%% Generated by Sphinx.
\def\sphinxdocclass{report}
\documentclass[letterpaper,10pt,english]{sphinxmanual}
\ifdefined\pdfpxdimen
   \let\sphinxpxdimen\pdfpxdimen\else\newdimen\sphinxpxdimen
\fi \sphinxpxdimen=.75bp\relax
\ifdefined\pdfimageresolution
    \pdfimageresolution= \numexpr \dimexpr1in\relax/\sphinxpxdimen\relax
\fi
%% let collapsible pdf bookmarks panel have high depth per default
\PassOptionsToPackage{bookmarksdepth=5}{hyperref}

\PassOptionsToPackage{booktabs}{sphinx}
\PassOptionsToPackage{colorrows}{sphinx}

\PassOptionsToPackage{warn}{textcomp}
\usepackage[utf8]{inputenc}
\ifdefined\DeclareUnicodeCharacter
% support both utf8 and utf8x syntaxes
  \ifdefined\DeclareUnicodeCharacterAsOptional
    \def\sphinxDUC#1{\DeclareUnicodeCharacter{"#1}}
  \else
    \let\sphinxDUC\DeclareUnicodeCharacter
  \fi
  \sphinxDUC{00A0}{\nobreakspace}
  \sphinxDUC{2500}{\sphinxunichar{2500}}
  \sphinxDUC{2502}{\sphinxunichar{2502}}
  \sphinxDUC{2514}{\sphinxunichar{2514}}
  \sphinxDUC{251C}{\sphinxunichar{251C}}
  \sphinxDUC{2572}{\textbackslash}
\fi
\usepackage{cmap}
\usepackage[T1]{fontenc}
\usepackage{amsmath,amssymb,amstext}
\usepackage{babel}



\usepackage{tgtermes}
\usepackage{tgheros}
\renewcommand{\ttdefault}{txtt}



\usepackage[Bjarne]{fncychap}
\usepackage{sphinx}

\fvset{fontsize=auto}
\usepackage{geometry}


% Include hyperref last.
\usepackage{hyperref}
% Fix anchor placement for figures with captions.
\usepackage{hypcap}% it must be loaded after hyperref.
% Set up styles of URL: it should be placed after hyperref.
\urlstyle{same}

\addto\captionsenglish{\renewcommand{\contentsname}{LAMMPS Documentation}}

\usepackage{sphinxmessages}



\setcounter{tocdepth}{2}
\setcounter{part}{-1}
\renewcommand{\sfdefault}{ptm}  % Use Times New Roman font for \textrm
\renewcommand{\sfdefault}{phv}  % Use Helvetica font for \textsf
\usepackage[columns=1]{idxlayout} % create index with only one column
% Set up math fonts to match text fonts
\DeclareSymbolFont{operators}   {OT1}{ptm}{m}{n}
\DeclareSymbolFont{bold}        {OT1}{ptm}{bx}{n}
\DeclareSymbolFont{italic}      {OT1}{ptm}{m}{it}
\DeclareSymbolFont{extraops}    {OT1}{ztmcm}{m}{n}
\DeclareSymbolFont{letters}     {OML}{ztmcm}{m}{it}
\DeclareSymbolFont{largesymbols}{OMX}{ztmcm}{m}{n}
% symbols (and \mathcal) are taken from computer modern.
% setup mappings
\DeclareSymbolFontAlphabet{\mathrm}    {operators}
\DeclareSymbolFontAlphabet{\mathnormal}{letters}
\DeclareMathAlphabet{\mathnormal}{OT1}{ptm}{m}{n}
\DeclareMathAlphabet{\mathrm}{OT1}{ptm}{m}{n}
\DeclareMathAlphabet{\mathbf}{OT1}{ptm}{bx}{n}
\DeclareMathAlphabet{\mathit}{OT1}{ptm}{m}{it}
\DeclareMathAlphabet{\mathtt}{OT1}{pcr}{m}{n}
\SetMathAlphabet\mathit{bold}{OT1}{ptm}{bx}{it}
% declare missing operators
\DeclareMathSymbol{\omicron}{0}{operators}{`\o}
\DeclareMathSymbol{\Gamma}{\mathalpha}{extraops}{"00}
\DeclareMathSymbol{\Delta}{\mathalpha}{extraops}{"01}
\DeclareMathSymbol{\Theta}{\mathalpha}{extraops}{"02}
\DeclareMathSymbol{\Lambda}{\mathalpha}{extraops}{"03}
\DeclareMathSymbol{\Xi}{\mathalpha}{extraops}{"04}
\DeclareMathSymbol{\Pi}{\mathalpha}{extraops}{"05}
\DeclareMathSymbol{\Sigma}{\mathalpha}{extraops}{"06}
\DeclareMathSymbol{\Upsilon}{\mathalpha}{extraops}{"07}
\DeclareMathSymbol{\Phi}{\mathalpha}{extraops}{"08}
\DeclareMathSymbol{\Psi}{\mathalpha}{extraops}{"09}
\DeclareMathSymbol{\Omega}{\mathalpha}{extraops}{"0A}

\renewcommand{\AA}{\mbox{\textrm{\r{A}}}}
% Make ToC number fields wider to accommodate sections with >= 100 subsections
%   or >= 10 subsections with >= 10 subsubsections
\makeatletter
% reset chapter counter for each part
\@addtoreset{chapter}{part}
\renewcommand*{\sphinxtableofcontentshook}{%
  \renewcommand*\l@section{\@dottedtocline{1}{1.5em}{3.1em}}
  \renewcommand*\l@subsection{\@dottedtocline{2}{4.6em}{4.5em}}
}
\makeatother


\title{LAMMPS Documentation}
\date{May 08, 2026}
\release{}
\author{The LAMMPS Developers}
\newcommand{\sphinxlogo}{\sphinxincludegraphics{lammps-logo-large.png}\par}
\renewcommand{\releasename}{}
\makeindex
\begin{document}

\ifdefined\shorthandoff
  \ifnum\catcode`\=\string=\active\shorthandoff{=}\fi
  \ifnum\catcode`\"=\active\shorthandoff{"}\fi
\fi

\pagestyle{empty}
\newcommand{\lammpsversion}{Release 29 Aug 2024 Update 1}

% customized titlepage
{%
  \hypersetup{pageanchor=false}% avoid duplicate destination warnings
  \begin{titlepage}%
    \sffamily\bfseries
    \begingroup % for PDF information dictionary
      \def\endgraf{ }\def\and{\& }%
      \pdfstringdefDisableCommands{\def\\{, }}% overwrite hyperref setup
      \hypersetup{pdfauthor={The LAMMPS Developers}, pdftitle={LAMMPS Documentation}}%
    \endgroup
    \noindent\rule{\textwidth}{4pt}\par
    \begin{center}%
      \sphinxlogo
      \vfill
      {\Huge LAMMPS Documentation \par}
      \vfill
      {\LARGE \lammpsversion \par}
      \vfill
      {\LARGE The LAMMPS Developers \par}
      {\Large developers@lammps.org $^*$ \par}
      \vfill\vfill\vfill
      {\normalsize ${}^*$ see
        \sphinxhref{https://www.lammps.org/authors.html}{https://www.lammps.org/authors.html}
        for details \par}
    \end{center}
    \noindent\rule{\textwidth}{4pt}\par
  \end{titlepage}%
  \setcounter{footnote}{0}%
  \clearpage
}

\pagestyle{plain}
\sphinxtableofcontents
\pagestyle{normal}
\phantomsection\label{\detokenize{Manual::doc}}



\part{About LAMMPS and this manual}
\label{\detokenize{Manual:about-lammps-and-this-manual}}
\sphinxAtStartPar
LAMMPS stands for \sphinxstylestrong{L}arge\sphinxhyphen{}scale \sphinxstylestrong{A}tomic/\sphinxstylestrong{M}olecular
\sphinxstylestrong{M}assively \sphinxstylestrong{P}arallel \sphinxstylestrong{S}imulator.

\sphinxAtStartPar
LAMMPS is a classical molecular dynamics simulation code focusing on
materials modeling.  It was designed to run efficiently on parallel
computers and to be easy to extend and modify.  Originally developed at
Sandia National Laboratories, a US Department of Energy facility, LAMMPS
now includes contributions from many research groups and individuals
from many institutions.  Most of the funding for LAMMPS has come from
the US Department of Energy (DOE).  LAMMPS is open\sphinxhyphen{}source software
distributed under the terms of the GNU Public License Version 2 (GPLv2).

\sphinxAtStartPar
The \sphinxhref{https://www.lammps.org}{LAMMPS website} has a variety of information about the
code.  It includes links to an online version of this manual, an
\sphinxhref{https://www.lammps.org/forum.html}{online forum} where users can post
questions and discuss LAMMPS, and a \sphinxhref{https://github.com/lammps/lammps}{GitHub site} where all LAMMPS development is
coordinated.


\bigskip\hrule\bigskip


\sphinxAtStartPar
The content for this manual is part of the LAMMPS distribution in its
doc directory.
\begin{itemize}
\item {} 
\sphinxAtStartPar
The version of the manual on the LAMMPS website corresponds to the
latest LAMMPS feature release.  It is available at:
\sphinxurl{https://docs.lammps.org/}.

\item {} 
\sphinxAtStartPar
A version of the manual corresponding to the latest LAMMPS stable
release (state of the \sphinxstyleemphasis{stable} branch on GitHub) is available online
at: \sphinxurl{https://docs.lammps.org/stable/}

\item {} 
\sphinxAtStartPar
A version of the manual with the features most recently added to
LAMMPS (state of the \sphinxstyleemphasis{develop} branch on GitHub) is available at:
\sphinxurl{https://docs.lammps.org/latest/}

\end{itemize}

\sphinxAtStartPar
If needed, you can build a copy on your local machine of the manual
(HTML pages or PDF file) for the version of LAMMPS you have
downloaded.  Follow the steps on the {\hyperref[\detokenize{Build_manual::doc}]{\sphinxcrossref{\DUrole{doc}{Build the LAMMPS documentation}}}} page.


\bigskip\hrule\bigskip


\sphinxAtStartPar
The manual is organized into three parts:
\begin{enumerate}
\sphinxsetlistlabels{\arabic}{enumi}{enumii}{}{.}%
\item {} 
\sphinxAtStartPar
The {\hyperref[\detokenize{Manual:user-documentation}]{\sphinxcrossref{\DUrole{std,std-ref}{User Guide}}}} with information about how
to obtain, configure, compile, install, and use LAMMPS,

\item {} 
\sphinxAtStartPar
the {\hyperref[\detokenize{Manual:programmer-documentation}]{\sphinxcrossref{\DUrole{std,std-ref}{Programmer Guide}}}} with
information about how to use the LAMMPS library interface from
different programming languages, how to modify and extend LAMMPS,
the program design, internal programming interfaces, and code
design conventions,

\item {} 
\sphinxAtStartPar
the {\hyperref[\detokenize{Manual:id1}]{\sphinxcrossref{\DUrole{std,std-ref}{Command Reference}}}} with detailed
descriptions of all input script commands available in LAMMPS.

\end{enumerate}


\bigskip\hrule\bigskip



\bigskip\hrule\bigskip



\part{User Guide}
\label{\detokenize{Manual:user-guide}}\label{\detokenize{Manual:user-documentation}}
\sphinxstepscope


\chapter{Introduction}
\label{\detokenize{Intro:introduction}}\label{\detokenize{Intro::doc}}
\sphinxAtStartPar
These pages provide a brief introduction to LAMMPS.

\sphinxstepscope


\section{Overview of LAMMPS}
\label{\detokenize{Intro_overview:overview-of-lammps}}\label{\detokenize{Intro_overview::doc}}
\sphinxAtStartPar
LAMMPS is a classical molecular dynamics (MD) code that models ensembles
of particles in a liquid, solid, or gaseous state.  It can model atomic,
polymeric, biological, solid\sphinxhyphen{}state (metals, ceramics, oxides), granular,
coarse\sphinxhyphen{}grained, or macroscopic systems using a variety of interatomic
potentials (force fields) and boundary conditions.  It can model 2d or
3d systems with sizes ranging from only a few particles up to billions.

\sphinxAtStartPar
LAMMPS can be built and run on single laptop or desktop machines, but is
designed for parallel computers.  It will run in serial and on any
parallel machine that supports the \sphinxhref{https://en.wikipedia.org/wiki/Message\_Passing\_Interface}{MPI} message\sphinxhyphen{}passing
library.  This includes shared\sphinxhyphen{}memory multicore, multi\sphinxhyphen{}CPU servers and
distributed\sphinxhyphen{}memory clusters and supercomputers. Parts of LAMMPS also
support \sphinxhref{https://www.openmp.org}{OpenMP multi\sphinxhyphen{}threading}, vectorization, and GPU
acceleration.

\sphinxAtStartPar
LAMMPS is written in C++ and requires a compiler that is at least
compatible with the C++\sphinxhyphen{}11 standard.  Earlier versions were written in
F77, F90, and C++\sphinxhyphen{}98.  See the \sphinxhref{https://www.lammps.org/history.html}{History page} of the website for details.  All
versions can be downloaded as source code from the \sphinxhref{https://www.lammps.org}{LAMMPS website}.

\sphinxAtStartPar
LAMMPS is designed to be easy to modify or extend with new capabilities,
such as new force fields, atom types, boundary conditions, or
diagnostics.  See the {\hyperref[\detokenize{Modify::doc}]{\sphinxcrossref{\DUrole{doc}{Modifying \& extending LAMMPS}}}} section of for more details.

\sphinxAtStartPar
In the most general sense, LAMMPS integrates Newton’s equations of
motion for a collection of interacting particles.  A single particle can
be an atom or molecule or electron, a coarse\sphinxhyphen{}grained cluster of atoms,
or a mesoscopic or macroscopic clump of material.  The interaction
models that LAMMPS includes are mostly short\sphinxhyphen{}ranged in nature; some
long\sphinxhyphen{}range models are included as well.

\sphinxAtStartPar
LAMMPS uses neighbor lists to keep track of nearby particles.  The lists
are optimized for systems with particles that are repulsive at short
distances, so that the local density of particles never becomes too
large.  This is in contrast to methods used for modeling plasma or
gravitational bodies (like galaxy formation).

\sphinxAtStartPar
On parallel machines, LAMMPS uses spatial\sphinxhyphen{}decomposition techniques with
MPI parallelization to partition the simulation domain into subdomains
of equal computational cost, one of which is assigned to each processor.
Processors communicate and store “ghost” atom information for atoms that
border their subdomain.  Multi\sphinxhyphen{}threading parallelization and GPU
acceleration with particle\sphinxhyphen{}decomposition can be used in addition.

\sphinxstepscope


\section{What does a LAMMPS version mean}
\label{\detokenize{Manual_version:what-does-a-lammps-version-mean}}\label{\detokenize{Manual_version::doc}}
\sphinxAtStartPar
The LAMMPS “version” is the date when it was released, such as 1 May
2014.  LAMMPS is updated continuously, and we aim to keep it working
correctly and reliably at all times.  Also, several variants of static
code analysis are run regularly to maintain or improve the overall code
quality, consistency, and compliance with programming standards, best
practices and style conventions.  You can follow its development in a
public \sphinxhref{https://github.com/lammps/lammps}{git repository on GitHub}.

\sphinxAtStartPar
Each version of LAMMPS contains all the documented \sphinxstyleemphasis{features} up to and
including its version date.  For recently added features, we add markers
to the documentation at which specific LAMMPS version a feature or
keyword was added or significantly changed.


\subsection{Identifying the Version}
\label{\detokenize{Manual_version:identifying-the-version}}
\sphinxAtStartPar
The version date is printed to the screen and log file every time you
run LAMMPS.  There also is an indication, if a LAMMPS binary was
compiled from version with modifications \sphinxstylestrong{after} a release.
It is also visible in the file src/version.h and in the LAMMPS directory
name created when you unpack a downloaded tarball.  And it is on the
first page of the {\hyperref[\detokenize{Manual::doc}]{\sphinxcrossref{\DUrole{doc}{manual}}}}.
\begin{itemize}
\item {} 
\sphinxAtStartPar
If you browse the HTML pages of the online version of the LAMMPS
manual, they will by default describe the most current feature release
version of LAMMPS.  In the navigation bar on the bottom left, there is
the option to view instead the documentation for the most recent
\sphinxstyleemphasis{stable} version or the documentation corresponding to the state of
the development branch.

\item {} 
\sphinxAtStartPar
If you browse the HTML pages included in your downloaded tarball, they
describe the version you have, which may be older than the online
version.

\end{itemize}


\subsection{LAMMPS releases, branches, and tags}
\label{\detokenize{Manual_version:lammps-releases-branches-and-tags}}
\begin{figure}[htbp]
\centering
\capstart

\noindent\sphinxincludegraphics{{lammps-releases}.png}
\caption{Relations between releases, main branches, and tags in the LAMMPS git repository}\label{\detokenize{Manual_version:id1}}\end{figure}


\subsubsection{Development}
\label{\detokenize{Manual_version:development}}
\sphinxAtStartPar
Modifications of the LAMMPS source code (like bug fixes, code
refactoring, updates to existing features, or addition of new features)
are organized into pull requests.  Pull requests will be merged into the
\sphinxstyleemphasis{develop} branch of the git repository after they pass automated testing
and code review by the LAMMPS developers.


\subsubsection{Feature Releases}
\label{\detokenize{Manual_version:feature-releases}}
\sphinxAtStartPar
When a sufficient number of new features and updates have accumulated
\sphinxstyleemphasis{and} the LAMMPS version on the \sphinxstyleemphasis{develop} branch passes an extended set
of automated tests, we release it as a \sphinxstyleemphasis{feature release}, which are
currently made every 4 to 8 weeks.  The \sphinxstyleemphasis{release} branch of the git
repository is updated with every such \sphinxstyleemphasis{feature release} and a tag in the
format \sphinxcode{\sphinxupquote{patch\_1May2014}} is added.  A summary of the most important
changes of these releases for the current year are posted on \sphinxhref{https://www.lammps.org/bug.html}{this
website page}.  More detailed release
notes are \sphinxhref{https://github.com/lammps/lammps/releases/}{available on GitHub}.


\subsubsection{Stable Releases}
\label{\detokenize{Manual_version:stable-releases}}
\sphinxAtStartPar
About once a year, we release a \sphinxstyleemphasis{stable release} version of LAMMPS.
This is done after a “stabilization period” where we apply only bug
fixes and small, non\sphinxhyphen{}intrusive changes to the \sphinxstyleemphasis{develop} branch but no
new features.  At the same time, the code is subjected to more detailed
and thorough manual testing than the default automated testing.
After such a \sphinxstyleemphasis{stable release}, both the \sphinxstyleemphasis{release} and the \sphinxstyleemphasis{stable}
branches are updated and two tags are applied, a \sphinxcode{\sphinxupquote{patch\_1May2014}} format
and a \sphinxcode{\sphinxupquote{stable\_1May2014}} format tag.


\subsubsection{Stable Release Updates}
\label{\detokenize{Manual_version:stable-release-updates}}
\sphinxAtStartPar
Between \sphinxstyleemphasis{stable releases}, we collect bug fixes and updates back\sphinxhyphen{}ported
from the \sphinxstyleemphasis{develop} branch in a branch called \sphinxstyleemphasis{maintenance}.  From the
\sphinxstyleemphasis{maintenance} branch we make occasional \sphinxstyleemphasis{stable update releases} and
update the \sphinxstyleemphasis{stable} branch accordingly.  The first update to the
\sphinxcode{\sphinxupquote{stable\_1May2014}} release would be tagged as
\sphinxcode{\sphinxupquote{stable\_1May2014\_update1}}.  These updates contain no new features.

\sphinxstepscope


\section{LAMMPS features}
\label{\detokenize{Intro_features:lammps-features}}\label{\detokenize{Intro_features::doc}}
\sphinxAtStartPar
LAMMPS is a classical molecular dynamics (MD) code with these general
classes of functionality:
\begin{enumerate}
\sphinxsetlistlabels{\arabic}{enumi}{enumii}{}{.}%
\item {} 
\sphinxAtStartPar
{\hyperref[\detokenize{Intro_features:general}]{\sphinxcrossref{\DUrole{std,std-ref}{General features}}}}

\item {} 
\sphinxAtStartPar
{\hyperref[\detokenize{Intro_features:particle}]{\sphinxcrossref{\DUrole{std,std-ref}{Particle and model types}}}}

\item {} 
\sphinxAtStartPar
{\hyperref[\detokenize{Intro_features:ff}]{\sphinxcrossref{\DUrole{std,std-ref}{Interatomic potentials (force fields)}}}}

\item {} 
\sphinxAtStartPar
{\hyperref[\detokenize{Intro_features:create}]{\sphinxcrossref{\DUrole{std,std-ref}{Atom creation}}}}

\item {} 
\sphinxAtStartPar
{\hyperref[\detokenize{Intro_features:ensemble}]{\sphinxcrossref{\DUrole{std,std-ref}{Ensembles, constraints, and boundary conditions}}}}

\item {} 
\sphinxAtStartPar
{\hyperref[\detokenize{Intro_features:integrate}]{\sphinxcrossref{\DUrole{std,std-ref}{Integrators}}}}

\item {} 
\sphinxAtStartPar
{\hyperref[\detokenize{Intro_features:diag}]{\sphinxcrossref{\DUrole{std,std-ref}{Diagnostics}}}}

\item {} 
\sphinxAtStartPar
{\hyperref[\detokenize{Intro_features:output}]{\sphinxcrossref{\DUrole{std,std-ref}{Output}}}}

\item {} 
\sphinxAtStartPar
{\hyperref[\detokenize{Intro_features:replica1}]{\sphinxcrossref{\DUrole{std,std-ref}{Multi\sphinxhyphen{}replica models}}}}

\item {} 
\sphinxAtStartPar
{\hyperref[\detokenize{Intro_features:prepost}]{\sphinxcrossref{\DUrole{std,std-ref}{Pre\sphinxhyphen{} and post\sphinxhyphen{}processing}}}}

\item {} 
\sphinxAtStartPar
{\hyperref[\detokenize{Intro_features:special}]{\sphinxcrossref{\DUrole{std,std-ref}{Specialized features (beyond MD itself)}}}}

\end{enumerate}


\bigskip\hrule\bigskip



\subsection{General features}
\label{\detokenize{Intro_features:general-features}}\label{\detokenize{Intro_features:general}}\begin{itemize}
\item {} 
\sphinxAtStartPar
runs on a single processor or in parallel

\item {} 
\sphinxAtStartPar
distributed memory message\sphinxhyphen{}passing parallelism (MPI)

\item {} 
\sphinxAtStartPar
shared memory multi\sphinxhyphen{}threading parallelism (OpenMP)

\item {} 
\sphinxAtStartPar
spatial decomposition of simulation domain for MPI parallelism

\item {} 
\sphinxAtStartPar
particle decomposition inside spatial decomposition for OpenMP and GPU parallelism

\item {} 
\sphinxAtStartPar
GPLv2 licensed open\sphinxhyphen{}source distribution

\item {} 
\sphinxAtStartPar
highly portable C++\sphinxhyphen{}11 (optional packages may require C++17)

\item {} 
\sphinxAtStartPar
modular code with most functionality in optional packages

\item {} 
\sphinxAtStartPar
only depends on MPI library for basic parallel functionality, MPI stub for serial compilation

\item {} 
\sphinxAtStartPar
other libraries are optional and only required for specific packages

\item {} 
\sphinxAtStartPar
GPU (CUDA, OpenCL, HIP, SYCL), Intel Xeon Phi, and OpenMP support for many code features

\item {} 
\sphinxAtStartPar
easy to extend with new features and functionality

\item {} 
\sphinxAtStartPar
runs from an input script

\item {} 
\sphinxAtStartPar
syntax for defining and using variables and formulas

\item {} 
\sphinxAtStartPar
syntax for looping over runs and breaking out of loops

\item {} 
\sphinxAtStartPar
run one or multiple simulations simultaneously (in parallel) from one script

\item {} 
\sphinxAtStartPar
build as library, invoke LAMMPS through library interface (from C, C++, Fortran) or provided Python wrapper or SWIG based wrappers

\item {} 
\sphinxAtStartPar
couple with other codes: LAMMPS calls other code, other code calls LAMMPS, umbrella code calls both, MDI coupling interface

\item {} 
\sphinxAtStartPar
call out to Python for computing forces, time integration, or other tasks

\item {} 
\sphinxAtStartPar
plugin interface for loading external features at runtime

\item {} 
\sphinxAtStartPar
large integrated collection of tests

\end{itemize}


\subsection{Particle and model types}
\label{\detokenize{Intro_features:particle-and-model-types}}\label{\detokenize{Intro_features:particle}}
\sphinxAtStartPar
(See {\hyperref[\detokenize{atom_style::doc}]{\sphinxcrossref{\DUrole{doc}{atom style}}}} command)
\begin{itemize}
\item {} 
\sphinxAtStartPar
atoms

\item {} 
\sphinxAtStartPar
coarse\sphinxhyphen{}grained particles (e.g. bead\sphinxhyphen{}spring polymers)

\item {} 
\sphinxAtStartPar
united\sphinxhyphen{}atom polymers or organic molecules

\item {} 
\sphinxAtStartPar
all\sphinxhyphen{}atom polymers, organic molecules, proteins, DNA

\item {} 
\sphinxAtStartPar
metals

\item {} 
\sphinxAtStartPar
metal oxides

\item {} 
\sphinxAtStartPar
granular materials

\item {} 
\sphinxAtStartPar
coarse\sphinxhyphen{}grained mesoscale models

\item {} 
\sphinxAtStartPar
finite\sphinxhyphen{}size spherical and ellipsoidal particles

\item {} 
\sphinxAtStartPar
finite\sphinxhyphen{}size line segment (2d) and triangle (3d) particles

\item {} 
\sphinxAtStartPar
finite\sphinxhyphen{}size rounded polygons (2d) and polyhedra (3d) particles

\item {} 
\sphinxAtStartPar
point dipole particles

\item {} 
\sphinxAtStartPar
particles with magnetic spin

\item {} 
\sphinxAtStartPar
rigid collections of n particles

\item {} 
\sphinxAtStartPar
hybrid combinations of these

\end{itemize}


\subsection{Interatomic potentials (force fields)}
\label{\detokenize{Intro_features:interatomic-potentials-force-fields}}\label{\detokenize{Intro_features:ff}}
\sphinxAtStartPar
(See {\hyperref[\detokenize{pair_style::doc}]{\sphinxcrossref{\DUrole{doc}{pair style}}}}, {\hyperref[\detokenize{bond_style::doc}]{\sphinxcrossref{\DUrole{doc}{bond style}}}},
{\hyperref[\detokenize{angle_style::doc}]{\sphinxcrossref{\DUrole{doc}{angle style}}}}, {\hyperref[\detokenize{dihedral_style::doc}]{\sphinxcrossref{\DUrole{doc}{dihedral style}}}},
{\hyperref[\detokenize{improper_style::doc}]{\sphinxcrossref{\DUrole{doc}{improper style}}}}, {\hyperref[\detokenize{kspace_style::doc}]{\sphinxcrossref{\DUrole{doc}{kspace style}}}}
commands)
\begin{itemize}
\item {} 
\sphinxAtStartPar
pairwise potentials: Lennard\sphinxhyphen{}Jones, Buckingham, Morse, Born\sphinxhyphen{}Mayer\sphinxhyphen{}Huggins, Yukawa, soft, Class II (COMPASS), hydrogen bond, harmonic, gaussian, tabulated, scripted

\item {} 
\sphinxAtStartPar
charged pairwise potentials: Coulombic, point\sphinxhyphen{}dipole

\item {} 
\sphinxAtStartPar
many\sphinxhyphen{}body potentials: EAM, Finnis/Sinclair, MEAM, MEAM+SW, EIM, EDIP, ADP, Stillinger\sphinxhyphen{}Weber, Tersoff, REBO, AIREBO, ReaxFF, COMB, Streitz\sphinxhyphen{}Mintmire, 3\sphinxhyphen{}body polymorphic, BOP, Vashishta

\item {} 
\sphinxAtStartPar
machine learning potentials: ACE, AGNI, GAP, Behler\sphinxhyphen{}Parrinello (N2P2), POD, RANN, SNAP

\item {} 
\sphinxAtStartPar
interfaces to ML potentials distributed by external groups: ANI, ChIMES, DeepPot, HIPNN, MTP

\item {} 
\sphinxAtStartPar
long\sphinxhyphen{}range interactions for charge, point\sphinxhyphen{}dipoles, and LJ dispersion:  Ewald, Wolf, PPPM (similar to particle\sphinxhyphen{}mesh Ewald), MSM, ScaFaCoS

\item {} 
\sphinxAtStartPar
polarization models: {\hyperref[\detokenize{fix_qeq::doc}]{\sphinxcrossref{\DUrole{doc}{QEq}}}}, {\hyperref[\detokenize{Howto_coreshell::doc}]{\sphinxcrossref{\DUrole{doc}{core/shell model}}}}, {\hyperref[\detokenize{Howto_drude::doc}]{\sphinxcrossref{\DUrole{doc}{Drude dipole model}}}}

\item {} 
\sphinxAtStartPar
charge equilibration (QEq via dynamic, point, shielded, Slater methods)

\item {} 
\sphinxAtStartPar
coarse\sphinxhyphen{}grained potentials: DPD, GayBerne, REsquared, colloidal, DLVO, oxDNA / oxRNA, SPICA

\item {} 
\sphinxAtStartPar
mesoscopic potentials: granular, Peridynamics, SPH, mesoscopic tubular potential (MESONT)

\item {} 
\sphinxAtStartPar
semi\sphinxhyphen{}empirical potentials: multi\sphinxhyphen{}ion generalized pseudopotential theory (MGPT), second moment tight binding + QEq (SMTB\sphinxhyphen{}Q)

\item {} 
\sphinxAtStartPar
electron force field (eFF, AWPMD)

\item {} 
\sphinxAtStartPar
bond potentials: harmonic, FENE, Morse, nonlinear, Class II (COMPASS), quartic (breakable), tabulated, scripted

\item {} 
\sphinxAtStartPar
angle potentials: harmonic, CHARMM, cosine, cosine/squared, cosine/periodic, Class II (COMPASS), tabulated, scripted

\item {} 
\sphinxAtStartPar
dihedral potentials: harmonic, CHARMM, multi\sphinxhyphen{}harmonic, helix, Class II (COMPASS), OPLS, tabulated, scripted

\item {} 
\sphinxAtStartPar
improper potentials: harmonic, cvff, umbrella, Class II (COMPASS), tabulated

\item {} 
\sphinxAtStartPar
polymer potentials: all\sphinxhyphen{}atom, united\sphinxhyphen{}atom, bead\sphinxhyphen{}spring, breakable

\item {} 
\sphinxAtStartPar
water potentials: TIP3P, TIP4P, SPC, SPC/E and variants

\item {} 
\sphinxAtStartPar
interlayer potentials for graphene and analogues, hetero\sphinxhyphen{}junctions

\item {} 
\sphinxAtStartPar
metal\sphinxhyphen{}organic framework potentials (QuickFF, MO\sphinxhyphen{}FF)

\item {} 
\sphinxAtStartPar
implicit solvent potentials: hydrodynamic lubrication, Debye

\item {} 
\sphinxAtStartPar
force\sphinxhyphen{}field compatibility with CHARMM, AMBER, DREIDING, OPLS, GROMACS, Class II (COMPASS), UFF, ClayFF, DREIDING, AMOEBA, INTERFACE

\item {} 
\sphinxAtStartPar
access to the \sphinxhref{https://openkim.org}{OpenKIM Repository} of potentials via the {\hyperref[\detokenize{kim_commands::doc}]{\sphinxcrossref{\DUrole{doc}{kim command}}}}

\item {} 
\sphinxAtStartPar
hybrid potentials: multiple pair, bond, angle, dihedral, improper potentials can be used in one simulation

\item {} 
\sphinxAtStartPar
overlaid potentials: superposition of multiple pair potentials (including many\sphinxhyphen{}body) with optional scale factor

\end{itemize}


\subsection{Atom creation}
\label{\detokenize{Intro_features:atom-creation}}\label{\detokenize{Intro_features:create}}
\sphinxAtStartPar
(See {\hyperref[\detokenize{read_data::doc}]{\sphinxcrossref{\DUrole{doc}{read\_data}}}}, {\hyperref[\detokenize{lattice::doc}]{\sphinxcrossref{\DUrole{doc}{lattice}}}},
{\hyperref[\detokenize{create_atoms::doc}]{\sphinxcrossref{\DUrole{doc}{create\_atoms}}}}, {\hyperref[\detokenize{delete_atoms::doc}]{\sphinxcrossref{\DUrole{doc}{delete\_atoms}}}},
{\hyperref[\detokenize{displace_atoms::doc}]{\sphinxcrossref{\DUrole{doc}{displace\_atoms}}}}, {\hyperref[\detokenize{replicate::doc}]{\sphinxcrossref{\DUrole{doc}{replicate}}}} commands)
\begin{itemize}
\item {} 
\sphinxAtStartPar
read in atom coordinates from files

\item {} 
\sphinxAtStartPar
create atoms on one or more lattices (e.g. grain boundaries)

\item {} 
\sphinxAtStartPar
delete geometric or logical groups of atoms (e.g. voids)

\item {} 
\sphinxAtStartPar
replicate existing atoms multiple times

\item {} 
\sphinxAtStartPar
displace atoms

\end{itemize}


\subsection{Ensembles, constraints, and boundary conditions}
\label{\detokenize{Intro_features:ensembles-constraints-and-boundary-conditions}}\label{\detokenize{Intro_features:ensemble}}
\sphinxAtStartPar
(See {\hyperref[\detokenize{fix::doc}]{\sphinxcrossref{\DUrole{doc}{fix}}}} command)
\begin{itemize}
\item {} 
\sphinxAtStartPar
2d or 3d systems

\item {} 
\sphinxAtStartPar
orthogonal or non\sphinxhyphen{}orthogonal (triclinic symmetry) simulation domains

\item {} 
\sphinxAtStartPar
constant NVE, NVT, NPT, NPH, Parrinello/Rahman integrators

\item {} 
\sphinxAtStartPar
thermostatting options for groups and geometric regions of atoms

\item {} 
\sphinxAtStartPar
pressure control via Nose/Hoover or Berendsen barostatting in 1 to 3 dimensions

\item {} 
\sphinxAtStartPar
simulation box deformation (tensile and shear)

\item {} 
\sphinxAtStartPar
harmonic (umbrella) constraint forces

\item {} 
\sphinxAtStartPar
rigid body constraints

\item {} 
\sphinxAtStartPar
SHAKE / RATTLE bond and angle constraints

\item {} 
\sphinxAtStartPar
motion constraints to manifold surfaces

\item {} 
\sphinxAtStartPar
Monte Carlo bond breaking, formation, swapping, template based reaction modeling

\item {} 
\sphinxAtStartPar
atom/molecule insertion and deletion

\item {} 
\sphinxAtStartPar
walls of various kinds, static and moving

\item {} 
\sphinxAtStartPar
non\sphinxhyphen{}equilibrium molecular dynamics (NEMD)

\item {} 
\sphinxAtStartPar
variety of additional boundary conditions and constraints

\end{itemize}


\subsection{Integrators}
\label{\detokenize{Intro_features:integrators}}\label{\detokenize{Intro_features:integrate}}
\sphinxAtStartPar
(See {\hyperref[\detokenize{run::doc}]{\sphinxcrossref{\DUrole{doc}{run}}}}, {\hyperref[\detokenize{run_style::doc}]{\sphinxcrossref{\DUrole{doc}{run\_style}}}}, {\hyperref[\detokenize{minimize::doc}]{\sphinxcrossref{\DUrole{doc}{minimize}}}} commands)
\begin{itemize}
\item {} 
\sphinxAtStartPar
velocity\sphinxhyphen{}Verlet integrator

\item {} 
\sphinxAtStartPar
Brownian dynamics

\item {} 
\sphinxAtStartPar
rigid body integration

\item {} 
\sphinxAtStartPar
energy minimization via conjugate gradient, steepest descent relaxation, or damped dynamics (FIRE, Quickmin)

\item {} 
\sphinxAtStartPar
rRESPA hierarchical timestepping

\item {} 
\sphinxAtStartPar
fixed or adaptive time step

\item {} 
\sphinxAtStartPar
rerun command for post\sphinxhyphen{}processing of dump files

\end{itemize}


\subsection{Diagnostics}
\label{\detokenize{Intro_features:diagnostics}}\label{\detokenize{Intro_features:diag}}\begin{itemize}
\item {} 
\sphinxAtStartPar
see various flavors of the {\hyperref[\detokenize{fix::doc}]{\sphinxcrossref{\DUrole{doc}{fix}}}} and {\hyperref[\detokenize{compute::doc}]{\sphinxcrossref{\DUrole{doc}{compute}}}} commands

\item {} 
\sphinxAtStartPar
introspection command for system, simulation, and compile time settings and configurations

\end{itemize}


\subsection{Output}
\label{\detokenize{Intro_features:output}}\label{\detokenize{Intro_features:id1}}
\sphinxAtStartPar
({\hyperref[\detokenize{dump::doc}]{\sphinxcrossref{\DUrole{doc}{dump}}}}, {\hyperref[\detokenize{restart::doc}]{\sphinxcrossref{\DUrole{doc}{restart}}}} commands)
\begin{itemize}
\item {} 
\sphinxAtStartPar
log file of thermodynamic info

\item {} 
\sphinxAtStartPar
text dump files of atom coordinates, velocities, other per\sphinxhyphen{}atom quantities

\item {} 
\sphinxAtStartPar
dump output on fixed and variable intervals, based timestep or simulated time

\item {} 
\sphinxAtStartPar
binary restart files

\item {} 
\sphinxAtStartPar
parallel I/O of dump and restart files

\item {} 
\sphinxAtStartPar
per\sphinxhyphen{}atom quantities (energy, stress, centro\sphinxhyphen{}symmetry parameter, CNA, etc.)

\item {} 
\sphinxAtStartPar
user\sphinxhyphen{}defined system\sphinxhyphen{}wide (log file) or per\sphinxhyphen{}atom (dump file) calculations

\item {} 
\sphinxAtStartPar
custom partitioning (chunks) for binning, and static or dynamic grouping of atoms for analysis

\item {} 
\sphinxAtStartPar
spatial, time, and per\sphinxhyphen{}chunk averaging of per\sphinxhyphen{}atom quantities

\item {} 
\sphinxAtStartPar
time averaging and histogramming of system\sphinxhyphen{}wide quantities

\item {} 
\sphinxAtStartPar
atom snapshots in native, XYZ, XTC, DCD, CFG, NetCDF, HDF5, ADIOS2, YAML formats

\item {} 
\sphinxAtStartPar
on\sphinxhyphen{}the\sphinxhyphen{}fly compression of output and decompression of read in files

\end{itemize}


\subsection{Multi\sphinxhyphen{}replica models}
\label{\detokenize{Intro_features:multi-replica-models}}\label{\detokenize{Intro_features:replica1}}\begin{itemize}
\item {} 
\sphinxAtStartPar
{\hyperref[\detokenize{neb::doc}]{\sphinxcrossref{\DUrole{doc}{nudged elastic band}}}}

\item {} 
\sphinxAtStartPar
{\hyperref[\detokenize{hyper::doc}]{\sphinxcrossref{\DUrole{doc}{hyperdynamics}}}}

\item {} 
\sphinxAtStartPar
{\hyperref[\detokenize{prd::doc}]{\sphinxcrossref{\DUrole{doc}{parallel replica dynamics}}}}

\item {} 
\sphinxAtStartPar
{\hyperref[\detokenize{tad::doc}]{\sphinxcrossref{\DUrole{doc}{temperature accelerated dynamics}}}}

\item {} 
\sphinxAtStartPar
{\hyperref[\detokenize{temper::doc}]{\sphinxcrossref{\DUrole{doc}{parallel tempering}}}}

\item {} 
\sphinxAtStartPar
path\sphinxhyphen{}integral MD: {\hyperref[\detokenize{fix_pimd::doc}]{\sphinxcrossref{\DUrole{doc}{first variant}}}}, {\hyperref[\detokenize{fix_ipi::doc}]{\sphinxcrossref{\DUrole{doc}{second variant}}}}

\item {} 
\sphinxAtStartPar
multi\sphinxhyphen{}walker collective variables with {\hyperref[\detokenize{fix_colvars::doc}]{\sphinxcrossref{\DUrole{doc}{Colvars}}}} and {\hyperref[\detokenize{fix_plumed::doc}]{\sphinxcrossref{\DUrole{doc}{Plumed}}}}

\end{itemize}


\subsection{Pre\sphinxhyphen{} and post\sphinxhyphen{}processing}
\label{\detokenize{Intro_features:pre-and-post-processing}}\label{\detokenize{Intro_features:prepost}}\begin{itemize}
\item {} 
\sphinxAtStartPar
A handful of pre\sphinxhyphen{} and post\sphinxhyphen{}processing tools are packaged with LAMMPS,
some of which can convert input and output files to/from formats used
by other codes; see the {\hyperref[\detokenize{Tools::doc}]{\sphinxcrossref{\DUrole{doc}{Tools}}}} page.

\item {} 
\sphinxAtStartPar
Our group has also written and released a separate toolkit called
\sphinxhref{https://lammps.github.io/pizza}{Pizza.py} which provides tools for doing setup, analysis,
plotting, and visualization for LAMMPS simulations.  Pizza.py is
written in \sphinxhref{https://www.python.org}{Python} and is available for download from \sphinxhref{https://lammps.github.io/pizza}{the Pizza.py WWW site}.

\end{itemize}


\subsection{Specialized features}
\label{\detokenize{Intro_features:specialized-features}}\label{\detokenize{Intro_features:special}}
\sphinxAtStartPar
LAMMPS can be built with optional packages which implement a variety of
additional capabilities.  See the {\hyperref[\detokenize{Packages::doc}]{\sphinxcrossref{\DUrole{doc}{Optional Packages}}}}
page for details.

\sphinxAtStartPar
These are LAMMPS capabilities which you may not think of as typical
classical MD options:
\begin{itemize}
\item {} 
\sphinxAtStartPar
{\hyperref[\detokenize{balance::doc}]{\sphinxcrossref{\DUrole{doc}{static}}}} and {\hyperref[\detokenize{fix_balance::doc}]{\sphinxcrossref{\DUrole{doc}{dynamic load\sphinxhyphen{}balancing}}}}, optional with recursive bisectioning decomposition

\item {} 
\sphinxAtStartPar
{\hyperref[\detokenize{Howto_body::doc}]{\sphinxcrossref{\DUrole{doc}{generalized aspherical particles}}}}

\item {} 
\sphinxAtStartPar
{\hyperref[\detokenize{fix_srd::doc}]{\sphinxcrossref{\DUrole{doc}{stochastic rotation dynamics (SRD)}}}}

\item {} 
\sphinxAtStartPar
{\hyperref[\detokenize{fix_imd::doc}]{\sphinxcrossref{\DUrole{doc}{real\sphinxhyphen{}time visualization and interactive MD}}}}, {\hyperref[\detokenize{dump_image::doc}]{\sphinxcrossref{\DUrole{doc}{built\sphinxhyphen{}in renderer for images and movies}}}}

\item {} 
\sphinxAtStartPar
calculate {\hyperref[\detokenize{compute_xrd::doc}]{\sphinxcrossref{\DUrole{doc}{virtual diffraction patterns}}}}

\item {} 
\sphinxAtStartPar
calculate {\hyperref[\detokenize{fix_phonon::doc}]{\sphinxcrossref{\DUrole{doc}{finite temperature phonon dispersion}}}} and the {\hyperref[\detokenize{dynamical_matrix::doc}]{\sphinxcrossref{\DUrole{doc}{dynamical matrix of minimized structures}}}}

\item {} 
\sphinxAtStartPar
{\hyperref[\detokenize{fix_atc::doc}]{\sphinxcrossref{\DUrole{doc}{atom\sphinxhyphen{}to\sphinxhyphen{}continuum coupling}}}} with finite elements

\item {} 
\sphinxAtStartPar
coupled rigid body integration via the {\hyperref[\detokenize{fix_poems::doc}]{\sphinxcrossref{\DUrole{doc}{POEMS}}}} library

\item {} 
\sphinxAtStartPar
{\hyperref[\detokenize{fix_qmmm::doc}]{\sphinxcrossref{\DUrole{doc}{QM/MM coupling}}}}

\item {} 
\sphinxAtStartPar
Monte Carlo via {\hyperref[\detokenize{fix_gcmc::doc}]{\sphinxcrossref{\DUrole{doc}{GCMC}}}} and {\hyperref[\detokenize{fix_tfmc::doc}]{\sphinxcrossref{\DUrole{doc}{tfMC}}}} and {\hyperref[\detokenize{fix_atom_swap::doc}]{\sphinxcrossref{\DUrole{doc}{atom swapping}}}}

\item {} 
\sphinxAtStartPar
{\hyperref[\detokenize{fix_ipi::doc}]{\sphinxcrossref{\DUrole{doc}{path\sphinxhyphen{}integral molecular dynamics (PIMD)}}}} and {\hyperref[\detokenize{fix_pimd::doc}]{\sphinxcrossref{\DUrole{doc}{this as well}}}}

\item {} 
\sphinxAtStartPar
{\hyperref[\detokenize{pair_dsmc::doc}]{\sphinxcrossref{\DUrole{doc}{Direct Simulation Monte Carlo}}}} for low\sphinxhyphen{}density fluids

\item {} 
\sphinxAtStartPar
{\hyperref[\detokenize{pair_peri::doc}]{\sphinxcrossref{\DUrole{doc}{Peridynamics modeling}}}}

\item {} 
\sphinxAtStartPar
{\hyperref[\detokenize{fix_lb_fluid::doc}]{\sphinxcrossref{\DUrole{doc}{Lattice Boltzmann fluid}}}}

\item {} 
\sphinxAtStartPar
{\hyperref[\detokenize{fix_tmd::doc}]{\sphinxcrossref{\DUrole{doc}{targeted}}}} and {\hyperref[\detokenize{fix_smd::doc}]{\sphinxcrossref{\DUrole{doc}{steered}}}} molecular dynamics

\item {} 
\sphinxAtStartPar
{\hyperref[\detokenize{fix_ttm::doc}]{\sphinxcrossref{\DUrole{doc}{two\sphinxhyphen{}temperature electron model}}}}

\end{itemize}

\sphinxstepscope


\section{LAMMPS non\sphinxhyphen{}features}
\label{\detokenize{Intro_nonfeatures:lammps-non-features}}\label{\detokenize{Intro_nonfeatures::doc}}
\sphinxAtStartPar
LAMMPS is designed to be a fast, parallel engine for molecular
dynamics (MD) simulations.  It provides only a modest amount of
functionality for setting up simulations and analyzing their output.

\sphinxAtStartPar
Originally, LAMMPS was not conceived and designed for:
\begin{itemize}
\item {} 
\sphinxAtStartPar
being run through a GUI

\item {} 
\sphinxAtStartPar
building molecular systems, or building molecular topologies

\item {} 
\sphinxAtStartPar
assign force\sphinxhyphen{}field coefficients automagically

\item {} 
\sphinxAtStartPar
perform sophisticated analysis of your MD simulation

\item {} 
\sphinxAtStartPar
visualize your MD simulation interactively

\item {} 
\sphinxAtStartPar
plot your output data

\end{itemize}

\sphinxAtStartPar
Over the years many of these limitations have been reduced or
removed. In part through features added to LAMMPS and in part
through external tools that either closely interface with LAMMPS
or extend LAMMPS.

\sphinxAtStartPar
Here are suggestions on how to perform these tasks:
\begin{itemize}
\item {} 
\sphinxAtStartPar
\sphinxstylestrong{GUI:} LAMMPS can be built as a library and a Python module that
wraps the library interface is provided.  Thus, GUI interfaces can be
written in Python or C/C++ that run LAMMPS and visualize or plot its
output.  Examples of this are provided in the python directory and
described on the {\hyperref[\detokenize{Python_head::doc}]{\sphinxcrossref{\DUrole{doc}{Python}}}} doc page.

\sphinxAtStartPar
Since version 2 August 2023 {\hyperref[\detokenize{Tools:lammps-gui}]{\sphinxcrossref{\DUrole{std,std-ref}{a LAMMPS\sphinxhyphen{}GUI tool}}}} is
included in LAMMPS.  Also, there are several external wrappers or GUI
front ends that are mentioned on the \sphinxhref{https://www.lammps.org/prepost.html}{Pre\sphinxhyphen{}/post\sphinxhyphen{}processing tools page} of the LAMMPS homepage.

\item {} 
\sphinxAtStartPar
\sphinxstylestrong{Builder:} Several pre\sphinxhyphen{}processing tools are packaged with LAMMPS.
Some of them convert input files in formats produced by other MD codes
such as CHARMM, AMBER, or Insight into LAMMPS input formats.  Some of
them are simple programs that will build simple molecular systems,
such as linear bead\sphinxhyphen{}spring polymer chains.  The moltemplate program is
a true molecular builder that will generate complex molecular models.
See the {\hyperref[\detokenize{Tools::doc}]{\sphinxcrossref{\DUrole{doc}{Tools}}}} page for details on tools packaged with
LAMMPS.  The \sphinxhref{https://www.lammps.org/prepost.html}{Pre\sphinxhyphen{}/post\sphinxhyphen{}processing tools page} of the LAMMPS homepage
describes a variety of third party tools for this task.  Furthermore,
some internal LAMMPS commands allow reconstructing, or selectively adding
topology information, as well as provide the option to insert molecule
templates instead of atoms for building bulk molecular systems.

\item {} 
\sphinxAtStartPar
\sphinxstylestrong{Force\sphinxhyphen{}field assignment:} The conversion tools described in the previous
bullet for CHARMM, AMBER, and Insight will also assign force field
coefficients in the LAMMPS format, assuming you provide CHARMM, AMBER,
or BIOVIA (formerly Accelrys) force field files. The tools
\sphinxhref{https://parmed.github.io/ParmEd/html/index.html}{ParmEd} and
\sphinxhref{https://github.com/shirtsgroup/InterMol}{InterMol} are particularly
powerful and flexible in converting force field and topology data
between various MD simulation programs.

\item {} 
\sphinxAtStartPar
\sphinxstylestrong{Simulation analysis:} If you want to perform analysis on\sphinxhyphen{}the\sphinxhyphen{}fly as
your simulation runs, see the {\hyperref[\detokenize{compute::doc}]{\sphinxcrossref{\DUrole{doc}{compute}}}} and {\hyperref[\detokenize{fix::doc}]{\sphinxcrossref{\DUrole{doc}{fix}}}} doc pages, which list commands that can be used in a LAMMPS
input script.  Also see the {\hyperref[\detokenize{Modify::doc}]{\sphinxcrossref{\DUrole{doc}{Modify}}}} page for info on
how to add your own analysis code or algorithms to LAMMPS.  For
post\sphinxhyphen{}processing, LAMMPS output such as {\hyperref[\detokenize{dump::doc}]{\sphinxcrossref{\DUrole{doc}{dump file snapshots}}}} can be converted into formats used by other MD or
post\sphinxhyphen{}processing codes.  To some degree, that conversion can be done
directly inside LAMMPS by interfacing to the VMD molfile plugins.  The
{\hyperref[\detokenize{rerun::doc}]{\sphinxcrossref{\DUrole{doc}{rerun}}}} command also allows post\sphinxhyphen{}processing of existing
trajectories, and through being able to read a variety of file
formats, this can also be used for analyzing trajectories from other
MD codes.  Some post\sphinxhyphen{}processing tools packaged with LAMMPS will do
these conversions.  Scripts provided in the tools/python directory can
extract and massage data in dump files to make it easier to import
into other programs.  See the {\hyperref[\detokenize{Tools::doc}]{\sphinxcrossref{\DUrole{doc}{Tools}}}} page for details on
these various options.

\sphinxAtStartPar
The \sphinxhref{https://www.lammps.org/prepost.html}{Pre\sphinxhyphen{}/post\sphinxhyphen{}processing page}
on the LAMMPS homepage lists some external packages for analysis of MD
simulation data, including data produced by LAMMPS.

\item {} 
\sphinxAtStartPar
\sphinxstylestrong{Visualization:} LAMMPS can produce NETPBM, JPG, or PNG format
snapshot images on\sphinxhyphen{}the\sphinxhyphen{}fly via its {\hyperref[\detokenize{dump_image::doc}]{\sphinxcrossref{\DUrole{doc}{dump image}}}}
command and pass them to an external program, \sphinxhref{https://ffmpeg.org/}{FFmpeg}, to generate movies from them.  The
{\hyperref[\detokenize{Tools:lammps-gui}]{\sphinxcrossref{\DUrole{std,std-ref}{LAMMPS\sphinxhyphen{}GUI tool}}}} has an \sphinxstyleemphasis{Snapshot Image Viewer}
which uses {\hyperref[\detokenize{dump_image::doc}]{\sphinxcrossref{\DUrole{doc}{dump image}}}} and allows to modify the
visualization settings interactively.  It also has a \sphinxstyleemphasis{Slide Show}
feature where images created by {\hyperref[\detokenize{dump_image::doc}]{\sphinxcrossref{\DUrole{doc}{dump image}}}} are
collected during a simulation and can be animated interactively or
exported to a movie with FFmpeg.

\sphinxAtStartPar
For high\sphinxhyphen{}quality, interactive visualization, there are many excellent
and free tools available.  See the \sphinxhref{https://www.lammps.org/viz.html}{Visualization Tools} page of the LAMMPS website for
visualization packages that can process LAMMPS output data.

\item {} 
\sphinxAtStartPar
\sphinxstylestrong{Plotting:} See the next bullet about Pizza.py as well as the
{\hyperref[\detokenize{Python_head::doc}]{\sphinxcrossref{\DUrole{doc}{Python}}}} page for examples of plotting LAMMPS
output.  Scripts provided with the \sphinxstyleemphasis{python} tool in the \sphinxcode{\sphinxupquote{tools}}
directory will extract and process data in log and dump files to make
it easier to analyze and plot.  See the {\hyperref[\detokenize{Tools::doc}]{\sphinxcrossref{\DUrole{doc}{Tools}}}} doc page
for more discussion of the various tools.

\sphinxAtStartPar
The {\hyperref[\detokenize{Tools:lammps-gui}]{\sphinxcrossref{\DUrole{std,std-ref}{LAMMPS\sphinxhyphen{}GUI tool}}}} has an \sphinxstyleemphasis{Chart Viewer} where
{\hyperref[\detokenize{thermo_style::doc}]{\sphinxcrossref{\DUrole{doc}{thermodynamic data}}}} computed by LAMMPS is
collected during the simulation and plotted immediately.

\item {} 
\sphinxAtStartPar
\sphinxstylestrong{Pizza.py:} Our group has also written a separate toolkit called
\sphinxhref{https://lammps.github.io/pizza}{Pizza.py} which can do certain kinds of
setup, analysis, plotting, and visualization (via OpenGL) for LAMMPS
simulations.  It thus provides some functionality for several of the
above bullets.  Pizza.py is written in \sphinxhref{https://www.python.org}{Python}
and is available for download from \sphinxhref{https://sjplimp.github.io/download.html}{this page}.

\end{itemize}

\sphinxstepscope


\section{LAMMPS portability and compatibility}
\label{\detokenize{Intro_portability:lammps-portability-and-compatibility}}\label{\detokenize{Intro_portability::doc}}
\sphinxAtStartPar
The primary form of distributing LAMMPS is through highly portable
source code.  But also several ways of obtaining LAMMPS as {\hyperref[\detokenize{Install::doc}]{\sphinxcrossref{\DUrole{doc}{precompiled
packages or through automated build mechanisms}}}} exist.  Most
of LAMMPS is written in C++, some support tools are written in Fortran
or Python or MATLAB.


\subsection{Programming language standards}
\label{\detokenize{Intro_portability:programming-language-standards}}
\sphinxAtStartPar
Most of the C++ code currently requires a compiler compatible with the
C++11 standard, the KOKKOS package currently requires C++17.  Most of
the Python code is written to be compatible with Python 3.5 or later or
Python 2.7.  Some Python scripts \sphinxstyleemphasis{require} Python 3 and a few others
still need to be ported from Python 2 to Python 3.


\subsection{Build systems}
\label{\detokenize{Intro_portability:build-systems}}
\sphinxAtStartPar
LAMMPS can be compiled from source code using a (traditional) build
system based on shell scripts, a few shell utilities (grep, sed, cat,
tr) and the GNU make program. This requires running within a Bourne
shell (\sphinxcode{\sphinxupquote{/bin/sh}}).  Alternatively, a build system with different back ends
can be created using CMake.  CMake must be at least version 3.16.


\subsection{Operating systems}
\label{\detokenize{Intro_portability:operating-systems}}
\sphinxAtStartPar
The primary development platform for LAMMPS is Linux.  Thus, the chances
for LAMMPS to compile without problems on Linux machines are the best.
Also, compilation and correct execution on macOS and Windows (using
Microsoft Visual C++) is checked automatically for largest part of the
source code.  Some (optional) features are not compatible with all
operating systems, either through limitations of the corresponding
LAMMPS source code or through source code or build system
incompatibilities of required libraries.

\sphinxAtStartPar
Executables for Windows may be created natively using either Cygwin or
Visual Studio or with a Linux to Windows MinGW cross\sphinxhyphen{}compiler.

\sphinxAtStartPar
Additionally, FreeBSD and Solaris have been tested successfully.


\subsection{Compilers}
\label{\detokenize{Intro_portability:compilers}}
\sphinxAtStartPar
The most commonly used compilers are the GNU compilers, but also Clang
and the Intel compilers have been successfully used on Linux, macOS, and
Windows.  Also, the Nvidia HPC SDK (formerly PGI compilers) will compile
LAMMPS (tested on Linux).


\subsection{CPU architectures}
\label{\detokenize{Intro_portability:cpu-architectures}}
\sphinxAtStartPar
The primary CPU architecture for running LAMMPS is 64\sphinxhyphen{}bit x86, but also
32\sphinxhyphen{}bit x86, and 64\sphinxhyphen{}bit ARM and PowerPC (64\sphinxhyphen{}bit, Little Endian) are
regularly tested.


\subsection{Portability compliance}
\label{\detokenize{Intro_portability:portability-compliance}}
\sphinxAtStartPar
Only a subset of the LAMMPS source code is \sphinxstyleemphasis{fully} compliant to \sphinxstyleemphasis{all}
of the above mentioned standards.  This is rather typical for projects
like LAMMPS that largely depend on contributions from the user community.
Not all contributors are trained as programmers and not all of them have
access to multiple platforms for testing.  As part of the continuous
integration process, however, all contributions are automatically tested
to compile, link, and pass some runtime tests on a selection of Linux
flavors, macOS, and Windows, and on Linux with different compilers.
Thus portability issues are often found before a pull request is merged.
Other platforms may be checked occasionally or when portability bugs are
reported.

\sphinxstepscope


\section{LAMMPS open\sphinxhyphen{}source license}
\label{\detokenize{Intro_opensource:lammps-open-source-license}}\label{\detokenize{Intro_opensource::doc}}

\subsection{GPL version of LAMMPS}
\label{\detokenize{Intro_opensource:gpl-version-of-lammps}}
\sphinxAtStartPar
LAMMPS is an open\sphinxhyphen{}source code, available free\sphinxhyphen{}of\sphinxhyphen{}charge, and distributed
under the terms of the \sphinxhref{https://github.com/lammps/lammps/blob/develop/LICENSE}{GNU Public License Version 2} (GPLv2),
which means you can use or modify the code however you wish for your own
purposes, but have to adhere to certain rules when redistributing it \sphinxhyphen{}
specifically in binary form \sphinxhyphen{} or are distributing software derived
from it or that includes parts of it.

\sphinxAtStartPar
LAMMPS comes with no warranty of any kind.

\sphinxAtStartPar
As each source file states in its header, it is a copyrighted code, and
thus not in the public domain. For more information about open\sphinxhyphen{}source
software and open\sphinxhyphen{}source distribution, see \sphinxhref{https://www.gnu.org}{www.gnu.org}
or \sphinxhref{https://www.opensource.org}{www.opensource.org}.  The legal text of the GPL as it
applies to LAMMPS is in the LICENSE file included in the LAMMPS distribution.

\sphinxAtStartPar
Here is a more specific summary of what the GPL means for LAMMPS users:

\sphinxAtStartPar
(1) Anyone is free to use, copy, modify, or extend LAMMPS in any way they
choose, including for commercial purposes.

\sphinxAtStartPar
(2) If you \sphinxstylestrong{distribute} a modified version of LAMMPS, it must remain
open\sphinxhyphen{}source, meaning you are required to distribute \sphinxstylestrong{all} of it under
the terms of the GPLv2.  You should \sphinxstylestrong{clearly} annotate such a modified
code as a derivative version of LAMMPS.  This is best done by changing
the name (example: LIGGGHTS is such a modified and extended version of
LAMMPS).

\sphinxAtStartPar
(3) If you release any code that includes or uses LAMMPS source code,
then it must also be open\sphinxhyphen{}sourced, meaning you distribute it under the
terms of the GPLv2.  You may write code that interfaces LAMMPS to a
differently licensed library.  In that case the code that provides the
interface must be licensed GPLv2, but not necessarily that library
unless you are distributing binaries that require the library to run.

\sphinxAtStartPar
(4) If you give LAMMPS files to someone else, the GPLv2 LICENSE file and
source file headers (including the copyright and GPLv2 notices) should
remain part of the code.


\subsection{LGPL version of LAMMPS}
\label{\detokenize{Intro_opensource:lgpl-version-of-lammps}}
\sphinxAtStartPar
We occasionally make stable LAMMPS releases available under the \sphinxhref{https://www.gnu.org/licenses/old-licenses/lgpl-2.1.html}{GNU
Lesser Public License v2.1}.  This is on request only and with
non\sphinxhyphen{}LGPL compliant files removed.  This allows uses linking non\sphinxhyphen{}GPL
compatible software with the (otherwise unmodified) LAMMPS library
or loading it dynamically at runtime.  Any \sphinxstylestrong{modifications} to
the LAMMPS code however, even with the LGPL licensed version, must still
be made available under the same open source terms as LAMMPS itself.

\sphinxstepscope


\section{Authors of LAMMPS}
\label{\detokenize{Intro_authors:authors-of-lammps}}\label{\detokenize{Intro_authors::doc}}
\sphinxAtStartPar
The current core LAMMPS developers are listed here (grouped by seniority
and sorted alphabetically by last name). You can email an individual
developer with code related questions for their area of expertise, or
send an email to all of them at this address: “developers at
lammps.org”.  General questions about LAMMPS should be posted in the
\sphinxhref{https://matsci.org/lammps/}{LAMMPS forum on MatSci}.

\small


\begin{savenotes}\sphinxattablestart
\sphinxthistablewithglobalstyle
\centering
\begin{tabular}[t]{\X{17}{100}\X{15}{100}\X{25}{100}\X{43}{100}}
\sphinxtoprule
\sphinxstyletheadfamily 
\sphinxAtStartPar
Name
&\sphinxstyletheadfamily 
\sphinxAtStartPar
Affiliation
&\sphinxstyletheadfamily 
\sphinxAtStartPar
Email
&\sphinxstyletheadfamily 
\sphinxAtStartPar
Areas of expertise
\\
\sphinxmidrule
\sphinxtableatstartofbodyhook
\sphinxAtStartPar
\sphinxhref{https://sites.google.com/site/akohlmey/}{Axel Kohlmeyer}
&
\sphinxAtStartPar
Temple U
&
\sphinxAtStartPar
akohlmey at gmail.com
&
\sphinxAtStartPar
OpenMP, library interfaces, LAMMPS\sphinxhyphen{}GUI, GitHub, MatSci forum, code maintenance, testing, releases
\\
\sphinxhline
\sphinxAtStartPar
\sphinxhref{https://sjplimp.github.io}{Steve Plimpton}
&
\sphinxAtStartPar
SNL (retired)
&
\sphinxAtStartPar
sjplimp at gmail.com
&
\sphinxAtStartPar
MD kernels, parallel algorithms \& scalability, code structure and design
\\
\sphinxhline
\sphinxAtStartPar
\sphinxhref{https://www2.sandia.gov/~athomps/}{Aidan Thompson}
&
\sphinxAtStartPar
SNL
&
\sphinxAtStartPar
athomps at sandia.gov
&
\sphinxAtStartPar
manybody potentials, machine learned potentials, materials science, statistical mechanics
\\
\sphinxhline&&&\\
\sphinxhline
\sphinxAtStartPar
\sphinxhref{https://rbberger.github.io/}{Richard Berger}
&
\sphinxAtStartPar
LANL
&
\sphinxAtStartPar
richard.berger at outlook.com
&
\sphinxAtStartPar
Python, HPC, DevOps
\\
\sphinxhline
\sphinxAtStartPar
\sphinxhref{https://enthalpiste.fr/}{Germain Clavier}
&
\sphinxAtStartPar
U Caen
&
\sphinxAtStartPar
germain.clavier at unicaen.fr
&
\sphinxAtStartPar
organic molecules, polymers, mechanical properties, surfaces, integrators, coarse\sphinxhyphen{}graining
\\
\sphinxhline
\sphinxAtStartPar
Joel Clemmer
&
\sphinxAtStartPar
SNL
&
\sphinxAtStartPar
jtclemm at sandia.gov
&
\sphinxAtStartPar
granular systems fluid/solid mechanics
\\
\sphinxhline
\sphinxAtStartPar
\sphinxhref{https://www.nanocipher.org/}{Jacob R. Gissinger}
&
\sphinxAtStartPar
Stevens Institute of Technology
&
\sphinxAtStartPar
jgissing at stevens.edu
&
\sphinxAtStartPar
reactive molecular dynamics, macro\sphinxhyphen{}molecular systems, type labels
\\
\sphinxhline
\sphinxAtStartPar
James Goff
&
\sphinxAtStartPar
SNL
&
\sphinxAtStartPar
jmgoff at sandia.gov
&
\sphinxAtStartPar
machine learned potentials, QEq solvers, Python
\\
\sphinxhline
\sphinxAtStartPar
Megan McCarthy
&
\sphinxAtStartPar
SNL
&
\sphinxAtStartPar
megmcca at sandia.gov
&
\sphinxAtStartPar
alloys, micro\sphinxhyphen{}structure, machine learned potentials
\\
\sphinxhline
\sphinxAtStartPar
Stan Moore
&
\sphinxAtStartPar
SNL
&
\sphinxAtStartPar
stamoor at sandia.gov
&
\sphinxAtStartPar
Kokkos, KSpace solvers, ReaxFF
\\
\sphinxhline
\sphinxAtStartPar
\sphinxhref{https://sites.google.com/site/ndtrung8/}{Trung Nguyen}
&
\sphinxAtStartPar
U Chicago
&
\sphinxAtStartPar
ndactrung at gmail.com
&
\sphinxAtStartPar
soft matter, GPU package
\\
\sphinxbottomrule
\end{tabular}
\sphinxtableafterendhook\par
\sphinxattableend\end{savenotes}

\normalsize

\sphinxAtStartPar
Past developers include Paul Crozier and Mark Stevens, both at SNL,
and Ray Shan, now at Materials Design.


\bigskip\hrule\bigskip


\sphinxAtStartPar
The \sphinxhref{https://www.lammps.org/authors.html}{Authors page} of the
\sphinxhref{https://www.lammps.org}{LAMMPS website} has a comprehensive list of all the individuals
who have contributed code for a new feature or command or tool to
LAMMPS.


\bigskip\hrule\bigskip


\sphinxAtStartPar
The following folks deserve special recognition.  Many of the packages
they have written are unique for an MD code and LAMMPS would not be as
general\sphinxhyphen{}purpose as it is without their expertise and efforts.
\begin{itemize}
\item {} 
\sphinxAtStartPar
Metin Aktulga (MSU), REAXFF package for C/C++ version of ReaxFF

\item {} 
\sphinxAtStartPar
Mike Brown (Intel), GPU and INTEL packages

\item {} 
\sphinxAtStartPar
Colin Denniston (U Western Ontario), LATBOLTZ package

\item {} 
\sphinxAtStartPar
Georg Ganzenmuller (EMI), MACHDYN and SPH packages

\item {} 
\sphinxAtStartPar
Andres Jaramillo\sphinxhyphen{}Botero (Caltech), EFF package for electron force field

\item {} 
\sphinxAtStartPar
Reese Jones (Sandia) and colleagues, ATC package for atom/continuum coupling

\item {} 
\sphinxAtStartPar
Christoph Kloss (DCS Computing), LIGGGHTS code for granular materials, built on top of LAMMPS

\item {} 
\sphinxAtStartPar
Rudra Mukherjee (JPL), POEMS package for articulated rigid body motion

\item {} 
\sphinxAtStartPar
Trung Ngyuen (U Chicago), GPU, RIGID, BODY, and DIELECTRIC packages

\item {} 
\sphinxAtStartPar
Mike Parks (Sandia), PERI package for Peridynamics

\item {} 
\sphinxAtStartPar
Roy Pollock (LLNL), Ewald and PPPM solvers

\item {} 
\sphinxAtStartPar
Julien Tranchida (CEA Cadarache), SPIN package

\item {} 
\sphinxAtStartPar
Christian Trott (Sandia), CUDA and KOKKOS packages

\item {} 
\sphinxAtStartPar
Ilya Valuev (JIHT), AWPMD package for wave packet MD

\item {} 
\sphinxAtStartPar
Greg Wagner (Northwestern U), MEAM package for MEAM potential

\end{itemize}


\bigskip\hrule\bigskip


\sphinxAtStartPar
As discussed on the \sphinxhref{https://www.lammps.org/history.html}{History page} of the website, LAMMPS
originated as a cooperative project between DOE labs and industrial
partners.  Folks involved in the design and testing of the original
version of LAMMPS were the following:
\begin{itemize}
\item {} 
\sphinxAtStartPar
John Carpenter (Mayo Clinic, formerly at Cray Research)

\item {} 
\sphinxAtStartPar
Terry Stouch (Lexicon Pharmaceuticals, formerly at Bristol Myers Squibb)

\item {} 
\sphinxAtStartPar
Steve Lustig (Dupont)

\item {} 
\sphinxAtStartPar
Jim Belak and Roy Pollock (LLNL)

\end{itemize}

\sphinxstepscope


\section{Citing LAMMPS}
\label{\detokenize{Intro_citing:citing-lammps}}\label{\detokenize{Intro_citing::doc}}

\subsection{Core Algorithms}
\label{\detokenize{Intro_citing:core-algorithms}}
\sphinxAtStartPar
The paper mentioned below is the best overview of LAMMPS, but there are
also publications describing particular models or algorithms implemented
in LAMMPS or complementary software that is has interfaces to.  Please
see below for how to cite contributions to LAMMPS.

\phantomsection\label{\detokenize{Intro_citing:lammps-paper}}
\sphinxAtStartPar
The latest canonical publication that describes the basic features, the
source code design, the program structure, the spatial decomposition
approach, the neighbor finding, basic communications algorithms, and how
users and developers have contributed to LAMMPS is:
\begin{quote}

\sphinxAtStartPar
\sphinxhref{https://doi.org/10.1016/j.cpc.2021.108171}{LAMMPS \sphinxhyphen{} A flexible simulation tool for particle\sphinxhyphen{}based materials modeling at the atomic, meso, and continuum scales, Comp. Phys. Comm. 271, 108171 (2022)}
\end{quote}

\sphinxAtStartPar
So a project using LAMMPS or a derivative application that uses LAMMPS
as a simulation engine should cite this paper.  The paper is expected to
be published in its final form under the same DOI in the first half
of 2022.  Please also give the URL of the LAMMPS website in your paper,
namely \sphinxurl{https://www.lammps.org}.

\sphinxAtStartPar
The original publication describing the parallel algorithms used in the
initial versions of LAMMPS is:
\begin{quote}

\sphinxAtStartPar
\sphinxhref{https://doi.org/10.1006/jcph.1995.1039}{S. Plimpton, Fast Parallel Algorithms for Short\sphinxhyphen{}Range Molecular Dynamics, J Comp Phys, 117, 1\sphinxhyphen{}19 (1995).}
\end{quote}


\subsection{DOI for the LAMMPS source code}
\label{\detokenize{Intro_citing:doi-for-the-lammps-source-code}}
\sphinxAtStartPar
The LAMMPS developers use the \sphinxhref{https://zenodo.org/}{Zenodo service at CERN}
to create digital object identifiers (DOI) for stable releases of the
LAMMPS source code.  There are two types of DOIs for the LAMMPS source code.

\sphinxAtStartPar
The canonical DOI for \sphinxstylestrong{all} versions of LAMMPS, which will always
point to the \sphinxstylestrong{latest} stable release version, is:
\begin{itemize}
\item {} 
\sphinxAtStartPar
DOI: \sphinxhref{https://dx.doi.org/10.5281/zenodo.3726416}{10.5281/zenodo.3726416}

\end{itemize}

\sphinxAtStartPar
In addition there are DOIs generated for individual stable releases:
\begin{itemize}
\item {} 
\sphinxAtStartPar
3 March 2020 version: \sphinxhref{https://dx.doi.org/10.5281/zenodo.3726417}{DOI:10.5281/zenodo.3726417}

\item {} 
\sphinxAtStartPar
29 October 2020 version: \sphinxhref{https://dx.doi.org/10.5281/zenodo.4157471}{DOI:10.5281/zenodo.4157471}

\item {} 
\sphinxAtStartPar
29 September 2021 version: \sphinxhref{https://dx.doi.org/10.5281/zenodo.6386596}{DOI:10.5281/zenodo.6386596}

\item {} 
\sphinxAtStartPar
23 June 2022 version: \sphinxhref{https://doi.org/10.5281/zenodo.10806836}{DOI:10.5281/zenodo.10806836}

\item {} 
\sphinxAtStartPar
2 August 2023 version: \sphinxhref{https://doi.org/10.5281/zenodo.10806852}{DOI:10.5281/zenodo.10806852}

\end{itemize}


\subsection{Home page}
\label{\detokenize{Intro_citing:home-page}}
\sphinxAtStartPar
The LAMMPS website at \sphinxhref{https://www.lammps.org}{https://www.lammps.org/} is the canonical location for information
about LAMMPS and its features.


\subsection{Citing contributions}
\label{\detokenize{Intro_citing:citing-contributions}}
\sphinxAtStartPar
LAMMPS has many features that use either previously published methods
and algorithms or novel features.  It also includes potential parameter
files for specific models.  Where available, a reminder about references
for optional features used in a specific run is printed to the screen
and log file.  Style and output location can be selected with the
{\hyperref[\detokenize{Run_options:cite}]{\sphinxcrossref{\DUrole{std,std-ref}{\sphinxhyphen{}cite command\sphinxhyphen{}line switch}}}}.  Additional references are
given in the documentation of the {\hyperref[\detokenize{Commands_all::doc}]{\sphinxcrossref{\DUrole{doc}{corresponding commands}}}} or in the {\hyperref[\detokenize{Howto::doc}]{\sphinxcrossref{\DUrole{doc}{Howto tutorials}}}}.  Please make
certain, that you provide the proper acknowledgments and citations in
any published works using LAMMPS.

\sphinxstepscope


\section{Additional website links}
\label{\detokenize{Intro_website:additional-website-links}}\label{\detokenize{Intro_website::doc}}
\sphinxAtStartPar
The \sphinxhref{https://www.lammps.org}{LAMMPS website} has a variety of additional info about
LAMMPS, beyond what is in this manual.  Some other useful resources
available online are listed below.
\begin{itemize}
\item {} 
\sphinxAtStartPar
\sphinxhref{https://github.com/lammps/lammps}{LAMMPS source code repository on GitHub}

\item {} 
\sphinxAtStartPar
\sphinxhref{https://matsci.org/lammps/}{LAMMPS forum on matsci.org}

\item {} 
\sphinxAtStartPar
\sphinxhref{https://www.lammps.org/bug.html}{Recent bug fixes and new features}

\item {} 
\sphinxAtStartPar
\sphinxhref{https://www.lammps.org/download.html}{Download info}

\item {} 
\sphinxAtStartPar
\sphinxhref{https://www.lammps.org/glossary.html}{Glossary of terms relevant to LAMMPS}

\item {} 
\sphinxAtStartPar
\sphinxhref{https://www.lammps.org/pictures.html}{LAMMPS highlights with images}

\item {} 
\sphinxAtStartPar
\sphinxhref{https://www.lammps.org/movies.html}{LAMMPS highlights with movies}

\item {} 
\sphinxAtStartPar
\sphinxhref{https://www.lammps.org/workshops.html}{Workshops}

\item {} 
\sphinxAtStartPar
\sphinxhref{https://www.lammps.org/tutorials.html}{Tutorials}

\item {} 
\sphinxAtStartPar
\sphinxhref{https://www.lammps.org/prepost.html}{Pre\sphinxhyphen{} and post\sphinxhyphen{}processing tools for LAMMPS}

\item {} 
\sphinxAtStartPar
\sphinxhref{https://www.lammps.org/external.html}{Other software usable with LAMMPS}

\item {} 
\sphinxAtStartPar
\sphinxhref{https://www.lammps.org/viz.html}{Viz tools usable with LAMMPS}

\item {} 
\sphinxAtStartPar
\sphinxhref{https://www.lammps.org/bench.html}{Benchmark performance}

\item {} 
\sphinxAtStartPar
\sphinxhref{https://www.lammps.org/papers.html}{Publications that have cited LAMMPS}

\item {} 
\sphinxAtStartPar
\sphinxhref{https://www.lammps.org/authors.html}{Authors of LAMMPS}

\item {} 
\sphinxAtStartPar
\sphinxhref{https://www.lammps.org/history.html}{History of LAMMPS development}

\item {} 
\sphinxAtStartPar
\sphinxhref{https://www.lammps.org/funding.html}{Funding for LAMMPS}

\end{itemize}

\sphinxstepscope


\chapter{Install LAMMPS}
\label{\detokenize{Install:install-lammps}}\label{\detokenize{Install::doc}}
\sphinxAtStartPar
You can download LAMMPS as an executable or as source code.

\sphinxAtStartPar
When downloading the LAMMPS source code, you also have to {\hyperref[\detokenize{Build::doc}]{\sphinxcrossref{\DUrole{doc}{build
LAMMPS}}}}.  But you have more flexibility as to what features to
include or exclude in the build.  When you download and install
pre\sphinxhyphen{}compiled LAMMPS executables, you are limited to install which
version of LAMMPS is available and which features are included of these
builds.  If you plan to {\hyperref[\detokenize{Modify::doc}]{\sphinxcrossref{\DUrole{doc}{modify or extend LAMMPS}}}}, then
you \sphinxstylestrong{must} build LAMMPS from the source code.

\begin{sphinxadmonition}{note}{Note:}
\sphinxAtStartPar
If you have questions about the pre\sphinxhyphen{}compiled LAMMPS executables, you
need to contact the people preparing those executables.  The LAMMPS
developers have no control over their choices of how they configure
and build their packages and when they update them.
\end{sphinxadmonition}


\bigskip\hrule\bigskip


\sphinxstepscope


\section{Download an executable for Linux}
\label{\detokenize{Install_linux:download-an-executable-for-linux}}\label{\detokenize{Install_linux::doc}}
\sphinxAtStartPar
Binaries are available for different versions of Linux:
\begin{itemize}
\item {} 
\sphinxAtStartPar
{\hyperref[\detokenize{Install_linux:static}]{\sphinxcrossref{\DUrole{std,std-ref}{Pre\sphinxhyphen{}built static Linux x86\_64 executables}}}}

\item {} 
\sphinxAtStartPar
{\hyperref[\detokenize{Install_linux:ubuntu}]{\sphinxcrossref{\DUrole{std,std-ref}{Pre\sphinxhyphen{}built Ubuntu and Debian Linux executables}}}}

\item {} 
\sphinxAtStartPar
{\hyperref[\detokenize{Install_linux:fedora}]{\sphinxcrossref{\DUrole{std,std-ref}{Pre\sphinxhyphen{}built Fedora Linux executables}}}}

\item {} 
\sphinxAtStartPar
{\hyperref[\detokenize{Install_linux:epel}]{\sphinxcrossref{\DUrole{std,std-ref}{Pre\sphinxhyphen{}built EPEL Linux executables (RHEL, CentOS)}}}}

\item {} 
\sphinxAtStartPar
{\hyperref[\detokenize{Install_linux:opensuse}]{\sphinxcrossref{\DUrole{std,std-ref}{Pre\sphinxhyphen{}built OpenSuse Linux executables}}}}

\item {} 
\sphinxAtStartPar
{\hyperref[\detokenize{Install_linux:gentoo}]{\sphinxcrossref{\DUrole{std,std-ref}{Gentoo Linux executable}}}}

\item {} 
\sphinxAtStartPar
{\hyperref[\detokenize{Install_linux:arch}]{\sphinxcrossref{\DUrole{std,std-ref}{Arch Linux build\sphinxhyphen{}script}}}}

\end{itemize}

\begin{sphinxadmonition}{note}{Note:}
\sphinxAtStartPar
If you have questions about these pre\sphinxhyphen{}compiled LAMMPS executables,
you need to contact the people preparing those packages.  The LAMMPS
developers have no control over how they configure and build their
packages and when they update them.  They may only provide packages
for stable release versions and not always update the packages in a
timely fashion after a new LAMMPS release is made.
\end{sphinxadmonition}


\bigskip\hrule\bigskip



\subsection{Pre\sphinxhyphen{}built static Linux x86\_64 executables}
\label{\detokenize{Install_linux:pre-built-static-linux-x86-64-executables}}\label{\detokenize{Install_linux:static}}
\sphinxAtStartPar
Pre\sphinxhyphen{}built LAMMPS executables for Linux, that are statically linked and
compiled for 64\sphinxhyphen{}bit x86 CPUs (x86\_64 or AMD64) are available for download
at \sphinxurl{https://download.lammps.org/static/}.
Because of that static linkage (and unlike the Linux distribution specific
packages listed below), they do not depend on any installed software and
thus should run on \sphinxstyleemphasis{any} 64\sphinxhyphen{}bit x86 machine with \sphinxstyleemphasis{any} Linux version.

\sphinxAtStartPar
These executable include most of the available packages and multi\sphinxhyphen{}thread
parallelization (via INTEL, KOKKOS, or OPENMP package).  They are
\sphinxstylestrong{not} compatible with MPI.  Several of the LAMMPS tools executables
(e.g. \sphinxcode{\sphinxupquote{msi2lmp}}) are included as well.  Because of the static linkage,
there is no \sphinxcode{\sphinxupquote{liblammps.so}} library file and thus also the LAMMPS
python module, which depends on it, is not included.

\sphinxAtStartPar
The compressed tar archives available for download have names following
the pattern \sphinxcode{\sphinxupquote{lammps\sphinxhyphen{}linux\sphinxhyphen{}x86\_64\sphinxhyphen{}\textless{}version\textgreater{}.tar.gz}} and will all unpack
into a \sphinxcode{\sphinxupquote{lammps\sphinxhyphen{}static}} folder.  The executables are then in the
\sphinxcode{\sphinxupquote{lammps\sphinxhyphen{}static/bin/}} folder.  Since they do not depend on any other
software, they may be freely moved or copied around.


\bigskip\hrule\bigskip



\subsection{Pre\sphinxhyphen{}built Ubuntu and Debian Linux executables}
\label{\detokenize{Install_linux:pre-built-ubuntu-and-debian-linux-executables}}\label{\detokenize{Install_linux:ubuntu}}
\sphinxAtStartPar
A pre\sphinxhyphen{}built LAMMPS executable, suitable for running on the latest Ubuntu
and Debian Linux versions, can be downloaded as a Debian package.  This
allows you to install LAMMPS with a single command, and stay (mostly)
up\sphinxhyphen{}to\sphinxhyphen{}date with the current stable version of LAMMPS by simply updating
your operating system.

\sphinxAtStartPar
To install LAMMPS do the following once:

\begin{sphinxVerbatim}[commandchars=\\\{\}]
sudo\PYG{+w}{ }apt\PYGZhy{}get\PYG{+w}{ }install\PYG{+w}{ }lammps
\end{sphinxVerbatim}

\sphinxAtStartPar
This downloads an executable named \sphinxcode{\sphinxupquote{lmp}} to your box and multiple
packages with supporting data, examples and libraries as well as any
missing dependencies.  For example, the LAMMPS binary in this package is
built with the {\hyperref[\detokenize{Build_extras:kim}]{\sphinxcrossref{\DUrole{std,std-ref}{KIM package}}}} enabled, which results in the
above command also installing the \sphinxcode{\sphinxupquote{kim\sphinxhyphen{}api}} binaries when LAMMPS is
installed, unless they were installed already.  In order to use
potentials from \sphinxhref{https://openkim.org}{openkim.org}, you can also install the
\sphinxcode{\sphinxupquote{openkim\sphinxhyphen{}models}} package:

\begin{sphinxVerbatim}[commandchars=\\\{\}]
sudo\PYG{+w}{ }apt\PYGZhy{}get\PYG{+w}{ }install\PYG{+w}{ }openkim\PYGZhy{}models
\end{sphinxVerbatim}

\sphinxAtStartPar
Or use the \sphinxhref{https://openkim.org/doc/usage/obtaining-models/\#installing\_api}{KIM\sphinxhyphen{}API commands}
to download and install individual models.

\sphinxAtStartPar
This LAMMPS executable can then be used in the usual way to run input
scripts:

\begin{sphinxVerbatim}[commandchars=\\\{\}]
lmp\PYG{+w}{ }\PYGZhy{}in\PYG{+w}{ }\PYG{k}{in}.lj
\end{sphinxVerbatim}

\sphinxAtStartPar
To update LAMMPS to the latest packaged version, do the following:

\begin{sphinxVerbatim}[commandchars=\\\{\}]
sudo\PYG{+w}{ }apt\PYGZhy{}get\PYG{+w}{ }update
\end{sphinxVerbatim}

\sphinxAtStartPar
This will also update other packages on your system.

\sphinxAtStartPar
To uninstall LAMMPS, do the following:

\begin{sphinxVerbatim}[commandchars=\\\{\}]
sudo\PYG{+w}{ }apt\PYGZhy{}get\PYG{+w}{ }remove\PYG{+w}{ }lammps
\end{sphinxVerbatim}

\sphinxAtStartPar
Please use \sphinxcode{\sphinxupquote{lmp \sphinxhyphen{}help}} to see which compilation options, packages,
and styles are included in the binary.

\sphinxAtStartPar
Thanks to Anton Gladky (gladky.anton at gmail.com) for setting up this
Ubuntu package capability.


\bigskip\hrule\bigskip



\subsection{Pre\sphinxhyphen{}built Fedora Linux executables}
\label{\detokenize{Install_linux:pre-built-fedora-linux-executables}}\label{\detokenize{Install_linux:fedora}}
\sphinxAtStartPar
Pre\sphinxhyphen{}built \sphinxhref{https://packages.fedoraproject.org/pkgs/lammps/}{LAMMPS packages for stable releases} are available in the
Fedora Linux distribution since Fedora version 28. The packages can be
installed via the dnf package manager. There are 3 basic varieties
(lammps = no MPI, lammps\sphinxhyphen{}mpich = MPICH MPI library, lammps\sphinxhyphen{}openmpi =
OpenMPI MPI library) and for each support for linking to the C library
interface (lammps\sphinxhyphen{}devel, lammps\sphinxhyphen{}mpich\sphinxhyphen{}devel, lammps\sphinxhyphen{}openmpi\sphinxhyphen{}devel), the
header for compiling programs using the C library interface
(lammps\sphinxhyphen{}headers), and the LAMMPS python module for Python 3. All
packages can be installed at the same time and the name of the LAMMPS
executable is \sphinxcode{\sphinxupquote{lmp}} and \sphinxcode{\sphinxupquote{lmp\_openmpi}} or \sphinxcode{\sphinxupquote{lmp\_mpich}} respectively.
By default, \sphinxcode{\sphinxupquote{lmp}} will refer to the serial executable, unless one of
the MPI environment modules is loaded (\sphinxcode{\sphinxupquote{module load mpi/mpich\sphinxhyphen{}x86\_64}}
or \sphinxcode{\sphinxupquote{module load mpi/openmpi\sphinxhyphen{}x86\_64}}).  Then the corresponding parallel
LAMMPS executable can be used.  The same mechanism applies when loading
the LAMMPS python module.

\sphinxAtStartPar
To install LAMMPS with OpenMPI and run an input \sphinxcode{\sphinxupquote{in.lj}} with 2 CPUs do:

\begin{sphinxVerbatim}[commandchars=\\\{\}]
dnf\PYG{+w}{ }install\PYG{+w}{ }lammps\PYGZhy{}openmpi
module\PYG{+w}{ }load\PYG{+w}{ }mpi/openmpi\PYGZhy{}x86\PYGZus{}64
mpirun\PYG{+w}{ }\PYGZhy{}np\PYG{+w}{ }\PYG{l+m}{2}\PYG{+w}{ }lmp\PYG{+w}{ }\PYGZhy{}in\PYG{+w}{ }\PYG{k}{in}.lj
\end{sphinxVerbatim}

\sphinxAtStartPar
The \sphinxcode{\sphinxupquote{dnf install}} command is needed only once.  In case of a new LAMMPS
stable release, \sphinxcode{\sphinxupquote{dnf update}} will automatically update to the newer
version as soon as the RPM files are built and uploaded to the download
mirrors. The \sphinxcode{\sphinxupquote{module load}} command is needed once per (shell) session
or shell terminal instance, unless it is automatically loaded from the
shell profile.

\sphinxAtStartPar
The LAMMPS binary is built with the {\hyperref[\detokenize{Build_extras:kim}]{\sphinxcrossref{\DUrole{std,std-ref}{KIM package}}}} which
results in the above command also installing the \sphinxtitleref{kim\sphinxhyphen{}api} binaries when LAMMPS
is installed.  In order to use potentials from \sphinxhref{https://openkim.org}{openkim.org}, you
can install the \sphinxtitleref{openkim\sphinxhyphen{}models} package

\begin{sphinxVerbatim}[commandchars=\\\{\}]
dnf\PYG{+w}{ }install\PYG{+w}{ }openkim\PYGZhy{}models
\end{sphinxVerbatim}

\sphinxAtStartPar
Please use \sphinxcode{\sphinxupquote{lmp \sphinxhyphen{}help}} to see which compilation options, packages,
and styles are included in the binary.

\sphinxAtStartPar
Thanks to Christoph Junghans (LANL) for making LAMMPS available in Fedora.


\bigskip\hrule\bigskip



\subsection{Pre\sphinxhyphen{}built EPEL Linux executable}
\label{\detokenize{Install_linux:pre-built-epel-linux-executable}}\label{\detokenize{Install_linux:epel}}
\sphinxAtStartPar
Pre\sphinxhyphen{}built LAMMPS (and KIM) packages for stable releases are available
in the \sphinxhref{https://docs.fedoraproject.org/en-US/epel/}{Extra Packages for Enterprise Linux (EPEL) repository}
for use with Red Hat Enterprise Linux (RHEL) or CentOS version 7.x
and compatible Linux distributions. Names of packages, executable,
and content are the same as described above for Fedora Linux.
But RHEL/CentOS 7.x uses the \sphinxcode{\sphinxupquote{yum}} package manager instead of \sphinxcode{\sphinxupquote{dnf}}
in Fedora 28.

\sphinxAtStartPar
Please use \sphinxcode{\sphinxupquote{lmp \sphinxhyphen{}help}} to see which compilation options, packages,
and styles are included in the binary.

\sphinxAtStartPar
Thanks to Christoph Junghans (LANL) for making LAMMPS available in EPEL.


\bigskip\hrule\bigskip



\subsection{Pre\sphinxhyphen{}built OpenSuse Linux executable}
\label{\detokenize{Install_linux:pre-built-opensuse-linux-executable}}\label{\detokenize{Install_linux:opensuse}}
\sphinxAtStartPar
A pre\sphinxhyphen{}built LAMMPS package for stable releases is available
in OpenSuse as of Leap 15.0. You can install the package with:

\begin{sphinxVerbatim}[commandchars=\\\{\}]
zypper\PYG{+w}{ }install\PYG{+w}{ }lammps
\end{sphinxVerbatim}

\sphinxAtStartPar
This includes support for OpenMPI. The name of the LAMMPS executable
is \sphinxcode{\sphinxupquote{lmp}}. To run an input in parallel on 2 CPUs you would do:

\begin{sphinxVerbatim}[commandchars=\\\{\}]
mpirun\PYG{+w}{ }\PYGZhy{}np\PYG{+w}{ }\PYG{l+m}{2}\PYG{+w}{ }lmp\PYG{+w}{ }\PYGZhy{}in\PYG{+w}{ }\PYG{k}{in}.lj
\end{sphinxVerbatim}

\sphinxAtStartPar
Please use \sphinxcode{\sphinxupquote{lmp \sphinxhyphen{}help}} to see which compilation options, packages,
and styles are included in the binary.

\sphinxAtStartPar
The LAMMPS binary is built with the {\hyperref[\detokenize{Build_extras:kim}]{\sphinxcrossref{\DUrole{std,std-ref}{KIM package}}}} which
results in the above command also installing the \sphinxtitleref{kim\sphinxhyphen{}api} binaries when LAMMPS
is installed.  In order to use potentials from \sphinxhref{https://openkim.org}{openkim.org}, you
can install the \sphinxtitleref{openkim\sphinxhyphen{}models} package

\begin{sphinxVerbatim}[commandchars=\\\{\}]
zypper\PYG{+w}{ }install\PYG{+w}{ }openkim\PYGZhy{}models
\end{sphinxVerbatim}

\sphinxAtStartPar
Thanks to Christoph Junghans (LANL) for making LAMMPS available in OpenSuse.


\bigskip\hrule\bigskip



\subsection{Gentoo Linux executable}
\label{\detokenize{Install_linux:gentoo-linux-executable}}\label{\detokenize{Install_linux:gentoo}}
\sphinxAtStartPar
LAMMPS is part of \sphinxhref{https://packages.gentoo.org/packages/sci-physics/lammps}{Gentoo’s main package tree} and can be
installed by typing:

\begin{sphinxVerbatim}[commandchars=\\\{\}]
emerge\PYG{+w}{ }\PYGZhy{}\PYGZhy{}ask\PYG{+w}{ }lammps
\end{sphinxVerbatim}

\sphinxAtStartPar
Note that in Gentoo the LAMMPS source code is downloaded and the package is
then compiled and installed on your machine.

\sphinxAtStartPar
Certain LAMMPS packages can be enabled via USE flags, type

\begin{sphinxVerbatim}[commandchars=\\\{\}]
equery\PYG{+w}{ }uses\PYG{+w}{ }lammps
\end{sphinxVerbatim}

\sphinxAtStartPar
for details.

\sphinxAtStartPar
Thanks to Nicolas Bock and Christoph Junghans (LANL) for setting up
this Gentoo capability.


\bigskip\hrule\bigskip



\subsection{Archlinux build\sphinxhyphen{}script}
\label{\detokenize{Install_linux:archlinux-build-script}}\label{\detokenize{Install_linux:arch}}
\sphinxAtStartPar
LAMMPS is available via Arch’s unofficial Arch User repository (AUR).
There are three scripts available, named \sphinxhref{https://aur.archlinux.org/packages/lammps}{lammps}, \sphinxhref{https://aur.archlinux.org/packages/lammps}{lammps\sphinxhyphen{}beta} and \sphinxhref{https://aur.archlinux.org/packages/lammps}{lammps\sphinxhyphen{}git}.  They respectively
package the stable, feature, and git releases.

\sphinxAtStartPar
To install, you will need to have the git package installed. You may use
any of the above names in\sphinxhyphen{}place of lammps.

\begin{sphinxVerbatim}[commandchars=\\\{\}]
git\PYG{+w}{ }clone\PYG{+w}{ }https://aur.archlinux.org/lammps.git
\PYG{n+nb}{cd}\PYG{+w}{ }lammps
makepkg\PYG{+w}{ }\PYGZhy{}s
makepkg\PYG{+w}{ }\PYGZhy{}i
\end{sphinxVerbatim}

\sphinxAtStartPar
To update LAMMPS, you may repeat the above, or change into the cloned
directory, and execute the following, after which, if there are any
changes, you may use makepkg as above.

\begin{sphinxVerbatim}[commandchars=\\\{\}]
git\PYG{+w}{ }pull
\end{sphinxVerbatim}

\sphinxAtStartPar
Alternatively, you may use an AUR helper to install these packages.

\sphinxAtStartPar
Note that the AUR provides build\sphinxhyphen{}scripts that download the source code
and then build and install the package on your machine.

\sphinxstepscope


\section{Download an executable for macOS}
\label{\detokenize{Install_mac:download-an-executable-for-macos}}\label{\detokenize{Install_mac::doc}}
\sphinxAtStartPar
LAMMPS can be downloaded, built, and configured for macOS with \sphinxhref{https://brew.sh}{Homebrew}.  (Alternatively, see the installation instructions for
{\hyperref[\detokenize{Install_conda::doc}]{\sphinxcrossref{\DUrole{doc}{downloading an executable via Conda}}}}.)  The
following LAMMPS packages are unavailable at this time because of
additional requirements not yet met: GPU, KOKKOS, MSCG, POEMS,
VORONOI.

\sphinxAtStartPar
After installing Homebrew, you can install LAMMPS on your system with
the following commands:

\begin{sphinxVerbatim}[commandchars=\\\{\}]
brew\PYG{+w}{ }install\PYG{+w}{ }lammps
\end{sphinxVerbatim}

\sphinxAtStartPar
This will install the executables “lammps\_serial” and “lammps\_mpi”, as
well as the LAMMPS “doc”, “potentials”, “tools”, “bench”, and “examples”
directories.

\sphinxAtStartPar
Once LAMMPS is installed, you can test the installation with the
Lennard\sphinxhyphen{}Jones benchmark file:

\begin{sphinxVerbatim}[commandchars=\\\{\}]
brew\PYG{+w}{ }\PYG{n+nb}{test}\PYG{+w}{ }lammps\PYG{+w}{ }\PYGZhy{}v
\end{sphinxVerbatim}

\sphinxAtStartPar
The LAMMPS binary is built with the {\hyperref[\detokenize{Build_extras:kim}]{\sphinxcrossref{\DUrole{std,std-ref}{KIM package}}}}, which
results in Homebrew also installing the \sphinxtitleref{kim\sphinxhyphen{}api} binaries when LAMMPS
is installed.  In order to use potentials from \sphinxhref{https://openkim.org}{openkim.org}, you can install the \sphinxtitleref{openkim\sphinxhyphen{}models} package

\begin{sphinxVerbatim}[commandchars=\\\{\}]
brew\PYG{+w}{ }install\PYG{+w}{ }openkim\PYGZhy{}models
\end{sphinxVerbatim}

\sphinxAtStartPar
If you have problems with the installation, you can post issues to
\sphinxhref{https://github.com/Homebrew/homebrew-core/issues}{this link}.

\sphinxAtStartPar
Thanks to Derek Thomas (derekt at cello.t.u\sphinxhyphen{}tokyo.ac.jp) for setting
up the Homebrew capability.

\sphinxstepscope


\section{Download an executable for Windows}
\label{\detokenize{Install_windows:download-an-executable-for-windows}}\label{\detokenize{Install_windows::doc}}
\sphinxAtStartPar
Pre\sphinxhyphen{}compiled Windows installers which install LAMMPS executables on a
Windows system can be downloaded from this site:
\begin{sphinxalltt}
\sphinxurl{https://packages.lammps.org/windows.html}
\end{sphinxalltt}

\sphinxAtStartPar
Note that each installer package has a date in its name, which
corresponds to the LAMMPS version of the same date.  Installers for
current and older versions of LAMMPS are available.  32\sphinxhyphen{}bit and 64\sphinxhyphen{}bit
installers are available, and each installer contains both a serial
and parallel executable.  The installer website also explains how to
install the Windows MPI package (MPICH2 from Argonne National Labs),
needed to run in parallel with MPI.

\sphinxAtStartPar
The LAMMPS binaries contain \sphinxstyleemphasis{all} {\hyperref[\detokenize{Packages::doc}]{\sphinxcrossref{\DUrole{doc}{optional packages}}}}
included in the source distribution except: ADIOS, H5MD, KIM, ML\sphinxhyphen{}PACE,
ML\sphinxhyphen{}QUIP, MSCG, NETCDF, QMMM, SCAFACOS, and VTK.  The serial version also
does not include the LATBOLTZ package.  The PYTHON package is only
available in the Python installers that bundle a Python runtime.  The
GPU package is compiled for OpenCL with mixed precision kernels.

\sphinxAtStartPar
The LAMMPS library is compiled as a shared library and the
{\hyperref[\detokenize{Python_module::doc}]{\sphinxcrossref{\DUrole{doc}{LAMMPS Python module}}}} is installed, so that
it is possible to load LAMMPS into a Python interpreter.

\sphinxAtStartPar
The installer site also has instructions on how to run LAMMPS under
Windows, once it is installed, in both serial and parallel.

\sphinxAtStartPar
When you download the installer package, you run it on your Windows
machine.  It will then prompt you with a dialog, where you can choose
the installation directory, unpack and copy several executables,
potential files, documentation PDFs, selected example files, etc.  It
will then update a few system settings (e.g. \sphinxcode{\sphinxupquote{PATH}}, \sphinxcode{\sphinxupquote{LAMMPS\_POTENTIALS}})
and add an entry into the Start Menu (with references to the
documentation, LAMMPS homepage and more).  From that menu, there is
also a link to an uninstaller that removes the files and undoes the
environment manipulations.

\sphinxAtStartPar
Note that to update to a newer version of LAMMPS, you should typically
uninstall the version you currently have, download a new installer, and
go through the installation procedure described above.  I.e. the same
procedure for installing/updating most Windows programs.  You can
install multiple versions of LAMMPS (in different directories), but only
the executable for the last\sphinxhyphen{}installed package will be found
automatically, so this should only be done for debugging purposes.

\sphinxstepscope


\section{Download an executable for Linux or macOS via Conda}
\label{\detokenize{Install_conda:download-an-executable-for-linux-or-macos-via-conda}}\label{\detokenize{Install_conda::doc}}
\sphinxAtStartPar
Pre\sphinxhyphen{}compiled LAMMPS binaries are available for macOS and Linux via the
\sphinxhref{https://docs.conda.io/en/latest/index.html}{Conda} package management system.

\sphinxAtStartPar
First, one must set up the Conda package manager on your system.  Follow
the instructions to install \sphinxhref{https://docs.conda.io/en/latest/miniconda.html}{Miniconda}, then
create a conda environment (named \sphinxtitleref{my\sphinxhyphen{}lammps\sphinxhyphen{}env} or whatever you
prefer) for your LAMMPS install:

\begin{sphinxVerbatim}[commandchars=\\\{\}]
conda\PYG{+w}{ }config\PYG{+w}{ }\PYGZhy{}\PYGZhy{}add\PYG{+w}{ }channels\PYG{+w}{ }conda\PYGZhy{}forge
conda\PYG{+w}{ }create\PYG{+w}{ }\PYGZhy{}n\PYG{+w}{ }my\PYGZhy{}lammps\PYGZhy{}env
\end{sphinxVerbatim}

\sphinxAtStartPar
Then, you can install LAMMPS on your system with the following command:

\begin{sphinxVerbatim}[commandchars=\\\{\}]
conda\PYG{+w}{ }activate\PYG{+w}{ }my\PYGZhy{}lammps\PYGZhy{}env
conda\PYG{+w}{ }install\PYG{+w}{ }lammps
\end{sphinxVerbatim}

\sphinxAtStartPar
The LAMMPS binary is built with the {\hyperref[\detokenize{Build_extras:kim}]{\sphinxcrossref{\DUrole{std,std-ref}{KIM package}}}}, which
results in Conda also installing the \sphinxtitleref{kim\sphinxhyphen{}api} binaries when LAMMPS is
installed.  In order to use potentials from \sphinxhref{https://openkim.org}{openkim.org},
you can install the \sphinxtitleref{openkim\sphinxhyphen{}models} package

\begin{sphinxVerbatim}[commandchars=\\\{\}]
conda\PYG{+w}{ }install\PYG{+w}{ }openkim\PYGZhy{}models
\end{sphinxVerbatim}

\sphinxAtStartPar
If you have problems with the installation, you can post issues to \sphinxhref{https://github.com/conda-forge/lammps-feedstock/issues}{this
link}.  Thanks to Jan Janssen
(Max\sphinxhyphen{}Planck\sphinxhyphen{}Institut fuer Eisenforschung) for setting up the Conda
capability.

\begin{sphinxadmonition}{note}{Note:}
\sphinxAtStartPar
If you have questions about these pre\sphinxhyphen{}compiled LAMMPS executables,
you need to contact the people preparing those packages.  The LAMMPS
developers have no control over their choices of how they configure
and build their packages and when they update them.
\end{sphinxadmonition}

\sphinxstepscope


\section{Download source and documentation as a tarball}
\label{\detokenize{Install_tarball:download-source-and-documentation-as-a-tarball}}\label{\detokenize{Install_tarball::doc}}
\sphinxAtStartPar
You can download a current LAMMPS tarball from the \sphinxhref{https://www.lammps.org/download.html}{download page}
of the \sphinxhref{https://www.lammps.org}{LAMMPS website} or from GitHub (see below).

\sphinxAtStartPar
You have two choices of tarballs, either the most recent stable release
or the most recent feature release.  Stable releases occur a few times
per year, and undergo more testing before release.  Also, between stable
releases bug fixes from the feature releases are back\sphinxhyphen{}ported and the
tarball occasionally updated.  Feature releases occur every 4 to 8
weeks.  The new contents in all feature releases are listed on the \sphinxhref{https://www.lammps.org/bug.html}{bug
and feature page} of the LAMMPS homepage.

\sphinxAtStartPar
Tarballs of older LAMMPS versions can also be downloaded from \sphinxhref{https://download.lammps.org/tars/}{this page}.

\sphinxAtStartPar
Tarballs downloaded from the LAMMPS homepage include the pre\sphinxhyphen{}translated
LAMMPS documentation (HTML and PDF files) corresponding to that version.

\sphinxAtStartPar
Once you have a tarball, uncompress and untar it with the following
command:

\begin{sphinxVerbatim}[commandchars=\\\{\}]
tar\PYG{+w}{ }\PYGZhy{}xzvf\PYG{+w}{ }lammps*.tar.gz
\end{sphinxVerbatim}

\sphinxAtStartPar
This will create a LAMMPS directory with the version date in its name,
e.g. \sphinxcode{\sphinxupquote{lammps\sphinxhyphen{}28Mar23}}.


\bigskip\hrule\bigskip


\sphinxAtStartPar
You can also download a compressed tar or zip archives from the
“Assets” sections of the \sphinxhref{https://github.com/lammps/lammps/releases}{LAMMPS GitHub releases site}.
The file name will be lammps\sphinxhyphen{}\textless{}version\textgreater{}.zip which can be unzipped
with the following command, to create a lammps\sphinxhyphen{}\textless{}version\textgreater{} directory:

\begin{sphinxVerbatim}[commandchars=\\\{\}]
unzip\PYG{+w}{ }lammps*.zip
\end{sphinxVerbatim}

\sphinxAtStartPar
This version corresponds to the selected LAMMPS feature or stable
release (as indicated by the matching git tag) and will only contain the
source code and no pre\sphinxhyphen{}built documentation.

\sphinxstepscope


\section{Download the LAMMPS source with git}
\label{\detokenize{Install_git:download-the-lammps-source-with-git}}\label{\detokenize{Install_git::doc}}
\sphinxAtStartPar
LAMMPS development is coordinated through the “LAMMPS GitHub site”.
If you clone the LAMMPS repository onto your local machine, it has
several advantages:
\begin{itemize}
\item {} 
\sphinxAtStartPar
You can stay current with changes to LAMMPS with a single git
command.

\item {} 
\sphinxAtStartPar
You can create your own development branches to add code to LAMMPS.

\item {} 
\sphinxAtStartPar
You can submit your new features back to GitHub for inclusion in
LAMMPS.  For that, you should first create your own {\hyperref[\detokenize{Howto_github::doc}]{\sphinxcrossref{\DUrole{doc}{fork on
GitHub}}}}, though.

\end{itemize}

\sphinxAtStartPar
You must have \sphinxhref{https://git-scm.com}{git} installed on your system to use the
commands explained below to communicate with the git servers on
GitHub.  For people still using subversion (svn), GitHub also
provides \sphinxhref{https://help.github.com/en/github/importing-your-projects-to-github/working-with-subversion-on-github}{limited support for subversion clients}.

\begin{sphinxadmonition}{note}{Note:}
\sphinxAtStartPar
As of October 2016, the official home of public LAMMPS development is
on GitHub.  The previously advertised LAMMPS git repositories on
git.lammps.org and bitbucket.org are now offline or deprecated.
\end{sphinxadmonition}

\sphinxAtStartPar
You can follow the LAMMPS development on 4 different git branches:
\begin{itemize}
\item {} 
\sphinxAtStartPar
\sphinxstylestrong{develop} : this branch follows the ongoing development and is
updated with every merge commit of a pull request

\item {} \begin{description}
\sphinxlineitem{\sphinxstylestrong{release}}{[}this branch is updated with every “feature release”;{]}
\sphinxAtStartPar
updates are always “fast\sphinxhyphen{}forward” merges from \sphinxstyleemphasis{develop}

\end{description}

\item {} 
\sphinxAtStartPar
\sphinxstylestrong{maintenance} : this branch collects back\sphinxhyphen{}ported bug fixes from the
\sphinxstyleemphasis{develop} branch to the \sphinxstyleemphasis{stable} branch.  It is used to update the
\sphinxstyleemphasis{stable} branch for “stable update releases”.

\item {} 
\sphinxAtStartPar
\sphinxstylestrong{stable} : this branch is updated from the \sphinxstyleemphasis{release} branch with
every “stable release” version and also has selected bug fixes with
every “update release” when the \sphinxstyleemphasis{maintenance} branch is merged into it

\end{itemize}

\sphinxAtStartPar
To access the git repositories on your box, use the clone command to
create a local copy of the LAMMPS repository with a command like:

\begin{sphinxVerbatim}[commandchars=\\\{\}]
git\PYG{+w}{ }clone\PYG{+w}{ }\PYGZhy{}b\PYG{+w}{ }release\PYG{+w}{ }https://github.com/lammps/lammps.git\PYG{+w}{ }mylammps
\end{sphinxVerbatim}

\sphinxAtStartPar
where “mylammps” is the name of the directory you wish to create on
your machine and “release” is one of the 3 branches listed above.
(Note that you actually download all 3 branches; you can switch
between them at any time using “git checkout \textless{}branch name\textgreater{}”.)

\begin{sphinxadmonition}{note}{Saving time and disk space when using \sphinxstyleliteralintitle{\sphinxupquote{git clone}}}

\sphinxAtStartPar
The complete git history of the LAMMPS project is quite large because
it contains the entire commit history of the project since fall 2006,
which includes the time when LAMMPS was managed with subversion.
This includes a few commits that have added and removed some large
files (mostly by accident).  If you do not need access to the entire
commit history (most people don’t), you can speed up the “cloning”
process and reduce local disk space requirements by using the
\sphinxcode{\sphinxupquote{\sphinxhyphen{}\sphinxhyphen{}depth}} git command line flag.  That will create a “shallow clone”
of the repository, which contains only a subset of the git history.
Using a depth of 1000 is usually sufficient to include the head
commits of the \sphinxstyleemphasis{develop}, the \sphinxstyleemphasis{release}, and the \sphinxstyleemphasis{maintenance}
branches.  To include the head commit of the \sphinxstyleemphasis{stable} branch you may
need a depth of up to 10000.  If you later need more of the git
history, you can always convert the shallow clone into a “full
clone”.
\end{sphinxadmonition}

\sphinxAtStartPar
Once the command completes, your directory will contain the same files
as if you unpacked a current LAMMPS tarball, with the exception, that
the HTML documentation files are not included. They can be generated
from the content provided in \sphinxcode{\sphinxupquote{doc/src}} by typing \sphinxcode{\sphinxupquote{make html}} from
the \sphinxcode{\sphinxupquote{doc}} directory.

\sphinxAtStartPar
After initial cloning, as bug fixes and new features are added to
LAMMPS you can stay up\sphinxhyphen{}to\sphinxhyphen{}date by typing the following git commands
from within the “mylammps” directory:

\begin{sphinxVerbatim}[commandchars=\\\{\}]
git\PYG{+w}{ }checkout\PYG{+w}{ }release\PYG{+w}{      }\PYG{c+c1}{\PYGZsh{} not needed if you always stay in this branch}
git\PYG{+w}{ }checkout\PYG{+w}{ }stable\PYG{+w}{       }\PYG{c+c1}{\PYGZsh{} use one of these 4 checkout commands}
git\PYG{+w}{ }checkout\PYG{+w}{ }develop\PYG{+w}{      }\PYG{c+c1}{\PYGZsh{} to choose the branch to follow}
git\PYG{+w}{ }checkout\PYG{+w}{ }maintenance
git\PYG{+w}{ }pull
\end{sphinxVerbatim}

\sphinxAtStartPar
Doing a “pull” will not change any files you have added to the LAMMPS
directory structure.  It will also not change any existing LAMMPS files
you have edited, unless those files have changed in the repository.  In
that case, git will attempt to merge the changes from the repository
file with your version of the file and tell you if there are any
conflicts.  See the git documentation for details.

\sphinxAtStartPar
If you want to access a particular previous release version of LAMMPS,
you can instead “check out” any version with a published tag. See the
output of \sphinxcode{\sphinxupquote{git tag \sphinxhyphen{}l}} for the list of tags.  The git command to do
this is as follows.

\begin{sphinxVerbatim}[commandchars=\\\{\}]
git\PYG{+w}{ }checkout\PYG{+w}{ }tagID
\end{sphinxVerbatim}

\sphinxAtStartPar
Stable versions and what tagID to use for a particular stable version
are discussed on \sphinxhref{https://www.lammps.org/bug.html\#version}{this page}.
Note that this command will print some warnings, because in order to get
back to the latest revision and to be able to update with \sphinxcode{\sphinxupquote{git pull}}
again, you will need to do \sphinxcode{\sphinxupquote{git checkout release}} (or
check out any other desired branch) first.

\sphinxAtStartPar
Once you have updated your local files with a \sphinxcode{\sphinxupquote{git pull}} (or \sphinxcode{\sphinxupquote{git
checkout}}), you still need to re\sphinxhyphen{}build LAMMPS if any source files have
changed.  How to do this depends on the build system you are using.
\begin{quote}

\begin{sphinxadmonition}{note}{CMake build}

\sphinxAtStartPar
Change to your build folder and type:

\begin{sphinxVerbatim}[commandchars=\\\{\}]
cmake\PYG{+w}{ }\PYGZhy{}\PYGZhy{}build\PYG{+w}{ }.
\end{sphinxVerbatim}

\sphinxAtStartPar
CMake should auto\sphinxhyphen{}detect whether it needs to re\sphinxhyphen{}run the CMake
configuration step and otherwise redo the build for all files
that have been changed or files that depend on changed files.
In case some build options have been changed or renamed, you
may have to update those by running:

\begin{sphinxVerbatim}[commandchars=\\\{\}]
cmake\PYG{+w}{ }.
\end{sphinxVerbatim}

\sphinxAtStartPar
and then rebuild.
\end{sphinxadmonition}

\begin{sphinxadmonition}{note}{Traditional make}

\sphinxAtStartPar
Switch to the src directory and type:

\begin{sphinxVerbatim}[commandchars=\\\{\}]
make\PYG{+w}{ }purge\PYG{+w}{             }\PYG{c+c1}{\PYGZsh{} remove any deprecated src files}
make\PYG{+w}{ }package\PYGZhy{}update\PYG{+w}{    }\PYG{c+c1}{\PYGZsh{} sync package files with src files}
make\PYG{+w}{ }foo\PYG{+w}{               }\PYG{c+c1}{\PYGZsh{} re\PYGZhy{}build for your machine (mpi, serial, etc)}
\end{sphinxVerbatim}

\sphinxAtStartPar
to enforce consistency of the source between the src folder
and package directories.  This is OK to do even if you don’t
use any packages. The \sphinxcode{\sphinxupquote{make purge}} command removes any deprecated
src files if they were removed by the update from a package
subdirectory.

\begin{sphinxadmonition}{warning}{Warning:}
\sphinxAtStartPar
If you wish to edit/change a src file that is from a package,
you should edit the version of the file inside the package
subdirectory with src, then re\sphinxhyphen{}install the package.  The
version in the source directory is merely a copy and will be
wiped out if you type “make package\sphinxhyphen{}update”.
\end{sphinxadmonition}
\end{sphinxadmonition}
\end{quote}

\begin{sphinxadmonition}{note}{Git protocols}

\sphinxAtStartPar
The servers at github.com support the “https” access protocol for
anonymous, read\sphinxhyphen{}only access.  If you have a suitably configured
GitHub account, you may also use SSH protocol with the URL
\sphinxcode{\sphinxupquote{git@github.com:lammps/lammps.git}}.
\end{sphinxadmonition}

\sphinxAtStartPar
The LAMMPS GitHub project is currently overseen by Axel Kohlmeyer
(Temple U, akohlmey at gmail.com).


\bigskip\hrule\bigskip


\sphinxAtStartPar
These are the files and subdirectories in the LAMMPS distribution:


\begin{savenotes}\sphinxattablestart
\sphinxthistablewithglobalstyle
\centering
\begin{tabulary}{\linewidth}[t]{TT}
\sphinxtoprule
\sphinxtableatstartofbodyhook
\sphinxAtStartPar
\sphinxcode{\sphinxupquote{README}}
&
\sphinxAtStartPar
Short description of the LAMMPS package
\\
\sphinxhline
\sphinxAtStartPar
\sphinxcode{\sphinxupquote{LICENSE}}
&
\sphinxAtStartPar
GNU General Public License (GPL)
\\
\sphinxhline
\sphinxAtStartPar
\sphinxcode{\sphinxupquote{SECURITY.md}}
&
\sphinxAtStartPar
Security policy for the LAMMPS package
\\
\sphinxhline
\sphinxAtStartPar
\sphinxcode{\sphinxupquote{bench}}
&
\sphinxAtStartPar
benchmark inputs
\\
\sphinxhline
\sphinxAtStartPar
\sphinxcode{\sphinxupquote{cmake}}
&
\sphinxAtStartPar
CMake build files
\\
\sphinxhline
\sphinxAtStartPar
\sphinxcode{\sphinxupquote{doc}}
&
\sphinxAtStartPar
documentation and tools to build the manual
\\
\sphinxhline
\sphinxAtStartPar
\sphinxcode{\sphinxupquote{examples}}
&
\sphinxAtStartPar
example input files
\\
\sphinxhline
\sphinxAtStartPar
\sphinxcode{\sphinxupquote{fortran}}
&
\sphinxAtStartPar
Fortran module for LAMMPS library interface
\\
\sphinxhline
\sphinxAtStartPar
\sphinxcode{\sphinxupquote{lib}}
&
\sphinxAtStartPar
additional provided or external libraries
\\
\sphinxhline
\sphinxAtStartPar
\sphinxcode{\sphinxupquote{potentials}}
&
\sphinxAtStartPar
selected interatomic potential files
\\
\sphinxhline
\sphinxAtStartPar
\sphinxcode{\sphinxupquote{python}}
&
\sphinxAtStartPar
Python module for LAMMPS library interface
\\
\sphinxhline
\sphinxAtStartPar
\sphinxcode{\sphinxupquote{src}}
&
\sphinxAtStartPar
LAMMPS source files
\\
\sphinxhline
\sphinxAtStartPar
\sphinxcode{\sphinxupquote{tools}}
&
\sphinxAtStartPar
pre\sphinxhyphen{} and post\sphinxhyphen{}processing tools
\\
\sphinxhline
\sphinxAtStartPar
\sphinxcode{\sphinxupquote{unittest}}
&
\sphinxAtStartPar
source code and inputs for testing LAMMPS
\\
\sphinxbottomrule
\end{tabulary}
\sphinxtableafterendhook\par
\sphinxattableend\end{savenotes}

\sphinxAtStartPar
You will have all of these if you downloaded the LAMMPS source code.
You will have only some of them if you downloaded executables, as
explained on the pages listed above.

\sphinxstepscope


\chapter{Build LAMMPS}
\label{\detokenize{Build:build-lammps}}\label{\detokenize{Build::doc}}
\sphinxAtStartPar
LAMMPS is built as a library and an executable from source code using
either traditional makefiles for use with GNU make (which may require
manual editing), or using a build environment generated by CMake (Unix
Makefiles, Ninja, Xcode, Visual Studio, KDevelop, CodeBlocks and more).

\sphinxAtStartPar
As an alternative, you can download a package with pre\sphinxhyphen{}built executables
or automated build trees, as described in the {\hyperref[\detokenize{Install::doc}]{\sphinxcrossref{\DUrole{doc}{Install}}}}
section of the manual.

\sphinxstepscope


\section{Build LAMMPS with CMake}
\label{\detokenize{Build_cmake:build-lammps-with-cmake}}\label{\detokenize{Build_cmake::doc}}
\sphinxAtStartPar
This page describes how to use \sphinxhref{https://cmake.org}{CMake} in general
to build LAMMPS.  Details for specific compile time settings and options
to enable and configure add\sphinxhyphen{}on packages are discussed with those
packages.  Links to those pages on the {\hyperref[\detokenize{Build::doc}]{\sphinxcrossref{\DUrole{doc}{Build overview}}}}
page.

\sphinxAtStartPar
The following text assumes some familiarity with CMake and focuses on
using the command line tool \sphinxcode{\sphinxupquote{cmake}} and what settings are supported
for building LAMMPS.  A more detailed tutorial on how to use CMake
itself, the text mode or graphical user interface, to change the
generated output files for different build tools and development
environments is on a {\hyperref[\detokenize{Howto_cmake::doc}]{\sphinxcrossref{\DUrole{doc}{separate page}}}}.

\begin{sphinxadmonition}{note}{Note:}
\sphinxAtStartPar
LAMMPS currently requires that CMake version 3.16 or later is available.
\end{sphinxadmonition}

\begin{sphinxadmonition}{warning}{Warning:}
\sphinxAtStartPar
You must not mix the {\hyperref[\detokenize{Build_make::doc}]{\sphinxcrossref{\DUrole{doc}{traditional make based}}}}
LAMMPS build procedure with using CMake.  No packages may be
installed or a build been previously attempted in the LAMMPS source
directory by using \sphinxcode{\sphinxupquote{make \textless{}machine\textgreater{}}}.  CMake will detect if this is
the case and generate an error.  To remove conflicting files from the
\sphinxcode{\sphinxupquote{src}} you can use the command \sphinxcode{\sphinxupquote{make no\sphinxhyphen{}all purge}} which will
uninstall all packages and delete all auto\sphinxhyphen{}generated files.
\end{sphinxadmonition}


\subsection{Advantages of using CMake}
\label{\detokenize{Build_cmake:advantages-of-using-cmake}}
\sphinxAtStartPar
CMake is an alternative to compiling LAMMPS in the traditional way
through {\hyperref[\detokenize{Build_make::doc}]{\sphinxcrossref{\DUrole{doc}{(manually customized) makefiles}}}}.  Using
CMake has multiple advantages that are specifically helpful for
people with limited experience in compiling software or for people
that want to modify or extend LAMMPS.
\begin{itemize}
\item {} 
\sphinxAtStartPar
CMake can detect available hardware, tools, features, and libraries
and adapt the LAMMPS default build configuration accordingly.

\item {} 
\sphinxAtStartPar
CMake can generate files for different build tools and integrated
development environments (IDE).

\item {} 
\sphinxAtStartPar
CMake supports customization of settings with a command line, text
mode, or graphical user interface.  No manual editing of files,
knowledge of file formats or complex command line syntax is required.

\item {} 
\sphinxAtStartPar
All enabled components are compiled in a single build operation.

\item {} 
\sphinxAtStartPar
Automated dependency tracking for all files and configuration options.

\item {} 
\sphinxAtStartPar
Support for true out\sphinxhyphen{}of\sphinxhyphen{}source compilation. Multiple configurations
and settings with different choices of LAMMPS packages, settings, or
compilers can be configured and built concurrently from the same
source tree.

\item {} 
\sphinxAtStartPar
Simplified packaging of LAMMPS for Linux distributions, environment
modules, or automated build tools like \sphinxhref{https://brew.sh/}{Homebrew}.

\item {} 
\sphinxAtStartPar
Integration of automated unit and regression testing (the LAMMPS side
of this is still under active development).

\end{itemize}


\subsection{Getting started}
\label{\detokenize{Build_cmake:getting-started}}\label{\detokenize{Build_cmake:cmake-build}}
\sphinxAtStartPar
Building LAMMPS with CMake is a two\sphinxhyphen{}step process.  In the first step,
you use CMake to generate a build environment in a new directory.  For
that purpose you can use either the command\sphinxhyphen{}line utility \sphinxcode{\sphinxupquote{cmake}} (or
\sphinxcode{\sphinxupquote{cmake3}}), the text\sphinxhyphen{}mode UI utility \sphinxcode{\sphinxupquote{ccmake}} (or \sphinxcode{\sphinxupquote{ccmake3}}) or the
graphical utility \sphinxcode{\sphinxupquote{cmake\sphinxhyphen{}gui}}, or use them interchangeably.  The
second step is then the compilation and linking of all objects,
libraries, and executables using the selected build tool.  Here is a
minimal example using the command line version of CMake to build LAMMPS
with no add\sphinxhyphen{}on packages enabled and no customization:

\begin{sphinxVerbatim}[commandchars=\\\{\}]
\PYG{n+nb}{cd}\PYG{+w}{ }lammps\PYG{+w}{                }\PYG{c+c1}{\PYGZsh{} change to the LAMMPS distribution directory}
mkdir\PYG{+w}{ }build\PYG{p}{;}\PYG{+w}{ }\PYG{n+nb}{cd}\PYG{+w}{ }build\PYG{+w}{    }\PYG{c+c1}{\PYGZsh{} create and use a build directory}
cmake\PYG{+w}{ }../cmake\PYG{+w}{           }\PYG{c+c1}{\PYGZsh{} configuration reading CMake scripts from ../cmake}
cmake\PYG{+w}{ }\PYGZhy{}\PYGZhy{}build\PYG{+w}{ }.\PYG{+w}{          }\PYG{c+c1}{\PYGZsh{} compilation (or type \PYGZdq{}make\PYGZdq{})}
\end{sphinxVerbatim}

\sphinxAtStartPar
This will create and change into a folder called \sphinxcode{\sphinxupquote{build}}, then run the
configuration step to generate build files for the default build command
and then launch that build command to compile LAMMPS.  During the
configuration step CMake will try to detect whether support for MPI,
OpenMP, FFTW, gzip, JPEG, PNG, and ffmpeg are available and enable the
corresponding configuration settings.  The progress of this
configuration can be followed on the screen and a summary of selected
options and settings will be printed at the end.  The \sphinxcode{\sphinxupquote{cmake \sphinxhyphen{}\sphinxhyphen{}build
.}} command will launch the compilation, which, if successful, will
ultimately produce a library \sphinxcode{\sphinxupquote{liblammps.a}} and the LAMMPS executable
\sphinxcode{\sphinxupquote{lmp}} inside the \sphinxcode{\sphinxupquote{build}} folder.

\sphinxAtStartPar
Compilation can take a long time, since LAMMPS is a large project with
many features. If your machine has multiple CPU cores (most do these
days), you can speed this up by compiling sources in parallel with
\sphinxcode{\sphinxupquote{make \sphinxhyphen{}j N}} (with N being the maximum number of concurrently executed
tasks).  Installation of the \sphinxhref{https://ccache.dev/}{ccache} (= Compiler
Cache) software may speed up repeated compilation even more, e.g. during
code development, especially when repeatedly switching between branches.

\sphinxAtStartPar
After the initial build, whenever you edit LAMMPS source files, enable
or disable packages, change compiler flags or build options, you must
re\sphinxhyphen{}compile and relink the LAMMPS executable with \sphinxcode{\sphinxupquote{cmake \sphinxhyphen{}\sphinxhyphen{}build .}} (or
\sphinxcode{\sphinxupquote{make}}).  If the compilation fails for some reason, try running
\sphinxcode{\sphinxupquote{cmake .}} and then compile again. The included dependency tracking
should make certain that only the necessary subset of files is
re\sphinxhyphen{}compiled.  You can also delete compiled objects, libraries, and
executables with \sphinxcode{\sphinxupquote{cmake \sphinxhyphen{}\sphinxhyphen{}build . \sphinxhyphen{}\sphinxhyphen{}target clean}} (or \sphinxcode{\sphinxupquote{make clean}}).

\sphinxAtStartPar
After compilation, you may optionally install the LAMMPS executable into
your system with:

\begin{sphinxVerbatim}[commandchars=\\\{\}]
make\PYG{+w}{ }install\PYG{+w}{    }\PYG{c+c1}{\PYGZsh{} optional, copy compiled files into installation location}
\end{sphinxVerbatim}

\sphinxAtStartPar
This will install the LAMMPS executable and library, some tools (if
configured) and additional files like LAMMPS API headers, manpages,
potential and force field files.  The location of the installation tree
defaults to \sphinxcode{\sphinxupquote{\$\{HOME\}/.local}}.


\subsection{Configuration and build options}
\label{\detokenize{Build_cmake:configuration-and-build-options}}\label{\detokenize{Build_cmake:cmake-options}}
\sphinxAtStartPar
The CMake commands have one mandatory argument: a folder containing a
file called \sphinxcode{\sphinxupquote{CMakeLists.txt}} (for LAMMPS it is located in the
\sphinxcode{\sphinxupquote{cmake}} folder) or a build folder containing a file called
\sphinxcode{\sphinxupquote{CMakeCache.txt}}, which is generated at the end of the CMake
configuration step.  The cache file contains all current CMake settings.

\sphinxAtStartPar
To modify settings, enable or disable features, you need to set
\sphinxstyleemphasis{variables} with either the \sphinxcode{\sphinxupquote{\sphinxhyphen{}D}} command line flag (\sphinxcode{\sphinxupquote{\sphinxhyphen{}D
VARIABLE1\_NAME=value}}) or change them in the text mode of the graphical
user interface.  The \sphinxcode{\sphinxupquote{\sphinxhyphen{}D}} flag can be used several times in one command.

\sphinxAtStartPar
For your convenience, we provide {\hyperref[\detokenize{Build_package:cmake-presets}]{\sphinxcrossref{\DUrole{std,std-ref}{CMake presets}}}}
that combine multiple settings to enable optional LAMMPS packages or use
a different compiler tool chain.  Those are loaded with the \sphinxcode{\sphinxupquote{\sphinxhyphen{}C}} flag
(\sphinxcode{\sphinxupquote{\sphinxhyphen{}C ../cmake/presets/basic.cmake}}).  This step would only be needed
once, as the settings from the preset files are stored in the
\sphinxcode{\sphinxupquote{CMakeCache.txt}} file. It is also possible to customize the build
by adding one or more \sphinxcode{\sphinxupquote{\sphinxhyphen{}D}} flags to the CMake command line.

\sphinxAtStartPar
Generating files for alternate build tools (e.g. Ninja) and project files
for IDEs like Eclipse, CodeBlocks, or Kate can be selected using the \sphinxcode{\sphinxupquote{\sphinxhyphen{}G}}
command line flag.  A list of available generator settings for your
specific CMake version is given when running \sphinxcode{\sphinxupquote{cmake \sphinxhyphen{}\sphinxhyphen{}help}}.


\subsection{Multi\sphinxhyphen{}configuration build systems}
\label{\detokenize{Build_cmake:multi-configuration-build-systems}}\label{\detokenize{Build_cmake:cmake-multiconfig}}
\sphinxAtStartPar
Throughout this manual, it is mostly assumed that LAMMPS is being built
on a Unix\sphinxhyphen{}like operating system with “make” as the underlying “builder”,
since this is the most common case.  In this case the build
“configuration” is chose using \sphinxcode{\sphinxupquote{\sphinxhyphen{}D CMAKE\_BUILD\_TYPE=\textless{}configuration\textgreater{}}}
with \sphinxcode{\sphinxupquote{\textless{}configuration\textgreater{}}} being one of “Release”, “Debug”,
“RelWithDebInfo”, or “MinSizeRel”.  Some build tools, however, can also
use or even require having a so\sphinxhyphen{}called multi\sphinxhyphen{}configuration build system
setup.  For a multi\sphinxhyphen{}configuration build, the built type (or
configuration) is selected at compile time using the same build
files. E.g.  with:

\begin{sphinxVerbatim}[commandchars=\\\{\}]
cmake\PYG{+w}{ }\PYGZhy{}\PYGZhy{}build\PYG{+w}{ }build\PYGZhy{}multi\PYG{+w}{ }\PYGZhy{}\PYGZhy{}config\PYG{+w}{ }Release
\end{sphinxVerbatim}

\sphinxAtStartPar
In that case the resulting binaries are not in the build folder directly
but in subdirectories corresponding to the build type (i.e. Release in
the example from above).  Similarly, for running unit tests the
configuration is selected with the \sphinxcode{\sphinxupquote{\sphinxhyphen{}C}} flag:

\begin{sphinxVerbatim}[commandchars=\\\{\}]
ctest\PYG{+w}{ }\PYGZhy{}C\PYG{+w}{ }Debug
\end{sphinxVerbatim}

\sphinxAtStartPar
The CMake scripts in LAMMPS have basic support for being compiled using
a multi\sphinxhyphen{}config build system, but not all of it has been ported.  This is
in particular applicable to compiling packages that require additional
libraries that would be downloaded and compiled by CMake.  The
\sphinxcode{\sphinxupquote{windows.cmake}} preset file tries to keep track of which packages can
be compiled natively with the MSVC compilers out\sphinxhyphen{}of\sphinxhyphen{}the box.  Not all of
the external libraries are portable to Windows, either.


\subsection{Installing CMake}
\label{\detokenize{Build_cmake:installing-cmake}}
\sphinxAtStartPar
Check if your machine already has CMake installed:

\begin{sphinxVerbatim}[commandchars=\\\{\}]
which\PYG{+w}{ }cmake\PYG{+w}{             }\PYG{c+c1}{\PYGZsh{} do you have it?}
which\PYG{+w}{ }cmake3\PYG{+w}{            }\PYG{c+c1}{\PYGZsh{} version 3 may have this name}
cmake\PYG{+w}{ }\PYGZhy{}\PYGZhy{}version\PYG{+w}{         }\PYG{c+c1}{\PYGZsh{} what specific version you have}
\end{sphinxVerbatim}

\sphinxAtStartPar
On clusters or supercomputers which use environment modules to manage
software packages, do this:

\begin{sphinxVerbatim}[commandchars=\\\{\}]
module\PYG{+w}{ }list\PYG{+w}{            }\PYG{c+c1}{\PYGZsh{} is a module for cmake already loaded?}
module\PYG{+w}{ }avail\PYG{+w}{           }\PYG{c+c1}{\PYGZsh{} is a module for cmake available?}
module\PYG{+w}{ }load\PYG{+w}{ }cmake\PYG{+w}{      }\PYG{c+c1}{\PYGZsh{} load cmake module with appropriate name}
\end{sphinxVerbatim}

\sphinxAtStartPar
Most Linux distributions offer pre\sphinxhyphen{}compiled cmake packages through their
package management system. If you do not have CMake or a recent enough
version (Note: for CentOS 7.x you need to enable the EPEL repository),
you can download the latest version from \sphinxurl{https://cmake.org/download/}.  Instructions on how to install it on
various platforms can be found \sphinxhref{https://cmake.org/install/}{on this page}.

\sphinxstepscope


\section{Build LAMMPS with make}
\label{\detokenize{Build_make:build-lammps-with-make}}\label{\detokenize{Build_make::doc}}
\sphinxAtStartPar
Building LAMMPS with traditional makefiles requires that you have a
\sphinxcode{\sphinxupquote{Makefile.\textless{}machine\textgreater{}}} file appropriate for your system in either the
\sphinxcode{\sphinxupquote{src/MAKE}}, \sphinxcode{\sphinxupquote{src/MAKE/MACHINES}}, \sphinxcode{\sphinxupquote{src/MAKE/OPTIONS}}, or
\sphinxcode{\sphinxupquote{src/MAKE/MINE}} directory (see below).  It can include various options
for customizing your LAMMPS build with a number of global compilation
options and features.


\subsection{Requirements}
\label{\detokenize{Build_make:requirements}}
\sphinxAtStartPar
Those makefiles are written for and tested with GNU make and may not
be compatible with other make programs.  In most cases, if the “make”
program is not GNU make, then there will be a GNU make program
available under the name “gmake”.  If GNU make or a compatible make is
not available, you may have to first install it or switch to building
with {\hyperref[\detokenize{Build_cmake::doc}]{\sphinxcrossref{\DUrole{doc}{CMake}}}}.  The makefiles of the traditional
make based build process and the scripts they are calling expect a few
additional tools to be available and functioning.
\begin{itemize}
\item {} 
\sphinxAtStartPar
A working C/C++ compiler toolchain supporting the C++11 standard; on
Linux, these are often the GNU compilers. Some older compiler versions
require adding flags like \sphinxcode{\sphinxupquote{\sphinxhyphen{}std=c++11}} to enable the C++11 mode.

\item {} 
\sphinxAtStartPar
A Bourne shell compatible “Unix” shell program (frequently this is \sphinxcode{\sphinxupquote{bash}})

\item {} 
\sphinxAtStartPar
A few shell utilities: \sphinxcode{\sphinxupquote{ls}}, \sphinxcode{\sphinxupquote{mv}}, \sphinxcode{\sphinxupquote{ln}}, \sphinxcode{\sphinxupquote{rm}}, \sphinxcode{\sphinxupquote{grep}}, \sphinxcode{\sphinxupquote{sed}}, \sphinxcode{\sphinxupquote{tr}}, \sphinxcode{\sphinxupquote{cat}}, \sphinxcode{\sphinxupquote{touch}}, \sphinxcode{\sphinxupquote{diff}}, \sphinxcode{\sphinxupquote{dirname}}

\item {} 
\sphinxAtStartPar
Python (optional, required for \sphinxcode{\sphinxupquote{make lib\sphinxhyphen{}\textless{}pkg\textgreater{}}} in the \sphinxcode{\sphinxupquote{src}}
folder).  Python scripts are currently tested with python 2.7 and
3.6 to 3.11. The procedure for {\hyperref[\detokenize{Build_manual::doc}]{\sphinxcrossref{\DUrole{doc}{building the documentation}}}} \sphinxstyleemphasis{requires} Python 3.5 or later.

\end{itemize}


\subsection{Getting started}
\label{\detokenize{Build_make:getting-started}}
\sphinxAtStartPar
To include LAMMPS packages (i.e. optional commands and styles) you must
enable (or “install”) them first, as discussed on the {\hyperref[\detokenize{Build_package::doc}]{\sphinxcrossref{\DUrole{doc}{Build
package}}}} page.  If a package requires (provided or
external) libraries, you must configure and build those libraries
\sphinxstylestrong{before} building LAMMPS itself and especially \sphinxstylestrong{before} enabling
such a package with \sphinxcode{\sphinxupquote{make yes\sphinxhyphen{}\textless{}package\textgreater{}}}.  {\hyperref[\detokenize{Build_cmake::doc}]{\sphinxcrossref{\DUrole{doc}{Building LAMMPS with
CMake}}}} can automate much of this for many types of
machines, especially workstations, desktops, and laptops, so we suggest
you try it first when building LAMMPS in those cases.

\sphinxAtStartPar
The commands below perform a default LAMMPS build, producing the LAMMPS
executable \sphinxcode{\sphinxupquote{lmp\_serial}} and \sphinxcode{\sphinxupquote{lmp\_mpi}} in \sphinxcode{\sphinxupquote{lammps/src}}:

\begin{sphinxVerbatim}[commandchars=\\\{\}]
\PYG{n+nb}{cd}\PYG{+w}{ }lammps/src\PYG{+w}{   }\PYG{c+c1}{\PYGZsh{} change to main LAMMPS source folder}
make\PYG{+w}{ }serial\PYG{+w}{     }\PYG{c+c1}{\PYGZsh{} build a serial LAMMPS executable using GNU g++}
make\PYG{+w}{ }mpi\PYG{+w}{        }\PYG{c+c1}{\PYGZsh{} build a parallel LAMMPS executable with MPI}
make\PYG{+w}{            }\PYG{c+c1}{\PYGZsh{} see a variety of make options}
\end{sphinxVerbatim}

\sphinxAtStartPar
Compilation can take a long time, since LAMMPS is a large project with
many features. If your machine has multiple CPU cores (most do these
days), you can speed this up by compiling sources in parallel with
\sphinxcode{\sphinxupquote{make \sphinxhyphen{}j N}} (with N being the maximum number of concurrently executed
tasks).  Installation of the \sphinxhref{https://ccache.dev/}{ccache} (= Compiler
Cache) software may speed up repeated compilation even more, e.g. during
code development, especially when repeatedly switching between branches.

\sphinxAtStartPar
After the initial build, whenever you edit LAMMPS source files, or add
or remove new files to the source directory (e.g. by installing or
uninstalling packages), you must re\sphinxhyphen{}compile and relink the LAMMPS
executable with the same \sphinxcode{\sphinxupquote{make \textless{}machine\textgreater{}}} command.  The makefile’s
dependency tracking should ensure that only the necessary subset of
files is re\sphinxhyphen{}compiled.  If you change settings in the makefile, you have
to recompile \sphinxstyleemphasis{everything}.  To delete all objects, you can use \sphinxcode{\sphinxupquote{make
clean\sphinxhyphen{}\textless{}machine\textgreater{}}}.

\begin{sphinxadmonition}{note}{Note:}
\sphinxAtStartPar
Before the actual compilation starts, LAMMPS will perform several
steps to collect information from the configuration and setup that is
then embedded into the executable.  When you build LAMMPS for the
first time, it will also compile a tool to quickly determine a list
of dependencies.  Those are required for the make program to
correctly detect, which files need to be recompiled or relinked
after changes were made to the sources.
\end{sphinxadmonition}


\subsection{Customized builds and alternate makefiles}
\label{\detokenize{Build_make:customized-builds-and-alternate-makefiles}}
\sphinxAtStartPar
The \sphinxcode{\sphinxupquote{src/MAKE}} directory tree contains the \sphinxcode{\sphinxupquote{Makefile.\textless{}machine\textgreater{}}}
files included in the LAMMPS distribution.  Typing \sphinxcode{\sphinxupquote{make example}} uses
\sphinxcode{\sphinxupquote{Makefile.example}} from one of those folders, if available.  The
\sphinxcode{\sphinxupquote{make serial}} and \sphinxcode{\sphinxupquote{make mpi}} lines above, for example, use
\sphinxcode{\sphinxupquote{src/MAKE/Makefile.serial}} and \sphinxcode{\sphinxupquote{src/MAKE/Makefile.mpi}},
respectively.  Other makefiles are in these directories:

\begin{sphinxVerbatim}[commandchars=\\\{\}]
OPTIONS\PYG{+w}{      }\PYG{c+c1}{\PYGZsh{} Makefiles which enable specific options}
MACHINES\PYG{+w}{     }\PYG{c+c1}{\PYGZsh{} Makefiles for specific machines}
MINE\PYG{+w}{         }\PYG{c+c1}{\PYGZsh{} customized Makefiles you create (you may need to create this folder)}
\end{sphinxVerbatim}

\sphinxAtStartPar
Simply typing \sphinxcode{\sphinxupquote{make}} lists all the available \sphinxcode{\sphinxupquote{Makefile.\textless{}machine\textgreater{}}}
files with a single line description toward the end of the output.  A
file with the same name can appear in multiple folders (not a good
idea).  The order the directories are searched is as follows:
\sphinxcode{\sphinxupquote{src/MAKE/MINE}}, \sphinxcode{\sphinxupquote{src/MAKE}}, \sphinxcode{\sphinxupquote{src/MAKE/OPTIONS}},
\sphinxcode{\sphinxupquote{src/MAKE/MACHINES}}.  This gives preference to a customized file you
put in \sphinxcode{\sphinxupquote{src/MAKE/MINE}}.  If you create your own custom makefile under
a new name, please edit the first line with the description and machine
name, so you will not confuse yourself, when looking at the machine
summary.

\sphinxAtStartPar
Makefiles you may wish to try out, include those listed below (some
require a package first be installed).  Many of these include specific
compiler flags for optimized performance.  Please note, however, that
some of these customized machine Makefile are contributed by users, and
thus may have modifications specific to the systems of those users.
Since compilers, OS configurations, and LAMMPS itself keep changing,
their settings may become outdated, too:

\begin{sphinxVerbatim}[commandchars=\\\{\}]
make\PYG{+w}{ }mac\PYG{+w}{             }\PYG{c+c1}{\PYGZsh{} build serial LAMMPS on macOS}
make\PYG{+w}{ }mac\PYGZus{}mpi\PYG{+w}{         }\PYG{c+c1}{\PYGZsh{} build parallel LAMMPS on macOS}
make\PYG{+w}{ }intel\PYGZus{}cpu\PYG{+w}{       }\PYG{c+c1}{\PYGZsh{} build with the INTEL package optimized for CPUs}
make\PYG{+w}{ }knl\PYG{+w}{             }\PYG{c+c1}{\PYGZsh{} build with the INTEL package optimized for KNLs}
make\PYG{+w}{ }opt\PYG{+w}{             }\PYG{c+c1}{\PYGZsh{} build with the OPT package optimized for CPUs}
make\PYG{+w}{ }omp\PYG{+w}{             }\PYG{c+c1}{\PYGZsh{} build with the OPENMP package optimized for OpenMP}
make\PYG{+w}{ }kokkos\PYGZus{}omp\PYG{+w}{      }\PYG{c+c1}{\PYGZsh{} build with the KOKKOS package for OpenMP}
make\PYG{+w}{ }kokkos\PYGZus{}cuda\PYGZus{}mpi\PYG{+w}{ }\PYG{c+c1}{\PYGZsh{} build with the KOKKOS package for GPUs}
make\PYG{+w}{ }kokkos\PYGZus{}phi\PYG{+w}{      }\PYG{c+c1}{\PYGZsh{} build with the KOKKOS package for KNLs}
\end{sphinxVerbatim}

\sphinxstepscope


\section{Link LAMMPS as a library to another code}
\label{\detokenize{Build_link:link-lammps-as-a-library-to-another-code}}\label{\detokenize{Build_link::doc}}
\sphinxAtStartPar
LAMMPS is designed as a library of C++ objects that can be integrated
into other applications, including Python scripts.  The files
\sphinxcode{\sphinxupquote{src/library.cpp}} and \sphinxcode{\sphinxupquote{src/library.h}} define a C\sphinxhyphen{}style API for using
LAMMPS as a library.  See the {\hyperref[\detokenize{Howto_library::doc}]{\sphinxcrossref{\DUrole{doc}{Library interface to LAMMPS}}}} page for a
description of the interface and how to use it for your needs.

\sphinxAtStartPar
The {\hyperref[\detokenize{Build_basics::doc}]{\sphinxcrossref{\DUrole{doc}{Basic build options}}}} page explains how to build LAMMPS as either a
shared or static library.  This results in a file in the compilation
folder called \sphinxcode{\sphinxupquote{liblammps.a}} or \sphinxcode{\sphinxupquote{liblammps\_\textless{}name\textgreater{}.a}} in case of
building a static library.  In case of a shared library, the name is the
same only that the suffix is going to be either \sphinxcode{\sphinxupquote{.so}} or \sphinxcode{\sphinxupquote{.dylib}} or
\sphinxcode{\sphinxupquote{.dll}} instead of \sphinxcode{\sphinxupquote{.a}} depending on the OS.  In some cases, the
\sphinxcode{\sphinxupquote{.so}} file may be a symbolic link to a file with the suffix \sphinxcode{\sphinxupquote{.so.0}}
(or some other number).

\begin{sphinxadmonition}{note}{Note:}
\sphinxAtStartPar
Care should be taken to use the same MPI library for the calling code
and the LAMMPS library, unless LAMMPS is to be compiled without (real)
MPI support using the included STUBS MPI library.
\end{sphinxadmonition}


\subsection{Link with LAMMPS as a static library}
\label{\detokenize{Build_link:link-with-lammps-as-a-static-library}}
\sphinxAtStartPar
The calling application can link to LAMMPS as a static library with
compilation and link commands, as in the examples shown below.  These
are examples for a code written in C in the file \sphinxcode{\sphinxupquote{caller.c}}.
The benefit of linking to a static library is, that the resulting
executable is independent of that library since all required
executable code from the library is copied into the calling executable.
\begin{quote}

\begin{sphinxadmonition}{note}{CMake build}

\sphinxAtStartPar
This assumes that LAMMPS has been configured without setting a
\sphinxcode{\sphinxupquote{LAMMPS\_MACHINE}} name, installed with \sphinxcode{\sphinxupquote{make install}}, and the
\sphinxcode{\sphinxupquote{PKG\_CONFIG\_PATH}} environment variable has been updated to
include the \sphinxcode{\sphinxupquote{liblammps.pc}} file installed into the configured
destination folder.  The commands to compile and link a coupled
executable are then:

\begin{sphinxVerbatim}[commandchars=\\\{\}]
mpicc\PYG{+w}{ }\PYGZhy{}c\PYG{+w}{ }\PYGZhy{}O\PYG{+w}{ }\PYG{k}{\PYGZdl{}(}pkg\PYGZhy{}config\PYG{+w}{ }\PYGZhy{}\PYGZhy{}cflags\PYG{+w}{ }liblammps\PYG{k}{)}\PYG{+w}{ }caller.c
mpicxx\PYG{+w}{ }\PYGZhy{}o\PYG{+w}{ }\PYG{n+nb}{caller}\PYG{+w}{ }caller.o\PYG{+w}{ }\PYGZhy{}\PYG{k}{\PYGZdl{}(}pkg\PYGZhy{}config\PYG{+w}{ }\PYGZhy{}\PYGZhy{}libs\PYG{+w}{ }liblammps\PYG{k}{)}
\end{sphinxVerbatim}
\end{sphinxadmonition}

\begin{sphinxadmonition}{note}{Traditional make}

\sphinxAtStartPar
This assumes that LAMMPS has been compiled in the folder
\sphinxcode{\sphinxupquote{\$\{HOME\}/lammps/src}} with “make mpi”. The commands to compile
and link a coupled executable are then:

\begin{sphinxVerbatim}[commandchars=\\\{\}]
mpicc\PYG{+w}{ }\PYGZhy{}c\PYG{+w}{ }\PYGZhy{}O\PYG{+w}{ }\PYGZhy{}I\PYG{l+s+si}{\PYGZdl{}\PYGZob{}}\PYG{n+nv}{HOME}\PYG{l+s+si}{\PYGZcb{}}/lammps/src\PYG{+w}{ }caller.c
mpicxx\PYG{+w}{ }\PYGZhy{}o\PYG{+w}{ }\PYG{n+nb}{caller}\PYG{+w}{ }caller.o\PYG{+w}{ }\PYGZhy{}L\PYG{l+s+si}{\PYGZdl{}\PYGZob{}}\PYG{n+nv}{HOME}\PYG{l+s+si}{\PYGZcb{}}/lammps/src\PYG{+w}{ }\PYGZhy{}llammps\PYGZus{}mpi
\end{sphinxVerbatim}

\sphinxAtStartPar
The \sphinxcode{\sphinxupquote{\sphinxhyphen{}I}} argument is the path to the location of the \sphinxcode{\sphinxupquote{library.h}}
header file containing the interface to the LAMMPS C\sphinxhyphen{}style library
interface.  The \sphinxcode{\sphinxupquote{\sphinxhyphen{}L}} argument is the path to where the
\sphinxcode{\sphinxupquote{liblammps\_mpi.a}} file is located.  The \sphinxcode{\sphinxupquote{\sphinxhyphen{}llammps\_mpi}} argument
is shorthand for telling the compiler to link the file
\sphinxcode{\sphinxupquote{liblammps\_mpi.a}}.  If LAMMPS has been built as a shared
library, then the linker will use \sphinxcode{\sphinxupquote{liblammps\_mpi.so}} instead.
If both files are available, the linker will usually prefer the
shared library.  In case of a shared library, you may need to
update the \sphinxcode{\sphinxupquote{LD\_LIBRARY\_PATH}} environment variable or running the
\sphinxcode{\sphinxupquote{caller}} executable will fail since it cannot find the shared
library at runtime.
\end{sphinxadmonition}
\end{quote}

\sphinxAtStartPar
However, it is only as simple as shown above for the case of a plain
LAMMPS library without any optional packages that depend on libraries
(bundled or external) or when using a shared library.  Otherwise, you
need to include all flags, libraries, and paths for the coupled
executable, that are also required to link the LAMMPS executable.
\begin{quote}

\begin{sphinxadmonition}{note}{CMake build}

\sphinxAtStartPar
When using CMake, additional libraries with sources in the lib
folder are built, but not included in \sphinxcode{\sphinxupquote{liblammps.a}} and
(currently) not installed with \sphinxcode{\sphinxupquote{make install}} and not included
in the \sphinxcode{\sphinxupquote{pkgconfig}} configuration file.  They can be found in the
top level build folder, but you have to determine the necessary
link flags manually.  It is therefore recommended to either use
the traditional make procedure to build and link with a static
library or build and link with a shared library instead.
\end{sphinxadmonition}

\begin{sphinxadmonition}{note}{Traditional make}

\sphinxAtStartPar
After you have compiled a static LAMMPS library using the
conventional build system for example with “make mode=static
serial”. And you also have installed the \sphinxcode{\sphinxupquote{POEMS}} package after
building its bundled library in \sphinxcode{\sphinxupquote{lib/poems}}. Then the commands
to build and link the coupled executable change to:

\begin{sphinxVerbatim}[commandchars=\\\{\}]
gcc\PYG{+w}{ }\PYGZhy{}c\PYG{+w}{ }\PYGZhy{}O\PYG{+w}{ }\PYGZhy{}I\PYG{l+s+si}{\PYGZdl{}\PYGZob{}}\PYG{n+nv}{HOME}\PYG{l+s+si}{\PYGZcb{}}/lammps/src\PYG{+w}{ }\PYGZhy{}caller.c
g++\PYG{+w}{ }\PYGZhy{}o\PYG{+w}{ }\PYG{n+nb}{caller}\PYG{+w}{ }caller.o\PYG{+w}{ }\PYGZhy{}L\PYG{l+s+si}{\PYGZdl{}\PYGZob{}}\PYG{n+nv}{HOME}\PYG{l+s+si}{\PYGZcb{}}/lammps/lib/poems\PYG{+w}{ }\PYG{l+s+se}{\PYGZbs{}}
\PYG{+w}{             }\PYGZhy{}L\PYG{l+s+si}{\PYGZdl{}\PYGZob{}}\PYG{n+nv}{HOME}\PYG{l+s+si}{\PYGZcb{}}/lammps/src/STUBS\PYG{+w}{ }\PYGZhy{}L\PYG{l+s+si}{\PYGZdl{}\PYGZob{}}\PYG{n+nv}{HOME}\PYG{l+s+si}{\PYGZcb{}}/lammps/src\PYG{+w}{ }\PYG{l+s+se}{\PYGZbs{}}
\PYG{+w}{             }\PYGZhy{}llammps\PYGZus{}serial\PYG{+w}{ }\PYGZhy{}lpoems\PYG{+w}{ }\PYGZhy{}lmpi\PYGZus{}stubs
\end{sphinxVerbatim}

\sphinxAtStartPar
Note, that you need to link with \sphinxcode{\sphinxupquote{g++}} instead of \sphinxcode{\sphinxupquote{gcc}} even
if you have written your code in C, since LAMMPS itself is C++
code.  You can display the currently applied settings for building
LAMMPS for the “serial” machine target by using the command:

\begin{sphinxVerbatim}[commandchars=\\\{\}]
make\PYG{+w}{ }\PYG{n+nv}{mode}\PYG{o}{=}print\PYG{+w}{ }serial
\end{sphinxVerbatim}

\sphinxAtStartPar
Which should output something like:

\begin{sphinxVerbatim}[commandchars=\\\{\}]
\PYG{c+c1}{\PYGZsh{} Compiler:}
\PYG{n+nv}{CXX}\PYG{o}{=}g++
\PYG{c+c1}{\PYGZsh{} Linker:}
\PYG{n+nv}{LD}\PYG{o}{=}g++
\PYG{c+c1}{\PYGZsh{} Compilation:}
\PYG{n+nv}{CXXFLAGS}\PYG{o}{=}\PYGZhy{}g\PYG{+w}{ }\PYGZhy{}O3\PYG{+w}{ }\PYGZhy{}DLAMMPS\PYGZus{}GZIP\PYG{+w}{ }\PYGZhy{}DLAMMPS\PYGZus{}MEMALIGN\PYG{o}{=}\PYG{l+m}{64}\PYG{+w}{ }\PYGZhy{}I\PYG{l+s+si}{\PYGZdl{}\PYGZob{}}\PYG{n+nv}{HOME}\PYG{l+s+si}{\PYGZcb{}}/compile/lammps/lib/poems\PYG{+w}{ }\PYGZhy{}I\PYG{l+s+si}{\PYGZdl{}\PYGZob{}}\PYG{n+nv}{HOME}\PYG{l+s+si}{\PYGZcb{}}/compile/lammps/src/STUBS
\PYG{c+c1}{\PYGZsh{} Linking:}
\PYG{n+nv}{LDFLAGS}\PYG{o}{=}\PYGZhy{}g\PYG{+w}{ }\PYGZhy{}O
\PYG{c+c1}{\PYGZsh{} Libraries:}
\PYG{n+nv}{LDLIBS}\PYG{o}{=}\PYGZhy{}L\PYG{l+s+si}{\PYGZdl{}\PYGZob{}}\PYG{n+nv}{HOME}\PYG{l+s+si}{\PYGZcb{}}/compile/lammps/src\PYG{+w}{ }\PYGZhy{}llammps\PYGZus{}serial\PYG{+w}{ }\PYGZhy{}L\PYG{l+s+si}{\PYGZdl{}\PYGZob{}}\PYG{n+nv}{HOME}\PYG{l+s+si}{\PYGZcb{}}/compile/lammps/lib/poems\PYG{+w}{ }\PYGZhy{}L\PYG{l+s+si}{\PYGZdl{}\PYGZob{}}\PYG{n+nv}{HOME}\PYG{l+s+si}{\PYGZcb{}}/compile/lammps/src/STUBS\PYG{+w}{ }\PYGZhy{}lpoems\PYG{+w}{ }\PYGZhy{}lmpi\PYGZus{}stubs
\end{sphinxVerbatim}

\sphinxAtStartPar
From this you can gather the necessary paths and flags.  With
makefiles for other \sphinxstyleemphasis{machine} configurations you need to do the
equivalent and replace “serial” with the corresponding “machine”
name of the makefile.
\end{sphinxadmonition}
\end{quote}


\subsection{Link with LAMMPS as a shared library}
\label{\detokenize{Build_link:link-with-lammps-as-a-shared-library}}
\sphinxAtStartPar
When linking to LAMMPS built as a shared library, the situation becomes
much simpler, as all dependent libraries and objects are either included
in the shared library or registered as a dependent library in the shared
library file.  Thus, those libraries need not be specified when linking
the calling executable.  Only the \sphinxcode{\sphinxupquote{\sphinxhyphen{}I}} flags are needed.  So the example
case from above of the serial version static LAMMPS library with the
POEMS package installed becomes:
\begin{quote}

\begin{sphinxadmonition}{note}{CMake build}

\sphinxAtStartPar
The commands with a shared LAMMPS library compiled with the CMake
build process are the same as for the static library.

\begin{sphinxVerbatim}[commandchars=\\\{\}]
mpicc\PYG{+w}{ }\PYGZhy{}c\PYG{+w}{ }\PYGZhy{}O\PYG{+w}{ }\PYG{k}{\PYGZdl{}(}pkg\PYGZhy{}config\PYG{+w}{ }\PYGZhy{}\PYGZhy{}cflags\PYG{+w}{ }liblammps\PYG{k}{)}\PYG{+w}{ }caller.c
mpicxx\PYG{+w}{ }\PYGZhy{}o\PYG{+w}{ }\PYG{n+nb}{caller}\PYG{+w}{ }caller.o\PYG{+w}{ }\PYGZhy{}\PYG{k}{\PYGZdl{}(}pkg\PYGZhy{}config\PYG{+w}{ }\PYGZhy{}\PYGZhy{}libs\PYG{+w}{ }liblammps\PYG{k}{)}
\end{sphinxVerbatim}
\end{sphinxadmonition}

\begin{sphinxadmonition}{note}{Traditional make}

\sphinxAtStartPar
The commands with a shared LAMMPS library compiled with the
traditional make build using \sphinxcode{\sphinxupquote{make mode=shared serial}} becomes:

\begin{sphinxVerbatim}[commandchars=\\\{\}]
gcc\PYG{+w}{ }\PYGZhy{}c\PYG{+w}{ }\PYGZhy{}O\PYG{+w}{ }\PYGZhy{}I\PYG{l+s+si}{\PYGZdl{}\PYGZob{}}\PYG{n+nv}{HOME}\PYG{l+s+si}{\PYGZcb{}}/lammps/src\PYG{+w}{ }\PYGZhy{}caller.c
g++\PYG{+w}{ }\PYGZhy{}o\PYG{+w}{ }\PYG{n+nb}{caller}\PYG{+w}{ }caller.o\PYG{+w}{ }\PYGZhy{}L\PYG{l+s+si}{\PYGZdl{}\PYGZob{}}\PYG{n+nv}{HOME}\PYG{l+s+si}{\PYGZcb{}}/lammps/src\PYG{+w}{ }\PYGZhy{}llammps\PYGZus{}serial
\end{sphinxVerbatim}
\end{sphinxadmonition}
\end{quote}


\subsubsection{Locating liblammps.so at runtime}
\label{\detokenize{Build_link:locating-liblammps-so-at-runtime}}
\sphinxAtStartPar
Unlike with a static link, now the \sphinxcode{\sphinxupquote{liblammps.so}} file is required at
runtime and needs to be in a folder, where the shared linker program of
the operating system can find it.  This would be either a folder like
\sphinxcode{\sphinxupquote{/usr/local/lib64}} or \sphinxcode{\sphinxupquote{\$\{HOME\}/.local/lib64}} or a folder pointed to
by the \sphinxcode{\sphinxupquote{LD\_LIBRARY\_PATH}} environment variable. You can type

\begin{sphinxVerbatim}[commandchars=\\\{\}]
printenv\PYG{+w}{ }LD\PYGZus{}LIBRARY\PYGZus{}PATH
\end{sphinxVerbatim}

\sphinxAtStartPar
to see what directories are in that list.

\sphinxAtStartPar
Or you can add the LAMMPS src directory or the directory you performed a
CMake style build in to your \sphinxcode{\sphinxupquote{LD\_LIBRARY\_PATH}} environment variable,
so that the current version of the shared library is always available to
programs that use it.

\sphinxAtStartPar
For the Bourne or Korn shells (/bin/sh, /bin/ksh, /bin/bash etc.), you
would add something like this to your \sphinxcode{\sphinxupquote{\$\{HOME\}/.profile}} file:

\begin{sphinxVerbatim}[commandchars=\\\{\}]
LD\PYGZus{}LIBRARY\PYGZus{}PATH\PYG{+w}{ }\PYG{l+s+si}{\PYGZdl{}\PYGZob{}}\PYG{n+nv}{LD\PYGZus{}LIBRARY\PYGZus{}PATH}\PYG{p}{\PYGZhy{}/usr/lib64}\PYG{l+s+si}{\PYGZcb{}}:\PYG{l+s+si}{\PYGZdl{}\PYGZob{}}\PYG{n+nv}{HOME}\PYG{l+s+si}{\PYGZcb{}}/lammps/src
\PYG{n+nb}{export}\PYG{+w}{ }LD\PYGZus{}LIBRARY\PYGZus{}PATH
\end{sphinxVerbatim}

\sphinxAtStartPar
For the csh or tcsh shells, you would equivalently add something like this
to your \sphinxcode{\sphinxupquote{\$\{HOME\}/.cshrc}} file:

\begin{sphinxVerbatim}[commandchars=\\\{\}]
\PYG{n+nb}{setenv }LD\PYGZus{}LIBRARY\PYGZus{}PATH \PYG{k}{\PYGZdl{}\PYGZob{}}\PYG{n+nv}{LD\PYGZus{}LIBRARY\PYGZus{}PATH}\PYG{k}{\PYGZcb{}}:\PYG{k}{\PYGZdl{}\PYGZob{}}\PYG{n+nv}{HOME}\PYG{k}{\PYGZcb{}}/lammps/src
\end{sphinxVerbatim}

\sphinxAtStartPar
You can verify whether all required shared libraries are found with the
\sphinxcode{\sphinxupquote{ldd}} tool.  Example:

\begin{sphinxVerbatim}[commandchars=\\\{\}]
\PYG{n+nv}{LD\PYGZus{}LIBRARY\PYGZus{}PATH}\PYG{o}{=}/home/user/lammps/src\PYG{+w}{ }ldd\PYG{+w}{ }\PYG{n+nb}{caller}
\PYG{+w}{     }linux\PYGZhy{}vdso.so.1\PYG{+w}{ }\PYG{o}{(}0x00007ffe729e0000\PYG{o}{)}
\PYG{+w}{     }liblammps.so\PYG{+w}{ }\PYG{o}{=}\PYGZgt{}\PYG{+w}{ }/home/user/lammps/src/liblammps.so\PYG{+w}{ }\PYG{o}{(}0x00007fc91bb9e000\PYG{o}{)}
\PYG{+w}{     }libstdc++.so.6\PYG{+w}{ }\PYG{o}{=}\PYGZgt{}\PYG{+w}{ }/lib64/libstdc++.so.6\PYG{+w}{ }\PYG{o}{(}0x00007fc91b984000\PYG{o}{)}
\PYG{+w}{     }libm.so.6\PYG{+w}{ }\PYG{o}{=}\PYGZgt{}\PYG{+w}{ }/lib64/libm.so.6\PYG{+w}{ }\PYG{o}{(}0x00007fc91b83e000\PYG{o}{)}
\PYG{+w}{     }libgcc\PYGZus{}s.so.1\PYG{+w}{ }\PYG{o}{=}\PYGZgt{}\PYG{+w}{ }/lib64/libgcc\PYGZus{}s.so.1\PYG{+w}{ }\PYG{o}{(}0x00007fc91b824000\PYG{o}{)}
\PYG{+w}{     }libc.so.6\PYG{+w}{ }\PYG{o}{=}\PYGZgt{}\PYG{+w}{ }/lib64/libc.so.6\PYG{+w}{ }\PYG{o}{(}0x00007fc91b65b000\PYG{o}{)}
\PYG{+w}{     }/lib64/ld\PYGZhy{}linux\PYGZhy{}x86\PYGZhy{}64.so.2\PYG{+w}{ }\PYG{o}{(}0x00007fc91c094000\PYG{o}{)}
\end{sphinxVerbatim}

\sphinxAtStartPar
If a required library is missing, you would get a ‘not found’ entry:

\begin{sphinxVerbatim}[commandchars=\\\{\}]
ldd\PYG{+w}{ }\PYG{n+nb}{caller}
\PYG{+w}{     }linux\PYGZhy{}vdso.so.1\PYG{+w}{ }\PYG{o}{(}0x00007ffd672fe000\PYG{o}{)}
\PYG{+w}{     }liblammps.so\PYG{+w}{ }\PYG{o}{=}\PYGZgt{}\PYG{+w}{ }not\PYG{+w}{ }found
\PYG{+w}{     }libstdc++.so.6\PYG{+w}{ }\PYG{o}{=}\PYGZgt{}\PYG{+w}{ }/usr/lib64/libstdc++.so.6\PYG{+w}{ }\PYG{o}{(}0x00007fb7c7e86000\PYG{o}{)}
\PYG{+w}{     }libm.so.6\PYG{+w}{ }\PYG{o}{=}\PYGZgt{}\PYG{+w}{ }/usr/lib64/libm.so.6\PYG{+w}{ }\PYG{o}{(}0x00007fb7c7d40000\PYG{o}{)}
\PYG{+w}{     }libgcc\PYGZus{}s.so.1\PYG{+w}{ }\PYG{o}{=}\PYGZgt{}\PYG{+w}{ }/usr/lib64/libgcc\PYGZus{}s.so.1\PYG{+w}{ }\PYG{o}{(}0x00007fb7c7d26000\PYG{o}{)}
\PYG{+w}{     }libc.so.6\PYG{+w}{ }\PYG{o}{=}\PYGZgt{}\PYG{+w}{ }/usr/lib64/libc.so.6\PYG{+w}{ }\PYG{o}{(}0x00007fb7c7b5d000\PYG{o}{)}
\PYG{+w}{     }/lib64/ld\PYGZhy{}linux\PYGZhy{}x86\PYGZhy{}64.so.2\PYG{+w}{ }\PYG{o}{(}0x00007fb7c80a2000\PYG{o}{)}
\end{sphinxVerbatim}

\sphinxstepscope


\section{Basic build options}
\label{\detokenize{Build_basics:basic-build-options}}\label{\detokenize{Build_basics::doc}}
\sphinxAtStartPar
The following topics are covered on this page, for building with both
CMake and make:
\begin{itemize}
\item {} 
\sphinxAtStartPar
{\hyperref[\detokenize{Build_basics:serial}]{\sphinxcrossref{\DUrole{std,std-ref}{Serial vs parallel build}}}}

\item {} 
\sphinxAtStartPar
{\hyperref[\detokenize{Build_basics:compile}]{\sphinxcrossref{\DUrole{std,std-ref}{Choice of compiler and compile/link options}}}}

\item {} 
\sphinxAtStartPar
{\hyperref[\detokenize{Build_basics:exe}]{\sphinxcrossref{\DUrole{std,std-ref}{Build the LAMMPS executable and library}}}}

\item {} 
\sphinxAtStartPar
{\hyperref[\detokenize{Build_basics:debug}]{\sphinxcrossref{\DUrole{std,std-ref}{Including and removing debug support}}}}

\item {} 
\sphinxAtStartPar
{\hyperref[\detokenize{Build_basics:install}]{\sphinxcrossref{\DUrole{std,std-ref}{Install LAMMPS after a build}}}}

\end{itemize}


\bigskip\hrule\bigskip



\subsection{Serial vs parallel build}
\label{\detokenize{Build_basics:serial-vs-parallel-build}}\label{\detokenize{Build_basics:serial}}
\sphinxAtStartPar
LAMMPS is written to use the ubiquitous \sphinxhref{https://en.wikipedia.org/wiki/Message\_Passing\_Interface}{MPI (Message Passing Interface)} library API
for distributed memory parallel computation.  You need to have such a
library installed for building and running LAMMPS in parallel using a
domain decomposition parallelization.  It is compatible with the MPI
standard version 2.x and later.  LAMMPS can also be built into a
“serial” executable for use with a single processor using the bundled
MPI STUBS library.

\sphinxAtStartPar
Independent of the distributed memory MPI parallelization, parts of
LAMMPS are also written with support for shared memory parallelization
using the \sphinxhref{https://en.wikipedia.org/wiki/OpenMP}{OpenMP} threading
standard. A more detailed discussion of that is below.
\begin{quote}

\begin{sphinxadmonition}{note}{CMake build}

\begin{sphinxVerbatim}[commandchars=\\\{\}]
\PYGZhy{}D\PYG{+w}{ }\PYG{n+nv}{BUILD\PYGZus{}MPI}\PYG{o}{=}value\PYG{+w}{        }\PYG{c+c1}{\PYGZsh{} yes or no, default is yes if CMake finds MPI}
\PYGZhy{}D\PYG{+w}{ }\PYG{n+nv}{BUILD\PYGZus{}OMP}\PYG{o}{=}value\PYG{+w}{        }\PYG{c+c1}{\PYGZsh{} yes or no, default is yes if a compatible}
\PYG{+w}{                          }\PYG{c+c1}{\PYGZsh{} compiler is detected}
\PYGZhy{}D\PYG{+w}{ }\PYG{n+nv}{LAMMPS\PYGZus{}MACHINE}\PYG{o}{=}name\PYG{+w}{    }\PYG{c+c1}{\PYGZsh{} name = mpi, serial, mybox, titan, laptop, etc}
\PYG{+w}{                          }\PYG{c+c1}{\PYGZsh{} no default value}
\end{sphinxVerbatim}

\sphinxAtStartPar
The executable created by CMake (after running make) is named
\sphinxcode{\sphinxupquote{lmp}} unless the \sphinxcode{\sphinxupquote{LAMMPS\_MACHINE}} option is set.  When setting
\sphinxcode{\sphinxupquote{LAMMPS\_MACHINE=name}}, the executable will be called
\sphinxcode{\sphinxupquote{lmp\_name}}.  Using \sphinxcode{\sphinxupquote{BUILD\_MPI=no}} will enforce building a
serial executable using the MPI STUBS library.
\end{sphinxadmonition}

\begin{sphinxadmonition}{note}{Traditional make}

\sphinxAtStartPar
The build with traditional makefiles has to be done inside the source folder \sphinxcode{\sphinxupquote{src}}.

\begin{sphinxVerbatim}[commandchars=\\\{\}]
make\PYG{+w}{ }mpi\PYG{+w}{      }\PYG{c+c1}{\PYGZsh{} parallel build, produces lmp\PYGZus{}mpi using Makefile.mpi}
make\PYG{+w}{ }serial\PYG{+w}{   }\PYG{c+c1}{\PYGZsh{} serial build, produces lmp\PYGZus{}serial using Makefile/serial}
make\PYG{+w}{ }mybox\PYG{+w}{    }\PYG{c+c1}{\PYGZsh{} uses Makefile.mybox to produce lmp\PYGZus{}mybox}
\end{sphinxVerbatim}

\sphinxAtStartPar
Any \sphinxcode{\sphinxupquote{make machine}} command will look up the make settings from a
file \sphinxcode{\sphinxupquote{Makefile.machine}} in the folder \sphinxcode{\sphinxupquote{src/MAKE}} or one of its
subdirectories \sphinxcode{\sphinxupquote{MINE}}, \sphinxcode{\sphinxupquote{MACHINES}}, or \sphinxcode{\sphinxupquote{OPTIONS}}, create a
folder \sphinxcode{\sphinxupquote{Obj\_machine}} with all objects and generated files and an
executable called \sphinxcode{\sphinxupquote{lmp\_machine}}.  The standard parallel build
with \sphinxcode{\sphinxupquote{make mpi}} assumes a standard MPI installation with MPI
compiler wrappers where all necessary compiler and linker flags to
get access and link with the suitable MPI headers and libraries
are set by the wrapper programs.  For other cases or the serial
build, you have to adjust the make file variables \sphinxcode{\sphinxupquote{MPI\_INC}},
\sphinxcode{\sphinxupquote{MPI\_PATH}}, \sphinxcode{\sphinxupquote{MPI\_LIB}} as well as \sphinxcode{\sphinxupquote{CC}} and \sphinxcode{\sphinxupquote{LINK}}.  To
enable OpenMP threading usually a compiler specific flag needs to
be added to the compile and link commands.  For the GNU compilers,
this is \sphinxcode{\sphinxupquote{\sphinxhyphen{}fopenmp}}, which can be added to the \sphinxcode{\sphinxupquote{CC}} and
\sphinxcode{\sphinxupquote{LINK}} makefile variables.

\sphinxAtStartPar
For the serial build the following make variables are set (see \sphinxcode{\sphinxupquote{src/MAKE/Makefile.serial}}):

\begin{sphinxVerbatim}[commandchars=\\\{\}]
\PYG{n+nv}{CC}\PYG{+w}{ }\PYG{o}{=}\PYG{+w}{       }g++
\PYG{n+nv}{LINK}\PYG{+w}{ }\PYG{o}{=}\PYG{+w}{     }g++
\PYG{n+nv}{MPI\PYGZus{}INC}\PYG{+w}{ }\PYG{o}{=}\PYG{+w}{  }\PYGZhy{}I../STUBS
\PYG{n+nv}{MPI\PYGZus{}PATH}\PYG{+w}{ }\PYG{o}{=}\PYG{+w}{ }\PYGZhy{}L../STUBS
\PYG{n+nv}{MPI\PYGZus{}LIB}\PYG{+w}{ }\PYG{o}{=}\PYG{+w}{  }\PYGZhy{}lmpi\PYGZus{}stubs
\end{sphinxVerbatim}

\sphinxAtStartPar
You also need to build the STUBS library for your platform before
making LAMMPS itself.  A \sphinxcode{\sphinxupquote{make serial}} build does this for you
automatically, otherwise, type \sphinxcode{\sphinxupquote{make mpi\sphinxhyphen{}stubs}} from the src
directory, or \sphinxcode{\sphinxupquote{make}} from the \sphinxcode{\sphinxupquote{src/STUBS}} dir.  If the build
fails, you may need to edit the \sphinxcode{\sphinxupquote{STUBS/Makefile}} for your
platform.  The stubs library does not provide MPI/IO functions
required by some LAMMPS packages, e.g. \sphinxcode{\sphinxupquote{LATBOLTZ}},
and thus is not compatible with those packages.

\begin{sphinxadmonition}{note}{Note:}
\sphinxAtStartPar
The file \sphinxcode{\sphinxupquote{src/STUBS/mpi.cpp}} provides a CPU timer function
called \sphinxcode{\sphinxupquote{MPI\_Wtime()}} that calls \sphinxcode{\sphinxupquote{gettimeofday()}}.  If your
operating system does not support \sphinxcode{\sphinxupquote{gettimeofday()}}, you will
need to insert code to call another timer.  Note that the
ANSI\sphinxhyphen{}standard function \sphinxcode{\sphinxupquote{clock()}} rolls over after an hour or
so, and is therefore insufficient for timing long LAMMPS
simulations.
\end{sphinxadmonition}
\end{sphinxadmonition}
\end{quote}


\subsubsection{MPI and OpenMP support in LAMMPS}
\label{\detokenize{Build_basics:mpi-and-openmp-support-in-lammps}}
\sphinxAtStartPar
If you are installing MPI yourself to build a parallel LAMMPS
executable, we recommend either MPICH or OpenMPI, which are regularly
used and tested with LAMMPS by the LAMMPS developers.  MPICH can be
downloaded from the \sphinxhref{https://www.mpich.org}{MPICH home page}, and
OpenMPI can be downloaded correspondingly from the \sphinxhref{https://www.open-mpi.org}{OpenMPI home page}.  Other MPI packages should also work.  No
specific vendor provided and standard compliant MPI library is currently
known to be incompatible with LAMMPS.  If you are running on a large
parallel machine, your system admins or the vendor should have already
installed a version of MPI, which is likely to be faster than a
self\sphinxhyphen{}installed MPICH or OpenMPI, so you should study the provided
documentation to find out how to build and link with it.

\sphinxAtStartPar
The majority of OpenMP (threading) support in LAMMPS is provided by the
\sphinxcode{\sphinxupquote{OPENMP}} package; see the {\hyperref[\detokenize{Speed_omp::doc}]{\sphinxcrossref{\DUrole{doc}{OPENMP package}}}}
page for details. The \sphinxcode{\sphinxupquote{INTEL}} package also includes OpenMP
threading (it is compatible with \sphinxcode{\sphinxupquote{OPENMP}} and will usually fall
back on styles from that package, if a \sphinxcode{\sphinxupquote{INTEL}} does not exist)
and adds vectorization support when compiled with compatible compilers,
in particular the Intel compilers on top of OpenMP. Also, the \sphinxcode{\sphinxupquote{KOKKOS}}
package can be compiled to include OpenMP threading.

\sphinxAtStartPar
In addition, there are a few commands in LAMMPS that have native OpenMP
support included as well.  These are commands in the \sphinxcode{\sphinxupquote{ML\sphinxhyphen{}SNAP}},
\sphinxcode{\sphinxupquote{DIFFRACTION}}, and \sphinxcode{\sphinxupquote{DPD\sphinxhyphen{}REACT}} packages.  Furthermore, some packages
support OpenMP threading indirectly through the libraries they interface
to: e.g. \sphinxcode{\sphinxupquote{KSPACE}}, and \sphinxcode{\sphinxupquote{COLVARS}}.  See the {\hyperref[\detokenize{Packages_details::doc}]{\sphinxcrossref{\DUrole{doc}{Packages details}}}} page for more info on these packages, and the pages
for their respective commands for OpenMP threading info.

\sphinxAtStartPar
For CMake, if you use \sphinxcode{\sphinxupquote{BUILD\_OMP=yes}}, you can use these packages
and turn on their native OpenMP support and turn on their native OpenMP
support at run time, by setting the \sphinxcode{\sphinxupquote{OMP\_NUM\_THREADS}} environment
variable before you launch LAMMPS.

\sphinxAtStartPar
For building via conventional make, the \sphinxcode{\sphinxupquote{CCFLAGS}} and \sphinxcode{\sphinxupquote{LINKFLAGS}}
variables in Makefile.machine need to include the compiler flag that
enables OpenMP. For the GNU compilers or Clang, it is \sphinxcode{\sphinxupquote{\sphinxhyphen{}fopenmp}}.
For (recent) Intel compilers, it is \sphinxcode{\sphinxupquote{\sphinxhyphen{}qopenmp}}.  If you are using a
different compiler, please refer to its documentation.


\subsubsection{OpenMP Compiler compatibility}
\label{\detokenize{Build_basics:openmp-compiler-compatibility}}\label{\detokenize{Build_basics:default-none-issues}}
\sphinxAtStartPar
Some compilers do not fully support the \sphinxcode{\sphinxupquote{default(none)}} directive and
others (e.g. GCC version 9 and beyond, Clang version 10 and later) may
implement strict OpenMP 4.0 and later semantics, which are incompatible
with the OpenMP 3.1 semantics used in LAMMPS for maximal compatibility
with compiler versions in use.  If compilation with OpenMP enabled fails
because of your compiler requiring strict OpenMP 4.0 semantics, you can
change the behavior by adding \sphinxcode{\sphinxupquote{\sphinxhyphen{}D LAMMPS\_OMP\_COMPAT=4}} to the
\sphinxcode{\sphinxupquote{LMP\_INC}} variable in your makefile, or add it to the command line
while configuring with CMake.  LAMMPS will auto\sphinxhyphen{}detect a suitable setting
for most GNU, Clang, and Intel compilers.


\bigskip\hrule\bigskip



\subsection{Choice of compiler and compile/link options}
\label{\detokenize{Build_basics:choice-of-compiler-and-compile-link-options}}\label{\detokenize{Build_basics:compile}}
\sphinxAtStartPar
The choice of compiler and compiler flags can be important for maximum
performance.  Vendor provided compilers for a specific hardware can
produce faster code than open\sphinxhyphen{}source compilers like the GNU compilers.
On the most common x86 hardware, the most popular C++ compilers are
quite similar in their ability to optimize regular C/C++ source code at
high optimization levels.  When using the \sphinxcode{\sphinxupquote{INTEL}} package, there is a
distinct advantage in using the \sphinxhref{https://software.intel.com/en-us/intel-compilers}{Intel C++ compiler} due to
much improved vectorization through SSE and AVX instructions on
compatible hardware.  The source code in that package conditionally
includes compiler specific directives to enable these high degrees of
vectorization.  This may change over time as equivalent vectorization
directives are included into the OpenMP standard and other compilers
adopt them.

\sphinxAtStartPar
On parallel clusters or supercomputers which use “environment modules”
for their compile/link environments, you can often access different
compilers by simply loading the appropriate module before building
LAMMPS.
\begin{quote}

\begin{sphinxadmonition}{note}{CMake build}

\sphinxAtStartPar
By default CMake will use the compiler it finds according to
internal preferences, and it will add optimization flags
appropriate to that compiler and any {\hyperref[\detokenize{Speed_packages::doc}]{\sphinxcrossref{\DUrole{doc}{accelerator packages}}}} you have included in the build.  CMake will
check if the detected or selected compiler is compatible with the
C++ support requirements of LAMMPS and stop with an error, if this
is not the case.

\sphinxAtStartPar
You can tell CMake to look for a specific compiler with setting
CMake variables (listed below) during configuration.  For a few
common choices, there are also presets in the \sphinxcode{\sphinxupquote{cmake/presets}}
folder.  For convenience, there is a \sphinxcode{\sphinxupquote{CMAKE\_TUNE\_FLAGS}} variable
that can be set to apply global compiler options (applied to
compilation only), to be used for adding compiler or host specific
optimization flags in addition to the “flags” variables listed
below. You may also specify the corresponding \sphinxcode{\sphinxupquote{CMAKE\_*\_FLAGS}}
variables individually, if you want to experiment with alternate
optimization flags.  You should specify all 3 compilers, so that
the (few) LAMMPS source files written in C or Fortran are built
with a compiler consistent with the one used for the C++ files:

\begin{sphinxVerbatim}[commandchars=\\\{\}]
\PYGZhy{}D\PYG{+w}{ }\PYG{n+nv}{CMAKE\PYGZus{}CXX\PYGZus{}COMPILER}\PYG{o}{=}name\PYG{+w}{            }\PYG{c+c1}{\PYGZsh{} name of C++ compiler}
\PYGZhy{}D\PYG{+w}{ }\PYG{n+nv}{CMAKE\PYGZus{}C\PYGZus{}COMPILER}\PYG{o}{=}name\PYG{+w}{              }\PYG{c+c1}{\PYGZsh{} name of C compiler}
\PYGZhy{}D\PYG{+w}{ }\PYG{n+nv}{CMAKE\PYGZus{}Fortran\PYGZus{}COMPILER}\PYG{o}{=}name\PYG{+w}{        }\PYG{c+c1}{\PYGZsh{} name of Fortran compiler}

\PYGZhy{}D\PYG{+w}{ }\PYG{n+nv}{CMAKE\PYGZus{}CXX\PYGZus{}FLAGS}\PYG{o}{=}string\PYG{+w}{             }\PYG{c+c1}{\PYGZsh{} flags to use with C++ compiler}
\PYGZhy{}D\PYG{+w}{ }\PYG{n+nv}{CMAKE\PYGZus{}C\PYGZus{}FLAGS}\PYG{o}{=}string\PYG{+w}{               }\PYG{c+c1}{\PYGZsh{} flags to use with C compiler}
\PYGZhy{}D\PYG{+w}{ }\PYG{n+nv}{CMAKE\PYGZus{}Fortran\PYGZus{}FLAGS}\PYG{o}{=}string\PYG{+w}{         }\PYG{c+c1}{\PYGZsh{} flags to use with Fortran compiler}
\end{sphinxVerbatim}

\sphinxAtStartPar
A few example command lines are:

\begin{sphinxVerbatim}[commandchars=\\\{\}]
\PYG{c+c1}{\PYGZsh{} Building with GNU Compilers:}
cmake\PYG{+w}{ }\PYGZhy{}DCMAKE\PYGZus{}C\PYGZus{}COMPILER\PYG{o}{=}gcc\PYG{+w}{ }\PYGZhy{}DCMAKE\PYGZus{}CXX\PYGZus{}COMPILER\PYG{o}{=}g++\PYG{+w}{ }\PYG{l+s+se}{\PYGZbs{}}
\PYG{+w}{      }\PYGZhy{}DCMAKE\PYGZus{}Fortran\PYGZus{}COMPILER\PYG{o}{=}gfortran\PYG{+w}{ }../cmake
\PYG{c+c1}{\PYGZsh{} Building with Intel Classic Compilers:}
cmake\PYG{+w}{ }\PYGZhy{}DCMAKE\PYGZus{}C\PYGZus{}COMPILER\PYG{o}{=}icc\PYG{+w}{ }\PYGZhy{}DCMAKE\PYGZus{}CXX\PYGZus{}COMPILER\PYG{o}{=}icpc\PYG{+w}{ }\PYG{l+s+se}{\PYGZbs{}}
\PYG{+w}{      }\PYGZhy{}DCMAKE\PYGZus{}Fortran\PYGZus{}COMPILER\PYG{o}{=}ifort\PYG{+w}{ }../cmake
\PYG{c+c1}{\PYGZsh{} Building with Intel oneAPI Compilers:}
cmake\PYG{+w}{ }\PYGZhy{}DCMAKE\PYGZus{}C\PYGZus{}COMPILER\PYG{o}{=}icx\PYG{+w}{ }\PYGZhy{}DCMAKE\PYGZus{}CXX\PYGZus{}COMPILER\PYG{o}{=}icpx\PYG{+w}{ }\PYG{l+s+se}{\PYGZbs{}}
\PYG{+w}{      }\PYGZhy{}DCMAKE\PYGZus{}Fortran\PYGZus{}COMPILER\PYG{o}{=}ifx\PYG{+w}{ }../cmake
\PYG{c+c1}{\PYGZsh{} Building with LLVM/Clang Compilers:}
cmake\PYG{+w}{ }\PYGZhy{}DCMAKE\PYGZus{}C\PYGZus{}COMPILER\PYG{o}{=}clang\PYG{+w}{ }\PYGZhy{}DCMAKE\PYGZus{}CXX\PYGZus{}COMPILER\PYG{o}{=}clang++\PYG{+w}{ }\PYG{l+s+se}{\PYGZbs{}}
\PYG{+w}{      }\PYGZhy{}DCMAKE\PYGZus{}Fortran\PYGZus{}COMPILER\PYG{o}{=}flang\PYG{+w}{ }../cmake
\PYG{c+c1}{\PYGZsh{} Building with PGI/Nvidia Compilers:}
cmake\PYG{+w}{ }\PYGZhy{}DCMAKE\PYGZus{}C\PYGZus{}COMPILER\PYG{o}{=}pgcc\PYG{+w}{ }\PYGZhy{}DCMAKE\PYGZus{}CXX\PYGZus{}COMPILER\PYG{o}{=}pgc++\PYG{+w}{ }\PYG{l+s+se}{\PYGZbs{}}
\PYG{+w}{      }\PYGZhy{}DCMAKE\PYGZus{}Fortran\PYGZus{}COMPILER\PYG{o}{=}pgfortran\PYG{+w}{ }../cmake
\PYG{c+c1}{\PYGZsh{} Building with the NVHPC Compilers:}
cmake\PYG{+w}{ }\PYGZhy{}DCMAKE\PYGZus{}C\PYGZus{}COMPILER\PYG{o}{=}nvc\PYG{+w}{ }\PYGZhy{}DCMAKE\PYGZus{}CXX\PYGZus{}COMPILER\PYG{o}{=}nvc++\PYG{+w}{ }\PYG{l+s+se}{\PYGZbs{}}
\PYG{+w}{      }\PYGZhy{}DCMAKE\PYGZus{}Fortran\PYGZus{}COMPILER\PYG{o}{=}nvfortran\PYG{+w}{ }../cmake
\end{sphinxVerbatim}

\sphinxAtStartPar
For compiling with the Clang/LLVM compilers a CMake preset is
provided that can be loaded with
\sphinxcode{\sphinxupquote{\sphinxhyphen{}C ../cmake/presets/clang.cmake}}.  Similarly,
\sphinxcode{\sphinxupquote{\sphinxhyphen{}C ../cmake/presets/intel.cmake}} should switch the compiler
toolchain to the legacy Intel compilers, \sphinxcode{\sphinxupquote{\sphinxhyphen{}C ../cmake/presets/oneapi.cmake}}
will switch to the LLVM based oneAPI Intel compilers,
\sphinxcode{\sphinxupquote{\sphinxhyphen{}C ../cmake/presets/pgi.cmake}} will switch the compiler to the PGI compilers,
and \sphinxcode{\sphinxupquote{\sphinxhyphen{}C ../cmake/presets/nvhpc.cmake}} will switch to the NVHPC compilers.

\sphinxAtStartPar
Furthermore, you can set \sphinxcode{\sphinxupquote{CMAKE\_TUNE\_FLAGS}} to specifically add
compiler flags to tune for optimal performance on given hosts.
This variable is empty by default.

\begin{sphinxadmonition}{note}{Note:}
\sphinxAtStartPar
When the cmake command completes, it prints a summary to the
screen which compilers it is using and what flags and settings
will be used for the compilation.  Note that if the top\sphinxhyphen{}level
compiler is \sphinxcode{\sphinxupquote{mpicxx}}, it is simply a wrapper on a real compiler.
The underlying compiler info is what CMake will try to
determine and report.  You should check to confirm you are
using the compiler and optimization flags you want.
\end{sphinxadmonition}
\end{sphinxadmonition}

\begin{sphinxadmonition}{note}{Makefile.machine settings for traditional make}

\sphinxAtStartPar
The “compiler/linker settings” section of a Makefile.machine lists
compiler and linker settings for your C++ compiler, including
optimization flags.  For a parallel build it is recommended to use
\sphinxcode{\sphinxupquote{mpicxx}} or \sphinxcode{\sphinxupquote{mpiCC}}, since these compiler wrappers will
include a variety of settings appropriate for your MPI
installation and thus avoiding the guesswork of finding the right
flags.

\sphinxAtStartPar
Parallel build (see \sphinxcode{\sphinxupquote{src/MAKE/Makefile.mpi}}):

\begin{sphinxVerbatim}[commandchars=\\\{\}]
\PYG{n+nv}{CC}\PYG{+w}{ }\PYG{o}{=}\PYG{+w}{        }mpicxx
\PYG{n+nv}{CCFLAGS}\PYG{+w}{ }\PYG{o}{=}\PYG{+w}{   }\PYGZhy{}g\PYG{+w}{ }\PYGZhy{}O3
\PYG{n+nv}{LINK}\PYG{+w}{ }\PYG{o}{=}\PYG{+w}{      }mpicxx
\PYG{n+nv}{LINKFLAGS}\PYG{+w}{ }\PYG{o}{=}\PYG{+w}{ }\PYGZhy{}g\PYG{+w}{ }\PYGZhy{}O
\end{sphinxVerbatim}

\sphinxAtStartPar
Serial build with GNU gcc (see \sphinxcode{\sphinxupquote{src/MAKE/Makefile.serial}}):

\begin{sphinxVerbatim}[commandchars=\\\{\}]
\PYG{n+nv}{CC}\PYG{+w}{ }\PYG{o}{=}\PYG{+w}{        }g++
\PYG{n+nv}{CCFLAGS}\PYG{+w}{ }\PYG{o}{=}\PYG{+w}{   }\PYGZhy{}g\PYG{+w}{ }\PYGZhy{}O3
\PYG{n+nv}{LINK}\PYG{+w}{ }\PYG{o}{=}\PYG{+w}{      }g++
\PYG{n+nv}{LINKFLAGS}\PYG{+w}{ }\PYG{o}{=}\PYG{+w}{ }\PYGZhy{}g\PYG{+w}{ }\PYGZhy{}O
\end{sphinxVerbatim}

\begin{sphinxadmonition}{note}{Note:}
\sphinxAtStartPar
If compilation stops with a message like the following:

\begin{sphinxVerbatim}[commandchars=\\\{\}]
\PYG{g+go}{g++ \PYGZhy{}g \PYGZhy{}O3  \PYGZhy{}DLAMMPS\PYGZus{}GZIP \PYGZhy{}DLAMMPS\PYGZus{}MEMALIGN=64    \PYGZhy{}I../STUBS     \PYGZhy{}c ../main.cpp}
\PYG{g+go}{In file included from ../pointers.h:24:0,}
\PYG{g+go}{           from ../input.h:17,}
\PYG{g+go}{           from ../main.cpp:16:}
\PYG{g+go}{../lmptype.h:34:2: error: \PYGZsh{}error LAMMPS requires a C++11 (or later) compliant compiler. Enable C++11 compatibility or upgrade the compiler.}
\end{sphinxVerbatim}

\sphinxAtStartPar
then you have either an unsupported (old) compiler or you have
to turn on C++11 mode.  The latter applies to GCC 4.8.x shipped
with RHEL 7.x and CentOS 7.x or GCC 5.4.x shipped with Ubuntu16.04.
For those compilers, you need to add the \sphinxcode{\sphinxupquote{\sphinxhyphen{}std=c++11}} flag.
If there is no compiler that supports this flag (or equivalent),
you would have to install a newer compiler that supports C++11;
either as a binary package or through compiling from source.

\sphinxAtStartPar
If you build LAMMPS with any {\hyperref[\detokenize{Speed_packages::doc}]{\sphinxcrossref{\DUrole{doc}{Accelerator packages}}}} included,
there may be specific compiler or linker flags that are either
required or recommended to enable required features and to
achieve optimal performance.  You need to include these in the
\sphinxcode{\sphinxupquote{CCFLAGS}} and \sphinxcode{\sphinxupquote{LINKFLAGS}} settings above.  For details, see the
documentation for the individual packages listed on the
{\hyperref[\detokenize{Speed_packages::doc}]{\sphinxcrossref{\DUrole{doc}{Accelerator packages}}}} page.  Or examine these files in the
\sphinxcode{\sphinxupquote{src/MAKE/OPTIONS}} directory.  They correspond to each of the 5
accelerator packages and their hardware variants:

\begin{sphinxVerbatim}[commandchars=\\\{\}]
Makefile.opt\PYG{+w}{                   }\PYG{c+c1}{\PYGZsh{} OPT package}
Makefile.omp\PYG{+w}{                   }\PYG{c+c1}{\PYGZsh{} OPENMP package}
Makefile.intel\PYGZus{}cpu\PYG{+w}{             }\PYG{c+c1}{\PYGZsh{} INTEL package for CPUs}
Makefile.intel\PYGZus{}coprocessor\PYG{+w}{     }\PYG{c+c1}{\PYGZsh{} INTEL package for KNLs}
Makefile.gpu\PYG{+w}{                   }\PYG{c+c1}{\PYGZsh{} GPU package}
Makefile.kokkos\PYGZus{}cuda\PYGZus{}mpi\PYG{+w}{       }\PYG{c+c1}{\PYGZsh{} KOKKOS package for GPUs}
Makefile.kokkos\PYGZus{}omp\PYG{+w}{            }\PYG{c+c1}{\PYGZsh{} KOKKOS package for CPUs (OpenMP)}
Makefile.kokkos\PYGZus{}phi\PYG{+w}{            }\PYG{c+c1}{\PYGZsh{} KOKKOS package for KNLs (OpenMP)}
\end{sphinxVerbatim}
\end{sphinxadmonition}
\end{sphinxadmonition}
\end{quote}


\bigskip\hrule\bigskip



\subsection{Build the LAMMPS executable and library}
\label{\detokenize{Build_basics:build-the-lammps-executable-and-library}}\label{\detokenize{Build_basics:library}}\label{\detokenize{Build_basics:exe}}
\sphinxAtStartPar
LAMMPS is always built as a library of C++ classes plus an executable.
The executable is a simple \sphinxcode{\sphinxupquote{main()}} function that sets up MPI and then
creates a LAMMPS class instance from the LAMMPS library, which
will then process commands provided via a file or from the console
input.  The LAMMPS library can also be called from another application
or a scripting language.  See the {\hyperref[\detokenize{Howto_couple::doc}]{\sphinxcrossref{\DUrole{doc}{Howto couple}}}} doc
page for more info on coupling LAMMPS to other codes.  See the
{\hyperref[\detokenize{Python_head::doc}]{\sphinxcrossref{\DUrole{doc}{Python}}}} page for more info on wrapping and
running LAMMPS from Python via its library interface.
\begin{quote}

\begin{sphinxadmonition}{note}{CMake build}

\sphinxAtStartPar
For CMake builds, you can select through setting CMake variables
between building a shared or a static LAMMPS library and what kind
of suffix is added to them (in case you want to concurrently
install multiple variants of binaries with different settings). If
none are set, defaults are applied.

\begin{sphinxVerbatim}[commandchars=\\\{\}]
\PYGZhy{}D\PYG{+w}{ }\PYG{n+nv}{BUILD\PYGZus{}SHARED\PYGZus{}LIBS}\PYG{o}{=}value\PYG{+w}{   }\PYG{c+c1}{\PYGZsh{} yes or no (default)}
\PYGZhy{}D\PYG{+w}{ }\PYG{n+nv}{LAMMPS\PYGZus{}MACHINE}\PYG{o}{=}name\PYG{+w}{       }\PYG{c+c1}{\PYGZsh{} name = mpi, serial, mybox, titan, laptop, etc}
\PYG{+w}{                             }\PYG{c+c1}{\PYGZsh{} no default value}
\end{sphinxVerbatim}

\sphinxAtStartPar
The compilation will always produce a LAMMPS library and an
executable linked to it.  By default, this will be a static
library named \sphinxcode{\sphinxupquote{liblammps.a}} and an executable named \sphinxcode{\sphinxupquote{lmp}}
Setting \sphinxcode{\sphinxupquote{BUILD\_SHARED\_LIBS=yes}} will instead produce a shared
library called \sphinxcode{\sphinxupquote{liblammps.so}} (or \sphinxcode{\sphinxupquote{liblammps.dylib}} or
\sphinxcode{\sphinxupquote{liblammps.dll}} depending on the platform) If
\sphinxcode{\sphinxupquote{LAMMPS\_MACHINE=name}} is set in addition, the name of the
generated libraries will be changed to either \sphinxcode{\sphinxupquote{liblammps\_name.a}}
or \sphinxcode{\sphinxupquote{liblammps\_name.so}}, respectively and the executable will
be called \sphinxcode{\sphinxupquote{lmp\_name}}.
\end{sphinxadmonition}

\begin{sphinxadmonition}{note}{Traditional make}

\sphinxAtStartPar
With the traditional makefile based build process, the choice of
the generated executable or library depends on the “mode” setting.
Several options are available and \sphinxcode{\sphinxupquote{mode=static}} is the default.

\begin{sphinxVerbatim}[commandchars=\\\{\}]
make\PYG{+w}{ }machine\PYG{+w}{               }\PYG{c+c1}{\PYGZsh{} build LAMMPS executable lmp\PYGZus{}machine}
make\PYG{+w}{ }\PYG{n+nv}{mode}\PYG{o}{=}static\PYG{+w}{ }machine\PYG{+w}{   }\PYG{c+c1}{\PYGZsh{} same as \PYGZdq{}make machine\PYGZdq{}}
make\PYG{+w}{ }\PYG{n+nv}{mode}\PYG{o}{=}shared\PYG{+w}{ }machine\PYG{+w}{   }\PYG{c+c1}{\PYGZsh{} build LAMMPS shared lib liblammps\PYGZus{}machine.so}
\PYG{+w}{                           }\PYG{c+c1}{\PYGZsh{} instead}
\end{sphinxVerbatim}

\sphinxAtStartPar
The “static” build will generate a static library called
\sphinxcode{\sphinxupquote{liblammps\_machine.a}} and an executable named \sphinxcode{\sphinxupquote{lmp\_machine}},
while the “shared” build will generate a shared library
\sphinxcode{\sphinxupquote{liblammps\_machine.so}} instead and \sphinxcode{\sphinxupquote{lmp\_machine}} will be
linked to it.  The build step will also create generic soft links,
named \sphinxcode{\sphinxupquote{liblammps.a}} and \sphinxcode{\sphinxupquote{liblammps.so}}, which point to the
specific \sphinxcode{\sphinxupquote{liblammps\_machine.a/so}} files.
\end{sphinxadmonition}
\end{quote}


\subsubsection{Additional information}
\label{\detokenize{Build_basics:additional-information}}
\sphinxAtStartPar
Note that for creating a shared library, all the libraries it depends on
must be compiled to be compatible with shared libraries.  This should be
the case for libraries included with LAMMPS, such as the dummy MPI
library in \sphinxcode{\sphinxupquote{src/STUBS}} or any package libraries in the \sphinxcode{\sphinxupquote{lib}}
directory, since they are always built in a shared library compatible
way using the \sphinxcode{\sphinxupquote{\sphinxhyphen{}fPIC}} compiler switch.  However, if an auxiliary
library (like MPI or FFTW) does not exist as a compatible format, the
shared library linking step may generate an error.  This means you will
need to install a compatible version of the auxiliary library.  The
build instructions for that library should tell you how to do this.

\sphinxAtStartPar
As an example, here is how to build and install the \sphinxhref{https://www.mpich.org}{MPICH library}, a popular open\sphinxhyphen{}source version of MPI, as a shared library
in the default /usr/local/lib location:

\begin{sphinxVerbatim}[commandchars=\\\{\}]
./configure\PYG{+w}{ }\PYGZhy{}\PYGZhy{}enable\PYGZhy{}shared
make
make\PYG{+w}{ }install
\end{sphinxVerbatim}

\sphinxAtStartPar
You may need to use \sphinxcode{\sphinxupquote{sudo make install}} in place of the last line if
you do not have write privileges for \sphinxcode{\sphinxupquote{/usr/local/lib}} or use the
\sphinxcode{\sphinxupquote{\sphinxhyphen{}\sphinxhyphen{}prefix}} configuration option to select an installation folder,
where you do have write access.  The end result should be the file
\sphinxcode{\sphinxupquote{/usr/local/lib/libmpich.so}}.  On many Linux installations, the folder
\sphinxcode{\sphinxupquote{\$\{HOME\}/.local}} is an alternative to using \sphinxcode{\sphinxupquote{/usr/local}} and does
not require superuser or sudo access.  In that case the configuration
step becomes:

\begin{sphinxVerbatim}[commandchars=\\\{\}]
./configure\PYG{+w}{ }\PYGZhy{}\PYGZhy{}enable\PYGZhy{}shared\PYG{+w}{ }\PYGZhy{}\PYGZhy{}prefix\PYG{o}{=}\PYG{l+s+si}{\PYGZdl{}\PYGZob{}}\PYG{n+nv}{HOME}\PYG{l+s+si}{\PYGZcb{}}/.local
\end{sphinxVerbatim}

\sphinxAtStartPar
Avoiding the use of “sudo” for custom software installation (i.e. from
source and not through a package manager tool provided by the OS) is
generally recommended to ensure the integrity of the system software
installation.


\bigskip\hrule\bigskip



\subsection{Including or removing debug support}
\label{\detokenize{Build_basics:including-or-removing-debug-support}}\label{\detokenize{Build_basics:debug}}
\sphinxAtStartPar
By default the compilation settings will include the \sphinxcode{\sphinxupquote{\sphinxhyphen{}g}} flag which
instructs the compiler to include debug information (e.g. which line of
source code a particular instruction correspond to).  This can be
extremely useful in case LAMMPS crashes and can help to provide crucial
information in {\hyperref[\detokenize{Errors_debug::doc}]{\sphinxcrossref{\DUrole{doc}{tracking down the origin of a crash}}}}
and help the LAMMPS developers fix bugs in the source code.  However,
this increases the storage requirements for object files, libraries, and
the executable 3\sphinxhyphen{}5 fold.

\sphinxAtStartPar
If this is a concern, you can change the compilation settings or remove
the debug information from the LAMMPS executable:
\begin{itemize}
\item {} 
\sphinxAtStartPar
\sphinxstylestrong{Traditional make}: edit your \sphinxcode{\sphinxupquote{Makefile.\textless{}machine\textgreater{}}} to remove the
\sphinxcode{\sphinxupquote{\sphinxhyphen{}g}} flag from the \sphinxcode{\sphinxupquote{CCFLAGS}} and \sphinxcode{\sphinxupquote{LINKFLAGS}} definitions

\item {} 
\sphinxAtStartPar
\sphinxstylestrong{CMake}: use \sphinxcode{\sphinxupquote{\sphinxhyphen{}D CMAKE\_BUILD\_TYPE=Release}} or explicitly reset
the applicable compiler flags (best done using the text mode or
graphical user interface).

\item {} 
\sphinxAtStartPar
\sphinxstylestrong{Remove debug info}: If you are only concerned about the executable
being too large, you can use the \sphinxcode{\sphinxupquote{strip}} tool (e.g. \sphinxcode{\sphinxupquote{strip
lmp\_serial}}) to remove the debug information from the executable file.
Do not strip libraries or object files, as that will render them unusable.

\end{itemize}


\bigskip\hrule\bigskip



\subsection{Build LAMMPS tools}
\label{\detokenize{Build_basics:build-lammps-tools}}\label{\detokenize{Build_basics:tools}}
\sphinxAtStartPar
Some tools described in {\hyperref[\detokenize{Tools::doc}]{\sphinxcrossref{\DUrole{doc}{Auxiliary tools}}}} can be built directly
using CMake or Make.
\begin{quote}

\begin{sphinxadmonition}{note}{CMake build}

\begin{sphinxVerbatim}[commandchars=\\\{\}]
\PYGZhy{}D\PYG{+w}{ }\PYG{n+nv}{BUILD\PYGZus{}TOOLS}\PYG{o}{=}value\PYG{+w}{         }\PYG{c+c1}{\PYGZsh{} yes or no (default). Build binary2txt,}
\PYG{+w}{                             }\PYG{c+c1}{\PYGZsh{} chain.x, micelle2d.x, msi2lmp, phana,}
\PYG{+w}{                             }\PYG{c+c1}{\PYGZsh{} stl\PYGZus{}bin2txt}
\PYGZhy{}D\PYG{+w}{ }\PYG{n+nv}{BUILD\PYGZus{}LAMMPS\PYGZus{}GUI}\PYG{o}{=}value\PYG{+w}{    }\PYG{c+c1}{\PYGZsh{} yes or no (default). Build LAMMPS\PYGZhy{}GUI}
\end{sphinxVerbatim}

\sphinxAtStartPar
The generated binaries will also become part of the LAMMPS installation
(see below).
\end{sphinxadmonition}

\begin{sphinxadmonition}{note}{Traditional make}

\begin{sphinxVerbatim}[commandchars=\\\{\}]
\PYG{n+nb}{cd}\PYG{+w}{ }lammps/tools
make\PYG{+w}{ }all\PYG{+w}{              }\PYG{c+c1}{\PYGZsh{} build all binaries of tools}
make\PYG{+w}{ }binary2txt\PYG{+w}{       }\PYG{c+c1}{\PYGZsh{} build only binary2txt tool}
make\PYG{+w}{ }chain\PYG{+w}{            }\PYG{c+c1}{\PYGZsh{} build only chain tool}
make\PYG{+w}{ }micelle2d\PYG{+w}{        }\PYG{c+c1}{\PYGZsh{} build only micelle2d tool}
\end{sphinxVerbatim}

\begin{sphinxadmonition}{note}{Note:}
\sphinxAtStartPar
Building the LAMMPS\sphinxhyphen{}GUI \sphinxstyleemphasis{requires} building LAMMPS with CMake.
\end{sphinxadmonition}
\end{sphinxadmonition}
\end{quote}


\bigskip\hrule\bigskip



\subsection{Install LAMMPS after a build}
\label{\detokenize{Build_basics:install-lammps-after-a-build}}\label{\detokenize{Build_basics:install}}
\sphinxAtStartPar
After building LAMMPS, you may wish to copy the LAMMPS executable or
library, along with other LAMMPS files (library header, doc files), to a
globally visible place on your system, for others to access.  Note that
you may need super\sphinxhyphen{}user privileges (e.g. sudo) if the directory you want
to copy files to is protected.
\begin{quote}

\begin{sphinxadmonition}{note}{CMake build}

\begin{sphinxVerbatim}[commandchars=\\\{\}]
cmake\PYG{+w}{ }\PYGZhy{}D\PYG{+w}{ }\PYG{n+nv}{CMAKE\PYGZus{}INSTALL\PYGZus{}PREFIX}\PYG{o}{=}path\PYG{+w}{ }\PYG{o}{[}options\PYG{+w}{ }...\PYG{o}{]}\PYG{+w}{ }../cmake
make\PYG{+w}{                        }\PYG{c+c1}{\PYGZsh{} perform make after CMake command}
make\PYG{+w}{ }install\PYG{+w}{                }\PYG{c+c1}{\PYGZsh{} perform the installation into prefix}
\end{sphinxVerbatim}

\sphinxAtStartPar
During the installation process CMake will by default remove any runtime
path settings for loading shared libraries.  Because of this you may
have to set or modify the \sphinxcode{\sphinxupquote{LD\_LIBRARY\_PATH}} (or \sphinxcode{\sphinxupquote{DYLD\_LIBRARY\_PATH}})
environment variable, if you are installing LAMMPS into a non\sphinxhyphen{}system
location and/or are linking to libraries in a non\sphinxhyphen{}system location that
depend on such runtime path settings.
As an alternative, you may set the CMake variable \sphinxcode{\sphinxupquote{LAMMPS\_INSTALL\_RPATH}}
to \sphinxcode{\sphinxupquote{on}} and then the runtime paths for any linked shared libraries
and the library installation folder for the LAMMPS library will be
embedded and thus the requirement to set environment variables is avoided.
The \sphinxcode{\sphinxupquote{off}} setting is usually preferred for packaged binaries or when
setting up environment modules, the \sphinxcode{\sphinxupquote{on}} setting is more convenient
for installing software into a non\sphinxhyphen{}system or personal folder.
\end{sphinxadmonition}

\begin{sphinxadmonition}{note}{Traditional make}

\sphinxAtStartPar
There is no “install” option in the \sphinxcode{\sphinxupquote{src/Makefile}} for LAMMPS.
If you wish to do this you will need to first build LAMMPS, then
manually copy the desired LAMMPS files to the appropriate system
directories.
\end{sphinxadmonition}
\end{quote}

\sphinxstepscope

\clearpage


\section{Optional build settings}
\label{\detokenize{Build_settings:optional-build-settings}}\label{\detokenize{Build_settings::doc}}
\sphinxAtStartPar
LAMMPS can be built with several optional settings.  Each subsection
explains how to do this for building both with CMake and make.
\begin{itemize}
\item {} 
\sphinxAtStartPar
{\hyperref[\detokenize{Build_settings:c-11-standard-compliance}]{\sphinxcrossref{C++11 standard compliance}}} when building all of LAMMPS

\item {} 
\sphinxAtStartPar
{\hyperref[\detokenize{Build_settings:fft-library}]{\sphinxcrossref{FFT library}}} for use with the {\hyperref[\detokenize{kspace_style::doc}]{\sphinxcrossref{\DUrole{doc}{kspace\_style pppm}}}} command

\item {} 
\sphinxAtStartPar
{\hyperref[\detokenize{Build_settings:size-of-lammps-integer-types-and-size-limits}]{\sphinxcrossref{Size of LAMMPS integer types and size limits}}}

\item {} 
\sphinxAtStartPar
{\hyperref[\detokenize{Build_settings:read-or-write-compressed-files}]{\sphinxcrossref{Read or write compressed files}}}

\item {} 
\sphinxAtStartPar
{\hyperref[\detokenize{Build_settings:output-of-jpeg-png-and-movie-files}]{\sphinxcrossref{Output of JPEG, PNG, and movie files}}} via the {\hyperref[\detokenize{dump_image::doc}]{\sphinxcrossref{\DUrole{doc}{dump image}}}} or {\hyperref[\detokenize{dump_image::doc}]{\sphinxcrossref{\DUrole{doc}{dump movie}}}} commands

\item {} 
\sphinxAtStartPar
{\hyperref[\detokenize{Build_settings:support-for-downloading-files}]{\sphinxcrossref{Support for downloading files}}}

\item {} 
\sphinxAtStartPar
{\hyperref[\detokenize{Build_settings:memory-allocation-alignment}]{\sphinxcrossref{Memory allocation alignment}}}

\item {} 
\sphinxAtStartPar
{\hyperref[\detokenize{Build_settings:workaround-for-long-long-integers}]{\sphinxcrossref{Workaround for long long integers}}}

\item {} 
\sphinxAtStartPar
{\hyperref[\detokenize{Build_settings:exception-handling-when-using-lammps-as-a-library}]{\sphinxcrossref{Exception handling when using LAMMPS as a library}}} to capture errors

\item {} 
\sphinxAtStartPar
{\hyperref[\detokenize{Build_settings:trigger-selected-floating-point-exceptions}]{\sphinxcrossref{Trigger selected floating\sphinxhyphen{}point exceptions}}}

\end{itemize}


\bigskip\hrule\bigskip



\subsection{C++11 standard compliance}
\label{\detokenize{Build_settings:c-11-standard-compliance}}\label{\detokenize{Build_settings:cxx11}}
\sphinxAtStartPar
A C++11 standard compatible compiler is a requirement for compiling LAMMPS.
LAMMPS version 3 March 2020 is the last version compatible with the previous
C++98 standard for the core code and most packages. Most currently used
C++ compilers are compatible with C++11, but some older ones may need extra
flags to enable C++11 compliance.  Example for GNU c++ 4.8.x:

\begin{sphinxVerbatim}[commandchars=\\\{\}]
\PYG{n+nv}{CCFLAGS}\PYG{+w}{ }\PYG{o}{=}\PYG{+w}{ }\PYGZhy{}g\PYG{+w}{ }\PYGZhy{}O3\PYG{+w}{ }\PYGZhy{}std\PYG{o}{=}c++11
\end{sphinxVerbatim}

\sphinxAtStartPar
Individual packages may require compliance with a later C++ standard
like C++14 or C++17.  These requirements will be documented with the
{\hyperref[\detokenize{Packages_details::doc}]{\sphinxcrossref{\DUrole{doc}{individual packages}}}}.


\bigskip\hrule\bigskip



\subsection{FFT library}
\label{\detokenize{Build_settings:fft-library}}\label{\detokenize{Build_settings:fft}}
\sphinxAtStartPar
When the KSPACE package is included in a LAMMPS build, the
{\hyperref[\detokenize{kspace_style::doc}]{\sphinxcrossref{\DUrole{doc}{kspace\_style pppm}}}} command performs 3d FFTs which
require use of an FFT library to compute 1d FFTs.  The KISS FFT
library is included with LAMMPS, but other libraries can be faster.
LAMMPS can use them if they are available on your system.

\sphinxAtStartPar
\DUrole{versionmodified,added}{Added in version 7Feb2024.}

\sphinxAtStartPar
Alternatively, LAMMPS can use the \sphinxhref{https://icl-utk-edu.github.io/heffte/}{heFFTe} library for the MPI
communication algorithms, which comes with many optimizations for
special cases, e.g. leveraging available 2D and 3D FFTs in the back end
libraries and better pipelining for packing and communication.
\begin{quote}

\begin{sphinxadmonition}{note}{CMake build}

\begin{sphinxVerbatim}[commandchars=\\\{\}]
\PYGZhy{}D\PYG{+w}{ }\PYG{n+nv}{FFT}\PYG{o}{=}value\PYG{+w}{              }\PYG{c+c1}{\PYGZsh{} FFTW3 or MKL or KISS, default is FFTW3 if found,}
\PYG{+w}{                          }\PYG{c+c1}{\PYGZsh{} else KISS}
\PYGZhy{}D\PYG{+w}{ }\PYG{n+nv}{FFT\PYGZus{}KOKKOS}\PYG{o}{=}value\PYG{+w}{       }\PYG{c+c1}{\PYGZsh{} FFTW3 or MKL or KISS or CUFFT or HIPFFT,}
\PYG{+w}{                          }\PYG{c+c1}{\PYGZsh{} default is KISS}
\PYGZhy{}D\PYG{+w}{ }\PYG{n+nv}{FFT\PYGZus{}SINGLE}\PYG{o}{=}value\PYG{+w}{       }\PYG{c+c1}{\PYGZsh{} yes or no (default), no = double precision}
\PYGZhy{}D\PYG{+w}{ }\PYG{n+nv}{FFT\PYGZus{}PACK}\PYG{o}{=}value\PYG{+w}{         }\PYG{c+c1}{\PYGZsh{} array (default) or pointer or memcpy}
\PYGZhy{}D\PYG{+w}{ }\PYG{n+nv}{FFT\PYGZus{}USE\PYGZus{}HEFFTE}\PYG{o}{=}value\PYG{+w}{   }\PYG{c+c1}{\PYGZsh{} yes or no (default), yes links to heFFTe}
\end{sphinxVerbatim}

\begin{sphinxadmonition}{note}{Note:}
\sphinxAtStartPar
When the Kokkos variant of a package is compiled and selected at run time,
the FFT library selected by the \sphinxcode{\sphinxupquote{FFT\_KOKKOS}} variable applies. Otherwise,
the FFT library selected by the FFT variable applies.
The same FFT settings apply to both. \sphinxcode{\sphinxupquote{FFT\_KOKKOS}} must be compatible with the
Kokkos back end \sphinxhyphen{} for example, when using the CUDA back end of Kokkos,
you must use either \sphinxcode{\sphinxupquote{CUFFT}} or \sphinxcode{\sphinxupquote{KISS}}.
\end{sphinxadmonition}

\sphinxAtStartPar
Usually these settings are all that is needed.  If FFTW3 is
selected, then CMake will try to detect, if threaded FFTW
libraries are available and enable them by default.  This setting
is independent of whether OpenMP threads are enabled and a package
like KOKKOS or OPENMP is used.  If CMake cannot detect the FFT
library, you can set these variables to assist:

\begin{sphinxVerbatim}[commandchars=\\\{\}]
\PYGZhy{}D\PYG{+w}{ }\PYG{n+nv}{FFTW3\PYGZus{}INCLUDE\PYGZus{}DIR}\PYG{o}{=}path\PYG{+w}{   }\PYG{c+c1}{\PYGZsh{} path to FFTW3 include files}
\PYGZhy{}D\PYG{+w}{ }\PYG{n+nv}{FFTW3\PYGZus{}LIBRARY}\PYG{o}{=}path\PYG{+w}{       }\PYG{c+c1}{\PYGZsh{} path to FFTW3 libraries}
\PYGZhy{}D\PYG{+w}{ }\PYG{n+nv}{FFTW3\PYGZus{}OMP\PYGZus{}LIBRARY}\PYG{o}{=}path\PYG{+w}{   }\PYG{c+c1}{\PYGZsh{} path to FFTW3 OpenMP wrapper libraries}
\PYGZhy{}D\PYG{+w}{ }\PYG{n+nv}{FFT\PYGZus{}FFTW\PYGZus{}THREADS}\PYG{o}{=}on\PYG{+w}{      }\PYG{c+c1}{\PYGZsh{} enable using OpenMP threaded FFTW3 libraries}
\PYGZhy{}D\PYG{+w}{ }\PYG{n+nv}{MKL\PYGZus{}INCLUDE\PYGZus{}DIR}\PYG{o}{=}path\PYG{+w}{     }\PYG{c+c1}{\PYGZsh{} ditto for Intel MKL library}
\PYGZhy{}D\PYG{+w}{ }\PYG{n+nv}{FFT\PYGZus{}MKL\PYGZus{}THREADS}\PYG{o}{=}on\PYG{+w}{       }\PYG{c+c1}{\PYGZsh{} enable using threaded FFTs with MKL libraries}
\PYGZhy{}D\PYG{+w}{ }\PYG{n+nv}{MKL\PYGZus{}LIBRARY}\PYG{o}{=}path\PYG{+w}{         }\PYG{c+c1}{\PYGZsh{} path to MKL libraries}
\PYGZhy{}D\PYG{+w}{ }\PYG{n+nv}{FFT\PYGZus{}HEFFTE\PYGZus{}BACKEND}\PYG{o}{=}value\PYG{+w}{ }\PYG{c+c1}{\PYGZsh{} FFTW or MKL or empty/undefined for the stock}
\PYG{+w}{                            }\PYG{c+c1}{\PYGZsh{} heFFTe back end}
\PYGZhy{}D\PYG{+w}{ }\PYG{n+nv}{Heffte\PYGZus{}ROOT}\PYG{o}{=}path\PYG{+w}{         }\PYG{c+c1}{\PYGZsh{} path to an existing heFFTe installation}
\end{sphinxVerbatim}

\begin{sphinxadmonition}{note}{Note:}
\sphinxAtStartPar
heFFTe comes with a builtin (= stock) back end for FFTs, i.e. a
default internal FFT implementation; however, this stock back
end is intended for testing purposes only and is not optimized
for production runs.
\end{sphinxadmonition}
\end{sphinxadmonition}

\begin{sphinxadmonition}{note}{Traditional make}

\sphinxAtStartPar
To change the FFT library to be used and its options, you have to edit
your machine Makefile. Below are examples how the makefile variables
could be changed.

\begin{sphinxVerbatim}[commandchars=\\\{\}]
\PYG{n+nv}{FFT\PYGZus{}INC}\PYG{+w}{ }\PYG{o}{=}\PYG{+w}{ }\PYGZhy{}DFFT\PYGZus{}\PYGZlt{}NAME\PYGZgt{}\PYG{+w}{        }\PYG{c+c1}{\PYGZsh{} where \PYGZlt{}NAME\PYGZgt{} is KISS (default), FFTW3,}
\PYG{c}{                              \PYGZsh{} FFTW (same as FFTW3), or MKL}
\PYG{n+nv}{FFT\PYGZus{}INC}\PYG{+w}{ }\PYG{o}{=}\PYG{+w}{ }\PYGZhy{}DFFT\PYGZus{}KOKKOS\PYGZus{}\PYGZlt{}NAME\PYGZgt{}\PYG{+w}{ }\PYG{c+c1}{\PYGZsh{} where \PYGZlt{}NAME\PYGZgt{} is KISS (default), FFTW3,}
\PYG{c}{                              \PYGZsh{} FFTW (same as FFTW3), MKL, CUFFT, or HIPFFT}
\PYG{n+nv}{FFT\PYGZus{}INC}\PYG{+w}{ }\PYG{o}{=}\PYG{+w}{ }\PYGZhy{}DFFT\PYGZus{}SINGLE\PYG{+w}{       }\PYG{c+c1}{\PYGZsh{} do not specify for double precision}
\PYG{n+nv}{FFT\PYGZus{}INC}\PYG{+w}{ }\PYG{o}{=}\PYG{+w}{ }\PYGZhy{}DFFT\PYGZus{}FFTW\PYGZus{}THREADS\PYG{+w}{ }\PYG{c+c1}{\PYGZsh{} enable using threaded FFTW3 libraries}
\PYG{n+nv}{FFT\PYGZus{}INC}\PYG{+w}{ }\PYG{o}{=}\PYG{+w}{ }\PYGZhy{}DFFT\PYGZus{}MKL\PYGZus{}THREADS\PYG{+w}{  }\PYG{c+c1}{\PYGZsh{} enable using threaded FFTs with MKL libraries}
\PYG{n+nv}{FFT\PYGZus{}INC}\PYG{+w}{ }\PYG{o}{=}\PYG{+w}{ }\PYGZhy{}DFFT\PYGZus{}PACK\PYGZus{}ARRAY\PYG{+w}{   }\PYG{c+c1}{\PYGZsh{} or \PYGZhy{}DFFT\PYGZus{}PACK\PYGZus{}POINTER or \PYGZhy{}DFFT\PYGZus{}PACK\PYGZus{}MEMCPY}
\PYG{c}{                             \PYGZsh{} default is FFT\PYGZus{}PACK\PYGZus{}ARRAY if not specified}
\end{sphinxVerbatim}

\begin{sphinxVerbatim}[commandchars=\\\{\}]
\PYG{n+nv}{FFT\PYGZus{}INC}\PYG{+w}{ }\PYG{o}{=}\PYG{+w}{  }\PYGZhy{}I/usr/local/include
\PYG{n+nv}{FFT\PYGZus{}PATH}\PYG{+w}{ }\PYG{o}{=}\PYG{+w}{ }\PYGZhy{}L/usr/local/lib

\PYG{c}{\PYGZsh{} hipFFT either precision}
\PYG{n+nv}{FFT\PYGZus{}LIB}\PYG{+w}{ }\PYG{o}{=}\PYG{+w}{  }\PYGZhy{}lhipfft

\PYG{c}{\PYGZsh{} cuFFT either precision}
\PYG{n+nv}{FFT\PYGZus{}LIB}\PYG{+w}{ }\PYG{o}{=}\PYG{+w}{  }\PYGZhy{}lcufft

\PYG{c}{\PYGZsh{} FFTW3 double precision}
\PYG{n+nv}{FFT\PYGZus{}LIB}\PYG{+w}{ }\PYG{o}{=}\PYG{+w}{  }\PYGZhy{}lfftw3

\PYG{c}{\PYGZsh{} FFTW3 double precision with threads (needs \PYGZhy{}DFFT\PYGZus{}FFTW\PYGZus{}THREADS)}
\PYG{n+nv}{FFT\PYGZus{}LIB}\PYG{+w}{ }\PYG{o}{=}\PYG{+w}{  }\PYGZhy{}lfftw3\PYG{+w}{ }\PYGZhy{}lfftw3\PYGZus{}omp

\PYG{c}{\PYGZsh{} FFTW3 single precision}
\PYG{n+nv}{FFT\PYGZus{}LIB}\PYG{+w}{ }\PYG{o}{=}\PYG{+w}{  }\PYGZhy{}lfftw3\PYG{+w}{ }\PYGZhy{}lfftw3f

\PYG{c}{\PYGZsh{} serial MKL with Intel compiler}
\PYG{n+nv}{FFT\PYGZus{}LIB}\PYG{+w}{ }\PYG{o}{=}\PYG{+w}{  }\PYGZhy{}lmkl\PYGZus{}intel\PYGZus{}lp64\PYG{+w}{ }\PYGZhy{}lmkl\PYGZus{}sequential\PYG{+w}{ }\PYGZhy{}lmkl\PYGZus{}core

\PYG{c}{\PYGZsh{} serial MKL with GNU compiler}
\PYG{n+nv}{FFT\PYGZus{}LIB}\PYG{+w}{ }\PYG{o}{=}\PYG{+w}{  }\PYGZhy{}lmkl\PYGZus{}gf\PYGZus{}lp64\PYG{+w}{ }\PYGZhy{}lmkl\PYGZus{}sequential\PYG{+w}{ }\PYGZhy{}lmkl\PYGZus{}core

\PYG{c}{\PYGZsh{} threaded MKL with Intel compiler}
\PYG{n+nv}{FFT\PYGZus{}LIB}\PYG{+w}{ }\PYG{o}{=}\PYG{+w}{  }\PYGZhy{}lmkl\PYGZus{}intel\PYGZus{}lp64\PYG{+w}{ }\PYGZhy{}lmkl\PYGZus{}intel\PYGZus{}thread\PYG{+w}{ }\PYGZhy{}lmkl\PYGZus{}core

\PYG{c}{\PYGZsh{} threaded MKL with GNU compiler}
\PYG{n+nv}{FFT\PYGZus{}LIB}\PYG{+w}{ }\PYG{o}{=}\PYG{+w}{  }\PYGZhy{}lmkl\PYGZus{}gf\PYGZus{}lp64\PYG{+w}{ }\PYGZhy{}lmkl\PYGZus{}gnu\PYGZus{}thread\PYG{+w}{ }\PYGZhy{}lmkl\PYGZus{}core

\PYG{c}{\PYGZsh{} MKL with automatic runtime selection of interface libs}
\PYG{n+nv}{FFT\PYGZus{}LIB}\PYG{+w}{ }\PYG{o}{=}\PYG{+w}{  }\PYGZhy{}lmkl\PYGZus{}rt
\end{sphinxVerbatim}

\sphinxAtStartPar
As with CMake, you do not need to set paths in \sphinxcode{\sphinxupquote{FFT\_INC}} or
\sphinxcode{\sphinxupquote{FFT\_PATH}}, if the compiler can find the FFT header and library
files in its default search path.  You must specify \sphinxcode{\sphinxupquote{FFT\_LIB}}
with the appropriate FFT libraries to include in the link.

\sphinxAtStartPar
Traditional make can also link to heFFTe using an existing installation

\begin{sphinxVerbatim}[commandchars=\\\{\}]
\PYG{c+cp}{include \PYGZlt{}path\PYGZhy{}to\PYGZhy{}heffte\PYGZhy{}installation\PYGZgt{}/share/heffte/HeffteMakefile.in}
\PYG{n+nv}{FFT\PYGZus{}INC}\PYG{+w}{ }\PYG{o}{=}\PYG{+w}{ }\PYGZhy{}DFFT\PYGZus{}HEFFTE\PYG{+w}{ }\PYGZhy{}DFFT\PYGZus{}HEFFTE\PYGZus{}FFTW\PYG{+w}{ }\PYG{k}{\PYGZdl{}(}heffte\PYGZus{}include\PYG{k}{)}
\PYG{n+nv}{FFT\PYGZus{}PATH}\PYG{+w}{ }\PYG{o}{=}
\PYG{n+nv}{FFT\PYGZus{}LIB}\PYG{+w}{ }\PYG{o}{=}\PYG{+w}{ }\PYG{k}{\PYGZdl{}(}heffte\PYGZus{}link\PYG{k}{)}\PYG{+w}{ }\PYG{k}{\PYGZdl{}(}heffte\PYGZus{}libs\PYG{k}{)}
\end{sphinxVerbatim}

\sphinxAtStartPar
The heFFTe install path will contain \sphinxcode{\sphinxupquote{HeffteMakefile.in}}.
which will define the \sphinxcode{\sphinxupquote{heffte\_}} include variables needed to link to heFFTe from
an external project using traditional make.
The \sphinxcode{\sphinxupquote{\sphinxhyphen{}DFFT\_HEFFTE}} is required to switch to using heFFTe, while the optional \sphinxcode{\sphinxupquote{\sphinxhyphen{}DFFT\_HEFFTE\_FFTW}}
selects the desired heFFTe back end, e.g., \sphinxcode{\sphinxupquote{\sphinxhyphen{}DFFT\_HEFFTE\_FFTW}} or \sphinxcode{\sphinxupquote{\sphinxhyphen{}DFFT\_HEFFTE\_MKL}},
omitting the variable will default to the \sphinxtitleref{stock} back end.
The heFFTe \sphinxtitleref{stock} back end is intended to be used for testing and debugging,
but is not performance optimized for large scale production runs.
\end{sphinxadmonition}
\end{quote}

\sphinxAtStartPar
The \sphinxhref{https://github.com/mborgerding/kissfft}{KISS FFT library} is
included in the LAMMPS distribution.  It is portable across all
platforms.  Depending on the size of the FFTs and the number of
processors used, the other libraries listed here can be faster.

\sphinxAtStartPar
However, note that long\sphinxhyphen{}range Coulombics are only a portion of the
per\sphinxhyphen{}timestep CPU cost, FFTs are only a portion of long\sphinxhyphen{}range Coulombics,
and 1d FFTs are only a portion of the FFT cost (parallel communication
can be costly).  A breakdown of these timings is printed to the screen
at the end of a run when using the {\hyperref[\detokenize{kspace_style::doc}]{\sphinxcrossref{\DUrole{doc}{kspace\_style pppm}}}} command. The {\hyperref[\detokenize{Run_output::doc}]{\sphinxcrossref{\DUrole{doc}{Screen and logfile output}}}} page gives more details.  A more detailed (and time
consuming) report of the FFT performance is generated with the
{\hyperref[\detokenize{kspace_modify::doc}]{\sphinxcrossref{\DUrole{doc}{kspace\_modify fftbench yes}}}} command.

\sphinxAtStartPar
FFTW is a fast, portable FFT library that should also work on any
platform and can be faster than the KISS FFT library.  You can download
it from \sphinxhref{https://www.fftw.org}{www.fftw.org}.  LAMMPS requires version
3.X; the legacy version 2.1.X is no longer supported.

\sphinxAtStartPar
Building FFTW for your box should be as simple as \sphinxcode{\sphinxupquote{./configure; make;
make install}}.  The install command typically requires root privileges
(e.g. invoke it via sudo), unless you specify a local directory with
the \sphinxcode{\sphinxupquote{\sphinxhyphen{}\sphinxhyphen{}prefix}} option of configure.  Type \sphinxcode{\sphinxupquote{./configure \sphinxhyphen{}\sphinxhyphen{}help}} to see
various options.

\sphinxAtStartPar
The Intel MKL math library is part of the Intel compiler suite.  It
can be used with the Intel or GNU compiler (see the \sphinxcode{\sphinxupquote{FFT\_LIB}} setting
above).

\sphinxAtStartPar
The cuFFT and hipFFT FFT libraries are packaged with NVIDIA’s CUDA and
AMD’s HIP installations, respectively. These FFT libraries require the
Kokkos acceleration package to be enabled and the Kokkos back end to be
GPU\sphinxhyphen{}resident (i.e., HIP or CUDA).

\sphinxAtStartPar
Performing 3d FFTs in parallel can be time\sphinxhyphen{}consuming due to data access
and required communication.  This cost can be reduced by performing
single\sphinxhyphen{}precision FFTs instead of double precision.  Single precision
means the real and imaginary parts of a complex datum are 4\sphinxhyphen{}byte floats.
Double precision means they are 8\sphinxhyphen{}byte doubles.  Note that Fourier
transform and related PPPM operations are somewhat less sensitive to
floating point truncation errors, and thus the resulting error is
generally less than the difference in precision. Using the
\sphinxcode{\sphinxupquote{\sphinxhyphen{}DFFT\_SINGLE}} setting trades off a little accuracy for reduced memory
use and parallel communication costs for transposing 3d FFT data.

\sphinxAtStartPar
When using \sphinxcode{\sphinxupquote{\sphinxhyphen{}DFFT\_SINGLE}} with FFTW3, you may need to ensure that
the FFTW3 installation includes support for single\sphinxhyphen{}precision.

\sphinxAtStartPar
When compiler FFTW3 from source, you can do the following, which should
produce the additional libraries \sphinxcode{\sphinxupquote{libfftw3f.a}} and/or \sphinxcode{\sphinxupquote{libfftw3f.so}}.

\begin{sphinxVerbatim}[commandchars=\\\{\}]
make\PYG{+w}{ }clean
./configure\PYG{+w}{ }\PYGZhy{}\PYGZhy{}enable\PYGZhy{}single\PYG{p}{;}\PYG{+w}{ }make\PYG{p}{;}\PYG{+w}{ }make\PYG{+w}{ }install
\end{sphinxVerbatim}

\sphinxAtStartPar
Performing 3d FFTs requires communication to transpose the 3d FFT
grid.  The data packing/unpacking for this can be done in one of 3
modes (ARRAY, POINTER, MEMCPY) as set by the \sphinxcode{\sphinxupquote{FFT\_PACK}} syntax above.
Depending on the machine, the size of the FFT grid, the number of
processors used, one option may be slightly faster.  The default is
ARRAY mode.

\sphinxAtStartPar
When using \sphinxcode{\sphinxupquote{\sphinxhyphen{}DFFT\_HEFFTE}} CMake will first look for an existing
install with hints provided by \sphinxcode{\sphinxupquote{\sphinxhyphen{}DHeffte\_ROOT}}, as recommended by the
CMake standard and note that the name is case sensitive. If CMake cannot
find a heFFTe installation with the correct back end (e.g., FFTW or
MKL), it will attempt to download and build the library automatically.
In this case, LAMMPS CMake will also accept all heFFTe specific
variables listed in the \sphinxhref{https://mkstoyanov.bitbucket.io/heffte/md\_doxygen\_installation.html}{heFFTe documentation}
and those variables will be passed into the heFFTe build.


\bigskip\hrule\bigskip


\clearpage


\subsection{Size of LAMMPS integer types and size limits}
\label{\detokenize{Build_settings:size-of-lammps-integer-types-and-size-limits}}\label{\detokenize{Build_settings:size}}
\sphinxAtStartPar
LAMMPS uses a few custom integer data types, which can be defined as
either 4\sphinxhyphen{}byte (= 32\sphinxhyphen{}bit) or 8\sphinxhyphen{}byte (= 64\sphinxhyphen{}bit) integers at compile time.
This has an impact on the size of a system that can be simulated, or how
large counters can become before “rolling over”.  The default setting of
“smallbig” is almost always adequate.
\begin{quote}

\begin{sphinxadmonition}{note}{CMake build}

\sphinxAtStartPar
With CMake the choice of integer types is made via setting a
variable during configuration.

\begin{sphinxVerbatim}[commandchars=\\\{\}]
\PYGZhy{}D\PYG{+w}{ }\PYG{n+nv}{LAMMPS\PYGZus{}SIZES}\PYG{o}{=}value\PYG{+w}{   }\PYG{c+c1}{\PYGZsh{} smallbig (default) or bigbig or smallsmall}
\end{sphinxVerbatim}

\sphinxAtStartPar
If the variable is not set explicitly, “smallbig” is used.
\end{sphinxadmonition}

\begin{sphinxadmonition}{note}{Traditional build}

\sphinxAtStartPar
If you want a setting different from the default, you need to edit the
\sphinxcode{\sphinxupquote{LMP\_INC}} variable setting your machine Makefile.

\begin{sphinxVerbatim}[commandchars=\\\{\}]
\PYG{n+nv}{LMP\PYGZus{}INC}\PYG{+w}{ }\PYG{o}{=}\PYG{+w}{ }\PYGZhy{}DLAMMPS\PYGZus{}SMALLBIG\PYG{+w}{    }\PYG{c+c1}{\PYGZsh{} or \PYGZhy{}DLAMMPS\PYGZus{}BIGBIG or \PYGZhy{}DLAMMPS\PYGZus{}SMALLSMALL}
\end{sphinxVerbatim}

\sphinxAtStartPar
The default setting is \sphinxcode{\sphinxupquote{\sphinxhyphen{}DLAMMPS\_SMALLBIG}} if nothing is specified
\end{sphinxadmonition}
\end{quote}


\subsubsection{LAMMPS system size restrictions}
\label{\detokenize{Build_settings:lammps-system-size-restrictions}}

\begin{savenotes}\sphinxattablestart
\sphinxthistablewithglobalstyle
\centering
\begin{tabular}[t]{\X{18}{100}\X{27}{100}\X{28}{100}\X{27}{100}}
\sphinxtoprule
\sphinxstyletheadfamily &\sphinxstyletheadfamily 
\sphinxAtStartPar
smallbig
&\sphinxstyletheadfamily 
\sphinxAtStartPar
bigbig
&\sphinxstyletheadfamily 
\sphinxAtStartPar
smallsmall
\\
\sphinxmidrule
\sphinxtableatstartofbodyhook
\sphinxAtStartPar
Total atom count
&
\sphinxAtStartPar
\(2^{63}\) atoms (= \(9.223 \cdot 10^{18}\))
&
\sphinxAtStartPar
\(2^{63}\) atoms (= \(9.223 \cdot 10^{18}\))
&
\sphinxAtStartPar
\(2^{31}\) atoms (= \(2.147 \cdot 10^9\))
\\
\sphinxhline
\sphinxAtStartPar
Total timesteps
&
\sphinxAtStartPar
\(2^{63}\) steps (= \(9.223 \cdot 10^{18}\))
&
\sphinxAtStartPar
\(2^{63}\) steps (= \(9.223 \cdot 10^{18}\))
&
\sphinxAtStartPar
\(2^{31}\) steps (= \(2.147 \cdot 10^9\))
\\
\sphinxhline
\sphinxAtStartPar
Atom ID values
&
\sphinxAtStartPar
\(1 \le i \le 2^{31} (= 2.147 \cdot 10^9)\)
&
\sphinxAtStartPar
\(1 \le i \le 2^{63} (= 9.223 \cdot 10^{18})\)
&
\sphinxAtStartPar
\(1 \le i \le 2^{31} (= 2.147 \cdot 10^9)\)
\\
\sphinxhline
\sphinxAtStartPar
Image flag values
&
\sphinxAtStartPar
\(-512 \le i \le 511\)
&
\sphinxAtStartPar
\(- 1\,048\,576 \le i \le 1\,048\,575\)
&
\sphinxAtStartPar
\(-512 \le i \le 511\)
\\
\sphinxbottomrule
\end{tabular}
\sphinxtableafterendhook\par
\sphinxattableend\end{savenotes}

\sphinxAtStartPar
The “bigbig” setting increases the size of image flags and atom IDs over
“smallbig” and the “smallsmall” setting is only needed if your machine
does not support 64\sphinxhyphen{}bit integers or incurs performance penalties when
using them.

\sphinxAtStartPar
These are limits for the core of the LAMMPS code, specific features or
some styles may impose additional limits.  The {\hyperref[\detokenize{Packages_details:pkg-atc}]{\sphinxcrossref{\DUrole{std,std-ref}{ATC}}}} package cannot be compiled with the “bigbig” setting.
Also, there are limitations when using the library interface where some
functions with known issues have been replaced by dummy calls printing a
corresponding error message rather than crashing randomly or corrupting
data.

\sphinxAtStartPar
Atom IDs are not required for atomic systems which do not store bond
topology information, though IDs are enabled by default.  The
{\hyperref[\detokenize{atom_modify::doc}]{\sphinxcrossref{\DUrole{doc}{atom\_modify id no}}}} command will turn them off.  Atom
IDs are required for molecular systems with bond topology (bonds,
angles, dihedrals, etc).  Similarly, some force or compute or fix styles
require atom IDs.  Thus, if you model a molecular system or use one of
those styles with more than 2 billion atoms, you need the “bigbig”
setting.

\sphinxAtStartPar
Regardless of the total system size limits, the maximum number of atoms
per MPI rank (local + ghost atoms) is limited to 2 billion for atomic
systems and 500 million for systems with bonds (the additional
restriction is due to using the 2 upper bits of the local atom index
in neighbor lists for storing special bonds info).

\sphinxAtStartPar
Image flags store 3 values per atom in a single integer, which count the
number of times an atom has moved through the periodic box in each
dimension.  See the {\hyperref[\detokenize{dump::doc}]{\sphinxcrossref{\DUrole{doc}{dump}}}} manual page for a discussion.  If
an atom moves through the periodic box more than this limit, the value
will “roll over”, e.g. from 511 to \sphinxhyphen{}512, which can cause diagnostics
like the mean\sphinxhyphen{}squared displacement, as calculated by the {\hyperref[\detokenize{compute_msd::doc}]{\sphinxcrossref{\DUrole{doc}{compute
msd}}}} command, to be faulty.

\sphinxAtStartPar
Also note that the GPU package requires its lib/gpu library to be
compiled with the same size setting, or the link will fail.  A CMake
build does this automatically.  When building with make, the setting
in whichever \sphinxcode{\sphinxupquote{lib/gpu/Makefile}} is used must be the same as above.


\bigskip\hrule\bigskip



\subsection{Output of JPEG, PNG, and movie files}
\label{\detokenize{Build_settings:output-of-jpeg-png-and-movie-files}}\label{\detokenize{Build_settings:graphics}}
\sphinxAtStartPar
The {\hyperref[\detokenize{dump_image::doc}]{\sphinxcrossref{\DUrole{doc}{dump image}}}} command has options to output JPEG or
PNG image files.  Likewise, the {\hyperref[\detokenize{dump_image::doc}]{\sphinxcrossref{\DUrole{doc}{dump movie}}}} command
outputs movie files in a variety of movie formats.  Using these options
requires the following settings:
\begin{quote}

\begin{sphinxadmonition}{note}{CMake build}

\begin{sphinxVerbatim}[commandchars=\\\{\}]
\PYGZhy{}D\PYG{+w}{ }\PYG{n+nv}{WITH\PYGZus{}JPEG}\PYG{o}{=}value\PYG{+w}{    }\PYG{c+c1}{\PYGZsh{} yes or no}
\PYG{+w}{                      }\PYG{c+c1}{\PYGZsh{} default = yes if CMake finds JPEG development files, else no}
\PYGZhy{}D\PYG{+w}{ }\PYG{n+nv}{WITH\PYGZus{}PNG}\PYG{o}{=}value\PYG{+w}{     }\PYG{c+c1}{\PYGZsh{} yes or no}
\PYG{+w}{                      }\PYG{c+c1}{\PYGZsh{} default = yes if CMake finds PNG and ZLIB development files,}
\PYG{+w}{                      }\PYG{c+c1}{\PYGZsh{} else no}
\PYGZhy{}D\PYG{+w}{ }\PYG{n+nv}{WITH\PYGZus{}FFMPEG}\PYG{o}{=}value\PYG{+w}{  }\PYG{c+c1}{\PYGZsh{} yes or no}
\PYG{+w}{                      }\PYG{c+c1}{\PYGZsh{} default = yes if CMake can find ffmpeg, else no}
\end{sphinxVerbatim}

\sphinxAtStartPar
Usually these settings are all that is needed.  If CMake cannot
find the graphics header, library, executable files, you can set
these variables:

\begin{sphinxVerbatim}[commandchars=\\\{\}]
\PYGZhy{}D\PYG{+w}{ }\PYG{n+nv}{JPEG\PYGZus{}INCLUDE\PYGZus{}DIR}\PYG{o}{=}path\PYG{+w}{    }\PYG{c+c1}{\PYGZsh{} path to jpeglib.h header file}
\PYGZhy{}D\PYG{+w}{ }\PYG{n+nv}{JPEG\PYGZus{}LIBRARY}\PYG{o}{=}path\PYG{+w}{        }\PYG{c+c1}{\PYGZsh{} path to libjpeg.a (.so) file}
\PYGZhy{}D\PYG{+w}{ }\PYG{n+nv}{PNG\PYGZus{}INCLUDE\PYGZus{}DIR}\PYG{o}{=}path\PYG{+w}{     }\PYG{c+c1}{\PYGZsh{} path to png.h header file}
\PYGZhy{}D\PYG{+w}{ }\PYG{n+nv}{PNG\PYGZus{}LIBRARY}\PYG{o}{=}path\PYG{+w}{         }\PYG{c+c1}{\PYGZsh{} path to libpng.a (.so) file}
\PYGZhy{}D\PYG{+w}{ }\PYG{n+nv}{ZLIB\PYGZus{}INCLUDE\PYGZus{}DIR}\PYG{o}{=}path\PYG{+w}{    }\PYG{c+c1}{\PYGZsh{} path to zlib.h header file}
\PYGZhy{}D\PYG{+w}{ }\PYG{n+nv}{ZLIB\PYGZus{}LIBRARY}\PYG{o}{=}path\PYG{+w}{        }\PYG{c+c1}{\PYGZsh{} path to libz.a (.so) file}
\PYGZhy{}D\PYG{+w}{ }\PYG{n+nv}{FFMPEG\PYGZus{}EXECUTABLE}\PYG{o}{=}path\PYG{+w}{   }\PYG{c+c1}{\PYGZsh{} path to ffmpeg executable}
\end{sphinxVerbatim}
\end{sphinxadmonition}

\begin{sphinxadmonition}{note}{Traditional make}

\begin{sphinxVerbatim}[commandchars=\\\{\}]
\PYG{n+nv}{LMP\PYGZus{}INC}\PYG{+w}{ }\PYG{o}{=}\PYG{+w}{ }\PYGZhy{}DLAMMPS\PYGZus{}JPEG\PYG{+w}{ }\PYGZhy{}DLAMMPS\PYGZus{}PNG\PYG{+w}{ }\PYGZhy{}DLAMMPS\PYGZus{}FFMPEG\PYG{+w}{  }\PYGZlt{}other\PYG{+w}{ }LMP\PYGZus{}INC\PYG{+w}{ }settings\PYGZgt{}

\PYG{n+nv}{JPG\PYGZus{}INC}\PYG{+w}{ }\PYG{o}{=}\PYG{+w}{ }\PYGZhy{}I/usr/local/include\PYG{+w}{   }\PYG{c+c1}{\PYGZsh{} path to jpeglib.h, png.h, zlib.h headers}
\PYG{c}{                                 \PYGZsh{} if make cannot find them}
\PYG{n+nv}{JPG\PYGZus{}PATH}\PYG{+w}{ }\PYG{o}{=}\PYG{+w}{ }\PYGZhy{}L/usr/lib\PYG{+w}{            }\PYG{c+c1}{\PYGZsh{} paths to libjpeg.a, libpng.a, libz.a (.so)}
\PYG{c}{                                 \PYGZsh{} files if make cannot find them}
\PYG{n+nv}{JPG\PYGZus{}LIB}\PYG{+w}{ }\PYG{o}{=}\PYG{+w}{ }\PYGZhy{}ljpeg\PYG{+w}{ }\PYGZhy{}lpng\PYG{+w}{ }\PYGZhy{}lz\PYG{+w}{       }\PYG{c+c1}{\PYGZsh{} library names}
\end{sphinxVerbatim}

\sphinxAtStartPar
As with CMake, you do not need to set \sphinxcode{\sphinxupquote{JPG\_INC}} or \sphinxcode{\sphinxupquote{JPG\_PATH}},
if make can find the graphics header and library files in their
default system locations.  You must specify \sphinxcode{\sphinxupquote{JPG\_LIB}} with a
list of graphics libraries to include in the link.  You must make
certain that the ffmpeg executable (or ffmpeg.exe on Windows) is
in a directory where LAMMPS can find it at runtime; that is
usually a directory list in your \sphinxcode{\sphinxupquote{PATH}} environment variable.
\end{sphinxadmonition}
\end{quote}

\sphinxAtStartPar
Using \sphinxcode{\sphinxupquote{ffmpeg}} to output movie files requires that your machine
supports the “popen” function in the standard runtime library.

\begin{sphinxadmonition}{note}{Note:}
\sphinxAtStartPar
On some clusters with high\sphinxhyphen{}speed networks, using the fork()
library call (required by popen()) can interfere with the fast
communication library and lead to simulations using ffmpeg to hang or
crash.
\end{sphinxadmonition}


\bigskip\hrule\bigskip



\subsection{Read or write compressed files}
\label{\detokenize{Build_settings:read-or-write-compressed-files}}\label{\detokenize{Build_settings:gzip}}
\sphinxAtStartPar
If this option is enabled, large files can be read or written with
compression by \sphinxcode{\sphinxupquote{gzip}} or similar tools by several LAMMPS commands,
including {\hyperref[\detokenize{read_data::doc}]{\sphinxcrossref{\DUrole{doc}{read\_data}}}}, {\hyperref[\detokenize{rerun::doc}]{\sphinxcrossref{\DUrole{doc}{rerun}}}}, and
{\hyperref[\detokenize{dump::doc}]{\sphinxcrossref{\DUrole{doc}{dump}}}}.  Supported compression tools and algorithms are currently
\sphinxcode{\sphinxupquote{gzip}}, \sphinxcode{\sphinxupquote{bzip2}}, \sphinxcode{\sphinxupquote{zstd}}, \sphinxcode{\sphinxupquote{xz}}, \sphinxcode{\sphinxupquote{lz4}}, and \sphinxcode{\sphinxupquote{lzma}} (via xz).
\begin{quote}

\begin{sphinxadmonition}{note}{CMake build}

\begin{sphinxVerbatim}[commandchars=\\\{\}]
\PYGZhy{}D\PYG{+w}{ }\PYG{n+nv}{WITH\PYGZus{}GZIP}\PYG{o}{=}value\PYG{+w}{  }\PYG{c+c1}{\PYGZsh{} yes or no}
\PYG{+w}{                    }\PYG{c+c1}{\PYGZsh{} default is yes if CMake can find the gzip program}
\end{sphinxVerbatim}
\end{sphinxadmonition}

\begin{sphinxadmonition}{note}{Traditional make}

\begin{sphinxVerbatim}[commandchars=\\\{\}]
\PYG{n+nv}{LMP\PYGZus{}INC}\PYG{+w}{ }\PYG{o}{=}\PYG{+w}{ }\PYGZhy{}DLAMMPS\PYGZus{}GZIP\PYG{+w}{   }\PYGZlt{}other\PYG{+w}{ }LMP\PYGZus{}INC\PYG{+w}{ }settings\PYGZgt{}
\end{sphinxVerbatim}
\end{sphinxadmonition}
\end{quote}

\sphinxAtStartPar
This option requires that your operating system fully supports the
“popen()” function in the standard runtime library and that a \sphinxcode{\sphinxupquote{gzip}}
or other executable can be found by LAMMPS in the standard search path
during a run.

\begin{sphinxadmonition}{note}{Note:}
\sphinxAtStartPar
On clusters with high\sphinxhyphen{}speed networks, using the “fork()” library call
(required by “popen()”) can interfere with the fast communication
library and lead to simulations using compressed output or input to
hang or crash. For selected operations, compressed file I/O is also
available using a compression library instead, which is what the
{\hyperref[\detokenize{Packages_details:pkg-compress}]{\sphinxcrossref{\DUrole{std,std-ref}{COMPRESS package}}}} enables.
\end{sphinxadmonition}


\bigskip\hrule\bigskip



\subsection{Support for downloading files}
\label{\detokenize{Build_settings:support-for-downloading-files}}\label{\detokenize{Build_settings:libcurl}}
\sphinxAtStartPar
\DUrole{versionmodified,added}{Added in version 29Aug2024.}

\sphinxAtStartPar
The {\hyperref[\detokenize{geturl::doc}]{\sphinxcrossref{\DUrole{doc}{geturl command}}}} command uses the \sphinxhref{https://curl.se/libcurl/}{the libcurl library} to download files.  This requires that
LAMMPS is compiled accordingly which needs the following settings:
\begin{quote}

\begin{sphinxadmonition}{note}{CMake build}

\begin{sphinxVerbatim}[commandchars=\\\{\}]
\PYGZhy{}D\PYG{+w}{ }\PYG{n+nv}{WITH\PYGZus{}CURL}\PYG{o}{=}value\PYG{+w}{      }\PYG{c+c1}{\PYGZsh{} yes or no}
\PYG{+w}{                        }\PYG{c+c1}{\PYGZsh{} default = yes if CMake finds CURL development files, else no}
\end{sphinxVerbatim}

\sphinxAtStartPar
Usually these settings are all that is needed.  If CMake cannot
find the graphics header, library, executable files, you can set
these variables:

\begin{sphinxVerbatim}[commandchars=\\\{\}]
\PYGZhy{}D\PYG{+w}{ }\PYG{n+nv}{CURL\PYGZus{}INCLUDE\PYGZus{}DIR}\PYG{o}{=}path\PYG{+w}{    }\PYG{c+c1}{\PYGZsh{} path to folder which contains curl.h header file}
\PYGZhy{}D\PYG{+w}{ }\PYG{n+nv}{CURL\PYGZus{}LIBRARY}\PYG{o}{=}path\PYG{+w}{        }\PYG{c+c1}{\PYGZsh{} path to libcurls.a (.so) file}
\end{sphinxVerbatim}
\end{sphinxadmonition}

\begin{sphinxadmonition}{note}{Traditional make}

\begin{sphinxVerbatim}[commandchars=\\\{\}]
\PYG{n+nv}{LMP\PYGZus{}INC}\PYG{+w}{ }\PYG{o}{=}\PYG{+w}{ }\PYGZhy{}DLAMMPS\PYGZus{}CURL\PYG{+w}{  }\PYGZlt{}other\PYG{+w}{ }LMP\PYGZus{}INC\PYG{+w}{ }settings\PYGZgt{}

\PYG{n+nv}{CURL\PYGZus{}INC}\PYG{+w}{ }\PYG{o}{=}\PYG{+w}{ }\PYGZhy{}I/usr/local/include\PYG{+w}{   }\PYG{c+c1}{\PYGZsh{} path to curl folder with curl.h}
\PYG{n+nv}{CURL\PYGZus{}PATH}\PYG{+w}{ }\PYG{o}{=}\PYG{+w}{ }\PYGZhy{}L/usr/lib\PYG{+w}{            }\PYG{c+c1}{\PYGZsh{} paths to libcurl.a(.so) if make cannot find it}
\PYG{n+nv}{CURL\PYGZus{}LIB}\PYG{+w}{ }\PYG{o}{=}\PYG{+w}{ }\PYGZhy{}lcurl\PYG{+w}{                 }\PYG{c+c1}{\PYGZsh{} library names}
\end{sphinxVerbatim}

\sphinxAtStartPar
As with CMake, you do not need to set \sphinxcode{\sphinxupquote{CURL\_INC}} or \sphinxcode{\sphinxupquote{CURL\_PATH}},
if make can find the libcurl header and library files in their
default system locations.  You must specify \sphinxcode{\sphinxupquote{CURL\_LIB}} with a
paths or linker flags to link to libcurl.
\end{sphinxadmonition}
\end{quote}


\bigskip\hrule\bigskip



\subsection{Memory allocation alignment}
\label{\detokenize{Build_settings:memory-allocation-alignment}}\label{\detokenize{Build_settings:align}}
\sphinxAtStartPar
This setting enables the use of the \sphinxcode{\sphinxupquote{posix\_memalign()}} call instead of
\sphinxcode{\sphinxupquote{malloc()}} when LAMMPS allocates large chunks of memory.  Vector
instructions on CPUs may become more efficient, if dynamically allocated
memory is aligned on larger\sphinxhyphen{}than\sphinxhyphen{}default byte boundaries.  On most
current operating systems, the \sphinxcode{\sphinxupquote{malloc()}} implementation returns
pointers that are aligned to 16\sphinxhyphen{}byte boundaries. Using SSE vector
instructions efficiently, however, requires memory blocks being aligned
on 64\sphinxhyphen{}byte boundaries.
\begin{quote}

\begin{sphinxadmonition}{note}{CMake build}

\begin{sphinxVerbatim}[commandchars=\\\{\}]
\PYGZhy{}D\PYG{+w}{ }\PYG{n+nv}{LAMMPS\PYGZus{}MEMALIGN}\PYG{o}{=}value\PYG{+w}{            }\PYG{c+c1}{\PYGZsh{} 0, 8, 16, 32, 64 (default)}
\end{sphinxVerbatim}

\sphinxAtStartPar
Use a \sphinxcode{\sphinxupquote{LAMMPS\_MEMALIGN}} value of 0 to disable using
\sphinxcode{\sphinxupquote{posix\_memalign()}} and revert to using the \sphinxcode{\sphinxupquote{malloc()}} C\sphinxhyphen{}library
function instead.  When compiling LAMMPS for Windows systems,
\sphinxcode{\sphinxupquote{malloc()}} will always be used and this setting is ignored.
\end{sphinxadmonition}

\begin{sphinxadmonition}{note}{Traditional make}

\begin{sphinxVerbatim}[commandchars=\\\{\}]
\PYG{n+nv}{LMP\PYGZus{}INC}\PYG{+w}{ }\PYG{o}{=}\PYG{+w}{ }\PYGZhy{}DLAMMPS\PYGZus{}MEMALIGN\PYG{o}{=}value\PYG{+w}{   }\PYG{c+c1}{\PYGZsh{} 8, 16, 32, 64}
\end{sphinxVerbatim}

\sphinxAtStartPar
Do not set \sphinxcode{\sphinxupquote{\sphinxhyphen{}DLAMMPS\_MEMALIGN}}, if you want to have memory
allocated with the \sphinxcode{\sphinxupquote{malloc()}} function call
instead. \sphinxcode{\sphinxupquote{\sphinxhyphen{}DLAMMPS\_MEMALIGN}} \sphinxstylestrong{cannot} be used on Windows, as
Windows different function calls with different semantics for
allocating aligned memory, that are not compatible with how LAMMPS
manages its dynamical memory.
\end{sphinxadmonition}
\end{quote}


\bigskip\hrule\bigskip



\subsection{Workaround for long long integers}
\label{\detokenize{Build_settings:workaround-for-long-long-integers}}\label{\detokenize{Build_settings:longlong}}
\sphinxAtStartPar
If your system or MPI version does not recognize “long long” data
types, the following setting will be needed.  It converts “long long”
to a “long” data type, which should be the desired 8\sphinxhyphen{}byte integer on
those systems:
\begin{quote}

\begin{sphinxadmonition}{note}{CMake build}

\begin{sphinxVerbatim}[commandchars=\\\{\}]
\PYGZhy{}D\PYG{+w}{ }\PYG{n+nv}{LAMMPS\PYGZus{}LONGLONG\PYGZus{}TO\PYGZus{}LONG}\PYG{o}{=}value\PYG{+w}{     }\PYG{c+c1}{\PYGZsh{} yes or no (default)}
\end{sphinxVerbatim}
\end{sphinxadmonition}

\begin{sphinxadmonition}{note}{Traditional make}

\begin{sphinxVerbatim}[commandchars=\\\{\}]
\PYG{n+nv}{LMP\PYGZus{}INC}\PYG{+w}{ }\PYG{o}{=}\PYG{+w}{ }\PYGZhy{}DLAMMPS\PYGZus{}LONGLONG\PYGZus{}TO\PYGZus{}LONG\PYG{+w}{  }\PYGZlt{}other\PYG{+w}{ }LMP\PYGZus{}INC\PYG{+w}{ }settings\PYGZgt{}
\end{sphinxVerbatim}
\end{sphinxadmonition}
\end{quote}


\bigskip\hrule\bigskip



\subsection{Exception handling when using LAMMPS as a library}
\label{\detokenize{Build_settings:exception-handling-when-using-lammps-as-a-library}}\label{\detokenize{Build_settings:exceptions}}
\sphinxAtStartPar
LAMMPS errors do not kill the calling code, but throw an exception.  In
the C\sphinxhyphen{}library interface, the call stack is unwound and control returns
to the caller, e.g. to Python or a code that is coupled to LAMMPS. The
error status can then be queried.  When using C++ directly, the calling
code has to be set up to \sphinxstyleemphasis{catch} exceptions thrown from within LAMMPS.

\begin{sphinxadmonition}{note}{Note:}
\sphinxAtStartPar
When LAMMPS is running in parallel, it is not always possible to
cleanly recover from an exception since not all parallel ranks may
throw an exception and thus other MPI ranks may get stuck waiting for
messages from the ones with errors.
\end{sphinxadmonition}


\bigskip\hrule\bigskip



\subsection{Trigger selected floating\sphinxhyphen{}point exceptions}
\label{\detokenize{Build_settings:trigger-selected-floating-point-exceptions}}\label{\detokenize{Build_settings:trap-fpe}}
\sphinxAtStartPar
Many kinds of CPUs have the capability to detect when a calculation
results in an invalid math operation, like a division by zero or calling
the square root with a negative argument.  The default behavior on
most operating systems is to continue and have values for \sphinxcode{\sphinxupquote{NaN}} (= not
a number) or \sphinxcode{\sphinxupquote{Inf}} (= infinity).  This allows software to detect and
recover from such conditions.  This behavior can be changed, however,
often through use of compiler flags.  On Linux systems (or more general
on systems using the GNU C library), these so\sphinxhyphen{}called floating\sphinxhyphen{}point traps
can also be selectively enabled through library calls.  LAMMPS supports
that by setting the \sphinxcode{\sphinxupquote{\sphinxhyphen{}DLAMMPS\_TRAP\_FPE}} pre\sphinxhyphen{}processor define.  As it is
done in the \sphinxcode{\sphinxupquote{main()}} function, this applies only to the standalone
executable, not the library.
\begin{quote}

\begin{sphinxadmonition}{note}{CMake build}

\begin{sphinxVerbatim}[commandchars=\\\{\}]
\PYGZhy{}D\PYG{+w}{ }\PYG{n+nv}{CMAKE\PYGZus{}TUNE\PYGZus{}FLAGS}\PYG{o}{=}\PYGZhy{}DLAMMPS\PYGZus{}TRAP\PYGZus{}FPE
\end{sphinxVerbatim}
\end{sphinxadmonition}

\begin{sphinxadmonition}{note}{Traditional make}

\begin{sphinxVerbatim}[commandchars=\\\{\}]
\PYG{n+nv}{LMP\PYGZus{}INC}\PYG{+w}{ }\PYG{o}{=}\PYG{+w}{ }\PYGZhy{}DLAMMPS\PYGZus{}TRAP\PYGZus{}FPE\PYG{+w}{  }\PYGZlt{}other\PYG{+w}{ }LMP\PYGZus{}INC\PYG{+w}{ }settings\PYGZgt{}
\end{sphinxVerbatim}
\end{sphinxadmonition}
\end{quote}

\sphinxAtStartPar
After compilation with this flag set, the LAMMPS executable will stop
and produce a core dump when a division by zero, overflow, illegal math
function argument or other invalid floating point operation is encountered.

\sphinxstepscope


\section{Include packages in build}
\label{\detokenize{Build_package:include-packages-in-build}}\label{\detokenize{Build_package::doc}}
\sphinxAtStartPar
In LAMMPS, a package is a group of files that enable a specific set of
features.  For example, force fields for molecular systems or
rigid\sphinxhyphen{}body constraints are in packages.  In the src directory, each
package is a subdirectory with the package name in capital letters.

\sphinxAtStartPar
An overview of packages is given on the {\hyperref[\detokenize{Packages::doc}]{\sphinxcrossref{\DUrole{doc}{Packages}}}} doc
page.  Brief overviews of each package are on the {\hyperref[\detokenize{Packages_details::doc}]{\sphinxcrossref{\DUrole{doc}{Packages details}}}} page.

\sphinxAtStartPar
When building LAMMPS, you can choose to include or exclude each
package.  Generally, there is no need to include a package if you
never plan to use its features.

\sphinxAtStartPar
If you get a run\sphinxhyphen{}time error that a LAMMPS command or style is
“unknown”, it is often because the command is contained in a package,
and your build did not include that package.  If the command or style
\sphinxstyleemphasis{is} available in a package included in the LAMMPS distribution,
the error message will indicate which package would be needed.
Running LAMMPS with the {\hyperref[\detokenize{Run_options::doc}]{\sphinxcrossref{\DUrole{doc}{\sphinxhyphen{}h command\sphinxhyphen{}line switch}}}}
will print \sphinxstyleemphasis{all} optional commands and packages that were enabled
when building that executable.

\sphinxAtStartPar
For the majority of packages, if you follow the single step below to
include it, you can then build LAMMPS exactly the same as you would
without any packages installed.  A few packages may require additional
steps, as explained on the {\hyperref[\detokenize{Build_extras::doc}]{\sphinxcrossref{\DUrole{doc}{Build extras}}}} page.

\sphinxAtStartPar
These links take you to the extra instructions for those select
packages:


\begin{savenotes}\sphinxattablestart
\sphinxthistablewithglobalstyle
\centering
\begin{tabulary}{\linewidth}[t]{TTTTTT}
\sphinxtoprule
\sphinxtableatstartofbodyhook
\sphinxAtStartPar
{\hyperref[\detokenize{Build_extras:adios}]{\sphinxcrossref{\DUrole{std,std-ref}{ADIOS}}}}
&
\sphinxAtStartPar
{\hyperref[\detokenize{Build_extras:atc}]{\sphinxcrossref{\DUrole{std,std-ref}{ATC}}}}
&
\sphinxAtStartPar
{\hyperref[\detokenize{Build_extras:awpmd}]{\sphinxcrossref{\DUrole{std,std-ref}{AWPMD}}}}
&
\sphinxAtStartPar
{\hyperref[\detokenize{Build_extras:colvar}]{\sphinxcrossref{\DUrole{std,std-ref}{COLVARS}}}}
&
\sphinxAtStartPar
{\hyperref[\detokenize{Build_extras:compress}]{\sphinxcrossref{\DUrole{std,std-ref}{COMPRESS}}}}
&
\sphinxAtStartPar
{\hyperref[\detokenize{Build_extras:electrode}]{\sphinxcrossref{\DUrole{std,std-ref}{ELECTRODE}}}}
\\
\sphinxhline
\sphinxAtStartPar
{\hyperref[\detokenize{Build_extras:gpu}]{\sphinxcrossref{\DUrole{std,std-ref}{GPU}}}}
&
\sphinxAtStartPar
{\hyperref[\detokenize{Build_extras:h5md}]{\sphinxcrossref{\DUrole{std,std-ref}{H5MD}}}}
&
\sphinxAtStartPar
{\hyperref[\detokenize{Build_extras:intel}]{\sphinxcrossref{\DUrole{std,std-ref}{INTEL}}}}
&
\sphinxAtStartPar
{\hyperref[\detokenize{Build_extras:kim}]{\sphinxcrossref{\DUrole{std,std-ref}{KIM}}}}
&
\sphinxAtStartPar
{\hyperref[\detokenize{Build_extras:kokkos}]{\sphinxcrossref{\DUrole{std,std-ref}{KOKKOS}}}}
&
\sphinxAtStartPar
{\hyperref[\detokenize{Build_extras:lepton}]{\sphinxcrossref{\DUrole{std,std-ref}{LEPTON}}}}
\\
\sphinxhline
\sphinxAtStartPar
{\hyperref[\detokenize{Build_extras:machdyn}]{\sphinxcrossref{\DUrole{std,std-ref}{MACHDYN}}}}
&
\sphinxAtStartPar
{\hyperref[\detokenize{Build_extras:mdi}]{\sphinxcrossref{\DUrole{std,std-ref}{MDI}}}}
&
\sphinxAtStartPar
{\hyperref[\detokenize{Build_extras:ml-hdnnp}]{\sphinxcrossref{\DUrole{std,std-ref}{ML\sphinxhyphen{}HDNNP}}}}
&
\sphinxAtStartPar
{\hyperref[\detokenize{Build_extras:mliap}]{\sphinxcrossref{\DUrole{std,std-ref}{ML\sphinxhyphen{}IAP}}}}
&
\sphinxAtStartPar
{\hyperref[\detokenize{Build_extras:ml-pace}]{\sphinxcrossref{\DUrole{std,std-ref}{ML\sphinxhyphen{}PACE}}}}
&
\sphinxAtStartPar
{\hyperref[\detokenize{Build_extras:ml-pod}]{\sphinxcrossref{\DUrole{std,std-ref}{ML\sphinxhyphen{}POD}}}}
\\
\sphinxhline
\sphinxAtStartPar
{\hyperref[\detokenize{Build_extras:ml-quip}]{\sphinxcrossref{\DUrole{std,std-ref}{ML\sphinxhyphen{}QUIP}}}}
&
\sphinxAtStartPar
{\hyperref[\detokenize{Build_extras:molfile}]{\sphinxcrossref{\DUrole{std,std-ref}{MOLFILE}}}}
&
\sphinxAtStartPar
{\hyperref[\detokenize{Build_extras:netcdf}]{\sphinxcrossref{\DUrole{std,std-ref}{NETCDF}}}}
&
\sphinxAtStartPar
{\hyperref[\detokenize{Build_extras:openmp}]{\sphinxcrossref{\DUrole{std,std-ref}{OPENMP}}}}
&
\sphinxAtStartPar
{\hyperref[\detokenize{Build_extras:opt}]{\sphinxcrossref{\DUrole{std,std-ref}{OPT}}}}
&
\sphinxAtStartPar
{\hyperref[\detokenize{Build_extras:plumed}]{\sphinxcrossref{\DUrole{std,std-ref}{PLUMED}}}}
\\
\sphinxhline
\sphinxAtStartPar
{\hyperref[\detokenize{Build_extras:poems}]{\sphinxcrossref{\DUrole{std,std-ref}{POEMS}}}}
&
\sphinxAtStartPar
{\hyperref[\detokenize{Build_extras:python}]{\sphinxcrossref{\DUrole{std,std-ref}{PYTHON}}}}
&
\sphinxAtStartPar
{\hyperref[\detokenize{Build_extras:qmmm}]{\sphinxcrossref{\DUrole{std,std-ref}{QMMM}}}}
&
\sphinxAtStartPar
{\hyperref[\detokenize{Build_extras:rheo}]{\sphinxcrossref{\DUrole{std,std-ref}{RHEO}}}}
&
\sphinxAtStartPar
{\hyperref[\detokenize{Build_extras:scafacos}]{\sphinxcrossref{\DUrole{std,std-ref}{SCAFACOS}}}}
&
\sphinxAtStartPar
{\hyperref[\detokenize{Build_extras:voronoi}]{\sphinxcrossref{\DUrole{std,std-ref}{VORONOI}}}}
\\
\sphinxhline
\sphinxAtStartPar
{\hyperref[\detokenize{Build_extras:vtk}]{\sphinxcrossref{\DUrole{std,std-ref}{VTK}}}}
&&&&&\\
\sphinxbottomrule
\end{tabulary}
\sphinxtableafterendhook\par
\sphinxattableend\end{savenotes}

\sphinxAtStartPar
The mechanism for including packages is simple but different for CMake
versus make.
\begin{quote}

\begin{sphinxadmonition}{note}{CMake build}

\begin{sphinxVerbatim}[commandchars=\\\{\}]
\PYGZhy{}D \PYG{n+nv}{PKG\PYGZus{}NAME}\PYG{o}{=}value          \PYG{c}{\PYGZsh{} yes or no (default)}
\end{sphinxVerbatim}

\sphinxAtStartPar
Examples:

\begin{sphinxVerbatim}[commandchars=\\\{\}]
\PYGZhy{}D \PYG{n+nv}{PKG\PYGZus{}MANYBODY}\PYG{o}{=}yes
\PYGZhy{}D \PYG{n+nv}{PKG\PYGZus{}INTEL}\PYG{o}{=}yes
\end{sphinxVerbatim}

\sphinxAtStartPar
All packages are included the same way.  See the shortcut section
below for how to install many packages at once with CMake.

\begin{sphinxadmonition}{note}{Note:}
\sphinxAtStartPar
If you switch between building with CMake and make builds, no
packages in the src directory can be installed when you invoke
\sphinxcode{\sphinxupquote{cmake}}.  CMake will give an error if that is not the case,
indicating how you can uninstall all packages in the src dir.
\end{sphinxadmonition}
\end{sphinxadmonition}

\begin{sphinxadmonition}{note}{Traditional make}

\begin{sphinxVerbatim}[commandchars=\\\{\}]
\PYG{n+nb}{cd}\PYG{+w}{ }lammps/src
make\PYG{+w}{ }ps\PYG{+w}{        }\PYG{c+c1}{\PYGZsh{} check which packages are currently installed}
make\PYG{+w}{ }yes\PYGZhy{}name\PYG{+w}{  }\PYG{c+c1}{\PYGZsh{} install a package with name}
make\PYG{+w}{ }no\PYGZhy{}name\PYG{+w}{   }\PYG{c+c1}{\PYGZsh{} uninstall a package with name}
make\PYG{+w}{ }mpi\PYG{+w}{       }\PYG{c+c1}{\PYGZsh{} build LAMMPS with whatever packages are now installed}
\end{sphinxVerbatim}

\sphinxAtStartPar
Examples:

\begin{sphinxVerbatim}[commandchars=\\\{\}]
make\PYG{+w}{ }no\PYGZhy{}rigid
make\PYG{+w}{ }yes\PYGZhy{}intel
\end{sphinxVerbatim}

\sphinxAtStartPar
All packages are included the same way.  See the shortcut section
below for how to install many packages at once with make.

\begin{sphinxadmonition}{note}{Note:}
\sphinxAtStartPar
You must always re\sphinxhyphen{}build LAMMPS (via make) after installing or
uninstalling a package, for the action to take effect. The
included dependency tracking will make certain only files that
are required to be rebuilt are recompiled.
\end{sphinxadmonition}

\begin{sphinxadmonition}{note}{Note:}
\sphinxAtStartPar
You cannot install or uninstall packages and build LAMMPS in a
single make command with multiple targets, e.g. \sphinxcode{\sphinxupquote{make
yes\sphinxhyphen{}colloid mpi}}.  This is because the make procedure creates
a list of source files that will be out\sphinxhyphen{}of\sphinxhyphen{}date for the build
if the package configuration changes within the same command.
You can include or exclude multiple packages in a single make
command, e.g. \sphinxcode{\sphinxupquote{make yes\sphinxhyphen{}colloid no\sphinxhyphen{}manybody}}.
\end{sphinxadmonition}
\end{sphinxadmonition}
\end{quote}


\subsection{Information for both build systems}
\label{\detokenize{Build_package:information-for-both-build-systems}}
\sphinxAtStartPar
Almost all packages can be included or excluded in a LAMMPS build,
independent of the other packages.  However, some packages include files
derived from files in other packages.  LAMMPS checks for this and does
the right thing.  Individual files are only included if their
dependencies are already included.  Likewise, if a package is excluded,
other files dependent on that package are also excluded.

\begin{sphinxadmonition}{note}{Note:}
\sphinxAtStartPar
By default no packages are installed.  Prior to August 2018, however,
if you downloaded a tarball, 3 packages (KSPACE, MANYBODY, MOLECULE)
were pre\sphinxhyphen{}installed via the traditional make procedure in the \sphinxcode{\sphinxupquote{src}}
directory.  That is no longer the case, so that CMake will build
as\sphinxhyphen{}is without needing to uninstall those packages.
\end{sphinxadmonition}


\bigskip\hrule\bigskip



\subsection{CMake presets for installing many packages}
\label{\detokenize{Build_package:cmake-presets-for-installing-many-packages}}\label{\detokenize{Build_package:cmake-presets}}
\sphinxAtStartPar
Instead of specifying all the CMake options via the command\sphinxhyphen{}line,
CMake allows initializing its settings cache using script files.
These are regular CMake files which can manipulate and set CMake
variables (which represent selected options), and can also contain
control flow constructs for more complex operations.

\sphinxAtStartPar
LAMMPS includes several of these files to define configuration
“presets”, similar to the options that exist for the Make based
system. Using these files, you can enable/disable portions of the
available packages in LAMMPS. If you need a custom preset, you can
make a copy of one of them and modify it to suit your needs.

\begin{sphinxVerbatim}[commandchars=\\\{\}]
\PYG{c+c1}{\PYGZsh{} enable just a few core packages}
cmake\PYG{+w}{ }\PYGZhy{}C\PYG{+w}{ }../cmake/presets/basic.cmake\PYG{+w}{    }\PYG{o}{[}OPTIONS\PYG{o}{]}\PYG{+w}{ }../cmake

\PYG{c+c1}{\PYGZsh{} enable most packages}
cmake\PYG{+w}{ }\PYGZhy{}C\PYG{+w}{ }../cmake/presets/most.cmake\PYG{+w}{     }\PYG{o}{[}OPTIONS\PYG{o}{]}\PYG{+w}{ }../cmake

\PYG{c+c1}{\PYGZsh{} enable packages which download sources or potential files}
cmake\PYG{+w}{ }\PYGZhy{}C\PYG{+w}{ }../cmake/presets/download.cmake\PYG{+w}{ }\PYG{o}{[}OPTIONS\PYG{o}{]}\PYG{+w}{ }../cmake

\PYG{c+c1}{\PYGZsh{} disable packages that do require extra libraries or tools}
cmake\PYG{+w}{ }\PYGZhy{}C\PYG{+w}{ }../cmake/presets/nolib.cmake\PYG{+w}{    }\PYG{o}{[}OPTIONS\PYG{o}{]}\PYG{+w}{ }../cmake

\PYG{c+c1}{\PYGZsh{} change settings to use the Clang compilers by default}
cmake\PYG{+w}{ }\PYGZhy{}C\PYG{+w}{ }../cmake/presets/clang.cmake\PYG{+w}{    }\PYG{o}{[}OPTIONS\PYG{o}{]}\PYG{+w}{ }../cmake

\PYG{c+c1}{\PYGZsh{} change settings to use the GNU compilers by default}
cmake\PYG{+w}{ }\PYGZhy{}C\PYG{+w}{ }../cmake/presets/gcc.cmake\PYG{+w}{      }\PYG{o}{[}OPTIONS\PYG{o}{]}\PYG{+w}{ }../cmake

\PYG{c+c1}{\PYGZsh{} change settings to use the Intel compilers by default}
cmake\PYG{+w}{ }\PYGZhy{}C\PYG{+w}{ }../cmake/presets/intel.cmake\PYG{+w}{    }\PYG{o}{[}OPTIONS\PYG{o}{]}\PYG{+w}{ }../cmake

\PYG{c+c1}{\PYGZsh{} change settings to use the PGI compilers by default}
cmake\PYG{+w}{ }\PYGZhy{}C\PYG{+w}{ }../cmake/presets/pgi.cmake\PYG{+w}{      }\PYG{o}{[}OPTIONS\PYG{o}{]}\PYG{+w}{ }../cmake

\PYG{c+c1}{\PYGZsh{} enable all packages}
cmake\PYG{+w}{ }\PYGZhy{}C\PYG{+w}{ }../cmake/presets/all\PYGZus{}on.cmake\PYG{+w}{   }\PYG{o}{[}OPTIONS\PYG{o}{]}\PYG{+w}{ }../cmake

\PYG{c+c1}{\PYGZsh{} disable all packages}
cmake\PYG{+w}{ }\PYGZhy{}C\PYG{+w}{ }../cmake/presets/all\PYGZus{}off.cmake\PYG{+w}{  }\PYG{o}{[}OPTIONS\PYG{o}{]}\PYG{+w}{ }../cmake

\PYG{c+c1}{\PYGZsh{}  compile with MinGW cross\PYGZhy{}compilers}
mingw64\PYGZhy{}cmake\PYG{+w}{ }\PYGZhy{}C\PYG{+w}{ }../cmake/presets/mingw\PYGZhy{}cross.cmake\PYG{+w}{ }\PYG{o}{[}OPTIONS\PYG{o}{]}\PYG{+w}{ }../cmake

\PYG{c+c1}{\PYGZsh{} compile serial multi\PYGZhy{}arch binaries on macOS}
cmake\PYG{+w}{ }\PYGZhy{}C\PYG{+w}{ }../cmake/presets/macos\PYGZhy{}multiarch.cmake\PYG{+w}{ }\PYG{o}{[}OPTIONS\PYG{o}{]}\PYG{+w}{ }../cmake
\end{sphinxVerbatim}

\sphinxAtStartPar
Presets that have names starting with “windows” are specifically for
compiling LAMMPS {\hyperref[\detokenize{Build_windows::doc}]{\sphinxcrossref{\DUrole{doc}{natively on Windows}}}} and
presets that have names starting with “kokkos” are specifically for
selecting configurations for compiling LAMMPS with {\hyperref[\detokenize{Build_extras:kokkos}]{\sphinxcrossref{\DUrole{std,std-ref}{KOKKOS}}}}.

\begin{sphinxadmonition}{note}{Note:}
\sphinxAtStartPar
Running cmake this way manipulates the CMake settings cache in your
current build directory.  You can combine multiple presets and options
in a single cmake run, or change settings incrementally by running
cmake with new flags.  If you use a present for selecting a set of
compilers, it will reset all settings from previous CMake runs.
\end{sphinxadmonition}


\subsubsection{Example}
\label{\detokenize{Build_package:example}}
\begin{sphinxVerbatim}[commandchars=\\\{\}]
\PYG{c+c1}{\PYGZsh{} build LAMMPS with most commonly used packages, but then remove}
\PYG{c+c1}{\PYGZsh{} those requiring additional library or tools, but still enable}
\PYG{c+c1}{\PYGZsh{} GPU package and configure it for using CUDA. You can run.}
mkdir\PYG{+w}{ }build
\PYG{n+nb}{cd}\PYG{+w}{ }build
cmake\PYG{+w}{ }\PYGZhy{}C\PYG{+w}{ }../cmake/presets/most.cmake\PYG{+w}{ }\PYGZhy{}C\PYG{+w}{ }../cmake/presets/nolib.cmake\PYG{+w}{ }\PYG{l+s+se}{\PYGZbs{}}
\PYG{+w}{      }\PYGZhy{}D\PYG{+w}{ }\PYG{n+nv}{PKG\PYGZus{}GPU}\PYG{o}{=}on\PYG{+w}{ }\PYGZhy{}D\PYG{+w}{ }\PYG{n+nv}{GPU\PYGZus{}API}\PYG{o}{=}cuda\PYG{+w}{ }../cmake

\PYG{c+c1}{\PYGZsh{} to add another package, say BODY to the previous configuration you can run:}
cmake\PYG{+w}{ }\PYGZhy{}D\PYG{+w}{ }\PYG{n+nv}{PKG\PYGZus{}BODY}\PYG{o}{=}on\PYG{+w}{ }.

\PYG{c+c1}{\PYGZsh{} to reset the package selection from above to the default of no packages}
\PYG{c+c1}{\PYGZsh{} but leaving all other settings untouched. You can run:}
cmake\PYG{+w}{ }\PYGZhy{}C\PYG{+w}{ }../cmake/presets/all\PYGZus{}off.cmake\PYG{+w}{ }.
\end{sphinxVerbatim}


\bigskip\hrule\bigskip



\subsection{Make shortcuts for installing many packages}
\label{\detokenize{Build_package:make-shortcuts-for-installing-many-packages}}
\sphinxAtStartPar
The following commands are useful for managing package source files
and their installation when building LAMMPS via traditional make.
Just type \sphinxcode{\sphinxupquote{make}} in lammps/src to see a one\sphinxhyphen{}line summary.

\sphinxAtStartPar
These commands install/uninstall sets of packages:

\begin{sphinxVerbatim}[commandchars=\\\{\}]
make\PYG{+w}{ }yes\PYGZhy{}all\PYG{+w}{                        }\PYG{c+c1}{\PYGZsh{} install all packages}
make\PYG{+w}{ }no\PYGZhy{}all\PYG{+w}{                         }\PYG{c+c1}{\PYGZsh{} check for changes and uninstall all packages}
make\PYG{+w}{ }no\PYGZhy{}installed\PYG{+w}{                   }\PYG{c+c1}{\PYGZsh{} only check and uninstall installed packages}
make\PYG{+w}{ }yes\PYGZhy{}basic\PYG{+w}{                      }\PYG{c+c1}{\PYGZsh{} install a few commonly used packages\PYGZsq{}}
make\PYG{+w}{ }no\PYGZhy{}basic\PYG{+w}{                       }\PYG{c+c1}{\PYGZsh{} remove a few commonly used packages\PYGZsq{}}
make\PYG{+w}{ }yes\PYGZhy{}most\PYG{+w}{                       }\PYG{c+c1}{\PYGZsh{} install most packages w/o libs\PYGZsq{}}
make\PYG{+w}{ }no\PYGZhy{}most\PYG{+w}{                        }\PYG{c+c1}{\PYGZsh{} remove most packages w/o libs\PYGZsq{}}
make\PYG{+w}{ }yes\PYGZhy{}lib\PYG{+w}{                        }\PYG{c+c1}{\PYGZsh{} install packages that require extra libraries}
make\PYG{+w}{ }no\PYGZhy{}lib\PYG{+w}{                         }\PYG{c+c1}{\PYGZsh{} uninstall packages that require extra libraries}
make\PYG{+w}{ }yes\PYGZhy{}ext\PYG{+w}{                        }\PYG{c+c1}{\PYGZsh{} install packages that require external libraries}
make\PYG{+w}{ }no\PYGZhy{}ext\PYG{+w}{                         }\PYG{c+c1}{\PYGZsh{} uninstall packages that require external libraries}
\end{sphinxVerbatim}

\sphinxAtStartPar
which install/uninstall various sets of packages.  Typing \sphinxcode{\sphinxupquote{make
package}} will list all the these commands.

\begin{sphinxadmonition}{note}{Note:}
\sphinxAtStartPar
Installing or uninstalling a package for the make based build process
works by simply copying files back and forth between the main source
directory src and the subdirectories with the package name (e.g.
src/KSPACE, src/ATC), so that the files are included or excluded
when LAMMPS is built.  Only source files in the src folder will be
compiled.
\end{sphinxadmonition}

\sphinxAtStartPar
The following make commands help manage files that exist in both the
src directory and in package subdirectories.  You do not normally
need to use these commands unless you are editing LAMMPS files or are
updating LAMMPS via git.

\sphinxAtStartPar
Type \sphinxcode{\sphinxupquote{make package\sphinxhyphen{}status}} or \sphinxcode{\sphinxupquote{make ps}} to show which packages are
currently installed.  For those that are installed, it will list any
files that are different in the src directory and package
subdirectory.

\sphinxAtStartPar
Type \sphinxcode{\sphinxupquote{make package\sphinxhyphen{}installed}} or \sphinxcode{\sphinxupquote{make pi}} to show which packages are
currently installed, without listing the status of packages that are
not installed.

\sphinxAtStartPar
Type \sphinxcode{\sphinxupquote{make package\sphinxhyphen{}update}} or \sphinxcode{\sphinxupquote{make pu}} to overwrite src files with
files from the package subdirectories if the package is installed.  It
should be used after the checkout has been {\hyperref[\detokenize{Install_git::doc}]{\sphinxcrossref{\DUrole{doc}{updated or changed
with git}}}}, this will only update the files in the package
subdirectories, but not the copies in the src folder.

\sphinxAtStartPar
Type \sphinxcode{\sphinxupquote{make package\sphinxhyphen{}overwrite}} to overwrite files in the package
subdirectories with src files.

\sphinxAtStartPar
Type \sphinxcode{\sphinxupquote{make package\sphinxhyphen{}diff}} to list all differences between pairs of
files in both the source directory and the package directory.

\sphinxstepscope


\section{Packages with extra build options}
\label{\detokenize{Build_extras:packages-with-extra-build-options}}\label{\detokenize{Build_extras::doc}}
\sphinxAtStartPar
When building with some packages, additional steps may be required,
in addition to


\begin{savenotes}\sphinxattablestart
\sphinxthistablewithglobalstyle
\centering
\begin{tabular}[t]{*{2}{\X{1}{2}}}
\sphinxtoprule
\sphinxstyletheadfamily 
\sphinxAtStartPar
CMake build
&\sphinxstyletheadfamily 
\sphinxAtStartPar
Traditional make
\\
\sphinxmidrule
\sphinxtableatstartofbodyhook
\begin{sphinxVerbatimintable}[commandchars=\\\{\}]
cmake\PYG{+w}{ }\PYGZhy{}D\PYG{+w}{ }\PYG{n+nv}{PKG\PYGZus{}NAME}\PYG{o}{=}yes
\end{sphinxVerbatimintable}
&
\begin{sphinxVerbatimintable}[commandchars=\\\{\}]
make\PYG{+w}{ }yes\PYGZhy{}name
\end{sphinxVerbatimintable}
\\
\sphinxbottomrule
\end{tabular}
\sphinxtableafterendhook\par
\sphinxattableend\end{savenotes}

\sphinxAtStartPar
as described on the {\hyperref[\detokenize{Build_package::doc}]{\sphinxcrossref{\DUrole{doc}{Build\_package}}}} page.

\sphinxAtStartPar
For a CMake build there may be additional optional or required
variables to set.  For a build with make, a provided library under the
lammps/lib directory may need to be built first.  Or an external
library may need to exist on your system or be downloaded and built.
You may need to tell LAMMPS where it is found on your system.

\sphinxAtStartPar
This is the list of packages that may require additional steps.


\begin{savenotes}\sphinxattablestart
\sphinxthistablewithglobalstyle
\centering
\begin{tabulary}{\linewidth}[t]{TTTTTT}
\sphinxtoprule
\sphinxtableatstartofbodyhook
\sphinxAtStartPar
{\hyperref[\detokenize{Build_extras:adios}]{\sphinxcrossref{\DUrole{std,std-ref}{ADIOS}}}}
&
\sphinxAtStartPar
{\hyperref[\detokenize{Build_extras:atc}]{\sphinxcrossref{\DUrole{std,std-ref}{ATC}}}}
&
\sphinxAtStartPar
{\hyperref[\detokenize{Build_extras:awpmd}]{\sphinxcrossref{\DUrole{std,std-ref}{AWPMD}}}}
&
\sphinxAtStartPar
{\hyperref[\detokenize{Build_extras:colvar}]{\sphinxcrossref{\DUrole{std,std-ref}{COLVARS}}}}
&
\sphinxAtStartPar
{\hyperref[\detokenize{Build_extras:compress}]{\sphinxcrossref{\DUrole{std,std-ref}{COMPRESS}}}}
&
\sphinxAtStartPar
{\hyperref[\detokenize{Build_extras:electrode}]{\sphinxcrossref{\DUrole{std,std-ref}{ELECTRODE}}}}
\\
\sphinxhline
\sphinxAtStartPar
{\hyperref[\detokenize{Build_extras:gpu}]{\sphinxcrossref{\DUrole{std,std-ref}{GPU}}}}
&
\sphinxAtStartPar
{\hyperref[\detokenize{Build_extras:h5md}]{\sphinxcrossref{\DUrole{std,std-ref}{H5MD}}}}
&
\sphinxAtStartPar
{\hyperref[\detokenize{Build_extras:intel}]{\sphinxcrossref{\DUrole{std,std-ref}{INTEL}}}}
&
\sphinxAtStartPar
{\hyperref[\detokenize{Build_extras:kim}]{\sphinxcrossref{\DUrole{std,std-ref}{KIM}}}}
&
\sphinxAtStartPar
{\hyperref[\detokenize{Build_extras:kokkos}]{\sphinxcrossref{\DUrole{std,std-ref}{KOKKOS}}}}
&
\sphinxAtStartPar
{\hyperref[\detokenize{Build_extras:lepton}]{\sphinxcrossref{\DUrole{std,std-ref}{LEPTON}}}}
\\
\sphinxhline
\sphinxAtStartPar
{\hyperref[\detokenize{Build_extras:machdyn}]{\sphinxcrossref{\DUrole{std,std-ref}{MACHDYN}}}}
&
\sphinxAtStartPar
{\hyperref[\detokenize{Build_extras:mdi}]{\sphinxcrossref{\DUrole{std,std-ref}{MDI}}}}
&
\sphinxAtStartPar
{\hyperref[\detokenize{Build_extras:ml-hdnnp}]{\sphinxcrossref{\DUrole{std,std-ref}{ML\sphinxhyphen{}HDNNP}}}}
&
\sphinxAtStartPar
{\hyperref[\detokenize{Build_extras:mliap}]{\sphinxcrossref{\DUrole{std,std-ref}{ML\sphinxhyphen{}IAP}}}}
&
\sphinxAtStartPar
{\hyperref[\detokenize{Build_extras:ml-pace}]{\sphinxcrossref{\DUrole{std,std-ref}{ML\sphinxhyphen{}PACE}}}}
&
\sphinxAtStartPar
{\hyperref[\detokenize{Build_extras:ml-pod}]{\sphinxcrossref{\DUrole{std,std-ref}{ML\sphinxhyphen{}POD}}}}
\\
\sphinxhline
\sphinxAtStartPar
{\hyperref[\detokenize{Build_extras:ml-quip}]{\sphinxcrossref{\DUrole{std,std-ref}{ML\sphinxhyphen{}QUIP}}}}
&
\sphinxAtStartPar
{\hyperref[\detokenize{Build_extras:molfile}]{\sphinxcrossref{\DUrole{std,std-ref}{MOLFILE}}}}
&
\sphinxAtStartPar
{\hyperref[\detokenize{Build_extras:netcdf}]{\sphinxcrossref{\DUrole{std,std-ref}{NETCDF}}}}
&
\sphinxAtStartPar
{\hyperref[\detokenize{Build_extras:openmp}]{\sphinxcrossref{\DUrole{std,std-ref}{OPENMP}}}}
&
\sphinxAtStartPar
{\hyperref[\detokenize{Build_extras:opt}]{\sphinxcrossref{\DUrole{std,std-ref}{OPT}}}}
&
\sphinxAtStartPar
{\hyperref[\detokenize{Build_extras:plumed}]{\sphinxcrossref{\DUrole{std,std-ref}{PLUMED}}}}
\\
\sphinxhline
\sphinxAtStartPar
{\hyperref[\detokenize{Build_extras:poems}]{\sphinxcrossref{\DUrole{std,std-ref}{POEMS}}}}
&
\sphinxAtStartPar
{\hyperref[\detokenize{Build_extras:python}]{\sphinxcrossref{\DUrole{std,std-ref}{PYTHON}}}}
&
\sphinxAtStartPar
{\hyperref[\detokenize{Build_extras:qmmm}]{\sphinxcrossref{\DUrole{std,std-ref}{QMMM}}}}
&
\sphinxAtStartPar
{\hyperref[\detokenize{Build_extras:rheo}]{\sphinxcrossref{\DUrole{std,std-ref}{RHEO}}}}
&
\sphinxAtStartPar
{\hyperref[\detokenize{Build_extras:scafacos}]{\sphinxcrossref{\DUrole{std,std-ref}{SCAFACOS}}}}
&
\sphinxAtStartPar
{\hyperref[\detokenize{Build_extras:voronoi}]{\sphinxcrossref{\DUrole{std,std-ref}{VORONOI}}}}
\\
\sphinxhline
\sphinxAtStartPar
{\hyperref[\detokenize{Build_extras:vtk}]{\sphinxcrossref{\DUrole{std,std-ref}{VTK}}}}
&&&&&\\
\sphinxbottomrule
\end{tabulary}
\sphinxtableafterendhook\par
\sphinxattableend\end{savenotes}


\bigskip\hrule\bigskip



\subsection{COMPRESS package}
\label{\detokenize{Build_extras:compress-package}}\label{\detokenize{Build_extras:compress}}
\sphinxAtStartPar
To build with this package you must have the \sphinxhref{https://zlib.net}{zlib compression library} available on your system to build dump styles with
a \sphinxcode{\sphinxupquote{/gz}} suffix.  There are also styles using the
\sphinxhref{https://facebook.github.io/zstd/}{Zstandard} library which have a
‘/zstd’ suffix.  The zstd library version must be at least 1.4.  Older
versions use an incompatible API and thus LAMMPS will fail to compile.
\begin{quote}

\begin{sphinxadmonition}{note}{CMake build}

\sphinxAtStartPar
If CMake cannot find the zlib library or include files, you can set
these variables:

\begin{sphinxVerbatim}[commandchars=\\\{\}]
\PYGZhy{}D\PYG{+w}{ }\PYG{n+nv}{ZLIB\PYGZus{}INCLUDE\PYGZus{}DIR}\PYG{o}{=}path\PYG{+w}{    }\PYG{c+c1}{\PYGZsh{} path to zlib.h header file}
\PYGZhy{}D\PYG{+w}{ }\PYG{n+nv}{ZLIB\PYGZus{}LIBRARY}\PYG{o}{=}path\PYG{+w}{        }\PYG{c+c1}{\PYGZsh{} path to libz.a (.so) file}
\end{sphinxVerbatim}

\sphinxAtStartPar
Support for Zstandard compression is auto\sphinxhyphen{}detected and for that
CMake depends on the \sphinxhref{https://www.freedesktop.org/wiki/Software/pkg-config/}{pkg\sphinxhyphen{}config} tool to
identify the necessary flags to compile with this library, so the
corresponding \sphinxcode{\sphinxupquote{libzstandard.pc}} file must be in a folder where
\sphinxcode{\sphinxupquote{pkg\sphinxhyphen{}config}} can find it, which may require adding it to the
\sphinxcode{\sphinxupquote{PKG\_CONFIG\_PATH}} environment variable.
\end{sphinxadmonition}

\begin{sphinxadmonition}{note}{Traditional make}

\sphinxAtStartPar
To include support for Zstandard compression, \sphinxcode{\sphinxupquote{\sphinxhyphen{}DLAMMPS\_ZSTD}}
must be added to the compiler flags.  If make cannot find the
libraries, you can edit the file \sphinxcode{\sphinxupquote{lib/compress/Makefile.lammps}}
to specify the paths and library names.  This must be done
\sphinxstylestrong{before} the package is installed.
\end{sphinxadmonition}
\end{quote}


\bigskip\hrule\bigskip



\subsection{GPU package}
\label{\detokenize{Build_extras:gpu-package}}\label{\detokenize{Build_extras:gpu}}
\sphinxAtStartPar
To build with this package, you must choose options for precision and
which GPU hardware to build for. The GPU package currently supports
three different types of backends: OpenCL, CUDA and HIP.


\subsubsection{CMake build}
\label{\detokenize{Build_extras:cmake-build}}
\begin{sphinxVerbatim}[commandchars=\\\{\}]
\PYGZhy{}D\PYG{+w}{ }\PYG{n+nv}{GPU\PYGZus{}API}\PYG{o}{=}value\PYG{+w}{             }\PYG{c+c1}{\PYGZsh{} value = opencl (default) or cuda or hip}
\PYGZhy{}D\PYG{+w}{ }\PYG{n+nv}{GPU\PYGZus{}PREC}\PYG{o}{=}value\PYG{+w}{            }\PYG{c+c1}{\PYGZsh{} precision setting}
\PYG{+w}{                             }\PYG{c+c1}{\PYGZsh{} value = double or mixed (default) or single}
\PYGZhy{}D\PYG{+w}{ }\PYG{n+nv}{GPU\PYGZus{}ARCH}\PYG{o}{=}value\PYG{+w}{            }\PYG{c+c1}{\PYGZsh{} primary GPU hardware choice for GPU\PYGZus{}API=cuda}
\PYG{+w}{                             }\PYG{c+c1}{\PYGZsh{} value = sm\PYGZus{}XX (see below, default is sm\PYGZus{}50)}
\PYGZhy{}D\PYG{+w}{ }\PYG{n+nv}{GPU\PYGZus{}DEBUG}\PYG{o}{=}value\PYG{+w}{           }\PYG{c+c1}{\PYGZsh{} enable debug code in the GPU package library,}
\PYG{+w}{                             }\PYG{c+c1}{\PYGZsh{} mostly useful for developers}
\PYG{+w}{                             }\PYG{c+c1}{\PYGZsh{} value = yes or no (default)}
\PYGZhy{}D\PYG{+w}{ }\PYG{n+nv}{HIP\PYGZus{}PATH}\PYG{o}{=}value\PYG{+w}{            }\PYG{c+c1}{\PYGZsh{} value = path to HIP installation. Must be set if}
\PYG{+w}{                             }\PYG{c+c1}{\PYGZsh{} GPU\PYGZus{}API=HIP}
\PYGZhy{}D\PYG{+w}{ }\PYG{n+nv}{HIP\PYGZus{}ARCH}\PYG{o}{=}value\PYG{+w}{            }\PYG{c+c1}{\PYGZsh{} primary GPU hardware choice for GPU\PYGZus{}API=hip}
\PYG{+w}{                             }\PYG{c+c1}{\PYGZsh{} value depends on selected HIP\PYGZus{}PLATFORM}
\PYG{+w}{                             }\PYG{c+c1}{\PYGZsh{} default is \PYGZsq{}gfx906\PYGZsq{} for HIP\PYGZus{}PLATFORM=amd and \PYGZsq{}sm\PYGZus{}50\PYGZsq{} for}
\PYG{+w}{                             }\PYG{c+c1}{\PYGZsh{} HIP\PYGZus{}PLATFORM=nvcc}
\PYGZhy{}D\PYG{+w}{ }\PYG{n+nv}{HIP\PYGZus{}USE\PYGZus{}DEVICE\PYGZus{}SORT}\PYG{o}{=}value\PYG{+w}{ }\PYG{c+c1}{\PYGZsh{} enables GPU sorting}
\PYG{+w}{                             }\PYG{c+c1}{\PYGZsh{} value = yes (default) or no}
\PYGZhy{}D\PYG{+w}{ }\PYG{n+nv}{CUDPP\PYGZus{}OPT}\PYG{o}{=}value\PYG{+w}{           }\PYG{c+c1}{\PYGZsh{} use GPU binning with CUDA (should be off for modern GPUs)}
\PYG{+w}{                             }\PYG{c+c1}{\PYGZsh{} enables CUDA Performance Primitives, must be \PYGZdq{}no\PYGZdq{} for}
\PYG{+w}{                             }\PYG{c+c1}{\PYGZsh{} CUDA\PYGZus{}MPS\PYGZus{}SUPPORT=yes}
\PYG{+w}{                             }\PYG{c+c1}{\PYGZsh{} value = yes or no (default)}
\PYGZhy{}D\PYG{+w}{ }\PYG{n+nv}{CUDA\PYGZus{}MPS\PYGZus{}SUPPORT}\PYG{o}{=}value\PYG{+w}{    }\PYG{c+c1}{\PYGZsh{} enables some tweaks required to run with active}
\PYG{+w}{                             }\PYG{c+c1}{\PYGZsh{} nvidia\PYGZhy{}cuda\PYGZhy{}mps daemon}
\PYG{+w}{                             }\PYG{c+c1}{\PYGZsh{} value = yes or no (default)}
\PYGZhy{}D\PYG{+w}{ }\PYG{n+nv}{CUDA\PYGZus{}BUILD\PYGZus{}MULTIARCH}\PYG{o}{=}value\PYG{+w}{  }\PYG{c+c1}{\PYGZsh{} enables building CUDA kernels for all supported GPU}
\PYG{+w}{                               }\PYG{c+c1}{\PYGZsh{} architectures}
\PYG{+w}{                               }\PYG{c+c1}{\PYGZsh{} value = yes (default) or no}
\PYGZhy{}D\PYG{+w}{ }\PYG{n+nv}{USE\PYGZus{}STATIC\PYGZus{}OPENCL\PYGZus{}LOADER}\PYG{o}{=}value\PYG{+w}{  }\PYG{c+c1}{\PYGZsh{} downloads/includes OpenCL ICD loader library,}
\PYG{+w}{                                   }\PYG{c+c1}{\PYGZsh{} no local OpenCL headers/libs needed}
\PYG{+w}{                                   }\PYG{c+c1}{\PYGZsh{} value = yes (default) or no}
\end{sphinxVerbatim}

\sphinxAtStartPar
\sphinxcode{\sphinxupquote{GPU\_ARCH}} settings for different GPU hardware is as follows:
\begin{itemize}
\item {} 
\sphinxAtStartPar
\sphinxcode{\sphinxupquote{sm\_30}} for Kepler (supported since CUDA 5 and until CUDA 10.x)

\item {} 
\sphinxAtStartPar
\sphinxcode{\sphinxupquote{sm\_35}} or \sphinxcode{\sphinxupquote{sm\_37}} for Kepler (supported since CUDA 5 and until CUDA 11.x)

\item {} 
\sphinxAtStartPar
\sphinxcode{\sphinxupquote{sm\_50}} or \sphinxcode{\sphinxupquote{sm\_52}} for Maxwell (supported since CUDA 6)

\item {} 
\sphinxAtStartPar
\sphinxcode{\sphinxupquote{sm\_60}} or \sphinxcode{\sphinxupquote{sm\_61}} for Pascal (supported since CUDA 8)

\item {} 
\sphinxAtStartPar
\sphinxcode{\sphinxupquote{sm\_70}} for Volta (supported since CUDA 9)

\item {} 
\sphinxAtStartPar
\sphinxcode{\sphinxupquote{sm\_75}} for Turing (supported since CUDA 10)

\item {} 
\sphinxAtStartPar
\sphinxcode{\sphinxupquote{sm\_80}} or sm\_86 for Ampere (supported since CUDA 11, sm\_86 since CUDA 11.1)

\item {} 
\sphinxAtStartPar
\sphinxcode{\sphinxupquote{sm\_89}} for Lovelace (supported since CUDA 11.8)

\item {} 
\sphinxAtStartPar
\sphinxcode{\sphinxupquote{sm\_90}} for Hopper (supported since CUDA 12.0)

\end{itemize}

\sphinxAtStartPar
A more detailed list can be found, for example,
at \sphinxhref{https://en.wikipedia.org/wiki/CUDA\#GPUs\_supported}{Wikipedia’s CUDA article}

\sphinxAtStartPar
CMake can detect which version of the CUDA toolkit is used and thus will
try to include support for \sphinxstylestrong{all} major GPU architectures supported by
this toolkit.  Thus the \sphinxcode{\sphinxupquote{GPU\_ARCH}} setting is merely an optimization, to
have code for the preferred GPU architecture directly included rather
than having to wait for the JIT compiler of the CUDA driver to translate
it.  This behavior can be turned off (e.g. to speed up compilation) by
setting \sphinxcode{\sphinxupquote{CUDA\_ENABLE\_MULTIARCH}} to \sphinxcode{\sphinxupquote{no}}.

\sphinxAtStartPar
When compiling for CUDA or HIP with CUDA, version 8.0 or later of the
CUDA toolkit is required and a GPU architecture of Kepler or later,
which must \sphinxstyleemphasis{also} be supported by the CUDA toolkit in use \sphinxstylestrong{and} the
CUDA driver in use.  When compiling for OpenCL, OpenCL version 1.2 or
later is required and the GPU must be supported by the GPU driver and
OpenCL runtime bundled with the driver.

\sphinxAtStartPar
When building with CMake, you \sphinxstylestrong{must NOT} build the GPU library in
\sphinxcode{\sphinxupquote{lib/gpu}} using the traditional build procedure. CMake will detect
files generated by that process and will terminate with an error and a
suggestion for how to remove them.

\sphinxAtStartPar
If you are compiling for OpenCL, the default setting is to download,
build, and link with a static OpenCL ICD loader library and standard
OpenCL headers.  This way no local OpenCL development headers or library
needs to be present and only OpenCL compatible drivers need to be
installed to use OpenCL.  If this is not desired, you can set
\sphinxcode{\sphinxupquote{USE\_STATIC\_OPENCL\_LOADER}} to \sphinxcode{\sphinxupquote{no}}.

\sphinxAtStartPar
The GPU library has some multi\sphinxhyphen{}thread support using OpenMP.  If LAMMPS
is built with \sphinxcode{\sphinxupquote{\sphinxhyphen{}D BUILD\_OMP=on}} this will also be enabled.

\sphinxAtStartPar
If you are compiling with HIP, note that before running CMake you will
have to set appropriate environment variables. Some variables such as
\sphinxcode{\sphinxupquote{HCC\_AMDGPU\_TARGET}} (for ROCm \textless{}= 4.0) or \sphinxcode{\sphinxupquote{CUDA\_PATH}} are
necessary for \sphinxcode{\sphinxupquote{hipcc}} and the linker to work correctly.

\sphinxAtStartPar
\DUrole{versionmodified,added}{Added in version 3Aug2022.}

\sphinxAtStartPar
Using the CHIP\sphinxhyphen{}SPV implementation of HIP is supported. It allows one to
run HIP code on Intel GPUs via the OpenCL or Level Zero backends. To use
CHIP\sphinxhyphen{}SPV, you must set \sphinxcode{\sphinxupquote{\sphinxhyphen{}DHIP\_USE\_DEVICE\_SORT=OFF}} in your CMake
command line as CHIP\sphinxhyphen{}SPV does not yet support hipCUB. As of Summer 2022,
the use of HIP for Intel GPUs is experimental. You should only use this
option in preparations to run on Aurora system at Argonne.

\begin{sphinxVerbatim}[commandchars=\\\{\}]
\PYG{c+c1}{\PYGZsh{} AMDGPU target (ROCm \PYGZlt{}= 4.0)}
\PYG{n+nb}{export}\PYG{+w}{ }\PYG{n+nv}{HIP\PYGZus{}PLATFORM}\PYG{o}{=}hcc
\PYG{n+nb}{export}\PYG{+w}{ }\PYG{n+nv}{HIP\PYGZus{}PATH}\PYG{o}{=}/path/to/HIP/install
\PYG{n+nb}{export}\PYG{+w}{ }\PYG{n+nv}{HCC\PYGZus{}AMDGPU\PYGZus{}TARGET}\PYG{o}{=}gfx906
cmake\PYG{+w}{ }\PYGZhy{}D\PYG{+w}{ }\PYG{n+nv}{PKG\PYGZus{}GPU}\PYG{o}{=}on\PYG{+w}{ }\PYGZhy{}D\PYG{+w}{ }\PYG{n+nv}{GPU\PYGZus{}API}\PYG{o}{=}HIP\PYG{+w}{ }\PYGZhy{}D\PYG{+w}{ }\PYG{n+nv}{HIP\PYGZus{}ARCH}\PYG{o}{=}gfx906\PYG{+w}{ }\PYGZhy{}D\PYG{+w}{ }\PYG{n+nv}{CMAKE\PYGZus{}CXX\PYGZus{}COMPILER}\PYG{o}{=}hipcc\PYG{+w}{ }..
make\PYG{+w}{ }\PYGZhy{}j\PYG{+w}{ }\PYG{l+m}{4}
\end{sphinxVerbatim}

\begin{sphinxVerbatim}[commandchars=\\\{\}]
\PYG{c+c1}{\PYGZsh{} AMDGPU target (ROCm \PYGZgt{}= 4.1)}
\PYG{n+nb}{export}\PYG{+w}{ }\PYG{n+nv}{HIP\PYGZus{}PLATFORM}\PYG{o}{=}amd
\PYG{n+nb}{export}\PYG{+w}{ }\PYG{n+nv}{HIP\PYGZus{}PATH}\PYG{o}{=}/path/to/HIP/install
cmake\PYG{+w}{ }\PYGZhy{}D\PYG{+w}{ }\PYG{n+nv}{PKG\PYGZus{}GPU}\PYG{o}{=}on\PYG{+w}{ }\PYGZhy{}D\PYG{+w}{ }\PYG{n+nv}{GPU\PYGZus{}API}\PYG{o}{=}HIP\PYG{+w}{ }\PYGZhy{}D\PYG{+w}{ }\PYG{n+nv}{HIP\PYGZus{}ARCH}\PYG{o}{=}gfx906\PYG{+w}{ }\PYGZhy{}D\PYG{+w}{ }\PYG{n+nv}{CMAKE\PYGZus{}CXX\PYGZus{}COMPILER}\PYG{o}{=}hipcc\PYG{+w}{ }..
make\PYG{+w}{ }\PYGZhy{}j\PYG{+w}{ }\PYG{l+m}{4}
\end{sphinxVerbatim}

\begin{sphinxVerbatim}[commandchars=\\\{\}]
\PYG{c+c1}{\PYGZsh{} CUDA target (not recommended, use GPU\PYGZus{}ARCH=cuda)}
\PYG{c+c1}{\PYGZsh{} !!! DO NOT set CMAKE\PYGZus{}CXX\PYGZus{}COMPILER !!!}
\PYG{n+nb}{export}\PYG{+w}{ }\PYG{n+nv}{HIP\PYGZus{}PLATFORM}\PYG{o}{=}nvcc
\PYG{n+nb}{export}\PYG{+w}{ }\PYG{n+nv}{HIP\PYGZus{}PATH}\PYG{o}{=}/path/to/HIP/install
\PYG{n+nb}{export}\PYG{+w}{ }\PYG{n+nv}{CUDA\PYGZus{}PATH}\PYG{o}{=}/usr/local/cuda
cmake\PYG{+w}{ }\PYGZhy{}D\PYG{+w}{ }\PYG{n+nv}{PKG\PYGZus{}GPU}\PYG{o}{=}on\PYG{+w}{ }\PYGZhy{}D\PYG{+w}{ }\PYG{n+nv}{GPU\PYGZus{}API}\PYG{o}{=}HIP\PYG{+w}{ }\PYGZhy{}D\PYG{+w}{ }\PYG{n+nv}{HIP\PYGZus{}ARCH}\PYG{o}{=}sm\PYGZus{}70\PYG{+w}{ }..
make\PYG{+w}{ }\PYGZhy{}j\PYG{+w}{ }\PYG{l+m}{4}
\end{sphinxVerbatim}

\begin{sphinxVerbatim}[commandchars=\\\{\}]
\PYG{c+c1}{\PYGZsh{} SPIR\PYGZhy{}V target (Intel GPUs)}
\PYG{n+nb}{export}\PYG{+w}{ }\PYG{n+nv}{HIP\PYGZus{}PLATFORM}\PYG{o}{=}spirv
\PYG{n+nb}{export}\PYG{+w}{ }\PYG{n+nv}{HIP\PYGZus{}PATH}\PYG{o}{=}/path/to/HIP/install
\PYG{n+nb}{export}\PYG{+w}{ }\PYG{n+nv}{CMAKE\PYGZus{}CXX\PYGZus{}COMPILER}\PYG{o}{=}\PYGZlt{}hipcc/clang++\PYGZgt{}
cmake\PYG{+w}{ }\PYGZhy{}D\PYG{+w}{ }\PYG{n+nv}{PKG\PYGZus{}GPU}\PYG{o}{=}on\PYG{+w}{ }\PYGZhy{}D\PYG{+w}{ }\PYG{n+nv}{GPU\PYGZus{}API}\PYG{o}{=}HIP\PYG{+w}{ }..
make\PYG{+w}{ }\PYGZhy{}j\PYG{+w}{ }\PYG{l+m}{4}
\end{sphinxVerbatim}


\subsubsection{Traditional make}
\label{\detokenize{Build_extras:traditional-make}}
\sphinxAtStartPar
Before building LAMMPS, you must build the GPU library in \sphinxcode{\sphinxupquote{lib/gpu}}.
You can do this manually if you prefer; follow the instructions in
\sphinxcode{\sphinxupquote{lib/gpu/README}}.  Note that the GPU library uses MPI calls, so you must
use the same MPI library (or the STUBS library) settings as the main
LAMMPS code.  This also applies to the \sphinxcode{\sphinxupquote{\sphinxhyphen{}DLAMMPS\_BIGBIG}},
\sphinxcode{\sphinxupquote{\sphinxhyphen{}DLAMMPS\_SMALLBIG}}, or \sphinxcode{\sphinxupquote{\sphinxhyphen{}DLAMMPS\_SMALLSMALL}} settings in whichever
Makefile you use.

\sphinxAtStartPar
You can also build the library in one step from the \sphinxcode{\sphinxupquote{lammps/src}} dir,
using a command like these, which simply invokes the \sphinxcode{\sphinxupquote{lib/gpu/Install.py}}
script with the specified args:

\begin{sphinxVerbatim}[commandchars=\\\{\}]
\PYG{c+c1}{\PYGZsh{} print help message}
make\PYG{+w}{ }lib\PYGZhy{}gpu

\PYG{c+c1}{\PYGZsh{} build GPU library with default Makefile.linux}
make\PYG{+w}{ }lib\PYGZhy{}gpu\PYG{+w}{ }\PYG{n+nv}{args}\PYG{o}{=}\PYG{l+s+s2}{\PYGZdq{}\PYGZhy{}b\PYGZdq{}}

\PYG{c+c1}{\PYGZsh{} create new Makefile.xk7.single, altered for single\PYGZhy{}precision}
make\PYG{+w}{ }lib\PYGZhy{}gpu\PYG{+w}{ }\PYG{n+nv}{args}\PYG{o}{=}\PYG{l+s+s2}{\PYGZdq{}\PYGZhy{}m xk7 \PYGZhy{}p single \PYGZhy{}o xk7.single\PYGZdq{}}

\PYG{c+c1}{\PYGZsh{} build GPU library with mixed precision and P100 using other settings in Makefile.mpi}
make\PYG{+w}{ }lib\PYGZhy{}gpu\PYG{+w}{ }\PYG{n+nv}{args}\PYG{o}{=}\PYG{l+s+s2}{\PYGZdq{}\PYGZhy{}m mpi \PYGZhy{}a sm\PYGZus{}60 \PYGZhy{}p mixed \PYGZhy{}b\PYGZdq{}}
\end{sphinxVerbatim}

\sphinxAtStartPar
Note that this procedure starts with a Makefile.machine in lib/gpu, as
specified by the \sphinxcode{\sphinxupquote{\sphinxhyphen{}m}} switch.  For your convenience, machine makefiles
for “mpi” and “serial” are provided, which have the same settings as
the corresponding machine makefiles in the main LAMMPS source
folder. In addition you can alter 4 important settings in the
Makefile.machine you start from via the corresponding \sphinxcode{\sphinxupquote{\sphinxhyphen{}c}}, \sphinxcode{\sphinxupquote{\sphinxhyphen{}a}}, \sphinxcode{\sphinxupquote{\sphinxhyphen{}p}}, \sphinxcode{\sphinxupquote{\sphinxhyphen{}e}}
switches (as in the examples above), and also save a copy of the new
Makefile if desired:
\begin{itemize}
\item {} 
\sphinxAtStartPar
\sphinxcode{\sphinxupquote{CUDA\_HOME}} = where NVIDIA CUDA software is installed on your system

\item {} 
\sphinxAtStartPar
\sphinxcode{\sphinxupquote{CUDA\_ARCH}} = \sphinxcode{\sphinxupquote{sm\_XX}}, what GPU hardware you have, same as CMake \sphinxcode{\sphinxupquote{GPU\_ARCH}} above

\item {} 
\sphinxAtStartPar
\sphinxcode{\sphinxupquote{CUDA\_PRECISION}} = precision (double, mixed, single)

\item {} 
\sphinxAtStartPar
\sphinxcode{\sphinxupquote{EXTRAMAKE}} = which \sphinxcode{\sphinxupquote{Makefile.lammps.*}} file to copy to Makefile.lammps

\end{itemize}

\sphinxAtStartPar
The file \sphinxcode{\sphinxupquote{Makefile.cuda}} is set up to include support for multiple
GPU architectures as supported by the CUDA toolkit in use. This is done
through using the \sphinxcode{\sphinxupquote{\sphinxhyphen{}\sphinxhyphen{}gencode}} flag, which can be used multiple times and
thus support all GPU architectures supported by your CUDA compiler.

\sphinxAtStartPar
To enable GPU binning via CUDA performance primitives set the Makefile variable
\sphinxcode{\sphinxupquote{CUDPP\_OPT = \sphinxhyphen{}DUSE\_CUDPP \sphinxhyphen{}Icudpp\_mini}}.  This should \sphinxstylestrong{not} be used with
most modern GPUs.

\sphinxAtStartPar
To support the CUDA multiprocessor server you can set the define
\sphinxcode{\sphinxupquote{\sphinxhyphen{}DCUDA\_MPS\_SUPPORT}}.  Please note that in this case you must \sphinxstylestrong{not} use
the CUDA performance primitives and thus set the variable \sphinxcode{\sphinxupquote{CUDPP\_OPT}}
to empty.

\sphinxAtStartPar
The GPU library has some multi\sphinxhyphen{}thread support using OpenMP.  You need to add
the compiler flag that enables OpenMP to the \sphinxcode{\sphinxupquote{CUDR\_OPTS}} Makefile variable.

\sphinxAtStartPar
If the library build is successful, 3 files should be created:
\sphinxcode{\sphinxupquote{lib/gpu/libgpu.a}}, \sphinxcode{\sphinxupquote{lib/gpu/nvc\_get\_devices}}, and
\sphinxcode{\sphinxupquote{lib/gpu/Makefile.lammps}}.  The latter has settings that enable LAMMPS
to link with CUDA libraries.  If the settings in \sphinxcode{\sphinxupquote{Makefile.lammps}} for
your machine are not correct, the LAMMPS build will fail, and
\sphinxcode{\sphinxupquote{lib/gpu/Makefile.lammps}} may need to be edited.

\begin{sphinxadmonition}{note}{Note:}
\sphinxAtStartPar
If you re\sphinxhyphen{}build the GPU library in \sphinxcode{\sphinxupquote{lib/gpu}}, you should always
uninstall the GPU package in \sphinxcode{\sphinxupquote{lammps/src}}, then re\sphinxhyphen{}install it and
re\sphinxhyphen{}build LAMMPS.  This is because the compilation of files in the GPU
package uses the library settings from the \sphinxcode{\sphinxupquote{lib/gpu/Makefile.machine}}
used to build the GPU library.
\end{sphinxadmonition}


\bigskip\hrule\bigskip



\subsection{KIM package}
\label{\detokenize{Build_extras:kim-package}}\label{\detokenize{Build_extras:kim}}
\sphinxAtStartPar
To build with this package, the KIM library with API v2 must be downloaded
and built on your system. It must include the KIM models that you want to
use with LAMMPS.

\sphinxAtStartPar
If you would like to use the {\hyperref[\detokenize{kim_commands::doc}]{\sphinxcrossref{\DUrole{doc}{kim query}}}}
command, you also need to have libcurl installed with the matching
development headers and the curl\sphinxhyphen{}config tool.

\sphinxAtStartPar
If you would like to use the {\hyperref[\detokenize{kim_commands::doc}]{\sphinxcrossref{\DUrole{doc}{kim property}}}}
command, you need to build LAMMPS with the PYTHON package installed
and linked to Python 3.6 or later. See the {\hyperref[\detokenize{Build_extras:python}]{\sphinxcrossref{\DUrole{std,std-ref}{PYTHON package build info}}}}
for more details on this. After successfully building LAMMPS with Python, you
also need to install the \sphinxcode{\sphinxupquote{kim\sphinxhyphen{}property}} Python package, which can be easily
done using \sphinxstyleemphasis{pip} as \sphinxcode{\sphinxupquote{pip install kim\sphinxhyphen{}property}}, or from the \sphinxstyleemphasis{conda\sphinxhyphen{}forge}
channel as \sphinxcode{\sphinxupquote{conda install kim\sphinxhyphen{}property}} if LAMMPS is built in Conda. More
detailed information is available at:
\sphinxhref{https://github.com/openkim/kim-property\#installing-kim-property}{kim\sphinxhyphen{}property installation}.

\sphinxAtStartPar
In addition to installing the KIM API, it is also necessary to install the
library of KIM models (interatomic potentials).
See \sphinxhref{https://openkim.org/doc/usage/obtaining-models}{Obtaining KIM Models} to
learn how to install a pre\sphinxhyphen{}build binary of the OpenKIM Repository of Models.
See the list of all KIM models here: \sphinxurl{https://openkim.org/browse/models}

\sphinxAtStartPar
(Also note that when downloading and installing from source
the KIM API library with all its models, may take a long time (tens of
minutes to hours) to build.  Of course you only need to do that once.)
\begin{quote}

\begin{sphinxadmonition}{note}{CMake build}

\begin{sphinxVerbatim}[commandchars=\\\{\}]
\PYGZhy{}D\PYG{+w}{ }\PYG{n+nv}{DOWNLOAD\PYGZus{}KIM}\PYG{o}{=}value\PYG{+w}{           }\PYG{c+c1}{\PYGZsh{} download OpenKIM API v2 for build}
\PYG{+w}{                                }\PYG{c+c1}{\PYGZsh{} value = no (default) or yes}
\PYGZhy{}D\PYG{+w}{ }\PYG{n+nv}{LMP\PYGZus{}DEBUG\PYGZus{}CURL}\PYG{o}{=}value\PYG{+w}{         }\PYG{c+c1}{\PYGZsh{} set libcurl verbose mode on/off}
\PYG{+w}{                                }\PYG{c+c1}{\PYGZsh{} value = off (default) or on}
\PYGZhy{}D\PYG{+w}{ }\PYG{n+nv}{LMP\PYGZus{}NO\PYGZus{}SSL\PYGZus{}CHECK}\PYG{o}{=}value\PYG{+w}{       }\PYG{c+c1}{\PYGZsh{} tell libcurl to not verify the peer}
\PYG{+w}{                                }\PYG{c+c1}{\PYGZsh{} value = no (default) or yes}
\PYGZhy{}D\PYG{+w}{ }\PYG{n+nv}{KIM\PYGZus{}EXTRA\PYGZus{}UNITTESTS}\PYG{o}{=}value\PYG{+w}{    }\PYG{c+c1}{\PYGZsh{} enables extra unit tests}
\PYG{+w}{                                }\PYG{c+c1}{\PYGZsh{} value = no (default) or yes}
\end{sphinxVerbatim}

\sphinxAtStartPar
If \sphinxcode{\sphinxupquote{DOWNLOAD\_KIM}} is set to \sphinxcode{\sphinxupquote{yes}} (or \sphinxcode{\sphinxupquote{on}}), the KIM API library
will be downloaded and built inside the CMake build directory.  If
the KIM library is already installed on your system (in a location
where CMake cannot find it), you may need to set the
\sphinxcode{\sphinxupquote{PKG\_CONFIG\_PATH}} environment variable so that libkim\sphinxhyphen{}api can be
found, or run the command \sphinxcode{\sphinxupquote{source kim\sphinxhyphen{}api\sphinxhyphen{}activate}}.

\sphinxAtStartPar
Extra unit tests can only be available if they are explicitly requested
(\sphinxcode{\sphinxupquote{KIM\_EXTRA\_UNITTESTS}} is set to \sphinxcode{\sphinxupquote{yes}} (or \sphinxcode{\sphinxupquote{on}})) and the prerequisites
are met. See {\hyperref[\detokenize{Build_extras:kim-extra-unittests}]{\sphinxcrossref{\DUrole{std,std-ref}{KIM Extra unit tests}}}} for
more details on this.
\end{sphinxadmonition}

\begin{sphinxadmonition}{note}{Traditional make}

\sphinxAtStartPar
You can download and build the KIM library manually if you prefer;
follow the instructions in \sphinxcode{\sphinxupquote{lib/kim/README}}.  You can also do
this in one step from the lammps/src directory, using a command like
these, which simply invokes the \sphinxcode{\sphinxupquote{lib/kim/Install.py}} script with
the specified args.

\begin{sphinxVerbatim}[commandchars=\\\{\}]
\PYG{c+c1}{\PYGZsh{} print help message}
make\PYG{+w}{ }lib\PYGZhy{}kim

\PYG{c+c1}{\PYGZsh{} (re\PYGZhy{})install KIM API lib with only example models}
make\PYG{+w}{ }lib\PYGZhy{}kim\PYG{+w}{ }\PYG{n+nv}{args}\PYG{o}{=}\PYG{l+s+s2}{\PYGZdq{}\PYGZhy{}b\PYGZdq{}}

\PYG{c+c1}{\PYGZsh{} ditto plus one model}
make\PYG{+w}{ }lib\PYGZhy{}kim\PYG{+w}{ }\PYG{n+nv}{args}\PYG{o}{=}\PYG{l+s+s2}{\PYGZdq{}\PYGZhy{}b \PYGZhy{}a Glue\PYGZus{}Ercolessi\PYGZus{}Adams\PYGZus{}Al\PYGZus{}\PYGZus{}MO\PYGZus{}324507536345\PYGZus{}001\PYGZdq{}}

\PYG{c+c1}{\PYGZsh{} install KIM API lib with all models}
make\PYG{+w}{ }lib\PYGZhy{}kim\PYG{+w}{ }\PYG{n+nv}{args}\PYG{o}{=}\PYG{l+s+s2}{\PYGZdq{}\PYGZhy{}b \PYGZhy{}a everything\PYGZdq{}}

\PYG{c+c1}{\PYGZsh{} add one model or model driver}
make\PYG{+w}{ }lib\PYGZhy{}kim\PYG{+w}{ }\PYG{n+nv}{args}\PYG{o}{=}\PYG{l+s+s2}{\PYGZdq{}\PYGZhy{}n \PYGZhy{}a EAM\PYGZus{}Dynamo\PYGZus{}Ackland\PYGZus{}W\PYGZus{}\PYGZus{}MO\PYGZus{}141627196590\PYGZus{}002\PYGZdq{}}

\PYG{c+c1}{\PYGZsh{} use an existing KIM API installation at the provided location}
make\PYG{+w}{ }lib\PYGZhy{}kim\PYG{+w}{ }\PYG{n+nv}{args}\PYG{o}{=}\PYG{l+s+s2}{\PYGZdq{}\PYGZhy{}p \PYGZlt{}prefix\PYGZgt{}\PYGZdq{}}

\PYG{c+c1}{\PYGZsh{} ditto but add one model or driver}
make\PYG{+w}{ }lib\PYGZhy{}kim\PYG{+w}{ }\PYG{n+nv}{args}\PYG{o}{=}\PYG{l+s+s2}{\PYGZdq{}\PYGZhy{}p \PYGZlt{}prefix\PYGZgt{} \PYGZhy{}a EAM\PYGZus{}Dynamo\PYGZus{}Ackland\PYGZus{}W\PYGZus{}\PYGZus{}MO\PYGZus{}141627196590\PYGZus{}002\PYGZdq{}}
\end{sphinxVerbatim}

\sphinxAtStartPar
When using the \sphinxcode{\sphinxupquote{\sphinxhyphen{}b}} option, the KIM library is built using its native
cmake build system.  The \sphinxcode{\sphinxupquote{lib/kim/Install.py}} script supports a
\sphinxcode{\sphinxupquote{CMAKE}} environment variable if the cmake executable is named other
than \sphinxcode{\sphinxupquote{cmake}} on your system.  Additional environment variables may be
provided on the command line for use by cmake.  For example, to use the
\sphinxcode{\sphinxupquote{cmake3}} executable and tell it to use the gnu version 11 compilers
to build KIM, one could use the following command line.

\begin{sphinxVerbatim}[commandchars=\\\{\}]
\PYG{c+c1}{\PYGZsh{} (re\PYGZhy{})install KIM API lib using cmake3 and gnu v11 compilers}
\PYG{c+c1}{\PYGZsh{} with only example models}
\PYG{n+nv}{CMAKE}\PYG{o}{=}cmake3\PYG{+w}{ }\PYG{n+nv}{CXX}\PYG{o}{=}g++\PYGZhy{}11\PYG{+w}{ }\PYG{n+nv}{CC}\PYG{o}{=}gcc\PYGZhy{}11\PYG{+w}{ }\PYG{n+nv}{FC}\PYG{o}{=}gfortran\PYGZhy{}11\PYG{+w}{ }make\PYG{+w}{ }lib\PYGZhy{}kim\PYG{+w}{ }\PYG{n+nv}{args}\PYG{o}{=}\PYG{l+s+s2}{\PYGZdq{}\PYGZhy{}b\PYGZdq{}}
\end{sphinxVerbatim}

\sphinxAtStartPar
Settings for debugging OpenKIM web queries discussed below need to
be applied by adding them to the \sphinxcode{\sphinxupquote{LMP\_INC}} variable through
editing the \sphinxcode{\sphinxupquote{Makefile.machine}} you are using.  For example:

\begin{sphinxVerbatim}[commandchars=\\\{\}]
\PYG{n+nv}{LMP\PYGZus{}INC}\PYG{+w}{ }\PYG{o}{=}\PYG{+w}{ }\PYGZhy{}DLMP\PYGZus{}NO\PYGZus{}SSL\PYGZus{}CHECK
\end{sphinxVerbatim}
\end{sphinxadmonition}
\end{quote}


\subsubsection{Debugging OpenKIM web queries in LAMMPS}
\label{\detokenize{Build_extras:debugging-openkim-web-queries-in-lammps}}
\sphinxAtStartPar
If \sphinxcode{\sphinxupquote{LMP\_DEBUG\_CURL}} is set, the libcurl verbose mode will be turned
on, and any libcurl calls within the KIM web query display a lot of
information about libcurl operations.  You hardly ever want this set in
production use, you will almost always want this when you debug or
report problems.

\sphinxAtStartPar
The libcurl library performs peer SSL certificate verification by
default.  This verification is done using a CA certificate store that
the SSL library can use to make sure the peer’s server certificate is
valid.  If SSL reports an error (“certificate verify failed”) during the
handshake and thus refuses further communicate with that server, you can
set \sphinxcode{\sphinxupquote{LMP\_NO\_SSL\_CHECK}} to override that behavior.  When LAMMPS is
compiled with \sphinxcode{\sphinxupquote{LMP\_NO\_SSL\_CHECK}} set, libcurl does not verify the peer
and connection attempts will succeed regardless of the names in the
certificate. This option is insecure.  As an alternative, you can
specify your own CA cert path by setting the environment variable
\sphinxcode{\sphinxupquote{CURL\_CA\_BUNDLE}} to the path of your choice.  A call to the KIM web
query would get this value from the environment variable.


\subsubsection{KIM Extra unit tests (CMake only)}
\label{\detokenize{Build_extras:kim-extra-unit-tests-cmake-only}}\label{\detokenize{Build_extras:kim-extra-unittests}}
\sphinxAtStartPar
During development, testing, or debugging, if
{\hyperref[\detokenize{Build_development::doc}]{\sphinxcrossref{\DUrole{doc}{unit testing}}}} is enabled in LAMMPS, one can also
enable extra tests on {\hyperref[\detokenize{kim_commands::doc}]{\sphinxcrossref{\DUrole{doc}{KIM commands}}}} by setting the
\sphinxcode{\sphinxupquote{KIM\_EXTRA\_UNITTESTS}} to \sphinxcode{\sphinxupquote{yes}} (or \sphinxcode{\sphinxupquote{on}}).

\sphinxAtStartPar
Enabling the extra unit tests have some requirements,
\begin{itemize}
\item {} 
\sphinxAtStartPar
It requires to have internet access.

\item {} 
\sphinxAtStartPar
It requires to have libcurl installed with the matching development headers
and the curl\sphinxhyphen{}config tool.

\item {} 
\sphinxAtStartPar
It requires to build LAMMPS with the PYTHON package installed and linked to
Python 3.6 or later. See the {\hyperref[\detokenize{Build_extras:python}]{\sphinxcrossref{\DUrole{std,std-ref}{PYTHON package build info}}}} for
more details on this.

\item {} 
\sphinxAtStartPar
It requires to have \sphinxcode{\sphinxupquote{kim\sphinxhyphen{}property}} Python package installed, which can be
easily done using \sphinxstyleemphasis{pip} as \sphinxcode{\sphinxupquote{pip install kim\sphinxhyphen{}property}}, or from the
\sphinxstyleemphasis{conda\sphinxhyphen{}forge} channel as \sphinxcode{\sphinxupquote{conda install kim\sphinxhyphen{}property}} if LAMMPS is built in
Conda. More detailed information is available at:
\sphinxhref{https://github.com/openkim/kim-property\#installing-kim-property}{kim\sphinxhyphen{}property installation}.

\item {} 
\sphinxAtStartPar
It is also necessary to install the following KIM models:
\begin{itemize}
\item {} 
\sphinxAtStartPar
\sphinxcode{\sphinxupquote{EAM\_Dynamo\_MendelevAckland\_2007v3\_Zr\_\_MO\_004835508849\_000}}

\item {} 
\sphinxAtStartPar
\sphinxcode{\sphinxupquote{EAM\_Dynamo\_ErcolessiAdams\_1994\_Al\_\_MO\_123629422045\_005}}

\item {} 
\sphinxAtStartPar
\sphinxcode{\sphinxupquote{LennardJones612\_UniversalShifted\_\_MO\_959249795837\_003}}

\end{itemize}

\sphinxAtStartPar
See \sphinxhref{https://openkim.org/doc/usage/obtaining-models}{Obtaining KIM Models}
to learn how to install a pre\sphinxhyphen{}built binary of the OpenKIM Repository of
Models or see
\sphinxhref{https://openkim.org/doc/usage/obtaining-models/\#installing\_models}{Installing KIM Models}
to learn how to install the specific KIM models.

\end{itemize}


\bigskip\hrule\bigskip



\subsection{KOKKOS package}
\label{\detokenize{Build_extras:kokkos-package}}\label{\detokenize{Build_extras:kokkos}}
\sphinxAtStartPar
Using the KOKKOS package requires choosing several settings.  You have
to select whether you want to compile with parallelization on the host
and whether you want to include offloading of calculations to a device
(e.g. a GPU).  The default setting is to have no host parallelization
and no device offloading.  In addition, you can select the hardware
architecture to select the instruction set.  Since most hardware is
backward compatible, you may choose settings for an older architecture
to have an executable that will run on this and newer architectures.

\begin{sphinxadmonition}{note}{Note:}
\sphinxAtStartPar
If you run Kokkos on a different GPU architecture than what LAMMPS
was compiled with, there will be a delay during device initialization
while the just\sphinxhyphen{}in\sphinxhyphen{}time compiler is recompiling all GPU kernels for
the new hardware.  This is, however, only supported for GPUs of the
\sphinxstylestrong{same} major hardware version and different minor hardware versions,
e.g. 5.0 and 5.2 but not 5.2 and 6.0.  LAMMPS will abort with an
error message indicating a mismatch, if that happens.
\end{sphinxadmonition}

\sphinxAtStartPar
The settings discussed below have been tested with LAMMPS and are
confirmed to work.  Kokkos is an active project with ongoing improvements
and projects working on including support for additional architectures.
More information on Kokkos can be found on the
\sphinxhref{https://github.com/kokkos}{Kokkos GitHub project}.


\subsubsection{Available Architecture settings}
\label{\detokenize{Build_extras:available-architecture-settings}}
\sphinxAtStartPar
These are the possible choices for the Kokkos architecture ID.
They must be specified in uppercase.


\begin{savenotes}
\sphinxatlongtablestart
\sphinxthistablewithglobalstyle
\makeatletter
  \LTleft \@totalleftmargin plus1fill
  \LTright\dimexpr\columnwidth-\@totalleftmargin-\linewidth\relax plus1fill
\makeatother
\begin{longtable}{lll}
\sphinxtoprule
\endfirsthead

\multicolumn{3}{c}{\sphinxnorowcolor
    \makebox[0pt]{\sphinxtablecontinued{\tablename\ \thetable{} \textendash{} continued from previous page}}%
}\\
\sphinxtoprule
\endhead

\sphinxbottomrule
\multicolumn{3}{r}{\sphinxnorowcolor
    \makebox[0pt][r]{\sphinxtablecontinued{continues on next page}}%
}\\
\endfoot

\endlastfoot
\sphinxtableatstartofbodyhook

\sphinxAtStartPar
\sphinxstylestrong{Arch\sphinxhyphen{}ID}
&
\sphinxAtStartPar
\sphinxstylestrong{HOST or GPU}
&
\sphinxAtStartPar
\sphinxstylestrong{Description}
\\
\sphinxhline
\sphinxAtStartPar
NATIVE
&
\sphinxAtStartPar
HOST
&
\sphinxAtStartPar
Local machine
\\
\sphinxhline
\sphinxAtStartPar
AMDAVX
&
\sphinxAtStartPar
HOST
&
\sphinxAtStartPar
AMD 64\sphinxhyphen{}bit x86 CPU (AVX 1)
\\
\sphinxhline
\sphinxAtStartPar
ZEN
&
\sphinxAtStartPar
HOST
&
\sphinxAtStartPar
AMD Zen class CPU (AVX 2)
\\
\sphinxhline
\sphinxAtStartPar
ZEN2
&
\sphinxAtStartPar
HOST
&
\sphinxAtStartPar
AMD Zen2 class CPU (AVX 2)
\\
\sphinxhline
\sphinxAtStartPar
ZEN3
&
\sphinxAtStartPar
HOST
&
\sphinxAtStartPar
AMD Zen3 class CPU (AVX 2)
\\
\sphinxhline
\sphinxAtStartPar
ARMV80
&
\sphinxAtStartPar
HOST
&
\sphinxAtStartPar
ARMv8.0 Compatible CPU
\\
\sphinxhline
\sphinxAtStartPar
ARMV81
&
\sphinxAtStartPar
HOST
&
\sphinxAtStartPar
ARMv8.1 Compatible CPU
\\
\sphinxhline
\sphinxAtStartPar
ARMV8\_THUNDERX
&
\sphinxAtStartPar
HOST
&
\sphinxAtStartPar
ARMv8 Cavium ThunderX CPU
\\
\sphinxhline
\sphinxAtStartPar
ARMV8\_THUNDERX2
&
\sphinxAtStartPar
HOST
&
\sphinxAtStartPar
ARMv8 Cavium ThunderX2 CPU
\\
\sphinxhline
\sphinxAtStartPar
A64FX
&
\sphinxAtStartPar
HOST
&
\sphinxAtStartPar
ARMv8.2 with SVE Support
\\
\sphinxhline
\sphinxAtStartPar
SNB
&
\sphinxAtStartPar
HOST
&
\sphinxAtStartPar
Intel Sandy/Ivy Bridge CPU (AVX 1)
\\
\sphinxhline
\sphinxAtStartPar
HSW
&
\sphinxAtStartPar
HOST
&
\sphinxAtStartPar
Intel Haswell CPU (AVX 2)
\\
\sphinxhline
\sphinxAtStartPar
BDW
&
\sphinxAtStartPar
HOST
&
\sphinxAtStartPar
Intel Broadwell Xeon E\sphinxhyphen{}class CPU (AVX 2 + transactional mem)
\\
\sphinxhline
\sphinxAtStartPar
SKL
&
\sphinxAtStartPar
HOST
&
\sphinxAtStartPar
Intel Skylake Client CPU
\\
\sphinxhline
\sphinxAtStartPar
SKX
&
\sphinxAtStartPar
HOST
&
\sphinxAtStartPar
Intel Skylake Xeon Server CPU (AVX512)
\\
\sphinxhline
\sphinxAtStartPar
ICL
&
\sphinxAtStartPar
HOST
&
\sphinxAtStartPar
Intel Ice Lake Client CPU (AVX512)
\\
\sphinxhline
\sphinxAtStartPar
ICX
&
\sphinxAtStartPar
HOST
&
\sphinxAtStartPar
Intel Ice Lake Xeon Server CPU (AVX512)
\\
\sphinxhline
\sphinxAtStartPar
SPR
&
\sphinxAtStartPar
HOST
&
\sphinxAtStartPar
Intel Sapphire Rapids Xeon Server CPU (AVX512)
\\
\sphinxhline
\sphinxAtStartPar
KNC
&
\sphinxAtStartPar
HOST
&
\sphinxAtStartPar
Intel Knights Corner Xeon Phi
\\
\sphinxhline
\sphinxAtStartPar
KNL
&
\sphinxAtStartPar
HOST
&
\sphinxAtStartPar
Intel Knights Landing Xeon Phi
\\
\sphinxhline
\sphinxAtStartPar
POWER8
&
\sphinxAtStartPar
HOST
&
\sphinxAtStartPar
IBM POWER8 CPU
\\
\sphinxhline
\sphinxAtStartPar
POWER9
&
\sphinxAtStartPar
HOST
&
\sphinxAtStartPar
IBM POWER9 CPU
\\
\sphinxhline
\sphinxAtStartPar
RISCV\_SG2042
&
\sphinxAtStartPar
HOST
&
\sphinxAtStartPar
SG2042 (RISC\sphinxhyphen{}V) CPU
\\
\sphinxhline
\sphinxAtStartPar
KEPLER30
&
\sphinxAtStartPar
GPU
&
\sphinxAtStartPar
NVIDIA Kepler generation CC 3.0 GPU
\\
\sphinxhline
\sphinxAtStartPar
KEPLER32
&
\sphinxAtStartPar
GPU
&
\sphinxAtStartPar
NVIDIA Kepler generation CC 3.2 GPU
\\
\sphinxhline
\sphinxAtStartPar
KEPLER35
&
\sphinxAtStartPar
GPU
&
\sphinxAtStartPar
NVIDIA Kepler generation CC 3.5 GPU
\\
\sphinxhline
\sphinxAtStartPar
KEPLER37
&
\sphinxAtStartPar
GPU
&
\sphinxAtStartPar
NVIDIA Kepler generation CC 3.7 GPU
\\
\sphinxhline
\sphinxAtStartPar
MAXWELL50
&
\sphinxAtStartPar
GPU
&
\sphinxAtStartPar
NVIDIA Maxwell generation CC 5.0 GPU
\\
\sphinxhline
\sphinxAtStartPar
MAXWELL52
&
\sphinxAtStartPar
GPU
&
\sphinxAtStartPar
NVIDIA Maxwell generation CC 5.2 GPU
\\
\sphinxhline
\sphinxAtStartPar
MAXWELL53
&
\sphinxAtStartPar
GPU
&
\sphinxAtStartPar
NVIDIA Maxwell generation CC 5.3 GPU
\\
\sphinxhline
\sphinxAtStartPar
PASCAL60
&
\sphinxAtStartPar
GPU
&
\sphinxAtStartPar
NVIDIA Pascal generation CC 6.0 GPU
\\
\sphinxhline
\sphinxAtStartPar
PASCAL61
&
\sphinxAtStartPar
GPU
&
\sphinxAtStartPar
NVIDIA Pascal generation CC 6.1 GPU
\\
\sphinxhline
\sphinxAtStartPar
VOLTA70
&
\sphinxAtStartPar
GPU
&
\sphinxAtStartPar
NVIDIA Volta generation CC 7.0 GPU
\\
\sphinxhline
\sphinxAtStartPar
VOLTA72
&
\sphinxAtStartPar
GPU
&
\sphinxAtStartPar
NVIDIA Volta generation CC 7.2 GPU
\\
\sphinxhline
\sphinxAtStartPar
TURING75
&
\sphinxAtStartPar
GPU
&
\sphinxAtStartPar
NVIDIA Turing generation CC 7.5 GPU
\\
\sphinxhline
\sphinxAtStartPar
AMPERE80
&
\sphinxAtStartPar
GPU
&
\sphinxAtStartPar
NVIDIA Ampere generation CC 8.0 GPU
\\
\sphinxhline
\sphinxAtStartPar
AMPERE86
&
\sphinxAtStartPar
GPU
&
\sphinxAtStartPar
NVIDIA Ampere generation CC 8.6 GPU
\\
\sphinxhline
\sphinxAtStartPar
ADA89
&
\sphinxAtStartPar
GPU
&
\sphinxAtStartPar
NVIDIA Ada Lovelace generation CC 8.9 GPU
\\
\sphinxhline
\sphinxAtStartPar
HOPPER90
&
\sphinxAtStartPar
GPU
&
\sphinxAtStartPar
NVIDIA Hopper generation CC 9.0 GPU
\\
\sphinxhline
\sphinxAtStartPar
AMD\_GFX906
&
\sphinxAtStartPar
GPU
&
\sphinxAtStartPar
AMD GPU MI50/MI60
\\
\sphinxhline
\sphinxAtStartPar
AMD\_GFX908
&
\sphinxAtStartPar
GPU
&
\sphinxAtStartPar
AMD GPU MI100
\\
\sphinxhline
\sphinxAtStartPar
AMD\_GFX90A
&
\sphinxAtStartPar
GPU
&
\sphinxAtStartPar
AMD GPU MI200
\\
\sphinxhline
\sphinxAtStartPar
AMD\_GFX942
&
\sphinxAtStartPar
GPU
&
\sphinxAtStartPar
AMD GPU MI300
\\
\sphinxhline
\sphinxAtStartPar
AMD\_GFX1030
&
\sphinxAtStartPar
GPU
&
\sphinxAtStartPar
AMD GPU V620/W6800
\\
\sphinxhline
\sphinxAtStartPar
AMD\_GFX1100
&
\sphinxAtStartPar
GPU
&
\sphinxAtStartPar
AMD GPU RX7900XTX
\\
\sphinxhline
\sphinxAtStartPar
AMD\_GFX1103
&
\sphinxAtStartPar
GPU
&
\sphinxAtStartPar
AMD Phoenix APU with Radeon 740M/760M/780M/880M/890M
\\
\sphinxhline
\sphinxAtStartPar
INTEL\_GEN
&
\sphinxAtStartPar
GPU
&
\sphinxAtStartPar
SPIR64\sphinxhyphen{}based devices, e.g. Intel GPUs, using JIT
\\
\sphinxhline
\sphinxAtStartPar
INTEL\_DG1
&
\sphinxAtStartPar
GPU
&
\sphinxAtStartPar
Intel Iris XeMAX GPU
\\
\sphinxhline
\sphinxAtStartPar
INTEL\_GEN9
&
\sphinxAtStartPar
GPU
&
\sphinxAtStartPar
Intel GPU Gen9
\\
\sphinxhline
\sphinxAtStartPar
INTEL\_GEN11
&
\sphinxAtStartPar
GPU
&
\sphinxAtStartPar
Intel GPU Gen11
\\
\sphinxhline
\sphinxAtStartPar
INTEL\_GEN12LP
&
\sphinxAtStartPar
GPU
&
\sphinxAtStartPar
Intel GPU Gen12LP
\\
\sphinxhline
\sphinxAtStartPar
INTEL\_XEHP
&
\sphinxAtStartPar
GPU
&
\sphinxAtStartPar
Intel GPU Xe\sphinxhyphen{}HP
\\
\sphinxhline
\sphinxAtStartPar
INTEL\_PVC
&
\sphinxAtStartPar
GPU
&
\sphinxAtStartPar
Intel GPU Ponte Vecchio
\\
\sphinxbottomrule
\end{longtable}
\sphinxtableafterendhook
\sphinxatlongtableend
\end{savenotes}

\sphinxAtStartPar
This list was last updated for version 4.3.0 of the Kokkos library.
\begin{quote}

\begin{sphinxadmonition}{note}{Basic CMake build settings:}

\sphinxAtStartPar
For multicore CPUs using OpenMP, set these 2 variables.

\begin{sphinxVerbatim}[commandchars=\\\{\}]
\PYGZhy{}D\PYG{+w}{ }\PYG{n+nv}{Kokkos\PYGZus{}ARCH\PYGZus{}HOSTARCH}\PYG{o}{=}yes\PYG{+w}{  }\PYG{c+c1}{\PYGZsh{} HOSTARCH = HOST from list above}
\PYGZhy{}D\PYG{+w}{ }\PYG{n+nv}{Kokkos\PYGZus{}ENABLE\PYGZus{}OPENMP}\PYG{o}{=}yes
\PYGZhy{}D\PYG{+w}{ }\PYG{n+nv}{BUILD\PYGZus{}OMP}\PYG{o}{=}yes
\end{sphinxVerbatim}

\sphinxAtStartPar
Please note that enabling OpenMP for KOKKOS requires that OpenMP is
also {\hyperref[\detokenize{Build_basics:serial}]{\sphinxcrossref{\DUrole{std,std-ref}{enabled for the rest of LAMMPS}}}}.

\sphinxAtStartPar
For Intel KNLs using OpenMP, set these variables:

\begin{sphinxVerbatim}[commandchars=\\\{\}]
\PYGZhy{}D\PYG{+w}{ }\PYG{n+nv}{Kokkos\PYGZus{}ARCH\PYGZus{}KNL}\PYG{o}{=}yes
\PYGZhy{}D\PYG{+w}{ }\PYG{n+nv}{Kokkos\PYGZus{}ENABLE\PYGZus{}OPENMP}\PYG{o}{=}yes
\end{sphinxVerbatim}

\sphinxAtStartPar
For NVIDIA GPUs using CUDA, set these variables:

\begin{sphinxVerbatim}[commandchars=\\\{\}]
\PYGZhy{}D\PYG{+w}{ }\PYG{n+nv}{Kokkos\PYGZus{}ARCH\PYGZus{}HOSTARCH}\PYG{o}{=}yes\PYG{+w}{   }\PYG{c+c1}{\PYGZsh{} HOSTARCH = HOST from list above}
\PYGZhy{}D\PYG{+w}{ }\PYG{n+nv}{Kokkos\PYGZus{}ARCH\PYGZus{}GPUARCH}\PYG{o}{=}yes\PYG{+w}{    }\PYG{c+c1}{\PYGZsh{} GPUARCH = GPU from list above}
\PYGZhy{}D\PYG{+w}{ }\PYG{n+nv}{Kokkos\PYGZus{}ENABLE\PYGZus{}CUDA}\PYG{o}{=}yes
\PYGZhy{}D\PYG{+w}{ }\PYG{n+nv}{Kokkos\PYGZus{}ENABLE\PYGZus{}OPENMP}\PYG{o}{=}yes
\end{sphinxVerbatim}

\sphinxAtStartPar
This will also enable executing FFTs on the GPU, either via the
internal KISSFFT library, or \sphinxhyphen{} by preference \sphinxhyphen{} with the cuFFT
library bundled with the CUDA toolkit, depending on whether CMake
can identify its location.

\sphinxAtStartPar
For AMD or NVIDIA GPUs using HIP, set these variables:

\begin{sphinxVerbatim}[commandchars=\\\{\}]
\PYGZhy{}D\PYG{+w}{ }\PYG{n+nv}{Kokkos\PYGZus{}ARCH\PYGZus{}HOSTARCH}\PYG{o}{=}yes\PYG{+w}{   }\PYG{c+c1}{\PYGZsh{} HOSTARCH = HOST from list above}
\PYGZhy{}D\PYG{+w}{ }\PYG{n+nv}{Kokkos\PYGZus{}ARCH\PYGZus{}GPUARCH}\PYG{o}{=}yes\PYG{+w}{    }\PYG{c+c1}{\PYGZsh{} GPUARCH = GPU from list above}
\PYGZhy{}D\PYG{+w}{ }\PYG{n+nv}{Kokkos\PYGZus{}ENABLE\PYGZus{}HIP}\PYG{o}{=}yes
\PYGZhy{}D\PYG{+w}{ }\PYG{n+nv}{Kokkos\PYGZus{}ENABLE\PYGZus{}OPENMP}\PYG{o}{=}yes
\end{sphinxVerbatim}

\sphinxAtStartPar
This will enable FFTs on the GPU, either by the internal KISSFFT library
or with the hipFFT wrapper library, which will call out to the
platform\sphinxhyphen{}appropriate vendor library: rocFFT on AMD GPUs or cuFFT on
NVIDIA GPUs.

\sphinxAtStartPar
To simplify compilation, five preset files are included in the
\sphinxcode{\sphinxupquote{cmake/presets}} folder, \sphinxcode{\sphinxupquote{kokkos\sphinxhyphen{}serial.cmake}},
\sphinxcode{\sphinxupquote{kokkos\sphinxhyphen{}openmp.cmake}}, \sphinxcode{\sphinxupquote{kokkos\sphinxhyphen{}cuda.cmake}},
\sphinxcode{\sphinxupquote{kokkos\sphinxhyphen{}hip.cmake}}, and \sphinxcode{\sphinxupquote{kokkos\sphinxhyphen{}sycl.cmake}}.  They will enable
the KOKKOS package and enable some hardware choices.  For GPU
support those preset files must be customized to match the
hardware used. So to compile with CUDA device parallelization with
some common packages enabled, you can do the following:

\begin{sphinxVerbatim}[commandchars=\\\{\}]
mkdir\PYG{+w}{ }build\PYGZhy{}kokkos\PYGZhy{}cuda
\PYG{n+nb}{cd}\PYG{+w}{ }build\PYGZhy{}kokkos\PYGZhy{}cuda
cmake\PYG{+w}{ }\PYGZhy{}C\PYG{+w}{ }../cmake/presets/basic.cmake\PYG{+w}{ }\PYG{l+s+se}{\PYGZbs{}}
\PYG{+w}{      }\PYGZhy{}C\PYG{+w}{ }../cmake/presets/kokkos\PYGZhy{}cuda.cmake\PYG{+w}{ }../cmake
cmake\PYG{+w}{ }\PYGZhy{}\PYGZhy{}build\PYG{+w}{ }.
\end{sphinxVerbatim}
\end{sphinxadmonition}

\begin{sphinxadmonition}{note}{Basic traditional make settings:}

\sphinxAtStartPar
Choose which hardware to support in \sphinxcode{\sphinxupquote{Makefile.machine}} via
\sphinxcode{\sphinxupquote{KOKKOS\_DEVICES}} and \sphinxcode{\sphinxupquote{KOKKOS\_ARCH}} settings.  See the
\sphinxcode{\sphinxupquote{src/MAKE/OPTIONS/Makefile.kokkos*}} files for examples.

\sphinxAtStartPar
For multicore CPUs using OpenMP:

\begin{sphinxVerbatim}[commandchars=\\\{\}]
\PYG{n+nv}{KOKKOS\PYGZus{}DEVICES}\PYG{+w}{ }\PYG{o}{=}\PYG{+w}{ }OpenMP
\PYG{n+nv}{KOKKOS\PYGZus{}ARCH}\PYG{+w}{ }\PYG{o}{=}\PYG{+w}{ }HOSTARCH\PYG{+w}{          }\PYG{c+c1}{\PYGZsh{} HOSTARCH = HOST from list above}
\end{sphinxVerbatim}

\sphinxAtStartPar
For Intel KNLs using OpenMP:

\begin{sphinxVerbatim}[commandchars=\\\{\}]
\PYG{n+nv}{KOKKOS\PYGZus{}DEVICES}\PYG{+w}{ }\PYG{o}{=}\PYG{+w}{ }OpenMP
\PYG{n+nv}{KOKKOS\PYGZus{}ARCH}\PYG{+w}{ }\PYG{o}{=}\PYG{+w}{ }KNL
\end{sphinxVerbatim}

\sphinxAtStartPar
For NVIDIA GPUs using CUDA:

\begin{sphinxVerbatim}[commandchars=\\\{\}]
\PYG{n+nv}{KOKKOS\PYGZus{}DEVICES}\PYG{+w}{ }\PYG{o}{=}\PYG{+w}{ }Cuda
\PYG{n+nv}{KOKKOS\PYGZus{}ARCH}\PYG{+w}{ }\PYG{o}{=}\PYG{+w}{ }HOSTARCH,GPUARCH\PYG{+w}{  }\PYG{c+c1}{\PYGZsh{} HOSTARCH = HOST from list above that is}
\PYG{c}{                                \PYGZsh{}            hosting the GPU}
\PYG{c}{                                \PYGZsh{} GPUARCH = GPU from list above}
\PYG{n+nv}{KOKKOS\PYGZus{}CUDA\PYGZus{}OPTIONS}\PYG{+w}{ }\PYG{o}{=}\PYG{+w}{ }\PYG{l+s+s2}{\PYGZdq{}enable\PYGZus{}lambda\PYGZdq{}}
\PYG{n+nv}{FFT\PYGZus{}INC}\PYG{+w}{ }\PYG{o}{=}\PYG{+w}{ }\PYGZhy{}DFFT\PYGZus{}CUFFT\PYG{+w}{           }\PYG{c+c1}{\PYGZsh{} enable use of cuFFT (optional)}
\PYG{n+nv}{FFT\PYGZus{}LIB}\PYG{+w}{ }\PYG{o}{=}\PYG{+w}{ }\PYGZhy{}lcufft\PYG{+w}{               }\PYG{c+c1}{\PYGZsh{} link to cuFFT library}
\end{sphinxVerbatim}

\sphinxAtStartPar
For GPUs, you also need the following lines in your
\sphinxcode{\sphinxupquote{Makefile.machine}} before the CC line is defined.  They tell
\sphinxcode{\sphinxupquote{mpicxx}} to use an \sphinxcode{\sphinxupquote{nvcc}} compiler wrapper, which will use
\sphinxcode{\sphinxupquote{nvcc}} for compiling CUDA files and a C++ compiler for
non\sphinxhyphen{}Kokkos, non\sphinxhyphen{}CUDA files.

\begin{sphinxVerbatim}[commandchars=\\\{\}]
\PYG{c}{\PYGZsh{} For OpenMPI}
\PYG{n+nv}{KOKKOS\PYGZus{}ABSOLUTE\PYGZus{}PATH}\PYG{+w}{ }\PYG{o}{=}\PYG{+w}{ }\PYG{k}{\PYGZdl{}(}shell\PYG{+w}{ }\PYG{n+nb}{cd}\PYG{+w}{ }\PYG{k}{\PYGZdl{}(}KOKKOS\PYGZus{}PATH\PYG{k}{)}\PYG{p}{;}\PYG{+w}{ }\PYG{n+nb}{pwd}\PYG{k}{)}
\PYG{k}{export }\PYG{n+nv}{OMPI\PYGZus{}CXX}\PYG{+w}{ }\PYG{o}{=}\PYG{+w}{ }\PYG{k}{\PYGZdl{}(}KOKKOS\PYGZus{}ABSOLUTE\PYGZus{}PATH\PYG{k}{)}/config/nvcc\PYGZus{}wrapper
\PYG{n+nv}{CC}\PYG{+w}{ }\PYG{o}{=}\PYG{+w}{ }mpicxx
\end{sphinxVerbatim}

\begin{sphinxVerbatim}[commandchars=\\\{\}]
\PYG{c}{\PYGZsh{} For MPICH and derivatives}
\PYG{n+nv}{KOKKOS\PYGZus{}ABSOLUTE\PYGZus{}PATH}\PYG{+w}{ }\PYG{o}{=}\PYG{+w}{ }\PYG{k}{\PYGZdl{}(}shell\PYG{+w}{ }\PYG{n+nb}{cd}\PYG{+w}{ }\PYG{k}{\PYGZdl{}(}KOKKOS\PYGZus{}PATH\PYG{k}{)}\PYG{p}{;}\PYG{+w}{ }\PYG{n+nb}{pwd}\PYG{k}{)}
\PYG{n+nv}{CC}\PYG{+w}{ }\PYG{o}{=}\PYG{+w}{ }mpicxx\PYG{+w}{ }\PYGZhy{}cxx\PYG{o}{=}\PYG{k}{\PYGZdl{}(}KOKKOS\PYGZus{}ABSOLUTE\PYGZus{}PATH\PYG{k}{)}/config/nvcc\PYGZus{}wrapper
\end{sphinxVerbatim}

\sphinxAtStartPar
For AMD or NVIDIA GPUs using HIP:

\begin{sphinxVerbatim}[commandchars=\\\{\}]
\PYG{n+nv}{KOKKOS\PYGZus{}DEVICES}\PYG{+w}{ }\PYG{o}{=}\PYG{+w}{ }HIP
\PYG{n+nv}{KOKKOS\PYGZus{}ARCH}\PYG{+w}{ }\PYG{o}{=}\PYG{+w}{ }HOSTARCH,GPUARCH\PYG{+w}{  }\PYG{c+c1}{\PYGZsh{} HOSTARCH = HOST from list above that is}
\PYG{c}{                                \PYGZsh{}            hosting the GPU}
\PYG{c}{                                \PYGZsh{} GPUARCH = GPU from list above}
\PYG{n+nv}{FFT\PYGZus{}INC}\PYG{+w}{ }\PYG{o}{=}\PYG{+w}{ }\PYGZhy{}DFFT\PYGZus{}HIPFFT\PYG{+w}{          }\PYG{c+c1}{\PYGZsh{} enable use of hipFFT (optional)}
\PYG{n+nv}{FFT\PYGZus{}LIB}\PYG{+w}{ }\PYG{o}{=}\PYG{+w}{ }\PYGZhy{}lhipfft\PYG{+w}{              }\PYG{c+c1}{\PYGZsh{} link to hipFFT library}
\end{sphinxVerbatim}
\end{sphinxadmonition}
\end{quote}


\subsubsection{Advanced KOKKOS compilation settings}
\label{\detokenize{Build_extras:advanced-kokkos-compilation-settings}}
\sphinxAtStartPar
There are other allowed options when building with the KOKKOS package
that can improve performance or assist in debugging or profiling. Below
are some examples that may be useful in combination with LAMMPS.  For
the full list (which keeps changing as the Kokkos package itself evolves),
please consult the Kokkos library documentation.

\sphinxAtStartPar
As alternative to using multi\sphinxhyphen{}threading via OpenMP
(\sphinxcode{\sphinxupquote{\sphinxhyphen{}DKokkos\_ENABLE\_OPENMP=on}} or \sphinxcode{\sphinxupquote{KOKKOS\_DEVICES=OpenMP}}) it is also
possible to use Posix threads directly (\sphinxcode{\sphinxupquote{\sphinxhyphen{}DKokkos\_ENABLE\_PTHREAD=on}}
or \sphinxcode{\sphinxupquote{KOKKOS\_DEVICES=Pthread}}).  While binding of threads to individual
or groups of CPU cores is managed in OpenMP with environment variables,
you need assistance from either the “hwloc” or “libnuma” library for the
Pthread thread parallelization option. To enable use with CMake:
\sphinxcode{\sphinxupquote{\sphinxhyphen{}DKokkos\_ENABLE\_HWLOC=on}} or \sphinxcode{\sphinxupquote{\sphinxhyphen{}DKokkos\_ENABLE\_LIBNUMA=on}}; and with
conventional make: \sphinxcode{\sphinxupquote{KOKKOS\_USE\_TPLS=hwloc}} or
\sphinxcode{\sphinxupquote{KOKKOS\_USE\_TPLS=libnuma}}.

\sphinxAtStartPar
The CMake option \sphinxcode{\sphinxupquote{\sphinxhyphen{}DKokkos\_ENABLE\_LIBRT=on}} or the makefile setting
\sphinxcode{\sphinxupquote{KOKKOS\_USE\_TPLS=librt}} enables the use of a more accurate timer
mechanism on many Unix\sphinxhyphen{}like platforms for internal profiling.

\sphinxAtStartPar
The CMake option \sphinxcode{\sphinxupquote{\sphinxhyphen{}DKokkos\_ENABLE\_DEBUG=on}} or the makefile setting
\sphinxcode{\sphinxupquote{KOKKOS\_DEBUG=yes}} enables printing of run\sphinxhyphen{}time
debugging information that can be useful. It also enables runtime
bounds checking on Kokkos data structures.  As to be expected, enabling
this option will negatively impact the performance and thus is only
recommended when developing a Kokkos\sphinxhyphen{}enabled style in LAMMPS.

\sphinxAtStartPar
The CMake option \sphinxcode{\sphinxupquote{\sphinxhyphen{}DKokkos\_ENABLE\_CUDA\_UVM=on}} or the makefile
setting \sphinxcode{\sphinxupquote{KOKKOS\_CUDA\_OPTIONS=enable\_lambda,force\_uvm}} enables the
use of CUDA “Unified Virtual Memory” (UVM) in Kokkos.  UVM allows to
transparently use RAM on the host to supplement the memory used on the
GPU (with some performance penalty) and thus enables running larger
problems that would otherwise not fit into the RAM on the GPU.

\sphinxAtStartPar
Please note, that the LAMMPS KOKKOS package must \sphinxstylestrong{always} be compiled
with the \sphinxstyleemphasis{enable\_lambda} option when using GPUs.  The CMake configuration
will thus always enable it.


\bigskip\hrule\bigskip



\subsection{LEPTON package}
\label{\detokenize{Build_extras:lepton-package}}\label{\detokenize{Build_extras:lepton}}
\sphinxAtStartPar
To build with this package, you must build the Lepton library which is
included in the LAMMPS source distribution in the \sphinxcode{\sphinxupquote{lib/lepton}} folder.
\begin{quote}

\begin{sphinxadmonition}{note}{CMake build}

\sphinxAtStartPar
This is the recommended build procedure for using Lepton in
LAMMPS. No additional settings are normally needed besides
\sphinxcode{\sphinxupquote{\sphinxhyphen{}D PKG\_LEPTON=yes}}.

\sphinxAtStartPar
On x86 hardware the Lepton library will also include a just\sphinxhyphen{}in\sphinxhyphen{}time
compiler for faster execution.  This is auto detected but can
be explicitly disabled by setting \sphinxcode{\sphinxupquote{\sphinxhyphen{}D LEPTON\_ENABLE\_JIT=no}}
(or enabled by setting it to yes).
\end{sphinxadmonition}

\begin{sphinxadmonition}{note}{Traditional make}

\sphinxAtStartPar
Before building LAMMPS, one must build the Lepton library in lib/lepton.

\sphinxAtStartPar
This can be done manually in the same folder by using or adapting
one of the provided Makefiles: for example, \sphinxcode{\sphinxupquote{Makefile.serial}} for
the GNU C++ compiler, or \sphinxcode{\sphinxupquote{Makefile.mpi}} for the MPI compiler wrapper.
The Lepton library is written in C++\sphinxhyphen{}11 and thus the C++ compiler
may need to be instructed to enable support for that.

\sphinxAtStartPar
In general, it is safer to use build setting consistent with the
rest of LAMMPS.  This is best carried out from the LAMMPS src
directory using a command like these, which simply invokes the
\sphinxcode{\sphinxupquote{lib/lepton/Install.py}} script with the specified args:

\begin{sphinxVerbatim}[commandchars=\\\{\}]
\PYG{c+c1}{\PYGZsh{} print help message}
make\PYG{+w}{ }lib\PYGZhy{}lepton

\PYG{c+c1}{\PYGZsh{} build with GNU g++ compiler (settings as with \PYGZdq{}make serial\PYGZdq{})}
make\PYG{+w}{ }lib\PYGZhy{}lepton\PYG{+w}{ }\PYG{n+nv}{args}\PYG{o}{=}\PYG{l+s+s2}{\PYGZdq{}\PYGZhy{}m serial\PYGZdq{}}

\PYG{c+c1}{\PYGZsh{} build with default MPI compiler (settings as with \PYGZdq{}make mpi\PYGZdq{})}
make\PYG{+w}{ }lib\PYGZhy{}lepton\PYG{+w}{ }\PYG{n+nv}{args}\PYG{o}{=}\PYG{l+s+s2}{\PYGZdq{}\PYGZhy{}m mpi\PYGZdq{}}
\end{sphinxVerbatim}

\sphinxAtStartPar
The “machine” argument of the \sphinxcode{\sphinxupquote{\sphinxhyphen{}m}} flag is used to find a
Makefile.machine to use as build recipe.

\sphinxAtStartPar
The build should produce a \sphinxcode{\sphinxupquote{build}} folder and the library \sphinxcode{\sphinxupquote{lib/lepton/liblmplepton.a}}
\end{sphinxadmonition}
\end{quote}


\bigskip\hrule\bigskip



\subsection{MACHDYN package}
\label{\detokenize{Build_extras:machdyn-package}}\label{\detokenize{Build_extras:machdyn}}
\sphinxAtStartPar
To build with this package, you must download the Eigen3 library.
Eigen3 is a template library, so you do not need to build it.
\begin{quote}

\begin{sphinxadmonition}{note}{CMake build}

\begin{sphinxVerbatim}[commandchars=\\\{\}]
\PYGZhy{}D\PYG{+w}{ }DOWNLOAD\PYGZus{}EIGEN3\PYG{+w}{            }\PYG{c+c1}{\PYGZsh{} download Eigen3, value = no (default) or yes}
\PYGZhy{}D\PYG{+w}{ }\PYG{n+nv}{EIGEN3\PYGZus{}INCLUDE\PYGZus{}DIR}\PYG{o}{=}path\PYG{+w}{    }\PYG{c+c1}{\PYGZsh{} path to Eigen library (only needed if a}
\PYG{+w}{                              }\PYG{c+c1}{\PYGZsh{} custom location)}
\end{sphinxVerbatim}

\sphinxAtStartPar
If \sphinxcode{\sphinxupquote{DOWNLOAD\_EIGEN3}} is set, the Eigen3 library will be
downloaded and inside the CMake build directory.  If the Eigen3
library is already on your system (in a location where CMake
cannot find it), set \sphinxcode{\sphinxupquote{EIGEN3\_INCLUDE\_DIR}} to the directory the
\sphinxcode{\sphinxupquote{Eigen3}} include file is in.
\end{sphinxadmonition}

\begin{sphinxadmonition}{note}{Traditional make}

\sphinxAtStartPar
You can download the Eigen3 library manually if you prefer; follow
the instructions in \sphinxcode{\sphinxupquote{lib/machdyn/README}}.  You can also do it in one
step from the \sphinxcode{\sphinxupquote{lammps/src}} dir, using a command like these,
which simply invokes the \sphinxcode{\sphinxupquote{lib/machdyn/Install.py}} script with the
specified args:

\begin{sphinxVerbatim}[commandchars=\\\{\}]
\PYG{c+c1}{\PYGZsh{} print help message}
make\PYG{+w}{ }lib\PYGZhy{}machdyn

\PYG{c+c1}{\PYGZsh{} download to lib/machdyn/eigen3}
make\PYG{+w}{ }lib\PYGZhy{}machdyn\PYG{+w}{ }\PYG{n+nv}{args}\PYG{o}{=}\PYG{l+s+s2}{\PYGZdq{}\PYGZhy{}b\PYGZdq{}}

\PYG{c+c1}{\PYGZsh{} use existing Eigen installation in /usr/include/eigen3}
make\PYG{+w}{ }lib\PYGZhy{}machdyn\PYG{+w}{ }\PYG{n+nv}{args}\PYG{o}{=}\PYG{l+s+s2}{\PYGZdq{}\PYGZhy{}p /usr/include/eigen3\PYGZdq{}}
\end{sphinxVerbatim}

\sphinxAtStartPar
Note that a symbolic (soft) link named \sphinxcode{\sphinxupquote{includelink}} is created
in \sphinxcode{\sphinxupquote{lib/machdyn}} to point to the Eigen dir.  When LAMMPS builds it
will use this link.  You should not need to edit the
\sphinxcode{\sphinxupquote{lib/machdyn/Makefile.lammps}} file.
\end{sphinxadmonition}
\end{quote}


\bigskip\hrule\bigskip



\subsection{ML\sphinxhyphen{}IAP package}
\label{\detokenize{Build_extras:ml-iap-package}}\label{\detokenize{Build_extras:mliap}}
\sphinxAtStartPar
Building the ML\sphinxhyphen{}IAP package requires including the {\hyperref[\detokenize{Packages_details:pkg-ml-snap}]{\sphinxcrossref{\DUrole{std,std-ref}{ML\sphinxhyphen{}SNAP}}}} package.  There will be an error message if this requirement
is not satisfied.  Using the \sphinxstyleemphasis{mliappy} model also requires enabling
Python support, which in turn requires to include the {\hyperref[\detokenize{Packages_details:pkg-python}]{\sphinxcrossref{\DUrole{std,std-ref}{PYTHON}}}} package \sphinxstylestrong{and} requires to have the \sphinxhref{https://cython.org}{cython} software installed and with it a working
\sphinxcode{\sphinxupquote{cythonize}} command.  This feature requires compiling LAMMPS with
Python version 3.6 or later.
\begin{quote}

\begin{sphinxadmonition}{note}{CMake build}

\begin{sphinxVerbatim}[commandchars=\\\{\}]
\PYGZhy{}D\PYG{+w}{ }\PYG{n+nv}{MLIAP\PYGZus{}ENABLE\PYGZus{}PYTHON}\PYG{o}{=}value\PYG{+w}{   }\PYG{c+c1}{\PYGZsh{} enable mliappy model (default is autodetect)}
\end{sphinxVerbatim}

\sphinxAtStartPar
Without this setting, CMake will check whether it can find a
suitable Python version and the \sphinxcode{\sphinxupquote{cythonize}} command and choose
the default accordingly.  During the build procedure the provided
.pyx file(s) will be automatically translated to C++ code and compiled.
Please do \sphinxstylestrong{not} run \sphinxcode{\sphinxupquote{cythonize}} manually in the \sphinxcode{\sphinxupquote{src/ML\sphinxhyphen{}IAP}} folder,
as that can lead to compilation errors if Python support is not enabled.
If you did it by accident, please remove the generated .cpp and .h files.
\end{sphinxadmonition}

\begin{sphinxadmonition}{note}{Traditional make}

\sphinxAtStartPar
The build uses the \sphinxcode{\sphinxupquote{lib/python/Makefile.mliap\_python}} file in the
compile/link process to add a rule to update the files generated by
the \sphinxcode{\sphinxupquote{cythonize}} command in case the corresponding .pyx file(s) were
modified.  You may need to modify \sphinxcode{\sphinxupquote{lib/python/Makefile.lammps}}
if the LAMMPS build fails.

\sphinxAtStartPar
To enable building the ML\sphinxhyphen{}IAP package with Python support enabled,
you need to add \sphinxcode{\sphinxupquote{\sphinxhyphen{}DMLIAP\_PYTHON}} to the \sphinxcode{\sphinxupquote{LMP\_INC}} variable in
your machine makefile.  You may have to manually run the
\sphinxcode{\sphinxupquote{cythonize}} command on .pyx file(s) in the \sphinxcode{\sphinxupquote{src}} folder, if
this is not automatically done during installing the ML\sphinxhyphen{}IAP
package.  Please do \sphinxstylestrong{not} run \sphinxcode{\sphinxupquote{cythonize}} in the \sphinxcode{\sphinxupquote{src/ML\sphinxhyphen{}IAP}}
folder, as that can lead to compilation errors if Python support
is not enabled.  If you did this by accident, please remove the
generated .cpp and .h files.
\end{sphinxadmonition}
\end{quote}


\bigskip\hrule\bigskip



\subsection{OPT package}
\label{\detokenize{Build_extras:opt-package}}\label{\detokenize{Build_extras:opt}}\begin{quote}

\begin{sphinxadmonition}{note}{CMake build}

\sphinxAtStartPar
No additional settings are needed besides \sphinxcode{\sphinxupquote{\sphinxhyphen{}D PKG\_OPT=yes}}
\end{sphinxadmonition}

\begin{sphinxadmonition}{note}{Traditional make}

\sphinxAtStartPar
The compiler flag \sphinxcode{\sphinxupquote{\sphinxhyphen{}restrict}} must be used to build LAMMPS with
the OPT package when using Intel compilers.  It should be added to
the \sphinxcode{\sphinxupquote{CCFLAGS}} line of your \sphinxcode{\sphinxupquote{Makefile.machine}}.  See
\sphinxcode{\sphinxupquote{src/MAKE/OPTIONS/Makefile.opt}} for an example.
\end{sphinxadmonition}
\end{quote}


\bigskip\hrule\bigskip



\subsection{POEMS package}
\label{\detokenize{Build_extras:poems-package}}\label{\detokenize{Build_extras:poems}}\begin{quote}

\begin{sphinxadmonition}{note}{CMake build}

\sphinxAtStartPar
No additional settings are needed besides \sphinxcode{\sphinxupquote{\sphinxhyphen{}D PKG\_OPT=yes}}
\end{sphinxadmonition}

\begin{sphinxadmonition}{note}{Traditional make}

\sphinxAtStartPar
Before building LAMMPS, you must build the POEMS library in
\sphinxcode{\sphinxupquote{lib/poems}}.  You can do this manually if you prefer; follow
the instructions in \sphinxcode{\sphinxupquote{lib/poems/README}}.  You can also do it in
one step from the \sphinxcode{\sphinxupquote{lammps/src}} dir, using a command like these,
which simply invokes the \sphinxcode{\sphinxupquote{lib/poems/Install.py}} script with the
specified args:

\begin{sphinxVerbatim}[commandchars=\\\{\}]
\PYG{c+c1}{\PYGZsh{} print help message}
make\PYG{+w}{ }lib\PYGZhy{}poems

\PYG{c+c1}{\PYGZsh{} build with GNU g++ compiler (settings as with \PYGZdq{}make serial\PYGZdq{})}
make\PYG{+w}{ }lib\PYGZhy{}poems\PYG{+w}{ }\PYG{n+nv}{args}\PYG{o}{=}\PYG{l+s+s2}{\PYGZdq{}\PYGZhy{}m serial\PYGZdq{}}

\PYG{c+c1}{\PYGZsh{} build with default MPI C++ compiler (settings as with \PYGZdq{}make mpi\PYGZdq{})}
make\PYG{+w}{ }lib\PYGZhy{}poems\PYG{+w}{ }\PYG{n+nv}{args}\PYG{o}{=}\PYG{l+s+s2}{\PYGZdq{}\PYGZhy{}m mpi\PYGZdq{}}

\PYG{c+c1}{\PYGZsh{} build with Intel Classic compiler}
make\PYG{+w}{ }lib\PYGZhy{}poems\PYG{+w}{ }\PYG{n+nv}{args}\PYG{o}{=}\PYG{l+s+s2}{\PYGZdq{}\PYGZhy{}m icc\PYGZdq{}}
\end{sphinxVerbatim}

\sphinxAtStartPar
The build should produce two files: \sphinxcode{\sphinxupquote{lib/poems/libpoems.a}} and
\sphinxcode{\sphinxupquote{lib/poems/Makefile.lammps}}.  The latter is copied from an
existing \sphinxcode{\sphinxupquote{Makefile.lammps.*}} and has settings needed to build
LAMMPS with the POEMS library (though typically the settings are
just blank).  If necessary, you can edit/create a new
\sphinxcode{\sphinxupquote{lib/poems/Makefile.machine}} file for your system, which should
define an \sphinxcode{\sphinxupquote{EXTRAMAKE}} variable to specify a corresponding
\sphinxcode{\sphinxupquote{Makefile.lammps.machine}} file.
\end{sphinxadmonition}
\end{quote}


\bigskip\hrule\bigskip



\subsection{PYTHON package}
\label{\detokenize{Build_extras:python-package}}\label{\detokenize{Build_extras:python}}
\sphinxAtStartPar
Building with the PYTHON package requires you have a the Python development
headers and library available on your system, which needs to be a Python 2.7
version or a Python 3.x version.  Since support for Python 2.x has ended,
using Python 3.x is strongly recommended. See \sphinxcode{\sphinxupquote{lib/python/README}} for
additional details.
\begin{quote}

\begin{sphinxadmonition}{note}{CMake build}

\begin{sphinxVerbatim}[commandchars=\\\{\}]
\PYGZhy{}D\PYG{+w}{ }\PYG{n+nv}{Python\PYGZus{}EXECUTABLE}\PYG{o}{=}path\PYG{+w}{   }\PYG{c+c1}{\PYGZsh{} path to Python executable to use}
\end{sphinxVerbatim}

\sphinxAtStartPar
Without this setting, CMake will guess the default Python version
on your system.  To use a different Python version, you can either
create a virtualenv, activate it and then run cmake.  Or you can
set the Python\_EXECUTABLE variable to specify which Python
interpreter should be used.  Note note that you will also need to
have the development headers installed for this version,
e.g. python2\sphinxhyphen{}devel.
\end{sphinxadmonition}

\begin{sphinxadmonition}{note}{Traditional make}

\sphinxAtStartPar
The build uses the \sphinxcode{\sphinxupquote{lib/python/Makefile.lammps}} file in the
compile/link process to find Python.  You should only need to
create a new \sphinxcode{\sphinxupquote{Makefile.lammps.*}} file (and copy it to
\sphinxcode{\sphinxupquote{Makefile.lammps}}) if the LAMMPS build fails.
\end{sphinxadmonition}
\end{quote}


\bigskip\hrule\bigskip



\subsection{VORONOI package}
\label{\detokenize{Build_extras:voronoi-package}}\label{\detokenize{Build_extras:voronoi}}
\sphinxAtStartPar
To build with this package, you must download and build the
\sphinxhref{https://math.lbl.gov/voro++/}{Voro++ library} or install a
binary package provided by your operating system.
\begin{quote}

\begin{sphinxadmonition}{note}{CMake build}

\begin{sphinxVerbatim}[commandchars=\\\{\}]
\PYGZhy{}D\PYG{+w}{ }\PYG{n+nv}{DOWNLOAD\PYGZus{}VORO}\PYG{o}{=}value\PYG{+w}{    }\PYG{c+c1}{\PYGZsh{} download Voro++ for build}
\PYG{+w}{                          }\PYG{c+c1}{\PYGZsh{} value = no (default) or yes}
\PYGZhy{}D\PYG{+w}{ }\PYG{n+nv}{VORO\PYGZus{}LIBRARY}\PYG{o}{=}path\PYG{+w}{      }\PYG{c+c1}{\PYGZsh{} Voro++ library file}
\PYG{+w}{                          }\PYG{c+c1}{\PYGZsh{} (only needed if at custom location)}
\PYGZhy{}D\PYG{+w}{ }\PYG{n+nv}{VORO\PYGZus{}INCLUDE\PYGZus{}DIR}\PYG{o}{=}path\PYG{+w}{  }\PYG{c+c1}{\PYGZsh{} Voro++ include directory}
\PYG{+w}{                          }\PYG{c+c1}{\PYGZsh{} (only needed if at custom location)}
\end{sphinxVerbatim}

\sphinxAtStartPar
If \sphinxcode{\sphinxupquote{DOWNLOAD\_VORO}} is set, the Voro++ library will be downloaded
and built inside the CMake build directory.  If the Voro++ library
is already on your system (in a location CMake cannot find it),
\sphinxcode{\sphinxupquote{VORO\_LIBRARY}} is the filename (plus path) of the Voro++ library
file, not the directory the library file is in.
\sphinxcode{\sphinxupquote{VORO\_INCLUDE\_DIR}} is the directory the Voro++ include file is
in.
\end{sphinxadmonition}

\begin{sphinxadmonition}{note}{Traditional make}

\sphinxAtStartPar
You can download and build the Voro++ library manually if you
prefer; follow the instructions in \sphinxcode{\sphinxupquote{lib/voronoi/README}}.  You
can also do it in one step from the \sphinxcode{\sphinxupquote{lammps/src}} dir, using a
command like these, which simply invokes the
\sphinxcode{\sphinxupquote{lib/voronoi/Install.py}} script with the specified args:

\begin{sphinxVerbatim}[commandchars=\\\{\}]
\PYG{c+c1}{\PYGZsh{} print help message}
make\PYG{+w}{ }lib\PYGZhy{}voronoi

\PYG{c+c1}{\PYGZsh{} download and build the default version in lib/voronoi/voro++\PYGZhy{}\PYGZlt{}version\PYGZgt{}}
make\PYG{+w}{ }lib\PYGZhy{}voronoi\PYG{+w}{ }\PYG{n+nv}{args}\PYG{o}{=}\PYG{l+s+s2}{\PYGZdq{}\PYGZhy{}b\PYGZdq{}}

\PYG{c+c1}{\PYGZsh{} use existing Voro++ installation in \PYGZdl{}HOME/voro++}
make\PYG{+w}{ }lib\PYGZhy{}voronoi\PYG{+w}{ }\PYG{n+nv}{args}\PYG{o}{=}\PYG{l+s+s2}{\PYGZdq{}}\PYG{l+s+s2}{\PYGZhy{}p }\PYG{n+nv}{\PYGZdl{}HOME}\PYG{l+s+s2}{/voro++}\PYG{l+s+s2}{\PYGZdq{}}

\PYG{c+c1}{\PYGZsh{} download and build the 0.4.6 version in lib/voronoi/voro++\PYGZhy{}0.4.6}
make\PYG{+w}{ }lib\PYGZhy{}voronoi\PYG{+w}{ }\PYG{n+nv}{args}\PYG{o}{=}\PYG{l+s+s2}{\PYGZdq{}\PYGZhy{}b \PYGZhy{}v voro++0.4.6\PYGZdq{}}
\end{sphinxVerbatim}

\sphinxAtStartPar
Note that two symbolic (soft) links, \sphinxcode{\sphinxupquote{includelink}} and
\sphinxcode{\sphinxupquote{liblink}}, are created in lib/voronoi to point to the Voro++
source dir.  When LAMMPS builds in \sphinxcode{\sphinxupquote{src}} it will use these
links.  You should not need to edit the
\sphinxcode{\sphinxupquote{lib/voronoi/Makefile.lammps}} file.
\end{sphinxadmonition}
\end{quote}


\bigskip\hrule\bigskip



\subsection{ADIOS package}
\label{\detokenize{Build_extras:adios-package}}\label{\detokenize{Build_extras:adios}}
\sphinxAtStartPar
The ADIOS package requires the \sphinxhref{https://github.com/ornladios/ADIOS2}{ADIOS I/O library}, version 2.3.1 or newer. Make
sure that you have ADIOS built either with or without MPI to match if
you build LAMMPS with or without MPI.  ADIOS compilation settings for
LAMMPS are automatically detected, if the PATH and LD\_LIBRARY\_PATH
environment variables have been updated for the local ADIOS installation
and the instructions below are followed for the respective build
systems.
\begin{quote}

\begin{sphinxadmonition}{note}{CMake build}

\begin{sphinxVerbatim}[commandchars=\\\{\}]
\PYGZhy{}D\PYG{+w}{ }\PYG{n+nv}{ADIOS2\PYGZus{}DIR}\PYG{o}{=}path\PYG{+w}{        }\PYG{c+c1}{\PYGZsh{} path is where ADIOS 2.x is installed}
\PYGZhy{}D\PYG{+w}{ }\PYG{n+nv}{PKG\PYGZus{}ADIOS}\PYG{o}{=}yes
\end{sphinxVerbatim}
\end{sphinxadmonition}

\begin{sphinxadmonition}{note}{Traditional make}

\sphinxAtStartPar
Turn on the ADIOS package before building LAMMPS. If the
ADIOS 2.x software is installed in PATH, there is nothing else to
do:

\begin{sphinxVerbatim}[commandchars=\\\{\}]
make\PYG{+w}{ }yes\PYGZhy{}adios
\end{sphinxVerbatim}

\sphinxAtStartPar
otherwise, set ADIOS2\_DIR environment variable when turning on the package:

\begin{sphinxVerbatim}[commandchars=\\\{\}]
\PYG{n+nv}{ADIOS2\PYGZus{}DIR}\PYG{o}{=}path\PYG{+w}{ }make\PYG{+w}{ }yes\PYGZhy{}adios\PYG{+w}{   }\PYG{c+c1}{\PYGZsh{} path is where ADIOS 2.x is installed}
\end{sphinxVerbatim}
\end{sphinxadmonition}
\end{quote}


\bigskip\hrule\bigskip



\subsection{ATC package}
\label{\detokenize{Build_extras:atc-package}}\label{\detokenize{Build_extras:atc}}
\sphinxAtStartPar
The ATC package requires the MANYBODY package also be installed.
\begin{quote}

\begin{sphinxadmonition}{note}{CMake build}

\sphinxAtStartPar
No additional settings are needed besides \sphinxcode{\sphinxupquote{\sphinxhyphen{}D PKG\_ATC=yes}}
and \sphinxcode{\sphinxupquote{\sphinxhyphen{}D PKG\_MANYBODY=yes}}.
\end{sphinxadmonition}

\begin{sphinxadmonition}{note}{Traditional make}

\sphinxAtStartPar
Before building LAMMPS, you must build the ATC library in
\sphinxcode{\sphinxupquote{lib/atc}}.  You can do this manually if you prefer; follow the
instructions in \sphinxcode{\sphinxupquote{lib/atc/README}}.  You can also do it in one
step from the \sphinxcode{\sphinxupquote{lammps/src}} dir, using a command like these,
which simply invokes the \sphinxcode{\sphinxupquote{lib/atc/Install.py}} script with the
specified args:

\begin{sphinxVerbatim}[commandchars=\\\{\}]
\PYG{c+c1}{\PYGZsh{} print help message}
make\PYG{+w}{ }lib\PYGZhy{}atc

\PYG{c+c1}{\PYGZsh{} build with GNU g++ compiler and MPI STUBS (settings as with \PYGZdq{}make serial\PYGZdq{})}
make\PYG{+w}{ }lib\PYGZhy{}atc\PYG{+w}{ }\PYG{n+nv}{args}\PYG{o}{=}\PYG{l+s+s2}{\PYGZdq{}\PYGZhy{}m serial\PYGZdq{}}

\PYG{c+c1}{\PYGZsh{} build with default MPI compiler (settings as with \PYGZdq{}make mpi\PYGZdq{})}
make\PYG{+w}{ }lib\PYGZhy{}atc\PYG{+w}{ }\PYG{n+nv}{args}\PYG{o}{=}\PYG{l+s+s2}{\PYGZdq{}\PYGZhy{}m mpi\PYGZdq{}}

\PYG{c+c1}{\PYGZsh{} build with Intel Classic compiler}
make\PYG{+w}{ }lib\PYGZhy{}atc\PYG{+w}{ }\PYG{n+nv}{args}\PYG{o}{=}\PYG{l+s+s2}{\PYGZdq{}\PYGZhy{}m icc\PYGZdq{}}
\end{sphinxVerbatim}

\sphinxAtStartPar
The build should produce two files: \sphinxcode{\sphinxupquote{lib/atc/libatc.a}} and
\sphinxcode{\sphinxupquote{lib/atc/Makefile.lammps}}.  The latter is copied from an
existing \sphinxcode{\sphinxupquote{Makefile.lammps.*}} and has settings needed to build
LAMMPS with the ATC library.  If necessary, you can edit/create a
new \sphinxcode{\sphinxupquote{lib/atc/Makefile.machine}} file for your system, which
should define an \sphinxcode{\sphinxupquote{EXTRAMAKE}} variable to specify a corresponding
\sphinxcode{\sphinxupquote{Makefile.lammps.\textless{}machine\textgreater{}}} file.

\sphinxAtStartPar
Note that the Makefile.lammps file has settings for the BLAS and
LAPACK linear algebra libraries.  As explained in
\sphinxcode{\sphinxupquote{lib/atc/README}} these can either exist on your system, or you
can use the files provided in \sphinxcode{\sphinxupquote{lib/linalg}}.  In the latter case
you also need to build the library in \sphinxcode{\sphinxupquote{lib/linalg}} with a
command like these:

\begin{sphinxVerbatim}[commandchars=\\\{\}]
\PYG{c+c1}{\PYGZsh{} print help message}
make\PYG{+w}{ }lib\PYGZhy{}linalg

\PYG{c+c1}{\PYGZsh{} build with GNU C++ compiler (settings as with \PYGZdq{}make serial\PYGZdq{})}
make\PYG{+w}{ }lib\PYGZhy{}linalg\PYG{+w}{ }\PYG{n+nv}{args}\PYG{o}{=}\PYG{l+s+s2}{\PYGZdq{}\PYGZhy{}m serial\PYGZdq{}}

\PYG{c+c1}{\PYGZsh{} build with default MPI C++ compiler (settings as with \PYGZdq{}make mpi\PYGZdq{})}
make\PYG{+w}{ }lib\PYGZhy{}linalg\PYG{+w}{ }\PYG{n+nv}{args}\PYG{o}{=}\PYG{l+s+s2}{\PYGZdq{}\PYGZhy{}m mpi\PYGZdq{}}

\PYG{c+c1}{\PYGZsh{} build with GNU Fortran compiler}
make\PYG{+w}{ }lib\PYGZhy{}linalg\PYG{+w}{ }\PYG{n+nv}{args}\PYG{o}{=}\PYG{l+s+s2}{\PYGZdq{}\PYGZhy{}m g++\PYGZdq{}}
\end{sphinxVerbatim}
\end{sphinxadmonition}
\end{quote}


\bigskip\hrule\bigskip



\subsection{AWPMD package}
\label{\detokenize{Build_extras:awpmd-package}}\label{\detokenize{Build_extras:awpmd}}\begin{quote}

\begin{sphinxadmonition}{note}{CMake build}

\sphinxAtStartPar
No additional settings are needed besides \sphinxcode{\sphinxupquote{\sphinxhyphen{}D PKG\_AQPMD=yes}}.
\end{sphinxadmonition}

\begin{sphinxadmonition}{note}{Traditional make}

\sphinxAtStartPar
Before building LAMMPS, you must build the AWPMD library in
\sphinxcode{\sphinxupquote{lib/awpmd}}.  You can do this manually if you prefer; follow the
instructions in \sphinxcode{\sphinxupquote{lib/awpmd/README}}.  You can also do it in one
step from the \sphinxcode{\sphinxupquote{lammps/src}} dir, using a command like these,
which simply invokes the \sphinxcode{\sphinxupquote{lib/awpmd/Install.py}} script with the
specified args:

\begin{sphinxVerbatim}[commandchars=\\\{\}]
\PYG{c+c1}{\PYGZsh{} print help message}
make\PYG{+w}{ }lib\PYGZhy{}awpmd

\PYG{c+c1}{\PYGZsh{} build with GNU g++ compiler and MPI STUBS (settings as with \PYGZdq{}make serial\PYGZdq{})}
make\PYG{+w}{ }lib\PYGZhy{}awpmd\PYG{+w}{ }\PYG{n+nv}{args}\PYG{o}{=}\PYG{l+s+s2}{\PYGZdq{}\PYGZhy{}m serial\PYGZdq{}}

\PYG{c+c1}{\PYGZsh{} build with default MPI compiler (settings as with \PYGZdq{}make mpi\PYGZdq{})}
make\PYG{+w}{ }lib\PYGZhy{}awpmd\PYG{+w}{ }\PYG{n+nv}{args}\PYG{o}{=}\PYG{l+s+s2}{\PYGZdq{}\PYGZhy{}m mpi\PYGZdq{}}

\PYG{c+c1}{\PYGZsh{} build with Intel Classic compiler}
make\PYG{+w}{ }lib\PYGZhy{}awpmd\PYG{+w}{ }\PYG{n+nv}{args}\PYG{o}{=}\PYG{l+s+s2}{\PYGZdq{}\PYGZhy{}m icc\PYGZdq{}}
\end{sphinxVerbatim}

\sphinxAtStartPar
The build should produce two files: \sphinxcode{\sphinxupquote{lib/awpmd/libawpmd.a}} and
\sphinxcode{\sphinxupquote{lib/awpmd/Makefile.lammps}}.  The latter is copied from an
existing \sphinxcode{\sphinxupquote{Makefile.lammps.*}} and has settings needed to build
LAMMPS with the AWPMD library.  If necessary, you can edit/create
a new \sphinxcode{\sphinxupquote{lib/awpmd/Makefile.machine}} file for your system, which
should define an \sphinxcode{\sphinxupquote{EXTRAMAKE}} variable to specify a corresponding
\sphinxcode{\sphinxupquote{Makefile.lammps.\textless{}machine\textgreater{}}} file.

\sphinxAtStartPar
Note that the \sphinxcode{\sphinxupquote{Makefile.lammps}} file has settings for the BLAS
and LAPACK linear algebra libraries.  As explained in
\sphinxcode{\sphinxupquote{lib/awpmd/README}} these can either exist on your system, or you
can use the files provided in \sphinxcode{\sphinxupquote{lib/linalg}}.  In the latter case
you also need to build the library in \sphinxcode{\sphinxupquote{lib/linalg}} with a
command like these:

\begin{sphinxVerbatim}[commandchars=\\\{\}]
\PYG{c+c1}{\PYGZsh{} print help message}
make\PYG{+w}{ }lib\PYGZhy{}linalg

\PYG{c+c1}{\PYGZsh{} build with GNU C++ compiler (settings as with \PYGZdq{}make serial\PYGZdq{})}
make\PYG{+w}{ }lib\PYGZhy{}linalg\PYG{+w}{ }\PYG{n+nv}{args}\PYG{o}{=}\PYG{l+s+s2}{\PYGZdq{}\PYGZhy{}m serial\PYGZdq{}}

\PYG{c+c1}{\PYGZsh{} build with default MPI C++ compiler (settings as with \PYGZdq{}make mpi\PYGZdq{})}
make\PYG{+w}{ }lib\PYGZhy{}linalg\PYG{+w}{ }\PYG{n+nv}{args}\PYG{o}{=}\PYG{l+s+s2}{\PYGZdq{}\PYGZhy{}m mpi\PYGZdq{}}

\PYG{c+c1}{\PYGZsh{} build with GNU C++ compiler}
make\PYG{+w}{ }lib\PYGZhy{}linalg\PYG{+w}{ }\PYG{n+nv}{args}\PYG{o}{=}\PYG{l+s+s2}{\PYGZdq{}\PYGZhy{}m g++\PYGZdq{}}
\end{sphinxVerbatim}
\end{sphinxadmonition}
\end{quote}


\bigskip\hrule\bigskip



\subsection{COLVARS package}
\label{\detokenize{Build_extras:colvars-package}}\label{\detokenize{Build_extras:colvar}}
\sphinxAtStartPar
This package enables the use of the \sphinxhref{https://colvars.github.io/}{Colvars}
module included in the LAMMPS source distribution.
\begin{quote}

\begin{sphinxadmonition}{note}{CMake build}

\sphinxAtStartPar
This is the recommended build procedure for using Colvars in
LAMMPS. No additional settings are normally needed besides
\sphinxcode{\sphinxupquote{\sphinxhyphen{}D PKG\_COLVARS=yes}}.
\end{sphinxadmonition}

\begin{sphinxadmonition}{note}{Traditional make}

\sphinxAtStartPar
As with other libraries distributed with LAMMPS, the Colvars library
needs to be built before building the LAMMPS program with the COLVARS
package enabled.

\sphinxAtStartPar
From the LAMMPS \sphinxcode{\sphinxupquote{src}} directory, this is most easily and safely done
via one of the following commands, which implicitly rely on the
\sphinxcode{\sphinxupquote{lib/colvars/Install.py}} script with optional arguments:

\begin{sphinxVerbatim}[commandchars=\\\{\}]
\PYG{c+c1}{\PYGZsh{} print help message}
make\PYG{+w}{ }lib\PYGZhy{}colvars

\PYG{c+c1}{\PYGZsh{} build with GNU g++ compiler (settings as with \PYGZdq{}make serial\PYGZdq{})}
make\PYG{+w}{ }lib\PYGZhy{}colvars\PYG{+w}{ }\PYG{n+nv}{args}\PYG{o}{=}\PYG{l+s+s2}{\PYGZdq{}\PYGZhy{}m serial\PYGZdq{}}

\PYG{c+c1}{\PYGZsh{} build with default MPI compiler (settings as with \PYGZdq{}make mpi\PYGZdq{})}
make\PYG{+w}{ }lib\PYGZhy{}colvars\PYG{+w}{ }\PYG{n+nv}{args}\PYG{o}{=}\PYG{l+s+s2}{\PYGZdq{}\PYGZhy{}m mpi\PYGZdq{}}

\PYG{c+c1}{\PYGZsh{} build with GNU g++ compiler and colvars debugging enabled}
make\PYG{+w}{ }lib\PYGZhy{}colvars\PYG{+w}{ }\PYG{n+nv}{args}\PYG{o}{=}\PYG{l+s+s2}{\PYGZdq{}\PYGZhy{}m g++\PYGZhy{}debug\PYGZdq{}}
\end{sphinxVerbatim}

\sphinxAtStartPar
The “machine” argument of the “\sphinxhyphen{}m” flag is used to find a
\sphinxcode{\sphinxupquote{Makefile.machine}} file to use as build recipe.  If such recipe does
not already exist in \sphinxcode{\sphinxupquote{lib/colvars}}, suitable settings will be
auto\sphinxhyphen{}generated consistent with those used in the core LAMMPS makefiles.

\sphinxAtStartPar
\DUrole{versionmodified,changed}{Changed in version 8Feb2023.}

\sphinxAtStartPar
Please note that Colvars uses the Lepton library, which is now
included with the LEPTON package; if you use anything other than
the \sphinxcode{\sphinxupquote{make lib\sphinxhyphen{}colvars}} command, please make sure to {\hyperref[\detokenize{Build_extras:lepton}]{\sphinxcrossref{\DUrole{std,std-ref}{build
Lepton beforehand}}}}.

\sphinxAtStartPar
Optional flags may be specified as environment variables:

\begin{sphinxVerbatim}[commandchars=\\\{\}]
\PYG{c+c1}{\PYGZsh{} Build with debug code (much slower)}
\PYG{n+nv}{COLVARS\PYGZus{}DEBUG}\PYG{o}{=}yes\PYG{+w}{ }make\PYG{+w}{ }lib\PYGZhy{}colvars\PYG{+w}{ }\PYG{n+nv}{args}\PYG{o}{=}\PYG{l+s+s2}{\PYGZdq{}\PYGZhy{}m machine\PYGZdq{}}

\PYG{c+c1}{\PYGZsh{} Build without Lepton (included otherwise)}
\PYG{n+nv}{COLVARS\PYGZus{}LEPTON}\PYG{o}{=}no\PYG{+w}{ }make\PYG{+w}{ }lib\PYGZhy{}colvars\PYG{+w}{ }\PYG{n+nv}{args}\PYG{o}{=}\PYG{l+s+s2}{\PYGZdq{}\PYGZhy{}m machine\PYGZdq{}}
\end{sphinxVerbatim}

\sphinxAtStartPar
The build should produce two files: the library
\sphinxcode{\sphinxupquote{lib/colvars/libcolvars.a}} and the specification file
\sphinxcode{\sphinxupquote{lib/colvars/Makefile.lammps}}.  The latter is auto\sphinxhyphen{}generated,
and normally does not need to be edited.
\end{sphinxadmonition}
\end{quote}


\bigskip\hrule\bigskip



\subsection{ELECTRODE package}
\label{\detokenize{Build_extras:electrode-package}}\label{\detokenize{Build_extras:electrode}}
\sphinxAtStartPar
This package depends on the KSPACE package.
\begin{quote}

\begin{sphinxadmonition}{note}{CMake build}

\begin{sphinxVerbatim}[commandchars=\\\{\}]
\PYGZhy{}D\PYG{+w}{ }\PYG{n+nv}{PKG\PYGZus{}ELECTRODE}\PYG{o}{=}yes\PYG{+w}{          }\PYG{c+c1}{\PYGZsh{} enable the package itself}
\PYGZhy{}D\PYG{+w}{ }\PYG{n+nv}{PKG\PYGZus{}KSPACE}\PYG{o}{=}yes\PYG{+w}{             }\PYG{c+c1}{\PYGZsh{} the ELECTRODE package requires KSPACE}
\PYGZhy{}D\PYG{+w}{ }\PYG{n+nv}{USE\PYGZus{}INTERNAL\PYGZus{}LINALG}\PYG{o}{=}value\PYG{+w}{  }\PYG{c+c1}{\PYGZsh{}}
\end{sphinxVerbatim}

\sphinxAtStartPar
Features in the ELECTRODE package are dependent on code in the
KSPACE package so the latter one \sphinxstyleemphasis{must} be enabled.

\sphinxAtStartPar
The ELECTRODE package also requires LAPACK (and BLAS) and CMake
can identify their locations and pass that info to the ELECTRODE
build script.  But on some systems this may cause problems when
linking or the dependency is not desired.  Try enabling
\sphinxcode{\sphinxupquote{USE\_INTERNAL\_LINALG}} in those cases to use the bundled linear
algebra library and work around the limitation.
\end{sphinxadmonition}

\begin{sphinxadmonition}{note}{Traditional make}

\sphinxAtStartPar
Before building LAMMPS, you must configure the ELECTRODE support
libraries and settings in \sphinxcode{\sphinxupquote{lib/electrode}}.  You can do this
manually, if you prefer, or do it in one step from the
\sphinxcode{\sphinxupquote{lammps/src}} dir, using a command like these, which simply
invokes the \sphinxcode{\sphinxupquote{lib/electrode/Install.py}} script with the specified
args:

\begin{sphinxVerbatim}[commandchars=\\\{\}]
\PYG{c+c1}{\PYGZsh{} print help message}
make\PYG{+w}{ }lib\PYGZhy{}electrode

\PYG{c+c1}{\PYGZsh{} build with GNU g++ compiler and MPI STUBS (settings as with \PYGZdq{}make serial\PYGZdq{})}
make\PYG{+w}{ }lib\PYGZhy{}electrode\PYG{+w}{ }\PYG{n+nv}{args}\PYG{o}{=}\PYG{l+s+s2}{\PYGZdq{}\PYGZhy{}m serial\PYGZdq{}}

\PYG{c+c1}{\PYGZsh{} build with default MPI compiler (settings as with \PYGZdq{}make mpi\PYGZdq{})}
make\PYG{+w}{ }lib\PYGZhy{}electrode\PYG{+w}{ }\PYG{n+nv}{args}\PYG{o}{=}\PYG{l+s+s2}{\PYGZdq{}\PYGZhy{}m mpi\PYGZdq{}}
\end{sphinxVerbatim}

\sphinxAtStartPar
Note that the \sphinxcode{\sphinxupquote{Makefile.lammps}} file has settings for the BLAS
and LAPACK linear algebra libraries.  These can either exist on
your system, or you can use the files provided in \sphinxcode{\sphinxupquote{lib/linalg}}.
In the latter case you also need to build the library in
\sphinxcode{\sphinxupquote{lib/linalg}} with a command like these:

\begin{sphinxVerbatim}[commandchars=\\\{\}]
\PYG{c+c1}{\PYGZsh{} print help message}
make\PYG{+w}{ }lib\PYGZhy{}linalg

\PYG{c+c1}{\PYGZsh{} build with GNU C++ compiler (settings as with \PYGZdq{}make serial\PYGZdq{})}
make\PYG{+w}{ }lib\PYGZhy{}linalg\PYG{+w}{ }\PYG{n+nv}{args}\PYG{o}{=}\PYG{l+s+s2}{\PYGZdq{}\PYGZhy{}m serial\PYGZdq{}}

\PYG{c+c1}{\PYGZsh{} build with default MPI C++ compiler (settings as with \PYGZdq{}make mpi\PYGZdq{})}
make\PYG{+w}{ }lib\PYGZhy{}linalg\PYG{+w}{ }\PYG{n+nv}{args}\PYG{o}{=}\PYG{l+s+s2}{\PYGZdq{}\PYGZhy{}m mpi\PYGZdq{}}

\PYG{c+c1}{\PYGZsh{} build with GNU C++ compiler}
make\PYG{+w}{ }lib\PYGZhy{}linalg\PYG{+w}{ }\PYG{n+nv}{args}\PYG{o}{=}\PYG{l+s+s2}{\PYGZdq{}\PYGZhy{}m g++\PYGZdq{}}
\end{sphinxVerbatim}

\sphinxAtStartPar
The package itself is activated with \sphinxcode{\sphinxupquote{make yes\sphinxhyphen{}KSPACE}} and
\sphinxcode{\sphinxupquote{make yes\sphinxhyphen{}ELECTRODE}}
\end{sphinxadmonition}
\end{quote}


\bigskip\hrule\bigskip



\subsection{ML\sphinxhyphen{}PACE package}
\label{\detokenize{Build_extras:ml-pace-package}}\label{\detokenize{Build_extras:ml-pace}}
\sphinxAtStartPar
This package requires a library that can be downloaded and built
in lib/pace or somewhere else, which must be done before building
LAMMPS with this package. The code for the library can be found
at: \sphinxurl{https://github.com/ICAMS/lammps-user-pace/}

\sphinxAtStartPar
Instead of including the ML\sphinxhyphen{}PACE package directly into LAMMPS, it
is also possible to skip this step and build the ML\sphinxhyphen{}PACE package as
a plugin using the CMake script files in the \sphinxcode{\sphinxupquote{examples/PACKAGE/pace/plugin}}
folder and then load this plugin at runtime with the {\hyperref[\detokenize{plugin::doc}]{\sphinxcrossref{\DUrole{doc}{plugin command}}}}.
\begin{quote}

\begin{sphinxadmonition}{note}{CMake build}

\sphinxAtStartPar
By default the library will be downloaded from the git repository
and built automatically when the ML\sphinxhyphen{}PACE package is enabled with
\sphinxcode{\sphinxupquote{\sphinxhyphen{}D PKG\_ML\sphinxhyphen{}PACE=yes}}.  The location for the sources may be
customized by setting the variable \sphinxcode{\sphinxupquote{PACELIB\_URL}} when
configuring with CMake (e.g. to use a local archive on machines
without internet access).  Since CMake checks the validity of the
archive with \sphinxcode{\sphinxupquote{md5sum}} you may also need to set \sphinxcode{\sphinxupquote{PACELIB\_MD5}}
if you provide a different library version than what is downloaded
automatically.
\end{sphinxadmonition}

\begin{sphinxadmonition}{note}{Traditional make}

\sphinxAtStartPar
You can download and build the ML\sphinxhyphen{}PACE library
in one step from the \sphinxcode{\sphinxupquote{lammps/src}} dir, using these commands,
which invoke the \sphinxcode{\sphinxupquote{lib/pace/Install.py}} script.

\begin{sphinxVerbatim}[commandchars=\\\{\}]
\PYG{c+c1}{\PYGZsh{} print help message}
make\PYG{+w}{ }lib\PYGZhy{}pace

\PYG{c+c1}{\PYGZsh{} download and build the default version in lib/pace}
make\PYG{+w}{ }lib\PYGZhy{}pace\PYG{+w}{ }\PYG{n+nv}{args}\PYG{o}{=}\PYG{l+s+s2}{\PYGZdq{}\PYGZhy{}b\PYGZdq{}}
\end{sphinxVerbatim}

\sphinxAtStartPar
You should not need to edit the \sphinxcode{\sphinxupquote{lib/pace/Makefile.lammps}} file.
\end{sphinxadmonition}
\end{quote}


\bigskip\hrule\bigskip



\subsection{ML\sphinxhyphen{}POD package}
\label{\detokenize{Build_extras:ml-pod-package}}\label{\detokenize{Build_extras:ml-pod}}\begin{quote}

\begin{sphinxadmonition}{note}{CMake build}

\sphinxAtStartPar
No additional settings are needed besides \sphinxcode{\sphinxupquote{\sphinxhyphen{}D PKG\_ML\sphinxhyphen{}POD=yes}}.
\end{sphinxadmonition}

\begin{sphinxadmonition}{note}{Traditional make}

\sphinxAtStartPar
Before building LAMMPS, you must configure the ML\sphinxhyphen{}POD support
settings in \sphinxcode{\sphinxupquote{lib/mlpod}}.  You can do this manually, if you
prefer, or do it in one step from the \sphinxcode{\sphinxupquote{lammps/src}} dir, using a
command like the following, which simply invoke the
\sphinxcode{\sphinxupquote{lib/mlpod/Install.py}} script with the specified args:

\begin{sphinxVerbatim}[commandchars=\\\{\}]
\PYG{c+c1}{\PYGZsh{} print help message}
make\PYG{+w}{ }lib\PYGZhy{}mlpod

\PYG{c+c1}{\PYGZsh{} build with GNU g++ compiler and MPI STUBS (settings as with \PYGZdq{}make serial\PYGZdq{})}
make\PYG{+w}{ }lib\PYGZhy{}mlpod\PYG{+w}{ }\PYG{n+nv}{args}\PYG{o}{=}\PYG{l+s+s2}{\PYGZdq{}\PYGZhy{}m serial\PYGZdq{}}

\PYG{c+c1}{\PYGZsh{} build with default MPI compiler (settings as with \PYGZdq{}make mpi\PYGZdq{})}
make\PYG{+w}{ }lib\PYGZhy{}mlpod\PYG{+w}{ }\PYG{n+nv}{args}\PYG{o}{=}\PYG{l+s+s2}{\PYGZdq{}\PYGZhy{}m mpi\PYGZdq{}}

\PYG{c+c1}{\PYGZsh{} same as above but use the bundled linalg lib}
make\PYG{+w}{ }lib\PYGZhy{}mlpod\PYG{+w}{ }\PYG{n+nv}{args}\PYG{o}{=}\PYG{l+s+s2}{\PYGZdq{}\PYGZhy{}m mpi \PYGZhy{}e linalg\PYGZdq{}}
\end{sphinxVerbatim}

\sphinxAtStartPar
Note that the \sphinxcode{\sphinxupquote{Makefile.lammps}} file has settings to use the BLAS
and LAPACK linear algebra libraries.  These can either exist on
your system, or you can use the files provided in \sphinxcode{\sphinxupquote{lib/linalg}}.
In the latter case you also need to build the library in
\sphinxcode{\sphinxupquote{lib/linalg}} with a command like these:

\begin{sphinxVerbatim}[commandchars=\\\{\}]
\PYG{c+c1}{\PYGZsh{} print help message}
make\PYG{+w}{ }lib\PYGZhy{}linalg

\PYG{c+c1}{\PYGZsh{} build with GNU C++ compiler (settings as with \PYGZdq{}make serial\PYGZdq{})}
make\PYG{+w}{ }lib\PYGZhy{}linalg\PYG{+w}{ }\PYG{n+nv}{args}\PYG{o}{=}\PYG{l+s+s2}{\PYGZdq{}\PYGZhy{}m serial\PYGZdq{}}

\PYG{c+c1}{\PYGZsh{} build with default MPI C++ compiler (settings as with \PYGZdq{}make mpi\PYGZdq{})}
make\PYG{+w}{ }lib\PYGZhy{}linalg\PYG{+w}{ }\PYG{n+nv}{args}\PYG{o}{=}\PYG{l+s+s2}{\PYGZdq{}\PYGZhy{}m mpi\PYGZdq{}}

\PYG{c+c1}{\PYGZsh{} build with GNU C++ compiler}
make\PYG{+w}{ }lib\PYGZhy{}linalg\PYG{+w}{ }\PYG{n+nv}{args}\PYG{o}{=}\PYG{l+s+s2}{\PYGZdq{}\PYGZhy{}m g++\PYGZdq{}}
\end{sphinxVerbatim}

\sphinxAtStartPar
The package itself is activated with \sphinxcode{\sphinxupquote{make yes\sphinxhyphen{}ML\sphinxhyphen{}POD}}.
\end{sphinxadmonition}
\end{quote}


\bigskip\hrule\bigskip



\subsection{ML\sphinxhyphen{}QUIP package}
\label{\detokenize{Build_extras:ml-quip-package}}\label{\detokenize{Build_extras:ml-quip}}
\sphinxAtStartPar
To build with this package, you must download and build the QUIP
library.  It can be obtained from GitHub.  For support of GAP
potentials, additional files with specific licensing conditions need
to be downloaded and configured.  The automatic download will from
within CMake will download the non\sphinxhyphen{}commercial use version.
\begin{quote}

\begin{sphinxadmonition}{note}{CMake build}

\begin{sphinxVerbatim}[commandchars=\\\{\}]
\PYGZhy{}D\PYG{+w}{ }\PYG{n+nv}{DOWNLOAD\PYGZus{}QUIP}\PYG{o}{=}value\PYG{+w}{       }\PYG{c+c1}{\PYGZsh{} download QUIP library for build}
\PYG{+w}{                             }\PYG{c+c1}{\PYGZsh{} value = no (default) or yes}
\PYGZhy{}D\PYG{+w}{ }\PYG{n+nv}{QUIP\PYGZus{}LIBRARY}\PYG{o}{=}path\PYG{+w}{         }\PYG{c+c1}{\PYGZsh{} path to libquip.a}
\PYG{+w}{                             }\PYG{c+c1}{\PYGZsh{} (only needed if a custom location)}
\PYGZhy{}D\PYG{+w}{ }\PYG{n+nv}{USE\PYGZus{}INTERNAL\PYGZus{}LINALG}\PYG{o}{=}value\PYG{+w}{ }\PYG{c+c1}{\PYGZsh{} Use the internal linear algebra library}
\PYG{+w}{                             }\PYG{c+c1}{\PYGZsh{} instead of LAPACK}
\PYG{+w}{                             }\PYG{c+c1}{\PYGZsh{} value = no (default) or yes}
\end{sphinxVerbatim}

\sphinxAtStartPar
CMake will try to download and build the QUIP library from GitHub,
if it is not found on the local machine. This requires to have git
installed. It will use the same compilers and flags as used for
compiling LAMMPS.  Currently this is only supported for the GNU
and the Intel compilers. Set the \sphinxcode{\sphinxupquote{QUIP\_LIBRARY}} variable if you
want to use a previously compiled and installed QUIP library and
CMake cannot find it.

\sphinxAtStartPar
The QUIP library requires LAPACK (and BLAS) and CMake can identify
their locations and pass that info to the QUIP build script. But
on some systems this triggers a (current) limitation of CMake and
the configuration will fail. Try enabling \sphinxcode{\sphinxupquote{USE\_INTERNAL\_LINALG}} in
those cases to use the bundled linear algebra library and work around
the limitation.
\end{sphinxadmonition}

\begin{sphinxadmonition}{note}{Traditional make}

\sphinxAtStartPar
The download/build procedure for the QUIP library, described in
\sphinxcode{\sphinxupquote{lib/quip/README}} file requires setting two environment
variables, \sphinxcode{\sphinxupquote{QUIP\_ROOT}} and \sphinxcode{\sphinxupquote{QUIP\_ARCH}}.  These are accessed by
the \sphinxcode{\sphinxupquote{lib/quip/Makefile.lammps}} file which is used when you
compile and link LAMMPS with this package.  You should only need
to edit \sphinxcode{\sphinxupquote{Makefile.lammps}} if the LAMMPS build can not use its
settings to successfully build on your system.
\end{sphinxadmonition}
\end{quote}


\bigskip\hrule\bigskip



\subsection{PLUMED package}
\label{\detokenize{Build_extras:plumed-package}}\label{\detokenize{Build_extras:plumed}}
\sphinxAtStartPar
Before building LAMMPS with this package, you must first build PLUMED.
PLUMED can be built as part of the LAMMPS build or installed separately
from LAMMPS using the generic \sphinxhref{https://plumed.github.io/doc-master/user-doc/html/\_installation.html}{PLUMED installation instructions}.
The PLUMED package has been tested to work with Plumed versions
2.4.x, 2.5.x, and 2.6.x and will error out, when trying to run calculations
with a different version of the Plumed kernel.

\sphinxAtStartPar
PLUMED can be linked into MD codes in three different modes: static,
shared, and runtime.  With the “static” mode, all the code that PLUMED
requires is linked statically into LAMMPS. LAMMPS is then fully
independent from the PLUMED installation, but you have to rebuild/relink
it in order to update the PLUMED code inside it.  With the “shared”
linkage mode, LAMMPS is linked to a shared library that contains the
PLUMED code.  This library should preferably be installed in a globally
accessible location. When PLUMED is linked in this way the same library
can be used by multiple MD packages.  Furthermore, the PLUMED library
LAMMPS uses can be updated without the need for a recompile of LAMMPS
for as long as the shared PLUMED library is ABI\sphinxhyphen{}compatible.

\sphinxAtStartPar
The third linkage mode is “runtime” which allows the user to specify
which PLUMED kernel should be used at runtime by using the PLUMED\_KERNEL
environment variable. This variable should point to the location of the
libplumedKernel.so dynamical shared object, which is then loaded at
runtime. This mode of linking is particularly convenient for doing
PLUMED development and comparing multiple PLUMED versions as these sorts
of comparisons can be done without recompiling the hosting MD code. All
three linkage modes are supported by LAMMPS on selected operating
systems (e.g. Linux) and using either CMake or traditional make
build. The “static” mode should be the most portable, while the
“runtime” mode support in LAMMPS makes the most assumptions about
operating system and compiler environment. If one mode does not work,
try a different one, switch to a different build system, consider a
global PLUMED installation or consider downloading PLUMED during the
LAMMPS build.

\sphinxAtStartPar
Instead of including the PLUMED package directly into LAMMPS, it
is also possible to skip this step and build the PLUMED package as
a plugin using the CMake script files in the \sphinxcode{\sphinxupquote{examples/PACKAGE/plumed/plugin}}
folder and then load this plugin at runtime with the {\hyperref[\detokenize{plugin::doc}]{\sphinxcrossref{\DUrole{doc}{plugin command}}}}.
\begin{quote}

\begin{sphinxadmonition}{note}{CMake build}

\sphinxAtStartPar
When the \sphinxcode{\sphinxupquote{\sphinxhyphen{}D PKG\_PLUMED=yes}} flag is included in the cmake
command you must ensure that \sphinxtitleref{the GNU Scientific Library (GSL)
\textless{}https://www.gnu.org/software/gsl/\textgreater{}} is installed in locations
that are accessible in your environment.  There are then two
additional variables that control the manner in which PLUMED is
obtained and linked into LAMMPS.

\begin{sphinxVerbatim}[commandchars=\\\{\}]
\PYGZhy{}D\PYG{+w}{ }\PYG{n+nv}{DOWNLOAD\PYGZus{}PLUMED}\PYG{o}{=}value\PYG{+w}{   }\PYG{c+c1}{\PYGZsh{} download PLUMED for build}
\PYG{+w}{                           }\PYG{c+c1}{\PYGZsh{} value = no (default) or yes}
\PYGZhy{}D\PYG{+w}{ }\PYG{n+nv}{PLUMED\PYGZus{}MODE}\PYG{o}{=}value\PYG{+w}{       }\PYG{c+c1}{\PYGZsh{} Linkage mode for PLUMED}
\PYG{+w}{                           }\PYG{c+c1}{\PYGZsh{} value = static (default), shared,}
\PYG{+w}{                           }\PYG{c+c1}{\PYGZsh{}         or runtime}
\end{sphinxVerbatim}

\sphinxAtStartPar
If \sphinxcode{\sphinxupquote{DOWNLOAD\_PLUMED}} is set to \sphinxcode{\sphinxupquote{yes}}, the PLUMED library will be
downloaded (the version of PLUMED that will be downloaded is
hard\sphinxhyphen{}coded to a vetted version of PLUMED, usually a recent stable
release version) and built inside the CMake build directory.  If
\sphinxcode{\sphinxupquote{DOWNLOAD\_PLUMED}} is set to “no” (the default), CMake will try
to detect and link to an installed version of PLUMED.  For this to
work, the PLUMED library has to be installed into a location where
the \sphinxcode{\sphinxupquote{pkg\sphinxhyphen{}config}} tool can find it or the \sphinxcode{\sphinxupquote{PKG\_CONFIG\_PATH}}
environment variable has to be set up accordingly.  PLUMED should
be installed in such a location if you compile it using the
default make; make install commands.

\sphinxAtStartPar
The \sphinxcode{\sphinxupquote{PLUMED\_MODE}} setting determines the linkage mode for the
PLUMED library.  The allowed values for this flag are “static”
(default), “shared”, or “runtime”.  If you want to switch the
linkage mode, just re\sphinxhyphen{}run CMake with a different setting. For a
discussion of PLUMED linkage modes, please see above.  When
\sphinxcode{\sphinxupquote{DOWNLOAD\_PLUMED}} is enabled the static linkage mode is
recommended.
\end{sphinxadmonition}

\begin{sphinxadmonition}{note}{Traditional make}

\sphinxAtStartPar
PLUMED needs to be installed before the PLUMED package is
installed so that LAMMPS can find the right settings when
compiling and linking the LAMMPS executable.  You can either
download and build PLUMED inside the LAMMPS plumed library folder
or use a previously installed PLUMED library and point LAMMPS to
its location. You also have to choose the linkage mode: “static”
(default), “shared” or “runtime”.  For a discussion of PLUMED
linkage modes, please see above.

\sphinxAtStartPar
Download/compilation/configuration of the plumed library can be done
from the src folder through the following make args:

\begin{sphinxVerbatim}[commandchars=\\\{\}]
\PYG{c+c1}{\PYGZsh{} print help message}
make\PYG{+w}{ }lib\PYGZhy{}plumed

\PYG{c+c1}{\PYGZsh{} download and build PLUMED in lib/plumed/plumed2}
make\PYG{+w}{ }lib\PYGZhy{}plumed\PYG{+w}{ }\PYG{n+nv}{args}\PYG{o}{=}\PYG{l+s+s2}{\PYGZdq{}\PYGZhy{}b\PYGZdq{}}

\PYG{c+c1}{\PYGZsh{} use existing PLUMED installation in \PYGZdl{}HOME/.local}
make\PYG{+w}{ }lib\PYGZhy{}plumed\PYG{+w}{ }\PYG{n+nv}{args}\PYG{o}{=}\PYG{l+s+s2}{\PYGZdq{}}\PYG{l+s+s2}{\PYGZhy{}p }\PYG{n+nv}{\PYGZdl{}HOME}\PYG{l+s+s2}{/.local}\PYG{l+s+s2}{\PYGZdq{}}

\PYG{c+c1}{\PYGZsh{} use existing PLUMED installation in /usr/local and}
\PYG{c+c1}{\PYGZsh{} use shared linkage mode}
make\PYG{+w}{ }lib\PYGZhy{}plumed\PYG{+w}{ }\PYG{n+nv}{args}\PYG{o}{=}\PYG{l+s+s2}{\PYGZdq{}\PYGZhy{}p /usr/local \PYGZhy{}m shared\PYGZdq{}}
\end{sphinxVerbatim}

\sphinxAtStartPar
Note that two symbolic (soft) links, \sphinxcode{\sphinxupquote{includelink}} and \sphinxcode{\sphinxupquote{liblink}}
are created in \sphinxcode{\sphinxupquote{lib/plumed}} that point to the location of the PLUMED
build to use. A new file \sphinxcode{\sphinxupquote{lib/plumed/Makefile.lammps}} is also
created with settings suitable for LAMMPS to compile and link
PLUMED using the desired linkage mode. After this step is
completed, you can install the PLUMED package and compile
LAMMPS in the usual manner:

\begin{sphinxVerbatim}[commandchars=\\\{\}]
make\PYG{+w}{ }yes\PYGZhy{}plumed
make\PYG{+w}{ }machine
\end{sphinxVerbatim}

\sphinxAtStartPar
Once this compilation completes you should be able to run LAMMPS
in the usual way.  For shared linkage mode, libplumed.so must be
found by the LAMMPS executable, which on many operating systems
means, you have to set the \sphinxcode{\sphinxupquote{LD\_LIBRARY\_PATH}} environment variable
accordingly.

\sphinxAtStartPar
Support for the different linkage modes in LAMMPS varies for
different operating systems, using the static linkage is expected
to be the most portable, and thus set to be the default.

\sphinxAtStartPar
If you want to change the linkage mode, you have to re\sphinxhyphen{}run \sphinxcode{\sphinxupquote{make
lib\sphinxhyphen{}plumed}} with the desired settings \sphinxstylestrong{and} do a re\sphinxhyphen{}install if
the PLUMED package with \sphinxcode{\sphinxupquote{make yes\sphinxhyphen{}plumed}} to update the
required makefile settings with the changes in the \sphinxcode{\sphinxupquote{lib/plumed}}
folder.
\end{sphinxadmonition}
\end{quote}


\bigskip\hrule\bigskip



\subsection{H5MD package}
\label{\detokenize{Build_extras:h5md-package}}\label{\detokenize{Build_extras:h5md}}
\sphinxAtStartPar
To build with this package you must have the HDF5 software package
installed on your system, which should include the h5cc compiler and
the HDF5 library.
\begin{quote}

\begin{sphinxadmonition}{note}{CMake build}

\sphinxAtStartPar
No additional settings are needed besides \sphinxcode{\sphinxupquote{\sphinxhyphen{}D PKG\_H5MD=yes}}.

\sphinxAtStartPar
This should auto\sphinxhyphen{}detect the H5MD library on your system.  Several
advanced CMake H5MD options exist if you need to specify where it
is installed.  Use the ccmake (terminal window) or cmake\sphinxhyphen{}gui
(graphical) tools to see these options and set them interactively
from their user interfaces.
\end{sphinxadmonition}

\begin{sphinxadmonition}{note}{Traditional make}

\sphinxAtStartPar
Before building LAMMPS, you must build the CH5MD library in
\sphinxcode{\sphinxupquote{lib/h5md}}.  You can do this manually if you prefer; follow the
instructions in \sphinxcode{\sphinxupquote{lib/h5md/README}}.  You can also do it in one
step from the \sphinxcode{\sphinxupquote{lammps/src}} dir, using a command like these,
which simply invokes the \sphinxcode{\sphinxupquote{lib/h5md/Install.py}} script with the
specified args:

\begin{sphinxVerbatim}[commandchars=\\\{\}]
make\PYG{+w}{ }lib\PYGZhy{}h5md\PYG{+w}{                     }\PYG{c+c1}{\PYGZsh{} print help message}
make\PYG{+w}{ }lib\PYGZhy{}h5md\PYG{+w}{ }\PYG{n+nv}{args}\PYG{o}{=}\PYG{l+s+s2}{\PYGZdq{}\PYGZhy{}m h5cc\PYGZdq{}}\PYG{+w}{      }\PYG{c+c1}{\PYGZsh{} build with h5cc compiler}
\end{sphinxVerbatim}

\sphinxAtStartPar
The build should produce two files: \sphinxcode{\sphinxupquote{lib/h5md/libch5md.a}} and
\sphinxcode{\sphinxupquote{lib/h5md/Makefile.lammps}}.  The latter is copied from an
existing \sphinxcode{\sphinxupquote{Makefile.lammps.*}} and has settings needed to build
LAMMPS with the system HDF5 library.  If necessary, you can
edit/create a new \sphinxcode{\sphinxupquote{lib/h5md/Makefile.machine}} file for your
system, which should define an EXTRAMAKE variable to specify a
corresponding \sphinxcode{\sphinxupquote{Makefile.lammps.\textless{}machine\textgreater{}}} file.
\end{sphinxadmonition}
\end{quote}


\bigskip\hrule\bigskip



\subsection{ML\sphinxhyphen{}HDNNP package}
\label{\detokenize{Build_extras:ml-hdnnp-package}}\label{\detokenize{Build_extras:ml-hdnnp}}
\sphinxAtStartPar
To build with the ML\sphinxhyphen{}HDNNP package it is required to download and build the
external \sphinxhref{https://github.com/CompPhysVienna/n2p2}{n2p2} library \sphinxcode{\sphinxupquote{v2.1.4}}
(or higher). The LAMMPS build process offers an automatic download and
compilation of \sphinxstyleemphasis{n2p2} or allows you to choose the installation directory of
\sphinxstyleemphasis{n2p2} manually. Please see the boxes below for the CMake and traditional build
system for detailed information.

\sphinxAtStartPar
In case of a manual installation of \sphinxstyleemphasis{n2p2} you only need to build the \sphinxstyleemphasis{n2p2} core
library \sphinxcode{\sphinxupquote{libnnp}} and interface library \sphinxcode{\sphinxupquote{libnnpif}}. When using GCC it should
suffice to execute \sphinxcode{\sphinxupquote{make libnnpif}} in the \sphinxstyleemphasis{n2p2} \sphinxcode{\sphinxupquote{src}} directory. For more
details please see \sphinxcode{\sphinxupquote{lib/hdnnp/README}} and the \sphinxhref{https://compphysvienna.github.io/n2p2/topics/build.html}{n2p2 build documentation}.
\begin{quote}

\begin{sphinxadmonition}{note}{CMake build}

\begin{sphinxVerbatim}[commandchars=\\\{\}]
\PYGZhy{}D\PYG{+w}{ }\PYG{n+nv}{DOWNLOAD\PYGZus{}N2P2}\PYG{o}{=}value\PYG{+w}{    }\PYG{c+c1}{\PYGZsh{} download n2p2 for build}
\PYG{+w}{                          }\PYG{c+c1}{\PYGZsh{} value = no (default) or yes}
\PYGZhy{}D\PYG{+w}{ }\PYG{n+nv}{N2P2\PYGZus{}DIR}\PYG{o}{=}path\PYG{+w}{          }\PYG{c+c1}{\PYGZsh{} n2p2 base directory}
\PYG{+w}{                          }\PYG{c+c1}{\PYGZsh{} (only needed if a custom location)}
\end{sphinxVerbatim}

\sphinxAtStartPar
If \sphinxcode{\sphinxupquote{DOWNLOAD\_N2P2}} is set, the \sphinxstyleemphasis{n2p2} library will be downloaded and
built inside the CMake build directory.  If the \sphinxstyleemphasis{n2p2} library is already
on your system (in a location CMake cannot find it), set the \sphinxcode{\sphinxupquote{N2P2\_DIR}}
to path where \sphinxstyleemphasis{n2p2} is located. If \sphinxstyleemphasis{n2p2} is located directly in
\sphinxcode{\sphinxupquote{lib/hdnnp/n2p2}} it will be automatically found by CMake.
\end{sphinxadmonition}

\begin{sphinxadmonition}{note}{Traditional make}

\sphinxAtStartPar
You can download and build the \sphinxstyleemphasis{n2p2} library manually if you prefer;
follow the instructions in \sphinxcode{\sphinxupquote{lib/hdnnp/README}}. You can also do it in
one step from the \sphinxcode{\sphinxupquote{lammps/src}} dir, using a command like these, which
simply invokes the \sphinxcode{\sphinxupquote{lib/hdnnp/Install.py}} script with the specified args:

\begin{sphinxVerbatim}[commandchars=\\\{\}]
\PYG{c+c1}{\PYGZsh{} print help message}
make\PYG{+w}{ }lib\PYGZhy{}hdnnp

\PYG{c+c1}{\PYGZsh{} download and build in lib/hdnnp/n2p2\PYGZhy{}...}
make\PYG{+w}{ }lib\PYGZhy{}hdnnp\PYG{+w}{ }\PYG{n+nv}{args}\PYG{o}{=}\PYG{l+s+s2}{\PYGZdq{}\PYGZhy{}b\PYGZdq{}}

\PYG{c+c1}{\PYGZsh{} download and build specific version}
make\PYG{+w}{ }lib\PYGZhy{}hdnnp\PYG{+w}{ }\PYG{n+nv}{args}\PYG{o}{=}\PYG{l+s+s2}{\PYGZdq{}\PYGZhy{}b \PYGZhy{}v 2.1.4\PYGZdq{}}

\PYG{c+c1}{\PYGZsh{} use the existing n2p2 installation in /usr/local/n2p2}
make\PYG{+w}{ }lib\PYGZhy{}hdnnp\PYG{+w}{ }\PYG{n+nv}{args}\PYG{o}{=}\PYG{l+s+s2}{\PYGZdq{}\PYGZhy{}p /usr/local/n2p2\PYGZdq{}}
\end{sphinxVerbatim}

\sphinxAtStartPar
Note that three symbolic (soft) links, \sphinxcode{\sphinxupquote{includelink}}, \sphinxcode{\sphinxupquote{liblink}} and
\sphinxcode{\sphinxupquote{Makefile.lammps}}, will be created in \sphinxcode{\sphinxupquote{lib/hdnnp}} to point to
\sphinxcode{\sphinxupquote{n2p2/include}}, \sphinxcode{\sphinxupquote{n2p2/lib}} and \sphinxcode{\sphinxupquote{n2p2/lib/Makefile.lammps\sphinxhyphen{}extra}},
respectively. When LAMMPS is built in \sphinxcode{\sphinxupquote{src}} it will use these links.
\end{sphinxadmonition}
\end{quote}


\bigskip\hrule\bigskip



\subsection{INTEL package}
\label{\detokenize{Build_extras:intel-package}}\label{\detokenize{Build_extras:intel}}
\sphinxAtStartPar
To build with this package, you must choose which hardware you want to
build for, either x86 CPUs or Intel KNLs in offload mode.  You should
also typically {\hyperref[\detokenize{Build_extras:openmp}]{\sphinxcrossref{\DUrole{std,std-ref}{install the OPENMP package}}}}, as it can be
used in tandem with the INTEL package to good effect, as explained
on the {\hyperref[\detokenize{Speed_intel::doc}]{\sphinxcrossref{\DUrole{doc}{INTEL package}}}} page.

\sphinxAtStartPar
When using Intel compilers version 16.0 or later is required.  You can
also use the GNU or Clang compilers and they will provide performance
improvements over regular styles and OPENMP styles, but less so than
with the Intel compilers.  Please also note, that some compilers have
been found to apply memory alignment constraints incompletely or
incorrectly and thus can cause segmentation faults in otherwise correct
code when using features from the INTEL package.
\begin{quote}

\begin{sphinxadmonition}{note}{CMake build}

\begin{sphinxVerbatim}[commandchars=\\\{\}]
\PYGZhy{}D\PYG{+w}{ }\PYG{n+nv}{INTEL\PYGZus{}ARCH}\PYG{o}{=}value\PYG{+w}{     }\PYG{c+c1}{\PYGZsh{} value = cpu (default) or knl}
\PYGZhy{}D\PYG{+w}{ }\PYG{n+nv}{INTEL\PYGZus{}LRT\PYGZus{}MODE}\PYG{o}{=}value\PYG{+w}{ }\PYG{c+c1}{\PYGZsh{} value = threads, none, or c++11}
\end{sphinxVerbatim}
\end{sphinxadmonition}

\begin{sphinxadmonition}{note}{Traditional make}

\sphinxAtStartPar
Choose which hardware to compile for in Makefile.machine via the
following settings.  See \sphinxcode{\sphinxupquote{src/MAKE/OPTIONS/Makefile.intel\_cpu*}}
and \sphinxcode{\sphinxupquote{Makefile.knl}} files for examples. and
\sphinxcode{\sphinxupquote{src/INTEL/README}} for additional information.

\sphinxAtStartPar
For CPUs:

\begin{sphinxVerbatim}[commandchars=\\\{\}]
\PYG{n+nv}{OPTFLAGS}\PYG{+w}{ }\PYG{o}{=}\PYG{+w}{  }\PYGZhy{}xHost\PYG{+w}{ }\PYGZhy{}O2\PYG{+w}{ }\PYGZhy{}fp\PYGZhy{}model\PYG{+w}{ }\PYG{n+nv}{fast}\PYG{o}{=}\PYG{l+m}{2}\PYG{+w}{ }\PYGZhy{}no\PYGZhy{}prec\PYGZhy{}div\PYG{+w}{ }\PYGZhy{}qoverride\PYGZhy{}limits\PYG{+w}{ }\PYGZhy{}qopt\PYGZhy{}zmm\PYGZhy{}usage\PYG{o}{=}high
\PYG{n+nv}{CCFLAGS}\PYG{+w}{ }\PYG{o}{=}\PYG{+w}{   }\PYGZhy{}g\PYG{+w}{ }\PYGZhy{}qopenmp\PYG{+w}{ }\PYGZhy{}DLAMMPS\PYGZus{}MEMALIGN\PYG{o}{=}\PYG{l+m}{64}\PYG{+w}{ }\PYGZhy{}no\PYGZhy{}offload\PYG{+w}{ }\PYGZhy{}fno\PYGZhy{}alias\PYG{+w}{ }\PYGZhy{}ansi\PYGZhy{}alias\PYG{+w}{ }\PYGZhy{}restrict\PYG{+w}{ }\PYG{k}{\PYGZdl{}(}OPTFLAGS\PYG{k}{)}
\PYG{n+nv}{LINKFLAGS}\PYG{+w}{ }\PYG{o}{=}\PYG{+w}{ }\PYGZhy{}g\PYG{+w}{ }\PYGZhy{}qopenmp\PYG{+w}{ }\PYG{k}{\PYGZdl{}(}OPTFLAGS\PYG{k}{)}
\PYG{n+nv}{LIB}\PYG{+w}{ }\PYG{o}{=}\PYG{+w}{       }\PYGZhy{}ltbbmalloc
\end{sphinxVerbatim}

\sphinxAtStartPar
For KNLs:

\begin{sphinxVerbatim}[commandchars=\\\{\}]
\PYG{n+nv}{OPTFLAGS}\PYG{+w}{ }\PYG{o}{=}\PYG{+w}{  }\PYGZhy{}xMIC\PYGZhy{}AVX512\PYG{+w}{ }\PYGZhy{}O2\PYG{+w}{ }\PYGZhy{}fp\PYGZhy{}model\PYG{+w}{ }\PYG{n+nv}{fast}\PYG{o}{=}\PYG{l+m}{2}\PYG{+w}{ }\PYGZhy{}no\PYGZhy{}prec\PYGZhy{}div\PYG{+w}{ }\PYGZhy{}qoverride\PYGZhy{}limits
\PYG{n+nv}{CCFLAGS}\PYG{+w}{ }\PYG{o}{=}\PYG{+w}{   }\PYGZhy{}g\PYG{+w}{ }\PYGZhy{}qopenmp\PYG{+w}{ }\PYGZhy{}DLAMMPS\PYGZus{}MEMALIGN\PYG{o}{=}\PYG{l+m}{64}\PYG{+w}{ }\PYGZhy{}no\PYGZhy{}offload\PYG{+w}{ }\PYGZhy{}fno\PYGZhy{}alias\PYG{+w}{ }\PYGZhy{}ansi\PYGZhy{}alias\PYG{+w}{ }\PYGZhy{}restrict\PYG{+w}{ }\PYG{k}{\PYGZdl{}(}OPTFLAGS\PYG{k}{)}
\PYG{n+nv}{LINKFLAGS}\PYG{+w}{ }\PYG{o}{=}\PYG{+w}{ }\PYGZhy{}g\PYG{+w}{ }\PYGZhy{}qopenmp\PYG{+w}{ }\PYG{k}{\PYGZdl{}(}OPTFLAGS\PYG{k}{)}
\PYG{n+nv}{LIB}\PYG{+w}{ }\PYG{o}{=}\PYG{+w}{       }\PYGZhy{}ltbbmalloc
\end{sphinxVerbatim}
\end{sphinxadmonition}
\end{quote}

\sphinxAtStartPar
In Long\sphinxhyphen{}range thread mode (LRT) a modified verlet style is used, that
operates the Kspace calculation in a separate thread concurrently to
other calculations. This has to be enabled in the {\hyperref[\detokenize{package::doc}]{\sphinxcrossref{\DUrole{doc}{package intel}}}} command at runtime. With the setting “threads” it used the
pthreads library, while “c++11” will use the built\sphinxhyphen{}in thread support
of C++11 compilers. The option “none” skips compilation of this
feature. The default is to use “threads” if pthreads is available and
otherwise “none”.

\sphinxAtStartPar
Best performance is achieved with Intel hardware, Intel compilers, as
well as the Intel TBB and MKL libraries. However, the code also
compiles, links, and runs with other compilers / hardware and without
TBB and MKL.


\bigskip\hrule\bigskip



\subsection{MDI package}
\label{\detokenize{Build_extras:mdi-package}}\label{\detokenize{Build_extras:mdi}}\begin{quote}

\begin{sphinxadmonition}{note}{CMake build}

\begin{sphinxVerbatim}[commandchars=\\\{\}]
\PYGZhy{}D\PYG{+w}{ }\PYG{n+nv}{DOWNLOAD\PYGZus{}MDI}\PYG{o}{=}value\PYG{+w}{    }\PYG{c+c1}{\PYGZsh{} download MDI Library for build}
\PYG{+w}{                         }\PYG{c+c1}{\PYGZsh{} value = no (default) or yes}
\end{sphinxVerbatim}
\end{sphinxadmonition}

\begin{sphinxadmonition}{note}{Traditional make}

\sphinxAtStartPar
Before building LAMMPS, you must build the MDI Library in
\sphinxcode{\sphinxupquote{lib/mdi}}.  You can do this by executing a command like one
of the following from the \sphinxcode{\sphinxupquote{lib/mdi}} directory:

\begin{sphinxVerbatim}[commandchars=\\\{\}]
python\PYG{+w}{ }Install.py\PYG{+w}{ }\PYGZhy{}m\PYG{+w}{ }gcc\PYG{+w}{       }\PYG{c+c1}{\PYGZsh{} build using gcc compiler}
python\PYG{+w}{ }Install.py\PYG{+w}{ }\PYGZhy{}m\PYG{+w}{ }icc\PYG{+w}{       }\PYG{c+c1}{\PYGZsh{} build using icc compiler}
\end{sphinxVerbatim}

\sphinxAtStartPar
The build should produce two files: \sphinxcode{\sphinxupquote{lib/mdi/includelink/mdi.h}}
and \sphinxcode{\sphinxupquote{lib/mdi/liblink/libmdi.so}}.
\end{sphinxadmonition}
\end{quote}


\bigskip\hrule\bigskip



\subsection{MOLFILE package}
\label{\detokenize{Build_extras:molfile-package}}\label{\detokenize{Build_extras:molfile}}\begin{quote}

\begin{sphinxadmonition}{note}{CMake build}

\begin{sphinxVerbatim}[commandchars=\\\{\}]
\PYGZhy{}D\PYG{+w}{ }\PYG{n+nv}{MOLFILE\PYGZus{}INCLUDE\PYGZus{}DIR}\PYG{o}{=}path\PYG{+w}{   }\PYG{c+c1}{\PYGZsh{} (optional) path where VMD molfile}
\PYG{+w}{                              }\PYG{c+c1}{\PYGZsh{} plugin headers are installed}
\PYGZhy{}D\PYG{+w}{ }\PYG{n+nv}{PKG\PYGZus{}MOLFILE}\PYG{o}{=}yes
\end{sphinxVerbatim}

\sphinxAtStartPar
Using \sphinxcode{\sphinxupquote{\sphinxhyphen{}D PKG\_MOLFILE=yes}} enables the package, and setting
\sphinxcode{\sphinxupquote{\sphinxhyphen{}D MOLFILE\_INCLUDE\_DIR}} allows to provide a custom location for
the molfile plugin header files. These should match the ABI of the
plugin files used, and thus one typically sets them to include
folder of the local VMD installation in use. LAMMPS ships with a
couple of default header files that correspond to a popular VMD
version, usually the latest release.
\end{sphinxadmonition}

\begin{sphinxadmonition}{note}{Traditional make}

\sphinxAtStartPar
The \sphinxcode{\sphinxupquote{lib/molfile/Makefile.lammps}} file has a setting for a
dynamic loading library libdl.a that is typically present on all
systems.  It is required for LAMMPS to link with this package.  If
the setting is not valid for your system, you will need to edit
the Makefile.lammps file.  See \sphinxcode{\sphinxupquote{lib/molfile/README}} and
\sphinxcode{\sphinxupquote{lib/molfile/Makefile.lammps}} for details. It is also possible
to configure a different folder with the VMD molfile plugin header
files. LAMMPS ships with a couple of default headers, but these
are not compatible with all VMD versions, so it is often best to
change this setting to the location of the same include files of
the local VMD installation in use.
\end{sphinxadmonition}
\end{quote}


\bigskip\hrule\bigskip



\subsection{NETCDF package}
\label{\detokenize{Build_extras:netcdf-package}}\label{\detokenize{Build_extras:netcdf}}
\sphinxAtStartPar
To build with this package you must have the NetCDF library installed
on your system.
\begin{quote}

\begin{sphinxadmonition}{note}{CMake build}

\sphinxAtStartPar
No additional settings are needed besides \sphinxcode{\sphinxupquote{\sphinxhyphen{}D PKG\_NETCDF=yes}}.

\sphinxAtStartPar
This should auto\sphinxhyphen{}detect the NETCDF library if it is installed on
your system at standard locations.  Several advanced CMake NETCDF
options exist if you need to specify where it was installed.  Use
the \sphinxcode{\sphinxupquote{ccmake}} (terminal window) or \sphinxcode{\sphinxupquote{cmake\sphinxhyphen{}gui}} (graphical)
tools to see these options and set them interactively from their
user interfaces.
\end{sphinxadmonition}

\begin{sphinxadmonition}{note}{Traditional make}

\sphinxAtStartPar
The \sphinxcode{\sphinxupquote{lib/netcdf/Makefile.lammps}} file has settings for NetCDF
include and library files which LAMMPS needs to build with this
package.  If the settings are not valid for your system, you will
need to edit the \sphinxcode{\sphinxupquote{Makefile.lammps}} file.  See
\sphinxcode{\sphinxupquote{lib/netcdf/README}} for details.
\end{sphinxadmonition}
\end{quote}


\bigskip\hrule\bigskip



\subsection{OPENMP package}
\label{\detokenize{Build_extras:openmp-package}}\label{\detokenize{Build_extras:openmp}}\begin{quote}

\begin{sphinxadmonition}{note}{CMake build}

\sphinxAtStartPar
No additional settings are required besides \sphinxcode{\sphinxupquote{\sphinxhyphen{}D
PKG\_OPENMP=yes}}.  If CMake detects OpenMP compiler support, the
OPENMP code will be compiled with multi\sphinxhyphen{}threading support
enabled, otherwise as optimized serial code.
\end{sphinxadmonition}

\begin{sphinxadmonition}{note}{Traditional make}

\sphinxAtStartPar
To enable multi\sphinxhyphen{}threading support in the OPENMP package (and
other styles supporting OpenMP) the following compile and link
flags must be added to your Makefile.machine file.  See
\sphinxcode{\sphinxupquote{src/MAKE/OPTIONS/Makefile.omp}} for an example.

\begin{sphinxVerbatim}[commandchars=\\\{\}]
CCFLAGS: \PYGZhy{}fopenmp               \PYGZsh{} for GNU and Clang Compilers
CCFLAGS: \PYGZhy{}qopenmp \PYGZhy{}restrict     \PYGZsh{} for Intel compilers on Linux
LINKFLAGS: \PYGZhy{}fopenmp             \PYGZsh{} for GNU and Clang Compilers
LINKFLAGS: \PYGZhy{}qopenmp             \PYGZsh{} for Intel compilers on Linux
\end{sphinxVerbatim}

\sphinxAtStartPar
For other platforms and compilers, please consult the
documentation about OpenMP support for your compiler.
\end{sphinxadmonition}
\end{quote}

\begin{sphinxadmonition}{note}{Adding OpenMP support on macOS}

\sphinxAtStartPar
Apple offers the \sphinxhref{https://developer.apple.com/xcode/}{Xcode package and IDE} for compiling software on
macOS, so you have likely installed it to compile LAMMPS.  Their
compiler is based on \sphinxhref{https://clang.llvm.org/}{Clang}, but while it
is capable of processing OpenMP directives, the necessary header
files and OpenMP runtime library are missing.  The \sphinxhref{https://www.r-project.org/}{R developers} have figured out a way to build those
in a compatible fashion. One can download them from
\sphinxurl{https://mac.r-project.org/openmp/}.  Simply adding those files as
instructed enables the Xcode C++ compiler to compile LAMMPS with \sphinxcode{\sphinxupquote{\sphinxhyphen{}D
BUILD\_OMP=yes}}.
\end{sphinxadmonition}


\bigskip\hrule\bigskip



\subsection{QMMM package}
\label{\detokenize{Build_extras:qmmm-package}}\label{\detokenize{Build_extras:qmmm}}
\sphinxAtStartPar
For using LAMMPS to do QM/MM simulations via the QMMM package you
need to build LAMMPS as a library.  A LAMMPS executable with {\hyperref[\detokenize{fix_qmmm::doc}]{\sphinxcrossref{\DUrole{doc}{fix
qmmm}}}} included can be built, but will not be able to do a
QM/MM simulation on as such.  You must also build a QM code \sphinxhyphen{} currently
only Quantum ESPRESSO (QE) is supported \sphinxhyphen{} and create a new executable
which links LAMMPS and the QM code together.  Details are given in the
\sphinxcode{\sphinxupquote{lib/qmmm/README}} file.  It is also recommended to read the
instructions for {\hyperref[\detokenize{Build_link::doc}]{\sphinxcrossref{\DUrole{doc}{linking with LAMMPS as a library}}}}
for background information.  This requires compatible Quantum Espresso
and LAMMPS versions.  The current interface and makefiles have last been
verified to work in February 2020 with Quantum Espresso versions 6.3 to
6.5.
\begin{quote}

\begin{sphinxadmonition}{note}{CMake build}

\sphinxAtStartPar
When using CMake, building a LAMMPS library is required and it is
recommended to build a shared library, since any libraries built
from the sources in the \sphinxstyleemphasis{lib} folder (including the essential
libqmmm.a) are not included in the static LAMMPS library and
(currently) not installed, while their code is included in the
shared LAMMPS library.  Thus a typical command line to configure
building LAMMPS for QMMM would be:

\begin{sphinxVerbatim}[commandchars=\\\{\}]
cmake\PYG{+w}{ }\PYGZhy{}C\PYG{+w}{ }../cmake/presets/basic.cmake\PYG{+w}{ }\PYGZhy{}D\PYG{+w}{ }\PYG{n+nv}{PKG\PYGZus{}QMMM}\PYG{o}{=}yes\PYG{+w}{ }\PYG{l+s+se}{\PYGZbs{}}
\PYG{+w}{    }\PYGZhy{}D\PYG{+w}{ }\PYG{n+nv}{BUILD\PYGZus{}LIB}\PYG{o}{=}yes\PYG{+w}{ }\PYGZhy{}DBUILD\PYGZus{}SHARED\PYGZus{}LIBS\PYG{o}{=}yes\PYG{+w}{ }../cmake
\end{sphinxVerbatim}

\sphinxAtStartPar
After completing the LAMMPS build and also configuring and
compiling Quantum ESPRESSO with external library support (via
“make couple”), go back to the \sphinxcode{\sphinxupquote{lib/qmmm}} folder and follow the
instructions on the README file to build the combined LAMMPS/QE
QM/MM executable (pwqmmm.x) in the \sphinxcode{\sphinxupquote{lib/qmmm}} folder.
\end{sphinxadmonition}

\begin{sphinxadmonition}{note}{Traditional make}

\sphinxAtStartPar
Before building LAMMPS, you must build the QMMM library in
\sphinxcode{\sphinxupquote{lib/qmmm}}.  You can do this manually if you prefer; follow the
first two steps explained in \sphinxcode{\sphinxupquote{lib/qmmm/README}}.  You can also do
it in one step from the \sphinxcode{\sphinxupquote{lammps/src}} dir, using a command like
these, which simply invokes the \sphinxcode{\sphinxupquote{lib/qmmm/Install.py}} script with
the specified args:

\begin{sphinxVerbatim}[commandchars=\\\{\}]
\PYG{c+c1}{\PYGZsh{} print help message}
make\PYG{+w}{ }lib\PYGZhy{}qmmm

\PYG{c+c1}{\PYGZsh{} build with GNU Fortran compiler (settings as in \PYGZdq{}make serial\PYGZdq{})}
make\PYG{+w}{ }lib\PYGZhy{}qmmm\PYG{+w}{ }\PYG{n+nv}{args}\PYG{o}{=}\PYG{l+s+s2}{\PYGZdq{}\PYGZhy{}m serial\PYGZdq{}}

\PYG{c+c1}{\PYGZsh{} build with default MPI compiler (settings as in \PYGZdq{}make mpi\PYGZdq{})}
make\PYG{+w}{ }lib\PYGZhy{}qmmm\PYG{+w}{ }\PYG{n+nv}{args}\PYG{o}{=}\PYG{l+s+s2}{\PYGZdq{}\PYGZhy{}m mpi\PYGZdq{}}

\PYG{c+c1}{\PYGZsh{} build with GNU Fortran compiler}
make\PYG{+w}{ }lib\PYGZhy{}qmmm\PYG{+w}{ }\PYG{n+nv}{args}\PYG{o}{=}\PYG{l+s+s2}{\PYGZdq{}\PYGZhy{}m gfortran\PYGZdq{}}
\end{sphinxVerbatim}

\sphinxAtStartPar
The build should produce two files: \sphinxcode{\sphinxupquote{lib/qmmm/libqmmm.a}} and
\sphinxcode{\sphinxupquote{lib/qmmm/Makefile.lammps}}.  The latter is copied from an
existing \sphinxcode{\sphinxupquote{Makefile.lammps.*}} and has settings needed to build
LAMMPS with the QMMM library (though typically the settings are
just blank).  If necessary, you can edit/create a new
\sphinxcode{\sphinxupquote{lib/qmmm/Makefile.\textless{}machine\textgreater{}}} file for your system, which should
define an \sphinxcode{\sphinxupquote{EXTRAMAKE}} variable to specify a corresponding
\sphinxcode{\sphinxupquote{Makefile.lammps.\textless{}machine\textgreater{}}} file.

\sphinxAtStartPar
You can then install QMMM package and build LAMMPS in the usual
manner.  After completing the LAMMPS build and compiling Quantum
ESPRESSO with external library support (via \sphinxcode{\sphinxupquote{make couple}}), go
back to the \sphinxcode{\sphinxupquote{lib/qmmm}} folder and follow the instructions in the
README file to build the combined LAMMPS/QE QM/MM executable
(\sphinxcode{\sphinxupquote{pwqmmm.x}}) in the \sphinxcode{\sphinxupquote{lib/qmmm}} folder.
\end{sphinxadmonition}
\end{quote}


\bigskip\hrule\bigskip



\subsection{RHEO package}
\label{\detokenize{Build_extras:rheo-package}}\label{\detokenize{Build_extras:rheo}}
\sphinxAtStartPar
To build with this package you must have the \sphinxtitleref{GNU Scientific Library
(GSL) \textless{}https://www.gnu.org/software/gsl/\textgreater{}} installed in locations that
are accessible in your environment.  The GSL library should be at least
version 2.7.
\begin{quote}

\begin{sphinxadmonition}{note}{CMake build}

\sphinxAtStartPar
If CMake cannot find the GSL library or include files, you can set:

\begin{sphinxVerbatim}[commandchars=\\\{\}]
\PYGZhy{}D\PYG{+w}{ }\PYG{n+nv}{GSL\PYGZus{}ROOT\PYGZus{}DIR}\PYG{o}{=}path\PYG{+w}{    }\PYG{c+c1}{\PYGZsh{} path to root of GSL installation}
\end{sphinxVerbatim}
\end{sphinxadmonition}

\begin{sphinxadmonition}{note}{Traditional make}

\sphinxAtStartPar
LAMMPS will try to auto\sphinxhyphen{}detect the GSL compiler and linker flags
from the corresponding \sphinxcode{\sphinxupquote{pkg\sphinxhyphen{}config}} file (\sphinxcode{\sphinxupquote{gsl.pc}}), otherwise
you can edit the file \sphinxcode{\sphinxupquote{lib/rheo/Makefile.lammps}}
to specify the paths and library names where indicated by comments.
This must be done \sphinxstylestrong{before} the package is installed.
\end{sphinxadmonition}
\end{quote}


\bigskip\hrule\bigskip



\subsection{SCAFACOS package}
\label{\detokenize{Build_extras:scafacos-package}}\label{\detokenize{Build_extras:scafacos}}
\sphinxAtStartPar
To build with this package, you must download and build the
\sphinxhref{http://www.scafacos.de}{ScaFaCoS Coulomb solver library}
\begin{quote}

\begin{sphinxadmonition}{note}{CMake build}

\begin{sphinxVerbatim}[commandchars=\\\{\}]
\PYGZhy{}D\PYG{+w}{ }\PYG{n+nv}{DOWNLOAD\PYGZus{}SCAFACOS}\PYG{o}{=}value\PYG{+w}{    }\PYG{c+c1}{\PYGZsh{} download ScaFaCoS for build, value = no (default) or yes}
\PYGZhy{}D\PYG{+w}{ }\PYG{n+nv}{SCAFACOS\PYGZus{}LIBRARY}\PYG{o}{=}path\PYG{+w}{      }\PYG{c+c1}{\PYGZsh{} ScaFaCos library file (only needed if at custom location)}
\PYGZhy{}D\PYG{+w}{ }\PYG{n+nv}{SCAFACOS\PYGZus{}INCLUDE\PYGZus{}DIR}\PYG{o}{=}path\PYG{+w}{  }\PYG{c+c1}{\PYGZsh{} ScaFaCoS include directory (only needed if at custom location)}
\end{sphinxVerbatim}

\sphinxAtStartPar
If \sphinxcode{\sphinxupquote{DOWNLOAD\_SCAFACOS}} is set, the ScaFaCoS library will be
downloaded and built inside the CMake build directory.  If the
ScaFaCoS library is already on your system (in a location CMake
cannot find it), \sphinxcode{\sphinxupquote{SCAFACOS\_LIBRARY}} is the filename (plus path) of
the ScaFaCoS library file, not the directory the library file is
in.  \sphinxcode{\sphinxupquote{SCAFACOS\_INCLUDE\_DIR}} is the directory the ScaFaCoS include
file is in.
\end{sphinxadmonition}

\begin{sphinxadmonition}{note}{Traditional make}

\sphinxAtStartPar
You can download and build the ScaFaCoS library manually if you
prefer; follow the instructions in \sphinxcode{\sphinxupquote{lib/scafacos/README}}.  You
can also do it in one step from the \sphinxcode{\sphinxupquote{lammps/src}} dir, using a
command like these, which simply invokes the
\sphinxcode{\sphinxupquote{lib/scafacos/Install.py}} script with the specified args:

\begin{sphinxVerbatim}[commandchars=\\\{\}]
\PYG{c+c1}{\PYGZsh{} print help message}
make\PYG{+w}{ }lib\PYGZhy{}scafacos

\PYG{c+c1}{\PYGZsh{} download and build in lib/scafacos/scafacos\PYGZhy{}\PYGZlt{}version\PYGZgt{}}
make\PYG{+w}{ }lib\PYGZhy{}scafacos\PYG{+w}{ }\PYG{n+nv}{args}\PYG{o}{=}\PYG{l+s+s2}{\PYGZdq{}\PYGZhy{}b\PYGZdq{}}

\PYG{c+c1}{\PYGZsh{} use existing ScaFaCoS installation in \PYGZdl{}HOME/scafacos}
make\PYG{+w}{ }lib\PYGZhy{}scafacos\PYG{+w}{ }\PYG{n+nv}{args}\PYG{o}{=}\PYG{l+s+s2}{\PYGZdq{}}\PYG{l+s+s2}{\PYGZhy{}p }\PYG{n+nv}{\PYGZdl{}HOME}\PYG{l+s+s2}{/scafacos}
\end{sphinxVerbatim}

\sphinxAtStartPar
Note that two symbolic (soft) links, \sphinxcode{\sphinxupquote{includelink}} and \sphinxcode{\sphinxupquote{liblink}}, are
created in \sphinxcode{\sphinxupquote{lib/scafacos}} to point to the ScaFaCoS src dir.  When LAMMPS
builds in src it will use these links.  You should not need to edit
the \sphinxcode{\sphinxupquote{lib/scafacos/Makefile.lammps}} file.
\end{sphinxadmonition}
\end{quote}


\bigskip\hrule\bigskip



\subsection{VTK package}
\label{\detokenize{Build_extras:vtk-package}}\label{\detokenize{Build_extras:vtk}}
\sphinxAtStartPar
To build with this package you must have the VTK library installed on
your system.
\begin{quote}

\begin{sphinxadmonition}{note}{CMake build}

\sphinxAtStartPar
No additional settings are needed besides \sphinxcode{\sphinxupquote{\sphinxhyphen{}D PKG\_VTK=yes}}.

\sphinxAtStartPar
This should auto\sphinxhyphen{}detect the VTK library if it is installed on your
system at standard locations.  Several advanced VTK options exist
if you need to specify where it was installed.  Use the \sphinxcode{\sphinxupquote{ccmake}}
(terminal window) or \sphinxcode{\sphinxupquote{cmake\sphinxhyphen{}gui}} (graphical) tools to see these
options and set them interactively from their user interfaces.
\end{sphinxadmonition}

\begin{sphinxadmonition}{note}{Traditional make}

\sphinxAtStartPar
The \sphinxcode{\sphinxupquote{lib/vtk/Makefile.lammps}} file has settings for accessing
VTK files and its library, which LAMMPS needs to build with this
package.  If the settings are not valid for your system, check if
one of the other \sphinxcode{\sphinxupquote{lib/vtk/Makefile.lammps.*}} files is compatible
and copy it to Makefile.lammps.  If none of the provided files
work, you will need to edit the \sphinxcode{\sphinxupquote{Makefile.lammps}} file.  See
\sphinxcode{\sphinxupquote{lib/vtk/README}} for details.
\end{sphinxadmonition}
\end{quote}

\sphinxstepscope


\section{Build the LAMMPS documentation}
\label{\detokenize{Build_manual:build-the-lammps-documentation}}\label{\detokenize{Build_manual::doc}}
\sphinxAtStartPar
Depending on how you obtained LAMMPS and whether you have built the
manual yourself, this directory has a number of subdirectories and
files. Here is a list with descriptions:

\begin{sphinxVerbatim}[commandchars=\\\{\}]
README\PYG{+w}{           }\PYG{c+c1}{\PYGZsh{} brief info about the documentation}
src\PYG{+w}{              }\PYG{c+c1}{\PYGZsh{} content files for LAMMPS documentation}
html\PYG{+w}{             }\PYG{c+c1}{\PYGZsh{} HTML version of the LAMMPS manual (see html/Manual.html)}
utils\PYG{+w}{            }\PYG{c+c1}{\PYGZsh{} tools and settings for building the documentation}
lammps.1\PYG{+w}{         }\PYG{c+c1}{\PYGZsh{} man page for the lammps command}
msi2lmp.1\PYG{+w}{        }\PYG{c+c1}{\PYGZsh{} man page for the msi2lmp command}
Manual.pdf\PYG{+w}{       }\PYG{c+c1}{\PYGZsh{} large PDF version of entire manual}
LAMMPS.epub\PYG{+w}{      }\PYG{c+c1}{\PYGZsh{} Manual in ePUB e\PYGZhy{}book format}
LAMMPS.mobi\PYG{+w}{      }\PYG{c+c1}{\PYGZsh{} Manual in MOBI e\PYGZhy{}book format}
docenv\PYG{+w}{           }\PYG{c+c1}{\PYGZsh{} virtualenv folder for processing the manual sources}
doctrees\PYG{+w}{         }\PYG{c+c1}{\PYGZsh{} temporary data from processing the manual}
doxygen\PYG{+w}{          }\PYG{c+c1}{\PYGZsh{} doxygen configuration and output}
.gitignore\PYG{+w}{       }\PYG{c+c1}{\PYGZsh{} list of files and folders to be ignored by git}
doxygen\PYGZhy{}warn.log\PYG{+w}{ }\PYG{c+c1}{\PYGZsh{} logfile with warnings from running doxygen}
github\PYGZhy{}development\PYGZhy{}workflow.md\PYG{+w}{   }\PYG{c+c1}{\PYGZsh{} notes on the LAMMPS development workflow}
\end{sphinxVerbatim}

\sphinxAtStartPar
If you downloaded LAMMPS as a tarball from \sphinxhref{https://www.lammps.org}{the LAMMPS website},
the html folder and the PDF files should be included.

\sphinxAtStartPar
If you downloaded LAMMPS from the public git repository, then the HTML
and PDF files are not included.  You can build the HTML or PDF files yourself,
by typing \sphinxcode{\sphinxupquote{make html}}  or \sphinxcode{\sphinxupquote{make pdf}} in the \sphinxcode{\sphinxupquote{doc}} folder.  This requires
various tools and files.  Some of them have to be installed (see below).  For
the rest the build process will attempt to download and install them into
a python virtual environment and local folders.

\sphinxAtStartPar
A current version of the manual (latest feature release, that is the state
of the \sphinxstyleemphasis{release} branch) is is available online at:
\sphinxurl{https://docs.lammps.org/}.
A version of the manual corresponding to the ongoing development (that is
the state of the \sphinxstyleemphasis{develop} branch) is available online at:
\sphinxurl{https://docs.lammps.org/latest/}
A version of the manual corresponding to the latest stable LAMMPS release
(that is the state of the \sphinxstyleemphasis{stable} branch) is available online at:
\sphinxurl{https://docs.lammps.org/stable/}


\subsection{Build using GNU make}
\label{\detokenize{Build_manual:build-using-gnu-make}}
\sphinxAtStartPar
The LAMMPS manual is written in \sphinxhref{https://www.sphinx-doc.org/en/master/usage/restructuredtext/index.html}{reStructuredText} format which
can be translated to different output format using the \sphinxhref{https://www.sphinx-doc.org/}{Sphinx} document generator tool.  It also
incorporates programmer documentation extracted from the LAMMPS C++
sources through the \sphinxhref{https://doxygen.nl/}{Doxygen} program.  Currently
the translation to HTML, PDF (via LaTeX), ePUB (for many e\sphinxhyphen{}book readers)
and MOBI (for Amazon Kindle readers) are supported.  For that to work a
Python interpreter version 3.8 or later, the \sphinxcode{\sphinxupquote{doxygen}} tools and
internet access to download additional files and tools are required.
This download is usually only required once or after the documentation
folder is returned to a pristine state with \sphinxcode{\sphinxupquote{make clean\sphinxhyphen{}all}}.

\sphinxAtStartPar
For the documentation build a python virtual environment is set up in
the folder \sphinxcode{\sphinxupquote{doc/docenv}} and various python packages are installed into
that virtual environment via the \sphinxcode{\sphinxupquote{pip}} tool.  For rendering embedded
LaTeX code also the \sphinxhref{https://www.mathjax.org/}{MathJax} JavaScript
engine needs to be downloaded.  If you need to pass additional options
to the pip commands to work (e.g. to use a web proxy or to point to
additional SSL certificates) you can set them via the \sphinxcode{\sphinxupquote{PIP\_OPTIONS}}
environment variable or uncomment and edit the \sphinxcode{\sphinxupquote{PIP\_OPTIONS}} setting
at beginning of the makefile.

\sphinxAtStartPar
The actual translation is then done via \sphinxcode{\sphinxupquote{make}} commands in the doc
folder.  The following \sphinxcode{\sphinxupquote{make}} commands are available:

\begin{sphinxVerbatim}[commandchars=\\\{\}]
make\PYG{+w}{ }html\PYG{+w}{          }\PYG{c+c1}{\PYGZsh{} generate HTML in html dir using Sphinx}
make\PYG{+w}{ }pdf\PYG{+w}{           }\PYG{c+c1}{\PYGZsh{} generate PDF  as Manual.pdf using Sphinx and PDFLaTeX}
make\PYG{+w}{ }epub\PYG{+w}{          }\PYG{c+c1}{\PYGZsh{} generate LAMMPS.epub in ePUB format using Sphinx}
make\PYG{+w}{ }mobi\PYG{+w}{          }\PYG{c+c1}{\PYGZsh{} generate LAMMPS.mobi in MOBI format using ebook\PYGZhy{}convert}

make\PYG{+w}{ }fasthtml\PYG{+w}{      }\PYG{c+c1}{\PYGZsh{} generate approximate HTML in fasthtml dir using Sphinx}
\PYG{+w}{                   }\PYG{c+c1}{\PYGZsh{} some Sphinx extensions do not work correctly with this}

make\PYG{+w}{ }clean\PYG{+w}{         }\PYG{c+c1}{\PYGZsh{} remove intermediate RST files created by HTML build}
make\PYG{+w}{ }clean\PYGZhy{}all\PYG{+w}{     }\PYG{c+c1}{\PYGZsh{} remove entire build folder and any cached data}

make\PYG{+w}{ }anchor\PYGZus{}check\PYG{+w}{  }\PYG{c+c1}{\PYGZsh{} check for duplicate anchor labels}
make\PYG{+w}{ }style\PYGZus{}check\PYG{+w}{   }\PYG{c+c1}{\PYGZsh{} check for complete and consistent style lists}
make\PYG{+w}{ }package\PYGZus{}check\PYG{+w}{ }\PYG{c+c1}{\PYGZsh{} check for complete and consistent package lists}
make\PYG{+w}{ }link\PYGZus{}check\PYG{+w}{    }\PYG{c+c1}{\PYGZsh{} check for broken or outdated URLs}
make\PYG{+w}{ }spelling\PYG{+w}{      }\PYG{c+c1}{\PYGZsh{} spell\PYGZhy{}check the manual}
\end{sphinxVerbatim}


\bigskip\hrule\bigskip



\subsection{Build using CMake}
\label{\detokenize{Build_manual:build-using-cmake}}
\sphinxAtStartPar
It is also possible to create the HTML version (and \sphinxstylestrong{only} the HTML
version) of the manual within the {\hyperref[\detokenize{Build_cmake::doc}]{\sphinxcrossref{\DUrole{doc}{CMake build directory}}}}.  The reason for this option is to include the
installation of the HTML manual pages into the “install” step when
installing LAMMPS after the CMake build via \sphinxcode{\sphinxupquote{cmake \sphinxhyphen{}\sphinxhyphen{}build . \sphinxhyphen{}\sphinxhyphen{}target
install}}.  The documentation build is included in the default build
target, but can also be requested independently with
\sphinxcode{\sphinxupquote{cmake \sphinxhyphen{}\sphinxhyphen{}build . \sphinxhyphen{}\sphinxhyphen{}target doc}}.  If you need to pass additional options
to the pip commands to work (e.g. to use a web proxy or to point to
additional SSL certificates) you can set them via the \sphinxcode{\sphinxupquote{PIP\_OPTIONS}}
environment variable.

\begin{sphinxVerbatim}[commandchars=\\\{\}]
\PYGZhy{}D\PYG{+w}{ }\PYG{n+nv}{BUILD\PYGZus{}DOC}\PYG{o}{=}value\PYG{+w}{       }\PYG{c+c1}{\PYGZsh{} yes or no (default)}
\end{sphinxVerbatim}


\bigskip\hrule\bigskip



\subsection{Prerequisites for HTML}
\label{\detokenize{Build_manual:prerequisites-for-html}}
\sphinxAtStartPar
To run the HTML documentation build toolchain, python 3, git, doxygen,
and virtualenv have to be installed locally.  Here are instructions for
common setups:
\begin{quote}

\begin{sphinxadmonition}{note}{Ubuntu}

\begin{sphinxVerbatim}[commandchars=\\\{\}]
sudo\PYG{+w}{ }apt\PYGZhy{}get\PYG{+w}{ }install\PYG{+w}{ }git\PYG{+w}{ }doxygen
\end{sphinxVerbatim}
\end{sphinxadmonition}

\begin{sphinxadmonition}{note}{RHEL or CentOS (Version 7.x)}

\begin{sphinxVerbatim}[commandchars=\\\{\}]
sudo\PYG{+w}{ }yum\PYG{+w}{ }install\PYG{+w}{ }git\PYG{+w}{ }doxygen
\end{sphinxVerbatim}
\end{sphinxadmonition}

\begin{sphinxadmonition}{note}{Fedora or RHEL/CentOS (8.x or later)}

\begin{sphinxVerbatim}[commandchars=\\\{\}]
sudo\PYG{+w}{ }dnf\PYG{+w}{ }install\PYG{+w}{ }git\PYG{+w}{ }doxygen
\end{sphinxVerbatim}
\end{sphinxadmonition}

\begin{sphinxadmonition}{note}{macOS}

\sphinxAtStartPar
\sphinxstyleemphasis{Python 3}

\sphinxAtStartPar
If Python 3 is not available on your macOS system, you can
download the latest Python 3 macOS package from
\sphinxurl{https://www.python.org} and install it.
This will install both Python 3 and pip3.
\end{sphinxadmonition}
\end{quote}


\subsection{Prerequisites for PDF}
\label{\detokenize{Build_manual:prerequisites-for-pdf}}
\sphinxAtStartPar
In addition to the tools needed for building the HTML format manual,
a working LaTeX installation with support for PDFLaTeX and a selection
of LaTeX styles/packages are required.  To run the PDFLaTeX translation
the \sphinxcode{\sphinxupquote{latexmk}} script needs to be installed as well.


\subsection{Prerequisites for ePUB and MOBI}
\label{\detokenize{Build_manual:prerequisites-for-epub-and-mobi}}
\sphinxAtStartPar
In addition to the tools needed for building the HTML format manual,
a working LaTeX installation with a few add\sphinxhyphen{}on LaTeX packages
as well as the \sphinxcode{\sphinxupquote{dvipng}} tool are required to convert embedded
math expressions transparently into embedded images.

\sphinxAtStartPar
For converting the generated ePUB file to a MOBI format file (for e\sphinxhyphen{}book
readers, like Kindle, that cannot read ePUB), you also need to have the
\sphinxcode{\sphinxupquote{ebook\sphinxhyphen{}convert}} tool from the “calibre” software
installed. \sphinxurl{https://calibre-ebook.com/}
Typing \sphinxcode{\sphinxupquote{make mobi}} will first create the ePUB file and then convert
it.  On the Kindle readers in particular, you also have support for PDF
files, so you could download and view the PDF version as an alternative.


\subsection{Instructions for Developers}
\label{\detokenize{Build_manual:instructions-for-developers}}
\sphinxAtStartPar
When adding new styles or options to the LAMMPS code, corresponding
documentation is required and either existing files in the \sphinxcode{\sphinxupquote{src}}
folder need to be updated or new files added. These files are written in
\sphinxhref{https://www.sphinx-doc.org/en/master/usage/restructuredtext/index.html}{reStructuredText} markup for translation with the Sphinx tool.

\sphinxAtStartPar
Before contributing any documentation, please check that both the HTML
and the PDF format documentation can translate without errors.  During
testing the html translation, you may use the \sphinxcode{\sphinxupquote{make fasthtml}} command
which does an approximate translation (i.e. not all Sphinx features and
extensions will work), but runs very fast because it will only translate
files that have been changed since the last \sphinxcode{\sphinxupquote{make fasthtml}} command.

\sphinxAtStartPar
Please also check the output to the console for any warnings or problems.  There will
be multiple tests run automatically:
\begin{itemize}
\item {} 
\sphinxAtStartPar
A test for correctness of all anchor labels and their references

\item {} 
\sphinxAtStartPar
A test that all LAMMPS packages (= folders with sources in
\sphinxcode{\sphinxupquote{lammps/src}}) are documented and listed.  A typical warning shows
the name of the folder with the suspected new package code and the
documentation files where they need to be listed:

\begin{sphinxVerbatim}[commandchars=\\\{\}]
Found 88 packages
Package NEWPACKAGE missing in Packages\PYGZus{}list.rst
Package NEWPACKAGE missing in Packages\PYGZus{}details.rst
\end{sphinxVerbatim}

\item {} 
\sphinxAtStartPar
A test that only standard, printable ASCII text characters are used.
This runs the command \sphinxcode{\sphinxupquote{env LC\_ALL=C grep \sphinxhyphen{}n \textquotesingle{}{[}\textasciicircum{} \sphinxhyphen{}\textasciitilde{}{]}\textquotesingle{} src/*.rst}} and
thus prints all offending lines with filename and line number
prepended to the screen.  Special characters like Greek letters
(\(\alpha~~\sigma~~\epsilon\)), super\sphinxhyphen{} or subscripts
(\(x^2~~\mathrm{U}_{LJ}\)), mathematical expressions
(\(\frac{1}{2}\mathrm{N}~~x\to\infty\)), or the Angstrom symbol
(\(\AA\)) should be typeset with embedded LaTeX (like this
\sphinxcode{\sphinxupquote{:math:\textasciigrave{}\textbackslash{}alpha \textbackslash{}sigma \textbackslash{}epsilon\textasciigrave{}}}, \sphinxcode{\sphinxupquote{:math:\textasciigrave{}x\textasciicircum{}2 \textbackslash{}mathrm\{E\}\_\{LJ\}\textasciigrave{}}},
\sphinxcode{\sphinxupquote{:math:\textasciigrave{}\textbackslash{}frac\{1\}\{2\}\textbackslash{}mathrm\{N\} x\textbackslash{}to\textbackslash{}infty\textasciigrave{}}}, or \sphinxcode{\sphinxupquote{:math:\textasciigrave{}\textbackslash{}AA\textasciigrave{}}}).

\item {} 
\sphinxAtStartPar
Embedded LaTeX is rendered in HTML output with \sphinxhref{https://www.mathjax.org/}{MathJax} and in PDF output by passing the embedded
text to LaTeX.  Some care has to be taken, though, since there are
limitations which macros and features can be used in either mode, so
it is recommended to always check whether any new or changed
documentation does translate and render correctly with either output.

\item {} 
\sphinxAtStartPar
A test whether all styles are documented and listed in their
respective overview pages.  A typical output with warnings looks like this:

\begin{sphinxVerbatim}[commandchars=\\\{\}]
Parsed style names w/o suffixes from C++ tree in ../src:
   Angle styles:      21    Atom styles:       24
   Body styles:        3    Bond styles:       17
   Command styles:    41    Compute styles:   143
   Dihedral styles:   16    Dump styles:       26
   Fix styles:       223    Improper styles:   13
   Integrate styles:   4    Kspace styles:     15
   Minimize styles:    9    Pair styles:      234
   Reader styles:      4    Region styles:      8
Compute style entry newcomp is missing or incomplete in Commands\PYGZus{}compute.rst
Compute style entry newcomp is missing or incomplete in compute.rst
Fix style entry newfix is missing or incomplete in Commands\PYGZus{}fix.rst
Fix style entry newfix is missing or incomplete in fix.rst
Pair style entry new is missing or incomplete in Commands\PYGZus{}pair.rst
Pair style entry new is missing or incomplete in pair\PYGZus{}style.rst
Found 6 issue(s) with style lists
\end{sphinxVerbatim}

\end{itemize}

\sphinxAtStartPar
In addition, there is the option to run a spellcheck on the entire
manual with \sphinxcode{\sphinxupquote{make spelling}}.  This requires \sphinxhref{https://github.com/AbiWord/enchant}{a library called enchant}.  To avoid printing out \sphinxstyleemphasis{false
positives} (e.g. keywords, names, abbreviations) those can be added to
the file \sphinxcode{\sphinxupquote{lammps/doc/utils/sphinx\sphinxhyphen{}config/false\_positives.txt}}.

\sphinxstepscope


\section{Notes for building LAMMPS on Windows}
\label{\detokenize{Build_windows:notes-for-building-lammps-on-windows}}\label{\detokenize{Build_windows::doc}}\begin{itemize}
\item {} 
\sphinxAtStartPar
{\hyperref[\detokenize{Build_windows:generic}]{\sphinxcrossref{\DUrole{std,std-ref}{General remarks}}}}

\item {} 
\sphinxAtStartPar
{\hyperref[\detokenize{Build_windows:linux}]{\sphinxcrossref{\DUrole{std,std-ref}{Running Linux on Windows}}}}

\item {} 
\sphinxAtStartPar
{\hyperref[\detokenize{Build_windows:gnu}]{\sphinxcrossref{\DUrole{std,std-ref}{Using GNU GCC ported to Windows}}}}

\item {} 
\sphinxAtStartPar
{\hyperref[\detokenize{Build_windows:msvc}]{\sphinxcrossref{\DUrole{std,std-ref}{Using Visual Studio}}}}

\item {} 
\sphinxAtStartPar
{\hyperref[\detokenize{Build_windows:oneapi}]{\sphinxcrossref{\DUrole{std,std-ref}{Using Intel oneAPI compilers and libraries}}}}

\item {} 
\sphinxAtStartPar
{\hyperref[\detokenize{Build_windows:cross}]{\sphinxcrossref{\DUrole{std,std-ref}{Using a cross\sphinxhyphen{}compiler}}}}

\end{itemize}


\bigskip\hrule\bigskip



\subsection{General remarks}
\label{\detokenize{Build_windows:general-remarks}}\label{\detokenize{Build_windows:generic}}
\sphinxAtStartPar
LAMMPS is developed and tested primarily on Linux machines.  The vast
majority of HPC clusters and supercomputers today run on Linux as well.
While portability to other platforms is desired, it is not always
achieved.  That is sometimes due to non\sphinxhyphen{}portable code in LAMMPS itself,
but more often due to portability limitations of external libraries and
tools required to build a specific feature or package.  The LAMMPS
developers are dependent on LAMMPS users giving feedback and providing
assistance in resolving portability issues.  This is particularly true
for compiling LAMMPS on Windows, since this platform has significant
differences in some low\sphinxhyphen{}level functionality.  As of LAMMPS version 14
December 2021, large parts of LAMMPS can be compiled natively with the
Microsoft Visual C++ Compilers.  As of LAMMPS version 31 May 2022, also
the Intel oneAPI compilers can compile large parts of LAMMPS natively on
Windows.  This is mostly facilitated by using the
{\hyperref[\detokenize{Developer_platform::doc}]{\sphinxcrossref{\DUrole{doc}{Platform abstraction functions}}}} in the \sphinxcode{\sphinxupquote{platform}} namespace and CMake.

\sphinxAtStartPar
Before trying to build LAMMPS on Windows yourself, please consider the
\sphinxhref{https://packages.lammps.org/windows.html}{pre\sphinxhyphen{}compiled Windows installer packages}
and see if they are sufficient for your needs.


\subsection{Running Linux on Windows}
\label{\detokenize{Build_windows:running-linux-on-windows}}\label{\detokenize{Build_windows:linux}}
\sphinxAtStartPar
If it is necessary for you to compile LAMMPS on a Windows machine
(e.g. because it is your main desktop), please also consider using a
virtual machine software and compile and run LAMMPS in a Linux virtual
machine, or \sphinxhyphen{} if you have a sufficiently up\sphinxhyphen{}to\sphinxhyphen{}date Windows 10 or
Windows 11 installation \sphinxhyphen{} consider using the Windows subsystem for
Linux.  This optional Windows feature allows you to run the bash shell
of a Linux system (Ubuntu by default) from within Windows and from there
on, you can pretty much use that shell like you are running on a regular
Ubuntu Linux machine (e.g. installing software via apt\sphinxhyphen{}get and more).
For more details on that, please see {\hyperref[\detokenize{Howto_wsl::doc}]{\sphinxcrossref{\DUrole{doc}{this tutorial}}}}.


\subsection{Using a GNU GCC ported to Windows}
\label{\detokenize{Build_windows:using-a-gnu-gcc-ported-to-windows}}\label{\detokenize{Build_windows:gnu}}
\sphinxAtStartPar
One option for compiling LAMMPS on Windows natively is to install a Bash
shell, Unix shell utilities, Perl, Python, GNU make, and a GNU compiler
ported to Windows.  The Cygwin package provides a unix/linux interface
to low\sphinxhyphen{}level Windows functions, so LAMMPS can be compiled on Windows.
The necessary (minor) modifications to LAMMPS are included, but may not
always up\sphinxhyphen{}to\sphinxhyphen{}date for recently added functionality and the corresponding
new code.  A machine makefile for using cygwin for the old build system
is provided.  Using CMake for this mode of compilation is untested and
not likely to work.

\sphinxAtStartPar
When compiling for Windows do \sphinxstylestrong{not} set the \sphinxcode{\sphinxupquote{\sphinxhyphen{}DLAMMPS\_MEMALIGN}}
define in the LMP\_INC makefile variable and add \sphinxcode{\sphinxupquote{\sphinxhyphen{}lwsock32 \sphinxhyphen{}lpsapi}} to
the linker flags in LIB makefile variable. Try adding \sphinxcode{\sphinxupquote{\sphinxhyphen{}static\sphinxhyphen{}libgcc}}
or \sphinxcode{\sphinxupquote{\sphinxhyphen{}static}} or both to the linker flags when your resulting LAMMPS
Windows executable complains about missing .dll files. The CMake
configuration should set this up automatically, but is untested.

\sphinxAtStartPar
In case of problems, you are recommended to contact somebody with
experience in using Cygwin.  If you do come across portability problems
requiring changes to the LAMMPS source code, or figure out corrections
yourself, please report them on the
\sphinxhref{https://matsci.org/c/lammps/lammps-development/}{LAMMPS forum at MatSci},
or file them as an issue or pull request on the LAMMPS GitHub project.


\subsection{Using Microsoft Visual Studio}
\label{\detokenize{Build_windows:using-microsoft-visual-studio}}\label{\detokenize{Build_windows:msvc}}
\sphinxAtStartPar
Following the integration of the {\hyperref[\detokenize{Developer_platform::doc}]{\sphinxcrossref{\DUrole{doc}{platform namespace}}}} into the LAMMPS code base, portability of LAMMPS
for native compilation on Windows using Visual Studio has been
significantly improved.  This has been tested with Visual Studio 2019
(aka version 16) and Visual Studio 2022 (aka version 17).  We strongly
recommend using Visual Studio 2022 version 17.1 or later.  Not all
features and packages in LAMMPS are currently supported out of the box,
but a preset \sphinxcode{\sphinxupquote{cmake/presets/windows.cmake}} is provided that contains
the packages that have been compiled successfully so far.  You \sphinxstylestrong{must}
use the CMake based build procedure, since there is no support for GNU
make or the Unix shell utilities required for the GNU make build
procedure.

\sphinxAtStartPar
It is possible to use both the integrated CMake support of the Visual
Studio IDE or use an external CMake installation (e.g. downloaded from
cmake.org) to create build files and compile LAMMPS from the command line.

\sphinxAtStartPar
Compilation via command line and unit tests are checked automatically
for the LAMMPS development branch through
\sphinxhref{https://github.com/lammps/lammps/actions/workflows/compile-msvc.yml}{GitHub Actions}.

\begin{sphinxadmonition}{note}{Note:}
\sphinxAtStartPar
Versions of Visual Studio before version 17.1 may scan the entire
LAMMPS source tree and likely miss the correct master
\sphinxcode{\sphinxupquote{CMakeLists.txt}} and get confused since there are multiple files
of that name in different folders but none in top level folder.
\end{sphinxadmonition}

\sphinxAtStartPar
Please note, that for either approach CMake will create a so\sphinxhyphen{}called
{\hyperref[\detokenize{Build_cmake:cmake-multiconfig}]{\sphinxcrossref{\DUrole{std,std-ref}{“multi\sphinxhyphen{}configuration” build environment}}}}, and
the command lines for building and testing LAMMPS must be adjusted
accordingly.

\sphinxAtStartPar
The LAMMPS cmake folder contains a \sphinxcode{\sphinxupquote{CMakeSettings.json}} file with
build configurations for MSVC compilers and the MS provided Clang
compiler package in Debug and Release mode.

\sphinxAtStartPar
To support running in parallel you can compile with OpenMP enabled using
the OPENMP package or install Microsoft MPI (including the SDK) and compile
LAMMPS with MPI enabled.

\begin{sphinxadmonition}{note}{Note:}
\sphinxAtStartPar
This is work in progress and you should contact the LAMMPS developers
via GitHub or the \sphinxhref{https://matsci.org/c/lammps/lammps-development/}{LAMMPS forum at MatSci},
if you have questions or LAMMPS specific problems.
\end{sphinxadmonition}


\subsection{Using Intel oneAPI Compilers and Libraries}
\label{\detokenize{Build_windows:using-intel-oneapi-compilers-and-libraries}}\label{\detokenize{Build_windows:oneapi}}
\sphinxAtStartPar
\DUrole{versionmodified,added}{Added in version 31May2022.}

\sphinxAtStartPar
After installing the \sphinxhref{https://www.intel.com/content/www/us/en/developer/tools/oneapi/toolkits.html}{Intel oneAPI}
base toolkit and the HPC toolkit, it is also possible to compile large
parts of LAMMPS natively on Windows using Intel compilers.  The HPC
toolkit provides two sets of C/C++ and Fortran compilers: the so\sphinxhyphen{}called
“classic” compilers (\sphinxcode{\sphinxupquote{icl.exe}} and \sphinxcode{\sphinxupquote{ifort.exe}}) and newer, LLVM
based compilers (\sphinxcode{\sphinxupquote{icx.exe}} and \sphinxcode{\sphinxupquote{ifx.exe}}).  In addition to the
compilers and their dependent modules, also the thread building blocks
(TBB) and the math kernel library (MKL) need to be installed.  Two
presets (\sphinxcode{\sphinxupquote{cmake/presets/windows\sphinxhyphen{}intel\sphinxhyphen{}llvm.cmake}} and
\sphinxcode{\sphinxupquote{cmake/presets/windows\sphinxhyphen{}intel\sphinxhyphen{}classic.cmake}}) are provided for
selecting the LLVM based or classic compilers, respectively. The preset
\sphinxcode{\sphinxupquote{cmake/presets/windows.cmake}} enables compatible packages that are not
dependent on additional features or libraries.  You \sphinxstylestrong{must} use the
CMake based build procedure and use Ninja as build tool.  For compiling
from the command prompt, thus both \sphinxhref{https://cmake.org}{CMake} and
\sphinxhref{https://ninja-build.org}{Ninja\sphinxhyphen{}build} binaries must be installed.  It
is also possible to use Visual Studio, if it is started (\sphinxcode{\sphinxupquote{devenv.exe}})
from a command prompt that has the Intel oneAPI compilers enabled.  The
Visual Studio settings file in the \sphinxcode{\sphinxupquote{cmake}} folder contains
configurations for both compiler variants in debug and release settings.
Those will use the CMake and Ninja binaries bundled with Visual Studio,
thus a separate installation is not required.

\begin{sphinxadmonition}{note}{Known Limitations}

\sphinxAtStartPar
In addition to portability issues with several packages and external
libraries, the classic Intel compilers are currently not able to
compile the googletest libraries and thus enabling the \sphinxcode{\sphinxupquote{\sphinxhyphen{}DENABLE\_TESTING}}
option will result in compilation failure.  The LLVM based compilers
are compatible.
\end{sphinxadmonition}

\begin{sphinxadmonition}{note}{Note:}
\sphinxAtStartPar
This is work in progress and you should contact the LAMMPS developers
via GitHub or the \sphinxhref{https://matsci.org/c/lammps/lammps-development/}{LAMMPS forum at MatSci},
if you have questions or LAMMPS specific problems.
\end{sphinxadmonition}


\subsection{Using a cross\sphinxhyphen{}compiler}
\label{\detokenize{Build_windows:using-a-cross-compiler}}\label{\detokenize{Build_windows:cross}}
\sphinxAtStartPar
If you need to provide custom LAMMPS binaries for Windows, but do not
need to do the compilation on Windows, please consider using a Linux to
Windows cross\sphinxhyphen{}compiler.  This is how currently the Windows binary
packages are created by the LAMMPS developers.  Because of that, this is
probably the currently best tested and supported way to build LAMMPS
executables for Windows.  A CMake preset selecting all packages
compatible with this cross\sphinxhyphen{}compilation build is provided.  The GPU
package can only be compiled with OpenCL support.  To compile with MPI
support, a pre\sphinxhyphen{}compiled library and the corresponding header files are
required.  When building with CMake the matching package will be
downloaded automatically, but MPI support has to be explicitly enabled
with \sphinxcode{\sphinxupquote{\sphinxhyphen{}DBUILD\_MPI=on}}.

\sphinxAtStartPar
Please keep in mind, though, that this only applies to \sphinxstylestrong{compiling} LAMMPS.
Whether the resulting binaries do work correctly is rarely tested by the
LAMMPS developers.  We instead rely on the feedback of the users
of these pre\sphinxhyphen{}compiled LAMMPS packages for Windows.  We will try to resolve
issues to the best of our abilities if we become aware of them. However
this is subject to time constraints and focus on HPC platforms.

\sphinxstepscope


\section{Notes for saving disk space when building LAMMPS from source}
\label{\detokenize{Build_diskspace:notes-for-saving-disk-space-when-building-lammps-from-source}}\label{\detokenize{Build_diskspace::doc}}
\sphinxAtStartPar
LAMMPS is a large software project with a large number of source files,
extensive documentation, and a large collection of example files.
When downloading LAMMPS by cloning the
\sphinxhref{https://github.com/lammps/lammps}{git repository from GitHub} this
will by default also download the entire commit history since September 2006.
Compiling LAMMPS will add the storage requirements of the compiled object
files and libraries to the tally.

\sphinxAtStartPar
In a user account on an HPC cluster with filesystem quotas or in other
environments with restricted disk space capacity it may be needed to
reduce the storage requirements. Here are some suggestions:
\begin{itemize}
\item {} 
\sphinxAtStartPar
Create a so\sphinxhyphen{}called shallow repository by cloning only the last commit
instead of the full project history by using \sphinxcode{\sphinxupquote{git clone git@github.com:lammps/lammps \sphinxhyphen{}\sphinxhyphen{}depth=1 \sphinxhyphen{}\sphinxhyphen{}branch=develop}}.
This reduces the downloaded size to about half.  With \sphinxcode{\sphinxupquote{\sphinxhyphen{}\sphinxhyphen{}depth=1}} it is not possible to check out different
versions/branches of LAMMPS, using \sphinxcode{\sphinxupquote{\sphinxhyphen{}\sphinxhyphen{}depth=1000}} will make multiple recent versions available at little
extra storage needs (the entire git history had nearly 30,000 commits in fall 2021).

\item {} 
\sphinxAtStartPar
Download a tar archive from either the \sphinxhref{https://www.lammps.org/download.html}{download section on the LAMMPS homepage}
or from the \sphinxhref{https://github.com/lammps/lammps/releases}{LAMMPS releases page on GitHub} these will not
contain the git history at all.

\item {} 
\sphinxAtStartPar
Build LAMMPS without the debug flag (remove \sphinxcode{\sphinxupquote{\sphinxhyphen{}g}} from the machine makefile or use \sphinxcode{\sphinxupquote{\sphinxhyphen{}DCMAKE\_BUILD\_TYPE=Release}})
or use the \sphinxcode{\sphinxupquote{strip}} command on the LAMMPS executable when no more debugging would be needed.  The strip command
may also be applied to the LAMMPS shared library. The static library may be deleted entirely.

\item {} 
\sphinxAtStartPar
Delete compiled object files and libraries after copying the LAMMPS executable to a permanent location.
When using the traditional build process, one may use \sphinxcode{\sphinxupquote{make clean\sphinxhyphen{}\textless{}machine\textgreater{}}} or \sphinxcode{\sphinxupquote{make clean\sphinxhyphen{}all}}
to delete object files in the src folder.  For CMake based builds, one may use \sphinxcode{\sphinxupquote{make clean}} or just
delete the entire build folder.

\item {} 
\sphinxAtStartPar
The folders containing the documentation tree (doc), the examples (examples) are not needed to build and
run LAMMPS and can be safely deleted.  Some files in the potentials folder are large and may be deleted,
if not needed.  The largest of those files (occupying about 120 MBytes combined) will only be downloaded on
demand, when the corresponding package is installed.

\item {} 
\sphinxAtStartPar
When using the CMake build procedure, the compilation can be done on a (local) scratch storage that will not
count toward the quota.  A local scratch file system may offer the additional benefit of speeding up creating
object files and linking with libraries compared to a networked file system.  Also with CMake (and unlike with
the traditional make) it is possible to compile LAMMPS executables with different settings and packages included
from the same source tree since all the configuration information is stored in the build folder.  So it is
not necessary to have multiple copies of LAMMPS.

\end{itemize}

\sphinxstepscope


\section{Development build options}
\label{\detokenize{Build_development:development-build-options}}\label{\detokenize{Build_development::doc}}
\sphinxAtStartPar
The build procedures in LAMMPS offers a few extra options which are
useful during development, testing or debugging.


\bigskip\hrule\bigskip



\subsection{Monitor compilation flags (CMake only)}
\label{\detokenize{Build_development:monitor-compilation-flags-cmake-only}}\label{\detokenize{Build_development:compilation}}
\sphinxAtStartPar
Sometimes it is necessary to verify the complete sequence of compilation flags
generated by the CMake build. To enable a more verbose output during
compilation you can use the following option.

\begin{sphinxVerbatim}[commandchars=\\\{\}]
\PYGZhy{}D\PYG{+w}{ }\PYG{n+nv}{CMAKE\PYGZus{}VERBOSE\PYGZus{}MAKEFILE}\PYG{o}{=}value\PYG{+w}{    }\PYG{c+c1}{\PYGZsh{} value = no (default) or yes}
\end{sphinxVerbatim}

\sphinxAtStartPar
Another way of doing this without reconfiguration is calling make with
variable VERBOSE set to 1:

\begin{sphinxVerbatim}[commandchars=\\\{\}]
make\PYG{+w}{ }\PYG{n+nv}{VERBOSE}\PYG{o}{=}\PYG{l+m}{1}
\end{sphinxVerbatim}


\bigskip\hrule\bigskip



\subsection{Enable static code analysis with clang\sphinxhyphen{}tidy (CMake only)}
\label{\detokenize{Build_development:enable-static-code-analysis-with-clang-tidy-cmake-only}}\label{\detokenize{Build_development:clang-tidy}}
\sphinxAtStartPar
The \sphinxhref{https://clang.llvm.org/extra/clang-tidy/}{clang\sphinxhyphen{}tidy tool} is a
static code analysis tool to diagnose (and potentially fix) typical
programming errors or coding style violations.  It has a modular framework
of tests that can be adjusted to help identifying problems before they
become bugs and also assist in modernizing large code bases (like LAMMPS).
It can be enabled for all C++ code with the following CMake flag

\begin{sphinxVerbatim}[commandchars=\\\{\}]
\PYGZhy{}D\PYG{+w}{ }\PYG{n+nv}{ENABLE\PYGZus{}CLANG\PYGZus{}TIDY}\PYG{o}{=}value\PYG{+w}{    }\PYG{c+c1}{\PYGZsh{} value = no (default) or yes}
\end{sphinxVerbatim}

\sphinxAtStartPar
With this flag enabled all source files will be processed twice, first to
be compiled and then to be analyzed. Please note that the analysis can be
significantly more time\sphinxhyphen{}consuming than the compilation itself.


\bigskip\hrule\bigskip



\subsection{Report missing and unneeded ‘\#include’ statements (CMake only)}
\label{\detokenize{Build_development:report-missing-and-unneeded-include-statements-cmake-only}}\label{\detokenize{Build_development:iwyu-processing}}
\sphinxAtStartPar
The conventions for how and when to use and order include statements in
LAMMPS are documented in {\hyperref[\detokenize{Modify_style::doc}]{\sphinxcrossref{\DUrole{doc}{LAMMPS programming style}}}}.  To assist with following
these conventions one can use the \sphinxhref{https://include-what-you-use.org/}{Include What You Use tool}.
This tool is still under development and for large and complex projects like LAMMPS
there are some false positives, so suggested changes need to be verified manually.
It is recommended to use at least version 0.16, which has much fewer incorrect
reports than earlier versions.  To install the IWYU toolkit, you need to have
the clang compiler \sphinxstylestrong{and} its development package installed.  Download the IWYU
version that matches the version of the clang compiler, configure, build, and
install it.

\sphinxAtStartPar
The necessary steps to generate the report can be enabled via a CMake variable
during CMake configuration.

\begin{sphinxVerbatim}[commandchars=\\\{\}]
\PYGZhy{}D\PYG{+w}{ }\PYG{n+nv}{ENABLE\PYGZus{}IWYU}\PYG{o}{=}value\PYG{+w}{    }\PYG{c+c1}{\PYGZsh{} value = no (default) or yes}
\end{sphinxVerbatim}

\sphinxAtStartPar
This will check if the required binary (include\sphinxhyphen{}what\sphinxhyphen{}you\sphinxhyphen{}use or iwyu)
and python script script (iwyu\sphinxhyphen{}tool or iwyu\_tool or iwyu\_tool.py) can
be found in the path.  The analysis can then be started with:

\begin{sphinxVerbatim}[commandchars=\\\{\}]
make\PYG{+w}{ }iwyu
\end{sphinxVerbatim}

\sphinxAtStartPar
This may first run some compilation, as the analysis is dependent
on recording all commands required to do the compilation.


\bigskip\hrule\bigskip



\subsection{Address, Leak, Undefined Behavior, and Thread Sanitizer Support (CMake only)}
\label{\detokenize{Build_development:address-leak-undefined-behavior-and-thread-sanitizer-support-cmake-only}}\label{\detokenize{Build_development:sanitizer}}
\sphinxAtStartPar
Compilers such as GCC and Clang support generating instrumented binaries
which use different sanitizer libraries to detect problems in the code
during run\sphinxhyphen{}time. They can detect issues like:
\begin{itemize}
\item {} 
\sphinxAtStartPar
\sphinxhref{https://clang.llvm.org/docs/AddressSanitizer.html\#memory-leak-detection}{memory leaks}

\item {} 
\sphinxAtStartPar
\sphinxhref{https://clang.llvm.org/docs/UndefinedBehaviorSanitizer.html}{undefined behavior}

\item {} 
\sphinxAtStartPar
\sphinxhref{https://clang.llvm.org/docs/ThreadSanitizer.html}{data races}

\end{itemize}

\sphinxAtStartPar
Please note that this kind of instrumentation usually comes with a
performance hit (but much less than using tools like \sphinxhref{https://valgrind.org}{Valgrind} with a more low level approach).  To enable
these features, additional compiler flags need to be added to the
compilation and linking stages.  This is done through setting the
\sphinxcode{\sphinxupquote{ENABLE\_SANITIZER}} variable during configuration. Examples:

\begin{sphinxVerbatim}[commandchars=\\\{\}]
\PYGZhy{}D\PYG{+w}{ }\PYG{n+nv}{ENABLE\PYGZus{}SANITIZER}\PYG{o}{=}none\PYG{+w}{       }\PYG{c+c1}{\PYGZsh{} no sanitizer active (default)}
\PYGZhy{}D\PYG{+w}{ }\PYG{n+nv}{ENABLE\PYGZus{}SANITIZER}\PYG{o}{=}address\PYG{+w}{    }\PYG{c+c1}{\PYGZsh{} enable address sanitizer / memory leak checker}
\PYGZhy{}D\PYG{+w}{ }\PYG{n+nv}{ENABLE\PYGZus{}SANITIZER}\PYG{o}{=}hwaddress\PYG{+w}{  }\PYG{c+c1}{\PYGZsh{} enable hardware assisted address sanitizer / memory leak checker}
\PYGZhy{}D\PYG{+w}{ }\PYG{n+nv}{ENABLE\PYGZus{}SANITIZER}\PYG{o}{=}leak\PYG{+w}{       }\PYG{c+c1}{\PYGZsh{} enable memory leak checker (only)}
\PYGZhy{}D\PYG{+w}{ }\PYG{n+nv}{ENABLE\PYGZus{}SANITIZER}\PYG{o}{=}undefined\PYG{+w}{  }\PYG{c+c1}{\PYGZsh{} enable undefined behavior sanitizer}
\PYGZhy{}D\PYG{+w}{ }\PYG{n+nv}{ENABLE\PYGZus{}SANITIZER}\PYG{o}{=}thread\PYG{+w}{     }\PYG{c+c1}{\PYGZsh{} enable thread sanitizer}
\end{sphinxVerbatim}


\bigskip\hrule\bigskip



\subsection{Code Coverage and Unit Testing (CMake only)}
\label{\detokenize{Build_development:code-coverage-and-unit-testing-cmake-only}}\label{\detokenize{Build_development:testing}}
\sphinxAtStartPar
The LAMMPS code is subject to multiple levels of automated testing
during development:
\begin{itemize}
\item {} 
\sphinxAtStartPar
Integration testing (i.e. whether the code compiles
on various platforms and with a variety of compilers and settings),

\item {} 
\sphinxAtStartPar
Unit testing (i.e. whether certain functions or classes of the code
produce the expected results for given inputs),

\item {} 
\sphinxAtStartPar
Run testing (i.e. whether selected input decks can run to completion
without crashing for multiple configurations),

\item {} 
\sphinxAtStartPar
Regression testing (i.e. whether selected input examples reproduce the
same results over a given number of steps and operations within a
given error margin).

\end{itemize}

\sphinxAtStartPar
The status of this automated testing can be viewed on \sphinxurl{https://ci.lammps.org}.

\sphinxAtStartPar
The scripts and inputs for integration, run, and regression testing
are maintained in a
\sphinxhref{https://github.com/lammps/lammps-testing}{separate repository}
of the LAMMPS project on GitHub.  A few tests are also run as GitHub
Actions and their configuration files are in the \sphinxcode{\sphinxupquote{.github/workflows/}}
folder of the LAMMPS git tree.

\sphinxAtStartPar
The unit testing facility is integrated into the CMake build process of
the LAMMPS source code distribution itself.  It can be enabled by
setting \sphinxcode{\sphinxupquote{\sphinxhyphen{}D ENABLE\_TESTING=on}} during the CMake configuration step.
It requires the \sphinxhref{https://pyyaml.org/}{YAML} library and matching
development headers to compile (if those are not found locally a recent
version of that library will be downloaded and compiled along with
LAMMPS and the test programs) and will download and compile a specific
version of the \sphinxhref{https://github.com/google/googletest/}{GoogleTest} C++
test framework that is used to implement the tests.  Those unit tests
may be combined with memory access and leak checking with valgrind
(see below for how to enable it).  In that case, running so\sphinxhyphen{}called
death tests will create a lot of false positives and thus they can be
disabled by configuring compilation with the additional setting
\sphinxcode{\sphinxupquote{\sphinxhyphen{}D SKIP\_DEATH\_TESTS=on}}.

\begin{sphinxadmonition}{note}{Software version and LAMMPS configuration requirements}

\sphinxAtStartPar
The compiler and library version requirements for the testing
framework are more strict than for the main part of LAMMPS.  For
example the default GNU C++ and Fortran compilers of RHEL/CentOS 7.x
(version 4.8.x) are not sufficient.  The CMake configuration will try
to detect incompatible versions and either skip incompatible tests or
stop with an error.  Also the number of available tests will depend on
installed LAMMPS packages, development environment, operating system,
and configuration settings.
\end{sphinxadmonition}

\sphinxAtStartPar
After compilation is complete, the unit testing is started in the build
folder using the \sphinxcode{\sphinxupquote{ctest}} command, which is part of the CMake software.
The output of this command will be looking something like this:

\begin{sphinxVerbatim}[commandchars=\\\{\}]
\PYG{g+gp}{   \PYGZdl{} }ctest
\PYG{g+go}{   Test project /home/akohlmey/compile/lammps/build\PYGZhy{}testing}
\PYG{g+go}{        Start   1: RunLammps}
\PYG{g+go}{  1/563 Test   \PYGZsh{}1: RunLammps ..................................   Passed    0.28 sec}
\PYG{g+go}{        Start   2: HelpMessage}
\PYG{g+go}{  2/563 Test   \PYGZsh{}2: HelpMessage ................................   Passed    0.06 sec}
\PYG{g+go}{        Start   3: InvalidFlag}
\PYG{g+go}{  3/563 Test   \PYGZsh{}3: InvalidFlag ................................   Passed    0.06 sec}
\PYG{g+go}{        Start   4: Tokenizer}
\PYG{g+go}{  4/563 Test   \PYGZsh{}4: Tokenizer ..................................   Passed    0.05 sec}
\PYG{g+go}{        Start   5: MemPool}
\PYG{g+go}{  5/563 Test   \PYGZsh{}5: MemPool ....................................   Passed    0.05 sec}
\PYG{g+go}{        Start   6: ArgUtils}
\PYG{g+go}{  6/563 Test   \PYGZsh{}6: ArgUtils ...................................   Passed    0.05 sec}
\PYG{g+go}{      [...]}
\PYG{g+go}{        Start 561: ImproperStyle:zero}
\PYG{g+go}{561/563 Test \PYGZsh{}561: ImproperStyle:zero .........................   Passed    0.07 sec}
\PYG{g+go}{        Start 562: TestMliapPyUnified}
\PYG{g+go}{562/563 Test \PYGZsh{}562: TestMliapPyUnified .........................   Passed    0.16 sec}
\PYG{g+go}{        Start 563: TestPairList}
\PYG{g+go}{563/563 Test \PYGZsh{}563: TestPairList ...............................   Passed    0.06 sec}

\PYG{g+go}{100\PYGZpc{} tests passed, 0 tests failed out of 563}

\PYG{g+go}{Label Time Summary:}
\PYG{g+go}{generated    =   0.85 sec*proc (3 tests)}
\PYG{g+go}{noWindows    =   4.16 sec*proc (2 tests)}
\PYG{g+go}{slow         =  78.33 sec*proc (67 tests)}
\PYG{g+go}{unstable     =  28.23 sec*proc (34 tests)}

\PYG{g+go}{Total Test time (real) = 132.34 sec}
\end{sphinxVerbatim}

\sphinxAtStartPar
The \sphinxcode{\sphinxupquote{ctest}} command has many options, the most important ones are:


\begin{savenotes}\sphinxattablestart
\sphinxthistablewithglobalstyle
\centering
\begin{tabular}[t]{\X{20}{100}\X{80}{100}}
\sphinxtoprule
\sphinxtableatstartofbodyhook
\sphinxAtStartPar
Option
&
\sphinxAtStartPar
Function
\\
\sphinxhline
\sphinxAtStartPar
\sphinxcode{\sphinxupquote{\sphinxhyphen{}V}}
&
\sphinxAtStartPar
verbose output: display output of individual test runs
\\
\sphinxhline
\sphinxAtStartPar
\sphinxcode{\sphinxupquote{\sphinxhyphen{}j \textless{}num\textgreater{}}}
&
\sphinxAtStartPar
parallel run: run \textless{}num\textgreater{} tests in parallel
\\
\sphinxhline
\sphinxAtStartPar
\sphinxcode{\sphinxupquote{\sphinxhyphen{}R \textless{}regex\textgreater{}}}
&
\sphinxAtStartPar
run subset of tests matching the regular expression \textless{}regex\textgreater{}
\\
\sphinxhline
\sphinxAtStartPar
\sphinxcode{\sphinxupquote{\sphinxhyphen{}E \textless{}regex\textgreater{}}}
&
\sphinxAtStartPar
exclude subset of tests matching the regular expression \textless{}regex\textgreater{}
\\
\sphinxhline
\sphinxAtStartPar
\sphinxcode{\sphinxupquote{\sphinxhyphen{}L \textless{}regex\textgreater{}}}
&
\sphinxAtStartPar
run subset of tests with a label matching the regular expression \textless{}regex\textgreater{}
\\
\sphinxhline
\sphinxAtStartPar
\sphinxcode{\sphinxupquote{\sphinxhyphen{}LE \textless{}regex\textgreater{}}}
&
\sphinxAtStartPar
exclude subset of tests with a label matching the regular expression \textless{}regex\textgreater{}
\\
\sphinxhline
\sphinxAtStartPar
\sphinxcode{\sphinxupquote{\sphinxhyphen{}N}}
&
\sphinxAtStartPar
dry\sphinxhyphen{}run: display list of tests without running them
\\
\sphinxhline
\sphinxAtStartPar
\sphinxcode{\sphinxupquote{\sphinxhyphen{}T memcheck}}
&
\sphinxAtStartPar
run tests with valgrind memory checker (if available)
\\
\sphinxbottomrule
\end{tabular}
\sphinxtableafterendhook\par
\sphinxattableend\end{savenotes}

\sphinxAtStartPar
In its full implementation, the unit test framework will consist of multiple
kinds of tests implemented in different programming languages (C++, C, Python,
Fortran) and testing different aspects of the LAMMPS software and its features.
The tests will adapt to the compilation settings of LAMMPS, so that tests
will be skipped if prerequisite features are not available in LAMMPS.

\begin{sphinxadmonition}{note}{Work in Progress}

\sphinxAtStartPar
The unit test framework was added in spring 2020 and is under active
development.  The coverage is not complete and will be expanded over
time.  Preference is given to parts of the code base that are easy to
test or commonly used.
\end{sphinxadmonition}

\sphinxAtStartPar
Tests as shown by the \sphinxcode{\sphinxupquote{ctest}} program are command lines defined in the
\sphinxcode{\sphinxupquote{CMakeLists.txt}} files in the \sphinxcode{\sphinxupquote{unittest}} directory tree.  A few
tests simply execute LAMMPS with specific command line flags and check
the output to the screen for expected content.  A large number of unit
tests are special tests programs using the \sphinxhref{https://github.com/google/googletest/}{GoogleTest framework} and linked to the LAMMPS
library that test individual functions or create a LAMMPS class
instance, execute one or more commands and check data inside the LAMMPS
class hierarchy.  There are also tests for the C\sphinxhyphen{}library, Fortran, and
Python module interfaces to LAMMPS.  The Python tests use the Python
“unittest” module in a similar fashion than the others use \sphinxtitleref{GoogleTest}.
These special test programs are structured to perform multiple
individual tests internally and each of those contains several checks
(aka assertions) for internal data being changed as expected.

\sphinxAtStartPar
Tests for force computing or modifying styles (e.g. styles for non\sphinxhyphen{}bonded
and bonded interactions and selected fixes) are run by using a more generic
test program that reads its input from files in YAML format. The YAML file
provides the information on how to customized the test program to test
a specific style and \sphinxhyphen{} if needed \sphinxhyphen{} with specific settings.
To add a test for another, similar style (e.g. a new pair style) it is
usually sufficient to add a suitable YAML file.  {\hyperref[\detokenize{Developer_unittest::doc}]{\sphinxcrossref{\DUrole{doc}{Detailed
instructions for adding tests}}}} are provided in the
Programmer Guide part of the manual.  A description of what happens
during the tests is given below.


\subsubsection{Unit tests for force styles}
\label{\detokenize{Build_development:unit-tests-for-force-styles}}
\sphinxAtStartPar
A large part of LAMMPS are different “styles” for computing non\sphinxhyphen{}bonded
and bonded interactions selected through the {\hyperref[\detokenize{pair_style::doc}]{\sphinxcrossref{\DUrole{doc}{pair\_style command}}}},
{\hyperref[\detokenize{bond_style::doc}]{\sphinxcrossref{\DUrole{doc}{bond\_style command}}}}, {\hyperref[\detokenize{angle_style::doc}]{\sphinxcrossref{\DUrole{doc}{angle\_style command}}}}, {\hyperref[\detokenize{dihedral_style::doc}]{\sphinxcrossref{\DUrole{doc}{dihedral\_style command}}}},
{\hyperref[\detokenize{improper_style::doc}]{\sphinxcrossref{\DUrole{doc}{improper\_style command}}}}, and {\hyperref[\detokenize{kspace_style::doc}]{\sphinxcrossref{\DUrole{doc}{kspace\_style command}}}}.  Since these all share
common interfaces, it is possible to write generic test programs that
will call those common interfaces for small test systems with less than
100 atoms and compare the results with pre\sphinxhyphen{}recorded reference results.
A test run is then a a collection multiple individual test runs each
with many comparisons to reference results based on template input
files, individual command settings, relative error margins, and
reference data stored in a YAML format file with \sphinxcode{\sphinxupquote{.yaml}}
suffix. Currently the programs \sphinxcode{\sphinxupquote{test\_pair\_style}}, \sphinxcode{\sphinxupquote{test\_bond\_style}},
\sphinxcode{\sphinxupquote{test\_angle\_style}}, \sphinxcode{\sphinxupquote{test\_dihedral\_style}}, and
\sphinxcode{\sphinxupquote{test\_improper\_style}} are implemented.  They will compare forces,
energies and (global) stress for all atoms after a \sphinxcode{\sphinxupquote{run 0}} calculation
and after a few steps of MD with {\hyperref[\detokenize{fix_nve::doc}]{\sphinxcrossref{\DUrole{doc}{fix nve}}}}, each in
multiple variants with different settings and also for multiple
accelerated styles. If a prerequisite style or package is missing, the
individual tests are skipped.  All force style tests will be executed on
a single MPI process, so using the CMake option \sphinxcode{\sphinxupquote{\sphinxhyphen{}D BUILD\_MPI=off}} can
significantly speed up testing, since this will skip the MPI
initialization for each test run.  Below is an example command and
output:

\begin{sphinxVerbatim}[commandchars=\\\{\}]
\PYG{g+gp}{\PYGZdl{} }test\PYGZus{}pair\PYGZus{}style\PYG{+w}{ }mol\PYGZhy{}pair\PYGZhy{}lj\PYGZus{}cut.yaml
\PYG{g+go}{[==========] Running 6 tests from 1 test suite.}
\PYG{g+go}{[\PYGZhy{}\PYGZhy{}\PYGZhy{}\PYGZhy{}\PYGZhy{}\PYGZhy{}\PYGZhy{}\PYGZhy{}\PYGZhy{}\PYGZhy{}] Global test environment set\PYGZhy{}up.}
\PYG{g+go}{[\PYGZhy{}\PYGZhy{}\PYGZhy{}\PYGZhy{}\PYGZhy{}\PYGZhy{}\PYGZhy{}\PYGZhy{}\PYGZhy{}\PYGZhy{}] 6 tests from PairStyle}
\PYG{g+go}{[ RUN      ] PairStyle.plain}
\PYG{g+go}{[       OK ] PairStyle.plain (24 ms)}
\PYG{g+go}{[ RUN      ] PairStyle.omp}
\PYG{g+go}{[       OK ] PairStyle.omp (18 ms)}
\PYG{g+go}{[ RUN      ] PairStyle.intel}
\PYG{g+go}{[       OK ] PairStyle.intel (6 ms)}
\PYG{g+go}{[ RUN      ] PairStyle.opt}
\PYG{g+go}{[  SKIPPED ] PairStyle.opt (0 ms)}
\PYG{g+go}{[ RUN      ] PairStyle.single}
\PYG{g+go}{[       OK ] PairStyle.single (7 ms)}
\PYG{g+go}{[ RUN      ] PairStyle.extract}
\PYG{g+go}{[       OK ] PairStyle.extract (6 ms)}
\PYG{g+go}{[\PYGZhy{}\PYGZhy{}\PYGZhy{}\PYGZhy{}\PYGZhy{}\PYGZhy{}\PYGZhy{}\PYGZhy{}\PYGZhy{}\PYGZhy{}] 6 tests from PairStyle (62 ms total)}

\PYG{g+go}{[\PYGZhy{}\PYGZhy{}\PYGZhy{}\PYGZhy{}\PYGZhy{}\PYGZhy{}\PYGZhy{}\PYGZhy{}\PYGZhy{}\PYGZhy{}] Global test environment tear\PYGZhy{}down}
\PYG{g+go}{[==========] 6 tests from 1 test suite ran. (63 ms total)}
\PYG{g+go}{[  PASSED  ] 5 tests.}
\PYG{g+go}{[  SKIPPED ] 1 test, listed below:}
\PYG{g+go}{[  SKIPPED ] PairStyle.opt}
\end{sphinxVerbatim}

\sphinxAtStartPar
In this particular case, 5 out of 6 sets of tests were conducted, the
tests for the \sphinxcode{\sphinxupquote{lj/cut/opt}} pair style was skipped, since the tests
executable did not include it.  To learn what individual tests are performed,
you (currently) need to read the source code.  You can use code coverage
recording (see next section) to confirm how well the tests cover the code
paths in the individual source files.

\sphinxAtStartPar
The force style test programs have a common set of options:


\begin{savenotes}\sphinxattablestart
\sphinxthistablewithglobalstyle
\centering
\begin{tabular}[t]{\X{25}{100}\X{75}{100}}
\sphinxtoprule
\sphinxtableatstartofbodyhook
\sphinxAtStartPar
Option
&
\sphinxAtStartPar
Function
\\
\sphinxhline
\sphinxAtStartPar
\sphinxcode{\sphinxupquote{\sphinxhyphen{}g \textless{}newfile\textgreater{}}}
&
\sphinxAtStartPar
regenerate reference data in new YAML file
\\
\sphinxhline
\sphinxAtStartPar
\sphinxcode{\sphinxupquote{\sphinxhyphen{}u}}
&
\sphinxAtStartPar
update reference data in the original YAML file
\\
\sphinxhline
\sphinxAtStartPar
\sphinxcode{\sphinxupquote{\sphinxhyphen{}s}}
&
\sphinxAtStartPar
print error statistics for each group of comparisons
\\
\sphinxhline
\sphinxAtStartPar
\sphinxcode{\sphinxupquote{\sphinxhyphen{}v}}
&
\sphinxAtStartPar
verbose output: also print the executed LAMMPS commands
\\
\sphinxbottomrule
\end{tabular}
\sphinxtableafterendhook\par
\sphinxattableend\end{savenotes}

\sphinxAtStartPar
The \sphinxcode{\sphinxupquote{ctest}} tool has no mechanism to directly pass flags to the individual
test programs, but a workaround has been implemented where these flags can be
set in an environment variable \sphinxcode{\sphinxupquote{TEST\_ARGS}}. Example:

\begin{sphinxVerbatim}[commandchars=\\\{\}]
env\PYG{+w}{ }\PYG{n+nv}{TEST\PYGZus{}ARGS}\PYG{o}{=}\PYGZhy{}s\PYG{+w}{ }ctest\PYG{+w}{ }\PYGZhy{}V\PYG{+w}{ }\PYGZhy{}R\PYG{+w}{ }BondStyle
\end{sphinxVerbatim}

\sphinxAtStartPar
To add a test for a style that is not yet covered, it is usually best
to copy a YAML file for a similar style to a new file, edit the details
of the style (how to call it, how to set its coefficients) and then
run test command with either the \sphinxcode{\sphinxupquote{\sphinxhyphen{}g}} and the replace the initial
test file with the regenerated one or the \sphinxcode{\sphinxupquote{\sphinxhyphen{}u}} option.  The \sphinxcode{\sphinxupquote{\sphinxhyphen{}u}} option
will destroy the original file, if the generation run does not complete,
so using \sphinxcode{\sphinxupquote{\sphinxhyphen{}g}} is recommended unless the YAML file is fully tested
and working.

\sphinxAtStartPar
Some of the force style tests are rather slow to run and some are very
sensitive to small differences like CPU architecture, compiler
toolchain, compiler optimization. Those tests are flagged with a “slow”
and/or “unstable” label, and thus those tests can be selectively
excluded with the \sphinxcode{\sphinxupquote{\sphinxhyphen{}LE}} flag or selected with the \sphinxcode{\sphinxupquote{\sphinxhyphen{}L}} flag.

\begin{sphinxadmonition}{note}{Recommendations and notes for YAML files}
\begin{itemize}
\item {} 
\sphinxAtStartPar
The reference results should be recorded without any code
optimization or related compiler flags enabled.

\item {} 
\sphinxAtStartPar
The \sphinxcode{\sphinxupquote{epsilon}} parameter defines the relative precision with which
the reference results must be met.  The test geometries often have
high and low energy parts and thus a significant impact from
floating\sphinxhyphen{}point math truncation errors is to be expected. Some
functional forms and potentials are more noisy than others, so this
parameter needs to be adjusted. Typically a value around 1.0e\sphinxhyphen{}13
can be used, but it may need to be as large as 1.0e\sphinxhyphen{}8 in some
cases.

\item {} 
\sphinxAtStartPar
The tests for pair styles from OPT, OPENMP and INTEL are
performed with automatically rescaled epsilon to account for
additional loss of precision from code optimizations and different
summation orders.

\item {} 
\sphinxAtStartPar
When compiling with (aggressive) compiler optimization, some tests
are likely to fail.  It is recommended to inspect the individual
tests in detail to decide, whether the specific error for a specific
property is acceptable (it often is), or this may be an indication
of mis\sphinxhyphen{}compiled code (or an undesired large loss of precision due
to significant reordering of operations and thus less error cancellation).

\end{itemize}
\end{sphinxadmonition}


\subsubsection{Unit tests for timestepping related fixes}
\label{\detokenize{Build_development:unit-tests-for-timestepping-related-fixes}}
\sphinxAtStartPar
A substantial subset of {\hyperref[\detokenize{fix::doc}]{\sphinxcrossref{\DUrole{doc}{fix styles}}}} are invoked regularly
during MD timestepping and manipulate per\sphinxhyphen{}atom properties like
positions, velocities, and forces.  For those fix styles, testing can be
done in a very similar fashion as for force fields and thus there is a
test program \sphinxtitleref{test\_fix\_timestep} that shares a lot of code, properties,
and command line flags with the force field style testers described in
the previous section.

\sphinxAtStartPar
This tester will set up a small molecular system run with verlet run
style for 4 MD steps, then write a binary restart and continue for
another 4 MD steps. At this point coordinates and velocities are
recorded and compared to reference data. Then the system is cleared,
restarted and running the second 4 MD steps again and the data is
compared to the same reference. That is followed by another restart
after which per atom type masses are replaced with per\sphinxhyphen{}atom masses and
the second 4 MD steps are repeated again and compared to the same
reference.  Also global scalar and vector data of the fix is recorded
and compared.  If the fix is a thermostat and thus the internal property
\sphinxcode{\sphinxupquote{t\_target}} can be extracted, then this is compared to the reference
data.  The tests are repeated with the respa run style.

\sphinxAtStartPar
If the fix has a multi\sphinxhyphen{}threaded version in the OPENMP package, then
the entire set of tests is repeated for that version as well.

\sphinxAtStartPar
For this to work, some additional conditions have to be met by the
YAML format test inputs.
\begin{itemize}
\item {} 
\sphinxAtStartPar
The fix to be tested (and only this fix), should be listed in the
\sphinxcode{\sphinxupquote{prerequisites:}} section

\item {} 
\sphinxAtStartPar
The fix to be tested must be specified in the \sphinxcode{\sphinxupquote{post\_commands:}}
section with the fix\sphinxhyphen{}ID \sphinxcode{\sphinxupquote{test}}.  This section may contain other
commands and other fixes (e.g. an instance of fix nve for testing
a thermostat or force manipulation fix)

\item {} 
\sphinxAtStartPar
For fixes that can tally contributions to the global virial, the
line \sphinxcode{\sphinxupquote{fix\_modify test virial yes}} should be included in the
\sphinxcode{\sphinxupquote{post\_commands:}} section of the test input.

\item {} 
\sphinxAtStartPar
For thermostat fixes the target temperature should be ramped from
an arbitrary value (e.g. 50K) to a pre\sphinxhyphen{}defined target temperature
entered as \sphinxcode{\sphinxupquote{\$\{t\_target\}}}.

\item {} 
\sphinxAtStartPar
For fixes that have thermostatting support included, but do not
have it enabled in the input (e.g. fix rigid with default settings),
the \sphinxcode{\sphinxupquote{post\_commands:}} section should contain the line
\sphinxcode{\sphinxupquote{variable t\_target delete}} to disable the target temperature ramp
check to avoid false positives.

\end{itemize}


\subsubsection{Use custom linker for faster link times when ENABLE\_TESTING is active}
\label{\detokenize{Build_development:use-custom-linker-for-faster-link-times-when-enable-testing-is-active}}
\sphinxAtStartPar
When compiling LAMMPS with enabled tests, most test executables will
need to be linked against the LAMMPS library.  Since this can be a very
large library with many C++ objects when many packages are enabled, link
times can become very long on machines that use the GNU BFD linker (e.g.
Linux systems).  Alternatives like the \sphinxcode{\sphinxupquote{mold}} linker, the \sphinxcode{\sphinxupquote{lld}}
linker of the LLVM project, or the \sphinxcode{\sphinxupquote{gold}} linker available with GNU
binutils can speed up this step substantially (in this order).  CMake
will by default test if any of the three can be enabled and use it when
\sphinxcode{\sphinxupquote{ENABLE\_TESTING}} is active.  It can also be selected manually through
the \sphinxcode{\sphinxupquote{CMAKE\_CUSTOM\_LINKER}} CMake variable.  Allowed values are
\sphinxcode{\sphinxupquote{mold}}, \sphinxcode{\sphinxupquote{lld}}, \sphinxcode{\sphinxupquote{gold}}, \sphinxcode{\sphinxupquote{bfd}}, or \sphinxcode{\sphinxupquote{default}}.  The \sphinxcode{\sphinxupquote{default}}
option will use the system default linker otherwise, the linker is
chosen explicitly.  This option is only available for the GNU or Clang
C++ compilers.


\subsubsection{Tests for other components and utility functions}
\label{\detokenize{Build_development:tests-for-other-components-and-utility-functions}}
\sphinxAtStartPar
Additional tests that validate utility functions or specific components
of LAMMPS are implemented as standalone executable which may, or may not
require creating a suitable LAMMPS instance.  These tests are more specific
and do not require YAML format input files.  To add a test, either an
existing source file needs to be extended or a new file added, which in turn
requires additions to the \sphinxcode{\sphinxupquote{CMakeLists.txt}} file in the source folder.


\subsubsection{Collect and visualize code coverage metrics}
\label{\detokenize{Build_development:collect-and-visualize-code-coverage-metrics}}
\sphinxAtStartPar
You can also collect code coverage metrics while running LAMMPS or the
tests by enabling code coverage support during the CMake configuration:

\begin{sphinxVerbatim}[commandchars=\\\{\}]
\PYGZhy{}D\PYG{+w}{ }\PYG{n+nv}{ENABLE\PYGZus{}COVERAGE}\PYG{o}{=}on\PYG{+w}{  }\PYG{c+c1}{\PYGZsh{} enable coverage measurements (off by default)}
\end{sphinxVerbatim}

\sphinxAtStartPar
This will instrument all object files to write information about which
lines of code were accessed during execution in files next to the
corresponding object files.  These can be post\sphinxhyphen{}processed to visually
show the degree of coverage and which code paths are accessed and which
are not taken.  When working on unit tests (see above), this can be
extremely helpful to determine which parts of the code are not executed
and thus what kind of tests are still missing. The coverage data is
cumulative, i.e. new data is added with each new run.

\sphinxAtStartPar
Enabling code coverage will also add the following build targets to
generate coverage reports after running the LAMMPS executable or the
unit tests:

\begin{sphinxVerbatim}[commandchars=\\\{\}]
make\PYG{+w}{ }gen\PYGZus{}coverage\PYGZus{}html\PYG{+w}{   }\PYG{c+c1}{\PYGZsh{} generate coverage report in HTML format}
make\PYG{+w}{ }gen\PYGZus{}coverage\PYGZus{}xml\PYG{+w}{    }\PYG{c+c1}{\PYGZsh{} generate coverage report in XML format}
make\PYG{+w}{ }clean\PYGZus{}coverage\PYGZus{}html\PYG{+w}{ }\PYG{c+c1}{\PYGZsh{} delete folder with HTML format coverage report}
make\PYG{+w}{ }reset\PYGZus{}coverage\PYG{+w}{      }\PYG{c+c1}{\PYGZsh{} delete all collected coverage data and HTML output}
\end{sphinxVerbatim}

\sphinxAtStartPar
These reports require \sphinxhref{https://gcovr.com/}{GCOVR} to be installed. The easiest way
to do this to install it via pip:

\begin{sphinxVerbatim}[commandchars=\\\{\}]
pip\PYG{+w}{ }install\PYG{+w}{ }git+https://github.com/gcovr/gcovr.git
\end{sphinxVerbatim}

\sphinxAtStartPar
After post\sphinxhyphen{}processing with \sphinxcode{\sphinxupquote{gen\_coverage\_html}} the results are in
a folder \sphinxcode{\sphinxupquote{coverage\_html}} and can be viewed with a web browser.
The images below illustrate how the data is presented.

\begin{figure}[htbp]
\centering
\capstart

\noindent\sphinxincludegraphics[width=0.600\linewidth]{{coverage-overview-top}.png}
\caption{Top of the overview page}\label{\detokenize{Build_development:id1}}\end{figure}

\begin{figure}[htbp]
\centering
\capstart

\noindent\sphinxincludegraphics[width=0.600\linewidth]{{coverage-overview-manybody}.png}
\caption{Styles with good coverage}\label{\detokenize{Build_development:id2}}\end{figure}

\begin{figure}[htbp]
\centering
\capstart

\noindent\sphinxincludegraphics[width=0.600\linewidth]{{coverage-file-top}.png}
\caption{Top of individual source page}\label{\detokenize{Build_development:id3}}\end{figure}

\begin{figure}[htbp]
\centering
\capstart

\noindent\sphinxincludegraphics[width=0.600\linewidth]{{coverage-file-branches}.png}
\caption{Source page with branches}\label{\detokenize{Build_development:id4}}\end{figure}


\subsection{Coding style utilities}
\label{\detokenize{Build_development:coding-style-utilities}}
\sphinxAtStartPar
To aid with enforcing some of the coding style conventions in LAMMPS
some additional build targets have been added. These require Python 3.5
or later and will only work properly on Unix\sphinxhyphen{}like operating and file systems.

\sphinxAtStartPar
The following options are available.

\begin{sphinxVerbatim}[commandchars=\\\{\}]
make\PYG{+w}{ }check\PYGZhy{}whitespace\PYG{+w}{    }\PYG{c+c1}{\PYGZsh{} search for files with whitespace issues}
make\PYG{+w}{ }fix\PYGZhy{}whitespace\PYG{+w}{      }\PYG{c+c1}{\PYGZsh{} correct whitespace issues in files}
make\PYG{+w}{ }check\PYGZhy{}homepage\PYG{+w}{      }\PYG{c+c1}{\PYGZsh{} search for files with old LAMMPS homepage URLs}
make\PYG{+w}{ }fix\PYGZhy{}homepage\PYG{+w}{        }\PYG{c+c1}{\PYGZsh{} correct LAMMPS homepage URLs in files}
make\PYG{+w}{ }check\PYGZhy{}errordocs\PYG{+w}{     }\PYG{c+c1}{\PYGZsh{} search for deprecated error docs in header files}
make\PYG{+w}{ }fix\PYGZhy{}errordocs\PYG{+w}{       }\PYG{c+c1}{\PYGZsh{} remove error docs in header files}
make\PYG{+w}{ }check\PYGZhy{}permissions\PYG{+w}{   }\PYG{c+c1}{\PYGZsh{} search for files with permissions issues}
make\PYG{+w}{ }fix\PYGZhy{}permissions\PYG{+w}{     }\PYG{c+c1}{\PYGZsh{} correct permissions issues in files}
make\PYG{+w}{ }check\PYGZhy{}docs\PYG{+w}{          }\PYG{c+c1}{\PYGZsh{} search for several issues in the manual}
make\PYG{+w}{ }check\PYGZhy{}version\PYG{+w}{       }\PYG{c+c1}{\PYGZsh{} list files with pending release version tags}
make\PYG{+w}{ }check\PYG{+w}{               }\PYG{c+c1}{\PYGZsh{} run all check targets from above}
\end{sphinxVerbatim}

\sphinxAtStartPar
These should help to make source and documentation files conforming
to some the coding style preferences of the LAMMPS developers.


\subsection{Clang\sphinxhyphen{}format support}
\label{\detokenize{Build_development:clang-format-support}}\label{\detokenize{Build_development:clang-format}}
\sphinxAtStartPar
For the code in the \sphinxcode{\sphinxupquote{unittest}} and \sphinxcode{\sphinxupquote{src}} trees we are transitioning
to use the \sphinxtitleref{clang\sphinxhyphen{}format} tool to assist with having a consistent source
code formatting style.  The \sphinxtitleref{clang\sphinxhyphen{}format} command bundled with Clang
version 8.0 or later is required.  The configuration is in files called
\sphinxcode{\sphinxupquote{.clang\sphinxhyphen{}format}} in the respective folders.  Since the modifications
from \sphinxtitleref{clang\sphinxhyphen{}format} can be significant and \sphinxhyphen{} especially for “legacy
style code” \sphinxhyphen{} they are not always improving readability, a large number
of files currently have a \sphinxcode{\sphinxupquote{// clang\sphinxhyphen{}format off}} at the top, which will
disable the processing.  As of fall 2021 all files have been either
“protected” this way or are enabled for full or partial \sphinxtitleref{clang\sphinxhyphen{}format}
processing.  Over time, the “protected” files will be refactored and
updated so that \sphinxtitleref{clang\sphinxhyphen{}format} may be applied to them as well.

\sphinxAtStartPar
It is recommended for all newly contributed files to use the clang\sphinxhyphen{}format
processing while writing the code or do the coding style processing
(including the scripts mentioned in the previous paragraph)

\sphinxAtStartPar
If \sphinxtitleref{clang\sphinxhyphen{}format} is available, files can be updated individually with
commands like the following:

\begin{sphinxVerbatim}[commandchars=\\\{\}]
clang\PYGZhy{}format\PYG{+w}{ }\PYGZhy{}i\PYG{+w}{ }some\PYGZus{}file.cpp
\end{sphinxVerbatim}

\sphinxAtStartPar
The following target are available for both, GNU make and CMake:

\begin{sphinxVerbatim}[commandchars=\\\{\}]
make\PYG{+w}{ }format\PYGZhy{}src\PYG{+w}{       }\PYG{c+c1}{\PYGZsh{} apply clang\PYGZhy{}format to all files in src and the package folders}
make\PYG{+w}{ }format\PYGZhy{}tests\PYG{+w}{     }\PYG{c+c1}{\PYGZsh{} apply clang\PYGZhy{}format to all files in the unittest tree}
\end{sphinxVerbatim}


\bigskip\hrule\bigskip



\subsection{GitHub command line interface}
\label{\detokenize{Build_development:github-command-line-interface}}\label{\detokenize{Build_development:gh-cli}}
\sphinxAtStartPar
GitHub is developing a \sphinxhref{https://cli.github.com}{tool for the command line} that interacts with the GitHub website via a
command called \sphinxcode{\sphinxupquote{gh}}.  This can be extremely convenient when working
with a Git repository hosted on GitHub (like LAMMPS).  It is thus highly
recommended to install it when doing LAMMPS development.

\sphinxAtStartPar
The capabilities of the \sphinxcode{\sphinxupquote{gh}} command is continually expanding, so
please see the documentation at \sphinxurl{https://cli.github.com/manual/}

\sphinxstepscope


\chapter{Run LAMMPS}
\label{\detokenize{Run_head:run-lammps}}\label{\detokenize{Run_head::doc}}
\sphinxAtStartPar
These pages explain how to run LAMMPS once you have {\hyperref[\detokenize{Install::doc}]{\sphinxcrossref{\DUrole{doc}{installed an executable}}}} or {\hyperref[\detokenize{Install::doc}]{\sphinxcrossref{\DUrole{doc}{downloaded the source code}}}}
and {\hyperref[\detokenize{Build::doc}]{\sphinxcrossref{\DUrole{doc}{built an executable}}}}.  The {\hyperref[\detokenize{Commands::doc}]{\sphinxcrossref{\DUrole{doc}{Commands}}}}
doc page describes how input scripts are structured and the commands
they can contain.

\sphinxstepscope


\section{Basics of running LAMMPS}
\label{\detokenize{Run_basics:basics-of-running-lammps}}\label{\detokenize{Run_basics::doc}}
\sphinxAtStartPar
LAMMPS is run from the command line, reading commands from a file via
the \sphinxcode{\sphinxupquote{\sphinxhyphen{}in}} command line flag, or from standard input.  Using the \sphinxcode{\sphinxupquote{\sphinxhyphen{}in
in.file}} variant is recommended (see note below).  The name of the
LAMMPS executable is either \sphinxcode{\sphinxupquote{lmp}} or \sphinxcode{\sphinxupquote{lmp\_\textless{}machine\textgreater{}}} with
\sphinxtitleref{\textless{}machine\textgreater{}} being the machine string used when compiling LAMMPS.  This
is required when compiling LAMMPS with the traditional build system
(e.g. with \sphinxcode{\sphinxupquote{make mpi}}), but optional when using CMake to configure and
build LAMMPS:

\begin{sphinxVerbatim}[commandchars=\\\{\}]
lmp\PYGZus{}serial\PYG{+w}{ }\PYGZhy{}in\PYG{+w}{ }\PYG{k}{in}.file
lmp\PYGZus{}serial\PYG{+w}{ }\PYGZlt{}\PYG{+w}{ }\PYG{k}{in}.file
lmp\PYG{+w}{ }\PYGZhy{}in\PYG{+w}{ }\PYG{k}{in}.file
lmp\PYG{+w}{ }\PYGZlt{}\PYG{+w}{ }\PYG{k}{in}.file
/path/to/lammps/src/lmp\PYGZus{}serial\PYG{+w}{ }\PYGZhy{}i\PYG{+w}{ }\PYG{k}{in}.file
mpirun\PYG{+w}{ }\PYGZhy{}np\PYG{+w}{ }\PYG{l+m}{4}\PYG{+w}{ }lmp\PYGZus{}mpi\PYG{+w}{ }\PYGZhy{}in\PYG{+w}{ }\PYG{k}{in}.file
mpiexec\PYG{+w}{ }\PYGZhy{}np\PYG{+w}{ }\PYG{l+m}{4}\PYG{+w}{ }lmp\PYG{+w}{ }\PYGZhy{}in\PYG{+w}{ }\PYG{k}{in}.file
mpirun\PYG{+w}{ }\PYGZhy{}np\PYG{+w}{ }\PYG{l+m}{8}\PYG{+w}{ }/path/to/lammps/src/lmp\PYGZus{}mpi\PYG{+w}{ }\PYGZhy{}in\PYG{+w}{ }\PYG{k}{in}.file
mpiexec\PYG{+w}{ }\PYGZhy{}n\PYG{+w}{ }\PYG{l+m}{6}\PYG{+w}{ }/usr/local/bin/lmp\PYG{+w}{ }\PYGZhy{}in\PYG{+w}{ }\PYG{k}{in}.file
\end{sphinxVerbatim}

\sphinxAtStartPar
You normally run the LAMMPS command in the directory where your input
script is located.  That is also where output files are produced by
default, unless you provide specific other paths in your input script or
on the command line.  As in some of the examples above, the LAMMPS
executable itself can be placed elsewhere.

\begin{sphinxadmonition}{note}{Note:}
\sphinxAtStartPar
The redirection operator “\textless{}” will not always work when running in
parallel with \sphinxcode{\sphinxupquote{mpirun}} or \sphinxcode{\sphinxupquote{mpiexec}}; for those systems the \sphinxhyphen{}in
form is required.
\end{sphinxadmonition}

\sphinxAtStartPar
As LAMMPS runs it prints info to the screen and a logfile named
\sphinxcode{\sphinxupquote{log.lammps}}.  More info about output is given on the {\hyperref[\detokenize{Run_output::doc}]{\sphinxcrossref{\DUrole{doc}{screen and
logfile output}}}} page.

\sphinxAtStartPar
If LAMMPS encounters errors in the input script or while running a
simulation it will print an ERROR message and stop or a WARNING
message and continue.  See the {\hyperref[\detokenize{Errors::doc}]{\sphinxcrossref{\DUrole{doc}{Common Problems}}}} page for a
discussion of the various kinds of errors LAMMPS can or can’t detect,
a list of all ERROR and WARNING messages, and what to do about them.


\bigskip\hrule\bigskip


\sphinxAtStartPar
LAMMPS can run the same problem on any number of processors, including a
single processor.  In theory you should get identical answers on any
number of processors and on any machine.  In practice, numerical
round\sphinxhyphen{}off due to using floating\sphinxhyphen{}point math can cause slight differences
and an eventual divergence of molecular dynamics trajectories.  See the
{\hyperref[\detokenize{Errors_common::doc}]{\sphinxcrossref{\DUrole{doc}{Errors common}}}} page for discussion of this.

\sphinxAtStartPar
LAMMPS can run as large a problem as will fit in the physical memory of
one or more processors.  If you run out of memory, you must run on more
processors or define a smaller problem.  The amount of memory needed and
how well it can be distributed across processors may vary based
on the models and settings and commands used.

\sphinxAtStartPar
If you run LAMMPS in parallel via mpirun, you should be aware of the
{\hyperref[\detokenize{processors::doc}]{\sphinxcrossref{\DUrole{doc}{processors}}}} command, which controls how MPI tasks are
mapped to the simulation box, as well as mpirun options that control how
MPI tasks are assigned to physical cores of the node(s) of the machine
you are running on.  These settings can improve performance, though the
defaults are often adequate.

\sphinxAtStartPar
For example, it is often important to bind MPI tasks (processes) to
physical cores (processor affinity), so that the operating system does
not migrate them during a simulation.  If this is not the default
behavior on your machine, the mpirun option \sphinxcode{\sphinxupquote{\sphinxhyphen{}\sphinxhyphen{}bind\sphinxhyphen{}to core}} (OpenMPI)
or \sphinxcode{\sphinxupquote{\sphinxhyphen{}bind\sphinxhyphen{}to core}} (MPICH) can be used.

\sphinxAtStartPar
If the LAMMPS command(s) you are using support multi\sphinxhyphen{}threading, you
can set the number of threads per MPI task via the environment
variable \sphinxcode{\sphinxupquote{OMP\_NUM\_THREADS}}, before you launch LAMMPS:

\begin{sphinxVerbatim}[commandchars=\\\{\}]
\PYG{n+nb}{export}\PYG{+w}{ }\PYG{n+nv}{OMP\PYGZus{}NUM\PYGZus{}THREADS}\PYG{o}{=}\PYG{l+m}{2}\PYG{+w}{     }\PYG{c+c1}{\PYGZsh{} bash}
setenv\PYG{+w}{ }OMP\PYGZus{}NUM\PYGZus{}THREADS\PYG{+w}{ }\PYG{l+m}{2}\PYG{+w}{     }\PYG{c+c1}{\PYGZsh{} csh or tcsh}
\end{sphinxVerbatim}

\sphinxAtStartPar
This can also be done via the {\hyperref[\detokenize{package::doc}]{\sphinxcrossref{\DUrole{doc}{package}}}} command or via
the {\hyperref[\detokenize{Run_options::doc}]{\sphinxcrossref{\DUrole{doc}{\sphinxhyphen{}pk command\sphinxhyphen{}line switch}}}} which invokes the
package command.  See the {\hyperref[\detokenize{package::doc}]{\sphinxcrossref{\DUrole{doc}{package}}}} command or
{\hyperref[\detokenize{Speed::doc}]{\sphinxcrossref{\DUrole{doc}{Speed}}}} doc pages for more details about which accelerator
packages and which commands support multi\sphinxhyphen{}threading.


\bigskip\hrule\bigskip


\sphinxAtStartPar
You can experiment with running LAMMPS using any of the input scripts
provided in the examples or bench directory.  Input scripts are named
\sphinxcode{\sphinxupquote{in.*}} and sample outputs are named \sphinxcode{\sphinxupquote{log.*.P}} where P is the number of
processors it was run on.

\sphinxAtStartPar
Some of the examples or benchmarks require LAMMPS to be built with
optional packages.

\sphinxstepscope


\section{Command\sphinxhyphen{}line options}
\label{\detokenize{Run_options:command-line-options}}\label{\detokenize{Run_options::doc}}
\sphinxAtStartPar
At run time, LAMMPS recognizes several optional command\sphinxhyphen{}line switches
which may be used in any order.  Either the full word or a one or two
letter abbreviation can be used:
\begin{itemize}
\item {} 
\sphinxAtStartPar
{\hyperref[\detokenize{Run_options:echo}]{\sphinxcrossref{\DUrole{std,std-ref}{\sphinxhyphen{}e or \sphinxhyphen{}echo}}}}

\item {} 
\sphinxAtStartPar
{\hyperref[\detokenize{Run_options:help}]{\sphinxcrossref{\DUrole{std,std-ref}{\sphinxhyphen{}h or \sphinxhyphen{}help}}}}

\item {} 
\sphinxAtStartPar
{\hyperref[\detokenize{Run_options:file}]{\sphinxcrossref{\DUrole{std,std-ref}{\sphinxhyphen{}i or \sphinxhyphen{}in}}}}

\item {} 
\sphinxAtStartPar
{\hyperref[\detokenize{Run_options:run-kokkos}]{\sphinxcrossref{\DUrole{std,std-ref}{\sphinxhyphen{}k or \sphinxhyphen{}kokkos}}}}

\item {} 
\sphinxAtStartPar
{\hyperref[\detokenize{Run_options:log}]{\sphinxcrossref{\DUrole{std,std-ref}{\sphinxhyphen{}l or \sphinxhyphen{}log}}}}

\item {} 
\sphinxAtStartPar
{\hyperref[\detokenize{Run_options:mdi-flags}]{\sphinxcrossref{\DUrole{std,std-ref}{\sphinxhyphen{}mdi}}}}

\item {} 
\sphinxAtStartPar
{\hyperref[\detokenize{Run_options:mpicolor}]{\sphinxcrossref{\DUrole{std,std-ref}{\sphinxhyphen{}m or \sphinxhyphen{}mpicolor}}}}

\item {} 
\sphinxAtStartPar
{\hyperref[\detokenize{Run_options:cite}]{\sphinxcrossref{\DUrole{std,std-ref}{\sphinxhyphen{}c or \sphinxhyphen{}cite}}}}

\item {} 
\sphinxAtStartPar
{\hyperref[\detokenize{Run_options:nocite}]{\sphinxcrossref{\DUrole{std,std-ref}{\sphinxhyphen{}nc or \sphinxhyphen{}nocite}}}}

\item {} 
\sphinxAtStartPar
{\hyperref[\detokenize{Run_options:nonbuf}]{\sphinxcrossref{\DUrole{std,std-ref}{\sphinxhyphen{}nb or \sphinxhyphen{}nonbuf}}}}

\item {} 
\sphinxAtStartPar
{\hyperref[\detokenize{Run_options:package}]{\sphinxcrossref{\DUrole{std,std-ref}{\sphinxhyphen{}pk or \sphinxhyphen{}package}}}}

\item {} 
\sphinxAtStartPar
{\hyperref[\detokenize{Run_options:partition}]{\sphinxcrossref{\DUrole{std,std-ref}{\sphinxhyphen{}p or \sphinxhyphen{}partition}}}}

\item {} 
\sphinxAtStartPar
{\hyperref[\detokenize{Run_options:plog}]{\sphinxcrossref{\DUrole{std,std-ref}{\sphinxhyphen{}pl or \sphinxhyphen{}plog}}}}

\item {} 
\sphinxAtStartPar
{\hyperref[\detokenize{Run_options:pscreen}]{\sphinxcrossref{\DUrole{std,std-ref}{\sphinxhyphen{}ps or \sphinxhyphen{}pscreen}}}}

\item {} 
\sphinxAtStartPar
{\hyperref[\detokenize{Run_options:reorder}]{\sphinxcrossref{\DUrole{std,std-ref}{\sphinxhyphen{}ro or \sphinxhyphen{}reorder}}}}

\item {} 
\sphinxAtStartPar
{\hyperref[\detokenize{Run_options:restart2data}]{\sphinxcrossref{\DUrole{std,std-ref}{\sphinxhyphen{}r2data or \sphinxhyphen{}restart2data}}}}

\item {} 
\sphinxAtStartPar
{\hyperref[\detokenize{Run_options:restart2dump}]{\sphinxcrossref{\DUrole{std,std-ref}{\sphinxhyphen{}r2dump or \sphinxhyphen{}restart2dump}}}}

\item {} 
\sphinxAtStartPar
{\hyperref[\detokenize{Run_options:restart2info}]{\sphinxcrossref{\DUrole{std,std-ref}{\sphinxhyphen{}r2info or \sphinxhyphen{}restart2info}}}}

\item {} 
\sphinxAtStartPar
{\hyperref[\detokenize{Run_options:screen}]{\sphinxcrossref{\DUrole{std,std-ref}{\sphinxhyphen{}sc or \sphinxhyphen{}screen}}}}

\item {} 
\sphinxAtStartPar
{\hyperref[\detokenize{Run_options:skiprun}]{\sphinxcrossref{\DUrole{std,std-ref}{\sphinxhyphen{}sr or skiprun}}}}

\item {} 
\sphinxAtStartPar
{\hyperref[\detokenize{Run_options:suffix}]{\sphinxcrossref{\DUrole{std,std-ref}{\sphinxhyphen{}sf or \sphinxhyphen{}suffix}}}}

\item {} 
\sphinxAtStartPar
{\hyperref[\detokenize{Run_options:var}]{\sphinxcrossref{\DUrole{std,std-ref}{\sphinxhyphen{}v or \sphinxhyphen{}var}}}}

\end{itemize}

\sphinxAtStartPar
For example, the lmp\_mpi executable might be launched as follows:

\begin{sphinxVerbatim}[commandchars=\\\{\}]
mpirun\PYG{+w}{ }\PYGZhy{}np\PYG{+w}{ }\PYG{l+m}{16}\PYG{+w}{ }lmp\PYGZus{}mpi\PYG{+w}{ }\PYGZhy{}v\PYG{+w}{ }f\PYG{+w}{ }tmp.out\PYG{+w}{ }\PYGZhy{}l\PYG{+w}{ }my.log\PYG{+w}{ }\PYGZhy{}sc\PYG{+w}{ }none\PYG{+w}{ }\PYGZhy{}i\PYG{+w}{ }\PYG{k}{in}.alloy
mpirun\PYG{+w}{ }\PYGZhy{}np\PYG{+w}{ }\PYG{l+m}{16}\PYG{+w}{ }lmp\PYGZus{}mpi\PYG{+w}{ }\PYGZhy{}var\PYG{+w}{ }f\PYG{+w}{ }tmp.out\PYG{+w}{ }\PYGZhy{}log\PYG{+w}{ }my.log\PYG{+w}{ }\PYGZhy{}screen\PYG{+w}{ }none\PYG{+w}{ }\PYGZhy{}in\PYG{+w}{ }\PYG{k}{in}.alloy
\end{sphinxVerbatim}


\bigskip\hrule\bigskip

\phantomsection\label{\detokenize{Run_options:echo}}
\sphinxAtStartPar
\sphinxstylestrong{\sphinxhyphen{}echo style}

\sphinxAtStartPar
Set the style of command echoing.  The style can be \sphinxstyleemphasis{none} or \sphinxstyleemphasis{screen}
or \sphinxstyleemphasis{log} or \sphinxstyleemphasis{both}.  Depending on the style, each command read from
the input script will be echoed to the screen and/or logfile.  This
can be useful to figure out which line of your script is causing an
input error.  The default value is \sphinxstyleemphasis{log}.  The echo style can also be
set by using the {\hyperref[\detokenize{echo::doc}]{\sphinxcrossref{\DUrole{doc}{echo}}}} command in the input script itself.


\bigskip\hrule\bigskip

\phantomsection\label{\detokenize{Run_options:help}}
\sphinxAtStartPar
\sphinxstylestrong{\sphinxhyphen{}help}

\sphinxAtStartPar
Print a brief help summary and a list of options compiled into this
executable for each LAMMPS style (atom\_style, fix, compute,
pair\_style, bond\_style, etc).  This can tell you if the command you
want to use was included via the appropriate package at compile time.
LAMMPS will print the info and immediately exit if this switch is
used.


\bigskip\hrule\bigskip

\phantomsection\label{\detokenize{Run_options:file}}
\sphinxAtStartPar
\sphinxstylestrong{\sphinxhyphen{}in file}

\sphinxAtStartPar
Specify a file to use as an input script.  This is an optional but
recommended switch when running LAMMPS in one\sphinxhyphen{}partition mode.  If it
is not specified, LAMMPS reads its script from standard input, typically
from a script via I/O redirection; e.g. lmp\_linux \textless{} in.run.
With many MPI implementations I/O redirection also works in parallel,
but using the \sphinxhyphen{}in flag will always work.

\sphinxAtStartPar
Note that this is a required switch when running LAMMPS in
multi\sphinxhyphen{}partition mode, since multiple processors cannot all read from
stdin concurrently.  The file name may be “none” for starting
multi\sphinxhyphen{}partition calculations without reading an initial input file
from the library interface.


\bigskip\hrule\bigskip

\phantomsection\label{\detokenize{Run_options:run-kokkos}}
\sphinxAtStartPar
\sphinxstylestrong{\sphinxhyphen{}kokkos on/off keyword/value …}

\sphinxAtStartPar
Explicitly enable or disable KOKKOS support, as provided by the KOKKOS
package.  Even if LAMMPS is built with this package, as described
in the {\hyperref[\detokenize{Speed_kokkos::doc}]{\sphinxcrossref{\DUrole{doc}{the KOKKOS package page}}}}, this switch must be set to enable
running with KOKKOS\sphinxhyphen{}enabled styles the package provides.  If the
switch is not set (the default), LAMMPS will operate as if the KOKKOS
package were not installed; i.e. you can run standard LAMMPS or with
the GPU or OPENMP packages, for testing or benchmarking purposes.

\sphinxAtStartPar
Additional optional keyword/value pairs can be specified which determine
how Kokkos will use the underlying hardware on your platform.  These
settings apply to each MPI task you launch via the \sphinxcode{\sphinxupquote{mpirun}} or
\sphinxcode{\sphinxupquote{mpiexec}} command.  You may choose to run one or more MPI tasks per
physical node.  Note that if you are running on a desktop machine, you
typically have one physical node.  On a cluster or supercomputer there
may be dozens or 1000s of physical nodes.

\sphinxAtStartPar
Either the full word or an abbreviation can be used for the keywords.
Note that the keywords do not use a leading minus sign.  I.e. the
keyword is “t”, not “\sphinxhyphen{}t”.  Also note that each of the keywords has a
default setting.  Examples of when to use these options and what
settings to use on different platforms is given on the {\hyperref[\detokenize{Speed_kokkos::doc}]{\sphinxcrossref{\DUrole{doc}{KOKKOS
package}}}} doc page.
\begin{itemize}
\item {} 
\sphinxAtStartPar
d or device

\item {} 
\sphinxAtStartPar
g or gpus

\item {} 
\sphinxAtStartPar
t or threads

\end{itemize}

\begin{sphinxVerbatim}[commandchars=\\\{\}]
device Nd
\end{sphinxVerbatim}

\sphinxAtStartPar
This option is only relevant if you built LAMMPS with CUDA=yes, you
have more than one GPU per node, and if you are running with only one
MPI task per node.  The Nd setting is the ID of the GPU on the node to
run on.  By default Nd = 0.  If you have multiple GPUs per node, they
have consecutive IDs numbered as 0,1,2,etc.  This setting allows you
to launch multiple independent jobs on the node, each with a single
MPI task per node, and assign each job to run on a different GPU.

\begin{sphinxVerbatim}[commandchars=\\\{\}]
gpus Ng Ns
\end{sphinxVerbatim}

\sphinxAtStartPar
This option is only relevant if you built LAMMPS with CUDA=yes, you
have more than one GPU per node, and you are running with multiple MPI
tasks per node (up to one per GPU).  The Ng setting is how many GPUs
you will use.  The Ns setting is optional.  If set, it is the ID of a
GPU to skip when assigning MPI tasks to GPUs.  This may be useful if
your desktop system reserves one GPU to drive the screen and the rest
are intended for computational work like running LAMMPS.  By default
Ng = 1 and Ns is not set.

\sphinxAtStartPar
Depending on which flavor of MPI you are running, LAMMPS will look for
one of these 4 environment variables

\begin{sphinxVerbatim}[commandchars=\\\{\}]
SLURM\PYGZus{}LOCALID (various MPI variants compiled with SLURM support)
MPT\PYGZus{}LRANK (HPE MPI)
MV2\PYGZus{}COMM\PYGZus{}WORLD\PYGZus{}LOCAL\PYGZus{}RANK (Mvapich)
OMPI\PYGZus{}COMM\PYGZus{}WORLD\PYGZus{}LOCAL\PYGZus{}RANK (OpenMPI)
\end{sphinxVerbatim}

\sphinxAtStartPar
which are initialized by the \sphinxcode{\sphinxupquote{srun}}, \sphinxcode{\sphinxupquote{mpirun}}, or \sphinxcode{\sphinxupquote{mpiexec}}
commands.  The environment variable setting for each MPI rank is used to
assign a unique GPU ID to the MPI task.

\begin{sphinxVerbatim}[commandchars=\\\{\}]
threads Nt
\end{sphinxVerbatim}

\sphinxAtStartPar
This option assigns Nt number of threads to each MPI task for
performing work when Kokkos is executing in OpenMP or pthreads mode.
The default is Nt = 1, which essentially runs in MPI\sphinxhyphen{}only mode.  If
there are Np MPI tasks per physical node, you generally want Np*Nt =
the number of physical cores per node, to use your available hardware
optimally.  This also sets the number of threads used by the host when
LAMMPS is compiled with CUDA=yes.

\sphinxAtStartPar
\DUrole{versionmodified,deprecated}{Deprecated since version 22Dec2022.}

\sphinxAtStartPar
Support for the “numa” or “n” option was removed as its functionality
was ignored in Kokkos for some time already.


\bigskip\hrule\bigskip

\phantomsection\label{\detokenize{Run_options:log}}
\sphinxAtStartPar
\sphinxstylestrong{\sphinxhyphen{}log file}

\sphinxAtStartPar
Specify a log file for LAMMPS to write status information to.  In
one\sphinxhyphen{}partition mode, if the switch is not used, LAMMPS writes to the
file log.lammps.  If this switch is used, LAMMPS writes to the
specified file.  In multi\sphinxhyphen{}partition mode, if the switch is not used, a
log.lammps file is created with high\sphinxhyphen{}level status information.  Each
partition also writes to a log.lammps.N file where N is the partition
ID.  If the switch is specified in multi\sphinxhyphen{}partition mode, the high\sphinxhyphen{}level
logfile is named “file” and each partition also logs information to a
file.N.  For both one\sphinxhyphen{}partition and multi\sphinxhyphen{}partition mode, if the
specified file is “none”, then no log files are created.  Using a
{\hyperref[\detokenize{log::doc}]{\sphinxcrossref{\DUrole{doc}{log}}}} command in the input script will override this setting.
Option \sphinxhyphen{}plog will override the name of the partition log files file.N.


\bigskip\hrule\bigskip

\phantomsection\label{\detokenize{Run_options:mdi-flags}}
\sphinxAtStartPar
\sphinxstylestrong{\sphinxhyphen{}mdi ‘multiple flags’}

\sphinxAtStartPar
This flag is only recognized and used when LAMMPS has support for the MolSSI
Driver Interface (MDI) included as part of the {\hyperref[\detokenize{Packages_details:pkg-mdi}]{\sphinxcrossref{\DUrole{std,std-ref}{MDI}}}}
package.  This flag is specific to the MDI library and controls how LAMMPS
interacts with MDI.  There are usually multiple flags that have to follow it
and those have to be placed in quotation marks.  For more information about
how to launch LAMMPS in MDI client/server mode please refer to the
{\hyperref[\detokenize{Howto_mdi::doc}]{\sphinxcrossref{\DUrole{doc}{MDI Howto}}}}.


\bigskip\hrule\bigskip

\phantomsection\label{\detokenize{Run_options:mpicolor}}
\sphinxAtStartPar
\sphinxstylestrong{\sphinxhyphen{}mpicolor color}

\sphinxAtStartPar
If used, this must be the first command\sphinxhyphen{}line argument after the LAMMPS
executable name.  It is only used when LAMMPS is launched by an mpirun
command which also launches another executable(s) at the same time.
(The other executable could be LAMMPS as well.)  The color is an
integer value which should be different for each executable (another
application may set this value in a different way).  LAMMPS and the
other executable(s) perform an MPI\_Comm\_split() with their own colors
to shrink the MPI\_COMM\_WORLD communication to be the subset of
processors they are actually running on.


\bigskip\hrule\bigskip

\phantomsection\label{\detokenize{Run_options:cite}}
\sphinxAtStartPar
\sphinxstylestrong{\sphinxhyphen{}cite style} or \sphinxstylestrong{file name}

\sphinxAtStartPar
Select how and where to output a reminder about citing contributions
to the LAMMPS code that were used during the run. Available keywords
for styles are “both”, “none”, “screen”, or “log”.  Any other keyword
will be considered a file name to write the detailed citation info to
instead of logfile or screen.  Default is the “log” style where there
is a short summary in the screen output and detailed citations
in BibTeX format in the logfile.  The option “both” selects the detailed
output for both, “none”, the short output for both, and “screen” will
write the detailed info to the screen and the short version to the log
file.  If a dedicated citation info file is requested, the screen and
log file output will be in the short format (same as with “none”).

\sphinxAtStartPar
See the {\hyperref[\detokenize{Intro_citing::doc}]{\sphinxcrossref{\DUrole{doc}{citation page}}}} for more details on
how to correctly reference and cite LAMMPS.


\bigskip\hrule\bigskip

\phantomsection\label{\detokenize{Run_options:nocite}}
\sphinxAtStartPar
\sphinxstylestrong{\sphinxhyphen{}nocite}

\sphinxAtStartPar
Disable generating a citation reminder (see above) at all.


\bigskip\hrule\bigskip

\phantomsection\label{\detokenize{Run_options:nonbuf}}
\sphinxAtStartPar
\sphinxstylestrong{\sphinxhyphen{}nonbuf}

\sphinxAtStartPar
\DUrole{versionmodified,added}{Added in version 15Sep2022.}

\sphinxAtStartPar
Turn off buffering for screen and logfile output.  For performance
reasons, output to the screen and logfile is usually buffered, i.e.
output is only written to a file if its buffer \sphinxhyphen{} typically 4096 bytes \sphinxhyphen{}
has been filled.  When LAMMPS crashes for some reason, however, that can
mean that there is important output missing.  With this flag the
buffering can be turned off (only for screen and logfile output) and any
output will be committed immediately.  Note that when running in
parallel with MPI, the screen output may still be buffered by the MPI
library and this cannot be changed by LAMMPS.  This flag should only be
used for debugging and not for production simulations as the performance
impact can be significant, especially for large parallel runs.


\bigskip\hrule\bigskip

\phantomsection\label{\detokenize{Run_options:package}}
\sphinxAtStartPar
\sphinxstylestrong{\sphinxhyphen{}package style args ….}

\sphinxAtStartPar
Invoke the {\hyperref[\detokenize{package::doc}]{\sphinxcrossref{\DUrole{doc}{package}}}} command with style and args.  The
syntax is the same as if the command appeared at the top of the input
script.  For example \sphinxcode{\sphinxupquote{\sphinxhyphen{}package gpu 2}} or \sphinxcode{\sphinxupquote{\sphinxhyphen{}pk gpu 2}} is the same as
{\hyperref[\detokenize{package::doc}]{\sphinxcrossref{\DUrole{doc}{package gpu 2}}}} in the input script.  The possible styles
and args are documented on the {\hyperref[\detokenize{package::doc}]{\sphinxcrossref{\DUrole{doc}{package}}}} doc page.  This
switch can be used multiple times, e.g. to set options for the
INTEL and OPENMP packages which can be used together.

\sphinxAtStartPar
Along with the \sphinxcode{\sphinxupquote{\sphinxhyphen{}suffix}} command\sphinxhyphen{}line switch, this is a convenient
mechanism for invoking accelerator packages and their options without
having to edit an input script.


\bigskip\hrule\bigskip

\phantomsection\label{\detokenize{Run_options:partition}}
\sphinxAtStartPar
\sphinxstylestrong{\sphinxhyphen{}partition 8x2 4 5 …}

\sphinxAtStartPar
Invoke LAMMPS in multi\sphinxhyphen{}partition mode.  When LAMMPS is run on P
processors and this switch is not used, LAMMPS runs in one partition,
i.e. all P processors run a single simulation.  If this switch is
used, the P processors are split into separate partitions and each
partition runs its own simulation.  The arguments to the switch
specify the number of processors in each partition.  Arguments of the
form MxN mean M partitions, each with N processors.  Arguments of the
form N mean a single partition with N processors.  The sum of
processors in all partitions must equal P.  Thus the command
\sphinxcode{\sphinxupquote{\sphinxhyphen{}partition 8x2 4 5}} has 10 partitions and runs on a total of 25
processors.

\sphinxAtStartPar
Running with multiple partitions can be useful for running
{\hyperref[\detokenize{Howto_replica::doc}]{\sphinxcrossref{\DUrole{doc}{multi\sphinxhyphen{}replica simulations}}}}, where each replica
runs on one or a few processors.  Note that with MPI installed on a
machine (e.g. your desktop), you can run on more (virtual) processors
than you have physical processors.

\sphinxAtStartPar
To run multiple independent simulations from one input script, using
multiple partitions, see the {\hyperref[\detokenize{Howto_multiple::doc}]{\sphinxcrossref{\DUrole{doc}{Howto multiple}}}}
page.  World\sphinxhyphen{} and universe\sphinxhyphen{}style {\hyperref[\detokenize{variable::doc}]{\sphinxcrossref{\DUrole{doc}{variables}}}} are useful
in this context.


\bigskip\hrule\bigskip

\phantomsection\label{\detokenize{Run_options:plog}}
\sphinxAtStartPar
\sphinxstylestrong{\sphinxhyphen{}plog file}

\sphinxAtStartPar
Specify the base name for the partition log files, so partition N
writes log information to file.N. If file is none, then no partition
log files are created.  This overrides the filename specified in the
\sphinxhyphen{}log command\sphinxhyphen{}line option.  This option is useful when working with
large numbers of partitions, allowing the partition log files to be
suppressed (\sphinxhyphen{}plog none) or placed in a subdirectory (\sphinxhyphen{}plog
replica\_files/log.lammps) If this option is not used the log file for
partition N is log.lammps.N or whatever is specified by the \sphinxhyphen{}log
command\sphinxhyphen{}line option.


\bigskip\hrule\bigskip

\phantomsection\label{\detokenize{Run_options:pscreen}}
\sphinxAtStartPar
\sphinxstylestrong{\sphinxhyphen{}pscreen file}

\sphinxAtStartPar
Specify the base name for the partition screen file, so partition N
writes screen information to file.N. If file is “none”, then no
partition screen files are created.  This overrides the filename
specified in the \sphinxhyphen{}screen command\sphinxhyphen{}line option.  This option is useful
when working with large numbers of partitions, allowing the partition
screen files to be suppressed (\sphinxhyphen{}pscreen none) or placed in a
subdirectory (\sphinxhyphen{}pscreen replica\_files/screen).  If this option is not
used the screen file for partition N is screen.N or whatever is
specified by the \sphinxhyphen{}screen command\sphinxhyphen{}line option.


\bigskip\hrule\bigskip

\phantomsection\label{\detokenize{Run_options:reorder}}
\sphinxAtStartPar
\sphinxstylestrong{\sphinxhyphen{}reorder}

\sphinxAtStartPar
This option has 2 forms:

\begin{sphinxVerbatim}[commandchars=\\\{\}]
\PYGZhy{}reorder nth N
\PYGZhy{}reorder custom filename
\end{sphinxVerbatim}

\sphinxAtStartPar
Reorder the processors in the MPI communicator used to instantiate
LAMMPS, in one of several ways.  The original MPI communicator ranks
all P processors from 0 to P\sphinxhyphen{}1.  The mapping of these ranks to
physical processors is done by MPI before LAMMPS begins.  It may be
useful in some cases to alter the rank order.  E.g. to ensure that
cores within each node are ranked in a desired order.  Or when using
the {\hyperref[\detokenize{run_style::doc}]{\sphinxcrossref{\DUrole{doc}{run\_style verlet/split}}}} command with 2 partitions
to ensure that a specific Kspace processor (in the second partition) is
matched up with a specific set of processors in the first partition.
See the {\hyperref[\detokenize{Speed_tips::doc}]{\sphinxcrossref{\DUrole{doc}{General tips}}}} page for more details.

\sphinxAtStartPar
If the keyword \sphinxstyleemphasis{nth} is used with a setting \sphinxstyleemphasis{N}, then it means every
Nth processor will be moved to the end of the ranking.  This is useful
when using the {\hyperref[\detokenize{run_style::doc}]{\sphinxcrossref{\DUrole{doc}{run\_style verlet/split}}}} command with 2
partitions via the \sphinxhyphen{}partition command\sphinxhyphen{}line switch.  The first set of
processors will be in the first partition, the second set in the second
partition.  The \sphinxhyphen{}reorder command\sphinxhyphen{}line switch can alter this so that
the first N procs in the first partition and one proc in the second partition
will be ordered consecutively, e.g. as the cores on one physical node.
This can boost performance.  For example, if you use \sphinxcode{\sphinxupquote{\sphinxhyphen{}reorder nth 4}}
and \sphinxcode{\sphinxupquote{\sphinxhyphen{}partition 9 3}} and you are running on 12 processors, the
processors will be reordered from

\begin{sphinxVerbatim}[commandchars=\\\{\}]
0 1 2 3 4 5 6 7 8 9 10 11
\end{sphinxVerbatim}

\sphinxAtStartPar
to

\begin{sphinxVerbatim}[commandchars=\\\{\}]
0 1 2 4 5 6 8 9 10 3 7 11
\end{sphinxVerbatim}

\sphinxAtStartPar
so that the processors in each partition will be

\begin{sphinxVerbatim}[commandchars=\\\{\}]
0 1 2 4 5 6 8 9 10
3 7 11
\end{sphinxVerbatim}

\sphinxAtStartPar
See the “processors” command for how to ensure processors from each
partition could then be grouped optimally for quad\sphinxhyphen{}core nodes.

\sphinxAtStartPar
If the keyword is \sphinxstyleemphasis{custom}, then a file that specifies a permutation
of the processor ranks is also specified.  The format of the reorder
file is as follows.  Any number of initial blank or comment lines
(starting with a “\#” character) can be present.  These should be
followed by P lines of the form:

\begin{sphinxVerbatim}[commandchars=\\\{\}]
I J
\end{sphinxVerbatim}

\sphinxAtStartPar
where P is the number of processors LAMMPS was launched with.  Note
that if running in multi\sphinxhyphen{}partition mode (see the \sphinxhyphen{}partition switch
above) P is the total number of processors in all partitions.  The I
and J values describe a permutation of the P processors.  Every I and
J should be values from 0 to P\sphinxhyphen{}1 inclusive.  In the set of P I values,
every proc ID should appear exactly once.  Ditto for the set of P J
values.  A single I,J pairing means that the physical processor with
rank I in the original MPI communicator will have rank J in the
reordered communicator.

\sphinxAtStartPar
Note that rank ordering can also be specified by many MPI
implementations, either by environment variables that specify how to
order physical processors, or by config files that specify what
physical processors to assign to each MPI rank.  The \sphinxhyphen{}reorder switch
simply gives you a portable way to do this without relying on MPI
itself.  See the {\hyperref[\detokenize{processors::doc}]{\sphinxcrossref{\DUrole{doc}{processors file}}}} command for how
to output info on the final assignment of physical processors to
the LAMMPS simulation domain.


\bigskip\hrule\bigskip

\phantomsection\label{\detokenize{Run_options:restart2data}}
\sphinxAtStartPar
\sphinxstylestrong{\sphinxhyphen{}restart2data restartfile datafile keyword value …}

\sphinxAtStartPar
Convert the restart file into a data file and immediately exit.  This
is the same operation as if the following 2\sphinxhyphen{}line input script were
run:

\begin{sphinxVerbatim}[commandchars=\\\{\}]
\PYG{k}{read\PYGZus{}restart}\PYG{+w}{ }\PYG{n}{restartfile}
\PYG{k}{write\PYGZus{}data }\PYG{n+nv+nvIdentifier}{datafile}\PYG{+w}{ }\PYG{n}{keyword}\PYG{+w}{ }\PYG{n}{value}\PYG{+w}{ }\PYG{n}{...}
\end{sphinxVerbatim}

\sphinxAtStartPar
The specified restartfile and/or datafile name may contain the wild\sphinxhyphen{}card
character “*”.  The restartfile name may also contain the wild\sphinxhyphen{}card
character “\%”.  The meaning of these characters is explained on the
{\hyperref[\detokenize{read_restart::doc}]{\sphinxcrossref{\DUrole{doc}{read\_restart}}}} and {\hyperref[\detokenize{write_data::doc}]{\sphinxcrossref{\DUrole{doc}{write\_data}}}} doc
pages.  The use of “\%” means that a parallel restart file can be read.
Note that a filename such as file.* may need to be enclosed in quotes or
the “*” character prefixed with a backslash (”") to avoid shell
expansion of the “*” character.

\sphinxAtStartPar
The syntax following restartfile, namely

\begin{sphinxVerbatim}[commandchars=\\\{\}]
datafile keyword value ...
\end{sphinxVerbatim}

\sphinxAtStartPar
is identical to the arguments of the {\hyperref[\detokenize{write_data::doc}]{\sphinxcrossref{\DUrole{doc}{write\_data}}}}
command.  See its documentation page for details.  This includes its
optional keyword/value settings.


\bigskip\hrule\bigskip

\phantomsection\label{\detokenize{Run_options:restart2dump}}
\sphinxAtStartPar
\sphinxstylestrong{\sphinxhyphen{}restart2dump restartfile group\sphinxhyphen{}ID dumpstyle dumpfile arg1 arg2 …}

\sphinxAtStartPar
Convert the restart file into a dump file and immediately exit.  This
is the same operation as if the following 2\sphinxhyphen{}line input script were
run:

\begin{sphinxVerbatim}[commandchars=\\\{\}]
\PYG{k}{read\PYGZus{}restart}\PYG{+w}{ }\PYG{n}{restartfile}
\PYG{k}{write\PYGZus{}dump }\PYG{n+nv+nvIdentifier}{group\PYGZhy{}ID}\PYG{+w}{ }\PYG{n}{dumpstyle}\PYG{+w}{ }\PYG{n}{dumpfile}\PYG{+w}{ }\PYG{n}{arg1}\PYG{+w}{ }\PYG{n}{arg2}\PYG{+w}{ }\PYG{n}{...}
\end{sphinxVerbatim}

\sphinxAtStartPar
Note that the specified restartfile and dumpfile names may contain
wild\sphinxhyphen{}card characters (”*” or “\%”) as explained on the
{\hyperref[\detokenize{read_restart::doc}]{\sphinxcrossref{\DUrole{doc}{read\_restart}}}} and {\hyperref[\detokenize{write_dump::doc}]{\sphinxcrossref{\DUrole{doc}{write\_dump}}}} doc
pages.  The use of “\%” means that a parallel restart file and/or
parallel dump file can be read and/or written.  Note that a filename
such as file.* may need to be enclosed in quotes or the “*” character
prefixed with a backslash (”") to avoid shell expansion of the “*”
character.

\sphinxAtStartPar
The syntax following restartfile, namely

\begin{sphinxVerbatim}[commandchars=\\\{\}]
\PYG{n}{group}\PYG{o}{\PYGZhy{}}\PYG{n}{ID}\PYG{+w}{ }\PYG{n}{dumpstyle}\PYG{+w}{ }\PYG{n}{dumpfile}\PYG{+w}{ }\PYG{n}{arg1}\PYG{+w}{ }\PYG{n}{arg2}\PYG{+w}{ }\PYG{n}{...}
\end{sphinxVerbatim}

\sphinxAtStartPar
is identical to the arguments of the {\hyperref[\detokenize{write_dump::doc}]{\sphinxcrossref{\DUrole{doc}{write\_dump}}}}
command.  See its documentation page for details.  This includes what
per\sphinxhyphen{}atom fields are written to the dump file and optional dump\_modify
settings, including ones that affect how parallel dump files are written,
e.g. the \sphinxstyleemphasis{nfile} and \sphinxstyleemphasis{fileper} keywords.  See the
{\hyperref[\detokenize{dump_modify::doc}]{\sphinxcrossref{\DUrole{doc}{dump\_modify}}}} page for details.


\bigskip\hrule\bigskip

\phantomsection\label{\detokenize{Run_options:restart2info}}
\sphinxAtStartPar
\sphinxstylestrong{\sphinxhyphen{}restart2info restartfile keyword …}

\sphinxAtStartPar
\DUrole{versionmodified,added}{Added in version 29Aug2024.}

\sphinxAtStartPar
Write out some info about the restart file and and immediately exit.
This is the same operation as if the following 2\sphinxhyphen{}line input script were
run:

\begin{sphinxVerbatim}[commandchars=\\\{\}]
\PYG{k}{read\PYGZus{}restart}\PYG{+w}{ }\PYG{n}{restartfile}
\PYG{k}{info}\PYG{+w}{ }\PYG{n}{system}\PYG{+w}{ }\PYG{k}{group }\PYG{n+nv+nvIdentifier}{computes}\PYG{+w}{ }\PYG{n}{fixes}
\end{sphinxVerbatim}

\sphinxAtStartPar
The specified restartfile name may contain the wild\sphinxhyphen{}card character “*”.
The restartfile name may also contain the wild\sphinxhyphen{}card character “\%”.  The
meaning of these characters is explained on the {\hyperref[\detokenize{read_restart::doc}]{\sphinxcrossref{\DUrole{doc}{read\_restart}}}} documentation.  The use of “\%” means that a parallel
restart file can be read.  Note that a filename such as file.* may need
to be enclosed in quotes or the “*” character prefixed with a backslash
(”") to avoid shell expansion of the “*” character.

\sphinxAtStartPar
Optional keywords may follow the restartfile argument.  These must be
valid keywords for the {\hyperref[\detokenize{info::doc}]{\sphinxcrossref{\DUrole{doc}{info command}}}}.  The most useful
ones \sphinxhyphen{} \sphinxstyleemphasis{system}, \sphinxstyleemphasis{group}, \sphinxstyleemphasis{computes}, and \sphinxstyleemphasis{fixes} \sphinxhyphen{} are already applied.
Appending keywords like \sphinxstyleemphasis{coeffs} or \sphinxstyleemphasis{communication} may provide
additional useful information stored in the restart file.


\bigskip\hrule\bigskip

\phantomsection\label{\detokenize{Run_options:screen}}
\sphinxAtStartPar
\sphinxstylestrong{\sphinxhyphen{}screen file}

\sphinxAtStartPar
Specify a file for LAMMPS to write its screen information to.  In
one\sphinxhyphen{}partition mode, if the switch is not used, LAMMPS writes to the
screen.  If this switch is used, LAMMPS writes to the specified file
instead and you will see no screen output.  In multi\sphinxhyphen{}partition mode,
if the switch is not used, high\sphinxhyphen{}level status information is written to
the screen.  Each partition also writes to a screen.N file where N is
the partition ID.  If the switch is specified in multi\sphinxhyphen{}partition mode,
the high\sphinxhyphen{}level screen dump is named “file” and each partition also
writes screen information to a file.N.  For both one\sphinxhyphen{}partition and
multi\sphinxhyphen{}partition mode, if the specified file is “none”, then no screen
output is performed. Option \sphinxhyphen{}pscreen will override the name of the
partition screen files file.N.


\bigskip\hrule\bigskip

\phantomsection\label{\detokenize{Run_options:skiprun}}
\sphinxAtStartPar
\sphinxstylestrong{\sphinxhyphen{}skiprun}

\sphinxAtStartPar
Insert the command {\hyperref[\detokenize{timer::doc}]{\sphinxcrossref{\DUrole{doc}{timer timeout 0 every 1}}}} at the
beginning of an input file or after a {\hyperref[\detokenize{clear::doc}]{\sphinxcrossref{\DUrole{doc}{clear}}}} command.
This has the effect that the entire LAMMPS input script is processed
without executing actual {\hyperref[\detokenize{run::doc}]{\sphinxcrossref{\DUrole{doc}{run}}}} or {\hyperref[\detokenize{minimize::doc}]{\sphinxcrossref{\DUrole{doc}{minimize}}}}
and similar commands (their main loops are skipped).  This can be
helpful and convenient to test input scripts of long running
calculations for correctness to avoid having them crash after a
long time due to a typo or syntax error in the middle or at the end.


\bigskip\hrule\bigskip

\phantomsection\label{\detokenize{Run_options:suffix}}
\sphinxAtStartPar
\sphinxstylestrong{\sphinxhyphen{}suffix style args}

\sphinxAtStartPar
Use variants of various styles if they exist.  The specified style can
be \sphinxstyleemphasis{gpu}, \sphinxstyleemphasis{intel}, \sphinxstyleemphasis{kk}, \sphinxstyleemphasis{omp}, \sphinxstyleemphasis{opt}, or \sphinxstyleemphasis{hybrid}.  These
refer to optional packages that LAMMPS can be built with, as described
in {\hyperref[\detokenize{Speed::doc}]{\sphinxcrossref{\DUrole{doc}{Accelerate performance}}}}.  The “gpu” style corresponds to the
GPU package, the “intel” style to the INTEL package, the “kk”
style to the KOKKOS package, the “opt” style to the OPT package, and
the “omp” style to the OPENMP package. The hybrid style is the only
style that accepts arguments. It allows for two packages to be
specified. The first package specified is the default and will be used
if it is available. If no style is available for the first package,
the style for the second package will be used if available. For
example, \sphinxcode{\sphinxupquote{\sphinxhyphen{}suffix hybrid intel omp}} will use styles from the
INTEL package if they are installed and available, but styles for
the OPENMP package otherwise.

\sphinxAtStartPar
Along with the \sphinxcode{\sphinxupquote{\sphinxhyphen{}package}} command\sphinxhyphen{}line switch, this is a convenient
mechanism for invoking accelerator packages and their options without
having to edit an input script.

\sphinxAtStartPar
As an example, all of the packages provide a {\hyperref[\detokenize{pair_lj::doc}]{\sphinxcrossref{\DUrole{doc}{pair\_style lj/cut}}}} variant, with style names lj/cut/gpu,
lj/cut/intel, lj/cut/kk, lj/cut/omp, and lj/cut/opt.  A variant style
can be specified explicitly in your input script, e.g. pair\_style
lj/cut/gpu.  If the \sphinxhyphen{}suffix switch is used the specified suffix
(gpu,intel,kk,omp,opt) is automatically appended whenever your input
script command creates a new {\hyperref[\detokenize{atom_style::doc}]{\sphinxcrossref{\DUrole{doc}{atom style}}}},
{\hyperref[\detokenize{pair_style::doc}]{\sphinxcrossref{\DUrole{doc}{pair style}}}}, {\hyperref[\detokenize{fix::doc}]{\sphinxcrossref{\DUrole{doc}{fix}}}},
{\hyperref[\detokenize{compute::doc}]{\sphinxcrossref{\DUrole{doc}{compute}}}}, or {\hyperref[\detokenize{run_style::doc}]{\sphinxcrossref{\DUrole{doc}{run style}}}}.  If the
variant version does not exist, the standard version is created.

\sphinxAtStartPar
For the GPU package, using this command\sphinxhyphen{}line switch also invokes the
default GPU settings, as if the command “package gpu 1” were used at
the top of your input script.  These settings can be changed by using
the \sphinxcode{\sphinxupquote{\sphinxhyphen{}package gpu}} command\sphinxhyphen{}line switch or the {\hyperref[\detokenize{package::doc}]{\sphinxcrossref{\DUrole{doc}{package gpu}}}} command in your script.

\sphinxAtStartPar
For the INTEL package, using this command\sphinxhyphen{}line switch also
invokes the default INTEL settings, as if the command “package
intel 1” were used at the top of your input script.  These settings
can be changed by using the \sphinxcode{\sphinxupquote{\sphinxhyphen{}package intel}} command\sphinxhyphen{}line switch or
the {\hyperref[\detokenize{package::doc}]{\sphinxcrossref{\DUrole{doc}{package intel}}}} command in your script. If the
OPENMP package is also installed, the hybrid style with “intel omp”
arguments can be used to make the omp suffix a second choice, if a
requested style is not available in the INTEL package.  It will
also invoke the default OPENMP settings, as if the command “package
omp 0” were used at the top of your input script.  These settings can
be changed by using the \sphinxcode{\sphinxupquote{\sphinxhyphen{}package omp}} command\sphinxhyphen{}line switch or the
{\hyperref[\detokenize{package::doc}]{\sphinxcrossref{\DUrole{doc}{package omp}}}} command in your script.

\sphinxAtStartPar
For the KOKKOS package, using this command\sphinxhyphen{}line switch also invokes
the default KOKKOS settings, as if the command “package kokkos” were
used at the top of your input script.  These settings can be changed
by using the \sphinxcode{\sphinxupquote{\sphinxhyphen{}package kokkos}} command\sphinxhyphen{}line switch or the {\hyperref[\detokenize{package::doc}]{\sphinxcrossref{\DUrole{doc}{package kokkos}}}} command in your script.

\sphinxAtStartPar
For the OMP package, using this command\sphinxhyphen{}line switch also invokes the
default OMP settings, as if the command “package omp 0” were used at
the top of your input script.  These settings can be changed by using
the \sphinxcode{\sphinxupquote{\sphinxhyphen{}package omp}} command\sphinxhyphen{}line switch or the {\hyperref[\detokenize{package::doc}]{\sphinxcrossref{\DUrole{doc}{package omp}}}} command in your script.

\sphinxAtStartPar
The {\hyperref[\detokenize{suffix::doc}]{\sphinxcrossref{\DUrole{doc}{suffix}}}} command can also be used within an input
script to set a suffix, or to turn off or back on any suffix setting
made via the command line.


\bigskip\hrule\bigskip

\phantomsection\label{\detokenize{Run_options:var}}
\sphinxAtStartPar
\sphinxstylestrong{\sphinxhyphen{}var name value1 value2 …}

\sphinxAtStartPar
Specify a variable that will be defined for substitution purposes when
the input script is read.  This switch can be used multiple times to
define multiple variables.  “Name” is the variable name which can be a
single character (referenced as \$x in the input script) or a full
string (referenced as \$\{abc\}).  An {\hyperref[\detokenize{variable::doc}]{\sphinxcrossref{\DUrole{doc}{index\sphinxhyphen{}style variable}}}} will be created and populated with the
subsequent values, e.g. a set of filenames.  Using this command\sphinxhyphen{}line
option is equivalent to putting the line “variable name index value1
value2 …”  at the beginning of the input script.  Defining an index
variable as a command\sphinxhyphen{}line argument overrides any setting for the same
index variable in the input script, since index variables cannot be
re\sphinxhyphen{}defined.

\sphinxAtStartPar
See the {\hyperref[\detokenize{variable::doc}]{\sphinxcrossref{\DUrole{doc}{variable}}}} command for more info on defining
index and other kinds of variables and the {\hyperref[\detokenize{Commands_parse::doc}]{\sphinxcrossref{\DUrole{doc}{Parsing rules}}}}
page for more info on using variables in input scripts.

\begin{sphinxadmonition}{note}{Note:}
\sphinxAtStartPar
Currently, the command\sphinxhyphen{}line parser looks for arguments that
start with “\sphinxhyphen{}” to indicate new switches.  Thus you cannot specify
multiple variable values if any of them start with a “\sphinxhyphen{}”, e.g. a
negative numeric value.  It is OK if the first value1 starts with a
“\sphinxhyphen{}”, since it is automatically skipped.
\end{sphinxadmonition}

\sphinxstepscope


\section{Screen and logfile output}
\label{\detokenize{Run_output:screen-and-logfile-output}}\label{\detokenize{Run_output::doc}}
\sphinxAtStartPar
As LAMMPS reads an input script, it prints information to both the
screen and a log file about significant actions it takes to setup a
simulation.  When the simulation is ready to begin, LAMMPS performs
various initializations, and prints info about the run it is about to
perform, including the amount of memory (in MBytes per processor) that
the simulation requires.  It also prints details of the initial
thermodynamic state of the system.  During the run itself,
thermodynamic information is printed periodically, every few
timesteps.  When the run concludes, LAMMPS prints the final
thermodynamic state and a total run time for the simulation.  It also
appends statistics about the CPU time and storage requirements for the
simulation.  An example set of statistics is shown here:
\begin{sphinxalltt}
Loop time of 0.942801 on 4 procs for 300 steps with 2004 atoms

Performance: 54.985 ns/day, 0.436 hours/ns, 318.201 timesteps/s, 637.674 katom\sphinxhyphen{}step/s
195.2\% CPU use with 2 MPI tasks x 2 OpenMP threads

MPI task timing breakdown:
Section |  min time  |  avg time  |  max time  |\%varavg| \%total
\sphinxhyphen{}\sphinxhyphen{}\sphinxhyphen{}\sphinxhyphen{}\sphinxhyphen{}\sphinxhyphen{}\sphinxhyphen{}\sphinxhyphen{}\sphinxhyphen{}\sphinxhyphen{}\sphinxhyphen{}\sphinxhyphen{}\sphinxhyphen{}\sphinxhyphen{}\sphinxhyphen{}\sphinxhyphen{}\sphinxhyphen{}\sphinxhyphen{}\sphinxhyphen{}\sphinxhyphen{}\sphinxhyphen{}\sphinxhyphen{}\sphinxhyphen{}\sphinxhyphen{}\sphinxhyphen{}\sphinxhyphen{}\sphinxhyphen{}\sphinxhyphen{}\sphinxhyphen{}\sphinxhyphen{}\sphinxhyphen{}\sphinxhyphen{}\sphinxhyphen{}\sphinxhyphen{}\sphinxhyphen{}\sphinxhyphen{}\sphinxhyphen{}\sphinxhyphen{}\sphinxhyphen{}\sphinxhyphen{}\sphinxhyphen{}\sphinxhyphen{}\sphinxhyphen{}\sphinxhyphen{}\sphinxhyphen{}\sphinxhyphen{}\sphinxhyphen{}\sphinxhyphen{}\sphinxhyphen{}\sphinxhyphen{}\sphinxhyphen{}\sphinxhyphen{}\sphinxhyphen{}\sphinxhyphen{}\sphinxhyphen{}\sphinxhyphen{}\sphinxhyphen{}\sphinxhyphen{}\sphinxhyphen{}\sphinxhyphen{}\sphinxhyphen{}\sphinxhyphen{}\sphinxhyphen{}
Pair    | 0.61419    | 0.62872    | 0.64325    |   1.8 | 66.69
Bond    | 0.0028608  | 0.0028899  | 0.002919   |   0.1 |  0.31
Kspace  | 0.12652    | 0.14048    | 0.15444    |   3.7 | 14.90
Neigh   | 0.10242    | 0.10242    | 0.10242    |   0.0 | 10.86
Comm    | 0.026753   | 0.027593   | 0.028434   |   0.5 |  2.93
Output  | 0.00018341 | 0.00030942 | 0.00043542 |   0.0 |  0.03
Modify  | 0.039117   | 0.039348   | 0.039579   |   0.1 |  4.17
Other   |            | 0.001041   |            |       |  0.11

Nlocal:           1002 ave        1006 max         998 min
Histogram: 1 0 0 0 0 0 0 0 0 1
Nghost:         8670.5 ave        8691 max        8650 min
Histogram: 1 0 0 0 0 0 0 0 0 1
Neighs:         354010 ave      357257 max      350763 min
Histogram: 1 0 0 0 0 0 0 0 0 1

Total \# of neighbors = 708020
Ave neighs/atom = 353.30339
Ave special neighs/atom = 2.3403194
Neighbor list builds = 26
Dangerous builds = 0
\end{sphinxalltt}


\bigskip\hrule\bigskip


\sphinxAtStartPar
The first section provides a global loop timing summary. The \sphinxstyleemphasis{loop time}
is the total wall\sphinxhyphen{}clock time for the MD steps of the simulation run,
excluding the time for initialization and setup (i.e. the parts that may
be skipped with {\hyperref[\detokenize{run::doc}]{\sphinxcrossref{\DUrole{doc}{run N pre no}}}}).  The \sphinxstyleemphasis{Performance} line is
provided for convenience to help predict how long it will take to run a
desired physical simulation and to have numbers useful for performance
comparison between different simulation settings or system sizes.  The
\sphinxstyleemphasis{CPU use} line provides the CPU utilization per MPI task; it should be
close to 100\% times the number of OpenMP threads (or 1 if not using
OpenMP).  Lower numbers correspond to delays due to file I/O or
insufficient thread utilization from parts of the code that have not
been multi\sphinxhyphen{}threaded.


\bigskip\hrule\bigskip


\sphinxAtStartPar
The \sphinxstyleemphasis{MPI task} section gives the breakdown of the CPU run time (in
seconds) into major categories:
\begin{itemize}
\item {} 
\sphinxAtStartPar
\sphinxstyleemphasis{Pair} = non\sphinxhyphen{}bonded force computations

\item {} 
\sphinxAtStartPar
\sphinxstyleemphasis{Bond} = bonded interactions: bonds, angles, dihedrals, impropers

\item {} 
\sphinxAtStartPar
\sphinxstyleemphasis{Kspace} = long\sphinxhyphen{}range interactions: Ewald, PPPM, MSM

\item {} 
\sphinxAtStartPar
\sphinxstyleemphasis{Neigh} = neighbor list construction

\item {} 
\sphinxAtStartPar
\sphinxstyleemphasis{Comm} = inter\sphinxhyphen{}processor communication of atoms and their properties

\item {} 
\sphinxAtStartPar
\sphinxstyleemphasis{Output} = output of thermodynamic info and dump files

\item {} 
\sphinxAtStartPar
\sphinxstyleemphasis{Modify} = fixes and computes invoked by fixes

\item {} 
\sphinxAtStartPar
\sphinxstyleemphasis{Other} = all the remaining time

\end{itemize}

\sphinxAtStartPar
For each category, there is a breakdown of the least, average and most
amount of wall time any processor spent on this category of
computation.  The “\%varavg” is the percentage by which the max or min
varies from the average.  This is an indication of load imbalance.  A
percentage close to 0 is perfect load balance.  A large percentage is
imbalance.  The final “\%total” column is the percentage of the total
loop time is spent in this category.

\sphinxAtStartPar
When using the {\hyperref[\detokenize{timer::doc}]{\sphinxcrossref{\DUrole{doc}{timer full}}}} setting, an additional column
is added that also prints the CPU utilization in percent. In addition,
when using \sphinxstyleemphasis{timer full} and the {\hyperref[\detokenize{package::doc}]{\sphinxcrossref{\DUrole{doc}{package omp}}}} command are
active, a similar timing summary of time spent in threaded regions to
monitor thread utilization and load balance is provided. A new \sphinxstyleemphasis{Thread
timings} section is also added, which lists the time spent in reducing
the per\sphinxhyphen{}thread data elements to the storage for non\sphinxhyphen{}threaded
computation. These thread timings are measured for the first MPI rank
only and thus, because the breakdown for MPI tasks can change from
MPI rank to MPI rank, this breakdown can be very different for
individual ranks. Here is an example output for this section:
\begin{sphinxalltt}
Thread timings breakdown (MPI rank 0):
Total threaded time 0.6846 / 90.6\%
Section |  min time  |  avg time  |  max time  |\%varavg| \%total
\sphinxhyphen{}\sphinxhyphen{}\sphinxhyphen{}\sphinxhyphen{}\sphinxhyphen{}\sphinxhyphen{}\sphinxhyphen{}\sphinxhyphen{}\sphinxhyphen{}\sphinxhyphen{}\sphinxhyphen{}\sphinxhyphen{}\sphinxhyphen{}\sphinxhyphen{}\sphinxhyphen{}\sphinxhyphen{}\sphinxhyphen{}\sphinxhyphen{}\sphinxhyphen{}\sphinxhyphen{}\sphinxhyphen{}\sphinxhyphen{}\sphinxhyphen{}\sphinxhyphen{}\sphinxhyphen{}\sphinxhyphen{}\sphinxhyphen{}\sphinxhyphen{}\sphinxhyphen{}\sphinxhyphen{}\sphinxhyphen{}\sphinxhyphen{}\sphinxhyphen{}\sphinxhyphen{}\sphinxhyphen{}\sphinxhyphen{}\sphinxhyphen{}\sphinxhyphen{}\sphinxhyphen{}\sphinxhyphen{}\sphinxhyphen{}\sphinxhyphen{}\sphinxhyphen{}\sphinxhyphen{}\sphinxhyphen{}\sphinxhyphen{}\sphinxhyphen{}\sphinxhyphen{}\sphinxhyphen{}\sphinxhyphen{}\sphinxhyphen{}\sphinxhyphen{}\sphinxhyphen{}\sphinxhyphen{}\sphinxhyphen{}\sphinxhyphen{}\sphinxhyphen{}\sphinxhyphen{}\sphinxhyphen{}\sphinxhyphen{}\sphinxhyphen{}\sphinxhyphen{}\sphinxhyphen{}
Pair    | 0.5127     | 0.5147     | 0.5167     |   0.3 | 75.18
Bond    | 0.0043139  | 0.0046779  | 0.0050418  |   0.5 |  0.68
Kspace  | 0.070572   | 0.074541   | 0.07851    |   1.5 | 10.89
Neigh   | 0.084778   | 0.086969   | 0.089161   |   0.7 | 12.70
Reduce  | 0.0036485  | 0.003737   | 0.0038254  |   0.1 |  0.55
\end{sphinxalltt}


\bigskip\hrule\bigskip


\sphinxAtStartPar
The third section above lists the number of owned atoms (Nlocal),
ghost atoms (Nghost), and pairwise neighbors stored per processor.
The max and min values give the spread of these values across
processors with a 10\sphinxhyphen{}bin histogram showing the distribution. The total
number of histogram counts is equal to the number of processors.


\bigskip\hrule\bigskip


\sphinxAtStartPar
The last section gives aggregate statistics (across all processors)
for pairwise neighbors and special neighbors that LAMMPS keeps track
of (see the {\hyperref[\detokenize{special_bonds::doc}]{\sphinxcrossref{\DUrole{doc}{special\_bonds}}}} command).  The number
of times neighbor lists were rebuilt is tallied, as is the number of
potentially \sphinxstyleemphasis{dangerous} rebuilds.  If atom movement triggered neighbor
list rebuilding (see the {\hyperref[\detokenize{neigh_modify::doc}]{\sphinxcrossref{\DUrole{doc}{neigh\_modify}}}} command),
then dangerous reneighborings are those that were triggered on the
first timestep atom movement was checked for.  If this count is
non\sphinxhyphen{}zero you may wish to reduce the delay factor to ensure no force
interactions are missed by atoms moving beyond the neighbor skin
distance before a rebuild takes place.


\bigskip\hrule\bigskip


\sphinxAtStartPar
If an energy minimization was performed via the
{\hyperref[\detokenize{minimize::doc}]{\sphinxcrossref{\DUrole{doc}{minimize}}}} command, additional information is printed,
e.g.

\begin{sphinxVerbatim}[commandchars=\\\{\}]
Minimization stats:
  Stopping criterion = linesearch alpha is zero
  Energy initial, next\PYGZhy{}to\PYGZhy{}last, final =
         \PYGZhy{}6372.3765206     \PYGZhy{}8328.46998942     \PYGZhy{}8328.46998942
  Force two\PYGZhy{}norm initial, final = 1059.36 5.36874
  Force max component initial, final = 58.6026 1.46872
  Final line search alpha, max atom move = 2.7842e\PYGZhy{}10 4.0892e\PYGZhy{}10
  Iterations, force evaluations = 701 1516
\end{sphinxVerbatim}

\sphinxAtStartPar
The first line prints the criterion that determined minimization was
converged. The next line lists the initial and final energy, as well
as the energy on the next\sphinxhyphen{}to\sphinxhyphen{}last iteration.  The next 2 lines give a
measure of the gradient of the energy (force on all atoms).  The
2\sphinxhyphen{}norm is the “length” of this 3N\sphinxhyphen{}component force vector; the largest
component (x, y, or z) of force (infinity\sphinxhyphen{}norm) is also given.  Then
information is provided about the line search and statistics on how
many iterations and force\sphinxhyphen{}evaluations the minimizer required.
Multiple force evaluations are typically done at each iteration to
perform a 1d line minimization in the search direction.  See the
{\hyperref[\detokenize{minimize::doc}]{\sphinxcrossref{\DUrole{doc}{minimize}}}} page for more details.


\bigskip\hrule\bigskip


\sphinxAtStartPar
If a {\hyperref[\detokenize{kspace_style::doc}]{\sphinxcrossref{\DUrole{doc}{kspace\_style}}}} long\sphinxhyphen{}range Coulombics solver
that performs FFTs was used during the run (PPPM, Ewald), then
additional information is printed, e.g.

\begin{sphinxVerbatim}[commandchars=\\\{\}]
FFT time (\PYGZpc{} of Kspce) = 0.200313 (8.34477)
FFT Gflps 3d 1d\PYGZhy{}only = 2.31074 9.19989
\end{sphinxVerbatim}

\sphinxAtStartPar
The first line is the time spent doing 3d FFTs (several per timestep)
and the fraction it represents of the total KSpace time (listed
above).  Each 3d FFT requires computation (3 sets of 1d FFTs) and
communication (transposes).  The total flops performed is 5Nlog\_2(N),
where N is the number of points in the 3d grid.  The FFTs are timed
with and without the communication and a Gflop rate is computed.  The
3d rate is with communication; the 1d rate is without (just the 1d
FFTs).  Thus you can estimate what fraction of your FFT time was spent
in communication, roughly 75\% in the example above.

\sphinxstepscope


\section{Running LAMMPS on Windows}
\label{\detokenize{Run_windows:running-lammps-on-windows}}\label{\detokenize{Run_windows::doc}}
\sphinxAtStartPar
To run a serial (non\sphinxhyphen{}MPI) executable, follow these steps:
\begin{itemize}
\item {} 
\sphinxAtStartPar
Install a LAMMPS installer package from \sphinxurl{https://packages.lammps.org/windows.html}

\item {} 
\sphinxAtStartPar
Open the “Command Prompt” or “Terminal” app.

\item {} 
\sphinxAtStartPar
Change to the directory where you have your input script,
(e.g. by typing: cd “Documents”).

\item {} 
\sphinxAtStartPar
At the command prompt, type “lmp \sphinxhyphen{}in in.file.lmp”, where
\sphinxcode{\sphinxupquote{in.file.lmp}} is the name of your LAMMPS input script.

\end{itemize}

\sphinxAtStartPar
Note that the serial executable includes support for multi\sphinxhyphen{}threading
parallelization from the styles in the OPENMP and KOKKOS packages.
To run with 4 threads, you can type this:

\begin{sphinxVerbatim}[commandchars=\\\{\}]
lmp\PYG{+w}{ }\PYGZhy{}in\PYG{+w}{ }\PYG{k}{in}.lj.lmp\PYG{+w}{ }\PYGZhy{}pk\PYG{+w}{ }omp\PYG{+w}{ }\PYG{l+m}{4}\PYG{+w}{ }\PYGZhy{}sf\PYG{+w}{ }omp
lmp\PYG{+w}{ }\PYGZhy{}in\PYG{+w}{ }\PYG{k}{in}.lj.lmp\PYG{+w}{ }\PYGZhy{}k\PYG{+w}{ }on\PYG{+w}{ }t\PYG{+w}{ }\PYG{l+m}{4}\PYG{+w}{ }\PYGZhy{}sf\PYG{+w}{ }kk
\end{sphinxVerbatim}

\sphinxAtStartPar
Alternately, you can also install a package with LAMMPS\sphinxhyphen{}GUI included and
open the LAMMPS\sphinxhyphen{}GUI app (the package includes the command line version
of LAMMPS as well) and open the input file in the GUI and run it from
there.  For details on LAMMPS\sphinxhyphen{}GUI, see {\hyperref[\detokenize{Howto_lammps_gui::doc}]{\sphinxcrossref{\DUrole{doc}{Using LAMMPS\sphinxhyphen{}GUI}}}}.


\bigskip\hrule\bigskip


\sphinxAtStartPar
For the MS\sphinxhyphen{}MPI executables, which allow you to run LAMMPS under Windows
in parallel using MPI rather than multi\sphinxhyphen{}threading, follow these steps.

\sphinxAtStartPar
Download and install the MS\sphinxhyphen{}MPI runtime package \sphinxcode{\sphinxupquote{msmpisetup.exe}} from
\sphinxurl{https://www.microsoft.com/en-us/download/details.aspx?id=105289} (Note
that the \sphinxcode{\sphinxupquote{msmpisdk.msi}} is \sphinxstylestrong{only} required for \sphinxstylestrong{compilation} of
LAMMPS from source on Windows using Microsoft Visual Studio). After
installation of MS\sphinxhyphen{}MPI perform a reboot.

\sphinxAtStartPar
Then you can run the executable in serial like in the example above
or in parallel using MPI with one of the following commands:

\begin{sphinxVerbatim}[commandchars=\\\{\}]
mpiexec\PYG{+w}{ }\PYGZhy{}localonly\PYG{+w}{ }\PYG{l+m}{4}\PYG{+w}{ }lmp\PYG{+w}{ }\PYGZhy{}in\PYG{+w}{ }\PYG{k}{in}.file.lmp
mpiexec\PYG{+w}{ }\PYGZhy{}np\PYG{+w}{ }\PYG{l+m}{4}\PYG{+w}{ }lmp\PYG{+w}{ }\PYGZhy{}in\PYG{+w}{ }\PYG{k}{in}.file.lmp
\end{sphinxVerbatim}

\sphinxAtStartPar
where \sphinxcode{\sphinxupquote{in.file.lmp}} is the name of your LAMMPS input script. For the
latter case, you may be prompted to enter the password that you set
during installation of the MPI library software.

\sphinxAtStartPar
In this mode, output may not immediately show up on the screen, so if
your input script takes a long time to execute, you may need to be
patient before the output shows up.

\sphinxAtStartPar
Note that the parallel executable also includes OpenMP multi\sphinxhyphen{}threading
through both the OPENMP and the KOKKOS package, which can be combined
with MPI using something like:

\begin{sphinxVerbatim}[commandchars=\\\{\}]
mpiexec\PYG{+w}{ }\PYGZhy{}localonly\PYG{+w}{ }\PYG{l+m}{2}\PYG{+w}{ }lmp\PYG{+w}{ }\PYGZhy{}in\PYG{+w}{ }\PYG{k}{in}.lj.lmp\PYG{+w}{ }\PYGZhy{}pk\PYG{+w}{ }omp\PYG{+w}{ }\PYG{l+m}{2}\PYG{+w}{ }\PYGZhy{}sf\PYG{+w}{ }omp
mpiexec\PYG{+w}{ }\PYGZhy{}localonly\PYG{+w}{ }\PYG{l+m}{2}\PYG{+w}{ }lmp\PYG{+w}{ }\PYGZhy{}in\PYG{+w}{ }\PYG{k}{in}.lj.lmp\PYG{+w}{ }\PYGZhy{}kokkos\PYG{+w}{ }on\PYG{+w}{ }t\PYG{+w}{ }\PYG{l+m}{2}\PYG{+w}{ }\PYGZhy{}sf\PYG{+w}{ }kk
\end{sphinxVerbatim}


\bigskip\hrule\bigskip


\sphinxAtStartPar
MPI parallelization will work for \sphinxstyleemphasis{all} functionality in LAMMPS and in
many cases the MPI parallelization is more efficient than
multi\sphinxhyphen{}threading since LAMMPS was designed from ground up for MPI
parallelization using domain decomposition.  Multi\sphinxhyphen{}threading is only
available for selected styles and implemented on top of the MPI
parallelization.  Multi\sphinxhyphen{}threading is most useful for systems with large
load imbalances when using domain decomposition and a smaller number
of threads (\textless{}= 8).

\sphinxstepscope


\chapter{Commands}
\label{\detokenize{Commands:commands}}\label{\detokenize{Commands::doc}}
\sphinxAtStartPar
These pages describe how a LAMMPS input script is formatted and the
commands in it are used to define a LAMMPS simulation.

\sphinxstepscope


\section{LAMMPS input scripts}
\label{\detokenize{Commands_input:lammps-input-scripts}}\label{\detokenize{Commands_input::doc}}
\sphinxAtStartPar
LAMMPS executes calculations by reading commands from a input script (text file), one
line at a time.  When the input script ends, LAMMPS exits.  This is different
from programs that read and process the entire input before starting a calculation.

\sphinxAtStartPar
Each command causes LAMMPS to take some immediate action without regard
for any commands that may be processed later. Commands may set an
internal variable, read in a file, or run a simulation.  These actions
can be grouped into three categories:
\begin{enumerate}
\sphinxsetlistlabels{\alph}{enumi}{enumii}{}{)}%
\item {} 
\sphinxAtStartPar
commands that change a global setting (examples: {\hyperref[\detokenize{timestep::doc}]{\sphinxcrossref{\DUrole{doc}{timestep}}}}, {\hyperref[\detokenize{newton::doc}]{\sphinxcrossref{\DUrole{doc}{newton}}}},
{\hyperref[\detokenize{echo::doc}]{\sphinxcrossref{\DUrole{doc}{echo}}}}, {\hyperref[\detokenize{log::doc}]{\sphinxcrossref{\DUrole{doc}{log}}}}, {\hyperref[\detokenize{thermo::doc}]{\sphinxcrossref{\DUrole{doc}{thermo}}}}, {\hyperref[\detokenize{restart::doc}]{\sphinxcrossref{\DUrole{doc}{restart}}}}),

\item {} 
\sphinxAtStartPar
commands that add, modify, remove, or replace “styles” that are
executed during a “run” (examples: {\hyperref[\detokenize{pair_style::doc}]{\sphinxcrossref{\DUrole{doc}{pair\_style}}}}, {\hyperref[\detokenize{fix::doc}]{\sphinxcrossref{\DUrole{doc}{fix}}}}, {\hyperref[\detokenize{compute::doc}]{\sphinxcrossref{\DUrole{doc}{compute}}}}, {\hyperref[\detokenize{dump::doc}]{\sphinxcrossref{\DUrole{doc}{dump}}}},
{\hyperref[\detokenize{thermo_style::doc}]{\sphinxcrossref{\DUrole{doc}{thermo\_style}}}}, {\hyperref[\detokenize{pair_modify::doc}]{\sphinxcrossref{\DUrole{doc}{pair\_modify}}}}), and

\item {} 
\sphinxAtStartPar
commands that execute a “run” or perform some other computation or
operation (examples: {\hyperref[\detokenize{print::doc}]{\sphinxcrossref{\DUrole{doc}{print}}}}, {\hyperref[\detokenize{run::doc}]{\sphinxcrossref{\DUrole{doc}{run}}}}, {\hyperref[\detokenize{minimize::doc}]{\sphinxcrossref{\DUrole{doc}{minimize}}}}, {\hyperref[\detokenize{temper::doc}]{\sphinxcrossref{\DUrole{doc}{temper}}}}, {\hyperref[\detokenize{write_dump::doc}]{\sphinxcrossref{\DUrole{doc}{write\_dump}}}}, {\hyperref[\detokenize{rerun::doc}]{\sphinxcrossref{\DUrole{doc}{rerun}}}},
{\hyperref[\detokenize{read_data::doc}]{\sphinxcrossref{\DUrole{doc}{read\_data}}}}, {\hyperref[\detokenize{read_restart::doc}]{\sphinxcrossref{\DUrole{doc}{read\_restart}}}})

\end{enumerate}

\sphinxAtStartPar
Commands in category a) have default settings, which means you only
need to use the command if you wish to change the defaults.

\sphinxAtStartPar
In many cases, the ordering of commands in an input script is not
important, but can have consequences when the global state is changed
between commands in the c) category. The following rules apply:
\begin{enumerate}
\sphinxsetlistlabels{\arabic}{enumi}{enumii}{(}{)}%
\item {} 
\sphinxAtStartPar
LAMMPS does not read your entire input script and then perform a
simulation with all the settings.  Rather, the input script is read
one line at a time and each command takes effect when it is read.
Thus this sequence of commands:

\begin{sphinxVerbatim}[commandchars=\\\{\}]
\PYG{k}{timestep}\PYG{+w}{ }\PYG{l+m}{0.5}
\PYG{k}{run}\PYG{+w}{      }\PYG{l+m}{100}
\PYG{k}{run}\PYG{+w}{      }\PYG{l+m}{100}
\end{sphinxVerbatim}

\sphinxAtStartPar
does something different than this sequence:

\begin{sphinxVerbatim}[commandchars=\\\{\}]
\PYG{k}{run}\PYG{+w}{      }\PYG{l+m}{100}
\PYG{k}{timestep}\PYG{+w}{ }\PYG{l+m}{0.5}
\PYG{k}{run}\PYG{+w}{      }\PYG{l+m}{100}
\end{sphinxVerbatim}

\sphinxAtStartPar
In the first case, the specified timestep (0.5 fs) is used for two
simulations of 100 timesteps each.  In the second case, the default
timestep (1.0 fs) is used for the first 100 step simulation and a
0.5 fs timestep is used for the second one.

\item {} 
\sphinxAtStartPar
Some commands are only valid when they follow other commands.  For
example you cannot set the temperature of a group of atoms until
atoms have been defined and a group command is used to define which
atoms belong to the group.

\item {} 
\sphinxAtStartPar
Sometimes command B will use values that can be set by command A.
This means command A must precede command B in the input script if
it is to have the desired effect.  For example, the {\hyperref[\detokenize{read_data::doc}]{\sphinxcrossref{\DUrole{doc}{read\_data}}}} command initializes the system by setting up the
simulation box and assigning atoms to processors.  If default values
are not desired, the {\hyperref[\detokenize{processors::doc}]{\sphinxcrossref{\DUrole{doc}{processors}}}} and
{\hyperref[\detokenize{boundary::doc}]{\sphinxcrossref{\DUrole{doc}{boundary}}}} commands need to be used before \sphinxcode{\sphinxupquote{read\_data}}
to tell LAMMPS how to map processors to the simulation box.

\end{enumerate}

\sphinxAtStartPar
Many input script errors are detected by LAMMPS and an ERROR or
WARNING message is printed.  The {\hyperref[\detokenize{Errors::doc}]{\sphinxcrossref{\DUrole{doc}{Errors}}}} page gives
more information on what errors mean.  The documentation for each
command lists restrictions on how the command can be used.

\sphinxAtStartPar
You can use the {\hyperref[\detokenize{Run_options:skiprun}]{\sphinxcrossref{\DUrole{std,std-ref}{\sphinxhyphen{}skiprun}}}} command line flag
to have LAMMPS skip the execution of any \sphinxcode{\sphinxupquote{run}}, \sphinxcode{\sphinxupquote{minimize}}, or similar
commands to check the entire input for correct syntax to avoid crashes
on typos or syntax errors in long runs.

\sphinxstepscope


\section{Parsing rules for input scripts}
\label{\detokenize{Commands_parse:parsing-rules-for-input-scripts}}\label{\detokenize{Commands_parse::doc}}
\sphinxAtStartPar
Each non\sphinxhyphen{}blank line in the input script is treated as a command.
LAMMPS commands are case sensitive.  Command names are lower\sphinxhyphen{}case, as
are specified command arguments.  Upper case letters may be used in
file names or user\sphinxhyphen{}chosen ID strings.

\sphinxAtStartPar
Here are 6 rules for how each line in the input script is parsed by
LAMMPS:

\phantomsection\label{\detokenize{Commands_parse:one}}\begin{enumerate}
\sphinxsetlistlabels{\arabic}{enumi}{enumii}{}{.}%
\item {} 
\sphinxAtStartPar
If the last printable character on the line is a “\&” character, the
command is assumed to continue on the next line.  The next line is
concatenated to the previous line by removing the “\&” character and
line break.  This allows long commands to be continued across two or
more lines.  See the discussion of triple quotes in {\hyperref[\detokenize{Commands_parse:six}]{\sphinxcrossref{\DUrole{std,std-ref}{6}}}}
for how to continue a command across multiple line without using “\&”
characters.

\end{enumerate}
\phantomsection\label{\detokenize{Commands_parse:two}}\begin{enumerate}
\sphinxsetlistlabels{\arabic}{enumi}{enumii}{}{.}%
\setcounter{enumi}{1}
\item {} 
\sphinxAtStartPar
All characters from the first “\#” character onward are treated as
comment and discarded.  The exception to this rule is described in
{\hyperref[\detokenize{Commands_parse:six}]{\sphinxcrossref{\DUrole{std,std-ref}{6}}}}. Note that a comment after a trailing “\&” character
will prevent the command from continuing on the next line.  Also note
that for multi\sphinxhyphen{}line commands a single leading “\#” will comment out
the entire command.

\begin{sphinxVerbatim}[commandchars=\\\{\}]
\PYG{c}{\PYGZsh{} this is a comment}
\PYG{k}{timestep}\PYG{+w}{  }\PYG{l+m}{1.0}\PYG{+w}{   }\PYG{c}{\PYGZsh{} this is also a comment}
\end{sphinxVerbatim}

\end{enumerate}
\phantomsection\label{\detokenize{Commands_parse:three}}\begin{enumerate}
\sphinxsetlistlabels{\arabic}{enumi}{enumii}{}{.}%
\setcounter{enumi}{2}
\item {} 
\sphinxAtStartPar
The line is searched repeatedly for \$ characters, which indicate
variables that are replaced with a text string.  The exception to
this rule is described in {\hyperref[\detokenize{Commands_parse:six}]{\sphinxcrossref{\DUrole{std,std-ref}{6}}}}.

\sphinxAtStartPar
If the \$ is followed by text in curly brackets ‘\{\}’, then the
variable name is the text inside the curly brackets.  If no curly
brackets follow the \$, then the variable name is the single character
immediately following the \$.  Thus \sphinxcode{\sphinxupquote{\$\{myTemp\}}} and \sphinxcode{\sphinxupquote{\$x}} refer to variables
named “myTemp” and “x”, while \sphinxcode{\sphinxupquote{\$xx}} will be interpreted as a variable
named “x” followed by an “x” character.

\sphinxAtStartPar
How the variable is converted to a text string depends on what style
of variable it is; see the {\hyperref[\detokenize{variable::doc}]{\sphinxcrossref{\DUrole{doc}{variable}}}} page for
details.  It can be a variable that stores multiple text strings, and
return one of them.  The returned text string can be multiple “words”
(space separated) which will then be interpreted as multiple
arguments in the input command.  The variable can also store a
numeric formula which will be evaluated and its numeric result
returned as a string.

\sphinxAtStartPar
As a special case, if the \$ is followed by parenthesis “()”, then the
text inside the parenthesis is treated as an “immediate” variable and
evaluated as an {\hyperref[\detokenize{variable::doc}]{\sphinxcrossref{\DUrole{doc}{equal\sphinxhyphen{}style variable}}}}.  This is a
way to use numeric formulas in an input script without having to
assign them to variable names.  For example, these 3 input script
lines:

\begin{sphinxVerbatim}[commandchars=\\\{\}]
\PYG{k}{variable }\PYG{n+nv+nvIdentifier}{X}\PYG{+w}{ }\PYG{n}{equal}\PYG{+w}{ }\PYG{n+nv}{(}\PYG{n+nv}{xlo+xhi}\PYG{n+nv}{)}\PYG{o}{/}\PYG{l+m}{2}\PYG{o}{+}\PYG{n}{sqrt}\PYG{n+nv}{(}\PYG{n+nv}{v\PYGZus{}area}\PYG{n+nv}{)}
\PYG{k}{region }\PYG{n+nv+nvIdentifier}{1}\PYG{+w}{ }\PYG{n}{block}\PYG{+w}{ }\PYG{n+nv}{\PYGZdl{}X}\PYG{+w}{ }\PYG{l+m}{2}\PYG{+w}{ }\PYG{n}{INF}\PYG{+w}{ }\PYG{n}{INF}\PYG{+w}{ }\PYG{n}{EDGE}\PYG{+w}{ }\PYG{n}{EDGE}
\PYG{k}{variable }\PYG{n+nv+nvIdentifier}{X}\PYG{+w}{ }\PYG{n}{delete}
\end{sphinxVerbatim}

\sphinxAtStartPar
can be replaced by:

\begin{sphinxVerbatim}[commandchars=\\\{\}]
\PYG{k}{region }\PYG{n+nv+nvIdentifier}{1}\PYG{+w}{ }\PYG{n}{block}\PYG{+w}{ }\PYG{n+nv}{\PYGZdl{}(}\PYG{n+nv}{(}\PYG{n+nv}{xlo+xhi}\PYG{n+nv}{)}\PYG{n+nv}{/2+sqrt}\PYG{n+nv}{(}\PYG{n+nv}{v\PYGZus{}area}\PYG{n+nv}{)}\PYG{n+nv}{)}\PYG{+w}{ }\PYG{l+m}{2}\PYG{+w}{ }\PYG{n}{INF}\PYG{+w}{ }\PYG{n}{INF}\PYG{+w}{ }\PYG{n}{EDGE}\PYG{+w}{ }\PYG{n}{EDGE}
\end{sphinxVerbatim}

\sphinxAtStartPar
so that you do not have to define (or discard) a temporary variable,
“X” in this case.

\sphinxAtStartPar
Additionally, the entire “immediate” variable expression may be
followed by a colon, followed by a C\sphinxhyphen{}style format string,
e.g. \sphinxcode{\sphinxupquote{:\%f}} or \sphinxcode{\sphinxupquote{:\%.10g}}.  The format string must be appropriate for
a double\sphinxhyphen{}precision floating\sphinxhyphen{}point value.  The format string is used
to output the result of the variable expression evaluation.  If a
format string is not specified, a high\sphinxhyphen{}precision \sphinxcode{\sphinxupquote{\%.20g}} is used as
the default format.

\sphinxAtStartPar
This can be useful for formatting print output to a desired precision:

\begin{sphinxVerbatim}[commandchars=\\\{\}]
\PYG{k}{print}\PYG{+w}{ }\PYG{l+s}{\PYGZdq{}}\PYG{l+s}{Final energy per atom: \PYGZdl{}(v\PYGZus{}ke\PYGZus{}per\PYGZus{}atom+v\PYGZus{}pe\PYGZus{}per\PYGZus{}atom:\PYGZpc{}10.3f) eV/atom}\PYG{l+s}{\PYGZdq{}}
\end{sphinxVerbatim}

\sphinxAtStartPar
Note that neither the curly\sphinxhyphen{}bracket or immediate form of variables
can contain nested \$ characters for other variables to substitute
for.  Thus you may \sphinxstylestrong{NOT} do this:

\begin{sphinxVerbatim}[commandchars=\\\{\}]
\PYG{k}{variable        }\PYG{n+nv+nvIdentifier}{a}\PYG{+w}{ }\PYG{n}{equal}\PYG{+w}{ }\PYG{l+m}{2}
\PYG{k}{variable        }\PYG{n+nv+nvIdentifier}{b2}\PYG{+w}{ }\PYG{n}{equal}\PYG{+w}{ }\PYG{l+m}{4}
\PYG{k}{print}\PYG{+w}{           }\PYG{l+s}{\PYGZdq{}}\PYG{l+s}{B2 = \PYGZdl{}\PYGZob{}b\PYGZdl{}a\PYGZcb{}}\PYG{l+s}{\PYGZdq{}}
\end{sphinxVerbatim}

\sphinxAtStartPar
Nor can you specify an expression like \sphinxcode{\sphinxupquote{\$(\$x\sphinxhyphen{}1.0)}} for an immediate
variable, but you could use \sphinxcode{\sphinxupquote{\$(v\_x\sphinxhyphen{}1.0)}}, since the latter is valid
syntax for an {\hyperref[\detokenize{variable::doc}]{\sphinxcrossref{\DUrole{doc}{equal\sphinxhyphen{}style variable}}}}.

\sphinxAtStartPar
See the {\hyperref[\detokenize{variable::doc}]{\sphinxcrossref{\DUrole{doc}{variable}}}} command for more details of how
strings are assigned to variables and evaluated, and how they can
be used in input script commands.

\end{enumerate}
\phantomsection\label{\detokenize{Commands_parse:four}}\begin{enumerate}
\sphinxsetlistlabels{\arabic}{enumi}{enumii}{}{.}%
\setcounter{enumi}{3}
\item {} 
\sphinxAtStartPar
The line is broken into “words” separated by white\sphinxhyphen{}space (tabs,
spaces).  Note that words can thus contain letters, digits,
underscores, or punctuation characters.

\end{enumerate}
\phantomsection\label{\detokenize{Commands_parse:five}}\begin{enumerate}
\sphinxsetlistlabels{\arabic}{enumi}{enumii}{}{.}%
\setcounter{enumi}{4}
\item {} 
\sphinxAtStartPar
The first word is the command name.  All successive words in the line
are arguments.

\end{enumerate}
\phantomsection\label{\detokenize{Commands_parse:six}}\begin{enumerate}
\sphinxsetlistlabels{\arabic}{enumi}{enumii}{}{.}%
\setcounter{enumi}{5}
\item {} 
\sphinxAtStartPar
If you want text with spaces to be treated as a single argument, it
can be enclosed in either single (’) or double (”) or triple (“””)
quotes.  A long single argument enclosed in single or double quotes
can span multiple lines if the “\&” character is used, as described
in {\hyperref[\detokenize{Commands_parse:one}]{\sphinxcrossref{\DUrole{std,std-ref}{1}}}} above.  When the lines are concatenated together
by LAMMPS (and the “\&” characters and line breaks removed), the
combined text will become a single line.  If you want multiple lines
of an argument to retain their line breaks, the text can be enclosed
in triple quotes, in which case “\&” characters are not needed and do
not function as line continuation character.  For example:

\begin{sphinxVerbatim}[commandchars=\\\{\}]
\PYG{k}{print}\PYG{+w}{ }\PYG{l+s}{\PYGZdq{}}\PYG{l+s}{Volume = \PYGZdl{}v}\PYG{l+s}{\PYGZdq{}}
\PYG{k}{print}\PYG{+w}{ }\PYG{l+s}{\PYGZsq{}}\PYG{l+s}{Volume = \PYGZdl{}v}\PYG{l+s}{\PYGZsq{}}
\PYG{k}{if}\PYG{+w}{ }\PYG{l+s}{\PYGZdq{}}\PYG{l+s}{\PYGZdl{}\PYGZob{}steps\PYGZcb{} \PYGZgt{} 1000}\PYG{l+s}{\PYGZdq{}}\PYG{+w}{ }\PYG{k}{then}\PYG{+w}{ }\PYG{n}{quit}
\PYG{k}{variable }\PYG{n+nv+nvIdentifier}{a}\PYG{+w}{ }\PYG{n}{string}\PYG{+w}{ }\PYG{l+s}{\PYGZdq{}}\PYG{l+s}{red green blue \PYGZam{}}
\PYG{l+s}{                purple orange cyan}\PYG{l+s}{\PYGZdq{}}
\PYG{k}{print}\PYG{+w}{ }\PYG{l+s}{\PYGZdq{}}\PYG{l+s}{\PYGZdq{}}\PYG{l+s}{\PYGZdq{}}
\PYG{l+s}{System volume = \PYGZdl{}v}
\PYG{l+s}{System temperature = \PYGZdl{}t}
\PYG{l+s}{\PYGZdq{}}\PYG{l+s}{\PYGZdq{}}\PYG{l+s}{\PYGZdq{}}
\end{sphinxVerbatim}

\sphinxAtStartPar
In each of these cases, the single, double, or triple quotes are
removed and the enclosed text stored internally as a single
argument.

\sphinxAtStartPar
See the {\hyperref[\detokenize{dump_modify::doc}]{\sphinxcrossref{\DUrole{doc}{dump modify format}}}}, {\hyperref[\detokenize{print::doc}]{\sphinxcrossref{\DUrole{doc}{print}}}}, {\hyperref[\detokenize{if::doc}]{\sphinxcrossref{\DUrole{doc}{if}}}}, and {\hyperref[\detokenize{python::doc}]{\sphinxcrossref{\DUrole{doc}{python}}}} commands for
examples.

\sphinxAtStartPar
A “\#” or “\$” character that is between quotes will not be treated as
a comment indicator in {\hyperref[\detokenize{Commands_parse:two}]{\sphinxcrossref{\DUrole{std,std-ref}{2}}}} or substituted for as a
variable in {\hyperref[\detokenize{Commands_parse:three}]{\sphinxcrossref{\DUrole{std,std-ref}{3}}}}.

\end{enumerate}

\begin{sphinxadmonition}{note}{Note:}
\sphinxAtStartPar
If the argument is itself a command that requires a quoted
argument (e.g. using a {\hyperref[\detokenize{print::doc}]{\sphinxcrossref{\DUrole{doc}{print}}}} command as part of an
{\hyperref[\detokenize{if::doc}]{\sphinxcrossref{\DUrole{doc}{if}}}} or {\hyperref[\detokenize{run::doc}]{\sphinxcrossref{\DUrole{doc}{run every}}}} command), then single, double, or
triple quotes can be nested in the usual manner.  See the doc pages
for those commands for examples.  Only one of level of nesting is
allowed, but that should be sufficient for most use cases.
\end{sphinxadmonition}

\begin{sphinxadmonition}{note}{ASCII versus UTF\sphinxhyphen{}8}

\sphinxAtStartPar
LAMMPS expects and processes 7\sphinxhyphen{}bit ASCII format text internally.
Many modern environments use UTF\sphinxhyphen{}8 encoding, which is a superset
of the 7\sphinxhyphen{}bit ASCII character table and thus mostly compatible.
However, there are several non\sphinxhyphen{}ASCII characters that can look
very similar to their ASCII equivalents or are invisible (so they
look like a blank), but are encoded differently.  Web browsers,
PDF viewers, document editors are known to sometimes replace one
with the other for a better looking output.  However, that can
lead to problems, for instance, when using cut\sphinxhyphen{}n\sphinxhyphen{}paste of input
file examples from web pages, or when using a document editor
(not a dedicated plain text editor) for writing LAMMPS inputs.
LAMMPS will try to detect this and substitute the non\sphinxhyphen{}ASCII
characters with their ASCII equivalents where known.  There also
is going to be a warning printed, if this occurs.  It is
recommended to avoid such characters altogether in LAMMPS input,
data and potential files.  The replacement tables are likely
incomplete and dependent on users reporting problems processing
correctly looking input containing UTF\sphinxhyphen{}8 encoded non\sphinxhyphen{}ASCII
characters.
\end{sphinxadmonition}

\sphinxstepscope


\section{Input script structure}
\label{\detokenize{Commands_structure:input-script-structure}}\label{\detokenize{Commands_structure::doc}}
\sphinxAtStartPar
This page describes the structure of a typical LAMMPS input script.
The examples directory in the LAMMPS distribution contains many sample
input scripts; it is discussed on the {\hyperref[\detokenize{Examples::doc}]{\sphinxcrossref{\DUrole{doc}{Examples}}}} doc
page.

\sphinxAtStartPar
A LAMMPS input script typically has 4 parts:
\begin{enumerate}
\sphinxsetlistlabels{\arabic}{enumi}{enumii}{}{.}%
\item {} 
\sphinxAtStartPar
{\hyperref[\detokenize{Commands_structure:init}]{\sphinxcrossref{\DUrole{std,std-ref}{Initialization}}}}

\item {} 
\sphinxAtStartPar
{\hyperref[\detokenize{Commands_structure:system}]{\sphinxcrossref{\DUrole{std,std-ref}{System definition}}}}

\item {} 
\sphinxAtStartPar
{\hyperref[\detokenize{Commands_structure:settings}]{\sphinxcrossref{\DUrole{std,std-ref}{Simulation settings}}}}

\item {} 
\sphinxAtStartPar
{\hyperref[\detokenize{Commands_structure:run}]{\sphinxcrossref{\DUrole{std,std-ref}{Run a simulation}}}}

\end{enumerate}

\sphinxAtStartPar
The last 2 parts can be repeated as many times as desired.  I.e. run a
simulation, change some settings, run some more, etc.  Each of the 4
parts is now described in more detail.  Remember that almost all
commands need only be used if a non\sphinxhyphen{}default value is desired.


\subsection{Initialization}
\label{\detokenize{Commands_structure:initialization}}\label{\detokenize{Commands_structure:init}}
\sphinxAtStartPar
Set parameters that need to be defined before atoms are created or
read\sphinxhyphen{}in from a file.

\sphinxAtStartPar
The relevant commands are {\hyperref[\detokenize{units::doc}]{\sphinxcrossref{\DUrole{doc}{units}}}},
{\hyperref[\detokenize{dimension::doc}]{\sphinxcrossref{\DUrole{doc}{dimension}}}}, {\hyperref[\detokenize{newton::doc}]{\sphinxcrossref{\DUrole{doc}{newton}}}},
{\hyperref[\detokenize{processors::doc}]{\sphinxcrossref{\DUrole{doc}{processors}}}}, {\hyperref[\detokenize{boundary::doc}]{\sphinxcrossref{\DUrole{doc}{boundary}}}},
{\hyperref[\detokenize{atom_style::doc}]{\sphinxcrossref{\DUrole{doc}{atom\_style}}}}, {\hyperref[\detokenize{atom_modify::doc}]{\sphinxcrossref{\DUrole{doc}{atom\_modify}}}}.

\sphinxAtStartPar
If force\sphinxhyphen{}field parameters appear in the files that will be read, these
commands tell LAMMPS what kinds of force fields are being used:
{\hyperref[\detokenize{pair_style::doc}]{\sphinxcrossref{\DUrole{doc}{pair\_style}}}}, {\hyperref[\detokenize{bond_style::doc}]{\sphinxcrossref{\DUrole{doc}{bond\_style}}}},
{\hyperref[\detokenize{angle_style::doc}]{\sphinxcrossref{\DUrole{doc}{angle\_style}}}}, {\hyperref[\detokenize{dihedral_style::doc}]{\sphinxcrossref{\DUrole{doc}{dihedral\_style}}}},
{\hyperref[\detokenize{improper_style::doc}]{\sphinxcrossref{\DUrole{doc}{improper\_style}}}}.


\subsection{System definition}
\label{\detokenize{Commands_structure:system-definition}}\label{\detokenize{Commands_structure:system}}
\sphinxAtStartPar
There are 3 ways to define the simulation cell and reserve space for
force field info and fill it with atoms in LAMMPS.  Read them in from
(1) a data file or (2) a restart file via the {\hyperref[\detokenize{read_data::doc}]{\sphinxcrossref{\DUrole{doc}{read\_data}}}} or {\hyperref[\detokenize{read_restart::doc}]{\sphinxcrossref{\DUrole{doc}{read\_restart}}}} commands,
respectively.  These files can also contain molecular topology
information.  Or (3) create a simulation cell and fill it with atoms on
a lattice (with no molecular topology), using these commands:
{\hyperref[\detokenize{lattice::doc}]{\sphinxcrossref{\DUrole{doc}{lattice}}}}, {\hyperref[\detokenize{region::doc}]{\sphinxcrossref{\DUrole{doc}{region}}}}, {\hyperref[\detokenize{create_box::doc}]{\sphinxcrossref{\DUrole{doc}{create\_box}}}}, {\hyperref[\detokenize{create_atoms::doc}]{\sphinxcrossref{\DUrole{doc}{create\_atoms}}}} or
{\hyperref[\detokenize{read_dump::doc}]{\sphinxcrossref{\DUrole{doc}{read\_dump}}}}.

\sphinxAtStartPar
The entire set of atoms can be duplicated to make a larger simulation
using the {\hyperref[\detokenize{replicate::doc}]{\sphinxcrossref{\DUrole{doc}{replicate}}}} command.


\subsection{Simulation settings}
\label{\detokenize{Commands_structure:simulation-settings}}\label{\detokenize{Commands_structure:settings}}
\sphinxAtStartPar
Once atoms and molecular topology are defined, a variety of settings
can be specified: force field coefficients, simulation parameters,
output options, and more.

\sphinxAtStartPar
Force field coefficients are set by these commands (they can also be
set in the read\sphinxhyphen{}in files): {\hyperref[\detokenize{pair_coeff::doc}]{\sphinxcrossref{\DUrole{doc}{pair\_coeff}}}},
{\hyperref[\detokenize{bond_coeff::doc}]{\sphinxcrossref{\DUrole{doc}{bond\_coeff}}}}, {\hyperref[\detokenize{angle_coeff::doc}]{\sphinxcrossref{\DUrole{doc}{angle\_coeff}}}},
{\hyperref[\detokenize{dihedral_coeff::doc}]{\sphinxcrossref{\DUrole{doc}{dihedral\_coeff}}}},
{\hyperref[\detokenize{improper_coeff::doc}]{\sphinxcrossref{\DUrole{doc}{improper\_coeff}}}},
{\hyperref[\detokenize{kspace_style::doc}]{\sphinxcrossref{\DUrole{doc}{kspace\_style}}}}, {\hyperref[\detokenize{dielectric::doc}]{\sphinxcrossref{\DUrole{doc}{dielectric}}}},
{\hyperref[\detokenize{special_bonds::doc}]{\sphinxcrossref{\DUrole{doc}{special\_bonds}}}}.

\sphinxAtStartPar
Various simulation parameters are set by these commands:
{\hyperref[\detokenize{neighbor::doc}]{\sphinxcrossref{\DUrole{doc}{neighbor}}}}, {\hyperref[\detokenize{neigh_modify::doc}]{\sphinxcrossref{\DUrole{doc}{neigh\_modify}}}},
{\hyperref[\detokenize{group::doc}]{\sphinxcrossref{\DUrole{doc}{group}}}}, {\hyperref[\detokenize{timestep::doc}]{\sphinxcrossref{\DUrole{doc}{timestep}}}},
{\hyperref[\detokenize{reset_timestep::doc}]{\sphinxcrossref{\DUrole{doc}{reset\_timestep}}}}, {\hyperref[\detokenize{run_style::doc}]{\sphinxcrossref{\DUrole{doc}{run\_style}}}},
{\hyperref[\detokenize{min_style::doc}]{\sphinxcrossref{\DUrole{doc}{min\_style}}}}, {\hyperref[\detokenize{min_modify::doc}]{\sphinxcrossref{\DUrole{doc}{min\_modify}}}}.

\sphinxAtStartPar
Fixes impose a variety of boundary conditions, time integration, and
diagnostic options.  The {\hyperref[\detokenize{fix::doc}]{\sphinxcrossref{\DUrole{doc}{fix}}}} command comes in many flavors.

\sphinxAtStartPar
Various computations can be specified for execution during a
simulation using the {\hyperref[\detokenize{compute::doc}]{\sphinxcrossref{\DUrole{doc}{compute}}}},
{\hyperref[\detokenize{compute_modify::doc}]{\sphinxcrossref{\DUrole{doc}{compute\_modify}}}}, and {\hyperref[\detokenize{variable::doc}]{\sphinxcrossref{\DUrole{doc}{variable}}}}
commands.

\sphinxAtStartPar
Output options are set by the {\hyperref[\detokenize{thermo::doc}]{\sphinxcrossref{\DUrole{doc}{thermo}}}}, {\hyperref[\detokenize{dump::doc}]{\sphinxcrossref{\DUrole{doc}{dump}}}},
and {\hyperref[\detokenize{restart::doc}]{\sphinxcrossref{\DUrole{doc}{restart}}}} commands.


\subsection{Run a simulation}
\label{\detokenize{Commands_structure:run-a-simulation}}\label{\detokenize{Commands_structure:run}}
\sphinxAtStartPar
A molecular dynamics simulation is run using the {\hyperref[\detokenize{run::doc}]{\sphinxcrossref{\DUrole{doc}{run}}}}
command.  Energy minimization (molecular statics) is performed using
the {\hyperref[\detokenize{minimize::doc}]{\sphinxcrossref{\DUrole{doc}{minimize}}}} command.  A parallel tempering
(replica\sphinxhyphen{}exchange) simulation can be run using the
{\hyperref[\detokenize{temper::doc}]{\sphinxcrossref{\DUrole{doc}{temper}}}} command.

\sphinxstepscope

\clearpage


\section{Commands by category}
\label{\detokenize{Commands_category:commands-by-category}}\label{\detokenize{Commands_category::doc}}
\sphinxAtStartPar
This page lists most of the LAMMPS commands, grouped by category.  The
{\hyperref[\detokenize{Commands_all::doc}]{\sphinxcrossref{\DUrole{doc}{General commands}}}} page lists all general commands
alphabetically.  Style options for entries like fix, compute, pair etc.
have their own pages where they are listed alphabetically.


\subsection{Initialization}
\label{\detokenize{Commands_category:initialization}}

\begin{savenotes}\sphinxattablestart
\sphinxthistablewithglobalstyle
\centering
\begin{tabulary}{\linewidth}[t]{TTTTT}
\sphinxtoprule
\sphinxtableatstartofbodyhook
\sphinxAtStartPar
{\hyperref[\detokenize{newton::doc}]{\sphinxcrossref{\DUrole{doc}{newton}}}}
&
\sphinxAtStartPar
{\hyperref[\detokenize{package::doc}]{\sphinxcrossref{\DUrole{doc}{package}}}}
&
\sphinxAtStartPar
{\hyperref[\detokenize{processors::doc}]{\sphinxcrossref{\DUrole{doc}{processors}}}}
&
\sphinxAtStartPar
{\hyperref[\detokenize{suffix::doc}]{\sphinxcrossref{\DUrole{doc}{suffix}}}}
&
\sphinxAtStartPar
{\hyperref[\detokenize{units::doc}]{\sphinxcrossref{\DUrole{doc}{units}}}}
\\
\sphinxbottomrule
\end{tabulary}
\sphinxtableafterendhook\par
\sphinxattableend\end{savenotes}


\subsection{Setup simulation box}
\label{\detokenize{Commands_category:setup-simulation-box}}

\begin{savenotes}\sphinxattablestart
\sphinxthistablewithglobalstyle
\centering
\begin{tabulary}{\linewidth}[t]{TTTT}
\sphinxtoprule
\sphinxtableatstartofbodyhook
\sphinxAtStartPar
{\hyperref[\detokenize{boundary::doc}]{\sphinxcrossref{\DUrole{doc}{boundary}}}}
&
\sphinxAtStartPar
{\hyperref[\detokenize{change_box::doc}]{\sphinxcrossref{\DUrole{doc}{change\_box}}}}
&
\sphinxAtStartPar
{\hyperref[\detokenize{create_box::doc}]{\sphinxcrossref{\DUrole{doc}{create\_box}}}}
&
\sphinxAtStartPar
{\hyperref[\detokenize{dimension::doc}]{\sphinxcrossref{\DUrole{doc}{dimension}}}}
\\
\sphinxhline
\sphinxAtStartPar
{\hyperref[\detokenize{lattice::doc}]{\sphinxcrossref{\DUrole{doc}{lattice}}}}
&
\sphinxAtStartPar
{\hyperref[\detokenize{region::doc}]{\sphinxcrossref{\DUrole{doc}{region}}}}
&&\\
\sphinxbottomrule
\end{tabulary}
\sphinxtableafterendhook\par
\sphinxattableend\end{savenotes}


\subsection{Setup atoms}
\label{\detokenize{Commands_category:setup-atoms}}

\begin{savenotes}\sphinxattablestart
\sphinxthistablewithglobalstyle
\centering
\begin{tabulary}{\linewidth}[t]{TTTT}
\sphinxtoprule
\sphinxtableatstartofbodyhook
\sphinxAtStartPar
{\hyperref[\detokenize{atom_modify::doc}]{\sphinxcrossref{\DUrole{doc}{atom\_modify}}}}
&
\sphinxAtStartPar
{\hyperref[\detokenize{atom_style::doc}]{\sphinxcrossref{\DUrole{doc}{atom\_style}}}}
&
\sphinxAtStartPar
{\hyperref[\detokenize{balance::doc}]{\sphinxcrossref{\DUrole{doc}{balance}}}}
&
\sphinxAtStartPar
{\hyperref[\detokenize{create_atoms::doc}]{\sphinxcrossref{\DUrole{doc}{create\_atoms}}}}
\\
\sphinxhline
\sphinxAtStartPar
{\hyperref[\detokenize{create_bonds::doc}]{\sphinxcrossref{\DUrole{doc}{create\_bonds}}}}
&
\sphinxAtStartPar
{\hyperref[\detokenize{delete_atoms::doc}]{\sphinxcrossref{\DUrole{doc}{delete\_atoms}}}}
&
\sphinxAtStartPar
{\hyperref[\detokenize{delete_bonds::doc}]{\sphinxcrossref{\DUrole{doc}{delete\_bonds}}}}
&
\sphinxAtStartPar
{\hyperref[\detokenize{displace_atoms::doc}]{\sphinxcrossref{\DUrole{doc}{displace\_atoms}}}}
\\
\sphinxhline
\sphinxAtStartPar
{\hyperref[\detokenize{group::doc}]{\sphinxcrossref{\DUrole{doc}{group}}}}
&
\sphinxAtStartPar
{\hyperref[\detokenize{mass::doc}]{\sphinxcrossref{\DUrole{doc}{mass}}}}
&
\sphinxAtStartPar
{\hyperref[\detokenize{molecule::doc}]{\sphinxcrossref{\DUrole{doc}{molecule}}}}
&
\sphinxAtStartPar
{\hyperref[\detokenize{read_data::doc}]{\sphinxcrossref{\DUrole{doc}{read\_data}}}}
\\
\sphinxhline
\sphinxAtStartPar
{\hyperref[\detokenize{read_dump::doc}]{\sphinxcrossref{\DUrole{doc}{read\_dump}}}}
&
\sphinxAtStartPar
{\hyperref[\detokenize{read_restart::doc}]{\sphinxcrossref{\DUrole{doc}{read\_restart}}}}
&
\sphinxAtStartPar
{\hyperref[\detokenize{replicate::doc}]{\sphinxcrossref{\DUrole{doc}{replicate}}}}
&
\sphinxAtStartPar
{\hyperref[\detokenize{set::doc}]{\sphinxcrossref{\DUrole{doc}{set}}}}
\\
\sphinxhline
\sphinxAtStartPar
{\hyperref[\detokenize{velocity::doc}]{\sphinxcrossref{\DUrole{doc}{velocity}}}}
&&&\\
\sphinxbottomrule
\end{tabulary}
\sphinxtableafterendhook\par
\sphinxattableend\end{savenotes}


\subsection{Force fields}
\label{\detokenize{Commands_category:force-fields}}

\begin{savenotes}\sphinxattablestart
\sphinxthistablewithglobalstyle
\centering
\begin{tabulary}{\linewidth}[t]{TTTT}
\sphinxtoprule
\sphinxtableatstartofbodyhook
\sphinxAtStartPar
{\hyperref[\detokenize{angle_coeff::doc}]{\sphinxcrossref{\DUrole{doc}{angle\_coeff}}}}
&
\sphinxAtStartPar
{\hyperref[\detokenize{angle_style::doc}]{\sphinxcrossref{\DUrole{doc}{angle\_style}}}}
&
\sphinxAtStartPar
{\hyperref[\detokenize{bond_coeff::doc}]{\sphinxcrossref{\DUrole{doc}{bond\_coeff}}}}
&
\sphinxAtStartPar
{\hyperref[\detokenize{bond_style::doc}]{\sphinxcrossref{\DUrole{doc}{bond\_style}}}}
\\
\sphinxhline
\sphinxAtStartPar
{\hyperref[\detokenize{bond_write::doc}]{\sphinxcrossref{\DUrole{doc}{bond\_write}}}}
&
\sphinxAtStartPar
{\hyperref[\detokenize{dielectric::doc}]{\sphinxcrossref{\DUrole{doc}{dielectric}}}}
&
\sphinxAtStartPar
{\hyperref[\detokenize{dihedral_coeff::doc}]{\sphinxcrossref{\DUrole{doc}{dihedral\_coeff}}}}
&
\sphinxAtStartPar
{\hyperref[\detokenize{dihedral_style::doc}]{\sphinxcrossref{\DUrole{doc}{dihedral\_style}}}}
\\
\sphinxhline
\sphinxAtStartPar
{\hyperref[\detokenize{improper_coeff::doc}]{\sphinxcrossref{\DUrole{doc}{improper\_coeff}}}}
&
\sphinxAtStartPar
{\hyperref[\detokenize{improper_style::doc}]{\sphinxcrossref{\DUrole{doc}{improper\_style}}}}
&
\sphinxAtStartPar
{\hyperref[\detokenize{kspace_modify::doc}]{\sphinxcrossref{\DUrole{doc}{kspace\_modify}}}}
&
\sphinxAtStartPar
{\hyperref[\detokenize{kspace_style::doc}]{\sphinxcrossref{\DUrole{doc}{kspace\_style}}}}
\\
\sphinxhline
\sphinxAtStartPar
{\hyperref[\detokenize{pair_coeff::doc}]{\sphinxcrossref{\DUrole{doc}{pair\_coeff}}}}
&
\sphinxAtStartPar
{\hyperref[\detokenize{pair_modify::doc}]{\sphinxcrossref{\DUrole{doc}{pair\_modify}}}}
&
\sphinxAtStartPar
{\hyperref[\detokenize{pair_style::doc}]{\sphinxcrossref{\DUrole{doc}{pair\_style}}}}
&
\sphinxAtStartPar
{\hyperref[\detokenize{pair_write::doc}]{\sphinxcrossref{\DUrole{doc}{pair\_write}}}}
\\
\sphinxhline
\sphinxAtStartPar
{\hyperref[\detokenize{special_bonds::doc}]{\sphinxcrossref{\DUrole{doc}{special\_bonds}}}}
&&&\\
\sphinxbottomrule
\end{tabulary}
\sphinxtableafterendhook\par
\sphinxattableend\end{savenotes}


\subsection{Settings}
\label{\detokenize{Commands_category:settings}}

\begin{savenotes}\sphinxattablestart
\sphinxthistablewithglobalstyle
\centering
\begin{tabulary}{\linewidth}[t]{TTTT}
\sphinxtoprule
\sphinxtableatstartofbodyhook
\sphinxAtStartPar
{\hyperref[\detokenize{comm_modify::doc}]{\sphinxcrossref{\DUrole{doc}{comm\_modify}}}}
&
\sphinxAtStartPar
{\hyperref[\detokenize{comm_style::doc}]{\sphinxcrossref{\DUrole{doc}{comm\_style}}}}
&
\sphinxAtStartPar
{\hyperref[\detokenize{info::doc}]{\sphinxcrossref{\DUrole{doc}{info}}}}
&
\sphinxAtStartPar
{\hyperref[\detokenize{min_modify::doc}]{\sphinxcrossref{\DUrole{doc}{min\_modify}}}}
\\
\sphinxhline
\sphinxAtStartPar
{\hyperref[\detokenize{min_style::doc}]{\sphinxcrossref{\DUrole{doc}{min\_style}}}}
&
\sphinxAtStartPar
{\hyperref[\detokenize{neigh_modify::doc}]{\sphinxcrossref{\DUrole{doc}{neigh\_modify}}}}
&
\sphinxAtStartPar
{\hyperref[\detokenize{neighbor::doc}]{\sphinxcrossref{\DUrole{doc}{neighbor}}}}
&
\sphinxAtStartPar
{\hyperref[\detokenize{partition::doc}]{\sphinxcrossref{\DUrole{doc}{partition}}}}
\\
\sphinxhline
\sphinxAtStartPar
{\hyperref[\detokenize{reset_timestep::doc}]{\sphinxcrossref{\DUrole{doc}{reset\_timestep}}}}
&
\sphinxAtStartPar
{\hyperref[\detokenize{run_style::doc}]{\sphinxcrossref{\DUrole{doc}{run\_style}}}}
&
\sphinxAtStartPar
{\hyperref[\detokenize{timer::doc}]{\sphinxcrossref{\DUrole{doc}{timer}}}}
&
\sphinxAtStartPar
{\hyperref[\detokenize{timestep::doc}]{\sphinxcrossref{\DUrole{doc}{timestep}}}}
\\
\sphinxbottomrule
\end{tabulary}
\sphinxtableafterendhook\par
\sphinxattableend\end{savenotes}


\subsection{Operations within timestepping (fixes) and diagnostics (computes)}
\label{\detokenize{Commands_category:operations-within-timestepping-fixes-and-diagnostics-computes}}

\begin{savenotes}\sphinxattablestart
\sphinxthistablewithglobalstyle
\centering
\begin{tabulary}{\linewidth}[t]{TTTT}
\sphinxtoprule
\sphinxtableatstartofbodyhook
\sphinxAtStartPar
{\hyperref[\detokenize{compute::doc}]{\sphinxcrossref{\DUrole{doc}{compute}}}}
&
\sphinxAtStartPar
{\hyperref[\detokenize{compute_modify::doc}]{\sphinxcrossref{\DUrole{doc}{compute\_modify}}}}
&
\sphinxAtStartPar
{\hyperref[\detokenize{fix::doc}]{\sphinxcrossref{\DUrole{doc}{fix}}}}
&
\sphinxAtStartPar
{\hyperref[\detokenize{fix_modify::doc}]{\sphinxcrossref{\DUrole{doc}{fix\_modify}}}}
\\
\sphinxhline
\sphinxAtStartPar
{\hyperref[\detokenize{uncompute::doc}]{\sphinxcrossref{\DUrole{doc}{uncompute}}}}
&
\sphinxAtStartPar
{\hyperref[\detokenize{unfix::doc}]{\sphinxcrossref{\DUrole{doc}{unfix}}}}
&&\\
\sphinxbottomrule
\end{tabulary}
\sphinxtableafterendhook\par
\sphinxattableend\end{savenotes}


\subsection{Output}
\label{\detokenize{Commands_category:output}}

\begin{savenotes}\sphinxattablestart
\sphinxthistablewithglobalstyle
\centering
\begin{tabulary}{\linewidth}[t]{TTTT}
\sphinxtoprule
\sphinxtableatstartofbodyhook
\sphinxAtStartPar
{\hyperref[\detokenize{dump_image::doc}]{\sphinxcrossref{\DUrole{doc}{dump image}}}}
&
\sphinxAtStartPar
{\hyperref[\detokenize{dump_image::doc}]{\sphinxcrossref{\DUrole{doc}{dump movie}}}}
&
\sphinxAtStartPar
{\hyperref[\detokenize{dump::doc}]{\sphinxcrossref{\DUrole{doc}{dump}}}}
&
\sphinxAtStartPar
{\hyperref[\detokenize{dump_modify::doc}]{\sphinxcrossref{\DUrole{doc}{dump\_modify}}}}
\\
\sphinxhline
\sphinxAtStartPar
{\hyperref[\detokenize{restart::doc}]{\sphinxcrossref{\DUrole{doc}{restart}}}}
&
\sphinxAtStartPar
{\hyperref[\detokenize{thermo::doc}]{\sphinxcrossref{\DUrole{doc}{thermo}}}}
&
\sphinxAtStartPar
{\hyperref[\detokenize{thermo_modify::doc}]{\sphinxcrossref{\DUrole{doc}{thermo\_modify}}}}
&
\sphinxAtStartPar
{\hyperref[\detokenize{thermo_style::doc}]{\sphinxcrossref{\DUrole{doc}{thermo\_style}}}}
\\
\sphinxhline
\sphinxAtStartPar
{\hyperref[\detokenize{undump::doc}]{\sphinxcrossref{\DUrole{doc}{undump}}}}
&
\sphinxAtStartPar
{\hyperref[\detokenize{write_coeff::doc}]{\sphinxcrossref{\DUrole{doc}{write\_coeff}}}}
&
\sphinxAtStartPar
{\hyperref[\detokenize{write_data::doc}]{\sphinxcrossref{\DUrole{doc}{write\_data}}}}
&
\sphinxAtStartPar
{\hyperref[\detokenize{write_dump::doc}]{\sphinxcrossref{\DUrole{doc}{write\_dump}}}}
\\
\sphinxhline
\sphinxAtStartPar
{\hyperref[\detokenize{write_restart::doc}]{\sphinxcrossref{\DUrole{doc}{write\_restart}}}}
&&&\\
\sphinxbottomrule
\end{tabulary}
\sphinxtableafterendhook\par
\sphinxattableend\end{savenotes}


\subsection{Actions}
\label{\detokenize{Commands_category:actions}}

\begin{savenotes}\sphinxattablestart
\sphinxthistablewithglobalstyle
\centering
\begin{tabulary}{\linewidth}[t]{TTTTTT}
\sphinxtoprule
\sphinxtableatstartofbodyhook
\sphinxAtStartPar
{\hyperref[\detokenize{minimize::doc}]{\sphinxcrossref{\DUrole{doc}{minimize}}}}
&
\sphinxAtStartPar
{\hyperref[\detokenize{neb::doc}]{\sphinxcrossref{\DUrole{doc}{neb}}}}
&
\sphinxAtStartPar
{\hyperref[\detokenize{neb_spin::doc}]{\sphinxcrossref{\DUrole{doc}{neb\_spin}}}}
&
\sphinxAtStartPar
{\hyperref[\detokenize{prd::doc}]{\sphinxcrossref{\DUrole{doc}{prd}}}}
&
\sphinxAtStartPar
{\hyperref[\detokenize{rerun::doc}]{\sphinxcrossref{\DUrole{doc}{rerun}}}}
&
\sphinxAtStartPar
{\hyperref[\detokenize{run::doc}]{\sphinxcrossref{\DUrole{doc}{run}}}}
\\
\sphinxhline
\sphinxAtStartPar
{\hyperref[\detokenize{tad::doc}]{\sphinxcrossref{\DUrole{doc}{tad}}}}
&
\sphinxAtStartPar
{\hyperref[\detokenize{temper::doc}]{\sphinxcrossref{\DUrole{doc}{temper}}}}
&&&&\\
\sphinxbottomrule
\end{tabulary}
\sphinxtableafterendhook\par
\sphinxattableend\end{savenotes}


\subsection{Input script control}
\label{\detokenize{Commands_category:input-script-control}}

\begin{savenotes}\sphinxattablestart
\sphinxthistablewithglobalstyle
\centering
\begin{tabulary}{\linewidth}[t]{TTTTTTT}
\sphinxtoprule
\sphinxtableatstartofbodyhook
\sphinxAtStartPar
{\hyperref[\detokenize{clear::doc}]{\sphinxcrossref{\DUrole{doc}{clear}}}}
&
\sphinxAtStartPar
{\hyperref[\detokenize{echo::doc}]{\sphinxcrossref{\DUrole{doc}{echo}}}}
&
\sphinxAtStartPar
{\hyperref[\detokenize{if::doc}]{\sphinxcrossref{\DUrole{doc}{if}}}}
&
\sphinxAtStartPar
{\hyperref[\detokenize{include::doc}]{\sphinxcrossref{\DUrole{doc}{include}}}}
&
\sphinxAtStartPar
{\hyperref[\detokenize{info::doc}]{\sphinxcrossref{\DUrole{doc}{info}}}}
&
\sphinxAtStartPar
{\hyperref[\detokenize{jump::doc}]{\sphinxcrossref{\DUrole{doc}{jump}}}}
&
\sphinxAtStartPar
{\hyperref[\detokenize{label::doc}]{\sphinxcrossref{\DUrole{doc}{label}}}}
\\
\sphinxhline
\sphinxAtStartPar
{\hyperref[\detokenize{log::doc}]{\sphinxcrossref{\DUrole{doc}{log}}}}
&
\sphinxAtStartPar
{\hyperref[\detokenize{next::doc}]{\sphinxcrossref{\DUrole{doc}{next}}}}
&
\sphinxAtStartPar
{\hyperref[\detokenize{print::doc}]{\sphinxcrossref{\DUrole{doc}{print}}}}
&
\sphinxAtStartPar
{\hyperref[\detokenize{python::doc}]{\sphinxcrossref{\DUrole{doc}{python}}}}
&
\sphinxAtStartPar
{\hyperref[\detokenize{quit::doc}]{\sphinxcrossref{\DUrole{doc}{quit}}}}
&
\sphinxAtStartPar
{\hyperref[\detokenize{shell::doc}]{\sphinxcrossref{\DUrole{doc}{shell}}}}
&
\sphinxAtStartPar
{\hyperref[\detokenize{variable::doc}]{\sphinxcrossref{\DUrole{doc}{variable}}}}
\\
\sphinxbottomrule
\end{tabulary}
\sphinxtableafterendhook\par
\sphinxattableend\end{savenotes}

\sphinxstepscope

\clearpage


\section{General commands}
\label{\detokenize{Commands_all:general-commands}}\label{\detokenize{Commands_all::doc}}
\sphinxAtStartPar
An alphabetic list of general LAMMPS commands.


\begin{savenotes}\sphinxattablestart
\sphinxthistablewithglobalstyle
\centering
\begin{tabulary}{\linewidth}[t]{TTTTTT}
\sphinxtoprule
\sphinxtableatstartofbodyhook
\sphinxAtStartPar
{\hyperref[\detokenize{angle_coeff::doc}]{\sphinxcrossref{\DUrole{doc}{angle\_coeff}}}}
&
\sphinxAtStartPar
{\hyperref[\detokenize{angle_style::doc}]{\sphinxcrossref{\DUrole{doc}{angle\_style}}}}
&
\sphinxAtStartPar
{\hyperref[\detokenize{angle_write::doc}]{\sphinxcrossref{\DUrole{doc}{angle\_write}}}}
&
\sphinxAtStartPar
{\hyperref[\detokenize{atom_modify::doc}]{\sphinxcrossref{\DUrole{doc}{atom\_modify}}}}
&
\sphinxAtStartPar
{\hyperref[\detokenize{atom_style::doc}]{\sphinxcrossref{\DUrole{doc}{atom\_style}}}}
&
\sphinxAtStartPar
{\hyperref[\detokenize{balance::doc}]{\sphinxcrossref{\DUrole{doc}{balance}}}}
\\
\sphinxhline
\sphinxAtStartPar
{\hyperref[\detokenize{bond_coeff::doc}]{\sphinxcrossref{\DUrole{doc}{bond\_coeff}}}}
&
\sphinxAtStartPar
{\hyperref[\detokenize{bond_style::doc}]{\sphinxcrossref{\DUrole{doc}{bond\_style}}}}
&
\sphinxAtStartPar
{\hyperref[\detokenize{bond_write::doc}]{\sphinxcrossref{\DUrole{doc}{bond\_write}}}}
&
\sphinxAtStartPar
{\hyperref[\detokenize{boundary::doc}]{\sphinxcrossref{\DUrole{doc}{boundary}}}}
&
\sphinxAtStartPar
{\hyperref[\detokenize{change_box::doc}]{\sphinxcrossref{\DUrole{doc}{change\_box}}}}
&
\sphinxAtStartPar
{\hyperref[\detokenize{clear::doc}]{\sphinxcrossref{\DUrole{doc}{clear}}}}
\\
\sphinxhline
\sphinxAtStartPar
{\hyperref[\detokenize{comm_modify::doc}]{\sphinxcrossref{\DUrole{doc}{comm\_modify}}}}
&
\sphinxAtStartPar
{\hyperref[\detokenize{comm_style::doc}]{\sphinxcrossref{\DUrole{doc}{comm\_style}}}}
&
\sphinxAtStartPar
{\hyperref[\detokenize{compute::doc}]{\sphinxcrossref{\DUrole{doc}{compute}}}}
&
\sphinxAtStartPar
{\hyperref[\detokenize{compute_modify::doc}]{\sphinxcrossref{\DUrole{doc}{compute\_modify}}}}
&
\sphinxAtStartPar
{\hyperref[\detokenize{create_atoms::doc}]{\sphinxcrossref{\DUrole{doc}{create\_atoms}}}}
&
\sphinxAtStartPar
{\hyperref[\detokenize{create_bonds::doc}]{\sphinxcrossref{\DUrole{doc}{create\_bonds}}}}
\\
\sphinxhline
\sphinxAtStartPar
{\hyperref[\detokenize{create_box::doc}]{\sphinxcrossref{\DUrole{doc}{create\_box}}}}
&
\sphinxAtStartPar
{\hyperref[\detokenize{delete_atoms::doc}]{\sphinxcrossref{\DUrole{doc}{delete\_atoms}}}}
&
\sphinxAtStartPar
{\hyperref[\detokenize{delete_bonds::doc}]{\sphinxcrossref{\DUrole{doc}{delete\_bonds}}}}
&
\sphinxAtStartPar
{\hyperref[\detokenize{dielectric::doc}]{\sphinxcrossref{\DUrole{doc}{dielectric}}}}
&
\sphinxAtStartPar
{\hyperref[\detokenize{dihedral_coeff::doc}]{\sphinxcrossref{\DUrole{doc}{dihedral\_coeff}}}}
&
\sphinxAtStartPar
{\hyperref[\detokenize{dihedral_style::doc}]{\sphinxcrossref{\DUrole{doc}{dihedral\_style}}}}
\\
\sphinxhline
\sphinxAtStartPar
{\hyperref[\detokenize{dihedral_write::doc}]{\sphinxcrossref{\DUrole{doc}{dihedral\_write}}}}
&
\sphinxAtStartPar
{\hyperref[\detokenize{dimension::doc}]{\sphinxcrossref{\DUrole{doc}{dimension}}}}
&
\sphinxAtStartPar
{\hyperref[\detokenize{displace_atoms::doc}]{\sphinxcrossref{\DUrole{doc}{displace\_atoms}}}}
&
\sphinxAtStartPar
{\hyperref[\detokenize{dump::doc}]{\sphinxcrossref{\DUrole{doc}{dump}}}}
&
\sphinxAtStartPar
{\hyperref[\detokenize{dump_modify::doc}]{\sphinxcrossref{\DUrole{doc}{dump\_modify}}}}
&
\sphinxAtStartPar
{\hyperref[\detokenize{echo::doc}]{\sphinxcrossref{\DUrole{doc}{echo}}}}
\\
\sphinxhline
\sphinxAtStartPar
{\hyperref[\detokenize{fix::doc}]{\sphinxcrossref{\DUrole{doc}{fix}}}}
&
\sphinxAtStartPar
{\hyperref[\detokenize{fix_modify::doc}]{\sphinxcrossref{\DUrole{doc}{fix\_modify}}}}
&
\sphinxAtStartPar
{\hyperref[\detokenize{geturl::doc}]{\sphinxcrossref{\DUrole{doc}{geturl}}}}
&
\sphinxAtStartPar
{\hyperref[\detokenize{group::doc}]{\sphinxcrossref{\DUrole{doc}{group}}}}
&
\sphinxAtStartPar
{\hyperref[\detokenize{if::doc}]{\sphinxcrossref{\DUrole{doc}{if}}}}
&
\sphinxAtStartPar
{\hyperref[\detokenize{improper_coeff::doc}]{\sphinxcrossref{\DUrole{doc}{improper\_coeff}}}}
\\
\sphinxhline
\sphinxAtStartPar
{\hyperref[\detokenize{improper_style::doc}]{\sphinxcrossref{\DUrole{doc}{improper\_style}}}}
&
\sphinxAtStartPar
{\hyperref[\detokenize{include::doc}]{\sphinxcrossref{\DUrole{doc}{include}}}}
&
\sphinxAtStartPar
{\hyperref[\detokenize{info::doc}]{\sphinxcrossref{\DUrole{doc}{info}}}}
&
\sphinxAtStartPar
{\hyperref[\detokenize{jump::doc}]{\sphinxcrossref{\DUrole{doc}{jump}}}}
&
\sphinxAtStartPar
{\hyperref[\detokenize{kspace_modify::doc}]{\sphinxcrossref{\DUrole{doc}{kspace\_modify}}}}
&
\sphinxAtStartPar
{\hyperref[\detokenize{kspace_style::doc}]{\sphinxcrossref{\DUrole{doc}{kspace\_style}}}}
\\
\sphinxhline
\sphinxAtStartPar
{\hyperref[\detokenize{label::doc}]{\sphinxcrossref{\DUrole{doc}{label}}}}
&
\sphinxAtStartPar
{\hyperref[\detokenize{labelmap::doc}]{\sphinxcrossref{\DUrole{doc}{labelmap}}}}
&
\sphinxAtStartPar
{\hyperref[\detokenize{lattice::doc}]{\sphinxcrossref{\DUrole{doc}{lattice}}}}
&
\sphinxAtStartPar
{\hyperref[\detokenize{log::doc}]{\sphinxcrossref{\DUrole{doc}{log}}}}
&
\sphinxAtStartPar
{\hyperref[\detokenize{mass::doc}]{\sphinxcrossref{\DUrole{doc}{mass}}}}
&
\sphinxAtStartPar
{\hyperref[\detokenize{minimize::doc}]{\sphinxcrossref{\DUrole{doc}{minimize}}}}
\\
\sphinxhline
\sphinxAtStartPar
{\hyperref[\detokenize{min_modify::doc}]{\sphinxcrossref{\DUrole{doc}{min\_modify}}}}
&
\sphinxAtStartPar
{\hyperref[\detokenize{min_style::doc}]{\sphinxcrossref{\DUrole{doc}{min\_style}}}}
&
\sphinxAtStartPar
{\hyperref[\detokenize{molecule::doc}]{\sphinxcrossref{\DUrole{doc}{molecule}}}}
&
\sphinxAtStartPar
{\hyperref[\detokenize{neigh_modify::doc}]{\sphinxcrossref{\DUrole{doc}{neigh\_modify}}}}
&
\sphinxAtStartPar
{\hyperref[\detokenize{neighbor::doc}]{\sphinxcrossref{\DUrole{doc}{neighbor}}}}
&
\sphinxAtStartPar
{\hyperref[\detokenize{newton::doc}]{\sphinxcrossref{\DUrole{doc}{newton}}}}
\\
\sphinxhline
\sphinxAtStartPar
{\hyperref[\detokenize{next::doc}]{\sphinxcrossref{\DUrole{doc}{next}}}}
&
\sphinxAtStartPar
{\hyperref[\detokenize{package::doc}]{\sphinxcrossref{\DUrole{doc}{package}}}}
&
\sphinxAtStartPar
{\hyperref[\detokenize{pair_coeff::doc}]{\sphinxcrossref{\DUrole{doc}{pair\_coeff}}}}
&
\sphinxAtStartPar
{\hyperref[\detokenize{pair_modify::doc}]{\sphinxcrossref{\DUrole{doc}{pair\_modify}}}}
&
\sphinxAtStartPar
{\hyperref[\detokenize{pair_style::doc}]{\sphinxcrossref{\DUrole{doc}{pair\_style}}}}
&
\sphinxAtStartPar
{\hyperref[\detokenize{pair_write::doc}]{\sphinxcrossref{\DUrole{doc}{pair\_write}}}}
\\
\sphinxhline
\sphinxAtStartPar
{\hyperref[\detokenize{partition::doc}]{\sphinxcrossref{\DUrole{doc}{partition}}}}
&
\sphinxAtStartPar
{\hyperref[\detokenize{print::doc}]{\sphinxcrossref{\DUrole{doc}{print}}}}
&
\sphinxAtStartPar
{\hyperref[\detokenize{processors::doc}]{\sphinxcrossref{\DUrole{doc}{processors}}}}
&
\sphinxAtStartPar
{\hyperref[\detokenize{quit::doc}]{\sphinxcrossref{\DUrole{doc}{quit}}}}
&
\sphinxAtStartPar
{\hyperref[\detokenize{read_data::doc}]{\sphinxcrossref{\DUrole{doc}{read\_data}}}}
&
\sphinxAtStartPar
{\hyperref[\detokenize{read_dump::doc}]{\sphinxcrossref{\DUrole{doc}{read\_dump}}}}
\\
\sphinxhline
\sphinxAtStartPar
{\hyperref[\detokenize{read_restart::doc}]{\sphinxcrossref{\DUrole{doc}{read\_restart}}}}
&
\sphinxAtStartPar
{\hyperref[\detokenize{region::doc}]{\sphinxcrossref{\DUrole{doc}{region}}}}
&
\sphinxAtStartPar
{\hyperref[\detokenize{replicate::doc}]{\sphinxcrossref{\DUrole{doc}{replicate}}}}
&
\sphinxAtStartPar
{\hyperref[\detokenize{rerun::doc}]{\sphinxcrossref{\DUrole{doc}{rerun}}}}
&
\sphinxAtStartPar
{\hyperref[\detokenize{reset_atoms::doc}]{\sphinxcrossref{\DUrole{doc}{reset\_atoms}}}}
&
\sphinxAtStartPar
{\hyperref[\detokenize{reset_timestep::doc}]{\sphinxcrossref{\DUrole{doc}{reset\_timestep}}}}
\\
\sphinxhline
\sphinxAtStartPar
{\hyperref[\detokenize{restart::doc}]{\sphinxcrossref{\DUrole{doc}{restart}}}}
&
\sphinxAtStartPar
{\hyperref[\detokenize{run::doc}]{\sphinxcrossref{\DUrole{doc}{run}}}}
&
\sphinxAtStartPar
{\hyperref[\detokenize{run_style::doc}]{\sphinxcrossref{\DUrole{doc}{run\_style}}}}
&
\sphinxAtStartPar
{\hyperref[\detokenize{set::doc}]{\sphinxcrossref{\DUrole{doc}{set}}}}
&
\sphinxAtStartPar
{\hyperref[\detokenize{shell::doc}]{\sphinxcrossref{\DUrole{doc}{shell}}}}
&
\sphinxAtStartPar
{\hyperref[\detokenize{special_bonds::doc}]{\sphinxcrossref{\DUrole{doc}{special\_bonds}}}}
\\
\sphinxhline
\sphinxAtStartPar
{\hyperref[\detokenize{suffix::doc}]{\sphinxcrossref{\DUrole{doc}{suffix}}}}
&
\sphinxAtStartPar
{\hyperref[\detokenize{thermo::doc}]{\sphinxcrossref{\DUrole{doc}{thermo}}}}
&
\sphinxAtStartPar
{\hyperref[\detokenize{thermo_modify::doc}]{\sphinxcrossref{\DUrole{doc}{thermo\_modify}}}}
&
\sphinxAtStartPar
{\hyperref[\detokenize{thermo_style::doc}]{\sphinxcrossref{\DUrole{doc}{thermo\_style}}}}
&
\sphinxAtStartPar
{\hyperref[\detokenize{timer::doc}]{\sphinxcrossref{\DUrole{doc}{timer}}}}
&
\sphinxAtStartPar
{\hyperref[\detokenize{timestep::doc}]{\sphinxcrossref{\DUrole{doc}{timestep}}}}
\\
\sphinxhline
\sphinxAtStartPar
{\hyperref[\detokenize{uncompute::doc}]{\sphinxcrossref{\DUrole{doc}{uncompute}}}}
&
\sphinxAtStartPar
{\hyperref[\detokenize{undump::doc}]{\sphinxcrossref{\DUrole{doc}{undump}}}}
&
\sphinxAtStartPar
{\hyperref[\detokenize{unfix::doc}]{\sphinxcrossref{\DUrole{doc}{unfix}}}}
&
\sphinxAtStartPar
{\hyperref[\detokenize{units::doc}]{\sphinxcrossref{\DUrole{doc}{units}}}}
&
\sphinxAtStartPar
{\hyperref[\detokenize{variable::doc}]{\sphinxcrossref{\DUrole{doc}{variable}}}}
&
\sphinxAtStartPar
{\hyperref[\detokenize{velocity::doc}]{\sphinxcrossref{\DUrole{doc}{velocity}}}}
\\
\sphinxhline
\sphinxAtStartPar
{\hyperref[\detokenize{write_coeff::doc}]{\sphinxcrossref{\DUrole{doc}{write\_coeff}}}}
&
\sphinxAtStartPar
{\hyperref[\detokenize{write_data::doc}]{\sphinxcrossref{\DUrole{doc}{write\_data}}}}
&
\sphinxAtStartPar
{\hyperref[\detokenize{write_dump::doc}]{\sphinxcrossref{\DUrole{doc}{write\_dump}}}}
&
\sphinxAtStartPar
{\hyperref[\detokenize{write_restart::doc}]{\sphinxcrossref{\DUrole{doc}{write\_restart}}}}
&&\\
\sphinxbottomrule
\end{tabulary}
\sphinxtableafterendhook\par
\sphinxattableend\end{savenotes}

\sphinxAtStartPar
Additional general LAMMPS commands provided by packages.  A few
commands have accelerated versions.  This is indicated by an
additional letter in parenthesis: k = KOKKOS.


\begin{savenotes}\sphinxattablestart
\sphinxthistablewithglobalstyle
\centering
\begin{tabulary}{\linewidth}[t]{TTTTTT}
\sphinxtoprule
\sphinxtableatstartofbodyhook
\sphinxAtStartPar
{\hyperref[\detokenize{dynamical_matrix::doc}]{\sphinxcrossref{\DUrole{doc}{dynamical\_matrix (k)}}}}
&
\sphinxAtStartPar
{\hyperref[\detokenize{group2ndx::doc}]{\sphinxcrossref{\DUrole{doc}{group2ndx}}}}
&
\sphinxAtStartPar
{\hyperref[\detokenize{hyper::doc}]{\sphinxcrossref{\DUrole{doc}{hyper}}}}
&
\sphinxAtStartPar
{\hyperref[\detokenize{kim_commands::doc}]{\sphinxcrossref{\DUrole{doc}{kim}}}}
&
\sphinxAtStartPar
{\hyperref[\detokenize{fitpod_command::doc}]{\sphinxcrossref{\DUrole{doc}{fitpod}}}}
&
\sphinxAtStartPar
{\hyperref[\detokenize{mdi::doc}]{\sphinxcrossref{\DUrole{doc}{mdi}}}}
\\
\sphinxhline
\sphinxAtStartPar
{\hyperref[\detokenize{group2ndx::doc}]{\sphinxcrossref{\DUrole{doc}{ndx2group}}}}
&
\sphinxAtStartPar
{\hyperref[\detokenize{neb::doc}]{\sphinxcrossref{\DUrole{doc}{neb}}}}
&
\sphinxAtStartPar
{\hyperref[\detokenize{neb_spin::doc}]{\sphinxcrossref{\DUrole{doc}{neb/spin}}}}
&
\sphinxAtStartPar
{\hyperref[\detokenize{plugin::doc}]{\sphinxcrossref{\DUrole{doc}{plugin}}}}
&
\sphinxAtStartPar
{\hyperref[\detokenize{prd::doc}]{\sphinxcrossref{\DUrole{doc}{prd}}}}
&
\sphinxAtStartPar
{\hyperref[\detokenize{python::doc}]{\sphinxcrossref{\DUrole{doc}{python}}}}
\\
\sphinxhline
\sphinxAtStartPar
{\hyperref[\detokenize{tad::doc}]{\sphinxcrossref{\DUrole{doc}{tad}}}}
&
\sphinxAtStartPar
{\hyperref[\detokenize{temper::doc}]{\sphinxcrossref{\DUrole{doc}{temper}}}}
&
\sphinxAtStartPar
{\hyperref[\detokenize{temper_grem::doc}]{\sphinxcrossref{\DUrole{doc}{temper/grem}}}}
&
\sphinxAtStartPar
{\hyperref[\detokenize{temper_npt::doc}]{\sphinxcrossref{\DUrole{doc}{temper/npt}}}}
&
\sphinxAtStartPar
{\hyperref[\detokenize{third_order::doc}]{\sphinxcrossref{\DUrole{doc}{third\_order (k)}}}}
&\\
\sphinxbottomrule
\end{tabulary}
\sphinxtableafterendhook\par
\sphinxattableend\end{savenotes}

\sphinxstepscope


\section{Fix styles}
\label{\detokenize{Commands_fix:fix-styles}}\label{\detokenize{Commands_fix::doc}}
\sphinxAtStartPar
An alphabetic list of all LAMMPS {\hyperref[\detokenize{fix::doc}]{\sphinxcrossref{\DUrole{doc}{fix}}}} commands.  Some styles
have accelerated versions.  This is indicated by additional letters in
parenthesis: g = GPU, i = INTEL, k = KOKKOS, o = OPENMP, t =
OPT.


\begin{savenotes}
\sphinxatlongtablestart
\sphinxthistablewithglobalstyle
\makeatletter
  \LTleft \@totalleftmargin plus1fill
  \LTright\dimexpr\columnwidth-\@totalleftmargin-\linewidth\relax plus1fill
\makeatother
\begin{longtable}{llll}
\sphinxtoprule
\endfirsthead

\multicolumn{4}{c}{\sphinxnorowcolor
    \makebox[0pt]{\sphinxtablecontinued{\tablename\ \thetable{} \textendash{} continued from previous page}}%
}\\
\sphinxtoprule
\endhead

\sphinxbottomrule
\multicolumn{4}{r}{\sphinxnorowcolor
    \makebox[0pt][r]{\sphinxtablecontinued{continues on next page}}%
}\\
\endfoot

\endlastfoot
\sphinxtableatstartofbodyhook

\sphinxAtStartPar
{\hyperref[\detokenize{fix_accelerate_cos::doc}]{\sphinxcrossref{\DUrole{doc}{accelerate/cos}}}}
&
\sphinxAtStartPar
{\hyperref[\detokenize{fix_acks2_reaxff::doc}]{\sphinxcrossref{\DUrole{doc}{acks2/reaxff (k)}}}}
&
\sphinxAtStartPar
{\hyperref[\detokenize{fix_adapt::doc}]{\sphinxcrossref{\DUrole{doc}{adapt}}}}
&
\sphinxAtStartPar
{\hyperref[\detokenize{fix_adapt_fep::doc}]{\sphinxcrossref{\DUrole{doc}{adapt/fep}}}}
\\
\sphinxhline
\sphinxAtStartPar
{\hyperref[\detokenize{fix_addforce::doc}]{\sphinxcrossref{\DUrole{doc}{addforce}}}}
&
\sphinxAtStartPar
{\hyperref[\detokenize{fix_add_heat::doc}]{\sphinxcrossref{\DUrole{doc}{add/heat}}}}
&
\sphinxAtStartPar
{\hyperref[\detokenize{fix_addtorque::doc}]{\sphinxcrossref{\DUrole{doc}{addtorque}}}}
&
\sphinxAtStartPar
{\hyperref[\detokenize{fix_alchemy::doc}]{\sphinxcrossref{\DUrole{doc}{alchemy}}}}
\\
\sphinxhline
\sphinxAtStartPar
{\hyperref[\detokenize{fix_amoeba_bitorsion::doc}]{\sphinxcrossref{\DUrole{doc}{amoeba/bitorsion}}}}
&
\sphinxAtStartPar
{\hyperref[\detokenize{fix_amoeba_pitorsion::doc}]{\sphinxcrossref{\DUrole{doc}{amoeba/pitorsion}}}}
&
\sphinxAtStartPar
{\hyperref[\detokenize{fix_append_atoms::doc}]{\sphinxcrossref{\DUrole{doc}{append/atoms}}}}
&
\sphinxAtStartPar
{\hyperref[\detokenize{fix_atc::doc}]{\sphinxcrossref{\DUrole{doc}{atc}}}}
\\
\sphinxhline
\sphinxAtStartPar
{\hyperref[\detokenize{fix_atom_swap::doc}]{\sphinxcrossref{\DUrole{doc}{atom/swap}}}}
&
\sphinxAtStartPar
{\hyperref[\detokenize{fix_ave_atom::doc}]{\sphinxcrossref{\DUrole{doc}{ave/atom}}}}
&
\sphinxAtStartPar
{\hyperref[\detokenize{fix_ave_chunk::doc}]{\sphinxcrossref{\DUrole{doc}{ave/chunk}}}}
&
\sphinxAtStartPar
{\hyperref[\detokenize{fix_ave_correlate::doc}]{\sphinxcrossref{\DUrole{doc}{ave/correlate}}}}
\\
\sphinxhline
\sphinxAtStartPar
{\hyperref[\detokenize{fix_ave_correlate_long::doc}]{\sphinxcrossref{\DUrole{doc}{ave/correlate/long}}}}
&
\sphinxAtStartPar
{\hyperref[\detokenize{fix_ave_grid::doc}]{\sphinxcrossref{\DUrole{doc}{ave/grid}}}}
&
\sphinxAtStartPar
{\hyperref[\detokenize{fix_ave_histo::doc}]{\sphinxcrossref{\DUrole{doc}{ave/histo}}}}
&
\sphinxAtStartPar
{\hyperref[\detokenize{fix_ave_histo::doc}]{\sphinxcrossref{\DUrole{doc}{ave/histo/weight}}}}
\\
\sphinxhline
\sphinxAtStartPar
{\hyperref[\detokenize{fix_ave_time::doc}]{\sphinxcrossref{\DUrole{doc}{ave/time}}}}
&
\sphinxAtStartPar
{\hyperref[\detokenize{fix_aveforce::doc}]{\sphinxcrossref{\DUrole{doc}{aveforce}}}}
&
\sphinxAtStartPar
{\hyperref[\detokenize{fix_balance::doc}]{\sphinxcrossref{\DUrole{doc}{balance}}}}
&
\sphinxAtStartPar
{\hyperref[\detokenize{fix_bocs::doc}]{\sphinxcrossref{\DUrole{doc}{bocs}}}}
\\
\sphinxhline
\sphinxAtStartPar
{\hyperref[\detokenize{fix_bond_break::doc}]{\sphinxcrossref{\DUrole{doc}{bond/break}}}}
&
\sphinxAtStartPar
{\hyperref[\detokenize{fix_bond_create::doc}]{\sphinxcrossref{\DUrole{doc}{bond/create}}}}
&
\sphinxAtStartPar
{\hyperref[\detokenize{fix_bond_create::doc}]{\sphinxcrossref{\DUrole{doc}{bond/create/angle}}}}
&
\sphinxAtStartPar
{\hyperref[\detokenize{fix_bond_react::doc}]{\sphinxcrossref{\DUrole{doc}{bond/react}}}}
\\
\sphinxhline
\sphinxAtStartPar
{\hyperref[\detokenize{fix_bond_swap::doc}]{\sphinxcrossref{\DUrole{doc}{bond/swap}}}}
&
\sphinxAtStartPar
{\hyperref[\detokenize{fix_box_relax::doc}]{\sphinxcrossref{\DUrole{doc}{box/relax}}}}
&
\sphinxAtStartPar
{\hyperref[\detokenize{fix_brownian::doc}]{\sphinxcrossref{\DUrole{doc}{brownian}}}}
&
\sphinxAtStartPar
{\hyperref[\detokenize{fix_brownian::doc}]{\sphinxcrossref{\DUrole{doc}{brownian/asphere}}}}
\\
\sphinxhline
\sphinxAtStartPar
{\hyperref[\detokenize{fix_brownian::doc}]{\sphinxcrossref{\DUrole{doc}{brownian/sphere}}}}
&
\sphinxAtStartPar
{\hyperref[\detokenize{fix_charge_regulation::doc}]{\sphinxcrossref{\DUrole{doc}{charge/regulation}}}}
&
\sphinxAtStartPar
{\hyperref[\detokenize{fix_cmap::doc}]{\sphinxcrossref{\DUrole{doc}{cmap}}}}
&
\sphinxAtStartPar
{\hyperref[\detokenize{fix_colvars::doc}]{\sphinxcrossref{\DUrole{doc}{colvars}}}}
\\
\sphinxhline
\sphinxAtStartPar
{\hyperref[\detokenize{fix_controller::doc}]{\sphinxcrossref{\DUrole{doc}{controller}}}}
&
\sphinxAtStartPar
{\hyperref[\detokenize{fix_damping_cundall::doc}]{\sphinxcrossref{\DUrole{doc}{damping/cundall}}}}
&
\sphinxAtStartPar
{\hyperref[\detokenize{fix_deform::doc}]{\sphinxcrossref{\DUrole{doc}{deform (k)}}}}
&
\sphinxAtStartPar
{\hyperref[\detokenize{fix_deform_pressure::doc}]{\sphinxcrossref{\DUrole{doc}{deform/pressure}}}}
\\
\sphinxhline
\sphinxAtStartPar
{\hyperref[\detokenize{fix_deposit::doc}]{\sphinxcrossref{\DUrole{doc}{deposit}}}}
&
\sphinxAtStartPar
{\hyperref[\detokenize{fix_dpd_energy::doc}]{\sphinxcrossref{\DUrole{doc}{dpd/energy (k)}}}}
&
\sphinxAtStartPar
{\hyperref[\detokenize{fix_drag::doc}]{\sphinxcrossref{\DUrole{doc}{drag}}}}
&
\sphinxAtStartPar
{\hyperref[\detokenize{fix_drude::doc}]{\sphinxcrossref{\DUrole{doc}{drude}}}}
\\
\sphinxhline
\sphinxAtStartPar
{\hyperref[\detokenize{fix_drude_transform::doc}]{\sphinxcrossref{\DUrole{doc}{drude/transform/direct}}}}
&
\sphinxAtStartPar
{\hyperref[\detokenize{fix_drude_transform::doc}]{\sphinxcrossref{\DUrole{doc}{drude/transform/inverse}}}}
&
\sphinxAtStartPar
{\hyperref[\detokenize{fix_dt_reset::doc}]{\sphinxcrossref{\DUrole{doc}{dt/reset (k)}}}}
&
\sphinxAtStartPar
{\hyperref[\detokenize{fix_dpd_source::doc}]{\sphinxcrossref{\DUrole{doc}{edpd/source}}}}
\\
\sphinxhline
\sphinxAtStartPar
{\hyperref[\detokenize{fix_efield::doc}]{\sphinxcrossref{\DUrole{doc}{efield (k)}}}}
&
\sphinxAtStartPar
{\hyperref[\detokenize{fix_efield::doc}]{\sphinxcrossref{\DUrole{doc}{efield/tip4p}}}}
&
\sphinxAtStartPar
{\hyperref[\detokenize{fix_ehex::doc}]{\sphinxcrossref{\DUrole{doc}{ehex}}}}
&
\sphinxAtStartPar
{\hyperref[\detokenize{fix_electrode::doc}]{\sphinxcrossref{\DUrole{doc}{electrode/conp (i)}}}}
\\
\sphinxhline
\sphinxAtStartPar
{\hyperref[\detokenize{fix_electrode::doc}]{\sphinxcrossref{\DUrole{doc}{electrode/conq (i)}}}}
&
\sphinxAtStartPar
{\hyperref[\detokenize{fix_electrode::doc}]{\sphinxcrossref{\DUrole{doc}{electrode/thermo (i)}}}}
&
\sphinxAtStartPar
{\hyperref[\detokenize{fix_electron_stopping::doc}]{\sphinxcrossref{\DUrole{doc}{electron/stopping}}}}
&
\sphinxAtStartPar
{\hyperref[\detokenize{fix_electron_stopping::doc}]{\sphinxcrossref{\DUrole{doc}{electron/stopping/fit}}}}
\\
\sphinxhline
\sphinxAtStartPar
{\hyperref[\detokenize{fix_enforce2d::doc}]{\sphinxcrossref{\DUrole{doc}{enforce2d (k)}}}}
&
\sphinxAtStartPar
{\hyperref[\detokenize{fix_eos_cv::doc}]{\sphinxcrossref{\DUrole{doc}{eos/cv}}}}
&
\sphinxAtStartPar
{\hyperref[\detokenize{fix_eos_table::doc}]{\sphinxcrossref{\DUrole{doc}{eos/table}}}}
&
\sphinxAtStartPar
{\hyperref[\detokenize{fix_eos_table_rx::doc}]{\sphinxcrossref{\DUrole{doc}{eos/table/rx (k)}}}}
\\
\sphinxhline
\sphinxAtStartPar
{\hyperref[\detokenize{fix_evaporate::doc}]{\sphinxcrossref{\DUrole{doc}{evaporate}}}}
&
\sphinxAtStartPar
{\hyperref[\detokenize{fix_external::doc}]{\sphinxcrossref{\DUrole{doc}{external}}}}
&
\sphinxAtStartPar
{\hyperref[\detokenize{fix_ffl::doc}]{\sphinxcrossref{\DUrole{doc}{ffl}}}}
&
\sphinxAtStartPar
{\hyperref[\detokenize{fix_filter_corotate::doc}]{\sphinxcrossref{\DUrole{doc}{filter/corotate}}}}
\\
\sphinxhline
\sphinxAtStartPar
{\hyperref[\detokenize{fix_flow_gauss::doc}]{\sphinxcrossref{\DUrole{doc}{flow/gauss}}}}
&
\sphinxAtStartPar
{\hyperref[\detokenize{fix_freeze::doc}]{\sphinxcrossref{\DUrole{doc}{freeze (k)}}}}
&
\sphinxAtStartPar
{\hyperref[\detokenize{fix_gcmc::doc}]{\sphinxcrossref{\DUrole{doc}{gcmc}}}}
&
\sphinxAtStartPar
{\hyperref[\detokenize{fix_gld::doc}]{\sphinxcrossref{\DUrole{doc}{gld}}}}
\\
\sphinxhline
\sphinxAtStartPar
{\hyperref[\detokenize{fix_gle::doc}]{\sphinxcrossref{\DUrole{doc}{gle}}}}
&
\sphinxAtStartPar
{\hyperref[\detokenize{fix_gravity::doc}]{\sphinxcrossref{\DUrole{doc}{gravity (ko)}}}}
&
\sphinxAtStartPar
{\hyperref[\detokenize{fix_grem::doc}]{\sphinxcrossref{\DUrole{doc}{grem}}}}
&
\sphinxAtStartPar
{\hyperref[\detokenize{fix_halt::doc}]{\sphinxcrossref{\DUrole{doc}{halt}}}}
\\
\sphinxhline
\sphinxAtStartPar
{\hyperref[\detokenize{fix_heat::doc}]{\sphinxcrossref{\DUrole{doc}{heat}}}}
&
\sphinxAtStartPar
{\hyperref[\detokenize{fix_heat_flow::doc}]{\sphinxcrossref{\DUrole{doc}{heat/flow}}}}
&
\sphinxAtStartPar
{\hyperref[\detokenize{fix_hyper_global::doc}]{\sphinxcrossref{\DUrole{doc}{hyper/global}}}}
&
\sphinxAtStartPar
{\hyperref[\detokenize{fix_hyper_local::doc}]{\sphinxcrossref{\DUrole{doc}{hyper/local}}}}
\\
\sphinxhline
\sphinxAtStartPar
{\hyperref[\detokenize{fix_imd::doc}]{\sphinxcrossref{\DUrole{doc}{imd}}}}
&
\sphinxAtStartPar
{\hyperref[\detokenize{fix_indent::doc}]{\sphinxcrossref{\DUrole{doc}{indent}}}}
&
\sphinxAtStartPar
{\hyperref[\detokenize{fix_ipi::doc}]{\sphinxcrossref{\DUrole{doc}{ipi}}}}
&
\sphinxAtStartPar
{\hyperref[\detokenize{fix_langevin::doc}]{\sphinxcrossref{\DUrole{doc}{langevin (k)}}}}
\\
\sphinxhline
\sphinxAtStartPar
{\hyperref[\detokenize{fix_langevin_drude::doc}]{\sphinxcrossref{\DUrole{doc}{langevin/drude}}}}
&
\sphinxAtStartPar
{\hyperref[\detokenize{fix_langevin_eff::doc}]{\sphinxcrossref{\DUrole{doc}{langevin/eff}}}}
&
\sphinxAtStartPar
{\hyperref[\detokenize{fix_langevin_spin::doc}]{\sphinxcrossref{\DUrole{doc}{langevin/spin}}}}
&
\sphinxAtStartPar
{\hyperref[\detokenize{fix_lb_fluid::doc}]{\sphinxcrossref{\DUrole{doc}{lb/fluid}}}}
\\
\sphinxhline
\sphinxAtStartPar
{\hyperref[\detokenize{fix_lb_momentum::doc}]{\sphinxcrossref{\DUrole{doc}{lb/momentum}}}}
&
\sphinxAtStartPar
{\hyperref[\detokenize{fix_lb_viscous::doc}]{\sphinxcrossref{\DUrole{doc}{lb/viscous}}}}
&
\sphinxAtStartPar
{\hyperref[\detokenize{fix_lineforce::doc}]{\sphinxcrossref{\DUrole{doc}{lineforce}}}}
&
\sphinxAtStartPar
{\hyperref[\detokenize{fix_manifoldforce::doc}]{\sphinxcrossref{\DUrole{doc}{manifoldforce}}}}
\\
\sphinxhline
\sphinxAtStartPar
{\hyperref[\detokenize{fix_mdi_qm::doc}]{\sphinxcrossref{\DUrole{doc}{mdi/qm}}}}
&
\sphinxAtStartPar
{\hyperref[\detokenize{fix_mdi_qmmm::doc}]{\sphinxcrossref{\DUrole{doc}{mdi/qmmm}}}}
&
\sphinxAtStartPar
{\hyperref[\detokenize{fix_meso_move::doc}]{\sphinxcrossref{\DUrole{doc}{meso/move}}}}
&
\sphinxAtStartPar
{\hyperref[\detokenize{fix_mol_swap::doc}]{\sphinxcrossref{\DUrole{doc}{mol/swap}}}}
\\
\sphinxhline
\sphinxAtStartPar
{\hyperref[\detokenize{fix_momentum::doc}]{\sphinxcrossref{\DUrole{doc}{momentum (k)}}}}
&
\sphinxAtStartPar
{\hyperref[\detokenize{fix_momentum::doc}]{\sphinxcrossref{\DUrole{doc}{momentum/chunk}}}}
&
\sphinxAtStartPar
{\hyperref[\detokenize{fix_move::doc}]{\sphinxcrossref{\DUrole{doc}{move}}}}
&
\sphinxAtStartPar
{\hyperref[\detokenize{fix_msst::doc}]{\sphinxcrossref{\DUrole{doc}{msst}}}}
\\
\sphinxhline
\sphinxAtStartPar
{\hyperref[\detokenize{fix_mvv_dpd::doc}]{\sphinxcrossref{\DUrole{doc}{mvv/dpd}}}}
&
\sphinxAtStartPar
{\hyperref[\detokenize{fix_mvv_dpd::doc}]{\sphinxcrossref{\DUrole{doc}{mvv/edpd}}}}
&
\sphinxAtStartPar
{\hyperref[\detokenize{fix_mvv_dpd::doc}]{\sphinxcrossref{\DUrole{doc}{mvv/tdpd}}}}
&
\sphinxAtStartPar
{\hyperref[\detokenize{fix_neb::doc}]{\sphinxcrossref{\DUrole{doc}{neb}}}}
\\
\sphinxhline
\sphinxAtStartPar
{\hyperref[\detokenize{fix_neb_spin::doc}]{\sphinxcrossref{\DUrole{doc}{neb/spin}}}}
&
\sphinxAtStartPar
{\hyperref[\detokenize{fix_nonaffine_displacement::doc}]{\sphinxcrossref{\DUrole{doc}{nonaffine/displacement}}}}
&
\sphinxAtStartPar
{\hyperref[\detokenize{fix_nh::doc}]{\sphinxcrossref{\DUrole{doc}{nph (ko)}}}}
&
\sphinxAtStartPar
{\hyperref[\detokenize{fix_nph_asphere::doc}]{\sphinxcrossref{\DUrole{doc}{nph/asphere (o)}}}}
\\
\sphinxhline
\sphinxAtStartPar
{\hyperref[\detokenize{fix_nph_body::doc}]{\sphinxcrossref{\DUrole{doc}{nph/body}}}}
&
\sphinxAtStartPar
{\hyperref[\detokenize{fix_nh_eff::doc}]{\sphinxcrossref{\DUrole{doc}{nph/eff}}}}
&
\sphinxAtStartPar
{\hyperref[\detokenize{fix_nph_sphere::doc}]{\sphinxcrossref{\DUrole{doc}{nph/sphere (o)}}}}
&
\sphinxAtStartPar
{\hyperref[\detokenize{fix_nphug::doc}]{\sphinxcrossref{\DUrole{doc}{nphug}}}}
\\
\sphinxhline
\sphinxAtStartPar
{\hyperref[\detokenize{fix_nh::doc}]{\sphinxcrossref{\DUrole{doc}{npt (giko)}}}}
&
\sphinxAtStartPar
{\hyperref[\detokenize{fix_npt_asphere::doc}]{\sphinxcrossref{\DUrole{doc}{npt/asphere (o)}}}}
&
\sphinxAtStartPar
{\hyperref[\detokenize{fix_npt_body::doc}]{\sphinxcrossref{\DUrole{doc}{npt/body}}}}
&
\sphinxAtStartPar
{\hyperref[\detokenize{fix_npt_cauchy::doc}]{\sphinxcrossref{\DUrole{doc}{npt/cauchy}}}}
\\
\sphinxhline
\sphinxAtStartPar
{\hyperref[\detokenize{fix_nh_eff::doc}]{\sphinxcrossref{\DUrole{doc}{npt/eff}}}}
&
\sphinxAtStartPar
{\hyperref[\detokenize{fix_npt_sphere::doc}]{\sphinxcrossref{\DUrole{doc}{npt/sphere (o)}}}}
&
\sphinxAtStartPar
{\hyperref[\detokenize{fix_nh_uef::doc}]{\sphinxcrossref{\DUrole{doc}{npt/uef}}}}
&
\sphinxAtStartPar
{\hyperref[\detokenize{fix_numdiff::doc}]{\sphinxcrossref{\DUrole{doc}{numdiff}}}}
\\
\sphinxhline
\sphinxAtStartPar
{\hyperref[\detokenize{fix_numdiff_virial::doc}]{\sphinxcrossref{\DUrole{doc}{numdiff/virial}}}}
&
\sphinxAtStartPar
{\hyperref[\detokenize{fix_nve::doc}]{\sphinxcrossref{\DUrole{doc}{nve (giko)}}}}
&
\sphinxAtStartPar
{\hyperref[\detokenize{fix_nve_asphere::doc}]{\sphinxcrossref{\DUrole{doc}{nve/asphere (gi)}}}}
&
\sphinxAtStartPar
{\hyperref[\detokenize{fix_nve_asphere_noforce::doc}]{\sphinxcrossref{\DUrole{doc}{nve/asphere/noforce}}}}
\\
\sphinxhline
\sphinxAtStartPar
{\hyperref[\detokenize{fix_nve_awpmd::doc}]{\sphinxcrossref{\DUrole{doc}{nve/awpmd}}}}
&
\sphinxAtStartPar
{\hyperref[\detokenize{fix_nve_body::doc}]{\sphinxcrossref{\DUrole{doc}{nve/body}}}}
&
\sphinxAtStartPar
{\hyperref[\detokenize{fix_nve_dot::doc}]{\sphinxcrossref{\DUrole{doc}{nve/dot}}}}
&
\sphinxAtStartPar
{\hyperref[\detokenize{fix_nve_dotc_langevin::doc}]{\sphinxcrossref{\DUrole{doc}{nve/dotc/langevin}}}}
\\
\sphinxhline
\sphinxAtStartPar
{\hyperref[\detokenize{fix_nve_eff::doc}]{\sphinxcrossref{\DUrole{doc}{nve/eff}}}}
&
\sphinxAtStartPar
{\hyperref[\detokenize{fix_nve_limit::doc}]{\sphinxcrossref{\DUrole{doc}{nve/limit}}}}
&
\sphinxAtStartPar
{\hyperref[\detokenize{fix_nve_line::doc}]{\sphinxcrossref{\DUrole{doc}{nve/line}}}}
&
\sphinxAtStartPar
{\hyperref[\detokenize{fix_nve_manifold_rattle::doc}]{\sphinxcrossref{\DUrole{doc}{nve/manifold/rattle}}}}
\\
\sphinxhline
\sphinxAtStartPar
{\hyperref[\detokenize{fix_nve_noforce::doc}]{\sphinxcrossref{\DUrole{doc}{nve/noforce}}}}
&
\sphinxAtStartPar
{\hyperref[\detokenize{fix_nve_sphere::doc}]{\sphinxcrossref{\DUrole{doc}{nve/sphere (ko)}}}}
&
\sphinxAtStartPar
{\hyperref[\detokenize{fix_nve_bpm_sphere::doc}]{\sphinxcrossref{\DUrole{doc}{nve/bpm/sphere}}}}
&
\sphinxAtStartPar
{\hyperref[\detokenize{fix_nve_spin::doc}]{\sphinxcrossref{\DUrole{doc}{nve/spin}}}}
\\
\sphinxhline
\sphinxAtStartPar
{\hyperref[\detokenize{fix_nve_tri::doc}]{\sphinxcrossref{\DUrole{doc}{nve/tri}}}}
&
\sphinxAtStartPar
{\hyperref[\detokenize{fix_nvk::doc}]{\sphinxcrossref{\DUrole{doc}{nvk}}}}
&
\sphinxAtStartPar
{\hyperref[\detokenize{fix_nh::doc}]{\sphinxcrossref{\DUrole{doc}{nvt (giko)}}}}
&
\sphinxAtStartPar
{\hyperref[\detokenize{fix_nvt_asphere::doc}]{\sphinxcrossref{\DUrole{doc}{nvt/asphere (o)}}}}
\\
\sphinxhline
\sphinxAtStartPar
{\hyperref[\detokenize{fix_nvt_body::doc}]{\sphinxcrossref{\DUrole{doc}{nvt/body}}}}
&
\sphinxAtStartPar
{\hyperref[\detokenize{fix_nh_eff::doc}]{\sphinxcrossref{\DUrole{doc}{nvt/eff}}}}
&
\sphinxAtStartPar
{\hyperref[\detokenize{fix_nvt_manifold_rattle::doc}]{\sphinxcrossref{\DUrole{doc}{nvt/manifold/rattle}}}}
&
\sphinxAtStartPar
{\hyperref[\detokenize{fix_nvt_sllod::doc}]{\sphinxcrossref{\DUrole{doc}{nvt/sllod (iko)}}}}
\\
\sphinxhline
\sphinxAtStartPar
{\hyperref[\detokenize{fix_nvt_sllod_eff::doc}]{\sphinxcrossref{\DUrole{doc}{nvt/sllod/eff}}}}
&
\sphinxAtStartPar
{\hyperref[\detokenize{fix_nvt_sphere::doc}]{\sphinxcrossref{\DUrole{doc}{nvt/sphere (o)}}}}
&
\sphinxAtStartPar
{\hyperref[\detokenize{fix_nh_uef::doc}]{\sphinxcrossref{\DUrole{doc}{nvt/uef}}}}
&
\sphinxAtStartPar
{\hyperref[\detokenize{fix_oneway::doc}]{\sphinxcrossref{\DUrole{doc}{oneway}}}}
\\
\sphinxhline
\sphinxAtStartPar
{\hyperref[\detokenize{fix_orient::doc}]{\sphinxcrossref{\DUrole{doc}{orient/bcc}}}}
&
\sphinxAtStartPar
{\hyperref[\detokenize{fix_orient::doc}]{\sphinxcrossref{\DUrole{doc}{orient/fcc}}}}
&
\sphinxAtStartPar
{\hyperref[\detokenize{fix_orient_eco::doc}]{\sphinxcrossref{\DUrole{doc}{orient/eco}}}}
&
\sphinxAtStartPar
{\hyperref[\detokenize{fix_pafi::doc}]{\sphinxcrossref{\DUrole{doc}{pafi}}}}
\\
\sphinxhline
\sphinxAtStartPar
{\hyperref[\detokenize{fix_pair::doc}]{\sphinxcrossref{\DUrole{doc}{pair}}}}
&
\sphinxAtStartPar
{\hyperref[\detokenize{fix_phonon::doc}]{\sphinxcrossref{\DUrole{doc}{phonon}}}}
&
\sphinxAtStartPar
{\hyperref[\detokenize{fix_pimd::doc}]{\sphinxcrossref{\DUrole{doc}{pimd/langevin}}}}
&
\sphinxAtStartPar
{\hyperref[\detokenize{fix_pimd::doc}]{\sphinxcrossref{\DUrole{doc}{pimd/nvt}}}}
\\
\sphinxhline
\sphinxAtStartPar
{\hyperref[\detokenize{fix_planeforce::doc}]{\sphinxcrossref{\DUrole{doc}{planeforce}}}}
&
\sphinxAtStartPar
{\hyperref[\detokenize{fix_plumed::doc}]{\sphinxcrossref{\DUrole{doc}{plumed}}}}
&
\sphinxAtStartPar
{\hyperref[\detokenize{fix_poems::doc}]{\sphinxcrossref{\DUrole{doc}{poems}}}}
&
\sphinxAtStartPar
{\hyperref[\detokenize{fix_polarize::doc}]{\sphinxcrossref{\DUrole{doc}{polarize/bem/gmres}}}}
\\
\sphinxhline
\sphinxAtStartPar
{\hyperref[\detokenize{fix_polarize::doc}]{\sphinxcrossref{\DUrole{doc}{polarize/bem/icc}}}}
&
\sphinxAtStartPar
{\hyperref[\detokenize{fix_polarize::doc}]{\sphinxcrossref{\DUrole{doc}{polarize/functional}}}}
&
\sphinxAtStartPar
{\hyperref[\detokenize{fix_pour::doc}]{\sphinxcrossref{\DUrole{doc}{pour}}}}
&
\sphinxAtStartPar
{\hyperref[\detokenize{fix_precession_spin::doc}]{\sphinxcrossref{\DUrole{doc}{precession/spin}}}}
\\
\sphinxhline
\sphinxAtStartPar
{\hyperref[\detokenize{fix_press_berendsen::doc}]{\sphinxcrossref{\DUrole{doc}{press/berendsen}}}}
&
\sphinxAtStartPar
{\hyperref[\detokenize{fix_press_langevin::doc}]{\sphinxcrossref{\DUrole{doc}{press/langevin}}}}
&
\sphinxAtStartPar
{\hyperref[\detokenize{fix_print::doc}]{\sphinxcrossref{\DUrole{doc}{print}}}}
&
\sphinxAtStartPar
{\hyperref[\detokenize{fix_propel_self::doc}]{\sphinxcrossref{\DUrole{doc}{propel/self}}}}
\\
\sphinxhline
\sphinxAtStartPar
{\hyperref[\detokenize{fix_property_atom::doc}]{\sphinxcrossref{\DUrole{doc}{property/atom (k)}}}}
&
\sphinxAtStartPar
{\hyperref[\detokenize{fix_python_invoke::doc}]{\sphinxcrossref{\DUrole{doc}{python/invoke}}}}
&
\sphinxAtStartPar
{\hyperref[\detokenize{fix_python_move::doc}]{\sphinxcrossref{\DUrole{doc}{python/move}}}}
&
\sphinxAtStartPar
{\hyperref[\detokenize{fix_qbmsst::doc}]{\sphinxcrossref{\DUrole{doc}{qbmsst}}}}
\\
\sphinxhline
\sphinxAtStartPar
{\hyperref[\detokenize{fix_qeq_comb::doc}]{\sphinxcrossref{\DUrole{doc}{qeq/comb (o)}}}}
&
\sphinxAtStartPar
{\hyperref[\detokenize{fix_qeq::doc}]{\sphinxcrossref{\DUrole{doc}{qeq/dynamic}}}}
&
\sphinxAtStartPar
{\hyperref[\detokenize{fix_qeq::doc}]{\sphinxcrossref{\DUrole{doc}{qeq/fire}}}}
&
\sphinxAtStartPar
{\hyperref[\detokenize{fix_qeq::doc}]{\sphinxcrossref{\DUrole{doc}{qeq/point}}}}
\\
\sphinxhline
\sphinxAtStartPar
{\hyperref[\detokenize{fix_qeq_reaxff::doc}]{\sphinxcrossref{\DUrole{doc}{qeq/reaxff (ko)}}}}
&
\sphinxAtStartPar
{\hyperref[\detokenize{fix_qeq::doc}]{\sphinxcrossref{\DUrole{doc}{qeq/shielded}}}}
&
\sphinxAtStartPar
{\hyperref[\detokenize{fix_qeq::doc}]{\sphinxcrossref{\DUrole{doc}{qeq/slater}}}}
&
\sphinxAtStartPar
{\hyperref[\detokenize{fix_qmmm::doc}]{\sphinxcrossref{\DUrole{doc}{qmmm}}}}
\\
\sphinxhline
\sphinxAtStartPar
{\hyperref[\detokenize{fix_qtb::doc}]{\sphinxcrossref{\DUrole{doc}{qtb}}}}
&
\sphinxAtStartPar
{\hyperref[\detokenize{fix_shake::doc}]{\sphinxcrossref{\DUrole{doc}{rattle}}}}
&
\sphinxAtStartPar
{\hyperref[\detokenize{fix_reaxff_bonds::doc}]{\sphinxcrossref{\DUrole{doc}{reaxff/bonds (k)}}}}
&
\sphinxAtStartPar
{\hyperref[\detokenize{fix_reaxff_species::doc}]{\sphinxcrossref{\DUrole{doc}{reaxff/species (k)}}}}
\\
\sphinxhline
\sphinxAtStartPar
{\hyperref[\detokenize{fix_recenter::doc}]{\sphinxcrossref{\DUrole{doc}{recenter}}}}
&
\sphinxAtStartPar
{\hyperref[\detokenize{fix_restrain::doc}]{\sphinxcrossref{\DUrole{doc}{restrain}}}}
&
\sphinxAtStartPar
{\hyperref[\detokenize{fix_rheo::doc}]{\sphinxcrossref{\DUrole{doc}{rheo}}}}
&
\sphinxAtStartPar
{\hyperref[\detokenize{fix_rheo_oxidation::doc}]{\sphinxcrossref{\DUrole{doc}{rheo/oxidation}}}}
\\
\sphinxhline
\sphinxAtStartPar
{\hyperref[\detokenize{fix_rheo_pressure::doc}]{\sphinxcrossref{\DUrole{doc}{rheo/pressure}}}}
&
\sphinxAtStartPar
{\hyperref[\detokenize{fix_rheo_thermal::doc}]{\sphinxcrossref{\DUrole{doc}{rheo/thermal}}}}
&
\sphinxAtStartPar
{\hyperref[\detokenize{fix_rheo_viscosity::doc}]{\sphinxcrossref{\DUrole{doc}{rheo/viscosity}}}}
&
\sphinxAtStartPar
{\hyperref[\detokenize{fix_rhok::doc}]{\sphinxcrossref{\DUrole{doc}{rhok}}}}
\\
\sphinxhline
\sphinxAtStartPar
{\hyperref[\detokenize{fix_rigid::doc}]{\sphinxcrossref{\DUrole{doc}{rigid (o)}}}}
&
\sphinxAtStartPar
{\hyperref[\detokenize{fix_rigid_meso::doc}]{\sphinxcrossref{\DUrole{doc}{rigid/meso}}}}
&
\sphinxAtStartPar
{\hyperref[\detokenize{fix_rigid::doc}]{\sphinxcrossref{\DUrole{doc}{rigid/nph (o)}}}}
&
\sphinxAtStartPar
{\hyperref[\detokenize{fix_rigid::doc}]{\sphinxcrossref{\DUrole{doc}{rigid/nph/small}}}}
\\
\sphinxhline
\sphinxAtStartPar
{\hyperref[\detokenize{fix_rigid::doc}]{\sphinxcrossref{\DUrole{doc}{rigid/npt (o)}}}}
&
\sphinxAtStartPar
{\hyperref[\detokenize{fix_rigid::doc}]{\sphinxcrossref{\DUrole{doc}{rigid/npt/small}}}}
&
\sphinxAtStartPar
{\hyperref[\detokenize{fix_rigid::doc}]{\sphinxcrossref{\DUrole{doc}{rigid/nve (o)}}}}
&
\sphinxAtStartPar
{\hyperref[\detokenize{fix_rigid::doc}]{\sphinxcrossref{\DUrole{doc}{rigid/nve/small}}}}
\\
\sphinxhline
\sphinxAtStartPar
{\hyperref[\detokenize{fix_rigid::doc}]{\sphinxcrossref{\DUrole{doc}{rigid/nvt (o)}}}}
&
\sphinxAtStartPar
{\hyperref[\detokenize{fix_rigid::doc}]{\sphinxcrossref{\DUrole{doc}{rigid/nvt/small}}}}
&
\sphinxAtStartPar
{\hyperref[\detokenize{fix_rigid::doc}]{\sphinxcrossref{\DUrole{doc}{rigid/small (o)}}}}
&
\sphinxAtStartPar
{\hyperref[\detokenize{fix_rx::doc}]{\sphinxcrossref{\DUrole{doc}{rx (k)}}}}
\\
\sphinxhline
\sphinxAtStartPar
{\hyperref[\detokenize{fix_saed_vtk::doc}]{\sphinxcrossref{\DUrole{doc}{saed/vtk}}}}
&
\sphinxAtStartPar
{\hyperref[\detokenize{fix_setforce::doc}]{\sphinxcrossref{\DUrole{doc}{setforce (k)}}}}
&
\sphinxAtStartPar
{\hyperref[\detokenize{fix_setforce::doc}]{\sphinxcrossref{\DUrole{doc}{setforce/spin}}}}
&
\sphinxAtStartPar
{\hyperref[\detokenize{fix_sgcmc::doc}]{\sphinxcrossref{\DUrole{doc}{sgcmc}}}}
\\
\sphinxhline
\sphinxAtStartPar
{\hyperref[\detokenize{fix_shake::doc}]{\sphinxcrossref{\DUrole{doc}{shake (k)}}}}
&
\sphinxAtStartPar
{\hyperref[\detokenize{fix_shardlow::doc}]{\sphinxcrossref{\DUrole{doc}{shardlow (k)}}}}
&
\sphinxAtStartPar
{\hyperref[\detokenize{fix_smd::doc}]{\sphinxcrossref{\DUrole{doc}{smd}}}}
&
\sphinxAtStartPar
{\hyperref[\detokenize{fix_smd_adjust_dt::doc}]{\sphinxcrossref{\DUrole{doc}{smd/adjust\_dt}}}}
\\
\sphinxhline
\sphinxAtStartPar
{\hyperref[\detokenize{fix_smd_integrate_tlsph::doc}]{\sphinxcrossref{\DUrole{doc}{smd/integrate\_tlsph}}}}
&
\sphinxAtStartPar
{\hyperref[\detokenize{fix_smd_integrate_ulsph::doc}]{\sphinxcrossref{\DUrole{doc}{smd/integrate\_ulsph}}}}
&
\sphinxAtStartPar
{\hyperref[\detokenize{fix_smd_move_triangulated_surface::doc}]{\sphinxcrossref{\DUrole{doc}{smd/move\_tri\_surf}}}}
&
\sphinxAtStartPar
{\hyperref[\detokenize{fix_smd_setvel::doc}]{\sphinxcrossref{\DUrole{doc}{smd/setvel}}}}
\\
\sphinxhline
\sphinxAtStartPar
{\hyperref[\detokenize{fix_smd_wall_surface::doc}]{\sphinxcrossref{\DUrole{doc}{smd/wall\_surface}}}}
&
\sphinxAtStartPar
{\hyperref[\detokenize{fix_sph::doc}]{\sphinxcrossref{\DUrole{doc}{sph}}}}
&
\sphinxAtStartPar
{\hyperref[\detokenize{fix_sph_stationary::doc}]{\sphinxcrossref{\DUrole{doc}{sph/stationary}}}}
&
\sphinxAtStartPar
{\hyperref[\detokenize{fix_spring::doc}]{\sphinxcrossref{\DUrole{doc}{spring}}}}
\\
\sphinxhline
\sphinxAtStartPar
{\hyperref[\detokenize{fix_spring_chunk::doc}]{\sphinxcrossref{\DUrole{doc}{spring/chunk}}}}
&
\sphinxAtStartPar
{\hyperref[\detokenize{fix_spring_rg::doc}]{\sphinxcrossref{\DUrole{doc}{spring/rg}}}}
&
\sphinxAtStartPar
{\hyperref[\detokenize{fix_spring_self::doc}]{\sphinxcrossref{\DUrole{doc}{spring/self (k)}}}}
&
\sphinxAtStartPar
{\hyperref[\detokenize{fix_srd::doc}]{\sphinxcrossref{\DUrole{doc}{srd}}}}
\\
\sphinxhline
\sphinxAtStartPar
{\hyperref[\detokenize{fix_store_force::doc}]{\sphinxcrossref{\DUrole{doc}{store/force}}}}
&
\sphinxAtStartPar
{\hyperref[\detokenize{fix_store_state::doc}]{\sphinxcrossref{\DUrole{doc}{store/state}}}}
&
\sphinxAtStartPar
{\hyperref[\detokenize{fix_dpd_source::doc}]{\sphinxcrossref{\DUrole{doc}{tdpd/source}}}}
&
\sphinxAtStartPar
{\hyperref[\detokenize{fix_temp_berendsen::doc}]{\sphinxcrossref{\DUrole{doc}{temp/berendsen (k)}}}}
\\
\sphinxhline
\sphinxAtStartPar
{\hyperref[\detokenize{fix_temp_csvr::doc}]{\sphinxcrossref{\DUrole{doc}{temp/csld}}}}
&
\sphinxAtStartPar
{\hyperref[\detokenize{fix_temp_csvr::doc}]{\sphinxcrossref{\DUrole{doc}{temp/csvr}}}}
&
\sphinxAtStartPar
{\hyperref[\detokenize{fix_temp_rescale::doc}]{\sphinxcrossref{\DUrole{doc}{temp/rescale (k)}}}}
&
\sphinxAtStartPar
{\hyperref[\detokenize{fix_temp_rescale_eff::doc}]{\sphinxcrossref{\DUrole{doc}{temp/rescale/eff}}}}
\\
\sphinxhline
\sphinxAtStartPar
{\hyperref[\detokenize{fix_tfmc::doc}]{\sphinxcrossref{\DUrole{doc}{tfmc}}}}
&
\sphinxAtStartPar
{\hyperref[\detokenize{fix_tgnh_drude::doc}]{\sphinxcrossref{\DUrole{doc}{tgnpt/drude}}}}
&
\sphinxAtStartPar
{\hyperref[\detokenize{fix_tgnh_drude::doc}]{\sphinxcrossref{\DUrole{doc}{tgnvt/drude}}}}
&
\sphinxAtStartPar
{\hyperref[\detokenize{fix_thermal_conductivity::doc}]{\sphinxcrossref{\DUrole{doc}{thermal/conductivity}}}}
\\
\sphinxhline
\sphinxAtStartPar
{\hyperref[\detokenize{fix_ti_spring::doc}]{\sphinxcrossref{\DUrole{doc}{ti/spring}}}}
&
\sphinxAtStartPar
{\hyperref[\detokenize{fix_tmd::doc}]{\sphinxcrossref{\DUrole{doc}{tmd}}}}
&
\sphinxAtStartPar
{\hyperref[\detokenize{fix_ttm::doc}]{\sphinxcrossref{\DUrole{doc}{ttm}}}}
&
\sphinxAtStartPar
{\hyperref[\detokenize{fix_ttm::doc}]{\sphinxcrossref{\DUrole{doc}{ttm/grid}}}}
\\
\sphinxhline
\sphinxAtStartPar
{\hyperref[\detokenize{fix_ttm::doc}]{\sphinxcrossref{\DUrole{doc}{ttm/mod}}}}
&
\sphinxAtStartPar
{\hyperref[\detokenize{fix_tune_kspace::doc}]{\sphinxcrossref{\DUrole{doc}{tune/kspace}}}}
&
\sphinxAtStartPar
{\hyperref[\detokenize{fix_vector::doc}]{\sphinxcrossref{\DUrole{doc}{vector}}}}
&
\sphinxAtStartPar
{\hyperref[\detokenize{fix_viscosity::doc}]{\sphinxcrossref{\DUrole{doc}{viscosity}}}}
\\
\sphinxhline
\sphinxAtStartPar
{\hyperref[\detokenize{fix_viscous::doc}]{\sphinxcrossref{\DUrole{doc}{viscous (k)}}}}
&
\sphinxAtStartPar
{\hyperref[\detokenize{fix_viscous_sphere::doc}]{\sphinxcrossref{\DUrole{doc}{viscous/sphere}}}}
&
\sphinxAtStartPar
{\hyperref[\detokenize{fix_wall_body_polygon::doc}]{\sphinxcrossref{\DUrole{doc}{wall/body/polygon}}}}
&
\sphinxAtStartPar
{\hyperref[\detokenize{fix_wall_body_polyhedron::doc}]{\sphinxcrossref{\DUrole{doc}{wall/body/polyhedron}}}}
\\
\sphinxhline
\sphinxAtStartPar
{\hyperref[\detokenize{fix_wall::doc}]{\sphinxcrossref{\DUrole{doc}{wall/colloid}}}}
&
\sphinxAtStartPar
{\hyperref[\detokenize{fix_wall_ees::doc}]{\sphinxcrossref{\DUrole{doc}{wall/ees}}}}
&
\sphinxAtStartPar
{\hyperref[\detokenize{fix_wall_flow::doc}]{\sphinxcrossref{\DUrole{doc}{wall/flow (k)}}}}
&
\sphinxAtStartPar
{\hyperref[\detokenize{fix_wall_gran::doc}]{\sphinxcrossref{\DUrole{doc}{wall/gran (k)}}}}
\\
\sphinxhline
\sphinxAtStartPar
{\hyperref[\detokenize{fix_wall_gran_region::doc}]{\sphinxcrossref{\DUrole{doc}{wall/gran/region}}}}
&
\sphinxAtStartPar
{\hyperref[\detokenize{fix_wall::doc}]{\sphinxcrossref{\DUrole{doc}{wall/harmonic}}}}
&
\sphinxAtStartPar
{\hyperref[\detokenize{fix_wall::doc}]{\sphinxcrossref{\DUrole{doc}{wall/lj1043}}}}
&
\sphinxAtStartPar
{\hyperref[\detokenize{fix_wall::doc}]{\sphinxcrossref{\DUrole{doc}{wall/lj126}}}}
\\
\sphinxhline
\sphinxAtStartPar
{\hyperref[\detokenize{fix_wall::doc}]{\sphinxcrossref{\DUrole{doc}{wall/lj93 (k)}}}}
&
\sphinxAtStartPar
{\hyperref[\detokenize{fix_wall::doc}]{\sphinxcrossref{\DUrole{doc}{wall/lepton}}}}
&
\sphinxAtStartPar
{\hyperref[\detokenize{fix_wall::doc}]{\sphinxcrossref{\DUrole{doc}{wall/morse}}}}
&
\sphinxAtStartPar
{\hyperref[\detokenize{fix_wall_piston::doc}]{\sphinxcrossref{\DUrole{doc}{wall/piston}}}}
\\
\sphinxhline
\sphinxAtStartPar
{\hyperref[\detokenize{fix_wall_reflect::doc}]{\sphinxcrossref{\DUrole{doc}{wall/reflect (k)}}}}
&
\sphinxAtStartPar
{\hyperref[\detokenize{fix_wall_reflect_stochastic::doc}]{\sphinxcrossref{\DUrole{doc}{wall/reflect/stochastic}}}}
&
\sphinxAtStartPar
{\hyperref[\detokenize{fix_wall_region::doc}]{\sphinxcrossref{\DUrole{doc}{wall/region}}}}
&
\sphinxAtStartPar
{\hyperref[\detokenize{fix_wall_ees::doc}]{\sphinxcrossref{\DUrole{doc}{wall/region/ees}}}}
\\
\sphinxhline
\sphinxAtStartPar
{\hyperref[\detokenize{fix_wall_srd::doc}]{\sphinxcrossref{\DUrole{doc}{wall/srd}}}}
&
\sphinxAtStartPar
{\hyperref[\detokenize{fix_wall::doc}]{\sphinxcrossref{\DUrole{doc}{wall/table}}}}
&
\sphinxAtStartPar
{\hyperref[\detokenize{fix_widom::doc}]{\sphinxcrossref{\DUrole{doc}{widom}}}}
&\\
\sphinxbottomrule
\end{longtable}
\sphinxtableafterendhook
\sphinxatlongtableend
\end{savenotes}

\sphinxstepscope


\section{Compute styles}
\label{\detokenize{Commands_compute:compute-styles}}\label{\detokenize{Commands_compute::doc}}
\sphinxAtStartPar
An alphabetic list of all LAMMPS {\hyperref[\detokenize{compute::doc}]{\sphinxcrossref{\DUrole{doc}{compute}}}} commands.
Some styles have accelerated versions.  This is indicated by
additional letters in parenthesis: g = GPU, i = INTEL, k =
KOKKOS, o = OPENMP, t = OPT.


\begin{savenotes}
\sphinxatlongtablestart
\sphinxthistablewithglobalstyle
\makeatletter
  \LTleft \@totalleftmargin plus1fill
  \LTright\dimexpr\columnwidth-\@totalleftmargin-\linewidth\relax plus1fill
\makeatother
\begin{longtable}{llll}
\sphinxtoprule
\endfirsthead

\multicolumn{4}{c}{\sphinxnorowcolor
    \makebox[0pt]{\sphinxtablecontinued{\tablename\ \thetable{} \textendash{} continued from previous page}}%
}\\
\sphinxtoprule
\endhead

\sphinxbottomrule
\multicolumn{4}{r}{\sphinxnorowcolor
    \makebox[0pt][r]{\sphinxtablecontinued{continues on next page}}%
}\\
\endfoot

\endlastfoot
\sphinxtableatstartofbodyhook

\sphinxAtStartPar
{\hyperref[\detokenize{compute_ackland_atom::doc}]{\sphinxcrossref{\DUrole{doc}{ackland/atom}}}}
&
\sphinxAtStartPar
{\hyperref[\detokenize{compute_adf::doc}]{\sphinxcrossref{\DUrole{doc}{adf}}}}
&
\sphinxAtStartPar
{\hyperref[\detokenize{compute_cluster_atom::doc}]{\sphinxcrossref{\DUrole{doc}{aggregate/atom}}}}
&
\sphinxAtStartPar
{\hyperref[\detokenize{compute_angle::doc}]{\sphinxcrossref{\DUrole{doc}{angle}}}}
\\
\sphinxhline
\sphinxAtStartPar
{\hyperref[\detokenize{compute_angle_local::doc}]{\sphinxcrossref{\DUrole{doc}{angle/local}}}}
&
\sphinxAtStartPar
{\hyperref[\detokenize{compute_angmom_chunk::doc}]{\sphinxcrossref{\DUrole{doc}{angmom/chunk}}}}
&
\sphinxAtStartPar
{\hyperref[\detokenize{compute_ave_sphere_atom::doc}]{\sphinxcrossref{\DUrole{doc}{ave/sphere/atom (k)}}}}
&
\sphinxAtStartPar
{\hyperref[\detokenize{compute_basal_atom::doc}]{\sphinxcrossref{\DUrole{doc}{basal/atom}}}}
\\
\sphinxhline
\sphinxAtStartPar
{\hyperref[\detokenize{compute_body_local::doc}]{\sphinxcrossref{\DUrole{doc}{body/local}}}}
&
\sphinxAtStartPar
{\hyperref[\detokenize{compute_bond::doc}]{\sphinxcrossref{\DUrole{doc}{bond}}}}
&
\sphinxAtStartPar
{\hyperref[\detokenize{compute_bond_local::doc}]{\sphinxcrossref{\DUrole{doc}{bond/local}}}}
&
\sphinxAtStartPar
{\hyperref[\detokenize{compute_born_matrix::doc}]{\sphinxcrossref{\DUrole{doc}{born/matrix}}}}
\\
\sphinxhline
\sphinxAtStartPar
{\hyperref[\detokenize{compute_centro_atom::doc}]{\sphinxcrossref{\DUrole{doc}{centro/atom}}}}
&
\sphinxAtStartPar
{\hyperref[\detokenize{compute_stress_atom::doc}]{\sphinxcrossref{\DUrole{doc}{centroid/stress/atom}}}}
&
\sphinxAtStartPar
{\hyperref[\detokenize{compute_chunk_atom::doc}]{\sphinxcrossref{\DUrole{doc}{chunk/atom}}}}
&
\sphinxAtStartPar
{\hyperref[\detokenize{compute_chunk_spread_atom::doc}]{\sphinxcrossref{\DUrole{doc}{chunk/spread/atom}}}}
\\
\sphinxhline
\sphinxAtStartPar
{\hyperref[\detokenize{compute_cluster_atom::doc}]{\sphinxcrossref{\DUrole{doc}{cluster/atom}}}}
&
\sphinxAtStartPar
{\hyperref[\detokenize{compute_cna_atom::doc}]{\sphinxcrossref{\DUrole{doc}{cna/atom}}}}
&
\sphinxAtStartPar
{\hyperref[\detokenize{compute_cnp_atom::doc}]{\sphinxcrossref{\DUrole{doc}{cnp/atom}}}}
&
\sphinxAtStartPar
{\hyperref[\detokenize{compute_com::doc}]{\sphinxcrossref{\DUrole{doc}{com}}}}
\\
\sphinxhline
\sphinxAtStartPar
{\hyperref[\detokenize{compute_com_chunk::doc}]{\sphinxcrossref{\DUrole{doc}{com/chunk}}}}
&
\sphinxAtStartPar
{\hyperref[\detokenize{compute_contact_atom::doc}]{\sphinxcrossref{\DUrole{doc}{contact/atom}}}}
&
\sphinxAtStartPar
{\hyperref[\detokenize{compute_coord_atom::doc}]{\sphinxcrossref{\DUrole{doc}{coord/atom (k)}}}}
&
\sphinxAtStartPar
{\hyperref[\detokenize{compute_count_type::doc}]{\sphinxcrossref{\DUrole{doc}{count/type}}}}
\\
\sphinxhline
\sphinxAtStartPar
{\hyperref[\detokenize{compute_damage_atom::doc}]{\sphinxcrossref{\DUrole{doc}{damage/atom}}}}
&
\sphinxAtStartPar
{\hyperref[\detokenize{compute_dihedral::doc}]{\sphinxcrossref{\DUrole{doc}{dihedral}}}}
&
\sphinxAtStartPar
{\hyperref[\detokenize{compute_dihedral_local::doc}]{\sphinxcrossref{\DUrole{doc}{dihedral/local}}}}
&
\sphinxAtStartPar
{\hyperref[\detokenize{compute_dilatation_atom::doc}]{\sphinxcrossref{\DUrole{doc}{dilatation/atom}}}}
\\
\sphinxhline
\sphinxAtStartPar
{\hyperref[\detokenize{compute_dipole::doc}]{\sphinxcrossref{\DUrole{doc}{dipole}}}}
&
\sphinxAtStartPar
{\hyperref[\detokenize{compute_dipole_chunk::doc}]{\sphinxcrossref{\DUrole{doc}{dipole/chunk}}}}
&
\sphinxAtStartPar
{\hyperref[\detokenize{compute_dipole::doc}]{\sphinxcrossref{\DUrole{doc}{dipole/tip4p}}}}
&
\sphinxAtStartPar
{\hyperref[\detokenize{compute_dipole_chunk::doc}]{\sphinxcrossref{\DUrole{doc}{dipole/tip4p/chunk}}}}
\\
\sphinxhline
\sphinxAtStartPar
{\hyperref[\detokenize{compute_displace_atom::doc}]{\sphinxcrossref{\DUrole{doc}{displace/atom}}}}
&
\sphinxAtStartPar
{\hyperref[\detokenize{compute_dpd::doc}]{\sphinxcrossref{\DUrole{doc}{dpd}}}}
&
\sphinxAtStartPar
{\hyperref[\detokenize{compute_dpd_atom::doc}]{\sphinxcrossref{\DUrole{doc}{dpd/atom}}}}
&
\sphinxAtStartPar
{\hyperref[\detokenize{compute_edpd_temp_atom::doc}]{\sphinxcrossref{\DUrole{doc}{edpd/temp/atom}}}}
\\
\sphinxhline
\sphinxAtStartPar
{\hyperref[\detokenize{compute_efield_atom::doc}]{\sphinxcrossref{\DUrole{doc}{efield/atom}}}}
&
\sphinxAtStartPar
{\hyperref[\detokenize{compute_efield_wolf_atom::doc}]{\sphinxcrossref{\DUrole{doc}{efield/wolf/atom}}}}
&
\sphinxAtStartPar
{\hyperref[\detokenize{compute_entropy_atom::doc}]{\sphinxcrossref{\DUrole{doc}{entropy/atom}}}}
&
\sphinxAtStartPar
{\hyperref[\detokenize{compute_erotate_asphere::doc}]{\sphinxcrossref{\DUrole{doc}{erotate/asphere}}}}
\\
\sphinxhline
\sphinxAtStartPar
{\hyperref[\detokenize{compute_erotate_rigid::doc}]{\sphinxcrossref{\DUrole{doc}{erotate/rigid}}}}
&
\sphinxAtStartPar
{\hyperref[\detokenize{compute_erotate_sphere::doc}]{\sphinxcrossref{\DUrole{doc}{erotate/sphere (k)}}}}
&
\sphinxAtStartPar
{\hyperref[\detokenize{compute_erotate_sphere_atom::doc}]{\sphinxcrossref{\DUrole{doc}{erotate/sphere/atom}}}}
&
\sphinxAtStartPar
{\hyperref[\detokenize{compute_event_displace::doc}]{\sphinxcrossref{\DUrole{doc}{event/displace}}}}
\\
\sphinxhline
\sphinxAtStartPar
{\hyperref[\detokenize{compute_fabric::doc}]{\sphinxcrossref{\DUrole{doc}{fabric}}}}
&
\sphinxAtStartPar
{\hyperref[\detokenize{compute_fep::doc}]{\sphinxcrossref{\DUrole{doc}{fep}}}}
&
\sphinxAtStartPar
{\hyperref[\detokenize{compute_fep_ta::doc}]{\sphinxcrossref{\DUrole{doc}{fep/ta}}}}
&
\sphinxAtStartPar
{\hyperref[\detokenize{compute_tally::doc}]{\sphinxcrossref{\DUrole{doc}{force/tally}}}}
\\
\sphinxhline
\sphinxAtStartPar
{\hyperref[\detokenize{compute_cluster_atom::doc}]{\sphinxcrossref{\DUrole{doc}{fragment/atom}}}}
&
\sphinxAtStartPar
{\hyperref[\detokenize{compute_global_atom::doc}]{\sphinxcrossref{\DUrole{doc}{global/atom}}}}
&
\sphinxAtStartPar
{\hyperref[\detokenize{compute_group_group::doc}]{\sphinxcrossref{\DUrole{doc}{group/group}}}}
&
\sphinxAtStartPar
{\hyperref[\detokenize{compute_gyration::doc}]{\sphinxcrossref{\DUrole{doc}{gyration}}}}
\\
\sphinxhline
\sphinxAtStartPar
{\hyperref[\detokenize{compute_gyration_chunk::doc}]{\sphinxcrossref{\DUrole{doc}{gyration/chunk}}}}
&
\sphinxAtStartPar
{\hyperref[\detokenize{compute_gyration_shape::doc}]{\sphinxcrossref{\DUrole{doc}{gyration/shape}}}}
&
\sphinxAtStartPar
{\hyperref[\detokenize{compute_gyration_shape_chunk::doc}]{\sphinxcrossref{\DUrole{doc}{gyration/shape/chunk}}}}
&
\sphinxAtStartPar
{\hyperref[\detokenize{compute_heat_flux::doc}]{\sphinxcrossref{\DUrole{doc}{heat/flux}}}}
\\
\sphinxhline
\sphinxAtStartPar
{\hyperref[\detokenize{compute_tally::doc}]{\sphinxcrossref{\DUrole{doc}{heat/flux/tally}}}}
&
\sphinxAtStartPar
{\hyperref[\detokenize{compute_tally::doc}]{\sphinxcrossref{\DUrole{doc}{heat/flux/virial/tally}}}}
&
\sphinxAtStartPar
{\hyperref[\detokenize{compute_hexorder_atom::doc}]{\sphinxcrossref{\DUrole{doc}{hexorder/atom}}}}
&
\sphinxAtStartPar
{\hyperref[\detokenize{compute_hma::doc}]{\sphinxcrossref{\DUrole{doc}{hma}}}}
\\
\sphinxhline
\sphinxAtStartPar
{\hyperref[\detokenize{compute_improper::doc}]{\sphinxcrossref{\DUrole{doc}{improper}}}}
&
\sphinxAtStartPar
{\hyperref[\detokenize{compute_improper_local::doc}]{\sphinxcrossref{\DUrole{doc}{improper/local}}}}
&
\sphinxAtStartPar
{\hyperref[\detokenize{compute_inertia_chunk::doc}]{\sphinxcrossref{\DUrole{doc}{inertia/chunk}}}}
&
\sphinxAtStartPar
{\hyperref[\detokenize{compute_ke::doc}]{\sphinxcrossref{\DUrole{doc}{ke}}}}
\\
\sphinxhline
\sphinxAtStartPar
{\hyperref[\detokenize{compute_ke_atom::doc}]{\sphinxcrossref{\DUrole{doc}{ke/atom}}}}
&
\sphinxAtStartPar
{\hyperref[\detokenize{compute_ke_atom_eff::doc}]{\sphinxcrossref{\DUrole{doc}{ke/atom/eff}}}}
&
\sphinxAtStartPar
{\hyperref[\detokenize{compute_ke_eff::doc}]{\sphinxcrossref{\DUrole{doc}{ke/eff}}}}
&
\sphinxAtStartPar
{\hyperref[\detokenize{compute_ke_rigid::doc}]{\sphinxcrossref{\DUrole{doc}{ke/rigid}}}}
\\
\sphinxhline
\sphinxAtStartPar
{\hyperref[\detokenize{compute_composition_atom::doc}]{\sphinxcrossref{\DUrole{doc}{composition/atom (k)}}}}
&
\sphinxAtStartPar
{\hyperref[\detokenize{compute_mliap::doc}]{\sphinxcrossref{\DUrole{doc}{mliap}}}}
&
\sphinxAtStartPar
{\hyperref[\detokenize{compute_momentum::doc}]{\sphinxcrossref{\DUrole{doc}{momentum}}}}
&
\sphinxAtStartPar
{\hyperref[\detokenize{compute_msd::doc}]{\sphinxcrossref{\DUrole{doc}{msd}}}}
\\
\sphinxhline
\sphinxAtStartPar
{\hyperref[\detokenize{compute_msd_chunk::doc}]{\sphinxcrossref{\DUrole{doc}{msd/chunk}}}}
&
\sphinxAtStartPar
{\hyperref[\detokenize{compute_msd_nongauss::doc}]{\sphinxcrossref{\DUrole{doc}{msd/nongauss}}}}
&
\sphinxAtStartPar
{\hyperref[\detokenize{compute_nbond_atom::doc}]{\sphinxcrossref{\DUrole{doc}{nbond/atom}}}}
&
\sphinxAtStartPar
{\hyperref[\detokenize{compute_omega_chunk::doc}]{\sphinxcrossref{\DUrole{doc}{omega/chunk}}}}
\\
\sphinxhline
\sphinxAtStartPar
{\hyperref[\detokenize{compute_orientorder_atom::doc}]{\sphinxcrossref{\DUrole{doc}{orientorder/atom (k)}}}}
&
\sphinxAtStartPar
{\hyperref[\detokenize{compute_pace::doc}]{\sphinxcrossref{\DUrole{doc}{pace}}}}
&
\sphinxAtStartPar
{\hyperref[\detokenize{compute_pair::doc}]{\sphinxcrossref{\DUrole{doc}{pair}}}}
&
\sphinxAtStartPar
{\hyperref[\detokenize{compute_pair_local::doc}]{\sphinxcrossref{\DUrole{doc}{pair/local}}}}
\\
\sphinxhline
\sphinxAtStartPar
{\hyperref[\detokenize{compute_pe::doc}]{\sphinxcrossref{\DUrole{doc}{pe}}}}
&
\sphinxAtStartPar
{\hyperref[\detokenize{compute_pe_atom::doc}]{\sphinxcrossref{\DUrole{doc}{pe/atom}}}}
&
\sphinxAtStartPar
{\hyperref[\detokenize{compute_tally::doc}]{\sphinxcrossref{\DUrole{doc}{pe/mol/tally}}}}
&
\sphinxAtStartPar
{\hyperref[\detokenize{compute_tally::doc}]{\sphinxcrossref{\DUrole{doc}{pe/tally}}}}
\\
\sphinxhline
\sphinxAtStartPar
{\hyperref[\detokenize{compute_plasticity_atom::doc}]{\sphinxcrossref{\DUrole{doc}{plasticity/atom}}}}
&
\sphinxAtStartPar
{\hyperref[\detokenize{compute_pod_atom::doc}]{\sphinxcrossref{\DUrole{doc}{pod/atom}}}}
&
\sphinxAtStartPar
{\hyperref[\detokenize{compute_pod_atom::doc}]{\sphinxcrossref{\DUrole{doc}{podd/atom}}}}
&
\sphinxAtStartPar
{\hyperref[\detokenize{compute_pod_atom::doc}]{\sphinxcrossref{\DUrole{doc}{pod/local}}}}
\\
\sphinxhline
\sphinxAtStartPar
{\hyperref[\detokenize{compute_pod_atom::doc}]{\sphinxcrossref{\DUrole{doc}{pod/global}}}}
&
\sphinxAtStartPar
{\hyperref[\detokenize{compute_pressure::doc}]{\sphinxcrossref{\DUrole{doc}{pressure}}}}
&
\sphinxAtStartPar
{\hyperref[\detokenize{compute_pressure_alchemy::doc}]{\sphinxcrossref{\DUrole{doc}{pressure/alchemy}}}}
&
\sphinxAtStartPar
{\hyperref[\detokenize{compute_pressure_uef::doc}]{\sphinxcrossref{\DUrole{doc}{pressure/uef}}}}
\\
\sphinxhline
\sphinxAtStartPar
{\hyperref[\detokenize{compute_property_atom::doc}]{\sphinxcrossref{\DUrole{doc}{property/atom}}}}
&
\sphinxAtStartPar
{\hyperref[\detokenize{compute_property_chunk::doc}]{\sphinxcrossref{\DUrole{doc}{property/chunk}}}}
&
\sphinxAtStartPar
{\hyperref[\detokenize{compute_property_grid::doc}]{\sphinxcrossref{\DUrole{doc}{property/grid}}}}
&
\sphinxAtStartPar
{\hyperref[\detokenize{compute_property_local::doc}]{\sphinxcrossref{\DUrole{doc}{property/local}}}}
\\
\sphinxhline
\sphinxAtStartPar
{\hyperref[\detokenize{compute_ptm_atom::doc}]{\sphinxcrossref{\DUrole{doc}{ptm/atom}}}}
&
\sphinxAtStartPar
{\hyperref[\detokenize{compute_rattlers_atom::doc}]{\sphinxcrossref{\DUrole{doc}{rattlers/atom}}}}
&
\sphinxAtStartPar
{\hyperref[\detokenize{compute_rdf::doc}]{\sphinxcrossref{\DUrole{doc}{rdf}}}}
&
\sphinxAtStartPar
{\hyperref[\detokenize{compute_reaxff_atom::doc}]{\sphinxcrossref{\DUrole{doc}{reaxff/atom (k)}}}}
\\
\sphinxhline
\sphinxAtStartPar
{\hyperref[\detokenize{compute_reduce::doc}]{\sphinxcrossref{\DUrole{doc}{reduce}}}}
&
\sphinxAtStartPar
{\hyperref[\detokenize{compute_reduce_chunk::doc}]{\sphinxcrossref{\DUrole{doc}{reduce/chunk}}}}
&
\sphinxAtStartPar
{\hyperref[\detokenize{compute_reduce::doc}]{\sphinxcrossref{\DUrole{doc}{reduce/region}}}}
&
\sphinxAtStartPar
{\hyperref[\detokenize{compute_rheo_property_atom::doc}]{\sphinxcrossref{\DUrole{doc}{rheo/property/atom}}}}
\\
\sphinxhline
\sphinxAtStartPar
{\hyperref[\detokenize{compute_rigid_local::doc}]{\sphinxcrossref{\DUrole{doc}{rigid/local}}}}
&
\sphinxAtStartPar
{\hyperref[\detokenize{compute_saed::doc}]{\sphinxcrossref{\DUrole{doc}{saed}}}}
&
\sphinxAtStartPar
{\hyperref[\detokenize{compute_slcsa_atom::doc}]{\sphinxcrossref{\DUrole{doc}{slcsa/atom}}}}
&
\sphinxAtStartPar
{\hyperref[\detokenize{compute_slice::doc}]{\sphinxcrossref{\DUrole{doc}{slice}}}}
\\
\sphinxhline
\sphinxAtStartPar
{\hyperref[\detokenize{compute_smd_contact_radius::doc}]{\sphinxcrossref{\DUrole{doc}{smd/contact/radius}}}}
&
\sphinxAtStartPar
{\hyperref[\detokenize{compute_smd_damage::doc}]{\sphinxcrossref{\DUrole{doc}{smd/damage}}}}
&
\sphinxAtStartPar
{\hyperref[\detokenize{compute_smd_hourglass_error::doc}]{\sphinxcrossref{\DUrole{doc}{smd/hourglass/error}}}}
&
\sphinxAtStartPar
{\hyperref[\detokenize{compute_smd_internal_energy::doc}]{\sphinxcrossref{\DUrole{doc}{smd/internal/energy}}}}
\\
\sphinxhline
\sphinxAtStartPar
{\hyperref[\detokenize{compute_smd_plastic_strain::doc}]{\sphinxcrossref{\DUrole{doc}{smd/plastic/strain}}}}
&
\sphinxAtStartPar
{\hyperref[\detokenize{compute_smd_plastic_strain_rate::doc}]{\sphinxcrossref{\DUrole{doc}{smd/plastic/strain/rate}}}}
&
\sphinxAtStartPar
{\hyperref[\detokenize{compute_smd_rho::doc}]{\sphinxcrossref{\DUrole{doc}{smd/rho}}}}
&
\sphinxAtStartPar
{\hyperref[\detokenize{compute_smd_tlsph_defgrad::doc}]{\sphinxcrossref{\DUrole{doc}{smd/tlsph/defgrad}}}}
\\
\sphinxhline
\sphinxAtStartPar
{\hyperref[\detokenize{compute_smd_tlsph_dt::doc}]{\sphinxcrossref{\DUrole{doc}{smd/tlsph/dt}}}}
&
\sphinxAtStartPar
{\hyperref[\detokenize{compute_smd_tlsph_num_neighs::doc}]{\sphinxcrossref{\DUrole{doc}{smd/tlsph/num/neighs}}}}
&
\sphinxAtStartPar
{\hyperref[\detokenize{compute_smd_tlsph_shape::doc}]{\sphinxcrossref{\DUrole{doc}{smd/tlsph/shape}}}}
&
\sphinxAtStartPar
{\hyperref[\detokenize{compute_smd_tlsph_strain::doc}]{\sphinxcrossref{\DUrole{doc}{smd/tlsph/strain}}}}
\\
\sphinxhline
\sphinxAtStartPar
{\hyperref[\detokenize{compute_smd_tlsph_strain_rate::doc}]{\sphinxcrossref{\DUrole{doc}{smd/tlsph/strain/rate}}}}
&
\sphinxAtStartPar
{\hyperref[\detokenize{compute_smd_tlsph_stress::doc}]{\sphinxcrossref{\DUrole{doc}{smd/tlsph/stress}}}}
&
\sphinxAtStartPar
{\hyperref[\detokenize{compute_smd_triangle_vertices::doc}]{\sphinxcrossref{\DUrole{doc}{smd/triangle/vertices}}}}
&
\sphinxAtStartPar
{\hyperref[\detokenize{compute_smd_ulsph_effm::doc}]{\sphinxcrossref{\DUrole{doc}{smd/ulsph/effm}}}}
\\
\sphinxhline
\sphinxAtStartPar
{\hyperref[\detokenize{compute_smd_ulsph_num_neighs::doc}]{\sphinxcrossref{\DUrole{doc}{smd/ulsph/num/neighs}}}}
&
\sphinxAtStartPar
{\hyperref[\detokenize{compute_smd_ulsph_strain::doc}]{\sphinxcrossref{\DUrole{doc}{smd/ulsph/strain}}}}
&
\sphinxAtStartPar
{\hyperref[\detokenize{compute_smd_ulsph_strain_rate::doc}]{\sphinxcrossref{\DUrole{doc}{smd/ulsph/strain/rate}}}}
&
\sphinxAtStartPar
{\hyperref[\detokenize{compute_smd_ulsph_stress::doc}]{\sphinxcrossref{\DUrole{doc}{smd/ulsph/stress}}}}
\\
\sphinxhline
\sphinxAtStartPar
{\hyperref[\detokenize{compute_smd_vol::doc}]{\sphinxcrossref{\DUrole{doc}{smd/vol}}}}
&
\sphinxAtStartPar
{\hyperref[\detokenize{compute_sna_atom::doc}]{\sphinxcrossref{\DUrole{doc}{snap}}}}
&
\sphinxAtStartPar
{\hyperref[\detokenize{compute_sna_atom::doc}]{\sphinxcrossref{\DUrole{doc}{sna/atom}}}}
&
\sphinxAtStartPar
{\hyperref[\detokenize{compute_sna_atom::doc}]{\sphinxcrossref{\DUrole{doc}{sna/grid}}}}
\\
\sphinxhline
\sphinxAtStartPar
{\hyperref[\detokenize{compute_sna_atom::doc}]{\sphinxcrossref{\DUrole{doc}{sna/grid/local}}}}
&
\sphinxAtStartPar
{\hyperref[\detokenize{compute_sna_atom::doc}]{\sphinxcrossref{\DUrole{doc}{snad/atom}}}}
&
\sphinxAtStartPar
{\hyperref[\detokenize{compute_sna_atom::doc}]{\sphinxcrossref{\DUrole{doc}{snav/atom}}}}
&
\sphinxAtStartPar
{\hyperref[\detokenize{compute_sph_e_atom::doc}]{\sphinxcrossref{\DUrole{doc}{sph/e/atom}}}}
\\
\sphinxhline
\sphinxAtStartPar
{\hyperref[\detokenize{compute_sph_rho_atom::doc}]{\sphinxcrossref{\DUrole{doc}{sph/rho/atom}}}}
&
\sphinxAtStartPar
{\hyperref[\detokenize{compute_sph_t_atom::doc}]{\sphinxcrossref{\DUrole{doc}{sph/t/atom}}}}
&
\sphinxAtStartPar
{\hyperref[\detokenize{compute_spin::doc}]{\sphinxcrossref{\DUrole{doc}{spin}}}}
&
\sphinxAtStartPar
{\hyperref[\detokenize{compute_stress_atom::doc}]{\sphinxcrossref{\DUrole{doc}{stress/atom}}}}
\\
\sphinxhline
\sphinxAtStartPar
{\hyperref[\detokenize{compute_stress_cartesian::doc}]{\sphinxcrossref{\DUrole{doc}{stress/cartesian}}}}
&
\sphinxAtStartPar
{\hyperref[\detokenize{compute_stress_curvilinear::doc}]{\sphinxcrossref{\DUrole{doc}{stress/cylinder}}}}
&
\sphinxAtStartPar
{\hyperref[\detokenize{compute_stress_mop::doc}]{\sphinxcrossref{\DUrole{doc}{stress/mop}}}}
&
\sphinxAtStartPar
{\hyperref[\detokenize{compute_stress_mop::doc}]{\sphinxcrossref{\DUrole{doc}{stress/mop/profile}}}}
\\
\sphinxhline
\sphinxAtStartPar
{\hyperref[\detokenize{compute_stress_curvilinear::doc}]{\sphinxcrossref{\DUrole{doc}{stress/spherical}}}}
&
\sphinxAtStartPar
{\hyperref[\detokenize{compute_tally::doc}]{\sphinxcrossref{\DUrole{doc}{stress/tally}}}}
&
\sphinxAtStartPar
{\hyperref[\detokenize{compute_tdpd_cc_atom::doc}]{\sphinxcrossref{\DUrole{doc}{tdpd/cc/atom}}}}
&
\sphinxAtStartPar
{\hyperref[\detokenize{compute_temp::doc}]{\sphinxcrossref{\DUrole{doc}{temp (k)}}}}
\\
\sphinxhline
\sphinxAtStartPar
{\hyperref[\detokenize{compute_temp_asphere::doc}]{\sphinxcrossref{\DUrole{doc}{temp/asphere}}}}
&
\sphinxAtStartPar
{\hyperref[\detokenize{compute_temp_body::doc}]{\sphinxcrossref{\DUrole{doc}{temp/body}}}}
&
\sphinxAtStartPar
{\hyperref[\detokenize{compute_temp_chunk::doc}]{\sphinxcrossref{\DUrole{doc}{temp/chunk}}}}
&
\sphinxAtStartPar
{\hyperref[\detokenize{compute_temp_com::doc}]{\sphinxcrossref{\DUrole{doc}{temp/com}}}}
\\
\sphinxhline
\sphinxAtStartPar
{\hyperref[\detokenize{compute_temp_cs::doc}]{\sphinxcrossref{\DUrole{doc}{temp/cs}}}}
&
\sphinxAtStartPar
{\hyperref[\detokenize{compute_temp_deform::doc}]{\sphinxcrossref{\DUrole{doc}{temp/deform (k)}}}}
&
\sphinxAtStartPar
{\hyperref[\detokenize{compute_temp_deform_eff::doc}]{\sphinxcrossref{\DUrole{doc}{temp/deform/eff}}}}
&
\sphinxAtStartPar
{\hyperref[\detokenize{compute_temp_drude::doc}]{\sphinxcrossref{\DUrole{doc}{temp/drude}}}}
\\
\sphinxhline
\sphinxAtStartPar
{\hyperref[\detokenize{compute_temp_eff::doc}]{\sphinxcrossref{\DUrole{doc}{temp/eff}}}}
&
\sphinxAtStartPar
{\hyperref[\detokenize{compute_temp_partial::doc}]{\sphinxcrossref{\DUrole{doc}{temp/partial}}}}
&
\sphinxAtStartPar
{\hyperref[\detokenize{compute_temp_profile::doc}]{\sphinxcrossref{\DUrole{doc}{temp/profile}}}}
&
\sphinxAtStartPar
{\hyperref[\detokenize{compute_temp_ramp::doc}]{\sphinxcrossref{\DUrole{doc}{temp/ramp}}}}
\\
\sphinxhline
\sphinxAtStartPar
{\hyperref[\detokenize{compute_temp_region::doc}]{\sphinxcrossref{\DUrole{doc}{temp/region}}}}
&
\sphinxAtStartPar
{\hyperref[\detokenize{compute_temp_region_eff::doc}]{\sphinxcrossref{\DUrole{doc}{temp/region/eff}}}}
&
\sphinxAtStartPar
{\hyperref[\detokenize{compute_temp_rotate::doc}]{\sphinxcrossref{\DUrole{doc}{temp/rotate}}}}
&
\sphinxAtStartPar
{\hyperref[\detokenize{compute_temp_sphere::doc}]{\sphinxcrossref{\DUrole{doc}{temp/sphere}}}}
\\
\sphinxhline
\sphinxAtStartPar
{\hyperref[\detokenize{compute_temp_uef::doc}]{\sphinxcrossref{\DUrole{doc}{temp/uef}}}}
&
\sphinxAtStartPar
{\hyperref[\detokenize{compute_ti::doc}]{\sphinxcrossref{\DUrole{doc}{ti}}}}
&
\sphinxAtStartPar
{\hyperref[\detokenize{compute_torque_chunk::doc}]{\sphinxcrossref{\DUrole{doc}{torque/chunk}}}}
&
\sphinxAtStartPar
{\hyperref[\detokenize{compute_vacf::doc}]{\sphinxcrossref{\DUrole{doc}{vacf}}}}
\\
\sphinxhline
\sphinxAtStartPar
{\hyperref[\detokenize{compute_vcm_chunk::doc}]{\sphinxcrossref{\DUrole{doc}{vcm/chunk}}}}
&
\sphinxAtStartPar
{\hyperref[\detokenize{compute_viscosity_cos::doc}]{\sphinxcrossref{\DUrole{doc}{viscosity/cos}}}}
&
\sphinxAtStartPar
{\hyperref[\detokenize{compute_voronoi_atom::doc}]{\sphinxcrossref{\DUrole{doc}{voronoi/atom}}}}
&
\sphinxAtStartPar
{\hyperref[\detokenize{compute_xrd::doc}]{\sphinxcrossref{\DUrole{doc}{xrd}}}}
\\
\sphinxbottomrule
\end{longtable}
\sphinxtableafterendhook
\sphinxatlongtableend
\end{savenotes}

\sphinxstepscope


\section{Pair styles}
\label{\detokenize{Commands_pair:pair-styles}}\label{\detokenize{Commands_pair::doc}}
\sphinxAtStartPar
All LAMMPS {\hyperref[\detokenize{pair_style::doc}]{\sphinxcrossref{\DUrole{doc}{pair\_style}}}} commands.  Some styles have
accelerated versions.  This is indicated by additional letters in
parenthesis: g = GPU, i = INTEL, k = KOKKOS, o = OPENMP, t =
OPT.


\begin{savenotes}
\sphinxatlongtablestart
\sphinxthistablewithglobalstyle
\makeatletter
  \LTleft \@totalleftmargin plus1fill
  \LTright\dimexpr\columnwidth-\@totalleftmargin-\linewidth\relax plus1fill
\makeatother
\begin{longtable}{lll}
\sphinxtoprule
\endfirsthead

\multicolumn{3}{c}{\sphinxnorowcolor
    \makebox[0pt]{\sphinxtablecontinued{\tablename\ \thetable{} \textendash{} continued from previous page}}%
}\\
\sphinxtoprule
\endhead

\sphinxbottomrule
\multicolumn{3}{r}{\sphinxnorowcolor
    \makebox[0pt][r]{\sphinxtablecontinued{continues on next page}}%
}\\
\endfoot

\endlastfoot
\sphinxtableatstartofbodyhook

\sphinxAtStartPar
{\hyperref[\detokenize{pair_none::doc}]{\sphinxcrossref{\DUrole{doc}{none}}}}
&
\sphinxAtStartPar
{\hyperref[\detokenize{pair_zero::doc}]{\sphinxcrossref{\DUrole{doc}{zero}}}}
&
\sphinxAtStartPar
{\hyperref[\detokenize{pair_hybrid::doc}]{\sphinxcrossref{\DUrole{doc}{hybrid (ko)}}}}
\\
\sphinxhline
\sphinxAtStartPar
{\hyperref[\detokenize{pair_hybrid::doc}]{\sphinxcrossref{\DUrole{doc}{hybrid/molecular (o)}}}}
&
\sphinxAtStartPar
{\hyperref[\detokenize{pair_hybrid::doc}]{\sphinxcrossref{\DUrole{doc}{hybrid/overlay (ko)}}}}
&
\sphinxAtStartPar
{\hyperref[\detokenize{pair_hybrid::doc}]{\sphinxcrossref{\DUrole{doc}{hybrid/scaled (o)}}}}
\\
\sphinxhline
\sphinxAtStartPar
{\hyperref[\detokenize{pair_kim::doc}]{\sphinxcrossref{\DUrole{doc}{kim}}}}
&
\sphinxAtStartPar
{\hyperref[\detokenize{pair_list::doc}]{\sphinxcrossref{\DUrole{doc}{list}}}}
&
\sphinxAtStartPar
{\hyperref[\detokenize{pair_tracker::doc}]{\sphinxcrossref{\DUrole{doc}{tracker}}}}
\\
\sphinxhline&&\\
\sphinxhline
\sphinxAtStartPar
{\hyperref[\detokenize{pair_adp::doc}]{\sphinxcrossref{\DUrole{doc}{adp (ko)}}}}
&
\sphinxAtStartPar
{\hyperref[\detokenize{pair_agni::doc}]{\sphinxcrossref{\DUrole{doc}{agni (o)}}}}
&
\sphinxAtStartPar
{\hyperref[\detokenize{pair_aip_water_2dm::doc}]{\sphinxcrossref{\DUrole{doc}{aip/water/2dm (t)}}}}
\\
\sphinxhline
\sphinxAtStartPar
{\hyperref[\detokenize{pair_airebo::doc}]{\sphinxcrossref{\DUrole{doc}{airebo (io)}}}}
&
\sphinxAtStartPar
{\hyperref[\detokenize{pair_airebo::doc}]{\sphinxcrossref{\DUrole{doc}{airebo/morse (io)}}}}
&
\sphinxAtStartPar
{\hyperref[\detokenize{pair_amoeba::doc}]{\sphinxcrossref{\DUrole{doc}{amoeba (g)}}}}
\\
\sphinxhline
\sphinxAtStartPar
{\hyperref[\detokenize{pair_atm::doc}]{\sphinxcrossref{\DUrole{doc}{atm}}}}
&
\sphinxAtStartPar
{\hyperref[\detokenize{pair_awpmd::doc}]{\sphinxcrossref{\DUrole{doc}{awpmd/cut}}}}
&
\sphinxAtStartPar
{\hyperref[\detokenize{pair_beck::doc}]{\sphinxcrossref{\DUrole{doc}{beck (go)}}}}
\\
\sphinxhline
\sphinxAtStartPar
{\hyperref[\detokenize{pair_body_nparticle::doc}]{\sphinxcrossref{\DUrole{doc}{body/nparticle}}}}
&
\sphinxAtStartPar
{\hyperref[\detokenize{pair_body_rounded_polygon::doc}]{\sphinxcrossref{\DUrole{doc}{body/rounded/polygon}}}}
&
\sphinxAtStartPar
{\hyperref[\detokenize{pair_body_rounded_polyhedron::doc}]{\sphinxcrossref{\DUrole{doc}{body/rounded/polyhedron}}}}
\\
\sphinxhline
\sphinxAtStartPar
{\hyperref[\detokenize{pair_bop::doc}]{\sphinxcrossref{\DUrole{doc}{bop}}}}
&
\sphinxAtStartPar
{\hyperref[\detokenize{pair_born::doc}]{\sphinxcrossref{\DUrole{doc}{born (go)}}}}
&
\sphinxAtStartPar
{\hyperref[\detokenize{pair_born::doc}]{\sphinxcrossref{\DUrole{doc}{born/coul/dsf}}}}
\\
\sphinxhline
\sphinxAtStartPar
{\hyperref[\detokenize{pair_cs::doc}]{\sphinxcrossref{\DUrole{doc}{born/coul/dsf/cs}}}}
&
\sphinxAtStartPar
{\hyperref[\detokenize{pair_born::doc}]{\sphinxcrossref{\DUrole{doc}{born/coul/long (go)}}}}
&
\sphinxAtStartPar
{\hyperref[\detokenize{pair_cs::doc}]{\sphinxcrossref{\DUrole{doc}{born/coul/long/cs (g)}}}}
\\
\sphinxhline
\sphinxAtStartPar
{\hyperref[\detokenize{pair_born::doc}]{\sphinxcrossref{\DUrole{doc}{born/coul/msm (o)}}}}
&
\sphinxAtStartPar
{\hyperref[\detokenize{pair_born::doc}]{\sphinxcrossref{\DUrole{doc}{born/coul/wolf (go)}}}}
&
\sphinxAtStartPar
{\hyperref[\detokenize{pair_cs::doc}]{\sphinxcrossref{\DUrole{doc}{born/coul/wolf/cs (g)}}}}
\\
\sphinxhline
\sphinxAtStartPar
{\hyperref[\detokenize{pair_born_gauss::doc}]{\sphinxcrossref{\DUrole{doc}{born/gauss}}}}
&
\sphinxAtStartPar
{\hyperref[\detokenize{pair_bpm_spring::doc}]{\sphinxcrossref{\DUrole{doc}{bpm/spring}}}}
&
\sphinxAtStartPar
{\hyperref[\detokenize{pair_brownian::doc}]{\sphinxcrossref{\DUrole{doc}{brownian (o)}}}}
\\
\sphinxhline
\sphinxAtStartPar
{\hyperref[\detokenize{pair_brownian::doc}]{\sphinxcrossref{\DUrole{doc}{brownian/poly (o)}}}}
&
\sphinxAtStartPar
{\hyperref[\detokenize{pair_buck::doc}]{\sphinxcrossref{\DUrole{doc}{buck (giko)}}}}
&
\sphinxAtStartPar
{\hyperref[\detokenize{pair_buck::doc}]{\sphinxcrossref{\DUrole{doc}{buck/coul/cut (giko)}}}}
\\
\sphinxhline
\sphinxAtStartPar
{\hyperref[\detokenize{pair_buck::doc}]{\sphinxcrossref{\DUrole{doc}{buck/coul/long (giko)}}}}
&
\sphinxAtStartPar
{\hyperref[\detokenize{pair_cs::doc}]{\sphinxcrossref{\DUrole{doc}{buck/coul/long/cs}}}}
&
\sphinxAtStartPar
{\hyperref[\detokenize{pair_buck::doc}]{\sphinxcrossref{\DUrole{doc}{buck/coul/msm (o)}}}}
\\
\sphinxhline
\sphinxAtStartPar
{\hyperref[\detokenize{pair_buck_long::doc}]{\sphinxcrossref{\DUrole{doc}{buck/long/coul/long (o)}}}}
&
\sphinxAtStartPar
{\hyperref[\detokenize{pair_mdf::doc}]{\sphinxcrossref{\DUrole{doc}{buck/mdf}}}}
&
\sphinxAtStartPar
{\hyperref[\detokenize{pair_buck6d_coul_gauss::doc}]{\sphinxcrossref{\DUrole{doc}{buck6d/coul/gauss/dsf}}}}
\\
\sphinxhline
\sphinxAtStartPar
{\hyperref[\detokenize{pair_buck6d_coul_gauss::doc}]{\sphinxcrossref{\DUrole{doc}{buck6d/coul/gauss/long}}}}
&
\sphinxAtStartPar
{\hyperref[\detokenize{pair_colloid::doc}]{\sphinxcrossref{\DUrole{doc}{colloid (go)}}}}
&
\sphinxAtStartPar
{\hyperref[\detokenize{pair_comb::doc}]{\sphinxcrossref{\DUrole{doc}{comb (o)}}}}
\\
\sphinxhline
\sphinxAtStartPar
{\hyperref[\detokenize{pair_comb::doc}]{\sphinxcrossref{\DUrole{doc}{comb3}}}}
&
\sphinxAtStartPar
{\hyperref[\detokenize{pair_cosine_squared::doc}]{\sphinxcrossref{\DUrole{doc}{cosine/squared}}}}
&
\sphinxAtStartPar
{\hyperref[\detokenize{pair_coul::doc}]{\sphinxcrossref{\DUrole{doc}{coul/cut (gko)}}}}
\\
\sphinxhline
\sphinxAtStartPar
{\hyperref[\detokenize{pair_dielectric::doc}]{\sphinxcrossref{\DUrole{doc}{coul/cut/dielectric}}}}
&
\sphinxAtStartPar
{\hyperref[\detokenize{pair_coul::doc}]{\sphinxcrossref{\DUrole{doc}{coul/cut/global (o)}}}}
&
\sphinxAtStartPar
{\hyperref[\detokenize{pair_fep_soft::doc}]{\sphinxcrossref{\DUrole{doc}{coul/cut/soft (o)}}}}
\\
\sphinxhline
\sphinxAtStartPar
{\hyperref[\detokenize{pair_coul::doc}]{\sphinxcrossref{\DUrole{doc}{coul/debye (gko)}}}}
&
\sphinxAtStartPar
{\hyperref[\detokenize{pair_coul_diel::doc}]{\sphinxcrossref{\DUrole{doc}{coul/diel (o)}}}}
&
\sphinxAtStartPar
{\hyperref[\detokenize{pair_coul::doc}]{\sphinxcrossref{\DUrole{doc}{coul/dsf (gko)}}}}
\\
\sphinxhline
\sphinxAtStartPar
{\hyperref[\detokenize{pair_coul::doc}]{\sphinxcrossref{\DUrole{doc}{coul/exclude}}}}
&
\sphinxAtStartPar
{\hyperref[\detokenize{pair_coul::doc}]{\sphinxcrossref{\DUrole{doc}{coul/long (gko)}}}}
&
\sphinxAtStartPar
{\hyperref[\detokenize{pair_cs::doc}]{\sphinxcrossref{\DUrole{doc}{coul/long/cs (g)}}}}
\\
\sphinxhline
\sphinxAtStartPar
{\hyperref[\detokenize{pair_dielectric::doc}]{\sphinxcrossref{\DUrole{doc}{coul/long/dielectric}}}}
&
\sphinxAtStartPar
{\hyperref[\detokenize{pair_fep_soft::doc}]{\sphinxcrossref{\DUrole{doc}{coul/long/soft (o)}}}}
&
\sphinxAtStartPar
{\hyperref[\detokenize{pair_coul::doc}]{\sphinxcrossref{\DUrole{doc}{coul/msm (o)}}}}
\\
\sphinxhline
\sphinxAtStartPar
{\hyperref[\detokenize{pair_coul_slater::doc}]{\sphinxcrossref{\DUrole{doc}{coul/slater/cut}}}}
&
\sphinxAtStartPar
{\hyperref[\detokenize{pair_coul_slater::doc}]{\sphinxcrossref{\DUrole{doc}{coul/slater/long (g)}}}}
&
\sphinxAtStartPar
{\hyperref[\detokenize{pair_coul_shield::doc}]{\sphinxcrossref{\DUrole{doc}{coul/shield}}}}
\\
\sphinxhline
\sphinxAtStartPar
{\hyperref[\detokenize{pair_coul::doc}]{\sphinxcrossref{\DUrole{doc}{coul/streitz}}}}
&
\sphinxAtStartPar
{\hyperref[\detokenize{pair_coul_tt::doc}]{\sphinxcrossref{\DUrole{doc}{coul/tt}}}}
&
\sphinxAtStartPar
{\hyperref[\detokenize{pair_coul::doc}]{\sphinxcrossref{\DUrole{doc}{coul/wolf (ko)}}}}
\\
\sphinxhline
\sphinxAtStartPar
{\hyperref[\detokenize{pair_cs::doc}]{\sphinxcrossref{\DUrole{doc}{coul/wolf/cs}}}}
&
\sphinxAtStartPar
{\hyperref[\detokenize{pair_dpd::doc}]{\sphinxcrossref{\DUrole{doc}{dpd (giko)}}}}
&
\sphinxAtStartPar
{\hyperref[\detokenize{pair_dpd_coul_slater_long::doc}]{\sphinxcrossref{\DUrole{doc}{dpd/coul/slater/long (g)}}}}
\\
\sphinxhline
\sphinxAtStartPar
{\hyperref[\detokenize{pair_dpd_ext::doc}]{\sphinxcrossref{\DUrole{doc}{dpd/ext (ko)}}}}
&
\sphinxAtStartPar
{\hyperref[\detokenize{pair_dpd_ext::doc}]{\sphinxcrossref{\DUrole{doc}{dpd/ext/tstat (ko)}}}}
&
\sphinxAtStartPar
{\hyperref[\detokenize{pair_dpd_fdt::doc}]{\sphinxcrossref{\DUrole{doc}{dpd/fdt}}}}
\\
\sphinxhline
\sphinxAtStartPar
{\hyperref[\detokenize{pair_dpd_fdt::doc}]{\sphinxcrossref{\DUrole{doc}{dpd/fdt/energy (k)}}}}
&
\sphinxAtStartPar
{\hyperref[\detokenize{pair_dpd::doc}]{\sphinxcrossref{\DUrole{doc}{dpd/tstat (gko)}}}}
&
\sphinxAtStartPar
{\hyperref[\detokenize{pair_dsmc::doc}]{\sphinxcrossref{\DUrole{doc}{dsmc}}}}
\\
\sphinxhline
\sphinxAtStartPar
{\hyperref[\detokenize{pair_e3b::doc}]{\sphinxcrossref{\DUrole{doc}{e3b}}}}
&
\sphinxAtStartPar
{\hyperref[\detokenize{pair_drip::doc}]{\sphinxcrossref{\DUrole{doc}{drip}}}}
&
\sphinxAtStartPar
{\hyperref[\detokenize{pair_eam::doc}]{\sphinxcrossref{\DUrole{doc}{eam (gikot)}}}}
\\
\sphinxhline
\sphinxAtStartPar
{\hyperref[\detokenize{pair_eam::doc}]{\sphinxcrossref{\DUrole{doc}{eam/alloy (gikot)}}}}
&
\sphinxAtStartPar
{\hyperref[\detokenize{pair_eam::doc}]{\sphinxcrossref{\DUrole{doc}{eam/cd}}}}
&
\sphinxAtStartPar
{\hyperref[\detokenize{pair_eam::doc}]{\sphinxcrossref{\DUrole{doc}{eam/cd/old}}}}
\\
\sphinxhline
\sphinxAtStartPar
{\hyperref[\detokenize{pair_eam::doc}]{\sphinxcrossref{\DUrole{doc}{eam/fs (gikot)}}}}
&
\sphinxAtStartPar
{\hyperref[\detokenize{pair_eam::doc}]{\sphinxcrossref{\DUrole{doc}{eam/he}}}}
&
\sphinxAtStartPar
{\hyperref[\detokenize{pair_edip::doc}]{\sphinxcrossref{\DUrole{doc}{edip (o)}}}}
\\
\sphinxhline
\sphinxAtStartPar
{\hyperref[\detokenize{pair_edip::doc}]{\sphinxcrossref{\DUrole{doc}{edip/multi}}}}
&
\sphinxAtStartPar
{\hyperref[\detokenize{pair_mesodpd::doc}]{\sphinxcrossref{\DUrole{doc}{edpd (g)}}}}
&
\sphinxAtStartPar
{\hyperref[\detokenize{pair_eff::doc}]{\sphinxcrossref{\DUrole{doc}{eff/cut}}}}
\\
\sphinxhline
\sphinxAtStartPar
{\hyperref[\detokenize{pair_eim::doc}]{\sphinxcrossref{\DUrole{doc}{eim (o)}}}}
&
\sphinxAtStartPar
{\hyperref[\detokenize{pair_exp6_rx::doc}]{\sphinxcrossref{\DUrole{doc}{exp6/rx (k)}}}}
&
\sphinxAtStartPar
{\hyperref[\detokenize{pair_extep::doc}]{\sphinxcrossref{\DUrole{doc}{extep}}}}
\\
\sphinxhline
\sphinxAtStartPar
{\hyperref[\detokenize{pair_gauss::doc}]{\sphinxcrossref{\DUrole{doc}{gauss (go)}}}}
&
\sphinxAtStartPar
{\hyperref[\detokenize{pair_gauss::doc}]{\sphinxcrossref{\DUrole{doc}{gauss/cut (o)}}}}
&
\sphinxAtStartPar
{\hyperref[\detokenize{pair_gayberne::doc}]{\sphinxcrossref{\DUrole{doc}{gayberne (gio)}}}}
\\
\sphinxhline
\sphinxAtStartPar
{\hyperref[\detokenize{pair_gran::doc}]{\sphinxcrossref{\DUrole{doc}{gran/hertz/history (o)}}}}
&
\sphinxAtStartPar
{\hyperref[\detokenize{pair_gran::doc}]{\sphinxcrossref{\DUrole{doc}{gran/hooke (o)}}}}
&
\sphinxAtStartPar
{\hyperref[\detokenize{pair_gran::doc}]{\sphinxcrossref{\DUrole{doc}{gran/hooke/history (ko)}}}}
\\
\sphinxhline
\sphinxAtStartPar
{\hyperref[\detokenize{pair_granular::doc}]{\sphinxcrossref{\DUrole{doc}{granular}}}}
&
\sphinxAtStartPar
{\hyperref[\detokenize{pair_gw::doc}]{\sphinxcrossref{\DUrole{doc}{gw}}}}
&
\sphinxAtStartPar
{\hyperref[\detokenize{pair_gw::doc}]{\sphinxcrossref{\DUrole{doc}{gw/zbl}}}}
\\
\sphinxhline
\sphinxAtStartPar
{\hyperref[\detokenize{pair_harmonic_cut::doc}]{\sphinxcrossref{\DUrole{doc}{harmonic/cut (o)}}}}
&
\sphinxAtStartPar
{\hyperref[\detokenize{pair_hbond_dreiding::doc}]{\sphinxcrossref{\DUrole{doc}{hbond/dreiding/lj (o)}}}}
&
\sphinxAtStartPar
{\hyperref[\detokenize{pair_hbond_dreiding::doc}]{\sphinxcrossref{\DUrole{doc}{hbond/dreiding/morse (o)}}}}
\\
\sphinxhline
\sphinxAtStartPar
{\hyperref[\detokenize{pair_hdnnp::doc}]{\sphinxcrossref{\DUrole{doc}{hdnnp}}}}
&
\sphinxAtStartPar
{\hyperref[\detokenize{pair_amoeba::doc}]{\sphinxcrossref{\DUrole{doc}{hippo (g)}}}}
&
\sphinxAtStartPar
{\hyperref[\detokenize{pair_ilp_graphene_hbn::doc}]{\sphinxcrossref{\DUrole{doc}{ilp/graphene/hbn (t)}}}}
\\
\sphinxhline
\sphinxAtStartPar
{\hyperref[\detokenize{pair_ilp_tmd::doc}]{\sphinxcrossref{\DUrole{doc}{ilp/tmd (t)}}}}
&
\sphinxAtStartPar
{\hyperref[\detokenize{pair_kolmogorov_crespi_full::doc}]{\sphinxcrossref{\DUrole{doc}{kolmogorov/crespi/full}}}}
&
\sphinxAtStartPar
{\hyperref[\detokenize{pair_kolmogorov_crespi_z::doc}]{\sphinxcrossref{\DUrole{doc}{kolmogorov/crespi/z}}}}
\\
\sphinxhline
\sphinxAtStartPar
{\hyperref[\detokenize{pair_lcbop::doc}]{\sphinxcrossref{\DUrole{doc}{lcbop}}}}
&
\sphinxAtStartPar
{\hyperref[\detokenize{pair_lebedeva_z::doc}]{\sphinxcrossref{\DUrole{doc}{lebedeva/z}}}}
&
\sphinxAtStartPar
{\hyperref[\detokenize{pair_mdf::doc}]{\sphinxcrossref{\DUrole{doc}{lennard/mdf}}}}
\\
\sphinxhline
\sphinxAtStartPar
{\hyperref[\detokenize{pair_lepton::doc}]{\sphinxcrossref{\DUrole{doc}{lepton (o)}}}}
&
\sphinxAtStartPar
{\hyperref[\detokenize{pair_lepton::doc}]{\sphinxcrossref{\DUrole{doc}{lepton/coul (o)}}}}
&
\sphinxAtStartPar
{\hyperref[\detokenize{pair_lepton::doc}]{\sphinxcrossref{\DUrole{doc}{lepton/sphere (o)}}}}
\\
\sphinxhline
\sphinxAtStartPar
{\hyperref[\detokenize{pair_line_lj::doc}]{\sphinxcrossref{\DUrole{doc}{line/lj}}}}
&
\sphinxAtStartPar
{\hyperref[\detokenize{pair_charmm::doc}]{\sphinxcrossref{\DUrole{doc}{lj/charmm/coul/charmm (giko)}}}}
&
\sphinxAtStartPar
{\hyperref[\detokenize{pair_charmm::doc}]{\sphinxcrossref{\DUrole{doc}{lj/charmm/coul/charmm/implicit (ko)}}}}
\\
\sphinxhline
\sphinxAtStartPar
{\hyperref[\detokenize{pair_charmm::doc}]{\sphinxcrossref{\DUrole{doc}{lj/charmm/coul/long (gikot)}}}}
&
\sphinxAtStartPar
{\hyperref[\detokenize{pair_fep_soft::doc}]{\sphinxcrossref{\DUrole{doc}{lj/charmm/coul/long/soft (o)}}}}
&
\sphinxAtStartPar
{\hyperref[\detokenize{pair_charmm::doc}]{\sphinxcrossref{\DUrole{doc}{lj/charmm/coul/msm (o)}}}}
\\
\sphinxhline
\sphinxAtStartPar
{\hyperref[\detokenize{pair_charmm::doc}]{\sphinxcrossref{\DUrole{doc}{lj/charmmfsw/coul/charmmfsh}}}}
&
\sphinxAtStartPar
{\hyperref[\detokenize{pair_charmm::doc}]{\sphinxcrossref{\DUrole{doc}{lj/charmmfsw/coul/long (k)}}}}
&
\sphinxAtStartPar
{\hyperref[\detokenize{pair_class2::doc}]{\sphinxcrossref{\DUrole{doc}{lj/class2 (gko)}}}}
\\
\sphinxhline
\sphinxAtStartPar
{\hyperref[\detokenize{pair_class2::doc}]{\sphinxcrossref{\DUrole{doc}{lj/class2/coul/cut (ko)}}}}
&
\sphinxAtStartPar
{\hyperref[\detokenize{pair_fep_soft::doc}]{\sphinxcrossref{\DUrole{doc}{lj/class2/coul/cut/soft}}}}
&
\sphinxAtStartPar
{\hyperref[\detokenize{pair_class2::doc}]{\sphinxcrossref{\DUrole{doc}{lj/class2/coul/long (gko)}}}}
\\
\sphinxhline
\sphinxAtStartPar
{\hyperref[\detokenize{pair_cs::doc}]{\sphinxcrossref{\DUrole{doc}{lj/class2/coul/long/cs}}}}
&
\sphinxAtStartPar
{\hyperref[\detokenize{pair_fep_soft::doc}]{\sphinxcrossref{\DUrole{doc}{lj/class2/coul/long/soft}}}}
&
\sphinxAtStartPar
{\hyperref[\detokenize{pair_fep_soft::doc}]{\sphinxcrossref{\DUrole{doc}{lj/class2/soft}}}}
\\
\sphinxhline
\sphinxAtStartPar
{\hyperref[\detokenize{pair_lj_cubic::doc}]{\sphinxcrossref{\DUrole{doc}{lj/cubic (go)}}}}
&
\sphinxAtStartPar
{\hyperref[\detokenize{pair_lj::doc}]{\sphinxcrossref{\DUrole{doc}{lj/cut (gikot)}}}}
&
\sphinxAtStartPar
{\hyperref[\detokenize{pair_lj_cut_coul::doc}]{\sphinxcrossref{\DUrole{doc}{lj/cut/coul/cut (gko)}}}}
\\
\sphinxhline
\sphinxAtStartPar
{\hyperref[\detokenize{pair_dielectric::doc}]{\sphinxcrossref{\DUrole{doc}{lj/cut/coul/cut/dielectric (o)}}}}
&
\sphinxAtStartPar
{\hyperref[\detokenize{pair_fep_soft::doc}]{\sphinxcrossref{\DUrole{doc}{lj/cut/coul/cut/soft (go)}}}}
&
\sphinxAtStartPar
{\hyperref[\detokenize{pair_lj_cut_coul::doc}]{\sphinxcrossref{\DUrole{doc}{lj/cut/coul/debye (gko)}}}}
\\
\sphinxhline
\sphinxAtStartPar
{\hyperref[\detokenize{pair_dielectric::doc}]{\sphinxcrossref{\DUrole{doc}{lj/cut/coul/debye/dielectric (o)}}}}
&
\sphinxAtStartPar
{\hyperref[\detokenize{pair_lj_cut_coul::doc}]{\sphinxcrossref{\DUrole{doc}{lj/cut/coul/dsf (gko)}}}}
&
\sphinxAtStartPar
{\hyperref[\detokenize{pair_lj_cut_coul::doc}]{\sphinxcrossref{\DUrole{doc}{lj/cut/coul/long (gikot)}}}}
\\
\sphinxhline
\sphinxAtStartPar
{\hyperref[\detokenize{pair_cs::doc}]{\sphinxcrossref{\DUrole{doc}{lj/cut/coul/long/cs}}}}
&
\sphinxAtStartPar
{\hyperref[\detokenize{pair_dielectric::doc}]{\sphinxcrossref{\DUrole{doc}{lj/cut/coul/long/dielectric (o)}}}}
&
\sphinxAtStartPar
{\hyperref[\detokenize{pair_fep_soft::doc}]{\sphinxcrossref{\DUrole{doc}{lj/cut/coul/long/soft (go)}}}}
\\
\sphinxhline
\sphinxAtStartPar
{\hyperref[\detokenize{pair_lj_cut_coul::doc}]{\sphinxcrossref{\DUrole{doc}{lj/cut/coul/msm (go)}}}}
&
\sphinxAtStartPar
{\hyperref[\detokenize{pair_dielectric::doc}]{\sphinxcrossref{\DUrole{doc}{lj/cut/coul/msm/dielectric}}}}
&
\sphinxAtStartPar
{\hyperref[\detokenize{pair_lj_cut_coul::doc}]{\sphinxcrossref{\DUrole{doc}{lj/cut/coul/wolf (o)}}}}
\\
\sphinxhline
\sphinxAtStartPar
{\hyperref[\detokenize{pair_dipole::doc}]{\sphinxcrossref{\DUrole{doc}{lj/cut/dipole/cut (gko)}}}}
&
\sphinxAtStartPar
{\hyperref[\detokenize{pair_dipole::doc}]{\sphinxcrossref{\DUrole{doc}{lj/cut/dipole/long (g)}}}}
&
\sphinxAtStartPar
{\hyperref[\detokenize{pair_dipole::doc}]{\sphinxcrossref{\DUrole{doc}{lj/cut/dipole/sf (go)}}}}
\\
\sphinxhline
\sphinxAtStartPar
{\hyperref[\detokenize{pair_fep_soft::doc}]{\sphinxcrossref{\DUrole{doc}{lj/cut/soft (o)}}}}
&
\sphinxAtStartPar
{\hyperref[\detokenize{pair_lj_cut_sphere::doc}]{\sphinxcrossref{\DUrole{doc}{lj/cut/sphere (o)}}}}
&
\sphinxAtStartPar
{\hyperref[\detokenize{pair_thole::doc}]{\sphinxcrossref{\DUrole{doc}{lj/cut/thole/long (o)}}}}
\\
\sphinxhline
\sphinxAtStartPar
{\hyperref[\detokenize{pair_lj_cut_tip4p::doc}]{\sphinxcrossref{\DUrole{doc}{lj/cut/tip4p/cut (o)}}}}
&
\sphinxAtStartPar
{\hyperref[\detokenize{pair_lj_cut_tip4p::doc}]{\sphinxcrossref{\DUrole{doc}{lj/cut/tip4p/long (got)}}}}
&
\sphinxAtStartPar
{\hyperref[\detokenize{pair_fep_soft::doc}]{\sphinxcrossref{\DUrole{doc}{lj/cut/tip4p/long/soft (o)}}}}
\\
\sphinxhline
\sphinxAtStartPar
{\hyperref[\detokenize{pair_lj_expand::doc}]{\sphinxcrossref{\DUrole{doc}{lj/expand (gko)}}}}
&
\sphinxAtStartPar
{\hyperref[\detokenize{pair_lj_expand::doc}]{\sphinxcrossref{\DUrole{doc}{lj/expand/coul/long (gk)}}}}
&
\sphinxAtStartPar
{\hyperref[\detokenize{pair_lj_expand_sphere::doc}]{\sphinxcrossref{\DUrole{doc}{lj/expand/sphere (o)}}}}
\\
\sphinxhline
\sphinxAtStartPar
{\hyperref[\detokenize{pair_gromacs::doc}]{\sphinxcrossref{\DUrole{doc}{lj/gromacs (gko)}}}}
&
\sphinxAtStartPar
{\hyperref[\detokenize{pair_gromacs::doc}]{\sphinxcrossref{\DUrole{doc}{lj/gromacs/coul/gromacs (ko)}}}}
&
\sphinxAtStartPar
{\hyperref[\detokenize{pair_lj_long::doc}]{\sphinxcrossref{\DUrole{doc}{lj/long/coul/long (iot)}}}}
\\
\sphinxhline
\sphinxAtStartPar
{\hyperref[\detokenize{pair_dielectric::doc}]{\sphinxcrossref{\DUrole{doc}{lj/long/coul/long/dielectric}}}}
&
\sphinxAtStartPar
{\hyperref[\detokenize{pair_dipole::doc}]{\sphinxcrossref{\DUrole{doc}{lj/long/dipole/long}}}}
&
\sphinxAtStartPar
{\hyperref[\detokenize{pair_lj_long::doc}]{\sphinxcrossref{\DUrole{doc}{lj/long/tip4p/long (o)}}}}
\\
\sphinxhline
\sphinxAtStartPar
{\hyperref[\detokenize{pair_mdf::doc}]{\sphinxcrossref{\DUrole{doc}{lj/mdf}}}}
&
\sphinxAtStartPar
{\hyperref[\detokenize{pair_lj_relres::doc}]{\sphinxcrossref{\DUrole{doc}{lj/relres (o)}}}}
&
\sphinxAtStartPar
{\hyperref[\detokenize{pair_spica::doc}]{\sphinxcrossref{\DUrole{doc}{lj/spica (gko)}}}}
\\
\sphinxhline
\sphinxAtStartPar
{\hyperref[\detokenize{pair_spica::doc}]{\sphinxcrossref{\DUrole{doc}{lj/spica/coul/long (gko)}}}}
&
\sphinxAtStartPar
{\hyperref[\detokenize{pair_spica::doc}]{\sphinxcrossref{\DUrole{doc}{lj/spica/coul/msm (o)}}}}
&
\sphinxAtStartPar
{\hyperref[\detokenize{pair_dipole::doc}]{\sphinxcrossref{\DUrole{doc}{lj/sf/dipole/sf (go)}}}}
\\
\sphinxhline
\sphinxAtStartPar
{\hyperref[\detokenize{pair_lj_smooth::doc}]{\sphinxcrossref{\DUrole{doc}{lj/smooth (go)}}}}
&
\sphinxAtStartPar
{\hyperref[\detokenize{pair_lj_smooth_linear::doc}]{\sphinxcrossref{\DUrole{doc}{lj/smooth/linear (o)}}}}
&
\sphinxAtStartPar
{\hyperref[\detokenize{pair_lj_switch3_coulgauss_long::doc}]{\sphinxcrossref{\DUrole{doc}{lj/switch3/coulgauss/long}}}}
\\
\sphinxhline
\sphinxAtStartPar
{\hyperref[\detokenize{pair_lj96::doc}]{\sphinxcrossref{\DUrole{doc}{lj96/cut (go)}}}}
&
\sphinxAtStartPar
{\hyperref[\detokenize{pair_local_density::doc}]{\sphinxcrossref{\DUrole{doc}{local/density}}}}
&
\sphinxAtStartPar
{\hyperref[\detokenize{pair_lubricate::doc}]{\sphinxcrossref{\DUrole{doc}{lubricate (o)}}}}
\\
\sphinxhline
\sphinxAtStartPar
{\hyperref[\detokenize{pair_lubricate::doc}]{\sphinxcrossref{\DUrole{doc}{lubricate/poly (o)}}}}
&
\sphinxAtStartPar
{\hyperref[\detokenize{pair_lubricateU::doc}]{\sphinxcrossref{\DUrole{doc}{lubricateU}}}}
&
\sphinxAtStartPar
{\hyperref[\detokenize{pair_lubricateU::doc}]{\sphinxcrossref{\DUrole{doc}{lubricateU/poly}}}}
\\
\sphinxhline
\sphinxAtStartPar
{\hyperref[\detokenize{pair_mesodpd::doc}]{\sphinxcrossref{\DUrole{doc}{mdpd (g)}}}}
&
\sphinxAtStartPar
{\hyperref[\detokenize{pair_mesodpd::doc}]{\sphinxcrossref{\DUrole{doc}{mdpd/rhosum}}}}
&
\sphinxAtStartPar
{\hyperref[\detokenize{pair_meam::doc}]{\sphinxcrossref{\DUrole{doc}{meam (k)}}}}
\\
\sphinxhline
\sphinxAtStartPar
{\hyperref[\detokenize{pair_meam::doc}]{\sphinxcrossref{\DUrole{doc}{meam/ms (k)}}}}
&
\sphinxAtStartPar
{\hyperref[\detokenize{pair_meam_spline::doc}]{\sphinxcrossref{\DUrole{doc}{meam/spline (o)}}}}
&
\sphinxAtStartPar
{\hyperref[\detokenize{pair_meam_sw_spline::doc}]{\sphinxcrossref{\DUrole{doc}{meam/sw/spline}}}}
\\
\sphinxhline
\sphinxAtStartPar
{\hyperref[\detokenize{pair_mesocnt::doc}]{\sphinxcrossref{\DUrole{doc}{mesocnt}}}}
&
\sphinxAtStartPar
{\hyperref[\detokenize{pair_mesocnt::doc}]{\sphinxcrossref{\DUrole{doc}{mesocnt/viscous}}}}
&
\sphinxAtStartPar
{\hyperref[\detokenize{pair_mgpt::doc}]{\sphinxcrossref{\DUrole{doc}{mgpt}}}}
\\
\sphinxhline
\sphinxAtStartPar
{\hyperref[\detokenize{pair_mie::doc}]{\sphinxcrossref{\DUrole{doc}{mie/cut (g)}}}}
&
\sphinxAtStartPar
{\hyperref[\detokenize{pair_mliap::doc}]{\sphinxcrossref{\DUrole{doc}{mliap (k)}}}}
&
\sphinxAtStartPar
{\hyperref[\detokenize{pair_lj_switch3_coulgauss_long::doc}]{\sphinxcrossref{\DUrole{doc}{mm3/switch3/coulgauss/long}}}}
\\
\sphinxhline
\sphinxAtStartPar
{\hyperref[\detokenize{pair_momb::doc}]{\sphinxcrossref{\DUrole{doc}{momb}}}}
&
\sphinxAtStartPar
{\hyperref[\detokenize{pair_morse::doc}]{\sphinxcrossref{\DUrole{doc}{morse (gkot)}}}}
&
\sphinxAtStartPar
{\hyperref[\detokenize{pair_morse::doc}]{\sphinxcrossref{\DUrole{doc}{morse/smooth/linear (o)}}}}
\\
\sphinxhline
\sphinxAtStartPar
{\hyperref[\detokenize{pair_fep_soft::doc}]{\sphinxcrossref{\DUrole{doc}{morse/soft}}}}
&
\sphinxAtStartPar
{\hyperref[\detokenize{pair_multi_lucy::doc}]{\sphinxcrossref{\DUrole{doc}{multi/lucy}}}}
&
\sphinxAtStartPar
{\hyperref[\detokenize{pair_multi_lucy_rx::doc}]{\sphinxcrossref{\DUrole{doc}{multi/lucy/rx (k)}}}}
\\
\sphinxhline
\sphinxAtStartPar
{\hyperref[\detokenize{pair_nb3b::doc}]{\sphinxcrossref{\DUrole{doc}{nb3b/harmonic}}}}
&
\sphinxAtStartPar
{\hyperref[\detokenize{pair_nb3b::doc}]{\sphinxcrossref{\DUrole{doc}{nb3b/screened}}}}
&
\sphinxAtStartPar
{\hyperref[\detokenize{pair_nm::doc}]{\sphinxcrossref{\DUrole{doc}{nm/cut (o)}}}}
\\
\sphinxhline
\sphinxAtStartPar
{\hyperref[\detokenize{pair_nm::doc}]{\sphinxcrossref{\DUrole{doc}{nm/cut/coul/cut (o)}}}}
&
\sphinxAtStartPar
{\hyperref[\detokenize{pair_nm::doc}]{\sphinxcrossref{\DUrole{doc}{nm/cut/coul/long (o)}}}}
&
\sphinxAtStartPar
{\hyperref[\detokenize{pair_nm::doc}]{\sphinxcrossref{\DUrole{doc}{nm/cut/split}}}}
\\
\sphinxhline
\sphinxAtStartPar
{\hyperref[\detokenize{pair_oxdna::doc}]{\sphinxcrossref{\DUrole{doc}{oxdna/coaxstk}}}}
&
\sphinxAtStartPar
{\hyperref[\detokenize{pair_oxdna::doc}]{\sphinxcrossref{\DUrole{doc}{oxdna/excv}}}}
&
\sphinxAtStartPar
{\hyperref[\detokenize{pair_oxdna::doc}]{\sphinxcrossref{\DUrole{doc}{oxdna/hbond}}}}
\\
\sphinxhline
\sphinxAtStartPar
{\hyperref[\detokenize{pair_oxdna::doc}]{\sphinxcrossref{\DUrole{doc}{oxdna/stk}}}}
&
\sphinxAtStartPar
{\hyperref[\detokenize{pair_oxdna::doc}]{\sphinxcrossref{\DUrole{doc}{oxdna/xstk}}}}
&
\sphinxAtStartPar
{\hyperref[\detokenize{pair_oxdna2::doc}]{\sphinxcrossref{\DUrole{doc}{oxdna2/coaxstk}}}}
\\
\sphinxhline
\sphinxAtStartPar
{\hyperref[\detokenize{pair_oxdna2::doc}]{\sphinxcrossref{\DUrole{doc}{oxdna2/dh}}}}
&
\sphinxAtStartPar
{\hyperref[\detokenize{pair_oxdna2::doc}]{\sphinxcrossref{\DUrole{doc}{oxdna2/excv}}}}
&
\sphinxAtStartPar
{\hyperref[\detokenize{pair_oxdna2::doc}]{\sphinxcrossref{\DUrole{doc}{oxdna2/hbond}}}}
\\
\sphinxhline
\sphinxAtStartPar
{\hyperref[\detokenize{pair_oxdna2::doc}]{\sphinxcrossref{\DUrole{doc}{oxdna2/stk}}}}
&
\sphinxAtStartPar
{\hyperref[\detokenize{pair_oxdna2::doc}]{\sphinxcrossref{\DUrole{doc}{oxdna2/xstk}}}}
&
\sphinxAtStartPar
{\hyperref[\detokenize{pair_oxrna2::doc}]{\sphinxcrossref{\DUrole{doc}{oxrna2/excv}}}}
\\
\sphinxhline
\sphinxAtStartPar
{\hyperref[\detokenize{pair_oxrna2::doc}]{\sphinxcrossref{\DUrole{doc}{oxrna2/hbond}}}}
&
\sphinxAtStartPar
{\hyperref[\detokenize{pair_oxrna2::doc}]{\sphinxcrossref{\DUrole{doc}{oxrna2/dh}}}}
&
\sphinxAtStartPar
{\hyperref[\detokenize{pair_oxrna2::doc}]{\sphinxcrossref{\DUrole{doc}{oxrna2/stk}}}}
\\
\sphinxhline
\sphinxAtStartPar
{\hyperref[\detokenize{pair_oxrna2::doc}]{\sphinxcrossref{\DUrole{doc}{oxrna2/xstk}}}}
&
\sphinxAtStartPar
{\hyperref[\detokenize{pair_oxrna2::doc}]{\sphinxcrossref{\DUrole{doc}{oxrna2/coaxstk}}}}
&
\sphinxAtStartPar
{\hyperref[\detokenize{pair_pace::doc}]{\sphinxcrossref{\DUrole{doc}{pace (k)}}}}
\\
\sphinxhline
\sphinxAtStartPar
{\hyperref[\detokenize{pair_pace::doc}]{\sphinxcrossref{\DUrole{doc}{pace/extrapolation (k)}}}}
&
\sphinxAtStartPar
{\hyperref[\detokenize{pair_pedone::doc}]{\sphinxcrossref{\DUrole{doc}{pedone (o)}}}}
&
\sphinxAtStartPar
{\hyperref[\detokenize{pair_pod::doc}]{\sphinxcrossref{\DUrole{doc}{pod (k)}}}}
\\
\sphinxhline
\sphinxAtStartPar
{\hyperref[\detokenize{pair_peri::doc}]{\sphinxcrossref{\DUrole{doc}{peri/eps}}}}
&
\sphinxAtStartPar
{\hyperref[\detokenize{pair_peri::doc}]{\sphinxcrossref{\DUrole{doc}{peri/lps (o)}}}}
&
\sphinxAtStartPar
{\hyperref[\detokenize{pair_peri::doc}]{\sphinxcrossref{\DUrole{doc}{peri/pmb (o)}}}}
\\
\sphinxhline
\sphinxAtStartPar
{\hyperref[\detokenize{pair_peri::doc}]{\sphinxcrossref{\DUrole{doc}{peri/ves}}}}
&
\sphinxAtStartPar
{\hyperref[\detokenize{pair_polymorphic::doc}]{\sphinxcrossref{\DUrole{doc}{polymorphic}}}}
&
\sphinxAtStartPar
{\hyperref[\detokenize{pair_python::doc}]{\sphinxcrossref{\DUrole{doc}{python}}}}
\\
\sphinxhline
\sphinxAtStartPar
{\hyperref[\detokenize{pair_quip::doc}]{\sphinxcrossref{\DUrole{doc}{quip}}}}
&
\sphinxAtStartPar
{\hyperref[\detokenize{pair_rann::doc}]{\sphinxcrossref{\DUrole{doc}{rann}}}}
&
\sphinxAtStartPar
{\hyperref[\detokenize{pair_reaxff::doc}]{\sphinxcrossref{\DUrole{doc}{reaxff (ko)}}}}
\\
\sphinxhline
\sphinxAtStartPar
{\hyperref[\detokenize{pair_airebo::doc}]{\sphinxcrossref{\DUrole{doc}{rebo (io)}}}}
&
\sphinxAtStartPar
{\hyperref[\detokenize{pair_rebomos::doc}]{\sphinxcrossref{\DUrole{doc}{rebomos (o)}}}}
&
\sphinxAtStartPar
{\hyperref[\detokenize{pair_resquared::doc}]{\sphinxcrossref{\DUrole{doc}{resquared (go)}}}}
\\
\sphinxhline
\sphinxAtStartPar
{\hyperref[\detokenize{pair_rheo::doc}]{\sphinxcrossref{\DUrole{doc}{rheo}}}}
&
\sphinxAtStartPar
{\hyperref[\detokenize{pair_rheo_solid::doc}]{\sphinxcrossref{\DUrole{doc}{rheo/solid}}}}
&
\sphinxAtStartPar
{\hyperref[\detokenize{pair_saip_metal::doc}]{\sphinxcrossref{\DUrole{doc}{saip/metal (t)}}}}
\\
\sphinxhline
\sphinxAtStartPar
{\hyperref[\detokenize{pair_sdpd_taitwater_isothermal::doc}]{\sphinxcrossref{\DUrole{doc}{sdpd/taitwater/isothermal}}}}
&
\sphinxAtStartPar
{\hyperref[\detokenize{pair_smatb::doc}]{\sphinxcrossref{\DUrole{doc}{smatb}}}}
&
\sphinxAtStartPar
{\hyperref[\detokenize{pair_smatb::doc}]{\sphinxcrossref{\DUrole{doc}{smatb/single}}}}
\\
\sphinxhline
\sphinxAtStartPar
{\hyperref[\detokenize{pair_smd_hertz::doc}]{\sphinxcrossref{\DUrole{doc}{smd/hertz}}}}
&
\sphinxAtStartPar
{\hyperref[\detokenize{pair_smd_tlsph::doc}]{\sphinxcrossref{\DUrole{doc}{smd/tlsph}}}}
&
\sphinxAtStartPar
{\hyperref[\detokenize{pair_smd_triangulated_surface::doc}]{\sphinxcrossref{\DUrole{doc}{smd/tri\_surface}}}}
\\
\sphinxhline
\sphinxAtStartPar
{\hyperref[\detokenize{pair_smd_ulsph::doc}]{\sphinxcrossref{\DUrole{doc}{smd/ulsph}}}}
&
\sphinxAtStartPar
{\hyperref[\detokenize{pair_smtbq::doc}]{\sphinxcrossref{\DUrole{doc}{smtbq}}}}
&
\sphinxAtStartPar
{\hyperref[\detokenize{pair_snap::doc}]{\sphinxcrossref{\DUrole{doc}{snap (ik)}}}}
\\
\sphinxhline
\sphinxAtStartPar
{\hyperref[\detokenize{pair_soft::doc}]{\sphinxcrossref{\DUrole{doc}{soft (gko)}}}}
&
\sphinxAtStartPar
{\hyperref[\detokenize{pair_sph_heatconduction::doc}]{\sphinxcrossref{\DUrole{doc}{sph/heatconduction (g)}}}}
&
\sphinxAtStartPar
{\hyperref[\detokenize{pair_sph_idealgas::doc}]{\sphinxcrossref{\DUrole{doc}{sph/idealgas}}}}
\\
\sphinxhline
\sphinxAtStartPar
{\hyperref[\detokenize{pair_sph_lj::doc}]{\sphinxcrossref{\DUrole{doc}{sph/lj (g)}}}}
&
\sphinxAtStartPar
{\hyperref[\detokenize{pair_sph_rhosum::doc}]{\sphinxcrossref{\DUrole{doc}{sph/rhosum}}}}
&
\sphinxAtStartPar
{\hyperref[\detokenize{pair_sph_taitwater::doc}]{\sphinxcrossref{\DUrole{doc}{sph/taitwater (g)}}}}
\\
\sphinxhline
\sphinxAtStartPar
{\hyperref[\detokenize{pair_sph_taitwater_morris::doc}]{\sphinxcrossref{\DUrole{doc}{sph/taitwater/morris}}}}
&
\sphinxAtStartPar
{\hyperref[\detokenize{pair_spin_dipole::doc}]{\sphinxcrossref{\DUrole{doc}{spin/dipole/cut}}}}
&
\sphinxAtStartPar
{\hyperref[\detokenize{pair_spin_dipole::doc}]{\sphinxcrossref{\DUrole{doc}{spin/dipole/long}}}}
\\
\sphinxhline
\sphinxAtStartPar
{\hyperref[\detokenize{pair_spin_dmi::doc}]{\sphinxcrossref{\DUrole{doc}{spin/dmi}}}}
&
\sphinxAtStartPar
{\hyperref[\detokenize{pair_spin_exchange::doc}]{\sphinxcrossref{\DUrole{doc}{spin/exchange}}}}
&
\sphinxAtStartPar
{\hyperref[\detokenize{pair_spin_exchange::doc}]{\sphinxcrossref{\DUrole{doc}{spin/exchange/biquadratic}}}}
\\
\sphinxhline
\sphinxAtStartPar
{\hyperref[\detokenize{pair_spin_magelec::doc}]{\sphinxcrossref{\DUrole{doc}{spin/magelec}}}}
&
\sphinxAtStartPar
{\hyperref[\detokenize{pair_spin_neel::doc}]{\sphinxcrossref{\DUrole{doc}{spin/neel}}}}
&
\sphinxAtStartPar
{\hyperref[\detokenize{pair_srp::doc}]{\sphinxcrossref{\DUrole{doc}{srp}}}}
\\
\sphinxhline
\sphinxAtStartPar
{\hyperref[\detokenize{pair_srp::doc}]{\sphinxcrossref{\DUrole{doc}{srp/react}}}}
&
\sphinxAtStartPar
{\hyperref[\detokenize{pair_sw::doc}]{\sphinxcrossref{\DUrole{doc}{sw (giko)}}}}
&
\sphinxAtStartPar
{\hyperref[\detokenize{pair_sw_angle_table::doc}]{\sphinxcrossref{\DUrole{doc}{sw/angle/table}}}}
\\
\sphinxhline
\sphinxAtStartPar
{\hyperref[\detokenize{pair_sw::doc}]{\sphinxcrossref{\DUrole{doc}{sw/mod (o)}}}}
&
\sphinxAtStartPar
{\hyperref[\detokenize{pair_table::doc}]{\sphinxcrossref{\DUrole{doc}{table (gko)}}}}
&
\sphinxAtStartPar
{\hyperref[\detokenize{pair_table_rx::doc}]{\sphinxcrossref{\DUrole{doc}{table/rx (k)}}}}
\\
\sphinxhline
\sphinxAtStartPar
{\hyperref[\detokenize{pair_mesodpd::doc}]{\sphinxcrossref{\DUrole{doc}{tdpd}}}}
&
\sphinxAtStartPar
{\hyperref[\detokenize{pair_tersoff::doc}]{\sphinxcrossref{\DUrole{doc}{tersoff (giko)}}}}
&
\sphinxAtStartPar
{\hyperref[\detokenize{pair_tersoff_mod::doc}]{\sphinxcrossref{\DUrole{doc}{tersoff/mod (gko)}}}}
\\
\sphinxhline
\sphinxAtStartPar
{\hyperref[\detokenize{pair_tersoff_mod::doc}]{\sphinxcrossref{\DUrole{doc}{tersoff/mod/c (o)}}}}
&
\sphinxAtStartPar
{\hyperref[\detokenize{pair_tersoff::doc}]{\sphinxcrossref{\DUrole{doc}{tersoff/table (o)}}}}
&
\sphinxAtStartPar
{\hyperref[\detokenize{pair_tersoff_zbl::doc}]{\sphinxcrossref{\DUrole{doc}{tersoff/zbl (gko)}}}}
\\
\sphinxhline
\sphinxAtStartPar
{\hyperref[\detokenize{pair_thole::doc}]{\sphinxcrossref{\DUrole{doc}{thole}}}}
&
\sphinxAtStartPar
{\hyperref[\detokenize{pair_threebody_table::doc}]{\sphinxcrossref{\DUrole{doc}{threebody/table}}}}
&
\sphinxAtStartPar
{\hyperref[\detokenize{pair_coul::doc}]{\sphinxcrossref{\DUrole{doc}{tip4p/cut (o)}}}}
\\
\sphinxhline
\sphinxAtStartPar
{\hyperref[\detokenize{pair_coul::doc}]{\sphinxcrossref{\DUrole{doc}{tip4p/long (o)}}}}
&
\sphinxAtStartPar
{\hyperref[\detokenize{pair_fep_soft::doc}]{\sphinxcrossref{\DUrole{doc}{tip4p/long/soft (o)}}}}
&
\sphinxAtStartPar
{\hyperref[\detokenize{pair_tri_lj::doc}]{\sphinxcrossref{\DUrole{doc}{tri/lj}}}}
\\
\sphinxhline
\sphinxAtStartPar
{\hyperref[\detokenize{pair_ufm::doc}]{\sphinxcrossref{\DUrole{doc}{ufm (got)}}}}
&
\sphinxAtStartPar
{\hyperref[\detokenize{pair_uf3::doc}]{\sphinxcrossref{\DUrole{doc}{uf3 (k)}}}}
&
\sphinxAtStartPar
{\hyperref[\detokenize{pair_vashishta::doc}]{\sphinxcrossref{\DUrole{doc}{vashishta (gko)}}}}
\\
\sphinxhline
\sphinxAtStartPar
{\hyperref[\detokenize{pair_vashishta::doc}]{\sphinxcrossref{\DUrole{doc}{vashishta/table (o)}}}}
&
\sphinxAtStartPar
{\hyperref[\detokenize{pair_wf_cut::doc}]{\sphinxcrossref{\DUrole{doc}{wf/cut}}}}
&
\sphinxAtStartPar
{\hyperref[\detokenize{pair_ylz::doc}]{\sphinxcrossref{\DUrole{doc}{ylz}}}}
\\
\sphinxhline
\sphinxAtStartPar
{\hyperref[\detokenize{pair_yukawa::doc}]{\sphinxcrossref{\DUrole{doc}{yukawa (gko)}}}}
&
\sphinxAtStartPar
{\hyperref[\detokenize{pair_yukawa_colloid::doc}]{\sphinxcrossref{\DUrole{doc}{yukawa/colloid (gko)}}}}
&
\sphinxAtStartPar
{\hyperref[\detokenize{pair_zbl::doc}]{\sphinxcrossref{\DUrole{doc}{zbl (gko)}}}}
\\
\sphinxbottomrule
\end{longtable}
\sphinxtableafterendhook
\sphinxatlongtableend
\end{savenotes}

\sphinxstepscope


\section{Bond styles}
\label{\detokenize{Commands_bond:bond-styles}}\label{\detokenize{Commands_bond:bond}}\label{\detokenize{Commands_bond::doc}}
\sphinxAtStartPar
All LAMMPS {\hyperref[\detokenize{bond_style::doc}]{\sphinxcrossref{\DUrole{doc}{bond\_style}}}} commands.  Some styles have
accelerated versions.  This is indicated by additional letters in
parenthesis: g = GPU, i = INTEL, k = KOKKOS, o = OPENMP, t =
OPT.


\begin{savenotes}\sphinxattablestart
\sphinxthistablewithglobalstyle
\centering
\begin{tabulary}{\linewidth}[t]{TTTTT}
\sphinxtoprule
\sphinxtableatstartofbodyhook
\sphinxAtStartPar
{\hyperref[\detokenize{bond_none::doc}]{\sphinxcrossref{\DUrole{doc}{none}}}}
&
\sphinxAtStartPar
{\hyperref[\detokenize{bond_zero::doc}]{\sphinxcrossref{\DUrole{doc}{zero}}}}
&
\sphinxAtStartPar
{\hyperref[\detokenize{bond_hybrid::doc}]{\sphinxcrossref{\DUrole{doc}{hybrid (k)}}}}
&&\\
\sphinxhline&&&&\\
\sphinxhline
\sphinxAtStartPar
{\hyperref[\detokenize{bond_bpm_rotational::doc}]{\sphinxcrossref{\DUrole{doc}{bpm/rotational}}}}
&
\sphinxAtStartPar
{\hyperref[\detokenize{bond_bpm_spring::doc}]{\sphinxcrossref{\DUrole{doc}{bpm/spring}}}}
&
\sphinxAtStartPar
{\hyperref[\detokenize{bond_class2::doc}]{\sphinxcrossref{\DUrole{doc}{class2 (ko)}}}}
&
\sphinxAtStartPar
{\hyperref[\detokenize{bond_fene::doc}]{\sphinxcrossref{\DUrole{doc}{fene (iko)}}}}
&
\sphinxAtStartPar
{\hyperref[\detokenize{bond_fene_expand::doc}]{\sphinxcrossref{\DUrole{doc}{fene/expand (o)}}}}
\\
\sphinxhline
\sphinxAtStartPar
{\hyperref[\detokenize{bond_fene::doc}]{\sphinxcrossref{\DUrole{doc}{fene/nm}}}}
&
\sphinxAtStartPar
{\hyperref[\detokenize{bond_gaussian::doc}]{\sphinxcrossref{\DUrole{doc}{gaussian}}}}
&
\sphinxAtStartPar
{\hyperref[\detokenize{bond_gromos::doc}]{\sphinxcrossref{\DUrole{doc}{gromos (o)}}}}
&
\sphinxAtStartPar
{\hyperref[\detokenize{bond_harmonic::doc}]{\sphinxcrossref{\DUrole{doc}{harmonic (iko)}}}}
&
\sphinxAtStartPar
{\hyperref[\detokenize{bond_harmonic_restrain::doc}]{\sphinxcrossref{\DUrole{doc}{harmonic/restrain}}}}
\\
\sphinxhline
\sphinxAtStartPar
{\hyperref[\detokenize{bond_harmonic_shift::doc}]{\sphinxcrossref{\DUrole{doc}{harmonic/shift (o)}}}}
&
\sphinxAtStartPar
{\hyperref[\detokenize{bond_harmonic_shift_cut::doc}]{\sphinxcrossref{\DUrole{doc}{harmonic/shift/cut (o)}}}}
&
\sphinxAtStartPar
{\hyperref[\detokenize{bond_lepton::doc}]{\sphinxcrossref{\DUrole{doc}{lepton (o)}}}}
&
\sphinxAtStartPar
{\hyperref[\detokenize{bond_mesocnt::doc}]{\sphinxcrossref{\DUrole{doc}{mesocnt}}}}
&
\sphinxAtStartPar
{\hyperref[\detokenize{bond_mm3::doc}]{\sphinxcrossref{\DUrole{doc}{mm3}}}}
\\
\sphinxhline
\sphinxAtStartPar
{\hyperref[\detokenize{bond_morse::doc}]{\sphinxcrossref{\DUrole{doc}{morse (o)}}}}
&
\sphinxAtStartPar
{\hyperref[\detokenize{bond_nonlinear::doc}]{\sphinxcrossref{\DUrole{doc}{nonlinear (o)}}}}
&
\sphinxAtStartPar
{\hyperref[\detokenize{bond_oxdna::doc}]{\sphinxcrossref{\DUrole{doc}{oxdna/fene}}}}
&
\sphinxAtStartPar
{\hyperref[\detokenize{bond_oxdna::doc}]{\sphinxcrossref{\DUrole{doc}{oxdna2/fene}}}}
&
\sphinxAtStartPar
{\hyperref[\detokenize{bond_oxdna::doc}]{\sphinxcrossref{\DUrole{doc}{oxrna2/fene}}}}
\\
\sphinxhline
\sphinxAtStartPar
{\hyperref[\detokenize{bond_quartic::doc}]{\sphinxcrossref{\DUrole{doc}{quartic (o)}}}}
&
\sphinxAtStartPar
{\hyperref[\detokenize{bond_rheo_shell::doc}]{\sphinxcrossref{\DUrole{doc}{rheo/shell}}}}
&
\sphinxAtStartPar
{\hyperref[\detokenize{bond_special::doc}]{\sphinxcrossref{\DUrole{doc}{special}}}}
&
\sphinxAtStartPar
{\hyperref[\detokenize{bond_table::doc}]{\sphinxcrossref{\DUrole{doc}{table (o)}}}}
&\\
\sphinxbottomrule
\end{tabulary}
\sphinxtableafterendhook\par
\sphinxattableend\end{savenotes}


\section{Angle styles}
\label{\detokenize{Commands_bond:angle-styles}}\label{\detokenize{Commands_bond:angle}}
\sphinxAtStartPar
All LAMMPS {\hyperref[\detokenize{angle_style::doc}]{\sphinxcrossref{\DUrole{doc}{angle\_style}}}} commands.  Some styles have
accelerated versions.  This is indicated by additional letters in
parenthesis: g = GPU, i = INTEL, k = KOKKOS, o = OPENMP, t =
OPT.


\begin{savenotes}\sphinxattablestart
\sphinxthistablewithglobalstyle
\centering
\begin{tabulary}{\linewidth}[t]{TTTTT}
\sphinxtoprule
\sphinxtableatstartofbodyhook
\sphinxAtStartPar
{\hyperref[\detokenize{angle_none::doc}]{\sphinxcrossref{\DUrole{doc}{none}}}}
&
\sphinxAtStartPar
{\hyperref[\detokenize{angle_zero::doc}]{\sphinxcrossref{\DUrole{doc}{zero}}}}
&
\sphinxAtStartPar
{\hyperref[\detokenize{angle_hybrid::doc}]{\sphinxcrossref{\DUrole{doc}{hybrid (k)}}}}
&&\\
\sphinxhline&&&&\\
\sphinxhline
\sphinxAtStartPar
{\hyperref[\detokenize{angle_amoeba::doc}]{\sphinxcrossref{\DUrole{doc}{amoeba}}}}
&
\sphinxAtStartPar
{\hyperref[\detokenize{angle_charmm::doc}]{\sphinxcrossref{\DUrole{doc}{charmm (iko)}}}}
&
\sphinxAtStartPar
{\hyperref[\detokenize{angle_class2::doc}]{\sphinxcrossref{\DUrole{doc}{class2 (ko)}}}}
&
\sphinxAtStartPar
{\hyperref[\detokenize{angle_class2::doc}]{\sphinxcrossref{\DUrole{doc}{class2/p6}}}}
&
\sphinxAtStartPar
{\hyperref[\detokenize{angle_cosine::doc}]{\sphinxcrossref{\DUrole{doc}{cosine (ko)}}}}
\\
\sphinxhline
\sphinxAtStartPar
{\hyperref[\detokenize{angle_cosine_buck6d::doc}]{\sphinxcrossref{\DUrole{doc}{cosine/buck6d}}}}
&
\sphinxAtStartPar
{\hyperref[\detokenize{angle_cosine_delta::doc}]{\sphinxcrossref{\DUrole{doc}{cosine/delta (o)}}}}
&
\sphinxAtStartPar
{\hyperref[\detokenize{angle_cosine_periodic::doc}]{\sphinxcrossref{\DUrole{doc}{cosine/periodic (o)}}}}
&
\sphinxAtStartPar
{\hyperref[\detokenize{angle_cosine_shift::doc}]{\sphinxcrossref{\DUrole{doc}{cosine/shift (o)}}}}
&
\sphinxAtStartPar
{\hyperref[\detokenize{angle_cosine_shift_exp::doc}]{\sphinxcrossref{\DUrole{doc}{cosine/shift/exp (o)}}}}
\\
\sphinxhline
\sphinxAtStartPar
{\hyperref[\detokenize{angle_cosine_squared::doc}]{\sphinxcrossref{\DUrole{doc}{cosine/squared (o)}}}}
&
\sphinxAtStartPar
{\hyperref[\detokenize{angle_cosine_squared_restricted::doc}]{\sphinxcrossref{\DUrole{doc}{cosine/squared/restricted (o)}}}}
&
\sphinxAtStartPar
{\hyperref[\detokenize{angle_cross::doc}]{\sphinxcrossref{\DUrole{doc}{cross}}}}
&
\sphinxAtStartPar
{\hyperref[\detokenize{angle_dipole::doc}]{\sphinxcrossref{\DUrole{doc}{dipole (o)}}}}
&
\sphinxAtStartPar
{\hyperref[\detokenize{angle_fourier::doc}]{\sphinxcrossref{\DUrole{doc}{fourier (o)}}}}
\\
\sphinxhline
\sphinxAtStartPar
{\hyperref[\detokenize{angle_fourier_simple::doc}]{\sphinxcrossref{\DUrole{doc}{fourier/simple (o)}}}}
&
\sphinxAtStartPar
{\hyperref[\detokenize{angle_gaussian::doc}]{\sphinxcrossref{\DUrole{doc}{gaussian}}}}
&
\sphinxAtStartPar
{\hyperref[\detokenize{angle_harmonic::doc}]{\sphinxcrossref{\DUrole{doc}{harmonic (iko)}}}}
&
\sphinxAtStartPar
{\hyperref[\detokenize{angle_lepton::doc}]{\sphinxcrossref{\DUrole{doc}{lepton (o)}}}}
&
\sphinxAtStartPar
{\hyperref[\detokenize{angle_mesocnt::doc}]{\sphinxcrossref{\DUrole{doc}{mesocnt}}}}
\\
\sphinxhline
\sphinxAtStartPar
{\hyperref[\detokenize{angle_mm3::doc}]{\sphinxcrossref{\DUrole{doc}{mm3}}}}
&
\sphinxAtStartPar
{\hyperref[\detokenize{angle_quartic::doc}]{\sphinxcrossref{\DUrole{doc}{quartic (o)}}}}
&
\sphinxAtStartPar
{\hyperref[\detokenize{angle_spica::doc}]{\sphinxcrossref{\DUrole{doc}{spica (ko)}}}}
&
\sphinxAtStartPar
{\hyperref[\detokenize{angle_table::doc}]{\sphinxcrossref{\DUrole{doc}{table (o)}}}}
&\\
\sphinxbottomrule
\end{tabulary}
\sphinxtableafterendhook\par
\sphinxattableend\end{savenotes}


\section{Dihedral styles}
\label{\detokenize{Commands_bond:dihedral-styles}}\label{\detokenize{Commands_bond:dihedral}}
\sphinxAtStartPar
All LAMMPS {\hyperref[\detokenize{dihedral_style::doc}]{\sphinxcrossref{\DUrole{doc}{dihedral\_style}}}} commands.  Some styles
have accelerated versions.  This is indicated by additional letters in
parenthesis: g = GPU, i = INTEL, k = KOKKOS, o = OPENMP, t =
OPT.


\begin{savenotes}\sphinxattablestart
\sphinxthistablewithglobalstyle
\centering
\begin{tabulary}{\linewidth}[t]{TTTTT}
\sphinxtoprule
\sphinxtableatstartofbodyhook
\sphinxAtStartPar
{\hyperref[\detokenize{dihedral_none::doc}]{\sphinxcrossref{\DUrole{doc}{none}}}}
&
\sphinxAtStartPar
{\hyperref[\detokenize{dihedral_zero::doc}]{\sphinxcrossref{\DUrole{doc}{zero}}}}
&
\sphinxAtStartPar
{\hyperref[\detokenize{dihedral_hybrid::doc}]{\sphinxcrossref{\DUrole{doc}{hybrid (k)}}}}
&&\\
\sphinxhline&&&&\\
\sphinxhline
\sphinxAtStartPar
{\hyperref[\detokenize{dihedral_charmm::doc}]{\sphinxcrossref{\DUrole{doc}{charmm (iko)}}}}
&
\sphinxAtStartPar
{\hyperref[\detokenize{dihedral_charmm::doc}]{\sphinxcrossref{\DUrole{doc}{charmmfsw (k)}}}}
&
\sphinxAtStartPar
{\hyperref[\detokenize{dihedral_class2::doc}]{\sphinxcrossref{\DUrole{doc}{class2 (ko)}}}}
&
\sphinxAtStartPar
{\hyperref[\detokenize{dihedral_cosine_shift_exp::doc}]{\sphinxcrossref{\DUrole{doc}{cosine/shift/exp (o)}}}}
&
\sphinxAtStartPar
{\hyperref[\detokenize{dihedral_cosine_squared_restricted::doc}]{\sphinxcrossref{\DUrole{doc}{cosine/squared/restricted}}}}
\\
\sphinxhline
\sphinxAtStartPar
{\hyperref[\detokenize{dihedral_fourier::doc}]{\sphinxcrossref{\DUrole{doc}{fourier (io)}}}}
&
\sphinxAtStartPar
{\hyperref[\detokenize{dihedral_harmonic::doc}]{\sphinxcrossref{\DUrole{doc}{harmonic (iko)}}}}
&
\sphinxAtStartPar
{\hyperref[\detokenize{dihedral_helix::doc}]{\sphinxcrossref{\DUrole{doc}{helix (o)}}}}
&
\sphinxAtStartPar
{\hyperref[\detokenize{dihedral_lepton::doc}]{\sphinxcrossref{\DUrole{doc}{lepton (o)}}}}
&
\sphinxAtStartPar
{\hyperref[\detokenize{dihedral_multi_harmonic::doc}]{\sphinxcrossref{\DUrole{doc}{multi/harmonic (o)}}}}
\\
\sphinxhline
\sphinxAtStartPar
{\hyperref[\detokenize{dihedral_nharmonic::doc}]{\sphinxcrossref{\DUrole{doc}{nharmonic (o)}}}}
&
\sphinxAtStartPar
{\hyperref[\detokenize{dihedral_opls::doc}]{\sphinxcrossref{\DUrole{doc}{opls (iko)}}}}
&
\sphinxAtStartPar
{\hyperref[\detokenize{dihedral_quadratic::doc}]{\sphinxcrossref{\DUrole{doc}{quadratic (o)}}}}
&
\sphinxAtStartPar
{\hyperref[\detokenize{dihedral_spherical::doc}]{\sphinxcrossref{\DUrole{doc}{spherical}}}}
&
\sphinxAtStartPar
{\hyperref[\detokenize{dihedral_table::doc}]{\sphinxcrossref{\DUrole{doc}{table (o)}}}}
\\
\sphinxhline
\sphinxAtStartPar
{\hyperref[\detokenize{dihedral_table::doc}]{\sphinxcrossref{\DUrole{doc}{table/cut}}}}
&&&&\\
\sphinxbottomrule
\end{tabulary}
\sphinxtableafterendhook\par
\sphinxattableend\end{savenotes}


\section{Improper styles}
\label{\detokenize{Commands_bond:improper-styles}}\label{\detokenize{Commands_bond:improper}}
\sphinxAtStartPar
All LAMMPS {\hyperref[\detokenize{improper_style::doc}]{\sphinxcrossref{\DUrole{doc}{improper\_style}}}} commands.  Some styles
have accelerated versions.  This is indicated by additional letters in
parenthesis: g = GPU, i = INTEL, k = KOKKOS, o = OPENMP, t =
OPT.


\begin{savenotes}\sphinxattablestart
\sphinxthistablewithglobalstyle
\centering
\begin{tabulary}{\linewidth}[t]{TTTTT}
\sphinxtoprule
\sphinxtableatstartofbodyhook
\sphinxAtStartPar
{\hyperref[\detokenize{improper_none::doc}]{\sphinxcrossref{\DUrole{doc}{none}}}}
&
\sphinxAtStartPar
{\hyperref[\detokenize{improper_zero::doc}]{\sphinxcrossref{\DUrole{doc}{zero}}}}
&
\sphinxAtStartPar
{\hyperref[\detokenize{improper_hybrid::doc}]{\sphinxcrossref{\DUrole{doc}{hybrid (k)}}}}
&&\\
\sphinxhline&&&&\\
\sphinxhline
\sphinxAtStartPar
{\hyperref[\detokenize{improper_amoeba::doc}]{\sphinxcrossref{\DUrole{doc}{amoeba}}}}
&
\sphinxAtStartPar
{\hyperref[\detokenize{improper_class2::doc}]{\sphinxcrossref{\DUrole{doc}{class2 (ko)}}}}
&
\sphinxAtStartPar
{\hyperref[\detokenize{improper_cossq::doc}]{\sphinxcrossref{\DUrole{doc}{cossq (o)}}}}
&
\sphinxAtStartPar
{\hyperref[\detokenize{improper_cvff::doc}]{\sphinxcrossref{\DUrole{doc}{cvff (io)}}}}
&
\sphinxAtStartPar
{\hyperref[\detokenize{improper_distance::doc}]{\sphinxcrossref{\DUrole{doc}{distance}}}}
\\
\sphinxhline
\sphinxAtStartPar
{\hyperref[\detokenize{improper_distharm::doc}]{\sphinxcrossref{\DUrole{doc}{distharm}}}}
&
\sphinxAtStartPar
{\hyperref[\detokenize{improper_fourier::doc}]{\sphinxcrossref{\DUrole{doc}{fourier (o)}}}}
&
\sphinxAtStartPar
{\hyperref[\detokenize{improper_harmonic::doc}]{\sphinxcrossref{\DUrole{doc}{harmonic (iko)}}}}
&
\sphinxAtStartPar
{\hyperref[\detokenize{improper_inversion_harmonic::doc}]{\sphinxcrossref{\DUrole{doc}{inversion/harmonic}}}}
&
\sphinxAtStartPar
{\hyperref[\detokenize{improper_ring::doc}]{\sphinxcrossref{\DUrole{doc}{ring (o)}}}}
\\
\sphinxhline
\sphinxAtStartPar
{\hyperref[\detokenize{improper_sqdistharm::doc}]{\sphinxcrossref{\DUrole{doc}{sqdistharm}}}}
&
\sphinxAtStartPar
{\hyperref[\detokenize{improper_umbrella::doc}]{\sphinxcrossref{\DUrole{doc}{umbrella (o)}}}}
&&&\\
\sphinxbottomrule
\end{tabulary}
\sphinxtableafterendhook\par
\sphinxattableend\end{savenotes}

\sphinxstepscope


\section{KSpace styles}
\label{\detokenize{Commands_kspace:kspace-styles}}\label{\detokenize{Commands_kspace::doc}}
\sphinxAtStartPar
All LAMMPS {\hyperref[\detokenize{kspace_style::doc}]{\sphinxcrossref{\DUrole{doc}{kspace\_style}}}} solvers.  Some styles have
accelerated versions.  This is indicated by additional letters in
parenthesis: g = GPU, i = INTEL, k = KOKKOS, o = OPENMP, t =
OPT.


\begin{savenotes}\sphinxattablestart
\sphinxthistablewithglobalstyle
\centering
\begin{tabulary}{\linewidth}[t]{TTTTT}
\sphinxtoprule
\sphinxtableatstartofbodyhook
\sphinxAtStartPar
{\hyperref[\detokenize{kspace_style::doc}]{\sphinxcrossref{\DUrole{doc}{ewald (o)}}}}
&
\sphinxAtStartPar
{\hyperref[\detokenize{kspace_style::doc}]{\sphinxcrossref{\DUrole{doc}{ewald/disp}}}}
&
\sphinxAtStartPar
{\hyperref[\detokenize{kspace_style::doc}]{\sphinxcrossref{\DUrole{doc}{ewald/disp/dipole}}}}
&
\sphinxAtStartPar
{\hyperref[\detokenize{kspace_style::doc}]{\sphinxcrossref{\DUrole{doc}{ewald/dipole}}}}
&
\sphinxAtStartPar
{\hyperref[\detokenize{kspace_style::doc}]{\sphinxcrossref{\DUrole{doc}{ewald/dipole/spin}}}}
\\
\sphinxhline
\sphinxAtStartPar
{\hyperref[\detokenize{kspace_style::doc}]{\sphinxcrossref{\DUrole{doc}{ewald/electrode}}}}
&
\sphinxAtStartPar
{\hyperref[\detokenize{kspace_style::doc}]{\sphinxcrossref{\DUrole{doc}{msm (o)}}}}
&
\sphinxAtStartPar
{\hyperref[\detokenize{kspace_style::doc}]{\sphinxcrossref{\DUrole{doc}{msm/cg (o)}}}}
&
\sphinxAtStartPar
{\hyperref[\detokenize{kspace_style::doc}]{\sphinxcrossref{\DUrole{doc}{msm/dielectric}}}}
&
\sphinxAtStartPar
{\hyperref[\detokenize{kspace_style::doc}]{\sphinxcrossref{\DUrole{doc}{pppm (giko)}}}}
\\
\sphinxhline
\sphinxAtStartPar
{\hyperref[\detokenize{kspace_style::doc}]{\sphinxcrossref{\DUrole{doc}{pppm/cg (o)}}}}
&
\sphinxAtStartPar
{\hyperref[\detokenize{kspace_style::doc}]{\sphinxcrossref{\DUrole{doc}{pppm/dipole}}}}
&
\sphinxAtStartPar
{\hyperref[\detokenize{kspace_style::doc}]{\sphinxcrossref{\DUrole{doc}{pppm/dipole/spin}}}}
&
\sphinxAtStartPar
{\hyperref[\detokenize{kspace_style::doc}]{\sphinxcrossref{\DUrole{doc}{pppm/dielectric}}}}
&
\sphinxAtStartPar
{\hyperref[\detokenize{kspace_style::doc}]{\sphinxcrossref{\DUrole{doc}{pppm/disp (io)}}}}
\\
\sphinxhline
\sphinxAtStartPar
{\hyperref[\detokenize{kspace_style::doc}]{\sphinxcrossref{\DUrole{doc}{pppm/disp/tip4p (o)}}}}
&
\sphinxAtStartPar
{\hyperref[\detokenize{kspace_style::doc}]{\sphinxcrossref{\DUrole{doc}{pppm/disp/dielectric}}}}
&
\sphinxAtStartPar
{\hyperref[\detokenize{kspace_style::doc}]{\sphinxcrossref{\DUrole{doc}{pppm/stagger}}}}
&
\sphinxAtStartPar
{\hyperref[\detokenize{kspace_style::doc}]{\sphinxcrossref{\DUrole{doc}{pppm/tip4p (o)}}}}
&
\sphinxAtStartPar
{\hyperref[\detokenize{kspace_style::doc}]{\sphinxcrossref{\DUrole{doc}{pppm/dielectric}}}}
\\
\sphinxhline
\sphinxAtStartPar
{\hyperref[\detokenize{kspace_style::doc}]{\sphinxcrossref{\DUrole{doc}{pppm/electrode (i)}}}}
&
\sphinxAtStartPar
{\hyperref[\detokenize{kspace_style::doc}]{\sphinxcrossref{\DUrole{doc}{scafacos}}}}
&&&\\
\sphinxbottomrule
\end{tabulary}
\sphinxtableafterendhook\par
\sphinxattableend\end{savenotes}

\sphinxstepscope


\section{Dump styles}
\label{\detokenize{Commands_dump:dump-styles}}\label{\detokenize{Commands_dump::doc}}
\sphinxAtStartPar
An alphabetic list of all LAMMPS {\hyperref[\detokenize{dump::doc}]{\sphinxcrossref{\DUrole{doc}{dump}}}} commands.


\begin{savenotes}\sphinxattablestart
\sphinxthistablewithglobalstyle
\centering
\begin{tabulary}{\linewidth}[t]{TTTTTT}
\sphinxtoprule
\sphinxtableatstartofbodyhook
\sphinxAtStartPar
{\hyperref[\detokenize{dump::doc}]{\sphinxcrossref{\DUrole{doc}{atom}}}}
&
\sphinxAtStartPar
{\hyperref[\detokenize{dump_adios::doc}]{\sphinxcrossref{\DUrole{doc}{atom/adios}}}}
&
\sphinxAtStartPar
{\hyperref[\detokenize{dump::doc}]{\sphinxcrossref{\DUrole{doc}{atom/gz}}}}
&
\sphinxAtStartPar
{\hyperref[\detokenize{dump::doc}]{\sphinxcrossref{\DUrole{doc}{atom/zstd}}}}
&
\sphinxAtStartPar
{\hyperref[\detokenize{dump::doc}]{\sphinxcrossref{\DUrole{doc}{cfg}}}}
&
\sphinxAtStartPar
{\hyperref[\detokenize{dump::doc}]{\sphinxcrossref{\DUrole{doc}{cfg/gz}}}}
\\
\sphinxhline
\sphinxAtStartPar
{\hyperref[\detokenize{dump_cfg_uef::doc}]{\sphinxcrossref{\DUrole{doc}{cfg/uef}}}}
&
\sphinxAtStartPar
{\hyperref[\detokenize{dump::doc}]{\sphinxcrossref{\DUrole{doc}{cfg/zstd}}}}
&
\sphinxAtStartPar
{\hyperref[\detokenize{dump::doc}]{\sphinxcrossref{\DUrole{doc}{custom}}}}
&
\sphinxAtStartPar
{\hyperref[\detokenize{dump_adios::doc}]{\sphinxcrossref{\DUrole{doc}{custom/adios}}}}
&
\sphinxAtStartPar
{\hyperref[\detokenize{dump::doc}]{\sphinxcrossref{\DUrole{doc}{custom/gz}}}}
&
\sphinxAtStartPar
{\hyperref[\detokenize{dump::doc}]{\sphinxcrossref{\DUrole{doc}{custom/zstd}}}}
\\
\sphinxhline
\sphinxAtStartPar
{\hyperref[\detokenize{dump::doc}]{\sphinxcrossref{\DUrole{doc}{dcd}}}}
&
\sphinxAtStartPar
{\hyperref[\detokenize{dump::doc}]{\sphinxcrossref{\DUrole{doc}{grid}}}}
&
\sphinxAtStartPar
{\hyperref[\detokenize{dump::doc}]{\sphinxcrossref{\DUrole{doc}{grid/vtk}}}}
&
\sphinxAtStartPar
{\hyperref[\detokenize{dump_h5md::doc}]{\sphinxcrossref{\DUrole{doc}{h5md}}}}
&
\sphinxAtStartPar
{\hyperref[\detokenize{dump_image::doc}]{\sphinxcrossref{\DUrole{doc}{image}}}}
&
\sphinxAtStartPar
{\hyperref[\detokenize{dump::doc}]{\sphinxcrossref{\DUrole{doc}{local}}}}
\\
\sphinxhline
\sphinxAtStartPar
{\hyperref[\detokenize{dump::doc}]{\sphinxcrossref{\DUrole{doc}{local/gz}}}}
&
\sphinxAtStartPar
{\hyperref[\detokenize{dump::doc}]{\sphinxcrossref{\DUrole{doc}{local/zstd}}}}
&
\sphinxAtStartPar
{\hyperref[\detokenize{dump_molfile::doc}]{\sphinxcrossref{\DUrole{doc}{molfile}}}}
&
\sphinxAtStartPar
{\hyperref[\detokenize{dump_image::doc}]{\sphinxcrossref{\DUrole{doc}{movie}}}}
&
\sphinxAtStartPar
{\hyperref[\detokenize{dump_netcdf::doc}]{\sphinxcrossref{\DUrole{doc}{netcdf}}}}
&
\sphinxAtStartPar
{\hyperref[\detokenize{dump::doc}]{\sphinxcrossref{\DUrole{doc}{netcdf/mpiio}}}}
\\
\sphinxhline
\sphinxAtStartPar
{\hyperref[\detokenize{dump_vtk::doc}]{\sphinxcrossref{\DUrole{doc}{vtk}}}}
&
\sphinxAtStartPar
{\hyperref[\detokenize{dump::doc}]{\sphinxcrossref{\DUrole{doc}{xtc}}}}
&
\sphinxAtStartPar
{\hyperref[\detokenize{dump::doc}]{\sphinxcrossref{\DUrole{doc}{xyz}}}}
&
\sphinxAtStartPar
{\hyperref[\detokenize{dump::doc}]{\sphinxcrossref{\DUrole{doc}{xyz/gz}}}}
&
\sphinxAtStartPar
{\hyperref[\detokenize{dump::doc}]{\sphinxcrossref{\DUrole{doc}{xyz/zstd}}}}
&
\sphinxAtStartPar
{\hyperref[\detokenize{dump::doc}]{\sphinxcrossref{\DUrole{doc}{yaml}}}}
\\
\sphinxbottomrule
\end{tabulary}
\sphinxtableafterendhook\par
\sphinxattableend\end{savenotes}

\sphinxstepscope


\section{Removed commands and packages}
\label{\detokenize{Commands_removed:removed-commands-and-packages}}\label{\detokenize{Commands_removed::doc}}
\sphinxAtStartPar
This page lists LAMMPS commands and packages that have been removed from
the distribution and provides suggestions for alternatives or
replacements.  LAMMPS has special dummy styles implemented, that will
stop LAMMPS and print a suitable error message in most cases, when a
style/command is used that has been removed or will replace the command
with the direct alternative (if available) and print a warning.


\subsection{restart2data tool}
\label{\detokenize{Commands_removed:restart2data-tool}}
\sphinxAtStartPar
\DUrole{versionmodified,changed}{Changed in version 23Nov2013.}

\sphinxAtStartPar
The functionality of the restart2data tool has been folded into the
LAMMPS executable directly instead of having a separate tool.  A
combination of the commands {\hyperref[\detokenize{read_restart::doc}]{\sphinxcrossref{\DUrole{doc}{read\_restart}}}} and
{\hyperref[\detokenize{write_data::doc}]{\sphinxcrossref{\DUrole{doc}{write\_data}}}} can be used to the same effect.  For
added convenience this conversion can also be triggered by
{\hyperref[\detokenize{Run_options::doc}]{\sphinxcrossref{\DUrole{doc}{command line flags}}}}


\subsection{Fix ave/spatial and fix ave/spatial/sphere}
\label{\detokenize{Commands_removed:fix-ave-spatial-and-fix-ave-spatial-sphere}}
\sphinxAtStartPar
\DUrole{versionmodified,deprecated}{Deprecated since version 11Dec2015.}

\sphinxAtStartPar
The fixes ave/spatial and ave/spatial/sphere have been removed from LAMMPS
since they were superseded by the more general and extensible “chunk
infrastructure”.  Here the system is partitioned in one of many possible
ways through the {\hyperref[\detokenize{compute_chunk_atom::doc}]{\sphinxcrossref{\DUrole{doc}{compute chunk/atom}}}} command
and then averaging is done using {\hyperref[\detokenize{fix_ave_chunk::doc}]{\sphinxcrossref{\DUrole{doc}{fix ave/chunk}}}}.
Please refer to the {\hyperref[\detokenize{Howto_chunk::doc}]{\sphinxcrossref{\DUrole{doc}{chunk HOWTO}}}} section for an overview.


\subsection{Box command}
\label{\detokenize{Commands_removed:box-command}}
\sphinxAtStartPar
\DUrole{versionmodified,deprecated}{Deprecated since version 22Dec2022.}

\sphinxAtStartPar
The \sphinxstyleemphasis{box} command has been removed and the LAMMPS code changed so it won’t
be needed.  If present, LAMMPS will ignore the command and print a warning.


\subsection{Reset\_ids, reset\_atom\_ids, reset\_mol\_ids commands}
\label{\detokenize{Commands_removed:reset-ids-reset-atom-ids-reset-mol-ids-commands}}
\sphinxAtStartPar
\DUrole{versionmodified,deprecated}{Deprecated since version 22Dec2022.}

\sphinxAtStartPar
The \sphinxstyleemphasis{reset\_ids}, \sphinxstyleemphasis{reset\_atom\_ids}, and \sphinxstyleemphasis{reset\_mol\_ids} commands have
been folded into the {\hyperref[\detokenize{reset_atoms::doc}]{\sphinxcrossref{\DUrole{doc}{reset\_atoms}}}} command.  If
present, LAMMPS will replace the commands accordingly and print a
warning.


\subsection{LATTE package}
\label{\detokenize{Commands_removed:latte-package}}
\sphinxAtStartPar
\DUrole{versionmodified,deprecated}{Deprecated since version 15Jun2023.}

\sphinxAtStartPar
The LATTE package with the fix latte command was removed from LAMMPS.
This functionality has been superseded by {\hyperref[\detokenize{fix_mdi_qm::doc}]{\sphinxcrossref{\DUrole{doc}{fix mdi/qm}}}}
and {\hyperref[\detokenize{fix_mdi_qmmm::doc}]{\sphinxcrossref{\DUrole{doc}{fix mdi/qmmm}}}} from the {\hyperref[\detokenize{Packages_details:pkg-mdi}]{\sphinxcrossref{\DUrole{std,std-ref}{MDI package}}}}.  These fixes are compatible with several quantum software
packages, including LATTE.  See the \sphinxcode{\sphinxupquote{examples/QUANTUM}} dir and the
{\hyperref[\detokenize{Howto_mdi::doc}]{\sphinxcrossref{\DUrole{doc}{MDI coupling HOWTO}}}} page.  MDI supports running LAMMPS
with LATTE as a plugin library (similar to the way fix latte worked), as
well as on a different set of MPI processors.


\subsection{MEAM package}
\label{\detokenize{Commands_removed:meam-package}}
\sphinxAtStartPar
The MEAM package in Fortran has been replaced by a C++ implementation.
The code in the {\hyperref[\detokenize{Packages_details:pkg-meam}]{\sphinxcrossref{\DUrole{std,std-ref}{MEAM package}}}} is a translation of the
Fortran code of MEAM into C++, which removes several restrictions
(e.g. there can be multiple instances in hybrid pair styles) and allows
for some optimizations leading to better performance.  The pair style
{\hyperref[\detokenize{pair_meam::doc}]{\sphinxcrossref{\DUrole{doc}{meam}}}} has the exact same syntax.  For a transition
period the C++ version of MEAM was called USER\sphinxhyphen{}MEAMC so it could
coexist with the Fortran version.


\subsection{Minimize style fire/old}
\label{\detokenize{Commands_removed:minimize-style-fire-old}}
\sphinxAtStartPar
\DUrole{versionmodified,deprecated}{Deprecated since version 8Feb2023.}

\sphinxAtStartPar
Minimize style \sphinxstyleemphasis{fire/old} has been removed. Its functionality can be
reproduced with \sphinxstyleemphasis{fire} with specific options. Please see the
{\hyperref[\detokenize{min_modify::doc}]{\sphinxcrossref{\DUrole{doc}{min\_modify command}}}} documentation for details.


\subsection{Pair style mesont/tpm, compute style mesont, atom style mesont}
\label{\detokenize{Commands_removed:pair-style-mesont-tpm-compute-style-mesont-atom-style-mesont}}
\sphinxAtStartPar
\DUrole{versionmodified,deprecated}{Deprecated since version 8Feb2023.}

\sphinxAtStartPar
Pair style \sphinxstyleemphasis{mesont/tpm}, compute style \sphinxstyleemphasis{mesont}, and atom style
\sphinxstyleemphasis{mesont} have been removed from the {\hyperref[\detokenize{Packages_details:pkg-mesont}]{\sphinxcrossref{\DUrole{std,std-ref}{MESONT package}}}}.
The same functionality is available through
{\hyperref[\detokenize{pair_mesocnt::doc}]{\sphinxcrossref{\DUrole{doc}{pair style mesocnt}}}},
{\hyperref[\detokenize{bond_mesocnt::doc}]{\sphinxcrossref{\DUrole{doc}{bond style mesocnt}}}} and
{\hyperref[\detokenize{angle_mesocnt::doc}]{\sphinxcrossref{\DUrole{doc}{angle style mesocnt}}}}.


\subsection{MPIIO package}
\label{\detokenize{Commands_removed:mpiio-package}}
\sphinxAtStartPar
\DUrole{versionmodified,deprecated}{Deprecated since version 21Nov2023.}

\sphinxAtStartPar
The MPIIO package has been removed from LAMMPS since it was unmaintained
for many years and thus not updated to incorporate required changes that
had been applied to the corresponding non\sphinxhyphen{}MPIIO commands. As a
consequence the MPIIO commands had become unreliable and sometimes
crashing LAMMPS or corrupting data.  Similar functionality is available
through the {\hyperref[\detokenize{Packages_details:pkg-adios}]{\sphinxcrossref{\DUrole{std,std-ref}{ADIOS package}}}} and the {\hyperref[\detokenize{Packages_details:pkg-netcdf}]{\sphinxcrossref{\DUrole{std,std-ref}{NETCDF
package}}}}.  Also, the {\hyperref[\detokenize{dump_modify::doc}]{\sphinxcrossref{\DUrole{doc}{dump\_modify nfile or dump\_modify
fileper}}}} keywords may be used for an efficient way of
writing out dump files when running on large numbers of processors.
Similarly, the “nfile” and “fileper” keywords exist for restarts:
see {\hyperref[\detokenize{restart::doc}]{\sphinxcrossref{\DUrole{doc}{restart}}}}, {\hyperref[\detokenize{read_restart::doc}]{\sphinxcrossref{\DUrole{doc}{read\_restart}}}},
{\hyperref[\detokenize{write_restart::doc}]{\sphinxcrossref{\DUrole{doc}{write\_restart}}}}.


\subsection{MSCG package}
\label{\detokenize{Commands_removed:mscg-package}}
\sphinxAtStartPar
\DUrole{versionmodified,deprecated}{Deprecated since version 21Nov2023.}

\sphinxAtStartPar
The MSCG package has been removed from LAMMPS since it was unmaintained
for many years and instead superseded by the \sphinxhref{https://software.rcc.uchicago.edu/mscg/}{OpenMSCG software} of the Voth group at the
University of Chicago, which can be used independent from LAMMPS.


\subsection{REAX package}
\label{\detokenize{Commands_removed:reax-package}}
\sphinxAtStartPar
The REAX package has been removed since it was superseded by the
{\hyperref[\detokenize{Packages_details:pkg-reaxff}]{\sphinxcrossref{\DUrole{std,std-ref}{REAXFF package}}}}.  The REAXFF package has been tested
to yield equivalent results to the REAX package, offers better
performance, supports OpenMP multi\sphinxhyphen{}threading via OPENMP, and GPU and
threading parallelization through KOKKOS.  The new pair styles are not
syntax compatible with the removed reax pair style, so input files will
have to be adapted.  The REAXFF package was originally called
USER\sphinxhyphen{}REAXC.


\subsection{USER\sphinxhyphen{}REAXC package}
\label{\detokenize{Commands_removed:user-reaxc-package}}
\sphinxAtStartPar
\DUrole{versionmodified,deprecated}{Deprecated since version 7Feb2024.}

\sphinxAtStartPar
The USER\sphinxhyphen{}REAXC package has been renamed to {\hyperref[\detokenize{Packages_details:pkg-reaxff}]{\sphinxcrossref{\DUrole{std,std-ref}{REAXFF}}}}.
In the process also the pair style and related fixes were renamed to use
the “reaxff” string instead of “reax/c”. For a while LAMMPS was maintaining
backward compatibility by providing aliases for the styles.  These have
been removed, so using “reaxff” is now \sphinxstyleemphasis{required}.


\subsection{USER\sphinxhyphen{}CUDA package}
\label{\detokenize{Commands_removed:user-cuda-package}}
\sphinxAtStartPar
The USER\sphinxhyphen{}CUDA package had been removed, since it had been unmaintained
for a long time and had known bugs and problems.  Significant parts of
the design were transferred to the
{\hyperref[\detokenize{Packages_details:pkg-kokkos}]{\sphinxcrossref{\DUrole{std,std-ref}{KOKKOS package}}}}, which has similar
performance characteristics on NVIDIA GPUs. Both, the KOKKOS
and the {\hyperref[\detokenize{Packages_details:pkg-gpu}]{\sphinxcrossref{\DUrole{std,std-ref}{GPU package}}}} are maintained
and allow running LAMMPS with GPU acceleration.


\subsection{i\sphinxhyphen{}PI tool}
\label{\detokenize{Commands_removed:i-pi-tool}}
\sphinxAtStartPar
\DUrole{versionmodified,changed}{Changed in version 27Jun2024.}

\sphinxAtStartPar
The i\sphinxhyphen{}PI tool has been removed from the LAMMPS distribution.  Instead,
instructions to install i\sphinxhyphen{}PI from PyPI via pip are provided.


\subsection{LAMMPS shell}
\label{\detokenize{Commands_removed:lammps-shell}}
\sphinxAtStartPar
\DUrole{versionmodified,changed}{Changed in version 29Aug2024.}

\sphinxAtStartPar
The LAMMPS shell has been removed from the LAMMPS distribution. Users
are encouraged to use the {\hyperref[\detokenize{Tools:lammps-gui}]{\sphinxcrossref{\DUrole{std,std-ref}{LAMMPS\sphinxhyphen{}GUI}}}} tool instead.

\sphinxstepscope


\chapter{Optional packages}
\label{\detokenize{Packages:optional-packages}}\label{\detokenize{Packages::doc}}
\sphinxAtStartPar
This section gives an overview of the optional packages that extend
LAMMPS functionality.  Packages are groups of files that enable a
specific set of features.  For example, force fields for molecular
systems or rigid\sphinxhyphen{}body constraints are in packages.  You can see the
list of all packages and “make” commands to manage them by typing
“make package” from within the src directory of the LAMMPS
distribution.  The {\hyperref[\detokenize{Build_package::doc}]{\sphinxcrossref{\DUrole{doc}{Build package}}}} page gives
general info on how to install and uninstall packages as part of the
LAMMPS build process.

\sphinxstepscope


\section{Available Packages}
\label{\detokenize{Packages_list:available-packages}}\label{\detokenize{Packages_list::doc}}
\sphinxAtStartPar
This is the list of packages included in LAMMPS.  The link for each
package name gives more details.

\sphinxAtStartPar
Packages are supported by either the LAMMPS developers or the
contributing authors and written in a syntax and style consistent with
the rest of LAMMPS.

\sphinxAtStartPar
The “Examples” column is a subdirectory in the examples directory of the
distribution which has one or more input scripts that use the package.
E.g. \sphinxcode{\sphinxupquote{peptide}} refers to the \sphinxcode{\sphinxupquote{examples/peptide}} directory; \sphinxcode{\sphinxupquote{PACKAGES/atc}} refers
to the \sphinxcode{\sphinxupquote{examples/PACKAGES/atc}} directory.  The “Lib” column indicates\textasciigrave{}\textasciigrave{}
whether an extra library is needed to build and use the package:
\begin{itemize}
\item {} 
\sphinxAtStartPar
no  = no library

\item {} 
\sphinxAtStartPar
sys = system library: you likely have it on your machine

\item {} 
\sphinxAtStartPar
int = internal library: provided with LAMMPS, but you may need to build it

\item {} 
\sphinxAtStartPar
ext = external library: you will need to download and install it on your machine

\end{itemize}


\begin{savenotes}
\sphinxatlongtablestart
\sphinxthistablewithglobalstyle
\makeatletter
  \LTleft \@totalleftmargin plus1fill
  \LTright\dimexpr\columnwidth-\@totalleftmargin-\linewidth\relax plus1fill
\makeatother
\begin{longtable}{\X{20}{100}\X{20}{100}\X{30}{100}\X{25}{100}\X{5}{100}}
\sphinxtoprule
\sphinxstyletheadfamily 
\sphinxAtStartPar
Package
&\sphinxstyletheadfamily 
\sphinxAtStartPar
Description
&\sphinxstyletheadfamily 
\sphinxAtStartPar
Doc page
&\sphinxstyletheadfamily 
\sphinxAtStartPar
Examples
&\sphinxstyletheadfamily 
\sphinxAtStartPar
Lib
\\
\sphinxmidrule
\endfirsthead

\multicolumn{5}{c}{\sphinxnorowcolor
    \makebox[0pt]{\sphinxtablecontinued{\tablename\ \thetable{} \textendash{} continued from previous page}}%
}\\
\sphinxtoprule
\sphinxstyletheadfamily 
\sphinxAtStartPar
Package
&\sphinxstyletheadfamily 
\sphinxAtStartPar
Description
&\sphinxstyletheadfamily 
\sphinxAtStartPar
Doc page
&\sphinxstyletheadfamily 
\sphinxAtStartPar
Examples
&\sphinxstyletheadfamily 
\sphinxAtStartPar
Lib
\\
\sphinxmidrule
\endhead

\sphinxbottomrule
\multicolumn{5}{r}{\sphinxnorowcolor
    \makebox[0pt][r]{\sphinxtablecontinued{continues on next page}}%
}\\
\endfoot

\endlastfoot
\sphinxtableatstartofbodyhook

\sphinxAtStartPar
{\hyperref[\detokenize{Packages_details:pkg-adios}]{\sphinxcrossref{\DUrole{std,std-ref}{ADIOS}}}}
&
\sphinxAtStartPar
dump output via ADIOS
&
\sphinxAtStartPar
{\hyperref[\detokenize{dump_adios::doc}]{\sphinxcrossref{\DUrole{doc}{dump adios}}}}
&
\sphinxAtStartPar
\sphinxcode{\sphinxupquote{PACKAGES/adios}}
&
\sphinxAtStartPar
ext
\\
\sphinxhline
\sphinxAtStartPar
{\hyperref[\detokenize{Packages_details:pkg-amoeba}]{\sphinxcrossref{\DUrole{std,std-ref}{AMOEBA}}}}
&
\sphinxAtStartPar
AMOEBA and HIPPO force fields
&
\sphinxAtStartPar
{\hyperref[\detokenize{Howto_amoeba::doc}]{\sphinxcrossref{\DUrole{doc}{AMOEBA and HIPPO howto}}}}
&
\sphinxAtStartPar
amoeba
&
\sphinxAtStartPar
no
\\
\sphinxhline
\sphinxAtStartPar
{\hyperref[\detokenize{Packages_details:pkg-asphere}]{\sphinxcrossref{\DUrole{std,std-ref}{ASPHERE}}}}
&
\sphinxAtStartPar
aspherical particle models
&
\sphinxAtStartPar
{\hyperref[\detokenize{Howto_spherical::doc}]{\sphinxcrossref{\DUrole{doc}{Howto spherical}}}}
&
\sphinxAtStartPar
ellipse
&
\sphinxAtStartPar
no
\\
\sphinxhline
\sphinxAtStartPar
{\hyperref[\detokenize{Packages_details:pkg-atc}]{\sphinxcrossref{\DUrole{std,std-ref}{ATC}}}}
&
\sphinxAtStartPar
Atom\sphinxhyphen{}to\sphinxhyphen{}Continuum coupling
&
\sphinxAtStartPar
{\hyperref[\detokenize{fix_atc::doc}]{\sphinxcrossref{\DUrole{doc}{fix atc}}}}
&
\sphinxAtStartPar
\sphinxcode{\sphinxupquote{PACKAGES/atc}}
&
\sphinxAtStartPar
int
\\
\sphinxhline
\sphinxAtStartPar
{\hyperref[\detokenize{Packages_details:pkg-awpmd}]{\sphinxcrossref{\DUrole{std,std-ref}{AWPMD}}}}
&
\sphinxAtStartPar
wave packet MD
&
\sphinxAtStartPar
{\hyperref[\detokenize{pair_awpmd::doc}]{\sphinxcrossref{\DUrole{doc}{pair\_style awpmd/cut}}}}
&
\sphinxAtStartPar
\sphinxcode{\sphinxupquote{PACKAGES/awpmd}}
&
\sphinxAtStartPar
int
\\
\sphinxhline
\sphinxAtStartPar
{\hyperref[\detokenize{Packages_details:pkg-bocs}]{\sphinxcrossref{\DUrole{std,std-ref}{BOCS}}}}
&
\sphinxAtStartPar
BOCS bottom up coarse graining
&
\sphinxAtStartPar
{\hyperref[\detokenize{fix_bocs::doc}]{\sphinxcrossref{\DUrole{doc}{fix bocs}}}}
&
\sphinxAtStartPar
\sphinxcode{\sphinxupquote{PACKAGES/bocs}}
&
\sphinxAtStartPar
no
\\
\sphinxhline
\sphinxAtStartPar
{\hyperref[\detokenize{Packages_details:pkg-body}]{\sphinxcrossref{\DUrole{std,std-ref}{BODY}}}}
&
\sphinxAtStartPar
body\sphinxhyphen{}style particles
&
\sphinxAtStartPar
{\hyperref[\detokenize{Howto_body::doc}]{\sphinxcrossref{\DUrole{doc}{Howto body}}}}
&
\sphinxAtStartPar
body
&
\sphinxAtStartPar
no
\\
\sphinxhline
\sphinxAtStartPar
{\hyperref[\detokenize{Packages_details:pkg-bpm}]{\sphinxcrossref{\DUrole{std,std-ref}{BPM}}}}
&
\sphinxAtStartPar
bonded particle models
&
\sphinxAtStartPar
{\hyperref[\detokenize{Howto_bpm::doc}]{\sphinxcrossref{\DUrole{doc}{Howto bpm}}}}
&
\sphinxAtStartPar
bpm
&
\sphinxAtStartPar
no
\\
\sphinxhline
\sphinxAtStartPar
{\hyperref[\detokenize{Packages_details:pkg-brownian}]{\sphinxcrossref{\DUrole{std,std-ref}{BROWNIAN}}}}
&
\sphinxAtStartPar
Brownian dynamics, self\sphinxhyphen{}propelled particles
&
\sphinxAtStartPar
{\hyperref[\detokenize{fix_brownian::doc}]{\sphinxcrossref{\DUrole{doc}{fix brownian}}}}, {\hyperref[\detokenize{fix_propel_self::doc}]{\sphinxcrossref{\DUrole{doc}{fix propel/self}}}}
&
\sphinxAtStartPar
\sphinxcode{\sphinxupquote{PACKAGES/brownian}}
&
\sphinxAtStartPar
no
\\
\sphinxhline
\sphinxAtStartPar
{\hyperref[\detokenize{Packages_details:pkg-cg-dna}]{\sphinxcrossref{\DUrole{std,std-ref}{CG\sphinxhyphen{}DNA}}}}
&
\sphinxAtStartPar
coarse\sphinxhyphen{}grained DNA force fields
&
\sphinxAtStartPar
\sphinxcode{\sphinxupquote{src/CG\sphinxhyphen{}DNA/README}}
&
\sphinxAtStartPar
\sphinxcode{\sphinxupquote{PACKAGES/cgdna}}
&
\sphinxAtStartPar
no
\\
\sphinxhline
\sphinxAtStartPar
{\hyperref[\detokenize{Packages_details:pkg-cg-spica}]{\sphinxcrossref{\DUrole{std,std-ref}{CG\sphinxhyphen{}SPICA}}}}
&
\sphinxAtStartPar
SPICA (SDK) coarse\sphinxhyphen{}graining model
&
\sphinxAtStartPar
{\hyperref[\detokenize{pair_spica::doc}]{\sphinxcrossref{\DUrole{doc}{pair\_style lj/spica}}}}
&
\sphinxAtStartPar
\sphinxcode{\sphinxupquote{PACKAGES/cgspica}}
&
\sphinxAtStartPar
no
\\
\sphinxhline
\sphinxAtStartPar
{\hyperref[\detokenize{Packages_details:pkg-class2}]{\sphinxcrossref{\DUrole{std,std-ref}{CLASS2}}}}
&
\sphinxAtStartPar
class 2 force fields
&
\sphinxAtStartPar
{\hyperref[\detokenize{pair_class2::doc}]{\sphinxcrossref{\DUrole{doc}{pair\_style lj/class2}}}}
&
\sphinxAtStartPar
n/a
&
\sphinxAtStartPar
no
\\
\sphinxhline
\sphinxAtStartPar
{\hyperref[\detokenize{Packages_details:pkg-colloid}]{\sphinxcrossref{\DUrole{std,std-ref}{COLLOID}}}}
&
\sphinxAtStartPar
colloidal particles
&
\sphinxAtStartPar
{\hyperref[\detokenize{atom_style::doc}]{\sphinxcrossref{\DUrole{doc}{atom\_style colloid}}}}
&
\sphinxAtStartPar
colloid
&
\sphinxAtStartPar
no
\\
\sphinxhline
\sphinxAtStartPar
{\hyperref[\detokenize{Packages_details:pkg-colvars}]{\sphinxcrossref{\DUrole{std,std-ref}{COLVARS}}}}
&
\sphinxAtStartPar
\sphinxhref{https://colvars.github.io/}{Colvars collective variables library}
&
\sphinxAtStartPar
{\hyperref[\detokenize{fix_colvars::doc}]{\sphinxcrossref{\DUrole{doc}{fix colvars}}}}
&
\sphinxAtStartPar
\sphinxcode{\sphinxupquote{PACKAGES/colvars}}
&
\sphinxAtStartPar
int
\\
\sphinxhline
\sphinxAtStartPar
{\hyperref[\detokenize{Packages_details:pkg-compress}]{\sphinxcrossref{\DUrole{std,std-ref}{COMPRESS}}}}
&
\sphinxAtStartPar
I/O compression
&
\sphinxAtStartPar
{\hyperref[\detokenize{dump::doc}]{\sphinxcrossref{\DUrole{doc}{dump */gz}}}}
&
\sphinxAtStartPar
n/a
&
\sphinxAtStartPar
sys
\\
\sphinxhline
\sphinxAtStartPar
{\hyperref[\detokenize{Packages_details:pkg-coreshell}]{\sphinxcrossref{\DUrole{std,std-ref}{CORESHELL}}}}
&
\sphinxAtStartPar
adiabatic core/shell model
&
\sphinxAtStartPar
{\hyperref[\detokenize{Howto_coreshell::doc}]{\sphinxcrossref{\DUrole{doc}{Howto coreshell}}}}
&
\sphinxAtStartPar
coreshell
&
\sphinxAtStartPar
no
\\
\sphinxhline
\sphinxAtStartPar
{\hyperref[\detokenize{Packages_details:pkg-dielectric}]{\sphinxcrossref{\DUrole{std,std-ref}{DIELECTRIC}}}}
&
\sphinxAtStartPar
dielectric boundary solvers and force styles
&
\sphinxAtStartPar
{\hyperref[\detokenize{compute_efield_atom::doc}]{\sphinxcrossref{\DUrole{doc}{compute efield/atom}}}}
&
\sphinxAtStartPar
\sphinxcode{\sphinxupquote{PACKAGES/dielectric}}
&
\sphinxAtStartPar
no
\\
\sphinxhline
\sphinxAtStartPar
{\hyperref[\detokenize{Packages_details:pkg-diffraction}]{\sphinxcrossref{\DUrole{std,std-ref}{DIFFRACTION}}}}
&
\sphinxAtStartPar
virtual x\sphinxhyphen{}ray and electron diffraction
&
\sphinxAtStartPar
{\hyperref[\detokenize{compute_xrd::doc}]{\sphinxcrossref{\DUrole{doc}{compute xrd}}}}
&
\sphinxAtStartPar
\sphinxcode{\sphinxupquote{PACKAGES/diffraction}}
&
\sphinxAtStartPar
no
\\
\sphinxhline
\sphinxAtStartPar
{\hyperref[\detokenize{Packages_details:pkg-dipole}]{\sphinxcrossref{\DUrole{std,std-ref}{DIPOLE}}}}
&
\sphinxAtStartPar
point dipole particles
&
\sphinxAtStartPar
{\hyperref[\detokenize{pair_dipole::doc}]{\sphinxcrossref{\DUrole{doc}{pair\_style lj/…/dipole}}}}
&
\sphinxAtStartPar
dipole
&
\sphinxAtStartPar
no
\\
\sphinxhline
\sphinxAtStartPar
{\hyperref[\detokenize{Packages_details:pkg-dpd-basic}]{\sphinxcrossref{\DUrole{std,std-ref}{DPD\sphinxhyphen{}BASIC}}}}
&
\sphinxAtStartPar
basic DPD models
&
\sphinxAtStartPar
{\hyperref[\detokenize{pair_dpd::doc}]{\sphinxcrossref{\DUrole{doc}{pair\_styles dpd}}}} {\hyperref[\detokenize{pair_dpd_ext::doc}]{\sphinxcrossref{\DUrole{doc}{dpd/ext}}}}
&
\sphinxAtStartPar
\sphinxcode{\sphinxupquote{PACKAGES/dpd\sphinxhyphen{}basic}}
&
\sphinxAtStartPar
no
\\
\sphinxhline
\sphinxAtStartPar
{\hyperref[\detokenize{Packages_details:pkg-dpd-meso}]{\sphinxcrossref{\DUrole{std,std-ref}{DPD\sphinxhyphen{}MESO}}}}
&
\sphinxAtStartPar
mesoscale DPD models
&
\sphinxAtStartPar
{\hyperref[\detokenize{pair_mesodpd::doc}]{\sphinxcrossref{\DUrole{doc}{pair\_style edpd}}}}
&
\sphinxAtStartPar
\sphinxcode{\sphinxupquote{PACKAGES/dpd\sphinxhyphen{}meso}}
&
\sphinxAtStartPar
no
\\
\sphinxhline
\sphinxAtStartPar
{\hyperref[\detokenize{Packages_details:pkg-dpd-react}]{\sphinxcrossref{\DUrole{std,std-ref}{DPD\sphinxhyphen{}REACT}}}}
&
\sphinxAtStartPar
reactive dissipative particle dynamics
&
\sphinxAtStartPar
\sphinxcode{\sphinxupquote{src/DPD\sphinxhyphen{}REACT/README}}
&
\sphinxAtStartPar
\sphinxcode{\sphinxupquote{PACKAGES/dpd\sphinxhyphen{}react}}
&
\sphinxAtStartPar
no
\\
\sphinxhline
\sphinxAtStartPar
{\hyperref[\detokenize{Packages_details:pkg-dpd-smooth}]{\sphinxcrossref{\DUrole{std,std-ref}{DPD\sphinxhyphen{}SMOOTH}}}}
&
\sphinxAtStartPar
smoothed dissipative particle dynamics
&
\sphinxAtStartPar
\sphinxcode{\sphinxupquote{src/DPD\sphinxhyphen{}SMOOTH/README}}
&
\sphinxAtStartPar
\sphinxcode{\sphinxupquote{PACKAGES/dpd\sphinxhyphen{}smooth}}
&
\sphinxAtStartPar
no
\\
\sphinxhline
\sphinxAtStartPar
{\hyperref[\detokenize{Packages_details:pkg-drude}]{\sphinxcrossref{\DUrole{std,std-ref}{DRUDE}}}}
&
\sphinxAtStartPar
Drude oscillators
&
\sphinxAtStartPar
{\hyperref[\detokenize{Howto_drude::doc}]{\sphinxcrossref{\DUrole{doc}{Howto drude}}}}
&
\sphinxAtStartPar
\sphinxcode{\sphinxupquote{PACKAGES/drude}}
&
\sphinxAtStartPar
no
\\
\sphinxhline
\sphinxAtStartPar
{\hyperref[\detokenize{Packages_details:pkg-eff}]{\sphinxcrossref{\DUrole{std,std-ref}{EFF}}}}
&
\sphinxAtStartPar
electron force field
&
\sphinxAtStartPar
{\hyperref[\detokenize{pair_eff::doc}]{\sphinxcrossref{\DUrole{doc}{pair\_style eff/cut}}}}
&
\sphinxAtStartPar
\sphinxcode{\sphinxupquote{PACKAGES/eff}}
&
\sphinxAtStartPar
no
\\
\sphinxhline
\sphinxAtStartPar
{\hyperref[\detokenize{Packages_details:pkg-electrode}]{\sphinxcrossref{\DUrole{std,std-ref}{ELECTRODE}}}}
&
\sphinxAtStartPar
electrode charges to match potential
&
\sphinxAtStartPar
{\hyperref[\detokenize{fix_electrode::doc}]{\sphinxcrossref{\DUrole{doc}{fix electrode/conp}}}}
&
\sphinxAtStartPar
\sphinxcode{\sphinxupquote{PACKAGES/electrode}}
&
\sphinxAtStartPar
no
\\
\sphinxhline
\sphinxAtStartPar
{\hyperref[\detokenize{Packages_details:pkg-extra-command}]{\sphinxcrossref{\DUrole{std,std-ref}{EXTRA\sphinxhyphen{}COMMAND}}}}
&
\sphinxAtStartPar
additional command styles
&
\sphinxAtStartPar
{\hyperref[\detokenize{Commands_all::doc}]{\sphinxcrossref{\DUrole{doc}{general commands}}}}
&
\sphinxAtStartPar
n/a
&
\sphinxAtStartPar
no
\\
\sphinxhline
\sphinxAtStartPar
{\hyperref[\detokenize{Packages_details:pkg-extra-compute}]{\sphinxcrossref{\DUrole{std,std-ref}{EXTRA\sphinxhyphen{}COMPUTE}}}}
&
\sphinxAtStartPar
additional compute styles
&
\sphinxAtStartPar
{\hyperref[\detokenize{compute::doc}]{\sphinxcrossref{\DUrole{doc}{compute}}}}
&
\sphinxAtStartPar
n/a
&
\sphinxAtStartPar
no
\\
\sphinxhline
\sphinxAtStartPar
{\hyperref[\detokenize{Packages_details:pkg-extra-dump}]{\sphinxcrossref{\DUrole{std,std-ref}{EXTRA\sphinxhyphen{}DUMP}}}}
&
\sphinxAtStartPar
additional dump styles
&
\sphinxAtStartPar
{\hyperref[\detokenize{dump::doc}]{\sphinxcrossref{\DUrole{doc}{dump}}}}
&
\sphinxAtStartPar
n/a
&
\sphinxAtStartPar
no
\\
\sphinxhline
\sphinxAtStartPar
{\hyperref[\detokenize{Packages_details:pkg-extra-fix}]{\sphinxcrossref{\DUrole{std,std-ref}{EXTRA\sphinxhyphen{}FIX}}}}
&
\sphinxAtStartPar
additional fix styles
&
\sphinxAtStartPar
{\hyperref[\detokenize{fix::doc}]{\sphinxcrossref{\DUrole{doc}{fix}}}}
&
\sphinxAtStartPar
n/a
&
\sphinxAtStartPar
no
\\
\sphinxhline
\sphinxAtStartPar
{\hyperref[\detokenize{Packages_details:pkg-extra-molecule}]{\sphinxcrossref{\DUrole{std,std-ref}{EXTRA\sphinxhyphen{}MOLECULE}}}}
&
\sphinxAtStartPar
additional molecular styles
&
\sphinxAtStartPar
{\hyperref[\detokenize{Commands_bond::doc}]{\sphinxcrossref{\DUrole{doc}{molecular styles}}}}
&
\sphinxAtStartPar
n/a
&
\sphinxAtStartPar
no
\\
\sphinxhline
\sphinxAtStartPar
{\hyperref[\detokenize{Packages_details:pkg-extra-pair}]{\sphinxcrossref{\DUrole{std,std-ref}{EXTRA\sphinxhyphen{}PAIR}}}}
&
\sphinxAtStartPar
additional pair styles
&
\sphinxAtStartPar
{\hyperref[\detokenize{pair_style::doc}]{\sphinxcrossref{\DUrole{doc}{pair\_style}}}}
&
\sphinxAtStartPar
n/a
&
\sphinxAtStartPar
no
\\
\sphinxhline
\sphinxAtStartPar
{\hyperref[\detokenize{Packages_details:pkg-fep}]{\sphinxcrossref{\DUrole{std,std-ref}{FEP}}}}
&
\sphinxAtStartPar
free energy perturbation
&
\sphinxAtStartPar
{\hyperref[\detokenize{compute_fep::doc}]{\sphinxcrossref{\DUrole{doc}{compute fep}}}}
&
\sphinxAtStartPar
\sphinxcode{\sphinxupquote{PACKAGES/fep}}
&
\sphinxAtStartPar
no
\\
\sphinxhline
\sphinxAtStartPar
{\hyperref[\detokenize{Packages_details:pkg-gpu}]{\sphinxcrossref{\DUrole{std,std-ref}{GPU}}}}
&
\sphinxAtStartPar
GPU\sphinxhyphen{}enabled styles
&
\sphinxAtStartPar
{\hyperref[\detokenize{Speed_gpu::doc}]{\sphinxcrossref{\DUrole{doc}{Section gpu}}}}
&
\sphinxAtStartPar
\sphinxhref{https://www.lammps.org/bench.html}{Benchmarks}
&
\sphinxAtStartPar
int
\\
\sphinxhline
\sphinxAtStartPar
{\hyperref[\detokenize{Packages_details:pkg-granular}]{\sphinxcrossref{\DUrole{std,std-ref}{GRANULAR}}}}
&
\sphinxAtStartPar
granular systems
&
\sphinxAtStartPar
{\hyperref[\detokenize{Howto_granular::doc}]{\sphinxcrossref{\DUrole{doc}{Howto granular}}}}
&
\sphinxAtStartPar
pour
&
\sphinxAtStartPar
no
\\
\sphinxhline
\sphinxAtStartPar
{\hyperref[\detokenize{Packages_details:pkg-h5md}]{\sphinxcrossref{\DUrole{std,std-ref}{H5MD}}}}
&
\sphinxAtStartPar
dump output via HDF5
&
\sphinxAtStartPar
{\hyperref[\detokenize{dump_h5md::doc}]{\sphinxcrossref{\DUrole{doc}{dump h5md}}}}
&
\sphinxAtStartPar
n/a
&
\sphinxAtStartPar
ext
\\
\sphinxhline
\sphinxAtStartPar
{\hyperref[\detokenize{Packages_details:pkg-intel}]{\sphinxcrossref{\DUrole{std,std-ref}{INTEL}}}}
&
\sphinxAtStartPar
optimized Intel CPU and KNL styles
&
\sphinxAtStartPar
{\hyperref[\detokenize{Speed_intel::doc}]{\sphinxcrossref{\DUrole{doc}{Speed intel}}}}
&
\sphinxAtStartPar
\sphinxhref{https://www.lammps.org/bench.html}{Benchmarks}
&
\sphinxAtStartPar
no
\\
\sphinxhline
\sphinxAtStartPar
{\hyperref[\detokenize{Packages_details:pkg-interlayer}]{\sphinxcrossref{\DUrole{std,std-ref}{INTERLAYER}}}}
&
\sphinxAtStartPar
Inter\sphinxhyphen{}layer pair potentials
&
\sphinxAtStartPar
{\hyperref[\detokenize{Commands_pair::doc}]{\sphinxcrossref{\DUrole{doc}{several pair styles}}}}
&
\sphinxAtStartPar
\sphinxcode{\sphinxupquote{PACKAGES/interlayer}}
&
\sphinxAtStartPar
no
\\
\sphinxhline
\sphinxAtStartPar
{\hyperref[\detokenize{Packages_details:pkg-kim}]{\sphinxcrossref{\DUrole{std,std-ref}{KIM}}}}
&
\sphinxAtStartPar
OpenKIM wrapper
&
\sphinxAtStartPar
{\hyperref[\detokenize{pair_kim::doc}]{\sphinxcrossref{\DUrole{doc}{pair\_style kim}}}}
&
\sphinxAtStartPar
kim
&
\sphinxAtStartPar
ext
\\
\sphinxhline
\sphinxAtStartPar
{\hyperref[\detokenize{Packages_details:pkg-kokkos}]{\sphinxcrossref{\DUrole{std,std-ref}{KOKKOS}}}}
&
\sphinxAtStartPar
Kokkos\sphinxhyphen{}enabled styles
&
\sphinxAtStartPar
{\hyperref[\detokenize{Speed_kokkos::doc}]{\sphinxcrossref{\DUrole{doc}{Speed kokkos}}}}
&
\sphinxAtStartPar
\sphinxhref{https://www.lammps.org/bench.html}{Benchmarks}
&
\sphinxAtStartPar
no
\\
\sphinxhline
\sphinxAtStartPar
{\hyperref[\detokenize{Packages_details:pkg-kspace}]{\sphinxcrossref{\DUrole{std,std-ref}{KSPACE}}}}
&
\sphinxAtStartPar
long\sphinxhyphen{}range Coulombic solvers
&
\sphinxAtStartPar
{\hyperref[\detokenize{kspace_style::doc}]{\sphinxcrossref{\DUrole{doc}{kspace\_style}}}}
&
\sphinxAtStartPar
peptide
&
\sphinxAtStartPar
no
\\
\sphinxhline
\sphinxAtStartPar
{\hyperref[\detokenize{Packages_details:pkg-latboltz}]{\sphinxcrossref{\DUrole{std,std-ref}{LATBOLTZ}}}}
&
\sphinxAtStartPar
Lattice Boltzmann fluid
&
\sphinxAtStartPar
{\hyperref[\detokenize{fix_lb_fluid::doc}]{\sphinxcrossref{\DUrole{doc}{fix lb/fluid}}}}
&
\sphinxAtStartPar
\sphinxcode{\sphinxupquote{PACKAGES/latboltz}}
&
\sphinxAtStartPar
no
\\
\sphinxhline
\sphinxAtStartPar
{\hyperref[\detokenize{Packages_details:pkg-lepton}]{\sphinxcrossref{\DUrole{std,std-ref}{LEPTON}}}}
&
\sphinxAtStartPar
evaluate strings as potential function
&
\sphinxAtStartPar
{\hyperref[\detokenize{pair_lepton::doc}]{\sphinxcrossref{\DUrole{doc}{pair\_style lepton}}}}
&
\sphinxAtStartPar
\sphinxcode{\sphinxupquote{PACKAGES/lepton}}
&
\sphinxAtStartPar
int
\\
\sphinxhline
\sphinxAtStartPar
{\hyperref[\detokenize{Packages_details:pkg-machdyn}]{\sphinxcrossref{\DUrole{std,std-ref}{MACHDYN}}}}
&
\sphinxAtStartPar
smoothed Mach dynamics
&
\sphinxAtStartPar
\sphinxhref{PDF/MACHDYN\_LAMMPS\_userguide.pdf}{SMD User Guide}
&
\sphinxAtStartPar
\sphinxcode{\sphinxupquote{PACKAGES/machdyn}}
&
\sphinxAtStartPar
ext
\\
\sphinxhline
\sphinxAtStartPar
{\hyperref[\detokenize{Packages_details:pkg-manifold}]{\sphinxcrossref{\DUrole{std,std-ref}{MANIFOLD}}}}
&
\sphinxAtStartPar
motion on 2d surfaces
&
\sphinxAtStartPar
{\hyperref[\detokenize{fix_manifoldforce::doc}]{\sphinxcrossref{\DUrole{doc}{fix manifoldforce}}}}
&
\sphinxAtStartPar
\sphinxcode{\sphinxupquote{PACKAGES/manifold}}
&
\sphinxAtStartPar
no
\\
\sphinxhline
\sphinxAtStartPar
{\hyperref[\detokenize{Packages_details:pkg-manybody}]{\sphinxcrossref{\DUrole{std,std-ref}{MANYBODY}}}}
&
\sphinxAtStartPar
many\sphinxhyphen{}body potentials
&
\sphinxAtStartPar
{\hyperref[\detokenize{pair_tersoff::doc}]{\sphinxcrossref{\DUrole{doc}{pair\_style tersoff}}}}
&
\sphinxAtStartPar
shear
&
\sphinxAtStartPar
no
\\
\sphinxhline
\sphinxAtStartPar
{\hyperref[\detokenize{Packages_details:pkg-mc}]{\sphinxcrossref{\DUrole{std,std-ref}{MC}}}}
&
\sphinxAtStartPar
Monte Carlo options
&
\sphinxAtStartPar
{\hyperref[\detokenize{fix_gcmc::doc}]{\sphinxcrossref{\DUrole{doc}{fix gcmc}}}}
&
\sphinxAtStartPar
n/a
&
\sphinxAtStartPar
no
\\
\sphinxhline
\sphinxAtStartPar
{\hyperref[\detokenize{Packages_details:pkg-mdi}]{\sphinxcrossref{\DUrole{std,std-ref}{MDI}}}}
&
\sphinxAtStartPar
client\sphinxhyphen{}server code coupling
&
\sphinxAtStartPar
{\hyperref[\detokenize{Howto_mdi::doc}]{\sphinxcrossref{\DUrole{doc}{MDI Howto}}}}
&
\sphinxAtStartPar
\sphinxcode{\sphinxupquote{PACKAGES/mdi}}
&
\sphinxAtStartPar
ext
\\
\sphinxhline
\sphinxAtStartPar
{\hyperref[\detokenize{Packages_details:pkg-meam}]{\sphinxcrossref{\DUrole{std,std-ref}{MEAM}}}}
&
\sphinxAtStartPar
modified EAM potential (C++)
&
\sphinxAtStartPar
{\hyperref[\detokenize{pair_meam::doc}]{\sphinxcrossref{\DUrole{doc}{pair\_style meam}}}}
&
\sphinxAtStartPar
meam
&
\sphinxAtStartPar
no
\\
\sphinxhline
\sphinxAtStartPar
{\hyperref[\detokenize{Packages_details:pkg-mesont}]{\sphinxcrossref{\DUrole{std,std-ref}{MESONT}}}}
&
\sphinxAtStartPar
mesoscopic tubular potential model
&
\sphinxAtStartPar
pair styles {\hyperref[\detokenize{pair_mesocnt::doc}]{\sphinxcrossref{\DUrole{doc}{mesocnt}}}}
&
\sphinxAtStartPar
\sphinxcode{\sphinxupquote{PACKAGES/mesont}}
&
\sphinxAtStartPar
no
\\
\sphinxhline
\sphinxAtStartPar
{\hyperref[\detokenize{Packages_details:pkg-mgpt}]{\sphinxcrossref{\DUrole{std,std-ref}{MGPT}}}}
&
\sphinxAtStartPar
fast MGPT multi\sphinxhyphen{}ion potentials
&
\sphinxAtStartPar
{\hyperref[\detokenize{pair_mgpt::doc}]{\sphinxcrossref{\DUrole{doc}{pair\_style mgpt}}}}
&
\sphinxAtStartPar
\sphinxcode{\sphinxupquote{PACKAGES/mgpt}}
&
\sphinxAtStartPar
no
\\
\sphinxhline
\sphinxAtStartPar
{\hyperref[\detokenize{Packages_details:pkg-misc}]{\sphinxcrossref{\DUrole{std,std-ref}{MISC}}}}
&
\sphinxAtStartPar
miscellaneous single\sphinxhyphen{}file commands
&
\sphinxAtStartPar
n/a
&
\sphinxAtStartPar
no
&
\sphinxAtStartPar
no
\\
\sphinxhline
\sphinxAtStartPar
{\hyperref[\detokenize{Packages_details:pkg-ml-hdnnp}]{\sphinxcrossref{\DUrole{std,std-ref}{ML\sphinxhyphen{}HDNNP}}}}
&
\sphinxAtStartPar
High\sphinxhyphen{}dimensional neural network potentials
&
\sphinxAtStartPar
{\hyperref[\detokenize{pair_hdnnp::doc}]{\sphinxcrossref{\DUrole{doc}{pair\_style hdnnp}}}}
&
\sphinxAtStartPar
\sphinxcode{\sphinxupquote{PACKAGES/hdnnp}}
&
\sphinxAtStartPar
ext
\\
\sphinxhline
\sphinxAtStartPar
{\hyperref[\detokenize{Packages_details:pkg-ml-iap}]{\sphinxcrossref{\DUrole{std,std-ref}{ML\sphinxhyphen{}IAP}}}}
&
\sphinxAtStartPar
multiple machine learning potentials
&
\sphinxAtStartPar
{\hyperref[\detokenize{pair_mliap::doc}]{\sphinxcrossref{\DUrole{doc}{pair\_style mliap}}}}
&
\sphinxAtStartPar
mliap
&
\sphinxAtStartPar
no
\\
\sphinxhline
\sphinxAtStartPar
{\hyperref[\detokenize{Packages_details:pkg-ml-pace}]{\sphinxcrossref{\DUrole{std,std-ref}{ML\sphinxhyphen{}PACE}}}}
&
\sphinxAtStartPar
Atomic Cluster Expansion potential
&
\sphinxAtStartPar
{\hyperref[\detokenize{pair_pace::doc}]{\sphinxcrossref{\DUrole{doc}{pair pace}}}}
&
\sphinxAtStartPar
\sphinxcode{\sphinxupquote{PACKAGES/pace}}
&
\sphinxAtStartPar
ext
\\
\sphinxhline
\sphinxAtStartPar
{\hyperref[\detokenize{Packages_details:pkg-ml-pod}]{\sphinxcrossref{\DUrole{std,std-ref}{ML\sphinxhyphen{}POD}}}}
&
\sphinxAtStartPar
Proper orthogonal decomposition potentials
&
\sphinxAtStartPar
{\hyperref[\detokenize{pair_pod::doc}]{\sphinxcrossref{\DUrole{doc}{pair pod}}}}
&
\sphinxAtStartPar
pod
&
\sphinxAtStartPar
ext
\\
\sphinxhline
\sphinxAtStartPar
{\hyperref[\detokenize{Packages_details:pkg-ml-quip}]{\sphinxcrossref{\DUrole{std,std-ref}{ML\sphinxhyphen{}QUIP}}}}
&
\sphinxAtStartPar
QUIP/libatoms interface
&
\sphinxAtStartPar
{\hyperref[\detokenize{pair_quip::doc}]{\sphinxcrossref{\DUrole{doc}{pair\_style quip}}}}
&
\sphinxAtStartPar
\sphinxcode{\sphinxupquote{PACKAGES/quip}}
&
\sphinxAtStartPar
ext
\\
\sphinxhline
\sphinxAtStartPar
{\hyperref[\detokenize{Packages_details:pkg-ml-rann}]{\sphinxcrossref{\DUrole{std,std-ref}{ML\sphinxhyphen{}RANN}}}}
&
\sphinxAtStartPar
Pair style for RANN potentials
&
\sphinxAtStartPar
{\hyperref[\detokenize{pair_rann::doc}]{\sphinxcrossref{\DUrole{doc}{pair rann}}}}
&
\sphinxAtStartPar
\sphinxcode{\sphinxupquote{PACKAGES/rann}}
&
\sphinxAtStartPar
no
\\
\sphinxhline
\sphinxAtStartPar
{\hyperref[\detokenize{Packages_details:pkg-ml-snap}]{\sphinxcrossref{\DUrole{std,std-ref}{ML\sphinxhyphen{}SNAP}}}}
&
\sphinxAtStartPar
quantum\sphinxhyphen{}fitted potential
&
\sphinxAtStartPar
{\hyperref[\detokenize{pair_snap::doc}]{\sphinxcrossref{\DUrole{doc}{pair\_style snap}}}}
&
\sphinxAtStartPar
snap
&
\sphinxAtStartPar
no
\\
\sphinxhline
\sphinxAtStartPar
{\hyperref[\detokenize{Packages_details:pkg-ml-uf3}]{\sphinxcrossref{\DUrole{std,std-ref}{ML\sphinxhyphen{}UF3}}}}
&
\sphinxAtStartPar
quantum\sphinxhyphen{}fitted ultra fast potentials
&
\sphinxAtStartPar
{\hyperref[\detokenize{pair_uf3::doc}]{\sphinxcrossref{\DUrole{doc}{pair\_style uf3}}}}
&
\sphinxAtStartPar
\sphinxcode{\sphinxupquote{PACKAGES/uf3}}
&
\sphinxAtStartPar
no
\\
\sphinxhline
\sphinxAtStartPar
{\hyperref[\detokenize{Packages_details:pkg-mofff}]{\sphinxcrossref{\DUrole{std,std-ref}{MOFFF}}}}
&
\sphinxAtStartPar
styles for \sphinxhref{https://www.mofplus.org/content/show/MOF-FF}{MOF\sphinxhyphen{}FF} force field
&
\sphinxAtStartPar
{\hyperref[\detokenize{pair_buck6d_coul_gauss::doc}]{\sphinxcrossref{\DUrole{doc}{pair\_style buck6d/coul/gauss}}}}
&
\sphinxAtStartPar
\sphinxcode{\sphinxupquote{PACKAGES/mofff}}
&
\sphinxAtStartPar
no
\\
\sphinxhline
\sphinxAtStartPar
{\hyperref[\detokenize{Packages_details:pkg-molecule}]{\sphinxcrossref{\DUrole{std,std-ref}{MOLECULE}}}}
&
\sphinxAtStartPar
molecular system force fields
&
\sphinxAtStartPar
{\hyperref[\detokenize{Howto_bioFF::doc}]{\sphinxcrossref{\DUrole{doc}{Howto bioFF}}}}
&
\sphinxAtStartPar
peptide
&
\sphinxAtStartPar
no
\\
\sphinxhline
\sphinxAtStartPar
{\hyperref[\detokenize{Packages_details:pkg-molfile}]{\sphinxcrossref{\DUrole{std,std-ref}{MOLFILE}}}}
&
\sphinxAtStartPar
\sphinxhref{https://www.ks.uiuc.edu/Research/vmd/}{VMD} molfile plug\sphinxhyphen{}ins
&
\sphinxAtStartPar
{\hyperref[\detokenize{dump_molfile::doc}]{\sphinxcrossref{\DUrole{doc}{dump molfile}}}}
&
\sphinxAtStartPar
n/a
&
\sphinxAtStartPar
ext
\\
\sphinxhline
\sphinxAtStartPar
{\hyperref[\detokenize{Packages_details:pkg-netcdf}]{\sphinxcrossref{\DUrole{std,std-ref}{NETCDF}}}}
&
\sphinxAtStartPar
dump output via NetCDF
&
\sphinxAtStartPar
{\hyperref[\detokenize{dump_netcdf::doc}]{\sphinxcrossref{\DUrole{doc}{dump netcdf}}}}
&
\sphinxAtStartPar
n/a
&
\sphinxAtStartPar
ext
\\
\sphinxhline
\sphinxAtStartPar
{\hyperref[\detokenize{Packages_details:pkg-openmp}]{\sphinxcrossref{\DUrole{std,std-ref}{OPENMP}}}}
&
\sphinxAtStartPar
OpenMP\sphinxhyphen{}enabled styles
&
\sphinxAtStartPar
{\hyperref[\detokenize{Speed_omp::doc}]{\sphinxcrossref{\DUrole{doc}{Speed omp}}}}
&
\sphinxAtStartPar
\sphinxhref{https://www.lammps.org/bench.html}{Benchmarks}
&
\sphinxAtStartPar
no
\\
\sphinxhline
\sphinxAtStartPar
{\hyperref[\detokenize{Packages_details:pkg-opt}]{\sphinxcrossref{\DUrole{std,std-ref}{OPT}}}}
&
\sphinxAtStartPar
optimized pair styles
&
\sphinxAtStartPar
{\hyperref[\detokenize{Speed_opt::doc}]{\sphinxcrossref{\DUrole{doc}{Speed opt}}}}
&
\sphinxAtStartPar
\sphinxhref{https://www.lammps.org/bench.html}{Benchmarks}
&
\sphinxAtStartPar
no
\\
\sphinxhline
\sphinxAtStartPar
{\hyperref[\detokenize{Packages_details:pkg-orient}]{\sphinxcrossref{\DUrole{std,std-ref}{ORIENT}}}}
&
\sphinxAtStartPar
fixes for orientation depended forces
&
\sphinxAtStartPar
{\hyperref[\detokenize{fix_orient::doc}]{\sphinxcrossref{\DUrole{doc}{fix orient/*}}}}
&
\sphinxAtStartPar
\sphinxcode{\sphinxupquote{PACKAGES/orient\_eco}}
&
\sphinxAtStartPar
no
\\
\sphinxhline
\sphinxAtStartPar
{\hyperref[\detokenize{Packages_details:pkg-peri}]{\sphinxcrossref{\DUrole{std,std-ref}{PERI}}}}
&
\sphinxAtStartPar
Peridynamics models
&
\sphinxAtStartPar
{\hyperref[\detokenize{pair_peri::doc}]{\sphinxcrossref{\DUrole{doc}{pair\_style peri}}}}
&
\sphinxAtStartPar
peri
&
\sphinxAtStartPar
no
\\
\sphinxhline
\sphinxAtStartPar
{\hyperref[\detokenize{Packages_details:pkg-phonon}]{\sphinxcrossref{\DUrole{std,std-ref}{PHONON}}}}
&
\sphinxAtStartPar
phonon dynamical matrix
&
\sphinxAtStartPar
{\hyperref[\detokenize{fix_phonon::doc}]{\sphinxcrossref{\DUrole{doc}{fix phonon}}}}
&
\sphinxAtStartPar
\sphinxcode{\sphinxupquote{PACKAGES/phonon}}
&
\sphinxAtStartPar
no
\\
\sphinxhline
\sphinxAtStartPar
{\hyperref[\detokenize{Packages_details:pkg-plugin}]{\sphinxcrossref{\DUrole{std,std-ref}{PLUGIN}}}}
&
\sphinxAtStartPar
Plugin loader command
&
\sphinxAtStartPar
{\hyperref[\detokenize{plugin::doc}]{\sphinxcrossref{\DUrole{doc}{plugin}}}}
&
\sphinxAtStartPar
plugins
&
\sphinxAtStartPar
no
\\
\sphinxhline
\sphinxAtStartPar
{\hyperref[\detokenize{Packages_details:pkg-plumed}]{\sphinxcrossref{\DUrole{std,std-ref}{PLUMED}}}}
&
\sphinxAtStartPar
\sphinxhref{https://www.plumed.org}{PLUMED free energy library}
&
\sphinxAtStartPar
{\hyperref[\detokenize{fix_plumed::doc}]{\sphinxcrossref{\DUrole{doc}{fix plumed}}}}
&
\sphinxAtStartPar
\sphinxcode{\sphinxupquote{PACKAGES/plumed}}
&
\sphinxAtStartPar
ext
\\
\sphinxhline
\sphinxAtStartPar
{\hyperref[\detokenize{Packages_details:pkg-poems}]{\sphinxcrossref{\DUrole{std,std-ref}{POEMS}}}}
&
\sphinxAtStartPar
coupled rigid body motion
&
\sphinxAtStartPar
{\hyperref[\detokenize{fix_poems::doc}]{\sphinxcrossref{\DUrole{doc}{fix poems}}}}
&
\sphinxAtStartPar
rigid
&
\sphinxAtStartPar
int
\\
\sphinxhline
\sphinxAtStartPar
{\hyperref[\detokenize{Packages_details:pkg-ptm}]{\sphinxcrossref{\DUrole{std,std-ref}{PTM}}}}
&
\sphinxAtStartPar
Polyhedral Template Matching
&
\sphinxAtStartPar
{\hyperref[\detokenize{compute_ptm_atom::doc}]{\sphinxcrossref{\DUrole{doc}{compute ptm/atom}}}}
&
\sphinxAtStartPar
n/a
&
\sphinxAtStartPar
no
\\
\sphinxhline
\sphinxAtStartPar
{\hyperref[\detokenize{Packages_details:pkg-python}]{\sphinxcrossref{\DUrole{std,std-ref}{PYTHON}}}}
&
\sphinxAtStartPar
embed Python code in an input script
&
\sphinxAtStartPar
{\hyperref[\detokenize{python::doc}]{\sphinxcrossref{\DUrole{doc}{python}}}}
&
\sphinxAtStartPar
python
&
\sphinxAtStartPar
sys
\\
\sphinxhline
\sphinxAtStartPar
{\hyperref[\detokenize{Packages_details:pkg-qeq}]{\sphinxcrossref{\DUrole{std,std-ref}{QEQ}}}}
&
\sphinxAtStartPar
QEq charge equilibration
&
\sphinxAtStartPar
{\hyperref[\detokenize{fix_qeq::doc}]{\sphinxcrossref{\DUrole{doc}{fix qeq}}}}
&
\sphinxAtStartPar
qeq
&
\sphinxAtStartPar
no
\\
\sphinxhline
\sphinxAtStartPar
{\hyperref[\detokenize{Packages_details:pkg-qmmm}]{\sphinxcrossref{\DUrole{std,std-ref}{QMMM}}}}
&
\sphinxAtStartPar
QM/MM coupling
&
\sphinxAtStartPar
{\hyperref[\detokenize{fix_qmmm::doc}]{\sphinxcrossref{\DUrole{doc}{fix qmmm}}}}
&
\sphinxAtStartPar
\sphinxcode{\sphinxupquote{PACKAGES/qmmm}}
&
\sphinxAtStartPar
ext
\\
\sphinxhline
\sphinxAtStartPar
{\hyperref[\detokenize{Packages_details:pkg-qtb}]{\sphinxcrossref{\DUrole{std,std-ref}{QTB}}}}
&
\sphinxAtStartPar
quantum nuclear effects
&
\sphinxAtStartPar
{\hyperref[\detokenize{fix_qtb::doc}]{\sphinxcrossref{\DUrole{doc}{fix qtb}}}} {\hyperref[\detokenize{fix_qbmsst::doc}]{\sphinxcrossref{\DUrole{doc}{fix qbmsst}}}}
&
\sphinxAtStartPar
qtb
&
\sphinxAtStartPar
no
\\
\sphinxhline
\sphinxAtStartPar
{\hyperref[\detokenize{Packages_details:pkg-rheo}]{\sphinxcrossref{\DUrole{std,std-ref}{RHEO}}}}
&
\sphinxAtStartPar
reproducing hydrodynamics and elastic objects
&
\sphinxAtStartPar
{\hyperref[\detokenize{Howto_rheo::doc}]{\sphinxcrossref{\DUrole{doc}{Howto rheo}}}}
&
\sphinxAtStartPar
rheo
&
\sphinxAtStartPar
no
\\
\sphinxhline
\sphinxAtStartPar
{\hyperref[\detokenize{Packages_details:pkg-reaction}]{\sphinxcrossref{\DUrole{std,std-ref}{REACTION}}}}
&
\sphinxAtStartPar
chemical reactions in classical MD
&
\sphinxAtStartPar
{\hyperref[\detokenize{fix_bond_react::doc}]{\sphinxcrossref{\DUrole{doc}{fix bond/react}}}}
&
\sphinxAtStartPar
\sphinxcode{\sphinxupquote{PACKAGES/reaction}}
&
\sphinxAtStartPar
no
\\
\sphinxhline
\sphinxAtStartPar
{\hyperref[\detokenize{Packages_details:pkg-reaxff}]{\sphinxcrossref{\DUrole{std,std-ref}{REAXFF}}}}
&
\sphinxAtStartPar
ReaxFF potential (C/C++)
&
\sphinxAtStartPar
{\hyperref[\detokenize{pair_reaxff::doc}]{\sphinxcrossref{\DUrole{doc}{pair\_style reaxff}}}}
&
\sphinxAtStartPar
reax
&
\sphinxAtStartPar
no
\\
\sphinxhline
\sphinxAtStartPar
{\hyperref[\detokenize{Packages_details:pkg-replica}]{\sphinxcrossref{\DUrole{std,std-ref}{REPLICA}}}}
&
\sphinxAtStartPar
multi\sphinxhyphen{}replica methods
&
\sphinxAtStartPar
{\hyperref[\detokenize{Howto_replica::doc}]{\sphinxcrossref{\DUrole{doc}{Howto replica}}}}
&
\sphinxAtStartPar
tad
&
\sphinxAtStartPar
no
\\
\sphinxhline
\sphinxAtStartPar
{\hyperref[\detokenize{Packages_details:pkg-rigid}]{\sphinxcrossref{\DUrole{std,std-ref}{RIGID}}}}
&
\sphinxAtStartPar
rigid bodies and constraints
&
\sphinxAtStartPar
{\hyperref[\detokenize{fix_rigid::doc}]{\sphinxcrossref{\DUrole{doc}{fix rigid}}}}
&
\sphinxAtStartPar
rigid
&
\sphinxAtStartPar
no
\\
\sphinxhline
\sphinxAtStartPar
{\hyperref[\detokenize{Packages_details:pkg-scafacos}]{\sphinxcrossref{\DUrole{std,std-ref}{SCAFACOS}}}}
&
\sphinxAtStartPar
wrapper for ScaFaCoS Kspace solver
&
\sphinxAtStartPar
{\hyperref[\detokenize{kspace_style::doc}]{\sphinxcrossref{\DUrole{doc}{kspace\_style scafacos}}}}
&
\sphinxAtStartPar
\sphinxcode{\sphinxupquote{PACKAGES/scafacos}}
&
\sphinxAtStartPar
ext
\\
\sphinxhline
\sphinxAtStartPar
{\hyperref[\detokenize{Packages_details:pkg-shock}]{\sphinxcrossref{\DUrole{std,std-ref}{SHOCK}}}}
&
\sphinxAtStartPar
shock loading methods
&
\sphinxAtStartPar
{\hyperref[\detokenize{fix_msst::doc}]{\sphinxcrossref{\DUrole{doc}{fix msst}}}}
&
\sphinxAtStartPar
n/a
&
\sphinxAtStartPar
no
\\
\sphinxhline
\sphinxAtStartPar
{\hyperref[\detokenize{Packages_details:pkg-smtbq}]{\sphinxcrossref{\DUrole{std,std-ref}{SMTBQ}}}}
&
\sphinxAtStartPar
second moment tight binding potentials
&
\sphinxAtStartPar
pair styles {\hyperref[\detokenize{pair_smtbq::doc}]{\sphinxcrossref{\DUrole{doc}{smtbq}}}}, {\hyperref[\detokenize{pair_smatb::doc}]{\sphinxcrossref{\DUrole{doc}{smatb}}}}
&
\sphinxAtStartPar
\sphinxcode{\sphinxupquote{PACKAGES/smtbq}}
&
\sphinxAtStartPar
no
\\
\sphinxhline
\sphinxAtStartPar
{\hyperref[\detokenize{Packages_details:pkg-sph}]{\sphinxcrossref{\DUrole{std,std-ref}{SPH}}}}
&
\sphinxAtStartPar
smoothed particle hydrodynamics
&
\sphinxAtStartPar
\sphinxhref{PDF/SPH\_LAMMPS\_userguide.pdf}{SPH User Guide}
&
\sphinxAtStartPar
\sphinxcode{\sphinxupquote{PACKAGES/sph}}
&
\sphinxAtStartPar
no
\\
\sphinxhline
\sphinxAtStartPar
{\hyperref[\detokenize{Packages_details:pkg-spin}]{\sphinxcrossref{\DUrole{std,std-ref}{SPIN}}}}
&
\sphinxAtStartPar
magnetic atomic spin dynamics
&
\sphinxAtStartPar
{\hyperref[\detokenize{Howto_spins::doc}]{\sphinxcrossref{\DUrole{doc}{Howto spins}}}}
&
\sphinxAtStartPar
SPIN
&
\sphinxAtStartPar
no
\\
\sphinxhline
\sphinxAtStartPar
{\hyperref[\detokenize{Packages_details:pkg-srd}]{\sphinxcrossref{\DUrole{std,std-ref}{SRD}}}}
&
\sphinxAtStartPar
stochastic rotation dynamics
&
\sphinxAtStartPar
{\hyperref[\detokenize{fix_srd::doc}]{\sphinxcrossref{\DUrole{doc}{fix srd}}}}
&
\sphinxAtStartPar
srd
&
\sphinxAtStartPar
no
\\
\sphinxhline
\sphinxAtStartPar
{\hyperref[\detokenize{Packages_details:pkg-tally}]{\sphinxcrossref{\DUrole{std,std-ref}{TALLY}}}}
&
\sphinxAtStartPar
pairwise tally computes
&
\sphinxAtStartPar
{\hyperref[\detokenize{compute_tally::doc}]{\sphinxcrossref{\DUrole{doc}{compute XXX/tally}}}}
&
\sphinxAtStartPar
\sphinxcode{\sphinxupquote{PACKAGES/tally}}
&
\sphinxAtStartPar
no
\\
\sphinxhline
\sphinxAtStartPar
{\hyperref[\detokenize{Packages_details:pkg-uef}]{\sphinxcrossref{\DUrole{std,std-ref}{UEF}}}}
&
\sphinxAtStartPar
extensional flow
&
\sphinxAtStartPar
{\hyperref[\detokenize{fix_nh_uef::doc}]{\sphinxcrossref{\DUrole{doc}{fix nvt/uef}}}}
&
\sphinxAtStartPar
\sphinxcode{\sphinxupquote{PACKAGES/uef}}
&
\sphinxAtStartPar
no
\\
\sphinxhline
\sphinxAtStartPar
{\hyperref[\detokenize{Packages_details:pkg-voronoi}]{\sphinxcrossref{\DUrole{std,std-ref}{VORONOI}}}}
&
\sphinxAtStartPar
Voronoi tesselation
&
\sphinxAtStartPar
{\hyperref[\detokenize{compute_voronoi_atom::doc}]{\sphinxcrossref{\DUrole{doc}{compute voronoi/atom}}}}
&
\sphinxAtStartPar
n/a
&
\sphinxAtStartPar
ext
\\
\sphinxhline
\sphinxAtStartPar
{\hyperref[\detokenize{Packages_details:pkg-vtk}]{\sphinxcrossref{\DUrole{std,std-ref}{VTK}}}}
&
\sphinxAtStartPar
dump output via VTK
&
\sphinxAtStartPar
{\hyperref[\detokenize{dump_vtk::doc}]{\sphinxcrossref{\DUrole{doc}{compute vtk}}}}
&
\sphinxAtStartPar
n/a
&
\sphinxAtStartPar
ext
\\
\sphinxhline
\sphinxAtStartPar
{\hyperref[\detokenize{Packages_details:pkg-yaff}]{\sphinxcrossref{\DUrole{std,std-ref}{YAFF}}}}
&
\sphinxAtStartPar
additional styles implemented in YAFF
&
\sphinxAtStartPar
{\hyperref[\detokenize{angle_cross::doc}]{\sphinxcrossref{\DUrole{doc}{angle\_style cross}}}}
&
\sphinxAtStartPar
\sphinxcode{\sphinxupquote{PACKAGES/yaff}}
&
\sphinxAtStartPar
no
\\
\sphinxbottomrule
\end{longtable}
\sphinxtableafterendhook
\sphinxatlongtableend
\end{savenotes}

\sphinxstepscope


\section{Package details}
\label{\detokenize{Packages_details:package-details}}\label{\detokenize{Packages_details::doc}}
\sphinxAtStartPar
Here is a brief description of all packages in LAMMPS.  It lists authors
(if applicable) and summarizes the package contents.  It has specific
instructions on how to install the package, including, if necessary,
info on how to download or build any extra library it requires.  It also
gives links to documentation, example scripts, and pictures/movies (if
available) that illustrate use of the package.

\sphinxAtStartPar
The majority of packages can be included in a LAMMPS build with a single
setting (\sphinxcode{\sphinxupquote{\sphinxhyphen{}D PKG\_\textless{}NAME\textgreater{}=on}} for CMake) or command (\sphinxcode{\sphinxupquote{make yes\sphinxhyphen{}\textless{}name\textgreater{}}}
for make).  See the {\hyperref[\detokenize{Build_package::doc}]{\sphinxcrossref{\DUrole{doc}{Build package}}}} page for more
info.  A few packages may require additional steps; this is indicated in
the descriptions below.  The {\hyperref[\detokenize{Build_extras::doc}]{\sphinxcrossref{\DUrole{doc}{Build extras}}}} page
gives those details.

\begin{sphinxadmonition}{note}{Note:}
\sphinxAtStartPar
To see the complete list of commands a package adds to LAMMPS,
you can examine the files in its src directory, e.g. \sphinxcode{\sphinxupquote{ls
src/GRANULAR}}.  Files with names that start with fix, compute, atom,
pair, bond, angle, etc correspond to commands with the same style name
as contained in the file name.
\end{sphinxadmonition}


\begin{savenotes}\sphinxattablestart
\sphinxthistablewithglobalstyle
\centering
\begin{tabulary}{\linewidth}[t]{TTTTTT}
\sphinxtoprule
\sphinxtableatstartofbodyhook
\sphinxAtStartPar
{\hyperref[\detokenize{Packages_details:pkg-adios}]{\sphinxcrossref{\DUrole{std,std-ref}{ADIOS}}}}
&
\sphinxAtStartPar
{\hyperref[\detokenize{Packages_details:pkg-amoeba}]{\sphinxcrossref{\DUrole{std,std-ref}{AMOEBA}}}}
&
\sphinxAtStartPar
{\hyperref[\detokenize{Packages_details:pkg-asphere}]{\sphinxcrossref{\DUrole{std,std-ref}{ASPHERE}}}}
&
\sphinxAtStartPar
{\hyperref[\detokenize{Packages_details:pkg-atc}]{\sphinxcrossref{\DUrole{std,std-ref}{ATC}}}}
&
\sphinxAtStartPar
{\hyperref[\detokenize{Packages_details:pkg-awpmd}]{\sphinxcrossref{\DUrole{std,std-ref}{AWPMD}}}}
&
\sphinxAtStartPar
{\hyperref[\detokenize{Packages_details:pkg-bocs}]{\sphinxcrossref{\DUrole{std,std-ref}{BOCS}}}}
\\
\sphinxhline
\sphinxAtStartPar
{\hyperref[\detokenize{Packages_details:pkg-body}]{\sphinxcrossref{\DUrole{std,std-ref}{BODY}}}}
&
\sphinxAtStartPar
{\hyperref[\detokenize{Packages_details:pkg-bpm}]{\sphinxcrossref{\DUrole{std,std-ref}{BPM}}}}
&
\sphinxAtStartPar
{\hyperref[\detokenize{Packages_details:pkg-brownian}]{\sphinxcrossref{\DUrole{std,std-ref}{BROWNIAN}}}}
&
\sphinxAtStartPar
{\hyperref[\detokenize{Packages_details:pkg-cg-dna}]{\sphinxcrossref{\DUrole{std,std-ref}{CG\sphinxhyphen{}DNA}}}}
&
\sphinxAtStartPar
{\hyperref[\detokenize{Packages_details:pkg-cg-spica}]{\sphinxcrossref{\DUrole{std,std-ref}{CG\sphinxhyphen{}SPICA}}}}
&
\sphinxAtStartPar
{\hyperref[\detokenize{Packages_details:pkg-class2}]{\sphinxcrossref{\DUrole{std,std-ref}{CLASS2}}}}
\\
\sphinxhline
\sphinxAtStartPar
{\hyperref[\detokenize{Packages_details:pkg-colloid}]{\sphinxcrossref{\DUrole{std,std-ref}{COLLOID}}}}
&
\sphinxAtStartPar
{\hyperref[\detokenize{Packages_details:pkg-colvars}]{\sphinxcrossref{\DUrole{std,std-ref}{COLVARS}}}}
&
\sphinxAtStartPar
{\hyperref[\detokenize{Packages_details:pkg-compress}]{\sphinxcrossref{\DUrole{std,std-ref}{COMPRESS}}}}
&
\sphinxAtStartPar
{\hyperref[\detokenize{Packages_details:pkg-coreshell}]{\sphinxcrossref{\DUrole{std,std-ref}{CORESHELL}}}}
&
\sphinxAtStartPar
{\hyperref[\detokenize{Packages_details:pkg-dielectric}]{\sphinxcrossref{\DUrole{std,std-ref}{DIELECTRIC}}}}
&
\sphinxAtStartPar
{\hyperref[\detokenize{Packages_details:pkg-diffraction}]{\sphinxcrossref{\DUrole{std,std-ref}{DIFFRACTION}}}}
\\
\sphinxhline
\sphinxAtStartPar
{\hyperref[\detokenize{Packages_details:pkg-dipole}]{\sphinxcrossref{\DUrole{std,std-ref}{DIPOLE}}}}
&
\sphinxAtStartPar
{\hyperref[\detokenize{Packages_details:pkg-dpd-basic}]{\sphinxcrossref{\DUrole{std,std-ref}{DPD\sphinxhyphen{}BASIC}}}}
&
\sphinxAtStartPar
{\hyperref[\detokenize{Packages_details:pkg-dpd-meso}]{\sphinxcrossref{\DUrole{std,std-ref}{DPD\sphinxhyphen{}MESO}}}}
&
\sphinxAtStartPar
{\hyperref[\detokenize{Packages_details:pkg-dpd-react}]{\sphinxcrossref{\DUrole{std,std-ref}{DPD\sphinxhyphen{}REACT}}}}
&
\sphinxAtStartPar
{\hyperref[\detokenize{Packages_details:pkg-dpd-smooth}]{\sphinxcrossref{\DUrole{std,std-ref}{DPD\sphinxhyphen{}SMOOTH}}}}
&
\sphinxAtStartPar
{\hyperref[\detokenize{Packages_details:pkg-drude}]{\sphinxcrossref{\DUrole{std,std-ref}{DRUDE}}}}
\\
\sphinxhline
\sphinxAtStartPar
{\hyperref[\detokenize{Packages_details:pkg-eff}]{\sphinxcrossref{\DUrole{std,std-ref}{EFF}}}}
&
\sphinxAtStartPar
{\hyperref[\detokenize{Packages_details:pkg-electrode}]{\sphinxcrossref{\DUrole{std,std-ref}{ELECTRODE}}}}
&
\sphinxAtStartPar
{\hyperref[\detokenize{Packages_details:pkg-extra-command}]{\sphinxcrossref{\DUrole{std,std-ref}{EXTRA\sphinxhyphen{}COMMAND}}}}
&
\sphinxAtStartPar
{\hyperref[\detokenize{Packages_details:pkg-extra-compute}]{\sphinxcrossref{\DUrole{std,std-ref}{EXTRA\sphinxhyphen{}COMPUTE}}}}
&
\sphinxAtStartPar
{\hyperref[\detokenize{Packages_details:pkg-extra-dump}]{\sphinxcrossref{\DUrole{std,std-ref}{EXTRA\sphinxhyphen{}DUMP}}}}
&
\sphinxAtStartPar
{\hyperref[\detokenize{Packages_details:pkg-extra-fix}]{\sphinxcrossref{\DUrole{std,std-ref}{EXTRA\sphinxhyphen{}FIX}}}}
\\
\sphinxhline
\sphinxAtStartPar
{\hyperref[\detokenize{Packages_details:pkg-extra-molecule}]{\sphinxcrossref{\DUrole{std,std-ref}{EXTRA\sphinxhyphen{}MOLECULE}}}}
&
\sphinxAtStartPar
{\hyperref[\detokenize{Packages_details:pkg-extra-pair}]{\sphinxcrossref{\DUrole{std,std-ref}{EXTRA\sphinxhyphen{}PAIR}}}}
&
\sphinxAtStartPar
{\hyperref[\detokenize{Packages_details:pkg-fep}]{\sphinxcrossref{\DUrole{std,std-ref}{FEP}}}}
&
\sphinxAtStartPar
{\hyperref[\detokenize{Packages_details:pkg-gpu}]{\sphinxcrossref{\DUrole{std,std-ref}{GPU}}}}
&
\sphinxAtStartPar
{\hyperref[\detokenize{Packages_details:pkg-granular}]{\sphinxcrossref{\DUrole{std,std-ref}{GRANULAR}}}}
&
\sphinxAtStartPar
{\hyperref[\detokenize{Packages_details:pkg-h5md}]{\sphinxcrossref{\DUrole{std,std-ref}{H5MD}}}}
\\
\sphinxhline
\sphinxAtStartPar
{\hyperref[\detokenize{Packages_details:pkg-intel}]{\sphinxcrossref{\DUrole{std,std-ref}{INTEL}}}}
&
\sphinxAtStartPar
{\hyperref[\detokenize{Packages_details:pkg-interlayer}]{\sphinxcrossref{\DUrole{std,std-ref}{INTERLAYER}}}}
&
\sphinxAtStartPar
{\hyperref[\detokenize{Packages_details:pkg-kim}]{\sphinxcrossref{\DUrole{std,std-ref}{KIM}}}}
&
\sphinxAtStartPar
{\hyperref[\detokenize{Packages_details:pkg-kokkos}]{\sphinxcrossref{\DUrole{std,std-ref}{KOKKOS}}}}
&
\sphinxAtStartPar
{\hyperref[\detokenize{Packages_details:pkg-kspace}]{\sphinxcrossref{\DUrole{std,std-ref}{KSPACE}}}}
&
\sphinxAtStartPar
{\hyperref[\detokenize{Packages_details:pkg-latboltz}]{\sphinxcrossref{\DUrole{std,std-ref}{LATBOLTZ}}}}
\\
\sphinxhline
\sphinxAtStartPar
{\hyperref[\detokenize{Packages_details:pkg-lepton}]{\sphinxcrossref{\DUrole{std,std-ref}{LEPTON}}}}
&
\sphinxAtStartPar
{\hyperref[\detokenize{Packages_details:pkg-machdyn}]{\sphinxcrossref{\DUrole{std,std-ref}{MACHDYN}}}}
&
\sphinxAtStartPar
{\hyperref[\detokenize{Packages_details:pkg-manifold}]{\sphinxcrossref{\DUrole{std,std-ref}{MANIFOLD}}}}
&
\sphinxAtStartPar
{\hyperref[\detokenize{Packages_details:pkg-manybody}]{\sphinxcrossref{\DUrole{std,std-ref}{MANYBODY}}}}
&
\sphinxAtStartPar
{\hyperref[\detokenize{Packages_details:pkg-mc}]{\sphinxcrossref{\DUrole{std,std-ref}{MC}}}}
&
\sphinxAtStartPar
{\hyperref[\detokenize{Packages_details:pkg-mdi}]{\sphinxcrossref{\DUrole{std,std-ref}{MDI}}}}
\\
\sphinxhline
\sphinxAtStartPar
{\hyperref[\detokenize{Packages_details:pkg-meam}]{\sphinxcrossref{\DUrole{std,std-ref}{MEAM}}}}
&
\sphinxAtStartPar
{\hyperref[\detokenize{Packages_details:pkg-mesont}]{\sphinxcrossref{\DUrole{std,std-ref}{MESONT}}}}
&
\sphinxAtStartPar
{\hyperref[\detokenize{Packages_details:pkg-mgpt}]{\sphinxcrossref{\DUrole{std,std-ref}{MGPT}}}}
&
\sphinxAtStartPar
{\hyperref[\detokenize{Packages_details:pkg-misc}]{\sphinxcrossref{\DUrole{std,std-ref}{MISC}}}}
&
\sphinxAtStartPar
{\hyperref[\detokenize{Packages_details:pkg-ml-hdnnp}]{\sphinxcrossref{\DUrole{std,std-ref}{ML\sphinxhyphen{}HDNNP}}}}
&
\sphinxAtStartPar
{\hyperref[\detokenize{Packages_details:pkg-ml-iap}]{\sphinxcrossref{\DUrole{std,std-ref}{ML\sphinxhyphen{}IAP}}}}
\\
\sphinxhline
\sphinxAtStartPar
{\hyperref[\detokenize{Packages_details:pkg-ml-pace}]{\sphinxcrossref{\DUrole{std,std-ref}{ML\sphinxhyphen{}PACE}}}}
&
\sphinxAtStartPar
{\hyperref[\detokenize{Packages_details:pkg-ml-pod}]{\sphinxcrossref{\DUrole{std,std-ref}{ML\sphinxhyphen{}POD}}}}
&
\sphinxAtStartPar
{\hyperref[\detokenize{Packages_details:pkg-ml-quip}]{\sphinxcrossref{\DUrole{std,std-ref}{ML\sphinxhyphen{}QUIP}}}}
&
\sphinxAtStartPar
{\hyperref[\detokenize{Packages_details:pkg-ml-rann}]{\sphinxcrossref{\DUrole{std,std-ref}{ML\sphinxhyphen{}RANN}}}}
&
\sphinxAtStartPar
{\hyperref[\detokenize{Packages_details:pkg-ml-snap}]{\sphinxcrossref{\DUrole{std,std-ref}{ML\sphinxhyphen{}SNAP}}}}
&
\sphinxAtStartPar
{\hyperref[\detokenize{Packages_details:pkg-ml-uf3}]{\sphinxcrossref{\DUrole{std,std-ref}{ML\sphinxhyphen{}UF3}}}}
\\
\sphinxhline
\sphinxAtStartPar
{\hyperref[\detokenize{Packages_details:pkg-mofff}]{\sphinxcrossref{\DUrole{std,std-ref}{MOFFF}}}}
&
\sphinxAtStartPar
{\hyperref[\detokenize{Packages_details:pkg-molecule}]{\sphinxcrossref{\DUrole{std,std-ref}{MOLECULE}}}}
&
\sphinxAtStartPar
{\hyperref[\detokenize{Packages_details:pkg-molfile}]{\sphinxcrossref{\DUrole{std,std-ref}{MOLFILE}}}}
&
\sphinxAtStartPar
{\hyperref[\detokenize{Packages_details:pkg-netcdf}]{\sphinxcrossref{\DUrole{std,std-ref}{NETCDF}}}}
&
\sphinxAtStartPar
{\hyperref[\detokenize{Packages_details:pkg-openmp}]{\sphinxcrossref{\DUrole{std,std-ref}{OPENMP}}}}
&
\sphinxAtStartPar
{\hyperref[\detokenize{Packages_details:pkg-opt}]{\sphinxcrossref{\DUrole{std,std-ref}{OPT}}}}
\\
\sphinxhline
\sphinxAtStartPar
{\hyperref[\detokenize{Packages_details:pkg-orient}]{\sphinxcrossref{\DUrole{std,std-ref}{ORIENT}}}}
&
\sphinxAtStartPar
{\hyperref[\detokenize{Packages_details:pkg-peri}]{\sphinxcrossref{\DUrole{std,std-ref}{PERI}}}}
&
\sphinxAtStartPar
{\hyperref[\detokenize{Packages_details:pkg-phonon}]{\sphinxcrossref{\DUrole{std,std-ref}{PHONON}}}}
&
\sphinxAtStartPar
{\hyperref[\detokenize{Packages_details:pkg-plugin}]{\sphinxcrossref{\DUrole{std,std-ref}{PLUGIN}}}}
&
\sphinxAtStartPar
{\hyperref[\detokenize{Packages_details:pkg-plumed}]{\sphinxcrossref{\DUrole{std,std-ref}{PLUMED}}}}
&
\sphinxAtStartPar
{\hyperref[\detokenize{Packages_details:pkg-poems}]{\sphinxcrossref{\DUrole{std,std-ref}{POEMS}}}}
\\
\sphinxhline
\sphinxAtStartPar
{\hyperref[\detokenize{Packages_details:pkg-ptm}]{\sphinxcrossref{\DUrole{std,std-ref}{PTM}}}}
&
\sphinxAtStartPar
{\hyperref[\detokenize{Packages_details:pkg-python}]{\sphinxcrossref{\DUrole{std,std-ref}{PYTHON}}}}
&
\sphinxAtStartPar
{\hyperref[\detokenize{Packages_details:pkg-qeq}]{\sphinxcrossref{\DUrole{std,std-ref}{QEQ}}}}
&
\sphinxAtStartPar
{\hyperref[\detokenize{Packages_details:pkg-qmmm}]{\sphinxcrossref{\DUrole{std,std-ref}{QMMM}}}}
&
\sphinxAtStartPar
{\hyperref[\detokenize{Packages_details:pkg-qtb}]{\sphinxcrossref{\DUrole{std,std-ref}{QTB}}}}
&
\sphinxAtStartPar
{\hyperref[\detokenize{Packages_details:pkg-rheo}]{\sphinxcrossref{\DUrole{std,std-ref}{RHEO}}}}
\\
\sphinxhline
\sphinxAtStartPar
{\hyperref[\detokenize{Packages_details:pkg-reaction}]{\sphinxcrossref{\DUrole{std,std-ref}{REACTION}}}}
&
\sphinxAtStartPar
{\hyperref[\detokenize{Packages_details:pkg-reaxff}]{\sphinxcrossref{\DUrole{std,std-ref}{REAXFF}}}}
&
\sphinxAtStartPar
{\hyperref[\detokenize{Packages_details:pkg-replica}]{\sphinxcrossref{\DUrole{std,std-ref}{REPLICA}}}}
&
\sphinxAtStartPar
{\hyperref[\detokenize{Packages_details:pkg-rigid}]{\sphinxcrossref{\DUrole{std,std-ref}{RIGID}}}}
&
\sphinxAtStartPar
{\hyperref[\detokenize{Packages_details:pkg-scafacos}]{\sphinxcrossref{\DUrole{std,std-ref}{SCAFACOS}}}}
&
\sphinxAtStartPar
{\hyperref[\detokenize{Packages_details:pkg-shock}]{\sphinxcrossref{\DUrole{std,std-ref}{SHOCK}}}}
\\
\sphinxhline
\sphinxAtStartPar
{\hyperref[\detokenize{Packages_details:pkg-smtbq}]{\sphinxcrossref{\DUrole{std,std-ref}{SMTBQ}}}}
&
\sphinxAtStartPar
{\hyperref[\detokenize{Packages_details:pkg-sph}]{\sphinxcrossref{\DUrole{std,std-ref}{SPH}}}}
&
\sphinxAtStartPar
{\hyperref[\detokenize{Packages_details:pkg-spin}]{\sphinxcrossref{\DUrole{std,std-ref}{SPIN}}}}
&
\sphinxAtStartPar
{\hyperref[\detokenize{Packages_details:pkg-srd}]{\sphinxcrossref{\DUrole{std,std-ref}{SRD}}}}
&
\sphinxAtStartPar
{\hyperref[\detokenize{Packages_details:pkg-tally}]{\sphinxcrossref{\DUrole{std,std-ref}{TALLY}}}}
&
\sphinxAtStartPar
{\hyperref[\detokenize{Packages_details:pkg-uef}]{\sphinxcrossref{\DUrole{std,std-ref}{UEF}}}}
\\
\sphinxhline
\sphinxAtStartPar
{\hyperref[\detokenize{Packages_details:pkg-voronoi}]{\sphinxcrossref{\DUrole{std,std-ref}{VORONOI}}}}
&
\sphinxAtStartPar
{\hyperref[\detokenize{Packages_details:pkg-vtk}]{\sphinxcrossref{\DUrole{std,std-ref}{VTK}}}}
&
\sphinxAtStartPar
{\hyperref[\detokenize{Packages_details:pkg-yaff}]{\sphinxcrossref{\DUrole{std,std-ref}{YAFF}}}}
&&&\\
\sphinxbottomrule
\end{tabulary}
\sphinxtableafterendhook\par
\sphinxattableend\end{savenotes}


\bigskip\hrule\bigskip



\subsection{ADIOS package}
\label{\detokenize{Packages_details:adios-package}}\label{\detokenize{Packages_details:pkg-adios}}
\sphinxAtStartPar
\sphinxstylestrong{Contents:}

\sphinxAtStartPar
ADIOS is a high\sphinxhyphen{}performance I/O library. This package implements the
{\hyperref[\detokenize{dump_adios::doc}]{\sphinxcrossref{\DUrole{doc}{dump atom/adios}}}}, {\hyperref[\detokenize{dump_adios::doc}]{\sphinxcrossref{\DUrole{doc}{dump custom/adios}}}} and
{\hyperref[\detokenize{read_dump::doc}]{\sphinxcrossref{\DUrole{doc}{read\_dump … format adios}}}}
commands to write and read data using the ADIOS library.

\sphinxAtStartPar
\sphinxstylestrong{Authors:} Norbert Podhorszki (ORNL) from the ADIOS developer team.

\sphinxAtStartPar
\DUrole{versionmodified,added}{Added in version 28Feb2019.}

\sphinxAtStartPar
\sphinxstylestrong{Install:}

\sphinxAtStartPar
This package has {\hyperref[\detokenize{Build_extras:adios}]{\sphinxcrossref{\DUrole{std,std-ref}{specific installation instructions}}}} on the {\hyperref[\detokenize{Build_extras::doc}]{\sphinxcrossref{\DUrole{doc}{Build extras}}}} page.

\sphinxAtStartPar
\sphinxstylestrong{Supporting info:}
\begin{itemize}
\item {} 
\sphinxAtStartPar
\sphinxcode{\sphinxupquote{src/ADIOS}}: filenames \sphinxhyphen{}\textgreater{} commands

\item {} 
\sphinxAtStartPar
\sphinxcode{\sphinxupquote{src/ADIOS/README}}

\item {} 
\sphinxAtStartPar
\sphinxcode{\sphinxupquote{examples/PACKAGES/adios}}

\item {} 
\sphinxAtStartPar
\sphinxurl{https://github.com/ornladios/ADIOS2}

\item {} 
\sphinxAtStartPar
{\hyperref[\detokenize{dump_adios::doc}]{\sphinxcrossref{\DUrole{doc}{dump atom/adios}}}}

\item {} 
\sphinxAtStartPar
{\hyperref[\detokenize{dump_adios::doc}]{\sphinxcrossref{\DUrole{doc}{dump custom/adios}}}}

\item {} 
\sphinxAtStartPar
{\hyperref[\detokenize{read_dump::doc}]{\sphinxcrossref{\DUrole{doc}{read\_dump}}}}

\end{itemize}


\bigskip\hrule\bigskip



\subsection{AMOEBA package}
\label{\detokenize{Packages_details:amoeba-package}}\label{\detokenize{Packages_details:pkg-amoeba}}
\sphinxAtStartPar
\sphinxstylestrong{Contents:}

\sphinxAtStartPar
Implementation of the AMOEBA and HIPPO polarized force fields
originally developed by Jay Ponder’s group at the U Washington at St
Louis.  The LAMMPS implementation is based on Fortran 90 code
provided by the Ponder group in their
\sphinxhref{https://dasher.wustl.edu/tinker/}{Tinker MD software}.
\begin{description}
\sphinxlineitem{\sphinxstylestrong{Authors:} Josh Rackers and Steve Plimpton (Sandia), Trung Nguyen (U}
\sphinxAtStartPar
Chicago)

\end{description}

\sphinxAtStartPar
\sphinxstylestrong{Supporting info:}
\begin{itemize}
\item {} 
\sphinxAtStartPar
\sphinxcode{\sphinxupquote{src/AMOEBA}}: filenames \sphinxhyphen{}\textgreater{} commands

\item {} 
\sphinxAtStartPar
{\hyperref[\detokenize{Howto_amoeba::doc}]{\sphinxcrossref{\DUrole{doc}{AMOEBA and HIPPO howto}}}}

\item {} 
\sphinxAtStartPar
{\hyperref[\detokenize{pair_amoeba::doc}]{\sphinxcrossref{\DUrole{doc}{pair\_style amoeba}}}}

\item {} 
\sphinxAtStartPar
{\hyperref[\detokenize{pair_amoeba::doc}]{\sphinxcrossref{\DUrole{doc}{pair\_style hippo}}}}

\item {} 
\sphinxAtStartPar
{\hyperref[\detokenize{atom_style::doc}]{\sphinxcrossref{\DUrole{doc}{atom\_style amoeba}}}}

\item {} 
\sphinxAtStartPar
{\hyperref[\detokenize{angle_amoeba::doc}]{\sphinxcrossref{\DUrole{doc}{angle\_style amoeba}}}}

\item {} 
\sphinxAtStartPar
{\hyperref[\detokenize{improper_amoeba::doc}]{\sphinxcrossref{\DUrole{doc}{improper\_style amoeba}}}}

\item {} 
\sphinxAtStartPar
{\hyperref[\detokenize{fix_amoeba_bitorsion::doc}]{\sphinxcrossref{\DUrole{doc}{fix amoeba/bitorsion}}}}

\item {} 
\sphinxAtStartPar
{\hyperref[\detokenize{fix_amoeba_pitorsion::doc}]{\sphinxcrossref{\DUrole{doc}{fix amoeba/pitorsion}}}}

\item {} 
\sphinxAtStartPar
tools/tinker/tinker2lmp.py

\item {} 
\sphinxAtStartPar
\sphinxcode{\sphinxupquote{examples/amoeba}}

\end{itemize}


\bigskip\hrule\bigskip



\subsection{ASPHERE package}
\label{\detokenize{Packages_details:asphere-package}}\label{\detokenize{Packages_details:pkg-asphere}}
\sphinxAtStartPar
\sphinxstylestrong{Contents:}

\sphinxAtStartPar
Computes, time\sphinxhyphen{}integration fixes, and pair styles for aspherical
particle models including ellipsoids, 2d lines, and 3d triangles.

\sphinxAtStartPar
\sphinxstylestrong{Supporting info:}
\begin{itemize}
\item {} 
\sphinxAtStartPar
\sphinxcode{\sphinxupquote{src/ASPHERE}}: filenames \sphinxhyphen{}\textgreater{} commands

\item {} 
\sphinxAtStartPar
{\hyperref[\detokenize{Howto_spherical::doc}]{\sphinxcrossref{\DUrole{doc}{Howto spherical}}}}

\item {} 
\sphinxAtStartPar
{\hyperref[\detokenize{pair_gayberne::doc}]{\sphinxcrossref{\DUrole{doc}{pair\_style gayberne}}}}

\item {} 
\sphinxAtStartPar
{\hyperref[\detokenize{pair_resquared::doc}]{\sphinxcrossref{\DUrole{doc}{pair\_style resquared}}}}

\item {} 
\sphinxAtStartPar
{\hyperref[\detokenize{pair_ylz::doc}]{\sphinxcrossref{\DUrole{doc}{pair\_style ylz}}}}

\item {} 
\sphinxAtStartPar
\sphinxhref{PDF/pair\_gayberne\_extra.pdf}{doc/PDF/pair\_gayberne\_extra.pdf}

\item {} 
\sphinxAtStartPar
\sphinxhref{PDF/pair\_resquared\_extra.pdf}{doc/PDF/pair\_resquared\_extra.pdf}

\item {} 
\sphinxAtStartPar
\sphinxcode{\sphinxupquote{examples/ASPHERE}}

\item {} 
\sphinxAtStartPar
\sphinxcode{\sphinxupquote{examples/ellipse}}

\item {} 
\sphinxAtStartPar
\sphinxurl{https://www.lammps.org/movies.html\#line}

\item {} 
\sphinxAtStartPar
\sphinxurl{https://www.lammps.org/movies.html\#tri}

\end{itemize}


\bigskip\hrule\bigskip



\subsection{ATC package}
\label{\detokenize{Packages_details:atc-package}}\label{\detokenize{Packages_details:pkg-atc}}
\sphinxAtStartPar
\sphinxstylestrong{Contents:}

\sphinxAtStartPar
ATC stands for atoms\sphinxhyphen{}to\sphinxhyphen{}continuum.  This package implements a
{\hyperref[\detokenize{fix_atc::doc}]{\sphinxcrossref{\DUrole{doc}{fix atc}}}} command to either couple molecular dynamics
with continuum finite element equations or perform on\sphinxhyphen{}the\sphinxhyphen{}fly
conversion of atomic information to continuum fields.

\sphinxAtStartPar
\sphinxstylestrong{Authors:} Reese Jones, Jeremy Templeton, Jon Zimmerman (Sandia).

\sphinxAtStartPar
\sphinxstylestrong{Install:}

\sphinxAtStartPar
This package has {\hyperref[\detokenize{Build_extras:atc}]{\sphinxcrossref{\DUrole{std,std-ref}{specific installation instructions}}}} on the
{\hyperref[\detokenize{Build_extras::doc}]{\sphinxcrossref{\DUrole{doc}{Build extras}}}} page.  The ATC package requires that
also the {\hyperref[\detokenize{Packages_details:pkg-manybody}]{\sphinxcrossref{\DUrole{std,std-ref}{MANYBODY}}}} package is installed.

\sphinxAtStartPar
\sphinxstylestrong{Supporting info:}
\begin{itemize}
\item {} 
\sphinxAtStartPar
\sphinxcode{\sphinxupquote{src/ATC}}: filenames \sphinxhyphen{}\textgreater{} commands

\item {} 
\sphinxAtStartPar
\sphinxcode{\sphinxupquote{src/ATC/README}}

\item {} 
\sphinxAtStartPar
{\hyperref[\detokenize{fix_atc::doc}]{\sphinxcrossref{\DUrole{doc}{fix atc}}}}

\item {} 
\sphinxAtStartPar
\sphinxcode{\sphinxupquote{examples/PACKAGES/atc}}

\item {} 
\sphinxAtStartPar
\sphinxurl{https://www.lammps.org/pictures.html\#atc}

\end{itemize}


\bigskip\hrule\bigskip



\subsection{AWPMD package}
\label{\detokenize{Packages_details:awpmd-package}}\label{\detokenize{Packages_details:pkg-awpmd}}
\sphinxAtStartPar
\sphinxstylestrong{Contents:}

\sphinxAtStartPar
AWPMD stands for Antisymmetrized Wave Packet Molecular Dynamics.  This
package implements an atom, pair, and fix style which allows electrons
to be treated as explicit particles in a classical molecular dynamics
model.

\sphinxAtStartPar
\sphinxstylestrong{Author:} Ilya Valuev (JIHT, Russia).

\sphinxAtStartPar
\sphinxstylestrong{Install:}

\sphinxAtStartPar
This package has {\hyperref[\detokenize{Build_extras:awpmd}]{\sphinxcrossref{\DUrole{std,std-ref}{specific installation instructions}}}} on the {\hyperref[\detokenize{Build_extras::doc}]{\sphinxcrossref{\DUrole{doc}{Build extras}}}} page.

\sphinxAtStartPar
\sphinxstylestrong{Supporting info:}
\begin{itemize}
\item {} 
\sphinxAtStartPar
\sphinxcode{\sphinxupquote{src/AWPMD}}: filenames \sphinxhyphen{}\textgreater{} commands

\item {} 
\sphinxAtStartPar
\sphinxcode{\sphinxupquote{src/AWPMD/README}}

\item {} 
\sphinxAtStartPar
{\hyperref[\detokenize{pair_awpmd::doc}]{\sphinxcrossref{\DUrole{doc}{pair\_style awpmd/cut}}}}

\item {} 
\sphinxAtStartPar
\sphinxcode{\sphinxupquote{examples/PACKAGES/awpmd}}

\end{itemize}


\bigskip\hrule\bigskip



\subsection{BOCS package}
\label{\detokenize{Packages_details:bocs-package}}\label{\detokenize{Packages_details:pkg-bocs}}
\sphinxAtStartPar
\sphinxstylestrong{Contents:}

\sphinxAtStartPar
This package provides {\hyperref[\detokenize{fix_bocs::doc}]{\sphinxcrossref{\DUrole{doc}{fix bocs}}}}, a modified version
of {\hyperref[\detokenize{fix_nh::doc}]{\sphinxcrossref{\DUrole{doc}{fix npt}}}} which includes the pressure correction to
the barostat as outlined in:

\sphinxAtStartPar
N. J. H. Dunn and W. G. Noid, “Bottom\sphinxhyphen{}up coarse\sphinxhyphen{}grained models that
accurately describe the structure, pressure, and compressibility of
molecular liquids”, J. Chem. Phys. 143, 243148 (2015).

\sphinxAtStartPar
\sphinxstylestrong{Authors:} Nicholas J. H. Dunn and Michael R. DeLyser (The
Pennsylvania State University)

\sphinxAtStartPar
\sphinxstylestrong{Supporting info:}

\sphinxAtStartPar
The BOCS package for LAMMPS is part of the BOCS software package:
\sphinxurl{https://github.com/noid-group/BOCS}

\sphinxAtStartPar
See the following reference for information about the entire package:

\sphinxAtStartPar
Dunn, NJH; Lebold, KM; DeLyser, MR; Rudzinski, JF; Noid, WG.
“BOCS: Bottom\sphinxhyphen{}Up Open\sphinxhyphen{}Source Coarse\sphinxhyphen{}Graining Software.”
J. Phys. Chem. B. 122, 13, 3363\sphinxhyphen{}3377 (2018).

\sphinxAtStartPar
Example inputs are in the \sphinxcode{\sphinxupquote{examples/PACKAGES/bocs}} folder.


\bigskip\hrule\bigskip



\subsection{BODY package}
\label{\detokenize{Packages_details:body-package}}\label{\detokenize{Packages_details:pkg-body}}
\sphinxAtStartPar
\sphinxstylestrong{Contents:}

\sphinxAtStartPar
Body\sphinxhyphen{}style particles with internal structure.  Computes,
time\sphinxhyphen{}integration fixes, pair styles, as well as the body styles
themselves.  See the {\hyperref[\detokenize{Howto_body::doc}]{\sphinxcrossref{\DUrole{doc}{Howto body}}}} page for an
overview.

\sphinxAtStartPar
\sphinxstylestrong{Supporting info:}
\begin{itemize}
\item {} 
\sphinxAtStartPar
\sphinxcode{\sphinxupquote{src/BODY}} filenames \sphinxhyphen{}\textgreater{} commands

\item {} 
\sphinxAtStartPar
{\hyperref[\detokenize{Howto_body::doc}]{\sphinxcrossref{\DUrole{doc}{Howto\_body}}}}

\item {} 
\sphinxAtStartPar
{\hyperref[\detokenize{atom_style::doc}]{\sphinxcrossref{\DUrole{doc}{atom\_style body}}}}

\item {} 
\sphinxAtStartPar
{\hyperref[\detokenize{fix_nve_body::doc}]{\sphinxcrossref{\DUrole{doc}{fix nve/body}}}}

\item {} 
\sphinxAtStartPar
{\hyperref[\detokenize{pair_body_nparticle::doc}]{\sphinxcrossref{\DUrole{doc}{pair\_style body/nparticle}}}}

\item {} 
\sphinxAtStartPar
\sphinxcode{\sphinxupquote{examples/body}}

\end{itemize}


\bigskip\hrule\bigskip



\subsection{BPM package}
\label{\detokenize{Packages_details:bpm-package}}\label{\detokenize{Packages_details:pkg-bpm}}
\sphinxAtStartPar
\sphinxstylestrong{Contents:}

\sphinxAtStartPar
Pair styles, bond styles, fixes, and computes for bonded particle
models for mesoscale simulations of solids and fracture.  See the
{\hyperref[\detokenize{Howto_bpm::doc}]{\sphinxcrossref{\DUrole{doc}{Howto bpm}}}} page for an overview.

\sphinxAtStartPar
\sphinxstylestrong{Authors:} Joel T. Clemmer (Sandia National Labs)

\sphinxAtStartPar
\DUrole{versionmodified,added}{Added in version 4May2022.}

\sphinxAtStartPar
\sphinxstylestrong{Supporting info:}
\begin{itemize}
\item {} 
\sphinxAtStartPar
\sphinxcode{\sphinxupquote{src/BPM}}: filenames \sphinxhyphen{}\textgreater{} commands

\item {} 
\sphinxAtStartPar
{\hyperref[\detokenize{Howto_bpm::doc}]{\sphinxcrossref{\DUrole{doc}{Howto\_bpm}}}}

\item {} 
\sphinxAtStartPar
{\hyperref[\detokenize{atom_style::doc}]{\sphinxcrossref{\DUrole{doc}{atom\_style bpm/sphere}}}}

\item {} 
\sphinxAtStartPar
{\hyperref[\detokenize{bond_bpm_rotational::doc}]{\sphinxcrossref{\DUrole{doc}{bond\_style bpm/rotational}}}}

\item {} 
\sphinxAtStartPar
{\hyperref[\detokenize{bond_bpm_spring::doc}]{\sphinxcrossref{\DUrole{doc}{bond\_style bpm/spring}}}}

\item {} 
\sphinxAtStartPar
{\hyperref[\detokenize{compute_nbond_atom::doc}]{\sphinxcrossref{\DUrole{doc}{compute nbond/atom}}}}

\item {} 
\sphinxAtStartPar
{\hyperref[\detokenize{fix_nve_bpm_sphere::doc}]{\sphinxcrossref{\DUrole{doc}{fix nve/bpm/sphere}}}}

\item {} 
\sphinxAtStartPar
{\hyperref[\detokenize{pair_bpm_spring::doc}]{\sphinxcrossref{\DUrole{doc}{pair\_style bpm/spring}}}}

\item {} 
\sphinxAtStartPar
\sphinxurl{https://www.lammps.org/movies.html\#bpmpackage}

\item {} 
\sphinxAtStartPar
\sphinxcode{\sphinxupquote{examples/bpm}}

\end{itemize}


\bigskip\hrule\bigskip



\subsection{BROWNIAN package}
\label{\detokenize{Packages_details:brownian-package}}\label{\detokenize{Packages_details:pkg-brownian}}
\sphinxAtStartPar
\sphinxstylestrong{Contents:}

\sphinxAtStartPar
This package provides {\hyperref[\detokenize{fix_brownian::doc}]{\sphinxcrossref{\DUrole{doc}{fix brownian, fix brownian/sphere, and
fix brownian/asphere}}}} as well as
{\hyperref[\detokenize{fix_propel_self::doc}]{\sphinxcrossref{\DUrole{doc}{fix propel/self}}}} which allow to do Brownian
Dynamics time integration of point, spherical and aspherical particles
and also support self\sphinxhyphen{}propelled particles.

\sphinxAtStartPar
\sphinxstylestrong{Authors:} Sam Cameron (University of Bristol),
Stefan Paquay (while at Brandeis University) (initial version of fix propel/self)

\sphinxAtStartPar
\DUrole{versionmodified,added}{Added in version 14May2021.}

\sphinxAtStartPar
Example inputs are in the \sphinxcode{\sphinxupquote{examples/PACKAGES/brownian}} folder.


\bigskip\hrule\bigskip



\subsection{CG\sphinxhyphen{}DNA package}
\label{\detokenize{Packages_details:cg-dna-package}}\label{\detokenize{Packages_details:pkg-cg-dna}}
\sphinxAtStartPar
\sphinxstylestrong{Contents:}

\sphinxAtStartPar
Several pair styles, bond styles, and integration fixes for coarse\sphinxhyphen{}grained
modelling of single\sphinxhyphen{} and double\sphinxhyphen{}stranded DNA and RNA based on the oxDNA and
oxRNA model of Doye, Louis and Ouldridge. The package includes Langevin\sphinxhyphen{}type
rigid\sphinxhyphen{}body integrators with improved stability.

\sphinxAtStartPar
\sphinxstylestrong{Author:} Oliver Henrich (University of Strathclyde, Glasgow).

\sphinxAtStartPar
\sphinxstylestrong{Install:}

\sphinxAtStartPar
The CG\sphinxhyphen{}DNA package requires that also the {\hyperref[\detokenize{Packages_details:pkg-molecule}]{\sphinxcrossref{\DUrole{std,std-ref}{MOLECULE}}}}
and {\hyperref[\detokenize{Packages_details:pkg-asphere}]{\sphinxcrossref{\DUrole{std,std-ref}{ASPHERE}}}} packages are installed.

\sphinxAtStartPar
\sphinxstylestrong{Supporting info:}
\begin{itemize}
\item {} 
\sphinxAtStartPar
\sphinxcode{\sphinxupquote{src/CG\sphinxhyphen{}DNA}}: filenames \sphinxhyphen{}\textgreater{} commands

\item {} 
\sphinxAtStartPar
\sphinxcode{\sphinxupquote{src/CG\sphinxhyphen{}DNA/README}}

\item {} 
\sphinxAtStartPar
{\hyperref[\detokenize{pair_oxdna::doc}]{\sphinxcrossref{\DUrole{doc}{pair\_style oxdna/*}}}}

\item {} 
\sphinxAtStartPar
{\hyperref[\detokenize{pair_oxdna2::doc}]{\sphinxcrossref{\DUrole{doc}{pair\_style oxdna2/*}}}}

\item {} 
\sphinxAtStartPar
{\hyperref[\detokenize{pair_oxrna2::doc}]{\sphinxcrossref{\DUrole{doc}{pair\_style oxrna2/*}}}}

\item {} 
\sphinxAtStartPar
{\hyperref[\detokenize{bond_oxdna::doc}]{\sphinxcrossref{\DUrole{doc}{bond\_style oxdna/*}}}}

\item {} 
\sphinxAtStartPar
{\hyperref[\detokenize{bond_oxdna::doc}]{\sphinxcrossref{\DUrole{doc}{bond\_style oxdna2/*}}}}

\item {} 
\sphinxAtStartPar
{\hyperref[\detokenize{bond_oxdna::doc}]{\sphinxcrossref{\DUrole{doc}{bond\_style oxrna2/*}}}}

\item {} 
\sphinxAtStartPar
{\hyperref[\detokenize{fix_nve_dotc_langevin::doc}]{\sphinxcrossref{\DUrole{doc}{fix nve/dotc/langevin}}}}

\item {} 
\sphinxAtStartPar
\sphinxcode{\sphinxupquote{examples/PACKAGES/cgdna}}

\end{itemize}


\bigskip\hrule\bigskip



\subsection{CG\sphinxhyphen{}SPICA package}
\label{\detokenize{Packages_details:cg-spica-package}}\label{\detokenize{Packages_details:pkg-cg-spica}}
\sphinxAtStartPar
\sphinxstylestrong{Contents:}

\sphinxAtStartPar
Several pair styles and an angle style which implement the
coarse\sphinxhyphen{}grained SPICA (formerly called SDK) model which enables
simulation of biological or soft material systems.

\sphinxAtStartPar
\sphinxstylestrong{Original Author:} Axel Kohlmeyer (Temple U).

\sphinxAtStartPar
\sphinxstylestrong{Maintainers:} Yusuke Miyazaki and Wataru Shinoda (Okayama U).

\sphinxAtStartPar
\sphinxstylestrong{Supporting info:}
\begin{itemize}
\item {} 
\sphinxAtStartPar
\sphinxcode{\sphinxupquote{src/CG\sphinxhyphen{}SPICA}}: filenames \sphinxhyphen{}\textgreater{} commands

\item {} 
\sphinxAtStartPar
\sphinxcode{\sphinxupquote{src/CG\sphinxhyphen{}SPICA/README}}

\item {} 
\sphinxAtStartPar
{\hyperref[\detokenize{pair_spica::doc}]{\sphinxcrossref{\DUrole{doc}{pair\_style lj/spica/*}}}}

\item {} 
\sphinxAtStartPar
{\hyperref[\detokenize{angle_spica::doc}]{\sphinxcrossref{\DUrole{doc}{angle\_style spica}}}}

\item {} 
\sphinxAtStartPar
\sphinxcode{\sphinxupquote{examples/PACKAGES/cgspica}}

\item {} 
\sphinxAtStartPar
\sphinxurl{https://www.lammps.org/pictures.html\#cg}

\item {} 
\sphinxAtStartPar
\sphinxurl{https://www.spica-ff.org/}

\end{itemize}


\bigskip\hrule\bigskip



\subsection{CLASS2 package}
\label{\detokenize{Packages_details:class2-package}}\label{\detokenize{Packages_details:pkg-class2}}
\sphinxAtStartPar
\sphinxstylestrong{Contents:}

\sphinxAtStartPar
Bond, angle, dihedral, improper, and pair styles for the COMPASS
CLASS2 molecular force field.

\sphinxAtStartPar
\sphinxstylestrong{Supporting info:}
\begin{itemize}
\item {} 
\sphinxAtStartPar
\sphinxcode{\sphinxupquote{src/CLASS2}}: filenames \sphinxhyphen{}\textgreater{} commands

\item {} 
\sphinxAtStartPar
{\hyperref[\detokenize{bond_class2::doc}]{\sphinxcrossref{\DUrole{doc}{bond\_style class2}}}}

\item {} 
\sphinxAtStartPar
{\hyperref[\detokenize{angle_class2::doc}]{\sphinxcrossref{\DUrole{doc}{angle\_style class2}}}}

\item {} 
\sphinxAtStartPar
{\hyperref[\detokenize{dihedral_class2::doc}]{\sphinxcrossref{\DUrole{doc}{dihedral\_style class2}}}}

\item {} 
\sphinxAtStartPar
{\hyperref[\detokenize{improper_class2::doc}]{\sphinxcrossref{\DUrole{doc}{improper\_style class2}}}}

\item {} 
\sphinxAtStartPar
{\hyperref[\detokenize{pair_class2::doc}]{\sphinxcrossref{\DUrole{doc}{pair\_style lj/class2}}}}

\end{itemize}


\bigskip\hrule\bigskip



\subsection{COLLOID package}
\label{\detokenize{Packages_details:colloid-package}}\label{\detokenize{Packages_details:pkg-colloid}}
\sphinxAtStartPar
\sphinxstylestrong{Contents:}

\sphinxAtStartPar
Coarse\sphinxhyphen{}grained finite\sphinxhyphen{}size colloidal particles.  Pair styles and fix
wall styles for colloidal interactions.  Includes the Fast Lubrication
Dynamics (FLD) method for hydrodynamic interactions, which is a
simplified approximation to Stokesian dynamics.

\sphinxAtStartPar
\sphinxstylestrong{Authors:} This package includes Fast Lubrication Dynamics pair styles
which were created by Amit Kumar and Michael Bybee from Jonathan
Higdon’s group at UIUC.

\sphinxAtStartPar
\sphinxstylestrong{Supporting info:}
\begin{itemize}
\item {} 
\sphinxAtStartPar
\sphinxcode{\sphinxupquote{src/COLLOID}}: filenames \sphinxhyphen{}\textgreater{} commands

\item {} 
\sphinxAtStartPar
{\hyperref[\detokenize{fix_wall::doc}]{\sphinxcrossref{\DUrole{doc}{fix wall/colloid}}}}

\item {} 
\sphinxAtStartPar
{\hyperref[\detokenize{pair_colloid::doc}]{\sphinxcrossref{\DUrole{doc}{pair\_style colloid}}}}

\item {} 
\sphinxAtStartPar
{\hyperref[\detokenize{pair_yukawa_colloid::doc}]{\sphinxcrossref{\DUrole{doc}{pair\_style yukawa/colloid}}}}

\item {} 
\sphinxAtStartPar
{\hyperref[\detokenize{pair_brownian::doc}]{\sphinxcrossref{\DUrole{doc}{pair\_style brownian}}}}

\item {} 
\sphinxAtStartPar
{\hyperref[\detokenize{pair_lubricate::doc}]{\sphinxcrossref{\DUrole{doc}{pair\_style lubricate}}}}

\item {} 
\sphinxAtStartPar
{\hyperref[\detokenize{pair_lubricateU::doc}]{\sphinxcrossref{\DUrole{doc}{pair\_style lubricateU}}}}

\item {} 
\sphinxAtStartPar
\sphinxcode{\sphinxupquote{examples/colloid}}

\item {} 
\sphinxAtStartPar
\sphinxcode{\sphinxupquote{examples/srd}}

\end{itemize}


\bigskip\hrule\bigskip



\subsection{COLVARS package}
\label{\detokenize{Packages_details:colvars-package}}\label{\detokenize{Packages_details:pkg-colvars}}
\sphinxAtStartPar
\sphinxstylestrong{Contents:}

\sphinxAtStartPar
Colvars stands for collective variables, which can be used to implement
various enhanced sampling methods, including Adaptive Biasing Force,
Metadynamics, Steered MD, Umbrella Sampling and Restraints.  A {\hyperref[\detokenize{fix_colvars::doc}]{\sphinxcrossref{\DUrole{doc}{fix
colvars}}}} command is implemented which wraps a COLVARS
library, which implements these methods.  simulations.

\sphinxAtStartPar
\sphinxstylestrong{Authors:} The COLVARS library is written and maintained by Giacomo
Fiorin (NIH, Bethesda, MD, USA) and Jerome Henin (CNRS, Paris, France),
originally for the NAMD MD code, but with portability in mind.  Axel
Kohlmeyer (Temple U) provided the interface to LAMMPS.

\sphinxAtStartPar
\sphinxstylestrong{Install:}

\sphinxAtStartPar
This package has {\hyperref[\detokenize{Build_extras:colvar}]{\sphinxcrossref{\DUrole{std,std-ref}{specific installation instructions}}}} on
the {\hyperref[\detokenize{Build_extras::doc}]{\sphinxcrossref{\DUrole{doc}{Build extras}}}} page.

\sphinxAtStartPar
\sphinxstylestrong{Supporting info:}
\begin{itemize}
\item {} 
\sphinxAtStartPar
\sphinxcode{\sphinxupquote{src/COLVARS}}: filenames \sphinxhyphen{}\textgreater{} commands

\item {} 
\sphinxAtStartPar
\sphinxhref{PDF/colvars-refman-lammps.pdf}{doc/PDF/colvars\sphinxhyphen{}refman\sphinxhyphen{}lammps.pdf}

\item {} 
\sphinxAtStartPar
\sphinxcode{\sphinxupquote{src/COLVARS/README}}

\item {} 
\sphinxAtStartPar
\sphinxcode{\sphinxupquote{lib/colvars/README}}

\item {} 
\sphinxAtStartPar
{\hyperref[\detokenize{fix_colvars::doc}]{\sphinxcrossref{\DUrole{doc}{fix colvars}}}}

\item {} 
\sphinxAtStartPar
{\hyperref[\detokenize{group2ndx::doc}]{\sphinxcrossref{\DUrole{doc}{group2ndx}}}}

\item {} 
\sphinxAtStartPar
{\hyperref[\detokenize{group2ndx::doc}]{\sphinxcrossref{\DUrole{doc}{ndx2group}}}}

\item {} 
\sphinxAtStartPar
\sphinxcode{\sphinxupquote{examples/PACKAGES/colvars}}

\end{itemize}


\bigskip\hrule\bigskip



\subsection{COMPRESS package}
\label{\detokenize{Packages_details:compress-package}}\label{\detokenize{Packages_details:pkg-compress}}
\sphinxAtStartPar
\sphinxstylestrong{Contents:}

\sphinxAtStartPar
Compressed output of dump files via the zlib compression library,
using dump styles with a “gz” in their style name.

\sphinxAtStartPar
To use this package you must have the zlib compression library
available on your system.

\sphinxAtStartPar
\sphinxstylestrong{Author:} Axel Kohlmeyer (Temple U).

\sphinxAtStartPar
\sphinxstylestrong{Install:}

\sphinxAtStartPar
This package has {\hyperref[\detokenize{Build_extras:compress}]{\sphinxcrossref{\DUrole{std,std-ref}{specific installation instructions}}}} on the {\hyperref[\detokenize{Build_extras::doc}]{\sphinxcrossref{\DUrole{doc}{Build extras}}}} page.

\sphinxAtStartPar
\sphinxstylestrong{Supporting info:}
\begin{itemize}
\item {} 
\sphinxAtStartPar
\sphinxcode{\sphinxupquote{src/COMPRESS}}: filenames \sphinxhyphen{}\textgreater{} commands

\item {} 
\sphinxAtStartPar
\sphinxcode{\sphinxupquote{src/COMPRESS/README}}

\item {} 
\sphinxAtStartPar
\sphinxcode{\sphinxupquote{lib/compress/README}}

\item {} 
\sphinxAtStartPar
{\hyperref[\detokenize{dump::doc}]{\sphinxcrossref{\DUrole{doc}{dump atom/gz}}}}

\item {} 
\sphinxAtStartPar
{\hyperref[\detokenize{dump::doc}]{\sphinxcrossref{\DUrole{doc}{dump cfg/gz}}}}

\item {} 
\sphinxAtStartPar
{\hyperref[\detokenize{dump::doc}]{\sphinxcrossref{\DUrole{doc}{dump custom/gz}}}}

\item {} 
\sphinxAtStartPar
{\hyperref[\detokenize{dump::doc}]{\sphinxcrossref{\DUrole{doc}{dump xyz/gz}}}}

\end{itemize}


\bigskip\hrule\bigskip



\subsection{CORESHELL package}
\label{\detokenize{Packages_details:coreshell-package}}\label{\detokenize{Packages_details:pkg-coreshell}}
\sphinxAtStartPar
\sphinxstylestrong{Contents:}

\sphinxAtStartPar
Compute and pair styles that implement the adiabatic core/shell model
for polarizability.  The pair styles augment Born, Buckingham, and
Lennard\sphinxhyphen{}Jones styles with core/shell capabilities.  The {\hyperref[\detokenize{compute_temp_cs::doc}]{\sphinxcrossref{\DUrole{doc}{compute
temp/cs}}}} command calculates the temperature of a
system with core/shell particles.  See the {\hyperref[\detokenize{Howto_coreshell::doc}]{\sphinxcrossref{\DUrole{doc}{Howto coreshell}}}} page for an overview of how to use this package.

\sphinxAtStartPar
\sphinxstylestrong{Author:} Hendrik Heenen (Technical U of Munich).

\sphinxAtStartPar
\sphinxstylestrong{Supporting info:}
\begin{itemize}
\item {} 
\sphinxAtStartPar
\sphinxcode{\sphinxupquote{src/CORESHELL}}: filenames \sphinxhyphen{}\textgreater{} commands

\item {} 
\sphinxAtStartPar
{\hyperref[\detokenize{Howto_coreshell::doc}]{\sphinxcrossref{\DUrole{doc}{Howto coreshell}}}}

\item {} 
\sphinxAtStartPar
{\hyperref[\detokenize{Howto_polarizable::doc}]{\sphinxcrossref{\DUrole{doc}{Howto polarizable}}}}

\item {} 
\sphinxAtStartPar
{\hyperref[\detokenize{compute_temp_cs::doc}]{\sphinxcrossref{\DUrole{doc}{compute temp/cs}}}}

\item {} 
\sphinxAtStartPar
{\hyperref[\detokenize{pair_cs::doc}]{\sphinxcrossref{\DUrole{doc}{pair\_style born/coul/long/cs}}}}

\item {} 
\sphinxAtStartPar
{\hyperref[\detokenize{pair_cs::doc}]{\sphinxcrossref{\DUrole{doc}{pair\_style buck/coul/long/cs}}}}

\item {} 
\sphinxAtStartPar
{\hyperref[\detokenize{pair_lj::doc}]{\sphinxcrossref{\DUrole{doc}{pair\_style lj/cut/coul/long/cs}}}}

\item {} 
\sphinxAtStartPar
\sphinxcode{\sphinxupquote{examples/coreshell}}

\end{itemize}


\bigskip\hrule\bigskip



\subsection{DIELECTRIC package}
\label{\detokenize{Packages_details:dielectric-package}}\label{\detokenize{Packages_details:pkg-dielectric}}
\sphinxAtStartPar
\sphinxstylestrong{Contents:}

\sphinxAtStartPar
An atom style, multiple pair styles, several fixes, Kspace styles and a
compute for simulating systems using boundary element solvers for
computing the induced charges at the interface between two media with
different dielectric constants.

\sphinxAtStartPar
\sphinxstylestrong{Install:}

\sphinxAtStartPar
To use this package, also the {\hyperref[\detokenize{Packages_details:pkg-kspace}]{\sphinxcrossref{\DUrole{std,std-ref}{KSPACE}}}} and
{\hyperref[\detokenize{Packages_details:pkg-extra-pair}]{\sphinxcrossref{\DUrole{std,std-ref}{EXTRA\sphinxhyphen{}PAIR}}}} packages need to be installed.

\sphinxAtStartPar
\sphinxstylestrong{Author:} Trung Nguyen and Monica Olvera de la Cruz (Northwestern U)

\sphinxAtStartPar
\DUrole{versionmodified,added}{Added in version 2Jul2021.}

\sphinxAtStartPar
\sphinxstylestrong{Supporting info:}
\begin{itemize}
\item {} 
\sphinxAtStartPar
\sphinxcode{\sphinxupquote{src/DIELECTRIC}}: filenames \sphinxhyphen{}\textgreater{} commands

\item {} 
\sphinxAtStartPar
{\hyperref[\detokenize{atom_style::doc}]{\sphinxcrossref{\DUrole{doc}{atom\_style dielectric}}}}

\item {} 
\sphinxAtStartPar
{\hyperref[\detokenize{pair_dielectric::doc}]{\sphinxcrossref{\DUrole{doc}{pair\_style coul/cut/dielectric}}}}

\item {} 
\sphinxAtStartPar
{\hyperref[\detokenize{pair_dielectric::doc}]{\sphinxcrossref{\DUrole{doc}{pair\_style coul/long/dielectric}}}}

\item {} 
\sphinxAtStartPar
{\hyperref[\detokenize{pair_dielectric::doc}]{\sphinxcrossref{\DUrole{doc}{pair\_style lj/cut/coul/cut/dielectric}}}}

\item {} 
\sphinxAtStartPar
{\hyperref[\detokenize{pair_dielectric::doc}]{\sphinxcrossref{\DUrole{doc}{pair\_style lj/cut/coul/debye/dielectric}}}}

\item {} 
\sphinxAtStartPar
{\hyperref[\detokenize{pair_dielectric::doc}]{\sphinxcrossref{\DUrole{doc}{pair\_style lj/cut/coul/long/dielectric}}}}

\item {} 
\sphinxAtStartPar
{\hyperref[\detokenize{pair_dielectric::doc}]{\sphinxcrossref{\DUrole{doc}{pair\_style lj/cut/coul/msm/dielectric}}}}

\item {} 
\sphinxAtStartPar
{\hyperref[\detokenize{kspace_style::doc}]{\sphinxcrossref{\DUrole{doc}{pair\_style pppm/dielectric}}}}

\item {} 
\sphinxAtStartPar
{\hyperref[\detokenize{kspace_style::doc}]{\sphinxcrossref{\DUrole{doc}{pair\_style pppm/disp/dielectric}}}}

\item {} 
\sphinxAtStartPar
{\hyperref[\detokenize{kspace_style::doc}]{\sphinxcrossref{\DUrole{doc}{pair\_style msm/dielectric}}}}

\item {} 
\sphinxAtStartPar
{\hyperref[\detokenize{fix_polarize::doc}]{\sphinxcrossref{\DUrole{doc}{fix\_style polarize/bem/icc}}}}

\item {} 
\sphinxAtStartPar
{\hyperref[\detokenize{fix_polarize::doc}]{\sphinxcrossref{\DUrole{doc}{fix\_style polarize/bem/gmres}}}}

\item {} 
\sphinxAtStartPar
{\hyperref[\detokenize{fix_polarize::doc}]{\sphinxcrossref{\DUrole{doc}{fix\_style polarize/functional}}}}

\item {} 
\sphinxAtStartPar
{\hyperref[\detokenize{compute_efield_atom::doc}]{\sphinxcrossref{\DUrole{doc}{compute efield/atom}}}}

\item {} 
\sphinxAtStartPar
\sphinxcode{\sphinxupquote{examples/PACKAGES/dielectric}}

\end{itemize}


\bigskip\hrule\bigskip



\subsection{DIFFRACTION package}
\label{\detokenize{Packages_details:diffraction-package}}\label{\detokenize{Packages_details:pkg-diffraction}}
\sphinxAtStartPar
\sphinxstylestrong{Contents:}

\sphinxAtStartPar
Two computes and a fix for calculating x\sphinxhyphen{}ray and electron diffraction
intensities based on kinematic diffraction theory.

\sphinxAtStartPar
\sphinxstylestrong{Author:} Shawn Coleman while at the U Arkansas.

\sphinxAtStartPar
\sphinxstylestrong{Supporting info:}
\begin{itemize}
\item {} 
\sphinxAtStartPar
\sphinxcode{\sphinxupquote{src/DIFFRACTION}}: filenames \sphinxhyphen{}\textgreater{} commands

\item {} 
\sphinxAtStartPar
{\hyperref[\detokenize{compute_saed::doc}]{\sphinxcrossref{\DUrole{doc}{compute saed}}}}

\item {} 
\sphinxAtStartPar
{\hyperref[\detokenize{compute_xrd::doc}]{\sphinxcrossref{\DUrole{doc}{compute xrd}}}}

\item {} 
\sphinxAtStartPar
{\hyperref[\detokenize{fix_saed_vtk::doc}]{\sphinxcrossref{\DUrole{doc}{fix saed/vtk}}}}

\item {} 
\sphinxAtStartPar
\sphinxcode{\sphinxupquote{examples/PACKAGES/diffraction}}

\end{itemize}


\bigskip\hrule\bigskip



\subsection{DIPOLE package}
\label{\detokenize{Packages_details:dipole-package}}\label{\detokenize{Packages_details:pkg-dipole}}
\sphinxAtStartPar
\sphinxstylestrong{Contents:}

\sphinxAtStartPar
An atom style and several pair styles for point dipole models with
short\sphinxhyphen{}range or long\sphinxhyphen{}range interactions.

\sphinxAtStartPar
\sphinxstylestrong{Supporting info:}
\begin{itemize}
\item {} 
\sphinxAtStartPar
\sphinxcode{\sphinxupquote{src/DIPOLE}}: filenames \sphinxhyphen{}\textgreater{} commands

\item {} 
\sphinxAtStartPar
{\hyperref[\detokenize{atom_style::doc}]{\sphinxcrossref{\DUrole{doc}{atom\_style dipole}}}}

\item {} 
\sphinxAtStartPar
{\hyperref[\detokenize{pair_dipole::doc}]{\sphinxcrossref{\DUrole{doc}{pair\_style lj/cut/dipole/cut}}}}

\item {} 
\sphinxAtStartPar
{\hyperref[\detokenize{pair_dipole::doc}]{\sphinxcrossref{\DUrole{doc}{pair\_style lj/cut/dipole/long}}}}

\item {} 
\sphinxAtStartPar
{\hyperref[\detokenize{pair_dipole::doc}]{\sphinxcrossref{\DUrole{doc}{pair\_style lj/long/dipole/long}}}}

\item {} 
\sphinxAtStartPar
{\hyperref[\detokenize{angle_dipole::doc}]{\sphinxcrossref{\DUrole{doc}{angle\_style dipole}}}}

\item {} 
\sphinxAtStartPar
\sphinxcode{\sphinxupquote{examples/dipole}}

\end{itemize}


\bigskip\hrule\bigskip



\subsection{DPD\sphinxhyphen{}BASIC package}
\label{\detokenize{Packages_details:dpd-basic-package}}\label{\detokenize{Packages_details:pkg-dpd-basic}}
\sphinxAtStartPar
\sphinxstylestrong{Contents:}

\sphinxAtStartPar
Pair styles for the basic dissipative particle dynamics (DPD) method
and DPD thermostatting.

\sphinxAtStartPar
Pair style {\hyperref[\detokenize{pair_dpd_coul_slater_long::doc}]{\sphinxcrossref{\DUrole{doc}{dpd/coul/slater/long}}}} also
includes smeared charges for coulomb interactions and thus requires the
{\hyperref[\detokenize{Packages_details:pkg-kspace}]{\sphinxcrossref{\DUrole{std,std-ref}{KSPACE}}}} package to be installed to handle the long\sphinxhyphen{}range
Coulomb part of the interactions.

\sphinxAtStartPar
\sphinxstylestrong{Authors:} Kurt Smith (U Pittsburgh), Martin Svoboda, Martin Lisal (ICPF and UJEP), Eddy Barraud (IFPEN)

\sphinxAtStartPar
\sphinxstylestrong{Supporting info:}
\begin{itemize}
\item {} 
\sphinxAtStartPar
\sphinxcode{\sphinxupquote{src/DPD\sphinxhyphen{}BASIC}}: filenames \sphinxhyphen{}\textgreater{} commands

\item {} 
\sphinxAtStartPar
{\hyperref[\detokenize{pair_dpd::doc}]{\sphinxcrossref{\DUrole{doc}{pair\_style dpd}}}}

\item {} 
\sphinxAtStartPar
{\hyperref[\detokenize{pair_dpd::doc}]{\sphinxcrossref{\DUrole{doc}{pair\_style dpd/tstat}}}}

\item {} 
\sphinxAtStartPar
{\hyperref[\detokenize{pair_dpd_ext::doc}]{\sphinxcrossref{\DUrole{doc}{pair\_style dpd/ext}}}}

\item {} 
\sphinxAtStartPar
{\hyperref[\detokenize{pair_dpd_ext::doc}]{\sphinxcrossref{\DUrole{doc}{pair\_style dpd/ext/tstat}}}}

\item {} 
\sphinxAtStartPar
{\hyperref[\detokenize{pair_dpd_coul_slater_long::doc}]{\sphinxcrossref{\DUrole{doc}{pair\_style dpd/coul/slater/long}}}}

\item {} 
\sphinxAtStartPar
\sphinxcode{\sphinxupquote{examples/PACKAGES/dpd\sphinxhyphen{}basic}}

\end{itemize}


\bigskip\hrule\bigskip



\subsection{DPD\sphinxhyphen{}MESO package}
\label{\detokenize{Packages_details:dpd-meso-package}}\label{\detokenize{Packages_details:pkg-dpd-meso}}
\sphinxAtStartPar
\sphinxstylestrong{Contents:}

\sphinxAtStartPar
Several extensions of the dissipative particle dynamics (DPD)
method.  Specifically, energy\sphinxhyphen{}conserving DPD (eDPD) that can model
non\sphinxhyphen{}isothermal processes, many\sphinxhyphen{}body DPD (mDPD) for simulating
vapor\sphinxhyphen{}liquid coexistence, and transport DPD (tDPD) for modeling
advection\sphinxhyphen{}diffusion\sphinxhyphen{}reaction systems. The equations of motion of these
DPD extensions are integrated through a modified velocity\sphinxhyphen{}Verlet (MVV)
algorithm.

\sphinxAtStartPar
\sphinxstylestrong{Author:} Zhen Li (Department of Mechanical Engineering, Clemson University)

\sphinxAtStartPar
\sphinxstylestrong{Supporting info:}
\begin{itemize}
\item {} 
\sphinxAtStartPar
\sphinxcode{\sphinxupquote{src/DPD\sphinxhyphen{}MESO}}: filenames \sphinxhyphen{}\textgreater{} commands

\item {} 
\sphinxAtStartPar
\sphinxcode{\sphinxupquote{src/DPD\sphinxhyphen{}MESO/README}}

\item {} 
\sphinxAtStartPar
{\hyperref[\detokenize{atom_style::doc}]{\sphinxcrossref{\DUrole{doc}{atom\_style edpd}}}}

\item {} 
\sphinxAtStartPar
{\hyperref[\detokenize{pair_mesodpd::doc}]{\sphinxcrossref{\DUrole{doc}{pair\_style edpd}}}}

\item {} 
\sphinxAtStartPar
{\hyperref[\detokenize{pair_mesodpd::doc}]{\sphinxcrossref{\DUrole{doc}{pair\_style mdpd}}}}

\item {} 
\sphinxAtStartPar
{\hyperref[\detokenize{pair_mesodpd::doc}]{\sphinxcrossref{\DUrole{doc}{pair\_style tdpd}}}}

\item {} 
\sphinxAtStartPar
{\hyperref[\detokenize{fix_mvv_dpd::doc}]{\sphinxcrossref{\DUrole{doc}{fix mvv/dpd}}}}

\item {} 
\sphinxAtStartPar
\sphinxcode{\sphinxupquote{examples/PACKAGES/mesodpd}}

\item {} 
\sphinxAtStartPar
\sphinxurl{https://www.lammps.org/movies.html\#mesodpd}

\end{itemize}


\bigskip\hrule\bigskip



\subsection{DPD\sphinxhyphen{}REACT package}
\label{\detokenize{Packages_details:dpd-react-package}}\label{\detokenize{Packages_details:pkg-dpd-react}}
\sphinxAtStartPar
\sphinxstylestrong{Contents:}

\sphinxAtStartPar
DPD stands for dissipative particle dynamics.  This package implements
coarse\sphinxhyphen{}grained DPD\sphinxhyphen{}based models for energetic, reactive molecular
crystalline materials.  It includes many pair styles specific to these
systems, including for reactive DPD, where each particle has internal
state for multiple species and a coupled set of chemical reaction ODEs
are integrated each timestep.  Highly accurate time integrators for
isothermal, isoenergetic, isobaric and isenthalpic conditions are
included.  These enable long timesteps via the Shardlow splitting
algorithm.

\sphinxAtStartPar
\sphinxstylestrong{Authors:} Jim Larentzos (ARL), Tim Mattox (Engility Corp), and John
Brennan (ARL).

\sphinxAtStartPar
\sphinxstylestrong{Supporting info:}
\begin{itemize}
\item {} 
\sphinxAtStartPar
\sphinxcode{\sphinxupquote{src/DPD\sphinxhyphen{}REACT}}: filenames \sphinxhyphen{}\textgreater{} commands

\item {} 
\sphinxAtStartPar
\sphinxcode{\sphinxupquote{src/DPD\sphinxhyphen{}REACT/README}}

\item {} 
\sphinxAtStartPar
{\hyperref[\detokenize{compute_dpd::doc}]{\sphinxcrossref{\DUrole{doc}{compute dpd}}}}

\item {} 
\sphinxAtStartPar
{\hyperref[\detokenize{compute_dpd_atom::doc}]{\sphinxcrossref{\DUrole{doc}{compute dpd/atom}}}}

\item {} 
\sphinxAtStartPar
{\hyperref[\detokenize{fix_eos_table::doc}]{\sphinxcrossref{\DUrole{doc}{fix eos/cv}}}}

\item {} 
\sphinxAtStartPar
{\hyperref[\detokenize{fix_eos_table::doc}]{\sphinxcrossref{\DUrole{doc}{fix eos/table}}}}

\item {} 
\sphinxAtStartPar
{\hyperref[\detokenize{fix_eos_table_rx::doc}]{\sphinxcrossref{\DUrole{doc}{fix eos/table/rx}}}}

\item {} 
\sphinxAtStartPar
{\hyperref[\detokenize{fix_shardlow::doc}]{\sphinxcrossref{\DUrole{doc}{fix shardlow}}}}

\item {} 
\sphinxAtStartPar
{\hyperref[\detokenize{fix_rx::doc}]{\sphinxcrossref{\DUrole{doc}{fix rx}}}}

\item {} 
\sphinxAtStartPar
{\hyperref[\detokenize{pair_table_rx::doc}]{\sphinxcrossref{\DUrole{doc}{pair\_style table/rx}}}}

\item {} 
\sphinxAtStartPar
{\hyperref[\detokenize{pair_dpd_fdt::doc}]{\sphinxcrossref{\DUrole{doc}{pair\_style dpd/fdt}}}}

\item {} 
\sphinxAtStartPar
{\hyperref[\detokenize{pair_dpd_fdt::doc}]{\sphinxcrossref{\DUrole{doc}{pair\_style dpd/fdt/energy}}}}

\item {} 
\sphinxAtStartPar
{\hyperref[\detokenize{pair_exp6_rx::doc}]{\sphinxcrossref{\DUrole{doc}{pair\_style exp6/rx}}}}

\item {} 
\sphinxAtStartPar
{\hyperref[\detokenize{pair_multi_lucy::doc}]{\sphinxcrossref{\DUrole{doc}{pair\_style multi/lucy}}}}

\item {} 
\sphinxAtStartPar
{\hyperref[\detokenize{pair_multi_lucy_rx::doc}]{\sphinxcrossref{\DUrole{doc}{pair\_style multi/lucy/rx}}}}

\item {} 
\sphinxAtStartPar
\sphinxcode{\sphinxupquote{examples/PACKAGES/dpd\sphinxhyphen{}react}}

\end{itemize}


\bigskip\hrule\bigskip



\subsection{DPD\sphinxhyphen{}SMOOTH package}
\label{\detokenize{Packages_details:dpd-smooth-package}}\label{\detokenize{Packages_details:pkg-dpd-smooth}}
\sphinxAtStartPar
\sphinxstylestrong{Contents:}

\sphinxAtStartPar
A pair style for smoothed dissipative particle dynamics (SDPD), which
is an extension of smoothed particle hydrodynamics (SPH) to mesoscale
where thermal fluctuations are important (see the
{\hyperref[\detokenize{Packages_details:pkg-sph}]{\sphinxcrossref{\DUrole{std,std-ref}{SPH package}}}}).
Also two fixes for moving and rigid body integration of SPH/SDPD particles
(particles of atom\_style meso).

\sphinxAtStartPar
\sphinxstylestrong{Author:} Morteza Jalalvand (Institute for Advanced Studies in Basic
Sciences, Iran).

\sphinxAtStartPar
\sphinxstylestrong{Supporting info:}
\begin{itemize}
\item {} 
\sphinxAtStartPar
\sphinxcode{\sphinxupquote{src/DPD\sphinxhyphen{}SMOOTH}}: filenames \sphinxhyphen{}\textgreater{} commands

\item {} 
\sphinxAtStartPar
\sphinxcode{\sphinxupquote{src/DPD\sphinxhyphen{}SMOOTH/README}}

\item {} 
\sphinxAtStartPar
{\hyperref[\detokenize{pair_sdpd_taitwater_isothermal::doc}]{\sphinxcrossref{\DUrole{doc}{pair\_style sdpd/taitwater/isothermal}}}}

\item {} 
\sphinxAtStartPar
{\hyperref[\detokenize{fix_meso_move::doc}]{\sphinxcrossref{\DUrole{doc}{fix meso/move}}}}

\item {} 
\sphinxAtStartPar
{\hyperref[\detokenize{fix_rigid_meso::doc}]{\sphinxcrossref{\DUrole{doc}{fix rigid/meso}}}}

\item {} 
\sphinxAtStartPar
\sphinxcode{\sphinxupquote{examples/PACKAGES/dpd\sphinxhyphen{}smooth}}

\end{itemize}


\bigskip\hrule\bigskip



\subsection{DRUDE package}
\label{\detokenize{Packages_details:drude-package}}\label{\detokenize{Packages_details:pkg-drude}}
\sphinxAtStartPar
\sphinxstylestrong{Contents:}

\sphinxAtStartPar
Fixes, pair styles, and a compute to simulate thermalized Drude
oscillators as a model of polarization.  See the {\hyperref[\detokenize{Howto_drude::doc}]{\sphinxcrossref{\DUrole{doc}{Howto drude}}}} and {\hyperref[\detokenize{Howto_drude2::doc}]{\sphinxcrossref{\DUrole{doc}{Howto drude2}}}} pages
for an overview of how to use the package.  There are auxiliary tools
for using this package in tools/drude.

\sphinxAtStartPar
\sphinxstylestrong{Authors:} Alain Dequidt (U Clermont Auvergne), Julien
Devemy (CNRS), and Agilio Padua (ENS de Lyon).

\sphinxAtStartPar
\sphinxstylestrong{Supporting info:}
\begin{itemize}
\item {} 
\sphinxAtStartPar
\sphinxcode{\sphinxupquote{src/DRUDE}}: filenames \sphinxhyphen{}\textgreater{} commands

\item {} 
\sphinxAtStartPar
{\hyperref[\detokenize{Howto_drude::doc}]{\sphinxcrossref{\DUrole{doc}{Howto drude}}}}

\item {} 
\sphinxAtStartPar
{\hyperref[\detokenize{Howto_drude2::doc}]{\sphinxcrossref{\DUrole{doc}{Howto drude2}}}}

\item {} 
\sphinxAtStartPar
{\hyperref[\detokenize{Howto_polarizable::doc}]{\sphinxcrossref{\DUrole{doc}{Howto polarizable}}}}

\item {} 
\sphinxAtStartPar
\sphinxcode{\sphinxupquote{src/DRUDE/README}}

\item {} 
\sphinxAtStartPar
{\hyperref[\detokenize{fix_drude::doc}]{\sphinxcrossref{\DUrole{doc}{fix drude}}}}

\item {} 
\sphinxAtStartPar
{\hyperref[\detokenize{fix_drude_transform::doc}]{\sphinxcrossref{\DUrole{doc}{fix drude/transform/*}}}}

\item {} 
\sphinxAtStartPar
{\hyperref[\detokenize{compute_temp_drude::doc}]{\sphinxcrossref{\DUrole{doc}{compute temp/drude}}}}

\item {} 
\sphinxAtStartPar
{\hyperref[\detokenize{pair_thole::doc}]{\sphinxcrossref{\DUrole{doc}{pair\_style thole}}}}

\item {} 
\sphinxAtStartPar
{\hyperref[\detokenize{pair_thole::doc}]{\sphinxcrossref{\DUrole{doc}{pair\_style lj/cut/thole/long}}}}

\item {} 
\sphinxAtStartPar
\sphinxcode{\sphinxupquote{examples/PACKAGES/drude}}

\item {} 
\sphinxAtStartPar
tools/drude

\end{itemize}


\bigskip\hrule\bigskip



\subsection{EFF package}
\label{\detokenize{Packages_details:eff-package}}\label{\detokenize{Packages_details:pkg-eff}}
\sphinxAtStartPar
\sphinxstylestrong{Contents:}

\sphinxAtStartPar
EFF stands for electron force field which allows a classical MD code
to model electrons as particles of variable radius.  This package
contains atom, pair, fix and compute styles which implement the eFF as
described in A. Jaramillo\sphinxhyphen{}Botero, J. Su, Q. An, and W.A. Goddard III,
JCC, 2010.  The eFF potential was first introduced by Su and Goddard,
in 2007.  There are auxiliary tools for using this package in
tools/eff; see its README file.

\sphinxAtStartPar
\sphinxstylestrong{Author:} Andres Jaramillo\sphinxhyphen{}Botero (CalTech).

\sphinxAtStartPar
\sphinxstylestrong{Supporting info:}
\begin{itemize}
\item {} 
\sphinxAtStartPar
\sphinxcode{\sphinxupquote{src/EFF}}: filenames \sphinxhyphen{}\textgreater{} commands

\item {} 
\sphinxAtStartPar
\sphinxcode{\sphinxupquote{src/EFF/README}}

\item {} 
\sphinxAtStartPar
{\hyperref[\detokenize{atom_style::doc}]{\sphinxcrossref{\DUrole{doc}{atom\_style electron}}}}

\item {} 
\sphinxAtStartPar
{\hyperref[\detokenize{fix_nve_eff::doc}]{\sphinxcrossref{\DUrole{doc}{fix nve/eff}}}}

\item {} 
\sphinxAtStartPar
{\hyperref[\detokenize{fix_nh_eff::doc}]{\sphinxcrossref{\DUrole{doc}{fix nvt/eff}}}}

\item {} 
\sphinxAtStartPar
{\hyperref[\detokenize{fix_nh_eff::doc}]{\sphinxcrossref{\DUrole{doc}{fix npt/eff}}}}

\item {} 
\sphinxAtStartPar
{\hyperref[\detokenize{fix_langevin_eff::doc}]{\sphinxcrossref{\DUrole{doc}{fix langevin/eff}}}}

\item {} 
\sphinxAtStartPar
{\hyperref[\detokenize{compute_temp_eff::doc}]{\sphinxcrossref{\DUrole{doc}{compute temp/eff}}}}

\item {} 
\sphinxAtStartPar
{\hyperref[\detokenize{pair_eff::doc}]{\sphinxcrossref{\DUrole{doc}{pair\_style eff/cut}}}}

\item {} 
\sphinxAtStartPar
{\hyperref[\detokenize{pair_eff::doc}]{\sphinxcrossref{\DUrole{doc}{pair\_style eff/inline}}}}

\item {} 
\sphinxAtStartPar
\sphinxcode{\sphinxupquote{examples/PACKAGES/eff}}

\item {} 
\sphinxAtStartPar
tools/eff/README

\item {} 
\sphinxAtStartPar
tools/eff

\item {} 
\sphinxAtStartPar
\sphinxurl{https://www.lammps.org/movies.html\#eff}

\end{itemize}


\bigskip\hrule\bigskip



\subsection{ELECTRODE package}
\label{\detokenize{Packages_details:electrode-package}}\label{\detokenize{Packages_details:pkg-electrode}}
\sphinxAtStartPar
\sphinxstylestrong{Contents:}

\sphinxAtStartPar
The ELECTRODE package allows the user to enforce a constant potential method for
groups of atoms that interact with the remaining atoms as electrolyte.

\sphinxAtStartPar
\sphinxstylestrong{Authors:} The ELECTRODE package is written and maintained by Ludwig
Ahrens\sphinxhyphen{}Iwers (TUHH, Hamburg, Germany), Shern Tee (UQ, Brisbane, Australia) and
Robert Meissner (TUHH, Hamburg, Germany).

\sphinxAtStartPar
\DUrole{versionmodified,added}{Added in version 4May2022.}

\sphinxAtStartPar
\sphinxstylestrong{Install:}

\sphinxAtStartPar
This package has {\hyperref[\detokenize{Build_extras:electrode}]{\sphinxcrossref{\DUrole{std,std-ref}{specific installation instructions}}}} on the
{\hyperref[\detokenize{Build_extras::doc}]{\sphinxcrossref{\DUrole{doc}{Build extras}}}} page.

\sphinxAtStartPar
\sphinxstylestrong{Supporting info:}
\begin{itemize}
\item {} 
\sphinxAtStartPar
{\hyperref[\detokenize{fix_electrode::doc}]{\sphinxcrossref{\DUrole{doc}{fix electrode/conp}}}}

\item {} 
\sphinxAtStartPar
{\hyperref[\detokenize{fix_electrode::doc}]{\sphinxcrossref{\DUrole{doc}{fix electrode/conq}}}}

\item {} 
\sphinxAtStartPar
{\hyperref[\detokenize{fix_electrode::doc}]{\sphinxcrossref{\DUrole{doc}{fix electrode/thermo}}}}

\end{itemize}


\bigskip\hrule\bigskip



\subsection{EXTRA\sphinxhyphen{}COMMAND package}
\label{\detokenize{Packages_details:extra-command-package}}\label{\detokenize{Packages_details:pkg-extra-command}}
\sphinxAtStartPar
\sphinxstylestrong{Contents:}

\sphinxAtStartPar
Additional command styles that are less commonly used.

\sphinxAtStartPar
\sphinxstylestrong{Supporting info:}
\begin{itemize}
\item {} 
\sphinxAtStartPar
\sphinxcode{\sphinxupquote{src/EXTRA\sphinxhyphen{}COMMAND}}: filenames \sphinxhyphen{}\textgreater{} commands

\item {} 
\sphinxAtStartPar
{\hyperref[\detokenize{Commands_all::doc}]{\sphinxcrossref{\DUrole{doc}{general commands}}}}

\end{itemize}


\bigskip\hrule\bigskip



\subsection{EXTRA\sphinxhyphen{}COMPUTE package}
\label{\detokenize{Packages_details:extra-compute-package}}\label{\detokenize{Packages_details:pkg-extra-compute}}
\sphinxAtStartPar
\sphinxstylestrong{Contents:}

\sphinxAtStartPar
Additional compute styles that are less commonly used.

\sphinxAtStartPar
\sphinxstylestrong{Supporting info:}
\begin{itemize}
\item {} 
\sphinxAtStartPar
\sphinxcode{\sphinxupquote{src/EXTRA\sphinxhyphen{}COMPUTE}}: filenames \sphinxhyphen{}\textgreater{} commands

\item {} 
\sphinxAtStartPar
{\hyperref[\detokenize{compute::doc}]{\sphinxcrossref{\DUrole{doc}{compute}}}}

\end{itemize}


\bigskip\hrule\bigskip



\subsection{EXTRA\sphinxhyphen{}DUMP package}
\label{\detokenize{Packages_details:extra-dump-package}}\label{\detokenize{Packages_details:pkg-extra-dump}}
\sphinxAtStartPar
\sphinxstylestrong{Contents:}

\sphinxAtStartPar
Additional dump styles that are less commonly used.

\sphinxAtStartPar
\sphinxstylestrong{Supporting info:}
\begin{itemize}
\item {} 
\sphinxAtStartPar
\sphinxcode{\sphinxupquote{src/EXTRA\sphinxhyphen{}DUMP}}: filenames \sphinxhyphen{}\textgreater{} commands

\item {} 
\sphinxAtStartPar
{\hyperref[\detokenize{dump::doc}]{\sphinxcrossref{\DUrole{doc}{dump}}}}

\end{itemize}


\bigskip\hrule\bigskip



\subsection{EXTRA\sphinxhyphen{}FIX package}
\label{\detokenize{Packages_details:extra-fix-package}}\label{\detokenize{Packages_details:pkg-extra-fix}}
\sphinxAtStartPar
\sphinxstylestrong{Contents:}

\sphinxAtStartPar
Additional fix styles that are less commonly used.

\sphinxAtStartPar
\sphinxstylestrong{Supporting info:}
\begin{itemize}
\item {} 
\sphinxAtStartPar
\sphinxcode{\sphinxupquote{src/EXTRA\sphinxhyphen{}FIX}}: filenames \sphinxhyphen{}\textgreater{} commands

\item {} 
\sphinxAtStartPar
{\hyperref[\detokenize{fix::doc}]{\sphinxcrossref{\DUrole{doc}{fix}}}}

\end{itemize}


\bigskip\hrule\bigskip



\subsection{EXTRA\sphinxhyphen{}MOLECULE package}
\label{\detokenize{Packages_details:extra-molecule-package}}\label{\detokenize{Packages_details:pkg-extra-molecule}}
\sphinxAtStartPar
\sphinxstylestrong{Contents:}

\sphinxAtStartPar
Additional bond, angle, dihedral, and improper styles that are less commonly used.

\sphinxAtStartPar
\sphinxstylestrong{Install:}

\sphinxAtStartPar
To use this package, also the {\hyperref[\detokenize{Packages_details:pkg-molecule}]{\sphinxcrossref{\DUrole{std,std-ref}{MOLECULE}}}} package needs to be installed.

\sphinxAtStartPar
\sphinxstylestrong{Supporting info:}
\begin{itemize}
\item {} 
\sphinxAtStartPar
\sphinxcode{\sphinxupquote{src/EXTRA\sphinxhyphen{}MOLECULE}}: filenames \sphinxhyphen{}\textgreater{} commands

\item {} 
\sphinxAtStartPar
{\hyperref[\detokenize{Commands_bond::doc}]{\sphinxcrossref{\DUrole{doc}{molecular styles}}}}

\end{itemize}


\bigskip\hrule\bigskip



\subsection{EXTRA\sphinxhyphen{}PAIR package}
\label{\detokenize{Packages_details:extra-pair-package}}\label{\detokenize{Packages_details:pkg-extra-pair}}
\sphinxAtStartPar
\sphinxstylestrong{Contents:}

\sphinxAtStartPar
Additional pair styles that are less commonly used.

\sphinxAtStartPar
\sphinxstylestrong{Supporting info:}
\begin{itemize}
\item {} 
\sphinxAtStartPar
\sphinxcode{\sphinxupquote{src/EXTRA\sphinxhyphen{}PAIR}}: filenames \sphinxhyphen{}\textgreater{} commands

\item {} 
\sphinxAtStartPar
{\hyperref[\detokenize{pair_style::doc}]{\sphinxcrossref{\DUrole{doc}{pair\_style}}}}

\end{itemize}


\bigskip\hrule\bigskip



\subsection{FEP package}
\label{\detokenize{Packages_details:fep-package}}\label{\detokenize{Packages_details:pkg-fep}}
\sphinxAtStartPar
\sphinxstylestrong{Contents:}

\sphinxAtStartPar
FEP stands for free energy perturbation.  This package provides methods
for performing FEP simulations by using a {\hyperref[\detokenize{fix_adapt_fep::doc}]{\sphinxcrossref{\DUrole{doc}{fix adapt/fep}}}} command with soft\sphinxhyphen{}core pair potentials, which have a
“soft” in their style name.  There are auxiliary tools for using this
package in \sphinxcode{\sphinxupquote{tools/fep}}; see its \sphinxcode{\sphinxupquote{README}} file.

\sphinxAtStartPar
\sphinxstylestrong{Author:} Agilio Padua (ENS de Lyon)

\sphinxAtStartPar
\sphinxstylestrong{Supporting info:}
\begin{itemize}
\item {} 
\sphinxAtStartPar
\sphinxcode{\sphinxupquote{src/FEP}}: filenames \sphinxhyphen{}\textgreater{} commands

\item {} 
\sphinxAtStartPar
\sphinxcode{\sphinxupquote{src/FEP/README}}

\item {} 
\sphinxAtStartPar
{\hyperref[\detokenize{fix_adapt_fep::doc}]{\sphinxcrossref{\DUrole{doc}{fix adapt/fep}}}}

\item {} 
\sphinxAtStartPar
{\hyperref[\detokenize{compute_fep::doc}]{\sphinxcrossref{\DUrole{doc}{compute fep}}}}

\item {} 
\sphinxAtStartPar
{\hyperref[\detokenize{pair_fep_soft::doc}]{\sphinxcrossref{\DUrole{doc}{pair\_style */soft}}}}

\item {} 
\sphinxAtStartPar
\sphinxcode{\sphinxupquote{examples/PACKAGES/fep}}

\item {} 
\sphinxAtStartPar
tools/fep/README

\item {} 
\sphinxAtStartPar
tools/fep

\end{itemize}


\bigskip\hrule\bigskip



\subsection{GPU package}
\label{\detokenize{Packages_details:gpu-package}}\label{\detokenize{Packages_details:pkg-gpu}}
\sphinxAtStartPar
\sphinxstylestrong{Contents:}

\sphinxAtStartPar
Dozens of pair styles and a version of the PPPM long\sphinxhyphen{}range Coulombic
solver optimized for GPUs.  All such styles have a “gpu” as a suffix
in their style name. The GPU code can be compiled with either CUDA or
OpenCL, however the OpenCL variants are no longer actively maintained
and only the CUDA versions are regularly tested.  The
{\hyperref[\detokenize{Speed_gpu::doc}]{\sphinxcrossref{\DUrole{doc}{GPU package}}}} page gives details of what hardware and GPU
software is required on your system, and details on how to build and
use this package.  Its styles can be invoked at run time via the \sphinxcode{\sphinxupquote{\sphinxhyphen{}sf
gpu}} or \sphinxcode{\sphinxupquote{\sphinxhyphen{}suffix gpu}} {\hyperref[\detokenize{Run_options::doc}]{\sphinxcrossref{\DUrole{doc}{command\sphinxhyphen{}line switches}}}}.  See
also the {\hyperref[\detokenize{Packages_details:pkg-kokkos}]{\sphinxcrossref{\DUrole{std,std-ref}{KOKKOS}}}} package, which has GPU\sphinxhyphen{}enabled styles.

\sphinxAtStartPar
\sphinxstylestrong{Authors:} Mike Brown (Intel) while at Sandia and ORNL and Trung Nguyen
(Northwestern U) while at ORNL and later. AMD HIP support by Evgeny
Kuznetsov, Vladimir Stegailov, and Vsevolod Nikolskiy (HSE University).

\sphinxAtStartPar
\sphinxstylestrong{Install:}

\sphinxAtStartPar
This package has {\hyperref[\detokenize{Build_extras:gpu}]{\sphinxcrossref{\DUrole{std,std-ref}{specific installation instructions}}}} on the
{\hyperref[\detokenize{Build_extras::doc}]{\sphinxcrossref{\DUrole{doc}{Build extras}}}} page.

\sphinxAtStartPar
\sphinxstylestrong{Supporting info:}
\begin{itemize}
\item {} 
\sphinxAtStartPar
\sphinxcode{\sphinxupquote{src/GPU}}: filenames \sphinxhyphen{}\textgreater{} commands

\item {} 
\sphinxAtStartPar
\sphinxcode{\sphinxupquote{src/GPU/README}}

\item {} 
\sphinxAtStartPar
\sphinxcode{\sphinxupquote{lib/gpu/README}}

\item {} 
\sphinxAtStartPar
{\hyperref[\detokenize{Speed_packages::doc}]{\sphinxcrossref{\DUrole{doc}{Accelerator packages}}}}

\item {} 
\sphinxAtStartPar
{\hyperref[\detokenize{Speed_gpu::doc}]{\sphinxcrossref{\DUrole{doc}{GPU package}}}}

\item {} 
\sphinxAtStartPar
{\hyperref[\detokenize{Run_options::doc}]{\sphinxcrossref{\DUrole{doc}{Section 2.6 \sphinxhyphen{}sf gpu}}}}

\item {} 
\sphinxAtStartPar
{\hyperref[\detokenize{Run_options::doc}]{\sphinxcrossref{\DUrole{doc}{Section 2.6 \sphinxhyphen{}pk gpu}}}}

\item {} 
\sphinxAtStartPar
{\hyperref[\detokenize{package::doc}]{\sphinxcrossref{\DUrole{doc}{package gpu}}}}

\item {} 
\sphinxAtStartPar
{\hyperref[\detokenize{Commands_all::doc}]{\sphinxcrossref{\DUrole{doc}{Commands}}}} pages ({\hyperref[\detokenize{Commands_pair::doc}]{\sphinxcrossref{\DUrole{doc}{pair}}}}, {\hyperref[\detokenize{Commands_kspace::doc}]{\sphinxcrossref{\DUrole{doc}{kspace}}}})
for styles followed by (g)

\item {} 
\sphinxAtStartPar
\sphinxhref{https://www.lammps.org/bench.html}{Benchmarks page} of website

\end{itemize}


\bigskip\hrule\bigskip



\subsection{GRANULAR package}
\label{\detokenize{Packages_details:granular-package}}\label{\detokenize{Packages_details:pkg-granular}}
\sphinxAtStartPar
\sphinxstylestrong{Contents:}

\sphinxAtStartPar
Pair styles and fixes for finite\sphinxhyphen{}size granular particles, which
interact with each other and boundaries via frictional and dissipative
potentials.

\sphinxAtStartPar
\sphinxstylestrong{Supporting info:}
\begin{itemize}
\item {} 
\sphinxAtStartPar
\sphinxcode{\sphinxupquote{src/GRANULAR}}: filenames \sphinxhyphen{}\textgreater{} commands

\item {} 
\sphinxAtStartPar
{\hyperref[\detokenize{Howto_granular::doc}]{\sphinxcrossref{\DUrole{doc}{Howto granular}}}}

\item {} 
\sphinxAtStartPar
{\hyperref[\detokenize{fix_pour::doc}]{\sphinxcrossref{\DUrole{doc}{fix pour}}}}

\item {} 
\sphinxAtStartPar
{\hyperref[\detokenize{fix_wall_gran::doc}]{\sphinxcrossref{\DUrole{doc}{fix wall/gran}}}}

\item {} 
\sphinxAtStartPar
{\hyperref[\detokenize{pair_gran::doc}]{\sphinxcrossref{\DUrole{doc}{pair\_style gran/hooke}}}}

\item {} 
\sphinxAtStartPar
{\hyperref[\detokenize{pair_gran::doc}]{\sphinxcrossref{\DUrole{doc}{pair\_style gran/hertz/history}}}}

\item {} 
\sphinxAtStartPar
\sphinxcode{\sphinxupquote{examples/granregion}}

\item {} 
\sphinxAtStartPar
\sphinxcode{\sphinxupquote{examples/pour}}

\item {} 
\sphinxAtStartPar
bench/in.chute

\item {} 
\sphinxAtStartPar
\sphinxurl{https://www.lammps.org/pictures.html\#jamming}

\item {} 
\sphinxAtStartPar
\sphinxurl{https://www.lammps.org/movies.html\#hopper}

\item {} 
\sphinxAtStartPar
\sphinxurl{https://www.lammps.org/movies.html\#dem}

\item {} 
\sphinxAtStartPar
\sphinxurl{https://www.lammps.org/movies.html\#brazil}

\item {} 
\sphinxAtStartPar
\sphinxurl{https://www.lammps.org/movies.html\#granregion}

\end{itemize}


\bigskip\hrule\bigskip



\subsection{H5MD package}
\label{\detokenize{Packages_details:h5md-package}}\label{\detokenize{Packages_details:pkg-h5md}}
\sphinxAtStartPar
\sphinxstylestrong{Contents:}

\sphinxAtStartPar
H5MD stands for HDF5 for MD.  \sphinxhref{https://www.hdfgroup.org/solutions/hdf5}{HDF5} is a portable, binary,
self\sphinxhyphen{}describing file format, used by many scientific simulations.
H5MD is a format for molecular simulations, built on top of HDF5.
This package implements a {\hyperref[\detokenize{dump_h5md::doc}]{\sphinxcrossref{\DUrole{doc}{dump h5md}}}} command to output
LAMMPS snapshots in this format.

\sphinxAtStartPar
To use this package you must have the HDF5 library available on your
system.

\sphinxAtStartPar
\sphinxstylestrong{Author:} Pierre de Buyl (KU Leuven) created both the package and the
H5MD format.

\sphinxAtStartPar
\sphinxstylestrong{Install:}

\sphinxAtStartPar
This package has {\hyperref[\detokenize{Build_extras:h5md}]{\sphinxcrossref{\DUrole{std,std-ref}{specific installation instructions}}}} on the {\hyperref[\detokenize{Build_extras::doc}]{\sphinxcrossref{\DUrole{doc}{Build extras}}}} page.

\sphinxAtStartPar
\sphinxstylestrong{Supporting info:}
\begin{itemize}
\item {} 
\sphinxAtStartPar
\sphinxcode{\sphinxupquote{src/H5MD}}: filenames \sphinxhyphen{}\textgreater{} commands

\item {} 
\sphinxAtStartPar
\sphinxcode{\sphinxupquote{src/H5MD/README}}

\item {} 
\sphinxAtStartPar
\sphinxcode{\sphinxupquote{lib/h5md/README}}

\item {} 
\sphinxAtStartPar
{\hyperref[\detokenize{dump_h5md::doc}]{\sphinxcrossref{\DUrole{doc}{dump h5md}}}}

\end{itemize}


\bigskip\hrule\bigskip



\subsection{INTEL package}
\label{\detokenize{Packages_details:intel-package}}\label{\detokenize{Packages_details:pkg-intel}}
\sphinxAtStartPar
\sphinxstylestrong{Contents:}

\sphinxAtStartPar
Dozens of pair, fix, bond, angle, dihedral, improper, and kspace styles
which are optimized for Intel CPUs and KNLs (Knights Landing).  All of
them have an “intel” in their style name.  The {\hyperref[\detokenize{Speed_intel::doc}]{\sphinxcrossref{\DUrole{doc}{INTEL package}}}} page gives details of what hardware and compilers are
required on your system, and how to build and use this package.  Its
styles can be invoked at run time via the \sphinxcode{\sphinxupquote{\sphinxhyphen{}sf intel}} or \sphinxcode{\sphinxupquote{\sphinxhyphen{}suffix intel}}
{\hyperref[\detokenize{Run_options::doc}]{\sphinxcrossref{\DUrole{doc}{command\sphinxhyphen{}line switches}}}}.  Also see the {\hyperref[\detokenize{Packages_details:pkg-kokkos}]{\sphinxcrossref{\DUrole{std,std-ref}{KOKKOS}}}}, {\hyperref[\detokenize{Packages_details:pkg-opt}]{\sphinxcrossref{\DUrole{std,std-ref}{OPT}}}}, and {\hyperref[\detokenize{Packages_details:pkg-openmp}]{\sphinxcrossref{\DUrole{std,std-ref}{OPENMP}}}}
packages, which have styles optimized for CPUs and KNLs.

\sphinxAtStartPar
You need to have an Intel compiler, version 14 or higher to take full
advantage of this package. While compilation with GNU compilers is
supported, performance will be sub\sphinxhyphen{}optimal.

\begin{sphinxadmonition}{note}{Note:}
\sphinxAtStartPar
the INTEL package contains styles that require using the
\sphinxhyphen{}restrict flag, when compiling with Intel compilers.
\end{sphinxadmonition}

\sphinxAtStartPar
\sphinxstylestrong{Author:} Mike Brown (Intel).

\sphinxAtStartPar
\sphinxstylestrong{Install:}

\sphinxAtStartPar
This package has {\hyperref[\detokenize{Build_extras:intel}]{\sphinxcrossref{\DUrole{std,std-ref}{specific installation instructions}}}} on the {\hyperref[\detokenize{Build_extras::doc}]{\sphinxcrossref{\DUrole{doc}{Build extras}}}} page.

\sphinxAtStartPar
\sphinxstylestrong{Supporting info:}
\begin{itemize}
\item {} 
\sphinxAtStartPar
\sphinxcode{\sphinxupquote{src/INTEL}}: filenames \sphinxhyphen{}\textgreater{} commands

\item {} 
\sphinxAtStartPar
\sphinxcode{\sphinxupquote{src/INTEL/README}}

\item {} 
\sphinxAtStartPar
{\hyperref[\detokenize{Speed_packages::doc}]{\sphinxcrossref{\DUrole{doc}{Accelerator packages}}}}

\item {} 
\sphinxAtStartPar
{\hyperref[\detokenize{Speed_intel::doc}]{\sphinxcrossref{\DUrole{doc}{INTEL package}}}}

\item {} 
\sphinxAtStartPar
{\hyperref[\detokenize{Run_options::doc}]{\sphinxcrossref{\DUrole{doc}{Section 2.6 \sphinxhyphen{}sf intel}}}}

\item {} 
\sphinxAtStartPar
{\hyperref[\detokenize{Run_options::doc}]{\sphinxcrossref{\DUrole{doc}{Section 2.6 \sphinxhyphen{}pk intel}}}}

\item {} 
\sphinxAtStartPar
{\hyperref[\detokenize{package::doc}]{\sphinxcrossref{\DUrole{doc}{package intel}}}}

\item {} 
\sphinxAtStartPar
Search the {\hyperref[\detokenize{Commands_all::doc}]{\sphinxcrossref{\DUrole{doc}{commands}}}} pages ({\hyperref[\detokenize{Commands_fix::doc}]{\sphinxcrossref{\DUrole{doc}{fix}}}}, {\hyperref[\detokenize{Commands_compute::doc}]{\sphinxcrossref{\DUrole{doc}{compute}}}},
{\hyperref[\detokenize{Commands_pair::doc}]{\sphinxcrossref{\DUrole{doc}{pair}}}}, {\hyperref[\detokenize{Commands_bond::doc}]{\sphinxcrossref{\DUrole{doc}{bond, angle, dihedral, improper}}}}, {\hyperref[\detokenize{Commands_kspace::doc}]{\sphinxcrossref{\DUrole{doc}{kspace}}}}) for styles followed by (i)

\item {} 
\sphinxAtStartPar
\sphinxcode{\sphinxupquote{src/INTEL/TEST}}

\item {} 
\sphinxAtStartPar
\sphinxhref{https://www.lammps.org/bench.html}{Benchmarks page} of website

\end{itemize}


\bigskip\hrule\bigskip



\subsection{INTERLAYER package}
\label{\detokenize{Packages_details:interlayer-package}}\label{\detokenize{Packages_details:pkg-interlayer}}
\sphinxAtStartPar
\sphinxstylestrong{Contents:}

\sphinxAtStartPar
A collection of pair styles specifically to be used for modeling layered
materials, most commonly graphene sheets (or equivalents).

\sphinxAtStartPar
\sphinxstylestrong{Supporting info:}
\begin{itemize}
\item {} 
\sphinxAtStartPar
\sphinxcode{\sphinxupquote{src/INTERLAYER}}: filenames \sphinxhyphen{}\textgreater{} commands

\item {} 
\sphinxAtStartPar
{\hyperref[\detokenize{Commands_pair::doc}]{\sphinxcrossref{\DUrole{doc}{Pair style}}}} page

\item {} 
\sphinxAtStartPar
\sphinxcode{\sphinxupquote{examples/PACKAGES/interlayer}}

\end{itemize}


\bigskip\hrule\bigskip



\subsection{KIM package}
\label{\detokenize{Packages_details:kim-package}}\label{\detokenize{Packages_details:pkg-kim}}
\sphinxAtStartPar
\sphinxstylestrong{Contents:}

\sphinxAtStartPar
This package contains a command with a set of sub\sphinxhyphen{}commands that serve as a
wrapper on the
\sphinxhref{https://openkim.org}{Open Knowledgebase of Interatomic Models (OpenKIM)}
repository of interatomic models (IMs) enabling compatible ones to be used in
LAMMPS simulations.

\sphinxAtStartPar
This includes {\hyperref[\detokenize{kim_commands::doc}]{\sphinxcrossref{\DUrole{doc}{kim init}}}}, and
{\hyperref[\detokenize{kim_commands::doc}]{\sphinxcrossref{\DUrole{doc}{kim interactions}}}} commands to select, initialize and
instantiate the IM, a {\hyperref[\detokenize{kim_commands::doc}]{\sphinxcrossref{\DUrole{doc}{kim query}}}} command to perform web
queries for material property predictions of OpenKIM IMs, a
{\hyperref[\detokenize{kim_commands::doc}]{\sphinxcrossref{\DUrole{doc}{kim param}}}} command to access KIM Model Parameters from
LAMMPS, and a {\hyperref[\detokenize{kim_commands::doc}]{\sphinxcrossref{\DUrole{doc}{kim property}}}} command to write material
properties computed in LAMMPS to standard KIM property instance format.

\sphinxAtStartPar
Support for KIM IMs that conform to the
\sphinxhref{https://openkim.org/kim-api/}{KIM Application Programming Interface (API)}
is provided by the {\hyperref[\detokenize{pair_kim::doc}]{\sphinxcrossref{\DUrole{doc}{pair\_style kim}}}} command.

\begin{sphinxadmonition}{note}{Note:}
\sphinxAtStartPar
The command \sphinxstyleemphasis{pair\_style kim} is called by \sphinxstyleemphasis{kim interactions} and is not
recommended to be directly used in input scripts.
\end{sphinxadmonition}

\sphinxAtStartPar
To use this package you must have the KIM API library available on your
system. The KIM API is available for download on the
\sphinxhref{https://openkim.org/kim-api/}{OpenKIM website}.
When installing LAMMPS from binary, the kim\sphinxhyphen{}api package
is a dependency that is automatically downloaded and installed.

\sphinxAtStartPar
Information about the KIM project can be found at its website:
\sphinxurl{https://openkim.org}.
The KIM project is led by Ellad Tadmor and Ryan Elliott (U Minnesota)
and is funded by the \sphinxhref{https://www.nsf.gov/}{National Science Foundation}.

\sphinxAtStartPar
\sphinxstylestrong{Authors:} Ryan Elliott (U Minnesota) is the main developer for the KIM
API and the \sphinxstyleemphasis{pair\_style kim} command. Yaser Afshar (U Minnesota),
Axel Kohlmeyer (Temple U), Ellad Tadmor (U Minnesota), and
Daniel Karls (U Minnesota) contributed to the
{\hyperref[\detokenize{kim_commands::doc}]{\sphinxcrossref{\DUrole{doc}{kim command}}}} interface in close collaboration with
Ryan Elliott.

\sphinxAtStartPar
\sphinxstylestrong{Install:}

\sphinxAtStartPar
This package has {\hyperref[\detokenize{Build_extras:kim}]{\sphinxcrossref{\DUrole{std,std-ref}{specific installation instructions}}}} on the
{\hyperref[\detokenize{Build_extras::doc}]{\sphinxcrossref{\DUrole{doc}{Build extras}}}} page.

\sphinxAtStartPar
\sphinxstylestrong{Supporting info:}
\begin{itemize}
\item {} 
\sphinxAtStartPar
{\hyperref[\detokenize{kim_commands::doc}]{\sphinxcrossref{\DUrole{doc}{kim command}}}}

\item {} 
\sphinxAtStartPar
{\hyperref[\detokenize{pair_kim::doc}]{\sphinxcrossref{\DUrole{doc}{pair\_style kim}}}}

\item {} 
\sphinxAtStartPar
\sphinxcode{\sphinxupquote{src/KIM}}: filenames \sphinxhyphen{}\textgreater{} commands

\item {} 
\sphinxAtStartPar
\sphinxcode{\sphinxupquote{src/KIM/README}}

\item {} 
\sphinxAtStartPar
\sphinxcode{\sphinxupquote{lib/kim/README}}

\item {} 
\sphinxAtStartPar
\sphinxcode{\sphinxupquote{examples/kim}}

\end{itemize}


\bigskip\hrule\bigskip



\subsection{KOKKOS package}
\label{\detokenize{Packages_details:kokkos-package}}\label{\detokenize{Packages_details:pkg-kokkos}}
\sphinxAtStartPar
\sphinxstylestrong{Contents:}

\sphinxAtStartPar
Dozens of atom, pair, bond, angle, dihedral, improper, fix, compute
styles adapted to compile using the Kokkos library which can convert
them to OpenMP or CUDA code so that they run efficiently on multicore
CPUs, KNLs, or GPUs.  All the styles have a “kk” as a suffix in their
style name.  The {\hyperref[\detokenize{Speed_kokkos::doc}]{\sphinxcrossref{\DUrole{doc}{KOKKOS package}}}} page gives details
of what hardware and software is required on your system, and how to
build and use this package.  Its styles can be invoked at run time via
the \sphinxcode{\sphinxupquote{\sphinxhyphen{}sf kk}} or \sphinxcode{\sphinxupquote{\sphinxhyphen{}suffix kk}} {\hyperref[\detokenize{Run_options::doc}]{\sphinxcrossref{\DUrole{doc}{command\sphinxhyphen{}line switches}}}}.
Also see the {\hyperref[\detokenize{Packages_details:pkg-gpu}]{\sphinxcrossref{\DUrole{std,std-ref}{GPU}}}}, {\hyperref[\detokenize{Packages_details:pkg-opt}]{\sphinxcrossref{\DUrole{std,std-ref}{OPT}}}}, {\hyperref[\detokenize{Packages_details:pkg-intel}]{\sphinxcrossref{\DUrole{std,std-ref}{INTEL}}}}, and {\hyperref[\detokenize{Packages_details:pkg-openmp}]{\sphinxcrossref{\DUrole{std,std-ref}{OPENMP}}}} packages, which have styles
optimized for CPUs, KNLs, and GPUs.

\sphinxAtStartPar
You must have a C++17 compatible compiler to use this package.
KOKKOS makes extensive use of advanced C++ features, which can
expose compiler bugs, especially when compiling for maximum
performance at high optimization levels. Please see the file
\sphinxcode{\sphinxupquote{lib/kokkos/README}} for a list of compilers and their respective
platforms, that are known to work.

\sphinxAtStartPar
\sphinxstylestrong{Authors:} The KOKKOS package was created primarily by Christian Trott
and Stan Moore (Sandia), with contributions from other folks as well.
It uses the open\sphinxhyphen{}source \sphinxhref{https://github.com/kokkos}{Kokkos library}
which was developed by Carter Edwards, Christian Trott, and others at
Sandia, and which is included in the LAMMPS distribution in
\sphinxcode{\sphinxupquote{lib/kokkos}}.

\sphinxAtStartPar
\sphinxstylestrong{Install:}

\sphinxAtStartPar
This package has {\hyperref[\detokenize{Build_extras:kokkos}]{\sphinxcrossref{\DUrole{std,std-ref}{specific installation instructions}}}} on the {\hyperref[\detokenize{Build_extras::doc}]{\sphinxcrossref{\DUrole{doc}{Build extras}}}} page.

\sphinxAtStartPar
\sphinxstylestrong{Supporting info:}
\begin{itemize}
\item {} 
\sphinxAtStartPar
\sphinxcode{\sphinxupquote{src/KOKKOS}}: filenames \sphinxhyphen{}\textgreater{} commands

\item {} 
\sphinxAtStartPar
\sphinxcode{\sphinxupquote{src/KOKKOS/README}}

\item {} 
\sphinxAtStartPar
\sphinxcode{\sphinxupquote{lib/kokkos/README}}

\item {} 
\sphinxAtStartPar
{\hyperref[\detokenize{Speed_packages::doc}]{\sphinxcrossref{\DUrole{doc}{Accelerator packages}}}}

\item {} 
\sphinxAtStartPar
{\hyperref[\detokenize{Speed_kokkos::doc}]{\sphinxcrossref{\DUrole{doc}{KOKKOS package}}}}

\item {} 
\sphinxAtStartPar
{\hyperref[\detokenize{Run_options::doc}]{\sphinxcrossref{\DUrole{doc}{Section 2.6 \sphinxhyphen{}k on …}}}}

\item {} 
\sphinxAtStartPar
{\hyperref[\detokenize{Run_options::doc}]{\sphinxcrossref{\DUrole{doc}{Section 2.6 \sphinxhyphen{}sf kk}}}}

\item {} 
\sphinxAtStartPar
{\hyperref[\detokenize{Run_options::doc}]{\sphinxcrossref{\DUrole{doc}{Section 2.6 \sphinxhyphen{}pk kokkos}}}}

\item {} 
\sphinxAtStartPar
{\hyperref[\detokenize{package::doc}]{\sphinxcrossref{\DUrole{doc}{package kokkos}}}}

\item {} 
\sphinxAtStartPar
Search the {\hyperref[\detokenize{Commands_all::doc}]{\sphinxcrossref{\DUrole{doc}{commands}}}} pages ({\hyperref[\detokenize{Commands_fix::doc}]{\sphinxcrossref{\DUrole{doc}{fix}}}}, {\hyperref[\detokenize{Commands_compute::doc}]{\sphinxcrossref{\DUrole{doc}{compute}}}},
{\hyperref[\detokenize{Commands_pair::doc}]{\sphinxcrossref{\DUrole{doc}{pair}}}}, {\hyperref[\detokenize{Commands_bond::doc}]{\sphinxcrossref{\DUrole{doc}{bond, angle, dihedral, improper}}}},
{\hyperref[\detokenize{Commands_kspace::doc}]{\sphinxcrossref{\DUrole{doc}{kspace}}}}) for styles followed by (k)

\item {} 
\sphinxAtStartPar
\sphinxhref{https://www.lammps.org/bench.html}{Benchmarks page} of website

\end{itemize}


\bigskip\hrule\bigskip



\subsection{KSPACE package}
\label{\detokenize{Packages_details:kspace-package}}\label{\detokenize{Packages_details:pkg-kspace}}
\sphinxAtStartPar
\sphinxstylestrong{Contents:}

\sphinxAtStartPar
A variety of long\sphinxhyphen{}range Coulombic solvers, as well as pair styles which
compute the corresponding short\sphinxhyphen{}range pairwise Coulombic interactions.
These include Ewald, particle\sphinxhyphen{}particle particle\sphinxhyphen{}mesh (PPPM), and
multilevel summation method (MSM) solvers.

\sphinxAtStartPar
\sphinxstylestrong{Install:}

\sphinxAtStartPar
Building with this package requires a 1d FFT library be present on your
system for use by the PPPM solvers.  This can be the KISS FFT library
provided with LAMMPS, third party libraries like FFTW, or a
vendor\sphinxhyphen{}supplied FFT library.  See the {\hyperref[\detokenize{Build_settings::doc}]{\sphinxcrossref{\DUrole{doc}{Build settings}}}} page for details on how to select different FFT
options for your LAMMPS build.

\sphinxAtStartPar
\sphinxstylestrong{Supporting info:}
\begin{itemize}
\item {} 
\sphinxAtStartPar
\sphinxcode{\sphinxupquote{src/KSPACE}}: filenames \sphinxhyphen{}\textgreater{} commands

\item {} 
\sphinxAtStartPar
{\hyperref[\detokenize{kspace_style::doc}]{\sphinxcrossref{\DUrole{doc}{kspace\_style}}}}

\item {} 
\sphinxAtStartPar
\sphinxhref{PDF/kspace.pdf}{doc/PDF/kspace.pdf}

\item {} 
\sphinxAtStartPar
{\hyperref[\detokenize{Howto_tip3p::doc}]{\sphinxcrossref{\DUrole{doc}{Howto tip3p}}}}

\item {} 
\sphinxAtStartPar
{\hyperref[\detokenize{Howto_tip4p::doc}]{\sphinxcrossref{\DUrole{doc}{Howto tip4p}}}}

\item {} 
\sphinxAtStartPar
{\hyperref[\detokenize{Howto_spc::doc}]{\sphinxcrossref{\DUrole{doc}{Howto spc}}}}

\item {} 
\sphinxAtStartPar
{\hyperref[\detokenize{pair_coul::doc}]{\sphinxcrossref{\DUrole{doc}{pair\_style coul}}}}

\item {} 
\sphinxAtStartPar
Search the {\hyperref[\detokenize{Commands_pair::doc}]{\sphinxcrossref{\DUrole{doc}{pair style}}}} page for styles with “long” or “msm” in name

\item {} 
\sphinxAtStartPar
\sphinxcode{\sphinxupquote{examples/peptide}}

\item {} 
\sphinxAtStartPar
bench/in.rhodo

\end{itemize}


\bigskip\hrule\bigskip



\subsection{LATBOLTZ package}
\label{\detokenize{Packages_details:latboltz-package}}\label{\detokenize{Packages_details:pkg-latboltz}}
\sphinxAtStartPar
\sphinxstylestrong{Contents:}

\sphinxAtStartPar
Fixes which implement a background Lattice\sphinxhyphen{}Boltzmann (LB) fluid, which
can be used to model MD particles influenced by hydrodynamic forces.

\sphinxAtStartPar
\sphinxstylestrong{Authors:} Frances Mackay and Colin Denniston (University of Western
Ontario).

\sphinxAtStartPar
\sphinxstylestrong{Install:}

\sphinxAtStartPar
The LATBOLTZ package requires that LAMMPS is build in {\hyperref[\detokenize{Build_basics:serial}]{\sphinxcrossref{\DUrole{std,std-ref}{MPI parallel mode}}}}.

\sphinxAtStartPar
\sphinxstylestrong{Supporting info:}
\begin{itemize}
\item {} 
\sphinxAtStartPar
\sphinxcode{\sphinxupquote{src/LATBOLTZ}}: filenames \sphinxhyphen{}\textgreater{} commands

\item {} 
\sphinxAtStartPar
\sphinxcode{\sphinxupquote{src/LATBOLTZ/README}}

\item {} 
\sphinxAtStartPar
{\hyperref[\detokenize{fix_lb_fluid::doc}]{\sphinxcrossref{\DUrole{doc}{fix lb/fluid}}}}

\item {} 
\sphinxAtStartPar
{\hyperref[\detokenize{fix_lb_momentum::doc}]{\sphinxcrossref{\DUrole{doc}{fix lb/momentum}}}}

\item {} 
\sphinxAtStartPar
{\hyperref[\detokenize{fix_lb_viscous::doc}]{\sphinxcrossref{\DUrole{doc}{fix lb/viscous}}}}

\item {} 
\sphinxAtStartPar
\sphinxcode{\sphinxupquote{examples/PACKAGES/latboltz}}

\end{itemize}


\bigskip\hrule\bigskip



\subsection{LEPTON package}
\label{\detokenize{Packages_details:lepton-package}}\label{\detokenize{Packages_details:pkg-lepton}}
\sphinxAtStartPar
\sphinxstylestrong{Contents:}

\sphinxAtStartPar
Styles for pair, bond, and angle forces that evaluate the potential
function from a string using the \sphinxhref{https://simtk.org/projects/lepton}{Lepton mathematical expression parser}.  Lepton is a C++ library that is
bundled with \sphinxhref{https://openmm.org/}{OpenMM} and can be used for
parsing, evaluating, differentiating, and analyzing mathematical
expressions.  This is a more lightweight and efficient alternative for
evaluating custom potential function to an embedded Python interpreter
as used in the {\hyperref[\detokenize{Packages_details:pkg-python}]{\sphinxcrossref{\DUrole{std,std-ref}{PYTHON package}}}}.  On the other hand,
since the potentials are evaluated form analytical expressions, they are
more precise than what can be done with {\hyperref[\detokenize{Tools:tabulate}]{\sphinxcrossref{\DUrole{std,std-ref}{tabulated potentials}}}}.

\sphinxAtStartPar
\sphinxstylestrong{Authors:} Axel Kohlmeyer (Temple U).  Lepton itself is developed
by Peter Eastman at Stanford University.

\sphinxAtStartPar
\DUrole{versionmodified,added}{Added in version 8Feb2023.}

\sphinxAtStartPar
\sphinxstylestrong{Install:}

\sphinxAtStartPar
This package has {\hyperref[\detokenize{Build_extras:lepton}]{\sphinxcrossref{\DUrole{std,std-ref}{specific installation instructions}}}} on
the {\hyperref[\detokenize{Build_extras::doc}]{\sphinxcrossref{\DUrole{doc}{Build extras}}}} page.

\sphinxAtStartPar
\sphinxstylestrong{Supporting info:}
\begin{itemize}
\item {} 
\sphinxAtStartPar
\sphinxcode{\sphinxupquote{src/LEPTON}}: filenames \sphinxhyphen{}\textgreater{} commands

\item {} 
\sphinxAtStartPar
\sphinxcode{\sphinxupquote{lib/lepton/README}}.md

\item {} 
\sphinxAtStartPar
{\hyperref[\detokenize{pair_lepton::doc}]{\sphinxcrossref{\DUrole{doc}{pair\_style lepton}}}}

\item {} 
\sphinxAtStartPar
{\hyperref[\detokenize{bond_lepton::doc}]{\sphinxcrossref{\DUrole{doc}{bond\_style lepton}}}}

\item {} 
\sphinxAtStartPar
{\hyperref[\detokenize{angle_lepton::doc}]{\sphinxcrossref{\DUrole{doc}{angle\_style lepton}}}}

\item {} 
\sphinxAtStartPar
{\hyperref[\detokenize{dihedral_lepton::doc}]{\sphinxcrossref{\DUrole{doc}{dihedral\_style lepton}}}}

\end{itemize}


\bigskip\hrule\bigskip



\subsection{MACHDYN package}
\label{\detokenize{Packages_details:machdyn-package}}\label{\detokenize{Packages_details:pkg-machdyn}}
\sphinxAtStartPar
\sphinxstylestrong{Contents:}

\sphinxAtStartPar
An atom style, fixes, computes, and several pair styles which
implements smoothed Mach dynamics (SMD) for solids, which is a model
related to smoothed particle hydrodynamics (SPH) for liquids (see the
{\hyperref[\detokenize{Packages_details:pkg-sph}]{\sphinxcrossref{\DUrole{std,std-ref}{SPH package}}}}).

\sphinxAtStartPar
This package solves solids mechanics problems via a state of the art
stabilized meshless method with hourglass control.  It can specify
hydrostatic interactions independently from material strength models,
i.e. pressure and deviatoric stresses are separated.  It provides many
material models (Johnson\sphinxhyphen{}Cook, plasticity with hardening,
Mie\sphinxhyphen{}Grueneisen, Polynomial EOS) and allows new material models to be
added.  It implements rigid boundary conditions (walls) which can be
specified as surface geometries from *.STL files.

\sphinxAtStartPar
\sphinxstylestrong{Author:} Georg Ganzenmuller (Fraunhofer\sphinxhyphen{}Institute for High\sphinxhyphen{}Speed
Dynamics, Ernst Mach Institute, Germany).

\sphinxAtStartPar
\sphinxstylestrong{Install:}

\sphinxAtStartPar
This package has {\hyperref[\detokenize{Build_extras:machdyn}]{\sphinxcrossref{\DUrole{std,std-ref}{specific installation instructions}}}} on the {\hyperref[\detokenize{Build_extras::doc}]{\sphinxcrossref{\DUrole{doc}{Build extras}}}} page.

\sphinxAtStartPar
\sphinxstylestrong{Supporting info:}
\begin{itemize}
\item {} 
\sphinxAtStartPar
\sphinxcode{\sphinxupquote{src/MACHDYN}}: filenames \sphinxhyphen{}\textgreater{} commands

\item {} 
\sphinxAtStartPar
\sphinxcode{\sphinxupquote{src/MACHDYN/README}}

\item {} 
\sphinxAtStartPar
\sphinxhref{PDF/MACHDYN\_LAMMPS\_userguide.pdf}{doc/PDF/MACHDYN\_LAMMPS\_userguide.pdf}

\item {} 
\sphinxAtStartPar
\sphinxcode{\sphinxupquote{examples/PACKAGES/machdyn}}

\item {} 
\sphinxAtStartPar
\sphinxurl{https://www.lammps.org/movies.html\#smd}

\end{itemize}


\bigskip\hrule\bigskip



\subsection{MANIFOLD package}
\label{\detokenize{Packages_details:manifold-package}}\label{\detokenize{Packages_details:pkg-manifold}}
\sphinxAtStartPar
\sphinxstylestrong{Contents:}

\sphinxAtStartPar
Several fixes and a “manifold” class which enable simulations of
particles constrained to a manifold (a 2D surface within the 3D
simulation box).  This is done by applying the RATTLE constraint
algorithm to formulate single\sphinxhyphen{}particle constraint functions
g(xi,yi,zi) = 0 and their derivative (i.e. the normal of the manifold)
n = grad(g).

\sphinxAtStartPar
\sphinxstylestrong{Author:} Stefan Paquay (until 2017: Eindhoven University of
Technology (TU/e), The Netherlands; since 2017: Brandeis University,
Waltham, MA, USA)

\sphinxAtStartPar
\sphinxstylestrong{Supporting info:}
\begin{itemize}
\item {} 
\sphinxAtStartPar
\sphinxcode{\sphinxupquote{src/MANIFOLD}}: filenames \sphinxhyphen{}\textgreater{} commands

\item {} 
\sphinxAtStartPar
\sphinxcode{\sphinxupquote{src/MANIFOLD/README}}

\item {} 
\sphinxAtStartPar
{\hyperref[\detokenize{Howto_manifold::doc}]{\sphinxcrossref{\DUrole{doc}{Howto manifold}}}}

\item {} 
\sphinxAtStartPar
{\hyperref[\detokenize{fix_manifoldforce::doc}]{\sphinxcrossref{\DUrole{doc}{fix manifoldforce}}}}

\item {} 
\sphinxAtStartPar
{\hyperref[\detokenize{fix_nve_manifold_rattle::doc}]{\sphinxcrossref{\DUrole{doc}{fix nve/manifold/rattle}}}}

\item {} 
\sphinxAtStartPar
{\hyperref[\detokenize{fix_nvt_manifold_rattle::doc}]{\sphinxcrossref{\DUrole{doc}{fix nvt/manifold/rattle}}}}

\item {} 
\sphinxAtStartPar
\sphinxcode{\sphinxupquote{examples/PACKAGES/manifold}}

\item {} 
\sphinxAtStartPar
\sphinxurl{https://www.lammps.org/movies.html\#manifold}

\end{itemize}


\bigskip\hrule\bigskip



\subsection{MANYBODY package}
\label{\detokenize{Packages_details:manybody-package}}\label{\detokenize{Packages_details:pkg-manybody}}
\sphinxAtStartPar
\sphinxstylestrong{Contents:}

\sphinxAtStartPar
A variety of many\sphinxhyphen{}body and bond\sphinxhyphen{}order potentials.  These include
(AI)REBO, BOP, EAM, EIM, Stillinger\sphinxhyphen{}Weber, and Tersoff potentials.

\sphinxAtStartPar
\sphinxstylestrong{Supporting info:}
\begin{itemize}
\item {} 
\sphinxAtStartPar
\sphinxcode{\sphinxupquote{src/MANYBODY}}: filenames \sphinxhyphen{}\textgreater{} commands

\item {} 
\sphinxAtStartPar
{\hyperref[\detokenize{Commands_pair::doc}]{\sphinxcrossref{\DUrole{doc}{Pair style}}}} page

\item {} 
\sphinxAtStartPar
\sphinxcode{\sphinxupquote{examples/comb}}

\item {} 
\sphinxAtStartPar
\sphinxcode{\sphinxupquote{examples/eim}}

\item {} 
\sphinxAtStartPar
\sphinxcode{\sphinxupquote{examples/nb3d}}

\item {} 
\sphinxAtStartPar
\sphinxcode{\sphinxupquote{examples/shear}}

\item {} 
\sphinxAtStartPar
\sphinxcode{\sphinxupquote{examples/streitz}}

\item {} 
\sphinxAtStartPar
\sphinxcode{\sphinxupquote{examples/vashishta}}

\item {} 
\sphinxAtStartPar
bench/in.eam

\end{itemize}


\bigskip\hrule\bigskip



\subsection{MC package}
\label{\detokenize{Packages_details:mc-package}}\label{\detokenize{Packages_details:pkg-mc}}
\sphinxAtStartPar
\sphinxstylestrong{Contents:}

\sphinxAtStartPar
Several fixes and a pair style that have Monte Carlo (MC) or MC\sphinxhyphen{}like
attributes.  These include fixes for creating, breaking, and swapping
bonds, for performing atomic swaps, and performing grand canonical
MC (GCMC), semi\sphinxhyphen{}grand canonical MC (SGCMC), or similar processes in
conjunction with molecular dynamics (MD).

\sphinxAtStartPar
\sphinxstylestrong{Supporting info:}
\begin{itemize}
\item {} 
\sphinxAtStartPar
\sphinxcode{\sphinxupquote{src/MC}}: filenames \sphinxhyphen{}\textgreater{} commands

\item {} 
\sphinxAtStartPar
{\hyperref[\detokenize{fix_atom_swap::doc}]{\sphinxcrossref{\DUrole{doc}{fix atom/swap}}}}

\item {} 
\sphinxAtStartPar
{\hyperref[\detokenize{fix_bond_break::doc}]{\sphinxcrossref{\DUrole{doc}{fix bond/break}}}}

\item {} 
\sphinxAtStartPar
{\hyperref[\detokenize{fix_bond_create::doc}]{\sphinxcrossref{\DUrole{doc}{fix bond/create}}}}

\item {} 
\sphinxAtStartPar
{\hyperref[\detokenize{fix_bond_create::doc}]{\sphinxcrossref{\DUrole{doc}{fix bond/create/angle}}}}

\item {} 
\sphinxAtStartPar
{\hyperref[\detokenize{fix_bond_swap::doc}]{\sphinxcrossref{\DUrole{doc}{fix bond/swap}}}}

\item {} 
\sphinxAtStartPar
{\hyperref[\detokenize{fix_charge_regulation::doc}]{\sphinxcrossref{\DUrole{doc}{fix charge/regulation}}}}

\item {} 
\sphinxAtStartPar
{\hyperref[\detokenize{fix_gcmc::doc}]{\sphinxcrossref{\DUrole{doc}{fix gcmc}}}}

\item {} 
\sphinxAtStartPar
{\hyperref[\detokenize{fix_sgcmc::doc}]{\sphinxcrossref{\DUrole{doc}{fix sgcmc}}}}

\item {} 
\sphinxAtStartPar
{\hyperref[\detokenize{fix_tfmc::doc}]{\sphinxcrossref{\DUrole{doc}{fix tfmc}}}}

\item {} 
\sphinxAtStartPar
{\hyperref[\detokenize{fix_widom::doc}]{\sphinxcrossref{\DUrole{doc}{fix widom}}}}

\item {} 
\sphinxAtStartPar
{\hyperref[\detokenize{pair_dsmc::doc}]{\sphinxcrossref{\DUrole{doc}{pair\_style dsmc}}}}

\item {} 
\sphinxAtStartPar
\sphinxurl{https://www.lammps.org/movies.html\#gcmc}

\end{itemize}


\bigskip\hrule\bigskip



\subsection{MDI package}
\label{\detokenize{Packages_details:mdi-package}}\label{\detokenize{Packages_details:pkg-mdi}}
\sphinxAtStartPar
\sphinxstylestrong{Contents:}

\sphinxAtStartPar
A LAMMPS command and fixes to allow client\sphinxhyphen{}server coupling of LAMMPS
to other atomic or molecular simulation codes or materials modeling
workflows via the \sphinxhref{https://molssi-mdi.github.io/MDI\_Library/html/index.html}{MolSSI Driver Interface
(MDI) library}.

\sphinxAtStartPar
\sphinxstylestrong{Author:} Taylor Barnes \sphinxhyphen{} MolSSI, taylor.a.barnes at gmail.com

\sphinxAtStartPar
\DUrole{versionmodified,added}{Added in version 14May2021.}

\sphinxAtStartPar
\sphinxstylestrong{Install:}

\sphinxAtStartPar
This package has {\hyperref[\detokenize{Build_extras:mdi}]{\sphinxcrossref{\DUrole{std,std-ref}{specific installation instructions}}}} on
the {\hyperref[\detokenize{Build_extras::doc}]{\sphinxcrossref{\DUrole{doc}{Build extras}}}} page.

\sphinxAtStartPar
\sphinxstylestrong{Supporting info:}
\begin{itemize}
\item {} 
\sphinxAtStartPar
\sphinxcode{\sphinxupquote{src/MDI/README}}

\item {} 
\sphinxAtStartPar
\sphinxcode{\sphinxupquote{lib/mdi/README}}

\item {} 
\sphinxAtStartPar
{\hyperref[\detokenize{Howto_mdi::doc}]{\sphinxcrossref{\DUrole{doc}{Howto MDI}}}}

\item {} 
\sphinxAtStartPar
{\hyperref[\detokenize{mdi::doc}]{\sphinxcrossref{\DUrole{doc}{mdi}}}}

\item {} 
\sphinxAtStartPar
{\hyperref[\detokenize{fix_mdi_qm::doc}]{\sphinxcrossref{\DUrole{doc}{fix mdi/qm}}}}

\item {} 
\sphinxAtStartPar
\sphinxcode{\sphinxupquote{examples/PACKAGES/mdi}}

\end{itemize}


\bigskip\hrule\bigskip



\subsection{MEAM package}
\label{\detokenize{Packages_details:meam-package}}\label{\detokenize{Packages_details:pkg-meam}}
\sphinxAtStartPar
\sphinxstylestrong{Contents:}

\sphinxAtStartPar
A pair style for the modified embedded atom (MEAM) potential
translated from the Fortran version in the (obsolete) MEAM package
to plain C++. The MEAM fully replaces the MEAM package, which
has been removed from LAMMPS after the 12 December 2018 version.

\sphinxAtStartPar
\sphinxstylestrong{Author:} Sebastian Huetter, (Otto\sphinxhyphen{}von\sphinxhyphen{}Guericke University Magdeburg)
based on the Fortran version of Greg Wagner (Northwestern U) while at
Sandia.

\sphinxAtStartPar
\sphinxstylestrong{Supporting info:}
\begin{itemize}
\item {} 
\sphinxAtStartPar
\sphinxcode{\sphinxupquote{src/MEAM}}: filenames \sphinxhyphen{}\textgreater{} commands

\item {} 
\sphinxAtStartPar
\sphinxcode{\sphinxupquote{src/MEAM/README}}

\item {} 
\sphinxAtStartPar
{\hyperref[\detokenize{pair_meam::doc}]{\sphinxcrossref{\DUrole{doc}{pair\_style meam}}}}

\item {} 
\sphinxAtStartPar
\sphinxcode{\sphinxupquote{examples/meam}}

\end{itemize}


\bigskip\hrule\bigskip



\subsection{MESONT package}
\label{\detokenize{Packages_details:mesont-package}}\label{\detokenize{Packages_details:pkg-mesont}}
\sphinxAtStartPar
\sphinxstylestrong{Contents:}

\sphinxAtStartPar
MESONT is a LAMMPS package for simulation of nanomechanics of nanotubes
(NTs).  The model is based on a coarse\sphinxhyphen{}grained representation of NTs as
“flexible cylinders” consisting of a variable number of
segments. Internal interactions within a NT and the van der Waals
interaction between the tubes are described by a mesoscopic force field
designed and parameterized based on the results of atomic\sphinxhyphen{}level
molecular dynamics simulations. The description of the force field is
provided in the papers listed in \sphinxcode{\sphinxupquote{src/MESONT/README}}.

\sphinxAtStartPar
This package used to have two independent implementations of this model:
the original implementation using a Fortran library written by the
developers of the model and a second implementation written in C++ by
Philipp Kloza (U Cambridge).  Since the C++ implementation offers the
same features as the original implementation with the addition of
friction, is typically faster, and easier to compile/install, the
Fortran library based implementation has since been obsoleted and
removed from the distribution. You have to download and compile
an older version of LAMMPS if you want to use those.

\sphinxAtStartPar
\sphinxstylestrong{Download of potential files:}

\sphinxAtStartPar
The potential files for these pair styles are \sphinxstyleemphasis{very} large and thus are
not included in the regular downloaded packages of LAMMPS or the git
repositories.  Instead, they will be automatically downloaded from a web
server when the package is installed for the first time.

\sphinxAtStartPar
\sphinxstylestrong{Authors of the obsoleted *mesont* styles:}

\sphinxAtStartPar
Maxim V. Shugaev (University of Virginia), Alexey N. Volkov (University
of Alabama), Leonid V. Zhigilei (University of Virginia)

\sphinxAtStartPar
\DUrole{versionmodified,deprecated}{Deprecated since version 8Feb2023.}

\sphinxAtStartPar
\sphinxstylestrong{Author of the C++ styles:}
Philipp Kloza (U Cambridge)

\sphinxAtStartPar
\DUrole{versionmodified,added}{Added in version 15Jun2020.}

\sphinxAtStartPar
\sphinxstylestrong{Supporting info:}
\begin{itemize}
\item {} 
\sphinxAtStartPar
\sphinxcode{\sphinxupquote{src/MESONT}}: filenames \sphinxhyphen{}\textgreater{} commands

\item {} 
\sphinxAtStartPar
\sphinxcode{\sphinxupquote{src/MESONT/README}}

\item {} 
\sphinxAtStartPar
{\hyperref[\detokenize{bond_mesocnt::doc}]{\sphinxcrossref{\DUrole{doc}{bond\_style mesocnt}}}}

\item {} 
\sphinxAtStartPar
{\hyperref[\detokenize{angle_mesocnt::doc}]{\sphinxcrossref{\DUrole{doc}{angle\_style mesocnt}}}}

\item {} 
\sphinxAtStartPar
{\hyperref[\detokenize{pair_mesocnt::doc}]{\sphinxcrossref{\DUrole{doc}{pair\_style mesocnt}}}}

\item {} 
\sphinxAtStartPar
\sphinxcode{\sphinxupquote{examples/PACKAGES/mesont}}

\end{itemize}


\bigskip\hrule\bigskip



\subsection{MGPT package}
\label{\detokenize{Packages_details:mgpt-package}}\label{\detokenize{Packages_details:pkg-mgpt}}
\sphinxAtStartPar
\sphinxstylestrong{Contents:}

\sphinxAtStartPar
A pair style which provides a fast implementation of the quantum\sphinxhyphen{}based
MGPT multi\sphinxhyphen{}ion potentials.  The MGPT or model GPT method derives from
first\sphinxhyphen{}principles DFT\sphinxhyphen{}based generalized pseudopotential theory (GPT)
through a series of systematic approximations valid for mid\sphinxhyphen{}period
transition metals with nearly half\sphinxhyphen{}filled d bands.  The MGPT method
was originally developed by John Moriarty at LLNL.  The pair style in
this package calculates forces and energies using an optimized
matrix\sphinxhyphen{}MGPT algorithm due to Tomas Oppelstrup at LLNL.

\sphinxAtStartPar
\sphinxstylestrong{Authors:} Tomas Oppelstrup and John Moriarty (LLNL).

\sphinxAtStartPar
\sphinxstylestrong{Supporting info:}
\begin{itemize}
\item {} 
\sphinxAtStartPar
\sphinxcode{\sphinxupquote{src/MGPT}}: filenames \sphinxhyphen{}\textgreater{} commands

\item {} 
\sphinxAtStartPar
\sphinxcode{\sphinxupquote{src/MGPT/README}}

\item {} 
\sphinxAtStartPar
{\hyperref[\detokenize{pair_mgpt::doc}]{\sphinxcrossref{\DUrole{doc}{pair\_style mgpt}}}}

\item {} 
\sphinxAtStartPar
\sphinxcode{\sphinxupquote{examples/PACKAGES/mgpt}}

\end{itemize}


\bigskip\hrule\bigskip



\subsection{MISC package}
\label{\detokenize{Packages_details:misc-package}}\label{\detokenize{Packages_details:pkg-misc}}
\sphinxAtStartPar
\sphinxstylestrong{Contents:}

\sphinxAtStartPar
A variety of compute, fix, pair, bond styles with specialized
capabilities that don’t align with other packages.  Do a directory
listing, \sphinxcode{\sphinxupquote{ls src/MISC}}, to see the list of commands.

\begin{sphinxadmonition}{note}{Note:}
\sphinxAtStartPar
the MISC package contains styles that require using the
\sphinxhyphen{}restrict flag, when compiling with Intel compilers.
\end{sphinxadmonition}

\sphinxAtStartPar
\sphinxstylestrong{Supporting info:}
\begin{itemize}
\item {} 
\sphinxAtStartPar
\sphinxcode{\sphinxupquote{src/MISC}}: filenames \sphinxhyphen{}\textgreater{} commands

\item {} 
\sphinxAtStartPar
{\hyperref[\detokenize{bond_special::doc}]{\sphinxcrossref{\DUrole{doc}{bond\_style special}}}}

\item {} 
\sphinxAtStartPar
{\hyperref[\detokenize{compute_viscosity_cos::doc}]{\sphinxcrossref{\DUrole{doc}{compute viscosity/cos}}}}

\item {} 
\sphinxAtStartPar
{\hyperref[\detokenize{fix_accelerate_cos::doc}]{\sphinxcrossref{\DUrole{doc}{fix accelerate/cos}}}}

\item {} 
\sphinxAtStartPar
{\hyperref[\detokenize{fix_imd::doc}]{\sphinxcrossref{\DUrole{doc}{fix imd}}}}

\item {} 
\sphinxAtStartPar
{\hyperref[\detokenize{fix_ipi::doc}]{\sphinxcrossref{\DUrole{doc}{fix ipi}}}}

\item {} 
\sphinxAtStartPar
{\hyperref[\detokenize{pair_agni::doc}]{\sphinxcrossref{\DUrole{doc}{pair\_style agni}}}}

\item {} 
\sphinxAtStartPar
{\hyperref[\detokenize{pair_list::doc}]{\sphinxcrossref{\DUrole{doc}{pair\_style list}}}}

\item {} 
\sphinxAtStartPar
{\hyperref[\detokenize{pair_srp::doc}]{\sphinxcrossref{\DUrole{doc}{pair\_style srp}}}}

\item {} 
\sphinxAtStartPar
{\hyperref[\detokenize{pair_tracker::doc}]{\sphinxcrossref{\DUrole{doc}{pair\_style tracker}}}}

\end{itemize}


\bigskip\hrule\bigskip



\subsection{ML\sphinxhyphen{}HDNNP package}
\label{\detokenize{Packages_details:ml-hdnnp-package}}\label{\detokenize{Packages_details:pkg-ml-hdnnp}}
\sphinxAtStartPar
\sphinxstylestrong{Contents:}

\sphinxAtStartPar
A {\hyperref[\detokenize{pair_hdnnp::doc}]{\sphinxcrossref{\DUrole{doc}{pair\_style hdnnp}}}} command which allows to use
high\sphinxhyphen{}dimensional neural network potentials (HDNNPs), a form of machine learning
potentials. HDNNPs must be carefully trained prior to their application in a
molecular dynamics simulation.

\sphinxAtStartPar
To use this package you must have the \sphinxhref{https://github.com/CompPhysVienna/n2p2}{n2p2} library installed and
compiled on your system.

\sphinxAtStartPar
\sphinxstylestrong{Author:} Andreas Singraber

\sphinxAtStartPar
\DUrole{versionmodified,added}{Added in version 27May2021.}

\sphinxAtStartPar
\sphinxstylestrong{Install:}

\sphinxAtStartPar
This package has {\hyperref[\detokenize{Build_extras:ml-hdnnp}]{\sphinxcrossref{\DUrole{std,std-ref}{specific installation instructions}}}} on the
{\hyperref[\detokenize{Build_extras::doc}]{\sphinxcrossref{\DUrole{doc}{Build extras}}}} page.

\sphinxAtStartPar
\sphinxstylestrong{Supporting info:}
\begin{itemize}
\item {} 
\sphinxAtStartPar
\sphinxcode{\sphinxupquote{src/ML\sphinxhyphen{}HDNNP}}: filenames \sphinxhyphen{}\textgreater{} commands

\item {} 
\sphinxAtStartPar
\sphinxcode{\sphinxupquote{src/ML\sphinxhyphen{}HDNNP/README}}

\item {} 
\sphinxAtStartPar
\sphinxcode{\sphinxupquote{lib/hdnnp/README}}

\item {} 
\sphinxAtStartPar
{\hyperref[\detokenize{pair_hdnnp::doc}]{\sphinxcrossref{\DUrole{doc}{pair\_style hdnnp}}}}

\item {} 
\sphinxAtStartPar
\sphinxcode{\sphinxupquote{examples/PACKAGES/hdnnp}}

\end{itemize}


\bigskip\hrule\bigskip



\subsection{ML\sphinxhyphen{}IAP package}
\label{\detokenize{Packages_details:ml-iap-package}}\label{\detokenize{Packages_details:pkg-ml-iap}}
\sphinxAtStartPar
\sphinxstylestrong{Contents:}

\sphinxAtStartPar
A general interface for machine\sphinxhyphen{}learning interatomic potentials, including PyTorch.

\sphinxAtStartPar
\sphinxstylestrong{Install:}

\sphinxAtStartPar
To use this package, also the {\hyperref[\detokenize{Packages_details:pkg-ml-snap}]{\sphinxcrossref{\DUrole{std,std-ref}{ML\sphinxhyphen{}SNAP}}}} package needs
to be installed.  To make the \sphinxstyleemphasis{mliappy} model available, also the
{\hyperref[\detokenize{Packages_details:pkg-python}]{\sphinxcrossref{\DUrole{std,std-ref}{PYTHON}}}} package needs to be installed, the version
of Python must be 3.6 or later, and the \sphinxhref{https://cython.org/}{cython} software
must be installed.

\sphinxAtStartPar
\sphinxstylestrong{Author:} Aidan Thompson (Sandia), Nicholas Lubbers (LANL).

\sphinxAtStartPar
\DUrole{versionmodified,added}{Added in version 30Jun2020.}

\sphinxAtStartPar
\sphinxstylestrong{Supporting info:}
\begin{itemize}
\item {} 
\sphinxAtStartPar
\sphinxcode{\sphinxupquote{src/ML\sphinxhyphen{}IAP}}: filenames \sphinxhyphen{}\textgreater{} commands

\item {} 
\sphinxAtStartPar
\sphinxcode{\sphinxupquote{src/ML\sphinxhyphen{}IAP/README.md}}

\item {} 
\sphinxAtStartPar
{\hyperref[\detokenize{pair_mliap::doc}]{\sphinxcrossref{\DUrole{doc}{pair\_style mliap}}}}

\item {} 
\sphinxAtStartPar
{\hyperref[\detokenize{compute_mliap::doc}]{\sphinxcrossref{\DUrole{doc}{compute\_style mliap}}}}

\item {} 
\sphinxAtStartPar
\sphinxcode{\sphinxupquote{examples/mliap}} (see README)

\end{itemize}

\sphinxAtStartPar
When built with the \sphinxstyleemphasis{mliappy} model this package includes an extension for
coupling with Python models, including PyTorch. In this case, the Python
interpreter linked to LAMMPS will need the \sphinxcode{\sphinxupquote{cython}} and \sphinxcode{\sphinxupquote{numpy}} modules
installed.  The provided examples build models with PyTorch, which would
therefore also needs to be installed to run those examples.


\bigskip\hrule\bigskip



\subsection{ML\sphinxhyphen{}PACE package}
\label{\detokenize{Packages_details:ml-pace-package}}\label{\detokenize{Packages_details:pkg-ml-pace}}
\sphinxAtStartPar
\sphinxstylestrong{Contents:}

\sphinxAtStartPar
A pair style for the Atomic Cluster Expansion potential (ACE).
ACE is a methodology for deriving a highly accurate classical potential
fit to a large archive of quantum mechanical (DFT) data. The ML\sphinxhyphen{}PACE
package provides an efficient implementation for running simulations
with ACE potentials.

\sphinxAtStartPar
\sphinxstylestrong{Authors:}

\sphinxAtStartPar
This package was written by Yury Lysogorskiy\textasciicircum{}1,
Cas van der Oord\textasciicircum{}2, Anton Bochkarev\textasciicircum{}1,
Sarath Menon\textasciicircum{}1, Matteo Rinaldi\textasciicircum{}1, Thomas Hammerschmidt\textasciicircum{}1, Matous Mrovec\textasciicircum{}1,
Aidan Thompson\textasciicircum{}3, Gabor Csanyi\textasciicircum{}2, Christoph Ortner\textasciicircum{}4, Ralf Drautz\textasciicircum{}1.
\begin{quote}

\sphinxAtStartPar
\textasciicircum{}1: Ruhr\sphinxhyphen{}University Bochum, Bochum, Germany

\sphinxAtStartPar
\textasciicircum{}2: University of Cambridge, Cambridge, United Kingdom

\sphinxAtStartPar
\textasciicircum{}3: Sandia National Laboratories, Albuquerque, New Mexico, USA

\sphinxAtStartPar
\textasciicircum{}4: University of British Columbia, Vancouver, BC, Canada
\end{quote}

\sphinxAtStartPar
\DUrole{versionmodified,added}{Added in version 14May2021.}

\sphinxAtStartPar
\sphinxstylestrong{Install:}

\sphinxAtStartPar
This package has {\hyperref[\detokenize{Build_extras:ml-pace}]{\sphinxcrossref{\DUrole{std,std-ref}{specific installation instructions}}}} on the
{\hyperref[\detokenize{Build_extras::doc}]{\sphinxcrossref{\DUrole{doc}{Build extras}}}} page.  This package may also be compiled
as a plugin to avoid licensing conflicts when distributing binaries.

\sphinxAtStartPar
\sphinxstylestrong{Supporting info:}
\begin{itemize}
\item {} 
\sphinxAtStartPar
\sphinxcode{\sphinxupquote{src/ML\sphinxhyphen{}PACE}}: filenames \sphinxhyphen{}\textgreater{} commands

\item {} 
\sphinxAtStartPar
{\hyperref[\detokenize{pair_pace::doc}]{\sphinxcrossref{\DUrole{doc}{pair\_style pace}}}}

\item {} 
\sphinxAtStartPar
\sphinxcode{\sphinxupquote{examples/PACKAGES/pace}}

\end{itemize}


\bigskip\hrule\bigskip



\subsection{ML\sphinxhyphen{}POD package}
\label{\detokenize{Packages_details:ml-pod-package}}\label{\detokenize{Packages_details:pkg-ml-pod}}
\sphinxAtStartPar
\sphinxstylestrong{Contents:}

\sphinxAtStartPar
A pair style and fitpod style for Proper Orthogonal Descriptors
(POD). POD is a methodology for deriving descriptors based on the proper
orthogonal decomposition. The ML\sphinxhyphen{}POD package provides an efficient
implementation for running simulations with POD potentials, along with
fitting the potentials natively in LAMMPS.

\sphinxAtStartPar
\sphinxstylestrong{Authors:}

\sphinxAtStartPar
Ngoc Cuong Nguyen (MIT), Andrew Rohskopf (Sandia)

\sphinxAtStartPar
\DUrole{versionmodified,added}{Added in version 22Dec2022.}

\sphinxAtStartPar
\sphinxstylestrong{Install:}

\sphinxAtStartPar
This package has {\hyperref[\detokenize{Build_extras:ml-pod}]{\sphinxcrossref{\DUrole{std,std-ref}{specific installation instructions}}}} on the
{\hyperref[\detokenize{Build_extras::doc}]{\sphinxcrossref{\DUrole{doc}{Build extras}}}} page.

\sphinxAtStartPar
\sphinxstylestrong{Supporting info:}
\begin{itemize}
\item {} 
\sphinxAtStartPar
\sphinxcode{\sphinxupquote{src/ML\sphinxhyphen{}POD}}: filenames \sphinxhyphen{}\textgreater{} commands

\item {} 
\sphinxAtStartPar
{\hyperref[\detokenize{pair_pod::doc}]{\sphinxcrossref{\DUrole{doc}{pair\_style pod}}}}

\item {} 
\sphinxAtStartPar
{\hyperref[\detokenize{fitpod_command::doc}]{\sphinxcrossref{\DUrole{doc}{command\_style fitpod}}}}

\item {} 
\sphinxAtStartPar
\sphinxcode{\sphinxupquote{examples/PACKAGES/pod}}

\end{itemize}


\bigskip\hrule\bigskip



\subsection{ML\sphinxhyphen{}QUIP package}
\label{\detokenize{Packages_details:ml-quip-package}}\label{\detokenize{Packages_details:pkg-ml-quip}}
\sphinxAtStartPar
\sphinxstylestrong{Contents:}

\sphinxAtStartPar
A {\hyperref[\detokenize{pair_quip::doc}]{\sphinxcrossref{\DUrole{doc}{pair\_style quip}}}} command which wraps the \sphinxhref{https://github.com/libAtoms/QUIP}{QUIP
libAtoms library}, which includes a variety of interatomic
potentials, including Gaussian Approximation Potential (GAP) models
developed by the Cambridge University group.

\sphinxAtStartPar
To use this package you must have the QUIP libAtoms library available
on your system.

\sphinxAtStartPar
\sphinxstylestrong{Author:} Albert Bartok (Cambridge University)

\sphinxAtStartPar
\sphinxstylestrong{Install:}

\sphinxAtStartPar
This package has {\hyperref[\detokenize{Build_extras:ml-quip}]{\sphinxcrossref{\DUrole{std,std-ref}{specific installation instructions}}}} on the
{\hyperref[\detokenize{Build_extras::doc}]{\sphinxcrossref{\DUrole{doc}{Build extras}}}} page.

\sphinxAtStartPar
\sphinxstylestrong{Supporting info:}
\begin{itemize}
\item {} 
\sphinxAtStartPar
\sphinxcode{\sphinxupquote{src/ML\sphinxhyphen{}QUIP}}: filenames \sphinxhyphen{}\textgreater{} commands

\item {} 
\sphinxAtStartPar
\sphinxcode{\sphinxupquote{src/ML\sphinxhyphen{}QUIP/README}}

\item {} 
\sphinxAtStartPar
{\hyperref[\detokenize{pair_quip::doc}]{\sphinxcrossref{\DUrole{doc}{pair\_style quip}}}}

\item {} 
\sphinxAtStartPar
\sphinxcode{\sphinxupquote{examples/PACKAGES/quip}}

\end{itemize}


\bigskip\hrule\bigskip



\subsection{ML\sphinxhyphen{}RANN package}
\label{\detokenize{Packages_details:ml-rann-package}}\label{\detokenize{Packages_details:pkg-ml-rann}}
\sphinxAtStartPar
\sphinxstylestrong{Contents:}

\sphinxAtStartPar
A pair style for using rapid atomistic neural network (RANN) potentials.
These neural network potentials work by first generating a series of symmetry
functions from the neighbor list and then using these values as the input layer
of a neural network.

\sphinxAtStartPar
\sphinxstylestrong{Authors:}

\sphinxAtStartPar
This package was written by Christopher Barrett
with contributions by Doyl Dickel, Mississippi State University.

\sphinxAtStartPar
\DUrole{versionmodified,added}{Added in version 27May2021.}

\sphinxAtStartPar
\sphinxstylestrong{Supporting info:}
\begin{itemize}
\item {} 
\sphinxAtStartPar
\sphinxcode{\sphinxupquote{src/ML\sphinxhyphen{}RANN}}: filenames \sphinxhyphen{}\textgreater{} commands

\item {} 
\sphinxAtStartPar
{\hyperref[\detokenize{pair_rann::doc}]{\sphinxcrossref{\DUrole{doc}{pair\_style rann}}}}

\item {} 
\sphinxAtStartPar
\sphinxcode{\sphinxupquote{examples/PACKAGES/rann}}

\end{itemize}


\bigskip\hrule\bigskip



\subsection{ML\sphinxhyphen{}SNAP package}
\label{\detokenize{Packages_details:ml-snap-package}}\label{\detokenize{Packages_details:pkg-ml-snap}}
\sphinxAtStartPar
\sphinxstylestrong{Contents:}

\sphinxAtStartPar
A pair style for the spectral neighbor analysis potential (SNAP).
SNAP is methodology for deriving a highly accurate classical potential
fit to a large archive of quantum mechanical (DFT) data. Also several
computes which analyze attributes of the potential.

\sphinxAtStartPar
\sphinxstylestrong{Author:} Aidan Thompson (Sandia).

\sphinxAtStartPar
\sphinxstylestrong{Supporting info:}
\begin{itemize}
\item {} 
\sphinxAtStartPar
\sphinxcode{\sphinxupquote{src/ML\sphinxhyphen{}SNAP}}: filenames \sphinxhyphen{}\textgreater{} commands

\item {} 
\sphinxAtStartPar
{\hyperref[\detokenize{pair_snap::doc}]{\sphinxcrossref{\DUrole{doc}{pair\_style snap}}}}

\item {} 
\sphinxAtStartPar
{\hyperref[\detokenize{compute_sna_atom::doc}]{\sphinxcrossref{\DUrole{doc}{compute sna/atom}}}}

\item {} 
\sphinxAtStartPar
{\hyperref[\detokenize{compute_sna_atom::doc}]{\sphinxcrossref{\DUrole{doc}{compute sna/grid}}}}

\item {} 
\sphinxAtStartPar
{\hyperref[\detokenize{compute_sna_atom::doc}]{\sphinxcrossref{\DUrole{doc}{compute sna/grid/local}}}}

\item {} 
\sphinxAtStartPar
{\hyperref[\detokenize{compute_sna_atom::doc}]{\sphinxcrossref{\DUrole{doc}{compute snad/atom}}}}

\item {} 
\sphinxAtStartPar
{\hyperref[\detokenize{compute_sna_atom::doc}]{\sphinxcrossref{\DUrole{doc}{compute snav/atom}}}}

\item {} 
\sphinxAtStartPar
\sphinxcode{\sphinxupquote{examples/snap}}

\end{itemize}


\bigskip\hrule\bigskip



\subsection{ML\sphinxhyphen{}UF3 package}
\label{\detokenize{Packages_details:ml-uf3-package}}\label{\detokenize{Packages_details:pkg-ml-uf3}}
\sphinxAtStartPar
\sphinxstylestrong{Contents:}

\sphinxAtStartPar
A pair style for the ultra\sphinxhyphen{}fast force field potentials (UF3). UF3 is a
methodology for deriving a highly accurate classical potential which is
fast to evaluate and is fitted to a large archives of quantum mechanical
(DFT) data.  The use of b\sphinxhyphen{}spline basis set in UF3 enables the rapid
evaluation of 2\sphinxhyphen{}body and 3\sphinxhyphen{}body interactions.

\sphinxAtStartPar
\sphinxstylestrong{Authors:} Ajinkya C Hire (University of Florida),
Hendrik Krass (University of Constance),
Matthias Rupp (Luxembourg Institute of Science and Technology),
Richard Hennig (University of Florida)

\sphinxAtStartPar
\sphinxstylestrong{Supporting info:}
\begin{itemize}
\item {} 
\sphinxAtStartPar
\sphinxcode{\sphinxupquote{src/ML\sphinxhyphen{}UF3}}: filenames \sphinxhyphen{}\textgreater{} commands

\item {} 
\sphinxAtStartPar
{\hyperref[\detokenize{pair_uf3::doc}]{\sphinxcrossref{\DUrole{doc}{pair\_style uf3}}}}

\item {} 
\sphinxAtStartPar
\sphinxcode{\sphinxupquote{examples/uf3}}

\item {} 
\sphinxAtStartPar
\sphinxurl{https://github.com/uf3/uf3}

\end{itemize}


\subsection{MOFFF package}
\label{\detokenize{Packages_details:mofff-package}}\label{\detokenize{Packages_details:pkg-mofff}}
\sphinxAtStartPar
\sphinxstylestrong{Contents:}

\sphinxAtStartPar
Pair, angle and improper styles needed to employ the MOF\sphinxhyphen{}FF
force field by Schmid and coworkers with LAMMPS.
MOF\sphinxhyphen{}FF is a first principles derived force field with the primary aim
to simulate MOFs and related porous framework materials, using spherical
Gaussian charges. It is described in S. Bureekaew et al., Phys. Stat. Sol. B
2013, 250, 1128\sphinxhyphen{}1141.
For the usage of MOF\sphinxhyphen{}FF see the example in the example directory as
well as the \sphinxhref{https://www.mofplus.org/content/show/MOF-FF}{MOF+} website.

\sphinxAtStartPar
\sphinxstylestrong{Author:} Hendrik Heenen (Technical U of Munich),
Rochus Schmid (Ruhr\sphinxhyphen{}University Bochum).

\sphinxAtStartPar
\sphinxstylestrong{Supporting info:}
\begin{itemize}
\item {} 
\sphinxAtStartPar
\sphinxcode{\sphinxupquote{src/MOFFF}}: filenames \sphinxhyphen{}\textgreater{} commands

\item {} 
\sphinxAtStartPar
\sphinxcode{\sphinxupquote{src/MOFFF/README}}

\item {} 
\sphinxAtStartPar
{\hyperref[\detokenize{pair_buck6d_coul_gauss::doc}]{\sphinxcrossref{\DUrole{doc}{pair\_style buck6d/coul/gauss}}}}

\item {} 
\sphinxAtStartPar
{\hyperref[\detokenize{angle_class2::doc}]{\sphinxcrossref{\DUrole{doc}{angle\_style class2}}}}

\item {} 
\sphinxAtStartPar
{\hyperref[\detokenize{angle_cosine_buck6d::doc}]{\sphinxcrossref{\DUrole{doc}{angle\_style cosine/buck6d}}}}

\item {} 
\sphinxAtStartPar
{\hyperref[\detokenize{improper_inversion_harmonic::doc}]{\sphinxcrossref{\DUrole{doc}{improper\_style inversion/harmonic}}}}

\item {} 
\sphinxAtStartPar
\sphinxcode{\sphinxupquote{examples/PACKAGES/mofff}}

\end{itemize}


\bigskip\hrule\bigskip



\subsection{MOLECULE package}
\label{\detokenize{Packages_details:molecule-package}}\label{\detokenize{Packages_details:pkg-molecule}}
\sphinxAtStartPar
\sphinxstylestrong{Contents:}

\sphinxAtStartPar
A large number of atom, pair, bond, angle, dihedral, improper styles
that are used to model molecular systems with fixed covalent bonds.
The pair styles include the Dreiding (hydrogen\sphinxhyphen{}bonding) and CHARMM
force fields, and a TIP4P water model.

\sphinxAtStartPar
\sphinxstylestrong{Supporting info:}
\begin{itemize}
\item {} 
\sphinxAtStartPar
\sphinxcode{\sphinxupquote{src/MOLECULE}}: filenames \sphinxhyphen{}\textgreater{} commands

\item {} 
\sphinxAtStartPar
{\hyperref[\detokenize{atom_style::doc}]{\sphinxcrossref{\DUrole{doc}{atom\_style}}}}

\item {} 
\sphinxAtStartPar
{\hyperref[\detokenize{bond_style::doc}]{\sphinxcrossref{\DUrole{doc}{bond\_style}}}}

\item {} 
\sphinxAtStartPar
{\hyperref[\detokenize{angle_style::doc}]{\sphinxcrossref{\DUrole{doc}{angle\_style}}}}

\item {} 
\sphinxAtStartPar
{\hyperref[\detokenize{dihedral_style::doc}]{\sphinxcrossref{\DUrole{doc}{dihedral\_style}}}}

\item {} 
\sphinxAtStartPar
{\hyperref[\detokenize{improper_style::doc}]{\sphinxcrossref{\DUrole{doc}{improper\_style}}}}

\item {} 
\sphinxAtStartPar
{\hyperref[\detokenize{pair_hbond_dreiding::doc}]{\sphinxcrossref{\DUrole{doc}{pair\_style hbond/dreiding/lj}}}}

\item {} 
\sphinxAtStartPar
{\hyperref[\detokenize{pair_charmm::doc}]{\sphinxcrossref{\DUrole{doc}{pair\_style lj/charmm/coul/charmm}}}}

\item {} 
\sphinxAtStartPar
{\hyperref[\detokenize{Howto_bioFF::doc}]{\sphinxcrossref{\DUrole{doc}{Howto bioFF}}}}

\item {} 
\sphinxAtStartPar
\sphinxcode{\sphinxupquote{examples/cmap}}

\item {} 
\sphinxAtStartPar
\sphinxcode{\sphinxupquote{examples/dreiding}}

\item {} 
\sphinxAtStartPar
\sphinxcode{\sphinxupquote{examples/micelle,}}

\item {} 
\sphinxAtStartPar
\sphinxcode{\sphinxupquote{examples/peptide}}

\item {} 
\sphinxAtStartPar
bench/in.chain

\item {} 
\sphinxAtStartPar
bench/in.rhodo

\end{itemize}


\bigskip\hrule\bigskip



\subsection{MOLFILE package}
\label{\detokenize{Packages_details:molfile-package}}\label{\detokenize{Packages_details:pkg-molfile}}
\sphinxAtStartPar
\sphinxstylestrong{Contents:}

\sphinxAtStartPar
A {\hyperref[\detokenize{dump_molfile::doc}]{\sphinxcrossref{\DUrole{doc}{dump molfile}}}} command which uses molfile plugins
that are bundled with the \sphinxhref{https://www.ks.uiuc.edu/Research/vmd/}{VMD}
molecular visualization and analysis program, to enable LAMMPS to dump
snapshots in formats compatible with various molecular simulation
tools.

\sphinxAtStartPar
To use this package you must have the desired VMD plugins available on
your system.

\sphinxAtStartPar
Note that this package only provides the interface code, not the
plugins themselves, which will be accessed when requesting a specific
plugin via the {\hyperref[\detokenize{dump_molfile::doc}]{\sphinxcrossref{\DUrole{doc}{dump molfile}}}} command.  Plugins can
be obtained from a VMD installation which has to match the platform
that you are using to compile LAMMPS for. By adding plugins to VMD,
support for new file formats can be added to LAMMPS (or VMD or other
programs that use them) without having to re\sphinxhyphen{}compile the application
itself.  More information about the VMD molfile plugins can be found
at
\sphinxurl{https://www.ks.uiuc.edu/Research/vmd/plugins/molfile}.

\sphinxAtStartPar
\sphinxstylestrong{Author:} Axel Kohlmeyer (Temple U).

\sphinxAtStartPar
\sphinxstylestrong{Install:}

\sphinxAtStartPar
This package has {\hyperref[\detokenize{Build_extras:molfile}]{\sphinxcrossref{\DUrole{std,std-ref}{specific installation instructions}}}} on the {\hyperref[\detokenize{Build_extras::doc}]{\sphinxcrossref{\DUrole{doc}{Build extras}}}} page.

\sphinxAtStartPar
\sphinxstylestrong{Supporting info:}
\begin{itemize}
\item {} 
\sphinxAtStartPar
\sphinxcode{\sphinxupquote{src/MOLFILE}}: filenames \sphinxhyphen{}\textgreater{} commands

\item {} 
\sphinxAtStartPar
\sphinxcode{\sphinxupquote{src/MOLFILE/README}}

\item {} 
\sphinxAtStartPar
\sphinxcode{\sphinxupquote{lib/molfile/README}}

\item {} 
\sphinxAtStartPar
{\hyperref[\detokenize{dump_molfile::doc}]{\sphinxcrossref{\DUrole{doc}{dump molfile}}}}

\end{itemize}


\bigskip\hrule\bigskip



\subsection{NETCDF package}
\label{\detokenize{Packages_details:netcdf-package}}\label{\detokenize{Packages_details:pkg-netcdf}}
\sphinxAtStartPar
\sphinxstylestrong{Contents:}

\sphinxAtStartPar
Dump styles for writing NetCDF formatted dump files.  NetCDF is a
portable, binary, self\sphinxhyphen{}describing file format developed on top of
HDF5. The file contents follow the AMBER NetCDF trajectory conventions
(\sphinxurl{https://ambermd.org/netcdf/nctraj.xhtml}), but include extensions.

\sphinxAtStartPar
To use this package you must have the NetCDF library available on your
system.

\sphinxAtStartPar
Note that NetCDF files can be directly visualized with the following
tools:
\begin{itemize}
\item {} 
\sphinxAtStartPar
\sphinxhref{https://www.ovito.org}{Ovito} (Ovito supports the AMBER convention and the extensions mentioned above)

\item {} 
\sphinxAtStartPar
\sphinxhref{https://www.ks.uiuc.edu/Research/vmd/}{VMD}

\end{itemize}

\sphinxAtStartPar
\sphinxstylestrong{Author:} Lars Pastewka (Karlsruhe Institute of Technology).

\sphinxAtStartPar
\sphinxstylestrong{Install:}

\sphinxAtStartPar
This package has {\hyperref[\detokenize{Build_extras:netcdf}]{\sphinxcrossref{\DUrole{std,std-ref}{specific installation instructions}}}} on the {\hyperref[\detokenize{Build_extras::doc}]{\sphinxcrossref{\DUrole{doc}{Build extras}}}} page.

\sphinxAtStartPar
\sphinxstylestrong{Supporting info:}
\begin{itemize}
\item {} 
\sphinxAtStartPar
\sphinxcode{\sphinxupquote{src/NETCDF}}: filenames \sphinxhyphen{}\textgreater{} commands

\item {} 
\sphinxAtStartPar
\sphinxcode{\sphinxupquote{src/NETCDF/README}}

\item {} 
\sphinxAtStartPar
\sphinxcode{\sphinxupquote{lib/netcdf/README}}

\item {} 
\sphinxAtStartPar
{\hyperref[\detokenize{dump_netcdf::doc}]{\sphinxcrossref{\DUrole{doc}{dump netcdf}}}}

\end{itemize}


\bigskip\hrule\bigskip



\subsection{OPENMP package}
\label{\detokenize{Packages_details:openmp-package}}\label{\detokenize{Packages_details:pkg-openmp}}
\sphinxAtStartPar
\sphinxstylestrong{Contents:}

\sphinxAtStartPar
Hundreds of pair, fix, compute, bond, angle, dihedral, improper, and
kspace styles which are altered to enable threading on many\sphinxhyphen{}core CPUs
via OpenMP directives.  All of them have an “omp” in their style name.
The {\hyperref[\detokenize{Speed_omp::doc}]{\sphinxcrossref{\DUrole{doc}{OPENMP package}}}} page gives details of what hardware
and compilers are required on your system, and how to build and use
this package.  Its styles can be invoked at run time via the \sphinxcode{\sphinxupquote{\sphinxhyphen{}sf omp}}
or \sphinxcode{\sphinxupquote{\sphinxhyphen{}suffix omp}} {\hyperref[\detokenize{Run_options::doc}]{\sphinxcrossref{\DUrole{doc}{command\sphinxhyphen{}line switches}}}}.  Also see
the {\hyperref[\detokenize{Packages_details:pkg-kokkos}]{\sphinxcrossref{\DUrole{std,std-ref}{KOKKOS}}}}, {\hyperref[\detokenize{Packages_details:pkg-opt}]{\sphinxcrossref{\DUrole{std,std-ref}{OPT}}}}, and {\hyperref[\detokenize{Packages_details:pkg-intel}]{\sphinxcrossref{\DUrole{std,std-ref}{INTEL}}}}
packages, which have styles optimized for CPUs.

\sphinxAtStartPar
\sphinxstylestrong{Author:} Axel Kohlmeyer (Temple U).

\begin{sphinxadmonition}{note}{Note:}
\sphinxAtStartPar
To enable multi\sphinxhyphen{}threading support the compile flag \sphinxcode{\sphinxupquote{\sphinxhyphen{}fopenmp}}
and the link flag \sphinxcode{\sphinxupquote{\sphinxhyphen{}fopenmp}} (for GNU compilers, you have to look up
the equivalent flags for other compilers) must be used to build LAMMPS.
When using Intel compilers, also the \sphinxcode{\sphinxupquote{\sphinxhyphen{}restrict}} flag is required.
The OPENMP package can be compiled without enabling OpenMP; then
all code will be compiled as serial and the only improvement over the
regular styles are some data access optimization. These flags should
be added to the CCFLAGS and LINKFLAGS lines of your Makefile.machine.
See \sphinxcode{\sphinxupquote{src/MAKE/OPTIONS/Makefile}}.omp for an example.
\end{sphinxadmonition}

\sphinxAtStartPar
Once you have an appropriate Makefile.machine, you can
install/uninstall the package and build LAMMPS in the usual manner:

\sphinxAtStartPar
\sphinxstylestrong{Install:}

\sphinxAtStartPar
This package has {\hyperref[\detokenize{Build_extras:openmp}]{\sphinxcrossref{\DUrole{std,std-ref}{specific installation instructions}}}} on
the {\hyperref[\detokenize{Build_extras::doc}]{\sphinxcrossref{\DUrole{doc}{Build extras}}}} page.

\sphinxAtStartPar
\sphinxstylestrong{Supporting info:}
\begin{itemize}
\item {} 
\sphinxAtStartPar
\sphinxcode{\sphinxupquote{src/OPENMP}}: filenames \sphinxhyphen{}\textgreater{} commands

\item {} 
\sphinxAtStartPar
\sphinxcode{\sphinxupquote{src/OPENMP/README}}

\item {} 
\sphinxAtStartPar
{\hyperref[\detokenize{Speed_packages::doc}]{\sphinxcrossref{\DUrole{doc}{Accelerator packages}}}}

\item {} 
\sphinxAtStartPar
{\hyperref[\detokenize{Speed_omp::doc}]{\sphinxcrossref{\DUrole{doc}{OPENMP package}}}}

\item {} 
\sphinxAtStartPar
{\hyperref[\detokenize{Run_options::doc}]{\sphinxcrossref{\DUrole{doc}{Command line option \sphinxhyphen{}suffix/\sphinxhyphen{}sf omp}}}}

\item {} 
\sphinxAtStartPar
{\hyperref[\detokenize{Run_options::doc}]{\sphinxcrossref{\DUrole{doc}{Command line option \sphinxhyphen{}package/\sphinxhyphen{}pk omp}}}}

\item {} 
\sphinxAtStartPar
{\hyperref[\detokenize{package::doc}]{\sphinxcrossref{\DUrole{doc}{package omp}}}}

\item {} 
\sphinxAtStartPar
Search the {\hyperref[\detokenize{Commands_all::doc}]{\sphinxcrossref{\DUrole{doc}{commands}}}} pages ({\hyperref[\detokenize{Commands_fix::doc}]{\sphinxcrossref{\DUrole{doc}{fix}}}}, {\hyperref[\detokenize{Commands_compute::doc}]{\sphinxcrossref{\DUrole{doc}{compute}}}},
{\hyperref[\detokenize{Commands_pair::doc}]{\sphinxcrossref{\DUrole{doc}{pair}}}}, {\hyperref[\detokenize{Commands_bond::doc}]{\sphinxcrossref{\DUrole{doc}{bond, angle, dihedral, improper}}}},
{\hyperref[\detokenize{Commands_kspace::doc}]{\sphinxcrossref{\DUrole{doc}{kspace}}}}) for styles followed by (o)

\item {} 
\sphinxAtStartPar
\sphinxhref{https://www.lammps.org/bench.html}{Benchmarks page} of website

\end{itemize}


\bigskip\hrule\bigskip



\subsection{OPT package}
\label{\detokenize{Packages_details:opt-package}}\label{\detokenize{Packages_details:pkg-opt}}
\sphinxAtStartPar
\sphinxstylestrong{Contents:}

\sphinxAtStartPar
A handful of pair styles which are optimized for improved CPU
performance on single or multiple cores.  These include EAM, LJ,
CHARMM, and Morse potentials.  The styles have an “opt” suffix in
their style name.  The {\hyperref[\detokenize{Speed_opt::doc}]{\sphinxcrossref{\DUrole{doc}{OPT package}}}} page gives
details of how to build and use this package.  Its styles can be
invoked at run time via the \sphinxcode{\sphinxupquote{\sphinxhyphen{}sf opt}} or \sphinxcode{\sphinxupquote{\sphinxhyphen{}suffix opt}} {\hyperref[\detokenize{Run_options::doc}]{\sphinxcrossref{\DUrole{doc}{command\sphinxhyphen{}line switches}}}}.  See also the {\hyperref[\detokenize{Packages_details:pkg-kokkos}]{\sphinxcrossref{\DUrole{std,std-ref}{KOKKOS}}}},
{\hyperref[\detokenize{Packages_details:pkg-intel}]{\sphinxcrossref{\DUrole{std,std-ref}{INTEL}}}}, and {\hyperref[\detokenize{Packages_details:pkg-openmp}]{\sphinxcrossref{\DUrole{std,std-ref}{OPENMP}}}} packages, which
have styles optimized for CPU performance.

\sphinxAtStartPar
\sphinxstylestrong{Authors:} James Fischer (High Performance Technologies), David Richie,
and Vincent Natoli (Stone Ridge Technology).

\sphinxAtStartPar
\sphinxstylestrong{Install:}

\sphinxAtStartPar
This package has {\hyperref[\detokenize{Build_extras:opt}]{\sphinxcrossref{\DUrole{std,std-ref}{specific installation instructions}}}} on the {\hyperref[\detokenize{Build_extras::doc}]{\sphinxcrossref{\DUrole{doc}{Build extras}}}} page.

\sphinxAtStartPar
\sphinxstylestrong{Supporting info:}
\begin{itemize}
\item {} 
\sphinxAtStartPar
\sphinxcode{\sphinxupquote{src/OPT}}: filenames \sphinxhyphen{}\textgreater{} commands

\item {} 
\sphinxAtStartPar
{\hyperref[\detokenize{Speed_packages::doc}]{\sphinxcrossref{\DUrole{doc}{Accelerator packages}}}}

\item {} 
\sphinxAtStartPar
{\hyperref[\detokenize{Speed_opt::doc}]{\sphinxcrossref{\DUrole{doc}{OPT package}}}}

\item {} 
\sphinxAtStartPar
{\hyperref[\detokenize{Run_options::doc}]{\sphinxcrossref{\DUrole{doc}{Section 2.6 \sphinxhyphen{}sf opt}}}}

\item {} 
\sphinxAtStartPar
Search the {\hyperref[\detokenize{Commands_pair::doc}]{\sphinxcrossref{\DUrole{doc}{pair style}}}} page for styles followed by (t)

\item {} 
\sphinxAtStartPar
\sphinxhref{https://www.lammps.org/bench.html}{Benchmarks page} of website

\end{itemize}


\subsection{ORIENT package}
\label{\detokenize{Packages_details:orient-package}}\label{\detokenize{Packages_details:pkg-orient}}
\sphinxAtStartPar
\sphinxstylestrong{Contents:}

\sphinxAtStartPar
A few fixes that apply orientation dependent forces for studying
grain boundary migration.

\sphinxAtStartPar
\sphinxstylestrong{Supporting info:}
\begin{itemize}
\item {} 
\sphinxAtStartPar
\sphinxcode{\sphinxupquote{src/ORIENT}}: filenames \sphinxhyphen{}\textgreater{} commands

\item {} 
\sphinxAtStartPar
{\hyperref[\detokenize{fix_orient::doc}]{\sphinxcrossref{\DUrole{doc}{fix orient/bcc}}}}

\item {} 
\sphinxAtStartPar
{\hyperref[\detokenize{fix_orient::doc}]{\sphinxcrossref{\DUrole{doc}{fix orient/fcc}}}}

\item {} 
\sphinxAtStartPar
{\hyperref[\detokenize{fix_orient_eco::doc}]{\sphinxcrossref{\DUrole{doc}{fix orient/eco}}}}

\end{itemize}


\bigskip\hrule\bigskip



\subsection{PERI package}
\label{\detokenize{Packages_details:peri-package}}\label{\detokenize{Packages_details:pkg-peri}}
\sphinxAtStartPar
\sphinxstylestrong{Contents:}

\sphinxAtStartPar
An atom style, several pair styles which implement different
Peridynamics materials models, and several computes which calculate
diagnostics.  Peridynamics is a particle\sphinxhyphen{}based meshless continuum
model.

\sphinxAtStartPar
\sphinxstylestrong{Authors:} The original package was created by Mike Parks (Sandia).
Additional Peridynamics models were added by Rezwanur Rahman and John
Foster (UTSA).

\sphinxAtStartPar
\sphinxstylestrong{Supporting info:}
\begin{itemize}
\item {} 
\sphinxAtStartPar
\sphinxcode{\sphinxupquote{src/PERI}}: filenames \sphinxhyphen{}\textgreater{} commands

\item {} 
\sphinxAtStartPar
{\hyperref[\detokenize{Howto_peri::doc}]{\sphinxcrossref{\DUrole{doc}{Peridynamics Howto}}}}

\item {} 
\sphinxAtStartPar
\sphinxhref{PDF/PDLammps\_overview.pdf}{doc/PDF/PDLammps\_overview.pdf}

\item {} 
\sphinxAtStartPar
\sphinxhref{PDF/PDLammps\_EPS.pdf}{doc/PDF/PDLammps\_EPS.pdf}

\item {} 
\sphinxAtStartPar
\sphinxhref{PDF/PDLammps\_VES.pdf}{doc/PDF/PDLammps\_VES.pdf}

\item {} 
\sphinxAtStartPar
{\hyperref[\detokenize{atom_style::doc}]{\sphinxcrossref{\DUrole{doc}{atom\_style peri}}}}

\item {} 
\sphinxAtStartPar
{\hyperref[\detokenize{pair_peri::doc}]{\sphinxcrossref{\DUrole{doc}{pair\_style peri/*}}}}

\item {} 
\sphinxAtStartPar
{\hyperref[\detokenize{compute_damage_atom::doc}]{\sphinxcrossref{\DUrole{doc}{compute damage/atom}}}}

\item {} 
\sphinxAtStartPar
{\hyperref[\detokenize{compute_plasticity_atom::doc}]{\sphinxcrossref{\DUrole{doc}{compute plasticity/atom}}}}

\item {} 
\sphinxAtStartPar
\sphinxcode{\sphinxupquote{examples/peri}}

\item {} 
\sphinxAtStartPar
\sphinxurl{https://www.lammps.org/movies.html\#peri}

\end{itemize}


\bigskip\hrule\bigskip



\subsection{PHONON package}
\label{\detokenize{Packages_details:phonon-package}}\label{\detokenize{Packages_details:pkg-phonon}}
\sphinxAtStartPar
\sphinxstylestrong{Contents:}

\sphinxAtStartPar
A {\hyperref[\detokenize{fix_phonon::doc}]{\sphinxcrossref{\DUrole{doc}{fix phonon}}}} command that calculates dynamical
matrices, which can then be used to compute phonon dispersion
relations, directly from molecular dynamics simulations.
And a {\hyperref[\detokenize{dynamical_matrix::doc}]{\sphinxcrossref{\DUrole{doc}{dynamical\_matrix}}}} as well as a
{\hyperref[\detokenize{third_order::doc}]{\sphinxcrossref{\DUrole{doc}{third\_order}}}} command to compute the dynamical matrix
and third order tensor from finite differences.

\sphinxAtStartPar
\sphinxstylestrong{Install:}

\sphinxAtStartPar
The fix phonon command also requires that the {\hyperref[\detokenize{Packages_details:pkg-kspace}]{\sphinxcrossref{\DUrole{std,std-ref}{KSPACE}}}}
package is installed.

\sphinxAtStartPar
\sphinxstylestrong{Authors:} Ling\sphinxhyphen{}Ti Kong (Shanghai Jiao Tong University) for “fix phonon”
and Charlie Sievers (UC Davis) for “dynamical\_matrix” and “third\_order”

\sphinxAtStartPar
\sphinxstylestrong{Supporting info:}
\begin{itemize}
\item {} 
\sphinxAtStartPar
\sphinxcode{\sphinxupquote{src/PHONON}}: filenames \sphinxhyphen{}\textgreater{} commands

\item {} 
\sphinxAtStartPar
\sphinxcode{\sphinxupquote{src/PHONON/README}}

\item {} 
\sphinxAtStartPar
{\hyperref[\detokenize{fix_phonon::doc}]{\sphinxcrossref{\DUrole{doc}{fix phonon}}}}

\item {} 
\sphinxAtStartPar
{\hyperref[\detokenize{dynamical_matrix::doc}]{\sphinxcrossref{\DUrole{doc}{dynamical\_matrix}}}}

\item {} 
\sphinxAtStartPar
{\hyperref[\detokenize{third_order::doc}]{\sphinxcrossref{\DUrole{doc}{third\_order}}}}

\item {} 
\sphinxAtStartPar
\sphinxcode{\sphinxupquote{examples/PACKAGES/phonon}}

\end{itemize}


\bigskip\hrule\bigskip



\subsection{PLUGIN package}
\label{\detokenize{Packages_details:plugin-package}}\label{\detokenize{Packages_details:pkg-plugin}}
\sphinxAtStartPar
\sphinxstylestrong{Contents:}

\sphinxAtStartPar
A {\hyperref[\detokenize{plugin::doc}]{\sphinxcrossref{\DUrole{doc}{plugin}}}} command that can load and unload several
kind of styles in LAMMPS from shared object files at runtime without
having to recompile and relink LAMMPS.

\sphinxAtStartPar
When the environment variable \sphinxcode{\sphinxupquote{LAMMPS\_PLUGIN\_PATH}} is set, then LAMMPS
will search the directory (or directories) listed in this path for files
with names that end in \sphinxcode{\sphinxupquote{plugin.so}} (e.g. \sphinxcode{\sphinxupquote{helloplugin.so}}) and will
try to load the contained plugins automatically at start\sphinxhyphen{}up.

\sphinxAtStartPar
\sphinxstylestrong{Authors:} Axel Kohlmeyer (Temple U)

\sphinxAtStartPar
\DUrole{versionmodified,added}{Added in version 8Apr2021.}

\sphinxAtStartPar
\sphinxstylestrong{Supporting info:}
\begin{itemize}
\item {} 
\sphinxAtStartPar
\sphinxcode{\sphinxupquote{src/PLUGIN}}: filenames \sphinxhyphen{}\textgreater{} commands

\item {} 
\sphinxAtStartPar
{\hyperref[\detokenize{plugin::doc}]{\sphinxcrossref{\DUrole{doc}{plugin command}}}}

\item {} 
\sphinxAtStartPar
{\hyperref[\detokenize{Developer_plugins::doc}]{\sphinxcrossref{\DUrole{doc}{Information on writing plugins}}}}

\item {} 
\sphinxAtStartPar
\sphinxcode{\sphinxupquote{examples/plugin}}

\end{itemize}


\bigskip\hrule\bigskip



\subsection{PLUMED package}
\label{\detokenize{Packages_details:plumed-package}}\label{\detokenize{Packages_details:pkg-plumed}}
\sphinxAtStartPar
\sphinxstylestrong{Contents:}

\sphinxAtStartPar
The fix plumed command allows you to use the PLUMED free energy plugin
for molecular dynamics to analyze and bias your LAMMPS trajectory on
the fly.  The PLUMED library is called from within the LAMMPS input
script by using the {\hyperref[\detokenize{fix_plumed::doc}]{\sphinxcrossref{\DUrole{doc}{fix plumed}}}} command.

\sphinxAtStartPar
\sphinxstylestrong{Authors:} The \sphinxhref{https://www.plumed.org}{PLUMED library} is written
and maintained by Massimilliano Bonomi, Giovanni Bussi, Carlo Camiloni,
and Gareth Tribello.

\sphinxAtStartPar
\sphinxstylestrong{Install:}

\sphinxAtStartPar
This package has {\hyperref[\detokenize{Build_extras:plumed}]{\sphinxcrossref{\DUrole{std,std-ref}{specific installation instructions}}}} on the
{\hyperref[\detokenize{Build_extras::doc}]{\sphinxcrossref{\DUrole{doc}{Build extras}}}} page.  This package may also be compiled
as a plugin to avoid licensing conflicts when distributing binaries.

\sphinxAtStartPar
\sphinxstylestrong{Supporting info:}
\begin{itemize}
\item {} 
\sphinxAtStartPar
\sphinxcode{\sphinxupquote{src/PLUMED/README}}

\item {} 
\sphinxAtStartPar
\sphinxcode{\sphinxupquote{lib/plumed/README}}

\item {} 
\sphinxAtStartPar
{\hyperref[\detokenize{fix_plumed::doc}]{\sphinxcrossref{\DUrole{doc}{fix plumed}}}}

\item {} 
\sphinxAtStartPar
\sphinxcode{\sphinxupquote{examples/PACKAGES/plumed}}

\end{itemize}


\bigskip\hrule\bigskip



\subsection{POEMS package}
\label{\detokenize{Packages_details:poems-package}}\label{\detokenize{Packages_details:pkg-poems}}
\sphinxAtStartPar
\sphinxstylestrong{Contents:}

\sphinxAtStartPar
A fix that wraps the Parallelizable Open source Efficient Multibody
Software (POEMS) library, which is able to simulate the dynamics of
articulated body systems.  These are systems with multiple rigid
bodies (collections of particles) whose motion is coupled by
connections at hinge points.

\sphinxAtStartPar
\sphinxstylestrong{Author:} Rudra Mukherjee (JPL) while at RPI.

\sphinxAtStartPar
\sphinxstylestrong{Install:}

\sphinxAtStartPar
This package has {\hyperref[\detokenize{Build_extras:poems}]{\sphinxcrossref{\DUrole{std,std-ref}{specific installation instructions}}}} on the {\hyperref[\detokenize{Build_extras::doc}]{\sphinxcrossref{\DUrole{doc}{Build extras}}}} page.

\sphinxAtStartPar
\sphinxstylestrong{Supporting info:}
\begin{itemize}
\item {} 
\sphinxAtStartPar
\sphinxcode{\sphinxupquote{src/POEMS}}: filenames \sphinxhyphen{}\textgreater{} commands

\item {} 
\sphinxAtStartPar
\sphinxcode{\sphinxupquote{src/POEMS/README}}

\item {} 
\sphinxAtStartPar
\sphinxcode{\sphinxupquote{lib/poems/README}}

\item {} 
\sphinxAtStartPar
{\hyperref[\detokenize{fix_poems::doc}]{\sphinxcrossref{\DUrole{doc}{fix poems}}}}

\item {} 
\sphinxAtStartPar
\sphinxcode{\sphinxupquote{examples/rigid}}

\end{itemize}


\bigskip\hrule\bigskip



\subsection{PTM package}
\label{\detokenize{Packages_details:ptm-package}}\label{\detokenize{Packages_details:pkg-ptm}}
\sphinxAtStartPar
\sphinxstylestrong{Contents:}

\sphinxAtStartPar
A {\hyperref[\detokenize{compute_ptm_atom::doc}]{\sphinxcrossref{\DUrole{doc}{compute ptm/atom}}}} command that calculates
local structure characterization using the Polyhedral Template
Matching methodology.

\sphinxAtStartPar
\sphinxstylestrong{Author:} Peter Mahler Larsen (MIT).

\sphinxAtStartPar
\sphinxstylestrong{Supporting info:}
\begin{itemize}
\item {} 
\sphinxAtStartPar
\sphinxcode{\sphinxupquote{src/PTM}}: filenames not starting with ptm\_ \sphinxhyphen{}\textgreater{} commands

\item {} 
\sphinxAtStartPar
\sphinxcode{\sphinxupquote{src/PTM}}: filenames starting with ptm\_ \sphinxhyphen{}\textgreater{} supporting code

\item {} 
\sphinxAtStartPar
\sphinxcode{\sphinxupquote{src/PTM/LICENSE}}

\item {} 
\sphinxAtStartPar
{\hyperref[\detokenize{compute_ptm_atom::doc}]{\sphinxcrossref{\DUrole{doc}{compute ptm/atom}}}}

\end{itemize}


\bigskip\hrule\bigskip



\subsection{PYTHON package}
\label{\detokenize{Packages_details:python-package}}\label{\detokenize{Packages_details:pkg-python}}
\sphinxAtStartPar
\sphinxstylestrong{Contents:}

\sphinxAtStartPar
A {\hyperref[\detokenize{python::doc}]{\sphinxcrossref{\DUrole{doc}{python}}}} command which allow you to execute Python code
from a LAMMPS input script.  The code can be in a separate file or
embedded in the input script itself.  See the {\hyperref[\detokenize{Python_call::doc}]{\sphinxcrossref{\DUrole{doc}{Python call}}}} page for an overview of using Python from LAMMPS in this
manner and all the {\hyperref[\detokenize{Python_head::doc}]{\sphinxcrossref{\DUrole{doc}{Python}}}} manual pages for other
ways to use LAMMPS and Python together.

\begin{sphinxadmonition}{note}{Note:}
\sphinxAtStartPar
Building with the PYTHON package assumes you have a Python development
environment (headers and libraries) available on your system, which needs
to be either Python version 2.7 or Python 3.5 and later.
\end{sphinxadmonition}

\sphinxAtStartPar
\sphinxstylestrong{Install:}

\sphinxAtStartPar
This package has {\hyperref[\detokenize{Build_extras:python}]{\sphinxcrossref{\DUrole{std,std-ref}{specific installation instructions}}}} on the {\hyperref[\detokenize{Build_extras::doc}]{\sphinxcrossref{\DUrole{doc}{Build extras}}}} page.

\sphinxAtStartPar
\sphinxstylestrong{Supporting info:}
\begin{itemize}
\item {} 
\sphinxAtStartPar
\sphinxcode{\sphinxupquote{src/PYTHON}}: filenames \sphinxhyphen{}\textgreater{} commands

\item {} 
\sphinxAtStartPar
{\hyperref[\detokenize{Python_head::doc}]{\sphinxcrossref{\DUrole{doc}{Python call}}}}

\item {} 
\sphinxAtStartPar
\sphinxcode{\sphinxupquote{lib/python/README}}

\item {} 
\sphinxAtStartPar
\sphinxcode{\sphinxupquote{examples/python}}

\end{itemize}


\bigskip\hrule\bigskip



\subsection{QEQ package}
\label{\detokenize{Packages_details:qeq-package}}\label{\detokenize{Packages_details:pkg-qeq}}
\sphinxAtStartPar
\sphinxstylestrong{Contents:}

\sphinxAtStartPar
Several fixes for performing charge equilibration (QEq) via different
algorithms.  These can be used with pair styles that perform QEq as
part of their formulation.

\sphinxAtStartPar
\sphinxstylestrong{Supporting info:}
\begin{itemize}
\item {} 
\sphinxAtStartPar
\sphinxcode{\sphinxupquote{src/QEQ}}: filenames \sphinxhyphen{}\textgreater{} commands

\item {} 
\sphinxAtStartPar
{\hyperref[\detokenize{fix_qeq::doc}]{\sphinxcrossref{\DUrole{doc}{fix qeq/*}}}}

\item {} 
\sphinxAtStartPar
\sphinxcode{\sphinxupquote{examples/qeq}}

\item {} 
\sphinxAtStartPar
\sphinxcode{\sphinxupquote{examples/streitz}}

\end{itemize}


\bigskip\hrule\bigskip



\subsection{QMMM package}
\label{\detokenize{Packages_details:qmmm-package}}\label{\detokenize{Packages_details:pkg-qmmm}}
\sphinxAtStartPar
\sphinxstylestrong{Contents:}

\sphinxAtStartPar
A {\hyperref[\detokenize{fix_qmmm::doc}]{\sphinxcrossref{\DUrole{doc}{fix qmmm}}}} command which allows LAMMPS to be used as
the MM code in a QM/MM simulation.  This is currently only available
in combination with the \sphinxhref{https://www.quantum-espresso.org}{Quantum ESPRESSO} package.

\sphinxAtStartPar
To use this package you must have Quantum ESPRESSO (QE) available on
your system and include its coupling library in the compilation and
then compile LAMMPS as a library.  For QM/MM calculations you then
build a custom binary with MPI support, that sets up 3 partitions with
MPI sub\sphinxhyphen{}communicators (for inter\sphinxhyphen{} and intra\sphinxhyphen{}partition communication)
and then calls the corresponding library interfaces on each partition
(2x LAMMPS and 1x QE).

\sphinxAtStartPar
The current implementation supports an ONIOM style mechanical coupling
and a multi\sphinxhyphen{}pole based electrostatic coupling to the Quantum ESPRESSO
plane wave DFT package.  The QM/MM interface has been written in a
manner that coupling to other QM codes should be possible without
changes to LAMMPS itself.

\sphinxAtStartPar
\sphinxstylestrong{Authors:} Axel Kohlmeyer (Temple U). Mariella Ippolito and Carlo Cavazzoni (CINECA, Italy)

\sphinxAtStartPar
\sphinxstylestrong{Install:}

\sphinxAtStartPar
This package has {\hyperref[\detokenize{Build_extras:qmmm}]{\sphinxcrossref{\DUrole{std,std-ref}{specific installation instructions}}}}
on the {\hyperref[\detokenize{Build_extras::doc}]{\sphinxcrossref{\DUrole{doc}{Build extras}}}} page.

\sphinxAtStartPar
\sphinxstylestrong{Supporting info:}
\begin{itemize}
\item {} 
\sphinxAtStartPar
\sphinxcode{\sphinxupquote{src/QMMM}}: filenames \sphinxhyphen{}\textgreater{} commands

\item {} 
\sphinxAtStartPar
\sphinxcode{\sphinxupquote{src/QMMM/README}}

\item {} 
\sphinxAtStartPar
\sphinxcode{\sphinxupquote{lib/qmmm/README}}

\item {} 
\sphinxAtStartPar
{\hyperref[\detokenize{fix_phonon::doc}]{\sphinxcrossref{\DUrole{doc}{fix phonon}}}}

\item {} 
\sphinxAtStartPar
\sphinxcode{\sphinxupquote{lib/qmmm/example\sphinxhyphen{}ec/README}}

\item {} 
\sphinxAtStartPar
\sphinxcode{\sphinxupquote{lib/qmmm/example\sphinxhyphen{}mc/README}}

\end{itemize}


\bigskip\hrule\bigskip



\subsection{QTB package}
\label{\detokenize{Packages_details:qtb-package}}\label{\detokenize{Packages_details:pkg-qtb}}
\sphinxAtStartPar
\sphinxstylestrong{Contents:}

\sphinxAtStartPar
Two fixes which provide a self\sphinxhyphen{}consistent quantum treatment of
vibrational modes in a classical molecular dynamics simulation.  By
coupling the MD simulation to a colored thermostat, it introduces zero
point energy into the system, altering the energy power spectrum and
the heat capacity to account for their quantum nature. This is useful
when modeling systems at temperatures lower than their classical
limits or when temperatures ramp across the classical limits in a
simulation.

\sphinxAtStartPar
\sphinxstylestrong{Author:} Yuan Shen (Stanford U).

\sphinxAtStartPar
\sphinxstylestrong{Supporting info:}
\begin{itemize}
\item {} 
\sphinxAtStartPar
\sphinxcode{\sphinxupquote{src/QTB}}: filenames \sphinxhyphen{}\textgreater{} commands

\item {} 
\sphinxAtStartPar
\sphinxcode{\sphinxupquote{src/QTB/README}}

\item {} 
\sphinxAtStartPar
{\hyperref[\detokenize{fix_qtb::doc}]{\sphinxcrossref{\DUrole{doc}{fix qtb}}}}

\item {} 
\sphinxAtStartPar
{\hyperref[\detokenize{fix_qbmsst::doc}]{\sphinxcrossref{\DUrole{doc}{fix qbmsst}}}}

\item {} 
\sphinxAtStartPar
\sphinxcode{\sphinxupquote{examples/PACKAGES/qtb}}

\end{itemize}


\bigskip\hrule\bigskip



\subsection{REACTION package}
\label{\detokenize{Packages_details:reaction-package}}\label{\detokenize{Packages_details:pkg-reaction}}
\sphinxAtStartPar
\sphinxstylestrong{Contents:}

\sphinxAtStartPar
This package implements the REACTER protocol, which allows for complex
bond topology changes (reactions) during a running MD simulation when
using classical force fields. Topology changes are defined in pre\sphinxhyphen{} and
post\sphinxhyphen{}reaction molecule templates and can include creation and deletion
of bonds, angles, dihedrals, impropers, atom types, bond types, angle
types, dihedral types, improper types, and/or atomic charges. Other
options currently available include reaction constraints (e.g., angle
and Arrhenius constraints), deletion of reaction byproducts or other
small molecules, creation of new atoms or molecules bonded to existing
atoms, and using LAMMPS variables for input parameters.

\sphinxAtStartPar
\sphinxstylestrong{Author:} Jacob R. Gissinger (NASA Langley Research Center).

\sphinxAtStartPar
\sphinxstylestrong{Supporting info:}
\begin{itemize}
\item {} 
\sphinxAtStartPar
\sphinxcode{\sphinxupquote{src/REACTION}}: filenames \sphinxhyphen{}\textgreater{} commands

\item {} 
\sphinxAtStartPar
\sphinxcode{\sphinxupquote{src/REACTION/README}}

\item {} 
\sphinxAtStartPar
{\hyperref[\detokenize{fix_bond_react::doc}]{\sphinxcrossref{\DUrole{doc}{fix bond/react}}}}

\item {} 
\sphinxAtStartPar
\sphinxcode{\sphinxupquote{examples/PACKAGES/reaction}}

\item {} 
\sphinxAtStartPar
\sphinxhref{https://www.lammps.org/workshops/Aug17/pdf/gissinger.pdf}{2017 LAMMPS Workshop}

\item {} 
\sphinxAtStartPar
\sphinxhref{https://www.lammps.org/workshops/Aug19/talk\_gissinger.pdf}{2019 LAMMPS Workshop}

\item {} 
\sphinxAtStartPar
\sphinxhref{https://www.lammps.org/workshops/Aug21/talk/jacob-gissinger/}{2021 LAMMPS Workshop}

\item {} 
\sphinxAtStartPar
\sphinxhref{https://www.reacter.org/}{REACTER website (reacter.org)}

\end{itemize}


\bigskip\hrule\bigskip



\subsection{REAXFF package}
\label{\detokenize{Packages_details:reaxff-package}}\label{\detokenize{Packages_details:pkg-reaxff}}
\sphinxAtStartPar
\sphinxstylestrong{Contents:}

\sphinxAtStartPar
A pair style which implements the ReaxFF potential in C/C++.  ReaxFF
is a universal reactive force field.  See the \sphinxcode{\sphinxupquote{src/REAXFF/README}} file
for more info on differences between the two packages.  Also two fixes
for monitoring molecules as bonds are created and destroyed.

\sphinxAtStartPar
\sphinxstylestrong{Author:} Hasan Metin Aktulga (MSU) while at Purdue University.

\sphinxAtStartPar
\sphinxstylestrong{Supporting info:}
\begin{itemize}
\item {} 
\sphinxAtStartPar
\sphinxcode{\sphinxupquote{src/REAXFF}}: filenames \sphinxhyphen{}\textgreater{} commands

\item {} 
\sphinxAtStartPar
\sphinxcode{\sphinxupquote{src/REAXFF/README}}

\item {} 
\sphinxAtStartPar
{\hyperref[\detokenize{pair_reaxff::doc}]{\sphinxcrossref{\DUrole{doc}{pair\_style reaxff}}}}

\item {} 
\sphinxAtStartPar
{\hyperref[\detokenize{fix_reaxff_bonds::doc}]{\sphinxcrossref{\DUrole{doc}{fix reaxff/bonds}}}}

\item {} 
\sphinxAtStartPar
{\hyperref[\detokenize{fix_reaxff_species::doc}]{\sphinxcrossref{\DUrole{doc}{fix reaxff/species}}}}

\item {} 
\sphinxAtStartPar
\sphinxcode{\sphinxupquote{examples/reaxff}}

\end{itemize}


\bigskip\hrule\bigskip



\subsection{REPLICA package}
\label{\detokenize{Packages_details:replica-package}}\label{\detokenize{Packages_details:pkg-replica}}
\sphinxAtStartPar
\sphinxstylestrong{Contents:}

\sphinxAtStartPar
A collection of multi\sphinxhyphen{}replica methods which can be used when running
multiple LAMMPS simulations (replicas).  See the {\hyperref[\detokenize{Howto_replica::doc}]{\sphinxcrossref{\DUrole{doc}{Howto replica}}}} page for an overview of how to run
multi\sphinxhyphen{}replica simulations in LAMMPS.  Methods in the package include
nudged elastic band (NEB), parallel replica dynamics (PRD),
temperature accelerated dynamics (TAD), parallel tempering, and a
verlet/split algorithm for performing long\sphinxhyphen{}range Coulombics on one set
of processors, and the remainder of the force field calculation on
another set.

\sphinxAtStartPar
\sphinxstylestrong{Supporting info:}
\begin{itemize}
\item {} 
\sphinxAtStartPar
\sphinxcode{\sphinxupquote{src/REPLICA}}: filenames \sphinxhyphen{}\textgreater{} commands

\item {} 
\sphinxAtStartPar
{\hyperref[\detokenize{Howto_replica::doc}]{\sphinxcrossref{\DUrole{doc}{Howto replica}}}}

\item {} 
\sphinxAtStartPar
{\hyperref[\detokenize{neb::doc}]{\sphinxcrossref{\DUrole{doc}{neb}}}}

\item {} 
\sphinxAtStartPar
{\hyperref[\detokenize{prd::doc}]{\sphinxcrossref{\DUrole{doc}{prd}}}}

\item {} 
\sphinxAtStartPar
{\hyperref[\detokenize{tad::doc}]{\sphinxcrossref{\DUrole{doc}{tad}}}}

\item {} 
\sphinxAtStartPar
{\hyperref[\detokenize{temper::doc}]{\sphinxcrossref{\DUrole{doc}{temper}}}},

\item {} 
\sphinxAtStartPar
{\hyperref[\detokenize{temper_npt::doc}]{\sphinxcrossref{\DUrole{doc}{temper/npt}}}},

\item {} 
\sphinxAtStartPar
{\hyperref[\detokenize{temper_grem::doc}]{\sphinxcrossref{\DUrole{doc}{temper/grem}}}},

\item {} 
\sphinxAtStartPar
{\hyperref[\detokenize{run_style::doc}]{\sphinxcrossref{\DUrole{doc}{run\_style verlet/split}}}}

\item {} 
\sphinxAtStartPar
\sphinxcode{\sphinxupquote{examples/neb}}

\item {} 
\sphinxAtStartPar
\sphinxcode{\sphinxupquote{examples/prd}}

\item {} 
\sphinxAtStartPar
\sphinxcode{\sphinxupquote{examples/tad}}

\item {} 
\sphinxAtStartPar
\sphinxcode{\sphinxupquote{examples/PACKAGES/grem}}

\end{itemize}


\bigskip\hrule\bigskip



\subsection{RHEO package}
\label{\detokenize{Packages_details:rheo-package}}\label{\detokenize{Packages_details:pkg-rheo}}
\sphinxAtStartPar
\sphinxstylestrong{Contents:}

\sphinxAtStartPar
Pair styles, bond styles, fixes, and computes for reproducing
hydrodynamics and elastic objects. See the {\hyperref[\detokenize{Howto_rheo::doc}]{\sphinxcrossref{\DUrole{doc}{Howto rheo}}}} page for an overview.

\sphinxAtStartPar
\sphinxstylestrong{Install:}

\sphinxAtStartPar
This package has {\hyperref[\detokenize{Build_extras:rheo}]{\sphinxcrossref{\DUrole{std,std-ref}{specific installation instructions}}}} on the {\hyperref[\detokenize{Build_extras::doc}]{\sphinxcrossref{\DUrole{doc}{Build extras}}}} page.

\sphinxAtStartPar
\sphinxstylestrong{Authors:} Joel T. Clemmer (Sandia National Labs),
Thomas C. O’Connor (Carnegie Mellon University)

\sphinxAtStartPar
\DUrole{versionmodified,added}{Added in version 29Aug2024.}

\sphinxAtStartPar
\sphinxstylestrong{Supporting info:}
\begin{itemize}
\item {} 
\sphinxAtStartPar
\sphinxcode{\sphinxupquote{src/RHEO}} filenames \sphinxhyphen{}\textgreater{} commands

\item {} 
\sphinxAtStartPar
{\hyperref[\detokenize{Howto_rheo::doc}]{\sphinxcrossref{\DUrole{doc}{Howto\_rheo}}}}

\item {} 
\sphinxAtStartPar
{\hyperref[\detokenize{atom_style::doc}]{\sphinxcrossref{\DUrole{doc}{atom\_style rheo}}}}

\item {} 
\sphinxAtStartPar
{\hyperref[\detokenize{atom_style::doc}]{\sphinxcrossref{\DUrole{doc}{atom\_style rheo/thermal}}}}

\item {} 
\sphinxAtStartPar
{\hyperref[\detokenize{bond_rheo_shell::doc}]{\sphinxcrossref{\DUrole{doc}{bond\_style rheo/shell}}}}

\item {} 
\sphinxAtStartPar
{\hyperref[\detokenize{compute_rheo_property_atom::doc}]{\sphinxcrossref{\DUrole{doc}{compute rheo/property/atom}}}}

\item {} 
\sphinxAtStartPar
{\hyperref[\detokenize{fix_rheo::doc}]{\sphinxcrossref{\DUrole{doc}{fix rheo}}}}

\item {} 
\sphinxAtStartPar
{\hyperref[\detokenize{fix_rheo_oxidation::doc}]{\sphinxcrossref{\DUrole{doc}{fix rheo/oxidation}}}}

\item {} 
\sphinxAtStartPar
{\hyperref[\detokenize{fix_rheo_pressure::doc}]{\sphinxcrossref{\DUrole{doc}{fix rheo/pressure}}}}

\item {} 
\sphinxAtStartPar
{\hyperref[\detokenize{fix_rheo_thermal::doc}]{\sphinxcrossref{\DUrole{doc}{fix rheo/thermal}}}}

\item {} 
\sphinxAtStartPar
{\hyperref[\detokenize{fix_rheo_viscosity::doc}]{\sphinxcrossref{\DUrole{doc}{fix rheo/viscosity}}}}

\item {} 
\sphinxAtStartPar
{\hyperref[\detokenize{pair_rheo::doc}]{\sphinxcrossref{\DUrole{doc}{pair\_style rheo}}}}

\item {} 
\sphinxAtStartPar
{\hyperref[\detokenize{pair_rheo_solid::doc}]{\sphinxcrossref{\DUrole{doc}{pair\_style rheo/solid}}}}

\item {} 
\sphinxAtStartPar
\sphinxurl{https://www.lammps.org/movies.html\#rheopackage}

\item {} 
\sphinxAtStartPar
\sphinxcode{\sphinxupquote{examples/rheo}}

\end{itemize}


\bigskip\hrule\bigskip



\subsection{RIGID package}
\label{\detokenize{Packages_details:rigid-package}}\label{\detokenize{Packages_details:pkg-rigid}}
\sphinxAtStartPar
\sphinxstylestrong{Contents:}

\sphinxAtStartPar
Fixes which enforce rigid constraints on collections of atoms or
particles.  This includes SHAKE and RATTLE, as well as various
rigid\sphinxhyphen{}body integrators for a few large bodies or many small bodies.
Also several computes which calculate properties of rigid bodies.

\sphinxAtStartPar
\sphinxstylestrong{Supporting info:}
\begin{itemize}
\item {} 
\sphinxAtStartPar
\sphinxcode{\sphinxupquote{src/RIGID}}: filenames \sphinxhyphen{}\textgreater{} commands

\item {} 
\sphinxAtStartPar
{\hyperref[\detokenize{compute_erotate_rigid::doc}]{\sphinxcrossref{\DUrole{doc}{compute erotate/rigid}}}}

\item {} 
\sphinxAtStartPar
{\hyperref[\detokenize{fix_shake::doc}]{\sphinxcrossref{\DUrole{doc}{fix shake}}}}

\item {} 
\sphinxAtStartPar
{\hyperref[\detokenize{fix_shake::doc}]{\sphinxcrossref{\DUrole{doc}{fix rattle}}}}

\item {} 
\sphinxAtStartPar
{\hyperref[\detokenize{fix_rigid::doc}]{\sphinxcrossref{\DUrole{doc}{fix rigid/*}}}}

\item {} 
\sphinxAtStartPar
\sphinxcode{\sphinxupquote{examples/ASPHERE}}

\item {} 
\sphinxAtStartPar
\sphinxcode{\sphinxupquote{examples/rigid}}

\item {} 
\sphinxAtStartPar
bench/in.rhodo

\item {} 
\sphinxAtStartPar
\sphinxurl{https://www.lammps.org/movies.html\#box}

\item {} 
\sphinxAtStartPar
\sphinxurl{https://www.lammps.org/movies.html\#star}

\end{itemize}


\bigskip\hrule\bigskip



\subsection{SCAFACOS package}
\label{\detokenize{Packages_details:scafacos-package}}\label{\detokenize{Packages_details:pkg-scafacos}}
\sphinxAtStartPar
\sphinxstylestrong{Contents:}

\sphinxAtStartPar
A KSpace style which wraps the \sphinxhref{http://www.scafacos.de}{ScaFaCoS Coulomb solver library} to compute long\sphinxhyphen{}range Coulombic
interactions.

\sphinxAtStartPar
To use this package you must have the ScaFaCoS library available on
your system.

\sphinxAtStartPar
\sphinxstylestrong{Author:} Rene Halver (JSC) wrote the scafacos LAMMPS command.

\sphinxAtStartPar
ScaFaCoS itself was developed by a consortium of German research
facilities with a BMBF (German Ministry of Science and Education)
funded project in 2009\sphinxhyphen{}2012. Participants of the consortium were the
Universities of Bonn, Chemnitz, Stuttgart, and Wuppertal as well as
the Forschungszentrum Juelich.

\sphinxAtStartPar
\sphinxstylestrong{Install:}

\sphinxAtStartPar
This package has {\hyperref[\detokenize{Build_extras:scafacos}]{\sphinxcrossref{\DUrole{std,std-ref}{specific installation instructions}}}} on the {\hyperref[\detokenize{Build_extras::doc}]{\sphinxcrossref{\DUrole{doc}{Build extras}}}} page.
The SCAFACOS package requires that LAMMPS is build in {\hyperref[\detokenize{Build_basics:serial}]{\sphinxcrossref{\DUrole{std,std-ref}{MPI parallel mode}}}}.

\sphinxAtStartPar
\sphinxstylestrong{Supporting info:}
\begin{itemize}
\item {} 
\sphinxAtStartPar
\sphinxcode{\sphinxupquote{src/SCAFACOS}}: filenames \sphinxhyphen{}\textgreater{} commands

\item {} 
\sphinxAtStartPar
\sphinxcode{\sphinxupquote{src/SCAFACOS/README}}

\item {} 
\sphinxAtStartPar
{\hyperref[\detokenize{kspace_style::doc}]{\sphinxcrossref{\DUrole{doc}{kspace\_style scafacos}}}}

\item {} 
\sphinxAtStartPar
{\hyperref[\detokenize{kspace_modify::doc}]{\sphinxcrossref{\DUrole{doc}{kspace\_modify}}}}

\item {} 
\sphinxAtStartPar
\sphinxcode{\sphinxupquote{examples/PACKAGES/scafacos}}

\end{itemize}


\bigskip\hrule\bigskip



\subsection{SHOCK package}
\label{\detokenize{Packages_details:shock-package}}\label{\detokenize{Packages_details:pkg-shock}}
\sphinxAtStartPar
\sphinxstylestrong{Contents:}

\sphinxAtStartPar
Fixes for running impact simulations where a shock\sphinxhyphen{}wave passes through
a material.

\sphinxAtStartPar
\sphinxstylestrong{Supporting info:}
\begin{itemize}
\item {} 
\sphinxAtStartPar
\sphinxcode{\sphinxupquote{src/SHOCK}}: filenames \sphinxhyphen{}\textgreater{} commands

\item {} 
\sphinxAtStartPar
{\hyperref[\detokenize{fix_append_atoms::doc}]{\sphinxcrossref{\DUrole{doc}{fix append/atoms}}}}

\item {} 
\sphinxAtStartPar
{\hyperref[\detokenize{fix_msst::doc}]{\sphinxcrossref{\DUrole{doc}{fix msst}}}}

\item {} 
\sphinxAtStartPar
{\hyperref[\detokenize{fix_nphug::doc}]{\sphinxcrossref{\DUrole{doc}{fix nphug}}}}

\item {} 
\sphinxAtStartPar
{\hyperref[\detokenize{fix_wall_piston::doc}]{\sphinxcrossref{\DUrole{doc}{fix wall/piston}}}}

\item {} 
\sphinxAtStartPar
\sphinxcode{\sphinxupquote{examples/hugoniostat}}

\item {} 
\sphinxAtStartPar
\sphinxcode{\sphinxupquote{examples/msst}}

\end{itemize}


\bigskip\hrule\bigskip



\subsection{SMTBQ package}
\label{\detokenize{Packages_details:smtbq-package}}\label{\detokenize{Packages_details:pkg-smtbq}}
\sphinxAtStartPar
\sphinxstylestrong{Contents:}

\sphinxAtStartPar
Pair styles which implement Second Moment Tight Binding models.
One with QEq charge equilibration (SMTBQ) for the description of
ionocovalent bonds in oxides, and two more as plain SMATB models.

\sphinxAtStartPar
\sphinxstylestrong{Authors:} SMTBQ: Nicolas Salles, Emile Maras, Olivier Politano, and Robert
Tetot (LAAS\sphinxhyphen{}CNRS, France);
SMATB: Daniele Rapetti (Politecnico di Torino)

\sphinxAtStartPar
\sphinxstylestrong{Supporting info:}
\begin{itemize}
\item {} 
\sphinxAtStartPar
\sphinxcode{\sphinxupquote{src/SMTBQ}}: filenames \sphinxhyphen{}\textgreater{} commands

\item {} 
\sphinxAtStartPar
\sphinxcode{\sphinxupquote{src/SMTBQ/README}}

\item {} 
\sphinxAtStartPar
{\hyperref[\detokenize{pair_smtbq::doc}]{\sphinxcrossref{\DUrole{doc}{pair\_style smtbq}}}}

\item {} 
\sphinxAtStartPar
{\hyperref[\detokenize{pair_smatb::doc}]{\sphinxcrossref{\DUrole{doc}{pair\_style smatb}}}}, {\hyperref[\detokenize{pair_smatb::doc}]{\sphinxcrossref{\DUrole{doc}{pair\_style smatb/single}}}}

\item {} 
\sphinxAtStartPar
\sphinxcode{\sphinxupquote{examples/PACKAGES/smtbq}}

\end{itemize}


\bigskip\hrule\bigskip



\subsection{SPH package}
\label{\detokenize{Packages_details:sph-package}}\label{\detokenize{Packages_details:pkg-sph}}
\sphinxAtStartPar
\sphinxstylestrong{Contents:}

\sphinxAtStartPar
An atom style, fixes, computes, and several pair styles which
implements smoothed particle hydrodynamics (SPH) for liquids.  See the
related {\hyperref[\detokenize{Packages_details:pkg-machdyn}]{\sphinxcrossref{\DUrole{std,std-ref}{MACHDYN package}}}} package for smooth Mach dynamics
(SMD) for solids.

\sphinxAtStartPar
This package contains ideal gas, Lennard\sphinxhyphen{}Jones equation of states,
Tait, and full support for complete (i.e. internal\sphinxhyphen{}energy dependent)
equations of state.  It allows for plain or Monaghans XSPH integration
of the equations of motion.  It has options for density continuity or
density summation to propagate the density field.  It has
{\hyperref[\detokenize{set::doc}]{\sphinxcrossref{\DUrole{doc}{set}}}} command options to set the internal energy and density
of particles from the input script and allows the same quantities to
be output with thermodynamic output or to dump files via the {\hyperref[\detokenize{compute_property_atom::doc}]{\sphinxcrossref{\DUrole{doc}{compute property/atom}}}} command.

\sphinxAtStartPar
\sphinxstylestrong{Author:} Georg Ganzenmuller (Fraunhofer\sphinxhyphen{}Institute for High\sphinxhyphen{}Speed
Dynamics, Ernst Mach Institute, Germany).

\sphinxAtStartPar
\sphinxstylestrong{Supporting info:}
\begin{itemize}
\item {} 
\sphinxAtStartPar
\sphinxcode{\sphinxupquote{src/SPH}}: filenames \sphinxhyphen{}\textgreater{} commands

\item {} 
\sphinxAtStartPar
\sphinxcode{\sphinxupquote{src/SPH/README}}

\item {} 
\sphinxAtStartPar
\sphinxhref{PDF/SPH\_LAMMPS\_userguide.pdf}{doc/PDF/SPH\_LAMMPS\_userguide.pdf}

\item {} 
\sphinxAtStartPar
\sphinxcode{\sphinxupquote{examples/PACKAGES/sph}}

\item {} 
\sphinxAtStartPar
\sphinxurl{https://www.lammps.org/movies.html\#sph}

\end{itemize}


\bigskip\hrule\bigskip



\subsection{SPIN package}
\label{\detokenize{Packages_details:spin-package}}\label{\detokenize{Packages_details:pkg-spin}}
\sphinxAtStartPar
\sphinxstylestrong{Contents:}

\sphinxAtStartPar
Model atomic magnetic spins classically, coupled to atoms moving in
the usual manner via MD.  Various pair, fix, and compute styles.

\sphinxAtStartPar
\sphinxstylestrong{Author:} Julien Tranchida (Sandia).

\sphinxAtStartPar
\sphinxstylestrong{Supporting info:}
\begin{itemize}
\item {} 
\sphinxAtStartPar
\sphinxcode{\sphinxupquote{src/SPIN}}: filenames \sphinxhyphen{}\textgreater{} commands

\item {} 
\sphinxAtStartPar
{\hyperref[\detokenize{Howto_spins::doc}]{\sphinxcrossref{\DUrole{doc}{Howto spins}}}}

\item {} 
\sphinxAtStartPar
{\hyperref[\detokenize{pair_spin_dipole::doc}]{\sphinxcrossref{\DUrole{doc}{pair\_style spin/dipole/cut}}}}

\item {} 
\sphinxAtStartPar
{\hyperref[\detokenize{pair_spin_dipole::doc}]{\sphinxcrossref{\DUrole{doc}{pair\_style spin/dipole/long}}}}

\item {} 
\sphinxAtStartPar
{\hyperref[\detokenize{pair_spin_dmi::doc}]{\sphinxcrossref{\DUrole{doc}{pair\_style spin/dmi}}}}

\item {} 
\sphinxAtStartPar
{\hyperref[\detokenize{pair_spin_exchange::doc}]{\sphinxcrossref{\DUrole{doc}{pair\_style spin/exchange}}}}

\item {} 
\sphinxAtStartPar
{\hyperref[\detokenize{pair_spin_exchange::doc}]{\sphinxcrossref{\DUrole{doc}{pair\_style spin/exchange/biquadratic}}}}

\item {} 
\sphinxAtStartPar
{\hyperref[\detokenize{pair_spin_magelec::doc}]{\sphinxcrossref{\DUrole{doc}{pair\_style spin/magelec}}}}

\item {} 
\sphinxAtStartPar
{\hyperref[\detokenize{pair_spin_neel::doc}]{\sphinxcrossref{\DUrole{doc}{pair\_style spin/neel}}}}

\item {} 
\sphinxAtStartPar
{\hyperref[\detokenize{fix_nve_spin::doc}]{\sphinxcrossref{\DUrole{doc}{fix nve/spin}}}}

\item {} 
\sphinxAtStartPar
{\hyperref[\detokenize{fix_langevin_spin::doc}]{\sphinxcrossref{\DUrole{doc}{fix langevin/spin}}}}

\item {} 
\sphinxAtStartPar
{\hyperref[\detokenize{fix_precession_spin::doc}]{\sphinxcrossref{\DUrole{doc}{fix precession/spin}}}}

\item {} 
\sphinxAtStartPar
{\hyperref[\detokenize{compute_spin::doc}]{\sphinxcrossref{\DUrole{doc}{compute spin}}}}

\item {} 
\sphinxAtStartPar
{\hyperref[\detokenize{neb_spin::doc}]{\sphinxcrossref{\DUrole{doc}{neb/spin}}}}

\item {} 
\sphinxAtStartPar
\sphinxcode{\sphinxupquote{examples/SPIN}}

\end{itemize}


\bigskip\hrule\bigskip



\subsection{SRD package}
\label{\detokenize{Packages_details:srd-package}}\label{\detokenize{Packages_details:pkg-srd}}
\sphinxAtStartPar
\sphinxstylestrong{Contents:}

\sphinxAtStartPar
A pair of fixes which implement the Stochastic Rotation Dynamics (SRD)
method for coarse\sphinxhyphen{}graining of a solvent, typically around large
colloidal particles.

\sphinxAtStartPar
\sphinxstylestrong{Supporting info:}
\begin{itemize}
\item {} 
\sphinxAtStartPar
\sphinxcode{\sphinxupquote{src/SRD}}: filenames \sphinxhyphen{}\textgreater{} commands

\item {} 
\sphinxAtStartPar
{\hyperref[\detokenize{fix_srd::doc}]{\sphinxcrossref{\DUrole{doc}{fix srd}}}}

\item {} 
\sphinxAtStartPar
{\hyperref[\detokenize{fix_wall_srd::doc}]{\sphinxcrossref{\DUrole{doc}{fix wall/srd}}}}

\item {} 
\sphinxAtStartPar
\sphinxcode{\sphinxupquote{examples/srd}}

\item {} 
\sphinxAtStartPar
\sphinxcode{\sphinxupquote{examples/ASPHERE}}

\item {} 
\sphinxAtStartPar
\sphinxurl{https://www.lammps.org/movies.html\#tri}

\item {} 
\sphinxAtStartPar
\sphinxurl{https://www.lammps.org/movies.html\#line}

\item {} 
\sphinxAtStartPar
\sphinxurl{https://www.lammps.org/movies.html\#poly}

\end{itemize}


\bigskip\hrule\bigskip



\subsection{TALLY package}
\label{\detokenize{Packages_details:tally-package}}\label{\detokenize{Packages_details:pkg-tally}}
\sphinxAtStartPar
\sphinxstylestrong{Contents:}

\sphinxAtStartPar
Several compute styles that can be called when pairwise interactions
are calculated to tally information (forces, heat flux, energy,
stress, etc) about individual interactions.

\sphinxAtStartPar
\sphinxstylestrong{Author:} Axel Kohlmeyer (Temple U).

\sphinxAtStartPar
\sphinxstylestrong{Supporting info:}
\begin{itemize}
\item {} 
\sphinxAtStartPar
\sphinxcode{\sphinxupquote{src/TALLY}}: filenames \sphinxhyphen{}\textgreater{} commands

\item {} 
\sphinxAtStartPar
\sphinxcode{\sphinxupquote{src/TALLY/README}}

\item {} 
\sphinxAtStartPar
{\hyperref[\detokenize{compute_tally::doc}]{\sphinxcrossref{\DUrole{doc}{compute */tally}}}}

\item {} 
\sphinxAtStartPar
\sphinxcode{\sphinxupquote{examples/PACKAGES/tally}}

\end{itemize}


\bigskip\hrule\bigskip



\subsection{UEF package}
\label{\detokenize{Packages_details:uef-package}}\label{\detokenize{Packages_details:pkg-uef}}
\sphinxAtStartPar
\sphinxstylestrong{Contents:}

\sphinxAtStartPar
A fix style for the integration of the equations of motion under
extensional flow with proper boundary conditions, as well as several
supporting compute styles and an output option.

\sphinxAtStartPar
\sphinxstylestrong{Author:} David Nicholson (MIT).

\sphinxAtStartPar
\sphinxstylestrong{Supporting info:}
\begin{itemize}
\item {} 
\sphinxAtStartPar
\sphinxcode{\sphinxupquote{src/UEF}}: filenames \sphinxhyphen{}\textgreater{} commands

\item {} 
\sphinxAtStartPar
\sphinxcode{\sphinxupquote{src/UEF/README}}

\item {} 
\sphinxAtStartPar
{\hyperref[\detokenize{fix_nh_uef::doc}]{\sphinxcrossref{\DUrole{doc}{fix nvt/uef}}}}

\item {} 
\sphinxAtStartPar
{\hyperref[\detokenize{fix_nh_uef::doc}]{\sphinxcrossref{\DUrole{doc}{fix npt/uef}}}}

\item {} 
\sphinxAtStartPar
{\hyperref[\detokenize{compute_pressure_uef::doc}]{\sphinxcrossref{\DUrole{doc}{compute pressure/uef}}}}

\item {} 
\sphinxAtStartPar
{\hyperref[\detokenize{compute_temp_uef::doc}]{\sphinxcrossref{\DUrole{doc}{compute temp/uef}}}}

\item {} 
\sphinxAtStartPar
{\hyperref[\detokenize{dump_cfg_uef::doc}]{\sphinxcrossref{\DUrole{doc}{dump cfg/uef}}}}

\item {} 
\sphinxAtStartPar
\sphinxcode{\sphinxupquote{examples/uef}}

\end{itemize}


\bigskip\hrule\bigskip



\subsection{VORONOI package}
\label{\detokenize{Packages_details:voronoi-package}}\label{\detokenize{Packages_details:pkg-voronoi}}
\sphinxAtStartPar
\sphinxstylestrong{Contents:}

\sphinxAtStartPar
A compute command which calculates the Voronoi tesselation of a
collection of atoms by wrapping the \sphinxhref{https://math.lbl.gov/voro++/}{Voro++ library}.  This can be used to calculate the
local volume or each atoms or its near neighbors.

\sphinxAtStartPar
To use this package you must have the Voro++ library available on your
system.

\sphinxAtStartPar
\sphinxstylestrong{Author:} Daniel Schwen (INL) while at LANL.  The open\sphinxhyphen{}source Voro++
library was written by Chris Rycroft (Harvard U) while at UC Berkeley
and LBNL.

\sphinxAtStartPar
\sphinxstylestrong{Install:}

\sphinxAtStartPar
This package has {\hyperref[\detokenize{Build_extras:voronoi}]{\sphinxcrossref{\DUrole{std,std-ref}{specific installation instructions}}}} on the {\hyperref[\detokenize{Build_extras::doc}]{\sphinxcrossref{\DUrole{doc}{Build extras}}}} page.

\sphinxAtStartPar
\sphinxstylestrong{Supporting info:}
\begin{itemize}
\item {} 
\sphinxAtStartPar
\sphinxcode{\sphinxupquote{src/VORONOI}}: filenames \sphinxhyphen{}\textgreater{} commands

\item {} 
\sphinxAtStartPar
\sphinxcode{\sphinxupquote{src/VORONOI/README}}

\item {} 
\sphinxAtStartPar
\sphinxcode{\sphinxupquote{lib/voronoi/README}}

\item {} 
\sphinxAtStartPar
{\hyperref[\detokenize{compute_voronoi_atom::doc}]{\sphinxcrossref{\DUrole{doc}{compute voronoi/atom}}}}

\item {} 
\sphinxAtStartPar
\sphinxcode{\sphinxupquote{examples/voronoi}}

\end{itemize}


\bigskip\hrule\bigskip



\subsection{VTK package}
\label{\detokenize{Packages_details:vtk-package}}\label{\detokenize{Packages_details:pkg-vtk}}
\sphinxAtStartPar
\sphinxstylestrong{Contents:}

\sphinxAtStartPar
A {\hyperref[\detokenize{dump_vtk::doc}]{\sphinxcrossref{\DUrole{doc}{dump vtk}}}} command which outputs snapshot info in the
\sphinxhref{https://www.vtk.org}{VTK format}, enabling visualization by \sphinxhref{https://www.paraview.org}{Paraview} or
other visualization packages.

\sphinxAtStartPar
To use this package you must have VTK library available on your
system.

\sphinxAtStartPar
\sphinxstylestrong{Authors:} Richard Berger (JKU) and Daniel Queteschiner (DCS Computing).

\sphinxAtStartPar
\sphinxstylestrong{Install:}

\sphinxAtStartPar
This package has {\hyperref[\detokenize{Build_extras:vtk}]{\sphinxcrossref{\DUrole{std,std-ref}{specific installation instructions}}}} on the {\hyperref[\detokenize{Build_extras::doc}]{\sphinxcrossref{\DUrole{doc}{Build extras}}}} page.

\sphinxAtStartPar
\sphinxstylestrong{Supporting info:}
\begin{itemize}
\item {} 
\sphinxAtStartPar
\sphinxcode{\sphinxupquote{src/VTK}}: filenames \sphinxhyphen{}\textgreater{} commands

\item {} 
\sphinxAtStartPar
\sphinxcode{\sphinxupquote{src/VTK/README}}

\item {} 
\sphinxAtStartPar
\sphinxcode{\sphinxupquote{lib/vtk/README}}

\item {} 
\sphinxAtStartPar
{\hyperref[\detokenize{dump_vtk::doc}]{\sphinxcrossref{\DUrole{doc}{dump vtk}}}}

\end{itemize}


\bigskip\hrule\bigskip



\subsection{YAFF package}
\label{\detokenize{Packages_details:yaff-package}}\label{\detokenize{Packages_details:pkg-yaff}}
\sphinxAtStartPar
\sphinxstylestrong{Contents:}

\sphinxAtStartPar
Some potentials that are also implemented in the Yet Another Force Field (\sphinxhref{https://github.com/molmod/yaff}{YAFF}) code.
The expressions and their use are discussed in the following papers
\begin{itemize}
\item {} 
\sphinxAtStartPar
Vanduyfhuys et al., J. Comput. Chem., 36 (13), 1015\sphinxhyphen{}1027 (2015) \sphinxhref{https://doi.org/10.1002/jcc.23877}{link}

\item {} 
\sphinxAtStartPar
Vanduyfhuys et al., J. Comput. Chem., 39 (16), 999\sphinxhyphen{}1011 (2018) \sphinxhref{https://doi.org/10.1002/jcc.25173}{link}

\end{itemize}

\sphinxAtStartPar
which discuss the \sphinxhref{https://molmod.github.io/QuickFF}{QuickFF} methodology.

\sphinxAtStartPar
\sphinxstylestrong{Author:} Steven Vandenbrande.

\sphinxAtStartPar
\DUrole{versionmodified,added}{Added in version 1Feb2019.}

\sphinxAtStartPar
\sphinxstylestrong{Supporting info:}
\begin{itemize}
\item {} 
\sphinxAtStartPar
\sphinxcode{\sphinxupquote{src/YAFF/README}}

\item {} 
\sphinxAtStartPar
{\hyperref[\detokenize{angle_cross::doc}]{\sphinxcrossref{\DUrole{doc}{angle\_style cross}}}}

\item {} 
\sphinxAtStartPar
{\hyperref[\detokenize{angle_mm3::doc}]{\sphinxcrossref{\DUrole{doc}{angle\_style mm3}}}}

\item {} 
\sphinxAtStartPar
{\hyperref[\detokenize{bond_mm3::doc}]{\sphinxcrossref{\DUrole{doc}{bond\_style mm3}}}}

\item {} 
\sphinxAtStartPar
{\hyperref[\detokenize{improper_distharm::doc}]{\sphinxcrossref{\DUrole{doc}{improper\_style distharm}}}}

\item {} 
\sphinxAtStartPar
{\hyperref[\detokenize{improper_sqdistharm::doc}]{\sphinxcrossref{\DUrole{doc}{improper\_style sqdistharm}}}}

\item {} 
\sphinxAtStartPar
{\hyperref[\detokenize{pair_lj_switch3_coulgauss_long::doc}]{\sphinxcrossref{\DUrole{doc}{pair\_style mm3/switch3/coulgauss/long}}}}

\item {} 
\sphinxAtStartPar
{\hyperref[\detokenize{pair_lj_switch3_coulgauss_long::doc}]{\sphinxcrossref{\DUrole{doc}{pair\_style lj/switch3/coulgauss/long}}}}

\item {} 
\sphinxAtStartPar
\sphinxcode{\sphinxupquote{examples/PACKAGES/yaff}}

\end{itemize}

\sphinxstepscope


\chapter{Accelerate performance}
\label{\detokenize{Speed:accelerate-performance}}\label{\detokenize{Speed::doc}}
\sphinxAtStartPar
This section describes various methods for improving LAMMPS
performance for different classes of problems running on different
kinds of machines.

\sphinxAtStartPar
There are two thrusts to the discussion that follows.  The first is
using code options that implement alternate algorithms that can
speed\sphinxhyphen{}up a simulation.  The second is to use one of the several
accelerator packages provided with LAMMPS that contain code optimized
for certain kinds of hardware, including multicore CPUs, GPUs, and
Intel Xeon Phi co\sphinxhyphen{}processors.

\sphinxAtStartPar
The \sphinxhref{https://www.lammps.org/bench.html}{Benchmark page} of the LAMMPS
website gives performance results for the various accelerator
packages discussed on the {\hyperref[\detokenize{Speed_packages::doc}]{\sphinxcrossref{\DUrole{doc}{Accelerator packages}}}}
page, for several of the standard LAMMPS benchmark problems, as a
function of problem size and number of compute nodes, on different
hardware platforms.

\sphinxstepscope


\section{Benchmarks}
\label{\detokenize{Speed_bench:benchmarks}}\label{\detokenize{Speed_bench::doc}}
\sphinxAtStartPar
Current LAMMPS performance is discussed on the \sphinxhref{https://www.lammps.org/bench.html}{Benchmarks page} of the \sphinxhref{https://www.lammps.org}{LAMMPS website}
where timings and parallel efficiency are listed.  The page has
several sections, which are briefly described below:
\begin{itemize}
\item {} 
\sphinxAtStartPar
CPU performance on 5 standard problems, strong and weak scaling

\item {} 
\sphinxAtStartPar
GPU and Xeon Phi performance on same and related problems

\item {} 
\sphinxAtStartPar
Comparison of cost of interatomic potentials

\item {} 
\sphinxAtStartPar
Performance of huge, billion\sphinxhyphen{}atom problems

\end{itemize}

\sphinxAtStartPar
The 5 standard problems are as follow:
\begin{enumerate}
\sphinxsetlistlabels{\arabic}{enumi}{enumii}{}{.}%
\item {} 
\sphinxAtStartPar
LJ = atomic fluid, Lennard\sphinxhyphen{}Jones potential with 2.5 sigma cutoff (55
neighbors per atom), NVE integration

\item {} 
\sphinxAtStartPar
Chain = bead\sphinxhyphen{}spring polymer melt of 100\sphinxhyphen{}mer chains, FENE bonds and LJ
pairwise interactions with a \(2^{\frac{1}{6}}\) sigma cutoff (5 neighbors per
atom), NVE integration

\item {} 
\sphinxAtStartPar
EAM = metallic solid, Cu EAM potential with 4.95 Angstrom cutoff (45
neighbors per atom), NVE integration

\item {} 
\sphinxAtStartPar
Chute = granular chute flow, frictional history potential with 1.1
sigma cutoff (7 neighbors per atom), NVE integration

\item {} 
\sphinxAtStartPar
Rhodo = rhodopsin protein in solvated lipid bilayer, CHARMM force
field with a 10 Angstrom LJ cutoff (440 neighbors per atom),
particle\sphinxhyphen{}particle particle\sphinxhyphen{}mesh (PPPM) for long\sphinxhyphen{}range Coulombics, NPT
integration

\end{enumerate}

\sphinxAtStartPar
Input files for these 5 problems are provided in the bench directory
of the LAMMPS distribution.  Each has 32,000 atoms and runs for 100
timesteps.  The size of the problem (number of atoms) can be varied
using command\sphinxhyphen{}line switches as described in the \sphinxcode{\sphinxupquote{bench/README}} file.
This is an easy way to test performance and either strong or weak
scalability on your machine.

\sphinxAtStartPar
The bench directory includes a few \sphinxcode{\sphinxupquote{log.*}} files that show performance
of these 5 problems on 1 or 4 cores of Linux desktop.  The \sphinxcode{\sphinxupquote{bench/FERMI}}
and \sphinxcode{\sphinxupquote{bench/KEPLER}} directories have input files and scripts and instructions
for running the same (or similar) problems using OpenMP or GPU or Xeon
Phi acceleration options.  See the \sphinxcode{\sphinxupquote{README}} files in those directories and the
{\hyperref[\detokenize{Speed_packages::doc}]{\sphinxcrossref{\DUrole{doc}{Accelerator packages}}}} pages for instructions on how
to build LAMMPS and run on that kind of hardware.

\sphinxAtStartPar
The \sphinxcode{\sphinxupquote{bench/POTENTIALS}} directory has input files which correspond to the
table of results on the
\sphinxhref{https://www.lammps.org/bench.html\#potentials}{Potentials} section of
the Benchmarks web page.  So you can also run those test problems on
your machine.

\sphinxAtStartPar
The \sphinxhref{https://www.lammps.org/bench.html\#billion}{billion\sphinxhyphen{}atom} section
of the Benchmarks web page has performance data for very large
benchmark runs of simple Lennard\sphinxhyphen{}Jones (LJ) models, which use the
\sphinxcode{\sphinxupquote{bench/in.lj}} input script.


\bigskip\hrule\bigskip


\sphinxAtStartPar
For all the benchmarks, a useful metric is the CPU cost per atom per
timestep.  Since performance scales roughly linearly with problem size
and timesteps for all LAMMPS models (i.e. interatomic or coarse\sphinxhyphen{}grained
potentials), the run time of any problem using the same model (atom
style, force field, cutoff, etc) can then be estimated.

\sphinxAtStartPar
Performance on a parallel machine can also be predicted from one\sphinxhyphen{}core
or one\sphinxhyphen{}node timings if the parallel efficiency can be estimated.  The
communication bandwidth and latency of a particular parallel machine
affects the efficiency.  On most machines LAMMPS will give a parallel
efficiency on these benchmarks above 50\% so long as the number of
atoms/core is a few 100 or greater, and closer to 100\% for large
numbers of atoms/core.  This is for all\sphinxhyphen{}MPI mode with one MPI task per
core.  For nodes with accelerator options or hardware (OpenMP, GPU,
Phi), you should first measure single node performance.  Then you can
estimate parallel performance for multi\sphinxhyphen{}node runs using the same logic
as for all\sphinxhyphen{}MPI mode, except that now you will typically need many more
atoms/node to achieve good scalability.

\sphinxstepscope


\section{Measuring performance}
\label{\detokenize{Speed_measure:measuring-performance}}\label{\detokenize{Speed_measure::doc}}
\sphinxAtStartPar
Before trying to make your simulation run faster, you should
understand how it currently performs and where the bottlenecks are.

\sphinxAtStartPar
The best way to do this is run the your system (actual number of
atoms) for a modest number of timesteps (say 100 steps) on several
different processor counts, including a single processor if possible.
Do this for an equilibrium version of your system, so that the
100\sphinxhyphen{}step timings are representative of a much longer run.  There is
typically no need to run for 1000s of timesteps to get accurate
timings; you can simply extrapolate from short runs.

\sphinxAtStartPar
For the set of runs, look at the timing data printed to the screen and
log file at the end of each LAMMPS run.  The
{\hyperref[\detokenize{Run_output::doc}]{\sphinxcrossref{\DUrole{doc}{screen and logfile output}}}} page gives an overview.

\sphinxAtStartPar
Running on one (or a few processors) should give a good estimate of
the serial performance and what portions of the timestep are taking
the most time.  Running the same problem on a few different processor
counts should give an estimate of parallel scalability.  I.e. if the
simulation runs 16x faster on 16 processors, its 100\% parallel
efficient; if it runs 8x faster on 16 processors, it’s 50\% efficient.

\sphinxAtStartPar
The most important data to look at in the timing info is the timing
breakdown and relative percentages.  For example, trying different
options for speeding up the long\sphinxhyphen{}range solvers will have little impact
if they only consume 10\% of the run time.  If the pairwise time is
dominating, you may want to look at GPU or OMP versions of the pair
style, as discussed below.  Comparing how the percentages change as
you increase the processor count gives you a sense of how different
operations within the timestep are scaling.  Note that if you are
running with a Kspace solver, there is additional output on the
breakdown of the Kspace time.  For PPPM, this includes the fraction
spent on FFTs, which can be communication intensive.

\sphinxAtStartPar
Another important detail in the timing info are the histograms of
atoms counts and neighbor counts.  If these vary widely across
processors, you have a load\sphinxhyphen{}imbalance issue.  This often results in
inaccurate relative timing data, because processors have to wait when
communication occurs for other processors to catch up.  Thus the
reported times for “Communication” or “Other” may be higher than they
really are, due to load\sphinxhyphen{}imbalance.  If this is an issue, you can
use the {\hyperref[\detokenize{timer::doc}]{\sphinxcrossref{\DUrole{doc}{timer sync}}}} command to obtain synchronized timings.

\sphinxstepscope


\section{General tips}
\label{\detokenize{Speed_tips:general-tips}}\label{\detokenize{Speed_tips::doc}}
\begin{sphinxadmonition}{note}{Note:}
\sphinxAtStartPar
this page is still a work in progress
\end{sphinxadmonition}

\sphinxAtStartPar
Here is a list of general ideas for improving simulation performance.
Most of them are only applicable to certain models and certain
bottlenecks in the current performance, so let the timing data you
generate be your guide.  It is hard, if not impossible, to predict how
much difference these options will make, since it is a function of
problem size, number of processors used, and your machine.  There is
no substitute for identifying performance bottlenecks, and trying out
various options.
\begin{itemize}
\item {} 
\sphinxAtStartPar
rRESPA

\item {} 
\sphinxAtStartPar
Two\sphinxhyphen{}FFT PPPM

\item {} 
\sphinxAtStartPar
Staggered PPPM

\item {} 
\sphinxAtStartPar
single vs double PPPM

\item {} 
\sphinxAtStartPar
partial charge PPPM

\item {} 
\sphinxAtStartPar
verlet/split run style

\item {} 
\sphinxAtStartPar
processor command for proc layout and numa layout

\item {} 
\sphinxAtStartPar
load\sphinxhyphen{}balancing: balance and fix balance

\end{itemize}

\sphinxAtStartPar
Two\sphinxhyphen{}FFT PPPM, also called \sphinxstyleemphasis{analytic differentiation} or \sphinxstyleemphasis{ad} PPPM,
uses 2 FFTs instead of the 4 FFTs used by the default \sphinxstyleemphasis{ik
differentiation} PPPM. However, 2\sphinxhyphen{}FFT PPPM also requires a slightly
larger mesh size to achieve the same accuracy as 4\sphinxhyphen{}FFT PPPM. For
problems where the FFT cost is the performance bottleneck (typically
large problems running on many processors), 2\sphinxhyphen{}FFT PPPM may be faster
than 4\sphinxhyphen{}FFT PPPM.

\sphinxAtStartPar
Staggered PPPM performs calculations using two different meshes, one
shifted slightly with respect to the other.  This can reduce force
aliasing errors and increase the accuracy of the method, but also
doubles the amount of work required. For high relative accuracy, using
staggered PPPM allows one to half the mesh size in each dimension as
compared to regular PPPM, which can give around a 4x speedup in the
kspace time. However, for low relative accuracy, using staggered PPPM
gives little benefit and can be up to 2x slower in the kspace
time. For example, the rhodopsin benchmark was run on a single
processor, and results for kspace time vs. relative accuracy for the
different methods are shown in the figure below.  For this system,
staggered PPPM (using ik differentiation) becomes useful when using a
relative accuracy of slightly greater than 1e\sphinxhyphen{}5 and above.

\noindent{\hspace*{\fill}\sphinxincludegraphics{{rhodo_staggered}.jpg}\hspace*{\fill}}

\begin{sphinxadmonition}{note}{Note:}
\sphinxAtStartPar
Using staggered PPPM may not give the same increase in accuracy
of energy and pressure as it does in forces, so some caution must be
used if energy and/or pressure are quantities of interest, such as
when using a barostat.
\end{sphinxadmonition}

\sphinxstepscope


\section{Accelerator packages}
\label{\detokenize{Speed_packages:accelerator-packages}}\label{\detokenize{Speed_packages::doc}}
\sphinxAtStartPar
Accelerated versions of various {\hyperref[\detokenize{pair_style::doc}]{\sphinxcrossref{\DUrole{doc}{pair\_style}}}},
{\hyperref[\detokenize{fix::doc}]{\sphinxcrossref{\DUrole{doc}{fixes}}}}, {\hyperref[\detokenize{compute::doc}]{\sphinxcrossref{\DUrole{doc}{computes}}}}, and other commands have
been added to LAMMPS, which will typically run faster than the
standard non\sphinxhyphen{}accelerated versions.  Some require appropriate hardware
to be present on your system, e.g. GPUs or Intel Xeon Phi
co\sphinxhyphen{}processors.

\sphinxAtStartPar
All of these commands are in packages provided with LAMMPS.  An
overview of packages is give on the {\hyperref[\detokenize{Packages::doc}]{\sphinxcrossref{\DUrole{doc}{Packages}}}} doc
pages.

\sphinxAtStartPar
These are the accelerator packages currently in LAMMPS:


\begin{savenotes}\sphinxattablestart
\sphinxthistablewithglobalstyle
\centering
\begin{tabulary}{\linewidth}[t]{TT}
\sphinxtoprule
\sphinxtableatstartofbodyhook
\sphinxAtStartPar
{\hyperref[\detokenize{Speed_gpu::doc}]{\sphinxcrossref{\DUrole{doc}{GPU Package}}}}
&
\sphinxAtStartPar
for GPUs via CUDA, OpenCL, or ROCm HIP
\\
\sphinxhline
\sphinxAtStartPar
{\hyperref[\detokenize{Speed_intel::doc}]{\sphinxcrossref{\DUrole{doc}{INTEL Package}}}}
&
\sphinxAtStartPar
for Intel CPUs and Intel Xeon Phi
\\
\sphinxhline
\sphinxAtStartPar
{\hyperref[\detokenize{Speed_kokkos::doc}]{\sphinxcrossref{\DUrole{doc}{KOKKOS Package}}}}
&
\sphinxAtStartPar
for NVIDIA GPUs, Intel Xeon Phi, and OpenMP threading
\\
\sphinxhline
\sphinxAtStartPar
{\hyperref[\detokenize{Speed_omp::doc}]{\sphinxcrossref{\DUrole{doc}{OPENMP Package}}}}
&
\sphinxAtStartPar
for OpenMP threading and generic CPU optimizations
\\
\sphinxhline
\sphinxAtStartPar
{\hyperref[\detokenize{Speed_opt::doc}]{\sphinxcrossref{\DUrole{doc}{OPT Package}}}}
&
\sphinxAtStartPar
generic CPU optimizations
\\
\sphinxbottomrule
\end{tabulary}
\sphinxtableafterendhook\par
\sphinxattableend\end{savenotes}

\sphinxstepscope


\subsection{GPU package}
\label{\detokenize{Speed_gpu:gpu-package}}\label{\detokenize{Speed_gpu::doc}}
\sphinxAtStartPar
The GPU package was developed by Mike Brown while at SNL and ORNL (now
at Intel Corp.) and his collaborators, particularly Trung Nguyen (now at
Northwestern).  Support for AMD GPUs via HIP was added by Vsevolod Nikolskiy
and coworkers at HSE University.

\sphinxAtStartPar
The GPU package provides GPU versions of many pair styles and for
parts of the {\hyperref[\detokenize{kspace_style::doc}]{\sphinxcrossref{\DUrole{doc}{kspace\_style pppm}}}} for long\sphinxhyphen{}range
Coulombics.  It has the following general features:
\begin{itemize}
\item {} 
\sphinxAtStartPar
It is designed to exploit common GPU hardware configurations where one
or more GPUs are coupled to many cores of one or more multicore CPUs,
e.g. within a node of a parallel machine.

\item {} 
\sphinxAtStartPar
Atom\sphinxhyphen{}based data (e.g. coordinates, forces) are moved back\sphinxhyphen{}and\sphinxhyphen{}forth
between the CPU(s) and GPU every timestep.

\item {} 
\sphinxAtStartPar
Neighbor lists can be built on the CPU or on the GPU

\item {} 
\sphinxAtStartPar
The charge assignment and force interpolation portions of PPPM can be
run on the GPU.  The FFT portion, which requires MPI communication
between processors, runs on the CPU.

\item {} 
\sphinxAtStartPar
Force computations of different style (pair vs. bond/angle/dihedral/improper)
can be performed concurrently on the GPU and CPU(s), respectively.

\item {} 
\sphinxAtStartPar
It allows for GPU computations to be performed in single or double
precision, or in mixed\sphinxhyphen{}mode precision, where pairwise forces are
computed in single precision, but accumulated into double\sphinxhyphen{}precision
force vectors.

\item {} 
\sphinxAtStartPar
LAMMPS\sphinxhyphen{}specific code is in the GPU package.  It makes calls to a
generic GPU library in the lib/gpu directory.  This library provides
either Nvidia support, AMD support, or more general OpenCL support
(for Nvidia GPUs, AMD GPUs, Intel GPUs, and multicore CPUs).
so that the same functionality is supported on a variety of hardware.

\end{itemize}

\sphinxAtStartPar
\sphinxstylestrong{Required hardware/software:}

\sphinxAtStartPar
To compile and use this package in CUDA mode, you currently need
to have an NVIDIA GPU and install the corresponding NVIDIA CUDA
toolkit software on your system (this is only tested on Linux
and unsupported on Windows):
\begin{itemize}
\item {} 
\sphinxAtStartPar
Check if you have an NVIDIA GPU: \sphinxcode{\sphinxupquote{cat /proc/driver/nvidia/gpus/\textbackslash{}*/information}}

\item {} 
\sphinxAtStartPar
Go to \sphinxurl{https://developer.nvidia.com/cuda-downloads}

\item {} 
\sphinxAtStartPar
Install a driver and toolkit appropriate for your system (SDK is not necessary)

\item {} 
\sphinxAtStartPar
Run \sphinxcode{\sphinxupquote{lammps/lib/gpu/nvc\_get\_devices}} (after building the GPU library, see below) to
list supported devices and properties

\end{itemize}

\sphinxAtStartPar
To compile and use this package in OpenCL mode, you currently need
to have the OpenCL headers and the (vendor neutral) OpenCL library installed.
In OpenCL mode, the acceleration depends on having an \sphinxhref{https://www.khronos.org/news/permalink/opencl-installable-client-driver-icd-loader}{OpenCL Installable Client Driver (ICD)}
installed. There can be multiple of them for the same or different hardware
(GPUs, CPUs, Accelerators) installed at the same time. OpenCL refers to those
as ‘platforms’.  The GPU library will try to auto\sphinxhyphen{}select the best suitable platform,
but this can be overridden using the platform option of the {\hyperref[\detokenize{package::doc}]{\sphinxcrossref{\DUrole{doc}{package}}}}
command. run \sphinxcode{\sphinxupquote{lammps/lib/gpu/ocl\_get\_devices}} to get a list of available
platforms and devices with a suitable ICD available.

\sphinxAtStartPar
To compile and use this package for Intel GPUs, OpenCL or the Intel oneAPI
HPC Toolkit can be installed using linux package managers. The latter also
provides optimized C++, MPI, and many other libraries and tools. See:
\begin{itemize}
\item {} 
\sphinxAtStartPar
\sphinxurl{https://software.intel.com/content/www/us/en/develop/tools/oneapi/hpc-toolkit/download.html}

\end{itemize}

\sphinxAtStartPar
If you do not have a discrete GPU card installed, this package can still provide
significant speedups on some CPUs that include integrated GPUs. Additionally, for
many macs, OpenCL is already included with the OS and Makefiles are available
in the \sphinxcode{\sphinxupquote{lib/gpu}} directory.

\sphinxAtStartPar
To compile and use this package in HIP mode, you have to have the AMD ROCm
software installed. Versions of ROCm older than 3.5 are currently deprecated
by AMD.

\sphinxAtStartPar
\sphinxstylestrong{Building LAMMPS with the GPU package:}

\sphinxAtStartPar
See the {\hyperref[\detokenize{Build_extras:gpu}]{\sphinxcrossref{\DUrole{std,std-ref}{Build extras}}}} page for
instructions.

\sphinxAtStartPar
\sphinxstylestrong{Run with the GPU package from the command line:}

\sphinxAtStartPar
The \sphinxcode{\sphinxupquote{mpirun}} or \sphinxcode{\sphinxupquote{mpiexec}} command sets the total number of MPI tasks
used by LAMMPS (one or multiple per compute node) and the number of MPI
tasks used per node.  E.g. the \sphinxcode{\sphinxupquote{mpirun}} command in MPICH does this via
its \sphinxcode{\sphinxupquote{\sphinxhyphen{}np}} and \sphinxcode{\sphinxupquote{\sphinxhyphen{}ppn}} switches.  Ditto for OpenMPI via \sphinxcode{\sphinxupquote{\sphinxhyphen{}np}} and
\sphinxcode{\sphinxupquote{\sphinxhyphen{}npernode}}.

\sphinxAtStartPar
When using the GPU package, you cannot assign more than one GPU to a
single MPI task.  However multiple MPI tasks can share the same GPU,
and in many cases it will be more efficient to run this way.  Likewise
it may be more efficient to use less MPI tasks/node than the available
\# of CPU cores.  Assignment of multiple MPI tasks to a GPU will happen
automatically if you create more MPI tasks/node than there are
GPUs/mode.  E.g. with 8 MPI tasks/node and 2 GPUs, each GPU will be
shared by 4 MPI tasks.

\sphinxAtStartPar
The GPU package also has limited support for OpenMP for both
multi\sphinxhyphen{}threading and vectorization of routines that are run on the CPUs.
This requires that the GPU library and LAMMPS are built with flags to
enable OpenMP support (e.g. \sphinxcode{\sphinxupquote{\sphinxhyphen{}fopenmp}}). Some styles for time integration
are also available in the GPU package. These run completely on the CPUs
in full double precision, but exploit multi\sphinxhyphen{}threading and vectorization
for faster performance.

\sphinxAtStartPar
Use the \sphinxcode{\sphinxupquote{\sphinxhyphen{}sf gpu}} {\hyperref[\detokenize{Run_options::doc}]{\sphinxcrossref{\DUrole{doc}{command\sphinxhyphen{}line switch}}}}, which will
automatically append “gpu” to styles that support it.  Use the \sphinxcode{\sphinxupquote{\sphinxhyphen{}pk
gpu Ng}} {\hyperref[\detokenize{Run_options::doc}]{\sphinxcrossref{\DUrole{doc}{command\sphinxhyphen{}line switch}}}} to set \sphinxcode{\sphinxupquote{Ng}} = \# of
GPUs/node to use. If \sphinxcode{\sphinxupquote{Ng}} is 0, the number is selected automatically as
the number of matching GPUs that have the highest number of compute
cores.

\begin{sphinxVerbatim}[commandchars=\\\{\}]
\PYG{c+c1}{\PYGZsh{} 1 MPI task uses 1 GPU}
lmp\PYGZus{}machine\PYG{+w}{ }\PYGZhy{}sf\PYG{+w}{ }gpu\PYG{+w}{ }\PYGZhy{}pk\PYG{+w}{ }gpu\PYG{+w}{ }\PYG{l+m}{1}\PYG{+w}{ }\PYGZhy{}in\PYG{+w}{ }\PYG{k}{in}.script

\PYG{c+c1}{\PYGZsh{} 12 MPI tasks share 2 GPUs on a single 16\PYGZhy{}core (or whatever) node}
mpirun\PYG{+w}{ }\PYGZhy{}np\PYG{+w}{ }\PYG{l+m}{12}\PYG{+w}{ }lmp\PYGZus{}machine\PYG{+w}{ }\PYGZhy{}sf\PYG{+w}{ }gpu\PYG{+w}{ }\PYGZhy{}pk\PYG{+w}{ }gpu\PYG{+w}{ }\PYG{l+m}{2}\PYG{+w}{ }\PYGZhy{}in\PYG{+w}{ }\PYG{k}{in}.script

\PYG{c+c1}{\PYGZsh{} ditto on 4 16\PYGZhy{}core nodes}
mpirun\PYG{+w}{ }\PYGZhy{}np\PYG{+w}{ }\PYG{l+m}{48}\PYG{+w}{ }\PYGZhy{}ppn\PYG{+w}{ }\PYG{l+m}{12}\PYG{+w}{ }lmp\PYGZus{}machine\PYG{+w}{ }\PYGZhy{}sf\PYG{+w}{ }gpu\PYG{+w}{ }\PYGZhy{}pk\PYG{+w}{ }gpu\PYG{+w}{ }\PYG{l+m}{2}\PYG{+w}{ }\PYGZhy{}in\PYG{+w}{ }\PYG{k}{in}.script
\end{sphinxVerbatim}

\sphinxAtStartPar
Note that if the \sphinxcode{\sphinxupquote{\sphinxhyphen{}sf gpu}} switch is used, it also issues a default
{\hyperref[\detokenize{package::doc}]{\sphinxcrossref{\DUrole{doc}{package gpu 0}}}} command, which will result in
automatic selection of the number of GPUs to use.

\sphinxAtStartPar
Using the \sphinxcode{\sphinxupquote{\sphinxhyphen{}pk}} switch explicitly allows for setting of the number of
GPUs/node to use and additional options.  Its syntax is the same as
the \sphinxcode{\sphinxupquote{package gpu}} command.  See the {\hyperref[\detokenize{package::doc}]{\sphinxcrossref{\DUrole{doc}{package}}}}
command page for details, including the default values used for
all its options if it is not specified.

\sphinxAtStartPar
Note that the default for the {\hyperref[\detokenize{package::doc}]{\sphinxcrossref{\DUrole{doc}{package gpu}}}} command is to
set the Newton flag to “off” pairwise interactions.  It does not
affect the setting for bonded interactions (LAMMPS default is “on”).
The “off” setting for pairwise interaction is currently required for
GPU package pair styles.

\sphinxAtStartPar
\sphinxstylestrong{Or run with the GPU package by editing an input script:}

\sphinxAtStartPar
The discussion above for the \sphinxcode{\sphinxupquote{mpirun}} or \sphinxcode{\sphinxupquote{mpiexec}} command, MPI
tasks/node, and use of multiple MPI tasks/GPU is the same.

\sphinxAtStartPar
Use the {\hyperref[\detokenize{suffix::doc}]{\sphinxcrossref{\DUrole{doc}{suffix gpu}}}} command, or you can explicitly add an
“gpu” suffix to individual styles in your input script, e.g.

\begin{sphinxVerbatim}[commandchars=\\\{\}]
\PYG{k}{pair\PYGZus{}style}\PYG{+w}{ }\PYG{n}{lj}\PYG{o}{/}\PYG{n}{cut}\PYG{o}{/}\PYG{n}{gpu}\PYG{+w}{ }\PYG{l+m}{2.5}
\end{sphinxVerbatim}

\sphinxAtStartPar
You must also use the {\hyperref[\detokenize{package::doc}]{\sphinxcrossref{\DUrole{doc}{package gpu}}}} command to enable the
GPU package, unless the \sphinxcode{\sphinxupquote{\sphinxhyphen{}sf gpu}} or \sphinxcode{\sphinxupquote{\sphinxhyphen{}pk gpu}} {\hyperref[\detokenize{Run_options::doc}]{\sphinxcrossref{\DUrole{doc}{command\sphinxhyphen{}line switches}}}} were used.  It specifies the number of
GPUs/node to use, as well as other options.

\sphinxAtStartPar
\sphinxstylestrong{Speed\sphinxhyphen{}ups to expect:}

\sphinxAtStartPar
The performance of a GPU versus a multicore CPU is a function of your
hardware, which pair style is used, the number of atoms/GPU, and the
precision used on the GPU (double, single, mixed). Using the GPU package
in OpenCL mode on CPUs (which uses vectorization and multithreading) is
usually resulting in inferior performance compared to using LAMMPS’ native
threading and vectorization support in the OPENMP and INTEL packages.

\sphinxAtStartPar
See the \sphinxhref{https://www.lammps.org/bench.html}{Benchmark page} of the
LAMMPS website for performance of the GPU package on various
hardware, including the Titan HPC platform at ORNL.

\sphinxAtStartPar
You should also experiment with how many MPI tasks per GPU to use to
give the best performance for your problem and machine.  This is also
a function of the problem size and the pair style being using.
Likewise, you should experiment with the precision setting for the GPU
library to see if single or mixed precision will give accurate
results, since they will typically be faster.

\sphinxAtStartPar
MPI parallelism typically outperforms OpenMP parallelism, but in some
cases using fewer MPI tasks and multiple OpenMP threads with the GPU
package can give better performance. 3\sphinxhyphen{}body potentials can often perform
better with multiple OMP threads because the inter\sphinxhyphen{}process communication
is higher for these styles with the GPU package in order to allow
deterministic results.

\sphinxAtStartPar
\sphinxstylestrong{Guidelines for best performance:}
\begin{itemize}
\item {} 
\sphinxAtStartPar
Using multiple MPI tasks per GPU will often give the best performance,
as allowed my most multicore CPU/GPU configurations.

\item {} 
\sphinxAtStartPar
If the number of particles per MPI task is small (e.g. 100s of
particles), it can be more efficient to run with fewer MPI tasks per
GPU, even if you do not use all the cores on the compute node.

\item {} 
\sphinxAtStartPar
The {\hyperref[\detokenize{package::doc}]{\sphinxcrossref{\DUrole{doc}{package gpu}}}} command has several options for tuning
performance.  Neighbor lists can be built on the GPU or CPU.  Force
calculations can be dynamically balanced across the CPU cores and
GPUs.  GPU\sphinxhyphen{}specific settings can be made which can be optimized
for different hardware.  See the {\hyperref[\detokenize{package::doc}]{\sphinxcrossref{\DUrole{doc}{package}}}} command
page for details.

\item {} 
\sphinxAtStartPar
As described by the {\hyperref[\detokenize{package::doc}]{\sphinxcrossref{\DUrole{doc}{package gpu}}}} command, GPU
accelerated pair styles can perform computations asynchronously with
CPU computations. The “Pair” time reported by LAMMPS will be the
maximum of the time required to complete the CPU pair style
computations and the time required to complete the GPU pair style
computations. Any time spent for GPU\sphinxhyphen{}enabled pair styles for
computations that run simultaneously with {\hyperref[\detokenize{bond_style::doc}]{\sphinxcrossref{\DUrole{doc}{bond}}}},
{\hyperref[\detokenize{angle_style::doc}]{\sphinxcrossref{\DUrole{doc}{angle}}}}, {\hyperref[\detokenize{dihedral_style::doc}]{\sphinxcrossref{\DUrole{doc}{dihedral}}}},
{\hyperref[\detokenize{improper_style::doc}]{\sphinxcrossref{\DUrole{doc}{improper}}}}, and {\hyperref[\detokenize{kspace_style::doc}]{\sphinxcrossref{\DUrole{doc}{long\sphinxhyphen{}range}}}}
calculations will not be included in the “Pair” time.

\item {} 
\sphinxAtStartPar
Since only part of the pppm kspace style is GPU accelerated, it
may be faster to only use GPU acceleration for Pair styles with
long\sphinxhyphen{}range electrostatics.  See the “pair/only” keyword of the
package command for a shortcut to do that.  The work between kspace
on the CPU and non\sphinxhyphen{}bonded interactions on the GPU can be balanced
through adjusting the coulomb cutoff without loss of accuracy.

\item {} 
\sphinxAtStartPar
When the \sphinxstyleemphasis{mode} setting for the package gpu command is force/neigh,
the time for neighbor list calculations on the GPU will be added into
the “Pair” time, not the “Neigh” time.  An additional breakdown of the
times required for various tasks on the GPU (data copy, neighbor
calculations, force computations, etc) are output only with the LAMMPS
screen output (not in the log file) at the end of each run.  These
timings represent total time spent on the GPU for each routine,
regardless of asynchronous CPU calculations.

\item {} 
\sphinxAtStartPar
The output section “GPU Time Info (average)” reports “Max Mem / Proc”.
This is the maximum memory used at one time on the GPU for data
storage by a single MPI process.

\end{itemize}


\subsubsection{Restrictions}
\label{\detokenize{Speed_gpu:restrictions}}
\sphinxAtStartPar
None.

\sphinxstepscope


\subsection{INTEL package}
\label{\detokenize{Speed_intel:intel-package}}\label{\detokenize{Speed_intel::doc}}
\sphinxAtStartPar
The INTEL package is maintained by Mike Brown at Intel
Corporation.  It provides two methods for accelerating simulations,
depending on the hardware you have.  The first is acceleration on
Intel CPUs by running in single, mixed, or double precision with
vectorization.  The second is acceleration on Intel Xeon Phi
co\sphinxhyphen{}processors via offloading neighbor list and non\sphinxhyphen{}bonded force
calculations to the Phi.  The same C++ code is used in both cases.
When offloading to a co\sphinxhyphen{}processor from a CPU, the same routine is run
twice, once on the CPU and once with an offload flag. This allows
LAMMPS to run on the CPU cores and co\sphinxhyphen{}processor cores simultaneously.


\subsubsection{Currently Available INTEL Styles}
\label{\detokenize{Speed_intel:currently-available-intel-styles}}\begin{itemize}
\item {} 
\sphinxAtStartPar
Angle Styles: charmm, harmonic

\item {} 
\sphinxAtStartPar
Bond Styles: fene, fourier, harmonic

\item {} 
\sphinxAtStartPar
Dihedral Styles: charmm, fourier, harmonic, opls

\item {} 
\sphinxAtStartPar
Fixes: nve, npt, nvt, nvt/sllod, nve/asphere

\item {} 
\sphinxAtStartPar
Improper Styles: cvff, harmonic

\item {} 
\sphinxAtStartPar
Pair Styles: airebo, airebo/morse, buck/coul/cut, buck/coul/long,
buck, dpd, eam, eam/alloy, eam/fs, gayberne, lj/charmm/coul/charmm,
lj/charmm/coul/long, lj/cut, lj/cut/coul/long, lj/long/coul/long,
rebo, sw, tersoff

\item {} 
\sphinxAtStartPar
K\sphinxhyphen{}Space Styles: pppm, pppm/disp

\end{itemize}

\begin{sphinxadmonition}{warning}{Warning:}
\sphinxAtStartPar
None of the styles in the INTEL package currently
support computing per\sphinxhyphen{}atom stress.  If any compute or fix in your
input requires it, LAMMPS will abort with an error message.
\end{sphinxadmonition}


\subsubsection{Speed\sphinxhyphen{}up to expect}
\label{\detokenize{Speed_intel:speed-up-to-expect}}
\sphinxAtStartPar
The speedup will depend on your simulation, the hardware, which
styles are used, the number of atoms, and the floating\sphinxhyphen{}point
precision mode. Performance improvements are shown compared to
LAMMPS \sphinxstyleemphasis{without using other acceleration packages} as these are
under active development (and subject to performance changes). The
measurements were performed using the input files available in
the \sphinxcode{\sphinxupquote{src/INTEL/TEST}} directory with the provided run script.
These are scalable in size; the results given are with 512K
particles (524K for Liquid Crystal). Most of the simulations are
standard LAMMPS benchmarks (indicated by the filename extension in
parenthesis) with modifications to the run length and to add a
warm\sphinxhyphen{}up run (for use with offload benchmarks).

\noindent{\hspace*{\fill}\sphinxincludegraphics{{user_intel}.png}\hspace*{\fill}}

\sphinxAtStartPar
Results are speedups obtained on Intel Xeon E5\sphinxhyphen{}2697v4 processors
(code\sphinxhyphen{}named Broadwell), Intel Xeon Phi 7250 processors (code\sphinxhyphen{}named
Knights Landing), and Intel Xeon Gold 6148 processors (code\sphinxhyphen{}named
Skylake) with “June 2017” LAMMPS built with Intel Parallel Studio
2017 update 2. Results are with 1 MPI task per physical core. See
\sphinxcode{\sphinxupquote{src/INTEL/TEST/README}} for the raw simulation rates and
instructions to reproduce.


\bigskip\hrule\bigskip



\subsubsection{Accuracy and order of operations}
\label{\detokenize{Speed_intel:accuracy-and-order-of-operations}}
\sphinxAtStartPar
In most molecular dynamics software, parallelization parameters
(\# of MPI, OpenMP, and vectorization) can change the results due
to changing the order of operations with finite\sphinxhyphen{}precision
calculations. The INTEL package is deterministic. This means
that the results should be reproducible from run to run with the
\sphinxstyleemphasis{same} parallel configurations and when using deterministic
libraries or library settings (MPI, OpenMP, FFT). However, there
are differences in the INTEL package that can change the
order of operations compared to LAMMPS without acceleration:
\begin{itemize}
\item {} 
\sphinxAtStartPar
Neighbor lists can be created in a different order

\item {} 
\sphinxAtStartPar
Bins used for sorting atoms can be oriented differently

\item {} 
\sphinxAtStartPar
The default stencil order for PPPM is 7. By default, LAMMPS will
calculate other PPPM parameters to fit the desired accuracy with
this order

\item {} 
\sphinxAtStartPar
The \sphinxstyleemphasis{newton} setting applies to all atoms, not just atoms shared
between MPI tasks

\item {} 
\sphinxAtStartPar
Vectorization can change the order for adding pairwise forces

\item {} 
\sphinxAtStartPar
When using the \sphinxcode{\sphinxupquote{\sphinxhyphen{}DLMP\_USE\_MKL\_RNG}} define (all included intel optimized
makefiles do) at build time, the random number generator for
dissipative particle dynamics (\sphinxcode{\sphinxupquote{pair style dpd/intel}}) uses the Mersenne
Twister generator included in the Intel MKL library (that should be
more robust than the default Masaglia random number generator)

\end{itemize}

\sphinxAtStartPar
The precision mode (described below) used with the INTEL
package can change the \sphinxstyleemphasis{accuracy} of the calculations. For the
default \sphinxstyleemphasis{mixed} precision option, calculations between pairs or
triplets of atoms are performed in single precision, intended to
be within the inherent error of MD simulations. All accumulation
is performed in double precision to prevent the error from growing
with the number of atoms in the simulation. \sphinxstyleemphasis{Single} precision
mode should not be used without appropriate validation.


\bigskip\hrule\bigskip



\subsubsection{Quick Start for Experienced Users}
\label{\detokenize{Speed_intel:quick-start-for-experienced-users}}
\sphinxAtStartPar
LAMMPS should be built with the INTEL package installed.
Simulations should be run with 1 MPI task per physical \sphinxstyleemphasis{core},
not \sphinxstyleemphasis{hardware thread}.
\begin{itemize}
\item {} 
\sphinxAtStartPar
Edit \sphinxcode{\sphinxupquote{src/MAKE/OPTIONS/Makefile.intel\_cpu\_intelmpi}} as necessary.

\item {} 
\sphinxAtStartPar
Set the environment variable \sphinxcode{\sphinxupquote{KMP\_BLOCKTIME=0}}

\item {} 
\sphinxAtStartPar
\sphinxcode{\sphinxupquote{\sphinxhyphen{}pk intel 0 omp \$t \sphinxhyphen{}sf intel}} added to LAMMPS command\sphinxhyphen{}line

\item {} 
\sphinxAtStartPar
\sphinxcode{\sphinxupquote{\$t}} should be 2 for Intel Xeon CPUs and 2 or 4 for Intel Xeon Phi

\item {} 
\sphinxAtStartPar
For some of the simple 2\sphinxhyphen{}body potentials without long\sphinxhyphen{}range
electrostatics, performance and scalability can be better with
the \sphinxcode{\sphinxupquote{newton off}} setting added to the input script

\item {} 
\sphinxAtStartPar
For simulations on higher node counts, add \sphinxcode{\sphinxupquote{processors * * * grid
numa}} to the beginning of the input script for better scalability

\item {} 
\sphinxAtStartPar
If using \sphinxcode{\sphinxupquote{kspace\_style pppm}} in the input script, add
\sphinxcode{\sphinxupquote{kspace\_modify diff ad}} for better performance

\end{itemize}

\sphinxAtStartPar
For Intel Xeon Phi CPUs:
\begin{itemize}
\item {} 
\sphinxAtStartPar
Runs should be performed using MCDRAM.

\end{itemize}

\sphinxAtStartPar
For simulations using \sphinxcode{\sphinxupquote{kspace\_style pppm}} on Intel CPUs supporting
AVX\sphinxhyphen{}512:
\begin{itemize}
\item {} 
\sphinxAtStartPar
Add \sphinxcode{\sphinxupquote{kspace\_modify diff ad}} to the input script

\item {} 
\sphinxAtStartPar
The command\sphinxhyphen{}line option should be changed to
\sphinxcode{\sphinxupquote{\sphinxhyphen{}pk intel 0 omp \$r lrt yes \sphinxhyphen{}sf intel}} where \sphinxcode{\sphinxupquote{\$r}} is the number of
threads minus 1.

\item {} 
\sphinxAtStartPar
Do not use thread affinity (set \sphinxcode{\sphinxupquote{KMP\_AFFINITY=none}})

\item {} 
\sphinxAtStartPar
The \sphinxcode{\sphinxupquote{newton off}} setting may provide better scalability

\end{itemize}

\sphinxAtStartPar
For Intel Xeon Phi co\sphinxhyphen{}processors (Offload):
\begin{itemize}
\item {} 
\sphinxAtStartPar
Edit \sphinxcode{\sphinxupquote{src/MAKE/OPTIONS/Makefile.intel\_co\sphinxhyphen{}processor}} as necessary

\item {} 
\sphinxAtStartPar
\sphinxcode{\sphinxupquote{\sphinxhyphen{}pk intel N omp 1}} added to command\sphinxhyphen{}line where \sphinxcode{\sphinxupquote{N}} is the number of
co\sphinxhyphen{}processors per node.

\end{itemize}


\bigskip\hrule\bigskip



\subsubsection{Required hardware/software}
\label{\detokenize{Speed_intel:required-hardware-software}}
\sphinxAtStartPar
When using Intel compilers version 16.0 or later is required.

\sphinxAtStartPar
In order to use offload to co\sphinxhyphen{}processors, an Intel Xeon Phi
co\sphinxhyphen{}processor and an Intel compiler are required.

\sphinxAtStartPar
Although any compiler can be used with the INTEL package,
currently, vectorization directives are disabled by default when
not using Intel compilers due to lack of standard support and
observations of decreased performance. The OpenMP standard now
supports directives for vectorization and we plan to transition the
code to this standard once it is available in most compilers. We
expect this to allow improved performance and support with other
compilers.

\sphinxAtStartPar
For Intel Xeon Phi x200 series processors (code\sphinxhyphen{}named Knights
Landing), there are multiple configuration options for the hardware.
For best performance, we recommend that the MCDRAM is configured in
“Flat” mode and with the cluster mode set to “Quadrant” or “SNC4”.
“Cache” mode can also be used, although the performance might be
slightly lower.


\subsubsection{Notes about Simultaneous Multithreading}
\label{\detokenize{Speed_intel:notes-about-simultaneous-multithreading}}
\sphinxAtStartPar
Modern CPUs often support Simultaneous Multithreading (SMT). On
Intel processors, this is called Hyper\sphinxhyphen{}Threading (HT) technology.
SMT is hardware support for running multiple threads efficiently on
a single core. \sphinxstyleemphasis{Hardware threads} or \sphinxstyleemphasis{logical cores} are often used
to refer to the number of threads that are supported in hardware.
For example, the Intel Xeon E5\sphinxhyphen{}2697v4 processor is described
as having 36 cores and 72 threads. This means that 36 MPI processes
or OpenMP threads can run simultaneously on separate cores, but that
up to 72 MPI processes or OpenMP threads can be running on the CPU
without costly operating system context switches.

\sphinxAtStartPar
Molecular dynamics simulations will often run faster when making use
of SMT. If a thread becomes stalled, for example because it is
waiting on data that has not yet arrived from memory, another thread
can start running so that the CPU pipeline is still being used
efficiently. Although benefits can be seen by launching a MPI task
for every hardware thread, for multinode simulations, we recommend
that OpenMP threads are used for SMT instead, either with the
INTEL package, {\hyperref[\detokenize{Speed_omp::doc}]{\sphinxcrossref{\DUrole{doc}{OPENMP package}}}}, or
{\hyperref[\detokenize{Speed_kokkos::doc}]{\sphinxcrossref{\DUrole{doc}{KOKKOS package}}}}. In the example above, up
to 36X speedups can be observed by using all 36 physical cores with
LAMMPS. By using all 72 hardware threads, an additional 10\sphinxhyphen{}30\%
performance gain can be achieved.

\sphinxAtStartPar
The BIOS on many platforms allows SMT to be disabled, however, we do
not recommend this on modern processors as there is little to no
benefit for any software package in most cases. The operating system
will report every hardware thread as a separate core allowing one to
determine the number of hardware threads available. On Linux systems,
this information can normally be obtained with:

\begin{sphinxVerbatim}[commandchars=\\\{\}]
cat\PYG{+w}{ }/proc/cpuinfo
\end{sphinxVerbatim}


\subsubsection{Building LAMMPS with the INTEL package}
\label{\detokenize{Speed_intel:building-lammps-with-the-intel-package}}
\sphinxAtStartPar
See the {\hyperref[\detokenize{Build_extras:intel}]{\sphinxcrossref{\DUrole{std,std-ref}{Build extras}}}} page for
instructions.  Some additional details are covered here.

\sphinxAtStartPar
For building with make, several example Makefiles for building with
the Intel compiler are included with LAMMPS in the \sphinxcode{\sphinxupquote{src/MAKE/OPTIONS/}}
directory:

\begin{sphinxVerbatim}[commandchars=\\\{\}]
Makefile.intel\PYGZus{}cpu\PYGZus{}intelmpi\PYG{+w}{ }\PYG{c+c1}{\PYGZsh{} Intel Compiler, Intel MPI, No Offload}
Makefile.knl\PYG{+w}{                }\PYG{c+c1}{\PYGZsh{} Intel Compiler, Intel MPI, No Offload}
Makefile.intel\PYGZus{}cpu\PYGZus{}mpich\PYG{+w}{    }\PYG{c+c1}{\PYGZsh{} Intel Compiler, MPICH, No Offload}
Makefile.intel\PYGZus{}cpu\PYGZus{}openpmi\PYG{+w}{  }\PYG{c+c1}{\PYGZsh{} Intel Compiler, OpenMPI, No Offload}
Makefile.intel\PYGZus{}co\PYGZhy{}processor\PYG{+w}{  }\PYG{c+c1}{\PYGZsh{} Intel Compiler, Intel MPI, Offload}
\end{sphinxVerbatim}

\sphinxAtStartPar
Makefile.knl is identical to Makefile.intel\_cpu\_intelmpi except that
it explicitly specifies that vectorization should be for Intel Xeon
Phi x200 processors making it easier to cross\sphinxhyphen{}compile. For users with
recent installations of Intel Parallel Studio, the process can be as
simple as:

\begin{sphinxVerbatim}[commandchars=\\\{\}]
make\PYG{+w}{ }yes\PYGZhy{}intel
\PYG{n+nb}{source}\PYG{+w}{ }/opt/intel/parallel\PYGZus{}studio\PYGZus{}xe\PYGZus{}2016.3.067/psxevars.sh
\PYG{c+c1}{\PYGZsh{} or psxevars.csh for C\PYGZhy{}shell}
make\PYG{+w}{ }intel\PYGZus{}cpu\PYGZus{}intelmpi
\end{sphinxVerbatim}

\sphinxAtStartPar
Note that if you build with support for a Phi co\sphinxhyphen{}processor, the same
binary can be used on nodes with or without co\sphinxhyphen{}processors installed.
However, if you do not have co\sphinxhyphen{}processors on your system, building
without offload support will produce a smaller binary.

\sphinxAtStartPar
The general requirements for Makefiles with the INTEL package
are as follows. When using Intel compilers, \sphinxcode{\sphinxupquote{\sphinxhyphen{}restrict}} is required
and \sphinxcode{\sphinxupquote{\sphinxhyphen{}qopenmp}} is highly recommended for \sphinxcode{\sphinxupquote{CCFLAGS}} and \sphinxcode{\sphinxupquote{LINKFLAGS}}.
\sphinxcode{\sphinxupquote{CCFLAGS}} should include \sphinxcode{\sphinxupquote{\sphinxhyphen{}DLMP\_INTEL\_USELRT}} (unless POSIX Threads
are not supported in the build environment) and \sphinxcode{\sphinxupquote{\sphinxhyphen{}DLMP\_USE\_MKL\_RNG}}
(unless Intel Math Kernel Library (MKL) is not available in the build
environment). For Intel compilers, \sphinxcode{\sphinxupquote{LIB}} should include \sphinxcode{\sphinxupquote{\sphinxhyphen{}ltbbmalloc}}
or if the library is not available, \sphinxcode{\sphinxupquote{\sphinxhyphen{}DLMP\_INTEL\_NO\_TBB}} can be added
to \sphinxcode{\sphinxupquote{CCFLAGS}}. For builds supporting offload, \sphinxcode{\sphinxupquote{\sphinxhyphen{}DLMP\_INTEL\_OFFLOAD}} is
required for \sphinxcode{\sphinxupquote{CCFLAGS}} and \sphinxcode{\sphinxupquote{\sphinxhyphen{}qoffload}} is required for \sphinxcode{\sphinxupquote{LINKFLAGS}}. Other
recommended \sphinxcode{\sphinxupquote{CCFLAG}} options for best performance are \sphinxcode{\sphinxupquote{\sphinxhyphen{}O2 \sphinxhyphen{}fno\sphinxhyphen{}alias
\sphinxhyphen{}ansi\sphinxhyphen{}alias \sphinxhyphen{}qoverride\sphinxhyphen{}limits fp\sphinxhyphen{}model fast=2 \sphinxhyphen{}no\sphinxhyphen{}prec\sphinxhyphen{}div}}.

\begin{sphinxadmonition}{note}{Note:}
\sphinxAtStartPar
See the \sphinxcode{\sphinxupquote{src/INTEL/README}} file for additional flags that
might be needed for best performance on Intel server processors
code\sphinxhyphen{}named “Skylake”.
\end{sphinxadmonition}

\begin{sphinxadmonition}{note}{Note:}
\sphinxAtStartPar
The vectorization and math capabilities can differ depending on
the CPU. For Intel compilers, the \sphinxcode{\sphinxupquote{\sphinxhyphen{}x}} flag specifies the type of
processor for which to optimize. \sphinxcode{\sphinxupquote{\sphinxhyphen{}xHost}} specifies that the compiler
should build for the processor used for compiling. For Intel Xeon Phi
x200 series processors, this option is \sphinxcode{\sphinxupquote{\sphinxhyphen{}xMIC\sphinxhyphen{}AVX512}}. For fourth
generation Intel Xeon (v4/Broadwell) processors, \sphinxcode{\sphinxupquote{\sphinxhyphen{}xCORE\sphinxhyphen{}AVX2}} should
be used. For older Intel Xeon processors, \sphinxcode{\sphinxupquote{\sphinxhyphen{}xAVX}} will perform best
in general for the different simulations in LAMMPS. The default
in most of the example Makefiles is to use \sphinxcode{\sphinxupquote{\sphinxhyphen{}xHost}}, however this
should not be used when cross\sphinxhyphen{}compiling.
\end{sphinxadmonition}


\subsubsection{Running LAMMPS with the INTEL package}
\label{\detokenize{Speed_intel:running-lammps-with-the-intel-package}}
\sphinxAtStartPar
Running LAMMPS with the INTEL package is similar to normal use
with the exceptions that one should 1) specify that LAMMPS should use
the INTEL package, 2) specify the number of OpenMP threads, and
3) optionally specify the specific LAMMPS styles that should use the
INTEL package. 1) and 2) can be performed from the command\sphinxhyphen{}line
or by editing the input script. 3) requires editing the input script.
Advanced performance tuning options are also described below to get
the best performance.

\sphinxAtStartPar
When running on a single node (including runs using offload to a
co\sphinxhyphen{}processor), best performance is normally obtained by using 1 MPI
task per physical core and additional OpenMP threads with SMT. For
Intel Xeon processors, 2 OpenMP threads should be used for SMT.
For Intel Xeon Phi CPUs, 2 or 4 OpenMP threads should be used
(best choice depends on the simulation). In cases where the user
specifies that LRT mode is used (described below), 1 or 3 OpenMP
threads should be used. For multi\sphinxhyphen{}node runs, using 1 MPI task per
physical core will often perform best, however, depending on the
machine and scale, users might get better performance by decreasing
the number of MPI tasks and using more OpenMP threads. For
performance, the product of the number of MPI tasks and OpenMP
threads should not exceed the number of available hardware threads in
almost all cases.

\begin{sphinxadmonition}{note}{Note:}
\sphinxAtStartPar
Setting core affinity is often used to pin MPI tasks and OpenMP
threads to a core or group of cores so that memory access can be
uniform. Unless disabled at build time, affinity for MPI tasks and
OpenMP threads on the host (CPU) will be set by default on the host
\sphinxstyleemphasis{when using offload to a co\sphinxhyphen{}processor}. In this case, it is unnecessary
to use other methods to control affinity (e.g. \sphinxcode{\sphinxupquote{taskset}}, \sphinxcode{\sphinxupquote{numactl}},
\sphinxcode{\sphinxupquote{I\_MPI\_PIN\_DOMAIN}}, etc.). This can be disabled with the \sphinxstyleemphasis{no\_affinity}
option to the {\hyperref[\detokenize{package::doc}]{\sphinxcrossref{\DUrole{doc}{package intel}}}} command or by disabling the
option at build time (by adding \sphinxcode{\sphinxupquote{\sphinxhyphen{}DINTEL\_OFFLOAD\_NOAFFINITY}} to the
\sphinxcode{\sphinxupquote{CCFLAGS}} line of your Makefile). Disabling this option is not
recommended, especially when running on a machine with Intel
Hyper\sphinxhyphen{}Threading technology disabled.
\end{sphinxadmonition}


\subsubsection{Run with the INTEL package from the command line}
\label{\detokenize{Speed_intel:run-with-the-intel-package-from-the-command-line}}
\sphinxAtStartPar
To enable INTEL optimizations for all available styles used in
the input script, the \sphinxcode{\sphinxupquote{\sphinxhyphen{}sf intel}} {\hyperref[\detokenize{Run_options::doc}]{\sphinxcrossref{\DUrole{doc}{command\sphinxhyphen{}line switch}}}} can be used without any requirement for
editing the input script. This switch will automatically append
“intel” to styles that support it. It also invokes a default command:
{\hyperref[\detokenize{package::doc}]{\sphinxcrossref{\DUrole{doc}{package intel 1}}}}. This package command is used to set
options for the INTEL package.  The default package command will
specify that INTEL calculations are performed in mixed precision,
that the number of OpenMP threads is specified by the OMP\_NUM\_THREADS
environment variable, and that if co\sphinxhyphen{}processors are present and the
binary was built with offload support, that 1 co\sphinxhyphen{}processor per node
will be used with automatic balancing of work between the CPU and the
co\sphinxhyphen{}processor.

\sphinxAtStartPar
You can specify different options for the INTEL package by using
the \sphinxcode{\sphinxupquote{\sphinxhyphen{}pk intel Nphi}} {\hyperref[\detokenize{Run_options::doc}]{\sphinxcrossref{\DUrole{doc}{command\sphinxhyphen{}line switch}}}} with
keyword/value pairs as specified in the documentation. Here, \sphinxcode{\sphinxupquote{Nphi}} = \#
of Xeon Phi co\sphinxhyphen{}processors/node (ignored without offload
support). Common options to the INTEL package include \sphinxstyleemphasis{omp} to
override any \sphinxcode{\sphinxupquote{OMP\_NUM\_THREADS}} setting and specify the number of OpenMP
threads, \sphinxstyleemphasis{mode} to set the floating\sphinxhyphen{}point precision mode, and \sphinxstyleemphasis{lrt} to
enable Long\sphinxhyphen{}Range Thread mode as described below. See the {\hyperref[\detokenize{package::doc}]{\sphinxcrossref{\DUrole{doc}{package intel}}}} command for details, including the default values
used for all its options if not specified, and how to set the number
of OpenMP threads via the \sphinxcode{\sphinxupquote{OMP\_NUM\_THREADS}} environment variable if
desired.

\sphinxAtStartPar
Examples (see documentation for your MPI/Machine for differences in
launching MPI applications):

\begin{sphinxVerbatim}[commandchars=\\\{\}]
\PYG{c+c1}{\PYGZsh{} 2 nodes, 36 MPI tasks/node, \PYGZdl{}OMP\PYGZus{}NUM\PYGZus{}THREADS OpenMP Threads}
mpirun\PYG{+w}{ }\PYGZhy{}np\PYG{+w}{ }\PYG{l+m}{72}\PYG{+w}{ }\PYGZhy{}ppn\PYG{+w}{ }\PYG{l+m}{36}\PYG{+w}{ }lmp\PYGZus{}machine\PYG{+w}{ }\PYGZhy{}sf\PYG{+w}{ }intel\PYG{+w}{ }\PYGZhy{}in\PYG{+w}{ }\PYG{k}{in}.script

\PYG{c+c1}{\PYGZsh{} Don\PYGZsq{}t use any co\PYGZhy{}processors that might be available,}
\PYG{c+c1}{\PYGZsh{} use 2 OpenMP threads for each task, use double precision}
mpirun\PYG{+w}{ }\PYGZhy{}np\PYG{+w}{ }\PYG{l+m}{72}\PYG{+w}{ }\PYGZhy{}ppn\PYG{+w}{ }\PYG{l+m}{36}\PYG{+w}{ }lmp\PYGZus{}machine\PYG{+w}{ }\PYGZhy{}sf\PYG{+w}{ }intel\PYG{+w}{ }\PYGZhy{}in\PYG{+w}{ }\PYG{k}{in}.script\PYG{+w}{ }\PYG{l+s+se}{\PYGZbs{}}
\PYG{+w}{       }\PYGZhy{}pk\PYG{+w}{ }intel\PYG{+w}{ }\PYG{l+m}{0}\PYG{+w}{ }omp\PYG{+w}{ }\PYG{l+m}{2}\PYG{+w}{ }mode\PYG{+w}{ }double
\end{sphinxVerbatim}


\subsubsection{Or run with the INTEL package by editing an input script}
\label{\detokenize{Speed_intel:or-run-with-the-intel-package-by-editing-an-input-script}}
\sphinxAtStartPar
As an alternative to adding command\sphinxhyphen{}line arguments, the input script
can be edited to enable the INTEL package. This requires adding
the {\hyperref[\detokenize{package::doc}]{\sphinxcrossref{\DUrole{doc}{package intel}}}} command to the top of the input
script. For the second example above, this would be:

\begin{sphinxVerbatim}[commandchars=\\\{\}]
\PYG{k}{package}\PYG{+w}{ }\PYG{n}{intel}\PYG{+w}{ }\PYG{l+m}{0}\PYG{+w}{ }\PYG{n}{omp}\PYG{+w}{ }\PYG{l+m}{2}\PYG{+w}{ }\PYG{n}{mode}\PYG{+w}{ }\PYG{n}{double}
\end{sphinxVerbatim}

\sphinxAtStartPar
To enable the INTEL package only for individual styles, you can
add an “intel” suffix to the individual style, e.g.:

\begin{sphinxVerbatim}[commandchars=\\\{\}]
\PYG{k}{pair\PYGZus{}style}\PYG{+w}{ }\PYG{n}{lj}\PYG{o}{/}\PYG{n}{cut}\PYG{o}{/}\PYG{n}{intel}\PYG{+w}{ }\PYG{l+m}{2.5}
\end{sphinxVerbatim}

\sphinxAtStartPar
Alternatively, the {\hyperref[\detokenize{suffix::doc}]{\sphinxcrossref{\DUrole{doc}{suffix intel}}}} command can be added to
the input script to enable INTEL styles for the commands that
follow in the input script.


\subsubsection{Tuning for Performance}
\label{\detokenize{Speed_intel:tuning-for-performance}}
\begin{sphinxadmonition}{note}{Note:}
\sphinxAtStartPar
The INTEL package will perform better with modifications
to the input script when {\hyperref[\detokenize{kspace_style::doc}]{\sphinxcrossref{\DUrole{doc}{PPPM}}}} is used:
{\hyperref[\detokenize{kspace_modify::doc}]{\sphinxcrossref{\DUrole{doc}{kspace\_modify diff ad}}}} should be added to the
input script.
\end{sphinxadmonition}

\sphinxAtStartPar
Long\sphinxhyphen{}Range Thread (LRT) mode is an option to the {\hyperref[\detokenize{package::doc}]{\sphinxcrossref{\DUrole{doc}{package intel}}}} command that can improve performance when using
{\hyperref[\detokenize{kspace_style::doc}]{\sphinxcrossref{\DUrole{doc}{PPPM}}}} for long\sphinxhyphen{}range electrostatics on processors
with SMT. It generates an extra pthread for each MPI task. The thread
is dedicated to performing some of the PPPM calculations and MPI
communications. This feature requires setting the pre\sphinxhyphen{}processor flag
\sphinxcode{\sphinxupquote{\sphinxhyphen{}DLMP\_INTEL\_USELRT}} in the makefile when compiling LAMMPS. It is unset
in the default makefiles (\sphinxcode{\sphinxupquote{Makefile.mpi}} and \sphinxcode{\sphinxupquote{Makefile.serial}}) but
it is set in all makefiles tuned for the INTEL package.  On Intel
Xeon Phi x200 series CPUs, the LRT feature will likely improve
performance, even on a single node. On Intel Xeon processors, using
this mode might result in better performance when using multiple nodes,
depending on the specific machine configuration. To enable LRT mode,
specify that the number of OpenMP threads is one less than would
normally be used for the run and add the \sphinxcode{\sphinxupquote{lrt yes}} option to the \sphinxcode{\sphinxupquote{\sphinxhyphen{}pk}}
command\sphinxhyphen{}line suffix or “package intel” command. For example, if a run
would normally perform best with “\sphinxhyphen{}pk intel 0 omp 4”, instead use
\sphinxcode{\sphinxupquote{\sphinxhyphen{}pk intel 0 omp 3 lrt yes}}. When using LRT, you should set the
environment variable \sphinxcode{\sphinxupquote{KMP\_AFFINITY=none}}. LRT mode is not supported
when using offload.

\begin{sphinxadmonition}{note}{Note:}
\sphinxAtStartPar
Changing the {\hyperref[\detokenize{newton::doc}]{\sphinxcrossref{\DUrole{doc}{newton}}}} setting to off can improve
performance and/or scalability for simple 2\sphinxhyphen{}body potentials such as
lj/cut or when using LRT mode on processors supporting AVX\sphinxhyphen{}512.
\end{sphinxadmonition}

\sphinxAtStartPar
Not all styles are supported in the INTEL package. You can mix
the INTEL package with styles from the {\hyperref[\detokenize{Speed_opt::doc}]{\sphinxcrossref{\DUrole{doc}{OPT}}}}
package or the {\hyperref[\detokenize{Speed_omp::doc}]{\sphinxcrossref{\DUrole{doc}{OPENMP package}}}}. Of course, this
requires that these packages were installed at build time. This can
performed automatically by using \sphinxcode{\sphinxupquote{\sphinxhyphen{}sf hybrid intel opt}} or \sphinxcode{\sphinxupquote{\sphinxhyphen{}sf hybrid
intel omp}} command\sphinxhyphen{}line options. Alternatively, the “opt” and “omp”
suffixes can be appended manually in the input script. For the latter,
the {\hyperref[\detokenize{package::doc}]{\sphinxcrossref{\DUrole{doc}{package omp}}}} command must be in the input script or
the \sphinxcode{\sphinxupquote{\sphinxhyphen{}pk omp Nt}} {\hyperref[\detokenize{Run_options::doc}]{\sphinxcrossref{\DUrole{doc}{command\sphinxhyphen{}line switch}}}} must be used
where \sphinxcode{\sphinxupquote{Nt}} is the number of OpenMP threads. The number of OpenMP threads
should not be set differently for the different packages. Note that
the {\hyperref[\detokenize{suffix::doc}]{\sphinxcrossref{\DUrole{doc}{suffix hybrid intel omp}}}} command can also be used
within the input script to automatically append the “omp” suffix to
styles when INTEL styles are not available.

\begin{sphinxadmonition}{note}{Note:}
\sphinxAtStartPar
For simulations on higher node counts, add {\hyperref[\detokenize{processors::doc}]{\sphinxcrossref{\DUrole{doc}{processors * * * grid numa}}}} to the beginning of the input script for
better scalability.
\end{sphinxadmonition}

\sphinxAtStartPar
When running on many nodes, performance might be better when using
fewer OpenMP threads and more MPI tasks. This will depend on the
simulation and the machine. Using the {\hyperref[\detokenize{run_style::doc}]{\sphinxcrossref{\DUrole{doc}{verlet/split}}}}
run style might also give better performance for simulations with
{\hyperref[\detokenize{kspace_style::doc}]{\sphinxcrossref{\DUrole{doc}{PPPM}}}} electrostatics. Note that this is an
alternative to LRT mode and the two cannot be used together.

\sphinxAtStartPar
Currently, when using Intel MPI with Intel Xeon Phi x200 series
CPUs, better performance might be obtained by setting the
environment variable \sphinxcode{\sphinxupquote{I\_MPI\_SHM\_LMT=shm}} for Linux kernels that do
not yet have full support for AVX\sphinxhyphen{}512. Runs on Intel Xeon Phi x200
series processors will always perform better using MCDRAM. Please
consult your system documentation for the best approach to specify
that MPI runs are performed in MCDRAM.


\subsubsection{Tuning for Offload Performance}
\label{\detokenize{Speed_intel:tuning-for-offload-performance}}
\sphinxAtStartPar
The default settings for offload should give good performance.

\sphinxAtStartPar
When using LAMMPS with offload to Intel co\sphinxhyphen{}processors, best performance
will typically be achieved with concurrent calculations performed on
both the CPU and the co\sphinxhyphen{}processor. This is achieved by offloading only
a fraction of the neighbor and pair computations to the co\sphinxhyphen{}processor or
using {\hyperref[\detokenize{pair_hybrid::doc}]{\sphinxcrossref{\DUrole{doc}{hybrid}}}} pair styles where only one style uses
the “intel” suffix. For simulations with long\sphinxhyphen{}range electrostatics or
bond, angle, dihedral, improper calculations, computation and data
transfer to the co\sphinxhyphen{}processor will run concurrently with computations
and MPI communications for these calculations on the host CPU. This
is illustrated in the figure below for the rhodopsin protein benchmark
running on E5\sphinxhyphen{}2697v2 processors with a Intel Xeon Phi 7120p
co\sphinxhyphen{}processor. In this plot, the vertical access is time and routines
running at the same time are running concurrently on both the host and
the co\sphinxhyphen{}processor.

\noindent{\hspace*{\fill}\sphinxincludegraphics{{offload_knc}.png}\hspace*{\fill}}

\sphinxAtStartPar
The fraction of the offloaded work is controlled by the \sphinxstyleemphasis{balance}
keyword in the {\hyperref[\detokenize{package::doc}]{\sphinxcrossref{\DUrole{doc}{package intel}}}} command. A balance of 0
runs all calculations on the CPU.  A balance of 1 runs all
supported calculations on the co\sphinxhyphen{}processor.  A balance of 0.5 runs half
of the calculations on the co\sphinxhyphen{}processor.  Setting the balance to \sphinxhyphen{}1
(the default) will enable dynamic load balancing that continuously
adjusts the fraction of offloaded work throughout the simulation.
Because data transfer cannot be timed, this option typically produces
results within 5 to 10 percent of the optimal fixed balance.

\sphinxAtStartPar
If running short benchmark runs with dynamic load balancing, adding a
short warm\sphinxhyphen{}up run (10\sphinxhyphen{}20 steps) will allow the load\sphinxhyphen{}balancer to find a
near\sphinxhyphen{}optimal setting that will carry over to additional runs.

\sphinxAtStartPar
The default for the {\hyperref[\detokenize{package::doc}]{\sphinxcrossref{\DUrole{doc}{package intel}}}} command is to have
all the MPI tasks on a given compute node use a single Xeon Phi
co\sphinxhyphen{}processor.  In general, running with a large number of MPI tasks on
each node will perform best with offload.  Each MPI task will
automatically get affinity to a subset of the hardware threads
available on the co\sphinxhyphen{}processor.  For example, if your card has 61 cores,
with 60 cores available for offload and 4 hardware threads per core
(240 total threads), running with 24 MPI tasks per node will cause
each MPI task to use a subset of 10 threads on the co\sphinxhyphen{}processor.  Fine
tuning of the number of threads to use per MPI task or the number of
threads to use per core can be accomplished with keyword settings of
the {\hyperref[\detokenize{package::doc}]{\sphinxcrossref{\DUrole{doc}{package intel}}}} command.

\sphinxAtStartPar
The INTEL package has two modes for deciding which atoms will be
handled by the co\sphinxhyphen{}processor.  This choice is controlled with the \sphinxstyleemphasis{ghost}
keyword of the {\hyperref[\detokenize{package::doc}]{\sphinxcrossref{\DUrole{doc}{package intel}}}} command.  When set to 0,
ghost atoms (atoms at the borders between MPI tasks) are not offloaded
to the card.  This allows for overlap of MPI communication of forces
with computation on the co\sphinxhyphen{}processor when the {\hyperref[\detokenize{newton::doc}]{\sphinxcrossref{\DUrole{doc}{newton}}}}
setting is “on”.  The default is dependent on the style being used,
however, better performance may be achieved by setting this option
explicitly.

\sphinxAtStartPar
When using offload with CPU Hyper\sphinxhyphen{}Threading disabled, it may help
performance to use fewer MPI tasks and OpenMP threads than available
cores.  This is due to the fact that additional threads are generated
internally to handle the asynchronous offload tasks.

\sphinxAtStartPar
If pair computations are being offloaded to an Intel Xeon Phi
co\sphinxhyphen{}processor, a diagnostic line is printed to the screen (not to the
log file), during the setup phase of a run, indicating that offload
mode is being used and indicating the number of co\sphinxhyphen{}processor threads
per MPI task.  Additionally, an offload timing summary is printed at
the end of each run.  When offloading, the frequency for {\hyperref[\detokenize{atom_modify::doc}]{\sphinxcrossref{\DUrole{doc}{atom sorting}}}} is changed to 1 so that the per\sphinxhyphen{}atom data is
effectively sorted at every rebuild of the neighbor lists. All the
available co\sphinxhyphen{}processor threads on each Phi will be divided among MPI
tasks, unless the \sphinxcode{\sphinxupquote{tptask}} option of the \sphinxcode{\sphinxupquote{\sphinxhyphen{}pk intel}} {\hyperref[\detokenize{Run_options::doc}]{\sphinxcrossref{\DUrole{doc}{command\sphinxhyphen{}line switch}}}} is used to limit the co\sphinxhyphen{}processor threads per
MPI task.


\subsubsection{Restrictions}
\label{\detokenize{Speed_intel:restrictions}}
\sphinxAtStartPar
When offloading to a co\sphinxhyphen{}processor, {\hyperref[\detokenize{pair_hybrid::doc}]{\sphinxcrossref{\DUrole{doc}{hybrid}}}} styles
that require skip lists for neighbor builds cannot be offloaded.
Using {\hyperref[\detokenize{pair_hybrid::doc}]{\sphinxcrossref{\DUrole{doc}{hybrid/overlay}}}} is allowed.  Only one intel
accelerated style may be used with hybrid styles when offloading.
{\hyperref[\detokenize{special_bonds::doc}]{\sphinxcrossref{\DUrole{doc}{Special\_bonds}}}} exclusion lists are not currently
supported with offload, however, the same effect can often be
accomplished by setting cutoffs for excluded atom types to 0.  None of
the pair styles in the INTEL package currently support the
“inner”, “middle”, “outer” options for rRESPA integration via the
{\hyperref[\detokenize{run_style::doc}]{\sphinxcrossref{\DUrole{doc}{run\_style respa}}}} command; only the “pair” option is
supported.


\subsubsection{References}
\label{\detokenize{Speed_intel:references}}\begin{itemize}
\item {} 
\sphinxAtStartPar
Brown, W.M., Carrillo, J.\sphinxhyphen{}M.Y., Mishra, B., Gavhane, N., Thakkar, F.M., De Kraker, A.R., Yamada, M., Ang, J.A., Plimpton, S.J., “Optimizing Classical Molecular Dynamics in LAMMPS”, in Intel Xeon Phi Processor High Performance Programming: Knights Landing Edition, J. Jeffers, J. Reinders, A. Sodani, Eds. Morgan Kaufmann.

\item {} 
\sphinxAtStartPar
Brown, W. M., Semin, A., Hebenstreit, M., Khvostov, S., Raman, K., Plimpton, S.J. \sphinxhref{https://dl.acm.org/citation.cfm?id=3014915}{Increasing Molecular Dynamics Simulation Rates with an 8\sphinxhyphen{}Fold Increase in Electrical Power Efficiency.} 2016 High Performance Computing, Networking, Storage and Analysis, SC16: International Conference (pp. 82\sphinxhyphen{}95).

\item {} 
\sphinxAtStartPar
Brown, W.M., Carrillo, J.\sphinxhyphen{}M.Y., Gavhane, N., Thakkar, F.M., Plimpton, S.J. Optimizing Legacy Molecular Dynamics Software with Directive\sphinxhyphen{}Based Offload. Computer Physics Communications. 2015. 195: p. 95\sphinxhyphen{}101.

\end{itemize}

\sphinxstepscope


\subsection{KOKKOS package}
\label{\detokenize{Speed_kokkos:kokkos-package}}\label{\detokenize{Speed_kokkos::doc}}
\sphinxAtStartPar
Kokkos is a templated C++ library that provides abstractions to allow
a single implementation of an application kernel (e.g. a pair style)
to run efficiently on different kinds of hardware, such as GPUs, Intel
Xeon Phis, or many\sphinxhyphen{}core CPUs. Kokkos maps the C++ kernel onto
different back end languages such as CUDA, OpenMP, or Pthreads.  The
Kokkos library also provides data abstractions to adjust (at compile
time) the memory layout of data structures like 2d and 3d arrays to
optimize performance on different hardware. For more information on
Kokkos, see \sphinxhref{https://github.com/kokkos/kokkos}{the Kokkos GitHub page}.

\sphinxAtStartPar
The LAMMPS KOKKOS package contains versions of pair, fix, and atom
styles that use data structures and macros provided by the Kokkos
library, which is included with LAMMPS in /lib/kokkos. The KOKKOS
package was developed primarily by Christian Trott (Sandia) and Stan
Moore (Sandia) with contributions of various styles by others,
including Sikandar Mashayak (UIUC), Ray Shan (Sandia), and Dan Ibanez
(Sandia). For more information on developing using Kokkos abstractions
see the \sphinxhref{https://github.com/kokkos/kokkos/wiki}{Kokkos Wiki}.

\begin{sphinxadmonition}{note}{Note:}
\sphinxAtStartPar
The Kokkos library is under active development and tracking the
availability of accelerator hardware, so is the KOKKOS package in
LAMMPS.  This means that only a certain range of versions of the
Kokkos library are compatible with the KOKKOS package of a certain
range of LAMMPS versions.  For that reason LAMMPS comes with a
bundled version of the Kokkos library that has been validated on
multiple platforms and may contain selected back\sphinxhyphen{}ported bug fixes
from upstream Kokkos versions.  While it is possible to build LAMMPS
with an external version of Kokkos, it is untested and may result in
incorrect execution or crashes.
\end{sphinxadmonition}

\sphinxAtStartPar
Kokkos currently provides full support for 4 modes of execution (per MPI
task). These are Serial (MPI\sphinxhyphen{}only for CPUs and Intel Phi), OpenMP
(threading for many\sphinxhyphen{}core CPUs and Intel Phi), CUDA (for NVIDIA GPUs) and
HIP (for AMD GPUs).  Additional modes (e.g. OpenMP target, Intel data
center GPUs) are under development.  You choose the mode at build time
to produce an executable compatible with a specific hardware.

\sphinxAtStartPar
The following compatibility notes have been last updated for LAMMPS
version 23 November 2023 and Kokkos version 4.2.

\begin{sphinxadmonition}{note}{C++17 support}

\sphinxAtStartPar
Kokkos requires using a compiler that supports the c++17 standard. For
some compilers, it may be necessary to add a flag to enable c++17 support.
For example, the GNU compiler uses the \sphinxcode{\sphinxupquote{\sphinxhyphen{}std=c++17}} flag. For a list of
compilers that have been tested with the Kokkos library, see the
\sphinxhref{https://kokkos.github.io/kokkos-core-wiki/requirements.html}{requirements document of the Kokkos Wiki}.
\end{sphinxadmonition}

\begin{sphinxadmonition}{note}{NVIDIA CUDA support}

\sphinxAtStartPar
To build with Kokkos support for NVIDIA GPUs, the NVIDIA CUDA toolkit
software version 11.0 or later must be installed on your system. See
the discussion for the {\hyperref[\detokenize{Speed_gpu::doc}]{\sphinxcrossref{\DUrole{doc}{GPU package}}}} for details of
how to check and do this.
\end{sphinxadmonition}

\begin{sphinxadmonition}{note}{AMD ROCm (HIP) support}

\sphinxAtStartPar
To build with Kokkos support for AMD GPUs, the AMD ROCm toolkit
software version 5.2.0 or later must be installed on your system.
\end{sphinxadmonition}

\begin{sphinxadmonition}{note}{CUDA and MPI library compatibility}

\sphinxAtStartPar
Kokkos with CUDA currently implicitly assumes that the MPI library is
GPU\sphinxhyphen{}aware.  This is not always the case, especially when using
pre\sphinxhyphen{}compiled MPI libraries provided by a Linux distribution. This is
not a problem when using only a single GPU with a single MPI
rank.  When running with multiple MPI ranks, you may see segmentation
faults without GPU\sphinxhyphen{}aware MPI support. These can be avoided by adding
the flags {\hyperref[\detokenize{Run_options::doc}]{\sphinxcrossref{\DUrole{doc}{\sphinxhyphen{}pk kokkos gpu/aware off}}}} to the
LAMMPS command line or by using the command {\hyperref[\detokenize{package::doc}]{\sphinxcrossref{\DUrole{doc}{package kokkos
gpu/aware off}}}} in the input file.
\end{sphinxadmonition}

\begin{sphinxadmonition}{note}{Intel Data Center GPU support}

\sphinxAtStartPar
Support for Kokkos with Intel Data Center GPU accelerators (formerly
known under the code name “Ponte Vecchio”) in LAMMPS is still a work
in progress.  Only a subset of the functionality works correctly.
Please contact the LAMMPS developers if you run into problems.
\end{sphinxadmonition}


\subsubsection{Building LAMMPS with the KOKKOS package}
\label{\detokenize{Speed_kokkos:building-lammps-with-the-kokkos-package}}
\sphinxAtStartPar
See the {\hyperref[\detokenize{Build_extras:kokkos}]{\sphinxcrossref{\DUrole{std,std-ref}{Build extras}}}} page for instructions.


\subsubsection{Running LAMMPS with the KOKKOS package}
\label{\detokenize{Speed_kokkos:running-lammps-with-the-kokkos-package}}
\sphinxAtStartPar
All Kokkos operations occur within the context of an individual MPI task
running on a single node of the machine. The total number of MPI tasks
used by LAMMPS (one or multiple per compute node) is set in the usual
manner via the \sphinxcode{\sphinxupquote{mpirun}} or \sphinxcode{\sphinxupquote{mpiexec}} commands, and is independent of
Kokkos. E.g. the mpirun command in OpenMPI does this via its \sphinxcode{\sphinxupquote{\sphinxhyphen{}np}} and
\sphinxcode{\sphinxupquote{\sphinxhyphen{}npernode}} switches. Ditto for MPICH via \sphinxcode{\sphinxupquote{\sphinxhyphen{}np}} and \sphinxcode{\sphinxupquote{\sphinxhyphen{}ppn}}.


\paragraph{Running on a multicore CPU}
\label{\detokenize{Speed_kokkos:running-on-a-multicore-cpu}}
\sphinxAtStartPar
Here is a quick overview of how to use the KOKKOS package
for CPU acceleration, assuming one or more 16\sphinxhyphen{}core nodes.

\begin{sphinxVerbatim}[commandchars=\\\{\}]
\PYG{c+c1}{\PYGZsh{} 1 node, 16 MPI tasks/node, no multi\PYGZhy{}threading}
mpirun\PYG{+w}{ }\PYGZhy{}np\PYG{+w}{ }\PYG{l+m}{16}\PYG{+w}{ }lmp\PYGZus{}kokkos\PYGZus{}mpi\PYGZus{}only\PYG{+w}{ }\PYGZhy{}k\PYG{+w}{ }on\PYG{+w}{ }\PYGZhy{}sf\PYG{+w}{ }kk\PYG{+w}{ }\PYGZhy{}in\PYG{+w}{ }\PYG{k}{in}.lj

\PYG{c+c1}{\PYGZsh{} 2 nodes, 1 MPI task/node, 16 threads/task}
mpirun\PYG{+w}{ }\PYGZhy{}np\PYG{+w}{ }\PYG{l+m}{2}\PYG{+w}{ }\PYGZhy{}ppn\PYG{+w}{ }\PYG{l+m}{1}\PYG{+w}{ }lmp\PYGZus{}kokkos\PYGZus{}omp\PYG{+w}{ }\PYGZhy{}k\PYG{+w}{ }on\PYG{+w}{ }t\PYG{+w}{ }\PYG{l+m}{16}\PYG{+w}{ }\PYGZhy{}sf\PYG{+w}{ }kk\PYG{+w}{ }\PYGZhy{}in\PYG{+w}{ }\PYG{k}{in}.lj

\PYG{c+c1}{\PYGZsh{} 1 node,  2 MPI tasks/node, 8 threads/task}
mpirun\PYG{+w}{ }\PYGZhy{}np\PYG{+w}{ }\PYG{l+m}{2}\PYG{+w}{ }lmp\PYGZus{}kokkos\PYGZus{}omp\PYG{+w}{ }\PYGZhy{}k\PYG{+w}{ }on\PYG{+w}{ }t\PYG{+w}{ }\PYG{l+m}{8}\PYG{+w}{ }\PYGZhy{}sf\PYG{+w}{ }kk\PYG{+w}{ }\PYGZhy{}in\PYG{+w}{ }\PYG{k}{in}.lj

\PYG{c+c1}{\PYGZsh{} 8 nodes, 4 MPI tasks/node, 4 threads/task}
mpirun\PYG{+w}{ }\PYGZhy{}np\PYG{+w}{ }\PYG{l+m}{32}\PYG{+w}{ }\PYGZhy{}ppn\PYG{+w}{ }\PYG{l+m}{4}\PYG{+w}{ }lmp\PYGZus{}kokkos\PYGZus{}omp\PYG{+w}{ }\PYGZhy{}k\PYG{+w}{ }on\PYG{+w}{ }t\PYG{+w}{ }\PYG{l+m}{4}\PYG{+w}{ }\PYGZhy{}sf\PYG{+w}{ }kk\PYG{+w}{ }\PYGZhy{}in\PYG{+w}{ }\PYG{k}{in}.lj
\end{sphinxVerbatim}

\sphinxAtStartPar
To run using the KOKKOS package, use the \sphinxcode{\sphinxupquote{\sphinxhyphen{}k on}}, \sphinxcode{\sphinxupquote{\sphinxhyphen{}sf kk}} and \sphinxcode{\sphinxupquote{\sphinxhyphen{}pk
kokkos}} {\hyperref[\detokenize{Run_options::doc}]{\sphinxcrossref{\DUrole{doc}{command\sphinxhyphen{}line switches}}}} in your \sphinxcode{\sphinxupquote{mpirun}}
command.  You must use the \sphinxcode{\sphinxupquote{\sphinxhyphen{}k on}} {\hyperref[\detokenize{Run_options::doc}]{\sphinxcrossref{\DUrole{doc}{command\sphinxhyphen{}line switch}}}} to enable the KOKKOS package. It takes
additional arguments for hardware settings appropriate to your system.
For OpenMP use:

\begin{sphinxVerbatim}[commandchars=\\\{\}]
\PYGZhy{}k on t Nt
\end{sphinxVerbatim}

\sphinxAtStartPar
The \sphinxcode{\sphinxupquote{t Nt}} option specifies how many OpenMP threads per MPI task to
use with a node. The default is \sphinxcode{\sphinxupquote{Nt}} = 1, which is MPI\sphinxhyphen{}only mode.  Note
that the product of MPI tasks * OpenMP threads/task should not exceed
the physical number of cores (on a node), otherwise performance will
suffer. If Hyper\sphinxhyphen{}Threading (HT) is enabled, then the product of MPI
tasks * OpenMP threads/task should not exceed the physical number of
cores * hardware threads.  The \sphinxcode{\sphinxupquote{\sphinxhyphen{}k on}} switch also issues a
\sphinxcode{\sphinxupquote{package kokkos}} command (with no additional arguments) which sets
various KOKKOS options to default values, as discussed on the
{\hyperref[\detokenize{package::doc}]{\sphinxcrossref{\DUrole{doc}{package}}}} command doc page.

\sphinxAtStartPar
The \sphinxcode{\sphinxupquote{\sphinxhyphen{}sf kk}} {\hyperref[\detokenize{Run_options::doc}]{\sphinxcrossref{\DUrole{doc}{command\sphinxhyphen{}line switch}}}} will automatically
append the “/kk” suffix to styles that support it.  In this manner no
modification to the input script is needed. Alternatively, one can run
with the KOKKOS package by editing the input script as described
below.

\begin{sphinxadmonition}{note}{Note:}
\sphinxAtStartPar
When using a single OpenMP thread, the Kokkos Serial back end (i.e.
\sphinxcode{\sphinxupquote{Makefile.kokkos\_mpi\_only}}) will give better performance than the OpenMP
back end (i.e. \sphinxcode{\sphinxupquote{Makefile.kokkos\_omp}}) because some of the overhead to make
the code thread\sphinxhyphen{}safe is removed.
\end{sphinxadmonition}

\begin{sphinxadmonition}{note}{Note:}
\sphinxAtStartPar
Use the \sphinxcode{\sphinxupquote{\sphinxhyphen{}pk kokkos}} {\hyperref[\detokenize{Run_options::doc}]{\sphinxcrossref{\DUrole{doc}{command\sphinxhyphen{}line switch}}}} to
change the default {\hyperref[\detokenize{package::doc}]{\sphinxcrossref{\DUrole{doc}{package kokkos}}}} options. See its doc
page for details and default settings. Experimenting with its options
can provide a speed\sphinxhyphen{}up for specific calculations. For example:
\end{sphinxadmonition}

\begin{sphinxVerbatim}[commandchars=\\\{\}]
\PYG{c+c1}{\PYGZsh{} Newton on, Half neighbor list, non\PYGZhy{}threaded comm}
mpirun\PYG{+w}{ }\PYGZhy{}np\PYG{+w}{ }\PYG{l+m}{16}\PYG{+w}{ }lmp\PYGZus{}kokkos\PYGZus{}mpi\PYGZus{}only\PYG{+w}{ }\PYGZhy{}k\PYG{+w}{ }on\PYG{+w}{ }\PYGZhy{}sf\PYG{+w}{ }kk\PYG{+w}{ }\PYG{l+s+se}{\PYGZbs{}}
\PYG{+w}{       }\PYGZhy{}pk\PYG{+w}{ }kokkos\PYG{+w}{ }newton\PYG{+w}{ }on\PYG{+w}{ }neigh\PYG{+w}{ }half\PYG{+w}{ }comm\PYG{+w}{ }no\PYG{+w}{ }\PYGZhy{}in\PYG{+w}{ }\PYG{k}{in}.lj
\end{sphinxVerbatim}

\sphinxAtStartPar
If the {\hyperref[\detokenize{newton::doc}]{\sphinxcrossref{\DUrole{doc}{newton}}}} command is used in the input
script, it can also override the Newton flag defaults.

\sphinxAtStartPar
For half neighbor lists and OpenMP, the KOKKOS package uses data
duplication (i.e. thread\sphinxhyphen{}private arrays) by default to avoid
thread\sphinxhyphen{}level write conflicts in the force arrays (and other data
structures as necessary). Data duplication is typically fastest for
small numbers of threads (i.e. 8 or less) but does increase memory
footprint and is not scalable to large numbers of threads. An
alternative to data duplication is to use thread\sphinxhyphen{}level atomic operations
which do not require data duplication. The use of atomic operations can
be enforced by compiling LAMMPS with the \sphinxcode{\sphinxupquote{\sphinxhyphen{}DLMP\_KOKKOS\_USE\_ATOMICS}}
pre\sphinxhyphen{}processor flag. Most but not all Kokkos\sphinxhyphen{}enabled pair\_styles support
data duplication. Alternatively, full neighbor lists avoid the need for
duplication or atomic operations but require more compute operations per
atom.  When using the Kokkos Serial back end or the OpenMP back end with
a single thread, no duplication or atomic operations are used. For CUDA
and half neighbor lists, the KOKKOS package always uses atomic operations.


\paragraph{CPU Cores, Sockets and Thread Affinity}
\label{\detokenize{Speed_kokkos:cpu-cores-sockets-and-thread-affinity}}
\sphinxAtStartPar
When using multi\sphinxhyphen{}threading, it is important for performance to bind
both MPI tasks to physical cores, and threads to physical cores, so
they do not migrate during a simulation.

\sphinxAtStartPar
If you are not certain MPI tasks are being bound (check the defaults
for your MPI installation), binding can be forced with these flags:

\begin{sphinxVerbatim}[commandchars=\\\{\}]
\PYG{c+c1}{\PYGZsh{} OpenMPI 1.8}
mpirun\PYG{+w}{ }\PYGZhy{}np\PYG{+w}{ }\PYG{l+m}{2}\PYG{+w}{ }\PYGZhy{}\PYGZhy{}bind\PYGZhy{}to\PYG{+w}{ }socket\PYG{+w}{ }\PYGZhy{}\PYGZhy{}map\PYGZhy{}by\PYG{+w}{ }socket\PYG{+w}{ }./lmp\PYGZus{}openmpi\PYG{+w}{ }...

\PYG{c+c1}{\PYGZsh{} Mvapich2 2.0}
mpiexec\PYG{+w}{ }\PYGZhy{}np\PYG{+w}{ }\PYG{l+m}{2}\PYG{+w}{ }\PYGZhy{}\PYGZhy{}bind\PYGZhy{}to\PYG{+w}{ }socket\PYG{+w}{ }\PYGZhy{}\PYGZhy{}map\PYGZhy{}by\PYG{+w}{ }socket\PYG{+w}{ }./lmp\PYGZus{}mvapich\PYG{+w}{ }...
\end{sphinxVerbatim}

\sphinxAtStartPar
For binding threads with KOKKOS OpenMP, use thread affinity environment
variables to force binding. With OpenMP 3.1 (gcc 4.7 or later, intel 12
or later) setting the environment variable \sphinxcode{\sphinxupquote{OMP\_PROC\_BIND=true}} should
be sufficient. In general, for best performance with OpenMP 4.0 or later
set \sphinxcode{\sphinxupquote{OMP\_PROC\_BIND=spread}} and \sphinxcode{\sphinxupquote{OMP\_PLACES=threads}}.  For binding
threads with the KOKKOS pthreads option, compile LAMMPS with the hwloc
or libnuma support enabled as described in the {\hyperref[\detokenize{Build_extras:kokkos}]{\sphinxcrossref{\DUrole{std,std-ref}{extra build options page}}}}.


\paragraph{Running on Knight’s Landing (KNL) Intel Xeon Phi}
\label{\detokenize{Speed_kokkos:running-on-knight-s-landing-knl-intel-xeon-phi}}
\sphinxAtStartPar
Here is a quick overview of how to use the KOKKOS package for the
Intel Knight’s Landing (KNL) Xeon Phi:

\sphinxAtStartPar
KNL Intel Phi chips have 68 physical cores. Typically 1 to 4 cores are
reserved for the OS, and only 64 or 66 cores are used. Each core has 4
Hyper\sphinxhyphen{}Threads,so there are effectively N = 256 (4*64) or N = 264 (4*66)
cores to run on. The product of MPI tasks * OpenMP threads/task should
not exceed this limit, otherwise performance will suffer. Note that
with the KOKKOS package you do not need to specify how many KNLs there
are per node; each KNL is simply treated as running some number of MPI
tasks.

\sphinxAtStartPar
Examples of mpirun commands that follow these rules are shown below.

\begin{sphinxVerbatim}[commandchars=\\\{\}]
\PYG{c+c1}{\PYGZsh{} Running on an Intel KNL node with 68 cores}
\PYG{c+c1}{\PYGZsh{} (272 threads/node via 4x hardware threading):}

\PYG{c+c1}{\PYGZsh{} 1 node, 64 MPI tasks/node, 4 threads/task}
mpirun\PYG{+w}{ }\PYGZhy{}np\PYG{+w}{ }\PYG{l+m}{64}\PYG{+w}{ }lmp\PYGZus{}kokkos\PYGZus{}phi\PYG{+w}{ }\PYGZhy{}k\PYG{+w}{ }on\PYG{+w}{ }t\PYG{+w}{ }\PYG{l+m}{4}\PYG{+w}{ }\PYGZhy{}sf\PYG{+w}{ }kk\PYG{+w}{ }\PYGZhy{}in\PYG{+w}{ }\PYG{k}{in}.lj

\PYG{c+c1}{\PYGZsh{} 1 node, 66 MPI tasks/node, 4 threads/task}
mpirun\PYG{+w}{ }\PYGZhy{}np\PYG{+w}{ }\PYG{l+m}{66}\PYG{+w}{ }lmp\PYGZus{}kokkos\PYGZus{}phi\PYG{+w}{ }\PYGZhy{}k\PYG{+w}{ }on\PYG{+w}{ }t\PYG{+w}{ }\PYG{l+m}{4}\PYG{+w}{ }\PYGZhy{}sf\PYG{+w}{ }kk\PYG{+w}{ }\PYGZhy{}in\PYG{+w}{ }\PYG{k}{in}.lj

\PYG{c+c1}{\PYGZsh{} 1 node, 32 MPI tasks/node, 8 threads/task}
mpirun\PYG{+w}{ }\PYGZhy{}np\PYG{+w}{ }\PYG{l+m}{32}\PYG{+w}{ }lmp\PYGZus{}kokkos\PYGZus{}phi\PYG{+w}{ }\PYGZhy{}k\PYG{+w}{ }on\PYG{+w}{ }t\PYG{+w}{ }\PYG{l+m}{8}\PYG{+w}{ }\PYGZhy{}sf\PYG{+w}{ }kk\PYG{+w}{ }\PYGZhy{}in\PYG{+w}{ }\PYG{k}{in}.lj

\PYG{c+c1}{\PYGZsh{} 8 nodes, 64 MPI tasks/node, 4 threads/task}
mpirun\PYG{+w}{ }\PYGZhy{}np\PYG{+w}{ }\PYG{l+m}{512}\PYG{+w}{ }\PYGZhy{}ppn\PYG{+w}{ }\PYG{l+m}{64}\PYG{+w}{ }lmp\PYGZus{}kokkos\PYGZus{}phi\PYG{+w}{ }\PYGZhy{}k\PYG{+w}{ }on\PYG{+w}{ }t\PYG{+w}{ }\PYG{l+m}{4}\PYG{+w}{ }\PYGZhy{}sf\PYG{+w}{ }kk\PYG{+w}{ }\PYGZhy{}in\PYG{+w}{ }\PYG{k}{in}.lj
\end{sphinxVerbatim}

\sphinxAtStartPar
The \sphinxcode{\sphinxupquote{\sphinxhyphen{}np}} setting of the mpirun command sets the number of MPI
tasks/node. The \sphinxcode{\sphinxupquote{\sphinxhyphen{}k on t Nt}} command\sphinxhyphen{}line switch sets the number of
threads/task as \sphinxcode{\sphinxupquote{Nt}}. The product of these two values should be N, i.e.
256 or 264.

\begin{sphinxadmonition}{note}{Note:}
\sphinxAtStartPar
The default for the {\hyperref[\detokenize{package::doc}]{\sphinxcrossref{\DUrole{doc}{package kokkos}}}} command when
running on KNL is to use “half” neighbor lists and set the Newton
flag to “on” for both pairwise and bonded interactions. This will
typically be best for many\sphinxhyphen{}body potentials. For simpler pairwise
potentials, it may be faster to use a “full” neighbor list with
Newton flag to “off”.  Use the \sphinxcode{\sphinxupquote{\sphinxhyphen{}pk kokkos}} {\hyperref[\detokenize{Run_options::doc}]{\sphinxcrossref{\DUrole{doc}{command\sphinxhyphen{}line switch}}}} to change the default {\hyperref[\detokenize{package::doc}]{\sphinxcrossref{\DUrole{doc}{package kokkos}}}}
options. See its documentation page for details and default
settings. Experimenting with its options can provide a speed\sphinxhyphen{}up for
specific calculations. For example:
\end{sphinxadmonition}

\begin{sphinxVerbatim}[commandchars=\\\{\}]
\PYG{c+c1}{\PYGZsh{}  Newton on, half neighbor list, threaded comm}
mpirun\PYG{+w}{ }\PYGZhy{}np\PYG{+w}{ }\PYG{l+m}{64}\PYG{+w}{ }lmp\PYGZus{}kokkos\PYGZus{}phi\PYG{+w}{ }\PYGZhy{}k\PYG{+w}{ }on\PYG{+w}{ }t\PYG{+w}{ }\PYG{l+m}{4}\PYG{+w}{ }\PYGZhy{}sf\PYG{+w}{ }kk\PYG{+w}{ }\PYGZhy{}pk\PYG{+w}{ }kokkos\PYG{+w}{ }comm\PYG{+w}{ }host\PYG{+w}{ }\PYGZhy{}in\PYG{+w}{ }\PYG{k}{in}.reax

\PYG{c+c1}{\PYGZsh{} Newton off, full neighbor list, non\PYGZhy{}threaded comm}
mpirun\PYG{+w}{ }\PYGZhy{}np\PYG{+w}{ }\PYG{l+m}{64}\PYG{+w}{ }lmp\PYGZus{}kokkos\PYGZus{}phi\PYG{+w}{ }\PYGZhy{}k\PYG{+w}{ }on\PYG{+w}{ }t\PYG{+w}{ }\PYG{l+m}{4}\PYG{+w}{ }\PYGZhy{}sf\PYG{+w}{ }kk\PYG{+w}{ }\PYG{l+s+se}{\PYGZbs{}}
\PYG{+w}{       }\PYGZhy{}pk\PYG{+w}{ }kokkos\PYG{+w}{ }newton\PYG{+w}{ }off\PYG{+w}{ }neigh\PYG{+w}{ }full\PYG{+w}{ }comm\PYG{+w}{ }no\PYG{+w}{ }\PYGZhy{}in\PYG{+w}{ }\PYG{k}{in}.lj
\end{sphinxVerbatim}

\begin{sphinxadmonition}{note}{Note:}
\sphinxAtStartPar
MPI tasks and threads should be bound to cores as described
above for CPUs.
\end{sphinxadmonition}

\begin{sphinxadmonition}{note}{Note:}
\sphinxAtStartPar
To build with Kokkos support for Intel Xeon Phi co\sphinxhyphen{}processors
such as Knight’s Corner (KNC), your system must be configured to use
them in “native” mode, not “offload” mode like the INTEL package
supports.
\end{sphinxadmonition}


\paragraph{Running on GPUs}
\label{\detokenize{Speed_kokkos:running-on-gpus}}
\sphinxAtStartPar
Use the \sphinxcode{\sphinxupquote{\sphinxhyphen{}k}} {\hyperref[\detokenize{Run_options::doc}]{\sphinxcrossref{\DUrole{doc}{command\sphinxhyphen{}line switch}}}} to specify the
number of GPUs per node. Typically the \sphinxcode{\sphinxupquote{\sphinxhyphen{}np}} setting of the \sphinxcode{\sphinxupquote{mpirun}} command
should set the number of MPI tasks/node to be equal to the number of
physical GPUs on the node. You can assign multiple MPI tasks to the same
GPU with the KOKKOS package, but this is usually only faster if some
portions of the input script have not been ported to use Kokkos. In this
case, also packing/unpacking communication buffers on the host may give
speedup (see the KOKKOS {\hyperref[\detokenize{package::doc}]{\sphinxcrossref{\DUrole{doc}{package}}}} command). Using CUDA MPS
is recommended in this scenario.

\sphinxAtStartPar
Using a GPU\sphinxhyphen{}aware MPI library is highly recommended. GPU\sphinxhyphen{}aware MPI use can be
avoided by using {\hyperref[\detokenize{package::doc}]{\sphinxcrossref{\DUrole{doc}{\sphinxhyphen{}pk kokkos gpu/aware off}}}}. As above for
multicore CPUs (and no GPU), if N is the number of physical cores/node,
then the number of MPI tasks/node should not exceed N.

\begin{sphinxVerbatim}[commandchars=\\\{\}]
\PYGZhy{}k on g Ng
\end{sphinxVerbatim}

\sphinxAtStartPar
Here are examples of how to use the KOKKOS package for GPUs, assuming
one or more nodes, each with two GPUs:

\begin{sphinxVerbatim}[commandchars=\\\{\}]
\PYG{c+c1}{\PYGZsh{} 1 node,   2 MPI tasks/node, 2 GPUs/node}
mpirun\PYG{+w}{ }\PYGZhy{}np\PYG{+w}{ }\PYG{l+m}{2}\PYG{+w}{ }lmp\PYGZus{}kokkos\PYGZus{}cuda\PYGZus{}openmpi\PYG{+w}{ }\PYGZhy{}k\PYG{+w}{ }on\PYG{+w}{ }g\PYG{+w}{ }\PYG{l+m}{2}\PYG{+w}{ }\PYGZhy{}sf\PYG{+w}{ }kk\PYG{+w}{ }\PYGZhy{}in\PYG{+w}{ }\PYG{k}{in}.lj

\PYG{c+c1}{\PYGZsh{} 16 nodes, 2 MPI tasks/node, 2 GPUs/node (32 GPUs total)}
mpirun\PYG{+w}{ }\PYGZhy{}np\PYG{+w}{ }\PYG{l+m}{32}\PYG{+w}{ }\PYGZhy{}ppn\PYG{+w}{ }\PYG{l+m}{2}\PYG{+w}{ }lmp\PYGZus{}kokkos\PYGZus{}cuda\PYGZus{}openmpi\PYG{+w}{ }\PYGZhy{}k\PYG{+w}{ }on\PYG{+w}{ }g\PYG{+w}{ }\PYG{l+m}{2}\PYG{+w}{ }\PYGZhy{}sf\PYG{+w}{ }kk\PYG{+w}{ }\PYGZhy{}in\PYG{+w}{ }\PYG{k}{in}.lj
\end{sphinxVerbatim}

\begin{sphinxadmonition}{note}{Note:}
\sphinxAtStartPar
The default for the {\hyperref[\detokenize{package::doc}]{\sphinxcrossref{\DUrole{doc}{package kokkos}}}} command when
running on GPUs is to use “full” neighbor lists and set the Newton
flag to “off” for both pairwise and bonded interactions, along with
threaded communication. When running on Maxwell or Kepler GPUs, this
will typically be best. For Pascal GPUs and beyond, using “half”
neighbor lists and setting the Newton flag to “on” may be faster. For
many pair styles, setting the neighbor binsize equal to twice the CPU
default value will give speedup, which is the default when running on
GPUs. Use the \sphinxcode{\sphinxupquote{\sphinxhyphen{}pk kokkos}} {\hyperref[\detokenize{Run_options::doc}]{\sphinxcrossref{\DUrole{doc}{command\sphinxhyphen{}line switch}}}}
to change the default {\hyperref[\detokenize{package::doc}]{\sphinxcrossref{\DUrole{doc}{package kokkos}}}} options. See
its documentation page for details and default
settings. Experimenting with its options can provide a speed\sphinxhyphen{}up for
specific calculations. For example:
\end{sphinxadmonition}

\begin{sphinxVerbatim}[commandchars=\\\{\}]
\PYG{c+c1}{\PYGZsh{} Newton on, half neighbor list, set binsize = neighbor ghost cutoff}
mpirun\PYG{+w}{ }\PYGZhy{}np\PYG{+w}{ }\PYG{l+m}{2}\PYG{+w}{ }lmp\PYGZus{}kokkos\PYGZus{}cuda\PYGZus{}openmpi\PYG{+w}{ }\PYGZhy{}k\PYG{+w}{ }on\PYG{+w}{ }g\PYG{+w}{ }\PYG{l+m}{2}\PYG{+w}{ }\PYGZhy{}sf\PYG{+w}{ }kk\PYG{+w}{ }\PYG{l+s+se}{\PYGZbs{}}
\PYG{+w}{       }\PYGZhy{}pk\PYG{+w}{ }kokkos\PYG{+w}{ }newton\PYG{+w}{ }on\PYG{+w}{ }neigh\PYG{+w}{ }half\PYG{+w}{ }binsize\PYG{+w}{ }\PYG{l+m}{2}.8\PYG{+w}{ }\PYGZhy{}in\PYG{+w}{ }\PYG{k}{in}.lj
\end{sphinxVerbatim}

\begin{sphinxadmonition}{note}{Note:}
\sphinxAtStartPar
The default binsize for {\hyperref[\detokenize{atom_modify::doc}]{\sphinxcrossref{\DUrole{doc}{atom sorting}}}} on GPUs
is equal to the default CPU neighbor binsize (i.e. 2x smaller than the
default GPU neighbor binsize). When running simple pair\sphinxhyphen{}wise
potentials like Lennard Jones on GPUs, using a 2x larger binsize for
atom sorting (equal to the default GPU neighbor binsize) and a more
frequent sorting than default (e.g. sorting every 100 time steps
instead of 1000) may improve performance.
\end{sphinxadmonition}

\begin{sphinxadmonition}{note}{Note:}
\sphinxAtStartPar
When running on GPUs with many MPI ranks (tens of thousands and
more), the creation of the atom map (required for molecular systems)
on the GPU can slow down significantly or run out of GPU memory and
thus slow down the whole calculation or cause a crash.  You can use
the \sphinxcode{\sphinxupquote{\sphinxhyphen{}pk kokkos atom/map no}} {\hyperref[\detokenize{Run_options::doc}]{\sphinxcrossref{\DUrole{doc}{command\sphinxhyphen{}line switch}}}}
of the {\hyperref[\detokenize{package::doc}]{\sphinxcrossref{\DUrole{doc}{package kokkos atom/map no}}}} command to create
the atom map on the CPU instead.
\end{sphinxadmonition}

\begin{sphinxadmonition}{note}{Note:}
\sphinxAtStartPar
When using a GPU, you will achieve the best performance if your
input script does not use fix or compute styles which are not yet
Kokkos\sphinxhyphen{}enabled. This allows data to stay on the GPU for multiple
timesteps, without being copied back to the host CPU. Invoking a
non\sphinxhyphen{}Kokkos fix or compute, or performing I/O for
{\hyperref[\detokenize{thermo_style::doc}]{\sphinxcrossref{\DUrole{doc}{thermo}}}} or {\hyperref[\detokenize{dump::doc}]{\sphinxcrossref{\DUrole{doc}{dump}}}} output will cause data
to be copied back to the CPU incurring a performance penalty.
\end{sphinxadmonition}

\begin{sphinxadmonition}{note}{Note:}
\sphinxAtStartPar
To get an accurate timing breakdown between time spend in pair,
kspace, etc., you must set the environment variable \sphinxcode{\sphinxupquote{CUDA\_LAUNCH\_BLOCKING=1}}.
However, this will reduce performance and is not recommended for production runs.
\end{sphinxadmonition}


\paragraph{Run with the KOKKOS package by editing an input script}
\label{\detokenize{Speed_kokkos:run-with-the-kokkos-package-by-editing-an-input-script}}
\sphinxAtStartPar
Alternatively the effect of the \sphinxcode{\sphinxupquote{\sphinxhyphen{}sf}} or \sphinxcode{\sphinxupquote{\sphinxhyphen{}pk}} switches can be
duplicated by adding the {\hyperref[\detokenize{package::doc}]{\sphinxcrossref{\DUrole{doc}{package kokkos}}}} or {\hyperref[\detokenize{suffix::doc}]{\sphinxcrossref{\DUrole{doc}{suffix kk}}}} commands to your input script.

\sphinxAtStartPar
The discussion above for building LAMMPS with the KOKKOS package, the
\sphinxcode{\sphinxupquote{mpirun}} or \sphinxcode{\sphinxupquote{mpiexec}} command, and setting appropriate thread
properties are the same.

\sphinxAtStartPar
You must still use the \sphinxcode{\sphinxupquote{\sphinxhyphen{}k on}} {\hyperref[\detokenize{Run_options::doc}]{\sphinxcrossref{\DUrole{doc}{command\sphinxhyphen{}line switch}}}}
to enable the KOKKOS package, and specify its additional arguments for
hardware options appropriate to your system, as documented above.

\sphinxAtStartPar
You can use the {\hyperref[\detokenize{suffix::doc}]{\sphinxcrossref{\DUrole{doc}{suffix kk}}}} command, or you can explicitly add a
“kk” suffix to individual styles in your input script, e.g.

\begin{sphinxVerbatim}[commandchars=\\\{\}]
\PYG{k}{pair\PYGZus{}style}\PYG{+w}{ }\PYG{n}{lj}\PYG{o}{/}\PYG{n}{cut}\PYG{o}{/}\PYG{n}{kk}\PYG{+w}{ }\PYG{l+m}{2.5}
\end{sphinxVerbatim}

\sphinxAtStartPar
You only need to use the {\hyperref[\detokenize{package::doc}]{\sphinxcrossref{\DUrole{doc}{package kokkos}}}} command if you
wish to change any of its option defaults, as set by the “\sphinxhyphen{}k on”
{\hyperref[\detokenize{Run_options::doc}]{\sphinxcrossref{\DUrole{doc}{command\sphinxhyphen{}line switch}}}}.

\sphinxAtStartPar
\sphinxstylestrong{Using OpenMP threading and CUDA together:}

\sphinxAtStartPar
With the KOKKOS package, both OpenMP multi\sphinxhyphen{}threading and GPUs can be
compiled and used together in a few special cases. In the makefile for
the conventional build, the \sphinxcode{\sphinxupquote{KOKKOS\_DEVICES}} variable must include both,
“Cuda” and “OpenMP”, as is the case for \sphinxcode{\sphinxupquote{/src/MAKE/OPTIONS/Makefile.kokkos\_cuda\_mpi}}.

\begin{sphinxVerbatim}[commandchars=\\\{\}]
\PYG{n+nv}{KOKKOS\PYGZus{}DEVICES}\PYG{o}{=}Cuda,OpenMP
\end{sphinxVerbatim}

\sphinxAtStartPar
When building with CMake you need to enable both features as it is done
in the \sphinxcode{\sphinxupquote{kokkos\sphinxhyphen{}cuda.cmake}} CMake preset file.

\begin{sphinxVerbatim}[commandchars=\\\{\}]
cmake\PYG{+w}{ }\PYGZhy{}DKokkos\PYGZus{}ENABLE\PYGZus{}CUDA\PYG{o}{=}yes\PYG{+w}{ }\PYGZhy{}DKokkos\PYGZus{}ENABLE\PYGZus{}OPENMP\PYG{o}{=}yes\PYG{+w}{ }../cmake
\end{sphinxVerbatim}

\sphinxAtStartPar
The suffix “/kk” is equivalent to “/kk/device”, and for Kokkos CUDA,
using the \sphinxcode{\sphinxupquote{\sphinxhyphen{}sf kk}} in the command line gives the default CUDA version
everywhere.  However, if the “/kk/host” suffix is added to a specific
style in the input script, the Kokkos OpenMP (CPU) version of that
specific style will be used instead.  Set the number of OpenMP threads
as \sphinxcode{\sphinxupquote{t Nt}} and the number of GPUs as \sphinxcode{\sphinxupquote{g Ng}}

\begin{sphinxVerbatim}[commandchars=\\\{\}]
\PYGZhy{}k on t Nt g Ng
\end{sphinxVerbatim}

\sphinxAtStartPar
For example, the command to run with 1 GPU and 8 OpenMP threads is then:

\begin{sphinxVerbatim}[commandchars=\\\{\}]
mpiexec\PYG{+w}{ }\PYGZhy{}np\PYG{+w}{ }\PYG{l+m}{1}\PYG{+w}{ }lmp\PYGZus{}kokkos\PYGZus{}cuda\PYGZus{}openmpi\PYG{+w}{ }\PYGZhy{}in\PYG{+w}{ }\PYG{k}{in}.lj\PYG{+w}{ }\PYGZhy{}k\PYG{+w}{ }on\PYG{+w}{ }g\PYG{+w}{ }\PYG{l+m}{1}\PYG{+w}{ }t\PYG{+w}{ }\PYG{l+m}{8}\PYG{+w}{ }\PYGZhy{}sf\PYG{+w}{ }kk
\end{sphinxVerbatim}

\sphinxAtStartPar
Conversely, if the \sphinxcode{\sphinxupquote{\sphinxhyphen{}sf kk/host}} is used in the command line and then
the “/kk” or “/kk/device” suffix is added to a specific style in your
input script, then only that specific style will run on the GPU while
everything else will run on the CPU in OpenMP mode. Note that the
execution of the CPU and GPU styles will NOT overlap, except for a
special case:

\sphinxAtStartPar
A kspace style and/or molecular topology (bonds, angles, etc.) running
on the host CPU can overlap with a pair style running on the
GPU. First compile with \sphinxcode{\sphinxupquote{\sphinxhyphen{}\sphinxhyphen{}default\sphinxhyphen{}stream per\sphinxhyphen{}thread}} added to \sphinxcode{\sphinxupquote{CCFLAGS}}
in the Kokkos CUDA Makefile.  Then explicitly use the “/kk/host”
suffix for kspace and bonds, angles, etc.  in the input file and the
“kk” suffix (equal to “kk/device”) on the command line.  Also make
sure the environment variable \sphinxcode{\sphinxupquote{CUDA\_LAUNCH\_BLOCKING}} is not set to “1”
so CPU/GPU overlap can occur.


\subsubsection{Performance to expect}
\label{\detokenize{Speed_kokkos:performance-to-expect}}
\sphinxAtStartPar
The performance of KOKKOS running in different modes is a function of
your hardware, which KOKKOS\sphinxhyphen{}enable styles are used, and the problem
size.

\sphinxAtStartPar
Generally speaking, the following rules of thumb apply:
\begin{itemize}
\item {} 
\sphinxAtStartPar
When running on CPUs only, with a single thread per MPI task,
performance of a KOKKOS style is somewhere between the standard
(un\sphinxhyphen{}accelerated) styles (MPI\sphinxhyphen{}only mode), and those provided by the
OPENMP package. However the difference between all 3 is small (less
than 20\%).

\item {} 
\sphinxAtStartPar
When running on CPUs only, with multiple threads per MPI task,
performance of a KOKKOS style is a bit slower than the OPENMP
package.

\item {} 
\sphinxAtStartPar
When running large number of atoms per GPU, KOKKOS is typically faster
than the GPU package when compiled for double precision.  The benefit
of using single or mixed precision with the GPU package depends
significantly on the hardware in use and the simulated system and pair
style.

\item {} 
\sphinxAtStartPar
When running on Intel Phi hardware, KOKKOS is not as fast as
the INTEL package, which is optimized for x86 hardware (not just
from Intel) and compilation with the Intel compilers.  The INTEL
package also can increase the vector length of vector instructions
by switching to single or mixed precision mode.

\item {} 
\sphinxAtStartPar
The KOKKOS package by default assumes that you are using exactly one
MPI rank per GPU. When trying to use multiple MPI ranks per GPU it is
mandatory to enable \sphinxhref{https://docs.nvidia.com/deploy/mps/index.html}{CUDA Multi\sphinxhyphen{}Process Service (MPS)} to get good
performance.  In this case it is better to not use all available
MPI ranks in order to avoid competing with the MPS daemon for
CPU resources.

\end{itemize}

\sphinxAtStartPar
See the \sphinxhref{https://www.lammps.org/bench.html}{Benchmark page} of the
LAMMPS website for performance of the KOKKOS package on different
hardware.


\subsubsection{Advanced Kokkos options}
\label{\detokenize{Speed_kokkos:advanced-kokkos-options}}
\sphinxAtStartPar
There are other allowed options when building with the KOKKOS package
that can improve performance or assist in debugging or profiling.
They are explained on the {\hyperref[\detokenize{Build_extras:kokkos}]{\sphinxcrossref{\DUrole{std,std-ref}{KOKKOS section of the build extras}}}} doc page,


\subsubsection{Restrictions}
\label{\detokenize{Speed_kokkos:restrictions}}
\sphinxAtStartPar
Currently, there are no precision options with the KOKKOS package. All
compilation and computation is performed in double precision.

\sphinxstepscope


\subsection{OPENMP package}
\label{\detokenize{Speed_omp:openmp-package}}\label{\detokenize{Speed_omp::doc}}
\sphinxAtStartPar
The OPENMP package was developed by Axel Kohlmeyer at Temple
University.  It provides optimized and multi\sphinxhyphen{}threaded versions
of many pair styles, nearly all bonded styles (bond, angle, dihedral,
improper), several Kspace styles, and a few fix styles.  It uses
the OpenMP interface for multi\sphinxhyphen{}threading, but can also be compiled
without OpenMP support, providing optimized serial styles in that case.


\subsubsection{Required hardware/software}
\label{\detokenize{Speed_omp:required-hardware-software}}
\sphinxAtStartPar
To enable multi\sphinxhyphen{}threading, your compiler must support the OpenMP interface.
You should have one or more multicore CPUs, as multiple threads can only be
launched by each MPI task on the local node (using shared memory).


\subsubsection{Building LAMMPS with the OPENMP package}
\label{\detokenize{Speed_omp:building-lammps-with-the-openmp-package}}
\sphinxAtStartPar
See the {\hyperref[\detokenize{Build_extras:openmp}]{\sphinxcrossref{\DUrole{std,std-ref}{Build extras}}}} page for
instructions.


\subsubsection{Run with the OPENMP package from the command line}
\label{\detokenize{Speed_omp:run-with-the-openmp-package-from-the-command-line}}
\sphinxAtStartPar
These examples assume one or more 16\sphinxhyphen{}core nodes.

\begin{sphinxVerbatim}[commandchars=\\\{\}]
\PYG{c+c1}{\PYGZsh{} 1 MPI task, 16 threads according to OMP\PYGZus{}NUM\PYGZus{}THREADS}
env\PYG{+w}{ }\PYG{n+nv}{OMP\PYGZus{}NUM\PYGZus{}THREADS}\PYG{o}{=}\PYG{l+m}{16}\PYG{+w}{ }lmp\PYGZus{}omp\PYG{+w}{ }\PYGZhy{}sf\PYG{+w}{ }omp\PYG{+w}{ }\PYGZhy{}in\PYG{+w}{ }\PYG{k}{in}.script

\PYG{c+c1}{\PYGZsh{} 1 MPI task, no threads, optimized kernels}
lmp\PYGZus{}mpi\PYG{+w}{ }\PYGZhy{}sf\PYG{+w}{ }omp\PYG{+w}{ }\PYGZhy{}in\PYG{+w}{ }\PYG{k}{in}.script

\PYG{c+c1}{\PYGZsh{} 4 MPI tasks, 4 threads/task}
mpirun\PYG{+w}{ }\PYGZhy{}np\PYG{+w}{ }\PYG{l+m}{4}\PYG{+w}{ }lmp\PYGZus{}omp\PYG{+w}{ }\PYGZhy{}sf\PYG{+w}{ }omp\PYG{+w}{ }\PYGZhy{}pk\PYG{+w}{ }omp\PYG{+w}{ }\PYG{l+m}{4}\PYG{+w}{ }\PYGZhy{}in\PYG{+w}{ }\PYG{k}{in}.script

\PYG{c+c1}{\PYGZsh{} 8 nodes, 4 MPI tasks/node, 4 threads/task}
mpirun\PYG{+w}{ }\PYGZhy{}np\PYG{+w}{ }\PYG{l+m}{32}\PYG{+w}{ }\PYGZhy{}ppn\PYG{+w}{ }\PYG{l+m}{4}\PYG{+w}{ }lmp\PYGZus{}omp\PYG{+w}{ }\PYGZhy{}sf\PYG{+w}{ }omp\PYG{+w}{ }\PYGZhy{}pk\PYG{+w}{ }omp\PYG{+w}{ }\PYG{l+m}{4}\PYG{+w}{ }\PYGZhy{}in\PYG{+w}{ }\PYG{k}{in}.script
\end{sphinxVerbatim}

\sphinxAtStartPar
The \sphinxcode{\sphinxupquote{mpirun}} or \sphinxcode{\sphinxupquote{mpiexec}} command sets the total number of MPI tasks
used by LAMMPS (one or multiple per compute node) and the number of MPI
tasks used per node.  E.g. the mpirun command in MPICH does this via
its \sphinxcode{\sphinxupquote{\sphinxhyphen{}np}} and \sphinxcode{\sphinxupquote{\sphinxhyphen{}ppn}} switches.  Ditto for OpenMPI via \sphinxcode{\sphinxupquote{\sphinxhyphen{}np}} and \sphinxcode{\sphinxupquote{\sphinxhyphen{}npernode}}.

\sphinxAtStartPar
You need to choose how many OpenMP threads per MPI task will be used
by the OPENMP package.  Note that the product of MPI tasks *
threads/task should not exceed the physical number of cores (on a
node), otherwise performance will suffer.

\sphinxAtStartPar
As in the lines above, use the \sphinxcode{\sphinxupquote{\sphinxhyphen{}sf omp}} {\hyperref[\detokenize{Run_options::doc}]{\sphinxcrossref{\DUrole{doc}{command\sphinxhyphen{}line switch}}}}, which will automatically append “omp” to
styles that support it.  The \sphinxcode{\sphinxupquote{\sphinxhyphen{}sf omp}} switch also issues a default
{\hyperref[\detokenize{package::doc}]{\sphinxcrossref{\DUrole{doc}{package omp 0}}}} command, which will set the number of
threads per MPI task via the \sphinxcode{\sphinxupquote{OMP\_NUM\_THREADS}} environment variable.

\sphinxAtStartPar
You can also use the \sphinxcode{\sphinxupquote{\sphinxhyphen{}pk omp Nt}} {\hyperref[\detokenize{Run_options::doc}]{\sphinxcrossref{\DUrole{doc}{command\sphinxhyphen{}line switch}}}}, to explicitly set \sphinxcode{\sphinxupquote{Nt}} = \# of OpenMP threads
per MPI task to use, as well as additional options.  Its syntax is the
same as the {\hyperref[\detokenize{package::doc}]{\sphinxcrossref{\DUrole{doc}{package omp}}}} command whose page gives
details, including the default values used if it is not specified.  It
also gives more details on how to set the number of threads via the
\sphinxcode{\sphinxupquote{OMP\_NUM\_THREADS}} environment variable.


\subsubsection{Or run with the OPENMP package by editing an input script}
\label{\detokenize{Speed_omp:or-run-with-the-openmp-package-by-editing-an-input-script}}
\sphinxAtStartPar
The discussion above for the \sphinxcode{\sphinxupquote{mpirun}} or \sphinxcode{\sphinxupquote{mpiexec}} command, MPI
tasks/node, and threads/MPI task is the same.

\sphinxAtStartPar
Use the {\hyperref[\detokenize{suffix::doc}]{\sphinxcrossref{\DUrole{doc}{suffix omp}}}} command, or you can explicitly add an
“omp” suffix to individual styles in your input script, e.g.

\begin{sphinxVerbatim}[commandchars=\\\{\}]
\PYG{k}{pair\PYGZus{}style}\PYG{+w}{ }\PYG{n}{lj}\PYG{o}{/}\PYG{n}{cut}\PYG{o}{/}\PYG{n}{omp}\PYG{+w}{ }\PYG{l+m}{2.5}
\end{sphinxVerbatim}

\sphinxAtStartPar
You must also use the {\hyperref[\detokenize{package::doc}]{\sphinxcrossref{\DUrole{doc}{package omp}}}} command to enable the
OPENMP package.  When you do this you also specify how many threads
per MPI task to use.  The command page explains other options and
how to set the number of threads via the \sphinxcode{\sphinxupquote{OMP\_NUM\_THREADS}} environment
variable.


\subsubsection{Speed\sphinxhyphen{}up to expect}
\label{\detokenize{Speed_omp:speed-up-to-expect}}
\sphinxAtStartPar
Depending on which styles are accelerated, you should look for a
reduction in the “Pair time”, “Bond time”, “KSpace time”, and “Loop
time” values printed at the end of a run.

\sphinxAtStartPar
You may see a small performance advantage (5 to 20\%) when running a
OPENMP style (in serial or parallel) with a single thread per MPI
task, versus running standard LAMMPS with its standard un\sphinxhyphen{}accelerated
styles (in serial or all\sphinxhyphen{}MPI parallelization with 1 task/core).  This
is because many of the OPENMP styles contain similar optimizations
to those used in the OPT package, described in
{\hyperref[\detokenize{Speed_opt::doc}]{\sphinxcrossref{\DUrole{doc}{the OPT package}}}} doc page.

\sphinxAtStartPar
With multiple threads/task, the optimal choice of number of MPI
tasks/node and OpenMP threads/task can vary a lot and should always be
tested via benchmark runs for a specific simulation running on a
specific machine, paying attention to guidelines discussed in the next
subsection.

\sphinxAtStartPar
A description of the multi\sphinxhyphen{}threading strategy used in the OPENMP
package and some performance examples are
\sphinxhref{https://drive.google.com/file/d/1d1gLK6Ru6aPYB50Ld2tO10Li8zgPVNB8/view?usp=sharing}{presented here}.


\subsubsection{Guidelines for best performance}
\label{\detokenize{Speed_omp:guidelines-for-best-performance}}
\sphinxAtStartPar
For many problems on current generation CPUs, running the OPENMP
package with a single thread/task is faster than running with multiple
threads/task.  This is because the MPI parallelization in LAMMPS is
often more efficient than multi\sphinxhyphen{}threading as implemented in the
OPENMP package.  The parallel efficiency (in a threaded sense) also
varies for different OPENMP styles.

\sphinxAtStartPar
Using multiple threads/task can be more effective under the following
circumstances:
\begin{itemize}
\item {} 
\sphinxAtStartPar
Individual compute nodes have a significant number of CPU cores but
the CPU itself has limited memory bandwidth, e.g. for Intel Xeon 53xx
(Clovertown) and 54xx (Harpertown) quad\sphinxhyphen{}core processors.  Running one
MPI task per CPU core will result in significant performance
degradation, so that running with 4 or even only 2 MPI tasks per node
is faster.  Running in hybrid MPI+OpenMP mode will reduce the
inter\sphinxhyphen{}node communication bandwidth contention in the same way, but
offers an additional speedup by utilizing the otherwise idle CPU
cores.

\item {} 
\sphinxAtStartPar
The interconnect used for MPI communication does not provide
sufficient bandwidth for a large number of MPI tasks per node.  For
example, this applies to running over gigabit ethernet or on Cray XT4
or XT5 series supercomputers.  As in the aforementioned case, this
effect worsens when using an increasing number of nodes.

\item {} 
\sphinxAtStartPar
The system has a spatially inhomogeneous particle density which does
not map well to the {\hyperref[\detokenize{processors::doc}]{\sphinxcrossref{\DUrole{doc}{domain decomposition scheme}}}} or
{\hyperref[\detokenize{balance::doc}]{\sphinxcrossref{\DUrole{doc}{load\sphinxhyphen{}balancing}}}} options that LAMMPS provides.  This is
because multi\sphinxhyphen{}threading achieves parallelism over the number of
particles, not via their distribution in space.

\item {} 
\sphinxAtStartPar
A machine is being used in “capability mode”, i.e. near the point
where MPI parallelism is maxed out.  For example, this can happen when
using the {\hyperref[\detokenize{kspace_style::doc}]{\sphinxcrossref{\DUrole{doc}{PPPM solver}}}} for long\sphinxhyphen{}range
electrostatics on large numbers of nodes.  The scaling of the KSpace
calculation (see the {\hyperref[\detokenize{kspace_style::doc}]{\sphinxcrossref{\DUrole{doc}{kspace\_style}}}} command) becomes
the performance\sphinxhyphen{}limiting factor.  Using multi\sphinxhyphen{}threading allows less
MPI tasks to be invoked and can speed\sphinxhyphen{}up the long\sphinxhyphen{}range solver, while
increasing overall performance by parallelizing the pairwise and
bonded calculations via OpenMP.  Likewise additional speedup can be
sometimes be achieved by increasing the length of the Coulombic cutoff
and thus reducing the work done by the long\sphinxhyphen{}range solver.  Using the
{\hyperref[\detokenize{run_style::doc}]{\sphinxcrossref{\DUrole{doc}{run\_style verlet/split}}}} command, which is compatible
with the OPENMP package, is an alternative way to reduce the number
of MPI tasks assigned to the KSpace calculation.

\end{itemize}

\sphinxAtStartPar
Additional performance tips are as follows:
\begin{itemize}
\item {} 
\sphinxAtStartPar
The best parallel efficiency from \sphinxstyleemphasis{omp} styles is typically achieved
when there is at least one MPI task per physical CPU chip, i.e. socket
or die.

\item {} 
\sphinxAtStartPar
It is usually most efficient to restrict threading to a single
socket, i.e. use one or more MPI task per socket.

\item {} 
\sphinxAtStartPar
NOTE: By default, several current MPI implementations use a processor
affinity setting that restricts each MPI task to a single CPU core.
Using multi\sphinxhyphen{}threading in this mode will force all threads to share the
one core and thus is likely to be counterproductive.  Instead, binding
MPI tasks to a (multicore) socket, should solve this issue.

\end{itemize}


\subsubsection{Restrictions}
\label{\detokenize{Speed_omp:restrictions}}
\sphinxAtStartPar
None.

\sphinxstepscope


\subsection{OPT package}
\label{\detokenize{Speed_opt:opt-package}}\label{\detokenize{Speed_opt::doc}}
\sphinxAtStartPar
The OPT package was developed by James Fischer (High Performance
Technologies), David Richie, and Vincent Natoli (Stone Ridge
Technologies).  It contains a handful of pair styles whose compute()
methods were rewritten in C++ templated form to reduce the overhead
due to if tests and other conditional code.


\subsubsection{Required hardware/software}
\label{\detokenize{Speed_opt:required-hardware-software}}
\sphinxAtStartPar
Any hardware. Any compiler.


\subsubsection{Building LAMMPS with the OPT package}
\label{\detokenize{Speed_opt:building-lammps-with-the-opt-package}}
\sphinxAtStartPar
See the {\hyperref[\detokenize{Build_extras:opt}]{\sphinxcrossref{\DUrole{std,std-ref}{Build extras}}}} page for instructions.


\subsubsection{Run with the OPT package from the command line}
\label{\detokenize{Speed_opt:run-with-the-opt-package-from-the-command-line}}
\begin{sphinxVerbatim}[commandchars=\\\{\}]
lmp\PYGZus{}mpi\PYG{+w}{ }\PYGZhy{}sf\PYG{+w}{ }opt\PYG{+w}{ }\PYGZhy{}in\PYG{+w}{ }\PYG{k}{in}.script\PYG{+w}{                }\PYG{c+c1}{\PYGZsh{} run in serial}
mpirun\PYG{+w}{ }\PYGZhy{}np\PYG{+w}{ }\PYG{l+m}{4}\PYG{+w}{ }lmp\PYGZus{}mpi\PYG{+w}{ }\PYGZhy{}sf\PYG{+w}{ }opt\PYG{+w}{ }\PYGZhy{}in\PYG{+w}{ }\PYG{k}{in}.script\PYG{+w}{   }\PYG{c+c1}{\PYGZsh{} run in parallel}
\end{sphinxVerbatim}

\sphinxAtStartPar
Use the “\sphinxhyphen{}sf opt” {\hyperref[\detokenize{Run_options::doc}]{\sphinxcrossref{\DUrole{doc}{command\sphinxhyphen{}line switch}}}}, which will
automatically append “opt” to styles that support it.


\subsubsection{Or run with the OPT package by editing an input script}
\label{\detokenize{Speed_opt:or-run-with-the-opt-package-by-editing-an-input-script}}
\sphinxAtStartPar
Use the {\hyperref[\detokenize{suffix::doc}]{\sphinxcrossref{\DUrole{doc}{suffix opt}}}} command, or you can explicitly add an
“opt” suffix to individual styles in your input script, e.g.

\begin{sphinxVerbatim}[commandchars=\\\{\}]
\PYG{k}{pair\PYGZus{}style}\PYG{+w}{ }\PYG{n}{lj}\PYG{o}{/}\PYG{n}{cut}\PYG{o}{/}\PYG{n}{opt}\PYG{+w}{ }\PYG{l+m}{2.5}
\end{sphinxVerbatim}


\subsubsection{Speed\sphinxhyphen{}up to expect}
\label{\detokenize{Speed_opt:speed-up-to-expect}}
\sphinxAtStartPar
You should see a reduction in the “Pair time” value printed at the end
of a run.  On most machines for reasonable problem sizes, it will be a
5 to 20\% savings.


\subsubsection{Guidelines for best performance}
\label{\detokenize{Speed_opt:guidelines-for-best-performance}}
\sphinxAtStartPar
Just try out an OPT pair style to see how it performs.


\subsubsection{Restrictions}
\label{\detokenize{Speed_opt:restrictions}}
\sphinxAtStartPar
None.

\sphinxAtStartPar
Inverting this list, LAMMPS currently has acceleration support for
three kinds of hardware, via the listed packages:


\begin{savenotes}\sphinxattablestart
\sphinxthistablewithglobalstyle
\centering
\begin{tabulary}{\linewidth}[t]{TT}
\sphinxtoprule
\sphinxtableatstartofbodyhook
\sphinxAtStartPar
Many\sphinxhyphen{}core CPUs
&
\sphinxAtStartPar
{\hyperref[\detokenize{Speed_intel::doc}]{\sphinxcrossref{\DUrole{doc}{INTEL}}}}, {\hyperref[\detokenize{Speed_kokkos::doc}]{\sphinxcrossref{\DUrole{doc}{KOKKOS}}}}, {\hyperref[\detokenize{Speed_omp::doc}]{\sphinxcrossref{\DUrole{doc}{OPENMP}}}}, {\hyperref[\detokenize{Speed_opt::doc}]{\sphinxcrossref{\DUrole{doc}{OPT}}}} packages
\\
\sphinxhline
\sphinxAtStartPar
GPUs
&
\sphinxAtStartPar
{\hyperref[\detokenize{Speed_gpu::doc}]{\sphinxcrossref{\DUrole{doc}{GPU}}}}, {\hyperref[\detokenize{Speed_kokkos::doc}]{\sphinxcrossref{\DUrole{doc}{KOKKOS}}}} packages
\\
\sphinxhline
\sphinxAtStartPar
Intel Phi/AVX
&
\sphinxAtStartPar
{\hyperref[\detokenize{Speed_intel::doc}]{\sphinxcrossref{\DUrole{doc}{INTEL}}}}, {\hyperref[\detokenize{Speed_kokkos::doc}]{\sphinxcrossref{\DUrole{doc}{KOKKOS}}}} packages
\\
\sphinxbottomrule
\end{tabulary}
\sphinxtableafterendhook\par
\sphinxattableend\end{savenotes}

\sphinxAtStartPar
Which package is fastest for your hardware may depend on the size
problem you are running and what commands (accelerated and
non\sphinxhyphen{}accelerated) are invoked by your input script.  While these doc
pages include performance guidelines, there is no substitute for
trying out the different packages appropriate to your hardware.

\sphinxAtStartPar
Any accelerated style has the same name as the corresponding standard
style, except that a suffix is appended.  Otherwise, the syntax for
the command that uses the style is identical, their functionality is
the same, and the numerical results it produces should also be the
same, except for precision and round\sphinxhyphen{}off effects.

\sphinxAtStartPar
For example, all of these styles are accelerated variants of the
Lennard\sphinxhyphen{}Jones {\hyperref[\detokenize{pair_lj::doc}]{\sphinxcrossref{\DUrole{doc}{pair\_style lj/cut}}}}:
\begin{itemize}
\item {} 
\sphinxAtStartPar
{\hyperref[\detokenize{pair_lj::doc}]{\sphinxcrossref{\DUrole{doc}{pair\_style lj/cut/gpu}}}}

\item {} 
\sphinxAtStartPar
{\hyperref[\detokenize{pair_lj::doc}]{\sphinxcrossref{\DUrole{doc}{pair\_style lj/cut/intel}}}}

\item {} 
\sphinxAtStartPar
{\hyperref[\detokenize{pair_lj::doc}]{\sphinxcrossref{\DUrole{doc}{pair\_style lj/cut/kk}}}}

\item {} 
\sphinxAtStartPar
{\hyperref[\detokenize{pair_lj::doc}]{\sphinxcrossref{\DUrole{doc}{pair\_style lj/cut/omp}}}}

\item {} 
\sphinxAtStartPar
{\hyperref[\detokenize{pair_lj::doc}]{\sphinxcrossref{\DUrole{doc}{pair\_style lj/cut/opt}}}}

\end{itemize}

\sphinxAtStartPar
To see what accelerate styles are currently available for a particular
style, find the style name in the {\hyperref[\detokenize{Commands_all::doc}]{\sphinxcrossref{\DUrole{doc}{Commands}}}}
style pages (fix,compute,pair,etc) and see what suffixes are listed
(g,i,k,o,t) with it.  The doc pages for individual commands
(e.g. {\hyperref[\detokenize{pair_lj::doc}]{\sphinxcrossref{\DUrole{doc}{pair lj/cut}}}} or {\hyperref[\detokenize{fix_nve::doc}]{\sphinxcrossref{\DUrole{doc}{fix nve}}}}) also list
any accelerated variants available for that style.

\sphinxAtStartPar
To use an accelerator package in LAMMPS, and one or more of the styles
it provides, follow these general steps.  Details vary from package to
package and are explained in the individual accelerator doc pages,
listed above:


\begin{savenotes}\sphinxattablestart
\sphinxthistablewithglobalstyle
\centering
\begin{tabulary}{\linewidth}[t]{TT}
\sphinxtoprule
\sphinxtableatstartofbodyhook
\sphinxAtStartPar
build the accelerator library
&
\sphinxAtStartPar
only for GPU package
\\
\sphinxhline
\sphinxAtStartPar
install the accelerator package
&
\sphinxAtStartPar
\sphinxcode{\sphinxupquote{make yes\sphinxhyphen{}opt}}, \sphinxcode{\sphinxupquote{make yes\sphinxhyphen{}intel}}, etc
\\
\sphinxhline
\sphinxAtStartPar
add compile/link flags to \sphinxcode{\sphinxupquote{Makefile.machine}}
in \sphinxcode{\sphinxupquote{src/MAKE}}
&
\sphinxAtStartPar
only for INTEL, KOKKOS, OPENMP,
OPT packages
\\
\sphinxhline
\sphinxAtStartPar
re\sphinxhyphen{}build LAMMPS
&
\sphinxAtStartPar
\sphinxcode{\sphinxupquote{make machine}}
\\
\sphinxhline
\sphinxAtStartPar
prepare and test a regular LAMMPS simulation
&
\sphinxAtStartPar
\sphinxcode{\sphinxupquote{lmp\_machine \sphinxhyphen{}in in.script;}}
\sphinxcode{\sphinxupquote{mpirun \sphinxhyphen{}np 32 lmp\_machine \sphinxhyphen{}in in.script}}
\\
\sphinxhline
\sphinxAtStartPar
enable specific accelerator support via \sphinxcode{\sphinxupquote{\sphinxhyphen{}k on}}
{\hyperref[\detokenize{Run_options::doc}]{\sphinxcrossref{\DUrole{doc}{command\sphinxhyphen{}line switch}}}}
&
\sphinxAtStartPar
only needed for KOKKOS package
\\
\sphinxhline
\sphinxAtStartPar
set any needed options for the package via \sphinxcode{\sphinxupquote{\sphinxhyphen{}pk}}
{\hyperref[\detokenize{Run_options::doc}]{\sphinxcrossref{\DUrole{doc}{command\sphinxhyphen{}line switch}}}} or
{\hyperref[\detokenize{package::doc}]{\sphinxcrossref{\DUrole{doc}{package}}}} command
&
\sphinxAtStartPar
only if defaults need to be changed
\\
\sphinxhline
\sphinxAtStartPar
use accelerated styles in your input via \sphinxcode{\sphinxupquote{\sphinxhyphen{}sf}}
{\hyperref[\detokenize{Run_options::doc}]{\sphinxcrossref{\DUrole{doc}{command\sphinxhyphen{}line switch}}}} or
{\hyperref[\detokenize{suffix::doc}]{\sphinxcrossref{\DUrole{doc}{suffix}}}} command
&
\sphinxAtStartPar
\sphinxcode{\sphinxupquote{lmp\_machine \sphinxhyphen{}in in.script \sphinxhyphen{}sf gpu}}
\\
\sphinxbottomrule
\end{tabulary}
\sphinxtableafterendhook\par
\sphinxattableend\end{savenotes}

\sphinxAtStartPar
Note that the first 4 steps can be done as a single command with
suitable make command invocations. This is discussed on the
{\hyperref[\detokenize{Packages::doc}]{\sphinxcrossref{\DUrole{doc}{Packages}}}} doc pages, and its use is illustrated in the
individual accelerator sections.  Typically these steps only need to
be done once, to create an executable that uses one or more
accelerator packages.

\sphinxAtStartPar
The last 4 steps can all be done from the command\sphinxhyphen{}line when LAMMPS is
launched, without changing your input script, as illustrated in the
individual accelerator sections.  Or you can add
{\hyperref[\detokenize{package::doc}]{\sphinxcrossref{\DUrole{doc}{package}}}} and {\hyperref[\detokenize{suffix::doc}]{\sphinxcrossref{\DUrole{doc}{suffix}}}} commands to your input
script.

\begin{sphinxadmonition}{note}{Note:}
\sphinxAtStartPar
With a few exceptions, you can build a single LAMMPS executable
with all its accelerator packages installed.  Note however that the
INTEL and KOKKOS packages require you to choose one of their
hardware options when building for a specific platform.  I.e. CPU or
Phi option for the INTEL package.  Or the OpenMP, CUDA, HIP, SYCL,
or Phi option for the KOKKOS package.  Or the OpenCL, HIP, or CUDA
option for the GPU package.
\end{sphinxadmonition}

\sphinxAtStartPar
These are the exceptions.  You cannot build a single executable with:
\begin{itemize}
\item {} 
\sphinxAtStartPar
both the INTEL Phi and KOKKOS Phi options

\item {} 
\sphinxAtStartPar
the INTEL Phi or Kokkos Phi option, and the GPU package

\end{itemize}

\sphinxAtStartPar
As mentioned above, the \sphinxhref{https://www.lammps.org/bench.html}{Benchmark page} of the LAMMPS website gives
performance results for the various accelerator packages for several
of the standard LAMMPS benchmark problems, as a function of problem
size and number of compute nodes, on different hardware platforms.

\sphinxAtStartPar
Here is a brief summary of what the various packages provide.  Details
are in the individual accelerator sections.
\begin{itemize}
\item {} 
\sphinxAtStartPar
Styles with a “gpu” suffix are part of the GPU package and can be run
on Intel, NVIDIA, or AMD GPUs.  The speed\sphinxhyphen{}up on a GPU depends on a
variety of factors, discussed in the accelerator sections.

\item {} 
\sphinxAtStartPar
Styles with an “intel” suffix are part of the INTEL
package. These styles support vectorized single and mixed precision
calculations, in addition to full double precision.  In extreme cases,
this can provide speedups over 3.5x on CPUs.  The package also
supports acceleration in “offload” mode to Intel(R) Xeon Phi(TM)
co\sphinxhyphen{}processors.  This can result in additional speedup over 2x depending
on the hardware configuration.

\item {} 
\sphinxAtStartPar
Styles with a “kk” suffix are part of the KOKKOS package, and can be
run using OpenMP on multicore CPUs, on an NVIDIA or AMD GPU, or on an
Intel Xeon Phi in “native” mode.  The speed\sphinxhyphen{}up depends on a variety of
factors, as discussed on the KOKKOS accelerator page.

\item {} 
\sphinxAtStartPar
Styles with an “omp” suffix are part of the OPENMP package and allow
a pair\sphinxhyphen{}style to be run in multi\sphinxhyphen{}threaded mode using OpenMP.  This can
be useful on nodes with high\sphinxhyphen{}core counts when using less MPI processes
than cores is advantageous, e.g. when running with PPPM so that FFTs
are run on fewer MPI processors or when the many MPI tasks would
overload the available bandwidth for communication.

\item {} 
\sphinxAtStartPar
Styles with an “opt” suffix are part of the OPT package and typically
speed\sphinxhyphen{}up the pairwise calculations of your simulation by 5\sphinxhyphen{}25\% on a
CPU.

\end{itemize}

\sphinxAtStartPar
The individual accelerator package doc pages explain:
\begin{itemize}
\item {} 
\sphinxAtStartPar
what hardware and software the accelerated package requires

\item {} 
\sphinxAtStartPar
how to build LAMMPS with the accelerated package

\item {} 
\sphinxAtStartPar
how to run with the accelerated package either via command\sphinxhyphen{}line switches or modifying the input script

\item {} 
\sphinxAtStartPar
speed\sphinxhyphen{}ups to expect

\item {} 
\sphinxAtStartPar
guidelines for best performance

\item {} 
\sphinxAtStartPar
restrictions

\end{itemize}

\sphinxstepscope


\section{Comparison of various accelerator packages}
\label{\detokenize{Speed_compare:comparison-of-various-accelerator-packages}}\label{\detokenize{Speed_compare::doc}}
\sphinxAtStartPar
The next section compares and contrasts the various accelerator
options, since there are multiple ways to perform OpenMP threading,
run on GPUs, optimize for vector units on CPUs and run on Intel
Xeon Phi (co\sphinxhyphen{})processors.

\sphinxAtStartPar
All of these packages can accelerate a LAMMPS calculation taking
advantage of hardware features, but they do it in different ways
and acceleration is not always guaranteed.

\sphinxAtStartPar
As a consequence, for a particular simulation on specific hardware,
one package may be faster than the other.  We give some guidelines
below, but the best way to determine which package is faster for your
input script is to try multiple of them on your machine and experiment
with available performance tuning settings.  See the benchmarking
section below for examples where this has been done.

\sphinxAtStartPar
\sphinxstylestrong{Guidelines for using each package optimally:}
\begin{itemize}
\item {} 
\sphinxAtStartPar
Both, the GPU and the KOKKOS package allows you to assign multiple
MPI ranks (= CPU cores) to the same GPU. For the GPU package, this
can lead to a speedup through better utilization of the GPU (by
overlapping computation and data transfer) and more efficient
computation of the non\sphinxhyphen{}GPU accelerated parts of LAMMPS through MPI
parallelization, as all system data is maintained and updated on
the host. For KOKKOS, there is less to no benefit from this, due
to its different memory management model, which tries to retain
data on the GPU.

\item {} 
\sphinxAtStartPar
The GPU package moves per\sphinxhyphen{}atom data (coordinates, forces, and
(optionally) neighbor list data, if not computed on the GPU) between
the CPU and GPU at every timestep.  The KOKKOS/CUDA package only does
this on timesteps when a CPU calculation is required (e.g. to invoke
a fix or compute that is non\sphinxhyphen{}GPU\sphinxhyphen{}ized). Hence, if you can formulate
your input script to only use GPU\sphinxhyphen{}ized fixes and computes, and avoid
doing I/O too often (thermo output, dump file snapshots, restart files),
then the data transfer cost of the KOKKOS/CUDA package can be very low,
causing it to run faster than the GPU package.

\item {} 
\sphinxAtStartPar
The GPU package is often faster than the KOKKOS/CUDA package, when the
number of atoms per GPU is on the smaller side.  The crossover point,
in terms of atoms/GPU at which the KOKKOS/CUDA package becomes faster
depends strongly on the pair style.  For example, for a simple Lennard Jones
system the crossover (in single precision) is often about 50K\sphinxhyphen{}100K
atoms per GPU.  When performing double precision calculations the
crossover point can be significantly smaller.

\item {} 
\sphinxAtStartPar
Both KOKKOS and GPU package compute bonded interactions (bonds, angles,
etc) on the CPU.  If the GPU package is running with several MPI processes
assigned to one GPU, the cost of computing the bonded interactions is
spread across more CPUs and hence the GPU package can run faster in these
cases.

\item {} 
\sphinxAtStartPar
When using LAMMPS with multiple MPI ranks assigned to the same GPU, its
performance depends to some extent on the available bandwidth between
the CPUs and the GPU. This can differ significantly based on the
available bus technology, capability of the host CPU and mainboard,
the wiring of the buses and whether switches are used to increase the
number of available bus slots, or if GPUs are housed in an external
enclosure.  This can become quite complex.

\item {} 
\sphinxAtStartPar
To achieve significant acceleration through GPUs, both KOKKOS and GPU
package require capable GPUs with fast on\sphinxhyphen{}device memory and efficient
data transfer rates. This requests capable upper mid\sphinxhyphen{}level to high\sphinxhyphen{}end
(desktop) GPUs. Using lower performance GPUs (e.g. on laptops) may
result in a slowdown instead.

\item {} 
\sphinxAtStartPar
For the GPU package, specifically when running in parallel with MPI,
if it often more efficient to exclude the PPPM kspace style from GPU
acceleration and instead run it \sphinxhyphen{} concurrently with a GPU accelerated
pair style \sphinxhyphen{} on the CPU. This can often be easily achieved with placing
a \sphinxstyleemphasis{suffix off} command before and a \sphinxstyleemphasis{suffix on} command after the
\sphinxstyleemphasis{kspace\_style pppm} command.

\item {} 
\sphinxAtStartPar
The KOKKOS/OpenMP and OPENMP package have different thread management
strategies, which should result in OPENMP being more efficient for a
small number of threads with increasing overhead as the number of threads
per MPI rank grows. The KOKKOS/OpenMP kernels have less overhead in that
case, but have lower performance with few threads.

\item {} 
\sphinxAtStartPar
The INTEL package contains many options and settings for achieving
additional performance on Intel hardware (CPU and accelerator cards), but
to unlock this potential, an Intel compiler is required. The package code
will compile with GNU gcc, but it will not be as efficient.

\end{itemize}

\sphinxAtStartPar
\sphinxstylestrong{Differences between the GPU and KOKKOS packages:}
\begin{itemize}
\item {} 
\sphinxAtStartPar
The GPU package accelerates only pair force, neighbor list, and (parts
of) PPPM calculations. The KOKKOS package attempts to run most of the
calculation on the GPU, but can transparently support non\sphinxhyphen{}accelerated
code (with a performance penalty due to having data transfers between
host and GPU).

\item {} 
\sphinxAtStartPar
The GPU package requires neighbor lists to be built on the CPU when using
exclusion lists, or a triclinic simulation box.

\item {} 
\sphinxAtStartPar
The GPU package can be compiled for CUDA or OpenCL and thus supports
both, NVIDIA and AMD GPUs well. On NVIDIA hardware, using CUDA is typically
resulting in equal or better performance over OpenCL.

\item {} 
\sphinxAtStartPar
OpenCL in the GPU package does theoretically also support Intel CPUs or
Intel Xeon Phi, but the native support for those in KOKKOS (or INTEL)
is superior.

\end{itemize}

\sphinxstepscope


\chapter{Howto discussions}
\label{\detokenize{Howto:howto-discussions}}\label{\detokenize{Howto::doc}}
\sphinxAtStartPar
These doc pages describe how to perform various tasks with LAMMPS,
both for users and developers.  The
\sphinxhref{https://www.lammps.org/glossary.html}{glossary} website page also lists MD
terminology, with links to corresponding LAMMPS manual pages.  The
example input scripts included in the \sphinxcode{\sphinxupquote{examples}} directory of the LAMMPS
source code distribution and highlighted on the {\hyperref[\detokenize{Examples::doc}]{\sphinxcrossref{\DUrole{doc}{Example scripts}}}} page
also show how to set up and run various kinds of simulations.


\section{General howto}
\label{\detokenize{Howto:general-howto}}
\sphinxstepscope


\subsection{Restart a simulation}
\label{\detokenize{Howto_restart:restart-a-simulation}}\label{\detokenize{Howto_restart::doc}}
\sphinxAtStartPar
There are 3 ways to continue a long LAMMPS simulation.  Multiple
{\hyperref[\detokenize{run::doc}]{\sphinxcrossref{\DUrole{doc}{run}}}} commands can be used in the same input script.  Each
run will continue from where the previous run left off.  Or binary
restart files can be saved to disk using the {\hyperref[\detokenize{restart::doc}]{\sphinxcrossref{\DUrole{doc}{restart}}}}
command.  At a later time, these binary files can be read via a
{\hyperref[\detokenize{read_restart::doc}]{\sphinxcrossref{\DUrole{doc}{read\_restart}}}} command in a new script.  Or they can
be converted to text data files using the {\hyperref[\detokenize{Run_options::doc}]{\sphinxcrossref{\DUrole{doc}{\sphinxhyphen{}r command\sphinxhyphen{}line switch}}}} and read by a {\hyperref[\detokenize{read_data::doc}]{\sphinxcrossref{\DUrole{doc}{read\_data}}}} command in a
new script.

\sphinxAtStartPar
Here we give examples of 2 scripts that read either a binary restart
file or a converted data file and then issue a new run command to
continue where the previous run left off.  They illustrate what
settings must be made in the new script.  Details are discussed in the
documentation for the {\hyperref[\detokenize{read_restart::doc}]{\sphinxcrossref{\DUrole{doc}{read\_restart}}}} and
{\hyperref[\detokenize{read_data::doc}]{\sphinxcrossref{\DUrole{doc}{read\_data}}}} commands.

\sphinxAtStartPar
Look at the \sphinxstyleemphasis{in.chain} input script provided in the \sphinxstyleemphasis{bench} directory
of the LAMMPS distribution to see the original script that these 2
scripts are based on.  If that script had the line

\begin{sphinxVerbatim}[commandchars=\\\{\}]
\PYG{k}{restart}\PYG{+w}{         }\PYG{l+m}{50}\PYG{+w}{ }\PYG{n}{tmp.restart}
\end{sphinxVerbatim}

\sphinxAtStartPar
added to it, it would produce two binary restart files (\sphinxcode{\sphinxupquote{tmp.restart.50}}
and \sphinxcode{\sphinxupquote{tmp.restart.100}}) as it ran.

\sphinxAtStartPar
This script could be used to read the first restart file and re\sphinxhyphen{}run the
last 50 timesteps:

\begin{sphinxVerbatim}[commandchars=\\\{\}]
\PYG{k}{read\PYGZus{}restart}\PYG{+w}{    }\PYG{n}{tmp.restart.50}

\PYG{k}{neighbor}\PYG{+w}{        }\PYG{l+m}{0.4}\PYG{+w}{ }\PYG{n}{bin}
\PYG{k}{neigh\PYGZus{}modify}\PYG{+w}{    }\PYG{n}{every}\PYG{+w}{ }\PYG{l+m}{1}\PYG{+w}{ }\PYG{n}{delay}\PYG{+w}{ }\PYG{l+m}{1}

\PYG{k}{fix             }\PYG{n+nv+nvIdentifier}{1}\PYG{+w}{ }\PYG{n+nv+nvIdentifier}{all}\PYG{+w}{ }\PYG{n}{nve}
\PYG{k}{fix             }\PYG{n+nv+nvIdentifier}{2}\PYG{+w}{ }\PYG{n+nv+nvIdentifier}{all}\PYG{+w}{ }\PYG{n}{langevin}\PYG{+w}{ }\PYG{l+m}{1.0}\PYG{+w}{ }\PYG{l+m}{1.0}\PYG{+w}{ }\PYG{l+m}{10.0}\PYG{+w}{ }\PYG{l+m}{904297}

\PYG{k}{timestep}\PYG{+w}{        }\PYG{l+m}{0.012}

\PYG{k}{run}\PYG{+w}{             }\PYG{l+m}{50}
\end{sphinxVerbatim}

\sphinxAtStartPar
Note that the following commands do not need to be repeated because
their settings are included in the restart file: \sphinxcode{\sphinxupquote{%
\PYG{k}{units}%
}}, \sphinxcode{\sphinxupquote{%
\PYG{k}{atom\PYGZus{}style}%
}},
\sphinxcode{\sphinxupquote{%
\PYG{k}{special\PYGZus{}bonds}%
}}, \sphinxcode{\sphinxupquote{%
\PYG{k}{pair\PYGZus{}style}%
}}, \sphinxcode{\sphinxupquote{%
\PYG{k}{bond\PYGZus{}style}%
}}.  However, these commands do
need to be used, since their settings are not in the restart file:
\sphinxcode{\sphinxupquote{%
\PYG{k}{neighbor}%
}}, \sphinxcode{\sphinxupquote{%
\PYG{k}{fix}%
}}, \sphinxcode{\sphinxupquote{%
\PYG{k}{timestep}%
}}.

\sphinxAtStartPar
If you actually use this script to perform a restarted run, you will
notice that the thermodynamic data match at step 50 (if you also put a
\sphinxcode{\sphinxupquote{%
\PYG{k}{thermo}\PYG{+w}{ }\PYG{l+m}{50}%
}} command in the original script), but do not match at step
100.  This is because the {\hyperref[\detokenize{fix_langevin::doc}]{\sphinxcrossref{\DUrole{doc}{fix langevin}}}} command
uses random numbers in a way that does not allow for perfect restarts.

\sphinxAtStartPar
As an alternate approach, the restart file could be converted to a data
file as follows:

\begin{sphinxVerbatim}[commandchars=\\\{\}]
lmp\PYGZus{}g++\PYG{+w}{ }\PYGZhy{}r\PYG{+w}{ }tmp.restart.50\PYG{+w}{ }tmp.restart.data
\end{sphinxVerbatim}

\sphinxAtStartPar
Then, this script could be used to re\sphinxhyphen{}run the last 50 steps:

\begin{sphinxVerbatim}[commandchars=\\\{\}]
\PYG{k}{units}\PYG{+w}{           }\PYG{n}{lj}
\PYG{k}{atom\PYGZus{}style}\PYG{+w}{      }\PYG{n}{bond}
\PYG{k}{pair\PYGZus{}style}\PYG{+w}{      }\PYG{n}{lj}\PYG{o}{/}\PYG{n}{cut}\PYG{+w}{ }\PYG{l+m}{1.12}
\PYG{k}{pair\PYGZus{}modify}\PYG{+w}{     }\PYG{n}{shift}\PYG{+w}{ }\PYG{n}{yes}
\PYG{k}{bond\PYGZus{}style}\PYG{+w}{      }\PYG{n}{fene}
\PYG{k}{special\PYGZus{}bonds}\PYG{+w}{   }\PYG{l+m}{0.0}\PYG{+w}{ }\PYG{l+m}{1.0}\PYG{+w}{ }\PYG{l+m}{1.0}

\PYG{k}{read\PYGZus{}data}\PYG{+w}{       }\PYG{n}{tmp.restart.data}

\PYG{k}{neighbor}\PYG{+w}{        }\PYG{l+m}{0.4}\PYG{+w}{ }\PYG{n}{bin}
\PYG{k}{neigh\PYGZus{}modify}\PYG{+w}{    }\PYG{n}{every}\PYG{+w}{ }\PYG{l+m}{1}\PYG{+w}{ }\PYG{n}{delay}\PYG{+w}{ }\PYG{l+m}{1}

\PYG{k}{fix             }\PYG{n+nv+nvIdentifier}{1}\PYG{+w}{ }\PYG{n+nv+nvIdentifier}{all}\PYG{+w}{ }\PYG{n}{nve}
\PYG{k}{fix             }\PYG{n+nv+nvIdentifier}{2}\PYG{+w}{ }\PYG{n+nv+nvIdentifier}{all}\PYG{+w}{ }\PYG{n}{langevin}\PYG{+w}{ }\PYG{l+m}{1.0}\PYG{+w}{ }\PYG{l+m}{1.0}\PYG{+w}{ }\PYG{l+m}{10.0}\PYG{+w}{ }\PYG{l+m}{904297}

\PYG{k}{timestep}\PYG{+w}{        }\PYG{l+m}{0.012}

\PYG{k}{reset\PYGZus{}timestep}\PYG{+w}{  }\PYG{l+m}{50}
\PYG{k}{run}\PYG{+w}{             }\PYG{l+m}{50}
\end{sphinxVerbatim}

\sphinxAtStartPar
Note that nearly all the settings specified in the original \sphinxcode{\sphinxupquote{in.chain}}
script must be repeated, except the \sphinxcode{\sphinxupquote{%
\PYG{k}{pair\PYGZus{}coeff}%
}} and \sphinxcode{\sphinxupquote{%
\PYG{k}{bond\PYGZus{}coeff}%
}}
commands, since the new data file lists the force field coefficients.
Also, the {\hyperref[\detokenize{reset_timestep::doc}]{\sphinxcrossref{\DUrole{doc}{reset\_timestep}}}} command is used to tell
LAMMPS the current timestep.  This value is stored in restart files, but
not in data files.

\sphinxstepscope


\subsection{Visualize LAMMPS snapshots}
\label{\detokenize{Howto_viz:visualize-lammps-snapshots}}\label{\detokenize{Howto_viz::doc}}
\sphinxAtStartPar
Snapshots from LAMMPS simulations can be viewed, visualized, and
analyzed in a variety of ways.

\sphinxAtStartPar
LAMMPS snapshots are created by the {\hyperref[\detokenize{dump::doc}]{\sphinxcrossref{\DUrole{doc}{dump}}}} command, which
can create files in several formats. The native LAMMPS dump format is a
text file (see \sphinxcode{\sphinxupquote{%
\PYG{k}{dump }\PYG{n+nv+nvIdentifier}{atom}%
}} or \sphinxcode{\sphinxupquote{%
\PYG{k}{dump }\PYG{n+nv+nvIdentifier}{custom}%
}}) which can be visualized by
\sphinxhref{https://www.lammps.org/viz.html}{several visualization tools} for MD
simulation trajectories.  \sphinxhref{https://www.ovito.org}{OVITO} and \sphinxhref{https://www.ks.uiuc.edu/Research/vmd}{VMD} seem to be the most popular
choices among them.

\sphinxAtStartPar
The {\hyperref[\detokenize{dump_image::doc}]{\sphinxcrossref{\DUrole{doc}{dump image}}}} and {\hyperref[\detokenize{dump_image::doc}]{\sphinxcrossref{\DUrole{doc}{dump movie}}}}
styles can output internally rendered images or convert them to a movie
during the MD run.  It is also possible to create visualizations from
LAMMPS inputs or restart file with the {\hyperref[\detokenize{Howto_lammps_gui::doc}]{\sphinxcrossref{\DUrole{doc}{LAMMPS\sphinxhyphen{}GUI}}}}, which uses the {\hyperref[\detokenize{dump_image::doc}]{\sphinxcrossref{\DUrole{doc}{dump image}}}}
command internally.  The \sphinxcode{\sphinxupquote{Snapshot Image Viewer}} can be used to
adjust the visualization of the system interactively and then export
the corresponding LAMMPS commands to the clipboard to be inserted
into input files.

\sphinxAtStartPar
Programs included with LAMMPS as auxiliary tools can convert
between LAMMPS format files and other formats.  See the {\hyperref[\detokenize{Tools::doc}]{\sphinxcrossref{\DUrole{doc}{Tools}}}} page for details.  These are rarely needed these days.

\sphinxstepscope


\subsection{Run multiple simulations from one input script}
\label{\detokenize{Howto_multiple:run-multiple-simulations-from-one-input-script}}\label{\detokenize{Howto_multiple::doc}}
\sphinxAtStartPar
This can be done in several ways.  See the documentation for
individual commands for more details on how these examples work.

\sphinxAtStartPar
If “multiple simulations” means to continue a previous simulation for
more timesteps, then you simply use the {\hyperref[\detokenize{run::doc}]{\sphinxcrossref{\DUrole{doc}{run}}}} command
multiple times.  For example, this script

\begin{sphinxVerbatim}[commandchars=\\\{\}]
\PYG{k}{units}\PYG{+w}{ }\PYG{n}{lj}
\PYG{k}{atom\PYGZus{}style}\PYG{+w}{ }\PYG{n}{atomic}
\PYG{k}{read\PYGZus{}data}\PYG{+w}{ }\PYG{n}{data.lj}
\PYG{k}{run}\PYG{+w}{ }\PYG{l+m}{10000}
\PYG{k}{run}\PYG{+w}{ }\PYG{l+m}{10000}
\PYG{k}{run}\PYG{+w}{ }\PYG{l+m}{10000}
\PYG{k}{run}\PYG{+w}{ }\PYG{l+m}{10000}
\PYG{k}{run}\PYG{+w}{ }\PYG{l+m}{10000}
\end{sphinxVerbatim}

\sphinxAtStartPar
would run 5 successive simulations of the same system for a total of
50,000 timesteps.

\sphinxAtStartPar
If you wish to run totally different simulations, one after the other,
the {\hyperref[\detokenize{clear::doc}]{\sphinxcrossref{\DUrole{doc}{clear}}}} command can be used in between them to
re\sphinxhyphen{}initialize LAMMPS.  For example, this script

\begin{sphinxVerbatim}[commandchars=\\\{\}]
\PYG{k}{units}\PYG{+w}{ }\PYG{n}{lj}
\PYG{k}{atom\PYGZus{}style}\PYG{+w}{ }\PYG{n}{atomic}
\PYG{k}{read\PYGZus{}data}\PYG{+w}{ }\PYG{n}{data.lj}
\PYG{k}{run}\PYG{+w}{ }\PYG{l+m}{10000}
\PYG{k}{clear}
\PYG{k}{units}\PYG{+w}{ }\PYG{n}{lj}
\PYG{k}{atom\PYGZus{}style}\PYG{+w}{ }\PYG{n}{atomic}
\PYG{k}{read\PYGZus{}data}\PYG{+w}{ }\PYG{n}{data.lj.new}
\PYG{k}{run}\PYG{+w}{ }\PYG{l+m}{10000}
\end{sphinxVerbatim}

\sphinxAtStartPar
would run 2 independent simulations, one after the other.

\sphinxAtStartPar
For large numbers of independent simulations, you can use
{\hyperref[\detokenize{variable::doc}]{\sphinxcrossref{\DUrole{doc}{variables}}}} and the {\hyperref[\detokenize{next::doc}]{\sphinxcrossref{\DUrole{doc}{next}}}} and
{\hyperref[\detokenize{jump::doc}]{\sphinxcrossref{\DUrole{doc}{jump}}}} commands to loop over the same input script
multiple times with different settings.  For example, this
script, named \sphinxcode{\sphinxupquote{in.polymer}}

\begin{sphinxVerbatim}[commandchars=\\\{\}]
\PYG{k}{variable }\PYG{n+nv+nvIdentifier}{d}\PYG{+w}{ }\PYG{n}{index}\PYG{+w}{ }\PYG{n}{run1}\PYG{+w}{ }\PYG{n}{run2}\PYG{+w}{ }\PYG{n}{run3}\PYG{+w}{ }\PYG{n}{run4}\PYG{+w}{ }\PYG{n}{run5}\PYG{+w}{ }\PYG{n}{run6}\PYG{+w}{ }\PYG{n}{run7}\PYG{+w}{ }\PYG{n}{run8}
\PYG{k}{shell}\PYG{+w}{ }\PYG{n}{cd}\PYG{+w}{ }\PYG{n+nv}{\PYGZdl{}d}
\PYG{k}{read\PYGZus{}data}\PYG{+w}{ }\PYG{n}{data.polymer}
\PYG{k}{run}\PYG{+w}{ }\PYG{l+m}{10000}
\PYG{k}{shell}\PYG{+w}{ }\PYG{n}{cd}\PYG{+w}{ }\PYG{n}{..}
\PYG{k}{clear}
\PYG{k}{next }\PYG{n+nv+nvIdentifier}{d}
\PYG{k}{jump }\PYG{l+s+sc}{in.polymer}
\end{sphinxVerbatim}

\sphinxAtStartPar
would run 8 simulations in different directories, using a \sphinxcode{\sphinxupquote{data.polymer}}
file in each directory.  The same concept could be used to run the
same system at 8 different temperatures, using a temperature variable
and storing the output in different log and dump files, for example

\begin{sphinxVerbatim}[commandchars=\\\{\}]
\PYG{k}{variable }\PYG{n+nv+nvIdentifier}{a}\PYG{+w}{ }\PYG{n}{loop}\PYG{+w}{ }\PYG{l+m}{8}
\PYG{k}{variable }\PYG{n+nv+nvIdentifier}{t}\PYG{+w}{ }\PYG{n}{index}\PYG{+w}{ }\PYG{l+m}{0.8}\PYG{+w}{ }\PYG{l+m}{0.85}\PYG{+w}{ }\PYG{l+m}{0.9}\PYG{+w}{ }\PYG{l+m}{0.95}\PYG{+w}{ }\PYG{l+m}{1.0}\PYG{+w}{ }\PYG{l+m}{1.05}\PYG{+w}{ }\PYG{l+m}{1.1}\PYG{+w}{ }\PYG{l+m}{1.15}
\PYG{k}{log}\PYG{+w}{ }\PYG{n}{log.}\PYG{n+nv}{\PYGZdl{}a}
\PYG{n}{read}\PYG{+w}{ }\PYG{n}{data.polymer}
\PYG{k}{velocity }\PYG{n+nv+nvIdentifier}{all}\PYG{+w}{ }\PYG{n}{create}\PYG{+w}{ }\PYG{n+nv}{\PYGZdl{}t}\PYG{+w}{ }\PYG{l+m}{352839}
\PYG{k}{fix }\PYG{n+nv+nvIdentifier}{1}\PYG{+w}{ }\PYG{n+nv+nvIdentifier}{all}\PYG{+w}{ }\PYG{n}{nvt}\PYG{+w}{ }\PYG{n+nv}{\PYGZdl{}t}\PYG{+w}{ }\PYG{n+nv}{\PYGZdl{}t}\PYG{+w}{ }\PYG{l+m}{100.0}
\PYG{k}{dump }\PYG{n+nv+nvIdentifier}{1}\PYG{+w}{ }\PYG{n+nv+nvIdentifier}{all}\PYG{+w}{ }\PYG{n}{atom}\PYG{+w}{ }\PYG{l+m}{1000}\PYG{+w}{ }\PYG{n}{dump.}\PYG{n+nv}{\PYGZdl{}a}
\PYG{k}{run}\PYG{+w}{ }\PYG{l+m}{100000}
\PYG{k}{clear}
\PYG{k}{next }\PYG{n+nv+nvIdentifier}{t}
\PYG{k}{next }\PYG{n+nv+nvIdentifier}{a}
\PYG{k}{jump }\PYG{l+s+sc}{in.polymer}
\end{sphinxVerbatim}

\sphinxAtStartPar
All of the above examples work whether you are running on 1 or
multiple processors, but assumed you are running LAMMPS on a single
partition of processors.  LAMMPS can be run on multiple partitions via
the {\hyperref[\detokenize{Run_options::doc}]{\sphinxcrossref{\DUrole{doc}{\sphinxhyphen{}partition command\sphinxhyphen{}line switch}}}}.

\sphinxAtStartPar
In the last 2 examples, if LAMMPS were run on 3 partitions, the same
scripts could be used if the \sphinxcode{\sphinxupquote{index}} and \sphinxcode{\sphinxupquote{loop}} variables were
replaced with \sphinxstyleemphasis{universe}\sphinxhyphen{}style variables, as described in the
{\hyperref[\detokenize{variable::doc}]{\sphinxcrossref{\DUrole{doc}{variable}}}} command.  Also, the \sphinxcode{\sphinxupquote{%
\PYG{k}{next }\PYG{n+nv+nvIdentifier}{t}%
}} and \sphinxcode{\sphinxupquote{%
\PYG{k}{next }\PYG{n+nv+nvIdentifier}{a}%
}}
commands would need to be replaced with a single \sphinxcode{\sphinxupquote{%
\PYG{k}{next }\PYG{n+nv+nvIdentifier}{a}\PYG{+w}{ }\PYG{n}{t}%
}} command.
With these modifications, the 8 simulations of each script would run
on the 3 partitions one after the other until all were finished.
Initially, 3 simulations would be started simultaneously, one on each
partition.  When one finished, that partition would then start
the fourth simulation, and so forth, until all 8 were completed.

\sphinxstepscope


\subsection{Multi\sphinxhyphen{}replica simulations}
\label{\detokenize{Howto_replica:multi-replica-simulations}}\label{\detokenize{Howto_replica::doc}}
\sphinxAtStartPar
Several commands in LAMMPS run multi\sphinxhyphen{}replica simulations, meaning
that multiple instances (replicas) of your simulation are run
simultaneously, with small amounts of data exchanged between replicas
periodically.

\sphinxAtStartPar
These are the relevant commands:
\begin{itemize}
\item {} 
\sphinxAtStartPar
{\hyperref[\detokenize{hyper::doc}]{\sphinxcrossref{\DUrole{doc}{hyper}}}} for bond boost hyperdynamics (HD)

\item {} 
\sphinxAtStartPar
{\hyperref[\detokenize{neb::doc}]{\sphinxcrossref{\DUrole{doc}{neb}}}} for nudged elastic band calculations (NEB)

\item {} 
\sphinxAtStartPar
{\hyperref[\detokenize{neb_spin::doc}]{\sphinxcrossref{\DUrole{doc}{neb\_spin}}}} for magnetic nudged elastic band calculations

\item {} 
\sphinxAtStartPar
{\hyperref[\detokenize{prd::doc}]{\sphinxcrossref{\DUrole{doc}{prd}}}} for parallel replica dynamics (PRD)

\item {} 
\sphinxAtStartPar
{\hyperref[\detokenize{tad::doc}]{\sphinxcrossref{\DUrole{doc}{tad}}}} for temperature accelerated dynamics (TAD)

\item {} 
\sphinxAtStartPar
{\hyperref[\detokenize{temper::doc}]{\sphinxcrossref{\DUrole{doc}{temper}}}} for parallel tempering with fixed volume

\item {} 
\sphinxAtStartPar
{\hyperref[\detokenize{temper_npt::doc}]{\sphinxcrossref{\DUrole{doc}{temper/npt}}}} for parallel tempering extended for NPT

\item {} 
\sphinxAtStartPar
{\hyperref[\detokenize{temper_grem::doc}]{\sphinxcrossref{\DUrole{doc}{temper/grem}}}} for parallel tempering with generalized replica exchange (gREM)

\item {} 
\sphinxAtStartPar
{\hyperref[\detokenize{fix_pimd::doc}]{\sphinxcrossref{\DUrole{doc}{fix pimd}}}} for path\sphinxhyphen{}integral molecular dynamics (PIMD)

\end{itemize}

\sphinxAtStartPar
NEB is a method for finding transition states and barrier potential energies.
HD, PRD, and TAD are methods for performing accelerated dynamics to find and
perform infrequent events.  Parallel tempering or replica exchange runs
different replicas at a series of temperature to facilitate rare\sphinxhyphen{}event
sampling.  PIMD runs different replicas whose individual particles in different
replicas are coupled together by springs to model a system of ring\sphinxhyphen{}polymers which
can represent the quantum nature of atom cores.

\sphinxAtStartPar
These commands can only be used if LAMMPS was built with the REPLICA
package.  See the {\hyperref[\detokenize{Build_package::doc}]{\sphinxcrossref{\DUrole{doc}{Build package}}}} page for more
info.

\sphinxAtStartPar
In all these cases, you must run with one or more processors per
replica.  The processors assigned to each replica are determined at
run\sphinxhyphen{}time by using the {\hyperref[\detokenize{Run_options::doc}]{\sphinxcrossref{\DUrole{doc}{\sphinxhyphen{}partition command\sphinxhyphen{}line switch}}}} to launch LAMMPS on multiple partitions, which in this
context are the same as replicas.  E.g.  these commands:

\begin{sphinxVerbatim}[commandchars=\\\{\}]
mpirun\PYG{+w}{ }\PYGZhy{}np\PYG{+w}{ }\PYG{l+m}{16}\PYG{+w}{ }lmp\PYGZus{}linux\PYG{+w}{ }\PYGZhy{}partition\PYG{+w}{ }8x2\PYG{+w}{ }\PYGZhy{}in\PYG{+w}{ }\PYG{k}{in}.temper
mpirun\PYG{+w}{ }\PYGZhy{}np\PYG{+w}{ }\PYG{l+m}{8}\PYG{+w}{ }lmp\PYGZus{}linux\PYG{+w}{ }\PYGZhy{}partition\PYG{+w}{ }8x1\PYG{+w}{ }\PYGZhy{}in\PYG{+w}{ }\PYG{k}{in}.neb
\end{sphinxVerbatim}

\sphinxAtStartPar
would each run 8 replicas, on either 16 or 8 processors.  Note the use
of the {\hyperref[\detokenize{Run_options::doc}]{\sphinxcrossref{\DUrole{doc}{\sphinxhyphen{}in command\sphinxhyphen{}line switch}}}} to specify the input
script which is required when running in multi\sphinxhyphen{}replica mode.

\sphinxAtStartPar
Also note that with MPI installed on a machine (e.g. your desktop), you
can run on more (virtual) processors than you have physical processors.
Thus, the above commands could be run on a single\sphinxhyphen{}processor (or
few\sphinxhyphen{}processor) desktop so that you can run a multi\sphinxhyphen{}replica simulation on
more replicas than you have physical processors. This is useful for
testing and debugging, since with most modern processors and MPI
libraries, the efficiency of a calculation can severely diminish when
oversubscribing processors.

\sphinxstepscope


\subsection{Library interface to LAMMPS}
\label{\detokenize{Howto_library:library-interface-to-lammps}}\label{\detokenize{Howto_library::doc}}
\sphinxAtStartPar
As described on the {\hyperref[\detokenize{Build_basics::doc}]{\sphinxcrossref{\DUrole{doc}{Build basics}}}} doc page, LAMMPS
can be built as a static or shared library, so that it can be called by
another code, used in a {\hyperref[\detokenize{Howto_couple::doc}]{\sphinxcrossref{\DUrole{doc}{coupled manner}}}} with other
codes, or driven through a {\hyperref[\detokenize{Python_head::doc}]{\sphinxcrossref{\DUrole{doc}{Python interface}}}}.

\sphinxAtStartPar
At the core of LAMMPS is the \sphinxcode{\sphinxupquote{LAMMPS}} class, which encapsulates the
state of the simulation program through the state of the various class
instances that it is composed of.  So a calculation using LAMMPS
requires creating an instance of the \sphinxcode{\sphinxupquote{LAMMPS}} class and then send it
(text) commands, either individually or from a file, or perform other
operations that modify the state stored inside that instance or drive
simulations.  This is essentially what the \sphinxcode{\sphinxupquote{src/main.cpp}} file does as
well for the standalone LAMMPS executable, reading commands either from
an input file or the standard input.

\sphinxAtStartPar
Creating a LAMMPS instance can be done by using C++ code directly or
through a C\sphinxhyphen{}style interface library to LAMMPS that is provided in the
files \sphinxcode{\sphinxupquote{src/library.cpp}} and \sphinxcode{\sphinxupquote{src/library.h}}.  This {\hyperref[\detokenize{Library:lammps-c-api}]{\sphinxcrossref{\DUrole{std,std-ref}{C language
API}}}}, can be used from C and C++, and is also the basis
for the {\hyperref[\detokenize{Python_module::doc}]{\sphinxcrossref{\DUrole{doc}{Python}}}} and {\hyperref[\detokenize{Fortran::doc}]{\sphinxcrossref{\DUrole{doc}{Fortran}}}}
interfaces or the {\hyperref[\detokenize{Tools:swig}]{\sphinxcrossref{\DUrole{std,std-ref}{SWIG based wrappers}}}} included in the
LAMMPS source code.

\sphinxAtStartPar
The \sphinxcode{\sphinxupquote{examples/COUPLE}} and \sphinxcode{\sphinxupquote{python/examples}} directories contain some
example programs written in C++, C, Fortran, and Python, which show how
a driver code can link to LAMMPS as a library, run LAMMPS on a subset of
processors (so the others are available to run some other code
concurrently), grab data from LAMMPS, change it, and send it back into
LAMMPS.

\sphinxAtStartPar
A detailed documentation of the available APIs and examples of how to
use them can be found in the {\hyperref[\detokenize{Manual:programmer-documentation}]{\sphinxcrossref{\DUrole{std,std-ref}{Programmer Guide}}}}
section of this manual.

\sphinxstepscope


\subsection{Coupling LAMMPS to other codes}
\label{\detokenize{Howto_couple:coupling-lammps-to-other-codes}}\label{\detokenize{Howto_couple::doc}}
\sphinxAtStartPar
LAMMPS is designed to support being coupled to other codes.  For
example, a quantum mechanics code might compute forces on a subset of
atoms and pass those forces to LAMMPS.  Or a continuum finite element
(FE) simulation might use atom positions as boundary conditions on FE
nodal points, compute a FE solution, and return interpolated forces on
MD atoms.

\sphinxAtStartPar
LAMMPS can be coupled to other codes in at least 4 different ways.  Each
has advantages and disadvantages, which you will have to think about in
the context of your application.
\begin{enumerate}
\sphinxsetlistlabels{\arabic}{enumi}{enumii}{}{.}%
\item {} 
\sphinxAtStartPar
Define a new {\hyperref[\detokenize{fix::doc}]{\sphinxcrossref{\DUrole{doc}{fix}}}} or {\hyperref[\detokenize{compute::doc}]{\sphinxcrossref{\DUrole{doc}{compute}}}} command
that calls the other code.  In this scenario, LAMMPS is the driver
code.  During timestepping, the fix or compute is invoked, and can
make library calls to the other code, which has been linked to LAMMPS
as a library.  This is the way the {\hyperref[\detokenize{Packages_details:pkg-voronoi}]{\sphinxcrossref{\DUrole{std,std-ref}{VORONOI}}}}
package, which computes Voronoi tesselations using the \sphinxhref{http://math.lbl.gov/voro++}{Voro++
library}, is interfaced to LAMMPS.  See
the {\hyperref[\detokenize{compute_voronoi_atom::doc}]{\sphinxcrossref{\DUrole{doc}{compute voronoi}}}} command for more
details.  Also see the {\hyperref[\detokenize{Modify::doc}]{\sphinxcrossref{\DUrole{doc}{Modify}}}} pages for information
on how to add a new fix or compute to LAMMPS.

\end{enumerate}
\begin{enumerate}
\sphinxsetlistlabels{\arabic}{enumi}{enumii}{}{.}%
\setcounter{enumi}{1}
\item {} 
\sphinxAtStartPar
Define a new LAMMPS command that calls the other code.  This is
conceptually similar to method (1), but in this case LAMMPS and the
other code are on a more equal footing.  Note that now the other code
is not called during the timestepping of a LAMMPS run, but between
runs.  The LAMMPS input script can be used to alternate LAMMPS runs
with calls to the other code, invoked via the new command.  The
{\hyperref[\detokenize{run::doc}]{\sphinxcrossref{\DUrole{doc}{run}}}} command facilitates this with its \sphinxstyleemphasis{every} option,
which makes it easy to run a few steps, invoke the command, run a few
steps, invoke the command, etc.

\sphinxAtStartPar
In this scenario, the other code can be called as a library, as in
1., or it could be a stand\sphinxhyphen{}alone code, invoked by a \sphinxcode{\sphinxupquote{system()}} call
made by the command (assuming your parallel machine allows one or
more processors to start up another program).  In the latter case the
stand\sphinxhyphen{}alone code could communicate with LAMMPS through files that the
command writes and reads.

\sphinxAtStartPar
See the {\hyperref[\detokenize{Modify_command::doc}]{\sphinxcrossref{\DUrole{doc}{Modify command}}}} page for information
on how to add a new command to LAMMPS.

\end{enumerate}
\begin{enumerate}
\sphinxsetlistlabels{\arabic}{enumi}{enumii}{}{.}%
\setcounter{enumi}{2}
\item {} 
\sphinxAtStartPar
Use LAMMPS as a library called by another code.  In this case, the
other code is the driver and calls LAMMPS as needed.  Alternately, a
wrapper code could link and call both LAMMPS and another code as
libraries.  Again, the {\hyperref[\detokenize{run::doc}]{\sphinxcrossref{\DUrole{doc}{run}}}} command has options that
allow it to be invoked with minimal overhead (no setup or clean\sphinxhyphen{}up)
if you wish to do multiple short runs, driven by another program.
Details about using the library interface are given in the
{\hyperref[\detokenize{Library::doc}]{\sphinxcrossref{\DUrole{doc}{library API}}}} documentation.

\end{enumerate}
\begin{enumerate}
\sphinxsetlistlabels{\arabic}{enumi}{enumii}{}{.}%
\setcounter{enumi}{3}
\item {} 
\sphinxAtStartPar
Couple LAMMPS with another code in a client/server fashion, using the
\sphinxhref{https://molssi-mdi.github.io/MDI\_Library/html/index.html}{MDI Library}
developed by the \sphinxhref{https://molssi.org}{Molecular Sciences Software Institute (MolSSI)} to run LAMMPS as either an MDI driver (client)
or an MDI engine (server).  The MDI driver issues commands to the MDI
server to exchange data between them.  See the {\hyperref[\detokenize{Howto_mdi::doc}]{\sphinxcrossref{\DUrole{doc}{Using LAMMPS with the MDI library for code coupling}}}} page for
more information about how LAMMPS can operate in either of these modes.

\end{enumerate}

\sphinxstepscope


\subsection{Using LAMMPS with the MDI library for code coupling}
\label{\detokenize{Howto_mdi:using-lammps-with-the-mdi-library-for-code-coupling}}\label{\detokenize{Howto_mdi::doc}}
\sphinxAtStartPar
Client/server coupling of two (or more) codes is where one code is the
“client” and sends request messages (data) to one (or more) “server”
code(s).  A server responds to each request with a reply message
(data).  This enables two (or more) codes to work in tandem to perform
a simulation.  In this context, LAMMPS can act as either a client or
server code.  It does this by using the \sphinxhref{https://molssi-mdi.github.io/MDI\_Library/html/index.html}{MolSSI Driver Interface (MDI)
library},
developed by the \sphinxhref{https://molssi.org}{Molecular Sciences Software Institute (MolSSI)}, which is supported by the {\hyperref[\detokenize{Packages_details:pkg-mdi}]{\sphinxcrossref{\DUrole{std,std-ref}{MDI}}}}
package.

\sphinxAtStartPar
Alternate methods for coupling codes with LAMMPS are described on the
{\hyperref[\detokenize{Howto_couple::doc}]{\sphinxcrossref{\DUrole{doc}{Coupling LAMMPS to other codes}}}} page.

\sphinxAtStartPar
Some advantages of client/server coupling are that the codes can run
as stand\sphinxhyphen{}alone executables; they need not be linked together.  Thus,
neither code needs to have a library interface.  This also makes it
easy to run the two codes on different numbers of processors.  If a
message protocol (format and content) is defined for a particular kind
of simulation, then in principle any code which implements the
client\sphinxhyphen{}side protocol can be used in tandem with any code which
implements the server\sphinxhyphen{}side protocol.  Neither code needs to know what
specific other code it is working with.

\sphinxAtStartPar
In MDI nomenclature, a client code is the “driver”, and a server code is
an “engine”.  One driver code can communicate with one or more instances
of one or more engine codes.  Driver and engine codes can be written in
any language: C, C++, Fortran, Python, etc.

\sphinxAtStartPar
In addition to allowing driver and engine(s) to run as stand\sphinxhyphen{}alone
executables, MDI also enables an engine to be a \sphinxstyleemphasis{plugin} to the client
code.  In this scenario, server code(s) are compiled as shared
libraries, and one (or more) instances of the server are instantiated
by the driver code.  If the driver code runs in parallel, it can split
its MPI communicator into multiple sub\sphinxhyphen{}communicators, and launch each
plugin engine instance on a sub\sphinxhyphen{}communicator.  Driver processors
within that sub\sphinxhyphen{}communicator exchange messages with the corresponding
engine instance, and can also send MPI messages to other processors in
the driver.  The driver code can also destroy engine instances and
re\sphinxhyphen{}instantiate them.  LAMMPS can operate as either a stand\sphinxhyphen{}alone or
plugin MDI engine.  When it operates as a driver, it can use either
stand\sphinxhyphen{}alone or plugin MDI engines.

\sphinxAtStartPar
The way in which an MDI driver communicates with an MDI engine is by
making MDI\_Send() and MDI\_Recv() calls, which are conceptually similar
to MPI\_Send() and MPI\_Recv() calls.  Each send or receive operation
uses a string to identify the command name, and optionally some data,
which can be a single value or vector of values of any data type.
Inside the MDI library, data is exchanged between the driver and
engine via MPI calls or sockets.  This is a run\sphinxhyphen{}time choice by the user.


\bigskip\hrule\bigskip


\sphinxAtStartPar
The {\hyperref[\detokenize{Packages_details:pkg-mdi}]{\sphinxcrossref{\DUrole{std,std-ref}{MDI}}}} package provides a {\hyperref[\detokenize{mdi::doc}]{\sphinxcrossref{\DUrole{doc}{mdi engine}}}}
command, which enables LAMMPS to operate as an MDI engine.  Its doc
page explains the variety of standard and custom MDI commands which
the LAMMPS engine recognizes and can respond to.

\sphinxAtStartPar
The package also provides a {\hyperref[\detokenize{mdi::doc}]{\sphinxcrossref{\DUrole{doc}{mdi plugin}}}} command, which
enables LAMMPS to operate as an MDI driver and load an MDI engine as a
plugin library.

\sphinxAtStartPar
The package furthermore includes a {\hyperref[\detokenize{fix_mdi_qm::doc}]{\sphinxcrossref{\DUrole{doc}{fix mdi/qm}}}}
command, in which LAMMPS operates as an MDI driver in conjunction with a
quantum mechanics code as an MDI engine.  The post\_force() method of the
\sphinxcode{\sphinxupquote{fix\_mdi\_qm.cpp}} file shows how a driver issues MDI commands to
another code.  This command can be used to couple to an MDI engine,
which is either a stand\sphinxhyphen{}alone code or a plugin library.

\sphinxAtStartPar
As explained in the {\hyperref[\detokenize{fix_mdi_qm::doc}]{\sphinxcrossref{\DUrole{doc}{fix mdi/qm}}}} command
documentation, it can be used to perform \sphinxstyleemphasis{ab initio} MD simulations or
energy minimizations, or to evaluate the quantum energy and forces for a
series of independent systems.  The \sphinxcode{\sphinxupquote{examples/mdi}} directory has
example input scripts for all of these use cases.

\sphinxAtStartPar
The package also has a {\hyperref[\detokenize{fix_mdi_qmmm::doc}]{\sphinxcrossref{\DUrole{doc}{fix mdi/qmmm}}}} command in
which LAMMPS operates as an MDI driver in conjunction with a quantum
mechanics code as an MDI engine to perform QM/MM simulations.  The
LAMMPS input script partitions the system into QM and MM (molecular
mechanics) atoms.  As described below the \sphinxcode{\sphinxupquote{examples/QUANTUM}} directory
has examples for coupling to 3 different quantum codes in this manner.


\bigskip\hrule\bigskip


\sphinxAtStartPar
The \sphinxcode{\sphinxupquote{examples/mdi}} directory contains Python scripts and LAMMPS input
script which use LAMMPS as either an MDI driver or engine, or both.
Currently, 5 example use cases are provided:
\begin{itemize}
\item {} 
\sphinxAtStartPar
Run ab initio MD (AIMD) using 2 instances of LAMMPS.  As a driver,
LAMMPS performs the timestepping in either NVE or NPT mode.  As an
engine, LAMMPS computes forces and is a surrogate for a quantum
code.

\item {} 
\sphinxAtStartPar
LAMMPS runs an MD simulation as a driver.  Every N steps it passes the
current snapshot to an MDI engine to evaluate the energy, virial, and
peratom forces.  As the engine, LAMMPS is a surrogate for a quantum
code.

\item {} 
\sphinxAtStartPar
LAMMPS loops over a series of data files and passes the configuration
to an MDI engine to evaluate the energy, virial, and peratom forces
and thus acts as a simulation driver.  As the engine, LAMMPS is used
as a surrogate for a quantum code.

\item {} 
\sphinxAtStartPar
A Python script driver invokes a sequence of unrelated LAMMPS
calculations.  Calculations can be single\sphinxhyphen{}point energy/force
evaluations, MD runs, or energy minimizations.

\item {} 
\sphinxAtStartPar
Run AIMD with a Python driver code and 2 LAMMPS instances as engines.
The first LAMMPS instance performs MD timestepping.  The second LAMMPS
instance acts as a surrogate QM code to compute forces.

\end{itemize}

\begin{sphinxadmonition}{note}{Note:}
\sphinxAtStartPar
In any of these examples where LAMMPS is used as an engine, an actual
QM code (provided it has support for MDI) could be used in its place,
without modifying the input scripts or launch commands, except to
specify the name of the QM code.
\end{sphinxadmonition}

\sphinxAtStartPar
The \sphinxcode{\sphinxupquote{examples/mdi/Run.sh}} file illustrates how to launch both driver
and engine codes so that they communicate using the MDI library via
either MPI or sockets, or using the engine as a stand\sphinxhyphen{}alone code, or
as a plugin library.


\bigskip\hrule\bigskip


\sphinxAtStartPar
As of March 2023, these are quantum codes with MDI support provided via
Python wrapper scripts included in the LAMMPS distribution.  These can
be used with the fix mdi/qm and fix mdi/qmmm commands to perform QM
calculations of an entire system (e.g. AIMD) or QM/MM simulations.  See
the \sphinxcode{\sphinxupquote{examples/QUANTUM}} sub\sphinxhyphen{}directories for more details:
\begin{itemize}
\item {} 
\sphinxAtStartPar
LATTE \sphinxhyphen{} AIMD only

\item {} 
\sphinxAtStartPar
PySCF \sphinxhyphen{} QM/MM only

\item {} 
\sphinxAtStartPar
NWChem \sphinxhyphen{} AIMD or QM/MM

\end{itemize}

\sphinxAtStartPar
There are also at least two quantum codes which have direct MDI
support, \sphinxhref{https://www.quantum-espresso.org/}{Quantum ESPRESSO (QE)}
and \sphinxhref{https://qsg.llnl.gov/node/101.html}{INQ}.  There are also
several QM codes which have indirect support through QCEngine or i\sphinxhyphen{}PI.
The former means they require a wrapper program (QCEngine) with MDI
support which writes/read files to pass data to the quantum code
itself.  The list of QCEngine\sphinxhyphen{}supported and i\sphinxhyphen{}PI\sphinxhyphen{}supported quantum
codes is on the \sphinxhref{https://molssi-mdi.github.io/MDI\_Library/html/index.html}{MDI webpage}.

\sphinxAtStartPar
These direct\sphinxhyphen{} and indirect\sphinxhyphen{}support codes should be usable for full
system calculations (e.g. AIMD).  Whether they support QM/MM models
depends on the individual QM code.

\sphinxstepscope


\subsection{Broken Bonds}
\label{\detokenize{Howto_broken_bonds:broken-bonds}}\label{\detokenize{Howto_broken_bonds::doc}}
\sphinxAtStartPar
Typically, molecular bond interactions persist for the duration of a
simulation in LAMMPS.  However, some commands break bonds dynamically,
including the following:
\begin{itemize}
\item {} 
\sphinxAtStartPar
{\hyperref[\detokenize{bond_quartic::doc}]{\sphinxcrossref{\DUrole{doc}{bond\_style quartic}}}}

\item {} 
\sphinxAtStartPar
{\hyperref[\detokenize{fix_bond_break::doc}]{\sphinxcrossref{\DUrole{doc}{fix bond/break}}}}

\item {} 
\sphinxAtStartPar
{\hyperref[\detokenize{fix_bond_react::doc}]{\sphinxcrossref{\DUrole{doc}{fix bond/react}}}}

\item {} 
\sphinxAtStartPar
{\hyperref[\detokenize{Howto_bpm::doc}]{\sphinxcrossref{\DUrole{doc}{BPM package}}}} bond styles

\end{itemize}

\sphinxAtStartPar
A bond can break if it is stretched beyond a user\sphinxhyphen{}defined threshold or
more generally if other criteria are met.

\sphinxAtStartPar
For the quartic bond style, when a bond is broken its bond type is set
to 0 to effectively break it and pairwise forces between the two atoms
in the broken bond are “turned on”.  Angles, dihedrals, etc cannot be
defined for a system when {\hyperref[\detokenize{bond_quartic::doc}]{\sphinxcrossref{\DUrole{doc}{bond\_style quartic}}}} is
used.

\sphinxAtStartPar
Similarly, bond styles in the BPM package are also incompatible with
angles, dihedrals, etc. and when a bond breaks its type is set to zero.
However, in the BPM package one can either turn off all pair interactions
between bonded particles or leave them on, overlaying pair forces on
top of bond forces. To remove pair forces, the special bond list is
dynamically updated.  More details can be found on the {\hyperref[\detokenize{Howto_bpm::doc}]{\sphinxcrossref{\DUrole{doc}{Howto BPM}}}} page.

\sphinxAtStartPar
The {\hyperref[\detokenize{fix_bond_break::doc}]{\sphinxcrossref{\DUrole{doc}{fix bond/break}}}} and {\hyperref[\detokenize{fix_bond_react::doc}]{\sphinxcrossref{\DUrole{doc}{fix bond/react}}}} commands allow breaking of bonds within a molecular
topology with may also define angles, dihedrals, etc.  These commands
update internal topology data structures to remove broken bonds, as
well as the appropriate angle, dihedral, etc interactions which
include the bond.  They also trigger a rebuild of the neighbor list
when this occurs, to turn on the appropriate pairwise forces.

\sphinxAtStartPar
Note that when bonds are dumped to a file via the {\hyperref[\detokenize{dump::doc}]{\sphinxcrossref{\DUrole{doc}{dump local}}}} command, bonds with type 0 are not included.

\sphinxAtStartPar
The {\hyperref[\detokenize{delete_bonds::doc}]{\sphinxcrossref{\DUrole{doc}{delete\_bonds}}}} command can be used to query
the status of broken bonds with type = 0 or permanently delete them,
e.g.:

\begin{sphinxVerbatim}[commandchars=\\\{\}]
\PYG{k}{delete\PYGZus{}bonds }\PYG{n+nv+nvIdentifier}{all}\PYG{+w}{ }\PYG{n}{stats}
\PYG{k}{delete\PYGZus{}bonds }\PYG{n+nv+nvIdentifier}{all}\PYG{+w}{ }\PYG{n}{bond}\PYG{+w}{ }\PYG{l+m}{0}\PYG{+w}{ }\PYG{n}{remove}
\end{sphinxVerbatim}

\sphinxAtStartPar
The compute {\hyperref[\detokenize{compute_count_type::doc}]{\sphinxcrossref{\DUrole{doc}{count/type}}}} command tallies the
current number of bonds (or angles, etc) for each bond (angle, etc)
type.  It also tallies broken bonds with type = 0.

\sphinxAtStartPar
The compute {\hyperref[\detokenize{compute_nbond_atom::doc}]{\sphinxcrossref{\DUrole{doc}{nbond/atom}}}} command tallies the
current number of bonds each atom is part of, excluding broken bonds
with type = 0.


\section{Settings howto}
\label{\detokenize{Howto:settings-howto}}
\sphinxstepscope


\subsection{2d simulations}
\label{\detokenize{Howto_2d:d-simulations}}\label{\detokenize{Howto_2d::doc}}
\sphinxAtStartPar
You must use the {\hyperref[\detokenize{dimension::doc}]{\sphinxcrossref{\DUrole{doc}{dimension}}}} command to specify a 2d
simulation.  The default is 3d.

\sphinxAtStartPar
A 2d simulation box must be periodic in z as set by the {\hyperref[\detokenize{boundary::doc}]{\sphinxcrossref{\DUrole{doc}{boundary}}}} command.  This is the default.

\sphinxAtStartPar
Simulation boxes in LAMMPS can be either orthogonal or triclinic in
shape.  Orthogonal boxes in 2d are a rectangle with 4 edges that are
each perpendicular to either the x or y coordinate axes.  Triclinic
boxes in 2d are a parallelogram with opposite pairs of faces parallel
to each other.  LAMMPS supports two forms of triclinic boxes,
restricted and general, which for 2d differ in how the box is oriented
with respect to the xy coordinate axes.  See the {\hyperref[\detokenize{Howto_triclinic::doc}]{\sphinxcrossref{\DUrole{doc}{Howto triclinic}}}} for a detailed description of all 3 kinds of
simulation boxes.

\sphinxAtStartPar
Here are examples of using the {\hyperref[\detokenize{create_box::doc}]{\sphinxcrossref{\DUrole{doc}{create\_box}}}} command
to define the simulation box for a 2d system.

\begin{sphinxVerbatim}[commandchars=\\\{\}]
\PYG{c}{\PYGZsh{} 2d orthogonal box using a block\PYGZhy{}style region}
\PYG{k}{region }\PYG{n+nv+nvIdentifier}{mybox}\PYG{+w}{ }\PYG{n}{block}\PYG{+w}{ }\PYG{o}{\PYGZhy{}}\PYG{l+m}{10}\PYG{+w}{ }\PYG{l+m}{10}\PYG{+w}{ }\PYG{l+m}{0}\PYG{+w}{ }\PYG{l+m}{10}\PYG{+w}{ }\PYG{o}{\PYGZhy{}}\PYG{l+m}{0.5}\PYG{+w}{ }\PYG{l+m}{0.5}
\PYG{k}{create\PYGZus{}box }\PYG{n+nv+nvIdentifier}{1}\PYG{+w}{ }\PYG{n+nv+nvIdentifier}{mybox}

\PYG{c}{\PYGZsh{} 2d restricted triclinic box using a prism\PYGZhy{}style region with only xy tilt}
\PYG{k}{region }\PYG{n+nv+nvIdentifier}{mybox}\PYG{+w}{ }\PYG{n}{prism}\PYG{+w}{ }\PYG{l+m}{0}\PYG{+w}{ }\PYG{l+m}{10}\PYG{+w}{ }\PYG{l+m}{0}\PYG{+w}{ }\PYG{l+m}{10}\PYG{+w}{ }\PYG{o}{\PYGZhy{}}\PYG{l+m}{0.5}\PYG{+w}{ }\PYG{l+m}{0.5}\PYG{+w}{ }\PYG{l+m}{2.0}\PYG{+w}{ }\PYG{l+m}{0.0}\PYG{+w}{ }\PYG{l+m}{0.0}
\PYG{k}{create\PYGZus{}box }\PYG{n+nv+nvIdentifier}{1}\PYG{+w}{ }\PYG{n+nv+nvIdentifier}{mybox}

\PYG{c}{\PYGZsh{} 2d general triclinic box using a primitive cell for a 2d hex lattice}
\PYG{k}{lattice}\PYG{+w}{       }\PYG{n}{custom}\PYG{+w}{ }\PYG{l+m}{1.0}\PYG{+w}{ }\PYG{n}{a1}\PYG{+w}{ }\PYG{l+m}{1.0}\PYG{+w}{ }\PYG{l+m}{0.0}\PYG{+w}{ }\PYG{l+m}{0.0}\PYG{+w}{ }\PYG{n}{a2}\PYG{+w}{ }\PYG{l+m}{0.5}\PYG{+w}{ }\PYG{l+m}{0.86602540378}\PYG{+w}{ }\PYG{l+m}{0.0}\PYG{+w}{ }\PYG{o}{\PYGZam{}}
\PYG{+w}{              }\PYG{n}{a3}\PYG{+w}{ }\PYG{l+m}{0.0}\PYG{+w}{ }\PYG{l+m}{0.0}\PYG{+w}{ }\PYG{l+m}{1.0}\PYG{+w}{ }\PYG{n}{basis}\PYG{+w}{ }\PYG{l+m}{0.0}\PYG{+w}{ }\PYG{l+m}{0.0}\PYG{+w}{ }\PYG{l+m}{0.0}\PYG{+w}{ }\PYG{n}{triclinic}\PYG{o}{/}\PYG{n}{general}
\PYG{k}{create\PYGZus{}box    }\PYG{n+nv+nvIdentifier}{1}\PYG{+w}{ }\PYG{n+nv+nvIdentifier}{NULL}\PYG{+w}{ }\PYG{l+m}{0}\PYG{+w}{ }\PYG{l+m}{5}\PYG{+w}{ }\PYG{l+m}{0}\PYG{+w}{ }\PYG{l+m}{5}\PYG{+w}{ }\PYG{o}{\PYGZhy{}}\PYG{l+m}{0.5}\PYG{+w}{ }\PYG{l+m}{0.5}
\end{sphinxVerbatim}

\sphinxAtStartPar
Note that for 2d orthogonal or restricted triclinic boxes, the box has
a 3rd dimension which must straddle z = 0.0 in the z dimension.
Typically the width of box in the z dimension should be narrow,
e.g. \sphinxhyphen{}0.5 to 0.5, but that is not required.  For a 2d general
triclinic box, the \sphinxstyleemphasis{a3} vector defined by the {\hyperref[\detokenize{lattice::doc}]{\sphinxcrossref{\DUrole{doc}{lattice}}}}
command must be (0.0,0.0,1.0), which is its default value.  Also the
\sphinxstyleemphasis{clo} and \sphinxstyleemphasis{chi} arguments of the {\hyperref[\detokenize{create_box::doc}]{\sphinxcrossref{\DUrole{doc}{create\_box}}}}
command must be \sphinxhyphen{}0.5 and 0.5.

\sphinxAtStartPar
Here are examples of using the {\hyperref[\detokenize{read_data::doc}]{\sphinxcrossref{\DUrole{doc}{read\_data}}}} command
to define the simulation box for a 2d system via keywords in the
header section of the data file.  These are the same boxes as the examples
for the {\hyperref[\detokenize{create_box::doc}]{\sphinxcrossref{\DUrole{doc}{create\_box}}}} command

\begin{sphinxVerbatim}[commandchars=\\\{\}]
\PYG{c}{\PYGZsh{} 2d orthogonal box}
\PYG{o}{\PYGZhy{}}\PYG{l+m}{10}\PYG{+w}{ }\PYG{l+m}{10}\PYG{+w}{   }\PYG{n}{xlo}\PYG{+w}{ }\PYG{n}{xhi}
\PYG{l+m}{0}\PYG{+w}{ }\PYG{l+m}{10}\PYG{+w}{     }\PYG{n}{ylo}\PYG{+w}{ }\PYG{n}{yhi}
\PYG{o}{\PYGZhy{}}\PYG{l+m}{0.5}\PYG{+w}{ }\PYG{l+m}{0.5}\PYG{+w}{ }\PYG{n}{zlo}\PYG{+w}{ }\PYG{n}{zhi}\PYG{+w}{       }\PYG{c}{\PYGZsh{} this is the default, so no need to specify}

\PYG{c}{\PYGZsh{} 2d restricted triclinic box with only xy tilt}
\PYG{o}{\PYGZhy{}}\PYG{l+m}{10}\PYG{+w}{ }\PYG{l+m}{10}\PYG{+w}{   }\PYG{n}{xlo}\PYG{+w}{ }\PYG{n}{xhi}
\PYG{l+m}{0}\PYG{+w}{ }\PYG{l+m}{10}\PYG{+w}{     }\PYG{n}{ylo}\PYG{+w}{ }\PYG{n}{yhi}
\PYG{o}{\PYGZhy{}}\PYG{l+m}{0.5}\PYG{+w}{ }\PYG{l+m}{0.5}\PYG{+w}{ }\PYG{n}{zlo}\PYG{+w}{ }\PYG{n}{zhi}\PYG{+w}{       }\PYG{c}{\PYGZsh{} this is the default, so no need to specify}
\PYG{l+m}{2.0}\PYG{+w}{ }\PYG{l+m}{0.0}\PYG{+w}{ }\PYG{l+m}{0.0}\PYG{+w}{ }\PYG{n}{xy}\PYG{+w}{ }\PYG{n}{xz}\PYG{+w}{ }\PYG{n}{yz}

\PYG{c}{\PYGZsh{} 3d general triclinic box using a primitive cell for a 2d hex lattice}
\PYG{l+m}{5}\PYG{+w}{ }\PYG{l+m}{0}\PYG{+w}{ }\PYG{l+m}{0}\PYG{+w}{              }\PYG{n}{avec}
\PYG{l+m}{2.5}\PYG{+w}{ }\PYG{l+m}{4.3301270189}\PYG{+w}{ }\PYG{l+m}{0}\PYG{+w}{ }\PYG{n}{bvec}
\PYG{l+m}{0}\PYG{+w}{ }\PYG{l+m}{0}\PYG{+w}{ }\PYG{l+m}{1}\PYG{+w}{              }\PYG{n}{cvec}\PYG{+w}{           }\PYG{c}{\PYGZsh{} this is the default, so no need to specify}
\PYG{l+m}{0}\PYG{+w}{ }\PYG{l+m}{0}\PYG{+w}{ }\PYG{o}{\PYGZhy{}}\PYG{l+m}{0.5}\PYG{+w}{           }\PYG{n}{abc}\PYG{+w}{ }\PYG{n}{origin}\PYG{+w}{     }\PYG{c}{\PYGZsh{} this is the default for 2d, so no need to specify}
\end{sphinxVerbatim}

\sphinxAtStartPar
Note that for 2d orthogonal or restricted triclinic boxes, the box has
a 3rd dimension specified by the \sphinxstyleemphasis{zlo zhi} values, which must straddle
z = 0.0.  Typically the width of box in the z dimension should be
narrow, e.g. \sphinxhyphen{}0.5 to 0.5, but that is not required.  For a 2d general
triclinic box, the z component of \sphinxstyleemphasis{avec} and \sphinxstyleemphasis{bvec} must be zero, and
\sphinxstyleemphasis{cvec} must be (0,0,1), which is the default.  The z component of \sphinxstyleemphasis{abc
origin} must also be \sphinxhyphen{}0.5, which is the default.

\sphinxAtStartPar
If using the {\hyperref[\detokenize{create_atoms::doc}]{\sphinxcrossref{\DUrole{doc}{create\_atoms}}}} command to create
atoms in the 2d simulation box, all the z coordinates of created atoms
will be zero.

\sphinxAtStartPar
If using the {\hyperref[\detokenize{read_data::doc}]{\sphinxcrossref{\DUrole{doc}{read\_data}}}} command to read in a data
file of atom coordinates for a 2d system, the z coordinates of all
atoms should be zero.  A value within epsilon of zero is also allowed
in case the data file was generated by another program with finite
numeric precision, in which case the z coord for the atom will be set
to zero.

\sphinxAtStartPar
Use the {\hyperref[\detokenize{fix_enforce2d::doc}]{\sphinxcrossref{\DUrole{doc}{fix enforce2d}}}} command as the last fix
defined in the input script.  It ensures that the z\sphinxhyphen{}components of
velocities and forces are zeroed out every timestep.  The reason to
make it the last fix is so that any forces added by other fixes will
also be zeroed out.

\sphinxAtStartPar
Many of the example input scripts included in the examples directory
are for 2d models.

\begin{sphinxadmonition}{note}{Note:}
\sphinxAtStartPar
Some models in LAMMPS treat particles as finite\sphinxhyphen{}size spheres, as
opposed to point particles.  See the {\hyperref[\detokenize{atom_style::doc}]{\sphinxcrossref{\DUrole{doc}{atom\_style sphere}}}} and {\hyperref[\detokenize{fix_nve_sphere::doc}]{\sphinxcrossref{\DUrole{doc}{fix nve/sphere}}}} commands
for details.  By default, for 2d simulations, such particles will
still be modeled as 3d spheres, not 2d discs (circles), meaning
their moment of inertia will be that of a sphere.  If you wish to
model them as 2d discs, see the {\hyperref[\detokenize{set::doc}]{\sphinxcrossref{\DUrole{doc}{set density/disc}}}}
command and the \sphinxstyleemphasis{disc} option for the {\hyperref[\detokenize{fix_nve_sphere::doc}]{\sphinxcrossref{\DUrole{doc}{fix nve/sphere}}}}, {\hyperref[\detokenize{fix_nvt_sphere::doc}]{\sphinxcrossref{\DUrole{doc}{fix nvt/sphere}}}},
{\hyperref[\detokenize{fix_nph_sphere::doc}]{\sphinxcrossref{\DUrole{doc}{fix nph/sphere}}}}, {\hyperref[\detokenize{fix_npt_sphere::doc}]{\sphinxcrossref{\DUrole{doc}{fix npt/sphere}}}} commands.
\end{sphinxadmonition}

\sphinxstepscope


\subsection{Type labels}
\label{\detokenize{Howto_type_labels:type-labels}}\label{\detokenize{Howto_type_labels::doc}}
\sphinxAtStartPar
\DUrole{versionmodified,added}{Added in version 15Sep2022.}

\sphinxAtStartPar
Each atom in LAMMPS has an associated numeric atom type. Similarly,
each bond, angle, dihedral, and improper is assigned a bond type,
angle type, and so on.  The primary use of these types is to map
potential (force field) parameters to the interactions of the atom,
bond, angle, dihedral, and improper.

\sphinxAtStartPar
By default, type values are entered as integers from 1 to Ntypes
wherever they appear in LAMMPS input or output files.  The total number
Ntypes for each interaction is “locked in” when the simulation box
is created.

\sphinxAtStartPar
A recent addition to LAMMPS is the option to use strings \sphinxhyphen{} referred to
as type labels \sphinxhyphen{} as an alternative.  Using type labels instead of
numeric types can be advantageous in various scenarios.  For example,
type labels can make inputs more readable and generic (i.e. usable
through the {\hyperref[\detokenize{include::doc}]{\sphinxcrossref{\DUrole{doc}{include command}}}} for different systems with
different numerical values assigned to types.  This generality also
applies to other inputs like data files read by {\hyperref[\detokenize{read_data::doc}]{\sphinxcrossref{\DUrole{doc}{read\_data}}}} or molecule template files read by the {\hyperref[\detokenize{molecule::doc}]{\sphinxcrossref{\DUrole{doc}{molecule}}}} command.  A discussion of the current type label support can
be found in {\hyperref[\detokenize{Howto_type_labels:typelabel24}]{\sphinxcrossref{\DUrole{std,std-ref}{(Gissinger)}}}}.  See below for a list of
other commands that can use type labels in different ways.

\sphinxAtStartPar
LAMMPS will \sphinxstyleemphasis{internally} continue to use numeric types, which means
that many previous restrictions still apply.  For example, the total
number of types is locked in when creating the simulation box, and
potential parameters for each type must be provided even if not used
by any interactions.

\sphinxAtStartPar
A collection of type labels for all type\sphinxhyphen{}kinds (atom types, bond types,
etc.) is stored as a “label map” which is simply a list of numeric types
and their associated type labels.  Within a type\sphinxhyphen{}kind, each type label
must be unique.  It can be assigned to only one numeric type.  To read
and write type labels to data files for a given type\sphinxhyphen{}kind, \sphinxstyleemphasis{all}
associated numeric types need have a type label assigned.  Partial
maps can be saved with the {\hyperref[\detokenize{labelmap::doc}]{\sphinxcrossref{\DUrole{doc}{labelmap write}}}} command
and read back with the {\hyperref[\detokenize{include::doc}]{\sphinxcrossref{\DUrole{doc}{include}}}} command.

\sphinxAtStartPar
Valid type labels can contain most ASCII characters, but cannot start
with a number, a ‘\#’, or a ‘*’.  Also, labels must not contain whitespace
characters.  When using the {\hyperref[\detokenize{labelmap::doc}]{\sphinxcrossref{\DUrole{doc}{labelmap command}}}} in the
LAMMPS input, if certain characters appear in the type label, such as
the single (’) or double (”) quote or the ‘\#’ character, the label
must be put in either double, single, or triple (“””) quotes.  Triple
quotes allow for the most generic type label strings, but they require
to have a leading and trailing blank space.  When defining type labels
the blanks will be ignored. Example:

\begin{sphinxVerbatim}[commandchars=\\\{\}]
\PYG{k}{labelmap}\PYG{+w}{ }\PYG{n}{angle}\PYG{+w}{ }\PYG{l+m}{1}\PYG{+w}{ }\PYG{l+s}{\PYGZdq{}}\PYG{l+s}{\PYGZdq{}}\PYG{l+s}{\PYGZdq{}}\PYG{l+s}{ C1\PYGZsq{}\PYGZhy{}C2}\PYG{l+s}{\PYGZdq{}}\PYG{o}{\PYGZhy{}}\PYG{n}{C3}\PYG{c}{\PYGZsh{} \PYGZdq{}\PYGZdq{}\PYGZdq{}}
\end{sphinxVerbatim}

\sphinxAtStartPar
This command will map the string \sphinxcode{\sphinxupquote{\textasciigrave{}C1\textquotesingle{}\sphinxhyphen{}C2"\sphinxhyphen{}C3\#\textasciigrave{}}} to the angle type 1.

\sphinxAtStartPar
There are two ways to define label maps.  One is via the {\hyperref[\detokenize{labelmap::doc}]{\sphinxcrossref{\DUrole{doc}{labelmap}}}} command.  The other is via the {\hyperref[\detokenize{read_data::doc}]{\sphinxcrossref{\DUrole{doc}{read\_data}}}}
command.  A data file can have sections such as \sphinxstyleemphasis{Atom Type Labels}, \sphinxstyleemphasis{Bond
Type Labels}, etc., which assign type labels to numeric types.  The
label map can be written out to data files by the {\hyperref[\detokenize{write_data::doc}]{\sphinxcrossref{\DUrole{doc}{write\_data}}}} command.  This map is also written to and read from
restart files, by the {\hyperref[\detokenize{write_restart::doc}]{\sphinxcrossref{\DUrole{doc}{write\_restart}}}} and
{\hyperref[\detokenize{read_restart::doc}]{\sphinxcrossref{\DUrole{doc}{read\_restart}}}} commands.


\bigskip\hrule\bigskip



\subsubsection{Use of type labels in LAMMPS input or output}
\label{\detokenize{Howto_type_labels:use-of-type-labels-in-lammps-input-or-output}}
\sphinxAtStartPar
Many LAMMPS input script commands that take a numeric type as an
argument can use the associated type label instead.  If a type label
is not defined for a particular numeric type, only its numeric type
can be used.

\sphinxAtStartPar
This example assigns labels to the atom types, and then uses the type
labels to redefine the pair coefficients.

\begin{sphinxVerbatim}[commandchars=\\\{\}]
\PYG{k}{pair\PYGZus{}coeff}\PYG{+w}{ }\PYG{l+m}{1}\PYG{+w}{ }\PYG{l+m}{2}\PYG{+w}{ }\PYG{l+m}{1.0}\PYG{+w}{ }\PYG{l+m}{1.0}\PYG{+w}{              }\PYG{c}{\PYGZsh{} numeric types}
\PYG{k}{labelmap}\PYG{+w}{ }\PYG{n}{atom}\PYG{+w}{ }\PYG{l+m}{1}\PYG{+w}{ }\PYG{n}{C}\PYG{+w}{ }\PYG{l+m}{2}\PYG{+w}{ }\PYG{n}{H}
\PYG{k}{pair\PYGZus{}coeff}\PYG{+w}{ }\PYG{n}{C}\PYG{+w}{ }\PYG{n}{H}\PYG{+w}{ }\PYG{l+m}{1.0}\PYG{+w}{ }\PYG{l+m}{1.0}\PYG{+w}{              }\PYG{c}{\PYGZsh{} type labels}
\end{sphinxVerbatim}

\sphinxAtStartPar
Adding support for type labels to various commands is an ongoing
project.  If an input script command (or a section in a file read by a
command) allows substituting a type label for a numeric type argument,
it will be explicitly mentioned in that command’s documentation page.

\sphinxAtStartPar
As a temporary measure, input script commands can take advantage of
variables and how they can be expanded during processing of the input.
The variables can use functions that will translate type label strings
to their respective number as defined in the current label map.  See the
{\hyperref[\detokenize{variable::doc}]{\sphinxcrossref{\DUrole{doc}{variable}}}} command for details.

\sphinxAtStartPar
For example, here is how the pair\_coeff command could be used with
type labels if it did not yet support them, either with an explicit
variable command or an implicit variable used in the pair\_coeff
command.

\begin{sphinxVerbatim}[commandchars=\\\{\}]
\PYG{k}{labelmap}\PYG{+w}{ }\PYG{n}{atom}\PYG{+w}{ }\PYG{l+m}{1}\PYG{+w}{ }\PYG{n}{C}\PYG{+w}{ }\PYG{l+m}{2}\PYG{+w}{ }\PYG{n}{H}
\PYG{k}{variable }\PYG{n+nv+nvIdentifier}{atom1}\PYG{+w}{ }\PYG{n}{equal}\PYG{+w}{ }\PYG{n}{label2type}\PYG{n+nv}{(}\PYG{n+nv}{atom,C}\PYG{n+nv}{)}
\PYG{k}{variable }\PYG{n+nv+nvIdentifier}{atom2}\PYG{+w}{ }\PYG{n}{equal}\PYG{+w}{ }\PYG{n}{label2type}\PYG{n+nv}{(}\PYG{n+nv}{atom,H}\PYG{n+nv}{)}
\PYG{k}{pair\PYGZus{}coeff}\PYG{+w}{ }\PYG{n+nv}{\PYGZdl{}\PYGZob{}}\PYG{n+nv}{atom1}\PYG{n+nv}{\PYGZcb{}}\PYG{+w}{ }\PYG{n+nv}{\PYGZdl{}\PYGZob{}}\PYG{n+nv}{atom2}\PYG{n+nv}{\PYGZcb{}}\PYG{+w}{ }\PYG{l+m}{1.0}\PYG{+w}{ }\PYG{l+m}{1.0}
\end{sphinxVerbatim}

\begin{sphinxVerbatim}[commandchars=\\\{\}]
\PYG{k}{labelmap}\PYG{+w}{ }\PYG{n}{atom}\PYG{+w}{ }\PYG{l+m}{1}\PYG{+w}{ }\PYG{n}{C}\PYG{+w}{ }\PYG{l+m}{2}\PYG{+w}{ }\PYG{n}{H}
\PYG{k}{pair\PYGZus{}coeff}\PYG{+w}{ }\PYG{n+nv}{\PYGZdl{}(}\PYG{n+nv}{label2type}\PYG{n+nv}{(}\PYG{n+nv}{atom,C}\PYG{n+nv}{)}\PYG{n+nv}{)}\PYG{+w}{ }\PYG{n+nv}{\PYGZdl{}(}\PYG{n+nv}{label2type}\PYG{n+nv}{(}\PYG{n+nv}{atom,H}\PYG{n+nv}{)}\PYG{n+nv}{)}\PYG{+w}{ }\PYG{l+m}{80.0}\PYG{+w}{ }\PYG{l+m}{1.2}
\end{sphinxVerbatim}


\bigskip\hrule\bigskip



\subsubsection{Commands that can use label types}
\label{\detokenize{Howto_type_labels:commands-that-can-use-label-types}}
\sphinxAtStartPar
Any workflow that involves reading multiple data files, molecule
templates or a combination of the two can be streamlined by using type
labels instead of numeric types, because types are automatically synced
between the files.  The creation of simulation\sphinxhyphen{}ready reaction templates
for {\hyperref[\detokenize{fix_bond_react::doc}]{\sphinxcrossref{\DUrole{doc}{fix bond/react}}}} is much simpler when using
type labels, and results in templates that can be used without
modification in multiple simulations or different systems.


\bigskip\hrule\bigskip

\phantomsection\label{\detokenize{Howto_type_labels:typelabel24}}
\sphinxAtStartPar
\sphinxstylestrong{(Gissinger)} J. R. Gissinger, I. Nikiforov, Y. Afshar, B. Waters, M. Choi, D. S. Karls, A. Stukowski, W. Im, H. Heinz, A. Kohlmeyer, and E. B. Tadmor, J Phys Chem B, 128, 3282\sphinxhyphen{}3297 (2024).

\sphinxstepscope


\subsection{Triclinic (non\sphinxhyphen{}orthogonal) simulation boxes}
\label{\detokenize{Howto_triclinic:triclinic-non-orthogonal-simulation-boxes}}\label{\detokenize{Howto_triclinic::doc}}
\sphinxAtStartPar
By default, LAMMPS uses an orthogonal simulation box to encompass the
particles.  The orthogonal box has its “origin” at (xlo,ylo,zlo) and
extends to (xhi,yhi,zhi).  Conceptually it is defined by 3 edge
vectors starting from the origin given by \sphinxstylestrong{A} = (xhi\sphinxhyphen{}xlo,0,0); \sphinxstylestrong{B}
= (0,yhi\sphinxhyphen{}ylo,0); \sphinxstylestrong{C} = (0,0,zhi\sphinxhyphen{}zlo).  The {\hyperref[\detokenize{boundary::doc}]{\sphinxcrossref{\DUrole{doc}{boundary}}}} command sets the boundary conditions for the 6 faces of
the box (periodic, non\sphinxhyphen{}periodic, etc).  The 6 parameters
(xlo,xhi,ylo,yhi,zlo,zhi) are defined at the time the simulation box
is created by one of these commands:
\begin{itemize}
\item {} 
\sphinxAtStartPar
{\hyperref[\detokenize{create_box::doc}]{\sphinxcrossref{\DUrole{doc}{create\_box}}}}

\item {} 
\sphinxAtStartPar
{\hyperref[\detokenize{read_data::doc}]{\sphinxcrossref{\DUrole{doc}{read\_data}}}}

\item {} 
\sphinxAtStartPar
{\hyperref[\detokenize{read_restart::doc}]{\sphinxcrossref{\DUrole{doc}{read\_restart}}}}

\item {} 
\sphinxAtStartPar
{\hyperref[\detokenize{read_dump::doc}]{\sphinxcrossref{\DUrole{doc}{read\_dump}}}}

\end{itemize}

\sphinxAtStartPar
Internally, LAMMPS defines box size parameters lx,ly,lz where lx =
xhi\sphinxhyphen{}xlo, and similarly in the y and z dimensions.  The 6 parameters, as
well as lx,ly,lz, can be output via the {\hyperref[\detokenize{thermo_style::doc}]{\sphinxcrossref{\DUrole{doc}{thermo\_style custom}}}} command.  See the {\hyperref[\detokenize{Howto_2d::doc}]{\sphinxcrossref{\DUrole{doc}{Howto 2d}}}} doc page
for info on how zlo and zhi are defined for 2d simulations.


\bigskip\hrule\bigskip



\subsubsection{Triclinic simulation boxes}
\label{\detokenize{Howto_triclinic:triclinic-simulation-boxes}}
\sphinxAtStartPar
LAMMPS also allows simulations to be performed using triclinic
(non\sphinxhyphen{}orthogonal) simulation boxes shaped as a 3d parallelepiped with
triclinic symmetry.  For 2d simulations a triclinic simulation box is
effectively a parallelogram; see the {\hyperref[\detokenize{Howto_2d::doc}]{\sphinxcrossref{\DUrole{doc}{Howto 2d}}}} doc
page for details.

\sphinxAtStartPar
One use of triclinic simulation boxes is to model solid\sphinxhyphen{}state crystals
with triclinic symmetry.  The {\hyperref[\detokenize{lattice::doc}]{\sphinxcrossref{\DUrole{doc}{lattice}}}} command can be
used with non\sphinxhyphen{}orthogonal basis vectors to define a lattice that will
tile a triclinic simulation box via the {\hyperref[\detokenize{create_atoms::doc}]{\sphinxcrossref{\DUrole{doc}{create\_atoms}}}} command.

\sphinxAtStartPar
A second use is to run Parrinello\sphinxhyphen{}Rahman dynamics via the {\hyperref[\detokenize{fix_nh::doc}]{\sphinxcrossref{\DUrole{doc}{fix
npt}}}} command, which will adjust the xy, xz, yz tilt factors
to compensate for off\sphinxhyphen{}diagonal components of the pressure tensor.  The
analog for an {\hyperref[\detokenize{minimize::doc}]{\sphinxcrossref{\DUrole{doc}{energy minimization}}}} is the {\hyperref[\detokenize{fix_box_relax::doc}]{\sphinxcrossref{\DUrole{doc}{fix
box/relax}}}} command.

\sphinxAtStartPar
A third use is to shear a bulk solid to study the response of the
material.  The {\hyperref[\detokenize{fix_deform::doc}]{\sphinxcrossref{\DUrole{doc}{fix deform}}}} command can be used for
this purpose.  It allows dynamic control of the xy, xz, yz tilt
factors as a simulation runs.  This is discussed in the {\hyperref[\detokenize{Howto_nemd::doc}]{\sphinxcrossref{\DUrole{doc}{Howto
NEMD}}}} doc page on non\sphinxhyphen{}equilibrium MD (NEMD) simulations.

\sphinxAtStartPar
Conceptually, a triclinic parallelepiped is defined with an “origin”
at (xlo,ylo,zhi) and 3 edge vectors \sphinxstylestrong{A} = (ax,ay,az), \sphinxstylestrong{B} =
(bx,by,bz), \sphinxstylestrong{C} = (cx,cy,cz) which can be arbitrary vectors, so long
as they are non\sphinxhyphen{}zero, distinct, and not co\sphinxhyphen{}planar.  In addition, they
must define a right\sphinxhyphen{}handed system, such that (\sphinxstylestrong{A} cross \sphinxstylestrong{B})
points in the direction of \sphinxstylestrong{C}.  Note that a left\sphinxhyphen{}handed system can
be converted to a right\sphinxhyphen{}handed system by simply swapping the order of
any pair of the \sphinxstylestrong{A}, \sphinxstylestrong{B}, \sphinxstylestrong{C} vectors.

\sphinxAtStartPar
The 4 commands listed above for defining orthogonal simulation boxes
have triclinic options which allow for specification of the origin and
edge vectors \sphinxstylestrong{A}, \sphinxstylestrong{B}, \sphinxstylestrong{C}.  For each command, this can be done
in one of two ways, for what LAMMPS calls a \sphinxstyleemphasis{general} triclinic box or
a \sphinxstyleemphasis{restricted} triclinic box.

\sphinxAtStartPar
A \sphinxstyleemphasis{general} triclinic box is specified by an origin (xlo, ylo, zlo)
and arbitrary edge vectors \sphinxstylestrong{A} = (ax,ay,az), \sphinxstylestrong{B} = (bx,by,bz), and
\sphinxstylestrong{C} = (cx,cy,cz).  So there are 12 parameters in total.

\sphinxAtStartPar
A \sphinxstyleemphasis{restricted} triclinic box also has an origin (xlo,ylo,zlo), but its
edge vectors are of the following restricted form: \sphinxstylestrong{A} =
(xhi\sphinxhyphen{}xlo,0,0), \sphinxstylestrong{B} = (xy,yhi\sphinxhyphen{}ylo,0), \sphinxstylestrong{C} = (xz,yz,zhi\sphinxhyphen{}zlo).  So
there are 9 parameters in total.  Note that the restricted form
requires \sphinxstylestrong{A} to be along the x\sphinxhyphen{}axis, \sphinxstylestrong{B} to be in the xy plane
with a y\sphinxhyphen{}component in the +y direction, and \sphinxstylestrong{C} to have its
z\sphinxhyphen{}component in the +z direction.  Note that a restricted triclinic box
is \sphinxstyleemphasis{right\sphinxhyphen{}handed} by construction since (\sphinxstylestrong{A} cross \sphinxstylestrong{B}) points in
the direction of \sphinxstylestrong{C}.

\sphinxAtStartPar
The \sphinxstyleemphasis{xy,xz,yz} values can be zero or positive or negative.  They are
called “tilt factors” because they are the amount of displacement
applied to edges of faces of an orthogonal box to change it into a
restricted triclinic parallelepiped.

\begin{sphinxadmonition}{note}{Note:}
\sphinxAtStartPar
Any right\sphinxhyphen{}handed general triclinic box (i.e. solid\sphinxhyphen{}state crystal
basis vectors) can be rotated in 3d around its origin in order to
conform to the LAMMPS definition of a restricted triclinic box.
See the discussion in the next sub\sphinxhyphen{}section about general triclinic
simulation boxes in LAMMPS.
\end{sphinxadmonition}

\sphinxAtStartPar
Note that the {\hyperref[\detokenize{thermo_style::doc}]{\sphinxcrossref{\DUrole{doc}{thermo\_style custom}}}} command has
keywords for outputting the various parameters that define the size
and shape of orthogonal, restricted triclinic, and general triclinic
simulation boxes.

\sphinxAtStartPar
For orthogonal boxes there 6 thermo keywords (xlo,ylo,zlo) and
(xhi,yhi,zhi).

\sphinxAtStartPar
For restricted triclinic boxes there are 9 thermo keywords for
(xlo,ylo,zlo), (xhi,yhi,zhi), and the (xy,xz,yz) tilt factors.

\sphinxAtStartPar
For general triclinic boxes there are 12 thermo keywords for
(xlo,ylo,zhi) and the components of the \sphinxstylestrong{A}, \sphinxstylestrong{B}, \sphinxstylestrong{C} edge
vectors, namely (avecx,avecy,avecz), (bvecx,bvecy,bvecz), and
(cvecx,cvecy,cvecz),

\sphinxAtStartPar
The remainder of this doc page explains (a) how LAMMPS operates with
general triclinic simulation boxes, (b) mathematical transformations
between general and restricted triclinic boxes which may be useful
when creating LAMMPS inputs or interpreting outputs for triclinic
simulations, and (c) how LAMMPS uses tilt factors for restricted
triclinic simulation boxes.


\bigskip\hrule\bigskip



\subsubsection{General triclinic simulation boxes in LAMMPS}
\label{\detokenize{Howto_triclinic:general-triclinic-simulation-boxes-in-lammps}}
\sphinxAtStartPar
LAMMPS allows specification of general triclinic simulation boxes with
their atoms as a convenience for users who may be converting data from
solid\sphinxhyphen{}state crystallographic representations or from DFT codes for
input to LAMMPS.  Likewise it allows output of dump files, data files,
and thermodynamic data (e.g. pressure tensor) in a general triclinic
format.

\sphinxAtStartPar
However internally, LAMMPS only uses restricted triclinic simulation
boxes.  This is for parallel efficiency and to formulate partitioning
of the simulation box across processors, neighbor list building, and
inter\sphinxhyphen{}processor communication of per\sphinxhyphen{}atom data with methods similar to
those used for orthogonal boxes.

\sphinxAtStartPar
This means 4 things which are important to understand:
\begin{itemize}
\item {} 
\sphinxAtStartPar
Input of a general triclinic system is immediately converted to a
restricted triclinic system.

\item {} 
\sphinxAtStartPar
If output of per\sphinxhyphen{}atom data for a general triclinic system is
requested (e.g. for atom coordinates in a dump file),
conversion from a restricted to general triclinic system is done at
the time of output.

\item {} 
\sphinxAtStartPar
The conversion of the simulation box and per\sphinxhyphen{}atom data from general
triclinic to restricted triclinic (and vice versa) is a 3d rotation
operation around an origin, which is the lower left corner of the
simulation box.  This means an input data file for a general
triclinic system should specify all per\sphinxhyphen{}atom quantities consistent
with the general triclinic box and its orientation relative to the
standard x,y,z coordinate axes.  For example, atom coordinates
should be inside the general triclinic simulation box defined by the
edge vectors \sphinxstylestrong{A}, \sphinxstylestrong{B}, \sphinxstylestrong{C} and its origin.  Likewise per\sphinxhyphen{}atom
velocities should be in directions consistent with the general
triclinic box orientation.  E.g. a velocity vector which will be in
the +x direction once LAMMPS converts from a general to restricted
triclinic box, should be specified in the data file in the direction
of the \sphinxstylestrong{A} edge vector.  See the {\hyperref[\detokenize{read_data::doc}]{\sphinxcrossref{\DUrole{doc}{read\_data}}}} doc
page for info on all the per\sphinxhyphen{}atom vector quantities to which this
rule applies when a data file for a general triclinic box is input.

\item {} 
\sphinxAtStartPar
If commands such as {\hyperref[\detokenize{write_data::doc}]{\sphinxcrossref{\DUrole{doc}{write\_data}}}} or {\hyperref[\detokenize{dump::doc}]{\sphinxcrossref{\DUrole{doc}{dump
custom}}}} are used to output general triclinic information, it
is effectively the inverse of the operation described in the
preceding bullet.

\item {} 
\sphinxAtStartPar
Other LAMMPS commands such as {\hyperref[\detokenize{region::doc}]{\sphinxcrossref{\DUrole{doc}{region}}}} or
{\hyperref[\detokenize{velocity::doc}]{\sphinxcrossref{\DUrole{doc}{velocity}}}} or {\hyperref[\detokenize{set::doc}]{\sphinxcrossref{\DUrole{doc}{set}}}}, operate on a
restricted triclinic system even if a general triclinic system was
defined initially.

\end{itemize}

\sphinxAtStartPar
This is the list of commands which have general triclinic options:
\begin{itemize}
\item {} 
\sphinxAtStartPar
{\hyperref[\detokenize{create_box::doc}]{\sphinxcrossref{\DUrole{doc}{create\_box}}}} \sphinxhyphen{} define a general triclinic box

\item {} 
\sphinxAtStartPar
{\hyperref[\detokenize{create_atoms::doc}]{\sphinxcrossref{\DUrole{doc}{create\_atoms}}}} \sphinxhyphen{} add atoms to a general triclinic box

\item {} 
\sphinxAtStartPar
{\hyperref[\detokenize{lattice::doc}]{\sphinxcrossref{\DUrole{doc}{lattice}}}} \sphinxhyphen{} define a custom lattice consistent with the \sphinxstylestrong{A}, \sphinxstylestrong{B}, \sphinxstylestrong{C} edge vectors of a general triclinic box

\item {} 
\sphinxAtStartPar
{\hyperref[\detokenize{read_data::doc}]{\sphinxcrossref{\DUrole{doc}{read\_data}}}} \sphinxhyphen{} read a data file for a general triclinic system

\item {} 
\sphinxAtStartPar
{\hyperref[\detokenize{write_data::doc}]{\sphinxcrossref{\DUrole{doc}{write\_data}}}} \sphinxhyphen{} write a data file for a general triclinic system

\item {} 
\sphinxAtStartPar
{\hyperref[\detokenize{dump::doc}]{\sphinxcrossref{\DUrole{doc}{dump atom, dump custom}}}} \sphinxhyphen{} output dump snapshots in general triclinic format

\item {} 
\sphinxAtStartPar
{\hyperref[\detokenize{dump_modify::doc}]{\sphinxcrossref{\DUrole{doc}{dump\_modify triclinic/general}}}} \sphinxhyphen{} select general triclinic format for dump output

\item {} 
\sphinxAtStartPar
{\hyperref[\detokenize{thermo_style::doc}]{\sphinxcrossref{\DUrole{doc}{thermo\_style}}}} \sphinxhyphen{} output the pressure tensor in
general triclinic format

\item {} 
\sphinxAtStartPar
{\hyperref[\detokenize{thermo_modify::doc}]{\sphinxcrossref{\DUrole{doc}{thermo\_modify triclinic/general}}}} \sphinxhyphen{} select general triclinic format for thermo output

\item {} 
\sphinxAtStartPar
{\hyperref[\detokenize{read_restart::doc}]{\sphinxcrossref{\DUrole{doc}{read\_restart}}}} \sphinxhyphen{} read a restart file for a general triclinic system

\item {} 
\sphinxAtStartPar
{\hyperref[\detokenize{read_restart::doc}]{\sphinxcrossref{\DUrole{doc}{write\_restart}}}} \sphinxhyphen{} write a restart file for a general triclinic system

\end{itemize}


\bigskip\hrule\bigskip



\subsubsection{Transformation from general to restricted triclinic boxes}
\label{\detokenize{Howto_triclinic:transformation-from-general-to-restricted-triclinic-boxes}}
\sphinxAtStartPar
Let \sphinxstylestrong{A},\sphinxstylestrong{B},\sphinxstylestrong{C} be the right\sphinxhyphen{}handed edge vectors of a
general triclinic simulation box.  The equivalent LAMMPS \sphinxstylestrong{a},\sphinxstylestrong{b},\sphinxstylestrong{c} for a restricted triclinic box are a 3d rotation of
\sphinxstylestrong{A}, \sphinxstylestrong{B}, and \sphinxstylestrong{C} and can be computed as follows:
\begin{equation*}
\begin{split}\begin{pmatrix} \mathbf{a}  & \mathbf{b}  & \mathbf{c} \end{pmatrix} = &
\begin{pmatrix}
  a_x & b_x & c_x \\
  0   & b_y & c_y \\
  0   & 0   & c_z \\
\end{pmatrix} \\
a_x = & A \\
b_x = & \mathbf{B} \cdot \mathbf{\hat{A}} \quad = \quad B \cos{\gamma} \\
b_y = & |\mathbf{\hat{A}} \times \mathbf{B}| \quad = \quad B \sin{\gamma} \quad =  \quad  \sqrt{B^2 - {b_x}^2} \\
c_x = & \mathbf{C} \cdot \mathbf{\hat{A}} \quad = \quad C \cos{\beta} \\
c_y = & \mathbf{C} \cdot \widehat{(\mathbf{A} \times \mathbf{B})} \times \mathbf{\hat{A}} \quad = \quad \frac{\mathbf{B} \cdot \mathbf{C} - b_x c_x}{b_y} \\
c_z = & |\mathbf{C} \cdot \widehat{(\mathbf{A} \times \mathbf{B})}|\quad = \quad \sqrt{C^2 - {c_x}^2 - {c_y}^2}\end{split}
\end{equation*}
\sphinxAtStartPar
where A = | \sphinxstylestrong{A} | indicates the scalar length of \sphinxstylestrong{A}. The hat
symbol (\textasciicircum{}) indicates the corresponding unit vector. \(\beta\) and
\(\gamma\) are angles between the \sphinxstylestrong{A}, \sphinxstylestrong{B}, \sphinxstylestrong{C} vectors
as described below.

\sphinxAtStartPar
For consistency, the same rotation applied to the triclinic box edge
vectors can also be applied to atom positions, velocities, and other
vector quantities.  This can be conveniently achieved by first
converting to fractional coordinates in the general triclinic
coordinates and then converting to coordinates in the restricted
triclinic basis.  The transformation is given by the following equation:
\begin{equation*}
\begin{split}\mathbf{x} = & \begin{pmatrix} \mathbf{a}  & \mathbf{b}  & \mathbf{c} \end{pmatrix} \cdot \frac{1}{V}
  \begin{pmatrix}
    \mathbf{B \times C}  \\
    \mathbf{C \times A}  \\
    \mathbf{A \times B}
  \end{pmatrix} \cdot \mathbf{X}\end{split}
\end{equation*}
\sphinxAtStartPar
where \sphinxstyleemphasis{V} is the volume of the box (same in either basis), \sphinxstylestrong{X} is
the fractional vector in the general triclinic basis and \sphinxstylestrong{x} is the
resulting vector in the restricted triclinic basis.


\bigskip\hrule\bigskip



\subsubsection{Crystallographic general triclinic representation of a simulation box}
\label{\detokenize{Howto_triclinic:crystallographic-general-triclinic-representation-of-a-simulation-box}}
\sphinxAtStartPar
General triclinic crystal structures are often defined using three
lattice constants \sphinxstyleemphasis{a}, \sphinxstyleemphasis{b}, and \sphinxstyleemphasis{c}, and three angles \(\alpha\),
\(\beta\), and \(\gamma\). Note that in this nomenclature, the
a, b, and c lattice constants are the scalar lengths of the edge
vectors \sphinxstylestrong{a}, \sphinxstylestrong{b}, and \sphinxstylestrong{c} defined above.  The relationship
between these 6 quantities (a, b, c, \(\alpha\), \(\beta\),
\(\gamma\)) and the LAMMPS restricted triclinic box sizes
(lx,ly,lz) = (xhi\sphinxhyphen{}xlo,yhi\sphinxhyphen{}ylo,zhi\sphinxhyphen{}zlo) and tilt factors (xy,xz,yz) is
as follows:
\begin{equation*}
\begin{split}a   = & {\rm lx} \\
b^2 = &  {\rm ly}^2 +  {\rm xy}^2 \\
c^2 = &  {\rm lz}^2 +  {\rm xz}^2 +  {\rm yz}^2 \\
\cos{\alpha} = & \frac{{\rm xy}*{\rm xz} + {\rm ly}*{\rm yz}}{b*c} \\
\cos{\beta}  = & \frac{\rm xz}{c} \\
\cos{\gamma} = & \frac{\rm xy}{b} \\\end{split}
\end{equation*}
\sphinxAtStartPar
The inverse relationship can be written as follows:
\begin{equation*}
\begin{split}{\rm lx}   = & a \\
{\rm xy}   = & b \cos{\gamma}  \\
{\rm xz}   = & c \cos{\beta}\\
{\rm ly}^2 = & b^2 - {\rm xy}^2 \\
{\rm yz}   = & \frac{b*c \cos{\alpha} - {\rm xy}*{\rm xz}}{\rm ly} \\
{\rm lz}^2 = & c^2 - {\rm xz}^2 - {\rm yz}^2 \\\end{split}
\end{equation*}
\sphinxAtStartPar
The values of \sphinxstyleemphasis{a}, \sphinxstyleemphasis{b}, \sphinxstyleemphasis{c}, \(\alpha\) , \(\beta\), and
\(\gamma\) can be printed out or accessed by computes using the
{\hyperref[\detokenize{thermo_style::doc}]{\sphinxcrossref{\DUrole{doc}{thermo\_style custom}}}} keywords \sphinxstyleemphasis{cella}, \sphinxstyleemphasis{cellb},
\sphinxstyleemphasis{cellc}, \sphinxstyleemphasis{cellalpha}, \sphinxstyleemphasis{cellbeta}, \sphinxstyleemphasis{cellgamma}, respectively.


\bigskip\hrule\bigskip



\subsubsection{Output of restricted and general triclinic boxes in a dump file}
\label{\detokenize{Howto_triclinic:output-of-restricted-and-general-triclinic-boxes-in-a-dump-file}}
\sphinxAtStartPar
As discussed on the {\hyperref[\detokenize{dump::doc}]{\sphinxcrossref{\DUrole{doc}{dump}}}} command doc page, when the BOX
BOUNDS for a snapshot is written to a dump file for a restricted
triclinic box, an orthogonal bounding box which encloses the triclinic
simulation box is output, along with the 3 tilt factors (xy, xz, yz) of
the restricted triclinic box, formatted as follows:

\begin{sphinxVerbatim}[commandchars=\\\{\}]
ITEM: BOX BOUNDS xy xz yz
xlo\PYGZus{}bound xhi\PYGZus{}bound xy
ylo\PYGZus{}bound yhi\PYGZus{}bound xz
zlo\PYGZus{}bound zhi\PYGZus{}bound yz
\end{sphinxVerbatim}

\sphinxAtStartPar
This bounding box is convenient for many visualization programs and is
calculated from the 9 restricted triclinic box parameters
(xlo,xhi,ylo,yhi,zlo,zhi,xy,xz,yz) as follows:

\begin{sphinxVerbatim}[commandchars=\\\{\}]
xlo\PYGZus{}bound = xlo + MIN(0.0,xy,xz,xy+xz)
xhi\PYGZus{}bound = xhi + MAX(0.0,xy,xz,xy+xz)
ylo\PYGZus{}bound = ylo + MIN(0.0,yz)
yhi\PYGZus{}bound = yhi + MAX(0.0,yz)
zlo\PYGZus{}bound = zlo
zhi\PYGZus{}bound = zhi
\end{sphinxVerbatim}

\sphinxAtStartPar
These formulas can be inverted if you need to convert the bounding box
back into the restricted triclinic box parameters, e.g. xlo =
xlo\_bound \sphinxhyphen{} MIN(0.0,xy,xz,xy+xz).


\bigskip\hrule\bigskip



\subsubsection{Periodicity and tilt factors for triclinic simulation boxes}
\label{\detokenize{Howto_triclinic:periodicity-and-tilt-factors-for-triclinic-simulation-boxes}}
\sphinxAtStartPar
There is no requirement that a triclinic box be periodic in any
dimension, though it typically should be in y or z if you wish to
enforce a shift in coordinates due to periodic boundary conditions
across the y or z boundaries.  See the doc page for the {\hyperref[\detokenize{boundary::doc}]{\sphinxcrossref{\DUrole{doc}{boundary}}}} command for an explanation of shifted coordinates for
restricted triclinic boxes which are periodic.

\sphinxAtStartPar
Some commands that work with triclinic boxes, e.g. the {\hyperref[\detokenize{fix_deform::doc}]{\sphinxcrossref{\DUrole{doc}{fix
deform}}}} and {\hyperref[\detokenize{fix_nh::doc}]{\sphinxcrossref{\DUrole{doc}{fix npt}}}} commands, require
periodicity or non\sphinxhyphen{}shrink\sphinxhyphen{}wrap boundary conditions in specific
dimensions.  See the command doc pages for details.

\sphinxAtStartPar
A restricted triclinic box can be defined with all 3 tilt factors =
0.0, so that it is initially orthogonal.  This is necessary if the box
will become non\sphinxhyphen{}orthogonal, e.g. due to use of the {\hyperref[\detokenize{fix_nh::doc}]{\sphinxcrossref{\DUrole{doc}{fix npt}}}} or {\hyperref[\detokenize{fix_deform::doc}]{\sphinxcrossref{\DUrole{doc}{fix deform}}}} commands.  Alternatively,
you can use the {\hyperref[\detokenize{change_box::doc}]{\sphinxcrossref{\DUrole{doc}{change\_box}}}} command to convert a
simulation box from orthogonal to restricted triclinic and vice versa.

\begin{sphinxadmonition}{note}{Note:}
\sphinxAtStartPar
Highly tilted restricted triclinic simulation boxes can be
computationally inefficient.  This is due to the large volume of
communication needed to acquire ghost atoms around a processor’s
irregular\sphinxhyphen{}shaped subdomain.  For extreme values of tilt, LAMMPS may
also lose atoms and generate an error.
\end{sphinxadmonition}

\sphinxAtStartPar
LAMMPS will issue a warning if you define a restricted triclinic box
with a tilt factor which skews the box more than half the distance of
the parallel box length, which is the first dimension in the tilt
factor (e.g. x for xz).

\sphinxAtStartPar
For example, if xlo = 2 and xhi = 12, then the x box length is 10 and
the xy tilt factor should be between \sphinxhyphen{}5 and 5 to avoid the warning.
Similarly, both xz and yz should be between \sphinxhyphen{}(xhi\sphinxhyphen{}xlo)/2 and
+(yhi\sphinxhyphen{}ylo)/2.  Note that these are not limitations, since if the
maximum tilt factor is 5 (as in this example), then simulations boxes
and atom configurations with tilt = …, \sphinxhyphen{}15, \sphinxhyphen{}5, 5, 15, 25, … are
all geometrically equivalent.

\sphinxAtStartPar
If the box tilt exceeds this limit during a dynamics run (e.g. due to
the {\hyperref[\detokenize{fix_deform::doc}]{\sphinxcrossref{\DUrole{doc}{fix deform}}}} command), then by default the box
is “flipped” to an equivalent shape with a tilt factor within the
warning bounds, and the run continues.  See the {\hyperref[\detokenize{fix_deform::doc}]{\sphinxcrossref{\DUrole{doc}{fix deform}}}} page for further details.  Box flips that would normally
occur using the {\hyperref[\detokenize{fix_deform::doc}]{\sphinxcrossref{\DUrole{doc}{fix deform}}}} or {\hyperref[\detokenize{fix_nh::doc}]{\sphinxcrossref{\DUrole{doc}{fix npt}}}} commands can be suppressed using the \sphinxstyleemphasis{flip no} option with
either of the commands.

\sphinxAtStartPar
One exception to box flipping is if the first dimension in the tilt
factor (e.g. x for xy) is non\sphinxhyphen{}periodic.  In that case, the limits on
the tilt factor are not enforced, since flipping the box in that
dimension would not change the atom positions due to non\sphinxhyphen{}periodicity.
In this mode, if the system tilts to large angles, the simulation will
simply become inefficient, due to the highly skewed simulation box.

\sphinxstepscope


\subsection{Thermostats}
\label{\detokenize{Howto_thermostat:thermostats}}\label{\detokenize{Howto_thermostat::doc}}
\sphinxAtStartPar
Thermostatting means controlling the temperature of particles in an MD
simulation.  {\hyperref[\detokenize{Howto_barostat::doc}]{\sphinxcrossref{\DUrole{doc}{Barostatting}}}} means controlling
the pressure.  Since the pressure includes a kinetic component due to
particle velocities, both these operations require calculation of the
temperature.  Typically a target temperature (T) and/or pressure (P)
is specified by the user, and the thermostat or barostat attempts to
equilibrate the system to the requested T and/or P.

\sphinxAtStartPar
Thermostatting in LAMMPS is performed by {\hyperref[\detokenize{fix::doc}]{\sphinxcrossref{\DUrole{doc}{fixes}}}}, or in one
case by a pair style.  Several thermostatting fixes are available:
Nose\sphinxhyphen{}Hoover (nvt), Berendsen, CSVR, Langevin, and direct rescaling
(temp/rescale).  Dissipative particle dynamics (DPD) thermostatting
can be invoked via the \sphinxstyleemphasis{dpd/tstat} pair style:
\begin{itemize}
\item {} 
\sphinxAtStartPar
{\hyperref[\detokenize{fix_nh::doc}]{\sphinxcrossref{\DUrole{doc}{fix nvt}}}}

\item {} 
\sphinxAtStartPar
{\hyperref[\detokenize{fix_nvt_sphere::doc}]{\sphinxcrossref{\DUrole{doc}{fix nvt/sphere}}}}

\item {} 
\sphinxAtStartPar
{\hyperref[\detokenize{fix_nvt_asphere::doc}]{\sphinxcrossref{\DUrole{doc}{fix nvt/asphere}}}}

\item {} 
\sphinxAtStartPar
{\hyperref[\detokenize{fix_nvt_sllod::doc}]{\sphinxcrossref{\DUrole{doc}{fix nvt/sllod}}}}

\item {} 
\sphinxAtStartPar
{\hyperref[\detokenize{fix_temp_berendsen::doc}]{\sphinxcrossref{\DUrole{doc}{fix temp/berendsen}}}}

\item {} 
\sphinxAtStartPar
{\hyperref[\detokenize{fix_temp_csvr::doc}]{\sphinxcrossref{\DUrole{doc}{fix temp/csvr}}}}

\item {} 
\sphinxAtStartPar
{\hyperref[\detokenize{fix_langevin::doc}]{\sphinxcrossref{\DUrole{doc}{fix langevin}}}}

\item {} 
\sphinxAtStartPar
{\hyperref[\detokenize{fix_temp_rescale::doc}]{\sphinxcrossref{\DUrole{doc}{fix temp/rescale}}}}

\item {} 
\sphinxAtStartPar
{\hyperref[\detokenize{pair_dpd::doc}]{\sphinxcrossref{\DUrole{doc}{pair\_style dpd/tstat}}}}

\end{itemize}

\sphinxAtStartPar
{\hyperref[\detokenize{fix_nh::doc}]{\sphinxcrossref{\DUrole{doc}{Fix nvt}}}} only thermostats the translational velocity of
particles.  {\hyperref[\detokenize{fix_nvt_sllod::doc}]{\sphinxcrossref{\DUrole{doc}{Fix nvt/sllod}}}} also does this,
except that it subtracts out a velocity bias due to a deforming box
and integrates the SLLOD equations of motion.  See the {\hyperref[\detokenize{Howto_nemd::doc}]{\sphinxcrossref{\DUrole{doc}{Howto
nemd}}}} page for further details.  {\hyperref[\detokenize{fix_nvt_sphere::doc}]{\sphinxcrossref{\DUrole{doc}{Fix nvt/sphere}}}} and {\hyperref[\detokenize{fix_nvt_asphere::doc}]{\sphinxcrossref{\DUrole{doc}{fix nvt/asphere}}}}
thermostat not only translation velocities but also rotational
velocities for spherical and aspherical particles.

\begin{sphinxadmonition}{note}{Note:}
\sphinxAtStartPar
A recent (2017) book by {\hyperref[\detokenize{Howto_thermostat:daivis-thermostat}]{\sphinxcrossref{\DUrole{std,std-ref}{(Daivis and Todd)}}}}
discusses use of the SLLOD method and non\sphinxhyphen{}equilibrium MD (NEMD)
thermostatting generally, for both simple and complex fluids,
e.g. molecular systems.  The latter can be tricky to do correctly.
\end{sphinxadmonition}

\sphinxAtStartPar
DPD thermostatting alters pairwise interactions in a manner analogous
to the per\sphinxhyphen{}particle thermostatting of {\hyperref[\detokenize{fix_langevin::doc}]{\sphinxcrossref{\DUrole{doc}{fix langevin}}}}.

\sphinxAtStartPar
Any of the thermostatting fixes can be instructed to use custom
temperature computes that remove bias which has two effects: first,
the current calculated temperature, which is compared to the requested
target temperature, is calculated with the velocity bias removed;
second, the thermostat adjusts only the thermal temperature component
of the particle’s velocities, which are the velocities with the bias
removed.  The removed bias is then added back to the adjusted
velocities.  See the doc pages for the individual fixes and for the
{\hyperref[\detokenize{fix_modify::doc}]{\sphinxcrossref{\DUrole{doc}{fix\_modify}}}} command for instructions on how to
assign a temperature compute to a thermostatting fix.

\sphinxAtStartPar
For example, you can apply a thermostat only to atoms in a spatial
region by using it in conjunction with {\hyperref[\detokenize{compute_temp_region::doc}]{\sphinxcrossref{\DUrole{doc}{compute temp/region}}}}.  Or you can apply a thermostat to only the x
and z components of velocity by using it with {\hyperref[\detokenize{compute_temp_partial::doc}]{\sphinxcrossref{\DUrole{doc}{compute
temp/partial}}}}.  Of you could thermostat only
the thermal temperature of a streaming flow of particles without
affecting the streaming velocity, by using {\hyperref[\detokenize{compute_temp_profile::doc}]{\sphinxcrossref{\DUrole{doc}{compute temp/profile}}}}.

\sphinxAtStartPar
Below is a list of custom temperature computes that can be used like
that:
\begin{itemize}
\item {} 
\sphinxAtStartPar
{\hyperref[\detokenize{compute_temp_asphere::doc}]{\sphinxcrossref{\DUrole{doc}{compute temp/asphere command}}}}

\item {} 
\sphinxAtStartPar
{\hyperref[\detokenize{compute_temp_body::doc}]{\sphinxcrossref{\DUrole{doc}{compute temp/body command}}}}

\item {} 
\sphinxAtStartPar
{\hyperref[\detokenize{compute_temp_chunk::doc}]{\sphinxcrossref{\DUrole{doc}{compute temp/chunk command}}}}

\item {} 
\sphinxAtStartPar
{\hyperref[\detokenize{compute_temp_com::doc}]{\sphinxcrossref{\DUrole{doc}{compute temp/com command}}}}

\item {} 
\sphinxAtStartPar
{\hyperref[\detokenize{compute_temp_deform::doc}]{\sphinxcrossref{\DUrole{doc}{compute temp/deform command}}}}

\item {} 
\sphinxAtStartPar
{\hyperref[\detokenize{compute_temp_partial::doc}]{\sphinxcrossref{\DUrole{doc}{compute temp/partial command}}}}

\item {} 
\sphinxAtStartPar
{\hyperref[\detokenize{compute_temp_profile::doc}]{\sphinxcrossref{\DUrole{doc}{compute temp/profile command}}}}

\item {} 
\sphinxAtStartPar
{\hyperref[\detokenize{compute_temp_ramp::doc}]{\sphinxcrossref{\DUrole{doc}{compute temp/ramp command}}}}

\item {} 
\sphinxAtStartPar
{\hyperref[\detokenize{compute_temp_region::doc}]{\sphinxcrossref{\DUrole{doc}{compute temp/region command}}}}

\item {} 
\sphinxAtStartPar
{\hyperref[\detokenize{compute_temp_rotate::doc}]{\sphinxcrossref{\DUrole{doc}{compute temp/rotate command}}}}

\item {} 
\sphinxAtStartPar
{\hyperref[\detokenize{compute_temp_sphere::doc}]{\sphinxcrossref{\DUrole{doc}{compute temp/sphere command}}}}

\end{itemize}

\begin{sphinxadmonition}{note}{Note:}
\sphinxAtStartPar
Only the nvt fixes perform time integration, meaning they update
the velocities and positions of particles due to forces and velocities
respectively.  The other thermostat fixes only adjust velocities; they
do NOT perform time integration updates.  Thus they should be used in
conjunction with a constant NVE integration fix such as these:
\end{sphinxadmonition}
\begin{itemize}
\item {} 
\sphinxAtStartPar
{\hyperref[\detokenize{fix_nve::doc}]{\sphinxcrossref{\DUrole{doc}{fix nve}}}}

\item {} 
\sphinxAtStartPar
{\hyperref[\detokenize{fix_nve_sphere::doc}]{\sphinxcrossref{\DUrole{doc}{fix nve/sphere}}}}

\item {} 
\sphinxAtStartPar
{\hyperref[\detokenize{fix_nve_asphere::doc}]{\sphinxcrossref{\DUrole{doc}{fix nve/asphere}}}}

\end{itemize}

\sphinxAtStartPar
Thermodynamic output, which can be setup via the {\hyperref[\detokenize{thermo_style::doc}]{\sphinxcrossref{\DUrole{doc}{thermo\_style}}}} command, often includes temperature values.  As
explained on the page for the {\hyperref[\detokenize{thermo_style::doc}]{\sphinxcrossref{\DUrole{doc}{thermo\_style}}}}
command, the default temperature is setup by the thermo command
itself.  It is NOT the temperature associated with any thermostatting
fix you have defined or with any compute you have defined that
calculates a temperature.  The doc pages for the thermostatting fixes
explain the ID of the temperature compute they create.  Thus if you
want to view these temperatures, you need to specify them explicitly
via the {\hyperref[\detokenize{thermo_style::doc}]{\sphinxcrossref{\DUrole{doc}{thermo\_style custom}}}} command.  Or you can
use the {\hyperref[\detokenize{thermo_modify::doc}]{\sphinxcrossref{\DUrole{doc}{thermo\_modify}}}} command to re\sphinxhyphen{}define what
temperature compute is used for default thermodynamic output.


\bigskip\hrule\bigskip

\phantomsection\label{\detokenize{Howto_thermostat:daivis-thermostat}}
\sphinxAtStartPar
\sphinxstylestrong{(Daivis and Todd)} Daivis and Todd, Nonequilibrium Molecular Dynamics (book),
Cambridge University Press, \sphinxurl{https://doi.org/10.1017/9781139017848}, (2017).

\sphinxstepscope


\subsection{Barostats}
\label{\detokenize{Howto_barostat:barostats}}\label{\detokenize{Howto_barostat::doc}}
\sphinxAtStartPar
Barostatting means controlling the pressure in an MD simulation.
{\hyperref[\detokenize{Howto_thermostat::doc}]{\sphinxcrossref{\DUrole{doc}{Thermostatting}}}} means controlling the
temperature of the particles.  Since the pressure includes a kinetic
component due to particle velocities, both these operations require
calculation of the temperature.  Typically a target temperature (T)
and/or pressure (P) is specified by the user, and the thermostat or
barostat attempts to equilibrate the system to the requested T and/or
P.

\sphinxAtStartPar
Barostatting in LAMMPS is performed by {\hyperref[\detokenize{fix::doc}]{\sphinxcrossref{\DUrole{doc}{fixes}}}}.  Two
barostatting methods are currently available: Nose\sphinxhyphen{}Hoover (npt and
nph) and Berendsen:
\begin{itemize}
\item {} 
\sphinxAtStartPar
{\hyperref[\detokenize{fix_nh::doc}]{\sphinxcrossref{\DUrole{doc}{fix npt}}}}

\item {} 
\sphinxAtStartPar
{\hyperref[\detokenize{fix_npt_sphere::doc}]{\sphinxcrossref{\DUrole{doc}{fix npt/sphere}}}}

\item {} 
\sphinxAtStartPar
{\hyperref[\detokenize{fix_npt_asphere::doc}]{\sphinxcrossref{\DUrole{doc}{fix npt/asphere}}}}

\item {} 
\sphinxAtStartPar
{\hyperref[\detokenize{fix_nh::doc}]{\sphinxcrossref{\DUrole{doc}{fix nph}}}}

\item {} 
\sphinxAtStartPar
{\hyperref[\detokenize{fix_press_berendsen::doc}]{\sphinxcrossref{\DUrole{doc}{fix press/berendsen}}}}

\end{itemize}

\sphinxAtStartPar
The {\hyperref[\detokenize{fix_nh::doc}]{\sphinxcrossref{\DUrole{doc}{fix npt}}}} commands include a Nose\sphinxhyphen{}Hoover thermostat
and barostat.  {\hyperref[\detokenize{fix_nh::doc}]{\sphinxcrossref{\DUrole{doc}{Fix nph}}}} is just a Nose/Hoover barostat;
it does no thermostatting.  Both {\hyperref[\detokenize{fix_nh::doc}]{\sphinxcrossref{\DUrole{doc}{fix nph}}}} and {\hyperref[\detokenize{fix_press_berendsen::doc}]{\sphinxcrossref{\DUrole{doc}{fix press/berendsen}}}} can be used in conjunction
with any of the thermostatting fixes.

\sphinxAtStartPar
As with the {\hyperref[\detokenize{Howto_thermostat::doc}]{\sphinxcrossref{\DUrole{doc}{thermostats}}}}, {\hyperref[\detokenize{fix_nh::doc}]{\sphinxcrossref{\DUrole{doc}{fix npt}}}}
and {\hyperref[\detokenize{fix_nh::doc}]{\sphinxcrossref{\DUrole{doc}{fix nph}}}} only use translational motion of the
particles in computing T and P and performing thermo/barostatting.
{\hyperref[\detokenize{fix_npt_sphere::doc}]{\sphinxcrossref{\DUrole{doc}{Fix npt/sphere}}}} and {\hyperref[\detokenize{fix_npt_asphere::doc}]{\sphinxcrossref{\DUrole{doc}{fix npt/asphere}}}} thermo/barostat using not only
translation velocities but also rotational velocities for spherical
and aspherical particles.

\sphinxAtStartPar
All of the barostatting fixes use the {\hyperref[\detokenize{compute_pressure::doc}]{\sphinxcrossref{\DUrole{doc}{compute pressure}}}} compute to calculate a current
pressure.  By default, this compute is created with a simple {\hyperref[\detokenize{compute_temp::doc}]{\sphinxcrossref{\DUrole{doc}{compute temp}}}} (see the last argument of the {\hyperref[\detokenize{compute_pressure::doc}]{\sphinxcrossref{\DUrole{doc}{compute pressure}}}} command), which is used to calculated
the kinetic component of the pressure.  The barostatting fixes can
also use temperature computes that remove bias for the purpose of
computing the kinetic component which contributes to the current
pressure.  See the doc pages for the individual fixes and for the
{\hyperref[\detokenize{fix_modify::doc}]{\sphinxcrossref{\DUrole{doc}{fix\_modify}}}} command for instructions on how to assign
a temperature or pressure compute to a barostatting fix.

\begin{sphinxadmonition}{note}{Note:}
\sphinxAtStartPar
As with the thermostats, the Nose/Hoover methods ({\hyperref[\detokenize{fix_nh::doc}]{\sphinxcrossref{\DUrole{doc}{fix npt}}}} and {\hyperref[\detokenize{fix_nh::doc}]{\sphinxcrossref{\DUrole{doc}{fix nph}}}}) perform time integration.
{\hyperref[\detokenize{fix_press_berendsen::doc}]{\sphinxcrossref{\DUrole{doc}{Fix press/berendsen}}}} does NOT, so it should
be used with one of the constant NVE fixes or with one of the NVT
fixes.
\end{sphinxadmonition}

\sphinxAtStartPar
Thermodynamic output, which can be setup via the
{\hyperref[\detokenize{thermo_style::doc}]{\sphinxcrossref{\DUrole{doc}{thermo\_style}}}} command, often includes pressure
values.  As explained on the page for the
{\hyperref[\detokenize{thermo_style::doc}]{\sphinxcrossref{\DUrole{doc}{thermo\_style}}}} command, the default pressure is
setup by the thermo command itself.  It is NOT the pressure associated
with any barostatting fix you have defined or with any compute you
have defined that calculates a pressure.  The doc pages for the
barostatting fixes explain the ID of the pressure compute they create.
Thus if you want to view these pressures, you need to specify them
explicitly via the {\hyperref[\detokenize{thermo_style::doc}]{\sphinxcrossref{\DUrole{doc}{thermo\_style custom}}}} command.
Or you can use the {\hyperref[\detokenize{thermo_modify::doc}]{\sphinxcrossref{\DUrole{doc}{thermo\_modify}}}} command to
re\sphinxhyphen{}define what pressure compute is used for default thermodynamic
output.

\sphinxstepscope


\subsection{Walls}
\label{\detokenize{Howto_walls:walls}}\label{\detokenize{Howto_walls::doc}}
\sphinxAtStartPar
Walls in an MD simulation are typically used to bound particle motion,
i.e. to serve as a boundary condition.

\sphinxAtStartPar
Walls in LAMMPS can be of rough (made of particles) or idealized
surfaces.  Ideal walls can be smooth, generating forces only in the
normal direction, or frictional, generating forces also in the
tangential direction.

\sphinxAtStartPar
Rough walls, built of particles, can be created in various ways.  The
particles themselves can be generated like any other particle, via the
{\hyperref[\detokenize{lattice::doc}]{\sphinxcrossref{\DUrole{doc}{lattice}}}} and {\hyperref[\detokenize{create_atoms::doc}]{\sphinxcrossref{\DUrole{doc}{create\_atoms}}}} commands,
or read in via the {\hyperref[\detokenize{read_data::doc}]{\sphinxcrossref{\DUrole{doc}{read\_data}}}} command.

\sphinxAtStartPar
Their motion can be constrained by many different commands, so that
they do not move at all, move together as a group at constant velocity
or in response to a net force acting on them, move in a prescribed
fashion (e.g. rotate around a point), etc.  Note that if a time
integration fix like {\hyperref[\detokenize{fix_nve::doc}]{\sphinxcrossref{\DUrole{doc}{fix nve}}}} or {\hyperref[\detokenize{fix_nh::doc}]{\sphinxcrossref{\DUrole{doc}{fix nvt}}}}
is not used with the group that contains wall particles, their
positions and velocities will not be updated.
\begin{itemize}
\item {} 
\sphinxAtStartPar
{\hyperref[\detokenize{fix_aveforce::doc}]{\sphinxcrossref{\DUrole{doc}{fix aveforce}}}} \sphinxhyphen{} set force on particles to average value, so they move together

\item {} 
\sphinxAtStartPar
{\hyperref[\detokenize{fix_setforce::doc}]{\sphinxcrossref{\DUrole{doc}{fix setforce}}}} \sphinxhyphen{} set force on particles to a value, e.g. 0.0

\item {} 
\sphinxAtStartPar
{\hyperref[\detokenize{fix_freeze::doc}]{\sphinxcrossref{\DUrole{doc}{fix freeze}}}} \sphinxhyphen{} freeze particles for use as granular walls

\item {} 
\sphinxAtStartPar
{\hyperref[\detokenize{fix_nve_noforce::doc}]{\sphinxcrossref{\DUrole{doc}{fix nve/noforce}}}} \sphinxhyphen{} advect particles by their velocity, but without force

\item {} 
\sphinxAtStartPar
{\hyperref[\detokenize{fix_move::doc}]{\sphinxcrossref{\DUrole{doc}{fix move}}}} \sphinxhyphen{} prescribe motion of particles by a linear velocity, oscillation, rotation, variable

\end{itemize}

\sphinxAtStartPar
The {\hyperref[\detokenize{fix_move::doc}]{\sphinxcrossref{\DUrole{doc}{fix move}}}} command offers the most generality, since
the motion of individual particles can be specified with
{\hyperref[\detokenize{variable::doc}]{\sphinxcrossref{\DUrole{doc}{variable}}}} formula which depends on time and/or the
particle position.

\sphinxAtStartPar
For rough walls, it may be useful to turn off pairwise interactions
between wall particles via the {\hyperref[\detokenize{neigh_modify::doc}]{\sphinxcrossref{\DUrole{doc}{neigh\_modify exclude}}}} command.

\sphinxAtStartPar
Rough walls can also be created by specifying frozen particles that do
not move and do not interact with mobile particles, and then tethering
other particles to the fixed particles, via a {\hyperref[\detokenize{bond_style::doc}]{\sphinxcrossref{\DUrole{doc}{bond}}}}.
The bonded particles do interact with other mobile particles.

\sphinxAtStartPar
Idealized walls can be specified via several fix commands.  {\hyperref[\detokenize{fix_wall_gran::doc}]{\sphinxcrossref{\DUrole{doc}{Fix wall/gran}}}} creates frictional walls for use with
granular particles; all the other commands create smooth walls.
\begin{itemize}
\item {} 
\sphinxAtStartPar
{\hyperref[\detokenize{fix_wall_reflect::doc}]{\sphinxcrossref{\DUrole{doc}{fix wall/reflect}}}} \sphinxhyphen{} reflective flat walls

\item {} 
\sphinxAtStartPar
{\hyperref[\detokenize{fix_wall::doc}]{\sphinxcrossref{\DUrole{doc}{fix wall/lj93}}}} \sphinxhyphen{} flat walls, with Lennard\sphinxhyphen{}Jones 9/3 potential

\item {} 
\sphinxAtStartPar
{\hyperref[\detokenize{fix_wall::doc}]{\sphinxcrossref{\DUrole{doc}{fix wall/lj126}}}} \sphinxhyphen{} flat walls, with Lennard\sphinxhyphen{}Jones 12/6 potential

\item {} 
\sphinxAtStartPar
{\hyperref[\detokenize{fix_wall::doc}]{\sphinxcrossref{\DUrole{doc}{fix wall/colloid}}}} \sphinxhyphen{} flat walls, with {\hyperref[\detokenize{pair_colloid::doc}]{\sphinxcrossref{\DUrole{doc}{pair\_style colloid}}}} potential

\item {} 
\sphinxAtStartPar
{\hyperref[\detokenize{fix_wall::doc}]{\sphinxcrossref{\DUrole{doc}{fix wall/harmonic}}}} \sphinxhyphen{} flat walls, with repulsive harmonic spring potential

\item {} 
\sphinxAtStartPar
{\hyperref[\detokenize{fix_wall::doc}]{\sphinxcrossref{\DUrole{doc}{fix wall/morse}}}} \sphinxhyphen{} flat walls, with Morse potential

\item {} 
\sphinxAtStartPar
{\hyperref[\detokenize{fix_wall_region::doc}]{\sphinxcrossref{\DUrole{doc}{fix wall/region}}}} \sphinxhyphen{} use region surface as wall

\item {} 
\sphinxAtStartPar
{\hyperref[\detokenize{fix_wall_gran::doc}]{\sphinxcrossref{\DUrole{doc}{fix wall/gran}}}} \sphinxhyphen{} flat or curved walls with {\hyperref[\detokenize{pair_gran::doc}]{\sphinxcrossref{\DUrole{doc}{pair\_style granular}}}} potential

\end{itemize}

\sphinxAtStartPar
The \sphinxstyleemphasis{lj93}, \sphinxstyleemphasis{lj126}, \sphinxstyleemphasis{colloid}, \sphinxstyleemphasis{harmonic}, and \sphinxstyleemphasis{morse} styles all
allow the flat walls to move with a constant velocity, or oscillate in
time.  The {\hyperref[\detokenize{fix_wall_region::doc}]{\sphinxcrossref{\DUrole{doc}{fix wall/region}}}} command offers the
most generality, since the region surface is treated as a wall, and
the geometry of the region can be a simple primitive volume (e.g. a
sphere, or cube, or plane), or a complex volume made from the union
and intersection of primitive volumes.  {\hyperref[\detokenize{region::doc}]{\sphinxcrossref{\DUrole{doc}{Regions}}}} can also
specify a volume “interior” or “exterior” to the specified primitive
shape or \sphinxstyleemphasis{union} or \sphinxstyleemphasis{intersection}.  {\hyperref[\detokenize{region::doc}]{\sphinxcrossref{\DUrole{doc}{Regions}}}} can also be
“dynamic” meaning they move with constant velocity, oscillate, or
rotate.

\sphinxAtStartPar
The only frictional idealized walls currently in LAMMPS are flat or
curved surfaces specified by the {\hyperref[\detokenize{fix_wall_gran::doc}]{\sphinxcrossref{\DUrole{doc}{fix wall/gran}}}}
command.  At some point we plan to allow region surfaces to be used as
frictional walls, as well as triangulated surfaces.

\sphinxstepscope


\subsection{NEMD simulations}
\label{\detokenize{Howto_nemd:nemd-simulations}}\label{\detokenize{Howto_nemd::doc}}
\sphinxAtStartPar
Non\sphinxhyphen{}equilibrium molecular dynamics or NEMD simulations are typically
used to measure a fluid’s rheological properties such as viscosity.
In LAMMPS, such simulations can be performed by first setting up a
non\sphinxhyphen{}orthogonal simulation box (see the preceding Howto section).

\sphinxAtStartPar
A shear strain can be applied to the simulation box at a desired
strain rate by using the {\hyperref[\detokenize{fix_deform::doc}]{\sphinxcrossref{\DUrole{doc}{fix deform}}}} command.  The
{\hyperref[\detokenize{fix_nvt_sllod::doc}]{\sphinxcrossref{\DUrole{doc}{fix nvt/sllod}}}} command can be used to thermostat
the sheared fluid and integrate the SLLOD equations of motion for the
system.  Fix nvt/sllod uses {\hyperref[\detokenize{compute_temp_deform::doc}]{\sphinxcrossref{\DUrole{doc}{compute temp/deform}}}} to compute a thermal temperature
by subtracting out the streaming velocity of the shearing atoms.  The
velocity profile or other properties of the fluid can be monitored via
the {\hyperref[\detokenize{fix_ave_chunk::doc}]{\sphinxcrossref{\DUrole{doc}{fix ave/chunk}}}} command.

\begin{sphinxadmonition}{note}{Note:}
\sphinxAtStartPar
A recent (2017) book by {\hyperref[\detokenize{Howto_nemd:daivis-nemd}]{\sphinxcrossref{\DUrole{std,std-ref}{(Daivis and Todd)}}}}
discusses use of the SLLOD method and non\sphinxhyphen{}equilibrium MD (NEMD)
thermostatting generally, for both simple and complex fluids,
e.g. molecular systems.  The latter can be tricky to do correctly.
\end{sphinxadmonition}

\sphinxAtStartPar
As discussed in the previous section on non\sphinxhyphen{}orthogonal simulation
boxes, the amount of tilt or skew that can be applied is limited by
LAMMPS for computational efficiency to be 1/2 of the parallel box
length.  However, {\hyperref[\detokenize{fix_deform::doc}]{\sphinxcrossref{\DUrole{doc}{fix deform}}}} can continuously strain
a box by an arbitrary amount.  As discussed in the {\hyperref[\detokenize{fix_deform::doc}]{\sphinxcrossref{\DUrole{doc}{fix deform}}}} command, when the tilt value reaches a limit,
the box is flipped to the opposite limit which is an equivalent tiling
of periodic space.  The strain rate can then continue to change as
before.  In a long NEMD simulation these box re\sphinxhyphen{}shaping events may
occur many times.

\sphinxAtStartPar
In a NEMD simulation, the “remap” option of {\hyperref[\detokenize{fix_deform::doc}]{\sphinxcrossref{\DUrole{doc}{fix deform}}}} should be set to “remap v”, since that is what
{\hyperref[\detokenize{fix_nvt_sllod::doc}]{\sphinxcrossref{\DUrole{doc}{fix nvt/sllod}}}} assumes to generate a velocity
profile consistent with the applied shear strain rate.

\sphinxAtStartPar
An alternative method for calculating viscosities is provided via the
{\hyperref[\detokenize{fix_viscosity::doc}]{\sphinxcrossref{\DUrole{doc}{fix viscosity}}}} command.

\sphinxAtStartPar
NEMD simulations can also be used to measure transport properties of a fluid
through a pore or channel. Simulations of steady\sphinxhyphen{}state flow can be performed
using the {\hyperref[\detokenize{fix_flow_gauss::doc}]{\sphinxcrossref{\DUrole{doc}{fix flow/gauss}}}} command.


\bigskip\hrule\bigskip

\phantomsection\label{\detokenize{Howto_nemd:daivis-nemd}}
\sphinxAtStartPar
\sphinxstylestrong{(Daivis and Todd)} Daivis and Todd, Nonequilibrium Molecular Dynamics (book),
Cambridge University Press, \sphinxurl{https://doi.org/10.1017/9781139017848}, (2017).

\sphinxstepscope


\subsection{Long\sphinxhyphen{}range dispersion settings}
\label{\detokenize{Howto_dispersion:long-range-dispersion-settings}}\label{\detokenize{Howto_dispersion::doc}}
\sphinxAtStartPar
The PPPM method computes interactions by splitting the pair potential
into two parts, one of which is computed in a normal pairwise fashion,
the so\sphinxhyphen{}called real\sphinxhyphen{}space part, and one of which is computed using the
Fourier transform, the so called reciprocal\sphinxhyphen{}space or kspace part.  For
both parts, the potential is not computed exactly but is approximated.
Thus, there is an error in both parts of the computation, the
real\sphinxhyphen{}space and the kspace error. The just mentioned facts are true
both for the PPPM for Coulomb as well as dispersion interactions. The
deciding difference \sphinxhyphen{} and also the reason why the parameters for
pppm/disp have to be selected with more care \sphinxhyphen{} is the impact of the
errors on the results: The kspace error of the PPPM for Coulomb and
dispersion interaction and the real\sphinxhyphen{}space error of the PPPM for
Coulomb interaction have the character of noise. In contrast, the
real\sphinxhyphen{}space error of the PPPM for dispersion has a clear physical
interpretation: the underprediction of cohesion. As a consequence, the
real\sphinxhyphen{}space error has a much stronger effect than the kspace error on
simulation results for pppm/disp.  Parameters must thus be chosen in a
way that this error is much smaller than the kspace error.

\sphinxAtStartPar
When using pppm/disp and not making any specifications on the PPPM
parameters via the kspace modify command, parameters will be tuned
such that the real\sphinxhyphen{}space error and the kspace error are equal.  This
will result in simulations that are either inaccurate or slow, both of
which is not desirable. For selecting parameters for the pppm/disp
that provide fast and accurate simulations, there are two approaches,
which both have their up\sphinxhyphen{} and downsides.

\sphinxAtStartPar
The first approach is to set desired real\sphinxhyphen{}space an kspace accuracies
via the \sphinxstyleemphasis{kspace\_modify force/disp/real} and \sphinxstyleemphasis{kspace\_modify
force/disp/kspace} commands. Note that the accuracies have to be
specified in force units and are thus dependent on the chosen unit
settings. For real units, 0.0001 and 0.002 seem to provide reasonable
accurate and efficient computations for the real\sphinxhyphen{}space and kspace
accuracies.  0.002 and 0.05 work well for most systems using lj
units. PPPM parameters will be generated based on the desired
accuracies. The upside of this approach is that it usually provides a
good set of parameters and will work for both the \sphinxstyleemphasis{kspace\_modify diff
ad} and \sphinxstyleemphasis{kspace\_modify diff ik} options.  The downside of the method
is that setting the PPPM parameters will take some time during the
initialization of the simulation.

\sphinxAtStartPar
The second approach is to set the parameters for the pppm/disp
explicitly using the \sphinxstyleemphasis{kspace\_modify mesh/disp}, \sphinxstyleemphasis{kspace\_modify
order/disp}, and \sphinxstyleemphasis{kspace\_modify gewald/disp} commands. This approach
requires a more experienced user who understands well the impact of
the choice of parameters on the simulation accuracy and
performance. This approach provides a fast initialization of the
simulation. However, it is sensitive to errors: A combination of
parameters that will perform well for one system might result in
far\sphinxhyphen{}from\sphinxhyphen{}optimal conditions for other simulations. For example,
parameters that provide accurate and fast computations for
all\sphinxhyphen{}atomistic force fields can provide insufficient accuracy or
united\sphinxhyphen{}atomistic force fields (which is related to that the latter
typically have larger dispersion coefficients).

\sphinxAtStartPar
To avoid inaccurate or inefficient simulations, the pppm/disp stops
simulations with an error message if no action is taken to control the
PPPM parameters. If the automatic parameter generation is desired and
real\sphinxhyphen{}space and kspace accuracies are desired to be equal, this error
message can be suppressed using the \sphinxstyleemphasis{kspace\_modify disp/auto yes}
command.

\sphinxAtStartPar
A reasonable approach that combines the upsides of both methods is to
make the first run using the \sphinxstyleemphasis{kspace\_modify force/disp/real} and
\sphinxstyleemphasis{kspace\_modify force/disp/kspace} commands, write down the PPPM
parameters from the output, and specify these parameters using the
second approach in subsequent runs (which have the same composition,
force field, and approximately the same volume).

\sphinxAtStartPar
Concerning the performance of the pppm/disp there are two more things
to consider. The first is that when using the pppm/disp, the cutoff
parameter does no longer affect the accuracy of the simulation
(subject to that gewald/disp is adjusted when changing the cutoff).
The performance can thus be increased by examining different values
for the cutoff parameter. A lower bound for the cutoff is only set by
the truncation error of the repulsive term of pair potentials.

\sphinxAtStartPar
The second is that the mixing rule of the pair style has an impact on
the computation time when using the pppm/disp. Fastest computations
are achieved when using the geometric mixing rule. Using the
arithmetic mixing rule substantially increases the computational cost.
The computational overhead can be reduced using the \sphinxstyleemphasis{kspace\_modify
mix/disp geom} and \sphinxstyleemphasis{kspace\_modify splittol} commands. The first
command simply enforces geometric mixing of the dispersion
coefficients in kspace computations.  This introduces some error in
the computations but will also significantly speed\sphinxhyphen{}up the
simulations. The second keyword sets the accuracy with which the
dispersion coefficients are approximated using a matrix factorization
approach.  This may result in better accuracy then using the first
command, but will usually also not provide an equally good increase of
efficiency.

\sphinxAtStartPar
Finally, pppm/disp can also be used when no mixing rules apply.
This can be achieved using the \sphinxstyleemphasis{kspace\_modify mix/disp none} command.
Note that the code does not check automatically whether any mixing
rule is fulfilled. If mixing rules do not apply, the user will have
to specify this command explicitly.


\section{Analysis howto}
\label{\detokenize{Howto:analysis-howto}}
\sphinxstepscope


\subsection{Output from LAMMPS (thermo, dumps, computes, fixes, variables)}
\label{\detokenize{Howto_output:output-from-lammps-thermo-dumps-computes-fixes-variables}}\label{\detokenize{Howto_output::doc}}
\sphinxAtStartPar
There are four basic forms of LAMMPS output:
\begin{itemize}
\item {} 
\sphinxAtStartPar
{\hyperref[\detokenize{thermo_style::doc}]{\sphinxcrossref{\DUrole{doc}{Thermodynamic output}}}}, which is a list of
quantities printed every few timesteps to the screen and logfile.

\item {} 
\sphinxAtStartPar
{\hyperref[\detokenize{dump::doc}]{\sphinxcrossref{\DUrole{doc}{Dump files}}}}, which contain snapshots of atoms and various
per\sphinxhyphen{}atom values and are written at a specified frequency.

\item {} 
\sphinxAtStartPar
Certain fixes can output user\sphinxhyphen{}specified quantities to files:
{\hyperref[\detokenize{fix_ave_time::doc}]{\sphinxcrossref{\DUrole{doc}{fix ave/time}}}} for time averaging,
{\hyperref[\detokenize{fix_ave_chunk::doc}]{\sphinxcrossref{\DUrole{doc}{fix ave/chunk}}}} for spatial or other averaging, and
{\hyperref[\detokenize{fix_print::doc}]{\sphinxcrossref{\DUrole{doc}{fix print}}}} for single\sphinxhyphen{}line output of
{\hyperref[\detokenize{variable::doc}]{\sphinxcrossref{\DUrole{doc}{variables}}}}.  Fix print can also output to the screen.

\item {} 
\sphinxAtStartPar
{\hyperref[\detokenize{restart::doc}]{\sphinxcrossref{\DUrole{doc}{Restart files}}}}.

\end{itemize}

\sphinxAtStartPar
A simulation prints one set of thermodynamic output and (optionally)
restart files.  It can generate any number of dump files and fix
output files, depending on what {\hyperref[\detokenize{dump::doc}]{\sphinxcrossref{\DUrole{doc}{dump}}}} and {\hyperref[\detokenize{fix::doc}]{\sphinxcrossref{\DUrole{doc}{fix}}}}
commands you specify.

\sphinxAtStartPar
As discussed below, LAMMPS gives you a variety of ways to determine
what quantities are calculated and printed when the thermodynamics,
dump, or fix commands listed above perform output.  Throughout this
discussion, note that users can also {\hyperref[\detokenize{Modify::doc}]{\sphinxcrossref{\DUrole{doc}{add their own computes and
fixes to LAMMPS}}}} which can generate values that can then be
output with these commands.

\sphinxAtStartPar
The following subsections discuss different LAMMPS commands related
to output and the kind of data they operate on and produce:
\begin{itemize}
\item {} 
\sphinxAtStartPar
{\hyperref[\detokenize{Howto_output:global}]{\sphinxcrossref{\DUrole{std,std-ref}{Global/per\sphinxhyphen{}atom/local/per\sphinxhyphen{}grid data}}}}

\item {} 
\sphinxAtStartPar
{\hyperref[\detokenize{Howto_output:scalar}]{\sphinxcrossref{\DUrole{std,std-ref}{Scalar/vector/array data}}}}

\item {} 
\sphinxAtStartPar
{\hyperref[\detokenize{Howto_output:disambiguation}]{\sphinxcrossref{\DUrole{std,std-ref}{Disambiguation}}}}

\item {} 
\sphinxAtStartPar
{\hyperref[\detokenize{Howto_output:thermo}]{\sphinxcrossref{\DUrole{std,std-ref}{Thermodynamic output}}}}

\item {} 
\sphinxAtStartPar
{\hyperref[\detokenize{Howto_output:dump}]{\sphinxcrossref{\DUrole{std,std-ref}{Dump file output}}}}

\item {} 
\sphinxAtStartPar
{\hyperref[\detokenize{Howto_output:fixoutput}]{\sphinxcrossref{\DUrole{std,std-ref}{Fixes that write output files}}}}

\item {} 
\sphinxAtStartPar
{\hyperref[\detokenize{Howto_output:computeoutput}]{\sphinxcrossref{\DUrole{std,std-ref}{Computes that process output quantities}}}}

\item {} 
\sphinxAtStartPar
{\hyperref[\detokenize{Howto_output:fixprocoutput}]{\sphinxcrossref{\DUrole{std,std-ref}{Fixes that process output quantities}}}}

\item {} 
\sphinxAtStartPar
{\hyperref[\detokenize{Howto_output:compute}]{\sphinxcrossref{\DUrole{std,std-ref}{Computes that generate values to output}}}}

\item {} 
\sphinxAtStartPar
{\hyperref[\detokenize{Howto_output:fix}]{\sphinxcrossref{\DUrole{std,std-ref}{Fixes that generate values to output}}}}

\item {} 
\sphinxAtStartPar
{\hyperref[\detokenize{Howto_output:variable}]{\sphinxcrossref{\DUrole{std,std-ref}{Variables that generate values to output}}}}

\item {} 
\sphinxAtStartPar
{\hyperref[\detokenize{Howto_output:table}]{\sphinxcrossref{\DUrole{std,std-ref}{Summary table of output options and data flow between commands}}}}

\end{itemize}


\subsubsection{Global/per\sphinxhyphen{}atom/local/per\sphinxhyphen{}grid data}
\label{\detokenize{Howto_output:global-per-atom-local-per-grid-data}}\label{\detokenize{Howto_output:global}}
\sphinxAtStartPar
Various output\sphinxhyphen{}related commands work with four different “styles” of
data: global, per\sphinxhyphen{}atom, local, and per\sphinxhyphen{}grid.  A global datum is one or
more system\sphinxhyphen{}wide values, e.g. the temperature of the system.  A
per\sphinxhyphen{}atom datum is one or more values per atom, e.g. the kinetic energy
of each atom.  Local datums are calculated by each processor based on
the atoms it owns, and there may be zero or more per atom, e.g. a list
of bond distances.

\sphinxAtStartPar
A per\sphinxhyphen{}grid datum is one or more values per grid cell, for a grid which
overlays the simulation domain.  Similar to atoms and per\sphinxhyphen{}atom data,
the grid cells and the data they store are distributed across
processors; each processor owns the grid cells whose center points
fall within its subdomain.


\subsubsection{Scalar/vector/array data}
\label{\detokenize{Howto_output:scalar-vector-array-data}}\label{\detokenize{Howto_output:scalar}}
\sphinxAtStartPar
Global, per\sphinxhyphen{}atom, local, and per\sphinxhyphen{}grid datums can come in three
“kinds”: a single scalar value, a vector of values, or a 2d array of
values.  More specifically these are the valid kinds for each style:
\begin{itemize}
\item {} 
\sphinxAtStartPar
global scalar

\item {} 
\sphinxAtStartPar
global vector

\item {} 
\sphinxAtStartPar
global array

\item {} 
\sphinxAtStartPar
per\sphinxhyphen{}atom vector

\item {} 
\sphinxAtStartPar
per\sphinxhyphen{}atom array

\item {} 
\sphinxAtStartPar
local vector

\item {} 
\sphinxAtStartPar
local array

\item {} 
\sphinxAtStartPar
per\sphinxhyphen{}grid vector

\item {} 
\sphinxAtStartPar
per\sphinxhyphen{}grid array

\end{itemize}

\sphinxAtStartPar
A per\sphinxhyphen{}atom vector means a single value per atom; the “vector” is the
length of the number of atoms.  A per\sphinxhyphen{}atom array means multiple values
per atom.  Similarly a local vector or array means one or multiple
values per entity (e.g. per bond in the system).  And a per\sphinxhyphen{}grid
vector or array means one or multiple values per grid cell.

\sphinxAtStartPar
The doc page for a compute or fix or variable that generates data will
specify both the styles and kinds of data it produces, e.g. a per\sphinxhyphen{}atom
vector.  Note that a compute or fix may generate multiple styles and
kinds of output.  However, for per\sphinxhyphen{}atom data only a vector or array is
output, never both.  Likewise for per\sphinxhyphen{}local and per\sphinxhyphen{}grid data.  An
example of a fix which generates multiple styles and kinds of data is
the {\hyperref[\detokenize{fix_mdi_qm::doc}]{\sphinxcrossref{\DUrole{doc}{fix mdi/qm}}}} command.  It outputs a global
scalar, global vector, and per\sphinxhyphen{}atom array for the quantum mechanical
energy and virial of the system and forces on each atom.

\sphinxAtStartPar
By contrast, different variable styles generate only a single kind of
data: a global scalar for an equal\sphinxhyphen{}style variable, global vector for a
vector\sphinxhyphen{}style variable, and a per\sphinxhyphen{}atom vector for an atom\sphinxhyphen{}style
variable.

\sphinxAtStartPar
When data is accessed by another command, as in many of the output
commands discussed below, it can be referenced via the following
bracket notation, where ID in this case is the ID of a compute.  The
leading “c\_” would be replaced by “f\_” for a fix, or “v\_” for a
variable (and ID would be the name of the variable):


\begin{savenotes}\sphinxattablestart
\sphinxthistablewithglobalstyle
\centering
\begin{tabulary}{\linewidth}[t]{TT}
\sphinxtoprule
\sphinxtableatstartofbodyhook
\sphinxAtStartPar
c\_ID
&
\sphinxAtStartPar
entire scalar, vector, or array
\\
\sphinxhline
\sphinxAtStartPar
c\_ID{[}I{]}
&
\sphinxAtStartPar
one element of vector, one column of array
\\
\sphinxhline
\sphinxAtStartPar
c\_ID{[}I{]}{[}J{]}
&
\sphinxAtStartPar
one element of array
\\
\sphinxbottomrule
\end{tabulary}
\sphinxtableafterendhook\par
\sphinxattableend\end{savenotes}

\sphinxAtStartPar
Note that using one bracket reduces the dimension of the data once
(vector \sphinxhyphen{}\textgreater{} scalar, array \sphinxhyphen{}\textgreater{} vector).  Using two brackets reduces the
dimension twice (array \sphinxhyphen{}\textgreater{} scalar).  Thus a command that uses scalar
values as input can also conceptually operate on an element of a
vector or array.

\sphinxAtStartPar
Per\sphinxhyphen{}grid vectors or arrays are accessed similarly, except that the ID
for the compute or fix includes a grid name and a data name.  This is
because a fix or compute can create multiple grids (of different
sizes) and multiple sets of data (for each grid).  The fix or compute
defines names for each grid and for each data set, so that all of them
can be accessed by other commands.  See the {\hyperref[\detokenize{Howto_grid::doc}]{\sphinxcrossref{\DUrole{doc}{Howto grid}}}} doc page for more details.


\subsubsection{Disambiguation}
\label{\detokenize{Howto_output:disambiguation}}\label{\detokenize{Howto_output:id1}}
\sphinxAtStartPar
When a compute or fix produces data in multiple styles, e.g. global
and per\sphinxhyphen{}atom, a reference to the data can sometimes be ambiguous.
Usually the context in which the input script references the data
determines which style is meant.

\sphinxAtStartPar
For example, if a compute outputs a global vector and a per\sphinxhyphen{}atom
array, an element of the global vector will be accessed by using
\sphinxcode{\sphinxupquote{c\_ID{[}I{]}}} in {\hyperref[\detokenize{thermo_style::doc}]{\sphinxcrossref{\DUrole{doc}{thermodynamic output}}}}, while a
column of the per\sphinxhyphen{}atom array will be accessed by using \sphinxcode{\sphinxupquote{c\_ID{[}I{]}}} in
a {\hyperref[\detokenize{dump::doc}]{\sphinxcrossref{\DUrole{doc}{dump custom}}}} command.

\sphinxAtStartPar
However, if a {\hyperref[\detokenize{variable::doc}]{\sphinxcrossref{\DUrole{doc}{atom\sphinxhyphen{}style variable}}}} references
\sphinxcode{\sphinxupquote{c\_ID{[}I{]}}}, then it could be intended to refer to a single element of
the global vector or a column of the per\sphinxhyphen{}atom array.  The doc page for
any command that has a potential ambiguity (variables are the most
common) will explain how to resolve the ambiguity.

\sphinxAtStartPar
In this case, an atom\sphinxhyphen{}style variables references per\sphinxhyphen{}atom data if it
exists.  If access to an element of a global vector is needed (as in
this example), an equal\sphinxhyphen{}style variable which references the value can
be defined and used in the atom\sphinxhyphen{}style variable formula instead.

\sphinxAtStartPar
Similarly, {\hyperref[\detokenize{thermo_style::doc}]{\sphinxcrossref{\DUrole{doc}{thermodynamic output}}}} can only
reference global data from a compute or fix.  But you can indirectly
access per\sphinxhyphen{}atom data as follows.  The reference \sphinxcode{\sphinxupquote{c\_ID{[}245{]}{[}2{]}}} for
the ID of a {\hyperref[\detokenize{compute_displace_atom::doc}]{\sphinxcrossref{\DUrole{doc}{compute displace/atom}}}}
command, refers to the y\sphinxhyphen{}component of displacement for the atom with
ID 245.  While you cannot use that reference directly in the
{\hyperref[\detokenize{thermo_style::doc}]{\sphinxcrossref{\DUrole{doc}{thermo\_style}}}} command, you can use it an
equal\sphinxhyphen{}style variable formula, and then reference the variable in
thermodynamic output.


\subsubsection{Thermodynamic output}
\label{\detokenize{Howto_output:thermodynamic-output}}\label{\detokenize{Howto_output:thermo}}
\sphinxAtStartPar
The frequency and format of thermodynamic output is set by the
{\hyperref[\detokenize{thermo::doc}]{\sphinxcrossref{\DUrole{doc}{thermo}}}}, {\hyperref[\detokenize{thermo_style::doc}]{\sphinxcrossref{\DUrole{doc}{thermo\_style}}}}, and
{\hyperref[\detokenize{thermo_modify::doc}]{\sphinxcrossref{\DUrole{doc}{thermo\_modify}}}} commands.  The {\hyperref[\detokenize{thermo_style::doc}]{\sphinxcrossref{\DUrole{doc}{thermo\_style}}}} command also specifies what values are calculated and
written out.  Pre\sphinxhyphen{}defined keywords can be specified (e.g. press, etotal,
etc).  Three additional kinds of keywords can also be specified (c\_ID,
f\_ID, v\_name), where a {\hyperref[\detokenize{compute::doc}]{\sphinxcrossref{\DUrole{doc}{compute}}}} or {\hyperref[\detokenize{fix::doc}]{\sphinxcrossref{\DUrole{doc}{fix}}}} or
{\hyperref[\detokenize{variable::doc}]{\sphinxcrossref{\DUrole{doc}{variable}}}} provides the value to be output.  In each
case, the compute, fix, or variable must generate global values for
input to the {\hyperref[\detokenize{dump::doc}]{\sphinxcrossref{\DUrole{doc}{thermo\_style custom}}}} command.

\sphinxAtStartPar
Note that thermodynamic output values can be “extensive” or
“intensive”.  The former scale with the number of atoms in the system
(e.g. total energy), the latter do not (e.g. temperature).  The
setting for {\hyperref[\detokenize{thermo_modify::doc}]{\sphinxcrossref{\DUrole{doc}{thermo\_modify norm}}}} determines whether
extensive quantities are normalized or not.  Computes and fixes
produce either extensive or intensive values; see their individual doc
pages for details.  {\hyperref[\detokenize{variable::doc}]{\sphinxcrossref{\DUrole{doc}{Equal\sphinxhyphen{}style variables}}}} produce only
intensive values; you can include a division by “natoms” in the
formula if desired, to make an extensive calculation produce an
intensive result.


\subsubsection{Dump file output}
\label{\detokenize{Howto_output:dump-file-output}}\label{\detokenize{Howto_output:dump}}
\sphinxAtStartPar
Dump file output is specified by the {\hyperref[\detokenize{dump::doc}]{\sphinxcrossref{\DUrole{doc}{dump}}}} and
{\hyperref[\detokenize{dump_modify::doc}]{\sphinxcrossref{\DUrole{doc}{dump\_modify}}}} commands.  There are several
pre\sphinxhyphen{}defined formats (dump atom, dump xtc, etc).

\sphinxAtStartPar
There is also a {\hyperref[\detokenize{dump::doc}]{\sphinxcrossref{\DUrole{doc}{dump custom}}}} format where the user
specifies what values are output with each atom.  Pre\sphinxhyphen{}defined atom
attributes can be specified (id, x, fx, etc).  Three additional kinds
of keywords can also be specified (c\_ID, f\_ID, v\_name), where a
{\hyperref[\detokenize{compute::doc}]{\sphinxcrossref{\DUrole{doc}{compute}}}} or {\hyperref[\detokenize{fix::doc}]{\sphinxcrossref{\DUrole{doc}{fix}}}} or {\hyperref[\detokenize{variable::doc}]{\sphinxcrossref{\DUrole{doc}{variable}}}} provides the values to be output.  In each case, the
compute, fix, or variable must generate per\sphinxhyphen{}atom values for input to
the {\hyperref[\detokenize{dump::doc}]{\sphinxcrossref{\DUrole{doc}{dump custom}}}} command.

\sphinxAtStartPar
There is also a {\hyperref[\detokenize{dump::doc}]{\sphinxcrossref{\DUrole{doc}{dump local}}}} format where the user specifies
what local values to output.  A pre\sphinxhyphen{}defined index keyword can be
specified to enumerate the local values.  Two additional kinds of
keywords can also be specified (c\_ID, f\_ID), where a
{\hyperref[\detokenize{compute::doc}]{\sphinxcrossref{\DUrole{doc}{compute}}}} or {\hyperref[\detokenize{fix::doc}]{\sphinxcrossref{\DUrole{doc}{fix}}}} or {\hyperref[\detokenize{variable::doc}]{\sphinxcrossref{\DUrole{doc}{variable}}}}
provides the values to be output.  In each case, the compute or fix
must generate local values for input to the {\hyperref[\detokenize{dump::doc}]{\sphinxcrossref{\DUrole{doc}{dump local}}}}
command.

\sphinxAtStartPar
There is also a {\hyperref[\detokenize{dump::doc}]{\sphinxcrossref{\DUrole{doc}{dump grid}}}} format where the user
specifies what per\sphinxhyphen{}grid values to output from computes or fixes that
generate per\sphinxhyphen{}grid data.


\subsubsection{Fixes that write output files}
\label{\detokenize{Howto_output:fixes-that-write-output-files}}\label{\detokenize{Howto_output:fixoutput}}
\sphinxAtStartPar
Several fixes take various quantities as input and can write output
files: {\hyperref[\detokenize{fix_ave_time::doc}]{\sphinxcrossref{\DUrole{doc}{fix ave/time}}}}, {\hyperref[\detokenize{fix_ave_chunk::doc}]{\sphinxcrossref{\DUrole{doc}{fix ave/chunk}}}}, {\hyperref[\detokenize{fix_ave_histo::doc}]{\sphinxcrossref{\DUrole{doc}{fix ave/histo}}}}, {\hyperref[\detokenize{fix_ave_correlate::doc}]{\sphinxcrossref{\DUrole{doc}{fix
ave/correlate}}}}, and {\hyperref[\detokenize{fix_print::doc}]{\sphinxcrossref{\DUrole{doc}{fix print}}}}.

\sphinxAtStartPar
The {\hyperref[\detokenize{fix_ave_time::doc}]{\sphinxcrossref{\DUrole{doc}{fix ave/time}}}} command enables direct output
to a file and/or time\sphinxhyphen{}averaging of global scalars or vectors.  The
user specifies one or more quantities as input.  These can be global
{\hyperref[\detokenize{compute::doc}]{\sphinxcrossref{\DUrole{doc}{compute}}}} values, global {\hyperref[\detokenize{fix::doc}]{\sphinxcrossref{\DUrole{doc}{fix}}}} values, or
{\hyperref[\detokenize{variable::doc}]{\sphinxcrossref{\DUrole{doc}{variables}}}} of any style except the atom style which
produces per\sphinxhyphen{}atom values.  Since a variable can refer to keywords used
by the {\hyperref[\detokenize{thermo_style::doc}]{\sphinxcrossref{\DUrole{doc}{thermo\_style custom}}}} command (like temp or
press) and individual per\sphinxhyphen{}atom values, a wide variety of quantities
can be time averaged and/or output in this way.  If the inputs are one
or more scalar values, then the fix generate a global scalar or vector
of output.  If the inputs are one or more vector values, then the fix
generates a global vector or array of output.  The time\sphinxhyphen{}averaged
output of this fix can also be used as input to other output commands.

\sphinxAtStartPar
The {\hyperref[\detokenize{fix_ave_chunk::doc}]{\sphinxcrossref{\DUrole{doc}{fix ave/chunk}}}} command enables direct output
to a file of chunk\sphinxhyphen{}averaged per\sphinxhyphen{}atom quantities like those output in
dump files.  Chunks can represent spatial bins or other collections of
atoms, e.g. individual molecules.  The per\sphinxhyphen{}atom quantities can be atom
density (mass or number) or atom attributes such as position,
velocity, force.  They can also be per\sphinxhyphen{}atom quantities calculated by a
{\hyperref[\detokenize{compute::doc}]{\sphinxcrossref{\DUrole{doc}{compute}}}}, by a {\hyperref[\detokenize{fix::doc}]{\sphinxcrossref{\DUrole{doc}{fix}}}}, or by an atom\sphinxhyphen{}style
{\hyperref[\detokenize{variable::doc}]{\sphinxcrossref{\DUrole{doc}{variable}}}}.  The chunk\sphinxhyphen{}averaged output of this fix is
global and can also be used as input to other output commands.

\sphinxAtStartPar
Note that the {\hyperref[\detokenize{fix_ave_grid::doc}]{\sphinxcrossref{\DUrole{doc}{fix ave/grid}}}} command can also
average the same per\sphinxhyphen{}atom quantities within spatial bins, but it does
this for a distributed grid whose grid cells are owned by different
processors.  It outputs per\sphinxhyphen{}grid data, not global data, so it is more
efficient for large numbers of averaging bins.

\sphinxAtStartPar
The {\hyperref[\detokenize{fix_ave_histo::doc}]{\sphinxcrossref{\DUrole{doc}{fix ave/histo}}}} command enables direct output
to a file of histogrammed quantities, which can be global or per\sphinxhyphen{}atom
or local quantities.  The histogram output of this fix can also be
used as input to other output commands.

\sphinxAtStartPar
The {\hyperref[\detokenize{fix_ave_correlate::doc}]{\sphinxcrossref{\DUrole{doc}{fix ave/correlate}}}} command enables direct
output to a file of time\sphinxhyphen{}correlated quantities, which can be global
values.  The correlation matrix output of this fix can also be used as
input to other output commands.

\sphinxAtStartPar
The {\hyperref[\detokenize{fix_print::doc}]{\sphinxcrossref{\DUrole{doc}{fix print}}}} command can generate a line of output
written to the screen and log file or to a separate file, periodically
during a running simulation.  The line can contain one or more
{\hyperref[\detokenize{variable::doc}]{\sphinxcrossref{\DUrole{doc}{variable}}}} values for any style variable except the
vector or atom styles).  As explained above, variables themselves can
contain references to global values generated by {\hyperref[\detokenize{thermo_style::doc}]{\sphinxcrossref{\DUrole{doc}{thermodynamic
keywords}}}}, {\hyperref[\detokenize{compute::doc}]{\sphinxcrossref{\DUrole{doc}{computes}}}}, {\hyperref[\detokenize{fix::doc}]{\sphinxcrossref{\DUrole{doc}{fixes}}}}, or other {\hyperref[\detokenize{variable::doc}]{\sphinxcrossref{\DUrole{doc}{variables}}}}, or to per\sphinxhyphen{}atom values
for a specific atom.  Thus the {\hyperref[\detokenize{fix_print::doc}]{\sphinxcrossref{\DUrole{doc}{fix print}}}} command is
a means to output a wide variety of quantities separate from normal
thermodynamic or dump file output.


\subsubsection{Computes that process output quantities}
\label{\detokenize{Howto_output:computes-that-process-output-quantities}}\label{\detokenize{Howto_output:computeoutput}}
\sphinxAtStartPar
The {\hyperref[\detokenize{compute_reduce::doc}]{\sphinxcrossref{\DUrole{doc}{compute reduce}}}} and {\hyperref[\detokenize{compute_reduce::doc}]{\sphinxcrossref{\DUrole{doc}{compute
reduce/region}}}} commands take one or more per\sphinxhyphen{}atom or
local vector quantities as inputs and “reduce” them (sum, min, max,
ave) to scalar quantities.  These are produced as output values which
can be used as input to other output commands.

\sphinxAtStartPar
The {\hyperref[\detokenize{compute_slice::doc}]{\sphinxcrossref{\DUrole{doc}{compute slice}}}} command take one or more
global vector or array quantities as inputs and extracts a subset of
their values to create a new vector or array.  These are produced as
output values which can be used as input to other output commands.

\sphinxAtStartPar
The {\hyperref[\detokenize{compute_property_atom::doc}]{\sphinxcrossref{\DUrole{doc}{compute property/atom}}}} command takes
a list of one or more pre\sphinxhyphen{}defined atom attributes (id, x, fx, etc) and
stores the values in a per\sphinxhyphen{}atom vector or array.  These are produced
as output values which can be used as input to other output commands.
The list of atom attributes is the same as for the {\hyperref[\detokenize{dump::doc}]{\sphinxcrossref{\DUrole{doc}{dump custom}}}} command.

\sphinxAtStartPar
The {\hyperref[\detokenize{compute_property_local::doc}]{\sphinxcrossref{\DUrole{doc}{compute property/local}}}} command
takes a list of one or more pre\sphinxhyphen{}defined local attributes (bond info,
angle info, etc) and stores the values in a local vector or array.
These are produced as output values which can be used as input to
other output commands.

\sphinxAtStartPar
The {\hyperref[\detokenize{compute_property_grid::doc}]{\sphinxcrossref{\DUrole{doc}{compute property/grid}}}} command takes
a list of one or more pre\sphinxhyphen{}defined per\sphinxhyphen{}grid attributes (id, grid cell
coords, etc) and stores the values in a per\sphinxhyphen{}grid vector or array.
These are produced as output values which can be used as input to the
{\hyperref[\detokenize{dump::doc}]{\sphinxcrossref{\DUrole{doc}{dump grid}}}} command.

\sphinxAtStartPar
The {\hyperref[\detokenize{compute_property_chunk::doc}]{\sphinxcrossref{\DUrole{doc}{compute property/chunk}}}} command
takes a list of one or more pre\sphinxhyphen{}defined chunk attributes (id, count,
coords for spatial bins) and stores the values in a global vector or
array.  These are produced as output values which can be used as input
to other output commands.


\subsubsection{Fixes that process output quantities}
\label{\detokenize{Howto_output:fixes-that-process-output-quantities}}\label{\detokenize{Howto_output:fixprocoutput}}
\sphinxAtStartPar
The {\hyperref[\detokenize{fix_vector::doc}]{\sphinxcrossref{\DUrole{doc}{fix vector}}}} command can create global vectors as
output from global scalars as input, accumulating them one element at
a time.

\sphinxAtStartPar
The {\hyperref[\detokenize{fix_ave_atom::doc}]{\sphinxcrossref{\DUrole{doc}{fix ave/atom}}}} command performs time\sphinxhyphen{}averaging
of per\sphinxhyphen{}atom vectors.  The per\sphinxhyphen{}atom quantities can be atom attributes
such as position, velocity, force.  They can also be per\sphinxhyphen{}atom
quantities calculated by a {\hyperref[\detokenize{compute::doc}]{\sphinxcrossref{\DUrole{doc}{compute}}}}, by a {\hyperref[\detokenize{fix::doc}]{\sphinxcrossref{\DUrole{doc}{fix}}}}, or by an atom\sphinxhyphen{}style {\hyperref[\detokenize{variable::doc}]{\sphinxcrossref{\DUrole{doc}{variable}}}}.  The
time\sphinxhyphen{}averaged per\sphinxhyphen{}atom output of this fix can be used as input to
other output commands.

\sphinxAtStartPar
The {\hyperref[\detokenize{fix_store_state::doc}]{\sphinxcrossref{\DUrole{doc}{fix store/state}}}} command can archive one
or more per\sphinxhyphen{}atom attributes at a particular time, so that the old
values can be used in a future calculation or output.  The list of
atom attributes is the same as for the {\hyperref[\detokenize{dump::doc}]{\sphinxcrossref{\DUrole{doc}{dump custom}}}}
command, including per\sphinxhyphen{}atom quantities calculated by a {\hyperref[\detokenize{compute::doc}]{\sphinxcrossref{\DUrole{doc}{compute}}}}, by a {\hyperref[\detokenize{fix::doc}]{\sphinxcrossref{\DUrole{doc}{fix}}}}, or by an atom\sphinxhyphen{}style {\hyperref[\detokenize{variable::doc}]{\sphinxcrossref{\DUrole{doc}{variable}}}}.  The output of this fix can be used as input to other
output commands.

\sphinxAtStartPar
The {\hyperref[\detokenize{fix_ave_grid::doc}]{\sphinxcrossref{\DUrole{doc}{fix ave/grid}}}} command performs time\sphinxhyphen{}averaging
of either per\sphinxhyphen{}atom or per\sphinxhyphen{}grid data.

\sphinxAtStartPar
For per\sphinxhyphen{}atom data it performs averaging for the atoms within each grid
cell, similar to the {\hyperref[\detokenize{fix_ave_chunk::doc}]{\sphinxcrossref{\DUrole{doc}{fix ave/chunk}}}} command when
its chunks are defined as regular 2d or 3d bins.  The per\sphinxhyphen{}atom
quantities can be atom density (mass or number) or atom attributes
such as position, velocity, force.  They can also be per\sphinxhyphen{}atom
quantities calculated by a {\hyperref[\detokenize{compute::doc}]{\sphinxcrossref{\DUrole{doc}{compute}}}}, by a {\hyperref[\detokenize{fix::doc}]{\sphinxcrossref{\DUrole{doc}{fix}}}}, or by an atom\sphinxhyphen{}style {\hyperref[\detokenize{variable::doc}]{\sphinxcrossref{\DUrole{doc}{variable}}}}.

\sphinxAtStartPar
The chief difference between the {\hyperref[\detokenize{fix_ave_grid::doc}]{\sphinxcrossref{\DUrole{doc}{fix ave/grid}}}}
and {\hyperref[\detokenize{fix_ave_chunk::doc}]{\sphinxcrossref{\DUrole{doc}{fix ave/chunk}}}} commands when used in this
context is that the former uses a distributed grid, while the latter
uses a global grid.  Distributed means that each processor owns the
subset of grid cells within its subdomain.  Global means that each
processor owns a copy of the entire grid.  The {\hyperref[\detokenize{fix_ave_grid::doc}]{\sphinxcrossref{\DUrole{doc}{fix ave/grid}}}} command is thus more efficient for large grids.

\sphinxAtStartPar
For per\sphinxhyphen{}grid data, the {\hyperref[\detokenize{fix_ave_grid::doc}]{\sphinxcrossref{\DUrole{doc}{fix ave/grid}}}} command
takes inputs for grid data produced by other computes or fixes and
averages the values for each grid point over time.


\subsubsection{Computes that generate values to output}
\label{\detokenize{Howto_output:computes-that-generate-values-to-output}}\label{\detokenize{Howto_output:compute}}
\sphinxAtStartPar
Every {\hyperref[\detokenize{compute::doc}]{\sphinxcrossref{\DUrole{doc}{compute}}}} in LAMMPS produces either global or
per\sphinxhyphen{}atom or local or per\sphinxhyphen{}grid values.  The values can be scalars or
vectors or arrays of data.  These values can be output using the other
commands described in this section.  The page for each compute command
describes what it produces.  Computes that produce per\sphinxhyphen{}atom or local
or per\sphinxhyphen{}grid values have the word “atom” or “local” or “grid as the
last word in their style name.  Computes without the word “atom” or
“local” or “grid” produce global values.


\subsubsection{Fixes that generate values to output}
\label{\detokenize{Howto_output:fixes-that-generate-values-to-output}}\label{\detokenize{Howto_output:fix}}
\sphinxAtStartPar
Some {\hyperref[\detokenize{fix::doc}]{\sphinxcrossref{\DUrole{doc}{fixes}}}} in LAMMPS produces either global or per\sphinxhyphen{}atom
or local or per\sphinxhyphen{}grid values which can be accessed by other commands.
The values can be scalars or vectors or arrays of data.  These values
can be output using the other commands described in this section.  The
page for each fix command tells whether it produces any output
quantities and describes them.


\subsubsection{Variables that generate values to output}
\label{\detokenize{Howto_output:variables-that-generate-values-to-output}}\label{\detokenize{Howto_output:variable}}
\sphinxAtStartPar
{\hyperref[\detokenize{variable::doc}]{\sphinxcrossref{\DUrole{doc}{Variables}}}} defined in an input script can store one or
more strings.  But equal\sphinxhyphen{}style, vector\sphinxhyphen{}style, and atom\sphinxhyphen{}style or
atomfile\sphinxhyphen{}style variables generate a global scalar value, global vector
or values, or a per\sphinxhyphen{}atom vector, respectively, when accessed.  The
formulas used to define these variables can contain references to the
thermodynamic keywords and to global and per\sphinxhyphen{}atom data generated by
computes, fixes, and other variables.  The values generated by
variables can be used as input to and thus output by the other
commands described in this section.

\sphinxAtStartPar
Per\sphinxhyphen{}grid variables have not (yet) been implemented.


\subsubsection{Summary table of output options and data flow between commands}
\label{\detokenize{Howto_output:summary-table-of-output-options-and-data-flow-between-commands}}\label{\detokenize{Howto_output:table}}
\sphinxAtStartPar
This table summarizes the various commands that can be used for
generating output from LAMMPS.  Each command produces output data of
some kind and/or writes data to a file.  Most of the commands can take
data from other commands as input.  Thus you can link many of these
commands together in pipeline form, where data produced by one command
is used as input to another command and eventually written to the
screen or to a file.  Note that to hook two commands together the
output and input data types must match, e.g. global/per\sphinxhyphen{}atom/local
data and scalar/vector/array data.

\sphinxAtStartPar
Also note that, as described above, when a command takes a scalar as
input, that could also be an element of a vector or array.  Likewise a
vector input could be a column of an array.


\begin{savenotes}\sphinxattablestart
\sphinxthistablewithglobalstyle
\centering
\begin{tabulary}{\linewidth}[t]{TTT}
\sphinxtoprule
\sphinxtableatstartofbodyhook
\sphinxAtStartPar
Command
&
\sphinxAtStartPar
Input
&
\sphinxAtStartPar
Output
\\
\sphinxhline
\sphinxAtStartPar
{\hyperref[\detokenize{thermo_style::doc}]{\sphinxcrossref{\DUrole{doc}{thermo\_style custom}}}}
&
\sphinxAtStartPar
global scalars
&
\sphinxAtStartPar
screen, log file
\\
\sphinxhline
\sphinxAtStartPar
{\hyperref[\detokenize{dump::doc}]{\sphinxcrossref{\DUrole{doc}{dump custom}}}}
&
\sphinxAtStartPar
per\sphinxhyphen{}atom vectors
&
\sphinxAtStartPar
dump file
\\
\sphinxhline
\sphinxAtStartPar
{\hyperref[\detokenize{dump::doc}]{\sphinxcrossref{\DUrole{doc}{dump local}}}}
&
\sphinxAtStartPar
local vectors
&
\sphinxAtStartPar
dump file
\\
\sphinxhline
\sphinxAtStartPar
{\hyperref[\detokenize{dump::doc}]{\sphinxcrossref{\DUrole{doc}{dump grid}}}}
&
\sphinxAtStartPar
per\sphinxhyphen{}grid vectors
&
\sphinxAtStartPar
dump file
\\
\sphinxhline
\sphinxAtStartPar
{\hyperref[\detokenize{fix_print::doc}]{\sphinxcrossref{\DUrole{doc}{fix print}}}}
&
\sphinxAtStartPar
global scalar from variable
&
\sphinxAtStartPar
screen, file
\\
\sphinxhline
\sphinxAtStartPar
{\hyperref[\detokenize{print::doc}]{\sphinxcrossref{\DUrole{doc}{print}}}}
&
\sphinxAtStartPar
global scalar from variable
&
\sphinxAtStartPar
screen
\\
\sphinxhline
\sphinxAtStartPar
{\hyperref[\detokenize{compute::doc}]{\sphinxcrossref{\DUrole{doc}{computes}}}}
&
\sphinxAtStartPar
N/A
&
\sphinxAtStartPar
global/per\sphinxhyphen{}atom/local/per\sphinxhyphen{}grid scalar/vector/array
\\
\sphinxhline
\sphinxAtStartPar
{\hyperref[\detokenize{fix::doc}]{\sphinxcrossref{\DUrole{doc}{fixes}}}}
&
\sphinxAtStartPar
N/A
&
\sphinxAtStartPar
global/per\sphinxhyphen{}atom/local/per\sphinxhyphen{}grid scalar/vector/array
\\
\sphinxhline
\sphinxAtStartPar
{\hyperref[\detokenize{variable::doc}]{\sphinxcrossref{\DUrole{doc}{variables}}}}
&
\sphinxAtStartPar
global scalars and vectors, per\sphinxhyphen{}atom vectors
&
\sphinxAtStartPar
global scalar and vector, per\sphinxhyphen{}atom vector
\\
\sphinxhline
\sphinxAtStartPar
{\hyperref[\detokenize{compute_reduce::doc}]{\sphinxcrossref{\DUrole{doc}{compute reduce}}}}
&
\sphinxAtStartPar
per\sphinxhyphen{}atom/local vectors
&
\sphinxAtStartPar
global scalar/vector
\\
\sphinxhline
\sphinxAtStartPar
{\hyperref[\detokenize{compute_slice::doc}]{\sphinxcrossref{\DUrole{doc}{compute slice}}}}
&
\sphinxAtStartPar
global vectors/arrays
&
\sphinxAtStartPar
global vector/array
\\
\sphinxhline
\sphinxAtStartPar
{\hyperref[\detokenize{compute_property_atom::doc}]{\sphinxcrossref{\DUrole{doc}{compute property/atom}}}}
&
\sphinxAtStartPar
N/A
&
\sphinxAtStartPar
per\sphinxhyphen{}atom vector/array
\\
\sphinxhline
\sphinxAtStartPar
{\hyperref[\detokenize{compute_property_local::doc}]{\sphinxcrossref{\DUrole{doc}{compute property/local}}}}
&
\sphinxAtStartPar
N/A
&
\sphinxAtStartPar
local vector/array
\\
\sphinxhline
\sphinxAtStartPar
{\hyperref[\detokenize{compute_property_grid::doc}]{\sphinxcrossref{\DUrole{doc}{compute property/grid}}}}
&
\sphinxAtStartPar
N/A
&
\sphinxAtStartPar
per\sphinxhyphen{}grid vector/array
\\
\sphinxhline
\sphinxAtStartPar
{\hyperref[\detokenize{compute_property_chunk::doc}]{\sphinxcrossref{\DUrole{doc}{compute property/chunk}}}}
&
\sphinxAtStartPar
N/A
&
\sphinxAtStartPar
global vector/array
\\
\sphinxhline
\sphinxAtStartPar
{\hyperref[\detokenize{fix_vector::doc}]{\sphinxcrossref{\DUrole{doc}{fix vector}}}}
&
\sphinxAtStartPar
global scalars
&
\sphinxAtStartPar
global vector
\\
\sphinxhline
\sphinxAtStartPar
{\hyperref[\detokenize{fix_ave_atom::doc}]{\sphinxcrossref{\DUrole{doc}{fix ave/atom}}}}
&
\sphinxAtStartPar
per\sphinxhyphen{}atom vectors
&
\sphinxAtStartPar
per\sphinxhyphen{}atom vector/array
\\
\sphinxhline
\sphinxAtStartPar
{\hyperref[\detokenize{fix_ave_time::doc}]{\sphinxcrossref{\DUrole{doc}{fix ave/time}}}}
&
\sphinxAtStartPar
global scalars/vectors
&
\sphinxAtStartPar
global scalar/vector/array, file
\\
\sphinxhline
\sphinxAtStartPar
{\hyperref[\detokenize{fix_ave_chunk::doc}]{\sphinxcrossref{\DUrole{doc}{fix ave/chunk}}}}
&
\sphinxAtStartPar
per\sphinxhyphen{}atom vectors
&
\sphinxAtStartPar
global array, file
\\
\sphinxhline
\sphinxAtStartPar
{\hyperref[\detokenize{fix_ave_grid::doc}]{\sphinxcrossref{\DUrole{doc}{fix ave/grid}}}}
&
\sphinxAtStartPar
per\sphinxhyphen{}atom vectors or per\sphinxhyphen{}grid vectors
&
\sphinxAtStartPar
per\sphinxhyphen{}grid vector/array
\\
\sphinxhline
\sphinxAtStartPar
{\hyperref[\detokenize{fix_ave_histo::doc}]{\sphinxcrossref{\DUrole{doc}{fix ave/histo}}}}
&
\sphinxAtStartPar
global/per\sphinxhyphen{}atom/local scalars and vectors
&
\sphinxAtStartPar
global array, file
\\
\sphinxhline
\sphinxAtStartPar
{\hyperref[\detokenize{fix_ave_correlate::doc}]{\sphinxcrossref{\DUrole{doc}{fix ave/correlate}}}}
&
\sphinxAtStartPar
global scalars
&
\sphinxAtStartPar
global array, file
\\
\sphinxhline
\sphinxAtStartPar
{\hyperref[\detokenize{fix_store_state::doc}]{\sphinxcrossref{\DUrole{doc}{fix store/state}}}}
&
\sphinxAtStartPar
per\sphinxhyphen{}atom vectors
&
\sphinxAtStartPar
per\sphinxhyphen{}atom vector/array
\\
\sphinxbottomrule
\end{tabulary}
\sphinxtableafterendhook\par
\sphinxattableend\end{savenotes}

\sphinxstepscope


\subsection{Use chunks to calculate system properties}
\label{\detokenize{Howto_chunk:use-chunks-to-calculate-system-properties}}\label{\detokenize{Howto_chunk::doc}}
\sphinxAtStartPar
In LAMMPS, “chunks” are collections of atoms, as defined by the
{\hyperref[\detokenize{compute_chunk_atom::doc}]{\sphinxcrossref{\DUrole{doc}{compute chunk/atom}}}} command, which assigns
each atom to a chunk ID (or to no chunk at all).  The number of chunks
and the assignment of chunk IDs to atoms can be static or change over
time.  Examples of “chunks” are molecules or spatial bins or atoms
with similar values (e.g. coordination number or potential energy).

\sphinxAtStartPar
The per\sphinxhyphen{}atom chunk IDs can be used as input to two other kinds of
commands, to calculate various properties of a system:
\begin{itemize}
\item {} 
\sphinxAtStartPar
{\hyperref[\detokenize{fix_ave_chunk::doc}]{\sphinxcrossref{\DUrole{doc}{fix ave/chunk}}}}

\item {} 
\sphinxAtStartPar
any of the {\hyperref[\detokenize{compute::doc}]{\sphinxcrossref{\DUrole{doc}{compute */chunk}}}} commands

\end{itemize}

\sphinxAtStartPar
Here a brief overview for each of the 4 kinds of chunk\sphinxhyphen{}related commands
is provided.  Then some examples are given of how to compute different
properties with chunk commands.


\subsubsection{Compute chunk/atom command:}
\label{\detokenize{Howto_chunk:compute-chunk-atom-command}}
\sphinxAtStartPar
This compute can assign atoms to chunks of various styles.  Only atoms
in the specified group and optional specified region are assigned to a
chunk.  Here are some possible chunk definitions:


\begin{savenotes}\sphinxattablestart
\sphinxthistablewithglobalstyle
\centering
\begin{tabulary}{\linewidth}[t]{TT}
\sphinxtoprule
\sphinxtableatstartofbodyhook
\sphinxAtStartPar
atoms in same molecule
&
\sphinxAtStartPar
chunk ID = molecule ID
\\
\sphinxhline
\sphinxAtStartPar
atoms of same atom type
&
\sphinxAtStartPar
chunk ID = atom type
\\
\sphinxhline
\sphinxAtStartPar
all atoms with same atom property (charge, radius, etc)
&
\sphinxAtStartPar
chunk ID = output of compute property/atom
\\
\sphinxhline
\sphinxAtStartPar
atoms in same cluster
&
\sphinxAtStartPar
chunk ID = output of {\hyperref[\detokenize{compute_cluster_atom::doc}]{\sphinxcrossref{\DUrole{doc}{compute cluster/atom}}}} command
\\
\sphinxhline
\sphinxAtStartPar
atoms in same spatial bin
&
\sphinxAtStartPar
chunk ID = bin ID
\\
\sphinxhline
\sphinxAtStartPar
atoms in same rigid body
&
\sphinxAtStartPar
chunk ID = molecule ID used to define rigid bodies
\\
\sphinxhline
\sphinxAtStartPar
atoms with similar potential energy
&
\sphinxAtStartPar
chunk ID = output of {\hyperref[\detokenize{compute_pe_atom::doc}]{\sphinxcrossref{\DUrole{doc}{compute pe/atom}}}}
\\
\sphinxhline
\sphinxAtStartPar
atoms with same local defect structure
&
\sphinxAtStartPar
chunk ID = output of {\hyperref[\detokenize{compute_centro_atom::doc}]{\sphinxcrossref{\DUrole{doc}{compute centro/atom}}}} or {\hyperref[\detokenize{compute_coord_atom::doc}]{\sphinxcrossref{\DUrole{doc}{compute coord/atom}}}} command
\\
\sphinxbottomrule
\end{tabulary}
\sphinxtableafterendhook\par
\sphinxattableend\end{savenotes}

\sphinxAtStartPar
Note that chunk IDs are integer values, so for atom properties or
computes that produce a floating point value, they will be truncated
to an integer.  You could also use the compute in a variable that
scales the floating point value to spread it across multiple integers.

\sphinxAtStartPar
Spatial bins can be of various kinds, e.g. 1d bins = slabs, 2d bins =
pencils, 3d bins = boxes, spherical bins, cylindrical bins.

\sphinxAtStartPar
This compute also calculates the number of chunks \sphinxstyleemphasis{Nchunk}, which is
used by other commands to tally per\sphinxhyphen{}chunk data.  \sphinxstyleemphasis{Nchunk} can be a
static value or change over time (e.g. the number of clusters).  The
chunk ID for an individual atom can also be static (e.g. a molecule
ID), or dynamic (e.g. what spatial bin an atom is in as it moves).

\sphinxAtStartPar
Note that this compute allows the per\sphinxhyphen{}atom output of other
{\hyperref[\detokenize{compute::doc}]{\sphinxcrossref{\DUrole{doc}{computes}}}}, {\hyperref[\detokenize{fix::doc}]{\sphinxcrossref{\DUrole{doc}{fixes}}}}, and
{\hyperref[\detokenize{variable::doc}]{\sphinxcrossref{\DUrole{doc}{variables}}}} to be used to define chunk IDs for each
atom.  This means you can write your own compute or fix to output a
per\sphinxhyphen{}atom quantity to use as chunk ID.  See the {\hyperref[\detokenize{Modify::doc}]{\sphinxcrossref{\DUrole{doc}{Modify}}}}
doc pages for info on how to do this.  You can also define a {\hyperref[\detokenize{variable::doc}]{\sphinxcrossref{\DUrole{doc}{per\sphinxhyphen{}atom variable}}}} in the input script that uses a formula to
generate a chunk ID for each atom.


\subsubsection{Fix ave/chunk command:}
\label{\detokenize{Howto_chunk:fix-ave-chunk-command}}
\sphinxAtStartPar
This fix takes the ID of a {\hyperref[\detokenize{compute_chunk_atom::doc}]{\sphinxcrossref{\DUrole{doc}{compute chunk/atom}}}} command as input.  For each chunk,
it then sums one or more specified per\sphinxhyphen{}atom values over the atoms in
each chunk.  The per\sphinxhyphen{}atom values can be any atom property, such as
velocity, force, charge, potential energy, kinetic energy, stress,
etc.  Additional keywords are defined for per\sphinxhyphen{}chunk properties like
density and temperature.  More generally any per\sphinxhyphen{}atom value generated
by other {\hyperref[\detokenize{compute::doc}]{\sphinxcrossref{\DUrole{doc}{computes}}}}, {\hyperref[\detokenize{fix::doc}]{\sphinxcrossref{\DUrole{doc}{fixes}}}}, and {\hyperref[\detokenize{variable::doc}]{\sphinxcrossref{\DUrole{doc}{per\sphinxhyphen{}atom variables}}}}, can be summed over atoms in each chunk.

\sphinxAtStartPar
Similar to other averaging fixes, this fix allows the summed per\sphinxhyphen{}chunk
values to be time\sphinxhyphen{}averaged in various ways, and output to a file.  The
fix produces a global array as output with one row of values per
chunk.


\subsubsection{Compute */chunk commands:}
\label{\detokenize{Howto_chunk:compute-chunk-commands}}
\sphinxAtStartPar
The following computes operate on chunks of atoms to produce per\sphinxhyphen{}chunk
values.  Any compute whose style name ends in “/chunk” is in this
category:
\begin{itemize}
\item {} 
\sphinxAtStartPar
{\hyperref[\detokenize{compute_com_chunk::doc}]{\sphinxcrossref{\DUrole{doc}{compute com/chunk}}}}

\item {} 
\sphinxAtStartPar
{\hyperref[\detokenize{compute_gyration_chunk::doc}]{\sphinxcrossref{\DUrole{doc}{compute gyration/chunk}}}}

\item {} 
\sphinxAtStartPar
{\hyperref[\detokenize{compute_inertia_chunk::doc}]{\sphinxcrossref{\DUrole{doc}{compute inertia/chunk}}}}

\item {} 
\sphinxAtStartPar
{\hyperref[\detokenize{compute_msd_chunk::doc}]{\sphinxcrossref{\DUrole{doc}{compute msd/chunk}}}}

\item {} 
\sphinxAtStartPar
{\hyperref[\detokenize{compute_property_chunk::doc}]{\sphinxcrossref{\DUrole{doc}{compute property/chunk}}}}

\item {} 
\sphinxAtStartPar
{\hyperref[\detokenize{compute_temp_chunk::doc}]{\sphinxcrossref{\DUrole{doc}{compute temp/chunk}}}}

\item {} 
\sphinxAtStartPar
{\hyperref[\detokenize{compute_vcm_chunk::doc}]{\sphinxcrossref{\DUrole{doc}{compute torque/chunk}}}}

\item {} 
\sphinxAtStartPar
{\hyperref[\detokenize{compute_vcm_chunk::doc}]{\sphinxcrossref{\DUrole{doc}{compute vcm/chunk}}}}

\end{itemize}

\sphinxAtStartPar
They each take the ID of a {\hyperref[\detokenize{compute_chunk_atom::doc}]{\sphinxcrossref{\DUrole{doc}{compute chunk/atom}}}} command as input.  As their names
indicate, they calculate the center\sphinxhyphen{}of\sphinxhyphen{}mass, radius of gyration,
moments of inertia, mean\sphinxhyphen{}squared displacement, temperature, torque,
and velocity of center\sphinxhyphen{}of\sphinxhyphen{}mass for each chunk of atoms.  The {\hyperref[\detokenize{compute_property_chunk::doc}]{\sphinxcrossref{\DUrole{doc}{compute property/chunk}}}} command can tally the
count of atoms in each chunk and extract other per\sphinxhyphen{}chunk properties.

\sphinxAtStartPar
The reason these various calculations are not part of the {\hyperref[\detokenize{fix_ave_chunk::doc}]{\sphinxcrossref{\DUrole{doc}{fix ave/chunk command}}}}, is that each requires a more
complicated operation than simply summing and averaging over per\sphinxhyphen{}atom
values in each chunk.  For example, many of them require calculation
of a center of mass, which requires summing mass*position over the
atoms and then dividing by summed mass.

\sphinxAtStartPar
All of these computes produce a global vector or global array as
output, with one or more values per chunk.  The output can be used in
various ways:
\begin{itemize}
\item {} 
\sphinxAtStartPar
As input to the {\hyperref[\detokenize{fix_ave_time::doc}]{\sphinxcrossref{\DUrole{doc}{fix ave/time}}}} command, which can
write the values to a file and optionally time average them.

\item {} 
\sphinxAtStartPar
As input to the {\hyperref[\detokenize{fix_ave_histo::doc}]{\sphinxcrossref{\DUrole{doc}{fix ave/histo}}}} command to
histogram values across chunks.  E.g. a histogram of cluster sizes or
molecule diffusion rates.

\item {} 
\sphinxAtStartPar
As input to special functions of {\hyperref[\detokenize{variable::doc}]{\sphinxcrossref{\DUrole{doc}{equal\sphinxhyphen{}style variables}}}}, like sum() and max() and ave().  E.g. to
find the largest cluster or fastest diffusing molecule or average
radius\sphinxhyphen{}of\sphinxhyphen{}gyration of a set of molecules (chunks).

\end{itemize}


\subsubsection{Other chunk commands:}
\label{\detokenize{Howto_chunk:other-chunk-commands}}\begin{itemize}
\item {} 
\sphinxAtStartPar
{\hyperref[\detokenize{compute_chunk_spread_atom::doc}]{\sphinxcrossref{\DUrole{doc}{compute chunk/spread/atom}}}}

\item {} 
\sphinxAtStartPar
{\hyperref[\detokenize{compute_reduce_chunk::doc}]{\sphinxcrossref{\DUrole{doc}{compute reduce/chunk}}}}

\end{itemize}

\sphinxAtStartPar
The {\hyperref[\detokenize{compute_chunk_spread_atom::doc}]{\sphinxcrossref{\DUrole{doc}{compute chunk/spread/atom}}}} command
spreads per\sphinxhyphen{}chunk values to each atom in the chunk, producing per\sphinxhyphen{}atom
values as its output.  This can be useful for outputting per\sphinxhyphen{}chunk
values to a per\sphinxhyphen{}atom {\hyperref[\detokenize{dump::doc}]{\sphinxcrossref{\DUrole{doc}{dump file}}}}.  Or for using an atom’s
associated chunk value in an {\hyperref[\detokenize{variable::doc}]{\sphinxcrossref{\DUrole{doc}{atom\sphinxhyphen{}style variable}}}}.  Or
as input to the {\hyperref[\detokenize{fix_ave_chunk::doc}]{\sphinxcrossref{\DUrole{doc}{fix ave/chunk}}}} command to
spatially average per\sphinxhyphen{}chunk values calculated by a per\sphinxhyphen{}chunk compute.

\sphinxAtStartPar
The {\hyperref[\detokenize{compute_reduce_chunk::doc}]{\sphinxcrossref{\DUrole{doc}{compute reduce/chunk}}}} command reduces a
peratom value across the atoms in each chunk to produce a value per
chunk.  When used with the {\hyperref[\detokenize{compute_chunk_spread_atom::doc}]{\sphinxcrossref{\DUrole{doc}{compute chunk/spread/atom}}}} command it can
create peratom values that induce a new set of chunks with a second
{\hyperref[\detokenize{compute_chunk_atom::doc}]{\sphinxcrossref{\DUrole{doc}{compute chunk/atom}}}} command.


\subsubsection{Example calculations with chunks}
\label{\detokenize{Howto_chunk:example-calculations-with-chunks}}
\sphinxAtStartPar
Here are examples using chunk commands to calculate various
properties:
\begin{enumerate}
\sphinxsetlistlabels{\arabic}{enumi}{enumii}{}{.}%
\item {} 
\sphinxAtStartPar
Average velocity in each of 1000 2d spatial bins:

\end{enumerate}

\begin{sphinxVerbatim}[commandchars=\\\{\}]
\PYG{k}{compute }\PYG{n+nv+nvIdentifier}{cc1}\PYG{+w}{ }\PYG{n+nv+nvIdentifier}{all}\PYG{+w}{ }\PYG{n}{chunk}\PYG{o}{/}\PYG{n}{atom}\PYG{+w}{ }\PYG{n}{bin}\PYG{o}{/}\PYG{l+m}{2}\PYG{n}{d}\PYG{+w}{ }\PYG{n}{x}\PYG{+w}{ }\PYG{l+m}{0.0}\PYG{+w}{ }\PYG{l+m}{0.1}\PYG{+w}{ }\PYG{n}{y}\PYG{+w}{ }\PYG{n}{lower}\PYG{+w}{ }\PYG{l+m}{0.01}\PYG{+w}{ }\PYG{n}{units}\PYG{+w}{ }\PYG{n}{reduced}
\PYG{k}{fix }\PYG{n+nv+nvIdentifier}{1}\PYG{+w}{ }\PYG{n+nv+nvIdentifier}{all}\PYG{+w}{ }\PYG{n}{ave}\PYG{o}{/}\PYG{n}{chunk}\PYG{+w}{ }\PYG{l+m}{100}\PYG{+w}{ }\PYG{l+m}{10}\PYG{+w}{ }\PYG{l+m}{1000}\PYG{+w}{ }\PYG{n}{cc1}\PYG{+w}{ }\PYG{n}{vx}\PYG{+w}{ }\PYG{n}{vy}\PYG{+w}{ }\PYG{n}{file}\PYG{+w}{ }\PYG{n}{tmp.out}
\end{sphinxVerbatim}

\sphinxAtStartPar
2. Temperature in each spatial bin, after subtracting a flow
velocity:

\begin{sphinxVerbatim}[commandchars=\\\{\}]
\PYG{k}{compute }\PYG{n+nv+nvIdentifier}{cc1}\PYG{+w}{ }\PYG{n+nv+nvIdentifier}{all}\PYG{+w}{ }\PYG{n}{chunk}\PYG{o}{/}\PYG{n}{atom}\PYG{+w}{ }\PYG{n}{bin}\PYG{o}{/}\PYG{l+m}{2}\PYG{n}{d}\PYG{+w}{ }\PYG{n}{x}\PYG{+w}{ }\PYG{l+m}{0.0}\PYG{+w}{ }\PYG{l+m}{0.1}\PYG{+w}{ }\PYG{n}{y}\PYG{+w}{ }\PYG{n}{lower}\PYG{+w}{ }\PYG{l+m}{0.1}\PYG{+w}{ }\PYG{n}{units}\PYG{+w}{ }\PYG{n}{reduced}
\PYG{k}{compute }\PYG{n+nv+nvIdentifier}{vbias}\PYG{+w}{ }\PYG{n+nv+nvIdentifier}{all}\PYG{+w}{ }\PYG{n}{temp}\PYG{o}{/}\PYG{n}{profile}\PYG{+w}{ }\PYG{l+m}{1}\PYG{+w}{ }\PYG{l+m}{0}\PYG{+w}{ }\PYG{l+m}{0}\PYG{+w}{ }\PYG{n}{y}\PYG{+w}{ }\PYG{l+m}{10}
\PYG{k}{fix }\PYG{n+nv+nvIdentifier}{1}\PYG{+w}{ }\PYG{n+nv+nvIdentifier}{all}\PYG{+w}{ }\PYG{n}{ave}\PYG{o}{/}\PYG{n}{chunk}\PYG{+w}{ }\PYG{l+m}{100}\PYG{+w}{ }\PYG{l+m}{10}\PYG{+w}{ }\PYG{l+m}{1000}\PYG{+w}{ }\PYG{n}{cc1}\PYG{+w}{ }\PYG{n}{temp}\PYG{+w}{ }\PYG{n}{bias}\PYG{+w}{ }\PYG{n}{vbias}\PYG{+w}{ }\PYG{n}{file}\PYG{+w}{ }\PYG{n}{tmp.out}
\end{sphinxVerbatim}
\begin{enumerate}
\sphinxsetlistlabels{\arabic}{enumi}{enumii}{}{.}%
\setcounter{enumi}{2}
\item {} 
\sphinxAtStartPar
Center of mass of each molecule:

\end{enumerate}

\begin{sphinxVerbatim}[commandchars=\\\{\}]
\PYG{k}{compute }\PYG{n+nv+nvIdentifier}{cc1}\PYG{+w}{ }\PYG{n+nv+nvIdentifier}{all}\PYG{+w}{ }\PYG{n}{chunk}\PYG{o}{/}\PYG{n}{atom}\PYG{+w}{ }\PYG{n}{molecule}
\PYG{k}{compute }\PYG{n+nv+nvIdentifier}{myChunk}\PYG{+w}{ }\PYG{n+nv+nvIdentifier}{all}\PYG{+w}{ }\PYG{n}{com}\PYG{o}{/}\PYG{n}{chunk}\PYG{+w}{ }\PYG{n}{cc1}
\PYG{k}{fix }\PYG{n+nv+nvIdentifier}{1}\PYG{+w}{ }\PYG{n+nv+nvIdentifier}{all}\PYG{+w}{ }\PYG{n}{ave}\PYG{o}{/}\PYG{n}{time}\PYG{+w}{ }\PYG{l+m}{100}\PYG{+w}{ }\PYG{l+m}{1}\PYG{+w}{ }\PYG{l+m}{100}\PYG{+w}{ }\PYG{n}{c\PYGZus{}myChunk[}\PYG{o}{*}\PYG{n}{]}\PYG{+w}{ }\PYG{n}{file}\PYG{+w}{ }\PYG{n}{tmp.out}\PYG{+w}{ }\PYG{n}{mode}\PYG{+w}{ }\PYG{n}{vector}
\end{sphinxVerbatim}
\begin{enumerate}
\sphinxsetlistlabels{\arabic}{enumi}{enumii}{}{.}%
\setcounter{enumi}{3}
\item {} 
\sphinxAtStartPar
Total force on each molecule and ave/max across all molecules:

\end{enumerate}

\begin{sphinxVerbatim}[commandchars=\\\{\}]
\PYG{k}{compute }\PYG{n+nv+nvIdentifier}{cc1}\PYG{+w}{ }\PYG{n+nv+nvIdentifier}{all}\PYG{+w}{ }\PYG{n}{chunk}\PYG{o}{/}\PYG{n}{atom}\PYG{+w}{ }\PYG{n}{molecule}
\PYG{k}{fix }\PYG{n+nv+nvIdentifier}{1}\PYG{+w}{ }\PYG{n+nv+nvIdentifier}{all}\PYG{+w}{ }\PYG{n}{ave}\PYG{o}{/}\PYG{n}{chunk}\PYG{+w}{ }\PYG{l+m}{1000}\PYG{+w}{ }\PYG{l+m}{1}\PYG{+w}{ }\PYG{l+m}{1000}\PYG{+w}{ }\PYG{n}{cc1}\PYG{+w}{ }\PYG{n}{fx}\PYG{+w}{ }\PYG{n}{fy}\PYG{+w}{ }\PYG{n}{fz}\PYG{+w}{ }\PYG{n}{file}\PYG{+w}{ }\PYG{n}{tmp.out}
\PYG{k}{variable }\PYG{n+nv+nvIdentifier}{xave}\PYG{+w}{ }\PYG{n}{equal}\PYG{+w}{ }\PYG{n}{ave}\PYG{n+nv}{(}\PYG{n+nv}{f\PYGZus{}1[2]}\PYG{n+nv}{)}
\PYG{k}{variable }\PYG{n+nv+nvIdentifier}{xmax}\PYG{+w}{ }\PYG{n}{equal}\PYG{+w}{ }\PYG{n}{max}\PYG{n+nv}{(}\PYG{n+nv}{f\PYGZus{}1[2]}\PYG{n+nv}{)}
\PYG{k}{thermo}\PYG{+w}{ }\PYG{l+m}{1000}
\PYG{k}{thermo\PYGZus{}style}\PYG{+w}{ }\PYG{n}{custom}\PYG{+w}{ }\PYG{n}{step}\PYG{+w}{ }\PYG{n}{temp}\PYG{+w}{ }\PYG{n}{v\PYGZus{}xave}\PYG{+w}{ }\PYG{n}{v\PYGZus{}xmax}
\end{sphinxVerbatim}
\begin{enumerate}
\sphinxsetlistlabels{\arabic}{enumi}{enumii}{}{.}%
\setcounter{enumi}{4}
\item {} 
\sphinxAtStartPar
Histogram of cluster sizes:

\end{enumerate}

\begin{sphinxVerbatim}[commandchars=\\\{\}]
\PYG{k}{compute }\PYG{n+nv+nvIdentifier}{cluster}\PYG{+w}{ }\PYG{n+nv+nvIdentifier}{all}\PYG{+w}{ }\PYG{n}{cluster}\PYG{o}{/}\PYG{n}{atom}\PYG{+w}{ }\PYG{l+m}{1.0}
\PYG{k}{compute }\PYG{n+nv+nvIdentifier}{cc1}\PYG{+w}{ }\PYG{n+nv+nvIdentifier}{all}\PYG{+w}{ }\PYG{n}{chunk}\PYG{o}{/}\PYG{n}{atom}\PYG{+w}{ }\PYG{n}{c\PYGZus{}cluster}\PYG{+w}{ }\PYG{n}{compress}\PYG{+w}{ }\PYG{n}{yes}
\PYG{k}{compute }\PYG{n+nv+nvIdentifier}{size}\PYG{+w}{ }\PYG{n+nv+nvIdentifier}{all}\PYG{+w}{ }\PYG{n}{property}\PYG{o}{/}\PYG{n}{chunk}\PYG{+w}{ }\PYG{n}{cc1}\PYG{+w}{ }\PYG{n}{count}
\PYG{k}{fix }\PYG{n+nv+nvIdentifier}{1}\PYG{+w}{ }\PYG{n+nv+nvIdentifier}{all}\PYG{+w}{ }\PYG{n}{ave}\PYG{o}{/}\PYG{n}{histo}\PYG{+w}{ }\PYG{l+m}{100}\PYG{+w}{ }\PYG{l+m}{1}\PYG{+w}{ }\PYG{l+m}{100}\PYG{+w}{ }\PYG{l+m}{0}\PYG{+w}{ }\PYG{l+m}{20}\PYG{+w}{ }\PYG{l+m}{20}\PYG{+w}{ }\PYG{n}{c\PYGZus{}size}\PYG{+w}{ }\PYG{n}{mode}\PYG{+w}{ }\PYG{n}{vector}\PYG{+w}{ }\PYG{n}{ave}\PYG{+w}{ }\PYG{n}{running}\PYG{+w}{ }\PYG{n}{beyond}\PYG{+w}{ }\PYG{n}{ignore}\PYG{+w}{ }\PYG{n}{file}\PYG{+w}{ }\PYG{n}{tmp.histo}
\end{sphinxVerbatim}

\sphinxAtStartPar
6. An example for using a per\sphinxhyphen{}chunk value to apply per\sphinxhyphen{}atom forces to
compress individual polymer chains (molecules) in a mixture, is
explained on the {\hyperref[\detokenize{compute_chunk_spread_atom::doc}]{\sphinxcrossref{\DUrole{doc}{compute chunk/spread/atom}}}} command doc page.

\sphinxAtStartPar
7. An example for using one set of per\sphinxhyphen{}chunk values for molecule
chunks, to create a second set of micelle\sphinxhyphen{}scale chunks (clustered
molecules, due to hydrophobicity), is explained on the
{\hyperref[\detokenize{compute_reduce_chunk::doc}]{\sphinxcrossref{\DUrole{doc}{compute reduce/chunk}}}} command doc page.

\sphinxAtStartPar
8. An example for using one set of per\sphinxhyphen{}chunk values (dipole moment
vectors) for molecule chunks, spreading the values to each atom in
each chunk, then defining a second set of chunks as spatial bins, and
using the {\hyperref[\detokenize{fix_ave_chunk::doc}]{\sphinxcrossref{\DUrole{doc}{fix ave/chunk}}}} command to calculate an
average dipole moment vector for each bin.  This example is explained
on the {\hyperref[\detokenize{compute_chunk_spread_atom::doc}]{\sphinxcrossref{\DUrole{doc}{compute chunk/spread/atom}}}}
command doc page.

\sphinxstepscope


\subsection{Using distributed grids}
\label{\detokenize{Howto_grid:using-distributed-grids}}\label{\detokenize{Howto_grid::doc}}
\sphinxAtStartPar
\DUrole{versionmodified,added}{Added in version 22Dec2022.}

\sphinxAtStartPar
LAMMPS has internal capabilities to create uniformly spaced grids
which overlay the simulation domain.  For 2d and 3d simulations these
are 2d and 3d grids respectively.  Conceptually a grid can be thought
of as a collection of grid cells.  Each grid cell can store one or
more values (data).

\sphinxAtStartPar
The grid cells and data they store are distributed across processors.
Each processor owns the grid cells (and data) whose center points lie
within the spatial subdomain of the processor.  If needed for its
computations, a processor may also store ghost grid cells with their
data.

\sphinxAtStartPar
Distributed grids can overlay orthogonal or triclinic simulation
boxes; see the {\hyperref[\detokenize{Howto_triclinic::doc}]{\sphinxcrossref{\DUrole{doc}{Howto triclinic}}}} doc page for
an explanation of the latter.  For a triclinic box, the grid cell
shape conforms to the shape of the simulation domain,
e.g. parallelograms instead of rectangles in 2d.

\sphinxAtStartPar
If the box size or shape changes during a simulation, the grid changes
with it, so that it always overlays the entire simulation domain.  For
non\sphinxhyphen{}periodic dimensions, the grid size in that dimension matches the
box size, as set by the {\hyperref[\detokenize{boundary::doc}]{\sphinxcrossref{\DUrole{doc}{boundary}}}} command for fixed
or shrink\sphinxhyphen{}wrapped boundaries.

\sphinxAtStartPar
If load\sphinxhyphen{}balancing is invoked by the {\hyperref[\detokenize{balance::doc}]{\sphinxcrossref{\DUrole{doc}{balance}}}} or
{\hyperref[\detokenize{fix_balance::doc}]{\sphinxcrossref{\DUrole{doc}{fix balance}}}} commands, then the subdomain owned
by a processor can change which may also change which grid cells they
own.

\sphinxAtStartPar
Post\sphinxhyphen{}processing and visualization of grid cell data can be enabled by
the {\hyperref[\detokenize{dump::doc}]{\sphinxcrossref{\DUrole{doc}{dump grid}}}}, {\hyperref[\detokenize{dump::doc}]{\sphinxcrossref{\DUrole{doc}{dump grid/vtk}}}}, and
{\hyperref[\detokenize{dump_image::doc}]{\sphinxcrossref{\DUrole{doc}{dump image}}}} commands.  The latter has an optional
\sphinxstyleemphasis{grid} keyword.  The \sphinxhref{https://www.ovito.org}{OVITO visualization tool} also plans (as of Nov 2022) to add support
for visualizing grid cell data (along with atoms) using {\hyperref[\detokenize{dump::doc}]{\sphinxcrossref{\DUrole{doc}{dump
grid}}}} output files as input.

\begin{sphinxadmonition}{note}{Note:}
\sphinxAtStartPar
For developers, distributed grids are implemented within the code via
two classes: Grid2d and Grid3d.  These partition the grid across
processors and have methods which allow forward and reverse
communication of ghost grid data as well as load balancing.  If you
write a new compute or fix which needs a distributed grid, these are
the classes to look at.  A new pair style could use a distributed
grid by having a fix define it.  Please see the section on
{\hyperref[\detokenize{Developer_grid::doc}]{\sphinxcrossref{\DUrole{doc}{using distributed grids within style classes}}}}
for a detailed description.
\end{sphinxadmonition}


\bigskip\hrule\bigskip


\sphinxAtStartPar
These are the commands which currently define or use distributed
grids:
\begin{itemize}
\item {} 
\sphinxAtStartPar
{\hyperref[\detokenize{fix_ttm::doc}]{\sphinxcrossref{\DUrole{doc}{fix ttm/grid}}}} \sphinxhyphen{} store electron temperature on grid

\item {} 
\sphinxAtStartPar
{\hyperref[\detokenize{fix_ave_grid::doc}]{\sphinxcrossref{\DUrole{doc}{fix ave/grid}}}} \sphinxhyphen{} time average per\sphinxhyphen{}atom or per\sphinxhyphen{}grid values

\item {} 
\sphinxAtStartPar
{\hyperref[\detokenize{compute_property_grid::doc}]{\sphinxcrossref{\DUrole{doc}{compute property/grid}}}} \sphinxhyphen{} generate grid IDs and coords

\item {} 
\sphinxAtStartPar
{\hyperref[\detokenize{dump::doc}]{\sphinxcrossref{\DUrole{doc}{dump grid}}}} \sphinxhyphen{} output per\sphinxhyphen{}grid values in LAMMPS format

\item {} 
\sphinxAtStartPar
{\hyperref[\detokenize{dump::doc}]{\sphinxcrossref{\DUrole{doc}{dump grid/vtk}}}} \sphinxhyphen{} output per\sphinxhyphen{}grid values in VTK format

\item {} 
\sphinxAtStartPar
{\hyperref[\detokenize{dump_image::doc}]{\sphinxcrossref{\DUrole{doc}{dump image grid}}}} \sphinxhyphen{} include colored grid in output images

\item {} 
\sphinxAtStartPar
{\hyperref[\detokenize{pair_amoeba::doc}]{\sphinxcrossref{\DUrole{doc}{pair\_style amoeba}}}} \sphinxhyphen{} FFT grids

\item {} 
\sphinxAtStartPar
{\hyperref[\detokenize{kspace_style::doc}]{\sphinxcrossref{\DUrole{doc}{kspace\_style pppm}}}} (and variants) \sphinxhyphen{} FFT grids

\item {} 
\sphinxAtStartPar
{\hyperref[\detokenize{kspace_style::doc}]{\sphinxcrossref{\DUrole{doc}{kspace\_style msm}}}} (and variants) \sphinxhyphen{} MSM grids

\end{itemize}

\sphinxAtStartPar
The grids used by the {\hyperref[\detokenize{kspace_style::doc}]{\sphinxcrossref{\DUrole{doc}{kspace\_style}}}} can not be
referenced by an input script.  However the grids and data created and
used by the other commands can be.

\sphinxAtStartPar
A compute or fix command may create one or more grids (of different
sizes).  Each grid can store one or more data fields.  A data field
can be a single value per grid point (per\sphinxhyphen{}grid vector) or multiple
values per grid point (per\sphinxhyphen{}grid array).  See the {\hyperref[\detokenize{Howto_output::doc}]{\sphinxcrossref{\DUrole{doc}{Howto output}}}} doc page for an explanation of how per\sphinxhyphen{}grid data can
be generated by some commands and used by other commands.

\sphinxAtStartPar
A command accesses grid data from a compute or fix using a \sphinxstyleemphasis{grid
reference} with the following syntax:
\begin{itemize}
\item {} 
\sphinxAtStartPar
c\_ID:gname:dname

\item {} 
\sphinxAtStartPar
c\_ID:gname:dname{[}I{]}

\item {} 
\sphinxAtStartPar
f\_ID:gname:dname

\item {} 
\sphinxAtStartPar
f\_ID:gname:dname{[}I{]}

\end{itemize}

\sphinxAtStartPar
The prefix “c\_” or “f\_” refers to the ID of the compute or fix; gname is
the name of the grid, which is assigned by the compute or fix; dname is
the name of the data field, which is also assigned by the compute or
fix.

\sphinxAtStartPar
If the data field is a per\sphinxhyphen{}grid vector (one value per grid point),
then no brackets are used to access the values.  If the data field is
a per\sphinxhyphen{}grid array (multiple values per grid point), then brackets are
used to specify the column I of the array.  I ranges from 1 to Ncol
inclusive, where Ncol is the number of columns in the array and is
defined by the compute or fix.

\sphinxAtStartPar
Currently, there are no per\sphinxhyphen{}grid variables implemented in LAMMPS.  We
may add this feature at some point.

\sphinxstepscope


\subsection{Calculate temperature}
\label{\detokenize{Howto_temperature:calculate-temperature}}\label{\detokenize{Howto_temperature::doc}}
\sphinxAtStartPar
Temperature is computed as kinetic energy divided by some number of
degrees of freedom (and the Boltzmann constant).  Since kinetic energy
is a function of particle velocity, there is often a need to
distinguish between a particle’s advection velocity (due to some
aggregate motion of particles) and its thermal velocity.  The sum of
the two is the particle’s total velocity, but the latter is often what
is wanted to compute a temperature.

\sphinxAtStartPar
LAMMPS has several options for computing temperatures, any of which
can be used in {\hyperref[\detokenize{Howto_thermostat::doc}]{\sphinxcrossref{\DUrole{doc}{thermostatting}}}} and
{\hyperref[\detokenize{Howto_barostat::doc}]{\sphinxcrossref{\DUrole{doc}{barostatting}}}}.  These {\hyperref[\detokenize{compute::doc}]{\sphinxcrossref{\DUrole{doc}{compute commands}}}} calculate temperature:
\begin{itemize}
\item {} 
\sphinxAtStartPar
{\hyperref[\detokenize{compute_temp::doc}]{\sphinxcrossref{\DUrole{doc}{compute temp}}}}

\item {} 
\sphinxAtStartPar
{\hyperref[\detokenize{compute_temp_sphere::doc}]{\sphinxcrossref{\DUrole{doc}{compute temp/sphere}}}}

\item {} 
\sphinxAtStartPar
{\hyperref[\detokenize{compute_temp_asphere::doc}]{\sphinxcrossref{\DUrole{doc}{compute temp/asphere}}}}

\item {} 
\sphinxAtStartPar
{\hyperref[\detokenize{compute_temp_com::doc}]{\sphinxcrossref{\DUrole{doc}{compute temp/com}}}}

\item {} 
\sphinxAtStartPar
{\hyperref[\detokenize{compute_temp_deform::doc}]{\sphinxcrossref{\DUrole{doc}{compute temp/deform}}}}

\item {} 
\sphinxAtStartPar
{\hyperref[\detokenize{compute_temp_partial::doc}]{\sphinxcrossref{\DUrole{doc}{compute temp/partial}}}}

\item {} 
\sphinxAtStartPar
{\hyperref[\detokenize{compute_temp_profile::doc}]{\sphinxcrossref{\DUrole{doc}{compute temp/profile}}}}

\item {} 
\sphinxAtStartPar
{\hyperref[\detokenize{compute_temp_ramp::doc}]{\sphinxcrossref{\DUrole{doc}{compute temp/ramp}}}}

\item {} 
\sphinxAtStartPar
{\hyperref[\detokenize{compute_temp_region::doc}]{\sphinxcrossref{\DUrole{doc}{compute temp/region}}}}

\end{itemize}

\sphinxAtStartPar
All but the first 3 calculate velocity biases directly (e.g. advection
velocities) that are removed when computing the thermal temperature.
{\hyperref[\detokenize{compute_temp_sphere::doc}]{\sphinxcrossref{\DUrole{doc}{Compute temp/sphere}}}} and {\hyperref[\detokenize{compute_temp_asphere::doc}]{\sphinxcrossref{\DUrole{doc}{compute temp/asphere}}}} compute kinetic energy for
finite\sphinxhyphen{}size particles that includes rotational degrees of freedom.
They both allow for velocity biases indirectly, via an optional extra
argument which is another temperature compute that subtracts a
velocity bias.  This allows the translational velocity of spherical or
aspherical particles to be adjusted in prescribed ways.

\sphinxstepscope


\subsection{Calculate elastic constants}
\label{\detokenize{Howto_elastic:calculate-elastic-constants}}\label{\detokenize{Howto_elastic::doc}}
\sphinxAtStartPar
Elastic constants characterize the stiffness of a material. The formal
definition is provided by the linear relation that holds between the
stress and strain tensors in the limit of infinitesimal deformation.
In tensor notation, this is expressed as
\begin{equation*}
\begin{split}s_{ij} = C_{ijkl} e_{kl}\end{split}
\end{equation*}
\sphinxAtStartPar
where
the repeated indices imply summation. \(s_{ij}\) are the elements of the
symmetric stress tensor. \(e_{kl}\) are the elements of the symmetric strain
tensor. \(C_{ijkl}\) are the elements of the fourth rank tensor of elastic
constants. In three dimensions, this tensor has \(3^4=81\) elements. Using
Voigt notation, the tensor can be written as a 6x6 matrix, where \(C_{ij}\)
is now the derivative of \(s_i\) w.r.t. \(e_j\). Because \(s_i\) is itself a
derivative w.r.t. \(e_i\), it follows that \(C_{ij}\) is also symmetric, with at
most \(\frac{7 \times 6}{2}\) = 21 distinct elements.

\sphinxAtStartPar
At zero temperature, it is easy to estimate these derivatives by
deforming the simulation box in one of the six directions using the
{\hyperref[\detokenize{change_box::doc}]{\sphinxcrossref{\DUrole{doc}{change\_box}}}} command and measuring the change in the
stress tensor. A general\sphinxhyphen{}purpose script that does this is given in the
\sphinxcode{\sphinxupquote{examples/ELASTIC}} directory described on the {\hyperref[\detokenize{Examples::doc}]{\sphinxcrossref{\DUrole{doc}{Examples}}}}
doc page.

\sphinxAtStartPar
Calculating elastic constants at finite temperature is more
challenging, because it is necessary to run a simulation that performs
time averages of differential properties. There are at least
3 ways to do this in LAMMPS. The most reliable way to do this is
by exploiting the relationship between elastic constants, stress
fluctuations, and the Born matrix, the second derivatives of energy
w.r.t. strain {\hyperref[\detokenize{Howto_elastic:ray}]{\sphinxcrossref{\DUrole{std,std-ref}{(Ray)}}}}.
The Born matrix calculation has been enabled by
the {\hyperref[\detokenize{compute_born_matrix::doc}]{\sphinxcrossref{\DUrole{doc}{compute born/matrix}}}} command,
which works for any bonded or non\sphinxhyphen{}bonded potential in LAMMPS.
The most expensive part of the calculation is the sampling of
the stress fluctuations. Several examples of this method are
provided in the \sphinxcode{\sphinxupquote{examples/ELASTIC\_T/BORN\_MATRIX}} directory
described on the {\hyperref[\detokenize{Examples::doc}]{\sphinxcrossref{\DUrole{doc}{Examples}}}} doc page.

\sphinxAtStartPar
A second way is to measure
the change in average stress tensor in an NVT simulations when
the cell volume undergoes a finite deformation. In order to balance
the systematic and statistical errors in this method, the magnitude of
the deformation must be chosen judiciously, and care must be taken to
fully equilibrate the deformed cell before sampling the stress
tensor. An example of this method is provided in the
\sphinxcode{\sphinxupquote{examples/ELASTIC\_T/DEFORMATION}} directory
described on the {\hyperref[\detokenize{Examples::doc}]{\sphinxcrossref{\DUrole{doc}{Examples}}}} doc page.

\sphinxAtStartPar
Another approach is to sample the triclinic cell fluctuations
that occur in an NPT simulation. This method can also be slow to
converge and requires careful post\sphinxhyphen{}processing {\hyperref[\detokenize{Howto_elastic:shinoda1}]{\sphinxcrossref{\DUrole{std,std-ref}{(Shinoda)}}}}.
We do not provide an example of this method.

\sphinxAtStartPar
A nice review of the advantages and disadvantages of all of these methods
is provided in the paper by Clavier et al. {\hyperref[\detokenize{Howto_elastic:clavier}]{\sphinxcrossref{\DUrole{std,std-ref}{(Clavier)}}}}.


\bigskip\hrule\bigskip

\phantomsection\label{\detokenize{Howto_elastic:ray}}
\sphinxAtStartPar
\sphinxstylestrong{(Ray)} J. R. Ray and A. Rahman, J Chem Phys, 80, 4423 (1984).

\phantomsection\label{\detokenize{Howto_elastic:shinoda1}}
\sphinxAtStartPar
\sphinxstylestrong{(Shinoda)} Shinoda, Shiga, and Mikami, Phys Rev B, 69, 134103 (2004).

\phantomsection\label{\detokenize{Howto_elastic:clavier}}
\sphinxAtStartPar
\sphinxstylestrong{(Clavier)} G. Clavier, N. Desbiens, E. Bourasseau, V. Lachet, N. Brusselle\sphinxhyphen{}Dupend and B. Rousseau, Mol Sim, 43, 1413 (2017).

\sphinxstepscope


\subsection{Calculate thermal conductivity}
\label{\detokenize{Howto_kappa:calculate-thermal-conductivity}}\label{\detokenize{Howto_kappa::doc}}
\sphinxAtStartPar
The thermal conductivity \(\kappa\) of a material can be measured in at
least 4 ways using various options in LAMMPS.  See the \sphinxcode{\sphinxupquote{examples/KAPPA}}
directory for scripts that implement the 4 methods discussed here for
a simple Lennard\sphinxhyphen{}Jones fluid model.  Also, see the {\hyperref[\detokenize{Howto_viscosity::doc}]{\sphinxcrossref{\DUrole{doc}{Howto viscosity}}}} page for an analogous discussion
for viscosity.

\sphinxAtStartPar
The thermal conductivity tensor \(\mathbf{\kappa}\) is a measure of the propensity
of a material to transmit heat energy in a diffusive manner as given
by Fourier’s law
\begin{equation*}
\begin{split}J = -\kappa \cdot \text{grad}(T)\end{split}
\end{equation*}
\sphinxAtStartPar
where \(J\) is the heat flux in units of energy per area per time and
\(\text{grad}(T)\) is the spatial gradient of temperature.  The thermal
conductivity thus has units of energy per distance per time per degree
K and is often approximated as an isotropic quantity, i.e. as a
scalar.

\sphinxAtStartPar
The first method is to setup two thermostatted regions at opposite
ends of a simulation box, or one in the middle and one at the end of a
periodic box.  By holding the two regions at different temperatures
with a {\hyperref[\detokenize{Howto_thermostat::doc}]{\sphinxcrossref{\DUrole{doc}{thermostatting fix}}}}, the energy added to
the hot region should equal the energy subtracted from the cold region
and be proportional to the heat flux moving between the regions.  See
the papers by {\hyperref[\detokenize{Howto_kappa:howto-ikeshoji}]{\sphinxcrossref{\DUrole{std,std-ref}{Ikeshoji and Hafskjold}}}} and
{\hyperref[\detokenize{Howto_kappa:howto-wirnsberger}]{\sphinxcrossref{\DUrole{std,std-ref}{Wirnsberger et al}}}} for details of this idea.  Note
that thermostatting fixes such as {\hyperref[\detokenize{fix_nh::doc}]{\sphinxcrossref{\DUrole{doc}{fix nvt}}}}, {\hyperref[\detokenize{fix_langevin::doc}]{\sphinxcrossref{\DUrole{doc}{fix langevin}}}}, and {\hyperref[\detokenize{fix_temp_rescale::doc}]{\sphinxcrossref{\DUrole{doc}{fix temp/rescale}}}} store the cumulative energy they
add/subtract.

\sphinxAtStartPar
Alternatively, as a second method, the {\hyperref[\detokenize{fix_heat::doc}]{\sphinxcrossref{\DUrole{doc}{fix heat}}}} or
{\hyperref[\detokenize{fix_ehex::doc}]{\sphinxcrossref{\DUrole{doc}{fix ehex}}}} commands can be used in place of thermostats
on each of two regions to add/subtract specified amounts of energy to
both regions.  In both cases, the resulting temperatures of the two
regions can be monitored with the “compute temp/region” command and
the temperature profile of the intermediate region can be monitored
with the {\hyperref[\detokenize{fix_ave_chunk::doc}]{\sphinxcrossref{\DUrole{doc}{fix ave/chunk}}}} and {\hyperref[\detokenize{compute_ke_atom::doc}]{\sphinxcrossref{\DUrole{doc}{compute ke/atom}}}} commands.

\sphinxAtStartPar
The third method is to perform a reverse non\sphinxhyphen{}equilibrium MD simulation
using the {\hyperref[\detokenize{fix_thermal_conductivity::doc}]{\sphinxcrossref{\DUrole{doc}{fix thermal/conductivity}}}}
command which implements the rNEMD algorithm of Muller\sphinxhyphen{}Plathe.
Kinetic energy is swapped between atoms in two different layers of the
simulation box.  This induces a temperature gradient between the two
layers which can be monitored with the {\hyperref[\detokenize{fix_ave_chunk::doc}]{\sphinxcrossref{\DUrole{doc}{fix ave/chunk}}}} and {\hyperref[\detokenize{compute_ke_atom::doc}]{\sphinxcrossref{\DUrole{doc}{compute ke/atom}}}} commands.  The fix tallies the
cumulative energy transfer that it performs.  See the {\hyperref[\detokenize{fix_thermal_conductivity::doc}]{\sphinxcrossref{\DUrole{doc}{fix thermal/conductivity}}}} command for
details.

\sphinxAtStartPar
The fourth method is based on the Green\sphinxhyphen{}Kubo (GK) formula which
relates the ensemble average of the auto\sphinxhyphen{}correlation of the heat flux
to \(\kappa\).  The heat flux can be calculated from the fluctuations of
per\sphinxhyphen{}atom potential and kinetic energies and per\sphinxhyphen{}atom stress tensor in
a steady\sphinxhyphen{}state equilibrated simulation.  This is in contrast to the
two preceding non\sphinxhyphen{}equilibrium methods, where energy flows continuously
between hot and cold regions of the simulation box.

\sphinxAtStartPar
The {\hyperref[\detokenize{compute_heat_flux::doc}]{\sphinxcrossref{\DUrole{doc}{compute heat/flux}}}} command can calculate
the needed heat flux and describes how to implement the Green\_Kubo
formalism using additional LAMMPS commands, such as the {\hyperref[\detokenize{fix_ave_correlate::doc}]{\sphinxcrossref{\DUrole{doc}{fix ave/correlate}}}} command to calculate the needed
auto\sphinxhyphen{}correlation.  See the page for the {\hyperref[\detokenize{compute_heat_flux::doc}]{\sphinxcrossref{\DUrole{doc}{compute heat/flux}}}} command for an example input script
that calculates the thermal conductivity of solid Ar via the GK
formalism.


\bigskip\hrule\bigskip

\phantomsection\label{\detokenize{Howto_kappa:howto-ikeshoji}}
\sphinxAtStartPar
\sphinxstylestrong{(Ikeshoji)} Ikeshoji and Hafskjold, Molecular Physics, 81, 251\sphinxhyphen{}261
(1994).

\phantomsection\label{\detokenize{Howto_kappa:howto-wirnsberger}}
\sphinxAtStartPar
\sphinxstylestrong{(Wirnsberger)} Wirnsberger, Frenkel, and Dellago, J Chem Phys, 143, 124104
(2015).

\sphinxstepscope


\subsection{Calculate viscosity}
\label{\detokenize{Howto_viscosity:calculate-viscosity}}\label{\detokenize{Howto_viscosity::doc}}
\sphinxAtStartPar
The shear viscosity \(\eta\) of a fluid can be measured in at least 6 ways
using various options in LAMMPS.  See the \sphinxcode{\sphinxupquote{examples/VISCOSITY}} directory
for scripts that implement the 5 methods discussed here for a simple
Lennard\sphinxhyphen{}Jones fluid model and 1 method for SPC/E water model.
Also, see the {\hyperref[\detokenize{Howto_kappa::doc}]{\sphinxcrossref{\DUrole{doc}{page on calculating thermal conductivity}}}}
for an analogous discussion for thermal conductivity.

\sphinxAtStartPar
\(\eta\) is a measure of the propensity of a fluid to transmit momentum in
a direction perpendicular to the direction of velocity or momentum
flow.  Alternatively it is the resistance the fluid has to being
sheared.  It is given by
\begin{equation*}
\begin{split}J = -\eta \cdot \text{grad}(V_{\text{stream}})\end{split}
\end{equation*}
\sphinxAtStartPar
where \(J\) is the momentum flux in units of momentum per area per time.
and \(\text{grad}(V_{\text{stream}})\) is the spatial gradient of the velocity of the fluid
moving in another direction, normal to the area through which the
momentum flows.  Viscosity thus has units of pressure\sphinxhyphen{}time.

\sphinxAtStartPar
The first method is to perform a non\sphinxhyphen{}equilibrium MD (NEMD) simulation
by shearing the simulation box via the {\hyperref[\detokenize{fix_deform::doc}]{\sphinxcrossref{\DUrole{doc}{fix deform}}}}
command, and using the {\hyperref[\detokenize{fix_nvt_sllod::doc}]{\sphinxcrossref{\DUrole{doc}{fix nvt/sllod}}}} command to
thermostat the fluid via the SLLOD equations of motion.
Alternatively, as a second method, one or more moving walls can be
used to shear the fluid in between them, again with some kind of
thermostat that modifies only the thermal (non\sphinxhyphen{}shearing) components of
velocity to prevent the fluid from heating up.

\begin{sphinxadmonition}{note}{Note:}
\sphinxAtStartPar
A recent (2017) book by {\hyperref[\detokenize{Howto_viscosity:daivis-viscosity}]{\sphinxcrossref{\DUrole{std,std-ref}{(Daivis and Todd)}}}}
discusses use of the SLLOD method and non\sphinxhyphen{}equilibrium MD (NEMD)
thermostatting generally, for both simple and complex fluids,
e.g. molecular systems.  The latter can be tricky to do correctly.
\end{sphinxadmonition}

\sphinxAtStartPar
In both cases, the velocity profile setup in the fluid by this
procedure can be monitored by the {\hyperref[\detokenize{fix_ave_chunk::doc}]{\sphinxcrossref{\DUrole{doc}{fix ave/chunk}}}}
command, which determines \(\text{grad}(V_{\text{stream}})\) in the equation above.
E.g. the derivative in the y\sphinxhyphen{}direction of the \(V_x\) component of fluid
motion or \(\text{grad}(V_{\text{stream}}) = \frac{\text{d} V_x}{\text{d} y}\).  The \(P_{xy}\) off\sphinxhyphen{}diagonal component of
the pressure or stress tensor, as calculated by the {\hyperref[\detokenize{compute_pressure::doc}]{\sphinxcrossref{\DUrole{doc}{compute pressure}}}} command, can also be monitored, which
is the \(J\) term in the equation above.  See the {\hyperref[\detokenize{Howto_nemd::doc}]{\sphinxcrossref{\DUrole{doc}{Howto nemd}}}} page for details on NEMD simulations.

\sphinxAtStartPar
The third method is to perform a reverse non\sphinxhyphen{}equilibrium MD simulation
using the {\hyperref[\detokenize{fix_viscosity::doc}]{\sphinxcrossref{\DUrole{doc}{fix viscosity}}}} command which implements
the rNEMD algorithm of Muller\sphinxhyphen{}Plathe.  Momentum in one dimension is
swapped between atoms in two different layers of the simulation box in
a different dimension.  This induces a velocity gradient which can be
monitored with the {\hyperref[\detokenize{fix_ave_chunk::doc}]{\sphinxcrossref{\DUrole{doc}{fix ave/chunk}}}} command.
The fix tallies the cumulative momentum transfer that it performs.
See the {\hyperref[\detokenize{fix_viscosity::doc}]{\sphinxcrossref{\DUrole{doc}{fix viscosity}}}} command for details.

\sphinxAtStartPar
The fourth method is based on the Green\sphinxhyphen{}Kubo (GK) formula which
relates the ensemble average of the auto\sphinxhyphen{}correlation of the
stress/pressure tensor to \(\eta\).  This can be done in a fully
equilibrated simulation which is in contrast to the two preceding
non\sphinxhyphen{}equilibrium methods, where momentum flows continuously through the
simulation box.

\sphinxAtStartPar
Here is an example input script that calculates the viscosity of
liquid Ar via the GK formalism:

\begin{sphinxVerbatim}[commandchars=\\\{\}]
\PYG{c}{\PYGZsh{} Sample LAMMPS input script for viscosity of liquid Ar}

\PYG{k}{units}\PYG{+w}{       }\PYG{n}{real}
\PYG{k}{variable    }\PYG{n+nv+nvIdentifier}{T}\PYG{+w}{ }\PYG{n}{equal}\PYG{+w}{ }\PYG{l+m}{200.0}\PYG{+w}{       }\PYG{c}{\PYGZsh{} run temperature}
\PYG{k}{variable    }\PYG{n+nv+nvIdentifier}{Tinit}\PYG{+w}{ }\PYG{n}{equal}\PYG{+w}{ }\PYG{l+m}{250.0}\PYG{+w}{   }\PYG{c}{\PYGZsh{} equilibration temperature}
\PYG{k}{variable    }\PYG{n+nv+nvIdentifier}{V}\PYG{+w}{ }\PYG{n}{equal}\PYG{+w}{ }\PYG{n}{vol}
\PYG{k}{variable    }\PYG{n+nv+nvIdentifier}{dt}\PYG{+w}{ }\PYG{n}{equal}\PYG{+w}{ }\PYG{l+m}{4.0}
\PYG{k}{variable    }\PYG{n+nv+nvIdentifier}{p}\PYG{+w}{ }\PYG{n}{equal}\PYG{+w}{ }\PYG{l+m}{400}\PYG{+w}{     }\PYG{c}{\PYGZsh{} correlation length}
\PYG{k}{variable    }\PYG{n+nv+nvIdentifier}{s}\PYG{+w}{ }\PYG{n}{equal}\PYG{+w}{ }\PYG{l+m}{5}\PYG{+w}{       }\PYG{c}{\PYGZsh{} sample interval}
\PYG{k}{variable    }\PYG{n+nv+nvIdentifier}{d}\PYG{+w}{ }\PYG{n}{equal}\PYG{+w}{ }\PYG{n+nv}{\PYGZdl{}p}\PYG{o}{*}\PYG{n+nv}{\PYGZdl{}s}\PYG{+w}{   }\PYG{c}{\PYGZsh{} dump interval}

\PYG{c}{\PYGZsh{} convert from LAMMPS real units to SI}

\PYG{k}{variable    }\PYG{n+nv+nvIdentifier}{kB}\PYG{+w}{ }\PYG{n}{equal}\PYG{+w}{ }\PYG{l+m}{1.3806504e\PYGZhy{}23}\PYG{+w}{    }\PYG{c}{\PYGZsh{} [J/K] Boltzmann}
\PYG{k}{variable    }\PYG{n+nv+nvIdentifier}{atm2Pa}\PYG{+w}{ }\PYG{n}{equal}\PYG{+w}{ }\PYG{l+m}{101325.0}
\PYG{k}{variable    }\PYG{n+nv+nvIdentifier}{A2m}\PYG{+w}{ }\PYG{n}{equal}\PYG{+w}{ }\PYG{l+m}{1.0e\PYGZhy{}10}
\PYG{k}{variable    }\PYG{n+nv+nvIdentifier}{fs2s}\PYG{+w}{ }\PYG{n}{equal}\PYG{+w}{ }\PYG{l+m}{1.0e\PYGZhy{}15}
\PYG{k}{variable    }\PYG{n+nv+nvIdentifier}{convert}\PYG{+w}{ }\PYG{n}{equal}\PYG{+w}{ }\PYG{n+nv}{\PYGZdl{}\PYGZob{}}\PYG{n+nv}{atm2Pa}\PYG{n+nv}{\PYGZcb{}}\PYG{o}{*}\PYG{n+nv}{\PYGZdl{}\PYGZob{}}\PYG{n+nv}{atm2Pa}\PYG{n+nv}{\PYGZcb{}}\PYG{o}{*}\PYG{n+nv}{\PYGZdl{}\PYGZob{}}\PYG{n+nv}{fs2s}\PYG{n+nv}{\PYGZcb{}}\PYG{o}{*}\PYG{n+nv}{\PYGZdl{}\PYGZob{}}\PYG{n+nv}{A2m}\PYG{n+nv}{\PYGZcb{}}\PYG{o}{*}\PYG{n+nv}{\PYGZdl{}\PYGZob{}}\PYG{n+nv}{A2m}\PYG{n+nv}{\PYGZcb{}}\PYG{o}{*}\PYG{n+nv}{\PYGZdl{}\PYGZob{}}\PYG{n+nv}{A2m}\PYG{n+nv}{\PYGZcb{}}

\PYG{c}{\PYGZsh{} setup problem}

\PYG{k}{dimension}\PYG{+w}{    }\PYG{l+m}{3}
\PYG{k}{boundary}\PYG{+w}{     }\PYG{n}{p}\PYG{+w}{ }\PYG{n}{p}\PYG{+w}{ }\PYG{n}{p}
\PYG{k}{lattice}\PYG{+w}{      }\PYG{n}{fcc}\PYG{+w}{ }\PYG{l+m}{5.376}\PYG{+w}{ }\PYG{n}{orient}\PYG{+w}{ }\PYG{n}{x}\PYG{+w}{ }\PYG{l+m}{1}\PYG{+w}{ }\PYG{l+m}{0}\PYG{+w}{ }\PYG{l+m}{0}\PYG{+w}{ }\PYG{n}{orient}\PYG{+w}{ }\PYG{n}{y}\PYG{+w}{ }\PYG{l+m}{0}\PYG{+w}{ }\PYG{l+m}{1}\PYG{+w}{ }\PYG{l+m}{0}\PYG{+w}{ }\PYG{n}{orient}\PYG{+w}{ }\PYG{n}{z}\PYG{+w}{ }\PYG{l+m}{0}\PYG{+w}{ }\PYG{l+m}{0}\PYG{+w}{ }\PYG{l+m}{1}
\PYG{k}{region       }\PYG{n+nv+nvIdentifier}{box}\PYG{+w}{ }\PYG{n}{block}\PYG{+w}{ }\PYG{l+m}{0}\PYG{+w}{ }\PYG{l+m}{4}\PYG{+w}{ }\PYG{l+m}{0}\PYG{+w}{ }\PYG{l+m}{4}\PYG{+w}{ }\PYG{l+m}{0}\PYG{+w}{ }\PYG{l+m}{4}
\PYG{k}{create\PYGZus{}box   }\PYG{n+nv+nvIdentifier}{1}\PYG{+w}{ }\PYG{n+nv+nvIdentifier}{box}
\PYG{k}{create\PYGZus{}atoms}\PYG{+w}{ }\PYG{l+m}{1}\PYG{+w}{ }\PYG{n}{box}
\PYG{k}{mass}\PYG{+w}{         }\PYG{l+m}{1}\PYG{+w}{ }\PYG{l+m}{39.948}
\PYG{k}{pair\PYGZus{}style}\PYG{+w}{   }\PYG{n}{lj}\PYG{o}{/}\PYG{n}{cut}\PYG{+w}{ }\PYG{l+m}{13.0}
\PYG{k}{pair\PYGZus{}coeff}\PYG{+w}{   }\PYG{o}{*}\PYG{+w}{ }\PYG{o}{*}\PYG{+w}{ }\PYG{l+m}{0.2381}\PYG{+w}{ }\PYG{l+m}{3.405}
\PYG{k}{timestep}\PYG{+w}{     }\PYG{n+nv}{\PYGZdl{}\PYGZob{}}\PYG{n+nv}{dt}\PYG{n+nv}{\PYGZcb{}}
\PYG{k}{thermo}\PYG{+w}{       }\PYG{n+nv}{\PYGZdl{}d}

\PYG{c}{\PYGZsh{} equilibration and thermalization}

\PYG{k}{velocity}\PYG{+w}{     }\PYG{n}{all}\PYG{+w}{ }\PYG{n}{create}\PYG{+w}{ }\PYG{n+nv}{\PYGZdl{}\PYGZob{}}\PYG{n+nv}{Tinit}\PYG{n+nv}{\PYGZcb{}}\PYG{+w}{ }\PYG{l+m}{102486}\PYG{+w}{ }\PYG{n}{mom}\PYG{+w}{ }\PYG{n}{yes}\PYG{+w}{ }\PYG{n}{rot}\PYG{+w}{ }\PYG{n}{yes}\PYG{+w}{ }\PYG{n}{dist}\PYG{+w}{ }\PYG{n}{gaussian}
\PYG{k}{fix          }\PYG{n+nv+nvIdentifier}{NVT}\PYG{+w}{ }\PYG{n+nv+nvIdentifier}{all}\PYG{+w}{ }\PYG{n}{nvt}\PYG{+w}{ }\PYG{n}{temp}\PYG{+w}{ }\PYG{n+nv}{\PYGZdl{}\PYGZob{}}\PYG{n+nv}{Tinit}\PYG{n+nv}{\PYGZcb{}}\PYG{+w}{ }\PYG{n+nv}{\PYGZdl{}\PYGZob{}}\PYG{n+nv}{Tinit}\PYG{n+nv}{\PYGZcb{}}\PYG{+w}{ }\PYG{l+m}{10}\PYG{+w}{ }\PYG{n}{drag}\PYG{+w}{ }\PYG{l+m}{0.2}
\PYG{k}{run}\PYG{+w}{          }\PYG{l+m}{8000}

\PYG{c}{\PYGZsh{} viscosity calculation, switch to NVE if desired}

\PYG{k}{velocity}\PYG{+w}{     }\PYG{n}{all}\PYG{+w}{ }\PYG{n}{create}\PYG{+w}{ }\PYG{n+nv}{\PYGZdl{}T}\PYG{+w}{ }\PYG{l+m}{102486}\PYG{+w}{ }\PYG{n}{mom}\PYG{+w}{ }\PYG{n}{yes}\PYG{+w}{ }\PYG{n}{rot}\PYG{+w}{ }\PYG{n}{yes}\PYG{+w}{ }\PYG{n}{dist}\PYG{+w}{ }\PYG{n}{gaussian}
\PYG{k}{fix          }\PYG{n+nv+nvIdentifier}{NVT}\PYG{+w}{ }\PYG{n+nv+nvIdentifier}{all}\PYG{+w}{ }\PYG{n}{nvt}\PYG{+w}{ }\PYG{n}{temp}\PYG{+w}{ }\PYG{n+nv}{\PYGZdl{}T}\PYG{+w}{ }\PYG{n+nv}{\PYGZdl{}T}\PYG{+w}{ }\PYG{l+m}{10}\PYG{+w}{ }\PYG{n}{drag}\PYG{+w}{ }\PYG{l+m}{0.2}
\PYG{c}{\PYGZsh{}unfix       NVT}
\PYG{c}{\PYGZsh{}fix         NVE all nve}

\PYG{k}{reset\PYGZus{}timestep}\PYG{+w}{ }\PYG{l+m}{0}
\PYG{k}{variable     }\PYG{n+nv+nvIdentifier}{pxy}\PYG{+w}{ }\PYG{n}{equal}\PYG{+w}{ }\PYG{n}{pxy}
\PYG{k}{variable     }\PYG{n+nv+nvIdentifier}{pxz}\PYG{+w}{ }\PYG{n}{equal}\PYG{+w}{ }\PYG{n}{pxz}
\PYG{k}{variable     }\PYG{n+nv+nvIdentifier}{pyz}\PYG{+w}{ }\PYG{n}{equal}\PYG{+w}{ }\PYG{n}{pyz}
\PYG{k}{fix          }\PYG{n+nv+nvIdentifier}{SS}\PYG{+w}{ }\PYG{n+nv+nvIdentifier}{all}\PYG{+w}{ }\PYG{n}{ave}\PYG{o}{/}\PYG{n}{correlate}\PYG{+w}{ }\PYG{n+nv}{\PYGZdl{}s}\PYG{+w}{ }\PYG{n+nv}{\PYGZdl{}p}\PYG{+w}{ }\PYG{n+nv}{\PYGZdl{}d}\PYG{+w}{ }\PYG{o}{\PYGZam{}}
\PYG{+w}{             }\PYG{n}{v\PYGZus{}pxy}\PYG{+w}{ }\PYG{n}{v\PYGZus{}pxz}\PYG{+w}{ }\PYG{n}{v\PYGZus{}pyz}\PYG{+w}{ }\PYG{n}{type}\PYG{+w}{ }\PYG{n}{auto}\PYG{+w}{ }\PYG{n}{file}\PYG{+w}{ }\PYG{n}{S0St.dat}\PYG{+w}{ }\PYG{n}{ave}\PYG{+w}{ }\PYG{n}{running}
\PYG{k}{variable     }\PYG{n+nv+nvIdentifier}{scale}\PYG{+w}{ }\PYG{n}{equal}\PYG{+w}{ }\PYG{n+nv}{\PYGZdl{}\PYGZob{}}\PYG{n+nv}{convert}\PYG{n+nv}{\PYGZcb{}}\PYG{o}{/}\PYG{n+nv}{(}\PYG{n+nv}{\PYGZdl{}\PYGZob{}kB\PYGZcb{}*\PYGZdl{}T}\PYG{n+nv}{)}\PYG{o}{*}\PYG{n+nv}{\PYGZdl{}V}\PYG{o}{*}\PYG{n+nv}{\PYGZdl{}s}\PYG{o}{*}\PYG{n+nv}{\PYGZdl{}\PYGZob{}}\PYG{n+nv}{dt}\PYG{n+nv}{\PYGZcb{}}
\PYG{k}{variable     }\PYG{n+nv+nvIdentifier}{v11}\PYG{+w}{ }\PYG{n}{equal}\PYG{+w}{ }\PYG{n}{trap}\PYG{n+nv}{(}\PYG{n+nv}{f\PYGZus{}SS[3]}\PYG{n+nv}{)}\PYG{o}{*}\PYG{n+nv}{\PYGZdl{}\PYGZob{}}\PYG{n+nv}{scale}\PYG{n+nv}{\PYGZcb{}}
\PYG{k}{variable     }\PYG{n+nv+nvIdentifier}{v22}\PYG{+w}{ }\PYG{n}{equal}\PYG{+w}{ }\PYG{n}{trap}\PYG{n+nv}{(}\PYG{n+nv}{f\PYGZus{}SS[4]}\PYG{n+nv}{)}\PYG{o}{*}\PYG{n+nv}{\PYGZdl{}\PYGZob{}}\PYG{n+nv}{scale}\PYG{n+nv}{\PYGZcb{}}
\PYG{k}{variable     }\PYG{n+nv+nvIdentifier}{v33}\PYG{+w}{ }\PYG{n}{equal}\PYG{+w}{ }\PYG{n}{trap}\PYG{n+nv}{(}\PYG{n+nv}{f\PYGZus{}SS[5]}\PYG{n+nv}{)}\PYG{o}{*}\PYG{n+nv}{\PYGZdl{}\PYGZob{}}\PYG{n+nv}{scale}\PYG{n+nv}{\PYGZcb{}}
\PYG{k}{thermo\PYGZus{}style}\PYG{+w}{ }\PYG{n}{custom}\PYG{+w}{ }\PYG{n}{step}\PYG{+w}{ }\PYG{n}{temp}\PYG{+w}{ }\PYG{n}{press}\PYG{+w}{ }\PYG{n}{v\PYGZus{}pxy}\PYG{+w}{ }\PYG{n}{v\PYGZus{}pxz}\PYG{+w}{ }\PYG{n}{v\PYGZus{}pyz}\PYG{+w}{ }\PYG{n}{v\PYGZus{}v11}\PYG{+w}{ }\PYG{n}{v\PYGZus{}v22}\PYG{+w}{ }\PYG{n}{v\PYGZus{}v33}
\PYG{k}{run}\PYG{+w}{          }\PYG{l+m}{100000}
\PYG{k}{variable     }\PYG{n+nv+nvIdentifier}{v}\PYG{+w}{ }\PYG{n}{equal}\PYG{+w}{ }\PYG{n+nv}{(}\PYG{n+nv}{v\PYGZus{}v11+v\PYGZus{}v22+v\PYGZus{}v33}\PYG{n+nv}{)}\PYG{o}{/}\PYG{l+m}{3.0}
\PYG{k}{variable     }\PYG{n+nv+nvIdentifier}{ndens}\PYG{+w}{ }\PYG{n}{equal}\PYG{+w}{ }\PYG{n}{count}\PYG{n+nv}{(}\PYG{n+nv}{all}\PYG{n+nv}{)}\PYG{o}{/}\PYG{n}{vol}
\PYG{k}{print}\PYG{+w}{        }\PYG{l+s}{\PYGZdq{}}\PYG{l+s}{average viscosity: \PYGZdl{}v [Pa.s] @ \PYGZdl{}T K, \PYGZdl{}\PYGZob{}ndens\PYGZcb{} atoms/A\PYGZca{}3}\PYG{l+s}{\PYGZdq{}}
\end{sphinxVerbatim}

\sphinxAtStartPar
The fifth method is related to the above Green\sphinxhyphen{}Kubo method,
but uses the Einstein formulation, analogous to the Einstein
mean\sphinxhyphen{}square\sphinxhyphen{}displacement formulation for self\sphinxhyphen{}diffusivity. The
time\sphinxhyphen{}integrated momentum fluxes play the role of Cartesian
coordinates, whose mean\sphinxhyphen{}square displacement increases linearly
with time at sufficiently long times.

\sphinxAtStartPar
The sixth is the periodic perturbation method, which is also a non\sphinxhyphen{}equilibrium MD method.
However, instead of measuring the momentum flux in response to an applied velocity gradient,
it measures the velocity profile in response to applied stress.
A cosine\sphinxhyphen{}shaped periodic acceleration is added to the system via the
{\hyperref[\detokenize{fix_accelerate_cos::doc}]{\sphinxcrossref{\DUrole{doc}{fix accelerate/cos}}}} command,
and the {\hyperref[\detokenize{compute_viscosity_cos::doc}]{\sphinxcrossref{\DUrole{doc}{compute viscosity/cos}}}} command is used to monitor the
generated velocity profile and remove the velocity bias before thermostatting.

\begin{sphinxadmonition}{note}{Note:}
\sphinxAtStartPar
An article by {\hyperref[\detokenize{Howto_viscosity:hess3}]{\sphinxcrossref{\DUrole{std,std-ref}{(Hess)}}}} discussed the accuracy and efficiency of these methods.
\end{sphinxadmonition}


\bigskip\hrule\bigskip

\phantomsection\label{\detokenize{Howto_viscosity:daivis-viscosity}}
\sphinxAtStartPar
\sphinxstylestrong{(Daivis and Todd)} Daivis and Todd, Nonequilibrium Molecular Dynamics (book),
Cambridge University Press, \sphinxurl{https://doi.org/10.1017/9781139017848}, (2017).

\phantomsection\label{\detokenize{Howto_viscosity:hess3}}
\sphinxAtStartPar
\sphinxstylestrong{(Hess)} Hess, B. The Journal of Chemical Physics 2002, 116 (1), 209\sphinxhyphen{}217.

\sphinxstepscope


\subsection{Calculate diffusion coefficients}
\label{\detokenize{Howto_diffusion:calculate-diffusion-coefficients}}\label{\detokenize{Howto_diffusion::doc}}
\sphinxAtStartPar
The diffusion coefficient \(D\) of a material can be measured in at least
2 ways using various options in LAMMPS.  See the \sphinxcode{\sphinxupquote{examples/DIFFUSE}}
directory for scripts that implement the 2 methods discussed here for
a simple Lennard\sphinxhyphen{}Jones fluid model.

\sphinxAtStartPar
The first method is to measure the mean\sphinxhyphen{}squared displacement (MSD) of
the system, via the {\hyperref[\detokenize{compute_msd::doc}]{\sphinxcrossref{\DUrole{doc}{compute msd}}}} command.  The slope
of the MSD versus time is proportional to the diffusion coefficient.
The instantaneous MSD values can be accumulated in a vector via the
{\hyperref[\detokenize{fix_vector::doc}]{\sphinxcrossref{\DUrole{doc}{fix vector}}}} command, and a line fit to the vector to
compute its slope via the {\hyperref[\detokenize{variable::doc}]{\sphinxcrossref{\DUrole{doc}{variable slope}}}} function, and
thus extract \(D\).

\sphinxAtStartPar
The second method is to measure the velocity auto\sphinxhyphen{}correlation function
(VACF) of the system, via the {\hyperref[\detokenize{compute_vacf::doc}]{\sphinxcrossref{\DUrole{doc}{compute vacf}}}}
command.  The time\sphinxhyphen{}integral of the VACF is proportional to the
diffusion coefficient.  The instantaneous VACF values can be
accumulated in a vector via the {\hyperref[\detokenize{fix_vector::doc}]{\sphinxcrossref{\DUrole{doc}{fix vector}}}} command,
and time integrated via the {\hyperref[\detokenize{variable::doc}]{\sphinxcrossref{\DUrole{doc}{variable trap}}}} function,
and thus extract \(D\).

\sphinxstepscope


\subsection{Output structured data from LAMMPS}
\label{\detokenize{Howto_structured_data:output-structured-data-from-lammps}}\label{\detokenize{Howto_structured_data::doc}}
\sphinxAtStartPar
LAMMPS can output structured data with the {\hyperref[\detokenize{print::doc}]{\sphinxcrossref{\DUrole{doc}{print}}}} and {\hyperref[\detokenize{fix_print::doc}]{\sphinxcrossref{\DUrole{doc}{fix
print}}}} command.  This gives you flexibility since you can build
custom data formats that contain system properties, thermo data, and variables
values. This output can be directed to the screen and/or to a file for post
processing.


\subsubsection{Writing the current system state, thermo data, variable values}
\label{\detokenize{Howto_structured_data:writing-the-current-system-state-thermo-data-variable-values}}
\sphinxAtStartPar
Use the {\hyperref[\detokenize{print::doc}]{\sphinxcrossref{\DUrole{doc}{print}}}} command to output the current system state, which
can include system properties, thermo data and variable values.


\paragraph{YAML}
\label{\detokenize{Howto_structured_data:yaml}}
\begin{sphinxVerbatim}[commandchars=\\\{\}]
\PYG{k}{print}\PYG{+w}{ }\PYG{l+s}{\PYGZdq{}}\PYG{l+s}{\PYGZdq{}}\PYG{l+s}{\PYGZdq{}}\PYG{l+s}{\PYGZhy{}\PYGZhy{}\PYGZhy{}}
\PYG{l+s}{timestep: \PYGZdl{}(step)}
\PYG{l+s}{pe: \PYGZdl{}(pe)}
\PYG{l+s}{ke: \PYGZdl{}(ke)}
\PYG{l+s}{...}\PYG{l+s}{\PYGZdq{}}\PYG{l+s}{\PYGZdq{}}\PYG{l+s}{\PYGZdq{}}\PYG{+w}{ }\PYG{n}{file}\PYG{+w}{ }\PYG{n}{current\PYGZus{}state.yaml}\PYG{+w}{ }\PYG{n}{screen}\PYG{+w}{ }\PYG{n}{no}
\end{sphinxVerbatim}
\sphinxSetupCaptionForVerbatim{current\_state.yaml}
\def\sphinxLiteralBlockLabel{\label{\detokenize{Howto_structured_data:id2}}}
\begin{sphinxVerbatim}[commandchars=\\\{\}]
\PYG{n+nn}{\PYGZhy{}\PYGZhy{}\PYGZhy{}}
\PYG{n+nt}{timestep}\PYG{p}{:}\PYG{+w}{ }\PYG{l+lScalar+lScalarPlain}{250}
\PYG{n+nt}{pe}\PYG{p}{:}\PYG{+w}{ }\PYG{l+lScalar+lScalarPlain}{\PYGZhy{}4.7774327356321810711}
\PYG{n+nt}{ke}\PYG{p}{:}\PYG{+w}{ }\PYG{l+lScalar+lScalarPlain}{2.4962152903997174569}
\end{sphinxVerbatim}


\paragraph{JSON}
\label{\detokenize{Howto_structured_data:json}}
\begin{sphinxVerbatim}[commandchars=\\\{\}]
\PYG{k}{print}\PYG{+w}{ }\PYG{l+s}{\PYGZdq{}}\PYG{l+s}{\PYGZdq{}}\PYG{l+s}{\PYGZdq{}}\PYG{l+s}{\PYGZob{}}
\PYG{l+s}{  }\PYG{l+s}{\PYGZdq{}}\PYG{n}{timestep}\PYG{l+s}{\PYGZdq{}}\PYG{l+s}{: \PYGZdl{}(step),}
\PYG{l+s}{  }\PYG{l+s}{\PYGZdq{}}\PYG{n}{pe}\PYG{l+s}{\PYGZdq{}}\PYG{l+s}{: \PYGZdl{}(pe),}
\PYG{l+s}{  }\PYG{l+s}{\PYGZdq{}}\PYG{n}{ke}\PYG{l+s}{\PYGZdq{}}\PYG{l+s}{: \PYGZdl{}(ke)}
\PYG{l+s}{\PYGZcb{}}\PYG{l+s}{\PYGZdq{}}\PYG{l+s}{\PYGZdq{}}\PYG{l+s}{\PYGZdq{}}\PYG{+w}{ }\PYG{n}{file}\PYG{+w}{ }\PYG{n}{current\PYGZus{}state.json}\PYG{+w}{ }\PYG{n}{screen}\PYG{+w}{ }\PYG{n}{no}
\end{sphinxVerbatim}
\sphinxSetupCaptionForVerbatim{current\_state.json}
\def\sphinxLiteralBlockLabel{\label{\detokenize{Howto_structured_data:id3}}}
\begin{sphinxVerbatim}[commandchars=\\\{\}]
\PYG{p}{\PYGZob{}}
\PYG{+w}{  }\PYG{n+nt}{\PYGZdq{}timestep\PYGZdq{}}\PYG{p}{:}\PYG{+w}{ }\PYG{l+m+mi}{250}\PYG{p}{,}
\PYG{+w}{  }\PYG{n+nt}{\PYGZdq{}pe\PYGZdq{}}\PYG{p}{:}\PYG{+w}{ }\PYG{l+m+mf}{\PYGZhy{}4.7774327356321810711}\PYG{p}{,}
\PYG{+w}{  }\PYG{n+nt}{\PYGZdq{}ke\PYGZdq{}}\PYG{p}{:}\PYG{+w}{ }\PYG{l+m+mf}{2.4962152903997174569}
\PYG{p}{\PYGZcb{}}
\end{sphinxVerbatim}


\subsubsection{YAML format thermo\_style or dump\_style output}
\label{\detokenize{Howto_structured_data:yaml-format-thermo-style-or-dump-style-output}}

\paragraph{Extracting data from log file}
\label{\detokenize{Howto_structured_data:extracting-data-from-log-file}}
\sphinxAtStartPar
\DUrole{versionmodified,added}{Added in version 24Mar2022.}

\sphinxAtStartPar
LAMMPS supports the thermo style “yaml” and for “custom” style
thermodynamic output the format can be changed to YAML with
{\hyperref[\detokenize{thermo_modify::doc}]{\sphinxcrossref{\DUrole{doc}{thermo\_modify line yaml}}}}.  This will produce a
block of output in a compact YAML format \sphinxhyphen{} one “document” per run \sphinxhyphen{} of
the following style:

\begin{sphinxVerbatim}[commandchars=\\\{\}]
\PYG{n+nn}{\PYGZhy{}\PYGZhy{}\PYGZhy{}}
\PYG{n+nt}{keywords}\PYG{p}{:}\PYG{+w}{ }\PYG{p+pIndicator}{[}\PYG{l+s}{\PYGZsq{}}\PYG{l+s}{Step}\PYG{l+s}{\PYGZsq{}}\PYG{p+pIndicator}{,}\PYG{+w}{ }\PYG{l+s}{\PYGZsq{}}\PYG{l+s}{Temp}\PYG{l+s}{\PYGZsq{}}\PYG{p+pIndicator}{,}\PYG{+w}{ }\PYG{l+s}{\PYGZsq{}}\PYG{l+s}{E\PYGZus{}pair}\PYG{l+s}{\PYGZsq{}}\PYG{p+pIndicator}{,}\PYG{+w}{ }\PYG{l+s}{\PYGZsq{}}\PYG{l+s}{E\PYGZus{}mol}\PYG{l+s}{\PYGZsq{}}\PYG{p+pIndicator}{,}\PYG{+w}{ }\PYG{l+s}{\PYGZsq{}}\PYG{l+s}{TotEng}\PYG{l+s}{\PYGZsq{}}\PYG{p+pIndicator}{,}\PYG{+w}{ }\PYG{l+s}{\PYGZsq{}}\PYG{l+s}{Press}\PYG{l+s}{\PYGZsq{}}\PYG{p+pIndicator}{,}\PYG{+w}{ }\PYG{p+pIndicator}{]}
\PYG{n+nt}{data}\PYG{p}{:}
\PYG{+w}{  }\PYG{p+pIndicator}{\PYGZhy{}}\PYG{+w}{ }\PYG{p+pIndicator}{[}\PYG{n+nv}{100}\PYG{p+pIndicator}{,}\PYG{+w}{ }\PYG{n+nv}{0.757453103239935}\PYG{p+pIndicator}{,}\PYG{+w}{ }\PYG{n+nv}{\PYGZhy{}5.7585054860159}\PYG{p+pIndicator}{,}\PYG{+w}{ }\PYG{n+nv}{0}\PYG{p+pIndicator}{,}\PYG{+w}{ }\PYG{n+nv}{\PYGZhy{}4.62236133677021}\PYG{p+pIndicator}{,}\PYG{+w}{ }\PYG{n+nv}{0.207261053624721}\PYG{p+pIndicator}{,}\PYG{+w}{ }\PYG{p+pIndicator}{]}
\PYG{+w}{  }\PYG{p+pIndicator}{\PYGZhy{}}\PYG{+w}{ }\PYG{p+pIndicator}{[}\PYG{n+nv}{110}\PYG{p+pIndicator}{,}\PYG{+w}{ }\PYG{n+nv}{0.759322359337036}\PYG{p+pIndicator}{,}\PYG{+w}{ }\PYG{n+nv}{\PYGZhy{}5.7614668389562}\PYG{p+pIndicator}{,}\PYG{+w}{ }\PYG{n+nv}{0}\PYG{p+pIndicator}{,}\PYG{+w}{ }\PYG{n+nv}{\PYGZhy{}4.62251889318624}\PYG{p+pIndicator}{,}\PYG{+w}{ }\PYG{n+nv}{0.194314975399602}\PYG{p+pIndicator}{,}\PYG{+w}{ }\PYG{p+pIndicator}{]}
\PYG{+w}{  }\PYG{p+pIndicator}{\PYGZhy{}}\PYG{+w}{ }\PYG{p+pIndicator}{[}\PYG{n+nv}{120}\PYG{p+pIndicator}{,}\PYG{+w}{ }\PYG{n+nv}{0.759372342462676}\PYG{p+pIndicator}{,}\PYG{+w}{ }\PYG{n+nv}{\PYGZhy{}5.76149365656489}\PYG{p+pIndicator}{,}\PYG{+w}{ }\PYG{n+nv}{0}\PYG{p+pIndicator}{,}\PYG{+w}{ }\PYG{n+nv}{\PYGZhy{}4.62247073844943}\PYG{p+pIndicator}{,}\PYG{+w}{ }\PYG{n+nv}{0.191600048851267}\PYG{p+pIndicator}{,}\PYG{+w}{ }\PYG{p+pIndicator}{]}
\PYG{+w}{  }\PYG{p+pIndicator}{\PYGZhy{}}\PYG{+w}{ }\PYG{p+pIndicator}{[}\PYG{n+nv}{130}\PYG{p+pIndicator}{,}\PYG{+w}{ }\PYG{n+nv}{0.756833027516501}\PYG{p+pIndicator}{,}\PYG{+w}{ }\PYG{n+nv}{\PYGZhy{}5.75777334823494}\PYG{p+pIndicator}{,}\PYG{+w}{ }\PYG{n+nv}{0}\PYG{p+pIndicator}{,}\PYG{+w}{ }\PYG{n+nv}{\PYGZhy{}4.62255928350835}\PYG{p+pIndicator}{,}\PYG{+w}{ }\PYG{n+nv}{0.208792327853067}\PYG{p+pIndicator}{,}\PYG{+w}{ }\PYG{p+pIndicator}{]}
\PYG{n+nn}{...}
\end{sphinxVerbatim}

\sphinxAtStartPar
This data can be extracted and parsed from a log file using python with:

\begin{sphinxVerbatim}[commandchars=\\\{\}]
\PYG{k+kn}{import}\PYG{+w}{ }\PYG{n+nn}{re}\PYG{o}{,}\PYG{+w}{ }\PYG{n+nn}{yaml}
\PYG{k}{try}\PYG{p}{:}
    \PYG{k+kn}{from}\PYG{+w}{ }\PYG{n+nn}{yaml}\PYG{+w}{ }\PYG{k+kn}{import} \PYG{n}{CSafeLoader} \PYG{k}{as} \PYG{n}{Loader}
\PYG{k}{except} \PYG{n+ne}{ImportError}\PYG{p}{:}
    \PYG{k+kn}{from}\PYG{+w}{ }\PYG{n+nn}{yaml}\PYG{+w}{ }\PYG{k+kn}{import} \PYG{n}{SafeLoader} \PYG{k}{as} \PYG{n}{Loader}

\PYG{n}{docs} \PYG{o}{=} \PYG{l+s+s2}{\PYGZdq{}}\PYG{l+s+s2}{\PYGZdq{}}
\PYG{k}{with} \PYG{n+nb}{open}\PYG{p}{(}\PYG{l+s+s2}{\PYGZdq{}}\PYG{l+s+s2}{log.lammps}\PYG{l+s+s2}{\PYGZdq{}}\PYG{p}{)} \PYG{k}{as} \PYG{n}{f}\PYG{p}{:}
    \PYG{k}{for} \PYG{n}{line} \PYG{o+ow}{in} \PYG{n}{f}\PYG{p}{:}
        \PYG{n}{m} \PYG{o}{=} \PYG{n}{re}\PYG{o}{.}\PYG{n}{search}\PYG{p}{(}\PYG{l+s+sa}{r}\PYG{l+s+s2}{\PYGZdq{}}\PYG{l+s+s2}{\PYGZca{}(keywords:.*\PYGZdl{}|data:\PYGZdl{}|\PYGZhy{}\PYGZhy{}\PYGZhy{}\PYGZdl{}|}\PYG{l+s+s2}{\PYGZbs{}}\PYG{l+s+s2}{.}\PYG{l+s+s2}{\PYGZbs{}}\PYG{l+s+s2}{.}\PYG{l+s+s2}{\PYGZbs{}}\PYG{l+s+s2}{.\PYGZdl{}|  \PYGZhy{} }\PYG{l+s+s2}{\PYGZbs{}}\PYG{l+s+s2}{[.*}\PYG{l+s+s2}{\PYGZbs{}}\PYG{l+s+s2}{]\PYGZdl{})}\PYG{l+s+s2}{\PYGZdq{}}\PYG{p}{,} \PYG{n}{line}\PYG{p}{)}
        \PYG{k}{if} \PYG{n}{m}\PYG{p}{:} \PYG{n}{docs} \PYG{o}{+}\PYG{o}{=} \PYG{n}{m}\PYG{o}{.}\PYG{n}{group}\PYG{p}{(}\PYG{l+m+mi}{0}\PYG{p}{)} \PYG{o}{+} \PYG{l+s+s1}{\PYGZsq{}}\PYG{l+s+se}{\PYGZbs{}n}\PYG{l+s+s1}{\PYGZsq{}}

\PYG{n}{thermo} \PYG{o}{=} \PYG{n+nb}{list}\PYG{p}{(}\PYG{n}{yaml}\PYG{o}{.}\PYG{n}{load\PYGZus{}all}\PYG{p}{(}\PYG{n}{docs}\PYG{p}{,} \PYG{n}{Loader}\PYG{o}{=}\PYG{n}{Loader}\PYG{p}{)}\PYG{p}{)}

\PYG{n+nb}{print}\PYG{p}{(}\PYG{l+s+s2}{\PYGZdq{}}\PYG{l+s+s2}{Number of runs: }\PYG{l+s+s2}{\PYGZdq{}}\PYG{p}{,} \PYG{n+nb}{len}\PYG{p}{(}\PYG{n}{thermo}\PYG{p}{)}\PYG{p}{)}
\PYG{n+nb}{print}\PYG{p}{(}\PYG{n}{thermo}\PYG{p}{[}\PYG{l+m+mi}{1}\PYG{p}{]}\PYG{p}{[}\PYG{l+s+s1}{\PYGZsq{}}\PYG{l+s+s1}{keywords}\PYG{l+s+s1}{\PYGZsq{}}\PYG{p}{]}\PYG{p}{[}\PYG{l+m+mi}{4}\PYG{p}{]}\PYG{p}{,} \PYG{l+s+s1}{\PYGZsq{}}\PYG{l+s+s1}{ = }\PYG{l+s+s1}{\PYGZsq{}}\PYG{p}{,} \PYG{n}{thermo}\PYG{p}{[}\PYG{l+m+mi}{1}\PYG{p}{]}\PYG{p}{[}\PYG{l+s+s1}{\PYGZsq{}}\PYG{l+s+s1}{data}\PYG{l+s+s1}{\PYGZsq{}}\PYG{p}{]}\PYG{p}{[}\PYG{l+m+mi}{2}\PYG{p}{]}\PYG{p}{[}\PYG{l+m+mi}{4}\PYG{p}{]}\PYG{p}{)}
\end{sphinxVerbatim}

\sphinxAtStartPar
After loading the YAML data, \sphinxtitleref{thermo} is a list containing a dictionary
for each “run” where the tag “keywords” maps to the list of thermo
header strings and the tag “data” has a list of lists where the outer
list represents the lines of output and the inner list the values of the
columns matching the header keywords for that step.  The second print()
command for example will print the header string for the fifth keyword
of the second run and the corresponding value for the third output line
of that run:

\begin{sphinxVerbatim}[commandchars=\\\{\}]
Number of runs:  2
TotEng  =  \PYGZhy{}4.62140097780047
\end{sphinxVerbatim}


\paragraph{Extracting data from dump file}
\label{\detokenize{Howto_structured_data:extracting-data-from-dump-file}}
\sphinxAtStartPar
\DUrole{versionmodified,added}{Added in version 4May2022.}

\sphinxAtStartPar
YAML format output has been added to multiple commands in LAMMPS,
for example {\hyperref[\detokenize{dump::doc}]{\sphinxcrossref{\DUrole{doc}{dump yaml}}}} or {\hyperref[\detokenize{fix_ave_time::doc}]{\sphinxcrossref{\DUrole{doc}{fix ave/time}}}}
Depending on the kind of data being written, organization of the data
or the specific syntax used may change, but the principles are very
similar and all files should be readable with a suitable YAML parser.
A simple example for this is given below:

\begin{sphinxVerbatim}[commandchars=\\\{\}]
\PYG{k+kn}{import}\PYG{+w}{ }\PYG{n+nn}{yaml}
\PYG{k}{try}\PYG{p}{:}
    \PYG{k+kn}{from}\PYG{+w}{ }\PYG{n+nn}{yaml}\PYG{+w}{ }\PYG{k+kn}{import} \PYG{n}{CSafeLoader} \PYG{k}{as} \PYG{n}{YamlLoader}
\PYG{k}{except} \PYG{n+ne}{ImportError}\PYG{p}{:}
    \PYG{k+kn}{from}\PYG{+w}{ }\PYG{n+nn}{yaml}\PYG{+w}{ }\PYG{k+kn}{import} \PYG{n}{SafeLoader} \PYG{k}{as} \PYG{n}{YamlLoader}

\PYG{n}{timesteps} \PYG{o}{=} \PYG{p}{[}\PYG{p}{]}
\PYG{k}{with} \PYG{n+nb}{open}\PYG{p}{(}\PYG{l+s+s2}{\PYGZdq{}}\PYG{l+s+s2}{dump.yaml}\PYG{l+s+s2}{\PYGZdq{}}\PYG{p}{,} \PYG{l+s+s2}{\PYGZdq{}}\PYG{l+s+s2}{r}\PYG{l+s+s2}{\PYGZdq{}}\PYG{p}{)} \PYG{k}{as} \PYG{n}{f}\PYG{p}{:}
    \PYG{n}{data} \PYG{o}{=} \PYG{n}{yaml}\PYG{o}{.}\PYG{n}{load\PYGZus{}all}\PYG{p}{(}\PYG{n}{f}\PYG{p}{,} \PYG{n}{Loader}\PYG{o}{=}\PYG{n}{YamlLoader}\PYG{p}{)}

    \PYG{k}{for} \PYG{n}{d} \PYG{o+ow}{in} \PYG{n}{data}\PYG{p}{:}
        \PYG{n+nb}{print}\PYG{p}{(}\PYG{l+s+s1}{\PYGZsq{}}\PYG{l+s+s1}{Processing timestep }\PYG{l+s+si}{\PYGZpc{}d}\PYG{l+s+s1}{\PYGZsq{}} \PYG{o}{\PYGZpc{}} \PYG{n}{d}\PYG{p}{[}\PYG{l+s+s1}{\PYGZsq{}}\PYG{l+s+s1}{timestep}\PYG{l+s+s1}{\PYGZsq{}}\PYG{p}{]}\PYG{p}{)}
        \PYG{n}{timesteps}\PYG{o}{.}\PYG{n}{append}\PYG{p}{(}\PYG{n}{d}\PYG{p}{)}

\PYG{n+nb}{print}\PYG{p}{(}\PYG{l+s+s1}{\PYGZsq{}}\PYG{l+s+s1}{Read }\PYG{l+s+si}{\PYGZpc{}d}\PYG{l+s+s1}{ timesteps from yaml dump}\PYG{l+s+s1}{\PYGZsq{}} \PYG{o}{\PYGZpc{}} \PYG{n+nb}{len}\PYG{p}{(}\PYG{n}{timesteps}\PYG{p}{)}\PYG{p}{)}
\PYG{n+nb}{print}\PYG{p}{(}\PYG{l+s+s1}{\PYGZsq{}}\PYG{l+s+s1}{Second timestep: }\PYG{l+s+s1}{\PYGZsq{}}\PYG{p}{,} \PYG{n}{timesteps}\PYG{p}{[}\PYG{l+m+mi}{1}\PYG{p}{]}\PYG{p}{[}\PYG{l+s+s1}{\PYGZsq{}}\PYG{l+s+s1}{timestep}\PYG{l+s+s1}{\PYGZsq{}}\PYG{p}{]}\PYG{p}{)}
\PYG{n+nb}{print}\PYG{p}{(}\PYG{l+s+s1}{\PYGZsq{}}\PYG{l+s+s1}{Box info: x: }\PYG{l+s+s1}{\PYGZsq{}} \PYG{p}{,} \PYG{n}{timesteps}\PYG{p}{[}\PYG{l+m+mi}{1}\PYG{p}{]}\PYG{p}{[}\PYG{l+s+s1}{\PYGZsq{}}\PYG{l+s+s1}{box}\PYG{l+s+s1}{\PYGZsq{}}\PYG{p}{]}\PYG{p}{[}\PYG{l+m+mi}{0}\PYG{p}{]}\PYG{p}{,} \PYG{l+s+s1}{\PYGZsq{}}\PYG{l+s+s1}{ y:}\PYG{l+s+s1}{\PYGZsq{}}\PYG{p}{,} \PYG{n}{timesteps}\PYG{p}{[}\PYG{l+m+mi}{1}\PYG{p}{]}\PYG{p}{[}\PYG{l+s+s1}{\PYGZsq{}}\PYG{l+s+s1}{box}\PYG{l+s+s1}{\PYGZsq{}}\PYG{p}{]}\PYG{p}{[}\PYG{l+m+mi}{1}\PYG{p}{]}\PYG{p}{,} \PYG{l+s+s1}{\PYGZsq{}}\PYG{l+s+s1}{ z:}\PYG{l+s+s1}{\PYGZsq{}}\PYG{p}{,}\PYG{n}{timesteps}\PYG{p}{[}\PYG{l+m+mi}{1}\PYG{p}{]}\PYG{p}{[}\PYG{l+s+s1}{\PYGZsq{}}\PYG{l+s+s1}{box}\PYG{l+s+s1}{\PYGZsq{}}\PYG{p}{]}\PYG{p}{[}\PYG{l+m+mi}{2}\PYG{p}{]}\PYG{p}{)}
\PYG{n+nb}{print}\PYG{p}{(}\PYG{l+s+s1}{\PYGZsq{}}\PYG{l+s+s1}{First 5 per\PYGZhy{}atom columns: }\PYG{l+s+s1}{\PYGZsq{}}\PYG{p}{,} \PYG{n}{timesteps}\PYG{p}{[}\PYG{l+m+mi}{1}\PYG{p}{]}\PYG{p}{[}\PYG{l+s+s1}{\PYGZsq{}}\PYG{l+s+s1}{keywords}\PYG{l+s+s1}{\PYGZsq{}}\PYG{p}{]}\PYG{p}{[}\PYG{l+m+mi}{0}\PYG{p}{:}\PYG{l+m+mi}{5}\PYG{p}{]}\PYG{p}{)}
\PYG{n+nb}{print}\PYG{p}{(}\PYG{l+s+s1}{\PYGZsq{}}\PYG{l+s+s1}{Corresponding 10th atom data: }\PYG{l+s+s1}{\PYGZsq{}}\PYG{p}{,} \PYG{n}{timesteps}\PYG{p}{[}\PYG{l+m+mi}{1}\PYG{p}{]}\PYG{p}{[}\PYG{l+s+s1}{\PYGZsq{}}\PYG{l+s+s1}{data}\PYG{l+s+s1}{\PYGZsq{}}\PYG{p}{]}\PYG{p}{[}\PYG{l+m+mi}{9}\PYG{p}{]}\PYG{p}{[}\PYG{l+m+mi}{0}\PYG{p}{:}\PYG{l+m+mi}{5}\PYG{p}{]}\PYG{p}{)}
\end{sphinxVerbatim}

\sphinxAtStartPar
The corresponding output for a YAML dump command added to the “melt” example is:

\begin{sphinxVerbatim}[commandchars=\\\{\}]
Processing timestep 0
Processing timestep 50
Processing timestep 100
Processing timestep 150
Processing timestep 200
Processing timestep 250
Read 6 timesteps from yaml dump
Second timestep:  50
Box info: x:  [0, 16.795961913825074]  y: [0, 16.795961913825074]  z: [0, 16.795961913825074]
First 5 per\PYGZhy{}atom columns:  [\PYGZsq{}id\PYGZsq{}, \PYGZsq{}type\PYGZsq{}, \PYGZsq{}x\PYGZsq{}, \PYGZsq{}y\PYGZsq{}, \PYGZsq{}z\PYGZsq{}]
Corresponding 10th atom data:  [10, 1, 4.43828, 0.968481, 0.108555]
\end{sphinxVerbatim}


\paragraph{Processing scalar data with Python}
\label{\detokenize{Howto_structured_data:processing-scalar-data-with-python}}

\begin{wrapfigure}{r}{0.330\linewidth}
\centering
\noindent\sphinxincludegraphics{{thermo_bondeng}.png}
\end{wrapfigure}

\mbox{}\par\vskip-\dimexpr\baselineskip+\parskip\relax

\sphinxAtStartPar
After reading and parsing the YAML format data, it can be easily
imported for further processing and visualization with the \sphinxhref{https://pandas.pydata.org/}{pandas} and \sphinxhref{https://matplotlib.org/}{matplotlib} Python modules.  Because of the organization
of the data in the YAML format thermo output, it needs to be told to
process only the ‘data’ part of the imported data to create a pandas
data frame, and one needs to set the column names from the ‘keywords’
entry.  The following Python script code example demonstrates this, and
creates the image shown on the right of a simple plot of various bonded
energy contributions versus the timestep from a run of the ‘peptide’
example input after changing the {\hyperref[\detokenize{thermo_style::doc}]{\sphinxcrossref{\DUrole{doc}{thermo style}}}} to
‘yaml’.  The properties to be used for x and y values can be
conveniently selected through the keywords.  Please note that those
keywords can be changed to custom strings with the {\hyperref[\detokenize{thermo_modify::doc}]{\sphinxcrossref{\DUrole{doc}{thermo\_modify
colname}}}} command.

\begin{sphinxVerbatim}[commandchars=\\\{\}]
\PYG{k+kn}{import}\PYG{+w}{ }\PYG{n+nn}{re}\PYG{o}{,}\PYG{+w}{ }\PYG{n+nn}{yaml}
\PYG{k+kn}{import}\PYG{+w}{ }\PYG{n+nn}{pandas}\PYG{+w}{ }\PYG{k}{as}\PYG{+w}{ }\PYG{n+nn}{pd}
\PYG{k+kn}{import}\PYG{+w}{ }\PYG{n+nn}{matplotlib}\PYG{n+nn}{.}\PYG{n+nn}{pyplot}\PYG{+w}{ }\PYG{k}{as}\PYG{+w}{ }\PYG{n+nn}{plt}

\PYG{k}{try}\PYG{p}{:}
    \PYG{k+kn}{from}\PYG{+w}{ }\PYG{n+nn}{yaml}\PYG{+w}{ }\PYG{k+kn}{import} \PYG{n}{CSafeLoader} \PYG{k}{as} \PYG{n}{Loader}
\PYG{k}{except} \PYG{n+ne}{ImportError}\PYG{p}{:}
    \PYG{k+kn}{from}\PYG{+w}{ }\PYG{n+nn}{yaml}\PYG{+w}{ }\PYG{k+kn}{import} \PYG{n}{SafeLoader} \PYG{k}{as} \PYG{n}{Loader}

\PYG{n}{docs} \PYG{o}{=} \PYG{l+s+s2}{\PYGZdq{}}\PYG{l+s+s2}{\PYGZdq{}}
\PYG{k}{with} \PYG{n+nb}{open}\PYG{p}{(}\PYG{l+s+s2}{\PYGZdq{}}\PYG{l+s+s2}{log.lammps}\PYG{l+s+s2}{\PYGZdq{}}\PYG{p}{)} \PYG{k}{as} \PYG{n}{f}\PYG{p}{:}
    \PYG{k}{for} \PYG{n}{line} \PYG{o+ow}{in} \PYG{n}{f}\PYG{p}{:}
        \PYG{n}{m} \PYG{o}{=} \PYG{n}{re}\PYG{o}{.}\PYG{n}{search}\PYG{p}{(}\PYG{l+s+sa}{r}\PYG{l+s+s2}{\PYGZdq{}}\PYG{l+s+s2}{\PYGZca{}(keywords:.*\PYGZdl{}|data:\PYGZdl{}|\PYGZhy{}\PYGZhy{}\PYGZhy{}\PYGZdl{}|}\PYG{l+s+s2}{\PYGZbs{}}\PYG{l+s+s2}{.}\PYG{l+s+s2}{\PYGZbs{}}\PYG{l+s+s2}{.}\PYG{l+s+s2}{\PYGZbs{}}\PYG{l+s+s2}{.\PYGZdl{}|  \PYGZhy{} }\PYG{l+s+s2}{\PYGZbs{}}\PYG{l+s+s2}{[.*}\PYG{l+s+s2}{\PYGZbs{}}\PYG{l+s+s2}{]\PYGZdl{})}\PYG{l+s+s2}{\PYGZdq{}}\PYG{p}{,} \PYG{n}{line}\PYG{p}{)}
        \PYG{k}{if} \PYG{n}{m}\PYG{p}{:} \PYG{n}{docs} \PYG{o}{+}\PYG{o}{=} \PYG{n}{m}\PYG{o}{.}\PYG{n}{group}\PYG{p}{(}\PYG{l+m+mi}{0}\PYG{p}{)} \PYG{o}{+} \PYG{l+s+s1}{\PYGZsq{}}\PYG{l+s+se}{\PYGZbs{}n}\PYG{l+s+s1}{\PYGZsq{}}

\PYG{n}{thermo} \PYG{o}{=} \PYG{n+nb}{list}\PYG{p}{(}\PYG{n}{yaml}\PYG{o}{.}\PYG{n}{load\PYGZus{}all}\PYG{p}{(}\PYG{n}{docs}\PYG{p}{,} \PYG{n}{Loader}\PYG{o}{=}\PYG{n}{Loader}\PYG{p}{)}\PYG{p}{)}

\PYG{n}{df} \PYG{o}{=} \PYG{n}{pd}\PYG{o}{.}\PYG{n}{DataFrame}\PYG{p}{(}\PYG{n}{data}\PYG{o}{=}\PYG{n}{thermo}\PYG{p}{[}\PYG{l+m+mi}{0}\PYG{p}{]}\PYG{p}{[}\PYG{l+s+s1}{\PYGZsq{}}\PYG{l+s+s1}{data}\PYG{l+s+s1}{\PYGZsq{}}\PYG{p}{]}\PYG{p}{,} \PYG{n}{columns}\PYG{o}{=}\PYG{n}{thermo}\PYG{p}{[}\PYG{l+m+mi}{0}\PYG{p}{]}\PYG{p}{[}\PYG{l+s+s1}{\PYGZsq{}}\PYG{l+s+s1}{keywords}\PYG{l+s+s1}{\PYGZsq{}}\PYG{p}{]}\PYG{p}{)}
\PYG{n}{fig} \PYG{o}{=} \PYG{n}{df}\PYG{o}{.}\PYG{n}{plot}\PYG{p}{(}\PYG{n}{x}\PYG{o}{=}\PYG{l+s+s1}{\PYGZsq{}}\PYG{l+s+s1}{Step}\PYG{l+s+s1}{\PYGZsq{}}\PYG{p}{,} \PYG{n}{y}\PYG{o}{=}\PYG{p}{[}\PYG{l+s+s1}{\PYGZsq{}}\PYG{l+s+s1}{E\PYGZus{}bond}\PYG{l+s+s1}{\PYGZsq{}}\PYG{p}{,} \PYG{l+s+s1}{\PYGZsq{}}\PYG{l+s+s1}{E\PYGZus{}angle}\PYG{l+s+s1}{\PYGZsq{}}\PYG{p}{,} \PYG{l+s+s1}{\PYGZsq{}}\PYG{l+s+s1}{E\PYGZus{}dihed}\PYG{l+s+s1}{\PYGZsq{}}\PYG{p}{,} \PYG{l+s+s1}{\PYGZsq{}}\PYG{l+s+s1}{E\PYGZus{}impro}\PYG{l+s+s1}{\PYGZsq{}}\PYG{p}{]}\PYG{p}{,} \PYG{n}{ylabel}\PYG{o}{=}\PYG{l+s+s1}{\PYGZsq{}}\PYG{l+s+s1}{Energy in kcal/mol}\PYG{l+s+s1}{\PYGZsq{}}\PYG{p}{)}
\PYG{n}{plt}\PYG{o}{.}\PYG{n}{savefig}\PYG{p}{(}\PYG{l+s+s1}{\PYGZsq{}}\PYG{l+s+s1}{thermo\PYGZus{}bondeng.png}\PYG{l+s+s1}{\PYGZsq{}}\PYG{p}{)}
\end{sphinxVerbatim}


\paragraph{Processing vector data with Python}
\label{\detokenize{Howto_structured_data:processing-vector-data-with-python}}
\sphinxAtStartPar
Global \sphinxstyleemphasis{vector} data as produced by {\hyperref[\detokenize{fix_ave_time::doc}]{\sphinxcrossref{\DUrole{doc}{fix ave/time}}}}
uses a slightly different organization of the data. You still have the
dictionary keys ‘keywords’ and ‘data’ for the column headers and the
data.  But the data is a dictionary indexed by the time step and for
each step there are multiple rows of values each with a list of the
averaged properties.  This requires a slightly different processing,
since the entire data cannot be directly imported into a single pandas
DataFrame class instance.  The following Python script example
demonstrates how to read such data.  The result will combine the data
for the different steps into one large “multi\sphinxhyphen{}index” table.  The pandas
IndexSlice class can then be used to select data from this combined data
frame.

\begin{sphinxVerbatim}[commandchars=\\\{\}]
\PYG{k+kn}{import}\PYG{+w}{ }\PYG{n+nn}{yaml}
\PYG{k+kn}{import}\PYG{+w}{ }\PYG{n+nn}{pandas}\PYG{+w}{ }\PYG{k}{as}\PYG{+w}{ }\PYG{n+nn}{pd}

\PYG{k}{try}\PYG{p}{:}
    \PYG{k+kn}{from}\PYG{+w}{ }\PYG{n+nn}{yaml}\PYG{+w}{ }\PYG{k+kn}{import} \PYG{n}{CSafeLoader} \PYG{k}{as} \PYG{n}{Loader}
\PYG{k}{except} \PYG{n+ne}{ImportError}\PYG{p}{:}
    \PYG{k+kn}{from}\PYG{+w}{ }\PYG{n+nn}{yaml}\PYG{+w}{ }\PYG{k+kn}{import} \PYG{n}{SafeLoader} \PYG{k}{as} \PYG{n}{Loader}

\PYG{k}{with} \PYG{n+nb}{open}\PYG{p}{(}\PYG{l+s+s2}{\PYGZdq{}}\PYG{l+s+s2}{ave.yaml}\PYG{l+s+s2}{\PYGZdq{}}\PYG{p}{)} \PYG{k}{as} \PYG{n}{f}\PYG{p}{:}
    \PYG{n}{ave} \PYG{o}{=} \PYG{n}{yaml}\PYG{o}{.}\PYG{n}{load}\PYG{p}{(}\PYG{n}{f}\PYG{p}{,} \PYG{n}{Loader}\PYG{o}{=}\PYG{n}{Loader}\PYG{p}{)}

\PYG{n}{keys} \PYG{o}{=} \PYG{n}{ave}\PYG{p}{[}\PYG{l+s+s1}{\PYGZsq{}}\PYG{l+s+s1}{keywords}\PYG{l+s+s1}{\PYGZsq{}}\PYG{p}{]}
\PYG{n}{df} \PYG{o}{=} \PYG{p}{\PYGZob{}}\PYG{p}{\PYGZcb{}}
\PYG{k}{for} \PYG{n}{k} \PYG{o+ow}{in} \PYG{n}{ave}\PYG{p}{[}\PYG{l+s+s1}{\PYGZsq{}}\PYG{l+s+s1}{data}\PYG{l+s+s1}{\PYGZsq{}}\PYG{p}{]}\PYG{o}{.}\PYG{n}{keys}\PYG{p}{(}\PYG{p}{)}\PYG{p}{:}
    \PYG{n}{df}\PYG{p}{[}\PYG{n}{k}\PYG{p}{]} \PYG{o}{=} \PYG{n}{pd}\PYG{o}{.}\PYG{n}{DataFrame}\PYG{p}{(}\PYG{n}{data}\PYG{o}{=}\PYG{n}{ave}\PYG{p}{[}\PYG{l+s+s1}{\PYGZsq{}}\PYG{l+s+s1}{data}\PYG{l+s+s1}{\PYGZsq{}}\PYG{p}{]}\PYG{p}{[}\PYG{n}{k}\PYG{p}{]}\PYG{p}{,} \PYG{n}{columns}\PYG{o}{=}\PYG{n}{keys}\PYG{p}{)}

\PYG{c+c1}{\PYGZsh{} create multi\PYGZhy{}index data frame}
\PYG{n}{df} \PYG{o}{=} \PYG{n}{pd}\PYG{o}{.}\PYG{n}{concat}\PYG{p}{(}\PYG{n}{df}\PYG{p}{)}

\PYG{c+c1}{\PYGZsh{} output only the first 3 value for steps 200 to 300 of the column Pressure}
\PYG{n}{idx} \PYG{o}{=} \PYG{n}{pd}\PYG{o}{.}\PYG{n}{IndexSlice}
\PYG{n+nb}{print}\PYG{p}{(}\PYG{n}{df}\PYG{p}{[}\PYG{l+s+s1}{\PYGZsq{}}\PYG{l+s+s1}{Pressure}\PYG{l+s+s1}{\PYGZsq{}}\PYG{p}{]}\PYG{o}{.}\PYG{n}{loc}\PYG{p}{[}\PYG{n}{idx}\PYG{p}{[}\PYG{l+m+mi}{200}\PYG{p}{:}\PYG{l+m+mi}{300}\PYG{p}{,} \PYG{l+m+mi}{0}\PYG{p}{:}\PYG{l+m+mi}{2}\PYG{p}{]}\PYG{p}{]}\PYG{p}{)}
\end{sphinxVerbatim}


\paragraph{Processing scalar data with Perl}
\label{\detokenize{Howto_structured_data:processing-scalar-data-with-perl}}
\sphinxAtStartPar
The ease of processing YAML data is not limited to Python. Here is an
example for extracting and processing a LAMMPS log file with Perl instead.

\begin{sphinxVerbatim}[commandchars=\\\{\}]
\PYG{k}{use}\PYG{+w}{ }\PYG{n+nn}{YAML::XS}\PYG{p}{;}

\PYG{n+nb}{open}\PYG{p}{(}\PYG{n}{LOG}\PYG{p}{,}\PYG{+w}{ }\PYG{l+s}{\PYGZdq{}log.lammps\PYGZdq{}}\PYG{p}{)}\PYG{+w}{ }\PYG{o+ow}{or}\PYG{+w}{ }\PYG{n+nb}{die}\PYG{p}{(}\PYG{l+s}{\PYGZdq{}could not open log.lammps: \PYGZdl{}!\PYGZdq{}}\PYG{p}{)}\PYG{p}{;}
\PYG{k}{my}\PYG{+w}{ }\PYG{n+nv}{\PYGZdl{}}\PYG{n+nv}{file}\PYG{+w}{ }\PYG{o}{=}\PYG{+w}{ }\PYG{l+s}{\PYGZdq{}\PYGZdq{}}\PYG{p}{;}
\PYG{k}{while}\PYG{p}{(}\PYG{k}{my}\PYG{+w}{ }\PYG{n+nv}{\PYGZdl{}}\PYG{n+nv}{line}\PYG{+w}{ }\PYG{o}{=}\PYG{+w}{ }\PYG{l+s+sr}{\PYGZlt{}LOG\PYGZgt{}}\PYG{p}{)}\PYG{+w}{ }\PYG{p}{\PYGZob{}}
\PYG{+w}{    }\PYG{k}{if}\PYG{+w}{ }\PYG{p}{(}\PYG{n+nv}{\PYGZdl{}}\PYG{n+nv}{line}\PYG{+w}{ }\PYG{o}{=\PYGZti{}}\PYG{l+s+sr}{ /\PYGZca{}(keywords:.*\PYGZdl{}|data:\PYGZdl{}|\PYGZhy{}\PYGZhy{}\PYGZhy{}\PYGZdl{}|\PYGZbs{}.\PYGZbs{}.\PYGZbs{}.\PYGZdl{}|  \PYGZhy{} \PYGZbs{}[.*\PYGZbs{}]\PYGZdl{})/}\PYG{p}{)}\PYG{+w}{ }\PYG{p}{\PYGZob{}}
\PYG{+w}{        }\PYG{n+nv}{\PYGZdl{}}\PYG{n+nv}{file}\PYG{+w}{ }\PYG{o}{.}\PYG{o}{=}\PYG{+w}{ }\PYG{n+nv}{\PYGZdl{}}\PYG{n+nv}{line}\PYG{p}{;}
\PYG{+w}{    }\PYG{p}{\PYGZcb{}}
\PYG{p}{\PYGZcb{}}
\PYG{n+nb}{close}\PYG{p}{(}\PYG{n}{LOG}\PYG{p}{)}\PYG{p}{;}

\PYG{c+c1}{\PYGZsh{} convert YAML to perl as nested hash and array references}
\PYG{k}{my}\PYG{+w}{ }\PYG{n+nv}{\PYGZdl{}}\PYG{n+nv}{thermo}\PYG{+w}{ }\PYG{o}{=}\PYG{+w}{ }\PYG{n}{Load}\PYG{+w}{ }\PYG{n+nv}{\PYGZdl{}}\PYG{n+nv}{file}\PYG{p}{;}

\PYG{c+c1}{\PYGZsh{} convert references to real arrays}
\PYG{k}{my}\PYG{+w}{ }\PYG{n+nv}{@}\PYG{n+nv}{keywords}\PYG{+w}{ }\PYG{o}{=}\PYG{+w}{ }\PYG{n+nv}{@}\PYG{p}{\PYGZob{}}\PYG{n+nv}{\PYGZdl{}}\PYG{n+nv}{thermo}\PYG{o}{\PYGZhy{}}\PYG{o}{\PYGZgt{}}\PYG{p}{\PYGZob{}}\PYG{l+s}{\PYGZsq{}keywords\PYGZsq{}}\PYG{p}{\PYGZcb{}}\PYG{p}{\PYGZcb{}}\PYG{p}{;}
\PYG{k}{my}\PYG{+w}{ }\PYG{n+nv}{@}\PYG{n+nv}{data}\PYG{+w}{ }\PYG{o}{=}\PYG{+w}{ }\PYG{n+nv}{@}\PYG{p}{\PYGZob{}}\PYG{n+nv}{\PYGZdl{}}\PYG{n+nv}{thermo}\PYG{o}{\PYGZhy{}}\PYG{o}{\PYGZgt{}}\PYG{p}{\PYGZob{}}\PYG{l+s}{\PYGZsq{}data\PYGZsq{}}\PYG{p}{\PYGZcb{}}\PYG{p}{\PYGZcb{}}\PYG{p}{;}

\PYG{c+c1}{\PYGZsh{} print first two columns}
\PYG{k}{print}\PYG{p}{(}\PYG{l+s}{\PYGZdq{}\PYGZdl{}keywords[0] \PYGZdl{}keywords[1]\PYGZbs{}n\PYGZdq{}}\PYG{p}{)}\PYG{p}{;}
\PYG{k}{foreach}\PYG{+w}{ }\PYG{p}{(}\PYG{n+nv}{@}\PYG{n+nv}{data}\PYG{p}{)}\PYG{+w}{ }\PYG{p}{\PYGZob{}}
\PYG{+w}{    }\PYG{k}{print}\PYG{p}{(}\PYG{l+s}{\PYGZdq{}\PYGZdl{}\PYGZob{}\PYGZdl{}\PYGZus{}\PYGZcb{}[0]  \PYGZdl{}\PYGZob{}\PYGZdl{}\PYGZus{}\PYGZcb{}[1]\PYGZbs{}n\PYGZdq{}}\PYG{p}{)}\PYG{p}{;}
\PYG{p}{\PYGZcb{}}
\end{sphinxVerbatim}


\subsubsection{Writing continuous data during a simulation}
\label{\detokenize{Howto_structured_data:writing-continuous-data-during-a-simulation}}
\sphinxAtStartPar
The {\hyperref[\detokenize{fix_print::doc}]{\sphinxcrossref{\DUrole{doc}{fix print}}}} command allows you to output an arbitrary string at defined times during a simulation run.


\paragraph{YAML}
\label{\detokenize{Howto_structured_data:id1}}
\begin{sphinxVerbatim}[commandchars=\\\{\}]
\PYG{k}{fix }\PYG{n+nv+nvIdentifier}{extra}\PYG{+w}{ }\PYG{n+nv+nvIdentifier}{all}\PYG{+w}{ }\PYG{n}{print}\PYG{+w}{ }\PYG{l+m}{50}\PYG{+w}{ }\PYG{l+s}{\PYGZdq{}}\PYG{l+s}{\PYGZdq{}}\PYG{l+s}{\PYGZdq{}}
\PYG{l+s}{\PYGZhy{} timestep: \PYGZdl{}(step)}
\PYG{l+s}{  pe: \PYGZdl{}(pe)}
\PYG{l+s}{  ke: \PYGZdl{}(ke)}\PYG{l+s}{\PYGZdq{}}\PYG{l+s}{\PYGZdq{}}\PYG{l+s}{\PYGZdq{}}\PYG{+w}{ }\PYG{n}{file}\PYG{+w}{ }\PYG{n}{output.yaml}\PYG{+w}{ }\PYG{n}{screen}\PYG{+w}{ }\PYG{n}{no}
\end{sphinxVerbatim}
\sphinxSetupCaptionForVerbatim{output.yaml}
\def\sphinxLiteralBlockLabel{\label{\detokenize{Howto_structured_data:id4}}}
\begin{sphinxVerbatim}[commandchars=\\\{\}]
\PYG{c+c1}{\PYGZsh{} Fix print output for fix extra}
\PYG{p+pIndicator}{\PYGZhy{}}\PYG{+w}{ }\PYG{n+nt}{timestep}\PYG{p}{:}\PYG{+w}{ }\PYG{l+lScalar+lScalarPlain}{0}
\PYG{+w}{  }\PYG{n+nt}{pe}\PYG{p}{:}\PYG{+w}{ }\PYG{l+lScalar+lScalarPlain}{\PYGZhy{}6.77336805325924729}
\PYG{+w}{  }\PYG{n+nt}{ke}\PYG{p}{:}\PYG{+w}{ }\PYG{l+lScalar+lScalarPlain}{4.4988750000000026219}

\PYG{p+pIndicator}{\PYGZhy{}}\PYG{+w}{ }\PYG{n+nt}{timestep}\PYG{p}{:}\PYG{+w}{ }\PYG{l+lScalar+lScalarPlain}{50}
\PYG{+w}{  }\PYG{n+nt}{pe}\PYG{p}{:}\PYG{+w}{ }\PYG{l+lScalar+lScalarPlain}{\PYGZhy{}4.8082494418323200591}
\PYG{+w}{  }\PYG{n+nt}{ke}\PYG{p}{:}\PYG{+w}{ }\PYG{l+lScalar+lScalarPlain}{2.5257981827119797558}

\PYG{p+pIndicator}{\PYGZhy{}}\PYG{+w}{ }\PYG{n+nt}{timestep}\PYG{p}{:}\PYG{+w}{ }\PYG{l+lScalar+lScalarPlain}{100}
\PYG{+w}{  }\PYG{n+nt}{pe}\PYG{p}{:}\PYG{+w}{ }\PYG{l+lScalar+lScalarPlain}{\PYGZhy{}4.7875608875581505686}
\PYG{+w}{  }\PYG{n+nt}{ke}\PYG{p}{:}\PYG{+w}{ }\PYG{l+lScalar+lScalarPlain}{2.5062598821985102582}

\PYG{p+pIndicator}{\PYGZhy{}}\PYG{+w}{ }\PYG{n+nt}{timestep}\PYG{p}{:}\PYG{+w}{ }\PYG{l+lScalar+lScalarPlain}{150}
\PYG{+w}{  }\PYG{n+nt}{pe}\PYG{p}{:}\PYG{+w}{ }\PYG{l+lScalar+lScalarPlain}{\PYGZhy{}4.7471033686005483787}
\PYG{+w}{  }\PYG{n+nt}{ke}\PYG{p}{:}\PYG{+w}{ }\PYG{l+lScalar+lScalarPlain}{2.466095925545450207}

\PYG{p+pIndicator}{\PYGZhy{}}\PYG{+w}{ }\PYG{n+nt}{timestep}\PYG{p}{:}\PYG{+w}{ }\PYG{l+lScalar+lScalarPlain}{200}
\PYG{+w}{  }\PYG{n+nt}{pe}\PYG{p}{:}\PYG{+w}{ }\PYG{l+lScalar+lScalarPlain}{\PYGZhy{}4.7509052858544134068}
\PYG{+w}{  }\PYG{n+nt}{ke}\PYG{p}{:}\PYG{+w}{ }\PYG{l+lScalar+lScalarPlain}{2.4701136792591693592}

\PYG{p+pIndicator}{\PYGZhy{}}\PYG{+w}{ }\PYG{n+nt}{timestep}\PYG{p}{:}\PYG{+w}{ }\PYG{l+lScalar+lScalarPlain}{250}
\PYG{+w}{  }\PYG{n+nt}{pe}\PYG{p}{:}\PYG{+w}{ }\PYG{l+lScalar+lScalarPlain}{\PYGZhy{}4.7774327356321810711}
\PYG{+w}{  }\PYG{n+nt}{ke}\PYG{p}{:}\PYG{+w}{ }\PYG{l+lScalar+lScalarPlain}{2.4962152903997174569}
\end{sphinxVerbatim}

\sphinxAtStartPar
Post\sphinxhyphen{}processing of YAML files can be easily be done with Python and other
scripting languages. In case of Python the \sphinxtitleref{yaml} package allows you to load the
data files and obtain a list of dictionaries.

\begin{sphinxVerbatim}[commandchars=\\\{\}]
\PYG{k+kn}{import}\PYG{+w}{ }\PYG{n+nn}{yaml}

\PYG{k}{with} \PYG{n+nb}{open}\PYG{p}{(}\PYG{l+s+s2}{\PYGZdq{}}\PYG{l+s+s2}{output.yaml}\PYG{l+s+s2}{\PYGZdq{}}\PYG{p}{)} \PYG{k}{as} \PYG{n}{f}\PYG{p}{:}
   \PYG{n}{data} \PYG{o}{=} \PYG{n}{yaml}\PYG{o}{.}\PYG{n}{load}\PYG{p}{(}\PYG{n}{f}\PYG{p}{,} \PYG{n}{Loader}\PYG{o}{=}\PYG{n}{yaml}\PYG{o}{.}\PYG{n}{FullLoader}\PYG{p}{)}

\PYG{n+nb}{print}\PYG{p}{(}\PYG{n}{data}\PYG{p}{)}
\end{sphinxVerbatim}

\begin{sphinxVerbatim}[commandchars=\\\{\}]
[\PYGZob{}\PYGZsq{}timestep\PYGZsq{}: 0, \PYGZsq{}pe\PYGZsq{}: \PYGZhy{}6.773368053259247, \PYGZsq{}ke\PYGZsq{}: 4.498875000000003\PYGZcb{},
 \PYGZob{}\PYGZsq{}timestep\PYGZsq{}: 50, \PYGZsq{}pe\PYGZsq{}: \PYGZhy{}4.80824944183232, \PYGZsq{}ke\PYGZsq{}: 2.5257981827119798\PYGZcb{},
 \PYGZob{}\PYGZsq{}timestep\PYGZsq{}: 100, \PYGZsq{}pe\PYGZsq{}: \PYGZhy{}4.787560887558151, \PYGZsq{}ke\PYGZsq{}: 2.5062598821985103\PYGZcb{},
 \PYGZob{}\PYGZsq{}timestep\PYGZsq{}: 150, \PYGZsq{}pe\PYGZsq{}: \PYGZhy{}4.747103368600548, \PYGZsq{}ke\PYGZsq{}: 2.46609592554545\PYGZcb{},
 \PYGZob{}\PYGZsq{}timestep\PYGZsq{}: 200, \PYGZsq{}pe\PYGZsq{}: \PYGZhy{}4.750905285854413, \PYGZsq{}ke\PYGZsq{}: 2.4701136792591694\PYGZcb{},
 \PYGZob{}\PYGZsq{}timestep\PYGZsq{}: 250, \PYGZsq{}pe\PYGZsq{}: \PYGZhy{}4.777432735632181, \PYGZsq{}ke\PYGZsq{}: 2.4962152903997175\PYGZcb{}]
\end{sphinxVerbatim}


\paragraph{Line Delimited JSON (LD\sphinxhyphen{}JSON)}
\label{\detokenize{Howto_structured_data:line-delimited-json-ld-json}}
\sphinxAtStartPar
The JSON format itself is very strict when it comes to delimiters. For continuous
output/streaming data it is beneficial use the \sphinxstyleemphasis{line delimited JSON} format.
Each line represents one JSON object.

\begin{sphinxVerbatim}[commandchars=\\\{\}]
\PYG{k}{fix }\PYG{n+nv+nvIdentifier}{extra}\PYG{+w}{ }\PYG{n+nv+nvIdentifier}{all}\PYG{+w}{ }\PYG{n}{print}\PYG{+w}{ }\PYG{l+m}{50}\PYG{+w}{ }\PYG{l+s}{\PYGZdq{}}\PYG{l+s}{\PYGZdq{}}\PYG{l+s}{\PYGZdq{}}\PYG{l+s}{\PYGZob{}}\PYG{l+s}{\PYGZdq{}}\PYG{n}{timestep}\PYG{l+s}{\PYGZdq{}}\PYG{l+s}{: \PYGZdl{}(step), }\PYG{l+s}{\PYGZdq{}}\PYG{n}{pe}\PYG{l+s}{\PYGZdq{}}\PYG{l+s}{: \PYGZdl{}(pe), }\PYG{l+s}{\PYGZdq{}}\PYG{n}{ke}\PYG{l+s}{\PYGZdq{}}\PYG{l+s}{: \PYGZdl{}(ke)\PYGZcb{}}\PYG{l+s}{\PYGZdq{}}\PYG{l+s}{\PYGZdq{}}\PYG{l+s}{\PYGZdq{}}\PYG{+w}{ }\PYG{o}{\PYGZam{}}
\PYG{+w}{    }\PYG{n}{title}\PYG{+w}{ }\PYG{l+s}{\PYGZdq{}}\PYG{l+s}{\PYGZdq{}}\PYG{+w}{ }\PYG{n}{file}\PYG{+w}{ }\PYG{n}{output.json}\PYG{+w}{ }\PYG{n}{screen}\PYG{+w}{ }\PYG{n}{no}
\end{sphinxVerbatim}
\sphinxSetupCaptionForVerbatim{output.json}
\def\sphinxLiteralBlockLabel{\label{\detokenize{Howto_structured_data:id5}}}
\begin{sphinxVerbatim}[commandchars=\\\{\}]
\PYG{p}{\PYGZob{}}\PYG{n+nt}{\PYGZdq{}timestep\PYGZdq{}}\PYG{p}{:}\PYG{+w}{ }\PYG{l+m+mi}{0}\PYG{p}{,}\PYG{+w}{ }\PYG{n+nt}{\PYGZdq{}pe\PYGZdq{}}\PYG{p}{:}\PYG{+w}{ }\PYG{l+m+mf}{\PYGZhy{}6.77336805325924729}\PYG{p}{,}\PYG{+w}{ }\PYG{n+nt}{\PYGZdq{}ke\PYGZdq{}}\PYG{p}{:}\PYG{+w}{ }\PYG{l+m+mf}{4.4988750000000026219}\PYG{p}{\PYGZcb{}}
\PYG{p}{\PYGZob{}}\PYG{n+nt}{\PYGZdq{}timestep\PYGZdq{}}\PYG{p}{:}\PYG{+w}{ }\PYG{l+m+mi}{50}\PYG{p}{,}\PYG{+w}{ }\PYG{n+nt}{\PYGZdq{}pe\PYGZdq{}}\PYG{p}{:}\PYG{+w}{ }\PYG{l+m+mf}{\PYGZhy{}4.8082494418323200591}\PYG{p}{,}\PYG{+w}{ }\PYG{n+nt}{\PYGZdq{}ke\PYGZdq{}}\PYG{p}{:}\PYG{+w}{ }\PYG{l+m+mf}{2.5257981827119797558}\PYG{p}{\PYGZcb{}}
\PYG{p}{\PYGZob{}}\PYG{n+nt}{\PYGZdq{}timestep\PYGZdq{}}\PYG{p}{:}\PYG{+w}{ }\PYG{l+m+mi}{100}\PYG{p}{,}\PYG{+w}{ }\PYG{n+nt}{\PYGZdq{}pe\PYGZdq{}}\PYG{p}{:}\PYG{+w}{ }\PYG{l+m+mf}{\PYGZhy{}4.7875608875581505686}\PYG{p}{,}\PYG{+w}{ }\PYG{n+nt}{\PYGZdq{}ke\PYGZdq{}}\PYG{p}{:}\PYG{+w}{ }\PYG{l+m+mf}{2.5062598821985102582}\PYG{p}{\PYGZcb{}}
\PYG{p}{\PYGZob{}}\PYG{n+nt}{\PYGZdq{}timestep\PYGZdq{}}\PYG{p}{:}\PYG{+w}{ }\PYG{l+m+mi}{150}\PYG{p}{,}\PYG{+w}{ }\PYG{n+nt}{\PYGZdq{}pe\PYGZdq{}}\PYG{p}{:}\PYG{+w}{ }\PYG{l+m+mf}{\PYGZhy{}4.7471033686005483787}\PYG{p}{,}\PYG{+w}{ }\PYG{n+nt}{\PYGZdq{}ke\PYGZdq{}}\PYG{p}{:}\PYG{+w}{ }\PYG{l+m+mf}{2.466095925545450207}\PYG{p}{\PYGZcb{}}
\PYG{p}{\PYGZob{}}\PYG{n+nt}{\PYGZdq{}timestep\PYGZdq{}}\PYG{p}{:}\PYG{+w}{ }\PYG{l+m+mi}{200}\PYG{p}{,}\PYG{+w}{ }\PYG{n+nt}{\PYGZdq{}pe\PYGZdq{}}\PYG{p}{:}\PYG{+w}{ }\PYG{l+m+mf}{\PYGZhy{}4.7509052858544134068}\PYG{p}{,}\PYG{+w}{ }\PYG{n+nt}{\PYGZdq{}ke\PYGZdq{}}\PYG{p}{:}\PYG{+w}{ }\PYG{l+m+mf}{2.4701136792591693592}\PYG{p}{\PYGZcb{}}
\PYG{p}{\PYGZob{}}\PYG{n+nt}{\PYGZdq{}timestep\PYGZdq{}}\PYG{p}{:}\PYG{+w}{ }\PYG{l+m+mi}{250}\PYG{p}{,}\PYG{+w}{ }\PYG{n+nt}{\PYGZdq{}pe\PYGZdq{}}\PYG{p}{:}\PYG{+w}{ }\PYG{l+m+mf}{\PYGZhy{}4.7774327356321810711}\PYG{p}{,}\PYG{+w}{ }\PYG{n+nt}{\PYGZdq{}ke\PYGZdq{}}\PYG{p}{:}\PYG{+w}{ }\PYG{l+m+mf}{2.4962152903997174569}\PYG{p}{\PYGZcb{}}
\end{sphinxVerbatim}

\sphinxAtStartPar
One simple way to load this data into a Python script is to use the \sphinxtitleref{pandas}
package. It can directly load these files into a data frame:

\begin{sphinxVerbatim}[commandchars=\\\{\}]
\PYG{k+kn}{import}\PYG{+w}{ }\PYG{n+nn}{pandas}\PYG{+w}{ }\PYG{k}{as}\PYG{+w}{ }\PYG{n+nn}{pd}

\PYG{n}{data} \PYG{o}{=} \PYG{n}{pd}\PYG{o}{.}\PYG{n}{read\PYGZus{}json}\PYG{p}{(}\PYG{l+s+s1}{\PYGZsq{}}\PYG{l+s+s1}{output.json}\PYG{l+s+s1}{\PYGZsq{}}\PYG{p}{,} \PYG{n}{lines}\PYG{o}{=}\PYG{k+kc}{True}\PYG{p}{)}
\PYG{n+nb}{print}\PYG{p}{(}\PYG{n}{data}\PYG{p}{)}
\end{sphinxVerbatim}

\begin{sphinxVerbatim}[commandchars=\\\{\}]
\PYG{+w}{   }timestep\PYG{+w}{        }pe\PYG{+w}{        }ke
\PYG{l+m}{0}\PYG{+w}{         }\PYG{l+m}{0}\PYG{+w}{ }\PYGZhy{}6.773368\PYG{+w}{  }\PYG{l+m}{4}.498875
\PYG{l+m}{1}\PYG{+w}{        }\PYG{l+m}{50}\PYG{+w}{ }\PYGZhy{}4.808249\PYG{+w}{  }\PYG{l+m}{2}.525798
\PYG{l+m}{2}\PYG{+w}{       }\PYG{l+m}{100}\PYG{+w}{ }\PYGZhy{}4.787561\PYG{+w}{  }\PYG{l+m}{2}.506260
\PYG{l+m}{3}\PYG{+w}{       }\PYG{l+m}{150}\PYG{+w}{ }\PYGZhy{}4.747103\PYG{+w}{  }\PYG{l+m}{2}.466096
\PYG{l+m}{4}\PYG{+w}{       }\PYG{l+m}{200}\PYG{+w}{ }\PYGZhy{}4.750905\PYG{+w}{  }\PYG{l+m}{2}.470114
\PYG{l+m}{5}\PYG{+w}{       }\PYG{l+m}{250}\PYG{+w}{ }\PYGZhy{}4.777433\PYG{+w}{  }\PYG{l+m}{2}.496215
\end{sphinxVerbatim}


\section{Force fields howto}
\label{\detokenize{Howto:force-fields-howto}}
\sphinxstepscope


\subsection{CHARMM, AMBER, COMPASS, and DREIDING force fields}
\label{\detokenize{Howto_bioFF:charmm-amber-compass-and-dreiding-force-fields}}\label{\detokenize{Howto_bioFF::doc}}
\sphinxAtStartPar
A compact summary of the concepts, definitions, and properties of
force fields with explicit bonded interactions (like the ones discussed
in this HowTo) is given in {\hyperref[\detokenize{Howto_bioFF:typelabel2}]{\sphinxcrossref{\DUrole{std,std-ref}{(Gissinger)}}}}.

\sphinxAtStartPar
A force field has 2 parts: the formulas that define it and the
coefficients used for a particular system.  Here we only discuss
formulas implemented in LAMMPS that correspond to formulas commonly used
in the CHARMM, AMBER, COMPASS, and DREIDING force fields.  Setting
coefficients is done either from special sections in an input data file
via the {\hyperref[\detokenize{read_data::doc}]{\sphinxcrossref{\DUrole{doc}{read\_data}}}} command or in the input script with
commands like {\hyperref[\detokenize{pair_coeff::doc}]{\sphinxcrossref{\DUrole{doc}{pair\_coeff}}}} or {\hyperref[\detokenize{bond_coeff::doc}]{\sphinxcrossref{\DUrole{doc}{bond\_coeff}}}} and so on.  See the {\hyperref[\detokenize{Tools::doc}]{\sphinxcrossref{\DUrole{doc}{Tools}}}} doc page for
additional tools that can use CHARMM, AMBER, or Materials Studio
generated files to assign force field coefficients and convert their
output into LAMMPS input. LAMMPS input scripts can also be generated by
\sphinxhref{https://charmm-gui.org/}{charmm\sphinxhyphen{}gui.org}.


\subsubsection{CHARMM and AMBER}
\label{\detokenize{Howto_bioFF:charmm-and-amber}}
\sphinxAtStartPar
The \sphinxhref{https://mackerell.umaryland.edu/charmm\_ff.shtml}{CHARMM force field} {\hyperref[\detokenize{Howto_bioFF:howto-mackerell}]{\sphinxcrossref{\DUrole{std,std-ref}{(MacKerell)}}}} and \sphinxhref{https://ambermd.org/AmberModels.php}{AMBER force field} {\hyperref[\detokenize{Howto_bioFF:howto-cornell}]{\sphinxcrossref{\DUrole{std,std-ref}{(Cornell)}}}}
have potential energy function of the form
\begin{equation*}
\begin{split}V & = \sum_{bonds} E_b + \sum_{angles} \!E_a + \!\overbrace{\sum_{dihedral} \!\!E_d}^{\substack{
       \text{charmm} \\
      \text{charmmfsw}
    }} +\!\!\! \sum_{impropers} \!\!\!E_i \\[.6em]
    & \quad + \!\!\!\!\!\!\!\!\!\!\underbrace{~\sum_{pairs} \left(E_{LJ}+E_{coul}\right)}_{\substack{
       \text{lj/charmm/coul/charmm} \\
      \text{lj/charmm/coul/charmm/implicit} \\
      \text{lj/charmm/coul/long} \\
      \text{lj/charmm/coul/msm} \\
       \text{lj/charmmfsw/coul/charmmfsh} \\
      \text{lj/charmmfsw/coul/long}
    }} \!\!\!\!\!\!\!\!+ \!\!\sum_{special}\! E_s + \!\!\!\!\sum_{residues} \!\!\!{\scriptstyle\mathrm{CMAP}(\phi,\psi)}\end{split}
\end{equation*}
\sphinxAtStartPar
The terms are computed by bond styles (relationship between two atoms),
angle styles (between 3 atoms) , dihedral/improper styles (between 4
atoms), pair styles (non\sphinxhyphen{}covalently bonded pair interactions) and
special bonds. The CMAP term (see {\hyperref[\detokenize{fix_cmap::doc}]{\sphinxcrossref{\DUrole{doc}{fix cmap}}}} command for
details) corrects for pairs of dihedral angles (“Correction MAP”) to
significantly improve the structural and dynamic properties of proteins
in crystalline and solution environments {\hyperref[\detokenize{Howto_bioFF:howto-brooks}]{\sphinxcrossref{\DUrole{std,std-ref}{(Brooks)}}}}. The AMBER force field does not include the CMAP term.

\sphinxAtStartPar
The interaction styles listed below compute force field formulas that
are consistent with common options in CHARMM or AMBER.  See each
command’s documentation for the formula it computes.
\begin{itemize}
\item {} 
\sphinxAtStartPar
{\hyperref[\detokenize{bond_harmonic::doc}]{\sphinxcrossref{\DUrole{doc}{bond\_style}}}} harmonic

\item {} 
\sphinxAtStartPar
{\hyperref[\detokenize{angle_charmm::doc}]{\sphinxcrossref{\DUrole{doc}{angle\_style}}}} charmm

\item {} 
\sphinxAtStartPar
{\hyperref[\detokenize{dihedral_charmm::doc}]{\sphinxcrossref{\DUrole{doc}{dihedral\_style}}}} charmmfsh

\item {} 
\sphinxAtStartPar
{\hyperref[\detokenize{dihedral_charmm::doc}]{\sphinxcrossref{\DUrole{doc}{dihedral\_style}}}} charmm

\item {} 
\sphinxAtStartPar
{\hyperref[\detokenize{pair_charmm::doc}]{\sphinxcrossref{\DUrole{doc}{pair\_style}}}} lj/charmmfsw/coul/charmmfsh

\item {} 
\sphinxAtStartPar
{\hyperref[\detokenize{pair_charmm::doc}]{\sphinxcrossref{\DUrole{doc}{pair\_style}}}} lj/charmmfsw/coul/long

\item {} 
\sphinxAtStartPar
{\hyperref[\detokenize{pair_charmm::doc}]{\sphinxcrossref{\DUrole{doc}{pair\_style}}}} lj/charmm/coul/charmm

\item {} 
\sphinxAtStartPar
{\hyperref[\detokenize{pair_charmm::doc}]{\sphinxcrossref{\DUrole{doc}{pair\_style}}}} lj/charmm/coul/charmm/implicit

\item {} 
\sphinxAtStartPar
{\hyperref[\detokenize{pair_charmm::doc}]{\sphinxcrossref{\DUrole{doc}{pair\_style}}}} lj/charmm/coul/long

\item {} 
\sphinxAtStartPar
{\hyperref[\detokenize{special_bonds::doc}]{\sphinxcrossref{\DUrole{doc}{special\_bonds}}}} charmm

\item {} 
\sphinxAtStartPar
{\hyperref[\detokenize{special_bonds::doc}]{\sphinxcrossref{\DUrole{doc}{special\_bonds}}}} amber

\end{itemize}

\sphinxAtStartPar
The pair styles compute Lennard Jones (LJ) and Coulombic interactions
with additional switching or shifting functions that ramp the energy
and/or force smoothly to zero between an inner \((a)\) and outer
\((b)\) cutoff. The older styles with \sphinxstyleemphasis{charmm} (not \sphinxstyleemphasis{charmmfsw} or
\sphinxstyleemphasis{charmmfsh}) in their name compute the LJ and Coulombic interactions
with an energy switching function (esw) \(S(r)\) which ramps the energy
smoothly to zero between the inner and outer cutoff. This can cause
irregularities in pairwise forces (due to the discontinuous second
derivative of energy at the boundaries of the switching region), which
in some cases can result in complications in energy minimization and
detectable artifacts in MD simulations.

\begin{sphinxuseclass}{sd-container-fluid}
\begin{sphinxuseclass}{sd-sphinx-override}
\begin{sphinxuseclass}{sd-mb-4}
\begin{sphinxuseclass}{sd-row}
\begin{sphinxuseclass}{sd-row-cols-1}
\begin{sphinxuseclass}{sd-row-cols-xs-1}
\begin{sphinxuseclass}{sd-row-cols-sm-1}
\begin{sphinxuseclass}{sd-row-cols-md-2}
\begin{sphinxuseclass}{sd-row-cols-lg-2}
\begin{sphinxuseclass}{sd-col}
\begin{sphinxuseclass}{sd-d-flex-column}\begin{equation*}
\begin{split}LJ(r) &= 4 \epsilon \left[ \left(\frac{\sigma}{r}\right)^{12} -
         \left(\frac{\sigma}{r}\right)^6 \right]\\[.6em]
C(r) &= \frac{C q_i q_j}{ \epsilon r}\\[.6em]
S(r) &=  \frac{ \left(b^2 - r^2\right)^2 \left(b^2 + 2r^2 - 3{a^2}\right)}
        { \left(b^2 - a^2\right)^3 }\\[.6em]
E_{LJ}(r) &=  \begin{cases}
  LJ(r), & r \leq a \\
  LJ(r) S(r), & a < r \leq b \\
  0, &r > b
\end{cases} \\[.6em]
E_{coul}(r) &=  \begin{cases}
  C(r), & r \leq a \\
  C(r) S(r), & a < r \leq b \\
  0, & r > b
\end{cases}\end{split}
\end{equation*}
\end{sphinxuseclass}
\end{sphinxuseclass}
\begin{sphinxuseclass}{sd-col}
\begin{sphinxuseclass}{sd-d-flex-column}
\noindent{\hspace*{\fill}\sphinxincludegraphics{{howto_charmm_ELJ}.png}\hspace*{\fill}}

\end{sphinxuseclass}
\end{sphinxuseclass}
\end{sphinxuseclass}
\end{sphinxuseclass}
\end{sphinxuseclass}
\end{sphinxuseclass}
\end{sphinxuseclass}
\end{sphinxuseclass}
\end{sphinxuseclass}
\end{sphinxuseclass}
\end{sphinxuseclass}
\sphinxAtStartPar
The newer styles with \sphinxstyleemphasis{charmmfsw} or \sphinxstyleemphasis{charmmfsh} in their name replace
energy switching with force switching (fsw) for LJ interactions and
force shifting (fsh) functions for Coulombic interactions
{\hyperref[\detokenize{Howto_bioFF:howto-steinbach}]{\sphinxcrossref{\DUrole{std,std-ref}{(Steinbach)}}}}

\begin{sphinxuseclass}{sd-container-fluid}
\begin{sphinxuseclass}{sd-sphinx-override}
\begin{sphinxuseclass}{sd-mb-4}
\begin{sphinxuseclass}{sd-row}
\begin{sphinxuseclass}{sd-row-cols-1}
\begin{sphinxuseclass}{sd-row-cols-xs-1}
\begin{sphinxuseclass}{sd-row-cols-sm-1}
\begin{sphinxuseclass}{sd-row-cols-md-2}
\begin{sphinxuseclass}{sd-row-cols-lg-2}
\begin{sphinxuseclass}{sd-col}
\begin{sphinxuseclass}{sd-d-flex-column}\begin{equation*}
\begin{split}    E_{LJ}(r) = & \begin{cases}
4  \epsilon \sigma^6  \left(\frac{\displaystyle\sigma
  ^6-r^6}{\displaystyle r^{12}}-\frac{\displaystyle\sigma ^6}{\displaystyle a^6
  b^6}+\frac{\displaystyle 1}{\displaystyle a^3 b^3}\right) & r\leq a \\
\frac{\displaystyle 4 \epsilon \sigma^6   \left(\sigma ^6
  \left(b^6-r^6\right)^2-b^3 r^6 \left(a^3+b^3\right)
  \left(b^3-r^3\right)^2\right)}{\displaystyle b^6 r^{12}
  \left(b^6-a^6\right)} & a<r \leq b\\
  0, & r>b
 \end{cases}\\[.6em]
  E_{coul}(r) & =  \begin{cases}
       C(r) \frac{\displaystyle (b-r)^2}{\displaystyle r b^2}, &  r \leq b \\
       0, & r > b
     \end{cases}\end{split}
\end{equation*}
\end{sphinxuseclass}
\end{sphinxuseclass}
\begin{sphinxuseclass}{sd-col}
\begin{sphinxuseclass}{sd-d-flex-column}
\noindent{\hspace*{\fill}\sphinxincludegraphics{{howto_charmmfsw_ELJ}.png}\hspace*{\fill}}

\end{sphinxuseclass}
\end{sphinxuseclass}
\end{sphinxuseclass}
\end{sphinxuseclass}
\end{sphinxuseclass}
\end{sphinxuseclass}
\end{sphinxuseclass}
\end{sphinxuseclass}
\end{sphinxuseclass}
\end{sphinxuseclass}
\end{sphinxuseclass}
\sphinxAtStartPar
These styles are used by LAMMPS input scripts generated by
\sphinxurl{https://charmm-gui.org/} {\hyperref[\detokenize{Howto_bioFF:howto-brooks}]{\sphinxcrossref{\DUrole{std,std-ref}{(Brooks)}}}}.

\begin{sphinxadmonition}{note}{Note:}
\sphinxAtStartPar
For CHARMM, newer \sphinxstyleemphasis{charmmfsw} or \sphinxstyleemphasis{charmmfsh} styles were released in
March 2017.  We recommend they be used instead of the older \sphinxstyleemphasis{charmm}
styles.  See discussion of the differences on the {\hyperref[\detokenize{pair_charmm::doc}]{\sphinxcrossref{\DUrole{doc}{pair charmm}}}} and {\hyperref[\detokenize{dihedral_charmm::doc}]{\sphinxcrossref{\DUrole{doc}{dihedral charmm}}}} doc
pages.
\end{sphinxadmonition}

\begin{sphinxadmonition}{note}{Note:}
\sphinxAtStartPar
The TIP3P water model is strongly recommended for use with the CHARMM
force field. In fact, \sphinxhref{https://matsci.org/t/using-spc-water-with-charmm-ff/24715}{“using the SPC model with CHARMM parameters is
a bad idea”} and \sphinxhref{https://matsci.org/t/hybrid-pair-styles-for-charmm-and-tip4p-ew/32609}{“to
enable TIP4P style water in CHARMM, you would have to write a new pair
style”}
. LAMMPS input scripts generated by Solution Builder on \sphinxurl{https://charmm-gui.org}
use TIP3P molecules for solvation.  Any other water model can and
probably will lead to false conclusions.
\end{sphinxadmonition}


\subsubsection{COMPASS}
\label{\detokenize{Howto_bioFF:compass}}
\sphinxAtStartPar
COMPASS is a general force field for atomistic simulation of common
organic molecules, inorganic small molecules, and polymers which was
developed using ab initio and empirical parameterization techniques
{\hyperref[\detokenize{Howto_bioFF:howto-sun}]{\sphinxcrossref{\DUrole{std,std-ref}{(Sun)}}}}.  See the {\hyperref[\detokenize{Tools::doc}]{\sphinxcrossref{\DUrole{doc}{Tools}}}} page for the
msi2lmp tool for creating LAMMPS template input and data files from
BIOVIA’s Materials Studio files.  Please note that the msi2lmp tool is
very old and largely unmaintained, so it does not support all features
of Materials Studio provided force field files, especially additions
during the last decade.  You should watch the output carefully and
compare results, where possible.  See {\hyperref[\detokenize{Howto_bioFF:howto-sun}]{\sphinxcrossref{\DUrole{std,std-ref}{(Sun)}}}} for a
description of the COMPASS force field.

\sphinxAtStartPar
These interaction styles listed below compute force field formulas that
are consistent with the COMPASS force field.  See each command’s
documentation for the formula it computes.
\begin{itemize}
\item {} 
\sphinxAtStartPar
{\hyperref[\detokenize{bond_class2::doc}]{\sphinxcrossref{\DUrole{doc}{bond\_style}}}} class2

\item {} 
\sphinxAtStartPar
{\hyperref[\detokenize{angle_class2::doc}]{\sphinxcrossref{\DUrole{doc}{angle\_style}}}} class2

\item {} 
\sphinxAtStartPar
{\hyperref[\detokenize{dihedral_class2::doc}]{\sphinxcrossref{\DUrole{doc}{dihedral\_style}}}} class2

\item {} 
\sphinxAtStartPar
{\hyperref[\detokenize{improper_class2::doc}]{\sphinxcrossref{\DUrole{doc}{improper\_style}}}} class2

\item {} 
\sphinxAtStartPar
{\hyperref[\detokenize{pair_class2::doc}]{\sphinxcrossref{\DUrole{doc}{pair\_style}}}} lj/class2

\item {} 
\sphinxAtStartPar
{\hyperref[\detokenize{pair_class2::doc}]{\sphinxcrossref{\DUrole{doc}{pair\_style}}}} lj/class2/coul/cut

\item {} 
\sphinxAtStartPar
{\hyperref[\detokenize{pair_class2::doc}]{\sphinxcrossref{\DUrole{doc}{pair\_style}}}} lj/class2/coul/long

\item {} 
\sphinxAtStartPar
{\hyperref[\detokenize{special_bonds::doc}]{\sphinxcrossref{\DUrole{doc}{special\_bonds}}}} lj/coul 0 0 1

\end{itemize}


\subsubsection{DREIDING}
\label{\detokenize{Howto_bioFF:dreiding}}
\sphinxAtStartPar
DREIDING is a generic force field developed by the \sphinxhref{http://www.wag.caltech.edu}{Goddard group} at Caltech and is useful for predicting
structures and dynamics of organic, biological and main\sphinxhyphen{}group inorganic
molecules.  The philosophy in DREIDING is to use general force constants
and geometry parameters based on simple hybridization considerations,
rather than individual force constants and geometric parameters that
depend on the particular combinations of atoms involved in the bond,
angle, or torsion terms.  DREIDING has an {\hyperref[\detokenize{pair_hbond_dreiding::doc}]{\sphinxcrossref{\DUrole{doc}{explicit hydrogen bond
term}}}} to describe interactions involving a
hydrogen atom on very electronegative atoms (N, O, F).  Unlike CHARMM
or AMBER, the DREIDING force field has not been parameterized for
considering solvents (like water).

\sphinxAtStartPar
See {\hyperref[\detokenize{Howto_bioFF:howto-mayo}]{\sphinxcrossref{\DUrole{std,std-ref}{(Mayo)}}}} for a description of the DREIDING force field

\sphinxAtStartPar
The interaction styles listed below compute force field formulas that
are consistent with the DREIDING force field.  See each command’s
documentation for the formula it computes.
\begin{itemize}
\item {} 
\sphinxAtStartPar
{\hyperref[\detokenize{bond_harmonic::doc}]{\sphinxcrossref{\DUrole{doc}{bond\_style}}}} harmonic

\item {} 
\sphinxAtStartPar
{\hyperref[\detokenize{bond_morse::doc}]{\sphinxcrossref{\DUrole{doc}{bond\_style}}}} morse

\item {} 
\sphinxAtStartPar
{\hyperref[\detokenize{angle_cosine_squared::doc}]{\sphinxcrossref{\DUrole{doc}{angle\_style}}}} cosine/squared

\item {} 
\sphinxAtStartPar
{\hyperref[\detokenize{angle_harmonic::doc}]{\sphinxcrossref{\DUrole{doc}{angle\_style}}}} harmonic

\item {} 
\sphinxAtStartPar
{\hyperref[\detokenize{angle_cosine::doc}]{\sphinxcrossref{\DUrole{doc}{angle\_style}}}} cosine

\item {} 
\sphinxAtStartPar
{\hyperref[\detokenize{angle_cosine_periodic::doc}]{\sphinxcrossref{\DUrole{doc}{angle\_style}}}} cosine/periodic

\item {} 
\sphinxAtStartPar
{\hyperref[\detokenize{dihedral_charmm::doc}]{\sphinxcrossref{\DUrole{doc}{dihedral\_style}}}} charmm

\item {} 
\sphinxAtStartPar
{\hyperref[\detokenize{improper_umbrella::doc}]{\sphinxcrossref{\DUrole{doc}{improper\_style}}}} umbrella

\item {} 
\sphinxAtStartPar
{\hyperref[\detokenize{pair_buck::doc}]{\sphinxcrossref{\DUrole{doc}{pair\_style}}}} buck

\item {} 
\sphinxAtStartPar
{\hyperref[\detokenize{pair_buck::doc}]{\sphinxcrossref{\DUrole{doc}{pair\_style}}}} buck/coul/cut

\item {} 
\sphinxAtStartPar
{\hyperref[\detokenize{pair_buck::doc}]{\sphinxcrossref{\DUrole{doc}{pair\_style}}}} buck/coul/long

\item {} 
\sphinxAtStartPar
{\hyperref[\detokenize{pair_lj::doc}]{\sphinxcrossref{\DUrole{doc}{pair\_style}}}} lj/cut

\item {} 
\sphinxAtStartPar
{\hyperref[\detokenize{pair_lj_cut_coul::doc}]{\sphinxcrossref{\DUrole{doc}{pair\_style}}}} lj/cut/coul/cut

\item {} 
\sphinxAtStartPar
{\hyperref[\detokenize{pair_lj_cut_coul::doc}]{\sphinxcrossref{\DUrole{doc}{pair\_style}}}} lj/cut/coul/long

\item {} 
\sphinxAtStartPar
{\hyperref[\detokenize{pair_hbond_dreiding::doc}]{\sphinxcrossref{\DUrole{doc}{pair\_style}}}} hbond/dreiding/lj

\item {} 
\sphinxAtStartPar
{\hyperref[\detokenize{pair_hbond_dreiding::doc}]{\sphinxcrossref{\DUrole{doc}{pair\_style}}}} hbond/dreiding/morse

\item {} 
\sphinxAtStartPar
{\hyperref[\detokenize{special_bonds::doc}]{\sphinxcrossref{\DUrole{doc}{special\_bonds}}}} dreiding

\end{itemize}


\bigskip\hrule\bigskip

\phantomsection\label{\detokenize{Howto_bioFF:typelabel2}}
\sphinxAtStartPar
\sphinxstylestrong{(Gissinger)} J. R. Gissinger, I. Nikiforov, Y. Afshar, B. Waters, M. Choi, D. S. Karls, A. Stukowski, W. Im, H. Heinz, A. Kohlmeyer, and E. B. Tadmor, J Phys Chem B, 128, 3282\sphinxhyphen{}3297 (2024).

\phantomsection\label{\detokenize{Howto_bioFF:howto-mackerell}}
\sphinxAtStartPar
\sphinxstylestrong{(MacKerell)} MacKerell, Bashford, Bellott, Dunbrack, Evanseck, Field, Fischer, Gao, Guo, Ha, et al (1998).  J Phys Chem, 102, 3586 . \sphinxurl{https://doi.org/10.1021/jp973084f}

\phantomsection\label{\detokenize{Howto_bioFF:howto-cornell}}
\sphinxAtStartPar
\sphinxstylestrong{(Cornell)} Cornell, Cieplak, Bayly, Gould, Merz, Ferguson, Spellmeyer, Fox, Caldwell, Kollman (1995).  JACS 117, 5179\sphinxhyphen{}5197. \sphinxurl{https://doi.org/10.1021/ja00124a002}

\phantomsection\label{\detokenize{Howto_bioFF:howto-steinbach}}
\sphinxAtStartPar
\sphinxstylestrong{(Steinbach)} Steinbach, Brooks (1994). J Comput Chem, 15, 667. \sphinxurl{https://doi.org/10.1002/jcc.540150702}

\phantomsection\label{\detokenize{Howto_bioFF:howto-brooks}}
\sphinxAtStartPar
\sphinxstylestrong{(Brooks)} Brooks, et al (2009). J Comput Chem, 30, 1545. \sphinxurl{https://onlinelibrary.wiley.com/doi/10.1002/jcc.21287}

\phantomsection\label{\detokenize{Howto_bioFF:howto-sun}}
\sphinxAtStartPar
\sphinxstylestrong{(Sun)} Sun (1998). J. Phys. Chem. B, 102, 7338\sphinxhyphen{}7364. \sphinxurl{https://doi.org/10.1021/jp980939v}

\phantomsection\label{\detokenize{Howto_bioFF:howto-mayo}}
\sphinxAtStartPar
\sphinxstylestrong{(Mayo)} Mayo, Olfason, Goddard III (1990). J Phys Chem, 94, 8897\sphinxhyphen{}8909. \sphinxurl{https://doi.org/10.1021/j100389a010}

\sphinxstepscope


\subsection{AMOEBA and HIPPO force fields}
\label{\detokenize{Howto_amoeba:amoeba-and-hippo-force-fields}}\label{\detokenize{Howto_amoeba::doc}}
\sphinxAtStartPar
The AMOEBA and HIPPO polarizable force fields were developed by Jay
Ponder’s group at the U Washington at St Louis.  The LAMMPS
implementation is based on Fortran 90 code provided by the Ponder
group in their \sphinxhref{https://dasher.wustl.edu/tinker/}{Tinker MD software}.

\sphinxAtStartPar
The current implementation (July 2022) of AMOEBA in LAMMPS matches the
version discussed in {\hyperref[\detokenize{pair_amoeba:amoeba-ponder}]{\sphinxcrossref{\DUrole{std,std-ref}{(Ponder)}}}}, {\hyperref[\detokenize{pair_amoeba:amoeba-ren}]{\sphinxcrossref{\DUrole{std,std-ref}{(Ren)}}}}, and {\hyperref[\detokenize{pair_amoeba:amoeba-shi}]{\sphinxcrossref{\DUrole{std,std-ref}{(Shi)}}}}.  Likewise the current
implementation of HIPPO in LAMMPS matches the version discussed in
{\hyperref[\detokenize{pair_amoeba:amoeba-rackers}]{\sphinxcrossref{\DUrole{std,std-ref}{(Rackers)}}}}.

\sphinxAtStartPar
These force fields can be used when polarization effects are desired
in simulations of water, organic molecules, and biomolecules including
proteins, provided that parameterizations (Tinker PRM force field
files) are available for the systems you are interested in.  Files in
the LAMMPS potentials directory with a “amoeba” or “hippo” suffix can
be used.  The Tinker distribution and website have additional force
field files as well:
\sphinxurl{https://github.com/TinkerTools/tinker/tree/release/params}.

\sphinxAtStartPar
Note that currently, HIPPO can only be used for water systems, but
HIPPO files for a variety of small organic and biomolecules are in
preparation by the Ponder group.  Those force field files will be
included in the LAMMPS distribution when available.

\sphinxAtStartPar
To use the AMOEBA or HIPPO force fields, a simulation must be 3d, and
fully periodic or fully non\sphinxhyphen{}periodic, and use an orthogonal (not
triclinic) simulation box.


\bigskip\hrule\bigskip


\sphinxAtStartPar
The AMOEBA and HIPPO force fields contain the following terms in their
energy (U) computation.  Further details for AMOEBA equations are in
{\hyperref[\detokenize{pair_amoeba:amoeba-ponder}]{\sphinxcrossref{\DUrole{std,std-ref}{(Ponder)}}}}, further details for the HIPPO
equations are in {\hyperref[\detokenize{pair_amoeba:amoeba-rackers}]{\sphinxcrossref{\DUrole{std,std-ref}{(Rackers)}}}}.
\begin{equation*}
\begin{split}U & = U_{intermolecular} + U_{intramolecular} \\
U_{intermolecular} & = U_{hal} + U_{repulsion} + U_{dispersion} + U_{multipole} + U_{polar} + U_{qxfer} \\
U_{intramolecular} & = U_{bond} + U_{angle} + U_{torsion} + U_{oop} + U_{b\theta} + U_{UB} + U_{pitorsion} + U_{bitorsion}\end{split}
\end{equation*}
\sphinxAtStartPar
For intermolecular terms, the AMOEBA force field includes only the
\(U_{hal}\), \(U_{multipole}\), \(U_{polar}\) terms.  The
HIPPO force field includes all but the \(U_{hal}\) term.  In
LAMMPS, these are all computed by the {\hyperref[\detokenize{pair_style::doc}]{\sphinxcrossref{\DUrole{doc}{pair\_style amoeba or hippo}}}} command.  Note that the \(U_{multipole}\) and
\(U_{polar}\) terms in this formula are not the same for the
AMOEBA and HIPPO force fields.

\sphinxAtStartPar
For intramolecular terms, the \(U_{bond}\), \(U_{angle}\),
\(U_{torsion}\), \(U_{oop}\) terms are computed by the
{\hyperref[\detokenize{bond_class2::doc}]{\sphinxcrossref{\DUrole{doc}{bond\_style class2}}}} {\hyperref[\detokenize{angle_amoeba::doc}]{\sphinxcrossref{\DUrole{doc}{angle\_style amoeba}}}}, {\hyperref[\detokenize{dihedral_fourier::doc}]{\sphinxcrossref{\DUrole{doc}{dihedral\_style fourier}}}}, and
{\hyperref[\detokenize{improper_amoeba::doc}]{\sphinxcrossref{\DUrole{doc}{improper\_style amoeba}}}} commands respectively.
The {\hyperref[\detokenize{angle_amoeba::doc}]{\sphinxcrossref{\DUrole{doc}{angle\_style amoeba}}}} command includes the
\(U_{b\theta}\) bond\sphinxhyphen{}angle cross term, and the \(U_{UB}\) term
for a Urey\sphinxhyphen{}Bradley bond contribution between the I,K atoms in the IJK
angle.

\sphinxAtStartPar
The \(U_{pitorsion}\) term is computed by the {\hyperref[\detokenize{fix_amoeba_pitorsion::doc}]{\sphinxcrossref{\DUrole{doc}{fix
amoeba/pitorsion}}}} command.  It computes 6\sphinxhyphen{}body
interaction between a pair of bonded atoms which each have 2
additional bond partners.

\sphinxAtStartPar
The \(U_{bitorsion}\) term is computed by the {\hyperref[\detokenize{fix_amoeba_bitorsion::doc}]{\sphinxcrossref{\DUrole{doc}{fix
amoeba/bitorsion}}}} command.  It computes 5\sphinxhyphen{}body
interaction between two 4\sphinxhyphen{}body torsions (dihedrals) which overlap,
having 3 atoms in common.

\sphinxAtStartPar
These command doc pages have additional details on the terms they
compute:
\begin{itemize}
\item {} 
\sphinxAtStartPar
{\hyperref[\detokenize{pair_amoeba::doc}]{\sphinxcrossref{\DUrole{doc}{pair\_style amoeba or hippo}}}}

\item {} 
\sphinxAtStartPar
{\hyperref[\detokenize{bond_class2::doc}]{\sphinxcrossref{\DUrole{doc}{bond\_style class2}}}}

\item {} 
\sphinxAtStartPar
{\hyperref[\detokenize{angle_amoeba::doc}]{\sphinxcrossref{\DUrole{doc}{angle\_style amoeba}}}}

\item {} 
\sphinxAtStartPar
{\hyperref[\detokenize{dihedral_fourier::doc}]{\sphinxcrossref{\DUrole{doc}{dihedral\_style fourier}}}}

\item {} 
\sphinxAtStartPar
{\hyperref[\detokenize{improper_amoeba::doc}]{\sphinxcrossref{\DUrole{doc}{improper\_style amoeba}}}}

\item {} 
\sphinxAtStartPar
{\hyperref[\detokenize{fix_amoeba_pitorsion::doc}]{\sphinxcrossref{\DUrole{doc}{fix amoeba/pitorsion}}}}

\item {} 
\sphinxAtStartPar
{\hyperref[\detokenize{fix_amoeba_bitorsion::doc}]{\sphinxcrossref{\DUrole{doc}{fix amoeba/bitorsion}}}}

\end{itemize}


\bigskip\hrule\bigskip


\sphinxAtStartPar
To use the AMOEBA or HIPPO force fields in LAMMPS, use commands like
the following appropriately in your input script.  The only change
needed for AMOEBA vs HIPPO simulation is for the {\hyperref[\detokenize{pair_style::doc}]{\sphinxcrossref{\DUrole{doc}{pair\_style}}}} and {\hyperref[\detokenize{pair_coeff::doc}]{\sphinxcrossref{\DUrole{doc}{pair\_coeff}}}} commands, as shown
below.  See examples/amoeba for example input scripts for both AMOEBA
and HIPPO.

\begin{sphinxVerbatim}[commandchars=\\\{\}]
\PYG{k}{units}\PYG{+w}{              }\PYG{n}{real}\PYG{+w}{                           }\PYG{c}{\PYGZsh{} required}
\PYG{k}{atom\PYGZus{}style}\PYG{+w}{         }\PYG{n}{amoeba}
\PYG{k}{bond\PYGZus{}style}\PYG{+w}{         }\PYG{n}{class2}\PYG{+w}{                         }\PYG{c}{\PYGZsh{} CLASS2 package}
\PYG{k}{angle\PYGZus{}style}\PYG{+w}{        }\PYG{n}{amoeba}
\PYG{k}{dihedral\PYGZus{}style}\PYG{+w}{     }\PYG{n}{fourier}\PYG{+w}{                        }\PYG{c}{\PYGZsh{} EXTRA\PYGZhy{}MOLECULE package}
\PYG{k}{improper\PYGZus{}style}\PYG{+w}{     }\PYG{n}{amoeba}
\PYG{+w}{                                                  }\PYG{c}{\PYGZsh{} required per\PYGZhy{}atom data}
\PYG{k}{fix                }\PYG{n+nv+nvIdentifier}{amtype}\PYG{+w}{ }\PYG{n+nv+nvIdentifier}{all}\PYG{+w}{ }\PYG{n}{property}\PYG{o}{/}\PYG{n}{atom}\PYG{+w}{ }\PYG{n}{i\PYGZus{}amtype}\PYG{+w}{ }\PYG{n}{ghost}\PYG{+w}{ }\PYG{n}{yes}
\PYG{k}{fix                }\PYG{n+nv+nvIdentifier}{extra}\PYG{+w}{ }\PYG{n+nv+nvIdentifier}{all}\PYG{+w}{ }\PYG{n}{property}\PYG{o}{/}\PYG{n}{atom}\PYG{+w}{ }\PYG{o}{\PYGZam{}}
\PYG{+w}{                   }\PYG{n}{i\PYGZus{}amgroup}\PYG{+w}{ }\PYG{n}{i\PYGZus{}ired}\PYG{+w}{ }\PYG{n}{i\PYGZus{}xaxis}\PYG{+w}{ }\PYG{n}{i\PYGZus{}yaxis}\PYG{+w}{ }\PYG{n}{i\PYGZus{}zaxis}\PYG{+w}{ }\PYG{n}{d\PYGZus{}pval}\PYG{+w}{ }\PYG{n}{ghost}\PYG{+w}{ }\PYG{n}{yes}
\PYG{k}{fix                }\PYG{n+nv+nvIdentifier}{polaxe}\PYG{+w}{ }\PYG{n+nv+nvIdentifier}{all}\PYG{+w}{ }\PYG{n}{property}\PYG{o}{/}\PYG{n}{atom}\PYG{+w}{ }\PYG{n}{i\PYGZus{}polaxe}

\PYG{k}{fix                }\PYG{n+nv+nvIdentifier}{pit}\PYG{+w}{ }\PYG{n+nv+nvIdentifier}{all}\PYG{+w}{ }\PYG{n}{amoeba}\PYG{o}{/}\PYG{n}{pitorsion}\PYG{+w}{       }\PYG{c}{\PYGZsh{} PiTorsion terms in FF}
\PYG{k}{fix\PYGZus{}modify         }\PYG{n+nv+nvIdentifier}{pit}\PYG{+w}{ }\PYG{n}{energy}\PYG{+w}{ }\PYG{n}{yes}
\PYG{+w}{                                                  }\PYG{c}{\PYGZsh{} Bitorsion terms in FF}
\PYG{k}{fix                }\PYG{n+nv+nvIdentifier}{bit}\PYG{+w}{ }\PYG{n+nv+nvIdentifier}{all}\PYG{+w}{ }\PYG{n}{amoeba}\PYG{o}{/}\PYG{n}{bitorsion}\PYG{+w}{ }\PYG{n}{bitorsion.ubiquitin.data}
\PYG{k}{fix\PYGZus{}modify         }\PYG{n+nv+nvIdentifier}{bit}\PYG{+w}{ }\PYG{n}{energy}\PYG{+w}{ }\PYG{n}{yes}

\PYG{k}{read\PYGZus{}data}\PYG{+w}{          }\PYG{n}{data.ubiquitin}\PYG{+w}{ }\PYG{k}{fix }\PYG{n+nv+nvIdentifier}{amtype}\PYG{+w}{ }\PYG{n+nv+nvIdentifier}{NULL}\PYG{+w}{ }\PYG{l+s}{\PYGZdq{}}\PYG{l+s}{Tinker Types}\PYG{l+s}{\PYGZdq{}}\PYG{+w}{ }\PYG{o}{\PYGZam{}}
\PYG{+w}{                   }\PYG{k}{fix }\PYG{n+nv+nvIdentifier}{pit}\PYG{+w}{ }\PYG{l+s}{\PYGZdq{}}\PYG{l+s}{pitorsion types}\PYG{l+s}{\PYGZdq{}}\PYG{+w}{ }\PYG{l+s}{\PYGZdq{}}\PYG{l+s}{PiTorsion Coeffs}\PYG{l+s}{\PYGZdq{}}\PYG{+w}{ }\PYG{o}{\PYGZam{}}
\PYG{+w}{                   }\PYG{k}{fix }\PYG{n+nv+nvIdentifier}{pit}\PYG{+w}{ }\PYG{n+nv+nvIdentifier}{pitorsions}\PYG{+w}{ }\PYG{n}{PiTorsions}\PYG{+w}{ }\PYG{o}{\PYGZam{}}
\PYG{+w}{                   }\PYG{k}{fix }\PYG{n+nv+nvIdentifier}{bit}\PYG{+w}{ }\PYG{n+nv+nvIdentifier}{bitorsions}\PYG{+w}{ }\PYG{n}{BiTorsions}

\PYG{k}{pair\PYGZus{}style}\PYG{+w}{         }\PYG{n}{amoeba}\PYG{+w}{                          }\PYG{c}{\PYGZsh{} AMOEBA FF}
\PYG{k}{pair\PYGZus{}coeff}\PYG{+w}{         }\PYG{o}{*}\PYG{+w}{ }\PYG{o}{*}\PYG{+w}{ }\PYG{n}{amoeba\PYGZus{}ubiquitin.prm}\PYG{+w}{ }\PYG{n}{amoeba\PYGZus{}ubiquitin.key}

\PYG{k}{pair\PYGZus{}style}\PYG{+w}{         }\PYG{n}{hippo}\PYG{+w}{                           }\PYG{c}{\PYGZsh{} HIPPO FF}
\PYG{k}{pair\PYGZus{}coeff}\PYG{+w}{         }\PYG{o}{*}\PYG{+w}{ }\PYG{o}{*}\PYG{+w}{ }\PYG{n}{hippo\PYGZus{}water.prm}\PYG{+w}{ }\PYG{n}{hippo\PYGZus{}water.key}

\PYG{k}{special\PYGZus{}bonds}\PYG{+w}{      }\PYG{n}{lj}\PYG{o}{/}\PYG{n}{coul}\PYG{+w}{ }\PYG{l+m}{0.5}\PYG{+w}{ }\PYG{l+m}{0.5}\PYG{+w}{ }\PYG{l+m}{0.5}\PYG{+w}{ }\PYG{n}{one}\PYG{o}{/}\PYG{n}{five}\PYG{+w}{ }\PYG{n}{yes}\PYG{+w}{     }\PYG{c}{\PYGZsh{} 1\PYGZhy{}5 neighbors}
\end{sphinxVerbatim}

\sphinxAtStartPar
The data file read by the {\hyperref[\detokenize{read_data::doc}]{\sphinxcrossref{\DUrole{doc}{read\_data}}}} command should
be created by the tools/tinker/tinker2lmp.py conversion program
described below.  It will create a section in the data file with the
header “Tinker Types”.  A {\hyperref[\detokenize{fix_property_atom::doc}]{\sphinxcrossref{\DUrole{doc}{fix property/atom}}}}
command for the data must be specified before the read\_data command.
In the example above the fix ID is \sphinxstyleemphasis{amtype}.

\sphinxAtStartPar
Similarly, if the system you are simulating defines AMOEBA/HIPPO
pitorsion or bitorsion interactions, there will be entries in the data
file for those interactions.  They require a {\hyperref[\detokenize{fix_amoeba_pitorsion::doc}]{\sphinxcrossref{\DUrole{doc}{fix
amoeba/pitortion}}}} and {\hyperref[\detokenize{fix_amoeba_bitorsion::doc}]{\sphinxcrossref{\DUrole{doc}{fix
amoeba/bitorsion}}}} command be defined.  In the
example above, the IDs for these two fixes are \sphinxstyleemphasis{pit} and \sphinxstyleemphasis{bit}.

\sphinxAtStartPar
Of course, if the system being modeled does not have one or more of
the following \textendash{} bond, angle, dihedral, improper, pitorsion,
bitorsion interactions \textendash{} then the corresponding style and fix
commands above do not need to be used.  See the example scripts in
examples/amoeba for water systems as examples; they are simpler than
what is listed above.

\sphinxAtStartPar
The two {\hyperref[\detokenize{fix_property_atom::doc}]{\sphinxcrossref{\DUrole{doc}{fix property/atom}}}} commands with IDs
(in the example above) \sphinxstyleemphasis{extra} and \sphinxstyleemphasis{polaxe} are also needed to define
internal per\sphinxhyphen{}atom quantities used by the AMOEBA and HIPPO force
fields.

\sphinxAtStartPar
The {\hyperref[\detokenize{pair_coeff::doc}]{\sphinxcrossref{\DUrole{doc}{pair\_coeff}}}} command used for either the AMOEBA
or HIPPO force field takes two arguments for Tinker force field files,
namely a PRM and KEY file.  The keyfile can be specified as NULL and
default values for a various settings will be used.  Note that these 2
files are meant to allow use of native Tinker files as\sphinxhyphen{}is.  However
LAMMPS does not support all the options which can be included
in a Tinker PRM or KEY file.  See specifics below.

\sphinxAtStartPar
A {\hyperref[\detokenize{special_bonds::doc}]{\sphinxcrossref{\DUrole{doc}{special\_bonds}}}} command with the \sphinxstyleemphasis{one/five}
option is required, since the AMOEBA/HIPPO force fields define
weighting factors for not only 1\sphinxhyphen{}2, 1\sphinxhyphen{}3, 1\sphinxhyphen{}4 interactions, but also
1\sphinxhyphen{}5 interactions.  This command will trigger a per\sphinxhyphen{}atom list of 1\sphinxhyphen{}5
neighbors to be generated.  The AMOEBA and HIPPO force fields define
their own custom weighting factors for all the 1\sphinxhyphen{}2, 1\sphinxhyphen{}3, 1\sphinxhyphen{}4, 1\sphinxhyphen{}5
terms which in the Tinker PRM and KEY files; they can be different for
different terms in the force field.

\sphinxAtStartPar
In addition to the list above, these command doc pages have additional
details:
\begin{itemize}
\item {} 
\sphinxAtStartPar
{\hyperref[\detokenize{atom_style::doc}]{\sphinxcrossref{\DUrole{doc}{atom\_style amoeba}}}}

\item {} 
\sphinxAtStartPar
{\hyperref[\detokenize{fix_property_atom::doc}]{\sphinxcrossref{\DUrole{doc}{fix property/atom}}}}

\item {} 
\sphinxAtStartPar
{\hyperref[\detokenize{special_bonds::doc}]{\sphinxcrossref{\DUrole{doc}{special\_bonds}}}}

\end{itemize}


\bigskip\hrule\bigskip


\sphinxAtStartPar
Tinker PRM and KEY files

\sphinxAtStartPar
A Tinker PRM file is composed of sections, each of which has multiple
lines.  This is the list of PRM sections LAMMPS knows how to parse and
use.  Any other sections are skipped:
\begin{itemize}
\item {} 
\sphinxAtStartPar
Angle Bending Parameters

\item {} 
\sphinxAtStartPar
Atom Type Definitions

\item {} 
\sphinxAtStartPar
Atomic Multipole Parameters

\item {} 
\sphinxAtStartPar
Bond Stretching Parameters

\item {} 
\sphinxAtStartPar
Charge Penetration Parameters

\item {} 
\sphinxAtStartPar
Charge Transfer Parameters

\item {} 
\sphinxAtStartPar
Dipole Polarizability Parameters

\item {} 
\sphinxAtStartPar
Dispersion Parameters

\item {} 
\sphinxAtStartPar
Force Field Definition

\item {} 
\sphinxAtStartPar
Literature References

\item {} 
\sphinxAtStartPar
Out\sphinxhyphen{}of\sphinxhyphen{}Plane Bend Parameters

\item {} 
\sphinxAtStartPar
Pauli Repulsion Parameters

\item {} 
\sphinxAtStartPar
Pi\sphinxhyphen{}Torsion Parameters

\item {} 
\sphinxAtStartPar
Stretch\sphinxhyphen{}Bend Parameters

\item {} 
\sphinxAtStartPar
Torsion\sphinxhyphen{}Torsion Parameters

\item {} 
\sphinxAtStartPar
Torsional Parameters

\item {} 
\sphinxAtStartPar
Urey\sphinxhyphen{}Bradley Parameters

\item {} 
\sphinxAtStartPar
Van der Waals Pair Parameters

\item {} 
\sphinxAtStartPar
Van der Waals Parameters

\end{itemize}

\sphinxAtStartPar
A Tinker KEY file is composed of lines, each of which has a keyword
followed by zero or more parameters.  This is the list of keywords
LAMMPS knows how to parse and use in the same manner Tinker does.  Any
other keywords are skipped.  The value in parenthesis is the default
value for the keyword if it is not specified, or if the keyfile in the
{\hyperref[\detokenize{pair_coeff::doc}]{\sphinxcrossref{\DUrole{doc}{pair\_coeff}}}} command is specified as NULL:
\begin{itemize}
\item {} 
\sphinxAtStartPar
a\sphinxhyphen{}axis (0.0)

\item {} 
\sphinxAtStartPar
b\sphinxhyphen{}axis (0.0)

\item {} 
\sphinxAtStartPar
c\sphinxhyphen{}axis (0.0)

\item {} 
\sphinxAtStartPar
ctrn\sphinxhyphen{}cutoff (6.0)

\item {} 
\sphinxAtStartPar
ctrn\sphinxhyphen{}taper (0.9 * ctrn\sphinxhyphen{}cutoff)

\item {} 
\sphinxAtStartPar
cutoff

\item {} 
\sphinxAtStartPar
delta\sphinxhyphen{}halgren (0.07)

\item {} 
\sphinxAtStartPar
dewald (no long\sphinxhyphen{}range dispersion unless specified)

\item {} 
\sphinxAtStartPar
dewald\sphinxhyphen{}alpha (0.4)

\item {} 
\sphinxAtStartPar
dewald\sphinxhyphen{}cutoff (7.0)

\item {} 
\sphinxAtStartPar
dispersion\sphinxhyphen{}cutoff (9.0)

\item {} 
\sphinxAtStartPar
dispersion\sphinxhyphen{}taper (9.0 * dispersion\sphinxhyphen{}cutoff)

\item {} 
\sphinxAtStartPar
dpme\sphinxhyphen{}grid

\item {} 
\sphinxAtStartPar
dpme\sphinxhyphen{}order (4)

\item {} 
\sphinxAtStartPar
ewald (no long\sphinxhyphen{}range electrostatics unless specified)

\item {} 
\sphinxAtStartPar
ewald\sphinxhyphen{}alpha (0.4)

\item {} 
\sphinxAtStartPar
ewald\sphinxhyphen{}cutoff (7.0)

\item {} 
\sphinxAtStartPar
gamma\sphinxhyphen{}halgren (0.12)

\item {} 
\sphinxAtStartPar
mpole\sphinxhyphen{}cutoff (9.0)

\item {} 
\sphinxAtStartPar
mpole\sphinxhyphen{}taper (0.65 * mpole\sphinxhyphen{}cutoff)

\item {} 
\sphinxAtStartPar
pcg\sphinxhyphen{}guess (enabled by default)

\item {} 
\sphinxAtStartPar
pcg\sphinxhyphen{}noguess (disable pcg\sphinxhyphen{}guess if specified)

\item {} 
\sphinxAtStartPar
pcg\sphinxhyphen{}noprecond (disable pcg\sphinxhyphen{}precond if specified)

\item {} 
\sphinxAtStartPar
pcg\sphinxhyphen{}peek (1.0)

\item {} 
\sphinxAtStartPar
pcg\sphinxhyphen{}precond (enabled by default)

\item {} 
\sphinxAtStartPar
pewald\sphinxhyphen{}alpha (0.4)

\item {} 
\sphinxAtStartPar
pme\sphinxhyphen{}grid

\item {} 
\sphinxAtStartPar
pme\sphinxhyphen{}order (5)

\item {} 
\sphinxAtStartPar
polar\sphinxhyphen{}eps (1.0e\sphinxhyphen{}6)

\item {} 
\sphinxAtStartPar
polar\sphinxhyphen{}iter (100)

\item {} 
\sphinxAtStartPar
polar\sphinxhyphen{}predict (no prediction operation unless specified)

\item {} 
\sphinxAtStartPar
ppme\sphinxhyphen{}order (5)

\item {} 
\sphinxAtStartPar
repulsion\sphinxhyphen{}cutoff (6.0)

\item {} 
\sphinxAtStartPar
repulsion\sphinxhyphen{}taper (0.9 * repulsion\sphinxhyphen{}cutoff)

\item {} 
\sphinxAtStartPar
taper

\item {} 
\sphinxAtStartPar
usolve\sphinxhyphen{}cutoff (4.5)

\item {} 
\sphinxAtStartPar
usolve\sphinxhyphen{}diag (2.0)

\item {} 
\sphinxAtStartPar
vdw\sphinxhyphen{}cutoff (9.0)

\item {} 
\sphinxAtStartPar
vdw\sphinxhyphen{}taper (0.9 * vdw\sphinxhyphen{}cutoff)

\end{itemize}


\bigskip\hrule\bigskip


\sphinxAtStartPar
Tinker2lmp.py tool

\sphinxAtStartPar
This conversion tool is found in the tools/tinker directory.
As shown in examples/amoeba/README, these commands produce
the data files found in examples/amoeba, and also illustrate
all the options available to use with the tinker2lmp.py script:

\begin{sphinxVerbatim}[commandchars=\\\{\}]
python\PYG{+w}{ }tinker2lmp.py\PYG{+w}{ }\PYGZhy{}xyz\PYG{+w}{ }water\PYGZus{}dimer.xyz\PYG{+w}{ }\PYGZhy{}amoeba\PYG{+w}{ }amoeba\PYGZus{}water.prm\PYG{+w}{ }\PYGZhy{}data\PYG{+w}{ }data.water\PYGZus{}dimer.amoeba\PYG{+w}{                }\PYG{c+c1}{\PYGZsh{} AMOEBA non\PYGZhy{}periodic system}
python\PYG{+w}{ }tinker2lmp.py\PYG{+w}{ }\PYGZhy{}xyz\PYG{+w}{ }water\PYGZus{}dimer.xyz\PYG{+w}{ }\PYGZhy{}hippo\PYG{+w}{ }hippo\PYGZus{}water.prm\PYG{+w}{ }\PYGZhy{}data\PYG{+w}{ }data.water\PYGZus{}dimer.hippo\PYG{+w}{                   }\PYG{c+c1}{\PYGZsh{} HIPPO non\PYGZhy{}periodic system}
python\PYG{+w}{ }tinker2lmp.py\PYG{+w}{ }\PYGZhy{}xyz\PYG{+w}{ }water\PYGZus{}box.xyz\PYG{+w}{ }\PYGZhy{}amoeba\PYG{+w}{ }amoeba\PYGZus{}water.prm\PYG{+w}{ }\PYGZhy{}data\PYG{+w}{ }data.water\PYGZus{}box.amoeba\PYG{+w}{ }\PYGZhy{}pbc\PYG{+w}{ }\PYG{l+m}{18}.643\PYG{+w}{ }\PYG{l+m}{18}.643\PYG{+w}{ }\PYG{l+m}{18}.643\PYG{+w}{    }\PYG{c+c1}{\PYGZsh{} AMOEBA periodic system}
python\PYG{+w}{ }tinker2lmp.py\PYG{+w}{ }\PYGZhy{}xyz\PYG{+w}{ }water\PYGZus{}box.xyz\PYG{+w}{ }\PYGZhy{}hippo\PYG{+w}{ }hippo\PYGZus{}water.prm\PYG{+w}{ }\PYGZhy{}data\PYG{+w}{ }data.water\PYGZus{}box.hippo\PYG{+w}{ }\PYGZhy{}pbc\PYG{+w}{ }\PYG{l+m}{18}.643\PYG{+w}{ }\PYG{l+m}{18}.643\PYG{+w}{ }\PYG{l+m}{18}.643\PYG{+w}{       }\PYG{c+c1}{\PYGZsh{} HIPPO periodic system}
python\PYG{+w}{ }tinker2lmp.py\PYG{+w}{ }\PYGZhy{}xyz\PYG{+w}{ }ubiquitin.xyz\PYG{+w}{ }\PYGZhy{}amoeba\PYG{+w}{ }amoeba\PYGZus{}ubiquitin.prm\PYG{+w}{ }\PYGZhy{}data\PYG{+w}{ }data.ubiquitin.new\PYG{+w}{ }\PYGZhy{}pbc\PYG{+w}{ }\PYG{l+m}{54}.99\PYG{+w}{ }\PYG{l+m}{41}.91\PYG{+w}{ }\PYG{l+m}{41}.91\PYG{+w}{ }\PYGZhy{}bitorsion\PYG{+w}{ }bitorsion.ubiquitin.data.new\PYG{+w}{   }\PYG{c+c1}{\PYGZsh{} system with bitorsions}
\end{sphinxVerbatim}

\sphinxAtStartPar
Switches and their arguments may be specified in any order.

\sphinxAtStartPar
The \sphinxhyphen{}xyz switch is required and specifies an input XYZ file as an
argument.  The format of this file is an extended XYZ format defined
and used by Tinker for its input.  Example *.xyz files are in the
examples/amoeba directory.  The file lists the atoms in the system.
Each atom has the following information: Tinker species name (ignored
by LAMMPS), xyz coordinates, Tinker numeric type, and a list of atom
IDs the atom is bonded to.

\sphinxAtStartPar
Here is more information about the extended XYZ format defined and
used by Tinker, and links to programs that convert standard PDB files
to the extended XYZ format:
\begin{itemize}
\item {} 
\sphinxAtStartPar
\sphinxurl{https://openbabel.org/docs/current/FileFormats/Tinker\_XYZ\_format.html}

\item {} 
\sphinxAtStartPar
\sphinxurl{https://github.com/emleddin/pdbxyz-xyzpdb}

\item {} 
\sphinxAtStartPar
\sphinxurl{https://github.com/TinkerTools/tinker/blob/release/source/pdbxyz.f}

\end{itemize}

\sphinxAtStartPar
The \sphinxhyphen{}amoeba or \sphinxhyphen{}hippo switch is required.  It specifies an input
AMOEBA or HIPPO PRM force field file as an argument.  This should be
the same file used by the {\hyperref[\detokenize{pair_style::doc}]{\sphinxcrossref{\DUrole{doc}{pair\_style}}}} command in
the input script.

\sphinxAtStartPar
The \sphinxhyphen{}data switch is required.  It specifies an output file name for
the LAMMPS data file that will be produced.

\sphinxAtStartPar
For periodic systems, the \sphinxhyphen{}pbc switch is required.  It specifies the
periodic box size for each dimension (x,y,z).  For a Tinker simulation
these are specified in the KEY file.

\sphinxAtStartPar
The \sphinxhyphen{}bitorsion switch is only needed if the system contains Tinker
bitorsion interactions.  The data for each type of bitorsion
interaction will be written to the specified file, and read by the
{\hyperref[\detokenize{fix_amoeba_bitorsion::doc}]{\sphinxcrossref{\DUrole{doc}{fix amoeba/bitorsion}}}} command.  The data
includes 2d arrays of values to which splines are fit, and thus is not
compatible with the LAMMPS data file format.


\bigskip\hrule\bigskip

\phantomsection\label{\detokenize{Howto_amoeba:howto-ponder}}
\sphinxAtStartPar
\sphinxstylestrong{(Ponder)} Ponder, Wu, Ren, Pande, Chodera, Schnieders, Haque,  Mobley, Lambrecht, DiStasio Jr, M. Head\sphinxhyphen{}Gordon, Clark,  Johnson, T. Head\sphinxhyphen{}Gordon, J Phys Chem B, 114, 2549\sphinxhyphen{}2564 (2010).

\phantomsection\label{\detokenize{Howto_amoeba:howto-rackers}}
\sphinxAtStartPar
\sphinxstylestrong{(Rackers)} Rackers, Silva, Wang, Ponder, J Chem Theory Comput, 17, 7056\sphinxhyphen{}7084 (2021).

\phantomsection\label{\detokenize{Howto_amoeba:howto-ren}}
\sphinxAtStartPar
\sphinxstylestrong{(Ren)} Ren and Ponder, J Phys Chem B, 107, 5933 (2003).

\phantomsection\label{\detokenize{Howto_amoeba:howto-shi}}
\sphinxAtStartPar
\sphinxstylestrong{(Shi)} Shi, Xia, Zhang, Best, Wu, Ponder, Ren, J Chem Theory Comp, 9, 4046, 2013.

\sphinxstepscope


\subsection{TIP3P water model}
\label{\detokenize{Howto_tip3p:tip3p-water-model}}\label{\detokenize{Howto_tip3p::doc}}
\sphinxAtStartPar
The TIP3P water model as implemented in CHARMM {\hyperref[\detokenize{Howto_tip3p:howto-tip3p}]{\sphinxcrossref{\DUrole{std,std-ref}{(MacKerell)}}}} specifies a 3\sphinxhyphen{}site rigid water molecule with charges and
Lennard\sphinxhyphen{}Jones parameters assigned to each of the three atoms.

\sphinxAtStartPar
A suitable pair style with cutoff Coulomb would be:
\begin{itemize}
\item {} 
\sphinxAtStartPar
{\hyperref[\detokenize{pair_lj_cut_coul::doc}]{\sphinxcrossref{\DUrole{doc}{pair\_style lj/cut/coul/cut}}}}

\end{itemize}

\sphinxAtStartPar
or these commands for a long\sphinxhyphen{}range Coulomb model:
\begin{itemize}
\item {} 
\sphinxAtStartPar
{\hyperref[\detokenize{pair_lj_cut_coul::doc}]{\sphinxcrossref{\DUrole{doc}{pair\_style lj/cut/coul/long}}}}

\item {} 
\sphinxAtStartPar
{\hyperref[\detokenize{pair_fep_soft::doc}]{\sphinxcrossref{\DUrole{doc}{pair\_style lj/cut/coul/long/soft}}}}

\item {} 
\sphinxAtStartPar
{\hyperref[\detokenize{kspace_style::doc}]{\sphinxcrossref{\DUrole{doc}{kspace\_style pppm}}}}

\item {} 
\sphinxAtStartPar
{\hyperref[\detokenize{kspace_style::doc}]{\sphinxcrossref{\DUrole{doc}{kspace\_style pppm/disp}}}}

\end{itemize}

\sphinxAtStartPar
In LAMMPS the {\hyperref[\detokenize{fix_shake::doc}]{\sphinxcrossref{\DUrole{doc}{fix shake or fix rattle}}}} command can be
used to hold the two O\sphinxhyphen{}H bonds and the H\sphinxhyphen{}O\sphinxhyphen{}H angle rigid.  A bond style
of {\hyperref[\detokenize{bond_harmonic::doc}]{\sphinxcrossref{\DUrole{doc}{harmonic}}}} and an angle style of {\hyperref[\detokenize{angle_harmonic::doc}]{\sphinxcrossref{\DUrole{doc}{harmonic}}}} or {\hyperref[\detokenize{angle_charmm::doc}]{\sphinxcrossref{\DUrole{doc}{charmm}}}} should also be used.
In case of rigid bonds also bond style {\hyperref[\detokenize{bond_zero::doc}]{\sphinxcrossref{\DUrole{doc}{zero}}}} and angle
style {\hyperref[\detokenize{angle_zero::doc}]{\sphinxcrossref{\DUrole{doc}{zero}}}} can be used.

\sphinxAtStartPar
The table below lists the force field parameters (in real {\hyperref[\detokenize{units::doc}]{\sphinxcrossref{\DUrole{doc}{units}}}}) to for the water molecule atoms to run a rigid or flexible
TIP3P\sphinxhyphen{}CHARMM model with a cutoff, the original 1983 TIP3P model
{\hyperref[\detokenize{Howto_tip3p:jorgensen1}]{\sphinxcrossref{\DUrole{std,std-ref}{(Jorgensen)}}}}, or a TIP3P model with parameters
optimized for a long\sphinxhyphen{}range Coulomb solver (e.g. Ewald or PPPM in LAMMPS)
{\hyperref[\detokenize{Howto_tip3p:price1}]{\sphinxcrossref{\DUrole{std,std-ref}{(Price)}}}}.   The K values can be used if a flexible TIP3P
model (without fix shake) is desired, for rigid bonds/angles they are
ignored.


\begin{savenotes}\sphinxattablestart
\sphinxthistablewithglobalstyle
\centering
\begin{tabular}[t]{\X{38}{100}\X{22}{100}\X{20}{100}\X{20}{100}}
\sphinxtoprule
\sphinxstyletheadfamily 
\sphinxAtStartPar
Parameter
&\sphinxstyletheadfamily 
\sphinxAtStartPar
TIP3P\sphinxhyphen{}CHARMM
&\sphinxstyletheadfamily 
\sphinxAtStartPar
TIP3P (original)
&\sphinxstyletheadfamily 
\sphinxAtStartPar
TIP3P (Ewald)
\\
\sphinxmidrule
\sphinxtableatstartofbodyhook
\sphinxAtStartPar
O mass (amu)
&
\sphinxAtStartPar
15.9994
&
\sphinxAtStartPar
15.9994
&
\sphinxAtStartPar
15.9994
\\
\sphinxhline
\sphinxAtStartPar
H mass (amu)
&
\sphinxAtStartPar
1.008
&
\sphinxAtStartPar
1.008
&
\sphinxAtStartPar
1.008
\\
\sphinxhline
\sphinxAtStartPar
O charge (\(e\))
&
\sphinxAtStartPar
\sphinxhyphen{}0.834
&
\sphinxAtStartPar
\sphinxhyphen{}0.834
&
\sphinxAtStartPar
\sphinxhyphen{}0.834
\\
\sphinxhline
\sphinxAtStartPar
H charge (\(e\))
&
\sphinxAtStartPar
0.417
&
\sphinxAtStartPar
0.417
&
\sphinxAtStartPar
0.417
\\
\sphinxhline
\sphinxAtStartPar
LJ \(\epsilon\) of OO (kcal/mole)
&
\sphinxAtStartPar
0.1521
&
\sphinxAtStartPar
0.1521
&
\sphinxAtStartPar
0.1020
\\
\sphinxhline
\sphinxAtStartPar
LJ \(\sigma\) of OO (\(\AA\))
&
\sphinxAtStartPar
3.1507
&
\sphinxAtStartPar
3.1507
&
\sphinxAtStartPar
3.188
\\
\sphinxhline
\sphinxAtStartPar
LJ \(\epsilon\) of HH (kcal/mole)
&
\sphinxAtStartPar
0.0460
&
\sphinxAtStartPar
0.0
&
\sphinxAtStartPar
0.0
\\
\sphinxhline
\sphinxAtStartPar
LJ \(\sigma\) of HH (\(\AA\))
&
\sphinxAtStartPar
0.4
&
\sphinxAtStartPar
1.0
&
\sphinxAtStartPar
1.0
\\
\sphinxhline
\sphinxAtStartPar
LJ \(\epsilon\) of OH (kcal/mole)
&
\sphinxAtStartPar
0.0836
&
\sphinxAtStartPar
0.0
&
\sphinxAtStartPar
0.0
\\
\sphinxhline
\sphinxAtStartPar
LJ \(\sigma\) of OH (\(\AA\))
&
\sphinxAtStartPar
1.7753
&
\sphinxAtStartPar
1.0
&
\sphinxAtStartPar
1.0
\\
\sphinxhline
\sphinxAtStartPar
K of OH bond (kcal/mole/\(\AA^2\))
&
\sphinxAtStartPar
450
&
\sphinxAtStartPar
450
&
\sphinxAtStartPar
450
\\
\sphinxhline
\sphinxAtStartPar
\(r_0\) of OH bond (\(\AA\))
&
\sphinxAtStartPar
0.9572
&
\sphinxAtStartPar
0.9572
&
\sphinxAtStartPar
0.9572
\\
\sphinxhline
\sphinxAtStartPar
K of HOH angle (kcal/mole)
&
\sphinxAtStartPar
55.0
&
\sphinxAtStartPar
55.0
&
\sphinxAtStartPar
55.0
\\
\sphinxhline
\sphinxAtStartPar
\(\theta_0\) of HOH angle
&
\sphinxAtStartPar
104.52\(^{\circ}\)
&
\sphinxAtStartPar
104.52\(^{\circ}\)
&
\sphinxAtStartPar
104.52\(^{\circ}\)
\\
\sphinxbottomrule
\end{tabular}
\sphinxtableafterendhook\par
\sphinxattableend\end{savenotes}

\sphinxAtStartPar
Below is the code for a LAMMPS input file and a molecule file
(\sphinxcode{\sphinxupquote{tip3p.mol}}) of TIP3P water for use with the {\hyperref[\detokenize{molecule::doc}]{\sphinxcrossref{\DUrole{doc}{molecule command}}}} demonstrating how to set up a small bulk water system for
TIP3P with rigid bonds.

\begin{sphinxVerbatim}[commandchars=\\\{\}]
\PYG{k}{units}\PYG{+w}{ }\PYG{n}{real}
\PYG{k}{atom\PYGZus{}style}\PYG{+w}{ }\PYG{n}{full}
\PYG{k}{region }\PYG{n+nv+nvIdentifier}{box}\PYG{+w}{ }\PYG{n}{block}\PYG{+w}{ }\PYG{o}{\PYGZhy{}}\PYG{l+m}{5}\PYG{+w}{ }\PYG{l+m}{5}\PYG{+w}{ }\PYG{o}{\PYGZhy{}}\PYG{l+m}{5}\PYG{+w}{ }\PYG{l+m}{5}\PYG{+w}{ }\PYG{o}{\PYGZhy{}}\PYG{l+m}{5}\PYG{+w}{ }\PYG{l+m}{5}
\PYG{k}{create\PYGZus{}box }\PYG{n+nv+nvIdentifier}{2}\PYG{+w}{ }\PYG{n+nv+nvIdentifier}{box}\PYG{+w}{ }\PYG{n}{bond}\PYG{o}{/}\PYG{n}{types}\PYG{+w}{ }\PYG{l+m}{1}\PYG{+w}{ }\PYG{n}{angle}\PYG{o}{/}\PYG{n}{types}\PYG{+w}{ }\PYG{l+m}{1}\PYG{+w}{ }\PYG{o}{\PYGZam{}}
\PYG{+w}{            }\PYG{n}{extra}\PYG{o}{/}\PYG{n}{bond}\PYG{o}{/}\PYG{n}{per}\PYG{o}{/}\PYG{n}{atom}\PYG{+w}{ }\PYG{l+m}{2}\PYG{+w}{ }\PYG{n}{extra}\PYG{o}{/}\PYG{n}{angle}\PYG{o}{/}\PYG{n}{per}\PYG{o}{/}\PYG{n}{atom}\PYG{+w}{ }\PYG{l+m}{1}\PYG{+w}{ }\PYG{n}{extra}\PYG{o}{/}\PYG{n}{special}\PYG{o}{/}\PYG{n}{per}\PYG{o}{/}\PYG{n}{atom}\PYG{+w}{ }\PYG{l+m}{2}

\PYG{k}{mass}\PYG{+w}{ }\PYG{l+m}{1}\PYG{+w}{ }\PYG{l+m}{15.9994}
\PYG{k}{mass}\PYG{+w}{ }\PYG{l+m}{2}\PYG{+w}{ }\PYG{l+m}{1.008}

\PYG{k}{pair\PYGZus{}style}\PYG{+w}{ }\PYG{n}{lj}\PYG{o}{/}\PYG{n}{cut}\PYG{o}{/}\PYG{n}{coul}\PYG{o}{/}\PYG{n}{cut}\PYG{+w}{ }\PYG{l+m}{8.0}
\PYG{k}{pair\PYGZus{}coeff}\PYG{+w}{ }\PYG{l+m}{1}\PYG{+w}{ }\PYG{l+m}{1}\PYG{+w}{ }\PYG{l+m}{0.1521}\PYG{+w}{ }\PYG{l+m}{3.1507}
\PYG{k}{pair\PYGZus{}coeff}\PYG{+w}{ }\PYG{l+m}{2}\PYG{+w}{ }\PYG{l+m}{2}\PYG{+w}{ }\PYG{l+m}{0.0}\PYG{+w}{    }\PYG{l+m}{1.0}

\PYG{k}{bond\PYGZus{}style}\PYG{+w}{ }\PYG{n}{zero}
\PYG{k}{bond\PYGZus{}coeff}\PYG{+w}{ }\PYG{l+m}{1}\PYG{+w}{ }\PYG{l+m}{0.9574}

\PYG{k}{angle\PYGZus{}style}\PYG{+w}{ }\PYG{n}{zero}
\PYG{k}{angle\PYGZus{}coeff}\PYG{+w}{ }\PYG{l+m}{1}\PYG{+w}{ }\PYG{l+m}{104.52}

\PYG{k}{molecule}\PYG{+w}{ }\PYG{n}{water}\PYG{+w}{ }\PYG{n}{tip3p.mol}
\PYG{k}{create\PYGZus{}atoms}\PYG{+w}{ }\PYG{l+m}{0}\PYG{+w}{ }\PYG{n}{random}\PYG{+w}{ }\PYG{l+m}{33}\PYG{+w}{ }\PYG{l+m}{34564}\PYG{+w}{ }\PYG{n}{NULL}\PYG{+w}{ }\PYG{n}{mol}\PYG{+w}{ }\PYG{n}{water}\PYG{+w}{ }\PYG{l+m}{25367}\PYG{+w}{ }\PYG{n}{overlap}\PYG{+w}{ }\PYG{l+m}{1.33}

\PYG{k}{fix }\PYG{n+nv+nvIdentifier}{rigid}\PYG{+w}{ }\PYG{n+nv+nvIdentifier}{all}\PYG{+w}{ }\PYG{n}{shake}\PYG{+w}{ }\PYG{l+m}{0.001}\PYG{+w}{ }\PYG{l+m}{10}\PYG{+w}{ }\PYG{l+m}{10000}\PYG{+w}{ }\PYG{n}{b}\PYG{+w}{ }\PYG{l+m}{1}\PYG{+w}{ }\PYG{n}{a}\PYG{+w}{ }\PYG{l+m}{1}
\PYG{k}{minimize}\PYG{+w}{ }\PYG{l+m}{0.0}\PYG{+w}{ }\PYG{l+m}{0.0}\PYG{+w}{ }\PYG{l+m}{1000}\PYG{+w}{ }\PYG{l+m}{10000}

\PYG{k}{reset\PYGZus{}timestep}\PYG{+w}{ }\PYG{l+m}{0}
\PYG{k}{timestep}\PYG{+w}{ }\PYG{l+m}{1.0}
\PYG{k}{velocity }\PYG{n+nv+nvIdentifier}{all}\PYG{+w}{ }\PYG{n}{create}\PYG{+w}{ }\PYG{l+m}{300.0}\PYG{+w}{ }\PYG{l+m}{5463576}
\PYG{k}{fix }\PYG{n+nv+nvIdentifier}{integrate}\PYG{+w}{ }\PYG{n+nv+nvIdentifier}{all}\PYG{+w}{ }\PYG{n}{nvt}\PYG{+w}{ }\PYG{n}{temp}\PYG{+w}{ }\PYG{l+m}{300}\PYG{+w}{ }\PYG{l+m}{300}\PYG{+w}{ }\PYG{l+m}{100.0}

\PYG{k}{thermo\PYGZus{}style}\PYG{+w}{ }\PYG{n}{custom}\PYG{+w}{ }\PYG{n}{step}\PYG{+w}{ }\PYG{n}{temp}\PYG{+w}{ }\PYG{n}{press}\PYG{+w}{ }\PYG{n}{etotal}\PYG{+w}{ }\PYG{n}{pe}

\PYG{k}{thermo}\PYG{+w}{ }\PYG{l+m}{1000}
\PYG{k}{run}\PYG{+w}{ }\PYG{l+m}{20000}
\PYG{k}{write\PYGZus{}data }\PYG{n+nv+nvIdentifier}{tip3p.data}\PYG{+w}{ }\PYG{n}{nocoeff}
\end{sphinxVerbatim}

\def\sphinxLiteralBlockLabel{\label{\detokenize{Howto_tip3p:tip3p-molecule}}}
\begin{sphinxVerbatim}[commandchars=\\\{\}]
\PYGZsh{} Water molecule. TIP3P geometry

3 atoms
2 bonds
1 angles

Coords

1    0.00000  \PYGZhy{}0.06556   0.00000
2    0.75695   0.52032   0.00000
3   \PYGZhy{}0.75695   0.52032   0.00000

Types

1        1   \PYGZsh{} O
2        2   \PYGZsh{} H
3        2   \PYGZsh{} H

Charges

1       \PYGZhy{}0.834
2        0.417
3        0.417

Bonds

1   1      1      2
2   1      1      3

Angles

1   1      2      1      3

Shake Flags

1 1
2 1
3 1

Shake Atoms

1 1 2 3
2 1 2 3
3 1 2 3

Shake Bond Types

1 1 1 1
2 1 1 1
3 1 1 1

Special Bond Counts

1 2 0 0
2 1 1 0
3 1 1 0

Special Bonds

1 2 3
2 1 3
3 1 2
\end{sphinxVerbatim}

\sphinxAtStartPar
Wikipedia also has a nice article on \sphinxhref{https://en.wikipedia.org/wiki/Water\_model}{water models}.


\bigskip\hrule\bigskip

\phantomsection\label{\detokenize{Howto_tip3p:howto-tip3p}}
\sphinxAtStartPar
\sphinxstylestrong{(MacKerell)} MacKerell, Bashford, Bellott, Dunbrack, Evanseck, Field,
Fischer, Gao, Guo, Ha, et al, J Phys Chem, 102, 3586 (1998).

\phantomsection\label{\detokenize{Howto_tip3p:jorgensen1}}
\sphinxAtStartPar
\sphinxstylestrong{(Jorgensen)} Jorgensen, Chandrasekhar, Madura, Impey, Klein, J Chem
Phys, 79, 926 (1983).

\phantomsection\label{\detokenize{Howto_tip3p:price1}}
\sphinxAtStartPar
\sphinxstylestrong{(Price)} Price and Brooks, J Chem Phys, 121, 10096 (2004).

\sphinxstepscope


\subsection{TIP4P water model}
\label{\detokenize{Howto_tip4p:tip4p-water-model}}\label{\detokenize{Howto_tip4p::doc}}
\sphinxAtStartPar
The four\sphinxhyphen{}point TIP4P rigid water model extends the traditional
{\hyperref[\detokenize{Howto_tip3p::doc}]{\sphinxcrossref{\DUrole{doc}{three\sphinxhyphen{}point TIP3P}}}} model by adding an additional
site M, usually massless, where the charge associated with the oxygen
atom is placed.  This site M is located at a fixed distance away from
the oxygen along the bisector of the HOH bond angle.  A bond style of
{\hyperref[\detokenize{bond_harmonic::doc}]{\sphinxcrossref{\DUrole{doc}{harmonic}}}} and an angle style of {\hyperref[\detokenize{angle_harmonic::doc}]{\sphinxcrossref{\DUrole{doc}{harmonic}}}} or {\hyperref[\detokenize{angle_charmm::doc}]{\sphinxcrossref{\DUrole{doc}{charmm}}}} should also be used.
In case of rigid bonds also bond style {\hyperref[\detokenize{bond_zero::doc}]{\sphinxcrossref{\DUrole{doc}{zero}}}} and angle
style {\hyperref[\detokenize{angle_zero::doc}]{\sphinxcrossref{\DUrole{doc}{zero}}}} can be used.

\sphinxAtStartPar
There are two ways to implement TIP4P water in LAMMPS:
\begin{enumerate}
\sphinxsetlistlabels{\arabic}{enumi}{enumii}{}{.}%
\item {} 
\sphinxAtStartPar
Use a specially written pair style that uses the {\hyperref[\detokenize{Howto_tip3p:tip3p-molecule}]{\sphinxcrossref{\DUrole{std,std-ref}{TIP3P geometry}}}} without the point M. The point M location is then
implicitly derived from the other atoms or each water molecule and
used during the force computation.  The forces on M are then
projected on the oxygen and the two hydrogen atoms.  This is
computationally very efficient, but the charge distribution in space
is only correct within the tip4p labeled styles.  So all other
computations using charges will “see” the negative charge incorrectly
on the oxygen atom.

\sphinxAtStartPar
This can be done with the following pair styles for Coulomb with a cutoff:
\begin{itemize}
\item {} 
\sphinxAtStartPar
{\hyperref[\detokenize{pair_coul::doc}]{\sphinxcrossref{\DUrole{doc}{pair\_style tip4p/cut}}}}

\item {} 
\sphinxAtStartPar
{\hyperref[\detokenize{pair_lj_cut_tip4p::doc}]{\sphinxcrossref{\DUrole{doc}{pair\_style lj/cut/tip4p/cut}}}}

\end{itemize}

\sphinxAtStartPar
or these commands for a long\sphinxhyphen{}range Coulomb treatment:
\begin{itemize}
\item {} 
\sphinxAtStartPar
{\hyperref[\detokenize{pair_coul::doc}]{\sphinxcrossref{\DUrole{doc}{pair\_style tip4p/long}}}}

\item {} 
\sphinxAtStartPar
{\hyperref[\detokenize{pair_lj_cut_tip4p::doc}]{\sphinxcrossref{\DUrole{doc}{pair\_style lj/cut/tip4p/long}}}}

\item {} 
\sphinxAtStartPar
{\hyperref[\detokenize{pair_lj_long::doc}]{\sphinxcrossref{\DUrole{doc}{pair\_style lj/long/tip4p/long}}}}

\item {} 
\sphinxAtStartPar
{\hyperref[\detokenize{pair_fep_soft::doc}]{\sphinxcrossref{\DUrole{doc}{pair\_style tip4p/long/soft}}}}

\item {} 
\sphinxAtStartPar
{\hyperref[\detokenize{pair_fep_soft::doc}]{\sphinxcrossref{\DUrole{doc}{pair\_style lj/cut/tip4p/long/soft}}}}

\item {} 
\sphinxAtStartPar
{\hyperref[\detokenize{kspace_style::doc}]{\sphinxcrossref{\DUrole{doc}{kspace\_style pppm/tip4p}}}}

\item {} 
\sphinxAtStartPar
{\hyperref[\detokenize{kspace_style::doc}]{\sphinxcrossref{\DUrole{doc}{kspace\_style pppm/disp/tip4p}}}}

\end{itemize}

\sphinxAtStartPar
The bond lengths and bond angles should be held fixed using the
{\hyperref[\detokenize{fix_shake::doc}]{\sphinxcrossref{\DUrole{doc}{fix shake}}}} or {\hyperref[\detokenize{fix_shake::doc}]{\sphinxcrossref{\DUrole{doc}{fix rattle}}}} command,
unless a parameterization for a flexible TIP4P model is used.  The
parameter sets listed below are all for rigid TIP4P model variants and
thus the bond and angle force constants are not used and can be set to
any legal value; only equilibrium length and angle are used.

\item {} 
\sphinxAtStartPar
Use an {\hyperref[\detokenize{Howto_tip4p:tip4p-molecule}]{\sphinxcrossref{\DUrole{std,std-ref}{explicit 4 point TIP4P geometry}}}} where
the oxygen atom carries no charge and the M point no Lennard\sphinxhyphen{}Jones
interactions.  Since {\hyperref[\detokenize{fix_shake::doc}]{\sphinxcrossref{\DUrole{doc}{fix shake}}}} or {\hyperref[\detokenize{fix_shake::doc}]{\sphinxcrossref{\DUrole{doc}{fix rattle}}}} may not be applied to this kind of geometry, {\hyperref[\detokenize{fix_rigid::doc}]{\sphinxcrossref{\DUrole{doc}{fix
rigid or fix rigid/small}}}} or its thermostatted variants
are required to maintain a rigid geometry.  This avoids some of the
issues with respect to analysis and non\sphinxhyphen{}tip4p styles, but it is a
more costly force computation (more atoms in the same volume and thus
more neighbors in the neighbor lists) and requires a much shorter
timestep for stable integration of the rigid body motion.  Since no
bonds or angles are required, they do not need to be defined and atom
style charge would be sufficient for a bulk TIP4P water system.  In
order to avoid that LAMMPS produces an error due to the massless M
site a tiny non\sphinxhyphen{}zero mass needs to be assigned.

\end{enumerate}

\sphinxAtStartPar
The table below lists the force field parameters (in real {\hyperref[\detokenize{units::doc}]{\sphinxcrossref{\DUrole{doc}{units}}}}) to for a selection of popular variants of the TIP4P model.
There is the rigid TIP4P model with a cutoff {\hyperref[\detokenize{Howto_tip4p:jorgensen5}]{\sphinxcrossref{\DUrole{std,std-ref}{(Jorgensen)}}}}, the TIP4/Ice model {\hyperref[\detokenize{Howto_tip4p:abascal1}]{\sphinxcrossref{\DUrole{std,std-ref}{(Abascal1)}}}}, the
TIP4P/2005 model {\hyperref[\detokenize{Howto_tip4p:abascal2}]{\sphinxcrossref{\DUrole{std,std-ref}{(Abascal2)}}}} and a version of TIP4P
parameters adjusted for use with a long\sphinxhyphen{}range Coulombic solver
(e.g. Ewald or PPPM in LAMMPS).  Note that for implicit TIP4P models the
OM distance is specified in the {\hyperref[\detokenize{pair_style::doc}]{\sphinxcrossref{\DUrole{doc}{pair\_style}}}} command,
not as part of the pair coefficients.


\begin{savenotes}\sphinxattablestart
\sphinxthistablewithglobalstyle
\centering
\begin{tabular}[t]{\X{36}{100}\X{19}{100}\X{13}{100}\X{15}{100}\X{17}{100}}
\sphinxtoprule
\sphinxstyletheadfamily 
\sphinxAtStartPar
Parameter
&\sphinxstyletheadfamily 
\sphinxAtStartPar
TIP4P (original)
&\sphinxstyletheadfamily 
\sphinxAtStartPar
TIP4P/Ice
&\sphinxstyletheadfamily 
\sphinxAtStartPar
TIP4P/2005
&\sphinxstyletheadfamily 
\sphinxAtStartPar
TIP4P (Ewald)
\\
\sphinxmidrule
\sphinxtableatstartofbodyhook
\sphinxAtStartPar
O mass (amu)
&
\sphinxAtStartPar
15.9994
&
\sphinxAtStartPar
15.9994
&
\sphinxAtStartPar
15.9994
&
\sphinxAtStartPar
15.9994
\\
\sphinxhline
\sphinxAtStartPar
H mass (amu)
&
\sphinxAtStartPar
1.008
&
\sphinxAtStartPar
1.008
&
\sphinxAtStartPar
1.008
&
\sphinxAtStartPar
1.008
\\
\sphinxhline
\sphinxAtStartPar
O or M charge (\(e\))
&
\sphinxAtStartPar
\sphinxhyphen{}1.040
&
\sphinxAtStartPar
\sphinxhyphen{}1.1794
&
\sphinxAtStartPar
\sphinxhyphen{}1.1128
&
\sphinxAtStartPar
\sphinxhyphen{}1.04844
\\
\sphinxhline
\sphinxAtStartPar
H charge (\(e\))
&
\sphinxAtStartPar
0.520
&
\sphinxAtStartPar
0.5897
&
\sphinxAtStartPar
0.5564
&
\sphinxAtStartPar
0.52422
\\
\sphinxhline
\sphinxAtStartPar
LJ \(\epsilon\) of OO (kcal/mole)
&
\sphinxAtStartPar
0.1550
&
\sphinxAtStartPar
0.21084
&
\sphinxAtStartPar
0.1852
&
\sphinxAtStartPar
0.16275
\\
\sphinxhline
\sphinxAtStartPar
LJ \(\sigma\) of OO (\(\AA\))
&
\sphinxAtStartPar
3.1536
&
\sphinxAtStartPar
3.1668
&
\sphinxAtStartPar
3.1589
&
\sphinxAtStartPar
3.16435
\\
\sphinxhline
\sphinxAtStartPar
LJ \(\epsilon\) of HH, MM, OH, OM, HM (kcal/mole)
&
\sphinxAtStartPar
0.0
&
\sphinxAtStartPar
0.0
&
\sphinxAtStartPar
0.0
&
\sphinxAtStartPar
0.0
\\
\sphinxhline
\sphinxAtStartPar
LJ \(\sigma\) of HH, MM, OH, OM, HM (\(\AA\))
&
\sphinxAtStartPar
1.0
&
\sphinxAtStartPar
1.0
&
\sphinxAtStartPar
1.0
&
\sphinxAtStartPar
1.0
\\
\sphinxhline
\sphinxAtStartPar
\(r_0\) of OH bond (\(\AA\))
&
\sphinxAtStartPar
0.9572
&
\sphinxAtStartPar
0.9572
&
\sphinxAtStartPar
0.9572
&
\sphinxAtStartPar
0.9572
\\
\sphinxhline
\sphinxAtStartPar
\(\theta_0\) of HOH angle
&
\sphinxAtStartPar
104.52\(^{\circ}\)
&
\sphinxAtStartPar
104.52\(^{\circ}\)
&
\sphinxAtStartPar
104.52\(^{\circ}\)
&
\sphinxAtStartPar
104.52\(^{\circ}\)
\\
\sphinxhline
\sphinxAtStartPar
OM distance (\(\AA\))
&
\sphinxAtStartPar
0.15
&
\sphinxAtStartPar
0.1577
&
\sphinxAtStartPar
0.1546
&
\sphinxAtStartPar
0.1250
\\
\sphinxbottomrule
\end{tabular}
\sphinxtableafterendhook\par
\sphinxattableend\end{savenotes}

\sphinxAtStartPar
Note that the when using the TIP4P pair style, the neighbor list cutoff
for Coulomb interactions is effectively extended by a distance 2 * (OM
distance), to account for the offset distance of the fictitious charges
on O atoms in water molecules.  Thus it is typically best in an
efficiency sense to use a LJ cutoff \textgreater{}= Coulomb cutoff + 2*(OM
distance), to shrink the size of the neighbor list.  This leads to
slightly larger cost for the long\sphinxhyphen{}range calculation, so you can test the
trade\sphinxhyphen{}off for your model.  The OM distance and the LJ and Coulombic
cutoffs are set in the {\hyperref[\detokenize{pair_lj_cut_tip4p::doc}]{\sphinxcrossref{\DUrole{doc}{pair\_style lj/cut/tip4p/long}}}} command.

\sphinxAtStartPar
Below is the code for a LAMMPS input file using the implicit method and
the {\hyperref[\detokenize{Howto_tip3p:tip3p-molecule}]{\sphinxcrossref{\DUrole{std,std-ref}{TIP3P molecule file}}}}.  Because the TIP4P
charges are different from TIP3P they need to be reset (or the molecule
file changed):

\begin{sphinxVerbatim}[commandchars=\\\{\}]
\PYG{k}{units}\PYG{+w}{ }\PYG{n}{real}
\PYG{k}{atom\PYGZus{}style}\PYG{+w}{ }\PYG{n}{full}
\PYG{k}{region }\PYG{n+nv+nvIdentifier}{box}\PYG{+w}{ }\PYG{n}{block}\PYG{+w}{ }\PYG{o}{\PYGZhy{}}\PYG{l+m}{5}\PYG{+w}{ }\PYG{l+m}{5}\PYG{+w}{ }\PYG{o}{\PYGZhy{}}\PYG{l+m}{5}\PYG{+w}{ }\PYG{l+m}{5}\PYG{+w}{ }\PYG{o}{\PYGZhy{}}\PYG{l+m}{5}\PYG{+w}{ }\PYG{l+m}{5}
\PYG{k}{create\PYGZus{}box }\PYG{n+nv+nvIdentifier}{2}\PYG{+w}{ }\PYG{n+nv+nvIdentifier}{box}\PYG{+w}{ }\PYG{n}{bond}\PYG{o}{/}\PYG{n}{types}\PYG{+w}{ }\PYG{l+m}{1}\PYG{+w}{ }\PYG{n}{angle}\PYG{o}{/}\PYG{n}{types}\PYG{+w}{ }\PYG{l+m}{1}\PYG{+w}{ }\PYG{o}{\PYGZam{}}
\PYG{+w}{            }\PYG{n}{extra}\PYG{o}{/}\PYG{n}{bond}\PYG{o}{/}\PYG{n}{per}\PYG{o}{/}\PYG{n}{atom}\PYG{+w}{ }\PYG{l+m}{2}\PYG{+w}{ }\PYG{n}{extra}\PYG{o}{/}\PYG{n}{angle}\PYG{o}{/}\PYG{n}{per}\PYG{o}{/}\PYG{n}{atom}\PYG{+w}{ }\PYG{l+m}{1}\PYG{+w}{ }\PYG{n}{extra}\PYG{o}{/}\PYG{n}{special}\PYG{o}{/}\PYG{n}{per}\PYG{o}{/}\PYG{n}{atom}\PYG{+w}{ }\PYG{l+m}{2}

\PYG{k}{mass}\PYG{+w}{ }\PYG{l+m}{1}\PYG{+w}{ }\PYG{l+m}{15.9994}
\PYG{k}{mass}\PYG{+w}{ }\PYG{l+m}{2}\PYG{+w}{ }\PYG{l+m}{1.008}

\PYG{k}{pair\PYGZus{}style}\PYG{+w}{ }\PYG{n}{lj}\PYG{o}{/}\PYG{n}{cut}\PYG{o}{/}\PYG{n}{tip4p}\PYG{o}{/}\PYG{n}{cut}\PYG{+w}{ }\PYG{l+m}{1}\PYG{+w}{ }\PYG{l+m}{2}\PYG{+w}{ }\PYG{l+m}{1}\PYG{+w}{ }\PYG{l+m}{1}\PYG{+w}{ }\PYG{l+m}{0.15}\PYG{+w}{ }\PYG{l+m}{8.0}
\PYG{k}{pair\PYGZus{}coeff}\PYG{+w}{ }\PYG{l+m}{1}\PYG{+w}{ }\PYG{l+m}{1}\PYG{+w}{ }\PYG{l+m}{0.1550}\PYG{+w}{ }\PYG{l+m}{3.1536}
\PYG{k}{pair\PYGZus{}coeff}\PYG{+w}{ }\PYG{l+m}{2}\PYG{+w}{ }\PYG{l+m}{2}\PYG{+w}{ }\PYG{l+m}{0.0}\PYG{+w}{    }\PYG{l+m}{1.0}

\PYG{k}{bond\PYGZus{}style}\PYG{+w}{ }\PYG{n}{zero}
\PYG{k}{bond\PYGZus{}coeff}\PYG{+w}{ }\PYG{l+m}{1}\PYG{+w}{ }\PYG{l+m}{0.9574}

\PYG{k}{angle\PYGZus{}style}\PYG{+w}{ }\PYG{n}{zero}
\PYG{k}{angle\PYGZus{}coeff}\PYG{+w}{ }\PYG{l+m}{1}\PYG{+w}{ }\PYG{l+m}{104.52}

\PYG{k}{molecule}\PYG{+w}{ }\PYG{n}{water}\PYG{+w}{ }\PYG{n}{tip3p.mol}\PYG{+w}{  }\PYG{c}{\PYGZsh{} this uses the TIP3P geometry}
\PYG{k}{create\PYGZus{}atoms}\PYG{+w}{ }\PYG{l+m}{0}\PYG{+w}{ }\PYG{n}{random}\PYG{+w}{ }\PYG{l+m}{33}\PYG{+w}{ }\PYG{l+m}{34564}\PYG{+w}{ }\PYG{n}{NULL}\PYG{+w}{ }\PYG{n}{mol}\PYG{+w}{ }\PYG{n}{water}\PYG{+w}{ }\PYG{l+m}{25367}\PYG{+w}{ }\PYG{n}{overlap}\PYG{+w}{ }\PYG{l+m}{1.33}
\PYG{c}{\PYGZsh{} must change charges for TIP4P}
\PYG{k}{set}\PYG{+w}{ }\PYG{n}{type}\PYG{+w}{ }\PYG{l+m}{1}\PYG{+w}{ }\PYG{n}{charge}\PYG{+w}{ }\PYG{o}{\PYGZhy{}}\PYG{l+m}{1.040}
\PYG{k}{set}\PYG{+w}{ }\PYG{n}{type}\PYG{+w}{ }\PYG{l+m}{2}\PYG{+w}{ }\PYG{n}{charge}\PYG{+w}{  }\PYG{l+m}{0.520}

\PYG{k}{fix }\PYG{n+nv+nvIdentifier}{rigid}\PYG{+w}{ }\PYG{n+nv+nvIdentifier}{all}\PYG{+w}{ }\PYG{n}{shake}\PYG{+w}{ }\PYG{l+m}{0.001}\PYG{+w}{ }\PYG{l+m}{10}\PYG{+w}{ }\PYG{l+m}{10000}\PYG{+w}{ }\PYG{n}{b}\PYG{+w}{ }\PYG{l+m}{1}\PYG{+w}{ }\PYG{n}{a}\PYG{+w}{ }\PYG{l+m}{1}
\PYG{k}{minimize}\PYG{+w}{ }\PYG{l+m}{0.0}\PYG{+w}{ }\PYG{l+m}{0.0}\PYG{+w}{ }\PYG{l+m}{1000}\PYG{+w}{ }\PYG{l+m}{10000}

\PYG{k}{reset\PYGZus{}timestep}\PYG{+w}{ }\PYG{l+m}{0}
\PYG{k}{timestep}\PYG{+w}{ }\PYG{l+m}{1.0}
\PYG{k}{velocity }\PYG{n+nv+nvIdentifier}{all}\PYG{+w}{ }\PYG{n}{create}\PYG{+w}{ }\PYG{l+m}{300.0}\PYG{+w}{ }\PYG{l+m}{5463576}
\PYG{k}{fix }\PYG{n+nv+nvIdentifier}{integrate}\PYG{+w}{ }\PYG{n+nv+nvIdentifier}{all}\PYG{+w}{ }\PYG{n}{nvt}\PYG{+w}{ }\PYG{n}{temp}\PYG{+w}{ }\PYG{l+m}{300}\PYG{+w}{ }\PYG{l+m}{300}\PYG{+w}{ }\PYG{l+m}{100.0}

\PYG{k}{thermo\PYGZus{}style}\PYG{+w}{ }\PYG{n}{custom}\PYG{+w}{ }\PYG{n}{step}\PYG{+w}{ }\PYG{n}{temp}\PYG{+w}{ }\PYG{n}{press}\PYG{+w}{ }\PYG{n}{etotal}\PYG{+w}{ }\PYG{n}{pe}

\PYG{k}{thermo}\PYG{+w}{ }\PYG{l+m}{1000}
\PYG{k}{run}\PYG{+w}{ }\PYG{l+m}{20000}
\PYG{k}{write\PYGZus{}data }\PYG{n+nv+nvIdentifier}{tip4p\PYGZhy{}implicit.data}\PYG{+w}{ }\PYG{n}{nocoeff}
\end{sphinxVerbatim}

\sphinxAtStartPar
Below is the code for a LAMMPS input file using the explicit method and
a TIP4P molecule file.  Because of using {\hyperref[\detokenize{fix_rigid::doc}]{\sphinxcrossref{\DUrole{doc}{fix rigid/small}}}} no bonds need to be defined and thus no extra storage needs
to be reserved for them, but we need to either switch to atom style full
or use {\hyperref[\detokenize{fix_property_atom::doc}]{\sphinxcrossref{\DUrole{doc}{fix property/atom mol}}}} so that fix
rigid/small can identify rigid bodies by their molecule ID.  Also a
{\hyperref[\detokenize{neigh_modify::doc}]{\sphinxcrossref{\DUrole{doc}{neigh\_modify exclude}}}} command is added to exclude
computing intramolecular non\sphinxhyphen{}bonded interactions, since those are
removed by the rigid fix anyway:

\begin{sphinxVerbatim}[commandchars=\\\{\}]
\PYG{k}{units}\PYG{+w}{ }\PYG{n}{real}
\PYG{k}{atom\PYGZus{}style}\PYG{+w}{ }\PYG{n}{charge}
\PYG{k}{atom\PYGZus{}modify}\PYG{+w}{ }\PYG{n}{map}\PYG{+w}{ }\PYG{n}{array}
\PYG{k}{region }\PYG{n+nv+nvIdentifier}{box}\PYG{+w}{ }\PYG{n}{block}\PYG{+w}{ }\PYG{o}{\PYGZhy{}}\PYG{l+m}{5}\PYG{+w}{ }\PYG{l+m}{5}\PYG{+w}{ }\PYG{o}{\PYGZhy{}}\PYG{l+m}{5}\PYG{+w}{ }\PYG{l+m}{5}\PYG{+w}{ }\PYG{o}{\PYGZhy{}}\PYG{l+m}{5}\PYG{+w}{ }\PYG{l+m}{5}
\PYG{k}{create\PYGZus{}box }\PYG{n+nv+nvIdentifier}{3}\PYG{+w}{ }\PYG{n+nv+nvIdentifier}{box}

\PYG{k}{mass}\PYG{+w}{ }\PYG{l+m}{1}\PYG{+w}{ }\PYG{l+m}{15.9994}
\PYG{k}{mass}\PYG{+w}{ }\PYG{l+m}{2}\PYG{+w}{ }\PYG{l+m}{1.008}
\PYG{k}{mass}\PYG{+w}{ }\PYG{l+m}{3}\PYG{+w}{ }\PYG{l+m}{1.0e\PYGZhy{}100}

\PYG{k}{pair\PYGZus{}style}\PYG{+w}{ }\PYG{n}{lj}\PYG{o}{/}\PYG{n}{cut}\PYG{o}{/}\PYG{n}{coul}\PYG{o}{/}\PYG{n}{cut}\PYG{+w}{ }\PYG{l+m}{8.0}
\PYG{k}{pair\PYGZus{}coeff}\PYG{+w}{ }\PYG{l+m}{1}\PYG{+w}{ }\PYG{l+m}{1}\PYG{+w}{ }\PYG{l+m}{0.1550}\PYG{+w}{ }\PYG{l+m}{3.1536}
\PYG{k}{pair\PYGZus{}coeff}\PYG{+w}{ }\PYG{l+m}{2}\PYG{+w}{ }\PYG{l+m}{2}\PYG{+w}{ }\PYG{l+m}{0.0}\PYG{+w}{    }\PYG{l+m}{1.0}
\PYG{k}{pair\PYGZus{}coeff}\PYG{+w}{ }\PYG{l+m}{3}\PYG{+w}{ }\PYG{l+m}{3}\PYG{+w}{ }\PYG{l+m}{0.0}\PYG{+w}{    }\PYG{l+m}{1.0}

\PYG{k}{fix }\PYG{n+nv+nvIdentifier}{mol}\PYG{+w}{ }\PYG{n+nv+nvIdentifier}{all}\PYG{+w}{ }\PYG{n}{property}\PYG{o}{/}\PYG{n}{atom}\PYG{+w}{ }\PYG{n}{mol}\PYG{+w}{ }\PYG{n}{ghost}\PYG{+w}{ }\PYG{n}{yes}
\PYG{k}{molecule}\PYG{+w}{ }\PYG{n}{water}\PYG{+w}{ }\PYG{n}{tip4p.mol}
\PYG{k}{create\PYGZus{}atoms}\PYG{+w}{ }\PYG{l+m}{0}\PYG{+w}{ }\PYG{n}{random}\PYG{+w}{ }\PYG{l+m}{33}\PYG{+w}{ }\PYG{l+m}{34564}\PYG{+w}{ }\PYG{n}{NULL}\PYG{+w}{ }\PYG{n}{mol}\PYG{+w}{ }\PYG{n}{water}\PYG{+w}{ }\PYG{l+m}{25367}\PYG{+w}{ }\PYG{n}{overlap}\PYG{+w}{ }\PYG{l+m}{1.33}
\PYG{k}{neigh\PYGZus{}modify}\PYG{+w}{ }\PYG{n}{exclude}\PYG{+w}{ }\PYG{n}{molecule}\PYG{o}{/}\PYG{n}{intra}\PYG{+w}{ }\PYG{n}{all}

\PYG{k}{timestep}\PYG{+w}{ }\PYG{l+m}{0.5}
\PYG{k}{fix }\PYG{n+nv+nvIdentifier}{integrate}\PYG{+w}{ }\PYG{n+nv+nvIdentifier}{all}\PYG{+w}{ }\PYG{n}{rigid}\PYG{o}{/}\PYG{n}{small}\PYG{+w}{ }\PYG{n}{molecule}\PYG{+w}{ }\PYG{n}{langevin}\PYG{+w}{ }\PYG{l+m}{300.0}\PYG{+w}{ }\PYG{l+m}{300.0}\PYG{+w}{ }\PYG{l+m}{100.0}\PYG{+w}{ }\PYG{l+m}{2345634}

\PYG{k}{thermo\PYGZus{}style}\PYG{+w}{ }\PYG{n}{custom}\PYG{+w}{ }\PYG{n}{step}\PYG{+w}{ }\PYG{n}{temp}\PYG{+w}{ }\PYG{n}{press}\PYG{+w}{ }\PYG{n}{etotal}\PYG{+w}{ }\PYG{n}{density}\PYG{+w}{ }\PYG{n}{pe}\PYG{+w}{ }\PYG{n}{ke}
\PYG{k}{thermo}\PYG{+w}{ }\PYG{l+m}{2000}
\PYG{k}{run}\PYG{+w}{ }\PYG{l+m}{40000}
\PYG{k}{write\PYGZus{}data }\PYG{n+nv+nvIdentifier}{tip4p\PYGZhy{}explicit.data}\PYG{+w}{ }\PYG{n}{nocoeff}
\end{sphinxVerbatim}

\def\sphinxLiteralBlockLabel{\label{\detokenize{Howto_tip4p:tip4p-molecule}}}
\begin{sphinxVerbatim}[commandchars=\\\{\}]
\PYGZsh{} Water molecule. Explicit TIP4P geometry for use with fix rigid

4 atoms

Coords

1    0.00000  \PYGZhy{}0.06556   0.00000
2    0.75695   0.52032   0.00000
3   \PYGZhy{}0.75695   0.52032   0.00000
4    0.00000   0.08444   0.00000

Types

1        1   \PYGZsh{} O
2        2   \PYGZsh{} H
3        2   \PYGZsh{} H
4        3   \PYGZsh{} M

Charges

1        0.000
2        0.520
3        0.520
4       \PYGZhy{}1.040
\end{sphinxVerbatim}

\sphinxAtStartPar
Wikipedia also has a nice article on \sphinxhref{https://en.wikipedia.org/wiki/Water\_model}{water models}.


\bigskip\hrule\bigskip

\phantomsection\label{\detokenize{Howto_tip4p:jorgensen5}}
\sphinxAtStartPar
\sphinxstylestrong{(Jorgensen)} Jorgensen, Chandrasekhar, Madura, Impey, Klein, J Chem
Phys, 79, 926 (1983).

\phantomsection\label{\detokenize{Howto_tip4p:abascal1}}\begin{description}
\sphinxlineitem{\sphinxstylestrong{(Abascal1)} Abascal, Sanz, Fernandez, Vega, J Chem Phys, 122, 234511 (2005)}
\sphinxAtStartPar
\sphinxurl{https://doi.org/10.1063/1.1931662}

\end{description}
\phantomsection\label{\detokenize{Howto_tip4p:abascal2}}\begin{description}
\sphinxlineitem{\sphinxstylestrong{(Abascal2)} Abascal, J Chem Phys, 123, 234505 (2005)}
\sphinxAtStartPar
\sphinxurl{https://doi.org/10.1063/1.2121687}

\end{description}

\sphinxstepscope


\subsection{TIP5P water model}
\label{\detokenize{Howto_tip5p:tip5p-water-model}}\label{\detokenize{Howto_tip5p::doc}}
\sphinxAtStartPar
The five\sphinxhyphen{}point TIP5P rigid water model extends the {\hyperref[\detokenize{Howto_tip3p::doc}]{\sphinxcrossref{\DUrole{doc}{three\sphinxhyphen{}point
TIP3P model}}}} by adding two additional sites L, usually
massless, where the charge associated with the oxygen atom is placed.
These sites L are located at a fixed distance away from the oxygen atom,
forming a tetrahedral angle that is rotated by 90 degrees from the HOH
plane.  Those sites thus somewhat approximate lone pairs of the oxygen
and consequently improve the water structure to become even more
“tetrahedral” in comparison to the {\hyperref[\detokenize{Howto_tip4p::doc}]{\sphinxcrossref{\DUrole{doc}{four\sphinxhyphen{}point TIP4P model}}}}.

\sphinxAtStartPar
A suitable pair style with cutoff Coulomb would be:
\begin{itemize}
\item {} 
\sphinxAtStartPar
{\hyperref[\detokenize{pair_lj_cut_coul::doc}]{\sphinxcrossref{\DUrole{doc}{pair\_style lj/cut/coul/cut}}}}

\end{itemize}

\sphinxAtStartPar
or these commands for a long\sphinxhyphen{}range model:
\begin{itemize}
\item {} 
\sphinxAtStartPar
{\hyperref[\detokenize{pair_lj_cut_coul::doc}]{\sphinxcrossref{\DUrole{doc}{pair\_style lj/cut/coul/long}}}}

\item {} 
\sphinxAtStartPar
{\hyperref[\detokenize{pair_fep_soft::doc}]{\sphinxcrossref{\DUrole{doc}{pair\_style lj/cut/coul/long/soft}}}}

\item {} 
\sphinxAtStartPar
{\hyperref[\detokenize{kspace_style::doc}]{\sphinxcrossref{\DUrole{doc}{kspace\_style pppm}}}}

\item {} 
\sphinxAtStartPar
{\hyperref[\detokenize{kspace_style::doc}]{\sphinxcrossref{\DUrole{doc}{kspace\_style pppm/disp}}}}

\end{itemize}

\sphinxAtStartPar
A TIP5P model \sphinxstyleemphasis{must} be run using a {\hyperref[\detokenize{fix_rigid::doc}]{\sphinxcrossref{\DUrole{doc}{rigid fix}}}} since
there is no other option to keep this kind of structure rigid in LAMMPS.
In order to avoid that LAMMPS produces an error due to the massless L
sites, those need to be assigned a tiny non\sphinxhyphen{}zero mass.

\sphinxAtStartPar
The table below lists the force field parameters (in real {\hyperref[\detokenize{units::doc}]{\sphinxcrossref{\DUrole{doc}{units}}}}) to for a the TIP5P model with a cutoff {\hyperref[\detokenize{Howto_tip5p:mahoney}]{\sphinxcrossref{\DUrole{std,std-ref}{(Mahoney)}}}} and the TIP5P\sphinxhyphen{}E model {\hyperref[\detokenize{Howto_tip5p:rick}]{\sphinxcrossref{\DUrole{std,std-ref}{(Rick)}}}} for use with a
long\sphinxhyphen{}range Coulombic solver (e.g. Ewald or PPPM in LAMMPS).


\begin{savenotes}\sphinxattablestart
\sphinxthistablewithglobalstyle
\centering
\begin{tabular}[t]{\X{50}{100}\X{25}{100}\X{25}{100}}
\sphinxtoprule
\sphinxstyletheadfamily 
\sphinxAtStartPar
Parameter
&\sphinxstyletheadfamily 
\sphinxAtStartPar
TIP5P
&\sphinxstyletheadfamily 
\sphinxAtStartPar
TIP5P\sphinxhyphen{}E
\\
\sphinxmidrule
\sphinxtableatstartofbodyhook
\sphinxAtStartPar
O mass (amu)
&
\sphinxAtStartPar
15.9994
&
\sphinxAtStartPar
15.9994
\\
\sphinxhline
\sphinxAtStartPar
H mass (amu)
&
\sphinxAtStartPar
1.008
&
\sphinxAtStartPar
1.008
\\
\sphinxhline
\sphinxAtStartPar
O charge (\(e\))
&
\sphinxAtStartPar
0.0
&
\sphinxAtStartPar
0.0
\\
\sphinxhline
\sphinxAtStartPar
L charge (\(e\))
&
\sphinxAtStartPar
\sphinxhyphen{}0.241
&
\sphinxAtStartPar
\sphinxhyphen{}0.241
\\
\sphinxhline
\sphinxAtStartPar
H charge (\(e\))
&
\sphinxAtStartPar
0.241
&
\sphinxAtStartPar
0.241
\\
\sphinxhline
\sphinxAtStartPar
LJ \(\epsilon\) of OO (kcal/mole)
&
\sphinxAtStartPar
0.1600
&
\sphinxAtStartPar
0.1780
\\
\sphinxhline
\sphinxAtStartPar
LJ \(\sigma\) of OO (\(\AA\))
&
\sphinxAtStartPar
3.1200
&
\sphinxAtStartPar
3.0970
\\
\sphinxhline
\sphinxAtStartPar
LJ \(\epsilon\) of HH, LL, OH, OL, HL (kcal/mole)
&
\sphinxAtStartPar
0.0
&
\sphinxAtStartPar
0.0
\\
\sphinxhline
\sphinxAtStartPar
LJ \(\sigma\) of HH, LL, OH, OL, HL (\(\AA\))
&
\sphinxAtStartPar
1.0
&
\sphinxAtStartPar
1.0
\\
\sphinxhline
\sphinxAtStartPar
\(r_0\) of OH bond (\(\AA\))
&
\sphinxAtStartPar
0.9572
&
\sphinxAtStartPar
0.9572
\\
\sphinxhline
\sphinxAtStartPar
\(\theta_0\) of HOH angle
&
\sphinxAtStartPar
104.52\(^{\circ}\)
&
\sphinxAtStartPar
104.52\(^{\circ}\)
\\
\sphinxhline
\sphinxAtStartPar
OL distance (\(\AA\))
&
\sphinxAtStartPar
0.70
&
\sphinxAtStartPar
0.70
\\
\sphinxhline
\sphinxAtStartPar
\(\theta_0\) of LOL angle
&
\sphinxAtStartPar
109.47\(^{\circ}\)
&
\sphinxAtStartPar
109.47\(^{\circ}\)
\\
\sphinxbottomrule
\end{tabular}
\sphinxtableafterendhook\par
\sphinxattableend\end{savenotes}

\sphinxAtStartPar
Below is the code for a LAMMPS input file for setting up a simulation of
TIP5P water with a molecule file.  Because of using {\hyperref[\detokenize{fix_rigid::doc}]{\sphinxcrossref{\DUrole{doc}{fix
rigid/small}}}} no bonds need to be defined and thus no extra
storage needs to be reserved for them, but we need to either switch to
atom style full or use {\hyperref[\detokenize{fix_property_atom::doc}]{\sphinxcrossref{\DUrole{doc}{fix property/atom mol}}}}
so that fix rigid/small can identify rigid bodies by their molecule ID.
Also a {\hyperref[\detokenize{neigh_modify::doc}]{\sphinxcrossref{\DUrole{doc}{neigh\_modify exclude}}}} command is added to
exclude computing intramolecular non\sphinxhyphen{}bonded interactions, since those
are removed by the rigid fix anyway:

\begin{sphinxVerbatim}[commandchars=\\\{\}]
\PYG{k}{units}\PYG{+w}{ }\PYG{n}{real}
\PYG{k}{atom\PYGZus{}style}\PYG{+w}{ }\PYG{n}{charge}
\PYG{k}{atom\PYGZus{}modify}\PYG{+w}{ }\PYG{n}{map}\PYG{+w}{ }\PYG{n}{array}
\PYG{k}{region }\PYG{n+nv+nvIdentifier}{box}\PYG{+w}{ }\PYG{n}{block}\PYG{+w}{ }\PYG{o}{\PYGZhy{}}\PYG{l+m}{5}\PYG{+w}{ }\PYG{l+m}{5}\PYG{+w}{ }\PYG{o}{\PYGZhy{}}\PYG{l+m}{5}\PYG{+w}{ }\PYG{l+m}{5}\PYG{+w}{ }\PYG{o}{\PYGZhy{}}\PYG{l+m}{5}\PYG{+w}{ }\PYG{l+m}{5}
\PYG{k}{create\PYGZus{}box }\PYG{n+nv+nvIdentifier}{3}\PYG{+w}{ }\PYG{n+nv+nvIdentifier}{box}

\PYG{k}{mass}\PYG{+w}{ }\PYG{l+m}{1}\PYG{+w}{ }\PYG{l+m}{15.9994}
\PYG{k}{mass}\PYG{+w}{ }\PYG{l+m}{2}\PYG{+w}{ }\PYG{l+m}{1.008}
\PYG{k}{mass}\PYG{+w}{ }\PYG{l+m}{3}\PYG{+w}{ }\PYG{l+m}{1.0e\PYGZhy{}100}

\PYG{k}{pair\PYGZus{}style}\PYG{+w}{ }\PYG{n}{lj}\PYG{o}{/}\PYG{n}{cut}\PYG{o}{/}\PYG{n}{coul}\PYG{o}{/}\PYG{n}{cut}\PYG{+w}{ }\PYG{l+m}{8.0}
\PYG{k}{pair\PYGZus{}coeff}\PYG{+w}{ }\PYG{l+m}{1}\PYG{+w}{ }\PYG{l+m}{1}\PYG{+w}{ }\PYG{l+m}{0.160}\PYG{+w}{  }\PYG{l+m}{3.12}
\PYG{k}{pair\PYGZus{}coeff}\PYG{+w}{ }\PYG{l+m}{2}\PYG{+w}{ }\PYG{l+m}{2}\PYG{+w}{ }\PYG{l+m}{0.0}\PYG{+w}{    }\PYG{l+m}{1.0}
\PYG{k}{pair\PYGZus{}coeff}\PYG{+w}{ }\PYG{l+m}{3}\PYG{+w}{ }\PYG{l+m}{3}\PYG{+w}{ }\PYG{l+m}{0.0}\PYG{+w}{    }\PYG{l+m}{1.0}

\PYG{k}{fix }\PYG{n+nv+nvIdentifier}{mol}\PYG{+w}{ }\PYG{n+nv+nvIdentifier}{all}\PYG{+w}{ }\PYG{n}{property}\PYG{o}{/}\PYG{n}{atom}\PYG{+w}{ }\PYG{n}{mol}
\PYG{k}{molecule}\PYG{+w}{ }\PYG{n}{water}\PYG{+w}{ }\PYG{n}{tip5p.mol}
\PYG{k}{create\PYGZus{}atoms}\PYG{+w}{ }\PYG{l+m}{0}\PYG{+w}{ }\PYG{n}{random}\PYG{+w}{ }\PYG{l+m}{33}\PYG{+w}{ }\PYG{l+m}{34564}\PYG{+w}{ }\PYG{n}{NULL}\PYG{+w}{ }\PYG{n}{mol}\PYG{+w}{ }\PYG{n}{water}\PYG{+w}{ }\PYG{l+m}{25367}\PYG{+w}{ }\PYG{n}{overlap}\PYG{+w}{ }\PYG{l+m}{1.33}
\PYG{k}{neigh\PYGZus{}modify}\PYG{+w}{ }\PYG{n}{exclude}\PYG{+w}{ }\PYG{n}{molecule}\PYG{o}{/}\PYG{n}{intra}\PYG{+w}{ }\PYG{n}{all}

\PYG{k}{timestep}\PYG{+w}{ }\PYG{l+m}{0.5}
\PYG{k}{fix }\PYG{n+nv+nvIdentifier}{integrate}\PYG{+w}{ }\PYG{n+nv+nvIdentifier}{all}\PYG{+w}{ }\PYG{n}{rigid}\PYG{o}{/}\PYG{n}{small}\PYG{+w}{ }\PYG{n}{molecule}\PYG{+w}{ }\PYG{n}{langevin}\PYG{+w}{ }\PYG{l+m}{300.0}\PYG{+w}{ }\PYG{l+m}{300.0}\PYG{+w}{ }\PYG{l+m}{50.0}\PYG{+w}{ }\PYG{l+m}{235664}
\PYG{k}{reset\PYGZus{}timestep}\PYG{+w}{ }\PYG{l+m}{0}

\PYG{k}{thermo\PYGZus{}style}\PYG{+w}{ }\PYG{n}{custom}\PYG{+w}{ }\PYG{n}{step}\PYG{+w}{ }\PYG{n}{temp}\PYG{+w}{ }\PYG{n}{press}\PYG{+w}{ }\PYG{n}{etotal}\PYG{+w}{ }\PYG{n}{density}\PYG{+w}{ }\PYG{n}{pe}\PYG{+w}{ }\PYG{n}{ke}
\PYG{k}{thermo}\PYG{+w}{ }\PYG{l+m}{1000}
\PYG{k}{run}\PYG{+w}{ }\PYG{l+m}{20000}
\PYG{k}{write\PYGZus{}data }\PYG{n+nv+nvIdentifier}{tip5p.data}\PYG{+w}{ }\PYG{n}{nocoeff}
\end{sphinxVerbatim}

\def\sphinxLiteralBlockLabel{\label{\detokenize{Howto_tip5p:tip5p-molecule}}}
\begin{sphinxVerbatim}[commandchars=\\\{\}]
\PYGZsh{} Water molecule. Explicit TIP5P geometry for use with fix rigid

5 atoms

Coords

1    0.00000  \PYGZhy{}0.06556   0.00000
2    0.75695   0.52032   0.00000
3   \PYGZhy{}0.75695   0.52032   0.00000
4    0.00000  \PYGZhy{}0.46971   0.57154
5    0.00000  \PYGZhy{}0.46971  \PYGZhy{}0.57154

Types

1        1   \PYGZsh{} O
2        2   \PYGZsh{} H
3        2   \PYGZsh{} H
4        3   \PYGZsh{} L
5        3   \PYGZsh{} L

Charges

1        0.000
2        0.241
3        0.241
4       \PYGZhy{}0.241
5       \PYGZhy{}0.241
\end{sphinxVerbatim}

\sphinxAtStartPar
Wikipedia also has a nice article on \sphinxhref{https://en.wikipedia.org/wiki/Water\_model}{water models}.


\bigskip\hrule\bigskip

\phantomsection\label{\detokenize{Howto_tip5p:mahoney}}
\sphinxAtStartPar
\sphinxstylestrong{(Mahoney)} Mahoney, Jorgensen, J Chem Phys 112, 8910 (2000)

\phantomsection\label{\detokenize{Howto_tip5p:rick}}
\sphinxAtStartPar
\sphinxstylestrong{(Rick)} Rick, J Chem Phys 120, 6085 (2004)

\sphinxstepscope


\subsection{SPC water model}
\label{\detokenize{Howto_spc:spc-water-model}}\label{\detokenize{Howto_spc::doc}}
\sphinxAtStartPar
The SPC water model specifies a 3\sphinxhyphen{}site rigid water molecule with
charges and Lennard\sphinxhyphen{}Jones parameters assigned to each of the three atoms.
In LAMMPS the {\hyperref[\detokenize{fix_shake::doc}]{\sphinxcrossref{\DUrole{doc}{fix shake}}}} command can be used to hold
the two O\sphinxhyphen{}H bonds and the H\sphinxhyphen{}O\sphinxhyphen{}H angle rigid.  A bond style of
\sphinxstyleemphasis{harmonic} and an angle style of \sphinxstyleemphasis{harmonic} or \sphinxstyleemphasis{charmm} should also be
used.

\sphinxAtStartPar
These are the additional parameters (in real units) to set for O and H
atoms and the water molecule to run a rigid SPC model.

\begin{DUlineblock}{0em}
\item[] O mass = 15.9994
\item[] H mass = 1.008
\item[] O charge = \sphinxhyphen{}0.820
\item[] H charge = 0.410
\item[] LJ \(\epsilon\) of OO = 0.1553
\item[] LJ \(\sigma\) of OO = 3.166
\item[] LJ \(\epsilon\), \(\sigma\) of OH, HH = 0.0
\item[] \(r_0\) of OH bond = 1.0
\item[] \(\theta_0\) of HOH angle = 109.47\(^{\circ}\)
\end{DUlineblock}

\sphinxAtStartPar
Note that as originally proposed, the SPC model was run with a 9
Angstrom cutoff for both LJ and Coulomb terms.  It can also be used
with long\sphinxhyphen{}range electrostatic solvers (e.g. Ewald or PPPM in LAMMPS)
without changing any of the parameters above, although it becomes
a different model in that mode of usage.

\sphinxAtStartPar
The SPC/E (extended) water model is the same, except
the partial charge assignments change:

\begin{DUlineblock}{0em}
\item[] O charge = \sphinxhyphen{}0.8476
\item[] H charge = 0.4238
\end{DUlineblock}

\sphinxAtStartPar
See the {\hyperref[\detokenize{Howto_spc:howto-berendsen}]{\sphinxcrossref{\DUrole{std,std-ref}{(Berendsen2)}}}} reference for more details on both
the SPC and SPC/E models.

\sphinxAtStartPar
Below is the code for a LAMMPS input file and a molecule file
(\sphinxcode{\sphinxupquote{spce.mol}}) of SPC/E water for use with the {\hyperref[\detokenize{molecule::doc}]{\sphinxcrossref{\DUrole{doc}{molecule command}}}} demonstrating how to set up a small bulk water system for
SPC/E with rigid bonds.

\begin{sphinxVerbatim}[commandchars=\\\{\}]
\PYG{k}{units}\PYG{+w}{ }\PYG{n}{real}
\PYG{k}{atom\PYGZus{}style}\PYG{+w}{ }\PYG{n}{full}
\PYG{k}{region }\PYG{n+nv+nvIdentifier}{box}\PYG{+w}{ }\PYG{n}{block}\PYG{+w}{ }\PYG{o}{\PYGZhy{}}\PYG{l+m}{5}\PYG{+w}{ }\PYG{l+m}{5}\PYG{+w}{ }\PYG{o}{\PYGZhy{}}\PYG{l+m}{5}\PYG{+w}{ }\PYG{l+m}{5}\PYG{+w}{ }\PYG{o}{\PYGZhy{}}\PYG{l+m}{5}\PYG{+w}{ }\PYG{l+m}{5}
\PYG{k}{create\PYGZus{}box }\PYG{n+nv+nvIdentifier}{2}\PYG{+w}{ }\PYG{n+nv+nvIdentifier}{box}\PYG{+w}{  }\PYG{n}{bond}\PYG{o}{/}\PYG{n}{types}\PYG{+w}{ }\PYG{l+m}{1}\PYG{+w}{ }\PYG{n}{angle}\PYG{o}{/}\PYG{n}{types}\PYG{+w}{ }\PYG{l+m}{1}\PYG{+w}{ }\PYG{o}{\PYGZam{}}
\PYG{+w}{                }\PYG{n}{extra}\PYG{o}{/}\PYG{n}{bond}\PYG{o}{/}\PYG{n}{per}\PYG{o}{/}\PYG{n}{atom}\PYG{+w}{ }\PYG{l+m}{2}\PYG{+w}{ }\PYG{n}{extra}\PYG{o}{/}\PYG{n}{angle}\PYG{o}{/}\PYG{n}{per}\PYG{o}{/}\PYG{n}{atom}\PYG{+w}{ }\PYG{l+m}{1}\PYG{+w}{ }\PYG{n}{extra}\PYG{o}{/}\PYG{n}{special}\PYG{o}{/}\PYG{n}{per}\PYG{o}{/}\PYG{n}{atom}\PYG{+w}{ }\PYG{l+m}{2}

\PYG{k}{mass}\PYG{+w}{ }\PYG{l+m}{1}\PYG{+w}{ }\PYG{l+m}{15.9994}
\PYG{k}{mass}\PYG{+w}{ }\PYG{l+m}{2}\PYG{+w}{ }\PYG{l+m}{1.008}

\PYG{k}{pair\PYGZus{}style}\PYG{+w}{ }\PYG{n}{lj}\PYG{o}{/}\PYG{n}{cut}\PYG{o}{/}\PYG{n}{coul}\PYG{o}{/}\PYG{n}{cut}\PYG{+w}{ }\PYG{l+m}{10.0}
\PYG{k}{pair\PYGZus{}coeff}\PYG{+w}{ }\PYG{l+m}{1}\PYG{+w}{ }\PYG{l+m}{1}\PYG{+w}{ }\PYG{l+m}{0.1553}\PYG{+w}{ }\PYG{l+m}{3.166}
\PYG{k}{pair\PYGZus{}coeff}\PYG{+w}{ }\PYG{l+m}{1}\PYG{+w}{ }\PYG{l+m}{2}\PYG{+w}{ }\PYG{l+m}{0.0}\PYG{+w}{    }\PYG{l+m}{1.0}
\PYG{k}{pair\PYGZus{}coeff}\PYG{+w}{ }\PYG{l+m}{2}\PYG{+w}{ }\PYG{l+m}{2}\PYG{+w}{ }\PYG{l+m}{0.0}\PYG{+w}{    }\PYG{l+m}{1.0}

\PYG{k}{bond\PYGZus{}style}\PYG{+w}{ }\PYG{n}{zero}
\PYG{k}{bond\PYGZus{}coeff}\PYG{+w}{ }\PYG{l+m}{1}\PYG{+w}{ }\PYG{l+m}{1.0}

\PYG{k}{angle\PYGZus{}style}\PYG{+w}{ }\PYG{n}{zero}
\PYG{k}{angle\PYGZus{}coeff}\PYG{+w}{ }\PYG{l+m}{1}\PYG{+w}{ }\PYG{l+m}{109.47}

\PYG{k}{molecule}\PYG{+w}{ }\PYG{n}{water}\PYG{+w}{ }\PYG{n}{spce.mol}
\PYG{k}{create\PYGZus{}atoms}\PYG{+w}{ }\PYG{l+m}{0}\PYG{+w}{ }\PYG{n}{random}\PYG{+w}{ }\PYG{l+m}{33}\PYG{+w}{ }\PYG{l+m}{34564}\PYG{+w}{ }\PYG{n}{NULL}\PYG{+w}{ }\PYG{n}{mol}\PYG{+w}{ }\PYG{n}{water}\PYG{+w}{ }\PYG{l+m}{25367}\PYG{+w}{ }\PYG{n}{overlap}\PYG{+w}{ }\PYG{l+m}{1.33}

\PYG{k}{timestep}\PYG{+w}{ }\PYG{l+m}{1.0}
\PYG{k}{fix }\PYG{n+nv+nvIdentifier}{rigid}\PYG{+w}{     }\PYG{n+nv+nvIdentifier}{all}\PYG{+w}{ }\PYG{n}{shake}\PYG{+w}{ }\PYG{l+m}{0.0001}\PYG{+w}{ }\PYG{l+m}{10}\PYG{+w}{ }\PYG{l+m}{10000}\PYG{+w}{ }\PYG{n}{b}\PYG{+w}{ }\PYG{l+m}{1}\PYG{+w}{ }\PYG{n}{a}\PYG{+w}{ }\PYG{l+m}{1}
\PYG{k}{minimize}\PYG{+w}{ }\PYG{l+m}{0.0}\PYG{+w}{ }\PYG{l+m}{0.0}\PYG{+w}{ }\PYG{l+m}{1000}\PYG{+w}{ }\PYG{l+m}{10000}
\PYG{k}{velocity }\PYG{n+nv+nvIdentifier}{all}\PYG{+w}{ }\PYG{n}{create}\PYG{+w}{ }\PYG{l+m}{300.0}\PYG{+w}{ }\PYG{l+m}{5463576}
\PYG{k}{fix }\PYG{n+nv+nvIdentifier}{integrate}\PYG{+w}{ }\PYG{n+nv+nvIdentifier}{all}\PYG{+w}{ }\PYG{n}{nvt}\PYG{+w}{ }\PYG{n}{temp}\PYG{+w}{ }\PYG{l+m}{300.0}\PYG{+w}{ }\PYG{l+m}{300.0}\PYG{+w}{ }\PYG{l+m}{100.0}

\PYG{k}{thermo\PYGZus{}style}\PYG{+w}{ }\PYG{n}{custom}\PYG{+w}{ }\PYG{n}{step}\PYG{+w}{ }\PYG{n}{temp}\PYG{+w}{ }\PYG{n}{press}\PYG{+w}{ }\PYG{n}{etotal}\PYG{+w}{ }\PYG{n}{density}\PYG{+w}{ }\PYG{n}{pe}\PYG{+w}{ }\PYG{n}{ke}
\PYG{k}{thermo}\PYG{+w}{ }\PYG{l+m}{1000}
\PYG{k}{run}\PYG{+w}{ }\PYG{l+m}{20000}\PYG{+w}{ }\PYG{n}{upto}
\PYG{k}{write\PYGZus{}data }\PYG{n+nv+nvIdentifier}{spce.data}\PYG{+w}{ }\PYG{n}{nocoeff}
\end{sphinxVerbatim}

\def\sphinxLiteralBlockLabel{\label{\detokenize{Howto_spc:spce-molecule}}}
\begin{sphinxVerbatim}[commandchars=\\\{\}]
\PYGZsh{} Water molecule. SPC/E geometry

3 atoms
2 bonds
1 angles

Coords

1    0.00000  \PYGZhy{}0.06461   0.00000
2    0.81649   0.51275   0.00000
3   \PYGZhy{}0.81649   0.51275   0.00000

Types

1        1   \PYGZsh{} O
2        2   \PYGZsh{} H
3        2   \PYGZsh{} H

Charges

1       \PYGZhy{}0.8476
2        0.4238
3        0.4238

Bonds

1   1      1      2
2   1      1      3

Angles

1   1      2      1      3

Shake Flags

1 1
2 1
3 1

Shake Atoms

1 1 2 3
2 1 2 3
3 1 2 3

Shake Bond Types

1 1 1 1
2 1 1 1
3 1 1 1

Special Bond Counts

1 2 0 0
2 1 1 0
3 1 1 0

Special Bonds

1 2 3
2 1 3
3 1 2
\end{sphinxVerbatim}

\sphinxAtStartPar
Wikipedia also has a nice article on \sphinxhref{https://en.wikipedia.org/wiki/Water\_model}{water models}.


\bigskip\hrule\bigskip

\phantomsection\label{\detokenize{Howto_spc:howto-berendsen}}
\sphinxAtStartPar
\sphinxstylestrong{(Berendsen2)} Berendsen, Grigera, Straatsma, J Phys Chem, 91, 6269\sphinxhyphen{}6271 (1987).


\section{Packages howto}
\label{\detokenize{Howto:packages-howto}}
\sphinxstepscope


\subsection{Finite\sphinxhyphen{}size spherical and aspherical particles}
\label{\detokenize{Howto_spherical:finite-size-spherical-and-aspherical-particles}}\label{\detokenize{Howto_spherical::doc}}
\sphinxAtStartPar
Typical MD models treat atoms or particles as point masses.  Sometimes
it is desirable to have a model with finite\sphinxhyphen{}size particles such as
spheroids or ellipsoids or generalized aspherical bodies.  The
difference is that such particles have a moment of inertia, rotational
energy, and angular momentum.  Rotation is induced by torque coming
from interactions with other particles.

\sphinxAtStartPar
LAMMPS has several options for running simulations with these kinds of
particles.  The following aspects are discussed in turn:
\begin{itemize}
\item {} 
\sphinxAtStartPar
atom styles

\item {} 
\sphinxAtStartPar
pair potentials

\item {} 
\sphinxAtStartPar
time integration

\item {} 
\sphinxAtStartPar
computes, thermodynamics, and dump output

\item {} 
\sphinxAtStartPar
rigid bodies composed of finite\sphinxhyphen{}size particles

\end{itemize}

\sphinxAtStartPar
Example input scripts for these kinds of models are in the body,
colloid, dipole, ellipse, line, peri, pour, and tri directories of the
{\hyperref[\detokenize{Examples::doc}]{\sphinxcrossref{\DUrole{doc}{examples directory}}}} in the LAMMPS distribution.


\subsubsection{Atom styles}
\label{\detokenize{Howto_spherical:atom-styles}}
\sphinxAtStartPar
There are several {\hyperref[\detokenize{atom_style::doc}]{\sphinxcrossref{\DUrole{doc}{atom styles}}}} that allow for
definition of finite\sphinxhyphen{}size particles: sphere, dipole, ellipsoid, line,
tri, peri, and body.

\sphinxAtStartPar
The sphere style defines particles that are spheroids and each
particle can have a unique diameter and mass (or density).  These
particles store an angular velocity (omega) and can be acted upon by
torque.  The “set” command can be used to modify the diameter and mass
of individual particles, after then are created.

\sphinxAtStartPar
The dipole style does not actually define finite\sphinxhyphen{}size particles, but
is often used in conjunction with spherical particles, via a command
like

\begin{sphinxVerbatim}[commandchars=\\\{\}]
\PYG{k}{atom\PYGZus{}style}\PYG{+w}{ }\PYG{n}{hybrid}\PYG{+w}{ }\PYG{n}{sphere}\PYG{+w}{ }\PYG{n}{dipole}
\end{sphinxVerbatim}

\sphinxAtStartPar
This is because when dipoles interact with each other, they induce
torques, and a particle must be finite\sphinxhyphen{}size (i.e. have a moment of
inertia) in order to respond and rotate.  See the {\hyperref[\detokenize{atom_style::doc}]{\sphinxcrossref{\DUrole{doc}{atom\_style dipole}}}} command for details.  The “set” command can be
used to modify the orientation and length of the dipole moment of
individual particles, after then are created.

\sphinxAtStartPar
The ellipsoid style defines particles that are ellipsoids and thus can
be aspherical.  Each particle has a shape, specified by 3 diameters,
and mass (or density).  These particles store an angular momentum and
their orientation (quaternion), and can be acted upon by torque.  They
do not store an angular velocity (omega), which can be in a different
direction than angular momentum, rather they compute it as needed.
The “set” command can be used to modify the diameter, orientation, and
mass of individual particles, after then are created.  It also has a
brief explanation of what quaternions are.

\sphinxAtStartPar
The line style defines line segment particles with two end points and
a mass (or density).  They can be used in 2d simulations, and they can
be joined together to form rigid bodies which represent arbitrary
polygons.

\sphinxAtStartPar
The tri style defines triangular particles with three corner points
and a mass (or density).  They can be used in 3d simulations, and they
can be joined together to form rigid bodies which represent arbitrary
particles with a triangulated surface.

\sphinxAtStartPar
The peri style is used with {\hyperref[\detokenize{pair_peri::doc}]{\sphinxcrossref{\DUrole{doc}{Peridynamic models}}}} and
defines particles as having a volume, that is used internally in the
{\hyperref[\detokenize{pair_peri::doc}]{\sphinxcrossref{\DUrole{doc}{pair\_style peri}}}} potentials.

\sphinxAtStartPar
The body style allows for definition of particles which can represent
complex entities, such as surface meshes of discrete points,
collections of sub\sphinxhyphen{}particles, deformable objects, etc.  The body style
is discussed in more detail on the {\hyperref[\detokenize{Howto_body::doc}]{\sphinxcrossref{\DUrole{doc}{Howto body}}}} doc
page.

\sphinxAtStartPar
Note that if one of these atom styles is used (or multiple styles via
the {\hyperref[\detokenize{atom_style::doc}]{\sphinxcrossref{\DUrole{doc}{atom\_style hybrid}}}} command), not all particles in
the system are required to be finite\sphinxhyphen{}size or aspherical.

\sphinxAtStartPar
For example, in the ellipsoid style, if the 3 shape parameters are set
to the same value, the particle will be a sphere rather than an
ellipsoid.  If the 3 shape parameters are all set to 0.0 or if the
diameter is set to 0.0, it will be a point particle.  In the line or
tri style, if the lineflag or triflag is specified as 0, then it
will be a point particle.

\sphinxAtStartPar
Some of the pair styles used to compute pairwise interactions between
finite\sphinxhyphen{}size particles also compute the correct interaction with point
particles as well, e.g. the interaction between a point particle and a
finite\sphinxhyphen{}size particle or between two point particles.  If necessary,
{\hyperref[\detokenize{pair_hybrid::doc}]{\sphinxcrossref{\DUrole{doc}{pair\_style hybrid}}}} can be used to ensure the correct
interactions are computed for the appropriate style of interactions.
Likewise, using groups to partition particles (ellipsoids versus
spheres versus point particles) will allow you to use the appropriate
time integrators and temperature computations for each class of
particles.  See the doc pages for various commands for details.

\sphinxAtStartPar
Also note that for {\hyperref[\detokenize{dimension::doc}]{\sphinxcrossref{\DUrole{doc}{2d simulations}}}}, atom styles sphere
and ellipsoid still use 3d particles, rather than as circular disks or
ellipses.  This means they have the same moment of inertia as the 3d
object.  When temperature is computed, the correct degrees of freedom
are used for rotation in a 2d versus 3d system.


\subsubsection{Pair potentials}
\label{\detokenize{Howto_spherical:pair-potentials}}
\sphinxAtStartPar
When a system with finite\sphinxhyphen{}size particles is defined, the particles
will only rotate and experience torque if the force field computes
such interactions.  These are the various {\hyperref[\detokenize{pair_style::doc}]{\sphinxcrossref{\DUrole{doc}{pair styles}}}} that generate torque:
\begin{itemize}
\item {} 
\sphinxAtStartPar
{\hyperref[\detokenize{pair_gran::doc}]{\sphinxcrossref{\DUrole{doc}{pair\_style gran/history}}}}

\item {} 
\sphinxAtStartPar
{\hyperref[\detokenize{pair_gran::doc}]{\sphinxcrossref{\DUrole{doc}{pair\_style gran/hertz}}}}

\item {} 
\sphinxAtStartPar
{\hyperref[\detokenize{pair_gran::doc}]{\sphinxcrossref{\DUrole{doc}{pair\_style gran/no\_history}}}}

\item {} 
\sphinxAtStartPar
{\hyperref[\detokenize{pair_dipole::doc}]{\sphinxcrossref{\DUrole{doc}{pair\_style dipole/cut}}}}

\item {} 
\sphinxAtStartPar
{\hyperref[\detokenize{pair_gayberne::doc}]{\sphinxcrossref{\DUrole{doc}{pair\_style gayberne}}}}

\item {} 
\sphinxAtStartPar
{\hyperref[\detokenize{pair_resquared::doc}]{\sphinxcrossref{\DUrole{doc}{pair\_style resquared}}}}

\item {} 
\sphinxAtStartPar
{\hyperref[\detokenize{pair_brownian::doc}]{\sphinxcrossref{\DUrole{doc}{pair\_style brownian}}}}

\item {} 
\sphinxAtStartPar
{\hyperref[\detokenize{pair_lubricate::doc}]{\sphinxcrossref{\DUrole{doc}{pair\_style lubricate}}}}

\item {} 
\sphinxAtStartPar
{\hyperref[\detokenize{pair_line_lj::doc}]{\sphinxcrossref{\DUrole{doc}{pair\_style line/lj}}}}

\item {} 
\sphinxAtStartPar
{\hyperref[\detokenize{pair_tri_lj::doc}]{\sphinxcrossref{\DUrole{doc}{pair\_style tri/lj}}}}

\item {} 
\sphinxAtStartPar
{\hyperref[\detokenize{pair_body_nparticle::doc}]{\sphinxcrossref{\DUrole{doc}{pair\_style body/nparticle}}}}

\end{itemize}

\sphinxAtStartPar
The granular pair styles are used with spherical particles.  The
dipole pair style is used with the dipole atom style, which could be
applied to spherical or ellipsoidal particles.  The GayBerne and
REsquared potentials require ellipsoidal particles, though they will
also work if the 3 shape parameters are the same (a sphere).  The
Brownian and lubrication potentials are used with spherical particles.
The line, tri, and body potentials are used with line segment,
triangular, and body particles respectively.


\subsubsection{Time integration}
\label{\detokenize{Howto_spherical:time-integration}}
\sphinxAtStartPar
There are several fixes that perform time integration on finite\sphinxhyphen{}size
spherical particles, meaning the integrators update the rotational
orientation and angular velocity or angular momentum of the particles:
\begin{itemize}
\item {} 
\sphinxAtStartPar
{\hyperref[\detokenize{fix_nve_sphere::doc}]{\sphinxcrossref{\DUrole{doc}{fix nve/sphere}}}}

\item {} 
\sphinxAtStartPar
{\hyperref[\detokenize{fix_nvt_sphere::doc}]{\sphinxcrossref{\DUrole{doc}{fix nvt/sphere}}}}

\item {} 
\sphinxAtStartPar
{\hyperref[\detokenize{fix_npt_sphere::doc}]{\sphinxcrossref{\DUrole{doc}{fix npt/sphere}}}}

\end{itemize}

\sphinxAtStartPar
Likewise, there are 3 fixes that perform time integration on
ellipsoidal particles:
\begin{itemize}
\item {} 
\sphinxAtStartPar
{\hyperref[\detokenize{fix_nve_asphere::doc}]{\sphinxcrossref{\DUrole{doc}{fix nve/asphere}}}}

\item {} 
\sphinxAtStartPar
{\hyperref[\detokenize{fix_nvt_asphere::doc}]{\sphinxcrossref{\DUrole{doc}{fix nvt/asphere}}}}

\item {} 
\sphinxAtStartPar
{\hyperref[\detokenize{fix_npt_asphere::doc}]{\sphinxcrossref{\DUrole{doc}{fix npt/asphere}}}}

\end{itemize}

\sphinxAtStartPar
The advantage of these fixes is that those which thermostat the
particles include the rotational degrees of freedom in the temperature
calculation and thermostatting.  The {\hyperref[\detokenize{fix_langevin::doc}]{\sphinxcrossref{\DUrole{doc}{fix langevin}}}}
command can also be used with its \sphinxstyleemphasis{omgea} or \sphinxstyleemphasis{angmom} options to
thermostat the rotational degrees of freedom for spherical or
ellipsoidal particles.  Other thermostatting fixes only operate on the
translational kinetic energy of finite\sphinxhyphen{}size particles.

\sphinxAtStartPar
These fixes perform constant NVE time integration on line segment,
triangular, and body particles:
\begin{itemize}
\item {} 
\sphinxAtStartPar
{\hyperref[\detokenize{fix_nve_line::doc}]{\sphinxcrossref{\DUrole{doc}{fix nve/line}}}}

\item {} 
\sphinxAtStartPar
{\hyperref[\detokenize{fix_nve_tri::doc}]{\sphinxcrossref{\DUrole{doc}{fix nve/tri}}}}

\item {} 
\sphinxAtStartPar
{\hyperref[\detokenize{fix_nve_body::doc}]{\sphinxcrossref{\DUrole{doc}{fix nve/body}}}}

\end{itemize}

\sphinxAtStartPar
Note that for mixtures of point and finite\sphinxhyphen{}size particles, these
integration fixes can only be used with {\hyperref[\detokenize{group::doc}]{\sphinxcrossref{\DUrole{doc}{groups}}}} which
contain finite\sphinxhyphen{}size particles.


\subsubsection{Computes, thermodynamics, and dump output}
\label{\detokenize{Howto_spherical:computes-thermodynamics-and-dump-output}}
\sphinxAtStartPar
There are several computes that calculate the temperature or
rotational energy of spherical or ellipsoidal particles:
\begin{itemize}
\item {} 
\sphinxAtStartPar
{\hyperref[\detokenize{compute_temp_sphere::doc}]{\sphinxcrossref{\DUrole{doc}{compute temp/sphere}}}}

\item {} 
\sphinxAtStartPar
{\hyperref[\detokenize{compute_temp_asphere::doc}]{\sphinxcrossref{\DUrole{doc}{compute temp/asphere}}}}

\item {} 
\sphinxAtStartPar
{\hyperref[\detokenize{compute_erotate_sphere::doc}]{\sphinxcrossref{\DUrole{doc}{compute erotate/sphere}}}}

\item {} 
\sphinxAtStartPar
{\hyperref[\detokenize{compute_erotate_asphere::doc}]{\sphinxcrossref{\DUrole{doc}{compute erotate/asphere}}}}

\end{itemize}

\sphinxAtStartPar
These include rotational degrees of freedom in their computation.  If
you wish the thermodynamic output of temperature or pressure to use
one of these computes (e.g. for a system entirely composed of
finite\sphinxhyphen{}size particles), then the compute can be defined and the
{\hyperref[\detokenize{thermo_modify::doc}]{\sphinxcrossref{\DUrole{doc}{thermo\_modify}}}} command used.  Note that by default
thermodynamic quantities will be calculated with a temperature that
only includes translational degrees of freedom.  See the
{\hyperref[\detokenize{thermo_style::doc}]{\sphinxcrossref{\DUrole{doc}{thermo\_style}}}} command for details.

\sphinxAtStartPar
These commands can be used to output various attributes of finite\sphinxhyphen{}size
particles:
\begin{itemize}
\item {} 
\sphinxAtStartPar
{\hyperref[\detokenize{dump::doc}]{\sphinxcrossref{\DUrole{doc}{dump custom}}}}

\item {} 
\sphinxAtStartPar
{\hyperref[\detokenize{compute_property_atom::doc}]{\sphinxcrossref{\DUrole{doc}{compute property/atom}}}}

\item {} 
\sphinxAtStartPar
{\hyperref[\detokenize{dump::doc}]{\sphinxcrossref{\DUrole{doc}{dump local}}}}

\item {} 
\sphinxAtStartPar
{\hyperref[\detokenize{compute_body_local::doc}]{\sphinxcrossref{\DUrole{doc}{compute body/local}}}}

\end{itemize}

\sphinxAtStartPar
Attributes include the dipole moment, the angular velocity, the
angular momentum, the quaternion, the torque, the end\sphinxhyphen{}point and
corner\sphinxhyphen{}point coordinates (for line and tri particles), and
sub\sphinxhyphen{}particle attributes of body particles.


\subsubsection{Rigid bodies composed of finite\sphinxhyphen{}size particles}
\label{\detokenize{Howto_spherical:rigid-bodies-composed-of-finite-size-particles}}
\sphinxAtStartPar
The {\hyperref[\detokenize{fix_rigid::doc}]{\sphinxcrossref{\DUrole{doc}{fix rigid}}}} command treats a collection of
particles as a rigid body, computes its inertia tensor, sums the total
force and torque on the rigid body each timestep due to forces on its
constituent particles, and integrates the motion of the rigid body.

\sphinxAtStartPar
If any of the constituent particles of a rigid body are finite\sphinxhyphen{}size
particles (spheres or ellipsoids or line segments or triangles), then
their contribution to the inertia tensor of the body is different than
if they were point particles.  This means the rotational dynamics of
the rigid body will be different.  Thus a model of a dimer is
different if the dimer consists of two point masses versus two
spheroids, even if the two particles have the same mass.  Finite\sphinxhyphen{}size
particles that experience torque due to their interaction with other
particles will also impart that torque to a rigid body they are part
of.

\sphinxAtStartPar
See the “fix rigid” command for example of complex rigid\sphinxhyphen{}body models
it is possible to define in LAMMPS.

\sphinxAtStartPar
Note that the {\hyperref[\detokenize{fix_shake::doc}]{\sphinxcrossref{\DUrole{doc}{fix shake}}}} command can also be used to
treat 2, 3, or 4 particles as a rigid body, but it always assumes the
particles are point masses.

\sphinxAtStartPar
Also note that body particles cannot be modeled with the {\hyperref[\detokenize{fix_rigid::doc}]{\sphinxcrossref{\DUrole{doc}{fix rigid}}}} command.  Body particles are treated by LAMMPS
as single particles, though they can store internal state, such as a
list of sub\sphinxhyphen{}particles.  Individual body particles are typically treated
as rigid bodies, and their motion integrated with a command like {\hyperref[\detokenize{fix_nve_body::doc}]{\sphinxcrossref{\DUrole{doc}{fix nve/body}}}}.  Interactions between pairs of body
particles are computed via a command like {\hyperref[\detokenize{pair_body_nparticle::doc}]{\sphinxcrossref{\DUrole{doc}{pair\_style body/nparticle}}}}.

\sphinxstepscope


\subsection{Granular models}
\label{\detokenize{Howto_granular:granular-models}}\label{\detokenize{Howto_granular::doc}}
\sphinxAtStartPar
Granular system are composed of spherical particles with a diameter,
as opposed to point particles.  This means they have an angular
velocity and torque can be imparted to them to cause them to rotate.

\sphinxAtStartPar
To run a simulation of a granular model, you will want to use
the following commands:
\begin{itemize}
\item {} 
\sphinxAtStartPar
{\hyperref[\detokenize{atom_style::doc}]{\sphinxcrossref{\DUrole{doc}{atom\_style sphere}}}}

\item {} 
\sphinxAtStartPar
{\hyperref[\detokenize{fix_nve_sphere::doc}]{\sphinxcrossref{\DUrole{doc}{fix nve/sphere}}}}

\item {} 
\sphinxAtStartPar
{\hyperref[\detokenize{fix_gravity::doc}]{\sphinxcrossref{\DUrole{doc}{fix gravity}}}}

\end{itemize}

\sphinxAtStartPar
This compute
\begin{itemize}
\item {} 
\sphinxAtStartPar
{\hyperref[\detokenize{compute_erotate_sphere::doc}]{\sphinxcrossref{\DUrole{doc}{compute erotate/sphere}}}}

\end{itemize}

\sphinxAtStartPar
calculates rotational kinetic energy which can be {\hyperref[\detokenize{Howto_output::doc}]{\sphinxcrossref{\DUrole{doc}{output with thermodynamic info}}}}.
The compute
\begin{itemize}
\item {} 
\sphinxAtStartPar
{\hyperref[\detokenize{compute_fabric::doc}]{\sphinxcrossref{\DUrole{doc}{compute fabric}}}}

\end{itemize}

\sphinxAtStartPar
calculates various versions of the fabric tensor for granular and non\sphinxhyphen{}granular pair styles.

\sphinxAtStartPar
Use one of these 4 pair potentials, which compute forces and torques
between interacting pairs of particles:
\begin{itemize}
\item {} 
\sphinxAtStartPar
{\hyperref[\detokenize{pair_gran::doc}]{\sphinxcrossref{\DUrole{doc}{pair\_style gran/history}}}}

\item {} 
\sphinxAtStartPar
{\hyperref[\detokenize{pair_gran::doc}]{\sphinxcrossref{\DUrole{doc}{pair\_style gran/no\_history}}}}

\item {} 
\sphinxAtStartPar
{\hyperref[\detokenize{pair_gran::doc}]{\sphinxcrossref{\DUrole{doc}{pair\_style gran/hertzian}}}}

\item {} 
\sphinxAtStartPar
{\hyperref[\detokenize{pair_granular::doc}]{\sphinxcrossref{\DUrole{doc}{pair\_style granular}}}}

\end{itemize}

\sphinxAtStartPar
These commands implement fix options specific to granular systems:
\begin{itemize}
\item {} 
\sphinxAtStartPar
{\hyperref[\detokenize{fix_freeze::doc}]{\sphinxcrossref{\DUrole{doc}{fix freeze}}}}

\item {} 
\sphinxAtStartPar
{\hyperref[\detokenize{fix_pour::doc}]{\sphinxcrossref{\DUrole{doc}{fix pour}}}}

\item {} 
\sphinxAtStartPar
{\hyperref[\detokenize{fix_viscous::doc}]{\sphinxcrossref{\DUrole{doc}{fix viscous}}}}

\item {} 
\sphinxAtStartPar
{\hyperref[\detokenize{fix_wall_gran::doc}]{\sphinxcrossref{\DUrole{doc}{fix wall/gran}}}}

\item {} 
\sphinxAtStartPar
{\hyperref[\detokenize{fix_wall_gran_region::doc}]{\sphinxcrossref{\DUrole{doc}{fix wall/gran/region}}}}

\end{itemize}

\sphinxAtStartPar
The fix style \sphinxstyleemphasis{freeze} zeroes both the force and torque of frozen
atoms, and should be used for granular system instead of the fix style
\sphinxstyleemphasis{setforce}.

\sphinxAtStartPar
To model heat conduction, one must add the temperature and heatflow
atom variables with:
\begin{itemize}
\item {} 
\sphinxAtStartPar
{\hyperref[\detokenize{fix_property_atom::doc}]{\sphinxcrossref{\DUrole{doc}{fix property/atom}}}}

\end{itemize}

\sphinxAtStartPar
a temperature integration fix
\begin{itemize}
\item {} 
\sphinxAtStartPar
{\hyperref[\detokenize{fix_heat_flow::doc}]{\sphinxcrossref{\DUrole{doc}{fix heat/flow}}}}

\end{itemize}

\sphinxAtStartPar
and a heat conduction option defined in both
\begin{itemize}
\item {} 
\sphinxAtStartPar
{\hyperref[\detokenize{pair_granular::doc}]{\sphinxcrossref{\DUrole{doc}{pair\_style granular}}}}

\item {} 
\sphinxAtStartPar
{\hyperref[\detokenize{fix_wall_gran::doc}]{\sphinxcrossref{\DUrole{doc}{fix wall/gran}}}}

\end{itemize}

\sphinxAtStartPar
For computational efficiency, you can eliminate needless pairwise
computations between frozen atoms by using this command:
\begin{itemize}
\item {} 
\sphinxAtStartPar
{\hyperref[\detokenize{neigh_modify::doc}]{\sphinxcrossref{\DUrole{doc}{neigh\_modify}}}} exclude

\end{itemize}

\begin{sphinxadmonition}{note}{Note:}
\sphinxAtStartPar
By default, for 2d systems, granular particles are still modeled
as 3d spheres, not 2d discs (circles), meaning their moment of inertia
will be the same as in 3d.  If you wish to model granular particles in
2d as 2d discs, see the note on this topic on the {\hyperref[\detokenize{Howto_2d::doc}]{\sphinxcrossref{\DUrole{doc}{Howto 2d}}}}
doc page, where 2d simulations are discussed.
\end{sphinxadmonition}

\sphinxAtStartPar
To add custom granular contact models, see the
{\hyperref[\detokenize{Modify_gran_sub_mod::doc}]{\sphinxcrossref{\DUrole{doc}{modifying granular sub\sphinxhyphen{}models page}}}}.

\sphinxstepscope


\subsection{Body particles}
\label{\detokenize{Howto_body:body-particles}}\label{\detokenize{Howto_body::doc}}
\sphinxAtStartPar
\sphinxstylestrong{Overview:}

\sphinxAtStartPar
In LAMMPS, body particles are generalized finite\sphinxhyphen{}size particles.
Individual body particles can represent complex entities, such as
surface meshes of discrete points, collections of sub\sphinxhyphen{}particles,
deformable objects, etc.  Note that other kinds of finite\sphinxhyphen{}size
spherical and aspherical particles are also supported by LAMMPS, such
as spheres, ellipsoids, line segments, and triangles, but they are
simpler entities than body particles.  See the {\hyperref[\detokenize{Howto_spherical::doc}]{\sphinxcrossref{\DUrole{doc}{Howto spherical}}}} page for a general overview of all these
particle types.

\sphinxAtStartPar
Body particles are used via the {\hyperref[\detokenize{atom_style::doc}]{\sphinxcrossref{\DUrole{doc}{atom\_style body}}}}
command.  It takes a body style as an argument.  The current body
styles supported by LAMMPS are as follows.  The name in the first
column is used as the \sphinxstyleemphasis{bstyle} argument for the {\hyperref[\detokenize{atom_style::doc}]{\sphinxcrossref{\DUrole{doc}{atom\_style body}}}} command.


\begin{savenotes}\sphinxattablestart
\sphinxthistablewithglobalstyle
\centering
\begin{tabulary}{\linewidth}[t]{TT}
\sphinxtoprule
\sphinxtableatstartofbodyhook
\sphinxAtStartPar
\sphinxstyleemphasis{nparticle}
&
\sphinxAtStartPar
rigid body with N sub\sphinxhyphen{}particles
\\
\sphinxhline
\sphinxAtStartPar
\sphinxstyleemphasis{rounded/polygon}
&
\sphinxAtStartPar
2d polygons with N vertices
\\
\sphinxhline
\sphinxAtStartPar
\sphinxstyleemphasis{rounded/polyhedron}
&
\sphinxAtStartPar
3d polyhedra with N vertices, E edges and F faces
\\
\sphinxbottomrule
\end{tabulary}
\sphinxtableafterendhook\par
\sphinxattableend\end{savenotes}

\sphinxAtStartPar
The body style determines what attributes are stored for each body and
thus how they can be used to compute pairwise body/body or
bond/non\sphinxhyphen{}body (point particle) interactions.  More details of each
style are described below.

\sphinxAtStartPar
More styles may be added in the future.  See the
{\hyperref[\detokenize{Modify_body::doc}]{\sphinxcrossref{\DUrole{doc}{page on creating new body styles}}}} for details on
how to add a new body style to the code.


\bigskip\hrule\bigskip


\sphinxAtStartPar
\sphinxstylestrong{When to use body particles:}

\sphinxAtStartPar
You should not use body particles to model a rigid body made of
simpler particles (e.g. point, sphere, ellipsoid, line segment,
triangular particles), if the interaction between pairs of rigid
bodies is just the summation of pairwise interactions between the
simpler particles.  LAMMPS already supports this kind of model via the
{\hyperref[\detokenize{fix_rigid::doc}]{\sphinxcrossref{\DUrole{doc}{fix rigid}}}} command.  Any of the numerous pair styles
that compute interactions between simpler particles can be used.  The
{\hyperref[\detokenize{fix_rigid::doc}]{\sphinxcrossref{\DUrole{doc}{fix rigid}}}} command time integrates the motion of the
rigid bodies.  All of the standard LAMMPS commands for thermostatting,
adding constraints, performing output, etc will operate as expected on
the simple particles.

\sphinxAtStartPar
By contrast, when body particles are used, LAMMPS treats an entire
body as a single particle for purposes of computing pairwise
interactions, building neighbor lists, migrating particles between
processors, output of particles to a dump file, etc.  This means that
interactions between pairs of bodies or between a body and non\sphinxhyphen{}body
(point) particle need to be encoded in an appropriate pair style.  If
such a pair style were to mimic the {\hyperref[\detokenize{fix_rigid::doc}]{\sphinxcrossref{\DUrole{doc}{fix rigid}}}}
model, it would need to loop over the entire collection of
interactions between pairs of simple particles within the two bodies,
each time a single body/body interaction was computed.

\sphinxAtStartPar
Thus it only makes sense to use body particles and develop such a pair
style, when particle/particle interactions are more complex than what
the {\hyperref[\detokenize{fix_rigid::doc}]{\sphinxcrossref{\DUrole{doc}{fix rigid}}}} command can already calculate.  For
example, consider particles with one or more of the following
attributes:
\begin{itemize}
\item {} 
\sphinxAtStartPar
represented by a surface mesh

\item {} 
\sphinxAtStartPar
represented by a collection of geometric entities (e.g. planes + spheres)

\item {} 
\sphinxAtStartPar
deformable

\item {} 
\sphinxAtStartPar
internal stress that induces fragmentation

\end{itemize}

\sphinxAtStartPar
For these models, the interaction between pairs of particles is likely
to be more complex than the summation of simple pairwise interactions.
An example is contact or frictional forces between particles with
planar surfaces that inter\sphinxhyphen{}penetrate.  Likewise, the body particle may
store internal state, such as a stress tensor used to compute a
fracture criterion.

\sphinxAtStartPar
These are additional LAMMPS commands that can be used with body
particles of different styles


\begin{savenotes}\sphinxattablestart
\sphinxthistablewithglobalstyle
\centering
\begin{tabulary}{\linewidth}[t]{TT}
\sphinxtoprule
\sphinxtableatstartofbodyhook
\sphinxAtStartPar
{\hyperref[\detokenize{fix_nve_body::doc}]{\sphinxcrossref{\DUrole{doc}{fix nve/body}}}}
&
\sphinxAtStartPar
integrate motion of a body particle in NVE ensemble
\\
\sphinxhline
\sphinxAtStartPar
{\hyperref[\detokenize{fix_nvt_body::doc}]{\sphinxcrossref{\DUrole{doc}{fix nvt/body}}}}
&
\sphinxAtStartPar
ditto for NVT ensemble
\\
\sphinxhline
\sphinxAtStartPar
{\hyperref[\detokenize{fix_npt_body::doc}]{\sphinxcrossref{\DUrole{doc}{fix npt/body}}}}
&
\sphinxAtStartPar
ditto for NPT ensemble
\\
\sphinxhline
\sphinxAtStartPar
{\hyperref[\detokenize{fix_nph_body::doc}]{\sphinxcrossref{\DUrole{doc}{fix nph/body}}}}
&
\sphinxAtStartPar
ditto for NPH ensemble
\\
\sphinxhline
\sphinxAtStartPar
{\hyperref[\detokenize{compute_body_local::doc}]{\sphinxcrossref{\DUrole{doc}{compute body/local}}}}
&
\sphinxAtStartPar
store sub\sphinxhyphen{}particle attributes of a body particle
\\
\sphinxhline
\sphinxAtStartPar
{\hyperref[\detokenize{compute_temp_body::doc}]{\sphinxcrossref{\DUrole{doc}{compute temp/body}}}}
&
\sphinxAtStartPar
compute temperature of body particles
\\
\sphinxhline
\sphinxAtStartPar
{\hyperref[\detokenize{dump::doc}]{\sphinxcrossref{\DUrole{doc}{dump local}}}}
&
\sphinxAtStartPar
output sub\sphinxhyphen{}particle attributes of a body particle
\\
\sphinxhline
\sphinxAtStartPar
{\hyperref[\detokenize{dump_image::doc}]{\sphinxcrossref{\DUrole{doc}{dump image}}}}
&
\sphinxAtStartPar
output body particle attributes as an image
\\
\sphinxbottomrule
\end{tabulary}
\sphinxtableafterendhook\par
\sphinxattableend\end{savenotes}

\sphinxAtStartPar
The pair styles currently defined for use with specific body styles
are listed in the sections below.

\sphinxAtStartPar
Note that for all the body styles, if the data file defines a general
triclinic box, then the orientation of the body particle and its
corresponding 6 moments of inertia and other orientation\sphinxhyphen{}dependent
values should reflect the fact the body is defined withing a general
triclinic box with edge vectors \sphinxstylestrong{A},**B**,**C**.  LAMMPS will rotate
the box to convert it to a restricted triclinic box.  This operation
will also rotate the orientation of the body particles.  See the
{\hyperref[\detokenize{Howto_triclinic::doc}]{\sphinxcrossref{\DUrole{doc}{Howto triclinic}}}} doc page for more details.
The sections below highlight the orientation\sphinxhyphen{}dependent values specific
to each body style.


\bigskip\hrule\bigskip


\sphinxAtStartPar
\sphinxstylestrong{Specifics of body style nparticle:}

\sphinxAtStartPar
The \sphinxstyleemphasis{nparticle} body style represents body particles as a rigid body
with a variable number N of sub\sphinxhyphen{}particles.  It is provided as a
vanilla, prototypical example of a body particle, although as
mentioned above, the {\hyperref[\detokenize{fix_rigid::doc}]{\sphinxcrossref{\DUrole{doc}{fix rigid}}}} command already
duplicates its functionality.

\sphinxAtStartPar
The atom\_style body command for this body style takes two additional
arguments:

\begin{sphinxVerbatim}[commandchars=\\\{\}]
atom\PYGZus{}style body nparticle Nmin Nmax
Nmin = minimum \PYGZsh{} of sub\PYGZhy{}particles in any body in the system
Nmax = maximum \PYGZsh{} of sub\PYGZhy{}particles in any body in the system
\end{sphinxVerbatim}

\sphinxAtStartPar
The Nmin and Nmax arguments are used to bound the size of data
structures used internally by each particle.

\sphinxAtStartPar
When the {\hyperref[\detokenize{read_data::doc}]{\sphinxcrossref{\DUrole{doc}{read\_data}}}} command reads a data file for this
body style, the following information must be provided for each entry
in the \sphinxstyleemphasis{Bodies} section of the data file:

\begin{sphinxVerbatim}[commandchars=\\\{\}]
atom\PYGZhy{}ID 1 M
N
ixx iyy izz ixy ixz iyz
x1 y1 z1
...
xN yN zN
\end{sphinxVerbatim}

\sphinxAtStartPar
where M = 6 + 3*N, and N is the number of sub\sphinxhyphen{}particles in the body
particle.

\sphinxAtStartPar
The integer line has a single value N.  The floating point line(s)
list 6 moments of inertia followed by the coordinates of the N
sub\sphinxhyphen{}particles (x1 to zN) as 3N values.  These values can be listed on
as many lines as you wish; see the {\hyperref[\detokenize{read_data::doc}]{\sphinxcrossref{\DUrole{doc}{read\_data}}}} command
for more details.

\sphinxAtStartPar
The 6 moments of inertia (ixx,iyy,izz,ixy,ixz,iyz) should be the
values consistent with the current orientation of the rigid body
around its center of mass.  The values are with respect to the
simulation box XYZ axes, not with respect to the principal axes of the
rigid body itself.  LAMMPS performs the latter calculation internally.

\sphinxAtStartPar
The coordinates of each sub\sphinxhyphen{}particle are specified as its x,y,z
displacement from the center\sphinxhyphen{}of\sphinxhyphen{}mass of the body particle.  The
center\sphinxhyphen{}of\sphinxhyphen{}mass position of the particle is specified by the x,y,z
values in the \sphinxstyleemphasis{Atoms} section of the data file, as is the total mass
of the body particle.

\sphinxAtStartPar
Note that if the data file defines a general triclinic simulation box,
these sub\sphinxhyphen{}particle displacements are orientation\sphinxhyphen{}dependent and, as
mentioned above, should reflect the body particle’s orientation within
the general triclinic box.

\sphinxAtStartPar
The {\hyperref[\detokenize{pair_body_nparticle::doc}]{\sphinxcrossref{\DUrole{doc}{pair\_style body/nparticle}}}} command can be used
with this body style to compute body/body and body/non\sphinxhyphen{}body interactions.


\bigskip\hrule\bigskip


\sphinxAtStartPar
\sphinxstylestrong{Specifics of body style rounded/polygon:}

\sphinxAtStartPar
The \sphinxstyleemphasis{rounded/polygon} body style represents body particles as a 2d
polygon with a variable number of N vertices.  This style can only be
used for 2d models; see the {\hyperref[\detokenize{boundary::doc}]{\sphinxcrossref{\DUrole{doc}{boundary}}}} command.  See the
{\hyperref[\detokenize{pair_body_rounded_polygon::doc}]{\sphinxcrossref{\DUrole{doc}{pair\_style body/rounded/polygon}}}} page for
a diagram of two squares with rounded circles at the vertices.  Special cases
for N = 1 (circle) and N = 2 (rod with rounded ends) can also be specified.

\sphinxAtStartPar
One use of this body style is for 2d discrete element models, as
described in {\hyperref[\detokenize{Howto_body:body-fraige}]{\sphinxcrossref{\DUrole{std,std-ref}{Fraige}}}}.

\sphinxAtStartPar
Similar to body style \sphinxstyleemphasis{nparticle}, the atom\_style body command for
this body style takes two additional arguments:

\begin{sphinxVerbatim}[commandchars=\\\{\}]
atom\PYGZus{}style body rounded/polygon Nmin Nmax
Nmin = minimum \PYGZsh{} of vertices in any body in the system
Nmax = maximum \PYGZsh{} of vertices in any body in the system
\end{sphinxVerbatim}

\sphinxAtStartPar
The Nmin and Nmax arguments are used to bound the size of data
structures used internally by each particle.

\sphinxAtStartPar
When the {\hyperref[\detokenize{read_data::doc}]{\sphinxcrossref{\DUrole{doc}{read\_data}}}} command reads a data file for this
body style, the following information must be provided for each body
in the \sphinxstyleemphasis{Bodies} section of the data file:

\begin{sphinxVerbatim}[commandchars=\\\{\}]
atom\PYGZhy{}ID 1 M
N
ixx iyy izz ixy ixz iyz
x1 y1 z1
...
xN yN zN
diameter
\end{sphinxVerbatim}

\sphinxAtStartPar
where M = 6 + 3*N + 1, and N is the number of vertices in the body
particle.

\sphinxAtStartPar
The integer line has a single value N.  The floating point line(s)
list 6 moments of inertia, followed by the coordinates of the N
vertices (x1 to zN) as 3N values (with z = 0.0 for each), followed by
a diameter value = the rounded diameter of the circle that surrounds
each vertex. The diameter value can be different for each body
particle. These floating\sphinxhyphen{}point values can be listed on as many lines
as you wish; see the {\hyperref[\detokenize{read_data::doc}]{\sphinxcrossref{\DUrole{doc}{read\_data}}}} command for more
details.

\begin{sphinxadmonition}{note}{Note:}
\sphinxAtStartPar
It is important that the vertices for each polygonal body particle be
listed in order around its perimeter, so that edges can be inferred.
LAMMPS does not check that this is the case.
\end{sphinxadmonition}

\sphinxAtStartPar
The 6 moments of inertia (ixx,iyy,izz,ixy,ixz,iyz) should be the
values consistent with the current orientation of the rigid body
around its center of mass.  The values are with respect to the
simulation box XYZ axes, not with respect to the principal axes of the
rigid body itself.  LAMMPS performs the latter calculation internally.

\sphinxAtStartPar
The coordinates of each vertex are specified as its x,y,z displacement
from the center\sphinxhyphen{}of\sphinxhyphen{}mass of the body particle.  The center\sphinxhyphen{}of\sphinxhyphen{}mass
position of the particle is specified by the x,y,z values in the
\sphinxstyleemphasis{Atoms} section of the data file.

\sphinxAtStartPar
For example, the following information would specify a square particle
whose edge length is sqrt(2) and rounded diameter is 1.0.  The
orientation of the square is aligned with the xy coordinate axes which
is consistent with the 6 moments of inertia: ixx iyy izz ixy ixz iyz =
1 1 4 0 0 0. Note that only Izz matters in 2D simulations.

\begin{sphinxVerbatim}[commandchars=\\\{\}]
3 1 19
4
1 1 4 0 0 0
\PYGZhy{}0.7071 \PYGZhy{}0.7071 0
\PYGZhy{}0.7071 0.7071 0
0.7071 0.7071 0
0.7071 \PYGZhy{}0.7071 0
1.0
\end{sphinxVerbatim}

\sphinxAtStartPar
A rod in 2D, whose length is 4.0, mass 1.0, rounded at two ends
by circles of diameter 0.5, is specified as follows:

\begin{sphinxVerbatim}[commandchars=\\\{\}]
1 1 13
2
1 1 1.33333 0 0 0
\PYGZhy{}2 0 0
2 0 0
0.5
\end{sphinxVerbatim}

\sphinxAtStartPar
A disk, whose diameter is 3.0, mass 1.0, is specified as follows:

\begin{sphinxVerbatim}[commandchars=\\\{\}]
1 1 10
1
1 1 4.5 0 0 0
0 0 0
3.0
\end{sphinxVerbatim}

\sphinxAtStartPar
Note that if the data file defines a general triclinic simulation box,
these polygon vertex displacements are orientation\sphinxhyphen{}dependent and, as
mentioned above, should reflect the body particle’s orientation within
the general triclinic box.

\sphinxAtStartPar
The {\hyperref[\detokenize{pair_body_rounded_polygon::doc}]{\sphinxcrossref{\DUrole{doc}{pair\_style body/rounded/polygon}}}}
command can be used with this body style to compute body/body
interactions.  The {\hyperref[\detokenize{fix_wall_body_polygon::doc}]{\sphinxcrossref{\DUrole{doc}{fix wall/body/polygon}}}}
command can be used with this body style to compute the interaction of
body particles with a wall.


\bigskip\hrule\bigskip


\sphinxAtStartPar
\sphinxstylestrong{Specifics of body style rounded/polyhedron:}

\sphinxAtStartPar
The \sphinxstyleemphasis{rounded/polyhedron} body style represents body particles as a 3d
polyhedron with a variable number of N vertices, E edges and F faces.
This style can only be used for 3d models; see the
{\hyperref[\detokenize{boundary::doc}]{\sphinxcrossref{\DUrole{doc}{boundary}}}} command.  See the “pair\_style
body/rounded/polygon” page for a diagram of a two 2d squares with
rounded circles at the vertices.  A 3d cube with rounded spheres at
the 8 vertices and 12 rounded edges would be similar.  Special cases
for N = 1 (sphere) and N = 2 (rod with rounded ends) can also be
specified.

\sphinxAtStartPar
This body style is for 3d discrete element models, as described in
{\hyperref[\detokenize{Howto_body:body-wang}]{\sphinxcrossref{\DUrole{std,std-ref}{Wang}}}}.

\sphinxAtStartPar
Similar to body style \sphinxstyleemphasis{rounded/polygon}, the atom\_style body command
for this body style takes two additional arguments:

\begin{sphinxVerbatim}[commandchars=\\\{\}]
atom\PYGZus{}style body rounded/polyhedron Nmin Nmax
Nmin = minimum \PYGZsh{} of vertices in any body in the system
Nmax = maximum \PYGZsh{} of vertices in any body in the system
\end{sphinxVerbatim}

\sphinxAtStartPar
The Nmin and Nmax arguments are used to bound the size of data
structures used internally by each particle.

\sphinxAtStartPar
When the {\hyperref[\detokenize{read_data::doc}]{\sphinxcrossref{\DUrole{doc}{read\_data}}}} command reads a data file for this
body style, the following information must be provided for each entry
in the \sphinxstyleemphasis{Bodies} section of the data file:

\begin{sphinxVerbatim}[commandchars=\\\{\}]
atom\PYGZhy{}ID 3 M
N E F
ixx iyy izz ixy ixz iyz
x1 y1 z1
...
xN yN zN
0 1
1 2
2 3
...
0 1 2 \PYGZhy{}1
0 2 3 \PYGZhy{}1
...
1 2 3 4
diameter
\end{sphinxVerbatim}

\sphinxAtStartPar
where M = 6 + 3*N + 2*E + 4*F + 1, and N is the number of vertices
in the body particle, E = number of edges, F = number of faces.  For N
= 1 or 2, the format is simpler.  E and F are ignored and no edges or
faces are listed, so that M = 6 + 3*N + 1.

\sphinxAtStartPar
The integer line has three values: number of vertices (N), number of
edges (E) and number of faces (F). The floating point line(s) list 6
moments of inertia followed by the coordinates of the N vertices (x1
to zN) as 3N values, followed by 2E vertex indices corresponding to
the end points of the E edges, then 4*F vertex indices defining F
faces.  The last value is the diameter value = the rounded diameter of
the sphere that surrounds each vertex. The diameter value can be
different for each body particle. These floating\sphinxhyphen{}point values can be
listed on as many lines as you wish; see the {\hyperref[\detokenize{read_data::doc}]{\sphinxcrossref{\DUrole{doc}{read\_data}}}} command for more details.

\sphinxAtStartPar
Note that vertices are numbered from 0 to N\sphinxhyphen{}1 inclusive.  The order of
the 2 vertices in each edge does not matter.  Faces can be triangles
or quadrilaterals.  In both cases 4 vertices must be specified.  For a
triangle the 4th vertex is \sphinxhyphen{}1.  The 4 vertices within each triangle or
quadrilateral face should be ordered by the right\sphinxhyphen{}hand rule so that
the normal vector of the face points outwards from the center of mass.
For polyhedron with faces with more than 4 vertices, you should split
the complex face into multiple simple faces, each of which is a
triangle or quadrilateral.

\begin{sphinxadmonition}{note}{Note:}
\sphinxAtStartPar
If a face is a quadrilateral then its 4 vertices must be co\sphinxhyphen{}planar.
LAMMPS does not check that this is the case.  If you have a quad\sphinxhyphen{}face
of a polyhedron that is not planar (e.g. a cube whose vertices have
been randomly displaced), then you should represent the single quad
face as two triangle faces instead.
\end{sphinxadmonition}

\sphinxAtStartPar
The 6 moments of inertia (ixx,iyy,izz,ixy,ixz,iyz) should be the
values consistent with the current orientation of the rigid body
around its center of mass.  The values are with respect to the
simulation box XYZ axes, not with respect to the principal axes of the
rigid body itself.  LAMMPS performs the latter calculation internally.

\sphinxAtStartPar
The coordinates of each vertex are specified as its x,y,z displacement
from the center\sphinxhyphen{}of\sphinxhyphen{}mass of the body particle.  The center\sphinxhyphen{}of\sphinxhyphen{}mass
position of the particle is specified by the x,y,z values in the
\sphinxstyleemphasis{Atoms} section of the data file.

\sphinxAtStartPar
For example, the following information would specify a cubic particle
whose edge length is 2.0 and rounded diameter is 0.5.
The orientation of the cube is aligned with the xyz coordinate axes
which is consistent with the 6 moments of inertia: ixx iyy izz ixy ixz
iyz = 0.667 0.667 0.667 0 0 0.

\begin{sphinxVerbatim}[commandchars=\\\{\}]
1 3 79
8 12 6
0.667 0.667 0.667 0 0 0
1 1 1
1 \PYGZhy{}1 1
\PYGZhy{}1 \PYGZhy{}1 1
\PYGZhy{}1 1 1
1 1 \PYGZhy{}1
1 \PYGZhy{}1 \PYGZhy{}1
\PYGZhy{}1 \PYGZhy{}1 \PYGZhy{}1
\PYGZhy{}1 1 \PYGZhy{}1
0 1
1 2
2 3
3 0
4 5
5 6
6 7
7 4
0 4
1 5
2 6
3 7
0 1 2 3
4 5 6 7
0 1 5 4
1 2 6 5
2 3 7 6
3 0 4 7
0.5
\end{sphinxVerbatim}

\sphinxAtStartPar
A rod in 3D, whose length is 4.0, mass 1.0 and rounded at two ends
by circles of diameter 0.5, is specified as follows:

\begin{sphinxVerbatim}[commandchars=\\\{\}]
1 3 13
2 1 1
0 1.33333 1.33333 0 0 0
\PYGZhy{}2 0 0
2 0 0
0.5
\end{sphinxVerbatim}

\sphinxAtStartPar
A sphere whose diameter is 3.0 and mass 1.0, is specified as follows:

\begin{sphinxVerbatim}[commandchars=\\\{\}]
1 3 10
1 1 1
0.9 0.9 0.9 0 0 0
0 0 0
3.0
\end{sphinxVerbatim}

\sphinxAtStartPar
The number of edges and faces for a rod or sphere must be listed,
but is ignored.

\sphinxAtStartPar
Note that if the data file defines a general triclinic simulation box,
these polyhedron vertex displacements are orientation\sphinxhyphen{}dependent and,
as mentioned above, should reflect the body particle’s orientation
within the general triclinic box.

\sphinxAtStartPar
The {\hyperref[\detokenize{pair_body_rounded_polyhedron::doc}]{\sphinxcrossref{\DUrole{doc}{pair\_style body/rounded/polhedron}}}} command can be used with this body
style to compute body/body interactions.  The {\hyperref[\detokenize{fix_wall_body_polygon::doc}]{\sphinxcrossref{\DUrole{doc}{fix
wall/body/polyhedron}}}} command can be used with
this body style to compute the interaction of body particles with a
wall.


\bigskip\hrule\bigskip


\sphinxAtStartPar
\sphinxstylestrong{Output specifics for all body styles:}

\sphinxAtStartPar
For the {\hyperref[\detokenize{compute_body_local::doc}]{\sphinxcrossref{\DUrole{doc}{compute body/local}}}} and {\hyperref[\detokenize{dump::doc}]{\sphinxcrossref{\DUrole{doc}{dump
local}}}} commands, all 3 of the body styles described on his page
produces one datum for each of the N vertices (of sub\sphinxhyphen{}particles) in a
body particle.  The datum has 3 values:

\begin{sphinxVerbatim}[commandchars=\\\{\}]
1 = x position of vertex (or sub\PYGZhy{}particle)
2 = y position of vertex
3 = z position of vertex
\end{sphinxVerbatim}

\sphinxAtStartPar
These values are the current position of the vertex within the
simulation domain, not a displacement from the center\sphinxhyphen{}of\sphinxhyphen{}mass (COM) of
the body particle itself.  These values are calculated using the
current COM and orientation of the body particle.

\sphinxAtStartPar
The {\hyperref[\detokenize{dump_image::doc}]{\sphinxcrossref{\DUrole{doc}{dump image}}}} command and its \sphinxstyleemphasis{body} keyword can
be used to render body particles.

\sphinxAtStartPar
For the \sphinxstyleemphasis{nparticle} body style, each body is drawn as a
collection of spheres, one for each sub\sphinxhyphen{}particle.  The size of each
sphere is determined by the \sphinxstyleemphasis{bflag1} parameter for the \sphinxstyleemphasis{body} keyword.
The \sphinxstyleemphasis{bflag2} argument is ignored.

\sphinxAtStartPar
For the \sphinxstyleemphasis{rounded/polygon} body style, each body is drawn as a polygon
with N line segments.  For the \sphinxstyleemphasis{rounded/polyhedron} body style, each
face of each body is drawn as a polygon with N line segments.  The
drawn diameter of each line segment is determined by the \sphinxstyleemphasis{bflag1}
parameter for the \sphinxstyleemphasis{body} keyword.  The \sphinxstyleemphasis{bflag2} argument is ignored.

\sphinxAtStartPar
Note that for both the \sphinxstyleemphasis{rounded/polygon} and \sphinxstyleemphasis{rounded/polyhedron}
styles, line segments are drawn between the pairs of vertices.
Depending on the diameters of the line segments this may be slightly
different than the physical extent of the body as calculated by the
{\hyperref[\detokenize{pair_body_rounded_polygon::doc}]{\sphinxcrossref{\DUrole{doc}{pair\_style rounded/polygon}}}} or
{\hyperref[\detokenize{pair_body_rounded_polyhedron::doc}]{\sphinxcrossref{\DUrole{doc}{pair\_style rounded/polyhedron}}}}
commands.  Conceptually, the pair styles define the surface of a 2d or
3d body by lines or planes that are tangent to the finite\sphinxhyphen{}size spheres
of specified diameter which are placed on each vertex position.


\bigskip\hrule\bigskip

\phantomsection\label{\detokenize{Howto_body:body-fraige}}
\sphinxAtStartPar
\sphinxstylestrong{(Fraige)} F. Y. Fraige, P. A. Langston, A. J. Matchett, J. Dodds,
Particuology, 6, 455 (2008).

\phantomsection\label{\detokenize{Howto_body:body-wang}}
\sphinxAtStartPar
\sphinxstylestrong{(Wang)} J. Wang, H. S. Yu, P. A. Langston, F. Y. Fraige, Granular
Matter, 13, 1 (2011).

\sphinxstepscope


\subsection{Bonded particle models}
\label{\detokenize{Howto_bpm:bonded-particle-models}}\label{\detokenize{Howto_bpm::doc}}
\sphinxAtStartPar
The BPM package implements bonded particle models which can be used to
simulate mesoscale solids.  Solids are constructed as a collection of
particles, which each represent a coarse\sphinxhyphen{}grained region of space much
larger than the atomistic scale.  Particles within a solid region are
then connected by a network of bonds to provide solid elasticity.

\sphinxAtStartPar
Unlike traditional bonds in molecular dynamics, the equilibrium bond
length can vary between bonds. Bonds store the reference state.  This
includes setting the equilibrium length equal to the initial distance
between the two particles, but can also include data on the bond
orientation for rotational models. This produces a stress\sphinxhyphen{}free initial
state. Furthermore, bonds are allowed to break under large strains,
producing fracture. The examples/bpm directory has sample input scripts
for simulations of the fragmentation of an impacted plate and the
pouring of extended, elastic bodies. See {\hyperref[\detokenize{Howto_bpm:howto-clemmer}]{\sphinxcrossref{\DUrole{std,std-ref}{(Clemmer)}}}}
for more general information on the approach and the LAMMPS implementation.
Example movies illustrating some of these capabilities are found at
\sphinxurl{https://www.lammps.org/movies.html\#bpmpackage}.


\bigskip\hrule\bigskip


\sphinxAtStartPar
Bonds can be created using a {\hyperref[\detokenize{read_data::doc}]{\sphinxcrossref{\DUrole{doc}{read data}}}} or
{\hyperref[\detokenize{create_bonds::doc}]{\sphinxcrossref{\DUrole{doc}{create bonds}}}} command. Alternatively, a
{\hyperref[\detokenize{molecule::doc}]{\sphinxcrossref{\DUrole{doc}{molecule}}}} template with bonds can be used with
{\hyperref[\detokenize{fix_deposit::doc}]{\sphinxcrossref{\DUrole{doc}{fix deposit}}}} or {\hyperref[\detokenize{fix_pour::doc}]{\sphinxcrossref{\DUrole{doc}{fix pour}}}} to create
solid grains.

\sphinxAtStartPar
In this implementation, bonds store their reference state when they are
first computed in the setup of the first simulation run. Data is then
preserved across run commands and is written to {\hyperref[\detokenize{restart::doc}]{\sphinxcrossref{\DUrole{doc}{binary restart
files}}}} such that restarting the system will not reset the
reference state of a bond. Bonds that are created midway into a run,
such as those created by pouring grains using {\hyperref[\detokenize{fix_pour::doc}]{\sphinxcrossref{\DUrole{doc}{fix pour}}}}, are initialized on that timestep.


\bigskip\hrule\bigskip


\sphinxAtStartPar
Currently, there are two types of bonds included in the BPM package. The
first bond style, {\hyperref[\detokenize{bond_bpm_spring::doc}]{\sphinxcrossref{\DUrole{doc}{bond bpm/spring}}}}, only applies
pairwise, central body forces. Point particles must have {\hyperref[\detokenize{atom_style::doc}]{\sphinxcrossref{\DUrole{doc}{bond atom
style}}}} and may be thought of as nodes in a spring
network. Alternatively, the second bond style, {\hyperref[\detokenize{bond_bpm_rotational::doc}]{\sphinxcrossref{\DUrole{doc}{bond bpm/rotational}}}}, resolves tangential forces and torques arising
with the shearing, bending, and twisting of the bond due to rotation or
displacement of particles.  Particles are similar to those used in the
{\hyperref[\detokenize{Howto_granular::doc}]{\sphinxcrossref{\DUrole{doc}{granular package}}}}, {\hyperref[\detokenize{atom_style::doc}]{\sphinxcrossref{\DUrole{doc}{atom style sphere}}}}. However, they must also track the current orientation of
particles and store bonds, and therefore use a {\hyperref[\detokenize{atom_style::doc}]{\sphinxcrossref{\DUrole{doc}{bpm/sphere atom
style}}}}.  This also requires a unique integrator {\hyperref[\detokenize{fix_nve_bpm_sphere::doc}]{\sphinxcrossref{\DUrole{doc}{fix
nve/bpm/sphere}}}} which numerically integrates
orientation similar to {\hyperref[\detokenize{fix_nve_asphere::doc}]{\sphinxcrossref{\DUrole{doc}{fix nve/asphere}}}}.

\sphinxAtStartPar
In addition to bond styles, a new pair style {\hyperref[\detokenize{pair_bpm_spring::doc}]{\sphinxcrossref{\DUrole{doc}{pair bpm/spring}}}} was added to accompany the bpm/spring bond
style. This pair style is simply a hookean repulsion with similar
velocity damping as its sister bond style.


\bigskip\hrule\bigskip


\sphinxAtStartPar
Bond data can be output using a combination of standard LAMMPS commands.
A list of IDs for bonded atoms can be generated using the
{\hyperref[\detokenize{compute_property_local::doc}]{\sphinxcrossref{\DUrole{doc}{compute property/local}}}} command.
Various properties of bonds can be computed using the
{\hyperref[\detokenize{compute_bond_local::doc}]{\sphinxcrossref{\DUrole{doc}{compute bond/local}}}} command. This
command allows one to access data saved to the bond’s history,
such as the reference length of the bond. More information on
bond history data can be found on the documentation pages for the specific
BPM bond styles. Finally, this data can be output using a {\hyperref[\detokenize{dump::doc}]{\sphinxcrossref{\DUrole{doc}{dump local}}}}
command. As one may output many columns from the same compute, the
{\hyperref[\detokenize{dump_modify::doc}]{\sphinxcrossref{\DUrole{doc}{dump modify}}}} \sphinxstyleemphasis{colname} option may be used to provide
more helpful column names. An example of this procedure is found in
/examples/bpm/pour/. External software, such as OVITO, can read these dump
files to render bond data.


\bigskip\hrule\bigskip


\sphinxAtStartPar
As bonds can be broken between neighbor list builds, the
{\hyperref[\detokenize{special_bonds::doc}]{\sphinxcrossref{\DUrole{doc}{special\_bonds}}}} command works differently for BPM
bond styles. There are two possible settings which determine how pair
interactions work between bonded particles.  First, one can overlay
pair forces with bond forces such that all bonded particles also
feel pair interactions. This can be accomplished by setting the \sphinxstyleemphasis{overlay/pair}
keyword present in all bpm bond styles to \sphinxstyleemphasis{yes} and requires using the
following special bond settings
\begin{quote}

\begin{sphinxVerbatim}[commandchars=\\\{\}]
\PYG{k}{special\PYGZus{}bonds}\PYG{+w}{ }\PYG{n}{lj}\PYG{o}{/}\PYG{n}{coul}\PYG{+w}{ }\PYG{l+m}{1}\PYG{+w}{ }\PYG{l+m}{1}\PYG{+w}{ }\PYG{l+m}{1}
\end{sphinxVerbatim}
\end{quote}

\sphinxAtStartPar
Alternatively, one can turn off all pair interactions between bonded
particles. Unlike {\hyperref[\detokenize{bond_quartic::doc}]{\sphinxcrossref{\DUrole{doc}{bond quartic}}}}, this is not done
by subtracting pair forces during the bond computation, but rather by
dynamically updating the special bond list. This is the default behavior
of BPM bond styles and is done by updating the 1\sphinxhyphen{}2 special bond list as
bonds break.  To do this, LAMMPS requires {\hyperref[\detokenize{newton::doc}]{\sphinxcrossref{\DUrole{doc}{newton}}}} bond off
such that all processors containing an atom know when a bond breaks.
Additionally, one must use the following special bond settings
\begin{quote}

\begin{sphinxVerbatim}[commandchars=\\\{\}]
\PYG{k}{special\PYGZus{}bonds}\PYG{+w}{ }\PYG{n}{lj}\PYG{+w}{ }\PYG{l+m}{0}\PYG{+w}{ }\PYG{l+m}{1}\PYG{+w}{ }\PYG{l+m}{1}\PYG{+w}{ }\PYG{n}{coul}\PYG{+w}{ }\PYG{l+m}{1}\PYG{+w}{ }\PYG{l+m}{1}\PYG{+w}{ }\PYG{l+m}{1}
\end{sphinxVerbatim}
\end{quote}

\sphinxAtStartPar
These settings accomplish two goals. First, they turn off 1\sphinxhyphen{}3 and 1\sphinxhyphen{}4
special bond lists, which are not currently supported for BPMs. As
BPMs often have dense bond networks, generating 1\sphinxhyphen{}3 and 1\sphinxhyphen{}4 special
bond lists is expensive.  By setting the lj weight for 1\sphinxhyphen{}2 bonds to
zero, this turns off pairwise interactions.  Even though there are no
charges in BPM models, setting a nonzero coul weight for 1\sphinxhyphen{}2 bonds
ensures all bonded neighbors are still included in the neighbor list
in case bonds break between neighbor list builds. If bond breakage is
disabled during a simulation run by setting the \sphinxstyleemphasis{break} keyword to \sphinxstyleemphasis{no},
a zero coul weight for 1\sphinxhyphen{}2 bonds can be used to exclude bonded atoms
from the neighbor list builds
\begin{quote}

\begin{sphinxVerbatim}[commandchars=\\\{\}]
\PYG{k}{special\PYGZus{}bonds}\PYG{+w}{ }\PYG{n}{lj}\PYG{+w}{ }\PYG{l+m}{0}\PYG{+w}{ }\PYG{l+m}{1}\PYG{+w}{ }\PYG{l+m}{1}\PYG{+w}{ }\PYG{n}{coul}\PYG{+w}{ }\PYG{l+m}{0}\PYG{+w}{ }\PYG{l+m}{1}\PYG{+w}{ }\PYG{l+m}{1}
\end{sphinxVerbatim}
\end{quote}

\sphinxAtStartPar
This can be useful for post\sphinxhyphen{}processing, or to determine pair interaction
properties between distinct bonded particles.

\sphinxAtStartPar
To monitor the fracture of bonds in the system, all BPM bond styles
have the ability to record instances of bond breakage to output using
the {\hyperref[\detokenize{dump::doc}]{\sphinxcrossref{\DUrole{doc}{dump local}}}} command. Since one may frequently output
a list of broken bonds and the time they broke, the
{\hyperref[\detokenize{dump_modify::doc}]{\sphinxcrossref{\DUrole{doc}{dump modify}}}} option \sphinxstyleemphasis{header no} may be useful to
avoid repeatedly printing the header of the dump file. An example of
this procedure is found in /examples/bpm/impact/. Additionally,
one can use {\hyperref[\detokenize{compute_nbond_atom::doc}]{\sphinxcrossref{\DUrole{doc}{compute nbond/atom}}}} to tally the
current number of bonds per atom.

\sphinxAtStartPar
See the {\hyperref[\detokenize{Howto_broken_bonds::doc}]{\sphinxcrossref{\DUrole{doc}{Howto}}}} page on broken bonds for
more information.


\bigskip\hrule\bigskip


\sphinxAtStartPar
While LAMMPS has many utilities to create and delete bonds, \sphinxstyleemphasis{only}
the following are currently compatible with BPM bond styles:
\begin{itemize}
\item {} 
\sphinxAtStartPar
{\hyperref[\detokenize{create_bonds::doc}]{\sphinxcrossref{\DUrole{doc}{create\_bonds}}}}

\item {} 
\sphinxAtStartPar
{\hyperref[\detokenize{delete_bonds::doc}]{\sphinxcrossref{\DUrole{doc}{delete\_bonds}}}}

\item {} 
\sphinxAtStartPar
{\hyperref[\detokenize{fix_bond_create::doc}]{\sphinxcrossref{\DUrole{doc}{fix bond/create}}}}

\item {} 
\sphinxAtStartPar
{\hyperref[\detokenize{fix_bond_break::doc}]{\sphinxcrossref{\DUrole{doc}{fix bond/break}}}}

\item {} 
\sphinxAtStartPar
{\hyperref[\detokenize{fix_bond_swap::doc}]{\sphinxcrossref{\DUrole{doc}{fix bond/swap}}}}

\end{itemize}

\begin{sphinxadmonition}{note}{Note:}
\sphinxAtStartPar
The {\hyperref[\detokenize{create_bonds::doc}]{\sphinxcrossref{\DUrole{doc}{create\_bonds}}}} command requires certain
{\hyperref[\detokenize{special_bonds::doc}]{\sphinxcrossref{\DUrole{doc}{special\_bonds}}}} settings.  To subtract pair
interactions, one will need to switch between different \sphinxstyleemphasis{special\_bonds}
settings in the input script. An example is found in
\sphinxcode{\sphinxupquote{examples/bpm/impact}}.
\end{sphinxadmonition}


\bigskip\hrule\bigskip

\phantomsection\label{\detokenize{Howto_bpm:howto-clemmer}}
\sphinxAtStartPar
\sphinxstylestrong{(Clemmer)} Clemmer, Monti, Lechman, Soft Matter, 20, 1702 (2024).

\sphinxstepscope


\subsection{Polarizable models}
\label{\detokenize{Howto_polarizable:polarizable-models}}\label{\detokenize{Howto_polarizable::doc}}
\sphinxAtStartPar
In polarizable force fields the charge distributions in molecules and
materials respond to their electrostatic environments. Polarizable
systems can be simulated in LAMMPS using three methods:
\begin{itemize}
\item {} 
\sphinxAtStartPar
the fluctuating charge method, implemented in the {\hyperref[\detokenize{fix_qeq::doc}]{\sphinxcrossref{\DUrole{doc}{QEQ}}}}
package,

\item {} 
\sphinxAtStartPar
the adiabatic core\sphinxhyphen{}shell method, implemented in the
{\hyperref[\detokenize{Howto_coreshell::doc}]{\sphinxcrossref{\DUrole{doc}{CORESHELL}}}} package,

\item {} 
\sphinxAtStartPar
the thermalized Drude dipole method, implemented in the
{\hyperref[\detokenize{Howto_drude::doc}]{\sphinxcrossref{\DUrole{doc}{DRUDE}}}} package.

\end{itemize}

\sphinxAtStartPar
The fluctuating charge method calculates instantaneous charges on
interacting atoms based on the electronegativity equalization
principle. It is implemented in the {\hyperref[\detokenize{fix_qeq::doc}]{\sphinxcrossref{\DUrole{doc}{fix qeq}}}} which is
available in several variants. It is a relatively efficient technique
since no additional particles are introduced. This method allows for
charge transfer between molecules or atom groups. However, because the
charges are located at the interaction sites, off\sphinxhyphen{}plane components of
polarization cannot be represented in planar molecules or atom groups.

\sphinxAtStartPar
The two other methods share the same basic idea: polarizable atoms are
split into one core atom and one satellite particle (called shell or
Drude particle) attached to it by a harmonic spring.  Both atoms bear
a charge and they represent collectively an induced electric dipole.
These techniques are computationally more expensive than the QEq
method because of additional particles and bonds. These two
charge\sphinxhyphen{}on\sphinxhyphen{}spring methods differ in certain features, with the
core\sphinxhyphen{}shell model being normally used for ionic/crystalline materials,
whereas the so\sphinxhyphen{}called Drude model is normally used for molecular
systems and fluid states.

\sphinxAtStartPar
The core\sphinxhyphen{}shell model is applicable to crystalline materials where the
high symmetry around each site leads to stable trajectories of the
core\sphinxhyphen{}shell pairs. However, bonded atoms in molecules can be so close
that a core would interact too strongly or even capture the Drude
particle of a neighbor. The Drude dipole model is relatively more
complex in order to remedy this and other issues. Specifically, the
Drude model includes specific thermostatting of the core\sphinxhyphen{}Drude pairs
and short\sphinxhyphen{}range damping of the induced dipoles.

\sphinxAtStartPar
The three polarization methods can be implemented through a
self\sphinxhyphen{}consistent calculation of charges or induced dipoles at each
timestep. In the fluctuating charge scheme this is done by the matrix
inversion method in {\hyperref[\detokenize{fix_qeq::doc}]{\sphinxcrossref{\DUrole{doc}{fix qeq/point}}}}, but for core\sphinxhyphen{}shell
or Drude\sphinxhyphen{}dipoles the relaxed\sphinxhyphen{}dipoles technique would require an slow
iterative procedure. These self\sphinxhyphen{}consistent solutions yield accurate
trajectories since the additional degrees of freedom representing
polarization are massless.  An alternative is to attribute a mass to
the additional degrees of freedom and perform time integration using
an extended Lagrangian technique. For the fluctuating charge scheme
this is done by {\hyperref[\detokenize{fix_qeq::doc}]{\sphinxcrossref{\DUrole{doc}{fix qeq/dynamic}}}}, and for the
charge\sphinxhyphen{}on\sphinxhyphen{}spring models by the methods outlined in the next two
sections. The assignment of masses to the additional degrees of
freedom can lead to unphysical trajectories if care is not exerted in
choosing the parameters of the polarizable models and the simulation
conditions.

\sphinxAtStartPar
In the core\sphinxhyphen{}shell model the vibration of the shells is kept faster
than the ionic vibrations to mimic the fast response of the
polarizable electrons.  But in molecular systems thermalizing the
core\sphinxhyphen{}Drude pairs at temperatures comparable to the rest of the
simulation leads to several problems (kinetic energy transfer, too
short a timestep, etc.) In order to avoid these problems the relative
motion of the Drude particles with respect to their cores is kept
“cold” so the vibration of the core\sphinxhyphen{}Drude pairs is very slow,
approaching the self\sphinxhyphen{}consistent regime.  In both models the
temperature is regulated using the velocities of the center of mass of
core+shell (or Drude) pairs, but in the Drude model the actual
relative core\sphinxhyphen{}Drude particle motion is thermostatted separately as
well.

\sphinxstepscope


\subsection{Adiabatic core/shell model}
\label{\detokenize{Howto_coreshell:adiabatic-core-shell-model}}\label{\detokenize{Howto_coreshell::doc}}
\sphinxAtStartPar
The adiabatic core\sphinxhyphen{}shell model by {\hyperref[\detokenize{Howto_coreshell:mitchellfincham}]{\sphinxcrossref{\DUrole{std,std-ref}{Mitchell and Fincham}}}} is a simple method for adding polarizability
to a system.  In order to mimic the electron shell of an ion, a
satellite particle is attached to it. This way the ions are split into
a core and a shell where the latter is meant to react to the
electrostatic environment inducing polarizability.  See the {\hyperref[\detokenize{Howto_polarizable::doc}]{\sphinxcrossref{\DUrole{doc}{Howto polarizable}}}} page for a discussion of all
the polarizable models available in LAMMPS.

\sphinxAtStartPar
Technically, shells are attached to the cores by a spring force f =
k*r where k is a parameterized spring constant and r is the distance
between the core and the shell. The charges of the core and the shell
add up to the ion charge, thus q(ion) = q(core) + q(shell). This
setup introduces the ion polarizability (alpha) given by
alpha = q(shell)\textasciicircum{}2 / k. In a
similar fashion the mass of the ion is distributed on the core and the
shell with the core having the larger mass.

\sphinxAtStartPar
To run this model in LAMMPS, {\hyperref[\detokenize{atom_style::doc}]{\sphinxcrossref{\DUrole{doc}{atom\_style}}}} \sphinxstyleemphasis{full} can
be used since atom charge and bonds are needed.  Each kind of
core/shell pair requires two atom types and a bond type.  The core and
shell of a core/shell pair should be bonded to each other with a
harmonic bond that provides the spring force. For example, a data file
for NaCl, as found in examples/coreshell, has this format:

\begin{sphinxVerbatim}[commandchars=\\\{\}]
432   atoms  \PYGZsh{} core and shell atoms
216   bonds  \PYGZsh{} number of core/shell springs

4     atom types  \PYGZsh{} 2 cores and 2 shells for Na and Cl
2     bond types

0.0 24.09597 xlo xhi
0.0 24.09597 ylo yhi
0.0 24.09597 zlo zhi

Masses       \PYGZsh{} core/shell mass ratio = 0.1

1 20.690784  \PYGZsh{} Na core
2 31.90500   \PYGZsh{} Cl core
3 2.298976   \PYGZsh{} Na shell
4 3.54500    \PYGZsh{} Cl shell

Atoms

1    1    2   1.5005    0.00000000   0.00000000   0.00000000 \PYGZsh{} core of core/shell pair 1
2    1    4  \PYGZhy{}2.5005    0.00000000   0.00000000   0.00000000 \PYGZsh{} shell of core/shell pair 1
3    2    1   1.5056    4.01599500   4.01599500   4.01599500 \PYGZsh{} core of core/shell pair 2
4    2    3  \PYGZhy{}0.5056    4.01599500   4.01599500   4.01599500 \PYGZsh{} shell of core/shell pair 2
(...)

Bonds   \PYGZsh{} Bond topology for spring forces

1     2     1     2   \PYGZsh{} spring for core/shell pair 1
2     2     3     4   \PYGZsh{} spring for core/shell pair 2
(...)
\end{sphinxVerbatim}

\sphinxAtStartPar
Non\sphinxhyphen{}Coulombic (e.g. Lennard\sphinxhyphen{}Jones) pairwise interactions are only
defined between the shells.  Coulombic interactions are defined
between all cores and shells.  If desired, additional bonds can be
specified between cores.

\sphinxAtStartPar
The {\hyperref[\detokenize{special_bonds::doc}]{\sphinxcrossref{\DUrole{doc}{special\_bonds}}}} command should be used to
turn\sphinxhyphen{}off the Coulombic interaction within core/shell pairs, since that
interaction is set by the bond spring.  This is done using the
{\hyperref[\detokenize{special_bonds::doc}]{\sphinxcrossref{\DUrole{doc}{special\_bonds}}}} command with a 1\sphinxhyphen{}2 weight = 0.0,
which is the default value.  It needs to be considered whether one has
to adjust the {\hyperref[\detokenize{special_bonds::doc}]{\sphinxcrossref{\DUrole{doc}{special\_bonds}}}} weighting according
to the molecular topology since the interactions of the shells are
bypassed over an extra bond.

\sphinxAtStartPar
Note that this core/shell implementation does not require all ions to
be polarized.  One can mix core/shell pairs and ions without a
satellite particle if desired.

\sphinxAtStartPar
Since the core/shell model permits distances of r = 0.0 between the
core and shell, a pair style with a “cs” suffix needs to be used to
implement a valid long\sphinxhyphen{}range Coulombic correction.  Several such pair
styles are provided in the CORESHELL package.  See {\hyperref[\detokenize{pair_cs::doc}]{\sphinxcrossref{\DUrole{doc}{this page}}}} for details.  All of the core/shell enabled pair
styles require the use of a long\sphinxhyphen{}range Coulombic solver, as specified
by the {\hyperref[\detokenize{kspace_style::doc}]{\sphinxcrossref{\DUrole{doc}{kspace\_style}}}} command.  Either the PPPM or
Ewald solvers can be used.

\sphinxAtStartPar
For the NaCL example problem, these pair style and bond style settings
are used:

\begin{sphinxVerbatim}[commandchars=\\\{\}]
\PYG{k}{pair\PYGZus{}style}\PYG{+w}{      }\PYG{n}{born}\PYG{o}{/}\PYG{n}{coul}\PYG{o}{/}\PYG{n}{long}\PYG{o}{/}\PYG{n}{cs}\PYG{+w}{ }\PYG{l+m}{20.0}\PYG{+w}{ }\PYG{l+m}{20.0}
\PYG{k}{pair\PYGZus{}coeff}\PYG{+w}{      }\PYG{o}{*}\PYG{+w}{ }\PYG{o}{*}\PYG{+w}{      }\PYG{l+m}{0.0}\PYG{+w}{ }\PYG{l+m}{1.000}\PYG{+w}{   }\PYG{l+m}{0.00}\PYG{+w}{  }\PYG{l+m}{0.00}\PYG{+w}{   }\PYG{l+m}{0.00}
\PYG{k}{pair\PYGZus{}coeff}\PYG{+w}{      }\PYG{l+m}{3}\PYG{+w}{ }\PYG{l+m}{3}\PYG{+w}{    }\PYG{l+m}{487.0}\PYG{+w}{ }\PYG{l+m}{0.23768}\PYG{+w}{ }\PYG{l+m}{0.00}\PYG{+w}{  }\PYG{l+m}{1.05}\PYG{+w}{   }\PYG{l+m}{0.50}\PYG{+w}{ }\PYG{c}{\PYGZsh{}Na\PYGZhy{}Na}
\PYG{k}{pair\PYGZus{}coeff}\PYG{+w}{      }\PYG{l+m}{3}\PYG{+w}{ }\PYG{l+m}{4}\PYG{+w}{ }\PYG{l+m}{145134.0}\PYG{+w}{ }\PYG{l+m}{0.23768}\PYG{+w}{ }\PYG{l+m}{0.00}\PYG{+w}{  }\PYG{l+m}{6.99}\PYG{+w}{   }\PYG{l+m}{8.70}\PYG{+w}{ }\PYG{c}{\PYGZsh{}Na\PYGZhy{}Cl}
\PYG{k}{pair\PYGZus{}coeff}\PYG{+w}{      }\PYG{l+m}{4}\PYG{+w}{ }\PYG{l+m}{4}\PYG{+w}{ }\PYG{l+m}{405774.0}\PYG{+w}{ }\PYG{l+m}{0.23768}\PYG{+w}{ }\PYG{l+m}{0.00}\PYG{+w}{ }\PYG{l+m}{72.40}\PYG{+w}{ }\PYG{l+m}{145.40}\PYG{+w}{ }\PYG{c}{\PYGZsh{}Cl\PYGZhy{}Cl}

\PYG{k}{bond\PYGZus{}style}\PYG{+w}{      }\PYG{n}{harmonic}
\PYG{k}{bond\PYGZus{}coeff}\PYG{+w}{      }\PYG{l+m}{1}\PYG{+w}{ }\PYG{l+m}{63.014}\PYG{+w}{ }\PYG{l+m}{0.0}
\PYG{k}{bond\PYGZus{}coeff}\PYG{+w}{      }\PYG{l+m}{2}\PYG{+w}{ }\PYG{l+m}{25.724}\PYG{+w}{ }\PYG{l+m}{0.0}
\end{sphinxVerbatim}

\sphinxAtStartPar
When running dynamics with the adiabatic core/shell model, the
following issues should be considered.  The relative motion of
the core and shell particles corresponds to the polarization,
hereby an instantaneous relaxation of the shells is approximated
and a fast core/shell spring frequency ensures a nearly constant
internal kinetic energy during the simulation.
Thermostats can alter this polarization behavior, by scaling the
internal kinetic energy, meaning the shell will not react freely to
its electrostatic environment.
Therefore it is typically desirable to decouple the relative motion of
the core/shell pair, which is an imaginary degree of freedom, from the
real physical system.  To do that, the {\hyperref[\detokenize{compute_temp_cs::doc}]{\sphinxcrossref{\DUrole{doc}{compute temp/cs}}}} command can be used, in conjunction with
any of the thermostat fixes, such as {\hyperref[\detokenize{fix_nh::doc}]{\sphinxcrossref{\DUrole{doc}{fix nvt}}}} or {\hyperref[\detokenize{fix_langevin::doc}]{\sphinxcrossref{\DUrole{doc}{fix langevin}}}}.  This compute uses the center\sphinxhyphen{}of\sphinxhyphen{}mass velocity
of the core/shell pairs to calculate a temperature, and ensures that
velocity is what is rescaled for thermostatting purposes.  This
compute also works for a system with both core/shell pairs and
non\sphinxhyphen{}polarized ions (ions without an attached satellite particle).  The
{\hyperref[\detokenize{compute_temp_cs::doc}]{\sphinxcrossref{\DUrole{doc}{compute temp/cs}}}} command requires input of two
groups, one for the core atoms, another for the shell atoms.
Non\sphinxhyphen{}polarized ions which might also be included in the treated system
should not be included into either of these groups, they are taken
into account by the \sphinxstyleemphasis{group\sphinxhyphen{}ID} (second argument) of the compute.  The
groups can be defined using the {\hyperref[\detokenize{group::doc}]{\sphinxcrossref{\DUrole{doc}{group *type*}}}} command.
Note that to perform thermostatting using this definition of
temperature, the {\hyperref[\detokenize{fix_modify::doc}]{\sphinxcrossref{\DUrole{doc}{fix modify temp}}}} command should be
used to assign the compute to the thermostat fix.  Likewise the
{\hyperref[\detokenize{thermo_modify::doc}]{\sphinxcrossref{\DUrole{doc}{thermo\_modify temp}}}} command can be used to make
this temperature be output for the overall system.

\sphinxAtStartPar
For the NaCl example, this can be done as follows:

\begin{sphinxVerbatim}[commandchars=\\\{\}]
\PYG{k}{group }\PYG{n+nv+nvIdentifier}{cores}\PYG{+w}{ }\PYG{n}{type}\PYG{+w}{ }\PYG{l+m}{1}\PYG{+w}{ }\PYG{l+m}{2}
\PYG{k}{group }\PYG{n+nv+nvIdentifier}{shells}\PYG{+w}{ }\PYG{n}{type}\PYG{+w}{ }\PYG{l+m}{3}\PYG{+w}{ }\PYG{l+m}{4}
\PYG{k}{compute }\PYG{n+nv+nvIdentifier}{CSequ}\PYG{+w}{ }\PYG{n+nv+nvIdentifier}{all}\PYG{+w}{ }\PYG{n}{temp}\PYG{o}{/}\PYG{n}{cs}\PYG{+w}{ }\PYG{n}{cores}\PYG{+w}{ }\PYG{n}{shells}
\PYG{k}{fix }\PYG{n+nv+nvIdentifier}{thermoberendsen}\PYG{+w}{ }\PYG{n+nv+nvIdentifier}{all}\PYG{+w}{ }\PYG{n}{temp}\PYG{o}{/}\PYG{n}{berendsen}\PYG{+w}{ }\PYG{l+m}{1427}\PYG{+w}{ }\PYG{l+m}{1427}\PYG{+w}{ }\PYG{l+m}{0.4}\PYG{+w}{    }\PYG{c}{\PYGZsh{} thermostat for the true physical system}
\PYG{k}{fix }\PYG{n+nv+nvIdentifier}{thermostatequ}\PYG{+w}{ }\PYG{n+nv+nvIdentifier}{all}\PYG{+w}{ }\PYG{n}{nve}\PYG{+w}{                               }\PYG{c}{\PYGZsh{} integrator as needed for the berendsen thermostat}
\PYG{k}{fix\PYGZus{}modify }\PYG{n+nv+nvIdentifier}{thermoberendsen}\PYG{+w}{ }\PYG{n}{temp}\PYG{+w}{ }\PYG{n}{CSequ}
\PYG{k}{thermo\PYGZus{}modify}\PYG{+w}{ }\PYG{n}{temp}\PYG{+w}{ }\PYG{n}{CSequ}\PYG{+w}{                                }\PYG{c}{\PYGZsh{} output of center\PYGZhy{}of\PYGZhy{}mass derived temperature}
\end{sphinxVerbatim}

\sphinxAtStartPar
The pressure for the core/shell system is computed via the regular
LAMMPS convention by {\hyperref[\detokenize{Howto_coreshell:mitchellfincham2}]{\sphinxcrossref{\DUrole{std,std-ref}{treating the cores and shells as individual particles}}}}. For the thermo output of the pressure
as well as for the application of a barostat, it is necessary to
use an additional {\hyperref[\detokenize{compute_pressure::doc}]{\sphinxcrossref{\DUrole{doc}{pressure}}}} compute based on
the default {\hyperref[\detokenize{compute_temp::doc}]{\sphinxcrossref{\DUrole{doc}{temperature}}}} and specifying it as a
second argument in {\hyperref[\detokenize{fix_modify::doc}]{\sphinxcrossref{\DUrole{doc}{fix modify}}}} and
{\hyperref[\detokenize{thermo_modify::doc}]{\sphinxcrossref{\DUrole{doc}{thermo\_modify}}}} resulting in:

\begin{sphinxVerbatim}[commandchars=\\\{\}]
\PYG{n+nv}{(}\PYG{n+nv}{...}\PYG{n+nv}{)}
\PYG{k}{compute }\PYG{n+nv+nvIdentifier}{CSequ}\PYG{+w}{ }\PYG{n+nv+nvIdentifier}{all}\PYG{+w}{ }\PYG{n}{temp}\PYG{o}{/}\PYG{n}{cs}\PYG{+w}{ }\PYG{n}{cores}\PYG{+w}{ }\PYG{n}{shells}
\PYG{k}{compute }\PYG{n+nv+nvIdentifier}{thermo\PYGZus{}press\PYGZus{}lmp}\PYG{+w}{ }\PYG{n+nv+nvIdentifier}{all}\PYG{+w}{ }\PYG{n}{pressure}\PYG{+w}{ }\PYG{n}{thermo\PYGZus{}temp}\PYG{+w}{       }\PYG{c}{\PYGZsh{} pressure for individual particles}
\PYG{k}{thermo\PYGZus{}modify}\PYG{+w}{ }\PYG{n}{temp}\PYG{+w}{ }\PYG{n}{CSequ}\PYG{+w}{ }\PYG{n}{press}\PYG{+w}{ }\PYG{n}{thermo\PYGZus{}press\PYGZus{}lmp}\PYG{+w}{         }\PYG{c}{\PYGZsh{} modify thermo to regular pressure}
\PYG{k}{fix }\PYG{n+nv+nvIdentifier}{press\PYGZus{}bar}\PYG{+w}{ }\PYG{n+nv+nvIdentifier}{all}\PYG{+w}{ }\PYG{n}{npt}\PYG{+w}{ }\PYG{n}{temp}\PYG{+w}{ }\PYG{l+m}{300}\PYG{+w}{ }\PYG{l+m}{300}\PYG{+w}{ }\PYG{l+m}{0.04}\PYG{+w}{ }\PYG{n}{iso}\PYG{+w}{ }\PYG{l+m}{0}\PYG{+w}{ }\PYG{l+m}{0}\PYG{+w}{ }\PYG{l+m}{0.4}
\PYG{k}{fix\PYGZus{}modify }\PYG{n+nv+nvIdentifier}{press\PYGZus{}bar}\PYG{+w}{ }\PYG{n}{temp}\PYG{+w}{ }\PYG{n}{CSequ}\PYG{+w}{ }\PYG{n}{press}\PYG{+w}{ }\PYG{n}{thermo\PYGZus{}press\PYGZus{}lmp}\PYG{+w}{  }\PYG{c}{\PYGZsh{} pressure modification for correct kinetic scalar}
\end{sphinxVerbatim}

\sphinxAtStartPar
If {\hyperref[\detokenize{compute_temp_cs::doc}]{\sphinxcrossref{\DUrole{doc}{compute temp/cs}}}} is used, the decoupled
relative motion of the core and the shell should in theory be
stable.  However numerical fluctuation can introduce a small
momentum to the system, which is noticeable over long trajectories.
Therefore it is recommendable to use the {\hyperref[\detokenize{fix_momentum::doc}]{\sphinxcrossref{\DUrole{doc}{fix momentum}}}} command in combination with {\hyperref[\detokenize{compute_temp_cs::doc}]{\sphinxcrossref{\DUrole{doc}{compute temp/cs}}}} when equilibrating the system to
prevent any drift.

\sphinxAtStartPar
When initializing the velocities of a system with core/shell pairs, it
is also desirable to not introduce energy into the relative motion of
the core/shell particles, but only assign a center\sphinxhyphen{}of\sphinxhyphen{}mass velocity to
the pairs.  This can be done by using the \sphinxstyleemphasis{bias} keyword of the
{\hyperref[\detokenize{velocity::doc}]{\sphinxcrossref{\DUrole{doc}{velocity create}}}} command and assigning the {\hyperref[\detokenize{compute_temp_cs::doc}]{\sphinxcrossref{\DUrole{doc}{compute temp/cs}}}} command to the \sphinxstyleemphasis{temp} keyword of the
{\hyperref[\detokenize{velocity::doc}]{\sphinxcrossref{\DUrole{doc}{velocity}}}} command, e.g.

\begin{sphinxVerbatim}[commandchars=\\\{\}]
\PYG{k}{velocity }\PYG{n+nv+nvIdentifier}{all}\PYG{+w}{ }\PYG{n}{create}\PYG{+w}{ }\PYG{l+m}{1427}\PYG{+w}{ }\PYG{l+m}{134}\PYG{+w}{ }\PYG{n}{bias}\PYG{+w}{ }\PYG{n}{yes}\PYG{+w}{ }\PYG{n}{temp}\PYG{+w}{ }\PYG{n}{CSequ}
\PYG{k}{velocity }\PYG{n+nv+nvIdentifier}{all}\PYG{+w}{ }\PYG{n}{scale}\PYG{+w}{ }\PYG{l+m}{1427}\PYG{+w}{ }\PYG{n}{temp}\PYG{+w}{ }\PYG{n}{CSequ}
\end{sphinxVerbatim}

\sphinxAtStartPar
To maintain the correct polarizability of the core/shell pairs, the
kinetic energy of the internal motion shall remain nearly constant.
Therefore the choice of spring force and mass ratio need to ensure
much faster relative motion of the two atoms within the core/shell pair
than their center\sphinxhyphen{}of\sphinxhyphen{}mass velocity. This allows the shells to
effectively react instantaneously to the electrostatic environment and
limits energy transfer to or from the core/shell oscillators.
This fast movement also dictates the timestep that can be used.

\sphinxAtStartPar
The primary literature of the adiabatic core/shell model suggests that
the fast relative motion of the core/shell pairs only allows negligible
energy transfer to the environment.
The mentioned energy transfer will typically lead to a small drift
in total energy over time.  This internal energy can be monitored
using the {\hyperref[\detokenize{compute_chunk_atom::doc}]{\sphinxcrossref{\DUrole{doc}{compute chunk/atom}}}} and {\hyperref[\detokenize{compute_temp_chunk::doc}]{\sphinxcrossref{\DUrole{doc}{compute temp/chunk}}}} commands.  The internal kinetic
energies of each core/shell pair can then be summed using the sum()
special function of the {\hyperref[\detokenize{variable::doc}]{\sphinxcrossref{\DUrole{doc}{variable}}}} command.  Or they can
be time/averaged and output using the {\hyperref[\detokenize{fix_ave_time::doc}]{\sphinxcrossref{\DUrole{doc}{fix ave/time}}}}
command.  To use these commands, each core/shell pair must be defined
as a “chunk”.  If each core/shell pair is defined as its own molecule,
the molecule ID can be used to define the chunks.  If cores are bonded
to each other to form larger molecules, the chunks can be identified
by the {\hyperref[\detokenize{fix_property_atom::doc}]{\sphinxcrossref{\DUrole{doc}{fix property/atom}}}} via assigning a
core/shell ID to each atom using a special field in the data file read
by the {\hyperref[\detokenize{read_data::doc}]{\sphinxcrossref{\DUrole{doc}{read\_data}}}} command.  This field can then be
accessed by the {\hyperref[\detokenize{compute_property_atom::doc}]{\sphinxcrossref{\DUrole{doc}{compute property/atom}}}}
command, to use as input to the {\hyperref[\detokenize{compute_chunk_atom::doc}]{\sphinxcrossref{\DUrole{doc}{compute chunk/atom}}}} command to define the core/shell
pairs as chunks.

\sphinxAtStartPar
For example if core/shell pairs are the only molecules:

\begin{sphinxVerbatim}[commandchars=\\\{\}]
\PYG{k}{read\PYGZus{}data}\PYG{+w}{ }\PYG{n}{NaCl\PYGZus{}CS\PYGZus{}x0.1\PYGZus{}prop.data}
\PYG{k}{compute }\PYG{n+nv+nvIdentifier}{prop}\PYG{+w}{ }\PYG{n+nv+nvIdentifier}{all}\PYG{+w}{ }\PYG{n}{property}\PYG{o}{/}\PYG{n}{atom}\PYG{+w}{ }\PYG{n}{molecule}
\PYG{k}{compute }\PYG{n+nv+nvIdentifier}{cs\PYGZus{}chunk}\PYG{+w}{ }\PYG{n+nv+nvIdentifier}{all}\PYG{+w}{ }\PYG{n}{chunk}\PYG{o}{/}\PYG{n}{atom}\PYG{+w}{ }\PYG{n}{c\PYGZus{}prop}
\PYG{k}{compute }\PYG{n+nv+nvIdentifier}{cstherm}\PYG{+w}{ }\PYG{n+nv+nvIdentifier}{all}\PYG{+w}{ }\PYG{n}{temp}\PYG{o}{/}\PYG{n}{chunk}\PYG{+w}{ }\PYG{n}{cs\PYGZus{}chunk}\PYG{+w}{ }\PYG{n}{temp}\PYG{+w}{ }\PYG{n}{internal}\PYG{+w}{ }\PYG{n}{com}\PYG{+w}{ }\PYG{n}{yes}\PYG{+w}{ }\PYG{n}{cdof}\PYG{+w}{ }\PYG{l+m}{3.0}\PYG{+w}{     }\PYG{c}{\PYGZsh{} note the chosen degrees of freedom for the core/shell pairs}
\PYG{k}{fix }\PYG{n+nv+nvIdentifier}{ave\PYGZus{}chunk}\PYG{+w}{ }\PYG{n+nv+nvIdentifier}{all}\PYG{+w}{ }\PYG{n}{ave}\PYG{o}{/}\PYG{n}{time}\PYG{+w}{ }\PYG{l+m}{10}\PYG{+w}{ }\PYG{l+m}{1}\PYG{+w}{ }\PYG{l+m}{10}\PYG{+w}{ }\PYG{n}{c\PYGZus{}cstherm}\PYG{+w}{ }\PYG{n}{file}\PYG{+w}{ }\PYG{n}{chunk.dump}\PYG{+w}{ }\PYG{n}{mode}\PYG{+w}{ }\PYG{n}{vector}
\end{sphinxVerbatim}

\sphinxAtStartPar
For example if core/shell pairs and other molecules are present:

\begin{sphinxVerbatim}[commandchars=\\\{\}]
\PYG{k}{fix }\PYG{n+nv+nvIdentifier}{csinfo}\PYG{+w}{ }\PYG{n+nv+nvIdentifier}{all}\PYG{+w}{ }\PYG{n}{property}\PYG{o}{/}\PYG{n}{atom}\PYG{+w}{ }\PYG{n}{i\PYGZus{}CSID}\PYG{+w}{                       }\PYG{c}{\PYGZsh{} property/atom command}
\PYG{k}{read\PYGZus{}data}\PYG{+w}{ }\PYG{n}{NaCl\PYGZus{}CS\PYGZus{}x0.1\PYGZus{}prop.data}\PYG{+w}{ }\PYG{k}{fix }\PYG{n+nv+nvIdentifier}{csinfo}\PYG{+w}{ }\PYG{n+nv+nvIdentifier}{NULL}\PYG{+w}{ }\PYG{n}{CS}\PYG{o}{\PYGZhy{}}\PYG{n}{Info}\PYG{+w}{  }\PYG{c}{\PYGZsh{} atom property added in the data\PYGZhy{}file}
\PYG{k}{compute }\PYG{n+nv+nvIdentifier}{prop}\PYG{+w}{ }\PYG{n+nv+nvIdentifier}{all}\PYG{+w}{ }\PYG{n}{property}\PYG{o}{/}\PYG{n}{atom}\PYG{+w}{ }\PYG{n}{i\PYGZus{}CSID}
\PYG{n+nv}{(}\PYG{n+nv}{...}\PYG{n+nv}{)}
\end{sphinxVerbatim}

\sphinxAtStartPar
The additional section in the date file would be formatted like this:

\begin{sphinxVerbatim}[commandchars=\\\{\}]
CS\PYGZhy{}Info         \PYGZsh{} header of additional section

1   1           \PYGZsh{} column 1 = atom ID, column 2 = core/shell ID
2   1
3   2
4   2
5   3
6   3
7   4
8   4
(...)
\end{sphinxVerbatim}


\bigskip\hrule\bigskip

\phantomsection\label{\detokenize{Howto_coreshell:mitchellfincham}}
\sphinxAtStartPar
\sphinxstylestrong{(Mitchell and Fincham)} Mitchell, Fincham, J Phys Condensed Matter,
5, 1031\sphinxhyphen{}1038 (1993).

\phantomsection\label{\detokenize{Howto_coreshell:mitchellfincham2}}
\sphinxAtStartPar
\sphinxstylestrong{(Fincham)} Fincham, Mackrodt and Mitchell, J Phys Condensed Matter,
6, 393\sphinxhyphen{}404 (1994).

\sphinxstepscope


\subsection{Drude induced dipoles}
\label{\detokenize{Howto_drude:drude-induced-dipoles}}\label{\detokenize{Howto_drude::doc}}
\sphinxAtStartPar
The thermalized Drude model represents induced dipoles by a pair of
charges (the core atom and the Drude particle) connected by a harmonic
spring.  See the {\hyperref[\detokenize{Howto_polarizable::doc}]{\sphinxcrossref{\DUrole{doc}{Howto polarizable}}}} doc page
for a discussion of all the polarizable models available in LAMMPS.

\sphinxAtStartPar
The Drude model has a number of features aimed at its use in
molecular systems ({\hyperref[\detokenize{Howto_drude:howto-lamoureux}]{\sphinxcrossref{\DUrole{std,std-ref}{Lamoureux and Roux}}}}):
\begin{itemize}
\item {} 
\sphinxAtStartPar
Thermostatting of the additional degrees of freedom associated with the
induced dipoles at very low temperature, in terms of the reduced
coordinates of the Drude particles with respect to their cores. This
makes the trajectory close to that of relaxed induced dipoles.

\item {} 
\sphinxAtStartPar
Consistent definition of 1\sphinxhyphen{}2 to 1\sphinxhyphen{}4 neighbors. A core\sphinxhyphen{}Drude particle
pair represents a single (polarizable) atom, so the special screening
factors in a covalent structure should be the same for the core and
the Drude particle.  Drude particles have to inherit the 1\sphinxhyphen{}2, 1\sphinxhyphen{}3, 1\sphinxhyphen{}4
special neighbor relations from their respective cores.

\item {} 
\sphinxAtStartPar
Stabilization of the interactions between induced dipoles. Drude
dipoles on covalently bonded atoms interact too strongly due to the
short distances, so an atom may capture the Drude particle of a
neighbor, or the induced dipoles within the same molecule may align
too much. To avoid this, damping at short range can be done by Thole
functions (for which there are physical grounds). This Thole damping
is applied to the point charges composing the induced dipole (the
charge of the Drude particle and the opposite charge on the core, not
to the total charge of the core atom).

\end{itemize}

\sphinxAtStartPar
A detailed tutorial covering the usage of Drude induced dipoles in
LAMMPS is on the {\hyperref[\detokenize{Howto_drude2::doc}]{\sphinxcrossref{\DUrole{doc}{here}}}}.

\sphinxAtStartPar
As with the core\sphinxhyphen{}shell model, the cores and Drude particles should
appear in the data file as standard atoms. The same holds for the
springs between them, which are described by standard harmonic bonds.
The nature of the atoms (core, Drude particle or non\sphinxhyphen{}polarizable) is
specified via the {\hyperref[\detokenize{fix_drude::doc}]{\sphinxcrossref{\DUrole{doc}{fix drude}}}} command.  The special
list of neighbors is automatically refactored to account for the
equivalence of core and Drude particles as regards special 1\sphinxhyphen{}2 to 1\sphinxhyphen{}4
screening. It may be necessary to use the \sphinxstyleemphasis{extra/special/per/atom}
keyword of the {\hyperref[\detokenize{read_data::doc}]{\sphinxcrossref{\DUrole{doc}{read\_data}}}} command. If using {\hyperref[\detokenize{fix_shake::doc}]{\sphinxcrossref{\DUrole{doc}{fix shake}}}}, make sure no Drude particle is in this fix
group.

\sphinxAtStartPar
There are three ways to thermostat the Drude particles at a low
temperature: use either {\hyperref[\detokenize{fix_langevin_drude::doc}]{\sphinxcrossref{\DUrole{doc}{fix langevin/drude}}}}
for a Langevin thermostat, or {\hyperref[\detokenize{fix_drude_transform::doc}]{\sphinxcrossref{\DUrole{doc}{fix drude/transform/*}}}} for a Nose\sphinxhyphen{}Hoover
thermostat, or {\hyperref[\detokenize{fix_tgnh_drude::doc}]{\sphinxcrossref{\DUrole{doc}{fix tgnvt/drude}}}} for a temperature\sphinxhyphen{}grouped Nose\sphinxhyphen{}Hoover thermostat.
The first and third require use of the command {\hyperref[\detokenize{comm_modify::doc}]{\sphinxcrossref{\DUrole{doc}{comm\_modify vel yes}}}}. The second requires two separate integration
fixes like \sphinxstyleemphasis{nvt} or \sphinxstyleemphasis{npt}. The correct temperatures of the reduced
degrees of freedom can be calculated using the {\hyperref[\detokenize{compute_temp_drude::doc}]{\sphinxcrossref{\DUrole{doc}{compute temp/drude}}}}. This requires also to use the
command \sphinxstyleemphasis{comm\_modify vel yes}.

\sphinxAtStartPar
Short\sphinxhyphen{}range damping of the induced dipole interactions can be achieved
using Thole functions through the {\hyperref[\detokenize{pair_thole::doc}]{\sphinxcrossref{\DUrole{doc}{pair style thole}}}} in {\hyperref[\detokenize{pair_hybrid::doc}]{\sphinxcrossref{\DUrole{doc}{pair\_style hybrid/overlay}}}}
with a Coulomb pair style. It may be useful to use \sphinxstyleemphasis{coul/long/cs} or
similar from the CORESHELL package if the core and Drude particle come
too close, which can cause numerical issues.


\bigskip\hrule\bigskip

\phantomsection\label{\detokenize{Howto_drude:howto-lamoureux}}
\sphinxAtStartPar
\sphinxstylestrong{(Lamoureux and Roux)} G. Lamoureux, B. Roux, J. Chem. Phys 119, 3025 (2003)

\sphinxstepscope


\subsection{Tutorial for Thermalized Drude oscillators in LAMMPS}
\label{\detokenize{Howto_drude2:tutorial-for-thermalized-drude-oscillators-in-lammps}}\label{\detokenize{Howto_drude2::doc}}
\sphinxAtStartPar
This tutorial explains how to use Drude oscillators in LAMMPS to
simulate polarizable systems using the DRUDE package. As an
illustration, the input files for a simulation of 250 phenol molecules
are documented. First of all, LAMMPS has to be compiled with the
DRUDE package activated. Then, the data file and input scripts
have to be modified to include the Drude dipoles and how to handle
them.

\sphinxAtStartPar
Example input scripts available: examples/PACKAGES/drude


\bigskip\hrule\bigskip


\sphinxAtStartPar
\sphinxstylestrong{Overview of Drude induced dipoles}

\sphinxAtStartPar
Polarizable atoms acquire an induced electric dipole moment under the
action of an external electric field, for example the electric field
created by the surrounding particles.  Drude oscillators represent
these dipoles by two fixed charges: the core (DC) and the Drude
particle (DP) bound by a harmonic potential. The Drude particle can be
thought of as the electron cloud whose center can be displaced from
the position of the corresponding nucleus.

\sphinxAtStartPar
The sum of the masses of a core\sphinxhyphen{}Drude pair should be the mass of the
initial (unsplit) atom, \(m_C + m_D = m\).  The sum of their charges
should be the charge of the initial (unsplit) atom, \(q_C + q_D = q\).
A harmonic potential between the core and Drude partners should be
present, with force constant \(k_D\) and an equilibrium distance of
zero. The (half\sphinxhyphen{})stiffness of the {\hyperref[\detokenize{bond_harmonic::doc}]{\sphinxcrossref{\DUrole{doc}{harmonic bond}}}}
\(K_D = k_D/2\) and the Drude charge \(q_D\) are related to the atom
polarizability \(\alpha\) by
\begin{equation*}
\begin{split}K_D = \frac 1 2\, \frac {q_D^2} \alpha\end{split}
\end{equation*}
\sphinxAtStartPar
Ideally, the mass of the Drude particle should be small, and the
stiffness of the harmonic bond should be large, so that the Drude
particle remains close to the core. The values of Drude mass, Drude
charge, and force constant can be chosen following different
strategies, as in the following examples of polarizable force
fields:
\begin{itemize}
\item {} 
\sphinxAtStartPar
{\hyperref[\detokenize{Howto_drude2:lamoureux2}]{\sphinxcrossref{\DUrole{std,std-ref}{Lamoureux and Roux}}}} suggest adopting a global half\sphinxhyphen{}stiffness, \(K_D\) = 500 kcal/(mol Ang \({}^2\)) \sphinxhyphen{} which corresponds to a force constant \(k_D\) = 4184 kJ/(mol Ang \({}^2\)) \sphinxhyphen{} for all types of core\sphinxhyphen{}Drude bond, a global mass \(m_D\) = 0.4 g/mol (or u) for all types of Drude particles, and to calculate the Drude charges for individual atom types from the atom polarizabilities using equation (1). This choice is followed in the polarizable CHARMM force field.

\item {} 
\sphinxAtStartPar
Alternately {\hyperref[\detokenize{Howto_drude2:schroeder}]{\sphinxcrossref{\DUrole{std,std-ref}{Schroeder and Steinhauser}}}} suggest adopting a global charge \(q_D\) = \sphinxhyphen{}1.0e and a global mass \(m_D\) = 0.1 g/mol (or u) for all Drude particles, and to calculate the force constant for each type of core\sphinxhyphen{}Drude bond from equation (1). The timesteps used by these authors are between 0.5 and 2 fs, with the degrees of freedom of the Drude oscillators kept cold at 1 K.

\item {} 
\sphinxAtStartPar
In both these force fields hydrogen atoms are treated as non\sphinxhyphen{}polarizable.

\end{itemize}

\sphinxAtStartPar
The motion of of the Drude particles can be calculated by minimizing
the energy of the induced dipoles at each timestep, by an iterative,
self\sphinxhyphen{}consistent procedure. The Drude particles can be massless and
therefore do not contribute to the kinetic energy. However, the
relaxed method is computational slow. An extended\sphinxhyphen{}lagrangian method
can be used to calculate the positions of the Drude particles, but
this requires them to have mass. It is important in this case to
decouple the degrees of freedom associated with the Drude oscillators
from those of the normal atoms. Thermalizing the Drude dipoles at
temperatures comparable to the rest of the simulation leads to several
problems (kinetic energy transfer, very short timestep, etc.), which
can be remedied by the “cold Drude” technique ({\hyperref[\detokenize{Howto_drude2:lamoureux2}]{\sphinxcrossref{\DUrole{std,std-ref}{Lamoureux and Roux}}}}).

\sphinxAtStartPar
Two closely related models are used to represent polarization through
“charges on a spring”: the core\sphinxhyphen{}shell model and the Drude
model. Although the basic idea is the same, the core\sphinxhyphen{}shell model is
normally used for ionic/crystalline materials, whereas the Drude model
is normally used for molecular systems and fluid states. In ionic
crystals the symmetry around each ion and the distance between them
are such that the core\sphinxhyphen{}shell model is sufficiently stable. But to be
applicable to molecular/covalent systems the Drude model includes two
important features:
\begin{enumerate}
\sphinxsetlistlabels{\arabic}{enumi}{enumii}{}{.}%
\item {} 
\sphinxAtStartPar
The possibility to thermostat the additional degrees of freedom associated with the induced dipoles at very low temperature, in terms of the reduced coordinates of the Drude particles with respect to their cores. This makes the trajectory close to that of relaxed induced dipoles.

\item {} 
\sphinxAtStartPar
The Drude dipoles on covalently bonded atoms interact too strongly due to the short distances, so an atom may capture the Drude particle (shell) of a neighbor, or the induced dipoles within the same molecule may align too much.  To avoid this, damping at short of the interactions between the point charges composing the induced dipole can be done by {\hyperref[\detokenize{Howto_drude2:thole2}]{\sphinxcrossref{\DUrole{std,std-ref}{Thole}}}} functions.

\end{enumerate}


\bigskip\hrule\bigskip


\sphinxAtStartPar
\sphinxstylestrong{Preparation of the data file}

\sphinxAtStartPar
The data file is similar to a standard LAMMPS data file for
\sphinxstyleemphasis{atom\_style full}.  The DPs and the \sphinxstyleemphasis{harmonic bonds} connecting them
to their DC should appear in the data file as normal atoms and bonds.

\sphinxAtStartPar
You can use the \sphinxstyleemphasis{polarizer} tool (Python script distributed with the
DRUDE package) to convert a non\sphinxhyphen{}polarizable data file (here
\sphinxstyleemphasis{data.102494.lmp}) to a polarizable data file (\sphinxstyleemphasis{data\sphinxhyphen{}p.lmp})

\begin{sphinxVerbatim}[commandchars=\\\{\}]
polarizer\PYG{+w}{ }\PYGZhy{}q\PYG{+w}{ }\PYGZhy{}f\PYG{+w}{ }phenol.dff\PYG{+w}{ }data.102494.lmp\PYG{+w}{ }data\PYGZhy{}p.lmp
\end{sphinxVerbatim}

\sphinxAtStartPar
This will automatically insert the new atoms and bonds.
The masses and charges of DCs and DPs are computed
from \sphinxstyleemphasis{phenol.dff}, as well as the DC\sphinxhyphen{}DP bond constants.  The file
\sphinxstyleemphasis{phenol.dff} contains the polarizabilities of the atom types
and the mass of the Drude particles, for instance:
\begin{sphinxalltt}
\# units: kJ/mol, A, deg
\# kforce is in the form k/2 r\_D\textasciicircum{}2
\# type  m\_D/u   q\_D/e    k\_D   alpha/A3  thole
OH      0.4    \sphinxhyphen{}1.0    4184.0   0.63     0.67
CA      0.4    \sphinxhyphen{}1.0    4184.0   1.36     2.51
CAI     0.4    \sphinxhyphen{}1.0    4184.0   1.09     2.51
\end{sphinxalltt}

\sphinxAtStartPar
The hydrogen atoms are absent from this file, so they will be treated
as non\sphinxhyphen{}polarizable atoms.  In the non\sphinxhyphen{}polarizable data file
\sphinxstyleemphasis{data.102494.lmp}, atom names corresponding to the atom type numbers
have to be specified as comments at the end of lines of the \sphinxstyleemphasis{Masses}
section.  You probably need to edit it to add these names. It should
look like

\begin{sphinxVerbatim}[commandchars=\\\{\}]
Masses

1 12.011 \PYGZsh{} CAI
2 12.011 \PYGZsh{} CA
3 15.999 \PYGZsh{} OH
4 1.008  \PYGZsh{} HA
5 1.008  \PYGZsh{} HO
\end{sphinxVerbatim}


\bigskip\hrule\bigskip


\sphinxAtStartPar
\sphinxstylestrong{Basic input file}

\sphinxAtStartPar
The atom style should be set to (or derive from) \sphinxstyleemphasis{full}, so that you
can define atomic charges and molecular bonds, angles, dihedrals…

\sphinxAtStartPar
The \sphinxstyleemphasis{polarizer} tool also outputs certain lines related to the input
script (the use of these lines will be explained below).  In order for
LAMMPS to recognize that you are using Drude oscillators, you should
use the fix \sphinxstyleemphasis{drude}. The command is

\begin{sphinxVerbatim}[commandchars=\\\{\}]
\PYG{k}{fix }\PYG{n+nv+nvIdentifier}{DRUDE}\PYG{+w}{ }\PYG{n+nv+nvIdentifier}{all}\PYG{+w}{ }\PYG{n}{drude}\PYG{+w}{ }\PYG{n}{C}\PYG{+w}{ }\PYG{n}{C}\PYG{+w}{ }\PYG{n}{C}\PYG{+w}{ }\PYG{n}{N}\PYG{+w}{ }\PYG{n}{N}\PYG{+w}{ }\PYG{n}{D}\PYG{+w}{ }\PYG{n}{D}\PYG{+w}{ }\PYG{n}{D}
\end{sphinxVerbatim}

\sphinxAtStartPar
The N, C, D following the \sphinxstyleemphasis{drude} keyword have the following meaning:
There is one tag for each atom type. This tag is C for DCs, D for DPs
and N for non\sphinxhyphen{}polarizable atoms.  Here the atom types 1 to 3 (C and O
atoms) are DC, atom types 4 and 5 (H atoms) are non\sphinxhyphen{}polarizable and
the atom types 6 to 8 are the newly created DPs.

\sphinxAtStartPar
By recognizing the fix \sphinxstyleemphasis{drude}, LAMMPS will find and store matching
DC\sphinxhyphen{}DP pairs and will treat DP as equivalent to their DC in the
\sphinxstyleemphasis{special bonds} relations.  It may be necessary to extend the space
for storing such special relations.  In this case extra space should
be reserved by using the \sphinxstyleemphasis{extra/special/per/atom} keyword of either
the {\hyperref[\detokenize{read_data::doc}]{\sphinxcrossref{\DUrole{doc}{read\_data}}}} or {\hyperref[\detokenize{create_box::doc}]{\sphinxcrossref{\DUrole{doc}{create\_box}}}}
command.  With our phenol, there is 1 more special neighbor for which
space is required.  Otherwise LAMMPS crashes and gives the required
value.

\begin{sphinxVerbatim}[commandchars=\\\{\}]
\PYG{k}{read\PYGZus{}data}\PYG{+w}{ }\PYG{n}{data}\PYG{o}{\PYGZhy{}}\PYG{n}{p.lmp}\PYG{+w}{ }\PYG{n}{extra}\PYG{o}{/}\PYG{n}{special}\PYG{o}{/}\PYG{n}{per}\PYG{o}{/}\PYG{n}{atom}\PYG{+w}{ }\PYG{l+m}{1}
\end{sphinxVerbatim}

\sphinxAtStartPar
Let us assume we want to run a simple NVT simulation at 300 K.  Note
that Drude oscillators need to be thermalized at a low temperature in
order to approximate a self\sphinxhyphen{}consistent field (SCF), therefore it is not
possible to simulate an NVE ensemble with this package.  Since dipoles
are approximated by a charged DC\sphinxhyphen{}DP pair, the \sphinxstyleemphasis{pair\_style} must
include Coulomb interactions, for instance \sphinxstyleemphasis{lj/cut/coul/long} with
\sphinxstyleemphasis{kspace\_style pppm}. For example, with a cutoff of 10. and a precision
1.e\sphinxhyphen{}4:

\begin{sphinxVerbatim}[commandchars=\\\{\}]
\PYG{k}{pair\PYGZus{}style}\PYG{+w}{ }\PYG{n}{lj}\PYG{o}{/}\PYG{n}{cut}\PYG{o}{/}\PYG{n}{coul}\PYG{o}{/}\PYG{n}{long}\PYG{+w}{ }\PYG{l+m}{10.0}
\PYG{k}{kspace\PYGZus{}style}\PYG{+w}{ }\PYG{n}{pppm}\PYG{+w}{ }\PYG{l+m}{1.0e\PYGZhy{}4}
\end{sphinxVerbatim}

\sphinxAtStartPar
As compared to the non\sphinxhyphen{}polarizable input file, \sphinxstyleemphasis{pair\_coeff} lines need
to be added for the DPs.  Since the DPs have no Lennard\sphinxhyphen{}Jones
interactions, their \(\epsilon\) is 0. so the only \sphinxstyleemphasis{pair\_coeff} line
that needs to be added is

\begin{sphinxVerbatim}[commandchars=\\\{\}]
\PYG{k}{pair\PYGZus{}coeff}\PYG{+w}{ }\PYG{o}{*}\PYG{+w}{ }\PYG{l+m}{6}\PYG{o}{*}\PYG{+w}{ }\PYG{l+m}{0.0}\PYG{+w}{ }\PYG{l+m}{0.0}\PYG{+w}{ }\PYG{c}{\PYGZsh{} All\PYGZhy{}DPs}
\end{sphinxVerbatim}

\sphinxAtStartPar
Now for the thermalization, the simplest choice is to use the {\hyperref[\detokenize{fix_langevin_drude::doc}]{\sphinxcrossref{\DUrole{doc}{fix langevin/drude}}}}.

\begin{sphinxVerbatim}[commandchars=\\\{\}]
\PYG{k}{fix }\PYG{n+nv+nvIdentifier}{LANG}\PYG{+w}{ }\PYG{n+nv+nvIdentifier}{all}\PYG{+w}{ }\PYG{n}{langevin}\PYG{o}{/}\PYG{n}{drude}\PYG{+w}{ }\PYG{l+m}{300}\PYG{n}{.}\PYG{+w}{ }\PYG{l+m}{100}\PYG{+w}{ }\PYG{l+m}{12435}\PYG{+w}{ }\PYG{l+m}{1}\PYG{n}{.}\PYG{+w}{ }\PYG{l+m}{20}\PYG{+w}{ }\PYG{l+m}{13977}
\end{sphinxVerbatim}

\sphinxAtStartPar
This applies a Langevin thermostat at temperature 300. to the centers
of mass of the DC\sphinxhyphen{}DP pairs, with relaxation time 100 and with random
seed 12345.  This fix applies also a Langevin thermostat at temperature
1. to the relative motion of the DPs around their DCs, with relaxation
time 20 and random seed 13977.  Only the DCs and non\sphinxhyphen{}polarizable
atoms need to be in this fix’s group.  LAMMPS will thermostat the DPs
together with their DC.  For this, ghost atoms need to know their
velocities. Thus you need to add the following command:

\begin{sphinxVerbatim}[commandchars=\\\{\}]
\PYG{k}{comm\PYGZus{}modify}\PYG{+w}{ }\PYG{n}{vel}\PYG{+w}{ }\PYG{n}{yes}
\end{sphinxVerbatim}

\sphinxAtStartPar
In order to avoid that the center of mass of the whole system
drifts due to the random forces of the Langevin thermostat on DCs, you
can add the \sphinxstyleemphasis{zero yes} option at the end of the fix line.

\sphinxAtStartPar
If the fix \sphinxstyleemphasis{shake} is used to constrain the C\sphinxhyphen{}H bonds, it should be
invoked after the fix \sphinxstyleemphasis{langevin/drude} for more accuracy.

\begin{sphinxVerbatim}[commandchars=\\\{\}]
\PYG{k}{fix }\PYG{n+nv+nvIdentifier}{SHAKE}\PYG{+w}{ }\PYG{n+nv+nvIdentifier}{ATOMS}\PYG{+w}{ }\PYG{n}{shake}\PYG{+w}{ }\PYG{l+m}{0.0001}\PYG{+w}{ }\PYG{l+m}{20}\PYG{+w}{ }\PYG{l+m}{0}\PYG{+w}{ }\PYG{n}{t}\PYG{+w}{ }\PYG{l+m}{4}\PYG{+w}{ }\PYG{l+m}{5}
\end{sphinxVerbatim}

\begin{sphinxadmonition}{note}{Note:}
\sphinxAtStartPar
The group of the fix \sphinxstyleemphasis{shake} must not include the DPs.  If the
group \sphinxstyleemphasis{ATOMS} is defined by non\sphinxhyphen{}DPs atom types, you could use
\end{sphinxadmonition}

\sphinxAtStartPar
Since the fix \sphinxstyleemphasis{langevin/drude} does not perform time integration (just
modification of forces but no position/velocity updates), the fix
\sphinxstyleemphasis{nve} should be used in conjunction.

\begin{sphinxVerbatim}[commandchars=\\\{\}]
\PYG{k}{fix }\PYG{n+nv+nvIdentifier}{NVE}\PYG{+w}{ }\PYG{n+nv+nvIdentifier}{all}\PYG{+w}{ }\PYG{n}{nve}
\end{sphinxVerbatim}

\sphinxAtStartPar
To avoid the flying ice cube artifact, where the atoms progressively freeze and the
center of mass of the whole system drifts faster and faster, the \sphinxstyleemphasis{fix momentum}
can be used. For instance:

\begin{sphinxVerbatim}[commandchars=\\\{\}]
\PYG{k}{fix }\PYG{n+nv+nvIdentifier}{MOMENTUM}\PYG{+w}{ }\PYG{n+nv+nvIdentifier}{all}\PYG{+w}{ }\PYG{n}{momentum}\PYG{+w}{ }\PYG{l+m}{100}\PYG{+w}{ }\PYG{n}{linear}\PYG{+w}{ }\PYG{l+m}{1}\PYG{+w}{ }\PYG{l+m}{1}\PYG{+w}{ }\PYG{l+m}{1}
\end{sphinxVerbatim}

\sphinxAtStartPar
Finally, do not forget to update the atom type elements if you use
them in a \sphinxstyleemphasis{dump\_modify … element …} command, by adding the element
type of the DPs. Here for instance

\begin{sphinxVerbatim}[commandchars=\\\{\}]
\PYG{k}{dump }\PYG{n+nv+nvIdentifier}{DUMP}\PYG{+w}{ }\PYG{n+nv+nvIdentifier}{all}\PYG{+w}{ }\PYG{n}{custom}\PYG{+w}{ }\PYG{l+m}{10}\PYG{+w}{ }\PYG{n}{dump.lammpstrj}\PYG{+w}{ }\PYG{n}{id}\PYG{+w}{ }\PYG{n}{mol}\PYG{+w}{ }\PYG{n}{type}\PYG{+w}{ }\PYG{n}{element}\PYG{+w}{ }\PYG{n}{x}\PYG{+w}{ }\PYG{n}{y}\PYG{+w}{ }\PYG{n}{z}\PYG{+w}{ }\PYG{n}{ix}\PYG{+w}{ }\PYG{n}{iy}\PYG{+w}{ }\PYG{n}{iz}
\PYG{k}{dump\PYGZus{}modify }\PYG{n+nv+nvIdentifier}{DUMP}\PYG{+w}{ }\PYG{n}{element}\PYG{+w}{ }\PYG{n}{C}\PYG{+w}{ }\PYG{n}{C}\PYG{+w}{ }\PYG{n}{O}\PYG{+w}{ }\PYG{n}{H}\PYG{+w}{ }\PYG{n}{H}\PYG{+w}{ }\PYG{n}{D}\PYG{+w}{ }\PYG{n}{D}\PYG{+w}{ }\PYG{n}{D}
\end{sphinxVerbatim}

\sphinxAtStartPar
The input file should now be ready for use!

\sphinxAtStartPar
You will notice that the global temperature \sphinxstyleemphasis{thermo\_temp} computed by
LAMMPS is not 300. K as wanted.  This is because LAMMPS treats DPs as
standard atoms in his default compute.  If you want to output the
temperatures of the DC\sphinxhyphen{}DP pair centers of mass and of the DPs relative
to their DCs, you should use the {\hyperref[\detokenize{compute_temp_drude::doc}]{\sphinxcrossref{\DUrole{doc}{compute temp\_drude}}}}

\begin{sphinxVerbatim}[commandchars=\\\{\}]
\PYG{k}{compute }\PYG{n+nv+nvIdentifier}{TDRUDE}\PYG{+w}{ }\PYG{n+nv+nvIdentifier}{all}\PYG{+w}{ }\PYG{n}{temp}\PYG{o}{/}\PYG{n}{drude}
\end{sphinxVerbatim}

\sphinxAtStartPar
And then output the correct temperatures of the Drude oscillators
using \sphinxstyleemphasis{thermo\_style custom} with respectively \sphinxstyleemphasis{c\_TDRUDE{[}1{]}} and
\sphinxstyleemphasis{c\_TDRUDE{[}2{]}}. These should be close to 300.0 and 1.0 on average.

\begin{sphinxVerbatim}[commandchars=\\\{\}]
\PYG{k}{thermo\PYGZus{}style}\PYG{+w}{ }\PYG{n}{custom}\PYG{+w}{ }\PYG{n}{step}\PYG{+w}{ }\PYG{n}{temp}\PYG{+w}{ }\PYG{n}{c\PYGZus{}TDRUDE[1]}\PYG{+w}{ }\PYG{n}{c\PYGZus{}TDRUDE[2]}
\end{sphinxVerbatim}


\bigskip\hrule\bigskip


\sphinxAtStartPar
\sphinxstylestrong{Thole screening}

\sphinxAtStartPar
Dipolar interactions represented by point charges on springs may not
be stable, for example if the atomic polarizability is too high for
instance, a DP can escape from its DC and be captured by another DC,
which makes the force and energy diverge and the simulation
crash. Even without reaching this extreme case, the correlation
between nearby dipoles on the same molecule may be exaggerated.  Often,
special bond relations prevent bonded neighboring atoms to see the
charge of each other’s DP, so that the problem does not always appear.
It is possible to use screened dipole\sphinxhyphen{}dipole interactions by using the
{\hyperref[\detokenize{pair_thole::doc}]{\sphinxcrossref{\DUrole{doc}{*pair\_style thole*}}}}.  This is implemented as a
correction to the Coulomb pair\_styles, which dampens at short distance
the interactions between the charges representing the induced dipoles.
It is to be used as \sphinxstyleemphasis{hybrid/overlay} with any standard \sphinxstyleemphasis{coul} pair
style.  In our example, we would use

\begin{sphinxVerbatim}[commandchars=\\\{\}]
\PYG{k}{pair\PYGZus{}style}\PYG{+w}{ }\PYG{n}{hybrid}\PYG{o}{/}\PYG{n}{overlay}\PYG{+w}{ }\PYG{n}{lj}\PYG{o}{/}\PYG{n}{cut}\PYG{o}{/}\PYG{n}{coul}\PYG{o}{/}\PYG{n}{long}\PYG{+w}{ }\PYG{l+m}{10.0}\PYG{+w}{ }\PYG{n}{thole}\PYG{+w}{ }\PYG{l+m}{2.6}\PYG{+w}{ }\PYG{l+m}{10.0}
\end{sphinxVerbatim}

\sphinxAtStartPar
This tells LAMMPS that we are using two pair\_styles.  The first one is
as above (\sphinxstyleemphasis{lj/cut/coul/long 10.0}).  The second one is a \sphinxstyleemphasis{thole}
pair\_style with default screening factor 2.6 ({\hyperref[\detokenize{Howto_drude2:noskov2}]{\sphinxcrossref{\DUrole{std,std-ref}{Noskov}}}}) and
cutoff 10.0.

\sphinxAtStartPar
Since \sphinxstyleemphasis{hybrid/overlay} does not support mixing rules, the interaction
coefficients of all the pairs of atom types with i \textless{} j should be
explicitly defined.  The output of the \sphinxstyleemphasis{polarizer} script can be used
to complete the \sphinxstyleemphasis{pair\_coeff} section of the input file.  In our
example, this will look like:

\begin{sphinxVerbatim}[commandchars=\\\{\}]
\PYG{k}{pair\PYGZus{}coeff}\PYG{+w}{    }\PYG{l+m}{1}\PYG{+w}{    }\PYG{l+m}{1}\PYG{+w}{ }\PYG{n}{lj}\PYG{o}{/}\PYG{n}{cut}\PYG{o}{/}\PYG{n}{coul}\PYG{o}{/}\PYG{n}{long}\PYG{+w}{    }\PYG{l+m}{0.0700}\PYG{+w}{   }\PYG{l+m}{3.550}
\PYG{k}{pair\PYGZus{}coeff}\PYG{+w}{    }\PYG{l+m}{1}\PYG{+w}{    }\PYG{l+m}{2}\PYG{+w}{ }\PYG{n}{lj}\PYG{o}{/}\PYG{n}{cut}\PYG{o}{/}\PYG{n}{coul}\PYG{o}{/}\PYG{n}{long}\PYG{+w}{    }\PYG{l+m}{0.0700}\PYG{+w}{   }\PYG{l+m}{3.550}
\PYG{k}{pair\PYGZus{}coeff}\PYG{+w}{    }\PYG{l+m}{1}\PYG{+w}{    }\PYG{l+m}{3}\PYG{+w}{ }\PYG{n}{lj}\PYG{o}{/}\PYG{n}{cut}\PYG{o}{/}\PYG{n}{coul}\PYG{o}{/}\PYG{n}{long}\PYG{+w}{    }\PYG{l+m}{0.1091}\PYG{+w}{   }\PYG{l+m}{3.310}
\PYG{k}{pair\PYGZus{}coeff}\PYG{+w}{    }\PYG{l+m}{1}\PYG{+w}{    }\PYG{l+m}{4}\PYG{+w}{ }\PYG{n}{lj}\PYG{o}{/}\PYG{n}{cut}\PYG{o}{/}\PYG{n}{coul}\PYG{o}{/}\PYG{n}{long}\PYG{+w}{    }\PYG{l+m}{0.0458}\PYG{+w}{   }\PYG{l+m}{2.985}
\PYG{k}{pair\PYGZus{}coeff}\PYG{+w}{    }\PYG{l+m}{2}\PYG{+w}{    }\PYG{l+m}{2}\PYG{+w}{ }\PYG{n}{lj}\PYG{o}{/}\PYG{n}{cut}\PYG{o}{/}\PYG{n}{coul}\PYG{o}{/}\PYG{n}{long}\PYG{+w}{    }\PYG{l+m}{0.0700}\PYG{+w}{   }\PYG{l+m}{3.550}
\PYG{k}{pair\PYGZus{}coeff}\PYG{+w}{    }\PYG{l+m}{2}\PYG{+w}{    }\PYG{l+m}{3}\PYG{+w}{ }\PYG{n}{lj}\PYG{o}{/}\PYG{n}{cut}\PYG{o}{/}\PYG{n}{coul}\PYG{o}{/}\PYG{n}{long}\PYG{+w}{    }\PYG{l+m}{0.1091}\PYG{+w}{   }\PYG{l+m}{3.310}
\PYG{k}{pair\PYGZus{}coeff}\PYG{+w}{    }\PYG{l+m}{2}\PYG{+w}{    }\PYG{l+m}{4}\PYG{+w}{ }\PYG{n}{lj}\PYG{o}{/}\PYG{n}{cut}\PYG{o}{/}\PYG{n}{coul}\PYG{o}{/}\PYG{n}{long}\PYG{+w}{    }\PYG{l+m}{0.0458}\PYG{+w}{   }\PYG{l+m}{2.985}
\PYG{k}{pair\PYGZus{}coeff}\PYG{+w}{    }\PYG{l+m}{3}\PYG{+w}{    }\PYG{l+m}{3}\PYG{+w}{ }\PYG{n}{lj}\PYG{o}{/}\PYG{n}{cut}\PYG{o}{/}\PYG{n}{coul}\PYG{o}{/}\PYG{n}{long}\PYG{+w}{    }\PYG{l+m}{0.1700}\PYG{+w}{   }\PYG{l+m}{3.070}
\PYG{k}{pair\PYGZus{}coeff}\PYG{+w}{    }\PYG{l+m}{3}\PYG{+w}{    }\PYG{l+m}{4}\PYG{+w}{ }\PYG{n}{lj}\PYG{o}{/}\PYG{n}{cut}\PYG{o}{/}\PYG{n}{coul}\PYG{o}{/}\PYG{n}{long}\PYG{+w}{    }\PYG{l+m}{0.0714}\PYG{+w}{   }\PYG{l+m}{2.745}
\PYG{k}{pair\PYGZus{}coeff}\PYG{+w}{    }\PYG{l+m}{4}\PYG{+w}{    }\PYG{l+m}{4}\PYG{+w}{ }\PYG{n}{lj}\PYG{o}{/}\PYG{n}{cut}\PYG{o}{/}\PYG{n}{coul}\PYG{o}{/}\PYG{n}{long}\PYG{+w}{    }\PYG{l+m}{0.0300}\PYG{+w}{   }\PYG{l+m}{2.420}
\PYG{k}{pair\PYGZus{}coeff}\PYG{+w}{    }\PYG{o}{*}\PYG{+w}{    }\PYG{l+m}{5}\PYG{+w}{ }\PYG{n}{lj}\PYG{o}{/}\PYG{n}{cut}\PYG{o}{/}\PYG{n}{coul}\PYG{o}{/}\PYG{n}{long}\PYG{+w}{    }\PYG{l+m}{0.0000}\PYG{+w}{   }\PYG{l+m}{0.000}
\PYG{k}{pair\PYGZus{}coeff}\PYG{+w}{    }\PYG{o}{*}\PYG{+w}{   }\PYG{l+m}{6}\PYG{o}{*}\PYG{+w}{ }\PYG{n}{lj}\PYG{o}{/}\PYG{n}{cut}\PYG{o}{/}\PYG{n}{coul}\PYG{o}{/}\PYG{n}{long}\PYG{+w}{    }\PYG{l+m}{0.0000}\PYG{+w}{   }\PYG{l+m}{0.000}
\PYG{k}{pair\PYGZus{}coeff}\PYG{+w}{    }\PYG{l+m}{1}\PYG{+w}{    }\PYG{l+m}{1}\PYG{+w}{ }\PYG{n}{thole}\PYG{+w}{   }\PYG{l+m}{1.090}\PYG{+w}{   }\PYG{l+m}{2.510}
\PYG{k}{pair\PYGZus{}coeff}\PYG{+w}{    }\PYG{l+m}{1}\PYG{+w}{    }\PYG{l+m}{2}\PYG{+w}{ }\PYG{n}{thole}\PYG{+w}{   }\PYG{l+m}{1.218}\PYG{+w}{   }\PYG{l+m}{2.510}
\PYG{k}{pair\PYGZus{}coeff}\PYG{+w}{    }\PYG{l+m}{1}\PYG{+w}{    }\PYG{l+m}{3}\PYG{+w}{ }\PYG{n}{thole}\PYG{+w}{   }\PYG{l+m}{0.829}\PYG{+w}{   }\PYG{l+m}{1.590}
\PYG{k}{pair\PYGZus{}coeff}\PYG{+w}{    }\PYG{l+m}{1}\PYG{+w}{    }\PYG{l+m}{6}\PYG{+w}{ }\PYG{n}{thole}\PYG{+w}{   }\PYG{l+m}{1.090}\PYG{+w}{   }\PYG{l+m}{2.510}
\PYG{k}{pair\PYGZus{}coeff}\PYG{+w}{    }\PYG{l+m}{1}\PYG{+w}{    }\PYG{l+m}{7}\PYG{+w}{ }\PYG{n}{thole}\PYG{+w}{   }\PYG{l+m}{1.218}\PYG{+w}{   }\PYG{l+m}{2.510}
\PYG{k}{pair\PYGZus{}coeff}\PYG{+w}{    }\PYG{l+m}{1}\PYG{+w}{    }\PYG{l+m}{8}\PYG{+w}{ }\PYG{n}{thole}\PYG{+w}{   }\PYG{l+m}{0.829}\PYG{+w}{   }\PYG{l+m}{1.590}
\PYG{k}{pair\PYGZus{}coeff}\PYG{+w}{    }\PYG{l+m}{2}\PYG{+w}{    }\PYG{l+m}{2}\PYG{+w}{ }\PYG{n}{thole}\PYG{+w}{   }\PYG{l+m}{1.360}\PYG{+w}{   }\PYG{l+m}{2.510}
\PYG{k}{pair\PYGZus{}coeff}\PYG{+w}{    }\PYG{l+m}{2}\PYG{+w}{    }\PYG{l+m}{3}\PYG{+w}{ }\PYG{n}{thole}\PYG{+w}{   }\PYG{l+m}{0.926}\PYG{+w}{   }\PYG{l+m}{1.590}
\PYG{k}{pair\PYGZus{}coeff}\PYG{+w}{    }\PYG{l+m}{2}\PYG{+w}{    }\PYG{l+m}{6}\PYG{+w}{ }\PYG{n}{thole}\PYG{+w}{   }\PYG{l+m}{1.218}\PYG{+w}{   }\PYG{l+m}{2.510}
\PYG{k}{pair\PYGZus{}coeff}\PYG{+w}{    }\PYG{l+m}{2}\PYG{+w}{    }\PYG{l+m}{7}\PYG{+w}{ }\PYG{n}{thole}\PYG{+w}{   }\PYG{l+m}{1.360}\PYG{+w}{   }\PYG{l+m}{2.510}
\PYG{k}{pair\PYGZus{}coeff}\PYG{+w}{    }\PYG{l+m}{2}\PYG{+w}{    }\PYG{l+m}{8}\PYG{+w}{ }\PYG{n}{thole}\PYG{+w}{   }\PYG{l+m}{0.926}\PYG{+w}{   }\PYG{l+m}{1.590}
\PYG{k}{pair\PYGZus{}coeff}\PYG{+w}{    }\PYG{l+m}{3}\PYG{+w}{    }\PYG{l+m}{3}\PYG{+w}{ }\PYG{n}{thole}\PYG{+w}{   }\PYG{l+m}{0.630}\PYG{+w}{   }\PYG{l+m}{0.670}
\PYG{k}{pair\PYGZus{}coeff}\PYG{+w}{    }\PYG{l+m}{3}\PYG{+w}{    }\PYG{l+m}{6}\PYG{+w}{ }\PYG{n}{thole}\PYG{+w}{   }\PYG{l+m}{0.829}\PYG{+w}{   }\PYG{l+m}{1.590}
\PYG{k}{pair\PYGZus{}coeff}\PYG{+w}{    }\PYG{l+m}{3}\PYG{+w}{    }\PYG{l+m}{7}\PYG{+w}{ }\PYG{n}{thole}\PYG{+w}{   }\PYG{l+m}{0.926}\PYG{+w}{   }\PYG{l+m}{1.590}
\PYG{k}{pair\PYGZus{}coeff}\PYG{+w}{    }\PYG{l+m}{3}\PYG{+w}{    }\PYG{l+m}{8}\PYG{+w}{ }\PYG{n}{thole}\PYG{+w}{   }\PYG{l+m}{0.630}\PYG{+w}{   }\PYG{l+m}{0.670}
\PYG{k}{pair\PYGZus{}coeff}\PYG{+w}{    }\PYG{l+m}{6}\PYG{+w}{    }\PYG{l+m}{6}\PYG{+w}{ }\PYG{n}{thole}\PYG{+w}{   }\PYG{l+m}{1.090}\PYG{+w}{   }\PYG{l+m}{2.510}
\PYG{k}{pair\PYGZus{}coeff}\PYG{+w}{    }\PYG{l+m}{6}\PYG{+w}{    }\PYG{l+m}{7}\PYG{+w}{ }\PYG{n}{thole}\PYG{+w}{   }\PYG{l+m}{1.218}\PYG{+w}{   }\PYG{l+m}{2.510}
\PYG{k}{pair\PYGZus{}coeff}\PYG{+w}{    }\PYG{l+m}{6}\PYG{+w}{    }\PYG{l+m}{8}\PYG{+w}{ }\PYG{n}{thole}\PYG{+w}{   }\PYG{l+m}{0.829}\PYG{+w}{   }\PYG{l+m}{1.590}
\PYG{k}{pair\PYGZus{}coeff}\PYG{+w}{    }\PYG{l+m}{7}\PYG{+w}{    }\PYG{l+m}{7}\PYG{+w}{ }\PYG{n}{thole}\PYG{+w}{   }\PYG{l+m}{1.360}\PYG{+w}{   }\PYG{l+m}{2.510}
\PYG{k}{pair\PYGZus{}coeff}\PYG{+w}{    }\PYG{l+m}{7}\PYG{+w}{    }\PYG{l+m}{8}\PYG{+w}{ }\PYG{n}{thole}\PYG{+w}{   }\PYG{l+m}{0.926}\PYG{+w}{   }\PYG{l+m}{1.590}
\PYG{k}{pair\PYGZus{}coeff}\PYG{+w}{    }\PYG{l+m}{8}\PYG{+w}{    }\PYG{l+m}{8}\PYG{+w}{ }\PYG{n}{thole}\PYG{+w}{   }\PYG{l+m}{0.630}\PYG{+w}{   }\PYG{l+m}{0.670}
\end{sphinxVerbatim}

\sphinxAtStartPar
For the \sphinxstyleemphasis{thole} pair style the coefficients are
\begin{enumerate}
\sphinxsetlistlabels{\arabic}{enumi}{enumii}{}{.}%
\item {} 
\sphinxAtStartPar
the atom polarizability in units of cubic length

\item {} 
\sphinxAtStartPar
the screening factor of the Thole function (optional, default value
specified by the pair\_style command)

\item {} 
\sphinxAtStartPar
the cutoff (optional, default value defined by the pair\_style command)

\end{enumerate}

\sphinxAtStartPar
The special neighbors have charge\sphinxhyphen{}charge and charge\sphinxhyphen{}dipole
interactions screened by the \sphinxstyleemphasis{coul} factors of the \sphinxstyleemphasis{special\_bonds}
command (0.0, 0.0, and 0.5 in the example above).  Without using the
pair\_style \sphinxstyleemphasis{thole}, dipole\sphinxhyphen{}dipole interactions are screened by the
same factor.  By using the pair\_style \sphinxstyleemphasis{thole}, dipole\sphinxhyphen{}dipole
interactions are screened by Thole’s function, whatever their special
relationship (except within each DC\sphinxhyphen{}DP pair of course).  Consider for
example 1\sphinxhyphen{}2 neighbors: using the pair\_style \sphinxstyleemphasis{thole}, their dipoles
will see each other (despite the \sphinxstyleemphasis{coul} factor being 0.) and the
interactions between these dipoles will be damped by Thole’s function.


\bigskip\hrule\bigskip


\sphinxAtStartPar
\sphinxstylestrong{Thermostats and barostats}

\sphinxAtStartPar
Using a Nose\sphinxhyphen{}Hoover barostat with the \sphinxstyleemphasis{langevin/drude} thermostat is
straightforward using fix \sphinxstyleemphasis{nph} instead of \sphinxstyleemphasis{nve}.  For example:

\begin{sphinxVerbatim}[commandchars=\\\{\}]
\PYG{k}{fix }\PYG{n+nv+nvIdentifier}{NPH}\PYG{+w}{ }\PYG{n+nv+nvIdentifier}{all}\PYG{+w}{ }\PYG{n}{nph}\PYG{+w}{ }\PYG{n}{iso}\PYG{+w}{ }\PYG{l+m}{1}\PYG{n}{.}\PYG{+w}{ }\PYG{l+m}{1}\PYG{n}{.}\PYG{+w}{ }\PYG{l+m}{500}
\end{sphinxVerbatim}

\sphinxAtStartPar
It is also possible to use a Nose\sphinxhyphen{}Hoover instead of a Langevin
thermostat.  This requires to use {\hyperref[\detokenize{fix_drude_transform::doc}]{\sphinxcrossref{\DUrole{doc}{*fix drude/transform*}}}} just before and after the
time integration fixes.  The \sphinxstyleemphasis{fix drude/transform/direct} converts the
atomic masses, positions, velocities and forces into a reduced
representation, where the DCs transform into the centers of mass of
the DC\sphinxhyphen{}DP pairs and the DPs transform into their relative position
with respect to their DC. The \sphinxstyleemphasis{fix drude/transform/inverse} performs
the reverse transformation.  For a NVT simulation, with the DCs and
atoms at 300 K and the DPs at 1 K relative to their DC one would use

\begin{sphinxVerbatim}[commandchars=\\\{\}]
\PYG{k}{fix }\PYG{n+nv+nvIdentifier}{DIRECT}\PYG{+w}{ }\PYG{n+nv+nvIdentifier}{all}\PYG{+w}{ }\PYG{n}{drude}\PYG{o}{/}\PYG{n}{transform}\PYG{o}{/}\PYG{n}{direct}
\PYG{k}{fix }\PYG{n+nv+nvIdentifier}{NVT1}\PYG{+w}{ }\PYG{n+nv+nvIdentifier}{ATOMS}\PYG{+w}{ }\PYG{n}{nvt}\PYG{+w}{ }\PYG{n}{temp}\PYG{+w}{ }\PYG{l+m}{300}\PYG{n}{.}\PYG{+w}{ }\PYG{l+m}{300}\PYG{n}{.}\PYG{+w}{ }\PYG{l+m}{100}
\PYG{k}{fix }\PYG{n+nv+nvIdentifier}{NVT2}\PYG{+w}{ }\PYG{n+nv+nvIdentifier}{DRUDES}\PYG{+w}{ }\PYG{n}{nvt}\PYG{+w}{ }\PYG{n}{temp}\PYG{+w}{ }\PYG{l+m}{1}\PYG{n}{.}\PYG{+w}{ }\PYG{l+m}{1}\PYG{n}{.}\PYG{+w}{ }\PYG{l+m}{20}
\PYG{k}{fix }\PYG{n+nv+nvIdentifier}{INVERSE}\PYG{+w}{ }\PYG{n+nv+nvIdentifier}{all}\PYG{+w}{ }\PYG{n}{drude}\PYG{o}{/}\PYG{n}{transform}\PYG{o}{/}\PYG{n}{inverse}
\end{sphinxVerbatim}

\sphinxAtStartPar
For our phenol example, the groups would be defined as

\begin{sphinxVerbatim}[commandchars=\\\{\}]
\PYG{k}{group }\PYG{n+nv+nvIdentifier}{ATOMS}\PYG{+w}{  }\PYG{n}{type}\PYG{+w}{ }\PYG{l+m}{1}\PYG{+w}{ }\PYG{l+m}{2}\PYG{+w}{ }\PYG{l+m}{3}\PYG{+w}{ }\PYG{l+m}{4}\PYG{+w}{ }\PYG{l+m}{5}\PYG{+w}{ }\PYG{c}{\PYGZsh{} DCs and non\PYGZhy{}polarizable atoms}
\PYG{k}{group }\PYG{n+nv+nvIdentifier}{CORES}\PYG{+w}{  }\PYG{n}{type}\PYG{+w}{ }\PYG{l+m}{1}\PYG{+w}{ }\PYG{l+m}{2}\PYG{+w}{ }\PYG{l+m}{3}\PYG{+w}{     }\PYG{c}{\PYGZsh{} DCs}
\PYG{k}{group }\PYG{n+nv+nvIdentifier}{DRUDES}\PYG{+w}{ }\PYG{n}{type}\PYG{+w}{ }\PYG{l+m}{6}\PYG{+w}{ }\PYG{l+m}{7}\PYG{+w}{ }\PYG{l+m}{8}\PYG{+w}{     }\PYG{c}{\PYGZsh{} DPs}
\end{sphinxVerbatim}

\sphinxAtStartPar
Note that with the fixes \sphinxstyleemphasis{drude/transform}, it is not required to
specify \sphinxstyleemphasis{comm\_modify vel yes} because the fixes do it anyway (several
times and for the forces also).

\sphinxAtStartPar
It is a bit more tricky to run a NPT simulation with Nose\sphinxhyphen{}Hoover
barostat and thermostat.  First, the volume should be integrated only
once. So the fix for DCs and atoms should be \sphinxstyleemphasis{npt} while the fix for
DPs should be \sphinxstyleemphasis{nvt} (or vice versa).  Second, the \sphinxstyleemphasis{fix npt} computes a
global pressure and thus a global temperature whatever the fix group.
We do want the pressure to correspond to the whole system, but we want
the temperature to correspond to the fix group only.  We must then use
the \sphinxstyleemphasis{fix\_modify} command for this.  In the end, the block of
instructions for thermostatting and barostatting will look like

\begin{sphinxVerbatim}[commandchars=\\\{\}]
\PYG{k}{compute }\PYG{n+nv+nvIdentifier}{TATOMS}\PYG{+w}{ }\PYG{n+nv+nvIdentifier}{ATOMS}\PYG{+w}{ }\PYG{n}{temp}
\PYG{k}{fix }\PYG{n+nv+nvIdentifier}{DIRECT}\PYG{+w}{ }\PYG{n+nv+nvIdentifier}{all}\PYG{+w}{ }\PYG{n}{drude}\PYG{o}{/}\PYG{n}{transform}\PYG{o}{/}\PYG{n}{direct}
\PYG{k}{fix }\PYG{n+nv+nvIdentifier}{NPT}\PYG{+w}{ }\PYG{n+nv+nvIdentifier}{ATOMS}\PYG{+w}{ }\PYG{n}{npt}\PYG{+w}{ }\PYG{n}{temp}\PYG{+w}{ }\PYG{l+m}{300}\PYG{n}{.}\PYG{+w}{ }\PYG{l+m}{300}\PYG{n}{.}\PYG{+w}{ }\PYG{l+m}{100}\PYG{+w}{ }\PYG{n}{iso}\PYG{+w}{ }\PYG{l+m}{1}\PYG{n}{.}\PYG{+w}{ }\PYG{l+m}{1}\PYG{n}{.}\PYG{+w}{ }\PYG{l+m}{500}
\PYG{k}{fix\PYGZus{}modify }\PYG{n+nv+nvIdentifier}{NPT}\PYG{+w}{ }\PYG{n}{temp}\PYG{+w}{ }\PYG{n}{TATOMS}\PYG{+w}{ }\PYG{n}{press}\PYG{+w}{ }\PYG{n}{thermo\PYGZus{}press}
\PYG{k}{fix }\PYG{n+nv+nvIdentifier}{NVT}\PYG{+w}{ }\PYG{n+nv+nvIdentifier}{DRUDES}\PYG{+w}{ }\PYG{n}{nvt}\PYG{+w}{ }\PYG{n}{temp}\PYG{+w}{ }\PYG{l+m}{1}\PYG{n}{.}\PYG{+w}{ }\PYG{l+m}{1}\PYG{n}{.}\PYG{+w}{ }\PYG{l+m}{20}
\PYG{k}{fix }\PYG{n+nv+nvIdentifier}{INVERSE}\PYG{+w}{ }\PYG{n+nv+nvIdentifier}{all}\PYG{+w}{ }\PYG{n}{drude}\PYG{o}{/}\PYG{n}{transform}\PYG{o}{/}\PYG{n}{inverse}
\end{sphinxVerbatim}

\sphinxAtStartPar
Another option for thermalizing the Drude model is to use the
temperature\sphinxhyphen{}grouped Nose\sphinxhyphen{}Hoover (TGNH) thermostat proposed by {\hyperref[\detokenize{Howto_drude2:tgnh-son}]{\sphinxcrossref{\DUrole{std,std-ref}{(Son)}}}}.
This is implemented as {\hyperref[\detokenize{fix_tgnh_drude::doc}]{\sphinxcrossref{\DUrole{doc}{fix tgnvt/drude}}}} and {\hyperref[\detokenize{fix_tgnh_drude::doc}]{\sphinxcrossref{\DUrole{doc}{fix tgnpt/drude}}}}.
It separates the kinetic energy into three contributions:
the molecular center of mass (COM) motion, the motion of atoms or atom\sphinxhyphen{}Drude pairs relative to molecular COMs,
and the relative motion of atom\sphinxhyphen{}Drude pairs.
An independent Nose\sphinxhyphen{}Hoover chain is applied to each type of motion.
When TGNH is used, the temperatures of molecular, atomic and Drude motion can be printed out with {\hyperref[\detokenize{thermo_style::doc}]{\sphinxcrossref{\DUrole{doc}{thermo\_style command}}}} command.

\sphinxAtStartPar
NVT simulation with TGNH thermostat

\begin{sphinxVerbatim}[commandchars=\\\{\}]
\PYG{k}{comm\PYGZus{}modify}\PYG{+w}{ }\PYG{n}{vel}\PYG{+w}{ }\PYG{n}{yes}
\PYG{k}{fix }\PYG{n+nv+nvIdentifier}{TGNVT}\PYG{+w}{ }\PYG{n+nv+nvIdentifier}{all}\PYG{+w}{ }\PYG{n}{tgnvt}\PYG{o}{/}\PYG{n}{drude}\PYG{+w}{ }\PYG{n}{temp}\PYG{+w}{ }\PYG{l+m}{300}\PYG{n}{.}\PYG{+w}{ }\PYG{l+m}{300}\PYG{n}{.}\PYG{+w}{ }\PYG{l+m}{100}\PYG{+w}{ }\PYG{l+m}{1}\PYG{n}{.}\PYG{+w}{ }\PYG{l+m}{20}
\PYG{k}{thermo\PYGZus{}style}\PYG{+w}{ }\PYG{n}{custom}\PYG{+w}{ }\PYG{n}{f\PYGZus{}TGNVT[1]}\PYG{+w}{ }\PYG{n}{f\PYGZus{}TGNVT[2]}\PYG{+w}{ }\PYG{n}{f\PYGZus{}TGNVT[3]}
\end{sphinxVerbatim}

\sphinxAtStartPar
NPT simulation with TGNH thermostat

\begin{sphinxVerbatim}[commandchars=\\\{\}]
\PYG{k}{comm\PYGZus{}modify}\PYG{+w}{ }\PYG{n}{vel}\PYG{+w}{ }\PYG{n}{yes}
\PYG{k}{fix }\PYG{n+nv+nvIdentifier}{TGNPT}\PYG{+w}{ }\PYG{n+nv+nvIdentifier}{all}\PYG{+w}{ }\PYG{n}{tgnpt}\PYG{o}{/}\PYG{n}{drude}\PYG{+w}{ }\PYG{n}{temp}\PYG{+w}{ }\PYG{l+m}{300}\PYG{n}{.}\PYG{+w}{ }\PYG{l+m}{300}\PYG{n}{.}\PYG{+w}{ }\PYG{l+m}{100}\PYG{+w}{ }\PYG{l+m}{1}\PYG{n}{.}\PYG{+w}{ }\PYG{l+m}{20}\PYG{+w}{ }\PYG{n}{iso}\PYG{+w}{ }\PYG{l+m}{1}\PYG{n}{.}\PYG{+w}{ }\PYG{l+m}{1}\PYG{n}{.}\PYG{+w}{ }\PYG{l+m}{500}
\PYG{k}{thermo\PYGZus{}style}\PYG{+w}{ }\PYG{n}{custom}\PYG{+w}{ }\PYG{n}{f\PYGZus{}TGNPT[1]}\PYG{+w}{ }\PYG{n}{f\PYGZus{}TGNPT[2]}\PYG{+w}{ }\PYG{n}{f\PYGZus{}TGNPT[3]}
\end{sphinxVerbatim}


\bigskip\hrule\bigskip


\sphinxAtStartPar
\sphinxstylestrong{Rigid bodies}

\sphinxAtStartPar
You may want to simulate molecules as rigid bodies (but polarizable).
Common cases are water models such as {\hyperref[\detokenize{Howto_drude2:swm4-ndp}]{\sphinxcrossref{\DUrole{std,std-ref}{SWM4\sphinxhyphen{}NDP}}}}, which is a
kind of polarizable TIP4P water.  The rigid bodies and the DPs should
be integrated separately, even with the Langevin thermostat.  Let us
review the different thermostats and ensemble combinations.

\sphinxAtStartPar
NVT ensemble using Langevin thermostat:

\begin{sphinxVerbatim}[commandchars=\\\{\}]
\PYG{k}{comm\PYGZus{}modify}\PYG{+w}{ }\PYG{n}{vel}\PYG{+w}{ }\PYG{n}{yes}
\PYG{k}{fix }\PYG{n+nv+nvIdentifier}{LANG}\PYG{+w}{ }\PYG{n+nv+nvIdentifier}{all}\PYG{+w}{ }\PYG{n}{langevin}\PYG{o}{/}\PYG{n}{drude}\PYG{+w}{ }\PYG{l+m}{300}\PYG{n}{.}\PYG{+w}{ }\PYG{l+m}{100}\PYG{+w}{ }\PYG{l+m}{12435}\PYG{+w}{ }\PYG{l+m}{1}\PYG{n}{.}\PYG{+w}{ }\PYG{l+m}{20}\PYG{+w}{ }\PYG{l+m}{13977}
\PYG{k}{fix }\PYG{n+nv+nvIdentifier}{RIGID}\PYG{+w}{ }\PYG{n+nv+nvIdentifier}{ATOMS}\PYG{+w}{ }\PYG{n}{rigid}\PYG{o}{/}\PYG{n}{nve}\PYG{o}{/}\PYG{n}{small}\PYG{+w}{ }\PYG{n}{molecule}
\PYG{k}{fix }\PYG{n+nv+nvIdentifier}{NVE}\PYG{+w}{ }\PYG{n+nv+nvIdentifier}{DRUDES}\PYG{+w}{ }\PYG{n}{nve}
\end{sphinxVerbatim}

\sphinxAtStartPar
NVT ensemble using Nose\sphinxhyphen{}Hoover thermostat:

\begin{sphinxVerbatim}[commandchars=\\\{\}]
\PYG{k}{fix }\PYG{n+nv+nvIdentifier}{DIRECT}\PYG{+w}{ }\PYG{n+nv+nvIdentifier}{all}\PYG{+w}{ }\PYG{n}{drude}\PYG{o}{/}\PYG{n}{transform}\PYG{o}{/}\PYG{n}{direct}
\PYG{k}{fix }\PYG{n+nv+nvIdentifier}{RIGID}\PYG{+w}{ }\PYG{n+nv+nvIdentifier}{ATOMS}\PYG{+w}{ }\PYG{n}{rigid}\PYG{o}{/}\PYG{n}{nvt}\PYG{o}{/}\PYG{n}{small}\PYG{+w}{ }\PYG{n}{molecule}\PYG{+w}{ }\PYG{n}{temp}\PYG{+w}{ }\PYG{l+m}{300}\PYG{n}{.}\PYG{+w}{ }\PYG{l+m}{300}\PYG{n}{.}\PYG{+w}{ }\PYG{l+m}{100}
\PYG{k}{fix }\PYG{n+nv+nvIdentifier}{NVT}\PYG{+w}{ }\PYG{n+nv+nvIdentifier}{DRUDES}\PYG{+w}{ }\PYG{n}{nvt}\PYG{+w}{ }\PYG{n}{temp}\PYG{+w}{ }\PYG{l+m}{1}\PYG{n}{.}\PYG{+w}{ }\PYG{l+m}{1}\PYG{n}{.}\PYG{+w}{ }\PYG{l+m}{20}
\PYG{k}{fix }\PYG{n+nv+nvIdentifier}{INVERSE}\PYG{+w}{ }\PYG{n+nv+nvIdentifier}{all}\PYG{+w}{ }\PYG{n}{drude}\PYG{o}{/}\PYG{n}{transform}\PYG{o}{/}\PYG{n}{inverse}
\end{sphinxVerbatim}

\sphinxAtStartPar
NPT ensemble with Langevin thermostat:

\begin{sphinxVerbatim}[commandchars=\\\{\}]
\PYG{k}{comm\PYGZus{}modify}\PYG{+w}{ }\PYG{n}{vel}\PYG{+w}{ }\PYG{n}{yes}
\PYG{k}{fix }\PYG{n+nv+nvIdentifier}{LANG}\PYG{+w}{ }\PYG{n+nv+nvIdentifier}{all}\PYG{+w}{ }\PYG{n}{langevin}\PYG{o}{/}\PYG{n}{drude}\PYG{+w}{ }\PYG{l+m}{300}\PYG{n}{.}\PYG{+w}{ }\PYG{l+m}{100}\PYG{+w}{ }\PYG{l+m}{12435}\PYG{+w}{ }\PYG{l+m}{1}\PYG{n}{.}\PYG{+w}{ }\PYG{l+m}{20}\PYG{+w}{ }\PYG{l+m}{13977}
\PYG{k}{fix }\PYG{n+nv+nvIdentifier}{RIGID}\PYG{+w}{ }\PYG{n+nv+nvIdentifier}{ATOMS}\PYG{+w}{ }\PYG{n}{rigid}\PYG{o}{/}\PYG{n}{nph}\PYG{o}{/}\PYG{n}{small}\PYG{+w}{ }\PYG{n}{molecule}\PYG{+w}{ }\PYG{n}{iso}\PYG{+w}{ }\PYG{l+m}{1}\PYG{n}{.}\PYG{+w}{ }\PYG{l+m}{1}\PYG{n}{.}\PYG{+w}{ }\PYG{l+m}{500}
\PYG{k}{fix }\PYG{n+nv+nvIdentifier}{NVE}\PYG{+w}{ }\PYG{n+nv+nvIdentifier}{DRUDES}\PYG{+w}{ }\PYG{n}{nve}
\end{sphinxVerbatim}

\sphinxAtStartPar
NPT ensemble using Nose\sphinxhyphen{}Hoover thermostat:

\begin{sphinxVerbatim}[commandchars=\\\{\}]
\PYG{k}{compute }\PYG{n+nv+nvIdentifier}{TATOM}\PYG{+w}{ }\PYG{n+nv+nvIdentifier}{ATOMS}\PYG{+w}{ }\PYG{n}{temp}
\PYG{k}{fix }\PYG{n+nv+nvIdentifier}{DIRECT}\PYG{+w}{ }\PYG{n+nv+nvIdentifier}{all}\PYG{+w}{ }\PYG{n}{drude}\PYG{o}{/}\PYG{n}{transform}\PYG{o}{/}\PYG{n}{direct}
\PYG{k}{fix }\PYG{n+nv+nvIdentifier}{RIGID}\PYG{+w}{ }\PYG{n+nv+nvIdentifier}{ATOMS}\PYG{+w}{ }\PYG{n}{rigid}\PYG{o}{/}\PYG{n}{npt}\PYG{o}{/}\PYG{n}{small}\PYG{+w}{ }\PYG{n}{molecule}\PYG{+w}{ }\PYG{n}{temp}\PYG{+w}{ }\PYG{l+m}{300}\PYG{n}{.}\PYG{+w}{ }\PYG{l+m}{300}\PYG{n}{.}\PYG{+w}{ }\PYG{l+m}{100}\PYG{+w}{ }\PYG{n}{iso}\PYG{+w}{ }\PYG{l+m}{1}\PYG{n}{.}\PYG{+w}{ }\PYG{l+m}{1}\PYG{n}{.}\PYG{+w}{ }\PYG{l+m}{500}
\PYG{k}{fix\PYGZus{}modify }\PYG{n+nv+nvIdentifier}{RIGID}\PYG{+w}{ }\PYG{n}{temp}\PYG{+w}{ }\PYG{n}{TATOM}\PYG{+w}{ }\PYG{n}{press}\PYG{+w}{ }\PYG{n}{thermo\PYGZus{}press}
\PYG{k}{fix }\PYG{n+nv+nvIdentifier}{NVT}\PYG{+w}{ }\PYG{n+nv+nvIdentifier}{DRUDES}\PYG{+w}{ }\PYG{n}{nvt}\PYG{+w}{ }\PYG{n}{temp}\PYG{+w}{ }\PYG{l+m}{1}\PYG{n}{.}\PYG{+w}{ }\PYG{l+m}{1}\PYG{n}{.}\PYG{+w}{ }\PYG{l+m}{20}
\PYG{k}{fix }\PYG{n+nv+nvIdentifier}{INVERSE}\PYG{+w}{ }\PYG{n+nv+nvIdentifier}{all}\PYG{+w}{ }\PYG{n}{drude}\PYG{o}{/}\PYG{n}{transform}\PYG{o}{/}\PYG{n}{inverse}
\end{sphinxVerbatim}


\bigskip\hrule\bigskip

\phantomsection\label{\detokenize{Howto_drude2:lamoureux2}}
\sphinxAtStartPar
\sphinxstylestrong{(Lamoureux and Roux)} Lamoureux and Roux, J Chem Phys, 119, 3025\sphinxhyphen{}3039 (2003)

\phantomsection\label{\detokenize{Howto_drude2:schroeder}}
\sphinxAtStartPar
\sphinxstylestrong{(Schroeder)}  Schroeder and Steinhauser, J Chem Phys, 133,
154511 (2010).

\phantomsection\label{\detokenize{Howto_drude2:thole2}}
\sphinxAtStartPar
\sphinxstylestrong{(Thole)} Chem Phys, 59, 341 (1981).

\phantomsection\label{\detokenize{Howto_drude2:noskov2}}
\sphinxAtStartPar
\sphinxstylestrong{(Noskov)} Noskov, Lamoureux and Roux, J Phys Chem B, 109, 6705 (2005).

\phantomsection\label{\detokenize{Howto_drude2:swm4-ndp}}
\sphinxAtStartPar
\sphinxstylestrong{(SWM4\sphinxhyphen{}NDP)} Lamoureux, Harder, Vorobyov, Roux, MacKerell, Chem Phys
Let, 418, 245\sphinxhyphen{}249 (2006)

\phantomsection\label{\detokenize{Howto_drude2:tgnh-son}}
\sphinxAtStartPar
\sphinxstylestrong{(Son)} Son, McDaniel, Cui and Yethiraj, J Phys Chem Lett, 10, 7523 (2019).

\sphinxstepscope


\subsection{Peridynamics with LAMMPS}
\label{\detokenize{Howto_peri:peridynamics-with-lammps}}\label{\detokenize{Howto_peri::doc}}
\sphinxAtStartPar
This Howto is based on the Sandia report 2010\sphinxhyphen{}5549 by Michael L. Parks,
Pablo Seleson, Steven J. Plimpton, Richard B. Lehoucq, and
Stewart A. Silling.


\subsubsection{Overview}
\label{\detokenize{Howto_peri:overview}}
\sphinxAtStartPar
Peridynamics is a nonlocal extension of classical continuum mechanics.
The discrete peridynamic model has the same computational structure as a
molecular dynamics model.  This Howto provides a brief overview of the
peridynamic model of a continuum, then discusses how the peridynamic
model is discretized within LAMMPS as described in the original article
{\hyperref[\detokenize{Howto_peri:parks2}]{\sphinxcrossref{\DUrole{std,std-ref}{(Parks)}}}}.  An example problem with comments is also
included.


\subsubsection{Quick Start}
\label{\detokenize{Howto_peri:quick-start}}
\sphinxAtStartPar
The peridynamics styles are included in the optional {\hyperref[\detokenize{Packages_details:pkg-peri}]{\sphinxcrossref{\DUrole{std,std-ref}{PERI package}}}}.  If your LAMMPS executable does not already include the
PERI package, you can see the {\hyperref[\detokenize{Build_package::doc}]{\sphinxcrossref{\DUrole{doc}{build instructions for packages}}}} for how to enable the package when compiling a custom
version of LAMMPS from source.

\sphinxAtStartPar
Here is a minimal example for setting up a peridynamics simulation.

\begin{sphinxVerbatim}[commandchars=\\\{\}]
\PYG{k}{units}\PYG{+w}{         }\PYG{n}{si}
\PYG{k}{boundary}\PYG{+w}{      }\PYG{n}{s}\PYG{+w}{ }\PYG{n}{s}\PYG{+w}{ }\PYG{n}{s}
\PYG{k}{lattice}\PYG{+w}{       }\PYG{n}{sc}\PYG{+w}{ }\PYG{l+m}{0.0005}
\PYG{k}{atom\PYGZus{}style}\PYG{+w}{    }\PYG{n}{peri}
\PYG{k}{atom\PYGZus{}modify}\PYG{+w}{   }\PYG{n}{map}\PYG{+w}{ }\PYG{n}{array}
\PYG{k}{neighbor}\PYG{+w}{      }\PYG{l+m}{0.0010}\PYG{+w}{ }\PYG{n}{bin}
\PYG{k}{region        }\PYG{n+nv+nvIdentifier}{target}\PYG{+w}{ }\PYG{n}{cylinder}\PYG{+w}{ }\PYG{n}{y}\PYG{+w}{ }\PYG{l+m}{0.0}\PYG{+w}{ }\PYG{l+m}{0.0}\PYG{+w}{ }\PYG{l+m}{0.0050}\PYG{+w}{ }\PYG{o}{\PYGZhy{}}\PYG{l+m}{0.0050}\PYG{+w}{ }\PYG{l+m}{0.0}\PYG{+w}{ }\PYG{n}{units}\PYG{+w}{ }\PYG{n}{box}
\PYG{k}{create\PYGZus{}box    }\PYG{n+nv+nvIdentifier}{1}\PYG{+w}{ }\PYG{n+nv+nvIdentifier}{target}
\PYG{k}{create\PYGZus{}atoms}\PYG{+w}{  }\PYG{l+m}{1}\PYG{+w}{ }\PYG{k}{region }\PYG{n+nv+nvIdentifier}{target}

\PYG{k}{pair\PYGZus{}style}\PYG{+w}{    }\PYG{n}{peri}\PYG{o}{/}\PYG{n}{pmb}
\PYG{k}{pair\PYGZus{}coeff}\PYG{+w}{    }\PYG{o}{*}\PYG{+w}{ }\PYG{o}{*}\PYG{+w}{ }\PYG{l+m}{1.6863e22}\PYG{+w}{ }\PYG{l+m}{0.0015001}\PYG{+w}{ }\PYG{l+m}{0.0005}\PYG{+w}{ }\PYG{l+m}{0.25}
\PYG{k}{set}\PYG{+w}{           }\PYG{k}{group }\PYG{n+nv+nvIdentifier}{all}\PYG{+w}{ }\PYG{n}{density}\PYG{+w}{ }\PYG{l+m}{2200}
\PYG{k}{set}\PYG{+w}{           }\PYG{k}{group }\PYG{n+nv+nvIdentifier}{all}\PYG{+w}{ }\PYG{n}{volume}\PYG{+w}{ }\PYG{l+m}{1.25e\PYGZhy{}10}
\PYG{k}{velocity      }\PYG{n+nv+nvIdentifier}{all}\PYG{+w}{ }\PYG{n}{set}\PYG{+w}{ }\PYG{l+m}{0.0}\PYG{+w}{ }\PYG{l+m}{0.0}\PYG{+w}{ }\PYG{l+m}{0.0}\PYG{+w}{ }\PYG{n}{sum}\PYG{+w}{ }\PYG{n}{no}\PYG{+w}{ }\PYG{n}{units}\PYG{+w}{ }\PYG{n}{box}
\PYG{k}{fix           }\PYG{n+nv+nvIdentifier}{1}\PYG{+w}{ }\PYG{n+nv+nvIdentifier}{all}\PYG{+w}{ }\PYG{n}{nve}
\PYG{k}{compute       }\PYG{n+nv+nvIdentifier}{1}\PYG{+w}{ }\PYG{n+nv+nvIdentifier}{all}\PYG{+w}{ }\PYG{n}{damage}\PYG{o}{/}\PYG{n}{atom}
\PYG{k}{timestep}\PYG{+w}{      }\PYG{l+m}{1.0e\PYGZhy{}7}
\end{sphinxVerbatim}

\sphinxAtStartPar
Some notes on this input example:
\begin{itemize}
\item {} 
\sphinxAtStartPar
peridynamics simulations typically use SI {\hyperref[\detokenize{units::doc}]{\sphinxcrossref{\DUrole{doc}{units}}}}

\item {} 
\sphinxAtStartPar
particles must be created on a {\hyperref[\detokenize{lattice::doc}]{\sphinxcrossref{\DUrole{doc}{simple cubic lattice}}}}

\item {} 
\sphinxAtStartPar
using the {\hyperref[\detokenize{atom_style::doc}]{\sphinxcrossref{\DUrole{doc}{atom style peri}}}} is required

\item {} 
\sphinxAtStartPar
an {\hyperref[\detokenize{atom_modify::doc}]{\sphinxcrossref{\DUrole{doc}{atom map}}}} is required for indexing particles

\item {} 
\sphinxAtStartPar
The {\hyperref[\detokenize{neighbor::doc}]{\sphinxcrossref{\DUrole{doc}{skin distance}}}} used when computing neighbor lists
should be defined appropriately for your choice of simulation
parameters. The \sphinxstyleemphasis{skin} should be set to a value such that the
peridynamic horizon plus the skin distance is larger than the maximum
possible distance between two bonded particles (before their bond
breaks). Here it is set to 0.001 meters.

\item {} 
\sphinxAtStartPar
a {\hyperref[\detokenize{pair_peri::doc}]{\sphinxcrossref{\DUrole{doc}{peridynamics pair style}}}} is required.  Available
choices are currently: \sphinxstyleemphasis{peri/eps}, \sphinxstyleemphasis{peri/lps}, \sphinxstyleemphasis{peri/pmb}, and
\sphinxstyleemphasis{peri/ves}.  The model parameters are set with a {\hyperref[\detokenize{pair_coeff::doc}]{\sphinxcrossref{\DUrole{doc}{pair\_coeff}}}} command.

\item {} 
\sphinxAtStartPar
the mass density and volume fraction for each particle must be
defined.  This is done with the two {\hyperref[\detokenize{set::doc}]{\sphinxcrossref{\DUrole{doc}{set}}}} commands for
\sphinxstyleemphasis{density} and \sphinxstyleemphasis{volume}.  For a simple cubic lattice, the volume of a
particle should be equal to the cube of the lattice constant, here
\(V_i = \Delta x ^3\).

\item {} 
\sphinxAtStartPar
with the {\hyperref[\detokenize{velocity::doc}]{\sphinxcrossref{\DUrole{doc}{velocity}}}} command all particles are initially at rest

\item {} 
\sphinxAtStartPar
a plain {\hyperref[\detokenize{fix_nve::doc}]{\sphinxcrossref{\DUrole{doc}{velocity\sphinxhyphen{}Verlet time integrator}}}} is used,
which is algebraically equivalent to a centered difference in time,
but numerically more stable

\item {} 
\sphinxAtStartPar
you can compute the damage at the location of each particle with
{\hyperref[\detokenize{compute_damage_atom::doc}]{\sphinxcrossref{\DUrole{doc}{compute damage/atom}}}}

\item {} 
\sphinxAtStartPar
finally, the timestep is set to 0.1 microseconds with the
{\hyperref[\detokenize{timestep::doc}]{\sphinxcrossref{\DUrole{doc}{timestep}}}} command.

\end{itemize}


\subsubsection{Peridynamic Model of a Continuum}
\label{\detokenize{Howto_peri:peridynamic-model-of-a-continuum}}
\sphinxAtStartPar
The following is not a complete overview of peridynamics, but a
discussion of only those details specific to the model we have
implemented within LAMMPS. For more on the peridynamic theory, the
reader is referred to {\hyperref[\detokenize{Howto_peri:silling2007-2}]{\sphinxcrossref{\DUrole{std,std-ref}{(Silling 2007)}}}}. To begin,
we define the notation we will use.


\paragraph{Basic Notation}
\label{\detokenize{Howto_peri:basic-notation}}
\sphinxAtStartPar
Within the peridynamic literature, the following notational conventions
are generally used. The position of a given point in the reference
configuration is \(\textbf{x}\). Let
\(\mathbf{u}(\mathbf{x},t)\) and \(\mathbf{y}(\mathbf{x},t)\)
denote the displacement and position, respectively, of the point
\(\mathbf{x}\) at time \(t\). Define the relative position and
displacement vectors of two bonded points \(\textbf{x}\) and
\(\textbf{x}^\prime\) as \(\mathbf{\xi} = \textbf{x}^\prime -
\textbf{x}\) and \(\mathbf{\eta} = \textbf{u}(\textbf{x}^\prime,t) -
\textbf{u}(\textbf{x},t)\), respectively. We note here that
\(\mathbf{\eta}\) is time\sphinxhyphen{}dependent, and that \(\mathbf{\xi}\)
is not. It follows that the relative position of the two bonded points
in the current configuration can be written as \(\boldsymbol{\xi} +
\boldsymbol{\eta} =
\mathbf{y}(\mathbf{x}^{\prime},t)-\mathbf{y}(\mathbf{x},t)\).

\sphinxAtStartPar
Peridynamic models are frequently written using \sphinxstyleemphasis{states}, which we
briefly describe here. For the purposes of our discussion, all states
are operators that act on vectors in \(\mathbb{R}^3\). For a more
complete discussion of states, see {\hyperref[\detokenize{Howto_peri:silling2007-2}]{\sphinxcrossref{\DUrole{std,std-ref}{(Silling 2007)}}}}. A \sphinxstyleemphasis{vector state} is an operator whose image is a
vector, and may be viewed as a generalization of a second\sphinxhyphen{}rank
tensor. Similarly, a \sphinxstyleemphasis{scalar state} is an operator whose image is a
scalar. Of particular interest is the vector force state
\(\underline{\mathbf{T}}\left[ \mathbf{x},t \right]\left<
\mathbf{x}^{\prime}-\mathbf{x} \right>\), which is a mapping, having
units of force per volume squared, of the vector
\(\mathbf{x}^{\prime}-\mathbf{x}\) to the force vector state
field. The vector state operator \(\underline{\mathbf{T}}\) may
itself be a function of \(\mathbf{x}\) and \(t\). The
constitutive model is completely contained within
\(\underline{\mathbf{T}}\).

\sphinxAtStartPar
In the peridynamic theory, the deformation at a point depends
collectively on all points interacting with that point. Using the
notation of {\hyperref[\detokenize{Howto_peri:silling2007-2}]{\sphinxcrossref{\DUrole{std,std-ref}{(Silling 2007)}}}}, we write the
peridynamic equation of motion as

\phantomsection\label{\detokenize{Howto_peri:perinewtonii}}\begin{equation*}
\begin{split}\rho(\mathbf{x}) \ddot{\mathbf{u}}(\mathbf{x},t) =
\int_{\mathcal{H}_{\mathbf{x}}} \left\{
\underline{\mathbf{T}}\left[\mathbf{x},t
\right]\left<\mathbf{x}^{\prime}-\mathbf{x} \right> -
\underline{\mathbf{T}}\left[\mathbf{x}^{\prime},t
\right]\left<\mathbf{x}-\mathbf{x}^{\prime} \right> \right\}
{d}V_{\mathbf{x}^\prime} + \mathbf{b}(\mathbf{x},t), \qquad\qquad\textrm{(1)}\end{split}
\end{equation*}
\sphinxAtStartPar
where \(\rho\) represents the mass density,
\(\underline{\mathbf{T}}\) the force vector state, and
\(\mathbf{b}\) an external body force density. A point
\(\mathbf{x}\) interacts with all the points
\(\mathbf{x}^{\prime}\) within the neighborhood
\(\mathcal{H}_{\mathbf{x}}\), assumed to be a spherical region of
radius \(\delta>0\) centered at \(\mathbf{x}\). \(\delta\)
is called the \sphinxstyleemphasis{horizon}, and is analogous to the cutoff radius used in
molecular dynamics. Conditions on \(\underline{\mathbf{T}}\) for
which {\hyperref[\detokenize{Howto_peri:perinewtonii}]{\sphinxcrossref{\DUrole{std,std-ref}{(1)}}}} satisfies the balance of linear and angular
momentum are given in {\hyperref[\detokenize{Howto_peri:silling2007-2}]{\sphinxcrossref{\DUrole{std,std-ref}{(Silling 2007)}}}}.

\sphinxAtStartPar
We consider only force vector states that can be written as
\begin{equation*}
\begin{split}\underline{\mathbf{T}} = \underline{t}\,\underline{\mathbf{M}},\end{split}
\end{equation*}
\sphinxAtStartPar
with \(\underline{t}\) a \sphinxstyleemphasis{scalar force state} and
\(\underline{\mathbf{M}}\) the \sphinxstyleemphasis{deformed direction vector
state}, defined by
\begin{equation*}
\begin{split}\underline{\mathbf{M}}\left< \boldsymbol{\xi} \right> =
\left\{ \begin{array}{cl}
\frac{\boldsymbol{\xi} + \boldsymbol{\eta}}{
\left\Vert \boldsymbol{\xi} + \boldsymbol{\eta} \right\Vert
} & \left\Vert \boldsymbol{\xi} + \boldsymbol{\eta} \right\Vert \neq 0 \\
\boldsymbol{0}  & \textrm{otherwise}
\end{array}
\right. . \qquad\qquad\textrm{(2)}\end{split}
\end{equation*}
\sphinxAtStartPar
Such force states correspond to so\sphinxhyphen{}called \sphinxstyleemphasis{ordinary} materials
{\hyperref[\detokenize{Howto_peri:silling2007-2}]{\sphinxcrossref{\DUrole{std,std-ref}{(Silling 2007)}}}}. These are the materials for which
the force between any two interacting points \(\textbf{x}\) and
\(\textbf{x}^\prime\) acts along the line between the points.


\paragraph{Linear Peridynamic Solid (LPS) Model}
\label{\detokenize{Howto_peri:linear-peridynamic-solid-lps-model}}
\sphinxAtStartPar
We summarize the linear peridynamic solid (LPS) material model. For more
on this model, the reader is referred to {\hyperref[\detokenize{Howto_peri:silling2007-2}]{\sphinxcrossref{\DUrole{std,std-ref}{(Silling 2007)}}}}.  This model is a nonlocal analogue to a classical
linear elastic isotropic material. The elastic properties of a a
classical linear elastic isotropic material are determined by (for
example) the bulk and shear moduli. For the LPS model, the elastic
properties are analogously determined by the bulk and shear moduli,
along with the horizon \(\delta\).

\sphinxAtStartPar
The LPS model has a force scalar state
\begin{equation*}
\begin{split}\underline{t} = \frac{3K\theta}{m}\underline{\omega}\,\underline{x} +
\alpha \underline{\omega}\,\underline{e}^{\rm d}, \qquad\qquad\textrm{(3)}\end{split}
\end{equation*}
\sphinxAtStartPar
with \(K\) the bulk modulus and \(\alpha\) related to the shear
modulus \(G\) as
\begin{equation*}
\begin{split}\alpha = \frac{15 G}{m}.\end{split}
\end{equation*}
\sphinxAtStartPar
The remaining components of the model are described as follows. Define
the reference position scalar state \(\underline{x}\) so that
\(\underline{x}\left<\boldsymbol{\xi} \right> = \left\Vert
\boldsymbol{\xi} \right\Vert\). Then, the weighted volume \(m\) is
defined as
\begin{equation*}
\begin{split}m\left[ \mathbf{x} \right] = \int_{\mathcal{H}_\mathbf{x}}
\underline{\omega} \left<\boldsymbol{\xi}\right>
\underline{x}\left<\boldsymbol{\xi} \right>
\underline{x}\left<\boldsymbol{\xi} \right>{d}V_{\boldsymbol{\xi} }.  \qquad\qquad\textrm{(4)}\end{split}
\end{equation*}
\sphinxAtStartPar
Let
\begin{equation*}
\begin{split}\underline{e}\left[ \mathbf{x},t \right] \left<\boldsymbol{\xi}
\right> = \left\Vert \boldsymbol{\xi} + \boldsymbol{\eta}
\right\Vert - \left\Vert \boldsymbol{\xi} \right\Vert\end{split}
\end{equation*}
\sphinxAtStartPar
be the extension scalar state, and
\begin{equation*}
\begin{split}\theta\left[ \mathbf{x}, t \right] = \frac{3}{m\left[ \mathbf{x}
\right]}\int_{\mathcal{H}_\mathbf{x}} \underline{\omega}
\left<\boldsymbol{\xi}\right> \underline{x}\left<\boldsymbol{\xi}
\right> \underline{e}\left[ \mathbf{x},t
\right]\left<\boldsymbol{\xi} \right>{d}V_{\boldsymbol{\xi}}\end{split}
\end{equation*}
\sphinxAtStartPar
be the dilatation. The isotropic and deviatoric parts of the extension
scalar state are defined, respectively, as
\begin{equation*}
\begin{split}\underline{e}^{\rm i}=\frac{\theta \underline{x}}{3}, \qquad
\underline{e}^{\rm d} = \underline{e}- \underline{e}^{\rm i},\end{split}
\end{equation*}
\sphinxAtStartPar
where the arguments of the state functions and the vectors on which they
operate are omitted for simplicity. We note that the LPS model is linear
in the dilatation \(\theta\), and in the deviatoric part of the
extension \(\underline{e}^{\rm d}\).

\begin{sphinxadmonition}{note}{Note:}
\sphinxAtStartPar
The weighted volume \(m\) is time\sphinxhyphen{}independent, and does
not change as bonds break. It is computed with respect to the bond
family defined at the reference (initial) configuration.
\end{sphinxadmonition}

\sphinxAtStartPar
The non\sphinxhyphen{}negative scalar state \(\underline{\omega}\) is an
\sphinxstyleemphasis{influence function} {\hyperref[\detokenize{Howto_peri:silling2007-2}]{\sphinxcrossref{\DUrole{std,std-ref}{(Silling 2007)}}}}. For more on
influence functions, see {\hyperref[\detokenize{Howto_peri:seleson2010}]{\sphinxcrossref{\DUrole{std,std-ref}{(Seleson 2010)}}}}. If an
influence function \(\underline{\omega}\) depends only upon the
scalar \(\left\Vert \boldsymbol{\xi} \right\Vert\), (i.e.,
\(\underline{\omega}\left<\boldsymbol{\xi}\right> =
\underline{\omega}\left<\left\Vert \boldsymbol{\xi} \right\Vert\right>\)), then \(\underline{\omega}\) is a spherical influence function.
For a spherical influence function, the LPS model is isotropic
{\hyperref[\detokenize{Howto_peri:silling2007-2}]{\sphinxcrossref{\DUrole{std,std-ref}{(Silling 2007)}}}}.

\begin{sphinxadmonition}{note}{Note:}
\sphinxAtStartPar
In the LAMMPS implementation of the LPS model, the influence function
\(\underline{\omega}\left<\left\Vert \boldsymbol{\xi}
\right\Vert\right> = 1 / \left\Vert \boldsymbol{\xi} \right\Vert\) is
used. However, the user can define their own influence function by
altering the method “influence\_function” in the file
\sphinxcode{\sphinxupquote{pair\_peri\_lps.cpp}}. The LAMMPS peridynamics code permits both
spherical and non\sphinxhyphen{}spherical influence functions (e.g., isotropic and
non\sphinxhyphen{}isotropic materials).
\end{sphinxadmonition}


\paragraph{Prototype Microelastic Brittle (PMB) Model}
\label{\detokenize{Howto_peri:prototype-microelastic-brittle-pmb-model}}
\sphinxAtStartPar
We summarize the prototype microelastic brittle (PMB) material
model. For more on this model, the reader is referred to
{\hyperref[\detokenize{Howto_peri:silling2000-2}]{\sphinxcrossref{\DUrole{std,std-ref}{(Silling 2000)}}}} and {\hyperref[\detokenize{Howto_peri:silling2005}]{\sphinxcrossref{\DUrole{std,std-ref}{(Silling 2005)}}}}.  This model is a special case of the LPS model; see
{\hyperref[\detokenize{Howto_peri:seleson2010}]{\sphinxcrossref{\DUrole{std,std-ref}{(Seleson 2010)}}}} for the derivation. The elastic
properties of the PMB model are determined by the bulk modulus \(K\)
and the horizon \(\delta\).

\sphinxAtStartPar
The PMB model is expressed using the scalar force state field

\phantomsection\label{\detokenize{Howto_peri:peripmbstate}}\begin{equation*}
\begin{split}\underline{t}\left[ \mathbf{x},t \right]\left< \boldsymbol{\xi} \right> = \frac{1}{2} f\left( \boldsymbol{\eta} ,\boldsymbol{\xi} \right), \qquad\qquad\textrm{(5)}\end{split}
\end{equation*}
\sphinxAtStartPar
with \(f\) a scalar\sphinxhyphen{}valued function. We assume that \(f\) takes
the form
\begin{equation*}
\begin{split}f = c s,\end{split}
\end{equation*}
\sphinxAtStartPar
where

\phantomsection\label{\detokenize{Howto_peri:peric}}\begin{equation*}
\begin{split}c = \frac{18K}{\pi \delta^4}, \qquad\qquad\textrm{(6)}\end{split}
\end{equation*}
\sphinxAtStartPar
with \(K\) the bulk modulus and \(\delta\) the horizon, and
\(s\) the bond stretch, defined as
\begin{equation*}
\begin{split}s(t,\mathbf{\eta},\mathbf{\xi}) = \frac{ \left\Vert {\mathbf{\eta}+\mathbf{\xi}} \right\Vert - \left\Vert {\mathbf{\xi}} \right\Vert }{\left\Vert {\mathbf{\xi}} \right\Vert}.\end{split}
\end{equation*}
\sphinxAtStartPar
Bond stretch is a unitless quantity, and identical to a one\sphinxhyphen{}dimensional
definition of strain. As such, we see that a bond at its equilibrium
length has stretch \(s=0\), and a bond at twice its equilibrium
length has stretch \(s=1\).  The constant \(c\) given above is
appropriate for 3D models only. For more on the origins of the constant
\(c\), see {\hyperref[\detokenize{Howto_peri:silling2005}]{\sphinxcrossref{\DUrole{std,std-ref}{(Silling 2005)}}}}. For the derivation
of \(c\) for 1D and 2D models, see {\hyperref[\detokenize{Howto_peri:emmrich2007}]{\sphinxcrossref{\DUrole{std,std-ref}{(Emmrich)}}}}.

\sphinxAtStartPar
Given {\hyperref[\detokenize{Howto_peri:peripmbstate}]{\sphinxcrossref{\DUrole{std,std-ref}{(5)}}}}, {\hyperref[\detokenize{Howto_peri:perinewtonii}]{\sphinxcrossref{\DUrole{std,std-ref}{(1)}}}} reduces to
\begin{equation*}
\begin{split}\rho(\mathbf{x}) \ddot{\mathbf{u}}(\mathbf{x},t) = \int_{\mathcal{H}_\mathbf{x}} \mathbf{f} \left( \boldsymbol{\eta},\boldsymbol{\xi} \right){d}V_{\boldsymbol{\xi}} + \mathbf{b}(\mathbf{x},t), \qquad\qquad\textrm{(7)}\end{split}
\end{equation*}
\sphinxAtStartPar
with
\begin{equation*}
\begin{split}\mathbf{f} \left( \boldsymbol{\eta},  \boldsymbol{\xi}\right) =f \left( \boldsymbol{\eta},  \boldsymbol{\xi}\right)  \frac{\boldsymbol{\xi}+ \boldsymbol{\eta}}{ \left\Vert {\boldsymbol{\xi} + \boldsymbol{\eta}} \right\Vert}.\end{split}
\end{equation*}
\sphinxAtStartPar
Unlike the LPS model, the PMB model has a Poisson ratio of
\(\nu=1/4\) in 3D, and \(\nu=1/3\) in 2D. This is reflected in
the input for the PMB model, which requires only the bulk modulus of the
material, whereas the LPS model requires both the bulk and shear moduli.


\paragraph{Damage}
\label{\detokenize{Howto_peri:damage}}\label{\detokenize{Howto_peri:peridamage}}
\sphinxAtStartPar
Bonds are made to break when they are stretched beyond a given
limit. Once a bond fails, it is failed forever {\hyperref[\detokenize{Howto_peri:silling2005}]{\sphinxcrossref{\DUrole{std,std-ref}{(Silling)}}}}. Further, new bonds are never created during the course
of a simulation. We discuss only one criterion for bond breaking, called
the \sphinxstyleemphasis{critical stretch} criterion.

\sphinxAtStartPar
Define \(\mu\) to be the history\sphinxhyphen{}dependent scalar
boolean function

\phantomsection\label{\detokenize{Howto_peri:perimu}}\begin{equation*}
\begin{split}\mu(t,\mathbf{\eta},\mathbf{\xi}) = \left\{
\begin{array}{cl}
1 & \mbox{if $s(t^\prime,\mathbf{\eta},\mathbf{\xi})< \min \left(s_0(t^\prime,\mathbf{\eta},\mathbf{\xi}) , s_0(t^\prime,\mathbf{\eta}^\prime,\mathbf{\xi}^\prime) \right)$ for all $0 \leq t^\prime \leq t$} \\
0 & \mbox{otherwise}
\end{array}\right\}.  \qquad\qquad\textrm{(8)}\end{split}
\end{equation*}
\sphinxAtStartPar
where \(\mathbf{\eta}^\prime = \textbf{u}(\textbf{x}^{\prime
\prime},t) - \textbf{u}(\textbf{x}^\prime,t)\) and
\(\mathbf{\xi}^\prime = \textbf{x}^{\prime \prime} -
\textbf{x}^\prime\). Here, \(s_0(t,\mathbf{\eta},\mathbf{\xi})\) is a
critical stretch defined as

\phantomsection\label{\detokenize{Howto_peri:peris0}}\begin{equation*}
\begin{split}s_0(t,\mathbf{\eta},\mathbf{\xi}) = s_{00} - \alpha s_{\min}(t,\mathbf{\eta},\mathbf{\xi}), \qquad s_{\min}(t) = \min_{\mathbf{\xi}} s(t,\mathbf{\eta},\mathbf{\xi}), \qquad\qquad\textrm{(9)}\end{split}
\end{equation*}
\sphinxAtStartPar
where \(s_{00}\) and \(\alpha\) are material\sphinxhyphen{}dependent
constants. The history function \(\mu\) breaks bonds when the
stretch \(s\) exceeds the critical stretch \(s_0\).

\sphinxAtStartPar
Although \(s_0(t,\mathbf{\eta},\mathbf{\xi})\) is expressed as a
property of a particle, bond breaking must be a symmetric operation for
all particle pairs sharing a bond. That is, particles \(\textbf{x}\)
and \(\textbf{x}^\prime\) must utilize the same test when deciding
to break their common bond. This can be done by any method that treats
the particles symmetrically. In the definition of \(\mu\) above, we
have chosen to take the minimum of the two \(s_0\) values for
particles \(\textbf{x}\) and \(\textbf{x}^\prime\) when
determining if the \(\textbf{x}\)\textendash{}\(\textbf{x}^\prime\) bond
should be broken.

\sphinxAtStartPar
Following {\hyperref[\detokenize{Howto_peri:silling2005}]{\sphinxcrossref{\DUrole{std,std-ref}{(Silling)}}}}, we can define the damage at a
point \(\textbf{x}\) as

\phantomsection\label{\detokenize{Howto_peri:peridamageeq}}\begin{equation*}
\begin{split}\varphi(\textbf{x}, t) = 1 - \frac{\int_{\mathcal{H}_{\textbf{x}}} \mu(t,\mathbf{\eta},\mathbf{\xi}) dV_{\textbf{x}^\prime}
}{ \int_{\mathcal{H}_{\textbf{x}}} dV_{\textbf{x}^\prime} }.  \qquad\qquad\textrm{(10)}\end{split}
\end{equation*}

\subsubsection{Discrete Peridynamic Model and LAMMPS Implementation}
\label{\detokenize{Howto_peri:discrete-peridynamic-model-and-lammps-implementation}}
\sphinxAtStartPar
In LAMMPS, instead of {\hyperref[\detokenize{Howto_peri:perinewtonii}]{\sphinxcrossref{\DUrole{std,std-ref}{(1)}}}}, we model this equation of
motion:
\begin{equation*}
\begin{split}\rho(\mathbf{x}) \ddot{\textbf{y}}(\mathbf{x},t) = \int_{\mathcal{H}_{\mathbf{x}}}
\left\{ \underline{\mathbf{T}}\left[ \mathbf{x},t \right]\left<\mathbf{x}^{\prime}-\mathbf{x} \right>
- \underline{\mathbf{T}}\left[\mathbf{x}^{\prime},t \right]\left<\mathbf{x}-\mathbf{x}^{\prime} \right> \right\}
  {d}V_{\mathbf{x}^\prime} + \mathbf{b}(\mathbf{x},t),\end{split}
\end{equation*}
\sphinxAtStartPar
where we explicitly track and store at each timestep the positions and
not the displacements of the particles. We observe that
\(\ddot{\textbf{y}}(\textbf{x}, t) = \ddot{\textbf{x}} +
\ddot{\textbf{u}}(\textbf{x}, t) = \ddot{\textbf{u}}(\textbf{x}, t)\), so
that this is equivalent to {\hyperref[\detokenize{Howto_peri:perinewtonii}]{\sphinxcrossref{\DUrole{std,std-ref}{(1)}}}}.


\paragraph{Spatial Discretization}
\label{\detokenize{Howto_peri:spatial-discretization}}
\sphinxAtStartPar
The region defining a peridynamic material is discretized into particles
forming a simple cubic lattice with lattice constant \(\Delta x\),
where each particle \(i\) is associated with some volume fraction
\(V_i\). For any particle \(i\), let \(\mathcal{F}_i\)
denote the family of particles for which particle \(i\) shares a
bond in the reference configuration.  That is,

\phantomsection\label{\detokenize{Howto_peri:peribondfamily}}\begin{equation*}
\begin{split}\mathcal{F}_i = \{ p ~ | ~ \left\Vert {\textbf{x}_p - \textbf{x}_i} \right\Vert \leq \delta \}. \qquad\qquad\textrm{(11)}\end{split}
\end{equation*}
\sphinxAtStartPar
The discretized equation of motion replaces {\hyperref[\detokenize{Howto_peri:perinewtonii}]{\sphinxcrossref{\DUrole{std,std-ref}{(1)}}}} with

\phantomsection\label{\detokenize{Howto_peri:peridiscretenewtonii}}\begin{equation*}
\begin{split}\rho \ddot{\textbf{y}}_i^n =
\sum_{p \in \mathcal{F}_i}
\left\{ \underline{\mathbf{T}}\left[ \textbf{x}_i,t \right]\left<\textbf{x}_p^{\prime}-\textbf{x}_i \right>
- \underline{\mathbf{T}}\left[\textbf{x}_p,t \right]\left<\textbf{x}_i-\textbf{x}_p \right> \right\}
  V_{p} + \textbf{b}_i^n, \qquad\qquad\textrm{(12)}\end{split}
\end{equation*}
\sphinxAtStartPar
where \(n\) is the timestep number and subscripts denote the particle number.


\paragraph{Short\sphinxhyphen{}Range Forces}
\label{\detokenize{Howto_peri:short-range-forces}}
\sphinxAtStartPar
In the model discussed so far, particles interact only through their
bond forces. A particle with no bonds becomes a free non\sphinxhyphen{}interacting
particle. To account for contact forces, short\sphinxhyphen{}range forces are
introduced {\hyperref[\detokenize{Howto_peri:silling2007-2}]{\sphinxcrossref{\DUrole{std,std-ref}{(Silling 2007)}}}}. We add to the force in
{\hyperref[\detokenize{Howto_peri:peridiscretenewtonii}]{\sphinxcrossref{\DUrole{std,std-ref}{(12)}}}} the following force
\begin{equation*}
\begin{split}\textbf{f}_S(\textbf{y}_p,\textbf{y}_i) = \min \left\{ 0, \frac{c_S}{\delta}(\left\Vert {\textbf{y}_p-\textbf{y}_i} \right\Vert - d_{pi}) \right\}
\frac{\textbf{y}_p-\textbf{y}_i}{\left\Vert {\textbf{y}_p-\textbf{y}_i} \right\Vert}, \qquad\qquad\textrm{(13)}\end{split}
\end{equation*}
\sphinxAtStartPar
where \(d_{pi}\) is the short\sphinxhyphen{}range interaction distance between
particles \(p\) and \(i\), and \(c_S\) is a multiple of the
constant \(c\) from {\hyperref[\detokenize{Howto_peri:peric}]{\sphinxcrossref{\DUrole{std,std-ref}{(6)}}}}. Note that the short\sphinxhyphen{}range force
is always repulsive, never attractive. In practice, we choose

\phantomsection\label{\detokenize{Howto_peri:perics}}\begin{equation*}
\begin{split}c_S = 15 \frac{18K}{\pi \delta^4}. \qquad\qquad\textrm{(14)}\end{split}
\end{equation*}
\sphinxAtStartPar
For the short\sphinxhyphen{}range interaction distance, we choose {\hyperref[\detokenize{Howto_peri:silling2007-2}]{\sphinxcrossref{\DUrole{std,std-ref}{(Silling 2007)}}}}
\begin{equation*}
\begin{split}d_{pi} = \min \left\{ 0.9 \left\Vert {\textbf{x}_p - \textbf{x}_i} \right\Vert, 1.35 (r_p + r_i) \right\}, \qquad\qquad\textrm{(15)}\end{split}
\end{equation*}
\sphinxAtStartPar
where \(r_i\) is called the \sphinxstyleemphasis{node radius} of particle
\(i\). Given a discrete lattice, we choose \(r_i\) to be half
the lattice constant.

\begin{sphinxadmonition}{note}{Note:}
\sphinxAtStartPar
For a simple cubic lattice, \(\Delta x = \Delta y = \Delta z\).
\end{sphinxadmonition}

\sphinxAtStartPar
Given this definition of \(d_{pi}\), contact forces appear only when
particles are under compression.

\sphinxAtStartPar
When accounting for short\sphinxhyphen{}range forces, it is convenient to define the
short\sphinxhyphen{}range family of particles
\begin{equation*}
\begin{split}\mathcal{F}^S_i = \{ p ~ | ~ \left\Vert {\textbf{y}_p - \textbf{y}_i} \right\Vert \leq d_{pi} \}.\end{split}
\end{equation*}

\paragraph{Modification to the Particle Volume}
\label{\detokenize{Howto_peri:modification-to-the-particle-volume}}
\sphinxAtStartPar
The right\sphinxhyphen{}hand side of {\hyperref[\detokenize{Howto_peri:peridiscretenewtonii}]{\sphinxcrossref{\DUrole{std,std-ref}{(12)}}}} may be thought of as
a midpoint quadrature of {\hyperref[\detokenize{Howto_peri:perinewtonii}]{\sphinxcrossref{\DUrole{std,std-ref}{(1)}}}}. To slightly improve the
accuracy of this quadrature, we discuss a modification to the particle
volume used in {\hyperref[\detokenize{Howto_peri:peridiscretenewtonii}]{\sphinxcrossref{\DUrole{std,std-ref}{(12)}}}}. In a situation where two
particles share a bond with \(\left\Vert { \textbf{x}_p -
\textbf{x}_i }\right\Vert = \delta\), for example, we suppose that only
approximately half the volume of each particle is “seen” by the other
{\hyperref[\detokenize{pair_peri:silling2007}]{\sphinxcrossref{\DUrole{std,std-ref}{(Silling 2007)}}}}. When computing the force of each
particle on the other we use \(V_p / 2\) rather than \(V_p\) in
{\hyperref[\detokenize{Howto_peri:peridiscretenewtonii}]{\sphinxcrossref{\DUrole{std,std-ref}{(12)}}}}. As such, we introduce a nodal volume
scaling function for all bonded particles where \(\delta - r_i \leq
\left\Vert { \textbf{x}_p - \textbf{x}_i } \right\Vert \leq \delta\) (see
the Figure below).

\sphinxAtStartPar
We choose to use a linear unitless nodal volume scaling function
\begin{equation*}
\begin{split}\nu(\textbf{x}_p - \textbf{x}_i) = \left\{
\begin{array}{cl}
-\frac{1}{2 r_i} \left\Vert {\textbf{x}_p - \textbf{x}_i} \right\Vert + \left( \frac{\delta}{2 r_i} + \frac{1}{2} \right) & \mbox{if }
\delta - r_i \leq \left\Vert {\textbf{x}_p - \textbf{x}_i } \right\Vert \leq \delta \\
1 & \mbox{if } \left\Vert {\textbf{x}_p - \textbf{x}_i } \right\Vert \leq \delta - r_i \\
0 & \mbox{otherwise}
\end{array}
\right\}\end{split}
\end{equation*}
\sphinxAtStartPar
If \(\left\Vert {\textbf{x}_p - \textbf{x}_i} \right\Vert = \delta\), \(\nu = 0.5\), and if
\(\left\Vert {\textbf{x}_p - \textbf{x}_i} \right\Vert = \delta - r_i\), \(\nu = 1.0\), for
example.

\begin{figure}[htbp]
\centering
\capstart

\noindent\sphinxincludegraphics{{pdlammps_fig1}.png}
\caption{Diagram showing horizon of a particular particle, demonstrating that
the volume associated with particles near the boundary of the horizon is
not completely contained within the horizon.}\label{\detokenize{Howto_peri:id2}}\end{figure}


\paragraph{Temporal Discretization}
\label{\detokenize{Howto_peri:temporal-discretization}}
\sphinxAtStartPar
When discretizing time in LAMMPS, we use a velocity\sphinxhyphen{}Verlet scheme, where
both the position and velocity of the particle are stored
explicitly. The velocity\sphinxhyphen{}Verlet scheme is generally expressed in three
steps. In {\hyperref[\detokenize{Howto_peri:algvelverlet}]{\sphinxcrossref{\DUrole{std,std-ref}{Algorithm 1}}}}, \(\rho_i\) denotes the
mass density of a particle and \(\widetilde{\textbf{f}}_i^n\)
denotes the the net force density on particle \(i\) at timestep
\(n\). The LAMMPS command {\hyperref[\detokenize{fix_nve::doc}]{\sphinxcrossref{\DUrole{doc}{fix nve}}}} performs a
velocity\sphinxhyphen{}Verlet integration.
\begin{quote}
\phantomsection\label{\detokenize{Howto_peri:algvelverlet}}
\begin{sphinxadmonition}{note}{Algorithm 1: Velocity Verlet}

\begin{DUlineblock}{0em}
\item[] 1: \(\textbf{v}_i^{n + 1/2} = \textbf{v}_i^n + \frac{\Delta t}{2 \rho_i} \widetilde{\textbf{f}}_i^n\)
\item[] 2: \(\textbf{y}_i^{n+1}     = \textbf{y}_i^n + \Delta t \textbf{v}_i^{n + 1/2}\)
\item[] 3: \(\textbf{v}_i^{n+1} = \textbf{v}_i^{n+1/2} + \frac{\Delta t}{2 \rho_i} \widetilde{\textbf{f}}_i^{n+1}\)
\end{DUlineblock}
\end{sphinxadmonition}
\end{quote}


\paragraph{Breaking Bonds}
\label{\detokenize{Howto_peri:breaking-bonds}}
\sphinxAtStartPar
During the course of simulation, it may be necessary to break bonds, as
described in the {\hyperref[\detokenize{Howto_peri:peridamage}]{\sphinxcrossref{\DUrole{std,std-ref}{Damage section}}}}. Bonds are recorded
as broken in a simulation by removing them from the bond family
\(\mathcal{F}_i\) (see {\hyperref[\detokenize{Howto_peri:peribondfamily}]{\sphinxcrossref{\DUrole{std,std-ref}{(11)}}}}).

\sphinxAtStartPar
A naive implementation would have us first loop over all bonds and
compute \(s_{min}\) in {\hyperref[\detokenize{Howto_peri:peris0}]{\sphinxcrossref{\DUrole{std,std-ref}{(9)}}}}, then loop over all bonds
again and break bonds with a stretch \(s > s0\) as in
{\hyperref[\detokenize{Howto_peri:perimu}]{\sphinxcrossref{\DUrole{std,std-ref}{(8)}}}}, and finally loop over all particles and compute forces
for the next step of {\hyperref[\detokenize{Howto_peri:algvelverlet}]{\sphinxcrossref{\DUrole{std,std-ref}{Algorithm 1}}}}. For reasons of
computational efficiency, we will utilize the values of \(s_0\) from
the \sphinxstyleemphasis{previous} timestep when deciding to break a bond.

\begin{sphinxadmonition}{note}{Note:}
\sphinxAtStartPar
For the first timestep, \(s_0\) is initialized to
\(\mathbf{\infty}\) for all nodes. This means that no bonds may
be broken until the second timestep. As such, it is recommended that
the first few timesteps of the peridynamic simulation not involve any
actions that might result in the breaking of bonds. As a practical
example, the projectile in the {\hyperref[\detokenize{Howto_peri:periexample}]{\sphinxcrossref{\DUrole{std,std-ref}{commented example below}}}} is placed such that it does not impact the target
brittle plate until several timesteps into the simulation.
\end{sphinxadmonition}


\paragraph{LPS Pseudocode}
\label{\detokenize{Howto_peri:lps-pseudocode}}
\sphinxAtStartPar
A sketch of the LPS model implementation in the PERI package appears in
{\hyperref[\detokenize{Howto_peri:algperilps}]{\sphinxcrossref{\DUrole{std,std-ref}{Algorithm 2}}}}. This algorithm makes use of the
routines in {\hyperref[\detokenize{Howto_peri:algperilpsm}]{\sphinxcrossref{\DUrole{std,std-ref}{Algorithm 3}}}} and {\hyperref[\detokenize{Howto_peri:algperilpstheta}]{\sphinxcrossref{\DUrole{std,std-ref}{Algorithm 4}}}}.
\begin{quote}
\phantomsection\label{\detokenize{Howto_peri:algperilps}}
\begin{sphinxadmonition}{note}{Algorithm 2: LPS Peridynamic Model Pseudocode}

\begin{DUlineblock}{0em}
\item[] Fix \(s_{00}\), \(\alpha\), horizon \(\delta\), bulk modulus \(K\), shear modulus \(G\), timestep \(\Delta t\), and generate initial lattice of particles with lattice constant \(\Delta x\). Let there be \(N\) particles. Define constant \(c_S\) for repulsive short\sphinxhyphen{}range forces.
\item[] Initialize bonds between all particles \(i \neq j\) where \(\left\Vert {\textbf{x}_j - \textbf{x}_i} \right\Vert \leq \delta\)
\item[] Initialize weighted volume \(m\) for all particles using {\hyperref[\detokenize{Howto_peri:algperilpsm}]{\sphinxcrossref{\DUrole{std,std-ref}{Algorithm 3}}}}
\item[] Initialize \(s_0 = \mathbf{\infty}\) \{\sphinxstyleemphasis{Initialize each entry to MAX\_DOUBLE}\}
\item[] \sphinxstylestrong{while} not done \sphinxstylestrong{do}
\item[]
\begin{DUlineblock}{\DUlineblockindent}
\item[] Perform step 1 of {\hyperref[\detokenize{Howto_peri:algvelverlet}]{\sphinxcrossref{\DUrole{std,std-ref}{Algorithm 1}}}}, updating velocities of all particles
\item[] Perform step 2 of {\hyperref[\detokenize{Howto_peri:algvelverlet}]{\sphinxcrossref{\DUrole{std,std-ref}{Algorithm 1}}}}, updating positions of all particles
\item[] \(\tilde{s}_0 = \mathbf{\infty}\) \{\sphinxstyleemphasis{Initialize each entry to MAX\_DOUBLE}\}
\item[] \sphinxstylestrong{for} \(i=1\) to \(N\) \sphinxstylestrong{do}
\item[]
\begin{DUlineblock}{\DUlineblockindent}
\item[] \{Compute short\sphinxhyphen{}range forces\}
\item[] \sphinxstylestrong{for all} particles \(j \in \mathcal{F}^S_i\) (the short\sphinxhyphen{}range family of nodes for particle \(i\)) \sphinxstylestrong{do}
\item[]
\begin{DUlineblock}{\DUlineblockindent}
\item[] \(r = \left\Vert {\textbf{y}_j - \textbf{y}_i} \right\Vert\)
\item[] \(dr = \min \{ 0, r - d \}\) \{\sphinxstyleemphasis{Short\sphinxhyphen{}range forces are only repulsive, never attractive}\}
\item[] \(k = \frac{c_S}{\delta} V_k dr\) \{\(c_S\) \sphinxstyleemphasis{defined in :ref:\textasciigrave{}(14) \textless{}pericS\textgreater{}\textasciigrave{}}\}
\item[] \(\textbf{f} = \textbf{f} + k \frac{\textbf{y}_j-\textbf{y}_i}{\left\Vert {\textbf{y}_j-\textbf{y}_i} \right\Vert}\)
\end{DUlineblock}
\item[] \sphinxstylestrong{end for}
\end{DUlineblock}
\item[] \sphinxstylestrong{end for}
\item[] Compute the dilatation for each particle using {\hyperref[\detokenize{Howto_peri:algperilpstheta}]{\sphinxcrossref{\DUrole{std,std-ref}{Algorithm 4}}}}
\item[] \sphinxstylestrong{for} \(i=1\) to \(N\) \sphinxstylestrong{do}
\item[]
\begin{DUlineblock}{\DUlineblockindent}
\item[] \{Compute bond forces\}
\item[] \sphinxstylestrong{for all} particles \(j\) sharing an unbroken bond with particle \(i\) \sphinxstylestrong{do}
\item[]
\begin{DUlineblock}{\DUlineblockindent}
\item[] \(e = \left\Vert {\textbf{y}_j - \textbf{y}_i} \right\Vert - \left\Vert {\textbf{x}_j - \textbf{x}_i} \right\Vert\)
\item[] \(\omega_+ = \underline{\omega}\left<\textbf{x}_j - \textbf{x}_i\right>\) \{\sphinxstyleemphasis{Influence function evaluation}\}
\item[] \(\omega_- = \underline{\omega}\left<\textbf{x}_i - \textbf{x}_j\right>\) \{\sphinxstyleemphasis{Influence function evaluation}\}
\item[] \(\hat{f} = \left[ (3K-5G)\left( \frac{\theta(i)}{m(i)}\omega_+ + \frac{\theta(j)}{m(j)}\omega_- \right) \left\Vert {\textbf{x}_j - \textbf{x}_i} \right\Vert + 15G \left( \frac{\omega_+}{m(i)} + \frac{\omega_-}{m(j)} \right) e \right] \nu(\textbf{x}_j - \textbf{x}_i) V_j\)
\item[] \(\textbf{f} = \textbf{f} + \hat{f} \frac{\textbf{y}_j-\textbf{y}_i}{\left\Vert {\textbf{y}_j-\textbf{y}_i} \right\Vert}\)
\item[] \sphinxstylestrong{if} \((dr / \left\Vert {\textbf{x}_j - \textbf{x}_i} \right\Vert) > \min(s_0(i), s_0(j))\) \sphinxstylestrong{then}
\item[]
\begin{DUlineblock}{\DUlineblockindent}
\item[] Break \(i\)’s bond with \(j\) \{\(j\) \sphinxstyleemphasis{‘s bond with} \(i\) \sphinxstyleemphasis{will be broken when this loop iterates on} \(j\)\}
\end{DUlineblock}
\item[] \sphinxstylestrong{end if}
\item[] \(\tilde{s}_0(i) = \min (\tilde{s}_0(i),s_{00}-\alpha(dr / \left\Vert {\textbf{x}_j - \textbf{x}_i} \right\Vert))\)
\end{DUlineblock}
\item[] \sphinxstylestrong{end for}
\end{DUlineblock}
\item[] \sphinxstylestrong{end for}
\item[] \(s_0 = \tilde{s}_0\) \{\sphinxstyleemphasis{Store for use in next timestep}\}
\item[] Perform step 3 of {\hyperref[\detokenize{Howto_peri:algvelverlet}]{\sphinxcrossref{\DUrole{std,std-ref}{Algorithm 1}}}}, updating velocities of all particles
\end{DUlineblock}
\item[] \sphinxstylestrong{end while}
\end{DUlineblock}
\end{sphinxadmonition}
\phantomsection\label{\detokenize{Howto_peri:algperilpsm}}
\begin{sphinxadmonition}{note}{Algorithm 3: Computation of Weighted Volume \sphinxstyleemphasis{m}}

\begin{DUlineblock}{0em}
\item[] \sphinxstylestrong{for} \(i=1\) to \(N\) \sphinxstylestrong{do}
\item[]
\begin{DUlineblock}{\DUlineblockindent}
\item[] \(m(i) = 0.0\)
\item[] \sphinxstylestrong{for all} particles \(j\) sharing a bond with particle \(i\) \sphinxstylestrong{do}
\item[]
\begin{DUlineblock}{\DUlineblockindent}
\item[] \(m(i) = m(i) + \underline{\omega}\left<\textbf{x}_j - \textbf{x}_i\right> \left\Vert {\textbf{x}_j - \textbf{x}_i} \right\Vert^2 \nu(\textbf{x}_j - \textbf{x}_i) V_j\)
\end{DUlineblock}
\item[] \sphinxstylestrong{end for}
\end{DUlineblock}
\item[] \sphinxstylestrong{end for}
\end{DUlineblock}
\end{sphinxadmonition}
\phantomsection\label{\detokenize{Howto_peri:algperilpstheta}}
\begin{sphinxadmonition}{note}{Algorithm 4: Computation of Dilatation \protect\(\theta\protect\)}

\begin{DUlineblock}{0em}
\item[] \sphinxstylestrong{for} \(i=1\) to \(N\) \sphinxstylestrong{do}
\item[]
\begin{DUlineblock}{\DUlineblockindent}
\item[] \(\theta(i) = 0.0\)
\item[] \sphinxstylestrong{for all} particles \(j\) sharing an unbroken bond with particle \(i\) \sphinxstylestrong{do}
\item[]
\begin{DUlineblock}{\DUlineblockindent}
\item[] \(e = \left\Vert {\textbf{y}_i - \textbf{y}_j} \right\Vert - \left\Vert {\textbf{x}_i - \textbf{x}_j} \right\Vert\)
\item[] \(\theta(i) = \theta(i) + \underline{\omega}\left<\textbf{x}_j - \textbf{x}_i\right> \left\Vert {\textbf{x}_j - \textbf{x}_i} \right\Vert e  \nu(\textbf{x}_j - \textbf{x}_i) V_j\)
\end{DUlineblock}
\item[] \sphinxstylestrong{end for}
\item[] \(\theta(i) = \frac{3}{m(i)}\theta(i)\)
\end{DUlineblock}
\item[] \sphinxstylestrong{end for}
\end{DUlineblock}
\end{sphinxadmonition}
\end{quote}


\paragraph{PMB Pseudocode}
\label{\detokenize{Howto_peri:pmb-pseudocode}}
\sphinxAtStartPar
A sketch of the PMB model implementation in the PERI package appears in
{\hyperref[\detokenize{Howto_peri:algperipmb}]{\sphinxcrossref{\DUrole{std,std-ref}{Algorithm 5}}}}.
\begin{quote}
\phantomsection\label{\detokenize{Howto_peri:algperipmb}}
\begin{sphinxadmonition}{note}{Algorithm 5: PMB Peridynamic Model Pseudocode}

\begin{DUlineblock}{0em}
\item[] Fix \(s_{00}\), \(\alpha\), horizon \(\delta\), spring constant \(c\), timestep \(\Delta t\), and generate initial lattice of particles with lattice constant \(\Delta x\). Let there be \(N\) particles.
\item[] Initialize bonds between all particles \(i \neq j\) where \(\left\Vert {\textbf{x}_j - \textbf{x}_i} \right\Vert \leq \delta\)
\item[] Initialize \(s_0 = \mathbf{\infty}\) \{\sphinxstyleemphasis{Initialize each entry to MAX\_DOUBLE}\}
\item[] \sphinxstylestrong{while} not done \sphinxstylestrong{do}
\item[]
\begin{DUlineblock}{\DUlineblockindent}
\item[] Perform step 1 of {\hyperref[\detokenize{Howto_peri:algvelverlet}]{\sphinxcrossref{\DUrole{std,std-ref}{Algorithm 1}}}}, updating velocities of all particles
\item[] Perform step 2 of {\hyperref[\detokenize{Howto_peri:algvelverlet}]{\sphinxcrossref{\DUrole{std,std-ref}{Algorithm 1}}}}, updating positions of all particles
\item[] \(\tilde{s}_0 = \mathbf{\infty}\) \{\sphinxstyleemphasis{Initialize each entry to MAX\_DOUBLE}\}
\item[] \sphinxstylestrong{for} \(i=1\) to \(N\) \sphinxstylestrong{do}
\item[]
\begin{DUlineblock}{\DUlineblockindent}
\item[] \{Compute short\sphinxhyphen{}range forces\}
\item[] \sphinxstylestrong{for all} particles \(j \in \mathcal{F}^S_i\) (the short\sphinxhyphen{}range family of nodes for particle \(i\)) \sphinxstylestrong{do}
\item[]
\begin{DUlineblock}{\DUlineblockindent}
\item[] \(r = \left\Vert {\textbf{y}_j - \textbf{y}_i} \right\Vert\)
\item[] \(dr = \min \{ 0, r - d \}\) \{\sphinxstyleemphasis{Short\sphinxhyphen{}range forces are only repulsive, never attractive}\}
\item[] \(k = \frac{c_S}{\delta} V_k dr\) \{\(c_S\) \sphinxstyleemphasis{defined in :ref:\textasciigrave{}(14) \textless{}pericS\textgreater{}\textasciigrave{}}\}
\item[] \(\textbf{f} = \textbf{f} + k \frac{\textbf{y}_j-\textbf{y}_i}{\left\Vert {\textbf{y}_j-\textbf{y}_i} \right\Vert}\)
\end{DUlineblock}
\item[] \sphinxstylestrong{end for}
\end{DUlineblock}
\item[] \sphinxstylestrong{end for}
\item[] \sphinxstylestrong{for} \(i=1\) to \(N\) \sphinxstylestrong{do}
\item[]
\begin{DUlineblock}{\DUlineblockindent}
\item[] \{Compute bond forces\}
\item[] \sphinxstylestrong{for all} particles \(j\) sharing an unbroken bond with particle \(i\) \sphinxstylestrong{do}
\item[]
\begin{DUlineblock}{\DUlineblockindent}
\item[] \(r = \left\Vert {\textbf{y}_j - \textbf{y}_i} \right\Vert\)
\item[] \(dr = r - \left\Vert {\textbf{x}_j - \textbf{x}_i} \right\Vert\)
\item[] \(k = \frac{c}{\left\Vert {\textbf{x}_i - \textbf{x}_j} \right\Vert} \nu(\textbf{x}_i - \textbf{x}_j) V_j dr\) \{\(c\) \sphinxstyleemphasis{defined in :ref:\textasciigrave{}(6) \textless{}peric\textgreater{}\textasciigrave{}}\}
\item[] \(\textbf{f} = \textbf{f} + k \frac{\textbf{y}_j-\textbf{y}_i}{\left\Vert {\textbf{y}_j-\textbf{y}_i} \right\Vert}\)
\item[] \sphinxstylestrong{if} \((dr / \left\Vert {\textbf{x}_j - \textbf{x}_i} \right\Vert) > \min(s_0(i), s_0(j))\) \sphinxstylestrong{then}
\item[]
\begin{DUlineblock}{\DUlineblockindent}
\item[] Break \(i\)’s bond with \(j\) \{\(j\)\sphinxstyleemphasis{‘s bond with} \(i\) \sphinxstyleemphasis{will be broken when this loop iterates on} \(j\)\}
\end{DUlineblock}
\item[] \sphinxstylestrong{end if}
\item[] \(\tilde{s}_0(i) = \min (\tilde{s}_0(i),s_{00}-\alpha(dr / \left\Vert {\textbf{x}_j - \textbf{x}_i} \right\Vert))\)
\end{DUlineblock}
\item[] \sphinxstylestrong{end for}
\end{DUlineblock}
\item[] \sphinxstylestrong{end for}
\item[] \(s_0 = \tilde{s}_0\) \{\sphinxstyleemphasis{Store for use in next timestep}\}
\item[] Perform step 3 of {\hyperref[\detokenize{Howto_peri:algvelverlet}]{\sphinxcrossref{\DUrole{std,std-ref}{Algorithm 1}}}}, updating velocities of all particles
\end{DUlineblock}
\item[] \sphinxstylestrong{end while}
\end{DUlineblock}
\end{sphinxadmonition}
\end{quote}


\paragraph{Damage}
\label{\detokenize{Howto_peri:id1}}
\sphinxAtStartPar
The damage associated with every particle (see {\hyperref[\detokenize{Howto_peri:peridamageeq}]{\sphinxcrossref{\DUrole{std,std-ref}{(10)}}}}) can
optionally be computed and output with a LAMMPS data dump. To do this,
your input script must contain the command {\hyperref[\detokenize{compute_damage_atom::doc}]{\sphinxcrossref{\DUrole{doc}{compute damage/atom}}}} This enables a LAMMPS per\sphinxhyphen{}atom compute to
calculate the damage associated with each particle every time a LAMMPS
{\hyperref[\detokenize{dump::doc}]{\sphinxcrossref{\DUrole{doc}{data dump}}}} frame is written.


\paragraph{Visualizing Simulation Results}
\label{\detokenize{Howto_peri:visualizing-simulation-results}}
\sphinxAtStartPar
There are multiple ways to visualize the simulation results. Typically,
you want to display the particles and color code them by the value
computed with the {\hyperref[\detokenize{compute_damage_atom::doc}]{\sphinxcrossref{\DUrole{doc}{compute damage/atom}}}}
command.

\sphinxAtStartPar
This can be done, for example, by using the built\sphinxhyphen{}in visualizer of the
{\hyperref[\detokenize{dump_image::doc}]{\sphinxcrossref{\DUrole{doc}{dump image or dump movie}}}} command to create snapshot
images or a movie. Below are example command lines for using dump image
with the {\hyperref[\detokenize{Howto_peri:periexample}]{\sphinxcrossref{\DUrole{std,std-ref}{example listed below}}}} and a set of images
created for steps 300, 600, and 2000 this way.

\begin{sphinxVerbatim}[commandchars=\\\{\}]
\PYG{k}{dump            }\PYG{n+nv+nvIdentifier}{D2}\PYG{+w}{ }\PYG{n+nv+nvIdentifier}{all}\PYG{+w}{ }\PYG{n}{image}\PYG{+w}{ }\PYG{l+m}{100}\PYG{+w}{ }\PYG{n}{dump.peri.}\PYG{o}{*}\PYG{n}{.png}\PYG{+w}{ }\PYG{n}{c\PYGZus{}C1}\PYG{+w}{ }\PYG{n}{type}\PYG{+w}{ }\PYG{n}{box}\PYG{+w}{ }\PYG{n}{no}\PYG{+w}{ }\PYG{l+m}{0.0}\PYG{+w}{ }\PYG{n}{view}\PYG{+w}{ }\PYG{l+m}{30}\PYG{+w}{ }\PYG{l+m}{60}\PYG{+w}{ }\PYG{n}{zoom}\PYG{+w}{ }\PYG{l+m}{1.5}\PYG{+w}{ }\PYG{n}{up}\PYG{+w}{ }\PYG{l+m}{0}\PYG{+w}{ }\PYG{l+m}{0}\PYG{+w}{ }\PYG{o}{\PYGZhy{}}\PYG{l+m}{1}\PYG{+w}{  }\PYG{n}{ssao}\PYG{+w}{ }\PYG{n}{yes}\PYG{+w}{ }\PYG{l+m}{4539}\PYG{+w}{ }\PYG{l+m}{0.6}
\PYG{k}{dump\PYGZus{}modify     }\PYG{n+nv+nvIdentifier}{D2}\PYG{+w}{ }\PYG{n}{pad}\PYG{+w}{ }\PYG{l+m}{5}\PYG{+w}{ }\PYG{n}{adiam}\PYG{+w}{ }\PYG{o}{*}\PYG{+w}{ }\PYG{l+m}{0.001}\PYG{+w}{ }\PYG{n}{amap}\PYG{+w}{ }\PYG{l+m}{0.0}\PYG{+w}{ }\PYG{l+m}{1.0}\PYG{+w}{ }\PYG{n}{ca}\PYG{+w}{ }\PYG{l+m}{0.1}\PYG{+w}{ }\PYG{l+m}{3}\PYG{+w}{ }\PYG{n}{min}\PYG{+w}{ }\PYG{n}{blue}\PYG{+w}{ }\PYG{l+m}{0.5}\PYG{+w}{ }\PYG{n}{yellow}\PYG{+w}{ }\PYG{n}{max}\PYG{+w}{ }\PYG{n}{red}
\end{sphinxVerbatim}

\sphinxAtStartPar
\sphinxincludegraphics[width=0.320\linewidth]{{dump.peri.300}.png}  \sphinxincludegraphics[width=0.320\linewidth]{{dump.peri.600}.png}  \sphinxincludegraphics[width=0.320\linewidth]{{dump.peri.2000}.png}

\sphinxAtStartPar
For interactive visualization, the \sphinxhref{https://ovito.org}{Ovito} is very
convenient to use. Below are steps to create a visualization of the
{\hyperref[\detokenize{Howto_peri:periexample}]{\sphinxcrossref{\DUrole{std,std-ref}{same example from below}}}} now using the generated
trajectory in the \sphinxcode{\sphinxupquote{dump.peri}} file.
\begin{itemize}
\item {} 
\sphinxAtStartPar
Launch Ovito

\item {} 
\sphinxAtStartPar
File \sphinxhyphen{}\textgreater{} Load File \sphinxhyphen{}\textgreater{} \sphinxcode{\sphinxupquote{dump.peri}}

\item {} 
\sphinxAtStartPar
Select “\sphinxhyphen{}\textgreater{} Particle types” and under “Appearance” set “Display radius:” to 0.0005

\item {} 
\sphinxAtStartPar
From the “Add modification:” drop down list select “Color coding”

\item {} 
\sphinxAtStartPar
Under “Color coding” select from the “Color gradient” drop down list “Jet”

\item {} 
\sphinxAtStartPar
Also under “Color coding” set “Start value:” to 0 and “End value:” to 1

\item {} 
\sphinxAtStartPar
You can improve the image quality by adding the “Ambient occlusion” modification

\end{itemize}

\begin{figure}[htbp]
\centering
\capstart

\noindent\sphinxincludegraphics{{ovito-peri-snap}.png}
\caption{Screenshot of visualizing a trajectory with Ovito}\label{\detokenize{Howto_peri:id3}}\end{figure}


\paragraph{Pitfalls}
\label{\detokenize{Howto_peri:pitfalls}}
\sphinxAtStartPar
\sphinxstylestrong{Parallel Scalability}

\sphinxAtStartPar
LAMMPS operates in parallel in a {\hyperref[\detokenize{Developer_par_part::doc}]{\sphinxcrossref{\DUrole{doc}{spatial\sphinxhyphen{}decomposition mode}}}}, where each processor owns a spatial subdomain of
the overall simulation domain and communicates with its neighboring
processors via distributed\sphinxhyphen{}memory message passing (MPI) to acquire ghost
atom information to allow forces on the atoms it owns to be
computed. LAMMPS also uses Verlet neighbor lists which are recomputed
every few timesteps as particles move. On these timesteps, particles
also migrate to new processors as needed. LAMMPS decomposes the overall
simulation domain so that spatial subdomains of nearly equal volume are
assigned to each processor. When each subdomain contains nearly the
same number of particles, this results in a reasonable load balance
among all processors. As is more typical with some peridynamic
simulations, some subdomains may contain many particles while other
subdomains contain few particles, resulting in a load imbalance that
impacts parallel scalability.

\sphinxAtStartPar
\sphinxstylestrong{Setting the “skin” distance}

\sphinxAtStartPar
The {\hyperref[\detokenize{neighbor::doc}]{\sphinxcrossref{\DUrole{doc}{neighbor}}}} command with LAMMPS is used to set the
so\sphinxhyphen{}called “skin” distance used when building neighbor lists. All atom
pairs within a cutoff distance equal to the horizon \(\delta\) plus
the skin distance are stored in the list. Unexpected crashes in LAMMPS
may be due to too small a skin distance. The skin should be set to a
value such that \(\delta\) plus the skin distance is larger than the
maximum possible distance between two bonded particles. For example, if
\(s_{00}\) is increased, the skin distance may also need to be
increased.

\sphinxAtStartPar
\sphinxstylestrong{“Lost” particles}

\sphinxAtStartPar
All particles are contained within the “simulation box” of LAMMPS. The
boundaries of this box may change with time, or not, depending on how
the LAMMPS {\hyperref[\detokenize{boundary::doc}]{\sphinxcrossref{\DUrole{doc}{boundary}}}} command has been set. If a
particle drifts outside the simulation box during the course of a
simulation, it is called \sphinxstyleemphasis{lost}.

\sphinxAtStartPar
As an option of the {\hyperref[\detokenize{thermo_modify::doc}]{\sphinxcrossref{\DUrole{doc}{themo\_modify}}}} command of
LAMMPS, the lost keyword determines whether LAMMPS checks for lost atoms
each time it computes thermodynamics and what it does if atoms are
lost. If the value is \sphinxstyleemphasis{ignore}, LAMMPS does not check for lost atoms. If
the value is \sphinxstyleemphasis{error} or \sphinxstyleemphasis{warn}, LAMMPS checks and either issues an error
or warning. The code will exit with an error and continue with a
warning. This can be a useful debugging option. The default behavior of
LAMMPS is to exit with an error if a particle is lost.

\sphinxAtStartPar
The peridynamic module within LAMMPS does not check for lost atoms. If a
particle with unbroken bonds is lost, those bonds are marked as broken
by the remaining particles.

\sphinxAtStartPar
\sphinxstylestrong{Defining the peridynamic horizon} \(\mathbf{\delta}\)

\sphinxAtStartPar
In the {\hyperref[\detokenize{pair_coeff::doc}]{\sphinxcrossref{\DUrole{doc}{pair\_coeff}}}} command, the user must specify the
horizon \(\delta\). This argument determines which particles are
bonded when the simulation is initialized. It is recommended that
\(\delta\) be set to a small fraction of a lattice constant larger than
desired.

\sphinxAtStartPar
For example, if the lattice constant is 0.0005 and you wish to set the
horizon to three times the lattice constant, then set \(\delta\) to
be 0.0015001, a value slightly larger than three times the lattice
constant. This guarantees that particles three lattice constants away
from each other are still bonded. If \(\delta\) is set to 0.0015,
for example, floating point error may result in some pairs of particles
three lattice constants apart not being bonded.

\sphinxAtStartPar
\sphinxstylestrong{Breaking bonds too early}

\sphinxAtStartPar
For technical reasons, the bonds in the simulation are not created until
the end of the first timestep of the simulation. Therefore, one should
not attempt to break bonds until at least the second step of the
simulation.


\paragraph{Bugs}
\label{\detokenize{Howto_peri:bugs}}
\sphinxAtStartPar
The user is cautioned that this code is a beta release. If you are
confident that you have found a bug in the peridynamic module, please
report it in a \sphinxtitleref{GitHub Issue \textless{}https://github.com/lammps/lammps/issues\textgreater{}}
or send an email to the LAMMPS developers. First, check the \sphinxhref{https://www.lammps.org/bug.html}{New
features and bug fixes} section of
the LAMMPS website site to see if the bug has already been reported or
fixed. If not, the most useful thing you can do for us is to isolate the
problem. Run it on the smallest number of atoms and fewest number of
processors and with the simplest input script that reproduces the
bug. In your message, describe the problem and any ideas you have as to
what is causing it or where in the code the problem might be.  We’ll
request your input script and data files if necessary.


\paragraph{Modifying and Extending the Peridynamic Module}
\label{\detokenize{Howto_peri:modifying-and-extending-the-peridynamic-module}}
\sphinxAtStartPar
To add new features or peridynamic potentials to the peridynamic module,
the user is referred to the {\hyperref[\detokenize{Modify::doc}]{\sphinxcrossref{\DUrole{doc}{Modifying \& extending LAMMPS}}}}
section. To develop a new bond\sphinxhyphen{}based material, start with the \sphinxstyleemphasis{peri/pmb}
pair style as a template. To develop a new state\sphinxhyphen{}based material, start
with the \sphinxstyleemphasis{peri/lps} pair style as a template.


\subsubsection{A Numerical Example}
\label{\detokenize{Howto_peri:a-numerical-example}}
\sphinxAtStartPar
To introduce the peridynamic implementation within LAMMPS, we replicate
a numerical experiment taken from section 6 of {\hyperref[\detokenize{Howto_peri:silling2005}]{\sphinxcrossref{\DUrole{std,std-ref}{(Silling 2005)}}}}.


\paragraph{Problem Description and Setup}
\label{\detokenize{Howto_peri:problem-description-and-setup}}
\sphinxAtStartPar
We consider the impact of a rigid sphere on a homogeneous disk of
brittle material. The sphere has diameter \(0.01\) m and velocity
100 m/s directed normal to the surface of the target. The target
material has density \(\rho = 2200\) kg/m:math:\sphinxtitleref{\textasciicircum{}3}. A PMB material
model is used with \(K = 14.9\) GPa and critical bond stretch
parameters given by \(s_{00} = 0.0005\) and \(\alpha = 0.25\). A
three\sphinxhyphen{}dimensional simple cubic lattice is constructed with lattice
constant \(0.0005\) m and horizon \(0.0015\) m. (The horizon is
three times the lattice constant.) The target is a cylinder of diameter
\(0.074\) m and thickness \(0.0025\) m, and the associated
lattice contains 103,110 particles. Each particle \(i\) has volume
fraction \(V_i = 1.25 \times 10^{-10} \textrm{m}^3\).

\sphinxAtStartPar
The spring constant in the PMB material model is (see {\hyperref[\detokenize{Howto_peri:peric}]{\sphinxcrossref{\DUrole{std,std-ref}{(6)}}}})
\begin{equation*}
\begin{split}c = \frac{18k}{\pi \delta^4} = \frac{18 (14.9 \times 10^9)}{ \pi (1.5 \times 10^{-3})^4} \approx 1.6863 \times 10^{22}.\end{split}
\end{equation*}
\sphinxAtStartPar
The CFL analysis from {\hyperref[\detokenize{Howto_peri:silling2005}]{\sphinxcrossref{\DUrole{std,std-ref}{(Silling2005)}}}} shows that a
timestep of \(1.0 \times 10^{-7}\) is safe.

\sphinxAtStartPar
We observe here that in IEEE double\sphinxhyphen{}precision floating point arithmetic
when computing the bond stretch \(s(t,\mathbf{\eta},\mathbf{\xi})\)
at each iteration where \(\left\Vert {\mathbf{\eta}+\mathbf{\xi}}
\right\Vert\) is computed during the iteration and \(\left\Vert
{\mathbf{\xi}} \right\Vert\) was computed and stored for the initial
lattice, it may be that \(fl(s) = \varepsilon\) with \(\left|
\varepsilon \right| \leq \varepsilon_{machine}\) for an unstretched
bond. Taking \(\varepsilon = 2.220446049250313 \times 10^{-16}\), we
see that the value \(c s V_i \approx 4.68 \times 10^{-4}\), computed
when determining \(f\), is perhaps larger than we would like,
especially when the true force should be zero. One simple way to avoid
this issue is to insert the following instructions in Algorithm
{\hyperref[\detokenize{Howto_peri:algperipmb}]{\sphinxcrossref{\DUrole{std,std-ref}{Algorithm 5}}}} after instruction 21 (and similarly for
Algorithm {\hyperref[\detokenize{Howto_peri:algperilps}]{\sphinxcrossref{\DUrole{std,std-ref}{Algorithm 2}}}}):
\begin{quote}

\begin{DUlineblock}{0em}
\item[] \sphinxstylestrong{if} \(\left| dr \right| < \varepsilon_{machine}\) \sphinxstylestrong{then}
\item[]
\begin{DUlineblock}{\DUlineblockindent}
\item[] \(dr = 0\)
\end{DUlineblock}
\item[] \sphinxstylestrong{end if}
\end{DUlineblock}
\end{quote}

\sphinxAtStartPar
Qualitatively, this says that displacements from equilibrium on the
order of \(10^{-16}\)m are taken to be exactly zero, a seemingly
reasonable assumption.


\paragraph{The Projectile}
\label{\detokenize{Howto_peri:the-projectile}}
\sphinxAtStartPar
The projectile used in the following experiments is not the one used in
{\hyperref[\detokenize{Howto_peri:silling2005}]{\sphinxcrossref{\DUrole{std,std-ref}{(Silling 2005)}}}}. The projectile used here exerts a
force
\begin{equation*}
\begin{split}F(r) = - k_s (r - R)^2\end{split}
\end{equation*}
\sphinxAtStartPar
on each atom where \(k_s\) is a specified force constant, \(r\) is
the distance from the atom to the center of the indenter, and \(R\)
is the radius of the projectile. The force is repulsive and \(F(r) =
0\) for \(r > R\). For our problem, the projectile radius \(R =
0.05\) m, and we have chosen \(k_s = 1.0 \times 10^{17}\) (compare
with {\hyperref[\detokenize{Howto_peri:peric}]{\sphinxcrossref{\DUrole{std,std-ref}{(6)}}}} above).


\paragraph{Writing the LAMMPS Input File}
\label{\detokenize{Howto_peri:writing-the-lammps-input-file}}
\sphinxAtStartPar
We discuss the example input script {\hyperref[\detokenize{Howto_peri:periexample}]{\sphinxcrossref{\DUrole{std,std-ref}{listed below}}}}.
In line 2 we specify that SI units are to be used. We specify the
dimension (3) and boundary conditions (“shrink\sphinxhyphen{}wrapped”) for the
computational domain in lines 3 and 4. In line 5 we specify that
peridynamic particles are to be used for this simulation. In line 7, we
set the “skin” distance used in building the LAMMPS neighbor list. In
line 8 we set the lattice constant (in meters) and in line 10 we define
the spatial region where the target will be placed. In line 12 we
specify a rectangular box enclosing the target region that defines the
simulation domain. Line 14 fills the target region with atoms. Lines 15
and 17 define the peridynamic material model, and lines 19 and 21 set
the particle density and particle volume, respectively. The particle
volume should be set to the cube of the lattice constant for a simple
cubic lattice. Line 23 sets the initial velocity of all particles to
zero. Line 25 instructs LAMMPS to integrate time with velocity\sphinxhyphen{}Verlet,
and lines 27\sphinxhyphen{}30 create the spherical projectile, sending it with a
velocity of 100 m/s towards the target. Line 32 declares a compute style
for the damage (percentage of broken bonds) associated with each
particle.  Line 33 sets the timestep, line 34 instructs LAMMPS to
provide a screen dump of thermodynamic quantities every 200 timesteps,
and line 35 instructs LAMMPS to create a data file (\sphinxcode{\sphinxupquote{dump.peri}}) with
a complete snapshot of the system every 100 timesteps. This file can be
used to create still images or movies. Finally, line 36 instructs LAMMPS
to run for 2000 timesteps.
\sphinxSetupCaptionForVerbatim{Peridynamics Example LAMMPS Input Script}
\def\sphinxLiteralBlockLabel{\label{\detokenize{Howto_peri:id4}}\label{\detokenize{Howto_peri:periexample}}}
\begin{sphinxVerbatim}[commandchars=\\\{\},numbers=left,firstnumber=1,stepnumber=1]
\PYG{c}{\PYGZsh{} 3D Peridynamic simulation with projectile\PYGZdq{}}
\PYG{k}{units}\PYG{+w}{           }\PYG{n}{si}
\PYG{k}{dimension}\PYG{+w}{       }\PYG{l+m}{3}
\PYG{k}{boundary}\PYG{+w}{        }\PYG{n}{s}\PYG{+w}{ }\PYG{n}{s}\PYG{+w}{ }\PYG{n}{s}
\PYG{k}{atom\PYGZus{}style}\PYG{+w}{      }\PYG{n}{peri}
\PYG{k}{atom\PYGZus{}modify}\PYG{+w}{     }\PYG{n}{map}\PYG{+w}{ }\PYG{n}{array}
\PYG{k}{neighbor}\PYG{+w}{        }\PYG{l+m}{0.0010}\PYG{+w}{ }\PYG{n}{bin}
\PYG{k}{lattice}\PYG{+w}{         }\PYG{n}{sc}\PYG{+w}{ }\PYG{l+m}{0.0005}
\PYG{c}{\PYGZsh{} Create desired target}
\PYG{k}{region          }\PYG{n+nv+nvIdentifier}{target}\PYG{+w}{ }\PYG{n}{cylinder}\PYG{+w}{ }\PYG{n}{y}\PYG{+w}{ }\PYG{l+m}{0.0}\PYG{+w}{ }\PYG{l+m}{0.0}\PYG{+w}{ }\PYG{l+m}{0.037}\PYG{+w}{ }\PYG{o}{\PYGZhy{}}\PYG{l+m}{0.0025}\PYG{+w}{ }\PYG{l+m}{0.0}\PYG{+w}{ }\PYG{n}{units}\PYG{+w}{ }\PYG{n}{box}
\PYG{c}{\PYGZsh{} Make 1 atom type}
\PYG{k}{create\PYGZus{}box      }\PYG{n+nv+nvIdentifier}{1}\PYG{+w}{ }\PYG{n+nv+nvIdentifier}{target}
\PYG{c}{\PYGZsh{} Create the atoms in the simulation region}
\PYG{k}{create\PYGZus{}atoms}\PYG{+w}{    }\PYG{l+m}{1}\PYG{+w}{ }\PYG{k}{region }\PYG{n+nv+nvIdentifier}{target}
\PYG{k}{pair\PYGZus{}style}\PYG{+w}{      }\PYG{n}{peri}\PYG{o}{/}\PYG{n}{pmb}
\PYG{c}{\PYGZsh{}               \PYGZlt{}type1\PYGZgt{} \PYGZlt{}type2\PYGZgt{}    \PYGZlt{}c\PYGZgt{}    \PYGZlt{}horizon\PYGZgt{}  \PYGZlt{}s00\PYGZgt{} \PYGZlt{}alpha\PYGZgt{}}
\PYG{k}{pair\PYGZus{}coeff}\PYG{+w}{         }\PYG{o}{*}\PYG{+w}{       }\PYG{o}{*}\PYG{+w}{    }\PYG{l+m}{1.6863e22}\PYG{+w}{ }\PYG{l+m}{0.0015001}\PYG{+w}{ }\PYG{l+m}{0.0005}\PYG{+w}{  }\PYG{l+m}{0.25}
\PYG{c}{\PYGZsh{} Set mass density}
\PYG{k}{set}\PYG{+w}{             }\PYG{k}{group }\PYG{n+nv+nvIdentifier}{all}\PYG{+w}{ }\PYG{n}{density}\PYG{+w}{ }\PYG{l+m}{2200}
\PYG{c}{\PYGZsh{} volume = lattice constant\PYGZca{}3}
\PYG{k}{set}\PYG{+w}{             }\PYG{k}{group }\PYG{n+nv+nvIdentifier}{all}\PYG{+w}{ }\PYG{n}{volume}\PYG{+w}{ }\PYG{l+m}{1.25e\PYGZhy{}10}
\PYG{c}{\PYGZsh{} Zero out velocities of particles}
\PYG{k}{velocity        }\PYG{n+nv+nvIdentifier}{all}\PYG{+w}{ }\PYG{n}{set}\PYG{+w}{ }\PYG{l+m}{0.0}\PYG{+w}{ }\PYG{l+m}{0.0}\PYG{+w}{ }\PYG{l+m}{0.0}\PYG{+w}{ }\PYG{n}{sum}\PYG{+w}{ }\PYG{n}{no}\PYG{+w}{ }\PYG{n}{units}\PYG{+w}{ }\PYG{n}{box}
\PYG{c}{\PYGZsh{} Use velocity\PYGZhy{}Verlet time integrator}
\PYG{k}{fix             }\PYG{n+nv+nvIdentifier}{F1}\PYG{+w}{ }\PYG{n+nv+nvIdentifier}{all}\PYG{+w}{ }\PYG{n}{nve}
\PYG{c}{\PYGZsh{} Construct spherical indenter to shatter target}
\PYG{k}{variable        }\PYG{n+nv+nvIdentifier}{y0}\PYG{+w}{ }\PYG{n}{equal}\PYG{+w}{ }\PYG{l+m}{0.00510}
\PYG{k}{variable        }\PYG{n+nv+nvIdentifier}{vy}\PYG{+w}{ }\PYG{n}{equal}\PYG{+w}{ }\PYG{o}{\PYGZhy{}}\PYG{l+m}{100}
\PYG{k}{variable        }\PYG{n+nv+nvIdentifier}{y}\PYG{+w}{ }\PYG{n}{equal}\PYG{+w}{ }\PYG{l+s}{\PYGZdq{}}\PYG{l+s}{v\PYGZus{}y0 + step*dt*v\PYGZus{}vy}\PYG{l+s}{\PYGZdq{}}
\PYG{k}{fix             }\PYG{n+nv+nvIdentifier}{F2}\PYG{+w}{ }\PYG{n+nv+nvIdentifier}{all}\PYG{+w}{ }\PYG{n}{indent}\PYG{+w}{ }\PYG{l+m}{1e17}\PYG{+w}{ }\PYG{n}{sphere}\PYG{+w}{ }\PYG{l+m}{0.0000}\PYG{+w}{ }\PYG{n}{v\PYGZus{}y}\PYG{+w}{ }\PYG{l+m}{0.0000}\PYG{+w}{ }\PYG{l+m}{0.0050}\PYG{+w}{ }\PYG{n}{units}\PYG{+w}{ }\PYG{n}{box}
\PYG{c}{\PYGZsh{} Compute damage for each particle}
\PYG{k}{compute         }\PYG{n+nv+nvIdentifier}{C1}\PYG{+w}{ }\PYG{n+nv+nvIdentifier}{all}\PYG{+w}{ }\PYG{n}{damage}\PYG{o}{/}\PYG{n}{atom}
\PYG{k}{timestep}\PYG{+w}{        }\PYG{l+m}{1.0e\PYGZhy{}7}
\PYG{k}{thermo}\PYG{+w}{          }\PYG{l+m}{200}
\PYG{k}{dump            }\PYG{n+nv+nvIdentifier}{D1}\PYG{+w}{ }\PYG{n+nv+nvIdentifier}{all}\PYG{+w}{ }\PYG{n}{custom}\PYG{+w}{ }\PYG{l+m}{100}\PYG{+w}{ }\PYG{n}{dump.peri}\PYG{+w}{ }\PYG{n}{id}\PYG{+w}{ }\PYG{n}{type}\PYG{+w}{ }\PYG{n}{x}\PYG{+w}{ }\PYG{n}{y}\PYG{+w}{ }\PYG{n}{z}\PYG{+w}{ }\PYG{n}{c\PYGZus{}C1}
\PYG{k}{run}\PYG{+w}{             }\PYG{l+m}{2000}
\end{sphinxVerbatim}

\begin{sphinxadmonition}{note}{Note:}
\sphinxAtStartPar
To use the LPS model, replace line 15 with {\hyperref[\detokenize{pair_peri::doc}]{\sphinxcrossref{\DUrole{doc}{pair\_style peri/lps}}}} and modify line 16 accordingly.
\end{sphinxadmonition}


\paragraph{Numerical Results and Discussion}
\label{\detokenize{Howto_peri:numerical-results-and-discussion}}
\sphinxAtStartPar
We ran the {\hyperref[\detokenize{Howto_peri:periexample}]{\sphinxcrossref{\DUrole{std,std-ref}{input script from above}}}}. Images of the
disk (projectile not shown) appear in Figure below.  The plot of damage
on the top monolayer was created by coloring each particle according to
its damage.

\sphinxAtStartPar
The symmetry in the computed solution arises because a “perfect” lattice
was used, and a because a perfectly spherical projectile impacted the
lattice at its geometric center. To break the symmetry in the solution,
the nodes in the peridynamic body may be perturbed slightly from the
lattice sites. To do this, the lattice of points can be slightly
perturbed using the {\hyperref[\detokenize{displace_atoms::doc}]{\sphinxcrossref{\DUrole{doc}{displace\_atoms}}}} command.

\begin{figure}[htbp]
\centering
\capstart

\noindent\sphinxincludegraphics{{pdlammps_fig2}.png}
\caption{Target during (a) and after (b,c) impact}\label{\detokenize{Howto_peri:id5}}\label{\detokenize{Howto_peri:figperitarget}}\end{figure}


\bigskip\hrule\bigskip

\phantomsection\label{\detokenize{Howto_peri:emmrich2007}}
\sphinxAtStartPar
\sphinxstylestrong{(Emmrich)} Emmrich, Weckner, Commun. Math. Sci., 5, 851\sphinxhyphen{}864 (2007),

\phantomsection\label{\detokenize{Howto_peri:parks2}}
\sphinxAtStartPar
\sphinxstylestrong{(Parks)} Parks, Lehoucq, Plimpton, Silling, Comp Phys Comm, 179(11), 777\sphinxhyphen{}783 (2008).

\phantomsection\label{\detokenize{Howto_peri:silling2000-2}}
\sphinxAtStartPar
\sphinxstylestrong{(Silling 2000)} Silling, J Mech Phys Solids, 48, 175\sphinxhyphen{}209 (2000).

\phantomsection\label{\detokenize{Howto_peri:silling2005}}
\sphinxAtStartPar
\sphinxstylestrong{(Silling 2005)} Silling Askari, Computer and Structures, 83, 1526\sphinxhyphen{}1535 (2005).

\phantomsection\label{\detokenize{Howto_peri:silling2007-2}}
\sphinxAtStartPar
\sphinxstylestrong{(Silling 2007)} Silling, Epton, Weckner, Xu, Askari, J Elasticity, 88, 151\sphinxhyphen{}184 (2007).

\phantomsection\label{\detokenize{Howto_peri:seleson2010}}
\sphinxAtStartPar
\sphinxstylestrong{(Seleson 2010)} Seleson, Parks, Int J Mult Comp Eng 9(6), pp. 689\sphinxhyphen{}706, 2011.

\sphinxstepscope


\subsection{Manifolds (surfaces)}
\label{\detokenize{Howto_manifold:manifolds-surfaces}}\label{\detokenize{Howto_manifold::doc}}
\sphinxAtStartPar
\sphinxstylestrong{Overview:}

\sphinxAtStartPar
This page is not about a LAMMPS input script command, but about
manifolds, which are generalized surfaces, as defined and used by the
MANIFOLD package, to track particle motion on the manifolds.  See
the src/MANIFOLD/README file for more details about the package
and its commands.

\sphinxAtStartPar
Below is a list of currently supported manifolds by the MANIFOLD
package, their parameters and a short description of them.  The
parameters listed here are in the same order as they should be passed
to the relevant fixes.


\begin{savenotes}\sphinxattablestart
\sphinxthistablewithglobalstyle
\centering
\begin{tabulary}{\linewidth}[t]{TTTT}
\sphinxtoprule
\sphinxtableatstartofbodyhook
\sphinxAtStartPar
\sphinxstyleemphasis{manifold}
&
\sphinxAtStartPar
\sphinxstyleemphasis{parameters}
&
\sphinxAtStartPar
\sphinxstyleemphasis{equation}
&
\sphinxAtStartPar
\sphinxstyleemphasis{description}
\\
\sphinxhline
\sphinxAtStartPar
cylinder
&
\sphinxAtStartPar
R
&
\sphinxAtStartPar
x\textasciicircum{}2 + y\textasciicircum{}2 \sphinxhyphen{} R\textasciicircum{}2 = 0
&
\sphinxAtStartPar
Cylinder along z\sphinxhyphen{}axis, axis going through (0,0,0)
\\
\sphinxhline
\sphinxAtStartPar
cylinder\_dent
&
\sphinxAtStartPar
R l a
&
\sphinxAtStartPar
x\textasciicircum{}2 + y\textasciicircum{}2 \sphinxhyphen{} r(z)\textasciicircum{}2 = 0, r(x) = R if | z | \textgreater{} l, r(z) = R \sphinxhyphen{} a*(1 + cos(z/l))/2 otherwise
&
\sphinxAtStartPar
A cylinder with a dent around z = 0
\\
\sphinxhline
\sphinxAtStartPar
dumbbell
&
\sphinxAtStartPar
a A B c
&
\sphinxAtStartPar
\sphinxhyphen{}( x\textasciicircum{}2 + y\textasciicircum{}2 ) + (a\textasciicircum{}2 \sphinxhyphen{} z\textasciicircum{}2/c\textasciicircum{}2) * ( 1 + (A*sin(B*z\textasciicircum{}2))\textasciicircum{}4) = 0
&
\sphinxAtStartPar
A dumbbell
\\
\sphinxhline
\sphinxAtStartPar
ellipsoid
&
\sphinxAtStartPar
a  b c
&
\sphinxAtStartPar
(x/a)\textasciicircum{}2 + (y/b)\textasciicircum{}2 + (z/c)\textasciicircum{}2 = 0
&
\sphinxAtStartPar
An ellipsoid
\\
\sphinxhline
\sphinxAtStartPar
gaussian\_bump
&
\sphinxAtStartPar
A l rc1 rc2
&
\sphinxAtStartPar
if( x \textless{} rc1) \sphinxhyphen{}z + A * exp( \sphinxhyphen{}x\textasciicircum{}2 / (2 l\textasciicircum{}2) ); else if( x \textless{} rc2 ) \sphinxhyphen{}z + a + b*x + c*x\textasciicircum{}2 + d*x\textasciicircum{}3; else z
&
\sphinxAtStartPar
A Gaussian bump at x = y = 0, smoothly tapered to a flat plane z = 0.
\\
\sphinxhline
\sphinxAtStartPar
plane
&
\sphinxAtStartPar
a b c x0 y0 z0
&
\sphinxAtStartPar
a*(x\sphinxhyphen{}x0) + b*(y\sphinxhyphen{}y0) + c*(z\sphinxhyphen{}z0) = 0
&
\sphinxAtStartPar
A plane with normal (a,b,c) going through point (x0,y0,z0)
\\
\sphinxhline
\sphinxAtStartPar
plane\_wiggle
&
\sphinxAtStartPar
a w
&
\sphinxAtStartPar
z \sphinxhyphen{} a*sin(w*x) = 0
&
\sphinxAtStartPar
A plane with a sinusoidal modulation on z along x.
\\
\sphinxhline
\sphinxAtStartPar
sphere
&
\sphinxAtStartPar
R
&
\sphinxAtStartPar
x\textasciicircum{}2 + y\textasciicircum{}2 + z\textasciicircum{}2 \sphinxhyphen{} R\textasciicircum{}2 = 0
&
\sphinxAtStartPar
A sphere of radius R
\\
\sphinxhline
\sphinxAtStartPar
supersphere
&
\sphinxAtStartPar
R q
&
\sphinxAtStartPar
| x |\textasciicircum{}q + | y |\textasciicircum{}q + | z |\textasciicircum{}q \sphinxhyphen{} R\textasciicircum{}q = 0
&
\sphinxAtStartPar
A supersphere of hyperradius R
\\
\sphinxhline
\sphinxAtStartPar
spine
&
\sphinxAtStartPar
a, A, B, B2, c
&
\sphinxAtStartPar
\sphinxhyphen{}(x\textasciicircum{}2 + y\textasciicircum{}2) + (a\textasciicircum{}2 \sphinxhyphen{} z\textasciicircum{}2/f(z)\textasciicircum{}2)*(1 + (A*sin(g(z)*z\textasciicircum{}2))\textasciicircum{}4), f(z) = c if z \textgreater{} 0, 1 otherwise; g(z) = B if z \textgreater{} 0, B2 otherwise
&
\sphinxAtStartPar
An approximation to a dendritic spine
\\
\sphinxhline
\sphinxAtStartPar
spine\_two
&
\sphinxAtStartPar
a, A, B, B2, c
&
\sphinxAtStartPar
\sphinxhyphen{}(x\textasciicircum{}2 + y\textasciicircum{}2) + (a\textasciicircum{}2 \sphinxhyphen{} z\textasciicircum{}2/f(z)\textasciicircum{}2)*(1 + (A*sin(g(z)*z\textasciicircum{}2))\textasciicircum{}2), f(z) = c if z \textgreater{} 0, 1 otherwise; g(z) = B if z \textgreater{} 0, B2 otherwise
&
\sphinxAtStartPar
Another approximation to a dendritic spine
\\
\sphinxhline
\sphinxAtStartPar
thylakoid
&
\sphinxAtStartPar
wB LB lB
&
\sphinxAtStartPar
Various, see {\hyperref[\detokenize{Howto_manifold:paquay1}]{\sphinxcrossref{\DUrole{std,std-ref}{(Paquay)}}}}
&
\sphinxAtStartPar
A model grana thylakoid consisting of two block\sphinxhyphen{}like compartments connected by a bridge of width wB, length LB and taper length lB
\\
\sphinxhline
\sphinxAtStartPar
torus
&
\sphinxAtStartPar
R r
&
\sphinxAtStartPar
(R \sphinxhyphen{} sqrt( x\textasciicircum{}2 + y\textasciicircum{}2 ) )\textasciicircum{}2 + z\textasciicircum{}2 \sphinxhyphen{} r\textasciicircum{}2
&
\sphinxAtStartPar
A torus with large radius R and small radius r, centered on (0,0,0)
\\
\sphinxbottomrule
\end{tabulary}
\sphinxtableafterendhook\par
\sphinxattableend\end{savenotes}
\phantomsection\label{\detokenize{Howto_manifold:paquay1}}
\sphinxAtStartPar
\sphinxstylestrong{(Paquay)} Paquay and Kusters, Biophys. J., 110, 6, (2016).
preprint available at \sphinxhref{https://arxiv.org/abs/1411.3019/}{arXiv:1411.3019}.

\sphinxstepscope


\subsection{Reproducing hydrodynamics and elastic objects (RHEO)}
\label{\detokenize{Howto_rheo:reproducing-hydrodynamics-and-elastic-objects-rheo}}\label{\detokenize{Howto_rheo::doc}}
\sphinxAtStartPar
The RHEO package is a hybrid implementation of smoothed particle
hydrodynamics (SPH) for fluid flow, which can couple to the {\hyperref[\detokenize{Howto_bpm::doc}]{\sphinxcrossref{\DUrole{doc}{BPM package}}}} to model solid elements. RHEO combines these methods to enable
mesh\sphinxhyphen{}free modeling of multi\sphinxhyphen{}phase material systems. Its SPH solver supports
many advanced options including reproducing kernels, particle shifting, free
surface identification, and solid surface reconstruction. To model fluid\sphinxhyphen{}solid
systems, the status of particles can dynamically change between a fluid and
solid state, e.g. during melting/solidification, which determines how they
interact and their physical behavior. The package is designed with modularity
in mind, so one can easily turn various features on/off, adjust physical
details of the system, or develop new capabilities. For instance, the numerics
associated with calculating gradients, reproducing kernels, etc. are separated
into distinct classes to simplify the development of new integration schemes
which can call these calculations. Additional numerical details can be found in
{\hyperref[\detokenize{Howto_rheo:howto-rheo-clemmer}]{\sphinxcrossref{\DUrole{std,std-ref}{(Clemmer)}}}}. Example movies illustrating some of these
capabilities are found at \sphinxurl{https://www.lammps.org/movies.html\#rheopackage}.

\sphinxAtStartPar
Note, if you simply want to run a traditional SPH simulation, the {\hyperref[\detokenize{Packages_details:pkg-sph}]{\sphinxcrossref{\DUrole{std,std-ref}{SPH package}}}} package is likely better suited for your application. It has fewer advanced
features and therefore benefits from improved performance. The {\hyperref[\detokenize{Packages_details:pkg-machdyn}]{\sphinxcrossref{\DUrole{std,std-ref}{MACHDYN}}}} package for solids may also be relevant for fluid\sphinxhyphen{}solid problems.


\bigskip\hrule\bigskip


\sphinxAtStartPar
At the core of the package is {\hyperref[\detokenize{fix_rheo::doc}]{\sphinxcrossref{\DUrole{doc}{fix rheo}}}} which integrates
particle trajectories and controls many optional features (e.g. the use
of reproducing kernels). In conjunction to fix rheo, one must specify an
instance of {\hyperref[\detokenize{fix_rheo_pressure::doc}]{\sphinxcrossref{\DUrole{doc}{fix rheo/pressure}}}} and
{\hyperref[\detokenize{fix_rheo_viscosity::doc}]{\sphinxcrossref{\DUrole{doc}{fix rheo/viscosity}}}} to define a pressure equation
of state and viscosity model, respectively. Optionally, one can model
a heat equation with {\hyperref[\detokenize{fix_rheo_thermal::doc}]{\sphinxcrossref{\DUrole{doc}{fix rheo/thermal}}}}, which also
allows the user to specify equations for a particle’s thermal conductivity,
specific heat, latent heat, and melting temperature. The ordering of these
fixes in an an input script matters. Fix rheo must be defined prior to all
other RHEO fixes.

\sphinxAtStartPar
Typically, RHEO requires atom style rheo. In addition to typical atom
properties like positions and forces, particles store a local density,
viscosity, pressure, and status. If thermal evolution is modeled, one must
use atom style rheo/thermal which also includes a local energy, temperature, and
conductivity. Note that the temperature is always derived from the energy.
This implies the \sphinxstyleemphasis{temperature} attribute of {\hyperref[\detokenize{set::doc}]{\sphinxcrossref{\DUrole{doc}{the set command}}}} does not
affect particles. Instead, one should use the \sphinxstyleemphasis{sph/e} attribute.

\sphinxAtStartPar
The status variable uses bit\sphinxhyphen{}masking to track various properties of a particle
such as its current state of matter (fluid or solid) and its location relative
to a surface. Some of these properties (and others) can be accessed using
{\hyperref[\detokenize{compute_rheo_property_atom::doc}]{\sphinxcrossref{\DUrole{doc}{compute rheo/property/atom}}}}. The \sphinxstyleemphasis{status}
attribute in {\hyperref[\detokenize{set::doc}]{\sphinxcrossref{\DUrole{doc}{the set command}}}} only allows control over the first bit
which sets the state of matter, 0 is fluid and 1 is solid.

\sphinxAtStartPar
Fluid interactions, including pressure forces, viscous forces, and heat exchange,
are calculated using {\hyperref[\detokenize{pair_rheo::doc}]{\sphinxcrossref{\DUrole{doc}{pair rheo}}}}. Unlike typical pair styles,
pair rheo ignores the {\hyperref[\detokenize{special_bonds::doc}]{\sphinxcrossref{\DUrole{doc}{special bond}}}} settings. Instead,
it determines whether to calculate forces based on the status of particles: e.g.,
hydrodynamic forces are only calculated if a fluid particle is involved.


\bigskip\hrule\bigskip


\sphinxAtStartPar
To model elastic objects, there are currently two mechanisms in RHEO, one designed
for bulk solid bodies and the other for thin shells. Both mechanisms rely on
introducing bonded forces between particles and therefore require a hybrid of atom
style bond and rheo (or rheo/thermal).

\sphinxAtStartPar
To create an elastic solid body, one has to (a) change the status of constituent
particles to solid (e.g. with the {\hyperref[\detokenize{set::doc}]{\sphinxcrossref{\DUrole{doc}{set}}}} command), (b) create bpm
bonds between the particles (see the {\hyperref[\detokenize{Howto_bpm::doc}]{\sphinxcrossref{\DUrole{doc}{bpm howto}}}} page for
more details), and (c) use {\hyperref[\detokenize{pair_rheo_solid::doc}]{\sphinxcrossref{\DUrole{doc}{pair rheo/solid}}}} to
apply repulsive contact forces between distinct solid bodies. Akin to pair rheo,
pair rheo/solid considers a particles fluid/solid phase to determine whether to
apply forces. However, unlike pair rheo, pair rheo/solid does obey special bond
settings such that contact forces do not have to be calculated between two bonded
solid particles in the same elastic body.

\sphinxAtStartPar
In systems with thermal evolution, fix rheo/thermal can optionally set a
melting/solidification temperature allowing particles to dynamically swap their
state between fluid and solid when the temperature exceeds or drops below the
critical temperature, respectively. Using the \sphinxstyleemphasis{react} option, one can specify a maximum
bond length and a bond type. Then, when solidifying, particles will search their
local neighbors and automatically create bonds with any neighboring solid particles
in range. For BPM bond styles, bonds will then use the immediate position of the two
particles to calculate a reference state. When melting, particles will delete any
bonds of the specified type when reverting to a fluid state. Special bonds are updated
as bonds are created/broken.

\sphinxAtStartPar
The other option for elastic objects is an elastic shell that is nominally much
thinner than a particle diameter, e.g. a oxide skin which gradually forms over time
on the surface of a fluid. Currently, this is implemented using
{\hyperref[\detokenize{fix_rheo_oxidation::doc}]{\sphinxcrossref{\DUrole{doc}{fix rheo/oxidation}}}} and bond style
{\hyperref[\detokenize{bond_rheo_shell::doc}]{\sphinxcrossref{\DUrole{doc}{rheo/shell}}}}. Essentially, fix rheo/oxidation creates candidate
bonds of a specified type between surface fluid particles within a specified distance.
a newly created rheo/shell bond will then start a timer. While the timer is counting
down, the bond will delete itself if particles move too far apart or move away from the
surface. However, if the timer reaches a user\sphinxhyphen{}defined threshold, then the bond will
activate and apply additional forces to the fluid particles. Bond style rheo/shell
then operates very similarly to a BPM bond style, storing a reference length and
breaking if stretched too far. Unlike the above method, this option does not remove
the underlying fluid interactions (although particle shifting is turned off) and does
not modify special bond settings of particles.

\sphinxAtStartPar
While these two options are not expected to be appropriate for every system,
either framework can be modified to create more suitable models (e.g. by changing the
criteria for creating/deleting a bond or altering force calculations).


\bigskip\hrule\bigskip

\phantomsection\label{\detokenize{Howto_rheo:howto-rheo-clemmer}}
\sphinxAtStartPar
\sphinxstylestrong{(Clemmer)} Clemmer, Pierce, O’Connor, Nevins, Jones, Lechman, Tencer, Appl. Math. Model., 130, 310\sphinxhyphen{}326 (2024).

\sphinxstepscope


\subsection{Magnetic spins}
\label{\detokenize{Howto_spins:magnetic-spins}}\label{\detokenize{Howto_spins::doc}}
\sphinxAtStartPar
The magnetic spin simulations are enabled by the SPIN package, whose
implementation is detailed in {\hyperref[\detokenize{Howto_spins:tranchida}]{\sphinxcrossref{\DUrole{std,std-ref}{Tranchida}}}}.

\sphinxAtStartPar
The model represents the simulation of atomic magnetic spins coupled
to lattice vibrations. The dynamics of those magnetic spins can be used
to simulate a broad range a phenomena related to magneto\sphinxhyphen{}elasticity, or
or to study the influence of defects on the magnetic properties of
materials.

\sphinxAtStartPar
The magnetic spins are interacting with each others and with the
lattice via pair interactions. Typically, the magnetic exchange
interaction can be defined using the
{\hyperref[\detokenize{pair_spin_exchange::doc}]{\sphinxcrossref{\DUrole{doc}{pair/spin/exchange}}}} command. This exchange
applies a magnetic torque to a given spin, considering the orientation
of its neighboring spins and their relative distances.
It also applies a force on the atoms as a function of the spin
orientations and their associated inter\sphinxhyphen{}atomic distances.

\sphinxAtStartPar
The command {\hyperref[\detokenize{fix_precession_spin::doc}]{\sphinxcrossref{\DUrole{doc}{fix precession/spin}}}} allows to
apply a constant magnetic torque on all the spins in the system. This
torque can be an external magnetic field (Zeeman interaction), and an
uniaxial or cubic magnetic anisotropy.

\sphinxAtStartPar
A Langevin thermostat can be applied to those magnetic spins using
{\hyperref[\detokenize{fix_langevin_spin::doc}]{\sphinxcrossref{\DUrole{doc}{fix langevin/spin}}}}. Typically, this thermostat
can be coupled to another Langevin thermostat applied to the atoms
using {\hyperref[\detokenize{fix_langevin::doc}]{\sphinxcrossref{\DUrole{doc}{fix langevin}}}} in order to simulate
thermostatted spin\sphinxhyphen{}lattice systems.

\sphinxAtStartPar
The magnetic damping can also be applied
using {\hyperref[\detokenize{fix_langevin_spin::doc}]{\sphinxcrossref{\DUrole{doc}{fix langevin/spin}}}}.
It allows to either dissipate the thermal energy of the Langevin
thermostat, or to perform a relaxation of the magnetic configuration
toward an equilibrium state.

\sphinxAtStartPar
The command {\hyperref[\detokenize{fix_setforce::doc}]{\sphinxcrossref{\DUrole{doc}{fix setforce/spin}}}} allows to set the
components of the magnetic precession vectors (while erasing and
replacing the previously computed magnetic precession vectors on
the atom).
This command can be used to freeze the magnetic moment of certain
atoms in the simulation by zeroing their precession vector.

\sphinxAtStartPar
The command {\hyperref[\detokenize{fix_nve_spin::doc}]{\sphinxcrossref{\DUrole{doc}{fix nve/spin}}}} can be used to
perform a symplectic integration of the combined dynamics of spins
and atomic motions.

\sphinxAtStartPar
The minimization style {\hyperref[\detokenize{min_spin::doc}]{\sphinxcrossref{\DUrole{doc}{min/spin}}}} can be applied
to the spins to perform a minimization of the spin configuration.

\sphinxAtStartPar
All the computed magnetic properties can be output by two main
commands. The first one is {\hyperref[\detokenize{compute_spin::doc}]{\sphinxcrossref{\DUrole{doc}{compute spin}}}}, that
enables to evaluate magnetic averaged quantities, such as the total
magnetization of the system along x, y, or z, the spin temperature, or
the magnetic energy. The second command
is {\hyperref[\detokenize{compute_property_atom::doc}]{\sphinxcrossref{\DUrole{doc}{compute property/atom}}}}.
It enables to output all the per atom magnetic quantities. Typically,
the orientation of a given magnetic spin, or the magnetic force
acting on this spin.


\bigskip\hrule\bigskip

\phantomsection\label{\detokenize{Howto_spins:tranchida}}
\sphinxAtStartPar
\sphinxstylestrong{(Tranchida)} Tranchida, Plimpton, Thibaudeau and Thompson,
Journal of Computational Physics, 372, 406\sphinxhyphen{}425, (2018).


\section{Tutorials howto}
\label{\detokenize{Howto:tutorials-howto}}
\sphinxstepscope


\subsection{Using CMake with LAMMPS}
\label{\detokenize{Howto_cmake:using-cmake-with-lammps}}\label{\detokenize{Howto_cmake::doc}}
\sphinxAtStartPar
The support for building LAMMPS with CMake is a recent addition to
LAMMPS thanks to the efforts of Christoph Junghans (LANL) and Richard
Berger (LANL).  One of the key strengths of CMake is that it is not
tied to a specific platform or build system. Instead it generates the
files necessary to build and develop for different build systems and on
different platforms.  Note, that this applies to the build system itself
not the LAMMPS code. In other words, without additional porting effort,
it is not possible \sphinxhyphen{} for example \sphinxhyphen{} to compile LAMMPS with Visual C++ on
Windows.  The build system output can also include support files
necessary to program LAMMPS as a project in integrated development
environments (IDE) like Eclipse, Visual Studio, QtCreator, Xcode,
CodeBlocks, Kate and others.

\sphinxAtStartPar
A second important feature of CMake is that it can detect and validate
available libraries, optimal settings, available support tools and so
on, so that by default LAMMPS will take advantage of available tools
without requiring to provide the details about how to enable/integrate
them.

\sphinxAtStartPar
The downside of this approach is, that there is some complexity
associated with running CMake itself and how to customize the building
of LAMMPS.  This tutorial will show how to manage this through some
selected examples.  Please see the chapter about {\hyperref[\detokenize{Build::doc}]{\sphinxcrossref{\DUrole{doc}{building LAMMPS}}}} for descriptions of specific flags and options for LAMMPS in
general and for specific packages.

\sphinxAtStartPar
CMake can be used through either the command\sphinxhyphen{}line interface (CLI)
program \sphinxcode{\sphinxupquote{cmake}} (or \sphinxcode{\sphinxupquote{cmake3}}), a text mode interactive user
interface (TUI) program \sphinxcode{\sphinxupquote{ccmake}} (or \sphinxcode{\sphinxupquote{ccmake3}}), or a graphical user
interface (GUI) program \sphinxcode{\sphinxupquote{cmake\sphinxhyphen{}gui}}.  All of them are portable
software available on all supported platforms and can be used
interchangeably.  As of LAMMPS version 2 August 2023, the minimum
required CMake version is 3.16.

\sphinxAtStartPar
All details about features and settings for CMake are in the \sphinxhref{https://cmake.org/documentation/}{CMake
online documentation}. We focus
below on the most important aspects with respect to compiling LAMMPS.


\subsubsection{Prerequisites}
\label{\detokenize{Howto_cmake:prerequisites}}
\sphinxAtStartPar
This tutorial assumes that you are operating in a command\sphinxhyphen{}line environment
using a shell like Bash or Zsh.
\begin{itemize}
\item {} 
\sphinxAtStartPar
Linux: any Terminal window will work or text console

\item {} 
\sphinxAtStartPar
macOS: launch the Terminal application

\item {} 
\sphinxAtStartPar
Windows 10 or 11: install and run the {\hyperref[\detokenize{Howto_wsl::doc}]{\sphinxcrossref{\DUrole{doc}{Windows Subsystem for Linux}}}}

\item {} 
\sphinxAtStartPar
other Unix\sphinxhyphen{}like operating systems like FreeBSD

\end{itemize}

\begin{sphinxadmonition}{note}{Note:}
\sphinxAtStartPar
It is also possible to use CMake on Windows 10 or 11 through either the Microsoft
Visual Studio IDE with the bundled CMake or from the Windows command prompt using
a separately installed CMake package, both using the native Microsoft Visual C++
compilers and (optionally) the Microsoft MPI SDK.  This tutorial, however, only
covers unix\sphinxhyphen{}like command line interfaces.
\end{sphinxadmonition}

\sphinxAtStartPar
We also assume that you have downloaded and unpacked a recent LAMMPS source code package
or used Git to create a clone of the LAMMPS sources on your compilation machine.

\sphinxAtStartPar
You should change into the top level directory of the LAMMPS source tree all
paths mentioned in the tutorial are relative to that.  Immediately after downloading
it should look like this:

\begin{sphinxVerbatim}[commandchars=\\\{\}]
\PYG{g+gp}{\PYGZdl{} }ls
\PYG{g+go}{bench  doc       lib      potentials  README  tools}
\PYG{g+go}{cmake  examples  LICENSE  python      src}
\end{sphinxVerbatim}


\subsubsection{Build versus source directory}
\label{\detokenize{Howto_cmake:build-versus-source-directory}}
\sphinxAtStartPar
When using CMake the build procedure is separated into multiple distinct phases:
\begin{enumerate}
\sphinxsetlistlabels{\arabic}{enumi}{enumii}{}{.}%
\item {} 
\sphinxAtStartPar
\sphinxstylestrong{Configuration:} detect or define which features and settings
should be enable and used and how LAMMPS should be compiled

\item {} 
\sphinxAtStartPar
\sphinxstylestrong{Compilation:} generate and compile all necessary source files
and build libraries and executables.

\item {} 
\sphinxAtStartPar
\sphinxstylestrong{Installation:} copy selected files from the compilation into
your file system, so they can be used without having to keep the
source and build tree around.

\end{enumerate}

\sphinxAtStartPar
The configuration and compilation of LAMMPS has to happen in a dedicated
\sphinxstyleemphasis{build directory} which must be different from the source directory.
Also the source directory (\sphinxcode{\sphinxupquote{src}}) must remain pristine, so it is not
allowed to “install” packages using the traditional make process and
after an compilation attempt all created source files must be removed.
This can be achieved with \sphinxcode{\sphinxupquote{make no\sphinxhyphen{}all purge}}.

\sphinxAtStartPar
You can pick \sphinxstylestrong{any} folder outside the source tree. We recommend to
create a folder \sphinxcode{\sphinxupquote{build}} in the top\sphinxhyphen{}level directory, or multiple
folders in case you want to have separate builds of LAMMPS with
different options (\sphinxcode{\sphinxupquote{build\sphinxhyphen{}parallel}}, \sphinxcode{\sphinxupquote{build\sphinxhyphen{}serial}}) or with
different compilers (\sphinxcode{\sphinxupquote{build\sphinxhyphen{}gnu}}, \sphinxcode{\sphinxupquote{build\sphinxhyphen{}clang}}, \sphinxcode{\sphinxupquote{build\sphinxhyphen{}intel}})
and so on.  All the auxiliary files created by one build process
(executable, object files, log files, etc) are stored in this directory
or subdirectories within it that CMake creates.


\subsubsection{Running CMake}
\label{\detokenize{Howto_cmake:running-cmake}}

\paragraph{CLI version}
\label{\detokenize{Howto_cmake:cli-version}}
\sphinxAtStartPar
In the (empty) \sphinxcode{\sphinxupquote{build}} directory, we now run the command \sphinxcode{\sphinxupquote{cmake
../cmake}}, which will start the configuration phase and you will see
the progress of the configuration printed to the screen followed by a
summary of the enabled features, options and compiler settings. A typical
summary screen will look like this:

\begin{sphinxVerbatim}[commandchars=\\\{\}]
\PYG{g+gp}{\PYGZdl{} }cmake\PYG{+w}{ }../cmake/
\PYG{g+go}{\PYGZhy{}\PYGZhy{} The CXX compiler identification is GNU 8.2.0}
\PYG{g+go}{\PYGZhy{}\PYGZhy{} Check for working CXX compiler: /opt/tools/gcc\PYGZhy{}8.2.0/bin/c++}
\PYG{g+go}{\PYGZhy{}\PYGZhy{} Check for working CXX compiler: /opt/tools/gcc\PYGZhy{}8.2.0/bin/c++ \PYGZhy{} works}
\PYG{g+go}{\PYGZhy{}\PYGZhy{} Detecting CXX compiler ABI info}
\PYG{g+go}{\PYGZhy{}\PYGZhy{} Detecting CXX compiler ABI info \PYGZhy{} done}
\PYG{g+go}{\PYGZhy{}\PYGZhy{} Detecting CXX compile features}
\PYG{g+go}{\PYGZhy{}\PYGZhy{} Detecting CXX compile features \PYGZhy{} done}
\PYG{g+go}{\PYGZhy{}\PYGZhy{} Found Git: /usr/bin/git (found version \PYGZdq{}2.25.2\PYGZdq{})}
\PYG{g+go}{\PYGZhy{}\PYGZhy{} Running check for auto\PYGZhy{}generated files from make\PYGZhy{}based build system}
\PYG{g+go}{\PYGZhy{}\PYGZhy{} Found MPI\PYGZus{}CXX: /usr/lib64/mpich/lib/libmpicxx.so (found version \PYGZdq{}3.1\PYGZdq{})}
\PYG{g+go}{\PYGZhy{}\PYGZhy{} Found MPI: TRUE (found version \PYGZdq{}3.1\PYGZdq{})}
\PYG{g+go}{\PYGZhy{}\PYGZhy{} Looking for C++ include omp.h}
\PYG{g+go}{\PYGZhy{}\PYGZhy{} Looking for C++ include omp.h \PYGZhy{} found}
\PYG{g+go}{\PYGZhy{}\PYGZhy{} Found OpenMP\PYGZus{}CXX: \PYGZhy{}fopenmp (found version \PYGZdq{}4.5\PYGZdq{})}
\PYG{g+go}{\PYGZhy{}\PYGZhy{} Found OpenMP: TRUE (found version \PYGZdq{}4.5\PYGZdq{})}
\PYG{g+go}{\PYGZhy{}\PYGZhy{} Found JPEG: /usr/lib64/libjpeg.so (found version \PYGZdq{}62\PYGZdq{})}
\PYG{g+go}{\PYGZhy{}\PYGZhy{} Found PNG: /usr/lib64/libpng.so (found version \PYGZdq{}1.6.37\PYGZdq{})}
\PYG{g+go}{\PYGZhy{}\PYGZhy{} Found ZLIB: /usr/lib64/libz.so (found version \PYGZdq{}1.2.11\PYGZdq{})}
\PYG{g+go}{\PYGZhy{}\PYGZhy{} Found GZIP: /usr/bin/gzip}
\PYG{g+go}{\PYGZhy{}\PYGZhy{} Found FFMPEG: /usr/bin/ffmpeg}
\PYG{g+go}{\PYGZhy{}\PYGZhy{} Performing Test COMPILER\PYGZus{}SUPPORTS\PYGZhy{}ffast\PYGZhy{}math}
\PYG{g+go}{\PYGZhy{}\PYGZhy{} Performing Test COMPILER\PYGZus{}SUPPORTS\PYGZhy{}ffast\PYGZhy{}math \PYGZhy{} Success}
\PYG{g+go}{\PYGZhy{}\PYGZhy{} Performing Test COMPILER\PYGZus{}SUPPORTS\PYGZhy{}march=native}
\PYG{g+go}{\PYGZhy{}\PYGZhy{} Performing Test COMPILER\PYGZus{}SUPPORTS\PYGZhy{}march=native \PYGZhy{} Success}
\PYG{g+go}{\PYGZhy{}\PYGZhy{} Looking for C++ include cmath}
\PYG{g+go}{\PYGZhy{}\PYGZhy{} Looking for C++ include cmath \PYGZhy{} found}
\PYG{g+go}{\PYGZhy{}\PYGZhy{} Generating style\PYGZus{}angle.h...}
\PYG{g+go}{[...]}
\PYG{g+go}{\PYGZhy{}\PYGZhy{} Generating lmpinstalledpkgs.h...}
\PYG{g+go}{\PYGZhy{}\PYGZhy{} The following tools and libraries have been found and configured:}
\PYG{g+go}{ * Git}
\PYG{g+go}{ * MPI}
\PYG{g+go}{ * OpenMP}
\PYG{g+go}{ * JPEG}
\PYG{g+go}{ * PNG}
\PYG{g+go}{ * ZLIB}

\PYG{g+go}{\PYGZhy{}\PYGZhy{} \PYGZlt{}\PYGZlt{}\PYGZlt{} Build configuration \PYGZgt{}\PYGZgt{}\PYGZgt{}}
\PYG{g+go}{   Build type:       RelWithDebInfo}
\PYG{g+go}{   Install path:     /home/akohlmey/.local}
\PYG{g+go}{   Generator:        Unix Makefiles using /usr/bin/gmake}
\PYG{g+go}{\PYGZhy{}\PYGZhy{} \PYGZlt{}\PYGZlt{}\PYGZlt{} Compilers and Flags: \PYGZgt{}\PYGZgt{}\PYGZgt{}}
\PYG{g+go}{\PYGZhy{}\PYGZhy{} C++ Compiler:     /opt/tools/gcc\PYGZhy{}8.2.0/bin/c++}
\PYG{g+go}{      Type:          GNU}
\PYG{g+go}{      Version:       8.2.0}
\PYG{g+go}{      C++ Flags:     \PYGZhy{}O2 \PYGZhy{}g \PYGZhy{}DNDEBUG}
\PYG{g+go}{      Defines:       LAMMPS\PYGZus{}SMALLBIG;LAMMPS\PYGZus{}MEMALIGN=64;LAMMPS\PYGZus{}JPEG;LAMMPS\PYGZus{}PNG;LAMMPS\PYGZus{}GZIP;LAMMPS\PYGZus{}FFMPEG}
\PYG{g+go}{      Options:       \PYGZhy{}ffast\PYGZhy{}math;\PYGZhy{}march=native}
\PYG{g+go}{\PYGZhy{}\PYGZhy{} \PYGZlt{}\PYGZlt{}\PYGZlt{} Linker flags: \PYGZgt{}\PYGZgt{}\PYGZgt{}}
\PYG{g+go}{\PYGZhy{}\PYGZhy{} Executable name:  lmp}
\PYG{g+go}{\PYGZhy{}\PYGZhy{} Static library flags:}
\PYG{g+go}{\PYGZhy{}\PYGZhy{} \PYGZlt{}\PYGZlt{}\PYGZlt{} MPI flags \PYGZgt{}\PYGZgt{}\PYGZgt{}}
\PYG{g+go}{\PYGZhy{}\PYGZhy{} MPI includes:     /usr/include/mpich\PYGZhy{}x86\PYGZus{}64}
\PYG{g+go}{\PYGZhy{}\PYGZhy{} MPI libraries:    /usr/lib64/mpich/lib/libmpicxx.so;/usr/lib64/mpich/lib/libmpi.so;}
\PYG{g+go}{\PYGZhy{}\PYGZhy{} Configuring done}
\PYG{g+go}{\PYGZhy{}\PYGZhy{} Generating done}
\PYG{g+go}{\PYGZhy{}\PYGZhy{} Build files have been written to: /home/akohlmey/compile/lammps/build}
\end{sphinxVerbatim}

\sphinxAtStartPar
The \sphinxcode{\sphinxupquote{cmake}} command has one mandatory argument, and that is a folder
with either the file \sphinxcode{\sphinxupquote{CMakeLists.txt}} or \sphinxcode{\sphinxupquote{CMakeCache.txt}}. The
\sphinxcode{\sphinxupquote{CMakeCache.txt}} file is created during the CMake configuration run
and contains all active settings, thus after a first run of CMake
all future runs in the build folder can use the folder \sphinxcode{\sphinxupquote{.}} and CMake
will know where to find the CMake scripts and reload the settings
from the previous step.  This means, that one can modify an existing
configuration by re\sphinxhyphen{}running CMake, but only needs to provide flags
indicating the desired change, everything else will be retained. One
can also mix compilation and configuration, i.e. start with a minimal
configuration and then, if needed, enable additional features and
recompile.

\sphinxAtStartPar
The steps above \sphinxstylestrong{will NOT compile the code}. The compilation can be
started in a portable fashion with \sphinxcode{\sphinxupquote{cmake \sphinxhyphen{}\sphinxhyphen{}build .}}, or you use the
selected built tool, e.g. \sphinxcode{\sphinxupquote{make}}.


\paragraph{TUI version}
\label{\detokenize{Howto_cmake:tui-version}}
\sphinxAtStartPar
For the text mode UI CMake program the basic principle is the same.
You start the command \sphinxcode{\sphinxupquote{ccmake ../cmake}} in the \sphinxcode{\sphinxupquote{build}} folder.


\begin{savenotes}\sphinxattablestart
\sphinxthistablewithglobalstyle
\centering
\begin{tabulary}{\linewidth}[t]{TTT}
\sphinxtoprule
\sphinxtableatstartofbodyhook

\begin{sphinxfigure-in-table}
\centering
\capstart
\noindent\sphinxincludegraphics[scale=0.33]{{ccmake-initial}.png}
\sphinxfigcaption{Initial \sphinxcode{\sphinxupquote{ccmake}} screen}\label{\detokenize{Howto_cmake:id1}}\end{sphinxfigure-in-table}\relax
&

\begin{sphinxfigure-in-table}
\centering
\capstart
\noindent\sphinxincludegraphics[scale=0.33]{{ccmake-config}.png}
\sphinxfigcaption{Configure output of \sphinxcode{\sphinxupquote{ccmake}}}\label{\detokenize{Howto_cmake:id2}}\end{sphinxfigure-in-table}\relax
&

\begin{sphinxfigure-in-table}
\centering
\capstart
\noindent\sphinxincludegraphics[scale=0.33]{{ccmake-options}.png}
\sphinxfigcaption{Options screen of \sphinxcode{\sphinxupquote{ccmake}}}\label{\detokenize{Howto_cmake:id3}}\end{sphinxfigure-in-table}\relax
\\
\sphinxbottomrule
\end{tabulary}
\sphinxtableafterendhook\par
\sphinxattableend\end{savenotes}

\sphinxAtStartPar
This will show you the initial screen (left image) with the empty
configuration cache. Now you type the ‘c’ key to run the configuration
step. That will do a first configuration run and show the summary
(center image). You exit the summary screen with ‘e’ and see now the
main screen with detected options and settings. You can now make changes
by moving and down with the arrow keys of the keyboard and modify
entries. For on/off settings, the enter key will toggle the state.
For others, hitting enter will allow you to modify the value and
you commit the change by hitting the enter key again or cancel using
the escape key.  All “new” settings will be marked with a star ‘*’
and for as long as one setting is marked like this, you have to
re\sphinxhyphen{}run the configuration by hitting the ‘c’ key again, sometimes
multiple times unless the TUI shows the word “generate” next to the
letter ‘g’ and by hitting the ‘g’ key the build files will be written
to the folder and the TUI exits.  You can quit without generating
build files by hitting ‘q’.


\paragraph{GUI version}
\label{\detokenize{Howto_cmake:gui-version}}
\sphinxAtStartPar
For the graphical CMake program the steps are similar to the TUI
version.  You can type the command \sphinxcode{\sphinxupquote{cmake\sphinxhyphen{}gui ../cmake}} in the
\sphinxcode{\sphinxupquote{build}} folder.  In this case the path to the CMake script folder is
not required, it can also be entered from the GUI.


\begin{savenotes}\sphinxattablestart
\sphinxthistablewithglobalstyle
\centering
\begin{tabulary}{\linewidth}[t]{TTT}
\sphinxtoprule
\sphinxtableatstartofbodyhook

\begin{sphinxfigure-in-table}
\centering
\capstart
\noindent\sphinxincludegraphics[scale=0.4]{{cmake-gui-initial}.png}
\sphinxfigcaption{Initial \sphinxcode{\sphinxupquote{cmake\sphinxhyphen{}gui}} screen}\label{\detokenize{Howto_cmake:id4}}\end{sphinxfigure-in-table}\relax
&

\begin{sphinxfigure-in-table}
\centering
\capstart
\noindent\sphinxincludegraphics[scale=0.6]{{cmake-gui-popup}.png}
\sphinxfigcaption{Generator selection in \sphinxcode{\sphinxupquote{cmake\sphinxhyphen{}gui}}}\label{\detokenize{Howto_cmake:id5}}\end{sphinxfigure-in-table}\relax
&

\begin{sphinxfigure-in-table}
\centering
\capstart
\noindent\sphinxincludegraphics[scale=0.4]{{cmake-gui-options}.png}
\sphinxfigcaption{Options screen of \sphinxcode{\sphinxupquote{cmake\sphinxhyphen{}gui}}}\label{\detokenize{Howto_cmake:id6}}\end{sphinxfigure-in-table}\relax
\\
\sphinxbottomrule
\end{tabulary}
\sphinxtableafterendhook\par
\sphinxattableend\end{savenotes}

\sphinxAtStartPar
Again, you start with an empty configuration cache (left image) and need
to start the configuration step.  For the very first configuration in a
folder, you will have a pop\sphinxhyphen{}up dialog (center image) asking to select
the desired build tool and some configuration settings (stick with the
default) and then you get the option screen with all new settings
highlighted in red.  You can modify them (or not) and click on the
“configure” button again until satisfied and click on the “generate”
button to write out the build files. You can exit the GUI from the
“File” menu or hit “ctrl\sphinxhyphen{}q”.


\subsubsection{Setting options}
\label{\detokenize{Howto_cmake:setting-options}}
\sphinxAtStartPar
Options that enable, disable or modify settings are modified by setting
the value of CMake variables. This is done on the command line with the
\sphinxstyleemphasis{\sphinxhyphen{}D} flag in the format \sphinxcode{\sphinxupquote{\sphinxhyphen{}D VARIABLE=value}}, e.g. \sphinxcode{\sphinxupquote{\sphinxhyphen{}D
CMAKE\_BUILD\_TYPE=Release}} or \sphinxcode{\sphinxupquote{\sphinxhyphen{}D BUILD\_MPI=on}}.  There is one quirk:
when used before the CMake directory, there may be a space between the
\sphinxstyleemphasis{\sphinxhyphen{}D} flag and the variable, after it must not be. Such CMake variables
can have boolean values (on/off, yes/no, or 1/0 are all valid) or are
strings representing a choice, or a path, or are free format. If the
string would contain whitespace, it must be put in quotes, for example
\sphinxcode{\sphinxupquote{\sphinxhyphen{}D CMAKE\_TUNE\_FLAGS="\sphinxhyphen{}ftree\sphinxhyphen{}vectorize \sphinxhyphen{}ffast\sphinxhyphen{}math"}}.

\sphinxAtStartPar
CMake variables fall into two categories: 1) common CMake variables that
are used by default for any CMake configuration setup and 2) project
specific variables, i.e. settings that are specific for LAMMPS.
Also CMake variables can be flagged as \sphinxstyleemphasis{advanced}, which means they are
not shown in the text mode or graphical CMake program in the overview
of all settings by default, but only when explicitly requested (by hitting
the ‘t’ key or clicking on the ‘Advanced’ check\sphinxhyphen{}box).


\paragraph{Some common CMake variables}
\label{\detokenize{Howto_cmake:some-common-cmake-variables}}

\begin{savenotes}\sphinxattablestart
\sphinxthistablewithglobalstyle
\centering
\begin{tabulary}{\linewidth}[t]{TT}
\sphinxtoprule
\sphinxstyletheadfamily 
\sphinxAtStartPar
Variable
&\sphinxstyletheadfamily 
\sphinxAtStartPar
Description
\\
\sphinxmidrule
\sphinxtableatstartofbodyhook
\sphinxAtStartPar
\sphinxcode{\sphinxupquote{CMAKE\_INSTALL\_PREFIX}}
&
\sphinxAtStartPar
root directory of install location for \sphinxcode{\sphinxupquote{make install}}  (default: \sphinxcode{\sphinxupquote{\$HOME/.local}})
\\
\sphinxhline
\sphinxAtStartPar
\sphinxcode{\sphinxupquote{LAMMPS\_INSTALL\_RPATH}}
&
\sphinxAtStartPar
set or remove runtime path setting from binaries for \sphinxcode{\sphinxupquote{make install}} (default: \sphinxcode{\sphinxupquote{off}})
\\
\sphinxhline
\sphinxAtStartPar
\sphinxcode{\sphinxupquote{CMAKE\_BUILD\_TYPE}}
&
\sphinxAtStartPar
controls compilation options:
one of \sphinxcode{\sphinxupquote{RelWithDebInfo}} (default), \sphinxcode{\sphinxupquote{Release}}, \sphinxcode{\sphinxupquote{Debug}}, \sphinxcode{\sphinxupquote{MinSizeRel}}
\\
\sphinxhline
\sphinxAtStartPar
\sphinxcode{\sphinxupquote{BUILD\_SHARED\_LIBS}}
&
\sphinxAtStartPar
if set to \sphinxcode{\sphinxupquote{on}} build the LAMMPS library as shared library (default: \sphinxcode{\sphinxupquote{off}})
\\
\sphinxhline
\sphinxAtStartPar
\sphinxcode{\sphinxupquote{CMAKE\_MAKE\_PROGRAM}}
&
\sphinxAtStartPar
name/path of the compilation command (default depends on \sphinxstyleemphasis{\sphinxhyphen{}G} option, usually \sphinxcode{\sphinxupquote{make}})
\\
\sphinxhline
\sphinxAtStartPar
\sphinxcode{\sphinxupquote{CMAKE\_VERBOSE\_MAKEFILE}}
&
\sphinxAtStartPar
if set to \sphinxcode{\sphinxupquote{on}} echo commands while executing during build (default: \sphinxcode{\sphinxupquote{off}})
\\
\sphinxhline
\sphinxAtStartPar
\sphinxcode{\sphinxupquote{CMAKE\_C\_COMPILER}}
&
\sphinxAtStartPar
C compiler to be used for compilation (default: system specific, \sphinxcode{\sphinxupquote{gcc}} on Linux)
\\
\sphinxhline
\sphinxAtStartPar
\sphinxcode{\sphinxupquote{CMAKE\_CXX\_COMPILER}}
&
\sphinxAtStartPar
C++ compiler to be used for compilation (default: system specific, \sphinxcode{\sphinxupquote{g++}} on Linux)
\\
\sphinxhline
\sphinxAtStartPar
\sphinxcode{\sphinxupquote{CMAKE\_Fortran\_COMPILER}}
&
\sphinxAtStartPar
Fortran compiler to be used for compilation (default: system specific, \sphinxcode{\sphinxupquote{gfortran}} on Linux)
\\
\sphinxhline
\sphinxAtStartPar
\sphinxcode{\sphinxupquote{CXX\_COMPILER\_LAUNCHER}}
&
\sphinxAtStartPar
tool to launch the C++ compiler, e.g. \sphinxcode{\sphinxupquote{ccache}} or \sphinxcode{\sphinxupquote{distcc}} for faster compilation (default: empty)
\\
\sphinxbottomrule
\end{tabulary}
\sphinxtableafterendhook\par
\sphinxattableend\end{savenotes}


\paragraph{Some common LAMMPS specific variables}
\label{\detokenize{Howto_cmake:some-common-lammps-specific-variables}}

\begin{savenotes}\sphinxattablestart
\sphinxthistablewithglobalstyle
\centering
\begin{tabulary}{\linewidth}[t]{TT}
\sphinxtoprule
\sphinxstyletheadfamily 
\sphinxAtStartPar
Variable
&\sphinxstyletheadfamily 
\sphinxAtStartPar
Description
\\
\sphinxmidrule
\sphinxtableatstartofbodyhook
\sphinxAtStartPar
\sphinxcode{\sphinxupquote{BUILD\_MPI}}
&
\sphinxAtStartPar
build LAMMPS with MPI support (default: \sphinxcode{\sphinxupquote{on}} if a working MPI available, else \sphinxcode{\sphinxupquote{off}})
\\
\sphinxhline
\sphinxAtStartPar
\sphinxcode{\sphinxupquote{BUILD\_OMP}}
&
\sphinxAtStartPar
build LAMMPS with OpenMP support (default: \sphinxcode{\sphinxupquote{on}} if compiler supports OpenMP fully, else \sphinxcode{\sphinxupquote{off}})
\\
\sphinxhline
\sphinxAtStartPar
\sphinxcode{\sphinxupquote{BUILD\_TOOLS}}
&
\sphinxAtStartPar
compile some additional executables from the \sphinxcode{\sphinxupquote{tools}} folder (default: \sphinxcode{\sphinxupquote{off}})
\\
\sphinxhline
\sphinxAtStartPar
\sphinxcode{\sphinxupquote{BUILD\_DOC}}
&
\sphinxAtStartPar
include building the HTML format documentation for packaging/installing (default: \sphinxcode{\sphinxupquote{off}})
\\
\sphinxhline
\sphinxAtStartPar
\sphinxcode{\sphinxupquote{CMAKE\_TUNE\_FLAGS}}
&
\sphinxAtStartPar
common compiler flags, for optimization or instrumentation (default:)
\\
\sphinxhline
\sphinxAtStartPar
\sphinxcode{\sphinxupquote{LAMMPS\_MACHINE}}
&
\sphinxAtStartPar
when set to \sphinxcode{\sphinxupquote{name}} the LAMMPS executable and library will be called \sphinxcode{\sphinxupquote{lmp\_name}} and \sphinxcode{\sphinxupquote{liblammps\_name.a}}
\\
\sphinxhline
\sphinxAtStartPar
\sphinxcode{\sphinxupquote{FFT}}
&
\sphinxAtStartPar
select which FFT library to use: \sphinxcode{\sphinxupquote{FFTW3}}, \sphinxcode{\sphinxupquote{MKL}}, \sphinxcode{\sphinxupquote{KISS}} (default, unless FFTW3 is found)
\\
\sphinxhline
\sphinxAtStartPar
\sphinxcode{\sphinxupquote{FFT\_KOKKOS}}
&
\sphinxAtStartPar
select which FFT library to use in Kokkos\sphinxhyphen{}enabled styles: \sphinxcode{\sphinxupquote{FFTW3}}, \sphinxcode{\sphinxupquote{MKL}}, \sphinxcode{\sphinxupquote{HIPFFT}}, \sphinxcode{\sphinxupquote{CUFFT}}, \sphinxcode{\sphinxupquote{KISS}} (default)
\\
\sphinxhline
\sphinxAtStartPar
\sphinxcode{\sphinxupquote{FFT\_SINGLE}}
&
\sphinxAtStartPar
select whether to use single precision FFTs (default: \sphinxcode{\sphinxupquote{off}})
\\
\sphinxhline
\sphinxAtStartPar
\sphinxcode{\sphinxupquote{WITH\_JPEG}}
&
\sphinxAtStartPar
whether to support JPEG format in {\hyperref[\detokenize{dump_image::doc}]{\sphinxcrossref{\DUrole{doc}{dump image}}}} (default: \sphinxcode{\sphinxupquote{on}} if found)
\\
\sphinxhline
\sphinxAtStartPar
\sphinxcode{\sphinxupquote{WITH\_PNG}}
&
\sphinxAtStartPar
whether to support PNG format in  {\hyperref[\detokenize{dump_image::doc}]{\sphinxcrossref{\DUrole{doc}{dump image}}}} (default: \sphinxcode{\sphinxupquote{on}} if found)
\\
\sphinxhline
\sphinxAtStartPar
\sphinxcode{\sphinxupquote{WITH\_GZIP}}
&
\sphinxAtStartPar
whether to support reading and writing compressed files (default: \sphinxcode{\sphinxupquote{on}} if found)
\\
\sphinxhline
\sphinxAtStartPar
\sphinxcode{\sphinxupquote{WITH\_FFMPEG}}
&
\sphinxAtStartPar
whether to support generating movies with {\hyperref[\detokenize{dump_image::doc}]{\sphinxcrossref{\DUrole{doc}{dump movie}}}} (default: \sphinxcode{\sphinxupquote{on}} if found)
\\
\sphinxbottomrule
\end{tabulary}
\sphinxtableafterendhook\par
\sphinxattableend\end{savenotes}


\paragraph{Enabling or disabling LAMMPS packages}
\label{\detokenize{Howto_cmake:enabling-or-disabling-lammps-packages}}
\sphinxAtStartPar
The LAMMPS software is organized into a common core that is always
included and a large number of {\hyperref[\detokenize{Packages::doc}]{\sphinxcrossref{\DUrole{doc}{add\sphinxhyphen{}on packages}}}} that
have to be enabled to be included into a LAMMPS executable.  Packages
are enabled through setting variables of the kind \sphinxcode{\sphinxupquote{PKG\_\textless{}NAME\textgreater{}}} to
\sphinxcode{\sphinxupquote{on}} and disabled by setting them to \sphinxcode{\sphinxupquote{off}} (or using \sphinxcode{\sphinxupquote{yes}},
\sphinxcode{\sphinxupquote{no}}, \sphinxcode{\sphinxupquote{1}}, \sphinxcode{\sphinxupquote{0}} correspondingly).  \sphinxcode{\sphinxupquote{\textless{}NAME\textgreater{}}} has to be replaced by
the name of the package, e.g. \sphinxcode{\sphinxupquote{MOLECULE}} or \sphinxcode{\sphinxupquote{EXTRA\sphinxhyphen{}PAIR}}.


\subsubsection{Using presets}
\label{\detokenize{Howto_cmake:using-presets}}
\sphinxAtStartPar
Since LAMMPS has a lot of optional features and packages, specifying
them all on the command line can be tedious. Or when selecting a
different compiler toolchain, multiple options have to be changed
consistently and that is rather error prone. Or when enabling certain
packages, they require consistent settings to be operated in a
particular mode.  For this purpose, we are providing a selection of
“preset files” for CMake in the folder \sphinxcode{\sphinxupquote{cmake/presets}}.  They
represent a way to pre\sphinxhyphen{}load or override the CMake configuration cache by
setting or changing CMake variables.  Preset files are loaded using the
\sphinxstyleemphasis{\sphinxhyphen{}C} command line flag. You can combine loading multiple preset files or
change some variables later with additional \sphinxstyleemphasis{\sphinxhyphen{}D} flags.  A few examples:

\begin{sphinxVerbatim}[commandchars=\\\{\}]
cmake\PYG{+w}{ }\PYGZhy{}C\PYG{+w}{ }../cmake/presets/basic.cmake\PYG{+w}{ }\PYGZhy{}D\PYG{+w}{ }\PYG{n+nv}{PKG\PYGZus{}MISC}\PYG{o}{=}on\PYG{+w}{ }../cmake
cmake\PYG{+w}{ }\PYGZhy{}C\PYG{+w}{ }../cmake/presets/clang.cmake\PYG{+w}{ }\PYGZhy{}C\PYG{+w}{ }../cmake/presets/most.cmake\PYG{+w}{ }../cmake
cmake\PYG{+w}{ }\PYGZhy{}C\PYG{+w}{ }../cmake/presets/basic.cmake\PYG{+w}{ }\PYGZhy{}D\PYG{+w}{ }\PYG{n+nv}{BUILD\PYGZus{}MPI}\PYG{o}{=}off\PYG{+w}{ }../cmake
\end{sphinxVerbatim}

\sphinxAtStartPar
The first command will install the packages \sphinxcode{\sphinxupquote{KSPACE}}, \sphinxcode{\sphinxupquote{MANYBODY}},
\sphinxcode{\sphinxupquote{MOLECULE}}, \sphinxcode{\sphinxupquote{RIGID}} and \sphinxcode{\sphinxupquote{MISC}}; the first four from the preset
file and the fifth from the explicit variable definition.  The second
command will first switch the compiler toolchain to use the Clang
compilers and install a large number of packages that are not depending
on any special external libraries or tools and are not very unusual.
The third command will enable the first four packages like above and
then enforce compiling LAMMPS as a serial program (using the MPI STUBS
library).

\sphinxAtStartPar
It is also possible to do this incrementally.

\begin{sphinxVerbatim}[commandchars=\\\{\}]
cmake\PYG{+w}{ }\PYGZhy{}C\PYG{+w}{ }../cmake/presets/basic.cmake\PYG{+w}{ }../cmake
cmake\PYG{+w}{ }\PYGZhy{}D\PYG{+w}{ }\PYG{n+nv}{PKG\PYGZus{}MISC}\PYG{o}{=}on\PYG{+w}{ }.
\end{sphinxVerbatim}

\sphinxAtStartPar
will achieve the same final configuration as in the first example above.
In this scenario it is particularly convenient to do the second
configuration step using either the text mode or graphical user
interface (\sphinxcode{\sphinxupquote{ccmake}} or \sphinxcode{\sphinxupquote{cmake\sphinxhyphen{}gui}}).

\begin{sphinxadmonition}{note}{Note:}
\sphinxAtStartPar
Using a preset to select a compiler package (\sphinxcode{\sphinxupquote{clang.cmake}},
\sphinxcode{\sphinxupquote{gcc.cmake}}, \sphinxcode{\sphinxupquote{intel.cmake}}, \sphinxcode{\sphinxupquote{oneapi.cmake}}, or \sphinxcode{\sphinxupquote{pgi.cmake}})
are an exception to the mechanism of updating the configuration
incrementally, as they will trigger a reset of cached internal CMake
settings and thus reset settings to their default values.
\end{sphinxadmonition}


\subsubsection{Compilation and build targets}
\label{\detokenize{Howto_cmake:compilation-and-build-targets}}
\sphinxAtStartPar
The actual compilation will be started by running the selected build
command (on Linux this is by default \sphinxcode{\sphinxupquote{make}}, see below how to select
alternatives).  You can also use the portable command \sphinxcode{\sphinxupquote{cmake \sphinxhyphen{}\sphinxhyphen{}build .}}
which will adapt to whatever the selected build command is.
This is particularly convenient, if you have set a custom build command
via the \sphinxcode{\sphinxupquote{CMAKE\_MAKE\_PROGRAM}} variable.

\sphinxAtStartPar
When calling the build program, you can also select which “target” is to
be build through appending the \sphinxcode{\sphinxupquote{\sphinxhyphen{}\sphinxhyphen{}target}} flag and the name of the target
to the build command.  When using \sphinxcode{\sphinxupquote{make}} as build tool, you can just append
the target name to the command. Example: \sphinxcode{\sphinxupquote{cmake \sphinxhyphen{}\sphinxhyphen{}build . \sphinxhyphen{}\sphinxhyphen{}target all}} or
\sphinxcode{\sphinxupquote{make all}}.  The following abstract targets are available:


\begin{savenotes}\sphinxattablestart
\sphinxthistablewithglobalstyle
\centering
\begin{tabulary}{\linewidth}[t]{TT}
\sphinxtoprule
\sphinxstyletheadfamily 
\sphinxAtStartPar
Target
&\sphinxstyletheadfamily 
\sphinxAtStartPar
Description
\\
\sphinxmidrule
\sphinxtableatstartofbodyhook
\sphinxAtStartPar
\sphinxcode{\sphinxupquote{all}}
&
\sphinxAtStartPar
build “everything” (default)
\\
\sphinxhline
\sphinxAtStartPar
\sphinxcode{\sphinxupquote{lammps}}
&
\sphinxAtStartPar
build the LAMMPS library and executable
\\
\sphinxhline
\sphinxAtStartPar
\sphinxcode{\sphinxupquote{doc}}
&
\sphinxAtStartPar
build the HTML documentation (if configured)
\\
\sphinxhline
\sphinxAtStartPar
\sphinxcode{\sphinxupquote{install}}
&
\sphinxAtStartPar
install all target files into folders in \sphinxcode{\sphinxupquote{CMAKE\_INSTALL\_PREFIX}}
\\
\sphinxhline
\sphinxAtStartPar
\sphinxcode{\sphinxupquote{test}}
&
\sphinxAtStartPar
run some tests (if configured with \sphinxcode{\sphinxupquote{\sphinxhyphen{}D ENABLE\_TESTING=on}})
\\
\sphinxhline
\sphinxAtStartPar
\sphinxcode{\sphinxupquote{clean}}
&
\sphinxAtStartPar
remove all generated files
\\
\sphinxbottomrule
\end{tabulary}
\sphinxtableafterendhook\par
\sphinxattableend\end{savenotes}


\subsubsection{Choosing generators}
\label{\detokenize{Howto_cmake:choosing-generators}}
\sphinxAtStartPar
While CMake usually defaults to creating makefiles to compile software
with the \sphinxcode{\sphinxupquote{make}} program, it supports multiple alternate build tools
(e.g. \sphinxcode{\sphinxupquote{ninja\sphinxhyphen{}build}} which tends to be faster and more efficient in
parallelizing builds than \sphinxcode{\sphinxupquote{make}}) and can generate project files for
integrated development environments (IDEs) like VisualStudio, Eclipse or
CodeBlocks.  This is specific to how the local CMake version was
configured and compiled. The list of available options can be seen at
the end of the output of \sphinxcode{\sphinxupquote{cmake \sphinxhyphen{}\sphinxhyphen{}help}}. Example on Fedora 31 this is:

\begin{sphinxVerbatim}[commandchars=\\\{\}]
Generators

The following generators are available on this platform (* marks default):
* Unix Makefiles               = Generates standard UNIX makefiles.
  Green Hills MULTI            = Generates Green Hills MULTI files
                                 (experimental, work\PYGZhy{}in\PYGZhy{}progress).
  Ninja                        = Generates build.ninja files.
  Ninja Multi\PYGZhy{}Config           = Generates build\PYGZhy{}\PYGZlt{}Config\PYGZgt{}.ninja files.
  Watcom WMake                 = Generates Watcom WMake makefiles.
  CodeBlocks \PYGZhy{} Ninja           = Generates CodeBlocks project files.
  CodeBlocks \PYGZhy{} Unix Makefiles  = Generates CodeBlocks project files.
  CodeLite \PYGZhy{} Ninja             = Generates CodeLite project files.
  CodeLite \PYGZhy{} Unix Makefiles    = Generates CodeLite project files.
  Sublime Text 2 \PYGZhy{} Ninja       = Generates Sublime Text 2 project files.
  Sublime Text 2 \PYGZhy{} Unix Makefiles
                               = Generates Sublime Text 2 project files.
  Kate \PYGZhy{} Ninja                 = Generates Kate project files.
  Kate \PYGZhy{} Unix Makefiles        = Generates Kate project files.
  Eclipse CDT4 \PYGZhy{} Ninja         = Generates Eclipse CDT 4.0 project files.
  Eclipse CDT4 \PYGZhy{} Unix Makefiles= Generates Eclipse CDT 4.0 project files.
\end{sphinxVerbatim}

\sphinxAtStartPar
Below is a screenshot of using the CodeBlocks IDE with the ninja build tool
after running CMake as follows:

\begin{sphinxVerbatim}[commandchars=\\\{\}]
cmake\PYG{+w}{ }\PYGZhy{}G\PYG{+w}{ }\PYG{l+s+s1}{\PYGZsq{}CodeBlocks \PYGZhy{} Ninja\PYGZsq{}}\PYG{+w}{ }../cmake/presets/most.cmake\PYG{+w}{ }../cmake/
\end{sphinxVerbatim}

\noindent{\hspace*{\fill}\sphinxincludegraphics{{cmake-codeblocks}.png}\hspace*{\fill}}

\sphinxstepscope


\subsection{LAMMPS GitHub tutorial}
\label{\detokenize{Howto_github:lammps-github-tutorial}}\label{\detokenize{Howto_github::doc}}
\sphinxAtStartPar
\sphinxstylestrong{written by Stefan Paquay}


\bigskip\hrule\bigskip


\sphinxAtStartPar
This document describes the process of how to use GitHub to integrate
changes or additions you have made to LAMMPS into the official LAMMPS
distribution.  It uses the process of updating this very tutorial as an
example to describe the individual steps and options.  You need to be
familiar with git and you may want to have a look at the \sphinxhref{https://git-scm.com/book/}{git book} to familiarize yourself with some of the
more advanced git features used below.

\sphinxAtStartPar
As of fall 2016, submitting contributions to LAMMPS via pull requests
on GitHub is the preferred option for integrating contributed features
or improvements to LAMMPS, as it significantly reduces the amount of
work required by the LAMMPS developers. Consequently, creating a pull
request will increase your chances to have your contribution included
and will reduce the time until the integration is complete. For more
information on the requirements to have your code included into LAMMPS
please see {\hyperref[\detokenize{Modify_contribute::doc}]{\sphinxcrossref{\DUrole{doc}{this page}}}}.


\bigskip\hrule\bigskip


\sphinxAtStartPar
\sphinxstylestrong{Making an account}

\sphinxAtStartPar
First of all, you need a GitHub account. This is fairly simple, just
go to \sphinxhref{https://github.com}{GitHub} and create an account by clicking
the “Sign up for GitHub” button. Once your account is created, you
can sign in by clicking the button in the top left and filling in your
username or e\sphinxhyphen{}mail address and password.


\bigskip\hrule\bigskip


\sphinxAtStartPar
\sphinxstylestrong{Forking the repository}

\sphinxAtStartPar
To get changes into LAMMPS, you need to first fork the \sphinxtitleref{lammps/lammps}
repository on GitHub. At the time of writing, \sphinxstyleemphasis{develop} is the preferred
target branch. Thus go to \sphinxhref{https://github.com/lammps/lammps}{LAMMPS on GitHub}
and make sure branch is set to “develop”, as shown in the figure below.

\noindent{\hspace*{\fill}\sphinxincludegraphics{{tutorial_branch}.png}\hspace*{\fill}}

\sphinxAtStartPar
If it is not, use the button to change it to \sphinxstyleemphasis{develop}. Once it is, use
the fork button to create a fork.

\noindent{\hspace*{\fill}\sphinxincludegraphics{{tutorial_fork}.png}\hspace*{\fill}}

\sphinxAtStartPar
This will create a fork (which is essentially a copy, but uses less
resources) of the LAMMPS repository under your own GitHub account. You
can make changes in this fork and later file \sphinxstyleemphasis{pull requests} to allow
the upstream repository to merge changes from your own fork into the one
we just forked from (or others that were forked from the same repository).
At the same time, you can set things up, so you can include changes from
upstream into your repository and thus keep it in sync with the ongoing
LAMMPS development.


\bigskip\hrule\bigskip


\sphinxAtStartPar
\sphinxstylestrong{Adding changes to your own fork}

\sphinxAtStartPar
Additions to the upstream version of LAMMPS are handled using \sphinxstyleemphasis{feature
branches}.  For every new feature, a so\sphinxhyphen{}called feature branch is
created, which contains only those modification relevant to one specific
feature. For example, adding a single fix would consist of creating a
branch with only the fix header and source file and nothing else.  It is
explained in more detail here: \sphinxhref{https://www.atlassian.com/git/tutorials/comparing-workflows/feature-branch-workflow}{feature branch workflow}.

\sphinxAtStartPar
\sphinxstylestrong{Feature branches}

\sphinxAtStartPar
First of all, create a clone of your version on GitHub on your local
machine via HTTPS:

\begin{sphinxVerbatim}[commandchars=\\\{\}]
git\PYG{+w}{ }clone\PYG{+w}{ }https://github.com/\PYGZlt{}your\PYG{+w}{ }user\PYG{+w}{ }name\PYGZgt{}/lammps.git\PYG{+w}{ }\PYGZlt{}some\PYG{+w}{ }name\PYGZgt{}
\end{sphinxVerbatim}

\sphinxAtStartPar
or, if you have set up your GitHub account for using SSH keys, via SSH:

\begin{sphinxVerbatim}[commandchars=\\\{\}]
git\PYG{+w}{ }clone\PYG{+w}{ }git@github.com:\PYGZlt{}your\PYG{+w}{ }user\PYG{+w}{ }name\PYGZgt{}/lammps.git
\end{sphinxVerbatim}

\sphinxAtStartPar
You can find the proper URL by clicking the “Clone or download”\sphinxhyphen{}button:

\noindent{\hspace*{\fill}\sphinxincludegraphics{{tutorial_https_block}.png}\hspace*{\fill}}

\sphinxAtStartPar
The above command copies (“clones”) the git repository to your local
machine to a directory with the name you chose. If none is given, it will
default to “lammps”. Typical names are “mylammps” or something similar.

\sphinxAtStartPar
You can use this local clone to make changes and test them without
interfering with the repository on GitHub.

\sphinxAtStartPar
To pull changes from upstream into this copy, you can go to the directory
and use git pull:

\begin{sphinxVerbatim}[commandchars=\\\{\}]
\PYG{n+nb}{cd}\PYG{+w}{ }mylammps
git\PYG{+w}{ }checkout\PYG{+w}{ }develop
git\PYG{+w}{ }pull\PYG{+w}{ }https://github.com/lammps/lammps\PYG{+w}{ }develop
\end{sphinxVerbatim}

\sphinxAtStartPar
You can also add this URL as a remote:

\begin{sphinxVerbatim}[commandchars=\\\{\}]
git\PYG{+w}{ }remote\PYG{+w}{ }add\PYG{+w}{ }upstream\PYG{+w}{ }https://www.github.com/lammps/lammps
\end{sphinxVerbatim}

\sphinxAtStartPar
From then on you can update your upstream branches with:

\begin{sphinxVerbatim}[commandchars=\\\{\}]
git\PYG{+w}{ }fetch\PYG{+w}{ }upstream
\end{sphinxVerbatim}

\sphinxAtStartPar
and then refer to the upstream repository branches with
\sphinxtitleref{upstream/develop} or \sphinxtitleref{upstream/release} and so on.

\sphinxAtStartPar
At this point, you typically make a feature branch from the updated
branch for the feature you want to work on. This tutorial contains the
workflow that updated this tutorial, and hence we will call the branch
“github\sphinxhyphen{}tutorial\sphinxhyphen{}update”:

\begin{sphinxVerbatim}[commandchars=\\\{\}]
git\PYG{+w}{ }fetch\PYG{+w}{ }upstream
git\PYG{+w}{ }checkout\PYG{+w}{ }\PYGZhy{}b\PYG{+w}{ }github\PYGZhy{}tutorial\PYGZhy{}update\PYG{+w}{ }upstream/develop
\end{sphinxVerbatim}

\sphinxAtStartPar
Now that we have changed branches, we can make our changes to our local
repository. Just remember that if you want to start working on another,
unrelated feature, you should switch branches!

\begin{sphinxadmonition}{note}{Note:}
\sphinxAtStartPar
Committing changes to the \sphinxstyleemphasis{develop}, \sphinxstyleemphasis{release}, or \sphinxstyleemphasis{stable} branches
is strongly discouraged.  While it may be convenient initially, it
will create more work in the long run.  Various texts and tutorials
on using git effectively discuss the motivation for using feature
branches instead.
\end{sphinxadmonition}

\sphinxAtStartPar
\sphinxstylestrong{After changes are made}

\sphinxAtStartPar
After everything is done, add the files to the branch and commit them:

\begin{sphinxVerbatim}[commandchars=\\\{\}]
git\PYG{+w}{ }add\PYG{+w}{ }doc/src/Howto\PYGZus{}github.txt
git\PYG{+w}{ }add\PYG{+w}{ }doc/src/JPG/tutorial*.png
\end{sphinxVerbatim}

\begin{sphinxadmonition}{warning}{Warning:}
\sphinxAtStartPar
Do not use \sphinxstyleemphasis{git commit \sphinxhyphen{}a} (or \sphinxstyleemphasis{git add \sphinxhyphen{}A}).  The \sphinxhyphen{}a flag (or \sphinxhyphen{}A
flag) will automatically include \sphinxstylestrong{all} modified \sphinxstylestrong{and} new files
and that is rarely the behavior you want.  It can easily lead to
accidentally adding unrelated and unwanted changes into the
repository.  Instead it is preferable to explicitly use \sphinxstyleemphasis{git add},
\sphinxstyleemphasis{git rm}, \sphinxstyleemphasis{git mv} for adding, removing, renaming individual files,
respectively, and then \sphinxstyleemphasis{git commit} to finalize the commit.
Carefully check all pending changes with \sphinxstyleemphasis{git status} before
committing them.  If you find doing this on the command line too
tedious, consider using a GUI, for example the one included in git
distributions written in Tk, i.e. use \sphinxstyleemphasis{git gui} (on some Linux
distributions it may be required to install an additional package to
use it).
\end{sphinxadmonition}

\sphinxAtStartPar
After adding all files, the change set can be committed with some
useful message that explains the change.

\begin{sphinxVerbatim}[commandchars=\\\{\}]
git\PYG{+w}{ }commit\PYG{+w}{ }\PYGZhy{}m\PYG{+w}{ }\PYG{l+s+s1}{\PYGZsq{}Finally updated the GitHub tutorial\PYGZsq{}}
\end{sphinxVerbatim}

\sphinxAtStartPar
After the commit, the changes can be pushed to the same branch on GitHub:

\begin{sphinxVerbatim}[commandchars=\\\{\}]
git\PYG{+w}{ }push
\end{sphinxVerbatim}

\sphinxAtStartPar
Git will ask you for your user name and password on GitHub if you have
not configured anything. If your local branch is not present on GitHub yet,
it will ask you to add it by running

\begin{sphinxVerbatim}[commandchars=\\\{\}]
git\PYG{+w}{ }push\PYG{+w}{ }\PYGZhy{}\PYGZhy{}set\PYGZhy{}upstream\PYG{+w}{ }origin\PYG{+w}{ }github\PYGZhy{}tutorial\PYGZhy{}update
\end{sphinxVerbatim}

\sphinxAtStartPar
If you correctly type your user name and
password, the feature branch should be added to your fork on GitHub.

\sphinxAtStartPar
If you want to make really sure you push to the right repository
(which is good practice), you can provide it explicitly:

\begin{sphinxVerbatim}[commandchars=\\\{\}]
git\PYG{+w}{ }push\PYG{+w}{ }origin
\end{sphinxVerbatim}

\sphinxAtStartPar
or using an explicit URL:

\begin{sphinxVerbatim}[commandchars=\\\{\}]
git\PYG{+w}{ }push\PYG{+w}{ }git@github.com:Pakketeretet2/lammps.git
\end{sphinxVerbatim}


\bigskip\hrule\bigskip


\sphinxAtStartPar
\sphinxstylestrong{Filing a pull request}

\sphinxAtStartPar
Up to this point in the tutorial, all changes were to \sphinxstyleemphasis{your} clones of
LAMMPS.  Eventually, however, you want this feature to be included into
the official LAMMPS version.  To do this, you will want to file a pull
request by clicking on the “New pull request” button:

\noindent{\hspace*{\fill}\sphinxincludegraphics{{tutorial_new_pull_request}.png}\hspace*{\fill}}

\sphinxAtStartPar
Make sure that the current branch is set to the correct one, which, in
this case, is “github\sphinxhyphen{}tutorial\sphinxhyphen{}update”. If done correctly, the only
changes you will see are those that were made on this branch.

\sphinxAtStartPar
This will open up a new window that lists changes made to the
repository. If you are just adding new files, there is not much to do,
but I suppose merge conflicts are to be resolved here if there are
changes in existing files. If all changes can automatically be merged,
green text at the top will say so and you can click the “Create pull
request” button, see image.

\noindent{\hspace*{\fill}\sphinxincludegraphics{{tutorial_create_new_pull_request1}.png}\hspace*{\fill}}

\sphinxAtStartPar
Before creating the pull request, make sure the short title is accurate
and add a comment with details about your pull request.  Here you write
what your modifications do and why they should be incorporated upstream.

\sphinxAtStartPar
Note the checkbox that says “Allow edits from maintainers”.
This is checked by default checkbox (although in my version of Firefox, only the checkmark is visible):

\noindent{\hspace*{\fill}\sphinxincludegraphics{{tutorial_edits_maintainers}.png}\hspace*{\fill}}

\sphinxAtStartPar
If it is checked, maintainers can immediately add their own edits to the
pull request.  This helps the inclusion of your branch significantly, as
simple/trivial changes can be added directly to your pull request branch
by the LAMMPS maintainers.  The alternative would be that they make
changes on their own version of the branch and file a reverse pull
request to you.  Just leave this box checked unless you have a very good
reason not to.

\sphinxAtStartPar
Now just write some nice comments and click on “Create pull request”.

\noindent{\hspace*{\fill}\sphinxincludegraphics{{tutorial_create_new_pull_request2}.png}\hspace*{\fill}}


\bigskip\hrule\bigskip


\sphinxAtStartPar
\sphinxstylestrong{After filing a pull request}

\begin{sphinxadmonition}{note}{Note:}
\sphinxAtStartPar
When you submit a pull request (or ask for a pull request) for the
first time, you will receive an invitation to become a LAMMPS project
collaborator. Please accept this invite as being a collaborator will
simplify certain administrative tasks and will probably speed up the
merging of your feature, too.
\end{sphinxadmonition}

\sphinxAtStartPar
You will notice that after filing the pull request, some checks are
performed automatically:

\noindent{\hspace*{\fill}\sphinxincludegraphics{{tutorial_automated_checks}.png}\hspace*{\fill}}

\sphinxAtStartPar
If all is fine, you will see this:

\noindent{\hspace*{\fill}\sphinxincludegraphics{{tutorial_automated_checks_passed}.png}\hspace*{\fill}}

\sphinxAtStartPar
If any of the checks are failing, your pull request will not be
processed, as your changes may break compilation for certain
configurations or may not merge cleanly. It is your responsibility
to remove the reason(s) for the failed test(s). If you need help
with this, please contact the LAMMPS developers by adding a comment
explaining your problems with resolving the failed tests.

\sphinxAtStartPar
A few further interesting things (can) happen to pull requests before
they are included.

\sphinxAtStartPar
\sphinxstylestrong{Additional changes}

\sphinxAtStartPar
First of all, any additional changes you push into your branch in your
repository will automatically become part of the pull request:

\noindent{\hspace*{\fill}\sphinxincludegraphics{{tutorial_additional_changes}.png}\hspace*{\fill}}

\sphinxAtStartPar
This means you can add changes that should be part of the feature after
filing the pull request, which is useful in case you have forgotten
them, or if a developer has requested that something needs to be changed
before the feature can be accepted into the official LAMMPS version.
After each push, the automated checks are run again.

\sphinxAtStartPar
\sphinxstylestrong{Labels}

\sphinxAtStartPar
LAMMPS developers may add labels to your pull request to assign it to
categories (mostly for bookkeeping purposes), but a few of them are
important: \sphinxstyleemphasis{needs\_work}, \sphinxstyleemphasis{work\_in\_progress}, \sphinxstyleemphasis{run\_tests},
\sphinxstyleemphasis{test\_for\_regression}, and \sphinxstyleemphasis{ready\_for\_merge}.  The first two indicate,
that your pull request is not considered to be complete. With
“needs\_work” the burden is on exclusively on you; while
“work\_in\_progress” can also mean, that a LAMMPS developer may want to
add changes. Please watch the comments to the pull requests. The two
“test” labels are used to trigger extended tests before the code is
merged. This is sometimes done by LAMMPS developers, if they suspect
that there may be some subtle side effects from your changes. It is not
done by default, because those tests are very time\sphinxhyphen{}consuming.  The
\sphinxstyleemphasis{ready\_for\_merge} label is usually attached when the LAMMPS developer
assigned to the pull request considers this request complete and to
trigger a final full test evaluation.

\sphinxAtStartPar
\sphinxstylestrong{Reviews}

\sphinxAtStartPar
As of Fall 2021, a pull request needs to pass all automatic tests and at
least 1 approving review from a LAMMPS developer with write access to
the repository before it is eligible for merging.  In case your changes
touch code that certain developers are associated with, they are
auto\sphinxhyphen{}requested by the GitHub software.  Those associations are set in
the file \sphinxhref{https://github.com/lammps/lammps/blob/develop/.github/CODEOWNERS}{.github/CODEOWNERS} Thus
if you want to be automatically notified to review when anybody changes
files or packages, that \sphinxstylestrong{you} have contributed to LAMMPS, you can add
suitable patterns to that file, or a LAMMPS developer may add you.

\sphinxAtStartPar
Otherwise, you can also manually request reviews from specific developers,
or LAMMPS developers \sphinxhyphen{} in their assessment of your pull request \sphinxhyphen{} may
determine who else should be reviewing your contribution and add that person.
Through reviews, LAMMPS developers also may request specific changes from you.
If those are not addressed, your pull requests cannot be merged.

\sphinxAtStartPar
\sphinxstylestrong{Assignees}

\sphinxAtStartPar
There is an assignee property for pull requests. If the request has not
been reviewed by any developer yet, it is not assigned to anyone. After
revision, a developer can choose to assign it to either a) you, b) a
LAMMPS developer (including him/herself) or c) Axel Kohlmeyer (akohlmey).
\begin{itemize}
\item {} 
\sphinxAtStartPar
Case a) happens if changes are required on your part

\item {} 
\sphinxAtStartPar
Case b) means that at the moment, it is being tested and reviewed by a
LAMMPS developer with the expectation that some changes would be required.
After the review, the developer can choose to implement changes directly
or suggest them to you.

\item {} 
\sphinxAtStartPar
Case c) means that the pull request has been assigned to the developer
overseeing the merging of pull requests into the \sphinxstyleemphasis{develop} branch.

\end{itemize}

\sphinxAtStartPar
In this case, Axel assigned the tutorial to Steve:

\noindent{\hspace*{\fill}\sphinxincludegraphics{{tutorial_steve_assignee}.png}\hspace*{\fill}}

\sphinxAtStartPar
\sphinxstylestrong{Edits from LAMMPS maintainers}

\sphinxAtStartPar
If you allowed edits from maintainers (the default), any LAMMPS
maintainer can add changes to your pull request.  In this case, both
Axel and Richard made changes to the tutorial:

\noindent{\hspace*{\fill}\sphinxincludegraphics{{tutorial_changes_others}.png}\hspace*{\fill}}

\sphinxAtStartPar
\sphinxstylestrong{Reverse pull requests}

\sphinxAtStartPar
Sometimes, however, you might not feel comfortable having other people
push changes into your own branch, or maybe the maintainers are not sure
their idea was the right one.  In such a case, they can make changes,
reassign you as the assignee, and file a “reverse pull request”, i.e.
file a pull request in \sphinxstylestrong{your} forked GitHub repository to include
changes in the branch, that you have submitted as a pull request
yourself.  In that case, you can choose to merge their changes back into
your branch, possibly make additional changes or corrections and proceed
from there.  It looks something like this:

\noindent{\hspace*{\fill}\sphinxincludegraphics{{tutorial_reverse_pull_request}.png}\hspace*{\fill}}

\sphinxAtStartPar
For some reason, the highlighted button did not work in my case, but I
can go to my own repository and merge the pull request from there:

\noindent{\hspace*{\fill}\sphinxincludegraphics{{tutorial_reverse_pull_request2}.png}\hspace*{\fill}}

\sphinxAtStartPar
Be sure to check the changes to see if you agree with them by clicking
on the tab button:

\noindent{\hspace*{\fill}\sphinxincludegraphics{{tutorial_reverse_pull_request3}.png}\hspace*{\fill}}

\sphinxAtStartPar
In this case, most of it is changes in the markup and a short rewrite of
Axel’s explanation of the “git gui” and “git add” commands.

\noindent{\hspace*{\fill}\sphinxincludegraphics{{tutorial_reverse_pull_request4}.png}\hspace*{\fill}}

\sphinxAtStartPar
Because the changes are OK with us, we are going to merge by clicking on
“Merge pull request”.  After a merge it looks like this:

\noindent{\hspace*{\fill}\sphinxincludegraphics{{tutorial_reverse_pull_request5}.png}\hspace*{\fill}}

\sphinxAtStartPar
Now, since in the meantime our local text for the tutorial also changed,
we need to pull Axel’s change back into our branch, and merge them:

\begin{sphinxVerbatim}[commandchars=\\\{\}]
git\PYG{+w}{ }add\PYG{+w}{ }Howto\PYGZus{}github.txt
git\PYG{+w}{ }add\PYG{+w}{ }JPG/tutorial\PYGZus{}reverse\PYGZus{}pull\PYGZus{}request*.png
git\PYG{+w}{ }commit\PYG{+w}{ }\PYGZhy{}m\PYG{+w}{ }\PYG{l+s+s2}{\PYGZdq{}Updated text and images on reverse pull requests\PYGZdq{}}
git\PYG{+w}{ }pull
\end{sphinxVerbatim}

\sphinxAtStartPar
In this case, the merge was painless because git could auto\sphinxhyphen{}merge:

\noindent{\hspace*{\fill}\sphinxincludegraphics{{tutorial_reverse_pull_request6}.png}\hspace*{\fill}}

\sphinxAtStartPar
With Axel’s changes merged in and some final text updates, our feature
branch is now perfect as far as we are concerned, so we are going to
commit and push again:

\begin{sphinxVerbatim}[commandchars=\\\{\}]
git\PYG{+w}{ }add\PYG{+w}{ }Howto\PYGZus{}github.txt
git\PYG{+w}{ }add\PYG{+w}{ }JPG/tutorial\PYGZus{}reverse\PYGZus{}pull\PYGZus{}request6.png
git\PYG{+w}{ }commit\PYG{+w}{ }\PYGZhy{}m\PYG{+w}{ }\PYG{l+s+s2}{\PYGZdq{}Merged Axel\PYGZsq{}s suggestions and updated text\PYGZdq{}}
git\PYG{+w}{ }push\PYG{+w}{ }git@github.com:Pakketeretet2/lammps
\end{sphinxVerbatim}

\sphinxAtStartPar
This merge also shows up on the lammps GitHub page:

\noindent{\hspace*{\fill}\sphinxincludegraphics{{tutorial_reverse_pull_request7}.png}\hspace*{\fill}}


\bigskip\hrule\bigskip


\sphinxAtStartPar
\sphinxstylestrong{After a merge}

\sphinxAtStartPar
When everything is fine, the feature branch is merged into the \sphinxstyleemphasis{develop} branch:

\noindent{\hspace*{\fill}\sphinxincludegraphics{{tutorial_merged}.png}\hspace*{\fill}}

\sphinxAtStartPar
Now one question remains: What to do with the feature branch that got
merged into upstream?

\sphinxAtStartPar
It is in principle safe to delete them from your own fork. This helps
keep it a bit more tidy. Note that you first have to switch to another
branch!

\begin{sphinxVerbatim}[commandchars=\\\{\}]
git\PYG{+w}{ }checkout\PYG{+w}{ }develop
git\PYG{+w}{ }pull\PYG{+w}{ }https://github.com/lammps/lammps\PYG{+w}{ }develop
git\PYG{+w}{ }branch\PYG{+w}{ }\PYGZhy{}d\PYG{+w}{ }github\PYGZhy{}tutorial\PYGZhy{}update
\end{sphinxVerbatim}

\sphinxAtStartPar
If you do not pull first, it is not really a problem but git will warn
you at the next statement that you are deleting a local branch that
was not yet fully merged into HEAD. This is because git does not yet
know your branch just got merged into LAMMPS upstream. If you
first delete and then pull, everything should still be fine.
You can display all branches that are fully merged by:

\sphinxAtStartPar
Finally, if you delete the branch locally, you might want to push this
to your remote(s) as well:

\begin{sphinxVerbatim}[commandchars=\\\{\}]
git\PYG{+w}{ }push\PYG{+w}{ }origin\PYG{+w}{ }:github\PYGZhy{}tutorial\PYGZhy{}update
\end{sphinxVerbatim}

\sphinxAtStartPar
\sphinxstylestrong{Recent changes in the workflow}

\sphinxAtStartPar
Some recent changes to the workflow are not captured in this tutorial.
For example, in addition to the \sphinxstyleemphasis{develop} branch, to which all new
features should be submitted, there is also a \sphinxstyleemphasis{release}, a \sphinxstyleemphasis{stable}, and
a \sphinxstyleemphasis{maintenance} branch; the \sphinxstyleemphasis{release} branch is updated from the
\sphinxstyleemphasis{develop} branch as part of a “feature release”, and \sphinxstyleemphasis{stable} (together
with \sphinxstyleemphasis{release}) are updated from \sphinxstyleemphasis{develop} when a “stable release” is
made. In between stable releases, selected bug fixes and infrastructure
updates are back\sphinxhyphen{}ported from the \sphinxstyleemphasis{develop} branch to the \sphinxstyleemphasis{maintenance}
branch and occasionally merged to \sphinxstyleemphasis{stable} as an update release.

\sphinxAtStartPar
Furthermore, the naming of the release tags now follow the pattern
“patch\_\textless{}Day\textgreater{}\textless{}Month\textgreater{}\textless{}Year\textgreater{}” to simplify comparisons between releases.
For stable releases additional “stable\_\textless{}Day\textgreater{}\textless{}Month\textgreater{}\textless{}Year\textgreater{}” tags are
applied and update releases are tagged with
“stable\_\textless{}Day\textgreater{}\textless{}Month\textgreater{}\textless{}Year\textgreater{}\_update\textless{}Number\textgreater{}”, Finally, all releases and
submissions are subject to automatic testing and code checks to make
sure they compile with a variety of compilers and popular operating
systems.  Some unit and regression testing is applied as well.

\sphinxAtStartPar
A detailed discussion of the LAMMPS developer GitHub workflow can be
found in the file \sphinxhref{https://github.com/lammps/lammps/blob/develop/doc/github-development-workflow.md}{doc/github\sphinxhyphen{}development\sphinxhyphen{}workflow.md}

\sphinxstepscope


\subsection{Using LAMMPS\sphinxhyphen{}GUI}
\label{\detokenize{Howto_lammps_gui:using-lammps-gui}}\label{\detokenize{Howto_lammps_gui::doc}}
\sphinxAtStartPar
This document describes \sphinxstylestrong{LAMMPS\sphinxhyphen{}GUI version 1.6}.


\bigskip\hrule\bigskip


\sphinxAtStartPar
LAMMPS\sphinxhyphen{}GUI is a graphical text editor customized for editing LAMMPS
input files that is linked to the {\hyperref[\detokenize{Library:lammps-c-api}]{\sphinxcrossref{\DUrole{std,std-ref}{LAMMPS library}}}}
and thus can run LAMMPS directly using the contents of the editor’s text
buffer as input.  It can retrieve and display information from LAMMPS
while it is running, display visualizations created with the {\hyperref[\detokenize{dump_image::doc}]{\sphinxcrossref{\DUrole{doc}{dump
image command}}}}, and is adapted specifically for editing
LAMMPS input files through text completion and reformatting, and linking
to the online LAMMPS documentation for known LAMMPS commands and styles.

\begin{sphinxadmonition}{note}{Note:}
\sphinxAtStartPar
Pre\sphinxhyphen{}compiled, ready\sphinxhyphen{}to\sphinxhyphen{}use LAMMPS\sphinxhyphen{}GUI executables for Linux x86\_64
(Ubuntu 20.04LTS or later and compatible), macOS (version 11 aka Big
Sur or later), and Windows (version 10 or later) {\hyperref[\detokenize{Tools:lammps-gui-install}]{\sphinxcrossref{\DUrole{std,std-ref}{are available}}}} for download.  Non\sphinxhyphen{}MPI LAMMPS executables (as
\sphinxcode{\sphinxupquote{lmp}}) for running LAMMPS from the command line and {\hyperref[\detokenize{Tools::doc}]{\sphinxcrossref{\DUrole{doc}{some
LAMMPS tools}}}} compiled executables are also included.

\sphinxAtStartPar
The source code for LAMMPS\sphinxhyphen{}GUI is included in the LAMMPS source code
distribution and can be found in the \sphinxcode{\sphinxupquote{tools/lammps\sphinxhyphen{}gui}} folder.  It
can be compiled alongside LAMMPS when {\hyperref[\detokenize{Build_cmake::doc}]{\sphinxcrossref{\DUrole{doc}{compiling with CMake}}}}.
\end{sphinxadmonition}

\sphinxAtStartPar
LAMMPS\sphinxhyphen{}GUI tries to provide an experience similar to what people
traditionally would have running LAMMPS using a command line window and
the console LAMMPS executable but just rolled into a single executable:
\begin{itemize}
\item {} 
\sphinxAtStartPar
writing \& editing LAMMPS input files with a text editor

\item {} 
\sphinxAtStartPar
run LAMMPS on those input file with selected command line flags

\item {} 
\sphinxAtStartPar
extract data from the created files and visualize it with and
external software

\end{itemize}

\sphinxAtStartPar
That procedure is quite effective for people proficient in using the
command line, as that allows them to use tools for the individual steps
that they are most comfortable with.  In fact, it is often \sphinxstyleemphasis{required} to
adopt this workflow when running LAMMPS simulations on high\sphinxhyphen{}performance
computing facilities.

\sphinxAtStartPar
The main benefit of using LAMMPS\sphinxhyphen{}GUI is that many basic tasks can be
done directly from the GUI \sphinxstylestrong{without} switching to a text console
window or using external programs, let alone writing scripts to extract
data from the generated output.  It also integrates well with graphical
desktop environments where the \sphinxtitleref{.lmp} filename extension can be
registered with LAMMPS\sphinxhyphen{}GUI as the executable to launch when double
clicking on such files.  Also, LAMMPS\sphinxhyphen{}GUI has support for drag\sphinxhyphen{}n\sphinxhyphen{}drop,
i.e.  an input file can be selected and then moved and dropped on the
LAMMPS\sphinxhyphen{}GUI executable, and LAMMPS\sphinxhyphen{}GUI will launch and read the file into
its buffer.  In many cases LAMMPS\sphinxhyphen{}GUI will be integrated into the
graphical desktop environment and can be launched like other
applications.

\sphinxAtStartPar
LAMMPS\sphinxhyphen{}GUI thus makes it easier for beginners to get started running
simple LAMMPS simulations.  It is very suitable for tutorials on LAMMPS
since you only need to learn how to use a single program for most tasks
and thus time can be saved and people can focus on learning LAMMPS.
The tutorials at \sphinxurl{https://lammpstutorials.github.io/} are specifically
updated for use with LAMMPS\sphinxhyphen{}GUI.

\sphinxAtStartPar
Another design goal is to keep the barrier low when replacing part of
the functionality of LAMMPS\sphinxhyphen{}GUI with external tools.  That said, LAMMPS\sphinxhyphen{}GUI
has some unique functionality that is not found elsewhere:
\begin{itemize}
\item {} 
\sphinxAtStartPar
auto\sphinxhyphen{}adapting to features available in the integrated LAMMPS library

\item {} 
\sphinxAtStartPar
interactive visualization using the {\hyperref[\detokenize{dump_image::doc}]{\sphinxcrossref{\DUrole{doc}{dump image}}}}
command with the option to copy\sphinxhyphen{}paste the resulting settings

\item {} 
\sphinxAtStartPar
automatic slide show generation from dump image out at runtime

\item {} 
\sphinxAtStartPar
automatic plotting of thermodynamics data at runtime

\item {} 
\sphinxAtStartPar
inspection of binary restart files

\end{itemize}

\sphinxAtStartPar
The following text provides a detailed tour of the features and
functionality of LAMMPS\sphinxhyphen{}GUI.  Suggestions for new features and
reports of bugs are always welcome.  You can use the {\hyperref[\detokenize{Errors_bugs::doc}]{\sphinxcrossref{\DUrole{doc}{the same
channels as for LAMMPS itself}}}} for that purpose.


\bigskip\hrule\bigskip



\subsubsection{Installing Pre\sphinxhyphen{}compiled LAMMPS\sphinxhyphen{}GUI Packages}
\label{\detokenize{Howto_lammps_gui:installing-pre-compiled-lammps-gui-packages}}
\sphinxAtStartPar
LAMMPS\sphinxhyphen{}GUI is available as pre\sphinxhyphen{}compiled binary packages for Linux
x86\_64, macOS 11 and later, and Windows 10 and later.  Alternately, it
can be compiled from source.


\paragraph{Windows 10 and later}
\label{\detokenize{Howto_lammps_gui:windows-10-and-later}}
\sphinxAtStartPar
After downloading the \sphinxcode{\sphinxupquote{LAMMPS\sphinxhyphen{}Win10\sphinxhyphen{}64bit\sphinxhyphen{}GUI\sphinxhyphen{}\textless{}version\textgreater{}.exe}} installer
package, you need to execute it, and start the installation process.
Since those packages are currently unsigned, you have to enable “Developer Mode”
in the Windows System Settings to run the installer.


\paragraph{MacOS 11 and later}
\label{\detokenize{Howto_lammps_gui:macos-11-and-later}}
\sphinxAtStartPar
After downloading the \sphinxcode{\sphinxupquote{LAMMPS\sphinxhyphen{}macOS\sphinxhyphen{}multiarch\sphinxhyphen{}GUI\sphinxhyphen{}\textless{}version\textgreater{}.dmg}}
installer package, you need to double\sphinxhyphen{}click it and then, in the window
that opens, drag the app bundle as indicated into the “Applications”
folder.  The follow the instructions in the “README.txt” file to
get access to the other included executables.


\paragraph{Linux on x86\_64}
\label{\detokenize{Howto_lammps_gui:linux-on-x86-64}}
\sphinxAtStartPar
For Linux with x86\_64 CPU there are currently two variants. The first
is compiled on Ubuntu 20.04LTS, is using some wrapper scripts, and
should be compatible with more recent Linux distributions.  After
downloading and unpacking the
\sphinxcode{\sphinxupquote{LAMMPS\sphinxhyphen{}Linux\sphinxhyphen{}x86\_64\sphinxhyphen{}GUI\sphinxhyphen{}\textless{}version\textgreater{}.tar.gz}} package.  You can switch
into the “LAMMPS\_GUI” folder and execute “./lammps\sphinxhyphen{}gui” directly.

\sphinxAtStartPar
The second variant uses \sphinxhref{https://www.flatpak.org}{flatpak} and
requires the flatpak management and runtime software to be installed.
After downloading the \sphinxcode{\sphinxupquote{LAMMPS\sphinxhyphen{}GUI\sphinxhyphen{}Linux\sphinxhyphen{}x86\_64\sphinxhyphen{}GUI\sphinxhyphen{}\textless{}version\textgreater{}.tar.gz}}
flatpak bundle, you can install it with \sphinxcode{\sphinxupquote{flatpak install \sphinxhyphen{}\sphinxhyphen{}user
LAMMPS\sphinxhyphen{}GUI\sphinxhyphen{}Linux\sphinxhyphen{}x86\_64\sphinxhyphen{}GUI\sphinxhyphen{}\textless{}version\textgreater{}.tar.gz}}.  After installation,
LAMMPS\sphinxhyphen{}GUI should be integrated into your desktop environment under
“Applications \textgreater{} Science” but also can be launched from the console with
\sphinxcode{\sphinxupquote{flatpak run org.lammps.lammps\sphinxhyphen{}gui}}.  The flatpak bundle also includes
the console LAMMPS executable \sphinxcode{\sphinxupquote{lmp}} which can be launched to run
simulations with, for example: \sphinxcode{\sphinxupquote{flatpak run \sphinxhyphen{}\sphinxhyphen{}command=lmp
org.lammps.lammps\sphinxhyphen{}gui \sphinxhyphen{}in in.melt}}.


\paragraph{Compiling from Source}
\label{\detokenize{Howto_lammps_gui:compiling-from-source}}
\sphinxAtStartPar
There also are instructions for {\hyperref[\detokenize{Tools:lammps-gui-compilation}]{\sphinxcrossref{\DUrole{std,std-ref}{compiling LAMMPS\sphinxhyphen{}GUI from source
code}}}} available elsewhere in the manual.
Compilation from source \sphinxstyleemphasis{requires} using CMake.


\bigskip\hrule\bigskip



\subsubsection{Starting LAMMPS\sphinxhyphen{}GUI}
\label{\detokenize{Howto_lammps_gui:starting-lammps-gui}}
\sphinxAtStartPar
When LAMMPS\sphinxhyphen{}GUI starts, it shows the main window, labeled \sphinxstyleemphasis{Editor}, with
either an empty buffer or the contents of the file used as argument. In
the latter case it may look like the following:

\sphinxAtStartPar
\sphinxincludegraphics[width=0.480\linewidth]{{lammps-gui-main}.png}  \sphinxincludegraphics[width=0.480\linewidth]{{lammps-gui-dark}.png}

\sphinxAtStartPar
There is the typical menu bar at the top, then the main editor buffer,
and a status bar at the bottom.  The input file contents are shown
with line numbers on the left and the input is colored according to
the LAMMPS input file syntax.  The status bar shows the status of
LAMMPS execution on the left (e.g. “Ready.” when idle) and the current
working directory on the right.  The name of the current file in the
buffer is shown in the window title; the word \sphinxtitleref{*modified*} is added if
the buffer edits have not yet saved to a file.  The geometry of the main
window is stored when exiting and restored when starting again.


\paragraph{Opening Files}
\label{\detokenize{Howto_lammps_gui:opening-files}}
\sphinxAtStartPar
The LAMMPS\sphinxhyphen{}GUI application can be launched without command line arguments
and then starts with an empty buffer in the \sphinxstyleemphasis{Editor} window.  If arguments
are given LAMMPS will use first command line argument as the file name for
the \sphinxstyleemphasis{Editor} buffer and reads its contents into the buffer, if the file
exists.  All further arguments are ignored.  Files can also be opened via
the \sphinxstyleemphasis{File} menu, the \sphinxtitleref{Ctrl\sphinxhyphen{}O} (\sphinxtitleref{Command\sphinxhyphen{}O} on macOS) keyboard shortcut
or by drag\sphinxhyphen{}and\sphinxhyphen{}drop of a file from a graphical file manager into the editor
window.  If a file extension (e.g. \sphinxcode{\sphinxupquote{.lmp}}) has been registered with the
graphical environment to launch LAMMPS\sphinxhyphen{}GUI, an existing input file can
be launched with LAMMPS\sphinxhyphen{}GUI through double clicking.

\sphinxAtStartPar
Only one file can be edited at a time, so opening a new file with a
filled buffer closes that buffer.  If the buffer has unsaved
modifications, you are asked to either cancel the operation, discard the
changes, or save them.  A buffer with modifications can be saved any
time from the “File” menu, by the keyboard shortcut \sphinxtitleref{Ctrl\sphinxhyphen{}S}
(\sphinxtitleref{Command\sphinxhyphen{}S} on macOS), or by clicking on the “Save” button at the very
left in the status bar.


\paragraph{Running LAMMPS}
\label{\detokenize{Howto_lammps_gui:running-lammps}}
\sphinxAtStartPar
From within the LAMMPS\sphinxhyphen{}GUI main window LAMMPS can be started either from
the \sphinxstyleemphasis{Run} menu using the \sphinxstyleemphasis{Run LAMMPS from Editor Buffer} entry, by
the keyboard shortcut \sphinxtitleref{Ctrl\sphinxhyphen{}Enter} (\sphinxtitleref{Command\sphinxhyphen{}Enter} on macOS), or by
clicking on the green “Run” button in the status bar.  All of these
operations causes LAMMPS to process the entire input script in the
editor buffer, which may contain multiple {\hyperref[\detokenize{run::doc}]{\sphinxcrossref{\DUrole{doc}{run}}}} or
{\hyperref[\detokenize{minimize::doc}]{\sphinxcrossref{\DUrole{doc}{minimize}}}} commands.

\sphinxAtStartPar
LAMMPS runs in a separate thread, so the GUI stays responsive and is
able to interact with the running calculation and access data it
produces.  It is important to note that running LAMMPS this way is
using the contents of the input buffer for the run (via the
{\hyperref[\detokenize{Library_execute:_CPPv422lammps_commands_stringPvPKc}]{\sphinxcrossref{\sphinxcode{\sphinxupquote{lammps\_commands\_string()}}}}} function of the LAMMPS C\sphinxhyphen{}library
interface), and \sphinxstylestrong{not} the original file it was read from.  Thus, if
there are unsaved changes in the buffer, they \sphinxstyleemphasis{will} be used.  As an
alternative, it is also possible to run LAMMPS by reading the contents
of a file from the \sphinxstyleemphasis{Run LAMMPS from File} menu entry or with
\sphinxtitleref{Ctrl\sphinxhyphen{}Shift\sphinxhyphen{}Enter}.  This option may be required in some rare cases
where the input uses some functionality that is not compatible with
running LAMMPS from a string buffer.  For consistency, any unsaved
changes in the buffer must be either saved to the file or undone
before LAMMPS can be run from a file.

\noindent{\hspace*{\fill}\sphinxincludegraphics[scale=0.75]{{lammps-gui-running}.png}\hspace*{\fill}}

\sphinxAtStartPar
While LAMMPS is running, the contents of the status bar change.  On
the left side there is a text indicating that LAMMPS is running, which
also indicates the number of active threads, when thread\sphinxhyphen{}parallel
acceleration was selected in the \sphinxstyleemphasis{Preferences} dialog.  On the right
side, a progress bar is shown that displays the estimated progress for
the current {\hyperref[\detokenize{run::doc}]{\sphinxcrossref{\DUrole{doc}{run}}}} or {\hyperref[\detokenize{minimize::doc}]{\sphinxcrossref{\DUrole{doc}{minimize}}}} command.

\sphinxAtStartPar
Also, the line number of the currently executed command is highlighted
in green.

\sphinxAtStartPar
If an error occurs (in the example below the command {\hyperref[\detokenize{label::doc}]{\sphinxcrossref{\DUrole{doc}{label}}}} was incorrectly capitalized as “Label”), an error message
dialog is shown and the line of the input which triggered the error is
highlighted.  The state of LAMMPS in the status bar is set to “Failed.”
instead of “Ready.”

\noindent{\hspace*{\fill}\sphinxincludegraphics[scale=0.75]{{lammps-gui-run-error}.png}\hspace*{\fill}}

\sphinxAtStartPar
Up to three additional windows may open during a run:
\begin{itemize}
\item {} 
\sphinxAtStartPar
an \sphinxstyleemphasis{Output} window with the captured screen output from LAMMPS

\item {} 
\sphinxAtStartPar
a \sphinxstyleemphasis{Charts} window with a line graph created from thermodynamic output of the run

\item {} 
\sphinxAtStartPar
a \sphinxstyleemphasis{Slide Show} window with images created by a {\hyperref[\detokenize{dump_image::doc}]{\sphinxcrossref{\DUrole{doc}{dump image command}}}}
in the input

\end{itemize}

\sphinxAtStartPar
More information on those windows and how to adjust their behavior and
contents is given below.

\sphinxAtStartPar
An active LAMMPS run can be stopped cleanly by using either the \sphinxstyleemphasis{Stop
LAMMPS} entry in the \sphinxstyleemphasis{Run} menu, the keyboard shortcut \sphinxtitleref{Ctrl\sphinxhyphen{}/}
(\sphinxtitleref{Command\sphinxhyphen{}/} on macOS), or by clicking on the red button in the status
bar.  This will cause the running LAMMPS process to complete the current
timestep (or iteration for energy minimization) and then complete the
processing of the buffer while skipping all run or minimize commands.
This is equivalent to the input script command {\hyperref[\detokenize{timer::doc}]{\sphinxcrossref{\DUrole{doc}{timer timeout 0}}}} and is implemented by calling the
{\hyperref[\detokenize{Library_utility:_CPPv420lammps_force_timeoutPv}]{\sphinxcrossref{\sphinxcode{\sphinxupquote{lammps\_force\_timeout()}}}}} function of the LAMMPS C\sphinxhyphen{}library
interface.  Please see the corresponding documentation pages to
understand the implications of this operation.


\subsubsection{Output Window}
\label{\detokenize{Howto_lammps_gui:output-window}}
\sphinxAtStartPar
By default, when starting a run, an \sphinxstyleemphasis{Output} window opens that displays
the screen output of the running LAMMPS calculation, as shown below.
This text would normally be seen in the command line window.

\noindent{\hspace*{\fill}\sphinxincludegraphics[scale=0.5]{{lammps-gui-log}.png}\hspace*{\fill}}

\sphinxAtStartPar
LAMMPS\sphinxhyphen{}GUI captures the screen output from LAMMPS as it is generated and
updates the \sphinxstyleemphasis{Output} window regularly during a run.

\sphinxAtStartPar
By default, the \sphinxstyleemphasis{Output} window is replaced each time a run is started.
The runs are counted and the run number for the current run is displayed
in the window title.  It is possible to change the behavior of
LAMMPS\sphinxhyphen{}GUI in the preferences dialog to create a \sphinxstyleemphasis{new} \sphinxstyleemphasis{Output} window
for every run or to not show the current \sphinxstyleemphasis{Output} window.  It is also
possible to show or hide the \sphinxstyleemphasis{current} \sphinxstyleemphasis{Output} window from the \sphinxstyleemphasis{View}
menu.

\sphinxAtStartPar
The text in the \sphinxstyleemphasis{Output} window is read\sphinxhyphen{}only and cannot be modified, but
keyboard shortcuts to select and copy all or parts of the text can be
used to transfer text to another program. Also, the keyboard shortcut
\sphinxtitleref{Ctrl\sphinxhyphen{}S} (\sphinxtitleref{Command\sphinxhyphen{}S} on macOS) is available to save the \sphinxstyleemphasis{Output} buffer to a
file.  The “Select All” and “Copy” functions, as well as a “Save Log to
File” option are also available from a context menu by clicking with the
right mouse button into the \sphinxstyleemphasis{Output} window text area.

\noindent{\hspace*{\fill}\sphinxincludegraphics[scale=0.5]{{lammps-gui-yaml}.png}\hspace*{\fill}}

\sphinxAtStartPar
Should the \sphinxstyleemphasis{Output} window contain embedded YAML format text (see above for a
demonstration), for example from using {\hyperref[\detokenize{thermo_style::doc}]{\sphinxcrossref{\DUrole{doc}{thermo\_style yaml}}}} or {\hyperref[\detokenize{thermo_modify::doc}]{\sphinxcrossref{\DUrole{doc}{thermo\_modify line yaml}}}}, the
keyboard shortcut \sphinxtitleref{Ctrl\sphinxhyphen{}Y} (\sphinxtitleref{Command\sphinxhyphen{}Y} on macOS) is available to save
only the YAML parts to a file.  This option is also available from a
context menu by clicking with the right mouse button into the \sphinxstyleemphasis{Output} window
text area.


\subsubsection{Charts Window}
\label{\detokenize{Howto_lammps_gui:charts-window}}
\sphinxAtStartPar
By default, when starting a run, a \sphinxstyleemphasis{Charts} window opens that displays a
plot of thermodynamic output of the LAMMPS calculation as shown below.

\noindent{\hspace*{\fill}\sphinxincludegraphics[scale=0.33]{{lammps-gui-chart}.png}\hspace*{\fill}}

\sphinxAtStartPar
The drop down menu on the top right allows selection of different
properties that are computed and written to thermo output.  Only one
property can be shown at a time.  The plots are updated regularly with
new data as the run progresses, so they can be used to visually monitor
the evolution of available properties.  The update interval can be set
in the \sphinxstyleemphasis{Preferences} dialog.  By default, the raw data for the selected
property is plotted as a blue graph. As soon as there are a sufficient
number of data points, there will be a second graph shown in red with a
smoothed version of the data.  From the drop down menu on the top left,
you can select whether to plot only the raw data, only the smoothed
data or both.  The smoothing uses a \sphinxhref{https://en.wikipedia.org/wiki/Savitzky\%E2\%80\%93Golay\_filter}{Savitzky\sphinxhyphen{}Golay convolution filter} The
window width (left) and order (right) parameters can be set in the boxes
next to the drop down menu.  Default settings are 10 and 4 which means
that the smoothing window includes 10 points each to the left and the
right of the current data point and a fourth order polynomial is fit to
the data in the window.

\sphinxAtStartPar
You can use the mouse to zoom into the graph (hold the left button and
drag to mark an area) or zoom out (right click) and you can reset the
view with a click to the “lens” button next to the data drop down menu.

\sphinxAtStartPar
The window title shows the current run number that this chart window
corresponds to.  Same as for the \sphinxstyleemphasis{Output} window, the chart window is
replaced on each new run, but the behavior can be changed in the
\sphinxstyleemphasis{Preferences} dialog.

\sphinxAtStartPar
From the \sphinxstyleemphasis{File} menu on the top left, it is possible to save an image
of the currently displayed plot or export the data in either plain text
columns (for use by plotting tools like \sphinxhref{http://www.gnuplot.info/}{gnuplot} or \sphinxhref{https://plasma-gate.weizmann.ac.il/Grace/}{grace}), as CSV data which can be
imported for further processing with Microsoft Excel \sphinxhref{https://www.libreoffice.org/}{LibreOffice Calc} or with Python via \sphinxhref{https://pandas.pydata.org/}{pandas}, or as YAML which can be imported into
Python with \sphinxhref{https://pyyaml.org/}{PyYAML} or pandas.

\sphinxAtStartPar
Thermo output data from successive run commands in the input script is
combined into a single data set unless the format, number, or names of
output columns are changed with a {\hyperref[\detokenize{thermo_style::doc}]{\sphinxcrossref{\DUrole{doc}{thermo\_style}}}} or
a {\hyperref[\detokenize{thermo_modify::doc}]{\sphinxcrossref{\DUrole{doc}{thermo\_modify}}}} command, or the current time step
is reset with {\hyperref[\detokenize{reset_timestep::doc}]{\sphinxcrossref{\DUrole{doc}{reset\_timestep}}}}, or if a
{\hyperref[\detokenize{clear::doc}]{\sphinxcrossref{\DUrole{doc}{clear}}}} command is issued.  This is where the YAML export
from the \sphinxstyleemphasis{Charts} window differs from that of the \sphinxstyleemphasis{Output} window:
here you get the compounded data set starting with the last change of
output fields or timestep setting, while the export from the log will
contain \sphinxstyleemphasis{all} YAML output but \sphinxstyleemphasis{segmented} into individual runs.


\subsubsection{Image Slide Show}
\label{\detokenize{Howto_lammps_gui:image-slide-show}}
\sphinxAtStartPar
By default, if the LAMMPS input contains a {\hyperref[\detokenize{dump_image::doc}]{\sphinxcrossref{\DUrole{doc}{dump image}}}} command, a “Slide Show” window opens which loads and
displays the images created by LAMMPS as they are written.  This is a
convenient way to visually monitor the progress of the simulation.

\noindent{\hspace*{\fill}\sphinxincludegraphics[scale=0.5]{{lammps-gui-slideshow}.png}\hspace*{\fill}}

\sphinxAtStartPar
The various buttons at the bottom right of the window allow single
stepping through the sequence of images or playing an animation (as a
continuous loop or once from first to last).  It is also possible to
zoom in or zoom out of the displayed images. The button on the very
left triggers an export of the slide show animation to a movie file,
provided the \sphinxhref{https://ffmpeg.org/}{FFmpeg program} is installed.

\sphinxAtStartPar
When clicking on the “garbage can” icon, all image files of the slide
show will be deleted.  Since their number can be large for long
simulations, this option enables to safely and quickly clean up the
clutter caused in the working directory by those image files without
risk of deleting other files by accident when using wildcards.


\subsubsection{Variable Info}
\label{\detokenize{Howto_lammps_gui:variable-info}}
\sphinxAtStartPar
During a run, it may be of interest to monitor the value of input script
variables, for example to monitor the progress of loops.  This can be
done by enabling the “Variables Window” in the \sphinxstyleemphasis{View} menu or by using
the \sphinxtitleref{Ctrl\sphinxhyphen{}Shift\sphinxhyphen{}W} keyboard shortcut.  This shows info similar to the
{\hyperref[\detokenize{info::doc}]{\sphinxcrossref{\DUrole{doc}{info variables}}}} command in a separate window as shown
below.

\noindent{\hspace*{\fill}\sphinxincludegraphics[scale=0.5]{{lammps-gui-variable-info}.png}\hspace*{\fill}}

\sphinxAtStartPar
Like for the \sphinxstyleemphasis{Output} and \sphinxstyleemphasis{Charts} windows, its content is continuously
updated during a run.  It will show “(none)” if there are no variables
defined.  Note that it is also possible to \sphinxstyleemphasis{set} {\hyperref[\detokenize{variable::doc}]{\sphinxcrossref{\DUrole{doc}{index style
variables}}}}, that would normally be set via command line
flags, via the “Set Variables…” dialog from the \sphinxstyleemphasis{Run} menu.
LAMMPS\sphinxhyphen{}GUI automatically defines the variable “gui\_run” to the current
value of the run counter.  That way it is possible to automatically
record a separate log for each run attempt by using the command

\begin{sphinxVerbatim}[commandchars=\\\{\}]
\PYG{k}{log}\PYG{+w}{ }\PYG{n}{logfile}\PYG{o}{\PYGZhy{}}\PYG{n+nv}{\PYGZdl{}\PYGZob{}}\PYG{n+nv}{gui\PYGZus{}run}\PYG{n+nv}{\PYGZcb{}}\PYG{n}{.txt}
\end{sphinxVerbatim}

\sphinxAtStartPar
at the beginning of an input file. That would record logs to files
\sphinxcode{\sphinxupquote{logfile\sphinxhyphen{}1.txt}}, \sphinxcode{\sphinxupquote{logfile\sphinxhyphen{}2.txt}}, and so on for successive runs.


\subsubsection{Snapshot Image Viewer}
\label{\detokenize{Howto_lammps_gui:snapshot-image-viewer}}\label{\detokenize{Howto_lammps_gui:snapshot-viewer}}
\sphinxAtStartPar
By selecting the \sphinxstyleemphasis{Create Image} entry in the \sphinxstyleemphasis{Run} menu, or by
hitting the \sphinxtitleref{Ctrl\sphinxhyphen{}I} (\sphinxtitleref{Command\sphinxhyphen{}I} on macOS) keyboard shortcut, or by
clicking on the “palette” button in the status bar of the \sphinxstyleemphasis{Editor}
window, LAMMPS\sphinxhyphen{}GUI sends a custom {\hyperref[\detokenize{dump_image::doc}]{\sphinxcrossref{\DUrole{doc}{write\_dump image}}}}
command to LAMMPS and reads back the resulting snapshot image with the
current state of the system into an image viewer.  This functionality is
\sphinxstyleemphasis{not} available \sphinxstyleemphasis{during} an ongoing run.  In case LAMMPS is not yet
initialized, LAMMPS\sphinxhyphen{}GUI tries to identify the line with the first run or
minimize command and execute all commands from the input buffer up to
that line, and then executes a “run 0” command.  This initializes the
system so an image of the initial state of the system can be rendered.
If there was an error in that process, the snapshot image viewer does
not appear.

\sphinxAtStartPar
When possible, LAMMPS\sphinxhyphen{}GUI tries to detect which elements the atoms
correspond to (via their mass) and then colorize them in the image and
set their atom diameters accordingly.  If this is not possible, for
instance when using reduced (= ‘lj’) {\hyperref[\detokenize{units::doc}]{\sphinxcrossref{\DUrole{doc}{units}}}}, then
LAMMPS\sphinxhyphen{}GUI will check the current pair style and if it is a
Lennard\sphinxhyphen{}Jones type potential, it will extract the \sphinxstyleemphasis{sigma} parameter
for each atom type and assign atom diameters from those numbers.
For cases where atom diameters are not auto\sphinxhyphen{}detected, the \sphinxstyleemphasis{Atom size} field
can be edited and a suitable value set manually. The default value
is inferred from the x\sphinxhyphen{}direction lattice spacing.

\sphinxAtStartPar
If elements cannot be detected the default sequence of colors of the
{\hyperref[\detokenize{dump_image::doc}]{\sphinxcrossref{\DUrole{doc}{dump image}}}} command is assigned to the different atom
types.

\sphinxAtStartPar
\sphinxincludegraphics[width=0.480\linewidth]{{lammps-gui-image}.png}  \sphinxincludegraphics[width=0.480\linewidth]{{lammps-gui-funnel}.png}

\sphinxAtStartPar
The default image size, some default image quality settings, the view
style and some colors can be changed in the \sphinxstyleemphasis{Preferences} dialog
window.  From the image viewer window further adjustments can be made:
actual image size, high\sphinxhyphen{}quality (SSAO) rendering, anti\sphinxhyphen{}aliasing, view
style, display of box or axes, zoom factor.  The view of the system can
be rotated horizontally and vertically.  It is also possible to only
display the atoms within a group defined in the input script (default is
“all”).  The image can also be re\sphinxhyphen{}centered on the center of mass of the
selected group.  After each change, the image is rendered again and the
display updated.  The small palette icon on the top left is colored
while LAMMPS is running to render the new image; it is grayed out when
LAMMPS is finished.  When there are many atoms to render and high
quality images with anti\sphinxhyphen{}aliasing are requested, re\sphinxhyphen{}rendering may take
several seconds.  From the \sphinxstyleemphasis{File} menu of the image window, the
current image can be saved to a file (keyboard shortcut \sphinxtitleref{Ctrl\sphinxhyphen{}S}) or
copied to the clipboard (keyboard shortcut \sphinxtitleref{Ctrl\sphinxhyphen{}C}) for pasting the
image into another application.

\sphinxAtStartPar
From the \sphinxstyleemphasis{File} menu it is also possible to copy the current
{\hyperref[\detokenize{dump_image::doc}]{\sphinxcrossref{\DUrole{doc}{dump image}}}} and {\hyperref[\detokenize{dump_image::doc}]{\sphinxcrossref{\DUrole{doc}{dump\_modify}}}}
commands to the clipboard so they can be pasted into a LAMMPS input file
so that the visualization settings of the snapshot image can be repeated
for the entire simulation (and thus be repeated in the slide show
viewer). This feature has the keyboard shortcut \sphinxtitleref{Ctrl\sphinxhyphen{}D}.


\subsubsection{Editor Window}
\label{\detokenize{Howto_lammps_gui:editor-window}}
\sphinxAtStartPar
The \sphinxstyleemphasis{Editor} window of LAMMPS\sphinxhyphen{}GUI has most of the usual functionality
that similar programs have: text selection via mouse or with cursor
moves while holding the Shift key, Cut (\sphinxtitleref{Ctrl\sphinxhyphen{}X}), Copy (\sphinxtitleref{Ctrl\sphinxhyphen{}C}),
Paste (\sphinxtitleref{Ctrl\sphinxhyphen{}V}), Undo (\sphinxtitleref{Ctrl\sphinxhyphen{}Z}), Redo (\sphinxtitleref{Ctrl\sphinxhyphen{}Shift\sphinxhyphen{}Z}), Select All
(\sphinxtitleref{Ctrl\sphinxhyphen{}A}).  When trying to exit the editor with a modified buffer, a
dialog will pop up asking whether to cancel the exit operation, or to
save or not save the buffer contents to a file.

\sphinxAtStartPar
The editor has an auto\sphinxhyphen{}save mode that can be enabled or disabled in the
\sphinxstyleemphasis{Preferences} dialog.  In auto\sphinxhyphen{}save mode, the editor buffer is
automatically saved before running LAMMPS or before exiting LAMMPS\sphinxhyphen{}GUI.


\paragraph{Context Specific Word Completion}
\label{\detokenize{Howto_lammps_gui:context-specific-word-completion}}
\sphinxAtStartPar
By default, LAMMPS\sphinxhyphen{}GUI displays a small pop\sphinxhyphen{}up frame with possible
choices for LAMMPS input script commands or styles after 2 characters of
a word have been typed.

\noindent{\hspace*{\fill}\sphinxincludegraphics[scale=0.75]{{lammps-gui-complete}.png}\hspace*{\fill}}

\sphinxAtStartPar
The word can then be completed through selecting an entry by scrolling
up and down with the cursor keys and selecting with the ‘Enter’ key or
by clicking on the entry with the mouse.  The automatic completion
pop\sphinxhyphen{}up can be disabled in the \sphinxstyleemphasis{Preferences} dialog, but the completion
can still be requested manually by either hitting the ‘Shift\sphinxhyphen{}TAB’ key or
by right\sphinxhyphen{}clicking with the mouse and selecting the option from the
context menu.  Most of the completion information is retrieved from the
active LAMMPS instance and thus it shows only available options that
have been enabled when compiling LAMMPS. That list, however, excludes
accelerated styles and commands; for improved clarity, only the
non\sphinxhyphen{}suffix version of styles are shown.


\paragraph{Line Reformatting}
\label{\detokenize{Howto_lammps_gui:line-reformatting}}
\sphinxAtStartPar
The editor supports reformatting lines according to the syntax in order
to have consistently aligned lines.  This primarily means adding
whitespace padding to commands, type specifiers, IDs and names.  This
reformatting is performed manually by hitting the ‘Tab’ key.  It is
also possible to have this done automatically when hitting the ‘Enter’
key to start a new line.  This feature can be turned on or off in the
\sphinxstyleemphasis{Preferences} dialog for \sphinxstyleemphasis{Editor Settings} with the
“Reformat with ‘Enter’” checkbox. The amount of padding for multiple
categories can be adjusted in the same dialog.

\sphinxAtStartPar
Internally this functionality is achieved by splitting the line into
“words” and then putting it back together with padding added where the
context can be detected; otherwise a single space is used between words.


\paragraph{Context Specific Help}
\label{\detokenize{Howto_lammps_gui:context-specific-help}}
\sphinxAtStartPar
\sphinxincludegraphics[width=0.480\linewidth]{{lammps-gui-popup-help}.png}  \sphinxincludegraphics[width=0.480\linewidth]{{lammps-gui-popup-view}.png}

\sphinxAtStartPar
A unique feature of LAMMPS\sphinxhyphen{}GUI is the option to look up the LAMMPS
documentation for the command in the current line.  This can be done by
either clicking the right mouse button or by using the \sphinxtitleref{Ctrl\sphinxhyphen{}?} keyboard
shortcut.  When using the mouse, there are additional entries in the
context menu that open the corresponding documentation page in the
online LAMMPS documentation in a web browser window.  When using the
keyboard, the first of those entries is chosen.

\sphinxAtStartPar
If the word under the cursor is a file, then additionally the context
menu has an entry to open the file in a read\sphinxhyphen{}only text viewer window.
If the file is a LAMMPS restart file, instead the menu entry offers to
{\hyperref[\detokenize{Howto_lammps_gui:inspect-restart}]{\sphinxcrossref{\DUrole{std,std-ref}{inspect the restart}}}}.

\sphinxAtStartPar
The text viewer is a convenient way to view the contents of files that
are referenced in the input.  The file viewer also supports on\sphinxhyphen{}the\sphinxhyphen{}fly
decompression based on the file name suffix in a {\hyperref[\detokenize{Build_settings:gzip}]{\sphinxcrossref{\DUrole{std,std-ref}{similar fashion
as available with LAMMPS}}}}.  If the necessary decompression
program is missing or the file cannot be decompressed, the viewer window
will contain a corresponding message.


\paragraph{Inspecting a Restart file}
\label{\detokenize{Howto_lammps_gui:inspecting-a-restart-file}}\label{\detokenize{Howto_lammps_gui:inspect-restart}}
\sphinxAtStartPar
When LAMMPS\sphinxhyphen{}GUI is asked to “Inspect a Restart”, it will read the
restart file into a LAMMPS instance and then open three different
windows.  The first window is a text viewer with the output of an
{\hyperref[\detokenize{info::doc}]{\sphinxcrossref{\DUrole{doc}{info command}}}} with system information stored in the
restart.  The second window is text viewer containing a data file
generated with a {\hyperref[\detokenize{write_data::doc}]{\sphinxcrossref{\DUrole{doc}{write\_data command}}}}.  The third
window is a {\hyperref[\detokenize{Howto_lammps_gui:snapshot-viewer}]{\sphinxcrossref{\DUrole{std,std-ref}{Snapshot Image Viewer}}}} containing a
visualization of the system in the restart.

\sphinxAtStartPar
If the restart file is larger than 250 MBytes, a dialog will ask
for confirmation before continuing, since large restart files
may require large amounts of RAM since the entire system must
be read into RAM.  Thus restart file for large simulations that
have been run on an HPC cluster may overload a laptop or local
workstation. The \sphinxstyleemphasis{Show Details…} button will display a rough
estimate of the additional memory required.


\subsubsection{Menu}
\label{\detokenize{Howto_lammps_gui:menu}}
\sphinxAtStartPar
The menu bar has entries \sphinxstyleemphasis{File}, \sphinxstyleemphasis{Edit}, \sphinxstyleemphasis{Run}, \sphinxstyleemphasis{View}, and
\sphinxstyleemphasis{About}.  Instead of using the mouse to click on them, the individual
menus can also be activated by hitting the \sphinxtitleref{Alt} key together with the
corresponding underlined letter, that is \sphinxtitleref{Alt\sphinxhyphen{}F} activates the
\sphinxstyleemphasis{File} menu.  For the corresponding activated sub\sphinxhyphen{}menus, the key
corresponding the underlined letters can be used to select entries
instead of using the mouse.


\paragraph{File}
\label{\detokenize{Howto_lammps_gui:file}}
\sphinxAtStartPar
The \sphinxstyleemphasis{File} menu offers the usual options:
\begin{itemize}
\item {} 
\sphinxAtStartPar
\sphinxstyleemphasis{New} clears the current buffer and resets the file name to \sphinxcode{\sphinxupquote{*unknown*}}

\item {} 
\sphinxAtStartPar
\sphinxstyleemphasis{Open} opens a dialog to select a new file for editing in the \sphinxstyleemphasis{Editor}

\item {} 
\sphinxAtStartPar
\sphinxstyleemphasis{View} opens a dialog to select a file for viewing in a \sphinxstyleemphasis{separate} window (read\sphinxhyphen{}only) with support for on\sphinxhyphen{}the\sphinxhyphen{}fly decompression as explained above.

\item {} 
\sphinxAtStartPar
\sphinxstyleemphasis{Inspect restart} opens a dialog to select a file.  If that file is a {\hyperref[\detokenize{write_restart::doc}]{\sphinxcrossref{\DUrole{doc}{LAMMPS restart}}}} three windows with {\hyperref[\detokenize{Howto_lammps_gui:inspect-restart}]{\sphinxcrossref{\DUrole{std,std-ref}{information about the file are opened}}}}.

\item {} 
\sphinxAtStartPar
\sphinxstyleemphasis{Save} saves the current file; if the file name is \sphinxcode{\sphinxupquote{*unknown*}}
a dialog will open to select a new file name

\item {} 
\sphinxAtStartPar
\sphinxstyleemphasis{Save As} opens a dialog to select and new file name (and folder, if
desired) and saves the buffer to it.  Writing the buffer to a
different folder will also switch the current working directory to
that folder.

\item {} 
\sphinxAtStartPar
\sphinxstyleemphasis{Quit} exits LAMMPS\sphinxhyphen{}GUI. If there are unsaved changes, a dialog will
appear to either cancel the operation, or to save, or to not save the
modified buffer.

\end{itemize}

\sphinxAtStartPar
In addition, up to 5 recent file names will be listed after the \sphinxstyleemphasis{Open}
entry that allows re\sphinxhyphen{}opening recently opened files.  This list is stored
when quitting and recovered when starting again.


\paragraph{Edit}
\label{\detokenize{Howto_lammps_gui:edit}}
\sphinxAtStartPar
The \sphinxstyleemphasis{Edit} menu offers the usual editor functions like \sphinxstyleemphasis{Undo}, \sphinxstyleemphasis{Redo},
\sphinxstyleemphasis{Cut}, \sphinxstyleemphasis{Copy}, \sphinxstyleemphasis{Paste}, and a \sphinxstyleemphasis{Find and Replace} dialog (keyboard
shortcut \sphinxtitleref{Ctrl\sphinxhyphen{}F}).  It can also open a \sphinxstyleemphasis{Preferences} dialog (keyboard
shortcut \sphinxtitleref{Ctrl\sphinxhyphen{}P}) and allows deleting all stored preferences and
settings, so they are reset to their default values.


\paragraph{Run}
\label{\detokenize{Howto_lammps_gui:run}}
\sphinxAtStartPar
The \sphinxstyleemphasis{Run} menu has options to start and stop a LAMMPS process.  Rather
than calling the LAMMPS executable as a separate executable, the
LAMMPS\sphinxhyphen{}GUI is linked to the LAMMPS library and thus can run LAMMPS
internally through the {\hyperref[\detokenize{Library:lammps-c-api}]{\sphinxcrossref{\DUrole{std,std-ref}{LAMMPS C\sphinxhyphen{}library interface}}}}
in a separate thread.

\sphinxAtStartPar
Specifically, a LAMMPS instance will be created by calling
{\hyperref[\detokenize{Library_create:_CPPv418lammps_open_no_mpiiPPcPPv}]{\sphinxcrossref{\sphinxcode{\sphinxupquote{lammps\_open\_no\_mpi()}}}}}.  The buffer contents are then executed by
calling {\hyperref[\detokenize{Library_execute:_CPPv422lammps_commands_stringPvPKc}]{\sphinxcrossref{\sphinxcode{\sphinxupquote{lammps\_commands\_string()}}}}}.  Certain commands and
features are only available after a LAMMPS instance is created.  Its
presence is indicated by a small LAMMPS \sphinxcode{\sphinxupquote{L}} logo in the status bar
at the bottom left of the main window.  As an alternative, it is also
possible to run LAMMPS using the contents of the edited file by
reading the file.  This is mainly provided as a fallback option in
case the input uses some feature that is not available when running
from a string buffer.

\sphinxAtStartPar
The LAMMPS calculations are run in a concurrent thread so that the GUI
can stay responsive and be updated during the run.  The GUI can retrieve
data from the running LAMMPS instance and tell it to stop at the next
timestep.  The \sphinxstyleemphasis{Stop LAMMPS} entry will do this by calling the
{\hyperref[\detokenize{Library_utility:_CPPv420lammps_force_timeoutPv}]{\sphinxcrossref{\sphinxcode{\sphinxupquote{lammps\_force\_timeout()}}}}} library function, which is equivalent
to a {\hyperref[\detokenize{timer::doc}]{\sphinxcrossref{\DUrole{doc}{timer timeout 0}}}} command.

\sphinxAtStartPar
The \sphinxstyleemphasis{Set Variables…} entry opens a dialog box where
{\hyperref[\detokenize{variable::doc}]{\sphinxcrossref{\DUrole{doc}{index style variables}}}} can be set. Those variables
are passed to the LAMMPS instance when it is created and are thus
set \sphinxstyleemphasis{before} a run is started.

\noindent{\hspace*{\fill}\sphinxincludegraphics[scale=0.5]{{lammps-gui-variables}.png}\hspace*{\fill}}

\sphinxAtStartPar
The \sphinxstyleemphasis{Set Variables} dialog will be pre\sphinxhyphen{}populated with entries that
are set as index variables in the input and any variables that are
used but not defined, if the built\sphinxhyphen{}in parser can detect them.  New
rows for additional variables can be added through the \sphinxstyleemphasis{Add Row}
button and existing rows can be deleted by clicking on the \sphinxstyleemphasis{X} icons
on the right.

\sphinxAtStartPar
The \sphinxstyleemphasis{Create Image} entry will send a {\hyperref[\detokenize{dump_image::doc}]{\sphinxcrossref{\DUrole{doc}{dump image}}}}
command to the LAMMPS instance, read the resulting file, and show it
in an \sphinxstyleemphasis{Image Viewer} window.

\sphinxAtStartPar
The \sphinxstyleemphasis{View in OVITO} entry will launch \sphinxhref{https://ovito.org}{OVITO}
with a {\hyperref[\detokenize{write_data::doc}]{\sphinxcrossref{\DUrole{doc}{data file}}}} containing the current state of
the system.  This option is only available if LAMMPS\sphinxhyphen{}GUI can find
the OVITO executable in the system path.

\sphinxAtStartPar
The \sphinxstyleemphasis{View in VMD} entry will launch VMD with a {\hyperref[\detokenize{write_data::doc}]{\sphinxcrossref{\DUrole{doc}{data file}}}} containing the current state of the system.  This option
is only available if LAMMPS\sphinxhyphen{}GUI can find the VMD executable in the
system path.


\paragraph{View}
\label{\detokenize{Howto_lammps_gui:view}}
\sphinxAtStartPar
The \sphinxstyleemphasis{View} menu offers to show or hide additional windows with log
output, charts, slide show, variables, or snapshot images.  The
default settings for their visibility can be changed in the
\sphinxstyleemphasis{Preferences} dialog.


\paragraph{About}
\label{\detokenize{Howto_lammps_gui:about}}
\sphinxAtStartPar
The \sphinxstyleemphasis{About} menu finally offers a couple of dialog windows and an
option to launch the LAMMPS online documentation in a web browser.  The
\sphinxstyleemphasis{About LAMMPS\sphinxhyphen{}GUI} entry displays a dialog with a summary of the
configuration settings of the LAMMPS library in use and the version
number of LAMMPS\sphinxhyphen{}GUI itself.  The \sphinxstyleemphasis{Quick Help} displays a dialog with
a minimal description of LAMMPS\sphinxhyphen{}GUI.  The \sphinxstyleemphasis{LAMMPS\sphinxhyphen{}GUI Howto} entry
will open this documentation page from the online documentation in a web
browser window.  The \sphinxstyleemphasis{LAMMPS Manual} entry will open the main page of
the LAMMPS online documentation in a web browser window.
The \sphinxstyleemphasis{LAMMPS Tutorial} entry will open the main page of the set of
LAMMPS tutorials authored and maintained by Simon Gravelle at
\sphinxurl{https://lammpstutorials.github.io/} in a web browser window.


\bigskip\hrule\bigskip



\subsubsection{Find and Replace}
\label{\detokenize{Howto_lammps_gui:find-and-replace}}
\noindent{\hspace*{\fill}\sphinxincludegraphics[scale=0.33]{{lammps-gui-find}.png}\hspace*{\fill}}

\sphinxAtStartPar
The \sphinxstyleemphasis{Find and Replace} dialog allows searching for and replacing
text in the \sphinxstyleemphasis{Editor} window.

\sphinxAtStartPar
The dialog can be opened either from the \sphinxstyleemphasis{Edit} menu or with the
keyboard shortcut \sphinxtitleref{Ctrl\sphinxhyphen{}F}. You can enter the text to search for.
Through three check\sphinxhyphen{}boxes the search behavior can be adjusted:
\begin{itemize}
\item {} 
\sphinxAtStartPar
If checked, “Match case” does a case sensitive search; otherwise
the search is case insensitive.

\item {} 
\sphinxAtStartPar
If checked, “Wrap around” starts searching from the start of the
document, if there is no match found from the current cursor position
until the end of the document; otherwise the search will stop.

\item {} 
\sphinxAtStartPar
If checked, the “Whole word” setting only finds full word matches
(white space and special characters are word boundaries).

\end{itemize}

\sphinxAtStartPar
Clicking on the \sphinxstyleemphasis{Next} button will search for the next occurrence of the
search text and select / highlight it. Clicking on the \sphinxstyleemphasis{Replace} button
will replace an already highlighted search text and find the next one.
If no text is selected, or the selected text does not match the
selection string, then the first click on the \sphinxstyleemphasis{Replace} button will
only search and highlight the next occurrence of the search string.
Clicking on the \sphinxstyleemphasis{Replace All} button will replace all occurrences from
the cursor position to the end of the file; if the \sphinxstyleemphasis{Wrap around} box is
checked, then it will replace \sphinxstylestrong{all} occurrences in the \sphinxstylestrong{entire}
document.  Clicking on the \sphinxstyleemphasis{Done} button will dismiss the dialog.


\bigskip\hrule\bigskip



\subsubsection{Preferences}
\label{\detokenize{Howto_lammps_gui:preferences}}
\sphinxAtStartPar
The \sphinxstyleemphasis{Preferences} dialog allows customization of the behavior and
look of LAMMPS\sphinxhyphen{}GUI.  The settings are grouped and each group is
displayed within a tab.

\sphinxAtStartPar
\sphinxincludegraphics[width=0.240\linewidth]{{lammps-gui-prefs-general}.png}  \sphinxincludegraphics[width=0.240\linewidth]{{lammps-gui-prefs-accel}.png}  \sphinxincludegraphics[width=0.240\linewidth]{{lammps-gui-prefs-image}.png}  \sphinxincludegraphics[width=0.240\linewidth]{{lammps-gui-prefs-editor}.png}


\paragraph{General Settings:}
\label{\detokenize{Howto_lammps_gui:general-settings}}\begin{itemize}
\item {} 
\sphinxAtStartPar
\sphinxstyleemphasis{Echo input to log:} when checked, all input commands, including
variable expansions, are echoed to the \sphinxstyleemphasis{Output} window. This is
equivalent to using \sphinxtitleref{\sphinxhyphen{}echo screen} at the command line.  There is no
log \sphinxstyleemphasis{file} produced by default, since LAMMPS\sphinxhyphen{}GUI uses \sphinxtitleref{\sphinxhyphen{}log none}.

\item {} 
\sphinxAtStartPar
\sphinxstyleemphasis{Include citation details:} when checked full citation info will be
included to the log window.  This is equivalent to using \sphinxtitleref{\sphinxhyphen{}cite
screen} on the command line.

\item {} 
\sphinxAtStartPar
\sphinxstyleemphasis{Show log window by default:} when checked, the screen output of a
LAMMPS run will be collected in a log window during the run

\item {} 
\sphinxAtStartPar
\sphinxstyleemphasis{Show chart window by default:} when checked, the thermodynamic
output of a LAMMPS run will be collected and displayed in a chart
window as line graphs.

\item {} 
\sphinxAtStartPar
\sphinxstyleemphasis{Show slide show window by default:} when checked, a slide show
window will be shown with images from a dump image command, if
present, in the LAMMPS input.

\item {} 
\sphinxAtStartPar
\sphinxstyleemphasis{Replace log window on new run:} when checked, an existing log
window will be replaced on a new LAMMPS run, otherwise each run will
create a new log window.

\item {} 
\sphinxAtStartPar
\sphinxstyleemphasis{Replace chart window on new run:} when checked, an existing chart
window will be replaced on a new LAMMPS run, otherwise each run will
create a new chart window.

\item {} 
\sphinxAtStartPar
\sphinxstyleemphasis{Replace image window on new render:} when checked, an existing
chart window will be replaced when a new snapshot image is requested,
otherwise each command will create a new image window.

\item {} 
\sphinxAtStartPar
\sphinxstyleemphasis{Path to LAMMPS Shared Library File:} this option is only visible
when LAMMPS\sphinxhyphen{}GUI was compiled to load the LAMMPS library at run time
instead of being linked to it directly.  With the \sphinxstyleemphasis{Browse..} button
or by changing the text, a different shared library file with a
different compilation of LAMMPS with different settings or from a
different version can be loaded.  After this setting was changed,
LAMMPS\sphinxhyphen{}GUI needs to be re\sphinxhyphen{}launched.

\item {} 
\sphinxAtStartPar
\sphinxstyleemphasis{Select Default Font:} Opens a font selection dialog where the type
and size for the default font (used for everything but the editor and
log) of the application can be set.

\item {} 
\sphinxAtStartPar
\sphinxstyleemphasis{Select Text Font:} Opens a font selection dialog where the type and
size for the text editor and log font of the application can be set.

\item {} 
\sphinxAtStartPar
\sphinxstyleemphasis{Data update interval:} Allows to set the time interval between data
updates during a LAMMPS run in milliseconds.  The default is to update
the data (for charts and output window) every 10 milliseconds.  This
is good for many cases.  Set this to 100 milliseconds or more if
LAMMPS\sphinxhyphen{}GUI consumes too many resources during a run.  For LAMMPS runs
that run \sphinxstyleemphasis{very} fast (for example in tutorial examples), however, data
may be missed and through lowering this interval, this can be
corrected.  However, this will make the GUI use more resources.  This
setting may be changed to a value between 1 and 1000 milliseconds.

\item {} 
\sphinxAtStartPar
\sphinxstyleemphasis{Charts update interval:} Allows to set the time interval between redrawing
the plots in the \sphinxstyleemphasis{Charts} window in milliseconds.  The default is to
redraw the plots every 500 milliseconds.  This is just for the drawing,
data collection is managed with the previous setting.

\end{itemize}


\paragraph{Accelerators:}
\label{\detokenize{Howto_lammps_gui:accelerators}}
\sphinxAtStartPar
This tab enables selection of an accelerator package for LAMMPS to use
and is equivalent to using the \sphinxtitleref{\sphinxhyphen{}suffix} and \sphinxtitleref{\sphinxhyphen{}package} flags on the
command line.  Only settings supported by the LAMMPS library and local
hardware are available.  The \sphinxtitleref{Number of threads} field allows setting
the maximum number of threads for the accelerator packages that use
threads.


\paragraph{Snapshot Image:}
\label{\detokenize{Howto_lammps_gui:snapshot-image}}
\sphinxAtStartPar
This tab allows setting defaults for the snapshot images displayed in
the \sphinxstyleemphasis{Image Viewer} window, such as its dimensions and the zoom factor
applied.  The \sphinxstyleemphasis{Antialias} switch will render images with twice the
number of pixels for width and height and then smoothly scale the image
back to the requested size.  This produces higher quality images with
smoother edges at the expense of requiring more CPU time to render the
image.  The \sphinxstyleemphasis{HQ Image mode} option turns on screen space ambient
occlusion (SSAO) mode when rendering images.  This is also more time
consuming, but produces a more ‘spatial’ representation of the system
shading of atoms by their depth.  The \sphinxstyleemphasis{Shiny Image mode} option will
render objects with a shiny surface when enabled.  Otherwise the
surfaces will be matted.  The \sphinxstyleemphasis{Show Box} option selects whether the
system box is drawn as a colored set of sticks.  Similarly, the \sphinxstyleemphasis{Show
Axes} option selects whether a representation of the three system axes
will be drawn as colored sticks. The \sphinxstyleemphasis{VDW Style} checkbox selects
whether atoms are represented by space filling spheres when checked or
by smaller spheres and sticks.  Finally there are a couple of drop down
lists to select the background and box colors.


\paragraph{Editor Settings:}
\label{\detokenize{Howto_lammps_gui:editor-settings}}
\sphinxAtStartPar
This tab allows tweaking settings of the editor window.  Specifically
the amount of padding to be added to LAMMPS commands, types or type
ranges, IDs (e.g. for fixes), and names (e.g. for groups).  The value
set is the minimum width for the text element and it can be chosen in
the range between 1 and 32.

\sphinxAtStartPar
The three settings which follow enable or disable the automatic
reformatting when hitting the ‘Enter’ key, the automatic display of
the completion pop\sphinxhyphen{}up window, and whether auto\sphinxhyphen{}save mode is enabled.
In auto\sphinxhyphen{}save mode the editor buffer is saved before a run or before
exiting LAMMPS\sphinxhyphen{}GUI.


\bigskip\hrule\bigskip



\subsubsection{Keyboard Shortcuts}
\label{\detokenize{Howto_lammps_gui:keyboard-shortcuts}}
\sphinxAtStartPar
Almost all functionality is accessible from the menu of the editor
window or through keyboard shortcuts.  The following shortcuts are
available (On macOS use the Command key instead of Ctrl/Control).


\begin{savenotes}\sphinxattablestart
\sphinxthistablewithglobalstyle
\centering
\begin{tabular}[t]{\X{16}{99}\X{19}{99}\X{13}{99}\X{16}{99}\X{13}{99}\X{22}{99}}
\sphinxtoprule
\sphinxstyletheadfamily 
\sphinxAtStartPar
Shortcut
&\sphinxstyletheadfamily 
\sphinxAtStartPar
Function
&\sphinxstyletheadfamily 
\sphinxAtStartPar
Shortcut
&\sphinxstyletheadfamily 
\sphinxAtStartPar
Function
&\sphinxstyletheadfamily 
\sphinxAtStartPar
Shortcut
&\sphinxstyletheadfamily 
\sphinxAtStartPar
Function
\\
\sphinxmidrule
\sphinxtableatstartofbodyhook
\sphinxAtStartPar
Ctrl+N
&
\sphinxAtStartPar
New File
&
\sphinxAtStartPar
Ctrl+Z
&
\sphinxAtStartPar
Undo edit
&
\sphinxAtStartPar
Ctrl+Enter
&
\sphinxAtStartPar
Run Input
\\
\sphinxhline
\sphinxAtStartPar
Ctrl+O
&
\sphinxAtStartPar
Open File
&
\sphinxAtStartPar
Ctrl+Shift+Z
&
\sphinxAtStartPar
Redo edit
&
\sphinxAtStartPar
Ctrl+/
&
\sphinxAtStartPar
Stop Active Run
\\
\sphinxhline
\sphinxAtStartPar
Ctrl+Shift+F
&
\sphinxAtStartPar
View File
&
\sphinxAtStartPar
Ctrl+C
&
\sphinxAtStartPar
Copy text
&
\sphinxAtStartPar
Ctrl+Shift+V
&
\sphinxAtStartPar
Set Variables
\\
\sphinxhline
\sphinxAtStartPar
Ctrl+S
&
\sphinxAtStartPar
Save File
&
\sphinxAtStartPar
Ctrl+X
&
\sphinxAtStartPar
Cut text
&
\sphinxAtStartPar
Ctrl+I
&
\sphinxAtStartPar
Snapshot Image
\\
\sphinxhline
\sphinxAtStartPar
Ctrl+Shift+S
&
\sphinxAtStartPar
Save File As
&
\sphinxAtStartPar
Ctrl+V
&
\sphinxAtStartPar
Paste text
&
\sphinxAtStartPar
Ctrl+L
&
\sphinxAtStartPar
Slide Show
\\
\sphinxhline
\sphinxAtStartPar
Ctrl+Q
&
\sphinxAtStartPar
Quit Application
&
\sphinxAtStartPar
Ctrl+A
&
\sphinxAtStartPar
Select All
&
\sphinxAtStartPar
Ctrl+F
&
\sphinxAtStartPar
Find and Replace
\\
\sphinxhline
\sphinxAtStartPar
Ctrl+W
&
\sphinxAtStartPar
Close Window
&
\sphinxAtStartPar
TAB
&
\sphinxAtStartPar
Reformat line
&
\sphinxAtStartPar
Shift+TAB
&
\sphinxAtStartPar
Show Completions
\\
\sphinxhline
\sphinxAtStartPar
Ctrl+Shift+Enter
&
\sphinxAtStartPar
Run File
&
\sphinxAtStartPar
Ctrl+Shift+W
&
\sphinxAtStartPar
Show Variables
&
\sphinxAtStartPar
Ctrl+P
&
\sphinxAtStartPar
Preferences
\\
\sphinxhline
\sphinxAtStartPar
Ctrl+Shift+A
&
\sphinxAtStartPar
About LAMMPS
&
\sphinxAtStartPar
Ctrl+Shift+H
&
\sphinxAtStartPar
Quick Help
&
\sphinxAtStartPar
Ctrl+Shift+G
&
\sphinxAtStartPar
LAMMPS\sphinxhyphen{}GUI Howto
\\
\sphinxhline
\sphinxAtStartPar
Ctrl+Shift+M
&
\sphinxAtStartPar
LAMMPS Manual
&
\sphinxAtStartPar
Ctrl+?
&
\sphinxAtStartPar
Context Help
&
\sphinxAtStartPar
Ctrl+Shift+T
&
\sphinxAtStartPar
LAMMPS Tutorial
\\
\sphinxbottomrule
\end{tabular}
\sphinxtableafterendhook\par
\sphinxattableend\end{savenotes}

\sphinxAtStartPar
Further editing keybindings \sphinxhref{https://doc.qt.io/qt-5/qplaintextedit.html\#editing-key-bindings}{are documented with the Qt documentation}.  In
case of conflicts the list above takes precedence.

\sphinxAtStartPar
All other windows only support a subset of keyboard shortcuts listed
above.  Typically, the shortcuts \sphinxtitleref{Ctrl\sphinxhyphen{}/} (Stop Run), \sphinxtitleref{Ctrl\sphinxhyphen{}W} (Close
Window), and \sphinxtitleref{Ctrl\sphinxhyphen{}Q} (Quit Application) are supported.

\sphinxstepscope


\subsection{Moltemplate Tutorial}
\label{\detokenize{Howto_moltemplate:moltemplate-tutorial}}\label{\detokenize{Howto_moltemplate::doc}}
\sphinxAtStartPar
In this tutorial, we are going to use the tool {\hyperref[\detokenize{Tools:moltemplate}]{\sphinxcrossref{\DUrole{std,std-ref}{Moltemplate}}}} to set up a classical molecular dynamic simulation using
the {\hyperref[\detokenize{Howto_moltemplate:oplsaa96}]{\sphinxcrossref{\DUrole{std,std-ref}{OPLS\sphinxhyphen{}AA force field}}}}. The first
task is to describe an organic compound and create a complete input deck
for LAMMPS. The second task is to map the OPLS\sphinxhyphen{}AA force field to a
molecular sample created with an external tool, e.g. PACKMOL, and
exported as a PDB file.  The files used in this tutorial can be found
in the \sphinxcode{\sphinxupquote{tools/moltemplate/tutorial\sphinxhyphen{}files}} folder of the LAMMPS
source code distribution.


\subsubsection{Simulating an organic solvent}
\label{\detokenize{Howto_moltemplate:simulating-an-organic-solvent}}
\sphinxAtStartPar
This example aims to create a cubic box of the organic solvent
formamide.

\sphinxAtStartPar
The first step is to create a molecular topology in the
LAMMPS\sphinxhyphen{}template (LT) file format representing a single molecule, which
will be stored in a Moltemplate object called \sphinxcode{\sphinxupquote{\_FAM inherits OPLSAA \{\}}}.
This command states that the object \sphinxcode{\sphinxupquote{\_FAM}} is based on an existing
object called \sphinxcode{\sphinxupquote{OPLSAA}}, which contains OPLS\sphinxhyphen{}AA parameters, atom type
definitions, partial charges, masses and bond\sphinxhyphen{}angle rules for many organic
and biological compounds.

\sphinxAtStartPar
The atomic structure is the starting point to populate the command
\sphinxcode{\sphinxupquote{write(\textquotesingle{}Data Atoms\textquotesingle{}) \{\}}}, which will write the \sphinxcode{\sphinxupquote{Atoms}} section in the
LAMMPS data file. The OPLS\sphinxhyphen{}AA force field uses the \sphinxcode{\sphinxupquote{atom\_style full}},
therefore, this column format is used:
\sphinxcode{\sphinxupquote{\# atomID molID atomType charge coordX coordY coordZ}}.
The \sphinxcode{\sphinxupquote{atomID}}s are replaced with Moltemplate \sphinxcode{\sphinxupquote{\$}}\sphinxhyphen{}type variables, which
are then substituted with unique numerical IDs. The same logic is applied
to the \sphinxcode{\sphinxupquote{molID}}, except that the same variable is used for the whole
molecule. The atom types are assigned using \sphinxcode{\sphinxupquote{@}}\sphinxhyphen{}type variables. The
assignment of atom types (e.g. \sphinxcode{\sphinxupquote{@atom:177}}, \sphinxcode{\sphinxupquote{@atom:178}}) is done using
the OPLS\sphinxhyphen{}AA atom types defined in the “In Charges” section of the file
\sphinxcode{\sphinxupquote{oplsaa.lt}}, looking for a reasonable match with the description of the atom.
The resulting file (\sphinxcode{\sphinxupquote{formamide.lt}}) follows:

\begin{sphinxVerbatim}[commandchars=\\\{\}]
\PYGZus{}FAM\PYG{+w}{ }inherits\PYG{+w}{ }OPLSAA\PYG{+w}{ }\PYG{o}{\PYGZob{}}

\PYG{+w}{  }\PYG{c+c1}{\PYGZsh{} atomID    molID  atomType  charge  coordX  coordY  coordZ}
\PYG{+w}{  }write\PYG{o}{(}\PYG{l+s+s1}{\PYGZsq{}Data Atoms\PYGZsq{}}\PYG{o}{)}\PYG{+w}{ }\PYG{o}{\PYGZob{}}
\PYG{+w}{    }\PYG{n+nv}{\PYGZdl{}atom}:C00\PYG{+w}{ }\PYG{n+nv}{\PYGZdl{}mol}\PYG{+w}{   }@atom:177\PYG{+w}{  }\PYG{l+m}{0}.00\PYG{+w}{    }\PYG{l+m}{0}.100\PYG{+w}{   }\PYG{l+m}{0}.490\PYG{+w}{   }\PYG{l+m}{0}.0
\PYG{+w}{    }\PYG{n+nv}{\PYGZdl{}atom}:O01\PYG{+w}{ }\PYG{n+nv}{\PYGZdl{}mol}\PYG{+w}{   }@atom:178\PYG{+w}{  }\PYG{l+m}{0}.00\PYG{+w}{    }\PYG{l+m}{1}.091\PYG{+w}{  }\PYGZhy{}0.250\PYG{+w}{   }\PYG{l+m}{0}.0
\PYG{+w}{    }\PYG{n+nv}{\PYGZdl{}atom}:N02\PYG{+w}{ }\PYG{n+nv}{\PYGZdl{}mol}\PYG{+w}{   }@atom:179\PYG{+w}{  }\PYG{l+m}{0}.00\PYG{+w}{   }\PYGZhy{}1.121\PYG{+w}{  }\PYGZhy{}0.181\PYG{+w}{   }\PYG{l+m}{0}.0
\PYG{+w}{    }\PYG{n+nv}{\PYGZdl{}atom}:H03\PYG{+w}{ }\PYG{n+nv}{\PYGZdl{}mol}\PYG{+w}{   }@atom:182\PYG{+w}{  }\PYG{l+m}{0}.00\PYG{+w}{   }\PYGZhy{}2.013\PYG{+w}{   }\PYG{l+m}{0}.272\PYG{+w}{   }\PYG{l+m}{0}.0
\PYG{+w}{    }\PYG{n+nv}{\PYGZdl{}atom}:H04\PYG{+w}{ }\PYG{n+nv}{\PYGZdl{}mol}\PYG{+w}{   }@atom:182\PYG{+w}{  }\PYG{l+m}{0}.00\PYG{+w}{   }\PYGZhy{}1.056\PYG{+w}{  }\PYGZhy{}1.190\PYG{+w}{   }\PYG{l+m}{0}.0
\PYG{+w}{    }\PYG{n+nv}{\PYGZdl{}atom}:H05\PYG{+w}{ }\PYG{n+nv}{\PYGZdl{}mol}\PYG{+w}{   }@atom:221\PYG{+w}{  }\PYG{l+m}{0}.00\PYG{+w}{    }\PYG{l+m}{0}.144\PYG{+w}{   }\PYG{l+m}{1}.570\PYG{+w}{   }\PYG{l+m}{0}.0
\PYG{+w}{  }\PYG{o}{\PYGZcb{}}

\PYG{+w}{  }\PYG{c+c1}{\PYGZsh{} A list of the bonds in the molecule:}
\PYG{+w}{  }\PYG{c+c1}{\PYGZsh{} BondID   AtomID1   AtomID2}
\PYG{+w}{  }write\PYG{o}{(}\PYG{l+s+s1}{\PYGZsq{}Data Bond List\PYGZsq{}}\PYG{o}{)}\PYG{+w}{ }\PYG{o}{\PYGZob{}}
\PYG{+w}{    }\PYG{n+nv}{\PYGZdl{}bond}:C1\PYG{+w}{ }\PYG{n+nv}{\PYGZdl{}atom}:C00\PYG{+w}{ }\PYG{n+nv}{\PYGZdl{}atom}:O01
\PYG{+w}{    }\PYG{n+nv}{\PYGZdl{}bond}:C2\PYG{+w}{ }\PYG{n+nv}{\PYGZdl{}atom}:C00\PYG{+w}{ }\PYG{n+nv}{\PYGZdl{}atom}:H05
\PYG{+w}{    }\PYG{n+nv}{\PYGZdl{}bond}:C3\PYG{+w}{ }\PYG{n+nv}{\PYGZdl{}atom}:C00\PYG{+w}{ }\PYG{n+nv}{\PYGZdl{}atom}:N02
\PYG{+w}{    }\PYG{n+nv}{\PYGZdl{}bond}:C4\PYG{+w}{ }\PYG{n+nv}{\PYGZdl{}atom}:N02\PYG{+w}{ }\PYG{n+nv}{\PYGZdl{}atom}:H03
\PYG{+w}{    }\PYG{n+nv}{\PYGZdl{}bond}:C5\PYG{+w}{ }\PYG{n+nv}{\PYGZdl{}atom}:N02\PYG{+w}{ }\PYG{n+nv}{\PYGZdl{}atom}:H04
\PYG{+w}{  }\PYG{o}{\PYGZcb{}}
\PYG{o}{\PYGZcb{}}
\end{sphinxVerbatim}

\sphinxAtStartPar
You don’t have to specify the charge in this example because they will
be assigned according to the atom type. Analogously, only a
“Data Bond List” section is needed as the atom type will determine the
bond type. The other bonded interactions (e.g. angles,
dihedrals, and impropers) will be automatically generated by
Moltemplate.

\sphinxAtStartPar
If the simulation is non\sphinxhyphen{}neutral, or Moltemplate complains that you have
missing bond, angle, or dihedral types, this means at least one of your
atom types is incorrect.

\sphinxAtStartPar
The second step is to create a master file with instructions to build a
starting structure and the LAMMPS commands to run an NPT simulation. The
master file (\sphinxcode{\sphinxupquote{solv\_01.lt}}) follows:

\begin{sphinxVerbatim}[commandchars=\\\{\}]
\PYG{c+c1}{\PYGZsh{} Import the force field.}
import\PYG{+w}{ }/usr/local/moltemplate/moltemplate/force\PYGZus{}fields/oplsaa.lt
import\PYG{+w}{ }formamide.lt\PYG{+w}{ }\PYG{c+c1}{\PYGZsh{} after oplsaa.lt, as it depends on it.}

\PYG{c+c1}{\PYGZsh{} Create the input sample.}
\PYG{n+nv}{solv}\PYG{+w}{ }\PYG{o}{=}\PYG{+w}{ }new\PYG{+w}{ }\PYGZus{}FAM\PYG{+w}{ }\PYG{o}{[}\PYG{l+m}{5}\PYG{o}{]}.move\PYG{o}{(}\PYG{+w}{ }\PYG{l+m}{4}.6,\PYG{+w}{ }\PYG{l+m}{0},\PYG{+w}{ }\PYG{l+m}{0}\PYG{o}{)}
\PYG{+w}{                }\PYG{o}{[}\PYG{l+m}{5}\PYG{o}{]}.move\PYG{o}{(}\PYG{+w}{ }\PYG{l+m}{0},\PYG{+w}{ }\PYG{l+m}{4}.6,\PYG{+w}{ }\PYG{l+m}{0}\PYG{o}{)}
\PYG{+w}{                }\PYG{o}{[}\PYG{l+m}{5}\PYG{o}{]}.move\PYG{o}{(}\PYG{+w}{ }\PYG{l+m}{0},\PYG{+w}{ }\PYG{l+m}{0},\PYG{+w}{ }\PYG{l+m}{4}.6\PYG{o}{)}
solv\PYG{o}{[}*\PYG{o}{]}\PYG{o}{[}*\PYG{o}{]}\PYG{o}{[}*\PYG{o}{]}.move\PYG{o}{(}\PYGZhy{}11.5,\PYG{+w}{ }\PYGZhy{}11.5,\PYG{+w}{ }\PYGZhy{}11.5\PYG{o}{)}

\PYG{c+c1}{\PYGZsh{} Set the simulation box.}
write\PYGZus{}once\PYG{o}{(}\PYG{l+s+s2}{\PYGZdq{}Data Boundary\PYGZdq{}}\PYG{o}{)}\PYG{+w}{ }\PYG{o}{\PYGZob{}}
\PYG{+w}{ }\PYGZhy{}11.5\PYG{+w}{ }\PYG{l+m}{11}.5\PYG{+w}{ }xlo\PYG{+w}{ }xhi
\PYG{+w}{ }\PYGZhy{}11.5\PYG{+w}{ }\PYG{l+m}{11}.5\PYG{+w}{ }ylo\PYG{+w}{ }yhi
\PYG{+w}{ }\PYGZhy{}11.5\PYG{+w}{ }\PYG{l+m}{11}.5\PYG{+w}{ }zlo\PYG{+w}{ }zhi
\PYG{o}{\PYGZcb{}}

\PYG{c+c1}{\PYGZsh{} Create an input deck for LAMMPS.}
write\PYGZus{}once\PYG{o}{(}\PYG{l+s+s2}{\PYGZdq{}In Init\PYGZdq{}}\PYG{o}{)}\PYG{o}{\PYGZob{}}
\PYG{+w}{ }\PYG{c+c1}{\PYGZsh{} Input variables.}
\PYG{+w}{ }variable\PYG{+w}{ }run\PYG{+w}{    }string\PYG{+w}{ }solv\PYGZus{}01\PYG{+w}{   }\PYG{c+c1}{\PYGZsh{} output name}
\PYG{+w}{ }variable\PYG{+w}{ }ts\PYG{+w}{     }equal\PYG{+w}{  }\PYG{l+m}{1}\PYG{+w}{         }\PYG{c+c1}{\PYGZsh{} timestep}
\PYG{+w}{ }variable\PYG{+w}{ }temp\PYG{+w}{   }equal\PYG{+w}{  }\PYG{l+m}{300}\PYG{+w}{       }\PYG{c+c1}{\PYGZsh{} equilibrium temperature}
\PYG{+w}{ }variable\PYG{+w}{ }p\PYG{+w}{      }equal\PYG{+w}{  }\PYG{l+m}{1}.\PYG{+w}{        }\PYG{c+c1}{\PYGZsh{} equilibrium pressure}
\PYG{+w}{ }variable\PYG{+w}{ }d\PYG{+w}{      }equal\PYG{+w}{  }\PYG{l+m}{1000}\PYG{+w}{      }\PYG{c+c1}{\PYGZsh{} output frequency}
\PYG{+w}{ }variable\PYG{+w}{ }equi\PYG{+w}{   }equal\PYG{+w}{  }\PYG{l+m}{5000}\PYG{+w}{      }\PYG{c+c1}{\PYGZsh{} Equilibration steps}
\PYG{+w}{ }variable\PYG{+w}{ }prod\PYG{+w}{   }equal\PYG{+w}{  }\PYG{l+m}{30000}\PYG{+w}{     }\PYG{c+c1}{\PYGZsh{} Production steps}

\PYG{+w}{ }\PYG{c+c1}{\PYGZsh{} PBC (set them before the creation of the box).}
\PYG{+w}{ }boundary\PYG{+w}{ }p\PYG{+w}{ }p\PYG{+w}{ }p
\PYG{o}{\PYGZcb{}}

\PYG{c+c1}{\PYGZsh{} Run an NPT simulation.}
write\PYGZus{}once\PYG{o}{(}\PYG{l+s+s2}{\PYGZdq{}In Run\PYGZdq{}}\PYG{o}{)}\PYG{o}{\PYGZob{}}
\PYG{+w}{ }\PYG{c+c1}{\PYGZsh{} Derived variables.}
\PYG{+w}{ }variable\PYG{+w}{ }tcouple\PYG{+w}{ }equal\PYG{+w}{ }\PYG{l+s+se}{\PYGZbs{}\PYGZdl{}}\PYG{l+s+se}{\PYGZbs{}\PYGZob{}}ts\PYG{l+s+se}{\PYGZbs{}\PYGZcb{}}*100
\PYG{+w}{ }variable\PYG{+w}{ }pcouple\PYG{+w}{ }equal\PYG{+w}{ }\PYG{l+s+se}{\PYGZbs{}\PYGZdl{}}\PYG{l+s+se}{\PYGZbs{}\PYGZob{}}ts\PYG{l+s+se}{\PYGZbs{}\PYGZcb{}}*1000

\PYG{+w}{ }\PYG{c+c1}{\PYGZsh{} Output.}
\PYG{+w}{ }thermo\PYG{+w}{          }\PYG{l+s+se}{\PYGZbs{}\PYGZdl{}}d
\PYG{+w}{ }thermo\PYGZus{}style\PYG{+w}{ }custom\PYG{+w}{ }step\PYG{+w}{ }etotal\PYG{+w}{ }evdwl\PYG{+w}{ }ecoul\PYG{+w}{ }elong\PYG{+w}{ }ebond\PYG{+w}{ }eangle\PYG{+w}{ }\PYG{p}{\PYGZam{}}
\PYG{+w}{ }edihed\PYG{+w}{ }eimp\PYG{+w}{ }ke\PYG{+w}{ }pe\PYG{+w}{ }temp\PYG{+w}{ }press\PYG{+w}{ }vol\PYG{+w}{ }density\PYG{+w}{ }cpu
\PYG{+w}{ }thermo\PYGZus{}modify\PYG{+w}{ }flush\PYG{+w}{ }yes

\PYG{+w}{ }\PYG{c+c1}{\PYGZsh{} Trajectory.}
\PYG{+w}{ }dump\PYG{+w}{ }TRJ\PYG{+w}{ }all\PYG{+w}{ }dcd\PYG{+w}{ }\PYG{l+s+se}{\PYGZbs{}\PYGZdl{}}d\PYG{+w}{ }\PYG{l+s+se}{\PYGZbs{}\PYGZdl{}}\PYG{l+s+se}{\PYGZbs{}\PYGZob{}}run\PYG{l+s+se}{\PYGZbs{}\PYGZcb{}}.dcd
\PYG{+w}{ }dump\PYGZus{}modify\PYG{+w}{ }TRJ\PYG{+w}{ }unwrap\PYG{+w}{ }yes

\PYG{+w}{ }\PYG{c+c1}{\PYGZsh{} Thermalisation and relaxation, NPT ensemble.}
\PYG{+w}{ }timestep\PYG{+w}{       }\PYG{l+s+se}{\PYGZbs{}\PYGZdl{}}\PYG{l+s+se}{\PYGZbs{}\PYGZob{}}ts\PYG{l+s+se}{\PYGZbs{}\PYGZcb{}}
\PYG{+w}{ }fix\PYG{+w}{             }NPT\PYG{+w}{ }all\PYG{+w}{ }npt\PYG{+w}{ }temp\PYG{+w}{ }\PYG{l+s+se}{\PYGZbs{}\PYGZdl{}}\PYG{l+s+se}{\PYGZbs{}\PYGZob{}}temp\PYG{l+s+se}{\PYGZbs{}\PYGZcb{}}\PYG{+w}{ }\PYG{l+s+se}{\PYGZbs{}\PYGZdl{}}\PYG{l+s+se}{\PYGZbs{}\PYGZob{}}temp\PYG{l+s+se}{\PYGZbs{}\PYGZcb{}}\PYG{+w}{ }\PYG{l+s+se}{\PYGZbs{}\PYGZdl{}}\PYG{l+s+se}{\PYGZbs{}\PYGZob{}}tcouple\PYG{l+s+se}{\PYGZbs{}\PYGZcb{}}\PYG{+w}{ }iso\PYG{+w}{ }\PYG{l+s+se}{\PYGZbs{}\PYGZdl{}}p\PYG{+w}{ }\PYG{l+s+se}{\PYGZbs{}\PYGZdl{}}p\PYG{+w}{ }\PYG{l+s+se}{\PYGZbs{}\PYGZdl{}}\PYG{l+s+se}{\PYGZbs{}\PYGZob{}}pcouple\PYG{l+s+se}{\PYGZbs{}\PYGZcb{}}
\PYG{+w}{ }velocity\PYG{+w}{ }all\PYG{+w}{ }create\PYG{+w}{ }\PYG{l+s+se}{\PYGZbs{}\PYGZdl{}}\PYG{l+s+se}{\PYGZbs{}\PYGZob{}}temp\PYG{l+s+se}{\PYGZbs{}\PYGZcb{}}\PYG{+w}{ }\PYG{l+m}{858096}\PYG{+w}{ }dist\PYG{+w}{ }gaussian
\PYG{+w}{ }\PYG{c+c1}{\PYGZsh{} Short runs to update the PPPM settings as the box shinks.}
\PYG{+w}{ }run\PYG{+w}{    }\PYG{l+s+se}{\PYGZbs{}\PYGZdl{}}\PYG{l+s+se}{\PYGZbs{}\PYGZob{}}equi\PYG{l+s+se}{\PYGZbs{}\PYGZcb{}}\PYG{+w}{ }post\PYG{+w}{ }no
\PYG{+w}{ }run\PYG{+w}{    }\PYG{l+s+se}{\PYGZbs{}\PYGZdl{}}\PYG{l+s+se}{\PYGZbs{}\PYGZob{}}equi\PYG{l+s+se}{\PYGZbs{}\PYGZcb{}}\PYG{+w}{ }post\PYG{+w}{ }no
\PYG{+w}{ }run\PYG{+w}{    }\PYG{l+s+se}{\PYGZbs{}\PYGZdl{}}\PYG{l+s+se}{\PYGZbs{}\PYGZob{}}equi\PYG{l+s+se}{\PYGZbs{}\PYGZcb{}}\PYG{+w}{ }post\PYG{+w}{ }no
\PYG{+w}{ }run\PYG{+w}{    }\PYG{l+s+se}{\PYGZbs{}\PYGZdl{}}\PYG{l+s+se}{\PYGZbs{}\PYGZob{}}equi\PYG{l+s+se}{\PYGZbs{}\PYGZcb{}}
\PYG{+w}{ }\PYG{c+c1}{\PYGZsh{} From now on, the density shouldn\PYGZsq{}t change too much.}
\PYG{+w}{ }run\PYG{+w}{    }\PYG{l+s+se}{\PYGZbs{}\PYGZdl{}}\PYG{l+s+se}{\PYGZbs{}\PYGZob{}}prod\PYG{l+s+se}{\PYGZbs{}\PYGZcb{}}
\PYG{+w}{ }unfix\PYG{+w}{ }NPT
\PYG{o}{\PYGZcb{}}
\end{sphinxVerbatim}

\sphinxAtStartPar
The first two commands insert the content of files \sphinxcode{\sphinxupquote{oplsaa.lt}} and
\sphinxcode{\sphinxupquote{formamide.lt}} into the master file. At this point, we can use the
command \sphinxcode{\sphinxupquote{solv = new \_FAM {[}N{]}}} to create N copies of a molecule of type
\sphinxcode{\sphinxupquote{\_FAM}}. In this case, we create an array of 5*5*5 molecules on a cubic
grid using the coordinate transformation command \sphinxcode{\sphinxupquote{.move( 4.6, 0, 0)}}.
See the Moltemplate documentation to learn more about the syntax. As
the sample was created from scratch, we also specify the simulation box
size in the “Data Boundary” section.

\sphinxAtStartPar
The LAMMPS setting for the force field are specified in the file
\sphinxcode{\sphinxupquote{oplsaa.lt}} and are written automatically in the input deck. We also
specify the boundary conditions and a set of variables in
the “In Init” section. The remaining commands to run an NPT simulation
are written in the “In Run” section. Note that in this script, LAMMPS
variables are protected with the escape character \sphinxcode{\sphinxupquote{\textbackslash{}}} to distinguish
them from Moltemplate variables, e.g. \sphinxcode{\sphinxupquote{\textbackslash{}\$\textbackslash{}\{run\textbackslash{}\}}} is a LAMMPS
variable that is written in the input deck as \sphinxcode{\sphinxupquote{\$\{run\}}}.

\sphinxAtStartPar
Compile the master file with:

\begin{sphinxVerbatim}[commandchars=\\\{\}]
moltemplate.sh\PYG{+w}{ }\PYGZhy{}overlay\PYGZhy{}all\PYG{+w}{ }solv\PYGZus{}01.lt
\end{sphinxVerbatim}

\sphinxAtStartPar
And execute the simulation with the following:

\begin{sphinxVerbatim}[commandchars=\\\{\}]
mpirun\PYG{+w}{ }\PYGZhy{}np\PYG{+w}{ }\PYG{l+m}{4}\PYG{+w}{ }lmp\PYG{+w}{ }\PYGZhy{}in\PYG{+w}{ }solv\PYGZus{}01.in\PYG{+w}{ }\PYGZhy{}l\PYG{+w}{ }solv\PYGZus{}01.log
\end{sphinxVerbatim}

\begin{figure}[htbp]
\centering
\capstart

\noindent\sphinxincludegraphics{{solv_01}.png}
\caption{Snapshot of the sample at the beginning and end of the simulation.
Rendered with Ovito.}\label{\detokenize{Howto_moltemplate:id1}}\end{figure}


\subsubsection{Mapping an existing structure}
\label{\detokenize{Howto_moltemplate:mapping-an-existing-structure}}
\sphinxAtStartPar
Another helpful way to use Moltemplate is mapping an existing molecular
sample to a force field. This is useful when a complex sample is
assembled from different simulations or created with specialized
software (e.g. PACKMOL). As in the previous example, all molecular
species in the sample must be defined using single\sphinxhyphen{}molecule Moltemplate
objects.  For this example, we use a short polymer in a box containing
water molecules and ions in the PDB file \sphinxcode{\sphinxupquote{model.pdb}}.

\sphinxAtStartPar
It is essential to understand that the order of atoms in the PDB file
and in the Moltemplate master script must match, as we are using the
coordinates from the PDB file in the order they appear. The order of
atoms and molecules in the PDB file provided is as follows:
\begin{itemize}
\item {} 
\sphinxAtStartPar
500 water molecules, with atoms ordered in this sequence:

\begin{sphinxVerbatim}[commandchars=\\\{\}]
ATOM      1  O   MOL D   1       5.901   7.384   1.103  0.00  0.00      DUM
ATOM      2  H   MOL D   1       6.047   8.238   0.581  0.00  0.00      DUM
ATOM      3  H   MOL D   1       6.188   7.533   2.057  0.00  0.00      DUM
\end{sphinxVerbatim}

\item {} 
\sphinxAtStartPar
1 polymer molecule.

\item {} 
\sphinxAtStartPar
1 Ca$^{\text{2+}}$ ion.

\item {} 
\sphinxAtStartPar
2 Cl$^{\text{\sphinxhyphen{}}}$ ions.

\end{itemize}

\sphinxAtStartPar
In the master LT file, this sequence of molecules is matched with the
following commands:

\begin{sphinxVerbatim}[commandchars=\\\{\}]
\PYG{c+c1}{\PYGZsh{} Create the sample.}
\PYG{n+nv}{wat}\PYG{o}{=}new\PYG{+w}{ }SPC\PYG{o}{[}\PYG{l+m}{500}\PYG{o}{]}
\PYG{n+nv}{pol}\PYG{o}{=}new\PYG{+w}{ }PolyNIPAM\PYG{o}{[}\PYG{l+m}{1}\PYG{o}{]}
\PYG{n+nv}{cat}\PYG{o}{=}new\PYG{+w}{ }Ca\PYG{o}{[}\PYG{l+m}{1}\PYG{o}{]}
\PYG{n+nv}{ani}\PYG{o}{=}new\PYG{+w}{ }Cl\PYG{o}{[}\PYG{l+m}{2}\PYG{o}{]}
\end{sphinxVerbatim}

\sphinxAtStartPar
Note that the first command would create 500 water molecules in the
same position in space, and the other commands will use the coordinates
specified in the corresponding molecular topology block. However, the
coordinates will be overwritten by rendering an external atomic
structure file. Note that if the same molecule species are scattered in
the input structure, it is recommended to reorder and group together
for molecule types to facilitate the creation of the input sample.

\sphinxAtStartPar
The molecular topology for the polymer is created as in the previous
example, with the atom types assigned as in the following schema:

\begin{figure}[htbp]
\centering
\capstart

\noindent\sphinxincludegraphics[scale=0.3]{{PolyNIPAM}.jpg}
\caption{Atom types assigned to the polymer’s repeating unit.}\label{\detokenize{Howto_moltemplate:id2}}\end{figure}

\sphinxAtStartPar
The molecular topology of the water and ions is stated directly into
the master file for the sake of space, but they could also be written
in a separate file(s) and imported before the sample is created.

\sphinxAtStartPar
The resulting master LT file defining short annealing at a fixed volume
(NVT) follows:

\begin{sphinxVerbatim}[commandchars=\\\{\}]
\PYG{c+c1}{\PYGZsh{} Use the OPLS\PYGZhy{}AA force field for all species.}
import\PYG{+w}{ }/usr/local/moltemplate/moltemplate/force\PYGZus{}fields/oplsaa.lt
import\PYG{+w}{ }PolyNIPAM.lt

\PYG{c+c1}{\PYGZsh{} Define the SPC water and ions as in the OPLS\PYGZhy{}AA}
Ca\PYG{+w}{ }inherits\PYG{+w}{ }OPLSAA\PYG{+w}{ }\PYG{o}{\PYGZob{}}
\PYG{+w}{  }write\PYG{o}{(}\PYG{l+s+s2}{\PYGZdq{}Data Atoms\PYGZdq{}}\PYG{o}{)}\PYG{o}{\PYGZob{}}
\PYG{+w}{    }\PYG{n+nv}{\PYGZdl{}atom}:a1\PYG{+w}{  }\PYG{n+nv}{\PYGZdl{}mol}:.\PYG{+w}{ }@atom:354\PYG{+w}{ }\PYG{l+m}{0}.0\PYG{+w}{  }\PYG{l+m}{0}.00000\PYG{+w}{ }\PYG{l+m}{0}.00000\PYG{+w}{ }\PYG{l+m}{0}.000000
\PYG{+w}{  }\PYG{o}{\PYGZcb{}}
\PYG{o}{\PYGZcb{}}
Cl\PYG{+w}{ }inherits\PYG{+w}{ }OPLSAA\PYG{+w}{ }\PYG{o}{\PYGZob{}}
\PYG{+w}{  }write\PYG{o}{(}\PYG{l+s+s2}{\PYGZdq{}Data Atoms\PYGZdq{}}\PYG{o}{)}\PYG{o}{\PYGZob{}}
\PYG{+w}{    }\PYG{n+nv}{\PYGZdl{}atom}:a1\PYG{+w}{  }\PYG{n+nv}{\PYGZdl{}mol}:.\PYG{+w}{ }@atom:344\PYG{+w}{ }\PYG{l+m}{0}.0\PYG{+w}{  }\PYG{l+m}{0}.00000\PYG{+w}{ }\PYG{l+m}{0}.00000\PYG{+w}{ }\PYG{l+m}{0}.000000
\PYG{+w}{  }\PYG{o}{\PYGZcb{}}
\PYG{o}{\PYGZcb{}}
SPC\PYG{+w}{ }inherits\PYG{+w}{ }OPLSAA\PYG{+w}{ }\PYG{o}{\PYGZob{}}
\PYG{+w}{  }write\PYG{o}{(}\PYG{l+s+s2}{\PYGZdq{}Data Atoms\PYGZdq{}}\PYG{o}{)}\PYG{o}{\PYGZob{}}
\PYG{+w}{    }\PYG{n+nv}{\PYGZdl{}atom}:O\PYG{+w}{  }\PYG{n+nv}{\PYGZdl{}mol}:.\PYG{+w}{ }@atom:76\PYG{+w}{ }\PYG{l+m}{0}.\PYG{+w}{  }\PYG{l+m}{0}.0000000\PYG{+w}{ }\PYG{l+m}{0}.00000\PYG{+w}{ }\PYG{l+m}{0}.000000
\PYG{+w}{    }\PYG{n+nv}{\PYGZdl{}atom}:H1\PYG{+w}{ }\PYG{n+nv}{\PYGZdl{}mol}:.\PYG{+w}{ }@atom:77\PYG{+w}{ }\PYG{l+m}{0}.\PYG{+w}{  }\PYG{l+m}{0}.8164904\PYG{+w}{ }\PYG{l+m}{0}.00000\PYG{+w}{ }\PYG{l+m}{0}.5773590
\PYG{+w}{    }\PYG{n+nv}{\PYGZdl{}atom}:H2\PYG{+w}{ }\PYG{n+nv}{\PYGZdl{}mol}:.\PYG{+w}{ }@atom:77\PYG{+w}{ }\PYG{l+m}{0}.\PYG{+w}{ }\PYGZhy{}0.8164904\PYG{+w}{ }\PYG{l+m}{0}.00000\PYG{+w}{ }\PYG{l+m}{0}.5773590
\PYG{+w}{  }\PYG{o}{\PYGZcb{}}
\PYG{+w}{  }write\PYG{o}{(}\PYG{l+s+s2}{\PYGZdq{}Data Bond List\PYGZdq{}}\PYG{o}{)}\PYG{+w}{ }\PYG{o}{\PYGZob{}}
\PYG{+w}{    }\PYG{n+nv}{\PYGZdl{}bond}:OH1\PYG{+w}{ }\PYG{n+nv}{\PYGZdl{}atom}:O\PYG{+w}{ }\PYG{n+nv}{\PYGZdl{}atom}:H1
\PYG{+w}{    }\PYG{n+nv}{\PYGZdl{}bond}:OH2\PYG{+w}{ }\PYG{n+nv}{\PYGZdl{}atom}:O\PYG{+w}{ }\PYG{n+nv}{\PYGZdl{}atom}:H2
\PYG{+w}{  }\PYG{o}{\PYGZcb{}}
\PYG{o}{\PYGZcb{}}

\PYG{c+c1}{\PYGZsh{} Create the sample.}
\PYG{n+nv}{wat}\PYG{o}{=}new\PYG{+w}{ }SPC\PYG{o}{[}\PYG{l+m}{500}\PYG{o}{]}
\PYG{n+nv}{pol}\PYG{o}{=}new\PYG{+w}{ }PolyNIPAM\PYG{o}{[}\PYG{l+m}{1}\PYG{o}{]}
\PYG{n+nv}{cat}\PYG{o}{=}new\PYG{+w}{ }Ca\PYG{o}{[}\PYG{l+m}{1}\PYG{o}{]}
\PYG{n+nv}{ani}\PYG{o}{=}new\PYG{+w}{ }Cl\PYG{o}{[}\PYG{l+m}{2}\PYG{o}{]}

\PYG{c+c1}{\PYGZsh{} Periodic boundary conditions:}
write\PYGZus{}once\PYG{o}{(}\PYG{l+s+s2}{\PYGZdq{}Data Boundary\PYGZdq{}}\PYG{o}{)}\PYG{o}{\PYGZob{}}
\PYG{+w}{  }\PYG{l+m}{0}\PYG{+w}{ }\PYG{l+m}{26}\PYG{+w}{ }xlo\PYG{+w}{ }xhi
\PYG{+w}{  }\PYG{l+m}{0}\PYG{+w}{ }\PYG{l+m}{26}\PYG{+w}{ }ylo\PYG{+w}{ }yhi
\PYG{+w}{  }\PYG{l+m}{0}\PYG{+w}{ }\PYG{l+m}{26}\PYG{+w}{ }zlo\PYG{+w}{ }zhi
\PYG{o}{\PYGZcb{}}

\PYG{c+c1}{\PYGZsh{} Define the input variables.}
write\PYGZus{}once\PYG{o}{(}\PYG{l+s+s2}{\PYGZdq{}In Init\PYGZdq{}}\PYG{o}{)}\PYG{o}{\PYGZob{}}
\PYG{+w}{  }\PYG{c+c1}{\PYGZsh{} Input variables.}
\PYG{+w}{  }variable\PYG{+w}{ }run\PYG{+w}{    }string\PYG{+w}{ }sample01\PYG{+w}{  }\PYG{c+c1}{\PYGZsh{} output name}
\PYG{+w}{  }variable\PYG{+w}{ }ts\PYG{+w}{     }equal\PYG{+w}{  }\PYG{l+m}{2}\PYG{+w}{         }\PYG{c+c1}{\PYGZsh{} timestep}
\PYG{+w}{  }variable\PYG{+w}{ }temp\PYG{+w}{   }equal\PYG{+w}{  }\PYG{l+m}{298}.15\PYG{+w}{    }\PYG{c+c1}{\PYGZsh{} equilibrium temperature}
\PYG{+w}{  }variable\PYG{+w}{ }p\PYG{+w}{      }equal\PYG{+w}{  }\PYG{l+m}{1}.\PYG{+w}{        }\PYG{c+c1}{\PYGZsh{} equilibrium pressure}
\PYG{+w}{  }variable\PYG{+w}{ }equi\PYG{+w}{   }equal\PYG{+w}{  }\PYG{l+m}{30000}\PYG{+w}{     }\PYG{c+c1}{\PYGZsh{} equilibration steps}

\PYG{+w}{  }\PYG{c+c1}{\PYGZsh{} PBC (set them before the creation of the box).}
\PYG{+w}{  }boundary\PYG{+w}{ }p\PYG{+w}{ }p\PYG{+w}{ }p
\PYG{+w}{  }neighbor\PYG{+w}{        }\PYG{l+m}{3}\PYG{+w}{ }bin
\PYG{o}{\PYGZcb{}}

\PYG{c+c1}{\PYGZsh{} Run an NVT simulation.}
write\PYGZus{}once\PYG{o}{(}\PYG{l+s+s2}{\PYGZdq{}In Run\PYGZdq{}}\PYG{o}{)}\PYG{o}{\PYGZob{}}
\PYG{+w}{  }\PYG{c+c1}{\PYGZsh{} Set the output.}
\PYG{+w}{  }thermo\PYG{+w}{          }\PYG{l+m}{1000}
\PYG{+w}{  }thermo\PYGZus{}style\PYG{+w}{    }custom\PYG{+w}{ }step\PYG{+w}{ }etotal\PYG{+w}{ }evdwl\PYG{+w}{ }ecoul\PYG{+w}{ }elong\PYG{+w}{ }ebond\PYG{+w}{ }eangle\PYG{+w}{ }\PYG{p}{\PYGZam{}}
\PYG{+w}{  }edihed\PYG{+w}{ }eimp\PYG{+w}{ }pe\PYG{+w}{ }ke\PYG{+w}{ }temp\PYG{+w}{ }press\PYG{+w}{ }atoms\PYG{+w}{ }vol\PYG{+w}{ }density\PYG{+w}{ }cpu
\PYG{+w}{  }thermo\PYGZus{}modify\PYG{+w}{ }flush\PYG{+w}{ }yes
\PYG{+w}{  }compute\PYG{+w}{ }pe1\PYG{+w}{ }all\PYG{+w}{ }pe/atom\PYG{+w}{ }pair
\PYG{+w}{  }dump\PYG{+w}{ }TRJ\PYG{+w}{ }all\PYG{+w}{ }custom\PYG{+w}{ }\PYG{l+m}{100}\PYG{+w}{ }\PYG{l+s+se}{\PYGZbs{}\PYGZdl{}}\PYG{l+s+se}{\PYGZbs{}\PYGZob{}}run\PYG{l+s+se}{\PYGZbs{}\PYGZcb{}}.dump\PYG{+w}{ }id\PYG{+w}{ }xu\PYG{+w}{ }yu\PYG{+w}{ }zu\PYG{+w}{ }c\PYGZus{}pe1

\PYG{+w}{  }\PYG{c+c1}{\PYGZsh{} Minimise the input structure, just in case.}
\PYG{+w}{  }minimize\PYG{+w}{        }.01\PYG{+w}{ }.001\PYG{+w}{ }\PYG{l+m}{1000}\PYG{+w}{ }\PYG{l+m}{100000}
\PYG{+w}{  }write\PYGZus{}data\PYG{+w}{ }\PYG{l+s+se}{\PYGZbs{}\PYGZdl{}}\PYG{l+s+se}{\PYGZbs{}\PYGZob{}}run\PYG{l+s+se}{\PYGZbs{}\PYGZcb{}}.min

\PYG{+w}{  }\PYG{c+c1}{\PYGZsh{} Set the constrains.}
\PYG{+w}{  }group\PYG{+w}{ }watergroup\PYG{+w}{ }\PYG{n+nb}{type}\PYG{+w}{ }@atom:76\PYG{+w}{ }@atom:77
\PYG{+w}{  }fix\PYG{+w}{ }\PYG{l+m}{0}\PYG{+w}{ }watergroup\PYG{+w}{ }shake\PYG{+w}{ }\PYG{l+m}{0}.0001\PYG{+w}{ }\PYG{l+m}{10}\PYG{+w}{ }\PYG{l+m}{0}\PYG{+w}{ }b\PYG{+w}{ }@bond:042\PYGZus{}043\PYG{+w}{ }a\PYG{+w}{ }@angle:043\PYGZus{}042\PYGZus{}043

\PYG{+w}{  }\PYG{c+c1}{\PYGZsh{} Short annealing.}
\PYG{+w}{  }timestep\PYG{+w}{        }\PYG{l+s+se}{\PYGZbs{}\PYGZdl{}}\PYG{l+s+se}{\PYGZbs{}\PYGZob{}}ts\PYG{l+s+se}{\PYGZbs{}\PYGZcb{}}
\PYG{+w}{  }fix\PYG{+w}{       }\PYG{l+m}{1}\PYG{+w}{ }all\PYG{+w}{ }nvt\PYG{+w}{ }temp\PYG{+w}{ }\PYG{l+s+se}{\PYGZbs{}\PYGZdl{}}\PYG{l+s+se}{\PYGZbs{}\PYGZob{}}temp\PYG{l+s+se}{\PYGZbs{}\PYGZcb{}}\PYG{+w}{ }\PYG{l+s+se}{\PYGZbs{}\PYGZdl{}}\PYG{l+s+se}{\PYGZbs{}\PYGZob{}}temp\PYG{l+s+se}{\PYGZbs{}\PYGZcb{}}\PYG{+w}{ }\PYG{l+s+se}{\PYGZbs{}\PYGZdl{}}\PYG{o}{(}\PYG{l+m}{100}*dt\PYG{o}{)}
\PYG{+w}{  }velocity\PYG{+w}{    }all\PYG{+w}{ }create\PYG{+w}{ }\PYG{l+s+se}{\PYGZbs{}\PYGZdl{}}\PYG{l+s+se}{\PYGZbs{}\PYGZob{}}temp\PYG{l+s+se}{\PYGZbs{}\PYGZcb{}}\PYG{+w}{ }\PYG{l+m}{315443}
\PYG{+w}{  }run\PYG{+w}{             }\PYG{l+s+se}{\PYGZbs{}\PYGZdl{}}\PYG{l+s+se}{\PYGZbs{}\PYGZob{}}equi\PYG{l+s+se}{\PYGZbs{}\PYGZcb{}}
\PYG{+w}{  }unfix\PYG{+w}{ }\PYG{l+m}{1}
\PYG{o}{\PYGZcb{}}
\end{sphinxVerbatim}

\sphinxAtStartPar
In this example, the water model is SPC and it is defined in the
\sphinxcode{\sphinxupquote{oplsaa.lt}} file with atom types \sphinxcode{\sphinxupquote{@atom:76}} and \sphinxcode{\sphinxupquote{@atom:77}}.  For
water we also use the \sphinxcode{\sphinxupquote{group}} and \sphinxcode{\sphinxupquote{fix shake}} commands with
Moltemplate \sphinxcode{\sphinxupquote{@}}\sphinxhyphen{}type variables, to ensure consistency with the
numerical values assigned during compilation. To identify the bond and
angle types, look for the extended \sphinxcode{\sphinxupquote{@atom}} IDs, which in this case
are:

\begin{sphinxVerbatim}[commandchars=\\\{\}]
replace\PYG{o}{\PYGZob{}}\PYG{+w}{ }@atom:76\PYG{+w}{ }@atom:76\PYGZus{}b042\PYGZus{}a042\PYGZus{}d042\PYGZus{}i042\PYG{+w}{ }\PYG{o}{\PYGZcb{}}
replace\PYG{o}{\PYGZob{}}\PYG{+w}{ }@atom:77\PYG{+w}{ }@atom:77\PYGZus{}b043\PYGZus{}a043\PYGZus{}d043\PYGZus{}i043\PYG{+w}{ }\PYG{o}{\PYGZcb{}}
\end{sphinxVerbatim}

\sphinxAtStartPar
From which we can identify the following “Data Bonds By Type”:
\sphinxcode{\sphinxupquote{@bond:042\_043 @atom:*\_b042*\_a*\_d*\_i* @atom:*\_b043*\_a*\_d*\_i*}} and
“Data Angles By Type”: \sphinxcode{\sphinxupquote{@angle:043\_042\_043 @atom:*\_b*\_a043*\_d*\_i*
@atom:*\_b*\_a042*\_d*\_i* @atom:*\_b*\_a043*\_d*\_i*}}

\sphinxAtStartPar
Compile the master file with:

\begin{sphinxVerbatim}[commandchars=\\\{\}]
moltemplate.sh\PYG{+w}{ }\PYGZhy{}overlay\PYGZhy{}all\PYG{+w}{ }\PYGZhy{}pdb\PYG{+w}{ }model.pdb\PYG{+w}{ }sample01.lt
\end{sphinxVerbatim}

\sphinxAtStartPar
And execute the simulation with the following:

\begin{sphinxVerbatim}[commandchars=\\\{\}]
mpirun\PYG{+w}{ }\PYGZhy{}np\PYG{+w}{ }\PYG{l+m}{4}\PYG{+w}{ }lmp\PYG{+w}{ }\PYGZhy{}in\PYG{+w}{ }sample01.in\PYG{+w}{ }\PYGZhy{}l\PYG{+w}{ }sample01.log
\end{sphinxVerbatim}

\begin{figure}[htbp]
\centering
\capstart

\noindent\sphinxincludegraphics{{sample01}.png}
\caption{Sample visualized with Ovito loading the trajectory into the DATA
file written after minimization.}\label{\detokenize{Howto_moltemplate:id3}}\end{figure}


\bigskip\hrule\bigskip

\phantomsection\label{\detokenize{Howto_moltemplate:oplsaa96}}
\sphinxAtStartPar
\sphinxstylestrong{(OPLS\sphinxhyphen{}AA)}  Jorgensen, Maxwell, Tirado\sphinxhyphen{}Rives, J Am Chem Soc, 118(45), 11225\sphinxhyphen{}11236 (1996).

\sphinxstepscope


\subsection{PyLammps Tutorial}
\label{\detokenize{Howto_pylammps:pylammps-tutorial}}\label{\detokenize{Howto_pylammps::doc}}
\begin{sphinxShadowBox}
\sphinxstyletopictitle{Contents}
\begin{itemize}
\item {} 
\sphinxAtStartPar
\phantomsection\label{\detokenize{Howto_pylammps:id2}}{\hyperref[\detokenize{Howto_pylammps:pylammps-tutorial}]{\sphinxcrossref{PyLammps Tutorial}}}
\begin{itemize}
\item {} 
\sphinxAtStartPar
\phantomsection\label{\detokenize{Howto_pylammps:id3}}{\hyperref[\detokenize{Howto_pylammps:overview}]{\sphinxcrossref{Overview}}}
\begin{itemize}
\item {} 
\sphinxAtStartPar
\phantomsection\label{\detokenize{Howto_pylammps:id4}}{\hyperref[\detokenize{Howto_pylammps:comparison-of-lammps-and-pylammps-interfaces}]{\sphinxcrossref{Comparison of lammps and PyLammps interfaces}}}
\begin{itemize}
\item {} 
\sphinxAtStartPar
\phantomsection\label{\detokenize{Howto_pylammps:id5}}{\hyperref[\detokenize{Howto_pylammps:lammps-lammps}]{\sphinxcrossref{lammps.lammps}}}

\item {} 
\sphinxAtStartPar
\phantomsection\label{\detokenize{Howto_pylammps:id6}}{\hyperref[\detokenize{Howto_pylammps:lammps-pylammps}]{\sphinxcrossref{lammps.PyLammps}}}

\end{itemize}

\end{itemize}

\item {} 
\sphinxAtStartPar
\phantomsection\label{\detokenize{Howto_pylammps:id7}}{\hyperref[\detokenize{Howto_pylammps:quick-start}]{\sphinxcrossref{Quick Start}}}
\begin{itemize}
\item {} 
\sphinxAtStartPar
\phantomsection\label{\detokenize{Howto_pylammps:id8}}{\hyperref[\detokenize{Howto_pylammps:system-wide-installation}]{\sphinxcrossref{System\sphinxhyphen{}wide Installation}}}
\begin{itemize}
\item {} 
\sphinxAtStartPar
\phantomsection\label{\detokenize{Howto_pylammps:id9}}{\hyperref[\detokenize{Howto_pylammps:step-1-building-lammps-as-a-shared-library}]{\sphinxcrossref{Step 1: Building LAMMPS as a shared library}}}

\item {} 
\sphinxAtStartPar
\phantomsection\label{\detokenize{Howto_pylammps:id10}}{\hyperref[\detokenize{Howto_pylammps:step-2-installing-the-lammps-python-package}]{\sphinxcrossref{Step 2: Installing the LAMMPS Python package}}}

\end{itemize}

\item {} 
\sphinxAtStartPar
\phantomsection\label{\detokenize{Howto_pylammps:id11}}{\hyperref[\detokenize{Howto_pylammps:installation-inside-of-a-virtualenv}]{\sphinxcrossref{Installation inside of a virtualenv}}}
\begin{itemize}
\item {} 
\sphinxAtStartPar
\phantomsection\label{\detokenize{Howto_pylammps:id12}}{\hyperref[\detokenize{Howto_pylammps:benefits-of-using-a-virtualenv}]{\sphinxcrossref{Benefits of using a virtualenv}}}

\item {} 
\sphinxAtStartPar
\phantomsection\label{\detokenize{Howto_pylammps:id13}}{\hyperref[\detokenize{Howto_pylammps:creating-a-virtualenv-with-lammps-installed}]{\sphinxcrossref{Creating a virtualenv with lammps installed}}}

\end{itemize}

\end{itemize}

\item {} 
\sphinxAtStartPar
\phantomsection\label{\detokenize{Howto_pylammps:id14}}{\hyperref[\detokenize{Howto_pylammps:creating-a-new-instance-of-pylammps}]{\sphinxcrossref{Creating a new instance of PyLammps}}}

\item {} 
\sphinxAtStartPar
\phantomsection\label{\detokenize{Howto_pylammps:id15}}{\hyperref[\detokenize{Howto_pylammps:commands}]{\sphinxcrossref{Commands}}}

\item {} 
\sphinxAtStartPar
\phantomsection\label{\detokenize{Howto_pylammps:id16}}{\hyperref[\detokenize{Howto_pylammps:system-state}]{\sphinxcrossref{System state}}}

\item {} 
\sphinxAtStartPar
\phantomsection\label{\detokenize{Howto_pylammps:id17}}{\hyperref[\detokenize{Howto_pylammps:working-with-lammps-variables}]{\sphinxcrossref{Working with LAMMPS variables}}}

\item {} 
\sphinxAtStartPar
\phantomsection\label{\detokenize{Howto_pylammps:id18}}{\hyperref[\detokenize{Howto_pylammps:retrieving-the-value-of-an-arbitrary-lammps-expressions}]{\sphinxcrossref{Retrieving the value of an arbitrary LAMMPS expressions}}}

\item {} 
\sphinxAtStartPar
\phantomsection\label{\detokenize{Howto_pylammps:id19}}{\hyperref[\detokenize{Howto_pylammps:accessing-atom-data}]{\sphinxcrossref{Accessing atom data}}}

\item {} 
\sphinxAtStartPar
\phantomsection\label{\detokenize{Howto_pylammps:id20}}{\hyperref[\detokenize{Howto_pylammps:evaluating-thermo-data}]{\sphinxcrossref{Evaluating thermo data}}}

\item {} 
\sphinxAtStartPar
\phantomsection\label{\detokenize{Howto_pylammps:id21}}{\hyperref[\detokenize{Howto_pylammps:error-handling-with-pylammps}]{\sphinxcrossref{Error handling with PyLammps}}}

\item {} 
\sphinxAtStartPar
\phantomsection\label{\detokenize{Howto_pylammps:id22}}{\hyperref[\detokenize{Howto_pylammps:using-pylammps-in-ipython-notebooks-and-jupyter}]{\sphinxcrossref{Using PyLammps in IPython notebooks and Jupyter}}}

\item {} 
\sphinxAtStartPar
\phantomsection\label{\detokenize{Howto_pylammps:id23}}{\hyperref[\detokenize{Howto_pylammps:ipylammps-examples}]{\sphinxcrossref{IPyLammps Examples}}}
\begin{itemize}
\item {} 
\sphinxAtStartPar
\phantomsection\label{\detokenize{Howto_pylammps:id24}}{\hyperref[\detokenize{Howto_pylammps:validating-a-dihedral-potential}]{\sphinxcrossref{Validating a dihedral potential}}}

\item {} 
\sphinxAtStartPar
\phantomsection\label{\detokenize{Howto_pylammps:id25}}{\hyperref[\detokenize{Howto_pylammps:running-a-monte-carlo-relaxation}]{\sphinxcrossref{Running a Monte Carlo relaxation}}}

\end{itemize}

\item {} 
\sphinxAtStartPar
\phantomsection\label{\detokenize{Howto_pylammps:id26}}{\hyperref[\detokenize{Howto_pylammps:using-pylammps-and-mpi4py-experimental}]{\sphinxcrossref{Using PyLammps and mpi4py (Experimental)}}}

\item {} 
\sphinxAtStartPar
\phantomsection\label{\detokenize{Howto_pylammps:id27}}{\hyperref[\detokenize{Howto_pylammps:feedback-and-contributing}]{\sphinxcrossref{Feedback and Contributing}}}

\end{itemize}

\end{itemize}
\end{sphinxShadowBox}


\subsubsection{Overview}
\label{\detokenize{Howto_pylammps:overview}}
\sphinxAtStartPar
{\hyperref[\detokenize{Python_module:lammps.PyLammps}]{\sphinxcrossref{\sphinxcode{\sphinxupquote{PyLammps}}}}} is a Python wrapper class for
LAMMPS which can be created on its own or use an existing
{\hyperref[\detokenize{Python_module:lammps.lammps}]{\sphinxcrossref{\sphinxcode{\sphinxupquote{lammps Python}}}}} object.  It creates a simpler,
more “pythonic” interface to common LAMMPS functionality, in contrast to
the {\hyperref[\detokenize{Python_module:lammps.lammps}]{\sphinxcrossref{\sphinxcode{\sphinxupquote{lammps}}}}} wrapper for the LAMMPS {\hyperref[\detokenize{Library:lammps-c-api}]{\sphinxcrossref{\DUrole{std,std-ref}{C
language library interface API}}}} which is written using
\sphinxhref{https://docs.python.org/3/library/ctypes.html}{Python ctypes}.  The {\hyperref[\detokenize{Python_module:lammps.lammps}]{\sphinxcrossref{\sphinxcode{\sphinxupquote{lammps}}}}}
wrapper is discussed on the {\hyperref[\detokenize{Python_head::doc}]{\sphinxcrossref{\DUrole{doc}{Use Python with LAMMPS}}}} doc page.

\sphinxAtStartPar
Unlike the flat \sphinxhref{https://docs.python.org/3/library/ctypes.html}{ctypes} interface, PyLammps exposes a
discoverable API.  It no longer requires knowledge of the underlying C++
code implementation.  Finally, the {\hyperref[\detokenize{Python_module:lammps.IPyLammps}]{\sphinxcrossref{\sphinxcode{\sphinxupquote{IPyLammps}}}}} wrapper builds on top of {\hyperref[\detokenize{Python_module:lammps.PyLammps}]{\sphinxcrossref{\sphinxcode{\sphinxupquote{PyLammps}}}}} and adds some additional features for \sphinxhref{https://ipython.org/}{IPython
integration} into \sphinxhref{https://jupyter.org/}{Jupyter notebooks}, e.g. for
embedded visualization output from {\hyperref[\detokenize{dump_image::doc}]{\sphinxcrossref{\DUrole{doc}{dump style image}}}}.


\paragraph{Comparison of lammps and PyLammps interfaces}
\label{\detokenize{Howto_pylammps:comparison-of-lammps-and-pylammps-interfaces}}

\subparagraph{lammps.lammps}
\label{\detokenize{Howto_pylammps:lammps-lammps}}\begin{itemize}
\item {} 
\sphinxAtStartPar
uses \sphinxhref{https://docs.python.org/3/library/ctypes.html}{ctypes}

\item {} 
\sphinxAtStartPar
direct memory access to native C++ data with optional support for NumPy arrays

\item {} 
\sphinxAtStartPar
provides functions to send and receive data to LAMMPS

\item {} 
\sphinxAtStartPar
interface modeled after the LAMMPS {\hyperref[\detokenize{Library:lammps-c-api}]{\sphinxcrossref{\DUrole{std,std-ref}{C language library interface API}}}}

\item {} 
\sphinxAtStartPar
requires knowledge of how LAMMPS internally works (C pointers, etc)

\item {} 
\sphinxAtStartPar
full support for running Python with MPI using \sphinxhref{https://mpi4py.readthedocs.io}{mpi4py}

\end{itemize}


\subparagraph{lammps.PyLammps}
\label{\detokenize{Howto_pylammps:lammps-pylammps}}\begin{itemize}
\item {} 
\sphinxAtStartPar
higher\sphinxhyphen{}level abstraction built on \sphinxstyleemphasis{top} of original {\hyperref[\detokenize{Python_module:lammps.lammps}]{\sphinxcrossref{\sphinxcode{\sphinxupquote{ctypes based interface}}}}}

\item {} 
\sphinxAtStartPar
manipulation of Python objects

\item {} 
\sphinxAtStartPar
communication with LAMMPS is hidden from API user

\item {} 
\sphinxAtStartPar
shorter, more concise Python

\item {} 
\sphinxAtStartPar
better IPython integration, designed for quick prototyping

\item {} 
\sphinxAtStartPar
designed for serial execution

\end{itemize}


\subsubsection{Quick Start}
\label{\detokenize{Howto_pylammps:quick-start}}

\paragraph{System\sphinxhyphen{}wide Installation}
\label{\detokenize{Howto_pylammps:system-wide-installation}}

\subparagraph{Step 1: Building LAMMPS as a shared library}
\label{\detokenize{Howto_pylammps:step-1-building-lammps-as-a-shared-library}}
\sphinxAtStartPar
To use LAMMPS inside of Python it has to be compiled as shared
library. This library is then loaded by the Python interface. In this
example we enable the MOLECULE package and compile LAMMPS with PNG, JPEG
and FFMPEG output support enabled.

\sphinxAtStartPar
Step 1a: For the CMake based build system, the steps are:

\begin{sphinxVerbatim}[commandchars=\\\{\}]
mkdir\PYG{+w}{ }\PYG{n+nv}{\PYGZdl{}LAMMPS\PYGZus{}DIR}/build\PYGZhy{}shared
\PYG{n+nb}{cd}\PYG{+w}{  }\PYG{n+nv}{\PYGZdl{}LAMMPS\PYGZus{}DIR}/build\PYGZhy{}shared

\PYG{c+c1}{\PYGZsh{} MPI, PNG, Jpeg, FFMPEG are auto\PYGZhy{}detected}
cmake\PYG{+w}{ }../cmake\PYG{+w}{ }\PYGZhy{}DPKG\PYGZus{}MOLECULE\PYG{o}{=}yes\PYG{+w}{ }\PYGZhy{}DBUILD\PYGZus{}LIB\PYG{o}{=}yes\PYG{+w}{ }\PYGZhy{}DBUILD\PYGZus{}SHARED\PYGZus{}LIBS\PYG{o}{=}yes
make
\end{sphinxVerbatim}

\sphinxAtStartPar
Step 1b: For the legacy, make based build system, the steps are:

\begin{sphinxVerbatim}[commandchars=\\\{\}]
\PYG{n+nb}{cd}\PYG{+w}{ }\PYG{n+nv}{\PYGZdl{}LAMMPS\PYGZus{}DIR}/src

\PYG{c+c1}{\PYGZsh{} add packages if necessary}
make\PYG{+w}{ }yes\PYGZhy{}MOLECULE

\PYG{c+c1}{\PYGZsh{} compile shared library using Makefile}
make\PYG{+w}{ }mpi\PYG{+w}{ }\PYG{n+nv}{mode}\PYG{o}{=}shlib\PYG{+w}{ }\PYG{n+nv}{LMP\PYGZus{}INC}\PYG{o}{=}\PYG{l+s+s2}{\PYGZdq{}\PYGZhy{}DLAMMPS\PYGZus{}PNG \PYGZhy{}DLAMMPS\PYGZus{}JPEG \PYGZhy{}DLAMMPS\PYGZus{}FFMPEG\PYGZdq{}}\PYG{+w}{ }\PYG{n+nv}{JPG\PYGZus{}LIB}\PYG{o}{=}\PYG{l+s+s2}{\PYGZdq{}\PYGZhy{}lpng \PYGZhy{}ljpeg\PYGZdq{}}
\end{sphinxVerbatim}


\subparagraph{Step 2: Installing the LAMMPS Python package}
\label{\detokenize{Howto_pylammps:step-2-installing-the-lammps-python-package}}
\sphinxAtStartPar
PyLammps is part of the lammps Python package. To install it simply install
that package into your current Python installation with:

\begin{sphinxVerbatim}[commandchars=\\\{\}]
make\PYG{+w}{ }install\PYGZhy{}python
\end{sphinxVerbatim}

\begin{sphinxadmonition}{note}{Note:}
\sphinxAtStartPar
Recompiling the shared library requires re\sphinxhyphen{}installing the Python package
\end{sphinxadmonition}


\paragraph{Installation inside of a virtualenv}
\label{\detokenize{Howto_pylammps:installation-inside-of-a-virtualenv}}
\sphinxAtStartPar
You can use virtualenv to create a custom Python environment specifically tuned
for your workflow.


\subparagraph{Benefits of using a virtualenv}
\label{\detokenize{Howto_pylammps:benefits-of-using-a-virtualenv}}\begin{itemize}
\item {} 
\sphinxAtStartPar
isolation of your system Python installation from your development installation

\item {} 
\sphinxAtStartPar
installation can happen in your user directory without root access (useful for HPC clusters)

\item {} 
\sphinxAtStartPar
installing packages through pip allows you to get newer versions of packages than e.g., through apt\sphinxhyphen{}get or yum package managers (and without root access)

\item {} 
\sphinxAtStartPar
you can even install specific old versions of a package if necessary

\end{itemize}

\sphinxAtStartPar
\sphinxstylestrong{Prerequisite (e.g. on Ubuntu)}

\begin{sphinxVerbatim}[commandchars=\\\{\}]
apt\PYGZhy{}get\PYG{+w}{ }install\PYG{+w}{ }python\PYGZhy{}virtualenv
\end{sphinxVerbatim}


\subparagraph{Creating a virtualenv with lammps installed}
\label{\detokenize{Howto_pylammps:creating-a-virtualenv-with-lammps-installed}}
\begin{sphinxVerbatim}[commandchars=\\\{\}]
\PYG{c+c1}{\PYGZsh{} create virtualenv named \PYGZsq{}testing\PYGZsq{}}
virtualenv\PYG{+w}{ }\PYG{n+nv}{\PYGZdl{}HOME}/python/testing

\PYG{c+c1}{\PYGZsh{} activate \PYGZsq{}testing\PYGZsq{} environment}
\PYG{n+nb}{source}\PYG{+w}{ }\PYG{n+nv}{\PYGZdl{}HOME}/python/testing/bin/activate
\end{sphinxVerbatim}

\sphinxAtStartPar
Now configure and compile the LAMMPS shared library as outlined above.
When using CMake and the shared library has already been build, you
need to re\sphinxhyphen{}run CMake to update the location of the python executable
to the location in the virtual environment with:

\begin{sphinxVerbatim}[commandchars=\\\{\}]
cmake\PYG{+w}{ }.\PYG{+w}{ }\PYGZhy{}DPython\PYGZus{}EXECUTABLE\PYG{o}{=}\PYG{k}{\PYGZdl{}(}which\PYG{+w}{ }python\PYG{k}{)}

\PYG{c+c1}{\PYGZsh{} install LAMMPS package in virtualenv}
\PYG{o}{(}testing\PYG{o}{)}\PYG{+w}{ }make\PYG{+w}{ }install\PYGZhy{}python

\PYG{c+c1}{\PYGZsh{} install other useful packages}
\PYG{o}{(}testing\PYG{o}{)}\PYG{+w}{ }pip\PYG{+w}{ }install\PYG{+w}{ }matplotlib\PYG{+w}{ }jupyter\PYG{+w}{ }mpi4py

...

\PYG{c+c1}{\PYGZsh{} return to original shell}
\PYG{o}{(}testing\PYG{o}{)}\PYG{+w}{ }deactivate
\end{sphinxVerbatim}


\subsubsection{Creating a new instance of PyLammps}
\label{\detokenize{Howto_pylammps:creating-a-new-instance-of-pylammps}}
\sphinxAtStartPar
To create a PyLammps object you need to first import the class from the lammps
module. By using the default constructor, a new \sphinxstyleemphasis{lammps} instance is created.

\begin{sphinxVerbatim}[commandchars=\\\{\}]
\PYG{k+kn}{from}\PYG{+w}{ }\PYG{n+nn}{lammps}\PYG{+w}{ }\PYG{k+kn}{import} \PYG{n}{PyLammps}
\PYG{n}{L} \PYG{o}{=} \PYG{n}{PyLammps}\PYG{p}{(}\PYG{p}{)}
\end{sphinxVerbatim}

\sphinxAtStartPar
You can also initialize PyLammps on top of this existing \sphinxstyleemphasis{lammps} object:

\begin{sphinxVerbatim}[commandchars=\\\{\}]
\PYG{k+kn}{from}\PYG{+w}{ }\PYG{n+nn}{lammps}\PYG{+w}{ }\PYG{k+kn}{import} \PYG{n}{lammps}\PYG{p}{,} \PYG{n}{PyLammps}
\PYG{n}{lmp} \PYG{o}{=} \PYG{n}{lammps}\PYG{p}{(}\PYG{p}{)}
\PYG{n}{L} \PYG{o}{=} \PYG{n}{PyLammps}\PYG{p}{(}\PYG{n}{ptr}\PYG{o}{=}\PYG{n}{lmp}\PYG{p}{)}
\end{sphinxVerbatim}


\subsubsection{Commands}
\label{\detokenize{Howto_pylammps:commands}}
\sphinxAtStartPar
Sending a LAMMPS command with the existing library interfaces is done using
the command method of the lammps object instance.

\sphinxAtStartPar
For instance, let’s take the following LAMMPS command:

\begin{sphinxVerbatim}[commandchars=\\\{\}]
\PYG{k}{region }\PYG{n+nv+nvIdentifier}{box}\PYG{+w}{ }\PYG{n}{block}\PYG{+w}{ }\PYG{l+m}{0}\PYG{+w}{ }\PYG{l+m}{10}\PYG{+w}{ }\PYG{l+m}{0}\PYG{+w}{ }\PYG{l+m}{5}\PYG{+w}{ }\PYG{o}{\PYGZhy{}}\PYG{l+m}{0.5}\PYG{+w}{ }\PYG{l+m}{0.5}
\end{sphinxVerbatim}

\sphinxAtStartPar
In the original interface this command can be executed with the following
Python code if \sphinxstyleemphasis{L} was a lammps instance:

\begin{sphinxVerbatim}[commandchars=\\\{\}]
\PYG{n}{L}\PYG{o}{.}\PYG{n}{command}\PYG{p}{(}\PYG{l+s+s2}{\PYGZdq{}}\PYG{l+s+s2}{region box block 0 10 0 5 \PYGZhy{}0.5 0.5}\PYG{l+s+s2}{\PYGZdq{}}\PYG{p}{)}
\end{sphinxVerbatim}

\sphinxAtStartPar
With the PyLammps interface, any command can be split up into arbitrary parts
separated by white\sphinxhyphen{}space, passed as individual arguments to a region method.

\begin{sphinxVerbatim}[commandchars=\\\{\}]
\PYG{n}{L}\PYG{o}{.}\PYG{n}{region}\PYG{p}{(}\PYG{l+s+s2}{\PYGZdq{}}\PYG{l+s+s2}{box block}\PYG{l+s+s2}{\PYGZdq{}}\PYG{p}{,} \PYG{l+m+mi}{0}\PYG{p}{,} \PYG{l+m+mi}{10}\PYG{p}{,} \PYG{l+m+mi}{0}\PYG{p}{,} \PYG{l+m+mi}{5}\PYG{p}{,} \PYG{o}{\PYGZhy{}}\PYG{l+m+mf}{0.5}\PYG{p}{,} \PYG{l+m+mf}{0.5}\PYG{p}{)}
\end{sphinxVerbatim}

\sphinxAtStartPar
Note that each parameter is set as Python literal floating\sphinxhyphen{}point number. In the
PyLammps interface, each command takes an arbitrary parameter list and transparently
merges it to a single command string, separating individual parameters by white\sphinxhyphen{}space.

\sphinxAtStartPar
The benefit of this approach is avoiding redundant command calls and easier
parameterization. In the original interface parameterization needed to be done
manually by creating formatted strings.

\begin{sphinxVerbatim}[commandchars=\\\{\}]
\PYG{n}{L}\PYG{o}{.}\PYG{n}{command}\PYG{p}{(}\PYG{l+s+s2}{\PYGZdq{}}\PYG{l+s+s2}{region box block }\PYG{l+s+si}{\PYGZpc{}f}\PYG{l+s+s2}{ }\PYG{l+s+si}{\PYGZpc{}f}\PYG{l+s+s2}{ }\PYG{l+s+si}{\PYGZpc{}f}\PYG{l+s+s2}{ }\PYG{l+s+si}{\PYGZpc{}f}\PYG{l+s+s2}{ }\PYG{l+s+si}{\PYGZpc{}f}\PYG{l+s+s2}{ }\PYG{l+s+si}{\PYGZpc{}f}\PYG{l+s+s2}{\PYGZdq{}} \PYG{o}{\PYGZpc{}} \PYG{p}{(}\PYG{n}{xlo}\PYG{p}{,} \PYG{n}{xhi}\PYG{p}{,} \PYG{n}{ylo}\PYG{p}{,} \PYG{n}{yhi}\PYG{p}{,} \PYG{n}{zlo}\PYG{p}{,} \PYG{n}{zhi}\PYG{p}{)}\PYG{p}{)}
\end{sphinxVerbatim}

\sphinxAtStartPar
In contrast, methods of PyLammps accept parameters directly and will convert
them automatically to a final command string.

\begin{sphinxVerbatim}[commandchars=\\\{\}]
\PYG{n}{L}\PYG{o}{.}\PYG{n}{region}\PYG{p}{(}\PYG{l+s+s2}{\PYGZdq{}}\PYG{l+s+s2}{box block}\PYG{l+s+s2}{\PYGZdq{}}\PYG{p}{,} \PYG{n}{xlo}\PYG{p}{,} \PYG{n}{xhi}\PYG{p}{,} \PYG{n}{ylo}\PYG{p}{,} \PYG{n}{yhi}\PYG{p}{,} \PYG{n}{zlo}\PYG{p}{,} \PYG{n}{zhi}\PYG{p}{)}
\end{sphinxVerbatim}


\subsubsection{System state}
\label{\detokenize{Howto_pylammps:system-state}}
\sphinxAtStartPar
In addition to dispatching commands directly through the PyLammps object, it
also provides several properties which allow you to query the system state.
\begin{description}
\sphinxlineitem{L.system}
\sphinxAtStartPar
Is a dictionary describing the system such as the bounding box or number of atoms

\sphinxlineitem{L.system.xlo, L.system.xhi}
\sphinxAtStartPar
bounding box limits along x\sphinxhyphen{}axis

\sphinxlineitem{L.system.ylo, L.system.yhi}
\sphinxAtStartPar
bounding box limits along y\sphinxhyphen{}axis

\sphinxlineitem{L.system.zlo, L.system.zhi}
\sphinxAtStartPar
bounding box limits along z\sphinxhyphen{}axis

\sphinxlineitem{L.communication}
\sphinxAtStartPar
configuration of communication subsystem, such as the number of threads or processors

\sphinxlineitem{L.communication.nthreads}
\sphinxAtStartPar
number of threads used by each LAMMPS process

\sphinxlineitem{L.communication.nprocs}
\sphinxAtStartPar
number of MPI processes used by LAMMPS

\sphinxlineitem{L.fixes}
\sphinxAtStartPar
List of fixes in the current system

\sphinxlineitem{L.computes}
\sphinxAtStartPar
List of active computes in the current system

\sphinxlineitem{L.dump}
\sphinxAtStartPar
List of active dumps in the current system

\sphinxlineitem{L.groups}
\sphinxAtStartPar
List of groups present in the current system

\end{description}


\subsubsection{Working with LAMMPS variables}
\label{\detokenize{Howto_pylammps:working-with-lammps-variables}}
\sphinxAtStartPar
LAMMPS variables can be both defined and accessed via the PyLammps interface.

\sphinxAtStartPar
To define a variable you can use the {\hyperref[\detokenize{variable::doc}]{\sphinxcrossref{\DUrole{doc}{variable}}}} command:

\begin{sphinxVerbatim}[commandchars=\\\{\}]
\PYG{n}{L}\PYG{o}{.}\PYG{n}{variable}\PYG{p}{(}\PYG{l+s+s2}{\PYGZdq{}}\PYG{l+s+s2}{a index 2}\PYG{l+s+s2}{\PYGZdq{}}\PYG{p}{)}
\end{sphinxVerbatim}

\sphinxAtStartPar
A dictionary of all variables is returned by L.variables

\sphinxAtStartPar
you can access an individual variable by retrieving a variable object from the
L.variables dictionary by name

\begin{sphinxVerbatim}[commandchars=\\\{\}]
\PYG{n}{a} \PYG{o}{=} \PYG{n}{L}\PYG{o}{.}\PYG{n}{variables}\PYG{p}{[}\PYG{l+s+s1}{\PYGZsq{}}\PYG{l+s+s1}{a}\PYG{l+s+s1}{\PYGZsq{}}\PYG{p}{]}
\end{sphinxVerbatim}

\sphinxAtStartPar
The variable value can then be easily read and written by accessing the value
property of this object.

\begin{sphinxVerbatim}[commandchars=\\\{\}]
\PYG{n+nb}{print}\PYG{p}{(}\PYG{n}{a}\PYG{o}{.}\PYG{n}{value}\PYG{p}{)}
\PYG{n}{a}\PYG{o}{.}\PYG{n}{value} \PYG{o}{=} \PYG{l+m+mi}{4}
\end{sphinxVerbatim}


\subsubsection{Retrieving the value of an arbitrary LAMMPS expressions}
\label{\detokenize{Howto_pylammps:retrieving-the-value-of-an-arbitrary-lammps-expressions}}
\sphinxAtStartPar
LAMMPS expressions can be immediately evaluated by using the eval method. The
passed string parameter can be any expression containing global thermo values,
variables, compute or fix data.

\begin{sphinxVerbatim}[commandchars=\\\{\}]
\PYG{n}{result} \PYG{o}{=} \PYG{n}{L}\PYG{o}{.}\PYG{n}{eval}\PYG{p}{(}\PYG{l+s+s2}{\PYGZdq{}}\PYG{l+s+s2}{ke}\PYG{l+s+s2}{\PYGZdq{}}\PYG{p}{)} \PYG{c+c1}{\PYGZsh{} kinetic energy}
\PYG{n}{result} \PYG{o}{=} \PYG{n}{L}\PYG{o}{.}\PYG{n}{eval}\PYG{p}{(}\PYG{l+s+s2}{\PYGZdq{}}\PYG{l+s+s2}{pe}\PYG{l+s+s2}{\PYGZdq{}}\PYG{p}{)} \PYG{c+c1}{\PYGZsh{} potential energy}

\PYG{n}{result} \PYG{o}{=} \PYG{n}{L}\PYG{o}{.}\PYG{n}{eval}\PYG{p}{(}\PYG{l+s+s2}{\PYGZdq{}}\PYG{l+s+s2}{v\PYGZus{}t/2.0}\PYG{l+s+s2}{\PYGZdq{}}\PYG{p}{)}
\end{sphinxVerbatim}


\subsubsection{Accessing atom data}
\label{\detokenize{Howto_pylammps:accessing-atom-data}}
\sphinxAtStartPar
All atoms in the current simulation can be accessed by using the L.atoms list.
Each element of this list is an object which exposes its properties (id, type,
position, velocity, force, etc.).

\begin{sphinxVerbatim}[commandchars=\\\{\}]
\PYG{c+c1}{\PYGZsh{} access first atom}
\PYG{n}{L}\PYG{o}{.}\PYG{n}{atoms}\PYG{p}{[}\PYG{l+m+mi}{0}\PYG{p}{]}\PYG{o}{.}\PYG{n}{id}
\PYG{n}{L}\PYG{o}{.}\PYG{n}{atoms}\PYG{p}{[}\PYG{l+m+mi}{0}\PYG{p}{]}\PYG{o}{.}\PYG{n}{type}

\PYG{c+c1}{\PYGZsh{} access second atom}
\PYG{n}{L}\PYG{o}{.}\PYG{n}{atoms}\PYG{p}{[}\PYG{l+m+mi}{1}\PYG{p}{]}\PYG{o}{.}\PYG{n}{position}
\PYG{n}{L}\PYG{o}{.}\PYG{n}{atoms}\PYG{p}{[}\PYG{l+m+mi}{1}\PYG{p}{]}\PYG{o}{.}\PYG{n}{velocity}
\PYG{n}{L}\PYG{o}{.}\PYG{n}{atoms}\PYG{p}{[}\PYG{l+m+mi}{1}\PYG{p}{]}\PYG{o}{.}\PYG{n}{force}
\end{sphinxVerbatim}

\sphinxAtStartPar
Some properties can also be used to set:

\begin{sphinxVerbatim}[commandchars=\\\{\}]
\PYG{c+c1}{\PYGZsh{} set position in 2D simulation}
\PYG{n}{L}\PYG{o}{.}\PYG{n}{atoms}\PYG{p}{[}\PYG{l+m+mi}{0}\PYG{p}{]}\PYG{o}{.}\PYG{n}{position} \PYG{o}{=} \PYG{p}{(}\PYG{l+m+mf}{1.0}\PYG{p}{,} \PYG{l+m+mf}{0.0}\PYG{p}{)}

\PYG{c+c1}{\PYGZsh{} set position in 3D simulation}
\PYG{n}{L}\PYG{o}{.}\PYG{n}{atoms}\PYG{p}{[}\PYG{l+m+mi}{0}\PYG{p}{]}\PYG{o}{.}\PYG{n}{position} \PYG{o}{=} \PYG{p}{(}\PYG{l+m+mf}{1.0}\PYG{p}{,} \PYG{l+m+mf}{0.0}\PYG{p}{,} \PYG{l+m+mf}{1.}\PYG{p}{)}
\end{sphinxVerbatim}


\subsubsection{Evaluating thermo data}
\label{\detokenize{Howto_pylammps:evaluating-thermo-data}}
\sphinxAtStartPar
Each simulation run usually produces thermo output based on system state,
computes, fixes or variables. The trajectories of these values can be queried
after a run via the L.runs list. This list contains a growing list of run data.
The first element is the output of the first run, the second element that of
the second run.

\begin{sphinxVerbatim}[commandchars=\\\{\}]
\PYG{n}{L}\PYG{o}{.}\PYG{n}{run}\PYG{p}{(}\PYG{l+m+mi}{1000}\PYG{p}{)}
\PYG{n}{L}\PYG{o}{.}\PYG{n}{runs}\PYG{p}{[}\PYG{l+m+mi}{0}\PYG{p}{]} \PYG{c+c1}{\PYGZsh{} data of first 1000 time steps}

\PYG{n}{L}\PYG{o}{.}\PYG{n}{run}\PYG{p}{(}\PYG{l+m+mi}{1000}\PYG{p}{)}
\PYG{n}{L}\PYG{o}{.}\PYG{n}{runs}\PYG{p}{[}\PYG{l+m+mi}{1}\PYG{p}{]} \PYG{c+c1}{\PYGZsh{} data of second 1000 time steps}
\end{sphinxVerbatim}

\sphinxAtStartPar
Each run contains a dictionary of all trajectories. Each trajectory is
accessible through its thermo name:

\begin{sphinxVerbatim}[commandchars=\\\{\}]
\PYG{n}{L}\PYG{o}{.}\PYG{n}{runs}\PYG{p}{[}\PYG{l+m+mi}{0}\PYG{p}{]}\PYG{o}{.}\PYG{n}{thermo}\PYG{o}{.}\PYG{n}{Step} \PYG{c+c1}{\PYGZsh{} list of time steps in first run}
\PYG{n}{L}\PYG{o}{.}\PYG{n}{runs}\PYG{p}{[}\PYG{l+m+mi}{0}\PYG{p}{]}\PYG{o}{.}\PYG{n}{thermo}\PYG{o}{.}\PYG{n}{Ke}   \PYG{c+c1}{\PYGZsh{} list of kinetic energy values in first run}
\end{sphinxVerbatim}

\sphinxAtStartPar
Together with matplotlib plotting data out of LAMMPS becomes simple:

\begin{sphinxVerbatim}[commandchars=\\\{\}]
\PYG{k+kn}{import}\PYG{+w}{ }\PYG{n+nn}{matplotlib}\PYG{n+nn}{.}\PYG{n+nn}{plot}\PYG{+w}{ }\PYG{k}{as}\PYG{+w}{ }\PYG{n+nn}{plt}
\PYG{n}{steps} \PYG{o}{=} \PYG{n}{L}\PYG{o}{.}\PYG{n}{runs}\PYG{p}{[}\PYG{l+m+mi}{0}\PYG{p}{]}\PYG{o}{.}\PYG{n}{thermo}\PYG{o}{.}\PYG{n}{Step}
\PYG{n}{ke}    \PYG{o}{=} \PYG{n}{L}\PYG{o}{.}\PYG{n}{runs}\PYG{p}{[}\PYG{l+m+mi}{0}\PYG{p}{]}\PYG{o}{.}\PYG{n}{thermo}\PYG{o}{.}\PYG{n}{Ke}
\PYG{n}{plt}\PYG{o}{.}\PYG{n}{plot}\PYG{p}{(}\PYG{n}{steps}\PYG{p}{,} \PYG{n}{ke}\PYG{p}{)}
\end{sphinxVerbatim}


\subsubsection{Error handling with PyLammps}
\label{\detokenize{Howto_pylammps:error-handling-with-pylammps}}
\sphinxAtStartPar
Using C++ exceptions in LAMMPS for errors allows capturing them on the
C++ side and rethrowing them on the Python side.  This way you can handle
LAMMPS errors through the Python exception handling mechanism.

\begin{sphinxadmonition}{warning}{Warning:}
\sphinxAtStartPar
Capturing a LAMMPS exception in Python can still mean that the
current LAMMPS process is in an illegal state and must be
terminated. It is advised to save your data and terminate the Python
instance as quickly as possible.
\end{sphinxadmonition}


\subsubsection{Using PyLammps in IPython notebooks and Jupyter}
\label{\detokenize{Howto_pylammps:using-pylammps-in-ipython-notebooks-and-jupyter}}
\sphinxAtStartPar
If the LAMMPS Python package is installed for the same Python interpreter as
IPython, you can use PyLammps directly inside of an IPython notebook inside of
Jupyter. Jupyter is a powerful integrated development environment (IDE) for
many dynamic languages like Python, Julia and others, which operates inside of
any web browser. Besides auto\sphinxhyphen{}completion and syntax highlighting it allows you
to create formatted documents using Markup, mathematical formulas, graphics and
animations intermixed with executable Python code. It is a great format for
tutorials and showcasing your latest research.

\sphinxAtStartPar
To launch an instance of Jupyter simply run the following command inside your
Python environment (this assumes you followed the Quick Start instructions):

\begin{sphinxVerbatim}[commandchars=\\\{\}]
jupyter\PYG{+w}{ }notebook
\end{sphinxVerbatim}


\subsubsection{IPyLammps Examples}
\label{\detokenize{Howto_pylammps:ipylammps-examples}}
\sphinxAtStartPar
Examples of IPython notebooks can be found in the python/examples/pylammps
subdirectory. To open these notebooks launch \sphinxstyleemphasis{jupyter notebook} inside this
directory and navigate to one of them. If you compiled and installed
a LAMMPS shared library with exceptions, PNG, JPEG and FFMPEG support
you should be able to rerun all of these notebooks.


\paragraph{Validating a dihedral potential}
\label{\detokenize{Howto_pylammps:validating-a-dihedral-potential}}
\sphinxAtStartPar
This example showcases how an IPython Notebook can be used to compare a simple
LAMMPS simulation of a harmonic dihedral potential to its analytical solution.
Four atoms are placed in the simulation and the dihedral potential is applied on
them using a datafile. Then one of the atoms is rotated along the central axis by
setting its position from Python, which changes the dihedral angle.

\begin{sphinxVerbatim}[commandchars=\\\{\}]
\PYG{n}{phi} \PYG{o}{=} \PYG{p}{[}\PYG{n}{d} \PYGZbs{}\PYG{o}{*} \PYG{n}{math}\PYG{o}{.}\PYG{n}{pi} \PYG{o}{/} \PYG{l+m+mi}{180} \PYG{k}{for} \PYG{n}{d} \PYG{o+ow}{in} \PYG{n+nb}{range}\PYG{p}{(}\PYG{l+m+mi}{360}\PYG{p}{)}\PYG{p}{]}

\PYG{n}{pos} \PYG{o}{=} \PYG{p}{[}\PYG{p}{(}\PYG{l+m+mf}{1.0}\PYG{p}{,} \PYG{n}{math}\PYG{o}{.}\PYG{n}{cos}\PYG{p}{(}\PYG{n}{p}\PYG{p}{)}\PYG{p}{,} \PYG{n}{math}\PYG{o}{.}\PYG{n}{sin}\PYG{p}{(}\PYG{n}{p}\PYG{p}{)}\PYG{p}{)} \PYG{k}{for} \PYG{n}{p} \PYG{o+ow}{in} \PYG{n}{phi}\PYG{p}{]}

\PYG{n}{pe} \PYG{o}{=} \PYG{p}{[}\PYG{p}{]}
\PYG{k}{for} \PYG{n}{p} \PYG{o+ow}{in} \PYG{n}{pos}\PYG{p}{:}
    \PYG{n}{L}\PYG{o}{.}\PYG{n}{atoms}\PYG{p}{[}\PYG{l+m+mi}{3}\PYG{p}{]}\PYG{o}{.}\PYG{n}{position} \PYG{o}{=} \PYG{n}{p}
    \PYG{n}{L}\PYG{o}{.}\PYG{n}{run}\PYG{p}{(}\PYG{l+m+mi}{0}\PYG{p}{)}
    \PYG{n}{pe}\PYG{o}{.}\PYG{n}{append}\PYG{p}{(}\PYG{n}{L}\PYG{o}{.}\PYG{n}{eval}\PYG{p}{(}\PYG{l+s+s2}{\PYGZdq{}}\PYG{l+s+s2}{pe}\PYG{l+s+s2}{\PYGZdq{}}\PYG{p}{)}\PYG{p}{)}
\end{sphinxVerbatim}

\sphinxAtStartPar
By evaluating the potential energy for each position we can verify that
trajectory with the analytical formula.  To compare both solutions, we plot
both trajectories over each other using matplotlib, which embeds the generated
plot inside the IPython notebook.

\noindent{\hspace*{\fill}\sphinxincludegraphics{{pylammps_dihedral}.jpg}\hspace*{\fill}}


\paragraph{Running a Monte Carlo relaxation}
\label{\detokenize{Howto_pylammps:running-a-monte-carlo-relaxation}}
\sphinxAtStartPar
This second example shows how to use PyLammps to create a 2D Monte Carlo Relaxation
simulation, computing and plotting energy terms and even embedding video output.

\sphinxAtStartPar
Initially, a 2D system is created in a state with minimal energy.

\noindent{\hspace*{\fill}\sphinxincludegraphics{{pylammps_mc_minimum}.jpg}\hspace*{\fill}}

\sphinxAtStartPar
It is then disordered by moving each atom by a random delta.

\begin{sphinxVerbatim}[commandchars=\\\{\}]
\PYG{n}{random}\PYG{o}{.}\PYG{n}{seed}\PYG{p}{(}\PYG{l+m+mi}{27848}\PYG{p}{)}
\PYG{n}{deltaperturb} \PYG{o}{=} \PYG{l+m+mf}{0.2}

\PYG{k}{for} \PYG{n}{i} \PYG{o+ow}{in} \PYG{n+nb}{range}\PYG{p}{(}\PYG{n}{L}\PYG{o}{.}\PYG{n}{system}\PYG{o}{.}\PYG{n}{natoms}\PYG{p}{)}\PYG{p}{:}
    \PYG{n}{x}\PYG{p}{,} \PYG{n}{y} \PYG{o}{=} \PYG{n}{L}\PYG{o}{.}\PYG{n}{atoms}\PYG{p}{[}\PYG{n}{i}\PYG{p}{]}\PYG{o}{.}\PYG{n}{position}
    \PYG{n}{dx} \PYG{o}{=} \PYG{n}{deltaperturb} \PYGZbs{}\PYG{o}{*} \PYG{n}{random}\PYG{o}{.}\PYG{n}{uniform}\PYG{p}{(}\PYG{o}{\PYGZhy{}}\PYG{l+m+mi}{1}\PYG{p}{,} \PYG{l+m+mi}{1}\PYG{p}{)}
    \PYG{n}{dy} \PYG{o}{=} \PYG{n}{deltaperturb} \PYGZbs{}\PYG{o}{*} \PYG{n}{random}\PYG{o}{.}\PYG{n}{uniform}\PYG{p}{(}\PYG{o}{\PYGZhy{}}\PYG{l+m+mi}{1}\PYG{p}{,} \PYG{l+m+mi}{1}\PYG{p}{)}
    \PYG{n}{L}\PYG{o}{.}\PYG{n}{atoms}\PYG{p}{[}\PYG{n}{i}\PYG{p}{]}\PYG{o}{.}\PYG{n}{position}  \PYG{o}{=} \PYG{p}{(}\PYG{n}{x}\PYG{o}{+}\PYG{n}{dx}\PYG{p}{,} \PYG{n}{y}\PYG{o}{+}\PYG{n}{dy}\PYG{p}{)}

\PYG{n}{L}\PYG{o}{.}\PYG{n}{run}\PYG{p}{(}\PYG{l+m+mi}{0}\PYG{p}{)}
\end{sphinxVerbatim}

\noindent{\hspace*{\fill}\sphinxincludegraphics{{pylammps_mc_disordered}.jpg}\hspace*{\fill}}

\sphinxAtStartPar
Finally, the Monte Carlo algorithm is implemented in Python. It continuously
moves random atoms by a random delta and only accepts certain moves.

\begin{sphinxVerbatim}[commandchars=\\\{\}]
\PYG{n}{estart} \PYG{o}{=} \PYG{n}{L}\PYG{o}{.}\PYG{n}{eval}\PYG{p}{(}\PYG{l+s+s2}{\PYGZdq{}}\PYG{l+s+s2}{pe}\PYG{l+s+s2}{\PYGZdq{}}\PYG{p}{)}
\PYG{n}{elast} \PYG{o}{=} \PYG{n}{estart}

\PYG{n}{naccept} \PYG{o}{=} \PYG{l+m+mi}{0}
\PYG{n}{energies} \PYG{o}{=} \PYG{p}{[}\PYG{n}{estart}\PYG{p}{]}

\PYG{n}{niterations} \PYG{o}{=} \PYG{l+m+mi}{3000}
\PYG{n}{deltamove} \PYG{o}{=} \PYG{l+m+mf}{0.1}
\PYG{n}{kT} \PYG{o}{=} \PYG{l+m+mf}{0.05}

\PYG{n}{natoms} \PYG{o}{=} \PYG{n}{L}\PYG{o}{.}\PYG{n}{system}\PYG{o}{.}\PYG{n}{natoms}

\PYG{k}{for} \PYG{n}{i} \PYG{o+ow}{in} \PYG{n+nb}{range}\PYG{p}{(}\PYG{n}{niterations}\PYG{p}{)}\PYG{p}{:}
    \PYG{n}{iatom} \PYG{o}{=} \PYG{n}{random}\PYG{o}{.}\PYG{n}{randrange}\PYG{p}{(}\PYG{l+m+mi}{0}\PYG{p}{,} \PYG{n}{natoms}\PYG{p}{)}
    \PYG{n}{current\PYGZus{}atom} \PYG{o}{=} \PYG{n}{L}\PYG{o}{.}\PYG{n}{atoms}\PYG{p}{[}\PYG{n}{iatom}\PYG{p}{]}

    \PYG{n}{x0}\PYG{p}{,} \PYG{n}{y0} \PYG{o}{=} \PYG{n}{current\PYGZus{}atom}\PYG{o}{.}\PYG{n}{position}

    \PYG{n}{dx} \PYG{o}{=} \PYG{n}{deltamove} \PYGZbs{}\PYG{o}{*} \PYG{n}{random}\PYG{o}{.}\PYG{n}{uniform}\PYG{p}{(}\PYG{o}{\PYGZhy{}}\PYG{l+m+mi}{1}\PYG{p}{,} \PYG{l+m+mi}{1}\PYG{p}{)}
    \PYG{n}{dy} \PYG{o}{=} \PYG{n}{deltamove} \PYGZbs{}\PYG{o}{*} \PYG{n}{random}\PYG{o}{.}\PYG{n}{uniform}\PYG{p}{(}\PYG{o}{\PYGZhy{}}\PYG{l+m+mi}{1}\PYG{p}{,} \PYG{l+m+mi}{1}\PYG{p}{)}

    \PYG{n}{current\PYGZus{}atom}\PYG{o}{.}\PYG{n}{position} \PYG{o}{=} \PYG{p}{(}\PYG{n}{x0}\PYG{o}{+}\PYG{n}{dx}\PYG{p}{,} \PYG{n}{y0}\PYG{o}{+}\PYG{n}{dy}\PYG{p}{)}

    \PYG{n}{L}\PYG{o}{.}\PYG{n}{run}\PYG{p}{(}\PYG{l+m+mi}{1}\PYG{p}{,} \PYG{l+s+s2}{\PYGZdq{}}\PYG{l+s+s2}{pre no post no}\PYG{l+s+s2}{\PYGZdq{}}\PYG{p}{)}

    \PYG{n}{e} \PYG{o}{=} \PYG{n}{L}\PYG{o}{.}\PYG{n}{eval}\PYG{p}{(}\PYG{l+s+s2}{\PYGZdq{}}\PYG{l+s+s2}{pe}\PYG{l+s+s2}{\PYGZdq{}}\PYG{p}{)}
    \PYG{n}{energies}\PYG{o}{.}\PYG{n}{append}\PYG{p}{(}\PYG{n}{e}\PYG{p}{)}

    \PYG{k}{if} \PYG{n}{e} \PYG{o}{\PYGZlt{}}\PYG{o}{=} \PYG{n}{elast}\PYG{p}{:}
        \PYG{n}{naccept} \PYG{o}{+}\PYG{o}{=} \PYG{l+m+mi}{1}
        \PYG{n}{elast} \PYG{o}{=} \PYG{n}{e}
    \PYG{k}{elif} \PYG{n}{random}\PYG{o}{.}\PYG{n}{random}\PYG{p}{(}\PYG{p}{)} \PYG{o}{\PYGZlt{}}\PYG{o}{=} \PYG{n}{math}\PYG{o}{.}\PYG{n}{exp}\PYG{p}{(}\PYG{n}{natoms}\PYGZbs{}\PYG{o}{*}\PYG{p}{(}\PYG{n}{elast}\PYG{o}{\PYGZhy{}}\PYG{n}{e}\PYG{p}{)}\PYG{o}{/}\PYG{n}{kT}\PYG{p}{)}\PYG{p}{:}
        \PYG{n}{naccept} \PYG{o}{+}\PYG{o}{=} \PYG{l+m+mi}{1}
        \PYG{n}{elast} \PYG{o}{=} \PYG{n}{e}
    \PYG{k}{else}\PYG{p}{:}
        \PYG{n}{current\PYGZus{}atom}\PYG{o}{.}\PYG{n}{position} \PYG{o}{=} \PYG{p}{(}\PYG{n}{x0}\PYG{p}{,} \PYG{n}{y0}\PYG{p}{)}
\end{sphinxVerbatim}

\sphinxAtStartPar
The energies of each iteration are collected in a Python list and finally plotted using matplotlib.

\noindent{\hspace*{\fill}\sphinxincludegraphics{{pylammps_mc_energies_plot}.jpg}\hspace*{\fill}}

\sphinxAtStartPar
The IPython notebook also shows how to use dump commands and embed video files
inside of the IPython notebook.


\subsubsection{Using PyLammps and mpi4py (Experimental)}
\label{\detokenize{Howto_pylammps:using-pylammps-and-mpi4py-experimental}}
\sphinxAtStartPar
PyLammps can be run in parallel using \sphinxhref{https://mpi4py.readthedocs.io}{mpi4py}. This python package can be installed
using

\begin{sphinxVerbatim}[commandchars=\\\{\}]
pip\PYG{+w}{ }install\PYG{+w}{ }mpi4py
\end{sphinxVerbatim}

\begin{sphinxadmonition}{warning}{Warning:}
\sphinxAtStartPar
Usually, any {\hyperref[\detokenize{Python_module:lammps.PyLammps}]{\sphinxcrossref{\sphinxcode{\sphinxupquote{PyLammps}}}}} command must be
executed by \sphinxstyleemphasis{all} MPI processes. However, evaluations and querying
the system state is only available on MPI rank 0.  Using these
functions from other MPI ranks will raise an exception.
\end{sphinxadmonition}

\sphinxAtStartPar
The following is a short example which reads in an existing LAMMPS input
file and executes it in parallel.  You can find in.melt in the
examples/melt folder.  Please take note that the
{\hyperref[\detokenize{Python_module:lammps.PyLammps.eval}]{\sphinxcrossref{\sphinxcode{\sphinxupquote{PyLammps.eval()}}}}} is called only from
MPI rank 0.

\begin{sphinxVerbatim}[commandchars=\\\{\}]
\PYG{k+kn}{from}\PYG{+w}{ }\PYG{n+nn}{mpi4py}\PYG{+w}{ }\PYG{k+kn}{import} \PYG{n}{MPI}
\PYG{k+kn}{from}\PYG{+w}{ }\PYG{n+nn}{lammps}\PYG{+w}{ }\PYG{k+kn}{import} \PYG{n}{PyLammps}

\PYG{n}{L} \PYG{o}{=} \PYG{n}{PyLammps}\PYG{p}{(}\PYG{p}{)}
\PYG{n}{L}\PYG{o}{.}\PYG{n}{file}\PYG{p}{(}\PYG{l+s+s2}{\PYGZdq{}}\PYG{l+s+s2}{in.melt}\PYG{l+s+s2}{\PYGZdq{}}\PYG{p}{)}

\PYG{k}{if} \PYG{n}{MPI}\PYG{o}{.}\PYG{n}{COMM\PYGZus{}WORLD}\PYG{o}{.}\PYG{n}{rank} \PYG{o}{==} \PYG{l+m+mi}{0}\PYG{p}{:}
    \PYG{n+nb}{print}\PYG{p}{(}\PYG{l+s+s2}{\PYGZdq{}}\PYG{l+s+s2}{Potential energy: }\PYG{l+s+s2}{\PYGZdq{}}\PYG{p}{,} \PYG{n}{L}\PYG{o}{.}\PYG{n}{eval}\PYG{p}{(}\PYG{l+s+s2}{\PYGZdq{}}\PYG{l+s+s2}{pe}\PYG{l+s+s2}{\PYGZdq{}}\PYG{p}{)}\PYG{p}{)}

\PYG{n}{MPI}\PYG{o}{.}\PYG{n}{Finalize}\PYG{p}{(}\PYG{p}{)}
\end{sphinxVerbatim}

\sphinxAtStartPar
To run this script (melt.py) in parallel using 4 MPI processes we invoke the
following mpirun command:

\begin{sphinxVerbatim}[commandchars=\\\{\}]
mpirun\PYG{+w}{ }\PYGZhy{}np\PYG{+w}{ }\PYG{l+m}{4}\PYG{+w}{ }python\PYG{+w}{ }melt.py
\end{sphinxVerbatim}


\subsubsection{Feedback and Contributing}
\label{\detokenize{Howto_pylammps:feedback-and-contributing}}
\sphinxAtStartPar
If you find this Python interface useful, please feel free to provide feedback
and ideas on how to improve it to Richard Berger (\sphinxhref{mailto:richard.berger@outlook.com}{richard.berger@outlook.com}). We also
want to encourage people to write tutorial style IPython notebooks showcasing LAMMPS usage
and maybe their latest research results.

\sphinxstepscope


\subsection{Using LAMMPS on Windows 10 with WSL}
\label{\detokenize{Howto_wsl:using-lammps-on-windows-10-with-wsl}}\label{\detokenize{Howto_wsl::doc}}
\sphinxAtStartPar
\sphinxstylestrong{written by Richard Berger}


\bigskip\hrule\bigskip


\sphinxAtStartPar
It’s always been tricky for us to have LAMMPS users and developers work on
Windows. We primarily develop LAMMPS to run on Linux clusters. To teach
LAMMPS in workshop settings, we had to redirect Windows users to
Linux Virtual Machines such as VirtualBox or Unix\sphinxhyphen{}like compilation with
Cygwin.

\sphinxAtStartPar
With the latest updates in Windows 10 (Version 2004, Build 19041 or
higher), Microsoft has added a new way to work on Linux\sphinxhyphen{}based code. The
\sphinxhref{https://learn.microsoft.com/en-us/windows/wsl/}{Windows Subsystem for Linux (WSL)}.  With WSL Version 2,
you now get a Linux Virtual Machine that transparently integrates into
Windows.  All you need is to ensure you have the latest Windows updates
installed and enable this new feature.  Linux VMs are then easily
installed using the Microsoft Store.

\sphinxAtStartPar
In this tutorial, I’ll show you how to set up and compile LAMMPS for both serial
and MPI usage in WSL2.


\subsubsection{Installation}
\label{\detokenize{Howto_wsl:installation}}

\paragraph{Upgrade to the latest Windows 10}
\label{\detokenize{Howto_wsl:upgrade-to-the-latest-windows-10}}
\sphinxAtStartPar
Type “Updates” in Windows Start and select “Check for Updates”.

\noindent\sphinxincludegraphics[scale=0.5]{{updates}.png}

\sphinxAtStartPar
Install all pending updates and reboot your system as many times as
necessary. Continue until your Windows installation is updated.

\noindent\sphinxincludegraphics[scale=0.5]{{windows_update}.png}

\sphinxAtStartPar
Verify your system has at least \sphinxstylestrong{version 2004 and build 19041 or later}. You
can find this information by clicking on “OS build info”.

\noindent\sphinxincludegraphics[scale=0.5]{{osinfo}.png}


\paragraph{Enable WSL}
\label{\detokenize{Howto_wsl:enable-wsl}}
\sphinxAtStartPar
Next, we must install two additional Windows features to enable WSL support.
Open a PowerShell window as an administrator. Type “PowerShell” in Windows
Start and select “Run as Administrator”.

\noindent\sphinxincludegraphics[scale=0.5]{{powershell}.png}

\sphinxAtStartPar
Windows will ask you for administrator access. After you accept a new command
line window will appear. Type in the following command to install WSL:

\begin{sphinxVerbatim}[commandchars=\\\{\}]
dism.exe /online /enable\PYGZhy{}feature /featurename:Microsoft\PYGZhy{}Windows\PYGZhy{}Subsystem\PYGZhy{}Linux /all /norestart
\end{sphinxVerbatim}

\noindent\sphinxincludegraphics{{wsl_install1}.png}

\sphinxAtStartPar
Next, enable the VirtualMachinePlatform feature using the following command:

\begin{sphinxVerbatim}[commandchars=\\\{\}]
dism.exe /online /enable\PYGZhy{}feature /featurename:VirtualMachinePlatform /all /norestart
\end{sphinxVerbatim}

\noindent\sphinxincludegraphics{{wsl_install2}.png}

\sphinxAtStartPar
Finally, reboot your system.


\paragraph{Update WSL kernel component}
\label{\detokenize{Howto_wsl:update-wsl-kernel-component}}
\sphinxAtStartPar
Download and install the \sphinxcode{\sphinxupquote{WSL Kernel Component Update}}.
Afterwards, reboot your system.


\paragraph{Set WSL2 as default}
\label{\detokenize{Howto_wsl:set-wsl2-as-default}}
\sphinxAtStartPar
Again, open PowerShell as administrator and run the following command:

\begin{sphinxVerbatim}[commandchars=\\\{\}]
\PYG{n}{wsl} \PYG{p}{\PYGZhy{}}\PYG{n}{\PYGZhy{}set}\PYG{n}{\PYGZhy{}default}\PYG{n}{\PYGZhy{}version} \PYG{n}{2}
\end{sphinxVerbatim}

\sphinxAtStartPar
This command ensures that all future Linux installations will use WSL version 2.

\noindent\sphinxincludegraphics{{wsl_install3}.png}


\paragraph{Install a Linux Distribution}
\label{\detokenize{Howto_wsl:install-a-linux-distribution}}
\sphinxAtStartPar
Next, we need to install a Linux distribution via the Microsoft Store.
Install \sphinxhref{ms-windows-store://pdp/?ProductId=9n6svws3rx71}{Ubuntu 20.04 LTS}.
Once installed, you can launch it like any other application from the Start
Menu.

\noindent\sphinxincludegraphics[scale=0.5]{{ubuntu_in_store}.png}


\subparagraph{Initial Setup}
\label{\detokenize{Howto_wsl:initial-setup}}
\sphinxAtStartPar
The first time you launch the Ubuntu Linux console, it will prompt you for a
UNIX username and password. You will need this password to perform \sphinxcode{\sphinxupquote{sudo}}
commands later. Once completed, your Linux shell is ready for use. All your
actions and commands will run as the Linux user you specified.

\noindent\sphinxincludegraphics[scale=0.5]{{first_login}.png}


\subsubsection{Windows Explorer / WSL integration}
\label{\detokenize{Howto_wsl:windows-explorer-wsl-integration}}
\sphinxAtStartPar
Your Linux installation will have its own Linux filesystem, which contains
the Ubuntu files. Your Linux user will have a regular Linux home directory in
\sphinxcode{\sphinxupquote{/home/\textless{}USERNAME\textgreater{}}}. This directory is different from your Windows User
directory. Windows and Linux filesystems are connected through WSL.

\sphinxAtStartPar
All hard drives in Windows are accessible in the \sphinxcode{\sphinxupquote{/mnt}} directory in Linux.
E.g., WSL maps the \sphinxcode{\sphinxupquote{C}} hard drive to the \sphinxcode{\sphinxupquote{/mnt/c}} directory. That means you
can access your Windows User directory in \sphinxcode{\sphinxupquote{/mnt/c/Users/\textless{}WINDOWS\_USERNAME\textgreater{}}}.

\sphinxAtStartPar
The Windows Explorer can also access the Linux filesystem. To illustrate this
integration, open an Ubuntu console and navigate to a directory of your
choice. To view this location in Windows Explorer, use the \sphinxcode{\sphinxupquote{explorer.exe .}}
command (do not forget the final dot!).

\noindent\sphinxincludegraphics[scale=0.5]{{wsl_integration}.png}


\bigskip\hrule\bigskip



\subsubsection{Compiling LAMMPS}
\label{\detokenize{Howto_wsl:compiling-lammps}}
\sphinxAtStartPar
You now have a fully functioning Ubuntu installation and can follow most
guides to install LAMMPS on a Linux system. Here are some of the essential
steps to follow:


\paragraph{Install prerequisite packages}
\label{\detokenize{Howto_wsl:install-prerequisite-packages}}
\sphinxAtStartPar
Before we can begin, we need to download the necessary compiler toolchain and
libraries to compile LAMMPS. In our Ubuntu\sphinxhyphen{}based Linux installation, we will
use the \sphinxcode{\sphinxupquote{apt}} package manager to install additional packages.

\sphinxAtStartPar
First, upgrade all existing packages using \sphinxcode{\sphinxupquote{apt update}} and \sphinxcode{\sphinxupquote{apt upgrade}}.

\begin{sphinxVerbatim}[commandchars=\\\{\}]
sudo\PYG{+w}{ }apt\PYG{+w}{ }update
sudo\PYG{+w}{ }apt\PYG{+w}{ }upgrade\PYG{+w}{ }\PYGZhy{}y
\end{sphinxVerbatim}

\sphinxAtStartPar
Next, install the following packages with \sphinxcode{\sphinxupquote{apt install}}:

\begin{sphinxVerbatim}[commandchars=\\\{\}]
sudo\PYG{+w}{ }apt\PYG{+w}{ }install\PYG{+w}{ }\PYGZhy{}y\PYG{+w}{ }cmake\PYG{+w}{ }build\PYGZhy{}essential\PYG{+w}{ }ccache\PYG{+w}{ }gfortran\PYG{+w}{ }openmpi\PYGZhy{}bin\PYG{+w}{ }libopenmpi\PYGZhy{}dev\PYG{+w}{ }\PYG{l+s+se}{\PYGZbs{}}
\PYG{+w}{                    }libfftw3\PYGZhy{}dev\PYG{+w}{ }libjpeg\PYGZhy{}dev\PYG{+w}{ }libpng\PYGZhy{}dev\PYG{+w}{ }python3\PYGZhy{}dev\PYG{+w}{ }python3\PYGZhy{}pip\PYG{+w}{ }\PYG{l+s+se}{\PYGZbs{}}
\PYG{+w}{                    }python3\PYGZhy{}virtualenv\PYG{+w}{ }libblas\PYGZhy{}dev\PYG{+w}{ }liblapack\PYGZhy{}dev\PYG{+w}{ }libhdf5\PYGZhy{}serial\PYGZhy{}dev\PYG{+w}{ }\PYG{l+s+se}{\PYGZbs{}}
\PYG{+w}{                    }hdf5\PYGZhy{}tools
\end{sphinxVerbatim}


\paragraph{Download LAMMPS}
\label{\detokenize{Howto_wsl:download-lammps}}
\sphinxAtStartPar
Obtain a copy of the LAMMPS source code and go into it using the \sphinxcode{\sphinxupquote{cd}} command.


\subparagraph{Option 1: Download a LAMMPS tarball using wget}
\label{\detokenize{Howto_wsl:option-1-download-a-lammps-tarball-using-wget}}
\begin{sphinxVerbatim}[commandchars=\\\{\}]
wget\PYG{+w}{ }https://github.com/lammps/lammps/archive/stable\PYGZus{}3Mar2020.tar.gz
tar\PYG{+w}{ }xvzf\PYG{+w}{ }stable\PYGZus{}3Mar2020.tar.gz
\PYG{n+nb}{cd}\PYG{+w}{ }lammps
\end{sphinxVerbatim}


\subparagraph{Option 2: Download a LAMMPS development version from GitHub}
\label{\detokenize{Howto_wsl:option-2-download-a-lammps-development-version-from-github}}
\begin{sphinxVerbatim}[commandchars=\\\{\}]
git\PYG{+w}{ }clone\PYG{+w}{ }\PYGZhy{}\PYGZhy{}depth\PYG{o}{=}\PYG{l+m}{1}\PYG{+w}{ }https://github.com/lammps/lammps.git
\PYG{n+nb}{cd}\PYG{+w}{ }lammps
\end{sphinxVerbatim}


\paragraph{Configure and Compile LAMMPS with CMake}
\label{\detokenize{Howto_wsl:configure-and-compile-lammps-with-cmake}}
\sphinxAtStartPar
A beginner\sphinxhyphen{}friendly way to compile LAMMPS is to use CMake. Create a \sphinxcode{\sphinxupquote{build}}
directory to compile LAMMPS and move into it. This directory will store the
build configuration and any binaries generated during compilation.

\begin{sphinxVerbatim}[commandchars=\\\{\}]
mkdir\PYG{+w}{ }build
\PYG{n+nb}{cd}\PYG{+w}{ }build
\end{sphinxVerbatim}

\sphinxAtStartPar
There are countless ways to compile LAMMPS. It is beyond the scope of this
tutorial. If you want to find out more about what can be enabled, please
consult the extensive \sphinxhref{https://docs.lammps.org/Build\_cmake.html}{documentation}.

\sphinxAtStartPar
To compile a minimal version of LAMMPS, we’re going to use a preset.
Presets are a way to specify a collection of CMake options using a file.

\begin{sphinxVerbatim}[commandchars=\\\{\}]
cmake\PYG{+w}{ }../cmake/presets/basic.cmake\PYG{+w}{ }../cmake
\end{sphinxVerbatim}

\sphinxAtStartPar
This command configures the build and generates the necessary Makefiles. To compile the binary, run the make command.

\begin{sphinxVerbatim}[commandchars=\\\{\}]
make\PYG{+w}{ }\PYGZhy{}j\PYG{+w}{ }\PYG{l+m}{4}
\end{sphinxVerbatim}

\sphinxAtStartPar
The \sphinxcode{\sphinxupquote{\sphinxhyphen{}j}} option specifies how many parallel processes will perform the
compilation. This option can significantly speed up compilation times. Use a
number that corresponds to the number of processors in your system.

\sphinxAtStartPar
After the compilation completes successfully, you will have an executable
called \sphinxcode{\sphinxupquote{lmp}} in the \sphinxcode{\sphinxupquote{build}} directory.

\noindent\sphinxincludegraphics[scale=0.5]{{compilation_result}.png}

\sphinxAtStartPar
Please take note of the absolute path of your \sphinxcode{\sphinxupquote{build}} directory. You will
need to know the location to execute the LAMMPS binary later.

\sphinxAtStartPar
One way of getting the absolute path of the current directory is through the
\sphinxcode{\sphinxupquote{\$PWD}} variable:

\begin{sphinxVerbatim}[commandchars=\\\{\}]
\PYG{c+c1}{\PYGZsh{} prints out the current value of the PWD variable}
\PYG{n+nb}{echo}\PYG{+w}{ }\PYG{n+nv}{\PYGZdl{}PWD}
\end{sphinxVerbatim}

\sphinxAtStartPar
Let us save this value in a temporary variable \sphinxcode{\sphinxupquote{LAMMPS\_BUILD\_DIR}} for future use:

\begin{sphinxVerbatim}[commandchars=\\\{\}]
\PYG{n+nv}{LAMMPS\PYGZus{}BUILD\PYGZus{}DIR}\PYG{o}{=}\PYG{n+nv}{\PYGZdl{}PWD}
\end{sphinxVerbatim}

\sphinxAtStartPar
The full path of the LAMMPS binary then is \sphinxcode{\sphinxupquote{\$LAMMPS\_BUILD\_DIR/lmp}}.


\bigskip\hrule\bigskip



\subsubsection{Running an example script}
\label{\detokenize{Howto_wsl:running-an-example-script}}
\sphinxAtStartPar
Now that we have a LAMMPS binary, we will run a script from the examples folder.

\sphinxAtStartPar
Switch into the \sphinxcode{\sphinxupquote{examples/melt}} folder:

\begin{sphinxVerbatim}[commandchars=\\\{\}]
cd ../examples/melt
\end{sphinxVerbatim}

\sphinxAtStartPar
To run this example in serial, use the following command line:

\begin{sphinxVerbatim}[commandchars=\\\{\}]
\PYGZdl{}LAMMPS\PYGZus{}BUILD\PYGZus{}DIR/lmp \PYGZhy{}in in.melt
\end{sphinxVerbatim}

\sphinxAtStartPar
To run the same script in parallel using MPI with 4 processes, do the following:

\begin{sphinxVerbatim}[commandchars=\\\{\}]
mpirun\PYG{+w}{ }\PYGZhy{}np\PYG{+w}{ }\PYG{l+m}{4}\PYG{+w}{ }\PYG{n+nv}{\PYGZdl{}LAMMPS\PYGZus{}BUILD\PYGZus{}DIR}/lmp\PYG{+w}{ }\PYGZhy{}in\PYG{+w}{ }\PYG{k}{in}.melt
\end{sphinxVerbatim}

\sphinxAtStartPar
If you run LAMMPS for the first time, the Windows Firewall might prompt you
to confirm access. LAMMPS is accessing the network stack to enable parallel
computation. Allow the access.

\noindent\sphinxincludegraphics[scale=0.75]{{windows_firewall}.png}

\sphinxAtStartPar
In either serial or MPI case, LAMMPS executes and will output something similar to this:

\begin{sphinxVerbatim}[commandchars=\\\{\}]
LAMMPS (30 Jun 2020)
...
...
...
Total \PYGZsh{} of neighbors = 151513
Ave neighs/atom = 37.878250
Neighbor list builds = 12
Dangerous builds not checked
Total wall time: 0:00:00
\end{sphinxVerbatim}

\sphinxAtStartPar
\sphinxstylestrong{Congratulations! You’ve successfully compiled and executed LAMMPS on WSL!}


\subsubsection{Final steps}
\label{\detokenize{Howto_wsl:final-steps}}
\sphinxAtStartPar
It is cumbersome to always specify the path of your LAMMPS binary. You can
avoid this by adding the absolute path of your \sphinxcode{\sphinxupquote{build}} directory to your PATH
environment variable.

\begin{sphinxVerbatim}[commandchars=\\\{\}]
\PYG{n+nb}{export}\PYG{+w}{ }\PYG{n+nv}{PATH}\PYG{o}{=}\PYG{n+nv}{\PYGZdl{}LAMMPS\PYGZus{}BUILD\PYGZus{}DIR}:\PYG{n+nv}{\PYGZdl{}PATH}
\end{sphinxVerbatim}

\sphinxAtStartPar
You can then run LAMMPS input scripts like this:

\begin{sphinxVerbatim}[commandchars=\\\{\}]
lmp\PYG{+w}{ }\PYGZhy{}in\PYG{+w}{ }\PYG{k}{in}.melt
\end{sphinxVerbatim}

\sphinxAtStartPar
or

\begin{sphinxVerbatim}[commandchars=\\\{\}]
mpirun\PYG{+w}{ }\PYGZhy{}np\PYG{+w}{ }\PYG{l+m}{4}\PYG{+w}{ }lmp\PYG{+w}{ }\PYGZhy{}in\PYG{+w}{ }\PYG{k}{in}.melt
\end{sphinxVerbatim}

\begin{sphinxadmonition}{note}{Note:}
\sphinxAtStartPar
The value of this \sphinxcode{\sphinxupquote{PATH}} variable will disappear once you close your
console window. To persist this setting edit the \sphinxcode{\sphinxupquote{\$HOME/.bashrc}} file using your
favorite text editor and add this line:

\begin{sphinxVerbatim}[commandchars=\\\{\}]
\PYG{n+nb}{export}\PYG{+w}{ }\PYG{n+nv}{PATH}\PYG{o}{=}/full/path/to/your/lammps/build:\PYG{n+nv}{\PYGZdl{}PATH}
\end{sphinxVerbatim}

\sphinxAtStartPar
\sphinxstylestrong{Example:}
If the LAMMPS executable \sphinxtitleref{lmp} has the following absolute path:

\begin{sphinxVerbatim}[commandchars=\\\{\}]
/home/\PYGZlt{}USERNAME\PYGZgt{}/lammps/build/lmp
\end{sphinxVerbatim}

\sphinxAtStartPar
the \sphinxcode{\sphinxupquote{PATH}} variable should be:

\begin{sphinxVerbatim}[commandchars=\\\{\}]
\PYG{n+nb}{export}\PYG{+w}{ }\PYG{n+nv}{PATH}\PYG{o}{=}/home/\PYGZlt{}USERNAME\PYGZgt{}/lammps/build:\PYG{n+nv}{\PYGZdl{}PATH}
\end{sphinxVerbatim}

\sphinxAtStartPar
Once set up, all your Ubuntu consoles will always have access to your \sphinxcode{\sphinxupquote{lmp}}
binary without having to specify its location.
\end{sphinxadmonition}


\subsubsection{Conclusion}
\label{\detokenize{Howto_wsl:conclusion}}
\sphinxAtStartPar
I hope this gives you good overview on how to start compiling and running LAMMPS on
Windows. WSL makes preparing and running scripts on Windows a much better
experience.

\sphinxAtStartPar
If you are completely new to Linux, I highly recommend investing some time in
studying Linux online tutorials. E.g., tutorials about Bash Shell and Basic
Unix commands (e.g., \sphinxhref{https://linuxjourney.com/}{Linux Journey}). Acquiring
these skills will make you much more productive in this environment.


\begin{sphinxseealso}{See also:}
\begin{itemize}
\item {} 
\sphinxAtStartPar
\sphinxhref{https://docs.microsoft.com/en-us/windows/wsl/}{Windows Subsystem for Linux Documentation}

\end{itemize}


\end{sphinxseealso}


\sphinxstepscope


\chapter{Example scripts}
\label{\detokenize{Examples:example-scripts}}\label{\detokenize{Examples::doc}}
\sphinxAtStartPar
The LAMMPS distribution includes an examples subdirectory with many
sample problems.  Many are 2d models that run quickly and are
straightforward to visualize, requiring at most a couple of minutes to
run on a desktop machine.  Each problem has an input script (in.*) and
produces a log file (log.*) when it runs.  Some use a data file
(data.*) of initial coordinates as additional input.  A few sample log
file run on different machines and different numbers of processors are
included in the directories to compare your answers to.  E.g. a log
file like log.date.crack.foo.P means the “crack” example was run on P
processors of machine “foo” on that date (i.e. with that version of
LAMMPS).

\sphinxAtStartPar
Many of the input files have commented\sphinxhyphen{}out lines for creating dump
files and image files.

\sphinxAtStartPar
If you uncomment the {\hyperref[\detokenize{dump::doc}]{\sphinxcrossref{\DUrole{doc}{dump}}}} command in the input script, a
text dump file will be produced, which can be animated by various
\sphinxhref{https://www.lammps.org/viz.html}{visualization programs}.

\sphinxAtStartPar
If you uncomment the {\hyperref[\detokenize{dump::doc}]{\sphinxcrossref{\DUrole{doc}{dump image}}}} command in the input
script, and assuming you have built LAMMPS with a JPG library, JPG
snapshot images will be produced when the simulation runs.  They can
be quickly post\sphinxhyphen{}processed into a movie using commands described on the
{\hyperref[\detokenize{dump_image::doc}]{\sphinxcrossref{\DUrole{doc}{dump image}}}} doc page.

\sphinxAtStartPar
Animations of many of the examples can be viewed on the Movies section
of the \sphinxhref{https://www.lammps.org/movies.html}{LAMMPS website}.

\sphinxAtStartPar
There are two kinds of subdirectories in the examples folder.  Lower
case named directories contain one or a few simple, quick\sphinxhyphen{}to\sphinxhyphen{}run
problems.  Upper case named directories contain up to several complex
scripts that illustrate a particular kind of simulation method or
model.  Some of these run for longer times, e.g. to measure a
particular quantity.

\sphinxAtStartPar
Lists of both kinds of directories are given below.


\bigskip\hrule\bigskip



\section{Lowercase directories}
\label{\detokenize{Examples:lowercase-directories}}

\begin{savenotes}
\sphinxatlongtablestart
\sphinxthistablewithglobalstyle
\makeatletter
  \LTleft \@totalleftmargin plus1fill
  \LTright\dimexpr\columnwidth-\@totalleftmargin-\linewidth\relax plus1fill
\makeatother
\begin{longtable}{ll}
\sphinxtoprule
\endfirsthead

\multicolumn{2}{c}{\sphinxnorowcolor
    \makebox[0pt]{\sphinxtablecontinued{\tablename\ \thetable{} \textendash{} continued from previous page}}%
}\\
\sphinxtoprule
\endhead

\sphinxbottomrule
\multicolumn{2}{r}{\sphinxnorowcolor
    \makebox[0pt][r]{\sphinxtablecontinued{continues on next page}}%
}\\
\endfoot

\endlastfoot
\sphinxtableatstartofbodyhook

\sphinxAtStartPar
accelerate
&
\sphinxAtStartPar
run with various acceleration options (OpenMP, GPU, Phi)
\\
\sphinxhline
\sphinxAtStartPar
airebo
&
\sphinxAtStartPar
polyethylene with AIREBO potential
\\
\sphinxhline
\sphinxAtStartPar
atm
&
\sphinxAtStartPar
Axilrod\sphinxhyphen{}Teller\sphinxhyphen{}Muto potential example
\\
\sphinxhline
\sphinxAtStartPar
balance
&
\sphinxAtStartPar
dynamic load balancing, 2d system
\\
\sphinxhline
\sphinxAtStartPar
body
&
\sphinxAtStartPar
body particles, 2d system
\\
\sphinxhline
\sphinxAtStartPar
bpm
&
\sphinxAtStartPar
BPM simulations of pouring elastic grains and plate impact
\\
\sphinxhline
\sphinxAtStartPar
cmap
&
\sphinxAtStartPar
CMAP 5\sphinxhyphen{}body contributions to CHARMM force field
\\
\sphinxhline
\sphinxAtStartPar
colloid
&
\sphinxAtStartPar
big colloid particles in a small particle solvent, 2d system
\\
\sphinxhline
\sphinxAtStartPar
comb
&
\sphinxAtStartPar
models using the COMB potential
\\
\sphinxhline
\sphinxAtStartPar
controller
&
\sphinxAtStartPar
use of fix controller as a thermostat
\\
\sphinxhline
\sphinxAtStartPar
coreshell
&
\sphinxAtStartPar
core/shell model using CORESHELL package
\\
\sphinxhline
\sphinxAtStartPar
crack
&
\sphinxAtStartPar
crack propagation in a 2d solid
\\
\sphinxhline
\sphinxAtStartPar
deposit
&
\sphinxAtStartPar
deposit atoms and molecules on a surface
\\
\sphinxhline
\sphinxAtStartPar
dipole
&
\sphinxAtStartPar
point dipolar particles, 2d system
\\
\sphinxhline
\sphinxAtStartPar
dreiding
&
\sphinxAtStartPar
methanol via Dreiding FF
\\
\sphinxhline
\sphinxAtStartPar
eim
&
\sphinxAtStartPar
NaCl using the EIM potential
\\
\sphinxhline
\sphinxAtStartPar
ellipse
&
\sphinxAtStartPar
ellipsoidal particles in spherical solvent, 2d system
\\
\sphinxhline
\sphinxAtStartPar
flow
&
\sphinxAtStartPar
Couette and Poiseuille flow in a 2d channel
\\
\sphinxhline
\sphinxAtStartPar
friction
&
\sphinxAtStartPar
frictional contact of spherical asperities between 2d surfaces
\\
\sphinxhline
\sphinxAtStartPar
mc
&
\sphinxAtStartPar
Monte Carlo features via fix gcmc, widom and other commands
\\
\sphinxhline
\sphinxAtStartPar
granregion
&
\sphinxAtStartPar
use of fix wall/region/gran as boundary on granular particles
\\
\sphinxhline
\sphinxAtStartPar
hugoniostat
&
\sphinxAtStartPar
Hugoniostat shock dynamics
\\
\sphinxhline
\sphinxAtStartPar
hyper
&
\sphinxAtStartPar
global and local hyperdynamics of diffusion on Pt surface
\\
\sphinxhline
\sphinxAtStartPar
indent
&
\sphinxAtStartPar
spherical indenter into a 2d solid
\\
\sphinxhline
\sphinxAtStartPar
kim
&
\sphinxAtStartPar
use of potentials from the \sphinxhref{https://openkim.org}{OpenKIM Repository}
\\
\sphinxhline
\sphinxAtStartPar
mdi
&
\sphinxAtStartPar
use of the MDI package and MolSSI MDI code coupling library
\\
\sphinxhline
\sphinxAtStartPar
meam
&
\sphinxAtStartPar
MEAM test for SiC and shear (same as shear examples)
\\
\sphinxhline
\sphinxAtStartPar
melt
&
\sphinxAtStartPar
rapid melt of 3d LJ system
\\
\sphinxhline
\sphinxAtStartPar
micelle
&
\sphinxAtStartPar
self\sphinxhyphen{}assembly of small lipid\sphinxhyphen{}like molecules into 2d bilayers
\\
\sphinxhline
\sphinxAtStartPar
min
&
\sphinxAtStartPar
energy minimization of 2d LJ melt
\\
\sphinxhline
\sphinxAtStartPar
msst
&
\sphinxAtStartPar
MSST shock dynamics
\\
\sphinxhline
\sphinxAtStartPar
multi
&
\sphinxAtStartPar
multi neighboring for systems with large interaction disparities
\\
\sphinxhline
\sphinxAtStartPar
nb3b
&
\sphinxAtStartPar
use of non\sphinxhyphen{}bonded 3\sphinxhyphen{}body harmonic pair style
\\
\sphinxhline
\sphinxAtStartPar
neb
&
\sphinxAtStartPar
nudged elastic band (NEB) calculation for barrier finding
\\
\sphinxhline
\sphinxAtStartPar
nemd
&
\sphinxAtStartPar
non\sphinxhyphen{}equilibrium MD of 2d sheared system
\\
\sphinxhline
\sphinxAtStartPar
obstacle
&
\sphinxAtStartPar
flow around two voids in a 2d channel
\\
\sphinxhline
\sphinxAtStartPar
peptide
&
\sphinxAtStartPar
dynamics of a small solvated peptide chain (5\sphinxhyphen{}mer)
\\
\sphinxhline
\sphinxAtStartPar
peri
&
\sphinxAtStartPar
Peridynamic model of cylinder impacted by indenter
\\
\sphinxhline
\sphinxAtStartPar
pour
&
\sphinxAtStartPar
pouring of granular particles into a 3d box, then chute flow
\\
\sphinxhline
\sphinxAtStartPar
prd
&
\sphinxAtStartPar
parallel replica dynamics of vacancy diffusion in bulk Si
\\
\sphinxhline
\sphinxAtStartPar
python
&
\sphinxAtStartPar
using embedded Python in a LAMMPS input script
\\
\sphinxhline
\sphinxAtStartPar
qeq
&
\sphinxAtStartPar
use of the QEQ package for charge equilibration
\\
\sphinxhline
\sphinxAtStartPar
rdf\sphinxhyphen{}adf
&
\sphinxAtStartPar
computing radial and angle distribution functions for water
\\
\sphinxhline
\sphinxAtStartPar
reax
&
\sphinxAtStartPar
RDX and TATB models using the ReaxFF
\\
\sphinxhline
\sphinxAtStartPar
rerun
&
\sphinxAtStartPar
use of rerun and read\_dump commands
\\
\sphinxhline
\sphinxAtStartPar
rheo
&
\sphinxAtStartPar
RHEO simulations of fluid flows and phase transitions
\\
\sphinxhline
\sphinxAtStartPar
rigid
&
\sphinxAtStartPar
rigid bodies modeled as independent or coupled
\\
\sphinxhline
\sphinxAtStartPar
shear
&
\sphinxAtStartPar
sideways shear applied to 2d solid, with and without a void
\\
\sphinxhline
\sphinxAtStartPar
snap
&
\sphinxAtStartPar
NVE dynamics for BCC tantalum crystal using SNAP potential
\\
\sphinxhline
\sphinxAtStartPar
srd
&
\sphinxAtStartPar
stochastic rotation dynamics (SRD) particles as solvent
\\
\sphinxhline
\sphinxAtStartPar
streitz
&
\sphinxAtStartPar
use of Streitz/Mintmire potential with charge equilibration
\\
\sphinxhline
\sphinxAtStartPar
tad
&
\sphinxAtStartPar
temperature\sphinxhyphen{}accelerated dynamics of vacancy diffusion in bulk Si
\\
\sphinxhline
\sphinxAtStartPar
threebody
&
\sphinxAtStartPar
regression test input for a variety of manybody potentials
\\
\sphinxhline
\sphinxAtStartPar
tracker
&
\sphinxAtStartPar
track interactions in LJ melt
\\
\sphinxhline
\sphinxAtStartPar
vashishta
&
\sphinxAtStartPar
use of the Vashishta potential
\\
\sphinxhline
\sphinxAtStartPar
voronoi
&
\sphinxAtStartPar
Voronoi tesselation via compute voronoi/atom command
\\
\sphinxbottomrule
\end{longtable}
\sphinxtableafterendhook
\sphinxatlongtableend
\end{savenotes}

\sphinxAtStartPar
Here is how you can run and visualize one of the sample problems:

\begin{sphinxVerbatim}[commandchars=\\\{\}]
\PYG{n+nb}{cd}\PYG{+w}{ }indent
cp\PYG{+w}{ }../../src/lmp\PYGZus{}linux\PYG{+w}{ }.\PYG{+w}{           }\PYG{c+c1}{\PYGZsh{} copy LAMMPS executable to this dir}
lmp\PYGZus{}linux\PYG{+w}{ }\PYGZhy{}in\PYG{+w}{ }\PYG{k}{in}.indent\PYG{+w}{            }\PYG{c+c1}{\PYGZsh{} run the problem}
\end{sphinxVerbatim}

\sphinxAtStartPar
Running the simulation produces the files \sphinxstyleemphasis{dump.indent} and
\sphinxstyleemphasis{log.lammps}.  You can visualize the dump file of snapshots with a
variety of third\sphinxhyphen{}party tools highlighted on the
\sphinxhref{https://www.lammps.org/viz.html}{Visualization} page of the LAMMPS
website.

\sphinxAtStartPar
If you uncomment the {\hyperref[\detokenize{dump_image::doc}]{\sphinxcrossref{\DUrole{doc}{dump image}}}} line(s) in the input
script a series of JPG images will be produced by the run (assuming
you built LAMMPS with JPG support; see the
{\hyperref[\detokenize{Build_settings::doc}]{\sphinxcrossref{\DUrole{doc}{Build\_settings}}}} page for details).  These can
be viewed individually or turned into a movie or animated by tools
like ImageMagick or QuickTime or various Windows\sphinxhyphen{}based tools.  See the
{\hyperref[\detokenize{dump_image::doc}]{\sphinxcrossref{\DUrole{doc}{dump image}}}} page for more details.  E.g. this
Imagemagick command would create a GIF file suitable for viewing in a
browser.

\begin{sphinxVerbatim}[commandchars=\\\{\}]
\PYGZpc{}\PYG{+w}{ }convert\PYG{+w}{ }\PYGZhy{}loop\PYG{+w}{ }\PYG{l+m}{1}\PYG{+w}{ }*.jpg\PYG{+w}{ }foo.gif
\end{sphinxVerbatim}


\bigskip\hrule\bigskip



\section{Uppercase directories}
\label{\detokenize{Examples:uppercase-directories}}

\begin{savenotes}\sphinxattablestart
\sphinxthistablewithglobalstyle
\centering
\begin{tabulary}{\linewidth}[t]{TT}
\sphinxtoprule
\sphinxtableatstartofbodyhook
\sphinxAtStartPar
ASPHERE
&
\sphinxAtStartPar
various aspherical particle models, using ellipsoids, rigid bodies, line/triangle particles, etc
\\
\sphinxhline
\sphinxAtStartPar
COUPLE
&
\sphinxAtStartPar
examples of how to use LAMMPS as a library
\\
\sphinxhline
\sphinxAtStartPar
DIFFUSE
&
\sphinxAtStartPar
compute diffusion coefficients via several methods
\\
\sphinxhline
\sphinxAtStartPar
ELASTIC
&
\sphinxAtStartPar
compute elastic constants at zero temperature
\\
\sphinxhline
\sphinxAtStartPar
ELASTIC\_T
&
\sphinxAtStartPar
compute elastic constants at finite temperature
\\
\sphinxhline
\sphinxAtStartPar
HEAT
&
\sphinxAtStartPar
compute thermal conductivity for LJ and water via fix ehex
\\
\sphinxhline
\sphinxAtStartPar
KAPPA
&
\sphinxAtStartPar
compute thermal conductivity via several methods
\\
\sphinxhline
\sphinxAtStartPar
MC\sphinxhyphen{}LOOP
&
\sphinxAtStartPar
using LAMMPS in a Monte Carlo mode to relax the energy of a system in a input script loop
\\
\sphinxhline
\sphinxAtStartPar
PACKAGES
&
\sphinxAtStartPar
examples for specific packages and contributed commands
\\
\sphinxhline
\sphinxAtStartPar
SPIN
&
\sphinxAtStartPar
examples for features of the SPIN package
\\
\sphinxhline
\sphinxAtStartPar
UNITS
&
\sphinxAtStartPar
examples that run the same simulation in lj, real, metal units
\\
\sphinxhline
\sphinxAtStartPar
VISCOSITY
&
\sphinxAtStartPar
compute viscosity via several methods
\\
\sphinxbottomrule
\end{tabulary}
\sphinxtableafterendhook\par
\sphinxattableend\end{savenotes}

\sphinxAtStartPar
Nearly all of these directories have README files which give more
details on how to understand and use their contents.

\sphinxAtStartPar
The PACKAGES directory has a large number of subdirectories which
correspond by name to specific packages.  They contain scripts that
illustrate how to use the command(s) provided in those packages.  Many
of the subdirectories have their own README files which give further
instructions.  See the {\hyperref[\detokenize{Packages_details::doc}]{\sphinxcrossref{\DUrole{doc}{Packages\_details}}}} doc
page for more info on specific packages.

\sphinxstepscope


\chapter{Auxiliary tools}
\label{\detokenize{Tools:auxiliary-tools}}\label{\detokenize{Tools::doc}}
\sphinxAtStartPar
LAMMPS is designed to be a computational kernel for performing
molecular dynamics computations.  Additional pre\sphinxhyphen{} and post\sphinxhyphen{}processing
steps are often necessary to setup and analyze a simulation.  A list
of such tools can be found on the \sphinxhref{https://www.lammps.org}{LAMMPS webpage} at these links:
\begin{itemize}
\item {} 
\sphinxAtStartPar
\sphinxhref{https://www.lammps.org/prepost.html}{Pre/Post processing}

\item {} 
\sphinxAtStartPar
\sphinxhref{https://www.lammps.org/external.html}{External LAMMPS packages \& tools}

\item {} 
\sphinxAtStartPar
\sphinxhref{https://lammps.github.io/pizza}{Pizza.py toolkit}

\end{itemize}

\sphinxAtStartPar
The last link for \sphinxhref{https://lammps.github.io/pizza}{Pizza.py} is a Python\sphinxhyphen{}based tool developed at
Sandia which provides tools for doing setup, analysis, plotting, and
visualization for LAMMPS simulations.

\sphinxAtStartPar
Additional tools included in the LAMMPS distribution are described on
this page.

\sphinxAtStartPar
Note that many users write their own setup or analysis tools or use
other existing codes and convert their output to a LAMMPS input format
or vice versa.  The tools listed here are included in the LAMMPS
distribution as examples of auxiliary tools.  Some of them are not
actively supported by the LAMMPS developers, as they were contributed
by LAMMPS users.  If you have problems using them, we can direct you
to the authors.

\sphinxAtStartPar
The source code for each of these codes is in the tools subdirectory
of the LAMMPS distribution.  There is a Makefile (which you may need
to edit for your platform) which will build several of the tools which
reside in that directory.  Most of them are larger packages in their
own subdirectories with their own Makefiles and/or README files.


\bigskip\hrule\bigskip



\section{Pre\sphinxhyphen{}processing tools}
\label{\detokenize{Tools:pre-processing-tools}}

\begin{savenotes}\sphinxattablestart
\sphinxthistablewithglobalstyle
\centering
\begin{tabulary}{\linewidth}[t]{TTTTTT}
\sphinxtoprule
\sphinxtableatstartofbodyhook
\sphinxAtStartPar
{\hyperref[\detokenize{Tools:amber}]{\sphinxcrossref{\DUrole{std,std-ref}{amber2lmp}}}}
&
\sphinxAtStartPar
{\hyperref[\detokenize{Tools:charmm}]{\sphinxcrossref{\DUrole{std,std-ref}{ch2lmp}}}}
&
\sphinxAtStartPar
{\hyperref[\detokenize{Tools:chain}]{\sphinxcrossref{\DUrole{std,std-ref}{chain}}}}
&
\sphinxAtStartPar
{\hyperref[\detokenize{Tools:createatoms}]{\sphinxcrossref{\DUrole{std,std-ref}{createatoms}}}}
&
\sphinxAtStartPar
{\hyperref[\detokenize{Tools:drude}]{\sphinxcrossref{\DUrole{std,std-ref}{drude}}}}
&
\sphinxAtStartPar
{\hyperref[\detokenize{Tools:eamdb}]{\sphinxcrossref{\DUrole{std,std-ref}{eam database}}}}
\\
\sphinxhline
\sphinxAtStartPar
{\hyperref[\detokenize{Tools:eamgn}]{\sphinxcrossref{\DUrole{std,std-ref}{eam generate}}}}
&
\sphinxAtStartPar
{\hyperref[\detokenize{Tools:eff}]{\sphinxcrossref{\DUrole{std,std-ref}{eff}}}}
&
\sphinxAtStartPar
{\hyperref[\detokenize{Tools:ipp}]{\sphinxcrossref{\DUrole{std,std-ref}{ipp}}}}
&
\sphinxAtStartPar
{\hyperref[\detokenize{Tools:micelle}]{\sphinxcrossref{\DUrole{std,std-ref}{micelle2d}}}}
&
\sphinxAtStartPar
{\hyperref[\detokenize{Tools:moltemplate}]{\sphinxcrossref{\DUrole{std,std-ref}{moltemplate}}}}
&
\sphinxAtStartPar
{\hyperref[\detokenize{Tools:msi}]{\sphinxcrossref{\DUrole{std,std-ref}{msi2lmp}}}}
\\
\sphinxhline
\sphinxAtStartPar
{\hyperref[\detokenize{Tools:polybond}]{\sphinxcrossref{\DUrole{std,std-ref}{polybond}}}}
&
\sphinxAtStartPar
{\hyperref[\detokenize{Tools:stlconvert}]{\sphinxcrossref{\DUrole{std,std-ref}{stl\_bin2txt}}}}
&
\sphinxAtStartPar
{\hyperref[\detokenize{Tools:tabulate}]{\sphinxcrossref{\DUrole{std,std-ref}{tabulate}}}}
&
\sphinxAtStartPar
{\hyperref[\detokenize{Tools:tinker}]{\sphinxcrossref{\DUrole{std,std-ref}{tinker}}}}
&&\\
\sphinxbottomrule
\end{tabulary}
\sphinxtableafterendhook\par
\sphinxattableend\end{savenotes}


\section{Post\sphinxhyphen{}processing tools}
\label{\detokenize{Tools:post-processing-tools}}

\begin{savenotes}\sphinxattablestart
\sphinxthistablewithglobalstyle
\centering
\begin{tabulary}{\linewidth}[t]{TTTTTT}
\sphinxtoprule
\sphinxtableatstartofbodyhook
\sphinxAtStartPar
{\hyperref[\detokenize{Tools:amber}]{\sphinxcrossref{\DUrole{std,std-ref}{amber2lmp}}}}
&
\sphinxAtStartPar
{\hyperref[\detokenize{Tools:binary}]{\sphinxcrossref{\DUrole{std,std-ref}{binary2txt}}}}
&
\sphinxAtStartPar
{\hyperref[\detokenize{Tools:charmm}]{\sphinxcrossref{\DUrole{std,std-ref}{ch2lmp}}}}
&
\sphinxAtStartPar
{\hyperref[\detokenize{Tools:colvars-tools}]{\sphinxcrossref{\DUrole{std,std-ref}{colvars}}}}
&
\sphinxAtStartPar
{\hyperref[\detokenize{Tools:eff}]{\sphinxcrossref{\DUrole{std,std-ref}{eff}}}}
&
\sphinxAtStartPar
{\hyperref[\detokenize{Tools:fep}]{\sphinxcrossref{\DUrole{std,std-ref}{fep}}}}
\\
\sphinxhline
\sphinxAtStartPar
{\hyperref[\detokenize{Tools:arc}]{\sphinxcrossref{\DUrole{std,std-ref}{lmp2arc}}}}
&
\sphinxAtStartPar
{\hyperref[\detokenize{Tools:cfg}]{\sphinxcrossref{\DUrole{std,std-ref}{lmp2cfg}}}}
&
\sphinxAtStartPar
{\hyperref[\detokenize{Tools:matlab}]{\sphinxcrossref{\DUrole{std,std-ref}{matlab}}}}
&
\sphinxAtStartPar
{\hyperref[\detokenize{Tools:phonon}]{\sphinxcrossref{\DUrole{std,std-ref}{phonon}}}}
&
\sphinxAtStartPar
{\hyperref[\detokenize{Tools:pymol}]{\sphinxcrossref{\DUrole{std,std-ref}{pymol\_asphere}}}}
&
\sphinxAtStartPar
{\hyperref[\detokenize{Tools:pythontools}]{\sphinxcrossref{\DUrole{std,std-ref}{python}}}}
\\
\sphinxhline
\sphinxAtStartPar
{\hyperref[\detokenize{Tools:replica}]{\sphinxcrossref{\DUrole{std,std-ref}{replica}}}}
&
\sphinxAtStartPar
{\hyperref[\detokenize{Tools:smd}]{\sphinxcrossref{\DUrole{std,std-ref}{smd}}}}
&
\sphinxAtStartPar
{\hyperref[\detokenize{Tools:spin}]{\sphinxcrossref{\DUrole{std,std-ref}{spin}}}}
&
\sphinxAtStartPar
{\hyperref[\detokenize{Tools:xmgrace}]{\sphinxcrossref{\DUrole{std,std-ref}{xmgrace}}}}
&&\\
\sphinxbottomrule
\end{tabulary}
\sphinxtableafterendhook\par
\sphinxattableend\end{savenotes}


\section{Miscellaneous tools}
\label{\detokenize{Tools:miscellaneous-tools}}

\begin{savenotes}\sphinxattablestart
\sphinxthistablewithglobalstyle
\centering
\begin{tabulary}{\linewidth}[t]{TTTTTT}
\sphinxtoprule
\sphinxtableatstartofbodyhook
\sphinxAtStartPar
{\hyperref[\detokenize{Tools:coding-standard}]{\sphinxcrossref{\DUrole{std,std-ref}{LAMMPS coding standards}}}}
&
\sphinxAtStartPar
{\hyperref[\detokenize{Tools:emacs}]{\sphinxcrossref{\DUrole{std,std-ref}{emacs}}}}
&
\sphinxAtStartPar
{\hyperref[\detokenize{Tools:ipi}]{\sphinxcrossref{\DUrole{std,std-ref}{i\sphinxhyphen{}PI}}}}
&
\sphinxAtStartPar
{\hyperref[\detokenize{Tools:kate}]{\sphinxcrossref{\DUrole{std,std-ref}{kate}}}}
&
\sphinxAtStartPar
{\hyperref[\detokenize{Tools:lammps-gui}]{\sphinxcrossref{\DUrole{std,std-ref}{LAMMPS\sphinxhyphen{}GUI}}}}
&
\sphinxAtStartPar
{\hyperref[\detokenize{Tools:magic}]{\sphinxcrossref{\DUrole{std,std-ref}{LAMMPS magic patterns for file(1)}}}}
\\
\sphinxhline
\sphinxAtStartPar
{\hyperref[\detokenize{Tools:offline}]{\sphinxcrossref{\DUrole{std,std-ref}{Offline build tool}}}}
&
\sphinxAtStartPar
{\hyperref[\detokenize{Tools:regression}]{\sphinxcrossref{\DUrole{std,std-ref}{Regression tester}}}}
&
\sphinxAtStartPar
{\hyperref[\detokenize{Tools:singularity-tool}]{\sphinxcrossref{\DUrole{std,std-ref}{singularity/apptainer}}}}
&
\sphinxAtStartPar
{\hyperref[\detokenize{Tools:swig}]{\sphinxcrossref{\DUrole{std,std-ref}{SWIG interface}}}}
&
\sphinxAtStartPar
{\hyperref[\detokenize{Tools:valgrind}]{\sphinxcrossref{\DUrole{std,std-ref}{valgrind}}}}
&
\sphinxAtStartPar
{\hyperref[\detokenize{Tools:vim}]{\sphinxcrossref{\DUrole{std,std-ref}{vim}}}}
\\
\sphinxbottomrule
\end{tabulary}
\sphinxtableafterendhook\par
\sphinxattableend\end{savenotes}


\bigskip\hrule\bigskip



\section{Tool descriptions}
\label{\detokenize{Tools:tool-descriptions}}

\subsection{amber2lmp tool}
\label{\detokenize{Tools:amber2lmp-tool}}\label{\detokenize{Tools:amber}}
\sphinxAtStartPar
The amber2lmp subdirectory contains three Python scripts for converting
files back\sphinxhyphen{}and\sphinxhyphen{}forth between the AMBER MD code and LAMMPS.  See the
README file in amber2lmp for more information.

\sphinxAtStartPar
These tools were written by Keir Novik while he was at Queen Mary
University of London.  Keir is no longer there and cannot support
these tools which are out\sphinxhyphen{}of\sphinxhyphen{}date with respect to the current LAMMPS
version (and maybe with respect to AMBER as well).  Since we don’t use
these tools at Sandia, you will need to experiment with them and make
necessary modifications yourself.


\bigskip\hrule\bigskip



\subsection{binary2txt tool}
\label{\detokenize{Tools:binary2txt-tool}}\label{\detokenize{Tools:binary}}
\sphinxAtStartPar
The file binary2txt.cpp converts one or more binary LAMMPS dump file
into ASCII text files.  The syntax for running the tool is

\begin{sphinxVerbatim}[commandchars=\\\{\}]
binary2txt\PYG{+w}{ }file1\PYG{+w}{ }file2\PYG{+w}{ }...
\end{sphinxVerbatim}

\sphinxAtStartPar
which creates file1.txt, file2.txt, etc.  This tool must be compiled
on a platform that can read the binary file created by a LAMMPS run,
since binary files are not compatible across all platforms.


\bigskip\hrule\bigskip



\subsection{ch2lmp tool}
\label{\detokenize{Tools:ch2lmp-tool}}\label{\detokenize{Tools:charmm}}
\sphinxAtStartPar
The ch2lmp subdirectory contains tools for converting files
back\sphinxhyphen{}and\sphinxhyphen{}forth between the CHARMM MD code and LAMMPS.

\sphinxAtStartPar
They are intended to make it easy to use CHARMM as a builder and as a
post\sphinxhyphen{}processor for LAMMPS. Using charmm2lammps.pl, you can convert a
PDB file with associated CHARMM info, including CHARMM force field
data, into its LAMMPS equivalent. Support for the CMAP correction of
CHARMM22 and later is available as an option. This tool can also add
solvent water molecules and Na+ or Cl\sphinxhyphen{} ions to the system.
Using lammps2pdb.pl you can convert LAMMPS atom dumps into PDB files.

\sphinxAtStartPar
See the README file in the ch2lmp subdirectory for more information.

\sphinxAtStartPar
These tools were created by Pieter in’t Veld (pjintve at sandia.gov)
and Paul Crozier (pscrozi at sandia.gov) at Sandia.

\sphinxAtStartPar
CMAP support added and tested by Xiaohu Hu (hux2 at ornl.gov) and
Robert A. Latour (latourr at clemson.edu), David Hyde\sphinxhyphen{}Volpe, and
Tigran Abramyan, (Clemson University) and
Chris Lorenz (chris.lorenz at kcl.ac.uk), King’s College London.


\bigskip\hrule\bigskip



\subsection{chain tool}
\label{\detokenize{Tools:chain-tool}}\label{\detokenize{Tools:chain}}
\sphinxAtStartPar
The file chain.f90 creates a LAMMPS data file containing bead\sphinxhyphen{}spring
polymer chains and/or monomer solvent atoms.  It uses a text file
containing chain definition parameters as an input.  The created
chains and solvent atoms can strongly overlap, so LAMMPS needs to run
the system initially with a “soft” pair potential to un\sphinxhyphen{}overlap it.
The syntax for running the tool is

\begin{sphinxVerbatim}[commandchars=\\\{\}]
chain\PYG{+w}{ }\PYGZlt{}\PYG{+w}{ }def.chain\PYG{+w}{ }\PYGZgt{}\PYG{+w}{ }data.file
\end{sphinxVerbatim}

\sphinxAtStartPar
See the def.chain or def.chain.ab files in the tools directory for
examples of definition files.  This tool was used to create the system
for the {\hyperref[\detokenize{Speed_bench::doc}]{\sphinxcrossref{\DUrole{doc}{chain benchmark}}}}.


\bigskip\hrule\bigskip



\subsection{LAMMPS coding standard}
\label{\detokenize{Tools:lammps-coding-standard}}\label{\detokenize{Tools:coding-standard}}
\sphinxAtStartPar
The \sphinxcode{\sphinxupquote{coding\_standard}} folder contains multiple python scripts to
check for and apply some LAMMPS coding conventions.  The following
scripts are available:

\begin{sphinxVerbatim}[commandchars=\\\{\}]
permissions.py   \PYGZsh{} detects if sources have executable permissions and scripts have not
whitespace.py    \PYGZsh{} detects TAB characters and trailing whitespace
homepage.py      \PYGZsh{} detects outdated LAMMPS homepage URLs (pointing to sandia.gov instead of lammps.org)
errordocs.py     \PYGZsh{} detects deprecated error docs in header files
versiontags.py   \PYGZsh{} detects .. versionadded:: or .. versionchanged:: with pending version date
\end{sphinxVerbatim}

\sphinxAtStartPar
The tools need to be given the main folder of the LAMMPS distribution
or individual file names as argument and will by default check them
and report any non\sphinxhyphen{}compliance.  With the optional \sphinxcode{\sphinxupquote{\sphinxhyphen{}f}} argument the
corresponding script will try to change the non\sphinxhyphen{}compliant file(s) to
match the conventions.

\sphinxAtStartPar
For convenience this scripts can also be invoked by the make file in
the \sphinxcode{\sphinxupquote{src}} folder with, \sphinxtitleref{make check\sphinxhyphen{}whitespace} or \sphinxtitleref{make fix\sphinxhyphen{}whitespace}
to either detect or edit the files.  Correspondingly for the other python
scripts. \sphinxtitleref{make check} will run all checks.


\bigskip\hrule\bigskip



\subsection{colvars tools}
\label{\detokenize{Tools:colvars-tools}}\label{\detokenize{Tools:id1}}
\sphinxAtStartPar
The colvars directory contains a collection of tools for post\sphinxhyphen{}processing
data produced by the colvars collective variable library.
To compile the tools, edit the makefile for your system and run “make”.

\sphinxAtStartPar
Please report problems and issues the colvars library and its tools
at: \sphinxurl{https://github.com/colvars/colvars/issues}

\sphinxAtStartPar
abf\_integrate:

\sphinxAtStartPar
MC\sphinxhyphen{}based integration of multidimensional free energy gradient
Version 20110511
\begin{sphinxalltt}
./abf\_integrate \textless{} filename \textgreater{} {[}\sphinxhyphen{}n \textless{} nsteps \textgreater{}{]} {[}\sphinxhyphen{}t \textless{} temp \textgreater{}{]} {[}\sphinxhyphen{}m {[}0|1{]} (metadynamics){]} {[}\sphinxhyphen{}h \textless{} hill\_height \textgreater{}{]} {[}\sphinxhyphen{}f \textless{} variable\_hill\_factor \textgreater{}{]}
\end{sphinxalltt}

\sphinxAtStartPar
The LAMMPS interface to the colvars collective variable library, as
well as these tools, were created by Axel Kohlmeyer (akohlmey at
gmail.com) while at ICTP, Italy.


\bigskip\hrule\bigskip



\subsection{createatoms tool}
\label{\detokenize{Tools:createatoms-tool}}\label{\detokenize{Tools:createatoms}}
\sphinxAtStartPar
The tools/createatoms directory contains a Fortran program called
createAtoms.f which can generate a variety of interesting crystal
structures and geometries and output the resulting list of atom
coordinates in LAMMPS or other formats.

\sphinxAtStartPar
See the included Manual.pdf for details.

\sphinxAtStartPar
The tool is authored by Xiaowang Zhou (Sandia), xzhou at sandia.gov.


\bigskip\hrule\bigskip



\subsection{drude tool}
\label{\detokenize{Tools:drude-tool}}\label{\detokenize{Tools:drude}}
\sphinxAtStartPar
The tools/drude directory contains a Python script called
polarizer.py which can add Drude oscillators to a LAMMPS
data file in the required format.

\sphinxAtStartPar
See the header of the polarizer.py file for details.

\sphinxAtStartPar
The tool is authored by Agilio Padua and Alain Dequidt: agilio.padua
at ens\sphinxhyphen{}lyon.fr, alain.dequidt at uca.fr


\bigskip\hrule\bigskip



\subsection{eam database tool}
\label{\detokenize{Tools:eam-database-tool}}\label{\detokenize{Tools:eamdb}}
\sphinxAtStartPar
The tools/eam\_database directory contains a Fortran and a Python program
that will generate EAM alloy setfl potential files for any combination
of the 17 elements: Cu, Ag, Au, Ni, Pd, Pt, Al, Pb, Fe, Mo, Ta, W, Mg,
Co, Ti, Zr, Cr.  The files can then be used with the {\hyperref[\detokenize{pair_eam::doc}]{\sphinxcrossref{\DUrole{doc}{pair\_style
eam/alloy}}}} command.

\sphinxAtStartPar
The Fortran version of the tool was authored by Xiaowang Zhou (Sandia),
xzhou at sandia.gov, with updates from Lucas Hale (NIST) lucas.hale at
nist.gov and is based on his paper:

\sphinxAtStartPar
X. W. Zhou, R. A. Johnson, and H. N. G. Wadley, Phys. Rev. B, 69,
144113 (2004).

\sphinxAtStartPar
The parameters for Cr were taken from:

\sphinxAtStartPar
Lin Z B, Johnson R A and Zhigilei L V, Phys. Rev. B 77 214108 (2008).

\sphinxAtStartPar
The Python version of the tool was authored  by Germain Clavier
(Unicaen) germain.clavier at unicaen.fr

\begin{sphinxadmonition}{note}{Note:}
\sphinxAtStartPar
The parameters in the database are only optimized for individual
elements. The mixed parameters for interactions between different
elements generated by this tool are derived from simple mixing rules
and are thus inferior to parameterizations that are specifically
optimized for specific mixtures and combinations of elements.
\end{sphinxadmonition}


\bigskip\hrule\bigskip



\subsection{eam generate tool}
\label{\detokenize{Tools:eam-generate-tool}}\label{\detokenize{Tools:eamgn}}
\sphinxAtStartPar
The tools/eam\_generate directory contains several one\sphinxhyphen{}file C programs
that convert an analytic formula into a tabulated {\hyperref[\detokenize{pair_eam::doc}]{\sphinxcrossref{\DUrole{doc}{embedded atom
method (EAM)}}}} setfl potential file.  The potentials they
produce are in the potentials directory, and can be used with the
{\hyperref[\detokenize{pair_eam::doc}]{\sphinxcrossref{\DUrole{doc}{pair\_style eam/alloy}}}} command.

\sphinxAtStartPar
The source files and potentials were provided by Gerolf Ziegenhain
(gerolf at ziegenhain.com).


\bigskip\hrule\bigskip



\subsection{eff tool}
\label{\detokenize{Tools:eff-tool}}\label{\detokenize{Tools:eff}}
\sphinxAtStartPar
The tools/eff directory contains various scripts for generating
structures and post\sphinxhyphen{}processing output for simulations using the
electron force field (eFF).

\sphinxAtStartPar
These tools were provided by Andres Jaramillo\sphinxhyphen{}Botero at CalTech
(ajaramil at wag.caltech.edu).


\bigskip\hrule\bigskip



\subsection{emacs tool}
\label{\detokenize{Tools:emacs-tool}}\label{\detokenize{Tools:emacs}}
\sphinxAtStartPar
The tools/emacs directory contains an Emacs Lisp add\sphinxhyphen{}on file for GNU Emacs
that enables a lammps\sphinxhyphen{}mode for editing input scripts when using GNU Emacs,
with various highlighting options set up.

\sphinxAtStartPar
These tools were provided by Aidan Thompson at Sandia
(athomps at sandia.gov).


\bigskip\hrule\bigskip



\subsection{fep tool}
\label{\detokenize{Tools:fep-tool}}\label{\detokenize{Tools:fep}}
\sphinxAtStartPar
The tools/fep directory contains Python scripts useful for
post\sphinxhyphen{}processing results from performing free\sphinxhyphen{}energy perturbation
simulations using the FEP package.

\sphinxAtStartPar
The scripts were contributed by Agilio Padua (ENS de Lyon), agilio.padua at ens\sphinxhyphen{}lyon.fr.

\sphinxAtStartPar
See README file in the tools/fep directory.


\bigskip\hrule\bigskip



\subsection{i\sphinxhyphen{}PI tool}
\label{\detokenize{Tools:i-pi-tool}}\label{\detokenize{Tools:ipi}}
\sphinxAtStartPar
\DUrole{versionmodified,changed}{Changed in version 27June2024.}

\sphinxAtStartPar
The tools/i\sphinxhyphen{}pi directory used to contain a bundled version of the i\sphinxhyphen{}PI
software package for use with LAMMPS.  This version, however, was
removed in 06/2024.

\sphinxAtStartPar
The i\sphinxhyphen{}PI package was created and is maintained by Michele Ceriotti,
michele.ceriotti at gmail.com, to interface to a variety of molecular
dynamics codes.

\sphinxAtStartPar
i\sphinxhyphen{}PI is now available via PyPI using the pip package manager at:
\sphinxurl{https://pypi.org/project/ipi/}

\sphinxAtStartPar
Here are the commands to set up a virtual environment and install
i\sphinxhyphen{}PI into it with all its dependencies.

\begin{sphinxVerbatim}[commandchars=\\\{\}]
python\PYG{+w}{ }\PYGZhy{}m\PYG{+w}{ }venv\PYG{+w}{ }ipienv
\PYG{n+nb}{source}\PYG{+w}{ }ipienv/bin/activate
pip\PYG{+w}{ }install\PYG{+w}{ }\PYGZhy{}\PYGZhy{}upgrade\PYG{+w}{ }pip
pip\PYG{+w}{ }install\PYG{+w}{ }ipi
\end{sphinxVerbatim}

\sphinxAtStartPar
To install the development version from GitHub, please use:

\begin{sphinxVerbatim}[commandchars=\\\{\}]
pip\PYG{+w}{ }install\PYG{+w}{ }git+https://github.com/i\PYGZhy{}pi/i\PYGZhy{}pi.git
\end{sphinxVerbatim}

\sphinxAtStartPar
For further information, please consult the {[}i\sphinxhyphen{}PI home
page{]}(\sphinxurl{https://ipi-code.org}).


\bigskip\hrule\bigskip



\subsection{ipp tool}
\label{\detokenize{Tools:ipp-tool}}\label{\detokenize{Tools:ipp}}
\sphinxAtStartPar
The tools/ipp directory contains a Perl script ipp which can be used
to facilitate the creation of a complicated file (say, a LAMMPS input
script or tools/createatoms input file) using a template file.

\sphinxAtStartPar
ipp was created and is maintained by Reese Jones (Sandia), rjones at
sandia.gov.

\sphinxAtStartPar
See two examples in the tools/ipp directory.  One of them is for the
tools/createatoms tool’s input file.


\bigskip\hrule\bigskip



\subsection{kate tool}
\label{\detokenize{Tools:kate-tool}}\label{\detokenize{Tools:kate}}
\sphinxAtStartPar
The file in the tools/kate directory is an add\sphinxhyphen{}on to the Kate editor
in the KDE suite that allow syntax highlighting of LAMMPS input
scripts.  See the README.txt file for details.

\sphinxAtStartPar
The file was provided by Alessandro Luigi Sellerio
(alessandro.sellerio at ieni.cnr.it).


\bigskip\hrule\bigskip



\subsection{LAMMPS\sphinxhyphen{}GUI}
\label{\detokenize{Tools:lammps-gui}}\label{\detokenize{Tools:id2}}
\sphinxAtStartPar
\DUrole{versionmodified,added}{Added in version 2Aug2023.}


\subsubsection{Overview}
\label{\detokenize{Tools:overview}}
\sphinxAtStartPar
LAMMPS\sphinxhyphen{}GUI is a graphical text editor customized for editing LAMMPS
input files that is linked to the {\hyperref[\detokenize{Library:lammps-c-api}]{\sphinxcrossref{\DUrole{std,std-ref}{LAMMPS C\sphinxhyphen{}library}}}}
and thus can run LAMMPS directly using the contents of the editor’s text
buffer as input.  It can retrieve and display information from LAMMPS
while it is running, display visualizations created with the {\hyperref[\detokenize{dump_image::doc}]{\sphinxcrossref{\DUrole{doc}{dump
image command}}}}, and is adapted specifically for editing
LAMMPS input files through syntax highlighting, text completion, and
reformatting, and linking to the online LAMMPS documentation for known
LAMMPS commands and styles.

\sphinxAtStartPar
This is similar to what people traditionally would do to run LAMMPS but
all rolled into a single application: that is, using a text editor,
plotting program, and a visualization program to edit the input, run
LAMMPS, process the output into graphs and visualizations from a command
line window.  This similarity is a design goal. While making it easy for
beginners to start with LAMMPS, it is also the expectation that
LAMMPS\sphinxhyphen{}GUI users will eventually transition to workflows that most
experienced LAMMPS users employ.

\sphinxAtStartPar
All features have been extensively exposed to keyboard shortcuts, so
that there is also appeal for experienced LAMMPS users for prototyping
and testing simulation setups.


\subsubsection{Features}
\label{\detokenize{Tools:features}}
\sphinxAtStartPar
A detailed discussion and explanation of all features and functionality
are in the {\hyperref[\detokenize{Howto_lammps_gui::doc}]{\sphinxcrossref{\DUrole{doc}{Using LAMMPS\sphinxhyphen{}GUI}}}} tutorial Howto page.

\sphinxAtStartPar
Here are a few highlights of LAMMPS\sphinxhyphen{}GUI
\begin{itemize}
\item {} 
\sphinxAtStartPar
Text editor with line numbers and syntax highlighting customized for LAMMPS

\item {} 
\sphinxAtStartPar
Text editor features command completion and auto\sphinxhyphen{}indentation for known commands and styles

\item {} 
\sphinxAtStartPar
Text editor will switch its working directory to folder of file in buffer

\item {} 
\sphinxAtStartPar
Many adjustable settings and preferences that are persistent including the 5 most recent files

\item {} 
\sphinxAtStartPar
Context specific LAMMPS command help via online documentation

\item {} 
\sphinxAtStartPar
LAMMPS is running in a concurrent thread, so the GUI remains responsive

\item {} 
\sphinxAtStartPar
Progress bar indicates how far a run command is completed

\item {} 
\sphinxAtStartPar
LAMMPS can be started and stopped with a mouse click or a hotkey

\item {} 
\sphinxAtStartPar
Screen output is captured in an \sphinxstyleemphasis{Output} Window

\item {} 
\sphinxAtStartPar
Thermodynamic output is captured and displayed as line graph in a \sphinxstyleemphasis{Chart} Window

\item {} 
\sphinxAtStartPar
Indicator for currently executed command

\item {} 
\sphinxAtStartPar
Indicator for line that caused an error

\item {} 
\sphinxAtStartPar
Visualization of current state in Image Viewer (via calling {\hyperref[\detokenize{dump_image::doc}]{\sphinxcrossref{\DUrole{doc}{write\_dump image}}}})

\item {} 
\sphinxAtStartPar
Capture of images created via {\hyperref[\detokenize{dump_image::doc}]{\sphinxcrossref{\DUrole{doc}{dump image}}}} in Slide show window

\item {} 
\sphinxAtStartPar
Dialog to set variables, similar to the LAMMPS command line flag ‘\sphinxhyphen{}v’ / ‘\sphinxhyphen{}var’

\item {} 
\sphinxAtStartPar
Support for GPU, INTEL, KOKKOS/OpenMP, OPENMAP, and OPT and accelerator packages

\end{itemize}


\subsubsection{Parallelization}
\label{\detokenize{Tools:parallelization}}
\sphinxAtStartPar
Due to its nature as a graphical application, it is not possible to use
the LAMMPS\sphinxhyphen{}GUI in parallel with MPI, but OpenMP multi\sphinxhyphen{}threading and GPU
acceleration is available and enabled by default.


\subsubsection{Prerequisites and portability}
\label{\detokenize{Tools:prerequisites-and-portability}}
\sphinxAtStartPar
LAMMPS\sphinxhyphen{}GUI is programmed in C++ based on the C++11 standard and using
the \sphinxhref{https://www.qt.io/product/framework}{Qt GUI framework}.
Currently, Qt version 5.12 or later is required; Qt 5.15LTS is
recommended; support for Qt version 6.x is available.  Building LAMMPS
with CMake is required.

\sphinxAtStartPar
The LAMMPS\sphinxhyphen{}GUI has been successfully compiled and tested on:
\begin{itemize}
\item {} 
\sphinxAtStartPar
Ubuntu Linux 20.04LTS x86\_64 using GCC 9, Qt version 5.12

\item {} 
\sphinxAtStartPar
Fedora Linux 40 x86\_64 using GCC 14 and Clang 17, Qt version 5.15LTS

\item {} 
\sphinxAtStartPar
Fedora Linux 40 x86\_64 using GCC 14, Qt version 6.7

\item {} 
\sphinxAtStartPar
Apple macOS 12 (Monterey) and macOS 13 (Ventura) with Xcode on arm64 and x86\_64, Qt version 5.15LTS

\item {} 
\sphinxAtStartPar
Windows 10 and 11 x86\_64 with Visual Studio 2022 and Visual C++ 14.36, Qt version 5.15LTS

\item {} 
\sphinxAtStartPar
Windows 10 and 11 x86\_64 with Visual Studio 2022 and Visual C++ 14.40, Qt version 6.7

\item {} 
\sphinxAtStartPar
Windows 10 and 11 x86\_64 with MinGW / GCC 10.0 cross\sphinxhyphen{}compiler on Fedora 38, Qt version 5.15LTS

\end{itemize}


\subsubsection{Pre\sphinxhyphen{}compiled executables}
\label{\detokenize{Tools:pre-compiled-executables}}\label{\detokenize{Tools:lammps-gui-install}}
\sphinxAtStartPar
Pre\sphinxhyphen{}compiled LAMMPS executable packages that include the GUI are
currently available from \sphinxurl{https://download.lammps.org/static} or
\sphinxurl{https://github.com/lammps/lammps/releases}.  For Windows, you need to
download and then run the application installer.  For macOS you download
and mount the disk image and then drag the application bundle to the
Applications folder.  For Linux (x86\_64) you currently have two
options: 1) you can download the tar.gz archive, unpack it and run the
GUI directly in place.  The \sphinxcode{\sphinxupquote{LAMMPS\_GUI}} folder may also be moved
around and added to the \sphinxcode{\sphinxupquote{PATH}} environment variable so the executables
will be found automatically.  2) you can download the \sphinxhref{https://www.flatpak.org/}{Flatpak file} and then install it locally with the
\sphinxstyleemphasis{flatpak} command: \sphinxcode{\sphinxupquote{flatpak install \sphinxhyphen{}\sphinxhyphen{}user
LAMMPS\sphinxhyphen{}Linux\sphinxhyphen{}x86\_64\sphinxhyphen{}GUI\sphinxhyphen{}\textless{}version\textgreater{}.flatpak}} and run it with \sphinxcode{\sphinxupquote{flatpak
run org.lammps.lammps\sphinxhyphen{}gui}}.  The flatpak bundle also includes the
command line version of LAMMPS and some LAMMPS tools like msi2lmp.  The
can be launched by using the \sphinxcode{\sphinxupquote{\sphinxhyphen{}\sphinxhyphen{}command}} flag. For example to run
LAMMPS directly on the \sphinxcode{\sphinxupquote{in.lj}} benchmark input you would type in the
\sphinxcode{\sphinxupquote{bench}} folder: \sphinxcode{\sphinxupquote{flatpak run \sphinxhyphen{}\sphinxhyphen{}command=lmp \sphinxhyphen{}in in.lj}} The flatpak
version should also appear in the applications menu of standard desktop
environments.  The LAMMPS\sphinxhyphen{}GUI executable is called \sphinxcode{\sphinxupquote{lammps\sphinxhyphen{}gui}} and
either takes no arguments or attempts to load the first argument as
LAMMPS input file.


\subsubsection{Compilation}
\label{\detokenize{Tools:compilation}}\label{\detokenize{Tools:lammps-gui-compilation}}
\sphinxAtStartPar
The source for the LAMMPS\sphinxhyphen{}GUI is included with the LAMMPS source code
distribution in the folder \sphinxcode{\sphinxupquote{tools/lammps\sphinxhyphen{}gui}} and thus it can be can
be built as part of a regular LAMMPS compilation.  {\hyperref[\detokenize{Howto_cmake::doc}]{\sphinxcrossref{\DUrole{doc}{Using CMake}}}} is required.  To enable its compilation, the CMake
variable \sphinxcode{\sphinxupquote{\sphinxhyphen{}D BUILD\_LAMMPS\_GUI=on}} must be set when creating the CMake
configuration.  All other settings (compiler, flags, compile type) for
LAMMPS\sphinxhyphen{}GUI are then inherited from the regular LAMMPS build.  If the Qt
library is packaged for Linux distributions, then its location is
typically auto\sphinxhyphen{}detected since the required CMake configuration files are
stored in a location where CMake can find them without additional help.
Otherwise, the location of the Qt library installation must be indicated
by setting \sphinxcode{\sphinxupquote{\sphinxhyphen{}D Qt5\_DIR=/path/to/qt5/lib/cmake/Qt5}}, which is a path to
a folder inside the Qt installation that contains the file
\sphinxcode{\sphinxupquote{Qt5Config.cmake}}. Similarly, for Qt6 the location of the Qt library
installation can be indicated by setting \sphinxcode{\sphinxupquote{\sphinxhyphen{}D
Qt6\_DIR=/path/to/qt6/lib/cmake/Qt6}}, if necessary.  When both, Qt5 and
Qt6 are available, Qt6 will be preferred unless \sphinxcode{\sphinxupquote{\sphinxhyphen{}D
LAMMPS\_GUI\_USE\_QT5=yes}} is set.

\sphinxAtStartPar
It is possible to build the LAMMPS\sphinxhyphen{}GUI as a standalone compilation
(e.g. when LAMMPS has been compiled with traditional make).  Then the
CMake configuration needs to be told where to find the LAMMPS headers
and the LAMMPS library, via \sphinxcode{\sphinxupquote{\sphinxhyphen{}D LAMMPS\_SOURCE\_DIR=/path/to/lammps/src}}.
CMake will try to guess a build folder with the LAMMPS library from that
path, but it can also be set with \sphinxcode{\sphinxupquote{\sphinxhyphen{}D LAMMPS\_LIB\_DIR=/path/to/lammps/lib}}.


\paragraph{Plugin version}
\label{\detokenize{Tools:plugin-version}}
\sphinxAtStartPar
Rather than linking to the LAMMPS library during compilation, it is also
possible to compile the GUI with a plugin loader that will load the
LAMMPS library dynamically at runtime during the start of the GUI from a
shared library; e.g. \sphinxcode{\sphinxupquote{liblammps.so}} or \sphinxcode{\sphinxupquote{liblammps.dylib}} or
\sphinxcode{\sphinxupquote{liblammps.dll}} (depending on the operating system).  This has the
advantage that the LAMMPS library can be built from updated or modified
LAMMPS source without having to recompile the GUI.  The ABI of the
LAMMPS C\sphinxhyphen{}library interface is very stable and generally backward
compatible.  This feature is enabled by setting \sphinxcode{\sphinxupquote{\sphinxhyphen{}D
LAMMPS\_GUI\_USE\_PLUGIN=on}} and then \sphinxcode{\sphinxupquote{\sphinxhyphen{}D
LAMMPS\_PLUGINLIB\_DIR=/path/to/lammps/plugin/loader}}. Typically, this
would be the \sphinxcode{\sphinxupquote{examples/COUPLE/plugin}} folder of the LAMMPS
distribution.

\sphinxAtStartPar
When compiling LAMMPS\sphinxhyphen{}GUI with plugin support, there is an additional
command line flag (\sphinxcode{\sphinxupquote{\sphinxhyphen{}p \textless{}path\textgreater{}}} or \sphinxcode{\sphinxupquote{\sphinxhyphen{}\sphinxhyphen{}pluginpath \textless{}path\textgreater{}}}) which
allows to override the path to LAMMPS shared library used by the GUI.
This is usually auto\sphinxhyphen{}detected on the first run and can be changed in the
LAMMPS\sphinxhyphen{}GUI \sphinxstyleemphasis{Preferences} dialog.  The command line flag allows to reset
this path to a valid value in case the original setting has become
invalid.  An empty path (“”) as argument restores the default setting.


\subsubsection{Platform notes}
\label{\detokenize{Tools:platform-notes}}

\paragraph{macOS}
\label{\detokenize{Tools:macos}}
\sphinxAtStartPar
When building on macOS, the build procedure will try to manufacture a
drag\sphinxhyphen{}n\sphinxhyphen{}drop installer, \sphinxcode{\sphinxupquote{LAMMPS\sphinxhyphen{}macOS\sphinxhyphen{}multiarch.dmg}}, when using the
‘dmg’ target (i.e. \sphinxcode{\sphinxupquote{cmake \sphinxhyphen{}\sphinxhyphen{}build \textless{}build dir\textgreater{} \sphinxhyphen{}\sphinxhyphen{}target dmg}} or \sphinxcode{\sphinxupquote{make dmg}}.

\sphinxAtStartPar
To build multi\sphinxhyphen{}arch executables that will run on both, arm64 and x86\_64
architectures natively, it is necessary to set the CMake variable \sphinxcode{\sphinxupquote{\sphinxhyphen{}D
CMAKE\_OSX\_ARCHITECTURES=arm64;x86\_64}}.  To achieve wide compatibility
with different macOS versions, you can also set \sphinxcode{\sphinxupquote{\sphinxhyphen{}D
CMAKE\_OSX\_DEPLOYMENT\_TARGET=11.0}} which will set compatibility to macOS
11 (Big Sur) and later, even if you are compiling on a more recent macOS
version.


\paragraph{Windows}
\label{\detokenize{Tools:windows}}
\sphinxAtStartPar
On Windows either native compilation from within Visual Studio 2022 with
Visual C++ is supported and tested, or compilation with the MinGW / GCC
cross\sphinxhyphen{}compiler environment on Fedora Linux.

\sphinxAtStartPar
\sphinxstylestrong{Visual Studio}

\sphinxAtStartPar
Using CMake and Ninja as build system are required.  Qt needs to be
installed, tested was a binary package downloaded from
\sphinxurl{https://www.qt.io}, which installs into the \sphinxcode{\sphinxupquote{C:\textbackslash{}\textbackslash{}Qt}} folder by default.
There is a custom \sphinxtitleref{x64\sphinxhyphen{}GUI\sphinxhyphen{}MSVC} build configuration provided in the
\sphinxcode{\sphinxupquote{CMakeSettings.json}} file that Visual Studio uses to store different
compilation settings for project.  Choosing this configuration will
activate building the \sphinxtitleref{lammps\sphinxhyphen{}gui.exe} executable in addition to LAMMPS
through importing package selection from the \sphinxcode{\sphinxupquote{windows.cmake}} preset
file and enabling building the LAMMPS\sphinxhyphen{}GUI and disabling building with MPI.
When requesting an installation from the \sphinxtitleref{Build} menu in Visual Studio,
it will create a compressed \sphinxcode{\sphinxupquote{LAMMPS\sphinxhyphen{}Win10\sphinxhyphen{}amd64.zip}} zip file with the
executables and required dependent .dll files.  This zip file can be
uncompressed and \sphinxcode{\sphinxupquote{lammps\sphinxhyphen{}gui.exe}} run directly from there.  The
uncompressed folder can be added to the \sphinxcode{\sphinxupquote{PATH}} environment and LAMMPS
and LAMMPS\sphinxhyphen{}GUI can be launched from anywhere from the command line.

\sphinxAtStartPar
\sphinxstylestrong{MinGW64 Cross\sphinxhyphen{}compiler}

\sphinxAtStartPar
The standard CMake build procedure can be applied and the
\sphinxcode{\sphinxupquote{mingw\sphinxhyphen{}cross.cmake}} preset used. By using \sphinxcode{\sphinxupquote{mingw64\sphinxhyphen{}cmake}} the CMake
command will automatically include a suitable CMake toolchain file (the
regular cmake command can be used after that to modify the configuration
settings, if needed).  After building the libraries and executables,
you can build the target ‘zip’ (i.e. \sphinxcode{\sphinxupquote{cmake \sphinxhyphen{}\sphinxhyphen{}build \textless{}build dir\textgreater{} \sphinxhyphen{}\sphinxhyphen{}target zip}}
or \sphinxcode{\sphinxupquote{make zip}} to stage all installed files into a LAMMPS\_GUI folder
and then run a script to copy all required dependencies, some other files,
and create a zip file from it.


\paragraph{Linux}
\label{\detokenize{Tools:linux}}
\sphinxAtStartPar
Version 5.12 or later of the Qt library is required. Those are provided
by, e.g., Ubuntu 20.04LTS.  Thus older Linux distributions are not
likely to be supported, while more recent ones will work, even for
pre\sphinxhyphen{}compiled executables (see above).  After compiling with
\sphinxcode{\sphinxupquote{cmake \sphinxhyphen{}\sphinxhyphen{}build \textless{}build folder\textgreater{}}}, use \sphinxcode{\sphinxupquote{cmake \sphinxhyphen{}\sphinxhyphen{}build \textless{}build
folder\textgreater{} \sphinxhyphen{}\sphinxhyphen{}target tgz}} or \sphinxcode{\sphinxupquote{make tgz}} to build a
\sphinxcode{\sphinxupquote{LAMMPS\sphinxhyphen{}Linux\sphinxhyphen{}amd64.tar.gz}} file with the executables and their
support libraries.

\sphinxAtStartPar
It is also possible to build a \sphinxhref{https://docs.flatpak.org/en/latest/single-file-bundles.html}{flatpak bundle} which is
a way to distribute applications in a way that is compatible with most
Linux distributions.  Use the “flatpak” target to trigger a compile
(\sphinxcode{\sphinxupquote{cmake \sphinxhyphen{}\sphinxhyphen{}build \textless{}build folder\textgreater{} \sphinxhyphen{}\sphinxhyphen{}target flatpak}} or \sphinxcode{\sphinxupquote{make flatpak}}).
Please note that this will not build from the local sources but from the
repository and branch listed in the \sphinxcode{\sphinxupquote{org.lammps.lammps\sphinxhyphen{}gui.yml}}
LAMMPS\sphinxhyphen{}GUI source folder.


\bigskip\hrule\bigskip



\subsection{lmp2arc tool}
\label{\detokenize{Tools:lmp2arc-tool}}\label{\detokenize{Tools:arc}}
\sphinxAtStartPar
The lmp2arc subdirectory contains a tool for converting LAMMPS output
files to the format for Accelrys’ Insight MD code (formerly
MSI/Biosym and its Discover MD code).  See the README file for more
information.

\sphinxAtStartPar
This tool was written by John Carpenter (Cray), Michael Peachey
(Cray), and Steve Lustig (Dupont).  John is now at the Mayo Clinic
(jec at mayo.edu), but still fields questions about the tool.

\sphinxAtStartPar
This tool was updated for the current LAMMPS C++ version by Jeff
Greathouse at Sandia (jagreat at sandia.gov).


\bigskip\hrule\bigskip



\subsection{lmp2cfg tool}
\label{\detokenize{Tools:lmp2cfg-tool}}\label{\detokenize{Tools:cfg}}
\sphinxAtStartPar
The lmp2cfg subdirectory contains a tool for converting LAMMPS output
files into a series of *.cfg files which can be read into the
\sphinxhref{http://li.mit.edu/Archive/Graphics/A/}{AtomEye} visualizer.  See
the README file for more information.

\sphinxAtStartPar
This tool was written by Ara Kooser at Sandia (askoose at sandia.gov).


\bigskip\hrule\bigskip



\subsection{Magic patterns for the “file” command}
\label{\detokenize{Tools:magic-patterns-for-the-file-command}}\label{\detokenize{Tools:magic}}
\sphinxAtStartPar
\DUrole{versionmodified,added}{Added in version 10Mar2021.}

\sphinxAtStartPar
The file \sphinxcode{\sphinxupquote{magic}} contains patterns that are used by the
\sphinxhref{https://en.wikipedia.org/wiki/File\_(command)}{file program}
available on most Unix\sphinxhyphen{}like operating systems which enables it
to detect various LAMMPS files and print some useful information
about them.  To enable these patterns, append or copy the contents
of the file to either the file \sphinxcode{\sphinxupquote{.magic}} in your home directory
or (as administrator) to \sphinxcode{\sphinxupquote{/etc/magic}} (for a system\sphinxhyphen{}wide
installation).  Afterwards the \sphinxcode{\sphinxupquote{file}} command should be able to
detect most LAMMPS restarts, dump, data and log files. Examples:

\begin{sphinxVerbatim}[commandchars=\\\{\}]
\PYG{g+gp}{\PYGZdl{} }file\PYG{+w}{ }*.*
\PYG{g+go}{dihedral\PYGZhy{}quadratic.restart:   LAMMPS binary restart file (rev 2), Version 10 Mar 2021, Little Endian}
\PYG{g+go}{mol\PYGZhy{}pair\PYGZhy{}wf\PYGZus{}cut.restart:      LAMMPS binary restart file (rev 2), Version 24 Dec 2020, Little Endian}
\PYG{g+go}{atom.bin:                     LAMMPS atom style binary dump (rev 2), Little Endian, First time step: 445570}
\PYG{g+go}{custom.bin:                   LAMMPS custom style binary dump (rev 2), Little Endian, First time step: 100}
\PYG{g+go}{bn1.lammpstrj:                LAMMPS text mode dump, First time step: 5000}
\PYG{g+go}{data.fourmol:                 LAMMPS data file written by LAMMPS}
\PYG{g+go}{pnc.data:                     LAMMPS data file written by msi2lmp}
\PYG{g+go}{data.spce:                    LAMMPS data file written by TopoTools}
\PYG{g+go}{B.data:                       LAMMPS data file written by OVITO}
\PYG{g+go}{log.lammps:                   LAMMPS log file written by version 10 Feb 2021}
\end{sphinxVerbatim}


\bigskip\hrule\bigskip



\subsection{matlab tool}
\label{\detokenize{Tools:matlab-tool}}\label{\detokenize{Tools:matlab}}
\sphinxAtStartPar
The matlab subdirectory contains several \sphinxhref{https://www.mathworks.com}{MATLAB} scripts for
post\sphinxhyphen{}processing LAMMPS output.  The scripts include readers for log
and dump files, a reader for EAM potential files, and a converter that
reads LAMMPS dump files and produces CFG files that can be visualized
with the \sphinxhref{http://li.mit.edu/Archive/Graphics/A/}{AtomEye}
visualizer.

\sphinxAtStartPar
See the README.pdf file for more information.

\sphinxAtStartPar
These scripts were written by Arun Subramaniyan at Purdue Univ
(asubrama at purdue.edu).


\bigskip\hrule\bigskip



\subsection{micelle2d tool}
\label{\detokenize{Tools:micelle2d-tool}}\label{\detokenize{Tools:micelle}}
\sphinxAtStartPar
The file micelle2d.f creates a LAMMPS data file containing short lipid
chains in a monomer solution.  It uses a text file containing lipid
definition parameters as an input.  The created molecules and solvent
atoms can strongly overlap, so LAMMPS needs to run the system
initially with a “soft” pair potential to un\sphinxhyphen{}overlap it.  The syntax
for running the tool is

\begin{sphinxVerbatim}[commandchars=\\\{\}]
micelle2d\PYG{+w}{ }\PYGZlt{}\PYG{+w}{ }def.micelle2d\PYG{+w}{ }\PYGZgt{}\PYG{+w}{ }data.file
\end{sphinxVerbatim}

\sphinxAtStartPar
See the def.micelle2d file in the tools directory for an example of a
definition file.  This tool was used to create the system for the
{\hyperref[\detokenize{Examples::doc}]{\sphinxcrossref{\DUrole{doc}{micelle example}}}}.


\bigskip\hrule\bigskip



\subsection{moltemplate tool}
\label{\detokenize{Tools:moltemplate-tool}}\label{\detokenize{Tools:moltemplate}}
\sphinxAtStartPar
The moltemplate subdirectory contains instructions for installing
moltemplate, a Python\sphinxhyphen{}based tool for building molecular systems based
on a text\sphinxhyphen{}file description, and creating LAMMPS data files that encode
their molecular topology as lists of bonds, angles, dihedrals, etc.
See the README.txt file for more information.

\sphinxAtStartPar
This tool was written by Andrew Jewett (jewett.aij at gmail.com), who
supports it.  It has its own WWW page at
\sphinxurl{https://moltemplate.org}.
The latest sources can be found \sphinxhref{https://github.com/jewettaij/moltemplate/releases}{on its GitHub page}


\bigskip\hrule\bigskip



\subsection{msi2lmp tool}
\label{\detokenize{Tools:msi2lmp-tool}}\label{\detokenize{Tools:msi}}
\sphinxAtStartPar
The msi2lmp subdirectory contains a tool for creating LAMMPS template
input and data files from BIOVIA’s Materias Studio files (formerly
Accelrys’ Insight MD code, formerly MSI/Biosym and its Discover MD code).

\sphinxAtStartPar
This tool was written by John Carpenter (Cray), Michael Peachey
(Cray), and Steve Lustig (Dupont). Several people contributed changes
to remove bugs and adapt its output to changes in LAMMPS.

\sphinxAtStartPar
This tool has several known limitations and is no longer under active
development, so there are no changes except for the occasional bug fix.

\sphinxAtStartPar
See the README file in the tools/msi2lmp folder for more information.


\bigskip\hrule\bigskip



\subsection{Scripts for building LAMMPS when offline}
\label{\detokenize{Tools:scripts-for-building-lammps-when-offline}}\label{\detokenize{Tools:offline}}
\sphinxAtStartPar
In some situations it might be necessary to build LAMMPS on a system
without direct internet access. The scripts in \sphinxcode{\sphinxupquote{tools/offline}} folder
allow you to pre\sphinxhyphen{}load external dependencies for both the documentation
build and for building LAMMPS with CMake.

\sphinxAtStartPar
It does so by
\begin{enumerate}
\sphinxsetlistlabels{\arabic}{enumi}{enumii}{}{.}%
\item {} 
\sphinxAtStartPar
downloading necessary \sphinxcode{\sphinxupquote{pip}} packages,

\item {} 
\sphinxAtStartPar
cloning \sphinxcode{\sphinxupquote{git}} repositories

\item {} 
\sphinxAtStartPar
downloading tarballs

\end{enumerate}

\sphinxAtStartPar
to a designated cache folder.

\sphinxAtStartPar
As of April 2021, all of these downloads make up around 600MB. By
default, the offline scripts will download everything into the
\sphinxcode{\sphinxupquote{\$HOME/.cache/lammps}} folder, but this can be changed by setting the
\sphinxcode{\sphinxupquote{LAMMPS\_CACHING\_DIR}} environment variable.

\sphinxAtStartPar
Once the caches have been initialized, they can be used for building the
LAMMPS documentation or compiling LAMMPS using CMake on an offline
system.

\sphinxAtStartPar
The \sphinxcode{\sphinxupquote{use\_caches.sh}} script must be sourced into the current shell
to initialize the offline build environment. Note that it must use
the same \sphinxcode{\sphinxupquote{LAMMPS\_CACHING\_DIR}}. This script does the following:
\begin{enumerate}
\sphinxsetlistlabels{\arabic}{enumi}{enumii}{}{.}%
\item {} 
\sphinxAtStartPar
Set up environment variables that modify the behavior of both,
\sphinxcode{\sphinxupquote{pip}} and \sphinxcode{\sphinxupquote{git}}

\item {} 
\sphinxAtStartPar
Start a simple local HTTP server using Python to host files for CMake

\end{enumerate}

\sphinxAtStartPar
Afterwards, it will print out instruction on how to modify the CMake
command line to make sure it uses the local HTTP server.

\sphinxAtStartPar
To undo the environment changes and shutdown the local HTTP server,
run the \sphinxcode{\sphinxupquote{deactivate\_caches}} command.


\subsubsection{Examples}
\label{\detokenize{Tools:examples}}
\sphinxAtStartPar
For all of the examples below, you first need to create the cache, which
requires an internet connection.

\begin{sphinxVerbatim}[commandchars=\\\{\}]
./tools/offline/init\PYGZus{}caches.sh
\end{sphinxVerbatim}

\sphinxAtStartPar
Afterwards, you can disconnect or copy the contents of the
\sphinxcode{\sphinxupquote{LAMMPS\_CACHING\_DIR}} folder to an offline system.


\subsubsection{Documentation Build}
\label{\detokenize{Tools:documentation-build}}
\sphinxAtStartPar
The documentation build will create a new virtual environment that
typically first installs dependencies from \sphinxcode{\sphinxupquote{pip}}. With the offline
environment loaded, these installations will instead grab the necessary
packages from your local cache.

\begin{sphinxVerbatim}[commandchars=\\\{\}]
\PYG{c+c1}{\PYGZsh{} if LAMMPS\PYGZus{}CACHING\PYGZus{}DIR is different from default, make sure to set it first}
\PYG{c+c1}{\PYGZsh{} export LAMMPS\PYGZus{}CACHING\PYGZus{}DIR=path/to/folder}
\PYG{n+nb}{source}\PYG{+w}{ }tools/offline/use\PYGZus{}caches.sh
\PYG{n+nb}{cd}\PYG{+w}{ }doc/
make\PYG{+w}{ }html

deactivate\PYGZus{}caches
\end{sphinxVerbatim}


\subsubsection{CMake Build}
\label{\detokenize{Tools:cmake-build}}
\sphinxAtStartPar
When compiling certain packages with external dependencies, the CMake
build system will download necessary files or sources from the web. For
more flexibility the CMake configuration allows users to specify the URL
of each of these dependencies.  What the \sphinxcode{\sphinxupquote{init\_caches.sh}} script does
is create a CMake “preset” file, which sets the URLs for all of the known
dependencies and redirects the download to the local cache.

\begin{sphinxVerbatim}[commandchars=\\\{\}]
\PYG{c+c1}{\PYGZsh{} if LAMMPS\PYGZus{}CACHING\PYGZus{}DIR is different from default, make sure to set it first}
\PYG{c+c1}{\PYGZsh{} export LAMMPS\PYGZus{}CACHING\PYGZus{}DIR=path/to/folder}
\PYG{n+nb}{source}\PYG{+w}{ }tools/offline/use\PYGZus{}caches.sh

mkdir\PYG{+w}{ }build
\PYG{n+nb}{cd}\PYG{+w}{ }build
cmake\PYG{+w}{ }\PYGZhy{}D\PYG{+w}{ }\PYG{n+nv}{LAMMPS\PYGZus{}DOWNLOADS\PYGZus{}URL}\PYG{o}{=}\PYG{l+s+si}{\PYGZdl{}\PYGZob{}}\PYG{n+nv}{HTTP\PYGZus{}CACHE\PYGZus{}URL}\PYG{l+s+si}{\PYGZcb{}}\PYG{+w}{ }\PYGZhy{}C\PYG{+w}{ }\PYG{l+s+s2}{\PYGZdq{}}\PYG{l+s+si}{\PYGZdl{}\PYGZob{}}\PYG{n+nv}{LAMMPS\PYGZus{}HTTP\PYGZus{}CACHE\PYGZus{}CONFIG}\PYG{l+s+si}{\PYGZcb{}}\PYG{l+s+s2}{\PYGZdq{}}\PYG{+w}{ }\PYGZhy{}C\PYG{+w}{ }../cmake/presets/most.cmake\PYG{+w}{ }../cmake
make\PYG{+w}{ }\PYGZhy{}j\PYG{+w}{ }\PYG{l+m}{8}

deactivate\PYGZus{}caches
\end{sphinxVerbatim}


\bigskip\hrule\bigskip



\subsection{phonon tool}
\label{\detokenize{Tools:phonon-tool}}\label{\detokenize{Tools:phonon}}
\sphinxAtStartPar
The phonon subdirectory contains a post\sphinxhyphen{}processing tool, \sphinxstyleemphasis{phana}, useful
for analyzing the output of the {\hyperref[\detokenize{fix_phonon::doc}]{\sphinxcrossref{\DUrole{doc}{fix phonon}}}} command
in the PHONON package.

\sphinxAtStartPar
See the README file for instruction on building the tool and what
library it needs.  And see the examples/PACKAGES/phonon directory
for example problems that can be post\sphinxhyphen{}processed with this tool.

\sphinxAtStartPar
This tool was written by Ling\sphinxhyphen{}Ti Kong at Shanghai Jiao Tong
University.


\bigskip\hrule\bigskip



\subsection{polybond tool}
\label{\detokenize{Tools:polybond-tool}}\label{\detokenize{Tools:polybond}}
\sphinxAtStartPar
The polybond subdirectory contains a Python\sphinxhyphen{}based tool useful for
performing “programmable polymer bonding”.  The Python file
lmpsdata.py provides a “Lmpsdata” class with various methods which can
be invoked by a user\sphinxhyphen{}written Python script to create data files with
complex bonding topologies.

\sphinxAtStartPar
See the Manual.pdf for details and example scripts.

\sphinxAtStartPar
This tool was written by Zachary Kraus at Georgia Tech.


\bigskip\hrule\bigskip



\subsection{pymol\_asphere tool}
\label{\detokenize{Tools:pymol-asphere-tool}}\label{\detokenize{Tools:pymol}}
\sphinxAtStartPar
The pymol\_asphere subdirectory contains a tool for converting a
LAMMPS dump file that contains orientation info for ellipsoidal
particles into an input file for the \sphinxhref{https://www.pymol.org}{PyMol visualization package} or its \sphinxhref{https://github.com/schrodinger/pymol-open-source}{open source variant}.

\sphinxAtStartPar
Specifically, the tool triangulates the ellipsoids so they can be
viewed as true ellipsoidal particles within PyMol.  See the README and
examples directory within pymol\_asphere for more information.

\sphinxAtStartPar
This tool was written by Mike Brown at Sandia.


\bigskip\hrule\bigskip



\subsection{python tool}
\label{\detokenize{Tools:python-tool}}\label{\detokenize{Tools:pythontools}}
\sphinxAtStartPar
The python subdirectory contains several Python scripts
that perform common LAMMPS post\sphinxhyphen{}processing tasks, such as:
\begin{itemize}
\item {} 
\sphinxAtStartPar
extract thermodynamic info from a log file as columns of numbers

\item {} 
\sphinxAtStartPar
plot two columns of thermodynamic info from a log file using GnuPlot

\item {} 
\sphinxAtStartPar
sort the snapshots in a dump file by atom ID

\item {} 
\sphinxAtStartPar
convert multiple {\hyperref[\detokenize{neb::doc}]{\sphinxcrossref{\DUrole{doc}{NEB}}}} dump files into one dump file for viz

\item {} 
\sphinxAtStartPar
convert dump files into XYZ, CFG, or PDB format for viz by other packages

\end{itemize}

\sphinxAtStartPar
These are simple scripts built on \sphinxhref{https://lammps.github.io/pizza}{Pizza.py} modules.  See the
README for more info on Pizza.py and how to use these scripts.


\bigskip\hrule\bigskip



\subsection{Regression tester tool}
\label{\detokenize{Tools:regression-tester-tool}}\label{\detokenize{Tools:regression}}
\sphinxAtStartPar
The regression\sphinxhyphen{}tests subdirectory contains a tool for performing
regression tests with a given LAMMPS binary.  The tool launches the
LAMMPS binary with any given input script under one of the \sphinxtitleref{examples}
subdirectories, and compares the thermo output in the generated log file
with those in the provided log file with the same number of processors
ub the same subdirectory. If the differences between the actual and
reference values are within specified tolerances, the test is considered
passed.  For each test batch, that is, a set of example input scripts,
the mpirun command, the LAMMPS command line arguments, and the
tolerances for individual thermo quantities can be specified in a
configuration file in YAML format.

\sphinxAtStartPar
The tool also reports if and how the run fails, and if a reference log file
is missing.  See the README file for more information.

\sphinxAtStartPar
This tool was written by Trung Nguyen at U of Chicago (ndactrung at gmail.com).


\bigskip\hrule\bigskip



\subsection{replica tool}
\label{\detokenize{Tools:replica-tool}}\label{\detokenize{Tools:replica}}
\sphinxAtStartPar
The tools/replica directory contains the reorder\_remd\_traj python script which
can be used to reorder the replica trajectories (resulting from the use of the
temper command) according to temperature. This will produce discontinuous
trajectories with all frames at the same temperature in each trajectory.
Additional options can be used to calculate the canonical configurational
log\sphinxhyphen{}weight for each frame at each temperature using the pymbar package. See
the README.md file for further details. Try out the peptide example provided.

\sphinxAtStartPar
This tool was written by (and is maintained by) Tanmoy Sanyal,
while at the Shell lab at UC Santa Barbara. (tanmoy dot 7989 at gmail.com)


\bigskip\hrule\bigskip



\subsection{smd tool}
\label{\detokenize{Tools:smd-tool}}\label{\detokenize{Tools:smd}}
\sphinxAtStartPar
The smd subdirectory contains a C++ file dump2vtk\_tris.cpp and
Makefile which can be compiled and used to convert triangle output
files created by the Smooth\sphinxhyphen{}Mach Dynamics (MACHDYN) package into a
VTK\sphinxhyphen{}compatible unstructured grid file.  It could then be read in and
visualized by VTK.

\sphinxAtStartPar
See the header of dump2vtk.cpp for more details.

\sphinxAtStartPar
This tool was written by the MACHDYN package author, Georg
Ganzenmuller at the Fraunhofer\sphinxhyphen{}Institute for High\sphinxhyphen{}Speed Dynamics,
Ernst Mach Institute in Germany (georg.ganzenmueller at emi.fhg.de).


\bigskip\hrule\bigskip



\subsection{spin tool}
\label{\detokenize{Tools:spin-tool}}\label{\detokenize{Tools:spin}}
\sphinxAtStartPar
The spin subdirectory contains a C file interpolate.c which can
be compiled and used to perform a cubic polynomial interpolation of
the MEP following a GNEB calculation.

\sphinxAtStartPar
See the README file in tools/spin/interpolate\_gneb for more details.

\sphinxAtStartPar
This tool was written by the SPIN package author, Julien
Tranchida at Sandia National Labs (jtranch at sandia.gov, and by Aleksei
Ivanov, at University of Iceland (ali5 at hi.is).


\bigskip\hrule\bigskip



\subsection{singularity/apptainer tool}
\label{\detokenize{Tools:singularity-apptainer-tool}}\label{\detokenize{Tools:singularity-tool}}
\sphinxAtStartPar
The singularity subdirectory contains container definitions files that
can be used to build container images for building and testing LAMMPS on
specific OS variants using the \sphinxhref{https://apptainer.org}{Apptainer} or
\sphinxhref{https://sylabs.io}{Singularity} container software. Contributions for
additional variants are welcome.  For more details please see the
README.md file in that folder.


\bigskip\hrule\bigskip



\subsection{stl\_bin2txt tool}
\label{\detokenize{Tools:stl-bin2txt-tool}}\label{\detokenize{Tools:stlconvert}}
\sphinxAtStartPar
The file stl\_bin2txt.cpp converts binary STL files \sphinxhyphen{} like they are
frequently offered for download on the web \sphinxhyphen{} into ASCII format STL files
that LAMMPS can read with the {\hyperref[\detokenize{create_atoms::doc}]{\sphinxcrossref{\DUrole{doc}{create\_atoms mesh}}}} or
the {\hyperref[\detokenize{fix_smd_wall_surface::doc}]{\sphinxcrossref{\DUrole{doc}{fix smd/wall\_surface}}}} commands.  The syntax
for running the tool is

\begin{sphinxVerbatim}[commandchars=\\\{\}]
stl\PYGZus{}bin2txt\PYG{+w}{ }infile.stl\PYG{+w}{ }outfile.stl
\end{sphinxVerbatim}

\sphinxAtStartPar
which creates outfile.stl from infile.stl.  This tool must be compiled
on a platform compatible with the byte\sphinxhyphen{}ordering that was used to create
the binary file.  This usually is a so\sphinxhyphen{}called little endian hardware
(like x86).


\bigskip\hrule\bigskip



\subsection{SWIG interface}
\label{\detokenize{Tools:swig-interface}}\label{\detokenize{Tools:swig}}
\sphinxAtStartPar
The \sphinxhref{https://swig.org}{SWIG tool} offers a mostly automated way to
incorporate compiled code modules into scripting languages.  It
processes the function prototypes in C and generates wrappers for a wide
variety of scripting languages from it.  Thus it can also be applied to
the {\hyperref[\detokenize{Library::doc}]{\sphinxcrossref{\DUrole{doc}{C language library interface}}}} of LAMMPS so that
build a wrapper that allows to call LAMMPS from programming languages
like: C\#/Mono, Lua, Java, JavaScript, Perl, Python, R, Ruby, Tcl, and
more.


\subsubsection{What is included}
\label{\detokenize{Tools:what-is-included}}
\sphinxAtStartPar
We provide here an “interface file”, \sphinxcode{\sphinxupquote{lammps.i}}, that has the content
of the \sphinxcode{\sphinxupquote{library.h}} file adapted so SWIG can process it.  That will
create wrappers for all the functions that are present in the LAMMPS C
library interface.  Please note that not all kinds of C functions can be
automatically translated, so you would have to add custom functions to
be able to utilize those where the automatic translation does not work.
A few functions for converting pointers and accessing arrays are
predefined.  We provide the file here on an “as is” basis to help people
getting started, but not as a fully tested and supported feature of the
LAMMPS distribution.  Any contributions to complete this are, of course,
welcome.  Please also note, that for the case of creating a Python wrapper,
a fully supported {\hyperref[\detokenize{Python_module::doc}]{\sphinxcrossref{\DUrole{doc}{Ctypes based lammps module}}}}
already exists.  That module is designed to be object\sphinxhyphen{}oriented while
SWIG will generate a 1:1 translation of the functions in the interface file.


\subsubsection{Building the wrapper}
\label{\detokenize{Tools:building-the-wrapper}}
\sphinxAtStartPar
When using CMake, the build steps for building a wrapper
module are integrated for the languages: Java, Lua,
Perl5, Python, Ruby, and Tcl.  These require that the
LAMMPS library is build as a shared library and all
necessary development headers and libraries are present.

\begin{sphinxVerbatim}[commandchars=\\\{\}]
\PYGZhy{}D\PYG{+w}{ }\PYG{n+nv}{WITH\PYGZus{}SWIG}\PYG{o}{=}on\PYG{+w}{          }\PYG{c+c1}{\PYGZsh{} to enable building any SWIG wrapper}
\PYGZhy{}D\PYG{+w}{ }\PYG{n+nv}{BUILD\PYGZus{}SWIG\PYGZus{}JAVA}\PYG{o}{=}on\PYG{+w}{    }\PYG{c+c1}{\PYGZsh{} to enable building the Java wrapper}
\PYGZhy{}D\PYG{+w}{ }\PYG{n+nv}{BUILD\PYGZus{}SWIG\PYGZus{}LUA}\PYG{o}{=}on\PYG{+w}{     }\PYG{c+c1}{\PYGZsh{} to enable building the Lua wrapper}
\PYGZhy{}D\PYG{+w}{ }\PYG{n+nv}{BUILD\PYGZus{}SWIG\PYGZus{}PERL5}\PYG{o}{=}on\PYG{+w}{   }\PYG{c+c1}{\PYGZsh{} to enable building the Perl 5.x wrapper}
\PYGZhy{}D\PYG{+w}{ }\PYG{n+nv}{BUILD\PYGZus{}SWIG\PYGZus{}PYTHON}\PYG{o}{=}on\PYG{+w}{  }\PYG{c+c1}{\PYGZsh{} to enable building the Python wrapper}
\PYGZhy{}D\PYG{+w}{ }\PYG{n+nv}{BUILD\PYGZus{}SWIG\PYGZus{}RUBY}\PYG{o}{=}on\PYG{+w}{    }\PYG{c+c1}{\PYGZsh{} to enable building the Ruby wrapper}
\PYGZhy{}D\PYG{+w}{ }\PYG{n+nv}{BUILD\PYGZus{}SWIG\PYGZus{}TCL}\PYG{o}{=}on\PYG{+w}{     }\PYG{c+c1}{\PYGZsh{} to enable building the Tcl wrapper}
\end{sphinxVerbatim}

\sphinxAtStartPar
Manual building allows a little more flexibility. E.g. one can choose
the name of the module and build and use a dynamically loaded object
for Tcl with:

\begin{sphinxVerbatim}[commandchars=\\\{\}]
swig\PYG{+w}{ }\PYGZhy{}tcl\PYG{+w}{ }\PYGZhy{}module\PYG{+w}{ }tcllammps\PYG{+w}{ }lammps.i
gcc\PYG{+w}{ }\PYGZhy{}fPIC\PYG{+w}{ }\PYGZhy{}shared\PYG{+w}{ }\PYG{k}{\PYGZdl{}(}pkg\PYGZhy{}config\PYG{+w}{ }tcl\PYG{+w}{ }\PYGZhy{}\PYGZhy{}cflags\PYG{k}{)}\PYG{+w}{ }\PYGZhy{}o\PYG{+w}{ }tcllammps.so\PYG{+w}{ }\PYG{l+s+se}{\PYGZbs{}}
\PYG{+w}{            }lammps\PYGZus{}wrap.c\PYG{+w}{ }\PYGZhy{}L\PYG{+w}{ }../src/\PYG{+w}{ }\PYGZhy{}llammps
tclsh
\end{sphinxVerbatim}

\sphinxAtStartPar
Or one can build an extended Tcl shell command with the wrapped
functions included with:

\begin{sphinxVerbatim}[commandchars=\\\{\}]
swig\PYG{+w}{ }\PYGZhy{}tcl\PYG{+w}{ }\PYGZhy{}module\PYG{+w}{ }tcllmps\PYG{+w}{ }lammps\PYGZus{}shell.i
gcc\PYG{+w}{ }\PYGZhy{}o\PYG{+w}{ }tcllmpsh\PYG{+w}{ }lammps\PYGZus{}wrap.c\PYG{+w}{ }\PYGZhy{}Xlinker\PYG{+w}{ }\PYGZhy{}export\PYGZhy{}dynamic\PYG{+w}{ }\PYG{l+s+se}{\PYGZbs{}}
\PYG{+w}{         }\PYGZhy{}DHAVE\PYGZus{}CONFIG\PYGZus{}H\PYG{+w}{ }\PYG{k}{\PYGZdl{}(}pkg\PYGZhy{}config\PYG{+w}{ }tcl\PYG{+w}{ }\PYGZhy{}\PYGZhy{}cflags\PYG{k}{)}\PYG{+w}{ }\PYG{l+s+se}{\PYGZbs{}}
\PYG{+w}{         }\PYG{k}{\PYGZdl{}(}pkg\PYGZhy{}config\PYG{+w}{ }tcl\PYG{+w}{ }\PYGZhy{}\PYGZhy{}libs\PYG{k}{)}\PYG{+w}{ }\PYGZhy{}L\PYG{+w}{ }../src\PYG{+w}{ }\PYGZhy{}llammps
\end{sphinxVerbatim}

\sphinxAtStartPar
In both cases it is assumed that the LAMMPS library was compiled
as a shared library in the \sphinxcode{\sphinxupquote{src}} folder. Otherwise the last
part of the commands needs to be adjusted.


\subsubsection{Utility functions}
\label{\detokenize{Tools:utility-functions}}
\sphinxAtStartPar
Definitions for several utility functions required to manage and access
data passed or returned as pointers are included in the \sphinxcode{\sphinxupquote{lammps.i}}
file.  So most of the functionality of the library interface should be
accessible.  What works and what does not depends a bit on the
individual language for which the wrappers are built and how well SWIG
supports those.  The \sphinxhref{https://swig.org/doc.html}{SWIG documentation}
has very detailed instructions and recommendations.


\subsubsection{Usage examples}
\label{\detokenize{Tools:usage-examples}}
\sphinxAtStartPar
The \sphinxcode{\sphinxupquote{tools/swig}} folder has multiple shell scripts, \sphinxcode{\sphinxupquote{run\_\textless{}name\textgreater{}\_example.sh}}
that will create a small example script and demonstrate how to load
the wrapper and run LAMMPS through it in the corresponding programming
language.

\sphinxAtStartPar
For illustration purposes below is a part of the Tcl example script.

\begin{sphinxVerbatim}[commandchars=\\\{\}]
\PYG{n+nb}{load}\PYG{+w}{ }.\PYG{o}{/}tcllammps.so
\PYG{k}{set}\PYG{+w}{ }lmp\PYG{+w}{ }\PYG{k}{[}\PYG{n+nv}{lammps\PYGZus{}open\PYGZus{}no\PYGZus{}mpi}\PYG{+w}{ }\PYG{l+m+mi}{0}\PYG{+w}{ }NULL\PYG{+w}{ }NULL\PYG{k}{]}
\PYG{n+nv}{lammps\PYGZus{}command}\PYG{+w}{ }\PYG{n+nv}{\PYGZdl{}lmp}\PYG{+w}{ }\PYG{l+s+s2}{\PYGZdq{}}\PYG{l+s+s2}{u}\PYG{l+s+s2}{n}\PYG{l+s+s2}{i}\PYG{l+s+s2}{t}\PYG{l+s+s2}{s}\PYG{l+s+s2}{ }\PYG{l+s+s2}{r}\PYG{l+s+s2}{e}\PYG{l+s+s2}{a}\PYG{l+s+s2}{l}\PYG{l+s+s2}{\PYGZdq{}}
\PYG{n+nv}{lammps\PYGZus{}command}\PYG{+w}{ }\PYG{n+nv}{\PYGZdl{}lmp}\PYG{+w}{ }\PYG{l+s+s2}{\PYGZdq{}}\PYG{l+s+s2}{l}\PYG{l+s+s2}{a}\PYG{l+s+s2}{t}\PYG{l+s+s2}{t}\PYG{l+s+s2}{i}\PYG{l+s+s2}{c}\PYG{l+s+s2}{e}\PYG{l+s+s2}{ }\PYG{l+s+s2}{f}\PYG{l+s+s2}{c}\PYG{l+s+s2}{c}\PYG{l+s+s2}{ }\PYG{l+s+s2}{2}\PYG{l+s+s2}{.}\PYG{l+s+s2}{5}\PYG{l+s+s2}{\PYGZdq{}}
\PYG{n+nv}{lammps\PYGZus{}command}\PYG{+w}{ }\PYG{n+nv}{\PYGZdl{}lmp}\PYG{+w}{ }\PYG{l+s+s2}{\PYGZdq{}}\PYG{l+s+s2}{r}\PYG{l+s+s2}{e}\PYG{l+s+s2}{g}\PYG{l+s+s2}{i}\PYG{l+s+s2}{o}\PYG{l+s+s2}{n}\PYG{l+s+s2}{ }\PYG{l+s+s2}{b}\PYG{l+s+s2}{o}\PYG{l+s+s2}{x}\PYG{l+s+s2}{ }\PYG{l+s+s2}{b}\PYG{l+s+s2}{l}\PYG{l+s+s2}{o}\PYG{l+s+s2}{c}\PYG{l+s+s2}{k}\PYG{l+s+s2}{ }\PYG{l+s+s2}{\PYGZhy{}}\PYG{l+s+s2}{5}\PYG{l+s+s2}{ }\PYG{l+s+s2}{5}\PYG{l+s+s2}{ }\PYG{l+s+s2}{\PYGZhy{}}\PYG{l+s+s2}{5}\PYG{l+s+s2}{ }\PYG{l+s+s2}{5}\PYG{l+s+s2}{ }\PYG{l+s+s2}{\PYGZhy{}}\PYG{l+s+s2}{5}\PYG{l+s+s2}{ }\PYG{l+s+s2}{5}\PYG{l+s+s2}{\PYGZdq{}}
\PYG{n+nv}{lammps\PYGZus{}command}\PYG{+w}{ }\PYG{n+nv}{\PYGZdl{}lmp}\PYG{+w}{ }\PYG{l+s+s2}{\PYGZdq{}}\PYG{l+s+s2}{c}\PYG{l+s+s2}{r}\PYG{l+s+s2}{e}\PYG{l+s+s2}{a}\PYG{l+s+s2}{t}\PYG{l+s+s2}{e}\PYG{l+s+s2}{\PYGZus{}}\PYG{l+s+s2}{b}\PYG{l+s+s2}{o}\PYG{l+s+s2}{x}\PYG{l+s+s2}{ }\PYG{l+s+s2}{1}\PYG{l+s+s2}{ }\PYG{l+s+s2}{b}\PYG{l+s+s2}{o}\PYG{l+s+s2}{x}\PYG{l+s+s2}{\PYGZdq{}}
\PYG{n+nv}{lammps\PYGZus{}command}\PYG{+w}{ }\PYG{n+nv}{\PYGZdl{}lmp}\PYG{+w}{ }\PYG{l+s+s2}{\PYGZdq{}}\PYG{l+s+s2}{c}\PYG{l+s+s2}{r}\PYG{l+s+s2}{e}\PYG{l+s+s2}{a}\PYG{l+s+s2}{t}\PYG{l+s+s2}{e}\PYG{l+s+s2}{\PYGZus{}}\PYG{l+s+s2}{a}\PYG{l+s+s2}{t}\PYG{l+s+s2}{o}\PYG{l+s+s2}{m}\PYG{l+s+s2}{s}\PYG{l+s+s2}{ }\PYG{l+s+s2}{1}\PYG{l+s+s2}{ }\PYG{l+s+s2}{b}\PYG{l+s+s2}{o}\PYG{l+s+s2}{x}\PYG{l+s+s2}{\PYGZdq{}}

\PYG{k}{set}\PYG{+w}{ }dt\PYG{+w}{ }\PYG{k}{[}\PYG{n+nv}{doublep\PYGZus{}value}\PYG{+w}{ }\PYG{k}{[}\PYG{n+nv}{voidp\PYGZus{}to\PYGZus{}doublep}\PYG{+w}{ }\PYG{k}{[}\PYG{n+nv}{lammps\PYGZus{}extract\PYGZus{}global}\PYG{+w}{ }\PYG{n+nv}{\PYGZdl{}lmp}\PYG{+w}{ }dt\PYG{k}{]}\PYG{k}{]}\PYG{k}{]}
\PYG{n+nb}{puts}\PYG{+w}{ }\PYG{l+s+s2}{\PYGZdq{}}\PYG{l+s+s2}{L}\PYG{l+s+s2}{A}\PYG{l+s+s2}{M}\PYG{l+s+s2}{M}\PYG{l+s+s2}{P}\PYG{l+s+s2}{S}\PYG{l+s+s2}{ }\PYG{l+s+s2}{v}\PYG{l+s+s2}{e}\PYG{l+s+s2}{r}\PYG{l+s+s2}{s}\PYG{l+s+s2}{i}\PYG{l+s+s2}{o}\PYG{l+s+s2}{n}\PYG{l+s+s2}{ }\PYG{l+s+s2}{\PYGZdl{}}\PYG{l+s+s2}{v}\PYG{l+s+s2}{e}\PYG{l+s+s2}{r}\PYG{l+s+s2}{\PYGZdq{}}
\PYG{n+nb}{puts}\PYG{+w}{ }\PYG{k}{[}\PYG{n+nb}{format}\PYG{+w}{ }\PYG{l+s+s2}{\PYGZdq{}}\PYG{l+s+s2}{N}\PYG{l+s+s2}{u}\PYG{l+s+s2}{m}\PYG{l+s+s2}{b}\PYG{l+s+s2}{e}\PYG{l+s+s2}{r}\PYG{l+s+s2}{ }\PYG{l+s+s2}{o}\PYG{l+s+s2}{f}\PYG{l+s+s2}{ }\PYG{l+s+s2}{c}\PYG{l+s+s2}{r}\PYG{l+s+s2}{e}\PYG{l+s+s2}{a}\PYG{l+s+s2}{t}\PYG{l+s+s2}{e}\PYG{l+s+s2}{d}\PYG{l+s+s2}{ }\PYG{l+s+s2}{a}\PYG{l+s+s2}{t}\PYG{l+s+s2}{o}\PYG{l+s+s2}{m}\PYG{l+s+s2}{s}\PYG{l+s+s2}{:}\PYG{l+s+s2}{ }\PYG{l+s+s2}{\PYGZpc{}}\PYG{l+s+s2}{g}\PYG{l+s+s2}{\PYGZdq{}}\PYG{+w}{ }\PYG{k}{[}\PYG{n+nv}{lammps\PYGZus{}get\PYGZus{}natoms}\PYG{+w}{ }\PYG{n+nv}{\PYGZdl{}lmp}\PYG{k}{]}\PYG{k}{]}
\PYG{n+nb}{puts}\PYG{+w}{ }\PYG{l+s+s2}{\PYGZdq{}}\PYG{l+s+s2}{C}\PYG{l+s+s2}{u}\PYG{l+s+s2}{r}\PYG{l+s+s2}{r}\PYG{l+s+s2}{e}\PYG{l+s+s2}{n}\PYG{l+s+s2}{t}\PYG{l+s+s2}{ }\PYG{l+s+s2}{s}\PYG{l+s+s2}{i}\PYG{l+s+s2}{z}\PYG{l+s+s2}{e}\PYG{l+s+s2}{ }\PYG{l+s+s2}{o}\PYG{l+s+s2}{f}\PYG{l+s+s2}{ }\PYG{l+s+s2}{t}\PYG{l+s+s2}{i}\PYG{l+s+s2}{m}\PYG{l+s+s2}{e}\PYG{l+s+s2}{s}\PYG{l+s+s2}{t}\PYG{l+s+s2}{e}\PYG{l+s+s2}{p}\PYG{l+s+s2}{:}\PYG{l+s+s2}{ }\PYG{l+s+s2}{\PYGZdl{}}\PYG{l+s+s2}{d}\PYG{l+s+s2}{t}\PYG{l+s+s2}{\PYGZdq{}}
\PYG{n+nb}{puts}\PYG{+w}{ }\PYG{l+s+s2}{\PYGZdq{}}\PYG{l+s+s2}{L}\PYG{l+s+s2}{A}\PYG{l+s+s2}{M}\PYG{l+s+s2}{M}\PYG{l+s+s2}{P}\PYG{l+s+s2}{S}\PYG{l+s+s2}{ }\PYG{l+s+s2}{v}\PYG{l+s+s2}{e}\PYG{l+s+s2}{r}\PYG{l+s+s2}{s}\PYG{l+s+s2}{i}\PYG{l+s+s2}{o}\PYG{l+s+s2}{n}\PYG{l+s+s2}{:}\PYG{l+s+s2}{ }\PYG{l+s+s2}{[}\PYG{l+s+s2}{l}\PYG{l+s+s2}{a}\PYG{l+s+s2}{m}\PYG{l+s+s2}{m}\PYG{l+s+s2}{p}\PYG{l+s+s2}{s}\PYG{l+s+s2}{\PYGZus{}}\PYG{l+s+s2}{v}\PYG{l+s+s2}{e}\PYG{l+s+s2}{r}\PYG{l+s+s2}{s}\PYG{l+s+s2}{i}\PYG{l+s+s2}{o}\PYG{l+s+s2}{n}\PYG{l+s+s2}{ }\PYG{l+s+s2}{\PYGZdl{}}\PYG{l+s+s2}{l}\PYG{l+s+s2}{m}\PYG{l+s+s2}{p}\PYG{l+s+s2}{]}\PYG{l+s+s2}{\PYGZdq{}}
\PYG{n+nv}{lammps\PYGZus{}close}\PYG{+w}{ }\PYG{n+nv}{\PYGZdl{}lmp}
\end{sphinxVerbatim}


\bigskip\hrule\bigskip



\subsection{tabulate tool}
\label{\detokenize{Tools:tabulate-tool}}\label{\detokenize{Tools:tabulate}}
\sphinxAtStartPar
\DUrole{versionmodified,added}{Added in version 22Dec2022.}

\sphinxAtStartPar
The \sphinxcode{\sphinxupquote{tabulate}} folder contains Python scripts scripts to generate tabulated
potential files for LAMMPS.  The bulk of the code is in the \sphinxcode{\sphinxupquote{tabulate}} module
in the \sphinxcode{\sphinxupquote{tabulate.py}} file.  Some example files demonstrating its use are
included.  See the README file for more information.


\bigskip\hrule\bigskip



\subsection{tinker tool}
\label{\detokenize{Tools:tinker-tool}}\label{\detokenize{Tools:tinker}}
\sphinxAtStartPar
The \sphinxcode{\sphinxupquote{tinker}} folder contains Python scripts scripts to convert Tinker input
files to LAMMPS.

\sphinxAtStartPar
See the README file for more information.

\sphinxAtStartPar
Those scripts were written by Steve Plimpton sjplimp at gmail.com


\bigskip\hrule\bigskip



\subsection{valgrind tool}
\label{\detokenize{Tools:valgrind-tool}}\label{\detokenize{Tools:valgrind}}
\sphinxAtStartPar
The \sphinxcode{\sphinxupquote{valgrind}} folder contains additional suppressions fur LAMMPS when using
valgrind’s memcheck tool to search for memory access violation and memory
leaks. These suppressions are automatically invoked when running tests through
CMake “ctest \sphinxhyphen{}T memcheck”. See the provided README file to add these
suppressions when running LAMMPS.


\bigskip\hrule\bigskip



\subsection{vim tool}
\label{\detokenize{Tools:vim-tool}}\label{\detokenize{Tools:vim}}
\sphinxAtStartPar
The files in the \sphinxcode{\sphinxupquote{tools/vim}} directory are add\sphinxhyphen{}ons to the VIM editor
that allow easier editing of LAMMPS input scripts.  See the \sphinxcode{\sphinxupquote{README.txt}}
file for details.

\sphinxAtStartPar
These files were provided by Gerolf Ziegenhain (gerolf at
ziegenhain.com)


\bigskip\hrule\bigskip



\subsection{xmgrace tool}
\label{\detokenize{Tools:xmgrace-tool}}\label{\detokenize{Tools:xmgrace}}
\sphinxAtStartPar
The files in the tools/xmgrace directory can be used to plot the
thermodynamic data in LAMMPS log files via the xmgrace plotting
package.  There are several tools in the directory that can be used in
post\sphinxhyphen{}processing mode.  The lammpsplot.cpp file can be compiled and
used to create plots from the current state of a running LAMMPS
simulation.

\sphinxAtStartPar
See the README file for details.

\sphinxAtStartPar
These files were provided by Vikas Varshney (vv0210 at gmail.com)

\sphinxstepscope


\chapter{Errors}
\label{\detokenize{Errors:errors}}\label{\detokenize{Errors::doc}}
\sphinxAtStartPar
These doc pages describe many of the error and warning message you can
encounter when using LAMMPS.  The common problems include conceptual
issues.  The messages and warnings doc pages give complete lists of all
the messages the code may generate, with additional details for many of
them.

\sphinxstepscope


\section{Common problems}
\label{\detokenize{Errors_common:common-problems}}\label{\detokenize{Errors_common::doc}}
\sphinxAtStartPar
If two LAMMPS runs do not produce the exact same answer on different
machines or different numbers of processors, this is typically not a
bug.  In theory you should get identical answers on any number of
processors and on any machine.  In practice, numerical round\sphinxhyphen{}off can
cause slight differences and eventual divergence of molecular dynamics
phase space trajectories within a few 100s or few 1000s of timesteps.
However, the statistical properties of the two runs (e.g. average
energy or temperature) should still be the same.

\sphinxAtStartPar
If the {\hyperref[\detokenize{velocity::doc}]{\sphinxcrossref{\DUrole{doc}{velocity}}}} command is used to set initial atom
velocities, a particular atom can be assigned a different velocity
when the problem is run on a different number of processors or on
different machines.  If this happens, the phase space trajectories of
the two simulations will rapidly diverge.  See the discussion of the
\sphinxstyleemphasis{loop} option in the {\hyperref[\detokenize{velocity::doc}]{\sphinxcrossref{\DUrole{doc}{velocity}}}} command for details and
options that avoid this issue.

\sphinxAtStartPar
Similarly, the {\hyperref[\detokenize{create_atoms::doc}]{\sphinxcrossref{\DUrole{doc}{create\_atoms}}}} command generates a
lattice of atoms.  For the same physical system, the ordering and
numbering of atoms by atom ID may be different depending on the number
of processors.

\sphinxAtStartPar
Some commands use random number generators which may be setup to
produce different random number streams on each processor and hence
will produce different effects when run on different numbers of
processors.  A commonly\sphinxhyphen{}used example is the {\hyperref[\detokenize{fix_langevin::doc}]{\sphinxcrossref{\DUrole{doc}{fix langevin}}}} command for thermostatting.

\sphinxAtStartPar
A LAMMPS simulation typically has two stages, setup and run.  Most
LAMMPS errors are detected at setup time; others like a bond
stretching too far may not occur until the middle of a run.

\sphinxAtStartPar
LAMMPS tries to flag errors and print informative error messages so
you can fix the problem.  For most errors it will also print the last
input script command that it was processing.  Of course, LAMMPS cannot
figure out your physics or numerical mistakes, like choosing too big a
timestep, specifying erroneous force field coefficients, or putting 2
atoms on top of each other!  If you run into errors that LAMMPS
does not catch that you think it should flag, please send an email to
the \sphinxhref{https://www.lammps.org/authors.html}{developers} or create an new
topic on the dedicated \sphinxhref{https://matsci.org/lammps/}{MatSci forum section}.

\sphinxAtStartPar
If you get an error message about an invalid command in your input
script, you can determine what command is causing the problem by
looking in the log.lammps file or using the {\hyperref[\detokenize{echo::doc}]{\sphinxcrossref{\DUrole{doc}{echo command}}}}
to see it on the screen.  If you get an error like “Invalid …
style”, with … being fix, compute, pair, etc, it means that you
mistyped the style name or that the command is part of an optional
package which was not compiled into your executable.  The list of
available styles in your executable can be listed by using
{\hyperref[\detokenize{Run_options::doc}]{\sphinxcrossref{\DUrole{doc}{the \sphinxhyphen{}h command\sphinxhyphen{}line switch}}}}.  The installation and
compilation of optional packages is explained on the
{\hyperref[\detokenize{Build_package::doc}]{\sphinxcrossref{\DUrole{doc}{Build packages}}}} doc page.

\sphinxAtStartPar
For a given command, LAMMPS expects certain arguments in a specified
order.  If you mess this up, LAMMPS will often flag the error, but it
may also simply read a bogus argument and assign a value that is
valid, but not what you wanted.  E.g. trying to read the string “abc”
as an integer value of 0.  Careful reading of the associated doc page
for the command should allow you to fix these problems. In most cases,
where LAMMPS expects to read a number, either integer or floating point,
it performs a stringent test on whether the provided input actually
is an integer or floating\sphinxhyphen{}point number, respectively, and reject the
input with an error message (for instance, when an integer is required,
but a floating\sphinxhyphen{}point number 1.0 is provided):

\begin{sphinxVerbatim}[commandchars=\\\{\}]
ERROR: Expected integer parameter instead of \PYGZsq{}1.0\PYGZsq{} in input script or data file
\end{sphinxVerbatim}

\sphinxAtStartPar
Some commands allow for using variable references in place of numeric
constants so that the value can be evaluated and may change over the
course of a run.  This is typically done with the syntax \sphinxstyleemphasis{v\_name} for a
parameter, where name is the name of the variable. On the other hand,
immediate variable expansion with the syntax \$\{name\} is performed while
reading the input and before parsing commands,

\begin{sphinxadmonition}{note}{Note:}
\sphinxAtStartPar
Using a variable reference (i.e. \sphinxstyleemphasis{v\_name}) is only allowed if
the documentation of the corresponding command explicitly says it is.
Otherwise, you will receive an error message of this kind:
\end{sphinxadmonition}

\begin{sphinxVerbatim}[commandchars=\\\{\}]
ERROR: Expected floating point parameter instead of \PYGZsq{}v\PYGZus{}name\PYGZsq{} in input script or data file
\end{sphinxVerbatim}

\sphinxAtStartPar
Generally, LAMMPS will print a message to the screen and logfile and
exit gracefully when it encounters a fatal error.  Sometimes it will
print a WARNING to the screen and logfile and continue on; you can
decide if the WARNING is important or not.  A WARNING message that is
generated in the middle of a run is only printed to the screen, not to
the logfile, to avoid cluttering up thermodynamic output.  If LAMMPS
crashes or hangs without spitting out an error message first then it
could be a bug (see {\hyperref[\detokenize{Errors_bugs::doc}]{\sphinxcrossref{\DUrole{doc}{this section}}}}) or one of the following
cases:

\sphinxAtStartPar
LAMMPS runs in the available memory a processor allows to be
allocated.  Most reasonable MD runs are compute limited, not memory
limited, so this should not be a bottleneck on most platforms.  Almost
all large memory allocations in the code are done via C\sphinxhyphen{}style malloc’s
which will generate an error message if you run out of memory.
Smaller chunks of memory are allocated via C++ “new” statements.  If
you are unlucky you could run out of memory just when one of these
small requests is made, in which case the code will crash or hang (in
parallel), since LAMMPS does not trap on those errors.

\sphinxAtStartPar
Illegal arithmetic can cause LAMMPS to run slow or crash.  This is
typically due to invalid physics and numerics that your simulation is
computing.  If you see wild thermodynamic values or NaN values in your
LAMMPS output, something is wrong with your simulation.  If you
suspect this is happening, it is a good idea to print out
thermodynamic info frequently (e.g. every timestep) via the
{\hyperref[\detokenize{thermo::doc}]{\sphinxcrossref{\DUrole{doc}{thermo}}}} so you can monitor what is happening.
Visualizing the atom movement is also a good idea to ensure your model
is behaving as you expect.

\sphinxAtStartPar
In parallel, one way LAMMPS can hang is due to how different MPI
implementations handle buffering of messages.  If the code hangs
without an error message, it may be that you need to specify an MPI
setting or two (usually via an environment variable) to enable
buffering or boost the sizes of messages that can be buffered.

\sphinxstepscope


\section{Error and warning details}
\label{\detokenize{Errors_details:error-and-warning-details}}\label{\detokenize{Errors_details::doc}}
\sphinxAtStartPar
Many errors or warnings are self\sphinxhyphen{}explanatory and thus straightforward to
resolve.  However, there are also cases, where there is no single cause
and explanation, where LAMMPS can only detect symptoms of an error but
not the exact cause, or where the explanation needs to be more detailed than
what can be fit into a message printed by the program.  The following are
discussions of such cases.


\subsection{Unknown identifier in data file}
\label{\detokenize{Errors_details:unknown-identifier-in-data-file}}\label{\detokenize{Errors_details:err0001}}
\sphinxAtStartPar
This error happens when LAMMPS encounters a line of text with an
unexpected keyword while {\hyperref[\detokenize{read_data::doc}]{\sphinxcrossref{\DUrole{doc}{reading a data file}}}}.  This
would be either header keywords or section header keywords.  This is
most commonly due to a mistyped keyword or due to a keyword that is
inconsistent with the {\hyperref[\detokenize{atom_style::doc}]{\sphinxcrossref{\DUrole{doc}{atom style}}}} used.

\sphinxAtStartPar
The header section informs LAMMPS how many entries or lines are expected
in the various sections (like Atoms, Masses, Pair Coeffs, \sphinxstyleemphasis{etc.}) of
the data file.  If there is a mismatch, LAMMPS will either keep reading
beyond the end of a section or stop reading before the section has
ended.  In that case the next line will not contain a recognized keyword.

\sphinxAtStartPar
Such a mismatch can also happen when the first line of the data
is \sphinxstyleemphasis{not} a comment as required by the format, but a line with a valid
header keyword.  That would result in LAMMPS expecting, for instance,
0 atoms because the “atoms” header line is the first line and thus
treated as a comment.

\sphinxAtStartPar
Another possibility to trigger this error is to have a keyword in the
data file that corresponds to a fix (e.g. {\hyperref[\detokenize{fix_cmap::doc}]{\sphinxcrossref{\DUrole{doc}{fix cmap}}}})
but the {\hyperref[\detokenize{read_data::doc}]{\sphinxcrossref{\DUrole{doc}{read\_data}}}} command is missing the (optional)
arguments that identify the fix and the header keyword and section
keyword or those arguments are inconsistent with the keywords in the
data file.


\subsection{Incorrect format in … section of data file}
\label{\detokenize{Errors_details:incorrect-format-in-section-of-data-file}}\label{\detokenize{Errors_details:err0002}}
\sphinxAtStartPar
This error happens when LAMMPS reads the contents of a section of a
{\hyperref[\detokenize{read_data::doc}]{\sphinxcrossref{\DUrole{doc}{data file}}}} and the number of parameters in the line
differs from what is expected.  This most commonly happens, when the
atom style is different from what is expected for a specific data file
since changing the atom style usually changes the format of the line.

\sphinxAtStartPar
This error can also happen when the number of entries indicated in the
header of a data file (e.g. the number of atoms) is larger than the
number of lines provided (e.g. in the corresponding Atoms section)
and then LAMMPS will continue reading into the next section and that
would have a completely different format.

\sphinxstepscope


\section{Reporting bugs}
\label{\detokenize{Errors_bugs:reporting-bugs}}\label{\detokenize{Errors_bugs::doc}}
\sphinxAtStartPar
If you are confident that you have found a bug in LAMMPS, please follow
the steps outlined below:
\begin{itemize}
\item {} 
\sphinxAtStartPar
Check the \sphinxhref{https://www.lammps.org/bug.html}{New features and bug fixes} section of the \sphinxhref{https://www.lammps.org}{LAMMPS WWW site} or the
\sphinxhref{https://github.com/lammps/lammps/releases}{GitHub Releases page} to
see if the bug has already been addressed in a patch release.

\item {} 
\sphinxAtStartPar
Check that your issue can be reproduced with the latest development
version of LAMMPS.

\item {} 
\sphinxAtStartPar
Check the manual carefully to verify that the unexpected behavior you
are observing is indeed in conflict with the documentation

\item {} 
\sphinxAtStartPar
Check the \sphinxhref{https://github.com/lammps/lammps/issues}{GitHub Issue page}
if your issue has already been reported and if it is still open.

\item {} 
\sphinxAtStartPar
Check the \sphinxhref{https://github.com/lammps/lammps/pulls}{GitHub Pull Requests page}
to see if there is already a fix for your bug pending.

\item {} 
\sphinxAtStartPar
Check the \sphinxhref{https://matsci.org/lammps/}{LAMMPS forum at MatSci}
to see if the issue has been discussed before.

\end{itemize}

\sphinxAtStartPar
If none of these steps yields any useful information, please file a new
bug report on the \sphinxhref{https://github.com/lammps/lammps/issues}{GitHub Issue page}.
The website will offer you to select a suitable template with explanations
and then you should replace those explanations with the information that
you can provide to reproduce your issue.

\sphinxAtStartPar
The most useful thing you can do to help us verify and fix a bug is to
isolate the problem.  Run it on the smallest number of atoms and fewest
number of processors with the simplest input script that reproduces the
bug.  Try to identify what command or combination of commands is causing
the problem and upload the complete input deck as a tar or zip archive.
Please avoid using binary restart files unless the issue requires it.
In the latter case you should also include an input deck to quickly
generate this restart from a data file or a simple additional input.
This input deck can be used with tools like a debugger or \sphinxhref{https://valgrind.org}{valgrind} to further {\hyperref[\detokenize{Errors_debug::doc}]{\sphinxcrossref{\DUrole{doc}{debug the crash}}}}.

\sphinxAtStartPar
You may also post a message in the \sphinxhref{https://matsci.org/c/lammps/lammps-development/}{development category of the LAMMPS
forum at MatSci}
describing the problem with the same kind of information.  The forum can
provide a faster response, especially if the bug reported is actually
expected behavior or other LAMMPS users have come across it before.

\sphinxstepscope


\section{Debugging crashes}
\label{\detokenize{Errors_debug:debugging-crashes}}\label{\detokenize{Errors_debug::doc}}
\sphinxAtStartPar
If LAMMPS crashes with a “segmentation fault” or a “bus error” or
similar message, then you can use the following two methods to further
narrow down the origin of the issue.  This will help the LAMMPS
developers (or yourself) to understand the reason for the crash and
apply a fix (either to the input script or the source code).
This requires that your LAMMPS executable includes the required
{\hyperref[\detokenize{Build_basics:debug}]{\sphinxcrossref{\DUrole{std,std-ref}{debug information}}}}. Otherwise it is not possible to
look up the names of functions or variables.

\sphinxAtStartPar
The following patch will introduce a bug into the code for pair style
{\hyperref[\detokenize{pair_lj::doc}]{\sphinxcrossref{\DUrole{doc}{lj/cut}}}} when using the \sphinxcode{\sphinxupquote{examples/melt/in.melt}} input.
We use it to show how to identify the origin of a segmentation fault.

\begin{sphinxVerbatim}[commandchars=\\\{\}]
\PYG{g+gd}{\PYGZhy{}\PYGZhy{}\PYGZhy{} a/src/pair\PYGZus{}lj\PYGZus{}cut.cpp}
\PYG{g+gi}{+++ b/src/pair\PYGZus{}lj\PYGZus{}cut.cpp}
\PYG{g+gu}{@@ \PYGZhy{}81,6 +81,7 @@ void PairLJCut::compute(int eflag, int vflag)}
\PYG{+w}{ }  int nlocal = atom\PYGZhy{}\PYGZgt{}nlocal;
\PYG{+w}{ }  double *special\PYGZus{}lj = force\PYGZhy{}\PYGZgt{}special\PYGZus{}lj;
\PYG{+w}{ }  int newton\PYGZus{}pair = force\PYGZhy{}\PYGZgt{}newton\PYGZus{}pair;
\PYG{g+gi}{+  double comx = 0.0;}

\PYG{+w}{ }  inum = list\PYGZhy{}\PYGZgt{}inum;
\PYG{+w}{ }  ilist = list\PYGZhy{}\PYGZgt{}ilist;
\PYG{g+gu}{@@ \PYGZhy{}134,8 +135,10 @@ void PairLJCut::compute(int eflag, int vflag)}
\PYG{+w}{ }                             evdwl,0.0,fpair,delx,dely,delz);
\PYG{+w}{ }      \PYGZcb{}
\PYG{+w}{ }    \PYGZcb{}
\PYG{g+gd}{\PYGZhy{}  \PYGZcb{}}

\PYG{g+gi}{+    comx += atom\PYGZhy{}\PYGZgt{}rmass[i]*x[i][0]; /* BUG */}
\PYG{g+gi}{+  \PYGZcb{}}
\PYG{g+gi}{+  printf(\PYGZdq{}comx = \PYGZpc{}g\PYGZbs{}n\PYGZdq{},comx);}
\PYG{+w}{ }  if (vflag\PYGZus{}fdotr) virial\PYGZus{}fdotr\PYGZus{}compute();
\PYG{+w}{ }\PYGZcb{}
\end{sphinxVerbatim}

\sphinxAtStartPar
After recompiling LAMMPS and running the input you should get something like this:

\begin{sphinxVerbatim}[commandchars=\\\{\}]
\PYG{g+gp}{\PYGZdl{} }./lmp\PYG{+w}{ }\PYGZhy{}in\PYG{+w}{ }\PYG{k}{in}.melt
\PYG{g+go}{LAMMPS (19 Mar 2020)}
\PYG{g+go}{  using 1 OpenMP thread(s) per MPI task}
\PYG{g+go}{Lattice spacing in x,y,z = 1.6796 1.6796 1.6796}
\PYG{g+go}{Created orthogonal box = (0 0 0) to (16.796 16.796 16.796)}
\PYG{g+go}{  1 by 1 by 1 MPI processor grid}
\PYG{g+go}{Created 4000 atoms}
\PYG{g+go}{  create\PYGZus{}atoms CPU = 0.000432253 secs}
\PYG{g+go}{Neighbor list info ...}
\PYG{g+go}{  update every 20 steps, delay 0 steps, check no}
\PYG{g+go}{  max neighbors/atom: 2000, page size: 100000}
\PYG{g+go}{  master list distance cutoff = 2.8}
\PYG{g+go}{  ghost atom cutoff = 2.8}
\PYG{g+go}{  binsize = 1.4, bins = 12 12 12}
\PYG{g+go}{  1 neighbor lists, perpetual/occasional/extra = 1 0 0}
\PYG{g+go}{  (1) pair lj/cut, perpetual}
\PYG{g+go}{      attributes: half, newton on}
\PYG{g+go}{      pair build: half/bin/atomonly/newton}
\PYG{g+go}{      stencil: half/bin/3d/newton}
\PYG{g+go}{      bin: standard}
\PYG{g+go}{Setting up Verlet run ...}
\PYG{g+go}{  Unit style    : lj}
\PYG{g+go}{  Current step  : 0}
\PYG{g+go}{  Time step     : 0.005}
\PYG{g+go}{Segmentation fault (core dumped)}
\end{sphinxVerbatim}


\subsection{Using the GDB debugger to get a stack trace}
\label{\detokenize{Errors_debug:using-the-gdb-debugger-to-get-a-stack-trace}}
\sphinxAtStartPar
There are two options to use the GDB debugger for identifying the origin
of the segmentation fault or similar crash. The GDB debugger has many
more features and options, as can be seen for example its \sphinxhref{https://www.sourceware.org/gdb/documentation/}{online
documentation}.


\subsubsection{Run LAMMPS from within the debugger}
\label{\detokenize{Errors_debug:run-lammps-from-within-the-debugger}}
\sphinxAtStartPar
Running LAMMPS under the control of the debugger as shown below only
works for a single MPI rank (for debugging a program running in parallel
you usually need a parallel debugger program).  A simple way to launch
GDB is to prefix the LAMMPS command line with \sphinxcode{\sphinxupquote{gdb \sphinxhyphen{}\sphinxhyphen{}args}} and then
type the command “run” at the GDB prompt.  This will launch the
debugger, load the LAMMPS executable and its debug info, and then run
it.  When it reaches the code causing the segmentation fault, it will
stop with a message why it stopped, print the current line of code, and
drop back to the GDB prompt.

\begin{sphinxVerbatim}[commandchars=\\\{\}]
\PYG{g+gp+gpVirtualEnv}{(gdb)} \PYG{g+go}{run}
\PYG{g+go}{[...]}
\PYG{g+go}{Setting up Verlet run ...}
\PYG{g+go}{  Unit style    : lj}
\PYG{g+go}{  Current step  : 0}
\PYG{g+go}{  Time step     : 0.005}

\PYG{g+go}{Program received signal SIGSEGV, Segmentation fault.}
\PYG{g+go}{0x00000000006653ab in LAMMPS\PYGZus{}NS::PairLJCut::compute (this=0x829740, eflag=1, vflag=\PYGZlt{}optimized out\PYGZgt{}) at /home/akohlmey/compile/lammps/src/pair\PYGZus{}lj\PYGZus{}cut.cpp:139}
\PYG{g+go}{139      comx += atom\PYGZhy{}\PYGZgt{}rmass[i]*x[i][0]; /* BUG */}
\PYG{g+gp+gpVirtualEnv}{(gdb)}
\end{sphinxVerbatim}

\sphinxAtStartPar
Now typing the command “where” will show the stack of functions starting from
the current function back to “main()”.

\begin{sphinxVerbatim}[commandchars=\\\{\}]
\PYG{g+gp+gpVirtualEnv}{(gdb)} \PYG{g+go}{where}
\PYG{g+gp}{\PYGZsh{}}\PYG{l+m}{0}\PYG{+w}{  }0x00000000006653ab\PYG{+w}{ }\PYG{k}{in}\PYG{+w}{ }LAMMPS\PYGZus{}NS::PairLJCut::compute\PYG{+w}{ }\PYG{o}{(}\PYG{n+nv}{this}\PYG{o}{=}0x829740,\PYG{+w}{ }\PYG{n+nv}{eflag}\PYG{o}{=}\PYG{l+m}{1},\PYG{+w}{ }\PYG{n+nv}{vflag}\PYG{o}{=}\PYGZlt{}optimized\PYG{+w}{ }out\PYGZgt{}\PYG{o}{)}\PYG{+w}{ }at\PYG{+w}{ }/home/akohlmey/compile/lammps/src/pair\PYGZus{}lj\PYGZus{}cut.cpp:139
\PYG{g+gp}{\PYGZsh{}}\PYG{l+m}{1}\PYG{+w}{  }0x00000000004cf0a2\PYG{+w}{ }\PYG{k}{in}\PYG{+w}{ }LAMMPS\PYGZus{}NS::Verlet::setup\PYG{+w}{ }\PYG{o}{(}\PYG{n+nv}{this}\PYG{o}{=}0x7e6c90,\PYG{+w}{ }\PYG{n+nv}{flag}\PYG{o}{=}\PYG{l+m}{1}\PYG{o}{)}\PYG{+w}{ }at\PYG{+w}{ }/home/akohlmey/compile/lammps/src/verlet.cpp:131
\PYG{g+gp}{\PYGZsh{}}\PYG{l+m}{2}\PYG{+w}{  }0x000000000049db42\PYG{+w}{ }\PYG{k}{in}\PYG{+w}{ }LAMMPS\PYGZus{}NS::Run::command\PYG{+w}{ }\PYG{o}{(}\PYG{n+nv}{this}\PYG{o}{=}this@entry\PYG{o}{=}0x7fffffffcca0,\PYG{+w}{ }\PYG{n+nv}{narg}\PYG{o}{=}narg@entry\PYG{o}{=}\PYG{l+m}{1},\PYG{+w}{ }\PYG{n+nv}{arg}\PYG{o}{=}arg@entry\PYG{o}{=}0x7e8750\PYG{o}{)}
\PYG{g+go}{    at /home/akohlmey/compile/lammps/src/run.cpp:177}
\PYG{g+gp}{\PYGZsh{}}\PYG{l+m}{3}\PYG{+w}{  }0x000000000041258a\PYG{+w}{ }\PYG{k}{in}\PYG{+w}{ }LAMMPS\PYGZus{}NS::Input::command\PYGZus{}creator\PYGZlt{}LAMMPS\PYGZus{}NS::Run\PYGZgt{}\PYG{+w}{ }\PYG{o}{(}\PYG{n+nv}{lmp}\PYG{o}{=}\PYGZlt{}optimized\PYG{+w}{ }out\PYGZgt{},\PYG{+w}{ }\PYG{n+nv}{narg}\PYG{o}{=}\PYG{l+m}{1},\PYG{+w}{ }\PYG{n+nv}{arg}\PYG{o}{=}0x7e8750\PYG{o}{)}
\PYG{g+go}{    at /home/akohlmey/compile/lammps/src/input.cpp:878}
\PYG{g+gp}{\PYGZsh{}}\PYG{l+m}{4}\PYG{+w}{  }0x0000000000410ad3\PYG{+w}{ }\PYG{k}{in}\PYG{+w}{ }LAMMPS\PYGZus{}NS::Input::execute\PYGZus{}command\PYG{+w}{ }\PYG{o}{(}\PYG{n+nv}{this}\PYG{o}{=}0x7d1410\PYG{o}{)}\PYG{+w}{ }at\PYG{+w}{ }/home/akohlmey/compile/lammps/src/input.cpp:864
\PYG{g+gp}{\PYGZsh{}}\PYG{l+m}{5}\PYG{+w}{  }0x00000000004111fb\PYG{+w}{ }\PYG{k}{in}\PYG{+w}{ }LAMMPS\PYGZus{}NS::Input::file\PYG{+w}{ }\PYG{o}{(}\PYG{n+nv}{this}\PYG{o}{=}0x7d1410\PYG{o}{)}\PYG{+w}{ }at\PYG{+w}{ }/home/akohlmey/compile/lammps/src/input.cpp:229
\PYG{g+gp}{\PYGZsh{}}\PYG{l+m}{6}\PYG{+w}{  }0x000000000040933a\PYG{+w}{ }\PYG{k}{in}\PYG{+w}{ }main\PYG{+w}{ }\PYG{o}{(}\PYG{n+nv}{argc}\PYG{o}{=}\PYGZlt{}optimized\PYG{+w}{ }out\PYGZgt{},\PYG{+w}{ }\PYG{n+nv}{argv}\PYG{o}{=}\PYGZlt{}optimized\PYG{+w}{ }out\PYGZgt{}\PYG{o}{)}\PYG{+w}{ }at\PYG{+w}{ }/home/akohlmey/compile/lammps/src/main.cpp:65
\PYG{g+gp+gpVirtualEnv}{(gdb)}
\end{sphinxVerbatim}

\sphinxAtStartPar
You can also print the value of variables and see if there is anything
unexpected.  Segmentation faults, for example, commonly happen when a
pointer variable is not assigned and still initialized to NULL.

\begin{sphinxVerbatim}[commandchars=\\\{\}]
\PYG{g+gp+gpVirtualEnv}{(gdb)} \PYG{g+go}{print x}
\PYG{g+gp}{\PYGZdl{}}\PYG{n+nv}{1}\PYG{+w}{ }\PYG{o}{=}\PYG{+w}{ }\PYG{o}{(}double\PYG{+w}{ }**\PYG{o}{)}\PYG{+w}{ }0x7ffff7ca1010
\PYG{g+gp+gpVirtualEnv}{(gdb)} \PYG{g+go}{print i}
\PYG{g+gp}{\PYGZdl{}}\PYG{n+nv}{2}\PYG{+w}{ }\PYG{o}{=}\PYG{+w}{ }\PYG{l+m}{0}
\PYG{g+gp+gpVirtualEnv}{(gdb)} \PYG{g+go}{print x[0]}
\PYG{g+gp}{\PYGZdl{}}\PYG{n+nv}{3}\PYG{+w}{ }\PYG{o}{=}\PYG{+w}{ }\PYG{o}{(}double\PYG{+w}{ }*\PYG{o}{)}\PYG{+w}{ }0x7ffff6d80010
\PYG{g+gp+gpVirtualEnv}{(gdb)} \PYG{g+go}{print x[0][0]}
\PYG{g+gp}{\PYGZdl{}}\PYG{n+nv}{4}\PYG{+w}{ }\PYG{o}{=}\PYG{+w}{ }\PYG{l+m}{0}
\PYG{g+gp+gpVirtualEnv}{(gdb)} \PYG{g+go}{print x[1][0]}
\PYG{g+gp}{\PYGZdl{}}\PYG{n+nv}{5}\PYG{+w}{ }\PYG{o}{=}\PYG{+w}{ }\PYG{l+m}{0}.83979809569125363
\PYG{g+gp+gpVirtualEnv}{(gdb)} \PYG{g+go}{print atom\PYGZhy{}\PYGZgt{}rmass}
\PYG{g+gp}{\PYGZdl{}}\PYG{n+nv}{6}\PYG{+w}{ }\PYG{o}{=}\PYG{+w}{ }\PYG{o}{(}double\PYG{+w}{ }*\PYG{o}{)}\PYG{+w}{ }0x0
\PYG{g+gp+gpVirtualEnv}{(gdb)}
\end{sphinxVerbatim}


\subsubsection{Inspect a core dump file with the debugger}
\label{\detokenize{Errors_debug:inspect-a-core-dump-file-with-the-debugger}}
\sphinxAtStartPar
When an executable crashes with a “core dumped” message, it creates a
file “core” or “core.\textless{}PID\#\textgreater{}” which contains the information about the
current state.  This file may be located in the folder where you ran
LAMMPS or in some hidden folder managed by the systemd daemon.  In the
latter case, you need to “extract” the core file with the \sphinxcode{\sphinxupquote{coredumpctl}}
utility to the current folder. Example: \sphinxcode{\sphinxupquote{coredumpctl \sphinxhyphen{}o core dump lmp}}.
Now you can launch the debugger to load the executable, its debug info
and the core dump and drop you to a prompt like before.

\begin{sphinxVerbatim}[commandchars=\\\{\}]
\PYG{g+gp}{\PYGZdl{} }gdb\PYG{+w}{ }lmp\PYG{+w}{ }core
\PYG{g+go}{Reading symbols from lmp...}
\PYG{g+go}{[New LWP 1928535]}
\PYG{g+go}{[Thread debugging using libthread\PYGZus{}db enabled]}
\PYG{g+go}{Using host libthread\PYGZus{}db library \PYGZdq{}/lib64/libthread\PYGZus{}db.so.1\PYGZdq{}.}
\PYG{g+go}{Core was generated by `./lmp \PYGZhy{}in in.melt\PYGZsq{}.}
\PYG{g+go}{Program terminated with signal SIGSEGV, Segmentation fault.}
\PYG{g+gp}{\PYGZsh{}}\PYG{l+m}{0}\PYG{+w}{  }0x00000000006653ab\PYG{+w}{ }\PYG{k}{in}\PYG{+w}{ }LAMMPS\PYGZus{}NS::PairLJCut::compute\PYG{+w}{ }\PYG{o}{(}\PYG{n+nv}{this}\PYG{o}{=}0x1b10740,\PYG{+w}{ }\PYG{n+nv}{eflag}\PYG{o}{=}\PYG{l+m}{1},\PYG{+w}{ }\PYG{n+nv}{vflag}\PYG{o}{=}\PYGZlt{}optimized\PYG{+w}{ }out\PYGZgt{}\PYG{o}{)}
\PYG{g+go}{    at /home/akohlmey/compile/lammps/src/pair\PYGZus{}lj\PYGZus{}cut.cpp:139}
\PYG{g+go}{139      comx += atom\PYGZhy{}\PYGZgt{}rmass[i]*x[i][0]; /* BUG */}
\PYG{g+gp+gpVirtualEnv}{(gdb)}
\end{sphinxVerbatim}

\sphinxAtStartPar
From here on, you use the same commands as shown before to get a stack
trace and print current values of (pointer) variables.


\subsection{Using valgrind to get a stack trace}
\label{\detokenize{Errors_debug:using-valgrind-to-get-a-stack-trace}}
\sphinxAtStartPar
The \sphinxhref{https://valgrind.org}{valgrind} suite of tools allows to closely
inspect the behavior of a compiled program by essentially emulating a
CPU and instrumenting the program while running.  This slows down
execution quite significantly, but can also report issues that are not
resulting in a crash.  The default valgrind tool is a memory checker and
you can use it by prefixing the normal command line with \sphinxcode{\sphinxupquote{valgrind}}.
Unlike GDB, this will also work for parallel execution, but it is
recommended to redirect the valgrind output to a file (e.g. with
\sphinxcode{\sphinxupquote{\sphinxhyphen{}\sphinxhyphen{}log\sphinxhyphen{}file=crash\sphinxhyphen{}\%p.txt}}, the \%p will be substituted with the
process ID) so that the messages of the multiple valgrind instances to
the console are not mixed.

\begin{sphinxVerbatim}[commandchars=\\\{\}]
\PYG{g+gp}{\PYGZdl{} }valgrind\PYG{+w}{ }./lmp\PYG{+w}{ }\PYGZhy{}in\PYG{+w}{ }\PYG{k}{in}.melt
\PYG{g+go}{==1933642== Memcheck, a memory error detector}
\PYG{g+go}{==1933642== Copyright (C) 2002\PYGZhy{}2017, and GNU GPL\PYGZsq{}d, by Julian Seward et al.}
\PYG{g+go}{==1933642== Using Valgrind\PYGZhy{}3.15.0 and LibVEX; rerun with \PYGZhy{}h for copyright info}
\PYG{g+go}{==1933642== Command: ./lmp \PYGZhy{}in in.melt}
\PYG{g+go}{==1933642==}
\PYG{g+go}{LAMMPS (19 Mar 2020)}
\PYG{g+go}{OMP\PYGZus{}NUM\PYGZus{}THREADS environment is not set. Defaulting to 1 thread. (src/comm.cpp:94)}
\PYG{g+go}{  using 1 OpenMP thread(s) per MPI task}
\PYG{g+go}{Lattice spacing in x,y,z = 1.6796 1.6796 1.6796}
\PYG{g+go}{Created orthogonal box = (0 0 0) to (16.796 16.796 16.796)}
\PYG{g+go}{  1 by 1 by 1 MPI processor grid}
\PYG{g+go}{Created 4000 atoms}
\PYG{g+go}{  create\PYGZus{}atoms CPU = 0.032964 secs}
\PYG{g+go}{Neighbor list info ...}
\PYG{g+go}{  update every 20 steps, delay 0 steps, check no}
\PYG{g+go}{  max neighbors/atom: 2000, page size: 100000}
\PYG{g+go}{  master list distance cutoff = 2.8}
\PYG{g+go}{  ghost atom cutoff = 2.8}
\PYG{g+go}{  binsize = 1.4, bins = 12 12 12}
\PYG{g+go}{  1 neighbor lists, perpetual/occasional/extra = 1 0 0}
\PYG{g+go}{  (1) pair lj/cut, perpetual}
\PYG{g+go}{      attributes: half, newton on}
\PYG{g+go}{      pair build: half/bin/atomonly/newton}
\PYG{g+go}{      stencil: half/bin/3d/newton}
\PYG{g+go}{      bin: standard}
\PYG{g+go}{Setting up Verlet run ...}
\PYG{g+go}{  Unit style    : lj}
\PYG{g+go}{  Current step  : 0}
\PYG{g+go}{  Time step     : 0.005}
\PYG{g+go}{==1933642== Invalid read of size 8}
\PYG{g+go}{==1933642==    at 0x6653AB: LAMMPS\PYGZus{}NS::PairLJCut::compute(int, int) (pair\PYGZus{}lj\PYGZus{}cut.cpp:139)}
\PYG{g+go}{==1933642==    by 0x4CF0A1: LAMMPS\PYGZus{}NS::Verlet::setup(int) (verlet.cpp:131)}
\PYG{g+go}{==1933642==    by 0x49DB41: LAMMPS\PYGZus{}NS::Run::command(int, char**) (run.cpp:177)}
\PYG{g+go}{==1933642==    by 0x412589: void LAMMPS\PYGZus{}NS::Input::command\PYGZus{}creator\PYGZlt{}LAMMPS\PYGZus{}NS::Run\PYGZgt{}(LAMMPS\PYGZus{}NS::LAMMPS*, int, char**) (input.cpp:881)}
\PYG{g+go}{==1933642==    by 0x410AD2: LAMMPS\PYGZus{}NS::Input::execute\PYGZus{}command() (input.cpp:864)}
\PYG{g+go}{==1933642==    by 0x4111FA: LAMMPS\PYGZus{}NS::Input::file() (input.cpp:229)}
\PYG{g+go}{==1933642==    by 0x409339: main (main.cpp:65)}
\PYG{g+go}{==1933642==  Address 0x0 is not stack\PYGZsq{}d, malloc\PYGZsq{}d or (recently) free\PYGZsq{}d}
\PYG{g+go}{==1933642==}
\end{sphinxVerbatim}

\sphinxAtStartPar
As you can see, the stack trace information is similar to that obtained
from GDB. In addition you get a more specific hint about what cause the
segmentation fault, i.e. that it is a NULL pointer dereference.  To find
out which pointer exactly was NULL, you need to use the debugger, though.

\sphinxstepscope


\section{Error messages}
\label{\detokenize{Errors_messages:error-messages}}\label{\detokenize{Errors_messages::doc}}
\sphinxAtStartPar
This is an alphabetic list of the ERROR messages LAMMPS prints out and
the reason why.  If the explanation here is not sufficient, the
documentation for the offending command may help.  Error messages also
list the source file and line number where the error was generated.
For example, a message like this:

\begin{sphinxVerbatim}[commandchars=\\\{\}]
ERROR: Illegal velocity command (velocity.cpp:78)
\end{sphinxVerbatim}

\sphinxAtStartPar
means that line \#78 in the file src/velocity.cpp generated the error.
Looking in the source code may help you figure out what went wrong.

\sphinxAtStartPar
Doc page with {\hyperref[\detokenize{Errors_warnings::doc}]{\sphinxcrossref{\DUrole{doc}{WARNING messages}}}}


\bigskip\hrule\bigskip

\begin{description}
\sphinxlineitem{\sphinxstyleemphasis{1\sphinxhyphen{}3 bond count is inconsistent}}
\sphinxAtStartPar
An inconsistency was detected when computing the number of 1\sphinxhyphen{}3
neighbors for each atom.  This likely means something is wrong with
the bond topologies you have defined.

\sphinxlineitem{\sphinxstyleemphasis{1\sphinxhyphen{}4 bond count is inconsistent}}
\sphinxAtStartPar
An inconsistency was detected when computing the number of 1\sphinxhyphen{}4
neighbors for each atom.  This likely means something is wrong with
the bond topologies you have defined.

\sphinxlineitem{\sphinxstyleemphasis{Accelerator sharing is not currently supported on system}}
\sphinxAtStartPar
Multiple MPI processes cannot share the accelerator on your
system. For NVIDIA GPUs, see the nvidia\sphinxhyphen{}smi command to change this
setting.

\sphinxlineitem{\sphinxstyleemphasis{All angle coeffs are not set}}
\sphinxAtStartPar
All angle coefficients must be set in the data file or by the
angle\_coeff command before running a simulation.

\sphinxlineitem{\sphinxstyleemphasis{All atom IDs = 0 but atom\_modify id = yes}}
\sphinxAtStartPar
Self\sphinxhyphen{}explanatory.

\sphinxlineitem{\sphinxstyleemphasis{All atoms of a swapped type must have same charge.}}
\sphinxAtStartPar
Self\sphinxhyphen{}explanatory.

\sphinxlineitem{\sphinxstyleemphasis{All atoms of a swapped type must have the same charge.}}
\sphinxAtStartPar
Self\sphinxhyphen{}explanatory.

\sphinxlineitem{\sphinxstyleemphasis{All bond coeffs are not set}}
\sphinxAtStartPar
All bond coefficients must be set in the data file or by the
bond\_coeff command before running a simulation.

\sphinxlineitem{\sphinxstyleemphasis{All dihedral coeffs are not set}}
\sphinxAtStartPar
All dihedral coefficients must be set in the data file or by the
dihedral\_coeff command before running a simulation.

\sphinxlineitem{\sphinxstyleemphasis{All improper coeffs are not set}}
\sphinxAtStartPar
All improper coefficients must be set in the data file or by the
improper\_coeff command before running a simulation.

\sphinxlineitem{\sphinxstyleemphasis{All masses are not set}}
\sphinxAtStartPar
For atom styles that define masses for each atom type, all masses must
be set in the data file or by the mass command before running a
simulation.  They must also be set before using the velocity
command.

\sphinxlineitem{\sphinxstyleemphasis{All mol IDs should be set for fix gcmc group atoms}}
\sphinxAtStartPar
The molecule flag is on, yet not all molecule ids in the fix group
have been set to non\sphinxhyphen{}zero positive values by the user. This is an
error since all atoms in the fix gcmc group are eligible for deletion,
rotation, and translation and therefore must have valid molecule ids.

\sphinxlineitem{\sphinxstyleemphasis{All pair coeffs are not set}}
\sphinxAtStartPar
All pair coefficients must be set in the data file or by the
pair\_coeff command before running a simulation.

\sphinxlineitem{\sphinxstyleemphasis{All read\_dump x,y,z fields must be specified for scaled, triclinic coords}}
\sphinxAtStartPar
For triclinic boxes and scaled coordinates you must specify all 3 of
the x,y,z fields, else LAMMPS cannot reconstruct the unscaled
coordinates.

\sphinxlineitem{\sphinxstyleemphasis{All universe/uloop variables must have same \# of values}}
\sphinxAtStartPar
Self\sphinxhyphen{}explanatory.

\sphinxlineitem{\sphinxstyleemphasis{All variables in next command must be same style}}
\sphinxAtStartPar
Self\sphinxhyphen{}explanatory.

\sphinxlineitem{\sphinxstyleemphasis{Angle atom missing in delete\_bonds}}
\sphinxAtStartPar
The delete\_bonds command cannot find one or more atoms in a particular
angle on a particular processor.  The pairwise cutoff is too short or
the atoms are too far apart to make a valid angle.

\sphinxlineitem{\sphinxstyleemphasis{Angle atom missing in set command}}
\sphinxAtStartPar
The set command cannot find one or more atoms in a particular angle on
a particular processor.  The pairwise cutoff is too short or the atoms
are too far apart to make a valid angle.

\sphinxlineitem{\sphinxstyleemphasis{Angle atoms \%d \%d \%d missing on proc \%d at step \%ld}}
\sphinxAtStartPar
One or more of three atoms needed to compute a particular angle are
missing on this processor.  Typically this is because the pairwise
cutoff is set too short or the angle has blown apart and an atom is
too far away.

\sphinxlineitem{\sphinxstyleemphasis{Angle atoms missing on proc \%d at step \%ld}}
\sphinxAtStartPar
One or more of three atoms needed to compute a particular angle are
missing on this processor.  Typically this is because the pairwise
cutoff is set too short or the angle has blown apart and an atom is
too far away.

\sphinxlineitem{\sphinxstyleemphasis{Angle coeff for hybrid has invalid style}}
\sphinxAtStartPar
Angle style hybrid uses another angle style as one of its
coefficients.  The angle style used in the angle\_coeff command or read
from a restart file is not recognized.

\sphinxlineitem{\sphinxstyleemphasis{Angle coeffs are not set}}
\sphinxAtStartPar
No angle coefficients have been assigned in the data file or via the
angle\_coeff command.

\sphinxlineitem{\sphinxstyleemphasis{Angle extent \textgreater{} half of periodic box length}}
\sphinxAtStartPar
This error was detected by the neigh\_modify check yes setting.  It is
an error because the angle atoms are so far apart it is ambiguous how
it should be defined.

\sphinxlineitem{\sphinxstyleemphasis{Angle potential must be defined for SHAKE}}
\sphinxAtStartPar
When shaking angles, an angle\_style potential must be used.

\sphinxlineitem{\sphinxstyleemphasis{Angle style hybrid cannot have hybrid as an argument}}
\sphinxAtStartPar
Self\sphinxhyphen{}explanatory.

\sphinxlineitem{\sphinxstyleemphasis{Angle style hybrid cannot have none as an argument}}
\sphinxAtStartPar
Self\sphinxhyphen{}explanatory.

\sphinxlineitem{\sphinxstyleemphasis{Angle style hybrid cannot use same angle style twice}}
\sphinxAtStartPar
Self\sphinxhyphen{}explanatory.

\sphinxlineitem{\sphinxstyleemphasis{Angle table must range from 0 to 180 degrees}}
\sphinxAtStartPar
Self\sphinxhyphen{}explanatory.

\sphinxlineitem{\sphinxstyleemphasis{Angle table parameters did not set N}}
\sphinxAtStartPar
List of angle table parameters must include N setting.

\sphinxlineitem{\sphinxstyleemphasis{Angle\_coeff command before angle\_style is defined}}
\sphinxAtStartPar
Coefficients cannot be set in the data file or via the angle\_coeff
command until an angle\_style has been assigned.

\sphinxlineitem{\sphinxstyleemphasis{Angle\_coeff command before simulation box is defined}}
\sphinxAtStartPar
The angle\_coeff command cannot be used before a read\_data,
read\_restart, or create\_box command.

\sphinxlineitem{\sphinxstyleemphasis{Angle\_coeff command when no angles allowed}}
\sphinxAtStartPar
The chosen atom style does not allow for angles to be defined.

\sphinxlineitem{\sphinxstyleemphasis{Angle\_style command when no angles allowed}}
\sphinxAtStartPar
The chosen atom style does not allow for angles to be defined.

\sphinxlineitem{\sphinxstyleemphasis{Angles assigned incorrectly}}
\sphinxAtStartPar
Angles read in from the data file were not assigned correctly to
atoms.  This means there is something invalid about the topology
definitions.

\sphinxlineitem{\sphinxstyleemphasis{Angles defined but no angle types}}
\sphinxAtStartPar
The data file header lists angles but no angle types.

\sphinxlineitem{\sphinxstyleemphasis{Append boundary must be shrink/minimum}}
\sphinxAtStartPar
The boundary style of the face where atoms are added
must be of type m (shrink/minimum).

\sphinxlineitem{\sphinxstyleemphasis{Arccos of invalid value in variable formula}}
\sphinxAtStartPar
Argument of arccos() must be between \sphinxhyphen{}1 and 1.

\sphinxlineitem{\sphinxstyleemphasis{Arcsin of invalid value in variable formula}}
\sphinxAtStartPar
Argument of arcsin() must be between \sphinxhyphen{}1 and 1.

\sphinxlineitem{\sphinxstyleemphasis{Assigning body parameters to non\sphinxhyphen{}body atom}}
\sphinxAtStartPar
Self\sphinxhyphen{}explanatory.

\sphinxlineitem{\sphinxstyleemphasis{Assigning ellipsoid parameters to non\sphinxhyphen{}ellipsoid atom}}
\sphinxAtStartPar
Self\sphinxhyphen{}explanatory.

\sphinxlineitem{\sphinxstyleemphasis{Assigning line parameters to non\sphinxhyphen{}line atom}}
\sphinxAtStartPar
Self\sphinxhyphen{}explanatory.

\sphinxlineitem{\sphinxstyleemphasis{Assigning quat to non\sphinxhyphen{}body atom}}
\sphinxAtStartPar
Self\sphinxhyphen{}explanatory.

\sphinxlineitem{\sphinxstyleemphasis{Assigning tri parameters to non\sphinxhyphen{}tri atom}}
\sphinxAtStartPar
Self\sphinxhyphen{}explanatory.

\sphinxlineitem{\sphinxstyleemphasis{At least one atom of each swapped type must be present to define charges.}}
\sphinxAtStartPar
Self\sphinxhyphen{}explanatory.

\sphinxlineitem{\sphinxstyleemphasis{Atom IDs must be consecutive for velocity create loop all}}
\sphinxAtStartPar
Self\sphinxhyphen{}explanatory.

\sphinxlineitem{\sphinxstyleemphasis{Atom IDs must be used for molecular systems}}
\sphinxAtStartPar
Atom IDs are used to identify and find partner atoms in bonds.

\sphinxlineitem{\sphinxstyleemphasis{Atom count changed in fix neb}}
\sphinxAtStartPar
This is not allowed in a NEB calculation.

\sphinxlineitem{\sphinxstyleemphasis{Atom count is inconsistent, cannot write data file}}
\sphinxAtStartPar
The sum of atoms across processors does not equal the global number
of atoms.  Probably some atoms have been lost.

\sphinxlineitem{\sphinxstyleemphasis{Atom count is inconsistent, cannot write restart file}}
\sphinxAtStartPar
Sum of atoms across processors does not equal initial total count.
This is probably because you have lost some atoms.

\sphinxlineitem{\sphinxstyleemphasis{Atom in too many rigid bodies \sphinxhyphen{} boost MAXBODY}}
\sphinxAtStartPar
Fix poems has a parameter MAXBODY (in fix\_poems.cpp) which determines
the maximum number of rigid bodies a single atom can belong to (i.e. a
multibody joint).  The bodies you have defined exceed this limit.

\sphinxlineitem{\sphinxstyleemphasis{Atom sort did not operate correctly}}
\sphinxAtStartPar
This is an internal LAMMPS error.  Please report it to the
developers.

\sphinxlineitem{\sphinxstyleemphasis{Atom style hybrid cannot have hybrid as an argument}}
\sphinxAtStartPar
Self\sphinxhyphen{}explanatory.

\sphinxlineitem{\sphinxstyleemphasis{Atom style hybrid cannot use same atom style twice}}
\sphinxAtStartPar
Self\sphinxhyphen{}explanatory.

\sphinxlineitem{\sphinxstyleemphasis{Atom style template molecule must have atom types}}
\sphinxAtStartPar
The defined molecule(s) does not specify atom types.

\sphinxlineitem{\sphinxstyleemphasis{Atom style was redefined after using fix property/atom}}
\sphinxAtStartPar
This is not allowed.

\sphinxlineitem{\sphinxstyleemphasis{Atom type must be zero in fix gcmc mol command}}
\sphinxAtStartPar
Self\sphinxhyphen{}explanatory.

\sphinxlineitem{\sphinxstyleemphasis{Atom vector in equal\sphinxhyphen{}style variable formula}}
\sphinxAtStartPar
Atom vectors generate one value per atom which is not allowed
in an equal\sphinxhyphen{}style variable.

\sphinxlineitem{\sphinxstyleemphasis{Atom\sphinxhyphen{}style variable in equal\sphinxhyphen{}style variable formula}}
\sphinxAtStartPar
Atom\sphinxhyphen{}style variables generate one value per atom which is not allowed
in an equal\sphinxhyphen{}style variable.

\sphinxlineitem{\sphinxstyleemphasis{Atom\_modify id command after simulation box is defined}}
\sphinxAtStartPar
The atom\_modify id command cannot be used after a read\_data,
read\_restart, or create\_box command.

\sphinxlineitem{\sphinxstyleemphasis{Atom\_modify map command after simulation box is defined}}
\sphinxAtStartPar
The atom\_modify map command cannot be used after a read\_data,
read\_restart, or create\_box command.

\sphinxlineitem{\sphinxstyleemphasis{Atom\_modify sort and first options cannot be used together}}
\sphinxAtStartPar
Self\sphinxhyphen{}explanatory.

\sphinxlineitem{\sphinxstyleemphasis{Atom\_style command after simulation box is defined}}
\sphinxAtStartPar
The atom\_style command cannot be used after a read\_data,
read\_restart, or create\_box command.

\sphinxlineitem{\sphinxstyleemphasis{Atom\_style line can only be used in 2d simulations}}
\sphinxAtStartPar
Self\sphinxhyphen{}explanatory.

\sphinxlineitem{\sphinxstyleemphasis{Atom\_style tri can only be used in 3d simulations}}
\sphinxAtStartPar
Self\sphinxhyphen{}explanatory.

\sphinxlineitem{\sphinxstyleemphasis{Atomfile variable could not read values}}
\sphinxAtStartPar
Check the file assigned to the variable.

\sphinxlineitem{\sphinxstyleemphasis{Atomfile variable in equal\sphinxhyphen{}style variable formula}}
\sphinxAtStartPar
Self\sphinxhyphen{}explanatory.

\sphinxlineitem{\sphinxstyleemphasis{Atomfile\sphinxhyphen{}style variable in equal\sphinxhyphen{}style variable formula}}
\sphinxAtStartPar
Self\sphinxhyphen{}explanatory.

\sphinxlineitem{\sphinxstyleemphasis{Attempt to pop empty stack in fix box/relax}}
\sphinxAtStartPar
Internal LAMMPS error.  Please report it to the developers.

\sphinxlineitem{\sphinxstyleemphasis{Attempt to push beyond stack limit in fix box/relax}}
\sphinxAtStartPar
Internal LAMMPS error.  Please report it to the developers.

\sphinxlineitem{\sphinxstyleemphasis{Attempting to rescale a 0.0 temperature}}
\sphinxAtStartPar
Cannot rescale a temperature that is already 0.0.

\sphinxlineitem{\sphinxstyleemphasis{Attempting to insert more particles than available lattice points}}
\sphinxAtStartPar
Self\sphinxhyphen{}explanatory.

\sphinxlineitem{\sphinxstyleemphasis{Bad FENE bond}}
\sphinxAtStartPar
Two atoms in a FENE bond have become so far apart that the bond cannot
be computed.

\sphinxlineitem{\sphinxstyleemphasis{Bad TIP4P angle type for PPPM/TIP4P}}
\sphinxAtStartPar
Specified angle type is not valid.

\sphinxlineitem{\sphinxstyleemphasis{Bad TIP4P angle type for PPPMDisp/TIP4P}}
\sphinxAtStartPar
Specified angle type is not valid.

\sphinxlineitem{\sphinxstyleemphasis{Bad TIP4P bond type for PPPM/TIP4P}}
\sphinxAtStartPar
Specified bond type is not valid.

\sphinxlineitem{\sphinxstyleemphasis{Bad TIP4P bond type for PPPMDisp/TIP4P}}
\sphinxAtStartPar
Specified bond type is not valid.

\sphinxlineitem{\sphinxstyleemphasis{Bad fix ID in fix append/atoms command}}
\sphinxAtStartPar
The value of the fix\_id for keyword spatial must start with “f\_”.

\sphinxlineitem{\sphinxstyleemphasis{Bad grid of processors}}
\sphinxAtStartPar
The 3d grid of processors defined by the processors command does not
match the number of processors LAMMPS is being run on.

\sphinxlineitem{\sphinxstyleemphasis{Bad kspace\_modify kmax/ewald parameter}}
\sphinxAtStartPar
Kspace\_modify values for the kmax/ewald keyword must be integers \textgreater{} 0

\sphinxlineitem{\sphinxstyleemphasis{Bad kspace\_modify slab parameter}}
\sphinxAtStartPar
Kspace\_modify value for the slab/volume keyword must be \textgreater{}= 2.0.

\sphinxlineitem{\sphinxstyleemphasis{Bad matrix inversion in mldivide3}}
\sphinxAtStartPar
This error should not occur unless the matrix is badly formed.

\sphinxlineitem{\sphinxstyleemphasis{Bad principal moments}}
\sphinxAtStartPar
Fix rigid did not compute the principal moments of inertia of a rigid
group of atoms correctly.

\sphinxlineitem{\sphinxstyleemphasis{Bad quadratic solve for particle/line collision}}
\sphinxAtStartPar
This is an internal error.  It should normally not occur.

\sphinxlineitem{\sphinxstyleemphasis{Bad quadratic solve for particle/tri collision}}
\sphinxAtStartPar
This is an internal error.  It should normally not occur.

\sphinxlineitem{\sphinxstyleemphasis{Bad real space Coulombic cutoff in fix tune/kspace}}
\sphinxAtStartPar
Fix tune/kspace tried to find the optimal real space Coulombic cutoff using
the Newton\sphinxhyphen{}Rhaphson method, but found a non\sphinxhyphen{}positive or NaN cutoff

\sphinxlineitem{\sphinxstyleemphasis{Balance command before simulation box is defined}}
\sphinxAtStartPar
The balance command cannot be used before a read\_data, read\_restart,
or create\_box command.

\sphinxlineitem{\sphinxstyleemphasis{Balance produced bad splits}}
\sphinxAtStartPar
This should not occur.  It means two or more cutting plane locations
are on top of each other or out of order.  Report the problem to the
developers.

\sphinxlineitem{\sphinxstyleemphasis{Balance rcb cannot be used with comm\_style brick}}
\sphinxAtStartPar
Comm\_style tiled must be used instead.

\sphinxlineitem{\sphinxstyleemphasis{Balance shift string is invalid}}
\sphinxAtStartPar
The string can only contain the characters “x”, “y”, or “z”.

\sphinxlineitem{\sphinxstyleemphasis{Bias compute does not calculate a velocity bias}}
\sphinxAtStartPar
The specified compute must compute a bias for temperature.

\sphinxlineitem{\sphinxstyleemphasis{Bias compute does not calculate temperature}}
\sphinxAtStartPar
The specified compute must compute temperature.

\sphinxlineitem{\sphinxstyleemphasis{Bias compute group does not match compute group}}
\sphinxAtStartPar
The specified compute must operate on the same group as the parent
compute.

\sphinxlineitem{\sphinxstyleemphasis{Big particle in fix srd cannot be point particle}}
\sphinxAtStartPar
Big particles must be extended spheroids or ellipsoids.

\sphinxlineitem{\sphinxstyleemphasis{Bigint setting in lmptype.h is invalid}}
\sphinxAtStartPar
Size of bigint is less than size of tagint.

\sphinxlineitem{\sphinxstyleemphasis{Bigint setting in lmptype.h is not compatible}}
\sphinxAtStartPar
Format of bigint stored in restart file is not consistent with LAMMPS
version you are running.  See the settings in src/lmptype.h

\sphinxlineitem{\sphinxstyleemphasis{Bitmapped lookup tables require int/float be same size}}
\sphinxAtStartPar
Cannot use pair tables on this machine, because of word sizes.  Use
the pair\_modify command with table 0 instead.

\sphinxlineitem{\sphinxstyleemphasis{Bitmapped table in file does not match requested table}}
\sphinxAtStartPar
Setting for bitmapped table in pair\_coeff command must match table
in file exactly.

\sphinxlineitem{\sphinxstyleemphasis{Bitmapped table is incorrect length in table file}}
\sphinxAtStartPar
Number of table entries is not a correct power of 2.

\sphinxlineitem{\sphinxstyleemphasis{Bond and angle potentials must be defined for TIP4P}}
\sphinxAtStartPar
Cannot use TIP4P pair potential unless bond and angle potentials
are defined.

\sphinxlineitem{\sphinxstyleemphasis{Bond atom missing in box size check}}
\sphinxAtStartPar
The second atom needed to compute a particular bond is missing on this
processor.  Typically this is because the pairwise cutoff is set too
short or the bond has blown apart and an atom is too far away.

\sphinxlineitem{\sphinxstyleemphasis{Bond atom missing in delete\_bonds}}
\sphinxAtStartPar
The delete\_bonds command cannot find one or more atoms in a particular
bond on a particular processor.  The pairwise cutoff is too short or
the atoms are too far apart to make a valid bond.

\sphinxlineitem{\sphinxstyleemphasis{Bond atom missing in image check}}
\sphinxAtStartPar
The second atom in a particular bond is missing on this processor.
Typically this is because the pairwise cutoff is set too short or the
bond has blown apart and an atom is too far away.

\sphinxlineitem{\sphinxstyleemphasis{Bond atom missing in set command}}
\sphinxAtStartPar
The set command cannot find one or more atoms in a particular bond on
a particular processor.  The pairwise cutoff is too short or the atoms
are too far apart to make a valid bond.

\sphinxlineitem{\sphinxstyleemphasis{Bond atoms \%d \%d missing on proc \%d at step \%ld}}
\sphinxAtStartPar
The second atom needed to compute a particular bond is missing on this
processor.  Typically this is because the pairwise cutoff is set too
short or the bond has blown apart and an atom is too far away.

\sphinxlineitem{\sphinxstyleemphasis{Bond atoms missing on proc \%d at step \%ld}}
\sphinxAtStartPar
The second atom needed to compute a particular bond is missing on this
processor.  Typically this is because the pairwise cutoff is set too
short or the bond has blown apart and an atom is too far away.

\sphinxlineitem{\sphinxstyleemphasis{Bond coeff for hybrid has invalid style}}
\sphinxAtStartPar
Bond style hybrid uses another bond style as one of its coefficients.
The bond style used in the bond\_coeff command or read from a restart
file is not recognized.

\sphinxlineitem{\sphinxstyleemphasis{Bond coeffs are not set}}
\sphinxAtStartPar
No bond coefficients have been assigned in the data file or via the
bond\_coeff command.

\sphinxlineitem{\sphinxstyleemphasis{Bond extent \textgreater{} half of periodic box length}}
\sphinxAtStartPar
This error was detected by the neigh\_modify check yes setting.  It is
an error because the bond atoms are so far apart it is ambiguous how
it should be defined.

\sphinxlineitem{\sphinxstyleemphasis{Bond potential must be defined for SHAKE}}
\sphinxAtStartPar
Cannot use fix shake unless bond potential is defined.

\sphinxlineitem{\sphinxstyleemphasis{Bond style hybrid cannot have hybrid as an argument}}
\sphinxAtStartPar
Self\sphinxhyphen{}explanatory.

\sphinxlineitem{\sphinxstyleemphasis{Bond style hybrid cannot have none as an argument}}
\sphinxAtStartPar
Self\sphinxhyphen{}explanatory.

\sphinxlineitem{\sphinxstyleemphasis{Bond style hybrid cannot use same bond style twice}}
\sphinxAtStartPar
Self\sphinxhyphen{}explanatory.

\sphinxlineitem{\sphinxstyleemphasis{Bond style quartic cannot be used with 3,4\sphinxhyphen{}body interactions}}
\sphinxAtStartPar
No angle, dihedral, or improper styles can be defined when using
bond style quartic.

\sphinxlineitem{\sphinxstyleemphasis{Bond style quartic cannot be used with atom style template}}
\sphinxAtStartPar
This bond style can change the bond topology which is not
allowed with this atom style.

\sphinxlineitem{\sphinxstyleemphasis{Bond style quartic requires special\_bonds = 1,1,1}}
\sphinxAtStartPar
This is a restriction of the current bond quartic implementation.

\sphinxlineitem{\sphinxstyleemphasis{Bond table parameters did not set N}}
\sphinxAtStartPar
List of bond table parameters must include N setting.

\sphinxlineitem{\sphinxstyleemphasis{Bond table values are not increasing}}
\sphinxAtStartPar
The values in the tabulated file must be monotonically increasing.

\sphinxlineitem{\sphinxstyleemphasis{BondAngle coeff for hybrid angle has invalid format}}
\sphinxAtStartPar
No “ba” field should appear in data file entry.

\sphinxlineitem{\sphinxstyleemphasis{BondBond coeff for hybrid angle has invalid format}}
\sphinxAtStartPar
No “bb” field should appear in data file entry.

\sphinxlineitem{\sphinxstyleemphasis{Bond\_coeff command before bond\_style is defined}}
\sphinxAtStartPar
Coefficients cannot be set in the data file or via the bond\_coeff
command until an bond\_style has been assigned.

\sphinxlineitem{\sphinxstyleemphasis{Bond\_coeff command before simulation box is defined}}
\sphinxAtStartPar
The bond\_coeff command cannot be used before a read\_data,
read\_restart, or create\_box command.

\sphinxlineitem{\sphinxstyleemphasis{Bond\_coeff command when no bonds allowed}}
\sphinxAtStartPar
The chosen atom style does not allow for bonds to be defined.

\sphinxlineitem{\sphinxstyleemphasis{Bond\_style command when no bonds allowed}}
\sphinxAtStartPar
The chosen atom style does not allow for bonds to be defined.

\sphinxlineitem{\sphinxstyleemphasis{Bonds assigned incorrectly}}
\sphinxAtStartPar
Bonds read in from the data file were not assigned correctly to atoms.
This means there is something invalid about the topology definitions.

\sphinxlineitem{\sphinxstyleemphasis{Bonds defined but no bond types}}
\sphinxAtStartPar
The data file header lists bonds but no bond types.

\sphinxlineitem{\sphinxstyleemphasis{Both restart files must use \% or neither}}
\sphinxAtStartPar
Self\sphinxhyphen{}explanatory.

\sphinxlineitem{\sphinxstyleemphasis{Both restart files must use MPI\sphinxhyphen{}IO or neither}}
\sphinxAtStartPar
Self\sphinxhyphen{}explanatory.

\sphinxlineitem{\sphinxstyleemphasis{Both sides of boundary must be periodic}}
\sphinxAtStartPar
Cannot specify a boundary as periodic only on the lo or hi side.  Must
be periodic on both sides.

\sphinxlineitem{\sphinxstyleemphasis{Boundary command after simulation box is defined}}
\sphinxAtStartPar
The boundary command cannot be used after a read\_data, read\_restart,
or create\_box command.

\sphinxlineitem{\sphinxstyleemphasis{Box bounds are invalid}}
\sphinxAtStartPar
The box boundaries specified in the read\_data file are invalid.  The
lo value must be less than the hi value for all 3 dimensions.

\sphinxlineitem{\sphinxstyleemphasis{Box command after simulation box is defined}}
\sphinxAtStartPar
The box command cannot be used after a read\_data, read\_restart, or
create\_box command.

\sphinxlineitem{\sphinxstyleemphasis{CPU neighbor lists must be used for ellipsoid/sphere mix.}}
\sphinxAtStartPar
When using Gay\sphinxhyphen{}Berne or RE\sphinxhyphen{}squared pair styles with both ellipsoidal and
spherical particles, the neighbor list must be built on the CPU

\sphinxlineitem{\sphinxstyleemphasis{Can not specify Pxy/Pxz/Pyz in fix box/relax with non\sphinxhyphen{}triclinic box}}
\sphinxAtStartPar
Only triclinic boxes can be used with off\sphinxhyphen{}diagonal pressure components.
See the region prism command for details.

\sphinxlineitem{\sphinxstyleemphasis{Can not specify Pxy/Pxz/Pyz in fix nvt/npt/nph with non\sphinxhyphen{}triclinic box}}
\sphinxAtStartPar
Only triclinic boxes can be used with off\sphinxhyphen{}diagonal pressure components.
See the region prism command for details.

\sphinxlineitem{\sphinxstyleemphasis{Can only use \sphinxhyphen{}plog with multiple partitions}}
\sphinxAtStartPar
Self\sphinxhyphen{}explanatory.  See page discussion of command\sphinxhyphen{}line switches.

\sphinxlineitem{\sphinxstyleemphasis{Can only use \sphinxhyphen{}pscreen with multiple partitions}}
\sphinxAtStartPar
Self\sphinxhyphen{}explanatory.  See page discussion of command\sphinxhyphen{}line switches.

\sphinxlineitem{\sphinxstyleemphasis{Can only use Kokkos supported regions with Kokkos package}}
\sphinxAtStartPar
Self\sphinxhyphen{}explanatory.

\sphinxlineitem{\sphinxstyleemphasis{Can only use NEB with 1\sphinxhyphen{}processor replicas}}
\sphinxAtStartPar
This is current restriction for NEB as implemented in LAMMPS.

\sphinxlineitem{\sphinxstyleemphasis{Can only use TAD with 1\sphinxhyphen{}processor replicas for NEB}}
\sphinxAtStartPar
This is current restriction for NEB as implemented in LAMMPS.

\sphinxlineitem{\sphinxstyleemphasis{Cannot (yet) do analytic differentiation with pppm/gpu}}
\sphinxAtStartPar
This is a current restriction of this command.

\sphinxlineitem{\sphinxstyleemphasis{Cannot (yet) request ghost atoms with Kokkos half neighbor list}}
\sphinxAtStartPar
This feature is not yet supported.

\sphinxlineitem{\sphinxstyleemphasis{Cannot (yet) use ‘electron’ units with dipoles}}
\sphinxAtStartPar
This feature is not yet supported.

\sphinxlineitem{\sphinxstyleemphasis{Cannot (yet) use Ewald with triclinic box and slab correction}}
\sphinxAtStartPar
This feature is not yet supported.

\sphinxlineitem{\sphinxstyleemphasis{Cannot (yet) use K\sphinxhyphen{}space slab correction with compute group/group for triclinic systems}}
\sphinxAtStartPar
This option is not yet supported.

\sphinxlineitem{\sphinxstyleemphasis{Cannot (yet) use MSM with 2d simulation}}
\sphinxAtStartPar
This feature is not yet supported.

\sphinxlineitem{\sphinxstyleemphasis{Cannot (yet) use PPPM with triclinic box and TIP4P}}
\sphinxAtStartPar
This feature is not yet supported.

\sphinxlineitem{\sphinxstyleemphasis{Cannot (yet) use PPPM with triclinic box and kspace\_modify diff ad}}
\sphinxAtStartPar
This feature is not yet supported.

\sphinxlineitem{\sphinxstyleemphasis{Cannot (yet) use PPPM with triclinic box and slab correction}}
\sphinxAtStartPar
This feature is not yet supported.

\sphinxlineitem{\sphinxstyleemphasis{Cannot (yet) use kspace slab correction with long\sphinxhyphen{}range dipoles and non\sphinxhyphen{}neutral systems or per\sphinxhyphen{}atom energy}}
\sphinxAtStartPar
This feature is not yet supported.

\sphinxlineitem{\sphinxstyleemphasis{Cannot (yet) use kspace\_modify diff ad with compute group/group}}
\sphinxAtStartPar
This option is not yet supported.

\sphinxlineitem{\sphinxstyleemphasis{Cannot (yet) use kspace\_style pppm/stagger with triclinic systems}}
\sphinxAtStartPar
This feature is not yet supported.

\sphinxlineitem{\sphinxstyleemphasis{Cannot (yet) use molecular templates with Kokkos}}
\sphinxAtStartPar
Self\sphinxhyphen{}explanatory.

\sphinxlineitem{\sphinxstyleemphasis{Cannot (yet) use respa with Kokkos}}
\sphinxAtStartPar
Self\sphinxhyphen{}explanatory.

\sphinxlineitem{\sphinxstyleemphasis{Cannot (yet) use rigid bodies with fix deform and Kokkos}}
\sphinxAtStartPar
Self\sphinxhyphen{}explanatory.

\sphinxlineitem{\sphinxstyleemphasis{Cannot (yet) use rigid bodies with fix nh and Kokkos}}
\sphinxAtStartPar
Self\sphinxhyphen{}explanatory.

\sphinxlineitem{\sphinxstyleemphasis{Cannot (yet) use single precision with MSM (remove \sphinxhyphen{}DFFT\_SINGLE from Makefile and re\sphinxhyphen{}compile)}}
\sphinxAtStartPar
Single precision cannot be used with MSM.

\sphinxlineitem{\sphinxstyleemphasis{Cannot add atoms to fix move variable}}
\sphinxAtStartPar
Atoms can not be added afterwards to this fix option.

\sphinxlineitem{\sphinxstyleemphasis{Cannot append atoms to a triclinic box}}
\sphinxAtStartPar
The simulation box must be defined with edges aligned with the
Cartesian axes.

\sphinxlineitem{\sphinxstyleemphasis{Cannot balance in z dimension for 2d simulation}}
\sphinxAtStartPar
Self\sphinxhyphen{}explanatory.

\sphinxlineitem{\sphinxstyleemphasis{Cannot change box ortho/triclinic with certain fixes defined}}
\sphinxAtStartPar
This is because those fixes store the shape of the box.  You need to
use unfix to discard the fix, change the box, then redefine a new
fix.

\sphinxlineitem{\sphinxstyleemphasis{Cannot change box ortho/triclinic with dumps defined}}
\sphinxAtStartPar
This is because some dumps store the shape of the box.  You need to
use undump to discard the dump, change the box, then redefine a new
dump.

\sphinxlineitem{\sphinxstyleemphasis{Cannot change box tilt factors for orthogonal box}}
\sphinxAtStartPar
Cannot use tilt factors unless the simulation box is non\sphinxhyphen{}orthogonal.

\sphinxlineitem{\sphinxstyleemphasis{Cannot change box to orthogonal when tilt is non\sphinxhyphen{}zero}}
\sphinxAtStartPar
Self\sphinxhyphen{}explanatory.

\sphinxlineitem{\sphinxstyleemphasis{Cannot change box z boundary to non\sphinxhyphen{}periodic for a 2d simulation}}
\sphinxAtStartPar
Self\sphinxhyphen{}explanatory.

\sphinxlineitem{\sphinxstyleemphasis{Cannot change dump\_modify every for dump dcd}}
\sphinxAtStartPar
The frequency of writing dump dcd snapshots cannot be changed.

\sphinxlineitem{\sphinxstyleemphasis{Cannot change dump\_modify every for dump xtc}}
\sphinxAtStartPar
The frequency of writing dump xtc snapshots cannot be changed.

\sphinxlineitem{\sphinxstyleemphasis{Cannot change timestep once fix srd is setup}}
\sphinxAtStartPar
This is because various SRD properties depend on the timestep
size.

\sphinxlineitem{\sphinxstyleemphasis{Cannot change timestep with fix pour}}
\sphinxAtStartPar
This is because fix pour pre\sphinxhyphen{}computes the time delay for particles to
fall out of the insertion volume due to gravity.

\sphinxlineitem{\sphinxstyleemphasis{Cannot change to comm\_style brick from tiled layout}}
\sphinxAtStartPar
Self\sphinxhyphen{}explanatory.

\sphinxlineitem{\sphinxstyleemphasis{Cannot change\_box after reading restart file with per\sphinxhyphen{}atom info}}
\sphinxAtStartPar
This is because the restart file info cannot be migrated with the
atoms.  You can get around this by performing a 0\sphinxhyphen{}timestep run which
will assign the restart file info to actual atoms.

\sphinxlineitem{\sphinxstyleemphasis{Cannot change\_box in xz or yz for 2d simulation}}
\sphinxAtStartPar
Self\sphinxhyphen{}explanatory.

\sphinxlineitem{\sphinxstyleemphasis{Cannot change\_box in z dimension for 2d simulation}}
\sphinxAtStartPar
Self\sphinxhyphen{}explanatory.

\sphinxlineitem{\sphinxstyleemphasis{Cannot clear group all}}
\sphinxAtStartPar
This operation is not allowed.

\sphinxlineitem{\sphinxstyleemphasis{Cannot close restart file \sphinxhyphen{} MPI error: \%s}}
\sphinxAtStartPar
This error was generated by MPI when reading/writing an MPI\sphinxhyphen{}IO restart
file.

\sphinxlineitem{\sphinxstyleemphasis{Cannot compute initial g\_ewald\_disp}}
\sphinxAtStartPar
LAMMPS failed to compute an initial guess for the PPPM\_disp g\_ewald\_6
factor that partitions the computation between real space and k\sphinxhyphen{}space
for Dispersion interactions.

\sphinxlineitem{\sphinxstyleemphasis{Cannot create an atom map unless atoms have IDs}}
\sphinxAtStartPar
The simulation requires a mapping from global atom IDs to local atoms,
but the atoms that have been defined have no IDs.

\sphinxlineitem{\sphinxstyleemphasis{Cannot create atoms with undefined lattice}}
\sphinxAtStartPar
Must use the lattice command before using the create\_atoms
command.

\sphinxlineitem{\sphinxstyleemphasis{Cannot create/grow a vector/array of pointers for \%s}}
\sphinxAtStartPar
LAMMPS code is making an illegal call to the templated memory
allocators, to create a vector or array of pointers.

\sphinxlineitem{\sphinxstyleemphasis{Cannot create\_atoms after reading restart file with per\sphinxhyphen{}atom info}}
\sphinxAtStartPar
The per\sphinxhyphen{}atom info was stored to be used when by a fix that you may
re\sphinxhyphen{}define.  If you add atoms before re\sphinxhyphen{}defining the fix, then there
will not be a correct amount of per\sphinxhyphen{}atom info.

\sphinxlineitem{\sphinxstyleemphasis{Cannot create\_box after simulation box is defined}}
\sphinxAtStartPar
A simulation box can only be defined once.

\sphinxlineitem{\sphinxstyleemphasis{Cannot currently use pair reax with pair hybrid}}
\sphinxAtStartPar
This is not yet supported.

\sphinxlineitem{\sphinxstyleemphasis{Cannot currently use pppm/gpu with fix balance.}}
\sphinxAtStartPar
Self\sphinxhyphen{}explanatory.

\sphinxlineitem{\sphinxstyleemphasis{Cannot delete group all}}
\sphinxAtStartPar
Self\sphinxhyphen{}explanatory.

\sphinxlineitem{\sphinxstyleemphasis{Cannot delete group currently used by a compute}}
\sphinxAtStartPar
Self\sphinxhyphen{}explanatory.

\sphinxlineitem{\sphinxstyleemphasis{Cannot delete group currently used by a dump}}
\sphinxAtStartPar
Self\sphinxhyphen{}explanatory.

\sphinxlineitem{\sphinxstyleemphasis{Cannot delete group currently used by a fix}}
\sphinxAtStartPar
Self\sphinxhyphen{}explanatory.

\sphinxlineitem{\sphinxstyleemphasis{Cannot delete group currently used by atom\_modify first}}
\sphinxAtStartPar
Self\sphinxhyphen{}explanatory.

\sphinxlineitem{\sphinxstyleemphasis{Cannot delete\_atoms bond yes for non\sphinxhyphen{}molecular systems}}
\sphinxAtStartPar
Self\sphinxhyphen{}explanatory.

\sphinxlineitem{\sphinxstyleemphasis{Cannot displace\_atoms after reading restart file with per\sphinxhyphen{}atom info}}
\sphinxAtStartPar
This is because the restart file info cannot be migrated with the
atoms.  You can get around this by performing a 0\sphinxhyphen{}timestep run which
will assign the restart file info to actual atoms.

\sphinxlineitem{\sphinxstyleemphasis{Cannot do GCMC on atoms in atom\_modify first group}}
\sphinxAtStartPar
This is a restriction due to the way atoms are organized in a list to
enable the atom\_modify first command.

\sphinxlineitem{\sphinxstyleemphasis{Cannot do atom/swap on atoms in atom\_modify first group}}
\sphinxAtStartPar
This is a restriction due to the way atoms are organized in a list to
enable the atom\_modify first command.

\sphinxlineitem{\sphinxstyleemphasis{Cannot dump sort on atom IDs with no atom IDs defined}}
\sphinxAtStartPar
Self\sphinxhyphen{}explanatory.

\sphinxlineitem{\sphinxstyleemphasis{Cannot dump sort when multiple dump files are written}}
\sphinxAtStartPar
In this mode, each processor dumps its atoms to a file, so
no sorting is allowed.

\sphinxlineitem{\sphinxstyleemphasis{Cannot embed Python when also extending Python with LAMMPS}}
\sphinxAtStartPar
When running LAMMPS via Python through the LAMMPS library interface
you cannot also user the input script python command.

\sphinxlineitem{\sphinxstyleemphasis{Cannot evaporate atoms in atom\_modify first group}}
\sphinxAtStartPar
This is a restriction due to the way atoms are organized in
a list to enable the atom\_modify first command.

\sphinxlineitem{\sphinxstyleemphasis{Cannot find create\_bonds group ID}}
\sphinxAtStartPar
Self\sphinxhyphen{}explanatory.

\sphinxlineitem{\sphinxstyleemphasis{Cannot find delete\_bonds group ID}}
\sphinxAtStartPar
Group ID used in the delete\_bonds command does not exist.

\sphinxlineitem{\sphinxstyleemphasis{Cannot find specified group ID for core particles}}
\sphinxAtStartPar
Self\sphinxhyphen{}explanatory.

\sphinxlineitem{\sphinxstyleemphasis{Cannot find specified group ID for shell particles}}
\sphinxAtStartPar
Self\sphinxhyphen{}explanatory.

\sphinxlineitem{\sphinxstyleemphasis{Cannot have both pair\_modify shift and tail set to yes}}
\sphinxAtStartPar
These 2 options are contradictory.

\sphinxlineitem{\sphinxstyleemphasis{Cannot intersect groups using a dynamic group}}
\sphinxAtStartPar
This operation is not allowed.

\sphinxlineitem{\sphinxstyleemphasis{Cannot mix molecular and molecule template atom styles}}
\sphinxAtStartPar
Self\sphinxhyphen{}explanatory.

\sphinxlineitem{\sphinxstyleemphasis{Cannot open \sphinxhyphen{}reorder file}}
\sphinxAtStartPar
Self\sphinxhyphen{}explanatory.

\sphinxlineitem{\sphinxstyleemphasis{Cannot open ADP potential file \%s}}
\sphinxAtStartPar
The specified ADP potential file cannot be opened.  Check that the
path and name are correct.

\sphinxlineitem{\sphinxstyleemphasis{Cannot open AIREBO potential file \%s}}
\sphinxAtStartPar
The specified AIREBO potential file cannot be opened.  Check that the
path and name are correct.

\sphinxlineitem{\sphinxstyleemphasis{Cannot open BOP potential file \%s}}
\sphinxAtStartPar
The specified BOP potential file cannot be opened.  Check that the
path and name are correct.

\sphinxlineitem{\sphinxstyleemphasis{Cannot open COMB potential file \%s}}
\sphinxAtStartPar
The specified COMB potential file cannot be opened.  Check that the
path and name are correct.

\sphinxlineitem{\sphinxstyleemphasis{Cannot open COMB3 lib.comb3 file}}
\sphinxAtStartPar
The COMB3 library file cannot be opened.  Check that the path and name
are correct.

\sphinxlineitem{\sphinxstyleemphasis{Cannot open COMB3 potential file \%s}}
\sphinxAtStartPar
The specified COMB3 potential file cannot be opened.  Check that the
path and name are correct.

\sphinxlineitem{\sphinxstyleemphasis{Cannot open EAM potential file \%s}}
\sphinxAtStartPar
The specified EAM potential file cannot be opened.  Check that the
path and name are correct.

\sphinxlineitem{\sphinxstyleemphasis{Cannot open EIM potential file \%s}}
\sphinxAtStartPar
The specified EIM potential file cannot be opened.  Check that the
path and name are correct.

\sphinxlineitem{\sphinxstyleemphasis{Cannot open LCBOP potential file \%s}}
\sphinxAtStartPar
The specified LCBOP potential file cannot be opened.  Check that the
path and name are correct.

\sphinxlineitem{\sphinxstyleemphasis{Cannot open MEAM potential file \%s}}
\sphinxAtStartPar
The specified MEAM potential file cannot be opened.  Check that the
path and name are correct.

\sphinxlineitem{\sphinxstyleemphasis{Cannot open SNAP coefficient file \%s}}
\sphinxAtStartPar
The specified SNAP coefficient file cannot be opened.  Check that the
path and name are correct.

\sphinxlineitem{\sphinxstyleemphasis{Cannot open SNAP parameter file \%s}}
\sphinxAtStartPar
The specified SNAP parameter file cannot be opened.  Check that the
path and name are correct.

\sphinxlineitem{\sphinxstyleemphasis{Cannot open Stillinger\sphinxhyphen{}Weber potential file \%s}}
\sphinxAtStartPar
The specified SW potential file cannot be opened.  Check that the path
and name are correct.

\sphinxlineitem{\sphinxstyleemphasis{Cannot open Tersoff potential file \%s}}
\sphinxAtStartPar
The specified potential file cannot be opened.  Check that the path
and name are correct.

\sphinxlineitem{\sphinxstyleemphasis{Cannot open Vashishta potential file \%s}}
\sphinxAtStartPar
The specified Vashishta potential file cannot be opened.  Check that the path
and name are correct.

\sphinxlineitem{\sphinxstyleemphasis{Cannot open balance output file}}
\sphinxAtStartPar
Self\sphinxhyphen{}explanatory.

\sphinxlineitem{\sphinxstyleemphasis{Cannot open coul/streitz potential file \%s}}
\sphinxAtStartPar
The specified coul/streitz potential file cannot be opened.  Check
that the path and name are correct.

\sphinxlineitem{\sphinxstyleemphasis{Cannot open custom file}}
\sphinxAtStartPar
Self\sphinxhyphen{}explanatory.

\sphinxlineitem{\sphinxstyleemphasis{Cannot open data file \%s}}
\sphinxAtStartPar
The specified file cannot be opened.  Check that the path and name are
correct.

\sphinxlineitem{\sphinxstyleemphasis{Cannot open dir to search for restart file}}
\sphinxAtStartPar
Using a “*” in the name of the restart file will open the current
directory to search for matching file names.

\sphinxlineitem{\sphinxstyleemphasis{Cannot open dump file}}
\sphinxAtStartPar
Self\sphinxhyphen{}explanatory.

\sphinxlineitem{\sphinxstyleemphasis{Cannot open dump file \%s}}
\sphinxAtStartPar
The output file for the dump command cannot be opened.  Check that the
path and name are correct.

\sphinxlineitem{\sphinxstyleemphasis{Cannot open file \%s}}
\sphinxAtStartPar
The specified file cannot be opened.  Check that the path and name are
correct. If the file is a compressed file, also check that the gzip
executable can be found and run.

\sphinxlineitem{\sphinxstyleemphasis{Cannot open file variable file \%s}}
\sphinxAtStartPar
The specified file cannot be opened.  Check that the path and name are
correct.

\sphinxlineitem{\sphinxstyleemphasis{Cannot open fix ave/chunk file \%s}}
\sphinxAtStartPar
The specified file cannot be opened.  Check that the path and name are
correct.

\sphinxlineitem{\sphinxstyleemphasis{Cannot open fix ave/correlate file \%s}}
\sphinxAtStartPar
The specified file cannot be opened.  Check that the path and name are
correct.

\sphinxlineitem{\sphinxstyleemphasis{Cannot open fix ave/histo file \%s}}
\sphinxAtStartPar
The specified file cannot be opened.  Check that the path and name are
correct.

\sphinxlineitem{\sphinxstyleemphasis{Cannot open fix ave/time file \%s}}
\sphinxAtStartPar
The specified file cannot be opened.  Check that the path and name are
correct.

\sphinxlineitem{\sphinxstyleemphasis{Cannot open fix balance output file}}
\sphinxAtStartPar
Self\sphinxhyphen{}explanatory.

\sphinxlineitem{\sphinxstyleemphasis{Cannot open fix poems file \%s}}
\sphinxAtStartPar
The specified file cannot be opened.  Check that the path and name are
correct.

\sphinxlineitem{\sphinxstyleemphasis{Cannot open fix print file \%s}}
\sphinxAtStartPar
The output file generated by the fix print command cannot be opened

\sphinxlineitem{\sphinxstyleemphasis{Cannot open fix qeq parameter file \%s}}
\sphinxAtStartPar
The specified file cannot be opened.  Check that the path and name are
correct.

\sphinxlineitem{\sphinxstyleemphasis{Cannot open fix qeq/comb file \%s}}
\sphinxAtStartPar
The output file for the fix qeq/combs command cannot be opened.
Check that the path and name are correct.

\sphinxlineitem{\sphinxstyleemphasis{Cannot open fix reax/bonds file \%s}}
\sphinxAtStartPar
The output file for the fix reax/bonds command cannot be opened.
Check that the path and name are correct.

\sphinxlineitem{\sphinxstyleemphasis{Cannot open fix rigid infile \%s}}
\sphinxAtStartPar
The specified file cannot be opened.  Check that the path and name are
correct.

\sphinxlineitem{\sphinxstyleemphasis{Cannot open fix rigid restart file \%s}}
\sphinxAtStartPar
The specified file cannot be opened.  Check that the path and name are
correct.

\sphinxlineitem{\sphinxstyleemphasis{Cannot open fix rigid/small infile \%s}}
\sphinxAtStartPar
The specified file cannot be opened.  Check that the path and name are
correct.

\sphinxlineitem{\sphinxstyleemphasis{Cannot open fix tmd file \%s}}
\sphinxAtStartPar
The output file for the fix tmd command cannot be opened.  Check that
the path and name are correct.

\sphinxlineitem{\sphinxstyleemphasis{Cannot open fix ttm file \%s}}
\sphinxAtStartPar
The output file for the fix ttm command cannot be opened.  Check that
the path and name are correct.

\sphinxlineitem{\sphinxstyleemphasis{Cannot open gzipped file}}
\sphinxAtStartPar
LAMMPS was compiled without support for reading and writing gzipped
files through a pipeline to the gzip program with \sphinxhyphen{}DLAMMPS\_GZIP.

\sphinxlineitem{\sphinxstyleemphasis{Cannot open input script \%s}}
\sphinxAtStartPar
Self\sphinxhyphen{}explanatory.

\sphinxlineitem{\sphinxstyleemphasis{Cannot open log.cite file}}
\sphinxAtStartPar
This file is created when you use some LAMMPS features, to indicate
what paper you should cite on behalf of those who implemented
the feature.  Check that you have write privileges into the directory
you are running in.

\sphinxlineitem{\sphinxstyleemphasis{Cannot open log.lammps for writing}}
\sphinxAtStartPar
The default LAMMPS log file cannot be opened.  Check that the
directory you are running in allows for files to be created.

\sphinxlineitem{\sphinxstyleemphasis{Cannot open logfile}}
\sphinxAtStartPar
The LAMMPS log file named in a command\sphinxhyphen{}line argument cannot be opened.
Check that the path and name are correct.

\sphinxlineitem{\sphinxstyleemphasis{Cannot open logfile \%s}}
\sphinxAtStartPar
The LAMMPS log file specified in the input script cannot be opened.
Check that the path and name are correct.

\sphinxlineitem{\sphinxstyleemphasis{Cannot open molecule file \%s}}
\sphinxAtStartPar
The specified file cannot be opened.  Check that the path and name are
correct.

\sphinxlineitem{\sphinxstyleemphasis{Cannot open nb3b/harmonic potential file \%s}}
\sphinxAtStartPar
The specified potential file cannot be opened.  Check that the path
and name are correct.

\sphinxlineitem{\sphinxstyleemphasis{Cannot open pair\_write file}}
\sphinxAtStartPar
The specified output file for pair energies and forces cannot be
opened.  Check that the path and name are correct.

\sphinxlineitem{\sphinxstyleemphasis{Cannot open polymorphic potential file \%s}}
\sphinxAtStartPar
The specified polymorphic potential file cannot be opened.  Check that
the path and name are correct.

\sphinxlineitem{\sphinxstyleemphasis{Cannot open print file \%s}}
\sphinxAtStartPar
Self\sphinxhyphen{}explanatory.

\sphinxlineitem{\sphinxstyleemphasis{Cannot open processors output file}}
\sphinxAtStartPar
Self\sphinxhyphen{}explanatory.

\sphinxlineitem{\sphinxstyleemphasis{Cannot open restart file \%s}}
\sphinxAtStartPar
Self\sphinxhyphen{}explanatory.

\sphinxlineitem{\sphinxstyleemphasis{Cannot open restart file for reading \sphinxhyphen{} MPI error: \%s}}
\sphinxAtStartPar
This error was generated by MPI when reading/writing an MPI\sphinxhyphen{}IO restart
file.

\sphinxlineitem{\sphinxstyleemphasis{Cannot open restart file for writing \sphinxhyphen{} MPI error: \%s}}
\sphinxAtStartPar
This error was generated by MPI when reading/writing an MPI\sphinxhyphen{}IO restart
file.

\sphinxlineitem{\sphinxstyleemphasis{Cannot open screen file}}
\sphinxAtStartPar
The screen file specified as a command\sphinxhyphen{}line argument cannot be
opened.  Check that the directory you are running in allows for files
to be created.

\sphinxlineitem{\sphinxstyleemphasis{Cannot open temporary file for world counter.}}
\sphinxAtStartPar
Self\sphinxhyphen{}explanatory.

\sphinxlineitem{\sphinxstyleemphasis{Cannot open universe log file}}
\sphinxAtStartPar
For a multi\sphinxhyphen{}partition run, the master log file cannot be opened.
Check that the directory you are running in allows for files to be
created.

\sphinxlineitem{\sphinxstyleemphasis{Cannot open universe screen file}}
\sphinxAtStartPar
For a multi\sphinxhyphen{}partition run, the master screen file cannot be opened.
Check that the directory you are running in allows for files to be
created.

\sphinxlineitem{\sphinxstyleemphasis{Cannot read from restart file \sphinxhyphen{} MPI error: \%s}}
\sphinxAtStartPar
This error was generated by MPI when reading/writing an MPI\sphinxhyphen{}IO restart
file.

\sphinxlineitem{\sphinxstyleemphasis{Cannot read\_data without add keyword after simulation box is defined}}
\sphinxAtStartPar
Self\sphinxhyphen{}explanatory.

\sphinxlineitem{\sphinxstyleemphasis{Cannot read\_restart after simulation box is defined}}
\sphinxAtStartPar
The read\_restart command cannot be used after a read\_data,
read\_restart, or create\_box command.

\sphinxlineitem{\sphinxstyleemphasis{Cannot redefine variable as a different style}}
\sphinxAtStartPar
An equal\sphinxhyphen{}style variable can be re\sphinxhyphen{}defined but only if it was
originally an equal\sphinxhyphen{}style variable.

\sphinxlineitem{\sphinxstyleemphasis{Cannot replicate 2d simulation in z dimension}}
\sphinxAtStartPar
The replicate command cannot replicate a 2d simulation in the z
dimension.

\sphinxlineitem{\sphinxstyleemphasis{Cannot replicate with fixes that store atom quantities}}
\sphinxAtStartPar
Either fixes are defined that create and store atom\sphinxhyphen{}based vectors or a
restart file was read which included atom\sphinxhyphen{}based vectors for fixes.
The replicate command cannot duplicate that information for new atoms.
You should use the replicate command before fixes are applied to the
system.

\sphinxlineitem{\sphinxstyleemphasis{Cannot reset timestep with a dynamic region defined}}
\sphinxAtStartPar
Dynamic regions (see the region command) have a time dependence.
Thus you cannot change the timestep when one or more of these
are defined.

\sphinxlineitem{\sphinxstyleemphasis{Cannot reset timestep with a time\sphinxhyphen{}dependent fix defined}}
\sphinxAtStartPar
You cannot reset the timestep when a fix that keeps track of elapsed
time is in place.

\sphinxlineitem{\sphinxstyleemphasis{Cannot run 2d simulation with non\sphinxhyphen{}periodic Z dimension}}
\sphinxAtStartPar
Use the boundary command to make the z dimension periodic in order to
run a 2d simulation.

\sphinxlineitem{\sphinxstyleemphasis{Cannot set bond topology types for atom style template}}
\sphinxAtStartPar
The bond, angle, etc types cannot be changed for this atom style since
they are static settings in the molecule template files.

\sphinxlineitem{\sphinxstyleemphasis{Cannot set both respa pair and inner/middle/outer}}
\sphinxAtStartPar
In the rRESPA integrator, you must compute pairwise potentials either
all together (pair), or in pieces (inner/middle/outer).  You can’t do
both.

\sphinxlineitem{\sphinxstyleemphasis{Cannot set cutoff/multi before simulation box is defined}}
\sphinxAtStartPar
Self\sphinxhyphen{}explanatory.

\sphinxlineitem{\sphinxstyleemphasis{Cannot set dpd/theta for this atom style}}
\sphinxAtStartPar
Self\sphinxhyphen{}explanatory.

\sphinxlineitem{\sphinxstyleemphasis{Cannot set dump\_modify flush for dump xtc}}
\sphinxAtStartPar
Self\sphinxhyphen{}explanatory.

\sphinxlineitem{\sphinxstyleemphasis{Cannot set mass for this atom style}}
\sphinxAtStartPar
This atom style does not support mass settings for each atom type.
Instead they are defined on a per\sphinxhyphen{}atom basis in the data file.

\sphinxlineitem{\sphinxstyleemphasis{Cannot set meso/cv for this atom style}}
\sphinxAtStartPar
Self\sphinxhyphen{}explanatory.

\sphinxlineitem{\sphinxstyleemphasis{Cannot set meso/e for this atom style}}
\sphinxAtStartPar
Self\sphinxhyphen{}explanatory.

\sphinxlineitem{\sphinxstyleemphasis{Cannot set meso/rho for this atom style}}
\sphinxAtStartPar
Self\sphinxhyphen{}explanatory.

\sphinxlineitem{\sphinxstyleemphasis{Cannot set non\sphinxhyphen{}zero image flag for non\sphinxhyphen{}periodic dimension}}
\sphinxAtStartPar
Self\sphinxhyphen{}explanatory.

\sphinxlineitem{\sphinxstyleemphasis{Cannot set non\sphinxhyphen{}zero z velocity for 2d simulation}}
\sphinxAtStartPar
Self\sphinxhyphen{}explanatory.

\sphinxlineitem{\sphinxstyleemphasis{Cannot set quaternion for atom that has none}}
\sphinxAtStartPar
Self\sphinxhyphen{}explanatory.

\sphinxlineitem{\sphinxstyleemphasis{Cannot set quaternion with xy components for 2d system}}
\sphinxAtStartPar
Self\sphinxhyphen{}explanatory.

\sphinxlineitem{\sphinxstyleemphasis{Cannot set respa hybrid and any of pair/inner/middle/outer}}
\sphinxAtStartPar
In the rRESPA integrator, you must compute pairwise potentials either
all together (pair), with different cutoff regions (inner/middle/outer),
or per hybrid sub\sphinxhyphen{}style (hybrid).  You cannot mix those.

\sphinxlineitem{\sphinxstyleemphasis{Cannot set respa middle without inner/outer}}
\sphinxAtStartPar
In the rRESPA integrator, you must define both a inner and outer
setting in order to use a middle setting.

\sphinxlineitem{\sphinxstyleemphasis{Cannot set restart file size \sphinxhyphen{} MPI error: \%s}}
\sphinxAtStartPar
This error was generated by MPI when reading/writing an MPI\sphinxhyphen{}IO restart
file.

\sphinxlineitem{\sphinxstyleemphasis{Cannot set smd/contact/radius for this atom style}}
\sphinxAtStartPar
Self\sphinxhyphen{}explanatory.

\sphinxlineitem{\sphinxstyleemphasis{Cannot set smd/mass/density for this atom style}}
\sphinxAtStartPar
Self\sphinxhyphen{}explanatory.

\sphinxlineitem{\sphinxstyleemphasis{Cannot set temperature for fix rigid/nph}}
\sphinxAtStartPar
The temp keyword cannot be specified.

\sphinxlineitem{\sphinxstyleemphasis{Cannot set theta for atom that is not a line}}
\sphinxAtStartPar
Self\sphinxhyphen{}explanatory.

\sphinxlineitem{\sphinxstyleemphasis{Cannot set this attribute for this atom style}}
\sphinxAtStartPar
The attribute being set does not exist for the defined atom style.

\sphinxlineitem{\sphinxstyleemphasis{Cannot set variable z velocity for 2d simulation}}
\sphinxAtStartPar
Self\sphinxhyphen{}explanatory.

\sphinxlineitem{\sphinxstyleemphasis{Cannot skew triclinic box in z for 2d simulation}}
\sphinxAtStartPar
Self\sphinxhyphen{}explanatory.

\sphinxlineitem{\sphinxstyleemphasis{Cannot subtract groups using a dynamic group}}
\sphinxAtStartPar
This operation is not allowed.

\sphinxlineitem{\sphinxstyleemphasis{Cannot union groups using a dynamic group}}
\sphinxAtStartPar
This operation is not allowed.

\sphinxlineitem{\sphinxstyleemphasis{Cannot use \sphinxhyphen{}kokkos on without KOKKOS installed}}
\sphinxAtStartPar
Self\sphinxhyphen{}explanatory.

\sphinxlineitem{\sphinxstyleemphasis{Cannot use \sphinxhyphen{}reorder after \sphinxhyphen{}partition}}
\sphinxAtStartPar
Self\sphinxhyphen{}explanatory.  See page discussion of command\sphinxhyphen{}line switches.

\sphinxlineitem{\sphinxstyleemphasis{Cannot use Ewald with 2d simulation}}
\sphinxAtStartPar
The kspace style ewald cannot be used in 2d simulations.  You can use
2d Ewald in a 3d simulation; see the kspace\_modify command.

\sphinxlineitem{\sphinxstyleemphasis{Cannot use Ewald/disp solver on system with no charge, dipole, or LJ particles}}
\sphinxAtStartPar
No atoms in system have a non\sphinxhyphen{}zero charge or dipole, or are LJ
particles.  Change charges/dipoles or change options of the kspace
solver/pair style.

\sphinxlineitem{\sphinxstyleemphasis{Cannot use EwaldDisp with 2d simulation}}
\sphinxAtStartPar
This is a current restriction of this command.

\sphinxlineitem{\sphinxstyleemphasis{Cannot use Kokkos pair style with rRESPA inner/middle}}
\sphinxAtStartPar
Self\sphinxhyphen{}explanatory.

\sphinxlineitem{\sphinxstyleemphasis{Cannot use NEB unless atom map exists}}
\sphinxAtStartPar
Use the atom\_modify command to create an atom map.

\sphinxlineitem{\sphinxstyleemphasis{Cannot use NEB with a single replica}}
\sphinxAtStartPar
Self\sphinxhyphen{}explanatory.

\sphinxlineitem{\sphinxstyleemphasis{Cannot use NEB with atom\_modify sort enabled}}
\sphinxAtStartPar
This is current restriction for NEB implemented in LAMMPS.

\sphinxlineitem{\sphinxstyleemphasis{Cannot use PPPM with 2d simulation}}
\sphinxAtStartPar
The kspace style pppm cannot be used in 2d simulations.  You can use
2d PPPM in a 3d simulation; see the kspace\_modify command.

\sphinxlineitem{\sphinxstyleemphasis{Cannot use PPPMDisp with 2d simulation}}
\sphinxAtStartPar
The kspace style pppm/disp cannot be used in 2d simulations.  You can
use 2d pppm/disp in a 3d simulation; see the kspace\_modify command.

\sphinxlineitem{\sphinxstyleemphasis{Cannot use PRD with a changing box}}
\sphinxAtStartPar
The current box dimensions are not copied between replicas

\sphinxlineitem{\sphinxstyleemphasis{Cannot use PRD with a time\sphinxhyphen{}dependent fix defined}}
\sphinxAtStartPar
PRD alters the timestep in ways that will mess up these fixes.

\sphinxlineitem{\sphinxstyleemphasis{Cannot use PRD with a time\sphinxhyphen{}dependent region defined}}
\sphinxAtStartPar
PRD alters the timestep in ways that will mess up these regions.

\sphinxlineitem{\sphinxstyleemphasis{Cannot use PRD with atom\_modify sort enabled}}
\sphinxAtStartPar
This is a current restriction of PRD.  You must turn off sorting,
which is enabled by default, via the atom\_modify command.

\sphinxlineitem{\sphinxstyleemphasis{Cannot use PRD with multi\sphinxhyphen{}processor replicas unless atom map exists}}
\sphinxAtStartPar
Use the atom\_modify command to create an atom map.

\sphinxlineitem{\sphinxstyleemphasis{Cannot use TAD unless atom map exists for NEB}}
\sphinxAtStartPar
See atom\_modify map command to set this.

\sphinxlineitem{\sphinxstyleemphasis{Cannot use TAD with a single replica for NEB}}
\sphinxAtStartPar
NEB requires multiple replicas.

\sphinxlineitem{\sphinxstyleemphasis{Cannot use TAD with atom\_modify sort enabled for NEB}}
\sphinxAtStartPar
This is a current restriction of NEB.

\sphinxlineitem{\sphinxstyleemphasis{Cannot use a damped dynamics min style with fix box/relax}}
\sphinxAtStartPar
This is a current restriction in LAMMPS.  Use another minimizer
style.

\sphinxlineitem{\sphinxstyleemphasis{Cannot use a damped dynamics min style with per\sphinxhyphen{}atom DOF}}
\sphinxAtStartPar
This is a current restriction in LAMMPS.  Use another minimizer
style.

\sphinxlineitem{\sphinxstyleemphasis{Cannot use append/atoms in periodic dimension}}
\sphinxAtStartPar
The boundary style of the face where atoms are added can not be of
type p (periodic).

\sphinxlineitem{\sphinxstyleemphasis{Cannot use atomfile\sphinxhyphen{}style variable unless atom map exists}}
\sphinxAtStartPar
Self\sphinxhyphen{}explanatory.  See the atom\_modify command to create a map.

\sphinxlineitem{\sphinxstyleemphasis{Cannot use both com and bias with compute temp/chunk}}
\sphinxAtStartPar
Self\sphinxhyphen{}explanatory.

\sphinxlineitem{\sphinxstyleemphasis{Cannot use chosen neighbor list style with buck/coul/cut/kk}}
\sphinxAtStartPar
Self\sphinxhyphen{}explanatory.

\sphinxlineitem{\sphinxstyleemphasis{Cannot use chosen neighbor list style with buck/coul/long/kk}}
\sphinxAtStartPar
Self\sphinxhyphen{}explanatory.

\sphinxlineitem{\sphinxstyleemphasis{Cannot use chosen neighbor list style with buck/kk}}
\sphinxAtStartPar
That style is not supported by Kokkos.

\sphinxlineitem{\sphinxstyleemphasis{Cannot use chosen neighbor list style with coul/cut/kk}}
\sphinxAtStartPar
That style is not supported by Kokkos.

\sphinxlineitem{\sphinxstyleemphasis{Cannot use chosen neighbor list style with coul/debye/kk}}
\sphinxAtStartPar
Self\sphinxhyphen{}explanatory.

\sphinxlineitem{\sphinxstyleemphasis{Cannot use chosen neighbor list style with coul/dsf/kk}}
\sphinxAtStartPar
That style is not supported by Kokkos.

\sphinxlineitem{\sphinxstyleemphasis{Cannot use chosen neighbor list style with coul/wolf/kk}}
\sphinxAtStartPar
That style is not supported by Kokkos.

\sphinxlineitem{\sphinxstyleemphasis{Cannot use chosen neighbor list style with lj/charmm/coul/charmm/implicit/kk}}
\sphinxAtStartPar
Self\sphinxhyphen{}explanatory.

\sphinxlineitem{\sphinxstyleemphasis{Cannot use chosen neighbor list style with lj/charmm/coul/charmm/kk}}
\sphinxAtStartPar
Self\sphinxhyphen{}explanatory.

\sphinxlineitem{\sphinxstyleemphasis{Cannot use chosen neighbor list style with lj/charmm/coul/long/kk}}
\sphinxAtStartPar
Self\sphinxhyphen{}explanatory.

\sphinxlineitem{\sphinxstyleemphasis{Cannot use chosen neighbor list style with lj/class2/coul/cut/kk}}
\sphinxAtStartPar
Self\sphinxhyphen{}explanatory.

\sphinxlineitem{\sphinxstyleemphasis{Cannot use chosen neighbor list style with lj/class2/coul/long/kk}}
\sphinxAtStartPar
Self\sphinxhyphen{}explanatory.

\sphinxlineitem{\sphinxstyleemphasis{Cannot use chosen neighbor list style with lj/class2/kk}}
\sphinxAtStartPar
Self\sphinxhyphen{}explanatory.

\sphinxlineitem{\sphinxstyleemphasis{Cannot use chosen neighbor list style with lj/cut/coul/cut/kk}}
\sphinxAtStartPar
That style is not supported by Kokkos.

\sphinxlineitem{\sphinxstyleemphasis{Cannot use chosen neighbor list style with lj/cut/coul/debye/kk}}
\sphinxAtStartPar
Self\sphinxhyphen{}explanatory.

\sphinxlineitem{\sphinxstyleemphasis{Cannot use chosen neighbor list style with lj/cut/coul/long/kk}}
\sphinxAtStartPar
That style is not supported by Kokkos.

\sphinxlineitem{\sphinxstyleemphasis{Cannot use chosen neighbor list style with lj/cut/kk}}
\sphinxAtStartPar
That style is not supported by Kokkos.

\sphinxlineitem{\sphinxstyleemphasis{Cannot use chosen neighbor list style with lj/expand/kk}}
\sphinxAtStartPar
Self\sphinxhyphen{}explanatory.

\sphinxlineitem{\sphinxstyleemphasis{Cannot use chosen neighbor list style with lj/gromacs/coul/gromacs/kk}}
\sphinxAtStartPar
Self\sphinxhyphen{}explanatory.

\sphinxlineitem{\sphinxstyleemphasis{Cannot use chosen neighbor list style with lj/gromacs/kk}}
\sphinxAtStartPar
Self\sphinxhyphen{}explanatory.

\sphinxlineitem{\sphinxstyleemphasis{Cannot use chosen neighbor list style with lj/spica/kk}}
\sphinxAtStartPar
That style is not supported by Kokkos.

\sphinxlineitem{\sphinxstyleemphasis{Cannot use chosen neighbor list style with pair eam/kk}}
\sphinxAtStartPar
That style is not supported by Kokkos.

\sphinxlineitem{\sphinxstyleemphasis{Cannot use chosen neighbor list style with pair eam/kk/alloy}}
\sphinxAtStartPar
Self\sphinxhyphen{}explanatory.

\sphinxlineitem{\sphinxstyleemphasis{Cannot use chosen neighbor list style with pair eam/kk/fs}}
\sphinxAtStartPar
Self\sphinxhyphen{}explanatory.

\sphinxlineitem{\sphinxstyleemphasis{Cannot use chosen neighbor list style with pair sw/kk}}
\sphinxAtStartPar
Self\sphinxhyphen{}explanatory.

\sphinxlineitem{\sphinxstyleemphasis{Cannot use chosen neighbor list style with tersoff/kk}}
\sphinxAtStartPar
Self\sphinxhyphen{}explanatory.

\sphinxlineitem{\sphinxstyleemphasis{Cannot use chosen neighbor list style with tersoff/zbl/kk}}
\sphinxAtStartPar
Self\sphinxhyphen{}explanatory.

\sphinxlineitem{\sphinxstyleemphasis{Cannot use compute chunk/atom bin z for 2d model}}
\sphinxAtStartPar
Self\sphinxhyphen{}explanatory.

\sphinxlineitem{\sphinxstyleemphasis{Cannot use compute cluster/atom unless atoms have IDs}}
\sphinxAtStartPar
Atom IDs are used to identify clusters.

\sphinxlineitem{\sphinxstyleemphasis{Cannot use create\_atoms rotate unless single style}}
\sphinxAtStartPar
Self\sphinxhyphen{}explanatory.

\sphinxlineitem{\sphinxstyleemphasis{Cannot use create\_bonds unless atoms have IDs}}
\sphinxAtStartPar
This command requires a mapping from global atom IDs to local atoms,
but the atoms that have been defined have no IDs.

\sphinxlineitem{\sphinxstyleemphasis{Cannot use create\_bonds with non\sphinxhyphen{}molecular system}}
\sphinxAtStartPar
Self\sphinxhyphen{}explanatory.

\sphinxlineitem{\sphinxstyleemphasis{Cannot use cwiggle in variable formula between runs}}
\sphinxAtStartPar
This is a function of elapsed time.

\sphinxlineitem{\sphinxstyleemphasis{Cannot use delete\_atoms bond yes with atom\_style template}}
\sphinxAtStartPar
This is because the bonds for that atom style are hardwired in the
molecule template.

\sphinxlineitem{\sphinxstyleemphasis{Cannot use delete\_atoms unless atoms have IDs}}
\sphinxAtStartPar
Your atoms do not have IDs, so the delete\_atoms command cannot be
used.

\sphinxlineitem{\sphinxstyleemphasis{Cannot use delete\_bonds with non\sphinxhyphen{}molecular system}}
\sphinxAtStartPar
Your choice of atom style does not have bonds.

\sphinxlineitem{\sphinxstyleemphasis{Cannot use dump\_modify fileper without \% in dump file name}}
\sphinxAtStartPar
Self\sphinxhyphen{}explanatory.

\sphinxlineitem{\sphinxstyleemphasis{Cannot use dump\_modify nfile without \% in dump file name}}
\sphinxAtStartPar
Self\sphinxhyphen{}explanatory.

\sphinxlineitem{\sphinxstyleemphasis{Cannot use dynamic group with fix adapt atom}}
\sphinxAtStartPar
This is not yet supported.

\sphinxlineitem{\sphinxstyleemphasis{Cannot use fix TMD unless atom map exists}}
\sphinxAtStartPar
Using this fix requires the ability to lookup an atom index, which is
provided by an atom map.  An atom map does not exist (by default) for
non\sphinxhyphen{}molecular problems.  Using the atom\_modify map command will force
an atom map to be created.

\sphinxlineitem{\sphinxstyleemphasis{Cannot use fix bond/break with non\sphinxhyphen{}molecular systems}}
\sphinxAtStartPar
Only systems with bonds that can be changed can be used.  Atom\_style
template does not qualify.

\sphinxlineitem{\sphinxstyleemphasis{Cannot use fix bond/create with non\sphinxhyphen{}molecular systems}}
\sphinxAtStartPar
Only systems with bonds that can be changed can be used.  Atom\_style
template does not qualify.

\sphinxlineitem{\sphinxstyleemphasis{Cannot use fix bond/swap with non\sphinxhyphen{}molecular systems}}
\sphinxAtStartPar
Only systems with bonds that can be changed can be used.  Atom\_style
template does not qualify.

\sphinxlineitem{\sphinxstyleemphasis{Cannot use fix box/relax on a 2nd non\sphinxhyphen{}periodic dimension}}
\sphinxAtStartPar
When specifying an off\sphinxhyphen{}diagonal pressure component, the second of the two
dimensions must be periodic.  E.g. if the xy component is specified,
then the y dimension must be periodic.

\sphinxlineitem{\sphinxstyleemphasis{Cannot use fix box/relax on a non\sphinxhyphen{}periodic dimension}}
\sphinxAtStartPar
When specifying a diagonal pressure component, the dimension must be
periodic.

\sphinxlineitem{\sphinxstyleemphasis{Cannot use fix box/relax with both relaxation and scaling on a tilt factor}}
\sphinxAtStartPar
When specifying scaling on a tilt factor component, that component can not
also be controlled by the barostat. E.g. if scalexy yes is specified and
also keyword tri or xy, this is wrong.

\sphinxlineitem{\sphinxstyleemphasis{Cannot use fix box/relax with tilt factor scaling on a 2nd non\sphinxhyphen{}periodic dimension}}
\sphinxAtStartPar
When specifying scaling on a tilt factor component, the second of the two
dimensions must be periodic.  E.g. if the xy component is specified,
then the y dimension must be periodic.

\sphinxlineitem{\sphinxstyleemphasis{Cannot use fix deform on a shrink\sphinxhyphen{}wrapped boundary}}
\sphinxAtStartPar
The x, y, z options cannot be applied to shrink\sphinxhyphen{}wrapped
dimensions.

\sphinxlineitem{\sphinxstyleemphasis{Cannot use fix deform tilt on a shrink\sphinxhyphen{}wrapped 2nd dim}}
\sphinxAtStartPar
This is because the shrink\sphinxhyphen{}wrapping will change the value
of the strain implied by the tilt factor.

\sphinxlineitem{\sphinxstyleemphasis{Cannot use fix deform trate on a box with zero tilt}}
\sphinxAtStartPar
The trate style alters the current strain.

\sphinxlineitem{\sphinxstyleemphasis{Cannot use fix deposit rigid and not molecule}}
\sphinxAtStartPar
Self\sphinxhyphen{}explanatory.

\sphinxlineitem{\sphinxstyleemphasis{Cannot use fix deposit rigid and shake}}
\sphinxAtStartPar
These two attributes are conflicting.

\sphinxlineitem{\sphinxstyleemphasis{Cannot use fix deposit shake and not molecule}}
\sphinxAtStartPar
Self\sphinxhyphen{}explanatory.

\sphinxlineitem{\sphinxstyleemphasis{Cannot use fix enforce2d with 3d simulation}}
\sphinxAtStartPar
Self\sphinxhyphen{}explanatory.

\sphinxlineitem{\sphinxstyleemphasis{Cannot use fix gcmc in a 2d simulation}}
\sphinxAtStartPar
Fix gcmc is set up to run in 3d only. No 2d simulations with fix gcmc
are allowed.

\sphinxlineitem{\sphinxstyleemphasis{Cannot use fix gcmc shake and not molecule}}
\sphinxAtStartPar
Self\sphinxhyphen{}explanatory.

\sphinxlineitem{\sphinxstyleemphasis{Cannot use fix msst without per\sphinxhyphen{}type mass defined}}
\sphinxAtStartPar
Self\sphinxhyphen{}explanatory.

\sphinxlineitem{\sphinxstyleemphasis{Cannot use fix npt and fix deform on same component of stress tensor}}
\sphinxAtStartPar
This would be changing the same box dimension twice.

\sphinxlineitem{\sphinxstyleemphasis{Cannot use fix nvt/npt/nph on a 2nd non\sphinxhyphen{}periodic dimension}}
\sphinxAtStartPar
When specifying an off\sphinxhyphen{}diagonal pressure component, the second of the two
dimensions must be periodic.  E.g. if the xy component is specified,
then the y dimension must be periodic.

\sphinxlineitem{\sphinxstyleemphasis{Cannot use fix nvt/npt/nph on a non\sphinxhyphen{}periodic dimension}}
\sphinxAtStartPar
When specifying a diagonal pressure component, the dimension must be
periodic.

\sphinxlineitem{\sphinxstyleemphasis{Cannot use fix nvt/npt/nph with both xy dynamics and xy scaling}}
\sphinxAtStartPar
Self\sphinxhyphen{}explanatory.

\sphinxlineitem{\sphinxstyleemphasis{Cannot use fix nvt/npt/nph with both xz dynamics and xz scaling}}
\sphinxAtStartPar
Self\sphinxhyphen{}explanatory.

\sphinxlineitem{\sphinxstyleemphasis{Cannot use fix nvt/npt/nph with both yz dynamics and yz scaling}}
\sphinxAtStartPar
Self\sphinxhyphen{}explanatory.

\sphinxlineitem{\sphinxstyleemphasis{Cannot use fix nvt/npt/nph with xy scaling when y is non\sphinxhyphen{}periodic dimension}}
\sphinxAtStartPar
The second dimension in the barostatted tilt factor must be periodic.

\sphinxlineitem{\sphinxstyleemphasis{Cannot use fix nvt/npt/nph with xz scaling when z is non\sphinxhyphen{}periodic dimension}}
\sphinxAtStartPar
The second dimension in the barostatted tilt factor must be periodic.

\sphinxlineitem{\sphinxstyleemphasis{Cannot use fix nvt/npt/nph with yz scaling when z is non\sphinxhyphen{}periodic dimension}}
\sphinxAtStartPar
The second dimension in the barostatted tilt factor must be periodic.

\sphinxlineitem{\sphinxstyleemphasis{Cannot use fix pour rigid and not molecule}}
\sphinxAtStartPar
Self\sphinxhyphen{}explanatory.

\sphinxlineitem{\sphinxstyleemphasis{Cannot use fix pour rigid and shake}}
\sphinxAtStartPar
These two attributes are conflicting.

\sphinxlineitem{\sphinxstyleemphasis{Cannot use fix pour shake and not molecule}}
\sphinxAtStartPar
Self\sphinxhyphen{}explanatory.

\sphinxlineitem{\sphinxstyleemphasis{Cannot use fix pour with triclinic box}}
\sphinxAtStartPar
This option is not yet supported.

\sphinxlineitem{\sphinxstyleemphasis{Cannot use fix press/berendsen and fix deform on same component of stress tensor}}
\sphinxAtStartPar
These commands both change the box size/shape, so you cannot use both
together.

\sphinxlineitem{\sphinxstyleemphasis{Cannot use fix press/berendsen on a non\sphinxhyphen{}periodic dimension}}
\sphinxAtStartPar
Self\sphinxhyphen{}explanatory.

\sphinxlineitem{\sphinxstyleemphasis{Cannot use fix press/berendsen with triclinic box}}
\sphinxAtStartPar
Self\sphinxhyphen{}explanatory.

\sphinxlineitem{\sphinxstyleemphasis{Cannot use fix reax/bonds without pair\_style reax}}
\sphinxAtStartPar
Self\sphinxhyphen{}explanatory.

\sphinxlineitem{\sphinxstyleemphasis{Cannot use fix rigid npt/nph and fix deform on same component of stress tensor}}
\sphinxAtStartPar
This would be changing the same box dimension twice.

\sphinxlineitem{\sphinxstyleemphasis{Cannot use fix rigid npt/nph on a non\sphinxhyphen{}periodic dimension}}
\sphinxAtStartPar
When specifying a diagonal pressure component, the dimension must be
periodic.

\sphinxlineitem{\sphinxstyleemphasis{Cannot use fix rigid/small npt/nph on a non\sphinxhyphen{}periodic dimension}}
\sphinxAtStartPar
When specifying a diagonal pressure component, the dimension must be
periodic.

\sphinxlineitem{\sphinxstyleemphasis{Cannot use fix shake with non\sphinxhyphen{}molecular system}}
\sphinxAtStartPar
Your choice of atom style does not have bonds.

\sphinxlineitem{\sphinxstyleemphasis{Cannot use fix ttm with 2d simulation}}
\sphinxAtStartPar
This is a current restriction of this fix due to the grid it creates.

\sphinxlineitem{\sphinxstyleemphasis{Cannot use fix ttm with triclinic box}}
\sphinxAtStartPar
This is a current restriction of this fix due to the grid it creates.

\sphinxlineitem{\sphinxstyleemphasis{Cannot use fix tune/kspace without a kspace style}}
\sphinxAtStartPar
Self\sphinxhyphen{}explanatory.

\sphinxlineitem{\sphinxstyleemphasis{Cannot use fix tune/kspace without a pair style}}
\sphinxAtStartPar
This fix (tune/kspace) can only be used when a pair style has been specified.

\sphinxlineitem{\sphinxstyleemphasis{Cannot use fix wall in periodic dimension}}
\sphinxAtStartPar
Self\sphinxhyphen{}explanatory.

\sphinxlineitem{\sphinxstyleemphasis{Cannot use fix wall zlo/zhi for a 2d simulation}}
\sphinxAtStartPar
Self\sphinxhyphen{}explanatory.

\sphinxlineitem{\sphinxstyleemphasis{Cannot use fix wall/reflect in periodic dimension}}
\sphinxAtStartPar
Self\sphinxhyphen{}explanatory.

\sphinxlineitem{\sphinxstyleemphasis{Cannot use fix wall/reflect zlo/zhi for a 2d simulation}}
\sphinxAtStartPar
Self\sphinxhyphen{}explanatory.

\sphinxlineitem{\sphinxstyleemphasis{Cannot use fix wall/srd in periodic dimension}}
\sphinxAtStartPar
Self\sphinxhyphen{}explanatory.

\sphinxlineitem{\sphinxstyleemphasis{Cannot use fix wall/srd more than once}}
\sphinxAtStartPar
Nor is their a need to since multiple walls can be specified
in one command.

\sphinxlineitem{\sphinxstyleemphasis{Cannot use fix wall/srd without fix srd}}
\sphinxAtStartPar
Self\sphinxhyphen{}explanatory.

\sphinxlineitem{\sphinxstyleemphasis{Cannot use fix wall/srd zlo/zhi for a 2d simulation}}
\sphinxAtStartPar
Self\sphinxhyphen{}explanatory.

\sphinxlineitem{\sphinxstyleemphasis{Cannot use fix\_deposit unless atoms have IDs}}
\sphinxAtStartPar
Self\sphinxhyphen{}explanatory.

\sphinxlineitem{\sphinxstyleemphasis{Cannot use fix\_pour unless atoms have IDs}}
\sphinxAtStartPar
Self\sphinxhyphen{}explanatory.

\sphinxlineitem{\sphinxstyleemphasis{Cannot use include command within an if command}}
\sphinxAtStartPar
Self\sphinxhyphen{}explanatory.

\sphinxlineitem{\sphinxstyleemphasis{Cannot use lines with fix srd unless overlap is set}}
\sphinxAtStartPar
This is because line segments are connected to each other.

\sphinxlineitem{\sphinxstyleemphasis{Cannot use multiple fix wall commands with pair brownian}}
\sphinxAtStartPar
Self\sphinxhyphen{}explanatory.

\sphinxlineitem{\sphinxstyleemphasis{Cannot use multiple fix wall commands with pair lubricate}}
\sphinxAtStartPar
Self\sphinxhyphen{}explanatory.

\sphinxlineitem{\sphinxstyleemphasis{Cannot use multiple fix wall commands with pair lubricate/poly}}
\sphinxAtStartPar
Self\sphinxhyphen{}explanatory.

\sphinxlineitem{\sphinxstyleemphasis{Cannot use multiple fix wall commands with pair lubricateU}}
\sphinxAtStartPar
Self\sphinxhyphen{}explanatory.

\sphinxlineitem{\sphinxstyleemphasis{Cannot use neigh\_modify exclude with GPU neighbor builds}}
\sphinxAtStartPar
This is a current limitation of the GPU implementation
in LAMMPS.

\sphinxlineitem{\sphinxstyleemphasis{Cannot use neighbor bins \sphinxhyphen{} box size \textless{}\textless{} cutoff}}
\sphinxAtStartPar
Too many neighbor bins will be created.  This typically happens when
the simulation box is very small in some dimension, compared to the
neighbor cutoff.  Use the “nsq” style instead of “bin” style.

\sphinxlineitem{\sphinxstyleemphasis{Cannot use newton pair with beck/gpu pair style}}
\sphinxAtStartPar
Self\sphinxhyphen{}explanatory.

\sphinxlineitem{\sphinxstyleemphasis{Cannot use newton pair with born/coul/long/gpu pair style}}
\sphinxAtStartPar
Self\sphinxhyphen{}explanatory.

\sphinxlineitem{\sphinxstyleemphasis{Cannot use newton pair with born/coul/wolf/gpu pair style}}
\sphinxAtStartPar
Self\sphinxhyphen{}explanatory.

\sphinxlineitem{\sphinxstyleemphasis{Cannot use newton pair with born/gpu pair style}}
\sphinxAtStartPar
Self\sphinxhyphen{}explanatory.

\sphinxlineitem{\sphinxstyleemphasis{Cannot use newton pair with buck/coul/cut/gpu pair style}}
\sphinxAtStartPar
Self\sphinxhyphen{}explanatory.

\sphinxlineitem{\sphinxstyleemphasis{Cannot use newton pair with buck/coul/long/gpu pair style}}
\sphinxAtStartPar
Self\sphinxhyphen{}explanatory.

\sphinxlineitem{\sphinxstyleemphasis{Cannot use newton pair with buck/gpu pair style}}
\sphinxAtStartPar
Self\sphinxhyphen{}explanatory.

\sphinxlineitem{\sphinxstyleemphasis{Cannot use newton pair with colloid/gpu pair style}}
\sphinxAtStartPar
Self\sphinxhyphen{}explanatory.

\sphinxlineitem{\sphinxstyleemphasis{Cannot use newton pair with coul/cut/gpu pair style}}
\sphinxAtStartPar
Self\sphinxhyphen{}explanatory.

\sphinxlineitem{\sphinxstyleemphasis{Cannot use newton pair with coul/debye/gpu pair style}}
\sphinxAtStartPar
Self\sphinxhyphen{}explanatory.

\sphinxlineitem{\sphinxstyleemphasis{Cannot use newton pair with coul/dsf/gpu pair style}}
\sphinxAtStartPar
Self\sphinxhyphen{}explanatory.

\sphinxlineitem{\sphinxstyleemphasis{Cannot use newton pair with coul/long/gpu pair style}}
\sphinxAtStartPar
Self\sphinxhyphen{}explanatory.

\sphinxlineitem{\sphinxstyleemphasis{Cannot use newton pair with dipole/cut/gpu pair style}}
\sphinxAtStartPar
Self\sphinxhyphen{}explanatory.

\sphinxlineitem{\sphinxstyleemphasis{Cannot use newton pair with dipole/sf/gpu pair style}}
\sphinxAtStartPar
Self\sphinxhyphen{}explanatory.

\sphinxlineitem{\sphinxstyleemphasis{Cannot use newton pair with dpd/gpu pair style}}
\sphinxAtStartPar
Self\sphinxhyphen{}explanatory.

\sphinxlineitem{\sphinxstyleemphasis{Cannot use newton pair with dpd/tstat/gpu pair style}}
\sphinxAtStartPar
Self\sphinxhyphen{}explanatory.

\sphinxlineitem{\sphinxstyleemphasis{Cannot use newton pair with eam/alloy/gpu pair style}}
\sphinxAtStartPar
Self\sphinxhyphen{}explanatory.

\sphinxlineitem{\sphinxstyleemphasis{Cannot use newton pair with eam/fs/gpu pair style}}
\sphinxAtStartPar
Self\sphinxhyphen{}explanatory.

\sphinxlineitem{\sphinxstyleemphasis{Cannot use newton pair with eam/gpu pair style}}
\sphinxAtStartPar
Self\sphinxhyphen{}explanatory.

\sphinxlineitem{\sphinxstyleemphasis{Cannot use newton pair with gauss/gpu pair style}}
\sphinxAtStartPar
Self\sphinxhyphen{}explanatory.

\sphinxlineitem{\sphinxstyleemphasis{Cannot use newton pair with gayberne/gpu pair style}}
\sphinxAtStartPar
Self\sphinxhyphen{}explanatory.

\sphinxlineitem{\sphinxstyleemphasis{Cannot use newton pair with lj/charmm/coul/long/gpu pair style}}
\sphinxAtStartPar
Self\sphinxhyphen{}explanatory.

\sphinxlineitem{\sphinxstyleemphasis{Cannot use newton pair with lj/class2/coul/long/gpu pair style}}
\sphinxAtStartPar
Self\sphinxhyphen{}explanatory.

\sphinxlineitem{\sphinxstyleemphasis{Cannot use newton pair with lj/class2/gpu pair style}}
\sphinxAtStartPar
Self\sphinxhyphen{}explanatory.

\sphinxlineitem{\sphinxstyleemphasis{Cannot use newton pair with lj/cubic/gpu pair style}}
\sphinxAtStartPar
Self\sphinxhyphen{}explanatory.

\sphinxlineitem{\sphinxstyleemphasis{Cannot use newton pair with lj/cut/coul/cut/gpu pair style}}
\sphinxAtStartPar
Self\sphinxhyphen{}explanatory.

\sphinxlineitem{\sphinxstyleemphasis{Cannot use newton pair with lj/cut/coul/debye/gpu pair style}}
\sphinxAtStartPar
Self\sphinxhyphen{}explanatory.

\sphinxlineitem{\sphinxstyleemphasis{Cannot use newton pair with lj/cut/coul/dsf/gpu pair style}}
\sphinxAtStartPar
Self\sphinxhyphen{}explanatory.

\sphinxlineitem{\sphinxstyleemphasis{Cannot use newton pair with lj/cut/coul/long/gpu pair style}}
\sphinxAtStartPar
Self\sphinxhyphen{}explanatory.

\sphinxlineitem{\sphinxstyleemphasis{Cannot use newton pair with lj/cut/coul/msm/gpu pair style}}
\sphinxAtStartPar
Self\sphinxhyphen{}explanatory.

\sphinxlineitem{\sphinxstyleemphasis{Cannot use newton pair with lj/cut/gpu pair style}}
\sphinxAtStartPar
Self\sphinxhyphen{}explanatory.

\sphinxlineitem{\sphinxstyleemphasis{Cannot use newton pair with lj/expand/gpu pair style}}
\sphinxAtStartPar
Self\sphinxhyphen{}explanatory.

\sphinxlineitem{\sphinxstyleemphasis{Cannot use newton pair with lj/gromacs/gpu pair style}}
\sphinxAtStartPar
Self\sphinxhyphen{}explanatory.

\sphinxlineitem{\sphinxstyleemphasis{Cannot use newton pair with lj/spica/coul/long/gpu pair style}}
\sphinxAtStartPar
Self\sphinxhyphen{}explanatory.

\sphinxlineitem{\sphinxstyleemphasis{Cannot use newton pair with lj/spica/gpu pair style}}
\sphinxAtStartPar
Self\sphinxhyphen{}explanatory.

\sphinxlineitem{\sphinxstyleemphasis{Cannot use newton pair with lj96/cut/gpu pair style}}
\sphinxAtStartPar
Self\sphinxhyphen{}explanatory.

\sphinxlineitem{\sphinxstyleemphasis{Cannot use newton pair with mie/cut/gpu pair style}}
\sphinxAtStartPar
Self\sphinxhyphen{}explanatory.

\sphinxlineitem{\sphinxstyleemphasis{Cannot use newton pair with morse/gpu pair style}}
\sphinxAtStartPar
Self\sphinxhyphen{}explanatory.

\sphinxlineitem{\sphinxstyleemphasis{Cannot use newton pair with resquared/gpu pair style}}
\sphinxAtStartPar
Self\sphinxhyphen{}explanatory.

\sphinxlineitem{\sphinxstyleemphasis{Cannot use newton pair with soft/gpu pair style}}
\sphinxAtStartPar
Self\sphinxhyphen{}explanatory.

\sphinxlineitem{\sphinxstyleemphasis{Cannot use newton pair with table/gpu pair style}}
\sphinxAtStartPar
Self\sphinxhyphen{}explanatory.

\sphinxlineitem{\sphinxstyleemphasis{Cannot use newton pair with yukawa/colloid/gpu pair style}}
\sphinxAtStartPar
Self\sphinxhyphen{}explanatory.

\sphinxlineitem{\sphinxstyleemphasis{Cannot use newton pair with yukawa/gpu pair style}}
\sphinxAtStartPar
Self\sphinxhyphen{}explanatory.

\sphinxlineitem{\sphinxstyleemphasis{Cannot use newton pair with zbl/gpu pair style}}
\sphinxAtStartPar
Self\sphinxhyphen{}explanatory.

\sphinxlineitem{\sphinxstyleemphasis{Cannot use non\sphinxhyphen{}zero forces in an energy minimization}}
\sphinxAtStartPar
Fix setforce cannot be used in this manner.  Use fix addforce
instead.

\sphinxlineitem{\sphinxstyleemphasis{Cannot use non\sphinxhyphen{}periodic boundares with fix ttm}}
\sphinxAtStartPar
This fix requires a fully periodic simulation box.

\sphinxlineitem{\sphinxstyleemphasis{Cannot use non\sphinxhyphen{}periodic boundaries with Ewald}}
\sphinxAtStartPar
For kspace style ewald, all 3 dimensions must have periodic boundaries
unless you use the kspace\_modify command to define a 2d slab with a
non\sphinxhyphen{}periodic z dimension.

\sphinxlineitem{\sphinxstyleemphasis{Cannot use non\sphinxhyphen{}periodic boundaries with EwaldDisp}}
\sphinxAtStartPar
For kspace style ewald/disp, all 3 dimensions must have periodic
boundaries unless you use the kspace\_modify command to define a 2d
slab with a non\sphinxhyphen{}periodic z dimension.

\sphinxlineitem{\sphinxstyleemphasis{Cannot use non\sphinxhyphen{}periodic boundaries with PPPM}}
\sphinxAtStartPar
For kspace style pppm, all 3 dimensions must have periodic boundaries
unless you use the kspace\_modify command to define a 2d slab with a
non\sphinxhyphen{}periodic z dimension.

\sphinxlineitem{\sphinxstyleemphasis{Cannot use non\sphinxhyphen{}periodic boundaries with PPPMDisp}}
\sphinxAtStartPar
For kspace style pppm/disp, all 3 dimensions must have periodic
boundaries unless you use the kspace\_modify command to define a 2d
slab with a non\sphinxhyphen{}periodic z dimension.

\sphinxlineitem{\sphinxstyleemphasis{Cannot use order greater than 8 with pppm/gpu.}}
\sphinxAtStartPar
Self\sphinxhyphen{}explanatory.

\sphinxlineitem{\sphinxstyleemphasis{Cannot use package gpu neigh yes with triclinic box}}
\sphinxAtStartPar
This is a current restriction in LAMMPS.

\sphinxlineitem{\sphinxstyleemphasis{Cannot use pair tail corrections with 2d simulations}}
\sphinxAtStartPar
The correction factors are only currently defined for 3d systems.

\sphinxlineitem{\sphinxstyleemphasis{Cannot use processors part command without using partitions}}
\sphinxAtStartPar
See the command\sphinxhyphen{}line \sphinxhyphen{}partition switch.

\sphinxlineitem{\sphinxstyleemphasis{Cannot use ramp in variable formula between runs}}
\sphinxAtStartPar
This is because the ramp() function is time dependent.

\sphinxlineitem{\sphinxstyleemphasis{Cannot use read\_data add before simulation box is defined}}
\sphinxAtStartPar
Self\sphinxhyphen{}explanatory.

\sphinxlineitem{\sphinxstyleemphasis{Cannot use read\_data extra with add flag}}
\sphinxAtStartPar
Self\sphinxhyphen{}explanatory.

\sphinxlineitem{\sphinxstyleemphasis{Cannot use read\_data offset without add flag}}
\sphinxAtStartPar
Self\sphinxhyphen{}explanatory.

\sphinxlineitem{\sphinxstyleemphasis{Cannot use read\_data shift without add flag}}
\sphinxAtStartPar
Self\sphinxhyphen{}explanatory.

\sphinxlineitem{\sphinxstyleemphasis{Cannot use region INF or EDGE when box does not exist}}
\sphinxAtStartPar
Regions that extend to the box boundaries can only be used after the
create\_box command has been used.

\sphinxlineitem{\sphinxstyleemphasis{Cannot use set atom with no atom IDs defined}}
\sphinxAtStartPar
Atom IDs are not defined, so they cannot be used to identify an atom.

\sphinxlineitem{\sphinxstyleemphasis{Cannot use set mol with no molecule IDs defined}}
\sphinxAtStartPar
Self\sphinxhyphen{}explanatory.

\sphinxlineitem{\sphinxstyleemphasis{Cannot use swiggle in variable formula between runs}}
\sphinxAtStartPar
This is a function of elapsed time.

\sphinxlineitem{\sphinxstyleemphasis{Cannot use tris with fix srd unless overlap is set}}
\sphinxAtStartPar
This is because triangles are connected to each other.

\sphinxlineitem{\sphinxstyleemphasis{Cannot use variable energy with constant efield in fix efield}}
\sphinxAtStartPar
LAMMPS computes the energy itself when the E\sphinxhyphen{}field is constant.

\sphinxlineitem{\sphinxstyleemphasis{Cannot use variable energy with constant force in fix addforce}}
\sphinxAtStartPar
This is because for constant force, LAMMPS can compute the change
in energy directly.

\sphinxlineitem{\sphinxstyleemphasis{Cannot use variable every setting for dump dcd}}
\sphinxAtStartPar
The format of DCD dump files requires snapshots be output
at a constant frequency.

\sphinxlineitem{\sphinxstyleemphasis{Cannot use variable every setting for dump xtc}}
\sphinxAtStartPar
The format of this file requires snapshots at regular intervals.

\sphinxlineitem{\sphinxstyleemphasis{Cannot use vdisplace in variable formula between runs}}
\sphinxAtStartPar
This is a function of elapsed time.

\sphinxlineitem{\sphinxstyleemphasis{Cannot use velocity bias command without temp keyword}}
\sphinxAtStartPar
Self\sphinxhyphen{}explanatory.

\sphinxlineitem{\sphinxstyleemphasis{Cannot use velocity create loop all unless atoms have IDs}}
\sphinxAtStartPar
Atoms in the simulation to do not have IDs, so this style
of velocity creation cannot be performed.

\sphinxlineitem{\sphinxstyleemphasis{Cannot use wall in periodic dimension}}
\sphinxAtStartPar
Self\sphinxhyphen{}explanatory.

\sphinxlineitem{\sphinxstyleemphasis{Cannot use write\_restart fileper without \% in restart file name}}
\sphinxAtStartPar
Self\sphinxhyphen{}explanatory.

\sphinxlineitem{\sphinxstyleemphasis{Cannot use write\_restart nfile without \% in restart file name}}
\sphinxAtStartPar
Self\sphinxhyphen{}explanatory.

\sphinxlineitem{\sphinxstyleemphasis{Cannot wiggle and shear fix wall/gran}}
\sphinxAtStartPar
Cannot specify both options at the same time.

\sphinxlineitem{\sphinxstyleemphasis{Cannot write to restart file \sphinxhyphen{} MPI error: \%s}}
\sphinxAtStartPar
This error was generated by MPI when reading/writing an MPI\sphinxhyphen{}IO restart
file.

\sphinxlineitem{\sphinxstyleemphasis{Cannot yet use KSpace solver with grid with comm style tiled}}
\sphinxAtStartPar
This is current restriction in LAMMPS.

\sphinxlineitem{\sphinxstyleemphasis{Cannot yet use comm\_style tiled with multi\sphinxhyphen{}mode comm}}
\sphinxAtStartPar
Self\sphinxhyphen{}explanatory.

\sphinxlineitem{\sphinxstyleemphasis{Cannot yet use comm\_style tiled with triclinic box}}
\sphinxAtStartPar
Self\sphinxhyphen{}explanatory.

\sphinxlineitem{\sphinxstyleemphasis{Cannot yet use compute tally with Kokkos}}
\sphinxAtStartPar
This feature is not yet supported.

\sphinxlineitem{\sphinxstyleemphasis{Cannot yet use fix bond/break with this improper style}}
\sphinxAtStartPar
This is a current restriction in LAMMPS.

\sphinxlineitem{\sphinxstyleemphasis{Cannot yet use fix bond/create with this improper style}}
\sphinxAtStartPar
This is a current restriction in LAMMPS.

\sphinxlineitem{\sphinxstyleemphasis{Cannot yet use minimize with Kokkos}}
\sphinxAtStartPar
This feature is not yet supported.

\sphinxlineitem{\sphinxstyleemphasis{Cannot yet use pair hybrid with Kokkos}}
\sphinxAtStartPar
This feature is not yet supported.

\sphinxlineitem{\sphinxstyleemphasis{Cannot zero Langevin force of 0 atoms}}
\sphinxAtStartPar
The group has zero atoms, so you cannot request its force
be zeroed.

\sphinxlineitem{\sphinxstyleemphasis{Cannot zero gld force for zero atoms}}
\sphinxAtStartPar
There are no atoms currently in the group.

\sphinxlineitem{\sphinxstyleemphasis{Cannot zero momentum of no atoms}}
\sphinxAtStartPar
Self\sphinxhyphen{}explanatory.

\sphinxlineitem{\sphinxstyleemphasis{Change\_box command before simulation box is defined}}
\sphinxAtStartPar
Self\sphinxhyphen{}explanatory.

\sphinxlineitem{\sphinxstyleemphasis{Change\_box volume used incorrectly}}
\sphinxAtStartPar
The “dim volume” option must be used immediately following one or two
settings for “dim1 …” (and optionally “dim2 …”) and must be for a
different dimension, i.e. dim != dim1 and dim != dim2.

\sphinxlineitem{\sphinxstyleemphasis{Chunk/atom compute does not exist for compute angmom/chunk}}
\sphinxAtStartPar
Self\sphinxhyphen{}explanatory.

\sphinxlineitem{\sphinxstyleemphasis{Chunk/atom compute does not exist for compute com/chunk}}
\sphinxAtStartPar
Self\sphinxhyphen{}explanatory.

\sphinxlineitem{\sphinxstyleemphasis{Chunk/atom compute does not exist for compute gyration/chunk}}
\sphinxAtStartPar
Self\sphinxhyphen{}explanatory.

\sphinxlineitem{\sphinxstyleemphasis{Chunk/atom compute does not exist for compute inertia/chunk}}
\sphinxAtStartPar
Self\sphinxhyphen{}explanatory.

\sphinxlineitem{\sphinxstyleemphasis{Chunk/atom compute does not exist for compute msd/chunk}}
\sphinxAtStartPar
Self\sphinxhyphen{}explanatory.

\sphinxlineitem{\sphinxstyleemphasis{Chunk/atom compute does not exist for compute omega/chunk}}
\sphinxAtStartPar
Self\sphinxhyphen{}explanatory.

\sphinxlineitem{\sphinxstyleemphasis{Chunk/atom compute does not exist for compute property/chunk}}
\sphinxAtStartPar
Self\sphinxhyphen{}explanatory.

\sphinxlineitem{\sphinxstyleemphasis{Chunk/atom compute does not exist for compute temp/chunk}}
\sphinxAtStartPar
Self\sphinxhyphen{}explanatory.

\sphinxlineitem{\sphinxstyleemphasis{Chunk/atom compute does not exist for compute torque/chunk}}
\sphinxAtStartPar
Self\sphinxhyphen{}explanatory.

\sphinxlineitem{\sphinxstyleemphasis{Chunk/atom compute does not exist for compute vcm/chunk}}
\sphinxAtStartPar
Self\sphinxhyphen{}explanatory.

\sphinxlineitem{\sphinxstyleemphasis{Chunk/atom compute does not exist for fix ave/chunk}}
\sphinxAtStartPar
Self\sphinxhyphen{}explanatory.

\sphinxlineitem{\sphinxstyleemphasis{Comm tiled invalid index in box drop brick}}
\sphinxAtStartPar
Internal error check in comm\_style tiled which should not occur.
Contact the developers.

\sphinxlineitem{\sphinxstyleemphasis{Comm tiled mis\sphinxhyphen{}match in box drop brick}}
\sphinxAtStartPar
Internal error check in comm\_style tiled which should not occur.
Contact the developers.

\sphinxlineitem{\sphinxstyleemphasis{Comm\_modify group != atom\_modify first group}}
\sphinxAtStartPar
Self\sphinxhyphen{}explanatory.

\sphinxlineitem{\sphinxstyleemphasis{Communication cutoff for comm\_style tiled cannot exceed periodic box length}}
\sphinxAtStartPar
Self\sphinxhyphen{}explanatory.

\sphinxlineitem{\sphinxstyleemphasis{Communication cutoff too small for SNAP micro load balancing}}
\sphinxAtStartPar
This can happen if you change the neighbor skin after your pair\_style
command or if your box dimensions grow during a run. You can set the
cutoff explicitly via the comm\_modify cutoff command.

\sphinxlineitem{\sphinxstyleemphasis{Compute \%s does not allow use of dynamic group}}
\sphinxAtStartPar
Dynamic groups have not yet been enabled for this compute.

\sphinxlineitem{\sphinxstyleemphasis{Compute for fix pafi does not calculate a local array}}
\sphinxAtStartPar
Self\sphinxhyphen{}explanatory.

\sphinxlineitem{\sphinxstyleemphasis{Compute for fix pafi must have 9 fields per atom}}
\sphinxAtStartPar
Self\sphinxhyphen{}explanatory.

\sphinxlineitem{\sphinxstyleemphasis{Compute ID for compute chunk /atom does not exist}}
\sphinxAtStartPar
Self\sphinxhyphen{}explanatory.

\sphinxlineitem{\sphinxstyleemphasis{Compute ID for compute chunk/atom does not exist}}
\sphinxAtStartPar
Self\sphinxhyphen{}explanatory.

\sphinxlineitem{\sphinxstyleemphasis{Compute gyration ID does not exist for compute gyration/shape}}
\sphinxAtStartPar
Self\sphinxhyphen{}explanatory. Provide a valid compute ID.

\sphinxlineitem{\sphinxstyleemphasis{Compute gyration/shape compute ID does not point to a gyration compute}}
\sphinxAtStartPar
Self\sphinxhyphen{}explanatory. Provide and ID of a compute gyration command.

\sphinxlineitem{\sphinxstyleemphasis{Compute ID for compute reduce does not exist}}
\sphinxAtStartPar
Self\sphinxhyphen{}explanatory.

\sphinxlineitem{\sphinxstyleemphasis{Compute ID for compute slice does not exist}}
\sphinxAtStartPar
Self\sphinxhyphen{}explanatory.

\sphinxlineitem{\sphinxstyleemphasis{Compute ID for fix ave/atom does not exist}}
\sphinxAtStartPar
Self\sphinxhyphen{}explanatory.

\sphinxlineitem{\sphinxstyleemphasis{Compute ID for fix ave/chunk does not exist}}
\sphinxAtStartPar
Self\sphinxhyphen{}explanatory.

\sphinxlineitem{\sphinxstyleemphasis{Compute ID for fix ave/correlate does not exist}}
\sphinxAtStartPar
Self\sphinxhyphen{}explanatory.

\sphinxlineitem{\sphinxstyleemphasis{Compute ID for fix ave/histo does not exist}}
\sphinxAtStartPar
Self\sphinxhyphen{}explanatory.

\sphinxlineitem{\sphinxstyleemphasis{Compute ID for fix ave/time does not exist}}
\sphinxAtStartPar
Self\sphinxhyphen{}explanatory.

\sphinxlineitem{\sphinxstyleemphasis{Compute ID for fix numdiff does not exist}}
\sphinxAtStartPar
Self\sphinxhyphen{}explanatory.

\sphinxlineitem{\sphinxstyleemphasis{Compute ID for fix numdiff/virial does not exist}}
\sphinxAtStartPar
Self\sphinxhyphen{}explanatory.

\sphinxlineitem{\sphinxstyleemphasis{Compute ID for fix store/state does not exist}}
\sphinxAtStartPar
Self\sphinxhyphen{}explanatory.

\sphinxlineitem{\sphinxstyleemphasis{Compute ID for fix vector does not exist}}
\sphinxAtStartPar
Self\sphinxhyphen{}explanatory.

\sphinxlineitem{\sphinxstyleemphasis{Compute ID must be alphanumeric or underscore characters}}
\sphinxAtStartPar
Self\sphinxhyphen{}explanatory.

\sphinxlineitem{\sphinxstyleemphasis{Compute angle/local used when angles are not allowed}}
\sphinxAtStartPar
The atom style does not support angles.

\sphinxlineitem{\sphinxstyleemphasis{Compute angmom/chunk does not use chunk/atom compute}}
\sphinxAtStartPar
The style of the specified compute is not chunk/atom.

\sphinxlineitem{\sphinxstyleemphasis{Compute body/local requires atom style body}}
\sphinxAtStartPar
Self\sphinxhyphen{}explanatory.

\sphinxlineitem{\sphinxstyleemphasis{Compute bond/local used when bonds are not allowed}}
\sphinxAtStartPar
The atom style does not support bonds.

\sphinxlineitem{\sphinxstyleemphasis{Compute centro/atom requires a pair style be defined}}
\sphinxAtStartPar
This is because the computation of the centro\sphinxhyphen{}symmetry values
uses a pairwise neighbor list.

\sphinxlineitem{\sphinxstyleemphasis{Compute chunk/atom bin/cylinder radius is too large for periodic box}}
\sphinxAtStartPar
Radius cannot be bigger than 1/2 of a non\sphinxhyphen{}axis  periodic dimension.

\sphinxlineitem{\sphinxstyleemphasis{Compute chunk/atom bin/sphere radius is too large for periodic box}}
\sphinxAtStartPar
Radius cannot be bigger than 1/2 of any periodic dimension.

\sphinxlineitem{\sphinxstyleemphasis{Compute chunk/atom compute array is accessed out\sphinxhyphen{}of\sphinxhyphen{}range}}
\sphinxAtStartPar
The index for the array is out of bounds.

\sphinxlineitem{\sphinxstyleemphasis{Compute chunk/atom compute does not calculate a per\sphinxhyphen{}atom array}}
\sphinxAtStartPar
Self\sphinxhyphen{}explanatory.

\sphinxlineitem{\sphinxstyleemphasis{Compute chunk/atom compute does not calculate a per\sphinxhyphen{}atom vector}}
\sphinxAtStartPar
Self\sphinxhyphen{}explanatory.

\sphinxlineitem{\sphinxstyleemphasis{Compute chunk/atom compute does not calculate per\sphinxhyphen{}atom values}}
\sphinxAtStartPar
Self\sphinxhyphen{}explanatory.

\sphinxlineitem{\sphinxstyleemphasis{Compute chunk/atom cylinder axis must be z for 2d}}
\sphinxAtStartPar
Self\sphinxhyphen{}explanatory.

\sphinxlineitem{\sphinxstyleemphasis{Compute chunk/atom fix array is accessed out\sphinxhyphen{}of\sphinxhyphen{}range}}
\sphinxAtStartPar
The index for the array is out of bounds.

\sphinxlineitem{\sphinxstyleemphasis{Compute chunk/atom fix does not calculate a per\sphinxhyphen{}atom array}}
\sphinxAtStartPar
Self\sphinxhyphen{}explanatory.

\sphinxlineitem{\sphinxstyleemphasis{Compute chunk/atom fix does not calculate a per\sphinxhyphen{}atom vector}}
\sphinxAtStartPar
Self\sphinxhyphen{}explanatory.

\sphinxlineitem{\sphinxstyleemphasis{Compute chunk/atom fix does not calculate per\sphinxhyphen{}atom values}}
\sphinxAtStartPar
Self\sphinxhyphen{}explanatory.

\sphinxlineitem{\sphinxstyleemphasis{Compute chunk/atom for triclinic boxes requires units reduced}}
\sphinxAtStartPar
Self\sphinxhyphen{}explanatory.

\sphinxlineitem{\sphinxstyleemphasis{Compute chunk/atom ids once but nchunk is not once}}
\sphinxAtStartPar
You cannot assign chunks IDs to atom permanently if the number of
chunks may change.

\sphinxlineitem{\sphinxstyleemphasis{Compute chunk/atom molecule for non\sphinxhyphen{}molecular system}}
\sphinxAtStartPar
Self\sphinxhyphen{}explanatory.

\sphinxlineitem{\sphinxstyleemphasis{Compute chunk/atom sphere z origin must be 0.0 for 2d}}
\sphinxAtStartPar
Self\sphinxhyphen{}explanatory.

\sphinxlineitem{\sphinxstyleemphasis{Compute chunk/atom stores no IDs for compute property/chunk}}
\sphinxAtStartPar
It will only store IDs if its compress option is enabled.

\sphinxlineitem{\sphinxstyleemphasis{Compute chunk/atom stores no coord1 for compute property/chunk}}
\sphinxAtStartPar
Only certain binning options for compute chunk/atom store coordinates.

\sphinxlineitem{\sphinxstyleemphasis{Compute chunk/atom stores no coord2 for compute property/chunk}}
\sphinxAtStartPar
Only certain binning options for compute chunk/atom store coordinates.

\sphinxlineitem{\sphinxstyleemphasis{Compute chunk/atom stores no coord3 for compute property/chunk}}
\sphinxAtStartPar
Only certain binning options for compute chunk/atom store coordinates.

\sphinxlineitem{\sphinxstyleemphasis{Compute chunk/atom variable is not atom\sphinxhyphen{}style variable}}
\sphinxAtStartPar
Self\sphinxhyphen{}explanatory.

\sphinxlineitem{\sphinxstyleemphasis{Compute chunk/atom without bins cannot use discard mixed}}
\sphinxAtStartPar
That discard option only applies to the binning styles.

\sphinxlineitem{\sphinxstyleemphasis{Compute cluster/atom cutoff is longer than pairwise cutoff}}
\sphinxAtStartPar
Cannot identify clusters beyond cutoff.

\sphinxlineitem{\sphinxstyleemphasis{Compute cluster/atom requires a pair style be defined}}
\sphinxAtStartPar
This is so that the pair style defines a cutoff distance which
is used to find clusters.

\sphinxlineitem{\sphinxstyleemphasis{Compute cna/atom cutoff is longer than pairwise cutoff}}
\sphinxAtStartPar
Self\sphinxhyphen{}explanatory.

\sphinxlineitem{\sphinxstyleemphasis{Compute cna/atom requires a pair style be defined}}
\sphinxAtStartPar
Self\sphinxhyphen{}explanatory.

\sphinxlineitem{\sphinxstyleemphasis{Compute com/chunk does not use chunk/atom compute}}
\sphinxAtStartPar
The style of the specified compute is not chunk/atom.

\sphinxlineitem{\sphinxstyleemphasis{Compute contact/atom requires a pair style be defined}}
\sphinxAtStartPar
Self\sphinxhyphen{}explanatory.

\sphinxlineitem{\sphinxstyleemphasis{Compute contact/atom requires atom style sphere}}
\sphinxAtStartPar
Self\sphinxhyphen{}explanatory.

\sphinxlineitem{\sphinxstyleemphasis{Compute coord/atom cutoff is longer than pairwise cutoff}}
\sphinxAtStartPar
Cannot compute coordination at distances longer than the pair cutoff,
since those atoms are not in the neighbor list.

\sphinxlineitem{\sphinxstyleemphasis{Compute coord/atom requires a pair style be defined}}
\sphinxAtStartPar
Self\sphinxhyphen{}explanatory.

\sphinxlineitem{\sphinxstyleemphasis{Compute damage/atom requires peridynamic potential}}
\sphinxAtStartPar
Damage is a Peridynamic\sphinxhyphen{}specific metric.  It requires you
to be running a Peridynamics simulation.

\sphinxlineitem{\sphinxstyleemphasis{Compute dihedral/local used when dihedrals are not allowed}}
\sphinxAtStartPar
The atom style does not support dihedrals.

\sphinxlineitem{\sphinxstyleemphasis{Compute dilatation/atom cannot be used with this pair style}}
\sphinxAtStartPar
Self\sphinxhyphen{}explanatory.

\sphinxlineitem{\sphinxstyleemphasis{Compute dilatation/atom requires Peridynamic pair style}}
\sphinxAtStartPar
Self\sphinxhyphen{}explanatory.

\sphinxlineitem{\sphinxstyleemphasis{Compute does not allow an extra compute or fix to be reset}}
\sphinxAtStartPar
This is an internal LAMMPS error.  Please report it to the
developers.

\sphinxlineitem{\sphinxstyleemphasis{Compute erotate/asphere requires atom style ellipsoid or line or tri}}
\sphinxAtStartPar
Self\sphinxhyphen{}explanatory.

\sphinxlineitem{\sphinxstyleemphasis{Compute erotate/asphere requires extended particles}}
\sphinxAtStartPar
This compute cannot be used with point particles.

\sphinxlineitem{\sphinxstyleemphasis{Compute erotate/rigid with non\sphinxhyphen{}rigid fix\sphinxhyphen{}ID}}
\sphinxAtStartPar
Self\sphinxhyphen{}explanatory.

\sphinxlineitem{\sphinxstyleemphasis{Compute erotate/sphere requires atom style sphere}}
\sphinxAtStartPar
Self\sphinxhyphen{}explanatory.

\sphinxlineitem{\sphinxstyleemphasis{Compute erotate/sphere/atom requires atom style sphere}}
\sphinxAtStartPar
Self\sphinxhyphen{}explanatory.

\sphinxlineitem{\sphinxstyleemphasis{Compute event/displace has invalid fix event assigned}}
\sphinxAtStartPar
This is an internal LAMMPS error.  Please report it to the
developers.

\sphinxlineitem{\sphinxstyleemphasis{Compute group/group group ID does not exist}}
\sphinxAtStartPar
Self\sphinxhyphen{}explanatory.

\sphinxlineitem{\sphinxstyleemphasis{Compute gyration/chunk does not use chunk/atom compute}}
\sphinxAtStartPar
The style of the specified compute is not chunk/atom.

\sphinxlineitem{\sphinxstyleemphasis{Compute heat/flux compute ID does not compute ke/atom}}
\sphinxAtStartPar
Self\sphinxhyphen{}explanatory.

\sphinxlineitem{\sphinxstyleemphasis{Compute heat/flux compute ID does not compute pe/atom}}
\sphinxAtStartPar
Self\sphinxhyphen{}explanatory.

\sphinxlineitem{\sphinxstyleemphasis{Compute heat/flux compute ID does not compute stress/atom}}
\sphinxAtStartPar
Self\sphinxhyphen{}explanatory.

\sphinxlineitem{\sphinxstyleemphasis{Compute hexorder/atom cutoff is longer than pairwise cutoff}}
\sphinxAtStartPar
Cannot compute order parameter beyond cutoff.

\sphinxlineitem{\sphinxstyleemphasis{Compute hexorder/atom requires a pair style be defined}}
\sphinxAtStartPar
Self\sphinxhyphen{}explanatory.

\sphinxlineitem{\sphinxstyleemphasis{Compute improper/local used when impropers are not allowed}}
\sphinxAtStartPar
The atom style does not support impropers.

\sphinxlineitem{\sphinxstyleemphasis{Compute inertia/chunk does not use chunk/atom compute}}
\sphinxAtStartPar
The style of the specified compute is not chunk/atom.

\sphinxlineitem{\sphinxstyleemphasis{Compute ke/rigid with non\sphinxhyphen{}rigid fix\sphinxhyphen{}ID}}
\sphinxAtStartPar
Self\sphinxhyphen{}explanatory.

\sphinxlineitem{\sphinxstyleemphasis{Compute msd/chunk does not use chunk/atom compute}}
\sphinxAtStartPar
The style of the specified compute is not chunk/atom.

\sphinxlineitem{\sphinxstyleemphasis{Compute msd/chunk nchunk is not static}}
\sphinxAtStartPar
This is required because the MSD cannot be computed consistently if
the number of chunks is changing.  Compute chunk/atom allows setting
nchunk to be static.

\sphinxlineitem{\sphinxstyleemphasis{Compute nve/asphere requires atom style ellipsoid}}
\sphinxAtStartPar
Self\sphinxhyphen{}explanatory.

\sphinxlineitem{\sphinxstyleemphasis{Compute nvt/nph/npt asphere requires atom style ellipsoid}}
\sphinxAtStartPar
Self\sphinxhyphen{}explanatory.

\sphinxlineitem{\sphinxstyleemphasis{Compute nvt/nph/npt body requires atom style body}}
\sphinxAtStartPar
Self\sphinxhyphen{}explanatory.

\sphinxlineitem{\sphinxstyleemphasis{Compute omega/chunk does not use chunk/atom compute}}
\sphinxAtStartPar
The style of the specified compute is not chunk/atom.

\sphinxlineitem{\sphinxstyleemphasis{Compute orientorder/atom cutoff is longer than pairwise cutoff}}
\sphinxAtStartPar
Cannot compute order parameter beyond cutoff.

\sphinxlineitem{\sphinxstyleemphasis{Compute orientorder/atom requires a pair style be defined}}
\sphinxAtStartPar
Self\sphinxhyphen{}explanatory.

\sphinxlineitem{\sphinxstyleemphasis{Compute pair must use group all}}
\sphinxAtStartPar
Pair styles accumulate energy on all atoms.

\sphinxlineitem{\sphinxstyleemphasis{Compute pe must use group all}}
\sphinxAtStartPar
Energies computed by potentials (pair, bond, etc) are computed on all
atoms.

\sphinxlineitem{\sphinxstyleemphasis{Compute plasticity/atom cannot be used with this pair style}}
\sphinxAtStartPar
Self\sphinxhyphen{}explanatory.

\sphinxlineitem{\sphinxstyleemphasis{Compute plasticity/atom requires Peridynamic pair style}}
\sphinxAtStartPar
Self\sphinxhyphen{}explanatory.

\sphinxlineitem{\sphinxstyleemphasis{Compute pressure must use group all}}
\sphinxAtStartPar
Virial contributions computed by potentials (pair, bond, etc) are
computed on all atoms.

\sphinxlineitem{\sphinxstyleemphasis{Compute pressure requires temperature ID to include kinetic energy}}
\sphinxAtStartPar
The keflag cannot be used unless a temperature compute is provided.

\sphinxlineitem{\sphinxstyleemphasis{Compute pressure temperature ID does not compute temperature}}
\sphinxAtStartPar
The compute ID assigned to a pressure computation must compute
temperature.

\sphinxlineitem{\sphinxstyleemphasis{Compute property/atom floating point vector does not exist}}
\sphinxAtStartPar
The command is accessing a vector added by the fix property/atom
command, that does not exist.

\sphinxlineitem{\sphinxstyleemphasis{Compute property/atom for atom property that is not allocated}}
\sphinxAtStartPar
Self\sphinxhyphen{}explanatory.

\sphinxlineitem{\sphinxstyleemphasis{Compute property/atom integer vector does not exist}}
\sphinxAtStartPar
The command is accessing a vector added by the fix property/atom
command, that does not exist.

\sphinxlineitem{\sphinxstyleemphasis{Compute property/chunk does not use chunk/atom compute}}
\sphinxAtStartPar
The style of the specified compute is not chunk/atom.

\sphinxlineitem{\sphinxstyleemphasis{Compute property/local cannot use these inputs together}}
\sphinxAtStartPar
Only inputs that generate the same number of datums can be used
together.  E.g. bond and angle quantities cannot be mixed.

\sphinxlineitem{\sphinxstyleemphasis{Compute property/local does not (yet) work with atom\_style template}}
\sphinxAtStartPar
Self\sphinxhyphen{}explanatory.

\sphinxlineitem{\sphinxstyleemphasis{Compute property/local for property that is not allocated}}
\sphinxAtStartPar
Self\sphinxhyphen{}explanatory.

\sphinxlineitem{\sphinxstyleemphasis{Compute rdf requires a pair style be defined}}
\sphinxAtStartPar
Self\sphinxhyphen{}explanatory.

\sphinxlineitem{\sphinxstyleemphasis{Compute reduce compute array is accessed out\sphinxhyphen{}of\sphinxhyphen{}range}}
\sphinxAtStartPar
An index for the array is out of bounds.

\sphinxlineitem{\sphinxstyleemphasis{Compute reduce compute calculates global values}}
\sphinxAtStartPar
A compute that calculates peratom or local values is required.

\sphinxlineitem{\sphinxstyleemphasis{Compute reduce compute does not calculate a local array}}
\sphinxAtStartPar
Self\sphinxhyphen{}explanatory.

\sphinxlineitem{\sphinxstyleemphasis{Compute reduce compute does not calculate a local vector}}
\sphinxAtStartPar
Self\sphinxhyphen{}explanatory.

\sphinxlineitem{\sphinxstyleemphasis{Compute reduce compute does not calculate a per\sphinxhyphen{}atom array}}
\sphinxAtStartPar
Self\sphinxhyphen{}explanatory.

\sphinxlineitem{\sphinxstyleemphasis{Compute reduce compute does not calculate a per\sphinxhyphen{}atom vector}}
\sphinxAtStartPar
Self\sphinxhyphen{}explanatory.

\sphinxlineitem{\sphinxstyleemphasis{Compute reduce fix array is accessed out\sphinxhyphen{}of\sphinxhyphen{}range}}
\sphinxAtStartPar
An index for the array is out of bounds.

\sphinxlineitem{\sphinxstyleemphasis{Compute reduce fix calculates global values}}
\sphinxAtStartPar
A fix that calculates peratom or local values is required.

\sphinxlineitem{\sphinxstyleemphasis{Compute reduce fix does not calculate a local array}}
\sphinxAtStartPar
Self\sphinxhyphen{}explanatory.

\sphinxlineitem{\sphinxstyleemphasis{Compute reduce fix does not calculate a local vector}}
\sphinxAtStartPar
Self\sphinxhyphen{}explanatory.

\sphinxlineitem{\sphinxstyleemphasis{Compute reduce fix does not calculate a per\sphinxhyphen{}atom array}}
\sphinxAtStartPar
Self\sphinxhyphen{}explanatory.

\sphinxlineitem{\sphinxstyleemphasis{Compute reduce fix does not calculate a per\sphinxhyphen{}atom vector}}
\sphinxAtStartPar
Self\sphinxhyphen{}explanatory.

\sphinxlineitem{\sphinxstyleemphasis{Compute reduce replace requires min or max mode}}
\sphinxAtStartPar
Self\sphinxhyphen{}explanatory.

\sphinxlineitem{\sphinxstyleemphasis{Compute reduce variable is not atom\sphinxhyphen{}style variable}}
\sphinxAtStartPar
Self\sphinxhyphen{}explanatory.

\sphinxlineitem{\sphinxstyleemphasis{Compute slice compute array is accessed out\sphinxhyphen{}of\sphinxhyphen{}range}}
\sphinxAtStartPar
An index for the array is out of bounds.

\sphinxlineitem{\sphinxstyleemphasis{Compute slice compute does not calculate a global array}}
\sphinxAtStartPar
Self\sphinxhyphen{}explanatory.

\sphinxlineitem{\sphinxstyleemphasis{Compute slice compute does not calculate a global vector}}
\sphinxAtStartPar
Self\sphinxhyphen{}explanatory.

\sphinxlineitem{\sphinxstyleemphasis{Compute slice compute does not calculate global vector or array}}
\sphinxAtStartPar
Self\sphinxhyphen{}explanatory.

\sphinxlineitem{\sphinxstyleemphasis{Compute slice compute vector is accessed out\sphinxhyphen{}of\sphinxhyphen{}range}}
\sphinxAtStartPar
The index for the vector is out of bounds.

\sphinxlineitem{\sphinxstyleemphasis{Compute slice fix array is accessed out\sphinxhyphen{}of\sphinxhyphen{}range}}
\sphinxAtStartPar
An index for the array is out of bounds.

\sphinxlineitem{\sphinxstyleemphasis{Compute slice fix does not calculate a global array}}
\sphinxAtStartPar
Self\sphinxhyphen{}explanatory.

\sphinxlineitem{\sphinxstyleemphasis{Compute slice fix does not calculate a global vector}}
\sphinxAtStartPar
Self\sphinxhyphen{}explanatory.

\sphinxlineitem{\sphinxstyleemphasis{Compute slice fix does not calculate global vector or array}}
\sphinxAtStartPar
Self\sphinxhyphen{}explanatory.

\sphinxlineitem{\sphinxstyleemphasis{Compute slice fix vector is accessed out\sphinxhyphen{}of\sphinxhyphen{}range}}
\sphinxAtStartPar
The index for the vector is out of bounds.

\sphinxlineitem{\sphinxstyleemphasis{Compute sna/atom cutoff is longer than pairwise cutoff}}
\sphinxAtStartPar
Self\sphinxhyphen{}explanatory.

\sphinxlineitem{\sphinxstyleemphasis{Compute sna/atom requires a pair style be defined}}
\sphinxAtStartPar
Self\sphinxhyphen{}explanatory.

\sphinxlineitem{\sphinxstyleemphasis{Compute snad/atom cutoff is longer than pairwise cutoff}}
\sphinxAtStartPar
Self\sphinxhyphen{}explanatory.

\sphinxlineitem{\sphinxstyleemphasis{Compute snad/atom requires a pair style be defined}}
\sphinxAtStartPar
Self\sphinxhyphen{}explanatory.

\sphinxlineitem{\sphinxstyleemphasis{Compute snav/atom cutoff is longer than pairwise cutoff}}
\sphinxAtStartPar
Self\sphinxhyphen{}explanatory.

\sphinxlineitem{\sphinxstyleemphasis{Compute snav/atom requires a pair style be defined}}
\sphinxAtStartPar
Self\sphinxhyphen{}explanatory.

\sphinxlineitem{\sphinxstyleemphasis{Compute stress/atom temperature ID does not compute temperature}}
\sphinxAtStartPar
The specified compute must compute temperature.

\sphinxlineitem{\sphinxstyleemphasis{Compute temp/asphere requires atom style ellipsoid}}
\sphinxAtStartPar
Self\sphinxhyphen{}explanatory.

\sphinxlineitem{\sphinxstyleemphasis{Compute temp/asphere requires extended particles}}
\sphinxAtStartPar
This compute cannot be used with point particles.

\sphinxlineitem{\sphinxstyleemphasis{Compute temp/body requires atom style body}}
\sphinxAtStartPar
Self\sphinxhyphen{}explanatory.

\sphinxlineitem{\sphinxstyleemphasis{Compute temp/body requires bodies}}
\sphinxAtStartPar
This compute can only be applied to body particles.

\sphinxlineitem{\sphinxstyleemphasis{Compute temp/chunk does not use chunk/atom compute}}
\sphinxAtStartPar
The style of the specified compute is not chunk/atom.

\sphinxlineitem{\sphinxstyleemphasis{Compute temp/cs requires ghost atoms store velocity}}
\sphinxAtStartPar
Use the comm\_modify vel yes command to enable this.

\sphinxlineitem{\sphinxstyleemphasis{Compute temp/cs used when bonds are not allowed}}
\sphinxAtStartPar
This compute only works on pairs of bonded particles.

\sphinxlineitem{\sphinxstyleemphasis{Compute temp/partial cannot use vz for 2d systemx}}
\sphinxAtStartPar
Self\sphinxhyphen{}explanatory.

\sphinxlineitem{\sphinxstyleemphasis{Compute temp/profile cannot bin z for 2d systems}}
\sphinxAtStartPar
Self\sphinxhyphen{}explanatory.

\sphinxlineitem{\sphinxstyleemphasis{Compute temp/profile cannot use vz for 2d systemx}}
\sphinxAtStartPar
Self\sphinxhyphen{}explanatory.

\sphinxlineitem{\sphinxstyleemphasis{Compute temp/sphere requires atom style sphere}}
\sphinxAtStartPar
Self\sphinxhyphen{}explanatory.

\sphinxlineitem{\sphinxstyleemphasis{Compute ti kspace style does not exist}}
\sphinxAtStartPar
Self\sphinxhyphen{}explanatory.

\sphinxlineitem{\sphinxstyleemphasis{Compute ti pair style does not exist}}
\sphinxAtStartPar
Self\sphinxhyphen{}explanatory.

\sphinxlineitem{\sphinxstyleemphasis{Compute ti tail when pair style does not compute tail corrections}}
\sphinxAtStartPar
Self\sphinxhyphen{}explanatory.

\sphinxlineitem{\sphinxstyleemphasis{Compute torque/chunk does not use chunk/atom compute}}
\sphinxAtStartPar
The style of the specified compute is not chunk/atom.

\sphinxlineitem{\sphinxstyleemphasis{Compute used in dump between runs is not current}}
\sphinxAtStartPar
The compute was not invoked on the current timestep, therefore it
cannot be used in a dump between runs.

\sphinxlineitem{\sphinxstyleemphasis{Compute used in variable between runs is not current}}
\sphinxAtStartPar
Computes cannot be invoked by a variable in between runs.  Thus they
must have been evaluated on the last timestep of the previous run in
order for their value(s) to be accessed.  See the page for the
variable command for more info.

\sphinxlineitem{\sphinxstyleemphasis{Compute used in variable thermo keyword between runs is not current}}
\sphinxAtStartPar
Some thermo keywords rely on a compute to calculate their value(s).
Computes cannot be invoked by a variable in between runs.  Thus they
must have been evaluated on the last timestep of the previous run in
order for their value(s) to be accessed.  See the page for the
variable command for more info.

\sphinxlineitem{\sphinxstyleemphasis{Compute vcm/chunk does not use chunk/atom compute}}
\sphinxAtStartPar
The style of the specified compute is not chunk/atom.

\sphinxlineitem{\sphinxstyleemphasis{Computed temperature for fix temp/berendsen cannot be 0.0}}
\sphinxAtStartPar
Self\sphinxhyphen{}explanatory.

\sphinxlineitem{\sphinxstyleemphasis{Computed temperature for fix temp/rescale cannot be 0.0}}
\sphinxAtStartPar
Cannot rescale the temperature to a new value if the current
temperature is 0.0.

\sphinxlineitem{\sphinxstyleemphasis{Core/shell partner atom not found}}
\sphinxAtStartPar
Could not find one of the atoms in the bond pair.

\sphinxlineitem{\sphinxstyleemphasis{Core/shell partners were not all found}}
\sphinxAtStartPar
Could not find or more atoms in the bond pairs.

\sphinxlineitem{\sphinxstyleemphasis{Could not adjust g\_ewald\_6}}
\sphinxAtStartPar
The Newton\sphinxhyphen{}Raphson solver failed to converge to a good value for
g\_ewald.  This error should not occur for typical problems.  Please
send an email to the developers.

\sphinxlineitem{\sphinxstyleemphasis{Could not compute g\_ewald}}
\sphinxAtStartPar
The Newton\sphinxhyphen{}Raphson solver failed to converge to a good value for
g\_ewald.  This error should not occur for typical problems.  Please
send an email to the developers.

\sphinxlineitem{\sphinxstyleemphasis{Could not compute grid size}}
\sphinxAtStartPar
The code is unable to compute a grid size consistent with the desired
accuracy.  This error should not occur for typical problems.  Please
send an email to the developers.

\sphinxlineitem{\sphinxstyleemphasis{Could not compute grid size for Coulomb interaction}}
\sphinxAtStartPar
The code is unable to compute a grid size consistent with the desired
accuracy.  This error should not occur for typical problems.  Please
send an email to the developers.

\sphinxlineitem{\sphinxstyleemphasis{Could not compute grid size for Dispersion}}
\sphinxAtStartPar
The code is unable to compute a grid size consistent with the desired
accuracy.  This error should not occur for typical problems.  Please
send an email to the developers.

\sphinxlineitem{\sphinxstyleemphasis{Could not create 3d FFT plan}}
\sphinxAtStartPar
The FFT setup for the PPPM solver failed, typically due
to lack of memory.  This is an unusual error.  Check the
size of the FFT grid you are requesting.

\sphinxlineitem{\sphinxstyleemphasis{Could not create 3d grid of processors}}
\sphinxAtStartPar
The specified constraints did not allow a Px by Py by Pz grid to be
created where Px * Py * Pz = P = total number of processors.

\sphinxlineitem{\sphinxstyleemphasis{Could not create 3d remap plan}}
\sphinxAtStartPar
The FFT setup in pppm failed.

\sphinxlineitem{\sphinxstyleemphasis{Could not create Python function arguments}}
\sphinxAtStartPar
This is an internal Python error, possibly because the number
of inputs to the function is too large.

\sphinxlineitem{\sphinxstyleemphasis{Could not create numa grid of processors}}
\sphinxAtStartPar
The specified constraints did not allow this style of grid to be
created.  Usually this is because the total processor count is not a
multiple of the cores/node or the user specified processor count is \textgreater{}
1 in one of the dimensions.

\sphinxlineitem{\sphinxstyleemphasis{Could not create twolevel 3d grid of processors}}
\sphinxAtStartPar
The specified constraints did not allow this style of grid to be
created.

\sphinxlineitem{\sphinxstyleemphasis{Could not evaluate Python function input variable}}
\sphinxAtStartPar
Self\sphinxhyphen{}explanatory.

\sphinxlineitem{\sphinxstyleemphasis{Could not find Python function}}
\sphinxAtStartPar
The provided Python code was run successfully, but it not
define a callable function with the required name.

\sphinxlineitem{\sphinxstyleemphasis{Could not find atom\_modify first group ID}}
\sphinxAtStartPar
Self\sphinxhyphen{}explanatory.

\sphinxlineitem{\sphinxstyleemphasis{Could not find change\_box group ID}}
\sphinxAtStartPar
Group ID used in the change\_box command does not exist.

\sphinxlineitem{\sphinxstyleemphasis{Could not find compute ID for PRD}}
\sphinxAtStartPar
Self\sphinxhyphen{}explanatory.

\sphinxlineitem{\sphinxstyleemphasis{Could not find compute ID for TAD}}
\sphinxAtStartPar
Self\sphinxhyphen{}explanatory.

\sphinxlineitem{\sphinxstyleemphasis{Could not find compute ID for temperature bias}}
\sphinxAtStartPar
Self\sphinxhyphen{}explanatory.

\sphinxlineitem{\sphinxstyleemphasis{Could not find compute ID to delete}}
\sphinxAtStartPar
Self\sphinxhyphen{}explanatory.

\sphinxlineitem{\sphinxstyleemphasis{Could not find compute displace/atom fix ID}}
\sphinxAtStartPar
Self\sphinxhyphen{}explanatory.

\sphinxlineitem{\sphinxstyleemphasis{Could not find compute event/displace fix ID}}
\sphinxAtStartPar
Self\sphinxhyphen{}explanatory.

\sphinxlineitem{\sphinxstyleemphasis{Could not find compute group ID}}
\sphinxAtStartPar
Self\sphinxhyphen{}explanatory.

\sphinxlineitem{\sphinxstyleemphasis{Could not find compute heat/flux compute ID}}
\sphinxAtStartPar
Self\sphinxhyphen{}explanatory.

\sphinxlineitem{\sphinxstyleemphasis{Could not find compute msd fix ID}}
\sphinxAtStartPar
Self\sphinxhyphen{}explanatory.

\sphinxlineitem{\sphinxstyleemphasis{Could not find compute msd/chunk fix ID}}
\sphinxAtStartPar
The compute creates an internal fix, which has been deleted.

\sphinxlineitem{\sphinxstyleemphasis{Could not find compute pressure temperature ID}}
\sphinxAtStartPar
The compute ID for calculating temperature does not exist.

\sphinxlineitem{\sphinxstyleemphasis{Could not find compute stress/atom temperature ID}}
\sphinxAtStartPar
Self\sphinxhyphen{}explanatory.

\sphinxlineitem{\sphinxstyleemphasis{Could not find compute vacf fix ID}}
\sphinxAtStartPar
Self\sphinxhyphen{}explanatory.

\sphinxlineitem{\sphinxstyleemphasis{Could not find compute/voronoi surface group ID}}
\sphinxAtStartPar
Self\sphinxhyphen{}explanatory.

\sphinxlineitem{\sphinxstyleemphasis{Could not find compute\_modify ID}}
\sphinxAtStartPar
Self\sphinxhyphen{}explanatory.

\sphinxlineitem{\sphinxstyleemphasis{Could not find custom per\sphinxhyphen{}atom property ID}}
\sphinxAtStartPar
Self\sphinxhyphen{}explanatory.

\sphinxlineitem{\sphinxstyleemphasis{Could not find delete\_atoms group ID}}
\sphinxAtStartPar
Group ID used in the delete\_atoms command does not exist.

\sphinxlineitem{\sphinxstyleemphasis{Could not find delete\_atoms region ID}}
\sphinxAtStartPar
Region ID used in the delete\_atoms command does not exist.

\sphinxlineitem{\sphinxstyleemphasis{Could not find displace\_atoms group ID}}
\sphinxAtStartPar
Group ID used in the displace\_atoms command does not exist.

\sphinxlineitem{\sphinxstyleemphasis{Could not find dump custom compute ID}}
\sphinxAtStartPar
Self\sphinxhyphen{}explanatory.

\sphinxlineitem{\sphinxstyleemphasis{Could not find dump custom fix ID}}
\sphinxAtStartPar
Self\sphinxhyphen{}explanatory.

\sphinxlineitem{\sphinxstyleemphasis{Could not find dump custom variable name}}
\sphinxAtStartPar
Self\sphinxhyphen{}explanatory.

\sphinxlineitem{\sphinxstyleemphasis{Could not find dump group ID}}
\sphinxAtStartPar
A group ID used in the dump command does not exist.

\sphinxlineitem{\sphinxstyleemphasis{Could not find dump local compute ID}}
\sphinxAtStartPar
Self\sphinxhyphen{}explanatory.

\sphinxlineitem{\sphinxstyleemphasis{Could not find dump local fix ID}}
\sphinxAtStartPar
Self\sphinxhyphen{}explanatory.

\sphinxlineitem{\sphinxstyleemphasis{Could not find dump modify compute ID}}
\sphinxAtStartPar
Self\sphinxhyphen{}explanatory.

\sphinxlineitem{\sphinxstyleemphasis{Could not find dump modify custom atom floating point property ID}}
\sphinxAtStartPar
Self\sphinxhyphen{}explanatory.

\sphinxlineitem{\sphinxstyleemphasis{Could not find dump modify custom atom integer property ID}}
\sphinxAtStartPar
Self\sphinxhyphen{}explanatory.

\sphinxlineitem{\sphinxstyleemphasis{Could not find dump modify fix ID}}
\sphinxAtStartPar
Self\sphinxhyphen{}explanatory.

\sphinxlineitem{\sphinxstyleemphasis{Could not find dump modify variable name}}
\sphinxAtStartPar
Self\sphinxhyphen{}explanatory.

\sphinxlineitem{\sphinxstyleemphasis{Could not find fix ID to delete}}
\sphinxAtStartPar
Self\sphinxhyphen{}explanatory.

\sphinxlineitem{\sphinxstyleemphasis{Could not find fix adapt storage fix ID}}
\sphinxAtStartPar
This should not happen unless you explicitly deleted
a secondary fix that fix adapt created internally.

\sphinxlineitem{\sphinxstyleemphasis{Could not find fix halt variable name}}
\sphinxAtStartPar
Self\sphinxhyphen{}explanatory.

\sphinxlineitem{\sphinxstyleemphasis{Could not find fix gcmc exclusion group ID}}
\sphinxAtStartPar
Self\sphinxhyphen{}explanatory.

\sphinxlineitem{\sphinxstyleemphasis{Could not find fix gcmc rotation group ID}}
\sphinxAtStartPar
Self\sphinxhyphen{}explanatory.

\sphinxlineitem{\sphinxstyleemphasis{Could not find fix group ID}}
\sphinxAtStartPar
A group ID used in the fix command does not exist.

\sphinxlineitem{\sphinxstyleemphasis{Could not find fix msst compute ID}}
\sphinxAtStartPar
Self\sphinxhyphen{}explanatory.

\sphinxlineitem{\sphinxstyleemphasis{Could not find fix poems group ID}}
\sphinxAtStartPar
A group ID used in the fix poems command does not exist.

\sphinxlineitem{\sphinxstyleemphasis{Could not find fix recenter group ID}}
\sphinxAtStartPar
A group ID used in the fix recenter command does not exist.

\sphinxlineitem{\sphinxstyleemphasis{Could not find fix rigid group ID}}
\sphinxAtStartPar
A group ID used in the fix rigid command does not exist.

\sphinxlineitem{\sphinxstyleemphasis{Could not find fix srd group ID}}
\sphinxAtStartPar
Self\sphinxhyphen{}explanatory.

\sphinxlineitem{\sphinxstyleemphasis{Could not find fix\_modify ID}}
\sphinxAtStartPar
A fix ID used in the fix\_modify command does not exist.

\sphinxlineitem{\sphinxstyleemphasis{Could not find fix\_modify pressure ID}}
\sphinxAtStartPar
The compute ID for computing pressure does not exist.

\sphinxlineitem{\sphinxstyleemphasis{Could not find fix\_modify temperature ID}}
\sphinxAtStartPar
The compute ID for computing temperature does not exist.

\sphinxlineitem{\sphinxstyleemphasis{Could not find group clear group ID}}
\sphinxAtStartPar
Self\sphinxhyphen{}explanatory.

\sphinxlineitem{\sphinxstyleemphasis{Could not find group delete group ID}}
\sphinxAtStartPar
Self\sphinxhyphen{}explanatory.

\sphinxlineitem{\sphinxstyleemphasis{Could not find pair fix ID}}
\sphinxAtStartPar
A fix is created internally by the pair style to store shear
history information.  You cannot delete it.

\sphinxlineitem{\sphinxstyleemphasis{Could not find set group ID}}
\sphinxAtStartPar
Group ID specified in set command does not exist.

\sphinxlineitem{\sphinxstyleemphasis{Could not find specified fix gcmc group ID}}
\sphinxAtStartPar
Self\sphinxhyphen{}explanatory.

\sphinxlineitem{\sphinxstyleemphasis{Could not find thermo compute ID}}
\sphinxAtStartPar
Compute ID specified in thermo\_style command does not exist.

\sphinxlineitem{\sphinxstyleemphasis{Could not find thermo custom compute ID}}
\sphinxAtStartPar
The compute ID needed by thermo style custom to compute a requested
quantity does not exist.

\sphinxlineitem{\sphinxstyleemphasis{Could not find thermo custom fix ID}}
\sphinxAtStartPar
The fix ID needed by thermo style custom to compute a requested
quantity does not exist.

\sphinxlineitem{\sphinxstyleemphasis{Could not find thermo custom variable name}}
\sphinxAtStartPar
Self\sphinxhyphen{}explanatory.

\sphinxlineitem{\sphinxstyleemphasis{Could not find thermo fix ID}}
\sphinxAtStartPar
Fix ID specified in thermo\_style command does not exist.

\sphinxlineitem{\sphinxstyleemphasis{Could not find thermo variable name}}
\sphinxAtStartPar
Self\sphinxhyphen{}explanatory.

\sphinxlineitem{\sphinxstyleemphasis{Could not find thermo\_modify pressure ID}}
\sphinxAtStartPar
The compute ID needed by thermo style custom to compute pressure does
not exist.

\sphinxlineitem{\sphinxstyleemphasis{Could not find thermo\_modify temperature ID}}
\sphinxAtStartPar
The compute ID needed by thermo style custom to compute temperature does
not exist.

\sphinxlineitem{\sphinxstyleemphasis{Could not find undump ID}}
\sphinxAtStartPar
A dump ID used in the undump command does not exist.

\sphinxlineitem{\sphinxstyleemphasis{Could not find velocity group ID}}
\sphinxAtStartPar
A group ID used in the velocity command does not exist.

\sphinxlineitem{\sphinxstyleemphasis{Could not find velocity temperature ID}}
\sphinxAtStartPar
The compute ID needed by the velocity command to compute temperature
does not exist.

\sphinxlineitem{\sphinxstyleemphasis{Could not find/initialize a specified accelerator device}}
\sphinxAtStartPar
Could not initialize at least one of the devices specified for the gpu
package

\sphinxlineitem{\sphinxstyleemphasis{Could not grab element entry from EIM potential file}}
\sphinxAtStartPar
Self\sphinxhyphen{}explanatory

\sphinxlineitem{\sphinxstyleemphasis{Could not grab global entry from EIM potential file}}
\sphinxAtStartPar
Self\sphinxhyphen{}explanatory.

\sphinxlineitem{\sphinxstyleemphasis{Could not grab pair entry from EIM potential file}}
\sphinxAtStartPar
Self\sphinxhyphen{}explanatory.

\sphinxlineitem{\sphinxstyleemphasis{Could not initialize embedded Python}}
\sphinxAtStartPar
The main module in Python was not accessible.

\sphinxlineitem{\sphinxstyleemphasis{Could not open Python file}}
\sphinxAtStartPar
The specified file of Python code cannot be opened.  Check that the
path and name are correct.

\sphinxlineitem{\sphinxstyleemphasis{Could not process Python file}}
\sphinxAtStartPar
The Python code in the specified file was not run successfully by
Python, probably due to errors in the Python code.

\sphinxlineitem{\sphinxstyleemphasis{Could not process Python string}}
\sphinxAtStartPar
The Python code in the here string was not run successfully by Python,
probably due to errors in the Python code.

\sphinxlineitem{\sphinxstyleemphasis{Coulomb PPPMDisp order has been reduced below minorder}}
\sphinxAtStartPar
The default minimum order is 2.  This can be reset by the
kspace\_modify minorder command.

\sphinxlineitem{\sphinxstyleemphasis{Coulombic cutoff not supported in pair\_style buck/long/coul/coul}}
\sphinxAtStartPar
Must use long\sphinxhyphen{}range Coulombic interactions.

\sphinxlineitem{\sphinxstyleemphasis{Coulombic cutoff not supported in pair\_style lj/long/coul/long}}
\sphinxAtStartPar
Must use long\sphinxhyphen{}range Coulombic interactions.

\sphinxlineitem{\sphinxstyleemphasis{Coulombic cutoff not supported in pair\_style lj/long/tip4p/long}}
\sphinxAtStartPar
Must use long\sphinxhyphen{}range Coulombic interactions.

\sphinxlineitem{\sphinxstyleemphasis{Coulombic cutoffs of pair hybrid sub\sphinxhyphen{}styles do not match}}
\sphinxAtStartPar
If using a Kspace solver, all Coulombic cutoffs of long pair styles must
be the same.

\sphinxlineitem{\sphinxstyleemphasis{Coulombic cut not supported in pair\_style lj/long/dipole/long}}
\sphinxAtStartPar
Must use long\sphinxhyphen{}range Coulombic interactions.

\sphinxlineitem{\sphinxstyleemphasis{Cound not find dump\_modify ID}}
\sphinxAtStartPar
Self\sphinxhyphen{}explanatory.

\sphinxlineitem{\sphinxstyleemphasis{Create\_atoms command before simulation box is defined}}
\sphinxAtStartPar
The create\_atoms command cannot be used before a read\_data,
read\_restart, or create\_box command.

\sphinxlineitem{\sphinxstyleemphasis{Create\_atoms molecule has atom IDs, but system does not}}
\sphinxAtStartPar
The atom\_style id command can be used to force atom IDs to be stored.

\sphinxlineitem{\sphinxstyleemphasis{Create\_atoms molecule must have atom types}}
\sphinxAtStartPar
The defined molecule does not specify atom types.

\sphinxlineitem{\sphinxstyleemphasis{Create\_atoms molecule must have coordinates}}
\sphinxAtStartPar
The defined molecule does not specify coordinates.

\sphinxlineitem{\sphinxstyleemphasis{Create\_atoms region ID does not exist}}
\sphinxAtStartPar
A region ID used in the create\_atoms command does not exist.

\sphinxlineitem{\sphinxstyleemphasis{Create\_bonds command before simulation box is defined}}
\sphinxAtStartPar
Self\sphinxhyphen{}explanatory.

\sphinxlineitem{\sphinxstyleemphasis{Create\_bonds command requires no kspace\_style be defined}}
\sphinxAtStartPar
This is so that atom pairs that are already bonded to not appear
in the neighbor list.

\sphinxlineitem{\sphinxstyleemphasis{Create\_bonds command requires special\_bonds 1\sphinxhyphen{}2 weights be 0.0}}
\sphinxAtStartPar
This is so that atom pairs that are already bonded to not appear in
the neighbor list.

\sphinxlineitem{\sphinxstyleemphasis{Create\_bonds max distance \textgreater{} neighbor cutoff}}
\sphinxAtStartPar
Can only create bonds for atom pairs that will be in neighbor list.

\sphinxlineitem{\sphinxstyleemphasis{Create\_bonds requires a pair style be defined}}
\sphinxAtStartPar
Self\sphinxhyphen{}explanatory.

\sphinxlineitem{\sphinxstyleemphasis{Create\_box region ID does not exist}}
\sphinxAtStartPar
Self\sphinxhyphen{}explanatory.

\sphinxlineitem{\sphinxstyleemphasis{Create\_box region does not support a bounding box}}
\sphinxAtStartPar
Not all regions represent bounded volumes.  You cannot use
such a region with the create\_box command.

\sphinxlineitem{\sphinxstyleemphasis{Custom floating point vector for fix store/state does not exist}}
\sphinxAtStartPar
The command is accessing a vector added by the fix property/atom
command, that does not exist.

\sphinxlineitem{\sphinxstyleemphasis{Custom integer vector for fix store/state does not exist}}
\sphinxAtStartPar
The command is accessing a vector added by the fix property/atom
command, that does not exist.

\sphinxlineitem{\sphinxstyleemphasis{Custom per\sphinxhyphen{}atom property ID is not floating point}}
\sphinxAtStartPar
Self\sphinxhyphen{}explanatory.

\sphinxlineitem{\sphinxstyleemphasis{Custom per\sphinxhyphen{}atom property ID is not integer}}
\sphinxAtStartPar
Self\sphinxhyphen{}explanatory.

\sphinxlineitem{\sphinxstyleemphasis{Cut\sphinxhyphen{}offs missing in pair\_style lj/long/dipole/long}}
\sphinxAtStartPar
Self\sphinxhyphen{}explanatory.

\sphinxlineitem{\sphinxstyleemphasis{Cutoffs missing in pair\_style buck/long/coul/long}}
\sphinxAtStartPar
Self\sphinxhyphen{}explanatory.

\sphinxlineitem{\sphinxstyleemphasis{Cutoffs missing in pair\_style lj/long/coul/long}}
\sphinxAtStartPar
Self\sphinxhyphen{}explanatory.

\sphinxlineitem{\sphinxstyleemphasis{Cyclic loop in joint connections}}
\sphinxAtStartPar
Fix poems cannot (yet) work with coupled bodies whose joints connect
the bodies in a ring (or cycle).

\sphinxlineitem{\sphinxstyleemphasis{Degenerate lattice primitive vectors}}
\sphinxAtStartPar
Invalid set of 3 lattice vectors for lattice command.

\sphinxlineitem{\sphinxstyleemphasis{Delete region ID does not exist}}
\sphinxAtStartPar
Self\sphinxhyphen{}explanatory.

\sphinxlineitem{\sphinxstyleemphasis{Delete\_atoms command before simulation box is defined}}
\sphinxAtStartPar
The delete\_atoms command cannot be used before a read\_data,
read\_restart, or create\_box command.

\sphinxlineitem{\sphinxstyleemphasis{Delete\_atoms cutoff \textgreater{} max neighbor cutoff}}
\sphinxAtStartPar
Can only delete atoms in atom pairs that will be in neighbor list.

\sphinxlineitem{\sphinxstyleemphasis{Delete\_atoms mol yes requires atom attribute molecule}}
\sphinxAtStartPar
Cannot use this option with a non\sphinxhyphen{}molecular system.

\sphinxlineitem{\sphinxstyleemphasis{Delete\_atoms requires a pair style be defined}}
\sphinxAtStartPar
This is because atom deletion within a cutoff uses a pairwise
neighbor list.

\sphinxlineitem{\sphinxstyleemphasis{Delete\_bonds command before simulation box is defined}}
\sphinxAtStartPar
The delete\_bonds command cannot be used before a read\_data,
read\_restart, or create\_box command.

\sphinxlineitem{\sphinxstyleemphasis{Delete\_bonds command with no atoms existing}}
\sphinxAtStartPar
No atoms are yet defined so the delete\_bonds command cannot be used.

\sphinxlineitem{\sphinxstyleemphasis{Deposition region extends outside simulation box}}
\sphinxAtStartPar
Self\sphinxhyphen{}explanatory.

\sphinxlineitem{\sphinxstyleemphasis{Did not assign all atoms correctly}}
\sphinxAtStartPar
Atoms read in from a data file were not assigned correctly to
processors.  This is likely due to some atom coordinates being
outside a non\sphinxhyphen{}periodic simulation box.

\sphinxlineitem{\sphinxstyleemphasis{Did not assign all restart atoms correctly}}
\sphinxAtStartPar
Atoms read in from the restart file were not assigned correctly to
processors.  This is likely due to some atom coordinates being outside
a non\sphinxhyphen{}periodic simulation box.  Normally this should not happen.  You
may wish to use the “remap” option on the read\_restart command to see
if this helps.

\sphinxlineitem{\sphinxstyleemphasis{Did not find all elements in MEAM library file}}
\sphinxAtStartPar
Some requested elements were not found in the MEAM file. Check spelling etc.

\sphinxlineitem{\sphinxstyleemphasis{Did not find fix shake partner info}}
\sphinxAtStartPar
Could not find bond partners implied by fix shake command.  This error
can be triggered if the delete\_bonds command was used before fix
shake, and it removed bonds without resetting the 1\sphinxhyphen{}2, 1\sphinxhyphen{}3, 1\sphinxhyphen{}4
weighting list via the special keyword.

\sphinxlineitem{\sphinxstyleemphasis{Did not find keyword in table file}}
\sphinxAtStartPar
Keyword used in pair\_coeff command was not found in table file.

\sphinxlineitem{\sphinxstyleemphasis{Did not set pressure for fix rigid/nph}}
\sphinxAtStartPar
The press keyword must be specified.

\sphinxlineitem{\sphinxstyleemphasis{Did not set temp for fix rigid/nvt/small}}
\sphinxAtStartPar
Self\sphinxhyphen{}explanatory.

\sphinxlineitem{\sphinxstyleemphasis{Did not set temp or press for fix rigid/npt/small}}
\sphinxAtStartPar
Self\sphinxhyphen{}explanatory.

\sphinxlineitem{\sphinxstyleemphasis{Did not set temperature for fix rigid/nvt}}
\sphinxAtStartPar
The temp keyword must be specified.

\sphinxlineitem{\sphinxstyleemphasis{Did not set temperature or pressure for fix rigid/npt}}
\sphinxAtStartPar
The temp and press keywords must be specified.

\sphinxlineitem{\sphinxstyleemphasis{Dihedral atom missing in delete\_bonds}}
\sphinxAtStartPar
The delete\_bonds command cannot find one or more atoms in a particular
dihedral on a particular processor.  The pairwise cutoff is too short
or the atoms are too far apart to make a valid dihedral.

\sphinxlineitem{\sphinxstyleemphasis{Dihedral atom missing in set command}}
\sphinxAtStartPar
The set command cannot find one or more atoms in a particular dihedral
on a particular processor.  The pairwise cutoff is too short or the
atoms are too far apart to make a valid dihedral.

\sphinxlineitem{\sphinxstyleemphasis{Dihedral atoms \%d \%d \%d \%d missing on proc \%d at step \%ld}}
\sphinxAtStartPar
One or more of 4 atoms needed to compute a particular dihedral are
missing on this processor.  Typically this is because the pairwise
cutoff is set too short or the dihedral has blown apart and an atom is
too far away.

\sphinxlineitem{\sphinxstyleemphasis{Dihedral atoms missing on proc \%d at step \%ld}}
\sphinxAtStartPar
One or more of 4 atoms needed to compute a particular dihedral are
missing on this processor.  Typically this is because the pairwise
cutoff is set too short or the dihedral has blown apart and an atom is
too far away.

\sphinxlineitem{\sphinxstyleemphasis{Dihedral charmm is incompatible with Pair style}}
\sphinxAtStartPar
Dihedral style charmm must be used with a pair style charmm
in order for the 1\sphinxhyphen{}4 epsilon/sigma parameters to be defined.

\sphinxlineitem{\sphinxstyleemphasis{Dihedral coeff for hybrid has invalid style}}
\sphinxAtStartPar
Dihedral style hybrid uses another dihedral style as one of its
coefficients.  The dihedral style used in the dihedral\_coeff command
or read from a restart file is not recognized.

\sphinxlineitem{\sphinxstyleemphasis{Dihedral coeffs are not set}}
\sphinxAtStartPar
No dihedral coefficients have been assigned in the data file or via
the dihedral\_coeff command.

\sphinxlineitem{\sphinxstyleemphasis{Dihedral style hybrid cannot have hybrid as an argument}}
\sphinxAtStartPar
Self\sphinxhyphen{}explanatory.

\sphinxlineitem{\sphinxstyleemphasis{Dihedral style hybrid cannot have none as an argument}}
\sphinxAtStartPar
Self\sphinxhyphen{}explanatory.

\sphinxlineitem{\sphinxstyleemphasis{Dihedral style hybrid cannot use same dihedral style twice}}
\sphinxAtStartPar
Self\sphinxhyphen{}explanatory.

\sphinxlineitem{\sphinxstyleemphasis{Dihedral/improper extent \textgreater{} half of periodic box length}}
\sphinxAtStartPar
This error was detected by the neigh\_modify check yes setting.  It is
an error because the dihedral atoms are so far apart it is ambiguous
how it should be defined.

\sphinxlineitem{\sphinxstyleemphasis{Dihedral\_coeff command before dihedral\_style is defined}}
\sphinxAtStartPar
Coefficients cannot be set in the data file or via the dihedral\_coeff
command until an dihedral\_style has been assigned.

\sphinxlineitem{\sphinxstyleemphasis{Dihedral\_coeff command before simulation box is defined}}
\sphinxAtStartPar
The dihedral\_coeff command cannot be used before a read\_data,
read\_restart, or create\_box command.

\sphinxlineitem{\sphinxstyleemphasis{Dihedral\_coeff command when no dihedrals allowed}}
\sphinxAtStartPar
The chosen atom style does not allow for dihedrals to be defined.

\sphinxlineitem{\sphinxstyleemphasis{Dihedral\_style command when no dihedrals allowed}}
\sphinxAtStartPar
The chosen atom style does not allow for dihedrals to be defined.

\sphinxlineitem{\sphinxstyleemphasis{Dihedrals assigned incorrectly}}
\sphinxAtStartPar
Dihedrals read in from the data file were not assigned correctly to
atoms.  This means there is something invalid about the topology
definitions.

\sphinxlineitem{\sphinxstyleemphasis{Dihedrals defined but no dihedral types}}
\sphinxAtStartPar
The data file header lists dihedrals but no dihedral types.

\sphinxlineitem{\sphinxstyleemphasis{Dimension command after simulation box is defined}}
\sphinxAtStartPar
The dimension command cannot be used after a read\_data,
read\_restart, or create\_box command.

\sphinxlineitem{\sphinxstyleemphasis{Disk limit not supported by OS or illegal path}}
\sphinxAtStartPar
Self\sphinxhyphen{}explanatory.

\sphinxlineitem{\sphinxstyleemphasis{Dispersion PPPMDisp order has been reduced below minorder}}
\sphinxAtStartPar
The default minimum order is 2.  This can be reset by the
kspace\_modify minorder command.

\sphinxlineitem{\sphinxstyleemphasis{Displace\_atoms command before simulation box is defined}}
\sphinxAtStartPar
The displace\_atoms command cannot be used before a read\_data,
read\_restart, or create\_box command.

\sphinxlineitem{\sphinxstyleemphasis{Distance must be \textgreater{} 0 for compute event/displace}}
\sphinxAtStartPar
Self\sphinxhyphen{}explanatory.

\sphinxlineitem{\sphinxstyleemphasis{Divide by 0 in influence function}}
\sphinxAtStartPar
This should not normally occur.  It is likely a problem with your
model.

\sphinxlineitem{\sphinxstyleemphasis{Divide by 0 in influence function of pair peri/lps}}
\sphinxAtStartPar
This should not normally occur.  It is likely a problem with your
model.

\sphinxlineitem{\sphinxstyleemphasis{Divide by 0 in variable formula}}
\sphinxAtStartPar
Self\sphinxhyphen{}explanatory.

\sphinxlineitem{\sphinxstyleemphasis{Domain too large for neighbor bins}}
\sphinxAtStartPar
The domain has become extremely large so that neighbor bins cannot be
used.  Most likely, one or more atoms have been blown out of the
simulation box to a great distance.

\sphinxlineitem{\sphinxstyleemphasis{Double precision is not supported on this accelerator}}
\sphinxAtStartPar
Self\sphinxhyphen{}explanatory

\sphinxlineitem{\sphinxstyleemphasis{Dump atom/gz only writes compressed files}}
\sphinxAtStartPar
The dump atom/gz output file name must have a .gz suffix.

\sphinxlineitem{\sphinxstyleemphasis{Dump cfg arguments can not mix xs|ys|zs with xsu|ysu|zsu}}
\sphinxAtStartPar
Self\sphinxhyphen{}explanatory.

\sphinxlineitem{\sphinxstyleemphasis{Dump cfg arguments must start with ‘mass type xs ys zs’ or ‘mass type xsu ysu zsu’}}
\sphinxAtStartPar
This is a requirement of the CFG output format.  See the dump cfg doc
page for more details.

\sphinxlineitem{\sphinxstyleemphasis{Dump cfg requires one snapshot per file}}
\sphinxAtStartPar
Use the wildcard “*” character in the filename.

\sphinxlineitem{\sphinxstyleemphasis{Dump cfg/gz only writes compressed files}}
\sphinxAtStartPar
The dump cfg/gz output file name must have a .gz suffix.

\sphinxlineitem{\sphinxstyleemphasis{Dump custom and fix not computed at compatible times}}
\sphinxAtStartPar
The fix must produce per\sphinxhyphen{}atom quantities on timesteps that dump custom
needs them.

\sphinxlineitem{\sphinxstyleemphasis{Dump custom compute does not calculate per\sphinxhyphen{}atom array}}
\sphinxAtStartPar
Self\sphinxhyphen{}explanatory.

\sphinxlineitem{\sphinxstyleemphasis{Dump custom compute does not calculate per\sphinxhyphen{}atom vector}}
\sphinxAtStartPar
Self\sphinxhyphen{}explanatory.

\sphinxlineitem{\sphinxstyleemphasis{Dump custom compute does not compute per\sphinxhyphen{}atom info}}
\sphinxAtStartPar
Self\sphinxhyphen{}explanatory.

\sphinxlineitem{\sphinxstyleemphasis{Dump custom compute vector is accessed out\sphinxhyphen{}of\sphinxhyphen{}range}}
\sphinxAtStartPar
Self\sphinxhyphen{}explanatory.

\sphinxlineitem{\sphinxstyleemphasis{Dump custom fix does not compute per\sphinxhyphen{}atom array}}
\sphinxAtStartPar
Self\sphinxhyphen{}explanatory.

\sphinxlineitem{\sphinxstyleemphasis{Dump custom fix does not compute per\sphinxhyphen{}atom info}}
\sphinxAtStartPar
Self\sphinxhyphen{}explanatory.

\sphinxlineitem{\sphinxstyleemphasis{Dump custom fix does not compute per\sphinxhyphen{}atom vector}}
\sphinxAtStartPar
Self\sphinxhyphen{}explanatory.

\sphinxlineitem{\sphinxstyleemphasis{Dump custom fix vector is accessed out\sphinxhyphen{}of\sphinxhyphen{}range}}
\sphinxAtStartPar
Self\sphinxhyphen{}explanatory.

\sphinxlineitem{\sphinxstyleemphasis{Dump custom variable is not atom\sphinxhyphen{}style variable}}
\sphinxAtStartPar
Only atom\sphinxhyphen{}style variables generate per\sphinxhyphen{}atom quantities, needed for
dump output.

\sphinxlineitem{\sphinxstyleemphasis{Dump custom/gz only writes compressed files}}
\sphinxAtStartPar
The dump custom/gz output file name must have a .gz suffix.

\sphinxlineitem{\sphinxstyleemphasis{Dump dcd of non\sphinxhyphen{}matching \# of atoms}}
\sphinxAtStartPar
Every snapshot written by dump dcd must contain the same \# of atoms.

\sphinxlineitem{\sphinxstyleemphasis{Dump dcd requires sorting by atom ID}}
\sphinxAtStartPar
Use the dump\_modify sort command to enable this.

\sphinxlineitem{\sphinxstyleemphasis{Dump every variable returned a bad timestep}}
\sphinxAtStartPar
The variable must return a timestep greater than the current timestep.

\sphinxlineitem{\sphinxstyleemphasis{Dump file MPI\sphinxhyphen{}IO output not allowed with \% in filename}}
\sphinxAtStartPar
This is because a \% signifies one file per processor and MPI\sphinxhyphen{}IO
creates one large file for all processors.

\sphinxlineitem{\sphinxstyleemphasis{Dump file does not contain requested snapshot}}
\sphinxAtStartPar
Self\sphinxhyphen{}explanatory.

\sphinxlineitem{\sphinxstyleemphasis{Dump file is incorrectly formatted}}
\sphinxAtStartPar
Self\sphinxhyphen{}explanatory.

\sphinxlineitem{\sphinxstyleemphasis{Dump image body yes requires atom style body}}
\sphinxAtStartPar
Self\sphinxhyphen{}explanatory.

\sphinxlineitem{\sphinxstyleemphasis{Dump image bond not allowed with no bond types}}
\sphinxAtStartPar
Self\sphinxhyphen{}explanatory.

\sphinxlineitem{\sphinxstyleemphasis{Dump image cannot perform sorting}}
\sphinxAtStartPar
Self\sphinxhyphen{}explanatory.

\sphinxlineitem{\sphinxstyleemphasis{Dump image line requires atom style line}}
\sphinxAtStartPar
Self\sphinxhyphen{}explanatory.

\sphinxlineitem{\sphinxstyleemphasis{Dump image requires one snapshot per file}}
\sphinxAtStartPar
Use a “*” in the filename.

\sphinxlineitem{\sphinxstyleemphasis{Dump image tri requires atom style tri}}
\sphinxAtStartPar
Self\sphinxhyphen{}explanatory.

\sphinxlineitem{\sphinxstyleemphasis{Dump local and fix not computed at compatible times}}
\sphinxAtStartPar
The fix must produce per\sphinxhyphen{}atom quantities on timesteps that dump local
needs them.

\sphinxlineitem{\sphinxstyleemphasis{Dump local attributes contain no compute or fix}}
\sphinxAtStartPar
Self\sphinxhyphen{}explanatory.

\sphinxlineitem{\sphinxstyleemphasis{Dump local cannot sort by atom ID}}
\sphinxAtStartPar
This is because dump local does not really dump per\sphinxhyphen{}atom info.

\sphinxlineitem{\sphinxstyleemphasis{Dump local compute does not calculate local array}}
\sphinxAtStartPar
Self\sphinxhyphen{}explanatory.

\sphinxlineitem{\sphinxstyleemphasis{Dump local compute does not calculate local vector}}
\sphinxAtStartPar
Self\sphinxhyphen{}explanatory.

\sphinxlineitem{\sphinxstyleemphasis{Dump local compute does not compute local info}}
\sphinxAtStartPar
Self\sphinxhyphen{}explanatory.

\sphinxlineitem{\sphinxstyleemphasis{Dump local compute vector is accessed out\sphinxhyphen{}of\sphinxhyphen{}range}}
\sphinxAtStartPar
Self\sphinxhyphen{}explanatory.

\sphinxlineitem{\sphinxstyleemphasis{Dump local count is not consistent across input fields}}
\sphinxAtStartPar
Every column of output must be the same length.

\sphinxlineitem{\sphinxstyleemphasis{Dump local fix does not compute local array}}
\sphinxAtStartPar
Self\sphinxhyphen{}explanatory.

\sphinxlineitem{\sphinxstyleemphasis{Dump local fix does not compute local info}}
\sphinxAtStartPar
Self\sphinxhyphen{}explanatory.

\sphinxlineitem{\sphinxstyleemphasis{Dump local fix does not compute local vector}}
\sphinxAtStartPar
Self\sphinxhyphen{}explanatory.

\sphinxlineitem{\sphinxstyleemphasis{Dump local fix vector is accessed out\sphinxhyphen{}of\sphinxhyphen{}range}}
\sphinxAtStartPar
Self\sphinxhyphen{}explanatory.

\sphinxlineitem{\sphinxstyleemphasis{Dump modify bcolor not allowed with no bond types}}
\sphinxAtStartPar
Self\sphinxhyphen{}explanatory.

\sphinxlineitem{\sphinxstyleemphasis{Dump modify bdiam not allowed with no bond types}}
\sphinxAtStartPar
Self\sphinxhyphen{}explanatory.

\sphinxlineitem{\sphinxstyleemphasis{Dump modify compute ID does not compute per\sphinxhyphen{}atom array}}
\sphinxAtStartPar
Self\sphinxhyphen{}explanatory.

\sphinxlineitem{\sphinxstyleemphasis{Dump modify compute ID does not compute per\sphinxhyphen{}atom info}}
\sphinxAtStartPar
Self\sphinxhyphen{}explanatory.

\sphinxlineitem{\sphinxstyleemphasis{Dump modify compute ID does not compute per\sphinxhyphen{}atom vector}}
\sphinxAtStartPar
Self\sphinxhyphen{}explanatory.

\sphinxlineitem{\sphinxstyleemphasis{Dump modify compute ID vector is not large enough}}
\sphinxAtStartPar
Self\sphinxhyphen{}explanatory.

\sphinxlineitem{\sphinxstyleemphasis{Dump modify element names do not match atom types}}
\sphinxAtStartPar
Number of element names must equal number of atom types.

\sphinxlineitem{\sphinxstyleemphasis{Dump modify fix ID does not compute per\sphinxhyphen{}atom array}}
\sphinxAtStartPar
Self\sphinxhyphen{}explanatory.

\sphinxlineitem{\sphinxstyleemphasis{Dump modify fix ID does not compute per\sphinxhyphen{}atom info}}
\sphinxAtStartPar
Self\sphinxhyphen{}explanatory.

\sphinxlineitem{\sphinxstyleemphasis{Dump modify fix ID does not compute per\sphinxhyphen{}atom vector}}
\sphinxAtStartPar
Self\sphinxhyphen{}explanatory.

\sphinxlineitem{\sphinxstyleemphasis{Dump modify fix ID vector is not large enough}}
\sphinxAtStartPar
Self\sphinxhyphen{}explanatory.

\sphinxlineitem{\sphinxstyleemphasis{Dump modify variable is not atom\sphinxhyphen{}style variable}}
\sphinxAtStartPar
Self\sphinxhyphen{}explanatory.

\sphinxlineitem{\sphinxstyleemphasis{Dump sort column is invalid}}
\sphinxAtStartPar
Self\sphinxhyphen{}explanatory.

\sphinxlineitem{\sphinxstyleemphasis{Dump xtc requires sorting by atom ID}}
\sphinxAtStartPar
Use the dump\_modify sort command to enable this.

\sphinxlineitem{\sphinxstyleemphasis{Dump xyz/gz only writes compressed files}}
\sphinxAtStartPar
The dump xyz/gz output file name must have a .gz suffix.

\sphinxlineitem{\sphinxstyleemphasis{Dump\_modify buffer yes not allowed for this style}}
\sphinxAtStartPar
Self\sphinxhyphen{}explanatory.

\sphinxlineitem{\sphinxstyleemphasis{Dump\_modify format string is too short}}
\sphinxAtStartPar
There are more fields to be dumped in a line of output than your
format string specifies.

\sphinxlineitem{\sphinxstyleemphasis{Dump\_modify region ID does not exist}}
\sphinxAtStartPar
Self\sphinxhyphen{}explanatory.

\sphinxlineitem{\sphinxstyleemphasis{Dumping an atom property that is not allocated}}
\sphinxAtStartPar
The chosen atom style does not define the per\sphinxhyphen{}atom quantity being
dumped.

\sphinxlineitem{\sphinxstyleemphasis{Duplicate atom IDs exist}}
\sphinxAtStartPar
Self\sphinxhyphen{}explanatory.

\sphinxlineitem{\sphinxstyleemphasis{Duplicate fields in read\_dump command}}
\sphinxAtStartPar
Self\sphinxhyphen{}explanatory.

\sphinxlineitem{\sphinxstyleemphasis{Duplicate particle in PeriDynamic bond \sphinxhyphen{} simulation box is too small}}
\sphinxAtStartPar
This is likely because your box length is shorter than 2 times
the bond length.

\sphinxlineitem{\sphinxstyleemphasis{Electronic temperature dropped below zero}}
\sphinxAtStartPar
Something has gone wrong with the fix ttm electron temperature model.

\sphinxlineitem{\sphinxstyleemphasis{Element not defined in potential file}}
\sphinxAtStartPar
The specified element is not in the potential file.

\sphinxlineitem{\sphinxstyleemphasis{Empty brackets in variable}}
\sphinxAtStartPar
There is no variable syntax that uses empty brackets.  Check
the variable doc page.

\sphinxlineitem{\sphinxstyleemphasis{Energy was not tallied on needed timestep}}
\sphinxAtStartPar
You are using a thermo keyword that requires potentials to
have tallied energy, but they did not on this timestep.  See the
variable page for ideas on how to make this work.

\sphinxlineitem{\sphinxstyleemphasis{Epsilon or sigma reference not set by pair style in PPPMDisp}}
\sphinxAtStartPar
Self\sphinxhyphen{}explanatory.

\sphinxlineitem{\sphinxstyleemphasis{Epsilon or sigma reference not set by pair style in ewald/n}}
\sphinxAtStartPar
The pair style is not providing the needed epsilon or sigma values.

\sphinxlineitem{\sphinxstyleemphasis{Error in MEAM parameter file: keyword \%s (further information)}}
\sphinxAtStartPar
Self\sphinxhyphen{}explanatory. Check the parameter file.

\sphinxlineitem{\sphinxstyleemphasis{Error in vdw spline: inner radius \textgreater{} outer radius}}
\sphinxAtStartPar
A pre\sphinxhyphen{}tabulated spline is invalid.  Likely a problem with the
potential parameters.

\sphinxlineitem{\sphinxstyleemphasis{Error writing averaged chunk data}}
\sphinxAtStartPar
Something in the output to the file triggered an error.

\sphinxlineitem{\sphinxstyleemphasis{Error writing file header}}
\sphinxAtStartPar
Something in the output to the file triggered an error.

\sphinxlineitem{\sphinxstyleemphasis{Error writing out correlation data}}
\sphinxAtStartPar
Something in the output to the file triggered an error.

\sphinxlineitem{\sphinxstyleemphasis{Error writing out histogram data}}
\sphinxAtStartPar
Something in the output to the file triggered an error.

\sphinxlineitem{\sphinxstyleemphasis{Error writing out time averaged data}}
\sphinxAtStartPar
Something in the output to the file triggered an error.

\sphinxlineitem{\sphinxstyleemphasis{Failed to allocate \%ld bytes for array \%s}}
\sphinxAtStartPar
Your LAMMPS simulation has run out of memory.  You need to run a
smaller simulation or on more processors.

\sphinxlineitem{\sphinxstyleemphasis{Failed to open FFmpeg pipeline to file \%s}}
\sphinxAtStartPar
The specified file cannot be opened.  Check that the path and name are
correct and writable and that the FFmpeg executable can be found and run.

\sphinxlineitem{\sphinxstyleemphasis{Failed to reallocate \%ld bytes for array \%s}}
\sphinxAtStartPar
Your LAMMPS simulation has run out of memory.  You need to run a
smaller simulation or on more processors.

\sphinxlineitem{\sphinxstyleemphasis{Fewer SRD bins than processors in some dimension}}
\sphinxAtStartPar
This is not allowed.  Make your SRD bin size smaller.

\sphinxlineitem{\sphinxstyleemphasis{File variable could not read value}}
\sphinxAtStartPar
Check the file assigned to the variable.

\sphinxlineitem{\sphinxstyleemphasis{Final box dimension due to fix deform is \textless{} 0.0}}
\sphinxAtStartPar
Self\sphinxhyphen{}explanatory.

\sphinxlineitem{\sphinxstyleemphasis{Fix \%s does not allow use of dynamic group}}
\sphinxAtStartPar
Dynamic groups have not yet been enabled for this fix.

\sphinxlineitem{\sphinxstyleemphasis{Fix ID for compute chunk/atom does not exist}}
\sphinxAtStartPar
Self\sphinxhyphen{}explanatory.

\sphinxlineitem{\sphinxstyleemphasis{Fix ID for compute erotate/rigid does not exist}}
\sphinxAtStartPar
Self\sphinxhyphen{}explanatory.

\sphinxlineitem{\sphinxstyleemphasis{Fix ID for compute ke/rigid does not exist}}
\sphinxAtStartPar
Self\sphinxhyphen{}explanatory.

\sphinxlineitem{\sphinxstyleemphasis{Fix ID for compute reduce does not exist}}
\sphinxAtStartPar
Self\sphinxhyphen{}explanatory.

\sphinxlineitem{\sphinxstyleemphasis{Fix ID for compute slice does not exist}}
\sphinxAtStartPar
Self\sphinxhyphen{}explanatory.

\sphinxlineitem{\sphinxstyleemphasis{Fix ID for fix ave/atom does not exist}}
\sphinxAtStartPar
Self\sphinxhyphen{}explanatory.

\sphinxlineitem{\sphinxstyleemphasis{Fix ID for fix ave/chunk does not exist}}
\sphinxAtStartPar
Self\sphinxhyphen{}explanatory.

\sphinxlineitem{\sphinxstyleemphasis{Fix ID for fix ave/correlate does not exist}}
\sphinxAtStartPar
Self\sphinxhyphen{}explanatory.

\sphinxlineitem{\sphinxstyleemphasis{Fix ID for fix ave/histo does not exist}}
\sphinxAtStartPar
Self\sphinxhyphen{}explanatory.

\sphinxlineitem{\sphinxstyleemphasis{Fix ID for fix ave/time does not exist}}
\sphinxAtStartPar
Self\sphinxhyphen{}explanatory.

\sphinxlineitem{\sphinxstyleemphasis{Fix ID for fix store/state does not exist}}
\sphinxAtStartPar
Self\sphinxhyphen{}explanatory

\sphinxlineitem{\sphinxstyleemphasis{Fix ID for fix vector does not exist}}
\sphinxAtStartPar
Self\sphinxhyphen{}explanatory.

\sphinxlineitem{\sphinxstyleemphasis{Fix ID for read\_data does not exist}}
\sphinxAtStartPar
Self\sphinxhyphen{}explanatory.

\sphinxlineitem{\sphinxstyleemphasis{Fix ID for velocity does not exist}}
\sphinxAtStartPar
Self\sphinxhyphen{}explanatory.

\sphinxlineitem{\sphinxstyleemphasis{Fix ID must be alphanumeric or underscore characters}}
\sphinxAtStartPar
Self\sphinxhyphen{}explanatory.

\sphinxlineitem{\sphinxstyleemphasis{Fix SRD: bad bin assignment for SRD advection}}
\sphinxAtStartPar
Something has gone wrong in your SRD model; try using more
conservative settings.

\sphinxlineitem{\sphinxstyleemphasis{Fix SRD: bad search bin assignment}}
\sphinxAtStartPar
Something has gone wrong in your SRD model; try using more
conservative settings.

\sphinxlineitem{\sphinxstyleemphasis{Fix SRD: bad stencil bin for big particle}}
\sphinxAtStartPar
Something has gone wrong in your SRD model; try using more
conservative settings.

\sphinxlineitem{\sphinxstyleemphasis{Fix SRD: too many big particles in bin}}
\sphinxAtStartPar
Reset the ATOMPERBIN parameter at the top of fix\_srd.cpp
to a larger value, and re\sphinxhyphen{}compile the code.

\sphinxlineitem{\sphinxstyleemphasis{Fix SRD: too many walls in bin}}
\sphinxAtStartPar
This should not happen unless your system has been setup incorrectly.

\sphinxlineitem{\sphinxstyleemphasis{Fix adapt interface to this pair style not supported}}
\sphinxAtStartPar
New coding for the pair style would need to be done.

\sphinxlineitem{\sphinxstyleemphasis{Fix adapt kspace style does not exist}}
\sphinxAtStartPar
Self\sphinxhyphen{}explanatory.

\sphinxlineitem{\sphinxstyleemphasis{Fix adapt pair style does not exist}}
\sphinxAtStartPar
Self\sphinxhyphen{}explanatory

\sphinxlineitem{\sphinxstyleemphasis{Fix adapt pair style param not supported}}
\sphinxAtStartPar
The pair style does not know about the parameter you specified.

\sphinxlineitem{\sphinxstyleemphasis{Fix adapt requires atom attribute charge}}
\sphinxAtStartPar
The atom style being used does not specify an atom charge.

\sphinxlineitem{\sphinxstyleemphasis{Fix adapt requires atom attribute diameter}}
\sphinxAtStartPar
The atom style being used does not specify an atom diameter.

\sphinxlineitem{\sphinxstyleemphasis{Fix adapt type pair range is not valid for pair hybrid sub\sphinxhyphen{}style}}
\sphinxAtStartPar
Self\sphinxhyphen{}explanatory.

\sphinxlineitem{\sphinxstyleemphasis{Fix append/atoms requires a lattice be defined}}
\sphinxAtStartPar
Use the lattice command for this purpose.

\sphinxlineitem{\sphinxstyleemphasis{Fix ave/atom compute array is accessed out\sphinxhyphen{}of\sphinxhyphen{}range}}
\sphinxAtStartPar
Self\sphinxhyphen{}explanatory.

\sphinxlineitem{\sphinxstyleemphasis{Fix ave/atom compute does not calculate a per\sphinxhyphen{}atom array}}
\sphinxAtStartPar
Self\sphinxhyphen{}explanatory.

\sphinxlineitem{\sphinxstyleemphasis{Fix ave/atom compute does not calculate a per\sphinxhyphen{}atom vector}}
\sphinxAtStartPar
A compute used by fix ave/atom must generate per\sphinxhyphen{}atom values.

\sphinxlineitem{\sphinxstyleemphasis{Fix ave/atom compute does not calculate per\sphinxhyphen{}atom values}}
\sphinxAtStartPar
A compute used by fix ave/atom must generate per\sphinxhyphen{}atom values.

\sphinxlineitem{\sphinxstyleemphasis{Fix ave/atom fix array is accessed out\sphinxhyphen{}of\sphinxhyphen{}range}}
\sphinxAtStartPar
Self\sphinxhyphen{}explanatory.

\sphinxlineitem{\sphinxstyleemphasis{Fix ave/atom fix does not calculate a per\sphinxhyphen{}atom array}}
\sphinxAtStartPar
Self\sphinxhyphen{}explanatory.

\sphinxlineitem{\sphinxstyleemphasis{Fix ave/atom fix does not calculate a per\sphinxhyphen{}atom vector}}
\sphinxAtStartPar
A fix used by fix ave/atom must generate per\sphinxhyphen{}atom values.

\sphinxlineitem{\sphinxstyleemphasis{Fix ave/atom fix does not calculate per\sphinxhyphen{}atom values}}
\sphinxAtStartPar
A fix used by fix ave/atom must generate per\sphinxhyphen{}atom values.

\sphinxlineitem{\sphinxstyleemphasis{Fix ave/atom variable is not atom\sphinxhyphen{}style variable}}
\sphinxAtStartPar
A variable used by fix ave/atom must generate per\sphinxhyphen{}atom values.

\sphinxlineitem{\sphinxstyleemphasis{Fix ave/chunk compute does not calculate a per\sphinxhyphen{}atom array}}
\sphinxAtStartPar
Self\sphinxhyphen{}explanatory.

\sphinxlineitem{\sphinxstyleemphasis{Fix ave/chunk compute does not calculate a per\sphinxhyphen{}atom vector}}
\sphinxAtStartPar
Self\sphinxhyphen{}explanatory.

\sphinxlineitem{\sphinxstyleemphasis{Fix ave/chunk compute does not calculate per\sphinxhyphen{}atom values}}
\sphinxAtStartPar
Self\sphinxhyphen{}explanatory.

\sphinxlineitem{\sphinxstyleemphasis{Fix ave/chunk compute vector is accessed out\sphinxhyphen{}of\sphinxhyphen{}range}}
\sphinxAtStartPar
Self\sphinxhyphen{}explanatory.

\sphinxlineitem{\sphinxstyleemphasis{Fix ave/chunk does not use chunk/atom compute}}
\sphinxAtStartPar
The specified compute is not for a compute chunk/atom command.

\sphinxlineitem{\sphinxstyleemphasis{Fix ave/chunk fix does not calculate a per\sphinxhyphen{}atom array}}
\sphinxAtStartPar
Self\sphinxhyphen{}explanatory.

\sphinxlineitem{\sphinxstyleemphasis{Fix ave/chunk fix does not calculate a per\sphinxhyphen{}atom vector}}
\sphinxAtStartPar
Self\sphinxhyphen{}explanatory.

\sphinxlineitem{\sphinxstyleemphasis{Fix ave/chunk fix does not calculate per\sphinxhyphen{}atom values}}
\sphinxAtStartPar
Self\sphinxhyphen{}explanatory.

\sphinxlineitem{\sphinxstyleemphasis{Fix ave/chunk fix vector is accessed out\sphinxhyphen{}of\sphinxhyphen{}range}}
\sphinxAtStartPar
Self\sphinxhyphen{}explanatory.

\sphinxlineitem{\sphinxstyleemphasis{Fix ave/chunk variable is not atom\sphinxhyphen{}style variable}}
\sphinxAtStartPar
Self\sphinxhyphen{}explanatory.

\sphinxlineitem{\sphinxstyleemphasis{Fix ave/correlate compute does not calculate a scalar}}
\sphinxAtStartPar
Self\sphinxhyphen{}explanatory.

\sphinxlineitem{\sphinxstyleemphasis{Fix ave/correlate compute does not calculate a vector}}
\sphinxAtStartPar
Self\sphinxhyphen{}explanatory.

\sphinxlineitem{\sphinxstyleemphasis{Fix ave/correlate compute vector is accessed out\sphinxhyphen{}of\sphinxhyphen{}range}}
\sphinxAtStartPar
The index for the vector is out of bounds.

\sphinxlineitem{\sphinxstyleemphasis{Fix ave/correlate fix does not calculate a scalar}}
\sphinxAtStartPar
Self\sphinxhyphen{}explanatory.

\sphinxlineitem{\sphinxstyleemphasis{Fix ave/correlate fix does not calculate a vector}}
\sphinxAtStartPar
Self\sphinxhyphen{}explanatory.

\sphinxlineitem{\sphinxstyleemphasis{Fix ave/correlate fix vector is accessed out\sphinxhyphen{}of\sphinxhyphen{}range}}
\sphinxAtStartPar
The index for the vector is out of bounds.

\sphinxlineitem{\sphinxstyleemphasis{Fix ave/correlate variable is not equal\sphinxhyphen{}style variable}}
\sphinxAtStartPar
Self\sphinxhyphen{}explanatory.

\sphinxlineitem{\sphinxstyleemphasis{Fix ave/histo cannot input local values in scalar mode}}
\sphinxAtStartPar
Self\sphinxhyphen{}explanatory.

\sphinxlineitem{\sphinxstyleemphasis{Fix ave/histo cannot input per\sphinxhyphen{}atom values in scalar mode}}
\sphinxAtStartPar
Self\sphinxhyphen{}explanatory.

\sphinxlineitem{\sphinxstyleemphasis{Fix ave/histo compute array is accessed out\sphinxhyphen{}of\sphinxhyphen{}range}}
\sphinxAtStartPar
Self\sphinxhyphen{}explanatory.

\sphinxlineitem{\sphinxstyleemphasis{Fix ave/histo compute does not calculate a global array}}
\sphinxAtStartPar
Self\sphinxhyphen{}explanatory.

\sphinxlineitem{\sphinxstyleemphasis{Fix ave/histo compute does not calculate a global scalar}}
\sphinxAtStartPar
Self\sphinxhyphen{}explanatory.

\sphinxlineitem{\sphinxstyleemphasis{Fix ave/histo compute does not calculate a global vector}}
\sphinxAtStartPar
Self\sphinxhyphen{}explanatory.

\sphinxlineitem{\sphinxstyleemphasis{Fix ave/histo compute does not calculate a local array}}
\sphinxAtStartPar
Self\sphinxhyphen{}explanatory.

\sphinxlineitem{\sphinxstyleemphasis{Fix ave/histo compute does not calculate a local vector}}
\sphinxAtStartPar
Self\sphinxhyphen{}explanatory.

\sphinxlineitem{\sphinxstyleemphasis{Fix ave/histo compute does not calculate a per\sphinxhyphen{}atom array}}
\sphinxAtStartPar
Self\sphinxhyphen{}explanatory.

\sphinxlineitem{\sphinxstyleemphasis{Fix ave/histo compute does not calculate a per\sphinxhyphen{}atom vector}}
\sphinxAtStartPar
Self\sphinxhyphen{}explanatory.

\sphinxlineitem{\sphinxstyleemphasis{Fix ave/histo compute does not calculate local values}}
\sphinxAtStartPar
Self\sphinxhyphen{}explanatory.

\sphinxlineitem{\sphinxstyleemphasis{Fix ave/histo compute does not calculate per\sphinxhyphen{}atom values}}
\sphinxAtStartPar
Self\sphinxhyphen{}explanatory.

\sphinxlineitem{\sphinxstyleemphasis{Fix ave/histo compute vector is accessed out\sphinxhyphen{}of\sphinxhyphen{}range}}
\sphinxAtStartPar
Self\sphinxhyphen{}explanatory.

\sphinxlineitem{\sphinxstyleemphasis{Fix ave/histo fix array is accessed out\sphinxhyphen{}of\sphinxhyphen{}range}}
\sphinxAtStartPar
Self\sphinxhyphen{}explanatory.

\sphinxlineitem{\sphinxstyleemphasis{Fix ave/histo fix does not calculate a global array}}
\sphinxAtStartPar
Self\sphinxhyphen{}explanatory.

\sphinxlineitem{\sphinxstyleemphasis{Fix ave/histo fix does not calculate a global scalar}}
\sphinxAtStartPar
Self\sphinxhyphen{}explanatory.

\sphinxlineitem{\sphinxstyleemphasis{Fix ave/histo fix does not calculate a global vector}}
\sphinxAtStartPar
Self\sphinxhyphen{}explanatory.

\sphinxlineitem{\sphinxstyleemphasis{Fix ave/histo fix does not calculate a local array}}
\sphinxAtStartPar
Self\sphinxhyphen{}explanatory.

\sphinxlineitem{\sphinxstyleemphasis{Fix ave/histo fix does not calculate a local vector}}
\sphinxAtStartPar
Self\sphinxhyphen{}explanatory.

\sphinxlineitem{\sphinxstyleemphasis{Fix ave/histo fix does not calculate a per\sphinxhyphen{}atom array}}
\sphinxAtStartPar
Self\sphinxhyphen{}explanatory.

\sphinxlineitem{\sphinxstyleemphasis{Fix ave/histo fix does not calculate a per\sphinxhyphen{}atom vector}}
\sphinxAtStartPar
Self\sphinxhyphen{}explanatory.

\sphinxlineitem{\sphinxstyleemphasis{Fix ave/histo fix does not calculate local values}}
\sphinxAtStartPar
Self\sphinxhyphen{}explanatory.

\sphinxlineitem{\sphinxstyleemphasis{Fix ave/histo fix does not calculate per\sphinxhyphen{}atom values}}
\sphinxAtStartPar
Self\sphinxhyphen{}explanatory.

\sphinxlineitem{\sphinxstyleemphasis{Fix ave/histo fix vector is accessed out\sphinxhyphen{}of\sphinxhyphen{}range}}
\sphinxAtStartPar
Self\sphinxhyphen{}explanatory.

\sphinxlineitem{\sphinxstyleemphasis{Fix ave/histo input is invalid compute}}
\sphinxAtStartPar
Self\sphinxhyphen{}explanatory.

\sphinxlineitem{\sphinxstyleemphasis{Fix ave/histo input is invalid fix}}
\sphinxAtStartPar
Self\sphinxhyphen{}explanatory.

\sphinxlineitem{\sphinxstyleemphasis{Fix ave/histo input is invalid variable}}
\sphinxAtStartPar
Self\sphinxhyphen{}explanatory.

\sphinxlineitem{\sphinxstyleemphasis{Fix ave/histo inputs are not all global, peratom, or local}}
\sphinxAtStartPar
All inputs in a single fix ave/histo command must be of the
same style.

\sphinxlineitem{\sphinxstyleemphasis{Fix ave/histo/weight value and weight vector lengths do not match}}
\sphinxAtStartPar
Self\sphinxhyphen{}explanatory.

\sphinxlineitem{\sphinxstyleemphasis{Fix ave/time cannot set output array intensive/extensive from these inputs}}
\sphinxAtStartPar
One of more of the vector inputs has individual elements which are
flagged as intensive or extensive.  Such an input cannot be flagged as
all intensive/extensive when turned into an array by fix ave/time.

\sphinxlineitem{\sphinxstyleemphasis{Fix ave/time cannot use variable with vector mode}}
\sphinxAtStartPar
Variables produce scalar values.

\sphinxlineitem{\sphinxstyleemphasis{Fix ave/time columns are inconsistent lengths}}
\sphinxAtStartPar
Self\sphinxhyphen{}explanatory.

\sphinxlineitem{\sphinxstyleemphasis{Fix ave/time compute array is accessed out\sphinxhyphen{}of\sphinxhyphen{}range}}
\sphinxAtStartPar
An index for the array is out of bounds.

\sphinxlineitem{\sphinxstyleemphasis{Fix ave/time compute does not calculate a scalar}}
\sphinxAtStartPar
Self\sphinxhyphen{}explanatory.

\sphinxlineitem{\sphinxstyleemphasis{Fix ave/time compute does not calculate a vector}}
\sphinxAtStartPar
Self\sphinxhyphen{}explanatory.

\sphinxlineitem{\sphinxstyleemphasis{Fix ave/time compute does not calculate an array}}
\sphinxAtStartPar
Self\sphinxhyphen{}explanatory.

\sphinxlineitem{\sphinxstyleemphasis{Fix ave/time compute vector is accessed out\sphinxhyphen{}of\sphinxhyphen{}range}}
\sphinxAtStartPar
The index for the vector is out of bounds.

\sphinxlineitem{\sphinxstyleemphasis{Fix ave/time fix array cannot be variable length}}
\sphinxAtStartPar
Self\sphinxhyphen{}explanatory.

\sphinxlineitem{\sphinxstyleemphasis{Fix ave/time fix array is accessed out\sphinxhyphen{}of\sphinxhyphen{}range}}
\sphinxAtStartPar
An index for the array is out of bounds.

\sphinxlineitem{\sphinxstyleemphasis{Fix ave/time fix does not calculate a scalar}}
\sphinxAtStartPar
Self\sphinxhyphen{}explanatory.

\sphinxlineitem{\sphinxstyleemphasis{Fix ave/time fix does not calculate a vector}}
\sphinxAtStartPar
Self\sphinxhyphen{}explanatory.

\sphinxlineitem{\sphinxstyleemphasis{Fix ave/time fix does not calculate an array}}
\sphinxAtStartPar
Self\sphinxhyphen{}explanatory.

\sphinxlineitem{\sphinxstyleemphasis{Fix ave/time fix vector cannot be variable length}}
\sphinxAtStartPar
Self\sphinxhyphen{}explanatory.

\sphinxlineitem{\sphinxstyleemphasis{Fix ave/time fix vector is accessed out\sphinxhyphen{}of\sphinxhyphen{}range}}
\sphinxAtStartPar
The index for the vector is out of bounds.

\sphinxlineitem{\sphinxstyleemphasis{Fix ave/time variable is not equal\sphinxhyphen{}style variable}}
\sphinxAtStartPar
Self\sphinxhyphen{}explanatory.

\sphinxlineitem{\sphinxstyleemphasis{Fix balance rcb cannot be used with comm\_style brick}}
\sphinxAtStartPar
Comm\_style tiled must be used instead.

\sphinxlineitem{\sphinxstyleemphasis{Fix balance shift string is invalid}}
\sphinxAtStartPar
The string can only contain the characters “x”, “y”, or “z”.

\sphinxlineitem{\sphinxstyleemphasis{Fix bond/break needs ghost atoms from further away}}
\sphinxAtStartPar
This is because the fix needs to walk bonds to a certain distance to
acquire needed info, The comm\_modify cutoff command can be used to
extend the communication range.

\sphinxlineitem{\sphinxstyleemphasis{Fix bond/create angle type is invalid}}
\sphinxAtStartPar
Self\sphinxhyphen{}explanatory.

\sphinxlineitem{\sphinxstyleemphasis{Fix bond/create cutoff is longer than pairwise cutoff}}
\sphinxAtStartPar
This is not allowed because bond creation is done using the
pairwise neighbor list.

\sphinxlineitem{\sphinxstyleemphasis{Fix bond/create dihedral type is invalid}}
\sphinxAtStartPar
Self\sphinxhyphen{}explanatory.

\sphinxlineitem{\sphinxstyleemphasis{Fix bond/create improper type is invalid}}
\sphinxAtStartPar
Self\sphinxhyphen{}explanatory.

\sphinxlineitem{\sphinxstyleemphasis{Fix bond/create induced too many angles/dihedrals/impropers per atom}}
\sphinxAtStartPar
See the read\_data command for info on using the “extra/angle/per/atom”,
(or dihedral, improper) keywords to allow for additional
angles, dihedrals, and impropers to be formed.

\sphinxlineitem{\sphinxstyleemphasis{Fix bond/create needs ghost atoms from further away}}
\sphinxAtStartPar
This is because the fix needs to walk bonds to a certain distance to
acquire needed info, The comm\_modify cutoff command can be used to
extend the communication range.

\sphinxlineitem{\sphinxstyleemphasis{Fix bond/react: Cannot use fix bond/react with non\sphinxhyphen{}molecular systems}}
\sphinxAtStartPar
Only systems with bonds that can be changed can be used. Atom\_style
template does not qualify.

\sphinxlineitem{\sphinxstyleemphasis{Fix bond/react: Invalid template atom ID in map file}}
\sphinxAtStartPar
Atom IDs in molecule templates range from 1 to the number of atoms in the template.

\sphinxlineitem{\sphinxstyleemphasis{Fix bond/react: Rmax cutoff is longer than pairwise cutoff}}
\sphinxAtStartPar
This is not allowed because bond creation is done using the pairwise
neighbor list.

\sphinxlineitem{\sphinxstyleemphasis{Fix bond/react: Molecule template ID for fix bond/react does not exist}}
\sphinxAtStartPar
A valid molecule template must have been created with the molecule
command.

\sphinxlineitem{\sphinxstyleemphasis{Fix bond/react: Reaction templates must contain the same number of atoms}}
\sphinxAtStartPar
There should be a one\sphinxhyphen{}to\sphinxhyphen{}one correspondence between atoms in the
pre\sphinxhyphen{}reacted and post\sphinxhyphen{}reacted templates, as specified by the map file.

\sphinxlineitem{\sphinxstyleemphasis{Fix bond/react: Unknown section in map file}}
\sphinxAtStartPar
Please ensure reaction map files are properly formatted.

\sphinxlineitem{\sphinxstyleemphasis{Fix bond/react: Atom/Bond type affected by reaction too close to template edge}}
\sphinxAtStartPar
This means an atom which changes type or connectivity during the
reaction is too close to an ‘edge’ atom defined in the map
file.  This could cause incorrect assignment of bonds, angle, etc.
Generally, this means you must include more atoms in your templates,
such that there are at least two atoms between each atom involved in
the reaction and an edge atom.

\sphinxlineitem{\sphinxstyleemphasis{Fix bond/react: Fix bond/react needs ghost atoms from farther away}}
\sphinxAtStartPar
This is because a processor needs to map the entire unreacted
molecule template onto simulation atoms it knows about. The
comm\_modify cutoff command can be used to extend the communication
range.

\sphinxlineitem{\sphinxstyleemphasis{Fix bond/react: A deleted atom cannot remain bonded to an atom that is not deleted}}
\sphinxAtStartPar
Self\sphinxhyphen{}explanatory.

\sphinxlineitem{\sphinxstyleemphasis{Fix bond/react: First neighbors of chiral atoms must be of mutually different types}}
\sphinxAtStartPar
Self\sphinxhyphen{}explanatory.

\sphinxlineitem{\sphinxstyleemphasis{Fix bond/react: Chiral atoms must have exactly four first neighbors}}
\sphinxAtStartPar
Self\sphinxhyphen{}explanatory.

\sphinxlineitem{\sphinxstyleemphasis{Fix bond/react: Molecule template ‘Coords’ section required for chiralIDs keyword}}
\sphinxAtStartPar
The coordinates of atoms in the pre\sphinxhyphen{}reacted template are used to determine
chirality.

\sphinxlineitem{\sphinxstyleemphasis{Fix bond/react special bond generation overflow}}
\sphinxAtStartPar
The number of special bonds per\sphinxhyphen{}atom created by a reaction exceeds the
system setting. See the read\_data or create\_box command for how to
specify this value.

\sphinxlineitem{\sphinxstyleemphasis{Fix bond/react topology/atom exceed system topology/atom}}
\sphinxAtStartPar
The number of bonds, angles etc per\sphinxhyphen{}atom created by a reaction exceeds
the system setting. See the read\_data or create\_box command for how to
specify this value.

\sphinxlineitem{\sphinxstyleemphasis{Fix bond/swap cannot use dihedral or improper styles}}
\sphinxAtStartPar
These styles cannot be defined when using this fix.

\sphinxlineitem{\sphinxstyleemphasis{Fix bond/swap requires pair and bond styles}}
\sphinxAtStartPar
Self\sphinxhyphen{}explanatory.

\sphinxlineitem{\sphinxstyleemphasis{Fix bond/swap requires special\_bonds = 0,1,1}}
\sphinxAtStartPar
Self\sphinxhyphen{}explanatory.

\sphinxlineitem{\sphinxstyleemphasis{Fix box/relax generated negative box length}}
\sphinxAtStartPar
The pressure being applied is likely too large.  Try applying
it incrementally, to build to the high pressure.

\sphinxlineitem{\sphinxstyleemphasis{Fix command before simulation box is defined}}
\sphinxAtStartPar
The fix command cannot be used before a read\_data, read\_restart, or
create\_box command.

\sphinxlineitem{\sphinxstyleemphasis{Fix deform cannot use yz variable with xy}}
\sphinxAtStartPar
The yz setting cannot be a variable if xy deformation is also
specified.  This is because LAMMPS cannot determine if the yz setting
will induce a box flip which would be invalid if xy is also changing.

\sphinxlineitem{\sphinxstyleemphasis{Fix deform is changing yz too much with xy}}
\sphinxAtStartPar
When both yz and xy are changing, it induces changes in xz if the
box must flip from one tilt extreme to another.  Thus it is not
allowed for yz to grow so much that a flip is induced.

\sphinxlineitem{\sphinxstyleemphasis{Fix deform tilt factors require triclinic box}}
\sphinxAtStartPar
Cannot deform the tilt factors of a simulation box unless it
is a triclinic (non\sphinxhyphen{}orthogonal) box.

\sphinxlineitem{\sphinxstyleemphasis{Fix deform volume setting is invalid}}
\sphinxAtStartPar
Cannot use volume style unless other dimensions are being controlled.

\sphinxlineitem{\sphinxstyleemphasis{Fix deposit and fix rigid/small not using same molecule template ID}}
\sphinxAtStartPar
Self\sphinxhyphen{}explanatory.

\sphinxlineitem{\sphinxstyleemphasis{Fix deposit and fix shake not using same molecule template ID}}
\sphinxAtStartPar
Self\sphinxhyphen{}explanatory.

\sphinxlineitem{\sphinxstyleemphasis{Fix deposit molecule must have atom types}}
\sphinxAtStartPar
The defined molecule does not specify atom types.

\sphinxlineitem{\sphinxstyleemphasis{Fix deposit molecule must have coordinates}}
\sphinxAtStartPar
The defined molecule does not specify coordinates.

\sphinxlineitem{\sphinxstyleemphasis{Fix deposit molecule template ID must be same as atom\_style template ID}}
\sphinxAtStartPar
When using atom\_style template, you cannot deposit molecules that are
not in that template.

\sphinxlineitem{\sphinxstyleemphasis{Fix deposit region cannot be dynamic}}
\sphinxAtStartPar
Only static regions can be used with fix deposit.

\sphinxlineitem{\sphinxstyleemphasis{Fix deposit region does not support a bounding box}}
\sphinxAtStartPar
Not all regions represent bounded volumes.  You cannot use
such a region with the fix deposit command.

\sphinxlineitem{\sphinxstyleemphasis{Fix deposit shake fix does not exist}}
\sphinxAtStartPar
Self\sphinxhyphen{}explanatory.

\sphinxlineitem{\sphinxstyleemphasis{Fix efield requires atom attribute q or mu}}
\sphinxAtStartPar
The atom style defined does not have this attribute.

\sphinxlineitem{\sphinxstyleemphasis{Fix efield with dipoles cannot use atom\sphinxhyphen{}style variables}}
\sphinxAtStartPar
This option is not supported.

\sphinxlineitem{\sphinxstyleemphasis{Fix evaporate molecule requires atom attribute molecule}}
\sphinxAtStartPar
The atom style being used does not define a molecule ID.

\sphinxlineitem{\sphinxstyleemphasis{Fix external callback function not set}}
\sphinxAtStartPar
This must be done by an external program in order to use this fix.

\sphinxlineitem{\sphinxstyleemphasis{Fix for fix ave/atom not computed at compatible time}}
\sphinxAtStartPar
Fixes generate their values on specific timesteps.  Fix ave/atom is
requesting a value on a non\sphinxhyphen{}allowed timestep.

\sphinxlineitem{\sphinxstyleemphasis{Fix for fix ave/chunk not computed at compatible time}}
\sphinxAtStartPar
Fixes generate their values on specific timesteps.  Fix ave/chunk is
requesting a value on a non\sphinxhyphen{}allowed timestep.

\sphinxlineitem{\sphinxstyleemphasis{Fix for fix ave/correlate not computed at compatible time}}
\sphinxAtStartPar
Fixes generate their values on specific timesteps.  Fix ave/correlate
is requesting a value on a non\sphinxhyphen{}allowed timestep.

\sphinxlineitem{\sphinxstyleemphasis{Fix for fix ave/histo not computed at compatible time}}
\sphinxAtStartPar
Fixes generate their values on specific timesteps.  Fix ave/histo is
requesting a value on a non\sphinxhyphen{}allowed timestep.

\sphinxlineitem{\sphinxstyleemphasis{Fix for fix ave/spatial not computed at compatible time}}
\sphinxAtStartPar
Fixes generate their values on specific timesteps.  Fix ave/spatial is
requesting a value on a non\sphinxhyphen{}allowed timestep.

\sphinxlineitem{\sphinxstyleemphasis{Fix for fix ave/time not computed at compatible time}}
\sphinxAtStartPar
Fixes generate their values on specific timesteps.  Fix ave/time
is requesting a value on a non\sphinxhyphen{}allowed timestep.

\sphinxlineitem{\sphinxstyleemphasis{Fix for fix store/state not computed at compatible time}}
\sphinxAtStartPar
Fixes generate their values on specific timesteps.  Fix store/state
is requesting a value on a non\sphinxhyphen{}allowed timestep.

\sphinxlineitem{\sphinxstyleemphasis{Fix for fix vector not computed at compatible time}}
\sphinxAtStartPar
Fixes generate their values on specific timesteps.  Fix vector is
requesting a value on a non\sphinxhyphen{}allowed timestep.

\sphinxlineitem{\sphinxstyleemphasis{Fix freeze requires atom attribute torque}}
\sphinxAtStartPar
The atom style defined does not have this attribute.

\sphinxlineitem{\sphinxstyleemphasis{Fix gcmc and fix shake not using same molecule template ID}}
\sphinxAtStartPar
Self\sphinxhyphen{}explanatory.

\sphinxlineitem{\sphinxstyleemphasis{Fix gcmc atom has charge, but atom style does not}}
\sphinxAtStartPar
Self\sphinxhyphen{}explanatory.

\sphinxlineitem{\sphinxstyleemphasis{Fix gcmc cannot exchange individual atoms belonging to a molecule}}
\sphinxAtStartPar
This is an error since you should not delete only one atom of a
molecule.  The user has specified atomic (non\sphinxhyphen{}molecular) gas
exchanges, but an atom belonging to a molecule could be deleted.

\sphinxlineitem{\sphinxstyleemphasis{Fix gcmc does not (yet) work with atom\_style template}}
\sphinxAtStartPar
Self\sphinxhyphen{}explanatory.

\sphinxlineitem{\sphinxstyleemphasis{Fix gcmc molecule command requires that atoms have molecule attributes}}
\sphinxAtStartPar
Should not choose the gcmc molecule feature if no molecules are being
simulated. The general molecule flag is off, but gcmc’s molecule flag
is on.

\sphinxlineitem{\sphinxstyleemphasis{Fix gcmc molecule has charges, but atom style does not}}
\sphinxAtStartPar
Self\sphinxhyphen{}explanatory.

\sphinxlineitem{\sphinxstyleemphasis{Fix gcmc molecule must have atom types}}
\sphinxAtStartPar
The defined molecule does not specify atom types.

\sphinxlineitem{\sphinxstyleemphasis{Fix gcmc molecule must have coordinates}}
\sphinxAtStartPar
The defined molecule does not specify coordinates.

\sphinxlineitem{\sphinxstyleemphasis{Fix gcmc molecule template ID must be same as atom\_style template ID}}
\sphinxAtStartPar
When using atom\_style template, you cannot insert molecules that are
not in that template.

\sphinxlineitem{\sphinxstyleemphasis{Fix gcmc put atom outside box}}
\sphinxAtStartPar
This should not normally happen.  Contact the developers.

\sphinxlineitem{\sphinxstyleemphasis{Fix gcmc ran out of available atom IDs}}
\sphinxAtStartPar
See the setting for tagint in the src/lmptype.h file.

\sphinxlineitem{\sphinxstyleemphasis{Fix gcmc ran out of available molecule IDs}}
\sphinxAtStartPar
See the setting for tagint in the src/lmptype.h file.

\sphinxlineitem{\sphinxstyleemphasis{Fix gcmc region cannot be dynamic}}
\sphinxAtStartPar
Only static regions can be used with fix gcmc.

\sphinxlineitem{\sphinxstyleemphasis{Fix gcmc region does not support a bounding box}}
\sphinxAtStartPar
Not all regions represent bounded volumes.  You cannot use
such a region with the fix gcmc command.

\sphinxlineitem{\sphinxstyleemphasis{Fix gcmc region extends outside simulation box}}
\sphinxAtStartPar
Self\sphinxhyphen{}explanatory.

\sphinxlineitem{\sphinxstyleemphasis{Fix gcmc shake fix does not exist}}
\sphinxAtStartPar
Self\sphinxhyphen{}explanatory.

\sphinxlineitem{\sphinxstyleemphasis{Fix gld c coefficients must be \textgreater{}= 0}}
\sphinxAtStartPar
Self\sphinxhyphen{}explanatory.

\sphinxlineitem{\sphinxstyleemphasis{Fix gld needs more prony series coefficients}}
\sphinxAtStartPar
Self\sphinxhyphen{}explanatory.

\sphinxlineitem{\sphinxstyleemphasis{Fix gld prony terms must be \textgreater{} 0}}
\sphinxAtStartPar
Self\sphinxhyphen{}explanatory.

\sphinxlineitem{\sphinxstyleemphasis{Fix gld series type must be pprony for now}}
\sphinxAtStartPar
Self\sphinxhyphen{}explanatory.

\sphinxlineitem{\sphinxstyleemphasis{Fix gld start temperature must be \textgreater{}= 0}}
\sphinxAtStartPar
Self\sphinxhyphen{}explanatory.

\sphinxlineitem{\sphinxstyleemphasis{Fix gld stop temperature must be \textgreater{}= 0}}
\sphinxAtStartPar
Self\sphinxhyphen{}explanatory.

\sphinxlineitem{\sphinxstyleemphasis{Fix gld tau coefficients must be \textgreater{} 0}}
\sphinxAtStartPar
Self\sphinxhyphen{}explanatory.

\sphinxlineitem{\sphinxstyleemphasis{Fix halt variable is not equal\sphinxhyphen{}style variable}}
\sphinxAtStartPar
Self\sphinxhyphen{}explanatory.

\sphinxlineitem{\sphinxstyleemphasis{Fix heat group has no atoms}}
\sphinxAtStartPar
Self\sphinxhyphen{}explanatory.

\sphinxlineitem{\sphinxstyleemphasis{Fix heat kinetic energy of an atom went negative}}
\sphinxAtStartPar
This will cause the velocity rescaling about to be performed by fix
heat to be invalid.

\sphinxlineitem{\sphinxstyleemphasis{Fix heat kinetic energy went negative}}
\sphinxAtStartPar
This will cause the velocity rescaling about to be performed by fix
heat to be invalid.

\sphinxlineitem{\sphinxstyleemphasis{Fix in variable not computed at compatible time}}
\sphinxAtStartPar
Fixes generate their values on specific timesteps.  The variable is
requesting the values on a non\sphinxhyphen{}allowed timestep.

\sphinxlineitem{\sphinxstyleemphasis{Fix langevin angmom is not yet implemented with kokkos}}
\sphinxAtStartPar
This option is not yet available.

\sphinxlineitem{\sphinxstyleemphasis{Fix langevin angmom requires atom style ellipsoid}}
\sphinxAtStartPar
Self\sphinxhyphen{}explanatory.

\sphinxlineitem{\sphinxstyleemphasis{Fix langevin angmom requires extended particles}}
\sphinxAtStartPar
This fix option cannot be used with point particles.

\sphinxlineitem{\sphinxstyleemphasis{Fix langevin gjf and respa are not compatible}}
\sphinxAtStartPar
Self\sphinxhyphen{}explanatory.

\sphinxlineitem{\sphinxstyleemphasis{Fix langevin gjf cannot have period equal to dt/2}}
\sphinxAtStartPar
If the period is equal to dt/2 then division by zero will happen.

\sphinxlineitem{\sphinxstyleemphasis{Fix langevin gjf should come before fix nve}}
\sphinxAtStartPar
Self\sphinxhyphen{}explanatory.

\sphinxlineitem{\sphinxstyleemphasis{Fix langevin gjf with tbias is not yet implemented with kokkos}}
\sphinxAtStartPar
This option is not yet available.

\sphinxlineitem{\sphinxstyleemphasis{Fix langevin omega is not yet implemented with kokkos}}
\sphinxAtStartPar
This option is not yet available.

\sphinxlineitem{\sphinxstyleemphasis{Fix langevin omega requires atom style sphere}}
\sphinxAtStartPar
Self\sphinxhyphen{}explanatory.

\sphinxlineitem{\sphinxstyleemphasis{Fix langevin omega requires extended particles}}
\sphinxAtStartPar
One of the particles has radius 0.0.

\sphinxlineitem{\sphinxstyleemphasis{Fix langevin period must be \textgreater{} 0.0}}
\sphinxAtStartPar
The time window for temperature relaxation must be \textgreater{} 0

\sphinxlineitem{\sphinxstyleemphasis{Fix langevin variable returned negative temperature}}
\sphinxAtStartPar
Self\sphinxhyphen{}explanatory.

\sphinxlineitem{\sphinxstyleemphasis{Fix momentum group has no atoms}}
\sphinxAtStartPar
Self\sphinxhyphen{}explanatory.

\sphinxlineitem{\sphinxstyleemphasis{Fix move cannot define z or vz variable for 2d problem}}
\sphinxAtStartPar
Self\sphinxhyphen{}explanatory.

\sphinxlineitem{\sphinxstyleemphasis{Fix move cannot rotate aroung non z\sphinxhyphen{}axis for 2d problem}}
\sphinxAtStartPar
Self\sphinxhyphen{}explanatory.

\sphinxlineitem{\sphinxstyleemphasis{Fix move cannot set linear z motion for 2d problem}}
\sphinxAtStartPar
Self\sphinxhyphen{}explanatory.

\sphinxlineitem{\sphinxstyleemphasis{Fix move cannot set wiggle z motion for 2d problem}}
\sphinxAtStartPar
Self\sphinxhyphen{}explanatory.

\sphinxlineitem{\sphinxstyleemphasis{Fix msst compute ID does not compute potential energy}}
\sphinxAtStartPar
Self\sphinxhyphen{}explanatory.

\sphinxlineitem{\sphinxstyleemphasis{Fix msst compute ID does not compute pressure}}
\sphinxAtStartPar
Self\sphinxhyphen{}explanatory.

\sphinxlineitem{\sphinxstyleemphasis{Fix msst compute ID does not compute temperature}}
\sphinxAtStartPar
Self\sphinxhyphen{}explanatory.

\sphinxlineitem{\sphinxstyleemphasis{Fix msst requires a periodic box}}
\sphinxAtStartPar
Self\sphinxhyphen{}explanatory.

\sphinxlineitem{\sphinxstyleemphasis{Fix msst tscale must satisfy 0 \textless{}= tscale \textless{} 1}}
\sphinxAtStartPar
Self\sphinxhyphen{}explanatory.

\sphinxlineitem{\sphinxstyleemphasis{Fix npt/nph has tilted box too far in one step \sphinxhyphen{} periodic cell is too far from equilibrium state}}
\sphinxAtStartPar
Self\sphinxhyphen{}explanatory.  The change in the box tilt is too extreme
on a short timescale.

\sphinxlineitem{\sphinxstyleemphasis{Fix numdiff requires an atom map, see atom\_modify}}
\sphinxAtStartPar
Self\sphinxhyphen{}explanatory. Efficient loop over all atoms for numerical
difference requires an atom map.

\sphinxlineitem{\sphinxstyleemphasis{Fix numdiff requires consecutive atom IDs}}
\sphinxAtStartPar
Self\sphinxhyphen{}explanatory. Efficient loop over all atoms for numerical
difference requires consecutive atom IDs.

\sphinxlineitem{\sphinxstyleemphasis{Fix numdiff/virial must use group all}}
\sphinxAtStartPar
Virial contributions computed by this fix are
computed on all atoms.

\sphinxlineitem{\sphinxstyleemphasis{Fix nve/asphere requires extended particles}}
\sphinxAtStartPar
This fix can only be used for particles with a shape setting.

\sphinxlineitem{\sphinxstyleemphasis{Fix nve/asphere/noforce requires atom style ellipsoid}}
\sphinxAtStartPar
Self\sphinxhyphen{}explanatory.

\sphinxlineitem{\sphinxstyleemphasis{Fix nve/asphere/noforce requires extended particles}}
\sphinxAtStartPar
One of the particles is not an ellipsoid.

\sphinxlineitem{\sphinxstyleemphasis{Fix nve/body requires atom style body}}
\sphinxAtStartPar
Self\sphinxhyphen{}explanatory.

\sphinxlineitem{\sphinxstyleemphasis{Fix nve/body requires bodies}}
\sphinxAtStartPar
This fix can only be used for particles that are bodies.

\sphinxlineitem{\sphinxstyleemphasis{Fix nve/line can only be used for 2d simulations}}
\sphinxAtStartPar
Self\sphinxhyphen{}explanatory.

\sphinxlineitem{\sphinxstyleemphasis{Fix nve/line requires atom style line}}
\sphinxAtStartPar
Self\sphinxhyphen{}explanatory.

\sphinxlineitem{\sphinxstyleemphasis{Fix nve/line requires line particles}}
\sphinxAtStartPar
Self\sphinxhyphen{}explanatory.

\sphinxlineitem{\sphinxstyleemphasis{Fix nve/sphere dipole requires atom attribute mu}}
\sphinxAtStartPar
An atom style with this attribute is needed.

\sphinxlineitem{\sphinxstyleemphasis{Fix nve/sphere requires atom style sphere}}
\sphinxAtStartPar
Self\sphinxhyphen{}explanatory.

\sphinxlineitem{\sphinxstyleemphasis{Fix nve/sphere requires extended particles}}
\sphinxAtStartPar
This fix can only be used for particles of a finite size.

\sphinxlineitem{\sphinxstyleemphasis{Fix nve/tri can only be used for 3d simulations}}
\sphinxAtStartPar
Self\sphinxhyphen{}explanatory.

\sphinxlineitem{\sphinxstyleemphasis{Fix nve/tri requires atom style tri}}
\sphinxAtStartPar
Self\sphinxhyphen{}explanatory.

\sphinxlineitem{\sphinxstyleemphasis{Fix nve/tri requires tri particles}}
\sphinxAtStartPar
Self\sphinxhyphen{}explanatory.

\sphinxlineitem{\sphinxstyleemphasis{Fix nvt/nph/npt asphere requires extended particles}}
\sphinxAtStartPar
The shape setting for a particle in the fix group has shape = 0.0,
which means it is a point particle.

\sphinxlineitem{\sphinxstyleemphasis{Fix nvt/nph/npt body requires bodies}}
\sphinxAtStartPar
Self\sphinxhyphen{}explanatory.

\sphinxlineitem{\sphinxstyleemphasis{Fix nvt/nph/npt sphere requires atom style sphere}}
\sphinxAtStartPar
Self\sphinxhyphen{}explanatory.

\sphinxlineitem{\sphinxstyleemphasis{Fix nvt/npt/nph damping parameters must be \textgreater{} 0.0}}
\sphinxAtStartPar
Self\sphinxhyphen{}explanatory.

\sphinxlineitem{\sphinxstyleemphasis{Fix nvt/npt/nph dilate group ID does not exist}}
\sphinxAtStartPar
Self\sphinxhyphen{}explanatory.

\sphinxlineitem{\sphinxstyleemphasis{Fix nvt/sphere requires extended particles}}
\sphinxAtStartPar
This fix can only be used for particles of a finite size.

\sphinxlineitem{\sphinxstyleemphasis{Fix orient/fcc file open failed}}
\sphinxAtStartPar
The fix orient/fcc command could not open a specified file.

\sphinxlineitem{\sphinxstyleemphasis{Fix orient/fcc file read failed}}
\sphinxAtStartPar
The fix orient/fcc command could not read the needed parameters from a
specified file.

\sphinxlineitem{\sphinxstyleemphasis{Fix orient/fcc found self twice}}
\sphinxAtStartPar
The neighbor lists used by fix orient/fcc are messed up.  If this
error occurs, it is likely a bug, so send an email to the
\sphinxhref{https://www.lammps.org/authors.html}{developers}.

\sphinxlineitem{\sphinxstyleemphasis{Fix peri neigh does not exist}}
\sphinxAtStartPar
Somehow a fix that the pair style defines has been deleted.

\sphinxlineitem{\sphinxstyleemphasis{Fix pour and fix rigid/small not using same molecule template ID}}
\sphinxAtStartPar
Self\sphinxhyphen{}explanatory.

\sphinxlineitem{\sphinxstyleemphasis{Fix pour and fix shake not using same molecule template ID}}
\sphinxAtStartPar
Self\sphinxhyphen{}explanatory.

\sphinxlineitem{\sphinxstyleemphasis{Fix pour insertion count per timestep is 0}}
\sphinxAtStartPar
Self\sphinxhyphen{}explanatory.

\sphinxlineitem{\sphinxstyleemphasis{Fix pour molecule must have atom types}}
\sphinxAtStartPar
The defined molecule does not specify atom types.

\sphinxlineitem{\sphinxstyleemphasis{Fix pour molecule must have coordinates}}
\sphinxAtStartPar
The defined molecule does not specify coordinates.

\sphinxlineitem{\sphinxstyleemphasis{Fix pour molecule template ID must be same as atom style template ID}}
\sphinxAtStartPar
When using atom\_style template, you cannot pour molecules that are
not in that template.

\sphinxlineitem{\sphinxstyleemphasis{Fix pour polydisperse fractions do not sum to 1.0}}
\sphinxAtStartPar
Self\sphinxhyphen{}explanatory.

\sphinxlineitem{\sphinxstyleemphasis{Fix pour region ID does not exist}}
\sphinxAtStartPar
Self\sphinxhyphen{}explanatory.

\sphinxlineitem{\sphinxstyleemphasis{Fix pour region cannot be dynamic}}
\sphinxAtStartPar
Only static regions can be used with fix pour.

\sphinxlineitem{\sphinxstyleemphasis{Fix pour region does not support a bounding box}}
\sphinxAtStartPar
Not all regions represent bounded volumes.  You cannot use
such a region with the fix pour command.

\sphinxlineitem{\sphinxstyleemphasis{Fix pour requires atom attributes radius, rmass}}
\sphinxAtStartPar
The atom style defined does not have these attributes.

\sphinxlineitem{\sphinxstyleemphasis{Fix pour rigid fix does not exist}}
\sphinxAtStartPar
Self\sphinxhyphen{}explanatory.

\sphinxlineitem{\sphinxstyleemphasis{Fix pour shake fix does not exist}}
\sphinxAtStartPar
Self\sphinxhyphen{}explanatory.

\sphinxlineitem{\sphinxstyleemphasis{Fix press/berendsen damping parameters must be \textgreater{} 0.0}}
\sphinxAtStartPar
Self\sphinxhyphen{}explanatory.

\sphinxlineitem{\sphinxstyleemphasis{Fix property/atom cannot specify mol twice}}
\sphinxAtStartPar
Self\sphinxhyphen{}explanatory.

\sphinxlineitem{\sphinxstyleemphasis{Fix property/atom cannot specify q twice}}
\sphinxAtStartPar
Self\sphinxhyphen{}explanatory.

\sphinxlineitem{\sphinxstyleemphasis{Fix property/atom mol when atom\_style already has molecule attribute}}
\sphinxAtStartPar
Self\sphinxhyphen{}explanatory.

\sphinxlineitem{\sphinxstyleemphasis{Fix property/atom q when atom\_style already has charge attribute}}
\sphinxAtStartPar
Self\sphinxhyphen{}explanatory.

\sphinxlineitem{\sphinxstyleemphasis{Fix property/atom vector name already exists}}
\sphinxAtStartPar
The name for an integer or floating\sphinxhyphen{}point vector must be unique.

\sphinxlineitem{\sphinxstyleemphasis{Fix qeq has negative upper Taper radius cutoff}}
\sphinxAtStartPar
Self\sphinxhyphen{}explanatory.

\sphinxlineitem{\sphinxstyleemphasis{Fix qeq/comb group has no atoms}}
\sphinxAtStartPar
Self\sphinxhyphen{}explanatory.

\sphinxlineitem{\sphinxstyleemphasis{Fix qeq/comb requires atom attribute q}}
\sphinxAtStartPar
An atom style with charge must be used to perform charge equilibration.

\sphinxlineitem{\sphinxstyleemphasis{Fix qeq/dynamic group has no atoms}}
\sphinxAtStartPar
Self\sphinxhyphen{}explanatory.

\sphinxlineitem{\sphinxstyleemphasis{Fix qeq/dynamic requires atom attribute q}}
\sphinxAtStartPar
Self\sphinxhyphen{}explanatory.

\sphinxlineitem{\sphinxstyleemphasis{Fix qeq/fire group has no atoms}}
\sphinxAtStartPar
Self\sphinxhyphen{}explanatory.

\sphinxlineitem{\sphinxstyleemphasis{Fix qeq/fire requires atom attribute q}}
\sphinxAtStartPar
Self\sphinxhyphen{}explanatory.

\sphinxlineitem{\sphinxstyleemphasis{Fix qeq/point group has no atoms}}
\sphinxAtStartPar
Self\sphinxhyphen{}explanatory.

\sphinxlineitem{\sphinxstyleemphasis{Fix qeq/point has insufficient QEq matrix size}}
\sphinxAtStartPar
Occurs when number of neighbor atoms for an atom increased too much
during a run.  Increase SAFE\_ZONE and MIN\_CAP in fix\_qeq.h and
re\sphinxhyphen{}compile.

\sphinxlineitem{\sphinxstyleemphasis{Fix qeq/point requires atom attribute q}}
\sphinxAtStartPar
Self\sphinxhyphen{}explanatory.

\sphinxlineitem{\sphinxstyleemphasis{Fix qeq/shielded group has no atoms}}
\sphinxAtStartPar
Self\sphinxhyphen{}explanatory.

\sphinxlineitem{\sphinxstyleemphasis{Fix qeq/shielded has insufficient QEq matrix size}}
\sphinxAtStartPar
Occurs when number of neighbor atoms for an atom increased too much
during a run.  Increase SAFE\_ZONE and MIN\_CAP in fix\_qeq.h and
re\sphinxhyphen{}compile.

\sphinxlineitem{\sphinxstyleemphasis{Fix qeq/shielded requires atom attribute q}}
\sphinxAtStartPar
Self\sphinxhyphen{}explanatory.

\sphinxlineitem{\sphinxstyleemphasis{Fix qeq/slater could not extract params from pair coul/streitz}}
\sphinxAtStartPar
This should not happen unless pair coul/streitz has been altered.

\sphinxlineitem{\sphinxstyleemphasis{Fix qeq/slater group has no atoms}}
\sphinxAtStartPar
Self\sphinxhyphen{}explanatory.

\sphinxlineitem{\sphinxstyleemphasis{Fix qeq/slater has insufficient QEq matrix size}}
\sphinxAtStartPar
Occurs when number of neighbor atoms for an atom increased too much
during a run.  Increase SAFE\_ZONE and MIN\_CAP in fix\_qeq.h and
re\sphinxhyphen{}compile.

\sphinxlineitem{\sphinxstyleemphasis{Fix qeq/slater requires atom attribute q}}
\sphinxAtStartPar
Self\sphinxhyphen{}explanatory.

\sphinxlineitem{\sphinxstyleemphasis{Fix reax/bonds numbonds \textgreater{} nsbmax\_most}}
\sphinxAtStartPar
The limit of the number of bonds expected by the ReaxFF force field
was exceeded.

\sphinxlineitem{\sphinxstyleemphasis{Fix recenter group has no atoms}}
\sphinxAtStartPar
Self\sphinxhyphen{}explanatory.

\sphinxlineitem{\sphinxstyleemphasis{Fix restrain requires an atom map, see atom\_modify}}
\sphinxAtStartPar
Self\sphinxhyphen{}explanatory.

\sphinxlineitem{\sphinxstyleemphasis{Fix rigid atom has non\sphinxhyphen{}zero image flag in a non\sphinxhyphen{}periodic dimension}}
\sphinxAtStartPar
Image flags for non\sphinxhyphen{}periodic dimensions should not be set.

\sphinxlineitem{\sphinxstyleemphasis{Fix rigid file has no lines}}
\sphinxAtStartPar
Self\sphinxhyphen{}explanatory.

\sphinxlineitem{\sphinxstyleemphasis{Fix rigid langevin period must be \textgreater{} 0.0}}
\sphinxAtStartPar
Self\sphinxhyphen{}explanatory.

\sphinxlineitem{\sphinxstyleemphasis{Fix rigid molecule requires atom attribute molecule}}
\sphinxAtStartPar
Self\sphinxhyphen{}explanatory.

\sphinxlineitem{\sphinxstyleemphasis{Fix rigid npt/nph dilate group ID does not exist}}
\sphinxAtStartPar
Self\sphinxhyphen{}explanatory.

\sphinxlineitem{\sphinxstyleemphasis{Fix rigid npt/nph does not yet allow triclinic box}}
\sphinxAtStartPar
This is a current restriction in LAMMPS.

\sphinxlineitem{\sphinxstyleemphasis{Fix rigid npt/nph period must be \textgreater{} 0.0}}
\sphinxAtStartPar
Self\sphinxhyphen{}explanatory.

\sphinxlineitem{\sphinxstyleemphasis{Fix rigid npt/small t\_chain should not be less than 1}}
\sphinxAtStartPar
Self\sphinxhyphen{}explanatory.

\sphinxlineitem{\sphinxstyleemphasis{Fix rigid npt/small t\_order must be 3 or 5}}
\sphinxAtStartPar
Self\sphinxhyphen{}explanatory.

\sphinxlineitem{\sphinxstyleemphasis{Fix rigid nvt/npt/nph damping parameters must be \textgreater{} 0.0}}
\sphinxAtStartPar
Self\sphinxhyphen{}explanatory.

\sphinxlineitem{\sphinxstyleemphasis{Fix rigid nvt/small t\_chain should not be less than 1}}
\sphinxAtStartPar
Self\sphinxhyphen{}explanatory.

\sphinxlineitem{\sphinxstyleemphasis{Fix rigid nvt/small t\_iter should not be less than 1}}
\sphinxAtStartPar
Self\sphinxhyphen{}explanatory.

\sphinxlineitem{\sphinxstyleemphasis{Fix rigid nvt/small t\_order must be 3 or 5}}
\sphinxAtStartPar
Self\sphinxhyphen{}explanatory.

\sphinxlineitem{\sphinxstyleemphasis{Fix rigid xy torque cannot be on for 2d simulation}}
\sphinxAtStartPar
Self\sphinxhyphen{}explanatory.

\sphinxlineitem{\sphinxstyleemphasis{Fix rigid z force cannot be on for 2d simulation}}
\sphinxAtStartPar
Self\sphinxhyphen{}explanatory.

\sphinxlineitem{\sphinxstyleemphasis{Fix rigid/npt period must be \textgreater{} 0.0}}
\sphinxAtStartPar
Self\sphinxhyphen{}explanatory.

\sphinxlineitem{\sphinxstyleemphasis{Fix rigid/npt temperature order must be 3 or 5}}
\sphinxAtStartPar
Self\sphinxhyphen{}explanatory.

\sphinxlineitem{\sphinxstyleemphasis{Fix rigid/npt/small period must be \textgreater{} 0.0}}
\sphinxAtStartPar
Self\sphinxhyphen{}explanatory.

\sphinxlineitem{\sphinxstyleemphasis{Fix rigid/nvt period must be \textgreater{} 0.0}}
\sphinxAtStartPar
Self\sphinxhyphen{}explanatory.

\sphinxlineitem{\sphinxstyleemphasis{Fix rigid/nvt temperature order must be 3 or 5}}
\sphinxAtStartPar
Self\sphinxhyphen{}explanatory.

\sphinxlineitem{\sphinxstyleemphasis{Fix rigid/nvt/small period must be \textgreater{} 0.0}}
\sphinxAtStartPar
Self\sphinxhyphen{}explanatory.

\sphinxlineitem{\sphinxstyleemphasis{Fix rigid/small atom has non\sphinxhyphen{}zero image flag in a non\sphinxhyphen{}periodic dimension}}
\sphinxAtStartPar
Image flags for non\sphinxhyphen{}periodic dimensions should not be set.

\sphinxlineitem{\sphinxstyleemphasis{Fix rigid/small langevin period must be \textgreater{} 0.0}}
\sphinxAtStartPar
Self\sphinxhyphen{}explanatory.

\sphinxlineitem{\sphinxstyleemphasis{Fix rigid/small molecule must have atom types}}
\sphinxAtStartPar
The defined molecule does not specify atom types.

\sphinxlineitem{\sphinxstyleemphasis{Fix rigid/small molecule must have coordinates}}
\sphinxAtStartPar
The defined molecule does not specify coordinates.

\sphinxlineitem{\sphinxstyleemphasis{Fix rigid/small npt/nph period must be \textgreater{} 0.0}}
\sphinxAtStartPar
Self\sphinxhyphen{}explanatory.

\sphinxlineitem{\sphinxstyleemphasis{Fix rigid/small nvt/npt/nph damping parameters must be \textgreater{} 0.0}}
\sphinxAtStartPar
Self\sphinxhyphen{}explanatory.

\sphinxlineitem{\sphinxstyleemphasis{Fix rigid/small nvt/npt/nph dilate group ID does not exist}}
\sphinxAtStartPar
Self\sphinxhyphen{}explanatory.

\sphinxlineitem{\sphinxstyleemphasis{Fix rigid/small requires an atom map, see atom\_modify}}
\sphinxAtStartPar
Self\sphinxhyphen{}explanatory.

\sphinxlineitem{\sphinxstyleemphasis{Fix rigid/small requires atom attribute molecule}}
\sphinxAtStartPar
Self\sphinxhyphen{}explanatory.

\sphinxlineitem{\sphinxstyleemphasis{Fix rigid: Bad principal moments}}
\sphinxAtStartPar
The principal moments of inertia computed for a rigid body
are not within the required tolerances.

\sphinxlineitem{\sphinxstyleemphasis{Fix shake cannot be used with minimization}}
\sphinxAtStartPar
Cannot use fix shake while doing an energy minimization since
it turns off bonds that should contribute to the energy.

\sphinxlineitem{\sphinxstyleemphasis{Fix shake molecule template must have shake info}}
\sphinxAtStartPar
The defined molecule does not specify SHAKE information.

\sphinxlineitem{\sphinxstyleemphasis{Fix spring couple group ID does not exist}}
\sphinxAtStartPar
Self\sphinxhyphen{}explanatory.

\sphinxlineitem{\sphinxstyleemphasis{Fix srd can only currently be used with comm\_style brick}}
\sphinxAtStartPar
This is a current restriction in LAMMPS.

\sphinxlineitem{\sphinxstyleemphasis{Fix srd lamda must be \textgreater{}= 0.6 of SRD grid size}}
\sphinxAtStartPar
This is a requirement for accuracy reasons.

\sphinxlineitem{\sphinxstyleemphasis{Fix srd no\sphinxhyphen{}slip requires atom attribute torque}}
\sphinxAtStartPar
This is because the SRD collisions will impart torque to the solute
particles.

\sphinxlineitem{\sphinxstyleemphasis{Fix srd requires SRD particles all have same mass}}
\sphinxAtStartPar
Self\sphinxhyphen{}explanatory.

\sphinxlineitem{\sphinxstyleemphasis{Fix srd requires ghost atoms store velocity}}
\sphinxAtStartPar
Use the comm\_modify vel yes command to enable this.

\sphinxlineitem{\sphinxstyleemphasis{Fix srd requires newton pair on}}
\sphinxAtStartPar
Self\sphinxhyphen{}explanatory.

\sphinxlineitem{\sphinxstyleemphasis{Fix store/state compute array is accessed out\sphinxhyphen{}of\sphinxhyphen{}range}}
\sphinxAtStartPar
Self\sphinxhyphen{}explanatory.

\sphinxlineitem{\sphinxstyleemphasis{Fix store/state compute does not calculate a per\sphinxhyphen{}atom array}}
\sphinxAtStartPar
The compute calculates a per\sphinxhyphen{}atom vector.

\sphinxlineitem{\sphinxstyleemphasis{Fix store/state compute does not calculate a per\sphinxhyphen{}atom vector}}
\sphinxAtStartPar
The compute calculates a per\sphinxhyphen{}atom vector.

\sphinxlineitem{\sphinxstyleemphasis{Fix store/state compute does not calculate per\sphinxhyphen{}atom values}}
\sphinxAtStartPar
Computes that calculate global or local quantities cannot be used
with fix store/state.

\sphinxlineitem{\sphinxstyleemphasis{Fix store/state fix array is accessed out\sphinxhyphen{}of\sphinxhyphen{}range}}
\sphinxAtStartPar
Self\sphinxhyphen{}explanatory.

\sphinxlineitem{\sphinxstyleemphasis{Fix store/state fix does not calculate a per\sphinxhyphen{}atom array}}
\sphinxAtStartPar
The fix calculates a per\sphinxhyphen{}atom vector.

\sphinxlineitem{\sphinxstyleemphasis{Fix store/state fix does not calculate a per\sphinxhyphen{}atom vector}}
\sphinxAtStartPar
The fix calculates a per\sphinxhyphen{}atom array.

\sphinxlineitem{\sphinxstyleemphasis{Fix store/state fix does not calculate per\sphinxhyphen{}atom values}}
\sphinxAtStartPar
Fixes that calculate global or local quantities cannot be used with
fix store/state.

\sphinxlineitem{\sphinxstyleemphasis{Fix store/state for atom property that is not allocated}}
\sphinxAtStartPar
Self\sphinxhyphen{}explanatory.

\sphinxlineitem{\sphinxstyleemphasis{Fix store/state variable is not atom\sphinxhyphen{}style variable}}
\sphinxAtStartPar
Only atom\sphinxhyphen{}style variables calculate per\sphinxhyphen{}atom quantities.

\sphinxlineitem{\sphinxstyleemphasis{Fix temp/berendsen period must be \textgreater{} 0.0}}
\sphinxAtStartPar
Self\sphinxhyphen{}explanatory.

\sphinxlineitem{\sphinxstyleemphasis{Fix temp/berendsen variable returned negative temperature}}
\sphinxAtStartPar
Self\sphinxhyphen{}explanatory.

\sphinxlineitem{\sphinxstyleemphasis{Fix temp/csld is not compatible with fix rattle or fix shake}}
\sphinxAtStartPar
These two commands cannot currently be used together with fix temp/csld.

\sphinxlineitem{\sphinxstyleemphasis{Fix temp/csld variable returned negative temperature}}
\sphinxAtStartPar
Self\sphinxhyphen{}explanatory.

\sphinxlineitem{\sphinxstyleemphasis{Fix temp/csvr variable returned negative temperature}}
\sphinxAtStartPar
Self\sphinxhyphen{}explanatory.

\sphinxlineitem{\sphinxstyleemphasis{Fix temp/rescale variable returned negative temperature}}
\sphinxAtStartPar
Self\sphinxhyphen{}explanatory.

\sphinxlineitem{\sphinxstyleemphasis{Fix tfmc displacement length must be \textgreater{} 0}}
\sphinxAtStartPar
Self\sphinxhyphen{}explanatory.

\sphinxlineitem{\sphinxstyleemphasis{Fix tfmc is not compatible with fix shake}}
\sphinxAtStartPar
These two commands cannot currently be used together.

\sphinxlineitem{\sphinxstyleemphasis{Fix tfmc temperature must be \textgreater{} 0}}
\sphinxAtStartPar
Self\sphinxhyphen{}explanatory.

\sphinxlineitem{\sphinxstyleemphasis{Fix thermal/conductivity swap value must be positive}}
\sphinxAtStartPar
Self\sphinxhyphen{}explanatory.

\sphinxlineitem{\sphinxstyleemphasis{Fix tmd must come after integration fixes}}
\sphinxAtStartPar
Any fix tmd command must appear in the input script after all time
integration fixes (nve, nvt, npt).  See the fix tmd documentation for
details.

\sphinxlineitem{\sphinxstyleemphasis{Fix ttm electron temperatures must be \textgreater{} 0.0}}
\sphinxAtStartPar
Self\sphinxhyphen{}explanatory.

\sphinxlineitem{\sphinxstyleemphasis{Fix ttm electronic\_density must be \textgreater{} 0.0}}
\sphinxAtStartPar
Self\sphinxhyphen{}explanatory.

\sphinxlineitem{\sphinxstyleemphasis{Fix ttm electronic\_specific\_heat must be \textgreater{} 0.0}}
\sphinxAtStartPar
Self\sphinxhyphen{}explanatory.

\sphinxlineitem{\sphinxstyleemphasis{Fix ttm electronic\_thermal\_conductivity must be \textgreater{}= 0.0}}
\sphinxAtStartPar
Self\sphinxhyphen{}explanatory.

\sphinxlineitem{\sphinxstyleemphasis{Fix ttm gamma\_p must be \textgreater{} 0.0}}
\sphinxAtStartPar
Self\sphinxhyphen{}explanatory.

\sphinxlineitem{\sphinxstyleemphasis{Fix ttm gamma\_s must be \textgreater{}= 0.0}}
\sphinxAtStartPar
Self\sphinxhyphen{}explanatory.

\sphinxlineitem{\sphinxstyleemphasis{Fix ttm number of nodes must be \textgreater{} 0}}
\sphinxAtStartPar
Self\sphinxhyphen{}explanatory.

\sphinxlineitem{\sphinxstyleemphasis{Fix ttm v\_0 must be \textgreater{}= 0.0}}
\sphinxAtStartPar
Self\sphinxhyphen{}explanatory.

\sphinxlineitem{\sphinxstyleemphasis{Fix used in compute chunk/atom not computed at compatible time}}
\sphinxAtStartPar
The chunk/atom compute cannot query the output of the fix on a timestep
it is needed.

\sphinxlineitem{\sphinxstyleemphasis{Fix used in compute reduce not computed at compatible time}}
\sphinxAtStartPar
Fixes generate their values on specific timesteps.  Compute reduce is
requesting a value on a non\sphinxhyphen{}allowed timestep.

\sphinxlineitem{\sphinxstyleemphasis{Fix used in compute slice not computed at compatible time}}
\sphinxAtStartPar
Fixes generate their values on specific timesteps.  Compute slice is
requesting a value on a non\sphinxhyphen{}allowed timestep.

\sphinxlineitem{\sphinxstyleemphasis{Fix vector cannot set output array intensive/extensive from these inputs}}
\sphinxAtStartPar
The inputs to the command have conflicting intensive/extensive attributes.
You need to use more than one fix vector command.

\sphinxlineitem{\sphinxstyleemphasis{Fix vector compute does not calculate a scalar}}
\sphinxAtStartPar
Self\sphinxhyphen{}explanatory.

\sphinxlineitem{\sphinxstyleemphasis{Fix vector compute does not calculate a vector}}
\sphinxAtStartPar
Self\sphinxhyphen{}explanatory.

\sphinxlineitem{\sphinxstyleemphasis{Fix vector compute vector is accessed out\sphinxhyphen{}of\sphinxhyphen{}range}}
\sphinxAtStartPar
Self\sphinxhyphen{}explanatory.

\sphinxlineitem{\sphinxstyleemphasis{Fix vector fix does not calculate a scalar}}
\sphinxAtStartPar
Self\sphinxhyphen{}explanatory.

\sphinxlineitem{\sphinxstyleemphasis{Fix vector fix does not calculate a vector}}
\sphinxAtStartPar
Self\sphinxhyphen{}explanatory.

\sphinxlineitem{\sphinxstyleemphasis{Fix vector fix vector is accessed out\sphinxhyphen{}of\sphinxhyphen{}range}}
\sphinxAtStartPar
Self\sphinxhyphen{}explanatory.

\sphinxlineitem{\sphinxstyleemphasis{Fix vector variable is not equal\sphinxhyphen{}style variable}}
\sphinxAtStartPar
Self\sphinxhyphen{}explanatory.

\sphinxlineitem{\sphinxstyleemphasis{Fix viscosity swap value must be positive}}
\sphinxAtStartPar
Self\sphinxhyphen{}explanatory.

\sphinxlineitem{\sphinxstyleemphasis{Fix viscosity vtarget value must be positive}}
\sphinxAtStartPar
Self\sphinxhyphen{}explanatory.

\sphinxlineitem{\sphinxstyleemphasis{Fix wall cutoff \textless{}= 0.0}}
\sphinxAtStartPar
Self\sphinxhyphen{}explanatory.

\sphinxlineitem{\sphinxstyleemphasis{Fix wall/colloid requires atom style sphere}}
\sphinxAtStartPar
Self\sphinxhyphen{}explanatory.

\sphinxlineitem{\sphinxstyleemphasis{Fix wall/colloid requires extended particles}}
\sphinxAtStartPar
One of the particles has radius 0.0.

\sphinxlineitem{\sphinxstyleemphasis{Fix wall/gran is incompatible with Pair style}}
\sphinxAtStartPar
Must use a granular pair style to define the parameters needed for
this fix.

\sphinxlineitem{\sphinxstyleemphasis{Fix wall/gran requires atom style sphere}}
\sphinxAtStartPar
Self\sphinxhyphen{}explanatory.

\sphinxlineitem{\sphinxstyleemphasis{Fix wall/piston command only available at zlo}}
\sphinxAtStartPar
The face keyword must be zlo.

\sphinxlineitem{\sphinxstyleemphasis{Fix wall/region colloid requires atom style sphere}}
\sphinxAtStartPar
Self\sphinxhyphen{}explanatory.

\sphinxlineitem{\sphinxstyleemphasis{Fix wall/region colloid requires extended particles}}
\sphinxAtStartPar
One of the particles has radius 0.0.

\sphinxlineitem{\sphinxstyleemphasis{Fix wall/region cutoff \textless{}= 0.0}}
\sphinxAtStartPar
Self\sphinxhyphen{}explanatory.

\sphinxlineitem{\sphinxstyleemphasis{Fix\_modify pressure ID does not compute pressure}}
\sphinxAtStartPar
The compute ID assigned to the fix must compute pressure.

\sphinxlineitem{\sphinxstyleemphasis{Fix\_modify temperature ID does not compute temperature}}
\sphinxAtStartPar
The compute ID assigned to the fix must compute temperature.

\sphinxlineitem{\sphinxstyleemphasis{For triclinic deformation, specified target stress must be hydrostatic}}
\sphinxAtStartPar
Triclinic pressure control is allowed using the tri keyword, but
non\sphinxhyphen{}hydrostatic pressure control can not be used in this case.

\sphinxlineitem{\sphinxstyleemphasis{Found no restart file matching pattern}}
\sphinxAtStartPar
When using a “*” in the restart file name, no matching file was found.

\sphinxlineitem{\sphinxstyleemphasis{GPU library not compiled for this accelerator}}
\sphinxAtStartPar
Self\sphinxhyphen{}explanatory.

\sphinxlineitem{\sphinxstyleemphasis{GPU package does not (yet) work with atom\_style template}}
\sphinxAtStartPar
Self\sphinxhyphen{}explanatory.

\sphinxlineitem{\sphinxstyleemphasis{GPU particle split must be set to 1 for this pair style.}}
\sphinxAtStartPar
For this pair style, you cannot run part of the force calculation on
the host.  See the package command.

\sphinxlineitem{\sphinxstyleemphasis{GPUs are requested but Kokkos has not been compiled for CUDA}}
\sphinxAtStartPar
Re\sphinxhyphen{}compile Kokkos with CUDA support to use GPUs.

\sphinxlineitem{\sphinxstyleemphasis{Ghost velocity forward comm not yet implemented with Kokkos}}
\sphinxAtStartPar
This is a current restriction.

\sphinxlineitem{\sphinxstyleemphasis{Gmask function in equal\sphinxhyphen{}style variable formula}}
\sphinxAtStartPar
Gmask is per\sphinxhyphen{}atom operation.

\sphinxlineitem{\sphinxstyleemphasis{Gravity changed since fix pour was created}}
\sphinxAtStartPar
The gravity vector defined by fix gravity must be static.

\sphinxlineitem{\sphinxstyleemphasis{Gravity must point in \sphinxhyphen{}y to use with fix pour in 2d}}
\sphinxAtStartPar
Self\sphinxhyphen{}explanatory.

\sphinxlineitem{\sphinxstyleemphasis{Gravity must point in \sphinxhyphen{}z to use with fix pour in 3d}}
\sphinxAtStartPar
Self\sphinxhyphen{}explanatory.

\sphinxlineitem{\sphinxstyleemphasis{Grmask function in equal\sphinxhyphen{}style variable formula}}
\sphinxAtStartPar
Grmask is per\sphinxhyphen{}atom operation.

\sphinxlineitem{\sphinxstyleemphasis{Group ID does not exist}}
\sphinxAtStartPar
A group ID used in the group command does not exist.

\sphinxlineitem{\sphinxstyleemphasis{Group ID in variable formula does not exist}}
\sphinxAtStartPar
Self\sphinxhyphen{}explanatory.

\sphinxlineitem{\sphinxstyleemphasis{Group all cannot be made dynamic}}
\sphinxAtStartPar
This operation is not allowed.

\sphinxlineitem{\sphinxstyleemphasis{Group command before simulation box is defined}}
\sphinxAtStartPar
The group command cannot be used before a read\_data, read\_restart, or
create\_box command.

\sphinxlineitem{\sphinxstyleemphasis{Group dynamic cannot reference itself}}
\sphinxAtStartPar
Self\sphinxhyphen{}explanatory.

\sphinxlineitem{\sphinxstyleemphasis{Group dynamic parent group cannot be dynamic}}
\sphinxAtStartPar
Self\sphinxhyphen{}explanatory.

\sphinxlineitem{\sphinxstyleemphasis{Group dynamic parent group does not exist}}
\sphinxAtStartPar
Self\sphinxhyphen{}explanatory.

\sphinxlineitem{\sphinxstyleemphasis{Group region ID does not exist}}
\sphinxAtStartPar
A region ID used in the group command does not exist.

\sphinxlineitem{\sphinxstyleemphasis{If read\_dump purges it cannot replace or trim}}
\sphinxAtStartPar
These operations are not compatible.  See the read\_dump doc
page for details.

\sphinxlineitem{\sphinxstyleemphasis{Illegal … command}}
\sphinxAtStartPar
Self\sphinxhyphen{}explanatory.  Check the input script syntax and compare to the
documentation for the command.  You can use \sphinxhyphen{}echo screen as a
command\sphinxhyphen{}line option when running LAMMPS to see the offending line.

\sphinxlineitem{\sphinxstyleemphasis{Illegal COMB parameter}}
\sphinxAtStartPar
One or more of the coefficients defined in the potential file is
invalid.

\sphinxlineitem{\sphinxstyleemphasis{Illegal COMB3 parameter}}
\sphinxAtStartPar
One or more of the coefficients defined in the potential file is
invalid.

\sphinxlineitem{\sphinxstyleemphasis{Illegal Stillinger\sphinxhyphen{}Weber parameter}}
\sphinxAtStartPar
One or more of the coefficients defined in the potential file is
invalid.

\sphinxlineitem{\sphinxstyleemphasis{Illegal Tersoff parameter}}
\sphinxAtStartPar
One or more of the coefficients defined in the potential file is
invalid.

\sphinxlineitem{\sphinxstyleemphasis{Illegal Vashishta parameter}}
\sphinxAtStartPar
One or more of the coefficients defined in the potential file is
invalid.

\sphinxlineitem{\sphinxstyleemphasis{Illegal compute voronoi/atom command (occupation and (surface or edges))}}
\sphinxAtStartPar
Self\sphinxhyphen{}explanatory.

\sphinxlineitem{\sphinxstyleemphasis{Illegal coul/streitz parameter}}
\sphinxAtStartPar
One or more of the coefficients defined in the potential file is
invalid.

\sphinxlineitem{\sphinxstyleemphasis{Illegal dump\_modify sfactor value (must be \textgreater{} 0.0)}}
\sphinxAtStartPar
Self\sphinxhyphen{}explanatory.

\sphinxlineitem{\sphinxstyleemphasis{Illegal dump\_modify tfactor value (must be \textgreater{} 0.0)}}
\sphinxAtStartPar
Self\sphinxhyphen{}explanatory.

\sphinxlineitem{\sphinxstyleemphasis{Illegal fix gcmc gas mass \textless{}= 0}}
\sphinxAtStartPar
The computed mass of the designated gas molecule or atom type was less
than or equal to zero.

\sphinxlineitem{\sphinxstyleemphasis{Illegal fix tfmc random seed}}
\sphinxAtStartPar
Seeds can only be nonzero positive integers.

\sphinxlineitem{\sphinxstyleemphasis{Illegal fix wall/piston velocity}}
\sphinxAtStartPar
The piston velocity must be positive.

\sphinxlineitem{\sphinxstyleemphasis{Illegal integrate style}}
\sphinxAtStartPar
Self\sphinxhyphen{}explanatory.

\sphinxlineitem{\sphinxstyleemphasis{Illegal nb3b/harmonic parameter}}
\sphinxAtStartPar
One or more of the coefficients defined in the potential file is
invalid.

\sphinxlineitem{\sphinxstyleemphasis{Illegal number of angle table entries}}
\sphinxAtStartPar
There must be at least 2 table entries.

\sphinxlineitem{\sphinxstyleemphasis{Illegal number of bond table entries}}
\sphinxAtStartPar
There must be at least 2 table entries.

\sphinxlineitem{\sphinxstyleemphasis{Illegal number of pair table entries}}
\sphinxAtStartPar
There must be at least 2 table entries.

\sphinxlineitem{\sphinxstyleemphasis{Illegal or unset periodicity in restart}}
\sphinxAtStartPar
This error should not normally occur unless the restart file is invalid.

\sphinxlineitem{\sphinxstyleemphasis{Illegal range increment value}}
\sphinxAtStartPar
The increment must be \textgreater{}= 1.

\sphinxlineitem{\sphinxstyleemphasis{Illegal simulation box}}
\sphinxAtStartPar
The lower bound of the simulation box is greater than the upper bound.

\sphinxlineitem{\sphinxstyleemphasis{Illegal size double vector read requested}}
\sphinxAtStartPar
This error should not normally occur unless the restart file is invalid.

\sphinxlineitem{\sphinxstyleemphasis{Illegal size integer vector read requested}}
\sphinxAtStartPar
This error should not normally occur unless the restart file is invalid.

\sphinxlineitem{\sphinxstyleemphasis{Illegal size string or corrupt restart}}
\sphinxAtStartPar
This error should not normally occur unless the restart file is invalid.

\sphinxlineitem{\sphinxstyleemphasis{Imageint setting in lmptype.h is invalid}}
\sphinxAtStartPar
Imageint must be as large or larger than smallint.

\sphinxlineitem{\sphinxstyleemphasis{Imageint setting in lmptype.h is not compatible}}
\sphinxAtStartPar
Format of imageint stored in restart file is not consistent with
LAMMPS version you are running.  See the settings in src/lmptype.h

\sphinxlineitem{\sphinxstyleemphasis{Improper atom missing in delete\_bonds}}
\sphinxAtStartPar
The delete\_bonds command cannot find one or more atoms in a particular
improper on a particular processor.  The pairwise cutoff is too short
or the atoms are too far apart to make a valid improper.

\sphinxlineitem{\sphinxstyleemphasis{Improper atom missing in set command}}
\sphinxAtStartPar
The set command cannot find one or more atoms in a particular improper
on a particular processor.  The pairwise cutoff is too short or the
atoms are too far apart to make a valid improper.

\sphinxlineitem{\sphinxstyleemphasis{Improper atoms \%d \%d \%d \%d missing on proc \%d at step \%ld}}
\sphinxAtStartPar
One or more of 4 atoms needed to compute a particular improper are
missing on this processor.  Typically this is because the pairwise
cutoff is set too short or the improper has blown apart and an atom is
too far away.

\sphinxlineitem{\sphinxstyleemphasis{Improper atoms missing on proc \%d at step \%ld}}
\sphinxAtStartPar
One or more of 4 atoms needed to compute a particular improper are
missing on this processor.  Typically this is because the pairwise
cutoff is set too short or the improper has blown apart and an atom is
too far away.

\sphinxlineitem{\sphinxstyleemphasis{Improper coeff for hybrid has invalid style}}
\sphinxAtStartPar
Improper style hybrid uses another improper style as one of its
coefficients.  The improper style used in the improper\_coeff command
or read from a restart file is not recognized.

\sphinxlineitem{\sphinxstyleemphasis{Improper coeffs are not set}}
\sphinxAtStartPar
No improper coefficients have been assigned in the data file or via
the improper\_coeff command.

\sphinxlineitem{\sphinxstyleemphasis{Improper style hybrid cannot have hybrid as an argument}}
\sphinxAtStartPar
Self\sphinxhyphen{}explanatory.

\sphinxlineitem{\sphinxstyleemphasis{Improper style hybrid cannot have none as an argument}}
\sphinxAtStartPar
Self\sphinxhyphen{}explanatory.

\sphinxlineitem{\sphinxstyleemphasis{Improper style hybrid cannot use same improper style twice}}
\sphinxAtStartPar
Self\sphinxhyphen{}explanatory.

\sphinxlineitem{\sphinxstyleemphasis{Improper\_coeff command before improper\_style is defined}}
\sphinxAtStartPar
Coefficients cannot be set in the data file or via the improper\_coeff
command until an improper\_style has been assigned.

\sphinxlineitem{\sphinxstyleemphasis{Improper\_coeff command before simulation box is defined}}
\sphinxAtStartPar
The improper\_coeff command cannot be used before a read\_data,
read\_restart, or create\_box command.

\sphinxlineitem{\sphinxstyleemphasis{Improper\_coeff command when no impropers allowed}}
\sphinxAtStartPar
The chosen atom style does not allow for impropers to be defined.

\sphinxlineitem{\sphinxstyleemphasis{Improper\_style command when no impropers allowed}}
\sphinxAtStartPar
The chosen atom style does not allow for impropers to be defined.

\sphinxlineitem{\sphinxstyleemphasis{Impropers assigned incorrectly}}
\sphinxAtStartPar
Impropers read in from the data file were not assigned correctly to
atoms.  This means there is something invalid about the topology
definitions.

\sphinxlineitem{\sphinxstyleemphasis{Impropers defined but no improper types}}
\sphinxAtStartPar
The data file header lists improper but no improper types.

\sphinxlineitem{\sphinxstyleemphasis{Incompatible KIM Simulator Model}}
\sphinxAtStartPar
The requested KIM Simulator Model was defined for a different MD code
and thus is not compatible with LAMMPS.

\sphinxlineitem{\sphinxstyleemphasis{Incompatible units for KIM Simulator Model}}
\sphinxAtStartPar
The selected unit style is not compatible with the requested KIM
Simulator Model.

\sphinxlineitem{\sphinxstyleemphasis{Incomplete use of variables in create\_atoms command}}
\sphinxAtStartPar
The var and set options must be used together.

\sphinxlineitem{\sphinxstyleemphasis{Inconsistent iparam/jparam values in fix bond/create command}}
\sphinxAtStartPar
If itype and jtype are the same, then their maxbond and newtype
settings must also be the same.

\sphinxlineitem{\sphinxstyleemphasis{Inconsistent line segment in data file}}
\sphinxAtStartPar
The end points of the line segment are not equal distances from the
center point which is the atom coordinate.

\sphinxlineitem{\sphinxstyleemphasis{Inconsistent triangle in data file}}
\sphinxAtStartPar
The centroid of the triangle as defined by the corner points is not
the atom coordinate.

\sphinxlineitem{\sphinxstyleemphasis{Inconsistent use of finite\sphinxhyphen{}size particles by molecule template molecules}}
\sphinxAtStartPar
Not all of the molecules define a radius for their constituent
particles.

\sphinxlineitem{\sphinxstyleemphasis{Incorrect \# of floating\sphinxhyphen{}point values in Bodies section of data file}}
\sphinxAtStartPar
See page for body style.

\sphinxlineitem{\sphinxstyleemphasis{Incorrect \# of integer values in Bodies section of data file}}
\sphinxAtStartPar
See page for body style.

\sphinxlineitem{\sphinxstyleemphasis{Incorrect \%s format in data file}}
\sphinxAtStartPar
A section of the data file being read by fix property/atom does
not have the correct number of values per line.

\sphinxlineitem{\sphinxstyleemphasis{Incorrect SNAP parameter file}}
\sphinxAtStartPar
The file cannot be parsed correctly, check its internal syntax.

\sphinxlineitem{\sphinxstyleemphasis{Incorrect args for angle coefficients}}
\sphinxAtStartPar
Self\sphinxhyphen{}explanatory.  Check the input script or data file.

\sphinxlineitem{\sphinxstyleemphasis{Incorrect args for bond coefficients}}
\sphinxAtStartPar
Self\sphinxhyphen{}explanatory.  Check the input script or data file.

\sphinxlineitem{\sphinxstyleemphasis{Incorrect args for dihedral coefficients}}
\sphinxAtStartPar
Self\sphinxhyphen{}explanatory.  Check the input script or data file.

\sphinxlineitem{\sphinxstyleemphasis{Incorrect args for improper coefficients}}
\sphinxAtStartPar
Self\sphinxhyphen{}explanatory.  Check the input script or data file.

\sphinxlineitem{\sphinxstyleemphasis{Incorrect args for pair coefficients}}
\sphinxAtStartPar
Self\sphinxhyphen{}explanatory.  Check the input script or data file.

\sphinxlineitem{\sphinxstyleemphasis{Incorrect args in pair\_style command}}
\sphinxAtStartPar
Self\sphinxhyphen{}explanatory.

\sphinxlineitem{\sphinxstyleemphasis{Incorrect atom format in data file}}
\sphinxAtStartPar
Number of values per atom line in the data file is not consistent with
the atom style.

\sphinxlineitem{\sphinxstyleemphasis{Incorrect atom format in neb file}}
\sphinxAtStartPar
The number of fields per line is not what expected.

\sphinxlineitem{\sphinxstyleemphasis{Incorrect bonus data format in data file}}
\sphinxAtStartPar
See the read\_data page for a description of how various kinds of
bonus data must be formatted for certain atom styles.

\sphinxlineitem{\sphinxstyleemphasis{Incorrect boundaries with slab Ewald}}
\sphinxAtStartPar
Must have periodic x,y dimensions and non\sphinxhyphen{}periodic z dimension to use
2d slab option with Ewald.

\sphinxlineitem{\sphinxstyleemphasis{Incorrect boundaries with slab EwaldDisp}}
\sphinxAtStartPar
Must have periodic x,y dimensions and non\sphinxhyphen{}periodic z dimension to use
2d slab option with Ewald.

\sphinxlineitem{\sphinxstyleemphasis{Incorrect boundaries with slab PPPM}}
\sphinxAtStartPar
Must have periodic x,y dimensions and non\sphinxhyphen{}periodic z dimension to use
2d slab option with PPPM.

\sphinxlineitem{\sphinxstyleemphasis{Incorrect boundaries with slab PPPMDisp}}
\sphinxAtStartPar
Must have periodic x,y dimensions and non\sphinxhyphen{}periodic z dimension to use
2d slab option with pppm/disp.

\sphinxlineitem{\sphinxstyleemphasis{Incorrect conversion in format string}}
\sphinxAtStartPar
A format style variable was not using either a \%f, a \%g, or a \%e conversion.
Or an immediate variable with format suffix was not using either
a \%f, a \%g or a \%e conversion in the format suffix.

\sphinxlineitem{\sphinxstyleemphasis{Incorrect element names in ADP potential file}}
\sphinxAtStartPar
The element names in the ADP file do not match those requested.

\sphinxlineitem{\sphinxstyleemphasis{Incorrect element names in EAM potential file}}
\sphinxAtStartPar
The element names in the EAM file do not match those requested.

\sphinxlineitem{\sphinxstyleemphasis{Incorrect format of … section in data file}}
\sphinxAtStartPar
Number or type of values per line in the given section of the data file
is not consistent with the requirements for this section.

\sphinxlineitem{\sphinxstyleemphasis{Incorrect format in COMB potential file}}
\sphinxAtStartPar
Incorrect number of words per line in the potential file.

\sphinxlineitem{\sphinxstyleemphasis{Incorrect format in COMB3 potential file}}
\sphinxAtStartPar
Incorrect number of words per line in the potential file.

\sphinxlineitem{\sphinxstyleemphasis{Incorrect format in MEAM library file}}
\sphinxAtStartPar
Incorrect number of words per line in the potential file.

\sphinxlineitem{\sphinxstyleemphasis{Incorrect format in SNAP coefficient file}}
\sphinxAtStartPar
Incorrect number of words per line in the coefficient file.

\sphinxlineitem{\sphinxstyleemphasis{Incorrect format in SNAP parameter file}}
\sphinxAtStartPar
Incorrect number of words per line in the parameter file.

\sphinxlineitem{\sphinxstyleemphasis{Incorrect format in Stillinger\sphinxhyphen{}Weber potential file}}
\sphinxAtStartPar
Incorrect number of words per line in the potential file.

\sphinxlineitem{\sphinxstyleemphasis{Incorrect format in TMD target file}}
\sphinxAtStartPar
Format of file read by fix tmd command is incorrect.

\sphinxlineitem{\sphinxstyleemphasis{Incorrect format in Tersoff potential file}}
\sphinxAtStartPar
Incorrect number of words per line in the potential file.

\sphinxlineitem{\sphinxstyleemphasis{Incorrect format in Vashishta potential file}}
\sphinxAtStartPar
Incorrect number of words per line in the potential file.

\sphinxlineitem{\sphinxstyleemphasis{Incorrect format in coul/streitz potential file}}
\sphinxAtStartPar
Incorrect number of words per line in the potential file.

\sphinxlineitem{\sphinxstyleemphasis{Incorrect format in nb3b/harmonic potential file}}
\sphinxAtStartPar
Incorrect number of words per line in the potential file.

\sphinxlineitem{\sphinxstyleemphasis{Incorrect integer value in Bodies section of data file}}
\sphinxAtStartPar
See page for body style.

\sphinxlineitem{\sphinxstyleemphasis{Incorrect multiplicity arg for dihedral coefficients}}
\sphinxAtStartPar
Self\sphinxhyphen{}explanatory.  Check the input script or data file.

\sphinxlineitem{\sphinxstyleemphasis{Incorrect number of elements in potential file}}
\sphinxAtStartPar
Self\sphinxhyphen{}explanatory.

\sphinxlineitem{\sphinxstyleemphasis{Incorrect rigid body format in fix rigid file}}
\sphinxAtStartPar
The number of fields per line is not what expected.

\sphinxlineitem{\sphinxstyleemphasis{Incorrect rigid body format in fix rigid/small file}}
\sphinxAtStartPar
The number of fields per line is not what expected.

\sphinxlineitem{\sphinxstyleemphasis{Incorrect sign arg for dihedral coefficients}}
\sphinxAtStartPar
Self\sphinxhyphen{}explanatory.  Check the input script or data file.

\sphinxlineitem{\sphinxstyleemphasis{Incorrect table format check for element types}}
\sphinxAtStartPar
Self\sphinxhyphen{}explanatory.

\sphinxlineitem{\sphinxstyleemphasis{Incorrect velocity format in data file}}
\sphinxAtStartPar
Each atom style defines a format for the Velocity section
of the data file.  The read\sphinxhyphen{}in lines do not match.

\sphinxlineitem{\sphinxstyleemphasis{Incorrect weight arg for dihedral coefficients}}
\sphinxAtStartPar
Self\sphinxhyphen{}explanatory.  Check the input script or data file.

\sphinxlineitem{\sphinxstyleemphasis{Index between variable brackets must be positive}}
\sphinxAtStartPar
Self\sphinxhyphen{}explanatory.

\sphinxlineitem{\sphinxstyleemphasis{Indexed per\sphinxhyphen{}atom vector in variable formula without atom map}}
\sphinxAtStartPar
Accessing a value from an atom vector requires the ability to lookup
an atom index, which is provided by an atom map.  An atom map does not
exist (by default) for non\sphinxhyphen{}molecular problems.  Using the atom\_modify
map command will force an atom map to be created.

\sphinxlineitem{\sphinxstyleemphasis{Initial temperatures not all set in fix ttm}}
\sphinxAtStartPar
Self\sphinxhyphen{}explanatory.

\sphinxlineitem{\sphinxstyleemphasis{Input line quote not followed by white\sphinxhyphen{}space}}
\sphinxAtStartPar
An end quote must be followed by white\sphinxhyphen{}space.

\sphinxlineitem{\sphinxstyleemphasis{Insertion region extends outside simulation box}}
\sphinxAtStartPar
Self\sphinxhyphen{}explanatory.

\sphinxlineitem{\sphinxstyleemphasis{Insufficient Jacobi rotations for POEMS body}}
\sphinxAtStartPar
Eigensolve for rigid body was not sufficiently accurate.

\sphinxlineitem{\sphinxstyleemphasis{Insufficient Jacobi rotations for body nparticle}}
\sphinxAtStartPar
Eigensolve for rigid body was not sufficiently accurate.

\sphinxlineitem{\sphinxstyleemphasis{Insufficient Jacobi rotations for rigid body}}
\sphinxAtStartPar
Eigensolve for rigid body was not sufficiently accurate.

\sphinxlineitem{\sphinxstyleemphasis{Insufficient Jacobi rotations for rigid molecule}}
\sphinxAtStartPar
Eigensolve for rigid body was not sufficiently accurate.

\sphinxlineitem{\sphinxstyleemphasis{Insufficient Jacobi rotations for triangle}}
\sphinxAtStartPar
The calculation of the inertia tensor of the triangle failed.  This
should not happen if it is a reasonably shaped triangle.

\sphinxlineitem{\sphinxstyleemphasis{Insufficient memory on accelerator}}
\sphinxAtStartPar
There is insufficient memory on one of the devices specified for the gpu
package

\sphinxlineitem{\sphinxstyleemphasis{Internal error in atom\_style body}}
\sphinxAtStartPar
This error should not occur.  Contact the developers.

\sphinxlineitem{\sphinxstyleemphasis{Invalid \sphinxhyphen{}reorder N value}}
\sphinxAtStartPar
Self\sphinxhyphen{}explanatory.

\sphinxlineitem{\sphinxstyleemphasis{Invalid Angles section in molecule file}}
\sphinxAtStartPar
Self\sphinxhyphen{}explanatory.

\sphinxlineitem{\sphinxstyleemphasis{Invalid Bonds section in molecule file}}
\sphinxAtStartPar
Self\sphinxhyphen{}explanatory.

\sphinxlineitem{\sphinxstyleemphasis{Invalid Boolean syntax in if command}}
\sphinxAtStartPar
Self\sphinxhyphen{}explanatory.

\sphinxlineitem{\sphinxstyleemphasis{Invalid Charges section in molecule file}}
\sphinxAtStartPar
Self\sphinxhyphen{}explanatory.

\sphinxlineitem{\sphinxstyleemphasis{Invalid Coords section in molecule file}}
\sphinxAtStartPar
Self\sphinxhyphen{}explanatory.

\sphinxlineitem{\sphinxstyleemphasis{Invalid Diameters section in molecule file}}
\sphinxAtStartPar
Self\sphinxhyphen{}explanatory.

\sphinxlineitem{\sphinxstyleemphasis{Invalid Dihedrals section in molecule file}}
\sphinxAtStartPar
Self\sphinxhyphen{}explanatory.

\sphinxlineitem{\sphinxstyleemphasis{Invalid Impropers section in molecule file}}
\sphinxAtStartPar
Self\sphinxhyphen{}explanatory.

\sphinxlineitem{\sphinxstyleemphasis{Invalid Kokkos command\sphinxhyphen{}line args}}
\sphinxAtStartPar
Self\sphinxhyphen{}explanatory.  See Section 2.7 of the manual for details.

\sphinxlineitem{\sphinxstyleemphasis{Invalid LAMMPS restart file}}
\sphinxAtStartPar
The file does not appear to be a LAMMPS restart file since
it does not contain the correct magic string at the beginning.

\sphinxlineitem{\sphinxstyleemphasis{Invalid Masses section in molecule file}}
\sphinxAtStartPar
Self\sphinxhyphen{}explanatory.

\sphinxlineitem{\sphinxstyleemphasis{Invalid molecule ID in molecule file}}
\sphinxAtStartPar
Molecule ID must be a non\sphinxhyphen{}zero positive integer.

\sphinxlineitem{\sphinxstyleemphasis{Invalid Molecules section in molecule file}}
\sphinxAtStartPar
Self\sphinxhyphen{}explanatory.

\sphinxlineitem{\sphinxstyleemphasis{Invalid REAX atom type}}
\sphinxAtStartPar
There is a mis\sphinxhyphen{}match between LAMMPS atom types and the elements
listed in the ReaxFF force field file.

\sphinxlineitem{\sphinxstyleemphasis{Invalid Special Bond Counts section in molecule file}}
\sphinxAtStartPar
Self\sphinxhyphen{}explanatory.

\sphinxlineitem{\sphinxstyleemphasis{Invalid Types section in molecule file}}
\sphinxAtStartPar
Self\sphinxhyphen{}explanatory.

\sphinxlineitem{\sphinxstyleemphasis{Invalid angle count in molecule file}}
\sphinxAtStartPar
Self\sphinxhyphen{}explanatory.

\sphinxlineitem{\sphinxstyleemphasis{Invalid angle table length}}
\sphinxAtStartPar
Length must be 2 or greater.

\sphinxlineitem{\sphinxstyleemphasis{Invalid angle type in Angles section of data file}}
\sphinxAtStartPar
Angle type must be positive integer and within range of specified angle
types.

\sphinxlineitem{\sphinxstyleemphasis{Invalid angle type in Angles section of molecule file}}
\sphinxAtStartPar
Self\sphinxhyphen{}explanatory.

\sphinxlineitem{\sphinxstyleemphasis{Invalid angle type index for fix shake}}
\sphinxAtStartPar
Self\sphinxhyphen{}explanatory.

\sphinxlineitem{\sphinxstyleemphasis{Invalid args for non\sphinxhyphen{}hybrid pair coefficients}}
\sphinxAtStartPar
“NULL” is only supported in pair\_coeff calls when using pair hybrid

\sphinxlineitem{\sphinxstyleemphasis{Invalid argument to factorial \%d}}
\sphinxAtStartPar
N must be \textgreater{}= 0 and \textless{}= 167, otherwise the factorial result is too
large.

\sphinxlineitem{\sphinxstyleemphasis{Invalid atom ID in \%s section of data file}}
\sphinxAtStartPar
An atom in a section of the data file being read by fix property/atom
has an invalid atom ID that is \textless{}= 0 or \textgreater{} the maximum existing atom ID.

\sphinxlineitem{\sphinxstyleemphasis{Invalid atom ID in Angles section of data file}}
\sphinxAtStartPar
Atom IDs must be positive integers and within range of defined
atoms.

\sphinxlineitem{\sphinxstyleemphasis{Invalid atom ID in Angles section of molecule file}}
\sphinxAtStartPar
Self\sphinxhyphen{}explanatory.

\sphinxlineitem{\sphinxstyleemphasis{Invalid atom ID in Atoms section of data file}}
\sphinxAtStartPar
Atom IDs must be positive integers.

\sphinxlineitem{\sphinxstyleemphasis{Invalid atom ID in Bodies section of data file}}
\sphinxAtStartPar
Atom IDs must be positive integers and within range of defined
atoms.

\sphinxlineitem{\sphinxstyleemphasis{Invalid atom ID in Bonds section of data file}}
\sphinxAtStartPar
Atom IDs must be positive integers and within range of defined
atoms.

\sphinxlineitem{\sphinxstyleemphasis{Invalid atom ID in Bonds section of molecule file}}
\sphinxAtStartPar
Self\sphinxhyphen{}explanatory.

\sphinxlineitem{\sphinxstyleemphasis{Invalid atom ID in Bonus section of data file}}
\sphinxAtStartPar
Atom IDs must be positive integers and within range of defined
atoms.

\sphinxlineitem{\sphinxstyleemphasis{Invalid atom ID in Dihedrals section of data file}}
\sphinxAtStartPar
Atom IDs must be positive integers and within range of defined
atoms.

\sphinxlineitem{\sphinxstyleemphasis{Invalid atom ID in Fragments section of molecule file}}
\sphinxAtStartPar
Self\sphinxhyphen{}explanatory.

\sphinxlineitem{\sphinxstyleemphasis{Invalid atom ID in Impropers section of data file}}
\sphinxAtStartPar
Atom IDs must be positive integers and within range of defined
atoms.

\sphinxlineitem{\sphinxstyleemphasis{Invalid atom ID in Velocities section of data file}}
\sphinxAtStartPar
Atom IDs must be positive integers and within range of defined
atoms.

\sphinxlineitem{\sphinxstyleemphasis{Invalid atom ID in dihedrals section of molecule file}}
\sphinxAtStartPar
Self\sphinxhyphen{}explanatory.

\sphinxlineitem{\sphinxstyleemphasis{Invalid atom ID in impropers section of molecule file}}
\sphinxAtStartPar
Self\sphinxhyphen{}explanatory.

\sphinxlineitem{\sphinxstyleemphasis{Invalid atom ID in variable file}}
\sphinxAtStartPar
Self\sphinxhyphen{}explanatory.

\sphinxlineitem{\sphinxstyleemphasis{Invalid atom IDs in neb file}}
\sphinxAtStartPar
An ID in the file was not found in the system.

\sphinxlineitem{\sphinxstyleemphasis{Invalid atom diameter in molecule file}}
\sphinxAtStartPar
Diameters must be \textgreater{}= 0.0.

\sphinxlineitem{\sphinxstyleemphasis{Invalid atom mass for fix shake}}
\sphinxAtStartPar
Mass specified in fix shake command must be \textgreater{} 0.0.

\sphinxlineitem{\sphinxstyleemphasis{Invalid atom mass in molecule file}}
\sphinxAtStartPar
Masses must be \textgreater{} 0.0.

\sphinxlineitem{\sphinxstyleemphasis{Invalid atom type in Atoms section of data file}}
\sphinxAtStartPar
Atom types must range from 1 to specified \# of types.

\sphinxlineitem{\sphinxstyleemphasis{Invalid atom type in create\_atoms command}}
\sphinxAtStartPar
The create\_box command specified the range of valid atom types.
An invalid type is being requested.

\sphinxlineitem{\sphinxstyleemphasis{Invalid atom type in create\_atoms mol command}}
\sphinxAtStartPar
The atom types in the defined molecule are added to the value
specified in the create\_atoms command, as an offset.  The final value
for each atom must be between 1 to N, where N is the number of atom
types.

\sphinxlineitem{\sphinxstyleemphasis{Invalid atom type in fix atom/swap command}}
\sphinxAtStartPar
The atom type specified in the atom/swap command does not exist.

\sphinxlineitem{\sphinxstyleemphasis{Invalid atom type in fix bond/create command}}
\sphinxAtStartPar
Self\sphinxhyphen{}explanatory.

\sphinxlineitem{\sphinxstyleemphasis{Invalid atom type in fix deposit command}}
\sphinxAtStartPar
Self\sphinxhyphen{}explanatory.

\sphinxlineitem{\sphinxstyleemphasis{Invalid atom type in fix deposit mol command}}
\sphinxAtStartPar
The atom types in the defined molecule are added to the value
specified in the create\_atoms command, as an offset.  The final value
for each atom must be between 1 to N, where N is the number of atom
types.

\sphinxlineitem{\sphinxstyleemphasis{Invalid atom type in fix gcmc command}}
\sphinxAtStartPar
The atom type specified in the gcmc command does not exist.

\sphinxlineitem{\sphinxstyleemphasis{Invalid atom type in fix pour command}}
\sphinxAtStartPar
Self\sphinxhyphen{}explanatory.

\sphinxlineitem{\sphinxstyleemphasis{Invalid atom type in fix pour mol command}}
\sphinxAtStartPar
The atom types in the defined molecule are added to the value
specified in the create\_atoms command, as an offset.  The final value
for each atom must be between 1 to N, where N is the number of atom
types.

\sphinxlineitem{\sphinxstyleemphasis{Invalid atom type in molecule file}}
\sphinxAtStartPar
Atom types must range from 1 to specified \# of types.

\sphinxlineitem{\sphinxstyleemphasis{Invalid atom type in neighbor exclusion list}}
\sphinxAtStartPar
Atom types must range from 1 to Ntypes inclusive.

\sphinxlineitem{\sphinxstyleemphasis{Invalid atom type index for fix shake}}
\sphinxAtStartPar
Atom types must range from 1 to Ntypes inclusive.

\sphinxlineitem{\sphinxstyleemphasis{Invalid atom types in pair\_write command}}
\sphinxAtStartPar
Atom types must range from 1 to Ntypes inclusive.

\sphinxlineitem{\sphinxstyleemphasis{Invalid atom vector in variable formula}}
\sphinxAtStartPar
The atom vector is not recognized.

\sphinxlineitem{\sphinxstyleemphasis{Invalid atom\_style body command}}
\sphinxAtStartPar
No body style argument was provided.

\sphinxlineitem{\sphinxstyleemphasis{Invalid atom\_style command}}
\sphinxAtStartPar
Self\sphinxhyphen{}explanatory.

\sphinxlineitem{\sphinxstyleemphasis{Invalid attribute in dump custom command}}
\sphinxAtStartPar
Self\sphinxhyphen{}explanatory.

\sphinxlineitem{\sphinxstyleemphasis{Invalid attribute in dump local command}}
\sphinxAtStartPar
Self\sphinxhyphen{}explanatory.

\sphinxlineitem{\sphinxstyleemphasis{Invalid attribute in dump modify command}}
\sphinxAtStartPar
Self\sphinxhyphen{}explanatory.

\sphinxlineitem{\sphinxstyleemphasis{Invalid basis setting in create\_atoms command}}
\sphinxAtStartPar
The basis index must be between 1 to N where N is the number of basis
atoms in the lattice.  The type index must be between 1 to N where N
is the number of atom types.

\sphinxlineitem{\sphinxstyleemphasis{Invalid basis setting in fix append/atoms command}}
\sphinxAtStartPar
The basis index must be between 1 to N where N is the number of basis
atoms in the lattice.  The type index must be between 1 to N where N
is the number of atom types.

\sphinxlineitem{\sphinxstyleemphasis{Invalid bin bounds in compute chunk/atom}}
\sphinxAtStartPar
The lo/hi values are inconsistent.

\sphinxlineitem{\sphinxstyleemphasis{Invalid bin bounds in fix ave/spatial}}
\sphinxAtStartPar
The lo/hi values are inconsistent.

\sphinxlineitem{\sphinxstyleemphasis{Invalid body nparticle command}}
\sphinxAtStartPar
Arguments in atom\sphinxhyphen{}style command are not correct.

\sphinxlineitem{\sphinxstyleemphasis{Invalid bond count in molecule file}}
\sphinxAtStartPar
Self\sphinxhyphen{}explanatory.

\sphinxlineitem{\sphinxstyleemphasis{Invalid bond table length}}
\sphinxAtStartPar
Length must be 2 or greater.

\sphinxlineitem{\sphinxstyleemphasis{Invalid bond type in Bonds section of data file}}
\sphinxAtStartPar
Bond type must be positive integer and within range of specified bond
types.

\sphinxlineitem{\sphinxstyleemphasis{Invalid bond type in Bonds section of molecule file}}
\sphinxAtStartPar
Self\sphinxhyphen{}explanatory.

\sphinxlineitem{\sphinxstyleemphasis{Invalid bond type in create\_bonds command}}
\sphinxAtStartPar
Self\sphinxhyphen{}explanatory.

\sphinxlineitem{\sphinxstyleemphasis{Invalid bond type in fix bond/break command}}
\sphinxAtStartPar
Self\sphinxhyphen{}explanatory.

\sphinxlineitem{\sphinxstyleemphasis{Invalid bond type in fix bond/create command}}
\sphinxAtStartPar
Self\sphinxhyphen{}explanatory.

\sphinxlineitem{\sphinxstyleemphasis{Invalid bond type index for fix shake}}
\sphinxAtStartPar
Self\sphinxhyphen{}explanatory.  Check the fix shake command in the input script.

\sphinxlineitem{\sphinxstyleemphasis{Invalid coeffs for this dihedral style}}
\sphinxAtStartPar
Cannot set class 2 coeffs in data file for this dihedral style.

\sphinxlineitem{\sphinxstyleemphasis{Invalid color in dump\_modify command}}
\sphinxAtStartPar
The specified color name was not in the list of recognized colors.
See the dump\_modify doc page.

\sphinxlineitem{\sphinxstyleemphasis{Invalid color map min/max values}}
\sphinxAtStartPar
The min/max values are not consistent with either each other or
with values in the color map.

\sphinxlineitem{\sphinxstyleemphasis{Invalid command\sphinxhyphen{}line argument}}
\sphinxAtStartPar
One or more command\sphinxhyphen{}line arguments is invalid.  Check the syntax of
the command you are using to launch LAMMPS.

\sphinxlineitem{\sphinxstyleemphasis{Invalid compute ID in variable formula}}
\sphinxAtStartPar
The compute is not recognized.

\sphinxlineitem{\sphinxstyleemphasis{Invalid create\_atoms rotation vector for 2d model}}
\sphinxAtStartPar
The rotation vector can only have a z component.

\sphinxlineitem{\sphinxstyleemphasis{Invalid custom OpenCL parameter string.}}
\sphinxAtStartPar
There are not enough or too many parameters in the custom string for package
GPU.

\sphinxlineitem{\sphinxstyleemphasis{Invalid cutoff in comm\_modify command}}
\sphinxAtStartPar
Specified cutoff must be \textgreater{}= 0.0.

\sphinxlineitem{\sphinxstyleemphasis{Invalid cutoffs in pair\_write command}}
\sphinxAtStartPar
Inner cutoff must be larger than 0.0 and less than outer cutoff.

\sphinxlineitem{\sphinxstyleemphasis{Invalid d1 or d2 value for pair colloid coeff}}
\sphinxAtStartPar
Neither d1 or d2 can be \textless{} 0.

\sphinxlineitem{\sphinxstyleemphasis{Invalid data file section: Angle Coeffs}}
\sphinxAtStartPar
Atom style does not allow angles.

\sphinxlineitem{\sphinxstyleemphasis{Invalid data file section: AngleAngle Coeffs}}
\sphinxAtStartPar
Atom style does not allow impropers.

\sphinxlineitem{\sphinxstyleemphasis{Invalid data file section: AngleAngleTorsion Coeffs}}
\sphinxAtStartPar
Atom style does not allow dihedrals.

\sphinxlineitem{\sphinxstyleemphasis{Invalid data file section: AngleTorsion Coeffs}}
\sphinxAtStartPar
Atom style does not allow dihedrals.

\sphinxlineitem{\sphinxstyleemphasis{Invalid data file section: Angles}}
\sphinxAtStartPar
Atom style does not allow angles.

\sphinxlineitem{\sphinxstyleemphasis{Invalid data file section: Bodies}}
\sphinxAtStartPar
Atom style does not allow bodies.

\sphinxlineitem{\sphinxstyleemphasis{Invalid data file section: Bond Coeffs}}
\sphinxAtStartPar
Atom style does not allow bonds.

\sphinxlineitem{\sphinxstyleemphasis{Invalid data file section: BondAngle Coeffs}}
\sphinxAtStartPar
Atom style does not allow angles.

\sphinxlineitem{\sphinxstyleemphasis{Invalid data file section: BondBond Coeffs}}
\sphinxAtStartPar
Atom style does not allow angles.

\sphinxlineitem{\sphinxstyleemphasis{Invalid data file section: BondBond13 Coeffs}}
\sphinxAtStartPar
Atom style does not allow dihedrals.

\sphinxlineitem{\sphinxstyleemphasis{Invalid data file section: Bonds}}
\sphinxAtStartPar
Atom style does not allow bonds.

\sphinxlineitem{\sphinxstyleemphasis{Invalid data file section: Dihedral Coeffs}}
\sphinxAtStartPar
Atom style does not allow dihedrals.

\sphinxlineitem{\sphinxstyleemphasis{Invalid data file section: Dihedrals}}
\sphinxAtStartPar
Atom style does not allow dihedrals.

\sphinxlineitem{\sphinxstyleemphasis{Invalid data file section: Ellipsoids}}
\sphinxAtStartPar
Atom style does not allow ellipsoids.

\sphinxlineitem{\sphinxstyleemphasis{Invalid data file section: EndBondTorsion Coeffs}}
\sphinxAtStartPar
Atom style does not allow dihedrals.

\sphinxlineitem{\sphinxstyleemphasis{Invalid data file section: Improper Coeffs}}
\sphinxAtStartPar
Atom style does not allow impropers.

\sphinxlineitem{\sphinxstyleemphasis{Invalid data file section: Impropers}}
\sphinxAtStartPar
Atom style does not allow impropers.

\sphinxlineitem{\sphinxstyleemphasis{Invalid data file section: Lines}}
\sphinxAtStartPar
Atom style does not allow lines.

\sphinxlineitem{\sphinxstyleemphasis{Invalid data file section: MiddleBondTorsion Coeffs}}
\sphinxAtStartPar
Atom style does not allow dihedrals.

\sphinxlineitem{\sphinxstyleemphasis{Invalid data file section: Triangles}}
\sphinxAtStartPar
Atom style does not allow triangles.

\sphinxlineitem{\sphinxstyleemphasis{Invalid delta\_conf in tad command}}
\sphinxAtStartPar
The value must be between 0 and 1 inclusive.

\sphinxlineitem{\sphinxstyleemphasis{Invalid density in Atoms section of data file}}
\sphinxAtStartPar
Density value cannot be \textless{}= 0.0.

\sphinxlineitem{\sphinxstyleemphasis{Invalid density in set command}}
\sphinxAtStartPar
Density must be \textgreater{} 0.0.

\sphinxlineitem{\sphinxstyleemphasis{Invalid diameter in set command}}
\sphinxAtStartPar
Self\sphinxhyphen{}explanatory.

\sphinxlineitem{\sphinxstyleemphasis{Invalid dihedral count in molecule file}}
\sphinxAtStartPar
Self\sphinxhyphen{}explanatory.

\sphinxlineitem{\sphinxstyleemphasis{Invalid dihedral type in Dihedrals section of data file}}
\sphinxAtStartPar
Dihedral type must be positive integer and within range of specified
dihedral types.

\sphinxlineitem{\sphinxstyleemphasis{Invalid dihedral type in dihedrals section of molecule file}}
\sphinxAtStartPar
Self\sphinxhyphen{}explanatory.

\sphinxlineitem{\sphinxstyleemphasis{Invalid dipole length in set command}}
\sphinxAtStartPar
Self\sphinxhyphen{}explanatory.

\sphinxlineitem{\sphinxstyleemphasis{Invalid displace\_atoms rotate axis for 2d}}
\sphinxAtStartPar
Axis must be in z direction.

\sphinxlineitem{\sphinxstyleemphasis{Invalid dump dcd filename}}
\sphinxAtStartPar
Filenames used with the dump dcd style cannot be binary or compressed
or cause multiple files to be written.

\sphinxlineitem{\sphinxstyleemphasis{Invalid dump frequency}}
\sphinxAtStartPar
Dump frequency must be 1 or greater.

\sphinxlineitem{\sphinxstyleemphasis{Invalid dump image element name}}
\sphinxAtStartPar
The specified element name was not in the standard list of elements.
See the dump\_modify doc page.

\sphinxlineitem{\sphinxstyleemphasis{Invalid dump image filename}}
\sphinxAtStartPar
The file produced by dump image cannot be binary and must
be for a single processor.

\sphinxlineitem{\sphinxstyleemphasis{Invalid dump image theta value}}
\sphinxAtStartPar
Theta must be between 0.0 and 180.0 inclusive.

\sphinxlineitem{\sphinxstyleemphasis{Invalid dump image zoom value}}
\sphinxAtStartPar
Zoom value must be \textgreater{} 0.0.

\sphinxlineitem{\sphinxstyleemphasis{Invalid dump movie filename}}
\sphinxAtStartPar
The file produced by dump movie cannot be binary or compressed
and must be a single file for a single processor.

\sphinxlineitem{\sphinxstyleemphasis{Invalid dump xtc filename}}
\sphinxAtStartPar
Filenames used with the dump xtc style cannot be binary or compressed
or cause multiple files to be written.

\sphinxlineitem{\sphinxstyleemphasis{Invalid dump xyz filename}}
\sphinxAtStartPar
Filenames used with the dump xyz style cannot be binary or cause files
to be written by each processor.

\sphinxlineitem{\sphinxstyleemphasis{Invalid dump\_modify threshold operator}}
\sphinxAtStartPar
Operator keyword used for threshold specification in not recognized.

\sphinxlineitem{\sphinxstyleemphasis{Invalid entry in \sphinxhyphen{}reorder file}}
\sphinxAtStartPar
Self\sphinxhyphen{}explanatory.

\sphinxlineitem{\sphinxstyleemphasis{Invalid fix ID in variable formula}}
\sphinxAtStartPar
The fix is not recognized.

\sphinxlineitem{\sphinxstyleemphasis{Invalid fix ave/time off column}}
\sphinxAtStartPar
Self\sphinxhyphen{}explanatory.

\sphinxlineitem{\sphinxstyleemphasis{Invalid fix box/relax command for a 2d simulation}}
\sphinxAtStartPar
Fix box/relax styles involving the z dimension cannot be used in
a 2d simulation.

\sphinxlineitem{\sphinxstyleemphasis{Invalid fix box/relax command pressure settings}}
\sphinxAtStartPar
If multiple dimensions are coupled, those dimensions must be specified.

\sphinxlineitem{\sphinxstyleemphasis{Invalid fix box/relax pressure settings}}
\sphinxAtStartPar
Settings for coupled dimensions must be the same.

\sphinxlineitem{\sphinxstyleemphasis{Invalid fix halt attribute}}
\sphinxAtStartPar
Self\sphinxhyphen{}explanatory.

\sphinxlineitem{\sphinxstyleemphasis{Invalid fix halt operator}}
\sphinxAtStartPar
Self\sphinxhyphen{}explanatory.

\sphinxlineitem{\sphinxstyleemphasis{Invalid fix nvt/npt/nph command for a 2d simulation}}
\sphinxAtStartPar
Cannot control z dimension in a 2d model.

\sphinxlineitem{\sphinxstyleemphasis{Invalid fix nvt/npt/nph command pressure settings}}
\sphinxAtStartPar
If multiple dimensions are coupled, those dimensions must be
specified.

\sphinxlineitem{\sphinxstyleemphasis{Invalid fix nvt/npt/nph pressure settings}}
\sphinxAtStartPar
Settings for coupled dimensions must be the same.

\sphinxlineitem{\sphinxstyleemphasis{Invalid fix press/berendsen for a 2d simulation}}
\sphinxAtStartPar
The z component of pressure cannot be controlled for a 2d model.

\sphinxlineitem{\sphinxstyleemphasis{Invalid fix press/berendsen pressure settings}}
\sphinxAtStartPar
Settings for coupled dimensions must be the same.

\sphinxlineitem{\sphinxstyleemphasis{Invalid fix qeq parameter file}}
\sphinxAtStartPar
Element index \textgreater{} number of atom types.

\sphinxlineitem{\sphinxstyleemphasis{Invalid fix rigid npt/nph command for a 2d simulation}}
\sphinxAtStartPar
Cannot control z dimension in a 2d model.

\sphinxlineitem{\sphinxstyleemphasis{Invalid fix rigid npt/nph command pressure settings}}
\sphinxAtStartPar
If multiple dimensions are coupled, those dimensions must be
specified.

\sphinxlineitem{\sphinxstyleemphasis{Invalid fix rigid/small npt/nph command for a 2d simulation}}
\sphinxAtStartPar
Cannot control z dimension in a 2d model.

\sphinxlineitem{\sphinxstyleemphasis{Invalid fix rigid/small npt/nph command pressure settings}}
\sphinxAtStartPar
If multiple dimensions are coupled, those dimensions must be
specified.

\sphinxlineitem{\sphinxstyleemphasis{Invalid flag in force field section of restart file}}
\sphinxAtStartPar
Unrecognized entry in restart file.

\sphinxlineitem{\sphinxstyleemphasis{Invalid flag in header section of restart file}}
\sphinxAtStartPar
Unrecognized entry in restart file.

\sphinxlineitem{\sphinxstyleemphasis{Invalid flag in peratom section of restart file}}
\sphinxAtStartPar
The format of this section of the file is not correct.

\sphinxlineitem{\sphinxstyleemphasis{Invalid flag in type arrays section of restart file}}
\sphinxAtStartPar
Unrecognized entry in restart file.

\sphinxlineitem{\sphinxstyleemphasis{Invalid frequency in temper command}}
\sphinxAtStartPar
Nevery must be \textgreater{} 0.

\sphinxlineitem{\sphinxstyleemphasis{Invalid group ID in neigh\_modify command}}
\sphinxAtStartPar
A group ID used in the neigh\_modify command does not exist.

\sphinxlineitem{\sphinxstyleemphasis{Invalid group function in variable formula}}
\sphinxAtStartPar
Group function is not recognized.

\sphinxlineitem{\sphinxstyleemphasis{Invalid group in comm\_modify command}}
\sphinxAtStartPar
Self\sphinxhyphen{}explanatory.

\sphinxlineitem{\sphinxstyleemphasis{Invalid image up vector}}
\sphinxAtStartPar
Up vector cannot be (0,0,0).

\sphinxlineitem{\sphinxstyleemphasis{Invalid immediate variable}}
\sphinxAtStartPar
Syntax of immediate value is incorrect.

\sphinxlineitem{\sphinxstyleemphasis{Invalid improper count in molecule file}}
\sphinxAtStartPar
Self\sphinxhyphen{}explanatory.

\sphinxlineitem{\sphinxstyleemphasis{Invalid improper type in Impropers section of data file}}
\sphinxAtStartPar
Improper type must be positive integer and within range of specified
improper types.

\sphinxlineitem{\sphinxstyleemphasis{Invalid improper type in impropers section of molecule file}}
\sphinxAtStartPar
Self\sphinxhyphen{}explanatory.

\sphinxlineitem{\sphinxstyleemphasis{Invalid index for non\sphinxhyphen{}body particles in compute body/local command}}
\sphinxAtStartPar
Only indices 1,2,3 can be used for non\sphinxhyphen{}body particles.

\sphinxlineitem{\sphinxstyleemphasis{Invalid index in compute body/local command}}
\sphinxAtStartPar
Self\sphinxhyphen{}explanatory.

\sphinxlineitem{\sphinxstyleemphasis{Invalid is\_active() function in variable formula}}
\sphinxAtStartPar
Self\sphinxhyphen{}explanatory.

\sphinxlineitem{\sphinxstyleemphasis{Invalid is\_available() function in variable formula}}
\sphinxAtStartPar
Self\sphinxhyphen{}explanatory.

\sphinxlineitem{\sphinxstyleemphasis{Invalid is\_defined() function in variable formula}}
\sphinxAtStartPar
Self\sphinxhyphen{}explanatory.

\sphinxlineitem{\sphinxstyleemphasis{Invalid keyword in angle table parameters}}
\sphinxAtStartPar
Self\sphinxhyphen{}explanatory.

\sphinxlineitem{\sphinxstyleemphasis{Invalid keyword in bond table parameters}}
\sphinxAtStartPar
Self\sphinxhyphen{}explanatory.

\sphinxlineitem{\sphinxstyleemphasis{Invalid keyword in compute angle/local command}}
\sphinxAtStartPar
Self\sphinxhyphen{}explanatory.

\sphinxlineitem{\sphinxstyleemphasis{Invalid keyword in compute bond/local command}}
\sphinxAtStartPar
Self\sphinxhyphen{}explanatory.

\sphinxlineitem{\sphinxstyleemphasis{Invalid keyword in compute dihedral/local command}}
\sphinxAtStartPar
Self\sphinxhyphen{}explanatory.

\sphinxlineitem{\sphinxstyleemphasis{Invalid keyword in compute improper/local command}}
\sphinxAtStartPar
Self\sphinxhyphen{}explanatory.

\sphinxlineitem{\sphinxstyleemphasis{Invalid keyword in compute pair/local command}}
\sphinxAtStartPar
Self\sphinxhyphen{}explanatory.

\sphinxlineitem{\sphinxstyleemphasis{Invalid keyword in compute property/atom command}}
\sphinxAtStartPar
Self\sphinxhyphen{}explanatory.

\sphinxlineitem{\sphinxstyleemphasis{Invalid keyword in compute property/chunk command}}
\sphinxAtStartPar
Self\sphinxhyphen{}explanatory.

\sphinxlineitem{\sphinxstyleemphasis{Invalid keyword in compute property/local command}}
\sphinxAtStartPar
Self\sphinxhyphen{}explanatory.

\sphinxlineitem{\sphinxstyleemphasis{Invalid keyword in dump cfg command}}
\sphinxAtStartPar
Self\sphinxhyphen{}explanatory.

\sphinxlineitem{\sphinxstyleemphasis{Invalid keyword in pair table parameters}}
\sphinxAtStartPar
Keyword used in list of table parameters is not recognized.

\sphinxlineitem{\sphinxstyleemphasis{Invalid length in set command}}
\sphinxAtStartPar
Self\sphinxhyphen{}explanatory.

\sphinxlineitem{\sphinxstyleemphasis{Invalid mass in set command}}
\sphinxAtStartPar
Self\sphinxhyphen{}explanatory.

\sphinxlineitem{\sphinxstyleemphasis{Invalid mass line in data file}}
\sphinxAtStartPar
Self\sphinxhyphen{}explanatory.

\sphinxlineitem{\sphinxstyleemphasis{Invalid mass value}}
\sphinxAtStartPar
Self\sphinxhyphen{}explanatory.

\sphinxlineitem{\sphinxstyleemphasis{Invalid math function in variable formula}}
\sphinxAtStartPar
Self\sphinxhyphen{}explanatory.

\sphinxlineitem{\sphinxstyleemphasis{Invalid math/group/special function in variable formula}}
\sphinxAtStartPar
Self\sphinxhyphen{}explanatory.

\sphinxlineitem{\sphinxstyleemphasis{Invalid option in lattice command for non\sphinxhyphen{}custom style}}
\sphinxAtStartPar
Certain lattice keywords are not supported unless the
lattice style is “custom”.

\sphinxlineitem{\sphinxstyleemphasis{Invalid order of forces within respa levels}}
\sphinxAtStartPar
For respa, ordering of force computations within respa levels must
obey certain rules.  E.g. bonds cannot be compute less frequently than
angles, pairwise forces cannot be computed less frequently than
kspace, etc.

\sphinxlineitem{\sphinxstyleemphasis{Invalid pair table cutoff}}
\sphinxAtStartPar
Cutoffs in pair\_coeff command are not valid with read\sphinxhyphen{}in pair table.

\sphinxlineitem{\sphinxstyleemphasis{Invalid pair table length}}
\sphinxAtStartPar
Length of read\sphinxhyphen{}in pair table is invalid

\sphinxlineitem{\sphinxstyleemphasis{Invalid param file for fix qeq/shielded}}
\sphinxAtStartPar
Invalid value of gamma.

\sphinxlineitem{\sphinxstyleemphasis{Invalid param file for fix qeq/slater}}
\sphinxAtStartPar
Zeta value is 0.0.

\sphinxlineitem{\sphinxstyleemphasis{Invalid partitions in processors part command}}
\sphinxAtStartPar
Valid partitions are numbered 1 to N and the sender and receiver
cannot be the same partition.

\sphinxlineitem{\sphinxstyleemphasis{Invalid python command}}
\sphinxAtStartPar
Self\sphinxhyphen{}explanatory.  Check the input script syntax and compare to the
documentation for the command.  You can use \sphinxhyphen{}echo screen as a
command\sphinxhyphen{}line option when running LAMMPS to see the offending line.

\sphinxlineitem{\sphinxstyleemphasis{Invalid radius in Atoms section of data file}}
\sphinxAtStartPar
Radius must be \textgreater{}= 0.0.

\sphinxlineitem{\sphinxstyleemphasis{Invalid random number seed in fix ttm command}}
\sphinxAtStartPar
Random number seed must be \textgreater{} 0.

\sphinxlineitem{\sphinxstyleemphasis{Invalid random number seed in set command}}
\sphinxAtStartPar
Random number seed must be \textgreater{} 0.

\sphinxlineitem{\sphinxstyleemphasis{Invalid replace values in compute reduce}}
\sphinxAtStartPar
Self\sphinxhyphen{}explanatory.

\sphinxlineitem{\sphinxstyleemphasis{Invalid rigid body ID in fix rigid file}}
\sphinxAtStartPar
The ID does not match the number of an existing ID of rigid bodies
that are defined by the fix rigid command.

\sphinxlineitem{\sphinxstyleemphasis{Invalid rigid body ID in fix rigid/small file}}
\sphinxAtStartPar
The ID does not match the number of an existing ID of rigid bodies
that are defined by the fix rigid/small command.

\sphinxlineitem{\sphinxstyleemphasis{Invalid run command N value}}
\sphinxAtStartPar
The number of timesteps must fit in a 32\sphinxhyphen{}bit integer.  If you want to
run for more steps than this, perform multiple shorter runs.

\sphinxlineitem{\sphinxstyleemphasis{Invalid run command start/stop value}}
\sphinxAtStartPar
Self\sphinxhyphen{}explanatory.

\sphinxlineitem{\sphinxstyleemphasis{Invalid run command upto value}}
\sphinxAtStartPar
Self\sphinxhyphen{}explanatory.

\sphinxlineitem{\sphinxstyleemphasis{Invalid seed for Marsaglia random \# generator}}
\sphinxAtStartPar
The initial seed for this random number generator must be a positive
integer less than or equal to 900 million.

\sphinxlineitem{\sphinxstyleemphasis{Invalid seed for Park random \# generator}}
\sphinxAtStartPar
The initial seed for this random number generator must be a positive
integer.

\sphinxlineitem{\sphinxstyleemphasis{Invalid shake angle type in molecule file}}
\sphinxAtStartPar
Self\sphinxhyphen{}explanatory.

\sphinxlineitem{\sphinxstyleemphasis{Invalid shake atom in molecule file}}
\sphinxAtStartPar
Self\sphinxhyphen{}explanatory.

\sphinxlineitem{\sphinxstyleemphasis{Invalid shake bond type in molecule file}}
\sphinxAtStartPar
Self\sphinxhyphen{}explanatory.

\sphinxlineitem{\sphinxstyleemphasis{Invalid shake flag in molecule file}}
\sphinxAtStartPar
Self\sphinxhyphen{}explanatory.

\sphinxlineitem{\sphinxstyleemphasis{Invalid shape in Ellipsoids section of data file}}
\sphinxAtStartPar
Self\sphinxhyphen{}explanatory.

\sphinxlineitem{\sphinxstyleemphasis{Invalid shape in Triangles section of data file}}
\sphinxAtStartPar
Two or more of the triangle corners are duplicate points.

\sphinxlineitem{\sphinxstyleemphasis{Invalid shape in set command}}
\sphinxAtStartPar
Self\sphinxhyphen{}explanatory.

\sphinxlineitem{\sphinxstyleemphasis{Invalid shear direction for fix wall/gran}}
\sphinxAtStartPar
Self\sphinxhyphen{}explanatory.

\sphinxlineitem{\sphinxstyleemphasis{Invalid special atom index in molecule file}}
\sphinxAtStartPar
Self\sphinxhyphen{}explanatory.

\sphinxlineitem{\sphinxstyleemphasis{Invalid special function in variable formula}}
\sphinxAtStartPar
Self\sphinxhyphen{}explanatory.

\sphinxlineitem{\sphinxstyleemphasis{Invalid style in pair\_write command}}
\sphinxAtStartPar
Self\sphinxhyphen{}explanatory.  Check the input script.

\sphinxlineitem{\sphinxstyleemphasis{Invalid syntax in variable formula}}
\sphinxAtStartPar
Self\sphinxhyphen{}explanatory.

\sphinxlineitem{\sphinxstyleemphasis{Invalid t\_event in prd command}}
\sphinxAtStartPar
Self\sphinxhyphen{}explanatory.

\sphinxlineitem{\sphinxstyleemphasis{Invalid t\_event in tad command}}
\sphinxAtStartPar
The value must be greater than 0.

\sphinxlineitem{\sphinxstyleemphasis{Invalid template atom in Atoms section of data file}}
\sphinxAtStartPar
The atom indices must be between 1 to N, where N is the number of
atoms in the template molecule the atom belongs to.

\sphinxlineitem{\sphinxstyleemphasis{Invalid template index in Atoms section of data file}}
\sphinxAtStartPar
The template indices must be between 1 to N, where N is the number of
molecules in the template.

\sphinxlineitem{\sphinxstyleemphasis{Invalid thermo keyword in variable formula}}
\sphinxAtStartPar
The keyword is not recognized.

\sphinxlineitem{\sphinxstyleemphasis{Invalid threads\_per\_atom specified.}}
\sphinxAtStartPar
For 3\sphinxhyphen{}body potentials on the GPU, the threads\_per\_atom setting cannot be
greater than 4 for NVIDIA GPUs.

\sphinxlineitem{\sphinxstyleemphasis{Invalid timestep reset for fix ave/atom}}
\sphinxAtStartPar
Resetting the timestep has invalidated the sequence of timesteps this
fix needs to process.

\sphinxlineitem{\sphinxstyleemphasis{Invalid timestep reset for fix ave/chunk}}
\sphinxAtStartPar
Resetting the timestep has invalidated the sequence of timesteps this
fix needs to process.

\sphinxlineitem{\sphinxstyleemphasis{Invalid timestep reset for fix ave/correlate}}
\sphinxAtStartPar
Resetting the timestep has invalidated the sequence of timesteps this
fix needs to process.

\sphinxlineitem{\sphinxstyleemphasis{Invalid timestep reset for fix ave/histo}}
\sphinxAtStartPar
Resetting the timestep has invalidated the sequence of timesteps this
fix needs to process.

\sphinxlineitem{\sphinxstyleemphasis{Invalid timestep reset for fix ave/spatial}}
\sphinxAtStartPar
Resetting the timestep has invalidated the sequence of timesteps this
fix needs to process.

\sphinxlineitem{\sphinxstyleemphasis{Invalid timestep reset for fix ave/time}}
\sphinxAtStartPar
Resetting the timestep has invalidated the sequence of timesteps this
fix needs to process.

\sphinxlineitem{\sphinxstyleemphasis{Invalid tmax in tad command}}
\sphinxAtStartPar
The value must be greater than 0.0.

\sphinxlineitem{\sphinxstyleemphasis{Invalid type for mass set}}
\sphinxAtStartPar
Mass command must set a type from 1\sphinxhyphen{}N where N is the number of atom
types.

\sphinxlineitem{\sphinxstyleemphasis{Invalid label2type() function syntax in variable formula}}
\sphinxAtStartPar
The first argument must be a label map kind (atom, bond, angle,
dihedral, or improper) and the second argument must be a valid type
label that has been assigned to a numeric type.

\sphinxlineitem{\sphinxstyleemphasis{Invalid use of library file() function}}
\sphinxAtStartPar
This function is called through the library interface.  This
error should not occur.  Contact the developers if it does.

\sphinxlineitem{\sphinxstyleemphasis{Invalid value in set command}}
\sphinxAtStartPar
The value specified for the setting is invalid, likely because it is
too small or too large.

\sphinxlineitem{\sphinxstyleemphasis{Invalid variable evaluation in variable formula}}
\sphinxAtStartPar
A variable used in a formula could not be evaluated.

\sphinxlineitem{\sphinxstyleemphasis{Invalid variable in next command}}
\sphinxAtStartPar
Self\sphinxhyphen{}explanatory.

\sphinxlineitem{\sphinxstyleemphasis{Invalid variable name}}
\sphinxAtStartPar
Variable name used in an input script line is invalid.

\sphinxlineitem{\sphinxstyleemphasis{Invalid variable name in variable formula}}
\sphinxAtStartPar
Variable name is not recognized.

\sphinxlineitem{\sphinxstyleemphasis{Invalid variable style in special function next}}
\sphinxAtStartPar
Only file\sphinxhyphen{}style or atomfile\sphinxhyphen{}style variables can be used with next().

\sphinxlineitem{\sphinxstyleemphasis{Invalid variable style with next command}}
\sphinxAtStartPar
Variable styles \sphinxstyleemphasis{equal} and \sphinxstyleemphasis{world} cannot be used in a next
command.

\sphinxlineitem{\sphinxstyleemphasis{Invalid volume in set command}}
\sphinxAtStartPar
Volume must be \textgreater{} 0.0.

\sphinxlineitem{\sphinxstyleemphasis{Invalid wiggle direction for fix wall/gran}}
\sphinxAtStartPar
Self\sphinxhyphen{}explanatory.

\sphinxlineitem{\sphinxstyleemphasis{Invoked angle equil angle on angle style none}}
\sphinxAtStartPar
Self\sphinxhyphen{}explanatory.

\sphinxlineitem{\sphinxstyleemphasis{Invoked angle single on angle style none}}
\sphinxAtStartPar
Self\sphinxhyphen{}explanatory.

\sphinxlineitem{\sphinxstyleemphasis{Invoked bond equil distance on bond style none}}
\sphinxAtStartPar
Self\sphinxhyphen{}explanatory.

\sphinxlineitem{\sphinxstyleemphasis{Invoked bond single on bond style none}}
\sphinxAtStartPar
Self\sphinxhyphen{}explanatory.

\sphinxlineitem{\sphinxstyleemphasis{Invoked pair single on pair style none}}
\sphinxAtStartPar
A command (e.g. a dump) attempted to invoke the single() function on a
pair style none, which is illegal.  You are probably attempting to
compute per\sphinxhyphen{}atom quantities with an undefined pair style.

\sphinxlineitem{\sphinxstyleemphasis{Invoking coulombic in pair style lj/coul requires atom attribute q}}
\sphinxAtStartPar
The atom style defined does not have this attribute.

\sphinxlineitem{\sphinxstyleemphasis{Invoking coulombic in pair style lj/long/dipole/long requires atom attribute q}}
\sphinxAtStartPar
The atom style defined does not have these attributes.

\sphinxlineitem{\sphinxstyleemphasis{KIM Simulator Model has no Model definition}}
\sphinxAtStartPar
There is no model definition (key: model\sphinxhyphen{}defn) in the KIM Simulator
Model.  Please contact the OpenKIM database maintainers to verify
and potentially correct this.

\sphinxlineitem{\sphinxstyleemphasis{KOKKOS package does not yet support comm\_style tiled}}
\sphinxAtStartPar
Self\sphinxhyphen{}explanatory.

\sphinxlineitem{\sphinxstyleemphasis{KOKKOS package requires a kokkos enabled atom\_style}}
\sphinxAtStartPar
Self\sphinxhyphen{}explanatory.

\sphinxlineitem{\sphinxstyleemphasis{KSpace accuracy must be \textgreater{} 0}}
\sphinxAtStartPar
The kspace accuracy designated in the input must be greater than zero.

\sphinxlineitem{\sphinxstyleemphasis{KSpace accuracy too large to estimate G vector}}
\sphinxAtStartPar
Reduce the accuracy request or specify gewald explicitly
via the kspace\_modify command.

\sphinxlineitem{\sphinxstyleemphasis{KSpace accuracy too low}}
\sphinxAtStartPar
Requested accuracy must be less than 1.0.

\sphinxlineitem{\sphinxstyleemphasis{KSpace solver requires a pair style}}
\sphinxAtStartPar
No pair style is defined.

\sphinxlineitem{\sphinxstyleemphasis{KSpace style does not yet support triclinic geometries}}
\sphinxAtStartPar
The specified kspace style does not allow for non\sphinxhyphen{}orthogonal
simulation boxes.

\sphinxlineitem{\sphinxstyleemphasis{KSpace style has not yet been set}}
\sphinxAtStartPar
Cannot use kspace\_modify command until a kspace style is set.

\sphinxlineitem{\sphinxstyleemphasis{KSpace style is incompatible with Pair style}}
\sphinxAtStartPar
Setting a kspace style requires that a pair style with matching
long\sphinxhyphen{}range Coulombic or dispersion components be used.

\sphinxlineitem{\sphinxstyleemphasis{Keyword \%s in MEAM parameter file not recognized}}
\sphinxAtStartPar
Self\sphinxhyphen{}explanatory.

\sphinxlineitem{\sphinxstyleemphasis{Kokkos has been compiled for CUDA but no GPUs are requested}}
\sphinxAtStartPar
One or more GPUs must be used when Kokkos is compiled for CUDA.

\sphinxlineitem{\sphinxstyleemphasis{Kspace\_modify mesh parameter must be all zero or all positive}}
\sphinxAtStartPar
Valid kspace mesh parameters are \textgreater{}0. The code will try to auto\sphinxhyphen{}detect
suitable values when all three mesh sizes are set to zero (the default).

\sphinxlineitem{\sphinxstyleemphasis{Kspace\_modify mesh/disp parameter must be all zero or all positive}}
\sphinxAtStartPar
Valid kspace mesh/disp parameters are \textgreater{}0. The code will try to auto\sphinxhyphen{}detect
suitable values when all three mesh sizes are set to zero \sphinxstylestrong{and}
the required accuracy via \sphinxstyleemphasis{force/disp/real} as well as
\sphinxstyleemphasis{force/disp/kspace} is set.

\sphinxlineitem{\sphinxstyleemphasis{Kspace style does not support compute group/group}}
\sphinxAtStartPar
Self\sphinxhyphen{}explanatory.

\sphinxlineitem{\sphinxstyleemphasis{Kspace style pppm/disp/tip4p requires newton on}}
\sphinxAtStartPar
Self\sphinxhyphen{}explanatory.

\sphinxlineitem{\sphinxstyleemphasis{Kspace style pppm/tip4p requires newton on}}
\sphinxAtStartPar
Self\sphinxhyphen{}explanatory.

\sphinxlineitem{\sphinxstyleemphasis{Kspace style requires atom attribute q}}
\sphinxAtStartPar
The atom style defined does not have these attributes.

\sphinxlineitem{\sphinxstyleemphasis{Kspace\_modify eigtol must be smaller than one}}
\sphinxAtStartPar
Self\sphinxhyphen{}explanatory.

\sphinxlineitem{\sphinxstyleemphasis{LAMMPS is not built with Python embedded}}
\sphinxAtStartPar
This is done by including the PYTHON package before LAMMPS is built.
This is required to use python\sphinxhyphen{}style variables.

\sphinxlineitem{\sphinxstyleemphasis{LAMMPS unit\_style lj not supported by KIM models}}
\sphinxAtStartPar
Self\sphinxhyphen{}explanatory. Check the input script or data file.

\sphinxlineitem{\sphinxstyleemphasis{LJ6 off not supported in pair\_style buck/long/coul/long}}
\sphinxAtStartPar
Self\sphinxhyphen{}explanatory.

\sphinxlineitem{\sphinxstyleemphasis{Label map is incomplete: all types must be assigned a unique type label}}
\sphinxAtStartPar
For a given type\sphinxhyphen{}kind (atom types, bond types, etc.) to be written to
the data file, all associated types must be assigned a type label, and
each type label can be assigned to only one numeric type.

\sphinxlineitem{\sphinxstyleemphasis{Label wasn’t found in input script}}
\sphinxAtStartPar
Self\sphinxhyphen{}explanatory.

\sphinxlineitem{\sphinxstyleemphasis{Labelmap command before simulation box is defined}}
\sphinxAtStartPar
The labelmap command cannot be used before a read\_data,
read\_restart, or create\_box command.

\sphinxlineitem{\sphinxstyleemphasis{Lattice orient vectors are not orthogonal}}
\sphinxAtStartPar
The three specified lattice orientation vectors must be mutually
orthogonal.

\sphinxlineitem{\sphinxstyleemphasis{Lattice orient vectors are not right\sphinxhyphen{}handed}}
\sphinxAtStartPar
The three specified lattice orientation vectors must create a
right\sphinxhyphen{}handed coordinate system such that a1 cross a2 = a3.

\sphinxlineitem{\sphinxstyleemphasis{Lattice primitive vectors are collinear}}
\sphinxAtStartPar
The specified lattice primitive vectors do not for a unit cell with
non\sphinxhyphen{}zero volume.

\sphinxlineitem{\sphinxstyleemphasis{Lattice settings are not compatible with 2d simulation}}
\sphinxAtStartPar
One or more of the specified lattice vectors has a non\sphinxhyphen{}zero z
component.

\sphinxlineitem{\sphinxstyleemphasis{Lattice spacings are invalid}}
\sphinxAtStartPar
Each x,y,z spacing must be \textgreater{} 0.

\sphinxlineitem{\sphinxstyleemphasis{Lattice style incompatible with simulation dimension}}
\sphinxAtStartPar
2d simulation can use sq, sq2, or hex lattice.  3d simulation can use
sc, bcc, or fcc lattice.

\sphinxlineitem{\sphinxstyleemphasis{Log of zero/negative value in variable formula}}
\sphinxAtStartPar
Self\sphinxhyphen{}explanatory.

\sphinxlineitem{\sphinxstyleemphasis{Lost atoms via balance: original \%ld current \%ld}}
\sphinxAtStartPar
This should not occur.  Report the problem to the developers.

\sphinxlineitem{\sphinxstyleemphasis{Lost atoms: original \%ld current \%ld}}
\sphinxAtStartPar
Lost atoms are checked for each time thermo output is done.  See the
thermo\_modify lost command for options.  Lost atoms usually indicate
bad dynamics, e.g. atoms have been blown far out of the simulation
box, or moved further than one processor’s subdomain away before
reneighboring.

\sphinxlineitem{\sphinxstyleemphasis{MEAM library error \%d}}
\sphinxAtStartPar
A call to the MEAM Fortran library returned an error.

\sphinxlineitem{\sphinxstyleemphasis{MPI\_LMP\_BIGINT and bigint in lmptype.h are not compatible}}
\sphinxAtStartPar
The size of the MPI datatype does not match the size of a bigint.

\sphinxlineitem{\sphinxstyleemphasis{MPI\_LMP\_TAGINT and tagint in lmptype.h are not compatible}}
\sphinxAtStartPar
The size of the MPI datatype does not match the size of a tagint.

\sphinxlineitem{\sphinxstyleemphasis{MSM can only currently be used with comm\_style brick}}
\sphinxAtStartPar
This is a current restriction in LAMMPS.

\sphinxlineitem{\sphinxstyleemphasis{MSM grid is too large}}
\sphinxAtStartPar
The global MSM grid is larger than OFFSET in one or more dimensions.
OFFSET is currently set to 16384.  You likely need to decrease the
requested accuracy.

\sphinxlineitem{\sphinxstyleemphasis{MSM order must be 4, 6, 8, or 10}}
\sphinxAtStartPar
This is a limitation of the MSM implementation in LAMMPS:
the MSM order can only be 4, 6, 8, or 10.

\sphinxlineitem{\sphinxstyleemphasis{Mass command before simulation box is defined}}
\sphinxAtStartPar
The mass command cannot be used before a read\_data, read\_restart, or
create\_box command.

\sphinxlineitem{\sphinxstyleemphasis{Matrix factorization to split dispersion coefficients failed}}
\sphinxAtStartPar
This should not normally happen.  Contact the developers.

\sphinxlineitem{\sphinxstyleemphasis{Min\_style command before simulation box is defined}}
\sphinxAtStartPar
The min\_style command cannot be used before a read\_data, read\_restart,
or create\_box command.

\sphinxlineitem{\sphinxstyleemphasis{Minimization could not find thermo\_pe compute}}
\sphinxAtStartPar
This compute is created by the thermo command.  It must have been
explicitly deleted by a uncompute command.

\sphinxlineitem{\sphinxstyleemphasis{Minimize command before simulation box is defined}}
\sphinxAtStartPar
The minimize command cannot be used before a read\_data, read\_restart,
or create\_box command.

\sphinxlineitem{\sphinxstyleemphasis{Mismatched brackets in variable}}
\sphinxAtStartPar
Self\sphinxhyphen{}explanatory.

\sphinxlineitem{\sphinxstyleemphasis{Mismatched compute in variable formula}}
\sphinxAtStartPar
A compute is referenced incorrectly or a compute that produces per\sphinxhyphen{}atom
values is used in an equal\sphinxhyphen{}style variable formula.

\sphinxlineitem{\sphinxstyleemphasis{Mismatched fix in variable formula}}
\sphinxAtStartPar
A fix is referenced incorrectly or a fix that produces per\sphinxhyphen{}atom
values is used in an equal\sphinxhyphen{}style variable formula.

\sphinxlineitem{\sphinxstyleemphasis{Mismatched parameter in MEAM library file: z!=lat}}
\sphinxAtStartPar
The coordination number and lattice do not match, check that consistent values are given.

\sphinxlineitem{\sphinxstyleemphasis{Mismatched variable in variable formula}}
\sphinxAtStartPar
A variable is referenced incorrectly or an atom\sphinxhyphen{}style variable that
produces per\sphinxhyphen{}atom values is used in an equal\sphinxhyphen{}style variable
formula.

\sphinxlineitem{\sphinxstyleemphasis{Modulo 0 in variable formula}}
\sphinxAtStartPar
Self\sphinxhyphen{}explanatory.

\sphinxlineitem{\sphinxstyleemphasis{Molecule IDs too large for compute chunk/atom}}
\sphinxAtStartPar
The IDs must not be larger than can be stored in a 32\sphinxhyphen{}bit integer
since chunk IDs are 32\sphinxhyphen{}bit integers.

\sphinxlineitem{\sphinxstyleemphasis{Molecule auto special bond generation overflow}}
\sphinxAtStartPar
Counts exceed maxspecial setting for other atoms in system.

\sphinxlineitem{\sphinxstyleemphasis{Molecule file has angles but no nangles setting}}
\sphinxAtStartPar
Self\sphinxhyphen{}explanatory.

\sphinxlineitem{\sphinxstyleemphasis{Molecule file has body params but no setting for them}}
\sphinxAtStartPar
Self\sphinxhyphen{}explanatory.

\sphinxlineitem{\sphinxstyleemphasis{Molecule file has bonds but no nbonds setting}}
\sphinxAtStartPar
Self\sphinxhyphen{}explanatory.

\sphinxlineitem{\sphinxstyleemphasis{Molecule file has dihedrals but no ndihedrals setting}}
\sphinxAtStartPar
Self\sphinxhyphen{}explanatory.

\sphinxlineitem{\sphinxstyleemphasis{Molecule file has fragments but no nfragments setting}}
\sphinxAtStartPar
Self\sphinxhyphen{}explanatory.

\sphinxlineitem{\sphinxstyleemphasis{Molecule file has impropers but no nimpropers setting}}
\sphinxAtStartPar
Self\sphinxhyphen{}explanatory.

\sphinxlineitem{\sphinxstyleemphasis{Molecule file has no Body Doubles section}}
\sphinxAtStartPar
Self\sphinxhyphen{}explanatory.

\sphinxlineitem{\sphinxstyleemphasis{Molecule file has no Body Integers section}}
\sphinxAtStartPar
Self\sphinxhyphen{}explanatory.

\sphinxlineitem{\sphinxstyleemphasis{Molecule file has no Fragments section}}
\sphinxAtStartPar
Self\sphinxhyphen{}explanatory.

\sphinxlineitem{\sphinxstyleemphasis{Molecule file has special flags but no bonds}}
\sphinxAtStartPar
Self\sphinxhyphen{}explanatory.

\sphinxlineitem{\sphinxstyleemphasis{Molecule file needs both Special Bond sections}}
\sphinxAtStartPar
Self\sphinxhyphen{}explanatory.

\sphinxlineitem{\sphinxstyleemphasis{Molecule file requires atom style body}}
\sphinxAtStartPar
Self\sphinxhyphen{}explanatory.

\sphinxlineitem{\sphinxstyleemphasis{Molecule file shake flags not before shake atoms}}
\sphinxAtStartPar
The order of the two sections is important.

\sphinxlineitem{\sphinxstyleemphasis{Molecule file shake flags not before shake bonds}}
\sphinxAtStartPar
The order of the two sections is important.

\sphinxlineitem{\sphinxstyleemphasis{Molecule file shake info is incomplete}}
\sphinxAtStartPar
All 3 SHAKE sections are needed.

\sphinxlineitem{\sphinxstyleemphasis{Molecule file special list does not match special count}}
\sphinxAtStartPar
The number of values in an atom’s special list does not match count.

\sphinxlineitem{\sphinxstyleemphasis{Molecule file z center\sphinxhyphen{}of\sphinxhyphen{}mass must be 0.0 for 2d}}
\sphinxAtStartPar
Self\sphinxhyphen{}explanatory.

\sphinxlineitem{\sphinxstyleemphasis{Molecule file z coord must be 0.0 for 2d}}
\sphinxAtStartPar
Self\sphinxhyphen{}explanatory.

\sphinxlineitem{\sphinxstyleemphasis{Molecule natoms must be 1 for body particle}}
\sphinxAtStartPar
Self\sphinxhyphen{}explanatory.

\sphinxlineitem{\sphinxstyleemphasis{Molecule sizescale must be 1.0 for body particle}}
\sphinxAtStartPar
Self\sphinxhyphen{}explanatory.

\sphinxlineitem{\sphinxstyleemphasis{Molecule template ID for atom\_style template does not exist}}
\sphinxAtStartPar
Self\sphinxhyphen{}explanatory.

\sphinxlineitem{\sphinxstyleemphasis{Molecule template ID for create\_atoms does not exist}}
\sphinxAtStartPar
Self\sphinxhyphen{}explanatory.

\sphinxlineitem{\sphinxstyleemphasis{Molecule template ID for fix deposit does not exist}}
\sphinxAtStartPar
Self\sphinxhyphen{}explanatory.

\sphinxlineitem{\sphinxstyleemphasis{Molecule template ID for fix gcmc does not exist}}
\sphinxAtStartPar
Self\sphinxhyphen{}explanatory.

\sphinxlineitem{\sphinxstyleemphasis{Molecule template ID for fix pour does not exist}}
\sphinxAtStartPar
Self\sphinxhyphen{}explanatory.

\sphinxlineitem{\sphinxstyleemphasis{Molecule template ID for fix rigid/small does not exist}}
\sphinxAtStartPar
Self\sphinxhyphen{}explanatory.

\sphinxlineitem{\sphinxstyleemphasis{Molecule template ID for fix shake does not exist}}
\sphinxAtStartPar
Self\sphinxhyphen{}explanatory.

\sphinxlineitem{\sphinxstyleemphasis{Molecule template ID must be alphanumeric or underscore characters}}
\sphinxAtStartPar
Self\sphinxhyphen{}explanatory.

\sphinxlineitem{\sphinxstyleemphasis{Molecule topology/atom exceeds system topology/atom}}
\sphinxAtStartPar
The number of bonds, angles, etc per\sphinxhyphen{}atom in the molecule exceeds the
system setting.  See the create\_box command for how to specify these
values.

\sphinxlineitem{\sphinxstyleemphasis{Molecule topology type exceeds system topology type}}
\sphinxAtStartPar
The number of bond, angle, etc types in the molecule exceeds the
system setting.  See the create\_box command for how to specify these
values.

\sphinxlineitem{\sphinxstyleemphasis{More than one fix deform}}
\sphinxAtStartPar
Only one fix deform can be defined at a time.

\sphinxlineitem{\sphinxstyleemphasis{More than one fix freeze}}
\sphinxAtStartPar
Only one of these fixes can be defined, since the granular pair
potentials access it.

\sphinxlineitem{\sphinxstyleemphasis{More than one fix shake}}
\sphinxAtStartPar
Only one fix shake can be defined.

\sphinxlineitem{\sphinxstyleemphasis{Mu not allowed when not using semi\sphinxhyphen{}grand in fix atom/swap command}}
\sphinxAtStartPar
Self\sphinxhyphen{}explanatory.

\sphinxlineitem{\sphinxstyleemphasis{Must define angle\_style before Angle Coeffs}}
\sphinxAtStartPar
Must use an angle\_style command before reading a data file that
defines Angle Coeffs.

\sphinxlineitem{\sphinxstyleemphasis{Must define angle\_style before BondAngle Coeffs}}
\sphinxAtStartPar
Must use an angle\_style command before reading a data file that
defines Angle Coeffs.

\sphinxlineitem{\sphinxstyleemphasis{Must define angle\_style before BondBond Coeffs}}
\sphinxAtStartPar
Must use an angle\_style command before reading a data file that
defines Angle Coeffs.

\sphinxlineitem{\sphinxstyleemphasis{Must define bond\_style before Bond Coeffs}}
\sphinxAtStartPar
Must use a bond\_style command before reading a data file that
defines Bond Coeffs.

\sphinxlineitem{\sphinxstyleemphasis{Must define dihedral\_style before AngleAngleTorsion Coeffs}}
\sphinxAtStartPar
Must use a dihedral\_style command before reading a data file that
defines AngleAngleTorsion Coeffs.

\sphinxlineitem{\sphinxstyleemphasis{Must define dihedral\_style before AngleTorsion Coeffs}}
\sphinxAtStartPar
Must use a dihedral\_style command before reading a data file that
defines AngleTorsion Coeffs.

\sphinxlineitem{\sphinxstyleemphasis{Must define dihedral\_style before BondBond13 Coeffs}}
\sphinxAtStartPar
Must use a dihedral\_style command before reading a data file that
defines BondBond13 Coeffs.

\sphinxlineitem{\sphinxstyleemphasis{Must define dihedral\_style before Dihedral Coeffs}}
\sphinxAtStartPar
Must use a dihedral\_style command before reading a data file that
defines Dihedral Coeffs.

\sphinxlineitem{\sphinxstyleemphasis{Must define dihedral\_style before EndBondTorsion Coeffs}}
\sphinxAtStartPar
Must use a dihedral\_style command before reading a data file that
defines EndBondTorsion Coeffs.

\sphinxlineitem{\sphinxstyleemphasis{Must define dihedral\_style before MiddleBondTorsion Coeffs}}
\sphinxAtStartPar
Must use a dihedral\_style command before reading a data file that
defines MiddleBondTorsion Coeffs.

\sphinxlineitem{\sphinxstyleemphasis{Must define improper\_style before AngleAngle Coeffs}}
\sphinxAtStartPar
Must use an improper\_style command before reading a data file that
defines AngleAngle Coeffs.

\sphinxlineitem{\sphinxstyleemphasis{Must define improper\_style before Improper Coeffs}}
\sphinxAtStartPar
Must use an improper\_style command before reading a data file that
defines Improper Coeffs.

\sphinxlineitem{\sphinxstyleemphasis{Must define pair\_style before Pair Coeffs}}
\sphinxAtStartPar
Must use a pair\_style command before reading a data file that defines
Pair Coeffs.

\sphinxlineitem{\sphinxstyleemphasis{Must define pair\_style before PairIJ Coeffs}}
\sphinxAtStartPar
Must use a pair\_style command before reading a data file that defines
PairIJ Coeffs.

\sphinxlineitem{\sphinxstyleemphasis{Must have more than one processor partition to temper}}
\sphinxAtStartPar
Cannot use the temper command with only one processor partition.  Use
the \sphinxhyphen{}partition command\sphinxhyphen{}line option.

\sphinxlineitem{\sphinxstyleemphasis{Must not have multiple fixes change box parameter …}}
\sphinxAtStartPar
Self\sphinxhyphen{}explanatory.

\sphinxlineitem{\sphinxstyleemphasis{Must read Angle Type Labels before Angles}}
\sphinxAtStartPar
An Angle Type Labels section of a data file must come before the Angles section.

\sphinxlineitem{\sphinxstyleemphasis{Must read Atom Type Labels before Atoms}}
\sphinxAtStartPar
An Atom Type Labels section of a data file must come before the Atoms section.

\sphinxlineitem{\sphinxstyleemphasis{Must read Atoms before Angles}}
\sphinxAtStartPar
The Atoms section of a data file must come before an Angles section.

\sphinxlineitem{\sphinxstyleemphasis{Must read Atoms before Bodies}}
\sphinxAtStartPar
The Atoms section of a data file must come before a Bodies section.

\sphinxlineitem{\sphinxstyleemphasis{Must read Atoms before Bonds}}
\sphinxAtStartPar
The Atoms section of a data file must come before a Bonds section.

\sphinxlineitem{\sphinxstyleemphasis{Must read Atoms before Dihedrals}}
\sphinxAtStartPar
The Atoms section of a data file must come before a Dihedrals section.

\sphinxlineitem{\sphinxstyleemphasis{Must read Atoms before Ellipsoids}}
\sphinxAtStartPar
The Atoms section of a data file must come before a Ellipsoids
section.

\sphinxlineitem{\sphinxstyleemphasis{Must read Atoms before Impropers}}
\sphinxAtStartPar
The Atoms section of a data file must come before an Impropers
section.

\sphinxlineitem{\sphinxstyleemphasis{Must read Atoms before Lines}}
\sphinxAtStartPar
The Atoms section of a data file must come before a Lines section.

\sphinxlineitem{\sphinxstyleemphasis{Must read Atoms before Triangles}}
\sphinxAtStartPar
The Atoms section of a data file must come before a Triangles section.

\sphinxlineitem{\sphinxstyleemphasis{Must read Atoms before Velocities}}
\sphinxAtStartPar
The Atoms section of a data file must come before a Velocities
section.

\sphinxlineitem{\sphinxstyleemphasis{Must read Bond Type Labels before Bonds}}
\sphinxAtStartPar
A Bond Type Labels section of a data file must come before the Bonds section.

\sphinxlineitem{\sphinxstyleemphasis{Must read Dihedral Type Labels before Dihedrals}}
\sphinxAtStartPar
An Dihedral Type Labels section of a data file must come before the Dihedrals section.

\sphinxlineitem{\sphinxstyleemphasis{Must read Improper Type Labels before Impropers}}
\sphinxAtStartPar
An Improper Type Labels section of a data file must come before the Impropers section.

\sphinxlineitem{\sphinxstyleemphasis{Must re\sphinxhyphen{}specify non\sphinxhyphen{}restarted pair style (xxx) after read\_restart}}
\sphinxAtStartPar
For pair styles, that do not store their settings in a restart file,
it must be defined with a new ‘pair\_style’ command after read\_restart.

\sphinxlineitem{\sphinxstyleemphasis{Must set both respa inner and outer}}
\sphinxAtStartPar
Cannot use just the inner or outer option with respa without using the
other.

\sphinxlineitem{\sphinxstyleemphasis{Must set number of threads via package omp command}}
\sphinxAtStartPar
Because you are using the OPENMP package, set the number of threads
via its settings, not by the pair\_style snap nthreads setting.

\sphinxlineitem{\sphinxstyleemphasis{Must shrink\sphinxhyphen{}wrap piston boundary}}
\sphinxAtStartPar
The boundary style of the face where the piston is applied must be of
type s (shrink\sphinxhyphen{}wrapped).

\sphinxlineitem{\sphinxstyleemphasis{Must specify a region in fix deposit}}
\sphinxAtStartPar
The region keyword must be specified with this fix.

\sphinxlineitem{\sphinxstyleemphasis{Must specify a region in fix pour}}
\sphinxAtStartPar
Self\sphinxhyphen{}explanatory.

\sphinxlineitem{\sphinxstyleemphasis{Must specify at least 2 types in fix atom/swap command}}
\sphinxAtStartPar
Self\sphinxhyphen{}explanatory.

\sphinxlineitem{\sphinxstyleemphasis{Must use ‘kim\_style init’ command before simulation box is defined}}
\sphinxAtStartPar
Self\sphinxhyphen{}explanatory.

\sphinxlineitem{\sphinxstyleemphasis{Must use ‘kim\_style define’ command after simulation box is defined}}
\sphinxAtStartPar
Self\sphinxhyphen{}explanatory.

\sphinxlineitem{\sphinxstyleemphasis{Must use ‘kim\_style init’ command before ‘kim\_style define’}}
\sphinxAtStartPar
Self\sphinxhyphen{}explanatory.

\sphinxlineitem{\sphinxstyleemphasis{Must use ‘kspace\_modify pressure/scalar no’ for rRESPA with kspace\_style MSM}}
\sphinxAtStartPar
The kspace scalar pressure option cannot (yet) be used with rRESPA.

\sphinxlineitem{\sphinxstyleemphasis{Must use ‘kspace\_modify pressure/scalar no’ for tensor components with kspace\_style msm}}
\sphinxAtStartPar
Otherwise MSM will compute only a scalar pressure.  See the kspace\_modify
command for details on this setting.

\sphinxlineitem{\sphinxstyleemphasis{Must use ‘kspace\_modify pressure/scalar no’ to obtain per\sphinxhyphen{}atom virial with kspace\_style MSM}}
\sphinxAtStartPar
The kspace scalar pressure option cannot be used to obtain per\sphinxhyphen{}atom virial.

\sphinxlineitem{\sphinxstyleemphasis{Must use ‘kspace\_modify pressure/scalar no’ with GPU MSM Pair styles}}
\sphinxAtStartPar
The kspace scalar pressure option is not (yet) compatible with GPU MSM Pair styles.

\sphinxlineitem{\sphinxstyleemphasis{Must use ‘kspace\_modify pressure/scalar no’ with kspace\_style msm/cg}}
\sphinxAtStartPar
The kspace scalar pressure option is not compatible with kspace\_style msm/cg.

\sphinxlineitem{\sphinxstyleemphasis{Must use \sphinxhyphen{}in switch with multiple partitions}}
\sphinxAtStartPar
A multi\sphinxhyphen{}partition simulation cannot read the input script from stdin.
The \sphinxhyphen{}in command\sphinxhyphen{}line option must be used to specify a file.

\sphinxlineitem{\sphinxstyleemphasis{Must use Kokkos half/thread or full neighbor list with threads or GPUs}}
\sphinxAtStartPar
Using Kokkos half\sphinxhyphen{}neighbor lists with threading is not allowed.

\sphinxlineitem{\sphinxstyleemphasis{Must use a block or cylinder region with fix pour}}
\sphinxAtStartPar
Self\sphinxhyphen{}explanatory.

\sphinxlineitem{\sphinxstyleemphasis{Must use a block region with fix pour for 2d simulations}}
\sphinxAtStartPar
Self\sphinxhyphen{}explanatory.

\sphinxlineitem{\sphinxstyleemphasis{Must use a bond style with TIP4P potential}}
\sphinxAtStartPar
TIP4P potentials assume bond lengths in water are constrained
by a fix shake command.

\sphinxlineitem{\sphinxstyleemphasis{Must use a molecular atom style with fix poems molecule}}
\sphinxAtStartPar
Self\sphinxhyphen{}explanatory.

\sphinxlineitem{\sphinxstyleemphasis{Must use a z\sphinxhyphen{}axis cylinder region with fix pour}}
\sphinxAtStartPar
Self\sphinxhyphen{}explanatory.

\sphinxlineitem{\sphinxstyleemphasis{Must use an angle style with TIP4P potential}}
\sphinxAtStartPar
TIP4P potentials assume angles in water are constrained by a fix shake
command.

\sphinxlineitem{\sphinxstyleemphasis{Must use atom map style array with Kokkos}}
\sphinxAtStartPar
See the atom\_modify map command.

\sphinxlineitem{\sphinxstyleemphasis{Must use atom style with molecule IDs with fix bond/swap}}
\sphinxAtStartPar
Self\sphinxhyphen{}explanatory.

\sphinxlineitem{\sphinxstyleemphasis{Must use pair\_style comb or comb3 with fix qeq/comb}}
\sphinxAtStartPar
Self\sphinxhyphen{}explanatory.

\sphinxlineitem{\sphinxstyleemphasis{Must use variable energy with fix addforce}}
\sphinxAtStartPar
Must define an energy variable when applying a dynamic
force during minimization.

\sphinxlineitem{\sphinxstyleemphasis{Must use variable energy with fix efield}}
\sphinxAtStartPar
You must define an energy when performing a minimization with a
variable E\sphinxhyphen{}field.

\sphinxlineitem{\sphinxstyleemphasis{NEB command before simulation box is defined}}
\sphinxAtStartPar
Self\sphinxhyphen{}explanatory.

\sphinxlineitem{\sphinxstyleemphasis{NEB requires damped dynamics minimizer}}
\sphinxAtStartPar
Use a different minimization style.

\sphinxlineitem{\sphinxstyleemphasis{NEB requires use of fix neb}}
\sphinxAtStartPar
Self\sphinxhyphen{}explanatory.

\sphinxlineitem{\sphinxstyleemphasis{NL ramp in wall/piston only implemented in zlo for now}}
\sphinxAtStartPar
The ramp keyword can only be used for piston applied to face zlo.

\sphinxlineitem{\sphinxstyleemphasis{Need nswaptypes mu values in fix atom/swap command}}
\sphinxAtStartPar
Self\sphinxhyphen{}explanatory.

\sphinxlineitem{\sphinxstyleemphasis{Needed bonus data not in data file}}
\sphinxAtStartPar
Some atom styles require bonus data.  See the read\_data page for
details.

\sphinxlineitem{\sphinxstyleemphasis{Needed molecular topology not in data file}}
\sphinxAtStartPar
The header of the data file indicated bonds, angles, etc would be
included, but they are not present.

\sphinxlineitem{\sphinxstyleemphasis{Neigh\_modify exclude molecule requires atom attribute molecule}}
\sphinxAtStartPar
Self\sphinxhyphen{}explanatory.

\sphinxlineitem{\sphinxstyleemphasis{Neigh\_modify include group != atom\_modify first group}}
\sphinxAtStartPar
Self\sphinxhyphen{}explanatory.

\sphinxlineitem{\sphinxstyleemphasis{Neighbor delay must be 0 or multiple of every setting}}
\sphinxAtStartPar
The delay and every parameters set via the neigh\_modify command are
inconsistent.  If the delay setting is non\sphinxhyphen{}zero, then it must be a
multiple of the every setting.

\sphinxlineitem{\sphinxstyleemphasis{Neighbor include group not allowed with ghost neighbors}}
\sphinxAtStartPar
This is a current restriction within LAMMPS.

\sphinxlineitem{\sphinxstyleemphasis{Neighbor list overflow, boost neigh\_modify one}}
\sphinxAtStartPar
There are too many neighbors of a single atom.  Use the neigh\_modify
command to increase the max number of neighbors allowed for one atom.
You may also want to boost the page size.

\sphinxlineitem{\sphinxstyleemphasis{Neighbor multi not yet enabled for ghost neighbors}}
\sphinxAtStartPar
This is a current restriction within LAMMPS.

\sphinxlineitem{\sphinxstyleemphasis{Neighbor multi not yet enabled for granular}}
\sphinxAtStartPar
Self\sphinxhyphen{}explanatory.

\sphinxlineitem{\sphinxstyleemphasis{Neighbor multi not yet enabled for rRESPA}}
\sphinxAtStartPar
Self\sphinxhyphen{}explanatory.

\sphinxlineitem{\sphinxstyleemphasis{Neighbor page size must be \textgreater{}= 10x the one atom setting}}
\sphinxAtStartPar
This is required to prevent wasting too much memory.

\sphinxlineitem{\sphinxstyleemphasis{New atom IDs exceed maximum allowed ID}}
\sphinxAtStartPar
See the setting for tagint in the src/lmptype.h file.

\sphinxlineitem{\sphinxstyleemphasis{New bond exceeded bonds per atom in create\_bonds}}
\sphinxAtStartPar
See the read\_data command for info on using the “extra/bond/per/atom”
keyword to allow for additional bonds to be formed

\sphinxlineitem{\sphinxstyleemphasis{New bond exceeded bonds per atom in fix bond/create}}
\sphinxAtStartPar
See the read\_data command for info on using the “extra/bond/per/atom”
keyword to allow for additional bonds to be formed

\sphinxlineitem{\sphinxstyleemphasis{New bond exceeded special list size in fix bond/create}}
\sphinxAtStartPar
See the “read\_data extra/special/per/atom” command
(or the “create\_box extra/special/per/atom” command)
for info on how to leave space in the special bonds
list to allow for additional bonds to be formed.

\sphinxlineitem{\sphinxstyleemphasis{Newton bond change after simulation box is defined}}
\sphinxAtStartPar
The newton command cannot be used to change the newton bond value
after a read\_data, read\_restart, or create\_box command.

\sphinxlineitem{\sphinxstyleemphasis{Next command must list all universe and uloop variables}}
\sphinxAtStartPar
This is to ensure they stay in sync.

\sphinxlineitem{\sphinxstyleemphasis{No Kspace style defined for compute group/group}}
\sphinxAtStartPar
Self\sphinxhyphen{}explanatory.

\sphinxlineitem{\sphinxstyleemphasis{No OpenMP support compiled in}}
\sphinxAtStartPar
An OpenMP flag is set, but LAMMPS was not built with
OpenMP support.

\sphinxlineitem{\sphinxstyleemphasis{No angle style is defined for compute angle/local}}
\sphinxAtStartPar
Self\sphinxhyphen{}explanatory.

\sphinxlineitem{\sphinxstyleemphasis{No angles allowed with this atom style}}
\sphinxAtStartPar
Self\sphinxhyphen{}explanatory.

\sphinxlineitem{\sphinxstyleemphasis{No atoms in data file}}
\sphinxAtStartPar
The header of the data file indicated that atoms would be included,
but they are not present.

\sphinxlineitem{\sphinxstyleemphasis{No basis atoms in lattice}}
\sphinxAtStartPar
Basis atoms must be defined for lattice style user.

\sphinxlineitem{\sphinxstyleemphasis{No bodies allowed with this atom style}}
\sphinxAtStartPar
Self\sphinxhyphen{}explanatory.  Check data file.

\sphinxlineitem{\sphinxstyleemphasis{No bond style is defined for compute bond/local}}
\sphinxAtStartPar
Self\sphinxhyphen{}explanatory.

\sphinxlineitem{\sphinxstyleemphasis{No bonds allowed with this atom style}}
\sphinxAtStartPar
Self\sphinxhyphen{}explanatory.

\sphinxlineitem{\sphinxstyleemphasis{No box information in dump. You have to use ‘box no’}}
\sphinxAtStartPar
Self\sphinxhyphen{}explanatory.

\sphinxlineitem{\sphinxstyleemphasis{No count or invalid atom count in molecule file}}
\sphinxAtStartPar
The number of atoms must be specified.

\sphinxlineitem{\sphinxstyleemphasis{No dihedral style is defined for compute dihedral/local}}
\sphinxAtStartPar
Self\sphinxhyphen{}explanatory.

\sphinxlineitem{\sphinxstyleemphasis{No dihedrals allowed with this atom style}}
\sphinxAtStartPar
Self\sphinxhyphen{}explanatory.

\sphinxlineitem{\sphinxstyleemphasis{No dump custom arguments specified}}
\sphinxAtStartPar
The dump custom command requires that atom quantities be specified to
output to dump file.

\sphinxlineitem{\sphinxstyleemphasis{No dump local arguments specified}}
\sphinxAtStartPar
Self\sphinxhyphen{}explanatory.

\sphinxlineitem{\sphinxstyleemphasis{No ellipsoids allowed with this atom style}}
\sphinxAtStartPar
Self\sphinxhyphen{}explanatory.  Check data file.

\sphinxlineitem{\sphinxstyleemphasis{No fix gravity defined for fix pour}}
\sphinxAtStartPar
Gravity is required to use fix pour.

\sphinxlineitem{\sphinxstyleemphasis{No improper style is defined for compute improper/local}}
\sphinxAtStartPar
Self\sphinxhyphen{}explanatory.

\sphinxlineitem{\sphinxstyleemphasis{No impropers allowed with this atom style}}
\sphinxAtStartPar
Self\sphinxhyphen{}explanatory.

\sphinxlineitem{\sphinxstyleemphasis{No input values for fix ave/spatial}}
\sphinxAtStartPar
Self\sphinxhyphen{}explanatory.

\sphinxlineitem{\sphinxstyleemphasis{No lines allowed with this atom style}}
\sphinxAtStartPar
Self\sphinxhyphen{}explanatory.  Check data file.

\sphinxlineitem{\sphinxstyleemphasis{No matching element in ADP potential file}}
\sphinxAtStartPar
The ADP potential file does not contain elements that match the
requested elements.

\sphinxlineitem{\sphinxstyleemphasis{No matching element in EAM potential file}}
\sphinxAtStartPar
The EAM potential file does not contain elements that match the
requested elements.

\sphinxlineitem{\sphinxstyleemphasis{No molecule topology allowed with atom style template}}
\sphinxAtStartPar
The data file cannot specify the number of bonds, angles, etc,
because this info if inferred from the molecule templates.

\sphinxlineitem{\sphinxstyleemphasis{No overlap of box and region for create\_atoms}}
\sphinxAtStartPar
Self\sphinxhyphen{}explanatory.

\sphinxlineitem{\sphinxstyleemphasis{No pair coul/streitz for fix qeq/slater}}
\sphinxAtStartPar
These commands must be used together.

\sphinxlineitem{\sphinxstyleemphasis{No pair hbond/dreiding coefficients set}}
\sphinxAtStartPar
Self\sphinxhyphen{}explanatory.

\sphinxlineitem{\sphinxstyleemphasis{No pair style defined for compute group/group}}
\sphinxAtStartPar
Cannot calculate group interactions without a pair style defined.

\sphinxlineitem{\sphinxstyleemphasis{No pair style is defined for compute pair/local}}
\sphinxAtStartPar
Self\sphinxhyphen{}explanatory.

\sphinxlineitem{\sphinxstyleemphasis{No pair style is defined for compute property/local}}
\sphinxAtStartPar
Self\sphinxhyphen{}explanatory.

\sphinxlineitem{\sphinxstyleemphasis{No rigid bodies defined}}
\sphinxAtStartPar
The fix specification did not end up defining any rigid bodies.

\sphinxlineitem{\sphinxstyleemphasis{No triangles allowed with this atom style}}
\sphinxAtStartPar
Self\sphinxhyphen{}explanatory.  Check data file.

\sphinxlineitem{\sphinxstyleemphasis{No values in fix ave/chunk command}}
\sphinxAtStartPar
Self\sphinxhyphen{}explanatory.

\sphinxlineitem{\sphinxstyleemphasis{No values in fix ave/time command}}
\sphinxAtStartPar
Self\sphinxhyphen{}explanatory.

\sphinxlineitem{\sphinxstyleemphasis{Non digit character between brackets in variable}}
\sphinxAtStartPar
Self\sphinxhyphen{}explanatory.

\sphinxlineitem{\sphinxstyleemphasis{Non integer \# of swaps in temper command}}
\sphinxAtStartPar
Swap frequency in temper command must evenly divide the total \# of
timesteps.

\sphinxlineitem{\sphinxstyleemphasis{Non\sphinxhyphen{}numeric box dimensions \sphinxhyphen{} simulation unstable}}
\sphinxAtStartPar
The box size has apparently blown up.

\sphinxlineitem{\sphinxstyleemphasis{Non\sphinxhyphen{}zero atom IDs with atom\_modify id = no}}
\sphinxAtStartPar
Self\sphinxhyphen{}explanatory.

\sphinxlineitem{\sphinxstyleemphasis{Non\sphinxhyphen{}zero read\_data shift z value for 2d simulation}}
\sphinxAtStartPar
Self\sphinxhyphen{}explanatory.

\sphinxlineitem{\sphinxstyleemphasis{Nprocs not a multiple of N for \sphinxhyphen{}reorder}}
\sphinxAtStartPar
Self\sphinxhyphen{}explanatory.

\sphinxlineitem{\sphinxstyleemphasis{Number of core atoms != number of shell atoms}}
\sphinxAtStartPar
There must be a one\sphinxhyphen{}to\sphinxhyphen{}one pairing of core and shell atoms.

\sphinxlineitem{\sphinxstyleemphasis{Numeric index is out of bounds}}
\sphinxAtStartPar
A command with an argument that specifies an integer or range of
integers is using a value that is less than 1 or greater than the
maximum allowed limit.

\sphinxlineitem{\sphinxstyleemphasis{One or more Atom IDs is negative}}
\sphinxAtStartPar
Atom IDs must be positive integers.

\sphinxlineitem{\sphinxstyleemphasis{One or more atom IDs is too big}}
\sphinxAtStartPar
The limit on atom IDs is set by the SMALLBIG, BIGBIG, SMALLSMALL
setting in your LAMMPS build.  See the {\hyperref[\detokenize{Build_settings::doc}]{\sphinxcrossref{\DUrole{doc}{Build settings}}}} page for more info.

\sphinxlineitem{\sphinxstyleemphasis{One or more atom IDs is zero}}
\sphinxAtStartPar
Either all atoms IDs must be zero or none of them.

\sphinxlineitem{\sphinxstyleemphasis{One or more atoms belong to multiple rigid bodies}}
\sphinxAtStartPar
Two or more rigid bodies defined by the fix rigid command cannot
contain the same atom.

\sphinxlineitem{\sphinxstyleemphasis{One or more rigid bodies are a single particle}}
\sphinxAtStartPar
Self\sphinxhyphen{}explanatory.

\sphinxlineitem{\sphinxstyleemphasis{One or zero atoms in rigid body}}
\sphinxAtStartPar
Any rigid body defined by the fix rigid command must contain 2 or more
atoms.

\sphinxlineitem{\sphinxstyleemphasis{Only 2 types allowed when not using semi\sphinxhyphen{}grand in fix atom/swap command}}
\sphinxAtStartPar
Self\sphinxhyphen{}explanatory.

\sphinxlineitem{\sphinxstyleemphasis{Only one cut\sphinxhyphen{}off allowed when requesting all long}}
\sphinxAtStartPar
Self\sphinxhyphen{}explanatory.

\sphinxlineitem{\sphinxstyleemphasis{Only one cutoff allowed when requesting all long}}
\sphinxAtStartPar
Self\sphinxhyphen{}explanatory.

\sphinxlineitem{\sphinxstyleemphasis{Only zhi currently implemented for fix append/atoms}}
\sphinxAtStartPar
Self\sphinxhyphen{}explanatory.

\sphinxlineitem{\sphinxstyleemphasis{Out of range atoms \sphinxhyphen{} cannot compute MSM}}
\sphinxAtStartPar
One or more atoms are attempting to map their charge to a MSM grid point
that is not owned by a processor.  This is likely for one of two
reasons, both of them bad.  First, it may mean that an atom near the
boundary of a processor’s subdomain has moved more than 1/2 the
{\hyperref[\detokenize{neighbor::doc}]{\sphinxcrossref{\DUrole{doc}{neighbor skin distance}}}} without neighbor lists being
rebuilt and atoms being migrated to new processors.  This also means
you may be missing pairwise interactions that need to be computed.
The solution is to change the re\sphinxhyphen{}neighboring criteria via the
{\hyperref[\detokenize{neigh_modify::doc}]{\sphinxcrossref{\DUrole{doc}{neigh\_modify}}}} command.  The safest settings are
“delay 0 every 1 check yes”.  Second, it may mean that an atom has
moved far outside a processor’s subdomain or even the entire
simulation box. This indicates bad physics, e.g. due to highly
overlapping atoms, too large a timestep, etc.

\sphinxlineitem{\sphinxstyleemphasis{Out of range atoms \sphinxhyphen{} cannot compute PPPM}}
\sphinxAtStartPar
One or more atoms are attempting to map their charge to a PPPM grid
point that is not owned by a processor.  This is likely for one of two
reasons, both of them bad.  First, it may mean that an atom near the
boundary of a processor’s subdomain has moved more than 1/2 the
{\hyperref[\detokenize{neighbor::doc}]{\sphinxcrossref{\DUrole{doc}{neighbor skin distance}}}} without neighbor lists being
rebuilt and atoms being migrated to new processors.  This also means
you may be missing pairwise interactions that need to be computed.
The solution is to change the re\sphinxhyphen{}neighboring criteria via the
{\hyperref[\detokenize{neigh_modify::doc}]{\sphinxcrossref{\DUrole{doc}{neigh\_modify}}}} command.  The safest settings are
“delay 0 every 1 check yes”.  Second, it may mean that an atom has
moved far outside a processor’s subdomain or even the entire
simulation box. This indicates bad physics, e.g. due to highly
overlapping atoms, too large a timestep, etc.

\sphinxlineitem{\sphinxstyleemphasis{Out of range atoms \sphinxhyphen{} cannot compute PPPMDisp}}
\sphinxAtStartPar
One or more atoms are attempting to map their charge to a PPPM grid
point that is not owned by a processor.  This is likely for one of two
reasons, both of them bad.  First, it may mean that an atom near the
boundary of a processor’s subdomain has moved more than 1/2 the
{\hyperref[\detokenize{neighbor::doc}]{\sphinxcrossref{\DUrole{doc}{neighbor skin distance}}}} without neighbor lists being
rebuilt and atoms being migrated to new processors.  This also means
you may be missing pairwise interactions that need to be computed.
The solution is to change the re\sphinxhyphen{}neighboring criteria via the
{\hyperref[\detokenize{neigh_modify::doc}]{\sphinxcrossref{\DUrole{doc}{neigh\_modify}}}} command.  The safest settings are
“delay 0 every 1 check yes”.  Second, it may mean that an atom has
moved far outside a processor’s subdomain or even the entire
simulation box. This indicates bad physics, e.g. due to highly
overlapping atoms, too large a timestep, etc.

\sphinxlineitem{\sphinxstyleemphasis{Overflow of allocated fix vector storage}}
\sphinxAtStartPar
This should not normally happen if the fix correctly calculated
how long the vector will grow to.  Contact the developers.

\sphinxlineitem{\sphinxstyleemphasis{Overlapping large/large in pair colloid}}
\sphinxAtStartPar
This potential is infinite when there is an overlap.

\sphinxlineitem{\sphinxstyleemphasis{Overlapping small/large in pair colloid}}
\sphinxAtStartPar
This potential is infinite when there is an overlap.

\sphinxlineitem{\sphinxstyleemphasis{POEMS fix must come before NPT/NPH fix}}
\sphinxAtStartPar
NPT/NPH fix must be defined in input script after all poems fixes,
else the fix contribution to the pressure virial is incorrect.

\sphinxlineitem{\sphinxstyleemphasis{PPPM can only currently be used with comm\_style brick}}
\sphinxAtStartPar
This is a current restriction in LAMMPS.

\sphinxlineitem{\sphinxstyleemphasis{PPPM grid is too large}}
\sphinxAtStartPar
The global PPPM grid is larger than OFFSET in one or more dimensions.
OFFSET is currently set to 4096.  You likely need to decrease the
requested accuracy.

\sphinxlineitem{\sphinxstyleemphasis{PPPM grid stencil extends beyond nearest neighbor processor}}
\sphinxAtStartPar
This is not allowed if the kspace\_modify overlap setting is no.

\sphinxlineitem{\sphinxstyleemphasis{PPPM order \textless{} minimum allowed order}}
\sphinxAtStartPar
The default minimum order is 2.  This can be reset by the
kspace\_modify minorder command.

\sphinxlineitem{\sphinxstyleemphasis{PPPM order cannot be \textless{} 2 or \textgreater{} than \%d}}
\sphinxAtStartPar
This is a limitation of the PPPM implementation in LAMMPS.

\sphinxlineitem{\sphinxstyleemphasis{PPPMDisp Coulomb grid is too large}}
\sphinxAtStartPar
The global PPPM grid is larger than OFFSET in one or more dimensions.
OFFSET is currently set to 4096.  You likely need to decrease the
requested accuracy.

\sphinxlineitem{\sphinxstyleemphasis{PPPMDisp Dispersion grid is too large}}
\sphinxAtStartPar
The global PPPM grid is larger than OFFSET in one or more dimensions.
OFFSET is currently set to 4096.  You likely need to decrease the
requested accuracy.

\sphinxlineitem{\sphinxstyleemphasis{PPPMDisp can only currently be used with comm\_style brick}}
\sphinxAtStartPar
This is a current restriction in LAMMPS.

\sphinxlineitem{\sphinxstyleemphasis{PPPMDisp coulomb order cannot be greater than \%d}}
\sphinxAtStartPar
This is a limitation of the PPPM implementation in LAMMPS.

\sphinxlineitem{\sphinxstyleemphasis{PPPMDisp used but no parameters set, for further information please see the pppm/disp documentation}}
\sphinxAtStartPar
An efficient and accurate usage of the pppm/disp requires settings via the kspace\_modify command. Please see the pppm/disp documentation for further instructions.

\sphinxlineitem{\sphinxstyleemphasis{PRD command before simulation box is defined}}
\sphinxAtStartPar
The prd command cannot be used before a read\_data,
read\_restart, or create\_box command.

\sphinxlineitem{\sphinxstyleemphasis{PRD nsteps must be multiple of t\_event}}
\sphinxAtStartPar
Self\sphinxhyphen{}explanatory.

\sphinxlineitem{\sphinxstyleemphasis{PRD t\_corr must be multiple of t\_event}}
\sphinxAtStartPar
Self\sphinxhyphen{}explanatory.

\sphinxlineitem{\sphinxstyleemphasis{Package command after simulation box is defined}}
\sphinxAtStartPar
The package command cannot be used after a read\_data, read\_restart, or
create\_box command.

\sphinxlineitem{\sphinxstyleemphasis{Package gpu command without GPU package installed}}
\sphinxAtStartPar
The GPU package must be installed via “make yes\sphinxhyphen{}gpu” before LAMMPS is
built.

\sphinxlineitem{\sphinxstyleemphasis{Package intel command without INTEL package installed}}
\sphinxAtStartPar
The INTEL package must be installed via “make yes\sphinxhyphen{}intel”
before LAMMPS is built.

\sphinxlineitem{\sphinxstyleemphasis{Package kokkos command without KOKKOS package enabled}}
\sphinxAtStartPar
The KOKKOS package must be installed via “make yes\sphinxhyphen{}kokkos” before
LAMMPS is built, and the “\sphinxhyphen{}k on” must be used to enable the package.

\sphinxlineitem{\sphinxstyleemphasis{Package omp command without OPENMP package installed}}
\sphinxAtStartPar
The OPENMP package must be installed via “make yes\sphinxhyphen{}openmp” before
LAMMPS is built.

\sphinxlineitem{\sphinxstyleemphasis{Pair body requires atom style body}}
\sphinxAtStartPar
Self\sphinxhyphen{}explanatory.

\sphinxlineitem{\sphinxstyleemphasis{Pair body requires body style nparticle}}
\sphinxAtStartPar
This pair style is specific to the nparticle body style.

\sphinxlineitem{\sphinxstyleemphasis{Pair brownian requires atom style sphere}}
\sphinxAtStartPar
Self\sphinxhyphen{}explanatory.

\sphinxlineitem{\sphinxstyleemphasis{Pair brownian requires extended particles}}
\sphinxAtStartPar
One of the particles has radius 0.0.

\sphinxlineitem{\sphinxstyleemphasis{Pair brownian requires monodisperse particles}}
\sphinxAtStartPar
All particles must be the same finite size.

\sphinxlineitem{\sphinxstyleemphasis{Pair brownian/poly requires atom style sphere}}
\sphinxAtStartPar
Self\sphinxhyphen{}explanatory.

\sphinxlineitem{\sphinxstyleemphasis{Pair brownian/poly requires extended particles}}
\sphinxAtStartPar
One of the particles has radius 0.0.

\sphinxlineitem{\sphinxstyleemphasis{Pair brownian/poly requires newton pair off}}
\sphinxAtStartPar
Self\sphinxhyphen{}explanatory.

\sphinxlineitem{\sphinxstyleemphasis{Pair coeff for hybrid has invalid style}}
\sphinxAtStartPar
Style in pair coeff must have been listed in pair\_style command.

\sphinxlineitem{\sphinxstyleemphasis{Pair coul/wolf requires atom attribute q}}
\sphinxAtStartPar
The atom style defined does not have this attribute.

\sphinxlineitem{\sphinxstyleemphasis{Pair cutoff \textless{} Respa interior cutoff}}
\sphinxAtStartPar
One or more pairwise cutoffs are too short to use with the specified
rRESPA cutoffs.

\sphinxlineitem{\sphinxstyleemphasis{Pair dipole/cut requires atom attributes q, mu, torque}}
\sphinxAtStartPar
The atom style defined does not have these attributes.

\sphinxlineitem{\sphinxstyleemphasis{Pair dipole/cut/gpu requires atom attributes q, mu, torque}}
\sphinxAtStartPar
The atom style defined does not have this attribute.

\sphinxlineitem{\sphinxstyleemphasis{Pair dipole/long requires atom attributes q, mu, torque}}
\sphinxAtStartPar
The atom style defined does not have these attributes.

\sphinxlineitem{\sphinxstyleemphasis{Pair dipole/sf/gpu requires atom attributes q, mu, torque}}
\sphinxAtStartPar
The atom style defined does not one or more of these attributes.

\sphinxlineitem{\sphinxstyleemphasis{Pair distance \textless{} table inner cutoff}}
\sphinxAtStartPar
Two atoms are closer together than the pairwise table allows.

\sphinxlineitem{\sphinxstyleemphasis{Pair distance \textgreater{} table outer cutoff}}
\sphinxAtStartPar
Two atoms are further apart than the pairwise table allows.

\sphinxlineitem{\sphinxstyleemphasis{Pair dpd requires ghost atoms store velocity}}
\sphinxAtStartPar
Use the comm\_modify vel yes command to enable this.

\sphinxlineitem{\sphinxstyleemphasis{Pair gayberne epsilon a,b,c coeffs are not all set}}
\sphinxAtStartPar
Each atom type involved in pair\_style gayberne must
have these 3 coefficients set at least once.

\sphinxlineitem{\sphinxstyleemphasis{Pair gayberne requires atom style ellipsoid}}
\sphinxAtStartPar
Self\sphinxhyphen{}explanatory.

\sphinxlineitem{\sphinxstyleemphasis{Pair gayberne requires atoms with same type have same shape}}
\sphinxAtStartPar
Self\sphinxhyphen{}explanatory.

\sphinxlineitem{\sphinxstyleemphasis{Pair gayberne/gpu requires atom style ellipsoid}}
\sphinxAtStartPar
Self\sphinxhyphen{}explanatory.

\sphinxlineitem{\sphinxstyleemphasis{Pair gayberne/gpu requires atoms with same type have same shape}}
\sphinxAtStartPar
Self\sphinxhyphen{}explanatory.

\sphinxlineitem{\sphinxstyleemphasis{Pair granular requires atom attributes radius, rmass}}
\sphinxAtStartPar
The atom style defined does not have these attributes.

\sphinxlineitem{\sphinxstyleemphasis{Pair granular requires ghost atoms store velocity}}
\sphinxAtStartPar
Use the comm\_modify vel yes command to enable this.

\sphinxlineitem{\sphinxstyleemphasis{Pair granular with shear history requires newton pair off}}
\sphinxAtStartPar
This is a current restriction of the implementation of pair
granular styles with history.

\sphinxlineitem{\sphinxstyleemphasis{Pair hybrid single calls do not support per sub\sphinxhyphen{}style special bond values}}
\sphinxAtStartPar
Self\sphinxhyphen{}explanatory.

\sphinxlineitem{\sphinxstyleemphasis{Pair hybrid sub\sphinxhyphen{}style does not support single call}}
\sphinxAtStartPar
You are attempting to invoke a single() call on a pair style
that does not support it.

\sphinxlineitem{\sphinxstyleemphasis{Pair hybrid sub\sphinxhyphen{}style is not used}}
\sphinxAtStartPar
No pair\_coeff command used a sub\sphinxhyphen{}style specified in the pair\_style
command.

\sphinxlineitem{\sphinxstyleemphasis{Pair inner cutoff \textless{} Respa interior cutoff}}
\sphinxAtStartPar
One or more pairwise cutoffs are too short to use with the specified
rRESPA cutoffs.

\sphinxlineitem{\sphinxstyleemphasis{Pair inner cutoff \textgreater{}= Pair outer cutoff}}
\sphinxAtStartPar
The specified cutoffs for the pair style are inconsistent.

\sphinxlineitem{\sphinxstyleemphasis{Pair line/lj requires atom style line}}
\sphinxAtStartPar
Self\sphinxhyphen{}explanatory.

\sphinxlineitem{\sphinxstyleemphasis{Pair lj/long/dipole/long requires atom attributes mu, torque}}
\sphinxAtStartPar
The atom style defined does not have these attributes.

\sphinxlineitem{\sphinxstyleemphasis{Pair lubricate requires atom style sphere}}
\sphinxAtStartPar
Self\sphinxhyphen{}explanatory.

\sphinxlineitem{\sphinxstyleemphasis{Pair lubricate requires ghost atoms store velocity}}
\sphinxAtStartPar
Use the comm\_modify vel yes command to enable this.

\sphinxlineitem{\sphinxstyleemphasis{Pair lubricate requires monodisperse particles}}
\sphinxAtStartPar
All particles must be the same finite size.

\sphinxlineitem{\sphinxstyleemphasis{Pair lubricate/poly requires atom style sphere}}
\sphinxAtStartPar
Self\sphinxhyphen{}explanatory.

\sphinxlineitem{\sphinxstyleemphasis{Pair lubricate/poly requires extended particles}}
\sphinxAtStartPar
One of the particles has radius 0.0.

\sphinxlineitem{\sphinxstyleemphasis{Pair lubricate/poly requires ghost atoms store velocity}}
\sphinxAtStartPar
Use the comm\_modify vel yes command to enable this.

\sphinxlineitem{\sphinxstyleemphasis{Pair lubricate/poly requires newton pair off}}
\sphinxAtStartPar
Self\sphinxhyphen{}explanatory.

\sphinxlineitem{\sphinxstyleemphasis{Pair lubricateU requires atom style sphere}}
\sphinxAtStartPar
Self\sphinxhyphen{}explanatory.

\sphinxlineitem{\sphinxstyleemphasis{Pair lubricateU requires ghost atoms store velocity}}
\sphinxAtStartPar
Use the comm\_modify vel yes command to enable this.

\sphinxlineitem{\sphinxstyleemphasis{Pair lubricateU requires monodisperse particles}}
\sphinxAtStartPar
All particles must be the same finite size.

\sphinxlineitem{\sphinxstyleemphasis{Pair lubricateU/poly requires ghost atoms store velocity}}
\sphinxAtStartPar
Use the comm\_modify vel yes command to enable this.

\sphinxlineitem{\sphinxstyleemphasis{Pair lubricateU/poly requires newton pair off}}
\sphinxAtStartPar
Self\sphinxhyphen{}explanatory.

\sphinxlineitem{\sphinxstyleemphasis{Pair peri lattice is not identical in x, y, and z}}
\sphinxAtStartPar
The lattice defined by the lattice command must be cubic.

\sphinxlineitem{\sphinxstyleemphasis{Pair peri requires a lattice be defined}}
\sphinxAtStartPar
Use the lattice command for this purpose.

\sphinxlineitem{\sphinxstyleemphasis{Pair peri requires an atom map, see atom\_modify}}
\sphinxAtStartPar
Even for atomic systems, an atom map is required to find Peridynamic
bonds.  Use the atom\_modify command to define one.

\sphinxlineitem{\sphinxstyleemphasis{Pair resquared epsilon a,b,c coeffs are not all set}}
\sphinxAtStartPar
Self\sphinxhyphen{}explanatory.

\sphinxlineitem{\sphinxstyleemphasis{Pair resquared epsilon and sigma coeffs are not all set}}
\sphinxAtStartPar
Self\sphinxhyphen{}explanatory.

\sphinxlineitem{\sphinxstyleemphasis{Pair resquared requires atom style ellipsoid}}
\sphinxAtStartPar
Self\sphinxhyphen{}explanatory.

\sphinxlineitem{\sphinxstyleemphasis{Pair resquared requires atoms with same type have same shape}}
\sphinxAtStartPar
Self\sphinxhyphen{}explanatory.

\sphinxlineitem{\sphinxstyleemphasis{Pair resquared/gpu requires atom style ellipsoid}}
\sphinxAtStartPar
Self\sphinxhyphen{}explanatory.

\sphinxlineitem{\sphinxstyleemphasis{Pair resquared/gpu requires atoms with same type have same shape}}
\sphinxAtStartPar
Self\sphinxhyphen{}explanatory.

\sphinxlineitem{\sphinxstyleemphasis{Pair style AIREBO requires atom IDs}}
\sphinxAtStartPar
This is a requirement to use the AIREBO potential.

\sphinxlineitem{\sphinxstyleemphasis{Pair style AIREBO requires newton pair on}}
\sphinxAtStartPar
See the newton command.  This is a restriction to use the AIREBO
potential.

\sphinxlineitem{\sphinxstyleemphasis{Pair style BOP requires atom IDs}}
\sphinxAtStartPar
This is a requirement to use the BOP potential.

\sphinxlineitem{\sphinxstyleemphasis{Pair style BOP requires newton pair on}}
\sphinxAtStartPar
See the newton command.  This is a restriction to use the BOP
potential.

\sphinxlineitem{\sphinxstyleemphasis{Pair style COMB requires atom IDs}}
\sphinxAtStartPar
This is a requirement to use the AIREBO potential.

\sphinxlineitem{\sphinxstyleemphasis{Pair style COMB requires atom attribute q}}
\sphinxAtStartPar
Self\sphinxhyphen{}explanatory.

\sphinxlineitem{\sphinxstyleemphasis{Pair style COMB requires newton pair on}}
\sphinxAtStartPar
See the newton command.  This is a restriction to use the COMB
potential.

\sphinxlineitem{\sphinxstyleemphasis{Pair style COMB3 requires atom IDs}}
\sphinxAtStartPar
This is a requirement to use the COMB3 potential.

\sphinxlineitem{\sphinxstyleemphasis{Pair style COMB3 requires atom attribute q}}
\sphinxAtStartPar
Self\sphinxhyphen{}explanatory.

\sphinxlineitem{\sphinxstyleemphasis{Pair style COMB3 requires newton pair on}}
\sphinxAtStartPar
See the newton command.  This is a restriction to use the COMB3
potential.

\sphinxlineitem{\sphinxstyleemphasis{Pair style LCBOP requires atom IDs}}
\sphinxAtStartPar
This is a requirement to use the LCBOP potential.

\sphinxlineitem{\sphinxstyleemphasis{Pair style LCBOP requires newton pair on}}
\sphinxAtStartPar
See the newton command.  This is a restriction to use the Tersoff
potential.

\sphinxlineitem{\sphinxstyleemphasis{Pair style MEAM requires newton pair on}}
\sphinxAtStartPar
See the newton command.  This is a restriction to use the MEAM
potential.

\sphinxlineitem{\sphinxstyleemphasis{Pair style SNAP requires newton pair on}}
\sphinxAtStartPar
See the newton command.  This is a restriction to use the SNAP
potential.

\sphinxlineitem{\sphinxstyleemphasis{Pair style Stillinger\sphinxhyphen{}Weber requires atom IDs}}
\sphinxAtStartPar
This is a requirement to use the SW potential.

\sphinxlineitem{\sphinxstyleemphasis{Pair style Stillinger\sphinxhyphen{}Weber requires newton pair on}}
\sphinxAtStartPar
See the newton command.  This is a restriction to use the SW
potential.

\sphinxlineitem{\sphinxstyleemphasis{Pair style Tersoff requires atom IDs}}
\sphinxAtStartPar
This is a requirement to use the Tersoff potential.

\sphinxlineitem{\sphinxstyleemphasis{Pair style Tersoff requires newton pair on}}
\sphinxAtStartPar
See the newton command.  This is a restriction to use the Tersoff
potential.

\sphinxlineitem{\sphinxstyleemphasis{Pair style Vashishta requires atom IDs}}
\sphinxAtStartPar
This is a requirement to use the Vashishta potential.

\sphinxlineitem{\sphinxstyleemphasis{Pair style Vashishta requires newton pair on}}
\sphinxAtStartPar
See the newton command.  This is a restriction to use the Vashishta
potential.

\sphinxlineitem{\sphinxstyleemphasis{Pair style bop requires comm ghost cutoff at least 3x larger than \%g}}
\sphinxAtStartPar
Use the communicate ghost command to set this.  See the pair bop
page for more details.

\sphinxlineitem{\sphinxstyleemphasis{Pair style born/coul/long requires atom attribute q}}
\sphinxAtStartPar
An atom style that defines this attribute must be used.

\sphinxlineitem{\sphinxstyleemphasis{Pair style born/coul/long/gpu requires atom attribute q}}
\sphinxAtStartPar
The atom style defined does not have this attribute.

\sphinxlineitem{\sphinxstyleemphasis{Pair style born/coul/wolf requires atom attribute q}}
\sphinxAtStartPar
The atom style defined does not have this attribute.

\sphinxlineitem{\sphinxstyleemphasis{Pair style buck/coul/cut requires atom attribute q}}
\sphinxAtStartPar
The atom style defined does not have this attribute.

\sphinxlineitem{\sphinxstyleemphasis{Pair style buck/coul/long requires atom attribute q}}
\sphinxAtStartPar
The atom style defined does not have these attributes.

\sphinxlineitem{\sphinxstyleemphasis{Pair style buck/coul/long/gpu requires atom attribute q}}
\sphinxAtStartPar
The atom style defined does not have this attribute.

\sphinxlineitem{\sphinxstyleemphasis{Pair style buck/long/coul/long requires atom attribute q}}
\sphinxAtStartPar
The atom style defined does not have this attribute.

\sphinxlineitem{\sphinxstyleemphasis{Pair style coul/cut requires atom attribute q}}
\sphinxAtStartPar
The atom style defined does not have these attributes.

\sphinxlineitem{\sphinxstyleemphasis{Pair style coul/cut/gpu requires atom attribute q}}
\sphinxAtStartPar
The atom style defined does not have this attribute.

\sphinxlineitem{\sphinxstyleemphasis{Pair style coul/debye/gpu requires atom attribute q}}
\sphinxAtStartPar
The atom style defined does not have this attribute.

\sphinxlineitem{\sphinxstyleemphasis{Pair style coul/dsf requires atom attribute q}}
\sphinxAtStartPar
The atom style defined does not have this attribute.

\sphinxlineitem{\sphinxstyleemphasis{Pair style coul/dsf/gpu requires atom attribute q}}
\sphinxAtStartPar
The atom style defined does not have this attribute.

\sphinxlineitem{\sphinxstyleemphasis{Pair style coul/long/gpu requires atom attribute q}}
\sphinxAtStartPar
The atom style defined does not have these attributes.

\sphinxlineitem{\sphinxstyleemphasis{Pair style coul/streitz requires atom attribute q}}
\sphinxAtStartPar
Self\sphinxhyphen{}explanatory.

\sphinxlineitem{\sphinxstyleemphasis{Pair style does not have extra field requested by compute pair/local}}
\sphinxAtStartPar
The pair style does not support the pN value requested by the compute
pair/local command.

\sphinxlineitem{\sphinxstyleemphasis{Pair style does not support bond\_style quartic}}
\sphinxAtStartPar
The pair style does not have a single() function, so it can
not be invoked by bond\_style quartic.

\sphinxlineitem{\sphinxstyleemphasis{Pair style does not support compute group/group}}
\sphinxAtStartPar
The pair\_style does not have a single() function, so it cannot be
invoked by the compute group/group command.

\sphinxlineitem{\sphinxstyleemphasis{Pair style does not support compute pair/local}}
\sphinxAtStartPar
The pair style does not have a single() function, so it can
not be invoked by compute pair/local.

\sphinxlineitem{\sphinxstyleemphasis{Pair style does not support compute property/local}}
\sphinxAtStartPar
The pair style does not have a single() function, so it can
not be invoked by fix bond/swap.

\sphinxlineitem{\sphinxstyleemphasis{Pair style does not support fix bond/swap}}
\sphinxAtStartPar
The pair style does not have a single() function, so it can
not be invoked by fix bond/swap.

\sphinxlineitem{\sphinxstyleemphasis{Pair style does not support pair\_write}}
\sphinxAtStartPar
The pair style does not have a single() function, so it can
not be invoked by pair write.

\sphinxlineitem{\sphinxstyleemphasis{Pair style does not support rRESPA inner/middle/outer}}
\sphinxAtStartPar
You are attempting to use rRESPA options with a pair style that
does not support them.

\sphinxlineitem{\sphinxstyleemphasis{Pair style granular with history requires atoms have IDs}}
\sphinxAtStartPar
Atoms in the simulation do not have IDs, so history effects
cannot be tracked by the granular pair potential.

\sphinxlineitem{\sphinxstyleemphasis{Pair style hbond/dreiding requires an atom map, see atom\_modify}}
\sphinxAtStartPar
Self\sphinxhyphen{}explanatory.

\sphinxlineitem{\sphinxstyleemphasis{Pair style hbond/dreiding requires atom IDs}}
\sphinxAtStartPar
Self\sphinxhyphen{}explanatory.

\sphinxlineitem{\sphinxstyleemphasis{Pair style hbond/dreiding requires molecular system}}
\sphinxAtStartPar
Self\sphinxhyphen{}explanatory.

\sphinxlineitem{\sphinxstyleemphasis{Pair style hbond/dreiding requires newton pair on}}
\sphinxAtStartPar
See the newton command for details.

\sphinxlineitem{\sphinxstyleemphasis{Pair style hybrid cannot have hybrid as an argument}}
\sphinxAtStartPar
Self\sphinxhyphen{}explanatory.

\sphinxlineitem{\sphinxstyleemphasis{Pair style hybrid cannot have none as an argument}}
\sphinxAtStartPar
Self\sphinxhyphen{}explanatory.

\sphinxlineitem{\sphinxstyleemphasis{Pair style is incompatible with KSpace style}}
\sphinxAtStartPar
If a pair style with a long\sphinxhyphen{}range Coulombic component is selected,
then a kspace style must also be used.

\sphinxlineitem{\sphinxstyleemphasis{Pair style is incompatible with TIP4P KSpace style}}
\sphinxAtStartPar
The pair style does not have the requires TIP4P settings.

\sphinxlineitem{\sphinxstyleemphasis{Pair style lj/charmm/coul/charmm requires atom attribute q}}
\sphinxAtStartPar
The atom style defined does not have these attributes.

\sphinxlineitem{\sphinxstyleemphasis{Pair style lj/charmm/coul/long requires atom attribute q}}
\sphinxAtStartPar
The atom style defined does not have these attributes.

\sphinxlineitem{\sphinxstyleemphasis{Pair style lj/charmm/coul/long/gpu requires atom attribute q}}
\sphinxAtStartPar
The atom style defined does not have this attribute.

\sphinxlineitem{\sphinxstyleemphasis{Pair style lj/class2/coul/cut requires atom attribute q}}
\sphinxAtStartPar
The atom style defined does not have this attribute.

\sphinxlineitem{\sphinxstyleemphasis{Pair style lj/class2/coul/long requires atom attribute q}}
\sphinxAtStartPar
The atom style defined does not have this attribute.

\sphinxlineitem{\sphinxstyleemphasis{Pair style lj/class2/coul/long/gpu requires atom attribute q}}
\sphinxAtStartPar
The atom style defined does not have this attribute.

\sphinxlineitem{\sphinxstyleemphasis{Pair style lj/cut/coul/cut requires atom attribute q}}
\sphinxAtStartPar
The atom style defined does not have this attribute.

\sphinxlineitem{\sphinxstyleemphasis{Pair style lj/cut/coul/cut/gpu requires atom attribute q}}
\sphinxAtStartPar
The atom style defined does not have this attribute.

\sphinxlineitem{\sphinxstyleemphasis{Pair style lj/cut/coul/debye/gpu requires atom attribute q}}
\sphinxAtStartPar
The atom style defined does not have this attribute.

\sphinxlineitem{\sphinxstyleemphasis{Pair style lj/cut/coul/dsf requires atom attribute q}}
\sphinxAtStartPar
The atom style defined does not have these attributes.

\sphinxlineitem{\sphinxstyleemphasis{Pair style lj/cut/coul/dsf/gpu requires atom attribute q}}
\sphinxAtStartPar
The atom style defined does not have this attribute.

\sphinxlineitem{\sphinxstyleemphasis{Pair style lj/cut/coul/long requires atom attribute q}}
\sphinxAtStartPar
The atom style defined does not have this attribute.

\sphinxlineitem{\sphinxstyleemphasis{Pair style lj/cut/coul/long/gpu requires atom attribute q}}
\sphinxAtStartPar
The atom style defined does not have this attribute.

\sphinxlineitem{\sphinxstyleemphasis{Pair style lj/cut/tip4p/cut requires atom IDs}}
\sphinxAtStartPar
This is a requirement to use this potential.

\sphinxlineitem{\sphinxstyleemphasis{Pair style lj/cut/tip4p/cut requires atom attribute q}}
\sphinxAtStartPar
The atom style defined does not have this attribute.

\sphinxlineitem{\sphinxstyleemphasis{Pair style lj/cut/tip4p/cut requires newton pair on}}
\sphinxAtStartPar
See the newton command.  This is a restriction to use this
potential.

\sphinxlineitem{\sphinxstyleemphasis{Pair style lj/cut/tip4p/long requires atom IDs}}
\sphinxAtStartPar
There are no atom IDs defined in the system and the TIP4P potential
requires them to find O,H atoms with a water molecule.

\sphinxlineitem{\sphinxstyleemphasis{Pair style lj/cut/tip4p/long requires atom attribute q}}
\sphinxAtStartPar
The atom style defined does not have these attributes.

\sphinxlineitem{\sphinxstyleemphasis{Pair style lj/cut/tip4p/long requires newton pair on}}
\sphinxAtStartPar
This is because the computation of constraint forces within a water
molecule adds forces to atoms owned by other processors.

\sphinxlineitem{\sphinxstyleemphasis{Pair style lj/gromacs/coul/gromacs requires atom attribute q}}
\sphinxAtStartPar
An atom\_style with this attribute is needed.

\sphinxlineitem{\sphinxstyleemphasis{Pair style lj/long/dipole/long does not currently support respa}}
\sphinxAtStartPar
This feature is not yet supported.

\sphinxlineitem{\sphinxstyleemphasis{Pair style lj/long/tip4p/long requires atom IDs}}
\sphinxAtStartPar
There are no atom IDs defined in the system and the TIP4P potential
requires them to find O,H atoms with a water molecule.

\sphinxlineitem{\sphinxstyleemphasis{Pair style lj/long/tip4p/long requires atom attribute q}}
\sphinxAtStartPar
The atom style defined does not have these attributes.

\sphinxlineitem{\sphinxstyleemphasis{Pair style lj/long/tip4p/long requires newton pair on}}
\sphinxAtStartPar
This is because the computation of constraint forces within a water
molecule adds forces to atoms owned by other processors.

\sphinxlineitem{\sphinxstyleemphasis{Pair style lj/spica/coul/long/gpu requires atom attribute q}}
\sphinxAtStartPar
The atom style defined does not have this attribute.

\sphinxlineitem{\sphinxstyleemphasis{Pair style nb3b/harmonic requires atom IDs}}
\sphinxAtStartPar
This is a requirement to use this potential.

\sphinxlineitem{\sphinxstyleemphasis{Pair style nb3b/harmonic requires newton pair on}}
\sphinxAtStartPar
See the newton command.  This is a restriction to use this potential.

\sphinxlineitem{\sphinxstyleemphasis{Pair style nm/cut/coul/cut requires atom attribute q}}
\sphinxAtStartPar
The atom style defined does not have this attribute.

\sphinxlineitem{\sphinxstyleemphasis{Pair style nm/cut/coul/long requires atom attribute q}}
\sphinxAtStartPar
The atom style defined does not have this attribute.

\sphinxlineitem{\sphinxstyleemphasis{Pair style peri requires atom style peri}}
\sphinxAtStartPar
Self\sphinxhyphen{}explanatory.

\sphinxlineitem{\sphinxstyleemphasis{Pair style polymorphic requires atom IDs}}
\sphinxAtStartPar
This is a requirement to use the polymorphic potential.

\sphinxlineitem{\sphinxstyleemphasis{Pair style polymorphic requires newton pair on}}
\sphinxAtStartPar
See the newton command.  This is a restriction to use the polymorphic
potential.

\sphinxlineitem{\sphinxstyleemphasis{Pair style reax requires atom IDs}}
\sphinxAtStartPar
This is a requirement to use the ReaxFF potential.

\sphinxlineitem{\sphinxstyleemphasis{Pair style reax requires atom attribute q}}
\sphinxAtStartPar
The atom style defined does not have this attribute.

\sphinxlineitem{\sphinxstyleemphasis{Pair style reax requires newton pair on}}
\sphinxAtStartPar
This is a requirement to use the ReaxFF potential.

\sphinxlineitem{\sphinxstyleemphasis{Pair style requires a KSpace style}}
\sphinxAtStartPar
No kspace style is defined.

\sphinxlineitem{\sphinxstyleemphasis{Pair style requires use of kspace\_style ewald/disp}}
\sphinxAtStartPar
Self\sphinxhyphen{}explanatory.

\sphinxlineitem{\sphinxstyleemphasis{Pair style sw/gpu requires atom IDs}}
\sphinxAtStartPar
This is a requirement to use this potential.

\sphinxlineitem{\sphinxstyleemphasis{Pair style sw/gpu requires newton pair off}}
\sphinxAtStartPar
See the newton command.  This is a restriction to use this potential.

\sphinxlineitem{\sphinxstyleemphasis{Pair style vashishta/gpu requires atom IDs}}
\sphinxAtStartPar
This is a requirement to use this potential.

\sphinxlineitem{\sphinxstyleemphasis{Pair style vashishta/gpu requires newton pair off}}
\sphinxAtStartPar
See the newton command.  This is a restriction to use this potential.

\sphinxlineitem{\sphinxstyleemphasis{Pair style tersoff/gpu requires atom IDs}}
\sphinxAtStartPar
This is a requirement to use the tersoff/gpu potential.

\sphinxlineitem{\sphinxstyleemphasis{Pair style tersoff/gpu requires newton pair off}}
\sphinxAtStartPar
See the newton command.  This is a restriction to use this pair style.

\sphinxlineitem{\sphinxstyleemphasis{Pair style tip4p/cut requires atom IDs}}
\sphinxAtStartPar
This is a requirement to use this potential.

\sphinxlineitem{\sphinxstyleemphasis{Pair style tip4p/cut requires atom attribute q}}
\sphinxAtStartPar
The atom style defined does not have this attribute.

\sphinxlineitem{\sphinxstyleemphasis{Pair style tip4p/cut requires newton pair on}}
\sphinxAtStartPar
See the newton command.  This is a restriction to use this potential.

\sphinxlineitem{\sphinxstyleemphasis{Pair style tip4p/long requires atom IDs}}
\sphinxAtStartPar
There are no atom IDs defined in the system and the TIP4P potential
requires them to find O,H atoms with a water molecule.

\sphinxlineitem{\sphinxstyleemphasis{Pair style tip4p/long requires atom attribute q}}
\sphinxAtStartPar
The atom style defined does not have these attributes.

\sphinxlineitem{\sphinxstyleemphasis{Pair style tip4p/long requires newton pair on}}
\sphinxAtStartPar
This is because the computation of constraint forces within a water
molecule adds forces to atoms owned by other processors.

\sphinxlineitem{\sphinxstyleemphasis{Pair table cutoffs must all be equal to use with KSpace}}
\sphinxAtStartPar
When using pair style table with a long\sphinxhyphen{}range KSpace solver, the
cutoffs for all atom type pairs must all be the same, since the
long\sphinxhyphen{}range solver starts at that cutoff.

\sphinxlineitem{\sphinxstyleemphasis{Pair table parameters did not set N}}
\sphinxAtStartPar
List of pair table parameters must include N setting.

\sphinxlineitem{\sphinxstyleemphasis{Pair tersoff/zbl requires metal or real units}}
\sphinxAtStartPar
This is a current restriction of this pair potential.

\sphinxlineitem{\sphinxstyleemphasis{Pair tersoff/zbl/kk requires metal or real units}}
\sphinxAtStartPar
This is a current restriction of this pair potential.

\sphinxlineitem{\sphinxstyleemphasis{Pair tri/lj requires atom style tri}}
\sphinxAtStartPar
Self\sphinxhyphen{}explanatory.

\sphinxlineitem{\sphinxstyleemphasis{Pair yukawa/colloid requires atom style sphere}}
\sphinxAtStartPar
Self\sphinxhyphen{}explanatory.

\sphinxlineitem{\sphinxstyleemphasis{Pair yukawa/colloid requires atoms with same type have same radius}}
\sphinxAtStartPar
Self\sphinxhyphen{}explanatory.

\sphinxlineitem{\sphinxstyleemphasis{Pair yukawa/colloid/gpu requires atom style sphere}}
\sphinxAtStartPar
Self\sphinxhyphen{}explanatory.

\sphinxlineitem{\sphinxstyleemphasis{PairKIM only works with 3D problems}}
\sphinxAtStartPar
This is a current limitation.

\sphinxlineitem{\sphinxstyleemphasis{Pair\_coeff command before pair\_style is defined}}
\sphinxAtStartPar
Self\sphinxhyphen{}explanatory.

\sphinxlineitem{\sphinxstyleemphasis{Pair\_coeff command before simulation box is defined}}
\sphinxAtStartPar
The pair\_coeff command cannot be used before a read\_data,
read\_restart, or create\_box command.

\sphinxlineitem{\sphinxstyleemphasis{Pair\_modify command before pair\_style is defined}}
\sphinxAtStartPar
Self\sphinxhyphen{}explanatory.

\sphinxlineitem{\sphinxstyleemphasis{Pair\_modify special setting for pair hybrid incompatible with global special\_bonds setting}}
\sphinxAtStartPar
Cannot override a setting of 0.0 or 1.0 or change a setting between
0.0 and 1.0.

\sphinxlineitem{\sphinxstyleemphasis{Pair\_write command before pair\_style is defined}}
\sphinxAtStartPar
Self\sphinxhyphen{}explanatory.

\sphinxlineitem{\sphinxstyleemphasis{Particle on or inside fix wall surface}}
\sphinxAtStartPar
Particles must be “exterior” to the wall in order for energy/force to
be calculated.

\sphinxlineitem{\sphinxstyleemphasis{Particle outside surface of region used in fix wall/region}}
\sphinxAtStartPar
Particles must be inside the region for energy/force to be calculated.
A particle outside the region generates an error.

\sphinxlineitem{\sphinxstyleemphasis{Per\sphinxhyphen{}atom compute in equal\sphinxhyphen{}style variable formula}}
\sphinxAtStartPar
Equal\sphinxhyphen{}style variables cannot use per\sphinxhyphen{}atom quantities.

\sphinxlineitem{\sphinxstyleemphasis{Per\sphinxhyphen{}atom energy was not tallied on needed timestep}}
\sphinxAtStartPar
You are using a thermo keyword that requires potentials to
have tallied energy, but they did not on this timestep.  See the
variable page for ideas on how to make this work.

\sphinxlineitem{\sphinxstyleemphasis{Per\sphinxhyphen{}atom fix in equal\sphinxhyphen{}style variable formula}}
\sphinxAtStartPar
Equal\sphinxhyphen{}style variables cannot use per\sphinxhyphen{}atom quantities.

\sphinxlineitem{\sphinxstyleemphasis{Per\sphinxhyphen{}atom virial was not tallied on needed timestep}}
\sphinxAtStartPar
You are using a thermo keyword that requires potentials to have
tallied the virial, but they did not on this timestep.  See the
variable page for ideas on how to make this work.

\sphinxlineitem{\sphinxstyleemphasis{Per\sphinxhyphen{}processor system is too big}}
\sphinxAtStartPar
The number of owned atoms plus ghost atoms on a single
processor must fit in 32\sphinxhyphen{}bit integer.

\sphinxlineitem{\sphinxstyleemphasis{Potential energy ID for fix neb does not exist}}
\sphinxAtStartPar
Self\sphinxhyphen{}explanatory.

\sphinxlineitem{\sphinxstyleemphasis{Potential energy ID for fix nvt/nph/npt does not exist}}
\sphinxAtStartPar
A compute for potential energy must be defined.

\sphinxlineitem{\sphinxstyleemphasis{Potential file has duplicate entry}}
\sphinxAtStartPar
The potential file has more than one entry for the same element.

\sphinxlineitem{\sphinxstyleemphasis{Potential file is missing an entry}}
\sphinxAtStartPar
The potential file does not have a needed entry.

\sphinxlineitem{\sphinxstyleemphasis{Power by 0 in variable formula}}
\sphinxAtStartPar
Self\sphinxhyphen{}explanatory.

\sphinxlineitem{\sphinxstyleemphasis{Pressure ID for fix box/relax does not exist}}
\sphinxAtStartPar
The compute ID needed to compute pressure for the fix does not
exist.

\sphinxlineitem{\sphinxstyleemphasis{Pressure ID for fix modify does not exist}}
\sphinxAtStartPar
Self\sphinxhyphen{}explanatory.

\sphinxlineitem{\sphinxstyleemphasis{Pressure ID for fix npt/nph does not exist}}
\sphinxAtStartPar
Self\sphinxhyphen{}explanatory.

\sphinxlineitem{\sphinxstyleemphasis{Pressure ID for fix press/berendsen does not exist}}
\sphinxAtStartPar
The compute ID needed to compute pressure for the fix does not
exist.

\sphinxlineitem{\sphinxstyleemphasis{Pressure ID for fix rigid npt/nph does not exist}}
\sphinxAtStartPar
Self\sphinxhyphen{}explanatory.

\sphinxlineitem{\sphinxstyleemphasis{Pressure ID for thermo does not exist}}
\sphinxAtStartPar
The compute ID needed to compute pressure for thermodynamics does not
exist.

\sphinxlineitem{\sphinxstyleemphasis{Pressure control can not be used with fix nvt}}
\sphinxAtStartPar
Self\sphinxhyphen{}explanatory.

\sphinxlineitem{\sphinxstyleemphasis{Pressure control can not be used with fix nvt/asphere}}
\sphinxAtStartPar
Self\sphinxhyphen{}explanatory.

\sphinxlineitem{\sphinxstyleemphasis{Pressure control can not be used with fix nvt/body}}
\sphinxAtStartPar
Self\sphinxhyphen{}explanatory.

\sphinxlineitem{\sphinxstyleemphasis{Pressure control can not be used with fix nvt/sllod}}
\sphinxAtStartPar
Self\sphinxhyphen{}explanatory.

\sphinxlineitem{\sphinxstyleemphasis{Pressure control can not be used with fix nvt/sphere}}
\sphinxAtStartPar
Self\sphinxhyphen{}explanatory.

\sphinxlineitem{\sphinxstyleemphasis{Pressure control must be used with fix nph}}
\sphinxAtStartPar
Self\sphinxhyphen{}explanatory.

\sphinxlineitem{\sphinxstyleemphasis{Pressure control must be used with fix nph/asphere}}
\sphinxAtStartPar
Self\sphinxhyphen{}explanatory.

\sphinxlineitem{\sphinxstyleemphasis{Pressure control must be used with fix nph/body}}
\sphinxAtStartPar
Self\sphinxhyphen{}explanatory.

\sphinxlineitem{\sphinxstyleemphasis{Pressure control must be used with fix nph/small}}
\sphinxAtStartPar
Self\sphinxhyphen{}explanatory.

\sphinxlineitem{\sphinxstyleemphasis{Pressure control must be used with fix nph/sphere}}
\sphinxAtStartPar
Self\sphinxhyphen{}explanatory.

\sphinxlineitem{\sphinxstyleemphasis{Pressure control must be used with fix nphug}}
\sphinxAtStartPar
A pressure control keyword (iso, aniso, tri, x, y, or z) must be
provided.

\sphinxlineitem{\sphinxstyleemphasis{Pressure control must be used with fix npt}}
\sphinxAtStartPar
Self\sphinxhyphen{}explanatory.

\sphinxlineitem{\sphinxstyleemphasis{Pressure control must be used with fix npt/asphere}}
\sphinxAtStartPar
Self\sphinxhyphen{}explanatory.

\sphinxlineitem{\sphinxstyleemphasis{Pressure control must be used with fix npt/body}}
\sphinxAtStartPar
Self\sphinxhyphen{}explanatory.

\sphinxlineitem{\sphinxstyleemphasis{Pressure control must be used with fix npt/sphere}}
\sphinxAtStartPar
Self\sphinxhyphen{}explanatory.

\sphinxlineitem{\sphinxstyleemphasis{Processor count in z must be 1 for 2d simulation}}
\sphinxAtStartPar
Self\sphinxhyphen{}explanatory.

\sphinxlineitem{\sphinxstyleemphasis{Processor partitions do not match number of allocated processors}}
\sphinxAtStartPar
The total number of processors in all partitions must match the number
of processors LAMMPS is running on.

\sphinxlineitem{\sphinxstyleemphasis{Processors command after simulation box is defined}}
\sphinxAtStartPar
The processors command cannot be used after a read\_data, read\_restart,
or create\_box command.

\sphinxlineitem{\sphinxstyleemphasis{Processors custom grid file is inconsistent}}
\sphinxAtStartPar
The vales in the custom file are not consistent with the number of
processors you are running on or the Px,Py,Pz settings of the
processors command.  Or there was not a setting for every processor.

\sphinxlineitem{\sphinxstyleemphasis{Processors grid numa and map style are incompatible}}
\sphinxAtStartPar
Using numa for gstyle in the processors command requires using
cart for the map option.

\sphinxlineitem{\sphinxstyleemphasis{Processors part option and grid style are incompatible}}
\sphinxAtStartPar
Cannot use gstyle numa or custom with the part option.

\sphinxlineitem{\sphinxstyleemphasis{Processors twogrid requires proc count be a multiple of core count}}
\sphinxAtStartPar
Self\sphinxhyphen{}explanatory.

\sphinxlineitem{\sphinxstyleemphasis{Pstart and Pstop must have the same value}}
\sphinxAtStartPar
Self\sphinxhyphen{}explanatory.

\sphinxlineitem{\sphinxstyleemphasis{Python function evaluation failed}}
\sphinxAtStartPar
The Python function did not run successfully and/or did not return a
value (if it is supposed to return a value).  This is probably due to
some error condition in the function.

\sphinxlineitem{\sphinxstyleemphasis{Python function is not callable}}
\sphinxAtStartPar
The provided Python code was run successfully, but it not
define a callable function with the required name.

\sphinxlineitem{\sphinxstyleemphasis{Python invoke of undefined function}}
\sphinxAtStartPar
Cannot invoke a function that has not been previously defined.

\sphinxlineitem{\sphinxstyleemphasis{Python variable does not match Python function}}
\sphinxAtStartPar
This matching is defined by the python\sphinxhyphen{}style variable and the python
command.

\sphinxlineitem{\sphinxstyleemphasis{Python variable has no function}}
\sphinxAtStartPar
No python command was used to define the function associated with the
python\sphinxhyphen{}style variable.

\sphinxlineitem{\sphinxstyleemphasis{QEQ with ‘newton pair off’ not supported}}
\sphinxAtStartPar
See the newton command.  This is a restriction to use the QEQ fixes.

\sphinxlineitem{\sphinxstyleemphasis{R0 \textless{} 0 for fix spring command}}
\sphinxAtStartPar
Equilibrium spring length is invalid.

\sphinxlineitem{\sphinxstyleemphasis{RATTLE coordinate constraints are not satisfied up to desired tolerance}}
\sphinxAtStartPar
Self\sphinxhyphen{}explanatory.

\sphinxlineitem{\sphinxstyleemphasis{RATTLE determinant = 0.0}}
\sphinxAtStartPar
The determinant of the matrix being solved for a single cluster
specified by the fix rattle command is numerically invalid.

\sphinxlineitem{\sphinxstyleemphasis{RATTLE failed}}
\sphinxAtStartPar
Certain constraints were not satisfied.

\sphinxlineitem{\sphinxstyleemphasis{RATTLE velocity constraints are not satisfied up to desired tolerance}}
\sphinxAtStartPar
Self\sphinxhyphen{}explanatory.

\sphinxlineitem{\sphinxstyleemphasis{Read data add offset is too big}}
\sphinxAtStartPar
It cannot be larger than the size of atom IDs, e.g. the maximum 32\sphinxhyphen{}bit
integer.

\sphinxlineitem{\sphinxstyleemphasis{Read dump of atom property that is not allocated}}
\sphinxAtStartPar
Self\sphinxhyphen{}explanatory.

\sphinxlineitem{\sphinxstyleemphasis{Read rerun dump file timestep \textgreater{} specified stop}}
\sphinxAtStartPar
Self\sphinxhyphen{}explanatory.

\sphinxlineitem{\sphinxstyleemphasis{Read restart MPI\sphinxhyphen{}IO input not allowed with \% in filename}}
\sphinxAtStartPar
This is because a \% signifies one file per processor and MPI\sphinxhyphen{}IO
creates one large file for all processors.

\sphinxlineitem{\sphinxstyleemphasis{Read\_data shrink wrap did not assign all atoms correctly}}
\sphinxAtStartPar
This is typically because the box\sphinxhyphen{}size specified in the data file is
large compared to the actual extent of atoms in a shrink\sphinxhyphen{}wrapped
dimension.  When LAMMPS shrink\sphinxhyphen{}wraps the box atoms will be lost if the
processor they are re\sphinxhyphen{}assigned to is too far away.  Choose a box
size closer to the actual extent of the atoms.

\sphinxlineitem{\sphinxstyleemphasis{Read\_dump command before simulation box is defined}}
\sphinxAtStartPar
The read\_dump command cannot be used before a read\_data, read\_restart,
or create\_box command.

\sphinxlineitem{\sphinxstyleemphasis{Read\_dump field not found in dump file}}
\sphinxAtStartPar
Self\sphinxhyphen{}explanatory.

\sphinxlineitem{\sphinxstyleemphasis{Read\_dump triclinic status does not match simulation}}
\sphinxAtStartPar
Both the dump snapshot and the current LAMMPS simulation must
be using either an orthogonal or triclinic box.

\sphinxlineitem{\sphinxstyleemphasis{Read\_dump xyz fields do not have consistent scaling/wrapping}}
\sphinxAtStartPar
Self\sphinxhyphen{}explanatory.

\sphinxlineitem{\sphinxstyleemphasis{Reax\_defs.h setting for NATDEF is too small}}
\sphinxAtStartPar
Edit the setting in the ReaxFF library and re\sphinxhyphen{}compile the
library and re\sphinxhyphen{}build LAMMPS.

\sphinxlineitem{\sphinxstyleemphasis{Reax\_defs.h setting for NNEIGHMAXDEF is too small}}
\sphinxAtStartPar
Edit the setting in the ReaxFF library and re\sphinxhyphen{}compile the
library and re\sphinxhyphen{}build LAMMPS.

\sphinxlineitem{\sphinxstyleemphasis{Receiving partition in processors part command is already a receiver}}
\sphinxAtStartPar
Cannot specify a partition to be a receiver twice.

\sphinxlineitem{\sphinxstyleemphasis{Region ID for compute chunk/atom does not exist}}
\sphinxAtStartPar
Self\sphinxhyphen{}explanatory.

\sphinxlineitem{\sphinxstyleemphasis{Region ID for compute reduce/region does not exist}}
\sphinxAtStartPar
Self\sphinxhyphen{}explanatory.

\sphinxlineitem{\sphinxstyleemphasis{Region ID for compute temp/region does not exist}}
\sphinxAtStartPar
Self\sphinxhyphen{}explanatory.

\sphinxlineitem{\sphinxstyleemphasis{Region ID for dump custom does not exist}}
\sphinxAtStartPar
Self\sphinxhyphen{}explanatory.

\sphinxlineitem{\sphinxstyleemphasis{Region ID for fix addforce does not exist}}
\sphinxAtStartPar
Self\sphinxhyphen{}explanatory.

\sphinxlineitem{\sphinxstyleemphasis{Region ID for fix atom/swap does not exist}}
\sphinxAtStartPar
Self\sphinxhyphen{}explanatory.

\sphinxlineitem{\sphinxstyleemphasis{Region ID for fix ave/spatial does not exist}}
\sphinxAtStartPar
Self\sphinxhyphen{}explanatory.

\sphinxlineitem{\sphinxstyleemphasis{Region ID for fix aveforce does not exist}}
\sphinxAtStartPar
Self\sphinxhyphen{}explanatory.

\sphinxlineitem{\sphinxstyleemphasis{Region ID for fix deposit does not exist}}
\sphinxAtStartPar
Self\sphinxhyphen{}explanatory.

\sphinxlineitem{\sphinxstyleemphasis{Region ID for fix efield does not exist}}
\sphinxAtStartPar
Self\sphinxhyphen{}explanatory.

\sphinxlineitem{\sphinxstyleemphasis{Region ID for fix evaporate does not exist}}
\sphinxAtStartPar
Self\sphinxhyphen{}explanatory.

\sphinxlineitem{\sphinxstyleemphasis{Region ID for fix gcmc does not exist}}
\sphinxAtStartPar
Self\sphinxhyphen{}explanatory.

\sphinxlineitem{\sphinxstyleemphasis{Region ID for fix heat does not exist}}
\sphinxAtStartPar
Self\sphinxhyphen{}explanatory.

\sphinxlineitem{\sphinxstyleemphasis{Region ID for fix setforce does not exist}}
\sphinxAtStartPar
Self\sphinxhyphen{}explanatory.

\sphinxlineitem{\sphinxstyleemphasis{Region ID for fix wall/region does not exist}}
\sphinxAtStartPar
Self\sphinxhyphen{}explanatory.

\sphinxlineitem{\sphinxstyleemphasis{Region ID for group dynamic does not exist}}
\sphinxAtStartPar
Self\sphinxhyphen{}explanatory.

\sphinxlineitem{\sphinxstyleemphasis{Region ID in variable formula does not exist}}
\sphinxAtStartPar
Self\sphinxhyphen{}explanatory.

\sphinxlineitem{\sphinxstyleemphasis{Region cannot have 0 length rotation vector}}
\sphinxAtStartPar
Self\sphinxhyphen{}explanatory.

\sphinxlineitem{\sphinxstyleemphasis{Region for fix oneway does not exist}}
\sphinxAtStartPar
Self\sphinxhyphen{}explanatory.

\sphinxlineitem{\sphinxstyleemphasis{Region intersect region ID does not exist}}
\sphinxAtStartPar
Self\sphinxhyphen{}explanatory.

\sphinxlineitem{\sphinxstyleemphasis{Region union or intersect cannot be dynamic}}
\sphinxAtStartPar
The sub\sphinxhyphen{}regions can be dynamic, but not the combined region.

\sphinxlineitem{\sphinxstyleemphasis{Region union region ID does not exist}}
\sphinxAtStartPar
One or more of the region IDs specified by the region union command
does not exist.

\sphinxlineitem{\sphinxstyleemphasis{Replacing a fix, but new style != old style}}
\sphinxAtStartPar
A fix ID can be used a second time, but only if the style matches the
previous fix.  In this case it is assumed you want to reset a fix’s
parameters.  This error may mean you are mistakenly re\sphinxhyphen{}using a fix ID
when you do not intend to.

\sphinxlineitem{\sphinxstyleemphasis{Replicate command before simulation box is defined}}
\sphinxAtStartPar
The replicate command cannot be used before a read\_data, read\_restart,
or create\_box command.

\sphinxlineitem{\sphinxstyleemphasis{Replicate did not assign all atoms correctly}}
\sphinxAtStartPar
Atoms replicated by the replicate command were not assigned correctly
to processors.  This is likely due to some atom coordinates being
outside a non\sphinxhyphen{}periodic simulation box.

\sphinxlineitem{\sphinxstyleemphasis{Replicated system atom IDs are too big}}
\sphinxAtStartPar
See the setting for tagint in the src/lmptype.h file.

\sphinxlineitem{\sphinxstyleemphasis{Replicated system is too big}}
\sphinxAtStartPar
See the setting for bigint in the src/lmptype.h file.

\sphinxlineitem{\sphinxstyleemphasis{Required border comm not yet implemented with Kokkos}}
\sphinxAtStartPar
There are various limitations in the communication options supported
by Kokkos.

\sphinxlineitem{\sphinxstyleemphasis{Rerun command before simulation box is defined}}
\sphinxAtStartPar
The rerun command cannot be used before a read\_data, read\_restart, or
create\_box command.

\sphinxlineitem{\sphinxstyleemphasis{Rerun dump file does not contain requested snapshot}}
\sphinxAtStartPar
Self\sphinxhyphen{}explanatory.

\sphinxlineitem{\sphinxstyleemphasis{Resetting timestep size is not allowed with fix move}}
\sphinxAtStartPar
This is because fix move is moving atoms based on elapsed time.

\sphinxlineitem{\sphinxstyleemphasis{Respa inner cutoffs are invalid}}
\sphinxAtStartPar
The first cutoff must be \textless{}= the second cutoff.

\sphinxlineitem{\sphinxstyleemphasis{Respa levels must be \textgreater{}= 1}}
\sphinxAtStartPar
Self\sphinxhyphen{}explanatory.

\sphinxlineitem{\sphinxstyleemphasis{Respa middle cutoffs are invalid}}
\sphinxAtStartPar
The first cutoff must be \textless{}= the second cutoff.

\sphinxlineitem{\sphinxstyleemphasis{Restart file MPI\sphinxhyphen{}IO output not allowed with \% in filename}}
\sphinxAtStartPar
This is because a \% signifies one file per processor and MPI\sphinxhyphen{}IO
creates one large file for all processors.

\sphinxlineitem{\sphinxstyleemphasis{Restart file byte ordering is not recognized}}
\sphinxAtStartPar
The file does not appear to be a LAMMPS restart file since it does not
contain a recognized byte\sphinxhyphen{}ordering flag at the beginning.

\sphinxlineitem{\sphinxstyleemphasis{Restart file byte ordering is swapped}}
\sphinxAtStartPar
The file was written on a machine with different byte\sphinxhyphen{}ordering than
the machine you are reading it on.  Convert it to a text data file
instead, on the machine you wrote it on.

\sphinxlineitem{\sphinxstyleemphasis{Restart file incompatible with current version}}
\sphinxAtStartPar
This is probably because you are trying to read a file created with a
version of LAMMPS that is too old compared to the current version.
Use your older version of LAMMPS and convert the restart file
to a data file.

\sphinxlineitem{\sphinxstyleemphasis{Restart file is a MPI\sphinxhyphen{}IO file}}
\sphinxAtStartPar
The file is inconsistent with the filename you specified for it.

\sphinxlineitem{\sphinxstyleemphasis{Restart file is a multi\sphinxhyphen{}proc file}}
\sphinxAtStartPar
The file is inconsistent with the filename you specified for it.

\sphinxlineitem{\sphinxstyleemphasis{Restart file is not a MPI\sphinxhyphen{}IO file}}
\sphinxAtStartPar
The file is inconsistent with the filename you specified for it.

\sphinxlineitem{\sphinxstyleemphasis{Restart file is not a multi\sphinxhyphen{}proc file}}
\sphinxAtStartPar
The file is inconsistent with the filename you specified for it.

\sphinxlineitem{\sphinxstyleemphasis{Restart variable returned a bad timestep}}
\sphinxAtStartPar
The variable must return a timestep greater than the current timestep.

\sphinxlineitem{\sphinxstyleemphasis{Restrain atoms \%d \%d \%d \%d missing on proc \%d at step \%ld}}
\sphinxAtStartPar
The 4 atoms in a restrain dihedral specified by the fix restrain
command are not all accessible to a processor.  This probably means an
atom has moved too far.

\sphinxlineitem{\sphinxstyleemphasis{Restrain atoms \%d \%d \%d missing on proc \%d at step \%ld}}
\sphinxAtStartPar
The three atoms in a restrain angle specified by the fix restrain
command are not all accessible to a processor.  This probably means an
atom has moved too far.

\sphinxlineitem{\sphinxstyleemphasis{Restrain atoms \%d \%d missing on proc \%d at step \%ld}}
\sphinxAtStartPar
The two atoms in a restrain bond specified by the fix restrain
command are not all accessible to a processor.  This probably means an
atom has moved too far.

\sphinxlineitem{\sphinxstyleemphasis{Reuse of compute ID}}
\sphinxAtStartPar
A compute ID cannot be used twice.

\sphinxlineitem{\sphinxstyleemphasis{Reuse of dump ID}}
\sphinxAtStartPar
A dump ID cannot be used twice.

\sphinxlineitem{\sphinxstyleemphasis{Reuse of molecule template ID}}
\sphinxAtStartPar
The template IDs must be unique.

\sphinxlineitem{\sphinxstyleemphasis{Reuse of region ID}}
\sphinxAtStartPar
A region ID cannot be used twice.

\sphinxlineitem{\sphinxstyleemphasis{Rigid body atoms \%d \%d missing on proc \%d at step \%ld}}
\sphinxAtStartPar
This means that an atom cannot find the atom that owns the rigid body
it is part of, or vice versa.  The solution is to use the communicate
cutoff command to ensure ghost atoms are acquired from far enough away
to encompass the max distance printed when the fix rigid/small command
was invoked.

\sphinxlineitem{\sphinxstyleemphasis{Rigid body has degenerate moment of inertia}}
\sphinxAtStartPar
Fix poems will only work with bodies (collections of atoms) that have
non\sphinxhyphen{}zero principal moments of inertia.  This means they must be 3 or
more non\sphinxhyphen{}collinear atoms, even with joint atoms removed.

\sphinxlineitem{\sphinxstyleemphasis{Rigid fix must come before NPT/NPH fix}}
\sphinxAtStartPar
NPT/NPH fix must be defined in input script after all rigid fixes,
else the rigid fix contribution to the pressure virial is
incorrect.

\sphinxlineitem{\sphinxstyleemphasis{Rmask function in equal\sphinxhyphen{}style variable formula}}
\sphinxAtStartPar
Rmask is per\sphinxhyphen{}atom operation.

\sphinxlineitem{\sphinxstyleemphasis{Run command before simulation box is defined}}
\sphinxAtStartPar
The run command cannot be used before a read\_data, read\_restart, or
create\_box command.

\sphinxlineitem{\sphinxstyleemphasis{Run command start value is after start of run}}
\sphinxAtStartPar
Self\sphinxhyphen{}explanatory.

\sphinxlineitem{\sphinxstyleemphasis{Run command stop value is before end of run}}
\sphinxAtStartPar
Self\sphinxhyphen{}explanatory.

\sphinxlineitem{\sphinxstyleemphasis{Run\_style command before simulation box is defined}}
\sphinxAtStartPar
The run\_style command cannot be used before a read\_data,
read\_restart, or create\_box command.

\sphinxlineitem{\sphinxstyleemphasis{SRD bin size for fix srd differs from user request}}
\sphinxAtStartPar
Fix SRD had to adjust the bin size to fit the simulation box.  See the
cubic keyword if you want this message to be an error vs warning.

\sphinxlineitem{\sphinxstyleemphasis{SRD bins for fix srd are not cubic enough}}
\sphinxAtStartPar
The bin shape is not within tolerance of cubic.  See the cubic
keyword if you want this message to be an error vs warning.

\sphinxlineitem{\sphinxstyleemphasis{SRD particle \%d started inside big particle \%d on step \%ld bounce \%d}}
\sphinxAtStartPar
See the inside keyword if you want this message to be an error vs
warning.

\sphinxlineitem{\sphinxstyleemphasis{SRD particle \%d started inside wall \%d on step \%ld bounce \%d}}
\sphinxAtStartPar
See the inside keyword if you want this message to be an error vs
warning.

\sphinxlineitem{\sphinxstyleemphasis{Same dimension twice in fix ave/spatial}}
\sphinxAtStartPar
Self\sphinxhyphen{}explanatory.

\sphinxlineitem{\sphinxstyleemphasis{Sending partition in processors part command is already a sender}}
\sphinxAtStartPar
Cannot specify a partition to be a sender twice.

\sphinxlineitem{\sphinxstyleemphasis{Set command before simulation box is defined}}
\sphinxAtStartPar
The set command cannot be used before a read\_data, read\_restart,
or create\_box command.

\sphinxlineitem{\sphinxstyleemphasis{Set command floating point vector does not exist}}
\sphinxAtStartPar
Self\sphinxhyphen{}explanatory.

\sphinxlineitem{\sphinxstyleemphasis{Set command integer vector does not exist}}
\sphinxAtStartPar
Self\sphinxhyphen{}explanatory.

\sphinxlineitem{\sphinxstyleemphasis{Set command with no atoms existing}}
\sphinxAtStartPar
No atoms are yet defined so the set command cannot be used.

\sphinxlineitem{\sphinxstyleemphasis{Set region ID does not exist}}
\sphinxAtStartPar
Region ID specified in set command does not exist.

\sphinxlineitem{\sphinxstyleemphasis{Shake angles have different bond types}}
\sphinxAtStartPar
All 3\sphinxhyphen{}atom angle\sphinxhyphen{}constrained SHAKE clusters specified by the fix shake
command that are the same angle type, must also have the same bond
types for the two bonds in the angle.

\sphinxlineitem{\sphinxstyleemphasis{Shake atoms \%d \%d \%d \%d missing on proc \%d at step \%ld}}
\sphinxAtStartPar
The 4 atoms in a single shake cluster specified by the fix shake
command are not all accessible to a processor.  This probably means
an atom has moved too far.

\sphinxlineitem{\sphinxstyleemphasis{Shake atoms \%d \%d \%d missing on proc \%d at step \%ld}}
\sphinxAtStartPar
The three atoms in a single shake cluster specified by the fix shake
command are not all accessible to a processor.  This probably means
an atom has moved too far.

\sphinxlineitem{\sphinxstyleemphasis{Shake atoms \%d \%d missing on proc \%d at step \%ld}}
\sphinxAtStartPar
The two atoms in a single shake cluster specified by the fix shake
command are not all accessible to a processor.  This probably means
an atom has moved too far.

\sphinxlineitem{\sphinxstyleemphasis{Shake cluster of more than 4 atoms}}
\sphinxAtStartPar
A single cluster specified by the fix shake command can have no more
than 4 atoms.

\sphinxlineitem{\sphinxstyleemphasis{Shake clusters are connected}}
\sphinxAtStartPar
A single cluster specified by the fix shake command must have a single
central atom with up to 3 other atoms bonded to it.

\sphinxlineitem{\sphinxstyleemphasis{Shake determinant = 0.0}}
\sphinxAtStartPar
The determinant of the matrix being solved for a single cluster
specified by the fix shake command is numerically invalid.

\sphinxlineitem{\sphinxstyleemphasis{Shake fix must come before NPT/NPH fix}}
\sphinxAtStartPar
NPT fix must be defined in input script after SHAKE fix, else the
SHAKE fix contribution to the pressure virial is incorrect.

\sphinxlineitem{\sphinxstyleemphasis{Shear history overflow, boost neigh\_modify one}}
\sphinxAtStartPar
There are too many neighbors of a single atom.  Use the neigh\_modify
command to increase the max number of neighbors allowed for one atom.
You may also want to boost the page size.

\sphinxlineitem{\sphinxstyleemphasis{Small to big integers are not sized correctly}}
\sphinxAtStartPar
This error occurs when the sizes of smallint, imageint, tagint, bigint,
as defined in src/lmptype.h are not what is expected.  Contact
the developers if this occurs.

\sphinxlineitem{\sphinxstyleemphasis{Smallint setting in lmptype.h is invalid}}
\sphinxAtStartPar
It has to be the size of an integer.

\sphinxlineitem{\sphinxstyleemphasis{Smallint setting in lmptype.h is not compatible}}
\sphinxAtStartPar
Smallint stored in restart file is not consistent with LAMMPS version
you are running.

\sphinxlineitem{\sphinxstyleemphasis{Special list size exceeded in fix bond/create}}
\sphinxAtStartPar
See the “read\_data extra/special/per/atom” command
(or the “create\_box extra/special/per/atom” command)
for info on how to leave space in the special bonds
list to allow for additional bonds to be formed.

\sphinxlineitem{\sphinxstyleemphasis{Species XXX is not supported by this KIM Simulator Model}}
\sphinxAtStartPar
The kim\_style define command was referencing a species that is not
present in the requested KIM Simulator Model.

\sphinxlineitem{\sphinxstyleemphasis{Specified processors != physical processors}}
\sphinxAtStartPar
The 3d grid of processors defined by the processors command does not
match the number of processors LAMMPS is being run on.

\sphinxlineitem{\sphinxstyleemphasis{Specified target stress must be uniaxial or hydrostatic}}
\sphinxAtStartPar
Self\sphinxhyphen{}explanatory.

\sphinxlineitem{\sphinxstyleemphasis{Sqrt of negative value in variable formula}}
\sphinxAtStartPar
Self\sphinxhyphen{}explanatory.

\sphinxlineitem{\sphinxstyleemphasis{Subsequent read data induced too many angles per atom}}
\sphinxAtStartPar
See the extra/angle/per/atom keyword for the create\_box
or the read\_data command to set this limit larger

\sphinxlineitem{\sphinxstyleemphasis{Subsequent read data induced too many bonds per atom}}
\sphinxAtStartPar
See the extra/bond/per/atom keyword for the create\_box
or the read\_data command to set this limit larger

\sphinxlineitem{\sphinxstyleemphasis{Subsequent read data induced too many dihedrals per atom}}
\sphinxAtStartPar
See the extra/dihedral/per/atom keyword for the create\_box
or the read\_data command to set this limit larger

\sphinxlineitem{\sphinxstyleemphasis{Subsequent read data induced too many impropers per atom}}
\sphinxAtStartPar
See the extra/improper/per/atom keyword for the create\_box
or the read\_data command to set this limit larger

\sphinxlineitem{\sphinxstyleemphasis{Substitution for illegal variable}}
\sphinxAtStartPar
Input script line contained a variable that could not be substituted
for.

\sphinxlineitem{\sphinxstyleemphasis{Support for writing images in JPEG format not included}}
\sphinxAtStartPar
LAMMPS was not built with the \sphinxhyphen{}DLAMMPS\_JPEG switch in the Makefile.

\sphinxlineitem{\sphinxstyleemphasis{Support for writing images in PNG format not included}}
\sphinxAtStartPar
LAMMPS was not built with the \sphinxhyphen{}DLAMMPS\_PNG switch in the Makefile.

\sphinxlineitem{\sphinxstyleemphasis{Support for writing movies not included}}
\sphinxAtStartPar
LAMMPS was not built with the \sphinxhyphen{}DLAMMPS\_FFMPEG switch in the Makefile

\sphinxlineitem{\sphinxstyleemphasis{System in data file is too big}}
\sphinxAtStartPar
See the setting for bigint in the src/lmptype.h file.

\sphinxlineitem{\sphinxstyleemphasis{System is not charge neutral, net charge = \%g}}
\sphinxAtStartPar
The total charge on all atoms on the system is not 0.0.
For some KSpace solvers this is an error.

\sphinxlineitem{\sphinxstyleemphasis{TAD nsteps must be multiple of t\_event}}
\sphinxAtStartPar
Self\sphinxhyphen{}explanatory.

\sphinxlineitem{\sphinxstyleemphasis{TIP4P hydrogen has incorrect atom type}}
\sphinxAtStartPar
The TIP4P pairwise computation found an H atom whose type does not
agree with the specified H type.

\sphinxlineitem{\sphinxstyleemphasis{TIP4P hydrogen is missing}}
\sphinxAtStartPar
The TIP4P pairwise computation failed to find the correct H atom
within a water molecule.

\sphinxlineitem{\sphinxstyleemphasis{TMD target file did not list all group atoms}}
\sphinxAtStartPar
The target file for the fix tmd command did not list all atoms in the
fix group.

\sphinxlineitem{\sphinxstyleemphasis{Tad command before simulation box is defined}}
\sphinxAtStartPar
Self\sphinxhyphen{}explanatory.

\sphinxlineitem{\sphinxstyleemphasis{Tagint setting in lmptype.h is invalid}}
\sphinxAtStartPar
Tagint must be as large or larger than smallint.

\sphinxlineitem{\sphinxstyleemphasis{Tagint setting in lmptype.h is not compatible}}
\sphinxAtStartPar
Format of tagint stored in restart file is not consistent with LAMMPS
version you are running.  See the settings in src/lmptype.h

\sphinxlineitem{\sphinxstyleemphasis{Target pressure for fix rigid/nph cannot be \textless{} 0.0}}
\sphinxAtStartPar
Self\sphinxhyphen{}explanatory.

\sphinxlineitem{\sphinxstyleemphasis{Target pressure for fix rigid/npt/small cannot be \textless{} 0.0}}
\sphinxAtStartPar
Self\sphinxhyphen{}explanatory.

\sphinxlineitem{\sphinxstyleemphasis{Target temperature for fix nvt/npt/nph cannot be 0.0}}
\sphinxAtStartPar
Self\sphinxhyphen{}explanatory.

\sphinxlineitem{\sphinxstyleemphasis{Target temperature for fix rigid/npt cannot be 0.0}}
\sphinxAtStartPar
Self\sphinxhyphen{}explanatory.

\sphinxlineitem{\sphinxstyleemphasis{Target temperature for fix rigid/npt/small cannot be 0.0}}
\sphinxAtStartPar
Self\sphinxhyphen{}explanatory.

\sphinxlineitem{\sphinxstyleemphasis{Target temperature for fix rigid/nvt cannot be 0.0}}
\sphinxAtStartPar
Self\sphinxhyphen{}explanatory.

\sphinxlineitem{\sphinxstyleemphasis{Target temperature for fix rigid/nvt/small cannot be 0.0}}
\sphinxAtStartPar
Self\sphinxhyphen{}explanatory.

\sphinxlineitem{\sphinxstyleemphasis{Temper command before simulation box is defined}}
\sphinxAtStartPar
The temper command cannot be used before a read\_data, read\_restart, or
create\_box command.

\sphinxlineitem{\sphinxstyleemphasis{Temperature ID for fix bond/swap does not exist}}
\sphinxAtStartPar
Self\sphinxhyphen{}explanatory.

\sphinxlineitem{\sphinxstyleemphasis{Temperature ID for fix box/relax does not exist}}
\sphinxAtStartPar
Self\sphinxhyphen{}explanatory.

\sphinxlineitem{\sphinxstyleemphasis{Temperature ID for fix nvt/npt does not exist}}
\sphinxAtStartPar
Self\sphinxhyphen{}explanatory.

\sphinxlineitem{\sphinxstyleemphasis{Temperature ID for fix press/berendsen does not exist}}
\sphinxAtStartPar
Self\sphinxhyphen{}explanatory.

\sphinxlineitem{\sphinxstyleemphasis{Temperature ID for fix rigid nvt/npt/nph does not exist}}
\sphinxAtStartPar
Self\sphinxhyphen{}explanatory.

\sphinxlineitem{\sphinxstyleemphasis{Temperature ID for fix temp/berendsen does not exist}}
\sphinxAtStartPar
Self\sphinxhyphen{}explanatory.

\sphinxlineitem{\sphinxstyleemphasis{Temperature ID for fix temp/csld does not exist}}
\sphinxAtStartPar
Self\sphinxhyphen{}explanatory.

\sphinxlineitem{\sphinxstyleemphasis{Temperature ID for fix temp/csvr does not exist}}
\sphinxAtStartPar
Self\sphinxhyphen{}explanatory.

\sphinxlineitem{\sphinxstyleemphasis{Temperature ID for fix temp/rescale does not exist}}
\sphinxAtStartPar
Self\sphinxhyphen{}explanatory.

\sphinxlineitem{\sphinxstyleemphasis{Temperature compute degrees of freedom \textless{} 0}}
\sphinxAtStartPar
This should not happen if you are calculating the temperature
on a valid set of atoms.

\sphinxlineitem{\sphinxstyleemphasis{Temperature control can not be used with fix nph}}
\sphinxAtStartPar
Self\sphinxhyphen{}explanatory.

\sphinxlineitem{\sphinxstyleemphasis{Temperature control can not be used with fix nph/asphere}}
\sphinxAtStartPar
Self\sphinxhyphen{}explanatory.

\sphinxlineitem{\sphinxstyleemphasis{Temperature control can not be used with fix nph/body}}
\sphinxAtStartPar
Self\sphinxhyphen{}explanatory.

\sphinxlineitem{\sphinxstyleemphasis{Temperature control can not be used with fix nph/sphere}}
\sphinxAtStartPar
Self\sphinxhyphen{}explanatory.

\sphinxlineitem{\sphinxstyleemphasis{Temperature control must be used with fix nphug}}
\sphinxAtStartPar
The temp keyword must be provided.

\sphinxlineitem{\sphinxstyleemphasis{Temperature control must be used with fix npt}}
\sphinxAtStartPar
Self\sphinxhyphen{}explanatory.

\sphinxlineitem{\sphinxstyleemphasis{Temperature control must be used with fix npt/asphere}}
\sphinxAtStartPar
Self\sphinxhyphen{}explanatory.

\sphinxlineitem{\sphinxstyleemphasis{Temperature control must be used with fix npt/body}}
\sphinxAtStartPar
Self\sphinxhyphen{}explanatory.

\sphinxlineitem{\sphinxstyleemphasis{Temperature control must be used with fix npt/sphere}}
\sphinxAtStartPar
Self\sphinxhyphen{}explanatory.

\sphinxlineitem{\sphinxstyleemphasis{Temperature control must be used with fix nvt}}
\sphinxAtStartPar
Self\sphinxhyphen{}explanatory.

\sphinxlineitem{\sphinxstyleemphasis{Temperature control must be used with fix nvt/asphere}}
\sphinxAtStartPar
Self\sphinxhyphen{}explanatory.

\sphinxlineitem{\sphinxstyleemphasis{Temperature control must be used with fix nvt/body}}
\sphinxAtStartPar
Self\sphinxhyphen{}explanatory.

\sphinxlineitem{\sphinxstyleemphasis{Temperature control must be used with fix nvt/sllod}}
\sphinxAtStartPar
Self\sphinxhyphen{}explanatory.

\sphinxlineitem{\sphinxstyleemphasis{Temperature control must be used with fix nvt/sphere}}
\sphinxAtStartPar
Self\sphinxhyphen{}explanatory.

\sphinxlineitem{\sphinxstyleemphasis{Temperature control must not be used with fix nph/small}}
\sphinxAtStartPar
Self\sphinxhyphen{}explanatory.

\sphinxlineitem{\sphinxstyleemphasis{Temperature for fix nvt/sllod does not have a bias}}
\sphinxAtStartPar
The specified compute must compute temperature with a bias.

\sphinxlineitem{\sphinxstyleemphasis{Tempering could not find thermo\_pe compute}}
\sphinxAtStartPar
This compute is created by the thermo command.  It must have been
explicitly deleted by a uncompute command.

\sphinxlineitem{\sphinxstyleemphasis{Tempering fix ID is not defined}}
\sphinxAtStartPar
The fix ID specified by the temper command does not exist.

\sphinxlineitem{\sphinxstyleemphasis{Tempering temperature fix is not valid}}
\sphinxAtStartPar
The fix specified by the temper command is not one that controls
temperature (nvt or langevin).

\sphinxlineitem{\sphinxstyleemphasis{Test\_descriptor\_string already allocated}}
\sphinxAtStartPar
This is an internal error.  Contact the developers.

\sphinxlineitem{\sphinxstyleemphasis{The package gpu command is required for gpu styles}}
\sphinxAtStartPar
Self\sphinxhyphen{}explanatory.

\sphinxlineitem{\sphinxstyleemphasis{Thermo and fix not computed at compatible times}}
\sphinxAtStartPar
Fixes generate values on specific timesteps.  The thermo output
does not match these timesteps.

\sphinxlineitem{\sphinxstyleemphasis{Thermo compute array is accessed out\sphinxhyphen{}of\sphinxhyphen{}range}}
\sphinxAtStartPar
Self\sphinxhyphen{}explanatory.

\sphinxlineitem{\sphinxstyleemphasis{Thermo compute does not compute array}}
\sphinxAtStartPar
Self\sphinxhyphen{}explanatory.

\sphinxlineitem{\sphinxstyleemphasis{Thermo compute does not compute scalar}}
\sphinxAtStartPar
Self\sphinxhyphen{}explanatory.

\sphinxlineitem{\sphinxstyleemphasis{Thermo compute does not compute vector}}
\sphinxAtStartPar
Self\sphinxhyphen{}explanatory.

\sphinxlineitem{\sphinxstyleemphasis{Thermo compute vector is accessed out\sphinxhyphen{}of\sphinxhyphen{}range}}
\sphinxAtStartPar
Self\sphinxhyphen{}explanatory.

\sphinxlineitem{\sphinxstyleemphasis{Thermo custom variable cannot be indexed}}
\sphinxAtStartPar
Self\sphinxhyphen{}explanatory.

\sphinxlineitem{\sphinxstyleemphasis{Thermo custom variable is not equal\sphinxhyphen{}style variable}}
\sphinxAtStartPar
Only equal\sphinxhyphen{}style variables can be output with thermodynamics, not
atom\sphinxhyphen{}style variables.

\sphinxlineitem{\sphinxstyleemphasis{Thermo every variable returned a bad timestep}}
\sphinxAtStartPar
The variable must return a timestep greater than the current timestep.

\sphinxlineitem{\sphinxstyleemphasis{Thermo fix array is accessed out\sphinxhyphen{}of\sphinxhyphen{}range}}
\sphinxAtStartPar
Self\sphinxhyphen{}explanatory.

\sphinxlineitem{\sphinxstyleemphasis{Thermo fix does not compute array}}
\sphinxAtStartPar
Self\sphinxhyphen{}explanatory.

\sphinxlineitem{\sphinxstyleemphasis{Thermo fix does not compute scalar}}
\sphinxAtStartPar
Self\sphinxhyphen{}explanatory.

\sphinxlineitem{\sphinxstyleemphasis{Thermo fix does not compute vector}}
\sphinxAtStartPar
Self\sphinxhyphen{}explanatory.

\sphinxlineitem{\sphinxstyleemphasis{Thermo fix vector is accessed out\sphinxhyphen{}of\sphinxhyphen{}range}}
\sphinxAtStartPar
Self\sphinxhyphen{}explanatory.

\sphinxlineitem{\sphinxstyleemphasis{Thermo keyword in variable requires thermo to use/init pe}}
\sphinxAtStartPar
You are using a thermo keyword in a variable that requires
potential energy to be calculated, but your thermo output
does not use it.  Add it to your thermo output.

\sphinxlineitem{\sphinxstyleemphasis{Thermo keyword in variable requires thermo to use/init press}}
\sphinxAtStartPar
You are using a thermo keyword in a variable that requires pressure to
be calculated, but your thermo output does not use it.  Add it to your
thermo output.

\sphinxlineitem{\sphinxstyleemphasis{Thermo keyword in variable requires thermo to use/init temp}}
\sphinxAtStartPar
You are using a thermo keyword in a variable that requires temperature
to be calculated, but your thermo output does not use it.  Add it to
your thermo output.

\sphinxlineitem{\sphinxstyleemphasis{Thermo style does not use press}}
\sphinxAtStartPar
Cannot use thermo\_modify to set this parameter since the thermo\_style
is not computing this quantity.

\sphinxlineitem{\sphinxstyleemphasis{Thermo style does not use temp}}
\sphinxAtStartPar
Cannot use thermo\_modify to set this parameter since the thermo\_style
is not computing this quantity.

\sphinxlineitem{\sphinxstyleemphasis{Thermo\_modify every variable returned a bad timestep}}
\sphinxAtStartPar
The returned timestep is less than or equal to the current timestep.

\sphinxlineitem{\sphinxstyleemphasis{Thermo\_modify int format does not contain d character}}
\sphinxAtStartPar
Self\sphinxhyphen{}explanatory.

\sphinxlineitem{\sphinxstyleemphasis{Thermo\_modify pressure ID does not compute pressure}}
\sphinxAtStartPar
The specified compute ID does not compute pressure.

\sphinxlineitem{\sphinxstyleemphasis{Thermo\_modify temperature ID does not compute temperature}}
\sphinxAtStartPar
The specified compute ID does not compute temperature.

\sphinxlineitem{\sphinxstyleemphasis{Thermo\_style command before simulation box is defined}}
\sphinxAtStartPar
The thermo\_style command cannot be used before a read\_data,
read\_restart, or create\_box command.

\sphinxlineitem{\sphinxstyleemphasis{This variable thermo keyword cannot be used between runs}}
\sphinxAtStartPar
Keywords that refer to time (such as cpu, elapsed) do not
make sense in between runs.

\sphinxlineitem{\sphinxstyleemphasis{Threshold for an atom property that is not allocated}}
\sphinxAtStartPar
A dump threshold has been requested on a quantity that is
not defined by the atom style used in this simulation.

\sphinxlineitem{\sphinxstyleemphasis{Timestep must be \textgreater{}= 0}}
\sphinxAtStartPar
Specified timestep is invalid.

\sphinxlineitem{\sphinxstyleemphasis{Too big a problem to use velocity create loop all}}
\sphinxAtStartPar
The system size must fit in a 32\sphinxhyphen{}bit integer to use this option.

\sphinxlineitem{\sphinxstyleemphasis{Too big a timestep for dump dcd}}
\sphinxAtStartPar
The timestep must fit in a 32\sphinxhyphen{}bit integer to use this dump style.

\sphinxlineitem{\sphinxstyleemphasis{Too big a timestep for dump xtc}}
\sphinxAtStartPar
The timestep must fit in a 32\sphinxhyphen{}bit integer to use this dump style.

\sphinxlineitem{\sphinxstyleemphasis{Too few bits for lookup table}}
\sphinxAtStartPar
Table size specified via pair\_modify command does not work with your
machine’s floating point representation.

\sphinxlineitem{\sphinxstyleemphasis{Too few lines in \%s section of data file}}
\sphinxAtStartPar
Self\sphinxhyphen{}explanatory.

\sphinxlineitem{\sphinxstyleemphasis{Too few values in body lines in data file}}
\sphinxAtStartPar
Self\sphinxhyphen{}explanatory.

\sphinxlineitem{\sphinxstyleemphasis{Too few values in body section of molecule file}}
\sphinxAtStartPar
Self\sphinxhyphen{}explanatory.

\sphinxlineitem{\sphinxstyleemphasis{Too many \sphinxhyphen{}pk arguments in command line}}
\sphinxAtStartPar
The string formed by concatenating the arguments is too long.  Use a
package command in the input script instead.

\sphinxlineitem{\sphinxstyleemphasis{Too many MSM grid levels}}
\sphinxAtStartPar
The max number of MSM grid levels is hardwired to 10.

\sphinxlineitem{\sphinxstyleemphasis{Too many args in variable function}}
\sphinxAtStartPar
More args are used than any variable function allows.

\sphinxlineitem{\sphinxstyleemphasis{Too many atom pairs for pair bop}}
\sphinxAtStartPar
The number of atomic pairs exceeds the expected number.  Check your
atomic structure to ensure that it is realistic.

\sphinxlineitem{\sphinxstyleemphasis{Too many atom sorting bins}}
\sphinxAtStartPar
This is likely due to an immense simulation box that has blown up
to a large size.

\sphinxlineitem{\sphinxstyleemphasis{Too many atom triplets for pair bop}}
\sphinxAtStartPar
The number of three atom groups for angle determinations exceeds the
expected number.  Check your atomic structure to ensure that it is
realistic.

\sphinxlineitem{\sphinxstyleemphasis{Too many atoms for dump dcd}}
\sphinxAtStartPar
The system size must fit in a 32\sphinxhyphen{}bit integer to use this dump
style.

\sphinxlineitem{\sphinxstyleemphasis{Too many atoms for dump xtc}}
\sphinxAtStartPar
The system size must fit in a 32\sphinxhyphen{}bit integer to use this dump
style.

\sphinxlineitem{\sphinxstyleemphasis{Too many elements extracted from MEAM library.}}
\sphinxAtStartPar
Increase ‘maxelt’ in meam.h and recompile.

\sphinxlineitem{\sphinxstyleemphasis{Too many exponent bits for lookup table}}
\sphinxAtStartPar
Table size specified via pair\_modify command does not work with your
machine’s floating point representation.

\sphinxlineitem{\sphinxstyleemphasis{Too many groups}}
\sphinxAtStartPar
The maximum number of atom groups (including the “all” group) is
given by MAX\_GROUP in group.cpp and is 32.

\sphinxlineitem{\sphinxstyleemphasis{Too many iterations}}
\sphinxAtStartPar
You must use a number of iterations that fit in a 32\sphinxhyphen{}bit integer
for minimization.

\sphinxlineitem{\sphinxstyleemphasis{Too many lines in one body in data file \sphinxhyphen{} boost MAXBODY}}
\sphinxAtStartPar
MAXBODY is a setting at the top of the src/read\_data.cpp file.
Set it larger and re\sphinxhyphen{}compile the code.

\sphinxlineitem{\sphinxstyleemphasis{Too many local+ghost atoms for neighbor list}}
\sphinxAtStartPar
The number of nlocal + nghost atoms on a processor
is limited by the size of a 32\sphinxhyphen{}bit integer with 2 bits
removed for masking 1\sphinxhyphen{}2, 1\sphinxhyphen{}3, 1\sphinxhyphen{}4 neighbors.

\sphinxlineitem{\sphinxstyleemphasis{Too many mantissa bits for lookup table}}
\sphinxAtStartPar
Table size specified via pair\_modify command does not work with your
machine’s floating point representation.

\sphinxlineitem{\sphinxstyleemphasis{Too many masses for fix shake}}
\sphinxAtStartPar
The fix shake command cannot list more masses than there are atom
types.

\sphinxlineitem{\sphinxstyleemphasis{Too many molecules for fix poems}}
\sphinxAtStartPar
The limit is 2\textasciicircum{}31 = \textasciitilde{}2 billion molecules.

\sphinxlineitem{\sphinxstyleemphasis{Too many molecules for fix rigid}}
\sphinxAtStartPar
The limit is 2\textasciicircum{}31 = \textasciitilde{}2 billion molecules.

\sphinxlineitem{\sphinxstyleemphasis{Too many neighbor bins}}
\sphinxAtStartPar
This is likely due to an immense simulation box that has blown up
to a large size.

\sphinxlineitem{\sphinxstyleemphasis{Too many timesteps}}
\sphinxAtStartPar
The cumulative timesteps must fit in a 64\sphinxhyphen{}bit integer.

\sphinxlineitem{\sphinxstyleemphasis{Too many timesteps for NEB}}
\sphinxAtStartPar
You must use a number of timesteps that fit in a 32\sphinxhyphen{}bit integer
for NEB.

\sphinxlineitem{\sphinxstyleemphasis{Too many total atoms}}
\sphinxAtStartPar
See the setting for bigint in the src/lmptype.h file.

\sphinxlineitem{\sphinxstyleemphasis{Too many total bits for bitmapped lookup table}}
\sphinxAtStartPar
Table size specified via pair\_modify command is too large.  Note that
a value of N generates a 2\textasciicircum{}N size table.

\sphinxlineitem{\sphinxstyleemphasis{Too many values in body lines in data file}}
\sphinxAtStartPar
Self\sphinxhyphen{}explanatory.

\sphinxlineitem{\sphinxstyleemphasis{Too many values in body section of molecule file}}
\sphinxAtStartPar
Self\sphinxhyphen{}explanatory.

\sphinxlineitem{\sphinxstyleemphasis{Too much buffered per\sphinxhyphen{}proc info for dump}}
\sphinxAtStartPar
The size of the buffered string must fit in a 32\sphinxhyphen{}bit integer for a
dump.

\sphinxlineitem{\sphinxstyleemphasis{Too much per\sphinxhyphen{}proc info for dump}}
\sphinxAtStartPar
Number of local atoms times number of columns must fit in a 32\sphinxhyphen{}bit
integer for dump.

\sphinxlineitem{\sphinxstyleemphasis{Topology type exceeds system topology type}}
\sphinxAtStartPar
The number of bond, angle, etc types exceeds the system setting. See
the create\_box or read\_data command for how to specify these values.

\sphinxlineitem{\sphinxstyleemphasis{Tree structure in joint connections}}
\sphinxAtStartPar
Fix poems cannot (yet) work with coupled bodies whose joints connect
the bodies in a tree structure.

\sphinxlineitem{\sphinxstyleemphasis{Tried to convert a double to int, but input\_double \textgreater{} INT\_MAX}}
\sphinxAtStartPar
Self\sphinxhyphen{}explanatory.

\sphinxlineitem{\sphinxstyleemphasis{Trying to build an occasional neighbor list before initialization completed}}
\sphinxAtStartPar
This is not allowed.  Source code caller needs to be modified.

\sphinxlineitem{\sphinxstyleemphasis{Two fix ave commands using same compute chunk/atom command in incompatible ways}}
\sphinxAtStartPar
They are both attempting to “lock” the chunk/atom command so that the
chunk assignments persist for some number of timesteps, but are doing
it in different ways.

\sphinxlineitem{\sphinxstyleemphasis{Two groups cannot be the same in fix spring couple}}
\sphinxAtStartPar
Self\sphinxhyphen{}explanatory.

\sphinxlineitem{\sphinxstyleemphasis{The \%s type label \%s is already in use for type \%s}}
\sphinxAtStartPar
For a given type\sphinxhyphen{}kind (atom types, bond types, etc.), a given type
label can be assigned to only one numeric type.

\sphinxlineitem{\sphinxstyleemphasis{Type label string \%s for \%s type \%s is invalid}}
\sphinxAtStartPar
See the labelmap command documentation for valid type labels.

\sphinxlineitem{\sphinxstyleemphasis{Unable to initialize accelerator for use}}
\sphinxAtStartPar
There was a problem initializing an accelerator for the gpu package

\sphinxlineitem{\sphinxstyleemphasis{Unbalanced quotes in input line}}
\sphinxAtStartPar
No matching end double quote was found following a leading double
quote.

\sphinxlineitem{\sphinxstyleemphasis{Unexpected end of \sphinxhyphen{}reorder file}}
\sphinxAtStartPar
Self\sphinxhyphen{}explanatory.

\sphinxlineitem{\sphinxstyleemphasis{Unexpected empty line in Angle Coeffs section}}
\sphinxAtStartPar
Read a blank line where there should be coefficient data.

\sphinxlineitem{\sphinxstyleemphasis{Unexpected empty line in Bond Coeffs section}}
\sphinxAtStartPar
Read a blank line where there should be coefficient data.

\sphinxlineitem{\sphinxstyleemphasis{Unexpected empty line in Dihedral Coeffs section}}
\sphinxAtStartPar
Read a blank line where there should be coefficient data.

\sphinxlineitem{\sphinxstyleemphasis{Unexpected empty line in Improper Coeffs section}}
\sphinxAtStartPar
Read a blank line where there should be coefficient data.

\sphinxlineitem{\sphinxstyleemphasis{Unexpected empty line in Pair Coeffs section}}
\sphinxAtStartPar
Read a blank line where there should be coefficient data.

\sphinxlineitem{\sphinxstyleemphasis{Unexpected end of custom file}}
\sphinxAtStartPar
Self\sphinxhyphen{}explanatory.

\sphinxlineitem{\sphinxstyleemphasis{Unexpected end of data file}}
\sphinxAtStartPar
LAMMPS hit the end of the data file while attempting to read a
section.  Something is wrong with the format of the data file.

\sphinxlineitem{\sphinxstyleemphasis{Unexpected end of dump file}}
\sphinxAtStartPar
A read operation from the file failed.

\sphinxlineitem{\sphinxstyleemphasis{Unexpected end of fix rigid file}}
\sphinxAtStartPar
A read operation from the file failed.

\sphinxlineitem{\sphinxstyleemphasis{Unexpected end of fix rigid/small file}}
\sphinxAtStartPar
A read operation from the file failed.

\sphinxlineitem{\sphinxstyleemphasis{Unexpected end of molecule file}}
\sphinxAtStartPar
Self\sphinxhyphen{}explanatory.

\sphinxlineitem{\sphinxstyleemphasis{Unexpected end of neb file}}
\sphinxAtStartPar
A read operation from the file failed.

\sphinxlineitem{\sphinxstyleemphasis{Units command after simulation box is defined}}
\sphinxAtStartPar
The units command cannot be used after a read\_data, read\_restart, or
create\_box command.

\sphinxlineitem{\sphinxstyleemphasis{Universe/uloop variable count \textless{} \# of partitions}}
\sphinxAtStartPar
A universe or uloop style variable must specify a number of values \textgreater{}= to the
number of processor partitions.

\sphinxlineitem{\sphinxstyleemphasis{Unrecognized angle style}}
\sphinxAtStartPar
The choice of angle style is unknown.

\sphinxlineitem{\sphinxstyleemphasis{Unrecognized atom style}}
\sphinxAtStartPar
The choice of atom style is unknown.

\sphinxlineitem{\sphinxstyleemphasis{Unrecognized body style}}
\sphinxAtStartPar
The choice of body style is unknown.

\sphinxlineitem{\sphinxstyleemphasis{Unrecognized bond style}}
\sphinxAtStartPar
The choice of bond style is unknown.

\sphinxlineitem{\sphinxstyleemphasis{Unknown category for info is\_active()}}
\sphinxAtStartPar
Self\sphinxhyphen{}explanatory.

\sphinxlineitem{\sphinxstyleemphasis{Unknown category for info is\_available()}}
\sphinxAtStartPar
Self\sphinxhyphen{}explanatory.

\sphinxlineitem{\sphinxstyleemphasis{Unknown category for info is\_defined()}}
\sphinxAtStartPar
Self\sphinxhyphen{}explanatory.

\sphinxlineitem{\sphinxstyleemphasis{Unrecognized command: \%s}}
\sphinxAtStartPar
The command is not known to LAMMPS.  Check the input script.

\sphinxlineitem{\sphinxstyleemphasis{Unrecognized compute style}}
\sphinxAtStartPar
The choice of compute style is unknown.

\sphinxlineitem{\sphinxstyleemphasis{Unrecognized dihedral style}}
\sphinxAtStartPar
The choice of dihedral style is unknown.

\sphinxlineitem{\sphinxstyleemphasis{Unrecognized dump reader style}}
\sphinxAtStartPar
The choice of dump reader style via the format keyword is unknown.

\sphinxlineitem{\sphinxstyleemphasis{Unrecognized dump style}}
\sphinxAtStartPar
The choice of dump style is unknown.

\sphinxlineitem{\sphinxstyleemphasis{Unknown error in GPU library}}
\sphinxAtStartPar
Self\sphinxhyphen{}explanatory.

\sphinxlineitem{\sphinxstyleemphasis{Unrecognized fix style}}
\sphinxAtStartPar
The choice of fix style is unknown.

\sphinxlineitem{\sphinxstyleemphasis{Unknown identifier in data file: \%s}}
\sphinxAtStartPar
A section of the data file cannot be read by LAMMPS.

\sphinxlineitem{\sphinxstyleemphasis{Unrecognized improper style}}
\sphinxAtStartPar
The choice of improper style is unknown.

\sphinxlineitem{\sphinxstyleemphasis{Unknown keyword in thermo\_style custom command}}
\sphinxAtStartPar
One or more specified keywords are not recognized.

\sphinxlineitem{\sphinxstyleemphasis{Unrecognized kspace style}}
\sphinxAtStartPar
The choice of kspace style is unknown.

\sphinxlineitem{\sphinxstyleemphasis{Unknown name for info newton category}}
\sphinxAtStartPar
Self\sphinxhyphen{}explanatory.

\sphinxlineitem{\sphinxstyleemphasis{Unknown name for info package category}}
\sphinxAtStartPar
Self\sphinxhyphen{}explanatory.

\sphinxlineitem{\sphinxstyleemphasis{Unknown name for info pair category}}
\sphinxAtStartPar
Self\sphinxhyphen{}explanatory.

\sphinxlineitem{\sphinxstyleemphasis{Unrecognized pair style}}
\sphinxAtStartPar
The choice of pair style is unknown.

\sphinxlineitem{\sphinxstyleemphasis{Unknown pair\_modify hybrid sub\sphinxhyphen{}style}}
\sphinxAtStartPar
The choice of sub\sphinxhyphen{}style is unknown.

\sphinxlineitem{\sphinxstyleemphasis{Unrecognized region style}}
\sphinxAtStartPar
The choice of region style is unknown.

\sphinxlineitem{\sphinxstyleemphasis{Unknown section in molecule file}}
\sphinxAtStartPar
Self\sphinxhyphen{}explanatory.

\sphinxlineitem{\sphinxstyleemphasis{Unknown table style in angle style table}}
\sphinxAtStartPar
Self\sphinxhyphen{}explanatory.

\sphinxlineitem{\sphinxstyleemphasis{Unknown table style in bond style table}}
\sphinxAtStartPar
Self\sphinxhyphen{}explanatory.

\sphinxlineitem{\sphinxstyleemphasis{Unknown table style in pair\_style command}}
\sphinxAtStartPar
Style of table is invalid for use with pair\_style table command.

\sphinxlineitem{\sphinxstyleemphasis{Unknown unit\_style}}
\sphinxAtStartPar
Self\sphinxhyphen{}explanatory. Check the input script or data file.

\sphinxlineitem{\sphinxstyleemphasis{Unrecognized lattice type in MEAM library file}}
\sphinxAtStartPar
The lattice type in an entry of the MEAM library file is not
valid.

\sphinxlineitem{\sphinxstyleemphasis{Unrecognized lattice type in MEAM parameter file}}
\sphinxAtStartPar
The lattice type in an entry of the MEAM parameter file is not
valid.

\sphinxlineitem{\sphinxstyleemphasis{Unrecognized pair style in compute pair command}}
\sphinxAtStartPar
Self\sphinxhyphen{}explanatory.

\sphinxlineitem{\sphinxstyleemphasis{Unsupported mixing rule in kspace\_style ewald/disp}}
\sphinxAtStartPar
Only geometric mixing is supported.

\sphinxlineitem{\sphinxstyleemphasis{Unsupported order in kspace\_style ewald/disp}}
\sphinxAtStartPar
Only 1/r\textasciicircum{}6 dispersion or dipole terms are supported.

\sphinxlineitem{\sphinxstyleemphasis{Unsupported order in kspace\_style pppm/disp, pair\_style \%s}}
\sphinxAtStartPar
Only pair styles with 1/r and 1/r\textasciicircum{}6 dependence are currently supported.

\sphinxlineitem{\sphinxstyleemphasis{Unsupported parameter in MEAM library file}}
\sphinxAtStartPar
Self\sphinxhyphen{}explanatory.

\sphinxlineitem{\sphinxstyleemphasis{Use cutoff keyword to set cutoff in single mode}}
\sphinxAtStartPar
Mode is single so cutoff/multi keyword cannot be used.

\sphinxlineitem{\sphinxstyleemphasis{Use cutoff/multi keyword to set cutoff in multi mode}}
\sphinxAtStartPar
Mode is multi so cutoff keyword cannot be used.

\sphinxlineitem{\sphinxstyleemphasis{Using fix nvt/sllod with inconsistent fix deform remap option}}
\sphinxAtStartPar
Fix nvt/sllod requires that deforming atoms have a velocity profile
provided by “remap v” as a fix deform option.

\sphinxlineitem{\sphinxstyleemphasis{Using fix nvt/sllod with no fix deform defined}}
\sphinxAtStartPar
Self\sphinxhyphen{}explanatory.

\sphinxlineitem{\sphinxstyleemphasis{Using fix srd with inconsistent fix deform remap option}}
\sphinxAtStartPar
When shearing the box in an SRD simulation, the remap v option for fix
deform needs to be used.

\sphinxlineitem{\sphinxstyleemphasis{Using pair lubricate with inconsistent fix deform remap option}}
\sphinxAtStartPar
Must use remap v option with fix deform with this pair style.

\sphinxlineitem{\sphinxstyleemphasis{Using pair lubricate/poly with inconsistent fix deform remap option}}
\sphinxAtStartPar
If fix deform is used, the remap v option is required.

\sphinxlineitem{\sphinxstyleemphasis{Using suffix gpu without GPU package installed}}
\sphinxAtStartPar
Self\sphinxhyphen{}explanatory.

\sphinxlineitem{\sphinxstyleemphasis{Using suffix intel without INTEL package installed}}
\sphinxAtStartPar
Self\sphinxhyphen{}explanatory.

\sphinxlineitem{\sphinxstyleemphasis{Using suffix kk without KOKKOS package enabled}}
\sphinxAtStartPar
Self\sphinxhyphen{}explanatory.

\sphinxlineitem{\sphinxstyleemphasis{Using suffix omp without OPENMP package installed}}
\sphinxAtStartPar
Self\sphinxhyphen{}explanatory.

\sphinxlineitem{\sphinxstyleemphasis{Using update dipole flag requires atom attribute mu}}
\sphinxAtStartPar
Self\sphinxhyphen{}explanatory.

\sphinxlineitem{\sphinxstyleemphasis{Using update dipole flag requires atom style sphere}}
\sphinxAtStartPar
Self\sphinxhyphen{}explanatory.

\sphinxlineitem{\sphinxstyleemphasis{Variable ID in variable formula does not exist}}
\sphinxAtStartPar
Self\sphinxhyphen{}explanatory.

\sphinxlineitem{\sphinxstyleemphasis{Variable atom ID is too large}}
\sphinxAtStartPar
Specified ID is larger than the maximum allowed atom ID.

\sphinxlineitem{\sphinxstyleemphasis{Variable evaluation before simulation box is defined}}
\sphinxAtStartPar
Cannot evaluate a compute or fix or atom\sphinxhyphen{}based value in a variable
before the simulation has been setup.

\sphinxlineitem{\sphinxstyleemphasis{Variable evaluation in fix wall gave bad value}}
\sphinxAtStartPar
The returned value for epsilon or sigma \textless{} 0.0.

\sphinxlineitem{\sphinxstyleemphasis{Variable evaluation in region gave bad value}}
\sphinxAtStartPar
Variable returned a radius \textless{} 0.0.

\sphinxlineitem{\sphinxstyleemphasis{Variable for compute ti is invalid style}}
\sphinxAtStartPar
Self\sphinxhyphen{}explanatory.

\sphinxlineitem{\sphinxstyleemphasis{Variable for create\_atoms is invalid style}}
\sphinxAtStartPar
The variables must be equal\sphinxhyphen{}style variables.

\sphinxlineitem{\sphinxstyleemphasis{Variable for displace\_atoms is invalid style}}
\sphinxAtStartPar
It must be an equal\sphinxhyphen{}style or atom\sphinxhyphen{}style variable.

\sphinxlineitem{\sphinxstyleemphasis{Variable for dump every is invalid style}}
\sphinxAtStartPar
Only equal\sphinxhyphen{}style variables can be used.

\sphinxlineitem{\sphinxstyleemphasis{Variable for dump image center is invalid style}}
\sphinxAtStartPar
Must be an equal\sphinxhyphen{}style variable.

\sphinxlineitem{\sphinxstyleemphasis{Variable for dump image phi is invalid style}}
\sphinxAtStartPar
Must be an equal\sphinxhyphen{}style variable.

\sphinxlineitem{\sphinxstyleemphasis{Variable for dump image theta is invalid style}}
\sphinxAtStartPar
Must be an equal\sphinxhyphen{}style variable.

\sphinxlineitem{\sphinxstyleemphasis{Variable for dump image zoom is invalid style}}
\sphinxAtStartPar
Must be an equal\sphinxhyphen{}style variable.

\sphinxlineitem{\sphinxstyleemphasis{Variable for fix adapt is invalid style}}
\sphinxAtStartPar
Only equal\sphinxhyphen{}style variables can be used.

\sphinxlineitem{\sphinxstyleemphasis{Variable for fix addforce is invalid style}}
\sphinxAtStartPar
Self\sphinxhyphen{}explanatory.

\sphinxlineitem{\sphinxstyleemphasis{Variable for fix aveforce is invalid style}}
\sphinxAtStartPar
Only equal\sphinxhyphen{}style variables can be used.

\sphinxlineitem{\sphinxstyleemphasis{Variable for fix deform is invalid style}}
\sphinxAtStartPar
The variable must be an equal\sphinxhyphen{}style variable.

\sphinxlineitem{\sphinxstyleemphasis{Variable for fix efield is invalid style}}
\sphinxAtStartPar
The variable must be an equal\sphinxhyphen{} or atom\sphinxhyphen{}style variable.

\sphinxlineitem{\sphinxstyleemphasis{Variable for fix gravity is invalid style}}
\sphinxAtStartPar
Only equal\sphinxhyphen{}style variables can be used.

\sphinxlineitem{\sphinxstyleemphasis{Variable for fix heat is invalid style}}
\sphinxAtStartPar
Only equal\sphinxhyphen{}style or atom\sphinxhyphen{}style variables can be used.

\sphinxlineitem{\sphinxstyleemphasis{Variable for fix indent is invalid style}}
\sphinxAtStartPar
Only equal\sphinxhyphen{}style variables can be used.

\sphinxlineitem{\sphinxstyleemphasis{Variable for fix indent is not equal style}}
\sphinxAtStartPar
Only equal\sphinxhyphen{}style variables can be used.

\sphinxlineitem{\sphinxstyleemphasis{Variable for fix langevin is invalid style}}
\sphinxAtStartPar
It must be an equal\sphinxhyphen{}style variable.

\sphinxlineitem{\sphinxstyleemphasis{Variable for fix move is invalid style}}
\sphinxAtStartPar
Only equal\sphinxhyphen{}style variables can be used.

\sphinxlineitem{\sphinxstyleemphasis{Variable for fix setforce is invalid style}}
\sphinxAtStartPar
Only equal\sphinxhyphen{}style variables can be used.

\sphinxlineitem{\sphinxstyleemphasis{Variable for fix temp/berendsen is invalid style}}
\sphinxAtStartPar
Only equal\sphinxhyphen{}style variables can be used.

\sphinxlineitem{\sphinxstyleemphasis{Variable for fix temp/csld is invalid style}}
\sphinxAtStartPar
Only equal\sphinxhyphen{}style variables can be used.

\sphinxlineitem{\sphinxstyleemphasis{Variable for fix temp/csvr is invalid style}}
\sphinxAtStartPar
Only equal\sphinxhyphen{}style variables can be used.

\sphinxlineitem{\sphinxstyleemphasis{Variable for fix temp/rescale is invalid style}}
\sphinxAtStartPar
Only equal\sphinxhyphen{}style variables can be used.

\sphinxlineitem{\sphinxstyleemphasis{Variable for fix wall is invalid style}}
\sphinxAtStartPar
Only equal\sphinxhyphen{}style variables can be used.

\sphinxlineitem{\sphinxstyleemphasis{Variable for fix wall/reflect is invalid style}}
\sphinxAtStartPar
Only equal\sphinxhyphen{}style variables can be used.

\sphinxlineitem{\sphinxstyleemphasis{Variable for fix wall/srd is invalid style}}
\sphinxAtStartPar
Only equal\sphinxhyphen{}style variables can be used.

\sphinxlineitem{\sphinxstyleemphasis{Variable for group dynamic is invalid style}}
\sphinxAtStartPar
The variable must be an atom\sphinxhyphen{}style variable.

\sphinxlineitem{\sphinxstyleemphasis{Variable for group is invalid style}}
\sphinxAtStartPar
Only atom\sphinxhyphen{}style variables can be used.

\sphinxlineitem{\sphinxstyleemphasis{Variable for region cylinder is invalid style}}
\sphinxAtStartPar
Only equal\sphinxhyphen{}style variables are allowed.

\sphinxlineitem{\sphinxstyleemphasis{Variable for region is invalid style}}
\sphinxAtStartPar
Only equal\sphinxhyphen{}style variables can be used.

\sphinxlineitem{\sphinxstyleemphasis{Variable for region is not equal style}}
\sphinxAtStartPar
Self\sphinxhyphen{}explanatory.

\sphinxlineitem{\sphinxstyleemphasis{Variable for region sphere is invalid style}}
\sphinxAtStartPar
Only equal\sphinxhyphen{}style variables are allowed.

\sphinxlineitem{\sphinxstyleemphasis{Variable for restart is invalid style}}
\sphinxAtStartPar
Only equal\sphinxhyphen{}style variables can be used.

\sphinxlineitem{\sphinxstyleemphasis{Variable for set command is invalid style}}
\sphinxAtStartPar
Only atom\sphinxhyphen{}style variables can be used.

\sphinxlineitem{\sphinxstyleemphasis{Variable for thermo every is invalid style}}
\sphinxAtStartPar
Only equal\sphinxhyphen{}style variables can be used.

\sphinxlineitem{\sphinxstyleemphasis{Variable for velocity set is invalid style}}
\sphinxAtStartPar
Only atom\sphinxhyphen{}style variables can be used.

\sphinxlineitem{\sphinxstyleemphasis{Variable for voronoi radius is not atom style}}
\sphinxAtStartPar
Self\sphinxhyphen{}explanatory.

\sphinxlineitem{\sphinxstyleemphasis{Variable formula compute array is accessed out\sphinxhyphen{}of\sphinxhyphen{}range}}
\sphinxAtStartPar
Self\sphinxhyphen{}explanatory.

\sphinxlineitem{\sphinxstyleemphasis{Variable formula compute vector is accessed out\sphinxhyphen{}of\sphinxhyphen{}range}}
\sphinxAtStartPar
Self\sphinxhyphen{}explanatory.

\sphinxlineitem{\sphinxstyleemphasis{Variable formula fix array is accessed out\sphinxhyphen{}of\sphinxhyphen{}range}}
\sphinxAtStartPar
Self\sphinxhyphen{}explanatory.

\sphinxlineitem{\sphinxstyleemphasis{Variable formula fix vector is accessed out\sphinxhyphen{}of\sphinxhyphen{}range}}
\sphinxAtStartPar
Self\sphinxhyphen{}explanatory.

\sphinxlineitem{\sphinxstyleemphasis{Variable has circular dependency}}
\sphinxAtStartPar
A circular dependency is when variable “a” in used by variable “b” and
variable “b” is also used by variable “a”.  Circular dependencies with
longer chains of dependence are also not allowed.

\sphinxlineitem{\sphinxstyleemphasis{Variable name between brackets must be alphanumeric or underscore characters}}
\sphinxAtStartPar
Self\sphinxhyphen{}explanatory.

\sphinxlineitem{\sphinxstyleemphasis{Variable name for compute chunk/atom does not exist}}
\sphinxAtStartPar
Self\sphinxhyphen{}explanatory.

\sphinxlineitem{\sphinxstyleemphasis{Variable name for compute reduce does not exist}}
\sphinxAtStartPar
Self\sphinxhyphen{}explanatory.

\sphinxlineitem{\sphinxstyleemphasis{Variable name for compute ti does not exist}}
\sphinxAtStartPar
Self\sphinxhyphen{}explanatory.

\sphinxlineitem{\sphinxstyleemphasis{Variable name for create\_atoms does not exist}}
\sphinxAtStartPar
Self\sphinxhyphen{}explanatory.

\sphinxlineitem{\sphinxstyleemphasis{Variable name for displace\_atoms does not exist}}
\sphinxAtStartPar
Self\sphinxhyphen{}explanatory.

\sphinxlineitem{\sphinxstyleemphasis{Variable name for dump every does not exist}}
\sphinxAtStartPar
Self\sphinxhyphen{}explanatory.

\sphinxlineitem{\sphinxstyleemphasis{Variable name for dump image center does not exist}}
\sphinxAtStartPar
Self\sphinxhyphen{}explanatory.

\sphinxlineitem{\sphinxstyleemphasis{Variable name for dump image phi does not exist}}
\sphinxAtStartPar
Self\sphinxhyphen{}explanatory.

\sphinxlineitem{\sphinxstyleemphasis{Variable name for dump image theta does not exist}}
\sphinxAtStartPar
Self\sphinxhyphen{}explanatory.

\sphinxlineitem{\sphinxstyleemphasis{Variable name for dump image zoom does not exist}}
\sphinxAtStartPar
Self\sphinxhyphen{}explanatory.

\sphinxlineitem{\sphinxstyleemphasis{Variable name for fix adapt does not exist}}
\sphinxAtStartPar
Self\sphinxhyphen{}explanatory.

\sphinxlineitem{\sphinxstyleemphasis{Variable name for fix addforce does not exist}}
\sphinxAtStartPar
Self\sphinxhyphen{}explanatory.

\sphinxlineitem{\sphinxstyleemphasis{Variable name for fix ave/atom does not exist}}
\sphinxAtStartPar
Self\sphinxhyphen{}explanatory.

\sphinxlineitem{\sphinxstyleemphasis{Variable name for fix ave/chunk does not exist}}
\sphinxAtStartPar
Self\sphinxhyphen{}explanatory.

\sphinxlineitem{\sphinxstyleemphasis{Variable name for fix ave/correlate does not exist}}
\sphinxAtStartPar
Self\sphinxhyphen{}explanatory.

\sphinxlineitem{\sphinxstyleemphasis{Variable name for fix ave/histo does not exist}}
\sphinxAtStartPar
Self\sphinxhyphen{}explanatory.

\sphinxlineitem{\sphinxstyleemphasis{Variable name for fix ave/spatial does not exist}}
\sphinxAtStartPar
Self\sphinxhyphen{}explanatory.

\sphinxlineitem{\sphinxstyleemphasis{Variable name for fix ave/time does not exist}}
\sphinxAtStartPar
Self\sphinxhyphen{}explanatory.

\sphinxlineitem{\sphinxstyleemphasis{Variable name for fix aveforce does not exist}}
\sphinxAtStartPar
Self\sphinxhyphen{}explanatory.

\sphinxlineitem{\sphinxstyleemphasis{Variable name for fix deform does not exist}}
\sphinxAtStartPar
Self\sphinxhyphen{}explanatory.

\sphinxlineitem{\sphinxstyleemphasis{Variable name for fix efield does not exist}}
\sphinxAtStartPar
Self\sphinxhyphen{}explanatory.

\sphinxlineitem{\sphinxstyleemphasis{Variable name for fix gravity does not exist}}
\sphinxAtStartPar
Self\sphinxhyphen{}explanatory.

\sphinxlineitem{\sphinxstyleemphasis{Variable name for fix heat does not exist}}
\sphinxAtStartPar
Self\sphinxhyphen{}explanatory.

\sphinxlineitem{\sphinxstyleemphasis{Variable name for fix indent does not exist}}
\sphinxAtStartPar
Self\sphinxhyphen{}explanatory.

\sphinxlineitem{\sphinxstyleemphasis{Variable name for fix langevin does not exist}}
\sphinxAtStartPar
Self\sphinxhyphen{}explanatory.

\sphinxlineitem{\sphinxstyleemphasis{Variable name for fix move does not exist}}
\sphinxAtStartPar
Self\sphinxhyphen{}explanatory.

\sphinxlineitem{\sphinxstyleemphasis{Variable name for fix setforce does not exist}}
\sphinxAtStartPar
Self\sphinxhyphen{}explanatory.

\sphinxlineitem{\sphinxstyleemphasis{Variable name for fix store/state does not exist}}
\sphinxAtStartPar
Self\sphinxhyphen{}explanatory.

\sphinxlineitem{\sphinxstyleemphasis{Variable name for fix temp/berendsen does not exist}}
\sphinxAtStartPar
Self\sphinxhyphen{}explanatory.

\sphinxlineitem{\sphinxstyleemphasis{Variable name for fix temp/csld does not exist}}
\sphinxAtStartPar
Self\sphinxhyphen{}explanatory.

\sphinxlineitem{\sphinxstyleemphasis{Variable name for fix temp/csvr does not exist}}
\sphinxAtStartPar
Self\sphinxhyphen{}explanatory.

\sphinxlineitem{\sphinxstyleemphasis{Variable name for fix temp/rescale does not exist}}
\sphinxAtStartPar
Self\sphinxhyphen{}explanatory.

\sphinxlineitem{\sphinxstyleemphasis{Variable name for fix vector does not exist}}
\sphinxAtStartPar
Self\sphinxhyphen{}explanatory.

\sphinxlineitem{\sphinxstyleemphasis{Variable name for fix wall does not exist}}
\sphinxAtStartPar
Self\sphinxhyphen{}explanatory.

\sphinxlineitem{\sphinxstyleemphasis{Variable name for fix wall/reflect does not exist}}
\sphinxAtStartPar
Self\sphinxhyphen{}explanatory.

\sphinxlineitem{\sphinxstyleemphasis{Variable name for fix wall/srd does not exist}}
\sphinxAtStartPar
Self\sphinxhyphen{}explanatory.

\sphinxlineitem{\sphinxstyleemphasis{Variable name for group does not exist}}
\sphinxAtStartPar
Self\sphinxhyphen{}explanatory.

\sphinxlineitem{\sphinxstyleemphasis{Variable name for group dynamic does not exist}}
\sphinxAtStartPar
Self\sphinxhyphen{}explanatory.

\sphinxlineitem{\sphinxstyleemphasis{Variable name for region cylinder does not exist}}
\sphinxAtStartPar
Self\sphinxhyphen{}explanatory.

\sphinxlineitem{\sphinxstyleemphasis{Variable name for region does not exist}}
\sphinxAtStartPar
Self\sphinxhyphen{}explanatory.

\sphinxlineitem{\sphinxstyleemphasis{Variable name for region sphere does not exist}}
\sphinxAtStartPar
Self\sphinxhyphen{}explanatory.

\sphinxlineitem{\sphinxstyleemphasis{Variable name for restart does not exist}}
\sphinxAtStartPar
Self\sphinxhyphen{}explanatory.

\sphinxlineitem{\sphinxstyleemphasis{Variable name for set command does not exist}}
\sphinxAtStartPar
Self\sphinxhyphen{}explanatory.

\sphinxlineitem{\sphinxstyleemphasis{Variable name for thermo every does not exist}}
\sphinxAtStartPar
Self\sphinxhyphen{}explanatory.

\sphinxlineitem{\sphinxstyleemphasis{Variable name for velocity set does not exist}}
\sphinxAtStartPar
Self\sphinxhyphen{}explanatory.

\sphinxlineitem{\sphinxstyleemphasis{Variable name for voronoi radius does not exist}}
\sphinxAtStartPar
Self\sphinxhyphen{}explanatory.

\sphinxlineitem{\sphinxstyleemphasis{Variable name must be alphanumeric or underscore characters}}
\sphinxAtStartPar
Self\sphinxhyphen{}explanatory.

\sphinxlineitem{\sphinxstyleemphasis{Variable uses atom property that is not allocated}}
\sphinxAtStartPar
Self\sphinxhyphen{}explanatory.

\sphinxlineitem{\sphinxstyleemphasis{Velocity command before simulation box is defined}}
\sphinxAtStartPar
The velocity command cannot be used before a read\_data, read\_restart,
or create\_box command.

\sphinxlineitem{\sphinxstyleemphasis{Velocity command with no atoms existing}}
\sphinxAtStartPar
A velocity command has been used, but no atoms yet exist.

\sphinxlineitem{\sphinxstyleemphasis{Velocity ramp in z for a 2d problem}}
\sphinxAtStartPar
Self\sphinxhyphen{}explanatory.

\sphinxlineitem{\sphinxstyleemphasis{Velocity rigid used with non\sphinxhyphen{}rigid fix\sphinxhyphen{}ID}}
\sphinxAtStartPar
Self\sphinxhyphen{}explanatory.

\sphinxlineitem{\sphinxstyleemphasis{Velocity temperature ID does calculate a velocity bias}}
\sphinxAtStartPar
The specified compute must compute a bias for temperature.

\sphinxlineitem{\sphinxstyleemphasis{Velocity temperature ID does not compute temperature}}
\sphinxAtStartPar
The compute ID given to the velocity command must compute
temperature.

\sphinxlineitem{\sphinxstyleemphasis{Verlet/split can only currently be used with comm\_style brick}}
\sphinxAtStartPar
This is a current restriction in LAMMPS.

\sphinxlineitem{\sphinxstyleemphasis{Verlet/split does not yet support TIP4P}}
\sphinxAtStartPar
This is a current limitation.

\sphinxlineitem{\sphinxstyleemphasis{Verlet/split requires 2 partitions}}
\sphinxAtStartPar
See the \sphinxhyphen{}partition command\sphinxhyphen{}line switch.

\sphinxlineitem{\sphinxstyleemphasis{Verlet/split requires Rspace partition layout be multiple of Kspace partition layout in each dim}}
\sphinxAtStartPar
This is controlled by the processors command.

\sphinxlineitem{\sphinxstyleemphasis{Verlet/split requires Rspace partition size be multiple of Kspace partition size}}
\sphinxAtStartPar
This is so there is an equal number of Rspace processors for every
Kspace processor.

\sphinxlineitem{\sphinxstyleemphasis{Virial was not tallied on needed timestep}}
\sphinxAtStartPar
You are using a thermo keyword that requires potentials to
have tallied the virial, but they did not on this timestep.  See the
variable page for ideas on how to make this work.

\sphinxlineitem{\sphinxstyleemphasis{Voro++ error: narea and neigh have a different size}}
\sphinxAtStartPar
This error is returned by the Voro++ library.

\sphinxlineitem{\sphinxstyleemphasis{Wall defined twice in fix wall command}}
\sphinxAtStartPar
Self\sphinxhyphen{}explanatory.

\sphinxlineitem{\sphinxstyleemphasis{Wall defined twice in fix wall/reflect command}}
\sphinxAtStartPar
Self\sphinxhyphen{}explanatory.

\sphinxlineitem{\sphinxstyleemphasis{Wall defined twice in fix wall/srd command}}
\sphinxAtStartPar
Self\sphinxhyphen{}explanatory.

\sphinxlineitem{\sphinxstyleemphasis{Water H epsilon must be 0.0 for pair style lj/cut/tip4p/cut}}
\sphinxAtStartPar
This is because LAMMPS does not compute the Lennard\sphinxhyphen{}Jones interactions
with these particles for efficiency reasons.

\sphinxlineitem{\sphinxstyleemphasis{Water H epsilon must be 0.0 for pair style lj/cut/tip4p/long}}
\sphinxAtStartPar
This is because LAMMPS does not compute the Lennard\sphinxhyphen{}Jones interactions
with these particles for efficiency reasons.

\sphinxlineitem{\sphinxstyleemphasis{Water H epsilon must be 0.0 for pair style lj/long/tip4p/long}}
\sphinxAtStartPar
This is because LAMMPS does not compute the Lennard\sphinxhyphen{}Jones interactions
with these particles for efficiency reasons.

\sphinxlineitem{\sphinxstyleemphasis{World variable count does not match \# of partitions}}
\sphinxAtStartPar
A world\sphinxhyphen{}style variable must specify a number of values equal to the
number of processor partitions.

\sphinxlineitem{\sphinxstyleemphasis{Write\_data command before simulation box is defined}}
\sphinxAtStartPar
Self\sphinxhyphen{}explanatory.

\sphinxlineitem{\sphinxstyleemphasis{Write\_restart command before simulation box is defined}}
\sphinxAtStartPar
The write\_restart command cannot be used before a read\_data,
read\_restart, or create\_box command.

\sphinxlineitem{\sphinxstyleemphasis{Zero length rotation vector with displace\_atoms}}
\sphinxAtStartPar
Self\sphinxhyphen{}explanatory.

\sphinxlineitem{\sphinxstyleemphasis{Zero length rotation vector with fix move}}
\sphinxAtStartPar
Self\sphinxhyphen{}explanatory.

\sphinxlineitem{\sphinxstyleemphasis{Zero\sphinxhyphen{}length lattice orient vector}}
\sphinxAtStartPar
Self\sphinxhyphen{}explanatory.

\end{description}

\sphinxstepscope


\section{Warning messages}
\label{\detokenize{Errors_warnings:warning-messages}}\label{\detokenize{Errors_warnings::doc}}
\sphinxAtStartPar
This is an alphabetic list of the WARNING messages LAMMPS prints out
and the reason why.  If the explanation here is not sufficient, the
documentation for the offending command may help.  Warning messages
also list the source file and line number where the warning was
generated.  For example, a message like this:

\begin{sphinxVerbatim}[commandchars=\\\{\}]
WARNING: Bond atom missing in box size check (domain.cpp:187)
\end{sphinxVerbatim}

\sphinxAtStartPar
means that line \#187 in the file src/domain.cpp generated the error.
Looking in the source code may help you figure out what went wrong.

\sphinxAtStartPar
Doc page with {\hyperref[\detokenize{Errors_messages::doc}]{\sphinxcrossref{\DUrole{doc}{ERROR messages}}}}


\bigskip\hrule\bigskip

\begin{description}
\sphinxlineitem{\sphinxstyleemphasis{Adjusting Coulombic cutoff for MSM, new cutoff = \%g}}
\sphinxAtStartPar
The adjust/cutoff command is turned on and the Coulombic cutoff has been
adjusted to match the user\sphinxhyphen{}specified accuracy.

\sphinxlineitem{\sphinxstyleemphasis{Angle atoms missing at step \%ld}}
\sphinxAtStartPar
One or more of three atoms needed to compute a particular angle are
missing on this processor.  Typically this is because the pairwise
cutoff is set too short or the angle has blown apart and an atom is
too far away.

\sphinxlineitem{\sphinxstyleemphasis{Angle style in data file differs from currently defined angle style}}
\sphinxAtStartPar
Self\sphinxhyphen{}explanatory.

\sphinxlineitem{\sphinxstyleemphasis{Angles are defined but no angle style is set}}
\sphinxAtStartPar
The topology contains angles, but there are no angle forces computed
since there was no angle\_style command.

\sphinxlineitem{\sphinxstyleemphasis{Atom style in data file differs from currently defined atom style}}
\sphinxAtStartPar
Self\sphinxhyphen{}explanatory.

\sphinxlineitem{\sphinxstyleemphasis{Bond atom missing in box size check}}
\sphinxAtStartPar
The second atom needed to compute a particular bond is missing on this
processor.  Typically this is because the pairwise cutoff is set too
short or the bond has blown apart and an atom is too far away.

\sphinxlineitem{\sphinxstyleemphasis{Bond atom missing in image check}}
\sphinxAtStartPar
The second atom in a particular bond is missing on this processor.
Typically this is because the pairwise cutoff is set too short or the
bond has blown apart and an atom is too far away.

\sphinxlineitem{\sphinxstyleemphasis{Bond atoms missing at step \%ld}}
\sphinxAtStartPar
The second atom needed to compute a particular bond is missing on this
processor.  Typically this is because the pairwise cutoff is set too
short or the bond has blown apart and an atom is too far away.

\sphinxlineitem{\sphinxstyleemphasis{Bond style in data file differs from currently defined bond style}}
\sphinxAtStartPar
Self\sphinxhyphen{}explanatory.

\sphinxlineitem{\sphinxstyleemphasis{Bonds are defined but no bond style is set}}
\sphinxAtStartPar
The topology contains bonds, but there are no bond forces computed
since there was no bond\_style command.

\sphinxlineitem{\sphinxstyleemphasis{Bond/angle/dihedral extent \textgreater{} half of periodic box length}}
\sphinxAtStartPar
This is a restriction because LAMMPS can be confused about which image
of an atom in the bonded interaction is the correct one to use.
“Extent” in this context means the maximum end\sphinxhyphen{}to\sphinxhyphen{}end length of the
bond/angle/dihedral.  LAMMPS computes this by taking the maximum bond
length, multiplying by the number of bonds in the interaction (e.g. 3
for a dihedral) and adding a small amount of stretch.

\sphinxlineitem{\sphinxstyleemphasis{Both groups in compute group/group have a net charge; the Kspace boundary correction to energy will be non\sphinxhyphen{}zero}}
\sphinxAtStartPar
Self\sphinxhyphen{}explanatory.

\sphinxlineitem{\sphinxstyleemphasis{Calling write\_dump before a full system init.}}
\sphinxAtStartPar
The write\_dump command is used before the system has been fully
initialized as part of a ‘run’ or ‘minimize’ command. Not all dump
styles and features are fully supported at this point and thus the
command may fail or produce incomplete or incorrect output. Insert
a “run 0” command, if a full system init is required.

\sphinxlineitem{\sphinxstyleemphasis{Cannot count rigid body degrees\sphinxhyphen{}of\sphinxhyphen{}freedom before bodies are fully initialized}}
\sphinxAtStartPar
This means the temperature associated with the rigid bodies may be
incorrect on this timestep.

\sphinxlineitem{\sphinxstyleemphasis{Cannot count rigid body degrees\sphinxhyphen{}of\sphinxhyphen{}freedom before bodies are initialized}}
\sphinxAtStartPar
This means the temperature associated with the rigid bodies may be
incorrect on this timestep.

\sphinxlineitem{\sphinxstyleemphasis{Cannot include log terms without 1/r terms; setting flagHI to 1}}
\sphinxAtStartPar
Self\sphinxhyphen{}explanatory.

\sphinxlineitem{\sphinxstyleemphasis{Cannot include log terms without 1/r terms; setting flagHI to 1.}}
\sphinxAtStartPar
Self\sphinxhyphen{}explanatory.

\sphinxlineitem{\sphinxstyleemphasis{Charges are set, but coulombic solver is not used}}
\sphinxAtStartPar
Self\sphinxhyphen{}explanatory.

\sphinxlineitem{\sphinxstyleemphasis{Charges did not converge at step \%ld: \%lg}}
\sphinxAtStartPar
Self\sphinxhyphen{}explanatory.

\sphinxlineitem{\sphinxstyleemphasis{Communication cutoff is 0.0. No ghost atoms will be generated. Atoms may get lost}}
\sphinxAtStartPar
The communication cutoff defaults to the maximum of what is inferred from
pair and bond styles (will be zero, if none are defined) and what is specified
via {\hyperref[\detokenize{comm_modify::doc}]{\sphinxcrossref{\DUrole{doc}{comm\_modify cutoff}}}} (defaults to 0.0).  If this results
to 0.0, no ghost atoms will be generated and LAMMPS may lose atoms or use
incorrect periodic images of atoms in interaction lists.  To avoid, either use
{\hyperref[\detokenize{pair_zero::doc}]{\sphinxcrossref{\DUrole{doc}{pair style zero}}}} with a suitable cutoff or use {\hyperref[\detokenize{comm_modify::doc}]{\sphinxcrossref{\DUrole{doc}{comm\_modify cutoff}}}}.

\sphinxlineitem{\sphinxstyleemphasis{Communication cutoff is shorter than a bond length based estimate. This may lead to errors.}}
\sphinxAtStartPar
Since LAMMPS stores topology data with individual atoms, all atoms
comprising a bond, angle, dihedral or improper must be present on any
subdomain that “owns” the atom with the information, either as a
local or a ghost atom. The communication cutoff is what determines up
to what distance from a subdomain boundary ghost atoms are created.
The communication cutoff is by default the largest non\sphinxhyphen{}bonded cutoff
plus the neighbor skin distance, but for short or non\sphinxhyphen{}bonded cutoffs
and/or long bonds, this may not be sufficient. This warning indicates
that there is an increased risk of a simulation stopping unexpectedly
because of Bond/Angle/Dihedral/Improper atoms missing.  It can be
silenced by manually setting the communication cutoff via
{\hyperref[\detokenize{comm_modify::doc}]{\sphinxcrossref{\DUrole{doc}{comm\_modify cutoff}}}}.  However, since the
heuristic used to determine the estimate is not always accurate, it
is not changed automatically and the warning may be ignored depending
on the specific system being simulated.

\sphinxlineitem{\sphinxstyleemphasis{Communication cutoff is too small for SNAP micro load balancing, increased to \%lf}}
\sphinxAtStartPar
Self\sphinxhyphen{}explanatory.

\sphinxlineitem{\sphinxstyleemphasis{Compute cna/atom cutoff may be too large to find ghost atom neighbors}}
\sphinxAtStartPar
The neighbor cutoff used may not encompass enough ghost atoms
to perform this operation correctly.

\sphinxlineitem{\sphinxstyleemphasis{Computing temperature of portions of rigid bodies}}
\sphinxAtStartPar
The group defined by the temperature compute does not encompass all
the atoms in one or more rigid bodies, so the change in
degrees\sphinxhyphen{}of\sphinxhyphen{}freedom for the atoms in those partial rigid bodies will
not be accounted for.

\sphinxlineitem{\sphinxstyleemphasis{Create\_bonds max distance \textgreater{} minimum neighbor cutoff}}
\sphinxAtStartPar
This means atom pairs for some atom types may not be in the neighbor
list and thus no bond can be created between them.

\sphinxlineitem{\sphinxstyleemphasis{Delete\_atoms cutoff \textgreater{} minimum neighbor cutoff}}
\sphinxAtStartPar
This means atom pairs for some atom types may not be in the neighbor
list and thus an atom in that pair cannot be deleted.

\sphinxlineitem{\sphinxstyleemphasis{Dihedral atoms missing at step \%ld}}
\sphinxAtStartPar
One or more of 4 atoms needed to compute a particular dihedral are
missing on this processor.  Typically this is because the pairwise
cutoff is set too short or the dihedral has blown apart and an atom is
too far away.

\sphinxlineitem{\sphinxstyleemphasis{Dihedral problem}}
\sphinxAtStartPar
Conformation of the 4 listed dihedral atoms is extreme; you may want
to check your simulation geometry.

\sphinxlineitem{\sphinxstyleemphasis{Dihedral problem: \%d \%ld \%d \%d \%d \%d}}
\sphinxAtStartPar
Conformation of the 4 listed dihedral atoms is extreme; you may want
to check your simulation geometry.

\sphinxlineitem{\sphinxstyleemphasis{Dihedral style in data file differs from currently defined dihedral style}}
\sphinxAtStartPar
Self\sphinxhyphen{}explanatory.

\sphinxlineitem{\sphinxstyleemphasis{Dihedrals are defined but no dihedral style is set}}
\sphinxAtStartPar
The topology contains dihedrals, but there are no dihedral forces computed
since there was no dihedral\_style command.

\sphinxlineitem{\sphinxstyleemphasis{Dump dcd/xtc timestamp may be wrong with fix dt/reset}}
\sphinxAtStartPar
If the fix changes the timestep, the dump dcd file will not
reflect the change.

\sphinxlineitem{\sphinxstyleemphasis{Energy due to X extra global DOFs will be included in minimizer energies}}
\sphinxAtStartPar
When using fixes like box/relax, the potential energy used by the minimizer
is augmented by an additional energy provided by the fix. Thus the printed
converged energy may be different from the total potential energy.

\sphinxlineitem{\sphinxstyleemphasis{Estimated error in splitting of dispersion coeffs is \%g}}
\sphinxAtStartPar
Error is greater than 0.0001 percent.

\sphinxlineitem{\sphinxstyleemphasis{Ewald/disp Newton solver failed, using old method to estimate g\_ewald}}
\sphinxAtStartPar
Self\sphinxhyphen{}explanatory. Choosing a different cutoff value may help.

\sphinxlineitem{\sphinxstyleemphasis{FENE bond too long}}
\sphinxAtStartPar
A FENE bond has stretched dangerously far.  It’s interaction strength
will be truncated to attempt to prevent the bond from blowing up.

\sphinxlineitem{\sphinxstyleemphasis{FENE bond too long: \%ld \%d \%d \%g}}
\sphinxAtStartPar
A FENE bond has stretched dangerously far.  It’s interaction strength
will be truncated to attempt to prevent the bond from blowing up.

\sphinxlineitem{\sphinxstyleemphasis{FENE bond too long: \%ld \%g}}
\sphinxAtStartPar
A FENE bond has stretched dangerously far.  It’s interaction strength
will be truncated to attempt to prevent the bond from blowing up.

\sphinxlineitem{\sphinxstyleemphasis{Fix halt condition for fix\sphinxhyphen{}id \%s met on step \%ld with value \%g}}
\sphinxAtStartPar
Self explanatory.

\sphinxlineitem{\sphinxstyleemphasis{Fix SRD walls overlap but fix srd overlap not set}}
\sphinxAtStartPar
You likely want to set this in your input script.

\sphinxlineitem{\sphinxstyleemphasis{Fix bond/create is used multiple times or with fix bond/break \sphinxhyphen{} may not work as expected}}
\sphinxAtStartPar
When using fix bond/create multiple times or in combination with
fix bond/break, the individual fix instances do not share information
about changes they made at the same time step and thus it may result
in unexpected behavior.

\sphinxlineitem{\sphinxstyleemphasis{Fix bond/react: Atom affected by reaction too close to template edge}}
\sphinxAtStartPar
This means an atom which changes type or connectivity during the
reaction is too close to an ‘edge’ atom defined in the superimpose
file. This could cause incorrect assignment of bonds, angle, etc.
Generally, this means you must include more atoms in your templates,
such that there are at least two atoms between each atom involved in
the reaction and an edge atom.

\sphinxlineitem{\sphinxstyleemphasis{Fix bond/swap will ignore defined angles}}
\sphinxAtStartPar
See the page for fix bond/swap for more info on this
restriction.

\sphinxlineitem{\sphinxstyleemphasis{Fix deposit near setting \textless{} possible overlap separation \%g}}
\sphinxAtStartPar
This test is performed for finite size particles with a diameter, not
for point particles.  The near setting is smaller than the particle
diameter which can lead to overlaps.

\sphinxlineitem{\sphinxstyleemphasis{Fix evaporate may delete atom with non\sphinxhyphen{}zero molecule ID}}
\sphinxAtStartPar
This is probably an error, since you should not delete only one atom
of a molecule.

\sphinxlineitem{\sphinxstyleemphasis{Fix gcmc using full\_energy option}}
\sphinxAtStartPar
Fix gcmc has automatically turned on the full\_energy option since it
is required for systems like the one specified by the user. User input
included one or more of the following: kspace, triclinic, a hybrid
pair style, an eam pair style, or no “single” function for the pair
style.

\sphinxlineitem{\sphinxstyleemphasis{Fix langevin gjf using random gaussians is not implemented with kokkos}}
\sphinxAtStartPar
This will most likely cause errors in kinetic fluctuations.

\sphinxlineitem{\sphinxstyleemphasis{Fix property/atom mol or charge w/out ghost communication}}
\sphinxAtStartPar
A model typically needs these properties defined for ghost atoms.

\sphinxlineitem{\sphinxstyleemphasis{Fix qeq CG convergence failed (\%g) after \%d iterations at \%ld step}}
\sphinxAtStartPar
Self\sphinxhyphen{}explanatory.

\sphinxlineitem{\sphinxstyleemphasis{Fix qeq has non\sphinxhyphen{}zero lower Taper radius cutoff}}
\sphinxAtStartPar
Absolute value must be \textless{}= 0.01.

\sphinxlineitem{\sphinxstyleemphasis{Fix qeq has very low Taper radius cutoff}}
\sphinxAtStartPar
Value should typically be \textgreater{}= 5.0.

\sphinxlineitem{\sphinxstyleemphasis{Fix qeq/dynamic tolerance may be too small for damped dynamics}}
\sphinxAtStartPar
Self\sphinxhyphen{}explanatory.

\sphinxlineitem{\sphinxstyleemphasis{Fix qeq/fire tolerance may be too small for damped fires}}
\sphinxAtStartPar
Self\sphinxhyphen{}explanatory.

\sphinxlineitem{\sphinxstyleemphasis{Fix rattle should come after all other integration fixes}}
\sphinxAtStartPar
This fix is designed to work after all other integration fixes change
atom positions.  Thus it should be the last integration fix specified.
If not, it will not satisfy the desired constraints as well as it
otherwise would.

\sphinxlineitem{\sphinxstyleemphasis{Fix recenter should come after all other integration fixes}}
\sphinxAtStartPar
Other fixes may change the position of the center\sphinxhyphen{}of\sphinxhyphen{}mass, so
fix recenter should come last.

\sphinxlineitem{\sphinxstyleemphasis{Fix srd SRD moves may trigger frequent reneighboring}}
\sphinxAtStartPar
This is because the SRD particles may move long distances.

\sphinxlineitem{\sphinxstyleemphasis{Fix srd grid size \textgreater{} 1/4 of big particle diameter}}
\sphinxAtStartPar
This may cause accuracy problems.

\sphinxlineitem{\sphinxstyleemphasis{Fix srd particle moved outside valid domain}}
\sphinxAtStartPar
This may indicate a problem with your simulation parameters.

\sphinxlineitem{\sphinxstyleemphasis{Fix srd particles may move \textgreater{} big particle diameter}}
\sphinxAtStartPar
This may cause accuracy problems.

\sphinxlineitem{\sphinxstyleemphasis{Fix srd viscosity \textless{} 0.0 due to low SRD density}}
\sphinxAtStartPar
This may cause accuracy problems.

\sphinxlineitem{\sphinxstyleemphasis{Fixes cannot send data in Kokkos communication, switching to classic communication}}
\sphinxAtStartPar
This is current restriction with Kokkos.

\sphinxlineitem{\sphinxstyleemphasis{For better accuracy use ‘pair\_modify table 0’}}
\sphinxAtStartPar
The user\sphinxhyphen{}specified force accuracy cannot be achieved unless the table
feature is disabled by using ‘pair\_modify table 0’.

\sphinxlineitem{\sphinxstyleemphasis{Geometric mixing assumed for 1/r\textasciicircum{}6 coefficients}}
\sphinxAtStartPar
Self\sphinxhyphen{}explanatory.

\sphinxlineitem{\sphinxstyleemphasis{Group for fix\_modify temp != fix group}}
\sphinxAtStartPar
The fix\_modify command is specifying a temperature computation that
computes a temperature on a different group of atoms than the fix
itself operates on.  This is probably not what you want to do.

\sphinxlineitem{\sphinxstyleemphasis{H matrix size has been exceeded: m\_fill=\%d H.m=\%dn}}
\sphinxAtStartPar
This is the size of the matrix.

\sphinxlineitem{\sphinxstyleemphasis{Ignoring unknown or incorrect info command flag}}
\sphinxAtStartPar
Self\sphinxhyphen{}explanatory.  An unknown argument was given to the info command.
Compare your input with the documentation.

\sphinxlineitem{\sphinxstyleemphasis{Improper atoms missing at step \%ld}}
\sphinxAtStartPar
One or more of 4 atoms needed to compute a particular improper are
missing on this processor.  Typically this is because the pairwise
cutoff is set too short or the improper has blown apart and an atom is
too far away.

\sphinxlineitem{\sphinxstyleemphasis{Improper problem: \%d \%ld \%d \%d \%d \%d}}
\sphinxAtStartPar
Conformation of the 4 listed improper atoms is extreme; you may want
to check your simulation geometry.

\sphinxlineitem{\sphinxstyleemphasis{Improper style in data file differs from currently defined improper style}}
\sphinxAtStartPar
Self\sphinxhyphen{}explanatory.

\sphinxlineitem{\sphinxstyleemphasis{Impropers are defined but no improper style is set}}
\sphinxAtStartPar
The topology contains impropers, but there are no improper forces computed
since there was no improper\_style command.

\sphinxlineitem{\sphinxstyleemphasis{Inconsistent image flags}}
\sphinxAtStartPar
The image flags for a pair on bonded atoms appear to be inconsistent.
Inconsistent means that when the coordinates of the two atoms are
unwrapped using the image flags, the two atoms are far apart.
Specifically they are further apart than half a periodic box length.
Or they are more than a box length apart in a non\sphinxhyphen{}periodic dimension.
This is usually due to the initial data file not having correct image
flags for the two atoms in a bond that straddles a periodic boundary.
They should be different by 1 in that case.  This is a warning because
inconsistent image flags will not cause problems for dynamics or most
LAMMPS simulations.  However they can cause problems when such atoms
are used with the fix rigid or replicate commands.  Note that if you
have an infinite periodic crystal with bonds then it is impossible to
have fully consistent image flags, since some bonds will cross
periodic boundaries and connect two atoms with the same image
flag.

\sphinxlineitem{\sphinxstyleemphasis{Increasing communication cutoff for GPU style}}
\sphinxAtStartPar
The pair style has increased the communication cutoff to be consistent with
the communication cutoff requirements for this pair style when run on the GPU.

\sphinxlineitem{\sphinxstyleemphasis{KIM Model does not provide ‘energy’; Potential energy will be zero}}
\sphinxAtStartPar
Self\sphinxhyphen{}explanatory.

\sphinxlineitem{\sphinxstyleemphasis{KIM Model does not provide ‘forces’; Forces will be zero}}
\sphinxAtStartPar
Self\sphinxhyphen{}explanatory.

\sphinxlineitem{\sphinxstyleemphasis{KIM Model does not provide ‘particleEnergy’; energy per atom will be zero}}
\sphinxAtStartPar
Self\sphinxhyphen{}explanatory.

\sphinxlineitem{\sphinxstyleemphasis{KIM Model does not provide ‘particleVirial’; virial per atom will be zero}}
\sphinxAtStartPar
Self\sphinxhyphen{}explanatory.

\sphinxlineitem{\sphinxstyleemphasis{Kspace\_modify slab param \textless{} 2.0 may cause unphysical behavior}}
\sphinxAtStartPar
The kspace\_modify slab parameter should be larger to ensure periodic
grids padded with empty space do not overlap.

\sphinxlineitem{\sphinxstyleemphasis{Less insertions than requested}}
\sphinxAtStartPar
The fix pour command was unsuccessful at finding open space
for as many particles as it tried to insert.

\sphinxlineitem{\sphinxstyleemphasis{Library error in lammps\_gather\_atoms}}
\sphinxAtStartPar
This library function cannot be used if atom IDs are not defined
or are not consecutively numbered.

\sphinxlineitem{\sphinxstyleemphasis{Library error in lammps\_scatter\_atoms}}
\sphinxAtStartPar
This library function cannot be used if atom IDs are not defined or
are not consecutively numbered, or if no atom map is defined.  See the
atom\_modify command for details about atom maps.

\sphinxlineitem{\sphinxstyleemphasis{Likewise 1\sphinxhyphen{}2 special neighbor interactions != 1.0}}
\sphinxAtStartPar
The topology contains bonds, but there is no bond style defined
and a 1\sphinxhyphen{}2 special neighbor scaling factor was not 1.0. This
means that pair style interactions may have scaled or missing
pairs in the neighbor list in expectation of interactions for
those pairs being computed from the bond style.

\sphinxlineitem{\sphinxstyleemphasis{Likewise 1\sphinxhyphen{}3 special neighbor interactions != 1.0}}
\sphinxAtStartPar
The topology contains angles, but there is no angle style defined
and a 1\sphinxhyphen{}3 special neighbor scaling factor was not 1.0. This
means that pair style interactions may have scaled or missing
pairs in the neighbor list in expectation of interactions for
those pairs being computed from the angle style.

\sphinxlineitem{\sphinxstyleemphasis{Likewise 1\sphinxhyphen{}4 special neighbor interactions != 1.0}}
\sphinxAtStartPar
The topology contains dihedrals, but there is no dihedral style defined
and a 1\sphinxhyphen{}4 special neighbor scaling factor was not 1.0. This
means that pair style interactions may have scaled or missing
pairs in the neighbor list in expectation of interactions for
those pairs being computed from the dihedral style.

\sphinxlineitem{\sphinxstyleemphasis{Lost atoms via change\_box: original \%ld current \%ld}}
\sphinxAtStartPar
The command options you have used caused atoms to be lost.

\sphinxlineitem{\sphinxstyleemphasis{Lost atoms via displace\_atoms: original \%ld current \%ld}}
\sphinxAtStartPar
The command options you have used caused atoms to be lost.

\sphinxlineitem{\sphinxstyleemphasis{Lost atoms: original \%ld current \%ld}}
\sphinxAtStartPar
Lost atoms are checked for each time thermo output is done.  See the
thermo\_modify lost command for options.  Lost atoms usually indicate
bad dynamics, e.g. atoms have been blown far out of the simulation
box, or moved further than one processor’s subdomain away before
reneighboring.

\sphinxlineitem{\sphinxstyleemphasis{MSM mesh too small, increasing to 2 points in each direction}}
\sphinxAtStartPar
Self\sphinxhyphen{}explanatory.

\sphinxlineitem{\sphinxstyleemphasis{Mismatch between velocity and compute groups}}
\sphinxAtStartPar
The temperature computation used by the velocity command will not be
on the same group of atoms that velocities are being set for.

\sphinxlineitem{\sphinxstyleemphasis{Mixing forced for lj coefficients}}
\sphinxAtStartPar
Self\sphinxhyphen{}explanatory.

\sphinxlineitem{\sphinxstyleemphasis{Molecule attributes do not match system attributes}}
\sphinxAtStartPar
An attribute is specified (e.g. diameter, charge) that is
not defined for the specified atom style.

\sphinxlineitem{\sphinxstyleemphasis{Molecule has bond topology but no special bond settings}}
\sphinxAtStartPar
This means the bonded atoms will not be excluded in pairwise
interactions.

\sphinxlineitem{\sphinxstyleemphasis{Molecule template for create\_atoms has multiple molecules}}
\sphinxAtStartPar
The create\_atoms command will only create molecules of a single type,
i.e. the first molecule in the template.

\sphinxlineitem{\sphinxstyleemphasis{Molecule template for fix gcmc has multiple molecules}}
\sphinxAtStartPar
The fix gcmc command will only create molecules of a single type,
i.e. the first molecule in the template.

\sphinxlineitem{\sphinxstyleemphasis{Molecule template for fix shake has multiple molecules}}
\sphinxAtStartPar
The fix shake command will only recognize molecules of a single
type, i.e. the first molecule in the template.

\sphinxlineitem{\sphinxstyleemphasis{More than one compute centro/atom}}
\sphinxAtStartPar
It is not efficient to use compute centro/atom more than once.

\sphinxlineitem{\sphinxstyleemphasis{More than one compute cluster/atom}}
\sphinxAtStartPar
It is not efficient to use compute cluster/atom  more than once.

\sphinxlineitem{\sphinxstyleemphasis{More than one compute cna/atom defined}}
\sphinxAtStartPar
It is not efficient to use compute cna/atom  more than once.

\sphinxlineitem{\sphinxstyleemphasis{More than one compute contact/atom}}
\sphinxAtStartPar
It is not efficient to use compute contact/atom more than once.

\sphinxlineitem{\sphinxstyleemphasis{More than one compute coord/atom}}
\sphinxAtStartPar
It is not efficient to use compute coord/atom more than once.

\sphinxlineitem{\sphinxstyleemphasis{More than one compute damage/atom}}
\sphinxAtStartPar
It is not efficient to use compute ke/atom more than once.

\sphinxlineitem{\sphinxstyleemphasis{More than one compute dilatation/atom}}
\sphinxAtStartPar
Self\sphinxhyphen{}explanatory.

\sphinxlineitem{\sphinxstyleemphasis{More than one compute erotate/sphere/atom}}
\sphinxAtStartPar
It is not efficient to use compute erorate/sphere/atom more than once.

\sphinxlineitem{\sphinxstyleemphasis{More than one compute hexorder/atom}}
\sphinxAtStartPar
It is not efficient to use compute hexorder/atom more than once.

\sphinxlineitem{\sphinxstyleemphasis{More than one compute ke/atom}}
\sphinxAtStartPar
It is not efficient to use compute ke/atom more than once.

\sphinxlineitem{\sphinxstyleemphasis{More than one compute orientorder/atom}}
\sphinxAtStartPar
It is not efficient to use compute orientorder/atom more than once.

\sphinxlineitem{\sphinxstyleemphasis{More than one compute plasticity/atom}}
\sphinxAtStartPar
Self\sphinxhyphen{}explanatory.

\sphinxlineitem{\sphinxstyleemphasis{More than one compute sna/atom}}
\sphinxAtStartPar
Self\sphinxhyphen{}explanatory.

\sphinxlineitem{\sphinxstyleemphasis{More than one compute sna/grid}}
\sphinxAtStartPar
Self\sphinxhyphen{}explanatory.

\sphinxlineitem{\sphinxstyleemphasis{More than one compute sna/grid/local}}
\sphinxAtStartPar
Self\sphinxhyphen{}explanatory.

\sphinxlineitem{\sphinxstyleemphasis{More than one compute snad/atom}}
\sphinxAtStartPar
Self\sphinxhyphen{}explanatory.

\sphinxlineitem{\sphinxstyleemphasis{More than one compute snav/atom}}
\sphinxAtStartPar
Self\sphinxhyphen{}explanatory.

\sphinxlineitem{\sphinxstyleemphasis{More than one fix poems}}
\sphinxAtStartPar
It is not efficient to use fix poems more than once.

\sphinxlineitem{\sphinxstyleemphasis{More than one fix rigid}}
\sphinxAtStartPar
It is not efficient to use fix rigid more than once.

\sphinxlineitem{\sphinxstyleemphasis{Neighbor exclusions used with KSpace solver may give inconsistent Coulombic energies}}
\sphinxAtStartPar
This is because excluding specific pair interactions also excludes
them from long\sphinxhyphen{}range interactions which may not be the desired effect.
The special\_bonds command handles this consistently by ensuring
excluded (or weighted) 1\sphinxhyphen{}2, 1\sphinxhyphen{}3, 1\sphinxhyphen{}4 interactions are treated
consistently by both the short\sphinxhyphen{}range pair style and the long\sphinxhyphen{}range
solver.  This is not done for exclusions of charged atom pairs via the
neigh\_modify exclude command.

\sphinxlineitem{\sphinxstyleemphasis{New thermo\_style command, previous thermo\_modify settings will be lost}}
\sphinxAtStartPar
If a thermo\_style command is used after a thermo\_modify command, the
settings changed by the thermo\_modify command will be reset to their
default values.  This is because the thermo\_modify command acts on
the currently defined thermo style, and a thermo\_style command creates
a new style.

\sphinxlineitem{\sphinxstyleemphasis{No Kspace calculation with verlet/split}}
\sphinxAtStartPar
The second partition performs a kspace calculation so the kspace\_style
command must be used.

\sphinxlineitem{\sphinxstyleemphasis{No automatic unit conversion to XTC file format conventions possible for units lj}}
\sphinxAtStartPar
This means no scaling will be performed.

\sphinxlineitem{\sphinxstyleemphasis{No fixes defined, atoms won’t move}}
\sphinxAtStartPar
If you are not using a fix like nve, nvt, npt then atom velocities and
coordinates will not be updated during timestepping.

\sphinxlineitem{\sphinxstyleemphasis{No joints between rigid bodies, use fix rigid instead}}
\sphinxAtStartPar
The bodies defined by fix poems are not connected by joints.  POEMS
will integrate the body motion, but it would be more efficient to use
fix rigid.

\sphinxlineitem{\sphinxstyleemphasis{Not using real units with pair reaxff}}
\sphinxAtStartPar
This is most likely an error, unless you have created your own ReaxFF
parameter file in a different set of units.

\sphinxlineitem{\sphinxstyleemphasis{Number of MSM mesh points changed to be a multiple of 2}}
\sphinxAtStartPar
MSM requires that the number of grid points in each direction be a multiple
of two and the number of grid points in one or more directions have been
adjusted to meet this requirement.

\sphinxlineitem{\sphinxstyleemphasis{OMP\_NUM\_THREADS environment is not set.}}
\sphinxAtStartPar
This environment variable must be set appropriately to use the
OPENMP package.

\sphinxlineitem{\sphinxstyleemphasis{One or more atoms are time integrated more than once}}
\sphinxAtStartPar
This is probably an error since you typically do not want to
advance the positions or velocities of an atom more than once
per timestep.

\sphinxlineitem{\sphinxstyleemphasis{One or more chunks do not contain all atoms in molecule}}
\sphinxAtStartPar
This may not be what you intended.

\sphinxlineitem{\sphinxstyleemphasis{One or more dynamic groups may not be updated at correct point in timestep}}
\sphinxAtStartPar
If there are other fixes that act immediately after the initial stage
of time integration within a timestep (i.e. after atoms move), then
the command that sets up the dynamic group should appear after those
fixes.  This will ensure that dynamic group assignments are made
after all atoms have moved.

\sphinxlineitem{\sphinxstyleemphasis{One or more respa levels compute no forces}}
\sphinxAtStartPar
This is computationally inefficient.

\sphinxlineitem{\sphinxstyleemphasis{Pair COMB charge \%.10f with force \%.10f hit max barrier}}
\sphinxAtStartPar
Something is possibly wrong with your model.

\sphinxlineitem{\sphinxstyleemphasis{Pair COMB charge \%.10f with force \%.10f hit min barrier}}
\sphinxAtStartPar
Something is possibly wrong with your model.

\sphinxlineitem{\sphinxstyleemphasis{Pair brownian needs newton pair on for momentum conservation}}
\sphinxAtStartPar
Self\sphinxhyphen{}explanatory.

\sphinxlineitem{\sphinxstyleemphasis{Pair dpd needs newton pair on for momentum conservation}}
\sphinxAtStartPar
Self\sphinxhyphen{}explanatory.

\sphinxlineitem{\sphinxstyleemphasis{Pair dsmc: num\_of\_collisions \textgreater{} number\_of\_A}}
\sphinxAtStartPar
Collision model in DSMC is breaking down.

\sphinxlineitem{\sphinxstyleemphasis{Pair dsmc: num\_of\_collisions \textgreater{} number\_of\_B}}
\sphinxAtStartPar
Collision model in DSMC is breaking down.

\sphinxlineitem{\sphinxstyleemphasis{Pair style in data file differs from currently defined pair style}}
\sphinxAtStartPar
Self\sphinxhyphen{}explanatory.

\sphinxlineitem{\sphinxstyleemphasis{Pair style restartinfo set but has no restart support}}
\sphinxAtStartPar
This pair style has a bug, where it does not support reading and
writing information to a restart file, but does not set the member
variable “restartinfo” to 0 as required in that case.

\sphinxlineitem{\sphinxstyleemphasis{Particle deposition was unsuccessful}}
\sphinxAtStartPar
The fix deposit command was not able to insert as many atoms as
needed.  The requested volume fraction may be too high, or other atoms
may be in the insertion region.

\sphinxlineitem{\sphinxstyleemphasis{Proc subdomain size \textless{} neighbor skin, could lead to lost atoms}}
\sphinxAtStartPar
The decomposition of the physical domain (likely due to load
balancing) has led to a processor’s subdomain being smaller than the
neighbor skin in one or more dimensions.  Since reneighboring is
triggered by atoms moving the skin distance, this may lead to lost
atoms, if an atom moves all the way across a neighboring processor’s
subdomain before reneighboring is triggered.

\sphinxlineitem{\sphinxstyleemphasis{Reducing PPPM order b/c stencil extends beyond nearest neighbor processor}}
\sphinxAtStartPar
This may lead to a larger grid than desired.  See the kspace\_modify overlap
command to prevent changing of the PPPM order.

\sphinxlineitem{\sphinxstyleemphasis{Reducing PPPMDisp Coulomb order b/c stencil extends beyond neighbor processor}}
\sphinxAtStartPar
This may lead to a larger grid than desired.  See the kspace\_modify overlap
command to prevent changing of the PPPM order.

\sphinxlineitem{\sphinxstyleemphasis{Reducing PPPMDisp dispersion order b/c stencil extends beyond neighbor processor}}
\sphinxAtStartPar
This may lead to a larger grid than desired.  See the kspace\_modify overlap
command to prevent changing of the PPPM order.

\sphinxlineitem{\sphinxstyleemphasis{Replacing a fix, but new group != old group}}
\sphinxAtStartPar
The ID and style of a fix match for a fix you are changing with a fix
command, but the new group you are specifying does not match the old
group.

\sphinxlineitem{\sphinxstyleemphasis{Replicating in a non\sphinxhyphen{}periodic dimension}}
\sphinxAtStartPar
The parameters for a replicate command will cause a non\sphinxhyphen{}periodic
dimension to be replicated; this may cause unwanted behavior.

\sphinxlineitem{\sphinxstyleemphasis{Resetting reneighboring criteria during PRD}}
\sphinxAtStartPar
A PRD simulation requires that neigh\_modify settings be delay = 0,
every = 1, check = yes.  Since these settings were not in place,
LAMMPS changed them and will restore them to their original values
after the PRD simulation.

\sphinxlineitem{\sphinxstyleemphasis{Resetting reneighboring criteria during TAD}}
\sphinxAtStartPar
A TAD simulation requires that neigh\_modify settings be delay = 0,
every = 1, check = yes.  Since these settings were not in place,
LAMMPS changed them and will restore them to their original values
after the PRD simulation.

\sphinxlineitem{\sphinxstyleemphasis{Resetting reneighboring criteria during minimization}}
\sphinxAtStartPar
Minimization requires that neigh\_modify settings be delay = 0, every =
1, check = yes.  Since these settings were not in place, LAMMPS
changed them and will restore them to their original values after the
minimization.

\sphinxlineitem{\sphinxstyleemphasis{Restart file used different \# of processors}}
\sphinxAtStartPar
The restart file was written out by a LAMMPS simulation running on a
different number of processors.  Due to round\sphinxhyphen{}off, the trajectories of
your restarted simulation may diverge a little more quickly than if
you ran on the same \# of processors.

\sphinxlineitem{\sphinxstyleemphasis{Restart file used different 3d processor grid}}
\sphinxAtStartPar
The restart file was written out by a LAMMPS simulation running on a
different 3d grid of processors.  Due to round\sphinxhyphen{}off, the trajectories
of your restarted simulation may diverge a little more quickly than if
you ran on the same \# of processors.

\sphinxlineitem{\sphinxstyleemphasis{Restart file used different boundary settings, using restart file values}}
\sphinxAtStartPar
Your input script cannot change these restart file settings.

\sphinxlineitem{\sphinxstyleemphasis{Restart file used different newton bond setting, using restart file value}}
\sphinxAtStartPar
The restart file value will override the setting in the input script.

\sphinxlineitem{\sphinxstyleemphasis{Restart file used different newton pair setting, using input script value}}
\sphinxAtStartPar
The input script value will override the setting in the restart file.

\sphinxlineitem{\sphinxstyleemphasis{Restrain problem: \%d \%ld \%d \%d \%d \%d}}
\sphinxAtStartPar
Conformation of the 4 listed dihedral atoms is extreme; you may want
to check your simulation geometry.

\sphinxlineitem{\sphinxstyleemphasis{Running PRD with only one replica}}
\sphinxAtStartPar
This is allowed, but you will get no parallel speed\sphinxhyphen{}up.

\sphinxlineitem{\sphinxstyleemphasis{SRD bin shifting turned on due to small lamda}}
\sphinxAtStartPar
This is done to try to preserve accuracy.

\sphinxlineitem{\sphinxstyleemphasis{SRD bin size for fix srd differs from user request}}
\sphinxAtStartPar
Fix SRD had to adjust the bin size to fit the simulation box.  See the
cubic keyword if you want this message to be an error vs warning.

\sphinxlineitem{\sphinxstyleemphasis{SRD bins for fix srd are not cubic enough}}
\sphinxAtStartPar
The bin shape is not within tolerance of cubic.  See the cubic
keyword if you want this message to be an error vs warning.

\sphinxlineitem{\sphinxstyleemphasis{SRD particle \%d started inside big particle \%d on step \%ld bounce \%d}}
\sphinxAtStartPar
See the inside keyword if you want this message to be an error vs
warning.

\sphinxlineitem{\sphinxstyleemphasis{SRD particle \%d started inside wall \%d on step \%ld bounce \%d}}
\sphinxAtStartPar
See the inside keyword if you want this message to be an error vs
warning.

\sphinxlineitem{\sphinxstyleemphasis{Shake determinant \textless{} 0.0}}
\sphinxAtStartPar
The determinant of the quadratic equation being solved for a single
cluster specified by the fix shake command is numerically suspect.  LAMMPS
will set it to 0.0 and continue.

\sphinxlineitem{\sphinxstyleemphasis{Shell command ‘\%s’ failed with error ‘\%s’}}
\sphinxAtStartPar
Self\sphinxhyphen{}explanatory.

\sphinxlineitem{\sphinxstyleemphasis{Shell command returned with non\sphinxhyphen{}zero status}}
\sphinxAtStartPar
This may indicate the shell command did not operate as expected.

\sphinxlineitem{\sphinxstyleemphasis{Should not allow rigid bodies to bounce off reflecting walls}}
\sphinxAtStartPar
LAMMPS allows this, but their dynamics are not computed correctly.

\sphinxlineitem{\sphinxstyleemphasis{Should not use fix nve/limit with fix shake or fix rattle}}
\sphinxAtStartPar
This will lead to invalid constraint forces in the SHAKE/RATTLE
computation.

\sphinxlineitem{\sphinxstyleemphasis{Simulations might be very slow because of large number of structure factors}}
\sphinxAtStartPar
Self\sphinxhyphen{}explanatory.

\sphinxlineitem{\sphinxstyleemphasis{Slab correction not needed for MSM}}
\sphinxAtStartPar
Slab correction is intended to be used with Ewald or PPPM and is not needed by MSM.

\sphinxlineitem{\sphinxstyleemphasis{Specifying an ‘subset’ value of ‘0’ is equivalent to no ‘subset’ keyword}}
\sphinxAtStartPar
Self\sphinxhyphen{}explanatory.

\sphinxlineitem{\sphinxstyleemphasis{System is not charge neutral, net charge = \%g}}
\sphinxAtStartPar
The total charge on all atoms on the system is not 0.0.
For some KSpace solvers this is only a warning.

\sphinxlineitem{\sphinxstyleemphasis{Table inner cutoff \textgreater{}= outer cutoff}}
\sphinxAtStartPar
You specified an inner cutoff for a Coulombic table that is longer
than the global cutoff.  Probably not what you wanted.

\sphinxlineitem{\sphinxstyleemphasis{Temperature for MSST is not for group all}}
\sphinxAtStartPar
User\sphinxhyphen{}assigned temperature to MSST fix does not compute temperature for
all atoms.  Since MSST computes a global pressure, the kinetic energy
contribution from the temperature is assumed to also be for all atoms.
Thus the pressure used by MSST could be inaccurate.

\sphinxlineitem{\sphinxstyleemphasis{Temperature for NPT is not for group all}}
\sphinxAtStartPar
User\sphinxhyphen{}assigned temperature to NPT fix does not compute temperature for
all atoms.  Since NPT computes a global pressure, the kinetic energy
contribution from the temperature is assumed to also be for all atoms.
Thus the pressure used by NPT could be inaccurate.

\sphinxlineitem{\sphinxstyleemphasis{Temperature for fix modify is not for group all}}
\sphinxAtStartPar
The temperature compute is being used with a pressure calculation
which does operate on group all, so this may be inconsistent.

\sphinxlineitem{\sphinxstyleemphasis{Temperature for thermo pressure is not for group all}}
\sphinxAtStartPar
User\sphinxhyphen{}assigned temperature to thermo via the thermo\_modify command does
not compute temperature for all atoms.  Since thermo computes a global
pressure, the kinetic energy contribution from the temperature is
assumed to also be for all atoms.  Thus the pressure printed by thermo
could be inaccurate.

\sphinxlineitem{\sphinxstyleemphasis{The fix ave/spatial command has been replaced by the more flexible fix ave/chunk and compute chunk/atom commands \textendash{} fix ave/spatial will be removed in the summer of 2015}}
\sphinxAtStartPar
Self\sphinxhyphen{}explanatory.

\sphinxlineitem{\sphinxstyleemphasis{The minimizer does not re\sphinxhyphen{}orient dipoles when using fix efield}}
\sphinxAtStartPar
This means that only the atom coordinates will be minimized,
not the orientation of the dipoles.

\sphinxlineitem{\sphinxstyleemphasis{Too many common neighbors in CNA \%d times}}
\sphinxAtStartPar
More than the maximum \# of neighbors was found multiple times.  This
was unexpected.

\sphinxlineitem{\sphinxstyleemphasis{Too many inner timesteps in fix ttm}}
\sphinxAtStartPar
Self\sphinxhyphen{}explanatory.

\sphinxlineitem{\sphinxstyleemphasis{Too many neighbors in CNA for \%d atoms}}
\sphinxAtStartPar
More than the maximum \# of neighbors was found multiple times.  This
was unexpected.

\sphinxlineitem{\sphinxstyleemphasis{Use special bonds = 0,1,1 with bond style fene}}
\sphinxAtStartPar
Most FENE models need this setting for the special\_bonds command.

\sphinxlineitem{\sphinxstyleemphasis{Use special bonds = 0,1,1 with bond style fene/expand}}
\sphinxAtStartPar
Most FENE models need this setting for the special\_bonds command.

\sphinxlineitem{\sphinxstyleemphasis{Using a many\sphinxhyphen{}body potential with bonds/angles/dihedrals and special\_bond exclusions}}
\sphinxAtStartPar
This is likely not what you want to do.  The exclusion settings will
eliminate neighbors in the neighbor list, which the many\sphinxhyphen{}body potential
needs to calculated its terms correctly.

\sphinxlineitem{\sphinxstyleemphasis{Using compute temp/deform with inconsistent fix deform remap option}}
\sphinxAtStartPar
Fix nvt/sllod assumes deforming atoms have a velocity profile provided
by “remap v” or “remap none” as a fix deform option.

\sphinxlineitem{\sphinxstyleemphasis{Using compute temp/deform with no fix deform defined}}
\sphinxAtStartPar
This is probably an error, since it makes little sense to use
compute temp/deform in this case.

\sphinxlineitem{\sphinxstyleemphasis{Using fix srd with box deformation but no SRD thermostat}}
\sphinxAtStartPar
The deformation will heat the SRD particles so this can
be dangerous.

\sphinxlineitem{\sphinxstyleemphasis{Using kspace solver on system with no charge}}
\sphinxAtStartPar
Self\sphinxhyphen{}explanatory.

\sphinxlineitem{\sphinxstyleemphasis{Using largest cut\sphinxhyphen{}off for lj/long/dipole/long long long}}
\sphinxAtStartPar
Self\sphinxhyphen{}explanatory.

\sphinxlineitem{\sphinxstyleemphasis{Using largest cutoff for buck/long/coul/long}}
\sphinxAtStartPar
Self\sphinxhyphen{}explanatory.

\sphinxlineitem{\sphinxstyleemphasis{Using largest cutoff for lj/long/coul/long}}
\sphinxAtStartPar
Self\sphinxhyphen{}explanatory.

\sphinxlineitem{\sphinxstyleemphasis{Using largest cutoff for pair\_style lj/long/tip4p/long}}
\sphinxAtStartPar
Self\sphinxhyphen{}explanatory.

\sphinxlineitem{\sphinxstyleemphasis{Using package gpu without any pair style defined}}
\sphinxAtStartPar
Self\sphinxhyphen{}explanatory.

\sphinxlineitem{\sphinxstyleemphasis{Using pair potential shift with pair\_modify compute no}}
\sphinxAtStartPar
The shift effects will thus not be computed.

\sphinxlineitem{\sphinxstyleemphasis{Using pair tail corrections with nonperiodic system}}
\sphinxAtStartPar
This is probably a bogus thing to do, since tail corrections are
computed by integrating the density of a periodic system out to
infinity.

\sphinxlineitem{\sphinxstyleemphasis{Using pair tail corrections with pair\_modify compute no}}
\sphinxAtStartPar
The tail corrections will thus not be computed.

\end{description}


\part{Programmer Guide}
\label{\detokenize{Manual:programmer-guide}}\label{\detokenize{Manual:programmer-documentation}}
\sphinxstepscope


\chapter{LAMMPS Library Interfaces}
\label{\detokenize{Library:lammps-library-interfaces}}\label{\detokenize{Library::doc}}
\sphinxAtStartPar
As described on the {\hyperref[\detokenize{Howto_library::doc}]{\sphinxcrossref{\DUrole{doc}{library interface to LAMMPS}}}}
page, LAMMPS can be built as a library (static or shared), so that it
can be called by another code, used in a {\hyperref[\detokenize{Howto_couple::doc}]{\sphinxcrossref{\DUrole{doc}{coupled manner}}}} with other codes, or driven through a {\hyperref[\detokenize{Python_head::doc}]{\sphinxcrossref{\DUrole{doc}{Python
script}}}}.  The LAMMPS standalone executable itself is
essentially a thin wrapper on top of the LAMMPS library, which creates a
LAMMPS instance, passes the input for processing to that instance, and
then exits.

\sphinxAtStartPar
Most of the APIs described below are based on C language wrapper
functions in the files \sphinxcode{\sphinxupquote{src/library.h}} and \sphinxcode{\sphinxupquote{src/library.cpp}}, but
it is also possible to use C++ directly.  The basic procedure is
always the same: you create one or more instances of
{\hyperref[\detokenize{Classes_lammps:_CPPv4N9LAMMPS_NS6LAMMPSE}]{\sphinxcrossref{\sphinxcode{\sphinxupquote{LAMMPS}}}}}, pass commands as strings or
from files to that LAMMPS instance to execute calculations, and/or
call functions that read, manipulate, and update data from the active
class instances inside LAMMPS to do analysis or perform operations
that are not possible with existing input script commands.

\phantomsection\label{\detokenize{Library:thread-safety}}
\begin{sphinxadmonition}{note}{Thread\sphinxhyphen{}safety}

\sphinxAtStartPar
LAMMPS was initially not conceived as a thread\sphinxhyphen{}safe program, but over
the years changes have been applied to replace operations that
collide with creating multiple LAMMPS instances from multiple\sphinxhyphen{}threads
of the same process with thread\sphinxhyphen{}safe alternatives.  This primarily
applies to the core LAMMPS code and less so on add\sphinxhyphen{}on packages,
especially when those packages require additional code in the \sphinxstyleemphasis{lib}
folder, interface LAMMPS to Fortran libraries, or the code uses
static variables (like the COLVARS package).

\sphinxAtStartPar
Another major issue to deal with is to correctly handle MPI.
Creating a LAMMPS instance requires passing an MPI communicator, or
it assumes the \sphinxcode{\sphinxupquote{MPI\_COMM\_WORLD}} communicator, which spans all MPI
processor ranks.  When creating multiple LAMMPS object instances
from different threads, this communicator has to be different for
each thread or else collisions can happen.  Or it has to be
guaranteed, that only one thread at a time is active.  MPI
communicators, however, are not a problem, if LAMMPS is compiled
with the MPI STUBS library, which implies that there is no MPI
communication and only 1 MPI rank.
\end{sphinxadmonition}


\bigskip\hrule\bigskip



\section{LAMMPS C Library API}
\label{\detokenize{Library:lammps-c-library-api}}\label{\detokenize{Library:lammps-c-api}}
\sphinxAtStartPar
The C library interface is the most commonly used path to manage LAMMPS
instances from a compiled code and it is the basis for the {\hyperref[\detokenize{Python_module::doc}]{\sphinxcrossref{\DUrole{doc}{Python}}}} and {\hyperref[\detokenize{Fortran::doc}]{\sphinxcrossref{\DUrole{doc}{Fortran}}}} modules.  Almost all
functions of the C language API require an argument containing a
“handle” in the form of a \sphinxcode{\sphinxupquote{void *}} type variable, which points to the
location of a LAMMPS class instance.

\sphinxAtStartPar
The \sphinxcode{\sphinxupquote{library.h}} header file by default does not include the \sphinxcode{\sphinxupquote{mpi.h}}
header file and thus hides the {\hyperref[\detokenize{Library_create:_CPPv411lammps_openiPPc8MPI_CommPPv}]{\sphinxcrossref{\sphinxcode{\sphinxupquote{lammps\_open()}}}}} function which
requires the declaration of the \sphinxcode{\sphinxupquote{MPI\_comm}} data type.  This is only
a problem when the communicator that would be passed is different from
\sphinxcode{\sphinxupquote{MPI\_COMM\_WORLD}}.  Otherwise calling {\hyperref[\detokenize{Library_create:_CPPv418lammps_open_no_mpiiPPcPPv}]{\sphinxcrossref{\sphinxcode{\sphinxupquote{lammps\_open\_no\_mpi()}}}}}
will work just as well.  To make {\hyperref[\detokenize{Library_create:_CPPv411lammps_openiPPc8MPI_CommPPv}]{\sphinxcrossref{\sphinxcode{\sphinxupquote{lammps\_open()}}}}} available,
you need to compile the code with \sphinxcode{\sphinxupquote{\sphinxhyphen{}DLAMMPS\_LIB\_MPI}} or add the line
\sphinxcode{\sphinxupquote{\#define LAMMPS\_LIB\_MPI}} before \sphinxcode{\sphinxupquote{\#include "library.h"}}.

\sphinxAtStartPar
Please note the \sphinxcode{\sphinxupquote{mpi.h}} file must usually be the same (and thus the
MPI library in use) for the LAMMPS code and library and the calling code.
The exception is when LAMMPS was compiled in serial mode using the
\sphinxcode{\sphinxupquote{STUBS}} MPI library.  In that case the calling code may be compiled
with a different MPI library so long as {\hyperref[\detokenize{Library_create:_CPPv418lammps_open_no_mpiiPPcPPv}]{\sphinxcrossref{\sphinxcode{\sphinxupquote{lammps\_open\_no\_mpi()}}}}}
is called to create a LAMMPS instance.  In that case each MPI rank will
run LAMMPS in serial mode.

\begin{sphinxadmonition}{note}{Errors versus exceptions}

\sphinxAtStartPar
If the LAMMPS executable encounters an error condition, it will abort
after printing an error message.  It does so by catching the
exceptions that LAMMPS could throw.  For a C library interface this
is usually not desirable since the calling code might lack the
ability to catch such exceptions.  Thus, the library functions will
catch those exceptions and return from the affected functions.  The
error status {\hyperref[\detokenize{Library_utility:_CPPv416lammps_has_errorPv}]{\sphinxcrossref{\sphinxcode{\sphinxupquote{can be queried}}}}} and an
{\hyperref[\detokenize{Library_utility:_CPPv429lammps_get_last_error_messagePvPci}]{\sphinxcrossref{\sphinxcode{\sphinxupquote{error message retrieved}}}}}.
This is, for example used by the {\hyperref[\detokenize{Python_module::doc}]{\sphinxcrossref{\DUrole{doc}{LAMMPS python module}}}} and then a suitable Python exception is thrown.
\end{sphinxadmonition}

\begin{sphinxadmonition}{note}{Using the C library interface as a plugin}

\sphinxAtStartPar
Rather than including the C library directly and link to the LAMMPS
library at compile time, you can use the \sphinxcode{\sphinxupquote{liblammpsplugin.h}} header
file and the \sphinxcode{\sphinxupquote{liblammpsplugin.c}} C code in the
\sphinxcode{\sphinxupquote{examples/COUPLE/plugin}} folder for an interface to LAMMPS that is
largely identical to the regular library interface, only that it will
load a LAMMPS shared library file at runtime.  This can be useful for
applications where the interface to LAMMPS would be an optional
feature.
\end{sphinxadmonition}

\begin{sphinxadmonition}{warning}{Warning:}
\sphinxAtStartPar
No checks are made on the arguments of the function calls of the C
library interface.  \sphinxstyleemphasis{All} function arguments must be non\sphinxhyphen{}NULL unless
\sphinxstyleemphasis{explicitly} allowed, and must point to consistent and valid data.  Buffers
for storing returned data must be allocated to a suitable size.
Passing invalid or unsuitable information will likely cause crashes
or corrupt data.
\end{sphinxadmonition}


\bigskip\hrule\bigskip


\sphinxstepscope


\subsection{Creating or deleting a LAMMPS object}
\label{\detokenize{Library_create:creating-or-deleting-a-lammps-object}}\label{\detokenize{Library_create::doc}}
\sphinxAtStartPar
This section documents the following functions:
\begin{itemize}
\item {} 
\sphinxAtStartPar
{\hyperref[\detokenize{Library_create:_CPPv411lammps_openiPPc8MPI_CommPPv}]{\sphinxcrossref{\sphinxcode{\sphinxupquote{lammps\_open()}}}}}

\item {} 
\sphinxAtStartPar
{\hyperref[\detokenize{Library_create:_CPPv418lammps_open_no_mpiiPPcPPv}]{\sphinxcrossref{\sphinxcode{\sphinxupquote{lammps\_open\_no\_mpi()}}}}}

\item {} 
\sphinxAtStartPar
{\hyperref[\detokenize{Library_create:_CPPv419lammps_open_fortraniPPci}]{\sphinxcrossref{\sphinxcode{\sphinxupquote{lammps\_open\_fortran()}}}}}

\item {} 
\sphinxAtStartPar
{\hyperref[\detokenize{Library_create:_CPPv412lammps_closePv}]{\sphinxcrossref{\sphinxcode{\sphinxupquote{lammps\_close()}}}}}

\item {} 
\sphinxAtStartPar
{\hyperref[\detokenize{Library_create:_CPPv415lammps_mpi_initv}]{\sphinxcrossref{\sphinxcode{\sphinxupquote{lammps\_mpi\_init()}}}}}

\item {} 
\sphinxAtStartPar
{\hyperref[\detokenize{Library_create:_CPPv419lammps_mpi_finalizev}]{\sphinxcrossref{\sphinxcode{\sphinxupquote{lammps\_mpi\_finalize()}}}}}

\item {} 
\sphinxAtStartPar
{\hyperref[\detokenize{Library_create:_CPPv422lammps_kokkos_finalizev}]{\sphinxcrossref{\sphinxcode{\sphinxupquote{lammps\_kokkos\_finalize()}}}}}

\item {} 
\sphinxAtStartPar
{\hyperref[\detokenize{Library_create:_CPPv422lammps_python_finalizev}]{\sphinxcrossref{\sphinxcode{\sphinxupquote{lammps\_python\_finalize()}}}}}

\item {} 
\sphinxAtStartPar
{\hyperref[\detokenize{Library_create:_CPPv412lammps_errorPviPKc}]{\sphinxcrossref{\sphinxcode{\sphinxupquote{lammps\_error()}}}}}

\end{itemize}


\bigskip\hrule\bigskip


\sphinxAtStartPar
The {\hyperref[\detokenize{Library_create:_CPPv411lammps_openiPPc8MPI_CommPPv}]{\sphinxcrossref{\sphinxcode{\sphinxupquote{lammps\_open()}}}}} and {\hyperref[\detokenize{Library_create:_CPPv418lammps_open_no_mpiiPPcPPv}]{\sphinxcrossref{\sphinxcode{\sphinxupquote{lammps\_open\_no\_mpi()}}}}} functions
are used to create and initialize a \sphinxcode{\sphinxupquote{LAMMPS()}} instance.  They
return a reference to this instance as a \sphinxcode{\sphinxupquote{void *}} pointer to be used
as the “handle” argument in subsequent function calls until that
instance is destroyed by calling {\hyperref[\detokenize{Library_create:_CPPv412lammps_closePv}]{\sphinxcrossref{\sphinxcode{\sphinxupquote{lammps\_close()}}}}}.  Here is a
simple example demonstrating its use:

\begin{sphinxVerbatim}[commandchars=\\\{\}]
\PYG{c+cp}{\PYGZsh{}}\PYG{c+cp}{include}\PYG{+w}{ }\PYG{c+cpf}{\PYGZdq{}library.h\PYGZdq{}}
\PYG{c+cp}{\PYGZsh{}}\PYG{c+cp}{include}\PYG{+w}{ }\PYG{c+cpf}{\PYGZlt{}stdio.h\PYGZgt{}}

\PYG{k+kt}{int}\PYG{+w}{ }\PYG{n+nf}{main}\PYG{p}{(}\PYG{k+kt}{int}\PYG{+w}{ }\PYG{n}{argc}\PYG{p}{,}\PYG{+w}{ }\PYG{k+kt}{char}\PYG{+w}{ }\PYG{o}{*}\PYG{o}{*}\PYG{n}{argv}\PYG{p}{)}
\PYG{p}{\PYGZob{}}
\PYG{+w}{  }\PYG{k+kt}{void}\PYG{+w}{ }\PYG{o}{*}\PYG{n}{handle}\PYG{p}{;}
\PYG{+w}{  }\PYG{k+kt}{int}\PYG{+w}{ }\PYG{n}{version}\PYG{p}{;}
\PYG{+w}{  }\PYG{k}{const}\PYG{+w}{ }\PYG{k+kt}{char}\PYG{+w}{ }\PYG{o}{*}\PYG{n}{lmpargv}\PYG{p}{[}\PYG{p}{]}\PYG{+w}{ }\PYG{o}{=}\PYG{+w}{ }\PYG{p}{\PYGZob{}}\PYG{+w}{ }\PYG{l+s}{\PYGZdq{}}\PYG{l+s}{liblammps}\PYG{l+s}{\PYGZdq{}}\PYG{p}{,}\PYG{+w}{ }\PYG{l+s}{\PYGZdq{}}\PYG{l+s}{\PYGZhy{}log}\PYG{l+s}{\PYGZdq{}}\PYG{p}{,}\PYG{+w}{ }\PYG{l+s}{\PYGZdq{}}\PYG{l+s}{none}\PYG{l+s}{\PYGZdq{}}\PYG{p}{\PYGZcb{}}\PYG{p}{;}
\PYG{+w}{  }\PYG{k+kt}{int}\PYG{+w}{ }\PYG{n}{lmpargc}\PYG{+w}{ }\PYG{o}{=}\PYG{+w}{ }\PYG{k}{sizeof}\PYG{p}{(}\PYG{n}{lmpargv}\PYG{p}{)}\PYG{o}{/}\PYG{k}{sizeof}\PYG{p}{(}\PYG{k}{const}\PYG{+w}{ }\PYG{k+kt}{char}\PYG{+w}{ }\PYG{o}{*}\PYG{p}{)}\PYG{p}{;}

\PYG{+w}{  }\PYG{c+cm}{/* create LAMMPS instance */}
\PYG{+w}{  }\PYG{n}{handle}\PYG{+w}{ }\PYG{o}{=}\PYG{+w}{ }\PYG{n}{lammps\PYGZus{}open\PYGZus{}no\PYGZus{}mpi}\PYG{p}{(}\PYG{n}{lmpargc}\PYG{p}{,}\PYG{+w}{ }\PYG{p}{(}\PYG{k+kt}{char}\PYG{+w}{ }\PYG{o}{*}\PYG{o}{*}\PYG{p}{)}\PYG{n}{lmpargv}\PYG{p}{,}\PYG{+w}{ }\PYG{n+nb}{NULL}\PYG{p}{)}\PYG{p}{;}
\PYG{+w}{  }\PYG{k}{if}\PYG{+w}{ }\PYG{p}{(}\PYG{n}{handle}\PYG{+w}{ }\PYG{o}{=}\PYG{o}{=}\PYG{+w}{ }\PYG{n+nb}{NULL}\PYG{p}{)}\PYG{+w}{ }\PYG{p}{\PYGZob{}}
\PYG{+w}{    }\PYG{n}{printf}\PYG{p}{(}\PYG{l+s}{\PYGZdq{}}\PYG{l+s}{LAMMPS initialization failed}\PYG{l+s}{\PYGZdq{}}\PYG{p}{)}\PYG{p}{;}
\PYG{+w}{    }\PYG{n}{lammps\PYGZus{}mpi\PYGZus{}finalize}\PYG{p}{(}\PYG{p}{)}\PYG{p}{;}
\PYG{+w}{    }\PYG{k}{return}\PYG{+w}{ }\PYG{l+m+mi}{1}\PYG{p}{;}
\PYG{+w}{  }\PYG{p}{\PYGZcb{}}

\PYG{+w}{  }\PYG{c+cm}{/* get and print numerical version code */}
\PYG{+w}{  }\PYG{n}{version}\PYG{+w}{ }\PYG{o}{=}\PYG{+w}{ }\PYG{n}{lammps\PYGZus{}version}\PYG{p}{(}\PYG{n}{handle}\PYG{p}{)}\PYG{p}{;}
\PYG{+w}{  }\PYG{n}{printf}\PYG{p}{(}\PYG{l+s}{\PYGZdq{}}\PYG{l+s}{LAMMPS Version: \PYGZpc{}d}\PYG{l+s+se}{\PYGZbs{}n}\PYG{l+s}{\PYGZdq{}}\PYG{p}{,}\PYG{n}{version}\PYG{p}{)}\PYG{p}{;}

\PYG{+w}{  }\PYG{c+cm}{/* delete LAMMPS instance and shut down MPI */}
\PYG{+w}{  }\PYG{n}{lammps\PYGZus{}close}\PYG{p}{(}\PYG{n}{handle}\PYG{p}{)}\PYG{p}{;}
\PYG{+w}{  }\PYG{n}{lammps\PYGZus{}mpi\PYGZus{}finalize}\PYG{p}{(}\PYG{p}{)}\PYG{p}{;}
\PYG{+w}{  }\PYG{k}{return}\PYG{+w}{ }\PYG{l+m+mi}{0}\PYG{p}{;}
\PYG{p}{\PYGZcb{}}
\end{sphinxVerbatim}

\sphinxAtStartPar
The LAMMPS library uses the MPI library it was compiled with and will
either run on all processors in the \sphinxcode{\sphinxupquote{MPI\_COMM\_WORLD}} communicator or
on the set of processors in the communicator passed as the \sphinxcode{\sphinxupquote{comm}}
argument of {\hyperref[\detokenize{Library_create:_CPPv411lammps_openiPPc8MPI_CommPPv}]{\sphinxcrossref{\sphinxcode{\sphinxupquote{lammps\_open()}}}}}.  This means the calling code can
run LAMMPS on all or a subset of processors.  For example, a wrapper
code might decide to alternate between LAMMPS and another code, allowing
them both to run on all the processors.  Or it might allocate part of
the processors to LAMMPS and the rest to the other code by creating a
custom communicator with \sphinxcode{\sphinxupquote{MPI\_Comm\_split()}} and running both codes
concurrently before syncing them up periodically.  Or it might
instantiate multiple instances of LAMMPS to perform different
calculations and either alternate between them, run them concurrently on
split communicators, or run them one after the other.  The
{\hyperref[\detokenize{Library_create:_CPPv411lammps_openiPPc8MPI_CommPPv}]{\sphinxcrossref{\sphinxcode{\sphinxupquote{lammps\_open()}}}}} function may be called multiple times for this
latter purpose.

\sphinxAtStartPar
The {\hyperref[\detokenize{Library_create:_CPPv412lammps_closePv}]{\sphinxcrossref{\sphinxcode{\sphinxupquote{lammps\_close()}}}}} function is used to shut down the
{\hyperref[\detokenize{Classes_lammps:_CPPv4N9LAMMPS_NS6LAMMPSE}]{\sphinxcrossref{\sphinxcode{\sphinxupquote{LAMMPS}}}}} class pointed to by the handle
passed as an argument and free all its memory. This has to be called
for every instance created with one of the {\hyperref[\detokenize{Library_create:_CPPv411lammps_openiPPc8MPI_CommPPv}]{\sphinxcrossref{\sphinxcode{\sphinxupquote{lammps\_open()}}}}}
functions.  It will, however, \sphinxstylestrong{not} call \sphinxcode{\sphinxupquote{MPI\_Finalize()}}, since
that may only be called once.  See {\hyperref[\detokenize{Library_create:_CPPv419lammps_mpi_finalizev}]{\sphinxcrossref{\sphinxcode{\sphinxupquote{lammps\_mpi\_finalize()}}}}} for
an alternative to invoking \sphinxcode{\sphinxupquote{MPI\_Finalize()}} explicitly from the
calling program.


\bigskip\hrule\bigskip

\index{lammps\_open (C++ function)@\spxentry{lammps\_open}\spxextra{C++ function}}

\begin{fulllineitems}
\phantomsection\label{\detokenize{Library_create:_CPPv411lammps_openiPPc8MPI_CommPPv}}\phantomsection\label{\detokenize{Library_create:_CPPv311lammps_openiPPc8MPI_CommPPv}}\phantomsection\label{\detokenize{Library_create:_CPPv211lammps_openiPPc8MPI_CommPPv}}\phantomsection\label{\detokenize{Library_create:lammps_open__i.cPP.MPI_Comm.voidPP}}
\pysigstartsignatures
\pysigstartmultiline
\pysiglinewithargsret{\phantomsection\label{\detokenize{Library_create:library_8h_1a72b42bbe7e60dc8851ef0e955cfad399}}\DUrole{kt}{void}\DUrole{w}{ }\DUrole{p}{*}\sphinxbfcode{\sphinxupquote{\DUrole{n}{lammps\_open}}}}{\DUrole{kt}{int}\DUrole{w}{ }\DUrole{n,sig-param}{argc}\sphinxparamcomma \DUrole{kt}{char}\DUrole{w}{ }\DUrole{p}{*}\DUrole{p}{*}\DUrole{n,sig-param}{argv}\sphinxparamcomma \DUrole{n}{MPI\_Comm}\DUrole{w}{ }\DUrole{n,sig-param}{comm}\sphinxparamcomma \DUrole{kt}{void}\DUrole{w}{ }\DUrole{p}{*}\DUrole{p}{*}\DUrole{n,sig-param}{ptr}}{}
\pysigstopmultiline
\pysigstopsignatures
\sphinxAtStartPar
Create instance of the LAMMPS class and return pointer to it.

\sphinxAtStartPar

\sphinxAtStartPar
The {\hyperref[\detokenize{Library_create:_CPPv411lammps_openiPPc8MPI_CommPPv}]{\sphinxcrossref{\sphinxcode{\sphinxupquote{lammps\_open()}}}}} function creates a new {\hyperref[\detokenize{Classes_lammps:_CPPv4N9LAMMPS_NS6LAMMPSE}]{\sphinxcrossref{\sphinxcode{\sphinxupquote{LAMMPS}}}}} class instance while passing in a list of strings
as if they were {\hyperref[\detokenize{Run_options::doc}]{\sphinxcrossref{\DUrole{doc}{command\sphinxhyphen{}line arguments}}}} for the
LAMMPS executable, and an MPI communicator for LAMMPS to run under.
Since the list of arguments is \sphinxstylestrong{exactly} as when called from the
command line, the first argument would be the name of the executable and
thus is otherwise ignored.  However \sphinxcode{\sphinxupquote{argc}} may be set to 0 and then
\sphinxcode{\sphinxupquote{argv}} may be \sphinxcode{\sphinxupquote{NULL}}.  If MPI is not yet initialized, \sphinxcode{\sphinxupquote{MPI\_Init()}}
will be called during creation of the LAMMPS class instance.

\sphinxAtStartPar
If for some reason the creation or initialization of the LAMMPS instance
fails a null pointer is returned.

\sphinxAtStartPar
\DUrole{versionmodified,changed}{Changed in version 18Sep2020: }This function now has the pointer to the created LAMMPS class
instance as return value.  For backward compatibility it is still
possible to provide the address of a pointer variable as final
argument \sphinxstyleemphasis{ptr}.

\sphinxAtStartPar
\DUrole{versionmodified,deprecated}{Deprecated since version 18Sep2020: }The \sphinxstyleemphasis{ptr} argument will be removed in a future release of LAMMPS.
It should be set to \sphinxcode{\sphinxupquote{NULL}} instead.
\begin{description}
\sphinxlineitem{\sphinxstyleemphasis{See also}}
\sphinxAtStartPar
{\hyperref[\detokenize{Library_create:_CPPv418lammps_open_no_mpiiPPcPPv}]{\sphinxcrossref{\sphinxcode{\sphinxupquote{lammps\_open\_no\_mpi()}}}}}, {\hyperref[\detokenize{Library_create:_CPPv419lammps_open_fortraniPPci}]{\sphinxcrossref{\sphinxcode{\sphinxupquote{lammps\_open\_fortran()}}}}}

\end{description}


\begin{sphinxadmonition}{note}{Note:}
\sphinxAtStartPar
This function is \sphinxstylestrong{only} declared when the code using the LAMMPS
\sphinxcode{\sphinxupquote{library.h}} include file is compiled with \sphinxcode{\sphinxupquote{\sphinxhyphen{}DLAMMPS\_LIB\_MPI}},
or contains a \sphinxcode{\sphinxupquote{\#define LAMMPS\_LIB\_MPI 1}} statement before
\sphinxcode{\sphinxupquote{\#include "library.h"}}.  Otherwise you can only use the
{\hyperref[\detokenize{Library_create:_CPPv418lammps_open_no_mpiiPPcPPv}]{\sphinxcrossref{\sphinxcode{\sphinxupquote{lammps\_open\_no\_mpi()}}}}} or {\hyperref[\detokenize{Library_create:_CPPv419lammps_open_fortraniPPci}]{\sphinxcrossref{\sphinxcode{\sphinxupquote{lammps\_open\_fortran()}}}}}
functions.
\end{sphinxadmonition}
\begin{quote}\begin{description}
\sphinxlineitem{Parameters}\begin{itemize}
\item {} 
\sphinxAtStartPar
\sphinxstyleliteralstrong{\sphinxupquote{argc}} \textendash{} number of command line arguments 

\item {} 
\sphinxAtStartPar
\sphinxstyleliteralstrong{\sphinxupquote{argv}} \textendash{} list of command line argument strings 

\item {} 
\sphinxAtStartPar
\sphinxstyleliteralstrong{\sphinxupquote{comm}} \textendash{} MPI communicator for this LAMMPS instance 

\item {} 
\sphinxAtStartPar
\sphinxstyleliteralstrong{\sphinxupquote{ptr}} \textendash{} pointer to a void pointer variable which serves as a handle; may be \sphinxcode{\sphinxupquote{NULL}}

\end{itemize}

\sphinxlineitem{Returns}
\sphinxAtStartPar
pointer to new LAMMPS instance cast to \sphinxcode{\sphinxupquote{void *}}

\end{description}\end{quote}

\end{fulllineitems}



\bigskip\hrule\bigskip

\index{lammps\_open\_no\_mpi (C++ function)@\spxentry{lammps\_open\_no\_mpi}\spxextra{C++ function}}

\begin{fulllineitems}
\phantomsection\label{\detokenize{Library_create:_CPPv418lammps_open_no_mpiiPPcPPv}}\phantomsection\label{\detokenize{Library_create:_CPPv318lammps_open_no_mpiiPPcPPv}}\phantomsection\label{\detokenize{Library_create:_CPPv218lammps_open_no_mpiiPPcPPv}}\phantomsection\label{\detokenize{Library_create:lammps_open_no_mpi__i.cPP.voidPP}}
\pysigstartsignatures
\pysigstartmultiline
\pysiglinewithargsret{\phantomsection\label{\detokenize{Library_create:library_8h_1a01c641f102a6d3fa96ed34750914a98a}}\DUrole{kt}{void}\DUrole{w}{ }\DUrole{p}{*}\sphinxbfcode{\sphinxupquote{\DUrole{n}{lammps\_open\_no\_mpi}}}}{\DUrole{kt}{int}\DUrole{w}{ }\DUrole{n,sig-param}{argc}\sphinxparamcomma \DUrole{kt}{char}\DUrole{w}{ }\DUrole{p}{*}\DUrole{p}{*}\DUrole{n,sig-param}{argv}\sphinxparamcomma \DUrole{kt}{void}\DUrole{w}{ }\DUrole{p}{*}\DUrole{p}{*}\DUrole{n,sig-param}{ptr}}{}
\pysigstopmultiline
\pysigstopsignatures
\sphinxAtStartPar
Variant of \sphinxcode{\sphinxupquote{lammps\_open()}} that implicitly uses \sphinxcode{\sphinxupquote{MPI\_COMM\_WORLD}}.

\sphinxAtStartPar

\sphinxAtStartPar
This function is a version of {\hyperref[\detokenize{Library_create:_CPPv411lammps_openiPPc8MPI_CommPPv}]{\sphinxcrossref{\sphinxcode{\sphinxupquote{lammps\_open()}}}}}, that is missing
the MPI communicator argument.  It will use \sphinxcode{\sphinxupquote{MPI\_COMM\_WORLD}} instead.
The type and purpose of arguments and return value are otherwise the
same.

\sphinxAtStartPar
Outside of the convenience, this function is useful, when the LAMMPS
library was compiled in serial mode, but the calling code runs in
parallel and the \sphinxcode{\sphinxupquote{MPI\_Comm}} data type of the STUBS library would not
be compatible with that of the calling code.

\sphinxAtStartPar
If for some reason the creation or initialization of the LAMMPS instance
fails a null pointer is returned.

\sphinxAtStartPar
\DUrole{versionmodified,changed}{Changed in version 18Sep2020: }This function now has the pointer to the created LAMMPS class
instance as return value.  For backward compatibility it is still
possible to provide the address of a pointer variable as final
argument \sphinxstyleemphasis{ptr}.

\sphinxAtStartPar
\DUrole{versionmodified,deprecated}{Deprecated since version 18Sep2020: }The \sphinxstyleemphasis{ptr} argument will be removed in a future release of LAMMPS.
It should be set to \sphinxcode{\sphinxupquote{NULL}} instead.
\begin{description}
\sphinxlineitem{\sphinxstyleemphasis{See also}}
\sphinxAtStartPar
{\hyperref[\detokenize{Library_create:_CPPv411lammps_openiPPc8MPI_CommPPv}]{\sphinxcrossref{\sphinxcode{\sphinxupquote{lammps\_open()}}}}}, {\hyperref[\detokenize{Library_create:_CPPv419lammps_open_fortraniPPci}]{\sphinxcrossref{\sphinxcode{\sphinxupquote{lammps\_open\_fortran()}}}}}

\end{description}

\begin{quote}\begin{description}
\sphinxlineitem{Parameters}\begin{itemize}
\item {} 
\sphinxAtStartPar
\sphinxstyleliteralstrong{\sphinxupquote{argc}} \textendash{} number of command line arguments 

\item {} 
\sphinxAtStartPar
\sphinxstyleliteralstrong{\sphinxupquote{argv}} \textendash{} list of command line argument strings 

\item {} 
\sphinxAtStartPar
\sphinxstyleliteralstrong{\sphinxupquote{ptr}} \textendash{} pointer to a void pointer variable which serves as a handle; may be \sphinxcode{\sphinxupquote{NULL}}

\end{itemize}

\sphinxlineitem{Returns}
\sphinxAtStartPar
pointer to new LAMMPS instance cast to \sphinxcode{\sphinxupquote{void *}}

\end{description}\end{quote}

\end{fulllineitems}



\bigskip\hrule\bigskip

\index{lammps\_open\_fortran (C++ function)@\spxentry{lammps\_open\_fortran}\spxextra{C++ function}}

\begin{fulllineitems}
\phantomsection\label{\detokenize{Library_create:_CPPv419lammps_open_fortraniPPci}}\phantomsection\label{\detokenize{Library_create:_CPPv319lammps_open_fortraniPPci}}\phantomsection\label{\detokenize{Library_create:_CPPv219lammps_open_fortraniPPci}}\phantomsection\label{\detokenize{Library_create:lammps_open_fortran__i.cPP.i}}
\pysigstartsignatures
\pysigstartmultiline
\pysiglinewithargsret{\phantomsection\label{\detokenize{Library_create:library_8h_1a4f33b973643d1fee38ecdd87c81bbb0f}}\DUrole{kt}{void}\DUrole{w}{ }\DUrole{p}{*}\sphinxbfcode{\sphinxupquote{\DUrole{n}{lammps\_open\_fortran}}}}{\DUrole{kt}{int}\DUrole{w}{ }\DUrole{n,sig-param}{argc}\sphinxparamcomma \DUrole{kt}{char}\DUrole{w}{ }\DUrole{p}{*}\DUrole{p}{*}\DUrole{n,sig-param}{argv}\sphinxparamcomma \DUrole{kt}{int}\DUrole{w}{ }\DUrole{n,sig-param}{f\_comm}}{}
\pysigstopmultiline
\pysigstopsignatures
\sphinxAtStartPar
Variant of \sphinxcode{\sphinxupquote{lammps\_open()}} using a Fortran MPI communicator.

\sphinxAtStartPar

\sphinxAtStartPar
\DUrole{versionmodified,added}{Added in version 18Sep2020.}

\sphinxAtStartPar
This function is a version of {\hyperref[\detokenize{Library_create:_CPPv411lammps_openiPPc8MPI_CommPPv}]{\sphinxcrossref{\sphinxcode{\sphinxupquote{lammps\_open()}}}}}, that uses an
integer for the MPI communicator as the MPI Fortran interface does.  It
is used in the \sphinxcode{\sphinxupquote{lammps()}} constructor of the LAMMPS Fortran
module.  Internally it converts the \sphinxstyleemphasis{f\_comm} argument into a C\sphinxhyphen{}style MPI
communicator with \sphinxcode{\sphinxupquote{MPI\_Comm\_f2c()}} and then calls
{\hyperref[\detokenize{Library_create:_CPPv411lammps_openiPPc8MPI_CommPPv}]{\sphinxcrossref{\sphinxcode{\sphinxupquote{lammps\_open()}}}}}.

\sphinxAtStartPar
If for some reason the creation or initialization of the LAMMPS instance
fails a null pointer is returned.
\begin{description}
\sphinxlineitem{\sphinxstyleemphasis{See also}}
\sphinxAtStartPar
{\hyperref[\detokenize{Library_create:_CPPv419lammps_open_fortraniPPci}]{\sphinxcrossref{\sphinxcode{\sphinxupquote{lammps\_open\_fortran()}}}}}, {\hyperref[\detokenize{Library_create:_CPPv418lammps_open_no_mpiiPPcPPv}]{\sphinxcrossref{\sphinxcode{\sphinxupquote{lammps\_open\_no\_mpi()}}}}}

\end{description}

\begin{quote}\begin{description}
\sphinxlineitem{Parameters}\begin{itemize}
\item {} 
\sphinxAtStartPar
\sphinxstyleliteralstrong{\sphinxupquote{argc}} \textendash{} number of command line arguments 

\item {} 
\sphinxAtStartPar
\sphinxstyleliteralstrong{\sphinxupquote{argv}} \textendash{} list of command line argument strings 

\item {} 
\sphinxAtStartPar
\sphinxstyleliteralstrong{\sphinxupquote{f\_comm}} \textendash{} Fortran style MPI communicator for this LAMMPS instance 

\end{itemize}

\sphinxlineitem{Returns}
\sphinxAtStartPar
pointer to new LAMMPS instance cast to \sphinxcode{\sphinxupquote{void *}}

\end{description}\end{quote}

\end{fulllineitems}



\bigskip\hrule\bigskip

\index{lammps\_close (C++ function)@\spxentry{lammps\_close}\spxextra{C++ function}}

\begin{fulllineitems}
\phantomsection\label{\detokenize{Library_create:_CPPv412lammps_closePv}}\phantomsection\label{\detokenize{Library_create:_CPPv312lammps_closePv}}\phantomsection\label{\detokenize{Library_create:_CPPv212lammps_closePv}}\phantomsection\label{\detokenize{Library_create:lammps_close__voidP}}
\pysigstartsignatures
\pysigstartmultiline
\pysiglinewithargsret{\phantomsection\label{\detokenize{Library_create:library_8h_1a75570a623be78c153a75eb533a0b65d2}}\DUrole{kt}{void}\DUrole{w}{ }\sphinxbfcode{\sphinxupquote{\DUrole{n}{lammps\_close}}}}{\DUrole{kt}{void}\DUrole{w}{ }\DUrole{p}{*}\DUrole{n,sig-param}{handle}}{}
\pysigstopmultiline
\pysigstopsignatures
\sphinxAtStartPar
Delete a LAMMPS instance created by lammps\_open() or its variants.

\sphinxAtStartPar
This function deletes the LAMMPS class instance pointed to by \sphinxcode{\sphinxupquote{handle}}
that was created by one of the {\hyperref[\detokenize{Library_create:_CPPv411lammps_openiPPc8MPI_CommPPv}]{\sphinxcrossref{\sphinxcode{\sphinxupquote{lammps\_open()}}}}} variants.  It
does \sphinxstylestrong{not} call \sphinxcode{\sphinxupquote{MPI\_Finalize()}} to allow creating and deleting
multiple LAMMPS instances concurrently or sequentially.  See
{\hyperref[\detokenize{Library_create:_CPPv419lammps_mpi_finalizev}]{\sphinxcrossref{\sphinxcode{\sphinxupquote{lammps\_mpi\_finalize()}}}}} for a function performing this operation.
\begin{quote}\begin{description}
\sphinxlineitem{Parameters}
\sphinxAtStartPar
\sphinxstyleliteralstrong{\sphinxupquote{handle}} \textendash{} pointer to a previously created LAMMPS instance 

\end{description}\end{quote}

\end{fulllineitems}



\bigskip\hrule\bigskip

\index{lammps\_mpi\_init (C++ function)@\spxentry{lammps\_mpi\_init}\spxextra{C++ function}}

\begin{fulllineitems}
\phantomsection\label{\detokenize{Library_create:_CPPv415lammps_mpi_initv}}\phantomsection\label{\detokenize{Library_create:_CPPv315lammps_mpi_initv}}\phantomsection\label{\detokenize{Library_create:_CPPv215lammps_mpi_initv}}\phantomsection\label{\detokenize{Library_create:lammps_mpi_init}}
\pysigstartsignatures
\pysigstartmultiline
\pysiglinewithargsret{\phantomsection\label{\detokenize{Library_create:library_8h_1afee0113d4bcf1cbe55a8f83176b39b59}}\DUrole{kt}{void}\DUrole{w}{ }\sphinxbfcode{\sphinxupquote{\DUrole{n}{lammps\_mpi\_init}}}}{}{}
\pysigstopmultiline
\pysigstopsignatures
\sphinxAtStartPar
Ensure the MPI environment is initialized.

\sphinxAtStartPar

\sphinxAtStartPar
\DUrole{versionmodified,added}{Added in version 18Sep2020.}

\sphinxAtStartPar
The MPI standard requires that any MPI application must call
\sphinxcode{\sphinxupquote{MPI\_Init()}} exactly once before performing any other MPI function
calls.  This function checks, whether MPI is already initialized and
calls \sphinxcode{\sphinxupquote{MPI\_Init()}} in case it is not.


\end{fulllineitems}



\bigskip\hrule\bigskip

\index{lammps\_mpi\_finalize (C++ function)@\spxentry{lammps\_mpi\_finalize}\spxextra{C++ function}}

\begin{fulllineitems}
\phantomsection\label{\detokenize{Library_create:_CPPv419lammps_mpi_finalizev}}\phantomsection\label{\detokenize{Library_create:_CPPv319lammps_mpi_finalizev}}\phantomsection\label{\detokenize{Library_create:_CPPv219lammps_mpi_finalizev}}\phantomsection\label{\detokenize{Library_create:lammps_mpi_finalize}}
\pysigstartsignatures
\pysigstartmultiline
\pysiglinewithargsret{\phantomsection\label{\detokenize{Library_create:library_8h_1a181557f91cb6d49ba67083a51cf110dc}}\DUrole{kt}{void}\DUrole{w}{ }\sphinxbfcode{\sphinxupquote{\DUrole{n}{lammps\_mpi\_finalize}}}}{}{}
\pysigstopmultiline
\pysigstopsignatures
\sphinxAtStartPar
Shut down the MPI infrastructure.

\sphinxAtStartPar

\sphinxAtStartPar
\DUrole{versionmodified,added}{Added in version 18Sep2020.}

\sphinxAtStartPar
The MPI standard requires that any MPI application calls
\sphinxcode{\sphinxupquote{MPI\_Finalize()}} before exiting.  Even if a calling program does not
do any MPI calls, MPI is still initialized internally to avoid errors
accessing any MPI functions.  This function should then be called right
before exiting the program to wait until all (parallel) tasks are
completed and then MPI is cleanly shut down.  After calling this
function no more MPI calls may be made.
\begin{description}
\sphinxlineitem{\sphinxstyleemphasis{See also}}
\sphinxAtStartPar
{\hyperref[\detokenize{Library_create:_CPPv422lammps_kokkos_finalizev}]{\sphinxcrossref{\sphinxcode{\sphinxupquote{lammps\_kokkos\_finalize()}}}}}, {\hyperref[\detokenize{Library_create:_CPPv422lammps_python_finalizev}]{\sphinxcrossref{\sphinxcode{\sphinxupquote{lammps\_python\_finalize()}}}}}

\end{description}


\end{fulllineitems}



\bigskip\hrule\bigskip

\index{lammps\_kokkos\_finalize (C++ function)@\spxentry{lammps\_kokkos\_finalize}\spxextra{C++ function}}

\begin{fulllineitems}
\phantomsection\label{\detokenize{Library_create:_CPPv422lammps_kokkos_finalizev}}\phantomsection\label{\detokenize{Library_create:_CPPv322lammps_kokkos_finalizev}}\phantomsection\label{\detokenize{Library_create:_CPPv222lammps_kokkos_finalizev}}\phantomsection\label{\detokenize{Library_create:lammps_kokkos_finalize}}
\pysigstartsignatures
\pysigstartmultiline
\pysiglinewithargsret{\phantomsection\label{\detokenize{Library_create:library_8h_1a295c98b87692b72b5bf594e7b291905b}}\DUrole{kt}{void}\DUrole{w}{ }\sphinxbfcode{\sphinxupquote{\DUrole{n}{lammps\_kokkos\_finalize}}}}{}{}
\pysigstopmultiline
\pysigstopsignatures
\sphinxAtStartPar
Shut down the Kokkos library environment.

\sphinxAtStartPar

\sphinxAtStartPar
\DUrole{versionmodified,added}{Added in version 2Jul2021.}

\sphinxAtStartPar
The Kokkos library may only be initialized once during the execution of
a process.  This is done automatically the first time Kokkos
functionality is used.  This requires that the Kokkos environment
must be explicitly shut down after any LAMMPS instance using it is
closed (to release associated resources).
After calling this function no Kokkos functionality may be used.
\begin{description}
\sphinxlineitem{\sphinxstyleemphasis{See also}}
\sphinxAtStartPar
{\hyperref[\detokenize{Library_create:_CPPv419lammps_mpi_finalizev}]{\sphinxcrossref{\sphinxcode{\sphinxupquote{lammps\_mpi\_finalize()}}}}}, {\hyperref[\detokenize{Library_create:_CPPv422lammps_python_finalizev}]{\sphinxcrossref{\sphinxcode{\sphinxupquote{lammps\_python\_finalize()}}}}}

\end{description}


\end{fulllineitems}



\bigskip\hrule\bigskip

\index{lammps\_python\_finalize (C++ function)@\spxentry{lammps\_python\_finalize}\spxextra{C++ function}}

\begin{fulllineitems}
\phantomsection\label{\detokenize{Library_create:_CPPv422lammps_python_finalizev}}\phantomsection\label{\detokenize{Library_create:_CPPv322lammps_python_finalizev}}\phantomsection\label{\detokenize{Library_create:_CPPv222lammps_python_finalizev}}\phantomsection\label{\detokenize{Library_create:lammps_python_finalize}}
\pysigstartsignatures
\pysigstartmultiline
\pysiglinewithargsret{\phantomsection\label{\detokenize{Library_create:library_8h_1a5d625667e3988ecae35e94490cb5aa8f}}\DUrole{kt}{void}\DUrole{w}{ }\sphinxbfcode{\sphinxupquote{\DUrole{n}{lammps\_python\_finalize}}}}{}{}
\pysigstopmultiline
\pysigstopsignatures
\sphinxAtStartPar
Clear the embedded Python environment

\sphinxAtStartPar

\sphinxAtStartPar
\DUrole{versionmodified,added}{Added in version 20Sep2021.}

\sphinxAtStartPar
This function resets and clears an embedded Python environment
by calling the \sphinxhref{https://docs.python.org/3/c-api/init.html\#c.Py\_FinalizeEx}{Py\_Finalize() function}
of the embedded Python library, if enabled.
This call would free up all allocated resources and release
loaded shared objects.

\sphinxAtStartPar
However, this is \sphinxstylestrong{not} done when a LAMMPS instance is deleted because
a) LAMMPS may have been used through the Python module and thus
the Python interpreter is external and not embedded into LAMMPS
and therefore may not be reset by LAMMPS b) some Python modules
and extensions, most notably NumPy, are not compatible with being
initialized multiple times, which would happen if additional
LAMMPS instances using Python would be created \sphinxstyleemphasis{after}
after calling Py\_Finalize().

\sphinxAtStartPar
This function can be called to explicitly clear the Python
environment in case it is safe to do so.
\begin{description}
\sphinxlineitem{\sphinxstyleemphasis{See also}}
\sphinxAtStartPar
{\hyperref[\detokenize{Library_create:_CPPv419lammps_mpi_finalizev}]{\sphinxcrossref{\sphinxcode{\sphinxupquote{lammps\_mpi\_finalize()}}}}}, {\hyperref[\detokenize{Library_create:_CPPv422lammps_kokkos_finalizev}]{\sphinxcrossref{\sphinxcode{\sphinxupquote{lammps\_kokkos\_finalize()}}}}}

\end{description}


\end{fulllineitems}



\bigskip\hrule\bigskip

\index{lammps\_error (C++ function)@\spxentry{lammps\_error}\spxextra{C++ function}}

\begin{fulllineitems}
\phantomsection\label{\detokenize{Library_create:_CPPv412lammps_errorPviPKc}}\phantomsection\label{\detokenize{Library_create:_CPPv312lammps_errorPviPKc}}\phantomsection\label{\detokenize{Library_create:_CPPv212lammps_errorPviPKc}}\phantomsection\label{\detokenize{Library_create:lammps_error__voidP.i.cCP}}
\pysigstartsignatures
\pysigstartmultiline
\pysiglinewithargsret{\phantomsection\label{\detokenize{Library_create:library_8h_1ad365a1950c1e4e6bf21e414cab86a57a}}\DUrole{kt}{void}\DUrole{w}{ }\sphinxbfcode{\sphinxupquote{\DUrole{n}{lammps\_error}}}}{\DUrole{kt}{void}\DUrole{w}{ }\DUrole{p}{*}\DUrole{n,sig-param}{handle}\sphinxparamcomma \DUrole{kt}{int}\DUrole{w}{ }\DUrole{n,sig-param}{error\_type}\sphinxparamcomma \DUrole{k}{const}\DUrole{w}{ }\DUrole{kt}{char}\DUrole{w}{ }\DUrole{p}{*}\DUrole{n,sig-param}{error\_text}}{}
\pysigstopmultiline
\pysigstopsignatures
\sphinxAtStartPar
Call a LAMMPS Error class function

\sphinxAtStartPar

\sphinxAtStartPar
\DUrole{versionmodified,added}{Added in version 3Nov2022.}

\sphinxAtStartPar
This function is a wrapper around functions in the \sphinxcode{\sphinxupquote{Error}} to print an
error message and then stop LAMMPS.

\sphinxAtStartPar
The \sphinxstyleemphasis{error\_type} parameter selects which function to call.  It is a sum
of constants from \sphinxcode{\sphinxupquote{\_LMP\_ERROR\_CONST}}.  If the value does not
match any valid combination of constants a warning is printed and the
function returns.

\begin{quote}\begin{description}
\sphinxlineitem{Parameters}\begin{itemize}
\item {} 
\sphinxAtStartPar
\sphinxstyleliteralstrong{\sphinxupquote{handle}} \textendash{} pointer to a previously created LAMMPS instance 

\item {} 
\sphinxAtStartPar
\sphinxstyleliteralstrong{\sphinxupquote{error\_type}} \textendash{} parameter to select function in the Error class 

\item {} 
\sphinxAtStartPar
\sphinxstyleliteralstrong{\sphinxupquote{error\_text}} \textendash{} error message 

\end{itemize}

\end{description}\end{quote}

\end{fulllineitems}


\sphinxstepscope


\subsection{Executing commands}
\label{\detokenize{Library_execute:executing-commands}}\label{\detokenize{Library_execute::doc}}
\sphinxAtStartPar
This section documents the following functions:
\begin{itemize}
\item {} 
\sphinxAtStartPar
{\hyperref[\detokenize{Library_execute:_CPPv411lammps_filePvPKc}]{\sphinxcrossref{\sphinxcode{\sphinxupquote{lammps\_file()}}}}}

\item {} 
\sphinxAtStartPar
{\hyperref[\detokenize{Library_execute:_CPPv414lammps_commandPvPKc}]{\sphinxcrossref{\sphinxcode{\sphinxupquote{lammps\_command()}}}}}

\item {} 
\sphinxAtStartPar
{\hyperref[\detokenize{Library_execute:_CPPv420lammps_commands_listPviPPKc}]{\sphinxcrossref{\sphinxcode{\sphinxupquote{lammps\_commands\_list()}}}}}

\item {} 
\sphinxAtStartPar
{\hyperref[\detokenize{Library_execute:_CPPv422lammps_commands_stringPvPKc}]{\sphinxcrossref{\sphinxcode{\sphinxupquote{lammps\_commands\_string()}}}}}

\end{itemize}


\bigskip\hrule\bigskip


\sphinxAtStartPar
Once a LAMMPS instance is created, there are multiple ways to “drive” a
simulation.  In most cases it is easiest to process single or multiple
LAMMPS commands like in an input file.  This can be done through reading
a file or passing single commands or lists of commands or blocks of
commands with the following functions.

\sphinxAtStartPar
Via these functions, the calling code can have LAMMPS act on a series
of {\hyperref[\detokenize{Commands_all::doc}]{\sphinxcrossref{\DUrole{doc}{input file commands}}}} that are either read from
a file or passed as strings.  For example, this allows setup of a
problem from an input script, and then running it in stages while
performing other operations in between or concurrently.  The caller
can interleave the LAMMPS function calls with operations it performs,
such as calls to extract information from or set information within
LAMMPS, or calls to another code’s library.

\sphinxAtStartPar
Just as with {\hyperref[\detokenize{Commands_parse::doc}]{\sphinxcrossref{\DUrole{doc}{input script parsing}}}} comments can
be included in the file or strings, and expansion of variables with
\sphinxcode{\sphinxupquote{\$\{name\}}} or \sphinxcode{\sphinxupquote{\$(expression)}} syntax is performed.
Below is a short example using some of these functions.

\begin{sphinxVerbatim}[commandchars=\\\{\}]
\PYG{c+cm}{/* define to make the otherwise hidden prototype for \PYGZdq{}lammps\PYGZus{}open()\PYGZdq{} visible */}
\PYG{c+cp}{\PYGZsh{}}\PYG{c+cp}{define LAMMPS\PYGZus{}LIB\PYGZus{}MPI}
\PYG{c+cp}{\PYGZsh{}}\PYG{c+cp}{include}\PYG{+w}{ }\PYG{c+cpf}{\PYGZdq{}library.h\PYGZdq{}}

\PYG{c+cp}{\PYGZsh{}}\PYG{c+cp}{include}\PYG{+w}{ }\PYG{c+cpf}{\PYGZlt{}mpi.h\PYGZgt{}}
\PYG{c+cp}{\PYGZsh{}}\PYG{c+cp}{include}\PYG{+w}{ }\PYG{c+cpf}{\PYGZlt{}stdio.h\PYGZgt{}}

\PYG{k+kt}{int}\PYG{+w}{ }\PYG{n+nf}{main}\PYG{p}{(}\PYG{k+kt}{int}\PYG{+w}{ }\PYG{n}{argc}\PYG{p}{,}\PYG{+w}{ }\PYG{k+kt}{char}\PYG{+w}{ }\PYG{o}{*}\PYG{o}{*}\PYG{n}{argv}\PYG{p}{)}
\PYG{p}{\PYGZob{}}
\PYG{+w}{  }\PYG{k+kt}{void}\PYG{+w}{ }\PYG{o}{*}\PYG{n}{handle}\PYG{p}{;}
\PYG{+w}{  }\PYG{k+kt}{int}\PYG{+w}{ }\PYG{n}{i}\PYG{p}{;}

\PYG{+w}{  }\PYG{n}{MPI\PYGZus{}Init}\PYG{p}{(}\PYG{o}{\PYGZam{}}\PYG{n}{argc}\PYG{p}{,}\PYG{+w}{ }\PYG{o}{\PYGZam{}}\PYG{n}{argv}\PYG{p}{)}\PYG{p}{;}
\PYG{+w}{  }\PYG{n}{handle}\PYG{+w}{ }\PYG{o}{=}\PYG{+w}{ }\PYG{n}{lammps\PYGZus{}open}\PYG{p}{(}\PYG{l+m+mi}{0}\PYG{p}{,}\PYG{+w}{ }\PYG{n+nb}{NULL}\PYG{p}{,}\PYG{+w}{ }\PYG{n}{MPI\PYGZus{}COMM\PYGZus{}WORLD}\PYG{p}{,}\PYG{+w}{ }\PYG{n+nb}{NULL}\PYG{p}{)}\PYG{p}{;}
\PYG{+w}{  }\PYG{n}{lammps\PYGZus{}file}\PYG{p}{(}\PYG{n}{handle}\PYG{p}{,}\PYG{l+s}{\PYGZdq{}}\PYG{l+s}{in.sysinit}\PYG{l+s}{\PYGZdq{}}\PYG{p}{)}\PYG{p}{;}
\PYG{+w}{  }\PYG{n}{lammps\PYGZus{}command}\PYG{p}{(}\PYG{n}{handle}\PYG{p}{,}\PYG{l+s}{\PYGZdq{}}\PYG{l+s}{run 1000 post no}\PYG{l+s}{\PYGZdq{}}\PYG{p}{)}\PYG{p}{;}

\PYG{+w}{  }\PYG{k}{for}\PYG{+w}{ }\PYG{p}{(}\PYG{n}{i}\PYG{o}{=}\PYG{l+m+mi}{0}\PYG{p}{;}\PYG{+w}{ }\PYG{n}{i}\PYG{+w}{ }\PYG{o}{\PYGZlt{}}\PYG{+w}{ }\PYG{l+m+mi}{100}\PYG{p}{;}\PYG{+w}{ }\PYG{o}{+}\PYG{o}{+}\PYG{n}{i}\PYG{p}{)}\PYG{+w}{ }\PYG{p}{\PYGZob{}}
\PYG{+w}{    }\PYG{n}{lammps\PYGZus{}commands\PYGZus{}string}\PYG{p}{(}\PYG{n}{handle}\PYG{p}{,}\PYG{l+s}{\PYGZdq{}}\PYG{l+s}{run 100 pre no post no}\PYG{l+s+se}{\PYGZbs{}n}\PYG{l+s}{\PYGZdq{}}
\PYG{+w}{                                  }\PYG{l+s}{\PYGZdq{}}\PYG{l+s}{print \PYGZsq{}PE = \PYGZdl{}(pe)\PYGZsq{}}\PYG{l+s+se}{\PYGZbs{}n}\PYG{l+s}{\PYGZdq{}}
\PYG{+w}{                                  }\PYG{l+s}{\PYGZdq{}}\PYG{l+s}{print \PYGZsq{}KE = \PYGZdl{}(ke)\PYGZsq{}}\PYG{l+s+se}{\PYGZbs{}n}\PYG{l+s}{\PYGZdq{}}\PYG{p}{)}\PYG{p}{;}
\PYG{+w}{  }\PYG{p}{\PYGZcb{}}
\PYG{+w}{  }\PYG{n}{lammps\PYGZus{}close}\PYG{p}{(}\PYG{n}{handle}\PYG{p}{)}\PYG{p}{;}
\PYG{+w}{  }\PYG{n}{MPI\PYGZus{}Finalize}\PYG{p}{(}\PYG{p}{)}\PYG{p}{;}
\PYG{+w}{  }\PYG{k}{return}\PYG{+w}{ }\PYG{l+m+mi}{0}\PYG{p}{;}
\PYG{p}{\PYGZcb{}}
\end{sphinxVerbatim}


\bigskip\hrule\bigskip

\index{lammps\_file (C++ function)@\spxentry{lammps\_file}\spxextra{C++ function}}

\begin{fulllineitems}
\phantomsection\label{\detokenize{Library_execute:_CPPv411lammps_filePvPKc}}\phantomsection\label{\detokenize{Library_execute:_CPPv311lammps_filePvPKc}}\phantomsection\label{\detokenize{Library_execute:_CPPv211lammps_filePvPKc}}\phantomsection\label{\detokenize{Library_execute:lammps_file__voidP.cCP}}
\pysigstartsignatures
\pysigstartmultiline
\pysiglinewithargsret{\phantomsection\label{\detokenize{Library_execute:library_8h_1a76c1ce1af98c09978578c7a94d778c0b}}\DUrole{kt}{void}\DUrole{w}{ }\sphinxbfcode{\sphinxupquote{\DUrole{n}{lammps\_file}}}}{\DUrole{kt}{void}\DUrole{w}{ }\DUrole{p}{*}\DUrole{n,sig-param}{handle}\sphinxparamcomma \DUrole{k}{const}\DUrole{w}{ }\DUrole{kt}{char}\DUrole{w}{ }\DUrole{p}{*}\DUrole{n,sig-param}{file}}{}
\pysigstopmultiline
\pysigstopsignatures
\sphinxAtStartPar
Process LAMMPS input from a file.

\sphinxAtStartPar

\sphinxAtStartPar
This function processes commands in the file pointed to by \sphinxstyleemphasis{filename}
line by line and thus functions very similar to the {\hyperref[\detokenize{include::doc}]{\sphinxcrossref{\DUrole{doc}{include}}}} command. The function returns when the end of the file is
reached and the commands have completed.

\sphinxAtStartPar
The actual work is done by the functions
{\hyperref[\detokenize{Classes_input:_CPPv4N9LAMMPS_NS5Input4fileEPKc}]{\sphinxcrossref{\sphinxcode{\sphinxupquote{Input::file(const char *)}}}}}
and {\hyperref[\detokenize{Classes_input:_CPPv4N9LAMMPS_NS5Input4fileEv}]{\sphinxcrossref{\sphinxcode{\sphinxupquote{Input::file()}}}}}.

\begin{quote}\begin{description}
\sphinxlineitem{Parameters}\begin{itemize}
\item {} 
\sphinxAtStartPar
\sphinxstyleliteralstrong{\sphinxupquote{handle}} \textendash{} pointer to a previously created LAMMPS instance 

\item {} 
\sphinxAtStartPar
\sphinxstyleliteralstrong{\sphinxupquote{filename}} \textendash{} name of a file with LAMMPS input 

\end{itemize}

\end{description}\end{quote}

\end{fulllineitems}



\bigskip\hrule\bigskip

\index{lammps\_command (C++ function)@\spxentry{lammps\_command}\spxextra{C++ function}}

\begin{fulllineitems}
\phantomsection\label{\detokenize{Library_execute:_CPPv414lammps_commandPvPKc}}\phantomsection\label{\detokenize{Library_execute:_CPPv314lammps_commandPvPKc}}\phantomsection\label{\detokenize{Library_execute:_CPPv214lammps_commandPvPKc}}\phantomsection\label{\detokenize{Library_execute:lammps_command__voidP.cCP}}
\pysigstartsignatures
\pysigstartmultiline
\pysiglinewithargsret{\phantomsection\label{\detokenize{Library_execute:library_8h_1a81fbef297272938327713bda40d40e0e}}\DUrole{kt}{char}\DUrole{w}{ }\DUrole{p}{*}\sphinxbfcode{\sphinxupquote{\DUrole{n}{lammps\_command}}}}{\DUrole{kt}{void}\DUrole{w}{ }\DUrole{p}{*}\DUrole{n,sig-param}{handle}\sphinxparamcomma \DUrole{k}{const}\DUrole{w}{ }\DUrole{kt}{char}\DUrole{w}{ }\DUrole{p}{*}\DUrole{n,sig-param}{cmd}}{}
\pysigstopmultiline
\pysigstopsignatures
\sphinxAtStartPar
Process a single LAMMPS input command from a string.

\sphinxAtStartPar

\sphinxAtStartPar
This function tells LAMMPS to execute the single command in the string
\sphinxstyleemphasis{cmd}.  The entire string is considered as command and need not have a
(final) newline character.  Newline characters in the body of the
string, however, will be treated as part of the command and will \sphinxstylestrong{not}
start a second command.  The function {\hyperref[\detokenize{Library_execute:_CPPv422lammps_commands_stringPvPKc}]{\sphinxcrossref{\sphinxcode{\sphinxupquote{lammps\_commands\_string()}}}}}
processes a string with multiple command lines.

\sphinxAtStartPar
The function returns the name of the command on success or \sphinxcode{\sphinxupquote{NULL}} when
passing a string without a command.

\begin{quote}\begin{description}
\sphinxlineitem{Parameters}\begin{itemize}
\item {} 
\sphinxAtStartPar
\sphinxstyleliteralstrong{\sphinxupquote{handle}} \textendash{} pointer to a previously created LAMMPS instance 

\item {} 
\sphinxAtStartPar
\sphinxstyleliteralstrong{\sphinxupquote{cmd}} \textendash{} string with a single LAMMPS command 

\end{itemize}

\sphinxlineitem{Returns}
\sphinxAtStartPar
string with parsed command name or \sphinxcode{\sphinxupquote{NULL}}

\end{description}\end{quote}

\end{fulllineitems}



\bigskip\hrule\bigskip

\index{lammps\_commands\_list (C++ function)@\spxentry{lammps\_commands\_list}\spxextra{C++ function}}

\begin{fulllineitems}
\phantomsection\label{\detokenize{Library_execute:_CPPv420lammps_commands_listPviPPKc}}\phantomsection\label{\detokenize{Library_execute:_CPPv320lammps_commands_listPviPPKc}}\phantomsection\label{\detokenize{Library_execute:_CPPv220lammps_commands_listPviPPKc}}\phantomsection\label{\detokenize{Library_execute:lammps_commands_list__voidP.i.cCPP}}
\pysigstartsignatures
\pysigstartmultiline
\pysiglinewithargsret{\phantomsection\label{\detokenize{Library_execute:library_8h_1ab542f51c42bcf8e39b64722ffa426890}}\DUrole{kt}{void}\DUrole{w}{ }\sphinxbfcode{\sphinxupquote{\DUrole{n}{lammps\_commands\_list}}}}{\DUrole{kt}{void}\DUrole{w}{ }\DUrole{p}{*}\DUrole{n,sig-param}{handle}\sphinxparamcomma \DUrole{kt}{int}\DUrole{w}{ }\DUrole{n,sig-param}{ncmd}\sphinxparamcomma \DUrole{k}{const}\DUrole{w}{ }\DUrole{kt}{char}\DUrole{w}{ }\DUrole{p}{*}\DUrole{p}{*}\DUrole{n,sig-param}{cmds}}{}
\pysigstopmultiline
\pysigstopsignatures
\sphinxAtStartPar
Process multiple LAMMPS input commands from list of strings.

\sphinxAtStartPar
This function processes multiple commands from a list of strings by
first concatenating the individual strings in \sphinxstyleemphasis{cmds} into a single
string, inserting newline characters as needed.  The combined string
is passed to {\hyperref[\detokenize{Library_execute:_CPPv422lammps_commands_stringPvPKc}]{\sphinxcrossref{\sphinxcode{\sphinxupquote{lammps\_commands\_string()}}}}} for processing.
\begin{quote}\begin{description}
\sphinxlineitem{Parameters}\begin{itemize}
\item {} 
\sphinxAtStartPar
\sphinxstyleliteralstrong{\sphinxupquote{handle}} \textendash{} pointer to a previously created LAMMPS instance 

\item {} 
\sphinxAtStartPar
\sphinxstyleliteralstrong{\sphinxupquote{ncmd}} \textendash{} number of lines in \sphinxstyleemphasis{cmds}

\item {} 
\sphinxAtStartPar
\sphinxstyleliteralstrong{\sphinxupquote{cmds}} \textendash{} list of strings with LAMMPS commands 

\end{itemize}

\end{description}\end{quote}

\end{fulllineitems}



\bigskip\hrule\bigskip

\index{lammps\_commands\_string (C++ function)@\spxentry{lammps\_commands\_string}\spxextra{C++ function}}

\begin{fulllineitems}
\phantomsection\label{\detokenize{Library_execute:_CPPv422lammps_commands_stringPvPKc}}\phantomsection\label{\detokenize{Library_execute:_CPPv322lammps_commands_stringPvPKc}}\phantomsection\label{\detokenize{Library_execute:_CPPv222lammps_commands_stringPvPKc}}\phantomsection\label{\detokenize{Library_execute:lammps_commands_string__voidP.cCP}}
\pysigstartsignatures
\pysigstartmultiline
\pysiglinewithargsret{\phantomsection\label{\detokenize{Library_execute:library_8h_1a82b8d3fb9db7e0a834f6da0996e61d6b}}\DUrole{kt}{void}\DUrole{w}{ }\sphinxbfcode{\sphinxupquote{\DUrole{n}{lammps\_commands\_string}}}}{\DUrole{kt}{void}\DUrole{w}{ }\DUrole{p}{*}\DUrole{n,sig-param}{handle}\sphinxparamcomma \DUrole{k}{const}\DUrole{w}{ }\DUrole{kt}{char}\DUrole{w}{ }\DUrole{p}{*}\DUrole{n,sig-param}{str}}{}
\pysigstopmultiline
\pysigstopsignatures
\sphinxAtStartPar
Process a block of LAMMPS input commands from a single string.

\sphinxAtStartPar

\sphinxAtStartPar
This function processes a multi\sphinxhyphen{}line string similar to a block of
commands from a file.  The string may have multiple lines (separated by
newline characters) and also single commands may be distributed over
multiple lines with continuation characters (’\&’).  Those lines are
combined by removing the ‘\&’ and the following newline character.  After
this processing the string is handed to LAMMPS for parsing and
executing.

\sphinxAtStartPar
\DUrole{versionmodified,added}{Added in version 21Nov2023: }The command is now able to process long strings with triple quotes and
loops using {\hyperref[\detokenize{jump::doc}]{\sphinxcrossref{\DUrole{doc}{jump SELF \textless{}label\textgreater{}}}}}.

\begin{quote}\begin{description}
\sphinxlineitem{Parameters}\begin{itemize}
\item {} 
\sphinxAtStartPar
\sphinxstyleliteralstrong{\sphinxupquote{handle}} \textendash{} pointer to a previously created LAMMPS instance 

\item {} 
\sphinxAtStartPar
\sphinxstyleliteralstrong{\sphinxupquote{str}} \textendash{} string with block of LAMMPS input commands 

\end{itemize}

\end{description}\end{quote}

\end{fulllineitems}


\sphinxstepscope


\subsection{System properties}
\label{\detokenize{Library_properties:system-properties}}\label{\detokenize{Library_properties::doc}}
\sphinxAtStartPar
This section documents the following functions:
\begin{itemize}
\item {} 
\sphinxAtStartPar
{\hyperref[\detokenize{Library_properties:_CPPv417lammps_get_natomsPv}]{\sphinxcrossref{\sphinxcode{\sphinxupquote{lammps\_get\_natoms()}}}}}

\item {} 
\sphinxAtStartPar
{\hyperref[\detokenize{Library_properties:_CPPv417lammps_get_thermoPvPKc}]{\sphinxcrossref{\sphinxcode{\sphinxupquote{lammps\_get\_thermo()}}}}}

\item {} 
\sphinxAtStartPar
{\hyperref[\detokenize{Library_properties:_CPPv418lammps_last_thermoPvPKci}]{\sphinxcrossref{\sphinxcode{\sphinxupquote{lammps\_last\_thermo()}}}}}

\item {} 
\sphinxAtStartPar
{\hyperref[\detokenize{Library_properties:_CPPv418lammps_extract_boxPvPdPdPdPdPdPiPi}]{\sphinxcrossref{\sphinxcode{\sphinxupquote{lammps\_extract\_box()}}}}}

\item {} 
\sphinxAtStartPar
{\hyperref[\detokenize{Library_properties:_CPPv416lammps_reset_boxPvPdPdddd}]{\sphinxcrossref{\sphinxcode{\sphinxupquote{lammps\_reset\_box()}}}}}

\item {} 
\sphinxAtStartPar
{\hyperref[\detokenize{Library_properties:_CPPv419lammps_memory_usagePvPd}]{\sphinxcrossref{\sphinxcode{\sphinxupquote{lammps\_memory\_usage()}}}}}

\item {} 
\sphinxAtStartPar
{\hyperref[\detokenize{Library_properties:_CPPv419lammps_get_mpi_commPv}]{\sphinxcrossref{\sphinxcode{\sphinxupquote{lammps\_get\_mpi\_comm()}}}}}

\item {} 
\sphinxAtStartPar
{\hyperref[\detokenize{Library_properties:_CPPv422lammps_extract_settingPvPKc}]{\sphinxcrossref{\sphinxcode{\sphinxupquote{lammps\_extract\_setting()}}}}}

\item {} 
\sphinxAtStartPar
{\hyperref[\detokenize{Library_properties:_CPPv430lammps_extract_global_datatypePvPKc}]{\sphinxcrossref{\sphinxcode{\sphinxupquote{lammps\_extract\_global\_datatype()}}}}}

\item {} 
\sphinxAtStartPar
{\hyperref[\detokenize{Library_properties:_CPPv421lammps_extract_globalPvPKc}]{\sphinxcrossref{\sphinxcode{\sphinxupquote{lammps\_extract\_global()}}}}}

\item {} 
\sphinxAtStartPar
{\hyperref[\detokenize{Library_properties:_CPPv429lammps_extract_pair_dimensionPvPKc}]{\sphinxcrossref{\sphinxcode{\sphinxupquote{lammps\_extract\_pair\_dimension()}}}}}

\item {} 
\sphinxAtStartPar
{\hyperref[\detokenize{Library_properties:_CPPv419lammps_extract_pairPvPKc}]{\sphinxcrossref{\sphinxcode{\sphinxupquote{lammps\_extract\_pair()}}}}}

\item {} 
\sphinxAtStartPar
{\hyperref[\detokenize{Library_properties:_CPPv415lammps_map_atomPvPKv}]{\sphinxcrossref{\sphinxcode{\sphinxupquote{lammps\_map\_atom()}}}}}

\end{itemize}


\bigskip\hrule\bigskip


\sphinxAtStartPar
The library interface allows the extraction of different kinds of
information about the active simulation instance and also \sphinxhyphen{} in some
cases \sphinxhyphen{} to apply modifications to it.  This enables combining of a
LAMMPS simulation with other processing and simulation methods computed
by the calling code, or by another code that is coupled to LAMMPS via
the library interface.  In some cases the data returned is direct
reference to the original data inside LAMMPS, cast to a void pointer.
In that case the data needs to be cast to a suitable pointer for the
calling program to access it, and you may need to know the correct
dimensions and lengths.  This also means you can directly change those
value(s) from the calling program (e.g., to modify atom positions).  Of
course, changing values should be done with care.  When accessing per\sphinxhyphen{}atom
data, please note that these data are the per\sphinxhyphen{}processor \sphinxstylestrong{local} data and are
indexed accordingly. Per\sphinxhyphen{}atom data can change sizes and ordering at
every neighbor list rebuild or atom sort event as atoms migrate between
subdomains and processors.

\begin{sphinxVerbatim}[commandchars=\\\{\}]
\PYG{c+cp}{\PYGZsh{}}\PYG{c+cp}{include}\PYG{+w}{ }\PYG{c+cpf}{\PYGZdq{}library.h\PYGZdq{}}
\PYG{c+cp}{\PYGZsh{}}\PYG{c+cp}{include}\PYG{+w}{ }\PYG{c+cpf}{\PYGZlt{}stdio.h\PYGZgt{}}

\PYG{k+kt}{int}\PYG{+w}{ }\PYG{n+nf}{main}\PYG{p}{(}\PYG{k+kt}{int}\PYG{+w}{ }\PYG{n}{argc}\PYG{p}{,}\PYG{+w}{ }\PYG{k+kt}{char}\PYG{+w}{ }\PYG{o}{*}\PYG{o}{*}\PYG{n}{argv}\PYG{p}{)}
\PYG{p}{\PYGZob{}}
\PYG{+w}{  }\PYG{k+kt}{void}\PYG{+w}{ }\PYG{o}{*}\PYG{n}{handle}\PYG{p}{;}
\PYG{+w}{  }\PYG{k+kt}{int}\PYG{+w}{ }\PYG{n}{i}\PYG{p}{;}

\PYG{+w}{  }\PYG{n}{handle}\PYG{+w}{ }\PYG{o}{=}\PYG{+w}{ }\PYG{n}{lammps\PYGZus{}open\PYGZus{}no\PYGZus{}mpi}\PYG{p}{(}\PYG{l+m+mi}{0}\PYG{p}{,}\PYG{+w}{ }\PYG{n+nb}{NULL}\PYG{p}{,}\PYG{+w}{ }\PYG{n+nb}{NULL}\PYG{p}{)}\PYG{p}{;}
\PYG{+w}{  }\PYG{n}{lammps\PYGZus{}file}\PYG{p}{(}\PYG{n}{handle}\PYG{p}{,}\PYG{l+s}{\PYGZdq{}}\PYG{l+s}{in.sysinit}\PYG{l+s}{\PYGZdq{}}\PYG{p}{)}\PYG{p}{;}
\PYG{+w}{  }\PYG{n}{printf}\PYG{p}{(}\PYG{l+s}{\PYGZdq{}}\PYG{l+s}{Running a simulation with \PYGZpc{}g atoms.}\PYG{l+s+se}{\PYGZbs{}n}\PYG{l+s}{\PYGZdq{}}\PYG{p}{,}
\PYG{+w}{         }\PYG{n}{lammps\PYGZus{}get\PYGZus{}natoms}\PYG{p}{(}\PYG{n}{handle}\PYG{p}{)}\PYG{p}{)}\PYG{p}{;}

\PYG{+w}{  }\PYG{n}{printf}\PYG{p}{(}\PYG{l+s}{\PYGZdq{}}\PYG{l+s}{ \PYGZpc{}d local and \PYGZpc{}d ghost atoms. \PYGZpc{}d atom types}\PYG{l+s+se}{\PYGZbs{}n}\PYG{l+s}{\PYGZdq{}}\PYG{p}{,}
\PYG{+w}{         }\PYG{n}{lammps\PYGZus{}extract\PYGZus{}setting}\PYG{p}{(}\PYG{n}{handle}\PYG{p}{,}\PYG{l+s}{\PYGZdq{}}\PYG{l+s}{nlocal}\PYG{l+s}{\PYGZdq{}}\PYG{p}{)}\PYG{p}{,}
\PYG{+w}{         }\PYG{n}{lammps\PYGZus{}extract\PYGZus{}setting}\PYG{p}{(}\PYG{n}{handle}\PYG{p}{,}\PYG{l+s}{\PYGZdq{}}\PYG{l+s}{nghost}\PYG{l+s}{\PYGZdq{}}\PYG{p}{)}\PYG{p}{,}
\PYG{+w}{         }\PYG{n}{lammps\PYGZus{}extract\PYGZus{}setting}\PYG{p}{(}\PYG{n}{handle}\PYG{p}{,}\PYG{l+s}{\PYGZdq{}}\PYG{l+s}{ntypes}\PYG{l+s}{\PYGZdq{}}\PYG{p}{)}\PYG{p}{)}\PYG{p}{;}

\PYG{+w}{  }\PYG{k+kt}{double}\PYG{+w}{  }\PYG{o}{*}\PYG{n}{dt}\PYG{+w}{ }\PYG{o}{=}\PYG{+w}{ }\PYG{p}{(}\PYG{k+kt}{double}\PYG{+w}{ }\PYG{o}{*}\PYG{p}{)}\PYG{n}{lammps\PYGZus{}extract\PYGZus{}global}\PYG{p}{(}\PYG{n}{handle}\PYG{p}{,}\PYG{l+s}{\PYGZdq{}}\PYG{l+s}{dt}\PYG{l+s}{\PYGZdq{}}\PYG{p}{)}\PYG{p}{;}
\PYG{+w}{  }\PYG{n}{printf}\PYG{p}{(}\PYG{l+s}{\PYGZdq{}}\PYG{l+s}{Changing timestep from \PYGZpc{}g to 0.5}\PYG{l+s+se}{\PYGZbs{}n}\PYG{l+s}{\PYGZdq{}}\PYG{p}{,}\PYG{+w}{ }\PYG{o}{*}\PYG{n}{dt}\PYG{p}{)}\PYG{p}{;}
\PYG{+w}{  }\PYG{o}{*}\PYG{n}{dt}\PYG{+w}{ }\PYG{o}{=}\PYG{+w}{ }\PYG{l+m+mf}{0.5}\PYG{p}{;}

\PYG{+w}{  }\PYG{n}{lammps\PYGZus{}command}\PYG{p}{(}\PYG{n}{handle}\PYG{p}{,}\PYG{l+s}{\PYGZdq{}}\PYG{l+s}{run 1000 post no}\PYG{l+s}{\PYGZdq{}}\PYG{p}{)}\PYG{p}{;}

\PYG{+w}{  }\PYG{k}{for}\PYG{+w}{ }\PYG{p}{(}\PYG{n}{i}\PYG{o}{=}\PYG{l+m+mi}{0}\PYG{p}{;}\PYG{+w}{ }\PYG{n}{i}\PYG{+w}{ }\PYG{o}{\PYGZlt{}}\PYG{+w}{ }\PYG{l+m+mi}{10}\PYG{p}{;}\PYG{+w}{ }\PYG{o}{+}\PYG{o}{+}\PYG{n}{i}\PYG{p}{)}\PYG{+w}{ }\PYG{p}{\PYGZob{}}
\PYG{+w}{    }\PYG{n}{lammps\PYGZus{}command}\PYG{p}{(}\PYG{n}{handle}\PYG{p}{,}\PYG{l+s}{\PYGZdq{}}\PYG{l+s}{run 100 pre no post no}\PYG{l+s}{\PYGZdq{}}\PYG{p}{)}\PYG{p}{;}
\PYG{+w}{    }\PYG{n}{printf}\PYG{p}{(}\PYG{l+s}{\PYGZdq{}}\PYG{l+s}{PE = \PYGZpc{}g}\PYG{l+s+se}{\PYGZbs{}n}\PYG{l+s}{KE = \PYGZpc{}g}\PYG{l+s+se}{\PYGZbs{}n}\PYG{l+s}{\PYGZdq{}}\PYG{p}{,}
\PYG{+w}{           }\PYG{n}{lammps\PYGZus{}get\PYGZus{}thermo}\PYG{p}{(}\PYG{n}{handle}\PYG{p}{,}\PYG{l+s}{\PYGZdq{}}\PYG{l+s}{pe}\PYG{l+s}{\PYGZdq{}}\PYG{p}{)}\PYG{p}{,}
\PYG{+w}{           }\PYG{n}{lammps\PYGZus{}get\PYGZus{}thermo}\PYG{p}{(}\PYG{n}{handle}\PYG{p}{,}\PYG{l+s}{\PYGZdq{}}\PYG{l+s}{ke}\PYG{l+s}{\PYGZdq{}}\PYG{p}{)}\PYG{p}{)}\PYG{p}{;}
\PYG{+w}{  }\PYG{p}{\PYGZcb{}}
\PYG{+w}{  }\PYG{n}{lammps\PYGZus{}close}\PYG{p}{(}\PYG{n}{handle}\PYG{p}{)}\PYG{p}{;}
\PYG{+w}{  }\PYG{k}{return}\PYG{+w}{ }\PYG{l+m+mi}{0}\PYG{p}{;}
\PYG{p}{\PYGZcb{}}
\end{sphinxVerbatim}


\bigskip\hrule\bigskip

\index{lammps\_get\_natoms (C++ function)@\spxentry{lammps\_get\_natoms}\spxextra{C++ function}}

\begin{fulllineitems}
\phantomsection\label{\detokenize{Library_properties:_CPPv417lammps_get_natomsPv}}\phantomsection\label{\detokenize{Library_properties:_CPPv317lammps_get_natomsPv}}\phantomsection\label{\detokenize{Library_properties:_CPPv217lammps_get_natomsPv}}\phantomsection\label{\detokenize{Library_properties:lammps_get_natoms__voidP}}
\pysigstartsignatures
\pysigstartmultiline
\pysiglinewithargsret{\phantomsection\label{\detokenize{Library_properties:library_8h_1aa7d4ad0dc261c21fa41cffa38ea3ac21}}\DUrole{kt}{double}\DUrole{w}{ }\sphinxbfcode{\sphinxupquote{\DUrole{n}{lammps\_get\_natoms}}}}{\DUrole{kt}{void}\DUrole{w}{ }\DUrole{p}{*}\DUrole{n,sig-param}{handle}}{}
\pysigstopmultiline
\pysigstopsignatures
\sphinxAtStartPar
Return the total number of atoms in the system.

\sphinxAtStartPar

\sphinxAtStartPar
This number may be very large when running large simulations across
multiple processes.  Depending on compile time choices, LAMMPS may be
using either 32\sphinxhyphen{}bit or a 64\sphinxhyphen{}bit integer to store this number. For
portability this function returns thus a double precision
floating point number, which can represent up to a 53\sphinxhyphen{}bit signed
integer exactly (\(\approx 10^{16}\)).

\sphinxAtStartPar
As an alternative, you can use {\hyperref[\detokenize{Library_properties:_CPPv421lammps_extract_globalPvPKc}]{\sphinxcrossref{\sphinxcode{\sphinxupquote{lammps\_extract\_global()}}}}}
and cast the resulting pointer to an integer pointer of the correct
size and dereference it.  The size of that integer (in bytes) can be
queried by calling {\hyperref[\detokenize{Library_properties:_CPPv422lammps_extract_settingPvPKc}]{\sphinxcrossref{\sphinxcode{\sphinxupquote{lammps\_extract\_setting()}}}}} to return
the size of a \sphinxcode{\sphinxupquote{bigint}} integer.

\sphinxAtStartPar
\DUrole{versionmodified,changed}{Changed in version 18Sep2020: }The type of the return value was changed from \sphinxcode{\sphinxupquote{int}} to \sphinxcode{\sphinxupquote{double}}
to accommodate reporting atom counts for larger systems that would
overflow a 32\sphinxhyphen{}bit int without having to depend on a 64\sphinxhyphen{}bit bit
integer type definition.

\begin{quote}\begin{description}
\sphinxlineitem{Parameters}
\sphinxAtStartPar
\sphinxstyleliteralstrong{\sphinxupquote{handle}} \textendash{} pointer to a previously created LAMMPS instance 

\sphinxlineitem{Returns}
\sphinxAtStartPar
total number of atoms or 0 if value is too large 

\end{description}\end{quote}

\end{fulllineitems}



\bigskip\hrule\bigskip

\index{lammps\_get\_thermo (C++ function)@\spxentry{lammps\_get\_thermo}\spxextra{C++ function}}

\begin{fulllineitems}
\phantomsection\label{\detokenize{Library_properties:_CPPv417lammps_get_thermoPvPKc}}\phantomsection\label{\detokenize{Library_properties:_CPPv317lammps_get_thermoPvPKc}}\phantomsection\label{\detokenize{Library_properties:_CPPv217lammps_get_thermoPvPKc}}\phantomsection\label{\detokenize{Library_properties:lammps_get_thermo__voidP.cCP}}
\pysigstartsignatures
\pysigstartmultiline
\pysiglinewithargsret{\phantomsection\label{\detokenize{Library_properties:library_8h_1a5f54fbac6f7be9c20605c0faf71ec77a}}\DUrole{kt}{double}\DUrole{w}{ }\sphinxbfcode{\sphinxupquote{\DUrole{n}{lammps\_get\_thermo}}}}{\DUrole{kt}{void}\DUrole{w}{ }\DUrole{p}{*}\DUrole{n,sig-param}{handle}\sphinxparamcomma \DUrole{k}{const}\DUrole{w}{ }\DUrole{kt}{char}\DUrole{w}{ }\DUrole{p}{*}\DUrole{n,sig-param}{keyword}}{}
\pysigstopmultiline
\pysigstopsignatures
\sphinxAtStartPar
Evaluate a thermo keyword.

\sphinxAtStartPar
This function returns the current value of a {\hyperref[\detokenize{thermo_style::doc}]{\sphinxcrossref{\DUrole{doc}{thermo keyword}}}}.
Unlike {\hyperref[\detokenize{Library_properties:_CPPv421lammps_extract_globalPvPKc}]{\sphinxcrossref{\sphinxcode{\sphinxupquote{lammps\_extract\_global()}}}}} it does not give access to the
storage of the desired data but returns its value as a \sphinxcode{\sphinxupquote{double}}, so it
can also return information that is computed on\sphinxhyphen{}the\sphinxhyphen{}fly.
Use {\hyperref[\detokenize{Library_properties:_CPPv418lammps_last_thermoPvPKci}]{\sphinxcrossref{\sphinxcode{\sphinxupquote{lammps\_last\_thermo()}}}}} to get access to the cached data from
the last thermo output.
\begin{quote}\begin{description}
\sphinxlineitem{Parameters}\begin{itemize}
\item {} 
\sphinxAtStartPar
\sphinxstyleliteralstrong{\sphinxupquote{handle}} \textendash{} pointer to a previously created LAMMPS instance 

\item {} 
\sphinxAtStartPar
\sphinxstyleliteralstrong{\sphinxupquote{keyword}} \textendash{} string with the name of the thermo keyword 

\end{itemize}

\sphinxlineitem{Returns}
\sphinxAtStartPar
value of the requested thermo property or 0.0 

\end{description}\end{quote}

\end{fulllineitems}



\bigskip\hrule\bigskip

\index{lammps\_last\_thermo (C++ function)@\spxentry{lammps\_last\_thermo}\spxextra{C++ function}}

\begin{fulllineitems}
\phantomsection\label{\detokenize{Library_properties:_CPPv418lammps_last_thermoPvPKci}}\phantomsection\label{\detokenize{Library_properties:_CPPv318lammps_last_thermoPvPKci}}\phantomsection\label{\detokenize{Library_properties:_CPPv218lammps_last_thermoPvPKci}}\phantomsection\label{\detokenize{Library_properties:lammps_last_thermo__voidP.cCP.i}}
\pysigstartsignatures
\pysigstartmultiline
\pysiglinewithargsret{\phantomsection\label{\detokenize{Library_properties:library_8h_1a657031de981d9a1ec7d973e28f9a3c13}}\DUrole{kt}{void}\DUrole{w}{ }\DUrole{p}{*}\sphinxbfcode{\sphinxupquote{\DUrole{n}{lammps\_last\_thermo}}}}{\DUrole{kt}{void}\DUrole{w}{ }\DUrole{p}{*}\DUrole{n,sig-param}{handle}\sphinxparamcomma \DUrole{k}{const}\DUrole{w}{ }\DUrole{kt}{char}\DUrole{w}{ }\DUrole{p}{*}\DUrole{n,sig-param}{what}\sphinxparamcomma \DUrole{kt}{int}\DUrole{w}{ }\DUrole{n,sig-param}{index}}{}
\pysigstopmultiline
\pysigstopsignatures
\sphinxAtStartPar
Access cached data from last thermo output

\sphinxAtStartPar

\sphinxAtStartPar
\DUrole{versionmodified,added}{Added in version 15Jun2023.}

\sphinxAtStartPar
This function provides access to cached data from the last thermo output.
This differs from {\hyperref[\detokenize{Library_properties:_CPPv417lammps_get_thermoPvPKc}]{\sphinxcrossref{\sphinxcode{\sphinxupquote{lammps\_get\_thermo()}}}}} in that it does \sphinxstylestrong{not}
trigger an evaluation.  Instead it provides direct access to a read\sphinxhyphen{}only
location of the last thermo output data and the corresponding keyword
strings.  How to handle the return value depends on the value of the \sphinxstyleemphasis{what}
argument string.


\begin{savenotes}\sphinxattablestart
\sphinxthistablewithglobalstyle
\centering
\begin{tabular}[t]{\X{14}{100}\X{51}{100}\X{25}{100}\X{10}{100}}
\sphinxtoprule
\sphinxstyletheadfamily 
\sphinxAtStartPar
Value of \sphinxstyleemphasis{what}
&\sphinxstyletheadfamily 
\sphinxAtStartPar
Description of return value
&\sphinxstyletheadfamily 
\sphinxAtStartPar
Data type
&\sphinxstyletheadfamily 
\sphinxAtStartPar
Uses index
\\
\sphinxmidrule
\sphinxtableatstartofbodyhook
\sphinxAtStartPar
setup
&
\sphinxAtStartPar
1 if setup is not completed and thus thermo data invalid, 0 otherwise
&
\sphinxAtStartPar
pointer to int
&
\sphinxAtStartPar
no
\\
\sphinxhline
\sphinxAtStartPar
line
&
\sphinxAtStartPar
line number (0\sphinxhyphen{}based) of current line in current file or buffer
&
\sphinxAtStartPar
pointer to int
&
\sphinxAtStartPar
no
\\
\sphinxhline
\sphinxAtStartPar
imagename
&
\sphinxAtStartPar
file name of the last {\hyperref[\detokenize{dump_image::doc}]{\sphinxcrossref{\DUrole{doc}{dump image}}}} file written
&
\sphinxAtStartPar
pointer to 0\sphinxhyphen{}terminated const char array
&
\sphinxAtStartPar
no
\\
\sphinxhline
\sphinxAtStartPar
step
&
\sphinxAtStartPar
timestep when the last thermo output was generated or \sphinxhyphen{}1
&
\sphinxAtStartPar
pointer to bigint
&
\sphinxAtStartPar
no
\\
\sphinxhline
\sphinxAtStartPar
num
&
\sphinxAtStartPar
number of fields in thermo output
&
\sphinxAtStartPar
pointer to int
&
\sphinxAtStartPar
no
\\
\sphinxhline
\sphinxAtStartPar
keyword
&
\sphinxAtStartPar
column keyword for thermo output
&
\sphinxAtStartPar
pointer to 0\sphinxhyphen{}terminated const char array
&
\sphinxAtStartPar
yes
\\
\sphinxhline
\sphinxAtStartPar
type
&
\sphinxAtStartPar
data type of thermo output column; see \sphinxcode{\sphinxupquote{\_LMP\_DATATYPE\_CONST}}
&
\sphinxAtStartPar
pointer to int
&
\sphinxAtStartPar
yes
\\
\sphinxhline
\sphinxAtStartPar
data
&
\sphinxAtStartPar
actual field data for column
&
\sphinxAtStartPar
pointer to int, int64\_t or double
&
\sphinxAtStartPar
yes
\\
\sphinxbottomrule
\end{tabular}
\sphinxtableafterendhook\par
\sphinxattableend\end{savenotes}


\begin{sphinxadmonition}{note}{Note:}
\sphinxAtStartPar
The \sphinxstyleemphasis{type} property points to a static location that is reassigned
with every call, so the returned pointer should be recast,
dereferenced, and assigned immediately. Otherwise, its value may be
changed with the next invocation of the function.
\end{sphinxadmonition}
\begin{quote}\begin{description}
\sphinxlineitem{Parameters}\begin{itemize}
\item {} 
\sphinxAtStartPar
\sphinxstyleliteralstrong{\sphinxupquote{handle}} \textendash{} pointer to a previously created LAMMPS instance 

\item {} 
\sphinxAtStartPar
\sphinxstyleliteralstrong{\sphinxupquote{what}} \textendash{} string with the kind of data requested 

\item {} 
\sphinxAtStartPar
\sphinxstyleliteralstrong{\sphinxupquote{index}} \textendash{} integer with index into data arrays, ignored for scalar data 

\end{itemize}

\sphinxlineitem{Returns}
\sphinxAtStartPar
pointer to location of requested data cast to void or NULL 

\end{description}\end{quote}

\end{fulllineitems}



\bigskip\hrule\bigskip

\index{lammps\_extract\_box (C++ function)@\spxentry{lammps\_extract\_box}\spxextra{C++ function}}

\begin{fulllineitems}
\phantomsection\label{\detokenize{Library_properties:_CPPv418lammps_extract_boxPvPdPdPdPdPdPiPi}}\phantomsection\label{\detokenize{Library_properties:_CPPv318lammps_extract_boxPvPdPdPdPdPdPiPi}}\phantomsection\label{\detokenize{Library_properties:_CPPv218lammps_extract_boxPvPdPdPdPdPdPiPi}}\phantomsection\label{\detokenize{Library_properties:lammps_extract_box__voidP.doubleP.doubleP.doubleP.doubleP.doubleP.iP.iP}}
\pysigstartsignatures
\pysigstartmultiline
\pysiglinewithargsret{\phantomsection\label{\detokenize{Library_properties:library_8h_1a5280ea3a305969261faa1a0258b9c7f9}}\DUrole{kt}{void}\DUrole{w}{ }\sphinxbfcode{\sphinxupquote{\DUrole{n}{lammps\_extract\_box}}}}{\DUrole{kt}{void}\DUrole{w}{ }\DUrole{p}{*}\DUrole{n,sig-param}{handle}\sphinxparamcomma \DUrole{kt}{double}\DUrole{w}{ }\DUrole{p}{*}\DUrole{n,sig-param}{boxlo}\sphinxparamcomma \DUrole{kt}{double}\DUrole{w}{ }\DUrole{p}{*}\DUrole{n,sig-param}{boxhi}\sphinxparamcomma \DUrole{kt}{double}\DUrole{w}{ }\DUrole{p}{*}\DUrole{n,sig-param}{xy}\sphinxparamcomma \DUrole{kt}{double}\DUrole{w}{ }\DUrole{p}{*}\DUrole{n,sig-param}{yz}\sphinxparamcomma \DUrole{kt}{double}\DUrole{w}{ }\DUrole{p}{*}\DUrole{n,sig-param}{xz}\sphinxparamcomma \DUrole{kt}{int}\DUrole{w}{ }\DUrole{p}{*}\DUrole{n,sig-param}{pflags}\sphinxparamcomma \DUrole{kt}{int}\DUrole{w}{ }\DUrole{p}{*}\DUrole{n,sig-param}{boxflag}}{}
\pysigstopmultiline
\pysigstopsignatures
\sphinxAtStartPar
Extract simulation box parameters.

\sphinxAtStartPar
This function (re\sphinxhyphen{})initializes the simulation box and boundary
information and then assign the designated data to the locations in the
pointers passed as arguments. Any argument (except the first) may be
a NULL pointer and then will not be assigned.
\begin{quote}\begin{description}
\sphinxlineitem{Parameters}\begin{itemize}
\item {} 
\sphinxAtStartPar
\sphinxstyleliteralstrong{\sphinxupquote{handle}} \textendash{} pointer to a previously created LAMMPS instance 

\item {} 
\sphinxAtStartPar
\sphinxstyleliteralstrong{\sphinxupquote{boxlo}} \textendash{} pointer to 3 doubles where the lower box boundary is stored 

\item {} 
\sphinxAtStartPar
\sphinxstyleliteralstrong{\sphinxupquote{boxhi}} \textendash{} pointer to 3 doubles where the upper box boundary is stored 

\item {} 
\sphinxAtStartPar
\sphinxstyleliteralstrong{\sphinxupquote{xy}} \textendash{} pointer to a double where the xy tilt factor is stored 

\item {} 
\sphinxAtStartPar
\sphinxstyleliteralstrong{\sphinxupquote{yz}} \textendash{} pointer to a double where the yz tilt factor is stored 

\item {} 
\sphinxAtStartPar
\sphinxstyleliteralstrong{\sphinxupquote{xz}} \textendash{} pointer to a double where the xz tilt factor is stored 

\item {} 
\sphinxAtStartPar
\sphinxstyleliteralstrong{\sphinxupquote{pflags}} \textendash{} pointer to 3 ints, set to 1 for periodic boundaries and 0 for non\sphinxhyphen{}periodic 

\item {} 
\sphinxAtStartPar
\sphinxstyleliteralstrong{\sphinxupquote{boxflag}} \textendash{} pointer to an int, which is set to 1 if the box will be changed during a simulation by a fix and 0 if not. 

\end{itemize}

\end{description}\end{quote}

\end{fulllineitems}



\bigskip\hrule\bigskip

\index{lammps\_reset\_box (C++ function)@\spxentry{lammps\_reset\_box}\spxextra{C++ function}}

\begin{fulllineitems}
\phantomsection\label{\detokenize{Library_properties:_CPPv416lammps_reset_boxPvPdPdddd}}\phantomsection\label{\detokenize{Library_properties:_CPPv316lammps_reset_boxPvPdPdddd}}\phantomsection\label{\detokenize{Library_properties:_CPPv216lammps_reset_boxPvPdPdddd}}\phantomsection\label{\detokenize{Library_properties:lammps_reset_box__voidP.doubleP.doubleP.double.double.double}}
\pysigstartsignatures
\pysigstartmultiline
\pysiglinewithargsret{\phantomsection\label{\detokenize{Library_properties:library_8h_1ada23ca4aacdcc9b5991af9968e0369e8}}\DUrole{kt}{void}\DUrole{w}{ }\sphinxbfcode{\sphinxupquote{\DUrole{n}{lammps\_reset\_box}}}}{\DUrole{kt}{void}\DUrole{w}{ }\DUrole{p}{*}\DUrole{n,sig-param}{handle}\sphinxparamcomma \DUrole{kt}{double}\DUrole{w}{ }\DUrole{p}{*}\DUrole{n,sig-param}{boxlo}\sphinxparamcomma \DUrole{kt}{double}\DUrole{w}{ }\DUrole{p}{*}\DUrole{n,sig-param}{boxhi}\sphinxparamcomma \DUrole{kt}{double}\DUrole{w}{ }\DUrole{n,sig-param}{xy}\sphinxparamcomma \DUrole{kt}{double}\DUrole{w}{ }\DUrole{n,sig-param}{yz}\sphinxparamcomma \DUrole{kt}{double}\DUrole{w}{ }\DUrole{n,sig-param}{xz}}{}
\pysigstopmultiline
\pysigstopsignatures
\sphinxAtStartPar
Reset simulation box parameters.

\sphinxAtStartPar
This function sets the simulation box dimensions (upper and lower bounds
and tilt factors) from the provided data and then re\sphinxhyphen{}initializes the box
information and all derived settings. It may only be called before atoms
are created.
\begin{quote}\begin{description}
\sphinxlineitem{Parameters}\begin{itemize}
\item {} 
\sphinxAtStartPar
\sphinxstyleliteralstrong{\sphinxupquote{handle}} \textendash{} pointer to a previously created LAMMPS instance 

\item {} 
\sphinxAtStartPar
\sphinxstyleliteralstrong{\sphinxupquote{boxlo}} \textendash{} pointer to 3 doubles containing the lower box boundary 

\item {} 
\sphinxAtStartPar
\sphinxstyleliteralstrong{\sphinxupquote{boxhi}} \textendash{} pointer to 3 doubles containing the upper box boundary 

\item {} 
\sphinxAtStartPar
\sphinxstyleliteralstrong{\sphinxupquote{xy}} \textendash{} xy tilt factor 

\item {} 
\sphinxAtStartPar
\sphinxstyleliteralstrong{\sphinxupquote{yz}} \textendash{} yz tilt factor 

\item {} 
\sphinxAtStartPar
\sphinxstyleliteralstrong{\sphinxupquote{xz}} \textendash{} xz tilt factor 

\end{itemize}

\end{description}\end{quote}

\end{fulllineitems}



\bigskip\hrule\bigskip

\index{lammps\_memory\_usage (C++ function)@\spxentry{lammps\_memory\_usage}\spxextra{C++ function}}

\begin{fulllineitems}
\phantomsection\label{\detokenize{Library_properties:_CPPv419lammps_memory_usagePvPd}}\phantomsection\label{\detokenize{Library_properties:_CPPv319lammps_memory_usagePvPd}}\phantomsection\label{\detokenize{Library_properties:_CPPv219lammps_memory_usagePvPd}}\phantomsection\label{\detokenize{Library_properties:lammps_memory_usage__voidP.doubleP}}
\pysigstartsignatures
\pysigstartmultiline
\pysiglinewithargsret{\phantomsection\label{\detokenize{Library_properties:library_8h_1ab1a5d06bc29dd1233cde03c857ff0821}}\DUrole{kt}{void}\DUrole{w}{ }\sphinxbfcode{\sphinxupquote{\DUrole{n}{lammps\_memory\_usage}}}}{\DUrole{kt}{void}\DUrole{w}{ }\DUrole{p}{*}\DUrole{n,sig-param}{handle}\sphinxparamcomma \DUrole{kt}{double}\DUrole{w}{ }\DUrole{p}{*}\DUrole{n,sig-param}{meminfo}}{}
\pysigstopmultiline
\pysigstopsignatures
\sphinxAtStartPar
Get memory usage information

\sphinxAtStartPar

\sphinxAtStartPar
\DUrole{versionmodified,added}{Added in version 18Sep2020.}

\sphinxAtStartPar
This function will retrieve memory usage information for the current
LAMMPS instance or process.  The \sphinxstyleemphasis{meminfo} buffer will be filled with
3 different numbers (if supported by the operating system).  The first
is the tally (in MBytes) of all large memory allocations made by LAMMPS.
This is a lower boundary of how much memory is requested and does not
account for memory allocated on the stack or allocations via \sphinxcode{\sphinxupquote{new}}.
The second number is the current memory allocation of the current process
as returned by a memory allocation reporting in the system library.  The
third number is the maximum amount of RAM (not swap) used by the process
so far. If any of the two latter parameters is not supported by the operating
system it will be set to zero.

\begin{quote}\begin{description}
\sphinxlineitem{Parameters}\begin{itemize}
\item {} 
\sphinxAtStartPar
\sphinxstyleliteralstrong{\sphinxupquote{handle}} \textendash{} pointer to a previously created LAMMPS instance 

\item {} 
\sphinxAtStartPar
\sphinxstyleliteralstrong{\sphinxupquote{meminfo}} \textendash{} buffer with space for at least 3 double to store data in. 

\end{itemize}

\end{description}\end{quote}

\end{fulllineitems}



\bigskip\hrule\bigskip

\index{lammps\_get\_mpi\_comm (C++ function)@\spxentry{lammps\_get\_mpi\_comm}\spxextra{C++ function}}

\begin{fulllineitems}
\phantomsection\label{\detokenize{Library_properties:_CPPv419lammps_get_mpi_commPv}}\phantomsection\label{\detokenize{Library_properties:_CPPv319lammps_get_mpi_commPv}}\phantomsection\label{\detokenize{Library_properties:_CPPv219lammps_get_mpi_commPv}}\phantomsection\label{\detokenize{Library_properties:lammps_get_mpi_comm__voidP}}
\pysigstartsignatures
\pysigstartmultiline
\pysiglinewithargsret{\phantomsection\label{\detokenize{Library_properties:library_8h_1a27501b4b00aa53f6621e996e3eeeb149}}\DUrole{kt}{int}\DUrole{w}{ }\sphinxbfcode{\sphinxupquote{\DUrole{n}{lammps\_get\_mpi\_comm}}}}{\DUrole{kt}{void}\DUrole{w}{ }\DUrole{p}{*}\DUrole{n,sig-param}{handle}}{}
\pysigstopmultiline
\pysigstopsignatures
\sphinxAtStartPar
Return current LAMMPS world communicator as integer

\sphinxAtStartPar

\sphinxAtStartPar
\DUrole{versionmodified,added}{Added in version 18Sep2020.}

\sphinxAtStartPar
This will take the LAMMPS “world” communicator and convert it to an
integer using \sphinxcode{\sphinxupquote{MPI\_Comm\_c2f()}}, so it is equivalent to the
corresponding MPI communicator in Fortran. This way it can be safely
passed around between different programming languages.  To convert it
to the C language representation use \sphinxcode{\sphinxupquote{MPI\_Comm\_f2c()}}.

\sphinxAtStartPar
If LAMMPS was compiled with MPI\_STUBS, this function returns \sphinxhyphen{}1.
\begin{description}
\sphinxlineitem{\sphinxstyleemphasis{See also}}
\sphinxAtStartPar
{\hyperref[\detokenize{Library_create:_CPPv419lammps_open_fortraniPPci}]{\sphinxcrossref{\sphinxcode{\sphinxupquote{lammps\_open\_fortran()}}}}}

\end{description}

\begin{quote}\begin{description}
\sphinxlineitem{Parameters}
\sphinxAtStartPar
\sphinxstyleliteralstrong{\sphinxupquote{handle}} \textendash{} pointer to a previously created LAMMPS instance 

\sphinxlineitem{Returns}
\sphinxAtStartPar
Fortran representation of the LAMMPS world communicator 

\end{description}\end{quote}

\end{fulllineitems}



\bigskip\hrule\bigskip

\index{lammps\_extract\_setting (C++ function)@\spxentry{lammps\_extract\_setting}\spxextra{C++ function}}

\begin{fulllineitems}
\phantomsection\label{\detokenize{Library_properties:_CPPv422lammps_extract_settingPvPKc}}\phantomsection\label{\detokenize{Library_properties:_CPPv322lammps_extract_settingPvPKc}}\phantomsection\label{\detokenize{Library_properties:_CPPv222lammps_extract_settingPvPKc}}\phantomsection\label{\detokenize{Library_properties:lammps_extract_setting__voidP.cCP}}
\pysigstartsignatures
\pysigstartmultiline
\pysiglinewithargsret{\phantomsection\label{\detokenize{Library_properties:library_8h_1aec668a9f577e41295afb673f7edc99da}}\DUrole{kt}{int}\DUrole{w}{ }\sphinxbfcode{\sphinxupquote{\DUrole{n}{lammps\_extract\_setting}}}}{\DUrole{kt}{void}\DUrole{w}{ }\DUrole{p}{*}\DUrole{n,sig-param}{handle}\sphinxparamcomma \DUrole{k}{const}\DUrole{w}{ }\DUrole{kt}{char}\DUrole{w}{ }\DUrole{p}{*}\DUrole{n,sig-param}{keyword}}{}
\pysigstopmultiline
\pysigstopsignatures
\sphinxAtStartPar
Query LAMMPS about global settings.

\sphinxAtStartPar

\sphinxAtStartPar
This function will retrieve or compute global properties. In contrast to
{\hyperref[\detokenize{Library_properties:_CPPv417lammps_get_thermoPvPKc}]{\sphinxcrossref{\sphinxcode{\sphinxupquote{lammps\_get\_thermo()}}}}} this function returns an \sphinxcode{\sphinxupquote{int}}.  The
following tables list the currently supported keyword.  If a keyword is
not recognized, the function returns \sphinxhyphen{}1.  The integer sizes functions may
be called without a valid LAMMPS object handle (it is ignored).
\begin{itemize}
\item {} 
\sphinxAtStartPar
{\hyperref[\detokenize{Library_properties:extract-integer-sizes}]{\sphinxcrossref{\DUrole{std,std-ref}{Integer sizes}}}}

\item {} 
\sphinxAtStartPar
{\hyperref[\detokenize{Library_properties:extract-image-masks}]{\sphinxcrossref{\DUrole{std,std-ref}{Image masks}}}}

\item {} 
\sphinxAtStartPar
{\hyperref[\detokenize{Library_properties:extract-system-status}]{\sphinxcrossref{\DUrole{std,std-ref}{System status}}}}

\item {} 
\sphinxAtStartPar
{\hyperref[\detokenize{Library_properties:extract-system-sizes}]{\sphinxcrossref{\DUrole{std,std-ref}{System sizes}}}}

\item {} 
\sphinxAtStartPar
{\hyperref[\detokenize{Library_properties:extract-atom-flags}]{\sphinxcrossref{\DUrole{std,std-ref}{Atom style flags}}}}

\end{itemize}
\phantomsection\label{\detokenize{Library_properties:extract-integer-sizes}}
\sphinxAtStartPar
\sphinxstylestrong{Integer sizes}


\begin{savenotes}\sphinxattablestart
\sphinxthistablewithglobalstyle
\centering
\begin{tabular}[t]{\X{17}{100}\X{83}{100}}
\sphinxtoprule
\sphinxstyletheadfamily 
\sphinxAtStartPar
Keyword
&\sphinxstyletheadfamily 
\sphinxAtStartPar
Description / Return value
\\
\sphinxmidrule
\sphinxtableatstartofbodyhook
\sphinxAtStartPar
bigint
&
\sphinxAtStartPar
size of the \sphinxcode{\sphinxupquote{bigint}} integer type, 4 or 8 bytes.
Set at {\hyperref[\detokenize{Build_settings:size}]{\sphinxcrossref{\DUrole{std,std-ref}{compile time}}}}.
\\
\sphinxhline
\sphinxAtStartPar
tagint
&
\sphinxAtStartPar
size of the \sphinxcode{\sphinxupquote{tagint}} integer type, 4 or 8 bytes.
Set at {\hyperref[\detokenize{Build_settings:size}]{\sphinxcrossref{\DUrole{std,std-ref}{compile time}}}}.
\\
\sphinxhline
\sphinxAtStartPar
imageint
&
\sphinxAtStartPar
size of the \sphinxcode{\sphinxupquote{imageint}} integer type, 4 or 8 bytes.
Set at {\hyperref[\detokenize{Build_settings:size}]{\sphinxcrossref{\DUrole{std,std-ref}{compile time}}}}.
\\
\sphinxbottomrule
\end{tabular}
\sphinxtableafterendhook\par
\sphinxattableend\end{savenotes}
\phantomsection\label{\detokenize{Library_properties:extract-image-masks}}
\sphinxAtStartPar
\sphinxstylestrong{Image masks}

\sphinxAtStartPar
These settings are related to how LAMMPS stores and interprets periodic images. The values are used
internally by the {\hyperref[\detokenize{Fortran::doc}]{\sphinxcrossref{\DUrole{doc}{Fortran interface}}}} and are not likely to be useful to users.


\begin{savenotes}\sphinxattablestart
\sphinxthistablewithglobalstyle
\centering
\begin{tabular}[t]{\X{17}{100}\X{83}{100}}
\sphinxtoprule
\sphinxstyletheadfamily 
\sphinxAtStartPar
Keyword
&\sphinxstyletheadfamily 
\sphinxAtStartPar
Description / Return value
\\
\sphinxmidrule
\sphinxtableatstartofbodyhook
\sphinxAtStartPar
IMGMASK
&
\sphinxAtStartPar
Bit\sphinxhyphen{}mask used to convert image flags to a single integer
\\
\sphinxhline
\sphinxAtStartPar
IMGMAX
&
\sphinxAtStartPar
Maximum allowed image number for a particular atom
\\
\sphinxhline
\sphinxAtStartPar
IMGBITS
&
\sphinxAtStartPar
Bits used in image counts
\\
\sphinxhline
\sphinxAtStartPar
IMG2BITS
&
\sphinxAtStartPar
Second bitmask used in image counts
\\
\sphinxbottomrule
\end{tabular}
\sphinxtableafterendhook\par
\sphinxattableend\end{savenotes}
\phantomsection\label{\detokenize{Library_properties:extract-system-status}}
\sphinxAtStartPar
\sphinxstylestrong{System status}


\begin{savenotes}\sphinxattablestart
\sphinxthistablewithglobalstyle
\centering
\begin{tabular}[t]{\X{17}{100}\X{83}{100}}
\sphinxtoprule
\sphinxstyletheadfamily 
\sphinxAtStartPar
Keyword
&\sphinxstyletheadfamily 
\sphinxAtStartPar
Description / Return value
\\
\sphinxmidrule
\sphinxtableatstartofbodyhook
\sphinxAtStartPar
dimension
&
\sphinxAtStartPar
Number of dimensions: 2 or 3. See {\hyperref[\detokenize{dimension::doc}]{\sphinxcrossref{\DUrole{doc}{dimension command}}}}.
\\
\sphinxhline
\sphinxAtStartPar
box\_exist
&
\sphinxAtStartPar
1 if the simulation box is defined, 0 if not.
See {\hyperref[\detokenize{create_box::doc}]{\sphinxcrossref{\DUrole{doc}{create\_box command}}}}.
\\
\sphinxhline
\sphinxAtStartPar
kokkos\_active
&
\sphinxAtStartPar
1 if the KOKKOS package is compiled in \sphinxstylestrong{and} activated, 0 if not.
See {\hyperref[\detokenize{Speed_kokkos::doc}]{\sphinxcrossref{\DUrole{doc}{KOKKOS package}}}}.
\\
\sphinxhline
\sphinxAtStartPar
kokkos\_nthreads
&
\sphinxAtStartPar
Number of Kokkos threads per MPI process, 0 if Kokkos is not active.
See {\hyperref[\detokenize{Speed_kokkos::doc}]{\sphinxcrossref{\DUrole{doc}{KOKKOS package}}}}.
\\
\sphinxhline
\sphinxAtStartPar
kokkos\_ngpus
&
\sphinxAtStartPar
Number of Kokkos gpus per physical node, 0 if Kokkos is not active or no GPU support.
See {\hyperref[\detokenize{Speed_kokkos::doc}]{\sphinxcrossref{\DUrole{doc}{KOKKOS package}}}}.
\\
\sphinxhline
\sphinxAtStartPar
nthreads
&
\sphinxAtStartPar
Number of requested OpenMP threads per MPI process for LAMMPS’ execution
\\
\sphinxhline
\sphinxAtStartPar
newton\_bond
&
\sphinxAtStartPar
1 if Newton’s 3rd law is applied to bonded interactions, 0 if not.
\\
\sphinxhline
\sphinxAtStartPar
newton\_pair
&
\sphinxAtStartPar
1 if Newton’s 3rd law is applied to non\sphinxhyphen{}bonded interactions, 0 if not.
\\
\sphinxhline
\sphinxAtStartPar
triclinic
&
\sphinxAtStartPar
1 if the the simulation box is triclinic, 0 if orthogonal.
See {\hyperref[\detokenize{change_box::doc}]{\sphinxcrossref{\DUrole{doc}{change\_box command}}}}.
\\
\sphinxbottomrule
\end{tabular}
\sphinxtableafterendhook\par
\sphinxattableend\end{savenotes}

\sphinxAtStartPar
\sphinxstylestrong{Communication status}


\begin{savenotes}\sphinxattablestart
\sphinxthistablewithglobalstyle
\centering
\begin{tabular}[t]{\X{15}{100}\X{85}{100}}
\sphinxtoprule
\sphinxstyletheadfamily 
\sphinxAtStartPar
Keyword
&\sphinxstyletheadfamily 
\sphinxAtStartPar
Description / Return value
\\
\sphinxmidrule
\sphinxtableatstartofbodyhook
\sphinxAtStartPar
universe\_rank
&
\sphinxAtStartPar
MPI rank on LAMMPS’ universe communicator (0 \textless{}= universe\_rank \textless{} universe\_size)
\\
\sphinxhline
\sphinxAtStartPar
universe\_size
&
\sphinxAtStartPar
Number of ranks on LAMMPS’ universe communicator (world\_size \textless{}= universe\_size)
\\
\sphinxhline
\sphinxAtStartPar
world\_rank
&
\sphinxAtStartPar
MPI rank on LAMMPS’ world communicator (0 \textless{}= world\_rank \textless{} world\_size, = comm\sphinxhyphen{}\textgreater{}me)
\\
\sphinxhline
\sphinxAtStartPar
world\_size
&
\sphinxAtStartPar
Number of ranks on LAMMPS’ world communicator (aka comm\sphinxhyphen{}\textgreater{}nprocs)
\\
\sphinxhline
\sphinxAtStartPar
comm\_style
&
\sphinxAtStartPar
communication style (0 = BRICK, 1 = TILED)
\\
\sphinxhline
\sphinxAtStartPar
comm\_layout
&
\sphinxAtStartPar
communication layout (0 = LAYOUT\_UNIFORM, 1 = LAYOUT\_NONUNIFORM, 2 = LAYOUT\_TILED)
\\
\sphinxhline
\sphinxAtStartPar
comm\_mode
&
\sphinxAtStartPar
communication mode (0 = SINGLE, 1 = MULTI, 2 = MULTIOLD)
\\
\sphinxhline
\sphinxAtStartPar
ghost\_velocity
&
\sphinxAtStartPar
whether velocities are communicated for ghost atoms (0 = no, 1 = yes)
\\
\sphinxbottomrule
\end{tabular}
\sphinxtableafterendhook\par
\sphinxattableend\end{savenotes}
\phantomsection\label{\detokenize{Library_properties:extract-system-sizes}}
\sphinxAtStartPar
\sphinxstylestrong{System sizes}


\begin{savenotes}\sphinxattablestart
\sphinxthistablewithglobalstyle
\centering
\begin{tabular}[t]{\X{17}{100}\X{83}{100}}
\sphinxtoprule
\sphinxstyletheadfamily 
\sphinxAtStartPar
Keyword
&\sphinxstyletheadfamily 
\sphinxAtStartPar
Description / Return value
\\
\sphinxmidrule
\sphinxtableatstartofbodyhook
\sphinxAtStartPar
nlocal
&
\sphinxAtStartPar
number of “owned” atoms of the current MPI rank.
\\
\sphinxhline
\sphinxAtStartPar
nghost
&
\sphinxAtStartPar
number of “ghost” atoms of the current MPI rank.
\\
\sphinxhline
\sphinxAtStartPar
nall
&
\sphinxAtStartPar
number of all “owned” and “ghost” atoms of the current MPI rank.
\\
\sphinxhline
\sphinxAtStartPar
nmax
&
\sphinxAtStartPar
maximum of nlocal+nghost across all MPI ranks (for per\sphinxhyphen{}atom data array size).
\\
\sphinxhline
\sphinxAtStartPar
ntypes
&
\sphinxAtStartPar
number of atom types
\\
\sphinxhline
\sphinxAtStartPar
nbondtypes
&
\sphinxAtStartPar
number of bond types
\\
\sphinxhline
\sphinxAtStartPar
nangletypes
&
\sphinxAtStartPar
number of angle types
\\
\sphinxhline
\sphinxAtStartPar
ndihedraltypes
&
\sphinxAtStartPar
number of dihedral types
\\
\sphinxhline
\sphinxAtStartPar
nimpropertypes
&
\sphinxAtStartPar
number of improper types
\\
\sphinxhline
\sphinxAtStartPar
nellipsoids
&
\sphinxAtStartPar
number of atoms that have ellipsoid data
\\
\sphinxhline
\sphinxAtStartPar
nlines
&
\sphinxAtStartPar
number of atoms that have line data (see {\hyperref[\detokenize{pair_line_lj::doc}]{\sphinxcrossref{\DUrole{doc}{pair style line/lj}}}})
\\
\sphinxhline
\sphinxAtStartPar
ntris
&
\sphinxAtStartPar
number of atoms that have triangle data (see {\hyperref[\detokenize{pair_tri_lj::doc}]{\sphinxcrossref{\DUrole{doc}{pair style tri/lj}}}})
\\
\sphinxhline
\sphinxAtStartPar
nbodies
&
\sphinxAtStartPar
number of atoms that have body data (see {\hyperref[\detokenize{Howto_body::doc}]{\sphinxcrossref{\DUrole{doc}{the Body particle HowTo}}}})
\\
\sphinxbottomrule
\end{tabular}
\sphinxtableafterendhook\par
\sphinxattableend\end{savenotes}
\phantomsection\label{\detokenize{Library_properties:extract-atom-flags}}
\sphinxAtStartPar
\sphinxstylestrong{Atom style flags}


\begin{savenotes}\sphinxattablestart
\sphinxthistablewithglobalstyle
\centering
\begin{tabular}[t]{\X{15}{100}\X{85}{100}}
\sphinxtoprule
\sphinxstyletheadfamily 
\sphinxAtStartPar
Keyword
&\sphinxstyletheadfamily 
\sphinxAtStartPar
Description / Return value
\\
\sphinxmidrule
\sphinxtableatstartofbodyhook
\sphinxAtStartPar
molecule\_flag
&
\sphinxAtStartPar
1 if the atom style includes molecular topology data. See {\hyperref[\detokenize{atom_style::doc}]{\sphinxcrossref{\DUrole{doc}{atom\_style command}}}}.
\\
\sphinxhline
\sphinxAtStartPar
q\_flag
&
\sphinxAtStartPar
1 if the atom style includes point charges. See {\hyperref[\detokenize{atom_style::doc}]{\sphinxcrossref{\DUrole{doc}{atom\_style command}}}}.
\\
\sphinxhline
\sphinxAtStartPar
mu\_flag
&
\sphinxAtStartPar
1 if the atom style includes point dipoles. See {\hyperref[\detokenize{atom_style::doc}]{\sphinxcrossref{\DUrole{doc}{atom\_style command}}}}.
\\
\sphinxhline
\sphinxAtStartPar
rmass\_flag
&
\sphinxAtStartPar
1 if the atom style includes per\sphinxhyphen{}atom masses, 0 if there are per\sphinxhyphen{}type masses. See {\hyperref[\detokenize{atom_style::doc}]{\sphinxcrossref{\DUrole{doc}{atom\_style command}}}}.
\\
\sphinxhline
\sphinxAtStartPar
radius\_flag
&
\sphinxAtStartPar
1 if the atom style includes a per\sphinxhyphen{}atom radius. See {\hyperref[\detokenize{atom_style::doc}]{\sphinxcrossref{\DUrole{doc}{atom\_style command}}}}.
\\
\sphinxhline
\sphinxAtStartPar
ellipsoid\_flag
&
\sphinxAtStartPar
1 if the atom style describes extended particles that may be ellipsoidal. See {\hyperref[\detokenize{atom_style::doc}]{\sphinxcrossref{\DUrole{doc}{atom\_style command}}}}.
\\
\sphinxhline
\sphinxAtStartPar
omega\_flag
&
\sphinxAtStartPar
1 if the atom style can store per\sphinxhyphen{}atom rotational velocities. See {\hyperref[\detokenize{atom_style::doc}]{\sphinxcrossref{\DUrole{doc}{atom\_style command}}}}.
\\
\sphinxhline
\sphinxAtStartPar
torque\_flag
&
\sphinxAtStartPar
1 if the atom style can store per\sphinxhyphen{}atom torques. See {\hyperref[\detokenize{atom_style::doc}]{\sphinxcrossref{\DUrole{doc}{atom\_style command}}}}.
\\
\sphinxhline
\sphinxAtStartPar
angmom\_flag
&
\sphinxAtStartPar
1 if the atom style can store per\sphinxhyphen{}atom angular momentum. See {\hyperref[\detokenize{atom_style::doc}]{\sphinxcrossref{\DUrole{doc}{atom\_style command}}}}.
\\
\sphinxbottomrule
\end{tabular}
\sphinxtableafterendhook\par
\sphinxattableend\end{savenotes}
\begin{description}
\sphinxlineitem{\sphinxstyleemphasis{See also}}
\sphinxAtStartPar
{\hyperref[\detokenize{Library_properties:_CPPv421lammps_extract_globalPvPKc}]{\sphinxcrossref{\sphinxcode{\sphinxupquote{lammps\_extract\_global()}}}}}

\end{description}

\begin{quote}\begin{description}
\sphinxlineitem{Parameters}\begin{itemize}
\item {} 
\sphinxAtStartPar
\sphinxstyleliteralstrong{\sphinxupquote{handle}} \textendash{} pointer to a previously created LAMMPS instance 

\item {} 
\sphinxAtStartPar
\sphinxstyleliteralstrong{\sphinxupquote{keyword}} \textendash{} string with the name of the thermo keyword 

\end{itemize}

\sphinxlineitem{Returns}
\sphinxAtStartPar
value of the queried setting or \sphinxhyphen{}1 if unknown 

\end{description}\end{quote}

\end{fulllineitems}



\bigskip\hrule\bigskip

\index{lammps\_extract\_global\_datatype (C++ function)@\spxentry{lammps\_extract\_global\_datatype}\spxextra{C++ function}}

\begin{fulllineitems}
\phantomsection\label{\detokenize{Library_properties:_CPPv430lammps_extract_global_datatypePvPKc}}\phantomsection\label{\detokenize{Library_properties:_CPPv330lammps_extract_global_datatypePvPKc}}\phantomsection\label{\detokenize{Library_properties:_CPPv230lammps_extract_global_datatypePvPKc}}\phantomsection\label{\detokenize{Library_properties:lammps_extract_global_datatype__voidP.cCP}}
\pysigstartsignatures
\pysigstartmultiline
\pysiglinewithargsret{\phantomsection\label{\detokenize{Library_properties:library_8h_1a7a5da49a662d7cc3e8eaa79d7cc344e0}}\DUrole{kt}{int}\DUrole{w}{ }\sphinxbfcode{\sphinxupquote{\DUrole{n}{lammps\_extract\_global\_datatype}}}}{\DUrole{kt}{void}\DUrole{w}{ }\DUrole{p}{*}\DUrole{n,sig-param}{handle}\sphinxparamcomma \DUrole{k}{const}\DUrole{w}{ }\DUrole{kt}{char}\DUrole{w}{ }\DUrole{p}{*}\DUrole{n,sig-param}{name}}{}
\pysigstopmultiline
\pysigstopsignatures
\sphinxAtStartPar
Get data type of internal global LAMMPS variables or arrays.

\sphinxAtStartPar

\sphinxAtStartPar
\DUrole{versionmodified,added}{Added in version 18Sep2020.}

\sphinxAtStartPar
This function returns an integer that encodes the data type of the global
property with the specified name. See \sphinxcode{\sphinxupquote{\_LMP\_DATATYPE\_CONST}} for valid
values. Callers of {\hyperref[\detokenize{Library_properties:_CPPv421lammps_extract_globalPvPKc}]{\sphinxcrossref{\sphinxcode{\sphinxupquote{lammps\_extract\_global()}}}}} can use this information
to then decide how to cast the \sphinxcode{\sphinxupquote{void *}} pointer and access the data.

\begin{quote}\begin{description}
\sphinxlineitem{Parameters}\begin{itemize}
\item {} 
\sphinxAtStartPar
\sphinxstyleliteralstrong{\sphinxupquote{handle}} \textendash{} pointer to a previously created LAMMPS instance (unused) 

\item {} 
\sphinxAtStartPar
\sphinxstyleliteralstrong{\sphinxupquote{name}} \textendash{} string with the name of the extracted property 

\end{itemize}

\sphinxlineitem{Returns}
\sphinxAtStartPar
integer constant encoding the data type of the property or \sphinxhyphen{}1 if not found. 

\end{description}\end{quote}

\end{fulllineitems}



\bigskip\hrule\bigskip

\index{lammps\_extract\_global (C++ function)@\spxentry{lammps\_extract\_global}\spxextra{C++ function}}

\begin{fulllineitems}
\phantomsection\label{\detokenize{Library_properties:_CPPv421lammps_extract_globalPvPKc}}\phantomsection\label{\detokenize{Library_properties:_CPPv321lammps_extract_globalPvPKc}}\phantomsection\label{\detokenize{Library_properties:_CPPv221lammps_extract_globalPvPKc}}\phantomsection\label{\detokenize{Library_properties:lammps_extract_global__voidP.cCP}}
\pysigstartsignatures
\pysigstartmultiline
\pysiglinewithargsret{\phantomsection\label{\detokenize{Library_properties:library_8h_1a67d3bbc188b284fc94a5df6f32d5b028}}\DUrole{kt}{void}\DUrole{w}{ }\DUrole{p}{*}\sphinxbfcode{\sphinxupquote{\DUrole{n}{lammps\_extract\_global}}}}{\DUrole{kt}{void}\DUrole{w}{ }\DUrole{p}{*}\DUrole{n,sig-param}{handle}\sphinxparamcomma \DUrole{k}{const}\DUrole{w}{ }\DUrole{kt}{char}\DUrole{w}{ }\DUrole{p}{*}\DUrole{n,sig-param}{name}}{}
\pysigstopmultiline
\pysigstopsignatures
\sphinxAtStartPar
Get pointer to internal global LAMMPS variables or arrays.

\sphinxAtStartPar

\sphinxAtStartPar
This function returns a pointer to the location of some global property
stored in one of the constituent classes of a LAMMPS instance.  The
returned pointer is cast to \sphinxcode{\sphinxupquote{void *}} and needs to be cast to a pointer
of the type that the entity represents. The pointers returned by this
function are generally persistent; therefore it is not necessary to call
the function again, unless a {\hyperref[\detokenize{clear::doc}]{\sphinxcrossref{\DUrole{doc}{clear command}}}} command is issued which wipes
out and recreates the contents of the {\hyperref[\detokenize{Classes_lammps:_CPPv4N9LAMMPS_NS6LAMMPSE}]{\sphinxcrossref{\sphinxcode{\sphinxupquote{LAMMPS}}}}} class.

\sphinxAtStartPar
Please also see {\hyperref[\detokenize{Library_properties:_CPPv422lammps_extract_settingPvPKc}]{\sphinxcrossref{\sphinxcode{\sphinxupquote{lammps\_extract\_setting()}}}}},
{\hyperref[\detokenize{Library_properties:_CPPv417lammps_get_thermoPvPKc}]{\sphinxcrossref{\sphinxcode{\sphinxupquote{lammps\_get\_thermo()}}}}}, and {\hyperref[\detokenize{Library_properties:_CPPv418lammps_extract_boxPvPdPdPdPdPdPiPi}]{\sphinxcrossref{\sphinxcode{\sphinxupquote{lammps\_extract\_box()}}}}}.

\sphinxAtStartPar
The following tables list the supported names, their data types, length
of the data area, and a short description.  The data type can also be
queried through calling {\hyperref[\detokenize{Library_properties:_CPPv430lammps_extract_global_datatypePvPKc}]{\sphinxcrossref{\sphinxcode{\sphinxupquote{lammps\_extract\_global\_datatype()}}}}}.
The \sphinxcode{\sphinxupquote{bigint}} type may be defined to be either an \sphinxcode{\sphinxupquote{int}} or an
\sphinxcode{\sphinxupquote{int64\_t}}.  This is set at {\hyperref[\detokenize{Build_settings:size}]{\sphinxcrossref{\DUrole{std,std-ref}{compile time}}}} of the LAMMPS
library and can be queried through calling
{\hyperref[\detokenize{Library_properties:_CPPv422lammps_extract_settingPvPKc}]{\sphinxcrossref{\sphinxcode{\sphinxupquote{lammps\_extract\_setting()}}}}}.
The function {\hyperref[\detokenize{Library_properties:_CPPv430lammps_extract_global_datatypePvPKc}]{\sphinxcrossref{\sphinxcode{\sphinxupquote{lammps\_extract\_global\_datatype()}}}}} will directly
report the “native” data type.  The following tables are provided:
\begin{itemize}
\item {} 
\sphinxAtStartPar
{\hyperref[\detokenize{Library_properties:extract-timestep-settings}]{\sphinxcrossref{\DUrole{std,std-ref}{Timestep settings}}}}

\item {} 
\sphinxAtStartPar
{\hyperref[\detokenize{Library_properties:extract-box-settings}]{\sphinxcrossref{\DUrole{std,std-ref}{Simulation box settings}}}}

\item {} 
\sphinxAtStartPar
{\hyperref[\detokenize{Library_properties:extract-system-settings}]{\sphinxcrossref{\DUrole{std,std-ref}{System property settings}}}}

\item {} 
\sphinxAtStartPar
{\hyperref[\detokenize{Library_properties:extract-git-settings}]{\sphinxcrossref{\DUrole{std,std-ref}{Git revision and version settings}}}}

\item {} 
\sphinxAtStartPar
{\hyperref[\detokenize{Library_properties:extract-unit-settings}]{\sphinxcrossref{\DUrole{std,std-ref}{Unit settings}}}}

\end{itemize}
\phantomsection\label{\detokenize{Library_properties:extract-timestep-settings}}
\sphinxAtStartPar
\sphinxstylestrong{Timestep settings}


\begin{savenotes}\sphinxattablestart
\sphinxthistablewithglobalstyle
\centering
\begin{tabular}[t]{\X{14}{100}\X{10}{100}\X{10}{100}\X{66}{100}}
\sphinxtoprule
\sphinxstyletheadfamily 
\sphinxAtStartPar
Name
&\sphinxstyletheadfamily 
\sphinxAtStartPar
Type
&\sphinxstyletheadfamily 
\sphinxAtStartPar
Length
&\sphinxstyletheadfamily 
\sphinxAtStartPar
Description
\\
\sphinxmidrule
\sphinxtableatstartofbodyhook
\sphinxAtStartPar
dt
&
\sphinxAtStartPar
double
&
\sphinxAtStartPar
1
&
\sphinxAtStartPar
length of the time step. See {\hyperref[\detokenize{timestep::doc}]{\sphinxcrossref{\DUrole{doc}{timestep command}}}}.
\\
\sphinxhline
\sphinxAtStartPar
ntimestep
&
\sphinxAtStartPar
bigint
&
\sphinxAtStartPar
1
&
\sphinxAtStartPar
current time step number. See {\hyperref[\detokenize{reset_timestep::doc}]{\sphinxcrossref{\DUrole{doc}{reset\_timestep command}}}}.
\\
\sphinxhline
\sphinxAtStartPar
atime
&
\sphinxAtStartPar
double
&
\sphinxAtStartPar
1
&
\sphinxAtStartPar
accumulated simulation time in time units.
\\
\sphinxhline
\sphinxAtStartPar
atimestep
&
\sphinxAtStartPar
bigint
&
\sphinxAtStartPar
1
&
\sphinxAtStartPar
the number of the timestep when “atime” was last updated.
\\
\sphinxhline
\sphinxAtStartPar
respa\_levels
&
\sphinxAtStartPar
int
&
\sphinxAtStartPar
1
&
\sphinxAtStartPar
\(N_{respa}\) = number of r\sphinxhyphen{}RESPA levels. See {\hyperref[\detokenize{run_style::doc}]{\sphinxcrossref{\DUrole{doc}{run\_style command}}}}.
\\
\sphinxhline
\sphinxAtStartPar
respa\_dt
&
\sphinxAtStartPar
double
&
\sphinxAtStartPar
\(N_{respa}\)
&
\sphinxAtStartPar
length of the time steps with r\sphinxhyphen{}RESPA. See {\hyperref[\detokenize{run_style::doc}]{\sphinxcrossref{\DUrole{doc}{run\_style command}}}}.
\\
\sphinxbottomrule
\end{tabular}
\sphinxtableafterendhook\par
\sphinxattableend\end{savenotes}
\phantomsection\label{\detokenize{Library_properties:extract-box-settings}}
\sphinxAtStartPar
\sphinxstylestrong{Simulation box settings}


\begin{savenotes}\sphinxattablestart
\sphinxthistablewithglobalstyle
\centering
\begin{tabular}[t]{\X{16}{100}\X{10}{100}\X{10}{100}\X{64}{100}}
\sphinxtoprule
\sphinxstyletheadfamily 
\sphinxAtStartPar
Name
&\sphinxstyletheadfamily 
\sphinxAtStartPar
Type
&\sphinxstyletheadfamily 
\sphinxAtStartPar
Length
&\sphinxstyletheadfamily 
\sphinxAtStartPar
Description
\\
\sphinxmidrule
\sphinxtableatstartofbodyhook
\sphinxAtStartPar
boxxhi
&
\sphinxAtStartPar
double
&
\sphinxAtStartPar
1
&
\sphinxAtStartPar
upper box boundary in x\sphinxhyphen{}direction; see {\hyperref[\detokenize{create_box::doc}]{\sphinxcrossref{\DUrole{doc}{create\_box command}}}}.
\\
\sphinxhline
\sphinxAtStartPar
boxylo
&
\sphinxAtStartPar
double
&
\sphinxAtStartPar
1
&
\sphinxAtStartPar
lower box boundary in y\sphinxhyphen{}direction; see {\hyperref[\detokenize{create_box::doc}]{\sphinxcrossref{\DUrole{doc}{create\_box command}}}}.
\\
\sphinxhline
\sphinxAtStartPar
boxyhi
&
\sphinxAtStartPar
double
&
\sphinxAtStartPar
1
&
\sphinxAtStartPar
upper box boundary in y\sphinxhyphen{}direction; see {\hyperref[\detokenize{create_box::doc}]{\sphinxcrossref{\DUrole{doc}{create\_box command}}}}.
\\
\sphinxhline
\sphinxAtStartPar
boxzlo
&
\sphinxAtStartPar
double
&
\sphinxAtStartPar
1
&
\sphinxAtStartPar
lower box boundary in z\sphinxhyphen{}direction; see {\hyperref[\detokenize{create_box::doc}]{\sphinxcrossref{\DUrole{doc}{create\_box command}}}}.
\\
\sphinxhline
\sphinxAtStartPar
boxzhi
&
\sphinxAtStartPar
double
&
\sphinxAtStartPar
1
&
\sphinxAtStartPar
upper box boundary in z\sphinxhyphen{}direction; see {\hyperref[\detokenize{create_box::doc}]{\sphinxcrossref{\DUrole{doc}{create\_box command}}}}.
\\
\sphinxhline
\sphinxAtStartPar
sublo
&
\sphinxAtStartPar
double
&
\sphinxAtStartPar
3
&
\sphinxAtStartPar
subbox lower boundaries
\\
\sphinxhline
\sphinxAtStartPar
subhi
&
\sphinxAtStartPar
double
&
\sphinxAtStartPar
3
&
\sphinxAtStartPar
subbox upper boundaries
\\
\sphinxhline
\sphinxAtStartPar
sublo\_lambda
&
\sphinxAtStartPar
double
&
\sphinxAtStartPar
3
&
\sphinxAtStartPar
subbox lower boundaries in fractional coordinates (for triclinic cells)
\\
\sphinxhline
\sphinxAtStartPar
subhi\_lambda
&
\sphinxAtStartPar
double
&
\sphinxAtStartPar
3
&
\sphinxAtStartPar
subbox upper boundaries in fractional coordinates (for triclinic cells)
\\
\sphinxhline
\sphinxAtStartPar
periodicity
&
\sphinxAtStartPar
int
&
\sphinxAtStartPar
3
&
\sphinxAtStartPar
0 if non\sphinxhyphen{}periodic, 1 if periodic for x, y, and z; see {\hyperref[\detokenize{boundary::doc}]{\sphinxcrossref{\DUrole{doc}{boundary command}}}}.
\\
\sphinxhline
\sphinxAtStartPar
triclinic
&
\sphinxAtStartPar
int
&
\sphinxAtStartPar
1
&
\sphinxAtStartPar
1 if box is triclinic, 0 if orthogonal; see {\hyperref[\detokenize{change_box::doc}]{\sphinxcrossref{\DUrole{doc}{change\_box command}}}}.
\\
\sphinxhline
\sphinxAtStartPar
xy
&
\sphinxAtStartPar
double
&
\sphinxAtStartPar
1
&
\sphinxAtStartPar
triclinic tilt factor; see {\hyperref[\detokenize{Howto_triclinic::doc}]{\sphinxcrossref{\DUrole{doc}{Triclinic (non\sphinxhyphen{}orthogonal) simulation boxes}}}}.
\\
\sphinxhline
\sphinxAtStartPar
yz
&
\sphinxAtStartPar
double
&
\sphinxAtStartPar
1
&
\sphinxAtStartPar
triclinic tilt factor; see {\hyperref[\detokenize{Howto_triclinic::doc}]{\sphinxcrossref{\DUrole{doc}{Triclinic (non\sphinxhyphen{}orthogonal) simulation boxes}}}}.
\\
\sphinxhline
\sphinxAtStartPar
xz
&
\sphinxAtStartPar
double
&
\sphinxAtStartPar
1
&
\sphinxAtStartPar
triclinic tilt factor; see {\hyperref[\detokenize{Howto_triclinic::doc}]{\sphinxcrossref{\DUrole{doc}{Triclinic (non\sphinxhyphen{}orthogonal) simulation boxes}}}}.
\\
\sphinxhline
\sphinxAtStartPar
xlattice
&
\sphinxAtStartPar
double
&
\sphinxAtStartPar
1
&
\sphinxAtStartPar
lattice spacing in x\sphinxhyphen{}direction; see {\hyperref[\detokenize{lattice::doc}]{\sphinxcrossref{\DUrole{doc}{lattice command}}}}.
\\
\sphinxhline
\sphinxAtStartPar
ylattice
&
\sphinxAtStartPar
double
&
\sphinxAtStartPar
1
&
\sphinxAtStartPar
lattice spacing in y\sphinxhyphen{}direction; see {\hyperref[\detokenize{lattice::doc}]{\sphinxcrossref{\DUrole{doc}{lattice command}}}}.
\\
\sphinxhline
\sphinxAtStartPar
zlattice
&
\sphinxAtStartPar
double
&
\sphinxAtStartPar
1
&
\sphinxAtStartPar
lattice spacing in z\sphinxhyphen{}direction; see {\hyperref[\detokenize{lattice::doc}]{\sphinxcrossref{\DUrole{doc}{lattice command}}}}.
\\
\sphinxhline
\sphinxAtStartPar
procgrid
&
\sphinxAtStartPar
int
&
\sphinxAtStartPar
3
&
\sphinxAtStartPar
processor count in x\sphinxhyphen{}, y\sphinxhyphen{}, and z\sphinxhyphen{} direction; see {\hyperref[\detokenize{processors::doc}]{\sphinxcrossref{\DUrole{doc}{processors command}}}}.
\\
\sphinxbottomrule
\end{tabular}
\sphinxtableafterendhook\par
\sphinxattableend\end{savenotes}
\phantomsection\label{\detokenize{Library_properties:extract-system-settings}}
\sphinxAtStartPar
\sphinxstylestrong{System property settings}


\begin{savenotes}\sphinxattablestart
\sphinxthistablewithglobalstyle
\centering
\begin{tabular}[t]{\X{18}{100}\X{12}{100}\X{12}{100}\X{58}{100}}
\sphinxtoprule
\sphinxstyletheadfamily 
\sphinxAtStartPar
Name
&\sphinxstyletheadfamily 
\sphinxAtStartPar
Type
&\sphinxstyletheadfamily 
\sphinxAtStartPar
Length
&\sphinxstyletheadfamily 
\sphinxAtStartPar
Description
\\
\sphinxmidrule
\sphinxtableatstartofbodyhook
\sphinxAtStartPar
ntypes
&
\sphinxAtStartPar
int
&
\sphinxAtStartPar
1
&
\sphinxAtStartPar
number of atom types
\\
\sphinxhline
\sphinxAtStartPar
nbonds
&
\sphinxAtStartPar
bigint
&
\sphinxAtStartPar
1
&
\sphinxAtStartPar
total number of bonds in the simulation.
\\
\sphinxhline
\sphinxAtStartPar
nangles
&
\sphinxAtStartPar
bigint
&
\sphinxAtStartPar
1
&
\sphinxAtStartPar
total number of angles in the simulation.
\\
\sphinxhline
\sphinxAtStartPar
ndihedrals
&
\sphinxAtStartPar
bigint
&
\sphinxAtStartPar
1
&
\sphinxAtStartPar
total number of dihedrals in the simulation.
\\
\sphinxhline
\sphinxAtStartPar
nimpropers
&
\sphinxAtStartPar
bigint
&
\sphinxAtStartPar
1
&
\sphinxAtStartPar
total number of impropers in the simulation.
\\
\sphinxhline
\sphinxAtStartPar
natoms
&
\sphinxAtStartPar
bigint
&
\sphinxAtStartPar
1
&
\sphinxAtStartPar
total number of atoms in the simulation.
\\
\sphinxhline
\sphinxAtStartPar
nlocal
&
\sphinxAtStartPar
int
&
\sphinxAtStartPar
1
&
\sphinxAtStartPar
number of “owned” atoms of the current MPI rank.
\\
\sphinxhline
\sphinxAtStartPar
nghost
&
\sphinxAtStartPar
int
&
\sphinxAtStartPar
1
&
\sphinxAtStartPar
number of “ghost” atoms of the current MPI rank.
\\
\sphinxhline
\sphinxAtStartPar
nmax
&
\sphinxAtStartPar
int
&
\sphinxAtStartPar
1
&
\sphinxAtStartPar
maximum of nlocal+nghost across all MPI ranks (for per\sphinxhyphen{}atom data array size).
\\
\sphinxhline
\sphinxAtStartPar
special\_lj
&
\sphinxAtStartPar
double
&
\sphinxAtStartPar
4
&
\sphinxAtStartPar
special {\hyperref[\detokenize{special_bonds::doc}]{\sphinxcrossref{\DUrole{doc}{pair weighting factors}}}} for LJ interactions (first element is always 1.0)
\\
\sphinxhline
\sphinxAtStartPar
special\_coul
&
\sphinxAtStartPar
double
&
\sphinxAtStartPar
4
&
\sphinxAtStartPar
special {\hyperref[\detokenize{special_bonds::doc}]{\sphinxcrossref{\DUrole{doc}{pair weighting factors}}}} for Coulomb interactions (first element is always 1.0)
\\
\sphinxhline
\sphinxAtStartPar
map\_style
&
\sphinxAtStartPar
int
&
\sphinxAtStartPar
1
&
\sphinxAtStartPar
{\hyperref[\detokenize{atom_modify::doc}]{\sphinxcrossref{\DUrole{doc}{atom map setting}}}}: 0 = none, 1 = array, 2 = hash, 3 = yes
\\
\sphinxhline
\sphinxAtStartPar
map\_tag\_max
&
\sphinxAtStartPar
int/bigint
&
\sphinxAtStartPar
1
&
\sphinxAtStartPar
largest atom ID that can be mapped to a local index (bigint with \sphinxhyphen{}DLAMMPS\_BIGBIG)
\\
\sphinxhline
\sphinxAtStartPar
sametag
&
\sphinxAtStartPar
int
&
\sphinxAtStartPar
variable
&
\sphinxAtStartPar
index of next local atom with the same ID in ascending order. \sphinxhyphen{}1 signals end.
\\
\sphinxhline
\sphinxAtStartPar
sortfreq
&
\sphinxAtStartPar
int
&
\sphinxAtStartPar
1
&
\sphinxAtStartPar
frequency of atom sorting. 0 means sorting is off.
\\
\sphinxhline
\sphinxAtStartPar
nextsort
&
\sphinxAtStartPar
bigint
&
\sphinxAtStartPar
1
&
\sphinxAtStartPar
timestep when atoms are sorted next
\\
\sphinxhline
\sphinxAtStartPar
q\_flag
&
\sphinxAtStartPar
int
&
\sphinxAtStartPar
1
&
\sphinxAtStartPar
\sphinxstylestrong{deprecated}. Use {\hyperref[\detokenize{Library_properties:_CPPv422lammps_extract_settingPvPKc}]{\sphinxcrossref{\sphinxcode{\sphinxupquote{lammps\_extract\_setting()}}}}} instead.
\\
\sphinxhline
\sphinxAtStartPar
atom\_style
&
\sphinxAtStartPar
char *
&
\sphinxAtStartPar
1
&
\sphinxAtStartPar
string with the current atom style.
\\
\sphinxhline
\sphinxAtStartPar
pair\_style
&
\sphinxAtStartPar
char *
&
\sphinxAtStartPar
1
&
\sphinxAtStartPar
string with the current pair style.
\\
\sphinxhline
\sphinxAtStartPar
bond\_style
&
\sphinxAtStartPar
char *
&
\sphinxAtStartPar
1
&
\sphinxAtStartPar
string with the current bond style.
\\
\sphinxhline
\sphinxAtStartPar
angle\_style
&
\sphinxAtStartPar
char *
&
\sphinxAtStartPar
1
&
\sphinxAtStartPar
string with the current angle style.
\\
\sphinxhline
\sphinxAtStartPar
dihedral\_style
&
\sphinxAtStartPar
char *
&
\sphinxAtStartPar
1
&
\sphinxAtStartPar
string with the current dihedral style.
\\
\sphinxhline
\sphinxAtStartPar
improper\_style
&
\sphinxAtStartPar
char *
&
\sphinxAtStartPar
1
&
\sphinxAtStartPar
string with the current improper style.
\\
\sphinxhline
\sphinxAtStartPar
kspace\_style
&
\sphinxAtStartPar
char *
&
\sphinxAtStartPar
1
&
\sphinxAtStartPar
string with the current KSpace style.
\\
\sphinxbottomrule
\end{tabular}
\sphinxtableafterendhook\par
\sphinxattableend\end{savenotes}
\phantomsection\label{\detokenize{Library_properties:extract-git-settings}}
\sphinxAtStartPar
\sphinxstylestrong{Git revision and version settings}


\begin{savenotes}\sphinxattablestart
\sphinxthistablewithglobalstyle
\centering
\begin{tabular}[t]{\X{16}{100}\X{14}{100}\X{10}{100}\X{60}{100}}
\sphinxtoprule
\sphinxstyletheadfamily 
\sphinxAtStartPar
Name
&\sphinxstyletheadfamily 
\sphinxAtStartPar
Type
&\sphinxstyletheadfamily 
\sphinxAtStartPar
Length
&\sphinxstyletheadfamily 
\sphinxAtStartPar
Description
\\
\sphinxmidrule
\sphinxtableatstartofbodyhook
\sphinxAtStartPar
git\_commit
&
\sphinxAtStartPar
const char *
&
\sphinxAtStartPar
1
&
\sphinxAtStartPar
Git commit hash for the LAMMPS version.
\\
\sphinxhline
\sphinxAtStartPar
git\_branch
&
\sphinxAtStartPar
const char *
&
\sphinxAtStartPar
1
&
\sphinxAtStartPar
Git branch for the LAMMPS version.
\\
\sphinxhline
\sphinxAtStartPar
git\_descriptor
&
\sphinxAtStartPar
const char *
&
\sphinxAtStartPar
1
&
\sphinxAtStartPar
Combined descriptor for the git revision
\\
\sphinxhline
\sphinxAtStartPar
lammps\_version
&
\sphinxAtStartPar
const char *
&
\sphinxAtStartPar
1
&
\sphinxAtStartPar
LAMMPS version string.
\\
\sphinxbottomrule
\end{tabular}
\sphinxtableafterendhook\par
\sphinxattableend\end{savenotes}
\phantomsection\label{\detokenize{Library_properties:extract-unit-settings}}
\sphinxAtStartPar
\sphinxstylestrong{Unit settings}


\begin{savenotes}\sphinxattablestart
\sphinxthistablewithglobalstyle
\centering
\begin{tabular}[t]{\X{16}{100}\X{12}{100}\X{10}{100}\X{62}{100}}
\sphinxtoprule
\sphinxstyletheadfamily 
\sphinxAtStartPar
Name
&\sphinxstyletheadfamily 
\sphinxAtStartPar
Type
&\sphinxstyletheadfamily 
\sphinxAtStartPar
Length
&\sphinxstyletheadfamily 
\sphinxAtStartPar
Description
\\
\sphinxmidrule
\sphinxtableatstartofbodyhook
\sphinxAtStartPar
units
&
\sphinxAtStartPar
char *
&
\sphinxAtStartPar
1
&
\sphinxAtStartPar
string with the current unit style. See {\hyperref[\detokenize{units::doc}]{\sphinxcrossref{\DUrole{doc}{units command}}}}.
\\
\sphinxhline
\sphinxAtStartPar
boltz
&
\sphinxAtStartPar
double
&
\sphinxAtStartPar
1
&
\sphinxAtStartPar
value of the “boltz” constant. See {\hyperref[\detokenize{units::doc}]{\sphinxcrossref{\DUrole{doc}{units command}}}}.
\\
\sphinxhline
\sphinxAtStartPar
hplanck
&
\sphinxAtStartPar
double
&
\sphinxAtStartPar
1
&
\sphinxAtStartPar
value of the “hplanck” constant. See {\hyperref[\detokenize{units::doc}]{\sphinxcrossref{\DUrole{doc}{units command}}}}.
\\
\sphinxhline
\sphinxAtStartPar
mvv2e
&
\sphinxAtStartPar
double
&
\sphinxAtStartPar
1
&
\sphinxAtStartPar
factor to convert \(\frac{1}{2}mv^2\) for a particle to
the current energy unit; See {\hyperref[\detokenize{units::doc}]{\sphinxcrossref{\DUrole{doc}{units command}}}}.
\\
\sphinxhline
\sphinxAtStartPar
ftm2v
&
\sphinxAtStartPar
double
&
\sphinxAtStartPar
1
&
\sphinxAtStartPar
(description missing) See {\hyperref[\detokenize{units::doc}]{\sphinxcrossref{\DUrole{doc}{units command}}}}.
\\
\sphinxhline
\sphinxAtStartPar
mv2d
&
\sphinxAtStartPar
double
&
\sphinxAtStartPar
1
&
\sphinxAtStartPar
(description missing) See {\hyperref[\detokenize{units::doc}]{\sphinxcrossref{\DUrole{doc}{units command}}}}.
\\
\sphinxhline
\sphinxAtStartPar
nktv2p
&
\sphinxAtStartPar
double
&
\sphinxAtStartPar
1
&
\sphinxAtStartPar
(description missing) See {\hyperref[\detokenize{units::doc}]{\sphinxcrossref{\DUrole{doc}{units command}}}}.
\\
\sphinxhline
\sphinxAtStartPar
qqr2e
&
\sphinxAtStartPar
double
&
\sphinxAtStartPar
1
&
\sphinxAtStartPar
factor to convert \(\frac{q_i q_j}{r}\) to energy units;
See {\hyperref[\detokenize{units::doc}]{\sphinxcrossref{\DUrole{doc}{units command}}}}.
\\
\sphinxhline
\sphinxAtStartPar
qe2f
&
\sphinxAtStartPar
double
&
\sphinxAtStartPar
1
&
\sphinxAtStartPar
(description missing) See {\hyperref[\detokenize{units::doc}]{\sphinxcrossref{\DUrole{doc}{units command}}}}.
\\
\sphinxhline
\sphinxAtStartPar
vxmu2f
&
\sphinxAtStartPar
double
&
\sphinxAtStartPar
1
&
\sphinxAtStartPar
(description missing) See {\hyperref[\detokenize{units::doc}]{\sphinxcrossref{\DUrole{doc}{units command}}}}.
\\
\sphinxhline
\sphinxAtStartPar
xxt2kmu
&
\sphinxAtStartPar
double
&
\sphinxAtStartPar
1
&
\sphinxAtStartPar
(description missing) See {\hyperref[\detokenize{units::doc}]{\sphinxcrossref{\DUrole{doc}{units command}}}}.
\\
\sphinxhline
\sphinxAtStartPar
dielectric
&
\sphinxAtStartPar
double
&
\sphinxAtStartPar
1
&
\sphinxAtStartPar
value of the dielectric constant. See {\hyperref[\detokenize{dielectric::doc}]{\sphinxcrossref{\DUrole{doc}{dielectric command}}}}.
\\
\sphinxhline
\sphinxAtStartPar
qqrd2e
&
\sphinxAtStartPar
double
&
\sphinxAtStartPar
1
&
\sphinxAtStartPar
(description missing) See {\hyperref[\detokenize{units::doc}]{\sphinxcrossref{\DUrole{doc}{units command}}}}.
\\
\sphinxhline
\sphinxAtStartPar
e\_mass
&
\sphinxAtStartPar
double
&
\sphinxAtStartPar
1
&
\sphinxAtStartPar
(description missing) See {\hyperref[\detokenize{units::doc}]{\sphinxcrossref{\DUrole{doc}{units command}}}}.
\\
\sphinxhline
\sphinxAtStartPar
hhmrr2e
&
\sphinxAtStartPar
double
&
\sphinxAtStartPar
1
&
\sphinxAtStartPar
(description missing) See {\hyperref[\detokenize{units::doc}]{\sphinxcrossref{\DUrole{doc}{units command}}}}.
\\
\sphinxhline
\sphinxAtStartPar
mvh2r
&
\sphinxAtStartPar
double
&
\sphinxAtStartPar
1
&
\sphinxAtStartPar
(description missing) See {\hyperref[\detokenize{units::doc}]{\sphinxcrossref{\DUrole{doc}{units command}}}}.
\\
\sphinxhline
\sphinxAtStartPar
angstrom
&
\sphinxAtStartPar
double
&
\sphinxAtStartPar
1
&
\sphinxAtStartPar
constant to convert current length unit to angstroms;
1.0 for reduced (aka “lj”) units. See {\hyperref[\detokenize{units::doc}]{\sphinxcrossref{\DUrole{doc}{units command}}}}.
\\
\sphinxhline
\sphinxAtStartPar
femtosecond
&
\sphinxAtStartPar
double
&
\sphinxAtStartPar
1
&
\sphinxAtStartPar
constant to convert current time unit to femtoseconds;
1.0 for reduced (aka “lj”) units
\\
\sphinxhline
\sphinxAtStartPar
qelectron
&
\sphinxAtStartPar
double
&
\sphinxAtStartPar
1
&
\sphinxAtStartPar
(description missing) See {\hyperref[\detokenize{units::doc}]{\sphinxcrossref{\DUrole{doc}{units command}}}}.
\\
\sphinxbottomrule
\end{tabular}
\sphinxtableafterendhook\par
\sphinxattableend\end{savenotes}


\begin{sphinxadmonition}{warning}{Warning:}
\sphinxAtStartPar
Modifying the data in the location pointed to by the returned pointer
may lead to inconsistent internal data and thus may cause failures or
crashes or bogus simulations.  In general it is thus usually better
to use a LAMMPS input command that sets or changes these parameters.
Those will take care of all side effects and necessary updates of
settings derived from such settings.  Where possible, a reference to
such a command or a relevant section of the manual is given below.
\end{sphinxadmonition}
\begin{quote}\begin{description}
\sphinxlineitem{Parameters}\begin{itemize}
\item {} 
\sphinxAtStartPar
\sphinxstyleliteralstrong{\sphinxupquote{handle}} \textendash{} pointer to a previously created LAMMPS instance 

\item {} 
\sphinxAtStartPar
\sphinxstyleliteralstrong{\sphinxupquote{name}} \textendash{} string with the name of the extracted property 

\end{itemize}

\sphinxlineitem{Returns}
\sphinxAtStartPar
pointer (cast to \sphinxcode{\sphinxupquote{void *}}) to the location of the requested property. NULL if name is not known. 

\end{description}\end{quote}

\end{fulllineitems}



\bigskip\hrule\bigskip

\index{lammps\_extract\_pair\_dimension (C++ function)@\spxentry{lammps\_extract\_pair\_dimension}\spxextra{C++ function}}

\begin{fulllineitems}
\phantomsection\label{\detokenize{Library_properties:_CPPv429lammps_extract_pair_dimensionPvPKc}}\phantomsection\label{\detokenize{Library_properties:_CPPv329lammps_extract_pair_dimensionPvPKc}}\phantomsection\label{\detokenize{Library_properties:_CPPv229lammps_extract_pair_dimensionPvPKc}}\phantomsection\label{\detokenize{Library_properties:lammps_extract_pair_dimension__voidP.cCP}}
\pysigstartsignatures
\pysigstartmultiline
\pysiglinewithargsret{\phantomsection\label{\detokenize{Library_properties:library_8h_1a8935ce652b4140c850625040c71db8c9}}\DUrole{kt}{int}\DUrole{w}{ }\sphinxbfcode{\sphinxupquote{\DUrole{n}{lammps\_extract\_pair\_dimension}}}}{\DUrole{kt}{void}\DUrole{w}{ }\DUrole{p}{*}\DUrole{n,sig-param}{handle}\sphinxparamcomma \DUrole{k}{const}\DUrole{w}{ }\DUrole{kt}{char}\DUrole{w}{ }\DUrole{p}{*}\DUrole{n,sig-param}{name}}{}
\pysigstopmultiline
\pysigstopsignatures
\sphinxAtStartPar
Get data dimension of pair style data accessible via Pair::extract().

\sphinxAtStartPar

\sphinxAtStartPar
\DUrole{versionmodified,added}{Added in version 29Aug2024.}

\sphinxAtStartPar
This function returns an integer that specified the dimensionality of
the data that can be extracted from the current pair style with \sphinxcode{\sphinxupquote{Pair::extract()}}.
Callers of {\hyperref[\detokenize{Library_properties:_CPPv419lammps_extract_pairPvPKc}]{\sphinxcrossref{\sphinxcode{\sphinxupquote{lammps\_extract\_pair()}}}}} can use this information
to then decide how to cast the \sphinxcode{\sphinxupquote{void *}} pointer and access the data.

\begin{quote}\begin{description}
\sphinxlineitem{Parameters}\begin{itemize}
\item {} 
\sphinxAtStartPar
\sphinxstyleliteralstrong{\sphinxupquote{handle}} \textendash{} pointer to a previously created LAMMPS instance 

\item {} 
\sphinxAtStartPar
\sphinxstyleliteralstrong{\sphinxupquote{name}} \textendash{} string with the name of the extracted property 

\end{itemize}

\sphinxlineitem{Returns}
\sphinxAtStartPar
integer constant encoding the dimensionality of the extractable pair style property or \sphinxhyphen{}1 if not found. 

\end{description}\end{quote}

\end{fulllineitems}



\bigskip\hrule\bigskip

\index{lammps\_extract\_pair (C++ function)@\spxentry{lammps\_extract\_pair}\spxextra{C++ function}}

\begin{fulllineitems}
\phantomsection\label{\detokenize{Library_properties:_CPPv419lammps_extract_pairPvPKc}}\phantomsection\label{\detokenize{Library_properties:_CPPv319lammps_extract_pairPvPKc}}\phantomsection\label{\detokenize{Library_properties:_CPPv219lammps_extract_pairPvPKc}}\phantomsection\label{\detokenize{Library_properties:lammps_extract_pair__voidP.cCP}}
\pysigstartsignatures
\pysigstartmultiline
\pysiglinewithargsret{\phantomsection\label{\detokenize{Library_properties:library_8h_1ab4982a565b8a85491b403e57f6e16d56}}\DUrole{kt}{void}\DUrole{w}{ }\DUrole{p}{*}\sphinxbfcode{\sphinxupquote{\DUrole{n}{lammps\_extract\_pair}}}}{\DUrole{kt}{void}\DUrole{w}{ }\DUrole{p}{*}\DUrole{n,sig-param}{handle}\sphinxparamcomma \DUrole{k}{const}\DUrole{w}{ }\DUrole{kt}{char}\DUrole{w}{ }\DUrole{p}{*}\DUrole{n,sig-param}{name}}{}
\pysigstopmultiline
\pysigstopsignatures
\sphinxAtStartPar
Get extract pair style data accessible via Pair::extract().

\sphinxAtStartPar

\sphinxAtStartPar
\DUrole{versionmodified,added}{Added in version 29Aug2024.}

\sphinxAtStartPar
This function returns a pointer to data available from the current pair
style with \sphinxcode{\sphinxupquote{Pair::extract()}}. The dimensionality of the returned
pointer can be determined with {\hyperref[\detokenize{Library_properties:_CPPv429lammps_extract_pair_dimensionPvPKc}]{\sphinxcrossref{\sphinxcode{\sphinxupquote{lammps\_extract\_pair\_dimension()}}}}}.

\begin{quote}\begin{description}
\sphinxlineitem{Parameters}\begin{itemize}
\item {} 
\sphinxAtStartPar
\sphinxstyleliteralstrong{\sphinxupquote{handle}} \textendash{} pointer to a previously created LAMMPS instance 

\item {} 
\sphinxAtStartPar
\sphinxstyleliteralstrong{\sphinxupquote{name}} \textendash{} string with the name of the extracted property 

\end{itemize}

\sphinxlineitem{Returns}
\sphinxAtStartPar
pointer (cast to \sphinxcode{\sphinxupquote{void *}}) to the location of the requested property. NULL if name is not known. 

\end{description}\end{quote}

\end{fulllineitems}



\bigskip\hrule\bigskip

\index{lammps\_map\_atom (C++ function)@\spxentry{lammps\_map\_atom}\spxextra{C++ function}}

\begin{fulllineitems}
\phantomsection\label{\detokenize{Library_properties:_CPPv415lammps_map_atomPvPKv}}\phantomsection\label{\detokenize{Library_properties:_CPPv315lammps_map_atomPvPKv}}\phantomsection\label{\detokenize{Library_properties:_CPPv215lammps_map_atomPvPKv}}\phantomsection\label{\detokenize{Library_properties:lammps_map_atom__voidP.voidCP}}
\pysigstartsignatures
\pysigstartmultiline
\pysiglinewithargsret{\phantomsection\label{\detokenize{Library_properties:library_8h_1a40e85a5c09ff9cc7763d119fea9603ec}}\DUrole{kt}{int}\DUrole{w}{ }\sphinxbfcode{\sphinxupquote{\DUrole{n}{lammps\_map\_atom}}}}{\DUrole{kt}{void}\DUrole{w}{ }\DUrole{p}{*}\DUrole{n,sig-param}{handle}\sphinxparamcomma \DUrole{k}{const}\DUrole{w}{ }\DUrole{kt}{void}\DUrole{w}{ }\DUrole{p}{*}\DUrole{n,sig-param}{id}}{}
\pysigstopmultiline
\pysigstopsignatures
\sphinxAtStartPar
Map global atom ID to local atom index

\sphinxAtStartPar

\sphinxAtStartPar
\DUrole{versionmodified,added}{Added in version 27June2024.}

\sphinxAtStartPar
This function returns an integer that corresponds to the local atom
index for an atom with the global atom ID \sphinxstyleemphasis{id}. The atom ID is passed
as a void pointer so that it can use the same interface for either a
32\sphinxhyphen{}bit or 64\sphinxhyphen{}bit tagint. The size of the tagint can be determined
using {\hyperref[\detokenize{Library_properties:_CPPv422lammps_extract_settingPvPKc}]{\sphinxcrossref{\sphinxcode{\sphinxupquote{lammps\_extract\_setting()}}}}}.

\begin{quote}\begin{description}
\sphinxlineitem{Parameters}\begin{itemize}
\item {} 
\sphinxAtStartPar
\sphinxstyleliteralstrong{\sphinxupquote{handle}} \textendash{} pointer to a previously created LAMMPS instance 

\item {} 
\sphinxAtStartPar
\sphinxstyleliteralstrong{\sphinxupquote{id}} \textendash{} void pointer to the atom ID (of data type tagint, i.e. 32\sphinxhyphen{}bit or 64\sphinxhyphen{}bit integer) 

\end{itemize}

\sphinxlineitem{Returns}
\sphinxAtStartPar
local atom index or \sphinxhyphen{}1 if the atom is not found or no map exists 

\end{description}\end{quote}

\end{fulllineitems}


\sphinxstepscope


\subsection{Per\sphinxhyphen{}atom properties}
\label{\detokenize{Library_atoms:per-atom-properties}}\label{\detokenize{Library_atoms::doc}}
\sphinxAtStartPar
This section documents the following functions:
\begin{itemize}
\item {} 
\sphinxAtStartPar
{\hyperref[\detokenize{Library_atoms:_CPPv428lammps_extract_atom_datatypePvPKc}]{\sphinxcrossref{\sphinxcode{\sphinxupquote{lammps\_extract\_atom\_datatype()}}}}}

\item {} 
\sphinxAtStartPar
{\hyperref[\detokenize{Library_atoms:_CPPv419lammps_extract_atomPvPKc}]{\sphinxcrossref{\sphinxcode{\sphinxupquote{lammps\_extract\_atom()}}}}}

\end{itemize}


\bigskip\hrule\bigskip

\index{lammps\_extract\_atom\_datatype (C++ function)@\spxentry{lammps\_extract\_atom\_datatype}\spxextra{C++ function}}

\begin{fulllineitems}
\phantomsection\label{\detokenize{Library_atoms:_CPPv428lammps_extract_atom_datatypePvPKc}}\phantomsection\label{\detokenize{Library_atoms:_CPPv328lammps_extract_atom_datatypePvPKc}}\phantomsection\label{\detokenize{Library_atoms:_CPPv228lammps_extract_atom_datatypePvPKc}}\phantomsection\label{\detokenize{Library_atoms:lammps_extract_atom_datatype__voidP.cCP}}
\pysigstartsignatures
\pysigstartmultiline
\pysiglinewithargsret{\phantomsection\label{\detokenize{Library_atoms:library_8h_1aed36a77c70d91c9497b5eee075d18aab}}\DUrole{kt}{int}\DUrole{w}{ }\sphinxbfcode{\sphinxupquote{\DUrole{n}{lammps\_extract\_atom\_datatype}}}}{\DUrole{kt}{void}\DUrole{w}{ }\DUrole{p}{*}\DUrole{n,sig-param}{handle}\sphinxparamcomma \DUrole{k}{const}\DUrole{w}{ }\DUrole{kt}{char}\DUrole{w}{ }\DUrole{p}{*}\DUrole{n,sig-param}{name}}{}
\pysigstopmultiline
\pysigstopsignatures
\sphinxAtStartPar
Get data type of a LAMMPS per\sphinxhyphen{}atom property

\sphinxAtStartPar

\sphinxAtStartPar
\DUrole{versionmodified,added}{Added in version 18Sep2020.}

\sphinxAtStartPar
This function returns an integer that encodes the data type of the per\sphinxhyphen{}atom
property with the specified name. See \sphinxcode{\sphinxupquote{\_LMP\_DATATYPE\_CONST}} for valid
values. Callers of {\hyperref[\detokenize{Library_atoms:_CPPv419lammps_extract_atomPvPKc}]{\sphinxcrossref{\sphinxcode{\sphinxupquote{lammps\_extract\_atom()}}}}} can use this information
to then decide how to cast the \sphinxcode{\sphinxupquote{void *}} pointer and access the data.

\begin{quote}\begin{description}
\sphinxlineitem{Parameters}\begin{itemize}
\item {} 
\sphinxAtStartPar
\sphinxstyleliteralstrong{\sphinxupquote{handle}} \textendash{} pointer to a previously created LAMMPS instance 

\item {} 
\sphinxAtStartPar
\sphinxstyleliteralstrong{\sphinxupquote{name}} \textendash{} string with the name of the extracted property 

\end{itemize}

\sphinxlineitem{Returns}
\sphinxAtStartPar
integer constant encoding the data type of the property or \sphinxhyphen{}1 if not found. 

\end{description}\end{quote}

\end{fulllineitems}



\bigskip\hrule\bigskip

\index{lammps\_extract\_atom (C++ function)@\spxentry{lammps\_extract\_atom}\spxextra{C++ function}}

\begin{fulllineitems}
\phantomsection\label{\detokenize{Library_atoms:_CPPv419lammps_extract_atomPvPKc}}\phantomsection\label{\detokenize{Library_atoms:_CPPv319lammps_extract_atomPvPKc}}\phantomsection\label{\detokenize{Library_atoms:_CPPv219lammps_extract_atomPvPKc}}\phantomsection\label{\detokenize{Library_atoms:lammps_extract_atom__voidP.cCP}}
\pysigstartsignatures
\pysigstartmultiline
\pysiglinewithargsret{\phantomsection\label{\detokenize{Library_atoms:library_8h_1ac7a2b3791b785a29e1cc383dec467819}}\DUrole{kt}{void}\DUrole{w}{ }\DUrole{p}{*}\sphinxbfcode{\sphinxupquote{\DUrole{n}{lammps\_extract\_atom}}}}{\DUrole{kt}{void}\DUrole{w}{ }\DUrole{p}{*}\DUrole{n,sig-param}{handle}\sphinxparamcomma \DUrole{k}{const}\DUrole{w}{ }\DUrole{kt}{char}\DUrole{w}{ }\DUrole{p}{*}\DUrole{n,sig-param}{name}}{}
\pysigstopmultiline
\pysigstopsignatures
\sphinxAtStartPar
Get pointer to a LAMMPS per\sphinxhyphen{}atom property.

\sphinxAtStartPar

\sphinxAtStartPar
This function returns a pointer to the location of per\sphinxhyphen{}atom properties
(and per\sphinxhyphen{}atom\sphinxhyphen{}type properties in the case of the ‘mass’ keyword).
Per\sphinxhyphen{}atom data is distributed across sub\sphinxhyphen{}domains and thus MPI ranks.  The
returned pointer is cast to \sphinxcode{\sphinxupquote{void *}} and needs to be cast to a pointer
of data type that the entity represents.

\sphinxAtStartPar
A table with supported keywords is included in the documentation
of the {\hyperref[\detokenize{Classes_atom:_CPPv4N9LAMMPS_NS4Atom7extractEPKc}]{\sphinxcrossref{\sphinxcode{\sphinxupquote{Atom::extract()}}}}} function.


\begin{sphinxadmonition}{warning}{Warning:}
\sphinxAtStartPar
The pointers returned by this function are generally not persistent
since per\sphinxhyphen{}atom data may be re\sphinxhyphen{}distributed, re\sphinxhyphen{}allocated, and
re\sphinxhyphen{}ordered at every re\sphinxhyphen{}neighboring operation.
\end{sphinxadmonition}
\begin{quote}\begin{description}
\sphinxlineitem{Parameters}\begin{itemize}
\item {} 
\sphinxAtStartPar
\sphinxstyleliteralstrong{\sphinxupquote{handle}} \textendash{} pointer to a previously created LAMMPS instance 

\item {} 
\sphinxAtStartPar
\sphinxstyleliteralstrong{\sphinxupquote{name}} \textendash{} string with the name of the extracted property 

\end{itemize}

\sphinxlineitem{Returns}
\sphinxAtStartPar
pointer (cast to \sphinxcode{\sphinxupquote{void *}}) to the location of the requested data or \sphinxcode{\sphinxupquote{NULL}} if not found. 

\end{description}\end{quote}

\end{fulllineitems}


\sphinxstepscope


\subsection{Compute, fixes, variables}
\label{\detokenize{Library_objects:compute-fixes-variables}}\label{\detokenize{Library_objects::doc}}
\sphinxAtStartPar
This section documents accessing or modifying data stored by computes,
fixes, or variables in LAMMPS using the following functions:
\begin{itemize}
\item {} 
\sphinxAtStartPar
{\hyperref[\detokenize{Library_objects:_CPPv422lammps_extract_computePvPKcii}]{\sphinxcrossref{\sphinxcode{\sphinxupquote{lammps\_extract\_compute()}}}}}

\item {} 
\sphinxAtStartPar
{\hyperref[\detokenize{Library_objects:_CPPv418lammps_extract_fixPvPKciiii}]{\sphinxcrossref{\sphinxcode{\sphinxupquote{lammps\_extract\_fix()}}}}}

\item {} 
\sphinxAtStartPar
{\hyperref[\detokenize{Library_objects:_CPPv432lammps_extract_variable_datatypePvPKc}]{\sphinxcrossref{\sphinxcode{\sphinxupquote{lammps\_extract\_variable\_datatype()}}}}}

\item {} 
\sphinxAtStartPar
{\hyperref[\detokenize{Library_objects:_CPPv423lammps_extract_variablePvPKcPKc}]{\sphinxcrossref{\sphinxcode{\sphinxupquote{lammps\_extract\_variable()}}}}}

\item {} 
\sphinxAtStartPar
{\hyperref[\detokenize{Library_objects:_CPPv419lammps_set_variablePvPKcPKc}]{\sphinxcrossref{\sphinxcode{\sphinxupquote{lammps\_set\_variable()}}}}}

\item {} 
\sphinxAtStartPar
{\hyperref[\detokenize{Library_objects:_CPPv426lammps_set_string_variablePvPKcPKc}]{\sphinxcrossref{\sphinxcode{\sphinxupquote{lammps\_set\_string\_variable()}}}}}

\item {} 
\sphinxAtStartPar
{\hyperref[\detokenize{Library_objects:_CPPv428lammps_set_internal_variablePvPKcd}]{\sphinxcrossref{\sphinxcode{\sphinxupquote{lammps\_set\_internal\_variable()}}}}}

\item {} 
\sphinxAtStartPar
{\hyperref[\detokenize{Library_objects:_CPPv420lammps_variable_infoPviPci}]{\sphinxcrossref{\sphinxcode{\sphinxupquote{lammps\_variable\_info()}}}}}

\end{itemize}


\bigskip\hrule\bigskip

\index{lammps\_extract\_compute (C++ function)@\spxentry{lammps\_extract\_compute}\spxextra{C++ function}}

\begin{fulllineitems}
\phantomsection\label{\detokenize{Library_objects:_CPPv422lammps_extract_computePvPKcii}}\phantomsection\label{\detokenize{Library_objects:_CPPv322lammps_extract_computePvPKcii}}\phantomsection\label{\detokenize{Library_objects:_CPPv222lammps_extract_computePvPKcii}}\phantomsection\label{\detokenize{Library_objects:lammps_extract_compute__voidP.cCP.i.i}}
\pysigstartsignatures
\pysigstartmultiline
\pysiglinewithargsret{\phantomsection\label{\detokenize{Library_objects:library_8h_1a494c3aa6290bd2f3c05b3e5f65862e32}}\DUrole{kt}{void}\DUrole{w}{ }\DUrole{p}{*}\sphinxbfcode{\sphinxupquote{\DUrole{n}{lammps\_extract\_compute}}}}{\DUrole{kt}{void}\DUrole{w}{ }\DUrole{p}{*}\DUrole{n,sig-param}{handle}\sphinxparamcomma \DUrole{k}{const}\DUrole{w}{ }\DUrole{kt}{char}\DUrole{w}{ }\DUrole{p}{*}\DUrole{n,sig-param}{id}\sphinxparamcomma \DUrole{kt}{int}\DUrole{w}{ }\DUrole{n,sig-param}{style}\sphinxparamcomma \DUrole{kt}{int}\DUrole{w}{ }\DUrole{n,sig-param}{type}}{}
\pysigstopmultiline
\pysigstopsignatures
\sphinxAtStartPar
Get pointer to data from a LAMMPS compute.

\sphinxAtStartPar

\sphinxAtStartPar
This function returns a pointer to the location of data provided by a
{\hyperref[\detokenize{compute::doc}]{\sphinxcrossref{\DUrole{doc}{compute command}}}} instance identified by the compute\sphinxhyphen{}ID.  Computes may
provide global, per\sphinxhyphen{}atom, or local data, and those may be a scalar, a
vector, or an array or they may provide the information about the
dimensions of the respective data.  Since computes may provide multiple
kinds of data, it is required to set style and type flags representing
what specific data is desired.  This also determines to what kind of
pointer the returned pointer needs to be cast to access the data
correctly.  The function returns \sphinxcode{\sphinxupquote{NULL}} if the compute ID is not found
or the requested data is not available or current. The following table
lists the available options.


\begin{savenotes}\sphinxattablestart
\sphinxthistablewithglobalstyle
\centering
\begin{tabular}[t]{\X{25}{100}\X{24}{100}\X{14}{100}\X{37}{100}}
\sphinxtoprule
\sphinxstyletheadfamily 
\sphinxAtStartPar
Style (see {\hyperref[\detokenize{Library_objects:_CPPv416_LMP_STYLE_CONST}]{\sphinxcrossref{\sphinxcode{\sphinxupquote{\_LMP\_STYLE\_CONST}}}}})
&\sphinxstyletheadfamily 
\sphinxAtStartPar
Type (see {\hyperref[\detokenize{Library_objects:_CPPv415_LMP_TYPE_CONST}]{\sphinxcrossref{\sphinxcode{\sphinxupquote{\_LMP\_TYPE\_CONST}}}}})
&\sphinxstyletheadfamily 
\sphinxAtStartPar
Returned type
&\sphinxstyletheadfamily 
\sphinxAtStartPar
Returned data
\\
\sphinxmidrule
\sphinxtableatstartofbodyhook
\sphinxAtStartPar
LMP\_STYLE\_GLOBAL
&
\sphinxAtStartPar
LMP\_TYPE\_SCALAR
&
\sphinxAtStartPar
\sphinxcode{\sphinxupquote{double *}}
&
\sphinxAtStartPar
Global scalar
\\
\sphinxhline
\sphinxAtStartPar
LMP\_STYLE\_GLOBAL
&
\sphinxAtStartPar
LMP\_TYPE\_VECTOR
&
\sphinxAtStartPar
\sphinxcode{\sphinxupquote{double *}}
&
\sphinxAtStartPar
Global vector
\\
\sphinxhline
\sphinxAtStartPar
LMP\_STYLE\_GLOBAL
&
\sphinxAtStartPar
LMP\_TYPE\_ARRAY
&
\sphinxAtStartPar
\sphinxcode{\sphinxupquote{double **}}
&
\sphinxAtStartPar
Global array
\\
\sphinxhline
\sphinxAtStartPar
LMP\_STYLE\_GLOBAL
&
\sphinxAtStartPar
LMP\_SIZE\_VECTOR
&
\sphinxAtStartPar
\sphinxcode{\sphinxupquote{int *}}
&
\sphinxAtStartPar
Length of global vector
\\
\sphinxhline
\sphinxAtStartPar
LMP\_STYLE\_GLOBAL
&
\sphinxAtStartPar
LMP\_SIZE\_ROWS
&
\sphinxAtStartPar
\sphinxcode{\sphinxupquote{int *}}
&
\sphinxAtStartPar
Rows of global array
\\
\sphinxhline
\sphinxAtStartPar
LMP\_STYLE\_GLOBAL
&
\sphinxAtStartPar
LMP\_SIZE\_COLS
&
\sphinxAtStartPar
\sphinxcode{\sphinxupquote{int *}}
&
\sphinxAtStartPar
Columns of global array
\\
\sphinxhline
\sphinxAtStartPar
LMP\_STYLE\_ATOM
&
\sphinxAtStartPar
LMP\_TYPE\_VECTOR
&
\sphinxAtStartPar
\sphinxcode{\sphinxupquote{double *}}
&
\sphinxAtStartPar
Per\sphinxhyphen{}atom value
\\
\sphinxhline
\sphinxAtStartPar
LMP\_STYLE\_ATOM
&
\sphinxAtStartPar
LMP\_TYPE\_ARRAY
&
\sphinxAtStartPar
\sphinxcode{\sphinxupquote{double **}}
&
\sphinxAtStartPar
Per\sphinxhyphen{}atom vector
\\
\sphinxhline
\sphinxAtStartPar
LMP\_STYLE\_ATOM
&
\sphinxAtStartPar
LMP\_SIZE\_COLS
&
\sphinxAtStartPar
\sphinxcode{\sphinxupquote{int *}}
&
\sphinxAtStartPar
Columns in per\sphinxhyphen{}atom array, 0 if vector
\\
\sphinxhline
\sphinxAtStartPar
LMP\_STYLE\_LOCAL
&
\sphinxAtStartPar
LMP\_TYPE\_VECTOR
&
\sphinxAtStartPar
\sphinxcode{\sphinxupquote{double *}}
&
\sphinxAtStartPar
Local data vector
\\
\sphinxhline
\sphinxAtStartPar
LMP\_STYLE\_LOCAL
&
\sphinxAtStartPar
LMP\_TYPE\_ARRAY
&
\sphinxAtStartPar
\sphinxcode{\sphinxupquote{double **}}
&
\sphinxAtStartPar
Local data array
\\
\sphinxhline
\sphinxAtStartPar
LMP\_STYLE\_LOCAL
&
\sphinxAtStartPar
LMP\_SIZE\_VECTOR
&
\sphinxAtStartPar
\sphinxcode{\sphinxupquote{int *}}
&
\sphinxAtStartPar
Alias for LMP\_SIZE\_ROWS
\\
\sphinxhline
\sphinxAtStartPar
LMP\_STYLE\_LOCAL
&
\sphinxAtStartPar
LMP\_SIZE\_ROWS
&
\sphinxAtStartPar
\sphinxcode{\sphinxupquote{int *}}
&
\sphinxAtStartPar
Number of local array rows or length of vector
\\
\sphinxhline
\sphinxAtStartPar
LMP\_STYLE\_LOCAL
&
\sphinxAtStartPar
LMP\_SIZE\_COLS
&
\sphinxAtStartPar
\sphinxcode{\sphinxupquote{int *}}
&
\sphinxAtStartPar
Number of local array columns, 0 if vector
\\
\sphinxbottomrule
\end{tabular}
\sphinxtableafterendhook\par
\sphinxattableend\end{savenotes}


\begin{sphinxadmonition}{note}{Note:}
\sphinxAtStartPar
If the compute’s data is not computed for the current step, the
compute will be invoked.  LAMMPS cannot easily check at that time, if
it is valid to invoke a compute, so it may fail with an error.  The
caller has to check to avoid such an error.
\end{sphinxadmonition}

\begin{sphinxadmonition}{warning}{Warning:}
\sphinxAtStartPar
The pointers returned by this function are generally not persistent
since the computed data may be re\sphinxhyphen{}distributed, re\sphinxhyphen{}allocated, and
re\sphinxhyphen{}ordered at every invocation. It is advisable to re\sphinxhyphen{}invoke this
function before the data is accessed, or make a copy if the data shall
be used after other LAMMPS commands have been issued.
\end{sphinxadmonition}
\begin{quote}\begin{description}
\sphinxlineitem{Parameters}\begin{itemize}
\item {} 
\sphinxAtStartPar
\sphinxstyleliteralstrong{\sphinxupquote{handle}} \textendash{} pointer to a previously created LAMMPS instance 

\item {} 
\sphinxAtStartPar
\sphinxstyleliteralstrong{\sphinxupquote{id}} \textendash{} string with ID of the compute 

\item {} 
\sphinxAtStartPar
\sphinxstyleliteralstrong{\sphinxupquote{style}} \textendash{} constant indicating the style of data requested (global, per\sphinxhyphen{}atom, or local) 

\item {} 
\sphinxAtStartPar
\sphinxstyleliteralstrong{\sphinxupquote{type}} \textendash{} constant indicating type of data (scalar, vector, or array) or size of rows or columns 

\end{itemize}

\sphinxlineitem{Returns}
\sphinxAtStartPar
pointer (cast to \sphinxcode{\sphinxupquote{void *}}) to the location of the requested data or \sphinxcode{\sphinxupquote{NULL}} if not found. 

\end{description}\end{quote}

\end{fulllineitems}



\bigskip\hrule\bigskip

\index{lammps\_extract\_fix (C++ function)@\spxentry{lammps\_extract\_fix}\spxextra{C++ function}}

\begin{fulllineitems}
\phantomsection\label{\detokenize{Library_objects:_CPPv418lammps_extract_fixPvPKciiii}}\phantomsection\label{\detokenize{Library_objects:_CPPv318lammps_extract_fixPvPKciiii}}\phantomsection\label{\detokenize{Library_objects:_CPPv218lammps_extract_fixPvPKciiii}}\phantomsection\label{\detokenize{Library_objects:lammps_extract_fix__voidP.cCP.i.i.i.i}}
\pysigstartsignatures
\pysigstartmultiline
\pysiglinewithargsret{\phantomsection\label{\detokenize{Library_objects:library_8h_1abf1f4c59970775f24893ad1eb41e35a2}}\DUrole{kt}{void}\DUrole{w}{ }\DUrole{p}{*}\sphinxbfcode{\sphinxupquote{\DUrole{n}{lammps\_extract\_fix}}}}{\DUrole{kt}{void}\DUrole{w}{ }\DUrole{p}{*}\DUrole{n,sig-param}{handle}\sphinxparamcomma \DUrole{k}{const}\DUrole{w}{ }\DUrole{kt}{char}\DUrole{w}{ }\DUrole{p}{*}\DUrole{n,sig-param}{id}\sphinxparamcomma \DUrole{kt}{int}\DUrole{w}{ }\DUrole{n,sig-param}{style}\sphinxparamcomma \DUrole{kt}{int}\DUrole{w}{ }\DUrole{n,sig-param}{type}\sphinxparamcomma \DUrole{kt}{int}\DUrole{w}{ }\DUrole{n,sig-param}{nrow}\sphinxparamcomma \DUrole{kt}{int}\DUrole{w}{ }\DUrole{n,sig-param}{ncol}}{}
\pysigstopmultiline
\pysigstopsignatures
\sphinxAtStartPar
Get pointer to data from a LAMMPS fix.

\sphinxAtStartPar

\sphinxAtStartPar
This function returns a pointer to data provided by a {\hyperref[\detokenize{fix::doc}]{\sphinxcrossref{\DUrole{doc}{fix command}}}}
instance identified by its fix\sphinxhyphen{}ID.  Fixes may provide global, per\sphinxhyphen{}atom,
or local data, and those may be a scalar, a vector, or an array, or they
may provide the information about the dimensions of the respective data.
Since individual fixes may provide multiple kinds of data, it is
required to set style and type flags representing what specific data is
desired.  This also determines to what kind of pointer the returned
pointer needs to be cast to access the data correctly.  The function
returns \sphinxcode{\sphinxupquote{NULL}} if the fix ID is not found or the requested data is not
available.

\sphinxAtStartPar
The following table lists the available options.


\begin{savenotes}\sphinxattablestart
\sphinxthistablewithglobalstyle
\centering
\begin{tabular}[t]{\X{25}{100}\X{24}{100}\X{14}{100}\X{37}{100}}
\sphinxtoprule
\sphinxstyletheadfamily 
\sphinxAtStartPar
Style (see {\hyperref[\detokenize{Library_objects:_CPPv416_LMP_STYLE_CONST}]{\sphinxcrossref{\sphinxcode{\sphinxupquote{\_LMP\_STYLE\_CONST}}}}})
&\sphinxstyletheadfamily 
\sphinxAtStartPar
Type (see {\hyperref[\detokenize{Library_objects:_CPPv415_LMP_TYPE_CONST}]{\sphinxcrossref{\sphinxcode{\sphinxupquote{\_LMP\_TYPE\_CONST}}}}})
&\sphinxstyletheadfamily 
\sphinxAtStartPar
Returned type
&\sphinxstyletheadfamily 
\sphinxAtStartPar
Returned data
\\
\sphinxmidrule
\sphinxtableatstartofbodyhook
\sphinxAtStartPar
LMP\_STYLE\_GLOBAL
&
\sphinxAtStartPar
LMP\_TYPE\_SCALAR
&
\sphinxAtStartPar
\sphinxcode{\sphinxupquote{double *}}
&
\sphinxAtStartPar
Copy of global scalar
\\
\sphinxhline
\sphinxAtStartPar
LMP\_STYLE\_GLOBAL
&
\sphinxAtStartPar
LMP\_TYPE\_VECTOR
&
\sphinxAtStartPar
\sphinxcode{\sphinxupquote{double *}}
&
\sphinxAtStartPar
Copy of global vector element at index nrow
\\
\sphinxhline
\sphinxAtStartPar
LMP\_STYLE\_GLOBAL
&
\sphinxAtStartPar
LMP\_TYPE\_ARRAY
&
\sphinxAtStartPar
\sphinxcode{\sphinxupquote{double *}}
&
\sphinxAtStartPar
Copy of global array element at nrow, ncol
\\
\sphinxhline
\sphinxAtStartPar
LMP\_STYLE\_GLOBAL
&
\sphinxAtStartPar
LMP\_SIZE\_VECTOR
&
\sphinxAtStartPar
\sphinxcode{\sphinxupquote{int *}}
&
\sphinxAtStartPar
Length of global vector
\\
\sphinxhline
\sphinxAtStartPar
LMP\_STYLE\_GLOBAL
&
\sphinxAtStartPar
LMP\_SIZE\_ROWS
&
\sphinxAtStartPar
\sphinxcode{\sphinxupquote{int *}}
&
\sphinxAtStartPar
Rows in global array
\\
\sphinxhline
\sphinxAtStartPar
LMP\_STYLE\_GLOBAL
&
\sphinxAtStartPar
LMP\_SIZE\_COLS
&
\sphinxAtStartPar
\sphinxcode{\sphinxupquote{int *}}
&
\sphinxAtStartPar
Columns in global array
\\
\sphinxhline
\sphinxAtStartPar
LMP\_STYLE\_ATOM
&
\sphinxAtStartPar
LMP\_TYPE\_VECTOR
&
\sphinxAtStartPar
\sphinxcode{\sphinxupquote{double *}}
&
\sphinxAtStartPar
Per\sphinxhyphen{}atom value
\\
\sphinxhline
\sphinxAtStartPar
LMP\_STYLE\_ATOM
&
\sphinxAtStartPar
LMP\_TYPE\_ARRAY
&
\sphinxAtStartPar
\sphinxcode{\sphinxupquote{double **}}
&
\sphinxAtStartPar
Per\sphinxhyphen{}atom vector
\\
\sphinxhline
\sphinxAtStartPar
LMP\_STYLE\_ATOM
&
\sphinxAtStartPar
LMP\_SIZE\_COLS
&
\sphinxAtStartPar
\sphinxcode{\sphinxupquote{int *}}
&
\sphinxAtStartPar
Columns of per\sphinxhyphen{}atom array, 0 if vector
\\
\sphinxhline
\sphinxAtStartPar
LMP\_STYLE\_LOCAL
&
\sphinxAtStartPar
LMP\_TYPE\_VECTOR
&
\sphinxAtStartPar
\sphinxcode{\sphinxupquote{double *}}
&
\sphinxAtStartPar
Local data vector
\\
\sphinxhline
\sphinxAtStartPar
LMP\_STYLE\_LOCAL
&
\sphinxAtStartPar
LMP\_TYPE\_ARRAY
&
\sphinxAtStartPar
\sphinxcode{\sphinxupquote{double **}}
&
\sphinxAtStartPar
Local data array
\\
\sphinxhline
\sphinxAtStartPar
LMP\_STYLE\_LOCAL
&
\sphinxAtStartPar
LMP\_SIZE\_ROWS
&
\sphinxAtStartPar
\sphinxcode{\sphinxupquote{int *}}
&
\sphinxAtStartPar
Number of local data rows
\\
\sphinxhline
\sphinxAtStartPar
LMP\_STYLE\_LOCAL
&
\sphinxAtStartPar
LMP\_SIZE\_COLS
&
\sphinxAtStartPar
\sphinxcode{\sphinxupquote{int *}}
&
\sphinxAtStartPar
Number of local data columns
\\
\sphinxbottomrule
\end{tabular}
\sphinxtableafterendhook\par
\sphinxattableend\end{savenotes}


\begin{sphinxadmonition}{note}{Note:}
\sphinxAtStartPar
When requesting global data, the fix data can only be accessed one
item at a time without access to the pointer itself.  Thus this
function will allocate storage for a single double value, copy the
returned value to it, and returns a pointer to the location of the
copy.  Therefore the allocated storage needs to be freed after its
use to avoid a memory leak. Example:

\begin{sphinxVerbatim}[commandchars=\\\{\}]
\PYG{k+kt}{double}\PYG{+w}{ }\PYG{o}{*}\PYG{n}{dptr}\PYG{+w}{ }\PYG{o}{=}\PYG{+w}{ }\PYG{p}{(}\PYG{k+kt}{double}\PYG{+w}{ }\PYG{o}{*}\PYG{p}{)}\PYG{+w}{ }\PYG{n}{lammps\PYGZus{}extract\PYGZus{}fix}\PYG{p}{(}\PYG{n}{handle}\PYG{p}{,}\PYG{+w}{ }\PYG{n}{name}\PYG{p}{,}
\PYG{+w}{                          }\PYG{n}{LMP\PYGZus{}STYLE\PYGZus{}GLOBAL}\PYG{p}{,}\PYG{+w}{ }\PYG{n}{LMP\PYGZus{}TYPE\PYGZus{}VECTOR}\PYG{p}{,}\PYG{+w}{ }\PYG{l+m+mi}{0}\PYG{p}{,}\PYG{+w}{ }\PYG{l+m+mi}{0}\PYG{p}{)}\PYG{p}{;}
\PYG{k+kt}{double}\PYG{+w}{ }\PYG{n}{value}\PYG{+w}{ }\PYG{o}{=}\PYG{+w}{ }\PYG{o}{*}\PYG{n}{dptr}\PYG{p}{;}
\PYG{n}{lammps\PYGZus{}free}\PYG{p}{(}\PYG{p}{(}\PYG{k+kt}{void}\PYG{+w}{ }\PYG{o}{*}\PYG{p}{)}\PYG{n}{dptr}\PYG{p}{)}\PYG{p}{;}
\end{sphinxVerbatim}
\end{sphinxadmonition}

\begin{sphinxadmonition}{note}{Note:}
\sphinxAtStartPar
LAMMPS cannot easily check if it is valid to access the data, so it
may fail with an error.  The caller has to avoid such an error.
\end{sphinxadmonition}

\begin{sphinxadmonition}{warning}{Warning:}
\sphinxAtStartPar
The pointers returned by this function for per\sphinxhyphen{}atom or local data are
generally not persistent, since the computed data may be re\sphinxhyphen{}distributed,
re\sphinxhyphen{}allocated, and re\sphinxhyphen{}ordered at every invocation of the fix.  It is thus
advisable to re\sphinxhyphen{}invoke this function before the data is accessed, or
make a copy, if the data shall be used after other LAMMPS commands have
been issued.
\end{sphinxadmonition}
\begin{quote}\begin{description}
\sphinxlineitem{Parameters}\begin{itemize}
\item {} 
\sphinxAtStartPar
\sphinxstyleliteralstrong{\sphinxupquote{handle}} \textendash{} pointer to a previously created LAMMPS instance 

\item {} 
\sphinxAtStartPar
\sphinxstyleliteralstrong{\sphinxupquote{id}} \textendash{} string with ID of the fix 

\item {} 
\sphinxAtStartPar
\sphinxstyleliteralstrong{\sphinxupquote{style}} \textendash{} constant indicating the style of data requested (global, per\sphinxhyphen{}atom, or local) 

\item {} 
\sphinxAtStartPar
\sphinxstyleliteralstrong{\sphinxupquote{type}} \textendash{} constant indicating type of data (scalar, vector, or array) or size of rows or columns 

\item {} 
\sphinxAtStartPar
\sphinxstyleliteralstrong{\sphinxupquote{nrow}} \textendash{} row index (only used for global vectors and arrays) 

\item {} 
\sphinxAtStartPar
\sphinxstyleliteralstrong{\sphinxupquote{ncol}} \textendash{} column index (only used for global arrays) 

\end{itemize}

\sphinxlineitem{Returns}
\sphinxAtStartPar
pointer (cast to \sphinxcode{\sphinxupquote{void *}}) to the location of the requested data or \sphinxcode{\sphinxupquote{NULL}} if not found. 

\end{description}\end{quote}

\end{fulllineitems}



\bigskip\hrule\bigskip

\index{lammps\_extract\_variable\_datatype (C++ function)@\spxentry{lammps\_extract\_variable\_datatype}\spxextra{C++ function}}

\begin{fulllineitems}
\phantomsection\label{\detokenize{Library_objects:_CPPv432lammps_extract_variable_datatypePvPKc}}\phantomsection\label{\detokenize{Library_objects:_CPPv332lammps_extract_variable_datatypePvPKc}}\phantomsection\label{\detokenize{Library_objects:_CPPv232lammps_extract_variable_datatypePvPKc}}\phantomsection\label{\detokenize{Library_objects:lammps_extract_variable_datatype__voidP.cCP}}
\pysigstartsignatures
\pysigstartmultiline
\pysiglinewithargsret{\phantomsection\label{\detokenize{Library_objects:library_8h_1a5725944733b98157d3f6993c2f6499ce}}\DUrole{kt}{int}\DUrole{w}{ }\sphinxbfcode{\sphinxupquote{\DUrole{n}{lammps\_extract\_variable\_datatype}}}}{\DUrole{kt}{void}\DUrole{w}{ }\DUrole{p}{*}\DUrole{n,sig-param}{handle}\sphinxparamcomma \DUrole{k}{const}\DUrole{w}{ }\DUrole{kt}{char}\DUrole{w}{ }\DUrole{p}{*}\DUrole{n,sig-param}{name}}{}
\pysigstopmultiline
\pysigstopsignatures
\sphinxAtStartPar
Get data type of a LAMMPS variable.

\sphinxAtStartPar

\sphinxAtStartPar
\DUrole{versionmodified,added}{Added in version 3Nov2022.}

\sphinxAtStartPar
This function returns an integer that encodes the data type of the variable
with the specified name. See {\hyperref[\detokenize{Library_objects:_CPPv414_LMP_VAR_CONST}]{\sphinxcrossref{\sphinxcode{\sphinxupquote{\_LMP\_VAR\_CONST}}}}} for valid values.
Callers of {\hyperref[\detokenize{Library_objects:_CPPv423lammps_extract_variablePvPKcPKc}]{\sphinxcrossref{\sphinxcode{\sphinxupquote{lammps\_extract\_variable()}}}}} can use this information to
decide how to cast the \sphinxcode{\sphinxupquote{void *}} pointer and access the data.

\begin{quote}\begin{description}
\sphinxlineitem{Parameters}\begin{itemize}
\item {} 
\sphinxAtStartPar
\sphinxstyleliteralstrong{\sphinxupquote{handle}} \textendash{} pointer to a previously created LAMMPS instance 

\item {} 
\sphinxAtStartPar
\sphinxstyleliteralstrong{\sphinxupquote{name}} \textendash{} string with the name of the extracted variable 

\end{itemize}

\sphinxlineitem{Returns}
\sphinxAtStartPar
integer constant encoding the data type of the property or \sphinxhyphen{}1 if not found. 

\end{description}\end{quote}

\end{fulllineitems}



\bigskip\hrule\bigskip

\index{lammps\_extract\_variable (C++ function)@\spxentry{lammps\_extract\_variable}\spxextra{C++ function}}

\begin{fulllineitems}
\phantomsection\label{\detokenize{Library_objects:_CPPv423lammps_extract_variablePvPKcPKc}}\phantomsection\label{\detokenize{Library_objects:_CPPv323lammps_extract_variablePvPKcPKc}}\phantomsection\label{\detokenize{Library_objects:_CPPv223lammps_extract_variablePvPKcPKc}}\phantomsection\label{\detokenize{Library_objects:lammps_extract_variable__voidP.cCP.cCP}}
\pysigstartsignatures
\pysigstartmultiline
\pysiglinewithargsret{\phantomsection\label{\detokenize{Library_objects:library_8h_1a93437a9fed4a4dec044d010e9df52d7f}}\DUrole{kt}{void}\DUrole{w}{ }\DUrole{p}{*}\sphinxbfcode{\sphinxupquote{\DUrole{n}{lammps\_extract\_variable}}}}{\DUrole{kt}{void}\DUrole{w}{ }\DUrole{p}{*}\DUrole{n,sig-param}{handle}\sphinxparamcomma \DUrole{k}{const}\DUrole{w}{ }\DUrole{kt}{char}\DUrole{w}{ }\DUrole{p}{*}\DUrole{n,sig-param}{name}\sphinxparamcomma \DUrole{k}{const}\DUrole{w}{ }\DUrole{kt}{char}\DUrole{w}{ }\DUrole{p}{*}\DUrole{n,sig-param}{group}}{}
\pysigstopmultiline
\pysigstopsignatures
\sphinxAtStartPar
Get pointer to data from a LAMMPS variable.

\sphinxAtStartPar

\sphinxAtStartPar
This function returns a pointer to data from a LAMMPS {\hyperref[\detokenize{variable::doc}]{\sphinxcrossref{\DUrole{doc}{variable command}}}}
identified by its name.  When the variable is either an \sphinxstyleemphasis{equal}\sphinxhyphen{}style
compatible variable, a \sphinxstyleemphasis{vector}\sphinxhyphen{}style variable, or an \sphinxstyleemphasis{atom}\sphinxhyphen{}style
variable, the variable is evaluated and the corresponding value(s) returned.
Variables of style \sphinxstyleemphasis{internal} are compatible with \sphinxstyleemphasis{equal}\sphinxhyphen{}style variables and
so are \sphinxstyleemphasis{python}\sphinxhyphen{}style variables, if they return a numeric value.  For other
variable styles, their string value is returned.  The function returns
\sphinxcode{\sphinxupquote{NULL}} when a variable of the provided \sphinxstyleemphasis{name} is not found or of an
incompatible style.  The \sphinxstyleemphasis{group} argument is only used for \sphinxstyleemphasis{atom}\sphinxhyphen{}style variables and ignored otherwise, with one exception: for style \sphinxstyleemphasis{vector},
if \sphinxstyleemphasis{group} is “GET\_VECTOR\_SIZE”, the returned pointer will yield the length
of the vector to be returned when dereferenced. This pointer must be
deallocated after the value is read to avoid a memory leak.
If \sphinxstyleemphasis{group} is set to \sphinxcode{\sphinxupquote{NULL}} when extracting data from an \sphinxstyleemphasis{atom}\sphinxhyphen{}style
variable, the group is assumed to be “all”.

\sphinxAtStartPar
When requesting data from an \sphinxstyleemphasis{equal}\sphinxhyphen{}style or compatible variable
this function allocates storage for a single double value, copies the
returned value to it, and returns a pointer to the location of the
copy.  Therefore the allocated storage needs to be freed after its
use to avoid a memory leak. Example:

\begin{sphinxVerbatim}[commandchars=\\\{\}]
\PYG{k+kt}{double}\PYG{+w}{ }\PYG{o}{*}\PYG{n}{dptr}\PYG{+w}{ }\PYG{o}{=}\PYG{+w}{ }\PYG{p}{(}\PYG{k+kt}{double}\PYG{+w}{ }\PYG{o}{*}\PYG{p}{)}\PYG{+w}{ }\PYG{n}{lammps\PYGZus{}extract\PYGZus{}variable}\PYG{p}{(}\PYG{n}{handle}\PYG{p}{,}\PYG{+w}{ }\PYG{n}{name}\PYG{p}{,}\PYG{+w}{ }\PYG{n+nb}{NULL}\PYG{p}{)}\PYG{p}{;}
\PYG{k+kt}{double}\PYG{+w}{ }\PYG{n}{value}\PYG{+w}{ }\PYG{o}{=}\PYG{+w}{ }\PYG{o}{*}\PYG{n}{dptr}\PYG{p}{;}
\PYG{n}{lammps\PYGZus{}free}\PYG{p}{(}\PYG{p}{(}\PYG{k+kt}{void}\PYG{+w}{ }\PYG{o}{*}\PYG{p}{)}\PYG{n}{dptr}\PYG{p}{)}\PYG{p}{;}
\end{sphinxVerbatim}

\sphinxAtStartPar
For \sphinxstyleemphasis{atom}\sphinxhyphen{}style variables, the return value is a pointer to an
allocated block of storage of double of the length \sphinxcode{\sphinxupquote{atom\sphinxhyphen{}\textgreater{}nlocal}}.
Since the data returned are a copy, the location will persist, but its
content will not be updated in case the variable is re\sphinxhyphen{}evaluated.
To avoid a memory leak, this pointer needs to be freed after use in
the calling program.

\sphinxAtStartPar
For \sphinxstyleemphasis{vector}\sphinxhyphen{}style variables, the returned pointer points to actual
LAMMPS data and thus it should \sphinxstylestrong{not} be deallocated.  Its length
depends on the variable, compute, or fix data used to construct the
\sphinxstyleemphasis{vector}\sphinxhyphen{}style variable.  This length can be fetched by calling this
function with \sphinxstyleemphasis{group} set to a non\sphinxhyphen{}NULL pointer (NULL returns the vector).
In that case it will return the vector length as an allocated int
pointer cast to a \sphinxcode{\sphinxupquote{void *}} pointer.  That pointer can be recast and
dereferenced to an integer yielding the length of the vector. This pointer
must be deallocated when finished with it to avoid memory leaks. Example:

\begin{sphinxVerbatim}[commandchars=\\\{\}]
\PYG{k+kt}{double}\PYG{+w}{ }\PYG{o}{*}\PYG{n}{vectvals}\PYG{+w}{ }\PYG{o}{=}\PYG{+w}{ }\PYG{p}{(}\PYG{k+kt}{double}\PYG{+w}{ }\PYG{o}{*}\PYG{p}{)}\PYG{+w}{ }\PYG{n}{lammps\PYGZus{}extract\PYGZus{}variable}\PYG{p}{(}\PYG{n}{handle}\PYG{p}{,}\PYG{+w}{ }\PYG{n}{name}\PYG{p}{,}\PYG{+w}{ }\PYG{n+nb}{NULL}\PYG{p}{)}\PYG{p}{;}
\PYG{k+kt}{int}\PYG{+w}{ }\PYG{o}{*}\PYG{n}{intptr}\PYG{+w}{ }\PYG{o}{=}\PYG{+w}{ }\PYG{p}{(}\PYG{k+kt}{int}\PYG{+w}{ }\PYG{o}{*}\PYG{p}{)}\PYG{+w}{ }\PYG{n}{lammps\PYGZus{}extract\PYGZus{}variable}\PYG{p}{(}\PYG{n}{handle}\PYG{p}{,}\PYG{+w}{ }\PYG{n}{name}\PYG{p}{,}\PYG{+w}{ }\PYG{l+m+mi}{1}\PYG{p}{)}\PYG{p}{;}
\PYG{k+kt}{int}\PYG{+w}{ }\PYG{n}{vectlen}\PYG{+w}{ }\PYG{o}{=}\PYG{+w}{ }\PYG{o}{*}\PYG{n}{intptr}\PYG{p}{;}
\PYG{n}{lammps\PYGZus{}free}\PYG{p}{(}\PYG{p}{(}\PYG{k+kt}{void}\PYG{+w}{ }\PYG{o}{*}\PYG{p}{)}\PYG{n}{intptr}\PYG{p}{)}\PYG{p}{;}
\end{sphinxVerbatim}

\sphinxAtStartPar
For other variable styles the returned pointer needs to be cast to
a char pointer and it should \sphinxstylestrong{not} be deallocated. Example:

\begin{sphinxVerbatim}[commandchars=\\\{\}]
\PYG{k}{const}\PYG{+w}{ }\PYG{k+kt}{char}\PYG{+w}{ }\PYG{o}{*}\PYG{n}{cptr}\PYG{+w}{ }\PYG{o}{=}\PYG{+w}{ }\PYG{p}{(}\PYG{k}{const}\PYG{+w}{ }\PYG{k+kt}{char}\PYG{+w}{ }\PYG{o}{*}\PYG{p}{)}\PYG{+w}{ }\PYG{n}{lammps\PYGZus{}extract\PYGZus{}variable}\PYG{p}{(}\PYG{n}{handle}\PYG{p}{,}\PYG{n}{name}\PYG{p}{,}\PYG{n+nb}{NULL}\PYG{p}{)}\PYG{p}{;}
\PYG{n}{printf}\PYG{p}{(}\PYG{l+s}{\PYGZdq{}}\PYG{l+s}{The value of variable \PYGZpc{}s is \PYGZpc{}s}\PYG{l+s+se}{\PYGZbs{}n}\PYG{l+s}{\PYGZdq{}}\PYG{p}{,}\PYG{+w}{ }\PYG{n}{name}\PYG{p}{,}\PYG{+w}{ }\PYG{n}{cptr}\PYG{p}{)}\PYG{p}{;}
\end{sphinxVerbatim}


\begin{sphinxadmonition}{note}{Note:}
\sphinxAtStartPar
LAMMPS cannot easily check if it is valid to access the data
referenced by the variables (e.g., computes, fixes, or thermodynamic
info), so it may fail with an error.  The caller has to make certain
that the data is extracted only when it safe to evaluate the variable
and thus an error or crash are avoided.
\end{sphinxadmonition}
\begin{quote}\begin{description}
\sphinxlineitem{Parameters}\begin{itemize}
\item {} 
\sphinxAtStartPar
\sphinxstyleliteralstrong{\sphinxupquote{handle}} \textendash{} pointer to a previously created LAMMPS instance 

\item {} 
\sphinxAtStartPar
\sphinxstyleliteralstrong{\sphinxupquote{name}} \textendash{} name of the variable 

\item {} 
\sphinxAtStartPar
\sphinxstyleliteralstrong{\sphinxupquote{group}} \textendash{} group\sphinxhyphen{}ID for atom style variable or \sphinxcode{\sphinxupquote{NULL}} or non\sphinxhyphen{}NULL to get vector length 

\end{itemize}

\sphinxlineitem{Returns}
\sphinxAtStartPar
pointer (cast to \sphinxcode{\sphinxupquote{void *}}) to the location of the requested data or \sphinxcode{\sphinxupquote{NULL}} if not found. 

\end{description}\end{quote}

\end{fulllineitems}



\bigskip\hrule\bigskip

\index{lammps\_set\_variable (C++ function)@\spxentry{lammps\_set\_variable}\spxextra{C++ function}}

\begin{fulllineitems}
\phantomsection\label{\detokenize{Library_objects:_CPPv419lammps_set_variablePvPKcPKc}}\phantomsection\label{\detokenize{Library_objects:_CPPv319lammps_set_variablePvPKcPKc}}\phantomsection\label{\detokenize{Library_objects:_CPPv219lammps_set_variablePvPKcPKc}}\phantomsection\label{\detokenize{Library_objects:lammps_set_variable__voidP.cCP.cCP}}
\pysigstartsignatures
\pysigstartmultiline
\pysiglinewithargsret{\phantomsection\label{\detokenize{Library_objects:library_8h_1ab3cc202eaec7472181870b4747b3d577}}\DUrole{kt}{int}\DUrole{w}{ }\sphinxbfcode{\sphinxupquote{\DUrole{n}{lammps\_set\_variable}}}}{\DUrole{kt}{void}\DUrole{w}{ }\DUrole{p}{*}\DUrole{n,sig-param}{handle}\sphinxparamcomma \DUrole{k}{const}\DUrole{w}{ }\DUrole{kt}{char}\DUrole{w}{ }\DUrole{p}{*}\DUrole{n,sig-param}{name}\sphinxparamcomma \DUrole{k}{const}\DUrole{w}{ }\DUrole{kt}{char}\DUrole{w}{ }\DUrole{p}{*}\DUrole{n,sig-param}{str}}{}
\pysigstopmultiline
\pysigstopsignatures
\sphinxAtStartPar
Set the value of a string\sphinxhyphen{}style variable. 
\sphinxAtStartPar

\sphinxAtStartPar
\DUrole{versionmodified,deprecated}{Deprecated since version 7Feb2024.}

\sphinxAtStartPar
This function assigns a new value from the string str to the
string\sphinxhyphen{}style variable \sphinxstyleemphasis{name}.  This is a way to directly change the
string value of a LAMMPS variable that was previous defined with a
{\hyperref[\detokenize{variable::doc}]{\sphinxcrossref{\DUrole{doc}{variable name string}}}} command without using any
LAMMPS commands to delete and redefine the variable.

\sphinxAtStartPar
Returns \sphinxhyphen{}1 if a variable of that name does not exist or if it is not
a string\sphinxhyphen{}style variable, otherwise 0.



\begin{sphinxadmonition}{warning}{Warning:}
\sphinxAtStartPar
This function is deprecated and {\hyperref[\detokenize{Library_objects:_CPPv426lammps_set_string_variablePvPKcPKc}]{\sphinxcrossref{\sphinxcode{\sphinxupquote{lammps\_set\_string\_variable()}}}}}
should be used instead.
\end{sphinxadmonition}
\begin{quote}\begin{description}
\sphinxlineitem{Parameters}\begin{itemize}
\item {} 
\sphinxAtStartPar
\sphinxstyleliteralstrong{\sphinxupquote{handle}} \textendash{} pointer to a previously created LAMMPS instance 

\item {} 
\sphinxAtStartPar
\sphinxstyleliteralstrong{\sphinxupquote{name}} \textendash{} name of the variable 

\item {} 
\sphinxAtStartPar
\sphinxstyleliteralstrong{\sphinxupquote{str}} \textendash{} new value of the variable 

\end{itemize}

\sphinxlineitem{Returns}
\sphinxAtStartPar
0 on success or \sphinxhyphen{}1 on failure 

\end{description}\end{quote}

\end{fulllineitems}



\bigskip\hrule\bigskip

\index{lammps\_set\_string\_variable (C++ function)@\spxentry{lammps\_set\_string\_variable}\spxextra{C++ function}}

\begin{fulllineitems}
\phantomsection\label{\detokenize{Library_objects:_CPPv426lammps_set_string_variablePvPKcPKc}}\phantomsection\label{\detokenize{Library_objects:_CPPv326lammps_set_string_variablePvPKcPKc}}\phantomsection\label{\detokenize{Library_objects:_CPPv226lammps_set_string_variablePvPKcPKc}}\phantomsection\label{\detokenize{Library_objects:lammps_set_string_variable__voidP.cCP.cCP}}
\pysigstartsignatures
\pysigstartmultiline
\pysiglinewithargsret{\phantomsection\label{\detokenize{Library_objects:library_8h_1a540ffa539718883760eb6fafab3627bc}}\DUrole{kt}{int}\DUrole{w}{ }\sphinxbfcode{\sphinxupquote{\DUrole{n}{lammps\_set\_string\_variable}}}}{\DUrole{kt}{void}\DUrole{w}{ }\DUrole{p}{*}\DUrole{n,sig-param}{handle}\sphinxparamcomma \DUrole{k}{const}\DUrole{w}{ }\DUrole{kt}{char}\DUrole{w}{ }\DUrole{p}{*}\DUrole{n,sig-param}{name}\sphinxparamcomma \DUrole{k}{const}\DUrole{w}{ }\DUrole{kt}{char}\DUrole{w}{ }\DUrole{p}{*}\DUrole{n,sig-param}{str}}{}
\pysigstopmultiline
\pysigstopsignatures
\sphinxAtStartPar
Set the value of a string\sphinxhyphen{}style variable. 
\sphinxAtStartPar

\sphinxAtStartPar
\DUrole{versionmodified,added}{Added in version 7Feb2024.}

\sphinxAtStartPar
This function assigns a new value from the string str to the
string\sphinxhyphen{}style variable \sphinxstyleemphasis{name}.  This is a way to directly change the
string value of a LAMMPS variable that was previous defined with a
{\hyperref[\detokenize{variable::doc}]{\sphinxcrossref{\DUrole{doc}{variable name string}}}} command without using any
LAMMPS commands to delete and redefine the variable.

\sphinxAtStartPar
Returns \sphinxhyphen{}1 if a variable of that name does not exist or if it is not
a string\sphinxhyphen{}style variable, otherwise 0.


\begin{quote}\begin{description}
\sphinxlineitem{Parameters}\begin{itemize}
\item {} 
\sphinxAtStartPar
\sphinxstyleliteralstrong{\sphinxupquote{handle}} \textendash{} pointer to a previously created LAMMPS instance 

\item {} 
\sphinxAtStartPar
\sphinxstyleliteralstrong{\sphinxupquote{name}} \textendash{} name of the variable 

\item {} 
\sphinxAtStartPar
\sphinxstyleliteralstrong{\sphinxupquote{str}} \textendash{} new value of the variable 

\end{itemize}

\sphinxlineitem{Returns}
\sphinxAtStartPar
0 on success or \sphinxhyphen{}1 on failure 

\end{description}\end{quote}

\end{fulllineitems}



\bigskip\hrule\bigskip

\index{lammps\_set\_internal\_variable (C++ function)@\spxentry{lammps\_set\_internal\_variable}\spxextra{C++ function}}

\begin{fulllineitems}
\phantomsection\label{\detokenize{Library_objects:_CPPv428lammps_set_internal_variablePvPKcd}}\phantomsection\label{\detokenize{Library_objects:_CPPv328lammps_set_internal_variablePvPKcd}}\phantomsection\label{\detokenize{Library_objects:_CPPv228lammps_set_internal_variablePvPKcd}}\phantomsection\label{\detokenize{Library_objects:lammps_set_internal_variable__voidP.cCP.double}}
\pysigstartsignatures
\pysigstartmultiline
\pysiglinewithargsret{\phantomsection\label{\detokenize{Library_objects:library_8h_1a250fd1b00f6b69daf98cc4a879c78059}}\DUrole{kt}{int}\DUrole{w}{ }\sphinxbfcode{\sphinxupquote{\DUrole{n}{lammps\_set\_internal\_variable}}}}{\DUrole{kt}{void}\DUrole{w}{ }\DUrole{p}{*}\DUrole{n,sig-param}{handle}\sphinxparamcomma \DUrole{k}{const}\DUrole{w}{ }\DUrole{kt}{char}\DUrole{w}{ }\DUrole{p}{*}\DUrole{n,sig-param}{name}\sphinxparamcomma \DUrole{kt}{double}\DUrole{w}{ }\DUrole{n,sig-param}{value}}{}
\pysigstopmultiline
\pysigstopsignatures
\sphinxAtStartPar
Set the value of an internal\sphinxhyphen{}style variable.

\sphinxAtStartPar

\sphinxAtStartPar
\DUrole{versionmodified,added}{Added in version 7Feb2024.}

\sphinxAtStartPar
This function assigns a new value from the floating point number \sphinxstyleemphasis{value}
to the internal\sphinxhyphen{}style variable \sphinxstyleemphasis{name}.  This is a way to directly change
the numerical value of such a LAMMPS variable that was previous defined
with a {\hyperref[\detokenize{variable::doc}]{\sphinxcrossref{\DUrole{doc}{variable name internal}}}} command without using
any LAMMPS commands to delete and redefine the variable.

\sphinxAtStartPar
Returns \sphinxhyphen{}1 if a variable of that name does not exist or is not an
internal\sphinxhyphen{}style variable, otherwise 0.

\begin{quote}\begin{description}
\sphinxlineitem{Parameters}\begin{itemize}
\item {} 
\sphinxAtStartPar
\sphinxstyleliteralstrong{\sphinxupquote{handle}} \textendash{} pointer to a previously created LAMMPS instance 

\item {} 
\sphinxAtStartPar
\sphinxstyleliteralstrong{\sphinxupquote{name}} \textendash{} name of the variable 

\item {} 
\sphinxAtStartPar
\sphinxstyleliteralstrong{\sphinxupquote{value}} \textendash{} new value of the variable 

\end{itemize}

\sphinxlineitem{Returns}
\sphinxAtStartPar
0 on success or \sphinxhyphen{}1 on failure 

\end{description}\end{quote}

\end{fulllineitems}



\bigskip\hrule\bigskip

\index{lammps\_variable\_info (C++ function)@\spxentry{lammps\_variable\_info}\spxextra{C++ function}}

\begin{fulllineitems}
\phantomsection\label{\detokenize{Library_objects:_CPPv420lammps_variable_infoPviPci}}\phantomsection\label{\detokenize{Library_objects:_CPPv320lammps_variable_infoPviPci}}\phantomsection\label{\detokenize{Library_objects:_CPPv220lammps_variable_infoPviPci}}\phantomsection\label{\detokenize{Library_objects:lammps_variable_info__voidP.i.cP.i}}
\pysigstartsignatures
\pysigstartmultiline
\pysiglinewithargsret{\phantomsection\label{\detokenize{Library_objects:library_8h_1a7d42a4fb49036f76e6042ca8302aa567}}\DUrole{kt}{int}\DUrole{w}{ }\sphinxbfcode{\sphinxupquote{\DUrole{n}{lammps\_variable\_info}}}}{\DUrole{kt}{void}\DUrole{w}{ }\DUrole{p}{*}\DUrole{n,sig-param}{handle}\sphinxparamcomma \DUrole{kt}{int}\DUrole{w}{ }\DUrole{n,sig-param}{idx}\sphinxparamcomma \DUrole{kt}{char}\DUrole{w}{ }\DUrole{p}{*}\DUrole{n,sig-param}{buf}\sphinxparamcomma \DUrole{kt}{int}\DUrole{w}{ }\DUrole{n,sig-param}{bufsize}}{}
\pysigstopmultiline
\pysigstopsignatures
\sphinxAtStartPar
Retrieve informational string for a variable.

\sphinxAtStartPar

\sphinxAtStartPar
\DUrole{versionmodified,added}{Added in version 21Nov2023.}

\sphinxAtStartPar
This function copies a string with human readable information about
a defined variable: name, style, current value(s) into the provided
C\sphinxhyphen{}style string buffer.  That is the same info as produced by the
{\hyperref[\detokenize{info::doc}]{\sphinxcrossref{\DUrole{doc}{info variables}}}} command. The length of the buffer must
be provided as \sphinxstyleemphasis{buf\_size} argument.  If the info exceeds the length
of the buffer, it will be truncated accordingly.  If the index is
out of range, the function returns 0 and \sphinxstyleemphasis{buffer} is set to an empty
string, otherwise 1.

\begin{quote}\begin{description}
\sphinxlineitem{Parameters}\begin{itemize}
\item {} 
\sphinxAtStartPar
\sphinxstyleliteralstrong{\sphinxupquote{handle}} \textendash{} pointer to a previously created LAMMPS instance cast to \sphinxcode{\sphinxupquote{void *}}. 

\item {} 
\sphinxAtStartPar
\sphinxstyleliteralstrong{\sphinxupquote{idx}} \textendash{} index of the variable (0 \textless{}= idx \textless{} nvar) 

\item {} 
\sphinxAtStartPar
\sphinxstyleliteralstrong{\sphinxupquote{buffer}} \textendash{} string buffer to copy the info to 

\item {} 
\sphinxAtStartPar
\sphinxstyleliteralstrong{\sphinxupquote{buf\_size}} \textendash{} size of the provided string buffer 

\end{itemize}

\sphinxlineitem{Returns}
\sphinxAtStartPar
1 if successful, otherwise 0 

\end{description}\end{quote}

\end{fulllineitems}



\bigskip\hrule\bigskip

\index{multitype::\_LMP\_DATATYPE\_CONST (C++ enum)@\spxentry{multitype::\_LMP\_DATATYPE\_CONST}\spxextra{C++ enum}}

\begin{fulllineitems}
\phantomsection\label{\detokenize{Library_objects:_CPPv4N9LAMMPS_NS9multitype19_LMP_DATATYPE_CONSTE}}\phantomsection\label{\detokenize{Library_objects:_CPPv3N9LAMMPS_NS9multitype19_LMP_DATATYPE_CONSTE}}\phantomsection\label{\detokenize{Library_objects:_CPPv2N9LAMMPS_NS9multitype19_LMP_DATATYPE_CONSTE}}
\pysigstartsignatures
\pysigstartmultiline
\pysigline{\phantomsection\label{\detokenize{Library_objects:structLAMMPS__NS_1_1multitype_1a8f60df5cd520b67959d22185afcdacc5}}\DUrole{k}{enum}\DUrole{w}{ }\sphinxcode{\sphinxupquote{\DUrole{n}{LAMMPS\_NS}\DUrole{p}{::}\DUrole{n}{multitype}\DUrole{p}{::}}}\sphinxbfcode{\sphinxupquote{\DUrole{n}{\_LMP\_DATATYPE\_CONST}}}}
\pysigstopmultiline
\pysigstopsignatures
\sphinxAtStartPar
Data type constants for extracting data from atoms, computes and fixes

\sphinxAtStartPar
This enum must be kept in sync with the corresponding enum or constants in \sphinxcode{\sphinxupquote{python/lammps/constants.py}}, \sphinxcode{\sphinxupquote{fortran/lammps.f90}}, \sphinxcode{\sphinxupquote{tools/swig/lammps.i}}, \sphinxcode{\sphinxupquote{src/library.h}}, and \sphinxcode{\sphinxupquote{examples/COUPLE/plugin/liblammpsplugin.h}}

\sphinxAtStartPar
\sphinxstyleemphasis{Values:}
\index{multitype::\_LMP\_DATATYPE\_CONST::LAMMPS\_NONE (C++ enumerator)@\spxentry{multitype::\_LMP\_DATATYPE\_CONST::LAMMPS\_NONE}\spxextra{C++ enumerator}}

\begin{fulllineitems}
\phantomsection\label{\detokenize{Library_objects:_CPPv4N9LAMMPS_NS9multitype19_LMP_DATATYPE_CONST11LAMMPS_NONEE}}\phantomsection\label{\detokenize{Library_objects:_CPPv3N9LAMMPS_NS9multitype19_LMP_DATATYPE_CONST11LAMMPS_NONEE}}\phantomsection\label{\detokenize{Library_objects:_CPPv2N9LAMMPS_NS9multitype19_LMP_DATATYPE_CONST11LAMMPS_NONEE}}
\pysigstartsignatures
\pysigstartmultiline
\pysigline{\phantomsection\label{\detokenize{Library_objects:structLAMMPS__NS_1_1multitype_1a8f60df5cd520b67959d22185afcdacc5a2092f040c8fd7fd46efa6591c4f03515}}\DUrole{k}{enumerator}\DUrole{w}{ }\sphinxbfcode{\sphinxupquote{\DUrole{n}{LAMMPS\_NONE}}}}
\pysigstopmultiline
\pysigstopsignatures
\sphinxAtStartPar
no data type assigned (yet) 

\end{fulllineitems}

\index{multitype::\_LMP\_DATATYPE\_CONST::LAMMPS\_INT (C++ enumerator)@\spxentry{multitype::\_LMP\_DATATYPE\_CONST::LAMMPS\_INT}\spxextra{C++ enumerator}}

\begin{fulllineitems}
\phantomsection\label{\detokenize{Library_objects:_CPPv4N9LAMMPS_NS9multitype19_LMP_DATATYPE_CONST10LAMMPS_INTE}}\phantomsection\label{\detokenize{Library_objects:_CPPv3N9LAMMPS_NS9multitype19_LMP_DATATYPE_CONST10LAMMPS_INTE}}\phantomsection\label{\detokenize{Library_objects:_CPPv2N9LAMMPS_NS9multitype19_LMP_DATATYPE_CONST10LAMMPS_INTE}}
\pysigstartsignatures
\pysigstartmultiline
\pysigline{\phantomsection\label{\detokenize{Library_objects:structLAMMPS__NS_1_1multitype_1a8f60df5cd520b67959d22185afcdacc5a6f4597ad5c533e63a78fe90151d4d348}}\DUrole{k}{enumerator}\DUrole{w}{ }\sphinxbfcode{\sphinxupquote{\DUrole{n}{LAMMPS\_INT}}}}
\pysigstopmultiline
\pysigstopsignatures
\sphinxAtStartPar
32\sphinxhyphen{}bit integer (array) 

\end{fulllineitems}

\index{multitype::\_LMP\_DATATYPE\_CONST::LAMMPS\_INT\_2D (C++ enumerator)@\spxentry{multitype::\_LMP\_DATATYPE\_CONST::LAMMPS\_INT\_2D}\spxextra{C++ enumerator}}

\begin{fulllineitems}
\phantomsection\label{\detokenize{Library_objects:_CPPv4N9LAMMPS_NS9multitype19_LMP_DATATYPE_CONST13LAMMPS_INT_2DE}}\phantomsection\label{\detokenize{Library_objects:_CPPv3N9LAMMPS_NS9multitype19_LMP_DATATYPE_CONST13LAMMPS_INT_2DE}}\phantomsection\label{\detokenize{Library_objects:_CPPv2N9LAMMPS_NS9multitype19_LMP_DATATYPE_CONST13LAMMPS_INT_2DE}}
\pysigstartsignatures
\pysigstartmultiline
\pysigline{\phantomsection\label{\detokenize{Library_objects:structLAMMPS__NS_1_1multitype_1a8f60df5cd520b67959d22185afcdacc5acdc5b9e5771b685c0f5c0377a06380da}}\DUrole{k}{enumerator}\DUrole{w}{ }\sphinxbfcode{\sphinxupquote{\DUrole{n}{LAMMPS\_INT\_2D}}}}
\pysigstopmultiline
\pysigstopsignatures
\sphinxAtStartPar
two\sphinxhyphen{}dimensional 32\sphinxhyphen{}bit integer array 

\end{fulllineitems}

\index{multitype::\_LMP\_DATATYPE\_CONST::LAMMPS\_DOUBLE (C++ enumerator)@\spxentry{multitype::\_LMP\_DATATYPE\_CONST::LAMMPS\_DOUBLE}\spxextra{C++ enumerator}}

\begin{fulllineitems}
\phantomsection\label{\detokenize{Library_objects:_CPPv4N9LAMMPS_NS9multitype19_LMP_DATATYPE_CONST13LAMMPS_DOUBLEE}}\phantomsection\label{\detokenize{Library_objects:_CPPv3N9LAMMPS_NS9multitype19_LMP_DATATYPE_CONST13LAMMPS_DOUBLEE}}\phantomsection\label{\detokenize{Library_objects:_CPPv2N9LAMMPS_NS9multitype19_LMP_DATATYPE_CONST13LAMMPS_DOUBLEE}}
\pysigstartsignatures
\pysigstartmultiline
\pysigline{\phantomsection\label{\detokenize{Library_objects:structLAMMPS__NS_1_1multitype_1a8f60df5cd520b67959d22185afcdacc5a5a6c3a30d227047ab64e210c8a2fd468}}\DUrole{k}{enumerator}\DUrole{w}{ }\sphinxbfcode{\sphinxupquote{\DUrole{n}{LAMMPS\_DOUBLE}}}}
\pysigstopmultiline
\pysigstopsignatures
\sphinxAtStartPar
64\sphinxhyphen{}bit double (array) 

\end{fulllineitems}

\index{multitype::\_LMP\_DATATYPE\_CONST::LAMMPS\_DOUBLE\_2D (C++ enumerator)@\spxentry{multitype::\_LMP\_DATATYPE\_CONST::LAMMPS\_DOUBLE\_2D}\spxextra{C++ enumerator}}

\begin{fulllineitems}
\phantomsection\label{\detokenize{Library_objects:_CPPv4N9LAMMPS_NS9multitype19_LMP_DATATYPE_CONST16LAMMPS_DOUBLE_2DE}}\phantomsection\label{\detokenize{Library_objects:_CPPv3N9LAMMPS_NS9multitype19_LMP_DATATYPE_CONST16LAMMPS_DOUBLE_2DE}}\phantomsection\label{\detokenize{Library_objects:_CPPv2N9LAMMPS_NS9multitype19_LMP_DATATYPE_CONST16LAMMPS_DOUBLE_2DE}}
\pysigstartsignatures
\pysigstartmultiline
\pysigline{\phantomsection\label{\detokenize{Library_objects:structLAMMPS__NS_1_1multitype_1a8f60df5cd520b67959d22185afcdacc5a515e6e960fa5de1f4e22596ede33fdeb}}\DUrole{k}{enumerator}\DUrole{w}{ }\sphinxbfcode{\sphinxupquote{\DUrole{n}{LAMMPS\_DOUBLE\_2D}}}}
\pysigstopmultiline
\pysigstopsignatures
\sphinxAtStartPar
two\sphinxhyphen{}dimensional 64\sphinxhyphen{}bit double array 

\end{fulllineitems}

\index{multitype::\_LMP\_DATATYPE\_CONST::LAMMPS\_INT64 (C++ enumerator)@\spxentry{multitype::\_LMP\_DATATYPE\_CONST::LAMMPS\_INT64}\spxextra{C++ enumerator}}

\begin{fulllineitems}
\phantomsection\label{\detokenize{Library_objects:_CPPv4N9LAMMPS_NS9multitype19_LMP_DATATYPE_CONST12LAMMPS_INT64E}}\phantomsection\label{\detokenize{Library_objects:_CPPv3N9LAMMPS_NS9multitype19_LMP_DATATYPE_CONST12LAMMPS_INT64E}}\phantomsection\label{\detokenize{Library_objects:_CPPv2N9LAMMPS_NS9multitype19_LMP_DATATYPE_CONST12LAMMPS_INT64E}}
\pysigstartsignatures
\pysigstartmultiline
\pysigline{\phantomsection\label{\detokenize{Library_objects:structLAMMPS__NS_1_1multitype_1a8f60df5cd520b67959d22185afcdacc5a3e0db9f4dfd611ebb971372cc245e80c}}\DUrole{k}{enumerator}\DUrole{w}{ }\sphinxbfcode{\sphinxupquote{\DUrole{n}{LAMMPS\_INT64}}}}
\pysigstopmultiline
\pysigstopsignatures
\sphinxAtStartPar
64\sphinxhyphen{}bit integer (array) 

\end{fulllineitems}

\index{multitype::\_LMP\_DATATYPE\_CONST::LAMMPS\_INT64\_2D (C++ enumerator)@\spxentry{multitype::\_LMP\_DATATYPE\_CONST::LAMMPS\_INT64\_2D}\spxextra{C++ enumerator}}

\begin{fulllineitems}
\phantomsection\label{\detokenize{Library_objects:_CPPv4N9LAMMPS_NS9multitype19_LMP_DATATYPE_CONST15LAMMPS_INT64_2DE}}\phantomsection\label{\detokenize{Library_objects:_CPPv3N9LAMMPS_NS9multitype19_LMP_DATATYPE_CONST15LAMMPS_INT64_2DE}}\phantomsection\label{\detokenize{Library_objects:_CPPv2N9LAMMPS_NS9multitype19_LMP_DATATYPE_CONST15LAMMPS_INT64_2DE}}
\pysigstartsignatures
\pysigstartmultiline
\pysigline{\phantomsection\label{\detokenize{Library_objects:structLAMMPS__NS_1_1multitype_1a8f60df5cd520b67959d22185afcdacc5abfc022018c51b292be00cbc87760ab60}}\DUrole{k}{enumerator}\DUrole{w}{ }\sphinxbfcode{\sphinxupquote{\DUrole{n}{LAMMPS\_INT64\_2D}}}}
\pysigstopmultiline
\pysigstopsignatures
\sphinxAtStartPar
two\sphinxhyphen{}dimensional 64\sphinxhyphen{}bit integer array 

\end{fulllineitems}

\index{multitype::\_LMP\_DATATYPE\_CONST::LAMMPS\_STRING (C++ enumerator)@\spxentry{multitype::\_LMP\_DATATYPE\_CONST::LAMMPS\_STRING}\spxextra{C++ enumerator}}

\begin{fulllineitems}
\phantomsection\label{\detokenize{Library_objects:_CPPv4N9LAMMPS_NS9multitype19_LMP_DATATYPE_CONST13LAMMPS_STRINGE}}\phantomsection\label{\detokenize{Library_objects:_CPPv3N9LAMMPS_NS9multitype19_LMP_DATATYPE_CONST13LAMMPS_STRINGE}}\phantomsection\label{\detokenize{Library_objects:_CPPv2N9LAMMPS_NS9multitype19_LMP_DATATYPE_CONST13LAMMPS_STRINGE}}
\pysigstartsignatures
\pysigstartmultiline
\pysigline{\phantomsection\label{\detokenize{Library_objects:structLAMMPS__NS_1_1multitype_1a8f60df5cd520b67959d22185afcdacc5ae8c997bf96560c2d658a3d2a01518d72}}\DUrole{k}{enumerator}\DUrole{w}{ }\sphinxbfcode{\sphinxupquote{\DUrole{n}{LAMMPS\_STRING}}}}
\pysigstopmultiline
\pysigstopsignatures
\sphinxAtStartPar
C\sphinxhyphen{}String 

\end{fulllineitems}


\end{fulllineitems}

\index{\_LMP\_STYLE\_CONST (C++ enum)@\spxentry{\_LMP\_STYLE\_CONST}\spxextra{C++ enum}}

\begin{fulllineitems}
\phantomsection\label{\detokenize{Library_objects:_CPPv416_LMP_STYLE_CONST}}\phantomsection\label{\detokenize{Library_objects:_CPPv316_LMP_STYLE_CONST}}\phantomsection\label{\detokenize{Library_objects:_CPPv216_LMP_STYLE_CONST}}
\pysigstartsignatures
\pysigstartmultiline
\pysigline{\phantomsection\label{\detokenize{Library_objects:library_8h_1af78e950523b27d0aa9b74ccf93666d36}}\DUrole{k}{enum}\DUrole{w}{ }\sphinxbfcode{\sphinxupquote{\DUrole{n}{\_LMP\_STYLE\_CONST}}}}
\pysigstopmultiline
\pysigstopsignatures
\sphinxAtStartPar
Style constants for extracting data from computes and fixes.

\sphinxAtStartPar
Must be kept in sync with the equivalent constants in \sphinxcode{\sphinxupquote{python/lammps/constants.py}}, \sphinxcode{\sphinxupquote{fortran/lammps.f90}}, \sphinxcode{\sphinxupquote{tools/swig/lammps.i}}, and \sphinxcode{\sphinxupquote{examples/COUPLE/plugin/liblammpsplugin.h}}

\sphinxAtStartPar
\sphinxstyleemphasis{Values:}
\index{LMP\_STYLE\_GLOBAL (C++ enumerator)@\spxentry{LMP\_STYLE\_GLOBAL}\spxextra{C++ enumerator}}

\begin{fulllineitems}
\phantomsection\label{\detokenize{Library_objects:_CPPv4N16_LMP_STYLE_CONST16LMP_STYLE_GLOBALE}}\phantomsection\label{\detokenize{Library_objects:_CPPv3N16_LMP_STYLE_CONST16LMP_STYLE_GLOBALE}}\phantomsection\label{\detokenize{Library_objects:_CPPv2N16_LMP_STYLE_CONST16LMP_STYLE_GLOBALE}}
\pysigstartsignatures
\pysigstartmultiline
\pysigline{\phantomsection\label{\detokenize{Library_objects:library_8h_1af78e950523b27d0aa9b74ccf93666d36a5d9f763d2a7c13c4aea83405aaf4b3f3}}\DUrole{k}{enumerator}\DUrole{w}{ }\sphinxbfcode{\sphinxupquote{\DUrole{n}{LMP\_STYLE\_GLOBAL}}}}
\pysigstopmultiline
\pysigstopsignatures
\sphinxAtStartPar
return global data 

\end{fulllineitems}

\index{LMP\_STYLE\_ATOM (C++ enumerator)@\spxentry{LMP\_STYLE\_ATOM}\spxextra{C++ enumerator}}

\begin{fulllineitems}
\phantomsection\label{\detokenize{Library_objects:_CPPv4N16_LMP_STYLE_CONST14LMP_STYLE_ATOME}}\phantomsection\label{\detokenize{Library_objects:_CPPv3N16_LMP_STYLE_CONST14LMP_STYLE_ATOME}}\phantomsection\label{\detokenize{Library_objects:_CPPv2N16_LMP_STYLE_CONST14LMP_STYLE_ATOME}}
\pysigstartsignatures
\pysigstartmultiline
\pysigline{\phantomsection\label{\detokenize{Library_objects:library_8h_1af78e950523b27d0aa9b74ccf93666d36aa1e72b0a28e28261262774733cdbadfa}}\DUrole{k}{enumerator}\DUrole{w}{ }\sphinxbfcode{\sphinxupquote{\DUrole{n}{LMP\_STYLE\_ATOM}}}}
\pysigstopmultiline
\pysigstopsignatures
\sphinxAtStartPar
return per\sphinxhyphen{}atom data 

\end{fulllineitems}

\index{LMP\_STYLE\_LOCAL (C++ enumerator)@\spxentry{LMP\_STYLE\_LOCAL}\spxextra{C++ enumerator}}

\begin{fulllineitems}
\phantomsection\label{\detokenize{Library_objects:_CPPv4N16_LMP_STYLE_CONST15LMP_STYLE_LOCALE}}\phantomsection\label{\detokenize{Library_objects:_CPPv3N16_LMP_STYLE_CONST15LMP_STYLE_LOCALE}}\phantomsection\label{\detokenize{Library_objects:_CPPv2N16_LMP_STYLE_CONST15LMP_STYLE_LOCALE}}
\pysigstartsignatures
\pysigstartmultiline
\pysigline{\phantomsection\label{\detokenize{Library_objects:library_8h_1af78e950523b27d0aa9b74ccf93666d36a9b74ce240cd632e21f8c1a5a1cf07ed0}}\DUrole{k}{enumerator}\DUrole{w}{ }\sphinxbfcode{\sphinxupquote{\DUrole{n}{LMP\_STYLE\_LOCAL}}}}
\pysigstopmultiline
\pysigstopsignatures
\sphinxAtStartPar
return local data 

\end{fulllineitems}


\end{fulllineitems}

\index{\_LMP\_TYPE\_CONST (C++ enum)@\spxentry{\_LMP\_TYPE\_CONST}\spxextra{C++ enum}}

\begin{fulllineitems}
\phantomsection\label{\detokenize{Library_objects:_CPPv415_LMP_TYPE_CONST}}\phantomsection\label{\detokenize{Library_objects:_CPPv315_LMP_TYPE_CONST}}\phantomsection\label{\detokenize{Library_objects:_CPPv215_LMP_TYPE_CONST}}
\pysigstartsignatures
\pysigstartmultiline
\pysigline{\phantomsection\label{\detokenize{Library_objects:library_8h_1a30cd6a0606720be7addec2816760879e}}\DUrole{k}{enum}\DUrole{w}{ }\sphinxbfcode{\sphinxupquote{\DUrole{n}{\_LMP\_TYPE\_CONST}}}}
\pysigstopmultiline
\pysigstopsignatures
\sphinxAtStartPar
Type and size constants for extracting data from computes and fixes.

\sphinxAtStartPar
Must be kept in sync with the equivalent constants in \sphinxcode{\sphinxupquote{python/lammps/constants.py}}, \sphinxcode{\sphinxupquote{fortran/lammps.f90}}, \sphinxcode{\sphinxupquote{tools/swig/lammps.i}}, and \sphinxcode{\sphinxupquote{examples/COUPLE/plugin/liblammpsplugin.h}}

\sphinxAtStartPar
\sphinxstyleemphasis{Values:}
\index{LMP\_TYPE\_SCALAR (C++ enumerator)@\spxentry{LMP\_TYPE\_SCALAR}\spxextra{C++ enumerator}}

\begin{fulllineitems}
\phantomsection\label{\detokenize{Library_objects:_CPPv4N15_LMP_TYPE_CONST15LMP_TYPE_SCALARE}}\phantomsection\label{\detokenize{Library_objects:_CPPv3N15_LMP_TYPE_CONST15LMP_TYPE_SCALARE}}\phantomsection\label{\detokenize{Library_objects:_CPPv2N15_LMP_TYPE_CONST15LMP_TYPE_SCALARE}}
\pysigstartsignatures
\pysigstartmultiline
\pysigline{\phantomsection\label{\detokenize{Library_objects:library_8h_1a30cd6a0606720be7addec2816760879ea3a0d5e6c49cc929d953b366ed58c821b}}\DUrole{k}{enumerator}\DUrole{w}{ }\sphinxbfcode{\sphinxupquote{\DUrole{n}{LMP\_TYPE\_SCALAR}}}}
\pysigstopmultiline
\pysigstopsignatures
\sphinxAtStartPar
return scalar 

\end{fulllineitems}

\index{LMP\_TYPE\_VECTOR (C++ enumerator)@\spxentry{LMP\_TYPE\_VECTOR}\spxextra{C++ enumerator}}

\begin{fulllineitems}
\phantomsection\label{\detokenize{Library_objects:_CPPv4N15_LMP_TYPE_CONST15LMP_TYPE_VECTORE}}\phantomsection\label{\detokenize{Library_objects:_CPPv3N15_LMP_TYPE_CONST15LMP_TYPE_VECTORE}}\phantomsection\label{\detokenize{Library_objects:_CPPv2N15_LMP_TYPE_CONST15LMP_TYPE_VECTORE}}
\pysigstartsignatures
\pysigstartmultiline
\pysigline{\phantomsection\label{\detokenize{Library_objects:library_8h_1a30cd6a0606720be7addec2816760879ead7beec11786a3dc637fcdd2ac62506e8}}\DUrole{k}{enumerator}\DUrole{w}{ }\sphinxbfcode{\sphinxupquote{\DUrole{n}{LMP\_TYPE\_VECTOR}}}}
\pysigstopmultiline
\pysigstopsignatures
\sphinxAtStartPar
return vector 

\end{fulllineitems}

\index{LMP\_TYPE\_ARRAY (C++ enumerator)@\spxentry{LMP\_TYPE\_ARRAY}\spxextra{C++ enumerator}}

\begin{fulllineitems}
\phantomsection\label{\detokenize{Library_objects:_CPPv4N15_LMP_TYPE_CONST14LMP_TYPE_ARRAYE}}\phantomsection\label{\detokenize{Library_objects:_CPPv3N15_LMP_TYPE_CONST14LMP_TYPE_ARRAYE}}\phantomsection\label{\detokenize{Library_objects:_CPPv2N15_LMP_TYPE_CONST14LMP_TYPE_ARRAYE}}
\pysigstartsignatures
\pysigstartmultiline
\pysigline{\phantomsection\label{\detokenize{Library_objects:library_8h_1a30cd6a0606720be7addec2816760879eaaacd44f7a32b76b06051cb01ff40451a}}\DUrole{k}{enumerator}\DUrole{w}{ }\sphinxbfcode{\sphinxupquote{\DUrole{n}{LMP\_TYPE\_ARRAY}}}}
\pysigstopmultiline
\pysigstopsignatures
\sphinxAtStartPar
return array 

\end{fulllineitems}

\index{LMP\_SIZE\_VECTOR (C++ enumerator)@\spxentry{LMP\_SIZE\_VECTOR}\spxextra{C++ enumerator}}

\begin{fulllineitems}
\phantomsection\label{\detokenize{Library_objects:_CPPv4N15_LMP_TYPE_CONST15LMP_SIZE_VECTORE}}\phantomsection\label{\detokenize{Library_objects:_CPPv3N15_LMP_TYPE_CONST15LMP_SIZE_VECTORE}}\phantomsection\label{\detokenize{Library_objects:_CPPv2N15_LMP_TYPE_CONST15LMP_SIZE_VECTORE}}
\pysigstartsignatures
\pysigstartmultiline
\pysigline{\phantomsection\label{\detokenize{Library_objects:library_8h_1a30cd6a0606720be7addec2816760879eab7e39a6a361f8bc210064c9d853b4ded}}\DUrole{k}{enumerator}\DUrole{w}{ }\sphinxbfcode{\sphinxupquote{\DUrole{n}{LMP\_SIZE\_VECTOR}}}}
\pysigstopmultiline
\pysigstopsignatures
\sphinxAtStartPar
return length of vector 

\end{fulllineitems}

\index{LMP\_SIZE\_ROWS (C++ enumerator)@\spxentry{LMP\_SIZE\_ROWS}\spxextra{C++ enumerator}}

\begin{fulllineitems}
\phantomsection\label{\detokenize{Library_objects:_CPPv4N15_LMP_TYPE_CONST13LMP_SIZE_ROWSE}}\phantomsection\label{\detokenize{Library_objects:_CPPv3N15_LMP_TYPE_CONST13LMP_SIZE_ROWSE}}\phantomsection\label{\detokenize{Library_objects:_CPPv2N15_LMP_TYPE_CONST13LMP_SIZE_ROWSE}}
\pysigstartsignatures
\pysigstartmultiline
\pysigline{\phantomsection\label{\detokenize{Library_objects:library_8h_1a30cd6a0606720be7addec2816760879eaa65d78d0215fbb5d1f8cf6dfa05ed485}}\DUrole{k}{enumerator}\DUrole{w}{ }\sphinxbfcode{\sphinxupquote{\DUrole{n}{LMP\_SIZE\_ROWS}}}}
\pysigstopmultiline
\pysigstopsignatures
\sphinxAtStartPar
return number of rows 

\end{fulllineitems}

\index{LMP\_SIZE\_COLS (C++ enumerator)@\spxentry{LMP\_SIZE\_COLS}\spxextra{C++ enumerator}}

\begin{fulllineitems}
\phantomsection\label{\detokenize{Library_objects:_CPPv4N15_LMP_TYPE_CONST13LMP_SIZE_COLSE}}\phantomsection\label{\detokenize{Library_objects:_CPPv3N15_LMP_TYPE_CONST13LMP_SIZE_COLSE}}\phantomsection\label{\detokenize{Library_objects:_CPPv2N15_LMP_TYPE_CONST13LMP_SIZE_COLSE}}
\pysigstartsignatures
\pysigstartmultiline
\pysigline{\phantomsection\label{\detokenize{Library_objects:library_8h_1a30cd6a0606720be7addec2816760879ea0213053b1c6d4283bb6d662a00fecc16}}\DUrole{k}{enumerator}\DUrole{w}{ }\sphinxbfcode{\sphinxupquote{\DUrole{n}{LMP\_SIZE\_COLS}}}}
\pysigstopmultiline
\pysigstopsignatures
\sphinxAtStartPar
return number of columns 

\end{fulllineitems}


\end{fulllineitems}

\index{\_LMP\_VAR\_CONST (C++ enum)@\spxentry{\_LMP\_VAR\_CONST}\spxextra{C++ enum}}

\begin{fulllineitems}
\phantomsection\label{\detokenize{Library_objects:_CPPv414_LMP_VAR_CONST}}\phantomsection\label{\detokenize{Library_objects:_CPPv314_LMP_VAR_CONST}}\phantomsection\label{\detokenize{Library_objects:_CPPv214_LMP_VAR_CONST}}
\pysigstartsignatures
\pysigstartmultiline
\pysigline{\phantomsection\label{\detokenize{Library_objects:library_8h_1a09d7d9d1e5512d11257c5b6fce67ee37}}\DUrole{k}{enum}\DUrole{w}{ }\sphinxbfcode{\sphinxupquote{\DUrole{n}{\_LMP\_VAR\_CONST}}}}
\pysigstopmultiline
\pysigstopsignatures
\sphinxAtStartPar
Variable style constants for extracting data from variables.

\sphinxAtStartPar
Must be kept in sync with the equivalent constants in \sphinxcode{\sphinxupquote{python/lammps/constants.py}}, \sphinxcode{\sphinxupquote{fortran/lammps.f90}}, \sphinxcode{\sphinxupquote{tools/swig/lammps.i}}, and \sphinxcode{\sphinxupquote{examples/COUPLE/plugin/liblammpsplugin.h}}

\sphinxAtStartPar
\sphinxstyleemphasis{Values:}
\index{LMP\_VAR\_EQUAL (C++ enumerator)@\spxentry{LMP\_VAR\_EQUAL}\spxextra{C++ enumerator}}

\begin{fulllineitems}
\phantomsection\label{\detokenize{Library_objects:_CPPv4N14_LMP_VAR_CONST13LMP_VAR_EQUALE}}\phantomsection\label{\detokenize{Library_objects:_CPPv3N14_LMP_VAR_CONST13LMP_VAR_EQUALE}}\phantomsection\label{\detokenize{Library_objects:_CPPv2N14_LMP_VAR_CONST13LMP_VAR_EQUALE}}
\pysigstartsignatures
\pysigstartmultiline
\pysigline{\phantomsection\label{\detokenize{Library_objects:library_8h_1a09d7d9d1e5512d11257c5b6fce67ee37aeb6a1c0a08cdfb496bcaef2cc37d5047}}\DUrole{k}{enumerator}\DUrole{w}{ }\sphinxbfcode{\sphinxupquote{\DUrole{n}{LMP\_VAR\_EQUAL}}}}
\pysigstopmultiline
\pysigstopsignatures
\sphinxAtStartPar
compatible with equal\sphinxhyphen{}style variables 

\end{fulllineitems}

\index{LMP\_VAR\_ATOM (C++ enumerator)@\spxentry{LMP\_VAR\_ATOM}\spxextra{C++ enumerator}}

\begin{fulllineitems}
\phantomsection\label{\detokenize{Library_objects:_CPPv4N14_LMP_VAR_CONST12LMP_VAR_ATOME}}\phantomsection\label{\detokenize{Library_objects:_CPPv3N14_LMP_VAR_CONST12LMP_VAR_ATOME}}\phantomsection\label{\detokenize{Library_objects:_CPPv2N14_LMP_VAR_CONST12LMP_VAR_ATOME}}
\pysigstartsignatures
\pysigstartmultiline
\pysigline{\phantomsection\label{\detokenize{Library_objects:library_8h_1a09d7d9d1e5512d11257c5b6fce67ee37a4a336ec13f20754d421f7d0d218cc1a4}}\DUrole{k}{enumerator}\DUrole{w}{ }\sphinxbfcode{\sphinxupquote{\DUrole{n}{LMP\_VAR\_ATOM}}}}
\pysigstopmultiline
\pysigstopsignatures
\sphinxAtStartPar
compatible with atom\sphinxhyphen{}style variables 

\end{fulllineitems}

\index{LMP\_VAR\_VECTOR (C++ enumerator)@\spxentry{LMP\_VAR\_VECTOR}\spxextra{C++ enumerator}}

\begin{fulllineitems}
\phantomsection\label{\detokenize{Library_objects:_CPPv4N14_LMP_VAR_CONST14LMP_VAR_VECTORE}}\phantomsection\label{\detokenize{Library_objects:_CPPv3N14_LMP_VAR_CONST14LMP_VAR_VECTORE}}\phantomsection\label{\detokenize{Library_objects:_CPPv2N14_LMP_VAR_CONST14LMP_VAR_VECTORE}}
\pysigstartsignatures
\pysigstartmultiline
\pysigline{\phantomsection\label{\detokenize{Library_objects:library_8h_1a09d7d9d1e5512d11257c5b6fce67ee37a381da764a4e3a5f1bc183cc9905be43c}}\DUrole{k}{enumerator}\DUrole{w}{ }\sphinxbfcode{\sphinxupquote{\DUrole{n}{LMP\_VAR\_VECTOR}}}}
\pysigstopmultiline
\pysigstopsignatures
\sphinxAtStartPar
compatible with vector\sphinxhyphen{}style variables 

\end{fulllineitems}

\index{LMP\_VAR\_STRING (C++ enumerator)@\spxentry{LMP\_VAR\_STRING}\spxextra{C++ enumerator}}

\begin{fulllineitems}
\phantomsection\label{\detokenize{Library_objects:_CPPv4N14_LMP_VAR_CONST14LMP_VAR_STRINGE}}\phantomsection\label{\detokenize{Library_objects:_CPPv3N14_LMP_VAR_CONST14LMP_VAR_STRINGE}}\phantomsection\label{\detokenize{Library_objects:_CPPv2N14_LMP_VAR_CONST14LMP_VAR_STRINGE}}
\pysigstartsignatures
\pysigstartmultiline
\pysigline{\phantomsection\label{\detokenize{Library_objects:library_8h_1a09d7d9d1e5512d11257c5b6fce67ee37a52551a127731cf3f7e193b9d1d83e6e8}}\DUrole{k}{enumerator}\DUrole{w}{ }\sphinxbfcode{\sphinxupquote{\DUrole{n}{LMP\_VAR\_STRING}}}}
\pysigstopmultiline
\pysigstopsignatures
\sphinxAtStartPar
return value will be a string (catch\sphinxhyphen{}all) 

\end{fulllineitems}


\end{fulllineitems}


\sphinxstepscope


\subsection{Scatter/gather operations}
\label{\detokenize{Library_scatter:scatter-gather-operations}}\label{\detokenize{Library_scatter::doc}}
\sphinxAtStartPar
This section has functions which gather per\sphinxhyphen{}atom data from one or more
processors into a contiguous global list ordered by atom ID.  The same
list is returned to all calling processors.  It also contains
functions which scatter per\sphinxhyphen{}atom data from a contiguous global list
across the processors that own those atom IDs.  It also has a
create\_atoms() function which can create new atoms by scattering them
appropriately to owning processors in the LAMMPS spatial
decomposition.

\sphinxAtStartPar
It documents the following functions:
\begin{itemize}
\item {} 
\sphinxAtStartPar
{\hyperref[\detokenize{Library_scatter:_CPPv419lammps_gather_atomsPvPKciiPv}]{\sphinxcrossref{\sphinxcode{\sphinxupquote{lammps\_gather\_atoms()}}}}}

\item {} 
\sphinxAtStartPar
{\hyperref[\detokenize{Library_scatter:_CPPv426lammps_gather_atoms_concatPvPKciiPv}]{\sphinxcrossref{\sphinxcode{\sphinxupquote{lammps\_gather\_atoms\_concat()}}}}}

\item {} 
\sphinxAtStartPar
{\hyperref[\detokenize{Library_scatter:_CPPv426lammps_gather_atoms_subsetPvPKciiiPiPv}]{\sphinxcrossref{\sphinxcode{\sphinxupquote{lammps\_gather\_atoms\_subset()}}}}}

\item {} 
\sphinxAtStartPar
{\hyperref[\detokenize{Library_scatter:_CPPv420lammps_scatter_atomsPvPKciiPv}]{\sphinxcrossref{\sphinxcode{\sphinxupquote{lammps\_scatter\_atoms()}}}}}

\item {} 
\sphinxAtStartPar
{\hyperref[\detokenize{Library_scatter:_CPPv427lammps_scatter_atoms_subsetPvPKciiiPiPv}]{\sphinxcrossref{\sphinxcode{\sphinxupquote{lammps\_scatter\_atoms\_subset()}}}}}

\item {} 
\sphinxAtStartPar
{\hyperref[\detokenize{Library_scatter:_CPPv419lammps_gather_bondsPvPv}]{\sphinxcrossref{\sphinxcode{\sphinxupquote{lammps\_gather\_bonds()}}}}}

\item {} 
\sphinxAtStartPar
{\hyperref[\detokenize{Library_scatter:_CPPv420lammps_gather_anglesPvPv}]{\sphinxcrossref{\sphinxcode{\sphinxupquote{lammps\_gather\_angles()}}}}}

\item {} 
\sphinxAtStartPar
{\hyperref[\detokenize{Library_scatter:_CPPv423lammps_gather_dihedralsPvPv}]{\sphinxcrossref{\sphinxcode{\sphinxupquote{lammps\_gather\_dihedrals()}}}}}

\item {} 
\sphinxAtStartPar
{\hyperref[\detokenize{Library_scatter:_CPPv423lammps_gather_impropersPvPv}]{\sphinxcrossref{\sphinxcode{\sphinxupquote{lammps\_gather\_impropers()}}}}}

\item {} 
\sphinxAtStartPar
{\hyperref[\detokenize{Library_scatter:_CPPv413lammps_gatherPvPKciiPv}]{\sphinxcrossref{\sphinxcode{\sphinxupquote{lammps\_gather()}}}}}

\item {} 
\sphinxAtStartPar
{\hyperref[\detokenize{Library_scatter:_CPPv420lammps_gather_concatPvPKciiPv}]{\sphinxcrossref{\sphinxcode{\sphinxupquote{lammps\_gather\_concat()}}}}}

\item {} 
\sphinxAtStartPar
{\hyperref[\detokenize{Library_scatter:_CPPv420lammps_gather_subsetPvPKciiiPiPv}]{\sphinxcrossref{\sphinxcode{\sphinxupquote{lammps\_gather\_subset()}}}}}

\item {} 
\sphinxAtStartPar
{\hyperref[\detokenize{Library_scatter:_CPPv414lammps_scatterPvPKciiPv}]{\sphinxcrossref{\sphinxcode{\sphinxupquote{lammps\_scatter()}}}}}

\item {} 
\sphinxAtStartPar
{\hyperref[\detokenize{Library_scatter:_CPPv421lammps_scatter_subsetPvPKciiiPiPv}]{\sphinxcrossref{\sphinxcode{\sphinxupquote{lammps\_scatter\_subset()}}}}}

\item {} 
\sphinxAtStartPar
{\hyperref[\detokenize{Library_scatter:_CPPv419lammps_create_atomsPviPKiPKiPKdPKdPKii}]{\sphinxcrossref{\sphinxcode{\sphinxupquote{lammps\_create\_atoms()}}}}}

\end{itemize}


\bigskip\hrule\bigskip

\index{lammps\_gather\_atoms (C++ function)@\spxentry{lammps\_gather\_atoms}\spxextra{C++ function}}

\begin{fulllineitems}
\phantomsection\label{\detokenize{Library_scatter:_CPPv419lammps_gather_atomsPvPKciiPv}}\phantomsection\label{\detokenize{Library_scatter:_CPPv319lammps_gather_atomsPvPKciiPv}}\phantomsection\label{\detokenize{Library_scatter:_CPPv219lammps_gather_atomsPvPKciiPv}}\phantomsection\label{\detokenize{Library_scatter:lammps_gather_atoms__voidP.cCP.i.i.voidP}}
\pysigstartsignatures
\pysigstartmultiline
\pysiglinewithargsret{\phantomsection\label{\detokenize{Library_scatter:library_8h_1ad3208ce7d03d54f74573a80565dc335d}}\DUrole{kt}{void}\DUrole{w}{ }\sphinxbfcode{\sphinxupquote{\DUrole{n}{lammps\_gather\_atoms}}}}{\DUrole{kt}{void}\DUrole{w}{ }\DUrole{p}{*}\DUrole{n,sig-param}{handle}\sphinxparamcomma \DUrole{k}{const}\DUrole{w}{ }\DUrole{kt}{char}\DUrole{w}{ }\DUrole{p}{*}\DUrole{n,sig-param}{name}\sphinxparamcomma \DUrole{kt}{int}\DUrole{w}{ }\DUrole{n,sig-param}{type}\sphinxparamcomma \DUrole{kt}{int}\DUrole{w}{ }\DUrole{n,sig-param}{count}\sphinxparamcomma \DUrole{kt}{void}\DUrole{w}{ }\DUrole{p}{*}\DUrole{n,sig-param}{data}}{}
\pysigstopmultiline
\pysigstopsignatures
\sphinxAtStartPar
Gather the named atom\sphinxhyphen{}based entity for all atoms across all processes, in order.

\sphinxAtStartPar

\sphinxAtStartPar
This subroutine gathers data for all atoms and stores them in a
one\sphinxhyphen{}dimensional array allocated by the user. The data will be ordered by
atom ID, which requires consecutive atom IDs (1 to \sphinxstyleemphasis{natoms}). If you need
a similar array but have non\sphinxhyphen{}consecutive atom IDs, see
{\hyperref[\detokenize{Library_scatter:_CPPv426lammps_gather_atoms_concatPvPKciiPv}]{\sphinxcrossref{\sphinxcode{\sphinxupquote{lammps\_gather\_atoms\_concat()}}}}}; for a similar array but for a subset
of atoms, see {\hyperref[\detokenize{Library_scatter:_CPPv426lammps_gather_atoms_subsetPvPKciiiPiPv}]{\sphinxcrossref{\sphinxcode{\sphinxupquote{lammps\_gather\_atoms\_subset()}}}}}.

\sphinxAtStartPar
The \sphinxstyleemphasis{data} array will be ordered in groups of \sphinxstyleemphasis{count} values, sorted by atom ID
(e.g., if \sphinxstyleemphasis{name} is \sphinxstyleemphasis{x} and \sphinxstyleemphasis{count} = 3, then \sphinxstyleemphasis{data} = x{[}0{]}{[}0{]}, x{[}0{]}{[}1{]},
x{[}0{]}{[}2{]}, x{[}1{]}{[}0{]}, x{[}1{]}{[}1{]}, x{[}1{]}{[}2{]}, x{[}2{]}{[}0{]}, \(\dots\));
\sphinxstyleemphasis{data} must be pre\sphinxhyphen{}allocated by the caller to length (\sphinxstyleemphasis{count} \(\times\)
\sphinxstyleemphasis{natoms}), as queried by {\hyperref[\detokenize{Library_properties:_CPPv417lammps_get_natomsPv}]{\sphinxcrossref{\sphinxcode{\sphinxupquote{lammps\_get\_natoms()}}}}},
{\hyperref[\detokenize{Library_properties:_CPPv421lammps_extract_globalPvPKc}]{\sphinxcrossref{\sphinxcode{\sphinxupquote{lammps\_extract\_global()}}}}}, or {\hyperref[\detokenize{Library_properties:_CPPv422lammps_extract_settingPvPKc}]{\sphinxcrossref{\sphinxcode{\sphinxupquote{lammps\_extract\_setting()}}}}}.

\begin{sphinxadmonition}{note}{Restrictions}

\sphinxAtStartPar
This function is not compatible with \sphinxcode{\sphinxupquote{\sphinxhyphen{}DLAMMPS\_BIGBIG}}.

\sphinxAtStartPar
Atom IDs must be defined and consecutive.

\sphinxAtStartPar
The total number of atoms must not be more than 2147483647 (max 32\sphinxhyphen{}bit signed int).
\end{sphinxadmonition}

\begin{quote}\begin{description}
\sphinxlineitem{Parameters}\begin{itemize}
\item {} 
\sphinxAtStartPar
\sphinxstyleliteralstrong{\sphinxupquote{handle}} \textendash{} pointer to a previously created LAMMPS instance 

\item {} 
\sphinxAtStartPar
\sphinxstyleliteralstrong{\sphinxupquote{name}} \textendash{} desired quantity (e.g., \sphinxstyleemphasis{x} or \sphinxstyleemphasis{charge}) 

\item {} 
\sphinxAtStartPar
\sphinxstyleliteralstrong{\sphinxupquote{type}} \textendash{} 0 for \sphinxcode{\sphinxupquote{int}} values, 1 for \sphinxcode{\sphinxupquote{double}} values 

\item {} 
\sphinxAtStartPar
\sphinxstyleliteralstrong{\sphinxupquote{count}} \textendash{} number of per\sphinxhyphen{}atom values (e.g., 1 for \sphinxstyleemphasis{type} or \sphinxstyleemphasis{charge}, 3 for \sphinxstyleemphasis{x} or \sphinxstyleemphasis{f}); use \sphinxstyleemphasis{count} = 3 with \sphinxstyleemphasis{image} if you want a single image flag unpacked into (\sphinxstyleemphasis{x},\sphinxstyleemphasis{y},\sphinxstyleemphasis{z}) components. 

\item {} 
\sphinxAtStartPar
\sphinxstyleliteralstrong{\sphinxupquote{data}} \textendash{} per\sphinxhyphen{}atom values packed in a 1\sphinxhyphen{}dimensional array of length \sphinxstyleemphasis{natoms} * \sphinxstyleemphasis{count}. 

\end{itemize}

\end{description}\end{quote}

\end{fulllineitems}



\bigskip\hrule\bigskip

\index{lammps\_gather\_atoms\_concat (C++ function)@\spxentry{lammps\_gather\_atoms\_concat}\spxextra{C++ function}}

\begin{fulllineitems}
\phantomsection\label{\detokenize{Library_scatter:_CPPv426lammps_gather_atoms_concatPvPKciiPv}}\phantomsection\label{\detokenize{Library_scatter:_CPPv326lammps_gather_atoms_concatPvPKciiPv}}\phantomsection\label{\detokenize{Library_scatter:_CPPv226lammps_gather_atoms_concatPvPKciiPv}}\phantomsection\label{\detokenize{Library_scatter:lammps_gather_atoms_concat__voidP.cCP.i.i.voidP}}
\pysigstartsignatures
\pysigstartmultiline
\pysiglinewithargsret{\phantomsection\label{\detokenize{Library_scatter:library_8h_1a5e433d99895ca2aa9d87843bc7ee49d5}}\DUrole{kt}{void}\DUrole{w}{ }\sphinxbfcode{\sphinxupquote{\DUrole{n}{lammps\_gather\_atoms\_concat}}}}{\DUrole{kt}{void}\DUrole{w}{ }\DUrole{p}{*}\DUrole{n,sig-param}{handle}\sphinxparamcomma \DUrole{k}{const}\DUrole{w}{ }\DUrole{kt}{char}\DUrole{w}{ }\DUrole{p}{*}\DUrole{n,sig-param}{name}\sphinxparamcomma \DUrole{kt}{int}\DUrole{w}{ }\DUrole{n,sig-param}{type}\sphinxparamcomma \DUrole{kt}{int}\DUrole{w}{ }\DUrole{n,sig-param}{count}\sphinxparamcomma \DUrole{kt}{void}\DUrole{w}{ }\DUrole{p}{*}\DUrole{n,sig-param}{data}}{}
\pysigstopmultiline
\pysigstopsignatures
\sphinxAtStartPar
Gather the named atom\sphinxhyphen{}based entity for all atoms across all processes, unordered.

\sphinxAtStartPar

\sphinxAtStartPar
This subroutine gathers data for all atoms and stores them in a
one\sphinxhyphen{}dimensional array allocated by the user. The data will be a concatenation
of chunks from each processor’s owned atoms, in whatever order the atoms are
in on each processor. This process has no requirement that the atom IDs be
consecutive. If you need the ID of each atom, you can do another
{\hyperref[\detokenize{Library_scatter:_CPPv426lammps_gather_atoms_concatPvPKciiPv}]{\sphinxcrossref{\sphinxcode{\sphinxupquote{lammps\_gather\_atoms\_concat()}}}}} call with \sphinxstyleemphasis{name} set to \sphinxcode{\sphinxupquote{id}}.
If you have consecutive IDs and want the data to be in order, use
{\hyperref[\detokenize{Library_scatter:_CPPv419lammps_gather_atomsPvPKciiPv}]{\sphinxcrossref{\sphinxcode{\sphinxupquote{lammps\_gather\_atoms()}}}}}; for a similar array but for a subset
of atoms, use {\hyperref[\detokenize{Library_scatter:_CPPv426lammps_gather_atoms_subsetPvPKciiiPiPv}]{\sphinxcrossref{\sphinxcode{\sphinxupquote{lammps\_gather\_atoms\_subset()}}}}}.

\sphinxAtStartPar
The \sphinxstyleemphasis{data} array will be in groups of \sphinxstyleemphasis{count} values, with \sphinxstyleemphasis{natoms}
groups total, but not in order by atom ID (e.g., if \sphinxstyleemphasis{name} is \sphinxstyleemphasis{x} and \sphinxstyleemphasis{count}
is 3, then \sphinxstyleemphasis{data} might be something like x{[}10{]}{[}0{]}, x{[}10{]}{[}1{]}, x{[}10{]}{[}2{]},
x{[}2{]}{[}0{]}, x{[}2{]}{[}1{]}, x{[}2{]}{[}2{]}, x{[}4{]}{[}0{]}, \(\dots\)); \sphinxstyleemphasis{data} must be
pre\sphinxhyphen{}allocated by the caller to length (\sphinxstyleemphasis{count} \(\times\) \sphinxstyleemphasis{natoms}), as
queried by {\hyperref[\detokenize{Library_properties:_CPPv417lammps_get_natomsPv}]{\sphinxcrossref{\sphinxcode{\sphinxupquote{lammps\_get\_natoms()}}}}}, {\hyperref[\detokenize{Library_properties:_CPPv421lammps_extract_globalPvPKc}]{\sphinxcrossref{\sphinxcode{\sphinxupquote{lammps\_extract\_global()}}}}},
or {\hyperref[\detokenize{Library_properties:_CPPv422lammps_extract_settingPvPKc}]{\sphinxcrossref{\sphinxcode{\sphinxupquote{lammps\_extract\_setting()}}}}}.

\begin{sphinxadmonition}{note}{Restrictions}

\sphinxAtStartPar
This function is not compatible with \sphinxcode{\sphinxupquote{\sphinxhyphen{}DLAMMPS\_BIGBIG}}.

\sphinxAtStartPar
Atom IDs must be defined.

\sphinxAtStartPar
The total number of atoms must not be more than 2147483647 (max 32\sphinxhyphen{}bit signed int).
\end{sphinxadmonition}

\begin{quote}\begin{description}
\sphinxlineitem{Parameters}\begin{itemize}
\item {} 
\sphinxAtStartPar
\sphinxstyleliteralstrong{\sphinxupquote{handle}} \textendash{} pointer to a previously created LAMMPS instance 

\item {} 
\sphinxAtStartPar
\sphinxstyleliteralstrong{\sphinxupquote{name}} \textendash{} desired quantity (e.g., \sphinxstyleemphasis{x} or \sphinxstyleemphasis{charge}\textbackslash{} ) 

\item {} 
\sphinxAtStartPar
\sphinxstyleliteralstrong{\sphinxupquote{type}} \textendash{} 0 for \sphinxcode{\sphinxupquote{int}} values, 1 for \sphinxcode{\sphinxupquote{double}} values 

\item {} 
\sphinxAtStartPar
\sphinxstyleliteralstrong{\sphinxupquote{count}} \textendash{} number of per\sphinxhyphen{}atom values (e.g., 1 for \sphinxstyleemphasis{type} or \sphinxstyleemphasis{charge}, 3 for \sphinxstyleemphasis{x} or \sphinxstyleemphasis{f}); use \sphinxstyleemphasis{count} = 3 with “image” if you want single image flags unpacked into (\sphinxstyleemphasis{x},\sphinxstyleemphasis{y},\sphinxstyleemphasis{z}) 

\item {} 
\sphinxAtStartPar
\sphinxstyleliteralstrong{\sphinxupquote{data}} \textendash{} per\sphinxhyphen{}atom values packed in a 1\sphinxhyphen{}dimensional array of length \sphinxstyleemphasis{natoms} * \sphinxstyleemphasis{count}. 

\end{itemize}

\end{description}\end{quote}

\end{fulllineitems}



\bigskip\hrule\bigskip

\index{lammps\_gather\_atoms\_subset (C++ function)@\spxentry{lammps\_gather\_atoms\_subset}\spxextra{C++ function}}

\begin{fulllineitems}
\phantomsection\label{\detokenize{Library_scatter:_CPPv426lammps_gather_atoms_subsetPvPKciiiPiPv}}\phantomsection\label{\detokenize{Library_scatter:_CPPv326lammps_gather_atoms_subsetPvPKciiiPiPv}}\phantomsection\label{\detokenize{Library_scatter:_CPPv226lammps_gather_atoms_subsetPvPKciiiPiPv}}\phantomsection\label{\detokenize{Library_scatter:lammps_gather_atoms_subset__voidP.cCP.i.i.i.iP.voidP}}
\pysigstartsignatures
\pysigstartmultiline
\pysiglinewithargsret{\phantomsection\label{\detokenize{Library_scatter:library_8h_1adb2aa14375aa277b1c32ad17a5f4a39f}}\DUrole{kt}{void}\DUrole{w}{ }\sphinxbfcode{\sphinxupquote{\DUrole{n}{lammps\_gather\_atoms\_subset}}}}{\DUrole{kt}{void}\DUrole{w}{ }\DUrole{p}{*}\DUrole{n,sig-param}{handle}\sphinxparamcomma \DUrole{k}{const}\DUrole{w}{ }\DUrole{kt}{char}\DUrole{w}{ }\DUrole{p}{*}\DUrole{n,sig-param}{name}\sphinxparamcomma \DUrole{kt}{int}\DUrole{w}{ }\DUrole{n,sig-param}{type}\sphinxparamcomma \DUrole{kt}{int}\DUrole{w}{ }\DUrole{n,sig-param}{count}\sphinxparamcomma \DUrole{kt}{int}\DUrole{w}{ }\DUrole{n,sig-param}{ndata}\sphinxparamcomma \DUrole{kt}{int}\DUrole{w}{ }\DUrole{p}{*}\DUrole{n,sig-param}{ids}\sphinxparamcomma \DUrole{kt}{void}\DUrole{w}{ }\DUrole{p}{*}\DUrole{n,sig-param}{data}}{}
\pysigstopmultiline
\pysigstopsignatures
\sphinxAtStartPar
Gather the named atom\sphinxhyphen{}based entity for a subset of atoms.

\sphinxAtStartPar

\sphinxAtStartPar
This subroutine gathers data for the requested atom IDs and stores them in a
one\sphinxhyphen{}dimensional array allocated by the user. The data will be ordered by atom
ID, but there is no requirement that the IDs be consecutive. If you wish to
return a similar array for \sphinxstyleemphasis{all} the atoms, use {\hyperref[\detokenize{Library_scatter:_CPPv419lammps_gather_atomsPvPKciiPv}]{\sphinxcrossref{\sphinxcode{\sphinxupquote{lammps\_gather\_atoms()}}}}}
or {\hyperref[\detokenize{Library_scatter:_CPPv426lammps_gather_atoms_concatPvPKciiPv}]{\sphinxcrossref{\sphinxcode{\sphinxupquote{lammps\_gather\_atoms\_concat()}}}}}.

\sphinxAtStartPar
The \sphinxstyleemphasis{data} array will be in groups of \sphinxstyleemphasis{count} values, sorted by atom ID
in the same order as the array \sphinxstyleemphasis{ids} (e.g., if \sphinxstyleemphasis{name} is \sphinxstyleemphasis{x}, \sphinxstyleemphasis{count} = 3, and
\sphinxstyleemphasis{ids} is \{100, 57, 210\}, then \sphinxstyleemphasis{data} might look like \{x{[}100{]}{[}0{]}, x{[}100{]}{[}1{]},
x{[}100{]}{[}2{]}, x{[}57{]}{[}0{]}, x{[}57{]}{[}1{]}, x{[}57{]}{[}2{]}, x{[}210{]}{[}0{]}, \(\dots\));
\sphinxstyleemphasis{ids} must be provided by the user with length \sphinxstyleemphasis{ndata}, and
\sphinxstyleemphasis{data} must be pre\sphinxhyphen{}allocated by the caller to length
(\sphinxstyleemphasis{count} \(\times\) \sphinxstyleemphasis{ndata}).

\begin{sphinxadmonition}{note}{Restrictions}

\sphinxAtStartPar
This function is not compatible with \sphinxcode{\sphinxupquote{\sphinxhyphen{}DLAMMPS\_BIGBIG}}.

\sphinxAtStartPar
Atom IDs must be defined and an {\hyperref[\detokenize{atom_modify::doc}]{\sphinxcrossref{\DUrole{doc}{atom map must be enabled}}}}

\sphinxAtStartPar
The total number of atoms must not be more than 2147483647 (max 32\sphinxhyphen{}bit signed int).
\end{sphinxadmonition}

\begin{quote}\begin{description}
\sphinxlineitem{Parameters}\begin{itemize}
\item {} 
\sphinxAtStartPar
\sphinxstyleliteralstrong{\sphinxupquote{handle}} \textendash{} pointer to a previously created LAMMPS instance 

\item {} 
\sphinxAtStartPar
\sphinxstyleliteralstrong{\sphinxupquote{name}} \textendash{} desired quantity (e.g., \sphinxstyleemphasis{x} or \sphinxstyleemphasis{charge}) 

\item {} 
\sphinxAtStartPar
\sphinxstyleliteralstrong{\sphinxupquote{type}} \textendash{} 0 for \sphinxcode{\sphinxupquote{int}} values, 1 for \sphinxcode{\sphinxupquote{double}} values 

\item {} 
\sphinxAtStartPar
\sphinxstyleliteralstrong{\sphinxupquote{count}} \textendash{} number of per\sphinxhyphen{}atom values (e.g., 1 for \sphinxstyleemphasis{type} or \sphinxstyleemphasis{charge}, 3 for \sphinxstyleemphasis{x} or \sphinxstyleemphasis{f}); use \sphinxstyleemphasis{count} = 3 with “image” if you want single image flags unpacked into (\sphinxstyleemphasis{x},\sphinxstyleemphasis{y},\sphinxstyleemphasis{z}) 

\item {} 
\sphinxAtStartPar
\sphinxstyleliteralstrong{\sphinxupquote{ndata}} \textendash{} number of atoms for which to return data (can be all of them) 

\item {} 
\sphinxAtStartPar
\sphinxstyleliteralstrong{\sphinxupquote{ids}} \textendash{} list of \sphinxstyleemphasis{ndata} atom IDs for which to return data 

\item {} 
\sphinxAtStartPar
\sphinxstyleliteralstrong{\sphinxupquote{data}} \textendash{} per\sphinxhyphen{}atom values packed in a 1\sphinxhyphen{}dimensional array of length \sphinxstyleemphasis{ndata} * \sphinxstyleemphasis{count}. 

\end{itemize}

\end{description}\end{quote}

\end{fulllineitems}



\bigskip\hrule\bigskip

\index{lammps\_scatter\_atoms (C++ function)@\spxentry{lammps\_scatter\_atoms}\spxextra{C++ function}}

\begin{fulllineitems}
\phantomsection\label{\detokenize{Library_scatter:_CPPv420lammps_scatter_atomsPvPKciiPv}}\phantomsection\label{\detokenize{Library_scatter:_CPPv320lammps_scatter_atomsPvPKciiPv}}\phantomsection\label{\detokenize{Library_scatter:_CPPv220lammps_scatter_atomsPvPKciiPv}}\phantomsection\label{\detokenize{Library_scatter:lammps_scatter_atoms__voidP.cCP.i.i.voidP}}
\pysigstartsignatures
\pysigstartmultiline
\pysiglinewithargsret{\phantomsection\label{\detokenize{Library_scatter:library_8h_1a912e450e211db31d9396ad5f421d500d}}\DUrole{kt}{void}\DUrole{w}{ }\sphinxbfcode{\sphinxupquote{\DUrole{n}{lammps\_scatter\_atoms}}}}{\DUrole{kt}{void}\DUrole{w}{ }\DUrole{p}{*}\DUrole{n,sig-param}{handle}\sphinxparamcomma \DUrole{k}{const}\DUrole{w}{ }\DUrole{kt}{char}\DUrole{w}{ }\DUrole{p}{*}\DUrole{n,sig-param}{name}\sphinxparamcomma \DUrole{kt}{int}\DUrole{w}{ }\DUrole{n,sig-param}{type}\sphinxparamcomma \DUrole{kt}{int}\DUrole{w}{ }\DUrole{n,sig-param}{count}\sphinxparamcomma \DUrole{kt}{void}\DUrole{w}{ }\DUrole{p}{*}\DUrole{n,sig-param}{data}}{}
\pysigstopmultiline
\pysigstopsignatures
\sphinxAtStartPar
Scatter the named atom\sphinxhyphen{}based entities in \sphinxstyleemphasis{data} to all processes.

\sphinxAtStartPar

\sphinxAtStartPar
This subroutine takes data stored in a one\sphinxhyphen{}dimensional array supplied by the
user and scatters them to all atoms on all processes. The data must be
ordered by atom ID, with the requirement that the IDs be consecutive.
Use {\hyperref[\detokenize{Library_scatter:_CPPv427lammps_scatter_atoms_subsetPvPKciiiPiPv}]{\sphinxcrossref{\sphinxcode{\sphinxupquote{lammps\_scatter\_atoms\_subset()}}}}} to scatter data for some (or all)
atoms, unordered.

\sphinxAtStartPar
The \sphinxstyleemphasis{data} array needs to be ordered in groups of \sphinxstyleemphasis{count} values, sorted by
atom ID (e.g., if \sphinxstyleemphasis{name} is \sphinxstyleemphasis{x} and \sphinxstyleemphasis{count} = 3, then
\sphinxstyleemphasis{data} = \{x{[}0{]}{[}0{]}, x{[}0{]}{[}1{]}, x{[}0{]}{[}2{]}, x{[}1{]}{[}0{]}, x{[}1{]}{[}1{]}, x{[}1{]}{[}2{]}, x{[}2{]}{[}0{]},
\(\dots\)\}); \sphinxstyleemphasis{data} must be of length (\sphinxstyleemphasis{count} \(\times\) \sphinxstyleemphasis{natoms}).

\begin{sphinxadmonition}{note}{Restrictions}

\sphinxAtStartPar
This function is not compatible with \sphinxcode{\sphinxupquote{\sphinxhyphen{}DLAMMPS\_BIGBIG}}.

\sphinxAtStartPar
Atom IDs must be defined, must be consecutive, and an
{\hyperref[\detokenize{atom_modify::doc}]{\sphinxcrossref{\DUrole{doc}{atom map must be enabled}}}}

\sphinxAtStartPar
The total number of atoms must not be more than 2147483647 (max 32\sphinxhyphen{}bit signed int).
\end{sphinxadmonition}

\begin{quote}\begin{description}
\sphinxlineitem{Parameters}\begin{itemize}
\item {} 
\sphinxAtStartPar
\sphinxstyleliteralstrong{\sphinxupquote{handle}} \textendash{} pointer to a previously created LAMMPS instance 

\item {} 
\sphinxAtStartPar
\sphinxstyleliteralstrong{\sphinxupquote{name}} \textendash{} desired quantity (e.g., \sphinxstyleemphasis{x} or \sphinxstyleemphasis{charge}) 

\item {} 
\sphinxAtStartPar
\sphinxstyleliteralstrong{\sphinxupquote{type}} \textendash{} 0 for \sphinxcode{\sphinxupquote{int}} values, 1 for \sphinxcode{\sphinxupquote{double}} values 

\item {} 
\sphinxAtStartPar
\sphinxstyleliteralstrong{\sphinxupquote{count}} \textendash{} number of per\sphinxhyphen{}atom values (e.g., 1 for \sphinxstyleemphasis{type} or \sphinxstyleemphasis{charge}, 3 for \sphinxstyleemphasis{x} or \sphinxstyleemphasis{f}); use \sphinxstyleemphasis{count} = 3 with \sphinxstyleemphasis{image} if you have a single image flag packed into (\sphinxstyleemphasis{x},\sphinxstyleemphasis{y},\sphinxstyleemphasis{z}) components. 

\item {} 
\sphinxAtStartPar
\sphinxstyleliteralstrong{\sphinxupquote{data}} \textendash{} per\sphinxhyphen{}atom values packed in a one\sphinxhyphen{}dimensional array of length \sphinxstyleemphasis{natoms} * \sphinxstyleemphasis{count}. 

\end{itemize}

\end{description}\end{quote}

\end{fulllineitems}



\bigskip\hrule\bigskip

\index{lammps\_scatter\_atoms\_subset (C++ function)@\spxentry{lammps\_scatter\_atoms\_subset}\spxextra{C++ function}}

\begin{fulllineitems}
\phantomsection\label{\detokenize{Library_scatter:_CPPv427lammps_scatter_atoms_subsetPvPKciiiPiPv}}\phantomsection\label{\detokenize{Library_scatter:_CPPv327lammps_scatter_atoms_subsetPvPKciiiPiPv}}\phantomsection\label{\detokenize{Library_scatter:_CPPv227lammps_scatter_atoms_subsetPvPKciiiPiPv}}\phantomsection\label{\detokenize{Library_scatter:lammps_scatter_atoms_subset__voidP.cCP.i.i.i.iP.voidP}}
\pysigstartsignatures
\pysigstartmultiline
\pysiglinewithargsret{\phantomsection\label{\detokenize{Library_scatter:library_8h_1a42726943d284b0edad9b8f3cd3129398}}\DUrole{kt}{void}\DUrole{w}{ }\sphinxbfcode{\sphinxupquote{\DUrole{n}{lammps\_scatter\_atoms\_subset}}}}{\DUrole{kt}{void}\DUrole{w}{ }\DUrole{p}{*}\DUrole{n,sig-param}{handle}\sphinxparamcomma \DUrole{k}{const}\DUrole{w}{ }\DUrole{kt}{char}\DUrole{w}{ }\DUrole{p}{*}\DUrole{n,sig-param}{name}\sphinxparamcomma \DUrole{kt}{int}\DUrole{w}{ }\DUrole{n,sig-param}{type}\sphinxparamcomma \DUrole{kt}{int}\DUrole{w}{ }\DUrole{n,sig-param}{count}\sphinxparamcomma \DUrole{kt}{int}\DUrole{w}{ }\DUrole{n,sig-param}{ndata}\sphinxparamcomma \DUrole{kt}{int}\DUrole{w}{ }\DUrole{p}{*}\DUrole{n,sig-param}{ids}\sphinxparamcomma \DUrole{kt}{void}\DUrole{w}{ }\DUrole{p}{*}\DUrole{n,sig-param}{data}}{}
\pysigstopmultiline
\pysigstopsignatures
\sphinxAtStartPar
Scatter the named atom\sphinxhyphen{}based entities in \sphinxstyleemphasis{data} from a subset of atoms to all processes.

\sphinxAtStartPar

\sphinxAtStartPar
This subroutine takes data stored in a one\sphinxhyphen{}dimensional array supplied by the
user and scatters them to a subset of atoms on all processes. The array
\sphinxstyleemphasis{data} contains data associated with atom IDs, but there is no requirement that
the IDs be consecutive, as they are provided in a separate array.
Use {\hyperref[\detokenize{Library_scatter:_CPPv420lammps_scatter_atomsPvPKciiPv}]{\sphinxcrossref{\sphinxcode{\sphinxupquote{lammps\_scatter\_atoms()}}}}} to scatter data for all atoms, in order.

\sphinxAtStartPar
The \sphinxstyleemphasis{data} array needs to be organized in groups of \sphinxstyleemphasis{count} values, with the
groups in the same order as the array \sphinxstyleemphasis{ids}. For example, if you want \sphinxstyleemphasis{data}
to be the array \{x{[}1{]}{[}0{]}, x{[}1{]}{[}1{]}, x{[}1{]}{[}2{]}, x{[}100{]}{[}0{]}, x{[}100{]}{[}1{]}, x{[}100{]}{[}2{]},
x{[}57{]}{[}0{]}, x{[}57{]}{[}1{]}, x{[}57{]}{[}2{]}\}, then \sphinxstyleemphasis{count} = 3, \sphinxstyleemphasis{ndata} = 3, and \sphinxstyleemphasis{ids} would
be \{1, 100, 57\}.

\begin{sphinxadmonition}{note}{Restrictions}

\sphinxAtStartPar
This function is not compatible with \sphinxcode{\sphinxupquote{\sphinxhyphen{}DLAMMPS\_BIGBIG}}.

\sphinxAtStartPar
Atom IDs must be defined and an {\hyperref[\detokenize{atom_modify::doc}]{\sphinxcrossref{\DUrole{doc}{atom map must be enabled}}}}

\sphinxAtStartPar
The total number of atoms must not be more than 2147483647 (max 32\sphinxhyphen{}bit signed int).
\end{sphinxadmonition}

\begin{quote}\begin{description}
\sphinxlineitem{Parameters}\begin{itemize}
\item {} 
\sphinxAtStartPar
\sphinxstyleliteralstrong{\sphinxupquote{handle}} \textendash{} pointer to a previously created LAMMPS instance 

\item {} 
\sphinxAtStartPar
\sphinxstyleliteralstrong{\sphinxupquote{name}} \textendash{} desired quantity (e.g., \sphinxstyleemphasis{x} or \sphinxstyleemphasis{charge}) 

\item {} 
\sphinxAtStartPar
\sphinxstyleliteralstrong{\sphinxupquote{type}} \textendash{} 0 for \sphinxcode{\sphinxupquote{int}} values, 1 for \sphinxcode{\sphinxupquote{double}} values 

\item {} 
\sphinxAtStartPar
\sphinxstyleliteralstrong{\sphinxupquote{count}} \textendash{} number of per\sphinxhyphen{}atom values (e.g., 1 for \sphinxstyleemphasis{type} or \sphinxstyleemphasis{charge}, 3 for \sphinxstyleemphasis{x} or \sphinxstyleemphasis{f}); use \sphinxstyleemphasis{count} = 3 with “image” if you have all the image flags packed into (\sphinxstyleemphasis{xyz}) 

\item {} 
\sphinxAtStartPar
\sphinxstyleliteralstrong{\sphinxupquote{ndata}} \textendash{} number of atoms listed in \sphinxstyleemphasis{ids} and \sphinxstyleemphasis{data} arrays 

\item {} 
\sphinxAtStartPar
\sphinxstyleliteralstrong{\sphinxupquote{ids}} \textendash{} list of \sphinxstyleemphasis{ndata} atom IDs to scatter data to 

\item {} 
\sphinxAtStartPar
\sphinxstyleliteralstrong{\sphinxupquote{data}} \textendash{} per\sphinxhyphen{}atom values packed in a 1\sphinxhyphen{}dimensional array of length \sphinxstyleemphasis{ndata} * \sphinxstyleemphasis{count}. 

\end{itemize}

\end{description}\end{quote}

\end{fulllineitems}



\bigskip\hrule\bigskip

\index{lammps\_gather\_bonds (C++ function)@\spxentry{lammps\_gather\_bonds}\spxextra{C++ function}}

\begin{fulllineitems}
\phantomsection\label{\detokenize{Library_scatter:_CPPv419lammps_gather_bondsPvPv}}\phantomsection\label{\detokenize{Library_scatter:_CPPv319lammps_gather_bondsPvPv}}\phantomsection\label{\detokenize{Library_scatter:_CPPv219lammps_gather_bondsPvPv}}\phantomsection\label{\detokenize{Library_scatter:lammps_gather_bonds__voidP.voidP}}
\pysigstartsignatures
\pysigstartmultiline
\pysiglinewithargsret{\phantomsection\label{\detokenize{Library_scatter:library_8h_1a360d274b56e28d7a9fd0e195ff2171af}}\DUrole{kt}{void}\DUrole{w}{ }\sphinxbfcode{\sphinxupquote{\DUrole{n}{lammps\_gather\_bonds}}}}{\DUrole{kt}{void}\DUrole{w}{ }\DUrole{p}{*}\DUrole{n,sig-param}{handle}\sphinxparamcomma \DUrole{kt}{void}\DUrole{w}{ }\DUrole{p}{*}\DUrole{n,sig-param}{data}}{}
\pysigstopmultiline
\pysigstopsignatures
\sphinxAtStartPar
Gather type and constituent atom info for all bonds

\sphinxAtStartPar

\sphinxAtStartPar
\DUrole{versionmodified,added}{Added in version 28Jul2021.}

\sphinxAtStartPar
This function copies the list of all bonds into a buffer provided by
the calling code. The buffer will be filled with bond type, bond atom 1,
bond atom 2 for each bond. Thus the buffer has to be allocated to the
dimension of 3 times the \sphinxstylestrong{total} number of bonds times the size of
the LAMMPS “tagint” type, which is either 4 or 8 bytes depending on
whether they are stored in 32\sphinxhyphen{}bit or 64\sphinxhyphen{}bit integers, respectively.
This size depends on the compile time settings used when compiling
the LAMMPS library and can be queried by calling
{\hyperref[\detokenize{Library_properties:_CPPv422lammps_extract_settingPvPKc}]{\sphinxcrossref{\sphinxcode{\sphinxupquote{lammps\_extract\_setting()}}}}} with the keyword “tagint”.

\sphinxAtStartPar
When running in parallel, the data buffer must be allocated on \sphinxstylestrong{all}
MPI ranks and will be filled with the information for \sphinxstylestrong{all} bonds
in the system.

\sphinxAtStartPar
Below is a brief C code demonstrating accessing this collected bond information.

\begin{sphinxVerbatim}[commandchars=\\\{\}]
\PYG{c+cp}{\PYGZsh{}}\PYG{c+cp}{include}\PYG{+w}{ }\PYG{c+cpf}{\PYGZdq{}library.h\PYGZdq{}}

\PYG{c+cp}{\PYGZsh{}}\PYG{c+cp}{include}\PYG{+w}{ }\PYG{c+cpf}{\PYGZlt{}stdint.h\PYGZgt{}}
\PYG{c+cp}{\PYGZsh{}}\PYG{c+cp}{include}\PYG{+w}{ }\PYG{c+cpf}{\PYGZlt{}stdio.h\PYGZgt{}}
\PYG{c+cp}{\PYGZsh{}}\PYG{c+cp}{include}\PYG{+w}{ }\PYG{c+cpf}{\PYGZlt{}stdlib.h\PYGZgt{}}

\PYG{k+kt}{int}\PYG{+w}{ }\PYG{n+nf}{main}\PYG{p}{(}\PYG{k+kt}{int}\PYG{+w}{ }\PYG{n}{argc}\PYG{p}{,}\PYG{+w}{ }\PYG{k+kt}{char}\PYG{+w}{ }\PYG{o}{*}\PYG{o}{*}\PYG{n}{argv}\PYG{p}{)}
\PYG{p}{\PYGZob{}}
\PYG{+w}{    }\PYG{k+kt}{int}\PYG{+w}{ }\PYG{n}{tagintsize}\PYG{p}{;}
\PYG{+w}{    }\PYG{k+kt}{int64\PYGZus{}t}\PYG{+w}{ }\PYG{n}{i}\PYG{p}{,}\PYG{+w}{ }\PYG{n}{nbonds}\PYG{p}{;}
\PYG{+w}{    }\PYG{k+kt}{void}\PYG{+w}{ }\PYG{o}{*}\PYG{n}{handle}\PYG{p}{,}\PYG{+w}{ }\PYG{o}{*}\PYG{n}{bonds}\PYG{p}{;}

\PYG{+w}{    }\PYG{n}{handle}\PYG{+w}{ }\PYG{o}{=}\PYG{+w}{ }\PYG{n}{lammps\PYGZus{}open\PYGZus{}no\PYGZus{}mpi}\PYG{p}{(}\PYG{l+m+mi}{0}\PYG{p}{,}\PYG{+w}{ }\PYG{n+nb}{NULL}\PYG{p}{,}\PYG{+w}{ }\PYG{n+nb}{NULL}\PYG{p}{)}\PYG{p}{;}
\PYG{+w}{    }\PYG{n}{lammps\PYGZus{}file}\PYG{p}{(}\PYG{n}{handle}\PYG{p}{,}\PYG{+w}{ }\PYG{l+s}{\PYGZdq{}}\PYG{l+s}{in.some\PYGZus{}input}\PYG{l+s}{\PYGZdq{}}\PYG{p}{)}\PYG{p}{;}

\PYG{+w}{    }\PYG{n}{tagintsize}\PYG{+w}{ }\PYG{o}{=}\PYG{+w}{ }\PYG{n}{lammps\PYGZus{}extract\PYGZus{}setting}\PYG{p}{(}\PYG{n}{handle}\PYG{p}{,}\PYG{+w}{ }\PYG{l+s}{\PYGZdq{}}\PYG{l+s}{tagint}\PYG{l+s}{\PYGZdq{}}\PYG{p}{)}\PYG{p}{;}
\PYG{+w}{    }\PYG{k}{if}\PYG{+w}{ }\PYG{p}{(}\PYG{n}{tagintsize}\PYG{+w}{ }\PYG{o}{=}\PYG{o}{=}\PYG{+w}{ }\PYG{l+m+mi}{4}\PYG{p}{)}
\PYG{+w}{        }\PYG{n}{nbonds}\PYG{+w}{ }\PYG{o}{=}\PYG{+w}{ }\PYG{o}{*}\PYG{p}{(}\PYG{k+kt}{int32\PYGZus{}t}\PYG{+w}{ }\PYG{o}{*}\PYG{p}{)}\PYG{n}{lammps\PYGZus{}extract\PYGZus{}global}\PYG{p}{(}\PYG{n}{handle}\PYG{p}{,}\PYG{+w}{ }\PYG{l+s}{\PYGZdq{}}\PYG{l+s}{nbonds}\PYG{l+s}{\PYGZdq{}}\PYG{p}{)}\PYG{p}{;}
\PYG{+w}{     }\PYG{k}{else}
\PYG{+w}{        }\PYG{n}{nbonds}\PYG{+w}{ }\PYG{o}{=}\PYG{+w}{ }\PYG{o}{*}\PYG{p}{(}\PYG{k+kt}{int64\PYGZus{}t}\PYG{+w}{ }\PYG{o}{*}\PYG{p}{)}\PYG{n}{lammps\PYGZus{}extract\PYGZus{}global}\PYG{p}{(}\PYG{n}{handle}\PYG{p}{,}\PYG{+w}{ }\PYG{l+s}{\PYGZdq{}}\PYG{l+s}{nbonds}\PYG{l+s}{\PYGZdq{}}\PYG{p}{)}\PYG{p}{;}
\PYG{+w}{    }\PYG{n}{bonds}\PYG{+w}{ }\PYG{o}{=}\PYG{+w}{ }\PYG{n}{malloc}\PYG{p}{(}\PYG{n}{nbonds}\PYG{+w}{ }\PYG{o}{*}\PYG{+w}{ }\PYG{l+m+mi}{3}\PYG{+w}{ }\PYG{o}{*}\PYG{+w}{ }\PYG{n}{tagintsize}\PYG{p}{)}\PYG{p}{;}

\PYG{+w}{    }\PYG{n}{lammps\PYGZus{}gather\PYGZus{}bonds}\PYG{p}{(}\PYG{n}{handle}\PYG{p}{,}\PYG{+w}{ }\PYG{n}{bonds}\PYG{p}{)}\PYG{p}{;}

\PYG{+w}{    }\PYG{k}{if}\PYG{+w}{ }\PYG{p}{(}\PYG{n}{lammps\PYGZus{}extract\PYGZus{}setting}\PYG{p}{(}\PYG{n}{handle}\PYG{p}{,}\PYG{+w}{ }\PYG{l+s}{\PYGZdq{}}\PYG{l+s}{world\PYGZus{}rank}\PYG{l+s}{\PYGZdq{}}\PYG{p}{)}\PYG{+w}{ }\PYG{o}{=}\PYG{o}{=}\PYG{+w}{ }\PYG{l+m+mi}{0}\PYG{p}{)}\PYG{+w}{ }\PYG{p}{\PYGZob{}}
\PYG{+w}{        }\PYG{k}{if}\PYG{+w}{ }\PYG{p}{(}\PYG{n}{tagintsize}\PYG{+w}{ }\PYG{o}{=}\PYG{o}{=}\PYG{+w}{ }\PYG{l+m+mi}{4}\PYG{p}{)}\PYG{+w}{ }\PYG{p}{\PYGZob{}}
\PYG{+w}{            }\PYG{k+kt}{int32\PYGZus{}t}\PYG{+w}{ }\PYG{o}{*}\PYG{n}{bonds\PYGZus{}real}\PYG{+w}{ }\PYG{o}{=}\PYG{+w}{ }\PYG{p}{(}\PYG{k+kt}{int32\PYGZus{}t}\PYG{+w}{ }\PYG{o}{*}\PYG{p}{)}\PYG{n}{bonds}\PYG{p}{;}
\PYG{+w}{            }\PYG{k}{for}\PYG{+w}{ }\PYG{p}{(}\PYG{n}{i}\PYG{+w}{ }\PYG{o}{=}\PYG{+w}{ }\PYG{l+m+mi}{0}\PYG{p}{;}\PYG{+w}{ }\PYG{n}{i}\PYG{+w}{ }\PYG{o}{\PYGZlt{}}\PYG{+w}{ }\PYG{n}{nbonds}\PYG{p}{;}\PYG{+w}{ }\PYG{o}{+}\PYG{o}{+}\PYG{n}{i}\PYG{p}{)}\PYG{+w}{ }\PYG{p}{\PYGZob{}}
\PYG{+w}{                }\PYG{n}{printf}\PYG{p}{(}\PYG{l+s}{\PYGZdq{}}\PYG{l+s}{bond \PYGZpc{} 4ld: type = \PYGZpc{}d, atoms: \PYGZpc{} 4d  \PYGZpc{} 4d}\PYG{l+s+se}{\PYGZbs{}n}\PYG{l+s}{\PYGZdq{}}\PYG{p}{,}\PYG{n}{i}\PYG{p}{,}
\PYG{+w}{                       }\PYG{n}{bonds\PYGZus{}real}\PYG{p}{[}\PYG{l+m+mi}{3}\PYG{o}{*}\PYG{n}{i}\PYG{p}{]}\PYG{p}{,}\PYG{+w}{ }\PYG{n}{bonds\PYGZus{}real}\PYG{p}{[}\PYG{l+m+mi}{3}\PYG{o}{*}\PYG{n}{i}\PYG{o}{+}\PYG{l+m+mi}{1}\PYG{p}{]}\PYG{p}{,}\PYG{+w}{ }\PYG{n}{bonds\PYGZus{}real}\PYG{p}{[}\PYG{l+m+mi}{3}\PYG{o}{*}\PYG{n}{i}\PYG{o}{+}\PYG{l+m+mi}{2}\PYG{p}{]}\PYG{p}{)}\PYG{p}{;}
\PYG{+w}{            }\PYG{p}{\PYGZcb{}}
\PYG{+w}{        }\PYG{p}{\PYGZcb{}}\PYG{+w}{ }\PYG{k}{else}\PYG{+w}{ }\PYG{p}{\PYGZob{}}
\PYG{+w}{            }\PYG{k+kt}{int64\PYGZus{}t}\PYG{+w}{ }\PYG{o}{*}\PYG{n}{bonds\PYGZus{}real}\PYG{+w}{ }\PYG{o}{=}\PYG{+w}{ }\PYG{p}{(}\PYG{k+kt}{int64\PYGZus{}t}\PYG{+w}{ }\PYG{o}{*}\PYG{p}{)}\PYG{n}{bonds}\PYG{p}{;}
\PYG{+w}{            }\PYG{k}{for}\PYG{+w}{ }\PYG{p}{(}\PYG{n}{i}\PYG{+w}{ }\PYG{o}{=}\PYG{+w}{ }\PYG{l+m+mi}{0}\PYG{p}{;}\PYG{+w}{ }\PYG{n}{i}\PYG{+w}{ }\PYG{o}{\PYGZlt{}}\PYG{+w}{ }\PYG{n}{nbonds}\PYG{p}{;}\PYG{+w}{ }\PYG{o}{+}\PYG{o}{+}\PYG{n}{i}\PYG{p}{)}\PYG{+w}{ }\PYG{p}{\PYGZob{}}
\PYG{+w}{                }\PYG{n}{printf}\PYG{p}{(}\PYG{l+s}{\PYGZdq{}}\PYG{l+s}{bond \PYGZpc{} 4ld: type = \PYGZpc{}ld, atoms: \PYGZpc{} 4ld  \PYGZpc{} 4ld}\PYG{l+s+se}{\PYGZbs{}n}\PYG{l+s}{\PYGZdq{}}\PYG{p}{,}\PYG{n}{i}\PYG{p}{,}
\PYG{+w}{                       }\PYG{n}{bonds\PYGZus{}real}\PYG{p}{[}\PYG{l+m+mi}{3}\PYG{o}{*}\PYG{n}{i}\PYG{p}{]}\PYG{p}{,}\PYG{+w}{ }\PYG{n}{bonds\PYGZus{}real}\PYG{p}{[}\PYG{l+m+mi}{3}\PYG{o}{*}\PYG{n}{i}\PYG{o}{+}\PYG{l+m+mi}{1}\PYG{p}{]}\PYG{p}{,}\PYG{+w}{ }\PYG{n}{bonds\PYGZus{}real}\PYG{p}{[}\PYG{l+m+mi}{3}\PYG{o}{*}\PYG{n}{i}\PYG{o}{+}\PYG{l+m+mi}{2}\PYG{p}{]}\PYG{p}{)}\PYG{p}{;}
\PYG{+w}{            }\PYG{p}{\PYGZcb{}}
\PYG{+w}{        }\PYG{p}{\PYGZcb{}}
\PYG{+w}{    }\PYG{p}{\PYGZcb{}}

\PYG{+w}{    }\PYG{n}{lammps\PYGZus{}close}\PYG{p}{(}\PYG{n}{handle}\PYG{p}{)}\PYG{p}{;}
\PYG{+w}{    }\PYG{n}{lammps\PYGZus{}mpi\PYGZus{}finalize}\PYG{p}{(}\PYG{p}{)}\PYG{p}{;}
\PYG{+w}{    }\PYG{n}{free}\PYG{p}{(}\PYG{n}{bonds}\PYG{p}{)}\PYG{p}{;}
\PYG{+w}{    }\PYG{k}{return}\PYG{+w}{ }\PYG{l+m+mi}{0}\PYG{p}{;}
\PYG{p}{\PYGZcb{}}
\end{sphinxVerbatim}

\begin{quote}\begin{description}
\sphinxlineitem{Parameters}\begin{itemize}
\item {} 
\sphinxAtStartPar
\sphinxstyleliteralstrong{\sphinxupquote{handle}} \textendash{} pointer to a previously created LAMMPS instance 

\item {} 
\sphinxAtStartPar
\sphinxstyleliteralstrong{\sphinxupquote{data}} \textendash{} pointer to data to copy the result to 

\end{itemize}

\end{description}\end{quote}

\end{fulllineitems}



\bigskip\hrule\bigskip

\index{lammps\_gather\_angles (C++ function)@\spxentry{lammps\_gather\_angles}\spxextra{C++ function}}

\begin{fulllineitems}
\phantomsection\label{\detokenize{Library_scatter:_CPPv420lammps_gather_anglesPvPv}}\phantomsection\label{\detokenize{Library_scatter:_CPPv320lammps_gather_anglesPvPv}}\phantomsection\label{\detokenize{Library_scatter:_CPPv220lammps_gather_anglesPvPv}}\phantomsection\label{\detokenize{Library_scatter:lammps_gather_angles__voidP.voidP}}
\pysigstartsignatures
\pysigstartmultiline
\pysiglinewithargsret{\phantomsection\label{\detokenize{Library_scatter:library_8h_1a6f8c39008d113b5c3450c2cffaadcf3c}}\DUrole{kt}{void}\DUrole{w}{ }\sphinxbfcode{\sphinxupquote{\DUrole{n}{lammps\_gather\_angles}}}}{\DUrole{kt}{void}\DUrole{w}{ }\DUrole{p}{*}\DUrole{n,sig-param}{handle}\sphinxparamcomma \DUrole{kt}{void}\DUrole{w}{ }\DUrole{p}{*}\DUrole{n,sig-param}{data}}{}
\pysigstopmultiline
\pysigstopsignatures
\sphinxAtStartPar
Gather type and constituent atom info for all angles

\sphinxAtStartPar

\sphinxAtStartPar
\DUrole{versionmodified,added}{Added in version 8Feb2023.}

\sphinxAtStartPar
This function copies the list of all angles into a buffer provided by
the calling code. The buffer will be filled with angle type, angle atom 1,
angle atom 2, angle atom 3 for each angle. Thus the buffer has to be allocated to the
dimension of 4 times the \sphinxstylestrong{total} number of angles times the size of
the LAMMPS “tagint” type, which is either 4 or 8 bytes depending on
whether they are stored in 32\sphinxhyphen{}bit or 64\sphinxhyphen{}bit integers, respectively.
This size depends on the compile time settings used when compiling
the LAMMPS library and can be queried by calling
{\hyperref[\detokenize{Library_properties:_CPPv422lammps_extract_settingPvPKc}]{\sphinxcrossref{\sphinxcode{\sphinxupquote{lammps\_extract\_setting()}}}}} with the keyword “tagint”.

\sphinxAtStartPar
When running in parallel, the data buffer must be allocated on \sphinxstylestrong{all}
MPI ranks and will be filled with the information for \sphinxstylestrong{all} angles
in the system.

\sphinxAtStartPar
Below is a brief C code demonstrating accessing this collected angle information.

\begin{sphinxVerbatim}[commandchars=\\\{\}]
\PYG{c+cp}{\PYGZsh{}}\PYG{c+cp}{include}\PYG{+w}{ }\PYG{c+cpf}{\PYGZdq{}library.h\PYGZdq{}}

\PYG{c+cp}{\PYGZsh{}}\PYG{c+cp}{include}\PYG{+w}{ }\PYG{c+cpf}{\PYGZlt{}stdint.h\PYGZgt{}}
\PYG{c+cp}{\PYGZsh{}}\PYG{c+cp}{include}\PYG{+w}{ }\PYG{c+cpf}{\PYGZlt{}stdio.h\PYGZgt{}}
\PYG{c+cp}{\PYGZsh{}}\PYG{c+cp}{include}\PYG{+w}{ }\PYG{c+cpf}{\PYGZlt{}stdlib.h\PYGZgt{}}

\PYG{k+kt}{int}\PYG{+w}{ }\PYG{n+nf}{main}\PYG{p}{(}\PYG{k+kt}{int}\PYG{+w}{ }\PYG{n}{argc}\PYG{p}{,}\PYG{+w}{ }\PYG{k+kt}{char}\PYG{+w}{ }\PYG{o}{*}\PYG{o}{*}\PYG{n}{argv}\PYG{p}{)}
\PYG{p}{\PYGZob{}}
\PYG{+w}{    }\PYG{k+kt}{int}\PYG{+w}{ }\PYG{n}{tagintsize}\PYG{p}{;}
\PYG{+w}{    }\PYG{k+kt}{int64\PYGZus{}t}\PYG{+w}{ }\PYG{n}{i}\PYG{p}{,}\PYG{+w}{ }\PYG{n}{nangles}\PYG{p}{;}
\PYG{+w}{    }\PYG{k+kt}{void}\PYG{+w}{ }\PYG{o}{*}\PYG{n}{handle}\PYG{p}{,}\PYG{+w}{ }\PYG{o}{*}\PYG{n}{angles}\PYG{p}{;}

\PYG{+w}{    }\PYG{n}{handle}\PYG{+w}{ }\PYG{o}{=}\PYG{+w}{ }\PYG{n}{lammps\PYGZus{}open\PYGZus{}no\PYGZus{}mpi}\PYG{p}{(}\PYG{l+m+mi}{0}\PYG{p}{,}\PYG{+w}{ }\PYG{n+nb}{NULL}\PYG{p}{,}\PYG{+w}{ }\PYG{n+nb}{NULL}\PYG{p}{)}\PYG{p}{;}
\PYG{+w}{    }\PYG{n}{lammps\PYGZus{}file}\PYG{p}{(}\PYG{n}{handle}\PYG{p}{,}\PYG{+w}{ }\PYG{l+s}{\PYGZdq{}}\PYG{l+s}{in.some\PYGZus{}input}\PYG{l+s}{\PYGZdq{}}\PYG{p}{)}\PYG{p}{;}

\PYG{+w}{    }\PYG{n}{tagintsize}\PYG{+w}{ }\PYG{o}{=}\PYG{+w}{ }\PYG{n}{lammps\PYGZus{}extract\PYGZus{}setting}\PYG{p}{(}\PYG{n}{handle}\PYG{p}{,}\PYG{+w}{ }\PYG{l+s}{\PYGZdq{}}\PYG{l+s}{tagint}\PYG{l+s}{\PYGZdq{}}\PYG{p}{)}\PYG{p}{;}
\PYG{+w}{    }\PYG{k}{if}\PYG{+w}{ }\PYG{p}{(}\PYG{n}{tagintsize}\PYG{+w}{ }\PYG{o}{=}\PYG{o}{=}\PYG{+w}{ }\PYG{l+m+mi}{4}\PYG{p}{)}
\PYG{+w}{        }\PYG{n}{nangles}\PYG{+w}{ }\PYG{o}{=}\PYG{+w}{ }\PYG{o}{*}\PYG{p}{(}\PYG{k+kt}{int32\PYGZus{}t}\PYG{+w}{ }\PYG{o}{*}\PYG{p}{)}\PYG{n}{lammps\PYGZus{}extract\PYGZus{}global}\PYG{p}{(}\PYG{n}{handle}\PYG{p}{,}\PYG{+w}{ }\PYG{l+s}{\PYGZdq{}}\PYG{l+s}{nangles}\PYG{l+s}{\PYGZdq{}}\PYG{p}{)}\PYG{p}{;}
\PYG{+w}{     }\PYG{k}{else}
\PYG{+w}{        }\PYG{n}{nangles}\PYG{+w}{ }\PYG{o}{=}\PYG{+w}{ }\PYG{o}{*}\PYG{p}{(}\PYG{k+kt}{int64\PYGZus{}t}\PYG{+w}{ }\PYG{o}{*}\PYG{p}{)}\PYG{n}{lammps\PYGZus{}extract\PYGZus{}global}\PYG{p}{(}\PYG{n}{handle}\PYG{p}{,}\PYG{+w}{ }\PYG{l+s}{\PYGZdq{}}\PYG{l+s}{nangles}\PYG{l+s}{\PYGZdq{}}\PYG{p}{)}\PYG{p}{;}
\PYG{+w}{    }\PYG{n}{angles}\PYG{+w}{ }\PYG{o}{=}\PYG{+w}{ }\PYG{n}{malloc}\PYG{p}{(}\PYG{n}{nangles}\PYG{+w}{ }\PYG{o}{*}\PYG{+w}{ }\PYG{l+m+mi}{4}\PYG{+w}{ }\PYG{o}{*}\PYG{+w}{ }\PYG{n}{tagintsize}\PYG{p}{)}\PYG{p}{;}

\PYG{+w}{    }\PYG{n}{lammps\PYGZus{}gather\PYGZus{}angles}\PYG{p}{(}\PYG{n}{handle}\PYG{p}{,}\PYG{+w}{ }\PYG{n}{angles}\PYG{p}{)}\PYG{p}{;}

\PYG{+w}{    }\PYG{k}{if}\PYG{+w}{ }\PYG{p}{(}\PYG{n}{lammps\PYGZus{}extract\PYGZus{}setting}\PYG{p}{(}\PYG{n}{handle}\PYG{p}{,}\PYG{+w}{ }\PYG{l+s}{\PYGZdq{}}\PYG{l+s}{world\PYGZus{}rank}\PYG{l+s}{\PYGZdq{}}\PYG{p}{)}\PYG{+w}{ }\PYG{o}{=}\PYG{o}{=}\PYG{+w}{ }\PYG{l+m+mi}{0}\PYG{p}{)}\PYG{+w}{ }\PYG{p}{\PYGZob{}}
\PYG{+w}{        }\PYG{k}{if}\PYG{+w}{ }\PYG{p}{(}\PYG{n}{tagintsize}\PYG{+w}{ }\PYG{o}{=}\PYG{o}{=}\PYG{+w}{ }\PYG{l+m+mi}{4}\PYG{p}{)}\PYG{+w}{ }\PYG{p}{\PYGZob{}}
\PYG{+w}{            }\PYG{k+kt}{int32\PYGZus{}t}\PYG{+w}{ }\PYG{o}{*}\PYG{n}{angles\PYGZus{}real}\PYG{+w}{ }\PYG{o}{=}\PYG{+w}{ }\PYG{p}{(}\PYG{k+kt}{int32\PYGZus{}t}\PYG{+w}{ }\PYG{o}{*}\PYG{p}{)}\PYG{n}{angles}\PYG{p}{;}
\PYG{+w}{            }\PYG{k}{for}\PYG{+w}{ }\PYG{p}{(}\PYG{n}{i}\PYG{+w}{ }\PYG{o}{=}\PYG{+w}{ }\PYG{l+m+mi}{0}\PYG{p}{;}\PYG{+w}{ }\PYG{n}{i}\PYG{+w}{ }\PYG{o}{\PYGZlt{}}\PYG{+w}{ }\PYG{n}{nangles}\PYG{p}{;}\PYG{+w}{ }\PYG{o}{+}\PYG{o}{+}\PYG{n}{i}\PYG{p}{)}\PYG{+w}{ }\PYG{p}{\PYGZob{}}
\PYG{+w}{                }\PYG{n}{printf}\PYG{p}{(}\PYG{l+s}{\PYGZdq{}}\PYG{l+s}{angle \PYGZpc{} 4ld: type = \PYGZpc{}d, atoms: \PYGZpc{} 4d  \PYGZpc{} 4d  \PYGZpc{} 4d}\PYG{l+s+se}{\PYGZbs{}n}\PYG{l+s}{\PYGZdq{}}\PYG{p}{,}\PYG{n}{i}\PYG{p}{,}
\PYG{+w}{                       }\PYG{n}{angles\PYGZus{}real}\PYG{p}{[}\PYG{l+m+mi}{4}\PYG{o}{*}\PYG{n}{i}\PYG{p}{]}\PYG{p}{,}\PYG{+w}{ }\PYG{n}{angles\PYGZus{}real}\PYG{p}{[}\PYG{l+m+mi}{4}\PYG{o}{*}\PYG{n}{i}\PYG{o}{+}\PYG{l+m+mi}{1}\PYG{p}{]}\PYG{p}{,}\PYG{+w}{ }\PYG{n}{angles\PYGZus{}real}\PYG{p}{[}\PYG{l+m+mi}{4}\PYG{o}{*}\PYG{n}{i}\PYG{o}{+}\PYG{l+m+mi}{2}\PYG{p}{]}\PYG{p}{,}\PYG{+w}{ }\PYG{n}{angles\PYGZus{}real}\PYG{p}{[}\PYG{l+m+mi}{4}\PYG{o}{*}\PYG{n}{i}\PYG{o}{+}\PYG{l+m+mi}{3}\PYG{p}{]}\PYG{p}{)}\PYG{p}{;}
\PYG{+w}{            }\PYG{p}{\PYGZcb{}}
\PYG{+w}{        }\PYG{p}{\PYGZcb{}}\PYG{+w}{ }\PYG{k}{else}\PYG{+w}{ }\PYG{p}{\PYGZob{}}
\PYG{+w}{            }\PYG{k+kt}{int64\PYGZus{}t}\PYG{+w}{ }\PYG{o}{*}\PYG{n}{angles\PYGZus{}real}\PYG{+w}{ }\PYG{o}{=}\PYG{+w}{ }\PYG{p}{(}\PYG{k+kt}{int64\PYGZus{}t}\PYG{+w}{ }\PYG{o}{*}\PYG{p}{)}\PYG{n}{angles}\PYG{p}{;}
\PYG{+w}{            }\PYG{k}{for}\PYG{+w}{ }\PYG{p}{(}\PYG{n}{i}\PYG{+w}{ }\PYG{o}{=}\PYG{+w}{ }\PYG{l+m+mi}{0}\PYG{p}{;}\PYG{+w}{ }\PYG{n}{i}\PYG{+w}{ }\PYG{o}{\PYGZlt{}}\PYG{+w}{ }\PYG{n}{nangles}\PYG{p}{;}\PYG{+w}{ }\PYG{o}{+}\PYG{o}{+}\PYG{n}{i}\PYG{p}{)}\PYG{+w}{ }\PYG{p}{\PYGZob{}}
\PYG{+w}{                }\PYG{n}{printf}\PYG{p}{(}\PYG{l+s}{\PYGZdq{}}\PYG{l+s}{angle \PYGZpc{} 4ld: type = \PYGZpc{}ld, atoms: \PYGZpc{} 4ld  \PYGZpc{} 4ld  \PYGZpc{} 4ld}\PYG{l+s+se}{\PYGZbs{}n}\PYG{l+s}{\PYGZdq{}}\PYG{p}{,}\PYG{n}{i}\PYG{p}{,}
\PYG{+w}{                       }\PYG{n}{angles\PYGZus{}real}\PYG{p}{[}\PYG{l+m+mi}{4}\PYG{o}{*}\PYG{n}{i}\PYG{p}{]}\PYG{p}{,}\PYG{+w}{ }\PYG{n}{angles\PYGZus{}real}\PYG{p}{[}\PYG{l+m+mi}{4}\PYG{o}{*}\PYG{n}{i}\PYG{o}{+}\PYG{l+m+mi}{1}\PYG{p}{]}\PYG{p}{,}\PYG{+w}{ }\PYG{n}{angles\PYGZus{}real}\PYG{p}{[}\PYG{l+m+mi}{4}\PYG{o}{*}\PYG{n}{i}\PYG{o}{+}\PYG{l+m+mi}{2}\PYG{p}{]}\PYG{p}{,}\PYG{+w}{ }\PYG{n}{angles\PYGZus{}real}\PYG{p}{[}\PYG{l+m+mi}{4}\PYG{o}{*}\PYG{n}{i}\PYG{o}{+}\PYG{l+m+mi}{3}\PYG{p}{]}\PYG{p}{)}\PYG{p}{;}
\PYG{+w}{            }\PYG{p}{\PYGZcb{}}
\PYG{+w}{        }\PYG{p}{\PYGZcb{}}
\PYG{+w}{    }\PYG{p}{\PYGZcb{}}

\PYG{+w}{    }\PYG{n}{lammps\PYGZus{}close}\PYG{p}{(}\PYG{n}{handle}\PYG{p}{)}\PYG{p}{;}
\PYG{+w}{    }\PYG{n}{lammps\PYGZus{}mpi\PYGZus{}finalize}\PYG{p}{(}\PYG{p}{)}\PYG{p}{;}
\PYG{+w}{    }\PYG{n}{free}\PYG{p}{(}\PYG{n}{angles}\PYG{p}{)}\PYG{p}{;}
\PYG{+w}{    }\PYG{k}{return}\PYG{+w}{ }\PYG{l+m+mi}{0}\PYG{p}{;}
\PYG{p}{\PYGZcb{}}
\end{sphinxVerbatim}

\begin{quote}\begin{description}
\sphinxlineitem{Parameters}\begin{itemize}
\item {} 
\sphinxAtStartPar
\sphinxstyleliteralstrong{\sphinxupquote{handle}} \textendash{} pointer to a previously created LAMMPS instance 

\item {} 
\sphinxAtStartPar
\sphinxstyleliteralstrong{\sphinxupquote{data}} \textendash{} pointer to data to copy the result to 

\end{itemize}

\end{description}\end{quote}

\end{fulllineitems}



\bigskip\hrule\bigskip

\index{lammps\_gather\_dihedrals (C++ function)@\spxentry{lammps\_gather\_dihedrals}\spxextra{C++ function}}

\begin{fulllineitems}
\phantomsection\label{\detokenize{Library_scatter:_CPPv423lammps_gather_dihedralsPvPv}}\phantomsection\label{\detokenize{Library_scatter:_CPPv323lammps_gather_dihedralsPvPv}}\phantomsection\label{\detokenize{Library_scatter:_CPPv223lammps_gather_dihedralsPvPv}}\phantomsection\label{\detokenize{Library_scatter:lammps_gather_dihedrals__voidP.voidP}}
\pysigstartsignatures
\pysigstartmultiline
\pysiglinewithargsret{\phantomsection\label{\detokenize{Library_scatter:library_8h_1a93dac1f42646ebfa89521fc77a38a785}}\DUrole{kt}{void}\DUrole{w}{ }\sphinxbfcode{\sphinxupquote{\DUrole{n}{lammps\_gather\_dihedrals}}}}{\DUrole{kt}{void}\DUrole{w}{ }\DUrole{p}{*}\DUrole{n,sig-param}{handle}\sphinxparamcomma \DUrole{kt}{void}\DUrole{w}{ }\DUrole{p}{*}\DUrole{n,sig-param}{data}}{}
\pysigstopmultiline
\pysigstopsignatures
\sphinxAtStartPar
Gather type and constituent atom info for all dihedrals

\sphinxAtStartPar

\sphinxAtStartPar
\DUrole{versionmodified,added}{Added in version 8Feb2023.}

\sphinxAtStartPar
This function copies the list of all dihedrals into a buffer provided by
the calling code. The buffer will be filled with dihedral type, dihedral atom 1,
dihedral atom 2, dihedral atom 3, dihedral atom 4 for each dihedral.
Thus the buffer has to be allocated to the
dimension of 5 times the \sphinxstylestrong{total} number of dihedrals times the size of
the LAMMPS “tagint” type, which is either 4 or 8 bytes depending on
whether they are stored in 32\sphinxhyphen{}bit or 64\sphinxhyphen{}bit integers, respectively.
This size depends on the compile time settings used when compiling
the LAMMPS library and can be queried by calling
{\hyperref[\detokenize{Library_properties:_CPPv422lammps_extract_settingPvPKc}]{\sphinxcrossref{\sphinxcode{\sphinxupquote{lammps\_extract\_setting()}}}}} with the keyword “tagint”.

\sphinxAtStartPar
When running in parallel, the data buffer must be allocated on \sphinxstylestrong{all}
MPI ranks and will be filled with the information for \sphinxstylestrong{all} dihedrals
in the system.

\sphinxAtStartPar
Below is a brief C code demonstrating accessing this collected dihedral information.

\begin{sphinxVerbatim}[commandchars=\\\{\}]
\PYG{c+cp}{\PYGZsh{}}\PYG{c+cp}{include}\PYG{+w}{ }\PYG{c+cpf}{\PYGZdq{}library.h\PYGZdq{}}

\PYG{c+cp}{\PYGZsh{}}\PYG{c+cp}{include}\PYG{+w}{ }\PYG{c+cpf}{\PYGZlt{}stdint.h\PYGZgt{}}
\PYG{c+cp}{\PYGZsh{}}\PYG{c+cp}{include}\PYG{+w}{ }\PYG{c+cpf}{\PYGZlt{}stdio.h\PYGZgt{}}
\PYG{c+cp}{\PYGZsh{}}\PYG{c+cp}{include}\PYG{+w}{ }\PYG{c+cpf}{\PYGZlt{}stdlib.h\PYGZgt{}}

\PYG{k+kt}{int}\PYG{+w}{ }\PYG{n+nf}{main}\PYG{p}{(}\PYG{k+kt}{int}\PYG{+w}{ }\PYG{n}{argc}\PYG{p}{,}\PYG{+w}{ }\PYG{k+kt}{char}\PYG{+w}{ }\PYG{o}{*}\PYG{o}{*}\PYG{n}{argv}\PYG{p}{)}
\PYG{p}{\PYGZob{}}
\PYG{+w}{    }\PYG{k+kt}{int}\PYG{+w}{ }\PYG{n}{tagintsize}\PYG{p}{;}
\PYG{+w}{    }\PYG{k+kt}{int64\PYGZus{}t}\PYG{+w}{ }\PYG{n}{i}\PYG{p}{,}\PYG{+w}{ }\PYG{n}{ndihedrals}\PYG{p}{;}
\PYG{+w}{    }\PYG{k+kt}{void}\PYG{+w}{ }\PYG{o}{*}\PYG{n}{handle}\PYG{p}{,}\PYG{+w}{ }\PYG{o}{*}\PYG{n}{dihedrals}\PYG{p}{;}

\PYG{+w}{    }\PYG{n}{handle}\PYG{+w}{ }\PYG{o}{=}\PYG{+w}{ }\PYG{n}{lammps\PYGZus{}open\PYGZus{}no\PYGZus{}mpi}\PYG{p}{(}\PYG{l+m+mi}{0}\PYG{p}{,}\PYG{+w}{ }\PYG{n+nb}{NULL}\PYG{p}{,}\PYG{+w}{ }\PYG{n+nb}{NULL}\PYG{p}{)}\PYG{p}{;}
\PYG{+w}{    }\PYG{n}{lammps\PYGZus{}file}\PYG{p}{(}\PYG{n}{handle}\PYG{p}{,}\PYG{+w}{ }\PYG{l+s}{\PYGZdq{}}\PYG{l+s}{in.some\PYGZus{}input}\PYG{l+s}{\PYGZdq{}}\PYG{p}{)}\PYG{p}{;}

\PYG{+w}{    }\PYG{n}{tagintsize}\PYG{+w}{ }\PYG{o}{=}\PYG{+w}{ }\PYG{n}{lammps\PYGZus{}extract\PYGZus{}setting}\PYG{p}{(}\PYG{n}{handle}\PYG{p}{,}\PYG{+w}{ }\PYG{l+s}{\PYGZdq{}}\PYG{l+s}{tagint}\PYG{l+s}{\PYGZdq{}}\PYG{p}{)}\PYG{p}{;}
\PYG{+w}{    }\PYG{k}{if}\PYG{+w}{ }\PYG{p}{(}\PYG{n}{tagintsize}\PYG{+w}{ }\PYG{o}{=}\PYG{o}{=}\PYG{+w}{ }\PYG{l+m+mi}{4}\PYG{p}{)}
\PYG{+w}{        }\PYG{n}{ndihedrals}\PYG{+w}{ }\PYG{o}{=}\PYG{+w}{ }\PYG{o}{*}\PYG{p}{(}\PYG{k+kt}{int32\PYGZus{}t}\PYG{+w}{ }\PYG{o}{*}\PYG{p}{)}\PYG{n}{lammps\PYGZus{}extract\PYGZus{}global}\PYG{p}{(}\PYG{n}{handle}\PYG{p}{,}\PYG{+w}{ }\PYG{l+s}{\PYGZdq{}}\PYG{l+s}{ndihedrals}\PYG{l+s}{\PYGZdq{}}\PYG{p}{)}\PYG{p}{;}
\PYG{+w}{     }\PYG{k}{else}
\PYG{+w}{        }\PYG{n}{ndihedrals}\PYG{+w}{ }\PYG{o}{=}\PYG{+w}{ }\PYG{o}{*}\PYG{p}{(}\PYG{k+kt}{int64\PYGZus{}t}\PYG{+w}{ }\PYG{o}{*}\PYG{p}{)}\PYG{n}{lammps\PYGZus{}extract\PYGZus{}global}\PYG{p}{(}\PYG{n}{handle}\PYG{p}{,}\PYG{+w}{ }\PYG{l+s}{\PYGZdq{}}\PYG{l+s}{ndihedrals}\PYG{l+s}{\PYGZdq{}}\PYG{p}{)}\PYG{p}{;}
\PYG{+w}{    }\PYG{n}{dihedrals}\PYG{+w}{ }\PYG{o}{=}\PYG{+w}{ }\PYG{n}{malloc}\PYG{p}{(}\PYG{n}{ndihedrals}\PYG{+w}{ }\PYG{o}{*}\PYG{+w}{ }\PYG{l+m+mi}{5}\PYG{+w}{ }\PYG{o}{*}\PYG{+w}{ }\PYG{n}{tagintsize}\PYG{p}{)}\PYG{p}{;}

\PYG{+w}{    }\PYG{n}{lammps\PYGZus{}gather\PYGZus{}dihedrals}\PYG{p}{(}\PYG{n}{handle}\PYG{p}{,}\PYG{+w}{ }\PYG{n}{dihedrals}\PYG{p}{)}\PYG{p}{;}

\PYG{+w}{    }\PYG{k}{if}\PYG{+w}{ }\PYG{p}{(}\PYG{n}{lammps\PYGZus{}extract\PYGZus{}setting}\PYG{p}{(}\PYG{n}{handle}\PYG{p}{,}\PYG{+w}{ }\PYG{l+s}{\PYGZdq{}}\PYG{l+s}{world\PYGZus{}rank}\PYG{l+s}{\PYGZdq{}}\PYG{p}{)}\PYG{+w}{ }\PYG{o}{=}\PYG{o}{=}\PYG{+w}{ }\PYG{l+m+mi}{0}\PYG{p}{)}\PYG{+w}{ }\PYG{p}{\PYGZob{}}
\PYG{+w}{        }\PYG{k}{if}\PYG{+w}{ }\PYG{p}{(}\PYG{n}{tagintsize}\PYG{+w}{ }\PYG{o}{=}\PYG{o}{=}\PYG{+w}{ }\PYG{l+m+mi}{4}\PYG{p}{)}\PYG{+w}{ }\PYG{p}{\PYGZob{}}
\PYG{+w}{            }\PYG{k+kt}{int32\PYGZus{}t}\PYG{+w}{ }\PYG{o}{*}\PYG{n}{dihedrals\PYGZus{}real}\PYG{+w}{ }\PYG{o}{=}\PYG{+w}{ }\PYG{p}{(}\PYG{k+kt}{int32\PYGZus{}t}\PYG{+w}{ }\PYG{o}{*}\PYG{p}{)}\PYG{n}{dihedrals}\PYG{p}{;}
\PYG{+w}{            }\PYG{k}{for}\PYG{+w}{ }\PYG{p}{(}\PYG{n}{i}\PYG{+w}{ }\PYG{o}{=}\PYG{+w}{ }\PYG{l+m+mi}{0}\PYG{p}{;}\PYG{+w}{ }\PYG{n}{i}\PYG{+w}{ }\PYG{o}{\PYGZlt{}}\PYG{+w}{ }\PYG{n}{ndihedrals}\PYG{p}{;}\PYG{+w}{ }\PYG{o}{+}\PYG{o}{+}\PYG{n}{i}\PYG{p}{)}\PYG{+w}{ }\PYG{p}{\PYGZob{}}
\PYG{+w}{                }\PYG{n}{printf}\PYG{p}{(}\PYG{l+s}{\PYGZdq{}}\PYG{l+s}{dihedral \PYGZpc{} 4ld: type = \PYGZpc{}d, atoms: \PYGZpc{} 4d  \PYGZpc{} 4d  \PYGZpc{} 4d  \PYGZpc{} 4d}\PYG{l+s+se}{\PYGZbs{}n}\PYG{l+s}{\PYGZdq{}}\PYG{p}{,}\PYG{n}{i}\PYG{p}{,}
\PYG{+w}{                       }\PYG{n}{dihedrals\PYGZus{}real}\PYG{p}{[}\PYG{l+m+mi}{5}\PYG{o}{*}\PYG{n}{i}\PYG{p}{]}\PYG{p}{,}\PYG{+w}{ }\PYG{n}{dihedrals\PYGZus{}real}\PYG{p}{[}\PYG{l+m+mi}{5}\PYG{o}{*}\PYG{n}{i}\PYG{o}{+}\PYG{l+m+mi}{1}\PYG{p}{]}\PYG{p}{,}\PYG{+w}{ }\PYG{n}{dihedrals\PYGZus{}real}\PYG{p}{[}\PYG{l+m+mi}{5}\PYG{o}{*}\PYG{n}{i}\PYG{o}{+}\PYG{l+m+mi}{2}\PYG{p}{]}\PYG{p}{,}\PYG{+w}{ }\PYG{n}{dihedrals\PYGZus{}real}\PYG{p}{[}\PYG{l+m+mi}{5}\PYG{o}{*}\PYG{n}{i}\PYG{o}{+}\PYG{l+m+mi}{3}\PYG{p}{]}\PYG{p}{,}\PYG{+w}{ }\PYG{n}{dihedrals\PYGZus{}real}\PYG{p}{[}\PYG{l+m+mi}{5}\PYG{o}{*}\PYG{n}{i}\PYG{o}{+}\PYG{l+m+mi}{4}\PYG{p}{]}\PYG{p}{)}\PYG{p}{;}
\PYG{+w}{            }\PYG{p}{\PYGZcb{}}
\PYG{+w}{        }\PYG{p}{\PYGZcb{}}\PYG{+w}{ }\PYG{k}{else}\PYG{+w}{ }\PYG{p}{\PYGZob{}}
\PYG{+w}{            }\PYG{k+kt}{int64\PYGZus{}t}\PYG{+w}{ }\PYG{o}{*}\PYG{n}{dihedrals\PYGZus{}real}\PYG{+w}{ }\PYG{o}{=}\PYG{+w}{ }\PYG{p}{(}\PYG{k+kt}{int64\PYGZus{}t}\PYG{+w}{ }\PYG{o}{*}\PYG{p}{)}\PYG{n}{dihedrals}\PYG{p}{;}
\PYG{+w}{            }\PYG{k}{for}\PYG{+w}{ }\PYG{p}{(}\PYG{n}{i}\PYG{+w}{ }\PYG{o}{=}\PYG{+w}{ }\PYG{l+m+mi}{0}\PYG{p}{;}\PYG{+w}{ }\PYG{n}{i}\PYG{+w}{ }\PYG{o}{\PYGZlt{}}\PYG{+w}{ }\PYG{n}{ndihedrals}\PYG{p}{;}\PYG{+w}{ }\PYG{o}{+}\PYG{o}{+}\PYG{n}{i}\PYG{p}{)}\PYG{+w}{ }\PYG{p}{\PYGZob{}}
\PYG{+w}{                }\PYG{n}{printf}\PYG{p}{(}\PYG{l+s}{\PYGZdq{}}\PYG{l+s}{dihedral \PYGZpc{} 4ld: type = \PYGZpc{}ld, atoms: \PYGZpc{} 4ld  \PYGZpc{} 4ld  \PYGZpc{} 4ld  \PYGZpc{} 4ld}\PYG{l+s+se}{\PYGZbs{}n}\PYG{l+s}{\PYGZdq{}}\PYG{p}{,}\PYG{n}{i}\PYG{p}{,}
\PYG{+w}{                       }\PYG{n}{dihedrals\PYGZus{}real}\PYG{p}{[}\PYG{l+m+mi}{5}\PYG{o}{*}\PYG{n}{i}\PYG{p}{]}\PYG{p}{,}\PYG{+w}{ }\PYG{n}{dihedrals\PYGZus{}real}\PYG{p}{[}\PYG{l+m+mi}{5}\PYG{o}{*}\PYG{n}{i}\PYG{o}{+}\PYG{l+m+mi}{1}\PYG{p}{]}\PYG{p}{,}\PYG{+w}{ }\PYG{n}{dihedrals\PYGZus{}real}\PYG{p}{[}\PYG{l+m+mi}{5}\PYG{o}{*}\PYG{n}{i}\PYG{o}{+}\PYG{l+m+mi}{2}\PYG{p}{]}\PYG{p}{,}\PYG{+w}{ }\PYG{n}{dihedrals\PYGZus{}real}\PYG{p}{[}\PYG{l+m+mi}{5}\PYG{o}{*}\PYG{n}{i}\PYG{o}{+}\PYG{l+m+mi}{3}\PYG{p}{]}\PYG{p}{,}\PYG{+w}{ }\PYG{n}{dihedrals\PYGZus{}real}\PYG{p}{[}\PYG{l+m+mi}{5}\PYG{o}{*}\PYG{n}{i}\PYG{o}{+}\PYG{l+m+mi}{4}\PYG{p}{]}\PYG{p}{)}\PYG{p}{;}
\PYG{+w}{            }\PYG{p}{\PYGZcb{}}
\PYG{+w}{        }\PYG{p}{\PYGZcb{}}
\PYG{+w}{    }\PYG{p}{\PYGZcb{}}

\PYG{+w}{    }\PYG{n}{lammps\PYGZus{}close}\PYG{p}{(}\PYG{n}{handle}\PYG{p}{)}\PYG{p}{;}
\PYG{+w}{    }\PYG{n}{lammps\PYGZus{}mpi\PYGZus{}finalize}\PYG{p}{(}\PYG{p}{)}\PYG{p}{;}
\PYG{+w}{    }\PYG{n}{free}\PYG{p}{(}\PYG{n}{dihedrals}\PYG{p}{)}\PYG{p}{;}
\PYG{+w}{    }\PYG{k}{return}\PYG{+w}{ }\PYG{l+m+mi}{0}\PYG{p}{;}
\PYG{p}{\PYGZcb{}}
\end{sphinxVerbatim}

\begin{quote}\begin{description}
\sphinxlineitem{Parameters}\begin{itemize}
\item {} 
\sphinxAtStartPar
\sphinxstyleliteralstrong{\sphinxupquote{handle}} \textendash{} pointer to a previously created LAMMPS instance 

\item {} 
\sphinxAtStartPar
\sphinxstyleliteralstrong{\sphinxupquote{data}} \textendash{} pointer to data to copy the result to 

\end{itemize}

\end{description}\end{quote}

\end{fulllineitems}



\bigskip\hrule\bigskip

\index{lammps\_gather\_impropers (C++ function)@\spxentry{lammps\_gather\_impropers}\spxextra{C++ function}}

\begin{fulllineitems}
\phantomsection\label{\detokenize{Library_scatter:_CPPv423lammps_gather_impropersPvPv}}\phantomsection\label{\detokenize{Library_scatter:_CPPv323lammps_gather_impropersPvPv}}\phantomsection\label{\detokenize{Library_scatter:_CPPv223lammps_gather_impropersPvPv}}\phantomsection\label{\detokenize{Library_scatter:lammps_gather_impropers__voidP.voidP}}
\pysigstartsignatures
\pysigstartmultiline
\pysiglinewithargsret{\phantomsection\label{\detokenize{Library_scatter:library_8h_1abf81ea4427e375c6605fb7362e995b92}}\DUrole{kt}{void}\DUrole{w}{ }\sphinxbfcode{\sphinxupquote{\DUrole{n}{lammps\_gather\_impropers}}}}{\DUrole{kt}{void}\DUrole{w}{ }\DUrole{p}{*}\DUrole{n,sig-param}{handle}\sphinxparamcomma \DUrole{kt}{void}\DUrole{w}{ }\DUrole{p}{*}\DUrole{n,sig-param}{data}}{}
\pysigstopmultiline
\pysigstopsignatures
\sphinxAtStartPar
Gather type and constituent atom info for all impropers

\sphinxAtStartPar

\sphinxAtStartPar
\DUrole{versionmodified,added}{Added in version 8Feb2023.}

\sphinxAtStartPar
This function copies the list of all impropers into a buffer provided by
the calling code. The buffer will be filled with improper type, improper atom 1,
improper atom 2, improper atom 3, improper atom 4 for each improper.
Thus the buffer has to be allocated to the
dimension of 5 times the \sphinxstylestrong{total} number of impropers times the size of
the LAMMPS “tagint” type, which is either 4 or 8 bytes depending on
whether they are stored in 32\sphinxhyphen{}bit or 64\sphinxhyphen{}bit integers, respectively.
This size depends on the compile time settings used when compiling
the LAMMPS library and can be queried by calling
{\hyperref[\detokenize{Library_properties:_CPPv422lammps_extract_settingPvPKc}]{\sphinxcrossref{\sphinxcode{\sphinxupquote{lammps\_extract\_setting()}}}}} with the keyword “tagint”.

\sphinxAtStartPar
When running in parallel, the data buffer must be allocated on \sphinxstylestrong{all}
MPI ranks and will be filled with the information for \sphinxstylestrong{all} impropers
in the system.

\sphinxAtStartPar
Below is a brief C code demonstrating accessing this collected improper information.

\begin{sphinxVerbatim}[commandchars=\\\{\}]
\PYG{c+cp}{\PYGZsh{}}\PYG{c+cp}{include}\PYG{+w}{ }\PYG{c+cpf}{\PYGZdq{}library.h\PYGZdq{}}

\PYG{c+cp}{\PYGZsh{}}\PYG{c+cp}{include}\PYG{+w}{ }\PYG{c+cpf}{\PYGZlt{}stdint.h\PYGZgt{}}
\PYG{c+cp}{\PYGZsh{}}\PYG{c+cp}{include}\PYG{+w}{ }\PYG{c+cpf}{\PYGZlt{}stdio.h\PYGZgt{}}
\PYG{c+cp}{\PYGZsh{}}\PYG{c+cp}{include}\PYG{+w}{ }\PYG{c+cpf}{\PYGZlt{}stdlib.h\PYGZgt{}}

\PYG{k+kt}{int}\PYG{+w}{ }\PYG{n+nf}{main}\PYG{p}{(}\PYG{k+kt}{int}\PYG{+w}{ }\PYG{n}{argc}\PYG{p}{,}\PYG{+w}{ }\PYG{k+kt}{char}\PYG{+w}{ }\PYG{o}{*}\PYG{o}{*}\PYG{n}{argv}\PYG{p}{)}
\PYG{p}{\PYGZob{}}
\PYG{+w}{    }\PYG{k+kt}{int}\PYG{+w}{ }\PYG{n}{tagintsize}\PYG{p}{;}
\PYG{+w}{    }\PYG{k+kt}{int64\PYGZus{}t}\PYG{+w}{ }\PYG{n}{i}\PYG{p}{,}\PYG{+w}{ }\PYG{n}{nimpropers}\PYG{p}{;}
\PYG{+w}{    }\PYG{k+kt}{void}\PYG{+w}{ }\PYG{o}{*}\PYG{n}{handle}\PYG{p}{,}\PYG{+w}{ }\PYG{o}{*}\PYG{n}{impropers}\PYG{p}{;}

\PYG{+w}{    }\PYG{n}{handle}\PYG{+w}{ }\PYG{o}{=}\PYG{+w}{ }\PYG{n}{lammps\PYGZus{}open\PYGZus{}no\PYGZus{}mpi}\PYG{p}{(}\PYG{l+m+mi}{0}\PYG{p}{,}\PYG{+w}{ }\PYG{n+nb}{NULL}\PYG{p}{,}\PYG{+w}{ }\PYG{n+nb}{NULL}\PYG{p}{)}\PYG{p}{;}
\PYG{+w}{    }\PYG{n}{lammps\PYGZus{}file}\PYG{p}{(}\PYG{n}{handle}\PYG{p}{,}\PYG{+w}{ }\PYG{l+s}{\PYGZdq{}}\PYG{l+s}{in.some\PYGZus{}input}\PYG{l+s}{\PYGZdq{}}\PYG{p}{)}\PYG{p}{;}

\PYG{+w}{    }\PYG{n}{tagintsize}\PYG{+w}{ }\PYG{o}{=}\PYG{+w}{ }\PYG{n}{lammps\PYGZus{}extract\PYGZus{}setting}\PYG{p}{(}\PYG{n}{handle}\PYG{p}{,}\PYG{+w}{ }\PYG{l+s}{\PYGZdq{}}\PYG{l+s}{tagint}\PYG{l+s}{\PYGZdq{}}\PYG{p}{)}\PYG{p}{;}
\PYG{+w}{    }\PYG{k}{if}\PYG{+w}{ }\PYG{p}{(}\PYG{n}{tagintsize}\PYG{+w}{ }\PYG{o}{=}\PYG{o}{=}\PYG{+w}{ }\PYG{l+m+mi}{4}\PYG{p}{)}
\PYG{+w}{        }\PYG{n}{nimpropers}\PYG{+w}{ }\PYG{o}{=}\PYG{+w}{ }\PYG{o}{*}\PYG{p}{(}\PYG{k+kt}{int32\PYGZus{}t}\PYG{+w}{ }\PYG{o}{*}\PYG{p}{)}\PYG{n}{lammps\PYGZus{}extract\PYGZus{}global}\PYG{p}{(}\PYG{n}{handle}\PYG{p}{,}\PYG{+w}{ }\PYG{l+s}{\PYGZdq{}}\PYG{l+s}{nimpropers}\PYG{l+s}{\PYGZdq{}}\PYG{p}{)}\PYG{p}{;}
\PYG{+w}{     }\PYG{k}{else}
\PYG{+w}{        }\PYG{n}{nimpropers}\PYG{+w}{ }\PYG{o}{=}\PYG{+w}{ }\PYG{o}{*}\PYG{p}{(}\PYG{k+kt}{int64\PYGZus{}t}\PYG{+w}{ }\PYG{o}{*}\PYG{p}{)}\PYG{n}{lammps\PYGZus{}extract\PYGZus{}global}\PYG{p}{(}\PYG{n}{handle}\PYG{p}{,}\PYG{+w}{ }\PYG{l+s}{\PYGZdq{}}\PYG{l+s}{nimpropers}\PYG{l+s}{\PYGZdq{}}\PYG{p}{)}\PYG{p}{;}
\PYG{+w}{    }\PYG{n}{impropers}\PYG{+w}{ }\PYG{o}{=}\PYG{+w}{ }\PYG{n}{malloc}\PYG{p}{(}\PYG{n}{nimpropers}\PYG{+w}{ }\PYG{o}{*}\PYG{+w}{ }\PYG{l+m+mi}{5}\PYG{+w}{ }\PYG{o}{*}\PYG{+w}{ }\PYG{n}{tagintsize}\PYG{p}{)}\PYG{p}{;}

\PYG{+w}{    }\PYG{n}{lammps\PYGZus{}gather\PYGZus{}impropers}\PYG{p}{(}\PYG{n}{handle}\PYG{p}{,}\PYG{+w}{ }\PYG{n}{impropers}\PYG{p}{)}\PYG{p}{;}

\PYG{+w}{    }\PYG{k}{if}\PYG{+w}{ }\PYG{p}{(}\PYG{n}{lammps\PYGZus{}extract\PYGZus{}setting}\PYG{p}{(}\PYG{n}{handle}\PYG{p}{,}\PYG{+w}{ }\PYG{l+s}{\PYGZdq{}}\PYG{l+s}{world\PYGZus{}rank}\PYG{l+s}{\PYGZdq{}}\PYG{p}{)}\PYG{+w}{ }\PYG{o}{=}\PYG{o}{=}\PYG{+w}{ }\PYG{l+m+mi}{0}\PYG{p}{)}\PYG{+w}{ }\PYG{p}{\PYGZob{}}
\PYG{+w}{        }\PYG{k}{if}\PYG{+w}{ }\PYG{p}{(}\PYG{n}{tagintsize}\PYG{+w}{ }\PYG{o}{=}\PYG{o}{=}\PYG{+w}{ }\PYG{l+m+mi}{4}\PYG{p}{)}\PYG{+w}{ }\PYG{p}{\PYGZob{}}
\PYG{+w}{            }\PYG{k+kt}{int32\PYGZus{}t}\PYG{+w}{ }\PYG{o}{*}\PYG{n}{impropers\PYGZus{}real}\PYG{+w}{ }\PYG{o}{=}\PYG{+w}{ }\PYG{p}{(}\PYG{k+kt}{int32\PYGZus{}t}\PYG{+w}{ }\PYG{o}{*}\PYG{p}{)}\PYG{n}{impropers}\PYG{p}{;}
\PYG{+w}{            }\PYG{k}{for}\PYG{+w}{ }\PYG{p}{(}\PYG{n}{i}\PYG{+w}{ }\PYG{o}{=}\PYG{+w}{ }\PYG{l+m+mi}{0}\PYG{p}{;}\PYG{+w}{ }\PYG{n}{i}\PYG{+w}{ }\PYG{o}{\PYGZlt{}}\PYG{+w}{ }\PYG{n}{nimpropers}\PYG{p}{;}\PYG{+w}{ }\PYG{o}{+}\PYG{o}{+}\PYG{n}{i}\PYG{p}{)}\PYG{+w}{ }\PYG{p}{\PYGZob{}}
\PYG{+w}{                }\PYG{n}{printf}\PYG{p}{(}\PYG{l+s}{\PYGZdq{}}\PYG{l+s}{improper \PYGZpc{} 4ld: type = \PYGZpc{}d, atoms: \PYGZpc{} 4d  \PYGZpc{} 4d  \PYGZpc{} 4d  \PYGZpc{} 4d}\PYG{l+s+se}{\PYGZbs{}n}\PYG{l+s}{\PYGZdq{}}\PYG{p}{,}\PYG{n}{i}\PYG{p}{,}
\PYG{+w}{                       }\PYG{n}{impropers\PYGZus{}real}\PYG{p}{[}\PYG{l+m+mi}{5}\PYG{o}{*}\PYG{n}{i}\PYG{p}{]}\PYG{p}{,}\PYG{+w}{ }\PYG{n}{impropers\PYGZus{}real}\PYG{p}{[}\PYG{l+m+mi}{5}\PYG{o}{*}\PYG{n}{i}\PYG{o}{+}\PYG{l+m+mi}{1}\PYG{p}{]}\PYG{p}{,}\PYG{+w}{ }\PYG{n}{impropers\PYGZus{}real}\PYG{p}{[}\PYG{l+m+mi}{5}\PYG{o}{*}\PYG{n}{i}\PYG{o}{+}\PYG{l+m+mi}{2}\PYG{p}{]}\PYG{p}{,}\PYG{+w}{ }\PYG{n}{impropers\PYGZus{}real}\PYG{p}{[}\PYG{l+m+mi}{5}\PYG{o}{*}\PYG{n}{i}\PYG{o}{+}\PYG{l+m+mi}{3}\PYG{p}{]}\PYG{p}{,}\PYG{+w}{ }\PYG{n}{impropers\PYGZus{}real}\PYG{p}{[}\PYG{l+m+mi}{5}\PYG{o}{*}\PYG{n}{i}\PYG{o}{+}\PYG{l+m+mi}{4}\PYG{p}{]}\PYG{p}{)}\PYG{p}{;}
\PYG{+w}{            }\PYG{p}{\PYGZcb{}}
\PYG{+w}{        }\PYG{p}{\PYGZcb{}}\PYG{+w}{ }\PYG{k}{else}\PYG{+w}{ }\PYG{p}{\PYGZob{}}
\PYG{+w}{            }\PYG{k+kt}{int64\PYGZus{}t}\PYG{+w}{ }\PYG{o}{*}\PYG{n}{impropers\PYGZus{}real}\PYG{+w}{ }\PYG{o}{=}\PYG{+w}{ }\PYG{p}{(}\PYG{k+kt}{int64\PYGZus{}t}\PYG{+w}{ }\PYG{o}{*}\PYG{p}{)}\PYG{n}{impropers}\PYG{p}{;}
\PYG{+w}{            }\PYG{k}{for}\PYG{+w}{ }\PYG{p}{(}\PYG{n}{i}\PYG{+w}{ }\PYG{o}{=}\PYG{+w}{ }\PYG{l+m+mi}{0}\PYG{p}{;}\PYG{+w}{ }\PYG{n}{i}\PYG{+w}{ }\PYG{o}{\PYGZlt{}}\PYG{+w}{ }\PYG{n}{nimpropers}\PYG{p}{;}\PYG{+w}{ }\PYG{o}{+}\PYG{o}{+}\PYG{n}{i}\PYG{p}{)}\PYG{+w}{ }\PYG{p}{\PYGZob{}}
\PYG{+w}{                }\PYG{n}{printf}\PYG{p}{(}\PYG{l+s}{\PYGZdq{}}\PYG{l+s}{improper \PYGZpc{} 4ld: type = \PYGZpc{}ld, atoms: \PYGZpc{} 4ld  \PYGZpc{} 4ld  \PYGZpc{} 4ld  \PYGZpc{} 4ld}\PYG{l+s+se}{\PYGZbs{}n}\PYG{l+s}{\PYGZdq{}}\PYG{p}{,}\PYG{n}{i}\PYG{p}{,}
\PYG{+w}{                       }\PYG{n}{impropers\PYGZus{}real}\PYG{p}{[}\PYG{l+m+mi}{5}\PYG{o}{*}\PYG{n}{i}\PYG{p}{]}\PYG{p}{,}\PYG{+w}{ }\PYG{n}{impropers\PYGZus{}real}\PYG{p}{[}\PYG{l+m+mi}{5}\PYG{o}{*}\PYG{n}{i}\PYG{o}{+}\PYG{l+m+mi}{1}\PYG{p}{]}\PYG{p}{,}\PYG{+w}{ }\PYG{n}{impropers\PYGZus{}real}\PYG{p}{[}\PYG{l+m+mi}{5}\PYG{o}{*}\PYG{n}{i}\PYG{o}{+}\PYG{l+m+mi}{2}\PYG{p}{]}\PYG{p}{,}\PYG{+w}{ }\PYG{n}{impropers\PYGZus{}real}\PYG{p}{[}\PYG{l+m+mi}{5}\PYG{o}{*}\PYG{n}{i}\PYG{o}{+}\PYG{l+m+mi}{3}\PYG{p}{]}\PYG{p}{,}\PYG{+w}{ }\PYG{n}{impropers\PYGZus{}real}\PYG{p}{[}\PYG{l+m+mi}{5}\PYG{o}{*}\PYG{n}{i}\PYG{o}{+}\PYG{l+m+mi}{4}\PYG{p}{]}\PYG{p}{)}\PYG{p}{;}
\PYG{+w}{            }\PYG{p}{\PYGZcb{}}
\PYG{+w}{        }\PYG{p}{\PYGZcb{}}
\PYG{+w}{    }\PYG{p}{\PYGZcb{}}

\PYG{+w}{    }\PYG{n}{lammps\PYGZus{}close}\PYG{p}{(}\PYG{n}{handle}\PYG{p}{)}\PYG{p}{;}
\PYG{+w}{    }\PYG{n}{lammps\PYGZus{}mpi\PYGZus{}finalize}\PYG{p}{(}\PYG{p}{)}\PYG{p}{;}
\PYG{+w}{    }\PYG{n}{free}\PYG{p}{(}\PYG{n}{impropers}\PYG{p}{)}\PYG{p}{;}
\PYG{+w}{    }\PYG{k}{return}\PYG{+w}{ }\PYG{l+m+mi}{0}\PYG{p}{;}
\PYG{p}{\PYGZcb{}}
\end{sphinxVerbatim}

\begin{quote}\begin{description}
\sphinxlineitem{Parameters}\begin{itemize}
\item {} 
\sphinxAtStartPar
\sphinxstyleliteralstrong{\sphinxupquote{handle}} \textendash{} pointer to a previously created LAMMPS instance 

\item {} 
\sphinxAtStartPar
\sphinxstyleliteralstrong{\sphinxupquote{data}} \textendash{} pointer to data to copy the result to 

\end{itemize}

\end{description}\end{quote}

\end{fulllineitems}



\bigskip\hrule\bigskip

\index{lammps\_gather (C++ function)@\spxentry{lammps\_gather}\spxextra{C++ function}}

\begin{fulllineitems}
\phantomsection\label{\detokenize{Library_scatter:_CPPv413lammps_gatherPvPKciiPv}}\phantomsection\label{\detokenize{Library_scatter:_CPPv313lammps_gatherPvPKciiPv}}\phantomsection\label{\detokenize{Library_scatter:_CPPv213lammps_gatherPvPKciiPv}}\phantomsection\label{\detokenize{Library_scatter:lammps_gather__voidP.cCP.i.i.voidP}}
\pysigstartsignatures
\pysigstartmultiline
\pysiglinewithargsret{\phantomsection\label{\detokenize{Library_scatter:library_8h_1aeaf5356f9fafd2fd9dbab0958a0a8ac3}}\DUrole{kt}{void}\DUrole{w}{ }\sphinxbfcode{\sphinxupquote{\DUrole{n}{lammps\_gather}}}}{\DUrole{kt}{void}\DUrole{w}{ }\DUrole{p}{*}\DUrole{n,sig-param}{handle}\sphinxparamcomma \DUrole{k}{const}\DUrole{w}{ }\DUrole{kt}{char}\DUrole{w}{ }\DUrole{p}{*}\DUrole{n,sig-param}{name}\sphinxparamcomma \DUrole{kt}{int}\DUrole{w}{ }\DUrole{n,sig-param}{type}\sphinxparamcomma \DUrole{kt}{int}\DUrole{w}{ }\DUrole{n,sig-param}{count}\sphinxparamcomma \DUrole{kt}{void}\DUrole{w}{ }\DUrole{p}{*}\DUrole{n,sig-param}{data}}{}
\pysigstopmultiline
\pysigstopsignatures
\sphinxAtStartPar
Gather the named per\sphinxhyphen{}atom, per\sphinxhyphen{}atom fix, per\sphinxhyphen{}atom compute, or fix property/atom\sphinxhyphen{}based entities from all processes, in order by atom ID.

\sphinxAtStartPar

\sphinxAtStartPar
This subroutine gathers data from all processes and stores them in a one\sphinxhyphen{}dimensional array
allocated by the user. The array \sphinxstyleemphasis{data} will be ordered by atom ID, which requires consecutive IDs
(1 to \sphinxstyleemphasis{natoms}). If you need a similar array but for non\sphinxhyphen{}consecutive atom IDs, see
{\hyperref[\detokenize{Library_scatter:_CPPv420lammps_gather_concatPvPKciiPv}]{\sphinxcrossref{\sphinxcode{\sphinxupquote{lammps\_gather\_concat()}}}}}; for a similar array but for a subset of atoms, see
{\hyperref[\detokenize{Library_scatter:_CPPv420lammps_gather_subsetPvPKciiiPiPv}]{\sphinxcrossref{\sphinxcode{\sphinxupquote{lammps\_gather\_subset()}}}}}.

\sphinxAtStartPar
The \sphinxstyleemphasis{data} array will be ordered in groups of \sphinxstyleemphasis{count} values, sorted by atom ID (e.g., if \sphinxstyleemphasis{name} is
\sphinxstyleemphasis{x}, then \sphinxstyleemphasis{data} is \{x{[}0{]}{[}0{]}, x{[}0{]}{[}1{]}, x{[}0{]}{[}2{]}, x{[}1{]}{[}0{]}, x{[}1{]}{[}1{]}, x{[}1{]}{[}2{]}, x{[}2{]}{[}0{]},
\(\dots\)\}); \sphinxstyleemphasis{data} must be pre\sphinxhyphen{}allocated by the caller to the correct length
(\sphinxstyleemphasis{count}\({}\times{}\)\sphinxstyleemphasis{natoms}), as queried by {\hyperref[\detokenize{Library_properties:_CPPv417lammps_get_natomsPv}]{\sphinxcrossref{\sphinxcode{\sphinxupquote{lammps\_get\_natoms()}}}}},
{\hyperref[\detokenize{Library_properties:_CPPv421lammps_extract_globalPvPKc}]{\sphinxcrossref{\sphinxcode{\sphinxupquote{lammps\_extract\_global()}}}}}, or {\hyperref[\detokenize{Library_properties:_CPPv422lammps_extract_settingPvPKc}]{\sphinxcrossref{\sphinxcode{\sphinxupquote{lammps\_extract\_setting()}}}}}.

\sphinxAtStartPar
This function will return an error if fix or compute data are requested and the fix or compute ID
given does not have per\sphinxhyphen{}atom data.

\begin{sphinxadmonition}{note}{Restrictions}

\sphinxAtStartPar
This function is not compatible with \sphinxcode{\sphinxupquote{\sphinxhyphen{}DLAMMPS\_BIGBIG}}.

\sphinxAtStartPar
Atom IDs must be defined and must be consecutive.

\sphinxAtStartPar
The total number of atoms must not be more than 2147483647 (max 32\sphinxhyphen{}bit signed int).
\end{sphinxadmonition}

\begin{quote}\begin{description}
\sphinxlineitem{Parameters}\begin{itemize}
\item {} 
\sphinxAtStartPar
\sphinxstyleliteralstrong{\sphinxupquote{handle}} \textendash{} pointer to a previously created LAMMPS instance 

\item {} 
\sphinxAtStartPar
\sphinxstyleliteralstrong{\sphinxupquote{name}} \textendash{} desired quantity (e.g., “x” or “f” for atom properties, “f\_id” for per\sphinxhyphen{}atom fix data, “c\_id” for per\sphinxhyphen{}atom compute data, “d\_name” or “i\_name” for fix property/atom vectors with \sphinxstyleemphasis{count} = 1, “d2\_name” or “i2\_name” for fix property/atom vectors with \sphinxstyleemphasis{count} \textgreater{} 1) 

\item {} 
\sphinxAtStartPar
\sphinxstyleliteralstrong{\sphinxupquote{type}} \textendash{} 0 for \sphinxcode{\sphinxupquote{int}} values, 1 for \sphinxcode{\sphinxupquote{double}} values 

\item {} 
\sphinxAtStartPar
\sphinxstyleliteralstrong{\sphinxupquote{count}} \textendash{} number of per\sphinxhyphen{}atom values (e.g., 1 for \sphinxstyleemphasis{type} or \sphinxstyleemphasis{charge}, 3 for \sphinxstyleemphasis{x} or \sphinxstyleemphasis{f}); use \sphinxstyleemphasis{count} = 3 with \sphinxstyleemphasis{image} if you want the image flags unpacked into (\sphinxstyleemphasis{x},\sphinxstyleemphasis{y},\sphinxstyleemphasis{z}) components. 

\item {} 
\sphinxAtStartPar
\sphinxstyleliteralstrong{\sphinxupquote{data}} \textendash{} per\sphinxhyphen{}atom values packed into a one\sphinxhyphen{}dimensional array of length \sphinxstyleemphasis{natoms} * \sphinxstyleemphasis{count}. 

\end{itemize}

\end{description}\end{quote}

\end{fulllineitems}



\bigskip\hrule\bigskip

\index{lammps\_gather\_concat (C++ function)@\spxentry{lammps\_gather\_concat}\spxextra{C++ function}}

\begin{fulllineitems}
\phantomsection\label{\detokenize{Library_scatter:_CPPv420lammps_gather_concatPvPKciiPv}}\phantomsection\label{\detokenize{Library_scatter:_CPPv320lammps_gather_concatPvPKciiPv}}\phantomsection\label{\detokenize{Library_scatter:_CPPv220lammps_gather_concatPvPKciiPv}}\phantomsection\label{\detokenize{Library_scatter:lammps_gather_concat__voidP.cCP.i.i.voidP}}
\pysigstartsignatures
\pysigstartmultiline
\pysiglinewithargsret{\phantomsection\label{\detokenize{Library_scatter:library_8h_1add32cd75d71123448c752d69d96a9009}}\DUrole{kt}{void}\DUrole{w}{ }\sphinxbfcode{\sphinxupquote{\DUrole{n}{lammps\_gather\_concat}}}}{\DUrole{kt}{void}\DUrole{w}{ }\DUrole{p}{*}\DUrole{n,sig-param}{handle}\sphinxparamcomma \DUrole{k}{const}\DUrole{w}{ }\DUrole{kt}{char}\DUrole{w}{ }\DUrole{p}{*}\DUrole{n,sig-param}{name}\sphinxparamcomma \DUrole{kt}{int}\DUrole{w}{ }\DUrole{n,sig-param}{type}\sphinxparamcomma \DUrole{kt}{int}\DUrole{w}{ }\DUrole{n,sig-param}{count}\sphinxparamcomma \DUrole{kt}{void}\DUrole{w}{ }\DUrole{p}{*}\DUrole{n,sig-param}{data}}{}
\pysigstopmultiline
\pysigstopsignatures
\sphinxAtStartPar
Gather the named per\sphinxhyphen{}atom, per\sphinxhyphen{}atom fix, per\sphinxhyphen{}atom compute, or fix property/atom\sphinxhyphen{}based entities from all processes, unordered.

\sphinxAtStartPar

\sphinxAtStartPar
This subroutine gathers data for all atoms and stores them in a one\sphinxhyphen{}dimensional array allocated by
the user. The data will be a concatenation of chunks from each processor’s owned atoms, in
whatever order the atoms are in on each processor. This process has no requirement that the atom
IDs be consecutive. If you need the ID of each atom, you can do another call to either
{\hyperref[\detokenize{Library_scatter:_CPPv426lammps_gather_atoms_concatPvPKciiPv}]{\sphinxcrossref{\sphinxcode{\sphinxupquote{lammps\_gather\_atoms\_concat()}}}}} or {\hyperref[\detokenize{Library_scatter:_CPPv420lammps_gather_concatPvPKciiPv}]{\sphinxcrossref{\sphinxcode{\sphinxupquote{lammps\_gather\_concat()}}}}} with \sphinxstyleemphasis{name} set to
\sphinxcode{\sphinxupquote{id}}. If you have consecutive IDs and want the data to be in order, use
{\hyperref[\detokenize{Library_scatter:_CPPv413lammps_gatherPvPKciiPv}]{\sphinxcrossref{\sphinxcode{\sphinxupquote{lammps\_gather()}}}}}; for a similar array but for a subset of atoms, use
{\hyperref[\detokenize{Library_scatter:_CPPv420lammps_gather_subsetPvPKciiiPiPv}]{\sphinxcrossref{\sphinxcode{\sphinxupquote{lammps\_gather\_subset()}}}}}.

\sphinxAtStartPar
The \sphinxstyleemphasis{data} array will be in groups of \sphinxstyleemphasis{count} values, with \sphinxstyleemphasis{natoms} groups total, but not in order
by atom ID (e.g., if \sphinxstyleemphasis{name} is \sphinxstyleemphasis{x} and \sphinxstyleemphasis{count} is 3, then \sphinxstyleemphasis{data} might be something like
\{x{[}10{]}{[}0{]}, x{[}10{]}{[}1{]}, x{[}10{]}{[}2{]}, x{[}2{]}{[}0{]}, x{[}2{]}{[}1{]}, x{[}2{]}{[}2{]}, x{[}4{]}{[}0{]}, \(\dots\)\}); \sphinxstyleemphasis{data} must be
pre\sphinxhyphen{}allocated by the caller to length (\sphinxstyleemphasis{count} \(\times\) \sphinxstyleemphasis{natoms}), as queried by
{\hyperref[\detokenize{Library_properties:_CPPv417lammps_get_natomsPv}]{\sphinxcrossref{\sphinxcode{\sphinxupquote{lammps\_get\_natoms()}}}}}, {\hyperref[\detokenize{Library_properties:_CPPv421lammps_extract_globalPvPKc}]{\sphinxcrossref{\sphinxcode{\sphinxupquote{lammps\_extract\_global()}}}}}, or
{\hyperref[\detokenize{Library_properties:_CPPv422lammps_extract_settingPvPKc}]{\sphinxcrossref{\sphinxcode{\sphinxupquote{lammps\_extract\_setting()}}}}}.

\begin{sphinxadmonition}{note}{Restrictions}

\sphinxAtStartPar
This function is not compatible with \sphinxcode{\sphinxupquote{\sphinxhyphen{}DLAMMPS\_BIGBIG}}.

\sphinxAtStartPar
Atom IDs must be defined.

\sphinxAtStartPar
The total number of atoms must be less than 2147483647 (max 32\sphinxhyphen{}bit signed int).
\end{sphinxadmonition}

\begin{quote}\begin{description}
\sphinxlineitem{Parameters}\begin{itemize}
\item {} 
\sphinxAtStartPar
\sphinxstyleliteralstrong{\sphinxupquote{handle}} \textendash{} pointer to a previously created LAMMPS instance 

\item {} 
\sphinxAtStartPar
\sphinxstyleliteralstrong{\sphinxupquote{name}} \textendash{} desired quantity (e.g., “x” or “f” for atom properties, “f\_id” for per\sphinxhyphen{}atom fix data, “c\_id” for per\sphinxhyphen{}atom compute data, “d\_name” or “i\_name” for fix property/atom vectors with count = 1, “d2\_name” or “i2\_name” for fix property/atom vectors with count \textgreater{} 1) 

\item {} 
\sphinxAtStartPar
\sphinxstyleliteralstrong{\sphinxupquote{type}} \textendash{} 0 for \sphinxcode{\sphinxupquote{int}} values, 1 for \sphinxcode{\sphinxupquote{double}} values 

\item {} 
\sphinxAtStartPar
\sphinxstyleliteralstrong{\sphinxupquote{count}} \textendash{} number of per\sphinxhyphen{}atom values (e.g., 1 for \sphinxstyleemphasis{type} or \sphinxstyleemphasis{charge}, 3 for \sphinxstyleemphasis{x} or \sphinxstyleemphasis{f}); use \sphinxstyleemphasis{count} = 3 with \sphinxstyleemphasis{image} if you want the image flags unpacked into (\sphinxstyleemphasis{x},\sphinxstyleemphasis{y},\sphinxstyleemphasis{z}) components. 

\item {} 
\sphinxAtStartPar
\sphinxstyleliteralstrong{\sphinxupquote{data}} \textendash{} per\sphinxhyphen{}atom values packed into a one\sphinxhyphen{}dimensional array of length \sphinxstyleemphasis{natoms} * \sphinxstyleemphasis{count}. 

\end{itemize}

\end{description}\end{quote}

\end{fulllineitems}



\bigskip\hrule\bigskip

\index{lammps\_gather\_subset (C++ function)@\spxentry{lammps\_gather\_subset}\spxextra{C++ function}}

\begin{fulllineitems}
\phantomsection\label{\detokenize{Library_scatter:_CPPv420lammps_gather_subsetPvPKciiiPiPv}}\phantomsection\label{\detokenize{Library_scatter:_CPPv320lammps_gather_subsetPvPKciiiPiPv}}\phantomsection\label{\detokenize{Library_scatter:_CPPv220lammps_gather_subsetPvPKciiiPiPv}}\phantomsection\label{\detokenize{Library_scatter:lammps_gather_subset__voidP.cCP.i.i.i.iP.voidP}}
\pysigstartsignatures
\pysigstartmultiline
\pysiglinewithargsret{\phantomsection\label{\detokenize{Library_scatter:library_8h_1ab098fa5d18996508a54430ff5e3a0faf}}\DUrole{kt}{void}\DUrole{w}{ }\sphinxbfcode{\sphinxupquote{\DUrole{n}{lammps\_gather\_subset}}}}{\DUrole{kt}{void}\DUrole{w}{ }\DUrole{p}{*}\DUrole{n,sig-param}{handle}\sphinxparamcomma \DUrole{k}{const}\DUrole{w}{ }\DUrole{kt}{char}\DUrole{w}{ }\DUrole{p}{*}\DUrole{n,sig-param}{name}\sphinxparamcomma \DUrole{kt}{int}\DUrole{w}{ }\DUrole{n,sig-param}{type}\sphinxparamcomma \DUrole{kt}{int}\DUrole{w}{ }\DUrole{n,sig-param}{count}\sphinxparamcomma \DUrole{kt}{int}\DUrole{w}{ }\DUrole{n,sig-param}{ndata}\sphinxparamcomma \DUrole{kt}{int}\DUrole{w}{ }\DUrole{p}{*}\DUrole{n,sig-param}{ids}\sphinxparamcomma \DUrole{kt}{void}\DUrole{w}{ }\DUrole{p}{*}\DUrole{n,sig-param}{data}}{}
\pysigstopmultiline
\pysigstopsignatures
\sphinxAtStartPar
Gather the named per\sphinxhyphen{}atom, per\sphinxhyphen{}atom fix, per\sphinxhyphen{}atom compute, or fix property/atom\sphinxhyphen{}based entities from all processes for a subset of atoms.

\sphinxAtStartPar

\sphinxAtStartPar
This subroutine gathers data for the requested atom IDs and stores them in a one\sphinxhyphen{}dimensional array
allocated by the user. The data will be ordered by atom ID, but there is no requirement that the
IDs be consecutive. If you wish to return a similar array for \sphinxstyleemphasis{all} the atoms, use
{\hyperref[\detokenize{Library_scatter:_CPPv413lammps_gatherPvPKciiPv}]{\sphinxcrossref{\sphinxcode{\sphinxupquote{lammps\_gather()}}}}} or {\hyperref[\detokenize{Library_scatter:_CPPv420lammps_gather_concatPvPKciiPv}]{\sphinxcrossref{\sphinxcode{\sphinxupquote{lammps\_gather\_concat()}}}}}.

\sphinxAtStartPar
The \sphinxstyleemphasis{data} array will be in groups of \sphinxstyleemphasis{count} values, sorted by atom ID in the same order as the
array \sphinxstyleemphasis{ids} (e.g., if \sphinxstyleemphasis{name} is \sphinxstyleemphasis{x}, \sphinxstyleemphasis{count} = 3, and \sphinxstyleemphasis{ids} is \{100, 57, 210\}, then \sphinxstyleemphasis{data} might
look like \{x{[}100{]}{[}0{]}, x{[}100{]}{[}1{]}, x{[}100{]}{[}2{]}, x{[}57{]}{[}0{]}, x{[}57{]}{[}1{]}, x{[}57{]}{[}2{]}, x{[}210{]}{[}0{]},
\(\dots\)\}); \sphinxstyleemphasis{ids} must be provided by the user with length \sphinxstyleemphasis{ndata}, and \sphinxstyleemphasis{data} must be
pre\sphinxhyphen{}allocated by the caller to length (\sphinxstyleemphasis{count}\({}\times{}\)\sphinxstyleemphasis{ndata}).

\begin{sphinxadmonition}{note}{Restrictions}

\sphinxAtStartPar
This function is not compatible with \sphinxcode{\sphinxupquote{\sphinxhyphen{}DLAMMPS\_BIGBIG}}.

\sphinxAtStartPar
Atom IDs must be defined and an {\hyperref[\detokenize{atom_modify::doc}]{\sphinxcrossref{\DUrole{doc}{atom map must be enabled}}}}

\sphinxAtStartPar
The total number of atoms must not be more than 2147483647 (max 32\sphinxhyphen{}bit signed int).
\end{sphinxadmonition}

\begin{quote}\begin{description}
\sphinxlineitem{Parameters}\begin{itemize}
\item {} 
\sphinxAtStartPar
\sphinxstyleliteralstrong{\sphinxupquote{handle}} \textendash{} pointer to a previously created LAMMPS instance 

\item {} 
\sphinxAtStartPar
\sphinxstyleliteralstrong{\sphinxupquote{name}} \textendash{} desired quantity (e.g., “x” or “f” for atom properties, “f\_id” for per\sphinxhyphen{}atom fix data, “c\_id” for per\sphinxhyphen{}atom compute data, “d\_name” or “i\_name” for fix property/atom vectors with \sphinxstyleemphasis{count} = 1, “d2\_name” or “i2\_name” for fix property/atom vectors with \sphinxstyleemphasis{count} \textgreater{} 1) 

\item {} 
\sphinxAtStartPar
\sphinxstyleliteralstrong{\sphinxupquote{type}} \textendash{} 0 for \sphinxcode{\sphinxupquote{int}} values, 1 for \sphinxcode{\sphinxupquote{double}} values 

\item {} 
\sphinxAtStartPar
\sphinxstyleliteralstrong{\sphinxupquote{count}} \textendash{} number of per\sphinxhyphen{}atom values (e.g., 1 for \sphinxstyleemphasis{type} or \sphinxstyleemphasis{charge}, 3 for \sphinxstyleemphasis{x} or \sphinxstyleemphasis{f}); use \sphinxstyleemphasis{count} = 3 with \sphinxstyleemphasis{image} if you want the image flags unpacked into (\sphinxstyleemphasis{x},\sphinxstyleemphasis{y},\sphinxstyleemphasis{z}) components. 

\item {} 
\sphinxAtStartPar
\sphinxstyleliteralstrong{\sphinxupquote{ndata}} \textendash{} number of atoms for which to return data (can be all of them) 

\item {} 
\sphinxAtStartPar
\sphinxstyleliteralstrong{\sphinxupquote{ids}} \textendash{} list of \sphinxstyleemphasis{ndata} atom IDs for which to return data 

\item {} 
\sphinxAtStartPar
\sphinxstyleliteralstrong{\sphinxupquote{data}} \textendash{} per\sphinxhyphen{}atom values packed into a one\sphinxhyphen{}dimensional array of length \sphinxstyleemphasis{ndata} * \sphinxstyleemphasis{count}. 

\end{itemize}

\end{description}\end{quote}

\end{fulllineitems}



\bigskip\hrule\bigskip

\index{lammps\_scatter (C++ function)@\spxentry{lammps\_scatter}\spxextra{C++ function}}

\begin{fulllineitems}
\phantomsection\label{\detokenize{Library_scatter:_CPPv414lammps_scatterPvPKciiPv}}\phantomsection\label{\detokenize{Library_scatter:_CPPv314lammps_scatterPvPKciiPv}}\phantomsection\label{\detokenize{Library_scatter:_CPPv214lammps_scatterPvPKciiPv}}\phantomsection\label{\detokenize{Library_scatter:lammps_scatter__voidP.cCP.i.i.voidP}}
\pysigstartsignatures
\pysigstartmultiline
\pysiglinewithargsret{\phantomsection\label{\detokenize{Library_scatter:library_8h_1a1251a6d15588762a6810ed4021bfce11}}\DUrole{kt}{void}\DUrole{w}{ }\sphinxbfcode{\sphinxupquote{\DUrole{n}{lammps\_scatter}}}}{\DUrole{kt}{void}\DUrole{w}{ }\DUrole{p}{*}\DUrole{n,sig-param}{handle}\sphinxparamcomma \DUrole{k}{const}\DUrole{w}{ }\DUrole{kt}{char}\DUrole{w}{ }\DUrole{p}{*}\DUrole{n,sig-param}{name}\sphinxparamcomma \DUrole{kt}{int}\DUrole{w}{ }\DUrole{n,sig-param}{type}\sphinxparamcomma \DUrole{kt}{int}\DUrole{w}{ }\DUrole{n,sig-param}{count}\sphinxparamcomma \DUrole{kt}{void}\DUrole{w}{ }\DUrole{p}{*}\DUrole{n,sig-param}{data}}{}
\pysigstopmultiline
\pysigstopsignatures
\sphinxAtStartPar
Scatter the named per\sphinxhyphen{}atom, per\sphinxhyphen{}atom fix, per\sphinxhyphen{}atom compute, or fix property/atom\sphinxhyphen{}based entity in \sphinxstyleemphasis{data} to all processes.

\sphinxAtStartPar

\sphinxAtStartPar
This subroutine takes data stored in a one\sphinxhyphen{}dimensional array supplied by the user and scatters
them to all atoms on all processes. The data must be ordered by atom ID, with the requirement that
the IDs be consecutive. Use {\hyperref[\detokenize{Library_scatter:_CPPv421lammps_scatter_subsetPvPKciiiPiPv}]{\sphinxcrossref{\sphinxcode{\sphinxupquote{lammps\_scatter\_subset()}}}}} to scatter data for some (or all)
atoms, unordered.

\sphinxAtStartPar
The \sphinxstyleemphasis{data} array needs to be ordered in groups of \sphinxstyleemphasis{count} values, sorted by atom ID (e.g., if
\sphinxstyleemphasis{name} is \sphinxstyleemphasis{x} and \sphinxstyleemphasis{count} = 3, then \sphinxstyleemphasis{data} = \{x{[}0{]}{[}0{]}, x{[}0{]}{[}1{]}, x{[}0{]}{[}2{]}, x{[}1{]}{[}0{]}, x{[}1{]}{[}1{]},
x{[}1{]}{[}2{]}, x{[}2{]}{[}0{]}, \(\dots\)\}); \sphinxstyleemphasis{data} must be of length (\sphinxstyleemphasis{count} \(\times\) \sphinxstyleemphasis{natoms}).

\begin{sphinxadmonition}{note}{Restrictions}

\sphinxAtStartPar
This function is not compatible with \sphinxcode{\sphinxupquote{\sphinxhyphen{}DLAMMPS\_BIGBIG}}.

\sphinxAtStartPar
Atom IDs must be defined, must be consecutive, and an
{\hyperref[\detokenize{atom_modify::doc}]{\sphinxcrossref{\DUrole{doc}{atom map must be enabled}}}}

\sphinxAtStartPar
The total number of atoms must not be more than 2147483647 (max 32\sphinxhyphen{}bit signed int).
\end{sphinxadmonition}

\begin{quote}\begin{description}
\sphinxlineitem{Parameters}\begin{itemize}
\item {} 
\sphinxAtStartPar
\sphinxstyleliteralstrong{\sphinxupquote{handle}} \textendash{} pointer to a previously created LAMMPS instance 

\item {} 
\sphinxAtStartPar
\sphinxstyleliteralstrong{\sphinxupquote{name}} \textendash{} desired quantity (e.g., “x” or “f” for atom properties, “f\_id” for per\sphinxhyphen{}atom fix data, “c\_id” for per\sphinxhyphen{}atom compute data, “d\_name” or “i\_name” for fix property/atom vectors with \sphinxstyleemphasis{count} = 1, “d2\_name” or “i2\_name” for fix property/atom vectors with \sphinxstyleemphasis{count} \textgreater{} 1) 

\item {} 
\sphinxAtStartPar
\sphinxstyleliteralstrong{\sphinxupquote{type}} \textendash{} 0 for \sphinxcode{\sphinxupquote{int}} values, 1 for \sphinxcode{\sphinxupquote{double}} values 

\item {} 
\sphinxAtStartPar
\sphinxstyleliteralstrong{\sphinxupquote{count}} \textendash{} number of per\sphinxhyphen{}atom values (e.g., 1 for \sphinxstyleemphasis{type} or \sphinxstyleemphasis{charge}, 3 for \sphinxstyleemphasis{x} or \sphinxstyleemphasis{f}); use \sphinxstyleemphasis{count} = 3 with \sphinxstyleemphasis{image} if you have a single image flag packed into (\sphinxstyleemphasis{x},\sphinxstyleemphasis{y},\sphinxstyleemphasis{z}) components. 

\item {} 
\sphinxAtStartPar
\sphinxstyleliteralstrong{\sphinxupquote{data}} \textendash{} per\sphinxhyphen{}atom values packed in a one\sphinxhyphen{}dimensional array of length \sphinxstyleemphasis{natoms} * \sphinxstyleemphasis{count}. 

\end{itemize}

\end{description}\end{quote}

\end{fulllineitems}



\bigskip\hrule\bigskip

\index{lammps\_scatter\_subset (C++ function)@\spxentry{lammps\_scatter\_subset}\spxextra{C++ function}}

\begin{fulllineitems}
\phantomsection\label{\detokenize{Library_scatter:_CPPv421lammps_scatter_subsetPvPKciiiPiPv}}\phantomsection\label{\detokenize{Library_scatter:_CPPv321lammps_scatter_subsetPvPKciiiPiPv}}\phantomsection\label{\detokenize{Library_scatter:_CPPv221lammps_scatter_subsetPvPKciiiPiPv}}\phantomsection\label{\detokenize{Library_scatter:lammps_scatter_subset__voidP.cCP.i.i.i.iP.voidP}}
\pysigstartsignatures
\pysigstartmultiline
\pysiglinewithargsret{\phantomsection\label{\detokenize{Library_scatter:library_8h_1ae379c37f2f8136e130030ddbe832695a}}\DUrole{kt}{void}\DUrole{w}{ }\sphinxbfcode{\sphinxupquote{\DUrole{n}{lammps\_scatter\_subset}}}}{\DUrole{kt}{void}\DUrole{w}{ }\DUrole{p}{*}\DUrole{n,sig-param}{handle}\sphinxparamcomma \DUrole{k}{const}\DUrole{w}{ }\DUrole{kt}{char}\DUrole{w}{ }\DUrole{p}{*}\DUrole{n,sig-param}{name}\sphinxparamcomma \DUrole{kt}{int}\DUrole{w}{ }\DUrole{n,sig-param}{type}\sphinxparamcomma \DUrole{kt}{int}\DUrole{w}{ }\DUrole{n,sig-param}{count}\sphinxparamcomma \DUrole{kt}{int}\DUrole{w}{ }\DUrole{n,sig-param}{ndata}\sphinxparamcomma \DUrole{kt}{int}\DUrole{w}{ }\DUrole{p}{*}\DUrole{n,sig-param}{ids}\sphinxparamcomma \DUrole{kt}{void}\DUrole{w}{ }\DUrole{p}{*}\DUrole{n,sig-param}{data}}{}
\pysigstopmultiline
\pysigstopsignatures
\sphinxAtStartPar
Scatter the named per\sphinxhyphen{}atom, per\sphinxhyphen{}atom fix, per\sphinxhyphen{}atom compute, or fix property/atom\sphinxhyphen{}based entities in \sphinxstyleemphasis{data} from a subset of atoms to all processes.

\sphinxAtStartPar

\sphinxAtStartPar
This subroutine takes data stored in a one\sphinxhyphen{}dimensional array supplied by the
user and scatters them to a subset of atoms on all processes. The array
\sphinxstyleemphasis{data} contains data associated with atom IDs, but there is no requirement that
the IDs be consecutive, as they are provided in a separate array.
Use {\hyperref[\detokenize{Library_scatter:_CPPv414lammps_scatterPvPKciiPv}]{\sphinxcrossref{\sphinxcode{\sphinxupquote{lammps\_scatter()}}}}} to scatter data for all atoms, in order.

\sphinxAtStartPar
The \sphinxstyleemphasis{data} array needs to be organized in groups of \sphinxstyleemphasis{count} values, with the
groups in the same order as the array \sphinxstyleemphasis{ids}. For example, if you want \sphinxstyleemphasis{data}
to be the array \{x{[}1{]}{[}0{]}, x{[}1{]}{[}1{]}, x{[}1{]}{[}2{]}, x{[}100{]}{[}0{]}, x{[}100{]}{[}1{]}, x{[}100{]}{[}2{]},
x{[}57{]}{[}0{]}, x{[}57{]}{[}1{]}, x{[}57{]}{[}2{]}\}, then \sphinxstyleemphasis{count} = 3, \sphinxstyleemphasis{ndata} = 3, and \sphinxstyleemphasis{ids} would
be \{1, 100, 57\}.

\begin{sphinxadmonition}{note}{Restrictions}

\sphinxAtStartPar
This function is not compatible with \sphinxcode{\sphinxupquote{\sphinxhyphen{}DLAMMPS\_BIGBIG}}.

\sphinxAtStartPar
Atom IDs must be defined and an {\hyperref[\detokenize{atom_modify::doc}]{\sphinxcrossref{\DUrole{doc}{atom map must be enabled}}}}

\sphinxAtStartPar
The total number of atoms must not be more than 2147483647 (max 32\sphinxhyphen{}bit signed int).
\end{sphinxadmonition}

\begin{quote}\begin{description}
\sphinxlineitem{Parameters}\begin{itemize}
\item {} 
\sphinxAtStartPar
\sphinxstyleliteralstrong{\sphinxupquote{handle}} \textendash{} pointer to a previously created LAMMPS instance 

\item {} 
\sphinxAtStartPar
\sphinxstyleliteralstrong{\sphinxupquote{name}} \textendash{} desired quantity (e.g., “x” or “f” for atom properties, “f\_id” for per\sphinxhyphen{}atom fix data, “c\_id” for per\sphinxhyphen{}atom compute data, “d\_name” or “i\_name” for fix property/atom vectors with \sphinxstyleemphasis{count} = 1, “d2\_name” or “i2\_name” for fix property/atom vectors with \sphinxstyleemphasis{count} \textgreater{} 1) 

\item {} 
\sphinxAtStartPar
\sphinxstyleliteralstrong{\sphinxupquote{type}} \textendash{} 0 for \sphinxcode{\sphinxupquote{int}} values, 1 for \sphinxcode{\sphinxupquote{double}} values 

\item {} 
\sphinxAtStartPar
\sphinxstyleliteralstrong{\sphinxupquote{count}} \textendash{} number of per\sphinxhyphen{}atom values (e.g., 1 for \sphinxstyleemphasis{type} or \sphinxstyleemphasis{charge}, 3 for \sphinxstyleemphasis{x} or \sphinxstyleemphasis{f}); use \sphinxstyleemphasis{count} = 3 with “image” if you want single image flags unpacked into (\sphinxstyleemphasis{x},\sphinxstyleemphasis{y},\sphinxstyleemphasis{z}) 

\item {} 
\sphinxAtStartPar
\sphinxstyleliteralstrong{\sphinxupquote{ndata}} \textendash{} number of atoms listed in \sphinxstyleemphasis{ids} and \sphinxstyleemphasis{data} arrays 

\item {} 
\sphinxAtStartPar
\sphinxstyleliteralstrong{\sphinxupquote{ids}} \textendash{} list of \sphinxstyleemphasis{ndata} atom IDs to scatter data to 

\item {} 
\sphinxAtStartPar
\sphinxstyleliteralstrong{\sphinxupquote{data}} \textendash{} per\sphinxhyphen{}atom values packed in a 1\sphinxhyphen{}dimensional array of length \sphinxstyleemphasis{ndata} * \sphinxstyleemphasis{count}. 

\end{itemize}

\end{description}\end{quote}

\end{fulllineitems}



\bigskip\hrule\bigskip

\index{lammps\_create\_atoms (C++ function)@\spxentry{lammps\_create\_atoms}\spxextra{C++ function}}

\begin{fulllineitems}
\phantomsection\label{\detokenize{Library_scatter:_CPPv419lammps_create_atomsPviPKiPKiPKdPKdPKii}}\phantomsection\label{\detokenize{Library_scatter:_CPPv319lammps_create_atomsPviPKiPKiPKdPKdPKii}}\phantomsection\label{\detokenize{Library_scatter:_CPPv219lammps_create_atomsPviPKiPKiPKdPKdPKii}}\phantomsection\label{\detokenize{Library_scatter:lammps_create_atoms__voidP.i.iCP.iCP.doubleCP.doubleCP.iCP.i}}
\pysigstartsignatures
\pysigstartmultiline
\pysiglinewithargsret{\phantomsection\label{\detokenize{Library_scatter:library_8h_1aeebd5ebadb0c014f80e2c25dd369ff44}}\DUrole{kt}{int}\DUrole{w}{ }\sphinxbfcode{\sphinxupquote{\DUrole{n}{lammps\_create\_atoms}}}}{\DUrole{kt}{void}\DUrole{w}{ }\DUrole{p}{*}\DUrole{n,sig-param}{handle}\sphinxparamcomma \DUrole{kt}{int}\DUrole{w}{ }\DUrole{n,sig-param}{n}\sphinxparamcomma \DUrole{k}{const}\DUrole{w}{ }\DUrole{kt}{int}\DUrole{w}{ }\DUrole{p}{*}\DUrole{n,sig-param}{id}\sphinxparamcomma \DUrole{k}{const}\DUrole{w}{ }\DUrole{kt}{int}\DUrole{w}{ }\DUrole{p}{*}\DUrole{n,sig-param}{type}\sphinxparamcomma \DUrole{k}{const}\DUrole{w}{ }\DUrole{kt}{double}\DUrole{w}{ }\DUrole{p}{*}\DUrole{n,sig-param}{x}\sphinxparamcomma \DUrole{k}{const}\DUrole{w}{ }\DUrole{kt}{double}\DUrole{w}{ }\DUrole{p}{*}\DUrole{n,sig-param}{v}\sphinxparamcomma \DUrole{k}{const}\DUrole{w}{ }\DUrole{kt}{int}\DUrole{w}{ }\DUrole{p}{*}\DUrole{n,sig-param}{image}\sphinxparamcomma \DUrole{kt}{int}\DUrole{w}{ }\DUrole{n,sig-param}{bexpand}}{}
\pysigstopmultiline
\pysigstopsignatures
\sphinxAtStartPar
Create N atoms from list of coordinates

\sphinxAtStartPar

\sphinxAtStartPar
The prototype for this function when compiling with \sphinxcode{\sphinxupquote{\sphinxhyphen{}DLAMMPS\_BIGBIG}}
is:

\begin{sphinxVerbatim}[commandchars=\\\{\}]
\PYG{k+kt}{int}\PYG{+w}{ }\PYG{n+nf}{lammps\PYGZus{}create\PYGZus{}atoms}\PYG{p}{(}\PYG{k+kt}{void}\PYG{+w}{ }\PYG{o}{*}\PYG{n}{handle}\PYG{p}{,}\PYG{+w}{ }\PYG{k+kt}{int}\PYG{+w}{ }\PYG{n}{n}\PYG{p}{,}\PYG{+w}{ }\PYG{k+kt}{int64\PYGZus{}t}\PYG{+w}{ }\PYG{o}{*}\PYG{n}{id}\PYG{p}{,}\PYG{+w}{ }\PYG{k+kt}{int}\PYG{+w}{ }\PYG{o}{*}\PYG{n}{type}\PYG{p}{,}\PYG{+w}{ }\PYG{k+kt}{double}\PYG{+w}{ }\PYG{o}{*}\PYG{n}{x}\PYG{p}{,}\PYG{+w}{ }\PYG{k+kt}{double}\PYG{+w}{ }\PYG{o}{*}\PYG{n}{v}\PYG{p}{,}\PYG{+w}{ }\PYG{k+kt}{int64\PYGZus{}t}\PYG{+w}{ }\PYG{o}{*}\PYG{n}{image}\PYG{p}{,}\PYG{+w}{ }\PYG{k+kt}{int}\PYG{+w}{ }\PYG{n}{bexpand}\PYG{p}{)}\PYG{p}{;}
\end{sphinxVerbatim}

\sphinxAtStartPar
This function creates additional atoms from a given list of coordinates
and a list of atom types.  Additionally the atom\sphinxhyphen{}IDs, velocities, and
image flags may be provided.  If atom\sphinxhyphen{}IDs are not provided, they will be
automatically created as a sequence following the largest existing
atom\sphinxhyphen{}ID.

\sphinxAtStartPar
This function is useful to add atoms to a simulation or \sphinxhyphen{} in tandem with
{\hyperref[\detokenize{Library_properties:_CPPv416lammps_reset_boxPvPdPdddd}]{\sphinxcrossref{\sphinxcode{\sphinxupquote{lammps\_reset\_box()}}}}} \sphinxhyphen{} to restore a previously extracted and
saved state of a simulation.  Additional properties for the new atoms
can then be assigned via the {\hyperref[\detokenize{Library_scatter:_CPPv420lammps_scatter_atomsPvPKciiPv}]{\sphinxcrossref{\sphinxcode{\sphinxupquote{lammps\_scatter\_atoms()}}}}}
{\hyperref[\detokenize{Library_atoms:_CPPv419lammps_extract_atomPvPKc}]{\sphinxcrossref{\sphinxcode{\sphinxupquote{lammps\_extract\_atom()}}}}} functions.

\sphinxAtStartPar
For non\sphinxhyphen{}periodic boundaries, atoms will \sphinxstylestrong{not} be created that have
coordinates outside the box unless it is a shrink\sphinxhyphen{}wrap boundary and the
shrinkexceed flag has been set to a non\sphinxhyphen{}zero value.  For periodic
boundaries atoms will be wrapped back into the simulation cell and its
image flags adjusted accordingly, unless explicit image flags are
provided.

\sphinxAtStartPar
The function returns the number of atoms created or \sphinxhyphen{}1 on failure (e.g.,
when called before as box has been created).

\sphinxAtStartPar
Coordinates and velocities have to be given in a 1d\sphinxhyphen{}array in the order
X(1), Y(1), Z(1), X(2), Y(2), Z(2), …, X(N), Y(N), Z(N).

\begin{quote}\begin{description}
\sphinxlineitem{Parameters}\begin{itemize}
\item {} 
\sphinxAtStartPar
\sphinxstyleliteralstrong{\sphinxupquote{handle}} \textendash{} pointer to a previously created LAMMPS instance 

\item {} 
\sphinxAtStartPar
\sphinxstyleliteralstrong{\sphinxupquote{n}} \textendash{} number of atoms, N, to be added to the system 

\item {} 
\sphinxAtStartPar
\sphinxstyleliteralstrong{\sphinxupquote{id}} \textendash{} pointer to N atom IDs; \sphinxcode{\sphinxupquote{NULL}} will generate IDs 

\item {} 
\sphinxAtStartPar
\sphinxstyleliteralstrong{\sphinxupquote{type}} \textendash{} pointer to N atom types (required) 

\item {} 
\sphinxAtStartPar
\sphinxstyleliteralstrong{\sphinxupquote{x}} \textendash{} pointer to 3N doubles with x\sphinxhyphen{},y\sphinxhyphen{},z\sphinxhyphen{} positions of the new atoms (required) 

\item {} 
\sphinxAtStartPar
\sphinxstyleliteralstrong{\sphinxupquote{v}} \textendash{} pointer to 3N doubles with x\sphinxhyphen{},y\sphinxhyphen{},z\sphinxhyphen{} velocities of the new atoms (set to 0.0 if \sphinxcode{\sphinxupquote{NULL}}) 

\item {} 
\sphinxAtStartPar
\sphinxstyleliteralstrong{\sphinxupquote{image}} \textendash{} pointer to N imageint sets of image flags, or \sphinxcode{\sphinxupquote{NULL}}

\item {} 
\sphinxAtStartPar
\sphinxstyleliteralstrong{\sphinxupquote{bexpand}} \textendash{} if 1, atoms outside of shrink\sphinxhyphen{}wrap boundaries will still be created and not dropped and the box extended 

\end{itemize}

\sphinxlineitem{Returns}
\sphinxAtStartPar
number of atoms created on success; \sphinxhyphen{}1 on failure (no box, no atom IDs, etc.) 

\end{description}\end{quote}

\end{fulllineitems}


\sphinxstepscope


\subsection{Neighbor list access}
\label{\detokenize{Library_neighbor:neighbor-list-access}}\label{\detokenize{Library_neighbor::doc}}
\sphinxAtStartPar
The following functions enable access to neighbor lists generated by
LAMMPS or querying of their properties:
\begin{itemize}
\item {} 
\sphinxAtStartPar
{\hyperref[\detokenize{Library_neighbor:_CPPv429lammps_find_compute_neighlistPvPKci}]{\sphinxcrossref{\sphinxcode{\sphinxupquote{lammps\_find\_compute\_neighlist()}}}}}

\item {} 
\sphinxAtStartPar
{\hyperref[\detokenize{Library_neighbor:_CPPv425lammps_find_fix_neighlistPvPKci}]{\sphinxcrossref{\sphinxcode{\sphinxupquote{lammps\_find\_fix\_neighlist()}}}}}

\item {} 
\sphinxAtStartPar
{\hyperref[\detokenize{Library_neighbor:_CPPv426lammps_find_pair_neighlistPvPKciii}]{\sphinxcrossref{\sphinxcode{\sphinxupquote{lammps\_find\_pair\_neighlist()}}}}}

\item {} 
\sphinxAtStartPar
{\hyperref[\detokenize{Library_neighbor:_CPPv429lammps_neighlist_num_elementsPvi}]{\sphinxcrossref{\sphinxcode{\sphinxupquote{lammps\_neighlist\_num\_elements()}}}}}

\item {} 
\sphinxAtStartPar
{\hyperref[\detokenize{Library_neighbor:_CPPv434lammps_neighlist_element_neighborsPviiPiPiPPi}]{\sphinxcrossref{\sphinxcode{\sphinxupquote{lammps\_neighlist\_element\_neighbors()}}}}}

\end{itemize}


\bigskip\hrule\bigskip

\index{lammps\_find\_compute\_neighlist (C++ function)@\spxentry{lammps\_find\_compute\_neighlist}\spxextra{C++ function}}

\begin{fulllineitems}
\phantomsection\label{\detokenize{Library_neighbor:_CPPv429lammps_find_compute_neighlistPvPKci}}\phantomsection\label{\detokenize{Library_neighbor:_CPPv329lammps_find_compute_neighlistPvPKci}}\phantomsection\label{\detokenize{Library_neighbor:_CPPv229lammps_find_compute_neighlistPvPKci}}\phantomsection\label{\detokenize{Library_neighbor:lammps_find_compute_neighlist__voidP.cCP.i}}
\pysigstartsignatures
\pysigstartmultiline
\pysiglinewithargsret{\phantomsection\label{\detokenize{Library_neighbor:library_8h_1a04b80d34fe95f6931d9f6906a2c7f869}}\DUrole{kt}{int}\DUrole{w}{ }\sphinxbfcode{\sphinxupquote{\DUrole{n}{lammps\_find\_compute\_neighlist}}}}{\DUrole{kt}{void}\DUrole{w}{ }\DUrole{p}{*}\DUrole{n,sig-param}{handle}\sphinxparamcomma \DUrole{k}{const}\DUrole{w}{ }\DUrole{kt}{char}\DUrole{w}{ }\DUrole{p}{*}\DUrole{n,sig-param}{id}\sphinxparamcomma \DUrole{kt}{int}\DUrole{w}{ }\DUrole{n,sig-param}{request}}{}
\pysigstopmultiline
\pysigstopsignatures
\sphinxAtStartPar
Find index of a neighbor list requested by a compute

\sphinxAtStartPar
The neighbor list request from a compute is identified by the compute ID and the request ID. The request ID is typically 0, but will be \begin{quote}

\sphinxAtStartPar
0 in case a compute has multiple neighbor list requests. 
\end{quote}

\begin{quote}\begin{description}
\sphinxlineitem{Parameters}\begin{itemize}
\item {} 
\sphinxAtStartPar
\sphinxstyleliteralstrong{\sphinxupquote{handle}} \textendash{} pointer to a previously created LAMMPS instance cast to \sphinxcode{\sphinxupquote{void *}}. 

\item {} 
\sphinxAtStartPar
\sphinxstyleliteralstrong{\sphinxupquote{id}} \textendash{} Identifier of compute instance 

\item {} 
\sphinxAtStartPar
\sphinxstyleliteralstrong{\sphinxupquote{reqid}} \textendash{} request id to identify neighbor list in case there are multiple requests from the same compute 

\end{itemize}

\sphinxlineitem{Returns}
\sphinxAtStartPar
return neighbor list index if found, otherwise \sphinxhyphen{}1 

\end{description}\end{quote}

\end{fulllineitems}



\bigskip\hrule\bigskip

\index{lammps\_find\_fix\_neighlist (C++ function)@\spxentry{lammps\_find\_fix\_neighlist}\spxextra{C++ function}}

\begin{fulllineitems}
\phantomsection\label{\detokenize{Library_neighbor:_CPPv425lammps_find_fix_neighlistPvPKci}}\phantomsection\label{\detokenize{Library_neighbor:_CPPv325lammps_find_fix_neighlistPvPKci}}\phantomsection\label{\detokenize{Library_neighbor:_CPPv225lammps_find_fix_neighlistPvPKci}}\phantomsection\label{\detokenize{Library_neighbor:lammps_find_fix_neighlist__voidP.cCP.i}}
\pysigstartsignatures
\pysigstartmultiline
\pysiglinewithargsret{\phantomsection\label{\detokenize{Library_neighbor:library_8h_1a51652806cc5a984712a87de083396a48}}\DUrole{kt}{int}\DUrole{w}{ }\sphinxbfcode{\sphinxupquote{\DUrole{n}{lammps\_find\_fix\_neighlist}}}}{\DUrole{kt}{void}\DUrole{w}{ }\DUrole{p}{*}\DUrole{n,sig-param}{handle}\sphinxparamcomma \DUrole{k}{const}\DUrole{w}{ }\DUrole{kt}{char}\DUrole{w}{ }\DUrole{p}{*}\DUrole{n,sig-param}{id}\sphinxparamcomma \DUrole{kt}{int}\DUrole{w}{ }\DUrole{n,sig-param}{request}}{}
\pysigstopmultiline
\pysigstopsignatures
\sphinxAtStartPar
Find index of a neighbor list requested by a fix

\sphinxAtStartPar
The neighbor list request from a fix is identified by the fix ID and the request ID. The request ID is typically 0, but will be \textgreater{} 0 in case a fix has multiple neighbor list requests.
\begin{quote}\begin{description}
\sphinxlineitem{Parameters}\begin{itemize}
\item {} 
\sphinxAtStartPar
\sphinxstyleliteralstrong{\sphinxupquote{handle}} \textendash{} pointer to a previously created LAMMPS instance cast to \sphinxcode{\sphinxupquote{void *}}. 

\item {} 
\sphinxAtStartPar
\sphinxstyleliteralstrong{\sphinxupquote{id}} \textendash{} Identifier of fix instance 

\item {} 
\sphinxAtStartPar
\sphinxstyleliteralstrong{\sphinxupquote{reqid}} \textendash{} request id to identify neighbor list in case there are multiple requests from the same fix 

\end{itemize}

\sphinxlineitem{Returns}
\sphinxAtStartPar
return neighbor list index if found, otherwise \sphinxhyphen{}1 

\end{description}\end{quote}

\end{fulllineitems}



\bigskip\hrule\bigskip

\index{lammps\_find\_pair\_neighlist (C++ function)@\spxentry{lammps\_find\_pair\_neighlist}\spxextra{C++ function}}

\begin{fulllineitems}
\phantomsection\label{\detokenize{Library_neighbor:_CPPv426lammps_find_pair_neighlistPvPKciii}}\phantomsection\label{\detokenize{Library_neighbor:_CPPv326lammps_find_pair_neighlistPvPKciii}}\phantomsection\label{\detokenize{Library_neighbor:_CPPv226lammps_find_pair_neighlistPvPKciii}}\phantomsection\label{\detokenize{Library_neighbor:lammps_find_pair_neighlist__voidP.cCP.i.i.i}}
\pysigstartsignatures
\pysigstartmultiline
\pysiglinewithargsret{\phantomsection\label{\detokenize{Library_neighbor:library_8h_1a859b206fea42bc2db03c992acc478cbc}}\DUrole{kt}{int}\DUrole{w}{ }\sphinxbfcode{\sphinxupquote{\DUrole{n}{lammps\_find\_pair\_neighlist}}}}{\DUrole{kt}{void}\DUrole{w}{ }\DUrole{p}{*}\DUrole{n,sig-param}{handle}\sphinxparamcomma \DUrole{k}{const}\DUrole{w}{ }\DUrole{kt}{char}\DUrole{w}{ }\DUrole{p}{*}\DUrole{n,sig-param}{style}\sphinxparamcomma \DUrole{kt}{int}\DUrole{w}{ }\DUrole{n,sig-param}{exact}\sphinxparamcomma \DUrole{kt}{int}\DUrole{w}{ }\DUrole{n,sig-param}{nsub}\sphinxparamcomma \DUrole{kt}{int}\DUrole{w}{ }\DUrole{n,sig-param}{request}}{}
\pysigstopmultiline
\pysigstopsignatures
\sphinxAtStartPar
Find index of a neighbor list requested by a pair style

\sphinxAtStartPar
This function determines which of the available neighbor lists for pair styles matches the given conditions. It first matches the style name. If exact is 1 the name must match exactly, if exact is 0, a regular expression or sub\sphinxhyphen{}string match is done. If the pair style is hybrid or hybrid/overlay the style is matched against the sub styles instead. If a the same pair style is used multiple times as a sub\sphinxhyphen{}style, the nsub argument must be \textgreater{} 0 and represents the nth instance of the sub\sphinxhyphen{}style (same as for the pair\_coeff command, for example). In that case nsub=0 will not produce a match and this function will return \sphinxhyphen{}1.

\sphinxAtStartPar
The final condition to be checked is the request ID (reqid). This will normally be 0, but some pair styles request multiple neighbor lists and set the request ID to a value \textgreater{} 0.
\begin{quote}\begin{description}
\sphinxlineitem{Parameters}\begin{itemize}
\item {} 
\sphinxAtStartPar
\sphinxstyleliteralstrong{\sphinxupquote{handle}} \textendash{} pointer to a previously created LAMMPS instance cast to \sphinxcode{\sphinxupquote{void *}}. 

\item {} 
\sphinxAtStartPar
\sphinxstyleliteralstrong{\sphinxupquote{style}} \textendash{} String used to search for pair style instance 

\item {} 
\sphinxAtStartPar
\sphinxstyleliteralstrong{\sphinxupquote{exact}} \textendash{} Flag to control whether style should match exactly or only a regular expression / sub\sphinxhyphen{}string match is applied. 

\item {} 
\sphinxAtStartPar
\sphinxstyleliteralstrong{\sphinxupquote{nsub}} \textendash{} match nsub\sphinxhyphen{}th hybrid sub\sphinxhyphen{}style instance of the same style 

\item {} 
\sphinxAtStartPar
\sphinxstyleliteralstrong{\sphinxupquote{reqid}} \textendash{} request id to identify neighbor list in case there are multiple requests from the same pair style instance 

\end{itemize}

\sphinxlineitem{Returns}
\sphinxAtStartPar
return neighbor list index if found, otherwise \sphinxhyphen{}1 

\end{description}\end{quote}

\end{fulllineitems}



\bigskip\hrule\bigskip

\index{lammps\_neighlist\_num\_elements (C++ function)@\spxentry{lammps\_neighlist\_num\_elements}\spxextra{C++ function}}

\begin{fulllineitems}
\phantomsection\label{\detokenize{Library_neighbor:_CPPv429lammps_neighlist_num_elementsPvi}}\phantomsection\label{\detokenize{Library_neighbor:_CPPv329lammps_neighlist_num_elementsPvi}}\phantomsection\label{\detokenize{Library_neighbor:_CPPv229lammps_neighlist_num_elementsPvi}}\phantomsection\label{\detokenize{Library_neighbor:lammps_neighlist_num_elements__voidP.i}}
\pysigstartsignatures
\pysigstartmultiline
\pysiglinewithargsret{\phantomsection\label{\detokenize{Library_neighbor:library_8h_1ab4cd0310c0638fad6352d0081194bdd8}}\DUrole{kt}{int}\DUrole{w}{ }\sphinxbfcode{\sphinxupquote{\DUrole{n}{lammps\_neighlist\_num\_elements}}}}{\DUrole{kt}{void}\DUrole{w}{ }\DUrole{p}{*}\DUrole{n,sig-param}{handle}\sphinxparamcomma \DUrole{kt}{int}\DUrole{w}{ }\DUrole{n,sig-param}{idx}}{}
\pysigstopmultiline
\pysigstopsignatures
\sphinxAtStartPar
Return the number of entries in the neighbor list with given index
\begin{quote}\begin{description}
\sphinxlineitem{Parameters}\begin{itemize}
\item {} 
\sphinxAtStartPar
\sphinxstyleliteralstrong{\sphinxupquote{handle}} \textendash{} pointer to a previously created LAMMPS instance cast to \sphinxcode{\sphinxupquote{void *}}. 

\item {} 
\sphinxAtStartPar
\sphinxstyleliteralstrong{\sphinxupquote{idx}} \textendash{} neighbor list index 

\end{itemize}

\sphinxlineitem{Returns}
\sphinxAtStartPar
return number of entries in neighbor list, \sphinxhyphen{}1 if idx is not a valid index 

\end{description}\end{quote}

\end{fulllineitems}



\bigskip\hrule\bigskip

\index{lammps\_neighlist\_element\_neighbors (C++ function)@\spxentry{lammps\_neighlist\_element\_neighbors}\spxextra{C++ function}}

\begin{fulllineitems}
\phantomsection\label{\detokenize{Library_neighbor:_CPPv434lammps_neighlist_element_neighborsPviiPiPiPPi}}\phantomsection\label{\detokenize{Library_neighbor:_CPPv334lammps_neighlist_element_neighborsPviiPiPiPPi}}\phantomsection\label{\detokenize{Library_neighbor:_CPPv234lammps_neighlist_element_neighborsPviiPiPiPPi}}\phantomsection\label{\detokenize{Library_neighbor:lammps_neighlist_element_neighbors__voidP.i.i.iP.iP.iPP}}
\pysigstartsignatures
\pysigstartmultiline
\pysiglinewithargsret{\phantomsection\label{\detokenize{Library_neighbor:library_8h_1a0bd03926cf3a3ebb93311dab400c792c}}\DUrole{kt}{void}\DUrole{w}{ }\sphinxbfcode{\sphinxupquote{\DUrole{n}{lammps\_neighlist\_element\_neighbors}}}}{\DUrole{kt}{void}\DUrole{w}{ }\DUrole{p}{*}\DUrole{n,sig-param}{handle}\sphinxparamcomma \DUrole{kt}{int}\DUrole{w}{ }\DUrole{n,sig-param}{idx}\sphinxparamcomma \DUrole{kt}{int}\DUrole{w}{ }\DUrole{n,sig-param}{element}\sphinxparamcomma \DUrole{kt}{int}\DUrole{w}{ }\DUrole{p}{*}\DUrole{n,sig-param}{iatom}\sphinxparamcomma \DUrole{kt}{int}\DUrole{w}{ }\DUrole{p}{*}\DUrole{n,sig-param}{numneigh}\sphinxparamcomma \DUrole{kt}{int}\DUrole{w}{ }\DUrole{p}{*}\DUrole{p}{*}\DUrole{n,sig-param}{neighbors}}{}
\pysigstopmultiline
\pysigstopsignatures
\sphinxAtStartPar
Return atom local index, number of neighbors, and array of neighbor local atom indices of neighbor list entry
\begin{quote}\begin{description}
\sphinxlineitem{Parameters}\begin{itemize}
\item {} 
\sphinxAtStartPar
\sphinxstyleliteralstrong{\sphinxupquote{handle}} \textendash{} pointer to a previously created LAMMPS instance cast to \sphinxcode{\sphinxupquote{void *}}. 

\item {} 
\sphinxAtStartPar
\sphinxstyleliteralstrong{\sphinxupquote{idx}} \textendash{} index of this neighbor list in the list of all neighbor lists 

\item {} 
\sphinxAtStartPar
\sphinxstyleliteralstrong{\sphinxupquote{element}} \textendash{} index of this neighbor list entry 

\item {} 
\sphinxAtStartPar
\sphinxstyleliteralstrong{\sphinxupquote{iatom}} \textendash{} \sphinxstylestrong{{[}out{]}} local atom index (i.e. in the range {[}0, nlocal + nghost), \sphinxhyphen{}1 if invalid idx or element value 

\item {} 
\sphinxAtStartPar
\sphinxstyleliteralstrong{\sphinxupquote{numneigh}} \textendash{} \sphinxstylestrong{{[}out{]}} number of neighbors of atom iatom or 0 

\item {} 
\sphinxAtStartPar
\sphinxstyleliteralstrong{\sphinxupquote{neighbors}} \textendash{} \sphinxstylestrong{{[}out{]}} pointer to array of neighbor atom local indices or NULL 

\end{itemize}

\end{description}\end{quote}

\end{fulllineitems}


\sphinxstepscope


\subsection{Configuration information}
\label{\detokenize{Library_config:configuration-information}}\label{\detokenize{Library_config::doc}}
\sphinxAtStartPar
This section documents the following functions:
\begin{itemize}
\item {} 
\sphinxAtStartPar
{\hyperref[\detokenize{Library_config:_CPPv414lammps_versionPv}]{\sphinxcrossref{\sphinxcode{\sphinxupquote{lammps\_version()}}}}}

\item {} 
\sphinxAtStartPar
{\hyperref[\detokenize{Library_config:_CPPv418lammps_get_os_infoPci}]{\sphinxcrossref{\sphinxcode{\sphinxupquote{lammps\_get\_os\_info()}}}}}

\item {} 
\sphinxAtStartPar
{\hyperref[\detokenize{Library_config:_CPPv429lammps_config_has_mpi_supportv}]{\sphinxcrossref{\sphinxcode{\sphinxupquote{lammps\_config\_has\_mpi\_support()}}}}}

\item {} 
\sphinxAtStartPar
{\hyperref[\detokenize{Library_config:_CPPv430lammps_config_has_gzip_supportv}]{\sphinxcrossref{\sphinxcode{\sphinxupquote{lammps\_config\_has\_gzip\_support()}}}}}

\item {} 
\sphinxAtStartPar
{\hyperref[\detokenize{Library_config:_CPPv429lammps_config_has_png_supportv}]{\sphinxcrossref{\sphinxcode{\sphinxupquote{lammps\_config\_has\_png\_support()}}}}}

\item {} 
\sphinxAtStartPar
{\hyperref[\detokenize{Library_config:_CPPv430lammps_config_has_jpeg_supportv}]{\sphinxcrossref{\sphinxcode{\sphinxupquote{lammps\_config\_has\_jpeg\_support()}}}}}

\item {} 
\sphinxAtStartPar
{\hyperref[\detokenize{Library_config:_CPPv432lammps_config_has_ffmpeg_supportv}]{\sphinxcrossref{\sphinxcode{\sphinxupquote{lammps\_config\_has\_ffmpeg\_support()}}}}}

\item {} 
\sphinxAtStartPar
{\hyperref[\detokenize{Library_config:_CPPv428lammps_config_has_exceptionsv}]{\sphinxcrossref{\sphinxcode{\sphinxupquote{lammps\_config\_has\_exceptions()}}}}}

\item {} 
\sphinxAtStartPar
{\hyperref[\detokenize{Library_config:_CPPv425lammps_config_has_packagePKc}]{\sphinxcrossref{\sphinxcode{\sphinxupquote{lammps\_config\_has\_package()}}}}}

\item {} 
\sphinxAtStartPar
{\hyperref[\detokenize{Library_config:_CPPv427lammps_config_package_countv}]{\sphinxcrossref{\sphinxcode{\sphinxupquote{lammps\_config\_package\_count()}}}}}

\item {} 
\sphinxAtStartPar
{\hyperref[\detokenize{Library_config:_CPPv426lammps_config_package_nameiPci}]{\sphinxcrossref{\sphinxcode{\sphinxupquote{lammps\_config\_package\_name()}}}}}

\item {} 
\sphinxAtStartPar
{\hyperref[\detokenize{Library_config:_CPPv425lammps_config_acceleratorPKcPKcPKc}]{\sphinxcrossref{\sphinxcode{\sphinxupquote{lammps\_config\_accelerator()}}}}}

\item {} 
\sphinxAtStartPar
{\hyperref[\detokenize{Library_config:_CPPv421lammps_has_gpu_devicev}]{\sphinxcrossref{\sphinxcode{\sphinxupquote{lammps\_has\_gpu\_device()}}}}}

\item {} 
\sphinxAtStartPar
\sphinxcode{\sphinxupquote{lammps\_gpu\_device\_info()}}

\item {} 
\sphinxAtStartPar
{\hyperref[\detokenize{Library_config:_CPPv416lammps_has_stylePvPKcPKc}]{\sphinxcrossref{\sphinxcode{\sphinxupquote{lammps\_has\_style()}}}}}

\item {} 
\sphinxAtStartPar
{\hyperref[\detokenize{Library_config:_CPPv418lammps_style_countPvPKc}]{\sphinxcrossref{\sphinxcode{\sphinxupquote{lammps\_style\_count()}}}}}

\item {} 
\sphinxAtStartPar
{\hyperref[\detokenize{Library_config:_CPPv417lammps_style_namePvPKciPci}]{\sphinxcrossref{\sphinxcode{\sphinxupquote{lammps\_style\_name()}}}}}

\item {} 
\sphinxAtStartPar
{\hyperref[\detokenize{Library_config:_CPPv413lammps_has_idPvPKcPKc}]{\sphinxcrossref{\sphinxcode{\sphinxupquote{lammps\_has\_id()}}}}}

\item {} 
\sphinxAtStartPar
{\hyperref[\detokenize{Library_config:_CPPv415lammps_id_countPvPKc}]{\sphinxcrossref{\sphinxcode{\sphinxupquote{lammps\_id\_count()}}}}}

\item {} 
\sphinxAtStartPar
{\hyperref[\detokenize{Library_config:_CPPv414lammps_id_namePvPKciPci}]{\sphinxcrossref{\sphinxcode{\sphinxupquote{lammps\_id\_name()}}}}}

\end{itemize}


\bigskip\hrule\bigskip


\sphinxAtStartPar
These library functions can be used to query the LAMMPS library for
compile time settings and included packages and styles.  This enables
programs that use the library interface to determine whether the
linked LAMMPS library is compatible with the requirements of the
application without crashing during the LAMMPS functions (e.g. due to
missing pair styles from packages) or to choose between different
options (e.g. whether to use \sphinxcode{\sphinxupquote{lj/cut}}, \sphinxcode{\sphinxupquote{lj/cut/opt}},
\sphinxcode{\sphinxupquote{lj/cut/omp}} or \sphinxcode{\sphinxupquote{lj/cut/intel}}).  Most of the functions can be
called directly without first creating a LAMMPS instance.  While
crashes within LAMMPS may be recovered from by enabling
{\hyperref[\detokenize{Build_settings:exceptions}]{\sphinxcrossref{\DUrole{std,std-ref}{exceptions}}}}, avoiding them proactively is a safer
approach.
\sphinxSetupCaptionForVerbatim{Example for using configuration settings functions}
\def\sphinxLiteralBlockLabel{\label{\detokenize{Library_config:id1}}}
\begin{sphinxVerbatim}[commandchars=\\\{\}]
\PYG{c+cp}{\PYGZsh{}}\PYG{c+cp}{include}\PYG{+w}{ }\PYG{c+cpf}{\PYGZdq{}library.h\PYGZdq{}}
\PYG{c+cp}{\PYGZsh{}}\PYG{c+cp}{include}\PYG{+w}{ }\PYG{c+cpf}{\PYGZlt{}stdio.h\PYGZgt{}}

\PYG{k+kt}{int}\PYG{+w}{ }\PYG{n+nf}{main}\PYG{p}{(}\PYG{k+kt}{int}\PYG{+w}{ }\PYG{n}{argc}\PYG{p}{,}\PYG{+w}{ }\PYG{k+kt}{char}\PYG{+w}{ }\PYG{o}{*}\PYG{o}{*}\PYG{n}{argv}\PYG{p}{)}
\PYG{p}{\PYGZob{}}
\PYG{+w}{  }\PYG{k+kt}{void}\PYG{+w}{ }\PYG{o}{*}\PYG{n}{handle}\PYG{p}{;}

\PYG{+w}{  }\PYG{n}{handle}\PYG{+w}{ }\PYG{o}{=}\PYG{+w}{ }\PYG{n}{lammps\PYGZus{}open\PYGZus{}no\PYGZus{}mpi}\PYG{p}{(}\PYG{l+m+mi}{0}\PYG{p}{,}\PYG{+w}{ }\PYG{n+nb}{NULL}\PYG{p}{,}\PYG{+w}{ }\PYG{n+nb}{NULL}\PYG{p}{)}\PYG{p}{;}
\PYG{+w}{  }\PYG{n}{lammps\PYGZus{}file}\PYG{p}{(}\PYG{n}{handle}\PYG{p}{,}\PYG{+w}{ }\PYG{l+s}{\PYGZdq{}}\PYG{l+s}{in.missing}\PYG{l+s}{\PYGZdq{}}\PYG{p}{)}\PYG{p}{;}
\PYG{+w}{  }\PYG{k}{if}\PYG{+w}{ }\PYG{p}{(}\PYG{n}{lammps\PYGZus{}has\PYGZus{}error}\PYG{p}{(}\PYG{n}{handle}\PYG{p}{)}\PYG{p}{)}\PYG{+w}{ }\PYG{p}{\PYGZob{}}
\PYG{+w}{    }\PYG{k+kt}{char}\PYG{+w}{ }\PYG{n}{errmsg}\PYG{p}{[}\PYG{l+m+mi}{256}\PYG{p}{]}\PYG{p}{;}
\PYG{+w}{    }\PYG{k+kt}{int}\PYG{+w}{ }\PYG{n}{errtype}\PYG{p}{;}
\PYG{+w}{    }\PYG{n}{errtype}\PYG{+w}{ }\PYG{o}{=}\PYG{+w}{ }\PYG{n}{lammps\PYGZus{}get\PYGZus{}last\PYGZus{}error\PYGZus{}message}\PYG{p}{(}\PYG{n}{handle}\PYG{p}{,}\PYG{+w}{ }\PYG{n}{errmsg}\PYG{p}{,}\PYG{+w}{ }\PYG{l+m+mi}{256}\PYG{p}{)}\PYG{p}{;}
\PYG{+w}{    }\PYG{n}{fprintf}\PYG{p}{(}\PYG{n}{stderr}\PYG{p}{,}\PYG{+w}{ }\PYG{l+s}{\PYGZdq{}}\PYG{l+s}{LAMMPS failed with error: \PYGZpc{}s}\PYG{l+s+se}{\PYGZbs{}n}\PYG{l+s}{\PYGZdq{}}\PYG{p}{,}\PYG{+w}{ }\PYG{n}{errmsg}\PYG{p}{)}\PYG{p}{;}
\PYG{+w}{    }\PYG{k}{return}\PYG{+w}{ }\PYG{l+m+mi}{1}\PYG{p}{;}
\PYG{+w}{  }\PYG{p}{\PYGZcb{}}
\PYG{+w}{  }\PYG{c+cm}{/* write compressed dump file depending on available of options */}
\PYG{+w}{  }\PYG{k}{if}\PYG{+w}{ }\PYG{p}{(}\PYG{n}{lammps\PYGZus{}has\PYGZus{}style}\PYG{p}{(}\PYG{n}{handle}\PYG{p}{,}\PYG{+w}{ }\PYG{l+s}{\PYGZdq{}}\PYG{l+s}{dump}\PYG{l+s}{\PYGZdq{}}\PYG{p}{,}\PYG{+w}{ }\PYG{l+s}{\PYGZdq{}}\PYG{l+s}{atom/zstd}\PYG{l+s}{\PYGZdq{}}\PYG{p}{)}\PYG{p}{)}\PYG{+w}{ }\PYG{p}{\PYGZob{}}
\PYG{+w}{    }\PYG{n}{lammps\PYGZus{}command}\PYG{p}{(}\PYG{n}{handle}\PYG{p}{,}\PYG{+w}{ }\PYG{l+s}{\PYGZdq{}}\PYG{l+s}{dump d1 all atom/zstd 100 dump.zst}\PYG{l+s}{\PYGZdq{}}\PYG{p}{)}\PYG{p}{;}
\PYG{+w}{  }\PYG{p}{\PYGZcb{}}\PYG{+w}{ }\PYG{k}{else}\PYG{+w}{ }\PYG{k}{if}\PYG{+w}{ }\PYG{p}{(}\PYG{n}{lammps\PYGZus{}has\PYGZus{}style}\PYG{p}{(}\PYG{n}{handle}\PYG{p}{,}\PYG{+w}{ }\PYG{l+s}{\PYGZdq{}}\PYG{l+s}{dump}\PYG{l+s}{\PYGZdq{}}\PYG{p}{,}\PYG{+w}{ }\PYG{l+s}{\PYGZdq{}}\PYG{l+s}{atom/gz}\PYG{l+s}{\PYGZdq{}}\PYG{p}{)}\PYG{p}{)}\PYG{+w}{ }\PYG{p}{\PYGZob{}}
\PYG{+w}{    }\PYG{n}{lammps\PYGZus{}command}\PYG{p}{(}\PYG{n}{handle}\PYG{p}{,}\PYG{+w}{ }\PYG{l+s}{\PYGZdq{}}\PYG{l+s}{dump d1 all atom/gz 100 dump.gz}\PYG{l+s}{\PYGZdq{}}\PYG{p}{)}\PYG{p}{;}
\PYG{+w}{  }\PYG{p}{\PYGZcb{}}\PYG{+w}{ }\PYG{k}{else}\PYG{+w}{ }\PYG{k}{if}\PYG{+w}{ }\PYG{p}{(}\PYG{n}{lammps\PYGZus{}config\PYGZus{}has\PYGZus{}gzip\PYGZus{}support}\PYG{p}{(}\PYG{p}{)}\PYG{p}{)}\PYG{+w}{ }\PYG{p}{\PYGZob{}}
\PYG{+w}{    }\PYG{n}{lammps\PYGZus{}command}\PYG{p}{(}\PYG{n}{handle}\PYG{p}{,}\PYG{+w}{ }\PYG{l+s}{\PYGZdq{}}\PYG{l+s}{dump d1 all atom 100 dump.gz}\PYG{l+s}{\PYGZdq{}}\PYG{p}{)}\PYG{p}{;}
\PYG{+w}{  }\PYG{p}{\PYGZcb{}}\PYG{+w}{ }\PYG{k}{else}\PYG{+w}{ }\PYG{p}{\PYGZob{}}
\PYG{+w}{    }\PYG{n}{lammps\PYGZus{}command}\PYG{p}{(}\PYG{n}{handle}\PYG{p}{,}\PYG{+w}{ }\PYG{l+s}{\PYGZdq{}}\PYG{l+s}{dump d1 all atom 100 dump}\PYG{l+s}{\PYGZdq{}}\PYG{p}{)}\PYG{p}{;}
\PYG{+w}{  }\PYG{p}{\PYGZcb{}}
\PYG{+w}{  }\PYG{n}{lammps\PYGZus{}close}\PYG{p}{(}\PYG{n}{handle}\PYG{p}{)}\PYG{p}{;}
\PYG{+w}{  }\PYG{k}{return}\PYG{+w}{ }\PYG{l+m+mi}{0}\PYG{p}{;}
\PYG{p}{\PYGZcb{}}
\end{sphinxVerbatim}


\bigskip\hrule\bigskip

\index{lammps\_version (C++ function)@\spxentry{lammps\_version}\spxextra{C++ function}}

\begin{fulllineitems}
\phantomsection\label{\detokenize{Library_config:_CPPv414lammps_versionPv}}\phantomsection\label{\detokenize{Library_config:_CPPv314lammps_versionPv}}\phantomsection\label{\detokenize{Library_config:_CPPv214lammps_versionPv}}\phantomsection\label{\detokenize{Library_config:lammps_version__voidP}}
\pysigstartsignatures
\pysigstartmultiline
\pysiglinewithargsret{\phantomsection\label{\detokenize{Library_config:library_8h_1a41f44a4eef975f8a189d3186259a8ef9}}\DUrole{kt}{int}\DUrole{w}{ }\sphinxbfcode{\sphinxupquote{\DUrole{n}{lammps\_version}}}}{\DUrole{kt}{void}\DUrole{w}{ }\DUrole{p}{*}\DUrole{n,sig-param}{handle}}{}
\pysigstopmultiline
\pysigstopsignatures
\sphinxAtStartPar
Get numerical representation of the LAMMPS version date.

\sphinxAtStartPar
The {\hyperref[\detokenize{Library_config:_CPPv414lammps_versionPv}]{\sphinxcrossref{\sphinxcode{\sphinxupquote{lammps\_version()}}}}} function returns an integer representing
the version of the LAMMPS code in the format YYYYMMDD.  This can be used
to implement backward compatibility in software using the LAMMPS library
interface.  The specific format guarantees, that this version number is
growing with every new LAMMPS release.
\begin{quote}\begin{description}
\sphinxlineitem{Parameters}
\sphinxAtStartPar
\sphinxstyleliteralstrong{\sphinxupquote{handle}} \textendash{} pointer to a previously created LAMMPS instance 

\sphinxlineitem{Returns}
\sphinxAtStartPar
an integer representing the version data in the format YYYYMMDD 

\end{description}\end{quote}

\end{fulllineitems}



\bigskip\hrule\bigskip

\index{lammps\_get\_os\_info (C++ function)@\spxentry{lammps\_get\_os\_info}\spxextra{C++ function}}

\begin{fulllineitems}
\phantomsection\label{\detokenize{Library_config:_CPPv418lammps_get_os_infoPci}}\phantomsection\label{\detokenize{Library_config:_CPPv318lammps_get_os_infoPci}}\phantomsection\label{\detokenize{Library_config:_CPPv218lammps_get_os_infoPci}}\phantomsection\label{\detokenize{Library_config:lammps_get_os_info__cP.i}}
\pysigstartsignatures
\pysigstartmultiline
\pysiglinewithargsret{\phantomsection\label{\detokenize{Library_config:library_8h_1ae4b6178f693ff83002ee8a1559edd6ed}}\DUrole{kt}{void}\DUrole{w}{ }\sphinxbfcode{\sphinxupquote{\DUrole{n}{lammps\_get\_os\_info}}}}{\DUrole{kt}{char}\DUrole{w}{ }\DUrole{p}{*}\DUrole{n,sig-param}{buffer}\sphinxparamcomma \DUrole{kt}{int}\DUrole{w}{ }\DUrole{n,sig-param}{buf\_size}}{}
\pysigstopmultiline
\pysigstopsignatures
\sphinxAtStartPar
Get operating system and architecture information

\sphinxAtStartPar

\sphinxAtStartPar
\DUrole{versionmodified,added}{Added in version 9Oct2020.}

\sphinxAtStartPar
The {\hyperref[\detokenize{Library_config:_CPPv418lammps_get_os_infoPci}]{\sphinxcrossref{\sphinxcode{\sphinxupquote{lammps\_get\_os\_info()}}}}} function can be used to retrieve
detailed information about the hosting operating system and
compiler/runtime.

\sphinxAtStartPar
A suitable buffer for a C\sphinxhyphen{}style string has to be provided and its length.
The assembled text will be truncated to not overflow this buffer. The
string is typically a few hundred bytes long.

\begin{quote}\begin{description}
\sphinxlineitem{Parameters}\begin{itemize}
\item {} 
\sphinxAtStartPar
\sphinxstyleliteralstrong{\sphinxupquote{buffer}} \textendash{} string buffer to copy the information to 

\item {} 
\sphinxAtStartPar
\sphinxstyleliteralstrong{\sphinxupquote{buf\_size}} \textendash{} size of the provided string buffer 

\end{itemize}

\end{description}\end{quote}

\end{fulllineitems}



\bigskip\hrule\bigskip

\index{lammps\_config\_has\_mpi\_support (C++ function)@\spxentry{lammps\_config\_has\_mpi\_support}\spxextra{C++ function}}

\begin{fulllineitems}
\phantomsection\label{\detokenize{Library_config:_CPPv429lammps_config_has_mpi_supportv}}\phantomsection\label{\detokenize{Library_config:_CPPv329lammps_config_has_mpi_supportv}}\phantomsection\label{\detokenize{Library_config:_CPPv229lammps_config_has_mpi_supportv}}\phantomsection\label{\detokenize{Library_config:lammps_config_has_mpi_support}}
\pysigstartsignatures
\pysigstartmultiline
\pysiglinewithargsret{\phantomsection\label{\detokenize{Library_config:library_8h_1ada45e7d079fddf8400e31cc98710b5f8}}\DUrole{kt}{int}\DUrole{w}{ }\sphinxbfcode{\sphinxupquote{\DUrole{n}{lammps\_config\_has\_mpi\_support}}}}{}{}
\pysigstopmultiline
\pysigstopsignatures
\sphinxAtStartPar
This function is used to query whether LAMMPS was compiled with a real MPI library or in serial. For the real MPI library it reports the size of the MPI communicator in bytes (4 or 8), which allows to check for compatibility with a hosting code.
\begin{quote}\begin{description}
\sphinxlineitem{Returns}
\sphinxAtStartPar
0 when compiled with MPI STUBS, otherwise the MPI\_Comm size in bytes 

\end{description}\end{quote}

\end{fulllineitems}



\bigskip\hrule\bigskip

\index{lammps\_config\_has\_gzip\_support (C++ function)@\spxentry{lammps\_config\_has\_gzip\_support}\spxextra{C++ function}}

\begin{fulllineitems}
\phantomsection\label{\detokenize{Library_config:_CPPv430lammps_config_has_gzip_supportv}}\phantomsection\label{\detokenize{Library_config:_CPPv330lammps_config_has_gzip_supportv}}\phantomsection\label{\detokenize{Library_config:_CPPv230lammps_config_has_gzip_supportv}}\phantomsection\label{\detokenize{Library_config:lammps_config_has_gzip_support}}
\pysigstartsignatures
\pysigstartmultiline
\pysiglinewithargsret{\phantomsection\label{\detokenize{Library_config:library_8h_1a0bd41d39a5dd5992d614e150b6d88a5f}}\DUrole{kt}{int}\DUrole{w}{ }\sphinxbfcode{\sphinxupquote{\DUrole{n}{lammps\_config\_has\_gzip\_support}}}}{}{}
\pysigstopmultiline
\pysigstopsignatures
\sphinxAtStartPar
Check if the LAMMPS library supports reading or writing compressed files via a pipe to gzip or similar compression programs

\sphinxAtStartPar
Several LAMMPS commands (e.g., {\hyperref[\detokenize{read_data::doc}]{\sphinxcrossref{\DUrole{doc}{read\_data command}}}}, {\hyperref[\detokenize{write_data::doc}]{\sphinxcrossref{\DUrole{doc}{write\_data command}}}},
{\hyperref[\detokenize{dump::doc}]{\sphinxcrossref{\DUrole{doc}{dump styles atom, custom, and xyz}}}}) support reading and
writing compressed files via creating a pipe to the \sphinxcode{\sphinxupquote{gzip}} program.
This function checks whether this feature was {\hyperref[\detokenize{Build_settings:gzip}]{\sphinxcrossref{\DUrole{std,std-ref}{enabled at compile
time}}}}. It does \sphinxstylestrong{not} check whether\textasciigrave{}\textasciigrave{}gzip\textasciigrave{}\textasciigrave{} or any other
supported compression programs themselves are installed and usable.
\begin{quote}\begin{description}
\sphinxlineitem{Returns}
\sphinxAtStartPar
1 if yes, otherwise 0 

\end{description}\end{quote}

\end{fulllineitems}



\bigskip\hrule\bigskip

\index{lammps\_config\_has\_png\_support (C++ function)@\spxentry{lammps\_config\_has\_png\_support}\spxextra{C++ function}}

\begin{fulllineitems}
\phantomsection\label{\detokenize{Library_config:_CPPv429lammps_config_has_png_supportv}}\phantomsection\label{\detokenize{Library_config:_CPPv329lammps_config_has_png_supportv}}\phantomsection\label{\detokenize{Library_config:_CPPv229lammps_config_has_png_supportv}}\phantomsection\label{\detokenize{Library_config:lammps_config_has_png_support}}
\pysigstartsignatures
\pysigstartmultiline
\pysiglinewithargsret{\phantomsection\label{\detokenize{Library_config:library_8h_1a9de61d7686a3db2344fb94a4a22c67e4}}\DUrole{kt}{int}\DUrole{w}{ }\sphinxbfcode{\sphinxupquote{\DUrole{n}{lammps\_config\_has\_png\_support}}}}{}{}
\pysigstopmultiline
\pysigstopsignatures
\sphinxAtStartPar
Check if the LAMMPS library supports writing PNG format images

\sphinxAtStartPar
The LAMMPS {\hyperref[\detokenize{dump_image::doc}]{\sphinxcrossref{\DUrole{doc}{dump style image}}}} supports writing multiple
image file formats.  Most of them, however, need support from an external
library, and using that has to be {\hyperref[\detokenize{Build_settings:graphics}]{\sphinxcrossref{\DUrole{std,std-ref}{enabled at compile time}}}}.
This function checks whether support for the \sphinxhref{https://en.wikipedia.org/wiki/Portable\_Network\_Graphics}{PNG image file format} is available
in the current LAMMPS library.
\begin{quote}\begin{description}
\sphinxlineitem{Returns}
\sphinxAtStartPar
1 if yes, otherwise 0 

\end{description}\end{quote}

\end{fulllineitems}



\bigskip\hrule\bigskip

\index{lammps\_config\_has\_jpeg\_support (C++ function)@\spxentry{lammps\_config\_has\_jpeg\_support}\spxextra{C++ function}}

\begin{fulllineitems}
\phantomsection\label{\detokenize{Library_config:_CPPv430lammps_config_has_jpeg_supportv}}\phantomsection\label{\detokenize{Library_config:_CPPv330lammps_config_has_jpeg_supportv}}\phantomsection\label{\detokenize{Library_config:_CPPv230lammps_config_has_jpeg_supportv}}\phantomsection\label{\detokenize{Library_config:lammps_config_has_jpeg_support}}
\pysigstartsignatures
\pysigstartmultiline
\pysiglinewithargsret{\phantomsection\label{\detokenize{Library_config:library_8h_1ac3835bbc1c6797e56f8aa47443d11af6}}\DUrole{kt}{int}\DUrole{w}{ }\sphinxbfcode{\sphinxupquote{\DUrole{n}{lammps\_config\_has\_jpeg\_support}}}}{}{}
\pysigstopmultiline
\pysigstopsignatures
\sphinxAtStartPar
Check if the LAMMPS library supports writing JPEG format images

\sphinxAtStartPar
The LAMMPS {\hyperref[\detokenize{dump_image::doc}]{\sphinxcrossref{\DUrole{doc}{dump style image}}}} supports writing multiple
image file formats.  Most of them, however, need support from an external
library, and using that has to be {\hyperref[\detokenize{Build_settings:graphics}]{\sphinxcrossref{\DUrole{std,std-ref}{enabled at compile time}}}}.
This function checks whether support for the \sphinxhref{https://jpeg.org/jpeg/}{JPEG image file format} is available in the current LAMMPS library.
\begin{quote}\begin{description}
\sphinxlineitem{Returns}
\sphinxAtStartPar
1 if yes, otherwise 0 

\end{description}\end{quote}

\end{fulllineitems}



\bigskip\hrule\bigskip

\index{lammps\_config\_has\_ffmpeg\_support (C++ function)@\spxentry{lammps\_config\_has\_ffmpeg\_support}\spxextra{C++ function}}

\begin{fulllineitems}
\phantomsection\label{\detokenize{Library_config:_CPPv432lammps_config_has_ffmpeg_supportv}}\phantomsection\label{\detokenize{Library_config:_CPPv332lammps_config_has_ffmpeg_supportv}}\phantomsection\label{\detokenize{Library_config:_CPPv232lammps_config_has_ffmpeg_supportv}}\phantomsection\label{\detokenize{Library_config:lammps_config_has_ffmpeg_support}}
\pysigstartsignatures
\pysigstartmultiline
\pysiglinewithargsret{\phantomsection\label{\detokenize{Library_config:library_8h_1acbbc33bf81d78002a4ed94e2b7ff8f06}}\DUrole{kt}{int}\DUrole{w}{ }\sphinxbfcode{\sphinxupquote{\DUrole{n}{lammps\_config\_has\_ffmpeg\_support}}}}{}{}
\pysigstopmultiline
\pysigstopsignatures
\sphinxAtStartPar
Check if the LAMMPS library supports creating movie files via a pipe to ffmpeg

\sphinxAtStartPar
The LAMMPS {\hyperref[\detokenize{dump_image::doc}]{\sphinxcrossref{\DUrole{doc}{dump style movie}}}} supports generating movies
from images on\sphinxhyphen{}the\sphinxhyphen{}fly via creating a pipe to the
\sphinxhref{https://ffmpeg.org/}{ffmpeg} program.
This function checks whether this feature was {\hyperref[\detokenize{Build_settings:graphics}]{\sphinxcrossref{\DUrole{std,std-ref}{enabled at compile time}}}}.
It does \sphinxstylestrong{not} check whether the \sphinxcode{\sphinxupquote{ffmpeg}} itself is installed and usable.
\begin{quote}\begin{description}
\sphinxlineitem{Returns}
\sphinxAtStartPar
1 if yes, otherwise 0 

\end{description}\end{quote}

\end{fulllineitems}



\bigskip\hrule\bigskip

\index{lammps\_config\_has\_exceptions (C++ function)@\spxentry{lammps\_config\_has\_exceptions}\spxextra{C++ function}}

\begin{fulllineitems}
\phantomsection\label{\detokenize{Library_config:_CPPv428lammps_config_has_exceptionsv}}\phantomsection\label{\detokenize{Library_config:_CPPv328lammps_config_has_exceptionsv}}\phantomsection\label{\detokenize{Library_config:_CPPv228lammps_config_has_exceptionsv}}\phantomsection\label{\detokenize{Library_config:lammps_config_has_exceptions}}
\pysigstartsignatures
\pysigstartmultiline
\pysiglinewithargsret{\phantomsection\label{\detokenize{Library_config:library_8h_1ade5f1cc597079d8cb32a9765cbfff44c}}\DUrole{kt}{int}\DUrole{w}{ }\sphinxbfcode{\sphinxupquote{\DUrole{n}{lammps\_config\_has\_exceptions}}}}{}{}
\pysigstopmultiline
\pysigstopsignatures
\sphinxAtStartPar
Check whether LAMMPS errors will throw C++ exceptions.

\sphinxAtStartPar

\sphinxAtStartPar

\sphinxAtStartPar
\DUrole{versionmodified,deprecated}{Deprecated since version 21Nov2023: }LAMMPS has now exceptions always enabled, so this function
will now always return 1 and can be removed from applications
using the library interface.

\sphinxAtStartPar
In case of an error, LAMMPS will either abort or throw a C++ exception.
The latter has to be {\hyperref[\detokenize{Build_settings:exceptions}]{\sphinxcrossref{\DUrole{std,std-ref}{enabled at compile time}}}}.
This function checks if exceptions were enabled.

\sphinxAtStartPar
When using the library interface with C++ exceptions enabled,
the library interface functions will “catch” them and the
error status can then be checked by calling
{\hyperref[\detokenize{Library_utility:_CPPv416lammps_has_errorPv}]{\sphinxcrossref{\sphinxcode{\sphinxupquote{lammps\_has\_error()}}}}} and the most recent error message
can be retrieved via {\hyperref[\detokenize{Library_utility:_CPPv429lammps_get_last_error_messagePvPci}]{\sphinxcrossref{\sphinxcode{\sphinxupquote{lammps\_get\_last\_error\_message()}}}}}.
This can allow to restart a calculation or delete and recreate
the LAMMPS instance when C++ exceptions are enabled.  One application
of using exceptions this way is the {\hyperref[\detokenize{Tools:lammps-gui}]{\sphinxcrossref{\DUrole{std,std-ref}{LAMMPS\sphinxhyphen{}GUI}}}}.  If C++
exceptions are disabled and an error happens during a call to
LAMMPS, the application will terminate.

 
\begin{quote}\begin{description}
\sphinxlineitem{Returns}
\sphinxAtStartPar
1 if yes, otherwise 0 

\end{description}\end{quote}

\end{fulllineitems}



\bigskip\hrule\bigskip

\index{lammps\_config\_has\_package (C++ function)@\spxentry{lammps\_config\_has\_package}\spxextra{C++ function}}

\begin{fulllineitems}
\phantomsection\label{\detokenize{Library_config:_CPPv425lammps_config_has_packagePKc}}\phantomsection\label{\detokenize{Library_config:_CPPv325lammps_config_has_packagePKc}}\phantomsection\label{\detokenize{Library_config:_CPPv225lammps_config_has_packagePKc}}\phantomsection\label{\detokenize{Library_config:lammps_config_has_package__cCP}}
\pysigstartsignatures
\pysigstartmultiline
\pysiglinewithargsret{\phantomsection\label{\detokenize{Library_config:library_8h_1adca5959e5b824ab7413c74ee05eb9f0d}}\DUrole{kt}{int}\DUrole{w}{ }\sphinxbfcode{\sphinxupquote{\DUrole{n}{lammps\_config\_has\_package}}}}{\DUrole{k}{const}\DUrole{w}{ }\DUrole{kt}{char}\DUrole{p}{*}}{}
\pysigstopmultiline
\pysigstopsignatures
\sphinxAtStartPar
Check whether a specific package has been included in LAMMPS

\sphinxAtStartPar
This function checks whether the LAMMPS library in use includes the
specific {\hyperref[\detokenize{Packages::doc}]{\sphinxcrossref{\DUrole{doc}{LAMMPS package}}}} provided as argument.
\begin{quote}\begin{description}
\sphinxlineitem{Parameters}
\sphinxAtStartPar
\sphinxstyleliteralstrong{\sphinxupquote{name}} \textendash{} string with the name of the package 

\sphinxlineitem{Returns}
\sphinxAtStartPar
1 if included, 0 if not. 

\end{description}\end{quote}

\end{fulllineitems}



\bigskip\hrule\bigskip

\index{lammps\_config\_package\_count (C++ function)@\spxentry{lammps\_config\_package\_count}\spxextra{C++ function}}

\begin{fulllineitems}
\phantomsection\label{\detokenize{Library_config:_CPPv427lammps_config_package_countv}}\phantomsection\label{\detokenize{Library_config:_CPPv327lammps_config_package_countv}}\phantomsection\label{\detokenize{Library_config:_CPPv227lammps_config_package_countv}}\phantomsection\label{\detokenize{Library_config:lammps_config_package_count}}
\pysigstartsignatures
\pysigstartmultiline
\pysiglinewithargsret{\phantomsection\label{\detokenize{Library_config:library_8h_1a233aefab279fd9290f14fae1e4c16be4}}\DUrole{kt}{int}\DUrole{w}{ }\sphinxbfcode{\sphinxupquote{\DUrole{n}{lammps\_config\_package\_count}}}}{}{}
\pysigstopmultiline
\pysigstopsignatures
\sphinxAtStartPar
Count the number of installed packages in the LAMMPS library.

\sphinxAtStartPar
This function counts how many {\hyperref[\detokenize{Packages::doc}]{\sphinxcrossref{\DUrole{doc}{LAMMPS packages}}}} are
included in the LAMMPS library in use.
\begin{quote}\begin{description}
\sphinxlineitem{Returns}
\sphinxAtStartPar
number of packages included 

\end{description}\end{quote}

\end{fulllineitems}



\bigskip\hrule\bigskip

\index{lammps\_config\_package\_name (C++ function)@\spxentry{lammps\_config\_package\_name}\spxextra{C++ function}}

\begin{fulllineitems}
\phantomsection\label{\detokenize{Library_config:_CPPv426lammps_config_package_nameiPci}}\phantomsection\label{\detokenize{Library_config:_CPPv326lammps_config_package_nameiPci}}\phantomsection\label{\detokenize{Library_config:_CPPv226lammps_config_package_nameiPci}}\phantomsection\label{\detokenize{Library_config:lammps_config_package_name__i.cP.i}}
\pysigstartsignatures
\pysigstartmultiline
\pysiglinewithargsret{\phantomsection\label{\detokenize{Library_config:library_8h_1aca1da31b8eb4342ee1eff02d721fa48d}}\DUrole{kt}{int}\DUrole{w}{ }\sphinxbfcode{\sphinxupquote{\DUrole{n}{lammps\_config\_package\_name}}}}{\DUrole{kt}{int}\sphinxparamcomma \DUrole{kt}{char}\DUrole{p}{*}\sphinxparamcomma \DUrole{kt}{int}}{}
\pysigstopmultiline
\pysigstopsignatures
\sphinxAtStartPar
Get the name of a package in the list of installed packages in the LAMMPS library.

\sphinxAtStartPar
This function copies the name of the package with the index \sphinxstyleemphasis{idx} into the
provided C\sphinxhyphen{}style string buffer.  The length of the buffer must be provided
as \sphinxstyleemphasis{buf\_size} argument.  If the name of the package exceeds the length of the
buffer, it will be truncated accordingly.  If the index is out of range,
the function returns 0 and \sphinxstyleemphasis{buffer} is set to an empty string, otherwise 1;
\begin{quote}\begin{description}
\sphinxlineitem{Parameters}\begin{itemize}
\item {} 
\sphinxAtStartPar
\sphinxstyleliteralstrong{\sphinxupquote{idx}} \textendash{} index of the package in the list of included packages (0 \textless{}= idx \textless{} package count) 

\item {} 
\sphinxAtStartPar
\sphinxstyleliteralstrong{\sphinxupquote{buffer}} \textendash{} string buffer to copy the name of the package to 

\item {} 
\sphinxAtStartPar
\sphinxstyleliteralstrong{\sphinxupquote{buf\_size}} \textendash{} size of the provided string buffer 

\end{itemize}

\sphinxlineitem{Returns}
\sphinxAtStartPar
1 if successful, otherwise 0 

\end{description}\end{quote}

\end{fulllineitems}



\bigskip\hrule\bigskip

\index{lammps\_config\_accelerator (C++ function)@\spxentry{lammps\_config\_accelerator}\spxextra{C++ function}}

\begin{fulllineitems}
\phantomsection\label{\detokenize{Library_config:_CPPv425lammps_config_acceleratorPKcPKcPKc}}\phantomsection\label{\detokenize{Library_config:_CPPv325lammps_config_acceleratorPKcPKcPKc}}\phantomsection\label{\detokenize{Library_config:_CPPv225lammps_config_acceleratorPKcPKcPKc}}\phantomsection\label{\detokenize{Library_config:lammps_config_accelerator__cCP.cCP.cCP}}
\pysigstartsignatures
\pysigstartmultiline
\pysiglinewithargsret{\phantomsection\label{\detokenize{Library_config:library_8h_1a503c6feaadbab30ac00b978fcf199399}}\DUrole{kt}{int}\DUrole{w}{ }\sphinxbfcode{\sphinxupquote{\DUrole{n}{lammps\_config\_accelerator}}}}{\DUrole{k}{const}\DUrole{w}{ }\DUrole{kt}{char}\DUrole{p}{*}\sphinxparamcomma \DUrole{k}{const}\DUrole{w}{ }\DUrole{kt}{char}\DUrole{p}{*}\sphinxparamcomma \DUrole{k}{const}\DUrole{w}{ }\DUrole{kt}{char}\DUrole{p}{*}}{}
\pysigstopmultiline
\pysigstopsignatures
\sphinxAtStartPar
Check for compile time settings in accelerator packages included in LAMMPS.

\sphinxAtStartPar
This function checks availability of compile time settings of included
{\hyperref[\detokenize{Speed_packages::doc}]{\sphinxcrossref{\DUrole{doc}{accelerator packages}}}} in LAMMPS.
Supported packages names are “GPU”, “KOKKOS”, “INTEL”, and “OPENMP”.
Supported categories are “api” with possible settings “cuda”, “hip”, “phi”,
“pthreads”, “opencl”, “openmp”, and “serial”, and “precision” with
possible settings “double”, “mixed”, and “single”.  If the combination
of package, category, and setting is available, the function returns 1,
otherwise 0.
\begin{quote}\begin{description}
\sphinxlineitem{Parameters}\begin{itemize}
\item {} 
\sphinxAtStartPar
\sphinxstyleliteralstrong{\sphinxupquote{package}} \textendash{} string with the name of the accelerator package 

\item {} 
\sphinxAtStartPar
\sphinxstyleliteralstrong{\sphinxupquote{category}} \textendash{} string with the category name of the setting 

\item {} 
\sphinxAtStartPar
\sphinxstyleliteralstrong{\sphinxupquote{setting}} \textendash{} string with the name of the specific setting 

\end{itemize}

\sphinxlineitem{Returns}
\sphinxAtStartPar
1 if available, 0 if not. 

\end{description}\end{quote}

\end{fulllineitems}



\bigskip\hrule\bigskip

\index{lammps\_has\_gpu\_device (C++ function)@\spxentry{lammps\_has\_gpu\_device}\spxextra{C++ function}}

\begin{fulllineitems}
\phantomsection\label{\detokenize{Library_config:_CPPv421lammps_has_gpu_devicev}}\phantomsection\label{\detokenize{Library_config:_CPPv321lammps_has_gpu_devicev}}\phantomsection\label{\detokenize{Library_config:_CPPv221lammps_has_gpu_devicev}}\phantomsection\label{\detokenize{Library_config:lammps_has_gpu_device}}
\pysigstartsignatures
\pysigstartmultiline
\pysiglinewithargsret{\phantomsection\label{\detokenize{Library_config:library_8h_1acd3fd9fa79e08b7fb530b51b9ca79c8e}}\DUrole{kt}{int}\DUrole{w}{ }\sphinxbfcode{\sphinxupquote{\DUrole{n}{lammps\_has\_gpu\_device}}}}{}{}
\pysigstopmultiline
\pysigstopsignatures
\sphinxAtStartPar
Check for presence of a viable GPU package device

\sphinxAtStartPar

\sphinxAtStartPar
\DUrole{versionmodified,added}{Added in version 14May2021.}

\sphinxAtStartPar
The {\hyperref[\detokenize{Library_config:_CPPv421lammps_has_gpu_devicev}]{\sphinxcrossref{\sphinxcode{\sphinxupquote{lammps\_has\_gpu\_device()}}}}} function checks at runtime if
an accelerator device is present that can be used with the
{\hyperref[\detokenize{Speed_gpu::doc}]{\sphinxcrossref{\DUrole{doc}{GPU package}}}}. If at least one suitable device is
present the function will return 1, otherwise 0.

\sphinxAtStartPar
More detailed information about the available device or devices can
be obtained by calling the
{\hyperref[\detokenize{Library_config:_CPPv426lammps_get_gpu_device_infoPci}]{\sphinxcrossref{\sphinxcode{\sphinxupquote{lammps\_get\_gpu\_device\_info()}}}}} function.

\begin{quote}\begin{description}
\sphinxlineitem{Returns}
\sphinxAtStartPar
1 if viable device is available, 0 if not. 

\end{description}\end{quote}

\end{fulllineitems}



\bigskip\hrule\bigskip

\index{lammps\_get\_gpu\_device\_info (C++ function)@\spxentry{lammps\_get\_gpu\_device\_info}\spxextra{C++ function}}

\begin{fulllineitems}
\phantomsection\label{\detokenize{Library_config:_CPPv426lammps_get_gpu_device_infoPci}}\phantomsection\label{\detokenize{Library_config:_CPPv326lammps_get_gpu_device_infoPci}}\phantomsection\label{\detokenize{Library_config:_CPPv226lammps_get_gpu_device_infoPci}}\phantomsection\label{\detokenize{Library_config:lammps_get_gpu_device_info__cP.i}}
\pysigstartsignatures
\pysigstartmultiline
\pysiglinewithargsret{\phantomsection\label{\detokenize{Library_config:library_8h_1a8439a498ff1b8117a45e705e2a636c18}}\DUrole{kt}{void}\DUrole{w}{ }\sphinxbfcode{\sphinxupquote{\DUrole{n}{lammps\_get\_gpu\_device\_info}}}}{\DUrole{kt}{char}\DUrole{w}{ }\DUrole{p}{*}\DUrole{n,sig-param}{buffer}\sphinxparamcomma \DUrole{kt}{int}\DUrole{w}{ }\DUrole{n,sig-param}{buf\_size}}{}
\pysigstopmultiline
\pysigstopsignatures
\sphinxAtStartPar
Get GPU package device information

\sphinxAtStartPar

\sphinxAtStartPar
\DUrole{versionmodified,added}{Added in version 14May2021.}

\sphinxAtStartPar
The {\hyperref[\detokenize{Library_config:_CPPv426lammps_get_gpu_device_infoPci}]{\sphinxcrossref{\sphinxcode{\sphinxupquote{lammps\_get\_gpu\_device\_info()}}}}} function can be used to retrieve
detailed information about any accelerator devices that are viable for use
with the {\hyperref[\detokenize{Speed_gpu::doc}]{\sphinxcrossref{\DUrole{doc}{GPU package}}}}.  It will produce a string that is
equivalent to the output of the \sphinxcode{\sphinxupquote{nvc\_get\_device}} or \sphinxcode{\sphinxupquote{ocl\_get\_device}} or
\sphinxcode{\sphinxupquote{hip\_get\_device}} tools that are compiled alongside LAMMPS if the GPU
package is enabled.

\sphinxAtStartPar
A suitable buffer for a C\sphinxhyphen{}style string has to be provided and its length.
The assembled text will be truncated to not overflow this buffer.  This
string can be several kilobytes long, if multiple devices are present.

\begin{quote}\begin{description}
\sphinxlineitem{Parameters}\begin{itemize}
\item {} 
\sphinxAtStartPar
\sphinxstyleliteralstrong{\sphinxupquote{buffer}} \textendash{} string buffer to copy the information to 

\item {} 
\sphinxAtStartPar
\sphinxstyleliteralstrong{\sphinxupquote{buf\_size}} \textendash{} size of the provided string buffer 

\end{itemize}

\end{description}\end{quote}

\end{fulllineitems}



\bigskip\hrule\bigskip

\index{lammps\_has\_style (C++ function)@\spxentry{lammps\_has\_style}\spxextra{C++ function}}

\begin{fulllineitems}
\phantomsection\label{\detokenize{Library_config:_CPPv416lammps_has_stylePvPKcPKc}}\phantomsection\label{\detokenize{Library_config:_CPPv316lammps_has_stylePvPKcPKc}}\phantomsection\label{\detokenize{Library_config:_CPPv216lammps_has_stylePvPKcPKc}}\phantomsection\label{\detokenize{Library_config:lammps_has_style__voidP.cCP.cCP}}
\pysigstartsignatures
\pysigstartmultiline
\pysiglinewithargsret{\phantomsection\label{\detokenize{Library_config:library_8h_1a74aa598ce1e183ed70ce5eabf7202bec}}\DUrole{kt}{int}\DUrole{w}{ }\sphinxbfcode{\sphinxupquote{\DUrole{n}{lammps\_has\_style}}}}{\DUrole{kt}{void}\DUrole{p}{*}\sphinxparamcomma \DUrole{k}{const}\DUrole{w}{ }\DUrole{kt}{char}\DUrole{p}{*}\sphinxparamcomma \DUrole{k}{const}\DUrole{w}{ }\DUrole{kt}{char}\DUrole{p}{*}}{}
\pysigstopmultiline
\pysigstopsignatures
\sphinxAtStartPar
Check if a specific style has been included in LAMMPS

\sphinxAtStartPar
This function checks if the LAMMPS library in use includes the
specific \sphinxstyleemphasis{style} of a specific \sphinxstyleemphasis{category} provided as an argument.
Valid categories are: \sphinxstyleemphasis{atom}, \sphinxstyleemphasis{integrate}, \sphinxstyleemphasis{minimize},
\sphinxstyleemphasis{pair}, \sphinxstyleemphasis{bond}, \sphinxstyleemphasis{angle}, \sphinxstyleemphasis{dihedral}, \sphinxstyleemphasis{improper}, \sphinxstyleemphasis{kspace},
\sphinxstyleemphasis{compute}, \sphinxstyleemphasis{fix}, \sphinxstyleemphasis{region}, \sphinxstyleemphasis{dump}, and \sphinxstyleemphasis{command}.
\begin{quote}\begin{description}
\sphinxlineitem{Parameters}\begin{itemize}
\item {} 
\sphinxAtStartPar
\sphinxstyleliteralstrong{\sphinxupquote{handle}} \textendash{} pointer to a previously created LAMMPS instance cast to \sphinxcode{\sphinxupquote{void *}}. 

\item {} 
\sphinxAtStartPar
\sphinxstyleliteralstrong{\sphinxupquote{category}} \textendash{} category of the style 

\item {} 
\sphinxAtStartPar
\sphinxstyleliteralstrong{\sphinxupquote{name}} \textendash{} name of the style 

\end{itemize}

\sphinxlineitem{Returns}
\sphinxAtStartPar
1 if included, 0 if not. 

\end{description}\end{quote}

\end{fulllineitems}



\bigskip\hrule\bigskip

\index{lammps\_style\_count (C++ function)@\spxentry{lammps\_style\_count}\spxextra{C++ function}}

\begin{fulllineitems}
\phantomsection\label{\detokenize{Library_config:_CPPv418lammps_style_countPvPKc}}\phantomsection\label{\detokenize{Library_config:_CPPv318lammps_style_countPvPKc}}\phantomsection\label{\detokenize{Library_config:_CPPv218lammps_style_countPvPKc}}\phantomsection\label{\detokenize{Library_config:lammps_style_count__voidP.cCP}}
\pysigstartsignatures
\pysigstartmultiline
\pysiglinewithargsret{\phantomsection\label{\detokenize{Library_config:library_8h_1aef8b5d9f8c06dc22c55deb9088dbafd6}}\DUrole{kt}{int}\DUrole{w}{ }\sphinxbfcode{\sphinxupquote{\DUrole{n}{lammps\_style\_count}}}}{\DUrole{kt}{void}\DUrole{p}{*}\sphinxparamcomma \DUrole{k}{const}\DUrole{w}{ }\DUrole{kt}{char}\DUrole{p}{*}}{}
\pysigstopmultiline
\pysigstopsignatures
\sphinxAtStartPar
Count the number of styles of category in the LAMMPS library.

\sphinxAtStartPar
This function counts how many styles in the provided \sphinxstyleemphasis{category}
are included in the LAMMPS library in use.
Please see {\hyperref[\detokenize{Library_config:_CPPv416lammps_has_stylePvPKcPKc}]{\sphinxcrossref{\sphinxcode{\sphinxupquote{lammps\_has\_style()}}}}} for a list of valid
categories.
\begin{quote}\begin{description}
\sphinxlineitem{Parameters}\begin{itemize}
\item {} 
\sphinxAtStartPar
\sphinxstyleliteralstrong{\sphinxupquote{handle}} \textendash{} pointer to a previously created LAMMPS instance cast to \sphinxcode{\sphinxupquote{void *}}. 

\item {} 
\sphinxAtStartPar
\sphinxstyleliteralstrong{\sphinxupquote{category}} \textendash{} category of styles 

\end{itemize}

\sphinxlineitem{Returns}
\sphinxAtStartPar
number of styles in category 

\end{description}\end{quote}

\end{fulllineitems}



\bigskip\hrule\bigskip

\index{lammps\_style\_name (C++ function)@\spxentry{lammps\_style\_name}\spxextra{C++ function}}

\begin{fulllineitems}
\phantomsection\label{\detokenize{Library_config:_CPPv417lammps_style_namePvPKciPci}}\phantomsection\label{\detokenize{Library_config:_CPPv317lammps_style_namePvPKciPci}}\phantomsection\label{\detokenize{Library_config:_CPPv217lammps_style_namePvPKciPci}}\phantomsection\label{\detokenize{Library_config:lammps_style_name__voidP.cCP.i.cP.i}}
\pysigstartsignatures
\pysigstartmultiline
\pysiglinewithargsret{\phantomsection\label{\detokenize{Library_config:library_8h_1ab91b5c7e97982f7a804d8d7cfa0f190a}}\DUrole{kt}{int}\DUrole{w}{ }\sphinxbfcode{\sphinxupquote{\DUrole{n}{lammps\_style\_name}}}}{\DUrole{kt}{void}\DUrole{p}{*}\sphinxparamcomma \DUrole{k}{const}\DUrole{w}{ }\DUrole{kt}{char}\DUrole{p}{*}\sphinxparamcomma \DUrole{kt}{int}\sphinxparamcomma \DUrole{kt}{char}\DUrole{p}{*}\sphinxparamcomma \DUrole{kt}{int}}{}
\pysigstopmultiline
\pysigstopsignatures
\sphinxAtStartPar
Look up the name of a style by index in the list of style of a given category in the LAMMPS library.

\sphinxAtStartPar
This function copies the name of the \sphinxstyleemphasis{category} style with the index \sphinxstyleemphasis{idx} into the provided C\sphinxhyphen{}style string buffer. The length of the buffer must be provided as \sphinxstyleemphasis{buf\_size} argument. If the name of the style exceeds the length of the buffer, it will be truncated accordingly. If the index is out of range, the function returns 0 and \sphinxstyleemphasis{buffer} is set to an empty string, otherwise 1.
\begin{quote}\begin{description}
\sphinxlineitem{Parameters}\begin{itemize}
\item {} 
\sphinxAtStartPar
\sphinxstyleliteralstrong{\sphinxupquote{handle}} \textendash{} pointer to a previously created LAMMPS instance cast to \sphinxcode{\sphinxupquote{void *}}. 

\item {} 
\sphinxAtStartPar
\sphinxstyleliteralstrong{\sphinxupquote{category}} \textendash{} category of styles 

\item {} 
\sphinxAtStartPar
\sphinxstyleliteralstrong{\sphinxupquote{idx}} \textendash{} index of the style in the list of \sphinxstyleemphasis{category} styles (0 \textless{}= idx \textless{} style count) 

\item {} 
\sphinxAtStartPar
\sphinxstyleliteralstrong{\sphinxupquote{buffer}} \textendash{} string buffer to copy the name of the style to 

\item {} 
\sphinxAtStartPar
\sphinxstyleliteralstrong{\sphinxupquote{buf\_size}} \textendash{} size of the provided string buffer 

\end{itemize}

\sphinxlineitem{Returns}
\sphinxAtStartPar
1 if successful, otherwise 0 

\end{description}\end{quote}

\end{fulllineitems}



\bigskip\hrule\bigskip

\index{lammps\_has\_id (C++ function)@\spxentry{lammps\_has\_id}\spxextra{C++ function}}

\begin{fulllineitems}
\phantomsection\label{\detokenize{Library_config:_CPPv413lammps_has_idPvPKcPKc}}\phantomsection\label{\detokenize{Library_config:_CPPv313lammps_has_idPvPKcPKc}}\phantomsection\label{\detokenize{Library_config:_CPPv213lammps_has_idPvPKcPKc}}\phantomsection\label{\detokenize{Library_config:lammps_has_id__voidP.cCP.cCP}}
\pysigstartsignatures
\pysigstartmultiline
\pysiglinewithargsret{\phantomsection\label{\detokenize{Library_config:library_8h_1ad1ee62546d5d560747379cba84bea178}}\DUrole{kt}{int}\DUrole{w}{ }\sphinxbfcode{\sphinxupquote{\DUrole{n}{lammps\_has\_id}}}}{\DUrole{kt}{void}\DUrole{p}{*}\sphinxparamcomma \DUrole{k}{const}\DUrole{w}{ }\DUrole{kt}{char}\DUrole{p}{*}\sphinxparamcomma \DUrole{k}{const}\DUrole{w}{ }\DUrole{kt}{char}\DUrole{p}{*}}{}
\pysigstopmultiline
\pysigstopsignatures
\sphinxAtStartPar
Check if a specific ID exists in the current LAMMPS instance

\sphinxAtStartPar

\sphinxAtStartPar
\DUrole{versionmodified,added}{Added in version 9Oct2020.}

\sphinxAtStartPar
This function checks if the current LAMMPS instance a \sphinxstyleemphasis{category} ID of
the given \sphinxstyleemphasis{name} exists.  Valid categories are: \sphinxstyleemphasis{compute}, \sphinxstyleemphasis{dump},
\sphinxstyleemphasis{fix}, \sphinxstyleemphasis{group}, \sphinxstyleemphasis{molecule}, \sphinxstyleemphasis{region}, and \sphinxstyleemphasis{variable}.

\begin{quote}\begin{description}
\sphinxlineitem{Parameters}\begin{itemize}
\item {} 
\sphinxAtStartPar
\sphinxstyleliteralstrong{\sphinxupquote{handle}} \textendash{} pointer to a previously created LAMMPS instance cast to \sphinxcode{\sphinxupquote{void *}}. 

\item {} 
\sphinxAtStartPar
\sphinxstyleliteralstrong{\sphinxupquote{category}} \textendash{} category of the id 

\item {} 
\sphinxAtStartPar
\sphinxstyleliteralstrong{\sphinxupquote{name}} \textendash{} name of the id 

\end{itemize}

\sphinxlineitem{Returns}
\sphinxAtStartPar
1 if included, 0 if not. 

\end{description}\end{quote}

\end{fulllineitems}



\bigskip\hrule\bigskip

\index{lammps\_id\_count (C++ function)@\spxentry{lammps\_id\_count}\spxextra{C++ function}}

\begin{fulllineitems}
\phantomsection\label{\detokenize{Library_config:_CPPv415lammps_id_countPvPKc}}\phantomsection\label{\detokenize{Library_config:_CPPv315lammps_id_countPvPKc}}\phantomsection\label{\detokenize{Library_config:_CPPv215lammps_id_countPvPKc}}\phantomsection\label{\detokenize{Library_config:lammps_id_count__voidP.cCP}}
\pysigstartsignatures
\pysigstartmultiline
\pysiglinewithargsret{\phantomsection\label{\detokenize{Library_config:library_8h_1aaa1311b9ce55c53bfae199c2ee9cb955}}\DUrole{kt}{int}\DUrole{w}{ }\sphinxbfcode{\sphinxupquote{\DUrole{n}{lammps\_id\_count}}}}{\DUrole{kt}{void}\DUrole{p}{*}\sphinxparamcomma \DUrole{k}{const}\DUrole{w}{ }\DUrole{kt}{char}\DUrole{p}{*}}{}
\pysigstopmultiline
\pysigstopsignatures
\sphinxAtStartPar
Count the number of IDs of a category.

\sphinxAtStartPar

\sphinxAtStartPar
\DUrole{versionmodified,added}{Added in version 9Oct2020.}

\sphinxAtStartPar
This function counts how many IDs in the provided \sphinxstyleemphasis{category}
are defined in the current LAMMPS instance.
Please see {\hyperref[\detokenize{Library_config:_CPPv413lammps_has_idPvPKcPKc}]{\sphinxcrossref{\sphinxcode{\sphinxupquote{lammps\_has\_id()}}}}} for a list of valid
categories.

\begin{quote}\begin{description}
\sphinxlineitem{Parameters}\begin{itemize}
\item {} 
\sphinxAtStartPar
\sphinxstyleliteralstrong{\sphinxupquote{handle}} \textendash{} pointer to a previously created LAMMPS instance cast to \sphinxcode{\sphinxupquote{void *}}. 

\item {} 
\sphinxAtStartPar
\sphinxstyleliteralstrong{\sphinxupquote{category}} \textendash{} category of IDs 

\end{itemize}

\sphinxlineitem{Returns}
\sphinxAtStartPar
number of IDs in category 

\end{description}\end{quote}

\end{fulllineitems}



\bigskip\hrule\bigskip

\index{lammps\_id\_name (C++ function)@\spxentry{lammps\_id\_name}\spxextra{C++ function}}

\begin{fulllineitems}
\phantomsection\label{\detokenize{Library_config:_CPPv414lammps_id_namePvPKciPci}}\phantomsection\label{\detokenize{Library_config:_CPPv314lammps_id_namePvPKciPci}}\phantomsection\label{\detokenize{Library_config:_CPPv214lammps_id_namePvPKciPci}}\phantomsection\label{\detokenize{Library_config:lammps_id_name__voidP.cCP.i.cP.i}}
\pysigstartsignatures
\pysigstartmultiline
\pysiglinewithargsret{\phantomsection\label{\detokenize{Library_config:library_8h_1a42ba710155bca355ba451fd45dcda0f9}}\DUrole{kt}{int}\DUrole{w}{ }\sphinxbfcode{\sphinxupquote{\DUrole{n}{lammps\_id\_name}}}}{\DUrole{kt}{void}\DUrole{p}{*}\sphinxparamcomma \DUrole{k}{const}\DUrole{w}{ }\DUrole{kt}{char}\DUrole{p}{*}\sphinxparamcomma \DUrole{kt}{int}\sphinxparamcomma \DUrole{kt}{char}\DUrole{p}{*}\sphinxparamcomma \DUrole{kt}{int}}{}
\pysigstopmultiline
\pysigstopsignatures
\sphinxAtStartPar
Look up the name of an ID by index in the list of IDs of a given category.

\sphinxAtStartPar

\sphinxAtStartPar
\DUrole{versionmodified,added}{Added in version 9Oct2020.}

\sphinxAtStartPar
This function copies the name of the \sphinxstyleemphasis{category} ID with the index
\sphinxstyleemphasis{idx} into the provided C\sphinxhyphen{}style string buffer.  The length of the buffer
must be provided as \sphinxstyleemphasis{buf\_size} argument.  If the name of the style
exceeds the length of the buffer, it will be truncated accordingly.
If the index is out of range, the function returns 0 and \sphinxstyleemphasis{buffer} is
set to an empty string, otherwise 1.

\begin{quote}\begin{description}
\sphinxlineitem{Parameters}\begin{itemize}
\item {} 
\sphinxAtStartPar
\sphinxstyleliteralstrong{\sphinxupquote{handle}} \textendash{} pointer to a previously created LAMMPS instance cast to \sphinxcode{\sphinxupquote{void *}}. 

\item {} 
\sphinxAtStartPar
\sphinxstyleliteralstrong{\sphinxupquote{category}} \textendash{} category of IDs 

\item {} 
\sphinxAtStartPar
\sphinxstyleliteralstrong{\sphinxupquote{idx}} \textendash{} index of the ID in the list of \sphinxstyleemphasis{category} styles (0 \textless{}= idx \textless{} count) 

\item {} 
\sphinxAtStartPar
\sphinxstyleliteralstrong{\sphinxupquote{buffer}} \textendash{} string buffer to copy the name of the style to 

\item {} 
\sphinxAtStartPar
\sphinxstyleliteralstrong{\sphinxupquote{buf\_size}} \textendash{} size of the provided string buffer 

\end{itemize}

\sphinxlineitem{Returns}
\sphinxAtStartPar
1 if successful, otherwise 0 

\end{description}\end{quote}

\end{fulllineitems}


\sphinxstepscope


\subsection{Utility functions}
\label{\detokenize{Library_utility:utility-functions}}\label{\detokenize{Library_utility::doc}}
\sphinxAtStartPar
To simplify some tasks, the library interface contains these utility
functions.  They do not directly call the LAMMPS library.
\begin{itemize}
\item {} 
\sphinxAtStartPar
{\hyperref[\detokenize{Library_utility:_CPPv425lammps_encode_image_flagsiii}]{\sphinxcrossref{\sphinxcode{\sphinxupquote{lammps\_encode\_image\_flags()}}}}}

\item {} 
\sphinxAtStartPar
{\hyperref[\detokenize{Library_utility:_CPPv425lammps_decode_image_flagsiPi}]{\sphinxcrossref{\sphinxcode{\sphinxupquote{lammps\_decode\_image\_flags()}}}}}

\item {} 
\sphinxAtStartPar
{\hyperref[\detokenize{Library_utility:_CPPv432lammps_set_fix_external_callbackPvPKc16FixExternalFnPtrPv}]{\sphinxcrossref{\sphinxcode{\sphinxupquote{lammps\_set\_fix\_external\_callback()}}}}}

\item {} 
\sphinxAtStartPar
{\hyperref[\detokenize{Library_utility:_CPPv429lammps_fix_external_get_forcePvPKc}]{\sphinxcrossref{\sphinxcode{\sphinxupquote{lammps\_fix\_external\_get\_force()}}}}}

\item {} 
\sphinxAtStartPar
{\hyperref[\detokenize{Library_utility:_CPPv437lammps_fix_external_set_energy_globalPvPKcd}]{\sphinxcrossref{\sphinxcode{\sphinxupquote{lammps\_fix\_external\_set\_energy\_global()}}}}}

\item {} 
\sphinxAtStartPar
{\hyperref[\detokenize{Library_utility:_CPPv438lammps_fix_external_set_energy_peratomPvPKcPd}]{\sphinxcrossref{\sphinxcode{\sphinxupquote{lammps\_fix\_external\_set\_energy\_peratom()}}}}}

\item {} 
\sphinxAtStartPar
{\hyperref[\detokenize{Library_utility:_CPPv437lammps_fix_external_set_virial_globalPvPKcPd}]{\sphinxcrossref{\sphinxcode{\sphinxupquote{lammps\_fix\_external\_set\_virial\_global()}}}}}

\item {} 
\sphinxAtStartPar
{\hyperref[\detokenize{Library_utility:_CPPv438lammps_fix_external_set_virial_peratomPvPKcPPd}]{\sphinxcrossref{\sphinxcode{\sphinxupquote{lammps\_fix\_external\_set\_virial\_peratom()}}}}}

\item {} 
\sphinxAtStartPar
{\hyperref[\detokenize{Library_utility:_CPPv437lammps_fix_external_set_vector_lengthPvPKci}]{\sphinxcrossref{\sphinxcode{\sphinxupquote{lammps\_fix\_external\_set\_vector\_length()}}}}}

\item {} 
\sphinxAtStartPar
{\hyperref[\detokenize{Library_utility:_CPPv430lammps_fix_external_set_vectorPvPKcid}]{\sphinxcrossref{\sphinxcode{\sphinxupquote{lammps\_fix\_external\_set\_vector()}}}}}

\item {} 
\sphinxAtStartPar
{\hyperref[\detokenize{Library_utility:_CPPv420lammps_flush_buffersPv}]{\sphinxcrossref{\sphinxcode{\sphinxupquote{lammps\_flush\_buffers()}}}}}

\item {} 
\sphinxAtStartPar
{\hyperref[\detokenize{Library_utility:_CPPv411lammps_freePv}]{\sphinxcrossref{\sphinxcode{\sphinxupquote{lammps\_free()}}}}}

\item {} 
\sphinxAtStartPar
{\hyperref[\detokenize{Library_utility:_CPPv417lammps_is_runningPv}]{\sphinxcrossref{\sphinxcode{\sphinxupquote{lammps\_is\_running()}}}}}

\item {} 
\sphinxAtStartPar
{\hyperref[\detokenize{Library_utility:_CPPv420lammps_force_timeoutPv}]{\sphinxcrossref{\sphinxcode{\sphinxupquote{lammps\_force\_timeout()}}}}}

\item {} 
\sphinxAtStartPar
{\hyperref[\detokenize{Library_utility:_CPPv416lammps_has_errorPv}]{\sphinxcrossref{\sphinxcode{\sphinxupquote{lammps\_has\_error()}}}}}

\item {} 
\sphinxAtStartPar
{\hyperref[\detokenize{Library_utility:_CPPv429lammps_get_last_error_messagePvPci}]{\sphinxcrossref{\sphinxcode{\sphinxupquote{lammps\_get\_last\_error\_message()}}}}}

\item {} 
\sphinxAtStartPar
{\hyperref[\detokenize{Library_utility:_CPPv425lammps_python_api_versionv}]{\sphinxcrossref{\sphinxcode{\sphinxupquote{lammps\_python\_api\_version()}}}}}

\end{itemize}

\sphinxAtStartPar
The {\hyperref[\detokenize{Library_utility:_CPPv411lammps_freePv}]{\sphinxcrossref{\sphinxcode{\sphinxupquote{lammps\_free()}}}}} function is a clean\sphinxhyphen{}up function to free
memory that the library had allocated previously via other function
calls.  Look for notes in the descriptions of the individual commands
where such memory buffers were allocated that require the use of
{\hyperref[\detokenize{Library_utility:_CPPv411lammps_freePv}]{\sphinxcrossref{\sphinxcode{\sphinxupquote{lammps\_free()}}}}}.


\bigskip\hrule\bigskip

\index{lammps\_encode\_image\_flags (C++ function)@\spxentry{lammps\_encode\_image\_flags}\spxextra{C++ function}}

\begin{fulllineitems}
\phantomsection\label{\detokenize{Library_utility:_CPPv425lammps_encode_image_flagsiii}}\phantomsection\label{\detokenize{Library_utility:_CPPv325lammps_encode_image_flagsiii}}\phantomsection\label{\detokenize{Library_utility:_CPPv225lammps_encode_image_flagsiii}}\phantomsection\label{\detokenize{Library_utility:lammps_encode_image_flags__i.i.i}}
\pysigstartsignatures
\pysigstartmultiline
\pysiglinewithargsret{\phantomsection\label{\detokenize{Library_utility:library_8h_1ad12aad3df456f3c85418ff8af6e55e10}}\DUrole{kt}{int}\DUrole{w}{ }\sphinxbfcode{\sphinxupquote{\DUrole{n}{lammps\_encode\_image\_flags}}}}{\DUrole{kt}{int}\DUrole{w}{ }\DUrole{n,sig-param}{ix}\sphinxparamcomma \DUrole{kt}{int}\DUrole{w}{ }\DUrole{n,sig-param}{iy}\sphinxparamcomma \DUrole{kt}{int}\DUrole{w}{ }\DUrole{n,sig-param}{iz}}{}
\pysigstopmultiline
\pysigstopsignatures
\sphinxAtStartPar
Encode three integer image flags into a single imageint.

\sphinxAtStartPar

\sphinxAtStartPar
The prototype for this function when compiling with \sphinxcode{\sphinxupquote{\sphinxhyphen{}DLAMMPS\_BIGBIG}}
is:

\begin{sphinxVerbatim}[commandchars=\\\{\}]
\PYG{k+kt}{int64\PYGZus{}t}\PYG{+w}{ }\PYG{n+nf}{lammps\PYGZus{}encode\PYGZus{}image\PYGZus{}flags}\PYG{p}{(}\PYG{k+kt}{int}\PYG{+w}{ }\PYG{n}{ix}\PYG{p}{,}\PYG{+w}{ }\PYG{k+kt}{int}\PYG{+w}{ }\PYG{n}{iy}\PYG{p}{,}\PYG{+w}{ }\PYG{k+kt}{int}\PYG{+w}{ }\PYG{n}{iz}\PYG{p}{)}\PYG{p}{;}
\end{sphinxVerbatim}

\sphinxAtStartPar
This function performs the bit\sphinxhyphen{}shift, addition, and bit\sphinxhyphen{}wise OR
operations necessary to combine the values of three integers
representing the image flags in x\sphinxhyphen{}, y\sphinxhyphen{}, and z\sphinxhyphen{}direction.  Unless
LAMMPS is compiled with \sphinxhyphen{}DLAMMPS\_BIGBIG, those integers are
limited 10\sphinxhyphen{}bit signed integers {[}\sphinxhyphen{}512, 511{]}.  Otherwise the return
type changes from \sphinxcode{\sphinxupquote{int}} to \sphinxcode{\sphinxupquote{int64\_t}} and the valid range for
the individual image flags becomes {[}\sphinxhyphen{}1048576,1048575{]},
i.e. that of a 21\sphinxhyphen{}bit signed integer.  There is no check on whether
the arguments conform to these requirements.

\begin{quote}\begin{description}
\sphinxlineitem{Parameters}\begin{itemize}
\item {} 
\sphinxAtStartPar
\sphinxstyleliteralstrong{\sphinxupquote{ix}} \textendash{} image flag value in x 

\item {} 
\sphinxAtStartPar
\sphinxstyleliteralstrong{\sphinxupquote{iy}} \textendash{} image flag value in y 

\item {} 
\sphinxAtStartPar
\sphinxstyleliteralstrong{\sphinxupquote{iz}} \textendash{} image flag value in z 

\end{itemize}

\sphinxlineitem{Returns}
\sphinxAtStartPar
encoded image flag integer 

\end{description}\end{quote}

\end{fulllineitems}



\bigskip\hrule\bigskip

\index{lammps\_decode\_image\_flags (C++ function)@\spxentry{lammps\_decode\_image\_flags}\spxextra{C++ function}}

\begin{fulllineitems}
\phantomsection\label{\detokenize{Library_utility:_CPPv425lammps_decode_image_flagsiPi}}\phantomsection\label{\detokenize{Library_utility:_CPPv325lammps_decode_image_flagsiPi}}\phantomsection\label{\detokenize{Library_utility:_CPPv225lammps_decode_image_flagsiPi}}\phantomsection\label{\detokenize{Library_utility:lammps_decode_image_flags__i.iP}}
\pysigstartsignatures
\pysigstartmultiline
\pysiglinewithargsret{\phantomsection\label{\detokenize{Library_utility:library_8h_1ae774bc2a40ffa11e7a7aecfac7a9c3c6}}\DUrole{kt}{void}\DUrole{w}{ }\sphinxbfcode{\sphinxupquote{\DUrole{n}{lammps\_decode\_image\_flags}}}}{\DUrole{kt}{int}\DUrole{w}{ }\DUrole{n,sig-param}{image}\sphinxparamcomma \DUrole{kt}{int}\DUrole{w}{ }\DUrole{p}{*}\DUrole{n,sig-param}{flags}}{}
\pysigstopmultiline
\pysigstopsignatures
\sphinxAtStartPar
Decode a single image flag integer into three regular integers

\sphinxAtStartPar

\sphinxAtStartPar
The prototype for this function when compiling with \sphinxcode{\sphinxupquote{\sphinxhyphen{}DLAMMPS\_BIGBIG}}
is:

\begin{sphinxVerbatim}[commandchars=\\\{\}]
\PYG{k+kt}{void}\PYG{+w}{ }\PYG{n+nf}{lammps\PYGZus{}decode\PYGZus{}image\PYGZus{}flags}\PYG{p}{(}\PYG{k+kt}{int64\PYGZus{}t}\PYG{+w}{ }\PYG{n}{image}\PYG{p}{,}\PYG{+w}{ }\PYG{k+kt}{int}\PYG{+w}{ }\PYG{o}{*}\PYG{n}{flags}\PYG{p}{)}\PYG{p}{;}
\end{sphinxVerbatim}

\sphinxAtStartPar
This function does the reverse operation of
{\hyperref[\detokenize{Library_utility:_CPPv425lammps_encode_image_flagsiii}]{\sphinxcrossref{\sphinxcode{\sphinxupquote{lammps\_encode\_image\_flags()}}}}} and takes an image flag integer
does the bit\sphinxhyphen{}shift and bit\sphinxhyphen{}masking operations to decode it and stores
the resulting three regular integers into the buffer pointed to by
\sphinxstyleemphasis{flags}.

\begin{quote}\begin{description}
\sphinxlineitem{Parameters}\begin{itemize}
\item {} 
\sphinxAtStartPar
\sphinxstyleliteralstrong{\sphinxupquote{image}} \textendash{} encoded image flag integer 

\item {} 
\sphinxAtStartPar
\sphinxstyleliteralstrong{\sphinxupquote{flags}} \textendash{} pointer to storage where the decoded image flags are stored. 

\end{itemize}

\end{description}\end{quote}

\end{fulllineitems}



\bigskip\hrule\bigskip

\index{lammps\_set\_fix\_external\_callback (C++ function)@\spxentry{lammps\_set\_fix\_external\_callback}\spxextra{C++ function}}

\begin{fulllineitems}
\phantomsection\label{\detokenize{Library_utility:_CPPv432lammps_set_fix_external_callbackPvPKc16FixExternalFnPtrPv}}\phantomsection\label{\detokenize{Library_utility:_CPPv332lammps_set_fix_external_callbackPvPKc16FixExternalFnPtrPv}}\phantomsection\label{\detokenize{Library_utility:_CPPv232lammps_set_fix_external_callbackPvPKc16FixExternalFnPtrPv}}\phantomsection\label{\detokenize{Library_utility:lammps_set_fix_external_callback__voidP.cCP.FixExternalFnPtr.voidP}}
\pysigstartsignatures
\pysigstartmultiline
\pysiglinewithargsret{\phantomsection\label{\detokenize{Library_utility:library_8h_1a2b2994fdf05652cb2f3fe2afc4d00d1a}}\DUrole{kt}{void}\DUrole{w}{ }\sphinxbfcode{\sphinxupquote{\DUrole{n}{lammps\_set\_fix\_external\_callback}}}}{\DUrole{kt}{void}\DUrole{w}{ }\DUrole{p}{*}\DUrole{n,sig-param}{handle}\sphinxparamcomma \DUrole{k}{const}\DUrole{w}{ }\DUrole{kt}{char}\DUrole{w}{ }\DUrole{p}{*}\DUrole{n,sig-param}{id}\sphinxparamcomma \DUrole{n}{FixExternalFnPtr}\DUrole{w}{ }\DUrole{n,sig-param}{funcptr}\sphinxparamcomma \DUrole{kt}{void}\DUrole{w}{ }\DUrole{p}{*}\DUrole{n,sig-param}{ptr}}{}
\pysigstopmultiline
\pysigstopsignatures
\sphinxAtStartPar
Set up the callback function for a fix external instance with the given ID.

\sphinxAtStartPar

\sphinxAtStartPar
Fix {\hyperref[\detokenize{fix_external::doc}]{\sphinxcrossref{\DUrole{doc}{external}}}} allows programs that are running LAMMPS through
its library interface to modify certain LAMMPS properties on specific
timesteps, similar to the way other fixes do.

\sphinxAtStartPar
This function sets the callback function for use with the “pf/callback”
mode. The function has to have C language bindings with the prototype:

\begin{sphinxVerbatim}[commandchars=\\\{\}]
\PYG{k+kt}{void}\PYG{+w}{ }\PYG{n+nf}{func}\PYG{p}{(}\PYG{k+kt}{void}\PYG{+w}{ }\PYG{o}{*}\PYG{n}{ptr}\PYG{p}{,}\PYG{+w}{ }\PYG{n}{bigint}\PYG{+w}{ }\PYG{n}{timestep}\PYG{p}{,}\PYG{+w}{ }\PYG{k+kt}{int}\PYG{+w}{ }\PYG{n}{nlocal}\PYG{p}{,}\PYG{+w}{ }\PYG{n}{tagint}\PYG{+w}{ }\PYG{o}{*}\PYG{n}{ids}\PYG{p}{,}\PYG{+w}{ }\PYG{k+kt}{double}\PYG{+w}{ }\PYG{o}{*}\PYG{o}{*}\PYG{n}{x}\PYG{p}{,}\PYG{+w}{ }\PYG{k+kt}{double}\PYG{+w}{ }\PYG{o}{*}\PYG{o}{*}\PYG{n}{fexternal}\PYG{p}{)}\PYG{p}{;}
\end{sphinxVerbatim}

\sphinxAtStartPar
The argument \sphinxstyleemphasis{ptr} to this function will be stored in fix external and
the passed as the first argument calling the callback function \sphinxtitleref{func()}.
This would usually be a pointer to the active LAMMPS instance, i.e. the same
pointer as the \sphinxstyleemphasis{handle} argument.  This would be needed to call
functions that set the global or per\sphinxhyphen{}atom energy or virial contributions
from within the callback function.

\sphinxAtStartPar
The callback mechanism is one of the two modes of how forces and can be
applied to a simulation with the help of fix external. The alternative
is the array mode where you call {\hyperref[\detokenize{Library_utility:_CPPv429lammps_fix_external_get_forcePvPKc}]{\sphinxcrossref{\sphinxcode{\sphinxupquote{lammps\_fix\_external\_get\_force()}}}}}.

\sphinxAtStartPar
Please see the documentation for {\hyperref[\detokenize{fix_external::doc}]{\sphinxcrossref{\DUrole{doc}{fix external}}}} for
more information about how to use the fix and how to couple it with an
external code.

\sphinxAtStartPar
\DUrole{versionmodified,changed}{Changed in version 28Jul2021.}

\begin{quote}\begin{description}
\sphinxlineitem{Parameters}\begin{itemize}
\item {} 
\sphinxAtStartPar
\sphinxstyleliteralstrong{\sphinxupquote{handle}} \textendash{} pointer to a previously created LAMMPS instance cast to \sphinxcode{\sphinxupquote{void *}}. 

\item {} 
\sphinxAtStartPar
\sphinxstyleliteralstrong{\sphinxupquote{id}} \textendash{} fix ID of fix external instance 

\item {} 
\sphinxAtStartPar
\sphinxstyleliteralstrong{\sphinxupquote{funcptr}} \textendash{} pointer to callback function 

\item {} 
\sphinxAtStartPar
\sphinxstyleliteralstrong{\sphinxupquote{ptr}} \textendash{} pointer to object in calling code, passed to callback function as first argument 

\end{itemize}

\end{description}\end{quote}

\end{fulllineitems}



\bigskip\hrule\bigskip

\index{lammps\_fix\_external\_get\_force (C++ function)@\spxentry{lammps\_fix\_external\_get\_force}\spxextra{C++ function}}

\begin{fulllineitems}
\phantomsection\label{\detokenize{Library_utility:_CPPv429lammps_fix_external_get_forcePvPKc}}\phantomsection\label{\detokenize{Library_utility:_CPPv329lammps_fix_external_get_forcePvPKc}}\phantomsection\label{\detokenize{Library_utility:_CPPv229lammps_fix_external_get_forcePvPKc}}\phantomsection\label{\detokenize{Library_utility:lammps_fix_external_get_force__voidP.cCP}}
\pysigstartsignatures
\pysigstartmultiline
\pysiglinewithargsret{\phantomsection\label{\detokenize{Library_utility:library_8h_1a944127eb18f7e86dc50fcf5024ebb390}}\DUrole{kt}{double}\DUrole{w}{ }\DUrole{p}{*}\DUrole{p}{*}\sphinxbfcode{\sphinxupquote{\DUrole{n}{lammps\_fix\_external\_get\_force}}}}{\DUrole{kt}{void}\DUrole{w}{ }\DUrole{p}{*}\DUrole{n,sig-param}{handle}\sphinxparamcomma \DUrole{k}{const}\DUrole{w}{ }\DUrole{kt}{char}\DUrole{w}{ }\DUrole{p}{*}\DUrole{n,sig-param}{id}}{}
\pysigstopmultiline
\pysigstopsignatures
\sphinxAtStartPar
Get pointer to the force array storage in a fix external instance with the given ID.

\sphinxAtStartPar

\sphinxAtStartPar
\DUrole{versionmodified,added}{Added in version 28Jul2021.}

\sphinxAtStartPar
Fix {\hyperref[\detokenize{fix_external::doc}]{\sphinxcrossref{\DUrole{doc}{external}}}} allows programs that are running
LAMMPS through its library interface to add or modify certain LAMMPS
properties on specific timesteps, similar to the way other fixes do.

\sphinxAtStartPar
This function provides access to the per\sphinxhyphen{}atom force storage in a fix
external instance with the given fix\sphinxhyphen{}ID to be added to the individual
atoms when using the “pf/array” mode.  The \sphinxstyleemphasis{fexternal} array can be
accessed like other “native” per\sphinxhyphen{}atom arrays accessible via the
{\hyperref[\detokenize{Library_atoms:_CPPv419lammps_extract_atomPvPKc}]{\sphinxcrossref{\sphinxcode{\sphinxupquote{lammps\_extract\_atom()}}}}} function.  Please note that the array
stores holds the forces for \sphinxstyleemphasis{local} atoms for each MPI ranks, in the
order determined by the neighbor list build.  Because the underlying
data structures can change as well as the order of atom as they migrate
between MPI processes because of the domain decomposition
parallelization, this function should be always called immediately
before the forces are going to be set to get an up\sphinxhyphen{}to\sphinxhyphen{}date pointer.
You can use, for example, {\hyperref[\detokenize{Library_properties:_CPPv422lammps_extract_settingPvPKc}]{\sphinxcrossref{\sphinxcode{\sphinxupquote{lammps\_extract\_setting()}}}}} to obtain
the number of local atoms \sphinxtitleref{nlocal} and then assume the dimensions of
the returned force array as \sphinxcode{\sphinxupquote{double force{[}nlocal{]}{[}3{]}}}.

\sphinxAtStartPar
This is an alternative to the callback mechanism in fix external set up
by {\hyperref[\detokenize{Library_utility:_CPPv432lammps_set_fix_external_callbackPvPKc16FixExternalFnPtrPv}]{\sphinxcrossref{\sphinxcode{\sphinxupquote{lammps\_set\_fix\_external\_callback()}}}}}. The main difference is
that this mechanism can be used when forces are be pre\sphinxhyphen{}computed and the
control alternates between LAMMPS and the external code, while the
callback mechanism can call the external code to compute the force when
the fix is triggered and needs them.

\sphinxAtStartPar
Please see the documentation for {\hyperref[\detokenize{fix_external::doc}]{\sphinxcrossref{\DUrole{doc}{fix external}}}} for
more information about how to use the fix and how to couple it with an
external code.

\begin{quote}\begin{description}
\sphinxlineitem{Parameters}\begin{itemize}
\item {} 
\sphinxAtStartPar
\sphinxstyleliteralstrong{\sphinxupquote{handle}} \textendash{} pointer to a previously created LAMMPS instance cast to \sphinxcode{\sphinxupquote{void *}}. 

\item {} 
\sphinxAtStartPar
\sphinxstyleliteralstrong{\sphinxupquote{id}} \textendash{} fix ID of fix external instance 

\end{itemize}

\sphinxlineitem{Returns}
\sphinxAtStartPar
a pointer to the per\sphinxhyphen{}atom force array allocated by the fix 

\end{description}\end{quote}

\end{fulllineitems}



\bigskip\hrule\bigskip

\index{lammps\_fix\_external\_set\_energy\_global (C++ function)@\spxentry{lammps\_fix\_external\_set\_energy\_global}\spxextra{C++ function}}

\begin{fulllineitems}
\phantomsection\label{\detokenize{Library_utility:_CPPv437lammps_fix_external_set_energy_globalPvPKcd}}\phantomsection\label{\detokenize{Library_utility:_CPPv337lammps_fix_external_set_energy_globalPvPKcd}}\phantomsection\label{\detokenize{Library_utility:_CPPv237lammps_fix_external_set_energy_globalPvPKcd}}\phantomsection\label{\detokenize{Library_utility:lammps_fix_external_set_energy_global__voidP.cCP.double}}
\pysigstartsignatures
\pysigstartmultiline
\pysiglinewithargsret{\phantomsection\label{\detokenize{Library_utility:library_8h_1a6ac4888e2c82846dcb9b1d013ef90f96}}\DUrole{kt}{void}\DUrole{w}{ }\sphinxbfcode{\sphinxupquote{\DUrole{n}{lammps\_fix\_external\_set\_energy\_global}}}}{\DUrole{kt}{void}\DUrole{w}{ }\DUrole{p}{*}\DUrole{n,sig-param}{handle}\sphinxparamcomma \DUrole{k}{const}\DUrole{w}{ }\DUrole{kt}{char}\DUrole{w}{ }\DUrole{p}{*}\DUrole{n,sig-param}{id}\sphinxparamcomma \DUrole{kt}{double}\DUrole{w}{ }\DUrole{n,sig-param}{eng}}{}
\pysigstopmultiline
\pysigstopsignatures
\sphinxAtStartPar
Set the global energy contribution for a fix external instance with the given ID.

\sphinxAtStartPar

\sphinxAtStartPar
\DUrole{versionmodified,added}{Added in version 28Jul2021.}

\sphinxAtStartPar
This is a companion function to {\hyperref[\detokenize{Library_utility:_CPPv432lammps_set_fix_external_callbackPvPKc16FixExternalFnPtrPv}]{\sphinxcrossref{\sphinxcode{\sphinxupquote{lammps\_set\_fix\_external\_callback()}}}}} and
{\hyperref[\detokenize{Library_utility:_CPPv429lammps_fix_external_get_forcePvPKc}]{\sphinxcrossref{\sphinxcode{\sphinxupquote{lammps\_fix\_external\_get\_force()}}}}} to also set the contribution
to the global energy from the external code.  The value of the \sphinxstyleemphasis{eng}
argument will be stored in the fix and applied on the current and all
following timesteps until changed by another call to this function.
The energy is in energy units as determined by the current {\hyperref[\detokenize{units::doc}]{\sphinxcrossref{\DUrole{doc}{units}}}}
settings and is the \sphinxstylestrong{total} energy of the contribution.  Thus when
running in parallel all MPI processes have to call this function with
the \sphinxstylestrong{same} value and this will be returned as scalar property of the
fix external instance when accessed in LAMMPS input commands or from
variables.

\sphinxAtStartPar
Please see the documentation for {\hyperref[\detokenize{fix_external::doc}]{\sphinxcrossref{\DUrole{doc}{fix external}}}} for
more information about how to use the fix and how to couple it with an
external code.

\begin{quote}\begin{description}
\sphinxlineitem{Parameters}\begin{itemize}
\item {} 
\sphinxAtStartPar
\sphinxstyleliteralstrong{\sphinxupquote{handle}} \textendash{} pointer to a previously created LAMMPS instance cast to \sphinxcode{\sphinxupquote{void *}}. 

\item {} 
\sphinxAtStartPar
\sphinxstyleliteralstrong{\sphinxupquote{id}} \textendash{} fix ID of fix external instance 

\item {} 
\sphinxAtStartPar
\sphinxstyleliteralstrong{\sphinxupquote{eng}} \textendash{} total energy to be added to the global energy 

\end{itemize}

\end{description}\end{quote}

\end{fulllineitems}



\bigskip\hrule\bigskip

\index{lammps\_fix\_external\_set\_energy\_peratom (C++ function)@\spxentry{lammps\_fix\_external\_set\_energy\_peratom}\spxextra{C++ function}}

\begin{fulllineitems}
\phantomsection\label{\detokenize{Library_utility:_CPPv438lammps_fix_external_set_energy_peratomPvPKcPd}}\phantomsection\label{\detokenize{Library_utility:_CPPv338lammps_fix_external_set_energy_peratomPvPKcPd}}\phantomsection\label{\detokenize{Library_utility:_CPPv238lammps_fix_external_set_energy_peratomPvPKcPd}}\phantomsection\label{\detokenize{Library_utility:lammps_fix_external_set_energy_peratom__voidP.cCP.doubleP}}
\pysigstartsignatures
\pysigstartmultiline
\pysiglinewithargsret{\phantomsection\label{\detokenize{Library_utility:library_8h_1aac2c2027c2779613dfe2ffa47205bbdf}}\DUrole{kt}{void}\DUrole{w}{ }\sphinxbfcode{\sphinxupquote{\DUrole{n}{lammps\_fix\_external\_set\_energy\_peratom}}}}{\DUrole{kt}{void}\DUrole{w}{ }\DUrole{p}{*}\DUrole{n,sig-param}{handle}\sphinxparamcomma \DUrole{k}{const}\DUrole{w}{ }\DUrole{kt}{char}\DUrole{w}{ }\DUrole{p}{*}\DUrole{n,sig-param}{id}\sphinxparamcomma \DUrole{kt}{double}\DUrole{w}{ }\DUrole{p}{*}\DUrole{n,sig-param}{eng}}{}
\pysigstopmultiline
\pysigstopsignatures
\sphinxAtStartPar
Set the per\sphinxhyphen{}atom energy contribution for a fix external instance with the given ID.

\sphinxAtStartPar

\sphinxAtStartPar
\DUrole{versionmodified,added}{Added in version 28Jul2021.}

\sphinxAtStartPar
This is a companion function to {\hyperref[\detokenize{Library_utility:_CPPv432lammps_set_fix_external_callbackPvPKc16FixExternalFnPtrPv}]{\sphinxcrossref{\sphinxcode{\sphinxupquote{lammps\_set\_fix\_external\_callback()}}}}}
to set the per\sphinxhyphen{}atom energy contribution due to the fix from the external code
as part of the callback function.  For this to work, the handle to the
LAMMPS object must be passed as the \sphinxstyleemphasis{ptr} argument when registering the
callback function.

\sphinxAtStartPar
Please see the documentation for {\hyperref[\detokenize{fix_external::doc}]{\sphinxcrossref{\DUrole{doc}{fix external}}}} for
more information about how to use the fix and how to couple it with an
external code.


\begin{sphinxadmonition}{note}{Note:}
\sphinxAtStartPar
This function is fully independent from {\hyperref[\detokenize{Library_utility:_CPPv437lammps_fix_external_set_energy_globalPvPKcd}]{\sphinxcrossref{\sphinxcode{\sphinxupquote{lammps\_fix\_external\_set\_energy\_global()}}}}}
and will \sphinxstylestrong{NOT} add any contributions to the global energy tally
and \sphinxstylestrong{NOT} check whether the sum of the contributions added here are
consistent with the global added energy.
\end{sphinxadmonition}
\begin{quote}\begin{description}
\sphinxlineitem{Parameters}\begin{itemize}
\item {} 
\sphinxAtStartPar
\sphinxstyleliteralstrong{\sphinxupquote{handle}} \textendash{} pointer to a previously created LAMMPS instance cast to \sphinxcode{\sphinxupquote{void *}}. 

\item {} 
\sphinxAtStartPar
\sphinxstyleliteralstrong{\sphinxupquote{id}} \textendash{} fix ID of fix external instance 

\item {} 
\sphinxAtStartPar
\sphinxstyleliteralstrong{\sphinxupquote{eng}} \textendash{} pointer to array of length nlocal with the energy to be added to the per\sphinxhyphen{}atom energy 

\end{itemize}

\end{description}\end{quote}

\end{fulllineitems}



\bigskip\hrule\bigskip

\index{lammps\_fix\_external\_set\_virial\_global (C++ function)@\spxentry{lammps\_fix\_external\_set\_virial\_global}\spxextra{C++ function}}

\begin{fulllineitems}
\phantomsection\label{\detokenize{Library_utility:_CPPv437lammps_fix_external_set_virial_globalPvPKcPd}}\phantomsection\label{\detokenize{Library_utility:_CPPv337lammps_fix_external_set_virial_globalPvPKcPd}}\phantomsection\label{\detokenize{Library_utility:_CPPv237lammps_fix_external_set_virial_globalPvPKcPd}}\phantomsection\label{\detokenize{Library_utility:lammps_fix_external_set_virial_global__voidP.cCP.doubleP}}
\pysigstartsignatures
\pysigstartmultiline
\pysiglinewithargsret{\phantomsection\label{\detokenize{Library_utility:library_8h_1af707a45ddf44addec9b8345546deb83e}}\DUrole{kt}{void}\DUrole{w}{ }\sphinxbfcode{\sphinxupquote{\DUrole{n}{lammps\_fix\_external\_set\_virial\_global}}}}{\DUrole{kt}{void}\DUrole{w}{ }\DUrole{p}{*}\DUrole{n,sig-param}{handle}\sphinxparamcomma \DUrole{k}{const}\DUrole{w}{ }\DUrole{kt}{char}\DUrole{w}{ }\DUrole{p}{*}\DUrole{n,sig-param}{id}\sphinxparamcomma \DUrole{kt}{double}\DUrole{w}{ }\DUrole{p}{*}\DUrole{n,sig-param}{virial}}{}
\pysigstopmultiline
\pysigstopsignatures
\sphinxAtStartPar
Set the global virial contribution for a fix external instance with the given ID.

\sphinxAtStartPar

\sphinxAtStartPar
\DUrole{versionmodified,added}{Added in version 28Jul2021.}

\sphinxAtStartPar
This is a companion function to {\hyperref[\detokenize{Library_utility:_CPPv432lammps_set_fix_external_callbackPvPKc16FixExternalFnPtrPv}]{\sphinxcrossref{\sphinxcode{\sphinxupquote{lammps\_set\_fix\_external\_callback()}}}}}
and {\hyperref[\detokenize{Library_utility:_CPPv429lammps_fix_external_get_forcePvPKc}]{\sphinxcrossref{\sphinxcode{\sphinxupquote{lammps\_fix\_external\_get\_force()}}}}} to set the contribution to
the global virial from the external code.

\sphinxAtStartPar
The 6 values of the \sphinxstyleemphasis{virial} array will be stored in the fix and applied
on the current and all following timesteps until changed by another call
to this function. The components of the virial need to be stored in the
order: \sphinxstyleemphasis{xx}, \sphinxstyleemphasis{yy}, \sphinxstyleemphasis{zz}, \sphinxstyleemphasis{xy}, \sphinxstyleemphasis{xz}, \sphinxstyleemphasis{yz}.  In LAMMPS the virial is
stored internally as \sphinxtitleref{stress*volume} in units of \sphinxtitleref{pressure*volume} as
determined by the current {\hyperref[\detokenize{units::doc}]{\sphinxcrossref{\DUrole{doc}{units}}}} settings and is the
\sphinxstylestrong{total} contribution.  Thus when running in parallel all MPI processes
have to call this function with the \sphinxstylestrong{same} value and this will then
be added by fix external.

\sphinxAtStartPar
Please see the documentation for {\hyperref[\detokenize{fix_external::doc}]{\sphinxcrossref{\DUrole{doc}{fix external}}}} for
more information about how to use the fix and how to couple it with an
external code.

\begin{quote}\begin{description}
\sphinxlineitem{Parameters}\begin{itemize}
\item {} 
\sphinxAtStartPar
\sphinxstyleliteralstrong{\sphinxupquote{handle}} \textendash{} pointer to a previously created LAMMPS instance cast to \sphinxcode{\sphinxupquote{void *}}. 

\item {} 
\sphinxAtStartPar
\sphinxstyleliteralstrong{\sphinxupquote{id}} \textendash{} fix ID of fix external instance 

\item {} 
\sphinxAtStartPar
\sphinxstyleliteralstrong{\sphinxupquote{virial}} \textendash{} the 6 global stress tensor components to be added to the global virial 

\end{itemize}

\end{description}\end{quote}

\end{fulllineitems}



\bigskip\hrule\bigskip

\index{lammps\_fix\_external\_set\_virial\_peratom (C++ function)@\spxentry{lammps\_fix\_external\_set\_virial\_peratom}\spxextra{C++ function}}

\begin{fulllineitems}
\phantomsection\label{\detokenize{Library_utility:_CPPv438lammps_fix_external_set_virial_peratomPvPKcPPd}}\phantomsection\label{\detokenize{Library_utility:_CPPv338lammps_fix_external_set_virial_peratomPvPKcPPd}}\phantomsection\label{\detokenize{Library_utility:_CPPv238lammps_fix_external_set_virial_peratomPvPKcPPd}}\phantomsection\label{\detokenize{Library_utility:lammps_fix_external_set_virial_peratom__voidP.cCP.doublePP}}
\pysigstartsignatures
\pysigstartmultiline
\pysiglinewithargsret{\phantomsection\label{\detokenize{Library_utility:library_8h_1ae85311d7c556c3c7aace6e329ab352a2}}\DUrole{kt}{void}\DUrole{w}{ }\sphinxbfcode{\sphinxupquote{\DUrole{n}{lammps\_fix\_external\_set\_virial\_peratom}}}}{\DUrole{kt}{void}\DUrole{w}{ }\DUrole{p}{*}\DUrole{n,sig-param}{handle}\sphinxparamcomma \DUrole{k}{const}\DUrole{w}{ }\DUrole{kt}{char}\DUrole{w}{ }\DUrole{p}{*}\DUrole{n,sig-param}{id}\sphinxparamcomma \DUrole{kt}{double}\DUrole{w}{ }\DUrole{p}{*}\DUrole{p}{*}\DUrole{n,sig-param}{virial}}{}
\pysigstopmultiline
\pysigstopsignatures
\sphinxAtStartPar
Set the per\sphinxhyphen{}atom virial contribution for a fix external instance with the given ID.

\sphinxAtStartPar

\sphinxAtStartPar
\DUrole{versionmodified,added}{Added in version 28Jul2021.}

\sphinxAtStartPar
This is a companion function to {\hyperref[\detokenize{Library_utility:_CPPv432lammps_set_fix_external_callbackPvPKc16FixExternalFnPtrPv}]{\sphinxcrossref{\sphinxcode{\sphinxupquote{lammps\_set\_fix\_external\_callback()}}}}}
to set the per\sphinxhyphen{}atom virial contribution due to the fix from the external code
as part of the callback function.  For this to work, the handle to the
LAMMPS object must be passed as the \sphinxstyleemphasis{ptr} argument when registering the
callback function.

\sphinxAtStartPar
The order and units of the per\sphinxhyphen{}atom stress tensor elements are the same
as for the global virial.  The code in fix external assumes the
dimensions of the per\sphinxhyphen{}atom virial array is \sphinxcode{\sphinxupquote{double virial{[}nlocal{]}{[}6{]}}}.

\sphinxAtStartPar
Please see the documentation for {\hyperref[\detokenize{fix_external::doc}]{\sphinxcrossref{\DUrole{doc}{fix external}}}} for
more information about how to use the fix and how to couple it with an
external code.


\begin{sphinxadmonition}{note}{Note:}
\sphinxAtStartPar
This function is fully independent from {\hyperref[\detokenize{Library_utility:_CPPv437lammps_fix_external_set_virial_globalPvPKcPd}]{\sphinxcrossref{\sphinxcode{\sphinxupquote{lammps\_fix\_external\_set\_virial\_global()}}}}}
and will \sphinxstylestrong{NOT} add any contributions to the global virial tally
and \sphinxstylestrong{NOT} check whether the sum of the contributions added here are
consistent with the global added virial.
\end{sphinxadmonition}
\begin{quote}\begin{description}
\sphinxlineitem{Parameters}\begin{itemize}
\item {} 
\sphinxAtStartPar
\sphinxstyleliteralstrong{\sphinxupquote{handle}} \textendash{} pointer to a previously created LAMMPS instance cast to \sphinxcode{\sphinxupquote{void *}}. 

\item {} 
\sphinxAtStartPar
\sphinxstyleliteralstrong{\sphinxupquote{id}} \textendash{} fix ID of fix external instance 

\item {} 
\sphinxAtStartPar
\sphinxstyleliteralstrong{\sphinxupquote{virial}} \textendash{} a list of nlocal entries with the 6 per\sphinxhyphen{}atom stress tensor components to be added to the per\sphinxhyphen{}atom virial 

\end{itemize}

\end{description}\end{quote}

\end{fulllineitems}



\bigskip\hrule\bigskip

\index{lammps\_fix\_external\_set\_vector\_length (C++ function)@\spxentry{lammps\_fix\_external\_set\_vector\_length}\spxextra{C++ function}}

\begin{fulllineitems}
\phantomsection\label{\detokenize{Library_utility:_CPPv437lammps_fix_external_set_vector_lengthPvPKci}}\phantomsection\label{\detokenize{Library_utility:_CPPv337lammps_fix_external_set_vector_lengthPvPKci}}\phantomsection\label{\detokenize{Library_utility:_CPPv237lammps_fix_external_set_vector_lengthPvPKci}}\phantomsection\label{\detokenize{Library_utility:lammps_fix_external_set_vector_length__voidP.cCP.i}}
\pysigstartsignatures
\pysigstartmultiline
\pysiglinewithargsret{\phantomsection\label{\detokenize{Library_utility:library_8h_1a65eccbce0a2c1154051c56296acc7598}}\DUrole{kt}{void}\DUrole{w}{ }\sphinxbfcode{\sphinxupquote{\DUrole{n}{lammps\_fix\_external\_set\_vector\_length}}}}{\DUrole{kt}{void}\DUrole{w}{ }\DUrole{p}{*}\DUrole{n,sig-param}{handle}\sphinxparamcomma \DUrole{k}{const}\DUrole{w}{ }\DUrole{kt}{char}\DUrole{w}{ }\DUrole{p}{*}\DUrole{n,sig-param}{id}\sphinxparamcomma \DUrole{kt}{int}\DUrole{w}{ }\DUrole{n,sig-param}{len}}{}
\pysigstopmultiline
\pysigstopsignatures
\sphinxAtStartPar
Set the vector length for a global vector stored with fix external for analysis

\sphinxAtStartPar

\sphinxAtStartPar
\DUrole{versionmodified,added}{Added in version 28Jul2021.}

\sphinxAtStartPar
This is a companion function to {\hyperref[\detokenize{Library_utility:_CPPv432lammps_set_fix_external_callbackPvPKc16FixExternalFnPtrPv}]{\sphinxcrossref{\sphinxcode{\sphinxupquote{lammps\_set\_fix\_external\_callback()}}}}} and
{\hyperref[\detokenize{Library_utility:_CPPv429lammps_fix_external_get_forcePvPKc}]{\sphinxcrossref{\sphinxcode{\sphinxupquote{lammps\_fix\_external\_get\_force()}}}}} to set the length of a global vector of
properties that will be stored with the fix via
{\hyperref[\detokenize{Library_utility:_CPPv430lammps_fix_external_set_vectorPvPKcid}]{\sphinxcrossref{\sphinxcode{\sphinxupquote{lammps\_fix\_external\_set\_vector()}}}}}.

\sphinxAtStartPar
This function needs to be called \sphinxstylestrong{before} a call to
{\hyperref[\detokenize{Library_utility:_CPPv430lammps_fix_external_set_vectorPvPKcid}]{\sphinxcrossref{\sphinxcode{\sphinxupquote{lammps\_fix\_external\_set\_vector()}}}}} and \sphinxstylestrong{before} a run or minimize
command. When running in parallel it must be called from \sphinxstylestrong{all} MPI
processes and with the same length parameter.

\sphinxAtStartPar
Please see the documentation for {\hyperref[\detokenize{fix_external::doc}]{\sphinxcrossref{\DUrole{doc}{fix external}}}} for
more information about how to use the fix and how to couple it with an
external code.

\begin{quote}\begin{description}
\sphinxlineitem{Parameters}\begin{itemize}
\item {} 
\sphinxAtStartPar
\sphinxstyleliteralstrong{\sphinxupquote{handle}} \textendash{} pointer to a previously created LAMMPS instance cast to \sphinxcode{\sphinxupquote{void *}}. 

\item {} 
\sphinxAtStartPar
\sphinxstyleliteralstrong{\sphinxupquote{id}} \textendash{} fix ID of fix external instance 

\item {} 
\sphinxAtStartPar
\sphinxstyleliteralstrong{\sphinxupquote{len}} \textendash{} length of the global vector to be stored with the fix 

\end{itemize}

\end{description}\end{quote}

\end{fulllineitems}



\bigskip\hrule\bigskip

\index{lammps\_fix\_external\_set\_vector (C++ function)@\spxentry{lammps\_fix\_external\_set\_vector}\spxextra{C++ function}}

\begin{fulllineitems}
\phantomsection\label{\detokenize{Library_utility:_CPPv430lammps_fix_external_set_vectorPvPKcid}}\phantomsection\label{\detokenize{Library_utility:_CPPv330lammps_fix_external_set_vectorPvPKcid}}\phantomsection\label{\detokenize{Library_utility:_CPPv230lammps_fix_external_set_vectorPvPKcid}}\phantomsection\label{\detokenize{Library_utility:lammps_fix_external_set_vector__voidP.cCP.i.double}}
\pysigstartsignatures
\pysigstartmultiline
\pysiglinewithargsret{\phantomsection\label{\detokenize{Library_utility:library_8h_1a627a37da0a0f1cd076cfac35ea488c18}}\DUrole{kt}{void}\DUrole{w}{ }\sphinxbfcode{\sphinxupquote{\DUrole{n}{lammps\_fix\_external\_set\_vector}}}}{\DUrole{kt}{void}\DUrole{w}{ }\DUrole{p}{*}\DUrole{n,sig-param}{handle}\sphinxparamcomma \DUrole{k}{const}\DUrole{w}{ }\DUrole{kt}{char}\DUrole{w}{ }\DUrole{p}{*}\DUrole{n,sig-param}{id}\sphinxparamcomma \DUrole{kt}{int}\DUrole{w}{ }\DUrole{n,sig-param}{idx}\sphinxparamcomma \DUrole{kt}{double}\DUrole{w}{ }\DUrole{n,sig-param}{val}}{}
\pysigstopmultiline
\pysigstopsignatures
\sphinxAtStartPar
Store a global vector value for a fix external instance with the given ID.

\sphinxAtStartPar

\sphinxAtStartPar
\DUrole{versionmodified,added}{Added in version 28Jul2021.}

\sphinxAtStartPar
This is a companion function to {\hyperref[\detokenize{Library_utility:_CPPv432lammps_set_fix_external_callbackPvPKc16FixExternalFnPtrPv}]{\sphinxcrossref{\sphinxcode{\sphinxupquote{lammps\_set\_fix\_external\_callback()}}}}} and
{\hyperref[\detokenize{Library_utility:_CPPv429lammps_fix_external_get_forcePvPKc}]{\sphinxcrossref{\sphinxcode{\sphinxupquote{lammps\_fix\_external\_get\_force()}}}}} to set the values of a global vector of
properties that will be stored with the fix.  And can be accessed from
within LAMMPS input commands (e.g., fix ave/time or variables) when used
in a vector context.

\sphinxAtStartPar
This function needs to be called \sphinxstylestrong{after} a call to
{\hyperref[\detokenize{Library_utility:_CPPv437lammps_fix_external_set_vector_lengthPvPKci}]{\sphinxcrossref{\sphinxcode{\sphinxupquote{lammps\_fix\_external\_set\_vector\_length()}}}}} and the  and \sphinxstylestrong{before} a run or minimize
command.  When running in parallel it must be called from \sphinxstylestrong{all} MPI
processes and with the \sphinxstylestrong{same} index and value parameters.  The value
is assumed to be extensive.

\sphinxAtStartPar
Please see the documentation for {\hyperref[\detokenize{fix_external::doc}]{\sphinxcrossref{\DUrole{doc}{fix external}}}} for
more information about how to use the fix and how to couple it with an
external code.


\begin{sphinxadmonition}{note}{Note:}
\sphinxAtStartPar
The index in the \sphinxstyleemphasis{idx} parameter is 1\sphinxhyphen{}based, i.e. the first element
is set with idx = 1 and the last element of the vector with idx = N,
where N is the value of the \sphinxstyleemphasis{len} parameter of the call to
{\hyperref[\detokenize{Library_utility:_CPPv437lammps_fix_external_set_vector_lengthPvPKci}]{\sphinxcrossref{\sphinxcode{\sphinxupquote{lammps\_fix\_external\_set\_vector\_length()}}}}}.
\end{sphinxadmonition}
\begin{quote}\begin{description}
\sphinxlineitem{Parameters}\begin{itemize}
\item {} 
\sphinxAtStartPar
\sphinxstyleliteralstrong{\sphinxupquote{handle}} \textendash{} pointer to a previously created LAMMPS instance cast to \sphinxcode{\sphinxupquote{void *}}. 

\item {} 
\sphinxAtStartPar
\sphinxstyleliteralstrong{\sphinxupquote{id}} \textendash{} fix ID of fix external instance 

\item {} 
\sphinxAtStartPar
\sphinxstyleliteralstrong{\sphinxupquote{idx}} \textendash{} 1\sphinxhyphen{}based index of in global vector 

\item {} 
\sphinxAtStartPar
\sphinxstyleliteralstrong{\sphinxupquote{val}} \textendash{} value to be stored in global vector 

\end{itemize}

\end{description}\end{quote}

\end{fulllineitems}



\bigskip\hrule\bigskip

\index{lammps\_flush\_buffers (C++ function)@\spxentry{lammps\_flush\_buffers}\spxextra{C++ function}}

\begin{fulllineitems}
\phantomsection\label{\detokenize{Library_utility:_CPPv420lammps_flush_buffersPv}}\phantomsection\label{\detokenize{Library_utility:_CPPv320lammps_flush_buffersPv}}\phantomsection\label{\detokenize{Library_utility:_CPPv220lammps_flush_buffersPv}}\phantomsection\label{\detokenize{Library_utility:lammps_flush_buffers__voidP}}
\pysigstartsignatures
\pysigstartmultiline
\pysiglinewithargsret{\phantomsection\label{\detokenize{Library_utility:library_8h_1af54eb973fef3aa2bfe7708a0cf1139ff}}\DUrole{kt}{void}\DUrole{w}{ }\sphinxbfcode{\sphinxupquote{\DUrole{n}{lammps\_flush\_buffers}}}}{\DUrole{kt}{void}\DUrole{w}{ }\DUrole{p}{*}\DUrole{n,sig-param}{ptr}}{}
\pysigstopmultiline
\pysigstopsignatures
\sphinxAtStartPar
Flush output buffers

\sphinxAtStartPar
This function can be used to flush buffered output to be written to screen
and logfile pointers to simplify capturing output from LAMMPS library calls.
\begin{quote}\begin{description}
\sphinxlineitem{Parameters}
\sphinxAtStartPar
\sphinxstyleliteralstrong{\sphinxupquote{handle}} \textendash{} pointer to a previously created LAMMPS instance cast to \sphinxcode{\sphinxupquote{void *}}. 

\end{description}\end{quote}

\end{fulllineitems}



\bigskip\hrule\bigskip

\index{lammps\_free (C++ function)@\spxentry{lammps\_free}\spxextra{C++ function}}

\begin{fulllineitems}
\phantomsection\label{\detokenize{Library_utility:_CPPv411lammps_freePv}}\phantomsection\label{\detokenize{Library_utility:_CPPv311lammps_freePv}}\phantomsection\label{\detokenize{Library_utility:_CPPv211lammps_freePv}}\phantomsection\label{\detokenize{Library_utility:lammps_free__voidP}}
\pysigstartsignatures
\pysigstartmultiline
\pysiglinewithargsret{\phantomsection\label{\detokenize{Library_utility:library_8h_1aa521d27718fc92e9a58a255ae9f9cf87}}\DUrole{kt}{void}\DUrole{w}{ }\sphinxbfcode{\sphinxupquote{\DUrole{n}{lammps\_free}}}}{\DUrole{kt}{void}\DUrole{w}{ }\DUrole{p}{*}\DUrole{n,sig-param}{ptr}}{}
\pysigstopmultiline
\pysigstopsignatures
\sphinxAtStartPar
Free memory buffer allocated by LAMMPS.

\sphinxAtStartPar
Some of the LAMMPS C library interface functions return data as pointer
to a buffer that has been allocated by LAMMPS or the library interface.
This function can be used to delete those in order to avoid memory
leaks.
\begin{quote}\begin{description}
\sphinxlineitem{Parameters}
\sphinxAtStartPar
\sphinxstyleliteralstrong{\sphinxupquote{ptr}} \textendash{} pointer to data allocated by LAMMPS 

\end{description}\end{quote}

\end{fulllineitems}



\bigskip\hrule\bigskip

\index{lammps\_is\_running (C++ function)@\spxentry{lammps\_is\_running}\spxextra{C++ function}}

\begin{fulllineitems}
\phantomsection\label{\detokenize{Library_utility:_CPPv417lammps_is_runningPv}}\phantomsection\label{\detokenize{Library_utility:_CPPv317lammps_is_runningPv}}\phantomsection\label{\detokenize{Library_utility:_CPPv217lammps_is_runningPv}}\phantomsection\label{\detokenize{Library_utility:lammps_is_running__voidP}}
\pysigstartsignatures
\pysigstartmultiline
\pysiglinewithargsret{\phantomsection\label{\detokenize{Library_utility:library_8h_1a96130bd0ac1fcfbb01f56fe92f8cf7da}}\DUrole{kt}{int}\DUrole{w}{ }\sphinxbfcode{\sphinxupquote{\DUrole{n}{lammps\_is\_running}}}}{\DUrole{kt}{void}\DUrole{w}{ }\DUrole{p}{*}\DUrole{n,sig-param}{handle}}{}
\pysigstopmultiline
\pysigstopsignatures
\sphinxAtStartPar
Check if LAMMPS is currently inside a run or minimization

\sphinxAtStartPar
This function can be used from signal handlers or multi\sphinxhyphen{}threaded applications to determine if the LAMMPS instance is currently active.
\begin{quote}\begin{description}
\sphinxlineitem{Parameters}
\sphinxAtStartPar
\sphinxstyleliteralstrong{\sphinxupquote{handle}} \textendash{} pointer to a previously created LAMMPS instance cast to \sphinxcode{\sphinxupquote{void *}}. 

\sphinxlineitem{Returns}
\sphinxAtStartPar
0 if idle or \textgreater{}0 if active 

\end{description}\end{quote}

\end{fulllineitems}



\bigskip\hrule\bigskip

\index{lammps\_force\_timeout (C++ function)@\spxentry{lammps\_force\_timeout}\spxextra{C++ function}}

\begin{fulllineitems}
\phantomsection\label{\detokenize{Library_utility:_CPPv420lammps_force_timeoutPv}}\phantomsection\label{\detokenize{Library_utility:_CPPv320lammps_force_timeoutPv}}\phantomsection\label{\detokenize{Library_utility:_CPPv220lammps_force_timeoutPv}}\phantomsection\label{\detokenize{Library_utility:lammps_force_timeout__voidP}}
\pysigstartsignatures
\pysigstartmultiline
\pysiglinewithargsret{\phantomsection\label{\detokenize{Library_utility:library_8h_1a298dbff8f43185d82bcbf0a8d76b9aa9}}\DUrole{kt}{void}\DUrole{w}{ }\sphinxbfcode{\sphinxupquote{\DUrole{n}{lammps\_force\_timeout}}}}{\DUrole{kt}{void}\DUrole{w}{ }\DUrole{p}{*}\DUrole{n,sig-param}{handle}}{}
\pysigstopmultiline
\pysigstopsignatures
\sphinxAtStartPar
Force a timeout to stop an ongoing run cleanly.

\sphinxAtStartPar
This function can be used from signal handlers or multi\sphinxhyphen{}threaded applications to cleanly terminate an ongoing run.
\begin{quote}\begin{description}
\sphinxlineitem{Parameters}
\sphinxAtStartPar
\sphinxstyleliteralstrong{\sphinxupquote{handle}} \textendash{} pointer to a previously created LAMMPS instance cast to \sphinxcode{\sphinxupquote{void *}}

\end{description}\end{quote}

\end{fulllineitems}



\bigskip\hrule\bigskip

\index{lammps\_has\_error (C++ function)@\spxentry{lammps\_has\_error}\spxextra{C++ function}}

\begin{fulllineitems}
\phantomsection\label{\detokenize{Library_utility:_CPPv416lammps_has_errorPv}}\phantomsection\label{\detokenize{Library_utility:_CPPv316lammps_has_errorPv}}\phantomsection\label{\detokenize{Library_utility:_CPPv216lammps_has_errorPv}}\phantomsection\label{\detokenize{Library_utility:lammps_has_error__voidP}}
\pysigstartsignatures
\pysigstartmultiline
\pysiglinewithargsret{\phantomsection\label{\detokenize{Library_utility:library_8h_1a6f2f23eeafb8d56983ef90ad55e0c37d}}\DUrole{kt}{int}\DUrole{w}{ }\sphinxbfcode{\sphinxupquote{\DUrole{n}{lammps\_has\_error}}}}{\DUrole{kt}{void}\DUrole{w}{ }\DUrole{p}{*}\DUrole{n,sig-param}{handle}}{}
\pysigstopmultiline
\pysigstopsignatures
\sphinxAtStartPar
Check if there is a (new) error message available

\sphinxAtStartPar
This function can be used to query if an error inside of LAMMPS
has thrown a {\hyperref[\detokenize{Build_settings:exceptions}]{\sphinxcrossref{\DUrole{std,std-ref}{C++ exception}}}}.

\begin{sphinxadmonition}{note}{Note:}
\sphinxAtStartPar
The \sphinxstyleemphasis{handle} pointer may be \sphinxcode{\sphinxupquote{NULL}} for this function, as would be
the case when a call to create a LAMMPS instance has failed.  Then
this function will not check the error status inside the LAMMPS
instance, but instead would check the global error buffer of the
library interface.
\end{sphinxadmonition}
\begin{quote}\begin{description}
\sphinxlineitem{Parameters}
\sphinxAtStartPar
\sphinxstyleliteralstrong{\sphinxupquote{handle}} \textendash{} pointer to a previously created LAMMPS instance cast to \sphinxcode{\sphinxupquote{void *}} or NULL 

\sphinxlineitem{Returns}
\sphinxAtStartPar
0 on no error, 1 on error. 

\end{description}\end{quote}

\end{fulllineitems}



\bigskip\hrule\bigskip

\index{lammps\_get\_last\_error\_message (C++ function)@\spxentry{lammps\_get\_last\_error\_message}\spxextra{C++ function}}

\begin{fulllineitems}
\phantomsection\label{\detokenize{Library_utility:_CPPv429lammps_get_last_error_messagePvPci}}\phantomsection\label{\detokenize{Library_utility:_CPPv329lammps_get_last_error_messagePvPci}}\phantomsection\label{\detokenize{Library_utility:_CPPv229lammps_get_last_error_messagePvPci}}\phantomsection\label{\detokenize{Library_utility:lammps_get_last_error_message__voidP.cP.i}}
\pysigstartsignatures
\pysigstartmultiline
\pysiglinewithargsret{\phantomsection\label{\detokenize{Library_utility:library_8h_1a1900bdb6311f7f9649b5c82e20257cff}}\DUrole{kt}{int}\DUrole{w}{ }\sphinxbfcode{\sphinxupquote{\DUrole{n}{lammps\_get\_last\_error\_message}}}}{\DUrole{kt}{void}\DUrole{w}{ }\DUrole{p}{*}\DUrole{n,sig-param}{handle}\sphinxparamcomma \DUrole{kt}{char}\DUrole{w}{ }\DUrole{p}{*}\DUrole{n,sig-param}{buffer}\sphinxparamcomma \DUrole{kt}{int}\DUrole{w}{ }\DUrole{n,sig-param}{buf\_size}}{}
\pysigstopmultiline
\pysigstopsignatures
\sphinxAtStartPar
Copy the last error message into the provided buffer

\sphinxAtStartPar
This function can be used to retrieve the error message that was set
in the event of an error inside of LAMMPS which resulted in a
{\hyperref[\detokenize{Build_settings:exceptions}]{\sphinxcrossref{\DUrole{std,std-ref}{C++ exception}}}}.  A suitable buffer for a C\sphinxhyphen{}style
string has to be provided and its length.  If the internally stored
error message is longer, it will be truncated accordingly.  If the
buffer is a NULL pointer, then nothing will be copied.  The return
value of the function corresponds to the kind of error: a “1” indicates
an error that occurred on all MPI ranks and is often recoverable, while
a “2” indicates an abort that would happen only in a single MPI rank
and thus may not be recoverable, as other MPI ranks may be waiting on
the failing MPI ranks to send messages.

\begin{sphinxadmonition}{note}{Note:}
\sphinxAtStartPar
The \sphinxstyleemphasis{handle} pointer may be \sphinxcode{\sphinxupquote{NULL}} for this function, as would be
the case when a call to create a LAMMPS instance has failed.  Then
this function will not check the error buffer inside the LAMMPS
instance, but instead would check the global error buffer of the
library interface.

\sphinxAtStartPar
The \sphinxstyleemphasis{buffer} pointer may be \sphinxcode{\sphinxupquote{NULL}}.  This will clear any error
status without copying the error message.
\end{sphinxadmonition}
\begin{quote}\begin{description}
\sphinxlineitem{Parameters}\begin{itemize}
\item {} 
\sphinxAtStartPar
\sphinxstyleliteralstrong{\sphinxupquote{handle}} \textendash{} pointer to a previously created LAMMPS instance cast to \sphinxcode{\sphinxupquote{void *}} or NULL. 

\item {} 
\sphinxAtStartPar
\sphinxstyleliteralstrong{\sphinxupquote{buffer}} \textendash{} string buffer to copy the error message to, may be NULL 

\item {} 
\sphinxAtStartPar
\sphinxstyleliteralstrong{\sphinxupquote{buf\_size}} \textendash{} size of the provided string buffer 

\end{itemize}

\sphinxlineitem{Returns}
\sphinxAtStartPar
1 when all ranks had the error, 2 on a single rank error. 

\end{description}\end{quote}

\end{fulllineitems}



\bigskip\hrule\bigskip

\index{lammps\_python\_api\_version (C++ function)@\spxentry{lammps\_python\_api\_version}\spxextra{C++ function}}

\begin{fulllineitems}
\phantomsection\label{\detokenize{Library_utility:_CPPv425lammps_python_api_versionv}}\phantomsection\label{\detokenize{Library_utility:_CPPv325lammps_python_api_versionv}}\phantomsection\label{\detokenize{Library_utility:_CPPv225lammps_python_api_versionv}}\phantomsection\label{\detokenize{Library_utility:lammps_python_api_version}}
\pysigstartsignatures
\pysigstartmultiline
\pysiglinewithargsret{\phantomsection\label{\detokenize{Library_utility:library_8h_1ac696ed37ed39b295adbf05577111da90}}\DUrole{kt}{int}\DUrole{w}{ }\sphinxbfcode{\sphinxupquote{\DUrole{n}{lammps\_python\_api\_version}}}}{}{}
\pysigstopmultiline
\pysigstopsignatures
\sphinxAtStartPar
Return API version of embedded Python interpreter

\sphinxAtStartPar

\sphinxAtStartPar
\DUrole{versionmodified,added}{Added in version 3Nov2022.}

\sphinxAtStartPar
This function is used by the ML\sphinxhyphen{}IAP python code (mliappy) to verify
the API version of the embedded python interpreter of the PYTHON
package.  It returns \sphinxhyphen{}1 if the PYTHON package is not enabled.

\begin{quote}\begin{description}
\sphinxlineitem{Returns}
\sphinxAtStartPar
PYTHON\_API\_VERSION constant of the python interpreter or \sphinxhyphen{}1 

\end{description}\end{quote}

\end{fulllineitems}


\sphinxstepscope


\subsection{Extending the C API}
\label{\detokenize{Library_add:extending-the-c-api}}\label{\detokenize{Library_add::doc}}
\sphinxAtStartPar
The functionality of the LAMMPS library interface has historically
been motivated by the needs of its users.  Functions have been added
or expanded as they were needed and used.  Contributions to the
interface are always welcome.  However with a refactoring of the
library interface and its documentation that started in Spring 2020,
there are now a few requirements for including new changes or extensions.
\begin{itemize}
\item {} 
\sphinxAtStartPar
New functions should be orthogonal to existing ones and not
implement functionality that can already be achieved with the
existing APIs.

\item {} 
\sphinxAtStartPar
All changes and additions should be documented with
\sphinxhref{https://doxygen.nl}{Doxygen} style comments and references
to those functions added to the corresponding files in the
\sphinxcode{\sphinxupquote{doc/src}} folder.

\item {} 
\sphinxAtStartPar
If possible, new unit tests to test those new features should
be added.

\item {} 
\sphinxAtStartPar
New features should also be implemented and documented not just
for the C interface, but also the Python and Fortran interfaces.

\item {} 
\sphinxAtStartPar
All additions should work and be compatible with \sphinxcode{\sphinxupquote{\sphinxhyphen{}DLAMMPS\_BIGBIG}},
\sphinxcode{\sphinxupquote{\sphinxhyphen{}DLAMMPS\_SMALLBIG}}, \sphinxcode{\sphinxupquote{\sphinxhyphen{}DLAMMPS\_SMALLSMALL}} as well as when
compiling with and without MPI support.

\item {} 
\sphinxAtStartPar
The \sphinxcode{\sphinxupquote{library.h}} file should be kept compatible to C code at
a level similar to C89. Its interfaces may not reference any
custom data types (e.g. \sphinxcode{\sphinxupquote{bigint}}, \sphinxcode{\sphinxupquote{tagint}}, and so on) that
are only known inside of LAMMPS; instead \sphinxcode{\sphinxupquote{int}} and \sphinxcode{\sphinxupquote{int64\_t}}
should be used.

\item {} 
\sphinxAtStartPar
only use C style comments, not C++ style

\end{itemize}

\sphinxAtStartPar
Please note that these are not \sphinxstylestrong{strict} requirements, but the LAMMPS
developers very much appreciate, if they are followed and can assist
with implementing what is missing.  It helps maintaining the code base
and keeping it consistent.


\bigskip\hrule\bigskip



\section{LAMMPS Python APIs}
\label{\detokenize{Library:lammps-python-apis}}\label{\detokenize{Library:lammps-python-api}}
\sphinxAtStartPar
The LAMMPS Python module enables calling the LAMMPS C library API from
Python by dynamically loading functions in the LAMMPS shared library through
the \sphinxhref{https://docs.python.org/3/library/ctypes.html}{Python ctypes module}.
Because of the dynamic loading, it is \sphinxstylestrong{required} that LAMMPS is compiled
in {\hyperref[\detokenize{Build_basics:exe}]{\sphinxcrossref{\DUrole{std,std-ref}{“shared” mode}}}}.  The Python interface is object\sphinxhyphen{}oriented, but
otherwise tries to be very similar to the C library API.  Three different
Python classes to run LAMMPS are available and they build on each other.
More information on this is in the {\hyperref[\detokenize{Python_head::doc}]{\sphinxcrossref{\DUrole{doc}{Use Python with LAMMPS}}}}
section of the manual.  Use of the LAMMPS Python module is described in
{\hyperref[\detokenize{Python_module::doc}]{\sphinxcrossref{\DUrole{doc}{The lammps Python module}}}}.


\bigskip\hrule\bigskip



\section{LAMMPS Fortran API}
\label{\detokenize{Library:lammps-fortran-api}}\label{\detokenize{Library:id1}}
\sphinxAtStartPar
The LAMMPS Fortran module is a wrapper around calling functions from the
LAMMPS C library API.  This is done using the ISO\_C\_BINDING feature in
Fortran 2003.  The interface is object\sphinxhyphen{}oriented but otherwise tries to
be very similar to the C library API and the basic Python module.

\sphinxstepscope


\subsection{The \sphinxstyleliteralintitle{\sphinxupquote{LIBLAMMPS}} Fortran Module}
\label{\detokenize{Fortran:the-liblammps-fortran-module}}\label{\detokenize{Fortran::doc}}
\sphinxAtStartPar
The \sphinxcode{\sphinxupquote{LIBLAMMPS}} module provides an interface to call LAMMPS from
Fortran.  It is based on the LAMMPS C library interface and requires a
fully Fortran 2003\sphinxhyphen{}compatible compiler to be compiled.  It is designed
to be self\sphinxhyphen{}contained and not require any support functions written in C,
C++, or Fortran other than those in the C library interface and the
LAMMPS Fortran module itself.

\sphinxAtStartPar
While C libraries have a defined binary interface (ABI) and can thus be
used from multiple compiler versions from different vendors as long as
they are compatible with the hosting operating system, the same is not
true for Fortran programs.  Thus, the LAMMPS Fortran module needs to be
compiled alongside the code using it from the source code in
\sphinxcode{\sphinxupquote{fortran/lammps.f90}} \sphinxstyleemphasis{and} with the same compiler used to build the
rest of the Fortran code that interfaces to LAMMPS.  When linking, you
also need to {\hyperref[\detokenize{Build_link::doc}]{\sphinxcrossref{\DUrole{doc}{link to the LAMMPS library}}}}.  A typical
command line for a simple program using the Fortran interface would be:

\begin{sphinxVerbatim}[commandchars=\\\{\}]
mpifort\PYG{+w}{ }\PYGZhy{}o\PYG{+w}{ }testlib.x\PYG{+w}{ }lammps.f90\PYG{+w}{ }testlib.f90\PYG{+w}{ }\PYGZhy{}L.\PYG{+w}{ }\PYGZhy{}llammps
\end{sphinxVerbatim}

\sphinxAtStartPar
Please note that the MPI compiler wrapper is only required when the
calling the library \sphinxstyleemphasis{from} an MPI\sphinxhyphen{}parallelized program.  Otherwise,
using the plain Fortran compiler (gfortran, ifort, flang, etc.) will
suffice, since there are no direct references to MPI library features,
definitions and subroutine calls; MPI communicators are referred to by
their integer index representation as required by the Fortran MPI
interface.  It may be necessary to link to additional libraries,
depending on how LAMMPS was configured and whether the LAMMPS library
{\hyperref[\detokenize{Build_link::doc}]{\sphinxcrossref{\DUrole{doc}{was compiled as a static or dynamic library}}}}.

\sphinxAtStartPar
If the LAMMPS library itself has been compiled with MPI support, the
resulting executable will be able to run LAMMPS in parallel with
\sphinxcode{\sphinxupquote{mpirun}}, \sphinxcode{\sphinxupquote{mpiexec}}, or equivalent.  This may be either on the
“world” communicator or a sub\sphinxhyphen{}communicator created by the calling
Fortran code.  If, on the other hand, the LAMMPS library has been
compiled \sphinxstylestrong{without} MPI support, each LAMMPS instance will run
independently using just one processor.

\sphinxAtStartPar
Please also note that the order of the source files matters: the
\sphinxcode{\sphinxupquote{lammps.f90}} file needs to be compiled first, since it provides the
\sphinxcode{\sphinxupquote{LIBLAMMPS}} module that would need to be imported by the calling
Fortran code in order to uses the Fortran interface.
A working example can be found together with equivalent examples in C and
C++ in the \sphinxcode{\sphinxupquote{examples/COUPLE/simple}} folder of the LAMMPS distribution.

\begin{sphinxadmonition}{note}{Fortran compiler compatibility}

\sphinxAtStartPar
A fully Fortran 2003 compatible Fortran compiler is required.
This means that currently only GNU Fortran 9 and later are
compatible and thus the default compilers of Red Hat or CentOS 7
and Ubuntu 18.04 LTS and not compatible.  Either newer compilers
need to be installed or the Linux updated.
\end{sphinxadmonition}


\bigskip\hrule\bigskip



\subsection{Creating or deleting a LAMMPS object}
\label{\detokenize{Fortran:creating-or-deleting-a-lammps-object}}
\sphinxAtStartPar
With the Fortran interface, the creation of a {\hyperref[\detokenize{Classes_lammps:_CPPv4N9LAMMPS_NS6LAMMPSE}]{\sphinxcrossref{\sphinxcode{\sphinxupquote{LAMMPS}}}}} instance is included in the constructor for
creating the \sphinxcode{\sphinxupquote{lammps()}} derived type.  To import the definition of
that type and its type\sphinxhyphen{}bound procedures, you need to add a \sphinxcode{\sphinxupquote{USE LIBLAMMPS}}
statement.  Internally, it will call either
{\hyperref[\detokenize{Library_create:_CPPv419lammps_open_fortraniPPci}]{\sphinxcrossref{\sphinxcode{\sphinxupquote{lammps\_open\_fortran()}}}}} or {\hyperref[\detokenize{Library_create:_CPPv418lammps_open_no_mpiiPPcPPv}]{\sphinxcrossref{\sphinxcode{\sphinxupquote{lammps\_open\_no\_mpi()}}}}} from
the C library API to create the class instance.  All arguments are
optional and {\hyperref[\detokenize{Library_create:_CPPv415lammps_mpi_initv}]{\sphinxcrossref{\sphinxcode{\sphinxupquote{lammps\_mpi\_init()}}}}} will be called automatically
if it is needed.  Similarly, a possible call to
{\hyperref[\detokenize{Library_create:_CPPv419lammps_mpi_finalizev}]{\sphinxcrossref{\sphinxcode{\sphinxupquote{lammps\_mpi\_finalize()}}}}} is integrated into the \sphinxcode{\sphinxupquote{close()}}
function and triggered with the optional logical argument set to
\sphinxcode{\sphinxupquote{.TRUE.}}. Here is a simple example:

\begin{sphinxVerbatim}[commandchars=\\\{\}]
\PYG{k}{PROGRAM }\PYG{n}{testlib}
\PYG{+w}{  }\PYG{k}{USE }\PYG{n}{LIBLAMMPS}\PYG{+w}{                 }\PYG{c}{! include the LAMMPS library interface}
\PYG{+w}{  }\PYG{k}{IMPLICIT }\PYG{k}{NONE}
\PYG{k}{  }\PYG{k}{TYPE}\PYG{p}{(}\PYG{n}{lammps}\PYG{p}{)}\PYG{+w}{ }\PYG{k+kd}{::}\PYG{+w}{ }\PYG{n}{lmp}\PYG{+w}{           }\PYG{c}{! derived type to hold LAMMPS instance}
\PYG{+w}{  }\PYG{k+kt}{CHARACTER}\PYG{p}{(}\PYG{n+nb}{LEN}\PYG{o}{=}\PYG{l+m+mi}{12}\PYG{p}{)}\PYG{p}{,}\PYG{+w}{ }\PYG{k}{PARAMETER}\PYG{+w}{ }\PYG{k+kd}{::}\PYG{+w}{ }\PYG{n}{args}\PYG{p}{(}\PYG{l+m+mi}{3}\PYG{p}{)}\PYG{+w}{ }\PYG{o}{=}\PYG{+w}{ }\PYG{p}{\PYGZam{}}
\PYG{+w}{      }\PYG{p}{[}\PYG{+w}{ }\PYG{k+kt}{CHARACTER}\PYG{p}{(}\PYG{n+nb}{LEN}\PYG{o}{=}\PYG{l+m+mi}{12}\PYG{p}{)}\PYG{+w}{ }\PYG{k+kd}{::}\PYG{+w}{ }\PYG{l+s+s1}{\PYGZsq{}liblammps\PYGZsq{}}\PYG{p}{,}\PYG{+w}{ }\PYG{l+s+s1}{\PYGZsq{}\PYGZhy{}log\PYGZsq{}}\PYG{p}{,}\PYG{+w}{ }\PYG{l+s+s1}{\PYGZsq{}none\PYGZsq{}}\PYG{+w}{ }\PYG{p}{]}

\PYG{+w}{  }\PYG{c}{! create a LAMMPS instance (and initialize MPI)}
\PYG{+w}{  }\PYG{n}{lmp}\PYG{+w}{ }\PYG{o}{=}\PYG{+w}{ }\PYG{n}{lammps}\PYG{p}{(}\PYG{n}{args}\PYG{p}{)}
\PYG{+w}{  }\PYG{c}{! get and print numerical version code}
\PYG{+w}{  }\PYG{k}{PRINT}\PYG{o}{*}\PYG{p}{,}\PYG{+w}{ }\PYG{l+s+s1}{\PYGZsq{}LAMMPS Version: \PYGZsq{}}\PYG{p}{,}\PYG{+w}{ }\PYG{n}{lmp}\PYG{p}{\PYGZpc{}}\PYG{n}{version}\PYG{p}{(}\PYG{p}{)}
\PYG{+w}{  }\PYG{c}{! delete LAMMPS instance (and shutdown MPI)}
\PYG{+w}{  }\PYG{k}{CALL }\PYG{n}{lmp}\PYG{p}{\PYGZpc{}}\PYG{k}{close}\PYG{p}{(}\PYG{p}{.}\PYG{n}{TRUE}\PYG{p}{.}\PYG{p}{)}
\PYG{k}{END }\PYG{k}{PROGRAM }\PYG{n}{testlib}
\end{sphinxVerbatim}

\sphinxAtStartPar
It is also possible to pass command line flags from Fortran to C/C++ and
thus make the resulting executable behave similarly to the standalone
executable (it will ignore the \sphinxtitleref{\sphinxhyphen{}in/\sphinxhyphen{}i} flag, though).  This allows
using the command line to configure accelerator and suffix settings,
configure screen and logfile output, or to set index style variables
from the command line and more.  Here is a correspondingly adapted
version of the previous example:

\begin{sphinxVerbatim}[commandchars=\\\{\}]
\PYG{k}{PROGRAM }\PYG{n}{testlib2}
\PYG{+w}{  }\PYG{k}{USE }\PYG{n}{LIBLAMMPS}\PYG{+w}{                 }\PYG{c}{! include the LAMMPS library interface}
\PYG{+w}{  }\PYG{k}{IMPLICIT }\PYG{k}{NONE}
\PYG{k}{  }\PYG{k}{TYPE}\PYG{p}{(}\PYG{n}{lammps}\PYG{p}{)}\PYG{+w}{ }\PYG{k+kd}{::}\PYG{+w}{ }\PYG{n}{lmp}\PYG{+w}{           }\PYG{c}{! derived type to hold LAMMPS instance}
\PYG{+w}{  }\PYG{k+kt}{CHARACTER}\PYG{p}{(}\PYG{n+nb}{LEN}\PYG{o}{=}\PYG{l+m+mi}{128}\PYG{p}{)}\PYG{p}{,}\PYG{+w}{ }\PYG{k}{ALLOCATABLE}\PYG{+w}{ }\PYG{k+kd}{::}\PYG{+w}{ }\PYG{n}{command\PYGZus{}args}\PYG{p}{(}\PYG{p}{:}\PYG{p}{)}
\PYG{+w}{  }\PYG{k+kt}{INTEGER}\PYG{+w}{ }\PYG{k+kd}{::}\PYG{+w}{ }\PYG{n}{i}\PYG{p}{,}\PYG{+w}{ }\PYG{n}{argc}

\PYG{+w}{  }\PYG{c}{! copy command line flags to `command\PYGZus{}args()`}
\PYG{+w}{  }\PYG{n}{argc}\PYG{+w}{ }\PYG{o}{=}\PYG{+w}{ }\PYG{n+nb}{COMMAND\PYGZus{}ARGUMENT\PYGZus{}COUNT}\PYG{p}{(}\PYG{p}{)}
\PYG{+w}{  }\PYG{k}{ALLOCATE}\PYG{p}{(}\PYG{n}{command\PYGZus{}args}\PYG{p}{(}\PYG{l+m+mi}{0}\PYG{p}{:}\PYG{n}{argc}\PYG{p}{)}\PYG{p}{)}
\PYG{+w}{  }\PYG{k}{DO }\PYG{n}{i}\PYG{o}{=}\PYG{l+m+mi}{0}\PYG{p}{,}\PYG{+w}{ }\PYG{n}{argc}
\PYG{+w}{    }\PYG{k}{CALL }\PYG{n+nb}{GET\PYGZus{}COMMAND\PYGZus{}ARGUMENT}\PYG{p}{(}\PYG{n}{i}\PYG{p}{,}\PYG{+w}{ }\PYG{n}{command\PYGZus{}args}\PYG{p}{(}\PYG{n}{i}\PYG{p}{)}\PYG{p}{)}
\PYG{+w}{  }\PYG{k}{END }\PYG{k}{DO}

\PYG{+w}{  }\PYG{c}{! create a LAMMPS instance (and initialize MPI)}
\PYG{+w}{  }\PYG{n}{lmp}\PYG{+w}{ }\PYG{o}{=}\PYG{+w}{ }\PYG{n}{lammps}\PYG{p}{(}\PYG{n}{command\PYGZus{}args}\PYG{p}{)}
\PYG{+w}{  }\PYG{c}{! get and print numerical version code}
\PYG{+w}{  }\PYG{k}{PRINT}\PYG{o}{*}\PYG{p}{,}\PYG{+w}{ }\PYG{l+s+s1}{\PYGZsq{}Program name:   \PYGZsq{}}\PYG{p}{,}\PYG{+w}{ }\PYG{n}{command\PYGZus{}args}\PYG{p}{(}\PYG{l+m+mi}{0}\PYG{p}{)}
\PYG{+w}{  }\PYG{k}{PRINT}\PYG{o}{*}\PYG{p}{,}\PYG{+w}{ }\PYG{l+s+s1}{\PYGZsq{}LAMMPS Version: \PYGZsq{}}\PYG{p}{,}\PYG{+w}{ }\PYG{n}{lmp}\PYG{p}{\PYGZpc{}}\PYG{n}{version}\PYG{p}{(}\PYG{p}{)}
\PYG{+w}{  }\PYG{c}{! delete LAMMPS instance (and shuts down MPI)}
\PYG{+w}{  }\PYG{k}{CALL }\PYG{n}{lmp}\PYG{p}{\PYGZpc{}}\PYG{k}{close}\PYG{p}{(}\PYG{p}{.}\PYG{n}{TRUE}\PYG{p}{.}\PYG{p}{)}
\PYG{+w}{  }\PYG{k}{DEALLOCATE}\PYG{p}{(}\PYG{n}{command\PYGZus{}args}\PYG{p}{)}
\PYG{k}{END }\PYG{k}{PROGRAM }\PYG{n}{testlib2}
\end{sphinxVerbatim}


\bigskip\hrule\bigskip



\subsection{Executing LAMMPS commands}
\label{\detokenize{Fortran:executing-lammps-commands}}
\sphinxAtStartPar
Once a LAMMPS instance is created, it is possible to “drive” the LAMMPS
simulation by telling LAMMPS to read commands from a file or to pass
individual or multiple commands from strings or lists of strings.  This
is done similarly to how it is implemented in the {\hyperref[\detokenize{Library_execute::doc}]{\sphinxcrossref{\DUrole{doc}{C library
interface}}}}. Before handing off the calls to the
C library interface, the corresponding Fortran versions of the calls
(\sphinxcode{\sphinxupquote{file()}}, \sphinxcode{\sphinxupquote{command()}}, \sphinxcode{\sphinxupquote{commands\_list()}}, and
\sphinxcode{\sphinxupquote{commands\_string()}}) have to make copies of the strings passed as
arguments so that they can be modified to be compatible with the
requirements of strings in C without affecting the original strings.
Those copies are automatically deleted after the functions return.
Below is a small demonstration of the uses of the different functions.

\begin{sphinxVerbatim}[commandchars=\\\{\}]
\PYG{k}{PROGRAM }\PYG{n}{testcmd}
\PYG{+w}{  }\PYG{k}{USE }\PYG{n}{LIBLAMMPS}
\PYG{+w}{  }\PYG{k}{TYPE}\PYG{p}{(}\PYG{n}{lammps}\PYG{p}{)}\PYG{+w}{ }\PYG{k+kd}{::}\PYG{+w}{ }\PYG{n}{lmp}
\PYG{+w}{  }\PYG{k+kt}{CHARACTER}\PYG{p}{(}\PYG{n+nb}{LEN}\PYG{o}{=}\PYG{l+m+mi}{512}\PYG{p}{)}\PYG{+w}{ }\PYG{k+kd}{::}\PYG{+w}{ }\PYG{n}{cmds}
\PYG{+w}{  }\PYG{k+kt}{CHARACTER}\PYG{p}{(}\PYG{n+nb}{LEN}\PYG{o}{=}\PYG{l+m+mi}{40}\PYG{p}{)}\PYG{p}{,}\PYG{+w}{ }\PYG{k}{ALLOCATABLE}\PYG{+w}{ }\PYG{k+kd}{::}\PYG{+w}{ }\PYG{n}{cmdlist}\PYG{p}{(}\PYG{p}{:}\PYG{p}{)}
\PYG{+w}{  }\PYG{k+kt}{CHARACTER}\PYG{p}{(}\PYG{n+nb}{LEN}\PYG{o}{=}\PYG{l+m+mi}{10}\PYG{p}{)}\PYG{+w}{ }\PYG{k+kd}{::}\PYG{+w}{ }\PYG{n}{trimmed}
\PYG{+w}{  }\PYG{k+kt}{INTEGER}\PYG{+w}{ }\PYG{k+kd}{::}\PYG{+w}{ }\PYG{n}{i}

\PYG{+w}{  }\PYG{n}{lmp}\PYG{+w}{ }\PYG{o}{=}\PYG{+w}{ }\PYG{n}{lammps}\PYG{p}{(}\PYG{p}{)}
\PYG{+w}{  }\PYG{k}{CALL }\PYG{n}{lmp}\PYG{p}{\PYGZpc{}}\PYG{k}{file}\PYG{p}{(}\PYG{l+s+s1}{\PYGZsq{}in.melt\PYGZsq{}}\PYG{p}{)}
\PYG{+w}{  }\PYG{k}{CALL }\PYG{n}{lmp}\PYG{p}{\PYGZpc{}}\PYG{n}{command}\PYG{p}{(}\PYG{l+s+s1}{\PYGZsq{}variable zpos index 1.0\PYGZsq{}}\PYG{p}{)}
\PYG{+w}{  }\PYG{c}{! define 10 groups of 10 atoms each}
\PYG{+w}{  }\PYG{k}{ALLOCATE}\PYG{p}{(}\PYG{n}{cmdlist}\PYG{p}{(}\PYG{l+m+mi}{10}\PYG{p}{)}\PYG{p}{)}
\PYG{+w}{  }\PYG{k}{DO }\PYG{n}{i}\PYG{o}{=}\PYG{l+m+mi}{1}\PYG{p}{,}\PYG{+w}{ }\PYG{l+m+mi}{10}
\PYG{+w}{    }\PYG{k}{WRITE}\PYG{p}{(}\PYG{n}{trimmed}\PYG{p}{,}\PYG{l+s+s1}{\PYGZsq{}(I10)\PYGZsq{}}\PYG{p}{)}\PYG{+w}{ }\PYG{l+m+mi}{10}\PYG{o}{*}\PYG{n}{i}
\PYG{+w}{    }\PYG{k}{WRITE}\PYG{p}{(}\PYG{n}{cmdlist}\PYG{p}{(}\PYG{n}{i}\PYG{p}{)}\PYG{p}{,}\PYG{l+s+s1}{\PYGZsq{}(A,I1,A,I10,A,A)\PYGZsq{}}\PYG{p}{)}\PYG{+w}{       }\PYG{p}{\PYGZam{}}
\PYG{+w}{        }\PYG{l+s+s1}{\PYGZsq{}group g\PYGZsq{}}\PYG{p}{,}\PYG{+w}{ }\PYG{n}{i}\PYG{o}{\PYGZhy{}}\PYG{l+m+mi}{1}\PYG{p}{,}\PYG{+w}{ }\PYG{l+s+s1}{\PYGZsq{} id \PYGZsq{}}\PYG{p}{,}\PYG{+w}{ }\PYG{l+m+mi}{10}\PYG{o}{*}\PYG{p}{(}\PYG{n}{i}\PYG{o}{\PYGZhy{}}\PYG{l+m+mi}{1}\PYG{p}{)}\PYG{o}{+}\PYG{l+m+mi}{1}\PYG{p}{,}\PYG{+w}{ }\PYG{l+s+s1}{\PYGZsq{}:\PYGZsq{}}\PYG{p}{,}\PYG{+w}{ }\PYG{n+nb}{ADJUSTL}\PYG{p}{(}\PYG{n}{trimmed}\PYG{p}{)}
\PYG{+w}{  }\PYG{k}{END }\PYG{k}{DO}
\PYG{k}{  }\PYG{k}{CALL }\PYG{n}{lmp}\PYG{p}{\PYGZpc{}}\PYG{n}{commands\PYGZus{}list}\PYG{p}{(}\PYG{n}{cmdlist}\PYG{p}{)}
\PYG{+w}{  }\PYG{c}{! run multiple commands from multi\PYGZhy{}line string}
\PYG{+w}{  }\PYG{n}{cmds}\PYG{+w}{ }\PYG{o}{=}\PYG{+w}{ }\PYG{l+s+s1}{\PYGZsq{}clear\PYGZsq{}}\PYG{+w}{ }\PYG{o}{/}\PYG{o}{/}\PYG{+w}{ }\PYG{n+nb}{NEW\PYGZus{}LINE}\PYG{p}{(}\PYG{l+s+s1}{\PYGZsq{}A\PYGZsq{}}\PYG{p}{)}\PYG{+w}{ }\PYG{o}{/}\PYG{o}{/}\PYG{+w}{                       }\PYG{p}{\PYGZam{}}
\PYG{+w}{      }\PYG{l+s+s1}{\PYGZsq{}region  box block 0 2 0 2 0 2\PYGZsq{}}\PYG{+w}{ }\PYG{o}{/}\PYG{o}{/}\PYG{+w}{ }\PYG{n+nb}{NEW\PYGZus{}LINE}\PYG{p}{(}\PYG{l+s+s1}{\PYGZsq{}A\PYGZsq{}}\PYG{p}{)}\PYG{+w}{ }\PYG{o}{/}\PYG{o}{/}\PYG{+w}{  }\PYG{p}{\PYGZam{}}
\PYG{+w}{      }\PYG{l+s+s1}{\PYGZsq{}create\PYGZus{}box 1 box\PYGZsq{}}\PYG{+w}{ }\PYG{o}{/}\PYG{o}{/}\PYG{+w}{ }\PYG{n+nb}{NEW\PYGZus{}LINE}\PYG{p}{(}\PYG{l+s+s1}{\PYGZsq{}A\PYGZsq{}}\PYG{p}{)}\PYG{+w}{ }\PYG{o}{/}\PYG{o}{/}\PYG{+w}{               }\PYG{p}{\PYGZam{}}
\PYG{+w}{      }\PYG{l+s+s1}{\PYGZsq{}create\PYGZus{}atoms 1 single 1.0 1.0 \PYGZdl{}\PYGZob{}zpos\PYGZcb{}\PYGZsq{}}
\PYG{+w}{  }\PYG{k}{CALL }\PYG{n}{lmp}\PYG{p}{\PYGZpc{}}\PYG{n}{commands\PYGZus{}string}\PYG{p}{(}\PYG{n}{cmds}\PYG{p}{)}
\PYG{+w}{  }\PYG{k}{CALL }\PYG{n}{lmp}\PYG{p}{\PYGZpc{}}\PYG{k}{close}\PYG{p}{(}\PYG{p}{.}\PYG{n}{TRUE}\PYG{p}{.}\PYG{p}{)}
\PYG{k}{END }\PYG{k}{PROGRAM }\PYG{n}{testcmd}
\end{sphinxVerbatim}


\bigskip\hrule\bigskip



\subsection{Accessing system properties}
\label{\detokenize{Fortran:accessing-system-properties}}
\sphinxAtStartPar
The C library interface allows the {\hyperref[\detokenize{Library_properties::doc}]{\sphinxcrossref{\DUrole{doc}{extraction of different kinds
of information}}}} about the active simulation
instance and also—in some cases—to apply modifications to it, and the
Fortran interface provides access to the same data using Fortran\sphinxhyphen{}style,
C\sphinxhyphen{}interoperable data types.  In some cases, the Fortran library interface makes
pointers to internal LAMMPS data structures accessible; when accessing them
through the library interfaces, special care is needed to avoid data corruption
and crashes.  Please see the documentation of the individual type\sphinxhyphen{}bound
procedures for details.

\sphinxAtStartPar
Below is an example demonstrating some of the possible uses.

\begin{sphinxVerbatim}[commandchars=\\\{\}]
\PYG{k}{PROGRAM }\PYG{n}{testprop}
\PYG{+w}{  }\PYG{k}{USE }\PYG{n}{LIBLAMMPS}
\PYG{+w}{  }\PYG{k}{USE}\PYG{p}{,}\PYG{+w}{ }\PYG{k}{INTRINSIC}\PYG{+w}{ }\PYG{k+kd}{::}\PYG{+w}{ }\PYG{n+nb}{ISO\PYGZus{}C\PYGZus{}BINDING}\PYG{p}{,}\PYG{+w}{ }\PYG{k}{ONLY}\PYG{+w}{ }\PYG{p}{:}\PYG{+w}{ }\PYG{k+kt}{c\PYGZus{}double}\PYG{p}{,}\PYG{+w}{ }\PYG{k+kt}{c\PYGZus{}int64\PYGZus{}t}\PYG{p}{,}\PYG{+w}{ }\PYG{k+kt}{c\PYGZus{}int}
\PYG{k+kt}{  }\PYG{k}{USE}\PYG{p}{,}\PYG{+w}{ }\PYG{k}{INTRINSIC}\PYG{+w}{ }\PYG{k+kd}{::}\PYG{+w}{ }\PYG{n}{ISO\PYGZus{}FORTRAN\PYGZus{}ENV}\PYG{p}{,}\PYG{+w}{ }\PYG{k}{ONLY}\PYG{+w}{ }\PYG{p}{:}\PYG{+w}{ }\PYG{n}{OUTPUT\PYGZus{}UNIT}
\PYG{+w}{  }\PYG{k}{TYPE}\PYG{p}{(}\PYG{n}{lammps}\PYG{p}{)}\PYG{+w}{ }\PYG{k+kd}{::}\PYG{+w}{ }\PYG{n}{lmp}
\PYG{+w}{  }\PYG{k+kt}{INTEGER}\PYG{p}{(}\PYG{n+nb}{KIND}\PYG{o}{=}\PYG{k+kt}{c\PYGZus{}int64\PYGZus{}t}\PYG{p}{)}\PYG{p}{,}\PYG{+w}{ }\PYG{k}{POINTER}\PYG{+w}{ }\PYG{k+kd}{::}\PYG{+w}{ }\PYG{n}{natoms}\PYG{p}{,}\PYG{+w}{ }\PYG{n}{ntimestep}\PYG{p}{,}\PYG{+w}{ }\PYG{n}{bval}
\PYG{+w}{  }\PYG{k+kt}{REAL}\PYG{p}{(}\PYG{n+nb}{KIND}\PYG{o}{=}\PYG{k+kt}{c\PYGZus{}double}\PYG{p}{)}\PYG{p}{,}\PYG{+w}{ }\PYG{k}{POINTER}\PYG{+w}{ }\PYG{k+kd}{::}\PYG{+w}{ }\PYG{n}{dt}\PYG{p}{,}\PYG{+w}{ }\PYG{n}{dval}
\PYG{+w}{  }\PYG{k+kt}{INTEGER}\PYG{p}{(}\PYG{n+nb}{KIND}\PYG{o}{=}\PYG{k+kt}{c\PYGZus{}int}\PYG{p}{)}\PYG{p}{,}\PYG{+w}{ }\PYG{k}{POINTER}\PYG{+w}{ }\PYG{k+kd}{::}\PYG{+w}{ }\PYG{n}{nfield}\PYG{p}{,}\PYG{+w}{ }\PYG{n}{typ}\PYG{p}{,}\PYG{+w}{ }\PYG{n}{ival}
\PYG{+w}{  }\PYG{k+kt}{INTEGER}\PYG{p}{(}\PYG{n+nb}{KIND}\PYG{o}{=}\PYG{k+kt}{c\PYGZus{}int}\PYG{p}{)}\PYG{+w}{ }\PYG{k+kd}{::}\PYG{+w}{ }\PYG{n}{i}
\PYG{+w}{  }\PYG{k+kt}{CHARACTER}\PYG{p}{(}\PYG{n+nb}{LEN}\PYG{o}{=}\PYG{l+m+mi}{11}\PYG{p}{)}\PYG{+w}{ }\PYG{k+kd}{::}\PYG{+w}{ }\PYG{n}{key}
\PYG{+w}{  }\PYG{k+kt}{REAL}\PYG{p}{(}\PYG{n+nb}{KIND}\PYG{o}{=}\PYG{k+kt}{c\PYGZus{}double}\PYG{p}{)}\PYG{+w}{ }\PYG{k+kd}{::}\PYG{+w}{ }\PYG{n}{pe}\PYG{p}{,}\PYG{+w}{ }\PYG{n}{ke}

\PYG{+w}{  }\PYG{n}{lmp}\PYG{+w}{ }\PYG{o}{=}\PYG{+w}{ }\PYG{n}{lammps}\PYG{p}{(}\PYG{p}{)}
\PYG{+w}{  }\PYG{k}{CALL }\PYG{n}{lmp}\PYG{p}{\PYGZpc{}}\PYG{k}{file}\PYG{p}{(}\PYG{l+s+s1}{\PYGZsq{}in.sysinit\PYGZsq{}}\PYG{p}{)}
\PYG{+w}{  }\PYG{n}{natoms}\PYG{+w}{ }\PYG{o}{=}\PYG{+w}{ }\PYG{n}{lmp}\PYG{p}{\PYGZpc{}}\PYG{n}{extract\PYGZus{}global}\PYG{p}{(}\PYG{l+s+s1}{\PYGZsq{}natoms\PYGZsq{}}\PYG{p}{)}
\PYG{+w}{  }\PYG{k}{WRITE}\PYG{p}{(}\PYG{n}{OUTPUT\PYGZus{}UNIT}\PYG{p}{,}\PYG{l+s+s1}{\PYGZsq{}(A,I0,A)\PYGZsq{}}\PYG{p}{)}\PYG{+w}{ }\PYG{l+s+s1}{\PYGZsq{}Running a simulation with \PYGZsq{}}\PYG{p}{,}\PYG{+w}{ }\PYG{n}{natoms}\PYG{p}{,}\PYG{+w}{ }\PYG{l+s+s1}{\PYGZsq{} atoms\PYGZsq{}}
\PYG{+w}{  }\PYG{k}{WRITE}\PYG{p}{(}\PYG{n}{OUTPUT\PYGZus{}UNIT}\PYG{p}{,}\PYG{l+s+s1}{\PYGZsq{}(I0,A,I0,A,I0,A)\PYGZsq{}}\PYG{p}{)}\PYG{+w}{ }\PYG{n}{lmp}\PYG{p}{\PYGZpc{}}\PYG{n}{extract\PYGZus{}setting}\PYG{p}{(}\PYG{l+s+s1}{\PYGZsq{}nlocal\PYGZsq{}}\PYG{p}{)}\PYG{p}{,}\PYG{+w}{ }\PYG{p}{\PYGZam{}}
\PYG{+w}{      }\PYG{l+s+s1}{\PYGZsq{} local and \PYGZsq{}}\PYG{p}{,}\PYG{+w}{ }\PYG{n}{lmp}\PYG{p}{\PYGZpc{}}\PYG{n}{extract\PYGZus{}setting}\PYG{p}{(}\PYG{l+s+s1}{\PYGZsq{}nghost\PYGZsq{}}\PYG{p}{)}\PYG{p}{,}\PYG{+w}{ }\PYG{l+s+s1}{\PYGZsq{} ghost atoms. \PYGZsq{}}\PYG{p}{,}\PYG{+w}{ }\PYG{p}{\PYGZam{}}
\PYG{+w}{      }\PYG{n}{lmp}\PYG{p}{\PYGZpc{}}\PYG{n}{extract\PYGZus{}setting}\PYG{p}{(}\PYG{l+s+s1}{\PYGZsq{}ntypes\PYGZsq{}}\PYG{p}{)}\PYG{p}{,}\PYG{+w}{ }\PYG{l+s+s1}{\PYGZsq{} atom types\PYGZsq{}}

\PYG{+w}{  }\PYG{k}{CALL }\PYG{n}{lmp}\PYG{p}{\PYGZpc{}}\PYG{n}{command}\PYG{p}{(}\PYG{l+s+s1}{\PYGZsq{}run 2 post no\PYGZsq{}}\PYG{p}{)}

\PYG{+w}{  }\PYG{n}{ntimestep}\PYG{+w}{ }\PYG{o}{=}\PYG{+w}{ }\PYG{n}{lmp}\PYG{p}{\PYGZpc{}}\PYG{n}{last\PYGZus{}thermo}\PYG{p}{(}\PYG{l+s+s1}{\PYGZsq{}step\PYGZsq{}}\PYG{p}{,}\PYG{+w}{ }\PYG{l+m+mi}{0}\PYG{p}{)}
\PYG{+w}{  }\PYG{n}{nfield}\PYG{+w}{ }\PYG{o}{=}\PYG{+w}{ }\PYG{n}{lmp}\PYG{p}{\PYGZpc{}}\PYG{n}{last\PYGZus{}thermo}\PYG{p}{(}\PYG{l+s+s1}{\PYGZsq{}num\PYGZsq{}}\PYG{p}{,}\PYG{+w}{ }\PYG{l+m+mi}{0}\PYG{p}{)}
\PYG{+w}{  }\PYG{k}{WRITE}\PYG{p}{(}\PYG{n}{OUTPUT\PYGZus{}UNIT}\PYG{p}{,}\PYG{l+s+s1}{\PYGZsq{}(A,I0,A,I0)\PYGZsq{}}\PYG{p}{)}\PYG{+w}{ }\PYG{l+s+s1}{\PYGZsq{}Last thermo output on step: \PYGZsq{}}\PYG{p}{,}\PYG{+w}{ }\PYG{n}{ntimestep}\PYG{p}{,}\PYG{+w}{ }\PYG{p}{\PYGZam{}}
\PYG{+w}{      }\PYG{l+s+s1}{\PYGZsq{},  number of fields: \PYGZsq{}}\PYG{p}{,}\PYG{+w}{ }\PYG{n}{nfield}
\PYG{+w}{  }\PYG{k}{DO }\PYG{n}{i}\PYG{o}{=}\PYG{l+m+mi}{1}\PYG{p}{,}\PYG{+w}{ }\PYG{n}{nfield}
\PYG{+w}{      }\PYG{n}{key}\PYG{+w}{ }\PYG{o}{=}\PYG{+w}{ }\PYG{n}{lmp}\PYG{p}{\PYGZpc{}}\PYG{n}{last\PYGZus{}thermo}\PYG{p}{(}\PYG{l+s+s1}{\PYGZsq{}keyword\PYGZsq{}}\PYG{p}{,}\PYG{n}{i}\PYG{p}{)}
\PYG{+w}{      }\PYG{n}{typ}\PYG{+w}{ }\PYG{o}{=}\PYG{+w}{ }\PYG{n}{lmp}\PYG{p}{\PYGZpc{}}\PYG{n}{last\PYGZus{}thermo}\PYG{p}{(}\PYG{l+s+s1}{\PYGZsq{}type\PYGZsq{}}\PYG{p}{,}\PYG{n}{i}\PYG{p}{)}
\PYG{+w}{      }\PYG{k}{IF}\PYG{+w}{ }\PYG{p}{(}\PYG{n}{typ}\PYG{+w}{ }\PYG{o}{==}\PYG{+w}{ }\PYG{n}{lmp}\PYG{p}{\PYGZpc{}}\PYG{n}{dtype}\PYG{p}{\PYGZpc{}}\PYG{n}{i32}\PYG{p}{)}\PYG{+w}{ }\PYG{k}{THEN}
\PYG{k}{          }\PYG{n}{ival}\PYG{+w}{ }\PYG{o}{=}\PYG{+w}{ }\PYG{n}{lmp}\PYG{p}{\PYGZpc{}}\PYG{n}{last\PYGZus{}thermo}\PYG{p}{(}\PYG{l+s+s1}{\PYGZsq{}data\PYGZsq{}}\PYG{p}{,}\PYG{n}{i}\PYG{p}{)}
\PYG{+w}{          }\PYG{k}{WRITE}\PYG{p}{(}\PYG{n}{OUTPUT\PYGZus{}UNIT}\PYG{p}{,}\PYG{o}{*}\PYG{p}{)}\PYG{+w}{ }\PYG{n}{key}\PYG{p}{,}\PYG{+w}{ }\PYG{l+s+s1}{\PYGZsq{}:\PYGZsq{}}\PYG{p}{,}\PYG{+w}{ }\PYG{n}{ival}
\PYG{+w}{      }\PYG{k}{ELSE }\PYG{k}{IF}\PYG{+w}{ }\PYG{p}{(}\PYG{n}{typ}\PYG{+w}{ }\PYG{o}{==}\PYG{+w}{ }\PYG{n}{lmp}\PYG{p}{\PYGZpc{}}\PYG{n}{dtype}\PYG{p}{\PYGZpc{}}\PYG{n}{i64}\PYG{p}{)}\PYG{+w}{ }\PYG{k}{THEN}
\PYG{k}{          }\PYG{n}{bval}\PYG{+w}{ }\PYG{o}{=}\PYG{+w}{ }\PYG{n}{lmp}\PYG{p}{\PYGZpc{}}\PYG{n}{last\PYGZus{}thermo}\PYG{p}{(}\PYG{l+s+s1}{\PYGZsq{}data\PYGZsq{}}\PYG{p}{,}\PYG{n}{i}\PYG{p}{)}
\PYG{+w}{          }\PYG{k}{WRITE}\PYG{p}{(}\PYG{n}{OUTPUT\PYGZus{}UNIT}\PYG{p}{,}\PYG{o}{*}\PYG{p}{)}\PYG{+w}{ }\PYG{n}{key}\PYG{p}{,}\PYG{+w}{ }\PYG{l+s+s1}{\PYGZsq{}:\PYGZsq{}}\PYG{p}{,}\PYG{+w}{ }\PYG{n}{bval}
\PYG{+w}{      }\PYG{k}{ELSE }\PYG{k}{IF}\PYG{+w}{ }\PYG{p}{(}\PYG{n}{typ}\PYG{+w}{ }\PYG{o}{==}\PYG{+w}{ }\PYG{n}{lmp}\PYG{p}{\PYGZpc{}}\PYG{n}{dtype}\PYG{p}{\PYGZpc{}}\PYG{n}{r64}\PYG{p}{)}\PYG{+w}{ }\PYG{k}{THEN}
\PYG{k}{          }\PYG{n}{dval}\PYG{+w}{ }\PYG{o}{=}\PYG{+w}{ }\PYG{n}{lmp}\PYG{p}{\PYGZpc{}}\PYG{n}{last\PYGZus{}thermo}\PYG{p}{(}\PYG{l+s+s1}{\PYGZsq{}data\PYGZsq{}}\PYG{p}{,}\PYG{n}{i}\PYG{p}{)}
\PYG{+w}{          }\PYG{k}{WRITE}\PYG{p}{(}\PYG{n}{OUTPUT\PYGZus{}UNIT}\PYG{p}{,}\PYG{o}{*}\PYG{p}{)}\PYG{+w}{ }\PYG{n}{key}\PYG{p}{,}\PYG{+w}{ }\PYG{l+s+s1}{\PYGZsq{}:\PYGZsq{}}\PYG{p}{,}\PYG{+w}{ }\PYG{n}{dval}
\PYG{+w}{      }\PYG{k}{END }\PYG{k}{IF}
\PYG{k}{  }\PYG{k}{END }\PYG{k}{DO}

\PYG{k}{  }\PYG{n}{dt}\PYG{+w}{ }\PYG{o}{=}\PYG{+w}{ }\PYG{n}{lmp}\PYG{p}{\PYGZpc{}}\PYG{n}{extract\PYGZus{}global}\PYG{p}{(}\PYG{l+s+s1}{\PYGZsq{}dt\PYGZsq{}}\PYG{p}{)}
\PYG{+w}{  }\PYG{n}{ntimestep}\PYG{+w}{ }\PYG{o}{=}\PYG{+w}{ }\PYG{n}{lmp}\PYG{p}{\PYGZpc{}}\PYG{n}{extract\PYGZus{}global}\PYG{p}{(}\PYG{l+s+s1}{\PYGZsq{}ntimestep\PYGZsq{}}\PYG{p}{)}
\PYG{+w}{  }\PYG{k}{WRITE}\PYG{p}{(}\PYG{n}{OUTPUT\PYGZus{}UNIT}\PYG{p}{,}\PYG{l+s+s1}{\PYGZsq{}(A,I0,A,F4.1,A)\PYGZsq{}}\PYG{p}{)}\PYG{+w}{ }\PYG{l+s+s1}{\PYGZsq{}At step: \PYGZsq{}}\PYG{p}{,}\PYG{+w}{ }\PYG{n}{ntimestep}\PYG{p}{,}\PYG{+w}{ }\PYG{p}{\PYGZam{}}
\PYG{+w}{      }\PYG{l+s+s1}{\PYGZsq{}  Changing timestep from\PYGZsq{}}\PYG{p}{,}\PYG{+w}{ }\PYG{n}{dt}\PYG{p}{,}\PYG{+w}{ }\PYG{l+s+s1}{\PYGZsq{} to 0.5\PYGZsq{}}
\PYG{+w}{  }\PYG{n}{dt}\PYG{+w}{ }\PYG{o}{=}\PYG{+w}{ }\PYG{l+m+mf}{0.5\PYGZus{}c\PYGZus{}double}
\PYG{+w}{  }\PYG{k}{CALL }\PYG{n}{lmp}\PYG{p}{\PYGZpc{}}\PYG{n}{command}\PYG{p}{(}\PYG{l+s+s1}{\PYGZsq{}run 2 post no\PYGZsq{}}\PYG{p}{)}

\PYG{+w}{  }\PYG{k}{WRITE}\PYG{p}{(}\PYG{n}{OUTPUT\PYGZus{}UNIT}\PYG{p}{,}\PYG{l+s+s1}{\PYGZsq{}(A,I0)\PYGZsq{}}\PYG{p}{)}\PYG{+w}{ }\PYG{l+s+s1}{\PYGZsq{}At step: \PYGZsq{}}\PYG{p}{,}\PYG{+w}{ }\PYG{n}{ntimestep}
\PYG{+w}{  }\PYG{n}{pe}\PYG{+w}{ }\PYG{o}{=}\PYG{+w}{ }\PYG{n}{lmp}\PYG{p}{\PYGZpc{}}\PYG{n}{get\PYGZus{}thermo}\PYG{p}{(}\PYG{l+s+s1}{\PYGZsq{}pe\PYGZsq{}}\PYG{p}{)}
\PYG{+w}{  }\PYG{n}{ke}\PYG{+w}{ }\PYG{o}{=}\PYG{+w}{ }\PYG{n}{lmp}\PYG{p}{\PYGZpc{}}\PYG{n}{get\PYGZus{}thermo}\PYG{p}{(}\PYG{l+s+s1}{\PYGZsq{}ke\PYGZsq{}}\PYG{p}{)}
\PYG{+w}{  }\PYG{k}{WRITE}\PYG{p}{(}\PYG{n}{OUTPUT\PYGZus{}UNIT}\PYG{p}{,}\PYG{o}{*}\PYG{p}{)}\PYG{+w}{ }\PYG{l+s+s1}{\PYGZsq{}PE = \PYGZsq{}}\PYG{p}{,}\PYG{+w}{ }\PYG{n}{pe}
\PYG{+w}{  }\PYG{k}{WRITE}\PYG{p}{(}\PYG{n}{OUTPUT\PYGZus{}UNIT}\PYG{p}{,}\PYG{o}{*}\PYG{p}{)}\PYG{+w}{ }\PYG{l+s+s1}{\PYGZsq{}KE = \PYGZsq{}}\PYG{p}{,}\PYG{+w}{ }\PYG{n}{ke}

\PYG{+w}{  }\PYG{k}{CALL }\PYG{n}{lmp}\PYG{p}{\PYGZpc{}}\PYG{k}{close}\PYG{p}{(}\PYG{p}{.}\PYG{n}{TRUE}\PYG{p}{.}\PYG{p}{)}
\PYG{k}{END }\PYG{k}{PROGRAM }\PYG{n}{testprop}
\end{sphinxVerbatim}


\bigskip\hrule\bigskip



\subsection{The \sphinxstyleliteralintitle{\sphinxupquote{LIBLAMMPS}} module API}
\label{\detokenize{Fortran:the-liblammps-module-api}}\index{LIBLAMMPS (module)@\spxentry{LIBLAMMPS}\spxextra{module}|spxpagem}\phantomsection\label{\detokenize{Fortran:f/LIBLAMMPS}}
\sphinxAtStartPar
Below are the detailed descriptions of definitions and interfaces
of the contents of the \sphinxcode{\sphinxupquote{LIBLAMMPS}} Fortran interface to LAMMPS.
\index{lammps (fortran type)@\spxentry{lammps}\spxextra{fortran type}|spxpagem}

\begin{fulllineitems}
\phantomsection\label{\detokenize{Fortran:f/_/lammps}}
\pysigstartsignatures
\pysigline{\sphinxbfcode{\sphinxupquote{type  }}\sphinxbfcode{\sphinxupquote{lammps}}}
\pysigstopsignatures
\sphinxAtStartPar
Derived type that is the general class of the Fortran interface.
It holds a reference to the {\hyperref[\detokenize{Classes_lammps:_CPPv4N9LAMMPS_NS6LAMMPSE}]{\sphinxcrossref{\sphinxcode{\sphinxupquote{LAMMPS}}}}}
class instance to which any of the included calls are forwarded.
\begin{quote}\begin{description}
\sphinxlineitem{Type fields}\begin{itemize}
\item {} 
\sphinxAtStartPar
\sphinxbfcode{\sphinxupquote{\%~}}\sphinxbfcode{\sphinxupquote{handle}}\sphinxstyleemphasis{ {[}}\sphinxstyleemphasis{c\_ptr}\sphinxstyleemphasis{{]}} :: reference to the LAMMPS class

\item {} 
\sphinxAtStartPar
\sphinxbfcode{\sphinxupquote{\%~}}\sphinxbfcode{\sphinxupquote{style}}\sphinxstyleemphasis{ {[}}\sphinxstyleemphasis{type(lammps\_style)}\sphinxstyleemphasis{{]}} :: derived type to access lammps style constants

\item {} 
\sphinxAtStartPar
\sphinxbfcode{\sphinxupquote{\%~}}\sphinxbfcode{\sphinxupquote{type}}\sphinxstyleemphasis{ {[}}\sphinxstyleemphasis{type(lammps\_type)}\sphinxstyleemphasis{{]}} :: derived type to access lammps type constants

\item {} 
\sphinxAtStartPar
\sphinxbfcode{\sphinxupquote{\%~}}\sphinxbfcode{\sphinxupquote{dtype}}\sphinxstyleemphasis{ {[}}\sphinxstyleemphasis{type(lammps\_dtype)}\sphinxstyleemphasis{{]}} :: derived type to access lammps data type constants

\item {} 
\sphinxAtStartPar
\sphinxbfcode{\sphinxupquote{\%~}}\sphinxbfcode{\sphinxupquote{close}}\sphinxstyleemphasis{ {[}}\sphinxstyleemphasis{subroutine}\sphinxstyleemphasis{{]}} :: \sphinxcode{\sphinxupquote{close()}}

\item {} 
\sphinxAtStartPar
\sphinxbfcode{\sphinxupquote{\%~}}\sphinxbfcode{\sphinxupquote{error}}\sphinxstyleemphasis{ {[}}\sphinxstyleemphasis{subroutine}\sphinxstyleemphasis{{]}} :: \sphinxcode{\sphinxupquote{error()}}

\item {} 
\sphinxAtStartPar
\sphinxbfcode{\sphinxupquote{\%~}}\sphinxbfcode{\sphinxupquote{file}}\sphinxstyleemphasis{ {[}}\sphinxstyleemphasis{subroutine}\sphinxstyleemphasis{{]}} :: \sphinxcode{\sphinxupquote{file()}}

\item {} 
\sphinxAtStartPar
\sphinxbfcode{\sphinxupquote{\%~}}\sphinxbfcode{\sphinxupquote{command}}\sphinxstyleemphasis{ {[}}\sphinxstyleemphasis{subroutine}\sphinxstyleemphasis{{]}} :: \sphinxcode{\sphinxupquote{command()}}

\item {} 
\sphinxAtStartPar
\sphinxbfcode{\sphinxupquote{\%~}}\sphinxbfcode{\sphinxupquote{commands\_list}}\sphinxstyleemphasis{ {[}}\sphinxstyleemphasis{subroutine}\sphinxstyleemphasis{{]}} :: \sphinxcode{\sphinxupquote{commands\_list()}}

\item {} 
\sphinxAtStartPar
\sphinxbfcode{\sphinxupquote{\%~}}\sphinxbfcode{\sphinxupquote{commands\_string}}\sphinxstyleemphasis{ {[}}\sphinxstyleemphasis{subroutine}\sphinxstyleemphasis{{]}} :: \sphinxcode{\sphinxupquote{commands\_string()}}

\item {} 
\sphinxAtStartPar
\sphinxbfcode{\sphinxupquote{\%~}}\sphinxbfcode{\sphinxupquote{get\_natoms}}\sphinxstyleemphasis{ {[}}\sphinxstyleemphasis{function}\sphinxstyleemphasis{{]}} :: \sphinxcode{\sphinxupquote{get\_natoms()}}

\item {} 
\sphinxAtStartPar
\sphinxbfcode{\sphinxupquote{\%~}}\sphinxbfcode{\sphinxupquote{get\_thermo}}\sphinxstyleemphasis{ {[}}\sphinxstyleemphasis{function}\sphinxstyleemphasis{{]}} :: \sphinxcode{\sphinxupquote{get\_thermo()}}

\item {} 
\sphinxAtStartPar
\sphinxbfcode{\sphinxupquote{\%~}}\sphinxbfcode{\sphinxupquote{last\_thermo}}\sphinxstyleemphasis{ {[}}\sphinxstyleemphasis{function}\sphinxstyleemphasis{{]}} :: \sphinxcode{\sphinxupquote{last\_thermo()}}

\item {} 
\sphinxAtStartPar
\sphinxbfcode{\sphinxupquote{\%~}}\sphinxbfcode{\sphinxupquote{extract\_box}}\sphinxstyleemphasis{ {[}}\sphinxstyleemphasis{subroutine}\sphinxstyleemphasis{{]}} :: \sphinxcode{\sphinxupquote{extract\_box()}}

\item {} 
\sphinxAtStartPar
\sphinxbfcode{\sphinxupquote{\%~}}\sphinxbfcode{\sphinxupquote{reset\_box}}\sphinxstyleemphasis{ {[}}\sphinxstyleemphasis{subroutine}\sphinxstyleemphasis{{]}} :: \sphinxcode{\sphinxupquote{reset\_box()}}

\item {} 
\sphinxAtStartPar
\sphinxbfcode{\sphinxupquote{\%~}}\sphinxbfcode{\sphinxupquote{memory\_usage}}\sphinxstyleemphasis{ {[}}\sphinxstyleemphasis{subroutine}\sphinxstyleemphasis{{]}} :: \sphinxcode{\sphinxupquote{memory\_usage()}}

\item {} 
\sphinxAtStartPar
\sphinxbfcode{\sphinxupquote{\%~}}\sphinxbfcode{\sphinxupquote{get\_mpi\_comm}}\sphinxstyleemphasis{ {[}}\sphinxstyleemphasis{function}\sphinxstyleemphasis{{]}} :: \sphinxcode{\sphinxupquote{get\_mpi\_comm()}}

\item {} 
\sphinxAtStartPar
\sphinxbfcode{\sphinxupquote{\%~}}\sphinxbfcode{\sphinxupquote{extract\_setting}}\sphinxstyleemphasis{ {[}}\sphinxstyleemphasis{function}\sphinxstyleemphasis{{]}} :: \sphinxcode{\sphinxupquote{extract\_setting()}}

\item {} 
\sphinxAtStartPar
\sphinxbfcode{\sphinxupquote{\%~}}\sphinxbfcode{\sphinxupquote{extract\_global}}\sphinxstyleemphasis{ {[}}\sphinxstyleemphasis{function}\sphinxstyleemphasis{{]}} :: \sphinxcode{\sphinxupquote{extract\_global()}}

\item {} 
\sphinxAtStartPar
\sphinxbfcode{\sphinxupquote{\%~}}\sphinxbfcode{\sphinxupquote{map\_atom}}\sphinxstyleemphasis{ {[}}\sphinxstyleemphasis{function}\sphinxstyleemphasis{{]}} :: \sphinxcode{\sphinxupquote{map\_atom()}}

\item {} 
\sphinxAtStartPar
\sphinxbfcode{\sphinxupquote{\%~}}\sphinxbfcode{\sphinxupquote{extract\_atom}}\sphinxstyleemphasis{ {[}}\sphinxstyleemphasis{function}\sphinxstyleemphasis{{]}} :: \sphinxcode{\sphinxupquote{extract\_atom()}}

\item {} 
\sphinxAtStartPar
\sphinxbfcode{\sphinxupquote{\%~}}\sphinxbfcode{\sphinxupquote{extract\_compute}}\sphinxstyleemphasis{ {[}}\sphinxstyleemphasis{function}\sphinxstyleemphasis{{]}} :: \sphinxcode{\sphinxupquote{extract\_compute()}}

\item {} 
\sphinxAtStartPar
\sphinxbfcode{\sphinxupquote{\%~}}\sphinxbfcode{\sphinxupquote{extract\_fix}}\sphinxstyleemphasis{ {[}}\sphinxstyleemphasis{function}\sphinxstyleemphasis{{]}} :: \sphinxcode{\sphinxupquote{extract\_fix()}}

\item {} 
\sphinxAtStartPar
\sphinxbfcode{\sphinxupquote{\%~}}\sphinxbfcode{\sphinxupquote{extract\_variable}}\sphinxstyleemphasis{ {[}}\sphinxstyleemphasis{function}\sphinxstyleemphasis{{]}} :: \sphinxcode{\sphinxupquote{extract\_variable()}}

\item {} 
\sphinxAtStartPar
\sphinxbfcode{\sphinxupquote{\%~}}\sphinxbfcode{\sphinxupquote{set\_variable}}\sphinxstyleemphasis{ {[}}\sphinxstyleemphasis{subroutine}\sphinxstyleemphasis{{]}} :: \sphinxcode{\sphinxupquote{set\_variable()}}

\item {} 
\sphinxAtStartPar
\sphinxbfcode{\sphinxupquote{\%~}}\sphinxbfcode{\sphinxupquote{set\_string\_variable}}\sphinxstyleemphasis{ {[}}\sphinxstyleemphasis{subroutine}\sphinxstyleemphasis{{]}} :: \sphinxcode{\sphinxupquote{set\_set\_string\_variable()}}

\item {} 
\sphinxAtStartPar
\sphinxbfcode{\sphinxupquote{\%~}}\sphinxbfcode{\sphinxupquote{set\_internal\_variable}}\sphinxstyleemphasis{ {[}}\sphinxstyleemphasis{subroutine}\sphinxstyleemphasis{{]}} :: \sphinxcode{\sphinxupquote{set\_internal\_variable()}}

\item {} 
\sphinxAtStartPar
\sphinxbfcode{\sphinxupquote{\%~}}\sphinxbfcode{\sphinxupquote{gather\_atoms}}\sphinxstyleemphasis{ {[}}\sphinxstyleemphasis{subroutine}\sphinxstyleemphasis{{]}} :: \sphinxcode{\sphinxupquote{gather\_atoms()}}

\item {} 
\sphinxAtStartPar
\sphinxbfcode{\sphinxupquote{\%~}}\sphinxbfcode{\sphinxupquote{gather\_atoms\_concat}}\sphinxstyleemphasis{ {[}}\sphinxstyleemphasis{subroutine}\sphinxstyleemphasis{{]}} :: \sphinxcode{\sphinxupquote{gather\_atoms\_concat()}}

\item {} 
\sphinxAtStartPar
\sphinxbfcode{\sphinxupquote{\%~}}\sphinxbfcode{\sphinxupquote{gather\_atoms\_subset}}\sphinxstyleemphasis{ {[}}\sphinxstyleemphasis{subroutine}\sphinxstyleemphasis{{]}} :: \sphinxcode{\sphinxupquote{gather\_atoms\_subset()}}

\item {} 
\sphinxAtStartPar
\sphinxbfcode{\sphinxupquote{\%~}}\sphinxbfcode{\sphinxupquote{scatter\_atoms}}\sphinxstyleemphasis{ {[}}\sphinxstyleemphasis{subroutine}\sphinxstyleemphasis{{]}} :: \sphinxcode{\sphinxupquote{scatter\_atoms()}}

\item {} 
\sphinxAtStartPar
\sphinxbfcode{\sphinxupquote{\%~}}\sphinxbfcode{\sphinxupquote{scatter\_atoms\_subset}}\sphinxstyleemphasis{ {[}}\sphinxstyleemphasis{subroutine}\sphinxstyleemphasis{{]}} :: \sphinxcode{\sphinxupquote{scatter\_atoms\_subset()}}

\item {} 
\sphinxAtStartPar
\sphinxbfcode{\sphinxupquote{\%~}}\sphinxbfcode{\sphinxupquote{gather\_bonds}}\sphinxstyleemphasis{ {[}}\sphinxstyleemphasis{subroutine}\sphinxstyleemphasis{{]}} :: \sphinxcode{\sphinxupquote{gather\_bonds()}}

\item {} 
\sphinxAtStartPar
\sphinxbfcode{\sphinxupquote{\%~}}\sphinxbfcode{\sphinxupquote{gather\_angles}}\sphinxstyleemphasis{ {[}}\sphinxstyleemphasis{subroutine}\sphinxstyleemphasis{{]}} :: \sphinxcode{\sphinxupquote{gather\_angles()}}

\item {} 
\sphinxAtStartPar
\sphinxbfcode{\sphinxupquote{\%~}}\sphinxbfcode{\sphinxupquote{gather\_dihedrals}}\sphinxstyleemphasis{ {[}}\sphinxstyleemphasis{subroutine}\sphinxstyleemphasis{{]}} :: \sphinxcode{\sphinxupquote{gather\_dihedrals()}}

\item {} 
\sphinxAtStartPar
\sphinxbfcode{\sphinxupquote{\%~}}\sphinxbfcode{\sphinxupquote{gather\_impropers}}\sphinxstyleemphasis{ {[}}\sphinxstyleemphasis{subroutine}\sphinxstyleemphasis{{]}} :: \sphinxcode{\sphinxupquote{gather\_impropers()}}

\item {} 
\sphinxAtStartPar
\sphinxbfcode{\sphinxupquote{\%~}}\sphinxbfcode{\sphinxupquote{gather}}\sphinxstyleemphasis{ {[}}\sphinxstyleemphasis{subroutine}\sphinxstyleemphasis{{]}} :: \sphinxcode{\sphinxupquote{gather()}}

\item {} 
\sphinxAtStartPar
\sphinxbfcode{\sphinxupquote{\%~}}\sphinxbfcode{\sphinxupquote{gather\_concat}}\sphinxstyleemphasis{ {[}}\sphinxstyleemphasis{subroutine}\sphinxstyleemphasis{{]}} :: \sphinxcode{\sphinxupquote{gather\_concat()}}

\item {} 
\sphinxAtStartPar
\sphinxbfcode{\sphinxupquote{\%~}}\sphinxbfcode{\sphinxupquote{gather\_subset}}\sphinxstyleemphasis{ {[}}\sphinxstyleemphasis{subroutine}\sphinxstyleemphasis{{]}} :: \sphinxcode{\sphinxupquote{gather\_subset()}}

\item {} 
\sphinxAtStartPar
\sphinxbfcode{\sphinxupquote{\%~}}\sphinxbfcode{\sphinxupquote{scatter}}\sphinxstyleemphasis{ {[}}\sphinxstyleemphasis{subroutine}\sphinxstyleemphasis{{]}} :: \sphinxcode{\sphinxupquote{scatter()}}

\item {} 
\sphinxAtStartPar
\sphinxbfcode{\sphinxupquote{\%~}}\sphinxbfcode{\sphinxupquote{scatter\_subset}}\sphinxstyleemphasis{ {[}}\sphinxstyleemphasis{subroutine}\sphinxstyleemphasis{{]}} :: \sphinxcode{\sphinxupquote{scatter\_subset()}}

\item {} 
\sphinxAtStartPar
\sphinxbfcode{\sphinxupquote{\%~}}\sphinxbfcode{\sphinxupquote{create\_atoms}}\sphinxstyleemphasis{ {[}}\sphinxstyleemphasis{subroutine}\sphinxstyleemphasis{{]}} :: \sphinxcode{\sphinxupquote{create\_atoms()}}

\item {} 
\sphinxAtStartPar
\sphinxbfcode{\sphinxupquote{\%~}}\sphinxbfcode{\sphinxupquote{find\_pair\_neighlist}}\sphinxstyleemphasis{ {[}}\sphinxstyleemphasis{function}\sphinxstyleemphasis{{]}} :: \sphinxcode{\sphinxupquote{find\_pair\_neighlist()}}

\item {} 
\sphinxAtStartPar
\sphinxbfcode{\sphinxupquote{\%~}}\sphinxbfcode{\sphinxupquote{find\_fix\_neighlist}}\sphinxstyleemphasis{ {[}}\sphinxstyleemphasis{function}\sphinxstyleemphasis{{]}} :: \sphinxcode{\sphinxupquote{find\_fix\_neighlist()}}

\item {} 
\sphinxAtStartPar
\sphinxbfcode{\sphinxupquote{\%~}}\sphinxbfcode{\sphinxupquote{find\_compute\_neighlist}}\sphinxstyleemphasis{ {[}}\sphinxstyleemphasis{function}\sphinxstyleemphasis{{]}} :: \sphinxcode{\sphinxupquote{find\_compute\_neighlist()}}

\item {} 
\sphinxAtStartPar
\sphinxbfcode{\sphinxupquote{\%~}}\sphinxbfcode{\sphinxupquote{neighlist\_num\_elements}}\sphinxstyleemphasis{ {[}}\sphinxstyleemphasis{function}\sphinxstyleemphasis{{]}} :: \sphinxcode{\sphinxupquote{neighlist\_num\_elements()}}

\item {} 
\sphinxAtStartPar
\sphinxbfcode{\sphinxupquote{\%~}}\sphinxbfcode{\sphinxupquote{neighlist\_element\_neighbors}}\sphinxstyleemphasis{ {[}}\sphinxstyleemphasis{subroutine}\sphinxstyleemphasis{{]}} :: \sphinxcode{\sphinxupquote{neighlist\_element\_neighbors()}}

\item {} 
\sphinxAtStartPar
\sphinxbfcode{\sphinxupquote{\%~}}\sphinxbfcode{\sphinxupquote{version}}\sphinxstyleemphasis{ {[}}\sphinxstyleemphasis{function}\sphinxstyleemphasis{{]}} :: \sphinxcode{\sphinxupquote{version()}}

\item {} 
\sphinxAtStartPar
\sphinxbfcode{\sphinxupquote{\%~}}\sphinxbfcode{\sphinxupquote{get\_os\_info}}\sphinxstyleemphasis{ {[}}\sphinxstyleemphasis{subroutine}\sphinxstyleemphasis{{]}} :: \sphinxcode{\sphinxupquote{get\_os\_info()}}

\item {} 
\sphinxAtStartPar
\sphinxbfcode{\sphinxupquote{\%~}}\sphinxbfcode{\sphinxupquote{config\_has\_mpi\_support}}\sphinxstyleemphasis{ {[}}\sphinxstyleemphasis{function}\sphinxstyleemphasis{{]}} :: \sphinxcode{\sphinxupquote{config\_has\_mpi\_support()}}

\item {} 
\sphinxAtStartPar
\sphinxbfcode{\sphinxupquote{\%~}}\sphinxbfcode{\sphinxupquote{config\_has\_gzip\_support}}\sphinxstyleemphasis{ {[}}\sphinxstyleemphasis{function}\sphinxstyleemphasis{{]}} :: \sphinxcode{\sphinxupquote{config\_has\_gzip\_support()}}

\item {} 
\sphinxAtStartPar
\sphinxbfcode{\sphinxupquote{\%~}}\sphinxbfcode{\sphinxupquote{config\_has\_png\_support}}\sphinxstyleemphasis{ {[}}\sphinxstyleemphasis{function}\sphinxstyleemphasis{{]}} :: \sphinxcode{\sphinxupquote{config\_has\_png\_support()}}

\item {} 
\sphinxAtStartPar
\sphinxbfcode{\sphinxupquote{\%~}}\sphinxbfcode{\sphinxupquote{config\_has\_jpeg\_support}}\sphinxstyleemphasis{ {[}}\sphinxstyleemphasis{function}\sphinxstyleemphasis{{]}} :: \sphinxcode{\sphinxupquote{config\_has\_jpeg\_support()}}

\item {} 
\sphinxAtStartPar
\sphinxbfcode{\sphinxupquote{\%~}}\sphinxbfcode{\sphinxupquote{config\_has\_ffmpeg\_support}}\sphinxstyleemphasis{ {[}}\sphinxstyleemphasis{function}\sphinxstyleemphasis{{]}} :: \sphinxcode{\sphinxupquote{config\_has\_ffmpeg\_support()}}

\item {} 
\sphinxAtStartPar
\sphinxbfcode{\sphinxupquote{\%~}}\sphinxbfcode{\sphinxupquote{config\_has\_exceptions}}\sphinxstyleemphasis{ {[}}\sphinxstyleemphasis{function}\sphinxstyleemphasis{{]}} :: \sphinxcode{\sphinxupquote{config\_has\_exceptions()}}

\item {} 
\sphinxAtStartPar
\sphinxbfcode{\sphinxupquote{\%~}}\sphinxbfcode{\sphinxupquote{config\_has\_package}}\sphinxstyleemphasis{ {[}}\sphinxstyleemphasis{function}\sphinxstyleemphasis{{]}} :: \sphinxcode{\sphinxupquote{config\_has\_package()}}

\item {} 
\sphinxAtStartPar
\sphinxbfcode{\sphinxupquote{\%~}}\sphinxbfcode{\sphinxupquote{config\_package\_count}}\sphinxstyleemphasis{ {[}}\sphinxstyleemphasis{function}\sphinxstyleemphasis{{]}} :: \sphinxcode{\sphinxupquote{config\_package\_count()}}

\item {} 
\sphinxAtStartPar
\sphinxbfcode{\sphinxupquote{\%~}}\sphinxbfcode{\sphinxupquote{config\_package\_name}}\sphinxstyleemphasis{ {[}}\sphinxstyleemphasis{function}\sphinxstyleemphasis{{]}} :: \sphinxcode{\sphinxupquote{config\_package\_name()}}

\item {} 
\sphinxAtStartPar
\sphinxbfcode{\sphinxupquote{\%~}}\sphinxbfcode{\sphinxupquote{installed\_packages}}\sphinxstyleemphasis{ {[}}\sphinxstyleemphasis{subroutine}\sphinxstyleemphasis{{]}} :: \sphinxcode{\sphinxupquote{installed\_packages()}}

\item {} 
\sphinxAtStartPar
\sphinxbfcode{\sphinxupquote{\%~}}\sphinxbfcode{\sphinxupquote{config\_accelerator}}\sphinxstyleemphasis{ {[}}\sphinxstyleemphasis{function}\sphinxstyleemphasis{{]}} :: \sphinxcode{\sphinxupquote{config\_accelerator()}}

\item {} 
\sphinxAtStartPar
\sphinxbfcode{\sphinxupquote{\%~}}\sphinxbfcode{\sphinxupquote{has\_gpu\_device}}\sphinxstyleemphasis{ {[}}\sphinxstyleemphasis{function}\sphinxstyleemphasis{{]}} :: \sphinxcode{\sphinxupquote{has\_gpu\_device()}}

\item {} 
\sphinxAtStartPar
\sphinxbfcode{\sphinxupquote{\%~}}\sphinxbfcode{\sphinxupquote{get\_gpu\_device\_info}}\sphinxstyleemphasis{ {[}}\sphinxstyleemphasis{subroutine}\sphinxstyleemphasis{{]}} :: \sphinxcode{\sphinxupquote{get\_gpu\_device\_info()}}

\item {} 
\sphinxAtStartPar
\sphinxbfcode{\sphinxupquote{\%~}}\sphinxbfcode{\sphinxupquote{has\_style}}\sphinxstyleemphasis{ {[}}\sphinxstyleemphasis{function}\sphinxstyleemphasis{{]}} :: \sphinxcode{\sphinxupquote{has\_style()}}

\item {} 
\sphinxAtStartPar
\sphinxbfcode{\sphinxupquote{\%~}}\sphinxbfcode{\sphinxupquote{style\_count}}\sphinxstyleemphasis{ {[}}\sphinxstyleemphasis{function}\sphinxstyleemphasis{{]}} :: \sphinxcode{\sphinxupquote{style\_count()}}

\item {} 
\sphinxAtStartPar
\sphinxbfcode{\sphinxupquote{\%~}}\sphinxbfcode{\sphinxupquote{style\_name}}\sphinxstyleemphasis{ {[}}\sphinxstyleemphasis{function}\sphinxstyleemphasis{{]}} :: \sphinxcode{\sphinxupquote{style\_name()}}

\item {} 
\sphinxAtStartPar
\sphinxbfcode{\sphinxupquote{\%~}}\sphinxbfcode{\sphinxupquote{has\_id}}\sphinxstyleemphasis{ {[}}\sphinxstyleemphasis{function}\sphinxstyleemphasis{{]}} :: \sphinxcode{\sphinxupquote{has\_id()}}

\item {} 
\sphinxAtStartPar
\sphinxbfcode{\sphinxupquote{\%~}}\sphinxbfcode{\sphinxupquote{id\_count}}\sphinxstyleemphasis{ {[}}\sphinxstyleemphasis{function}\sphinxstyleemphasis{{]}} :: \sphinxcode{\sphinxupquote{id\_count()}}

\item {} 
\sphinxAtStartPar
\sphinxbfcode{\sphinxupquote{\%~}}\sphinxbfcode{\sphinxupquote{id\_name}}\sphinxstyleemphasis{ {[}}\sphinxstyleemphasis{subroutine}\sphinxstyleemphasis{{]}} :: \sphinxcode{\sphinxupquote{id\_name()}}

\item {} 
\sphinxAtStartPar
\sphinxbfcode{\sphinxupquote{\%~}}\sphinxbfcode{\sphinxupquote{plugin\_count}}\sphinxstyleemphasis{ {[}}\sphinxstyleemphasis{subroutine}\sphinxstyleemphasis{{]}} :: \sphinxcode{\sphinxupquote{plugin\_count()}}

\item {} 
\sphinxAtStartPar
\sphinxbfcode{\sphinxupquote{\%~}}\sphinxbfcode{\sphinxupquote{plugin\_name}} :: \sphinxcode{\sphinxupquote{plugin\_name()}}

\item {} 
\sphinxAtStartPar
\sphinxbfcode{\sphinxupquote{\%~}}\sphinxbfcode{\sphinxupquote{encode\_image\_flags}}\sphinxstyleemphasis{ {[}}\sphinxstyleemphasis{function}\sphinxstyleemphasis{{]}} :: \sphinxcode{\sphinxupquote{encode\_image\_flags()}}

\item {} 
\sphinxAtStartPar
\sphinxbfcode{\sphinxupquote{\%~}}\sphinxbfcode{\sphinxupquote{decode\_image\_flags}}\sphinxstyleemphasis{ {[}}\sphinxstyleemphasis{subroutine}\sphinxstyleemphasis{{]}} :: \sphinxcode{\sphinxupquote{decode\_image\_flags()}}

\item {} 
\sphinxAtStartPar
\sphinxbfcode{\sphinxupquote{\%~}}\sphinxbfcode{\sphinxupquote{set\_fix\_external\_callback}}\sphinxstyleemphasis{ {[}}\sphinxstyleemphasis{subroutine}\sphinxstyleemphasis{{]}} :: \sphinxcode{\sphinxupquote{set\_fix\_external\_callback()}}

\item {} 
\sphinxAtStartPar
\sphinxbfcode{\sphinxupquote{\%~}}\sphinxbfcode{\sphinxupquote{fix\_external\_get\_force}}\sphinxstyleemphasis{ {[}}\sphinxstyleemphasis{function}\sphinxstyleemphasis{{]}} :: \sphinxcode{\sphinxupquote{fix\_external\_get\_force()}}

\item {} 
\sphinxAtStartPar
\sphinxbfcode{\sphinxupquote{\%~}}\sphinxbfcode{\sphinxupquote{fix\_external\_set\_energy\_global}}\sphinxstyleemphasis{ {[}}\sphinxstyleemphasis{subroutine}\sphinxstyleemphasis{{]}} :: \sphinxcode{\sphinxupquote{fix\_external\_set\_energy\_global()}}

\item {} 
\sphinxAtStartPar
\sphinxbfcode{\sphinxupquote{\%~}}\sphinxbfcode{\sphinxupquote{fix\_external\_set\_virial\_global}}\sphinxstyleemphasis{ {[}}\sphinxstyleemphasis{subroutine}\sphinxstyleemphasis{{]}} :: \sphinxcode{\sphinxupquote{fix\_external\_set\_virial\_global()}}

\item {} 
\sphinxAtStartPar
\sphinxbfcode{\sphinxupquote{\%~}}\sphinxbfcode{\sphinxupquote{fix\_external\_set\_energy\_peratom}}\sphinxstyleemphasis{ {[}}\sphinxstyleemphasis{subroutine}\sphinxstyleemphasis{{]}} :: \sphinxcode{\sphinxupquote{fix\_external\_set\_energy\_peratom()}}

\item {} 
\sphinxAtStartPar
\sphinxbfcode{\sphinxupquote{\%~}}\sphinxbfcode{\sphinxupquote{fix\_external\_set\_virial\_peratom}}\sphinxstyleemphasis{ {[}}\sphinxstyleemphasis{subroutine}\sphinxstyleemphasis{{]}} :: \sphinxcode{\sphinxupquote{fix\_external\_set\_virial\_peratom()}}

\item {} 
\sphinxAtStartPar
\sphinxbfcode{\sphinxupquote{\%~}}\sphinxbfcode{\sphinxupquote{fix\_external\_set\_vector\_length}}\sphinxstyleemphasis{ {[}}\sphinxstyleemphasis{subroutine}\sphinxstyleemphasis{{]}} :: \sphinxcode{\sphinxupquote{fix\_external\_set\_vector\_length()}}

\item {} 
\sphinxAtStartPar
\sphinxbfcode{\sphinxupquote{\%~}}\sphinxbfcode{\sphinxupquote{fix\_external\_set\_vector}}\sphinxstyleemphasis{ {[}}\sphinxstyleemphasis{subroutine}\sphinxstyleemphasis{{]}} :: \sphinxcode{\sphinxupquote{fix\_external\_set\_vector()}}

\item {} 
\sphinxAtStartPar
\sphinxbfcode{\sphinxupquote{\%~}}\sphinxbfcode{\sphinxupquote{flush\_buffers}}\sphinxstyleemphasis{ {[}}\sphinxstyleemphasis{subroutine}\sphinxstyleemphasis{{]}} :: \sphinxcode{\sphinxupquote{flush\_buffers()}}

\item {} 
\sphinxAtStartPar
\sphinxbfcode{\sphinxupquote{\%~}}\sphinxbfcode{\sphinxupquote{is\_running}}\sphinxstyleemphasis{ {[}}\sphinxstyleemphasis{function}\sphinxstyleemphasis{{]}} :: \sphinxcode{\sphinxupquote{is\_running()}}

\item {} 
\sphinxAtStartPar
\sphinxbfcode{\sphinxupquote{\%~}}\sphinxbfcode{\sphinxupquote{force\_timeout}}\sphinxstyleemphasis{ {[}}\sphinxstyleemphasis{subroutine}\sphinxstyleemphasis{{]}} :: \sphinxcode{\sphinxupquote{force\_timeout()}}

\item {} 
\sphinxAtStartPar
\sphinxbfcode{\sphinxupquote{\%~}}\sphinxbfcode{\sphinxupquote{has\_error}}\sphinxstyleemphasis{ {[}}\sphinxstyleemphasis{function}\sphinxstyleemphasis{{]}} :: \sphinxcode{\sphinxupquote{has\_error()}}

\item {} 
\sphinxAtStartPar
\sphinxbfcode{\sphinxupquote{\%~}}\sphinxbfcode{\sphinxupquote{get\_last\_error\_message}}\sphinxstyleemphasis{ {[}}\sphinxstyleemphasis{subroutine}\sphinxstyleemphasis{{]}} :: \sphinxcode{\sphinxupquote{get\_last\_error\_message()}}

\end{itemize}

\end{description}\end{quote}

\end{fulllineitems}



\bigskip\hrule\bigskip

\index{lammps() (fortran function)@\spxentry{lammps()}\spxextra{fortran function}|spxpagem}

\begin{fulllineitems}

\pysigstartsignatures
\pysiglinewithargsret{\sphinxbfcode{\sphinxupquote{function  }}\sphinxbfcode{\sphinxupquote{lammps}}}{\sphinxparam{{[}args{]}{[},comm{]}}}{}
\pysigstopsignatures
\sphinxAtStartPar
This is the constructor for the Fortran class and will forward
the arguments to a call to either {\hyperref[\detokenize{Library_create:_CPPv419lammps_open_fortraniPPci}]{\sphinxcrossref{\sphinxcode{\sphinxupquote{lammps\_open\_fortran()}}}}}
or {\hyperref[\detokenize{Library_create:_CPPv418lammps_open_no_mpiiPPcPPv}]{\sphinxcrossref{\sphinxcode{\sphinxupquote{lammps\_open\_no\_mpi()}}}}}. If the LAMMPS library has been
compiled with MPI support, it will also initialize MPI, if it has
not already been initialized before.

\sphinxAtStartPar
The \sphinxstyleemphasis{args} argument with the list of command line parameters is
optional and so it the \sphinxstyleemphasis{comm} argument with the MPI communicator.
If \sphinxstyleemphasis{comm} is not provided, \sphinxcode{\sphinxupquote{MPI\_COMM\_WORLD}} is assumed. For
more details please see the documentation of {\hyperref[\detokenize{Library_create:_CPPv411lammps_openiPPc8MPI_CommPPv}]{\sphinxcrossref{\sphinxcode{\sphinxupquote{lammps\_open()}}}}}.
\begin{quote}\begin{description}
\sphinxlineitem{Options}\begin{itemize}
\item {} 
\sphinxAtStartPar
\sphinxstylestrong{args}\sphinxstyleemphasis{ {[}}\sphinxstyleemphasis{character(len=*),dimension(:),optional}\sphinxstyleemphasis{{]}} :: arguments as list of strings

\item {} 
\sphinxAtStartPar
\sphinxstylestrong{comm}\sphinxstyleemphasis{ {[}}\sphinxstyleemphasis{integer,optional}\sphinxstyleemphasis{{]}} :: MPI communicator

\end{itemize}

\sphinxlineitem{Call to}
\sphinxAtStartPar
{\hyperref[\detokenize{Library_create:_CPPv419lammps_open_fortraniPPci}]{\sphinxcrossref{\sphinxcode{\sphinxupquote{lammps\_open\_fortran()}}}}} {\hyperref[\detokenize{Library_create:_CPPv418lammps_open_no_mpiiPPcPPv}]{\sphinxcrossref{\sphinxcode{\sphinxupquote{lammps\_open\_no\_mpi()}}}}}

\sphinxlineitem{Return}
\sphinxAtStartPar
\sphinxstylestrong{lammps} :: an instance of the \sphinxcode{\sphinxupquote{lammps}} derived type

\end{description}\end{quote}

\begin{sphinxadmonition}{note}{Note:}
\sphinxAtStartPar
The \sphinxcode{\sphinxupquote{MPI\_F08}} module, which defines Fortran 2008 bindings for MPI,
is not directly supported by this interface due to the complexities of
supporting both the \sphinxcode{\sphinxupquote{MPI\_F08}} and \sphinxcode{\sphinxupquote{MPI}} modules at the same
time. However, you should be able to use the \sphinxcode{\sphinxupquote{MPI\_VAL}} member of the
\sphinxcode{\sphinxupquote{MPI\_comm}} derived type to access the integer value of the
communicator, such as in

\begin{sphinxVerbatim}[commandchars=\\\{\}]
\PYG{k}{PROGRAM }\PYG{n}{testmpi}
\PYG{+w}{  }\PYG{k}{USE }\PYG{n}{LIBLAMMPS}
\PYG{+w}{  }\PYG{k}{USE }\PYG{n}{MPI\PYGZus{}F08}
\PYG{+w}{  }\PYG{k}{TYPE}\PYG{p}{(}\PYG{n}{lammps}\PYG{p}{)}\PYG{+w}{ }\PYG{k+kd}{::}\PYG{+w}{ }\PYG{n}{lmp}
\PYG{+w}{  }\PYG{n}{lmp}\PYG{+w}{ }\PYG{o}{=}\PYG{+w}{ }\PYG{n}{lammps}\PYG{p}{(}\PYG{n}{comm}\PYG{o}{=}\PYG{n}{MPI\PYGZus{}COMM\PYGZus{}SELF}\PYG{p}{\PYGZpc{}}\PYG{n}{MPI\PYGZus{}VAL}\PYG{p}{)}
\PYG{k}{END }\PYG{k}{PROGRAM }\PYG{n}{testmpi}
\end{sphinxVerbatim}
\end{sphinxadmonition}

\end{fulllineitems}

\index{lammps\_style (fortran type)@\spxentry{lammps\_style}\spxextra{fortran type}|spxpagem}

\begin{fulllineitems}
\phantomsection\label{\detokenize{Fortran:f/_/lammps_style}}
\pysigstartsignatures
\pysigline{\sphinxbfcode{\sphinxupquote{type  }}\sphinxbfcode{\sphinxupquote{lammps\_style}}}
\pysigstopsignatures
\sphinxAtStartPar
This derived type is there to provide a convenient interface for the style
constants used with \sphinxcode{\sphinxupquote{extract\_compute()}}, \sphinxcode{\sphinxupquote{extract\_fix()}}, and
\sphinxcode{\sphinxupquote{extract\_variable()}}. Assuming your LAMMPS instance is called \sphinxcode{\sphinxupquote{lmp}},
these constants will be \sphinxcode{\sphinxupquote{lmp\%style\%global}}, \sphinxcode{\sphinxupquote{lmp\%style\%atom}},
and \sphinxcode{\sphinxupquote{lmp\%style\%local}}. These values are identical to the values described
in {\hyperref[\detokenize{Library_objects:_CPPv416_LMP_STYLE_CONST}]{\sphinxcrossref{\sphinxcode{\sphinxupquote{\_LMP\_STYLE\_CONST}}}}} for the C library interface.
\begin{quote}\begin{description}
\sphinxlineitem{Type fields}\begin{itemize}
\item {} 
\sphinxAtStartPar
\sphinxbfcode{\sphinxupquote{\%~}}\sphinxbfcode{\sphinxupquote{global}}\sphinxstyleemphasis{ {[}}\sphinxstyleemphasis{integer(c\_int)}\sphinxstyleemphasis{{]}} :: used to request global data

\item {} 
\sphinxAtStartPar
\sphinxbfcode{\sphinxupquote{\%~}}\sphinxbfcode{\sphinxupquote{atom}}\sphinxstyleemphasis{ {[}}\sphinxstyleemphasis{integer(c\_int)}\sphinxstyleemphasis{{]}} :: used to request per\sphinxhyphen{}atom data

\item {} 
\sphinxAtStartPar
\sphinxbfcode{\sphinxupquote{\%~}}\sphinxbfcode{\sphinxupquote{local}}\sphinxstyleemphasis{ {[}}\sphinxstyleemphasis{integer(c\_int)}\sphinxstyleemphasis{{]}} :: used to request local data

\end{itemize}

\end{description}\end{quote}

\end{fulllineitems}

\index{lammps\_type (fortran type)@\spxentry{lammps\_type}\spxextra{fortran type}|spxpagem}

\begin{fulllineitems}
\phantomsection\label{\detokenize{Fortran:f/_/lammps_type}}
\pysigstartsignatures
\pysigline{\sphinxbfcode{\sphinxupquote{type  }}\sphinxbfcode{\sphinxupquote{lammps\_type}}}
\pysigstopsignatures
\sphinxAtStartPar
This derived type is there to provide a convenient interface for the type
constants used with \sphinxcode{\sphinxupquote{extract\_compute()}}, \sphinxcode{\sphinxupquote{extract\_fix()}}, and
\sphinxcode{\sphinxupquote{extract\_variable()}}. Assuming your LAMMPS instance is called \sphinxcode{\sphinxupquote{lmp}},
these constants will be \sphinxcode{\sphinxupquote{lmp\%type\%scalar}}, \sphinxcode{\sphinxupquote{lmp\%type\%vector}}, and
\sphinxcode{\sphinxupquote{lmp\%type\%array}}. These values are identical to the values described
in {\hyperref[\detokenize{Library_objects:_CPPv415_LMP_TYPE_CONST}]{\sphinxcrossref{\sphinxcode{\sphinxupquote{\_LMP\_TYPE\_CONST}}}}} for the C library interface.
\begin{quote}\begin{description}
\sphinxlineitem{Type fields}\begin{itemize}
\item {} 
\sphinxAtStartPar
\sphinxbfcode{\sphinxupquote{\%~}}\sphinxbfcode{\sphinxupquote{scalar}}\sphinxstyleemphasis{ {[}}\sphinxstyleemphasis{integer(c\_int)}\sphinxstyleemphasis{{]}} :: used to request scalars

\item {} 
\sphinxAtStartPar
\sphinxbfcode{\sphinxupquote{\%~}}\sphinxbfcode{\sphinxupquote{vector}}\sphinxstyleemphasis{ {[}}\sphinxstyleemphasis{integer(c\_int)}\sphinxstyleemphasis{{]}} :: used to request vectors

\item {} 
\sphinxAtStartPar
\sphinxbfcode{\sphinxupquote{\%~}}\sphinxbfcode{\sphinxupquote{array}}\sphinxstyleemphasis{ {[}}\sphinxstyleemphasis{integer(c\_int)}\sphinxstyleemphasis{{]}} :: used to request arrays (matrices)

\end{itemize}

\end{description}\end{quote}

\end{fulllineitems}



\subsubsection{Procedures Bound to the \sphinxstyleliteralintitle{\sphinxupquote{lammps}} Derived Type}
\label{\detokenize{Fortran:procedures-bound-to-the-lammps-derived-type}}\index{close() (fortran subroutine)@\spxentry{close()}\spxextra{fortran subroutine}|spxpagem}

\begin{fulllineitems}
\phantomsection\label{\detokenize{Fortran:f/_/close}}
\pysigstartsignatures
\pysiglinewithargsret{\sphinxbfcode{\sphinxupquote{subroutine  }}\sphinxbfcode{\sphinxupquote{close}}}{\sphinxoptional{\sphinxparam{finalize}}}{}
\pysigstopsignatures
\sphinxAtStartPar
This method will close down the LAMMPS instance through calling
{\hyperref[\detokenize{Library_create:_CPPv412lammps_closePv}]{\sphinxcrossref{\sphinxcode{\sphinxupquote{lammps\_close()}}}}}.  If the \sphinxstyleemphasis{finalize} argument is present and
has a value of \sphinxcode{\sphinxupquote{.TRUE.}}, then this subroutine also calls
{\hyperref[\detokenize{Library_create:_CPPv422lammps_kokkos_finalizev}]{\sphinxcrossref{\sphinxcode{\sphinxupquote{lammps\_kokkos\_finalize()}}}}} and
{\hyperref[\detokenize{Library_create:_CPPv419lammps_mpi_finalizev}]{\sphinxcrossref{\sphinxcode{\sphinxupquote{lammps\_mpi\_finalize()}}}}}.
\begin{quote}\begin{description}
\sphinxlineitem{Options}
\sphinxAtStartPar
\sphinxstylestrong{finalize}\sphinxstyleemphasis{ {[}}\sphinxstyleemphasis{logical,optional}\sphinxstyleemphasis{{]}} :: shut down the MPI environment of the LAMMPS
library if \sphinxcode{\sphinxupquote{.TRUE.}}.

\sphinxlineitem{Call to}
\sphinxAtStartPar
{\hyperref[\detokenize{Library_create:_CPPv412lammps_closePv}]{\sphinxcrossref{\sphinxcode{\sphinxupquote{lammps\_close()}}}}} {\hyperref[\detokenize{Library_create:_CPPv419lammps_mpi_finalizev}]{\sphinxcrossref{\sphinxcode{\sphinxupquote{lammps\_mpi\_finalize()}}}}} {\hyperref[\detokenize{Library_create:_CPPv422lammps_kokkos_finalizev}]{\sphinxcrossref{\sphinxcode{\sphinxupquote{lammps\_kokkos\_finalize()}}}}}

\end{description}\end{quote}

\end{fulllineitems}



\bigskip\hrule\bigskip

\index{error() (fortran subroutine)@\spxentry{error()}\spxextra{fortran subroutine}|spxpagem}

\begin{fulllineitems}
\phantomsection\label{\detokenize{Fortran:f/_/error}}
\pysigstartsignatures
\pysiglinewithargsret{\sphinxbfcode{\sphinxupquote{subroutine  }}\sphinxbfcode{\sphinxupquote{error}}}{\sphinxparam{error\_type}\sphinxparamcomma \sphinxparam{error\_text}}{}
\pysigstopsignatures
\sphinxAtStartPar
This method is a wrapper around the {\hyperref[\detokenize{Library_create:_CPPv412lammps_errorPviPKc}]{\sphinxcrossref{\sphinxcode{\sphinxupquote{lammps\_error()}}}}} function and
will dispatch an error through the LAMMPS Error class.

\sphinxAtStartPar
\DUrole{versionmodified,added}{Added in version 3Nov2022.}
\begin{quote}\begin{description}
\sphinxlineitem{Parameters}\begin{itemize}
\item {} 
\sphinxAtStartPar
\sphinxstylestrong{error\_type}\sphinxstyleemphasis{ {[}}\sphinxstyleemphasis{integer(c\_int)}\sphinxstyleemphasis{{]}} :: constant to select which Error class function to call

\item {} 
\sphinxAtStartPar
\sphinxstylestrong{error\_text}\sphinxstyleemphasis{ {[}}\sphinxstyleemphasis{character(len=*)}\sphinxstyleemphasis{{]}} :: error message

\end{itemize}

\sphinxlineitem{Call to}
\sphinxAtStartPar
{\hyperref[\detokenize{Library_create:_CPPv412lammps_errorPviPKc}]{\sphinxcrossref{\sphinxcode{\sphinxupquote{lammps\_error()}}}}}

\end{description}\end{quote}

\end{fulllineitems}



\bigskip\hrule\bigskip

\index{file() (fortran subroutine)@\spxentry{file()}\spxextra{fortran subroutine}|spxpagem}

\begin{fulllineitems}
\phantomsection\label{\detokenize{Fortran:f/_/file}}
\pysigstartsignatures
\pysiglinewithargsret{\sphinxbfcode{\sphinxupquote{subroutine  }}\sphinxbfcode{\sphinxupquote{file}}}{\sphinxparam{filename}}{}
\pysigstopsignatures
\sphinxAtStartPar
This method will call {\hyperref[\detokenize{Library_execute:_CPPv411lammps_filePvPKc}]{\sphinxcrossref{\sphinxcode{\sphinxupquote{lammps\_file()}}}}} to have LAMMPS read
and process commands from a file.
\begin{quote}\begin{description}
\sphinxlineitem{Parameters}
\sphinxAtStartPar
\sphinxstylestrong{filename}\sphinxstyleemphasis{ {[}}\sphinxstyleemphasis{character(len=*)}\sphinxstyleemphasis{{]}} :: name of file with LAMMPS commands

\sphinxlineitem{Call to}
\sphinxAtStartPar
{\hyperref[\detokenize{Library_execute:_CPPv411lammps_filePvPKc}]{\sphinxcrossref{\sphinxcode{\sphinxupquote{lammps\_file()}}}}}

\end{description}\end{quote}

\end{fulllineitems}



\bigskip\hrule\bigskip

\index{command() (fortran subroutine)@\spxentry{command()}\spxextra{fortran subroutine}|spxpagem}

\begin{fulllineitems}
\phantomsection\label{\detokenize{Fortran:f/_/command}}
\pysigstartsignatures
\pysiglinewithargsret{\sphinxbfcode{\sphinxupquote{subroutine  }}\sphinxbfcode{\sphinxupquote{command}}}{\sphinxparam{cmd}}{}
\pysigstopsignatures
\sphinxAtStartPar
This method will call {\hyperref[\detokenize{Library_execute:_CPPv414lammps_commandPvPKc}]{\sphinxcrossref{\sphinxcode{\sphinxupquote{lammps\_command()}}}}} to have LAMMPS
execute a single command.
\begin{quote}\begin{description}
\sphinxlineitem{Parameters}
\sphinxAtStartPar
\sphinxstylestrong{cmd}\sphinxstyleemphasis{ {[}}\sphinxstyleemphasis{character(len=*)}\sphinxstyleemphasis{{]}} :: single LAMMPS command

\sphinxlineitem{Call to}
\sphinxAtStartPar
{\hyperref[\detokenize{Library_execute:_CPPv414lammps_commandPvPKc}]{\sphinxcrossref{\sphinxcode{\sphinxupquote{lammps\_command()}}}}}

\end{description}\end{quote}

\end{fulllineitems}



\bigskip\hrule\bigskip

\index{commands\_list() (fortran subroutine)@\spxentry{commands\_list()}\spxextra{fortran subroutine}|spxpagem}

\begin{fulllineitems}
\phantomsection\label{\detokenize{Fortran:f/_/commands_list}}
\pysigstartsignatures
\pysiglinewithargsret{\sphinxbfcode{\sphinxupquote{subroutine  }}\sphinxbfcode{\sphinxupquote{commands\_list}}}{\sphinxparam{cmds}}{}
\pysigstopsignatures
\sphinxAtStartPar
This method will call {\hyperref[\detokenize{Library_execute:_CPPv420lammps_commands_listPviPPKc}]{\sphinxcrossref{\sphinxcode{\sphinxupquote{lammps\_commands\_list()}}}}} to have LAMMPS
execute a list of input lines.
\begin{quote}\begin{description}
\sphinxlineitem{Parameters}
\sphinxAtStartPar
\sphinxstylestrong{cmd}\sphinxstyleemphasis{ {[}}\sphinxstyleemphasis{character(len=*)}\sphinxstyleemphasis{,}\sphinxstyleemphasis{dimension(:)}\sphinxstyleemphasis{{]}} :: list of LAMMPS input lines

\sphinxlineitem{Call to}
\sphinxAtStartPar
{\hyperref[\detokenize{Library_execute:_CPPv420lammps_commands_listPviPPKc}]{\sphinxcrossref{\sphinxcode{\sphinxupquote{lammps\_commands\_list()}}}}}

\end{description}\end{quote}

\end{fulllineitems}



\bigskip\hrule\bigskip

\index{commands\_string() (fortran subroutine)@\spxentry{commands\_string()}\spxextra{fortran subroutine}|spxpagem}

\begin{fulllineitems}
\phantomsection\label{\detokenize{Fortran:f/_/commands_string}}
\pysigstartsignatures
\pysiglinewithargsret{\sphinxbfcode{\sphinxupquote{subroutine  }}\sphinxbfcode{\sphinxupquote{commands\_string}}}{\sphinxparam{str}}{}
\pysigstopsignatures
\sphinxAtStartPar
This method will call {\hyperref[\detokenize{Library_execute:_CPPv422lammps_commands_stringPvPKc}]{\sphinxcrossref{\sphinxcode{\sphinxupquote{lammps\_commands\_string()}}}}} to have LAMMPS
execute a block of commands from a string.
\begin{quote}\begin{description}
\sphinxlineitem{Parameters}
\sphinxAtStartPar
\sphinxstylestrong{str}\sphinxstyleemphasis{ {[}}\sphinxstyleemphasis{character(len=*)}\sphinxstyleemphasis{{]}} :: LAMMPS input in string

\sphinxlineitem{Call to}
\sphinxAtStartPar
{\hyperref[\detokenize{Library_execute:_CPPv422lammps_commands_stringPvPKc}]{\sphinxcrossref{\sphinxcode{\sphinxupquote{lammps\_commands\_string()}}}}}

\end{description}\end{quote}

\end{fulllineitems}



\bigskip\hrule\bigskip

\index{get\_natoms() (fortran function)@\spxentry{get\_natoms()}\spxextra{fortran function}|spxpagem}

\begin{fulllineitems}
\phantomsection\label{\detokenize{Fortran:f/_/get_natoms}}
\pysigstartsignatures
\pysiglinewithargsret{\sphinxbfcode{\sphinxupquote{function  }}\sphinxbfcode{\sphinxupquote{get\_natoms}}}{}{}
\pysigstopsignatures
\sphinxAtStartPar
This function will call {\hyperref[\detokenize{Library_properties:_CPPv417lammps_get_natomsPv}]{\sphinxcrossref{\sphinxcode{\sphinxupquote{lammps\_get\_natoms()}}}}} and return the number
of atoms in the system.
\begin{quote}\begin{description}
\sphinxlineitem{Call to}
\sphinxAtStartPar
{\hyperref[\detokenize{Library_properties:_CPPv417lammps_get_natomsPv}]{\sphinxcrossref{\sphinxcode{\sphinxupquote{lammps\_get\_natoms()}}}}}

\sphinxlineitem{Return}
\sphinxAtStartPar
\sphinxstylestrong{natoms}\sphinxstyleemphasis{ {[}}\sphinxstyleemphasis{real(c\_double)}\sphinxstyleemphasis{{]}} :: number of atoms

\end{description}\end{quote}

\begin{sphinxadmonition}{note}{Note:}
\sphinxAtStartPar
If you would prefer to get the number of atoms in its native format
(i.e., as a 32\sphinxhyphen{} or 64\sphinxhyphen{}bit integer, depending on how LAMMPS was compiled),
this can be extracted with \sphinxcode{\sphinxupquote{extract\_global()}}.
\end{sphinxadmonition}

\end{fulllineitems}



\bigskip\hrule\bigskip

\index{get\_thermo() (fortran function)@\spxentry{get\_thermo()}\spxextra{fortran function}|spxpagem}

\begin{fulllineitems}
\phantomsection\label{\detokenize{Fortran:f/_/get_thermo}}
\pysigstartsignatures
\pysiglinewithargsret{\sphinxbfcode{\sphinxupquote{function  }}\sphinxbfcode{\sphinxupquote{get\_thermo}}}{\sphinxparam{name}}{}
\pysigstopsignatures
\sphinxAtStartPar
This function will call {\hyperref[\detokenize{Library_properties:_CPPv417lammps_get_thermoPvPKc}]{\sphinxcrossref{\sphinxcode{\sphinxupquote{lammps\_get\_thermo()}}}}} and return the value
of the corresponding thermodynamic keyword.

\sphinxAtStartPar
\DUrole{versionmodified,added}{Added in version 3Nov2022.}
\begin{quote}\begin{description}
\sphinxlineitem{Parameters}
\sphinxAtStartPar
\sphinxstylestrong{name}\sphinxstyleemphasis{ {[}}\sphinxstyleemphasis{character(len=*)}\sphinxstyleemphasis{{]}} :: string with the name of the thermo keyword

\sphinxlineitem{Call to}
\sphinxAtStartPar
{\hyperref[\detokenize{Library_properties:_CPPv417lammps_get_thermoPvPKc}]{\sphinxcrossref{\sphinxcode{\sphinxupquote{lammps\_get\_thermo()}}}}}

\sphinxlineitem{Return}
\sphinxAtStartPar
\sphinxstylestrong{value}\sphinxstyleemphasis{ {[}}\sphinxstyleemphasis{real(c\_double)}\sphinxstyleemphasis{{]}} :: value of the requested thermo property or \sphinxtitleref{0.0\_c\_double}

\end{description}\end{quote}

\end{fulllineitems}



\bigskip\hrule\bigskip

\index{last\_thermo() (fortran function)@\spxentry{last\_thermo()}\spxextra{fortran function}|spxpagem}

\begin{fulllineitems}
\phantomsection\label{\detokenize{Fortran:f/_/last_thermo}}
\pysigstartsignatures
\pysiglinewithargsret{\sphinxbfcode{\sphinxupquote{function  }}\sphinxbfcode{\sphinxupquote{last\_thermo}}}{\sphinxparam{what}\sphinxparamcomma \sphinxparam{index}}{}
\pysigstopsignatures
\sphinxAtStartPar
This function will call {\hyperref[\detokenize{Library_properties:_CPPv418lammps_last_thermoPvPKci}]{\sphinxcrossref{\sphinxcode{\sphinxupquote{lammps\_last\_thermo()}}}}} and returns
either a string or a pointer to a cached copy of LAMMPS last thermodynamic
output, depending on the data requested through \sphinxstyleemphasis{what}.  Note that \sphinxstyleemphasis{index}
uses 1\sphinxhyphen{}based indexing to access thermo output columns.

\sphinxAtStartPar
\DUrole{versionmodified,added}{Added in version 15Jun2023.}

\sphinxAtStartPar
Note that this function actually does not return a value, but rather
associates the pointer on the left side of the assignment to point to
internal LAMMPS data (with the exception of string data, which are
copied and returned as ordinary Fortran strings).  Pointers must be
of the correct data type to point to said data (typically
\sphinxcode{\sphinxupquote{INTEGER(c\_int)}}, \sphinxcode{\sphinxupquote{INTEGER(c\_int64\_t)}}, or \sphinxcode{\sphinxupquote{REAL(c\_double)}}).
The pointer being associated with LAMMPS data is type\sphinxhyphen{}checked at
run\sphinxhyphen{}time via an overloaded assignment operator.  The pointers
returned by this function point to temporary, read\sphinxhyphen{}only data that may
be overwritten at any time, so their target values need to be copied
to local storage if they are supposed to persist.

\sphinxAtStartPar
For example,

\begin{sphinxVerbatim}[commandchars=\\\{\}]
\PYG{k}{PROGRAM }\PYG{n}{thermo}
\PYG{+w}{  }\PYG{k}{USE }\PYG{n}{LIBLAMMPS}
\PYG{+w}{  }\PYG{k}{USE}\PYG{p}{,}\PYG{+w}{ }\PYG{k}{INTRINSIC}\PYG{+w}{ }\PYG{k+kd}{::}\PYG{+w}{ }\PYG{n+nb}{ISO\PYGZus{}C\PYGZus{}BINDING}\PYG{p}{,}\PYG{+w}{ }\PYG{k}{ONLY}\PYG{+w}{ }\PYG{p}{:}\PYG{+w}{ }\PYG{k+kt}{c\PYGZus{}double}\PYG{p}{,}\PYG{+w}{ }\PYG{k+kt}{c\PYGZus{}int64\PYGZus{}t}\PYG{p}{,}\PYG{+w}{ }\PYG{k+kt}{c\PYGZus{}int}
\PYG{k+kt}{  }\PYG{k}{TYPE}\PYG{p}{(}\PYG{n}{lammps}\PYG{p}{)}\PYG{+w}{ }\PYG{k+kd}{::}\PYG{+w}{ }\PYG{n}{lmp}
\PYG{+w}{  }\PYG{k+kt}{INTEGER}\PYG{p}{(}\PYG{n+nb}{KIND}\PYG{o}{=}\PYG{k+kt}{c\PYGZus{}int64\PYGZus{}t}\PYG{p}{)}\PYG{p}{,}\PYG{+w}{ }\PYG{k}{POINTER}\PYG{+w}{ }\PYG{k+kd}{::}\PYG{+w}{ }\PYG{n}{ntimestep}\PYG{p}{,}\PYG{+w}{ }\PYG{n}{bval}
\PYG{+w}{  }\PYG{k+kt}{REAL}\PYG{p}{(}\PYG{n+nb}{KIND}\PYG{o}{=}\PYG{k+kt}{c\PYGZus{}double}\PYG{p}{)}\PYG{p}{,}\PYG{+w}{ }\PYG{k}{POINTER}\PYG{+w}{ }\PYG{k+kd}{::}\PYG{+w}{ }\PYG{n}{dval}
\PYG{+w}{  }\PYG{k+kt}{INTEGER}\PYG{p}{(}\PYG{n+nb}{KIND}\PYG{o}{=}\PYG{k+kt}{c\PYGZus{}int}\PYG{p}{)}\PYG{p}{,}\PYG{+w}{ }\PYG{k}{POINTER}\PYG{+w}{ }\PYG{k+kd}{::}\PYG{+w}{ }\PYG{n}{nfield}\PYG{p}{,}\PYG{+w}{ }\PYG{n}{typ}\PYG{p}{,}\PYG{+w}{ }\PYG{n}{ival}
\PYG{+w}{  }\PYG{k+kt}{INTEGER}\PYG{p}{(}\PYG{n+nb}{KIND}\PYG{o}{=}\PYG{k+kt}{c\PYGZus{}int}\PYG{p}{)}\PYG{+w}{ }\PYG{k+kd}{::}\PYG{+w}{ }\PYG{n}{i}
\PYG{+w}{  }\PYG{k+kt}{CHARACTER}\PYG{p}{(}\PYG{n+nb}{LEN}\PYG{o}{=}\PYG{l+m+mi}{11}\PYG{p}{)}\PYG{+w}{ }\PYG{k+kd}{::}\PYG{+w}{ }\PYG{n}{key}

\PYG{+w}{  }\PYG{n}{lmp}\PYG{+w}{ }\PYG{o}{=}\PYG{+w}{ }\PYG{n}{lammps}\PYG{p}{(}\PYG{p}{)}
\PYG{+w}{  }\PYG{k}{CALL }\PYG{n}{lmp}\PYG{p}{\PYGZpc{}}\PYG{k}{file}\PYG{p}{(}\PYG{l+s+s1}{\PYGZsq{}in.sysinit\PYGZsq{}}\PYG{p}{)}

\PYG{+w}{  }\PYG{n}{ntimestep}\PYG{+w}{ }\PYG{o}{=}\PYG{+w}{ }\PYG{n}{lmp}\PYG{p}{\PYGZpc{}}\PYG{n}{last\PYGZus{}thermo}\PYG{p}{(}\PYG{l+s+s1}{\PYGZsq{}step\PYGZsq{}}\PYG{p}{,}\PYG{+w}{ }\PYG{l+m+mi}{0}\PYG{p}{)}
\PYG{+w}{  }\PYG{n}{nfield}\PYG{+w}{ }\PYG{o}{=}\PYG{+w}{ }\PYG{n}{lmp}\PYG{p}{\PYGZpc{}}\PYG{n}{last\PYGZus{}thermo}\PYG{p}{(}\PYG{l+s+s1}{\PYGZsq{}num\PYGZsq{}}\PYG{p}{,}\PYG{+w}{ }\PYG{l+m+mi}{0}\PYG{p}{)}
\PYG{+w}{  }\PYG{k}{PRINT}\PYG{o}{*}\PYG{p}{,}\PYG{+w}{ }\PYG{l+s+s1}{\PYGZsq{}Last thermo output on step: \PYGZsq{}}\PYG{p}{,}\PYG{+w}{ }\PYG{n}{ntimestep}\PYG{p}{,}\PYG{+w}{ }\PYG{l+s+s1}{\PYGZsq{}  Number of fields: \PYGZsq{}}\PYG{p}{,}\PYG{+w}{ }\PYG{n}{nfield}
\PYG{+w}{  }\PYG{k}{DO }\PYG{n}{i}\PYG{o}{=}\PYG{l+m+mi}{1}\PYG{p}{,}\PYG{+w}{ }\PYG{n}{nfield}
\PYG{+w}{      }\PYG{n}{key}\PYG{+w}{ }\PYG{o}{=}\PYG{+w}{ }\PYG{n}{lmp}\PYG{p}{\PYGZpc{}}\PYG{n}{last\PYGZus{}thermo}\PYG{p}{(}\PYG{l+s+s1}{\PYGZsq{}keyword\PYGZsq{}}\PYG{p}{,}\PYG{n}{i}\PYG{p}{)}
\PYG{+w}{      }\PYG{n}{typ}\PYG{+w}{ }\PYG{o}{=}\PYG{+w}{ }\PYG{n}{lmp}\PYG{p}{\PYGZpc{}}\PYG{n}{last\PYGZus{}thermo}\PYG{p}{(}\PYG{l+s+s1}{\PYGZsq{}type\PYGZsq{}}\PYG{p}{,}\PYG{n}{i}\PYG{p}{)}
\PYG{+w}{      }\PYG{k}{IF}\PYG{+w}{ }\PYG{p}{(}\PYG{n}{typ}\PYG{+w}{ }\PYG{o}{==}\PYG{+w}{ }\PYG{n}{lmp}\PYG{p}{\PYGZpc{}}\PYG{n}{dtype}\PYG{p}{\PYGZpc{}}\PYG{n}{i32}\PYG{p}{)}\PYG{+w}{ }\PYG{k}{THEN}
\PYG{k}{          }\PYG{n}{ival}\PYG{+w}{ }\PYG{o}{=}\PYG{+w}{ }\PYG{n}{lmp}\PYG{p}{\PYGZpc{}}\PYG{n}{last\PYGZus{}thermo}\PYG{p}{(}\PYG{l+s+s1}{\PYGZsq{}data\PYGZsq{}}\PYG{p}{,}\PYG{n}{i}\PYG{p}{)}
\PYG{+w}{          }\PYG{k}{PRINT}\PYG{o}{*}\PYG{p}{,}\PYG{+w}{ }\PYG{n}{key}\PYG{p}{,}\PYG{+w}{ }\PYG{l+s+s1}{\PYGZsq{}:\PYGZsq{}}\PYG{p}{,}\PYG{+w}{ }\PYG{n}{ival}
\PYG{+w}{      }\PYG{k}{ELSE }\PYG{k}{IF}\PYG{+w}{ }\PYG{p}{(}\PYG{n}{typ}\PYG{+w}{ }\PYG{o}{==}\PYG{+w}{ }\PYG{n}{lmp}\PYG{p}{\PYGZpc{}}\PYG{n}{dtype}\PYG{p}{\PYGZpc{}}\PYG{n}{i64}\PYG{p}{)}\PYG{+w}{ }\PYG{k}{THEN}
\PYG{k}{          }\PYG{n}{bval}\PYG{+w}{ }\PYG{o}{=}\PYG{+w}{ }\PYG{n}{lmp}\PYG{p}{\PYGZpc{}}\PYG{n}{last\PYGZus{}thermo}\PYG{p}{(}\PYG{l+s+s1}{\PYGZsq{}data\PYGZsq{}}\PYG{p}{,}\PYG{n}{i}\PYG{p}{)}
\PYG{+w}{          }\PYG{k}{PRINT}\PYG{o}{*}\PYG{p}{,}\PYG{+w}{ }\PYG{n}{key}\PYG{p}{,}\PYG{+w}{ }\PYG{l+s+s1}{\PYGZsq{}:\PYGZsq{}}\PYG{p}{,}\PYG{+w}{ }\PYG{n}{bval}
\PYG{+w}{      }\PYG{k}{ELSE }\PYG{k}{IF}\PYG{+w}{ }\PYG{p}{(}\PYG{n}{typ}\PYG{+w}{ }\PYG{o}{==}\PYG{+w}{ }\PYG{n}{lmp}\PYG{p}{\PYGZpc{}}\PYG{n}{dtype}\PYG{p}{\PYGZpc{}}\PYG{n}{r64}\PYG{p}{)}\PYG{+w}{ }\PYG{k}{THEN}
\PYG{k}{          }\PYG{n}{dval}\PYG{+w}{ }\PYG{o}{=}\PYG{+w}{ }\PYG{n}{lmp}\PYG{p}{\PYGZpc{}}\PYG{n}{last\PYGZus{}thermo}\PYG{p}{(}\PYG{l+s+s1}{\PYGZsq{}data\PYGZsq{}}\PYG{p}{,}\PYG{n}{i}\PYG{p}{)}
\PYG{+w}{          }\PYG{k}{PRINT}\PYG{o}{*}\PYG{p}{,}\PYG{+w}{ }\PYG{n}{key}\PYG{p}{,}\PYG{+w}{ }\PYG{l+s+s1}{\PYGZsq{}:\PYGZsq{}}\PYG{p}{,}\PYG{+w}{ }\PYG{n}{dval}
\PYG{+w}{      }\PYG{k}{END }\PYG{k}{IF}
\PYG{k}{  }\PYG{k}{END }\PYG{k}{DO}
\PYG{k}{  }\PYG{k}{CALL }\PYG{n}{lmp}\PYG{p}{\PYGZpc{}}\PYG{k}{close}\PYG{p}{(}\PYG{p}{.}\PYG{n}{TRUE}\PYG{p}{.}\PYG{p}{)}
\PYG{k}{END }\PYG{k}{PROGRAM }\PYG{n}{thermo}
\end{sphinxVerbatim}

\sphinxAtStartPar
would extract the last timestep where thermo output was done and the number
of columns it printed.  Then it loops over the columns to print out column
header keywords and the corresponding data.

\begin{sphinxadmonition}{note}{Note:}
\sphinxAtStartPar
If \sphinxcode{\sphinxupquote{last\_thermo()}} returns a string, the string must have a length
greater than or equal to the length of the string (not including the
terminal \sphinxcode{\sphinxupquote{NULL}} character) that LAMMPS returns. If the variable’s
length is too short, the string will be truncated.  As usual in Fortran,
strings are padded with spaces at the end.  If you use an allocatable
string, the string \sphinxstylestrong{must be allocated} prior to calling this function.
\end{sphinxadmonition}
\begin{quote}\begin{description}
\sphinxlineitem{Parameters}\begin{itemize}
\item {} 
\sphinxAtStartPar
\sphinxstylestrong{what}\sphinxstyleemphasis{ {[}}\sphinxstyleemphasis{character(len=*)}\sphinxstyleemphasis{{]}} :: string with the name of the thermo keyword

\item {} 
\sphinxAtStartPar
\sphinxstylestrong{index}\sphinxstyleemphasis{ {[}}\sphinxstyleemphasis{integer(c\_int)}\sphinxstyleemphasis{{]}} :: 1\sphinxhyphen{}based column index

\end{itemize}

\sphinxlineitem{Call to}
\sphinxAtStartPar
{\hyperref[\detokenize{Library_properties:_CPPv418lammps_last_thermoPvPKci}]{\sphinxcrossref{\sphinxcode{\sphinxupquote{lammps\_last\_thermo()}}}}}

\sphinxlineitem{Return}
\sphinxAtStartPar
\sphinxstylestrong{pointer}\sphinxstyleemphasis{ {[}}\sphinxstyleemphasis{polymorphic}\sphinxstyleemphasis{{]}} :: pointer to LAMMPS data. The left\sphinxhyphen{}hand side of the
assignment should be either a string (if expecting string data) or a
C\sphinxhyphen{}compatible pointer (e.g., \sphinxcode{\sphinxupquote{INTEGER(c\_int), POINTER :: nlocal}}) to the
extracted property.

\end{description}\end{quote}

\begin{sphinxadmonition}{warning}{Warning:}
\sphinxAtStartPar
Modifying the data in the location pointed to by the returned pointer
may lead to inconsistent internal data and thus may cause failures,
crashes, or bogus simulations.  In general, it is much better
to use a LAMMPS input command that sets or changes these parameters.
Using an input command will take care of all side effects and necessary
updates of settings derived from such settings.
\end{sphinxadmonition}

\end{fulllineitems}



\bigskip\hrule\bigskip

\index{extract\_box() (fortran subroutine)@\spxentry{extract\_box()}\spxextra{fortran subroutine}|spxpagem}

\begin{fulllineitems}
\phantomsection\label{\detokenize{Fortran:f/_/extract_box}}
\pysigstartsignatures
\pysiglinewithargsret{\sphinxbfcode{\sphinxupquote{subroutine  }}\sphinxbfcode{\sphinxupquote{extract\_box}}}{\sphinxparam{{[}boxlo{]}{[}, boxhi{]}{[}, xy{]}{[}, yz{]}{[}, xz{]}{[}, pflags{]}{[}, boxflag{]}}}{}
\pysigstopsignatures
\sphinxAtStartPar
This subroutine will call {\hyperref[\detokenize{Library_properties:_CPPv418lammps_extract_boxPvPdPdPdPdPdPiPi}]{\sphinxcrossref{\sphinxcode{\sphinxupquote{lammps\_extract\_box()}}}}}. All
parameters are optional, though obviously at least one should be
present. The parameters \sphinxstyleemphasis{pflags} and \sphinxstyleemphasis{boxflag} are stored in LAMMPS
as integers, but should be declared as \sphinxcode{\sphinxupquote{LOGICAL}} variables when
calling from Fortran.

\sphinxAtStartPar
\DUrole{versionmodified,added}{Added in version 3Nov2022.}
\begin{quote}\begin{description}
\sphinxlineitem{Options}\begin{itemize}
\item {} 
\sphinxAtStartPar
\sphinxstylestrong{boxlo}\sphinxstyleemphasis{ {[}}\sphinxstyleemphasis{real(c\_double)}\sphinxstyleemphasis{,}\sphinxstyleemphasis{dimension(3),optional}\sphinxstyleemphasis{{]}} :: vector in which to store
lower\sphinxhyphen{}bounds of simulation box

\item {} 
\sphinxAtStartPar
\sphinxstylestrong{boxhi}\sphinxstyleemphasis{ {[}}\sphinxstyleemphasis{real(c\_double)}\sphinxstyleemphasis{,}\sphinxstyleemphasis{dimension(3),optional}\sphinxstyleemphasis{{]}} :: vector in which to store
upper\sphinxhyphen{}bounds of simulation box

\item {} 
\sphinxAtStartPar
\sphinxstylestrong{xy}\sphinxstyleemphasis{ {[}}\sphinxstyleemphasis{real(c\_double)}\sphinxstyleemphasis{,}\sphinxstyleemphasis{optional}\sphinxstyleemphasis{{]}} :: variable in which to store \sphinxstyleemphasis{xy} tilt factor

\item {} 
\sphinxAtStartPar
\sphinxstylestrong{yz}\sphinxstyleemphasis{ {[}}\sphinxstyleemphasis{real(c\_double)}\sphinxstyleemphasis{,}\sphinxstyleemphasis{optional}\sphinxstyleemphasis{{]}} :: variable in which to store \sphinxstyleemphasis{yz} tilt factor

\item {} 
\sphinxAtStartPar
\sphinxstylestrong{xz}\sphinxstyleemphasis{ {[}}\sphinxstyleemphasis{real(c\_double)}\sphinxstyleemphasis{,}\sphinxstyleemphasis{optional}\sphinxstyleemphasis{{]}} :: variable in which to store \sphinxstyleemphasis{xz} tilt factor

\item {} 
\sphinxAtStartPar
\sphinxstylestrong{pflags}\sphinxstyleemphasis{ {[}}\sphinxstyleemphasis{logical,dimension(3),optional}\sphinxstyleemphasis{{]}} :: vector in which to store
periodicity flags (\sphinxcode{\sphinxupquote{.TRUE.}} means periodic in that dimension)

\item {} 
\sphinxAtStartPar
\sphinxstylestrong{boxflag}\sphinxstyleemphasis{ {[}}\sphinxstyleemphasis{logical,optional}\sphinxstyleemphasis{{]}} :: variable in which to store boolean denoting
whether the box will change during a simulation
(\sphinxcode{\sphinxupquote{.TRUE.}} means box will change)

\end{itemize}

\sphinxlineitem{Call to}
\sphinxAtStartPar
{\hyperref[\detokenize{Library_properties:_CPPv418lammps_extract_boxPvPdPdPdPdPdPiPi}]{\sphinxcrossref{\sphinxcode{\sphinxupquote{lammps\_extract\_box()}}}}}

\end{description}\end{quote}

\end{fulllineitems}


\begin{sphinxadmonition}{note}{Note:}
\sphinxAtStartPar
Note that a frequent use case of this function is to extract only one or
more of the options rather than all seven. For example, assuming “lmp”
represents a properly\sphinxhyphen{}initialized LAMMPS instance, the following code will
extract the periodic box settings into the variable “periodic”:

\begin{sphinxVerbatim}[commandchars=\\\{\}]
\PYG{c}{! code to start up}
\PYG{k+kt}{LOGICAL}\PYG{+w}{ }\PYG{k+kd}{::}\PYG{+w}{ }\PYG{n}{periodic}\PYG{p}{(}\PYG{l+m+mi}{3}\PYG{p}{)}
\PYG{c}{! code to initialize LAMMPS / run things / etc.}
\PYG{k}{CALL }\PYG{n}{lmp}\PYG{p}{\PYGZpc{}}\PYG{n}{extract\PYGZus{}box}\PYG{p}{(}\PYG{n}{pflags}\PYG{+w}{ }\PYG{o}{=}\PYG{+w}{ }\PYG{n}{periodic}\PYG{p}{)}
\end{sphinxVerbatim}
\end{sphinxadmonition}


\bigskip\hrule\bigskip

\index{reset\_box() (fortran subroutine)@\spxentry{reset\_box()}\spxextra{fortran subroutine}|spxpagem}

\begin{fulllineitems}
\phantomsection\label{\detokenize{Fortran:f/_/reset_box}}
\pysigstartsignatures
\pysiglinewithargsret{\sphinxbfcode{\sphinxupquote{subroutine  }}\sphinxbfcode{\sphinxupquote{reset\_box}}}{\sphinxparam{boxlo}\sphinxparamcomma \sphinxparam{boxhi}\sphinxparamcomma \sphinxparam{xy}\sphinxparamcomma \sphinxparam{yz}\sphinxparamcomma \sphinxparam{xz}}{}
\pysigstopsignatures
\sphinxAtStartPar
This subroutine will call {\hyperref[\detokenize{Library_properties:_CPPv416lammps_reset_boxPvPdPdddd}]{\sphinxcrossref{\sphinxcode{\sphinxupquote{lammps\_reset\_box()}}}}}. All parameters
are required.

\sphinxAtStartPar
\DUrole{versionmodified,added}{Added in version 3Nov2022.}
\begin{quote}\begin{description}
\sphinxlineitem{Parameters}\begin{itemize}
\item {} 
\sphinxAtStartPar
\sphinxstylestrong{boxlo}\sphinxstyleemphasis{ {[}}\sphinxstyleemphasis{real(c\_double)}\sphinxstyleemphasis{,}\sphinxstyleemphasis{dimension(3)}\sphinxstyleemphasis{{]}} :: vector of three doubles containing
the lower box boundary

\item {} 
\sphinxAtStartPar
\sphinxstylestrong{boxhi}\sphinxstyleemphasis{ {[}}\sphinxstyleemphasis{real(c\_double)}\sphinxstyleemphasis{,}\sphinxstyleemphasis{dimension(3)}\sphinxstyleemphasis{{]}} :: vector of three doubles containing
the upper box boundary

\item {} 
\sphinxAtStartPar
\sphinxstylestrong{xy}\sphinxstyleemphasis{ {[}}\sphinxstyleemphasis{real(c\_double)}\sphinxstyleemphasis{{]}} :: \sphinxstyleemphasis{x\textendash{}y} tilt factor

\item {} 
\sphinxAtStartPar
\sphinxstylestrong{yz}\sphinxstyleemphasis{ {[}}\sphinxstyleemphasis{real(c\_double)}\sphinxstyleemphasis{{]}} :: \sphinxstyleemphasis{y\textendash{}z} tilt factor

\item {} 
\sphinxAtStartPar
\sphinxstylestrong{xz}\sphinxstyleemphasis{ {[}}\sphinxstyleemphasis{real(c\_double)}\sphinxstyleemphasis{{]}} :: \sphinxstyleemphasis{x\textendash{}z} tilt factor

\end{itemize}

\sphinxlineitem{Call to}
\sphinxAtStartPar
{\hyperref[\detokenize{Library_properties:_CPPv416lammps_reset_boxPvPdPdddd}]{\sphinxcrossref{\sphinxcode{\sphinxupquote{lammps\_reset\_box()}}}}}

\end{description}\end{quote}

\end{fulllineitems}



\bigskip\hrule\bigskip

\index{memory\_usage() (fortran subroutine)@\spxentry{memory\_usage()}\spxextra{fortran subroutine}|spxpagem}

\begin{fulllineitems}
\phantomsection\label{\detokenize{Fortran:f/_/memory_usage}}
\pysigstartsignatures
\pysiglinewithargsret{\sphinxbfcode{\sphinxupquote{subroutine  }}\sphinxbfcode{\sphinxupquote{memory\_usage}}}{\sphinxparam{meminfo}}{}
\pysigstopsignatures
\sphinxAtStartPar
This subroutine will call {\hyperref[\detokenize{Library_properties:_CPPv419lammps_memory_usagePvPd}]{\sphinxcrossref{\sphinxcode{\sphinxupquote{lammps\_memory\_usage()}}}}} and store the
result in the three\sphinxhyphen{}element array \sphinxstyleemphasis{meminfo}.

\sphinxAtStartPar
\DUrole{versionmodified,added}{Added in version 3Nov2022.}
\begin{quote}\begin{description}
\sphinxlineitem{Parameters}
\sphinxAtStartPar
\sphinxstylestrong{meminfo}\sphinxstyleemphasis{ {[}}\sphinxstyleemphasis{real(c\_double)}\sphinxstyleemphasis{,}\sphinxstyleemphasis{dimension(3)}\sphinxstyleemphasis{{]}} :: vector of three doubles in which
to store memory usage data

\sphinxlineitem{Call to}
\sphinxAtStartPar
{\hyperref[\detokenize{Library_properties:_CPPv419lammps_memory_usagePvPd}]{\sphinxcrossref{\sphinxcode{\sphinxupquote{lammps\_memory\_usage()}}}}}

\end{description}\end{quote}

\end{fulllineitems}



\bigskip\hrule\bigskip

\index{get\_mpi\_comm() (fortran function)@\spxentry{get\_mpi\_comm()}\spxextra{fortran function}|spxpagem}

\begin{fulllineitems}
\phantomsection\label{\detokenize{Fortran:f/_/get_mpi_comm}}
\pysigstartsignatures
\pysiglinewithargsret{\sphinxbfcode{\sphinxupquote{function  }}\sphinxbfcode{\sphinxupquote{get\_mpi\_comm}}}{}{}
\pysigstopsignatures
\sphinxAtStartPar
This function returns a Fortran representation of the LAMMPS “world”
communicator.

\sphinxAtStartPar
\DUrole{versionmodified,added}{Added in version 3Nov2022.}
\begin{quote}\begin{description}
\sphinxlineitem{Call to}
\sphinxAtStartPar
{\hyperref[\detokenize{Library_properties:_CPPv419lammps_get_mpi_commPv}]{\sphinxcrossref{\sphinxcode{\sphinxupquote{lammps\_get\_mpi\_comm()}}}}}

\sphinxlineitem{Return}
\sphinxAtStartPar
\sphinxstylestrong{comm}\sphinxstyleemphasis{ {[}}\sphinxstyleemphasis{integer}\sphinxstyleemphasis{{]}} :: Fortran integer equivalent to the MPI communicator LAMMPS is
using

\end{description}\end{quote}

\begin{sphinxadmonition}{note}{Note:}
\sphinxAtStartPar
The C library interface currently returns type \sphinxcode{\sphinxupquote{int}} instead of
type \sphinxcode{\sphinxupquote{MPI\_Fint}}, which is the C type corresponding to Fortran
\sphinxcode{\sphinxupquote{INTEGER}} types of the default kind.  On most compilers, these
are the same anyway, but this interface exchanges values this way
to avoid warning messages.
\end{sphinxadmonition}

\begin{sphinxadmonition}{note}{Note:}
\sphinxAtStartPar
The \sphinxcode{\sphinxupquote{MPI\_F08}} module, which defines Fortran 2008 bindings for MPI,
is not directly supported by this function.  However, you should be
able to convert between the two using the \sphinxtitleref{MPI\_VAL} member of the
communicator.  For example,

\begin{sphinxVerbatim}[commandchars=\\\{\}]
\PYG{k}{USE }\PYG{n}{MPI\PYGZus{}F08}
\PYG{k}{USE }\PYG{n}{LIBLAMMPS}
\PYG{k}{TYPE}\PYG{p}{(}\PYG{n}{lammps}\PYG{p}{)}\PYG{+w}{ }\PYG{k+kd}{::}\PYG{+w}{ }\PYG{n}{lmp}
\PYG{k}{TYPE}\PYG{p}{(}\PYG{n}{MPI\PYGZus{}Comm}\PYG{p}{)}\PYG{+w}{ }\PYG{k+kd}{::}\PYG{+w}{ }\PYG{n}{comm}
\PYG{c}{! ... [commands to set up LAMMPS/etc.]}
\PYG{n}{comm}\PYG{p}{\PYGZpc{}}\PYG{n}{MPI\PYGZus{}VAL}\PYG{+w}{ }\PYG{o}{=}\PYG{+w}{ }\PYG{n}{lmp}\PYG{p}{\PYGZpc{}}\PYG{n}{get\PYGZus{}mpi\PYGZus{}comm}\PYG{p}{(}\PYG{p}{)}
\end{sphinxVerbatim}

\sphinxAtStartPar
should assign an \sphinxcode{\sphinxupquote{MPI\_F08}} communicator properly.
\end{sphinxadmonition}

\end{fulllineitems}



\bigskip\hrule\bigskip

\index{extract\_setting() (fortran function)@\spxentry{extract\_setting()}\spxextra{fortran function}|spxpagem}

\begin{fulllineitems}
\phantomsection\label{\detokenize{Fortran:f/_/extract_setting}}
\pysigstartsignatures
\pysiglinewithargsret{\sphinxbfcode{\sphinxupquote{function  }}\sphinxbfcode{\sphinxupquote{extract\_setting}}}{\sphinxparam{keyword}}{}
\pysigstopsignatures
\sphinxAtStartPar
Query LAMMPS about global settings. See the documentation for the
{\hyperref[\detokenize{Library_properties:_CPPv422lammps_extract_settingPvPKc}]{\sphinxcrossref{\sphinxcode{\sphinxupquote{lammps\_extract\_setting()}}}}} function from the C library.

\sphinxAtStartPar
\DUrole{versionmodified,added}{Added in version 3Nov2022.}
\begin{quote}\begin{description}
\sphinxlineitem{Parameters}
\sphinxAtStartPar
\sphinxstylestrong{keyword}\sphinxstyleemphasis{ {[}}\sphinxstyleemphasis{character(len=*)}\sphinxstyleemphasis{{]}} :: string containing the name of the thermo keyword

\sphinxlineitem{Call to}
\sphinxAtStartPar
{\hyperref[\detokenize{Library_properties:_CPPv422lammps_extract_settingPvPKc}]{\sphinxcrossref{\sphinxcode{\sphinxupquote{lammps\_extract\_setting()}}}}}

\sphinxlineitem{Return}
\sphinxAtStartPar
\sphinxstylestrong{setting}\sphinxstyleemphasis{ {[}}\sphinxstyleemphasis{integer(c\_int)}\sphinxstyleemphasis{{]}} :: value of the queried setting or \(-1\) if
unknown

\end{description}\end{quote}

\end{fulllineitems}



\bigskip\hrule\bigskip

\index{extract\_global() (fortran function)@\spxentry{extract\_global()}\spxextra{fortran function}|spxpagem}

\begin{fulllineitems}
\phantomsection\label{\detokenize{Fortran:f/_/extract_global}}
\pysigstartsignatures
\pysiglinewithargsret{\sphinxbfcode{\sphinxupquote{function  }}\sphinxbfcode{\sphinxupquote{extract\_global}}}{\sphinxparam{name}}{}
\pysigstopsignatures
\sphinxAtStartPar
This function calls {\hyperref[\detokenize{Library_properties:_CPPv421lammps_extract_globalPvPKc}]{\sphinxcrossref{\sphinxcode{\sphinxupquote{lammps\_extract\_global()}}}}} and returns
either a string or a pointer to internal global LAMMPS data,
depending on the data requested through \sphinxstyleemphasis{name}.

\sphinxAtStartPar
\DUrole{versionmodified,added}{Added in version 3Nov2022.}

\sphinxAtStartPar
Note that this function actually does not return a value, but rather
associates the pointer on the left side of the assignment to point to
internal LAMMPS data (with the exception of string data, which are
copied and returned as ordinary Fortran strings). Pointers must be of
the correct data type to point to said data (typically
\sphinxcode{\sphinxupquote{INTEGER(c\_int)}}, \sphinxcode{\sphinxupquote{INTEGER(c\_int64\_t)}}, or \sphinxcode{\sphinxupquote{REAL(c\_double)}})
and have compatible kind and rank.  The pointer being associated with
LAMMPS data is type\sphinxhyphen{}, kind\sphinxhyphen{}, and rank\sphinxhyphen{}checked at run\sphinxhyphen{}time via an
overloaded assignment operator.  The pointers returned by this
function are generally persistent; therefore it is not necessary to
call the function again unless a {\hyperref[\detokenize{clear::doc}]{\sphinxcrossref{\DUrole{doc}{clear command}}}} command has been
issued, which wipes out and recreates the contents of the
{\hyperref[\detokenize{Classes_lammps:_CPPv4N9LAMMPS_NS6LAMMPSE}]{\sphinxcrossref{\sphinxcode{\sphinxupquote{LAMMPS}}}}} class.

\sphinxAtStartPar
For example,

\begin{sphinxVerbatim}[commandchars=\\\{\}]
\PYG{k}{PROGRAM }\PYG{n}{demo}
\PYG{+w}{  }\PYG{k}{USE}\PYG{p}{,}\PYG{+w}{ }\PYG{k}{INTRINSIC}\PYG{+w}{ }\PYG{k+kd}{::}\PYG{+w}{ }\PYG{n+nb}{ISO\PYGZus{}C\PYGZus{}BINDING}\PYG{p}{,}\PYG{+w}{ }\PYG{k}{ONLY}\PYG{+w}{ }\PYG{p}{:}\PYG{+w}{ }\PYG{k+kt}{c\PYGZus{}int64\PYGZus{}t}\PYG{p}{,}\PYG{+w}{ }\PYG{k+kt}{c\PYGZus{}int}\PYG{p}{,}\PYG{+w}{ }\PYG{k+kt}{c\PYGZus{}double}
\PYG{k+kt}{  }\PYG{k}{USE }\PYG{n}{LIBLAMMPS}
\PYG{+w}{  }\PYG{k}{TYPE}\PYG{p}{(}\PYG{n}{lammps}\PYG{p}{)}\PYG{+w}{ }\PYG{k+kd}{::}\PYG{+w}{ }\PYG{n}{lmp}
\PYG{+w}{  }\PYG{k+kt}{INTEGER}\PYG{p}{(}\PYG{k+kt}{c\PYGZus{}int}\PYG{p}{)}\PYG{p}{,}\PYG{+w}{ }\PYG{k}{POINTER}\PYG{+w}{ }\PYG{k+kd}{::}\PYG{+w}{ }\PYG{n}{nlocal}\PYG{+w}{ }\PYG{o}{=}\PYG{o}{\PYGZgt{}}\PYG{+w}{ }\PYG{n+nb}{NULL}\PYG{p}{(}\PYG{p}{)}
\PYG{+w}{  }\PYG{k+kt}{INTEGER}\PYG{p}{(}\PYG{k+kt}{c\PYGZus{}int64\PYGZus{}t}\PYG{p}{)}\PYG{p}{,}\PYG{+w}{ }\PYG{k}{POINTER}\PYG{+w}{ }\PYG{k+kd}{::}\PYG{+w}{ }\PYG{n}{ntimestep}\PYG{+w}{ }\PYG{o}{=}\PYG{o}{\PYGZgt{}}\PYG{+w}{ }\PYG{n+nb}{NULL}\PYG{p}{(}\PYG{p}{)}
\PYG{+w}{  }\PYG{k+kt}{REAL}\PYG{p}{(}\PYG{k+kt}{c\PYGZus{}double}\PYG{p}{)}\PYG{p}{,}\PYG{+w}{ }\PYG{k}{POINTER}\PYG{+w}{ }\PYG{k+kd}{::}\PYG{+w}{ }\PYG{n}{dt}\PYG{+w}{ }\PYG{o}{=}\PYG{o}{\PYGZgt{}}\PYG{+w}{ }\PYG{n+nb}{NULL}\PYG{p}{(}\PYG{p}{)}
\PYG{+w}{  }\PYG{k+kt}{CHARACTER}\PYG{p}{(}\PYG{n+nb}{LEN}\PYG{o}{=}\PYG{l+m+mi}{10}\PYG{p}{)}\PYG{+w}{ }\PYG{k+kd}{::}\PYG{+w}{ }\PYG{n}{units}
\PYG{+w}{  }\PYG{n}{lmp}\PYG{+w}{ }\PYG{o}{=}\PYG{+w}{ }\PYG{n}{lammps}\PYG{p}{(}\PYG{p}{)}
\PYG{+w}{  }\PYG{c}{! other commands}
\PYG{+w}{  }\PYG{n}{nlocal}\PYG{+w}{ }\PYG{o}{=}\PYG{+w}{ }\PYG{n}{lmp}\PYG{p}{\PYGZpc{}}\PYG{n}{extract\PYGZus{}global}\PYG{p}{(}\PYG{l+s+s1}{\PYGZsq{}nlocal\PYGZsq{}}\PYG{p}{)}
\PYG{+w}{  }\PYG{n}{ntimestep}\PYG{+w}{ }\PYG{o}{=}\PYG{+w}{ }\PYG{n}{lmp}\PYG{p}{\PYGZpc{}}\PYG{n}{extract\PYGZus{}global}\PYG{p}{(}\PYG{l+s+s1}{\PYGZsq{}ntimestep\PYGZsq{}}\PYG{p}{)}
\PYG{+w}{  }\PYG{n}{dt}\PYG{+w}{ }\PYG{o}{=}\PYG{+w}{ }\PYG{n}{lmp}\PYG{p}{\PYGZpc{}}\PYG{n}{extract\PYGZus{}global}\PYG{p}{(}\PYG{l+s+s1}{\PYGZsq{}dt\PYGZsq{}}\PYG{p}{)}
\PYG{+w}{  }\PYG{n}{units}\PYG{+w}{ }\PYG{o}{=}\PYG{+w}{ }\PYG{n}{lmp}\PYG{p}{\PYGZpc{}}\PYG{n}{extract\PYGZus{}global}\PYG{p}{(}\PYG{l+s+s1}{\PYGZsq{}units\PYGZsq{}}\PYG{p}{)}
\PYG{+w}{  }\PYG{c}{! more commands}
\PYG{+w}{  }\PYG{n}{lmp}\PYG{p}{.}\PYG{k}{close}\PYG{p}{(}\PYG{p}{.}\PYG{n}{TRUE}\PYG{p}{.}\PYG{p}{)}
\PYG{k}{END }\PYG{k}{PROGRAM }\PYG{n}{demo}
\end{sphinxVerbatim}

\sphinxAtStartPar
would extract the number of atoms on this processor, the current time step,
the size of the current time step, and the units being used into the
variables \sphinxstyleemphasis{nlocal}, \sphinxstyleemphasis{ntimestep}, \sphinxstyleemphasis{dt}, and \sphinxstyleemphasis{units}, respectively.

\begin{sphinxadmonition}{note}{Note:}
\sphinxAtStartPar
If \sphinxcode{\sphinxupquote{extract\_global()}} returns a string, the string must have
a length greater than or equal to the length of the string (not
including the terminal \sphinxcode{\sphinxupquote{NULL}} character) that LAMMPS returns. If
the variable’s length is too short, the string will be
truncated. As usual in Fortran, strings are padded with spaces at
the end. If you use an allocatable string, the string \sphinxstylestrong{must be
allocated} prior to calling this function, but you can
automatically reallocate it to the correct length after the
function returns, viz.,

\begin{sphinxVerbatim}[commandchars=\\\{\}]
\PYG{k}{PROGRAM }\PYG{n}{test}
\PYG{+w}{  }\PYG{k}{USE }\PYG{n}{LIBLAMMPS}
\PYG{+w}{  }\PYG{k}{TYPE}\PYG{p}{(}\PYG{n}{lammps}\PYG{p}{)}\PYG{+w}{ }\PYG{k+kd}{::}\PYG{+w}{ }\PYG{n}{lmp}
\PYG{+w}{  }\PYG{k+kt}{CHARACTER}\PYG{p}{(}\PYG{n+nb}{LEN}\PYG{o}{=}\PYG{p}{:}\PYG{p}{)}\PYG{p}{,}\PYG{+w}{ }\PYG{k}{ALLOCATABLE}\PYG{+w}{ }\PYG{k+kd}{::}\PYG{+w}{ }\PYG{n}{str}
\PYG{+w}{  }\PYG{n}{lmp}\PYG{+w}{ }\PYG{o}{=}\PYG{+w}{ }\PYG{n}{lammps}\PYG{p}{(}\PYG{p}{)}
\PYG{+w}{  }\PYG{k}{CALL }\PYG{n}{lmp}\PYG{p}{\PYGZpc{}}\PYG{n}{command}\PYG{p}{(}\PYG{l+s+s1}{\PYGZsq{}units metal\PYGZsq{}}\PYG{p}{)}
\PYG{+w}{  }\PYG{k}{ALLOCATE}\PYG{p}{(}\PYG{k+kt}{CHARACTER}\PYG{p}{(}\PYG{n+nb}{LEN}\PYG{o}{=}\PYG{l+m+mi}{80}\PYG{p}{)}\PYG{+w}{ }\PYG{k+kd}{::}\PYG{+w}{ }\PYG{n}{str}\PYG{p}{)}
\PYG{+w}{  }\PYG{n}{str}\PYG{+w}{ }\PYG{o}{=}\PYG{+w}{ }\PYG{n}{lmp}\PYG{p}{\PYGZpc{}}\PYG{n}{extract\PYGZus{}global}\PYG{p}{(}\PYG{l+s+s1}{\PYGZsq{}units\PYGZsq{}}\PYG{p}{)}
\PYG{+w}{  }\PYG{n}{str}\PYG{+w}{ }\PYG{o}{=}\PYG{+w}{ }\PYG{n+nb}{TRIM}\PYG{p}{(}\PYG{n}{str}\PYG{p}{)}\PYG{+w}{ }\PYG{c}{! re\PYGZhy{}allocates to length len\PYGZus{}trim(str) here}
\PYG{+w}{  }\PYG{k}{PRINT}\PYG{o}{*}\PYG{p}{,}\PYG{+w}{ }\PYG{n+nb}{LEN}\PYG{p}{(}\PYG{n}{str}\PYG{p}{)}\PYG{p}{,}\PYG{+w}{ }\PYG{n+nb}{LEN\PYGZus{}TRIM}\PYG{p}{(}\PYG{n}{str}\PYG{p}{)}
\PYG{k}{END }\PYG{k}{PROGRAM }\PYG{n}{test}
\end{sphinxVerbatim}

\sphinxAtStartPar
will print the number 5 (the length of the word “metal”) twice.
\end{sphinxadmonition}
\begin{quote}\begin{description}
\sphinxlineitem{Parameters}
\sphinxAtStartPar
\sphinxstylestrong{name}\sphinxstyleemphasis{ {[}}\sphinxstyleemphasis{character(len=*)}\sphinxstyleemphasis{{]}} :: string with the name of the property to extract

\sphinxlineitem{Call to}
\sphinxAtStartPar
{\hyperref[\detokenize{Library_properties:_CPPv421lammps_extract_globalPvPKc}]{\sphinxcrossref{\sphinxcode{\sphinxupquote{lammps\_extract\_global()}}}}}

\sphinxlineitem{Return}
\sphinxAtStartPar
\sphinxstylestrong{pointer}\sphinxstyleemphasis{ {[}}\sphinxstyleemphasis{polymorphic}\sphinxstyleemphasis{{]}} :: pointer to LAMMPS data. The left\sphinxhyphen{}hand side of the
assignment should be either a string (if expecting string data) or a
C\sphinxhyphen{}compatible pointer (e.g., \sphinxcode{\sphinxupquote{INTEGER(c\_int), POINTER :: nlocal}}) to the
extracted property. If expecting vector data, the pointer should have
dimension “:”.

\end{description}\end{quote}

\begin{sphinxadmonition}{warning}{Warning:}
\sphinxAtStartPar
Modifying the data in the location pointed to by the returned pointer
may lead to inconsistent internal data and thus may cause failures,
crashes, or bogus simulations.  In general, it is much better
to use a LAMMPS input command that sets or changes these parameters.
Using an input command will take care of all side effects and necessary
updates of settings derived from such settings.
\end{sphinxadmonition}

\end{fulllineitems}



\bigskip\hrule\bigskip

\index{extract\_atom() (fortran function)@\spxentry{extract\_atom()}\spxextra{fortran function}|spxpagem}

\begin{fulllineitems}
\phantomsection\label{\detokenize{Fortran:f/_/extract_atom}}
\pysigstartsignatures
\pysiglinewithargsret{\sphinxbfcode{\sphinxupquote{function  }}\sphinxbfcode{\sphinxupquote{extract\_atom}}}{\sphinxparam{name}}{}
\pysigstopsignatures
\sphinxAtStartPar
This function calls {\hyperref[\detokenize{Library_atoms:_CPPv419lammps_extract_atomPvPKc}]{\sphinxcrossref{\sphinxcode{\sphinxupquote{lammps\_extract\_atom()}}}}} and returns a pointer to
LAMMPS data tied to the \sphinxcode{\sphinxupquote{Atom}} class, depending on the data
requested through \sphinxstyleemphasis{name}.

\sphinxAtStartPar
\DUrole{versionmodified,added}{Added in version 3Nov2022.}

\sphinxAtStartPar
Note that this function actually does not return a pointer, but rather
associates the pointer on the left side of the assignment to point
to internal LAMMPS data. Pointers must be of the correct type, kind, and
rank (e.g., \sphinxcode{\sphinxupquote{INTEGER(c\_int), DIMENSION(:)}} for “type”, “mask”, or “id”;
\sphinxcode{\sphinxupquote{INTEGER(c\_int64\_t), DIMENSION(:)}} for “id” if LAMMPS was compiled
with the \sphinxcode{\sphinxupquote{\sphinxhyphen{}DLAMMPS\_BIGBIG}} flag; \sphinxcode{\sphinxupquote{REAL(c\_double), DIMENSION(:,:)}} for
“x”, “v”, or “f”; and so forth). The pointer being associated with LAMMPS
data is type\sphinxhyphen{}, kind\sphinxhyphen{}, and rank\sphinxhyphen{}checked at run\sphinxhyphen{}time.
\begin{quote}\begin{description}
\sphinxlineitem{Parameters}
\sphinxAtStartPar
\sphinxstylestrong{name}\sphinxstyleemphasis{ {[}}\sphinxstyleemphasis{character(len=*)}\sphinxstyleemphasis{{]}} :: string with the name of the property to extract

\sphinxlineitem{Call to}
\sphinxAtStartPar
{\hyperref[\detokenize{Library_atoms:_CPPv419lammps_extract_atomPvPKc}]{\sphinxcrossref{\sphinxcode{\sphinxupquote{lammps\_extract\_atom()}}}}}

\sphinxlineitem{Return}
\sphinxAtStartPar
\sphinxstylestrong{pointer}\sphinxstyleemphasis{ {[}}\sphinxstyleemphasis{polymorphic}\sphinxstyleemphasis{{]}} :: pointer to LAMMPS data. The left\sphinxhyphen{}hand side of the
assignment should be a C\sphinxhyphen{}interoperable pointer of appropriate kind and rank
(e.g., \sphinxcode{\sphinxupquote{INTEGER(c\_int), POINTER :: mask(:)}}) to the extracted
property. If expecting vector data, the pointer should have dimension “:”;
if expecting matrix data, the pointer should have dimension “:,:”.

\end{description}\end{quote}

\begin{sphinxadmonition}{warning}{Warning:}
\sphinxAtStartPar
Pointers returned by this function are generally not persistent, as
per\sphinxhyphen{}atom data may be redistributed, reallocated, and reordered at every
re\sphinxhyphen{}neighboring operation. It is advisable to re\sphinxhyphen{}bind pointers using
\sphinxcode{\sphinxupquote{extract\_atom()}} between runs.
\end{sphinxadmonition}

\begin{sphinxadmonition}{note}{Array index order}

\sphinxAtStartPar
Two\sphinxhyphen{}dimensional arrays returned from \sphinxcode{\sphinxupquote{extract\_atom()}} will be
\sphinxstylestrong{transposed} from equivalent arrays in C, and they will be indexed
from 1 instead of 0. For example, in C,

\begin{sphinxVerbatim}[commandchars=\\\{\}]
\PYG{k+kt}{void}\PYG{+w}{ }\PYG{o}{*}\PYG{n}{lmp}\PYG{p}{;}
\PYG{k+kt}{double}\PYG{+w}{ }\PYG{o}{*}\PYG{o}{*}\PYG{n}{x}\PYG{p}{;}
\PYG{c+cm}{/* more code to setup, etc. */}
\PYG{n}{x}\PYG{+w}{ }\PYG{o}{=}\PYG{+w}{ }\PYG{n}{lammps\PYGZus{}extract\PYGZus{}atom}\PYG{p}{(}\PYG{n}{lmp}\PYG{p}{,}\PYG{+w}{ }\PYG{l+s}{\PYGZdq{}}\PYG{l+s}{x}\PYG{l+s}{\PYGZdq{}}\PYG{p}{)}\PYG{p}{;}
\PYG{n}{printf}\PYG{p}{(}\PYG{l+s}{\PYGZdq{}}\PYG{l+s}{\PYGZpc{}f}\PYG{l+s+se}{\PYGZbs{}n}\PYG{l+s}{\PYGZdq{}}\PYG{p}{,}\PYG{+w}{ }\PYG{n}{x}\PYG{p}{[}\PYG{l+m+mi}{5}\PYG{p}{]}\PYG{p}{[}\PYG{l+m+mi}{1}\PYG{p}{]}\PYG{p}{)}\PYG{p}{;}
\end{sphinxVerbatim}

\sphinxAtStartPar
will print the \sphinxstyleemphasis{y}\sphinxhyphen{}coordinate of the sixth atom on this processor.
Conversely,

\begin{sphinxVerbatim}[commandchars=\\\{\}]
\PYG{k}{TYPE}\PYG{p}{(}\PYG{n}{lammps}\PYG{p}{)}\PYG{+w}{ }\PYG{k+kd}{::}\PYG{+w}{ }\PYG{n}{lmp}
\PYG{k+kt}{REAL}\PYG{p}{(}\PYG{k+kt}{c\PYGZus{}double}\PYG{p}{)}\PYG{p}{,}\PYG{+w}{ }\PYG{k}{DIMENSION}\PYG{p}{(}\PYG{p}{:}\PYG{p}{,}\PYG{p}{:}\PYG{p}{)}\PYG{p}{,}\PYG{+w}{ }\PYG{k}{POINTER}\PYG{+w}{ }\PYG{k+kd}{::}\PYG{+w}{ }\PYG{n}{x}\PYG{+w}{ }\PYG{o}{=}\PYG{o}{\PYGZgt{}}\PYG{+w}{ }\PYG{n+nb}{NULL}\PYG{p}{(}\PYG{p}{)}
\PYG{c}{! more code to setup, etc.}
\PYG{n}{x}\PYG{+w}{ }\PYG{o}{=}\PYG{+w}{ }\PYG{n}{lmp}\PYG{p}{\PYGZpc{}}\PYG{n}{extract\PYGZus{}atom}\PYG{p}{(}\PYG{l+s+s2}{\PYGZdq{}x\PYGZdq{}}\PYG{p}{)}
\PYG{k}{PRINT}\PYG{+w}{ }\PYG{l+s+s1}{\PYGZsq{}(f0.6)\PYGZsq{}}\PYG{p}{,}\PYG{+w}{ }\PYG{n}{x}\PYG{p}{(}\PYG{l+m+mi}{2}\PYG{p}{,}\PYG{l+m+mi}{6}\PYG{p}{)}
\end{sphinxVerbatim}

\sphinxAtStartPar
will print the \sphinxstyleemphasis{y}\sphinxhyphen{}coordinate of the sixth atom on this processor
(note the transposition of the two indices). This is not a choice, but
rather a consequence of the different conventions adopted by the Fortran
and C standards decades ago: in C, the block of data

\begin{sphinxVerbatim}[commandchars=\\\{\}]
1 2 3 4 5 6 7 8 9 10 11 12 13 14 15 16
\end{sphinxVerbatim}

\sphinxAtStartPar
interpreted as a \(4\times4\) matrix would be
\begin{equation*}
\begin{split}\begin{bmatrix}
  1 & 2 & 3 & 4 \\
  5 & 6 & 7 & 8 \\
  9 & 10 & 11 & 12 \\
  13 & 14 & 15 & 16
\end{bmatrix},\end{split}
\end{equation*}
\sphinxAtStartPar
that is, in row\sphinxhyphen{}major order. In Fortran, the same block of data is
interpreted in column\sphinxhyphen{}major order, namely,
\begin{equation*}
\begin{split}\begin{bmatrix}
  1 & 5 & 9  & 13 \\
  2 & 6 & 10 & 14 \\
  3 & 7 & 11 & 15 \\
  4 & 8 & 12 & 16
\end{bmatrix}.\end{split}
\end{equation*}
\sphinxAtStartPar
This difference in interpretation of the same block of data by the two
languages means, in effect, that matrices from C or C++ will be
transposed when interpreted in Fortran.
\end{sphinxadmonition}

\begin{sphinxadmonition}{note}{Note:}
\sphinxAtStartPar
If you would like the indices to start at 0 instead of 1 (which follows
typical notation in C and C++, but not Fortran), you can create another
pointer and associate it thus:

\begin{sphinxVerbatim}[commandchars=\\\{\}]
\PYG{k+kt}{REAL}\PYG{p}{(}\PYG{k+kt}{c\PYGZus{}double}\PYG{p}{)}\PYG{p}{,}\PYG{+w}{ }\PYG{k}{DIMENSION}\PYG{p}{(}\PYG{p}{:}\PYG{p}{,}\PYG{p}{:}\PYG{p}{)}\PYG{p}{,}\PYG{+w}{ }\PYG{k}{POINTER}\PYG{+w}{ }\PYG{k+kd}{::}\PYG{+w}{ }\PYG{n}{x}\PYG{p}{,}\PYG{+w}{ }\PYG{n}{x0}
\PYG{n}{x}\PYG{+w}{ }\PYG{o}{=}\PYG{+w}{ }\PYG{n}{lmp}\PYG{p}{\PYGZpc{}}\PYG{n}{extract\PYGZus{}atom}\PYG{p}{(}\PYG{l+s+s2}{\PYGZdq{}x\PYGZdq{}}\PYG{p}{)}
\PYG{n}{x0}\PYG{p}{(}\PYG{l+m+mi}{0}\PYG{p}{:}\PYG{p}{,}\PYG{l+m+mi}{0}\PYG{p}{:}\PYG{p}{)}\PYG{+w}{ }\PYG{o}{=}\PYG{o}{\PYGZgt{}}\PYG{+w}{ }\PYG{n}{x}
\end{sphinxVerbatim}

\sphinxAtStartPar
The above would cause the dimensions of \sphinxstyleemphasis{x} to be (1:3, 1:nmax)
and those of \sphinxstyleemphasis{x0} to be (0:2, 0:nmax\(-\)1).
\end{sphinxadmonition}

\end{fulllineitems}



\bigskip\hrule\bigskip

\index{extract\_compute() (fortran function)@\spxentry{extract\_compute()}\spxextra{fortran function}|spxpagem}

\begin{fulllineitems}
\phantomsection\label{\detokenize{Fortran:f/_/extract_compute}}
\pysigstartsignatures
\pysiglinewithargsret{\sphinxbfcode{\sphinxupquote{function  }}\sphinxbfcode{\sphinxupquote{extract\_compute}}}{\sphinxparam{id}\sphinxparamcomma \sphinxparam{style}\sphinxparamcomma \sphinxparam{type}}{}
\pysigstopsignatures
\sphinxAtStartPar
This function calls {\hyperref[\detokenize{Library_objects:_CPPv422lammps_extract_computePvPKcii}]{\sphinxcrossref{\sphinxcode{\sphinxupquote{lammps\_extract\_compute()}}}}} and returns a pointer
to LAMMPS data tied to the \sphinxcode{\sphinxupquote{Compute}} class, specifically data
provided by the compute identified by \sphinxstyleemphasis{id}. Computes may provide global,
per\sphinxhyphen{}atom, or local data, and those data may be a scalar, a vector, or an
array. Since computes may provide multiple kinds of data, the user is
required to specify which set of data is to be returned through the
\sphinxstyleemphasis{style} and \sphinxstyleemphasis{type} variables.

\sphinxAtStartPar
\DUrole{versionmodified,added}{Added in version 3Nov2022.}

\sphinxAtStartPar
Note that this function actually does not return a value, but rather
associates the pointer on the left side of the assignment to point to
internal LAMMPS data. Pointers must be of the correct data type to point to
said data (i.e., \sphinxcode{\sphinxupquote{REAL(c\_double)}}) and have compatible rank.  The pointer
being associated with LAMMPS data is type\sphinxhyphen{}, kind\sphinxhyphen{}, and rank\sphinxhyphen{}checked at
run\sphinxhyphen{}time via an overloaded assignment operator.

\sphinxAtStartPar
For example,

\begin{sphinxVerbatim}[commandchars=\\\{\}]
\PYG{k}{TYPE}\PYG{p}{(}\PYG{n}{lammps}\PYG{p}{)}\PYG{+w}{ }\PYG{k+kd}{::}\PYG{+w}{ }\PYG{n}{lmp}
\PYG{k+kt}{REAL}\PYG{p}{(}\PYG{k+kt}{c\PYGZus{}double}\PYG{p}{)}\PYG{p}{,}\PYG{+w}{ }\PYG{k}{DIMENSION}\PYG{p}{(}\PYG{p}{:}\PYG{p}{)}\PYG{p}{,}\PYG{+w}{ }\PYG{k}{POINTER}\PYG{+w}{ }\PYG{k+kd}{::}\PYG{+w}{ }\PYG{n}{COM}
\PYG{c}{! code to setup, create atoms, etc.}
\PYG{k}{CALL }\PYG{n}{lmp}\PYG{p}{\PYGZpc{}}\PYG{n}{command}\PYG{p}{(}\PYG{l+s+s1}{\PYGZsq{}compute COM all com\PYGZsq{}}\PYG{p}{)}
\PYG{n}{COM}\PYG{+w}{ }\PYG{o}{=}\PYG{+w}{ }\PYG{n}{lmp}\PYG{p}{\PYGZpc{}}\PYG{n}{extract\PYGZus{}compute}\PYG{p}{(}\PYG{l+s+s1}{\PYGZsq{}COM\PYGZsq{}}\PYG{p}{,}\PYG{+w}{ }\PYG{n}{lmp}\PYG{p}{\PYGZpc{}}\PYG{n}{style}\PYG{p}{\PYGZpc{}}\PYG{n}{global}\PYG{p}{,}\PYG{+w}{ }\PYG{n}{lmp}\PYG{p}{\PYGZpc{}}\PYG{n}{style}\PYG{p}{\PYGZpc{}}\PYG{k}{type}\PYG{p}{)}
\end{sphinxVerbatim}

\sphinxAtStartPar
will bind the variable \sphinxstyleemphasis{COM} to the center of mass of the atoms created in
your simulation. The vector in this case has length 3; the length (or, in
the case of array data, the number of rows and columns) is determined for
you based on data from the \sphinxcode{\sphinxupquote{Compute}} class.

\begin{sphinxadmonition}{note}{Array index order}

\sphinxAtStartPar
Two\sphinxhyphen{}dimensional arrays returned from \sphinxcode{\sphinxupquote{extract\_compute()}} will be
\sphinxstylestrong{transposed} from equivalent arrays in C, and they will be indexed
from 1 instead of 0. See the note at \sphinxcode{\sphinxupquote{extract\_atom()}} for
further details.
\end{sphinxadmonition}

\sphinxAtStartPar
The following combinations are possible (assuming \sphinxcode{\sphinxupquote{lmp}} is the name of
your LAMMPS instance):


\begin{savenotes}\sphinxattablestart
\sphinxthistablewithglobalstyle
\centering
\begin{tabular}[t]{\X{21}{100}\X{20}{100}\X{40}{100}\X{19}{100}}
\sphinxtoprule
\sphinxstyletheadfamily 
\sphinxAtStartPar
Style
&\sphinxstyletheadfamily 
\sphinxAtStartPar
Type
&\sphinxstyletheadfamily 
\sphinxAtStartPar
Type to assign to
&\sphinxstyletheadfamily 
\sphinxAtStartPar
Returned data
\\
\sphinxmidrule
\sphinxtableatstartofbodyhook
\sphinxAtStartPar
\sphinxcode{\sphinxupquote{lmp\%style\%global}}
&
\sphinxAtStartPar
\sphinxcode{\sphinxupquote{lmp\%type\%scalar}}
&
\sphinxAtStartPar
\sphinxcode{\sphinxupquote{REAL(c\_double), POINTER}}
&
\sphinxAtStartPar
Global scalar
\\
\sphinxhline
\sphinxAtStartPar
\sphinxcode{\sphinxupquote{lmp\%style\%global}}
&
\sphinxAtStartPar
\sphinxcode{\sphinxupquote{lmp\%type\%vector}}
&
\sphinxAtStartPar
\sphinxcode{\sphinxupquote{REAL(c\_double), DIMENSION(:), POINTER}}
&
\sphinxAtStartPar
Global vector
\\
\sphinxhline
\sphinxAtStartPar
\sphinxcode{\sphinxupquote{lmp\%style\%global}}
&
\sphinxAtStartPar
\sphinxcode{\sphinxupquote{lmp\%type\%array}}
&
\sphinxAtStartPar
\sphinxcode{\sphinxupquote{REAL(c\_double), DIMENSION(:,:), POINTER}}
&
\sphinxAtStartPar
Global array
\\
\sphinxhline
\sphinxAtStartPar
\sphinxcode{\sphinxupquote{lmp\%style\%atom}}
&
\sphinxAtStartPar
\sphinxcode{\sphinxupquote{lmp\%type\%vector}}
&
\sphinxAtStartPar
\sphinxcode{\sphinxupquote{REAL(c\_double), DIMENSION(:), POINTER}}
&
\sphinxAtStartPar
Per\sphinxhyphen{}atom vector
\\
\sphinxhline
\sphinxAtStartPar
\sphinxcode{\sphinxupquote{lmp\%style\%atom}}
&
\sphinxAtStartPar
\sphinxcode{\sphinxupquote{lmp\%type\%array}}
&
\sphinxAtStartPar
\sphinxcode{\sphinxupquote{REAL(c\_double), DIMENSION(:,:), POINTER}}
&
\sphinxAtStartPar
Per\sphinxhyphen{}atom array
\\
\sphinxhline
\sphinxAtStartPar
\sphinxcode{\sphinxupquote{lmp\%style\%local}}
&
\sphinxAtStartPar
\sphinxcode{\sphinxupquote{lmp\%type\%vector}}
&
\sphinxAtStartPar
\sphinxcode{\sphinxupquote{REAL(c\_double), DIMENSION(:), POINTER}}
&
\sphinxAtStartPar
Local vector
\\
\sphinxhline
\sphinxAtStartPar
\sphinxcode{\sphinxupquote{lmp\%style\%local}}
&
\sphinxAtStartPar
\sphinxcode{\sphinxupquote{lmp\%type\%array}}
&
\sphinxAtStartPar
\sphinxcode{\sphinxupquote{REAL(c\_double), DIMENSION(:,:), POINTER}}
&
\sphinxAtStartPar
Local array
\\
\sphinxbottomrule
\end{tabular}
\sphinxtableafterendhook\par
\sphinxattableend\end{savenotes}
\begin{quote}\begin{description}
\sphinxlineitem{Parameters}\begin{itemize}
\item {} 
\sphinxAtStartPar
\sphinxstylestrong{id}\sphinxstyleemphasis{ {[}}\sphinxstyleemphasis{character(len=*)}\sphinxstyleemphasis{{]}} :: compute ID from which to extract data

\item {} 
\sphinxAtStartPar
\sphinxstylestrong{style}\sphinxstyleemphasis{ {[}}\sphinxstyleemphasis{integer(c\_int)}\sphinxstyleemphasis{{]}} :: value indicating the style of data to extract
(global, per\sphinxhyphen{}atom, or local)

\item {} 
\sphinxAtStartPar
\sphinxstylestrong{type}\sphinxstyleemphasis{ {[}}\sphinxstyleemphasis{integer(c\_int)}\sphinxstyleemphasis{{]}} :: value indicating the type of data to extract
(scalar, vector, or array)

\end{itemize}

\sphinxlineitem{Call to}
\sphinxAtStartPar
{\hyperref[\detokenize{Library_objects:_CPPv422lammps_extract_computePvPKcii}]{\sphinxcrossref{\sphinxcode{\sphinxupquote{lammps\_extract\_compute()}}}}}

\sphinxlineitem{Return}
\sphinxAtStartPar
\sphinxstylestrong{pointer}\sphinxstyleemphasis{ {[}}\sphinxstyleemphasis{polymorphic}\sphinxstyleemphasis{{]}} :: pointer to LAMMPS data. The left\sphinxhyphen{}hand side of the assignment
should be a C\sphinxhyphen{}compatible pointer (e.g., \sphinxcode{\sphinxupquote{REAL(c\_double), POINTER :: x}})
to the extracted property. If expecting vector data, the pointer should
have dimension “:”; if expecting array (matrix) data, the pointer should
have dimension “:,:”.

\end{description}\end{quote}

\begin{sphinxadmonition}{note}{Note:}
\sphinxAtStartPar
If the compute’s data are not already computed for the current step, the
compute will be invoked. LAMMPS cannot easily check at that time if it is
valid to invoke a compute, so it may fail with an error. The caller has
to check to avoid such an error.
\end{sphinxadmonition}

\begin{sphinxadmonition}{warning}{Warning:}
\sphinxAtStartPar
The pointers returned by this function are generally not persistent,
since the computed data may be re\sphinxhyphen{}distributed, re\sphinxhyphen{}allocated, and
re\sphinxhyphen{}ordered at every invocation. It is advisable to re\sphinxhyphen{}invoke this
function before the data are accessed or make a copy if the data are to
be used after other LAMMPS commands have been issued. Do \sphinxstylestrong{not} modify
the data returned by this function.
\end{sphinxadmonition}

\end{fulllineitems}



\bigskip\hrule\bigskip

\index{extract\_fix() (fortran function)@\spxentry{extract\_fix()}\spxextra{fortran function}|spxpagem}

\begin{fulllineitems}
\phantomsection\label{\detokenize{Fortran:f/_/extract_fix}}
\pysigstartsignatures
\pysiglinewithargsret{\sphinxbfcode{\sphinxupquote{function  }}\sphinxbfcode{\sphinxupquote{extract\_fix}}}{\sphinxparam{id, style, type{[}, nrow{]}{[}, ncol{]}}}{}
\pysigstopsignatures
\sphinxAtStartPar
This function calls {\hyperref[\detokenize{Library_objects:_CPPv418lammps_extract_fixPvPKciiii}]{\sphinxcrossref{\sphinxcode{\sphinxupquote{lammps\_extract\_fix()}}}}} and returns a pointer to
LAMMPS data tied to the \sphinxcode{\sphinxupquote{Fix}} class, specifically data provided
by the fix identified by \sphinxstyleemphasis{id}. Fixes may provide global, per\sphinxhyphen{}atom, or
local data, and those data may be a scalar, a vector, or an array. Since
many fixes provide multiple kinds of data, the user is required to specify
which set of data is to be returned through the \sphinxstyleemphasis{style} and \sphinxstyleemphasis{type}
variables.

\sphinxAtStartPar
\DUrole{versionmodified,added}{Added in version 3Nov2022.}

\sphinxAtStartPar
Global data are calculated at the time they are requested and are only
available element\sphinxhyphen{}by\sphinxhyphen{}element. As such, the user is expected to provide
the \sphinxstyleemphasis{nrow} variable to specify which element of a global vector or the
\sphinxstyleemphasis{nrow} and \sphinxstyleemphasis{ncol} variables to specify which element of a global array the
user wishes LAMMPS to return. The \sphinxstyleemphasis{ncol} variable is optional for global
scalar or vector data, and both \sphinxstyleemphasis{nrow} and \sphinxstyleemphasis{ncol} are optional when a
global scalar is requested, as well as when per\sphinxhyphen{}atom or local data are
requested. The following combinations are possible (assuming \sphinxcode{\sphinxupquote{lmp}} is the
name of your LAMMPS instance):


\begin{savenotes}\sphinxattablestart
\sphinxthistablewithglobalstyle
\centering
\begin{tabular}[t]{\X{20}{100}\X{19}{100}\X{11}{100}\X{11}{100}\X{21}{100}\X{18}{100}}
\sphinxtoprule
\sphinxstyletheadfamily 
\sphinxAtStartPar
Style
&\sphinxstyletheadfamily 
\sphinxAtStartPar
Type
&\sphinxstyletheadfamily 
\sphinxAtStartPar
nrow
&\sphinxstyletheadfamily 
\sphinxAtStartPar
ncol
&\sphinxstyletheadfamily 
\sphinxAtStartPar
Type to assign to
&\sphinxstyletheadfamily 
\sphinxAtStartPar
Returned data
\\
\sphinxmidrule
\sphinxtableatstartofbodyhook
\sphinxAtStartPar
\sphinxcode{\sphinxupquote{lmp\%style\%global}}
&
\sphinxAtStartPar
\sphinxcode{\sphinxupquote{lmp\%type\%scalar}}
&
\sphinxAtStartPar
Ignored
&
\sphinxAtStartPar
Ignored
&
\sphinxAtStartPar
\sphinxcode{\sphinxupquote{REAL(c\_double)}}
&
\sphinxAtStartPar
Global scalar
\\
\sphinxhline
\sphinxAtStartPar
\sphinxcode{\sphinxupquote{lmp\%style\%global}}
&
\sphinxAtStartPar
\sphinxcode{\sphinxupquote{lmp\%type\%vector}}
&
\sphinxAtStartPar
Required
&
\sphinxAtStartPar
Ignored
&
\sphinxAtStartPar
\sphinxcode{\sphinxupquote{REAL(c\_double)}}
&
\sphinxAtStartPar
Element of global vector
\\
\sphinxhline
\sphinxAtStartPar
\sphinxcode{\sphinxupquote{lmp\%style\%global}}
&
\sphinxAtStartPar
\sphinxcode{\sphinxupquote{lmp\%type\%array}}
&
\sphinxAtStartPar
Required
&
\sphinxAtStartPar
Required
&
\sphinxAtStartPar
\sphinxcode{\sphinxupquote{REAL(c\_double)}}
&
\sphinxAtStartPar
Element of global array
\\
\sphinxhline
\sphinxAtStartPar
\sphinxcode{\sphinxupquote{lmp\%style\%atom}}
&
\sphinxAtStartPar
\sphinxcode{\sphinxupquote{lmp\%type\%scalar}}
&&&&
\sphinxAtStartPar
(not allowed)
\\
\sphinxhline
\sphinxAtStartPar
\sphinxcode{\sphinxupquote{lmp\%style\%atom}}
&
\sphinxAtStartPar
\sphinxcode{\sphinxupquote{lmp\%type\%vector}}
&
\sphinxAtStartPar
Ignored
&
\sphinxAtStartPar
Ignored
&
\sphinxAtStartPar
\sphinxcode{\sphinxupquote{REAL(c\_double), DIMENSION(:), POINTER}}
&
\sphinxAtStartPar
Per\sphinxhyphen{}atom vector
\\
\sphinxhline
\sphinxAtStartPar
\sphinxcode{\sphinxupquote{lmp\%style\%atom}}
&
\sphinxAtStartPar
\sphinxcode{\sphinxupquote{lmp\%type\%array}}
&
\sphinxAtStartPar
Ignored
&
\sphinxAtStartPar
Ignored
&
\sphinxAtStartPar
\sphinxcode{\sphinxupquote{REAL(c\_double), DIMENSION(:,:), POINTER}}
&
\sphinxAtStartPar
Per\sphinxhyphen{}atom array
\\
\sphinxhline
\sphinxAtStartPar
\sphinxcode{\sphinxupquote{lmp\%style\%local}}
&
\sphinxAtStartPar
\sphinxcode{\sphinxupquote{lmp\%type\%scalar}}
&&&&
\sphinxAtStartPar
(not allowed)
\\
\sphinxhline
\sphinxAtStartPar
\sphinxcode{\sphinxupquote{lmp\%style\%local}}
&
\sphinxAtStartPar
\sphinxcode{\sphinxupquote{lmp\%type\%vector}}
&
\sphinxAtStartPar
Ignored
&
\sphinxAtStartPar
Ignored
&
\sphinxAtStartPar
\sphinxcode{\sphinxupquote{REAL(c\_double), DIMENSION(:), POINTER}}
&
\sphinxAtStartPar
Per\sphinxhyphen{}atom vector
\\
\sphinxhline
\sphinxAtStartPar
\sphinxcode{\sphinxupquote{lmp\%style\%local}}
&
\sphinxAtStartPar
\sphinxcode{\sphinxupquote{lmp\%type\%array}}
&
\sphinxAtStartPar
Ignored
&
\sphinxAtStartPar
Ignored
&
\sphinxAtStartPar
\sphinxcode{\sphinxupquote{REAL(c\_double), DIMENSION(:,:), POINTER}}
&
\sphinxAtStartPar
Per\sphinxhyphen{}atom array
\\
\sphinxbottomrule
\end{tabular}
\sphinxtableafterendhook\par
\sphinxattableend\end{savenotes}

\sphinxAtStartPar
In the case of global data, this function returns a value of type
\sphinxcode{\sphinxupquote{REAL(c\_double)}}. For per\sphinxhyphen{}atom or local data, this function does not
return a value but instead associates the pointer on the left side of the
assignment to point to internal LAMMPS data. Pointers must be of the correct
type and kind to point to said data (i.e., \sphinxcode{\sphinxupquote{REAL(c\_double)}}) and have
compatible rank.  The pointer being associated with LAMMPS data is type\sphinxhyphen{},
kind\sphinxhyphen{}, and rank\sphinxhyphen{}checked at run\sphinxhyphen{}time via an overloaded assignment operator.

\sphinxAtStartPar
For example,

\begin{sphinxVerbatim}[commandchars=\\\{\}]
\PYG{k}{TYPE}\PYG{p}{(}\PYG{n}{lammps}\PYG{p}{)}\PYG{+w}{ }\PYG{k+kd}{::}\PYG{+w}{ }\PYG{n}{lmp}
\PYG{k+kt}{REAL}\PYG{p}{(}\PYG{k+kt}{c\PYGZus{}double}\PYG{p}{)}\PYG{+w}{ }\PYG{k+kd}{::}\PYG{+w}{ }\PYG{n}{dr}\PYG{p}{,}\PYG{+w}{ }\PYG{n}{dx}\PYG{p}{,}\PYG{+w}{ }\PYG{n}{dy}\PYG{p}{,}\PYG{+w}{ }\PYG{n}{dz}
\PYG{c}{! more code to set up, etc.}
\PYG{n}{lmp}\PYG{p}{\PYGZpc{}}\PYG{n}{command}\PYG{p}{(}\PYG{l+s+s1}{\PYGZsq{}fix george all recenter 2 2 2\PYGZsq{}}\PYG{p}{)}
\PYG{c}{! more code}
\PYG{n}{dr}\PYG{+w}{ }\PYG{o}{=}\PYG{+w}{ }\PYG{n}{lmp}\PYG{p}{\PYGZpc{}}\PYG{n}{extract\PYGZus{}fix}\PYG{p}{(}\PYG{l+s+s2}{\PYGZdq{}george\PYGZdq{}}\PYG{p}{,}\PYG{+w}{ }\PYG{n}{lmp}\PYG{p}{\PYGZpc{}}\PYG{n}{style}\PYG{p}{\PYGZpc{}}\PYG{n}{global}\PYG{p}{,}\PYG{+w}{ }\PYG{n}{lmp}\PYG{p}{\PYGZpc{}}\PYG{n}{style}\PYG{p}{\PYGZpc{}}\PYG{n}{scalar}\PYG{p}{)}
\PYG{n}{dx}\PYG{+w}{ }\PYG{o}{=}\PYG{+w}{ }\PYG{n}{lmp}\PYG{p}{\PYGZpc{}}\PYG{n}{extract\PYGZus{}fix}\PYG{p}{(}\PYG{l+s+s2}{\PYGZdq{}george\PYGZdq{}}\PYG{p}{,}\PYG{+w}{ }\PYG{n}{lmp}\PYG{p}{\PYGZpc{}}\PYG{n}{style}\PYG{p}{\PYGZpc{}}\PYG{n}{global}\PYG{p}{,}\PYG{+w}{ }\PYG{n}{lmp}\PYG{p}{\PYGZpc{}}\PYG{n}{style}\PYG{p}{\PYGZpc{}}\PYG{n}{vector}\PYG{p}{,}\PYG{+w}{ }\PYG{l+m+mi}{1}\PYG{p}{)}
\PYG{n}{dy}\PYG{+w}{ }\PYG{o}{=}\PYG{+w}{ }\PYG{n}{lmp}\PYG{p}{\PYGZpc{}}\PYG{n}{extract\PYGZus{}fix}\PYG{p}{(}\PYG{l+s+s2}{\PYGZdq{}george\PYGZdq{}}\PYG{p}{,}\PYG{+w}{ }\PYG{n}{lmp}\PYG{p}{\PYGZpc{}}\PYG{n}{style}\PYG{p}{\PYGZpc{}}\PYG{n}{global}\PYG{p}{,}\PYG{+w}{ }\PYG{n}{lmp}\PYG{p}{\PYGZpc{}}\PYG{n}{style}\PYG{p}{\PYGZpc{}}\PYG{n}{vector}\PYG{p}{,}\PYG{+w}{ }\PYG{l+m+mi}{2}\PYG{p}{)}
\PYG{n}{dz}\PYG{+w}{ }\PYG{o}{=}\PYG{+w}{ }\PYG{n}{lmp}\PYG{p}{\PYGZpc{}}\PYG{n}{extract\PYGZus{}fix}\PYG{p}{(}\PYG{l+s+s2}{\PYGZdq{}george\PYGZdq{}}\PYG{p}{,}\PYG{+w}{ }\PYG{n}{lmp}\PYG{p}{\PYGZpc{}}\PYG{n}{style}\PYG{p}{\PYGZpc{}}\PYG{n}{global}\PYG{p}{,}\PYG{+w}{ }\PYG{n}{lmp}\PYG{p}{\PYGZpc{}}\PYG{n}{style}\PYG{p}{\PYGZpc{}}\PYG{n}{vector}\PYG{p}{,}\PYG{+w}{ }\PYG{l+m+mi}{3}\PYG{p}{)}
\end{sphinxVerbatim}

\sphinxAtStartPar
will extract the global scalar calculated by
{\hyperref[\detokenize{fix_recenter::doc}]{\sphinxcrossref{\DUrole{doc}{fix recenter}}}} into the variable \sphinxstyleemphasis{dr} and the
three elements of the global vector calculated by fix recenter into the
variables \sphinxstyleemphasis{dx}, \sphinxstyleemphasis{dy}, and \sphinxstyleemphasis{dz}, respectively.

\sphinxAtStartPar
If asked for per\sphinxhyphen{}atom or local data, \sphinxcode{\sphinxupquote{extract\_fix()}} returns a
pointer to actual LAMMPS data.  The pointer returned will have the
appropriate size to match the internal data, and will be
type/kind/rank\sphinxhyphen{}checked at the time of the assignment. For example,

\begin{sphinxVerbatim}[commandchars=\\\{\}]
\PYG{k}{TYPE}\PYG{p}{(}\PYG{n}{lammps}\PYG{p}{)}\PYG{+w}{ }\PYG{k+kd}{::}\PYG{+w}{ }\PYG{n}{lmp}
\PYG{k+kt}{REAL}\PYG{p}{(}\PYG{k+kt}{c\PYGZus{}double}\PYG{p}{)}\PYG{p}{,}\PYG{+w}{ }\PYG{k}{DIMENSION}\PYG{p}{(}\PYG{p}{:}\PYG{p}{)}\PYG{p}{,}\PYG{+w}{ }\PYG{k}{POINTER}\PYG{+w}{ }\PYG{k+kd}{::}\PYG{+w}{ }\PYG{n}{r}
\PYG{c}{! more code to set up, etc.}
\PYG{n}{lmp}\PYG{p}{\PYGZpc{}}\PYG{n}{command}\PYG{p}{(}\PYG{l+s+s1}{\PYGZsq{}fix state all store/state 0 x y z\PYGZsq{}}\PYG{p}{)}
\PYG{c}{! more code}
\PYG{n}{r}\PYG{+w}{ }\PYG{o}{=}\PYG{+w}{ }\PYG{n}{lmp}\PYG{p}{\PYGZpc{}}\PYG{n}{extract\PYGZus{}fix}\PYG{p}{(}\PYG{l+s+s1}{\PYGZsq{}state\PYGZsq{}}\PYG{p}{,}\PYG{+w}{ }\PYG{n}{lmp}\PYG{p}{\PYGZpc{}}\PYG{n}{style}\PYG{p}{\PYGZpc{}}\PYG{n}{atom}\PYG{p}{,}\PYG{+w}{ }\PYG{n}{lmp}\PYG{p}{\PYGZpc{}}\PYG{k}{type}\PYG{p}{\PYGZpc{}}\PYG{k}{array}\PYG{p}{)}
\end{sphinxVerbatim}

\sphinxAtStartPar
will bind the pointer \sphinxstyleemphasis{r} to internal LAMMPS data representing the per\sphinxhyphen{}atom
array computed by {\hyperref[\detokenize{fix_store_state::doc}]{\sphinxcrossref{\DUrole{doc}{fix store/state}}}} when three
inputs are specified. Similarly,

\begin{sphinxVerbatim}[commandchars=\\\{\}]
\PYG{k}{TYPE}\PYG{p}{(}\PYG{n}{lammps}\PYG{p}{)}\PYG{+w}{ }\PYG{k+kd}{::}\PYG{+w}{ }\PYG{n}{lmp}
\PYG{k+kt}{REAL}\PYG{p}{(}\PYG{k+kt}{c\PYGZus{}double}\PYG{p}{)}\PYG{p}{,}\PYG{+w}{ }\PYG{k}{DIMENSION}\PYG{p}{(}\PYG{p}{:}\PYG{p}{)}\PYG{p}{,}\PYG{+w}{ }\PYG{k}{POINTER}\PYG{+w}{ }\PYG{k+kd}{::}\PYG{+w}{ }\PYG{n}{x}
\PYG{c}{! more code to set up, etc.}
\PYG{n}{lmp}\PYG{p}{\PYGZpc{}}\PYG{n}{command}\PYG{p}{(}\PYG{l+s+s1}{\PYGZsq{}fix state all store/state 0 x\PYGZsq{}}\PYG{p}{)}
\PYG{c}{! more code}
\PYG{n}{x}\PYG{+w}{ }\PYG{o}{=}\PYG{+w}{ }\PYG{n}{lmp}\PYG{p}{\PYGZpc{}}\PYG{n}{extract\PYGZus{}fix}\PYG{p}{(}\PYG{l+s+s1}{\PYGZsq{}state\PYGZsq{}}\PYG{p}{,}\PYG{+w}{ }\PYG{n}{lmp}\PYG{p}{\PYGZpc{}}\PYG{n}{style}\PYG{p}{\PYGZpc{}}\PYG{n}{atom}\PYG{p}{,}\PYG{+w}{ }\PYG{n}{lmp}\PYG{p}{\PYGZpc{}}\PYG{k}{type}\PYG{p}{\PYGZpc{}}\PYG{n}{vector}\PYG{p}{)}
\end{sphinxVerbatim}

\sphinxAtStartPar
will associate the pointer \sphinxstyleemphasis{x} with internal LAMMPS data corresponding to
the per\sphinxhyphen{}atom vector computed by {\hyperref[\detokenize{fix_store_state::doc}]{\sphinxcrossref{\DUrole{doc}{fix store/state}}}}
when only one input is specified. Similar examples with \sphinxcode{\sphinxupquote{lmp\%style\%atom}}
replaced by \sphinxcode{\sphinxupquote{lmp\%style\%local}} will extract local data from fixes that
define local vectors and/or arrays.

\begin{sphinxadmonition}{warning}{Warning:}
\sphinxAtStartPar
The pointers returned by this function for per\sphinxhyphen{}atom or local data are
generally not persistent, since the computed data may be redistributed,
reallocated, and reordered at every invocation of the fix.  It is thus
advisable to re\sphinxhyphen{}invoke this function before the data are accessed or to
make a copy if the data are to be used after other LAMMPS commands have
been issued.
\end{sphinxadmonition}

\begin{sphinxadmonition}{note}{Note:}
\sphinxAtStartPar
LAMMPS cannot easily check if it is valid to access the data, so it
may fail with an error.  The caller has to avoid such an error.
\end{sphinxadmonition}
\begin{quote}\begin{description}
\sphinxlineitem{Parameters}\begin{itemize}
\item {} 
\sphinxAtStartPar
\sphinxstylestrong{id}\sphinxstyleemphasis{ {[}}\sphinxstyleemphasis{character(len=*)}\sphinxstyleemphasis{{]}} :: string with the name of the fix from which
to extract data

\item {} 
\sphinxAtStartPar
\sphinxstylestrong{style}\sphinxstyleemphasis{ {[}}\sphinxstyleemphasis{integer(c\_int)}\sphinxstyleemphasis{{]}} :: value indicating the style of data to extract
(global, per\sphinxhyphen{}atom, or local)

\item {} 
\sphinxAtStartPar
\sphinxstylestrong{type}\sphinxstyleemphasis{ {[}}\sphinxstyleemphasis{integer(c\_int)}\sphinxstyleemphasis{{]}} :: value indicating the type of data to extract
(scalar, vector, or array)

\item {} 
\sphinxAtStartPar
\sphinxstylestrong{nrow}\sphinxstyleemphasis{ {[}}\sphinxstyleemphasis{integer(c\_int)}\sphinxstyleemphasis{{]}} :: row index (used only for global vectors and arrays)

\item {} 
\sphinxAtStartPar
\sphinxstylestrong{ncol}\sphinxstyleemphasis{ {[}}\sphinxstyleemphasis{integer(c\_int)}\sphinxstyleemphasis{{]}} :: column index (only used for global arrays)

\end{itemize}

\sphinxlineitem{Call to}
\sphinxAtStartPar
{\hyperref[\detokenize{Library_objects:_CPPv418lammps_extract_fixPvPKciiii}]{\sphinxcrossref{\sphinxcode{\sphinxupquote{lammps\_extract\_fix()}}}}}

\sphinxlineitem{Return}
\sphinxAtStartPar
\sphinxstylestrong{data}\sphinxstyleemphasis{ {[}}\sphinxstyleemphasis{polymorphic}\sphinxstyleemphasis{{]}} :: LAMMPS data (for global data) or a pointer to LAMMPS data
(for per\sphinxhyphen{}atom or local data). The left\sphinxhyphen{}hand side of the assignment should
be of type \sphinxcode{\sphinxupquote{REAL(c\_double)}} and have appropriate rank (i.e.,
\sphinxcode{\sphinxupquote{DIMENSION(:)}} if expecting per\sphinxhyphen{}atom or local vector data and
\sphinxcode{\sphinxupquote{DIMENSION(:,:)}} if expecting per\sphinxhyphen{}atom or local array data). If expecting
local or per\sphinxhyphen{}atom data, it should have the \sphinxcode{\sphinxupquote{POINTER}} attribute, but
if expecting global data, it should be an ordinary (non\sphinxhyphen{}\sphinxcode{\sphinxupquote{POINTER}})
variable.

\end{description}\end{quote}

\begin{sphinxadmonition}{note}{Array index order}

\sphinxAtStartPar
Two\sphinxhyphen{}dimensional global, per\sphinxhyphen{}atom, or local array data from
\sphinxcode{\sphinxupquote{extract\_fix()}} will be \sphinxstylestrong{transposed} from equivalent arrays in
C (or in the ordinary LAMMPS interface accessed through thermodynamic
output), and they will be indexed from 1, not 0. This is true even for
global data, which are returned as scalars—this is done primarily so
the interface is consistent, as there is no choice but to transpose the
indices for per\sphinxhyphen{}atom or local array data. See the similar note under
\sphinxcode{\sphinxupquote{extract\_atom()}} for further details.
\end{sphinxadmonition}

\end{fulllineitems}



\bigskip\hrule\bigskip

\index{extract\_variable() (fortran function)@\spxentry{extract\_variable()}\spxextra{fortran function}|spxpagem}

\begin{fulllineitems}
\phantomsection\label{\detokenize{Fortran:f/_/extract_variable}}
\pysigstartsignatures
\pysiglinewithargsret{\sphinxbfcode{\sphinxupquote{function  }}\sphinxbfcode{\sphinxupquote{extract\_variable}}}{\sphinxparam{name}\sphinxoptional{\sphinxparamcomma \sphinxparam{group}}}{}
\pysigstopsignatures
\sphinxAtStartPar
This function calls {\hyperref[\detokenize{Library_objects:_CPPv423lammps_extract_variablePvPKcPKc}]{\sphinxcrossref{\sphinxcode{\sphinxupquote{lammps\_extract\_variable()}}}}} and returns a
scalar, vector, or string containing the value of the variable identified by
\sphinxstyleemphasis{name}. When the variable is an \sphinxstyleemphasis{equal}\sphinxhyphen{}style variable (or one compatible
with that style such as \sphinxstyleemphasis{internal}), the variable is evaluated and the
corresponding value returned. When the variable is an \sphinxstyleemphasis{atom}\sphinxhyphen{}style variable,
the variable is evaluated and a vector of values is returned. With all
other variables, a string is returned. The \sphinxstyleemphasis{group} argument is only used
for \sphinxstyleemphasis{atom} style variables and is ignored otherwise. If \sphinxstyleemphasis{group} is absent
for \sphinxstyleemphasis{atom}\sphinxhyphen{}style variables, the group is assumed to be “all”.

\sphinxAtStartPar
\DUrole{versionmodified,added}{Added in version 3Nov2022.}

\sphinxAtStartPar
This function returns the values of the variables, not pointers to them.
Vectors pointing to \sphinxstyleemphasis{atom}\sphinxhyphen{}style variables should be of type
\sphinxcode{\sphinxupquote{REAL(c\_double)}}, be of rank 1 (i.e., \sphinxcode{\sphinxupquote{DIMENSION(:)}}), and have the
\sphinxcode{\sphinxupquote{ALLOCATABLE}} attribute.

\begin{sphinxadmonition}{note}{Note:}
\sphinxAtStartPar
Unlike the C library interface, the Fortran interface does not require
you to deallocate memory when you are through; this is done for you,
behind the scenes.
\end{sphinxadmonition}

\sphinxAtStartPar
For example,

\begin{sphinxVerbatim}[commandchars=\\\{\}]
\PYG{k}{TYPE}\PYG{p}{(}\PYG{n}{lammps}\PYG{p}{)}\PYG{+w}{ }\PYG{k+kd}{::}\PYG{+w}{ }\PYG{n}{lmp}
\PYG{k+kt}{REAL}\PYG{p}{(}\PYG{k+kt}{c\PYGZus{}double}\PYG{p}{)}\PYG{+w}{ }\PYG{k+kd}{::}\PYG{+w}{ }\PYG{n}{area}
\PYG{c}{! more code to set up, etc.}
\PYG{n}{lmp}\PYG{p}{\PYGZpc{}}\PYG{n}{command}\PYG{p}{(}\PYG{l+s+s1}{\PYGZsq{}variable A equal lx*ly\PYGZsq{}}\PYG{p}{)}
\PYG{c}{! more code}
\PYG{n}{area}\PYG{+w}{ }\PYG{o}{=}\PYG{+w}{ }\PYG{n}{lmp}\PYG{p}{\PYGZpc{}}\PYG{n}{extract\PYGZus{}variable}\PYG{p}{(}\PYG{l+s+s2}{\PYGZdq{}A\PYGZdq{}}\PYG{p}{)}
\end{sphinxVerbatim}

\sphinxAtStartPar
will extract the \sphinxstyleemphasis{x}\textendash{}\sphinxstyleemphasis{y} cross\sphinxhyphen{}sectional area of the simulation into the
variable \sphinxstyleemphasis{area}.
\begin{quote}\begin{description}
\sphinxlineitem{Parameters}
\sphinxAtStartPar
\sphinxstylestrong{name}\sphinxstyleemphasis{ {[}}\sphinxstyleemphasis{character(len=*)}\sphinxstyleemphasis{{]}} :: variable name to evaluate

\sphinxlineitem{Options}
\sphinxAtStartPar
\sphinxstylestrong{group}\sphinxstyleemphasis{ {[}}\sphinxstyleemphasis{character(len=*)}\sphinxstyleemphasis{,}\sphinxstyleemphasis{optional}\sphinxstyleemphasis{{]}} :: group for which to extract per\sphinxhyphen{}atom
data (if absent, use “all”)

\sphinxlineitem{Call to}
\sphinxAtStartPar
{\hyperref[\detokenize{Library_objects:_CPPv423lammps_extract_variablePvPKcPKc}]{\sphinxcrossref{\sphinxcode{\sphinxupquote{lammps\_extract\_variable()}}}}}

\sphinxlineitem{Return}
\sphinxAtStartPar
\sphinxstylestrong{data}\sphinxstyleemphasis{ {[}}\sphinxstyleemphasis{polymorphic}\sphinxstyleemphasis{{]}} :: scalar of type \sphinxcode{\sphinxupquote{REAL(c\_double)}} (for \sphinxstyleemphasis{equal}\sphinxhyphen{}style
variables and others that are \sphinxstyleemphasis{equal}\sphinxhyphen{}compatible), vector of type
\sphinxcode{\sphinxupquote{REAL(c\_double), DIMENSION(:), ALLOCATABLE}} for \sphinxstyleemphasis{atom}\sphinxhyphen{} or \sphinxstyleemphasis{vector}\sphinxhyphen{}style
variables, or \sphinxcode{\sphinxupquote{CHARACTER(LEN=*)}} for \sphinxstyleemphasis{string}\sphinxhyphen{}style and compatible
variables. Strings whose length is too short to hold the result will be
truncated. Allocatable strings must be allocated before this function is
called; see note at \sphinxcode{\sphinxupquote{extract\_global()}} regarding allocatable strings.
Allocatable arrays (for \sphinxstyleemphasis{atom}\sphinxhyphen{} and \sphinxstyleemphasis{vector}\sphinxhyphen{}style data) will be
reallocated on assignment.

\end{description}\end{quote}

\end{fulllineitems}


\begin{sphinxadmonition}{note}{Note:}
\sphinxAtStartPar
LAMMPS cannot easily check if it is valid to access the data
referenced by the variables (e.g., computes, fixes, or thermodynamic
info), so it may fail with an error.  The caller has to make certain
that the data are extracted only when it is safe to evaluate the variable
and thus an error and crash are avoided.
\end{sphinxadmonition}


\bigskip\hrule\bigskip

\index{set\_variable() (fortran subroutine)@\spxentry{set\_variable()}\spxextra{fortran subroutine}|spxpagem}

\begin{fulllineitems}
\phantomsection\label{\detokenize{Fortran:f/_/set_variable}}
\pysigstartsignatures
\pysiglinewithargsret{\sphinxbfcode{\sphinxupquote{subroutine  }}\sphinxbfcode{\sphinxupquote{set\_variable}}}{\sphinxparam{name}\sphinxparamcomma \sphinxparam{str}}{}
\pysigstopsignatures
\sphinxAtStartPar
Set the value of a string\sphinxhyphen{}style variable.

\sphinxAtStartPar
\DUrole{versionmodified,deprecated}{Deprecated since version 7Feb2024.}

\sphinxAtStartPar
This function assigns a new value from the string \sphinxstyleemphasis{str} to the string\sphinxhyphen{}style
variable \sphinxstyleemphasis{name}. If \sphinxstyleemphasis{name} does not exist or is not a string\sphinxhyphen{}style
variable, an error is generated.

\begin{sphinxadmonition}{warning}{Warning:}
\sphinxAtStartPar
This subroutine is deprecated and \sphinxcode{\sphinxupquote{set\_string\_variable()}}
should be used instead.
\end{sphinxadmonition}
\begin{quote}\begin{description}
\sphinxlineitem{Parameters}\begin{itemize}
\item {} 
\sphinxAtStartPar
\sphinxstylestrong{name}\sphinxstyleemphasis{ {[}}\sphinxstyleemphasis{character(len=*)}\sphinxstyleemphasis{{]}} :: name of the variable

\item {} 
\sphinxAtStartPar
\sphinxstylestrong{str}\sphinxstyleemphasis{ {[}}\sphinxstyleemphasis{character(len=*)}\sphinxstyleemphasis{{]}} :: new value to assign to the variable

\end{itemize}

\sphinxlineitem{Call to}
\sphinxAtStartPar
{\hyperref[\detokenize{Library_objects:_CPPv419lammps_set_variablePvPKcPKc}]{\sphinxcrossref{\sphinxcode{\sphinxupquote{lammps\_set\_variable()}}}}}

\end{description}\end{quote}

\end{fulllineitems}



\bigskip\hrule\bigskip

\index{set\_string\_variable() (fortran subroutine)@\spxentry{set\_string\_variable()}\spxextra{fortran subroutine}|spxpagem}

\begin{fulllineitems}
\phantomsection\label{\detokenize{Fortran:f/_/set_string_variable}}
\pysigstartsignatures
\pysiglinewithargsret{\sphinxbfcode{\sphinxupquote{subroutine  }}\sphinxbfcode{\sphinxupquote{set\_string\_variable}}}{\sphinxparam{name}\sphinxparamcomma \sphinxparam{str}}{}
\pysigstopsignatures
\sphinxAtStartPar
Set the value of a string\sphinxhyphen{}style variable.

\sphinxAtStartPar
\DUrole{versionmodified,added}{Added in version 7Feb2024.}

\sphinxAtStartPar
This function assigns a new value from the string \sphinxstyleemphasis{str} to the string\sphinxhyphen{}style
variable \sphinxstyleemphasis{name}. If \sphinxstyleemphasis{name} does not exist or is not a string\sphinxhyphen{}style
variable, an error is generated.
\begin{quote}\begin{description}
\sphinxlineitem{Parameters}\begin{itemize}
\item {} 
\sphinxAtStartPar
\sphinxstylestrong{name}\sphinxstyleemphasis{ {[}}\sphinxstyleemphasis{character(len=*)}\sphinxstyleemphasis{{]}} :: name of the variable

\item {} 
\sphinxAtStartPar
\sphinxstylestrong{str}\sphinxstyleemphasis{ {[}}\sphinxstyleemphasis{character(len=*)}\sphinxstyleemphasis{{]}} :: new value to assign to the variable

\end{itemize}

\sphinxlineitem{Call to}
\sphinxAtStartPar
{\hyperref[\detokenize{Library_objects:_CPPv426lammps_set_string_variablePvPKcPKc}]{\sphinxcrossref{\sphinxcode{\sphinxupquote{lammps\_set\_string\_variable()}}}}}

\end{description}\end{quote}

\end{fulllineitems}



\bigskip\hrule\bigskip

\index{set\_internal\_variable() (fortran subroutine)@\spxentry{set\_internal\_variable()}\spxextra{fortran subroutine}|spxpagem}

\begin{fulllineitems}
\phantomsection\label{\detokenize{Fortran:f/_/set_internal_variable}}
\pysigstartsignatures
\pysiglinewithargsret{\sphinxbfcode{\sphinxupquote{subroutine  }}\sphinxbfcode{\sphinxupquote{set\_internal\_variable}}}{\sphinxparam{name}\sphinxparamcomma \sphinxparam{val}}{}
\pysigstopsignatures
\sphinxAtStartPar
Set the value of a internal\sphinxhyphen{}style variable.

\sphinxAtStartPar
\DUrole{versionmodified,added}{Added in version 7Feb2024.}

\sphinxAtStartPar
This function assigns a new value from the floating\sphinxhyphen{}point number \sphinxstyleemphasis{val} to
the internal\sphinxhyphen{}style variable \sphinxstyleemphasis{name}. If \sphinxstyleemphasis{name} does not exist or is not
an internal\sphinxhyphen{}style variable, an error is generated.
\begin{quote}\begin{description}
\sphinxlineitem{Parameters}\begin{itemize}
\item {} 
\sphinxAtStartPar
\sphinxstylestrong{name}\sphinxstyleemphasis{ {[}}\sphinxstyleemphasis{character(len=*)}\sphinxstyleemphasis{{]}} :: name of the variable

\item {} 
\sphinxAtStartPar
\sphinxstylestrong{val}\sphinxstyleemphasis{ {[}}\sphinxstyleemphasis{read(c\_double)}\sphinxstyleemphasis{{]}} :: new value to assign to the variable

\end{itemize}

\sphinxlineitem{Call to}
\sphinxAtStartPar
{\hyperref[\detokenize{Library_objects:_CPPv428lammps_set_internal_variablePvPKcd}]{\sphinxcrossref{\sphinxcode{\sphinxupquote{lammps\_set\_internal\_variable()}}}}}

\end{description}\end{quote}

\end{fulllineitems}



\bigskip\hrule\bigskip

\index{gather\_atoms() (fortran subroutine)@\spxentry{gather\_atoms()}\spxextra{fortran subroutine}|spxpagem}

\begin{fulllineitems}
\phantomsection\label{\detokenize{Fortran:f/_/gather_atoms}}
\pysigstartsignatures
\pysiglinewithargsret{\sphinxbfcode{\sphinxupquote{subroutine  }}\sphinxbfcode{\sphinxupquote{gather\_atoms}}}{\sphinxparam{name}\sphinxparamcomma \sphinxparam{count}\sphinxparamcomma \sphinxparam{data}}{}
\pysigstopsignatures
\sphinxAtStartPar
This function calls {\hyperref[\detokenize{Library_scatter:_CPPv419lammps_gather_atomsPvPKciiPv}]{\sphinxcrossref{\sphinxcode{\sphinxupquote{lammps\_gather\_atoms()}}}}} to gather the named
atom\sphinxhyphen{}based entity for all atoms on all processors and return it in the
vector \sphinxstyleemphasis{data}. The vector \sphinxstyleemphasis{data} will be ordered by atom
ID, which requires consecutive atom IDs (1 to \sphinxstyleemphasis{natoms}).

\sphinxAtStartPar
\DUrole{versionmodified,added}{Added in version 3Nov2022.}

\sphinxAtStartPar
If you need a similar array but have non\sphinxhyphen{}consecutive atom IDs, see
\sphinxcode{\sphinxupquote{gather\_atoms\_concat()}}; for a similar array but for a subset
of atoms, see \sphinxcode{\sphinxupquote{gather\_atoms\_subset()}}.

\sphinxAtStartPar
The \sphinxstyleemphasis{data} array will be ordered in groups of \sphinxstyleemphasis{count} values, sorted by atom
ID (e.g., if \sphinxstyleemphasis{name} is \sphinxstyleemphasis{x} and \sphinxstyleemphasis{count} = 3, then \sphinxstyleemphasis{data} = {[}\sphinxstyleemphasis{x}(1,1),
\sphinxstyleemphasis{x}(2,1), \sphinxstyleemphasis{x}(3,1), \sphinxstyleemphasis{x}(1,2), \sphinxstyleemphasis{x}(2,2), \sphinxstyleemphasis{x}(3,2), \sphinxstyleemphasis{x}(1,3),
\(\dots\){]}); \sphinxstyleemphasis{data} must be \sphinxcode{\sphinxupquote{ALLOCATABLE}} and will be allocated to
length (\sphinxstyleemphasis{count} \(\times\) \sphinxstyleemphasis{natoms}), as queried by
\sphinxcode{\sphinxupquote{get\_natoms()}}.

\sphinxAtStartPar
This function is not compatible with \sphinxcode{\sphinxupquote{\sphinxhyphen{}DLAMMPS\_BIGBIG}}.
\begin{quote}\begin{description}
\sphinxlineitem{Parameters}\begin{itemize}
\item {} 
\sphinxAtStartPar
\sphinxstylestrong{name}\sphinxstyleemphasis{ {[}}\sphinxstyleemphasis{character(len=*)}\sphinxstyleemphasis{{]}} :: desired quantity (e.g., \sphinxstyleemphasis{x} or \sphinxstyleemphasis{mask})

\item {} 
\sphinxAtStartPar
\sphinxstylestrong{count}\sphinxstyleemphasis{ {[}}\sphinxstyleemphasis{integer(c\_int)}\sphinxstyleemphasis{{]}} :: number of per\sphinxhyphen{}atom values you expect per atom
(e.g., 1 for \sphinxstyleemphasis{type}, \sphinxstyleemphasis{mask}, or \sphinxstyleemphasis{charge}; 3 for \sphinxstyleemphasis{x}, \sphinxstyleemphasis{v}, or \sphinxstyleemphasis{f}). Use
\sphinxstyleemphasis{count} = 3 with \sphinxstyleemphasis{image} if you want a single image flag unpacked into
\sphinxstyleemphasis{x}/\sphinxstyleemphasis{y}/\sphinxstyleemphasis{z} components.

\item {} 
\sphinxAtStartPar
\sphinxstylestrong{data}\sphinxstyleemphasis{ {[}}\sphinxstyleemphasis{polymorphic,dimension(:),allocatable}\sphinxstyleemphasis{{]}} :: array into which to store
the data. Array \sphinxstyleemphasis{must} have the \sphinxcode{\sphinxupquote{ALLOCATABLE}} attribute and be of rank 1
(i.e., \sphinxcode{\sphinxupquote{DIMENSION(:)}}). If this array is already allocated, it will be
reallocated to fit the length of the incoming data. It should have type
\sphinxcode{\sphinxupquote{INTEGER(c\_int)}} if expecting integer data and \sphinxcode{\sphinxupquote{REAL(c\_double)}} if
expecting floating\sphinxhyphen{}point data.

\end{itemize}

\sphinxlineitem{Call to}
\sphinxAtStartPar
{\hyperref[\detokenize{Library_scatter:_CPPv419lammps_gather_atomsPvPKciiPv}]{\sphinxcrossref{\sphinxcode{\sphinxupquote{lammps\_gather\_atoms()}}}}}

\end{description}\end{quote}

\begin{sphinxadmonition}{note}{Note:}
\sphinxAtStartPar
If you want data from this function to be accessible as a two\sphinxhyphen{}dimensional
array, you can declare a rank\sphinxhyphen{}2 pointer and reassign it, like so:

\begin{sphinxVerbatim}[commandchars=\\\{\}]
\PYG{k}{USE}\PYG{p}{,}\PYG{+w}{ }\PYG{k}{INTRINSIC}\PYG{+w}{ }\PYG{k+kd}{::}\PYG{+w}{ }\PYG{n+nb}{ISO\PYGZus{}C\PYGZus{}BINDING}
\PYG{k}{USE }\PYG{n}{LIBLAMMPS}
\PYG{k}{TYPE}\PYG{p}{(}\PYG{n}{lammps}\PYG{p}{)}\PYG{+w}{ }\PYG{k+kd}{::}\PYG{+w}{ }\PYG{n}{lmp}
\PYG{k+kt}{REAL}\PYG{p}{(}\PYG{k+kt}{c\PYGZus{}double}\PYG{p}{)}\PYG{p}{,}\PYG{+w}{ }\PYG{k}{DIMENSION}\PYG{p}{(}\PYG{p}{:}\PYG{p}{)}\PYG{p}{,}\PYG{+w}{ }\PYG{k}{ALLOCATABLE}\PYG{p}{,}\PYG{+w}{ }\PYG{k}{TARGET}\PYG{+w}{ }\PYG{k+kd}{::}\PYG{+w}{ }\PYG{n}{xdata}
\PYG{k+kt}{REAL}\PYG{p}{(}\PYG{k+kt}{c\PYGZus{}double}\PYG{p}{)}\PYG{p}{,}\PYG{+w}{ }\PYG{k}{DIMENSION}\PYG{p}{(}\PYG{p}{:}\PYG{p}{,}\PYG{p}{:}\PYG{p}{)}\PYG{p}{,}\PYG{+w}{ }\PYG{k}{POINTER}\PYG{+w}{ }\PYG{k+kd}{::}\PYG{+w}{ }\PYG{n}{x}
\PYG{c}{! other code to set up, etc.}
\PYG{k}{CALL }\PYG{n}{lmp}\PYG{p}{\PYGZpc{}}\PYG{n}{gather\PYGZus{}atoms}\PYG{p}{(}\PYG{l+s+s1}{\PYGZsq{}x\PYGZsq{}}\PYG{p}{,}\PYG{l+m+mi}{3}\PYG{p}{,}\PYG{n}{xdata}\PYG{p}{)}
\PYG{n}{x}\PYG{p}{(}\PYG{l+m+mi}{1}\PYG{p}{:}\PYG{l+m+mi}{3}\PYG{p}{,}\PYG{l+m+mi}{1}\PYG{p}{:}\PYG{n}{size}\PYG{p}{(}\PYG{n}{xdata}\PYG{p}{)}\PYG{o}{/}\PYG{l+m+mi}{3}\PYG{p}{)}\PYG{+w}{ }\PYG{o}{=}\PYG{o}{\PYGZgt{}}\PYG{+w}{ }\PYG{n}{xdata}
\end{sphinxVerbatim}

\sphinxAtStartPar
You can then access the \sphinxstyleemphasis{y}\sphinxhyphen{}component of atom 3 with \sphinxcode{\sphinxupquote{x(2,3)}}.
See the note about array index order at \sphinxcode{\sphinxupquote{extract\_atom()}}.
\end{sphinxadmonition}

\end{fulllineitems}



\bigskip\hrule\bigskip

\index{gather\_atoms\_concat() (fortran subroutine)@\spxentry{gather\_atoms\_concat()}\spxextra{fortran subroutine}|spxpagem}

\begin{fulllineitems}
\phantomsection\label{\detokenize{Fortran:f/_/gather_atoms_concat}}
\pysigstartsignatures
\pysiglinewithargsret{\sphinxbfcode{\sphinxupquote{subroutine  }}\sphinxbfcode{\sphinxupquote{gather\_atoms\_concat}}}{\sphinxparam{name}\sphinxparamcomma \sphinxparam{count}\sphinxparamcomma \sphinxparam{data}}{}
\pysigstopsignatures
\sphinxAtStartPar
This function calls {\hyperref[\detokenize{Library_scatter:_CPPv426lammps_gather_atoms_concatPvPKciiPv}]{\sphinxcrossref{\sphinxcode{\sphinxupquote{lammps\_gather\_atoms\_concat()}}}}} to gather the
named atom\sphinxhyphen{}based entity for all atoms on all processors and return it in the
vector \sphinxstyleemphasis{data}.

\sphinxAtStartPar
\DUrole{versionmodified,added}{Added in version 3Nov2022.}

\sphinxAtStartPar
The vector \sphinxstyleemphasis{data} will not be ordered by atom ID, and there is no
restriction on the IDs being consecutive. If you need the IDs, you can do
another \sphinxcode{\sphinxupquote{gather\_atoms\_concat()}} with \sphinxstyleemphasis{name} set to \sphinxcode{\sphinxupquote{id}}.

\sphinxAtStartPar
If you need a similar array but have consecutive atom IDs, see
\sphinxcode{\sphinxupquote{gather\_atoms()}}; for a similar array but for a subset of atoms, see
\sphinxcode{\sphinxupquote{gather\_atoms\_subset()}}.

\sphinxAtStartPar
This function is not compatible with \sphinxcode{\sphinxupquote{\sphinxhyphen{}DLAMMPS\_BIGBIG}}.
\begin{quote}\begin{description}
\sphinxlineitem{Parameters}\begin{itemize}
\item {} 
\sphinxAtStartPar
\sphinxstylestrong{name}\sphinxstyleemphasis{ {[}}\sphinxstyleemphasis{character(len=*)}\sphinxstyleemphasis{{]}} :: desired quantity (e.g., \sphinxstyleemphasis{x} or \sphinxstyleemphasis{mask})

\item {} 
\sphinxAtStartPar
\sphinxstylestrong{count}\sphinxstyleemphasis{ {[}}\sphinxstyleemphasis{integer(c\_int)}\sphinxstyleemphasis{{]}} :: number of per\sphinxhyphen{}atom values you expect per atom
(e.g., 1 for \sphinxstyleemphasis{type}, \sphinxstyleemphasis{mask}, or \sphinxstyleemphasis{charge}; 3 for \sphinxstyleemphasis{x}, \sphinxstyleemphasis{v}, or \sphinxstyleemphasis{f}). Use
\sphinxstyleemphasis{count} = 3 with \sphinxstyleemphasis{image} if you want a single image flag unpacked into
\sphinxstyleemphasis{x}/\sphinxstyleemphasis{y}/\sphinxstyleemphasis{z} components.

\item {} 
\sphinxAtStartPar
\sphinxstylestrong{data}\sphinxstyleemphasis{ {[}}\sphinxstyleemphasis{polymorphic,dimension(:),allocatable}\sphinxstyleemphasis{{]}} :: array into which to store
the data. Array \sphinxstyleemphasis{must} have the \sphinxcode{\sphinxupquote{ALLOCATABLE}} attribute and be of rank 1
(i.e., \sphinxcode{\sphinxupquote{DIMENSION(:)}}). If this array is already allocated, it will be
reallocated to fit the length of the incoming data. It should have type
\sphinxcode{\sphinxupquote{INTEGER(c\_int)}} if expecting integer data and \sphinxcode{\sphinxupquote{REAL(c\_double)}} if
expecting floating\sphinxhyphen{}point data.

\end{itemize}

\sphinxlineitem{Call to}
\sphinxAtStartPar
{\hyperref[\detokenize{Library_scatter:_CPPv426lammps_gather_atoms_concatPvPKciiPv}]{\sphinxcrossref{\sphinxcode{\sphinxupquote{lammps\_gather\_atoms\_concat()}}}}}

\end{description}\end{quote}

\end{fulllineitems}



\bigskip\hrule\bigskip

\index{gather\_atoms\_subset() (fortran subroutine)@\spxentry{gather\_atoms\_subset()}\spxextra{fortran subroutine}|spxpagem}

\begin{fulllineitems}
\phantomsection\label{\detokenize{Fortran:f/_/gather_atoms_subset}}
\pysigstartsignatures
\pysiglinewithargsret{\sphinxbfcode{\sphinxupquote{subroutine  }}\sphinxbfcode{\sphinxupquote{gather\_atoms\_subset}}}{\sphinxparam{name}\sphinxparamcomma \sphinxparam{count}\sphinxparamcomma \sphinxparam{ids}\sphinxparamcomma \sphinxparam{data}}{}
\pysigstopsignatures
\sphinxAtStartPar
This function calls {\hyperref[\detokenize{Library_scatter:_CPPv426lammps_gather_atoms_subsetPvPKciiiPiPv}]{\sphinxcrossref{\sphinxcode{\sphinxupquote{lammps\_gather\_atoms\_subset()}}}}} to gather the
named atom\sphinxhyphen{}based entity for the atoms in the array \sphinxstyleemphasis{ids} from all processors
and return it in the vector \sphinxstyleemphasis{data}.

\sphinxAtStartPar
\DUrole{versionmodified,added}{Added in version 3Nov2022.}

\sphinxAtStartPar
This subroutine gathers data for the requested atom IDs and stores them in a
one\sphinxhyphen{}dimensional allocatable array. The data will be ordered by
atom ID, but there is no requirement that the IDs be consecutive. If you
wish to return a similar array for \sphinxstyleemphasis{all} the atoms, use
\sphinxcode{\sphinxupquote{gather\_atoms()}} or \sphinxcode{\sphinxupquote{gather\_atoms\_concat()}}.

\sphinxAtStartPar
The \sphinxstyleemphasis{data} array will be in groups of \sphinxstyleemphasis{count} values, sorted by atom ID
in the same order as the array \sphinxstyleemphasis{ids} (e.g., if \sphinxstyleemphasis{name} is \sphinxstyleemphasis{x}, \sphinxstyleemphasis{count} = 3,
and \sphinxstyleemphasis{ids} is {[}100, 57, 210{]}, then \sphinxstyleemphasis{data} might look like
{[}\sphinxstyleemphasis{x}(1,100), \sphinxstyleemphasis{x}(2,100), \sphinxstyleemphasis{x}(3,100), \sphinxstyleemphasis{x}(1,57), \sphinxstyleemphasis{x}(2,57),
\sphinxstyleemphasis{x}(3,57), \sphinxstyleemphasis{x}(1,210), \(\dots\){]}; \sphinxstyleemphasis{ids} must be provided by the
user, and \sphinxstyleemphasis{data} must be of rank 1 (i.e., \sphinxcode{\sphinxupquote{DIMENSION(:)}}) and have the
\sphinxcode{\sphinxupquote{ALLOCATABLE}} attribute.

\sphinxAtStartPar
This function is not compatible with \sphinxcode{\sphinxupquote{\sphinxhyphen{}DLAMMPS\_BIGBIG}}.
\begin{quote}\begin{description}
\sphinxlineitem{Parameters}\begin{itemize}
\item {} 
\sphinxAtStartPar
\sphinxstylestrong{name}\sphinxstyleemphasis{ {[}}\sphinxstyleemphasis{character(len=*)}\sphinxstyleemphasis{{]}} :: desired quantity (e.g., \sphinxstyleemphasis{x} or \sphinxstyleemphasis{mask})

\item {} 
\sphinxAtStartPar
\sphinxstylestrong{count}\sphinxstyleemphasis{ {[}}\sphinxstyleemphasis{integer(c\_int)}\sphinxstyleemphasis{{]}} :: number of per\sphinxhyphen{}atom values you expect per atom
(e.g., 1 for \sphinxstyleemphasis{type}, \sphinxstyleemphasis{mask}, or \sphinxstyleemphasis{charge}; 3 for \sphinxstyleemphasis{x}, \sphinxstyleemphasis{v}, or \sphinxstyleemphasis{f}). Use
\sphinxstyleemphasis{count} = 3 with \sphinxstyleemphasis{image} if you want a single image flag unpacked into
\sphinxstyleemphasis{x}/\sphinxstyleemphasis{y}/\sphinxstyleemphasis{z} components.

\item {} 
\sphinxAtStartPar
\sphinxstylestrong{ids}\sphinxstyleemphasis{ {[}}\sphinxstyleemphasis{integer(c\_int)}\sphinxstyleemphasis{,}\sphinxstyleemphasis{dimension(:)}\sphinxstyleemphasis{{]}} :: atom IDs corresponding to the atoms
to be gathered

\item {} 
\sphinxAtStartPar
\sphinxstylestrong{data}\sphinxstyleemphasis{ {[}}\sphinxstyleemphasis{polymorphic,dimension(:),allocatable}\sphinxstyleemphasis{{]}} :: array into which to store
the data. Array \sphinxstyleemphasis{must} have the \sphinxcode{\sphinxupquote{ALLOCATABLE}} attribute and be of rank 1
(i.e., \sphinxcode{\sphinxupquote{DIMENSION(:)}}). If this array is already allocated, it will be
reallocated to fit the length of the incoming data. It should have type
\sphinxcode{\sphinxupquote{INTEGER(c\_int)}} if expecting integer data and \sphinxcode{\sphinxupquote{REAL(c\_double)}} if
expecting floating\sphinxhyphen{}point data.

\end{itemize}

\sphinxlineitem{Call to}
\sphinxAtStartPar
{\hyperref[\detokenize{Library_scatter:_CPPv426lammps_gather_atoms_subsetPvPKciiiPiPv}]{\sphinxcrossref{\sphinxcode{\sphinxupquote{lammps\_gather\_atoms\_subset()}}}}}

\end{description}\end{quote}

\end{fulllineitems}



\bigskip\hrule\bigskip

\index{scatter\_atoms() (fortran subroutine)@\spxentry{scatter\_atoms()}\spxextra{fortran subroutine}|spxpagem}

\begin{fulllineitems}
\phantomsection\label{\detokenize{Fortran:f/_/scatter_atoms}}
\pysigstartsignatures
\pysiglinewithargsret{\sphinxbfcode{\sphinxupquote{subroutine  }}\sphinxbfcode{\sphinxupquote{scatter\_atoms}}}{\sphinxparam{name}\sphinxparamcomma \sphinxparam{data}}{}
\pysigstopsignatures
\sphinxAtStartPar
This function calls {\hyperref[\detokenize{Library_scatter:_CPPv420lammps_scatter_atomsPvPKciiPv}]{\sphinxcrossref{\sphinxcode{\sphinxupquote{lammps\_scatter\_atoms()}}}}} to scatter the named
atom\sphinxhyphen{}based entities in \sphinxstyleemphasis{data} to all processors.

\sphinxAtStartPar
\DUrole{versionmodified,added}{Added in version 3Nov2022.}

\sphinxAtStartPar
This subroutine takes data stored in a one\sphinxhyphen{}dimensional array supplied by the
user and scatters them to all atoms on all processors. The data must be
ordered by atom ID, with the requirement that the IDs be consecutive.
Use \sphinxcode{\sphinxupquote{scatter\_atoms\_subset()}} to scatter data for some (or all)
atoms, in any order.

\sphinxAtStartPar
The \sphinxstyleemphasis{data} array needs to be ordered in groups of \sphinxstyleemphasis{count} values, sorted by
atom ID (e.g., if \sphinxstyleemphasis{name} is \sphinxstyleemphasis{x} and \sphinxstyleemphasis{count} = 3, then
\sphinxstyleemphasis{data} = {[}\sphinxstyleemphasis{x}(1,1), \sphinxstyleemphasis{x}(2,1), \sphinxstyleemphasis{x}(3,1), \sphinxstyleemphasis{x}(1,2), \sphinxstyleemphasis{x}(2,2),
\sphinxstyleemphasis{x}(3,2), \sphinxstyleemphasis{x}(1,3), \(\dots\){]}); \sphinxstyleemphasis{data} must be of length \sphinxstyleemphasis{natoms}
or 3*\sphinxstyleemphasis{natoms}.
\begin{quote}\begin{description}
\sphinxlineitem{Parameters}\begin{itemize}
\item {} 
\sphinxAtStartPar
\sphinxstylestrong{name}\sphinxstyleemphasis{ {[}}\sphinxstyleemphasis{character(len=*)}\sphinxstyleemphasis{{]}} :: quantity to be scattered (e.g., \sphinxstyleemphasis{x} or \sphinxstyleemphasis{charge})

\item {} 
\sphinxAtStartPar
\sphinxstylestrong{data}\sphinxstyleemphasis{ {[}}\sphinxstyleemphasis{polymorphic,dimension(:)}\sphinxstyleemphasis{{]}} :: per\sphinxhyphen{}atom values packed in a one\sphinxhyphen{}dimensional array
containing the data to be scattered. This array must have length \sphinxstyleemphasis{natoms}
(e.g., for \sphinxstyleemphasis{type} or \sphinxstyleemphasis{charge}) or length \sphinxstyleemphasis{natoms}\({}\times 3\)
(e.g., for \sphinxstyleemphasis{x} or \sphinxstyleemphasis{f}). The array \sphinxstyleemphasis{data} must be rank 1 (i.e.,
\sphinxcode{\sphinxupquote{DIMENSION(:)}}) and be of type \sphinxcode{\sphinxupquote{INTEGER(c\_int)}} (e.g., for \sphinxstyleemphasis{mask} or
\sphinxstyleemphasis{type}) or of type \sphinxcode{\sphinxupquote{REAL(c\_double)}} (e.g., for \sphinxstyleemphasis{x} or \sphinxstyleemphasis{charge} or \sphinxstyleemphasis{f}).

\end{itemize}

\sphinxlineitem{Call to}
\sphinxAtStartPar
{\hyperref[\detokenize{Library_scatter:_CPPv420lammps_scatter_atomsPvPKciiPv}]{\sphinxcrossref{\sphinxcode{\sphinxupquote{lammps\_scatter\_atoms()}}}}}

\end{description}\end{quote}

\end{fulllineitems}



\bigskip\hrule\bigskip

\index{scatter\_atoms\_subset() (fortran subroutine)@\spxentry{scatter\_atoms\_subset()}\spxextra{fortran subroutine}|spxpagem}

\begin{fulllineitems}
\phantomsection\label{\detokenize{Fortran:f/_/scatter_atoms_subset}}
\pysigstartsignatures
\pysiglinewithargsret{\sphinxbfcode{\sphinxupquote{subroutine  }}\sphinxbfcode{\sphinxupquote{scatter\_atoms\_subset}}}{\sphinxparam{name}\sphinxparamcomma \sphinxparam{ids}\sphinxparamcomma \sphinxparam{data}}{}
\pysigstopsignatures
\sphinxAtStartPar
This function calls {\hyperref[\detokenize{Library_scatter:_CPPv427lammps_scatter_atoms_subsetPvPKciiiPiPv}]{\sphinxcrossref{\sphinxcode{\sphinxupquote{lammps\_scatter\_atoms\_subset()}}}}} to scatter the
named atom\sphinxhyphen{}based entities in \sphinxstyleemphasis{data} to all processors.

\sphinxAtStartPar
\DUrole{versionmodified,added}{Added in version 3Nov2022.}

\sphinxAtStartPar
This subroutine takes data stored in a one\sphinxhyphen{}dimensional array supplied by the
user and scatters them to a subset of atoms on all processors. The array
\sphinxstyleemphasis{data} contains data associated with atom IDs, but there is no requirement
that the IDs be consecutive, as they are provided in a separate array,
\sphinxstyleemphasis{ids}. Use \sphinxcode{\sphinxupquote{scatter\_atoms()}} to scatter data for all atoms, in order.

\sphinxAtStartPar
The \sphinxstyleemphasis{data} array needs to be organized in groups of 1 or 3 values,
depending on which quantity is being scattered, with the groups in the same
order as the array \sphinxstyleemphasis{ids}. For example, if you want \sphinxstyleemphasis{data} to be the array
{[}\sphinxstyleemphasis{x}(1,1), \sphinxstyleemphasis{x}(2,1), \sphinxstyleemphasis{x}(3,1), \sphinxstyleemphasis{x}(1,100), \sphinxstyleemphasis{x}(2,100),
\sphinxstyleemphasis{x}(3,100), \sphinxstyleemphasis{x}(1,57), \sphinxstyleemphasis{x}(2,57), \sphinxstyleemphasis{x}(3,57){]}, then \sphinxstyleemphasis{ids} would be
{[}1, 100, 57{]} and \sphinxstyleemphasis{name} would be \sphinxstyleemphasis{x}.
\begin{quote}\begin{description}
\sphinxlineitem{Parameters}\begin{itemize}
\item {} 
\sphinxAtStartPar
\sphinxstylestrong{name}\sphinxstyleemphasis{ {[}}\sphinxstyleemphasis{character(len=*)}\sphinxstyleemphasis{{]}} :: quantity to be scattered (e.g., \sphinxstyleemphasis{x} or \sphinxstyleemphasis{charge})

\item {} 
\sphinxAtStartPar
\sphinxstylestrong{ids}\sphinxstyleemphasis{ {[}}\sphinxstyleemphasis{integer(c\_int)}\sphinxstyleemphasis{,}\sphinxstyleemphasis{dimension(:)}\sphinxstyleemphasis{{]}} :: atom IDs corresponding to the atoms
being scattered

\item {} 
\sphinxAtStartPar
\sphinxstylestrong{data}\sphinxstyleemphasis{ {[}}\sphinxstyleemphasis{polymorphic,dimension(:)}\sphinxstyleemphasis{{]}} :: per\sphinxhyphen{}atom values packed into a
one\sphinxhyphen{}dimensional array containing the data to be scattered. This array must
have either the same length as \sphinxstyleemphasis{ids} (for \sphinxstyleemphasis{mask}, \sphinxstyleemphasis{type}, etc.) or three
times its length (for \sphinxstyleemphasis{x}, \sphinxstyleemphasis{f}, etc.); the array must be rank 1
and be of type \sphinxcode{\sphinxupquote{INTEGER(c\_int)}} (e.g., for \sphinxstyleemphasis{mask} or \sphinxstyleemphasis{type}) or of type
\sphinxcode{\sphinxupquote{REAL(c\_double)}} (e.g., for \sphinxstyleemphasis{charge}, \sphinxstyleemphasis{x}, or \sphinxstyleemphasis{f}).

\end{itemize}

\sphinxlineitem{Call to}
\sphinxAtStartPar
{\hyperref[\detokenize{Library_scatter:_CPPv427lammps_scatter_atoms_subsetPvPKciiiPiPv}]{\sphinxcrossref{\sphinxcode{\sphinxupquote{lammps\_scatter\_atoms\_subset()}}}}}

\end{description}\end{quote}

\end{fulllineitems}



\bigskip\hrule\bigskip

\index{gather\_bonds() (fortran subroutine)@\spxentry{gather\_bonds()}\spxextra{fortran subroutine}|spxpagem}

\begin{fulllineitems}
\phantomsection\label{\detokenize{Fortran:f/_/gather_bonds}}
\pysigstartsignatures
\pysiglinewithargsret{\sphinxbfcode{\sphinxupquote{subroutine  }}\sphinxbfcode{\sphinxupquote{gather\_bonds}}}{\sphinxparam{data}}{}
\pysigstopsignatures
\sphinxAtStartPar
Gather type and constituent atom information for all bonds.

\sphinxAtStartPar
\DUrole{versionmodified,added}{Added in version 3Nov2022.}

\sphinxAtStartPar
This function copies the list of all bonds into an allocatable array.
The array will be filled with (bond type, bond atom 1, bond atom 2) for each
bond. The array is allocated to the right length (i.e., three times the
number of bonds). The array \sphinxstyleemphasis{data} must be of the same type as the LAMMPS
\sphinxcode{\sphinxupquote{tagint}} type, which is equivalent to either \sphinxcode{\sphinxupquote{INTEGER(c\_int)}} or
\sphinxcode{\sphinxupquote{INTEGER(c\_int64\_t)}}, depending on whether \sphinxcode{\sphinxupquote{\sphinxhyphen{}DLAMMPS\_BIGBIG}} was used
when LAMMPS was built. If the supplied array does not match, an error will
result at run\sphinxhyphen{}time.

\sphinxAtStartPar
An example of how to use this routine is below:

\begin{sphinxVerbatim}[commandchars=\\\{\}]
\PYG{k}{PROGRAM }\PYG{n}{bonds}
\PYG{+w}{  }\PYG{k}{USE}\PYG{p}{,}\PYG{+w}{ }\PYG{k}{INTRINSIC}\PYG{+w}{ }\PYG{k+kd}{::}\PYG{+w}{ }\PYG{n+nb}{ISO\PYGZus{}C\PYGZus{}BINDING}\PYG{p}{,}\PYG{+w}{ }\PYG{k}{ONLY}\PYG{+w}{ }\PYG{p}{:}\PYG{+w}{ }\PYG{k+kt}{c\PYGZus{}int}
\PYG{k+kt}{  }\PYG{k}{USE}\PYG{p}{,}\PYG{+w}{ }\PYG{k}{INTRINSIC}\PYG{+w}{ }\PYG{k+kd}{::}\PYG{+w}{ }\PYG{n}{ISO\PYGZus{}FORTRAN\PYGZus{}ENV}\PYG{p}{,}\PYG{+w}{ }\PYG{k}{ONLY}\PYG{+w}{ }\PYG{p}{:}\PYG{+w}{ }\PYG{n}{OUTPUT\PYGZus{}UNIT}
\PYG{+w}{  }\PYG{k}{USE }\PYG{n}{LIBLAMMPS}
\PYG{+w}{  }\PYG{k}{IMPLICIT }\PYG{k}{NONE}
\PYG{k}{  }\PYG{k+kt}{INTEGER}\PYG{p}{(}\PYG{k+kt}{c\PYGZus{}int}\PYG{p}{)}\PYG{p}{,}\PYG{+w}{ }\PYG{k}{DIMENSION}\PYG{p}{(}\PYG{p}{:}\PYG{p}{)}\PYG{p}{,}\PYG{+w}{ }\PYG{k}{ALLOCATABLE}\PYG{p}{,}\PYG{+w}{ }\PYG{k}{TARGET}\PYG{+w}{ }\PYG{k+kd}{::}\PYG{+w}{ }\PYG{n}{bonds}
\PYG{+w}{  }\PYG{k+kt}{INTEGER}\PYG{p}{(}\PYG{k+kt}{c\PYGZus{}int}\PYG{p}{)}\PYG{p}{,}\PYG{+w}{ }\PYG{k}{DIMENSION}\PYG{p}{(}\PYG{p}{:}\PYG{p}{,}\PYG{p}{:}\PYG{p}{)}\PYG{p}{,}\PYG{+w}{ }\PYG{k}{POINTER}\PYG{+w}{ }\PYG{k+kd}{::}\PYG{+w}{ }\PYG{n}{bonds\PYGZus{}array}
\PYG{+w}{  }\PYG{k}{TYPE}\PYG{p}{(}\PYG{n}{lammps}\PYG{p}{)}\PYG{+w}{ }\PYG{k+kd}{::}\PYG{+w}{ }\PYG{n}{lmp}
\PYG{+w}{  }\PYG{k+kt}{INTEGER}\PYG{+w}{ }\PYG{k+kd}{::}\PYG{+w}{ }\PYG{n}{i}
\PYG{+w}{  }\PYG{c}{! other commands to initialize LAMMPS, create bonds, etc.}
\PYG{+w}{  }\PYG{k}{CALL }\PYG{n}{lmp}\PYG{p}{\PYGZpc{}}\PYG{n}{gather\PYGZus{}bonds}\PYG{p}{(}\PYG{n}{bonds}\PYG{p}{)}
\PYG{+w}{  }\PYG{n}{bonds\PYGZus{}array}\PYG{p}{(}\PYG{l+m+mi}{1}\PYG{p}{:}\PYG{l+m+mi}{3}\PYG{p}{,}\PYG{+w}{ }\PYG{l+m+mi}{1}\PYG{p}{:}\PYG{n}{SIZE}\PYG{p}{(}\PYG{n}{bonds}\PYG{p}{)}\PYG{o}{/}\PYG{l+m+mi}{3}\PYG{p}{)}\PYG{+w}{ }\PYG{o}{=}\PYG{o}{\PYGZgt{}}\PYG{+w}{ }\PYG{n}{bonds}
\PYG{+w}{  }\PYG{k}{DO }\PYG{n}{i}\PYG{+w}{ }\PYG{o}{=}\PYG{+w}{ }\PYG{l+m+mi}{1}\PYG{p}{,}\PYG{+w}{ }\PYG{n}{SIZE}\PYG{p}{(}\PYG{n}{bonds}\PYG{p}{)}\PYG{o}{/}\PYG{l+m+mi}{3}
\PYG{+w}{    }\PYG{k}{WRITE}\PYG{p}{(}\PYG{n}{OUTPUT\PYGZus{}UNIT}\PYG{p}{,}\PYG{l+s+s1}{\PYGZsq{}(A,1X,I4,A,I4,1X,I4)\PYGZsq{}}\PYG{p}{)}\PYG{+w}{ }\PYG{l+s+s1}{\PYGZsq{}bond\PYGZsq{}}\PYG{p}{,}\PYG{+w}{ }\PYG{n}{bonds\PYGZus{}array}\PYG{p}{(}\PYG{l+m+mi}{1}\PYG{p}{,}\PYG{n}{i}\PYG{p}{)}\PYG{p}{,}\PYG{+w}{ }\PYG{p}{\PYGZam{}}
\PYG{+w}{      }\PYG{l+s+s1}{\PYGZsq{}; type = \PYGZsq{}}\PYG{p}{,}\PYG{+w}{ }\PYG{n}{bonds\PYGZus{}array}\PYG{p}{(}\PYG{l+m+mi}{2}\PYG{p}{,}\PYG{n}{i}\PYG{p}{)}\PYG{p}{,}\PYG{+w}{ }\PYG{n}{bonds\PYGZus{}array}\PYG{p}{(}\PYG{l+m+mi}{3}\PYG{p}{,}\PYG{n}{i}\PYG{p}{)}
\PYG{+w}{  }\PYG{k}{END }\PYG{k}{DO}
\PYG{k}{END }\PYG{k}{PROGRAM }\PYG{n}{bonds}
\end{sphinxVerbatim}
\begin{quote}\begin{description}
\sphinxlineitem{Parameters}
\sphinxAtStartPar
\sphinxstylestrong{data}\sphinxstyleemphasis{ {[}}\sphinxstyleemphasis{integer(kind=*),allocatable}\sphinxstyleemphasis{{]}} :: array into which to copy the result. *The \sphinxcode{\sphinxupquote{KIND}} parameter is
either \sphinxcode{\sphinxupquote{c\_int}} or, if LAMMPS was compiled with \sphinxcode{\sphinxupquote{\sphinxhyphen{}DLAMMPS\_BIGBIG}},
kind \sphinxcode{\sphinxupquote{c\_int64\_t}}.

\sphinxlineitem{Call to}
\sphinxAtStartPar
{\hyperref[\detokenize{Library_scatter:_CPPv419lammps_gather_bondsPvPv}]{\sphinxcrossref{\sphinxcode{\sphinxupquote{lammps\_gather\_bonds()}}}}}

\end{description}\end{quote}

\end{fulllineitems}



\bigskip\hrule\bigskip

\index{gather\_angles() (fortran subroutine)@\spxentry{gather\_angles()}\spxextra{fortran subroutine}|spxpagem}

\begin{fulllineitems}
\phantomsection\label{\detokenize{Fortran:f/_/gather_angles}}
\pysigstartsignatures
\pysiglinewithargsret{\sphinxbfcode{\sphinxupquote{subroutine  }}\sphinxbfcode{\sphinxupquote{gather\_angles}}}{\sphinxparam{data}}{}
\pysigstopsignatures
\sphinxAtStartPar
Gather type and constituent atom information for all angles.

\sphinxAtStartPar
\DUrole{versionmodified,added}{Added in version 8Feb2023.}

\sphinxAtStartPar
This function copies the list of all angles into an allocatable array.
The array will be filled with (angle type, angle atom 1, angle atom 2, angle atom 3)
for each angle. The array is allocated to the right length (i.e., four times the
number of angles). The array \sphinxstyleemphasis{data} must be of the same type as the LAMMPS
\sphinxcode{\sphinxupquote{tagint}} type, which is equivalent to either \sphinxcode{\sphinxupquote{INTEGER(c\_int)}} or
\sphinxcode{\sphinxupquote{INTEGER(c\_int64\_t)}}, depending on whether \sphinxcode{\sphinxupquote{\sphinxhyphen{}DLAMMPS\_BIGBIG}} was used
when LAMMPS was built. If the supplied array does not match, an error will
result at run\sphinxhyphen{}time.

\sphinxAtStartPar
An example of how to use this routine is below:

\begin{sphinxVerbatim}[commandchars=\\\{\}]
\PYG{k}{PROGRAM }\PYG{n}{angles}
\PYG{+w}{  }\PYG{k}{USE}\PYG{p}{,}\PYG{+w}{ }\PYG{k}{INTRINSIC}\PYG{+w}{ }\PYG{k+kd}{::}\PYG{+w}{ }\PYG{n+nb}{ISO\PYGZus{}C\PYGZus{}BINDING}\PYG{p}{,}\PYG{+w}{ }\PYG{k}{ONLY}\PYG{+w}{ }\PYG{p}{:}\PYG{+w}{ }\PYG{k+kt}{c\PYGZus{}int}
\PYG{k+kt}{  }\PYG{k}{USE}\PYG{p}{,}\PYG{+w}{ }\PYG{k}{INTRINSIC}\PYG{+w}{ }\PYG{k+kd}{::}\PYG{+w}{ }\PYG{n}{ISO\PYGZus{}FORTRAN\PYGZus{}ENV}\PYG{p}{,}\PYG{+w}{ }\PYG{k}{ONLY}\PYG{+w}{ }\PYG{p}{:}\PYG{+w}{ }\PYG{n}{OUTPUT\PYGZus{}UNIT}
\PYG{+w}{  }\PYG{k}{USE }\PYG{n}{LIBLAMMPS}
\PYG{+w}{  }\PYG{k}{IMPLICIT }\PYG{k}{NONE}
\PYG{k}{  }\PYG{k+kt}{INTEGER}\PYG{p}{(}\PYG{k+kt}{c\PYGZus{}int}\PYG{p}{)}\PYG{p}{,}\PYG{+w}{ }\PYG{k}{DIMENSION}\PYG{p}{(}\PYG{p}{:}\PYG{p}{)}\PYG{p}{,}\PYG{+w}{ }\PYG{k}{ALLOCATABLE}\PYG{p}{,}\PYG{+w}{ }\PYG{k}{TARGET}\PYG{+w}{ }\PYG{k+kd}{::}\PYG{+w}{ }\PYG{n}{angles}
\PYG{+w}{  }\PYG{k+kt}{INTEGER}\PYG{p}{(}\PYG{k+kt}{c\PYGZus{}int}\PYG{p}{)}\PYG{p}{,}\PYG{+w}{ }\PYG{k}{DIMENSION}\PYG{p}{(}\PYG{p}{:}\PYG{p}{,}\PYG{p}{:}\PYG{p}{)}\PYG{p}{,}\PYG{+w}{ }\PYG{k}{POINTER}\PYG{+w}{ }\PYG{k+kd}{::}\PYG{+w}{ }\PYG{n}{angles\PYGZus{}array}
\PYG{+w}{  }\PYG{k}{TYPE}\PYG{p}{(}\PYG{n}{lammps}\PYG{p}{)}\PYG{+w}{ }\PYG{k+kd}{::}\PYG{+w}{ }\PYG{n}{lmp}
\PYG{+w}{  }\PYG{k+kt}{INTEGER}\PYG{+w}{ }\PYG{k+kd}{::}\PYG{+w}{ }\PYG{n}{i}
\PYG{+w}{  }\PYG{c}{! other commands to initialize LAMMPS, create angles, etc.}
\PYG{+w}{  }\PYG{k}{CALL }\PYG{n}{lmp}\PYG{p}{\PYGZpc{}}\PYG{n}{gather\PYGZus{}angles}\PYG{p}{(}\PYG{n}{angles}\PYG{p}{)}
\PYG{+w}{  }\PYG{n}{angles\PYGZus{}array}\PYG{p}{(}\PYG{l+m+mi}{1}\PYG{p}{:}\PYG{l+m+mi}{4}\PYG{p}{,}\PYG{+w}{ }\PYG{l+m+mi}{1}\PYG{p}{:}\PYG{n}{SIZE}\PYG{p}{(}\PYG{n}{angles}\PYG{p}{)}\PYG{o}{/}\PYG{l+m+mi}{4}\PYG{p}{)}\PYG{+w}{ }\PYG{o}{=}\PYG{o}{\PYGZgt{}}\PYG{+w}{ }\PYG{n}{angles}
\PYG{+w}{  }\PYG{k}{DO }\PYG{n}{i}\PYG{+w}{ }\PYG{o}{=}\PYG{+w}{ }\PYG{l+m+mi}{1}\PYG{p}{,}\PYG{+w}{ }\PYG{n}{SIZE}\PYG{p}{(}\PYG{n}{angles}\PYG{p}{)}\PYG{o}{/}\PYG{l+m+mi}{4}
\PYG{+w}{    }\PYG{k}{WRITE}\PYG{p}{(}\PYG{n}{OUTPUT\PYGZus{}UNIT}\PYG{p}{,}\PYG{l+s+s1}{\PYGZsq{}(A,1X,I4,A,I4,1X,I4,1X,I4)\PYGZsq{}}\PYG{p}{)}\PYG{+w}{ }\PYG{l+s+s1}{\PYGZsq{}angle\PYGZsq{}}\PYG{p}{,}\PYG{+w}{ }\PYG{n}{angles\PYGZus{}array}\PYG{p}{(}\PYG{l+m+mi}{1}\PYG{p}{,}\PYG{n}{i}\PYG{p}{)}\PYG{p}{,}\PYG{+w}{ }\PYG{p}{\PYGZam{}}
\PYG{+w}{      }\PYG{l+s+s1}{\PYGZsq{}; type = \PYGZsq{}}\PYG{p}{,}\PYG{+w}{ }\PYG{n}{angles\PYGZus{}array}\PYG{p}{(}\PYG{l+m+mi}{2}\PYG{p}{,}\PYG{n}{i}\PYG{p}{)}\PYG{p}{,}\PYG{+w}{ }\PYG{n}{angles\PYGZus{}array}\PYG{p}{(}\PYG{l+m+mi}{3}\PYG{p}{,}\PYG{n}{i}\PYG{p}{)}\PYG{p}{,}\PYG{+w}{ }\PYG{n}{angles\PYGZus{}array}\PYG{p}{(}\PYG{l+m+mi}{4}\PYG{p}{,}\PYG{n}{i}\PYG{p}{)}
\PYG{+w}{  }\PYG{k}{END }\PYG{k}{DO}
\PYG{k}{END }\PYG{k}{PROGRAM }\PYG{n}{angles}
\end{sphinxVerbatim}
\begin{quote}\begin{description}
\sphinxlineitem{Parameters}
\sphinxAtStartPar
\sphinxstylestrong{data}\sphinxstyleemphasis{ {[}}\sphinxstyleemphasis{integer(kind=*),allocatable}\sphinxstyleemphasis{{]}} :: array into which to copy the result. *The \sphinxcode{\sphinxupquote{KIND}} parameter is
either \sphinxcode{\sphinxupquote{c\_int}} or, if LAMMPS was compiled with \sphinxcode{\sphinxupquote{\sphinxhyphen{}DLAMMPS\_BIGBIG}},
kind \sphinxcode{\sphinxupquote{c\_int64\_t}}.

\sphinxlineitem{Call to}
\sphinxAtStartPar
{\hyperref[\detokenize{Library_scatter:_CPPv420lammps_gather_anglesPvPv}]{\sphinxcrossref{\sphinxcode{\sphinxupquote{lammps\_gather\_angles()}}}}}

\end{description}\end{quote}

\end{fulllineitems}



\bigskip\hrule\bigskip

\index{gather() (fortran subroutine)@\spxentry{gather()}\spxextra{fortran subroutine}|spxpagem}

\begin{fulllineitems}
\phantomsection\label{\detokenize{Fortran:f/_/gather}}
\pysigstartsignatures
\pysiglinewithargsret{\sphinxbfcode{\sphinxupquote{subroutine  }}\sphinxbfcode{\sphinxupquote{gather}}}{\sphinxparam{self}\sphinxparamcomma \sphinxparam{name}\sphinxparamcomma \sphinxparam{count}\sphinxparamcomma \sphinxparam{data}}{}
\pysigstopsignatures
\sphinxAtStartPar
Gather the named per\sphinxhyphen{}atom, per\sphinxhyphen{}atom fix, per\sphinxhyphen{}atom compute, or fix
property/atom\sphinxhyphen{}based entities from all processes, in order by atom ID.

\sphinxAtStartPar
\DUrole{versionmodified,added}{Added in version 22Dec2022.}

\sphinxAtStartPar
This subroutine gathers data from all processes and stores them in a
one\sphinxhyphen{}dimensional allocatable array. The array \sphinxstyleemphasis{data} will be
ordered by atom ID, which requires consecutive IDs (1 to \sphinxstyleemphasis{natoms}). If you
need a similar array but for non\sphinxhyphen{}consecutive atom IDs, see
{\hyperref[\detokenize{Library_scatter:_CPPv420lammps_gather_concatPvPKciiPv}]{\sphinxcrossref{\sphinxcode{\sphinxupquote{lammps\_gather\_concat()}}}}}; for a similar array but for a subset of
atoms, see {\hyperref[\detokenize{Library_scatter:_CPPv420lammps_gather_subsetPvPKciiiPiPv}]{\sphinxcrossref{\sphinxcode{\sphinxupquote{lammps\_gather\_subset()}}}}}.

\sphinxAtStartPar
The \sphinxstyleemphasis{data} array will be ordered in groups of \sphinxstyleemphasis{count} values, sorted by atom
ID (e.g., if \sphinxstyleemphasis{name} is \sphinxstyleemphasis{x}, then \sphinxstyleemphasis{data} is {[}x(1,1), x(2,1), x(3,1), x(1,2),
x(2,2), x(3,2), x(1,3), \(\dots\){]}); \sphinxstyleemphasis{data} must be \sphinxcode{\sphinxupquote{ALLOCATABLE}} and
will be allocated to length (\sphinxstyleemphasis{count}\({}\times{}\)\sphinxstyleemphasis{natoms}), as
queried by \sphinxcode{\sphinxupquote{get\_natoms()}}.

\sphinxAtStartPar
This function will return an error if fix or compute data are requested and
the fix or compute ID given does not have per\sphinxhyphen{}atom data. See the note about
re\sphinxhyphen{}interpreting the vector as a matrix at \sphinxcode{\sphinxupquote{gather\_atoms()}}.

\sphinxAtStartPar
This function is not compatible with \sphinxcode{\sphinxupquote{\sphinxhyphen{}DLAMMPS\_BIGBIG}}.
\begin{quote}\begin{description}
\sphinxlineitem{Parameters}\begin{itemize}
\item {} 
\sphinxAtStartPar
\sphinxstylestrong{name}\sphinxstyleemphasis{ {[}}\sphinxstyleemphasis{character(len=*)}\sphinxstyleemphasis{{]}} :: desired quantity (e.g., “x” or “mask” for atom
properties, “f\_id” for per\sphinxhyphen{}atom fix data, “c\_id” for per\sphinxhyphen{}atom compute data,
“d\_name” or “i\_name” for fix property/atom vectors with \sphinxstyleemphasis{count} = 1,
“d2\_name” or “i2\_name” for fix property/atom vectors with
\sphinxstyleemphasis{count}\({}> 1\))

\item {} 
\sphinxAtStartPar
\sphinxstylestrong{count}\sphinxstyleemphasis{ {[}}\sphinxstyleemphasis{integer(c\_int)}\sphinxstyleemphasis{{]}} :: number of per\sphinxhyphen{}atom values (e.g., 1 for \sphinxstyleemphasis{type} or
\sphinxstyleemphasis{charge}, 3 for \sphinxstyleemphasis{x} or \sphinxstyleemphasis{f}); use \sphinxstyleemphasis{count} = 3 with \sphinxstyleemphasis{image} if you want the
image flags unpacked into (\sphinxstyleemphasis{x},\sphinxstyleemphasis{y},\sphinxstyleemphasis{z}) components.

\item {} 
\sphinxAtStartPar
\sphinxstylestrong{data}\sphinxstyleemphasis{ {[}}\sphinxstyleemphasis{real(c\_double)}\sphinxstyleemphasis{,}\sphinxstyleemphasis{dimension(:),allocatable}\sphinxstyleemphasis{{]}} :: array into which to store
the data. Array \sphinxstyleemphasis{must} have the \sphinxcode{\sphinxupquote{ALLOCATABLE}} attribute and be of rank 1
(i.e., \sphinxcode{\sphinxupquote{DIMENSION(:)}}). If this array is already allocated, it will be
reallocated to fit the length of the incoming data.

\end{itemize}

\sphinxlineitem{Call to}
\sphinxAtStartPar
{\hyperref[\detokenize{Library_scatter:_CPPv413lammps_gatherPvPKciiPv}]{\sphinxcrossref{\sphinxcode{\sphinxupquote{lammps\_gather()}}}}}

\end{description}\end{quote}

\end{fulllineitems}



\bigskip\hrule\bigskip

\index{gather\_dihedrals() (fortran subroutine)@\spxentry{gather\_dihedrals()}\spxextra{fortran subroutine}|spxpagem}

\begin{fulllineitems}
\phantomsection\label{\detokenize{Fortran:f/_/gather_dihedrals}}
\pysigstartsignatures
\pysiglinewithargsret{\sphinxbfcode{\sphinxupquote{subroutine  }}\sphinxbfcode{\sphinxupquote{gather\_dihedrals}}}{\sphinxparam{data}}{}
\pysigstopsignatures
\sphinxAtStartPar
Gather type and constituent atom information for all dihedrals.

\sphinxAtStartPar
\DUrole{versionmodified,added}{Added in version 8Feb2023.}

\sphinxAtStartPar
This function copies the list of all dihedrals into an allocatable array.
The array will be filled with (dihedral type, dihedral atom 1, dihedral atom 2,
dihedral atom 3, dihedral atom 4) for each dihedral. The array is allocated to
the right length (i.e., five times the number of dihedrals).
The array \sphinxstyleemphasis{data} must be of the same type as the LAMMPS
\sphinxcode{\sphinxupquote{tagint}} type, which is equivalent to either \sphinxcode{\sphinxupquote{INTEGER(c\_int)}} or
\sphinxcode{\sphinxupquote{INTEGER(c\_int64\_t)}}, depending on whether \sphinxcode{\sphinxupquote{\sphinxhyphen{}DLAMMPS\_BIGBIG}} was used
when LAMMPS was built. If the supplied array does not match, an error will
result at run\sphinxhyphen{}time.

\sphinxAtStartPar
An example of how to use this routine is below:

\begin{sphinxVerbatim}[commandchars=\\\{\}]
\PYG{k}{PROGRAM }\PYG{n}{dihedrals}
\PYG{+w}{  }\PYG{k}{USE}\PYG{p}{,}\PYG{+w}{ }\PYG{k}{INTRINSIC}\PYG{+w}{ }\PYG{k+kd}{::}\PYG{+w}{ }\PYG{n+nb}{ISO\PYGZus{}C\PYGZus{}BINDING}\PYG{p}{,}\PYG{+w}{ }\PYG{k}{ONLY}\PYG{+w}{ }\PYG{p}{:}\PYG{+w}{ }\PYG{k+kt}{c\PYGZus{}int}
\PYG{k+kt}{  }\PYG{k}{USE}\PYG{p}{,}\PYG{+w}{ }\PYG{k}{INTRINSIC}\PYG{+w}{ }\PYG{k+kd}{::}\PYG{+w}{ }\PYG{n}{ISO\PYGZus{}FORTRAN\PYGZus{}ENV}\PYG{p}{,}\PYG{+w}{ }\PYG{k}{ONLY}\PYG{+w}{ }\PYG{p}{:}\PYG{+w}{ }\PYG{n}{OUTPUT\PYGZus{}UNIT}
\PYG{+w}{  }\PYG{k}{USE }\PYG{n}{LIBLAMMPS}
\PYG{+w}{  }\PYG{k}{IMPLICIT }\PYG{k}{NONE}
\PYG{k}{  }\PYG{k+kt}{INTEGER}\PYG{p}{(}\PYG{k+kt}{c\PYGZus{}int}\PYG{p}{)}\PYG{p}{,}\PYG{+w}{ }\PYG{k}{DIMENSION}\PYG{p}{(}\PYG{p}{:}\PYG{p}{)}\PYG{p}{,}\PYG{+w}{ }\PYG{k}{ALLOCATABLE}\PYG{p}{,}\PYG{+w}{ }\PYG{k}{TARGET}\PYG{+w}{ }\PYG{k+kd}{::}\PYG{+w}{ }\PYG{n}{dihedrals}
\PYG{+w}{  }\PYG{k+kt}{INTEGER}\PYG{p}{(}\PYG{k+kt}{c\PYGZus{}int}\PYG{p}{)}\PYG{p}{,}\PYG{+w}{ }\PYG{k}{DIMENSION}\PYG{p}{(}\PYG{p}{:}\PYG{p}{,}\PYG{p}{:}\PYG{p}{)}\PYG{p}{,}\PYG{+w}{ }\PYG{k}{POINTER}\PYG{+w}{ }\PYG{k+kd}{::}\PYG{+w}{ }\PYG{n}{dihedrals\PYGZus{}array}
\PYG{+w}{  }\PYG{k}{TYPE}\PYG{p}{(}\PYG{n}{lammps}\PYG{p}{)}\PYG{+w}{ }\PYG{k+kd}{::}\PYG{+w}{ }\PYG{n}{lmp}
\PYG{+w}{  }\PYG{k+kt}{INTEGER}\PYG{+w}{ }\PYG{k+kd}{::}\PYG{+w}{ }\PYG{n}{i}
\PYG{+w}{  }\PYG{c}{! other commands to initialize LAMMPS, create dihedrals, etc.}
\PYG{+w}{  }\PYG{k}{CALL }\PYG{n}{lmp}\PYG{p}{\PYGZpc{}}\PYG{n}{gather\PYGZus{}dihedrals}\PYG{p}{(}\PYG{n}{dihedrals}\PYG{p}{)}
\PYG{+w}{  }\PYG{n}{dihedrals\PYGZus{}array}\PYG{p}{(}\PYG{l+m+mi}{1}\PYG{p}{:}\PYG{l+m+mi}{5}\PYG{p}{,}\PYG{+w}{ }\PYG{l+m+mi}{1}\PYG{p}{:}\PYG{n}{SIZE}\PYG{p}{(}\PYG{n}{dihedrals}\PYG{p}{)}\PYG{o}{/}\PYG{l+m+mi}{5}\PYG{p}{)}\PYG{+w}{ }\PYG{o}{=}\PYG{o}{\PYGZgt{}}\PYG{+w}{ }\PYG{n}{dihedrals}
\PYG{+w}{  }\PYG{k}{DO }\PYG{n}{i}\PYG{+w}{ }\PYG{o}{=}\PYG{+w}{ }\PYG{l+m+mi}{1}\PYG{p}{,}\PYG{+w}{ }\PYG{n}{SIZE}\PYG{p}{(}\PYG{n}{dihedrals}\PYG{p}{)}\PYG{o}{/}\PYG{l+m+mi}{5}
\PYG{+w}{    }\PYG{k}{WRITE}\PYG{p}{(}\PYG{n}{OUTPUT\PYGZus{}UNIT}\PYG{p}{,}\PYG{l+s+s1}{\PYGZsq{}(A,1X,I4,A,I4,1X,I4,1X,I4,1X,I4)\PYGZsq{}}\PYG{p}{)}\PYG{+w}{ }\PYG{l+s+s1}{\PYGZsq{}dihedral\PYGZsq{}}\PYG{p}{,}\PYG{+w}{ }\PYG{n}{dihedrals\PYGZus{}array}\PYG{p}{(}\PYG{l+m+mi}{1}\PYG{p}{,}\PYG{n}{i}\PYG{p}{)}\PYG{p}{,}\PYG{+w}{ }\PYG{p}{\PYGZam{}}
\PYG{+w}{      }\PYG{l+s+s1}{\PYGZsq{}; type = \PYGZsq{}}\PYG{p}{,}\PYG{+w}{ }\PYG{n}{dihedrals\PYGZus{}array}\PYG{p}{(}\PYG{l+m+mi}{2}\PYG{p}{,}\PYG{n}{i}\PYG{p}{)}\PYG{p}{,}\PYG{+w}{ }\PYG{n}{dihedrals\PYGZus{}array}\PYG{p}{(}\PYG{l+m+mi}{3}\PYG{p}{,}\PYG{n}{i}\PYG{p}{)}\PYG{p}{,}\PYG{+w}{ }\PYG{n}{dihedrals\PYGZus{}array}\PYG{p}{(}\PYG{l+m+mi}{4}\PYG{p}{,}\PYG{n}{i}\PYG{p}{)}\PYG{p}{,}\PYG{+w}{ }\PYG{n}{dihedrals\PYGZus{}array}\PYG{p}{(}\PYG{l+m+mi}{5}\PYG{p}{,}\PYG{n}{i}\PYG{p}{)}
\PYG{+w}{  }\PYG{k}{END }\PYG{k}{DO}
\PYG{k}{END }\PYG{k}{PROGRAM }\PYG{n}{dihedrals}
\end{sphinxVerbatim}
\begin{quote}\begin{description}
\sphinxlineitem{Parameters}
\sphinxAtStartPar
\sphinxstylestrong{data}\sphinxstyleemphasis{ {[}}\sphinxstyleemphasis{integer(kind=*),allocatable}\sphinxstyleemphasis{{]}} :: array into which to copy the result. *The \sphinxcode{\sphinxupquote{KIND}} parameter is
either \sphinxcode{\sphinxupquote{c\_int}} or, if LAMMPS was compiled with \sphinxcode{\sphinxupquote{\sphinxhyphen{}DLAMMPS\_BIGBIG}},
kind \sphinxcode{\sphinxupquote{c\_int64\_t}}.

\sphinxlineitem{Call to}
\sphinxAtStartPar
{\hyperref[\detokenize{Library_scatter:_CPPv423lammps_gather_dihedralsPvPv}]{\sphinxcrossref{\sphinxcode{\sphinxupquote{lammps\_gather\_dihedrals()}}}}}

\end{description}\end{quote}

\end{fulllineitems}



\bigskip\hrule\bigskip

\index{gather\_impropers() (fortran subroutine)@\spxentry{gather\_impropers()}\spxextra{fortran subroutine}|spxpagem}

\begin{fulllineitems}
\phantomsection\label{\detokenize{Fortran:f/_/gather_impropers}}
\pysigstartsignatures
\pysiglinewithargsret{\sphinxbfcode{\sphinxupquote{subroutine  }}\sphinxbfcode{\sphinxupquote{gather\_impropers}}}{\sphinxparam{data}}{}
\pysigstopsignatures
\sphinxAtStartPar
Gather type and constituent atom information for all impropers.

\sphinxAtStartPar
\DUrole{versionmodified,added}{Added in version 8Feb2023.}

\sphinxAtStartPar
This function copies the list of all impropers into an allocatable array.
The array will be filled with (improper type, improper atom 1, improper atom 2,
improper atom 3, improper atom 4) for each improper. The array is allocated to
the right length (i.e., five times the number of impropers).
The array \sphinxstyleemphasis{data} must be of the same type as the LAMMPS
\sphinxcode{\sphinxupquote{tagint}} type, which is equivalent to either \sphinxcode{\sphinxupquote{INTEGER(c\_int)}} or
\sphinxcode{\sphinxupquote{INTEGER(c\_int64\_t)}}, depending on whether \sphinxcode{\sphinxupquote{\sphinxhyphen{}DLAMMPS\_BIGBIG}} was used
when LAMMPS was built. If the supplied array does not match, an error will
result at run\sphinxhyphen{}time.

\sphinxAtStartPar
An example of how to use this routine is below:

\begin{sphinxVerbatim}[commandchars=\\\{\}]
\PYG{k}{PROGRAM }\PYG{n}{impropers}
\PYG{+w}{  }\PYG{k}{USE}\PYG{p}{,}\PYG{+w}{ }\PYG{k}{INTRINSIC}\PYG{+w}{ }\PYG{k+kd}{::}\PYG{+w}{ }\PYG{n+nb}{ISO\PYGZus{}C\PYGZus{}BINDING}\PYG{p}{,}\PYG{+w}{ }\PYG{k}{ONLY}\PYG{+w}{ }\PYG{p}{:}\PYG{+w}{ }\PYG{k+kt}{c\PYGZus{}int}
\PYG{k+kt}{  }\PYG{k}{USE}\PYG{p}{,}\PYG{+w}{ }\PYG{k}{INTRINSIC}\PYG{+w}{ }\PYG{k+kd}{::}\PYG{+w}{ }\PYG{n}{ISO\PYGZus{}FORTRAN\PYGZus{}ENV}\PYG{p}{,}\PYG{+w}{ }\PYG{k}{ONLY}\PYG{+w}{ }\PYG{p}{:}\PYG{+w}{ }\PYG{n}{OUTPUT\PYGZus{}UNIT}
\PYG{+w}{  }\PYG{k}{USE }\PYG{n}{LIBLAMMPS}
\PYG{+w}{  }\PYG{k}{IMPLICIT }\PYG{k}{NONE}
\PYG{k}{  }\PYG{k+kt}{INTEGER}\PYG{p}{(}\PYG{k+kt}{c\PYGZus{}int}\PYG{p}{)}\PYG{p}{,}\PYG{+w}{ }\PYG{k}{DIMENSION}\PYG{p}{(}\PYG{p}{:}\PYG{p}{)}\PYG{p}{,}\PYG{+w}{ }\PYG{k}{ALLOCATABLE}\PYG{p}{,}\PYG{+w}{ }\PYG{k}{TARGET}\PYG{+w}{ }\PYG{k+kd}{::}\PYG{+w}{ }\PYG{n}{impropers}
\PYG{+w}{  }\PYG{k+kt}{INTEGER}\PYG{p}{(}\PYG{k+kt}{c\PYGZus{}int}\PYG{p}{)}\PYG{p}{,}\PYG{+w}{ }\PYG{k}{DIMENSION}\PYG{p}{(}\PYG{p}{:}\PYG{p}{,}\PYG{p}{:}\PYG{p}{)}\PYG{p}{,}\PYG{+w}{ }\PYG{k}{POINTER}\PYG{+w}{ }\PYG{k+kd}{::}\PYG{+w}{ }\PYG{n}{impropers\PYGZus{}array}
\PYG{+w}{  }\PYG{k}{TYPE}\PYG{p}{(}\PYG{n}{lammps}\PYG{p}{)}\PYG{+w}{ }\PYG{k+kd}{::}\PYG{+w}{ }\PYG{n}{lmp}
\PYG{+w}{  }\PYG{k+kt}{INTEGER}\PYG{+w}{ }\PYG{k+kd}{::}\PYG{+w}{ }\PYG{n}{i}
\PYG{+w}{  }\PYG{c}{! other commands to initialize LAMMPS, create impropers, etc.}
\PYG{+w}{  }\PYG{k}{CALL }\PYG{n}{lmp}\PYG{p}{\PYGZpc{}}\PYG{n}{gather\PYGZus{}impropers}\PYG{p}{(}\PYG{n}{impropers}\PYG{p}{)}
\PYG{+w}{  }\PYG{n}{impropers\PYGZus{}array}\PYG{p}{(}\PYG{l+m+mi}{1}\PYG{p}{:}\PYG{l+m+mi}{5}\PYG{p}{,}\PYG{+w}{ }\PYG{l+m+mi}{1}\PYG{p}{:}\PYG{n}{SIZE}\PYG{p}{(}\PYG{n}{impropers}\PYG{p}{)}\PYG{o}{/}\PYG{l+m+mi}{5}\PYG{p}{)}\PYG{+w}{ }\PYG{o}{=}\PYG{o}{\PYGZgt{}}\PYG{+w}{ }\PYG{n}{impropers}
\PYG{+w}{  }\PYG{k}{DO }\PYG{n}{i}\PYG{+w}{ }\PYG{o}{=}\PYG{+w}{ }\PYG{l+m+mi}{1}\PYG{p}{,}\PYG{+w}{ }\PYG{n}{SIZE}\PYG{p}{(}\PYG{n}{impropers}\PYG{p}{)}\PYG{o}{/}\PYG{l+m+mi}{5}
\PYG{+w}{    }\PYG{k}{WRITE}\PYG{p}{(}\PYG{n}{OUTPUT\PYGZus{}UNIT}\PYG{p}{,}\PYG{l+s+s1}{\PYGZsq{}(A,1X,I4,A,I4,1X,I4,1X,I4,1X,I4)\PYGZsq{}}\PYG{p}{)}\PYG{+w}{ }\PYG{l+s+s1}{\PYGZsq{}improper\PYGZsq{}}\PYG{p}{,}\PYG{+w}{ }\PYG{n}{impropers\PYGZus{}array}\PYG{p}{(}\PYG{l+m+mi}{1}\PYG{p}{,}\PYG{n}{i}\PYG{p}{)}\PYG{p}{,}\PYG{+w}{ }\PYG{p}{\PYGZam{}}
\PYG{+w}{      }\PYG{l+s+s1}{\PYGZsq{}; type = \PYGZsq{}}\PYG{p}{,}\PYG{+w}{ }\PYG{n}{impropers\PYGZus{}array}\PYG{p}{(}\PYG{l+m+mi}{2}\PYG{p}{,}\PYG{n}{i}\PYG{p}{)}\PYG{p}{,}\PYG{+w}{ }\PYG{n}{impropers\PYGZus{}array}\PYG{p}{(}\PYG{l+m+mi}{3}\PYG{p}{,}\PYG{n}{i}\PYG{p}{)}\PYG{p}{,}\PYG{+w}{ }\PYG{n}{impropers\PYGZus{}array}\PYG{p}{(}\PYG{l+m+mi}{4}\PYG{p}{,}\PYG{n}{i}\PYG{p}{)}\PYG{p}{,}\PYG{+w}{ }\PYG{n}{impropers\PYGZus{}array}\PYG{p}{(}\PYG{l+m+mi}{5}\PYG{p}{,}\PYG{n}{i}\PYG{p}{)}
\PYG{+w}{  }\PYG{k}{END }\PYG{k}{DO}
\PYG{k}{END }\PYG{k}{PROGRAM }\PYG{n}{impropers}
\end{sphinxVerbatim}
\begin{quote}\begin{description}
\sphinxlineitem{Parameters}
\sphinxAtStartPar
\sphinxstylestrong{data}\sphinxstyleemphasis{ {[}}\sphinxstyleemphasis{integer(kind=*),allocatable}\sphinxstyleemphasis{{]}} :: array into which to copy the result. *The \sphinxcode{\sphinxupquote{KIND}} parameter is
either \sphinxcode{\sphinxupquote{c\_int}} or, if LAMMPS was compiled with \sphinxcode{\sphinxupquote{\sphinxhyphen{}DLAMMPS\_BIGBIG}},
kind \sphinxcode{\sphinxupquote{c\_int64\_t}}.

\sphinxlineitem{Call to}
\sphinxAtStartPar
{\hyperref[\detokenize{Library_scatter:_CPPv423lammps_gather_impropersPvPv}]{\sphinxcrossref{\sphinxcode{\sphinxupquote{lammps\_gather\_impropers()}}}}}

\end{description}\end{quote}

\end{fulllineitems}



\bigskip\hrule\bigskip

\index{gather\_concat() (fortran subroutine)@\spxentry{gather\_concat()}\spxextra{fortran subroutine}|spxpagem}

\begin{fulllineitems}
\phantomsection\label{\detokenize{Fortran:f/_/gather_concat}}
\pysigstartsignatures
\pysiglinewithargsret{\sphinxbfcode{\sphinxupquote{subroutine  }}\sphinxbfcode{\sphinxupquote{gather\_concat}}}{\sphinxparam{self}\sphinxparamcomma \sphinxparam{name}\sphinxparamcomma \sphinxparam{count}\sphinxparamcomma \sphinxparam{data}}{}
\pysigstopsignatures
\sphinxAtStartPar
Gather the named per\sphinxhyphen{}atom, per\sphinxhyphen{}atom fix, per\sphinxhyphen{}atom compute, or fix
property/atom\sphinxhyphen{}based entities from all processes, unordered.

\sphinxAtStartPar
\DUrole{versionmodified,added}{Added in version 22Dec2022.}

\sphinxAtStartPar
This subroutine gathers data for all atoms and stores them in a
one\sphinxhyphen{}dimensional allocatable array. The data will be a
concatenation of chunks from each processor’s owned atoms, in whatever order
the atoms are in on each processor. This process has no requirement that the
atom IDs be consecutive. If you need the ID of each atom, you can do another
call to either \sphinxcode{\sphinxupquote{gather\_atoms\_concat()}} or \sphinxcode{\sphinxupquote{gather\_concat()}} with
\sphinxstyleemphasis{name} set to \sphinxcode{\sphinxupquote{id}}. If you have consecutive IDs and want the data to be in
order, use \sphinxcode{\sphinxupquote{gather()}}; for a similar array but for a subset of
atoms, use \sphinxcode{\sphinxupquote{gather\_subset()}}.

\sphinxAtStartPar
The \sphinxstyleemphasis{data} array will be in groups of \sphinxstyleemphasis{count} values, with \sphinxstyleemphasis{natoms} groups
total, but not in order by atom ID (e.g., if \sphinxstyleemphasis{name} is \sphinxstyleemphasis{x} and \sphinxstyleemphasis{count} is 3,
then \sphinxstyleemphasis{data} might be something like {[}x(1,11), x(2,11), x(3,11), x(1,3),
x(2,3), x(3,3), x(1,5), \(\dots\){]}); \sphinxstyleemphasis{data} must be \sphinxcode{\sphinxupquote{ALLOCATABLE}} and
will be allocated to length (\sphinxstyleemphasis{count} \(\times\) \sphinxstyleemphasis{natoms}), as queried by
\sphinxcode{\sphinxupquote{get\_natoms()}}.

\sphinxAtStartPar
This function is not compatible with \sphinxcode{\sphinxupquote{\sphinxhyphen{}DLAMMPS\_BIGBIG}}.
\begin{quote}\begin{description}
\sphinxlineitem{Parameters}\begin{itemize}
\item {} 
\sphinxAtStartPar
\sphinxstylestrong{name}\sphinxstyleemphasis{ {[}}\sphinxstyleemphasis{character(len=*)}\sphinxstyleemphasis{{]}} :: desired quantity (e.g., “x” or “mask” for atom
properties, “f\_id” for per\sphinxhyphen{}atom fix data, “c\_id” for per\sphinxhyphen{}atom compute data,
“d\_name” or “i\_name” for fix property/atom vectors with \sphinxstyleemphasis{count} = 1,
“d2\_name” or “i2\_name” for fix property/atom vectors with
\sphinxstyleemphasis{count}\({}> 1\))

\item {} 
\sphinxAtStartPar
\sphinxstylestrong{count}\sphinxstyleemphasis{ {[}}\sphinxstyleemphasis{integer(c\_int)}\sphinxstyleemphasis{{]}} :: number of per\sphinxhyphen{}atom values you expect per atom
(e.g., 1 for \sphinxstyleemphasis{type}, \sphinxstyleemphasis{mask}, or \sphinxstyleemphasis{charge}; 3 for \sphinxstyleemphasis{x}, \sphinxstyleemphasis{v}, or \sphinxstyleemphasis{f}). Use
\sphinxstyleemphasis{count} = 3 with \sphinxstyleemphasis{image} if you want a single image flag unpacked into
\sphinxstyleemphasis{x}/\sphinxstyleemphasis{y}/\sphinxstyleemphasis{z} components.

\item {} 
\sphinxAtStartPar
\sphinxstylestrong{data}\sphinxstyleemphasis{ {[}}\sphinxstyleemphasis{polymorphic,dimension(:),allocatable}\sphinxstyleemphasis{{]}} :: array into which to store
the data. Array \sphinxstyleemphasis{must} have the \sphinxcode{\sphinxupquote{ALLOCATABLE}} attribute and be of rank 1
(i.e., \sphinxcode{\sphinxupquote{DIMENSION(:)}}). If this array is already allocated, it will be
reallocated to fit the length of the incoming data. It should have type
\sphinxcode{\sphinxupquote{INTEGER(c\_int)}} if expecting integer data and \sphinxcode{\sphinxupquote{REAL(c\_double)}} if
expecting floating\sphinxhyphen{}point data.

\end{itemize}

\sphinxlineitem{Call to}
\sphinxAtStartPar
{\hyperref[\detokenize{Library_scatter:_CPPv420lammps_gather_concatPvPKciiPv}]{\sphinxcrossref{\sphinxcode{\sphinxupquote{lammps\_gather\_concat()}}}}}

\end{description}\end{quote}

\end{fulllineitems}



\bigskip\hrule\bigskip

\index{gather\_subset() (fortran subroutine)@\spxentry{gather\_subset()}\spxextra{fortran subroutine}|spxpagem}

\begin{fulllineitems}
\phantomsection\label{\detokenize{Fortran:f/_/gather_subset}}
\pysigstartsignatures
\pysiglinewithargsret{\sphinxbfcode{\sphinxupquote{subroutine  }}\sphinxbfcode{\sphinxupquote{gather\_subset}}}{\sphinxparam{name}\sphinxparamcomma \sphinxparam{count}\sphinxparamcomma \sphinxparam{ids}\sphinxparamcomma \sphinxparam{data}}{}
\pysigstopsignatures
\sphinxAtStartPar
Gather the named per\sphinxhyphen{}atom, per\sphinxhyphen{}atom fix, per\sphinxhyphen{}atom compute, or fix
property/atom\sphinxhyphen{}based entities from all processes for a subset of atoms.

\sphinxAtStartPar
\DUrole{versionmodified,added}{Added in version 22Dec2022.}

\sphinxAtStartPar
This subroutine gathers data for the requested atom IDs and stores them in a
one\sphinxhyphen{}dimensional allocatable array. The data will be ordered by atom ID, but
there is no requirement that the IDs be consecutive. If you wish to return a
similar array for \sphinxstyleemphasis{all} the atoms, use \sphinxcode{\sphinxupquote{gather()}} or
\sphinxcode{\sphinxupquote{gather\_concat()}}.

\sphinxAtStartPar
The \sphinxstyleemphasis{data} array will be in groups of \sphinxstyleemphasis{count} values, sorted by atom ID in
the same order as the array \sphinxstyleemphasis{ids} (e.g., if \sphinxstyleemphasis{name} is \sphinxstyleemphasis{x}, \sphinxstyleemphasis{count} = 3, and
\sphinxstyleemphasis{ids} is {[}100, 57, 210{]}, then \sphinxstyleemphasis{data} might look like {[}\sphinxstyleemphasis{x}(1,100),
\sphinxstyleemphasis{x}(2,100), \sphinxstyleemphasis{x}(3,100), \sphinxstyleemphasis{x}(1,57), \sphinxstyleemphasis{x}(2,57), \sphinxstyleemphasis{x}(3,57),
\sphinxstyleemphasis{x}(1,210), \(\dots\){]}); \sphinxstyleemphasis{ids} must be provided by the user, and
\sphinxstyleemphasis{data} must have the \sphinxcode{\sphinxupquote{ALLOCATABLE}} attribute and be of rank 1 (i.e.,
\sphinxcode{\sphinxupquote{DIMENSION(:)}}). If \sphinxstyleemphasis{data} is already allocated, it will be reallocated to
fit the length of the incoming data.

\sphinxAtStartPar
This function is not compatible with \sphinxcode{\sphinxupquote{\sphinxhyphen{}DLAMMPS\_BIGBIG}}.
\begin{quote}\begin{description}
\sphinxlineitem{Parameters}\begin{itemize}
\item {} 
\sphinxAtStartPar
\sphinxstylestrong{name}\sphinxstyleemphasis{ {[}}\sphinxstyleemphasis{character(len=*)}\sphinxstyleemphasis{{]}} :: quantity to be scattered

\item {} 
\sphinxAtStartPar
\sphinxstylestrong{ids}\sphinxstyleemphasis{ {[}}\sphinxstyleemphasis{integer(c\_int)}\sphinxstyleemphasis{,}\sphinxstyleemphasis{dimension(:)}\sphinxstyleemphasis{{]}} :: atom IDs corresponding to the atoms
being scattered (e.g., “x” or “f” for atom properties, “f\_id” for per\sphinxhyphen{}atom
fix data, “c\_id” for per\sphinxhyphen{}atom compute data, “d\_name” or “i\_name” for fix
property/atom vectors with \sphinxstyleemphasis{count} = 1, “d2\_name” or “i2\_name” for fix
property/atom vectors with \sphinxstyleemphasis{count}\({} > 1\))

\item {} 
\sphinxAtStartPar
\sphinxstylestrong{count}\sphinxstyleemphasis{ {[}}\sphinxstyleemphasis{integer(c\_int)}\sphinxstyleemphasis{{]}} :: number of per\sphinxhyphen{}atom values you expect per atom
(e.g., 1 for \sphinxstyleemphasis{type}, \sphinxstyleemphasis{mask}, or \sphinxstyleemphasis{charge}; 3 for \sphinxstyleemphasis{x}, \sphinxstyleemphasis{v}, or \sphinxstyleemphasis{f}). Use
\sphinxstyleemphasis{count} = 3 with \sphinxstyleemphasis{image} if you want a single image flag unpacked into
\sphinxstyleemphasis{x}/\sphinxstyleemphasis{y}/\sphinxstyleemphasis{z} components.

\item {} 
\sphinxAtStartPar
\sphinxstylestrong{data}\sphinxstyleemphasis{ {[}}\sphinxstyleemphasis{polymorphic,dimension(:),allocatable}\sphinxstyleemphasis{{]}} :: per\sphinxhyphen{}atom values packed into a one\sphinxhyphen{}dimensional array containing the
data to be scattered. This array must have the \sphinxcode{\sphinxupquote{ALLOCATABLE}} attribute
and will be allocated either to the same length as \sphinxstyleemphasis{ids}
(for \sphinxstyleemphasis{mask}, \sphinxstyleemphasis{type}, etc.) or to three times its length (for \sphinxstyleemphasis{x}, \sphinxstyleemphasis{f},
etc.); the array must be rank 1 and be of type \sphinxcode{\sphinxupquote{INTEGER(c\_int)}} (e.g.,
for \sphinxstyleemphasis{mask} or \sphinxstyleemphasis{type}) or of type \sphinxcode{\sphinxupquote{REAL(c\_double)}} (e.g., for \sphinxstyleemphasis{charge},
\sphinxstyleemphasis{x}, or \sphinxstyleemphasis{f}).

\end{itemize}

\sphinxlineitem{Call to}
\sphinxAtStartPar
{\hyperref[\detokenize{Library_scatter:_CPPv420lammps_gather_subsetPvPKciiiPiPv}]{\sphinxcrossref{\sphinxcode{\sphinxupquote{lammps\_gather\_subset()}}}}}

\end{description}\end{quote}

\end{fulllineitems}



\bigskip\hrule\bigskip

\index{scatter() (fortran subroutine)@\spxentry{scatter()}\spxextra{fortran subroutine}|spxpagem}

\begin{fulllineitems}
\phantomsection\label{\detokenize{Fortran:f/_/scatter}}
\pysigstartsignatures
\pysiglinewithargsret{\sphinxbfcode{\sphinxupquote{subroutine  }}\sphinxbfcode{\sphinxupquote{scatter}}}{\sphinxparam{name}\sphinxparamcomma \sphinxparam{data}}{}
\pysigstopsignatures
\sphinxAtStartPar
This function calls {\hyperref[\detokenize{Library_scatter:_CPPv414lammps_scatterPvPKciiPv}]{\sphinxcrossref{\sphinxcode{\sphinxupquote{lammps\_scatter()}}}}} to scatter the named
per\sphinxhyphen{}atom, per\sphinxhyphen{}atom fix, per\sphinxhyphen{}atom compute, or fix property/atom\sphinxhyphen{}based entity
in \sphinxstyleemphasis{data} to all processes.

\sphinxAtStartPar
\DUrole{versionmodified,added}{Added in version 22Dec2022.}

\sphinxAtStartPar
This subroutine takes data stored in a one\sphinxhyphen{}dimensional array supplied by the
user and scatters them to all atoms on all processes. The data must be
ordered by atom ID, with the requirement that the IDs be consecutive. Use
\sphinxcode{\sphinxupquote{scatter\_subset()}} to scatter data for some (or all) atoms, unordered.

\sphinxAtStartPar
The \sphinxstyleemphasis{data} array needs to be ordered in groups of \sphinxstyleemphasis{count} values, sorted by
atom ID (e.g., if \sphinxstyleemphasis{name} is \sphinxstyleemphasis{x} and \sphinxstyleemphasis{count} = 3, then \sphinxstyleemphasis{data} = {[}\sphinxstyleemphasis{x}(1,1),
\sphinxstyleemphasis{x}(2,1), \sphinxstyleemphasis{x}(3,1), \sphinxstyleemphasis{x}(1,2), \sphinxstyleemphasis{x}(2,2), \sphinxstyleemphasis{x}(3,2), \sphinxstyleemphasis{x}(1,3),
\(\dots\){]}); \sphinxstyleemphasis{data} must be of length (\sphinxstyleemphasis{count} \(\times\) \sphinxstyleemphasis{natoms}).

\sphinxAtStartPar
This function is not compatible with \sphinxcode{\sphinxupquote{\sphinxhyphen{}DLAMMPS\_BIGBIG}}.
\begin{quote}\begin{description}
\sphinxlineitem{Parameters}\begin{itemize}
\item {} 
\sphinxAtStartPar
\sphinxstylestrong{name}\sphinxstyleemphasis{ {[}}\sphinxstyleemphasis{character(len=*)}\sphinxstyleemphasis{{]}} :: desired quantity (e.g., “x” or “f” for atom
properties, “f\_id” for per\sphinxhyphen{}atom fix data, “c\_id” for per\sphinxhyphen{}atom compute data,
“d\_name” or “i\_name” for fix property/atom vectors with \sphinxstyleemphasis{count} = 1,
“d2\_name” or “i2\_name” for fix property/atom vectors with \sphinxstyleemphasis{count}\({} > 1\))

\item {} 
\sphinxAtStartPar
\sphinxstylestrong{data}\sphinxstyleemphasis{ {[}}\sphinxstyleemphasis{polymorphic,dimension(:)}\sphinxstyleemphasis{{]}} :: per\sphinxhyphen{}atom values packed in a one\sphinxhyphen{}dimensional array; \sphinxstyleemphasis{data} should be
of type \sphinxcode{\sphinxupquote{INTEGER(c\_int)}} or \sphinxcode{\sphinxupquote{REAL(c\_double)}}, depending on the type of
data being scattered, and be of rank 1 (i.e., \sphinxcode{\sphinxupquote{DIMENSION(:)}}).

\end{itemize}

\sphinxlineitem{Call to}
\sphinxAtStartPar
{\hyperref[\detokenize{Library_scatter:_CPPv414lammps_scatterPvPKciiPv}]{\sphinxcrossref{\sphinxcode{\sphinxupquote{lammps\_scatter()}}}}}

\end{description}\end{quote}

\end{fulllineitems}



\bigskip\hrule\bigskip

\index{scatter\_subset() (fortran subroutine)@\spxentry{scatter\_subset()}\spxextra{fortran subroutine}|spxpagem}

\begin{fulllineitems}
\phantomsection\label{\detokenize{Fortran:f/_/scatter_subset}}
\pysigstartsignatures
\pysiglinewithargsret{\sphinxbfcode{\sphinxupquote{subroutine  }}\sphinxbfcode{\sphinxupquote{scatter\_subset}}}{\sphinxparam{name}\sphinxparamcomma \sphinxparam{ids}\sphinxparamcomma \sphinxparam{data}}{}
\pysigstopsignatures
\sphinxAtStartPar
This function calls {\hyperref[\detokenize{Library_scatter:_CPPv421lammps_scatter_subsetPvPKciiiPiPv}]{\sphinxcrossref{\sphinxcode{\sphinxupquote{lammps\_scatter\_subset()}}}}} to scatter the named
per\sphinxhyphen{}atom, per\sphinxhyphen{}atom fix, per\sphinxhyphen{}atom compute, or fix property/atom\sphinxhyphen{}based
entities in \sphinxstyleemphasis{data} from a subset of atoms to all processes.

\sphinxAtStartPar
\DUrole{versionmodified,added}{Added in version 22Dec2022.}

\sphinxAtStartPar
This subroutine takes data stored in a one\sphinxhyphen{}dimensional array supplied by the
user and scatters them to a subset of atoms on all processes. The array
\sphinxstyleemphasis{data} contains data associated with atom IDs, but there is no requirement
that the IDs be consecutive, as they are provided in a separate array.
Use \sphinxcode{\sphinxupquote{scatter()}} to scatter data for all atoms, in order.

\sphinxAtStartPar
The \sphinxstyleemphasis{data} array needs to be organized in groups of \sphinxstyleemphasis{count} values, with the
groups in the same order as the array \sphinxstyleemphasis{ids}. For example, if you want \sphinxstyleemphasis{data}
to be the array {[}x(1,1), x(2,1), x(3,1), x(1,100), x(2,100), x(3,100),
x(1,57), x(2,57), x(3,57){]}, then \sphinxstyleemphasis{count} = 3 and \sphinxstyleemphasis{ids} = {[}1, 100, 57{]}.

\sphinxAtStartPar
This function is not compatible with \sphinxcode{\sphinxupquote{\sphinxhyphen{}DLAMMPS\_BIGBIG}}.
\begin{quote}\begin{description}
\sphinxlineitem{Parameters}\begin{itemize}
\item {} 
\sphinxAtStartPar
\sphinxstylestrong{name}\sphinxstyleemphasis{ {[}}\sphinxstyleemphasis{character(len=*)}\sphinxstyleemphasis{{]}} :: desired quantity (e.g., “x” or “mask” for atom
properties, “f\_id” for per\sphinxhyphen{}atom fix data, “c\_id” for per\sphinxhyphen{}atom compute data,
“d\_name” or “i\_name” for fix property/atom vectors with \sphinxstyleemphasis{count} = 1,
“d2\_name” or “i2\_name” for fix property/atom vectors with
\sphinxstyleemphasis{count}\({}> 1\))

\item {} 
\sphinxAtStartPar
\sphinxstylestrong{ids}\sphinxstyleemphasis{ {[}}\sphinxstyleemphasis{integer(c\_int)}\sphinxstyleemphasis{{]}} :: list of atom IDs to scatter data for

\item {} 
\sphinxAtStartPar
\sphinxstylestrong{data}\sphinxstyleemphasis{ {[}}\sphinxstyleemphasis{polymorphic }\sphinxstyleemphasis{,}\sphinxstyleemphasis{dimension(:)}\sphinxstyleemphasis{{]}} :: per\sphinxhyphen{}atom values packed in a
one\sphinxhyphen{}dimensional array of length \sphinxstyleemphasis{size(ids)} * \sphinxstyleemphasis{count}.

\end{itemize}

\sphinxlineitem{Call to}
\sphinxAtStartPar
{\hyperref[\detokenize{Library_scatter:_CPPv421lammps_scatter_subsetPvPKciiiPiPv}]{\sphinxcrossref{\sphinxcode{\sphinxupquote{lammps\_scatter\_subset()}}}}}

\end{description}\end{quote}

\end{fulllineitems}



\bigskip\hrule\bigskip

\index{create\_atoms() (fortran subroutine)@\spxentry{create\_atoms()}\spxextra{fortran subroutine}|spxpagem}

\begin{fulllineitems}
\phantomsection\label{\detokenize{Fortran:f/_/create_atoms}}
\pysigstartsignatures
\pysiglinewithargsret{\sphinxbfcode{\sphinxupquote{subroutine  }}\sphinxbfcode{\sphinxupquote{create\_atoms}}}{\sphinxparam{{[}id,{]} type, x, {[}v,{]} {[}image,{]} {[}bexpand{]}}}{}
\pysigstopsignatures
\sphinxAtStartPar
This method calls {\hyperref[\detokenize{Library_scatter:_CPPv419lammps_create_atomsPviPKiPKiPKdPKdPKii}]{\sphinxcrossref{\sphinxcode{\sphinxupquote{lammps\_create\_atoms()}}}}} to create additional atoms
from a given list of coordinates and a list of atom types. Additionally,
the atom IDs, velocities, and image flags may be provided.

\sphinxAtStartPar
\DUrole{versionmodified,added}{Added in version 3Nov2022.}
\begin{quote}\begin{description}
\sphinxlineitem{Parameters}\begin{itemize}
\item {} 
\sphinxAtStartPar
\sphinxstylestrong{type}\sphinxstyleemphasis{ {[}}\sphinxstyleemphasis{integer(c\_int)}\sphinxstyleemphasis{,}\sphinxstyleemphasis{dimension(N)}\sphinxstyleemphasis{{]}} :: vector of \(N\) atom types
(required/see note below)

\item {} 
\sphinxAtStartPar
\sphinxstylestrong{x}\sphinxstyleemphasis{ {[}}\sphinxstyleemphasis{real(c\_double)}\sphinxstyleemphasis{,}\sphinxstyleemphasis{dimension(3N)}\sphinxstyleemphasis{{]}} :: vector of \(3N\ x/y/z\) positions
of the new atoms, arranged as \([x_1,y_1,z_1,x_2,y_2,\dotsc]\)
(required/see note below)

\end{itemize}

\sphinxlineitem{Options}\begin{itemize}
\item {} 
\sphinxAtStartPar
\sphinxstylestrong{id}\sphinxstyleemphasis{ {[}}\sphinxstyleemphasis{integer(kind=*)}\sphinxstyleemphasis{,}\sphinxstyleemphasis{dimension(N),optional}\sphinxstyleemphasis{{]}} :: vector of \(N\) atom
IDs; if absent, LAMMPS will generate them for you. *The \sphinxcode{\sphinxupquote{KIND}} parameter
should be \sphinxcode{\sphinxupquote{c\_int}} unless LAMMPS was compiled with \sphinxcode{\sphinxupquote{\sphinxhyphen{}DLAMMPS\_BIGBIG}}, in
which case it should be \sphinxcode{\sphinxupquote{c\_int64\_t}}.

\item {} 
\sphinxAtStartPar
\sphinxstylestrong{v}\sphinxstyleemphasis{ {[}}\sphinxstyleemphasis{real(c\_double)}\sphinxstyleemphasis{,}\sphinxstyleemphasis{dimension(3N),optional}\sphinxstyleemphasis{{]}} :: vector of \(3N\)
\sphinxstyleemphasis{x}/\sphinxstyleemphasis{y}/\sphinxstyleemphasis{z} velocities of the new atoms, arranged as
\([v_{1,x},v_{1,y},v_{1,z},v_{2,x}, \dotsc]\); if absent, they will be
set to zero

\item {} 
\sphinxAtStartPar
\sphinxstylestrong{image}\sphinxstyleemphasis{ {[}}\sphinxstyleemphasis{integer(kind=*)}\sphinxstyleemphasis{,}\sphinxstyleemphasis{dimension(N),optional}\sphinxstyleemphasis{{]}} :: vector of \(N\) image
flags; if absent, they are set to zero. *The \sphinxcode{\sphinxupquote{KIND}} parameter should be
\sphinxcode{\sphinxupquote{c\_int}} unless LAMMPS was compiled with \sphinxcode{\sphinxupquote{\sphinxhyphen{}DLAMMPS\_BIGBIG}}, in which
case it should be \sphinxcode{\sphinxupquote{c\_int64\_t}}. See note below.

\item {} 
\sphinxAtStartPar
\sphinxstylestrong{bexpand}\sphinxstyleemphasis{ {[}}\sphinxstyleemphasis{logical,optional}\sphinxstyleemphasis{{]}} :: if \sphinxcode{\sphinxupquote{.TRUE.}}, atoms outside of shrink\sphinxhyphen{}wrap boundaries
will be created, not dropped, and the box dimensions will be extended.
Default is \sphinxcode{\sphinxupquote{.FALSE.}}

\end{itemize}

\sphinxlineitem{Call to}
\sphinxAtStartPar
{\hyperref[\detokenize{Library_scatter:_CPPv419lammps_create_atomsPviPKiPKiPKdPKdPKii}]{\sphinxcrossref{\sphinxcode{\sphinxupquote{lammps\_create\_atoms()}}}}}

\end{description}\end{quote}

\begin{sphinxadmonition}{note}{Note:}
\sphinxAtStartPar
The \sphinxstyleemphasis{type} and \sphinxstyleemphasis{x} arguments are required, but they are declared
\sphinxcode{\sphinxupquote{OPTIONAL}} in the module because making them mandatory would require
\sphinxstyleemphasis{id} to be present as well. To have LAMMPS generate the ids for you,
use a call something like

\begin{sphinxVerbatim}[commandchars=\\\{\}]
\PYG{n}{lmp}\PYG{p}{\PYGZpc{}}\PYG{n}{create\PYGZus{}atoms}\PYG{p}{(}\PYG{k}{type}\PYG{o}{=}\PYG{n}{new\PYGZus{}types}\PYG{p}{,}\PYG{+w}{ }\PYG{n}{x}\PYG{o}{=}\PYG{n}{new\PYGZus{}xs}\PYG{p}{)}
\end{sphinxVerbatim}
\end{sphinxadmonition}

\begin{sphinxadmonition}{note}{Note:}
\sphinxAtStartPar
When LAMMPS has been compiled with \sphinxcode{\sphinxupquote{\sphinxhyphen{}DLAMMPS\_BIGBIG}}, it is not
possible to include the \sphinxstyleemphasis{image} parameter but omit the \sphinxstyleemphasis{id} parameter.
Either \sphinxstyleemphasis{id} must be present, or both \sphinxstyleemphasis{id} and \sphinxstyleemphasis{image} must be absent.
This is required because having all arguments be optional in both
generic functions creates an ambiguous interface. This limitation does
not exist if LAMMPS was not compiled with \sphinxcode{\sphinxupquote{\sphinxhyphen{}DLAMMPS\_BIGBIG}}.
\end{sphinxadmonition}

\end{fulllineitems}



\bigskip\hrule\bigskip

\index{find\_pair\_neighlist() (fortran function)@\spxentry{find\_pair\_neighlist()}\spxextra{fortran function}|spxpagem}

\begin{fulllineitems}
\phantomsection\label{\detokenize{Fortran:f/_/find_pair_neighlist}}
\pysigstartsignatures
\pysiglinewithargsret{\sphinxbfcode{\sphinxupquote{function  }}\sphinxbfcode{\sphinxupquote{find\_pair\_neighlist}}}{\sphinxparam{style{[}, exact{]}{[}, nsub{]}{[}, reqid{]}}}{}
\pysigstopsignatures
\sphinxAtStartPar
Find index of a neighbor list requested by a pair style.

\sphinxAtStartPar
\DUrole{versionmodified,added}{Added in version 3Nov2022.}

\sphinxAtStartPar
This function determines which of the available neighbor lists for pair
styles matches the given conditions.  It first matches the style name.
If \sphinxstyleemphasis{exact} is \sphinxcode{\sphinxupquote{.TRUE.}}, the name must match exactly; if \sphinxcode{\sphinxupquote{.FALSE.}}, a
regular expression or sub\sphinxhyphen{}string match is done.  If the pair style is
\sphinxstyleemphasis{hybrid} or \sphinxstyleemphasis{hybrid/overlay}, the style is matched against the sub\sphinxhyphen{}styles
instead. If the same pair style is used multiple times as a sub\sphinxhyphen{}style, the
\sphinxstyleemphasis{nsub} argument must be \(> 0\); this argument represents the \sphinxstyleemphasis{n}th
instance of the sub\sphinxhyphen{}style (same as for the {\hyperref[\detokenize{pair_coeff::doc}]{\sphinxcrossref{\DUrole{doc}{pair\_coeff}}}}
command, for example). In that case, \sphinxstyleemphasis{nsub}\({} = 0\) will not
produce a match, and the function will return \(-1\).

\sphinxAtStartPar
The final condition to be checked is the request ID (\sphinxstyleemphasis{reqid}). This
will usually be zero, but some pair styles request multiple neighbor
lists and set the request ID to a value greater than zero.
\begin{quote}\begin{description}
\sphinxlineitem{Parameters}
\sphinxAtStartPar
\sphinxstylestrong{style}\sphinxstyleemphasis{ {[}}\sphinxstyleemphasis{character(len=*)}\sphinxstyleemphasis{{]}} :: String used to search for pair style instance.

\sphinxlineitem{Options}\begin{itemize}
\item {} 
\sphinxAtStartPar
\sphinxstylestrong{exact}\sphinxstyleemphasis{ {[}}\sphinxstyleemphasis{logical,optional}\sphinxstyleemphasis{{]}} :: Flag to control whether style should match exactly or only a
regular expression/sub\sphinxhyphen{}string match is applied. Default: \sphinxcode{\sphinxupquote{.TRUE.}}.

\item {} 
\sphinxAtStartPar
\sphinxstylestrong{nsub}\sphinxstyleemphasis{ {[}}\sphinxstyleemphasis{integer(c\_int)}\sphinxstyleemphasis{,}\sphinxstyleemphasis{optional}\sphinxstyleemphasis{{]}} :: Match \sphinxstyleemphasis{nsub}th hybrid sub\sphinxhyphen{}style
instance of the same style. Default: 0.

\item {} 
\sphinxAtStartPar
\sphinxstylestrong{reqid}\sphinxstyleemphasis{ {[}}\sphinxstyleemphasis{integer(c\_int)}\sphinxstyleemphasis{,}\sphinxstyleemphasis{optional}\sphinxstyleemphasis{{]}} :: Request ID to identify the neighbor list
in case there are multiple requests from the same pair style instance.
Default: 0.

\end{itemize}

\sphinxlineitem{Call to}
\sphinxAtStartPar
{\hyperref[\detokenize{Library_neighbor:_CPPv426lammps_find_pair_neighlistPvPKciii}]{\sphinxcrossref{\sphinxcode{\sphinxupquote{lammps\_find\_pair\_neighlist()}}}}}

\sphinxlineitem{Return}
\sphinxAtStartPar
\sphinxstylestrong{index}\sphinxstyleemphasis{ {[}}\sphinxstyleemphasis{integer(c\_int)}\sphinxstyleemphasis{{]}} :: Neighbor list index if found, otherwise
\(-1\).

\end{description}\end{quote}

\end{fulllineitems}



\bigskip\hrule\bigskip

\index{find\_fix\_neighlist() (fortran function)@\spxentry{find\_fix\_neighlist()}\spxextra{fortran function}|spxpagem}

\begin{fulllineitems}
\phantomsection\label{\detokenize{Fortran:f/_/find_fix_neighlist}}
\pysigstartsignatures
\pysiglinewithargsret{\sphinxbfcode{\sphinxupquote{function  }}\sphinxbfcode{\sphinxupquote{find\_fix\_neighlist}}}{\sphinxparam{id}\sphinxoptional{\sphinxparamcomma \sphinxparam{reqid}}}{}
\pysigstopsignatures
\sphinxAtStartPar
Find index of a neighbor list requested by a fix.

\sphinxAtStartPar
\DUrole{versionmodified,added}{Added in version 3Nov2022.}

\sphinxAtStartPar
The neighbor list request from a fix is identified by the fix ID and the
request ID. The request ID is typically zero, but will be \(>0\) for
fixes with multiple neighbor list requests.
\begin{quote}\begin{description}
\sphinxlineitem{Parameters}
\sphinxAtStartPar
\sphinxstylestrong{id}\sphinxstyleemphasis{ {[}}\sphinxstyleemphasis{character(len=*)}\sphinxstyleemphasis{{]}} :: Identifier of fix instance

\sphinxlineitem{Options}
\sphinxAtStartPar
\sphinxstylestrong{reqid}\sphinxstyleemphasis{ {[}}\sphinxstyleemphasis{integer(c\_int)}\sphinxstyleemphasis{,}\sphinxstyleemphasis{optional}\sphinxstyleemphasis{{]}} :: request ID to identify the neighbor list
in cases in which there are multiple requests from the same fix.
Default: 0.

\sphinxlineitem{Call to}
\sphinxAtStartPar
{\hyperref[\detokenize{Library_neighbor:_CPPv425lammps_find_fix_neighlistPvPKci}]{\sphinxcrossref{\sphinxcode{\sphinxupquote{lammps\_find\_fix\_neighlist()}}}}}

\sphinxlineitem{Return}
\sphinxAtStartPar
\sphinxstylestrong{index}\sphinxstyleemphasis{ {[}}\sphinxstyleemphasis{integer(c\_int)}\sphinxstyleemphasis{{]}} :: neighbor list index if found, otherwise \(-1\)

\end{description}\end{quote}

\end{fulllineitems}



\bigskip\hrule\bigskip

\index{find\_compute\_neighlist() (fortran function)@\spxentry{find\_compute\_neighlist()}\spxextra{fortran function}|spxpagem}

\begin{fulllineitems}
\phantomsection\label{\detokenize{Fortran:f/_/find_compute_neighlist}}
\pysigstartsignatures
\pysiglinewithargsret{\sphinxbfcode{\sphinxupquote{function  }}\sphinxbfcode{\sphinxupquote{find\_compute\_neighlist}}}{\sphinxparam{id}\sphinxoptional{\sphinxparamcomma \sphinxparam{reqid}}}{}
\pysigstopsignatures
\sphinxAtStartPar
Find index of a neighbor list requested by a compute.

\sphinxAtStartPar
\DUrole{versionmodified,added}{Added in version 3Nov2022.}

\sphinxAtStartPar
The neighbor list request from a compute is identified by the compute ID and
the request ID.  The request ID is typically zero, but will be \(> 0\)
in case a compute has multiple neighbor list requests.
\begin{quote}\begin{description}
\sphinxlineitem{Parameters}
\sphinxAtStartPar
\sphinxstylestrong{id}\sphinxstyleemphasis{ {[}}\sphinxstyleemphasis{character(len=*)}\sphinxstyleemphasis{{]}} :: Identifier of compute instance.

\sphinxlineitem{Options}
\sphinxAtStartPar
\sphinxstylestrong{reqid}\sphinxstyleemphasis{ {[}}\sphinxstyleemphasis{integer(c\_int)}\sphinxstyleemphasis{,}\sphinxstyleemphasis{optional}\sphinxstyleemphasis{{]}} :: request ID to identify the neighbor list
in cases in which there are multiple requests from the same compute.
Default: 0.

\sphinxlineitem{Call to}
\sphinxAtStartPar
{\hyperref[\detokenize{Library_neighbor:_CPPv429lammps_find_compute_neighlistPvPKci}]{\sphinxcrossref{\sphinxcode{\sphinxupquote{lammps\_find\_compute\_neighlist()}}}}}

\sphinxlineitem{Return}
\sphinxAtStartPar
\sphinxstylestrong{index}\sphinxstyleemphasis{ {[}}\sphinxstyleemphasis{integer(c\_int)}\sphinxstyleemphasis{{]}} :: neighbor list index if found, otherwise \(-1\).

\end{description}\end{quote}

\end{fulllineitems}



\bigskip\hrule\bigskip

\index{neighlist\_num\_elements() (fortran function)@\spxentry{neighlist\_num\_elements()}\spxextra{fortran function}|spxpagem}

\begin{fulllineitems}
\phantomsection\label{\detokenize{Fortran:f/_/neighlist_num_elements}}
\pysigstartsignatures
\pysiglinewithargsret{\sphinxbfcode{\sphinxupquote{function  }}\sphinxbfcode{\sphinxupquote{neighlist\_num\_elements}}}{\sphinxparam{idx}}{}
\pysigstopsignatures
\sphinxAtStartPar
Return the number of entries in the neighbor list with the given index.

\sphinxAtStartPar
\DUrole{versionmodified,added}{Added in version 3Nov2022.}
\begin{quote}\begin{description}
\sphinxlineitem{Parameters}
\sphinxAtStartPar
\sphinxstylestrong{idx}\sphinxstyleemphasis{ {[}}\sphinxstyleemphasis{integer(c\_int)}\sphinxstyleemphasis{{]}} :: neighbor list index

\sphinxlineitem{Call to}
\sphinxAtStartPar
{\hyperref[\detokenize{Library_neighbor:_CPPv429lammps_neighlist_num_elementsPvi}]{\sphinxcrossref{\sphinxcode{\sphinxupquote{lammps\_neighlist\_num\_elements()}}}}} {\hyperref[\detokenize{Library_neighbor:_CPPv429lammps_neighlist_num_elementsPvi}]{\sphinxcrossref{\sphinxcode{\sphinxupquote{lammps\_neighlist\_num\_elements()}}}}}

\sphinxlineitem{Return}
\sphinxAtStartPar
\sphinxstylestrong{inum}\sphinxstyleemphasis{ {[}}\sphinxstyleemphasis{integer(c\_int)}\sphinxstyleemphasis{{]}} :: number of entries in neighbor list, or \(-1\) if \sphinxstyleemphasis{idx} is not
a valid index.

\end{description}\end{quote}

\end{fulllineitems}



\bigskip\hrule\bigskip

\index{neighlist\_element\_neighbors() (fortran subroutine)@\spxentry{neighlist\_element\_neighbors()}\spxextra{fortran subroutine}|spxpagem}

\begin{fulllineitems}
\phantomsection\label{\detokenize{Fortran:f/_/neighlist_element_neighbors}}
\pysigstartsignatures
\pysiglinewithargsret{\sphinxbfcode{\sphinxupquote{subroutine  }}\sphinxbfcode{\sphinxupquote{neighlist\_element\_neighbors}}}{\sphinxparam{idx}\sphinxparamcomma \sphinxparam{element}\sphinxparamcomma \sphinxparam{iatom}\sphinxparamcomma \sphinxparam{neighbors}}{}
\pysigstopsignatures
\sphinxAtStartPar
Return atom local index, number of neighbors, and array of neighbor local
atom indices of a neighbor list entry.

\sphinxAtStartPar
\DUrole{versionmodified,added}{Added in version 3Nov2022.}
\begin{quote}\begin{description}
\sphinxlineitem{Parameters}\begin{itemize}
\item {} 
\sphinxAtStartPar
\sphinxstylestrong{idx}\sphinxstyleemphasis{ {[}}\sphinxstyleemphasis{integer(c\_int)}\sphinxstyleemphasis{{]}} :: index of this neighbor list in the list of all
neighbor lists

\item {} 
\sphinxAtStartPar
\sphinxstylestrong{element}\sphinxstyleemphasis{ {[}}\sphinxstyleemphasis{integer(c\_int)}\sphinxstyleemphasis{{]}} :: index of this neighbor list entry

\item {} 
\sphinxAtStartPar
\sphinxstylestrong{iatom}\sphinxstyleemphasis{ {[}}\sphinxstyleemphasis{integer(c\_int)}\sphinxstyleemphasis{{]}} :: local atom index (i.e., in the range
{[}1,nlocal+nghost{]}; \sphinxhyphen{}1 if invalid or element value

\item {} 
\sphinxAtStartPar
\sphinxstylestrong{neighbors}\sphinxstyleemphasis{ {[}}\sphinxstyleemphasis{integer(c\_int)}\sphinxstyleemphasis{,}\sphinxstyleemphasis{dimension(:),pointer}\sphinxstyleemphasis{{]}} :: pointer to an array of
neighboring atom local indices

\end{itemize}

\sphinxlineitem{Call to}
\sphinxAtStartPar
{\hyperref[\detokenize{Library_neighbor:_CPPv434lammps_neighlist_element_neighborsPviiPiPiPPi}]{\sphinxcrossref{\sphinxcode{\sphinxupquote{lammps\_neighlist\_element\_neighbors()}}}}}

\end{description}\end{quote}

\end{fulllineitems}



\bigskip\hrule\bigskip

\index{version() (fortran function)@\spxentry{version()}\spxextra{fortran function}|spxpagem}

\begin{fulllineitems}
\phantomsection\label{\detokenize{Fortran:f/_/version}}
\pysigstartsignatures
\pysiglinewithargsret{\sphinxbfcode{\sphinxupquote{function  }}\sphinxbfcode{\sphinxupquote{version}}}{}{}
\pysigstopsignatures
\sphinxAtStartPar
This method returns the numeric LAMMPS version like
{\hyperref[\detokenize{Library_config:_CPPv414lammps_versionPv}]{\sphinxcrossref{\sphinxcode{\sphinxupquote{lammps\_version()}}}}} does.
\begin{quote}\begin{description}
\sphinxlineitem{Call to}
\sphinxAtStartPar
{\hyperref[\detokenize{Library_config:_CPPv414lammps_versionPv}]{\sphinxcrossref{\sphinxcode{\sphinxupquote{lammps\_version()}}}}}

\sphinxlineitem{Return}
\sphinxAtStartPar
\sphinxstylestrong{version}\sphinxstyleemphasis{ {[}}\sphinxstyleemphasis{integer}\sphinxstyleemphasis{{]}} :: LAMMPS version

\end{description}\end{quote}

\end{fulllineitems}



\bigskip\hrule\bigskip

\index{get\_os\_info() (fortran subroutine)@\spxentry{get\_os\_info()}\spxextra{fortran subroutine}|spxpagem}

\begin{fulllineitems}
\phantomsection\label{\detokenize{Fortran:f/_/get_os_info}}
\pysigstartsignatures
\pysiglinewithargsret{\sphinxbfcode{\sphinxupquote{subroutine  }}\sphinxbfcode{\sphinxupquote{get\_os\_info}}}{\sphinxparam{buffer}}{}
\pysigstopsignatures
\sphinxAtStartPar
This function can be used to retrieve detailed information about the hosting
operating system and compiler/runtime environment.

\sphinxAtStartPar
\DUrole{versionmodified,added}{Added in version 3Nov2022.}

\sphinxAtStartPar
A suitable buffer has to be provided. The assembled text will be truncated
so as not to overflow this buffer. The string is typically a few hundred
bytes long.
\begin{quote}\begin{description}
\sphinxlineitem{Parameters}
\sphinxAtStartPar
\sphinxstylestrong{buffer}\sphinxstyleemphasis{ {[}}\sphinxstyleemphasis{character(len=*)}\sphinxstyleemphasis{{]}} :: string that will house the information.

\sphinxlineitem{Call to}
\sphinxAtStartPar
{\hyperref[\detokenize{Library_config:_CPPv418lammps_get_os_infoPci}]{\sphinxcrossref{\sphinxcode{\sphinxupquote{lammps\_get\_os\_info()}}}}}

\end{description}\end{quote}

\end{fulllineitems}



\bigskip\hrule\bigskip

\index{config\_has\_mpi\_support() (fortran function)@\spxentry{config\_has\_mpi\_support()}\spxextra{fortran function}|spxpagem}

\begin{fulllineitems}
\phantomsection\label{\detokenize{Fortran:f/_/config_has_mpi_support}}
\pysigstartsignatures
\pysiglinewithargsret{\sphinxbfcode{\sphinxupquote{function  }}\sphinxbfcode{\sphinxupquote{config\_has\_mpi\_support}}}{}{}
\pysigstopsignatures
\sphinxAtStartPar
This function is used to query whether LAMMPS was compiled with a real MPI
library or in serial.

\sphinxAtStartPar
\DUrole{versionmodified,added}{Added in version 3Nov2022.}
\begin{quote}\begin{description}
\sphinxlineitem{Call to}
\sphinxAtStartPar
{\hyperref[\detokenize{Library_config:_CPPv429lammps_config_has_mpi_supportv}]{\sphinxcrossref{\sphinxcode{\sphinxupquote{lammps\_config\_has\_mpi\_support()}}}}}

\sphinxlineitem{Return}
\sphinxAtStartPar
\sphinxstylestrong{has\_mpi}\sphinxstyleemphasis{ {[}}\sphinxstyleemphasis{logical}\sphinxstyleemphasis{{]}} :: \sphinxcode{\sphinxupquote{.FALSE.}} when compiled with STUBS, \sphinxcode{\sphinxupquote{.TRUE.}} if
complied with MPI.

\end{description}\end{quote}

\end{fulllineitems}



\bigskip\hrule\bigskip

\index{config\_has\_gzip\_support() (fortran function)@\spxentry{config\_has\_gzip\_support()}\spxextra{fortran function}|spxpagem}

\begin{fulllineitems}
\phantomsection\label{\detokenize{Fortran:f/_/config_has_gzip_support}}
\pysigstartsignatures
\pysiglinewithargsret{\sphinxbfcode{\sphinxupquote{function  }}\sphinxbfcode{\sphinxupquote{config\_has\_gzip\_support}}}{}{}
\pysigstopsignatures
\sphinxAtStartPar
Check if the LAMMPS library supports reading or writing compressed
files via a pipe to gzip or similar compression programs.

\sphinxAtStartPar
\DUrole{versionmodified,added}{Added in version 3Nov2022.}

\sphinxAtStartPar
Several LAMMPS commands (e.g., {\hyperref[\detokenize{read_data::doc}]{\sphinxcrossref{\DUrole{doc}{read\_data command}}}}, {\hyperref[\detokenize{write_data::doc}]{\sphinxcrossref{\DUrole{doc}{write\_data command}}}},
{\hyperref[\detokenize{dump::doc}]{\sphinxcrossref{\DUrole{doc}{dump styles atom, custom, and xyz}}}}) support reading and writing
compressed files via creating a pipe to the \sphinxcode{\sphinxupquote{gzip}} program.  This function
checks whether this feature was {\hyperref[\detokenize{Build_settings:gzip}]{\sphinxcrossref{\DUrole{std,std-ref}{enabled at compile time}}}}.
It does \sphinxstylestrong{not} check whether \sphinxcode{\sphinxupquote{gzip}} or any other supported compression
programs themselves are installed and usable.
\begin{quote}\begin{description}
\sphinxlineitem{Call to}
\sphinxAtStartPar
{\hyperref[\detokenize{Library_config:_CPPv430lammps_config_has_gzip_supportv}]{\sphinxcrossref{\sphinxcode{\sphinxupquote{lammps\_config\_has\_gzip\_support()}}}}}

\sphinxlineitem{Return}
\sphinxAtStartPar
\sphinxstylestrong{has\_gzip}\sphinxstyleemphasis{ {[}}\sphinxstyleemphasis{logical}\sphinxstyleemphasis{{]}}

\end{description}\end{quote}

\end{fulllineitems}



\bigskip\hrule\bigskip

\index{config\_has\_png\_support() (fortran function)@\spxentry{config\_has\_png\_support()}\spxextra{fortran function}|spxpagem}

\begin{fulllineitems}
\phantomsection\label{\detokenize{Fortran:f/_/config_has_png_support}}
\pysigstartsignatures
\pysiglinewithargsret{\sphinxbfcode{\sphinxupquote{function  }}\sphinxbfcode{\sphinxupquote{config\_has\_png\_support}}}{}{}
\pysigstopsignatures
\sphinxAtStartPar
Check if the LAMMPS library supports writing PNG format images.

\sphinxAtStartPar
\DUrole{versionmodified,added}{Added in version 3Nov2022.}

\sphinxAtStartPar
The LAMMPS {\hyperref[\detokenize{dump_image::doc}]{\sphinxcrossref{\DUrole{doc}{dump style image}}}} supports writing multiple
image file formats.  Most of them, however, need support from an external
library, and using that has to be {\hyperref[\detokenize{Build_settings:graphics}]{\sphinxcrossref{\DUrole{std,std-ref}{enabled at compile time}}}}.
This function checks whether support for the \sphinxhref{https://en.wikipedia.org/wiki/Portable\_Network\_Graphics}{PNG image file format} is available
in the current LAMMPS library.
\begin{quote}\begin{description}
\sphinxlineitem{Call to}
\sphinxAtStartPar
{\hyperref[\detokenize{Library_config:_CPPv429lammps_config_has_png_supportv}]{\sphinxcrossref{\sphinxcode{\sphinxupquote{lammps\_config\_has\_png\_support()}}}}}

\sphinxlineitem{Return}
\sphinxAtStartPar
\sphinxstylestrong{has\_png}\sphinxstyleemphasis{ {[}}\sphinxstyleemphasis{logical}\sphinxstyleemphasis{{]}}

\end{description}\end{quote}

\end{fulllineitems}



\bigskip\hrule\bigskip

\index{config\_has\_jpeg\_support() (fortran function)@\spxentry{config\_has\_jpeg\_support()}\spxextra{fortran function}|spxpagem}

\begin{fulllineitems}
\phantomsection\label{\detokenize{Fortran:f/_/config_has_jpeg_support}}
\pysigstartsignatures
\pysiglinewithargsret{\sphinxbfcode{\sphinxupquote{function  }}\sphinxbfcode{\sphinxupquote{config\_has\_jpeg\_support}}}{}{}
\pysigstopsignatures
\sphinxAtStartPar
Check if the LAMMPS library supports writing JPEG format images.

\sphinxAtStartPar
\DUrole{versionmodified,added}{Added in version 3Nov2022.}

\sphinxAtStartPar
The LAMMPS {\hyperref[\detokenize{dump_image::doc}]{\sphinxcrossref{\DUrole{doc}{dump style image}}}} supports writing multiple
image file formats.  Most of them, however, need support from an external
library, and using that has to be {\hyperref[\detokenize{Build_settings:graphics}]{\sphinxcrossref{\DUrole{std,std-ref}{enabled at compile time}}}}.
This function checks whether support for the \sphinxhref{https://jpeg.org/jpeg/}{JPEG image file format} is available in the current LAMMPS library.
\begin{quote}\begin{description}
\sphinxlineitem{Call to}
\sphinxAtStartPar
{\hyperref[\detokenize{Library_config:_CPPv430lammps_config_has_jpeg_supportv}]{\sphinxcrossref{\sphinxcode{\sphinxupquote{lammps\_config\_has\_jpeg\_support()}}}}}

\sphinxlineitem{Return}
\sphinxAtStartPar
\sphinxstylestrong{has\_jpeg}\sphinxstyleemphasis{ {[}}\sphinxstyleemphasis{logical}\sphinxstyleemphasis{{]}}

\end{description}\end{quote}

\end{fulllineitems}



\bigskip\hrule\bigskip

\index{config\_has\_ffmpeg\_support() (fortran function)@\spxentry{config\_has\_ffmpeg\_support()}\spxextra{fortran function}|spxpagem}

\begin{fulllineitems}
\phantomsection\label{\detokenize{Fortran:f/_/config_has_ffmpeg_support}}
\pysigstartsignatures
\pysiglinewithargsret{\sphinxbfcode{\sphinxupquote{function  }}\sphinxbfcode{\sphinxupquote{config\_has\_ffmpeg\_support}}}{}{}
\pysigstopsignatures
\sphinxAtStartPar
Check if the LAMMPS library supports creating movie files via a pipe to
ffmpeg.

\sphinxAtStartPar
\DUrole{versionmodified,added}{Added in version 3Nov2022.}

\sphinxAtStartPar
The LAMMPS {\hyperref[\detokenize{dump_image::doc}]{\sphinxcrossref{\DUrole{doc}{dump style movie}}}} supports generating movies
from images on\sphinxhyphen{}the\sphinxhyphen{}fly via creating a pipe to the
\sphinxhref{https://ffmpeg.org/}{ffmpeg} program.
This function checks whether this feature was
{\hyperref[\detokenize{Build_settings:graphics}]{\sphinxcrossref{\DUrole{std,std-ref}{enabled at compile time}}}}.
It does \sphinxstylestrong{not} check whether the \sphinxcode{\sphinxupquote{ffmpeg}} itself is installed and usable.
\begin{quote}\begin{description}
\sphinxlineitem{Call to}
\sphinxAtStartPar
{\hyperref[\detokenize{Library_config:_CPPv432lammps_config_has_ffmpeg_supportv}]{\sphinxcrossref{\sphinxcode{\sphinxupquote{lammps\_config\_has\_ffmpeg\_support()}}}}}

\sphinxlineitem{Return}
\sphinxAtStartPar
\sphinxstylestrong{has\_ffmpeg}\sphinxstyleemphasis{ {[}}\sphinxstyleemphasis{logical}\sphinxstyleemphasis{{]}}

\end{description}\end{quote}

\end{fulllineitems}



\bigskip\hrule\bigskip

\index{config\_has\_exceptions() (fortran function)@\spxentry{config\_has\_exceptions()}\spxextra{fortran function}|spxpagem}

\begin{fulllineitems}
\phantomsection\label{\detokenize{Fortran:f/_/config_has_exceptions}}
\pysigstartsignatures
\pysiglinewithargsret{\sphinxbfcode{\sphinxupquote{function  }}\sphinxbfcode{\sphinxupquote{config\_has\_exceptions}}}{}{}
\pysigstopsignatures
\sphinxAtStartPar
Check whether LAMMPS errors will throw C++ exceptions.

\sphinxAtStartPar
\DUrole{versionmodified,added}{Added in version 3Nov2022.}

\sphinxAtStartPar
When using the library interface, the library interface functions
will “catch” exceptions, and then the error status can be checked by
calling \sphinxcode{\sphinxupquote{has\_error()}}.  The most recent error message can be
retrieved via \sphinxcode{\sphinxupquote{get\_last\_error\_message()}}.  This allows to
restart a calculation or delete and recreate the LAMMPS instance when
a C++ exception occurs.  One application of using exceptions this way
is the {\hyperref[\detokenize{Tools:lammps-gui}]{\sphinxcrossref{\DUrole{std,std-ref}{LAMMPS\sphinxhyphen{}GUI}}}}.
\begin{quote}\begin{description}
\sphinxlineitem{Call to}
\sphinxAtStartPar
{\hyperref[\detokenize{Library_config:_CPPv428lammps_config_has_exceptionsv}]{\sphinxcrossref{\sphinxcode{\sphinxupquote{lammps\_config\_has\_exceptions()}}}}}

\sphinxlineitem{Return}
\sphinxAtStartPar
\sphinxstylestrong{has\_exceptions}\sphinxstyleemphasis{ {[}}\sphinxstyleemphasis{logical}\sphinxstyleemphasis{{]}}

\end{description}\end{quote}

\end{fulllineitems}



\bigskip\hrule\bigskip

\index{config\_has\_package() (fortran function)@\spxentry{config\_has\_package()}\spxextra{fortran function}|spxpagem}

\begin{fulllineitems}
\phantomsection\label{\detokenize{Fortran:f/_/config_has_package}}
\pysigstartsignatures
\pysiglinewithargsret{\sphinxbfcode{\sphinxupquote{function  }}\sphinxbfcode{\sphinxupquote{config\_has\_package}}}{\sphinxparam{name}}{}
\pysigstopsignatures
\sphinxAtStartPar
Check whether a specific package has been included in LAMMPS

\sphinxAtStartPar
\DUrole{versionmodified,added}{Added in version 3Nov2022.}

\sphinxAtStartPar
This function checks whether the LAMMPS library in use includes the specific
{\hyperref[\detokenize{Packages::doc}]{\sphinxcrossref{\DUrole{doc}{LAMMPS package}}}} provided as argument.
\begin{quote}\begin{description}
\sphinxlineitem{Call to}
\sphinxAtStartPar
{\hyperref[\detokenize{Library_config:_CPPv425lammps_config_has_packagePKc}]{\sphinxcrossref{\sphinxcode{\sphinxupquote{lammps\_config\_has\_package()}}}}}

\sphinxlineitem{Return}
\sphinxAtStartPar
\sphinxstylestrong{has\_package}\sphinxstyleemphasis{ {[}}\sphinxstyleemphasis{logical}\sphinxstyleemphasis{{]}}

\end{description}\end{quote}

\end{fulllineitems}



\bigskip\hrule\bigskip

\index{config\_package\_count() (fortran function)@\spxentry{config\_package\_count()}\spxextra{fortran function}|spxpagem}

\begin{fulllineitems}
\phantomsection\label{\detokenize{Fortran:f/_/config_package_count}}
\pysigstartsignatures
\pysiglinewithargsret{\sphinxbfcode{\sphinxupquote{function  }}\sphinxbfcode{\sphinxupquote{config\_package\_count}}}{}{}
\pysigstopsignatures
\sphinxAtStartPar
Count the number of installed packages in the LAMMPS library.

\sphinxAtStartPar
\DUrole{versionmodified,added}{Added in version 3Nov2022.}

\sphinxAtStartPar
This function counts how many {\hyperref[\detokenize{Packages::doc}]{\sphinxcrossref{\DUrole{doc}{LAMMPS packages}}}} are
included in the LAMMPS library in use. It directly calls the C library
function {\hyperref[\detokenize{Library_config:_CPPv427lammps_config_package_countv}]{\sphinxcrossref{\sphinxcode{\sphinxupquote{lammps\_config\_package\_count()}}}}}.
\begin{quote}\begin{description}
\sphinxlineitem{Call to}
\sphinxAtStartPar
{\hyperref[\detokenize{Library_config:_CPPv427lammps_config_package_countv}]{\sphinxcrossref{\sphinxcode{\sphinxupquote{lammps\_config\_package\_count()}}}}}

\sphinxlineitem{Return}
\sphinxAtStartPar
\sphinxstylestrong{npackages}\sphinxstyleemphasis{ {[}}\sphinxstyleemphasis{integer(c\_int)}\sphinxstyleemphasis{{]}} :: number of packages installed

\end{description}\end{quote}

\end{fulllineitems}



\bigskip\hrule\bigskip

\index{config\_package\_name() (fortran subroutine)@\spxentry{config\_package\_name()}\spxextra{fortran subroutine}|spxpagem}

\begin{fulllineitems}
\phantomsection\label{\detokenize{Fortran:f/_/config_package_name}}
\pysigstartsignatures
\pysiglinewithargsret{\sphinxbfcode{\sphinxupquote{subroutine  }}\sphinxbfcode{\sphinxupquote{config\_package\_name}}}{\sphinxparam{idx}\sphinxparamcomma \sphinxparam{buffer}}{}
\pysigstopsignatures
\sphinxAtStartPar
Get the name of a package in the list of installed packages in the LAMMPS
library.

\sphinxAtStartPar
\DUrole{versionmodified,added}{Added in version 3Nov2022.}

\sphinxAtStartPar
This subroutine copies the name of the package with the index \sphinxstyleemphasis{idx} into the
provided string \sphinxstyleemphasis{buffer}. If the name of the package exceeds the length of
the buffer, it will be truncated accordingly.  If the index is out of range,
\sphinxstyleemphasis{buffer} is set to an empty string.
\begin{quote}\begin{description}
\sphinxlineitem{Parameters}\begin{itemize}
\item {} 
\sphinxAtStartPar
\sphinxstylestrong{idx}\sphinxstyleemphasis{ {[}}\sphinxstyleemphasis{integer(c\_int)}\sphinxstyleemphasis{{]}} :: index of the package in the list of included packages
\((0 \le idx < \text{package count})\)

\item {} 
\sphinxAtStartPar
\sphinxstylestrong{buffer}\sphinxstyleemphasis{ {[}}\sphinxstyleemphasis{character(len=*)}\sphinxstyleemphasis{{]}} :: string to hold the name of the package

\end{itemize}

\sphinxlineitem{Call to}
\sphinxAtStartPar
{\hyperref[\detokenize{Library_config:_CPPv426lammps_config_package_nameiPci}]{\sphinxcrossref{\sphinxcode{\sphinxupquote{lammps\_config\_package\_name()}}}}}

\end{description}\end{quote}

\end{fulllineitems}



\bigskip\hrule\bigskip

\index{installed\_packages() (fortran subroutine)@\spxentry{installed\_packages()}\spxextra{fortran subroutine}|spxpagem}

\begin{fulllineitems}
\phantomsection\label{\detokenize{Fortran:f/_/installed_packages}}
\pysigstartsignatures
\pysiglinewithargsret{\sphinxbfcode{\sphinxupquote{subroutine  }}\sphinxbfcode{\sphinxupquote{installed\_packages}}}{\sphinxparam{package}\sphinxoptional{\sphinxparamcomma \sphinxparam{length}}}{}
\pysigstopsignatures
\sphinxAtStartPar
Obtain a list of the names of enabled packages in the LAMMPS shared library
and store it in \sphinxstyleemphasis{package}.

\sphinxAtStartPar
\DUrole{versionmodified,added}{Added in version 3Nov2022.}

\sphinxAtStartPar
This function is analogous to the {\hyperref[\detokenize{Python_module:lammps.lammps.installed_packages}]{\sphinxcrossref{\sphinxcode{\sphinxupquote{installed\_packages}}}}} function in the Python API.
The optional argument \sphinxstyleemphasis{length} sets the length of each string in the vector
\sphinxstyleemphasis{package} (default: 31).
\begin{quote}\begin{description}
\sphinxlineitem{Parameters}
\sphinxAtStartPar
\sphinxstylestrong{package}\sphinxstyleemphasis{ {[}}\sphinxstyleemphasis{character(len=:)}\sphinxstyleemphasis{,}\sphinxstyleemphasis{dimension(:),allocatable}\sphinxstyleemphasis{{]}} :: list of packages;
\sphinxstyleemphasis{must} have the \sphinxcode{\sphinxupquote{ALLOCATABLE}} attribute and be of rank 1
(i.e., \sphinxcode{\sphinxupquote{DIMENSION(:)}}) with allocatable length.

\sphinxlineitem{Options}
\sphinxAtStartPar
\sphinxstylestrong{length}\sphinxstyleemphasis{ {[}}\sphinxstyleemphasis{integer,optional}\sphinxstyleemphasis{{]}} :: length of each string in the list. Default: 31.

\sphinxlineitem{Call to}
\sphinxAtStartPar
{\hyperref[\detokenize{Library_config:_CPPv427lammps_config_package_countv}]{\sphinxcrossref{\sphinxcode{\sphinxupquote{lammps\_config\_package\_count()}}}}} {\hyperref[\detokenize{Library_config:_CPPv426lammps_config_package_nameiPci}]{\sphinxcrossref{\sphinxcode{\sphinxupquote{lammps\_config\_package\_name()}}}}}

\end{description}\end{quote}

\end{fulllineitems}



\bigskip\hrule\bigskip

\index{config\_accelerator() (fortran function)@\spxentry{config\_accelerator()}\spxextra{fortran function}|spxpagem}

\begin{fulllineitems}
\phantomsection\label{\detokenize{Fortran:f/_/config_accelerator}}
\pysigstartsignatures
\pysiglinewithargsret{\sphinxbfcode{\sphinxupquote{function  }}\sphinxbfcode{\sphinxupquote{config\_accelerator}}}{\sphinxparam{package}\sphinxparamcomma \sphinxparam{category}\sphinxparamcomma \sphinxparam{setting}}{}
\pysigstopsignatures
\sphinxAtStartPar
This function calls {\hyperref[\detokenize{Library_config:_CPPv425lammps_config_acceleratorPKcPKcPKc}]{\sphinxcrossref{\sphinxcode{\sphinxupquote{lammps\_config\_accelerator()}}}}} to check the
availability of compile time settings of included
{\hyperref[\detokenize{Speed_packages::doc}]{\sphinxcrossref{\DUrole{doc}{accelerator packages}}}} in LAMMPS.

\sphinxAtStartPar
\DUrole{versionmodified,added}{Added in version 3Nov2022.}

\sphinxAtStartPar
Supported packages names are “GPU”, “KOKKOS”, “INTEL”, and “OPENMP”.
Supported categories are “api” with possible settings “cuda”, “hip”, “phi”,
“pthreads”, “opencl”, “openmp”, and “serial”; and “precision” with
possible settings “double”, “mixed”, and “single”.
\begin{quote}\begin{description}
\sphinxlineitem{Parameters}\begin{itemize}
\item {} 
\sphinxAtStartPar
\sphinxstylestrong{package}\sphinxstyleemphasis{ {[}}\sphinxstyleemphasis{character(len=*)}\sphinxstyleemphasis{{]}} :: string with the name of the accelerator
package

\item {} 
\sphinxAtStartPar
\sphinxstylestrong{category}\sphinxstyleemphasis{ {[}}\sphinxstyleemphasis{character(len=*)}\sphinxstyleemphasis{{]}} :: string with the name of the setting

\item {} 
\sphinxAtStartPar
\sphinxstylestrong{setting}\sphinxstyleemphasis{ {[}}\sphinxstyleemphasis{character(len=*)}\sphinxstyleemphasis{{]}} :: string with the name of the specific
setting

\end{itemize}

\sphinxlineitem{Call to}
\sphinxAtStartPar
{\hyperref[\detokenize{Library_config:_CPPv425lammps_config_acceleratorPKcPKcPKc}]{\sphinxcrossref{\sphinxcode{\sphinxupquote{lammps\_config\_accelerator()}}}}}

\sphinxlineitem{Return}
\sphinxAtStartPar
\sphinxstylestrong{available}\sphinxstyleemphasis{ {[}}\sphinxstyleemphasis{logical}\sphinxstyleemphasis{{]}} :: \sphinxcode{\sphinxupquote{.TRUE.}} if the combination of package, category,
and setting is available, otherwise \sphinxcode{\sphinxupquote{.FALSE.}}.

\end{description}\end{quote}

\end{fulllineitems}



\bigskip\hrule\bigskip

\index{has\_gpu\_device() (fortran function)@\spxentry{has\_gpu\_device()}\spxextra{fortran function}|spxpagem}

\begin{fulllineitems}
\phantomsection\label{\detokenize{Fortran:f/_/has_gpu_device}}
\pysigstartsignatures
\pysiglinewithargsret{\sphinxbfcode{\sphinxupquote{function  }}\sphinxbfcode{\sphinxupquote{has\_gpu\_device}}}{}{}
\pysigstopsignatures
\sphinxAtStartPar
Checks for the presence of a viable GPU package device.

\sphinxAtStartPar
\DUrole{versionmodified,added}{Added in version 3Nov2022.}

\sphinxAtStartPar
This function calls {\hyperref[\detokenize{Library_config:_CPPv421lammps_has_gpu_devicev}]{\sphinxcrossref{\sphinxcode{\sphinxupquote{lammps\_has\_gpu\_device()}}}}}, which checks at
runtime whether an accelerator device is present that can be used with the
{\hyperref[\detokenize{Speed_gpu::doc}]{\sphinxcrossref{\DUrole{doc}{GPU package}}}}.

\sphinxAtStartPar
More detailed information about the available device or devices can
be obtained by calling the
\sphinxcode{\sphinxupquote{get\_gpu\_device\_info()}} subroutine.
\begin{quote}\begin{description}
\sphinxlineitem{Call to}
\sphinxAtStartPar
{\hyperref[\detokenize{Library_config:_CPPv421lammps_has_gpu_devicev}]{\sphinxcrossref{\sphinxcode{\sphinxupquote{lammps\_has\_gpu\_device()}}}}}

\sphinxlineitem{Return}
\sphinxAtStartPar
\sphinxstylestrong{available}\sphinxstyleemphasis{ {[}}\sphinxstyleemphasis{logical}\sphinxstyleemphasis{{]}} :: \sphinxcode{\sphinxupquote{.TRUE.}} if a viable device is available, \sphinxcode{\sphinxupquote{.FALSE.}} if not.

\end{description}\end{quote}

\end{fulllineitems}



\bigskip\hrule\bigskip

\index{get\_gpu\_device\_info() (fortran subroutine)@\spxentry{get\_gpu\_device\_info()}\spxextra{fortran subroutine}|spxpagem}

\begin{fulllineitems}
\phantomsection\label{\detokenize{Fortran:f/_/get_gpu_device_info}}
\pysigstartsignatures
\pysiglinewithargsret{\sphinxbfcode{\sphinxupquote{subroutine  }}\sphinxbfcode{\sphinxupquote{get\_gpu\_device\_info}}}{\sphinxparam{buffer}}{}
\pysigstopsignatures
\sphinxAtStartPar
Get GPU package device information.

\sphinxAtStartPar
\DUrole{versionmodified,added}{Added in version 3Nov2022.}

\sphinxAtStartPar
Calls {\hyperref[\detokenize{Library_config:_CPPv426lammps_get_gpu_device_infoPci}]{\sphinxcrossref{\sphinxcode{\sphinxupquote{lammps\_get\_gpu\_device\_info()}}}}} to retrieve detailed
information about any accelerator devices that are viable for use with the
{\hyperref[\detokenize{Speed_gpu::doc}]{\sphinxcrossref{\DUrole{doc}{GPU package}}}}. It will fill \sphinxstyleemphasis{buffer} with a string that is
equivalent to the output of the \sphinxcode{\sphinxupquote{nvc\_get\_device}} or \sphinxcode{\sphinxupquote{ocl\_get\_device}} or
\sphinxcode{\sphinxupquote{hip\_get\_device}} tools that are compiled alongside LAMMPS if the GPU
package is enabled.

\sphinxAtStartPar
A suitable\sphinxhyphen{}length Fortran string has to be provided. The assembled text will
be truncated so as not to overflow this buffer.  This string can be several
kilobytes long if multiple devices are present.
\begin{quote}\begin{description}
\sphinxlineitem{Parameters}
\sphinxAtStartPar
\sphinxstylestrong{buffer}\sphinxstyleemphasis{ {[}}\sphinxstyleemphasis{character(len=*)}\sphinxstyleemphasis{{]}} :: string into which to copy the information.

\sphinxlineitem{Call to}
\sphinxAtStartPar
{\hyperref[\detokenize{Library_config:_CPPv426lammps_get_gpu_device_infoPci}]{\sphinxcrossref{\sphinxcode{\sphinxupquote{lammps\_get\_gpu\_device\_info()}}}}}

\end{description}\end{quote}

\end{fulllineitems}



\bigskip\hrule\bigskip

\index{has\_style() (fortran function)@\spxentry{has\_style()}\spxextra{fortran function}|spxpagem}

\begin{fulllineitems}
\phantomsection\label{\detokenize{Fortran:f/_/has_style}}
\pysigstartsignatures
\pysiglinewithargsret{\sphinxbfcode{\sphinxupquote{function  }}\sphinxbfcode{\sphinxupquote{has\_style}}}{\sphinxparam{category}\sphinxparamcomma \sphinxparam{name}}{}
\pysigstopsignatures
\sphinxAtStartPar
Check whether a specific style has been included in LAMMPS.

\sphinxAtStartPar
\DUrole{versionmodified,added}{Added in version 3Nov2022.}

\sphinxAtStartPar
This function calls {\hyperref[\detokenize{Library_config:_CPPv416lammps_has_stylePvPKcPKc}]{\sphinxcrossref{\sphinxcode{\sphinxupquote{lammps\_has\_style()}}}}} to check whether the
LAMMPS library in use includes the specific style \sphinxstyleemphasis{name} associated with a
specific \sphinxstyleemphasis{category} provided as arguments.  Please see
{\hyperref[\detokenize{Library_config:_CPPv416lammps_has_stylePvPKcPKc}]{\sphinxcrossref{\sphinxcode{\sphinxupquote{lammps\_has\_style()}}}}} for a list of valid categories.
\begin{quote}\begin{description}
\sphinxlineitem{Parameters}\begin{itemize}
\item {} 
\sphinxAtStartPar
\sphinxstylestrong{category}\sphinxstyleemphasis{ {[}}\sphinxstyleemphasis{character(len=*)}\sphinxstyleemphasis{{]}} :: category of the style

\item {} 
\sphinxAtStartPar
\sphinxstylestrong{name}\sphinxstyleemphasis{ {[}}\sphinxstyleemphasis{character(len=*)}\sphinxstyleemphasis{{]}} :: name of the style

\end{itemize}

\sphinxlineitem{Call to}
\sphinxAtStartPar
{\hyperref[\detokenize{Library_config:_CPPv416lammps_has_stylePvPKcPKc}]{\sphinxcrossref{\sphinxcode{\sphinxupquote{lammps\_has\_style()}}}}}

\sphinxlineitem{Return}
\sphinxAtStartPar
\sphinxstylestrong{has\_style}\sphinxstyleemphasis{ {[}}\sphinxstyleemphasis{logical}\sphinxstyleemphasis{{]}} :: \sphinxcode{\sphinxupquote{.TRUE.}} if included, \sphinxcode{\sphinxupquote{.FALSE.}} if not.

\end{description}\end{quote}

\end{fulllineitems}



\bigskip\hrule\bigskip

\index{style\_count() (fortran function)@\spxentry{style\_count()}\spxextra{fortran function}|spxpagem}

\begin{fulllineitems}
\phantomsection\label{\detokenize{Fortran:f/_/style_count}}
\pysigstartsignatures
\pysiglinewithargsret{\sphinxbfcode{\sphinxupquote{function  }}\sphinxbfcode{\sphinxupquote{style\_count}}}{\sphinxparam{category}}{}
\pysigstopsignatures
\sphinxAtStartPar
Count the number of styles of \sphinxstyleemphasis{category} in the LAMMPS library.

\sphinxAtStartPar
\DUrole{versionmodified,added}{Added in version 3Nov2022.}

\sphinxAtStartPar
This function counts how many styles in the provided \sphinxstyleemphasis{category} are
included in the LAMMPS library currently in use. Please see
{\hyperref[\detokenize{Library_config:_CPPv416lammps_has_stylePvPKcPKc}]{\sphinxcrossref{\sphinxcode{\sphinxupquote{lammps\_has\_style()}}}}} for a list of valid categories.
\begin{quote}\begin{description}
\sphinxlineitem{Parameters}
\sphinxAtStartPar
\sphinxstylestrong{category}\sphinxstyleemphasis{ {[}}\sphinxstyleemphasis{character(len=*)}\sphinxstyleemphasis{{]}} :: category of styles to count

\sphinxlineitem{Call to}
\sphinxAtStartPar
{\hyperref[\detokenize{Library_config:_CPPv418lammps_style_countPvPKc}]{\sphinxcrossref{\sphinxcode{\sphinxupquote{lammps\_style\_count()}}}}}

\sphinxlineitem{Return}
\sphinxAtStartPar
\sphinxstylestrong{count}\sphinxstyleemphasis{ {[}}\sphinxstyleemphasis{integer(c\_int)}\sphinxstyleemphasis{{]}} :: number of styles in \sphinxstyleemphasis{category}

\end{description}\end{quote}

\end{fulllineitems}



\bigskip\hrule\bigskip

\index{style\_name() (fortran subroutine)@\spxentry{style\_name()}\spxextra{fortran subroutine}|spxpagem}

\begin{fulllineitems}
\phantomsection\label{\detokenize{Fortran:f/_/style_name}}
\pysigstartsignatures
\pysiglinewithargsret{\sphinxbfcode{\sphinxupquote{subroutine  }}\sphinxbfcode{\sphinxupquote{style\_name}}}{\sphinxparam{category}\sphinxparamcomma \sphinxparam{idx}\sphinxparamcomma \sphinxparam{buffer}}{}
\pysigstopsignatures
\sphinxAtStartPar
Look up the name of a style by index in the list of styles of a given
category in the LAMMPS library.

\sphinxAtStartPar
\DUrole{versionmodified,added}{Added in version 3Nov2022.}

\sphinxAtStartPar
This function calls {\hyperref[\detokenize{Library_config:_CPPv417lammps_style_namePvPKciPci}]{\sphinxcrossref{\sphinxcode{\sphinxupquote{lammps\_style\_name()}}}}} and copies the name of
the \sphinxstyleemphasis{category} style with index \sphinxstyleemphasis{idx} into the provided string \sphinxstyleemphasis{buffer}.
The length of \sphinxstyleemphasis{buffer} must be long enough to contain the name of the
style; if it is too short, the name will be truncated accordingly.
If \sphinxstyleemphasis{idx} is out of range, \sphinxstyleemphasis{buffer} will be the empty string and a warning
will be issued.
\begin{quote}\begin{description}
\sphinxlineitem{Parameters}\begin{itemize}
\item {} 
\sphinxAtStartPar
\sphinxstylestrong{category}\sphinxstyleemphasis{ {[}}\sphinxstyleemphasis{character(len=*)}\sphinxstyleemphasis{{]}} :: category of styles

\item {} 
\sphinxAtStartPar
\sphinxstylestrong{idx}\sphinxstyleemphasis{ {[}}\sphinxstyleemphasis{integer(c\_int)}\sphinxstyleemphasis{{]}} :: index of the style in the list of \sphinxstyleemphasis{category}
styles \((1 \leq idx \leq \text{style count})\)

\item {} 
\sphinxAtStartPar
\sphinxstylestrong{buffer}\sphinxstyleemphasis{ {[}}\sphinxstyleemphasis{character(len*)}\sphinxstyleemphasis{{]}} :: string buffer to copy the name of the style
into

\end{itemize}

\sphinxlineitem{Call to}
\sphinxAtStartPar
{\hyperref[\detokenize{Library_config:_CPPv417lammps_style_namePvPKciPci}]{\sphinxcrossref{\sphinxcode{\sphinxupquote{lammps\_style\_name()}}}}}

\end{description}\end{quote}

\end{fulllineitems}



\bigskip\hrule\bigskip

\index{has\_id() (fortran function)@\spxentry{has\_id()}\spxextra{fortran function}|spxpagem}

\begin{fulllineitems}
\phantomsection\label{\detokenize{Fortran:f/_/has_id}}
\pysigstartsignatures
\pysiglinewithargsret{\sphinxbfcode{\sphinxupquote{function  }}\sphinxbfcode{\sphinxupquote{has\_id}}}{\sphinxparam{category}\sphinxparamcomma \sphinxparam{name}}{}
\pysigstopsignatures
\sphinxAtStartPar
This function checks if the current LAMMPS instance a \sphinxstyleemphasis{category} ID of
the given \sphinxstyleemphasis{name} exists.  Valid categories are: \sphinxstyleemphasis{compute}, \sphinxstyleemphasis{dump},
\sphinxstyleemphasis{fix}, \sphinxstyleemphasis{group}, \sphinxstyleemphasis{molecule}, \sphinxstyleemphasis{region}, and \sphinxstyleemphasis{variable}.

\sphinxAtStartPar
\DUrole{versionmodified,added}{Added in version 3Nov2022.}
\begin{quote}\begin{description}
\sphinxlineitem{Parameters}\begin{itemize}
\item {} 
\sphinxAtStartPar
\sphinxstylestrong{category}\sphinxstyleemphasis{ {[}}\sphinxstyleemphasis{character(len=*)}\sphinxstyleemphasis{{]}} :: category of the ID

\item {} 
\sphinxAtStartPar
\sphinxstylestrong{name}\sphinxstyleemphasis{ {[}}\sphinxstyleemphasis{character(len=*)}\sphinxstyleemphasis{{]}} :: name of the ID

\end{itemize}

\sphinxlineitem{Call to}
\sphinxAtStartPar
{\hyperref[\detokenize{Library_config:_CPPv413lammps_has_idPvPKcPKc}]{\sphinxcrossref{\sphinxcode{\sphinxupquote{lammps\_has\_id()}}}}}

\sphinxlineitem{Return}
\sphinxAtStartPar
\sphinxstylestrong{has\_id}\sphinxstyleemphasis{ {[}}\sphinxstyleemphasis{logical}\sphinxstyleemphasis{{]}} :: \sphinxcode{\sphinxupquote{.TRUE.}} if \sphinxstyleemphasis{category} style \sphinxstyleemphasis{name} exists, \sphinxcode{\sphinxupquote{.FALSE.}} if not.

\end{description}\end{quote}

\end{fulllineitems}



\bigskip\hrule\bigskip

\index{id\_count() (fortran function)@\spxentry{id\_count()}\spxextra{fortran function}|spxpagem}

\begin{fulllineitems}
\phantomsection\label{\detokenize{Fortran:f/_/id_count}}
\pysigstartsignatures
\pysiglinewithargsret{\sphinxbfcode{\sphinxupquote{function  }}\sphinxbfcode{\sphinxupquote{id\_count}}}{\sphinxparam{category}}{}
\pysigstopsignatures
\sphinxAtStartPar
This function counts how many IDs in the provided \sphinxstyleemphasis{category} are defined in
the current LAMMPS instance. Please see \sphinxcode{\sphinxupquote{has\_id()}} for a list of
valid categories.

\sphinxAtStartPar
\DUrole{versionmodified,added}{Added in version 3Nov2022.}
\begin{quote}\begin{description}
\sphinxlineitem{Parameters}
\sphinxAtStartPar
\sphinxstylestrong{category}\sphinxstyleemphasis{ {[}}\sphinxstyleemphasis{character(len=*)}\sphinxstyleemphasis{{]}} :: category of the ID

\sphinxlineitem{Call to}
\sphinxAtStartPar
{\hyperref[\detokenize{Library_config:_CPPv415lammps_id_countPvPKc}]{\sphinxcrossref{\sphinxcode{\sphinxupquote{lammps\_id\_count()}}}}}

\sphinxlineitem{Return}
\sphinxAtStartPar
\sphinxstylestrong{count}\sphinxstyleemphasis{ {[}}\sphinxstyleemphasis{integer(c\_int)}\sphinxstyleemphasis{{]}} :: number of IDs in \sphinxstyleemphasis{category}

\end{description}\end{quote}

\end{fulllineitems}



\bigskip\hrule\bigskip

\index{id\_name() (fortran subroutine)@\spxentry{id\_name()}\spxextra{fortran subroutine}|spxpagem}

\begin{fulllineitems}
\phantomsection\label{\detokenize{Fortran:f/_/id_name}}
\pysigstartsignatures
\pysiglinewithargsret{\sphinxbfcode{\sphinxupquote{subroutine  }}\sphinxbfcode{\sphinxupquote{id\_name}}}{\sphinxparam{category}\sphinxparamcomma \sphinxparam{idx}\sphinxparamcomma \sphinxparam{buffer}}{}
\pysigstopsignatures
\sphinxAtStartPar
Look up the name of an ID by index in the list of IDs of a given category.

\sphinxAtStartPar
\DUrole{versionmodified,added}{Added in version 3Nov2022.}

\sphinxAtStartPar
This function copies the name of the \sphinxstyleemphasis{category} ID with the index \sphinxstyleemphasis{idx} into
the provided string \sphinxstyleemphasis{buffer}.  The length of the buffer must be long
enough to hold the string; if the name of the style exceeds the length of
the buffer, it will be truncated accordingly. If \sphinxstyleemphasis{buffer} is
\sphinxcode{\sphinxupquote{ALLOCATABLE}}, it must be allocated \sphinxstyleemphasis{before} the function is called.
If \sphinxstyleemphasis{idx} is out of range, \sphinxstyleemphasis{buffer} is set to an empty string and a warning
is issued.
\begin{quote}\begin{description}
\sphinxlineitem{Parameters}\begin{itemize}
\item {} 
\sphinxAtStartPar
\sphinxstylestrong{category}\sphinxstyleemphasis{ {[}}\sphinxstyleemphasis{character(len=*)}\sphinxstyleemphasis{{]}} :: category of IDs

\item {} 
\sphinxAtStartPar
\sphinxstylestrong{idx}\sphinxstyleemphasis{ {[}}\sphinxstyleemphasis{integer(c\_int)}\sphinxstyleemphasis{{]}} :: index of the ID in the list of \sphinxstyleemphasis{category}
styles (\(0 \leq idx < count\))

\item {} 
\sphinxAtStartPar
\sphinxstylestrong{buffer}\sphinxstyleemphasis{ {[}}\sphinxstyleemphasis{character(len=*)}\sphinxstyleemphasis{{]}} :: string into which to copy the name of the
style

\end{itemize}

\sphinxlineitem{Call to}
\sphinxAtStartPar
{\hyperref[\detokenize{Library_config:_CPPv414lammps_id_namePvPKciPci}]{\sphinxcrossref{\sphinxcode{\sphinxupquote{lammps\_id\_name()}}}}}

\end{description}\end{quote}

\end{fulllineitems}



\bigskip\hrule\bigskip

\index{plugin\_count() (fortran function)@\spxentry{plugin\_count()}\spxextra{fortran function}|spxpagem}

\begin{fulllineitems}
\phantomsection\label{\detokenize{Fortran:f/_/plugin_count}}
\pysigstartsignatures
\pysiglinewithargsret{\sphinxbfcode{\sphinxupquote{function  }}\sphinxbfcode{\sphinxupquote{plugin\_count}}}{}{}
\pysigstopsignatures
\sphinxAtStartPar
This function counts the number of loaded plugins.

\sphinxAtStartPar
\DUrole{versionmodified,added}{Added in version 3Nov2022.}
\begin{quote}\begin{description}
\sphinxlineitem{Call to}
\sphinxAtStartPar
\sphinxcode{\sphinxupquote{lammps\_plugin\_count()}}

\sphinxlineitem{Return}
\sphinxAtStartPar
\sphinxstylestrong{n}\sphinxstyleemphasis{ {[}}\sphinxstyleemphasis{integer(c\_int)}\sphinxstyleemphasis{{]}} :: number of loaded plugins

\end{description}\end{quote}

\end{fulllineitems}



\bigskip\hrule\bigskip

\index{plugin\_name() (fortran subroutine)@\spxentry{plugin\_name()}\spxextra{fortran subroutine}|spxpagem}

\begin{fulllineitems}
\phantomsection\label{\detokenize{Fortran:f/_/plugin_name}}
\pysigstartsignatures
\pysiglinewithargsret{\sphinxbfcode{\sphinxupquote{subroutine  }}\sphinxbfcode{\sphinxupquote{plugin\_name}}}{\sphinxparam{idx}\sphinxparamcomma \sphinxparam{stylebuf}\sphinxparamcomma \sphinxparam{namebuf}}{}
\pysigstopsignatures
\sphinxAtStartPar
Look up the style and name of a plugin by its index in the list of plugins.

\sphinxAtStartPar
\DUrole{versionmodified,added}{Added in version 3Nov2022.}

\sphinxAtStartPar
This function copies the name of the \sphinxstyleemphasis{style} plugin with the index \sphinxstyleemphasis{idx}
into the provided C\sphinxhyphen{}style string buffer.  The length of the buffer must be
provided as \sphinxstyleemphasis{buf\_size} argument.  If the name of the style exceeds the
length of the buffer, it will be truncated accordingly.  If the index is out
of range, both strings are set to the empty string and a warning is
printed.
\begin{quote}\begin{description}
\sphinxlineitem{Parameters}\begin{itemize}
\item {} 
\sphinxAtStartPar
\sphinxstylestrong{idx}\sphinxstyleemphasis{ {[}}\sphinxstyleemphasis{integer(c\_int)}\sphinxstyleemphasis{{]}} :: index of the plugin in the list all or
\sphinxstyleemphasis{style} plugins

\item {} 
\sphinxAtStartPar
\sphinxstylestrong{stylebuf}\sphinxstyleemphasis{ {[}}\sphinxstyleemphasis{character(len=*)}\sphinxstyleemphasis{{]}} :: string into which to copy the style of the
plugin

\item {} 
\sphinxAtStartPar
\sphinxstylestrong{namebuf}\sphinxstyleemphasis{ {[}}\sphinxstyleemphasis{character(len=*)}\sphinxstyleemphasis{{]}} :: string into which to copy the style of the
plugin

\end{itemize}

\sphinxlineitem{Call to}
\sphinxAtStartPar
\sphinxcode{\sphinxupquote{lammps\_plugin\_name()}}

\end{description}\end{quote}

\end{fulllineitems}



\bigskip\hrule\bigskip

\index{encode\_image\_flags() (fortran function)@\spxentry{encode\_image\_flags()}\spxextra{fortran function}|spxpagem}

\begin{fulllineitems}
\phantomsection\label{\detokenize{Fortran:f/_/encode_image_flags}}
\pysigstartsignatures
\pysiglinewithargsret{\sphinxbfcode{\sphinxupquote{function  }}\sphinxbfcode{\sphinxupquote{encode\_image\_flags}}}{\sphinxparam{ix}\sphinxparamcomma \sphinxparam{iy}\sphinxparamcomma \sphinxparam{iz}}{}
\pysigstopsignatures
\sphinxAtStartPar
Encodes three integer image flags into a single imageint.

\sphinxAtStartPar
\DUrole{versionmodified,added}{Added in version 3Nov2022.}

\sphinxAtStartPar
This function performs the bit\sphinxhyphen{}shift, addition, and bit\sphinxhyphen{}wise OR operations
necessary to combine the values of three integers representing the image
flags in the \(x\)\sphinxhyphen{}, \(y\)\sphinxhyphen{}, and \(z\)\sphinxhyphen{}directions. Unless LAMMPS
is compiled with \sphinxcode{\sphinxupquote{\sphinxhyphen{}DLAMMPS\_BIGBIG}}, those integers are limited to 10\sphinxhyphen{}bit
signed integers \([-512,512)\). If \sphinxcode{\sphinxupquote{\sphinxhyphen{}DLAMMPS\_BIGBIG}} was used when
compiling, then the return value is of kind \sphinxcode{\sphinxupquote{c\_int64\_t}} instead of
kind \sphinxcode{\sphinxupquote{c\_int}}, and the valid range for the individual image flags becomes
\([-1048576,1048575)\) (i.e., the range of a 21\sphinxhyphen{}bit signed integer).
There is no check on whether the arguments conform to these requirements;
values out of range will simply be wrapped back into the interval.
\begin{quote}\begin{description}
\sphinxlineitem{Parameters}\begin{itemize}
\item {} 
\sphinxAtStartPar
\sphinxstylestrong{ix}\sphinxstyleemphasis{ {[}}\sphinxstyleemphasis{integer(c\_int)}\sphinxstyleemphasis{{]}} :: image flag in \(x\)\sphinxhyphen{}direction

\item {} 
\sphinxAtStartPar
\sphinxstylestrong{iy}\sphinxstyleemphasis{ {[}}\sphinxstyleemphasis{integer(c\_int)}\sphinxstyleemphasis{{]}} :: image flag in \(y\)\sphinxhyphen{}direction

\item {} 
\sphinxAtStartPar
\sphinxstylestrong{iz}\sphinxstyleemphasis{ {[}}\sphinxstyleemphasis{integer(c\_int)}\sphinxstyleemphasis{{]}} :: image flag in \(z\)\sphinxhyphen{}direction

\end{itemize}

\sphinxlineitem{Return}
\sphinxAtStartPar
\sphinxstylestrong{imageint}\sphinxstyleemphasis{ {[}}\sphinxstyleemphasis{integer(kind=*)}\sphinxstyleemphasis{{]}} :: encoded image flag. *The \sphinxcode{\sphinxupquote{KIND}} parameter
is \sphinxcode{\sphinxupquote{c\_int}} unless LAMMPS was built with \sphinxcode{\sphinxupquote{\sphinxhyphen{}DLAMMPS\_BIGBIG}}, in which
case it is \sphinxcode{\sphinxupquote{c\_int64\_t}}.

\end{description}\end{quote}

\begin{sphinxadmonition}{note}{Note:}
\sphinxAtStartPar
The fact that the programmer does not know the \sphinxcode{\sphinxupquote{KIND}} parameter of the
return value until compile time means that it is impossible to define an
interface that works for both sizes of \sphinxcode{\sphinxupquote{imageint}}. One side effect of
this is that you must assign the return value of this function to a
variable; it cannot be used as the argument to another function or as part
of an array constructor. For example,

\begin{sphinxVerbatim}[commandchars=\\\{\}]
\PYG{n}{my\PYGZus{}images}\PYG{+w}{ }\PYG{o}{=}\PYG{+w}{ }\PYG{p}{[}\PYG{n}{lmp}\PYG{p}{\PYGZpc{}}\PYG{n}{encode\PYGZus{}image\PYGZus{}flags}\PYG{p}{(}\PYG{l+m+mi}{0}\PYG{p}{,}\PYG{l+m+mi}{0}\PYG{p}{,}\PYG{l+m+mi}{0}\PYG{p}{)}\PYG{p}{,}\PYG{+w}{ }\PYG{n}{lmp}\PYG{p}{\PYGZpc{}}\PYG{n}{encode\PYGZus{}image\PYGZus{}flags}\PYG{p}{(}\PYG{l+m+mi}{1}\PYG{p}{,}\PYG{l+m+mi}{0}\PYG{p}{,}\PYG{l+m+mi}{0}\PYG{p}{)}\PYG{p}{]}
\end{sphinxVerbatim}

\sphinxAtStartPar
will \sphinxstyleemphasis{not} work; instead, do something like

\begin{sphinxVerbatim}[commandchars=\\\{\}]
\PYG{n}{my\PYGZus{}images}\PYG{p}{(}\PYG{l+m+mi}{1}\PYG{p}{)}\PYG{+w}{ }\PYG{o}{=}\PYG{+w}{ }\PYG{n}{lmp}\PYG{p}{\PYGZpc{}}\PYG{n}{encode\PYGZus{}image\PYGZus{}flags}\PYG{p}{(}\PYG{l+m+mi}{0}\PYG{p}{,}\PYG{l+m+mi}{0}\PYG{p}{,}\PYG{l+m+mi}{0}\PYG{p}{)}
\PYG{n}{my\PYGZus{}images}\PYG{p}{(}\PYG{l+m+mi}{2}\PYG{p}{)}\PYG{+w}{ }\PYG{o}{=}\PYG{+w}{ }\PYG{n}{lmp}\PYG{p}{\PYGZpc{}}\PYG{n}{encode\PYGZus{}image\PYGZus{}flags}\PYG{p}{(}\PYG{l+m+mi}{1}\PYG{p}{,}\PYG{l+m+mi}{0}\PYG{p}{,}\PYG{l+m+mi}{0}\PYG{p}{)}
\end{sphinxVerbatim}
\end{sphinxadmonition}

\end{fulllineitems}



\bigskip\hrule\bigskip

\index{decode\_image\_flags() (fortran subroutine)@\spxentry{decode\_image\_flags()}\spxextra{fortran subroutine}|spxpagem}

\begin{fulllineitems}
\phantomsection\label{\detokenize{Fortran:f/_/decode_image_flags}}
\pysigstartsignatures
\pysiglinewithargsret{\sphinxbfcode{\sphinxupquote{subroutine  }}\sphinxbfcode{\sphinxupquote{decode\_image\_flags}}}{\sphinxparam{image}\sphinxparamcomma \sphinxparam{flags}}{}
\pysigstopsignatures
\sphinxAtStartPar
This function does the reverse operation of \sphinxcode{\sphinxupquote{encode\_image\_flags()}}:
it takes the image flag and performs the bit\sphinxhyphen{}shift and bit\sphinxhyphen{}masking
operations to decode it and stores the resulting three integers into the
array \sphinxstyleemphasis{flags}.

\sphinxAtStartPar
\DUrole{versionmodified,added}{Added in version 3Nov2022.}
\begin{quote}\begin{description}
\sphinxlineitem{Parameters}\begin{itemize}
\item {} 
\sphinxAtStartPar
\sphinxstylestrong{image}\sphinxstyleemphasis{ {[}}\sphinxstyleemphasis{integer(kind=*)}\sphinxstyleemphasis{{]}} :: encoded image flag. *The \sphinxcode{\sphinxupquote{KIND}} parameter is
either \sphinxcode{\sphinxupquote{c\_int}} or, if LAMMPS was compiled with \sphinxcode{\sphinxupquote{\sphinxhyphen{}DLAMMPS\_BIGBIG}},
\sphinxcode{\sphinxupquote{c\_int64\_t}}. Kind compatibility is checked at run\sphinxhyphen{}time.

\item {} 
\sphinxAtStartPar
\sphinxstylestrong{flags}\sphinxstyleemphasis{ {[}}\sphinxstyleemphasis{integer(c\_int)}\sphinxstyleemphasis{,}\sphinxstyleemphasis{dimension(3)}\sphinxstyleemphasis{{]}} :: three\sphinxhyphen{}element vector where the
decoded image flags will be stored.

\end{itemize}

\end{description}\end{quote}

\end{fulllineitems}



\bigskip\hrule\bigskip

\index{set\_fix\_external\_callback() (fortran subroutine)@\spxentry{set\_fix\_external\_callback()}\spxextra{fortran subroutine}|spxpagem}

\begin{fulllineitems}
\phantomsection\label{\detokenize{Fortran:f/_/set_fix_external_callback}}
\pysigstartsignatures
\pysiglinewithargsret{\sphinxbfcode{\sphinxupquote{subroutine  }}\sphinxbfcode{\sphinxupquote{set\_fix\_external\_callback}}}{\sphinxparam{id}\sphinxparamcomma \sphinxparam{callback}\sphinxparamcomma \sphinxparam{caller}}{}
\pysigstopsignatures
\sphinxAtStartPar
Set the callback function for a {\hyperref[\detokenize{fix_external::doc}]{\sphinxcrossref{\DUrole{doc}{fix external}}}} instance
with the given ID.

\sphinxAtStartPar
\DUrole{versionmodified,added}{Added in version 22Dec2022.}

\sphinxAtStartPar
Fix {\hyperref[\detokenize{fix_external::doc}]{\sphinxcrossref{\DUrole{doc}{external}}}} allows programs that are running LAMMPS
through its library interface to modify certain LAMMPS properties on
specific time steps, similar to the way other fixes do.

\sphinxAtStartPar
This subroutine sets the callback function for use with the “pf/callback”
mode. The function should have Fortran language bindings with the following
interface, which depends on how LAMMPS was compiled:

\begin{sphinxVerbatim}[commandchars=\\\{\}]
\PYG{k}{ABSTRACT }\PYG{k}{INTERFACE}
\PYG{k}{  }\PYG{k}{SUBROUTINE }\PYG{n}{external\PYGZus{}callback}\PYG{p}{(}\PYG{n}{caller}\PYG{p}{,}\PYG{+w}{ }\PYG{n}{timestep}\PYG{p}{,}\PYG{+w}{ }\PYG{n}{ids}\PYG{p}{,}\PYG{+w}{ }\PYG{n}{x}\PYG{p}{,}\PYG{+w}{ }\PYG{n}{fexternal}\PYG{p}{)}
\PYG{+w}{    }\PYG{k}{USE}\PYG{p}{,}\PYG{+w}{ }\PYG{k}{INTRINSIC}\PYG{+w}{ }\PYG{k+kd}{::}\PYG{+w}{ }\PYG{n+nb}{ISO\PYGZus{}C\PYGZus{}BINDING}\PYG{p}{,}\PYG{+w}{ }\PYG{k}{ONLY}\PYG{+w}{ }\PYG{p}{:}\PYG{+w}{ }\PYG{k+kt}{c\PYGZus{}int}\PYG{p}{,}\PYG{+w}{ }\PYG{k+kt}{c\PYGZus{}double}\PYG{p}{,}\PYG{+w}{ }\PYG{k+kt}{c\PYGZus{}int64\PYGZus{}t}
\PYG{k+kt}{    }\PYG{k}{CLASS}\PYG{p}{(}\PYG{o}{*}\PYG{p}{)}\PYG{p}{,}\PYG{+w}{ }\PYG{k}{INTENT}\PYG{p}{(}\PYG{n}{INOUT}\PYG{p}{)}\PYG{+w}{ }\PYG{k+kd}{::}\PYG{+w}{ }\PYG{n}{caller}
\PYG{+w}{    }\PYG{k+kt}{INTEGER}\PYG{p}{(}\PYG{n}{c\PYGZus{}bigint}\PYG{p}{)}\PYG{p}{,}\PYG{+w}{ }\PYG{k}{INTENT}\PYG{p}{(}\PYG{n}{IN}\PYG{p}{)}\PYG{+w}{ }\PYG{k+kd}{::}\PYG{+w}{ }\PYG{n}{timestep}
\PYG{+w}{    }\PYG{k+kt}{INTEGER}\PYG{p}{(}\PYG{n}{c\PYGZus{}tagint}\PYG{p}{)}\PYG{p}{,}\PYG{+w}{ }\PYG{k}{DIMENSION}\PYG{p}{(}\PYG{p}{:}\PYG{p}{)}\PYG{p}{,}\PYG{+w}{ }\PYG{k}{INTENT}\PYG{p}{(}\PYG{n}{IN}\PYG{p}{)}\PYG{+w}{ }\PYG{k+kd}{::}\PYG{+w}{ }\PYG{n}{ids}
\PYG{+w}{    }\PYG{k+kt}{REAL}\PYG{p}{(}\PYG{k+kt}{c\PYGZus{}double}\PYG{p}{)}\PYG{p}{,}\PYG{+w}{ }\PYG{k}{DIMENSION}\PYG{p}{(}\PYG{p}{:}\PYG{p}{,}\PYG{p}{:}\PYG{p}{)}\PYG{p}{,}\PYG{+w}{ }\PYG{k}{INTENT}\PYG{p}{(}\PYG{n}{IN}\PYG{p}{)}\PYG{+w}{ }\PYG{k+kd}{::}\PYG{+w}{ }\PYG{n}{x}
\PYG{+w}{    }\PYG{k+kt}{REAL}\PYG{p}{(}\PYG{k+kt}{c\PYGZus{}double}\PYG{p}{)}\PYG{p}{,}\PYG{+w}{ }\PYG{k}{DIMENSION}\PYG{p}{(}\PYG{p}{:}\PYG{p}{,}\PYG{p}{:}\PYG{p}{)}\PYG{p}{,}\PYG{+w}{ }\PYG{k}{INTENT}\PYG{p}{(}\PYG{n}{OUT}\PYG{p}{)}\PYG{+w}{ }\PYG{k+kd}{::}\PYG{+w}{ }\PYG{n}{fexternal}
\PYG{+w}{  }\PYG{k}{END }\PYG{k}{SUBROUTINE }\PYG{n}{external\PYGZus{}callback}
\PYG{k}{END }\PYG{k}{INTERFACE}
\end{sphinxVerbatim}

\sphinxAtStartPar
where \sphinxcode{\sphinxupquote{c\_bigint}} is \sphinxcode{\sphinxupquote{c\_int}} if \sphinxcode{\sphinxupquote{\sphinxhyphen{}DLAMMPS\_SMALLSMALL}} was used and
\sphinxcode{\sphinxupquote{c\_int64\_t}} otherwise; and \sphinxcode{\sphinxupquote{c\_tagint}} is \sphinxcode{\sphinxupquote{c\_int64\_t}} if
\sphinxcode{\sphinxupquote{\sphinxhyphen{}DLAMMPS\_BIGBIG}} was used and \sphinxcode{\sphinxupquote{c\_int}} otherwise.

\sphinxAtStartPar
The argument \sphinxstyleemphasis{caller} to \sphinxcode{\sphinxupquote{set\_fix\_external\_callback()}} is unlimited
polymorphic (i.e., it can be any Fortran object you want to pass to the
calling function) and will be available as the first argument to the
callback function. It can be your LAMMPS instance, which you might need if
the callback function needs access to the library interface.
The argument must be a scalar; to pass non\sphinxhyphen{}scalar data, wrap those data in
a derived type and pass an instance of the derived type to \sphinxstyleemphasis{caller}.

\sphinxAtStartPar
The array \sphinxstyleemphasis{ids} is an array of length \sphinxstyleemphasis{nlocal} (as accessed from the
\sphinxcode{\sphinxupquote{Atom}} class or through \sphinxcode{\sphinxupquote{extract\_global()}}). The arrays
\sphinxstyleemphasis{x} and \sphinxstyleemphasis{fexternal} are \(3 \times {}\)\sphinxstyleemphasis{nlocal} arrays; these are
transposed from what they would look like in C (see note about array index
order at \sphinxcode{\sphinxupquote{extract\_atom()}}).

\sphinxAtStartPar
The callback mechanism is one of two ways that forces can be applied to a
simulation with the help of {\hyperref[\detokenize{fix_external::doc}]{\sphinxcrossref{\DUrole{doc}{fix external}}}}. The
alternative is \sphinxstyleemphasis{array} mode, where one calls
\sphinxcode{\sphinxupquote{fix\_external\_get\_force()}}.

\sphinxAtStartPar
Please see the documentation for {\hyperref[\detokenize{fix_external::doc}]{\sphinxcrossref{\DUrole{doc}{fix external}}}} for
more information about how to use the fix and couple it with external
programs.
\begin{quote}\begin{description}
\sphinxlineitem{Parameters}\begin{itemize}
\item {} 
\sphinxAtStartPar
\sphinxstylestrong{id}\sphinxstyleemphasis{ {[}}\sphinxstyleemphasis{character(len=*)}\sphinxstyleemphasis{{]}} :: ID of {\hyperref[\detokenize{fix_external::doc}]{\sphinxcrossref{\DUrole{doc}{fix external}}}} instance

\item {} 
\sphinxAtStartPar
\sphinxstylestrong{callback}\sphinxstyleemphasis{ {[}}\sphinxstyleemphasis{external}\sphinxstyleemphasis{{]}} :: subroutine {\hyperref[\detokenize{fix_external::doc}]{\sphinxcrossref{\DUrole{doc}{fix external}}}} should call

\item {} 
\sphinxAtStartPar
\sphinxstylestrong{caller}\sphinxstyleemphasis{ {[}}\sphinxstyleemphasis{class(*)}\sphinxstyleemphasis{,}\sphinxstyleemphasis{optional}\sphinxstyleemphasis{{]}} :: object you wish to pass to the callback
procedure (must be a scalar; see note)

\end{itemize}

\sphinxlineitem{Call to}
\sphinxAtStartPar
{\hyperref[\detokenize{Library_utility:_CPPv432lammps_set_fix_external_callbackPvPKc16FixExternalFnPtrPv}]{\sphinxcrossref{\sphinxcode{\sphinxupquote{lammps\_set\_fix\_external\_callback()}}}}}

\end{description}\end{quote}

\begin{sphinxadmonition}{note}{Note:}
\sphinxAtStartPar
The interface for your callback function must match types precisely
with the abstract interface block given above. \sphinxstylestrong{The compiler probably
will not be able to check this for you.} In particular, the first
argument (“caller”) must be of type \sphinxcode{\sphinxupquote{CLASS(*)}} or you will probably
get a segmentation fault or at least a misinterpretation of whatever is
in memory there. You can resolve the object using the \sphinxcode{\sphinxupquote{SELECT TYPE}}
construct. An example callback function (assuming LAMMPS was compiled
with \sphinxcode{\sphinxupquote{\sphinxhyphen{}DLAMMPS\_SMALLBIG}}) that applies something akin to Hooke’s Law
(with each atom having a different \sphinxstyleemphasis{k} value) is shown below.

\begin{sphinxVerbatim}[commandchars=\\\{\}]
\PYG{k}{MODULE }\PYG{n}{stuff}
\PYG{+w}{  }\PYG{k}{USE}\PYG{p}{,}\PYG{+w}{ }\PYG{k}{INTRINSIC}\PYG{+w}{ }\PYG{k+kd}{::}\PYG{+w}{ }\PYG{n+nb}{ISO\PYGZus{}C\PYGZus{}BINDING}\PYG{p}{,}\PYG{+w}{ }\PYG{k}{ONLY}\PYG{+w}{ }\PYG{p}{:}\PYG{+w}{ }\PYG{k+kt}{c\PYGZus{}int}\PYG{p}{,}\PYG{+w}{ }\PYG{k+kt}{c\PYGZus{}double}\PYG{p}{,}\PYG{+w}{ }\PYG{k+kt}{c\PYGZus{}int64\PYGZus{}t}
\PYG{k+kt}{  }\PYG{k}{USE}\PYG{p}{,}\PYG{+w}{ }\PYG{k}{INTRINSIC}\PYG{+w}{ }\PYG{k+kd}{::}\PYG{+w}{ }\PYG{n}{ISO\PYGZus{}FORTRAN\PYGZus{}ENV}\PYG{p}{,}\PYG{+w}{ }\PYG{k}{ONLY}\PYG{+w}{ }\PYG{p}{:}\PYG{+w}{ }\PYG{n}{error\PYGZus{}unit}
\PYG{+w}{  }\PYG{k}{IMPLICIT }\PYG{k}{NONE}

\PYG{k}{  }\PYG{k}{TYPE }\PYG{n}{shield}
\PYG{+w}{     }\PYG{k+kt}{REAL}\PYG{p}{(}\PYG{k+kt}{c\PYGZus{}double}\PYG{p}{)}\PYG{p}{,}\PYG{+w}{ }\PYG{k}{DIMENSION}\PYG{p}{(}\PYG{p}{:}\PYG{p}{,}\PYG{p}{:}\PYG{p}{)}\PYG{p}{,}\PYG{+w}{ }\PYG{k}{ALLOCATABLE}\PYG{+w}{ }\PYG{k+kd}{::}\PYG{+w}{ }\PYG{n}{k}
\PYG{+w}{     }\PYG{c}{! assume k gets allocated to dimension(3,nlocal) at some point}
\PYG{+w}{     }\PYG{c}{! and assigned values}
\PYG{+w}{  }\PYG{k}{END }\PYG{k}{TYPE }\PYG{n}{shield}

\PYG{+w}{  }\PYG{k}{SUBROUTINE }\PYG{n}{my\PYGZus{}callback}\PYG{p}{(}\PYG{n}{caller}\PYG{p}{,}\PYG{+w}{ }\PYG{n}{timestep}\PYG{p}{,}\PYG{+w}{ }\PYG{n}{ids}\PYG{p}{,}\PYG{+w}{ }\PYG{n}{x}\PYG{p}{,}\PYG{+w}{ }\PYG{n}{fexternal}\PYG{p}{)}
\PYG{+w}{    }\PYG{k}{CLASS}\PYG{p}{(}\PYG{o}{*}\PYG{p}{)}\PYG{p}{,}\PYG{+w}{ }\PYG{k}{INTENT}\PYG{p}{(}\PYG{n}{INOUT}\PYG{p}{)}\PYG{+w}{ }\PYG{k+kd}{::}\PYG{+w}{ }\PYG{n}{caller}
\PYG{+w}{    }\PYG{k+kt}{INTEGER}\PYG{p}{(}\PYG{k+kt}{c\PYGZus{}int}\PYG{p}{)}\PYG{p}{,}\PYG{+w}{ }\PYG{k}{INTENT}\PYG{p}{(}\PYG{n}{IN}\PYG{p}{)}\PYG{+w}{ }\PYG{k+kd}{::}\PYG{+w}{ }\PYG{n}{timestep}
\PYG{+w}{    }\PYG{k+kt}{INTEGER}\PYG{p}{(}\PYG{k+kt}{c\PYGZus{}int64\PYGZus{}t}\PYG{p}{)}\PYG{p}{,}\PYG{+w}{ }\PYG{k}{INTENT}\PYG{p}{(}\PYG{n}{IN}\PYG{p}{)}\PYG{+w}{ }\PYG{k+kd}{::}\PYG{+w}{ }\PYG{n}{ids}
\PYG{+w}{    }\PYG{k+kt}{REAL}\PYG{p}{(}\PYG{k+kt}{c\PYGZus{}double}\PYG{p}{)}\PYG{p}{,}\PYG{+w}{ }\PYG{k}{INTENT}\PYG{p}{(}\PYG{n}{IN}\PYG{p}{)}\PYG{+w}{ }\PYG{k+kd}{::}\PYG{+w}{ }\PYG{n}{x}\PYG{p}{(}\PYG{p}{:}\PYG{p}{,}\PYG{p}{:}\PYG{p}{)}
\PYG{+w}{    }\PYG{k+kt}{REAL}\PYG{p}{(}\PYG{k+kt}{c\PYGZus{}double}\PYG{p}{)}\PYG{p}{,}\PYG{+w}{ }\PYG{k}{INTENT}\PYG{p}{(}\PYG{n}{OUT}\PYG{p}{)}\PYG{+w}{ }\PYG{k+kd}{::}\PYG{+w}{ }\PYG{n}{fexternal}\PYG{p}{(}\PYG{p}{:}\PYG{p}{,}\PYG{p}{:}\PYG{p}{)}

\PYG{+w}{    }\PYG{k}{SELECT }\PYG{k}{TYPE}\PYG{+w}{ }\PYG{p}{(}\PYG{n}{caller}\PYG{p}{)}
\PYG{+w}{      }\PYG{k}{TYPE }\PYG{k}{IS}\PYG{+w}{ }\PYG{p}{(}\PYG{n}{shield}\PYG{p}{)}
\PYG{+w}{         }\PYG{n}{fexternal}\PYG{+w}{ }\PYG{o}{=}\PYG{+w}{ }\PYG{o}{\PYGZhy{}}\PYG{+w}{ }\PYG{n}{caller}\PYG{p}{\PYGZpc{}}\PYG{n}{k}\PYG{+w}{ }\PYG{o}{*}\PYG{+w}{ }\PYG{n}{x}
\PYG{+w}{      }\PYG{k}{CLASS }\PYG{n}{DEFAULT}
\PYG{+w}{         }\PYG{k}{WRITE}\PYG{p}{(}\PYG{n}{error\PYGZus{}unit}\PYG{p}{,}\PYG{o}{*}\PYG{p}{)}\PYG{+w}{ }\PYG{l+s+s1}{\PYGZsq{}UH OH...\PYGZsq{}}
\PYG{+w}{    }\PYG{k}{END }\PYG{k}{SELECT}
\PYG{k}{  }\PYG{k}{END }\PYG{k}{SUBROUTINE }\PYG{n}{my\PYGZus{}callback}
\PYG{k}{END }\PYG{k}{MODULE }\PYG{n}{stuff}

\PYG{c}{! then, when assigning the callback function, do this:}
\PYG{k}{PROGRAM }\PYG{n}{example}
\PYG{+w}{  }\PYG{k}{USE }\PYG{n}{LIBLAMMPS}
\PYG{+w}{  }\PYG{k}{USE }\PYG{n}{stuff}
\PYG{+w}{  }\PYG{k}{TYPE}\PYG{p}{(}\PYG{n}{lammps}\PYG{p}{)}\PYG{+w}{ }\PYG{k+kd}{::}\PYG{+w}{ }\PYG{n}{lmp}
\PYG{+w}{  }\PYG{k}{TYPE}\PYG{p}{(}\PYG{n}{shield}\PYG{p}{)}\PYG{+w}{ }\PYG{k+kd}{::}\PYG{+w}{ }\PYG{n}{my\PYGZus{}shield}
\PYG{+w}{  }\PYG{n}{lmp}\PYG{+w}{ }\PYG{o}{=}\PYG{+w}{ }\PYG{n}{lammps}\PYG{p}{(}\PYG{p}{)}
\PYG{+w}{  }\PYG{k}{CALL }\PYG{n}{lmp}\PYG{p}{\PYGZpc{}}\PYG{n}{command}\PYG{p}{(}\PYG{l+s+s1}{\PYGZsq{}fix ext all external pf/callback 1 1\PYGZsq{}}\PYG{p}{)}
\PYG{+w}{  }\PYG{k}{CALL }\PYG{n}{lmp}\PYG{p}{\PYGZpc{}}\PYG{n}{set\PYGZus{}fix\PYGZus{}external\PYGZus{}callback}\PYG{p}{(}\PYG{l+s+s1}{\PYGZsq{}ext\PYGZsq{}}\PYG{p}{,}\PYG{+w}{ }\PYG{n}{my\PYGZus{}callback}\PYG{p}{,}\PYG{+w}{ }\PYG{n}{my\PYGZus{}shield}\PYG{p}{)}
\PYG{k}{END }\PYG{k}{PROGRAM }\PYG{n}{example}
\end{sphinxVerbatim}
\end{sphinxadmonition}

\end{fulllineitems}



\bigskip\hrule\bigskip

\index{fix\_external\_get\_force() (fortran function)@\spxentry{fix\_external\_get\_force()}\spxextra{fortran function}|spxpagem}

\begin{fulllineitems}
\phantomsection\label{\detokenize{Fortran:f/_/fix_external_get_force}}
\pysigstartsignatures
\pysiglinewithargsret{\sphinxbfcode{\sphinxupquote{function  }}\sphinxbfcode{\sphinxupquote{fix\_external\_get\_force}}}{\sphinxparam{id}}{}
\pysigstopsignatures
\sphinxAtStartPar
Get pointer to the force array storage in a fix external instance with the
given ID.

\sphinxAtStartPar
\DUrole{versionmodified,added}{Added in version 22Dec2022.}

\sphinxAtStartPar
Fix {\hyperref[\detokenize{fix_external::doc}]{\sphinxcrossref{\DUrole{doc}{external}}}} allows programs that are running LAMMPS
through its library interfaces to add or modify certain LAMMPS properties on
specific time steps, similar to the way other fixes do.

\sphinxAtStartPar
This function provides access to the per\sphinxhyphen{}atom force storage in a fix
external instance with the given fix\sphinxhyphen{}ID to be added to the individual
atoms when using the “pf/array” mode.  The \sphinxstyleemphasis{fexternal} array can be
accessed like other “native” per\sphinxhyphen{}atom arrays accessible via the
\sphinxcode{\sphinxupquote{extract\_atom()}} function.  Please note that the array
stores the forces for \sphinxstyleemphasis{local} atoms for each MPI rank, in the order
determined by the neighbor list build.  Because the underlying data
structures can change as well as the order of atom as they migrate between
MPI processes because of the domain decomposition parallelization, this
function should be always called immediately before the forces are going to
be set to get an up\sphinxhyphen{}to\sphinxhyphen{}date pointer.  You can use, for example,
\sphinxcode{\sphinxupquote{extract\_setting()}} to obtain the number of local atoms \sphinxtitleref{nlocal} and
then assume the dimensions of the returned force array as
\sphinxcode{\sphinxupquote{REAL(c\_double) :: force(3,nlocal)}}.

\sphinxAtStartPar
This function is an alternative to the callback mechanism in fix external
set up by \sphinxcode{\sphinxupquote{set\_fix\_external\_callback()}}. The main difference is that
this mechanism can be used when forces are to be pre\sphinxhyphen{}computed and the
control alternates between LAMMPS and the external driver, while the
callback mechanism can call an external subroutine to compute the force when
the fix is triggered and needs them.

\sphinxAtStartPar
Please see the documentation for {\hyperref[\detokenize{fix_external::doc}]{\sphinxcrossref{\DUrole{doc}{fix external}}}} for
more information about how to use the fix and how to couple it with an
external program.
\begin{quote}\begin{description}
\sphinxlineitem{Parameters}
\sphinxAtStartPar
\sphinxstylestrong{id}\sphinxstyleemphasis{ {[}}\sphinxstyleemphasis{character(len=*)}\sphinxstyleemphasis{{]}} :: ID of {\hyperref[\detokenize{fix_external::doc}]{\sphinxcrossref{\DUrole{doc}{fix external}}}} instance

\sphinxlineitem{Call to}
\sphinxAtStartPar
{\hyperref[\detokenize{Library_utility:_CPPv429lammps_fix_external_get_forcePvPKc}]{\sphinxcrossref{\sphinxcode{\sphinxupquote{lammps\_fix\_external\_get\_force()}}}}}

\sphinxlineitem{Return}
\sphinxAtStartPar
\sphinxstylestrong{fexternal}\sphinxstyleemphasis{ {[}}\sphinxstyleemphasis{real(c\_double),dimension(3,nlocal)}\sphinxstyleemphasis{{]}} :: pointer to the per\sphinxhyphen{}atom force array allocated by the fix

\end{description}\end{quote}

\end{fulllineitems}



\bigskip\hrule\bigskip

\index{fix\_external\_set\_energy\_global() (fortran subroutine)@\spxentry{fix\_external\_set\_energy\_global()}\spxextra{fortran subroutine}|spxpagem}

\begin{fulllineitems}
\phantomsection\label{\detokenize{Fortran:f/_/fix_external_set_energy_global}}
\pysigstartsignatures
\pysiglinewithargsret{\sphinxbfcode{\sphinxupquote{subroutine  }}\sphinxbfcode{\sphinxupquote{fix\_external\_set\_energy\_global}}}{\sphinxparam{id}\sphinxparamcomma \sphinxparam{eng}}{}
\pysigstopsignatures
\sphinxAtStartPar
Set the global energy contribution for a {\hyperref[\detokenize{fix_external::doc}]{\sphinxcrossref{\DUrole{doc}{fix external}}}}
instance with the given ID.

\sphinxAtStartPar
\DUrole{versionmodified,added}{Added in version 22Dec2022.}

\sphinxAtStartPar
This is a companion function to \sphinxcode{\sphinxupquote{set\_fix\_external\_callback()}}
and \sphinxcode{\sphinxupquote{fix\_external\_get\_force()}} that also sets the contribution to the
global energy from the external program.  The value of the \sphinxstyleemphasis{eng} argument
will be stored in the fix and applied on the current and all following
time steps until changed by another call to this function. The energy is in
energy units as determined by the current {\hyperref[\detokenize{units::doc}]{\sphinxcrossref{\DUrole{doc}{units}}}} settings and
is the \sphinxstylestrong{total} energy of the contribution.  Thus, when running in
parallel, all MPI processes have to call this function with the \sphinxstylestrong{same}
value, and this will be returned as a scalar property of the fix external
instance when accessed in LAMMPS input commands or from variables.

\sphinxAtStartPar
Please see the documentation for {\hyperref[\detokenize{fix_external::doc}]{\sphinxcrossref{\DUrole{doc}{fix external}}}} for more
information about how to use the fix and how to couple it with an external
program.
\begin{quote}\begin{description}
\sphinxlineitem{Parameters}\begin{itemize}
\item {} 
\sphinxAtStartPar
\sphinxstylestrong{id}\sphinxstyleemphasis{ {[}}\sphinxstyleemphasis{character(len=*)}\sphinxstyleemphasis{{]}} :: fix ID of fix external instance

\item {} 
\sphinxAtStartPar
\sphinxstylestrong{eng}\sphinxstyleemphasis{ {[}}\sphinxstyleemphasis{real(c\_double)}\sphinxstyleemphasis{{]}} :: total energy to be added to the global energy

\end{itemize}

\sphinxlineitem{Call to}
\sphinxAtStartPar
{\hyperref[\detokenize{Library_utility:_CPPv437lammps_fix_external_set_energy_globalPvPKcd}]{\sphinxcrossref{\sphinxcode{\sphinxupquote{lammps\_fix\_external\_set\_energy\_global()}}}}}

\end{description}\end{quote}

\end{fulllineitems}



\bigskip\hrule\bigskip

\index{fix\_external\_set\_virial\_global() (fortran subroutine)@\spxentry{fix\_external\_set\_virial\_global()}\spxextra{fortran subroutine}|spxpagem}

\begin{fulllineitems}
\phantomsection\label{\detokenize{Fortran:f/_/fix_external_set_virial_global}}
\pysigstartsignatures
\pysiglinewithargsret{\sphinxbfcode{\sphinxupquote{subroutine  }}\sphinxbfcode{\sphinxupquote{fix\_external\_set\_virial\_global}}}{\sphinxparam{id}\sphinxparamcomma \sphinxparam{virial}}{}
\pysigstopsignatures
\sphinxAtStartPar
Set the global virial contribution for a fix external instance with the
given ID.

\sphinxAtStartPar
\DUrole{versionmodified,added}{Added in version 22Dec2022.}

\sphinxAtStartPar
This is a companion function to \sphinxcode{\sphinxupquote{set\_fix\_external\_callback()}}
and \sphinxcode{\sphinxupquote{fix\_external\_get\_force()}} to set the contribution to the global
virial from an external program.

\sphinxAtStartPar
The six values of the \sphinxstyleemphasis{virial} array will be stored in the fix and applied
on the current and all following time steps until changed by another call
to this function. The components of the virial need to be stored in the
following order: \sphinxstyleemphasis{xx}, \sphinxstyleemphasis{yy}, \sphinxstyleemphasis{zz}, \sphinxstyleemphasis{xy}, \sphinxstyleemphasis{xz}, \sphinxstyleemphasis{yz}.  In LAMMPS, the virial
is stored internally as \sphinxtitleref{stress*volume} in units of \sphinxtitleref{pressure*volume} as
determined by the current {\hyperref[\detokenize{units::doc}]{\sphinxcrossref{\DUrole{doc}{units}}}} settings and is the
\sphinxstylestrong{total} contribution.  Thus, when running in parallel, all MPI processes
have to call this function with the \sphinxstylestrong{same} value, and this will then be
added by fix external.

\sphinxAtStartPar
Please see the documentation for {\hyperref[\detokenize{fix_external::doc}]{\sphinxcrossref{\DUrole{doc}{fix external}}}} for
more information about how to use the fix and how to couple it with an
external code.
\begin{quote}\begin{description}
\sphinxlineitem{Parameters}\begin{itemize}
\item {} 
\sphinxAtStartPar
\sphinxstylestrong{id}\sphinxstyleemphasis{ {[}}\sphinxstyleemphasis{character(len=*)}\sphinxstyleemphasis{{]}} :: fix ID of fix external instance

\item {} 
\sphinxAtStartPar
\sphinxstylestrong{virial}\sphinxstyleemphasis{ {[}}\sphinxstyleemphasis{real(c\_double)}\sphinxstyleemphasis{,}\sphinxstyleemphasis{dimension(6)}\sphinxstyleemphasis{{]}} :: the six global stress tensor
components to be added to the global virial

\end{itemize}

\sphinxlineitem{Call to}
\sphinxAtStartPar
{\hyperref[\detokenize{Library_utility:_CPPv437lammps_fix_external_set_virial_globalPvPKcPd}]{\sphinxcrossref{\sphinxcode{\sphinxupquote{lammps\_fix\_external\_set\_virial\_global()}}}}}

\end{description}\end{quote}

\end{fulllineitems}



\bigskip\hrule\bigskip

\index{fix\_external\_set\_energy\_peratom() (fortran subroutine)@\spxentry{fix\_external\_set\_energy\_peratom()}\spxextra{fortran subroutine}|spxpagem}

\begin{fulllineitems}
\phantomsection\label{\detokenize{Fortran:f/_/fix_external_set_energy_peratom}}
\pysigstartsignatures
\pysiglinewithargsret{\sphinxbfcode{\sphinxupquote{subroutine  }}\sphinxbfcode{\sphinxupquote{fix\_external\_set\_energy\_peratom}}}{\sphinxparam{id}\sphinxparamcomma \sphinxparam{eng}}{}
\pysigstopsignatures
\sphinxAtStartPar
Set the per\sphinxhyphen{}atom energy contribution for a fix external instance with the
given ID.

\sphinxAtStartPar
\DUrole{versionmodified,added}{Added in version 22Dec2022.}

\sphinxAtStartPar
This is a companion function to \sphinxcode{\sphinxupquote{set\_fix\_external\_callback()}} to set
the per\sphinxhyphen{}atom energy contribution due to the fix from the external program as
part of the callback function.  For this to work, the LAMMPS object must be
passed as part of the \sphinxstyleemphasis{caller} argument when registering the callback
function, or the callback function must otherwise have access to the
LAMMPS object, such as through a module\sphinxhyphen{}based pointer.

\begin{sphinxadmonition}{note}{Note:}
\sphinxAtStartPar
This function is fully independent from
\sphinxcode{\sphinxupquote{fix\_external\_set\_energy\_global()}} and will \sphinxstylestrong{NOT} add any
contributions to the global energy tally and will \sphinxstylestrong{NOT} check whether
the sum of the contributions added here are consistent with the global
added energy.
\end{sphinxadmonition}

\sphinxAtStartPar
Please see the documentation for {\hyperref[\detokenize{fix_external::doc}]{\sphinxcrossref{\DUrole{doc}{fix external}}}} for
more information about how to use the fix and how to couple it with an
external code.
\begin{quote}\begin{description}
\sphinxlineitem{Parameters}\begin{itemize}
\item {} 
\sphinxAtStartPar
\sphinxstylestrong{id}\sphinxstyleemphasis{ {[}}\sphinxstyleemphasis{character(len=*)}\sphinxstyleemphasis{{]}} :: fix ID of the fix external instance

\item {} 
\sphinxAtStartPar
\sphinxstylestrong{eng}\sphinxstyleemphasis{ {[}}\sphinxstyleemphasis{real(c\_double)}\sphinxstyleemphasis{,}\sphinxstyleemphasis{dimension(:)}\sphinxstyleemphasis{{]}} :: array of length \sphinxstyleemphasis{nlocal} containing
the energy to add to the per\sphinxhyphen{}atom energy

\end{itemize}

\sphinxlineitem{Call to}
\sphinxAtStartPar
{\hyperref[\detokenize{Library_utility:_CPPv438lammps_fix_external_set_energy_peratomPvPKcPd}]{\sphinxcrossref{\sphinxcode{\sphinxupquote{lammps\_fix\_external\_set\_energy\_peratom()}}}}}

\end{description}\end{quote}

\end{fulllineitems}



\bigskip\hrule\bigskip

\index{set\_fix\_external\_set\_virial\_peratom() (fortran subroutine)@\spxentry{set\_fix\_external\_set\_virial\_peratom()}\spxextra{fortran subroutine}|spxpagem}

\begin{fulllineitems}
\phantomsection\label{\detokenize{Fortran:f/_/set_fix_external_set_virial_peratom}}
\pysigstartsignatures
\pysiglinewithargsret{\sphinxbfcode{\sphinxupquote{subroutine  }}\sphinxbfcode{\sphinxupquote{set\_fix\_external\_set\_virial\_peratom}}}{\sphinxparam{id}\sphinxparamcomma \sphinxparam{virial}}{}
\pysigstopsignatures
\sphinxAtStartPar
This is a companion function to \sphinxcode{\sphinxupquote{set\_fix\_external\_callback()}} to set
the per\sphinxhyphen{}atom virial contribution due to the fix from the external program as
part of the callback function.  For this to work, the LAMMPS object must be
passed as the \sphinxstyleemphasis{caller} argument when registering the callback function.

\sphinxAtStartPar
\DUrole{versionmodified,added}{Added in version 22Dec2022.}

\begin{sphinxadmonition}{note}{Note:}
\sphinxAtStartPar
This function is fully independent from
\sphinxcode{\sphinxupquote{fix\_external\_set\_virial\_global()}} and will \sphinxstylestrong{NOT} add any
contributions to the global virial tally and \sphinxstylestrong{NOT} check whether the
sum of the contributions added here are consistent with the global added
virial.
\end{sphinxadmonition}

\sphinxAtStartPar
The order and units of the per\sphinxhyphen{}atom stress tensor elements are the same
as for the global virial.  The type and dimensions of the per\sphinxhyphen{}atom virial
array must be \sphinxcode{\sphinxupquote{REAL(c\_double), DIMENSION(6,nlocal)}}.

\sphinxAtStartPar
Please see the documentation for {\hyperref[\detokenize{fix_external::doc}]{\sphinxcrossref{\DUrole{doc}{fix external}}}} for
more information about how to use the fix and how to couple it with an
external program.
\begin{quote}\begin{description}
\sphinxlineitem{Parameters}\begin{itemize}
\item {} 
\sphinxAtStartPar
\sphinxstylestrong{id}\sphinxstyleemphasis{ {[}}\sphinxstyleemphasis{character(len=*)}\sphinxstyleemphasis{{]}} :: fix ID of fix external instance

\item {} 
\sphinxAtStartPar
\sphinxstylestrong{virial}\sphinxstyleemphasis{ {[}}\sphinxstyleemphasis{real(c\_double)}\sphinxstyleemphasis{,}\sphinxstyleemphasis{dimension(:,:)}\sphinxstyleemphasis{{]}} :: an array of \(6 \times{}\)\sphinxstyleemphasis{nlocal} components to be added to the per\sphinxhyphen{}atom virial

\end{itemize}

\sphinxlineitem{Call to}
\sphinxAtStartPar
\sphinxcode{\sphinxupquote{lammps\_set\_virial\_peratom()}}

\end{description}\end{quote}

\end{fulllineitems}



\bigskip\hrule\bigskip

\index{fix\_external\_set\_vector\_length() (fortran subroutine)@\spxentry{fix\_external\_set\_vector\_length()}\spxextra{fortran subroutine}|spxpagem}

\begin{fulllineitems}
\phantomsection\label{\detokenize{Fortran:f/_/fix_external_set_vector_length}}
\pysigstartsignatures
\pysiglinewithargsret{\sphinxbfcode{\sphinxupquote{subroutine  }}\sphinxbfcode{\sphinxupquote{fix\_external\_set\_vector\_length}}}{\sphinxparam{id}\sphinxparamcomma \sphinxparam{length}}{}
\pysigstopsignatures
\sphinxAtStartPar
Set the vector length for a global vector stored with fix external for
analysis.

\sphinxAtStartPar
\DUrole{versionmodified,added}{Added in version 22Dec2022.}

\sphinxAtStartPar
This is a companion function to \sphinxcode{\sphinxupquote{set\_fix\_external\_callback()}} and
\sphinxcode{\sphinxupquote{fix\_external\_get\_force()}} to set the length of a global vector of
properties that will be stored with the fix via
\sphinxcode{\sphinxupquote{fix\_external\_set\_vector()}}.

\sphinxAtStartPar
This function needs to be called \sphinxstylestrong{before} a call to
\sphinxcode{\sphinxupquote{fix\_external\_set\_vector()}} and \sphinxstylestrong{before} a run or minimize command.
When running in parallel, it must be called from \sphinxstylestrong{all} MPI
processes with the same length argument.

\sphinxAtStartPar
Please see the documentation for {\hyperref[\detokenize{fix_external::doc}]{\sphinxcrossref{\DUrole{doc}{fix external}}}} for
more information about how to use the fix and how to couple it with an
external program.
\begin{quote}\begin{description}
\sphinxlineitem{Parameters}\begin{itemize}
\item {} 
\sphinxAtStartPar
\sphinxstylestrong{id}\sphinxstyleemphasis{ {[}}\sphinxstyleemphasis{character(len=*)}\sphinxstyleemphasis{{]}} :: fix ID of fix external instance

\item {} 
\sphinxAtStartPar
\sphinxstylestrong{length}\sphinxstyleemphasis{ {[}}\sphinxstyleemphasis{integer(c\_int)}\sphinxstyleemphasis{{]}} :: length of the global vector to be stored with the
fix

\end{itemize}

\sphinxlineitem{Call to}
\sphinxAtStartPar
{\hyperref[\detokenize{Library_utility:_CPPv437lammps_fix_external_set_vector_lengthPvPKci}]{\sphinxcrossref{\sphinxcode{\sphinxupquote{lammps\_fix\_external\_set\_vector\_length()}}}}}

\end{description}\end{quote}

\end{fulllineitems}



\bigskip\hrule\bigskip

\index{fix\_external\_set\_vector() (fortran subroutine)@\spxentry{fix\_external\_set\_vector()}\spxextra{fortran subroutine}|spxpagem}

\begin{fulllineitems}
\phantomsection\label{\detokenize{Fortran:f/_/fix_external_set_vector}}
\pysigstartsignatures
\pysiglinewithargsret{\sphinxbfcode{\sphinxupquote{subroutine  }}\sphinxbfcode{\sphinxupquote{fix\_external\_set\_vector}}}{\sphinxparam{id}\sphinxparamcomma \sphinxparam{idx}\sphinxparamcomma \sphinxparam{val}}{}
\pysigstopsignatures
\sphinxAtStartPar
Store a global vector value for a fix external instance with the given ID.

\sphinxAtStartPar
\DUrole{versionmodified,added}{Added in version 22Dec2022.}

\sphinxAtStartPar
This is a companion function to \sphinxcode{\sphinxupquote{set\_fix\_external\_callback()}} and
\sphinxcode{\sphinxupquote{fix\_external\_get\_force()}} to set the values of a global vector of
properties that will be stored with the fix and can be accessed from
within LAMMPS input commands (e.g., fix ave/time or variables) when used
in a vector context.

\sphinxAtStartPar
This function needs to be called \sphinxstylestrong{after} a call to
\sphinxcode{\sphinxupquote{fix\_external\_set\_vector\_length()}} and \sphinxstylestrong{before} a run or minimize
command.  When running in parallel, it must be called from \sphinxstylestrong{all} MPI
processes with the \sphinxstylestrong{same} \sphinxstyleemphasis{idx} and \sphinxstyleemphasis{val} parameters.  The variable
\sphinxstyleemphasis{val} is assumed to be extensive.

\begin{sphinxadmonition}{note}{Note:}
\sphinxAtStartPar
The index in the \sphinxstyleemphasis{idx} parameter is 1\sphinxhyphen{}based (i.e., the first element
is set with \sphinxstyleemphasis{idx}\({} = 1\), and the last element of the vector
with \sphinxstyleemphasis{idx}\({} = N\), where \(N\) is the value of the \sphinxstyleemphasis{length}
parameter of the call to \sphinxcode{\sphinxupquote{fix\_external\_set\_vector\_length()}}).
\end{sphinxadmonition}

\sphinxAtStartPar
Please see the documentation for {\hyperref[\detokenize{fix_external::doc}]{\sphinxcrossref{\DUrole{doc}{fix external}}}} for
more information about how to use the fix and how to couple it with an
external code.
\begin{quote}\begin{description}
\sphinxlineitem{Parameters}\begin{itemize}
\item {} 
\sphinxAtStartPar
\sphinxstylestrong{id}\sphinxstyleemphasis{ {[}}\sphinxstyleemphasis{character(len=*)}\sphinxstyleemphasis{{]}} :: ID of fix external instance

\item {} 
\sphinxAtStartPar
\sphinxstylestrong{idx}\sphinxstyleemphasis{ {[}}\sphinxstyleemphasis{integer(c\_int)}\sphinxstyleemphasis{{]}} :: 1\sphinxhyphen{}based index in global vector

\item {} 
\sphinxAtStartPar
\sphinxstylestrong{val}\sphinxstyleemphasis{ {[}}\sphinxstyleemphasis{integer(c\_int)}\sphinxstyleemphasis{{]}} :: value to be stored in global vector at index \sphinxstyleemphasis{idx}

\end{itemize}

\sphinxlineitem{Call to}
\sphinxAtStartPar
{\hyperref[\detokenize{Library_utility:_CPPv430lammps_fix_external_set_vectorPvPKcid}]{\sphinxcrossref{\sphinxcode{\sphinxupquote{lammps\_fix\_external\_set\_vector()}}}}}

\end{description}\end{quote}

\end{fulllineitems}



\bigskip\hrule\bigskip

\index{flush\_buffers() (fortran subroutine)@\spxentry{flush\_buffers()}\spxextra{fortran subroutine}|spxpagem}

\begin{fulllineitems}
\phantomsection\label{\detokenize{Fortran:f/_/flush_buffers}}
\pysigstartsignatures
\pysiglinewithargsret{\sphinxbfcode{\sphinxupquote{subroutine  }}\sphinxbfcode{\sphinxupquote{flush\_buffers}}}{}{}
\pysigstopsignatures
\sphinxAtStartPar
This function calls {\hyperref[\detokenize{Library_utility:_CPPv420lammps_flush_buffersPv}]{\sphinxcrossref{\sphinxcode{\sphinxupquote{lammps\_flush\_buffers()}}}}}, which flushes buffered
output to be written to screen and logfile. This can simplify capturing
output from LAMMPS library calls.

\sphinxAtStartPar
\DUrole{versionmodified,added}{Added in version 3Nov2022.}
\begin{quote}\begin{description}
\sphinxlineitem{Call to}
\sphinxAtStartPar
{\hyperref[\detokenize{Library_utility:_CPPv420lammps_flush_buffersPv}]{\sphinxcrossref{\sphinxcode{\sphinxupquote{lammps\_flush\_buffers()}}}}}

\end{description}\end{quote}

\end{fulllineitems}



\bigskip\hrule\bigskip

\index{is\_running() (fortran function)@\spxentry{is\_running()}\spxextra{fortran function}|spxpagem}

\begin{fulllineitems}
\phantomsection\label{\detokenize{Fortran:f/_/is_running}}
\pysigstartsignatures
\pysiglinewithargsret{\sphinxbfcode{\sphinxupquote{function  }}\sphinxbfcode{\sphinxupquote{is\_running}}}{}{}
\pysigstopsignatures
\sphinxAtStartPar
Check if LAMMPS is currently inside a run or minimization.

\sphinxAtStartPar
\DUrole{versionmodified,added}{Added in version 3Nov2022.}

\sphinxAtStartPar
This function can be used from signal handlers or multi\sphinxhyphen{}threaded
applications to determine if the LAMMPS instance is currently active.
\begin{quote}\begin{description}
\sphinxlineitem{Call to}
\sphinxAtStartPar
{\hyperref[\detokenize{Library_utility:_CPPv417lammps_is_runningPv}]{\sphinxcrossref{\sphinxcode{\sphinxupquote{lammps\_is\_running()}}}}}

\sphinxlineitem{Return}
\sphinxAtStartPar
\sphinxstylestrong{is\_running}\sphinxstyleemphasis{ {[}}\sphinxstyleemphasis{logical}\sphinxstyleemphasis{{]}} :: \sphinxcode{\sphinxupquote{.FALSE.}} if idle or \sphinxcode{\sphinxupquote{.TRUE.}} if active

\end{description}\end{quote}

\end{fulllineitems}



\bigskip\hrule\bigskip

\index{force\_timeout() (fortran subroutine)@\spxentry{force\_timeout()}\spxextra{fortran subroutine}|spxpagem}

\begin{fulllineitems}
\phantomsection\label{\detokenize{Fortran:f/_/force_timeout}}
\pysigstartsignatures
\pysiglinewithargsret{\sphinxbfcode{\sphinxupquote{subroutine  }}\sphinxbfcode{\sphinxupquote{force\_timeout}}}{}{}
\pysigstopsignatures
\sphinxAtStartPar
Force a timeout to stop an ongoing run cleanly.

\sphinxAtStartPar
\DUrole{versionmodified,added}{Added in version 3Nov2022.}

\sphinxAtStartPar
This function can be used from signal handlers or multi\sphinxhyphen{}threaded
applications to terminate an ongoing run cleanly.
\begin{quote}\begin{description}
\sphinxlineitem{Call to}
\sphinxAtStartPar
{\hyperref[\detokenize{Library_utility:_CPPv420lammps_force_timeoutPv}]{\sphinxcrossref{\sphinxcode{\sphinxupquote{lammps\_force\_timeout()}}}}}

\end{description}\end{quote}

\end{fulllineitems}



\bigskip\hrule\bigskip

\index{has\_error() (fortran function)@\spxentry{has\_error()}\spxextra{fortran function}|spxpagem}

\begin{fulllineitems}
\phantomsection\label{\detokenize{Fortran:f/_/has_error}}
\pysigstartsignatures
\pysiglinewithargsret{\sphinxbfcode{\sphinxupquote{function  }}\sphinxbfcode{\sphinxupquote{has\_error}}}{}{}
\pysigstopsignatures
\sphinxAtStartPar
Check if there is a (new) error message available.

\sphinxAtStartPar
\DUrole{versionmodified,added}{Added in version 3Nov2022.}

\sphinxAtStartPar
This function can be used to query if an error inside of LAMMPS
has thrown a {\hyperref[\detokenize{Build_settings:exceptions}]{\sphinxcrossref{\DUrole{std,std-ref}{C++ exception}}}}.
\begin{quote}\begin{description}
\sphinxlineitem{Call to}
\sphinxAtStartPar
{\hyperref[\detokenize{Library_utility:_CPPv416lammps_has_errorPv}]{\sphinxcrossref{\sphinxcode{\sphinxupquote{lammps\_has\_error()}}}}}

\sphinxlineitem{Return}
\sphinxAtStartPar
\sphinxstylestrong{has\_error}\sphinxstyleemphasis{ {[}}\sphinxstyleemphasis{logical}\sphinxstyleemphasis{{]}} :: \sphinxcode{\sphinxupquote{.TRUE.}} if there is an error.

\end{description}\end{quote}

\end{fulllineitems}



\bigskip\hrule\bigskip

\index{get\_last\_error\_message() (fortran subroutine)@\spxentry{get\_last\_error\_message()}\spxextra{fortran subroutine}|spxpagem}

\begin{fulllineitems}
\phantomsection\label{\detokenize{Fortran:f/_/get_last_error_message}}
\pysigstartsignatures
\pysiglinewithargsret{\sphinxbfcode{\sphinxupquote{subroutine  }}\sphinxbfcode{\sphinxupquote{get\_last\_error\_message}}}{\sphinxparam{buffer}\sphinxoptional{\sphinxparamcomma \sphinxparam{status}}}{}
\pysigstopsignatures
\sphinxAtStartPar
Copy the last error message into the provided buffer.

\sphinxAtStartPar
\DUrole{versionmodified,added}{Added in version 3Nov2022.}

\sphinxAtStartPar
This function can be used to retrieve the error message that was set
in the event of an error inside of LAMMPS that resulted in a
{\hyperref[\detokenize{Build_settings:exceptions}]{\sphinxcrossref{\DUrole{std,std-ref}{C++ exception}}}}.  A suitable buffer for a string has
to be provided.  If the internally\sphinxhyphen{}stored error message is longer than the
string, it will be truncated accordingly.  The optional argument \sphinxstyleemphasis{status}
indicates the kind of error: a “1” indicates an error that occurred on all
MPI ranks and is often recoverable, while a “2” indicates an abort that
would happen only in a single MPI rank and thus may not be recoverable, as
other MPI ranks may be waiting on the failing MPI rank(s) to send messages.
\begin{quote}\begin{description}
\sphinxlineitem{Parameters}
\sphinxAtStartPar
\sphinxstylestrong{buffer}\sphinxstyleemphasis{ {[}}\sphinxstyleemphasis{character(len=*)}\sphinxstyleemphasis{{]}} :: string buffer to copy the error message into

\sphinxlineitem{Options}
\sphinxAtStartPar
\sphinxstylestrong{status}\sphinxstyleemphasis{ {[}}\sphinxstyleemphasis{integer(c\_int)}\sphinxstyleemphasis{,}\sphinxstyleemphasis{optional}\sphinxstyleemphasis{{]}} :: 1 when all ranks had the error,
2 on a single\sphinxhyphen{}rank error.

\sphinxlineitem{Call to}
\sphinxAtStartPar
{\hyperref[\detokenize{Library_utility:_CPPv429lammps_get_last_error_messagePvPci}]{\sphinxcrossref{\sphinxcode{\sphinxupquote{lammps\_get\_last\_error\_message()}}}}}

\end{description}\end{quote}

\end{fulllineitems}



\bigskip\hrule\bigskip



\section{LAMMPS C++ API}
\label{\detokenize{Library:lammps-cplusplus-api}}\label{\detokenize{Library:id2}}
\sphinxAtStartPar
It is also possible to invoke the LAMMPS C++ API directly in your code.
It lacks some of the convenience of the C library API, but it allows
more direct access to simulation data and thus more low\sphinxhyphen{}level manipulations.
The following links provide some examples and references to the C++ API.

\sphinxstepscope


\subsection{Using the C++ API directly}
\label{\detokenize{Cplusplus:using-the-c-api-directly}}\label{\detokenize{Cplusplus::doc}}
\sphinxAtStartPar
Using the C++ classes of the LAMMPS library is lacking some of the
convenience of the C library API, but it allows a more direct access to
simulation data and thus more low\sphinxhyphen{}level manipulations and tighter
integration of LAMMPS into another code.  While for the complete C
library API is provided in the \sphinxcode{\sphinxupquote{library.h}} header file, for using
the C++ API it is required to include the individual header files
defining the individual classes in use.  Typically the name of the
class and the name of the header follow some simple rule.  Examples
are given below.


\subsection{Creating or deleting a LAMMPS object}
\label{\detokenize{Cplusplus:creating-or-deleting-a-lammps-object}}
\sphinxAtStartPar
When using the LAMMPS library interfaces, the core task is to create an
instance of the {\hyperref[\detokenize{Classes_lammps:_CPPv4N9LAMMPS_NS6LAMMPSE}]{\sphinxcrossref{\sphinxcode{\sphinxupquote{LAMMPS\_NS::LAMMPS}}}}} class.  In C++ this can
be done directly through the \sphinxcode{\sphinxupquote{new}} operator.  All further operations
are then initiated through calling member functions of some of the
components of the LAMMPS class or accessing their data members.  The
destruction of the LAMMPS instance is correspondingly initiated by using
the \sphinxcode{\sphinxupquote{delete}} operator.  Here is a simple example:

\begin{sphinxVerbatim}[commandchars=\\\{\}]
\PYG{c+cp}{\PYGZsh{}}\PYG{c+cp}{include}\PYG{+w}{ }\PYG{c+cpf}{\PYGZdq{}lammps.h\PYGZdq{}}

\PYG{c+cp}{\PYGZsh{}}\PYG{c+cp}{include}\PYG{+w}{ }\PYG{c+cpf}{\PYGZlt{}mpi.h\PYGZgt{}}
\PYG{c+cp}{\PYGZsh{}}\PYG{c+cp}{include}\PYG{+w}{ }\PYG{c+cpf}{\PYGZlt{}iostream\PYGZgt{}}

\PYG{k+kt}{int}\PYG{+w}{ }\PYG{n+nf}{main}\PYG{p}{(}\PYG{k+kt}{int}\PYG{+w}{ }\PYG{n}{argc}\PYG{p}{,}\PYG{+w}{ }\PYG{k+kt}{char}\PYG{+w}{ }\PYG{o}{*}\PYG{o}{*}\PYG{n}{argv}\PYG{p}{)}
\PYG{p}{\PYGZob{}}
\PYG{+w}{    }\PYG{n}{LAMMPS\PYGZus{}NS}\PYG{o}{:}\PYG{o}{:}\PYG{n}{LAMMPS}\PYG{+w}{ }\PYG{o}{*}\PYG{n}{lmp}\PYG{p}{;}
\PYG{+w}{    }\PYG{c+c1}{// custom argument vector for LAMMPS library}
\PYG{+w}{    }\PYG{k}{const}\PYG{+w}{ }\PYG{k+kt}{char}\PYG{+w}{ }\PYG{o}{*}\PYG{n}{lmpargv}\PYG{p}{[}\PYG{p}{]}\PYG{+w}{ }\PYG{p}{\PYGZob{}}\PYG{l+s}{\PYGZdq{}}\PYG{l+s}{liblammps}\PYG{l+s}{\PYGZdq{}}\PYG{p}{,}\PYG{+w}{ }\PYG{l+s}{\PYGZdq{}}\PYG{l+s}{\PYGZhy{}log}\PYG{l+s}{\PYGZdq{}}\PYG{p}{,}\PYG{+w}{ }\PYG{l+s}{\PYGZdq{}}\PYG{l+s}{none}\PYG{l+s}{\PYGZdq{}}\PYG{p}{\PYGZcb{}}\PYG{p}{;}
\PYG{+w}{    }\PYG{k+kt}{int}\PYG{+w}{ }\PYG{n}{lmpargc}\PYG{+w}{ }\PYG{o}{=}\PYG{+w}{ }\PYG{k}{sizeof}\PYG{p}{(}\PYG{n}{lmpargv}\PYG{p}{)}\PYG{o}{/}\PYG{k}{sizeof}\PYG{p}{(}\PYG{k}{const}\PYG{+w}{ }\PYG{k+kt}{char}\PYG{+w}{ }\PYG{o}{*}\PYG{p}{)}\PYG{p}{;}

\PYG{+w}{    }\PYG{c+c1}{// explicitly initialize MPI}
\PYG{+w}{    }\PYG{n}{MPI\PYGZus{}Init}\PYG{p}{(}\PYG{o}{\PYGZam{}}\PYG{n}{argc}\PYG{p}{,}\PYG{+w}{ }\PYG{o}{\PYGZam{}}\PYG{n}{argv}\PYG{p}{)}\PYG{p}{;}

\PYG{+w}{    }\PYG{c+c1}{// create LAMMPS instance}
\PYG{+w}{    }\PYG{n}{lmp}\PYG{+w}{ }\PYG{o}{=}\PYG{+w}{ }\PYG{k}{new}\PYG{+w}{ }\PYG{n}{LAMMPS\PYGZus{}NS}\PYG{o}{:}\PYG{o}{:}\PYG{n}{LAMMPS}\PYG{p}{(}\PYG{n}{lmpargc}\PYG{p}{,}\PYG{+w}{ }\PYG{p}{(}\PYG{k+kt}{char}\PYG{+w}{ }\PYG{o}{*}\PYG{o}{*}\PYG{p}{)}\PYG{n}{lmpargv}\PYG{p}{,}\PYG{+w}{ }\PYG{n}{MPI\PYGZus{}COMM\PYGZus{}WORLD}\PYG{p}{)}\PYG{p}{;}
\PYG{+w}{    }\PYG{c+c1}{// output numerical version string}
\PYG{+w}{    }\PYG{n}{std}\PYG{o}{:}\PYG{o}{:}\PYG{n}{cout}\PYG{+w}{ }\PYG{o}{\PYGZlt{}}\PYG{o}{\PYGZlt{}}\PYG{+w}{ }\PYG{l+s}{\PYGZdq{}}\PYG{l+s}{LAMMPS version ID: }\PYG{l+s}{\PYGZdq{}}\PYG{+w}{ }\PYG{o}{\PYGZlt{}}\PYG{o}{\PYGZlt{}}\PYG{+w}{ }\PYG{n}{lmp}\PYG{o}{\PYGZhy{}}\PYG{o}{\PYGZgt{}}\PYG{n}{num\PYGZus{}ver}\PYG{+w}{ }\PYG{o}{\PYGZlt{}}\PYG{o}{\PYGZlt{}}\PYG{+w}{ }\PYG{n}{std}\PYG{o}{:}\PYG{o}{:}\PYG{n}{endl}\PYG{p}{;}
\PYG{+w}{    }\PYG{c+c1}{// delete LAMMPS instance}
\PYG{+w}{    }\PYG{k}{delete}\PYG{+w}{ }\PYG{n}{lmp}\PYG{p}{;}

\PYG{+w}{    }\PYG{c+c1}{// stop MPI environment}
\PYG{+w}{    }\PYG{n}{MPI\PYGZus{}Finalize}\PYG{p}{(}\PYG{p}{)}\PYG{p}{;}
\PYG{+w}{    }\PYG{k}{return}\PYG{+w}{ }\PYG{l+m+mi}{0}\PYG{p}{;}
\PYG{p}{\PYGZcb{}}
\end{sphinxVerbatim}

\sphinxAtStartPar
This minimal example only requires to include the \sphinxcode{\sphinxupquote{lammps.h}} header
file since it only accesses a non\sphinxhyphen{}pointer member of the LAMMPS class.


\subsection{Executing LAMMPS commands}
\label{\detokenize{Cplusplus:executing-lammps-commands}}
\sphinxAtStartPar
Once a LAMMPS instance is created by your C++ code, you need to set up a
simulation and that is most conveniently done by “driving” it through
issuing commands like you would do when running a LAMMPS simulation from
an input script. Processing of input in LAMMPS is handled by the
{\hyperref[\detokenize{Classes_input:_CPPv4N9LAMMPS_NS5InputE}]{\sphinxcrossref{\sphinxcode{\sphinxupquote{Input}}}}} class an instance of which is a
member of the {\hyperref[\detokenize{Classes_lammps:_CPPv4N9LAMMPS_NS6LAMMPSE}]{\sphinxcrossref{\sphinxcode{\sphinxupquote{LAMMPS}}}}} class.  You have
two options: reading commands from a file, or executing a single
command from a string. See below for a small example:

\begin{sphinxVerbatim}[commandchars=\\\{\}]
\PYG{c+cp}{\PYGZsh{}}\PYG{c+cp}{include}\PYG{+w}{ }\PYG{c+cpf}{\PYGZdq{}lammps.h\PYGZdq{}}
\PYG{c+cp}{\PYGZsh{}}\PYG{c+cp}{include}\PYG{+w}{ }\PYG{c+cpf}{\PYGZdq{}input.h\PYGZdq{}}
\PYG{c+cp}{\PYGZsh{}}\PYG{c+cp}{include}\PYG{+w}{ }\PYG{c+cpf}{\PYGZlt{}mpi.h\PYGZgt{}}

\PYG{k}{using}\PYG{+w}{ }\PYG{k}{namespace}\PYG{+w}{ }\PYG{n+nn}{LAMMPS\PYGZus{}NS}\PYG{p}{;}

\PYG{k+kt}{int}\PYG{+w}{ }\PYG{n+nf}{main}\PYG{p}{(}\PYG{k+kt}{int}\PYG{+w}{ }\PYG{n}{argc}\PYG{p}{,}\PYG{+w}{ }\PYG{k+kt}{char}\PYG{+w}{ }\PYG{o}{*}\PYG{o}{*}\PYG{n}{argv}\PYG{p}{)}
\PYG{p}{\PYGZob{}}
\PYG{+w}{    }\PYG{k}{const}\PYG{+w}{ }\PYG{k+kt}{char}\PYG{+w}{ }\PYG{o}{*}\PYG{n}{lmpargv}\PYG{p}{[}\PYG{p}{]}\PYG{+w}{ }\PYG{p}{\PYGZob{}}\PYG{l+s}{\PYGZdq{}}\PYG{l+s}{liblammps}\PYG{l+s}{\PYGZdq{}}\PYG{p}{,}\PYG{+w}{ }\PYG{l+s}{\PYGZdq{}}\PYG{l+s}{\PYGZhy{}log}\PYG{l+s}{\PYGZdq{}}\PYG{p}{,}\PYG{+w}{ }\PYG{l+s}{\PYGZdq{}}\PYG{l+s}{none}\PYG{l+s}{\PYGZdq{}}\PYG{p}{\PYGZcb{}}\PYG{p}{;}
\PYG{+w}{    }\PYG{k+kt}{int}\PYG{+w}{ }\PYG{n}{lmpargc}\PYG{+w}{ }\PYG{o}{=}\PYG{+w}{ }\PYG{k}{sizeof}\PYG{p}{(}\PYG{n}{lmpargv}\PYG{p}{)}\PYG{o}{/}\PYG{k}{sizeof}\PYG{p}{(}\PYG{k}{const}\PYG{+w}{ }\PYG{k+kt}{char}\PYG{+w}{ }\PYG{o}{*}\PYG{p}{)}\PYG{p}{;}

\PYG{+w}{    }\PYG{n}{MPI\PYGZus{}Init}\PYG{p}{(}\PYG{o}{\PYGZam{}}\PYG{n}{argc}\PYG{p}{,}\PYG{+w}{ }\PYG{o}{\PYGZam{}}\PYG{n}{argv}\PYG{p}{)}\PYG{p}{;}
\PYG{+w}{    }\PYG{n}{LAMMPS}\PYG{+w}{ }\PYG{o}{*}\PYG{n}{lmp}\PYG{+w}{ }\PYG{o}{=}\PYG{+w}{ }\PYG{k}{new}\PYG{+w}{ }\PYG{n}{LAMMPS}\PYG{p}{(}\PYG{n}{lmpargc}\PYG{p}{,}\PYG{+w}{ }\PYG{p}{(}\PYG{k+kt}{char}\PYG{+w}{ }\PYG{o}{*}\PYG{o}{*}\PYG{p}{)}\PYG{n}{lmpargv}\PYG{p}{,}\PYG{+w}{ }\PYG{n}{MPI\PYGZus{}COMM\PYGZus{}WORLD}\PYG{p}{)}\PYG{p}{;}
\PYG{+w}{    }\PYG{n}{lmp}\PYG{o}{\PYGZhy{}}\PYG{o}{\PYGZgt{}}\PYG{n}{input}\PYG{o}{\PYGZhy{}}\PYG{o}{\PYGZgt{}}\PYG{n}{file}\PYG{p}{(}\PYG{l+s}{\PYGZdq{}}\PYG{l+s}{in.melt}\PYG{l+s}{\PYGZdq{}}\PYG{p}{)}\PYG{p}{;}
\PYG{+w}{    }\PYG{n}{lmp}\PYG{o}{\PYGZhy{}}\PYG{o}{\PYGZgt{}}\PYG{n}{input}\PYG{o}{\PYGZhy{}}\PYG{o}{\PYGZgt{}}\PYG{n}{one}\PYG{p}{(}\PYG{l+s}{\PYGZdq{}}\PYG{l+s}{run 100 post no}\PYG{l+s}{\PYGZdq{}}\PYG{p}{)}\PYG{p}{;}
\PYG{+w}{    }\PYG{k}{delete}\PYG{+w}{ }\PYG{n}{lmp}\PYG{p}{;}
\PYG{+w}{    }\PYG{k}{return}\PYG{+w}{ }\PYG{l+m+mi}{0}\PYG{p}{;}
\PYG{p}{\PYGZcb{}}
\end{sphinxVerbatim}

\sphinxstepscope


\chapter{Use Python with LAMMPS}
\label{\detokenize{Python_head:use-python-with-lammps}}\label{\detokenize{Python_head::doc}}
\sphinxAtStartPar
These pages describe various ways that LAMMPS and Python can be used
together.

\sphinxstepscope


\section{Overview}
\label{\detokenize{Python_overview:overview}}\label{\detokenize{Python_overview::doc}}
\sphinxAtStartPar
The LAMMPS distribution includes a \sphinxcode{\sphinxupquote{python}} directory with the Python
code needed to run LAMMPS from Python.  The \sphinxcode{\sphinxupquote{python/lammps}} package
contains {\hyperref[\detokenize{Python_module::doc}]{\sphinxcrossref{\DUrole{doc}{the “lammps” Python module}}}} that wraps the
LAMMPS C\sphinxhyphen{}library interface.  This module makes it is possible to do the
following either from a Python script, or interactively from a Python
prompt:
\begin{itemize}
\item {} 
\sphinxAtStartPar
create one or more instances of LAMMPS

\item {} 
\sphinxAtStartPar
invoke LAMMPS commands or read them from an input script

\item {} 
\sphinxAtStartPar
run LAMMPS incrementally

\item {} 
\sphinxAtStartPar
extract LAMMPS results

\item {} 
\sphinxAtStartPar
and modify internal LAMMPS data structures.

\end{itemize}

\sphinxAtStartPar
From a Python script you can do this in serial or in parallel.  Running
Python interactively in parallel does not generally work, unless you
have a version of Python that extends Python to enable multiple
instances of Python to read what you type.

\sphinxAtStartPar
To do all of this, you must build LAMMPS in {\hyperref[\detokenize{Build_basics:exe}]{\sphinxcrossref{\DUrole{std,std-ref}{“shared” mode}}}}
and make certain that your Python interpreter can find the \sphinxcode{\sphinxupquote{lammps}}
Python package and the LAMMPS shared library file.

\sphinxAtStartPar
The Python wrapper for LAMMPS uses the \sphinxhref{https://docs.python.org/3/library/ctypes.html}{ctypes} package in
Python, which auto\sphinxhyphen{}generates the interface code needed between Python
and a set of C\sphinxhyphen{}style library functions.  Ctypes has been part of the
standard Python distribution since version 2.5.  You can check which
version of Python you have by simply typing “python” at a shell prompt.
Below is an example output for Python version 3.8.5.

\begin{sphinxVerbatim}[commandchars=\\\{\}]
\PYG{g+gp}{\PYGZdl{} }python
\PYG{g+go}{Python 3.8.5 (default, Aug 12 2020, 00:00:00)}
\PYG{g+go}{[GCC 10.2.1 20200723 (Red Hat 10.2.1\PYGZhy{}1)] on linux}
\PYG{g+go}{Type \PYGZdq{}help\PYGZdq{}, \PYGZdq{}copyright\PYGZdq{}, \PYGZdq{}credits\PYGZdq{} or \PYGZdq{}license\PYGZdq{} for more information.}
\PYG{g+go}{\PYGZgt{}\PYGZgt{}\PYGZgt{}}
\end{sphinxVerbatim}

\begin{sphinxadmonition}{warning}{Warning:}
\sphinxAtStartPar
The options described in this section of the manual for using Python
with LAMMPS currently support either Python 2 or 3.  Specifically
version 2.7 or later and 3.6 or later.  Since the Python community no
longer maintains Python 2 (see \sphinxhref{https://www.python.org/doc/sunset-python-2/}{this notice}), we recommend use of
Python 3 with LAMMPS.  While Python 2 code should continue to work,
that is not something we can guarantee long\sphinxhyphen{}term.  If you notice
Python code in the LAMMPS distribution that is not compatible with
Python 3, please contact the LAMMPS developers or submit \sphinxhref{https://github.com/lammps/lammps/issues}{and issue
on GitHub}
\end{sphinxadmonition}


\bigskip\hrule\bigskip


\sphinxAtStartPar
LAMMPS can work together with Python in three ways.  First, Python can
wrap LAMMPS through the its {\hyperref[\detokenize{Library::doc}]{\sphinxcrossref{\DUrole{doc}{library interface}}}}, so
that a Python script can create one or more instances of LAMMPS and
launch one or more simulations.  In Python terms, this is referred to as
“extending” Python with a LAMMPS module.

\begin{figure}[htbp]
\centering
\capstart

\noindent\sphinxincludegraphics{{python-invoke-lammps}.png}
\caption{Launching LAMMPS via Python}\label{\detokenize{Python_overview:id1}}\end{figure}

\sphinxAtStartPar
Second, the lower\sphinxhyphen{}level Python interface in the {\hyperref[\detokenize{Python_module:lammps.lammps}]{\sphinxcrossref{\sphinxcode{\sphinxupquote{lammps Python
class}}}}} can be used indirectly through the provided
{\hyperref[\detokenize{Python_module:lammps.PyLammps}]{\sphinxcrossref{\sphinxcode{\sphinxupquote{PyLammps}}}}} and {\hyperref[\detokenize{Python_module:lammps.IPyLammps}]{\sphinxcrossref{\sphinxcode{\sphinxupquote{IPyLammps}}}}} wrapper classes, also written in Python.  These
wrappers try to simplify the usage of LAMMPS in Python by providing a
more object\sphinxhyphen{}based interface to common LAMMPS functionality.  They also
reduce the amount of code necessary to parameterize LAMMPS scripts
through Python and make variables and computes directly accessible.

\begin{figure}[htbp]
\centering
\capstart

\noindent\sphinxincludegraphics{{pylammps-invoke-lammps}.png}
\caption{Using the PyLammps / IPyLammps wrappers}\label{\detokenize{Python_overview:id2}}\end{figure}

\sphinxAtStartPar
Third, LAMMPS can use the Python interpreter, so that a LAMMPS input
script or styles can invoke Python code directly, and pass information
back\sphinxhyphen{}and\sphinxhyphen{}forth between the input script and Python functions you write.
This Python code can also call back to LAMMPS to query or change its
attributes through the LAMMPS Python module mentioned above.  In Python
terms, this is called “embedding” Python into LAMMPS.  When used in this
mode, Python can perform script operations that the simple LAMMPS input
script syntax can not.

\begin{figure}[htbp]
\centering
\capstart

\noindent\sphinxincludegraphics{{lammps-invoke-python}.png}
\caption{Calling Python code from LAMMPS}\label{\detokenize{Python_overview:id3}}\end{figure}

\sphinxstepscope


\section{Installation}
\label{\detokenize{Python_install:installation}}\label{\detokenize{Python_install::doc}}
\sphinxAtStartPar
The LAMMPS Python module enables calling the {\hyperref[\detokenize{Library:lammps-c-api}]{\sphinxcrossref{\DUrole{std,std-ref}{LAMMPS C library API}}}} from Python by dynamically loading functions in the
LAMMPS shared library through the Python \sphinxhref{https://docs.python.org/3/library/ctypes.html}{ctypes}
module.  Because of the dynamic loading, it is required that LAMMPS is
compiled in {\hyperref[\detokenize{Build_basics:exe}]{\sphinxcrossref{\DUrole{std,std-ref}{“shared” mode}}}}.

\sphinxAtStartPar
Two components are necessary for Python to be able to invoke LAMMPS code:
\begin{itemize}
\item {} 
\sphinxAtStartPar
The LAMMPS Python Package (\sphinxcode{\sphinxupquote{lammps}}) from the \sphinxcode{\sphinxupquote{python}} folder

\item {} 
\sphinxAtStartPar
The LAMMPS Shared Library (\sphinxcode{\sphinxupquote{liblammps.so}}, \sphinxcode{\sphinxupquote{liblammps.dylib}} or
\sphinxcode{\sphinxupquote{liblammps.dll}}) from the folder where you compiled LAMMPS.

\end{itemize}


\subsection{Installing the LAMMPS Python Module and Shared Library}
\label{\detokenize{Python_install:installing-the-lammps-python-module-and-shared-library}}\label{\detokenize{Python_install:python-install-guides}}
\sphinxAtStartPar
Making LAMMPS usable within Python and vice versa requires putting the
LAMMPS Python package (\sphinxcode{\sphinxupquote{lammps}}) into a location where the Python
interpreter can find it and installing the LAMMPS shared library into a
folder that the dynamic loader searches or inside of the installed
\sphinxcode{\sphinxupquote{lammps}} package folder.  There are multiple ways to achieve this.
\begin{enumerate}
\sphinxsetlistlabels{\arabic}{enumi}{enumii}{}{.}%
\item {} 
\sphinxAtStartPar
Install both components into a Python \sphinxcode{\sphinxupquote{site\sphinxhyphen{}packages}} folder, either
system\sphinxhyphen{}wide or in the corresponding user\sphinxhyphen{}specific folder. This way no
additional environment variables need to be set, but the shared
library is otherwise not accessible.

\item {} 
\sphinxAtStartPar
Do an installation into a virtual environment.

\item {} 
\sphinxAtStartPar
Leave the files where they are in the source/development tree and
adjust some environment variables.

\begin{sphinxadmonition}{note}{Python package}

\sphinxAtStartPar
Compile LAMMPS with either {\hyperref[\detokenize{Build_cmake::doc}]{\sphinxcrossref{\DUrole{doc}{CMake}}}} or the
{\hyperref[\detokenize{Build_make::doc}]{\sphinxcrossref{\DUrole{doc}{traditional make}}}} procedure in {\hyperref[\detokenize{Build_basics:exe}]{\sphinxcrossref{\DUrole{std,std-ref}{shared
mode}}}}.  After compilation has finished, type (in the
compilation folder):

\begin{sphinxVerbatim}[commandchars=\\\{\}]
make\PYG{+w}{ }install\PYGZhy{}python
\end{sphinxVerbatim}

\sphinxAtStartPar
This will try to build a so\sphinxhyphen{}called (binary) wheel file, a
compressed binary python package and then install it with the
python package manager ‘pip’.  Installation will be attempted into
a system\sphinxhyphen{}wide \sphinxcode{\sphinxupquote{site\sphinxhyphen{}packages}} folder and if that fails into the
corresponding folder in the user’s home directory.  For a
system\sphinxhyphen{}wide installation you usually would have to gain superuser
privilege first, e.g. though \sphinxcode{\sphinxupquote{sudo}}


\begin{savenotes}\sphinxattablestart
\sphinxthistablewithglobalstyle
\centering
\begin{tabular}[t]{*{3}{\X{1}{3}}}
\sphinxtoprule
\sphinxstyletheadfamily 
\sphinxAtStartPar
File
&\sphinxstyletheadfamily 
\sphinxAtStartPar
Location
&\sphinxstyletheadfamily 
\sphinxAtStartPar
Notes
\\
\sphinxmidrule
\sphinxtableatstartofbodyhook
\sphinxAtStartPar
LAMMPS Python package
&\begin{itemize}
\item {} 
\sphinxAtStartPar
\sphinxcode{\sphinxupquote{\$HOME/.local/lib/pythonX.Y/site\sphinxhyphen{}packages/lammps}}

\end{itemize}
&
\sphinxAtStartPar
\sphinxcode{\sphinxupquote{X.Y}} depends on the installed Python version
\\
\sphinxhline
\sphinxAtStartPar
LAMMPS shared library
&\begin{itemize}
\item {} 
\sphinxAtStartPar
\sphinxcode{\sphinxupquote{\$HOME/.local/lib/pythonX.Y/site\sphinxhyphen{}packages/lammps}}

\end{itemize}
&
\sphinxAtStartPar
\sphinxcode{\sphinxupquote{X.Y}} depends on the installed Python version
\\
\sphinxbottomrule
\end{tabular}
\sphinxtableafterendhook\par
\sphinxattableend\end{savenotes}

\sphinxAtStartPar
For a system\sphinxhyphen{}wide installation those folders would then become.


\begin{savenotes}\sphinxattablestart
\sphinxthistablewithglobalstyle
\centering
\begin{tabular}[t]{*{3}{\X{1}{3}}}
\sphinxtoprule
\sphinxstyletheadfamily 
\sphinxAtStartPar
File
&\sphinxstyletheadfamily 
\sphinxAtStartPar
Location
&\sphinxstyletheadfamily 
\sphinxAtStartPar
Notes
\\
\sphinxmidrule
\sphinxtableatstartofbodyhook
\sphinxAtStartPar
LAMMPS Python package
&\begin{itemize}
\item {} 
\sphinxAtStartPar
\sphinxcode{\sphinxupquote{/usr/lib/pythonX.Y/site\sphinxhyphen{}packages/lammps}}

\end{itemize}
&
\sphinxAtStartPar
\sphinxcode{\sphinxupquote{X.Y}} depends on the installed Python version
\\
\sphinxhline
\sphinxAtStartPar
LAMMPS shared library
&\begin{itemize}
\item {} 
\sphinxAtStartPar
\sphinxcode{\sphinxupquote{/usr/lib/pythonX.Y/site\sphinxhyphen{}packages/lammps}}

\end{itemize}
&
\sphinxAtStartPar
\sphinxcode{\sphinxupquote{X.Y}} depends on the installed Python version
\\
\sphinxbottomrule
\end{tabular}
\sphinxtableafterendhook\par
\sphinxattableend\end{savenotes}

\sphinxAtStartPar
No environment variables need to be set for those, as those
folders are searched by default by Python or the LAMMPS Python
package.

\sphinxAtStartPar
\DUrole{versionmodified,changed}{Changed in version 24Mar2022.}

\begin{sphinxadmonition}{note}{Note:}
\sphinxAtStartPar
If there is an existing installation of the LAMMPS python
module, \sphinxcode{\sphinxupquote{make install\sphinxhyphen{}python}} will try to update it.
However, that will fail if the older version of the module
was installed by LAMMPS versions until 17Feb2022.  Those
were using the distutils package, which does not create a
“manifest” that allows a clean uninstall.  The \sphinxcode{\sphinxupquote{make
install\sphinxhyphen{}python}} command will always produce a
lammps\sphinxhyphen{}\textless{}version\textgreater{}\sphinxhyphen{}\textless{}python\textgreater{}\sphinxhyphen{}\textless{}abi\textgreater{}\sphinxhyphen{}\textless{}os\textgreater{}\sphinxhyphen{}\textless{}arch\textgreater{}.whl file (the
‘wheel’). And this file can be later installed directly with
\sphinxcode{\sphinxupquote{python \sphinxhyphen{}m pip install \textless{}wheel file\textgreater{}.whl}} without having to
type \sphinxcode{\sphinxupquote{make install\sphinxhyphen{}python}} again and repeating the build
step, too.
\end{sphinxadmonition}

\sphinxAtStartPar
For the traditional make process you can override the python
version to version x.y when calling \sphinxcode{\sphinxupquote{make}} with
\sphinxcode{\sphinxupquote{PYTHON=pythonX.Y}}.  For a CMake based compilation this choice
has to be made during the CMake configuration step.

\sphinxAtStartPar
If the default settings of \sphinxcode{\sphinxupquote{make install\sphinxhyphen{}python}} are not what you want,
you can invoke \sphinxcode{\sphinxupquote{install.py}} from the python directory manually as

\begin{sphinxVerbatim}[commandchars=\\\{\}]
python\PYG{+w}{ }install.py\PYG{+w}{ }\PYGZhy{}p\PYG{+w}{ }\PYGZlt{}python\PYG{+w}{ }package\PYGZgt{}\PYG{+w}{ }\PYGZhy{}l\PYG{+w}{ }\PYGZlt{}shared\PYG{+w}{ }library\PYGZgt{}\PYG{+w}{ }\PYGZhy{}v\PYG{+w}{ }\PYGZlt{}version.h\PYG{+w}{ }file\PYGZgt{}\PYG{+w}{ }\PYG{o}{[}\PYGZhy{}n\PYG{o}{]}
\end{sphinxVerbatim}
\begin{itemize}
\item {} 
\sphinxAtStartPar
The \sphinxcode{\sphinxupquote{\sphinxhyphen{}p}} flag points to the \sphinxcode{\sphinxupquote{lammps}} Python package folder to be installed,

\item {} 
\sphinxAtStartPar
the \sphinxcode{\sphinxupquote{\sphinxhyphen{}l}} flag points to the LAMMPS shared library file to be installed,

\item {} 
\sphinxAtStartPar
the \sphinxcode{\sphinxupquote{\sphinxhyphen{}v}} flag points to the LAMMPS version header file to extract the version date,

\item {} 
\sphinxAtStartPar
and the optional \sphinxcode{\sphinxupquote{\sphinxhyphen{}n}} instructs the script to only build a wheel file
but not attempt to install it.

\end{itemize}
\end{sphinxadmonition}

\begin{sphinxadmonition}{note}{Virtual environment}

\sphinxAtStartPar
A virtual environment is a minimal Python installation inside of a
folder.  It allows isolating and customizing a Python environment
that is mostly independent from a user or system installation.
For the core Python environment, it uses symbolic links to the
system installation and thus it can be set up quickly and will not
take up much disk space.  This gives you the flexibility to
install (newer/different) versions of Python packages that would
potentially conflict with already installed system packages.  It
also does not requite any superuser privileges. See \sphinxhref{https://peps.python.org/pep-0405/}{PEP 405:
Python Virtual Environments} for more
information.

\sphinxAtStartPar
To create a virtual environment in the folder \sphinxcode{\sphinxupquote{\$HOME/myenv}},
use the \sphinxhref{https://docs.python.org/3/library/venv.html}{venv} module as follows.

\begin{sphinxVerbatim}[commandchars=\\\{\}]
\PYG{c+c1}{\PYGZsh{} create virtual environment in folder \PYGZdl{}HOME/myenv}
python3\PYG{+w}{ }\PYGZhy{}m\PYG{+w}{ }venv\PYG{+w}{ }\PYG{n+nv}{\PYGZdl{}HOME}/myenv
\end{sphinxVerbatim}

\sphinxAtStartPar
For Python versions prior 3.3 you can use \sphinxhref{https://packaging.python.org/en/latest/key\_projects/\#virtualenv}{virtualenv}
command instead of “python3 \sphinxhyphen{}m venv”.  This step has to be done
only once.

\sphinxAtStartPar
To activate the virtual environment type:

\begin{sphinxVerbatim}[commandchars=\\\{\}]
\PYG{n+nb}{source}\PYG{+w}{ }\PYG{n+nv}{\PYGZdl{}HOME}/myenv/bin/activate
\end{sphinxVerbatim}

\sphinxAtStartPar
This has to be done every time you log in or open a new terminal
window and after you turn off the virtual environment with the
\sphinxcode{\sphinxupquote{deactivate}} command.

\sphinxAtStartPar
When using CMake to build LAMMPS, you need to set
\sphinxcode{\sphinxupquote{CMAKE\_INSTALL\_PREFIX}} to the value of the \sphinxcode{\sphinxupquote{\$VIRTUAL\_ENV}}
environment variable during the configuration step. For the
traditional make procedure, no additional steps are needed.
After compiling LAMMPS you can do a “Python package only”
installation with \sphinxcode{\sphinxupquote{make install\sphinxhyphen{}python}} and the LAMMPS Python
package and the shared library file are installed into the
following locations:


\begin{savenotes}\sphinxattablestart
\sphinxthistablewithglobalstyle
\centering
\begin{tabular}[t]{*{3}{\X{1}{3}}}
\sphinxtoprule
\sphinxstyletheadfamily 
\sphinxAtStartPar
File
&\sphinxstyletheadfamily 
\sphinxAtStartPar
Location
&\sphinxstyletheadfamily 
\sphinxAtStartPar
Notes
\\
\sphinxmidrule
\sphinxtableatstartofbodyhook
\sphinxAtStartPar
LAMMPS Python Module
&\begin{itemize}
\item {} 
\sphinxAtStartPar
\sphinxcode{\sphinxupquote{\$VIRTUAL\_ENV/lib/pythonX.Y/site\sphinxhyphen{}packages/lammps}}

\end{itemize}
&
\sphinxAtStartPar
\sphinxcode{\sphinxupquote{X.Y}} depends on the installed Python version
\\
\sphinxhline
\sphinxAtStartPar
LAMMPS shared library
&\begin{itemize}
\item {} 
\sphinxAtStartPar
\sphinxcode{\sphinxupquote{\$VIRTUAL\_ENV/lib/pythonX.Y/site\sphinxhyphen{}packages/lammps}}

\end{itemize}
&
\sphinxAtStartPar
\sphinxcode{\sphinxupquote{X.Y}} depends on the installed Python version
\\
\sphinxbottomrule
\end{tabular}
\sphinxtableafterendhook\par
\sphinxattableend\end{savenotes}
\end{sphinxadmonition}

\begin{sphinxadmonition}{note}{In place usage}

\sphinxAtStartPar
You can also {\hyperref[\detokenize{Build::doc}]{\sphinxcrossref{\DUrole{doc}{compile LAMMPS}}}} as usual in
{\hyperref[\detokenize{Build_basics:exe}]{\sphinxcrossref{\DUrole{std,std-ref}{“shared” mode}}}} leave the shared library and Python
package inside the source/compilation folders. Instead of
copying the files where they can be found, you need to set the environment
variables \sphinxcode{\sphinxupquote{PYTHONPATH}} (for the Python package) and
\sphinxcode{\sphinxupquote{LD\_LIBRARY\_PATH}} (or \sphinxcode{\sphinxupquote{DYLD\_LIBRARY\_PATH}} on macOS

\sphinxAtStartPar
For Bourne shells (bash, ksh and similar) the commands are:

\begin{sphinxVerbatim}[commandchars=\\\{\}]
\PYG{n+nb}{export}\PYG{+w}{ }\PYG{n+nv}{PYTHONPATH}\PYG{o}{=}\PYG{l+s+si}{\PYGZdl{}\PYGZob{}}\PYG{n+nv}{PYTHONPATH}\PYG{l+s+si}{\PYGZcb{}}:\PYG{l+s+si}{\PYGZdl{}\PYGZob{}}\PYG{n+nv}{HOME}\PYG{l+s+si}{\PYGZcb{}}/lammps/python
\PYG{n+nb}{export}\PYG{+w}{ }\PYG{n+nv}{LD\PYGZus{}LIBRARY\PYGZus{}PATH}\PYG{o}{=}\PYG{l+s+si}{\PYGZdl{}\PYGZob{}}\PYG{n+nv}{LD\PYGZus{}LIBRARY\PYGZus{}PATH}\PYG{l+s+si}{\PYGZcb{}}:\PYG{l+s+si}{\PYGZdl{}\PYGZob{}}\PYG{n+nv}{HOME}\PYG{l+s+si}{\PYGZcb{}}/lammps/src
\end{sphinxVerbatim}

\sphinxAtStartPar
For the C\sphinxhyphen{}shells like csh or tcsh the commands are:

\begin{sphinxVerbatim}[commandchars=\\\{\}]
\PYG{n+nb}{setenv }PYTHONPATH \PYG{k}{\PYGZdl{}\PYGZob{}}\PYG{n+nv}{PYTHONPATH}\PYG{k}{\PYGZcb{}}:\PYG{k}{\PYGZdl{}\PYGZob{}}\PYG{n+nv}{HOME}\PYG{k}{\PYGZcb{}}/lammps/python
\PYG{n+nb}{setenv }LD\PYGZus{}LIBRARY\PYGZus{}PATH \PYG{k}{\PYGZdl{}\PYGZob{}}\PYG{n+nv}{LD\PYGZus{}LIBRARY\PYGZus{}PATH}\PYG{k}{\PYGZcb{}}:\PYG{k}{\PYGZdl{}\PYGZob{}}\PYG{n+nv}{HOME}\PYG{k}{\PYGZcb{}}/lammps/src
\end{sphinxVerbatim}

\sphinxAtStartPar
On macOS you may also need to set \sphinxcode{\sphinxupquote{DYLD\_LIBRARY\_PATH}} accordingly.
You can make those changes permanent by editing your \sphinxcode{\sphinxupquote{\$HOME/.bashrc}}
or \sphinxcode{\sphinxupquote{\$HOME/.login}} files, respectively.

\begin{sphinxadmonition}{note}{Note:}
\sphinxAtStartPar
The \sphinxcode{\sphinxupquote{PYTHONPATH}} needs to point to the parent folder that contains the \sphinxcode{\sphinxupquote{lammps}} package!
\end{sphinxadmonition}
\end{sphinxadmonition}

\end{enumerate}

\sphinxAtStartPar
To verify if LAMMPS can be successfully started from Python, start the
Python interpreter, load the \sphinxcode{\sphinxupquote{lammps}} Python module and create a
LAMMPS instance.  This should not generate an error message and produce
output similar to the following:
\begin{quote}

\begin{sphinxVerbatim}[commandchars=\\\{\}]
\PYG{g+gp}{\PYGZdl{} }python
\PYG{g+go}{Python 3.8.5 (default, Sep  5 2020, 10:50:12)}
\PYG{g+go}{[GCC 10.2.0] on linux}
\PYG{g+go}{Type \PYGZdq{}help\PYGZdq{}, \PYGZdq{}copyright\PYGZdq{}, \PYGZdq{}credits\PYGZdq{} or \PYGZdq{}license\PYGZdq{} for more information.}
\PYG{g+go}{\PYGZgt{}\PYGZgt{}\PYGZgt{} import lammps}
\PYG{g+go}{\PYGZgt{}\PYGZgt{}\PYGZgt{} lmp = lammps.lammps()}
\PYG{g+go}{LAMMPS (18 Sep 2020)}
\PYG{g+go}{using 1 OpenMP thread(s) per MPI task}
\PYG{g+go}{\PYGZgt{}\PYGZgt{}\PYGZgt{}}
\end{sphinxVerbatim}
\end{quote}

\begin{sphinxadmonition}{note}{Note:}
\sphinxAtStartPar
Unless you opted for “In place use”, you will have to rerun the installation
any time you recompile LAMMPS to ensure the latest Python package and shared
library are installed and used.
\end{sphinxadmonition}

\begin{sphinxadmonition}{note}{Note:}
\sphinxAtStartPar
If you want Python to be able to load different versions of the
LAMMPS shared library with different settings, you will need to
manually copy the files under different names
(e.g. \sphinxcode{\sphinxupquote{liblammps\_mpi.so}} or \sphinxcode{\sphinxupquote{liblammps\_gpu.so}}) into the
appropriate folder as indicated above. You can then select the
desired library through the \sphinxstyleemphasis{name} argument of the LAMMPS object
constructor (see {\hyperref[\detokenize{Python_create:python-create-lammps}]{\sphinxcrossref{\DUrole{std,std-ref}{Creating or deleting a LAMMPS object}}}}).
\end{sphinxadmonition}


\subsection{Extending Python to run in parallel}
\label{\detokenize{Python_install:extending-python-to-run-in-parallel}}\label{\detokenize{Python_install:python-install-mpi4py}}
\sphinxAtStartPar
If you wish to run LAMMPS in parallel from Python, you need to extend
your Python with an interface to MPI.  This also allows you to
make MPI calls directly from Python in your script, if you desire.

\sphinxAtStartPar
We have tested this with \sphinxhref{https://mpi4py.readthedocs.io/}{MPI for Python}
(aka mpi4py) and you will find installation instruction for it below.

\sphinxAtStartPar
Installation of mpi4py (version 3.0.3 as of Sep 2020) can be done as
follows:
\begin{itemize}
\item {} 
\sphinxAtStartPar
Via \sphinxcode{\sphinxupquote{pip}} into a local user folder with:

\begin{sphinxVerbatim}[commandchars=\\\{\}]
pip\PYG{+w}{ }install\PYG{+w}{ }\PYGZhy{}\PYGZhy{}user\PYG{+w}{ }mpi4py
\end{sphinxVerbatim}

\item {} 
\sphinxAtStartPar
Via \sphinxcode{\sphinxupquote{dnf}} into a system folder for RedHat/Fedora systems:

\begin{sphinxVerbatim}[commandchars=\\\{\}]
\PYG{c+c1}{\PYGZsh{} for use with OpenMPI}
sudo\PYG{+w}{ }dnf\PYG{+w}{ }install\PYG{+w}{ }python3\PYGZhy{}mpi4py\PYGZhy{}openmpi
\PYG{c+c1}{\PYGZsh{} for use with MPICH}
sudo\PYG{+w}{ }dnf\PYG{+w}{ }install\PYG{+w}{ }python3\PYGZhy{}mpi4py\PYGZhy{}openmpi
\end{sphinxVerbatim}

\item {} 
\sphinxAtStartPar
Via \sphinxcode{\sphinxupquote{pip}} into a virtual environment (see above):

\begin{sphinxVerbatim}[commandchars=\\\{\}]
\PYG{g+gp}{\PYGZdl{} }\PYG{n+nb}{source}\PYG{+w}{ }\PYG{n+nv}{\PYGZdl{}HOME}/myenv/activate
\PYG{g+gp+gpVirtualEnv}{(myenv)}\PYG{g+gp}{\PYGZdl{} }pip\PYG{+w}{ }install\PYG{+w}{ }mpi4py
\end{sphinxVerbatim}

\item {} 
\sphinxAtStartPar
Via \sphinxcode{\sphinxupquote{pip}} into a system folder (not recommended):

\begin{sphinxVerbatim}[commandchars=\\\{\}]
sudo\PYG{+w}{ }pip\PYG{+w}{ }install\PYG{+w}{ }mpi4py
\end{sphinxVerbatim}

\end{itemize}

\sphinxAtStartPar
For more detailed installation instructions and additional options,
please see the \sphinxhref{https://mpi4py.readthedocs.io/en/stable/install.html}{mpi4py installation} page.

\sphinxAtStartPar
To use \sphinxcode{\sphinxupquote{mpi4py}} and LAMMPS in parallel from Python, you \sphinxstylestrong{must} make
certain that \sphinxstylestrong{both} are using the \sphinxstylestrong{same} implementation and version
of MPI library.  If you only have one MPI library installed on your
system this is not an issue, but it can be if you have multiple MPI
installations (e.g. on an HPC cluster to be selected through environment
modules).  Your LAMMPS build is explicit about which MPI it is using,
since it is either detected during CMake configuration or in the
traditional make build system you specify the details in your low\sphinxhyphen{}level
\sphinxcode{\sphinxupquote{src/MAKE/Makefile.foo}} file. The installation process of \sphinxcode{\sphinxupquote{mpi4py}}
uses the \sphinxcode{\sphinxupquote{mpicc}} command to find information about the MPI it uses to
build against.  And it tries to load “libmpi.so” from the
\sphinxcode{\sphinxupquote{LD\_LIBRARY\_PATH}}.  This may or may not find the MPI library that
LAMMPS is using.  If you have problems running both mpi4py and LAMMPS
together, this is an issue you may need to address, e.g. by loading the
module for different MPI installation so that mpi4py finds the right
one.

\sphinxAtStartPar
If you have successfully installed mpi4py, you should be able to run
Python and type

\begin{sphinxVerbatim}[commandchars=\\\{\}]
\PYG{k+kn}{from}\PYG{+w}{ }\PYG{n+nn}{mpi4py}\PYG{+w}{ }\PYG{k+kn}{import} \PYG{n}{MPI}
\end{sphinxVerbatim}

\sphinxAtStartPar
without error.  You should also be able to run Python in parallel
on a simple test script

\begin{sphinxVerbatim}[commandchars=\\\{\}]
mpirun\PYG{+w}{ }\PYGZhy{}np\PYG{+w}{ }\PYG{l+m}{4}\PYG{+w}{ }python3\PYG{+w}{ }test.py
\end{sphinxVerbatim}

\sphinxAtStartPar
where \sphinxcode{\sphinxupquote{test.py}} contains the lines

\begin{sphinxVerbatim}[commandchars=\\\{\}]
\PYG{k+kn}{from}\PYG{+w}{ }\PYG{n+nn}{mpi4py}\PYG{+w}{ }\PYG{k+kn}{import} \PYG{n}{MPI}
\PYG{n}{comm} \PYG{o}{=} \PYG{n}{MPI}\PYG{o}{.}\PYG{n}{COMM\PYGZus{}WORLD}
\PYG{n+nb}{print}\PYG{p}{(}\PYG{l+s+s2}{\PYGZdq{}}\PYG{l+s+s2}{Proc }\PYG{l+s+si}{\PYGZpc{}d}\PYG{l+s+s2}{ out of }\PYG{l+s+si}{\PYGZpc{}d}\PYG{l+s+s2}{ procs}\PYG{l+s+s2}{\PYGZdq{}} \PYG{o}{\PYGZpc{}} \PYG{p}{(}\PYG{n}{comm}\PYG{o}{.}\PYG{n}{Get\PYGZus{}rank}\PYG{p}{(}\PYG{p}{)}\PYG{p}{,}\PYG{n}{comm}\PYG{o}{.}\PYG{n}{Get\PYGZus{}size}\PYG{p}{(}\PYG{p}{)}\PYG{p}{)}\PYG{p}{)}
\end{sphinxVerbatim}

\sphinxAtStartPar
and see one line of output for each processor you run on.  Please note
that the order of the lines is not deterministic

\begin{sphinxVerbatim}[commandchars=\\\{\}]
\PYG{g+gp}{\PYGZdl{} }mpirun\PYG{+w}{ }\PYGZhy{}np\PYG{+w}{ }\PYG{l+m}{4}\PYG{+w}{ }python3\PYG{+w}{ }test.py
\PYG{g+go}{Proc 0 out of 4 procs}
\PYG{g+go}{Proc 1 out of 4 procs}
\PYG{g+go}{Proc 2 out of 4 procs}
\PYG{g+go}{Proc 3 out of 4 procs}
\end{sphinxVerbatim}

\sphinxstepscope


\section{Run LAMMPS from Python}
\label{\detokenize{Python_run:run-lammps-from-python}}\label{\detokenize{Python_run::doc}}
\sphinxAtStartPar
After compiling the LAMMPS shared library and making it ready to use,
you can now write and run Python scripts that import the LAMMPS Python
module and launch LAMMPS simulations through Python scripts.
The following pages take you through the various steps necessary,
what functionality is available and give some examples how to use it.


\bigskip\hrule\bigskip


\sphinxstepscope


\subsection{Running LAMMPS and Python in serial}
\label{\detokenize{Python_launch:running-lammps-and-python-in-serial}}\label{\detokenize{Python_launch::doc}}
\sphinxAtStartPar
To run a LAMMPS in serial, type these lines into Python
interactively from the \sphinxcode{\sphinxupquote{bench}} directory:

\begin{sphinxVerbatim}[commandchars=\\\{\}]
\PYG{k+kn}{from}\PYG{+w}{ }\PYG{n+nn}{lammps}\PYG{+w}{ }\PYG{k+kn}{import} \PYG{n}{lammps}
\PYG{n}{lmp} \PYG{o}{=} \PYG{n}{lammps}\PYG{p}{(}\PYG{p}{)}
\PYG{n}{lmp}\PYG{o}{.}\PYG{n}{file}\PYG{p}{(}\PYG{l+s+s2}{\PYGZdq{}}\PYG{l+s+s2}{in.lj}\PYG{l+s+s2}{\PYGZdq{}}\PYG{p}{)}
\end{sphinxVerbatim}

\sphinxAtStartPar
Or put the same lines in the file \sphinxcode{\sphinxupquote{test.py}} and run it as

\begin{sphinxVerbatim}[commandchars=\\\{\}]
python3\PYG{+w}{ }test.py
\end{sphinxVerbatim}

\sphinxAtStartPar
Either way, you should see the results of running the \sphinxcode{\sphinxupquote{in.lj}} benchmark
on a single processor appear on the screen, the same as if you had
typed something like:

\begin{sphinxVerbatim}[commandchars=\\\{\}]
lmp\PYGZus{}serial\PYG{+w}{ }\PYGZhy{}in\PYG{+w}{ }\PYG{k}{in}.lj
\end{sphinxVerbatim}


\subsection{Running LAMMPS and Python in parallel with MPI}
\label{\detokenize{Python_launch:running-lammps-and-python-in-parallel-with-mpi}}
\sphinxAtStartPar
To run LAMMPS in parallel, assuming you have installed the
\sphinxhref{https://mpi4py.readthedocs.io}{mpi4py} package as discussed
{\hyperref[\detokenize{Python_install:python-install-mpi4py}]{\sphinxcrossref{\DUrole{std,std-ref}{Extending Python to run in parallel}}}}, create a \sphinxcode{\sphinxupquote{test.py}} file containing these lines:

\begin{sphinxVerbatim}[commandchars=\\\{\}]
\PYG{k+kn}{from}\PYG{+w}{ }\PYG{n+nn}{mpi4py}\PYG{+w}{ }\PYG{k+kn}{import} \PYG{n}{MPI}
\PYG{k+kn}{from}\PYG{+w}{ }\PYG{n+nn}{lammps}\PYG{+w}{ }\PYG{k+kn}{import} \PYG{n}{lammps}
\PYG{n}{lmp} \PYG{o}{=} \PYG{n}{lammps}\PYG{p}{(}\PYG{p}{)}
\PYG{n}{lmp}\PYG{o}{.}\PYG{n}{file}\PYG{p}{(}\PYG{l+s+s2}{\PYGZdq{}}\PYG{l+s+s2}{in.lj}\PYG{l+s+s2}{\PYGZdq{}}\PYG{p}{)}
\PYG{n}{me} \PYG{o}{=} \PYG{n}{MPI}\PYG{o}{.}\PYG{n}{COMM\PYGZus{}WORLD}\PYG{o}{.}\PYG{n}{Get\PYGZus{}rank}\PYG{p}{(}\PYG{p}{)}
\PYG{n}{nprocs} \PYG{o}{=} \PYG{n}{MPI}\PYG{o}{.}\PYG{n}{COMM\PYGZus{}WORLD}\PYG{o}{.}\PYG{n}{Get\PYGZus{}size}\PYG{p}{(}\PYG{p}{)}
\PYG{n+nb}{print}\PYG{p}{(}\PYG{l+s+s2}{\PYGZdq{}}\PYG{l+s+s2}{Proc }\PYG{l+s+si}{\PYGZpc{}d}\PYG{l+s+s2}{ out of }\PYG{l+s+si}{\PYGZpc{}d}\PYG{l+s+s2}{ procs has}\PYG{l+s+s2}{\PYGZdq{}} \PYG{o}{\PYGZpc{}} \PYG{p}{(}\PYG{n}{me}\PYG{p}{,}\PYG{n}{nprocs}\PYG{p}{)}\PYG{p}{,}\PYG{n}{lmp}\PYG{p}{)}
\PYG{n}{MPI}\PYG{o}{.}\PYG{n}{Finalize}\PYG{p}{(}\PYG{p}{)}
\end{sphinxVerbatim}

\sphinxAtStartPar
You can run the script in parallel as:

\begin{sphinxVerbatim}[commandchars=\\\{\}]
mpirun\PYG{+w}{ }\PYGZhy{}np\PYG{+w}{ }\PYG{l+m}{4}\PYG{+w}{ }python3\PYG{+w}{ }test.py
\end{sphinxVerbatim}

\sphinxAtStartPar
and you should see the same output as if you had typed

\begin{sphinxVerbatim}[commandchars=\\\{\}]
mpirun\PYG{+w}{ }\PYGZhy{}np\PYG{+w}{ }\PYG{l+m}{4}\PYG{+w}{ }lmp\PYGZus{}mpi\PYG{+w}{ }\PYGZhy{}in\PYG{+w}{ }\PYG{k}{in}.lj
\end{sphinxVerbatim}

\sphinxAtStartPar
Note that without the mpi4py specific lines from \sphinxcode{\sphinxupquote{test.py}}

\begin{sphinxVerbatim}[commandchars=\\\{\}]
\PYG{k+kn}{from}\PYG{+w}{ }\PYG{n+nn}{lammps}\PYG{+w}{ }\PYG{k+kn}{import} \PYG{n}{lammps}
\PYG{n}{lmp} \PYG{o}{=} \PYG{n}{lammps}\PYG{p}{(}\PYG{p}{)}
\PYG{n}{lmp}\PYG{o}{.}\PYG{n}{file}\PYG{p}{(}\PYG{l+s+s2}{\PYGZdq{}}\PYG{l+s+s2}{in.lj}\PYG{l+s+s2}{\PYGZdq{}}\PYG{p}{)}
\end{sphinxVerbatim}

\sphinxAtStartPar
running the script with \sphinxcode{\sphinxupquote{mpirun}} on \(P\) processors would lead to
\(P\) independent simulations to run parallel, each with a single
processor. Therefore, if you use the mpi4py lines and you see multiple LAMMPS
single processor outputs, mpi4py is not working correctly.

\sphinxAtStartPar
Also note that once you import the mpi4py module, mpi4py initializes MPI
for you, and you can use MPI calls directly in your Python script, as
described in the mpi4py documentation.  The last line of your Python
script should be \sphinxcode{\sphinxupquote{MPI.finalize()}}, to ensure MPI is shut down
correctly.


\subsection{Running Python scripts}
\label{\detokenize{Python_launch:running-python-scripts}}
\sphinxAtStartPar
Note that any Python script (not just for LAMMPS) can be invoked in
one of several ways:

\begin{sphinxVerbatim}[commandchars=\\\{\}]
python\PYG{+w}{ }script.py
python\PYG{+w}{ }\PYGZhy{}i\PYG{+w}{ }script.py
./script.py
\end{sphinxVerbatim}

\sphinxAtStartPar
The last command requires that the first line of the script be
something like this:

\begin{sphinxVerbatim}[commandchars=\\\{\}]
\PYG{c+ch}{\PYGZsh{}!/usr/bin/python}
\end{sphinxVerbatim}

\sphinxAtStartPar
or

\begin{sphinxVerbatim}[commandchars=\\\{\}]
\PYG{c+ch}{\PYGZsh{}!/usr/bin/env python}
\end{sphinxVerbatim}

\sphinxAtStartPar
where the path in the first case needs to point to where you have Python
installed (the second option is workaround for when this may change),
and that you have made the script file executable:

\begin{sphinxVerbatim}[commandchars=\\\{\}]
chmod\PYG{+w}{ }+x\PYG{+w}{ }script.py
\end{sphinxVerbatim}

\sphinxAtStartPar
Without the \sphinxcode{\sphinxupquote{\sphinxhyphen{}i}} flag, Python will exit when the script finishes.
With the \sphinxcode{\sphinxupquote{\sphinxhyphen{}i}} flag, you will be left in the Python interpreter when
the script finishes, so you can type subsequent commands.  As mentioned
above, you can only run Python interactively when running Python on a
single processor, not in parallel.


\bigskip\hrule\bigskip


\sphinxstepscope


\subsection{Creating or deleting a LAMMPS object}
\label{\detokenize{Python_create:creating-or-deleting-a-lammps-object}}\label{\detokenize{Python_create:python-create-lammps}}\label{\detokenize{Python_create::doc}}
\sphinxAtStartPar
With the Python interface the creation of a {\hyperref[\detokenize{Classes_lammps:_CPPv4N9LAMMPS_NS6LAMMPSE}]{\sphinxcrossref{\sphinxcode{\sphinxupquote{LAMMPS}}}}} instance is included in the constructors for the
{\hyperref[\detokenize{Python_module:lammps.lammps}]{\sphinxcrossref{\sphinxcode{\sphinxupquote{lammps}}}}}, {\hyperref[\detokenize{Python_module:lammps.PyLammps}]{\sphinxcrossref{\sphinxcode{\sphinxupquote{PyLammps}}}}},
and {\hyperref[\detokenize{Python_module:lammps.IPyLammps}]{\sphinxcrossref{\sphinxcode{\sphinxupquote{IPyLammps}}}}} classes.
Internally it will call either {\hyperref[\detokenize{Library_create:_CPPv411lammps_openiPPc8MPI_CommPPv}]{\sphinxcrossref{\sphinxcode{\sphinxupquote{lammps\_open()}}}}} or {\hyperref[\detokenize{Library_create:_CPPv418lammps_open_no_mpiiPPcPPv}]{\sphinxcrossref{\sphinxcode{\sphinxupquote{lammps\_open\_no\_mpi()}}}}} from the C
library API to create the class instance.

\sphinxAtStartPar
All arguments are optional.  The \sphinxstyleemphasis{name} argument allows loading a
LAMMPS shared library that is named \sphinxcode{\sphinxupquote{liblammps\_machine.so}} instead of
the default name of \sphinxcode{\sphinxupquote{liblammps.so}}.  In most cases the latter will be
installed or used.  The \sphinxstyleemphasis{ptr} argument is for use of the
{\hyperref[\detokenize{Python_module:module-lammps}]{\sphinxcrossref{\sphinxcode{\sphinxupquote{lammps}}}}} module from inside a LAMMPS instance, e.g. with the
{\hyperref[\detokenize{python::doc}]{\sphinxcrossref{\DUrole{doc}{python}}}} command, where a pointer to the already existing
{\hyperref[\detokenize{Classes_lammps:_CPPv4N9LAMMPS_NS6LAMMPSE}]{\sphinxcrossref{\sphinxcode{\sphinxupquote{LAMMPS}}}}} class instance can be passed
to the Python class and used instead of creating a new instance.  The
\sphinxstyleemphasis{comm} argument may be used in combination with the \sphinxhref{https://mpi4py.readthedocs.io/}{mpi4py}
module to pass an MPI communicator to LAMMPS and thus it is possible
to run the Python module like the library interface on a subset of the
MPI ranks after splitting the communicator.

\sphinxAtStartPar
Here are simple examples using all three Python interfaces:
\begin{quote}

\begin{sphinxadmonition}{note}{lammps API}

\begin{sphinxVerbatim}[commandchars=\\\{\}]
\PYG{k+kn}{from}\PYG{+w}{ }\PYG{n+nn}{lammps}\PYG{+w}{ }\PYG{k+kn}{import} \PYG{n}{lammps}

\PYG{c+c1}{\PYGZsh{} NOTE: argv[0] is set by the lammps class constructor}
\PYG{n}{args} \PYG{o}{=} \PYG{p}{[}\PYG{l+s+s2}{\PYGZdq{}}\PYG{l+s+s2}{\PYGZhy{}log}\PYG{l+s+s2}{\PYGZdq{}}\PYG{p}{,} \PYG{l+s+s2}{\PYGZdq{}}\PYG{l+s+s2}{none}\PYG{l+s+s2}{\PYGZdq{}}\PYG{p}{]}

\PYG{c+c1}{\PYGZsh{} create LAMMPS instance}
\PYG{n}{lmp} \PYG{o}{=} \PYG{n}{lammps}\PYG{p}{(}\PYG{n}{cmdargs}\PYG{o}{=}\PYG{n}{args}\PYG{p}{)}

\PYG{c+c1}{\PYGZsh{} get and print numerical version code}
\PYG{n+nb}{print}\PYG{p}{(}\PYG{l+s+s2}{\PYGZdq{}}\PYG{l+s+s2}{LAMMPS Version: }\PYG{l+s+s2}{\PYGZdq{}}\PYG{p}{,} \PYG{n}{lmp}\PYG{o}{.}\PYG{n}{version}\PYG{p}{(}\PYG{p}{)}\PYG{p}{)}

\PYG{c+c1}{\PYGZsh{} explicitly close and delete LAMMPS instance (optional)}
\PYG{n}{lmp}\PYG{o}{.}\PYG{n}{close}\PYG{p}{(}\PYG{p}{)}
\end{sphinxVerbatim}
\end{sphinxadmonition}

\begin{sphinxadmonition}{note}{PyLammps API}

\sphinxAtStartPar
The {\hyperref[\detokenize{Python_module:lammps.PyLammps}]{\sphinxcrossref{\sphinxcode{\sphinxupquote{PyLammps}}}}} class is a wrapper around the
{\hyperref[\detokenize{Python_module:lammps.lammps}]{\sphinxcrossref{\sphinxcode{\sphinxupquote{lammps}}}}} class and all of its lower level functions.
By default, it will create a new instance of {\hyperref[\detokenize{Python_module:lammps.lammps}]{\sphinxcrossref{\sphinxcode{\sphinxupquote{lammps}}}}} passing
along all arguments to the constructor of {\hyperref[\detokenize{Python_module:lammps.lammps}]{\sphinxcrossref{\sphinxcode{\sphinxupquote{lammps}}}}}.

\begin{sphinxVerbatim}[commandchars=\\\{\}]
\PYG{k+kn}{from}\PYG{+w}{ }\PYG{n+nn}{lammps}\PYG{+w}{ }\PYG{k+kn}{import} \PYG{n}{PyLammps}

\PYG{c+c1}{\PYGZsh{} NOTE: argv[0] is set by the lammps class constructor}
\PYG{n}{args} \PYG{o}{=} \PYG{p}{[}\PYG{l+s+s2}{\PYGZdq{}}\PYG{l+s+s2}{\PYGZhy{}log}\PYG{l+s+s2}{\PYGZdq{}}\PYG{p}{,} \PYG{l+s+s2}{\PYGZdq{}}\PYG{l+s+s2}{none}\PYG{l+s+s2}{\PYGZdq{}}\PYG{p}{]}

\PYG{c+c1}{\PYGZsh{} create LAMMPS instance}
\PYG{n}{L} \PYG{o}{=} \PYG{n}{PyLammps}\PYG{p}{(}\PYG{n}{cmdargs}\PYG{o}{=}\PYG{n}{args}\PYG{p}{)}

\PYG{c+c1}{\PYGZsh{} get and print numerical version code}
\PYG{n+nb}{print}\PYG{p}{(}\PYG{l+s+s2}{\PYGZdq{}}\PYG{l+s+s2}{LAMMPS Version: }\PYG{l+s+s2}{\PYGZdq{}}\PYG{p}{,} \PYG{n}{L}\PYG{o}{.}\PYG{n}{version}\PYG{p}{(}\PYG{p}{)}\PYG{p}{)}

\PYG{c+c1}{\PYGZsh{} explicitly close and delete LAMMPS instance (optional)}
\PYG{n}{L}\PYG{o}{.}\PYG{n}{close}\PYG{p}{(}\PYG{p}{)}
\end{sphinxVerbatim}

\sphinxAtStartPar
{\hyperref[\detokenize{Python_module:lammps.PyLammps}]{\sphinxcrossref{\sphinxcode{\sphinxupquote{PyLammps}}}}} objects can also be created on top of an existing
{\hyperref[\detokenize{Python_module:lammps.lammps}]{\sphinxcrossref{\sphinxcode{\sphinxupquote{lammps}}}}} object:

\begin{sphinxVerbatim}[commandchars=\\\{\}]
\PYG{k+kn}{from}\PYG{+w}{ }\PYG{n+nn}{lammps}\PYG{+w}{ }\PYG{k+kn}{import} \PYG{n}{lammps}\PYG{p}{,} \PYG{n}{PyLammps}
\PYG{o}{.}\PYG{o}{.}\PYG{o}{.}
\PYG{c+c1}{\PYGZsh{} create LAMMPS instance}
\PYG{n}{lmp} \PYG{o}{=} \PYG{n}{lammps}\PYG{p}{(}\PYG{n}{cmdargs}\PYG{o}{=}\PYG{n}{args}\PYG{p}{)}

\PYG{c+c1}{\PYGZsh{} create PyLammps instance using previously created LAMMPS instance}
\PYG{n}{L} \PYG{o}{=} \PYG{n}{PyLammps}\PYG{p}{(}\PYG{n}{ptr}\PYG{o}{=}\PYG{n}{lmp}\PYG{p}{)}
\end{sphinxVerbatim}

\sphinxAtStartPar
This is useful if you have to create the {\hyperref[\detokenize{Python_module:lammps.lammps}]{\sphinxcrossref{\sphinxcode{\sphinxupquote{lammps}}}}}
instance is a specific way, but want to take advantage of the
{\hyperref[\detokenize{Python_module:lammps.PyLammps}]{\sphinxcrossref{\sphinxcode{\sphinxupquote{PyLammps}}}}} interface.
\end{sphinxadmonition}

\begin{sphinxadmonition}{note}{IPyLammps API}

\sphinxAtStartPar
The {\hyperref[\detokenize{Python_module:lammps.IPyLammps}]{\sphinxcrossref{\sphinxcode{\sphinxupquote{IPyLammps}}}}} class is an extension of the
{\hyperref[\detokenize{Python_module:lammps.PyLammps}]{\sphinxcrossref{\sphinxcode{\sphinxupquote{PyLammps}}}}} class. It has the same construction behavior. By
default, it will create a new instance of {\hyperref[\detokenize{Python_module:module-lammps}]{\sphinxcrossref{\sphinxcode{\sphinxupquote{lammps}}}}} passing
along all arguments to the constructor of {\hyperref[\detokenize{Python_module:module-lammps}]{\sphinxcrossref{\sphinxcode{\sphinxupquote{lammps}}}}}.

\begin{sphinxVerbatim}[commandchars=\\\{\}]
\PYG{k+kn}{from}\PYG{+w}{ }\PYG{n+nn}{lammps}\PYG{+w}{ }\PYG{k+kn}{import} \PYG{n}{IPyLammps}

\PYG{c+c1}{\PYGZsh{} NOTE: argv[0] is set by the lammps class constructor}
\PYG{n}{args} \PYG{o}{=} \PYG{p}{[}\PYG{l+s+s2}{\PYGZdq{}}\PYG{l+s+s2}{\PYGZhy{}log}\PYG{l+s+s2}{\PYGZdq{}}\PYG{p}{,} \PYG{l+s+s2}{\PYGZdq{}}\PYG{l+s+s2}{none}\PYG{l+s+s2}{\PYGZdq{}}\PYG{p}{]}

\PYG{c+c1}{\PYGZsh{} create LAMMPS instance}
\PYG{n}{L} \PYG{o}{=} \PYG{n}{IPyLammps}\PYG{p}{(}\PYG{n}{cmdargs}\PYG{o}{=}\PYG{n}{args}\PYG{p}{)}

\PYG{c+c1}{\PYGZsh{} get and print numerical version code}
\PYG{n+nb}{print}\PYG{p}{(}\PYG{l+s+s2}{\PYGZdq{}}\PYG{l+s+s2}{LAMMPS Version: }\PYG{l+s+s2}{\PYGZdq{}}\PYG{p}{,} \PYG{n}{L}\PYG{o}{.}\PYG{n}{version}\PYG{p}{(}\PYG{p}{)}\PYG{p}{)}

\PYG{c+c1}{\PYGZsh{} explicitly close and delete LAMMPS instance (optional)}
\PYG{n}{L}\PYG{o}{.}\PYG{n}{close}\PYG{p}{(}\PYG{p}{)}
\end{sphinxVerbatim}

\sphinxAtStartPar
You can also initialize IPyLammps on top of an existing {\hyperref[\detokenize{Python_module:module-lammps}]{\sphinxcrossref{\sphinxcode{\sphinxupquote{lammps}}}}} or \sphinxcode{\sphinxupquote{PyLammps}} object:

\begin{sphinxVerbatim}[commandchars=\\\{\}]
\PYG{k+kn}{from}\PYG{+w}{ }\PYG{n+nn}{lammps}\PYG{+w}{ }\PYG{k+kn}{import} \PYG{n}{lammps}\PYG{p}{,} \PYG{n}{IPyLammps}
\PYG{o}{.}\PYG{o}{.}\PYG{o}{.}
\PYG{c+c1}{\PYGZsh{} create LAMMPS instance}
\PYG{n}{lmp} \PYG{o}{=} \PYG{n}{lammps}\PYG{p}{(}\PYG{n}{cmdargs}\PYG{o}{=}\PYG{n}{args}\PYG{p}{)}

\PYG{c+c1}{\PYGZsh{} create PyLammps instance using previously created LAMMPS instance}
\PYG{n}{L} \PYG{o}{=} \PYG{n}{PyLammps}\PYG{p}{(}\PYG{n}{ptr}\PYG{o}{=}\PYG{n}{lmp}\PYG{p}{)}
\end{sphinxVerbatim}

\sphinxAtStartPar
This is useful if you have to create the {\hyperref[\detokenize{Python_module:lammps.lammps}]{\sphinxcrossref{\sphinxcode{\sphinxupquote{lammps}}}}}
instance is a specific way, but want to take advantage of the
{\hyperref[\detokenize{Python_module:lammps.IPyLammps}]{\sphinxcrossref{\sphinxcode{\sphinxupquote{IPyLammps}}}}} interface.
\end{sphinxadmonition}
\end{quote}

\sphinxAtStartPar
In all of the above cases, same as with the {\hyperref[\detokenize{Library:lammps-c-api}]{\sphinxcrossref{\DUrole{std,std-ref}{C library API}}}}, this will use the
\sphinxcode{\sphinxupquote{MPI\_COMM\_WORLD}} communicator for the MPI library that LAMMPS was
compiled with.

\sphinxAtStartPar
The {\hyperref[\detokenize{Python_module:lammps.lammps.close}]{\sphinxcrossref{\sphinxcode{\sphinxupquote{lmp.close()}}}}} call is
optional since the LAMMPS class instance will also be deleted
automatically during the {\hyperref[\detokenize{Python_module:lammps.lammps}]{\sphinxcrossref{\sphinxcode{\sphinxupquote{lammps}}}}} class
destructor.  Instead of {\hyperref[\detokenize{Python_module:lammps.lammps.close}]{\sphinxcrossref{\sphinxcode{\sphinxupquote{lmp.close()}}}}}
it is also possible to call {\hyperref[\detokenize{Python_module:lammps.lammps.finalize}]{\sphinxcrossref{\sphinxcode{\sphinxupquote{lmp.finalize()}}}}};
this will destruct the LAMMPS instance, but also finalized and release
the MPI and/or Kokkos environment if enabled and active.

\sphinxAtStartPar
Note that you can create multiple LAMMPS objects in your Python
script, and coordinate and run multiple simulations, e.g.

\begin{sphinxVerbatim}[commandchars=\\\{\}]
\PYG{k+kn}{from}\PYG{+w}{ }\PYG{n+nn}{lammps}\PYG{+w}{ }\PYG{k+kn}{import} \PYG{n}{lammps}
\PYG{n}{lmp1} \PYG{o}{=} \PYG{n}{lammps}\PYG{p}{(}\PYG{p}{)}
\PYG{n}{lmp2} \PYG{o}{=} \PYG{n}{lammps}\PYG{p}{(}\PYG{p}{)}
\PYG{n}{lmp1}\PYG{o}{.}\PYG{n}{file}\PYG{p}{(}\PYG{l+s+s2}{\PYGZdq{}}\PYG{l+s+s2}{in.file1}\PYG{l+s+s2}{\PYGZdq{}}\PYG{p}{)}
\PYG{n}{lmp2}\PYG{o}{.}\PYG{n}{file}\PYG{p}{(}\PYG{l+s+s2}{\PYGZdq{}}\PYG{l+s+s2}{in.file2}\PYG{l+s+s2}{\PYGZdq{}}\PYG{p}{)}
\end{sphinxVerbatim}

\sphinxstepscope


\subsection{Executing commands}
\label{\detokenize{Python_execute:executing-commands}}\label{\detokenize{Python_execute::doc}}
\sphinxAtStartPar
Once an instance of the {\hyperref[\detokenize{Python_module:lammps.lammps}]{\sphinxcrossref{\sphinxcode{\sphinxupquote{lammps}}}}},
{\hyperref[\detokenize{Python_module:lammps.PyLammps}]{\sphinxcrossref{\sphinxcode{\sphinxupquote{PyLammps}}}}}, or
{\hyperref[\detokenize{Python_module:lammps.IPyLammps}]{\sphinxcrossref{\sphinxcode{\sphinxupquote{IPyLammps}}}}} class is created, there are
multiple ways to “feed” it commands. In a way that is not very different from
running a LAMMPS input script, except that Python has many more facilities
for structured programming than the LAMMPS input script syntax. Furthermore
it is possible to “compute” what the next LAMMPS command should be.
\begin{quote}

\begin{sphinxadmonition}{note}{lammps API}

\sphinxAtStartPar
Same as in the equivalent
{\hyperref[\detokenize{Library_execute::doc}]{\sphinxcrossref{\DUrole{doc}{C library functions}}}}, commands can be read from a file, a
single string, a list of strings and a block of commands in a single
multi\sphinxhyphen{}line string. They are processed under the same boundary conditions
as the C library counterparts.  The example below demonstrates the use
of \sphinxcode{\sphinxupquote{lammps.file()}}, \sphinxcode{\sphinxupquote{lammps.command()}},
\sphinxcode{\sphinxupquote{lammps.commands\_list()}}, and \sphinxcode{\sphinxupquote{lammps.commands\_string()}}:

\begin{sphinxVerbatim}[commandchars=\\\{\}]
\PYG{k+kn}{from}\PYG{+w}{ }\PYG{n+nn}{lammps}\PYG{+w}{ }\PYG{k+kn}{import} \PYG{n}{lammps}
\PYG{n}{lmp} \PYG{o}{=} \PYG{n}{lammps}\PYG{p}{(}\PYG{p}{)}

\PYG{c+c1}{\PYGZsh{} read commands from file \PYGZsq{}in.melt\PYGZsq{}}
\PYG{n}{lmp}\PYG{o}{.}\PYG{n}{file}\PYG{p}{(}\PYG{l+s+s1}{\PYGZsq{}}\PYG{l+s+s1}{in.melt}\PYG{l+s+s1}{\PYGZsq{}}\PYG{p}{)}

\PYG{c+c1}{\PYGZsh{} issue a single command}
\PYG{n}{lmp}\PYG{o}{.}\PYG{n}{command}\PYG{p}{(}\PYG{l+s+s1}{\PYGZsq{}}\PYG{l+s+s1}{variable zpos index 1.0}\PYG{l+s+s1}{\PYGZsq{}}\PYG{p}{)}

\PYG{c+c1}{\PYGZsh{} create 10 groups with 10 atoms each}
\PYG{n}{cmds} \PYG{o}{=} \PYG{p}{[}\PYG{l+s+s2}{\PYGZdq{}}\PYG{l+s+s2}{group g}\PYG{l+s+si}{\PYGZob{}\PYGZcb{}}\PYG{l+s+s2}{ id }\PYG{l+s+si}{\PYGZob{}\PYGZcb{}}\PYG{l+s+s2}{:}\PYG{l+s+si}{\PYGZob{}\PYGZcb{}}\PYG{l+s+s2}{\PYGZdq{}}\PYG{o}{.}\PYG{n}{format}\PYG{p}{(}\PYG{n}{i}\PYG{p}{,}\PYG{l+m+mi}{10}\PYG{o}{*}\PYG{n}{i}\PYG{o}{+}\PYG{l+m+mi}{1}\PYG{p}{,}\PYG{l+m+mi}{10}\PYG{o}{*}\PYG{p}{(}\PYG{n}{i}\PYG{o}{+}\PYG{l+m+mi}{1}\PYG{p}{)}\PYG{p}{)} \PYG{k}{for} \PYG{n}{i} \PYG{o+ow}{in} \PYG{n+nb}{range}\PYG{p}{(}\PYG{l+m+mi}{10}\PYG{p}{)}\PYG{p}{]}
\PYG{n}{lmp}\PYG{o}{.}\PYG{n}{commands\PYGZus{}list}\PYG{p}{(}\PYG{n}{cmds}\PYG{p}{)}

\PYG{c+c1}{\PYGZsh{} run commands from a multi\PYGZhy{}line string}
\PYG{n}{block} \PYG{o}{=} \PYG{l+s+s2}{\PYGZdq{}\PYGZdq{}\PYGZdq{}}
\PYG{l+s+s2}{clear}
\PYG{l+s+s2}{region  box block 0 2 0 2 0 2}
\PYG{l+s+s2}{create\PYGZus{}box 1 box}
\PYG{l+s+s2}{create\PYGZus{}atoms 1 single 1.0 1.0 \PYGZdl{}}\PYG{l+s+si}{\PYGZob{}zpos\PYGZcb{}}
\PYG{l+s+s2}{\PYGZdq{}\PYGZdq{}\PYGZdq{}}
\PYG{n}{lmp}\PYG{o}{.}\PYG{n}{commands\PYGZus{}string}\PYG{p}{(}\PYG{n}{block}\PYG{p}{)}
\end{sphinxVerbatim}
\end{sphinxadmonition}

\begin{sphinxadmonition}{note}{PyLammps/IPyLammps API}

\sphinxAtStartPar
Unlike the lammps API, the PyLammps/IPyLammps APIs allow running LAMMPS
commands by calling equivalent member functions of {\hyperref[\detokenize{Python_module:lammps.PyLammps}]{\sphinxcrossref{\sphinxcode{\sphinxupquote{PyLammps}}}}}
and {\hyperref[\detokenize{Python_module:lammps.IPyLammps}]{\sphinxcrossref{\sphinxcode{\sphinxupquote{IPyLammps}}}}} instances.

\sphinxAtStartPar
For instance, the following LAMMPS command

\begin{sphinxVerbatim}[commandchars=\\\{\}]
\PYG{k}{region }\PYG{n+nv+nvIdentifier}{box}\PYG{+w}{ }\PYG{n}{block}\PYG{+w}{ }\PYG{l+m}{0}\PYG{+w}{ }\PYG{l+m}{10}\PYG{+w}{ }\PYG{l+m}{0}\PYG{+w}{ }\PYG{l+m}{5}\PYG{+w}{ }\PYG{o}{\PYGZhy{}}\PYG{l+m}{0.5}\PYG{+w}{ }\PYG{l+m}{0.5}
\end{sphinxVerbatim}

\sphinxAtStartPar
can be executed using with the lammps API with the following Python code if \sphinxcode{\sphinxupquote{lmp}} is an
instance of {\hyperref[\detokenize{Python_module:lammps.lammps}]{\sphinxcrossref{\sphinxcode{\sphinxupquote{lammps}}}}}:

\begin{sphinxVerbatim}[commandchars=\\\{\}]
\PYG{k+kn}{from}\PYG{+w}{ }\PYG{n+nn}{lammps}\PYG{+w}{ }\PYG{k+kn}{import} \PYG{n}{lammps}

\PYG{n}{lmp} \PYG{o}{=} \PYG{n}{lammps}\PYG{p}{(}\PYG{p}{)}
\PYG{n}{lmp}\PYG{o}{.}\PYG{n}{command}\PYG{p}{(}\PYG{l+s+s2}{\PYGZdq{}}\PYG{l+s+s2}{region box block 0 10 0 5 \PYGZhy{}0.5 0.5}\PYG{l+s+s2}{\PYGZdq{}}\PYG{p}{)}
\end{sphinxVerbatim}

\sphinxAtStartPar
With the PyLammps interface, any LAMMPS command can be split up into arbitrary parts.
These parts are then passed to a member function with the name of the {\hyperref[\detokenize{Commands_all::doc}]{\sphinxcrossref{\DUrole{doc}{command}}}}.
For the {\hyperref[\detokenize{region::doc}]{\sphinxcrossref{\DUrole{doc}{region}}}} command that means the \sphinxcode{\sphinxupquote{region()}} method can be called.
The arguments of the command can be passed as one string, or
individually.

\begin{sphinxVerbatim}[commandchars=\\\{\}]
\PYG{k+kn}{from}\PYG{+w}{ }\PYG{n+nn}{lammps}\PYG{+w}{ }\PYG{k+kn}{import} \PYG{n}{PyLammps}

\PYG{n}{L} \PYG{o}{=} \PYG{n}{PyLammps}\PYG{p}{(}\PYG{p}{)}

\PYG{c+c1}{\PYGZsh{} pass command parameters as one string}
\PYG{n}{L}\PYG{o}{.}\PYG{n}{region}\PYG{p}{(}\PYG{l+s+s2}{\PYGZdq{}}\PYG{l+s+s2}{box block 0 10 0 5 \PYGZhy{}0.5 0.5}\PYG{l+s+s2}{\PYGZdq{}}\PYG{p}{)}

\PYG{c+c1}{\PYGZsh{} OR pass them individually}
\PYG{n}{L}\PYG{o}{.}\PYG{n}{region}\PYG{p}{(}\PYG{l+s+s2}{\PYGZdq{}}\PYG{l+s+s2}{box block}\PYG{l+s+s2}{\PYGZdq{}}\PYG{p}{,} \PYG{l+m+mi}{0}\PYG{p}{,} \PYG{l+m+mi}{10}\PYG{p}{,} \PYG{l+m+mi}{0}\PYG{p}{,} \PYG{l+m+mi}{5}\PYG{p}{,} \PYG{o}{\PYGZhy{}}\PYG{l+m+mf}{0.5}\PYG{p}{,} \PYG{l+m+mf}{0.5}\PYG{p}{)}
\end{sphinxVerbatim}

\sphinxAtStartPar
In the latter example, all parameters except the first are Python floating\sphinxhyphen{}point literals. The
member function takes the entire parameter list and transparently merges it to a single command
string.

\sphinxAtStartPar
The benefit of this approach is avoiding redundant command calls and easier
parameterization. In the lammps API parameterization needed to be done
manually by creating formatted command strings.

\begin{sphinxVerbatim}[commandchars=\\\{\}]
\PYG{n}{lmp}\PYG{o}{.}\PYG{n}{command}\PYG{p}{(}\PYG{l+s+s2}{\PYGZdq{}}\PYG{l+s+s2}{region box block }\PYG{l+s+si}{\PYGZpc{}f}\PYG{l+s+s2}{ }\PYG{l+s+si}{\PYGZpc{}f}\PYG{l+s+s2}{ }\PYG{l+s+si}{\PYGZpc{}f}\PYG{l+s+s2}{ }\PYG{l+s+si}{\PYGZpc{}f}\PYG{l+s+s2}{ }\PYG{l+s+si}{\PYGZpc{}f}\PYG{l+s+s2}{ }\PYG{l+s+si}{\PYGZpc{}f}\PYG{l+s+s2}{\PYGZdq{}} \PYG{o}{\PYGZpc{}} \PYG{p}{(}\PYG{n}{xlo}\PYG{p}{,} \PYG{n}{xhi}\PYG{p}{,} \PYG{n}{ylo}\PYG{p}{,} \PYG{n}{yhi}\PYG{p}{,} \PYG{n}{zlo}\PYG{p}{,} \PYG{n}{zhi}\PYG{p}{)}\PYG{p}{)}
\end{sphinxVerbatim}

\sphinxAtStartPar
In contrast, methods of PyLammps accept parameters directly and will convert
them automatically to a final command string.

\begin{sphinxVerbatim}[commandchars=\\\{\}]
\PYG{n}{L}\PYG{o}{.}\PYG{n}{region}\PYG{p}{(}\PYG{l+s+s2}{\PYGZdq{}}\PYG{l+s+s2}{box block}\PYG{l+s+s2}{\PYGZdq{}}\PYG{p}{,} \PYG{n}{xlo}\PYG{p}{,} \PYG{n}{xhi}\PYG{p}{,} \PYG{n}{ylo}\PYG{p}{,} \PYG{n}{yhi}\PYG{p}{,} \PYG{n}{zlo}\PYG{p}{,} \PYG{n}{zhi}\PYG{p}{)}
\end{sphinxVerbatim}

\sphinxAtStartPar
Using these facilities, the example shown for the lammps API can be rewritten as follows:

\begin{sphinxVerbatim}[commandchars=\\\{\}]
\PYG{k+kn}{from}\PYG{+w}{ }\PYG{n+nn}{lammps}\PYG{+w}{ }\PYG{k+kn}{import} \PYG{n}{PyLammps}
\PYG{n}{L} \PYG{o}{=} \PYG{n}{PyLammps}\PYG{p}{(}\PYG{p}{)}

\PYG{c+c1}{\PYGZsh{} read commands from file \PYGZsq{}in.melt\PYGZsq{}}
\PYG{n}{L}\PYG{o}{.}\PYG{n}{file}\PYG{p}{(}\PYG{l+s+s1}{\PYGZsq{}}\PYG{l+s+s1}{in.melt}\PYG{l+s+s1}{\PYGZsq{}}\PYG{p}{)}

\PYG{c+c1}{\PYGZsh{} issue a single command}
\PYG{n}{L}\PYG{o}{.}\PYG{n}{variable}\PYG{p}{(}\PYG{l+s+s1}{\PYGZsq{}}\PYG{l+s+s1}{zpos}\PYG{l+s+s1}{\PYGZsq{}}\PYG{p}{,} \PYG{l+s+s1}{\PYGZsq{}}\PYG{l+s+s1}{index}\PYG{l+s+s1}{\PYGZsq{}}\PYG{p}{,} \PYG{l+m+mf}{1.0}\PYG{p}{)}

\PYG{c+c1}{\PYGZsh{} create 10 groups with 10 atoms each}
\PYG{k}{for} \PYG{n}{i} \PYG{o+ow}{in} \PYG{n+nb}{range}\PYG{p}{(}\PYG{l+m+mi}{10}\PYG{p}{)}\PYG{p}{:}
   \PYG{n}{L}\PYG{o}{.}\PYG{n}{group}\PYG{p}{(}\PYG{l+s+sa}{f}\PYG{l+s+s2}{\PYGZdq{}}\PYG{l+s+s2}{g}\PYG{l+s+si}{\PYGZob{}}\PYG{n}{i}\PYG{l+s+si}{\PYGZcb{}}\PYG{l+s+s2}{\PYGZdq{}}\PYG{p}{,} \PYG{l+s+s2}{\PYGZdq{}}\PYG{l+s+s2}{id}\PYG{l+s+s2}{\PYGZdq{}}\PYG{p}{,} \PYG{l+s+sa}{f}\PYG{l+s+s2}{\PYGZdq{}}\PYG{l+s+si}{\PYGZob{}}\PYG{l+m+mi}{10}\PYG{o}{*}\PYG{n}{i}\PYG{o}{+}\PYG{l+m+mi}{1}\PYG{l+s+si}{\PYGZcb{}}\PYG{l+s+s2}{:}\PYG{l+s+si}{\PYGZob{}}\PYG{l+m+mi}{10}\PYG{o}{*}\PYG{p}{(}\PYG{n}{i}\PYG{o}{+}\PYG{l+m+mi}{1}\PYG{p}{)}\PYG{l+s+si}{\PYGZcb{}}\PYG{l+s+s2}{\PYGZdq{}}\PYG{p}{)}

\PYG{n}{L}\PYG{o}{.}\PYG{n}{clear}\PYG{p}{(}\PYG{p}{)}
\PYG{n}{L}\PYG{o}{.}\PYG{n}{region}\PYG{p}{(}\PYG{l+s+s2}{\PYGZdq{}}\PYG{l+s+s2}{box block}\PYG{l+s+s2}{\PYGZdq{}}\PYG{p}{,} \PYG{l+m+mi}{0}\PYG{p}{,} \PYG{l+m+mi}{2}\PYG{p}{,} \PYG{l+m+mi}{0}\PYG{p}{,} \PYG{l+m+mi}{2}\PYG{p}{,} \PYG{l+m+mi}{0}\PYG{p}{,} \PYG{l+m+mi}{2}\PYG{p}{)}
\PYG{n}{L}\PYG{o}{.}\PYG{n}{create\PYGZus{}box}\PYG{p}{(}\PYG{l+m+mi}{1}\PYG{p}{,} \PYG{l+s+s2}{\PYGZdq{}}\PYG{l+s+s2}{box}\PYG{l+s+s2}{\PYGZdq{}}\PYG{p}{)}
\PYG{n}{L}\PYG{o}{.}\PYG{n}{create\PYGZus{}atoms}\PYG{p}{(}\PYG{l+m+mi}{1}\PYG{p}{,} \PYG{l+s+s2}{\PYGZdq{}}\PYG{l+s+s2}{single}\PYG{l+s+s2}{\PYGZdq{}}\PYG{p}{,} \PYG{l+m+mf}{1.0}\PYG{p}{,} \PYG{l+m+mf}{1.0}\PYG{p}{,} \PYG{l+s+s2}{\PYGZdq{}}\PYG{l+s+s2}{\PYGZdl{}}\PYG{l+s+si}{\PYGZob{}zpos\PYGZcb{}}\PYG{l+s+s2}{\PYGZdq{}}\PYG{p}{)}
\end{sphinxVerbatim}
\end{sphinxadmonition}
\end{quote}

\sphinxstepscope


\subsection{System properties}
\label{\detokenize{Python_properties:system-properties}}\label{\detokenize{Python_properties::doc}}
\sphinxAtStartPar
Similar to what is described in {\hyperref[\detokenize{Library_properties::doc}]{\sphinxcrossref{\DUrole{doc}{System properties}}}}, the instances of
{\hyperref[\detokenize{Python_module:lammps.lammps}]{\sphinxcrossref{\sphinxcode{\sphinxupquote{lammps}}}}}, {\hyperref[\detokenize{Python_module:lammps.PyLammps}]{\sphinxcrossref{\sphinxcode{\sphinxupquote{PyLammps}}}}}, or
{\hyperref[\detokenize{Python_module:lammps.IPyLammps}]{\sphinxcrossref{\sphinxcode{\sphinxupquote{IPyLammps}}}}} can be used to extract different kinds
of information about the active LAMMPS instance and also to modify some of it. The
main difference between the interfaces is how the information is exposed.

\sphinxAtStartPar
While the {\hyperref[\detokenize{Python_module:lammps.lammps}]{\sphinxcrossref{\sphinxcode{\sphinxupquote{lammps}}}}} is just a thin layer that wraps C API calls,
{\hyperref[\detokenize{Python_module:lammps.PyLammps}]{\sphinxcrossref{\sphinxcode{\sphinxupquote{PyLammps}}}}} and {\hyperref[\detokenize{Python_module:lammps.IPyLammps}]{\sphinxcrossref{\sphinxcode{\sphinxupquote{IPyLammps}}}}} expose
information as objects and properties.

\sphinxAtStartPar
In some cases the data returned is a direct reference to the original data
inside LAMMPS cast to \sphinxcode{\sphinxupquote{ctypes}} pointers. Where possible, the wrappers will
determine the \sphinxcode{\sphinxupquote{ctypes}} data type and cast pointers accordingly. If
\sphinxcode{\sphinxupquote{numpy}} is installed arrays can also be extracted as numpy arrays, which
will access the C arrays directly and have the correct dimensions to protect
against invalid accesses.

\begin{sphinxadmonition}{warning}{Warning:}
\sphinxAtStartPar
When accessing per\sphinxhyphen{}atom data,
please note that this data is the per\sphinxhyphen{}processor local data and indexed
accordingly. These arrays can change sizes and order at every neighbor list
rebuild and atom sort event as atoms are migrating between subdomains.

\begin{sphinxadmonition}{note}{lammps API}

\begin{sphinxVerbatim}[commandchars=\\\{\}]
\PYG{k+kn}{from}\PYG{+w}{ }\PYG{n+nn}{lammps}\PYG{+w}{ }\PYG{k+kn}{import} \PYG{n}{lammps}

\PYG{n}{lmp} \PYG{o}{=} \PYG{n}{lammps}\PYG{p}{(}\PYG{p}{)}
\PYG{n}{lmp}\PYG{o}{.}\PYG{n}{file}\PYG{p}{(}\PYG{l+s+s2}{\PYGZdq{}}\PYG{l+s+s2}{in.sysinit}\PYG{l+s+s2}{\PYGZdq{}}\PYG{p}{)}

\PYG{n}{natoms} \PYG{o}{=} \PYG{n}{lmp}\PYG{o}{.}\PYG{n}{get\PYGZus{}natoms}\PYG{p}{(}\PYG{p}{)}
\PYG{n+nb}{print}\PYG{p}{(}\PYG{l+s+sa}{f}\PYG{l+s+s2}{\PYGZdq{}}\PYG{l+s+s2}{running simulation with }\PYG{l+s+si}{\PYGZob{}}\PYG{n}{natoms}\PYG{l+s+si}{\PYGZcb{}}\PYG{l+s+s2}{ atoms}\PYG{l+s+s2}{\PYGZdq{}}\PYG{p}{)}

\PYG{n}{lmp}\PYG{o}{.}\PYG{n}{command}\PYG{p}{(}\PYG{l+s+s2}{\PYGZdq{}}\PYG{l+s+s2}{run 1000 post no}\PYG{l+s+s2}{\PYGZdq{}}\PYG{p}{)}\PYG{p}{;}

\PYG{k}{for} \PYG{n}{i} \PYG{o+ow}{in} \PYG{n+nb}{range}\PYG{p}{(}\PYG{l+m+mi}{10}\PYG{p}{)}\PYG{p}{:}
   \PYG{n}{lmp}\PYG{o}{.}\PYG{n}{command}\PYG{p}{(}\PYG{l+s+s2}{\PYGZdq{}}\PYG{l+s+s2}{run 100 pre no post no}\PYG{l+s+s2}{\PYGZdq{}}\PYG{p}{)}
   \PYG{n}{pe} \PYG{o}{=} \PYG{n}{lmp}\PYG{o}{.}\PYG{n}{get\PYGZus{}thermo}\PYG{p}{(}\PYG{l+s+s2}{\PYGZdq{}}\PYG{l+s+s2}{pe}\PYG{l+s+s2}{\PYGZdq{}}\PYG{p}{)}
   \PYG{n}{ke} \PYG{o}{=} \PYG{n}{lmp}\PYG{o}{.}\PYG{n}{get\PYGZus{}thermo}\PYG{p}{(}\PYG{l+s+s2}{\PYGZdq{}}\PYG{l+s+s2}{ke}\PYG{l+s+s2}{\PYGZdq{}}\PYG{p}{)}
   \PYG{n+nb}{print}\PYG{p}{(}\PYG{l+s+sa}{f}\PYG{l+s+s2}{\PYGZdq{}}\PYG{l+s+s2}{PE = }\PYG{l+s+si}{\PYGZob{}}\PYG{n}{pe}\PYG{l+s+si}{\PYGZcb{}}\PYG{l+s+se}{\PYGZbs{}n}\PYG{l+s+s2}{KE = }\PYG{l+s+si}{\PYGZob{}}\PYG{n}{ke}\PYG{l+s+si}{\PYGZcb{}}\PYG{l+s+s2}{\PYGZdq{}}\PYG{p}{)}

\PYG{n}{lmp}\PYG{o}{.}\PYG{n}{close}\PYG{p}{(}\PYG{p}{)}
\end{sphinxVerbatim}

\sphinxAtStartPar
\sphinxstylestrong{Methods}:
\begin{itemize}
\item {} 
\sphinxAtStartPar
{\hyperref[\detokenize{Python_module:lammps.lammps.version}]{\sphinxcrossref{\sphinxcode{\sphinxupquote{version()}}}}}: return the numerical version id, e.g. LAMMPS 2 Sep 2015 \sphinxhyphen{}\textgreater{} 20150902

\item {} 
\sphinxAtStartPar
{\hyperref[\detokenize{Python_module:lammps.lammps.get_thermo}]{\sphinxcrossref{\sphinxcode{\sphinxupquote{get\_thermo()}}}}}: return current value of a thermo keyword

\item {} 
\sphinxAtStartPar
{\hyperref[\detokenize{Python_module:lammps.lammps.last_thermo}]{\sphinxcrossref{\sphinxcode{\sphinxupquote{last\_thermo()}}}}}: return a dictionary of the last thermodynamic output

\item {} 
\sphinxAtStartPar
{\hyperref[\detokenize{Python_module:lammps.lammps.get_natoms}]{\sphinxcrossref{\sphinxcode{\sphinxupquote{get\_natoms()}}}}}: total \# of atoms as int

\item {} 
\sphinxAtStartPar
{\hyperref[\detokenize{Python_module:lammps.lammps.reset_box}]{\sphinxcrossref{\sphinxcode{\sphinxupquote{reset\_box()}}}}}: reset the simulation box size

\item {} 
\sphinxAtStartPar
{\hyperref[\detokenize{Python_module:lammps.lammps.extract_setting}]{\sphinxcrossref{\sphinxcode{\sphinxupquote{extract\_setting()}}}}}: return a global setting

\item {} 
\sphinxAtStartPar
{\hyperref[\detokenize{Python_module:lammps.lammps.extract_global}]{\sphinxcrossref{\sphinxcode{\sphinxupquote{extract\_global()}}}}}: extract a global quantity

\item {} 
\sphinxAtStartPar
{\hyperref[\detokenize{Python_module:lammps.lammps.extract_box}]{\sphinxcrossref{\sphinxcode{\sphinxupquote{extract\_box()}}}}}: extract box info

\item {} 
\sphinxAtStartPar
{\hyperref[\detokenize{Python_module:lammps.lammps.create_atoms}]{\sphinxcrossref{\sphinxcode{\sphinxupquote{create\_atoms()}}}}}: create N atoms with IDs, types, x, v, and image flags

\end{itemize}

\sphinxAtStartPar
\sphinxstylestrong{Properties}:
\begin{itemize}
\item {} 
\sphinxAtStartPar
{\hyperref[\detokenize{Python_module:lammps.lammps.last_thermo_step}]{\sphinxcrossref{\sphinxcode{\sphinxupquote{last\_thermo\_step}}}}}: the last timestep thermodynamic output was computed

\end{itemize}
\end{sphinxadmonition}

\begin{sphinxadmonition}{note}{PyLammps/IPyLammps API}

\sphinxAtStartPar
In addition to the functions provided by {\hyperref[\detokenize{Python_module:lammps.lammps}]{\sphinxcrossref{\sphinxcode{\sphinxupquote{lammps}}}}}, {\hyperref[\detokenize{Python_module:lammps.PyLammps}]{\sphinxcrossref{\sphinxcode{\sphinxupquote{PyLammps}}}}} objects
have several properties which allow you to query the system state:
\begin{description}
\sphinxlineitem{L.system}
\sphinxAtStartPar
Is a dictionary describing the system such as the bounding box or number of atoms

\sphinxlineitem{L.system.xlo, L.system.xhi}
\sphinxAtStartPar
bounding box limits along x\sphinxhyphen{}axis

\sphinxlineitem{L.system.ylo, L.system.yhi}
\sphinxAtStartPar
bounding box limits along y\sphinxhyphen{}axis

\sphinxlineitem{L.system.zlo, L.system.zhi}
\sphinxAtStartPar
bounding box limits along z\sphinxhyphen{}axis

\sphinxlineitem{L.communication}
\sphinxAtStartPar
configuration of communication subsystem, such as the number of threads or processors

\sphinxlineitem{L.communication.nthreads}
\sphinxAtStartPar
number of threads used by each LAMMPS process

\sphinxlineitem{L.communication.nprocs}
\sphinxAtStartPar
number of MPI processes used by LAMMPS

\sphinxlineitem{L.fixes}
\sphinxAtStartPar
List of fixes in the current system

\sphinxlineitem{L.computes}
\sphinxAtStartPar
List of active computes in the current system

\sphinxlineitem{L.dump}
\sphinxAtStartPar
List of active dumps in the current system

\sphinxlineitem{L.groups}
\sphinxAtStartPar
List of groups present in the current system

\end{description}

\sphinxAtStartPar
\sphinxstylestrong{Retrieving the value of an arbitrary LAMMPS expressions}

\sphinxAtStartPar
LAMMPS expressions can be immediately evaluated by using the \sphinxcode{\sphinxupquote{eval}} method. The
passed string parameter can be any expression containing global {\hyperref[\detokenize{thermo::doc}]{\sphinxcrossref{\DUrole{doc}{thermo command}}}} values,
variables, compute or fix data (see {\hyperref[\detokenize{Howto_output::doc}]{\sphinxcrossref{\DUrole{doc}{Output from LAMMPS (thermo, dumps, computes, fixes, variables)}}}}):

\begin{sphinxVerbatim}[commandchars=\\\{\}]
\PYG{n}{result} \PYG{o}{=} \PYG{n}{L}\PYG{o}{.}\PYG{n}{eval}\PYG{p}{(}\PYG{l+s+s2}{\PYGZdq{}}\PYG{l+s+s2}{ke}\PYG{l+s+s2}{\PYGZdq{}}\PYG{p}{)} \PYG{c+c1}{\PYGZsh{} kinetic energy}
\PYG{n}{result} \PYG{o}{=} \PYG{n}{L}\PYG{o}{.}\PYG{n}{eval}\PYG{p}{(}\PYG{l+s+s2}{\PYGZdq{}}\PYG{l+s+s2}{pe}\PYG{l+s+s2}{\PYGZdq{}}\PYG{p}{)} \PYG{c+c1}{\PYGZsh{} potential energy}

\PYG{n}{result} \PYG{o}{=} \PYG{n}{L}\PYG{o}{.}\PYG{n}{eval}\PYG{p}{(}\PYG{l+s+s2}{\PYGZdq{}}\PYG{l+s+s2}{v\PYGZus{}t/2.0}\PYG{l+s+s2}{\PYGZdq{}}\PYG{p}{)}
\end{sphinxVerbatim}

\sphinxAtStartPar
\sphinxstylestrong{Example}

\begin{sphinxVerbatim}[commandchars=\\\{\}]
\PYG{k+kn}{from}\PYG{+w}{ }\PYG{n+nn}{lammps}\PYG{+w}{ }\PYG{k+kn}{import} \PYG{n}{PyLammps}

\PYG{n}{L} \PYG{o}{=} \PYG{n}{PyLammps}\PYG{p}{(}\PYG{p}{)}
\PYG{n}{L}\PYG{o}{.}\PYG{n}{file}\PYG{p}{(}\PYG{l+s+s2}{\PYGZdq{}}\PYG{l+s+s2}{in.sysinit}\PYG{l+s+s2}{\PYGZdq{}}\PYG{p}{)}

\PYG{n+nb}{print}\PYG{p}{(}\PYG{l+s+sa}{f}\PYG{l+s+s2}{\PYGZdq{}}\PYG{l+s+s2}{running simulation with }\PYG{l+s+si}{\PYGZob{}}\PYG{n}{L}\PYG{o}{.}\PYG{n}{system}\PYG{o}{.}\PYG{n}{natoms}\PYG{l+s+si}{\PYGZcb{}}\PYG{l+s+s2}{ atoms}\PYG{l+s+s2}{\PYGZdq{}}\PYG{p}{)}

\PYG{n}{L}\PYG{o}{.}\PYG{n}{run}\PYG{p}{(}\PYG{l+m+mi}{1000}\PYG{p}{,} \PYG{l+s+s2}{\PYGZdq{}}\PYG{l+s+s2}{post no}\PYG{l+s+s2}{\PYGZdq{}}\PYG{p}{)}\PYG{p}{;}

\PYG{k}{for} \PYG{n}{i} \PYG{o+ow}{in} \PYG{n+nb}{range}\PYG{p}{(}\PYG{l+m+mi}{10}\PYG{p}{)}\PYG{p}{:}
   \PYG{n}{L}\PYG{o}{.}\PYG{n}{run}\PYG{p}{(}\PYG{l+m+mi}{100}\PYG{p}{,} \PYG{l+s+s2}{\PYGZdq{}}\PYG{l+s+s2}{pre no post no}\PYG{l+s+s2}{\PYGZdq{}}\PYG{p}{)}
   \PYG{n}{pe} \PYG{o}{=} \PYG{n}{L}\PYG{o}{.}\PYG{n}{eval}\PYG{p}{(}\PYG{l+s+s2}{\PYGZdq{}}\PYG{l+s+s2}{pe}\PYG{l+s+s2}{\PYGZdq{}}\PYG{p}{)}
   \PYG{n}{ke} \PYG{o}{=} \PYG{n}{L}\PYG{o}{.}\PYG{n}{eval}\PYG{p}{(}\PYG{l+s+s2}{\PYGZdq{}}\PYG{l+s+s2}{ke}\PYG{l+s+s2}{\PYGZdq{}}\PYG{p}{)}
   \PYG{n+nb}{print}\PYG{p}{(}\PYG{l+s+sa}{f}\PYG{l+s+s2}{\PYGZdq{}}\PYG{l+s+s2}{PE = }\PYG{l+s+si}{\PYGZob{}}\PYG{n}{pe}\PYG{l+s+si}{\PYGZcb{}}\PYG{l+s+se}{\PYGZbs{}n}\PYG{l+s+s2}{KE = }\PYG{l+s+si}{\PYGZob{}}\PYG{n}{ke}\PYG{l+s+si}{\PYGZcb{}}\PYG{l+s+s2}{\PYGZdq{}}\PYG{p}{)}
\end{sphinxVerbatim}
\end{sphinxadmonition}
\end{sphinxadmonition}

\sphinxstepscope


\subsection{Per\sphinxhyphen{}atom properties}
\label{\detokenize{Python_atoms:per-atom-properties}}\label{\detokenize{Python_atoms::doc}}
\sphinxAtStartPar
Similar to what is described in {\hyperref[\detokenize{Library_atoms::doc}]{\sphinxcrossref{\DUrole{doc}{Per\sphinxhyphen{}atom properties}}}}, the instances of
{\hyperref[\detokenize{Python_module:lammps.lammps}]{\sphinxcrossref{\sphinxcode{\sphinxupquote{lammps}}}}}, {\hyperref[\detokenize{Python_module:lammps.PyLammps}]{\sphinxcrossref{\sphinxcode{\sphinxupquote{PyLammps}}}}}, or
{\hyperref[\detokenize{Python_module:lammps.IPyLammps}]{\sphinxcrossref{\sphinxcode{\sphinxupquote{IPyLammps}}}}} can be used to extract atom quantities
and modify some of them. The main difference between the interfaces is how the information
is exposed.

\sphinxAtStartPar
While the {\hyperref[\detokenize{Python_module:lammps.lammps}]{\sphinxcrossref{\sphinxcode{\sphinxupquote{lammps}}}}} is just a thin layer that wraps C API calls,
{\hyperref[\detokenize{Python_module:lammps.PyLammps}]{\sphinxcrossref{\sphinxcode{\sphinxupquote{PyLammps}}}}} and {\hyperref[\detokenize{Python_module:lammps.IPyLammps}]{\sphinxcrossref{\sphinxcode{\sphinxupquote{IPyLammps}}}}} expose
information as objects and properties.

\sphinxAtStartPar
In some cases the data returned is a direct reference to the original data
inside LAMMPS cast to \sphinxcode{\sphinxupquote{ctypes}} pointers. Where possible, the wrappers will
determine the \sphinxcode{\sphinxupquote{ctypes}} data type and cast pointers accordingly. If
\sphinxcode{\sphinxupquote{numpy}} is installed arrays can also be extracted as numpy arrays, which
will access the C arrays directly and have the correct dimensions to protect
against invalid accesses.

\begin{sphinxadmonition}{warning}{Warning:}
\sphinxAtStartPar
When accessing per\sphinxhyphen{}atom data,
please note that this data is the per\sphinxhyphen{}processor local data and indexed
accordingly. These arrays can change sizes and order at every neighbor list
rebuild and atom sort event as atoms are migrating between subdomains.

\begin{sphinxadmonition}{note}{lammps API}

\begin{sphinxVerbatim}[commandchars=\\\{\}]
\PYG{k+kn}{from}\PYG{+w}{ }\PYG{n+nn}{lammps}\PYG{+w}{ }\PYG{k+kn}{import} \PYG{n}{lammps}

\PYG{n}{lmp} \PYG{o}{=} \PYG{n}{lammps}\PYG{p}{(}\PYG{p}{)}
\PYG{n}{lmp}\PYG{o}{.}\PYG{n}{file}\PYG{p}{(}\PYG{l+s+s2}{\PYGZdq{}}\PYG{l+s+s2}{in.sysinit}\PYG{l+s+s2}{\PYGZdq{}}\PYG{p}{)}

\PYG{n}{nlocal} \PYG{o}{=} \PYG{n}{lmp}\PYG{o}{.}\PYG{n}{extract\PYGZus{}global}\PYG{p}{(}\PYG{l+s+s2}{\PYGZdq{}}\PYG{l+s+s2}{nlocal}\PYG{l+s+s2}{\PYGZdq{}}\PYG{p}{)}
\PYG{n}{x} \PYG{o}{=} \PYG{n}{lmp}\PYG{o}{.}\PYG{n}{extract\PYGZus{}atom}\PYG{p}{(}\PYG{l+s+s2}{\PYGZdq{}}\PYG{l+s+s2}{x}\PYG{l+s+s2}{\PYGZdq{}}\PYG{p}{)}

\PYG{k}{for} \PYG{n}{i} \PYG{o+ow}{in} \PYG{n+nb}{range}\PYG{p}{(}\PYG{n}{nlocal}\PYG{p}{)}\PYG{p}{:}
   \PYG{n+nb}{print}\PYG{p}{(}\PYG{l+s+s2}{\PYGZdq{}}\PYG{l+s+s2}{(x,y,z) = (}\PYG{l+s+s2}{\PYGZdq{}}\PYG{p}{,} \PYG{n}{x}\PYG{p}{[}\PYG{n}{i}\PYG{p}{]}\PYG{p}{[}\PYG{l+m+mi}{0}\PYG{p}{]}\PYG{p}{,} \PYG{n}{x}\PYG{p}{[}\PYG{n}{i}\PYG{p}{]}\PYG{p}{[}\PYG{l+m+mi}{1}\PYG{p}{]}\PYG{p}{,} \PYG{n}{x}\PYG{p}{[}\PYG{n}{i}\PYG{p}{]}\PYG{p}{[}\PYG{l+m+mi}{2}\PYG{p}{]}\PYG{p}{,} \PYG{l+s+s2}{\PYGZdq{}}\PYG{l+s+s2}{)}\PYG{l+s+s2}{\PYGZdq{}}\PYG{p}{)}

\PYG{n}{lmp}\PYG{o}{.}\PYG{n}{close}\PYG{p}{(}\PYG{p}{)}
\end{sphinxVerbatim}

\sphinxAtStartPar
\sphinxstylestrong{Methods}:
\begin{itemize}
\item {} 
\sphinxAtStartPar
{\hyperref[\detokenize{Python_module:lammps.lammps.extract_atom}]{\sphinxcrossref{\sphinxcode{\sphinxupquote{extract\_atom()}}}}}: extract a per\sphinxhyphen{}atom quantity

\end{itemize}

\sphinxAtStartPar
\sphinxstylestrong{Numpy Methods}:
\begin{itemize}
\item {} 
\sphinxAtStartPar
{\hyperref[\detokenize{Python_module:lammps.numpy_wrapper.numpy_wrapper.extract_atom}]{\sphinxcrossref{\sphinxcode{\sphinxupquote{numpy.extract\_atom()}}}}}: extract a per\sphinxhyphen{}atom quantity as numpy array

\end{itemize}
\end{sphinxadmonition}

\begin{sphinxadmonition}{note}{PyLammps/IPyLammps API}

\sphinxAtStartPar
All atoms in the current simulation can be accessed by using the {\hyperref[\detokenize{Python_module:lammps.PyLammps.atoms}]{\sphinxcrossref{\sphinxcode{\sphinxupquote{PyLammps.atoms}}}}} property.
Each element of this list is a {\hyperref[\detokenize{Python_module:lammps.Atom}]{\sphinxcrossref{\sphinxcode{\sphinxupquote{Atom}}}}} or {\hyperref[\detokenize{Python_module:lammps.Atom2D}]{\sphinxcrossref{\sphinxcode{\sphinxupquote{Atom2D}}}}} object. The attributes of
these objects provide access to their data (id, type, position, velocity, force, etc.):

\begin{sphinxVerbatim}[commandchars=\\\{\}]
\PYG{c+c1}{\PYGZsh{} access first atom}
\PYG{n}{L}\PYG{o}{.}\PYG{n}{atoms}\PYG{p}{[}\PYG{l+m+mi}{0}\PYG{p}{]}\PYG{o}{.}\PYG{n}{id}
\PYG{n}{L}\PYG{o}{.}\PYG{n}{atoms}\PYG{p}{[}\PYG{l+m+mi}{0}\PYG{p}{]}\PYG{o}{.}\PYG{n}{type}

\PYG{c+c1}{\PYGZsh{} access second atom}
\PYG{n}{L}\PYG{o}{.}\PYG{n}{atoms}\PYG{p}{[}\PYG{l+m+mi}{1}\PYG{p}{]}\PYG{o}{.}\PYG{n}{position}
\PYG{n}{L}\PYG{o}{.}\PYG{n}{atoms}\PYG{p}{[}\PYG{l+m+mi}{1}\PYG{p}{]}\PYG{o}{.}\PYG{n}{velocity}
\PYG{n}{L}\PYG{o}{.}\PYG{n}{atoms}\PYG{p}{[}\PYG{l+m+mi}{1}\PYG{p}{]}\PYG{o}{.}\PYG{n}{force}
\end{sphinxVerbatim}

\sphinxAtStartPar
Some attributes can be changed:

\begin{sphinxVerbatim}[commandchars=\\\{\}]
\PYG{c+c1}{\PYGZsh{} set position in 2D simulation}
\PYG{n}{L}\PYG{o}{.}\PYG{n}{atoms}\PYG{p}{[}\PYG{l+m+mi}{0}\PYG{p}{]}\PYG{o}{.}\PYG{n}{position} \PYG{o}{=} \PYG{p}{(}\PYG{l+m+mf}{1.0}\PYG{p}{,} \PYG{l+m+mf}{0.0}\PYG{p}{)}

\PYG{c+c1}{\PYGZsh{} set position in 3D simulation}
\PYG{n}{L}\PYG{o}{.}\PYG{n}{atoms}\PYG{p}{[}\PYG{l+m+mi}{0}\PYG{p}{]}\PYG{o}{.}\PYG{n}{position} \PYG{o}{=} \PYG{p}{(}\PYG{l+m+mf}{1.0}\PYG{p}{,} \PYG{l+m+mf}{0.0}\PYG{p}{,} \PYG{l+m+mf}{1.0}\PYG{p}{)}
\end{sphinxVerbatim}
\end{sphinxadmonition}
\end{sphinxadmonition}

\sphinxstepscope


\subsection{Compute, fixes, variables}
\label{\detokenize{Python_objects:compute-fixes-variables}}\label{\detokenize{Python_objects::doc}}
\sphinxAtStartPar
This section documents accessing or modifying data from objects like
computes, fixes, or variables in LAMMPS using the {\hyperref[\detokenize{Python_module:module-lammps}]{\sphinxcrossref{\sphinxcode{\sphinxupquote{lammps}}}}} module.
\begin{quote}

\begin{sphinxadmonition}{note}{lammps API}

\sphinxAtStartPar
For {\hyperref[\detokenize{Python_module:lammps.lammps.extract_compute}]{\sphinxcrossref{\sphinxcode{\sphinxupquote{lammps.extract\_compute()}}}}} and
{\hyperref[\detokenize{Python_module:lammps.lammps.extract_fix}]{\sphinxcrossref{\sphinxcode{\sphinxupquote{lammps.extract\_fix()}}}}}, the global, per\sphinxhyphen{}atom,
or local data calculated by the compute or fix can be accessed. What is returned
depends on whether the compute or fix calculates a scalar or vector or array.
For a scalar, a single double value is returned.  If the compute or fix calculates
a vector or array, a pointer to the internal LAMMPS data is returned, which you can
use via normal Python subscripting.

\sphinxAtStartPar
The one exception is that for a fix that calculates a
global vector or array, a single double value from the vector or array
is returned, indexed by I (vector) or I and J (array).  I,J are
zero\sphinxhyphen{}based indices.
See the {\hyperref[\detokenize{Howto_output::doc}]{\sphinxcrossref{\DUrole{doc}{Howto output}}}} page for a discussion of
global, per\sphinxhyphen{}atom, and local data, and of scalar, vector, and array
data types.  See the doc pages for individual {\hyperref[\detokenize{compute::doc}]{\sphinxcrossref{\DUrole{doc}{computes}}}}
and {\hyperref[\detokenize{fix::doc}]{\sphinxcrossref{\DUrole{doc}{fixes}}}} for a description of what they calculate and
store.

\sphinxAtStartPar
For {\hyperref[\detokenize{Python_module:lammps.lammps.extract_variable}]{\sphinxcrossref{\sphinxcode{\sphinxupquote{lammps.extract\_variable()}}}}},
an {\hyperref[\detokenize{variable::doc}]{\sphinxcrossref{\DUrole{doc}{equal\sphinxhyphen{}style or atom\sphinxhyphen{}style variable}}}} is evaluated and
its result returned.

\sphinxAtStartPar
For equal\sphinxhyphen{}style variables a single \sphinxcode{\sphinxupquote{c\_double}} value is returned and the
group argument is ignored.  For atom\sphinxhyphen{}style variables, a vector of
\sphinxcode{\sphinxupquote{c\_double}} is returned, one value per atom, which you can use via normal
Python subscripting. The values will be zero for atoms not in the
specified group.

\sphinxAtStartPar
{\hyperref[\detokenize{Python_module:lammps.numpy_wrapper.numpy_wrapper.extract_compute}]{\sphinxcrossref{\sphinxcode{\sphinxupquote{lammps.numpy.extract\_compute()}}}}},
{\hyperref[\detokenize{Python_module:lammps.numpy_wrapper.numpy_wrapper.extract_fix}]{\sphinxcrossref{\sphinxcode{\sphinxupquote{lammps.numpy.extract\_fix()}}}}}, and
{\hyperref[\detokenize{Python_module:lammps.numpy_wrapper.numpy_wrapper.extract_variable}]{\sphinxcrossref{\sphinxcode{\sphinxupquote{lammps.numpy.extract\_variable()}}}}} are
equivalent NumPy implementations that return NumPy arrays instead of \sphinxcode{\sphinxupquote{ctypes}} pointers.

\sphinxAtStartPar
The {\hyperref[\detokenize{Python_module:lammps.lammps.set_variable}]{\sphinxcrossref{\sphinxcode{\sphinxupquote{lammps.set\_variable()}}}}} method sets an
existing string\sphinxhyphen{}style variable to a new string value, so that subsequent LAMMPS
commands can access the variable.

\sphinxAtStartPar
\sphinxstylestrong{Methods}:
\begin{itemize}
\item {} 
\sphinxAtStartPar
{\hyperref[\detokenize{Python_module:lammps.lammps.extract_compute}]{\sphinxcrossref{\sphinxcode{\sphinxupquote{lammps.extract\_compute()}}}}}: extract value(s) from a compute

\item {} 
\sphinxAtStartPar
{\hyperref[\detokenize{Python_module:lammps.lammps.extract_fix}]{\sphinxcrossref{\sphinxcode{\sphinxupquote{lammps.extract\_fix()}}}}}: extract value(s) from a fix

\item {} 
\sphinxAtStartPar
{\hyperref[\detokenize{Python_module:lammps.lammps.extract_variable}]{\sphinxcrossref{\sphinxcode{\sphinxupquote{lammps.extract\_variable()}}}}}: extract value(s) from a variable

\item {} 
\sphinxAtStartPar
{\hyperref[\detokenize{Python_module:lammps.lammps.set_variable}]{\sphinxcrossref{\sphinxcode{\sphinxupquote{lammps.set\_variable()}}}}}: set existing named string\sphinxhyphen{}style variable to value

\end{itemize}

\sphinxAtStartPar
\sphinxstylestrong{NumPy Methods}:
\begin{itemize}
\item {} 
\sphinxAtStartPar
{\hyperref[\detokenize{Python_module:lammps.numpy_wrapper.numpy_wrapper.extract_compute}]{\sphinxcrossref{\sphinxcode{\sphinxupquote{lammps.numpy.extract\_compute()}}}}}: extract value(s) from a compute, return arrays as numpy arrays

\item {} 
\sphinxAtStartPar
{\hyperref[\detokenize{Python_module:lammps.numpy_wrapper.numpy_wrapper.extract_fix}]{\sphinxcrossref{\sphinxcode{\sphinxupquote{lammps.numpy.extract\_fix()}}}}}: extract value(s) from a fix, return arrays as numpy arrays

\item {} 
\sphinxAtStartPar
{\hyperref[\detokenize{Python_module:lammps.numpy_wrapper.numpy_wrapper.extract_variable}]{\sphinxcrossref{\sphinxcode{\sphinxupquote{lammps.numpy.extract\_variable()}}}}}: extract value(s) from a variable, return arrays as numpy arrays

\end{itemize}
\end{sphinxadmonition}

\begin{sphinxadmonition}{note}{PyLammps/IPyLammps API}

\sphinxAtStartPar
PyLammps and IPyLammps classes currently do not add any additional ways of
retrieving information out of computes and fixes. This information can still be accessed by using the lammps API:

\begin{sphinxVerbatim}[commandchars=\\\{\}]
\PYG{n}{L}\PYG{o}{.}\PYG{n}{lmp}\PYG{o}{.}\PYG{n}{extract\PYGZus{}compute}\PYG{p}{(}\PYG{o}{.}\PYG{o}{.}\PYG{o}{.}\PYG{p}{)}
\PYG{n}{L}\PYG{o}{.}\PYG{n}{lmp}\PYG{o}{.}\PYG{n}{extract\PYGZus{}fix}\PYG{p}{(}\PYG{o}{.}\PYG{o}{.}\PYG{o}{.}\PYG{p}{)}
\PYG{c+c1}{\PYGZsh{} OR}
\PYG{n}{L}\PYG{o}{.}\PYG{n}{lmp}\PYG{o}{.}\PYG{n}{numpy}\PYG{o}{.}\PYG{n}{extract\PYGZus{}compute}\PYG{p}{(}\PYG{o}{.}\PYG{o}{.}\PYG{o}{.}\PYG{p}{)}
\PYG{n}{L}\PYG{o}{.}\PYG{n}{lmp}\PYG{o}{.}\PYG{n}{numpy}\PYG{o}{.}\PYG{n}{extract\PYGZus{}fix}\PYG{p}{(}\PYG{o}{.}\PYG{o}{.}\PYG{o}{.}\PYG{p}{)}
\end{sphinxVerbatim}

\sphinxAtStartPar
LAMMPS variables can be both defined and accessed via the {\hyperref[\detokenize{Python_module:lammps.PyLammps}]{\sphinxcrossref{\sphinxcode{\sphinxupquote{PyLammps}}}}} interface.

\sphinxAtStartPar
To define a variable you can use the {\hyperref[\detokenize{variable::doc}]{\sphinxcrossref{\DUrole{doc}{variable}}}} command:

\begin{sphinxVerbatim}[commandchars=\\\{\}]
\PYG{n}{L}\PYG{o}{.}\PYG{n}{variable}\PYG{p}{(}\PYG{l+s+s2}{\PYGZdq{}}\PYG{l+s+s2}{a index 2}\PYG{l+s+s2}{\PYGZdq{}}\PYG{p}{)}
\end{sphinxVerbatim}

\sphinxAtStartPar
A dictionary of all variables is returned by the {\hyperref[\detokenize{Python_module:lammps.PyLammps.variables}]{\sphinxcrossref{\sphinxcode{\sphinxupquote{PyLammps.variables}}}}} property:

\sphinxAtStartPar
you can access an individual variable by retrieving a variable object from the
\sphinxcode{\sphinxupquote{L.variables}} dictionary by name

\begin{sphinxVerbatim}[commandchars=\\\{\}]
\PYG{n}{a} \PYG{o}{=} \PYG{n}{L}\PYG{o}{.}\PYG{n}{variables}\PYG{p}{[}\PYG{l+s+s1}{\PYGZsq{}}\PYG{l+s+s1}{a}\PYG{l+s+s1}{\PYGZsq{}}\PYG{p}{]}
\end{sphinxVerbatim}

\sphinxAtStartPar
The variable value can then be easily read and written by accessing the value
property of this object.

\begin{sphinxVerbatim}[commandchars=\\\{\}]
\PYG{n+nb}{print}\PYG{p}{(}\PYG{n}{a}\PYG{o}{.}\PYG{n}{value}\PYG{p}{)}
\PYG{n}{a}\PYG{o}{.}\PYG{n}{value} \PYG{o}{=} \PYG{l+m+mi}{4}
\end{sphinxVerbatim}
\end{sphinxadmonition}
\end{quote}

\sphinxstepscope


\subsection{Scatter/gather operations}
\label{\detokenize{Python_scatter:scatter-gather-operations}}\label{\detokenize{Python_scatter::doc}}
\begin{sphinxVerbatim}[commandchars=\\\{\}]
\PYG{n}{data} \PYG{o}{=} \PYG{n}{lmp}\PYG{o}{.}\PYG{n}{gather\PYGZus{}atoms}\PYG{p}{(}\PYG{n}{name}\PYG{p}{,}\PYG{n+nb}{type}\PYG{p}{,}\PYG{n}{count}\PYG{p}{)}  \PYG{c+c1}{\PYGZsh{} return per\PYGZhy{}atom property of all atoms gathered into data, ordered by atom ID}
                                          \PYG{c+c1}{\PYGZsh{} name = \PYGZdq{}x\PYGZdq{}, \PYGZdq{}charge\PYGZdq{}, \PYGZdq{}type\PYGZdq{}, etc}
\PYG{n}{data} \PYG{o}{=} \PYG{n}{lmp}\PYG{o}{.}\PYG{n}{gather\PYGZus{}atoms\PYGZus{}concat}\PYG{p}{(}\PYG{n}{name}\PYG{p}{,}\PYG{n+nb}{type}\PYG{p}{,}\PYG{n}{count}\PYG{p}{)}  \PYG{c+c1}{\PYGZsh{} ditto, but concatenated atom values from each proc (unordered)}
\PYG{n}{data} \PYG{o}{=} \PYG{n}{lmp}\PYG{o}{.}\PYG{n}{gather\PYGZus{}atoms\PYGZus{}subset}\PYG{p}{(}\PYG{n}{name}\PYG{p}{,}\PYG{n+nb}{type}\PYG{p}{,}\PYG{n}{count}\PYG{p}{,}\PYG{n}{ndata}\PYG{p}{,}\PYG{n}{ids}\PYG{p}{)}  \PYG{c+c1}{\PYGZsh{} ditto, but for subset of Ndata atoms with IDs}

\PYG{n}{lmp}\PYG{o}{.}\PYG{n}{scatter\PYGZus{}atoms}\PYG{p}{(}\PYG{n}{name}\PYG{p}{,}\PYG{n+nb}{type}\PYG{p}{,}\PYG{n}{count}\PYG{p}{,}\PYG{n}{data}\PYG{p}{)}   \PYG{c+c1}{\PYGZsh{} scatter per\PYGZhy{}atom property to all atoms from data, ordered by atom ID}
                                          \PYG{c+c1}{\PYGZsh{} name = \PYGZdq{}x\PYGZdq{}, \PYGZdq{}charge\PYGZdq{}, \PYGZdq{}type\PYGZdq{}, etc}
                                          \PYG{c+c1}{\PYGZsh{} count = \PYGZsh{} of per\PYGZhy{}atom values, 1 or 3, etc}

\PYG{n}{lmp}\PYG{o}{.}\PYG{n}{scatter\PYGZus{}atoms\PYGZus{}subset}\PYG{p}{(}\PYG{n}{name}\PYG{p}{,}\PYG{n+nb}{type}\PYG{p}{,}\PYG{n}{count}\PYG{p}{,}\PYG{n}{ndata}\PYG{p}{,}\PYG{n}{ids}\PYG{p}{,}\PYG{n}{data}\PYG{p}{)}  \PYG{c+c1}{\PYGZsh{} ditto, but for subset of Ndata atoms with IDs}
\end{sphinxVerbatim}

\sphinxAtStartPar
The gather methods collect peratom info of the requested type (atom
coords, atom types, forces, etc) from all processors, and returns the
same vector of values to each calling processor.  The scatter
functions do the inverse.  They distribute a vector of peratom values,
passed by all calling processors, to individual atoms, which may be
owned by different processors.

\sphinxAtStartPar
Note that the data returned by the gather methods,
e.g. \sphinxcode{\sphinxupquote{gather\_atoms("x")}}, is
different from the data structure returned by
{\hyperref[\detokenize{Python_module:lammps.lammps.extract_atom}]{\sphinxcrossref{\sphinxcode{\sphinxupquote{extract\_atom("x")}}}}} in four ways.
(1) \sphinxcode{\sphinxupquote{gather\_atoms()}} returns a vector which you index as x{[}i{]};
\sphinxcode{\sphinxupquote{extract\_atom()}} returns an array which you index as x{[}i{]}{[}j{]}.
(2) \sphinxcode{\sphinxupquote{gather\_atoms()}} orders the atoms by atom ID while
\sphinxcode{\sphinxupquote{extract\_atom()}} does not.  (3) \sphinxcode{\sphinxupquote{gather\_atoms()}} returns
a list of all atoms in the simulation; \sphinxcode{\sphinxupquote{extract\_atoms()}} returns just
the atoms local to each processor.  (4) Finally, the \sphinxcode{\sphinxupquote{gather\_atoms()}}
data structure is a copy of the atom coords stored internally in
LAMMPS, whereas \sphinxcode{\sphinxupquote{extract\_atom()}} returns an array that effectively
points directly to the internal data.  This means you can change
values inside LAMMPS from Python by assigning a new values to the
\sphinxcode{\sphinxupquote{extract\_atom()}} array.  To do this with the \sphinxcode{\sphinxupquote{gather\_atoms()}} vector, you
need to change values in the vector, then invoke the
\sphinxcode{\sphinxupquote{scatter\_atoms("x")}} method.

\sphinxAtStartPar
For the scatter methods, the array of coordinates passed to must be a
ctypes vector of ints or doubles, allocated and initialized something
like this:

\begin{sphinxVerbatim}[commandchars=\\\{\}]
\PYG{k+kn}{from}\PYG{+w}{ }\PYG{n+nn}{ctypes}\PYG{+w}{ }\PYG{k+kn}{import} \PYG{n}{c\PYGZus{}double}
\PYG{n}{natoms} \PYG{o}{=} \PYG{n}{lmp}\PYG{o}{.}\PYG{n}{get\PYGZus{}natoms}\PYG{p}{(}\PYG{p}{)}
\PYG{n}{n3} \PYG{o}{=} \PYG{l+m+mi}{3}\PYG{o}{*}\PYG{n}{natoms}
\PYG{n}{x} \PYG{o}{=} \PYG{p}{(}\PYG{n}{n3}\PYG{o}{*}\PYG{n}{c\PYGZus{}double}\PYG{p}{)}\PYG{p}{(}\PYG{p}{)}
\PYG{n}{x}\PYG{p}{[}\PYG{l+m+mi}{0}\PYG{p}{]} \PYG{o}{=} \PYG{n}{x} \PYG{n}{coord} \PYG{n}{of} \PYG{n}{atom} \PYG{k}{with} \PYG{n}{ID} \PYG{l+m+mi}{1}
\PYG{n}{x}\PYG{p}{[}\PYG{l+m+mi}{1}\PYG{p}{]} \PYG{o}{=} \PYG{n}{y} \PYG{n}{coord} \PYG{n}{of} \PYG{n}{atom} \PYG{k}{with} \PYG{n}{ID} \PYG{l+m+mi}{1}
\PYG{n}{x}\PYG{p}{[}\PYG{l+m+mi}{2}\PYG{p}{]} \PYG{o}{=} \PYG{n}{z} \PYG{n}{coord} \PYG{n}{of} \PYG{n}{atom} \PYG{k}{with} \PYG{n}{ID} \PYG{l+m+mi}{1}
\PYG{n}{x}\PYG{p}{[}\PYG{l+m+mi}{3}\PYG{p}{]} \PYG{o}{=} \PYG{n}{x} \PYG{n}{coord} \PYG{n}{of} \PYG{n}{atom} \PYG{k}{with} \PYG{n}{ID} \PYG{l+m+mi}{2}
\PYG{o}{.}\PYG{o}{.}\PYG{o}{.}
\PYG{n}{x}\PYG{p}{[}\PYG{n}{n3}\PYG{o}{\PYGZhy{}}\PYG{l+m+mi}{1}\PYG{p}{]} \PYG{o}{=} \PYG{n}{z} \PYG{n}{coord} \PYG{n}{of} \PYG{n}{atom} \PYG{k}{with} \PYG{n}{ID} \PYG{n}{natoms}
\PYG{n}{lmp}\PYG{o}{.}\PYG{n}{scatter\PYGZus{}atoms}\PYG{p}{(}\PYG{l+s+s2}{\PYGZdq{}}\PYG{l+s+s2}{x}\PYG{l+s+s2}{\PYGZdq{}}\PYG{p}{,} \PYG{l+m+mi}{1}\PYG{p}{,} \PYG{l+m+mi}{3}\PYG{p}{,} \PYG{n}{x}\PYG{p}{)}
\end{sphinxVerbatim}

\sphinxAtStartPar
The coordinates can also be provided as arguments to the initializer of x:

\begin{sphinxVerbatim}[commandchars=\\\{\}]
\PYG{k+kn}{from}\PYG{+w}{ }\PYG{n+nn}{ctypes}\PYG{+w}{ }\PYG{k+kn}{import} \PYG{n}{c\PYGZus{}double}
\PYG{n}{natoms} \PYG{o}{=} \PYG{l+m+mi}{2}
\PYG{n}{n3} \PYG{o}{=} \PYG{l+m+mi}{3}\PYG{o}{*}\PYG{n}{natoms}
\PYG{c+c1}{\PYGZsh{} init in constructor}
\PYG{n}{x} \PYG{o}{=} \PYG{p}{(}\PYG{n}{n3}\PYG{o}{*}\PYG{n}{c\PYGZus{}double}\PYG{p}{)}\PYG{p}{(}\PYG{l+m+mf}{0.0}\PYG{p}{,} \PYG{l+m+mf}{0.0}\PYG{p}{,} \PYG{l+m+mf}{0.0}\PYG{p}{,} \PYG{l+m+mf}{1.0}\PYG{p}{,} \PYG{l+m+mf}{1.0}\PYG{p}{,} \PYG{l+m+mf}{1.0}\PYG{p}{)}
\PYG{n}{lmp}\PYG{o}{.}\PYG{n}{scatter\PYGZus{}atoms}\PYG{p}{(}\PYG{l+s+s2}{\PYGZdq{}}\PYG{l+s+s2}{x}\PYG{l+s+s2}{\PYGZdq{}}\PYG{p}{,} \PYG{l+m+mi}{1}\PYG{p}{,} \PYG{l+m+mi}{3}\PYG{p}{,} \PYG{n}{x}\PYG{p}{)}
\PYG{c+c1}{\PYGZsh{} or using a list}
\PYG{n}{coords} \PYG{o}{=} \PYG{p}{[}\PYG{l+m+mf}{1.0}\PYG{p}{,} \PYG{l+m+mf}{2.0}\PYG{p}{,} \PYG{l+m+mf}{3.0}\PYG{p}{,} \PYG{o}{\PYGZhy{}}\PYG{l+m+mf}{3.0}\PYG{p}{,} \PYG{o}{\PYGZhy{}}\PYG{l+m+mf}{2.0}\PYG{p}{,} \PYG{o}{\PYGZhy{}}\PYG{l+m+mf}{1.0}\PYG{p}{]}
\PYG{n}{x} \PYG{o}{=} \PYG{p}{(}\PYG{n}{c\PYGZus{}double}\PYG{o}{*}\PYG{n+nb}{len}\PYG{p}{(}\PYG{n}{coords}\PYG{p}{)}\PYG{p}{)}\PYG{p}{(}\PYG{o}{*}\PYG{n}{coords}\PYG{p}{)}
\end{sphinxVerbatim}

\sphinxAtStartPar
Alternatively, you can just change values in the vector returned by
the gather methods, since they are also ctypes vectors.

\sphinxstepscope


\subsection{Neighbor list access}
\label{\detokenize{Python_neighbor:neighbor-list-access}}\label{\detokenize{Python_neighbor::doc}}
\sphinxAtStartPar
Access to neighbor lists is handled through a couple of wrapper classes
that allows one to treat it like either a python list or a NumPy array.  The
access procedure is similar to that of the C\sphinxhyphen{}library interface: use one
of the “find” functions to look up the index of the neighbor list in the
global table of neighbor lists and then get access to the neighbor list
data. The code sample below demonstrates reading the neighbor list data
using the NumPy access method.

\begin{sphinxVerbatim}[commandchars=\\\{\}]
\PYG{k+kn}{from}\PYG{+w}{ }\PYG{n+nn}{lammps}\PYG{+w}{ }\PYG{k+kn}{import} \PYG{n}{lammps}
\PYG{k+kn}{import}\PYG{+w}{ }\PYG{n+nn}{numpy}\PYG{+w}{ }\PYG{k}{as}\PYG{+w}{ }\PYG{n+nn}{np}

\PYG{n}{lmp} \PYG{o}{=} \PYG{n}{lammps}\PYG{p}{(}\PYG{p}{)}
\PYG{n}{lmp}\PYG{o}{.}\PYG{n}{commands\PYGZus{}string}\PYG{p}{(}\PYG{l+s+s2}{\PYGZdq{}\PYGZdq{}\PYGZdq{}}
\PYG{l+s+s2}{region box block \PYGZhy{}2 2 \PYGZhy{}2 2 \PYGZhy{}2 2}
\PYG{l+s+s2}{lattice fcc 1.0}
\PYG{l+s+s2}{create\PYGZus{}box 1 box}
\PYG{l+s+s2}{create\PYGZus{}atoms 1 box}
\PYG{l+s+s2}{mass 1 1.0}
\PYG{l+s+s2}{pair\PYGZus{}style lj/cut 2.5}
\PYG{l+s+s2}{pair\PYGZus{}coeff 1 1 1.0 1.0}
\PYG{l+s+s2}{run 0 post no}\PYG{l+s+s2}{\PYGZdq{}\PYGZdq{}\PYGZdq{}}\PYG{p}{)}

\PYG{c+c1}{\PYGZsh{} look up the neighbor list}
\PYG{n}{nlidx} \PYG{o}{=} \PYG{n}{lmp}\PYG{o}{.}\PYG{n}{find\PYGZus{}pair\PYGZus{}neighlist}\PYG{p}{(}\PYG{l+s+s1}{\PYGZsq{}}\PYG{l+s+s1}{lj/cut}\PYG{l+s+s1}{\PYGZsq{}}\PYG{p}{)}
\PYG{n}{nl} \PYG{o}{=} \PYG{n}{lmp}\PYG{o}{.}\PYG{n}{numpy}\PYG{o}{.}\PYG{n}{get\PYGZus{}neighlist}\PYG{p}{(}\PYG{n}{nlidx}\PYG{p}{)}
\PYG{n}{tags} \PYG{o}{=} \PYG{n}{lmp}\PYG{o}{.}\PYG{n}{extract\PYGZus{}atom}\PYG{p}{(}\PYG{l+s+s1}{\PYGZsq{}}\PYG{l+s+s1}{id}\PYG{l+s+s1}{\PYGZsq{}}\PYG{p}{)}
\PYG{n+nb}{print}\PYG{p}{(}\PYG{l+s+s2}{\PYGZdq{}}\PYG{l+s+s2}{half neighbor list with }\PYG{l+s+si}{\PYGZob{}\PYGZcb{}}\PYG{l+s+s2}{ entries}\PYG{l+s+s2}{\PYGZdq{}}\PYG{o}{.}\PYG{n}{format}\PYG{p}{(}\PYG{n}{nl}\PYG{o}{.}\PYG{n}{size}\PYG{p}{)}\PYG{p}{)}
\PYG{c+c1}{\PYGZsh{} print neighbor list contents}
\PYG{k}{for} \PYG{n}{i} \PYG{o+ow}{in} \PYG{n+nb}{range}\PYG{p}{(}\PYG{l+m+mi}{0}\PYG{p}{,}\PYG{n}{nl}\PYG{o}{.}\PYG{n}{size}\PYG{p}{)}\PYG{p}{:}
    \PYG{n}{idx}\PYG{p}{,} \PYG{n}{nlist}  \PYG{o}{=} \PYG{n}{nl}\PYG{o}{.}\PYG{n}{get}\PYG{p}{(}\PYG{n}{i}\PYG{p}{)}
    \PYG{n+nb}{print}\PYG{p}{(}\PYG{l+s+s2}{\PYGZdq{}}\PYG{l+s+se}{\PYGZbs{}n}\PYG{l+s+s2}{atom }\PYG{l+s+si}{\PYGZob{}\PYGZcb{}}\PYG{l+s+s2}{ with ID }\PYG{l+s+si}{\PYGZob{}\PYGZcb{}}\PYG{l+s+s2}{ has }\PYG{l+s+si}{\PYGZob{}\PYGZcb{}}\PYG{l+s+s2}{ neighbors:}\PYG{l+s+s2}{\PYGZdq{}}\PYG{o}{.}\PYG{n}{format}\PYG{p}{(}\PYG{n}{idx}\PYG{p}{,}\PYG{n}{tags}\PYG{p}{[}\PYG{n}{idx}\PYG{p}{]}\PYG{p}{,}\PYG{n}{nlist}\PYG{o}{.}\PYG{n}{size}\PYG{p}{)}\PYG{p}{)}
    \PYG{k}{if} \PYG{n}{nlist}\PYG{o}{.}\PYG{n}{size} \PYG{o}{\PYGZgt{}} \PYG{l+m+mi}{0}\PYG{p}{:}
        \PYG{k}{for} \PYG{n}{n} \PYG{o+ow}{in} \PYG{n}{np}\PYG{o}{.}\PYG{n}{nditer}\PYG{p}{(}\PYG{n}{nlist}\PYG{p}{)}\PYG{p}{:}
            \PYG{n+nb}{print}\PYG{p}{(}\PYG{l+s+s2}{\PYGZdq{}}\PYG{l+s+s2}{  atom }\PYG{l+s+si}{\PYGZob{}\PYGZcb{}}\PYG{l+s+s2}{ with ID }\PYG{l+s+si}{\PYGZob{}\PYGZcb{}}\PYG{l+s+s2}{\PYGZdq{}}\PYG{o}{.}\PYG{n}{format}\PYG{p}{(}\PYG{n}{n}\PYG{p}{,}\PYG{n}{tags}\PYG{p}{[}\PYG{n}{n}\PYG{p}{]}\PYG{p}{)}\PYG{p}{)}
\end{sphinxVerbatim}

\sphinxAtStartPar
Another example for extracting a full neighbor list without evaluating a
potential is shown below.

\begin{sphinxVerbatim}[commandchars=\\\{\}]
\PYG{k+kn}{from}\PYG{+w}{ }\PYG{n+nn}{lammps}\PYG{+w}{ }\PYG{k+kn}{import} \PYG{n}{lammps}
\PYG{k+kn}{import}\PYG{+w}{ }\PYG{n+nn}{numpy}\PYG{+w}{ }\PYG{k}{as}\PYG{+w}{ }\PYG{n+nn}{np}

\PYG{n}{lmp} \PYG{o}{=} \PYG{n}{lammps}\PYG{p}{(}\PYG{p}{)}
\PYG{n}{lmp}\PYG{o}{.}\PYG{n}{commands\PYGZus{}string}\PYG{p}{(}\PYG{l+s+s2}{\PYGZdq{}\PYGZdq{}\PYGZdq{}}
\PYG{l+s+s2}{newton off}
\PYG{l+s+s2}{region box block \PYGZhy{}2 2 \PYGZhy{}2 2 \PYGZhy{}2 2}
\PYG{l+s+s2}{lattice fcc 1.0}
\PYG{l+s+s2}{create\PYGZus{}box 1 box}
\PYG{l+s+s2}{create\PYGZus{}atoms 1 box}
\PYG{l+s+s2}{mass 1 1.0}
\PYG{l+s+s2}{pair\PYGZus{}style zero 1.0 full}
\PYG{l+s+s2}{pair\PYGZus{}coeff * *}
\PYG{l+s+s2}{run 0 post no}\PYG{l+s+s2}{\PYGZdq{}\PYGZdq{}\PYGZdq{}}\PYG{p}{)}

\PYG{c+c1}{\PYGZsh{} look up the neighbor list}
\PYG{n}{nlidx} \PYG{o}{=} \PYG{n}{lmp}\PYG{o}{.}\PYG{n}{find\PYGZus{}pair\PYGZus{}neighlist}\PYG{p}{(}\PYG{l+s+s1}{\PYGZsq{}}\PYG{l+s+s1}{zero}\PYG{l+s+s1}{\PYGZsq{}}\PYG{p}{)}
\PYG{n}{nl} \PYG{o}{=} \PYG{n}{lmp}\PYG{o}{.}\PYG{n}{numpy}\PYG{o}{.}\PYG{n}{get\PYGZus{}neighlist}\PYG{p}{(}\PYG{n}{nlidx}\PYG{p}{)}
\PYG{n}{tags} \PYG{o}{=} \PYG{n}{lmp}\PYG{o}{.}\PYG{n}{extract\PYGZus{}atom}\PYG{p}{(}\PYG{l+s+s1}{\PYGZsq{}}\PYG{l+s+s1}{id}\PYG{l+s+s1}{\PYGZsq{}}\PYG{p}{)}
\PYG{n+nb}{print}\PYG{p}{(}\PYG{l+s+s2}{\PYGZdq{}}\PYG{l+s+s2}{full neighbor list with }\PYG{l+s+si}{\PYGZob{}\PYGZcb{}}\PYG{l+s+s2}{ entries}\PYG{l+s+s2}{\PYGZdq{}}\PYG{o}{.}\PYG{n}{format}\PYG{p}{(}\PYG{n}{nl}\PYG{o}{.}\PYG{n}{size}\PYG{p}{)}\PYG{p}{)}
\PYG{c+c1}{\PYGZsh{} print neighbor list contents}
\PYG{k}{for} \PYG{n}{i} \PYG{o+ow}{in} \PYG{n+nb}{range}\PYG{p}{(}\PYG{l+m+mi}{0}\PYG{p}{,}\PYG{n}{nl}\PYG{o}{.}\PYG{n}{size}\PYG{p}{)}\PYG{p}{:}
    \PYG{n}{idx}\PYG{p}{,} \PYG{n}{nlist}  \PYG{o}{=} \PYG{n}{nl}\PYG{o}{.}\PYG{n}{get}\PYG{p}{(}\PYG{n}{i}\PYG{p}{)}
    \PYG{n+nb}{print}\PYG{p}{(}\PYG{l+s+s2}{\PYGZdq{}}\PYG{l+s+se}{\PYGZbs{}n}\PYG{l+s+s2}{atom }\PYG{l+s+si}{\PYGZob{}\PYGZcb{}}\PYG{l+s+s2}{ with ID }\PYG{l+s+si}{\PYGZob{}\PYGZcb{}}\PYG{l+s+s2}{ has }\PYG{l+s+si}{\PYGZob{}\PYGZcb{}}\PYG{l+s+s2}{ neighbors:}\PYG{l+s+s2}{\PYGZdq{}}\PYG{o}{.}\PYG{n}{format}\PYG{p}{(}\PYG{n}{idx}\PYG{p}{,}\PYG{n}{tags}\PYG{p}{[}\PYG{n}{idx}\PYG{p}{]}\PYG{p}{,}\PYG{n}{nlist}\PYG{o}{.}\PYG{n}{size}\PYG{p}{)}\PYG{p}{)}
    \PYG{k}{if} \PYG{n}{nlist}\PYG{o}{.}\PYG{n}{size} \PYG{o}{\PYGZgt{}} \PYG{l+m+mi}{0}\PYG{p}{:}
        \PYG{k}{for} \PYG{n}{n} \PYG{o+ow}{in} \PYG{n}{np}\PYG{o}{.}\PYG{n}{nditer}\PYG{p}{(}\PYG{n}{nlist}\PYG{p}{)}\PYG{p}{:}
            \PYG{k}{pass}
            \PYG{n+nb}{print}\PYG{p}{(}\PYG{l+s+s2}{\PYGZdq{}}\PYG{l+s+s2}{  atom }\PYG{l+s+si}{\PYGZob{}\PYGZcb{}}\PYG{l+s+s2}{ with ID }\PYG{l+s+si}{\PYGZob{}\PYGZcb{}}\PYG{l+s+s2}{\PYGZdq{}}\PYG{o}{.}\PYG{n}{format}\PYG{p}{(}\PYG{n}{n}\PYG{p}{,}\PYG{n}{tags}\PYG{p}{[}\PYG{n}{n}\PYG{p}{]}\PYG{p}{)}\PYG{p}{)}
\end{sphinxVerbatim}

\sphinxAtStartPar
\sphinxstylestrong{Methods:}
\begin{itemize}
\item {} 
\sphinxAtStartPar
{\hyperref[\detokenize{Python_module:lammps.lammps.get_neighlist}]{\sphinxcrossref{\sphinxcode{\sphinxupquote{lammps.get\_neighlist()}}}}}: Get neighbor list for given index

\item {} 
\sphinxAtStartPar
\sphinxcode{\sphinxupquote{lammps.get\_neighlist\_size()}}: Get number of elements in neighbor list

\item {} 
\sphinxAtStartPar
\sphinxcode{\sphinxupquote{lammps.get\_neighlist\_element\_neighbors()}}: Get element in neighbor list and its neighbors

\item {} 
\sphinxAtStartPar
{\hyperref[\detokenize{Python_module:lammps.lammps.find_pair_neighlist}]{\sphinxcrossref{\sphinxcode{\sphinxupquote{lammps.find\_pair\_neighlist()}}}}}: Find neighbor list of pair style

\item {} 
\sphinxAtStartPar
{\hyperref[\detokenize{Python_module:lammps.lammps.find_pair_neighlist}]{\sphinxcrossref{\sphinxcode{\sphinxupquote{lammps.find\_fix\_neighlist()}}}}}: Find neighbor list of fix style

\item {} 
\sphinxAtStartPar
{\hyperref[\detokenize{Python_module:lammps.lammps.find_pair_neighlist}]{\sphinxcrossref{\sphinxcode{\sphinxupquote{lammps.find\_compute\_neighlist()}}}}}: Find neighbor list of compute style

\end{itemize}

\sphinxAtStartPar
\sphinxstylestrong{NumPy Methods:}
\begin{itemize}
\item {} 
\sphinxAtStartPar
{\hyperref[\detokenize{Python_module:lammps.numpy_wrapper.numpy_wrapper.get_neighlist}]{\sphinxcrossref{\sphinxcode{\sphinxupquote{lammps.numpy.get\_neighlist()}}}}}: Get neighbor list for given index, which uses NumPy arrays for its element neighbor arrays

\item {} 
\sphinxAtStartPar
{\hyperref[\detokenize{Python_module:lammps.numpy_wrapper.numpy_wrapper.get_neighlist_element_neighbors}]{\sphinxcrossref{\sphinxcode{\sphinxupquote{lammps.numpy.get\_neighlist\_element\_neighbors()}}}}}: Get element in neighbor list and its neighbors (as a numpy array)

\end{itemize}

\sphinxstepscope


\subsection{Configuration information}
\label{\detokenize{Python_config:configuration-information}}\label{\detokenize{Python_config::doc}}
\sphinxAtStartPar
The following methods can be used to query the LAMMPS library
about compile time settings and included packages and styles.
\sphinxSetupCaptionForVerbatim{Example for using configuration settings functions}
\def\sphinxLiteralBlockLabel{\label{\detokenize{Python_config:id1}}}
\begin{sphinxVerbatim}[commandchars=\\\{\}]
\PYG{k+kn}{from}\PYG{+w}{ }\PYG{n+nn}{lammps}\PYG{+w}{ }\PYG{k+kn}{import} \PYG{n}{lammps}

\PYG{n}{lmp} \PYG{o}{=} \PYG{n}{lammps}\PYG{p}{(}\PYG{p}{)}

\PYG{k}{try}\PYG{p}{:}
    \PYG{n}{lmp}\PYG{o}{.}\PYG{n}{file}\PYG{p}{(}\PYG{l+s+s2}{\PYGZdq{}}\PYG{l+s+s2}{in.missing}\PYG{l+s+s2}{\PYGZdq{}}\PYG{p}{)}
\PYG{k}{except} \PYG{n+ne}{Exception} \PYG{k}{as} \PYG{n}{e}\PYG{p}{:}
    \PYG{n+nb}{print}\PYG{p}{(}\PYG{l+s+s2}{\PYGZdq{}}\PYG{l+s+s2}{LAMMPS failed with error:}\PYG{l+s+s2}{\PYGZdq{}}\PYG{p}{,} \PYG{n}{e}\PYG{p}{)}

\PYG{c+c1}{\PYGZsh{} write compressed dump file depending on available of options}

\PYG{k}{if} \PYG{n}{lmp}\PYG{o}{.}\PYG{n}{has\PYGZus{}style}\PYG{p}{(}\PYG{l+s+s2}{\PYGZdq{}}\PYG{l+s+s2}{dump}\PYG{l+s+s2}{\PYGZdq{}}\PYG{p}{,} \PYG{l+s+s2}{\PYGZdq{}}\PYG{l+s+s2}{atom/zstd}\PYG{l+s+s2}{\PYGZdq{}}\PYG{p}{)}\PYG{p}{:}
    \PYG{n}{lmp}\PYG{o}{.}\PYG{n}{command}\PYG{p}{(}\PYG{l+s+s2}{\PYGZdq{}}\PYG{l+s+s2}{dump d1 all atom/zstd 100 dump.zst}\PYG{l+s+s2}{\PYGZdq{}}\PYG{p}{)}
\PYG{k}{elif} \PYG{n}{lmp}\PYG{o}{.}\PYG{n}{has\PYGZus{}style}\PYG{p}{(}\PYG{l+s+s2}{\PYGZdq{}}\PYG{l+s+s2}{dump}\PYG{l+s+s2}{\PYGZdq{}}\PYG{p}{,} \PYG{l+s+s2}{\PYGZdq{}}\PYG{l+s+s2}{atom/gz}\PYG{l+s+s2}{\PYGZdq{}}\PYG{p}{)}\PYG{p}{:}
    \PYG{n}{lmp}\PYG{o}{.}\PYG{n}{command}\PYG{p}{(}\PYG{l+s+s2}{\PYGZdq{}}\PYG{l+s+s2}{dump d1 all atom/gz 100 dump.gz}\PYG{l+s+s2}{\PYGZdq{}}\PYG{p}{)}
\PYG{k}{elif} \PYG{n}{lmp}\PYG{o}{.}\PYG{n}{has\PYGZus{}gzip\PYGZus{}support}\PYG{p}{(}\PYG{p}{)}\PYG{p}{:}
    \PYG{n}{lmp}\PYG{o}{.}\PYG{n}{command}\PYG{p}{(}\PYG{l+s+s2}{\PYGZdq{}}\PYG{l+s+s2}{dump d1 all atom 100 dump.gz}\PYG{l+s+s2}{\PYGZdq{}}\PYG{p}{)}
\PYG{k}{else}\PYG{p}{:}
    \PYG{n}{lmp}\PYG{o}{.}\PYG{n}{command}\PYG{p}{(}\PYG{l+s+s2}{\PYGZdq{}}\PYG{l+s+s2}{dump d1 all atom 100 dump}\PYG{l+s+s2}{\PYGZdq{}}\PYG{p}{)}
\end{sphinxVerbatim}


\bigskip\hrule\bigskip


\sphinxAtStartPar
\sphinxstylestrong{Methods:}
\begin{itemize}
\item {} 
\sphinxAtStartPar
{\hyperref[\detokenize{Python_module:lammps.lammps.has_mpi_support}]{\sphinxcrossref{\sphinxcode{\sphinxupquote{lammps.has\_mpi\_support}}}}}

\item {} 
\sphinxAtStartPar
{\hyperref[\detokenize{Python_module:lammps.lammps.has_exceptions}]{\sphinxcrossref{\sphinxcode{\sphinxupquote{lammps.has\_exceptions}}}}}

\item {} 
\sphinxAtStartPar
{\hyperref[\detokenize{Python_module:lammps.lammps.has_gzip_support}]{\sphinxcrossref{\sphinxcode{\sphinxupquote{lammps.has\_gzip\_support}}}}}

\item {} 
\sphinxAtStartPar
{\hyperref[\detokenize{Python_module:lammps.lammps.has_png_support}]{\sphinxcrossref{\sphinxcode{\sphinxupquote{lammps.has\_png\_support}}}}}

\item {} 
\sphinxAtStartPar
{\hyperref[\detokenize{Python_module:lammps.lammps.has_jpeg_support}]{\sphinxcrossref{\sphinxcode{\sphinxupquote{lammps.has\_jpeg\_support}}}}}

\item {} 
\sphinxAtStartPar
{\hyperref[\detokenize{Python_module:lammps.lammps.has_ffmpeg_support}]{\sphinxcrossref{\sphinxcode{\sphinxupquote{lammps.has\_ffmpeg\_support}}}}}

\item {} 
\sphinxAtStartPar
{\hyperref[\detokenize{Python_module:lammps.lammps.installed_packages}]{\sphinxcrossref{\sphinxcode{\sphinxupquote{lammps.installed\_packages}}}}}

\item {} 
\sphinxAtStartPar
{\hyperref[\detokenize{Python_module:lammps.lammps.accelerator_config}]{\sphinxcrossref{\sphinxcode{\sphinxupquote{lammps.get\_accelerator\_config}}}}}

\item {} 
\sphinxAtStartPar
{\hyperref[\detokenize{Python_module:lammps.lammps.has_style}]{\sphinxcrossref{\sphinxcode{\sphinxupquote{lammps.has\_style()}}}}}

\item {} 
\sphinxAtStartPar
{\hyperref[\detokenize{Python_module:lammps.lammps.available_styles}]{\sphinxcrossref{\sphinxcode{\sphinxupquote{lammps.available\_styles()}}}}}

\end{itemize}

\sphinxstepscope


\section{The \sphinxstyleliteralintitle{\sphinxupquote{lammps}} Python module}
\label{\detokenize{Python_module:module-lammps}}\label{\detokenize{Python_module:the-lammps-python-module}}\label{\detokenize{Python_module::doc}}\index{module@\spxentry{module}!lammps@\spxentry{lammps}}\index{lammps@\spxentry{lammps}!module@\spxentry{module}}
\sphinxAtStartPar
The LAMMPS Python interface is implemented as a module called {\hyperref[\detokenize{Python_module:module-lammps}]{\sphinxcrossref{\sphinxcode{\sphinxupquote{lammps}}}}}
which is defined in the \sphinxcode{\sphinxupquote{lammps}} package in the \sphinxcode{\sphinxupquote{python}} folder of the
LAMMPS source code distribution.  After compilation of LAMMPS, the module can
be installed into a Python system folder or a user folder with \sphinxcode{\sphinxupquote{make
install\sphinxhyphen{}python}}.  Components of the module can then loaded into a Python
session with the \sphinxcode{\sphinxupquote{import}} command.

\sphinxAtStartPar
There are multiple Python interface classes in the {\hyperref[\detokenize{Python_module:module-lammps}]{\sphinxcrossref{\sphinxcode{\sphinxupquote{lammps}}}}} module:
\begin{itemize}
\item {} 
\sphinxAtStartPar
the {\hyperref[\detokenize{Python_module:lammps.lammps}]{\sphinxcrossref{\sphinxcode{\sphinxupquote{lammps}}}}} class. This is a wrapper around
the C\sphinxhyphen{}library interface and its member functions try to replicate the
{\hyperref[\detokenize{Library:lammps-c-api}]{\sphinxcrossref{\DUrole{std,std-ref}{C\sphinxhyphen{}library API}}}} closely.  This is the most
feature\sphinxhyphen{}complete Python API.

\item {} 
\sphinxAtStartPar
the {\hyperref[\detokenize{Python_module:lammps.PyLammps}]{\sphinxcrossref{\sphinxcode{\sphinxupquote{PyLammps}}}}} class. This is a more high\sphinxhyphen{}level
and more Python style class implemented on top of the
{\hyperref[\detokenize{Python_module:lammps.lammps}]{\sphinxcrossref{\sphinxcode{\sphinxupquote{lammps}}}}} class.

\item {} 
\sphinxAtStartPar
the {\hyperref[\detokenize{Python_module:lammps.IPyLammps}]{\sphinxcrossref{\sphinxcode{\sphinxupquote{IPyLammps}}}}} class is derived from
{\hyperref[\detokenize{Python_module:lammps.PyLammps}]{\sphinxcrossref{\sphinxcode{\sphinxupquote{PyLammps}}}}} and adds embedded graphics
features to conveniently include LAMMPS into \sphinxhref{https://jupyter.org/}{Jupyter} notebooks.

\end{itemize}

\begin{sphinxadmonition}{note}{Version check}

\sphinxAtStartPar
The {\hyperref[\detokenize{Python_module:module-lammps}]{\sphinxcrossref{\sphinxcode{\sphinxupquote{lammps}}}}} module stores the version number of the LAMMPS
version it is installed from.  When initializing the
{\hyperref[\detokenize{Python_module:lammps.lammps}]{\sphinxcrossref{\sphinxcode{\sphinxupquote{lammps}}}}} class, this version is checked to
be the same as the result from {\hyperref[\detokenize{Python_module:lammps.lammps.version}]{\sphinxcrossref{\sphinxcode{\sphinxupquote{lammps.version()}}}}}, the version
of the LAMMPS shared library that the module interfaces to.  If the
they are not the same an AttributeError exception is raised since a
mismatch of versions (e.g.  due to incorrect use of the
\sphinxcode{\sphinxupquote{LD\_LIBRARY\_PATH}} or \sphinxcode{\sphinxupquote{PYTHONPATH}} environment variables can lead
to crashes or data corruption and otherwise incorrect behavior.
\end{sphinxadmonition}

\sphinxAtStartPar
LAMMPS module global members:
\index{\_\_version\_\_ (in module lammps)@\spxentry{\_\_version\_\_}\spxextra{in module lammps}}

\begin{fulllineitems}
\phantomsection\label{\detokenize{Python_module:lammps.__version__}}
\pysigstartsignatures
\pysigline{\sphinxcode{\sphinxupquote{lammps.}}\sphinxbfcode{\sphinxupquote{\_\_version\_\_}}}
\pysigstopsignatures
\sphinxAtStartPar
Numerical representation of the LAMMPS version this
module was taken from.  Has the same format as the
result of {\hyperref[\detokenize{Python_module:lammps.lammps.version}]{\sphinxcrossref{\sphinxcode{\sphinxupquote{lammps.version()}}}}}.

\end{fulllineitems}



\bigskip\hrule\bigskip



\subsection{The \sphinxstyleliteralintitle{\sphinxupquote{lammps}} class API}
\label{\detokenize{Python_module:the-lammps-class-api}}
\sphinxAtStartPar
The {\hyperref[\detokenize{Python_module:lammps.lammps}]{\sphinxcrossref{\sphinxcode{\sphinxupquote{lammps}}}}} class is the core of the LAMMPS
Python interfaces.  It is a wrapper around the {\hyperref[\detokenize{Library:lammps-c-api}]{\sphinxcrossref{\DUrole{std,std-ref}{LAMMPS C library
API}}}} using the \sphinxhref{https://docs.python.org/3/library/ctypes.html}{Python ctypes module} and a shared library
compiled from the LAMMPS sources code.  The individual methods in this
class try to closely follow the corresponding C functions.  The handle
argument that needs to be passed to the C functions is stored internally
in the class and automatically added when calling the C library
functions. Below is a detailed documentation of the API.
\index{lammps (class in lammps)@\spxentry{lammps}\spxextra{class in lammps}}

\begin{fulllineitems}
\phantomsection\label{\detokenize{Python_module:lammps.lammps}}
\pysigstartsignatures
\pysiglinewithargsret{\sphinxbfcode{\sphinxupquote{class\DUrole{w}{ }}}\sphinxcode{\sphinxupquote{lammps.}}\sphinxbfcode{\sphinxupquote{lammps}}}{\sphinxparam{\DUrole{n}{name}\DUrole{o}{=}\DUrole{default_value}{\textquotesingle{}\textquotesingle{}}}\sphinxparamcomma \sphinxparam{\DUrole{n}{cmdargs}\DUrole{o}{=}\DUrole{default_value}{None}}\sphinxparamcomma \sphinxparam{\DUrole{n}{ptr}\DUrole{o}{=}\DUrole{default_value}{None}}\sphinxparamcomma \sphinxparam{\DUrole{n}{comm}\DUrole{o}{=}\DUrole{default_value}{None}}}{}
\pysigstopsignatures
\sphinxAtStartPar
Create an instance of the LAMMPS Python class.

\sphinxAtStartPar
This is a Python wrapper class that exposes the LAMMPS C\sphinxhyphen{}library
interface to Python.  It either requires that LAMMPS has been compiled
as shared library which is then dynamically loaded via the ctypes
Python module or that this module called from a Python function that
is called from a Python interpreter embedded into a LAMMPS executable,
for example through the {\hyperref[\detokenize{python::doc}]{\sphinxcrossref{\DUrole{doc}{python invoke}}}} command.
When the class is instantiated it calls the {\hyperref[\detokenize{Library_create:_CPPv411lammps_openiPPc8MPI_CommPPv}]{\sphinxcrossref{\sphinxcode{\sphinxupquote{lammps\_open()}}}}}
function of the LAMMPS C\sphinxhyphen{}library interface, which in
turn will create an instance of the {\hyperref[\detokenize{Classes_lammps:_CPPv4N9LAMMPS_NS6LAMMPSE}]{\sphinxcrossref{\sphinxcode{\sphinxupquote{LAMMPS}}}}}
C++ class.  The handle to this C++ class is stored internally
and automatically passed to the calls to the C library interface.
\begin{quote}\begin{description}
\sphinxlineitem{Parameters}\begin{itemize}
\item {} 
\sphinxAtStartPar
\sphinxstyleliteralstrong{\sphinxupquote{name}} (\sphinxstyleliteralemphasis{\sphinxupquote{string}}) \textendash{} “machine” name of the shared LAMMPS library (“mpi” loads \sphinxcode{\sphinxupquote{liblammps\_mpi.so}}, “” loads \sphinxcode{\sphinxupquote{liblammps.so}})

\item {} 
\sphinxAtStartPar
\sphinxstyleliteralstrong{\sphinxupquote{cmdargs}} (\sphinxstyleliteralemphasis{\sphinxupquote{list}}) \textendash{} list of command line arguments to be passed to the {\hyperref[\detokenize{Library_create:_CPPv411lammps_openiPPc8MPI_CommPPv}]{\sphinxcrossref{\sphinxcode{\sphinxupquote{lammps\_open()}}}}} function.  The executable name is automatically added.

\item {} 
\sphinxAtStartPar
\sphinxstyleliteralstrong{\sphinxupquote{ptr}} (\sphinxstyleliteralemphasis{\sphinxupquote{pointer}}) \textendash{} pointer to a LAMMPS C++ class instance when called from an embedded Python interpreter.  None means load symbols from shared library.

\item {} 
\sphinxAtStartPar
\sphinxstyleliteralstrong{\sphinxupquote{comm}} (\sphinxstyleliteralemphasis{\sphinxupquote{MPI\_Comm}}) \textendash{} MPI communicator (as provided by \sphinxhref{https://mpi4py.readthedocs.io/}{mpi4py}). \sphinxcode{\sphinxupquote{None}} means use \sphinxcode{\sphinxupquote{MPI\_COMM\_WORLD}} implicitly.

\end{itemize}

\end{description}\end{quote}
\index{numpy (lammps.lammps property)@\spxentry{numpy}\spxextra{lammps.lammps property}}

\begin{fulllineitems}
\phantomsection\label{\detokenize{Python_module:lammps.lammps.numpy}}
\pysigstartsignatures
\pysigline{\sphinxbfcode{\sphinxupquote{property\DUrole{w}{ }}}\sphinxbfcode{\sphinxupquote{numpy}}}
\pysigstopsignatures
\sphinxAtStartPar
Return object to access numpy versions of API

\sphinxAtStartPar
It provides alternative implementations of API functions that
return numpy arrays instead of ctypes pointers. If numpy is not installed,
accessing this property will lead to an ImportError.
\begin{quote}\begin{description}
\sphinxlineitem{Returns}
\sphinxAtStartPar
instance of numpy wrapper object

\sphinxlineitem{Return type}
\sphinxAtStartPar
{\hyperref[\detokenize{Python_module:lammps.numpy_wrapper.numpy_wrapper}]{\sphinxcrossref{numpy\_wrapper}}}

\end{description}\end{quote}

\end{fulllineitems}

\index{close() (lammps.lammps method)@\spxentry{close()}\spxextra{lammps.lammps method}}

\begin{fulllineitems}
\phantomsection\label{\detokenize{Python_module:lammps.lammps.close}}
\pysigstartsignatures
\pysiglinewithargsret{\sphinxbfcode{\sphinxupquote{close}}}{}{}
\pysigstopsignatures
\sphinxAtStartPar
Explicitly delete a LAMMPS instance through the C\sphinxhyphen{}library interface.

\sphinxAtStartPar
This is a wrapper around the {\hyperref[\detokenize{Library_create:_CPPv412lammps_closePv}]{\sphinxcrossref{\sphinxcode{\sphinxupquote{lammps\_close()}}}}} function of the C\sphinxhyphen{}library interface.

\end{fulllineitems}

\index{finalize() (lammps.lammps method)@\spxentry{finalize()}\spxextra{lammps.lammps method}}

\begin{fulllineitems}
\phantomsection\label{\detokenize{Python_module:lammps.lammps.finalize}}
\pysigstartsignatures
\pysiglinewithargsret{\sphinxbfcode{\sphinxupquote{finalize}}}{}{}
\pysigstopsignatures
\sphinxAtStartPar
Shut down the MPI communication and Kokkos environment (if active) through the
library interface by  calling {\hyperref[\detokenize{Library_create:_CPPv419lammps_mpi_finalizev}]{\sphinxcrossref{\sphinxcode{\sphinxupquote{lammps\_mpi\_finalize()}}}}} and
{\hyperref[\detokenize{Library_create:_CPPv422lammps_kokkos_finalizev}]{\sphinxcrossref{\sphinxcode{\sphinxupquote{lammps\_kokkos\_finalize()}}}}}.

\sphinxAtStartPar
You cannot create or use any LAMMPS instances after this function is called
unless LAMMPS was compiled without MPI and without Kokkos support.

\end{fulllineitems}

\index{error() (lammps.lammps method)@\spxentry{error()}\spxextra{lammps.lammps method}}

\begin{fulllineitems}
\phantomsection\label{\detokenize{Python_module:lammps.lammps.error}}
\pysigstartsignatures
\pysiglinewithargsret{\sphinxbfcode{\sphinxupquote{error}}}{\sphinxparam{\DUrole{n}{error\_type}}\sphinxparamcomma \sphinxparam{\DUrole{n}{error\_text}}}{}
\pysigstopsignatures
\sphinxAtStartPar
Forward error to the LAMMPS Error class.

\sphinxAtStartPar
\DUrole{versionmodified,added}{Added in version 3Nov2022.}

\sphinxAtStartPar
This is a wrapper around the {\hyperref[\detokenize{Library_create:_CPPv412lammps_errorPviPKc}]{\sphinxcrossref{\sphinxcode{\sphinxupquote{lammps\_error()}}}}} function of the C\sphinxhyphen{}library interface.
\begin{quote}\begin{description}
\sphinxlineitem{Parameters}\begin{itemize}
\item {} 
\sphinxAtStartPar
\sphinxstyleliteralstrong{\sphinxupquote{error\_type}} (\sphinxstyleliteralemphasis{\sphinxupquote{int}})

\item {} 
\sphinxAtStartPar
\sphinxstyleliteralstrong{\sphinxupquote{error\_text}} (\sphinxstyleliteralemphasis{\sphinxupquote{string}})

\end{itemize}

\end{description}\end{quote}

\end{fulllineitems}

\index{version() (lammps.lammps method)@\spxentry{version()}\spxextra{lammps.lammps method}}

\begin{fulllineitems}
\phantomsection\label{\detokenize{Python_module:lammps.lammps.version}}
\pysigstartsignatures
\pysiglinewithargsret{\sphinxbfcode{\sphinxupquote{version}}}{}{}
\pysigstopsignatures
\sphinxAtStartPar
Return a numerical representation of the LAMMPS version in use.

\sphinxAtStartPar
This is a wrapper around the {\hyperref[\detokenize{Library_config:_CPPv414lammps_versionPv}]{\sphinxcrossref{\sphinxcode{\sphinxupquote{lammps\_version()}}}}} function of the C\sphinxhyphen{}library interface.
\begin{quote}\begin{description}
\sphinxlineitem{Returns}
\sphinxAtStartPar
version number

\sphinxlineitem{Return type}
\sphinxAtStartPar
int

\end{description}\end{quote}

\end{fulllineitems}

\index{get\_os\_info() (lammps.lammps method)@\spxentry{get\_os\_info()}\spxextra{lammps.lammps method}}

\begin{fulllineitems}
\phantomsection\label{\detokenize{Python_module:lammps.lammps.get_os_info}}
\pysigstartsignatures
\pysiglinewithargsret{\sphinxbfcode{\sphinxupquote{get\_os\_info}}}{}{}
\pysigstopsignatures
\sphinxAtStartPar
Return a string with information about the OS and compiler runtime

\sphinxAtStartPar
This is a wrapper around the {\hyperref[\detokenize{Library_config:_CPPv418lammps_get_os_infoPci}]{\sphinxcrossref{\sphinxcode{\sphinxupquote{lammps\_get\_os\_info()}}}}} function of the C\sphinxhyphen{}library interface.
\begin{quote}\begin{description}
\sphinxlineitem{Returns}
\sphinxAtStartPar
OS info string

\sphinxlineitem{Return type}
\sphinxAtStartPar
string

\end{description}\end{quote}

\end{fulllineitems}

\index{get\_mpi\_comm() (lammps.lammps method)@\spxentry{get\_mpi\_comm()}\spxextra{lammps.lammps method}}

\begin{fulllineitems}
\phantomsection\label{\detokenize{Python_module:lammps.lammps.get_mpi_comm}}
\pysigstartsignatures
\pysiglinewithargsret{\sphinxbfcode{\sphinxupquote{get\_mpi\_comm}}}{}{}
\pysigstopsignatures
\sphinxAtStartPar
Get the MPI communicator in use by the current LAMMPS instance

\sphinxAtStartPar
This is a wrapper around the {\hyperref[\detokenize{Library_properties:_CPPv419lammps_get_mpi_commPv}]{\sphinxcrossref{\sphinxcode{\sphinxupquote{lammps\_get\_mpi\_comm()}}}}} function
of the C\sphinxhyphen{}library interface.  It will return \sphinxcode{\sphinxupquote{None}} if either the
LAMMPS library was compiled without MPI support or the mpi4py
Python module is not available.
\begin{quote}\begin{description}
\sphinxlineitem{Returns}
\sphinxAtStartPar
MPI communicator

\sphinxlineitem{Return type}
\sphinxAtStartPar
MPI\_Comm

\end{description}\end{quote}

\end{fulllineitems}

\index{file() (lammps.lammps method)@\spxentry{file()}\spxextra{lammps.lammps method}}

\begin{fulllineitems}
\phantomsection\label{\detokenize{Python_module:lammps.lammps.file}}
\pysigstartsignatures
\pysiglinewithargsret{\sphinxbfcode{\sphinxupquote{file}}}{\sphinxparam{\DUrole{n}{path}}}{}
\pysigstopsignatures
\sphinxAtStartPar
Read LAMMPS commands from a file.

\sphinxAtStartPar
This is a wrapper around the {\hyperref[\detokenize{Library_execute:_CPPv411lammps_filePvPKc}]{\sphinxcrossref{\sphinxcode{\sphinxupquote{lammps\_file()}}}}} function of the C\sphinxhyphen{}library interface.
It will open the file with the name/path \sphinxtitleref{file} and process the LAMMPS commands line by line until
the end. The function will return when the end of the file is reached.
\begin{quote}\begin{description}
\sphinxlineitem{Parameters}
\sphinxAtStartPar
\sphinxstyleliteralstrong{\sphinxupquote{path}} (\sphinxstyleliteralemphasis{\sphinxupquote{string}}) \textendash{} Name of the file/path with LAMMPS commands

\end{description}\end{quote}

\end{fulllineitems}

\index{command() (lammps.lammps method)@\spxentry{command()}\spxextra{lammps.lammps method}}

\begin{fulllineitems}
\phantomsection\label{\detokenize{Python_module:lammps.lammps.command}}
\pysigstartsignatures
\pysiglinewithargsret{\sphinxbfcode{\sphinxupquote{command}}}{\sphinxparam{\DUrole{n}{cmd}}}{}
\pysigstopsignatures
\sphinxAtStartPar
Process a single LAMMPS input command from a string.

\sphinxAtStartPar
This is a wrapper around the {\hyperref[\detokenize{Library_execute:_CPPv414lammps_commandPvPKc}]{\sphinxcrossref{\sphinxcode{\sphinxupquote{lammps\_command()}}}}}
function of the C\sphinxhyphen{}library interface.
\begin{quote}\begin{description}
\sphinxlineitem{Parameters}
\sphinxAtStartPar
\sphinxstyleliteralstrong{\sphinxupquote{cmd}} (\sphinxstyleliteralemphasis{\sphinxupquote{string}}) \textendash{} a single lammps command

\end{description}\end{quote}

\end{fulllineitems}

\index{commands\_list() (lammps.lammps method)@\spxentry{commands\_list()}\spxextra{lammps.lammps method}}

\begin{fulllineitems}
\phantomsection\label{\detokenize{Python_module:lammps.lammps.commands_list}}
\pysigstartsignatures
\pysiglinewithargsret{\sphinxbfcode{\sphinxupquote{commands\_list}}}{\sphinxparam{\DUrole{n}{cmdlist}}}{}
\pysigstopsignatures
\sphinxAtStartPar
Process multiple LAMMPS input commands from a list of strings.

\sphinxAtStartPar
This is a wrapper around the
{\hyperref[\detokenize{Library_execute:_CPPv420lammps_commands_listPviPPKc}]{\sphinxcrossref{\sphinxcode{\sphinxupquote{lammps\_commands\_list()}}}}} function of
the C\sphinxhyphen{}library interface.
\begin{quote}\begin{description}
\sphinxlineitem{Parameters}
\sphinxAtStartPar
\sphinxstyleliteralstrong{\sphinxupquote{cmdlist}} (\sphinxstyleliteralemphasis{\sphinxupquote{list}}\sphinxstyleliteralemphasis{\sphinxupquote{ of }}\sphinxstyleliteralemphasis{\sphinxupquote{strings}}) \textendash{} a single lammps command

\end{description}\end{quote}

\end{fulllineitems}

\index{commands\_string() (lammps.lammps method)@\spxentry{commands\_string()}\spxextra{lammps.lammps method}}

\begin{fulllineitems}
\phantomsection\label{\detokenize{Python_module:lammps.lammps.commands_string}}
\pysigstartsignatures
\pysiglinewithargsret{\sphinxbfcode{\sphinxupquote{commands\_string}}}{\sphinxparam{\DUrole{n}{multicmd}}}{}
\pysigstopsignatures
\sphinxAtStartPar
Process a block of LAMMPS input commands from a string.

\sphinxAtStartPar
This is a wrapper around the
{\hyperref[\detokenize{Library_execute:_CPPv422lammps_commands_stringPvPKc}]{\sphinxcrossref{\sphinxcode{\sphinxupquote{lammps\_commands\_string()}}}}}
function of the C\sphinxhyphen{}library interface.
\begin{quote}\begin{description}
\sphinxlineitem{Parameters}
\sphinxAtStartPar
\sphinxstyleliteralstrong{\sphinxupquote{multicmd}} (\sphinxstyleliteralemphasis{\sphinxupquote{string}}) \textendash{} text block of lammps commands

\end{description}\end{quote}

\end{fulllineitems}

\index{get\_natoms() (lammps.lammps method)@\spxentry{get\_natoms()}\spxextra{lammps.lammps method}}

\begin{fulllineitems}
\phantomsection\label{\detokenize{Python_module:lammps.lammps.get_natoms}}
\pysigstartsignatures
\pysiglinewithargsret{\sphinxbfcode{\sphinxupquote{get\_natoms}}}{}{}
\pysigstopsignatures
\sphinxAtStartPar
Get the total number of atoms in the LAMMPS instance.

\sphinxAtStartPar
Will be precise up to 53\sphinxhyphen{}bit signed integer due to the
underlying {\hyperref[\detokenize{Library_properties:_CPPv417lammps_get_natomsPv}]{\sphinxcrossref{\sphinxcode{\sphinxupquote{lammps\_get\_natoms()}}}}} function returning a double.
\begin{quote}\begin{description}
\sphinxlineitem{Returns}
\sphinxAtStartPar
number of atoms

\sphinxlineitem{Return type}
\sphinxAtStartPar
int

\end{description}\end{quote}

\end{fulllineitems}

\index{extract\_box() (lammps.lammps method)@\spxentry{extract\_box()}\spxextra{lammps.lammps method}}

\begin{fulllineitems}
\phantomsection\label{\detokenize{Python_module:lammps.lammps.extract_box}}
\pysigstartsignatures
\pysiglinewithargsret{\sphinxbfcode{\sphinxupquote{extract\_box}}}{}{}
\pysigstopsignatures
\sphinxAtStartPar
Extract simulation box parameters

\sphinxAtStartPar
This is a wrapper around the {\hyperref[\detokenize{Library_properties:_CPPv418lammps_extract_boxPvPdPdPdPdPdPiPi}]{\sphinxcrossref{\sphinxcode{\sphinxupquote{lammps\_extract\_box()}}}}} function
of the C\sphinxhyphen{}library interface.  Unlike in the C function, the result is
returned as a list.
\begin{quote}\begin{description}
\sphinxlineitem{Returns}
\sphinxAtStartPar
list of the extracted data: boxlo, boxhi, xy, yz, xz, periodicity, box\_change

\sphinxlineitem{Return type}
\sphinxAtStartPar
{[} 3*double, 3*double, double, double, 3*int, int{]}

\end{description}\end{quote}

\end{fulllineitems}

\index{reset\_box() (lammps.lammps method)@\spxentry{reset\_box()}\spxextra{lammps.lammps method}}

\begin{fulllineitems}
\phantomsection\label{\detokenize{Python_module:lammps.lammps.reset_box}}
\pysigstartsignatures
\pysiglinewithargsret{\sphinxbfcode{\sphinxupquote{reset\_box}}}{\sphinxparam{\DUrole{n}{boxlo}}\sphinxparamcomma \sphinxparam{\DUrole{n}{boxhi}}\sphinxparamcomma \sphinxparam{\DUrole{n}{xy}}\sphinxparamcomma \sphinxparam{\DUrole{n}{yz}}\sphinxparamcomma \sphinxparam{\DUrole{n}{xz}}}{}
\pysigstopsignatures
\sphinxAtStartPar
Reset simulation box parameters

\sphinxAtStartPar
This is a wrapper around the {\hyperref[\detokenize{Library_properties:_CPPv416lammps_reset_boxPvPdPdddd}]{\sphinxcrossref{\sphinxcode{\sphinxupquote{lammps\_reset\_box()}}}}} function
of the C\sphinxhyphen{}library interface.
\begin{quote}\begin{description}
\sphinxlineitem{Parameters}\begin{itemize}
\item {} 
\sphinxAtStartPar
\sphinxstyleliteralstrong{\sphinxupquote{boxlo}} (\sphinxstyleliteralemphasis{\sphinxupquote{list}}\sphinxstyleliteralemphasis{\sphinxupquote{ of }}\sphinxstyleliteralemphasis{\sphinxupquote{3 floating point numbers}}) \textendash{} new lower box boundaries

\item {} 
\sphinxAtStartPar
\sphinxstyleliteralstrong{\sphinxupquote{boxhi}} (\sphinxstyleliteralemphasis{\sphinxupquote{list}}\sphinxstyleliteralemphasis{\sphinxupquote{ of }}\sphinxstyleliteralemphasis{\sphinxupquote{3 floating point numbers}}) \textendash{} new upper box boundaries

\item {} 
\sphinxAtStartPar
\sphinxstyleliteralstrong{\sphinxupquote{xy}} (\sphinxstyleliteralemphasis{\sphinxupquote{float}}) \textendash{} xy tilt factor

\item {} 
\sphinxAtStartPar
\sphinxstyleliteralstrong{\sphinxupquote{yz}} (\sphinxstyleliteralemphasis{\sphinxupquote{float}}) \textendash{} yz tilt factor

\item {} 
\sphinxAtStartPar
\sphinxstyleliteralstrong{\sphinxupquote{xz}} (\sphinxstyleliteralemphasis{\sphinxupquote{float}}) \textendash{} xz tilt factor

\end{itemize}

\end{description}\end{quote}

\end{fulllineitems}

\index{get\_thermo() (lammps.lammps method)@\spxentry{get\_thermo()}\spxextra{lammps.lammps method}}

\begin{fulllineitems}
\phantomsection\label{\detokenize{Python_module:lammps.lammps.get_thermo}}
\pysigstartsignatures
\pysiglinewithargsret{\sphinxbfcode{\sphinxupquote{get\_thermo}}}{\sphinxparam{\DUrole{n}{name}}}{}
\pysigstopsignatures
\sphinxAtStartPar
Get current value of a thermo keyword

\sphinxAtStartPar
This is a wrapper around the {\hyperref[\detokenize{Library_properties:_CPPv417lammps_get_thermoPvPKc}]{\sphinxcrossref{\sphinxcode{\sphinxupquote{lammps\_get\_thermo()}}}}}
function of the C\sphinxhyphen{}library interface.
\begin{quote}\begin{description}
\sphinxlineitem{Parameters}
\sphinxAtStartPar
\sphinxstyleliteralstrong{\sphinxupquote{name}} (\sphinxstyleliteralemphasis{\sphinxupquote{string}}) \textendash{} name of thermo keyword

\sphinxlineitem{Returns}
\sphinxAtStartPar
value of thermo keyword

\sphinxlineitem{Return type}
\sphinxAtStartPar
double or None

\end{description}\end{quote}

\end{fulllineitems}

\index{last\_thermo\_step (lammps.lammps property)@\spxentry{last\_thermo\_step}\spxextra{lammps.lammps property}}

\begin{fulllineitems}
\phantomsection\label{\detokenize{Python_module:lammps.lammps.last_thermo_step}}
\pysigstartsignatures
\pysigline{\sphinxbfcode{\sphinxupquote{property\DUrole{w}{ }}}\sphinxbfcode{\sphinxupquote{last\_thermo\_step}}}
\pysigstopsignatures
\sphinxAtStartPar
Get the last timestep where thermodynamic data was computed
\begin{quote}\begin{description}
\sphinxlineitem{Returns}
\sphinxAtStartPar
the timestep or a negative number if there has not been any thermo output yet

\sphinxlineitem{Return type}
\sphinxAtStartPar
int

\end{description}\end{quote}

\end{fulllineitems}

\index{last\_thermo() (lammps.lammps method)@\spxentry{last\_thermo()}\spxextra{lammps.lammps method}}

\begin{fulllineitems}
\phantomsection\label{\detokenize{Python_module:lammps.lammps.last_thermo}}
\pysigstartsignatures
\pysiglinewithargsret{\sphinxbfcode{\sphinxupquote{last\_thermo}}}{}{}
\pysigstopsignatures
\sphinxAtStartPar
Get a dictionary of the last thermodynamic output

\sphinxAtStartPar
This is a wrapper around the {\hyperref[\detokenize{Library_properties:_CPPv418lammps_last_thermoPvPKci}]{\sphinxcrossref{\sphinxcode{\sphinxupquote{lammps\_last\_thermo()}}}}}
function of the C\sphinxhyphen{}library interface.  It collects the cached thermo
data from the last timestep into a dictionary.  The return value
is None, if there has not been any thermo output yet.
\begin{quote}\begin{description}
\sphinxlineitem{Returns}
\sphinxAtStartPar
a dictionary containing the last computed thermo output values

\sphinxlineitem{Return type}
\sphinxAtStartPar
dict or None

\end{description}\end{quote}

\end{fulllineitems}

\index{extract\_setting() (lammps.lammps method)@\spxentry{extract\_setting()}\spxextra{lammps.lammps method}}

\begin{fulllineitems}
\phantomsection\label{\detokenize{Python_module:lammps.lammps.extract_setting}}
\pysigstartsignatures
\pysiglinewithargsret{\sphinxbfcode{\sphinxupquote{extract\_setting}}}{\sphinxparam{\DUrole{n}{name}}}{}
\pysigstopsignatures
\sphinxAtStartPar
Query LAMMPS about global settings that can be expressed as an integer.

\sphinxAtStartPar
This is a wrapper around the {\hyperref[\detokenize{Library_properties:_CPPv422lammps_extract_settingPvPKc}]{\sphinxcrossref{\sphinxcode{\sphinxupquote{lammps\_extract\_setting()}}}}}
function of the C\sphinxhyphen{}library interface.  Its documentation includes
a list of the supported keywords.
\begin{quote}\begin{description}
\sphinxlineitem{Parameters}
\sphinxAtStartPar
\sphinxstyleliteralstrong{\sphinxupquote{name}} (\sphinxstyleliteralemphasis{\sphinxupquote{string}}) \textendash{} name of the setting

\sphinxlineitem{Returns}
\sphinxAtStartPar
value of the setting

\sphinxlineitem{Return type}
\sphinxAtStartPar
int

\end{description}\end{quote}

\end{fulllineitems}

\index{extract\_global\_datatype() (lammps.lammps method)@\spxentry{extract\_global\_datatype()}\spxextra{lammps.lammps method}}

\begin{fulllineitems}
\phantomsection\label{\detokenize{Python_module:lammps.lammps.extract_global_datatype}}
\pysigstartsignatures
\pysiglinewithargsret{\sphinxbfcode{\sphinxupquote{extract\_global\_datatype}}}{\sphinxparam{\DUrole{n}{name}}}{}
\pysigstopsignatures
\sphinxAtStartPar
Retrieve global property datatype from LAMMPS

\sphinxAtStartPar
This is a wrapper around the {\hyperref[\detokenize{Library_properties:_CPPv430lammps_extract_global_datatypePvPKc}]{\sphinxcrossref{\sphinxcode{\sphinxupquote{lammps\_extract\_global\_datatype()}}}}}
function of the C\sphinxhyphen{}library interface. Its documentation includes a
list of the supported keywords.
This function returns \sphinxcode{\sphinxupquote{None}} if the keyword is not
recognized. Otherwise it will return a positive integer value that
corresponds to one of the {\hyperref[\detokenize{Python_module:py-datatype-constants}]{\sphinxcrossref{\DUrole{std,std-ref}{data type}}}}
constants defined in the {\hyperref[\detokenize{Python_module:module-lammps}]{\sphinxcrossref{\sphinxcode{\sphinxupquote{lammps}}}}} module.
\begin{quote}\begin{description}
\sphinxlineitem{Parameters}
\sphinxAtStartPar
\sphinxstyleliteralstrong{\sphinxupquote{name}} (\sphinxstyleliteralemphasis{\sphinxupquote{string}}) \textendash{} name of the property

\sphinxlineitem{Returns}
\sphinxAtStartPar
data type of global property, see {\hyperref[\detokenize{Python_module:py-datatype-constants}]{\sphinxcrossref{\DUrole{std,std-ref}{Data Types}}}}

\sphinxlineitem{Return type}
\sphinxAtStartPar
int

\end{description}\end{quote}

\end{fulllineitems}

\index{extract\_global() (lammps.lammps method)@\spxentry{extract\_global()}\spxextra{lammps.lammps method}}

\begin{fulllineitems}
\phantomsection\label{\detokenize{Python_module:lammps.lammps.extract_global}}
\pysigstartsignatures
\pysiglinewithargsret{\sphinxbfcode{\sphinxupquote{extract\_global}}}{\sphinxparam{\DUrole{n}{name}}\sphinxparamcomma \sphinxparam{\DUrole{n}{dtype}\DUrole{o}{=}\DUrole{default_value}{None}}}{}
\pysigstopsignatures
\sphinxAtStartPar
Query LAMMPS about global settings of different types.

\sphinxAtStartPar
This is a wrapper around the {\hyperref[\detokenize{Library_properties:_CPPv421lammps_extract_globalPvPKc}]{\sphinxcrossref{\sphinxcode{\sphinxupquote{lammps\_extract\_global()}}}}} function
of the C\sphinxhyphen{}library interface.  Since there are no pointers in Python, this
method will \sphinxhyphen{} unlike the C function \sphinxhyphen{} return the value or a list of
values.  The {\hyperref[\detokenize{Library_properties:_CPPv421lammps_extract_globalPvPKc}]{\sphinxcrossref{\sphinxcode{\sphinxupquote{lammps\_extract\_global()}}}}} documentation includes a
list of the supported keywords and their data types.
Since Python needs to know the data type to be able to interpret
the result, by default, this function will try to auto\sphinxhyphen{}detect the data type
by asking the library. You can also force a specific data type.  For that
purpose the {\hyperref[\detokenize{Python_module:module-lammps}]{\sphinxcrossref{\sphinxcode{\sphinxupquote{lammps}}}}} module contains {\hyperref[\detokenize{Python_module:py-datatype-constants}]{\sphinxcrossref{\DUrole{std,std-ref}{data type}}}}
constants. This function returns \sphinxcode{\sphinxupquote{None}} if either the keyword is not recognized,
or an invalid data type constant is used.
\begin{quote}\begin{description}
\sphinxlineitem{Parameters}\begin{itemize}
\item {} 
\sphinxAtStartPar
\sphinxstyleliteralstrong{\sphinxupquote{name}} (\sphinxstyleliteralemphasis{\sphinxupquote{string}}) \textendash{} name of the property

\item {} 
\sphinxAtStartPar
\sphinxstyleliteralstrong{\sphinxupquote{dtype}} (\sphinxstyleliteralemphasis{\sphinxupquote{int}}\sphinxstyleliteralemphasis{\sphinxupquote{, }}\sphinxstyleliteralemphasis{\sphinxupquote{optional}}) \textendash{} data type of the returned data (see {\hyperref[\detokenize{Python_module:py-datatype-constants}]{\sphinxcrossref{\DUrole{std,std-ref}{Data Types}}}})

\end{itemize}

\sphinxlineitem{Returns}
\sphinxAtStartPar
value of the property or list of values or None

\sphinxlineitem{Return type}
\sphinxAtStartPar
int, float, list, or NoneType

\end{description}\end{quote}

\end{fulllineitems}

\index{extract\_pair\_dimension() (lammps.lammps method)@\spxentry{extract\_pair\_dimension()}\spxextra{lammps.lammps method}}

\begin{fulllineitems}
\phantomsection\label{\detokenize{Python_module:lammps.lammps.extract_pair_dimension}}
\pysigstartsignatures
\pysiglinewithargsret{\sphinxbfcode{\sphinxupquote{extract\_pair\_dimension}}}{\sphinxparam{\DUrole{n}{name}}}{}
\pysigstopsignatures
\sphinxAtStartPar
Retrieve pair style  property dimensionality from LAMMPS

\sphinxAtStartPar
\DUrole{versionmodified,added}{Added in version 29Aug2024.}

\sphinxAtStartPar
This is a wrapper around the {\hyperref[\detokenize{Library_properties:_CPPv429lammps_extract_pair_dimensionPvPKc}]{\sphinxcrossref{\sphinxcode{\sphinxupquote{lammps\_extract\_pair\_dimension()}}}}}
function of the C\sphinxhyphen{}library interface. The list of supported keywords
depends on the pair style. This function returns \sphinxcode{\sphinxupquote{None}} if the keyword
is not recognized.
\begin{quote}\begin{description}
\sphinxlineitem{Parameters}
\sphinxAtStartPar
\sphinxstyleliteralstrong{\sphinxupquote{name}} (\sphinxstyleliteralemphasis{\sphinxupquote{string}}) \textendash{} name of the property

\sphinxlineitem{Returns}
\sphinxAtStartPar
dimensionality of the extractable data (typically 0, 1, or 2)

\sphinxlineitem{Return type}
\sphinxAtStartPar
int

\end{description}\end{quote}

\end{fulllineitems}

\index{extract\_pair() (lammps.lammps method)@\spxentry{extract\_pair()}\spxextra{lammps.lammps method}}

\begin{fulllineitems}
\phantomsection\label{\detokenize{Python_module:lammps.lammps.extract_pair}}
\pysigstartsignatures
\pysiglinewithargsret{\sphinxbfcode{\sphinxupquote{extract\_pair}}}{\sphinxparam{\DUrole{n}{name}}}{}
\pysigstopsignatures
\sphinxAtStartPar
Extract pair style data from LAMMPS.

\sphinxAtStartPar
\DUrole{versionmodified,added}{Added in version 29Aug2024.}

\sphinxAtStartPar
This is a wrapper around the {\hyperref[\detokenize{Library_properties:_CPPv419lammps_extract_pairPvPKc}]{\sphinxcrossref{\sphinxcode{\sphinxupquote{lammps\_extract\_pair()}}}}} function
of the C\sphinxhyphen{}library interface.  Since there are no pointers in Python, this
method will \sphinxhyphen{} unlike the C function \sphinxhyphen{} return the value or a list of
values.
Since Python needs to know the dimensionality to be able to interpret
the result, this function will detect the dimensionality by asking the library.
This function returns \sphinxcode{\sphinxupquote{None}} if the keyword is not recognized.
\begin{quote}\begin{description}
\sphinxlineitem{Parameters}
\sphinxAtStartPar
\sphinxstyleliteralstrong{\sphinxupquote{name}} (\sphinxstyleliteralemphasis{\sphinxupquote{string}}) \textendash{} name of the property

\sphinxlineitem{Returns}
\sphinxAtStartPar
value of the property or list of values or None

\sphinxlineitem{Return type}
\sphinxAtStartPar
float, list of float, list of list of floats, or NoneType

\end{description}\end{quote}

\end{fulllineitems}

\index{map\_atom() (lammps.lammps method)@\spxentry{map\_atom()}\spxextra{lammps.lammps method}}

\begin{fulllineitems}
\phantomsection\label{\detokenize{Python_module:lammps.lammps.map_atom}}
\pysigstartsignatures
\pysiglinewithargsret{\sphinxbfcode{\sphinxupquote{map\_atom}}}{\sphinxparam{\DUrole{n}{id}}}{}
\pysigstopsignatures
\sphinxAtStartPar
Map a global atom ID (aka tag) to the local atom index

\sphinxAtStartPar
This is a wrapper around the {\hyperref[\detokenize{Library_properties:_CPPv415lammps_map_atomPvPKv}]{\sphinxcrossref{\sphinxcode{\sphinxupquote{lammps\_map\_atom()}}}}}
function of the C\sphinxhyphen{}library interface.
\begin{quote}\begin{description}
\sphinxlineitem{Parameters}
\sphinxAtStartPar
\sphinxstyleliteralstrong{\sphinxupquote{id}} (\sphinxstyleliteralemphasis{\sphinxupquote{int}}) \textendash{} atom ID

\sphinxlineitem{Returns}
\sphinxAtStartPar
local index

\sphinxlineitem{Return type}
\sphinxAtStartPar
int

\end{description}\end{quote}

\end{fulllineitems}

\index{extract\_atom\_datatype() (lammps.lammps method)@\spxentry{extract\_atom\_datatype()}\spxextra{lammps.lammps method}}

\begin{fulllineitems}
\phantomsection\label{\detokenize{Python_module:lammps.lammps.extract_atom_datatype}}
\pysigstartsignatures
\pysiglinewithargsret{\sphinxbfcode{\sphinxupquote{extract\_atom\_datatype}}}{\sphinxparam{\DUrole{n}{name}}}{}
\pysigstopsignatures
\sphinxAtStartPar
Retrieve per\sphinxhyphen{}atom property datatype from LAMMPS

\sphinxAtStartPar
This is a wrapper around the {\hyperref[\detokenize{Library_atoms:_CPPv428lammps_extract_atom_datatypePvPKc}]{\sphinxcrossref{\sphinxcode{\sphinxupquote{lammps\_extract\_atom\_datatype()}}}}}
function of the C\sphinxhyphen{}library interface. Its documentation includes a
list of the supported keywords.
This function returns \sphinxcode{\sphinxupquote{None}} if the keyword is not
recognized. Otherwise it will return an integer value that
corresponds to one of the {\hyperref[\detokenize{Python_module:py-datatype-constants}]{\sphinxcrossref{\DUrole{std,std-ref}{data type}}}} constants
defined in the {\hyperref[\detokenize{Python_module:module-lammps}]{\sphinxcrossref{\sphinxcode{\sphinxupquote{lammps}}}}} module.
\begin{quote}\begin{description}
\sphinxlineitem{Parameters}
\sphinxAtStartPar
\sphinxstyleliteralstrong{\sphinxupquote{name}} (\sphinxstyleliteralemphasis{\sphinxupquote{string}}) \textendash{} name of the property

\sphinxlineitem{Returns}
\sphinxAtStartPar
data type of per\sphinxhyphen{}atom property (see {\hyperref[\detokenize{Python_module:py-datatype-constants}]{\sphinxcrossref{\DUrole{std,std-ref}{Data Types}}}})

\sphinxlineitem{Return type}
\sphinxAtStartPar
int

\end{description}\end{quote}

\end{fulllineitems}

\index{extract\_atom() (lammps.lammps method)@\spxentry{extract\_atom()}\spxextra{lammps.lammps method}}

\begin{fulllineitems}
\phantomsection\label{\detokenize{Python_module:lammps.lammps.extract_atom}}
\pysigstartsignatures
\pysiglinewithargsret{\sphinxbfcode{\sphinxupquote{extract\_atom}}}{\sphinxparam{\DUrole{n}{name}}\sphinxparamcomma \sphinxparam{\DUrole{n}{dtype}\DUrole{o}{=}\DUrole{default_value}{None}}}{}
\pysigstopsignatures
\sphinxAtStartPar
Retrieve per\sphinxhyphen{}atom properties from LAMMPS

\sphinxAtStartPar
This is a wrapper around the {\hyperref[\detokenize{Library_atoms:_CPPv419lammps_extract_atomPvPKc}]{\sphinxcrossref{\sphinxcode{\sphinxupquote{lammps\_extract\_atom()}}}}}
function of the C\sphinxhyphen{}library interface. Its documentation includes a
list of the supported keywords and their data types.
Since Python needs to know the data type to be able to interpret
the result, by default, this function will try to auto\sphinxhyphen{}detect the data type
by asking the library. You can also force a specific data type by setting \sphinxcode{\sphinxupquote{dtype}}
to one of the {\hyperref[\detokenize{Python_module:py-datatype-constants}]{\sphinxcrossref{\DUrole{std,std-ref}{data type}}}} constants defined in the
{\hyperref[\detokenize{Python_module:module-lammps}]{\sphinxcrossref{\sphinxcode{\sphinxupquote{lammps}}}}} module.
This function returns \sphinxcode{\sphinxupquote{None}} if either the keyword is not
recognized, or an invalid data type constant is used.

\begin{sphinxadmonition}{note}{Note:}
\sphinxAtStartPar
While the returned arrays of per\sphinxhyphen{}atom data are dimensioned
for the range {[}0:nmax{]} \sphinxhyphen{} as is the underlying storage \sphinxhyphen{}
the data is usually only valid for the range of {[}0:nlocal{]},
unless the property of interest is also updated for ghost
atoms.  In some cases, this depends on a LAMMPS setting, see
for example {\hyperref[\detokenize{comm_modify::doc}]{\sphinxcrossref{\DUrole{doc}{comm\_modify vel yes}}}}.
\end{sphinxadmonition}
\begin{quote}\begin{description}
\sphinxlineitem{Parameters}\begin{itemize}
\item {} 
\sphinxAtStartPar
\sphinxstyleliteralstrong{\sphinxupquote{name}} (\sphinxstyleliteralemphasis{\sphinxupquote{string}}) \textendash{} name of the property

\item {} 
\sphinxAtStartPar
\sphinxstyleliteralstrong{\sphinxupquote{dtype}} (\sphinxstyleliteralemphasis{\sphinxupquote{int}}\sphinxstyleliteralemphasis{\sphinxupquote{, }}\sphinxstyleliteralemphasis{\sphinxupquote{optional}}) \textendash{} data type of the returned data (see {\hyperref[\detokenize{Python_module:py-datatype-constants}]{\sphinxcrossref{\DUrole{std,std-ref}{Data Types}}}})

\end{itemize}

\sphinxlineitem{Returns}
\sphinxAtStartPar
requested data or \sphinxcode{\sphinxupquote{None}}

\sphinxlineitem{Return type}
\sphinxAtStartPar
ctypes.POINTER(ctypes.c\_int32), ctypes.POINTER(ctypes.POINTER(ctypes.c\_int32)),
ctypes.POINTER(ctypes.c\_int64), ctypes.POINTER(ctypes.POINTER(ctypes.c\_int64)),
ctypes.POINTER(ctypes.c\_double), ctypes.POINTER(ctypes.POINTER(ctypes.c\_double)),
or NoneType

\end{description}\end{quote}

\end{fulllineitems}

\index{extract\_compute() (lammps.lammps method)@\spxentry{extract\_compute()}\spxextra{lammps.lammps method}}

\begin{fulllineitems}
\phantomsection\label{\detokenize{Python_module:lammps.lammps.extract_compute}}
\pysigstartsignatures
\pysiglinewithargsret{\sphinxbfcode{\sphinxupquote{extract\_compute}}}{\sphinxparam{\DUrole{n}{cid}}\sphinxparamcomma \sphinxparam{\DUrole{n}{cstyle}}\sphinxparamcomma \sphinxparam{\DUrole{n}{ctype}}}{}
\pysigstopsignatures
\sphinxAtStartPar
Retrieve data from a LAMMPS compute

\sphinxAtStartPar
This is a wrapper around the {\hyperref[\detokenize{Library_objects:_CPPv422lammps_extract_computePvPKcii}]{\sphinxcrossref{\sphinxcode{\sphinxupquote{lammps\_extract\_compute()}}}}}
function of the C\sphinxhyphen{}library interface.
This function returns \sphinxcode{\sphinxupquote{None}} if either the compute id is not
recognized, or an invalid combination of {\hyperref[\detokenize{Python_module:py-style-constants}]{\sphinxcrossref{\DUrole{std,std-ref}{cstyle}}}}
and {\hyperref[\detokenize{Python_module:py-type-constants}]{\sphinxcrossref{\DUrole{std,std-ref}{ctype}}}} constants is used. The
names and functionality of the constants are the same as for
the corresponding C\sphinxhyphen{}library function.  For requests to return
a scalar or a size, the value is returned, otherwise a pointer.
\begin{quote}\begin{description}
\sphinxlineitem{Parameters}\begin{itemize}
\item {} 
\sphinxAtStartPar
\sphinxstyleliteralstrong{\sphinxupquote{cid}} (\sphinxstyleliteralemphasis{\sphinxupquote{string}}) \textendash{} compute ID

\item {} 
\sphinxAtStartPar
\sphinxstyleliteralstrong{\sphinxupquote{cstyle}} (\sphinxstyleliteralemphasis{\sphinxupquote{int}}) \textendash{} style of the data retrieve (global, atom, or local), see {\hyperref[\detokenize{Python_module:py-style-constants}]{\sphinxcrossref{\DUrole{std,std-ref}{Style Constants}}}}

\item {} 
\sphinxAtStartPar
\sphinxstyleliteralstrong{\sphinxupquote{ctype}} (\sphinxstyleliteralemphasis{\sphinxupquote{int}}) \textendash{} type or size of the returned data (scalar, vector, or array), see {\hyperref[\detokenize{Python_module:py-type-constants}]{\sphinxcrossref{\DUrole{std,std-ref}{Type Constants}}}}

\end{itemize}

\sphinxlineitem{Returns}
\sphinxAtStartPar
requested data as scalar, pointer to 1d or 2d double array, or None

\sphinxlineitem{Return type}
\sphinxAtStartPar
c\_double, ctypes.POINTER(c\_double), ctypes.POINTER(ctypes.POINTER(c\_double)), or NoneType

\end{description}\end{quote}

\end{fulllineitems}

\index{extract\_fix() (lammps.lammps method)@\spxentry{extract\_fix()}\spxextra{lammps.lammps method}}

\begin{fulllineitems}
\phantomsection\label{\detokenize{Python_module:lammps.lammps.extract_fix}}
\pysigstartsignatures
\pysiglinewithargsret{\sphinxbfcode{\sphinxupquote{extract\_fix}}}{\sphinxparam{\DUrole{n}{fid}}\sphinxparamcomma \sphinxparam{\DUrole{n}{fstyle}}\sphinxparamcomma \sphinxparam{\DUrole{n}{ftype}}\sphinxparamcomma \sphinxparam{\DUrole{n}{nrow}\DUrole{o}{=}\DUrole{default_value}{0}}\sphinxparamcomma \sphinxparam{\DUrole{n}{ncol}\DUrole{o}{=}\DUrole{default_value}{0}}}{}
\pysigstopsignatures
\sphinxAtStartPar
Retrieve data from a LAMMPS fix

\sphinxAtStartPar
This is a wrapper around the {\hyperref[\detokenize{Library_objects:_CPPv418lammps_extract_fixPvPKciiii}]{\sphinxcrossref{\sphinxcode{\sphinxupquote{lammps\_extract\_fix()}}}}} function
of the C\sphinxhyphen{}library interface.  This function returns \sphinxcode{\sphinxupquote{None}} if
either the fix id is not recognized, or an invalid combination of
{\hyperref[\detokenize{Python_module:py-style-constants}]{\sphinxcrossref{\DUrole{std,std-ref}{fstyle}}}} and {\hyperref[\detokenize{Python_module:py-type-constants}]{\sphinxcrossref{\DUrole{std,std-ref}{ftype}}}} constants is used. The names and functionality
of the constants are the same as for the corresponding C\sphinxhyphen{}library
function.  For requests to return a scalar or a size, the value is
returned, also when accessing global vectors or arrays, otherwise a
pointer.

\begin{sphinxadmonition}{note}{Note:}
\sphinxAtStartPar
When requesting global data, the fix data can only be accessed
one item at a time without access to the whole vector or array.
Thus this function will always return a scalar.  To access vector
or array elements the “nrow” and “ncol” arguments need to be set
accordingly (they default to 0).
\end{sphinxadmonition}
\begin{quote}\begin{description}
\sphinxlineitem{Parameters}\begin{itemize}
\item {} 
\sphinxAtStartPar
\sphinxstyleliteralstrong{\sphinxupquote{fid}} (\sphinxstyleliteralemphasis{\sphinxupquote{string}}) \textendash{} fix ID

\item {} 
\sphinxAtStartPar
\sphinxstyleliteralstrong{\sphinxupquote{fstyle}} (\sphinxstyleliteralemphasis{\sphinxupquote{int}}) \textendash{} style of the data retrieve (global, atom, or local), see {\hyperref[\detokenize{Python_module:py-style-constants}]{\sphinxcrossref{\DUrole{std,std-ref}{Style Constants}}}}

\item {} 
\sphinxAtStartPar
\sphinxstyleliteralstrong{\sphinxupquote{ftype}} (\sphinxstyleliteralemphasis{\sphinxupquote{int}}) \textendash{} type or size of the returned data (scalar, vector, or array), see {\hyperref[\detokenize{Python_module:py-type-constants}]{\sphinxcrossref{\DUrole{std,std-ref}{Type Constants}}}}

\item {} 
\sphinxAtStartPar
\sphinxstyleliteralstrong{\sphinxupquote{nrow}} (\sphinxstyleliteralemphasis{\sphinxupquote{int}}) \textendash{} index of global vector element or row index of global array element

\item {} 
\sphinxAtStartPar
\sphinxstyleliteralstrong{\sphinxupquote{ncol}} (\sphinxstyleliteralemphasis{\sphinxupquote{int}}) \textendash{} column index of global array element

\end{itemize}

\sphinxlineitem{Returns}
\sphinxAtStartPar
requested data or None

\sphinxlineitem{Return type}
\sphinxAtStartPar
c\_double, ctypes.POINTER(c\_double), ctypes.POINTER(ctypes.POINTER(c\_double)), or NoneType

\end{description}\end{quote}

\end{fulllineitems}

\index{extract\_variable() (lammps.lammps method)@\spxentry{extract\_variable()}\spxextra{lammps.lammps method}}

\begin{fulllineitems}
\phantomsection\label{\detokenize{Python_module:lammps.lammps.extract_variable}}
\pysigstartsignatures
\pysiglinewithargsret{\sphinxbfcode{\sphinxupquote{extract\_variable}}}{\sphinxparam{\DUrole{n}{name}}\sphinxparamcomma \sphinxparam{\DUrole{n}{group}\DUrole{o}{=}\DUrole{default_value}{None}}\sphinxparamcomma \sphinxparam{\DUrole{n}{vartype}\DUrole{o}{=}\DUrole{default_value}{None}}}{}
\pysigstopsignatures
\sphinxAtStartPar
Evaluate a LAMMPS variable and return its data

\sphinxAtStartPar
This function is a wrapper around the function
{\hyperref[\detokenize{Library_objects:_CPPv423lammps_extract_variablePvPKcPKc}]{\sphinxcrossref{\sphinxcode{\sphinxupquote{lammps\_extract\_variable()}}}}} of the C library interface,
evaluates variable name and returns a copy of the computed data.
The memory temporarily allocated by the C\sphinxhyphen{}interface is deleted
after the data is copied to a Python variable or list.
The variable must be either an equal\sphinxhyphen{}style (or equivalent)
variable or an atom\sphinxhyphen{}style variable. The variable type can be
provided as the \sphinxcode{\sphinxupquote{vartype}} parameter, which may be one of several
constants: \sphinxcode{\sphinxupquote{LMP\_VAR\_EQUAL}}, \sphinxcode{\sphinxupquote{LMP\_VAR\_ATOM}}, \sphinxcode{\sphinxupquote{LMP\_VAR\_VECTOR}},
or \sphinxcode{\sphinxupquote{LMP\_VAR\_STRING}}. If omitted or \sphinxcode{\sphinxupquote{None}}, LAMMPS will determine its
value for you based on a call to
{\hyperref[\detokenize{Library_objects:_CPPv432lammps_extract_variable_datatypePvPKc}]{\sphinxcrossref{\sphinxcode{\sphinxupquote{lammps\_extract\_variable\_datatype()}}}}} from the C library interface.
The group parameter is only used for atom\sphinxhyphen{}style variables and defaults to
the group “all”.
\begin{quote}\begin{description}
\sphinxlineitem{Parameters}\begin{itemize}
\item {} 
\sphinxAtStartPar
\sphinxstyleliteralstrong{\sphinxupquote{name}} (\sphinxstyleliteralemphasis{\sphinxupquote{string}}) \textendash{} name of the variable to execute

\item {} 
\sphinxAtStartPar
\sphinxstyleliteralstrong{\sphinxupquote{group}} (\sphinxstyleliteralemphasis{\sphinxupquote{string}}\sphinxstyleliteralemphasis{\sphinxupquote{, }}\sphinxstyleliteralemphasis{\sphinxupquote{only for atom\sphinxhyphen{}style variables}}) \textendash{} name of group for atom\sphinxhyphen{}style variable

\item {} 
\sphinxAtStartPar
\sphinxstyleliteralstrong{\sphinxupquote{vartype}} (\sphinxstyleliteralemphasis{\sphinxupquote{int}}) \textendash{} type of variable, see {\hyperref[\detokenize{Python_module:py-vartype-constants}]{\sphinxcrossref{\DUrole{std,std-ref}{Variable Type Constants}}}}

\end{itemize}

\sphinxlineitem{Returns}
\sphinxAtStartPar
the requested data

\sphinxlineitem{Return type}
\sphinxAtStartPar
c\_double, (c\_double), or NoneType

\end{description}\end{quote}

\end{fulllineitems}

\index{flush\_buffers() (lammps.lammps method)@\spxentry{flush\_buffers()}\spxextra{lammps.lammps method}}

\begin{fulllineitems}
\phantomsection\label{\detokenize{Python_module:lammps.lammps.flush_buffers}}
\pysigstartsignatures
\pysiglinewithargsret{\sphinxbfcode{\sphinxupquote{flush\_buffers}}}{}{}
\pysigstopsignatures
\sphinxAtStartPar
Flush output buffers

\sphinxAtStartPar
This is a wrapper around the {\hyperref[\detokenize{Library_utility:_CPPv420lammps_flush_buffersPv}]{\sphinxcrossref{\sphinxcode{\sphinxupquote{lammps\_flush\_buffers()}}}}}
function of the C\sphinxhyphen{}library interface.

\end{fulllineitems}

\index{set\_variable() (lammps.lammps method)@\spxentry{set\_variable()}\spxextra{lammps.lammps method}}

\begin{fulllineitems}
\phantomsection\label{\detokenize{Python_module:lammps.lammps.set_variable}}
\pysigstartsignatures
\pysiglinewithargsret{\sphinxbfcode{\sphinxupquote{set\_variable}}}{\sphinxparam{\DUrole{n}{name}}\sphinxparamcomma \sphinxparam{\DUrole{n}{value}}}{}
\pysigstopsignatures
\sphinxAtStartPar
Set a new value for a LAMMPS string style variable

\sphinxAtStartPar
\DUrole{versionmodified,deprecated}{Deprecated since version 7Feb2024.}

\sphinxAtStartPar
This is a wrapper around the {\hyperref[\detokenize{Library_objects:_CPPv419lammps_set_variablePvPKcPKc}]{\sphinxcrossref{\sphinxcode{\sphinxupquote{lammps\_set\_variable()}}}}}
function of the C\sphinxhyphen{}library interface.
\begin{quote}\begin{description}
\sphinxlineitem{Parameters}\begin{itemize}
\item {} 
\sphinxAtStartPar
\sphinxstyleliteralstrong{\sphinxupquote{name}} (\sphinxstyleliteralemphasis{\sphinxupquote{string}}) \textendash{} name of the variable

\item {} 
\sphinxAtStartPar
\sphinxstyleliteralstrong{\sphinxupquote{value}} (\sphinxstyleliteralemphasis{\sphinxupquote{any. will be converted to a string}}) \textendash{} new variable value

\end{itemize}

\sphinxlineitem{Returns}
\sphinxAtStartPar
either 0 on success or \sphinxhyphen{}1 on failure

\sphinxlineitem{Return type}
\sphinxAtStartPar
int

\end{description}\end{quote}

\end{fulllineitems}

\index{set\_string\_variable() (lammps.lammps method)@\spxentry{set\_string\_variable()}\spxextra{lammps.lammps method}}

\begin{fulllineitems}
\phantomsection\label{\detokenize{Python_module:lammps.lammps.set_string_variable}}
\pysigstartsignatures
\pysiglinewithargsret{\sphinxbfcode{\sphinxupquote{set\_string\_variable}}}{\sphinxparam{\DUrole{n}{name}}\sphinxparamcomma \sphinxparam{\DUrole{n}{value}}}{}
\pysigstopsignatures
\sphinxAtStartPar
Set a new value for a LAMMPS string style variable

\sphinxAtStartPar
\DUrole{versionmodified,added}{Added in version 7Feb2024.}

\sphinxAtStartPar
This is a wrapper around the {\hyperref[\detokenize{Library_objects:_CPPv426lammps_set_string_variablePvPKcPKc}]{\sphinxcrossref{\sphinxcode{\sphinxupquote{lammps\_set\_string\_variable()}}}}}
function of the C\sphinxhyphen{}library interface.
\begin{quote}\begin{description}
\sphinxlineitem{Parameters}\begin{itemize}
\item {} 
\sphinxAtStartPar
\sphinxstyleliteralstrong{\sphinxupquote{name}} (\sphinxstyleliteralemphasis{\sphinxupquote{string}}) \textendash{} name of the variable

\item {} 
\sphinxAtStartPar
\sphinxstyleliteralstrong{\sphinxupquote{value}} (\sphinxstyleliteralemphasis{\sphinxupquote{any. will be converted to a string}}) \textendash{} new variable value

\end{itemize}

\sphinxlineitem{Returns}
\sphinxAtStartPar
either 0 on success or \sphinxhyphen{}1 on failure

\sphinxlineitem{Return type}
\sphinxAtStartPar
int

\end{description}\end{quote}

\end{fulllineitems}

\index{set\_internal\_variable() (lammps.lammps method)@\spxentry{set\_internal\_variable()}\spxextra{lammps.lammps method}}

\begin{fulllineitems}
\phantomsection\label{\detokenize{Python_module:lammps.lammps.set_internal_variable}}
\pysigstartsignatures
\pysiglinewithargsret{\sphinxbfcode{\sphinxupquote{set\_internal\_variable}}}{\sphinxparam{\DUrole{n}{name}}\sphinxparamcomma \sphinxparam{\DUrole{n}{value}}}{}
\pysigstopsignatures
\sphinxAtStartPar
Set a new value for a LAMMPS internal style variable

\sphinxAtStartPar
\DUrole{versionmodified,added}{Added in version 7Feb2024.}

\sphinxAtStartPar
This is a wrapper around the {\hyperref[\detokenize{Library_objects:_CPPv428lammps_set_internal_variablePvPKcd}]{\sphinxcrossref{\sphinxcode{\sphinxupquote{lammps\_set\_internal\_variable()}}}}}
function of the C\sphinxhyphen{}library interface.
\begin{quote}\begin{description}
\sphinxlineitem{Parameters}\begin{itemize}
\item {} 
\sphinxAtStartPar
\sphinxstyleliteralstrong{\sphinxupquote{name}} (\sphinxstyleliteralemphasis{\sphinxupquote{string}}) \textendash{} name of the variable

\item {} 
\sphinxAtStartPar
\sphinxstyleliteralstrong{\sphinxupquote{value}} (\sphinxstyleliteralemphasis{\sphinxupquote{float}}\sphinxstyleliteralemphasis{\sphinxupquote{ or }}\sphinxstyleliteralemphasis{\sphinxupquote{compatible. will be converted to float}}) \textendash{} new variable value

\end{itemize}

\sphinxlineitem{Returns}
\sphinxAtStartPar
either 0 on success or \sphinxhyphen{}1 on failure

\sphinxlineitem{Return type}
\sphinxAtStartPar
int

\end{description}\end{quote}

\end{fulllineitems}

\index{gather\_bonds() (lammps.lammps method)@\spxentry{gather\_bonds()}\spxextra{lammps.lammps method}}

\begin{fulllineitems}
\phantomsection\label{\detokenize{Python_module:lammps.lammps.gather_bonds}}
\pysigstartsignatures
\pysiglinewithargsret{\sphinxbfcode{\sphinxupquote{gather\_bonds}}}{}{}
\pysigstopsignatures
\sphinxAtStartPar
Retrieve global list of bonds

\sphinxAtStartPar
\DUrole{versionmodified,added}{Added in version 28Jul2021.}

\sphinxAtStartPar
This is a wrapper around the {\hyperref[\detokenize{Library_scatter:_CPPv419lammps_gather_bondsPvPv}]{\sphinxcrossref{\sphinxcode{\sphinxupquote{lammps\_gather\_bonds()}}}}}
function of the C\sphinxhyphen{}library interface.

\sphinxAtStartPar
This function returns a tuple with the number of bonds and a
flat list of ctypes integer values with the bond type, bond atom1,
bond atom2 for each bond.
\begin{quote}\begin{description}
\sphinxlineitem{Returns}
\sphinxAtStartPar
a tuple with the number of bonds and a list of c\_int or c\_long

\sphinxlineitem{Return type}
\sphinxAtStartPar
(int, 3*nbonds*c\_tagint)

\end{description}\end{quote}

\end{fulllineitems}

\index{gather\_angles() (lammps.lammps method)@\spxentry{gather\_angles()}\spxextra{lammps.lammps method}}

\begin{fulllineitems}
\phantomsection\label{\detokenize{Python_module:lammps.lammps.gather_angles}}
\pysigstartsignatures
\pysiglinewithargsret{\sphinxbfcode{\sphinxupquote{gather\_angles}}}{}{}
\pysigstopsignatures
\sphinxAtStartPar
Retrieve global list of angles

\sphinxAtStartPar
\DUrole{versionmodified,added}{Added in version 8Feb2023.}

\sphinxAtStartPar
This is a wrapper around the {\hyperref[\detokenize{Library_scatter:_CPPv420lammps_gather_anglesPvPv}]{\sphinxcrossref{\sphinxcode{\sphinxupquote{lammps\_gather\_angles()}}}}}
function of the C\sphinxhyphen{}library interface.

\sphinxAtStartPar
This function returns a tuple with the number of angles and a
flat list of ctypes integer values with the angle type, angle atom1,
angle atom2, angle atom3 for each angle.
\begin{quote}\begin{description}
\sphinxlineitem{Returns}
\sphinxAtStartPar
a tuple with the number of angles and a list of c\_int or c\_long

\sphinxlineitem{Return type}
\sphinxAtStartPar
(int, 4*nangles*c\_tagint)

\end{description}\end{quote}

\end{fulllineitems}

\index{gather\_dihedrals() (lammps.lammps method)@\spxentry{gather\_dihedrals()}\spxextra{lammps.lammps method}}

\begin{fulllineitems}
\phantomsection\label{\detokenize{Python_module:lammps.lammps.gather_dihedrals}}
\pysigstartsignatures
\pysiglinewithargsret{\sphinxbfcode{\sphinxupquote{gather\_dihedrals}}}{}{}
\pysigstopsignatures
\sphinxAtStartPar
Retrieve global list of dihedrals

\sphinxAtStartPar
\DUrole{versionmodified,added}{Added in version 8Feb2023.}

\sphinxAtStartPar
This is a wrapper around the {\hyperref[\detokenize{Library_scatter:_CPPv423lammps_gather_dihedralsPvPv}]{\sphinxcrossref{\sphinxcode{\sphinxupquote{lammps\_gather\_dihedrals()}}}}}
function of the C\sphinxhyphen{}library interface.

\sphinxAtStartPar
This function returns a tuple with the number of dihedrals and a
flat list of ctypes integer values with the dihedral type, dihedral atom1,
dihedral atom2, dihedral atom3, dihedral atom4 for each dihedral.
\begin{quote}\begin{description}
\sphinxlineitem{Returns}
\sphinxAtStartPar
a tuple with the number of dihedrals and a list of c\_int or c\_long

\sphinxlineitem{Return type}
\sphinxAtStartPar
(int, 5*ndihedrals*c\_tagint)

\end{description}\end{quote}

\end{fulllineitems}

\index{gather\_impropers() (lammps.lammps method)@\spxentry{gather\_impropers()}\spxextra{lammps.lammps method}}

\begin{fulllineitems}
\phantomsection\label{\detokenize{Python_module:lammps.lammps.gather_impropers}}
\pysigstartsignatures
\pysiglinewithargsret{\sphinxbfcode{\sphinxupquote{gather\_impropers}}}{}{}
\pysigstopsignatures
\sphinxAtStartPar
Retrieve global list of impropers

\sphinxAtStartPar
\DUrole{versionmodified,added}{Added in version 8Feb2023.}

\sphinxAtStartPar
This is a wrapper around the {\hyperref[\detokenize{Library_scatter:_CPPv423lammps_gather_impropersPvPv}]{\sphinxcrossref{\sphinxcode{\sphinxupquote{lammps\_gather\_impropers()}}}}}
function of the C\sphinxhyphen{}library interface.

\sphinxAtStartPar
This function returns a tuple with the number of impropers and a
flat list of ctypes integer values with the improper type, improper atom1,
improper atom2, improper atom3, improper atom4 for each improper.
\begin{quote}\begin{description}
\sphinxlineitem{Returns}
\sphinxAtStartPar
a tuple with the number of impropers and a list of c\_int or c\_long

\sphinxlineitem{Return type}
\sphinxAtStartPar
(int, 5*nimpropers*c\_tagint)

\end{description}\end{quote}

\end{fulllineitems}

\index{encode\_image\_flags() (lammps.lammps method)@\spxentry{encode\_image\_flags()}\spxextra{lammps.lammps method}}

\begin{fulllineitems}
\phantomsection\label{\detokenize{Python_module:lammps.lammps.encode_image_flags}}
\pysigstartsignatures
\pysiglinewithargsret{\sphinxbfcode{\sphinxupquote{encode\_image\_flags}}}{\sphinxparam{\DUrole{n}{ix}}\sphinxparamcomma \sphinxparam{\DUrole{n}{iy}}\sphinxparamcomma \sphinxparam{\DUrole{n}{iz}}}{}
\pysigstopsignatures
\sphinxAtStartPar
convert 3 integers with image flags for x\sphinxhyphen{}, y\sphinxhyphen{}, and z\sphinxhyphen{}direction
into a single integer like it is used internally in LAMMPS

\sphinxAtStartPar
This method is a wrapper around the {\hyperref[\detokenize{Library_utility:_CPPv425lammps_encode_image_flagsiii}]{\sphinxcrossref{\sphinxcode{\sphinxupquote{lammps\_encode\_image\_flags()}}}}}
function of library interface.
\begin{quote}\begin{description}
\sphinxlineitem{Parameters}\begin{itemize}
\item {} 
\sphinxAtStartPar
\sphinxstyleliteralstrong{\sphinxupquote{ix}} (\sphinxstyleliteralemphasis{\sphinxupquote{int}}) \textendash{} x\sphinxhyphen{}direction image flag

\item {} 
\sphinxAtStartPar
\sphinxstyleliteralstrong{\sphinxupquote{iy}} (\sphinxstyleliteralemphasis{\sphinxupquote{int}}) \textendash{} y\sphinxhyphen{}direction image flag

\item {} 
\sphinxAtStartPar
\sphinxstyleliteralstrong{\sphinxupquote{iz}} (\sphinxstyleliteralemphasis{\sphinxupquote{int}}) \textendash{} z\sphinxhyphen{}direction image flag

\end{itemize}

\sphinxlineitem{Returns}
\sphinxAtStartPar
encoded image flags

\sphinxlineitem{Return type}
\sphinxAtStartPar
lammps.c\_imageint

\end{description}\end{quote}

\end{fulllineitems}

\index{decode\_image\_flags() (lammps.lammps method)@\spxentry{decode\_image\_flags()}\spxextra{lammps.lammps method}}

\begin{fulllineitems}
\phantomsection\label{\detokenize{Python_module:lammps.lammps.decode_image_flags}}
\pysigstartsignatures
\pysiglinewithargsret{\sphinxbfcode{\sphinxupquote{decode\_image\_flags}}}{\sphinxparam{\DUrole{n}{image}}}{}
\pysigstopsignatures
\sphinxAtStartPar
Convert encoded image flag integer into list of three regular integers.

\sphinxAtStartPar
This method is a wrapper around the {\hyperref[\detokenize{Library_utility:_CPPv425lammps_decode_image_flagsiPi}]{\sphinxcrossref{\sphinxcode{\sphinxupquote{lammps\_decode\_image\_flags()}}}}}
function of library interface.
\begin{quote}\begin{description}
\sphinxlineitem{Parameters}
\sphinxAtStartPar
\sphinxstyleliteralstrong{\sphinxupquote{image}} (\sphinxstyleliteralemphasis{\sphinxupquote{lammps.c\_imageint}}) \textendash{} encoded image flags

\sphinxlineitem{Returns}
\sphinxAtStartPar
list of three image flags in x\sphinxhyphen{}, y\sphinxhyphen{}, and z\sphinxhyphen{} direction

\sphinxlineitem{Return type}
\sphinxAtStartPar
list of 3 int

\end{description}\end{quote}

\end{fulllineitems}

\index{create\_atoms() (lammps.lammps method)@\spxentry{create\_atoms()}\spxextra{lammps.lammps method}}

\begin{fulllineitems}
\phantomsection\label{\detokenize{Python_module:lammps.lammps.create_atoms}}
\pysigstartsignatures
\pysiglinewithargsret{\sphinxbfcode{\sphinxupquote{create\_atoms}}}{\sphinxparam{\DUrole{n}{n}}\sphinxparamcomma \sphinxparam{\DUrole{n}{id}}\sphinxparamcomma \sphinxparam{\DUrole{n}{type}}\sphinxparamcomma \sphinxparam{\DUrole{n}{x}}\sphinxparamcomma \sphinxparam{\DUrole{n}{v}\DUrole{o}{=}\DUrole{default_value}{None}}\sphinxparamcomma \sphinxparam{\DUrole{n}{image}\DUrole{o}{=}\DUrole{default_value}{None}}\sphinxparamcomma \sphinxparam{\DUrole{n}{shrinkexceed}\DUrole{o}{=}\DUrole{default_value}{False}}}{}
\pysigstopsignatures
\sphinxAtStartPar
Create N atoms from list of coordinates and properties

\sphinxAtStartPar
This function is a wrapper around the {\hyperref[\detokenize{Library_scatter:_CPPv419lammps_create_atomsPviPKiPKiPKdPKdPKii}]{\sphinxcrossref{\sphinxcode{\sphinxupquote{lammps\_create\_atoms()}}}}}
function of the C\sphinxhyphen{}library interface, and the behavior is similar except
that the \sphinxstyleemphasis{v}, \sphinxstyleemphasis{image}, and \sphinxstyleemphasis{shrinkexceed} arguments are optional and
default to \sphinxstyleemphasis{None}, \sphinxstyleemphasis{None}, and \sphinxstyleemphasis{False}, respectively. With \sphinxstyleemphasis{None} being
equivalent to a \sphinxcode{\sphinxupquote{NULL}} pointer in C.

\sphinxAtStartPar
The lists of coordinates, types, atom IDs, velocities, image flags can
be provided in any format that may be converted into the required
internal data types.  Also the list may contain more than \sphinxstyleemphasis{N} entries,
but not fewer.  In the latter case, the function will return without
attempting to create atoms.  You may use the {\hyperref[\detokenize{Python_module:lammps.lammps.encode_image_flags}]{\sphinxcrossref{\sphinxcode{\sphinxupquote{encode\_image\_flags}}}}} method to properly combine three integers
with image flags into a single integer.
\begin{quote}\begin{description}
\sphinxlineitem{Parameters}\begin{itemize}
\item {} 
\sphinxAtStartPar
\sphinxstyleliteralstrong{\sphinxupquote{n}} (\sphinxstyleliteralemphasis{\sphinxupquote{int}}) \textendash{} number of atoms for which data is provided

\item {} 
\sphinxAtStartPar
\sphinxstyleliteralstrong{\sphinxupquote{id}} (\sphinxstyleliteralemphasis{\sphinxupquote{list}}\sphinxstyleliteralemphasis{\sphinxupquote{ of }}\sphinxstyleliteralemphasis{\sphinxupquote{lammps.tagint}}) \textendash{} list of atom IDs with at least n elements or None

\item {} 
\sphinxAtStartPar
\sphinxstyleliteralstrong{\sphinxupquote{type}} (\sphinxstyleliteralemphasis{\sphinxupquote{list}}\sphinxstyleliteralemphasis{\sphinxupquote{ of }}\sphinxstyleliteralemphasis{\sphinxupquote{int}}) \textendash{} list of atom types

\item {} 
\sphinxAtStartPar
\sphinxstyleliteralstrong{\sphinxupquote{x}} (\sphinxstyleliteralemphasis{\sphinxupquote{list}}\sphinxstyleliteralemphasis{\sphinxupquote{ of }}\sphinxstyleliteralemphasis{\sphinxupquote{float}}) \textendash{} list of coordinates for x\sphinxhyphen{}, y\sphinxhyphen{}, and z (flat list of 3n entries)

\item {} 
\sphinxAtStartPar
\sphinxstyleliteralstrong{\sphinxupquote{v}} (\sphinxstyleliteralemphasis{\sphinxupquote{list}}\sphinxstyleliteralemphasis{\sphinxupquote{ of }}\sphinxstyleliteralemphasis{\sphinxupquote{float}}) \textendash{} list of velocities for x\sphinxhyphen{}, y\sphinxhyphen{}, and z (flat list of 3n entries) or None (optional)

\item {} 
\sphinxAtStartPar
\sphinxstyleliteralstrong{\sphinxupquote{image}} (\sphinxstyleliteralemphasis{\sphinxupquote{list}}\sphinxstyleliteralemphasis{\sphinxupquote{ of }}\sphinxstyleliteralemphasis{\sphinxupquote{lammps.imageint}}) \textendash{} list of encoded image flags (optional)

\item {} 
\sphinxAtStartPar
\sphinxstyleliteralstrong{\sphinxupquote{shrinkexceed}} (\sphinxstyleliteralemphasis{\sphinxupquote{bool}}) \textendash{} whether to expand shrink\sphinxhyphen{}wrap boundaries if atoms are outside the box (optional)

\end{itemize}

\sphinxlineitem{Returns}
\sphinxAtStartPar
number of atoms created. 0 if insufficient or invalid data

\sphinxlineitem{Return type}
\sphinxAtStartPar
int

\end{description}\end{quote}

\end{fulllineitems}

\index{has\_mpi\_support (lammps.lammps property)@\spxentry{has\_mpi\_support}\spxextra{lammps.lammps property}}

\begin{fulllineitems}
\phantomsection\label{\detokenize{Python_module:lammps.lammps.has_mpi_support}}
\pysigstartsignatures
\pysigline{\sphinxbfcode{\sphinxupquote{property\DUrole{w}{ }}}\sphinxbfcode{\sphinxupquote{has\_mpi\_support}}}
\pysigstopsignatures
\sphinxAtStartPar
Report whether the LAMMPS shared library was compiled with a
real MPI library or in serial.

\sphinxAtStartPar
This is a wrapper around the {\hyperref[\detokenize{Library_config:_CPPv429lammps_config_has_mpi_supportv}]{\sphinxcrossref{\sphinxcode{\sphinxupquote{lammps\_config\_has\_mpi\_support()}}}}}
function of the library interface.
\begin{quote}\begin{description}
\sphinxlineitem{Returns}
\sphinxAtStartPar
False when compiled with MPI STUBS, otherwise True

\sphinxlineitem{Return type}
\sphinxAtStartPar
bool

\end{description}\end{quote}

\end{fulllineitems}

\index{is\_running (lammps.lammps property)@\spxentry{is\_running}\spxextra{lammps.lammps property}}

\begin{fulllineitems}
\phantomsection\label{\detokenize{Python_module:lammps.lammps.is_running}}
\pysigstartsignatures
\pysigline{\sphinxbfcode{\sphinxupquote{property\DUrole{w}{ }}}\sphinxbfcode{\sphinxupquote{is\_running}}}
\pysigstopsignatures
\sphinxAtStartPar
Report whether being called from a function during a run or a minimization

\sphinxAtStartPar
\DUrole{versionmodified,added}{Added in version 9Oct2020.}

\sphinxAtStartPar
Various LAMMPS commands must not be called during an ongoing
run or minimization.  This property allows to check for that.
This is a wrapper around the {\hyperref[\detokenize{Library_utility:_CPPv417lammps_is_runningPv}]{\sphinxcrossref{\sphinxcode{\sphinxupquote{lammps\_is\_running()}}}}}
function of the library interface.
\begin{quote}\begin{description}
\sphinxlineitem{Returns}
\sphinxAtStartPar
True when called during a run otherwise false

\sphinxlineitem{Return type}
\sphinxAtStartPar
bool

\end{description}\end{quote}

\end{fulllineitems}

\index{force\_timeout() (lammps.lammps method)@\spxentry{force\_timeout()}\spxextra{lammps.lammps method}}

\begin{fulllineitems}
\phantomsection\label{\detokenize{Python_module:lammps.lammps.force_timeout}}
\pysigstartsignatures
\pysiglinewithargsret{\sphinxbfcode{\sphinxupquote{force\_timeout}}}{}{}
\pysigstopsignatures
\sphinxAtStartPar
Trigger an immediate timeout, i.e. a “soft stop” of a run.

\sphinxAtStartPar
\DUrole{versionmodified,added}{Added in version 9Oct2020.}

\sphinxAtStartPar
This function allows to cleanly stop an ongoing run or minimization
at the next loop iteration.
This is a wrapper around the {\hyperref[\detokenize{Library_utility:_CPPv420lammps_force_timeoutPv}]{\sphinxcrossref{\sphinxcode{\sphinxupquote{lammps\_force\_timeout()}}}}}
function of the library interface.

\end{fulllineitems}

\index{has\_exceptions (lammps.lammps property)@\spxentry{has\_exceptions}\spxextra{lammps.lammps property}}

\begin{fulllineitems}
\phantomsection\label{\detokenize{Python_module:lammps.lammps.has_exceptions}}
\pysigstartsignatures
\pysigline{\sphinxbfcode{\sphinxupquote{property\DUrole{w}{ }}}\sphinxbfcode{\sphinxupquote{has\_exceptions}}}
\pysigstopsignatures
\sphinxAtStartPar
Report whether the LAMMPS shared library was compiled with C++
exceptions handling enabled

\sphinxAtStartPar
This is a wrapper around the {\hyperref[\detokenize{Library_config:_CPPv428lammps_config_has_exceptionsv}]{\sphinxcrossref{\sphinxcode{\sphinxupquote{lammps\_config\_has\_exceptions()}}}}}
function of the library interface.
\begin{quote}\begin{description}
\sphinxlineitem{Returns}
\sphinxAtStartPar
state of C++ exception support

\sphinxlineitem{Return type}
\sphinxAtStartPar
bool

\end{description}\end{quote}

\end{fulllineitems}

\index{has\_gzip\_support (lammps.lammps property)@\spxentry{has\_gzip\_support}\spxextra{lammps.lammps property}}

\begin{fulllineitems}
\phantomsection\label{\detokenize{Python_module:lammps.lammps.has_gzip_support}}
\pysigstartsignatures
\pysigline{\sphinxbfcode{\sphinxupquote{property\DUrole{w}{ }}}\sphinxbfcode{\sphinxupquote{has\_gzip\_support}}}
\pysigstopsignatures
\sphinxAtStartPar
Report whether the LAMMPS shared library was compiled with support
for reading and writing compressed files through \sphinxcode{\sphinxupquote{gzip}}.

\sphinxAtStartPar
This is a wrapper around the {\hyperref[\detokenize{Library_config:_CPPv430lammps_config_has_gzip_supportv}]{\sphinxcrossref{\sphinxcode{\sphinxupquote{lammps\_config\_has\_gzip\_support()}}}}}
function of the library interface.
\begin{quote}\begin{description}
\sphinxlineitem{Returns}
\sphinxAtStartPar
state of gzip support

\sphinxlineitem{Return type}
\sphinxAtStartPar
bool

\end{description}\end{quote}

\end{fulllineitems}

\index{has\_png\_support (lammps.lammps property)@\spxentry{has\_png\_support}\spxextra{lammps.lammps property}}

\begin{fulllineitems}
\phantomsection\label{\detokenize{Python_module:lammps.lammps.has_png_support}}
\pysigstartsignatures
\pysigline{\sphinxbfcode{\sphinxupquote{property\DUrole{w}{ }}}\sphinxbfcode{\sphinxupquote{has\_png\_support}}}
\pysigstopsignatures
\sphinxAtStartPar
Report whether the LAMMPS shared library was compiled with support
for writing images in PNG format.

\sphinxAtStartPar
This is a wrapper around the {\hyperref[\detokenize{Library_config:_CPPv429lammps_config_has_png_supportv}]{\sphinxcrossref{\sphinxcode{\sphinxupquote{lammps\_config\_has\_png\_support()}}}}}
function of the library interface.
\begin{quote}\begin{description}
\sphinxlineitem{Returns}
\sphinxAtStartPar
state of PNG support

\sphinxlineitem{Return type}
\sphinxAtStartPar
bool

\end{description}\end{quote}

\end{fulllineitems}

\index{has\_jpeg\_support (lammps.lammps property)@\spxentry{has\_jpeg\_support}\spxextra{lammps.lammps property}}

\begin{fulllineitems}
\phantomsection\label{\detokenize{Python_module:lammps.lammps.has_jpeg_support}}
\pysigstartsignatures
\pysigline{\sphinxbfcode{\sphinxupquote{property\DUrole{w}{ }}}\sphinxbfcode{\sphinxupquote{has\_jpeg\_support}}}
\pysigstopsignatures
\sphinxAtStartPar
Report whether the LAMMPS shared library was compiled with support
for writing images in JPEG format.

\sphinxAtStartPar
This is a wrapper around the {\hyperref[\detokenize{Library_config:_CPPv430lammps_config_has_jpeg_supportv}]{\sphinxcrossref{\sphinxcode{\sphinxupquote{lammps\_config\_has\_jpeg\_support()}}}}}
function of the library interface.
\begin{quote}\begin{description}
\sphinxlineitem{Returns}
\sphinxAtStartPar
state of JPEG support

\sphinxlineitem{Return type}
\sphinxAtStartPar
bool

\end{description}\end{quote}

\end{fulllineitems}

\index{has\_ffmpeg\_support (lammps.lammps property)@\spxentry{has\_ffmpeg\_support}\spxextra{lammps.lammps property}}

\begin{fulllineitems}
\phantomsection\label{\detokenize{Python_module:lammps.lammps.has_ffmpeg_support}}
\pysigstartsignatures
\pysigline{\sphinxbfcode{\sphinxupquote{property\DUrole{w}{ }}}\sphinxbfcode{\sphinxupquote{has\_ffmpeg\_support}}}
\pysigstopsignatures
\sphinxAtStartPar
State of support for writing movies with \sphinxcode{\sphinxupquote{ffmpeg}} in the LAMMPS shared library

\sphinxAtStartPar
This is a wrapper around the {\hyperref[\detokenize{Library_config:_CPPv432lammps_config_has_ffmpeg_supportv}]{\sphinxcrossref{\sphinxcode{\sphinxupquote{lammps\_config\_has\_ffmpeg\_support()}}}}}
function of the library interface.
\begin{quote}\begin{description}
\sphinxlineitem{Returns}
\sphinxAtStartPar
state of ffmpeg support

\sphinxlineitem{Return type}
\sphinxAtStartPar
bool

\end{description}\end{quote}

\end{fulllineitems}

\index{has\_curl\_support (lammps.lammps property)@\spxentry{has\_curl\_support}\spxextra{lammps.lammps property}}

\begin{fulllineitems}
\phantomsection\label{\detokenize{Python_module:lammps.lammps.has_curl_support}}
\pysigstartsignatures
\pysigline{\sphinxbfcode{\sphinxupquote{property\DUrole{w}{ }}}\sphinxbfcode{\sphinxupquote{has\_curl\_support}}}
\pysigstopsignatures
\sphinxAtStartPar
Report whether the LAMMPS shared library was compiled with support
for downloading files through libcurl.

\sphinxAtStartPar
This is a wrapper around the \sphinxcode{\sphinxupquote{lammps\_config\_has\_curl\_support()}}
function of the library interface.
\begin{quote}\begin{description}
\sphinxlineitem{Returns}
\sphinxAtStartPar
state of CURL support

\sphinxlineitem{Return type}
\sphinxAtStartPar
bool

\end{description}\end{quote}

\end{fulllineitems}

\index{has\_package() (lammps.lammps method)@\spxentry{has\_package()}\spxextra{lammps.lammps method}}

\begin{fulllineitems}
\phantomsection\label{\detokenize{Python_module:lammps.lammps.has_package}}
\pysigstartsignatures
\pysiglinewithargsret{\sphinxbfcode{\sphinxupquote{has\_package}}}{\sphinxparam{\DUrole{n}{name}}}{}
\pysigstopsignatures
\sphinxAtStartPar
Report if the named package has been enabled in the LAMMPS shared library.

\sphinxAtStartPar
\DUrole{versionmodified,added}{Added in version 3Nov2022.}

\sphinxAtStartPar
This is a wrapper around the {\hyperref[\detokenize{Library_config:_CPPv425lammps_config_has_packagePKc}]{\sphinxcrossref{\sphinxcode{\sphinxupquote{lammps\_config\_has\_package()}}}}}
function of the library interface.
\begin{quote}\begin{description}
\sphinxlineitem{Parameters}
\sphinxAtStartPar
\sphinxstyleliteralstrong{\sphinxupquote{name}} (\sphinxstyleliteralemphasis{\sphinxupquote{string}}) \textendash{} name of the package

\sphinxlineitem{Returns}
\sphinxAtStartPar
state of package availability

\sphinxlineitem{Return type}
\sphinxAtStartPar
bool

\end{description}\end{quote}

\end{fulllineitems}

\index{accelerator\_config (lammps.lammps property)@\spxentry{accelerator\_config}\spxextra{lammps.lammps property}}

\begin{fulllineitems}
\phantomsection\label{\detokenize{Python_module:lammps.lammps.accelerator_config}}
\pysigstartsignatures
\pysigline{\sphinxbfcode{\sphinxupquote{property\DUrole{w}{ }}}\sphinxbfcode{\sphinxupquote{accelerator\_config}}}
\pysigstopsignatures
\sphinxAtStartPar
Return table with available accelerator configuration settings.

\sphinxAtStartPar
This is a wrapper around the {\hyperref[\detokenize{Library_config:_CPPv425lammps_config_acceleratorPKcPKcPKc}]{\sphinxcrossref{\sphinxcode{\sphinxupquote{lammps\_config\_accelerator()}}}}}
function of the library interface which loops over all known packages
and categories and returns enabled features as a nested dictionary
with all enabled settings as list of strings.
\begin{quote}\begin{description}
\sphinxlineitem{Returns}
\sphinxAtStartPar
nested dictionary with all known enabled settings as list of strings

\sphinxlineitem{Return type}
\sphinxAtStartPar
dictionary

\end{description}\end{quote}

\end{fulllineitems}

\index{has\_gpu\_device (lammps.lammps property)@\spxentry{has\_gpu\_device}\spxextra{lammps.lammps property}}

\begin{fulllineitems}
\phantomsection\label{\detokenize{Python_module:lammps.lammps.has_gpu_device}}
\pysigstartsignatures
\pysigline{\sphinxbfcode{\sphinxupquote{property\DUrole{w}{ }}}\sphinxbfcode{\sphinxupquote{has\_gpu\_device}}}
\pysigstopsignatures
\sphinxAtStartPar
Availability of GPU package compatible device

\sphinxAtStartPar
This is a wrapper around the {\hyperref[\detokenize{Library_config:_CPPv421lammps_has_gpu_devicev}]{\sphinxcrossref{\sphinxcode{\sphinxupquote{lammps\_has\_gpu\_device()}}}}}
function of the C library interface.
\begin{quote}\begin{description}
\sphinxlineitem{Returns}
\sphinxAtStartPar
True if a GPU package compatible device is present, otherwise False

\sphinxlineitem{Return type}
\sphinxAtStartPar
bool

\end{description}\end{quote}

\end{fulllineitems}

\index{get\_gpu\_device\_info() (lammps.lammps method)@\spxentry{get\_gpu\_device\_info()}\spxextra{lammps.lammps method}}

\begin{fulllineitems}
\phantomsection\label{\detokenize{Python_module:lammps.lammps.get_gpu_device_info}}
\pysigstartsignatures
\pysiglinewithargsret{\sphinxbfcode{\sphinxupquote{get\_gpu\_device\_info}}}{}{}
\pysigstopsignatures
\sphinxAtStartPar
Return a string with detailed information about any devices that are
usable by the GPU package.

\sphinxAtStartPar
This is a wrapper around the {\hyperref[\detokenize{Library_config:_CPPv426lammps_get_gpu_device_infoPci}]{\sphinxcrossref{\sphinxcode{\sphinxupquote{lammps\_get\_gpu\_device\_info()}}}}}
function of the C\sphinxhyphen{}library interface.
\begin{quote}\begin{description}
\sphinxlineitem{Returns}
\sphinxAtStartPar
GPU device info string

\sphinxlineitem{Return type}
\sphinxAtStartPar
string

\end{description}\end{quote}

\end{fulllineitems}

\index{installed\_packages (lammps.lammps property)@\spxentry{installed\_packages}\spxextra{lammps.lammps property}}

\begin{fulllineitems}
\phantomsection\label{\detokenize{Python_module:lammps.lammps.installed_packages}}
\pysigstartsignatures
\pysigline{\sphinxbfcode{\sphinxupquote{property\DUrole{w}{ }}}\sphinxbfcode{\sphinxupquote{installed\_packages}}}
\pysigstopsignatures
\sphinxAtStartPar
List of the names of enabled packages in the LAMMPS shared library

\sphinxAtStartPar
This is a wrapper around the functions {\hyperref[\detokenize{Library_config:_CPPv427lammps_config_package_countv}]{\sphinxcrossref{\sphinxcode{\sphinxupquote{lammps\_config\_package\_count()}}}}}
and {\hyperref[\detokenize{Library_config:_CPPv426lammps_config_package_nameiPci}]{\sphinxcrossref{\sphinxcode{\sphinxupquote{lammps\_config\_package\_name()}}}}} of the library interface.

\sphinxAtStartPar
:return

\end{fulllineitems}

\index{has\_style() (lammps.lammps method)@\spxentry{has\_style()}\spxextra{lammps.lammps method}}

\begin{fulllineitems}
\phantomsection\label{\detokenize{Python_module:lammps.lammps.has_style}}
\pysigstartsignatures
\pysiglinewithargsret{\sphinxbfcode{\sphinxupquote{has\_style}}}{\sphinxparam{\DUrole{n}{category}}\sphinxparamcomma \sphinxparam{\DUrole{n}{name}}}{}
\pysigstopsignatures
\sphinxAtStartPar
Returns whether a given style name is available in a given category

\sphinxAtStartPar
This is a wrapper around the function {\hyperref[\detokenize{Library_config:_CPPv416lammps_has_stylePvPKcPKc}]{\sphinxcrossref{\sphinxcode{\sphinxupquote{lammps\_has\_style()}}}}}
of the library interface.
\begin{quote}\begin{description}
\sphinxlineitem{Parameters}\begin{itemize}
\item {} 
\sphinxAtStartPar
\sphinxstyleliteralstrong{\sphinxupquote{category}} (\sphinxstyleliteralemphasis{\sphinxupquote{string}}) \textendash{} name of category

\item {} 
\sphinxAtStartPar
\sphinxstyleliteralstrong{\sphinxupquote{name}} (\sphinxstyleliteralemphasis{\sphinxupquote{string}}) \textendash{} name of the style

\end{itemize}

\sphinxlineitem{Returns}
\sphinxAtStartPar
true if style is available in given category

\sphinxlineitem{Return type}
\sphinxAtStartPar
bool

\end{description}\end{quote}

\end{fulllineitems}

\index{available\_styles() (lammps.lammps method)@\spxentry{available\_styles()}\spxextra{lammps.lammps method}}

\begin{fulllineitems}
\phantomsection\label{\detokenize{Python_module:lammps.lammps.available_styles}}
\pysigstartsignatures
\pysiglinewithargsret{\sphinxbfcode{\sphinxupquote{available\_styles}}}{\sphinxparam{\DUrole{n}{category}}}{}
\pysigstopsignatures
\sphinxAtStartPar
Returns a list of styles available for a given category

\sphinxAtStartPar
This is a wrapper around the functions {\hyperref[\detokenize{Library_config:_CPPv418lammps_style_countPvPKc}]{\sphinxcrossref{\sphinxcode{\sphinxupquote{lammps\_style\_count()}}}}}
and {\hyperref[\detokenize{Library_config:_CPPv417lammps_style_namePvPKciPci}]{\sphinxcrossref{\sphinxcode{\sphinxupquote{lammps\_style\_name()}}}}} of the library interface.
\begin{quote}\begin{description}
\sphinxlineitem{Parameters}
\sphinxAtStartPar
\sphinxstyleliteralstrong{\sphinxupquote{category}} (\sphinxstyleliteralemphasis{\sphinxupquote{string}}) \textendash{} name of category

\sphinxlineitem{Returns}
\sphinxAtStartPar
list of style names in given category

\sphinxlineitem{Return type}
\sphinxAtStartPar
list

\end{description}\end{quote}

\end{fulllineitems}

\index{has\_id() (lammps.lammps method)@\spxentry{has\_id()}\spxextra{lammps.lammps method}}

\begin{fulllineitems}
\phantomsection\label{\detokenize{Python_module:lammps.lammps.has_id}}
\pysigstartsignatures
\pysiglinewithargsret{\sphinxbfcode{\sphinxupquote{has\_id}}}{\sphinxparam{\DUrole{n}{category}}\sphinxparamcomma \sphinxparam{\DUrole{n}{name}}}{}
\pysigstopsignatures
\sphinxAtStartPar
Returns whether a given ID name is available in a given category

\sphinxAtStartPar
\DUrole{versionmodified,added}{Added in version 9Oct2020.}

\sphinxAtStartPar
This is a wrapper around the function {\hyperref[\detokenize{Library_config:_CPPv413lammps_has_idPvPKcPKc}]{\sphinxcrossref{\sphinxcode{\sphinxupquote{lammps\_has\_id()}}}}}
of the library interface.
\begin{quote}\begin{description}
\sphinxlineitem{Parameters}\begin{itemize}
\item {} 
\sphinxAtStartPar
\sphinxstyleliteralstrong{\sphinxupquote{category}} (\sphinxstyleliteralemphasis{\sphinxupquote{string}}) \textendash{} name of category

\item {} 
\sphinxAtStartPar
\sphinxstyleliteralstrong{\sphinxupquote{name}} (\sphinxstyleliteralemphasis{\sphinxupquote{string}}) \textendash{} name of the ID

\end{itemize}

\sphinxlineitem{Returns}
\sphinxAtStartPar
true if ID is available in given category

\sphinxlineitem{Return type}
\sphinxAtStartPar
bool

\end{description}\end{quote}

\end{fulllineitems}

\index{available\_ids() (lammps.lammps method)@\spxentry{available\_ids()}\spxextra{lammps.lammps method}}

\begin{fulllineitems}
\phantomsection\label{\detokenize{Python_module:lammps.lammps.available_ids}}
\pysigstartsignatures
\pysiglinewithargsret{\sphinxbfcode{\sphinxupquote{available\_ids}}}{\sphinxparam{\DUrole{n}{category}}}{}
\pysigstopsignatures
\sphinxAtStartPar
Returns a list of IDs available for a given category

\sphinxAtStartPar
\DUrole{versionmodified,added}{Added in version 9Oct2020.}

\sphinxAtStartPar
This is a wrapper around the functions {\hyperref[\detokenize{Library_config:_CPPv415lammps_id_countPvPKc}]{\sphinxcrossref{\sphinxcode{\sphinxupquote{lammps\_id\_count()}}}}}
and {\hyperref[\detokenize{Library_config:_CPPv414lammps_id_namePvPKciPci}]{\sphinxcrossref{\sphinxcode{\sphinxupquote{lammps\_id\_name()}}}}} of the library interface.
\begin{quote}\begin{description}
\sphinxlineitem{Parameters}
\sphinxAtStartPar
\sphinxstyleliteralstrong{\sphinxupquote{category}} (\sphinxstyleliteralemphasis{\sphinxupquote{string}}) \textendash{} name of category

\sphinxlineitem{Returns}
\sphinxAtStartPar
list of id names in given category

\sphinxlineitem{Return type}
\sphinxAtStartPar
list

\end{description}\end{quote}

\end{fulllineitems}

\index{available\_plugins() (lammps.lammps method)@\spxentry{available\_plugins()}\spxextra{lammps.lammps method}}

\begin{fulllineitems}
\phantomsection\label{\detokenize{Python_module:lammps.lammps.available_plugins}}
\pysigstartsignatures
\pysiglinewithargsret{\sphinxbfcode{\sphinxupquote{available\_plugins}}}{\sphinxparam{\DUrole{n}{category}}}{}
\pysigstopsignatures
\sphinxAtStartPar
Returns a list of plugins available for a given category

\sphinxAtStartPar
\DUrole{versionmodified,added}{Added in version 10Mar2021.}

\sphinxAtStartPar
This is a wrapper around the functions \sphinxcode{\sphinxupquote{lammps\_plugin\_count()}}
and \sphinxcode{\sphinxupquote{lammps\_plugin\_name()}} of the library interface.
\begin{quote}\begin{description}
\sphinxlineitem{Returns}
\sphinxAtStartPar
list of style/name pairs of loaded plugins

\sphinxlineitem{Return type}
\sphinxAtStartPar
list

\end{description}\end{quote}

\end{fulllineitems}

\index{set\_fix\_external\_callback() (lammps.lammps method)@\spxentry{set\_fix\_external\_callback()}\spxextra{lammps.lammps method}}

\begin{fulllineitems}
\phantomsection\label{\detokenize{Python_module:lammps.lammps.set_fix_external_callback}}
\pysigstartsignatures
\pysiglinewithargsret{\sphinxbfcode{\sphinxupquote{set\_fix\_external\_callback}}}{\sphinxparam{\DUrole{n}{fix\_id}}\sphinxparamcomma \sphinxparam{\DUrole{n}{callback}}\sphinxparamcomma \sphinxparam{\DUrole{n}{caller}\DUrole{o}{=}\DUrole{default_value}{None}}}{}
\pysigstopsignatures
\sphinxAtStartPar
Set the callback function for a fix external instance with a given fix ID.

\sphinxAtStartPar
Optionally also set a reference to the calling object.

\sphinxAtStartPar
This is a wrapper around the {\hyperref[\detokenize{Library_utility:_CPPv432lammps_set_fix_external_callbackPvPKc16FixExternalFnPtrPv}]{\sphinxcrossref{\sphinxcode{\sphinxupquote{lammps\_set\_fix\_external\_callback()}}}}} function
of the C\sphinxhyphen{}library interface.  However this is set up to call a Python function with
the following arguments.
\begin{itemize}
\item {} 
\sphinxAtStartPar
object is the value of the “caller” argument

\item {} 
\sphinxAtStartPar
ntimestep is the current timestep

\item {} 
\sphinxAtStartPar
nlocal is the number of local atoms on the current MPI process

\item {} 
\sphinxAtStartPar
tag is a 1d NumPy array of integers representing the atom IDs of the local atoms

\item {} 
\sphinxAtStartPar
x is a 2d NumPy array of doubles of the coordinates of the local atoms

\item {} 
\sphinxAtStartPar
f is a 2d NumPy array of doubles of the forces on the local atoms that will be added

\end{itemize}

\sphinxAtStartPar
\DUrole{versionmodified,changed}{Changed in version 28Jul2021.}
\begin{quote}\begin{description}
\sphinxlineitem{Parameters}\begin{itemize}
\item {} 
\sphinxAtStartPar
\sphinxstyleliteralstrong{\sphinxupquote{fix\_id}} \textendash{} Fix\sphinxhyphen{}ID of a fix external instance

\item {} 
\sphinxAtStartPar
\sphinxstyleliteralstrong{\sphinxupquote{callback}} \textendash{} Python function that will be called from fix external

\item {} 
\sphinxAtStartPar
\sphinxstyleliteralstrong{\sphinxupquote{caller}} \textendash{} reference to some object passed to the callback function

\end{itemize}

\sphinxlineitem{Type}
\sphinxAtStartPar
string

\sphinxlineitem{Type}
\sphinxAtStartPar
function

\sphinxlineitem{Type}
\sphinxAtStartPar
object, optional

\end{description}\end{quote}

\end{fulllineitems}

\index{fix\_external\_get\_force() (lammps.lammps method)@\spxentry{fix\_external\_get\_force()}\spxextra{lammps.lammps method}}

\begin{fulllineitems}
\phantomsection\label{\detokenize{Python_module:lammps.lammps.fix_external_get_force}}
\pysigstartsignatures
\pysiglinewithargsret{\sphinxbfcode{\sphinxupquote{fix\_external\_get\_force}}}{\sphinxparam{\DUrole{n}{fix\_id}}}{}
\pysigstopsignatures
\sphinxAtStartPar
Get access to the array with per\sphinxhyphen{}atom forces of a fix external instance with a given fix ID.

\sphinxAtStartPar
\DUrole{versionmodified,added}{Added in version 28Jul2021.}

\sphinxAtStartPar
This is a wrapper around the {\hyperref[\detokenize{Library_utility:_CPPv429lammps_fix_external_get_forcePvPKc}]{\sphinxcrossref{\sphinxcode{\sphinxupquote{lammps\_fix\_external\_get\_force()}}}}} function
of the C\sphinxhyphen{}library interface.
\begin{quote}\begin{description}
\sphinxlineitem{Parameters}
\sphinxAtStartPar
\sphinxstyleliteralstrong{\sphinxupquote{fix\_id}} \textendash{} Fix\sphinxhyphen{}ID of a fix external instance

\sphinxlineitem{Type}
\sphinxAtStartPar
string

\sphinxlineitem{Returns}
\sphinxAtStartPar
requested data

\sphinxlineitem{Return type}
\sphinxAtStartPar
ctypes.POINTER(ctypes.POINTER(ctypes.double))

\end{description}\end{quote}

\end{fulllineitems}

\index{fix\_external\_set\_energy\_global() (lammps.lammps method)@\spxentry{fix\_external\_set\_energy\_global()}\spxextra{lammps.lammps method}}

\begin{fulllineitems}
\phantomsection\label{\detokenize{Python_module:lammps.lammps.fix_external_set_energy_global}}
\pysigstartsignatures
\pysiglinewithargsret{\sphinxbfcode{\sphinxupquote{fix\_external\_set\_energy\_global}}}{\sphinxparam{\DUrole{n}{fix\_id}}\sphinxparamcomma \sphinxparam{\DUrole{n}{eng}}}{}
\pysigstopsignatures
\sphinxAtStartPar
Set the global energy contribution for a fix external instance with the given ID.

\sphinxAtStartPar
\DUrole{versionmodified,added}{Added in version 28Jul2021.}

\sphinxAtStartPar
This is a wrapper around the {\hyperref[\detokenize{Library_utility:_CPPv437lammps_fix_external_set_energy_globalPvPKcd}]{\sphinxcrossref{\sphinxcode{\sphinxupquote{lammps\_fix\_external\_set\_energy\_global()}}}}} function
of the C\sphinxhyphen{}library interface.
\begin{quote}\begin{description}
\sphinxlineitem{Parameters}\begin{itemize}
\item {} 
\sphinxAtStartPar
\sphinxstyleliteralstrong{\sphinxupquote{fix\_id}} \textendash{} Fix\sphinxhyphen{}ID of a fix external instance

\item {} 
\sphinxAtStartPar
\sphinxstyleliteralstrong{\sphinxupquote{eng}} \textendash{} potential energy value to be added by fix external

\end{itemize}

\sphinxlineitem{Type}
\sphinxAtStartPar
string

\sphinxlineitem{Type}
\sphinxAtStartPar
float

\end{description}\end{quote}

\end{fulllineitems}

\index{fix\_external\_set\_virial\_global() (lammps.lammps method)@\spxentry{fix\_external\_set\_virial\_global()}\spxextra{lammps.lammps method}}

\begin{fulllineitems}
\phantomsection\label{\detokenize{Python_module:lammps.lammps.fix_external_set_virial_global}}
\pysigstartsignatures
\pysiglinewithargsret{\sphinxbfcode{\sphinxupquote{fix\_external\_set\_virial\_global}}}{\sphinxparam{\DUrole{n}{fix\_id}}\sphinxparamcomma \sphinxparam{\DUrole{n}{virial}}}{}
\pysigstopsignatures
\sphinxAtStartPar
Set the global virial contribution for a fix external instance with the given ID.

\sphinxAtStartPar
\DUrole{versionmodified,added}{Added in version 28Jul2021.}

\sphinxAtStartPar
This is a wrapper around the {\hyperref[\detokenize{Library_utility:_CPPv437lammps_fix_external_set_virial_globalPvPKcPd}]{\sphinxcrossref{\sphinxcode{\sphinxupquote{lammps\_fix\_external\_set\_virial\_global()}}}}} function
of the C\sphinxhyphen{}library interface.
\begin{quote}\begin{description}
\sphinxlineitem{Parameters}\begin{itemize}
\item {} 
\sphinxAtStartPar
\sphinxstyleliteralstrong{\sphinxupquote{fix\_id}} \textendash{} Fix\sphinxhyphen{}ID of a fix external instance

\item {} 
\sphinxAtStartPar
\sphinxstyleliteralstrong{\sphinxupquote{eng}} \textendash{} list of 6 floating point numbers with the virial to be added by fix external

\end{itemize}

\sphinxlineitem{Type}
\sphinxAtStartPar
string

\sphinxlineitem{Type}
\sphinxAtStartPar
float

\end{description}\end{quote}

\end{fulllineitems}

\index{fix\_external\_set\_energy\_peratom() (lammps.lammps method)@\spxentry{fix\_external\_set\_energy\_peratom()}\spxextra{lammps.lammps method}}

\begin{fulllineitems}
\phantomsection\label{\detokenize{Python_module:lammps.lammps.fix_external_set_energy_peratom}}
\pysigstartsignatures
\pysiglinewithargsret{\sphinxbfcode{\sphinxupquote{fix\_external\_set\_energy\_peratom}}}{\sphinxparam{\DUrole{n}{fix\_id}}\sphinxparamcomma \sphinxparam{\DUrole{n}{eatom}}}{}
\pysigstopsignatures
\sphinxAtStartPar
Set the per\sphinxhyphen{}atom energy contribution for a fix external instance with the given ID.

\sphinxAtStartPar
\DUrole{versionmodified,added}{Added in version 28Jul2021.}

\sphinxAtStartPar
This is a wrapper around the {\hyperref[\detokenize{Library_utility:_CPPv438lammps_fix_external_set_energy_peratomPvPKcPd}]{\sphinxcrossref{\sphinxcode{\sphinxupquote{lammps\_fix\_external\_set\_energy\_peratom()}}}}} function
of the C\sphinxhyphen{}library interface.
\begin{quote}\begin{description}
\sphinxlineitem{Parameters}\begin{itemize}
\item {} 
\sphinxAtStartPar
\sphinxstyleliteralstrong{\sphinxupquote{fix\_id}} \textendash{} Fix\sphinxhyphen{}ID of a fix external instance

\item {} 
\sphinxAtStartPar
\sphinxstyleliteralstrong{\sphinxupquote{eatom}} \textendash{} list of potential energy values for local atoms to be added by fix external

\end{itemize}

\sphinxlineitem{Type}
\sphinxAtStartPar
string

\sphinxlineitem{Type}
\sphinxAtStartPar
float

\end{description}\end{quote}

\end{fulllineitems}

\index{fix\_external\_set\_virial\_peratom() (lammps.lammps method)@\spxentry{fix\_external\_set\_virial\_peratom()}\spxextra{lammps.lammps method}}

\begin{fulllineitems}
\phantomsection\label{\detokenize{Python_module:lammps.lammps.fix_external_set_virial_peratom}}
\pysigstartsignatures
\pysiglinewithargsret{\sphinxbfcode{\sphinxupquote{fix\_external\_set\_virial\_peratom}}}{\sphinxparam{\DUrole{n}{fix\_id}}\sphinxparamcomma \sphinxparam{\DUrole{n}{vatom}}}{}
\pysigstopsignatures
\sphinxAtStartPar
Set the per\sphinxhyphen{}atom virial contribution for a fix external instance with the given ID.

\sphinxAtStartPar
\DUrole{versionmodified,added}{Added in version 28Jul2021.}

\sphinxAtStartPar
This is a wrapper around the {\hyperref[\detokenize{Library_utility:_CPPv438lammps_fix_external_set_virial_peratomPvPKcPPd}]{\sphinxcrossref{\sphinxcode{\sphinxupquote{lammps\_fix\_external\_set\_virial\_peratom()}}}}} function
of the C\sphinxhyphen{}library interface.
\begin{quote}\begin{description}
\sphinxlineitem{Parameters}\begin{itemize}
\item {} 
\sphinxAtStartPar
\sphinxstyleliteralstrong{\sphinxupquote{fix\_id}} \textendash{} Fix\sphinxhyphen{}ID of a fix external instance

\item {} 
\sphinxAtStartPar
\sphinxstyleliteralstrong{\sphinxupquote{vatom}} \textendash{} list of natoms lists with 6 floating point numbers to be added by fix external

\end{itemize}

\sphinxlineitem{Type}
\sphinxAtStartPar
string

\sphinxlineitem{Type}
\sphinxAtStartPar
float

\end{description}\end{quote}

\end{fulllineitems}

\index{fix\_external\_set\_vector\_length() (lammps.lammps method)@\spxentry{fix\_external\_set\_vector\_length()}\spxextra{lammps.lammps method}}

\begin{fulllineitems}
\phantomsection\label{\detokenize{Python_module:lammps.lammps.fix_external_set_vector_length}}
\pysigstartsignatures
\pysiglinewithargsret{\sphinxbfcode{\sphinxupquote{fix\_external\_set\_vector\_length}}}{\sphinxparam{\DUrole{n}{fix\_id}}\sphinxparamcomma \sphinxparam{\DUrole{n}{length}}}{}
\pysigstopsignatures
\sphinxAtStartPar
Set the vector length for a global vector stored with fix external for analysis

\sphinxAtStartPar
\DUrole{versionmodified,added}{Added in version 28Jul2021.}

\sphinxAtStartPar
This is a wrapper around the {\hyperref[\detokenize{Library_utility:_CPPv437lammps_fix_external_set_vector_lengthPvPKci}]{\sphinxcrossref{\sphinxcode{\sphinxupquote{lammps\_fix\_external\_set\_vector\_length()}}}}} function
of the C\sphinxhyphen{}library interface.
\begin{quote}\begin{description}
\sphinxlineitem{Parameters}\begin{itemize}
\item {} 
\sphinxAtStartPar
\sphinxstyleliteralstrong{\sphinxupquote{fix\_id}} \textendash{} Fix\sphinxhyphen{}ID of a fix external instance

\item {} 
\sphinxAtStartPar
\sphinxstyleliteralstrong{\sphinxupquote{length}} \textendash{} length of the global vector

\end{itemize}

\sphinxlineitem{Type}
\sphinxAtStartPar
string

\sphinxlineitem{Type}
\sphinxAtStartPar
int

\end{description}\end{quote}

\end{fulllineitems}

\index{fix\_external\_set\_vector() (lammps.lammps method)@\spxentry{fix\_external\_set\_vector()}\spxextra{lammps.lammps method}}

\begin{fulllineitems}
\phantomsection\label{\detokenize{Python_module:lammps.lammps.fix_external_set_vector}}
\pysigstartsignatures
\pysiglinewithargsret{\sphinxbfcode{\sphinxupquote{fix\_external\_set\_vector}}}{\sphinxparam{\DUrole{n}{fix\_id}}\sphinxparamcomma \sphinxparam{\DUrole{n}{idx}}\sphinxparamcomma \sphinxparam{\DUrole{n}{val}}}{}
\pysigstopsignatures
\sphinxAtStartPar
Store a global vector value for a fix external instance with the given ID.

\sphinxAtStartPar
\DUrole{versionmodified,added}{Added in version 28Jul2021.}

\sphinxAtStartPar
This is a wrapper around the {\hyperref[\detokenize{Library_utility:_CPPv430lammps_fix_external_set_vectorPvPKcid}]{\sphinxcrossref{\sphinxcode{\sphinxupquote{lammps\_fix\_external\_set\_vector()}}}}} function
of the C\sphinxhyphen{}library interface.
\begin{quote}\begin{description}
\sphinxlineitem{Parameters}\begin{itemize}
\item {} 
\sphinxAtStartPar
\sphinxstyleliteralstrong{\sphinxupquote{fix\_id}} \textendash{} Fix\sphinxhyphen{}ID of a fix external instance

\item {} 
\sphinxAtStartPar
\sphinxstyleliteralstrong{\sphinxupquote{idx}} \textendash{} 1\sphinxhyphen{}based index of the value in the global vector

\item {} 
\sphinxAtStartPar
\sphinxstyleliteralstrong{\sphinxupquote{val}} \textendash{} value to be stored in the global vector

\end{itemize}

\sphinxlineitem{Type}
\sphinxAtStartPar
string

\sphinxlineitem{Type}
\sphinxAtStartPar
int

\sphinxlineitem{Type}
\sphinxAtStartPar
float

\end{description}\end{quote}

\end{fulllineitems}

\index{get\_neighlist() (lammps.lammps method)@\spxentry{get\_neighlist()}\spxextra{lammps.lammps method}}

\begin{fulllineitems}
\phantomsection\label{\detokenize{Python_module:lammps.lammps.get_neighlist}}
\pysigstartsignatures
\pysiglinewithargsret{\sphinxbfcode{\sphinxupquote{get\_neighlist}}}{\sphinxparam{\DUrole{n}{idx}}}{}
\pysigstopsignatures
\sphinxAtStartPar
Returns an instance of {\hyperref[\detokenize{Python_module:lammps.NeighList}]{\sphinxcrossref{\sphinxcode{\sphinxupquote{NeighList}}}}} which wraps access to the neighbor list with the given index

\sphinxAtStartPar
See {\hyperref[\detokenize{Python_module:lammps.numpy_wrapper.numpy_wrapper.get_neighlist}]{\sphinxcrossref{\sphinxcode{\sphinxupquote{lammps.numpy.get\_neighlist()}}}}} if you want to use
NumPy arrays instead of \sphinxcode{\sphinxupquote{c\_int}} pointers.
\begin{quote}\begin{description}
\sphinxlineitem{Parameters}
\sphinxAtStartPar
\sphinxstyleliteralstrong{\sphinxupquote{idx}} (\sphinxstyleliteralemphasis{\sphinxupquote{int}}) \textendash{} index of neighbor list

\sphinxlineitem{Returns}
\sphinxAtStartPar
an instance of {\hyperref[\detokenize{Python_module:lammps.NeighList}]{\sphinxcrossref{\sphinxcode{\sphinxupquote{NeighList}}}}} wrapping access to neighbor list data

\sphinxlineitem{Return type}
\sphinxAtStartPar
{\hyperref[\detokenize{Python_module:lammps.NeighList}]{\sphinxcrossref{NeighList}}}

\end{description}\end{quote}

\end{fulllineitems}

\index{get\_neighlist\_size() (lammps.lammps method)@\spxentry{get\_neighlist\_size()}\spxextra{lammps.lammps method}}

\begin{fulllineitems}
\phantomsection\label{\detokenize{Python_module:lammps.lammps.get_neighlist_size}}
\pysigstartsignatures
\pysiglinewithargsret{\sphinxbfcode{\sphinxupquote{get\_neighlist\_size}}}{\sphinxparam{\DUrole{n}{idx}}}{}
\pysigstopsignatures
\sphinxAtStartPar
Return the number of elements in neighbor list with the given index
\begin{quote}\begin{description}
\sphinxlineitem{Parameters}
\sphinxAtStartPar
\sphinxstyleliteralstrong{\sphinxupquote{idx}} (\sphinxstyleliteralemphasis{\sphinxupquote{int}}) \textendash{} neighbor list index

\sphinxlineitem{Returns}
\sphinxAtStartPar
number of elements in neighbor list with index idx

\sphinxlineitem{Return type}
\sphinxAtStartPar
int

\end{description}\end{quote}

\end{fulllineitems}

\index{get\_neighlist\_element\_neighbors() (lammps.lammps method)@\spxentry{get\_neighlist\_element\_neighbors()}\spxextra{lammps.lammps method}}

\begin{fulllineitems}
\phantomsection\label{\detokenize{Python_module:lammps.lammps.get_neighlist_element_neighbors}}
\pysigstartsignatures
\pysiglinewithargsret{\sphinxbfcode{\sphinxupquote{get\_neighlist\_element\_neighbors}}}{\sphinxparam{\DUrole{n}{idx}}\sphinxparamcomma \sphinxparam{\DUrole{n}{element}}}{}
\pysigstopsignatures
\sphinxAtStartPar
Return data of neighbor list entry
\begin{quote}\begin{description}
\sphinxlineitem{Parameters}\begin{itemize}
\item {} 
\sphinxAtStartPar
\sphinxstyleliteralstrong{\sphinxupquote{element}} (\sphinxstyleliteralemphasis{\sphinxupquote{int}}) \textendash{} neighbor list index

\item {} 
\sphinxAtStartPar
\sphinxstyleliteralstrong{\sphinxupquote{element}} \textendash{} neighbor list element index

\end{itemize}

\sphinxlineitem{Returns}
\sphinxAtStartPar
tuple with atom local index, number of neighbors and array of neighbor local atom indices

\sphinxlineitem{Return type}
\sphinxAtStartPar
(int, int, POINTER(c\_int))

\end{description}\end{quote}

\end{fulllineitems}

\index{find\_pair\_neighlist() (lammps.lammps method)@\spxentry{find\_pair\_neighlist()}\spxextra{lammps.lammps method}}

\begin{fulllineitems}
\phantomsection\label{\detokenize{Python_module:lammps.lammps.find_pair_neighlist}}
\pysigstartsignatures
\pysiglinewithargsret{\sphinxbfcode{\sphinxupquote{find\_pair\_neighlist}}}{\sphinxparam{\DUrole{n}{style}}\sphinxparamcomma \sphinxparam{\DUrole{n}{exact}\DUrole{o}{=}\DUrole{default_value}{True}}\sphinxparamcomma \sphinxparam{\DUrole{n}{nsub}\DUrole{o}{=}\DUrole{default_value}{0}}\sphinxparamcomma \sphinxparam{\DUrole{n}{reqid}\DUrole{o}{=}\DUrole{default_value}{0}}}{}
\pysigstopsignatures
\sphinxAtStartPar
Find neighbor list index of pair style neighbor list

\sphinxAtStartPar
Search for a neighbor list requested by a pair style instance that
matches “style”.  If exact is True, the pair style name must match
exactly. If exact is False, the pair style name is matched against
“style” as regular expression or sub\sphinxhyphen{}string. If the pair style is a
hybrid pair style, the style is instead matched against the hybrid
sub\sphinxhyphen{}styles. If the same pair style is used as sub\sphinxhyphen{}style multiple
types, you must set nsub to a value n \textgreater{} 0 which indicates the nth
instance of that sub\sphinxhyphen{}style to be used (same as for the pair\_coeff
command). The default value of 0 will fail to match in that case.

\sphinxAtStartPar
Once the pair style instance has been identified, it may have
requested multiple neighbor lists. Those are uniquely identified by
a request ID \textgreater{} 0 as set by the pair style. Otherwise the request
ID is 0.
\begin{quote}\begin{description}
\sphinxlineitem{Parameters}\begin{itemize}
\item {} 
\sphinxAtStartPar
\sphinxstyleliteralstrong{\sphinxupquote{style}} (\sphinxstyleliteralemphasis{\sphinxupquote{string}}) \textendash{} name of pair style that should be searched for

\item {} 
\sphinxAtStartPar
\sphinxstyleliteralstrong{\sphinxupquote{exact}} (\sphinxstyleliteralemphasis{\sphinxupquote{bool}}\sphinxstyleliteralemphasis{\sphinxupquote{, }}\sphinxstyleliteralemphasis{\sphinxupquote{optional}}) \textendash{} controls whether style should match exactly or only must be contained in pair style name, defaults to True

\item {} 
\sphinxAtStartPar
\sphinxstyleliteralstrong{\sphinxupquote{nsub}} (\sphinxstyleliteralemphasis{\sphinxupquote{int}}\sphinxstyleliteralemphasis{\sphinxupquote{, }}\sphinxstyleliteralemphasis{\sphinxupquote{optional}}) \textendash{} match nsub\sphinxhyphen{}th hybrid sub\sphinxhyphen{}style, defaults to 0

\item {} 
\sphinxAtStartPar
\sphinxstyleliteralstrong{\sphinxupquote{reqid}} (\sphinxstyleliteralemphasis{\sphinxupquote{int}}\sphinxstyleliteralemphasis{\sphinxupquote{, }}\sphinxstyleliteralemphasis{\sphinxupquote{optional}}) \textendash{} list request id, \textgreater{} 0 in case there are more than one, defaults to 0

\end{itemize}

\sphinxlineitem{Returns}
\sphinxAtStartPar
neighbor list index if found, otherwise \sphinxhyphen{}1

\sphinxlineitem{Return type}
\sphinxAtStartPar
int

\end{description}\end{quote}

\end{fulllineitems}

\index{find\_fix\_neighlist() (lammps.lammps method)@\spxentry{find\_fix\_neighlist()}\spxextra{lammps.lammps method}}

\begin{fulllineitems}
\phantomsection\label{\detokenize{Python_module:lammps.lammps.find_fix_neighlist}}
\pysigstartsignatures
\pysiglinewithargsret{\sphinxbfcode{\sphinxupquote{find\_fix\_neighlist}}}{\sphinxparam{\DUrole{n}{fixid}}\sphinxparamcomma \sphinxparam{\DUrole{n}{reqid}\DUrole{o}{=}\DUrole{default_value}{0}}}{}
\pysigstopsignatures
\sphinxAtStartPar
Find neighbor list index of fix neighbor list

\sphinxAtStartPar
The fix instance requesting the neighbor list is uniquely identified
by the fix ID.  In case the fix has requested multiple neighbor
lists, those are uniquely identified by a request ID \textgreater{} 0 as set by
the fix.  Otherwise the request ID is 0 (the default).
\begin{quote}\begin{description}
\sphinxlineitem{Parameters}\begin{itemize}
\item {} 
\sphinxAtStartPar
\sphinxstyleliteralstrong{\sphinxupquote{fixid}} (\sphinxstyleliteralemphasis{\sphinxupquote{string}}) \textendash{} name of fix

\item {} 
\sphinxAtStartPar
\sphinxstyleliteralstrong{\sphinxupquote{reqid}} (\sphinxstyleliteralemphasis{\sphinxupquote{int}}\sphinxstyleliteralemphasis{\sphinxupquote{, }}\sphinxstyleliteralemphasis{\sphinxupquote{optional}}) \textendash{} id of neighbor list request, in case there are more than one request, defaults to 0

\end{itemize}

\sphinxlineitem{Returns}
\sphinxAtStartPar
neighbor list index if found, otherwise \sphinxhyphen{}1

\sphinxlineitem{Return type}
\sphinxAtStartPar
int

\end{description}\end{quote}

\end{fulllineitems}

\index{find\_compute\_neighlist() (lammps.lammps method)@\spxentry{find\_compute\_neighlist()}\spxextra{lammps.lammps method}}

\begin{fulllineitems}
\phantomsection\label{\detokenize{Python_module:lammps.lammps.find_compute_neighlist}}
\pysigstartsignatures
\pysiglinewithargsret{\sphinxbfcode{\sphinxupquote{find\_compute\_neighlist}}}{\sphinxparam{\DUrole{n}{computeid}}\sphinxparamcomma \sphinxparam{\DUrole{n}{reqid}\DUrole{o}{=}\DUrole{default_value}{0}}}{}
\pysigstopsignatures
\sphinxAtStartPar
Find neighbor list index of compute neighbor list

\sphinxAtStartPar
The compute instance requesting the neighbor list is uniquely
identified by the compute ID.  In case the compute has requested
multiple neighbor lists, those are uniquely identified by a request
ID \textgreater{} 0 as set by the compute.  Otherwise the request ID is 0 (the
default).
\begin{quote}\begin{description}
\sphinxlineitem{Parameters}\begin{itemize}
\item {} 
\sphinxAtStartPar
\sphinxstyleliteralstrong{\sphinxupquote{computeid}} (\sphinxstyleliteralemphasis{\sphinxupquote{string}}) \textendash{} name of compute

\item {} 
\sphinxAtStartPar
\sphinxstyleliteralstrong{\sphinxupquote{reqid}} (\sphinxstyleliteralemphasis{\sphinxupquote{int}}\sphinxstyleliteralemphasis{\sphinxupquote{, }}\sphinxstyleliteralemphasis{\sphinxupquote{optional}}) \textendash{} index of neighbor list request, in case there are more than one request, defaults to 0

\end{itemize}

\sphinxlineitem{Returns}
\sphinxAtStartPar
neighbor list index if found, otherwise \sphinxhyphen{}1

\sphinxlineitem{Return type}
\sphinxAtStartPar
int

\end{description}\end{quote}

\end{fulllineitems}


\end{fulllineitems}

\index{numpy\_wrapper (class in lammps.numpy\_wrapper)@\spxentry{numpy\_wrapper}\spxextra{class in lammps.numpy\_wrapper}}

\begin{fulllineitems}
\phantomsection\label{\detokenize{Python_module:lammps.numpy_wrapper.numpy_wrapper}}
\pysigstartsignatures
\pysiglinewithargsret{\sphinxbfcode{\sphinxupquote{class\DUrole{w}{ }}}\sphinxcode{\sphinxupquote{lammps.numpy\_wrapper.}}\sphinxbfcode{\sphinxupquote{numpy\_wrapper}}}{\sphinxparam{\DUrole{n}{lmp}}}{}
\pysigstopsignatures
\sphinxAtStartPar
lammps API NumPy Wrapper

\sphinxAtStartPar
This is a wrapper class that provides additional methods on top of an
existing {\hyperref[\detokenize{Python_module:module-lammps}]{\sphinxcrossref{\sphinxcode{\sphinxupquote{lammps}}}}} instance. The methods transform raw ctypes
pointers into NumPy arrays, which give direct access to the
original data while protecting against out\sphinxhyphen{}of\sphinxhyphen{}bounds accesses.

\sphinxAtStartPar
There is no need to explicitly instantiate this class. Each instance
of {\hyperref[\detokenize{Python_module:module-lammps}]{\sphinxcrossref{\sphinxcode{\sphinxupquote{lammps}}}}} has a \sphinxcode{\sphinxupquote{numpy}} property
that returns an instance.
\begin{quote}\begin{description}
\sphinxlineitem{Parameters}
\sphinxAtStartPar
\sphinxstyleliteralstrong{\sphinxupquote{lmp}} ({\hyperref[\detokenize{Python_module:lammps.lammps}]{\sphinxcrossref{\sphinxstyleliteralemphasis{\sphinxupquote{lammps}}}}}) \textendash{} instance of the {\hyperref[\detokenize{Python_module:module-lammps}]{\sphinxcrossref{\sphinxcode{\sphinxupquote{lammps}}}}} class

\end{description}\end{quote}
\index{extract\_atom() (lammps.numpy\_wrapper.numpy\_wrapper method)@\spxentry{extract\_atom()}\spxextra{lammps.numpy\_wrapper.numpy\_wrapper method}}

\begin{fulllineitems}
\phantomsection\label{\detokenize{Python_module:lammps.numpy_wrapper.numpy_wrapper.extract_atom}}
\pysigstartsignatures
\pysiglinewithargsret{\sphinxbfcode{\sphinxupquote{extract\_atom}}}{\sphinxparam{\DUrole{n}{name}}\sphinxparamcomma \sphinxparam{\DUrole{n}{dtype}\DUrole{o}{=}\DUrole{default_value}{None}}\sphinxparamcomma \sphinxparam{\DUrole{n}{nelem}\DUrole{o}{=}\DUrole{default_value}{None}}\sphinxparamcomma \sphinxparam{\DUrole{n}{dim}\DUrole{o}{=}\DUrole{default_value}{None}}}{}
\pysigstopsignatures
\sphinxAtStartPar
Retrieve per\sphinxhyphen{}atom properties from LAMMPS as NumPy arrays

\sphinxAtStartPar
This is a wrapper around the \sphinxcode{\sphinxupquote{lammps.extract\_atom()}} method.
It behaves the same as the original method, but returns NumPy arrays
instead of \sphinxcode{\sphinxupquote{ctypes}} pointers.

\begin{sphinxadmonition}{note}{Note:}
\sphinxAtStartPar
The returned arrays of per\sphinxhyphen{}atom data are by default dimensioned
for the range {[}0:nlocal{]} since that data is \sphinxstyleemphasis{always} valid.  The
underlying storage for the data, however, is typically allocated
for the range of {[}0:nmax{]}. Whether there is valid data in the range
{[}nlocal:nlocal+nghost{]} depends on whether the property of interest
is also updated for ghost atoms.  This is not often the case.  In
some cases, it depends on a LAMMPS setting, see for example
{\hyperref[\detokenize{comm_modify::doc}]{\sphinxcrossref{\DUrole{doc}{comm\_modify vel yes}}}}.  By using the optional
\sphinxstyleemphasis{nelem} parameter the size of the returned NumPy can be overridden.
There is no check whether the number of elements chosen is valid.
\end{sphinxadmonition}
\begin{quote}\begin{description}
\sphinxlineitem{Parameters}\begin{itemize}
\item {} 
\sphinxAtStartPar
\sphinxstyleliteralstrong{\sphinxupquote{name}} (\sphinxstyleliteralemphasis{\sphinxupquote{string}}) \textendash{} name of the property

\item {} 
\sphinxAtStartPar
\sphinxstyleliteralstrong{\sphinxupquote{dtype}} (\sphinxstyleliteralemphasis{\sphinxupquote{int}}\sphinxstyleliteralemphasis{\sphinxupquote{, }}\sphinxstyleliteralemphasis{\sphinxupquote{optional}}) \textendash{} type of the returned data (see {\hyperref[\detokenize{Python_module:py-datatype-constants}]{\sphinxcrossref{\DUrole{std,std-ref}{Data Types}}}})

\item {} 
\sphinxAtStartPar
\sphinxstyleliteralstrong{\sphinxupquote{nelem}} (\sphinxstyleliteralemphasis{\sphinxupquote{int}}\sphinxstyleliteralemphasis{\sphinxupquote{, }}\sphinxstyleliteralemphasis{\sphinxupquote{optional}}) \textendash{} number of elements in array

\item {} 
\sphinxAtStartPar
\sphinxstyleliteralstrong{\sphinxupquote{dim}} (\sphinxstyleliteralemphasis{\sphinxupquote{int}}\sphinxstyleliteralemphasis{\sphinxupquote{, }}\sphinxstyleliteralemphasis{\sphinxupquote{optional}}) \textendash{} dimension of each element

\end{itemize}

\sphinxlineitem{Returns}
\sphinxAtStartPar
requested data as NumPy array with direct access to C data or None

\sphinxlineitem{Return type}
\sphinxAtStartPar
numpy.array or NoneType

\end{description}\end{quote}

\end{fulllineitems}

\index{extract\_compute() (lammps.numpy\_wrapper.numpy\_wrapper method)@\spxentry{extract\_compute()}\spxextra{lammps.numpy\_wrapper.numpy\_wrapper method}}

\begin{fulllineitems}
\phantomsection\label{\detokenize{Python_module:lammps.numpy_wrapper.numpy_wrapper.extract_compute}}
\pysigstartsignatures
\pysiglinewithargsret{\sphinxbfcode{\sphinxupquote{extract\_compute}}}{\sphinxparam{\DUrole{n}{cid}}\sphinxparamcomma \sphinxparam{\DUrole{n}{cstyle}}\sphinxparamcomma \sphinxparam{\DUrole{n}{ctype}}}{}
\pysigstopsignatures
\sphinxAtStartPar
Retrieve data from a LAMMPS compute

\sphinxAtStartPar
This is a wrapper around the
{\hyperref[\detokenize{Python_module:lammps.lammps.extract_compute}]{\sphinxcrossref{\sphinxcode{\sphinxupquote{lammps.extract\_compute()}}}}} method.
It behaves the same as the original method, but returns NumPy arrays
instead of \sphinxcode{\sphinxupquote{ctypes}} pointers.
\begin{quote}\begin{description}
\sphinxlineitem{Parameters}\begin{itemize}
\item {} 
\sphinxAtStartPar
\sphinxstyleliteralstrong{\sphinxupquote{cid}} (\sphinxstyleliteralemphasis{\sphinxupquote{string}}) \textendash{} compute ID

\item {} 
\sphinxAtStartPar
\sphinxstyleliteralstrong{\sphinxupquote{cstyle}} (\sphinxstyleliteralemphasis{\sphinxupquote{int}}) \textendash{} style of the data retrieve (global, atom, or local), see {\hyperref[\detokenize{Python_module:py-style-constants}]{\sphinxcrossref{\DUrole{std,std-ref}{Style Constants}}}}

\item {} 
\sphinxAtStartPar
\sphinxstyleliteralstrong{\sphinxupquote{ctype}} (\sphinxstyleliteralemphasis{\sphinxupquote{int}}) \textendash{} type of the returned data (scalar, vector, or array), see {\hyperref[\detokenize{Python_module:py-type-constants}]{\sphinxcrossref{\DUrole{std,std-ref}{Type Constants}}}}

\end{itemize}

\sphinxlineitem{Returns}
\sphinxAtStartPar
requested data either as float, as NumPy array with direct access to C data, or None

\sphinxlineitem{Return type}
\sphinxAtStartPar
float, numpy.array, or NoneType

\end{description}\end{quote}

\end{fulllineitems}

\index{extract\_fix() (lammps.numpy\_wrapper.numpy\_wrapper method)@\spxentry{extract\_fix()}\spxextra{lammps.numpy\_wrapper.numpy\_wrapper method}}

\begin{fulllineitems}
\phantomsection\label{\detokenize{Python_module:lammps.numpy_wrapper.numpy_wrapper.extract_fix}}
\pysigstartsignatures
\pysiglinewithargsret{\sphinxbfcode{\sphinxupquote{extract\_fix}}}{\sphinxparam{\DUrole{n}{fid}}\sphinxparamcomma \sphinxparam{\DUrole{n}{fstyle}}\sphinxparamcomma \sphinxparam{\DUrole{n}{ftype}}\sphinxparamcomma \sphinxparam{\DUrole{n}{nrow}\DUrole{o}{=}\DUrole{default_value}{0}}\sphinxparamcomma \sphinxparam{\DUrole{n}{ncol}\DUrole{o}{=}\DUrole{default_value}{0}}}{}
\pysigstopsignatures
\sphinxAtStartPar
Retrieve data from a LAMMPS fix

\sphinxAtStartPar
This is a wrapper around the {\hyperref[\detokenize{Python_module:lammps.lammps.extract_fix}]{\sphinxcrossref{\sphinxcode{\sphinxupquote{lammps.extract\_fix()}}}}} method.
It behaves the same as the original method, but returns NumPy arrays
instead of \sphinxcode{\sphinxupquote{ctypes}} pointers.

\begin{sphinxadmonition}{note}{Note:}
\sphinxAtStartPar
When requesting global data, the fix data can only be accessed one
item at a time without access to the whole vector or array.  Thus this
function will always return a scalar.  To access vector or array elements
the “nrow” and “ncol” arguments need to be set accordingly (they default to 0).
\end{sphinxadmonition}
\begin{quote}\begin{description}
\sphinxlineitem{Parameters}\begin{itemize}
\item {} 
\sphinxAtStartPar
\sphinxstyleliteralstrong{\sphinxupquote{fid}} (\sphinxstyleliteralemphasis{\sphinxupquote{string}}) \textendash{} fix ID

\item {} 
\sphinxAtStartPar
\sphinxstyleliteralstrong{\sphinxupquote{fstyle}} (\sphinxstyleliteralemphasis{\sphinxupquote{int}}) \textendash{} style of the data retrieve (global, atom, or local), see {\hyperref[\detokenize{Python_module:py-style-constants}]{\sphinxcrossref{\DUrole{std,std-ref}{Style Constants}}}}

\item {} 
\sphinxAtStartPar
\sphinxstyleliteralstrong{\sphinxupquote{ftype}} (\sphinxstyleliteralemphasis{\sphinxupquote{int}}) \textendash{} type or size of the returned data (scalar, vector, or array), see {\hyperref[\detokenize{Python_module:py-type-constants}]{\sphinxcrossref{\DUrole{std,std-ref}{Type Constants}}}}

\item {} 
\sphinxAtStartPar
\sphinxstyleliteralstrong{\sphinxupquote{nrow}} (\sphinxstyleliteralemphasis{\sphinxupquote{int}}) \textendash{} index of global vector element or row index of global array element

\item {} 
\sphinxAtStartPar
\sphinxstyleliteralstrong{\sphinxupquote{ncol}} (\sphinxstyleliteralemphasis{\sphinxupquote{int}}) \textendash{} column index of global array element

\end{itemize}

\sphinxlineitem{Returns}
\sphinxAtStartPar
requested data

\sphinxlineitem{Return type}
\sphinxAtStartPar
integer or double value, pointer to 1d or 2d double array  or None

\end{description}\end{quote}

\end{fulllineitems}

\index{extract\_variable() (lammps.numpy\_wrapper.numpy\_wrapper method)@\spxentry{extract\_variable()}\spxextra{lammps.numpy\_wrapper.numpy\_wrapper method}}

\begin{fulllineitems}
\phantomsection\label{\detokenize{Python_module:lammps.numpy_wrapper.numpy_wrapper.extract_variable}}
\pysigstartsignatures
\pysiglinewithargsret{\sphinxbfcode{\sphinxupquote{extract\_variable}}}{\sphinxparam{\DUrole{n}{name}}\sphinxparamcomma \sphinxparam{\DUrole{n}{group}\DUrole{o}{=}\DUrole{default_value}{None}}\sphinxparamcomma \sphinxparam{\DUrole{n}{vartype}\DUrole{o}{=}\DUrole{default_value}{0}}}{}
\pysigstopsignatures
\sphinxAtStartPar
Evaluate a LAMMPS variable and return its data

\sphinxAtStartPar
This function is a wrapper around the function
{\hyperref[\detokenize{Python_module:lammps.lammps.extract_variable}]{\sphinxcrossref{\sphinxcode{\sphinxupquote{lammps.extract\_variable()}}}}}
method. It behaves the same as the original method, but returns NumPy arrays
instead of \sphinxcode{\sphinxupquote{ctypes}} pointers.
\begin{quote}\begin{description}
\sphinxlineitem{Parameters}\begin{itemize}
\item {} 
\sphinxAtStartPar
\sphinxstyleliteralstrong{\sphinxupquote{name}} (\sphinxstyleliteralemphasis{\sphinxupquote{string}}) \textendash{} name of the variable to execute

\item {} 
\sphinxAtStartPar
\sphinxstyleliteralstrong{\sphinxupquote{group}} (\sphinxstyleliteralemphasis{\sphinxupquote{string}}) \textendash{} name of group for atom\sphinxhyphen{}style variable (ignored for equal\sphinxhyphen{}style variables)

\item {} 
\sphinxAtStartPar
\sphinxstyleliteralstrong{\sphinxupquote{vartype}} (\sphinxstyleliteralemphasis{\sphinxupquote{int}}) \textendash{} type of variable, see {\hyperref[\detokenize{Python_module:py-vartype-constants}]{\sphinxcrossref{\DUrole{std,std-ref}{Variable Type Constants}}}}

\end{itemize}

\sphinxlineitem{Returns}
\sphinxAtStartPar
the requested data or None

\sphinxlineitem{Return type}
\sphinxAtStartPar
c\_double, numpy.array, or NoneType

\end{description}\end{quote}

\end{fulllineitems}

\index{gather\_bonds() (lammps.numpy\_wrapper.numpy\_wrapper method)@\spxentry{gather\_bonds()}\spxextra{lammps.numpy\_wrapper.numpy\_wrapper method}}

\begin{fulllineitems}
\phantomsection\label{\detokenize{Python_module:lammps.numpy_wrapper.numpy_wrapper.gather_bonds}}
\pysigstartsignatures
\pysiglinewithargsret{\sphinxbfcode{\sphinxupquote{gather\_bonds}}}{}{}
\pysigstopsignatures
\sphinxAtStartPar
Retrieve global list of bonds as a NumPy array

\sphinxAtStartPar
\DUrole{versionmodified,added}{Added in version 28Jul2021.}

\sphinxAtStartPar
This is a wrapper around {\hyperref[\detokenize{Python_module:lammps.lammps.gather_bonds}]{\sphinxcrossref{\sphinxcode{\sphinxupquote{lammps.gather\_bonds()}}}}}.
It behaves the same as the original method, but returns a NumPy array instead
of a \sphinxcode{\sphinxupquote{ctypes}} list.
\begin{quote}\begin{description}
\sphinxlineitem{Returns}
\sphinxAtStartPar
the requested data as a 2d\sphinxhyphen{}integer numpy array

\sphinxlineitem{Return type}
\sphinxAtStartPar
numpy.array(nbonds,3)

\end{description}\end{quote}

\end{fulllineitems}

\index{gather\_angles() (lammps.numpy\_wrapper.numpy\_wrapper method)@\spxentry{gather\_angles()}\spxextra{lammps.numpy\_wrapper.numpy\_wrapper method}}

\begin{fulllineitems}
\phantomsection\label{\detokenize{Python_module:lammps.numpy_wrapper.numpy_wrapper.gather_angles}}
\pysigstartsignatures
\pysiglinewithargsret{\sphinxbfcode{\sphinxupquote{gather\_angles}}}{}{}
\pysigstopsignatures
\sphinxAtStartPar
Retrieve global list of angles as a NumPy array

\sphinxAtStartPar
\DUrole{versionmodified,added}{Added in version 8Feb2023.}

\sphinxAtStartPar
This is a wrapper around {\hyperref[\detokenize{Python_module:lammps.lammps.gather_angles}]{\sphinxcrossref{\sphinxcode{\sphinxupquote{lammps.gather\_angles()}}}}}.
It behaves the same as the original method, but returns a NumPy array instead
of a \sphinxcode{\sphinxupquote{ctypes}} list.
\begin{quote}\begin{description}
\sphinxlineitem{Returns}
\sphinxAtStartPar
the requested data as a 2d\sphinxhyphen{}integer numpy array

\sphinxlineitem{Return type}
\sphinxAtStartPar
numpy.array(nangles,4)

\end{description}\end{quote}

\end{fulllineitems}

\index{gather\_dihedrals() (lammps.numpy\_wrapper.numpy\_wrapper method)@\spxentry{gather\_dihedrals()}\spxextra{lammps.numpy\_wrapper.numpy\_wrapper method}}

\begin{fulllineitems}
\phantomsection\label{\detokenize{Python_module:lammps.numpy_wrapper.numpy_wrapper.gather_dihedrals}}
\pysigstartsignatures
\pysiglinewithargsret{\sphinxbfcode{\sphinxupquote{gather\_dihedrals}}}{}{}
\pysigstopsignatures
\sphinxAtStartPar
Retrieve global list of dihedrals as a NumPy array

\sphinxAtStartPar
\DUrole{versionmodified,added}{Added in version 8Feb2023.}

\sphinxAtStartPar
This is a wrapper around {\hyperref[\detokenize{Python_module:lammps.lammps.gather_dihedrals}]{\sphinxcrossref{\sphinxcode{\sphinxupquote{lammps.gather\_dihedrals()}}}}}.
It behaves the same as the original method, but returns a NumPy array instead
of a \sphinxcode{\sphinxupquote{ctypes}} list.
\begin{quote}\begin{description}
\sphinxlineitem{Returns}
\sphinxAtStartPar
the requested data as a 2d\sphinxhyphen{}integer numpy array

\sphinxlineitem{Return type}
\sphinxAtStartPar
numpy.array(ndihedrals,5)

\end{description}\end{quote}

\end{fulllineitems}

\index{gather\_impropers() (lammps.numpy\_wrapper.numpy\_wrapper method)@\spxentry{gather\_impropers()}\spxextra{lammps.numpy\_wrapper.numpy\_wrapper method}}

\begin{fulllineitems}
\phantomsection\label{\detokenize{Python_module:lammps.numpy_wrapper.numpy_wrapper.gather_impropers}}
\pysigstartsignatures
\pysiglinewithargsret{\sphinxbfcode{\sphinxupquote{gather\_impropers}}}{}{}
\pysigstopsignatures
\sphinxAtStartPar
Retrieve global list of impropers as a NumPy array

\sphinxAtStartPar
\DUrole{versionmodified,added}{Added in version 8Feb2023.}

\sphinxAtStartPar
This is a wrapper around {\hyperref[\detokenize{Python_module:lammps.lammps.gather_impropers}]{\sphinxcrossref{\sphinxcode{\sphinxupquote{lammps.gather\_impropers()}}}}}.
It behaves the same as the original method, but returns a NumPy array instead
of a \sphinxcode{\sphinxupquote{ctypes}} list.
\begin{quote}\begin{description}
\sphinxlineitem{Returns}
\sphinxAtStartPar
the requested data as a 2d\sphinxhyphen{}integer numpy array

\sphinxlineitem{Return type}
\sphinxAtStartPar
numpy.array(nimpropers,5)

\end{description}\end{quote}

\end{fulllineitems}

\index{fix\_external\_get\_force() (lammps.numpy\_wrapper.numpy\_wrapper method)@\spxentry{fix\_external\_get\_force()}\spxextra{lammps.numpy\_wrapper.numpy\_wrapper method}}

\begin{fulllineitems}
\phantomsection\label{\detokenize{Python_module:lammps.numpy_wrapper.numpy_wrapper.fix_external_get_force}}
\pysigstartsignatures
\pysiglinewithargsret{\sphinxbfcode{\sphinxupquote{fix\_external\_get\_force}}}{\sphinxparam{\DUrole{n}{fix\_id}}}{}
\pysigstopsignatures
\sphinxAtStartPar
Get access to the array with per\sphinxhyphen{}atom forces of a fix external instance with a given fix ID.

\sphinxAtStartPar
\DUrole{versionmodified,changed}{Changed in version 28Jul2021.}

\sphinxAtStartPar
This function is a wrapper around the
{\hyperref[\detokenize{Python_module:lammps.lammps.fix_external_get_force}]{\sphinxcrossref{\sphinxcode{\sphinxupquote{lammps.fix\_external\_get\_force()}}}}}
method.  It behaves the same as the original method, but returns a NumPy array instead
of a \sphinxcode{\sphinxupquote{ctypes}} pointer.
\begin{quote}\begin{description}
\sphinxlineitem{Parameters}
\sphinxAtStartPar
\sphinxstyleliteralstrong{\sphinxupquote{fix\_id}} \textendash{} Fix\sphinxhyphen{}ID of a fix external instance

\sphinxlineitem{Type}
\sphinxAtStartPar
string

\sphinxlineitem{Returns}
\sphinxAtStartPar
requested data

\sphinxlineitem{Return type}
\sphinxAtStartPar
numpy.array

\end{description}\end{quote}

\end{fulllineitems}

\index{fix\_external\_set\_energy\_peratom() (lammps.numpy\_wrapper.numpy\_wrapper method)@\spxentry{fix\_external\_set\_energy\_peratom()}\spxextra{lammps.numpy\_wrapper.numpy\_wrapper method}}

\begin{fulllineitems}
\phantomsection\label{\detokenize{Python_module:lammps.numpy_wrapper.numpy_wrapper.fix_external_set_energy_peratom}}
\pysigstartsignatures
\pysiglinewithargsret{\sphinxbfcode{\sphinxupquote{fix\_external\_set\_energy\_peratom}}}{\sphinxparam{\DUrole{n}{fix\_id}}\sphinxparamcomma \sphinxparam{\DUrole{n}{eatom}}}{}
\pysigstopsignatures
\sphinxAtStartPar
Set the per\sphinxhyphen{}atom energy contribution for a fix external instance with the given ID.

\sphinxAtStartPar
\DUrole{versionmodified,added}{Added in version 28Jul2021.}

\sphinxAtStartPar
This function is an alternative to
{\hyperref[\detokenize{Python_module:lammps.lammps.fix_external_set_energy_peratom}]{\sphinxcrossref{\sphinxcode{\sphinxupquote{lammps.fix\_external\_set\_energy\_peratom()}}}}}
method.  It behaves the same as the original method, but accepts a NumPy array
instead of a list as argument.
\begin{quote}\begin{description}
\sphinxlineitem{Parameters}\begin{itemize}
\item {} 
\sphinxAtStartPar
\sphinxstyleliteralstrong{\sphinxupquote{fix\_id}} \textendash{} Fix\sphinxhyphen{}ID of a fix external instance

\item {} 
\sphinxAtStartPar
\sphinxstyleliteralstrong{\sphinxupquote{eatom}} \textendash{} per\sphinxhyphen{}atom potential energy

\end{itemize}

\sphinxlineitem{Type}
\sphinxAtStartPar
string

\sphinxlineitem{Type}
\sphinxAtStartPar
numpy.array

\end{description}\end{quote}

\end{fulllineitems}

\index{fix\_external\_set\_virial\_peratom() (lammps.numpy\_wrapper.numpy\_wrapper method)@\spxentry{fix\_external\_set\_virial\_peratom()}\spxextra{lammps.numpy\_wrapper.numpy\_wrapper method}}

\begin{fulllineitems}
\phantomsection\label{\detokenize{Python_module:lammps.numpy_wrapper.numpy_wrapper.fix_external_set_virial_peratom}}
\pysigstartsignatures
\pysiglinewithargsret{\sphinxbfcode{\sphinxupquote{fix\_external\_set\_virial\_peratom}}}{\sphinxparam{\DUrole{n}{fix\_id}}\sphinxparamcomma \sphinxparam{\DUrole{n}{vatom}}}{}
\pysigstopsignatures
\sphinxAtStartPar
Set the per\sphinxhyphen{}atom virial contribution for a fix external instance with the given ID.

\sphinxAtStartPar
\DUrole{versionmodified,added}{Added in version 28Jul2021.}

\sphinxAtStartPar
This function is an alternative to
{\hyperref[\detokenize{Python_module:lammps.lammps.fix_external_set_virial_peratom}]{\sphinxcrossref{\sphinxcode{\sphinxupquote{lammps.fix\_external\_set\_virial\_peratom()}}}}}
method.  It behaves the same as the original method, but accepts a NumPy array
instead of a list as argument.
\begin{quote}\begin{description}
\sphinxlineitem{Parameters}\begin{itemize}
\item {} 
\sphinxAtStartPar
\sphinxstyleliteralstrong{\sphinxupquote{fix\_id}} \textendash{} Fix\sphinxhyphen{}ID of a fix external instance

\item {} 
\sphinxAtStartPar
\sphinxstyleliteralstrong{\sphinxupquote{eatom}} \textendash{} per\sphinxhyphen{}atom potential energy

\end{itemize}

\sphinxlineitem{Type}
\sphinxAtStartPar
string

\sphinxlineitem{Type}
\sphinxAtStartPar
numpy.array

\end{description}\end{quote}

\end{fulllineitems}

\index{get\_neighlist() (lammps.numpy\_wrapper.numpy\_wrapper method)@\spxentry{get\_neighlist()}\spxextra{lammps.numpy\_wrapper.numpy\_wrapper method}}

\begin{fulllineitems}
\phantomsection\label{\detokenize{Python_module:lammps.numpy_wrapper.numpy_wrapper.get_neighlist}}
\pysigstartsignatures
\pysiglinewithargsret{\sphinxbfcode{\sphinxupquote{get\_neighlist}}}{\sphinxparam{\DUrole{n}{idx}}}{}
\pysigstopsignatures
\sphinxAtStartPar
Returns an instance of {\hyperref[\detokenize{Python_module:lammps.numpy_wrapper.NumPyNeighList}]{\sphinxcrossref{\sphinxcode{\sphinxupquote{NumPyNeighList}}}}} which wraps access to the neighbor list with the given index
\begin{quote}\begin{description}
\sphinxlineitem{Parameters}
\sphinxAtStartPar
\sphinxstyleliteralstrong{\sphinxupquote{idx}} (\sphinxstyleliteralemphasis{\sphinxupquote{int}}) \textendash{} index of neighbor list

\sphinxlineitem{Returns}
\sphinxAtStartPar
an instance of {\hyperref[\detokenize{Python_module:lammps.numpy_wrapper.NumPyNeighList}]{\sphinxcrossref{\sphinxcode{\sphinxupquote{NumPyNeighList}}}}} wrapping access to neighbor list data

\sphinxlineitem{Return type}
\sphinxAtStartPar
{\hyperref[\detokenize{Python_module:lammps.numpy_wrapper.NumPyNeighList}]{\sphinxcrossref{NumPyNeighList}}}

\end{description}\end{quote}

\end{fulllineitems}

\index{get\_neighlist\_element\_neighbors() (lammps.numpy\_wrapper.numpy\_wrapper method)@\spxentry{get\_neighlist\_element\_neighbors()}\spxextra{lammps.numpy\_wrapper.numpy\_wrapper method}}

\begin{fulllineitems}
\phantomsection\label{\detokenize{Python_module:lammps.numpy_wrapper.numpy_wrapper.get_neighlist_element_neighbors}}
\pysigstartsignatures
\pysiglinewithargsret{\sphinxbfcode{\sphinxupquote{get\_neighlist\_element\_neighbors}}}{\sphinxparam{\DUrole{n}{idx}}\sphinxparamcomma \sphinxparam{\DUrole{n}{element}}}{}
\pysigstopsignatures
\sphinxAtStartPar
Return data of neighbor list entry

\sphinxAtStartPar
This function is a wrapper around the function
{\hyperref[\detokenize{Python_module:lammps.lammps.get_neighlist_element_neighbors}]{\sphinxcrossref{\sphinxcode{\sphinxupquote{lammps.get\_neighlist\_element\_neighbors()}}}}}
method. It behaves the same as the original method, but returns a NumPy array containing the neighbors
instead of a \sphinxcode{\sphinxupquote{ctypes}} pointer.
\begin{quote}\begin{description}
\sphinxlineitem{Parameters}\begin{itemize}
\item {} 
\sphinxAtStartPar
\sphinxstyleliteralstrong{\sphinxupquote{element}} (\sphinxstyleliteralemphasis{\sphinxupquote{int}}) \textendash{} neighbor list index

\item {} 
\sphinxAtStartPar
\sphinxstyleliteralstrong{\sphinxupquote{element}} \textendash{} neighbor list element index

\end{itemize}

\sphinxlineitem{Returns}
\sphinxAtStartPar
tuple with atom local index and numpy array of neighbor local atom indices

\sphinxlineitem{Return type}
\sphinxAtStartPar
(int, numpy.array)

\end{description}\end{quote}

\end{fulllineitems}


\end{fulllineitems}



\bigskip\hrule\bigskip



\subsection{The \sphinxstyleliteralintitle{\sphinxupquote{PyLammps}} class API}
\label{\detokenize{Python_module:the-pylammps-class-api}}
\sphinxAtStartPar
The {\hyperref[\detokenize{Python_module:lammps.PyLammps}]{\sphinxcrossref{\sphinxcode{\sphinxupquote{PyLammps}}}}} class is a wrapper that creates a
simpler, more “Pythonic” interface to common LAMMPS functionality. LAMMPS
data structures are exposed through objects and properties. This makes Python
scripts shorter and more concise. See the {\hyperref[\detokenize{Howto_pylammps::doc}]{\sphinxcrossref{\DUrole{doc}{PyLammps Tutorial}}}} for an introduction on how to use this interface.
\index{PyLammps (class in lammps)@\spxentry{PyLammps}\spxextra{class in lammps}}

\begin{fulllineitems}
\phantomsection\label{\detokenize{Python_module:lammps.PyLammps}}
\pysigstartsignatures
\pysiglinewithargsret{\sphinxbfcode{\sphinxupquote{class\DUrole{w}{ }}}\sphinxcode{\sphinxupquote{lammps.}}\sphinxbfcode{\sphinxupquote{PyLammps}}}{\sphinxparam{\DUrole{n}{name}\DUrole{o}{=}\DUrole{default_value}{\textquotesingle{}\textquotesingle{}}}\sphinxparamcomma \sphinxparam{\DUrole{n}{cmdargs}\DUrole{o}{=}\DUrole{default_value}{None}}\sphinxparamcomma \sphinxparam{\DUrole{n}{ptr}\DUrole{o}{=}\DUrole{default_value}{None}}\sphinxparamcomma \sphinxparam{\DUrole{n}{comm}\DUrole{o}{=}\DUrole{default_value}{None}}\sphinxparamcomma \sphinxparam{\DUrole{n}{verbose}\DUrole{o}{=}\DUrole{default_value}{False}}}{}
\pysigstopsignatures
\sphinxAtStartPar
This is a Python wrapper class around the lower\sphinxhyphen{}level
{\hyperref[\detokenize{Python_module:module-lammps}]{\sphinxcrossref{\sphinxcode{\sphinxupquote{lammps}}}}} class, exposing a more Python\sphinxhyphen{}like,
object\sphinxhyphen{}oriented interface for prototyping systems inside of IPython and
Jupyter notebooks.

\sphinxAtStartPar
It either creates its own instance of {\hyperref[\detokenize{Python_module:module-lammps}]{\sphinxcrossref{\sphinxcode{\sphinxupquote{lammps}}}}} or can be
initialized with an existing instance. The arguments are the same of the
lower\sphinxhyphen{}level interface. The original interface can still be accessed via
\sphinxcode{\sphinxupquote{PyLammps.lmp}}.
\begin{quote}\begin{description}
\sphinxlineitem{Parameters}\begin{itemize}
\item {} 
\sphinxAtStartPar
\sphinxstyleliteralstrong{\sphinxupquote{name}} (\sphinxstyleliteralemphasis{\sphinxupquote{string}}) \textendash{} “machine” name of the shared LAMMPS library (“mpi” loads \sphinxcode{\sphinxupquote{liblammps\_mpi.so}}, “” loads \sphinxcode{\sphinxupquote{liblammps.so}})

\item {} 
\sphinxAtStartPar
\sphinxstyleliteralstrong{\sphinxupquote{cmdargs}} (\sphinxstyleliteralemphasis{\sphinxupquote{list}}) \textendash{} list of command line arguments to be passed to the {\hyperref[\detokenize{Library_create:_CPPv411lammps_openiPPc8MPI_CommPPv}]{\sphinxcrossref{\sphinxcode{\sphinxupquote{lammps\_open()}}}}} function.  The executable name is automatically added.

\item {} 
\sphinxAtStartPar
\sphinxstyleliteralstrong{\sphinxupquote{ptr}} (\sphinxstyleliteralemphasis{\sphinxupquote{pointer}}) \textendash{} pointer to a LAMMPS C++ class instance when called from an embedded Python interpreter.  None means load symbols from shared library.

\item {} 
\sphinxAtStartPar
\sphinxstyleliteralstrong{\sphinxupquote{comm}} (\sphinxstyleliteralemphasis{\sphinxupquote{MPI\_Comm}}) \textendash{} MPI communicator (as provided by \sphinxhref{https://mpi4py.readthedocs.io/}{mpi4py}). \sphinxcode{\sphinxupquote{None}} means use \sphinxcode{\sphinxupquote{MPI\_COMM\_WORLD}} implicitly.

\item {} 
\sphinxAtStartPar
\sphinxstyleliteralstrong{\sphinxupquote{verbose}} (\sphinxstyleliteralemphasis{\sphinxupquote{bool}}) \textendash{} print all LAMMPS output to stdout

\end{itemize}

\sphinxlineitem{Variables}\begin{itemize}
\item {} 
\sphinxAtStartPar
\sphinxstyleliteralstrong{\sphinxupquote{lmp}} ({\hyperref[\detokenize{Python_module:module-lammps}]{\sphinxcrossref{\sphinxcode{\sphinxupquote{lammps}}}}}) \textendash{} instance of original LAMMPS Python interface

\item {} 
\sphinxAtStartPar
\sphinxstyleliteralstrong{\sphinxupquote{runs}} \textendash{} list of completed runs, each storing the thermo output

\end{itemize}

\end{description}\end{quote}
\index{close() (lammps.PyLammps method)@\spxentry{close()}\spxextra{lammps.PyLammps method}}

\begin{fulllineitems}
\phantomsection\label{\detokenize{Python_module:lammps.PyLammps.close}}
\pysigstartsignatures
\pysiglinewithargsret{\sphinxbfcode{\sphinxupquote{close}}}{}{}
\pysigstopsignatures
\sphinxAtStartPar
Explicitly delete a LAMMPS instance

\sphinxAtStartPar
This is a wrapper around the {\hyperref[\detokenize{Python_module:lammps.lammps.close}]{\sphinxcrossref{\sphinxcode{\sphinxupquote{lammps.close()}}}}} of the Python interface.

\end{fulllineitems}

\index{version() (lammps.PyLammps method)@\spxentry{version()}\spxextra{lammps.PyLammps method}}

\begin{fulllineitems}
\phantomsection\label{\detokenize{Python_module:lammps.PyLammps.version}}
\pysigstartsignatures
\pysiglinewithargsret{\sphinxbfcode{\sphinxupquote{version}}}{}{}
\pysigstopsignatures
\sphinxAtStartPar
Return a numerical representation of the LAMMPS version in use.

\sphinxAtStartPar
This is a wrapper around the {\hyperref[\detokenize{Python_module:lammps.lammps.version}]{\sphinxcrossref{\sphinxcode{\sphinxupquote{lammps.version()}}}}} function of the Python interface.
\begin{quote}\begin{description}
\sphinxlineitem{Returns}
\sphinxAtStartPar
version number

\sphinxlineitem{Return type}
\sphinxAtStartPar
int

\end{description}\end{quote}

\end{fulllineitems}

\index{file() (lammps.PyLammps method)@\spxentry{file()}\spxextra{lammps.PyLammps method}}

\begin{fulllineitems}
\phantomsection\label{\detokenize{Python_module:lammps.PyLammps.file}}
\pysigstartsignatures
\pysiglinewithargsret{\sphinxbfcode{\sphinxupquote{file}}}{\sphinxparam{\DUrole{n}{file}}}{}
\pysigstopsignatures
\sphinxAtStartPar
Read LAMMPS commands from a file.

\sphinxAtStartPar
This is a wrapper around the {\hyperref[\detokenize{Python_module:lammps.lammps.file}]{\sphinxcrossref{\sphinxcode{\sphinxupquote{lammps.file()}}}}} function of the Python interface.
\begin{quote}\begin{description}
\sphinxlineitem{Parameters}
\sphinxAtStartPar
\sphinxstyleliteralstrong{\sphinxupquote{path}} (\sphinxstyleliteralemphasis{\sphinxupquote{string}}) \textendash{} Name of the file/path with LAMMPS commands

\end{description}\end{quote}

\end{fulllineitems}

\index{enable\_cmd\_history (lammps.PyLammps property)@\spxentry{enable\_cmd\_history}\spxextra{lammps.PyLammps property}}

\begin{fulllineitems}
\phantomsection\label{\detokenize{Python_module:lammps.PyLammps.enable_cmd_history}}
\pysigstartsignatures
\pysigline{\sphinxbfcode{\sphinxupquote{property\DUrole{w}{ }}}\sphinxbfcode{\sphinxupquote{enable\_cmd\_history}}}
\pysigstopsignatures\begin{quote}\begin{description}
\sphinxlineitem{Getter}
\sphinxAtStartPar
Return whether command history is saved

\sphinxlineitem{Setter}
\sphinxAtStartPar
Set if command history should be saved

\sphinxlineitem{Type}
\sphinxAtStartPar
bool

\end{description}\end{quote}

\end{fulllineitems}

\index{write\_script() (lammps.PyLammps method)@\spxentry{write\_script()}\spxextra{lammps.PyLammps method}}

\begin{fulllineitems}
\phantomsection\label{\detokenize{Python_module:lammps.PyLammps.write_script}}
\pysigstartsignatures
\pysiglinewithargsret{\sphinxbfcode{\sphinxupquote{write\_script}}}{\sphinxparam{\DUrole{n}{filepath}}}{}
\pysigstopsignatures
\sphinxAtStartPar
Write LAMMPS script file containing all commands executed up until now
\begin{quote}\begin{description}
\sphinxlineitem{Parameters}
\sphinxAtStartPar
\sphinxstyleliteralstrong{\sphinxupquote{filepath}} (\sphinxstyleliteralemphasis{\sphinxupquote{string}}) \textendash{} path to script file that should be written

\end{description}\end{quote}

\end{fulllineitems}

\index{clear\_cmd\_history() (lammps.PyLammps method)@\spxentry{clear\_cmd\_history()}\spxextra{lammps.PyLammps method}}

\begin{fulllineitems}
\phantomsection\label{\detokenize{Python_module:lammps.PyLammps.clear_cmd_history}}
\pysigstartsignatures
\pysiglinewithargsret{\sphinxbfcode{\sphinxupquote{clear\_cmd\_history}}}{}{}
\pysigstopsignatures
\sphinxAtStartPar
Clear LAMMPS command history up to this point

\end{fulllineitems}

\index{append\_cmd\_history() (lammps.PyLammps method)@\spxentry{append\_cmd\_history()}\spxextra{lammps.PyLammps method}}

\begin{fulllineitems}
\phantomsection\label{\detokenize{Python_module:lammps.PyLammps.append_cmd_history}}
\pysigstartsignatures
\pysiglinewithargsret{\sphinxbfcode{\sphinxupquote{append\_cmd\_history}}}{\sphinxparam{\DUrole{n}{cmd}}}{}
\pysigstopsignatures
\sphinxAtStartPar
Commands will be added to the command history but not executed.

\sphinxAtStartPar
Add \sphinxtitleref{commands} only to the command history, but do not execute them, so that you can
conveniently create LAMMPS input files, using
{\hyperref[\detokenize{Python_module:lammps.PyLammps.write_script}]{\sphinxcrossref{\sphinxcode{\sphinxupquote{PyLammps.write\_script()}}}}}.

\end{fulllineitems}

\index{command() (lammps.PyLammps method)@\spxentry{command()}\spxextra{lammps.PyLammps method}}

\begin{fulllineitems}
\phantomsection\label{\detokenize{Python_module:lammps.PyLammps.command}}
\pysigstartsignatures
\pysiglinewithargsret{\sphinxbfcode{\sphinxupquote{command}}}{\sphinxparam{\DUrole{n}{cmd}}}{}
\pysigstopsignatures
\sphinxAtStartPar
Execute LAMMPS command

\sphinxAtStartPar
If {\hyperref[\detokenize{Python_module:lammps.PyLammps.enable_cmd_history}]{\sphinxcrossref{\sphinxcode{\sphinxupquote{PyLammps.enable\_cmd\_history}}}}} is set to \sphinxcode{\sphinxupquote{True}}, commands executed
will be recorded. The entire command history can be written to a file using
{\hyperref[\detokenize{Python_module:lammps.PyLammps.write_script}]{\sphinxcrossref{\sphinxcode{\sphinxupquote{PyLammps.write\_script()}}}}}. To clear the command history, use
{\hyperref[\detokenize{Python_module:lammps.PyLammps.clear_cmd_history}]{\sphinxcrossref{\sphinxcode{\sphinxupquote{PyLammps.clear\_cmd\_history()}}}}}.
\begin{quote}\begin{description}
\sphinxlineitem{Parameters}
\sphinxAtStartPar
\sphinxstyleliteralstrong{\sphinxupquote{cmd}} \textendash{} command string that should be executed

\sphinxlineitem{Type}
\sphinxAtStartPar
cmd: string

\end{description}\end{quote}

\end{fulllineitems}

\index{run() (lammps.PyLammps method)@\spxentry{run()}\spxextra{lammps.PyLammps method}}

\begin{fulllineitems}
\phantomsection\label{\detokenize{Python_module:lammps.PyLammps.run}}
\pysigstartsignatures
\pysiglinewithargsret{\sphinxbfcode{\sphinxupquote{run}}}{\sphinxparam{\DUrole{o}{*}\DUrole{n}{args}}\sphinxparamcomma \sphinxparam{\DUrole{o}{**}\DUrole{n}{kwargs}}}{}
\pysigstopsignatures
\sphinxAtStartPar
Execute LAMMPS run command with given arguments

\sphinxAtStartPar
Thermo data of the run is recorded and saved as new entry in
\sphinxcode{\sphinxupquote{PyLammps.runs}}. The latest run can be retrieved by
{\hyperref[\detokenize{Python_module:lammps.PyLammps.last_run}]{\sphinxcrossref{\sphinxcode{\sphinxupquote{PyLammps.last\_run}}}}}.

\sphinxAtStartPar
Note, for recording of all thermo steps during a run, the PYTHON package
needs to be enabled in LAMMPS. Otherwise, it will only capture the final
timestep.

\end{fulllineitems}

\index{last\_run (lammps.PyLammps property)@\spxentry{last\_run}\spxextra{lammps.PyLammps property}}

\begin{fulllineitems}
\phantomsection\label{\detokenize{Python_module:lammps.PyLammps.last_run}}
\pysigstartsignatures
\pysigline{\sphinxbfcode{\sphinxupquote{property\DUrole{w}{ }}}\sphinxbfcode{\sphinxupquote{last\_run}}}
\pysigstopsignatures
\sphinxAtStartPar
Return data produced of last completed run command
\begin{quote}\begin{description}
\sphinxlineitem{Getter}
\sphinxAtStartPar
Returns an object containing information about the last run command

\sphinxlineitem{Type}
\sphinxAtStartPar
dict

\end{description}\end{quote}

\end{fulllineitems}

\index{atoms (lammps.PyLammps property)@\spxentry{atoms}\spxextra{lammps.PyLammps property}}

\begin{fulllineitems}
\phantomsection\label{\detokenize{Python_module:lammps.PyLammps.atoms}}
\pysigstartsignatures
\pysigline{\sphinxbfcode{\sphinxupquote{property\DUrole{w}{ }}}\sphinxbfcode{\sphinxupquote{atoms}}}
\pysigstopsignatures
\sphinxAtStartPar
All atoms of this LAMMPS instance
\begin{quote}\begin{description}
\sphinxlineitem{Getter}
\sphinxAtStartPar
Returns a list of atoms currently in the system

\sphinxlineitem{Type}
\sphinxAtStartPar
{\hyperref[\detokenize{Python_module:lammps.AtomList}]{\sphinxcrossref{AtomList}}}

\end{description}\end{quote}

\end{fulllineitems}

\index{system (lammps.PyLammps property)@\spxentry{system}\spxextra{lammps.PyLammps property}}

\begin{fulllineitems}
\phantomsection\label{\detokenize{Python_module:lammps.PyLammps.system}}
\pysigstartsignatures
\pysigline{\sphinxbfcode{\sphinxupquote{property\DUrole{w}{ }}}\sphinxbfcode{\sphinxupquote{system}}}
\pysigstopsignatures
\sphinxAtStartPar
The system state of this LAMMPS instance
\begin{quote}\begin{description}
\sphinxlineitem{Getter}
\sphinxAtStartPar
Returns an object with properties storing the current system state

\sphinxlineitem{Type}
\sphinxAtStartPar
namedtuple

\end{description}\end{quote}

\end{fulllineitems}

\index{communication (lammps.PyLammps property)@\spxentry{communication}\spxextra{lammps.PyLammps property}}

\begin{fulllineitems}
\phantomsection\label{\detokenize{Python_module:lammps.PyLammps.communication}}
\pysigstartsignatures
\pysigline{\sphinxbfcode{\sphinxupquote{property\DUrole{w}{ }}}\sphinxbfcode{\sphinxupquote{communication}}}
\pysigstopsignatures
\sphinxAtStartPar
The communication state of this LAMMPS instance
\begin{quote}\begin{description}
\sphinxlineitem{Getter}
\sphinxAtStartPar
Returns an object with properties storing the current communication state

\sphinxlineitem{Type}
\sphinxAtStartPar
namedtuple

\end{description}\end{quote}

\end{fulllineitems}

\index{computes (lammps.PyLammps property)@\spxentry{computes}\spxextra{lammps.PyLammps property}}

\begin{fulllineitems}
\phantomsection\label{\detokenize{Python_module:lammps.PyLammps.computes}}
\pysigstartsignatures
\pysigline{\sphinxbfcode{\sphinxupquote{property\DUrole{w}{ }}}\sphinxbfcode{\sphinxupquote{computes}}}
\pysigstopsignatures
\sphinxAtStartPar
The list of active computes of this LAMMPS instance
\begin{quote}\begin{description}
\sphinxlineitem{Getter}
\sphinxAtStartPar
Returns a list of computes that are currently active in this LAMMPS instance

\sphinxlineitem{Type}
\sphinxAtStartPar
list

\end{description}\end{quote}

\end{fulllineitems}

\index{dumps (lammps.PyLammps property)@\spxentry{dumps}\spxextra{lammps.PyLammps property}}

\begin{fulllineitems}
\phantomsection\label{\detokenize{Python_module:lammps.PyLammps.dumps}}
\pysigstartsignatures
\pysigline{\sphinxbfcode{\sphinxupquote{property\DUrole{w}{ }}}\sphinxbfcode{\sphinxupquote{dumps}}}
\pysigstopsignatures
\sphinxAtStartPar
The list of active dumps of this LAMMPS instance
\begin{quote}\begin{description}
\sphinxlineitem{Getter}
\sphinxAtStartPar
Returns a list of dumps that are currently active in this LAMMPS instance

\sphinxlineitem{Type}
\sphinxAtStartPar
list

\end{description}\end{quote}

\end{fulllineitems}

\index{fixes (lammps.PyLammps property)@\spxentry{fixes}\spxextra{lammps.PyLammps property}}

\begin{fulllineitems}
\phantomsection\label{\detokenize{Python_module:lammps.PyLammps.fixes}}
\pysigstartsignatures
\pysigline{\sphinxbfcode{\sphinxupquote{property\DUrole{w}{ }}}\sphinxbfcode{\sphinxupquote{fixes}}}
\pysigstopsignatures
\sphinxAtStartPar
The list of active fixes of this LAMMPS instance
\begin{quote}\begin{description}
\sphinxlineitem{Getter}
\sphinxAtStartPar
Returns a list of fixes that are currently active in this LAMMPS instance

\sphinxlineitem{Type}
\sphinxAtStartPar
list

\end{description}\end{quote}

\end{fulllineitems}

\index{groups (lammps.PyLammps property)@\spxentry{groups}\spxextra{lammps.PyLammps property}}

\begin{fulllineitems}
\phantomsection\label{\detokenize{Python_module:lammps.PyLammps.groups}}
\pysigstartsignatures
\pysigline{\sphinxbfcode{\sphinxupquote{property\DUrole{w}{ }}}\sphinxbfcode{\sphinxupquote{groups}}}
\pysigstopsignatures
\sphinxAtStartPar
The list of active atom groups of this LAMMPS instance
\begin{quote}\begin{description}
\sphinxlineitem{Getter}
\sphinxAtStartPar
Returns a list of atom groups that are currently active in this LAMMPS instance

\sphinxlineitem{Type}
\sphinxAtStartPar
list

\end{description}\end{quote}

\end{fulllineitems}

\index{variables (lammps.PyLammps property)@\spxentry{variables}\spxextra{lammps.PyLammps property}}

\begin{fulllineitems}
\phantomsection\label{\detokenize{Python_module:lammps.PyLammps.variables}}
\pysigstartsignatures
\pysigline{\sphinxbfcode{\sphinxupquote{property\DUrole{w}{ }}}\sphinxbfcode{\sphinxupquote{variables}}}
\pysigstopsignatures
\sphinxAtStartPar
Returns a dictionary of all variables defined in the current LAMMPS instance
\begin{quote}\begin{description}
\sphinxlineitem{Getter}
\sphinxAtStartPar
Returns a dictionary of all variables that are defined in this LAMMPS instance

\sphinxlineitem{Type}
\sphinxAtStartPar
dict

\end{description}\end{quote}

\end{fulllineitems}

\index{eval() (lammps.PyLammps method)@\spxentry{eval()}\spxextra{lammps.PyLammps method}}

\begin{fulllineitems}
\phantomsection\label{\detokenize{Python_module:lammps.PyLammps.eval}}
\pysigstartsignatures
\pysiglinewithargsret{\sphinxbfcode{\sphinxupquote{eval}}}{\sphinxparam{\DUrole{n}{expr}}}{}
\pysigstopsignatures
\sphinxAtStartPar
Evaluate LAMMPS input file expression.

\sphinxAtStartPar
This is equivalent to using immediate variable expressions in the format “\$(…)”
in the LAMMPS input and will return the result of that expression.

\begin{sphinxadmonition}{warning}{Warning:}
\sphinxAtStartPar
This function is only supported on MPI rank 0.  Calling it from a different
MPI rank will raise an exception.
\end{sphinxadmonition}
\begin{quote}\begin{description}
\sphinxlineitem{Parameters}
\sphinxAtStartPar
\sphinxstyleliteralstrong{\sphinxupquote{expr}} (\sphinxstyleliteralemphasis{\sphinxupquote{string}}) \textendash{} the expression string that should be evaluated inside of LAMMPS

\sphinxlineitem{Returns}
\sphinxAtStartPar
the value of the evaluated expression

\sphinxlineitem{Return type}
\sphinxAtStartPar
float if numeric, string otherwise

\end{description}\end{quote}

\end{fulllineitems}

\index{lmp\_print() (lammps.PyLammps method)@\spxentry{lmp\_print()}\spxextra{lammps.PyLammps method}}

\begin{fulllineitems}
\phantomsection\label{\detokenize{Python_module:lammps.PyLammps.lmp_print}}
\pysigstartsignatures
\pysiglinewithargsret{\sphinxbfcode{\sphinxupquote{lmp\_print}}}{\sphinxparam{\DUrole{n}{s}}}{}
\pysigstopsignatures
\sphinxAtStartPar
needed for Python2 compatibility, since print is a reserved keyword

\end{fulllineitems}


\end{fulllineitems}

\index{AtomList (class in lammps)@\spxentry{AtomList}\spxextra{class in lammps}}

\begin{fulllineitems}
\phantomsection\label{\detokenize{Python_module:lammps.AtomList}}
\pysigstartsignatures
\pysiglinewithargsret{\sphinxbfcode{\sphinxupquote{class\DUrole{w}{ }}}\sphinxcode{\sphinxupquote{lammps.}}\sphinxbfcode{\sphinxupquote{AtomList}}}{\sphinxparam{\DUrole{n}{pylammps\_instance}}}{}
\pysigstopsignatures
\sphinxAtStartPar
A dynamic list of atoms that returns either an {\hyperref[\detokenize{Python_module:lammps.Atom}]{\sphinxcrossref{\sphinxcode{\sphinxupquote{Atom}}}}} or
{\hyperref[\detokenize{Python_module:lammps.Atom2D}]{\sphinxcrossref{\sphinxcode{\sphinxupquote{Atom2D}}}}} instance for each atom. Instances are only allocated
when accessed.
\begin{quote}\begin{description}
\sphinxlineitem{Variables}\begin{itemize}
\item {} 
\sphinxAtStartPar
\sphinxstyleliteralstrong{\sphinxupquote{natoms}} \textendash{} total number of atoms

\item {} 
\sphinxAtStartPar
\sphinxstyleliteralstrong{\sphinxupquote{dimensions}} \textendash{} number of dimensions in system

\end{itemize}

\end{description}\end{quote}

\end{fulllineitems}

\index{Atom (class in lammps)@\spxentry{Atom}\spxextra{class in lammps}}

\begin{fulllineitems}
\phantomsection\label{\detokenize{Python_module:lammps.Atom}}
\pysigstartsignatures
\pysiglinewithargsret{\sphinxbfcode{\sphinxupquote{class\DUrole{w}{ }}}\sphinxcode{\sphinxupquote{lammps.}}\sphinxbfcode{\sphinxupquote{Atom}}}{\sphinxparam{\DUrole{n}{pylammps\_instance}}\sphinxparamcomma \sphinxparam{\DUrole{n}{index}}}{}
\pysigstopsignatures
\sphinxAtStartPar
A wrapper class then represents a single atom inside of LAMMPS

\sphinxAtStartPar
It provides access to properties of the atom and allows you to change some of them.
\index{id (lammps.Atom property)@\spxentry{id}\spxextra{lammps.Atom property}}

\begin{fulllineitems}
\phantomsection\label{\detokenize{Python_module:lammps.Atom.id}}
\pysigstartsignatures
\pysigline{\sphinxbfcode{\sphinxupquote{property\DUrole{w}{ }}}\sphinxbfcode{\sphinxupquote{id}}}
\pysigstopsignatures
\sphinxAtStartPar
Return the atom ID
\begin{quote}\begin{description}
\sphinxlineitem{Type}
\sphinxAtStartPar
int

\end{description}\end{quote}

\end{fulllineitems}

\index{type (lammps.Atom property)@\spxentry{type}\spxextra{lammps.Atom property}}

\begin{fulllineitems}
\phantomsection\label{\detokenize{Python_module:lammps.Atom.type}}
\pysigstartsignatures
\pysigline{\sphinxbfcode{\sphinxupquote{property\DUrole{w}{ }}}\sphinxbfcode{\sphinxupquote{type}}}
\pysigstopsignatures
\sphinxAtStartPar
Return the atom type
\begin{quote}\begin{description}
\sphinxlineitem{Type}
\sphinxAtStartPar
int

\end{description}\end{quote}

\end{fulllineitems}

\index{mol (lammps.Atom property)@\spxentry{mol}\spxextra{lammps.Atom property}}

\begin{fulllineitems}
\phantomsection\label{\detokenize{Python_module:lammps.Atom.mol}}
\pysigstartsignatures
\pysigline{\sphinxbfcode{\sphinxupquote{property\DUrole{w}{ }}}\sphinxbfcode{\sphinxupquote{mol}}}
\pysigstopsignatures
\sphinxAtStartPar
Return the atom molecule index
\begin{quote}\begin{description}
\sphinxlineitem{Type}
\sphinxAtStartPar
int

\end{description}\end{quote}

\end{fulllineitems}

\index{mass (lammps.Atom property)@\spxentry{mass}\spxextra{lammps.Atom property}}

\begin{fulllineitems}
\phantomsection\label{\detokenize{Python_module:lammps.Atom.mass}}
\pysigstartsignatures
\pysigline{\sphinxbfcode{\sphinxupquote{property\DUrole{w}{ }}}\sphinxbfcode{\sphinxupquote{mass}}}
\pysigstopsignatures
\sphinxAtStartPar
Return the atom mass as a per\sphinxhyphen{}atom property.
This returns either the per\sphinxhyphen{}type mass or the per\sphinxhyphen{}atom
mass (AKA ‘rmass’) depending on what is available with
preference for the per\sphinxhyphen{}atom mass.

\sphinxAtStartPar
\DUrole{versionmodified,changed}{Changed in version 17Apr2024: }Support both per\sphinxhyphen{}type and per\sphinxhyphen{}atom masses. With
per\sphinxhyphen{}type return “mass{[}type{[}i{]}{]}” else return “rmass{[}i{]}”.
Per\sphinxhyphen{}atom mass is preferred if available.
\begin{quote}\begin{description}
\sphinxlineitem{Type}
\sphinxAtStartPar
float

\end{description}\end{quote}

\end{fulllineitems}

\index{radius (lammps.Atom property)@\spxentry{radius}\spxextra{lammps.Atom property}}

\begin{fulllineitems}
\phantomsection\label{\detokenize{Python_module:lammps.Atom.radius}}
\pysigstartsignatures
\pysigline{\sphinxbfcode{\sphinxupquote{property\DUrole{w}{ }}}\sphinxbfcode{\sphinxupquote{radius}}}
\pysigstopsignatures
\sphinxAtStartPar
Return the particle radius
\begin{quote}\begin{description}
\sphinxlineitem{Type}
\sphinxAtStartPar
float

\end{description}\end{quote}

\end{fulllineitems}

\index{position (lammps.Atom property)@\spxentry{position}\spxextra{lammps.Atom property}}

\begin{fulllineitems}
\phantomsection\label{\detokenize{Python_module:lammps.Atom.position}}
\pysigstartsignatures
\pysigline{\sphinxbfcode{\sphinxupquote{property\DUrole{w}{ }}}\sphinxbfcode{\sphinxupquote{position}}}
\pysigstopsignatures\begin{quote}\begin{description}
\sphinxlineitem{Getter}
\sphinxAtStartPar
Return position of atom

\sphinxlineitem{Setter}
\sphinxAtStartPar
Set position of atom

\sphinxlineitem{Type}
\sphinxAtStartPar
numpy.array (float, float, float)

\end{description}\end{quote}

\end{fulllineitems}

\index{velocity (lammps.Atom property)@\spxentry{velocity}\spxextra{lammps.Atom property}}

\begin{fulllineitems}
\phantomsection\label{\detokenize{Python_module:lammps.Atom.velocity}}
\pysigstartsignatures
\pysigline{\sphinxbfcode{\sphinxupquote{property\DUrole{w}{ }}}\sphinxbfcode{\sphinxupquote{velocity}}}
\pysigstopsignatures\begin{quote}\begin{description}
\sphinxlineitem{Getter}
\sphinxAtStartPar
Return velocity of atom

\sphinxlineitem{Setter}
\sphinxAtStartPar
Set velocity of atom

\sphinxlineitem{Type}
\sphinxAtStartPar
numpy.array (float, float, float)

\end{description}\end{quote}

\end{fulllineitems}

\index{force (lammps.Atom property)@\spxentry{force}\spxextra{lammps.Atom property}}

\begin{fulllineitems}
\phantomsection\label{\detokenize{Python_module:lammps.Atom.force}}
\pysigstartsignatures
\pysigline{\sphinxbfcode{\sphinxupquote{property\DUrole{w}{ }}}\sphinxbfcode{\sphinxupquote{force}}}
\pysigstopsignatures
\sphinxAtStartPar
Return the total force acting on the atom
\begin{quote}\begin{description}
\sphinxlineitem{Type}
\sphinxAtStartPar
numpy.array (float, float, float)

\end{description}\end{quote}

\end{fulllineitems}

\index{omega (lammps.Atom property)@\spxentry{omega}\spxextra{lammps.Atom property}}

\begin{fulllineitems}
\phantomsection\label{\detokenize{Python_module:lammps.Atom.omega}}
\pysigstartsignatures
\pysigline{\sphinxbfcode{\sphinxupquote{property\DUrole{w}{ }}}\sphinxbfcode{\sphinxupquote{omega}}}
\pysigstopsignatures
\sphinxAtStartPar
Return the rotational velocity of the particle
\begin{quote}\begin{description}
\sphinxlineitem{Type}
\sphinxAtStartPar
numpy.array (float, float, float)

\end{description}\end{quote}

\end{fulllineitems}

\index{torque (lammps.Atom property)@\spxentry{torque}\spxextra{lammps.Atom property}}

\begin{fulllineitems}
\phantomsection\label{\detokenize{Python_module:lammps.Atom.torque}}
\pysigstartsignatures
\pysigline{\sphinxbfcode{\sphinxupquote{property\DUrole{w}{ }}}\sphinxbfcode{\sphinxupquote{torque}}}
\pysigstopsignatures
\sphinxAtStartPar
Return the total torque acting on the particle
\begin{quote}\begin{description}
\sphinxlineitem{Type}
\sphinxAtStartPar
numpy.array (float, float, float)

\end{description}\end{quote}

\end{fulllineitems}

\index{angular\_momentum (lammps.Atom property)@\spxentry{angular\_momentum}\spxextra{lammps.Atom property}}

\begin{fulllineitems}
\phantomsection\label{\detokenize{Python_module:lammps.Atom.angular_momentum}}
\pysigstartsignatures
\pysigline{\sphinxbfcode{\sphinxupquote{property\DUrole{w}{ }}}\sphinxbfcode{\sphinxupquote{angular\_momentum}}}
\pysigstopsignatures
\sphinxAtStartPar
Return the angular momentum of the particle
\begin{quote}\begin{description}
\sphinxlineitem{Type}
\sphinxAtStartPar
numpy.array (float, float, float)

\end{description}\end{quote}

\end{fulllineitems}

\index{charge (lammps.Atom property)@\spxentry{charge}\spxextra{lammps.Atom property}}

\begin{fulllineitems}
\phantomsection\label{\detokenize{Python_module:lammps.Atom.charge}}
\pysigstartsignatures
\pysigline{\sphinxbfcode{\sphinxupquote{property\DUrole{w}{ }}}\sphinxbfcode{\sphinxupquote{charge}}}
\pysigstopsignatures
\sphinxAtStartPar
Return the atom charge
\begin{quote}\begin{description}
\sphinxlineitem{Type}
\sphinxAtStartPar
float

\end{description}\end{quote}

\end{fulllineitems}


\end{fulllineitems}

\index{Atom2D (class in lammps)@\spxentry{Atom2D}\spxextra{class in lammps}}

\begin{fulllineitems}
\phantomsection\label{\detokenize{Python_module:lammps.Atom2D}}
\pysigstartsignatures
\pysiglinewithargsret{\sphinxbfcode{\sphinxupquote{class\DUrole{w}{ }}}\sphinxcode{\sphinxupquote{lammps.}}\sphinxbfcode{\sphinxupquote{Atom2D}}}{\sphinxparam{\DUrole{n}{pylammps\_instance}}\sphinxparamcomma \sphinxparam{\DUrole{n}{index}}}{}
\pysigstopsignatures
\sphinxAtStartPar
A wrapper class then represents a single 2D atom inside of LAMMPS

\sphinxAtStartPar
Inherits all properties from the {\hyperref[\detokenize{Python_module:lammps.Atom}]{\sphinxcrossref{\sphinxcode{\sphinxupquote{Atom}}}}} class, but returns 2D versions
of position, velocity, and force.

\sphinxAtStartPar
It provides access to properties of the atom and allows you to change some of them.
\index{position (lammps.Atom2D property)@\spxentry{position}\spxextra{lammps.Atom2D property}}

\begin{fulllineitems}
\phantomsection\label{\detokenize{Python_module:lammps.Atom2D.position}}
\pysigstartsignatures
\pysigline{\sphinxbfcode{\sphinxupquote{property\DUrole{w}{ }}}\sphinxbfcode{\sphinxupquote{position}}}
\pysigstopsignatures
\sphinxAtStartPar
Access to coordinates of an atom
\begin{quote}\begin{description}
\sphinxlineitem{Getter}
\sphinxAtStartPar
Return position of atom

\sphinxlineitem{Setter}
\sphinxAtStartPar
Set position of atom

\sphinxlineitem{Type}
\sphinxAtStartPar
numpy.array (float, float)

\end{description}\end{quote}

\end{fulllineitems}

\index{velocity (lammps.Atom2D property)@\spxentry{velocity}\spxextra{lammps.Atom2D property}}

\begin{fulllineitems}
\phantomsection\label{\detokenize{Python_module:lammps.Atom2D.velocity}}
\pysigstartsignatures
\pysigline{\sphinxbfcode{\sphinxupquote{property\DUrole{w}{ }}}\sphinxbfcode{\sphinxupquote{velocity}}}
\pysigstopsignatures
\sphinxAtStartPar
Access to velocity of an atom
\begin{quote}\begin{description}
\sphinxlineitem{Getter}
\sphinxAtStartPar
Return velocity of atom

\sphinxlineitem{Setter}
\sphinxAtStartPar
Set velocity of atom

\sphinxlineitem{Type}
\sphinxAtStartPar
numpy.array (float, float)

\end{description}\end{quote}

\end{fulllineitems}

\index{force (lammps.Atom2D property)@\spxentry{force}\spxextra{lammps.Atom2D property}}

\begin{fulllineitems}
\phantomsection\label{\detokenize{Python_module:lammps.Atom2D.force}}
\pysigstartsignatures
\pysigline{\sphinxbfcode{\sphinxupquote{property\DUrole{w}{ }}}\sphinxbfcode{\sphinxupquote{force}}}
\pysigstopsignatures
\sphinxAtStartPar
Access to force of an atom
\begin{quote}\begin{description}
\sphinxlineitem{Getter}
\sphinxAtStartPar
Return force of atom

\sphinxlineitem{Setter}
\sphinxAtStartPar
Set force of atom

\sphinxlineitem{Type}
\sphinxAtStartPar
numpy.array (float, float)

\end{description}\end{quote}

\end{fulllineitems}


\end{fulllineitems}



\bigskip\hrule\bigskip



\subsection{The \sphinxstyleliteralintitle{\sphinxupquote{IPyLammps}} class API}
\label{\detokenize{Python_module:the-ipylammps-class-api}}
\sphinxAtStartPar
The {\hyperref[\detokenize{Python_module:lammps.PyLammps}]{\sphinxcrossref{\sphinxcode{\sphinxupquote{IPyLammps}}}}} class is an extension of
{\hyperref[\detokenize{Python_module:lammps.PyLammps}]{\sphinxcrossref{\sphinxcode{\sphinxupquote{PyLammps}}}}}, adding additional functions to
quickly display visualizations such as images and videos inside of IPython.
See the {\hyperref[\detokenize{Howto_pylammps::doc}]{\sphinxcrossref{\DUrole{doc}{PyLammps Tutorial}}}} for examples.
\index{IPyLammps (class in lammps)@\spxentry{IPyLammps}\spxextra{class in lammps}}

\begin{fulllineitems}
\phantomsection\label{\detokenize{Python_module:lammps.IPyLammps}}
\pysigstartsignatures
\pysiglinewithargsret{\sphinxbfcode{\sphinxupquote{class\DUrole{w}{ }}}\sphinxcode{\sphinxupquote{lammps.}}\sphinxbfcode{\sphinxupquote{IPyLammps}}}{\sphinxparam{\DUrole{n}{name}\DUrole{o}{=}\DUrole{default_value}{\textquotesingle{}\textquotesingle{}}}\sphinxparamcomma \sphinxparam{\DUrole{n}{cmdargs}\DUrole{o}{=}\DUrole{default_value}{None}}\sphinxparamcomma \sphinxparam{\DUrole{n}{ptr}\DUrole{o}{=}\DUrole{default_value}{None}}\sphinxparamcomma \sphinxparam{\DUrole{n}{comm}\DUrole{o}{=}\DUrole{default_value}{None}}}{}
\pysigstopsignatures
\sphinxAtStartPar
IPython wrapper for LAMMPS which adds embedded graphics capabilities to PyLammps interface

\sphinxAtStartPar
It either creates its own instance of {\hyperref[\detokenize{Python_module:module-lammps}]{\sphinxcrossref{\sphinxcode{\sphinxupquote{lammps}}}}} or can be
initialized with an existing instance. The arguments are the same of the
lower\sphinxhyphen{}level interface. The original interface can still be accessed via
\sphinxcode{\sphinxupquote{PyLammps.lmp}}.
\begin{quote}\begin{description}
\sphinxlineitem{Parameters}\begin{itemize}
\item {} 
\sphinxAtStartPar
\sphinxstyleliteralstrong{\sphinxupquote{name}} (\sphinxstyleliteralemphasis{\sphinxupquote{string}}) \textendash{} “machine” name of the shared LAMMPS library (“mpi” loads \sphinxcode{\sphinxupquote{liblammps\_mpi.so}}, “” loads \sphinxcode{\sphinxupquote{liblammps.so}})

\item {} 
\sphinxAtStartPar
\sphinxstyleliteralstrong{\sphinxupquote{cmdargs}} (\sphinxstyleliteralemphasis{\sphinxupquote{list}}) \textendash{} list of command line arguments to be passed to the {\hyperref[\detokenize{Library_create:_CPPv411lammps_openiPPc8MPI_CommPPv}]{\sphinxcrossref{\sphinxcode{\sphinxupquote{lammps\_open()}}}}} function.  The executable name is automatically added.

\item {} 
\sphinxAtStartPar
\sphinxstyleliteralstrong{\sphinxupquote{ptr}} (\sphinxstyleliteralemphasis{\sphinxupquote{pointer}}) \textendash{} pointer to a LAMMPS C++ class instance when called from an embedded Python interpreter.  None means load symbols from shared library.

\item {} 
\sphinxAtStartPar
\sphinxstyleliteralstrong{\sphinxupquote{comm}} (\sphinxstyleliteralemphasis{\sphinxupquote{MPI\_Comm}}) \textendash{} MPI communicator (as provided by \sphinxhref{https://mpi4py.readthedocs.io/}{mpi4py}). \sphinxcode{\sphinxupquote{None}} means use \sphinxcode{\sphinxupquote{MPI\_COMM\_WORLD}} implicitly.

\end{itemize}

\end{description}\end{quote}
\index{image() (lammps.IPyLammps method)@\spxentry{image()}\spxextra{lammps.IPyLammps method}}

\begin{fulllineitems}
\phantomsection\label{\detokenize{Python_module:lammps.IPyLammps.image}}
\pysigstartsignatures
\pysiglinewithargsret{\sphinxbfcode{\sphinxupquote{image}}}{\sphinxparam{\DUrole{n}{filename}\DUrole{o}{=}\DUrole{default_value}{\textquotesingle{}snapshot.png\textquotesingle{}}}\sphinxparamcomma \sphinxparam{\DUrole{n}{group}\DUrole{o}{=}\DUrole{default_value}{\textquotesingle{}all\textquotesingle{}}}\sphinxparamcomma \sphinxparam{\DUrole{n}{color}\DUrole{o}{=}\DUrole{default_value}{\textquotesingle{}type\textquotesingle{}}}\sphinxparamcomma \sphinxparam{\DUrole{n}{diameter}\DUrole{o}{=}\DUrole{default_value}{\textquotesingle{}type\textquotesingle{}}}\sphinxparamcomma \sphinxparam{\DUrole{n}{size}\DUrole{o}{=}\DUrole{default_value}{None}}\sphinxparamcomma \sphinxparam{\DUrole{n}{view}\DUrole{o}{=}\DUrole{default_value}{None}}\sphinxparamcomma \sphinxparam{\DUrole{n}{center}\DUrole{o}{=}\DUrole{default_value}{None}}\sphinxparamcomma \sphinxparam{\DUrole{n}{up}\DUrole{o}{=}\DUrole{default_value}{None}}\sphinxparamcomma \sphinxparam{\DUrole{n}{zoom}\DUrole{o}{=}\DUrole{default_value}{1.0}}\sphinxparamcomma \sphinxparam{\DUrole{n}{background\_color}\DUrole{o}{=}\DUrole{default_value}{\textquotesingle{}white\textquotesingle{}}}}{}
\pysigstopsignatures
\sphinxAtStartPar
Generate image using write\_dump command and display it

\sphinxAtStartPar
See {\hyperref[\detokenize{dump_image::doc}]{\sphinxcrossref{\DUrole{doc}{dump image}}}} for more information.
\begin{quote}\begin{description}
\sphinxlineitem{Parameters}\begin{itemize}
\item {} 
\sphinxAtStartPar
\sphinxstyleliteralstrong{\sphinxupquote{filename}} (\sphinxstyleliteralemphasis{\sphinxupquote{string}}) \textendash{} Name of the image file that should be generated. The extension determines whether it is PNG or JPEG

\item {} 
\sphinxAtStartPar
\sphinxstyleliteralstrong{\sphinxupquote{group}} (\sphinxstyleliteralemphasis{\sphinxupquote{string}}) \textendash{} the group of atoms write\_image should use

\item {} 
\sphinxAtStartPar
\sphinxstyleliteralstrong{\sphinxupquote{color}} (\sphinxstyleliteralemphasis{\sphinxupquote{string}}) \textendash{} name of property used to determine color

\item {} 
\sphinxAtStartPar
\sphinxstyleliteralstrong{\sphinxupquote{diameter}} (\sphinxstyleliteralemphasis{\sphinxupquote{string}}) \textendash{} name of property used to determine atom diameter

\item {} 
\sphinxAtStartPar
\sphinxstyleliteralstrong{\sphinxupquote{size}} (\sphinxstyleliteralemphasis{\sphinxupquote{tuple}}\sphinxstyleliteralemphasis{\sphinxupquote{ (}}\sphinxstyleliteralemphasis{\sphinxupquote{width}}\sphinxstyleliteralemphasis{\sphinxupquote{, }}\sphinxstyleliteralemphasis{\sphinxupquote{height}}\sphinxstyleliteralemphasis{\sphinxupquote{)}}) \textendash{} dimensions of image

\item {} 
\sphinxAtStartPar
\sphinxstyleliteralstrong{\sphinxupquote{view}} (\sphinxstyleliteralemphasis{\sphinxupquote{tuple}}\sphinxstyleliteralemphasis{\sphinxupquote{ (}}\sphinxstyleliteralemphasis{\sphinxupquote{theta}}\sphinxstyleliteralemphasis{\sphinxupquote{, }}\sphinxstyleliteralemphasis{\sphinxupquote{phi}}\sphinxstyleliteralemphasis{\sphinxupquote{)}}) \textendash{} view parameters

\item {} 
\sphinxAtStartPar
\sphinxstyleliteralstrong{\sphinxupquote{center}} (\sphinxstyleliteralemphasis{\sphinxupquote{tuple}}\sphinxstyleliteralemphasis{\sphinxupquote{ (}}\sphinxstyleliteralemphasis{\sphinxupquote{flag}}\sphinxstyleliteralemphasis{\sphinxupquote{, }}\sphinxstyleliteralemphasis{\sphinxupquote{center\_x}}\sphinxstyleliteralemphasis{\sphinxupquote{, }}\sphinxstyleliteralemphasis{\sphinxupquote{center\_y}}\sphinxstyleliteralemphasis{\sphinxupquote{, }}\sphinxstyleliteralemphasis{\sphinxupquote{center\_z}}\sphinxstyleliteralemphasis{\sphinxupquote{)}}) \textendash{} center parameters

\item {} 
\sphinxAtStartPar
\sphinxstyleliteralstrong{\sphinxupquote{up}} (\sphinxstyleliteralemphasis{\sphinxupquote{tuple}}\sphinxstyleliteralemphasis{\sphinxupquote{ (}}\sphinxstyleliteralemphasis{\sphinxupquote{up\_x}}\sphinxstyleliteralemphasis{\sphinxupquote{, }}\sphinxstyleliteralemphasis{\sphinxupquote{up\_y}}\sphinxstyleliteralemphasis{\sphinxupquote{, }}\sphinxstyleliteralemphasis{\sphinxupquote{up\_z}}\sphinxstyleliteralemphasis{\sphinxupquote{)}}) \textendash{} vector pointing to up direction

\item {} 
\sphinxAtStartPar
\sphinxstyleliteralstrong{\sphinxupquote{zoom}} (\sphinxstyleliteralemphasis{\sphinxupquote{float}}) \textendash{} zoom factor

\item {} 
\sphinxAtStartPar
\sphinxstyleliteralstrong{\sphinxupquote{background\_color}} (\sphinxstyleliteralemphasis{\sphinxupquote{string}}) \textendash{} background color of scene

\end{itemize}

\sphinxlineitem{Returns}
\sphinxAtStartPar
Image instance used to display image in notebook

\sphinxlineitem{Return type}
\sphinxAtStartPar
\sphinxcode{\sphinxupquote{IPython.core.display.Image}}

\end{description}\end{quote}

\end{fulllineitems}

\index{video() (lammps.IPyLammps method)@\spxentry{video()}\spxextra{lammps.IPyLammps method}}

\begin{fulllineitems}
\phantomsection\label{\detokenize{Python_module:lammps.IPyLammps.video}}
\pysigstartsignatures
\pysiglinewithargsret{\sphinxbfcode{\sphinxupquote{video}}}{\sphinxparam{\DUrole{n}{filename}}}{}
\pysigstopsignatures
\sphinxAtStartPar
Load video from file

\sphinxAtStartPar
Can be used to visualize videos from {\hyperref[\detokenize{dump_image::doc}]{\sphinxcrossref{\DUrole{doc}{dump movie}}}}.
\begin{quote}\begin{description}
\sphinxlineitem{Parameters}
\sphinxAtStartPar
\sphinxstyleliteralstrong{\sphinxupquote{filename}} (\sphinxstyleliteralemphasis{\sphinxupquote{string}}) \textendash{} Path to video file

\sphinxlineitem{Returns}
\sphinxAtStartPar
HTML Video Tag used by notebook to embed a video

\sphinxlineitem{Return type}
\sphinxAtStartPar
\sphinxcode{\sphinxupquote{IPython.display.HTML}}

\end{description}\end{quote}

\end{fulllineitems}


\end{fulllineitems}



\bigskip\hrule\bigskip



\subsection{Additional components of the \sphinxstyleliteralintitle{\sphinxupquote{lammps}} module}
\label{\detokenize{Python_module:additional-components-of-the-lammps-module}}
\sphinxAtStartPar
The {\hyperref[\detokenize{Python_module:module-lammps}]{\sphinxcrossref{\sphinxcode{\sphinxupquote{lammps}}}}} module additionally contains several constants
and the {\hyperref[\detokenize{Python_module:lammps.NeighList}]{\sphinxcrossref{\sphinxcode{\sphinxupquote{NeighList}}}}} class:


\subsubsection{Data Types}
\label{\detokenize{Python_module:data-types}}\label{\detokenize{Python_module:py-datatype-constants}}

\begin{fulllineitems}

\pysigstartsignatures
\pysigline{\sphinxbfcode{\sphinxupquote{LAMMPS\_INT,~LAMMPS\_INT\_2D,~LAMMPS\_DOUBLE,~LAMMPS\_DOUBLE\_2D,~LAMMPS\_INT64,~LAMMPS\_INT64\_2D,~LAMMPS\_STRING}}}
\pysigstopsignatures
\sphinxAtStartPar
Constants in the {\hyperref[\detokenize{Python_module:module-lammps}]{\sphinxcrossref{\sphinxcode{\sphinxupquote{lammps}}}}} module to indicate how to
cast data when the C library function returns a void pointer.
Used in {\hyperref[\detokenize{Python_module:lammps.lammps.extract_global}]{\sphinxcrossref{\sphinxcode{\sphinxupquote{lammps.extract\_global()}}}}} and {\hyperref[\detokenize{Python_module:lammps.lammps.extract_atom}]{\sphinxcrossref{\sphinxcode{\sphinxupquote{lammps.extract\_atom()}}}}}.
See \sphinxcode{\sphinxupquote{\_LMP\_DATATYPE\_CONST}} for the equivalent constants in the
C library interface.

\end{fulllineitems}



\subsubsection{Style Constants}
\label{\detokenize{Python_module:style-constants}}\label{\detokenize{Python_module:py-style-constants}}

\begin{fulllineitems}

\pysigstartsignatures
\pysigline{\sphinxbfcode{\sphinxupquote{LMP\_STYLE\_GLOBAL,~LMP\_STYLE\_ATOM,~LMP\_STYLE\_LOCAL}}}
\pysigstopsignatures
\sphinxAtStartPar
Constants in the {\hyperref[\detokenize{Python_module:module-lammps}]{\sphinxcrossref{\sphinxcode{\sphinxupquote{lammps}}}}} module to select what style of data
to request from computes or fixes. See {\hyperref[\detokenize{Library_objects:_CPPv416_LMP_STYLE_CONST}]{\sphinxcrossref{\sphinxcode{\sphinxupquote{\_LMP\_STYLE\_CONST}}}}}
for the equivalent constants in the C library interface. Used in
{\hyperref[\detokenize{Python_module:lammps.lammps.extract_compute}]{\sphinxcrossref{\sphinxcode{\sphinxupquote{lammps.extract\_compute()}}}}}, {\hyperref[\detokenize{Python_module:lammps.lammps.extract_fix}]{\sphinxcrossref{\sphinxcode{\sphinxupquote{lammps.extract\_fix()}}}}}, and their NumPy variants
{\hyperref[\detokenize{Python_module:lammps.numpy_wrapper.numpy_wrapper.extract_compute}]{\sphinxcrossref{\sphinxcode{\sphinxupquote{lammps.numpy.extract\_compute()}}}}} and
{\hyperref[\detokenize{Python_module:lammps.numpy_wrapper.numpy_wrapper.extract_fix}]{\sphinxcrossref{\sphinxcode{\sphinxupquote{lammps.numpy.extract\_fix()}}}}}.

\end{fulllineitems}



\subsubsection{Type Constants}
\label{\detokenize{Python_module:type-constants}}\label{\detokenize{Python_module:py-type-constants}}

\begin{fulllineitems}

\pysigstartsignatures
\pysigline{\sphinxbfcode{\sphinxupquote{LMP\_TYPE\_SCALAR,~LMP\_TYPE\_VECTOR,~LMP\_TYPE\_ARRAY,~LMP\_SIZE\_VECTOR,~LMP\_SIZE\_ROWS,~LMP\_SIZE\_COLS}}}
\pysigstopsignatures
\sphinxAtStartPar
Constants in the {\hyperref[\detokenize{Python_module:module-lammps}]{\sphinxcrossref{\sphinxcode{\sphinxupquote{lammps}}}}} module to select what type of data
to request  from computes  or fixes.  See {\hyperref[\detokenize{Library_objects:_CPPv415_LMP_TYPE_CONST}]{\sphinxcrossref{\sphinxcode{\sphinxupquote{\_LMP\_TYPE\_CONST}}}}}
for the equivalent constants in the C library interface. Used in
{\hyperref[\detokenize{Python_module:lammps.lammps.extract_compute}]{\sphinxcrossref{\sphinxcode{\sphinxupquote{lammps.extract\_compute()}}}}}, {\hyperref[\detokenize{Python_module:lammps.lammps.extract_fix}]{\sphinxcrossref{\sphinxcode{\sphinxupquote{lammps.extract\_fix()}}}}}, and their NumPy variants
{\hyperref[\detokenize{Python_module:lammps.numpy_wrapper.numpy_wrapper.extract_compute}]{\sphinxcrossref{\sphinxcode{\sphinxupquote{lammps.numpy.extract\_compute()}}}}} and
{\hyperref[\detokenize{Python_module:lammps.numpy_wrapper.numpy_wrapper.extract_fix}]{\sphinxcrossref{\sphinxcode{\sphinxupquote{lammps.numpy.extract\_fix()}}}}}.

\end{fulllineitems}



\subsubsection{Variable Type Constants}
\label{\detokenize{Python_module:variable-type-constants}}\label{\detokenize{Python_module:py-vartype-constants}}

\begin{fulllineitems}

\pysigstartsignatures
\pysigline{\sphinxbfcode{\sphinxupquote{LMP\_VAR\_EQUAL,~LMP\_VAR\_ATOM}}}
\pysigstopsignatures
\sphinxAtStartPar
Constants in the {\hyperref[\detokenize{Python_module:module-lammps}]{\sphinxcrossref{\sphinxcode{\sphinxupquote{lammps}}}}} module to select what type of
variable to query when calling {\hyperref[\detokenize{Python_module:lammps.lammps.extract_variable}]{\sphinxcrossref{\sphinxcode{\sphinxupquote{lammps.extract\_variable()}}}}}. See also: {\hyperref[\detokenize{variable::doc}]{\sphinxcrossref{\DUrole{doc}{variable command}}}}.

\end{fulllineitems}



\subsubsection{Classes representing internal objects}
\label{\detokenize{Python_module:classes-representing-internal-objects}}\index{NeighList (class in lammps)@\spxentry{NeighList}\spxextra{class in lammps}}

\begin{fulllineitems}
\phantomsection\label{\detokenize{Python_module:lammps.NeighList}}
\pysigstartsignatures
\pysiglinewithargsret{\sphinxbfcode{\sphinxupquote{class\DUrole{w}{ }}}\sphinxcode{\sphinxupquote{lammps.}}\sphinxbfcode{\sphinxupquote{NeighList}}}{\sphinxparam{\DUrole{n}{lmp}}\sphinxparamcomma \sphinxparam{\DUrole{n}{idx}}}{}
\pysigstopsignatures
\sphinxAtStartPar
This is a wrapper class that exposes the contents of a neighbor list.

\sphinxAtStartPar
It can be used like a regular Python list. Each element is a tuple of:
\begin{itemize}
\item {} 
\sphinxAtStartPar
the atom local index

\item {} 
\sphinxAtStartPar
its number of neighbors

\item {} 
\sphinxAtStartPar
and a pointer to an c\_int array containing local atom indices of its
neighbors

\end{itemize}

\sphinxAtStartPar
Internally it uses the lower\sphinxhyphen{}level LAMMPS C\sphinxhyphen{}library interface.
\begin{quote}\begin{description}
\sphinxlineitem{Parameters}\begin{itemize}
\item {} 
\sphinxAtStartPar
\sphinxstyleliteralstrong{\sphinxupquote{lmp}} ({\hyperref[\detokenize{Python_module:lammps.lammps}]{\sphinxcrossref{\sphinxstyleliteralemphasis{\sphinxupquote{lammps}}}}}) \textendash{} reference to instance of {\hyperref[\detokenize{Python_module:module-lammps}]{\sphinxcrossref{\sphinxcode{\sphinxupquote{lammps}}}}}

\item {} 
\sphinxAtStartPar
\sphinxstyleliteralstrong{\sphinxupquote{idx}} (\sphinxstyleliteralemphasis{\sphinxupquote{int}}) \textendash{} neighbor list index

\end{itemize}

\end{description}\end{quote}
\index{size (lammps.NeighList property)@\spxentry{size}\spxextra{lammps.NeighList property}}

\begin{fulllineitems}
\phantomsection\label{\detokenize{Python_module:lammps.NeighList.size}}
\pysigstartsignatures
\pysigline{\sphinxbfcode{\sphinxupquote{property\DUrole{w}{ }}}\sphinxbfcode{\sphinxupquote{size}}}
\pysigstopsignatures\begin{quote}\begin{description}
\sphinxlineitem{Returns}
\sphinxAtStartPar
number of elements in neighbor list

\end{description}\end{quote}

\end{fulllineitems}

\index{get() (lammps.NeighList method)@\spxentry{get()}\spxextra{lammps.NeighList method}}

\begin{fulllineitems}
\phantomsection\label{\detokenize{Python_module:lammps.NeighList.get}}
\pysigstartsignatures
\pysiglinewithargsret{\sphinxbfcode{\sphinxupquote{get}}}{\sphinxparam{\DUrole{n}{element}}}{}
\pysigstopsignatures
\sphinxAtStartPar
Access a specific neighbor list entry. “element” must be a number from 0 to the size\sphinxhyphen{}1 of the list
\begin{quote}\begin{description}
\sphinxlineitem{Returns}
\sphinxAtStartPar
tuple with atom local index, number of neighbors and ctypes pointer to neighbor’s local atom indices

\sphinxlineitem{Return type}
\sphinxAtStartPar
(int, int, ctypes.POINTER(c\_int))

\end{description}\end{quote}

\end{fulllineitems}

\index{find() (lammps.NeighList method)@\spxentry{find()}\spxextra{lammps.NeighList method}}

\begin{fulllineitems}
\phantomsection\label{\detokenize{Python_module:lammps.NeighList.find}}
\pysigstartsignatures
\pysiglinewithargsret{\sphinxbfcode{\sphinxupquote{find}}}{\sphinxparam{\DUrole{n}{iatom}}}{}
\pysigstopsignatures
\sphinxAtStartPar
Find the neighbor list for a specific (local) atom iatom.
If there is no list for iatom, (\sphinxhyphen{}1, None) is returned.
\begin{quote}\begin{description}
\sphinxlineitem{Returns}
\sphinxAtStartPar
tuple with number of neighbors and ctypes pointer to neighbor’s local atom indices

\sphinxlineitem{Return type}
\sphinxAtStartPar
(int, ctypes.POINTER(c\_int))

\end{description}\end{quote}

\end{fulllineitems}


\end{fulllineitems}

\index{NumPyNeighList (class in lammps.numpy\_wrapper)@\spxentry{NumPyNeighList}\spxextra{class in lammps.numpy\_wrapper}}

\begin{fulllineitems}
\phantomsection\label{\detokenize{Python_module:lammps.numpy_wrapper.NumPyNeighList}}
\pysigstartsignatures
\pysiglinewithargsret{\sphinxbfcode{\sphinxupquote{class\DUrole{w}{ }}}\sphinxcode{\sphinxupquote{lammps.numpy\_wrapper.}}\sphinxbfcode{\sphinxupquote{NumPyNeighList}}}{\sphinxparam{\DUrole{n}{lmp}}\sphinxparamcomma \sphinxparam{\DUrole{n}{idx}}}{}
\pysigstopsignatures
\sphinxAtStartPar
This is a wrapper class that exposes the contents of a neighbor list.

\sphinxAtStartPar
It can be used like a regular Python list. Each element is a tuple of:
\begin{itemize}
\item {} 
\sphinxAtStartPar
the atom local index

\item {} 
\sphinxAtStartPar
a NumPy array containing the local atom indices of its neighbors

\end{itemize}

\sphinxAtStartPar
Internally it uses the lower\sphinxhyphen{}level LAMMPS C\sphinxhyphen{}library interface.
\begin{quote}\begin{description}
\sphinxlineitem{Parameters}\begin{itemize}
\item {} 
\sphinxAtStartPar
\sphinxstyleliteralstrong{\sphinxupquote{lmp}} ({\hyperref[\detokenize{Python_module:lammps.lammps}]{\sphinxcrossref{\sphinxstyleliteralemphasis{\sphinxupquote{lammps}}}}}) \textendash{} reference to instance of {\hyperref[\detokenize{Python_module:module-lammps}]{\sphinxcrossref{\sphinxcode{\sphinxupquote{lammps}}}}}

\item {} 
\sphinxAtStartPar
\sphinxstyleliteralstrong{\sphinxupquote{idx}} (\sphinxstyleliteralemphasis{\sphinxupquote{int}}) \textendash{} neighbor list index

\end{itemize}

\end{description}\end{quote}
\index{get() (lammps.numpy\_wrapper.NumPyNeighList method)@\spxentry{get()}\spxextra{lammps.numpy\_wrapper.NumPyNeighList method}}

\begin{fulllineitems}
\phantomsection\label{\detokenize{Python_module:lammps.numpy_wrapper.NumPyNeighList.get}}
\pysigstartsignatures
\pysiglinewithargsret{\sphinxbfcode{\sphinxupquote{get}}}{\sphinxparam{\DUrole{n}{element}}}{}
\pysigstopsignatures
\sphinxAtStartPar
Access a specific neighbor list entry. “element” must be a number from 0 to the size\sphinxhyphen{}1 of the list
\begin{quote}\begin{description}
\sphinxlineitem{Returns}
\sphinxAtStartPar
tuple with atom local index, numpy array of neighbor local atom indices

\sphinxlineitem{Return type}
\sphinxAtStartPar
(int, numpy.array)

\end{description}\end{quote}

\end{fulllineitems}

\index{find() (lammps.numpy\_wrapper.NumPyNeighList method)@\spxentry{find()}\spxextra{lammps.numpy\_wrapper.NumPyNeighList method}}

\begin{fulllineitems}
\phantomsection\label{\detokenize{Python_module:lammps.numpy_wrapper.NumPyNeighList.find}}
\pysigstartsignatures
\pysiglinewithargsret{\sphinxbfcode{\sphinxupquote{find}}}{\sphinxparam{\DUrole{n}{iatom}}}{}
\pysigstopsignatures
\sphinxAtStartPar
Find the neighbor list for a specific (local) atom iatom.
If there is no list for iatom, None is returned.
\begin{quote}\begin{description}
\sphinxlineitem{Returns}
\sphinxAtStartPar
numpy array of neighbor local atom indices

\sphinxlineitem{Return type}
\sphinxAtStartPar
numpy.array or None

\end{description}\end{quote}

\end{fulllineitems}


\end{fulllineitems}


\sphinxstepscope


\section{Extending the Python interface}
\label{\detokenize{Python_ext:extending-the-python-interface}}\label{\detokenize{Python_ext::doc}}
\sphinxAtStartPar
As noted previously, most of the {\hyperref[\detokenize{Python_module:lammps.lammps}]{\sphinxcrossref{\sphinxcode{\sphinxupquote{lammps}}}}}
Python class methods correspond one\sphinxhyphen{}to\sphinxhyphen{}one with the functions in the
LAMMPS library interface in \sphinxcode{\sphinxupquote{src/library.cpp}} and \sphinxcode{\sphinxupquote{library.h}}.
This means you can extend the Python wrapper by following these steps:
\begin{itemize}
\item {} 
\sphinxAtStartPar
Add a new interface function to \sphinxcode{\sphinxupquote{src/library.cpp}} and
\sphinxcode{\sphinxupquote{src/library.h}}.

\item {} 
\sphinxAtStartPar
Rebuild LAMMPS as a shared library.

\item {} 
\sphinxAtStartPar
Add a wrapper method to \sphinxcode{\sphinxupquote{python/lammps/core.py}} for this interface
function.

\item {} 
\sphinxAtStartPar
Define the corresponding \sphinxcode{\sphinxupquote{argtypes}} list and \sphinxcode{\sphinxupquote{restype}}
in the \sphinxcode{\sphinxupquote{lammps.\_\_init\_\_()}} function.

\item {} 
\sphinxAtStartPar
Re\sphinxhyphen{}install the shared library and the python module, if needed

\item {} 
\sphinxAtStartPar
You should now be able to invoke the new interface function from a
Python script.

\end{itemize}

\sphinxstepscope


\section{Calling Python from LAMMPS}
\label{\detokenize{Python_call:calling-python-from-lammps}}\label{\detokenize{Python_call::doc}}
\sphinxAtStartPar
LAMMPS has several commands which can be used to invoke Python
code directly from an input script:
\begin{itemize}
\item {} 
\sphinxAtStartPar
{\hyperref[\detokenize{python::doc}]{\sphinxcrossref{\DUrole{doc}{python}}}}

\item {} 
\sphinxAtStartPar
{\hyperref[\detokenize{variable::doc}]{\sphinxcrossref{\DUrole{doc}{variable python}}}}

\item {} 
\sphinxAtStartPar
{\hyperref[\detokenize{fix_python_invoke::doc}]{\sphinxcrossref{\DUrole{doc}{fix python/invoke}}}}

\item {} 
\sphinxAtStartPar
{\hyperref[\detokenize{pair_python::doc}]{\sphinxcrossref{\DUrole{doc}{pair\_style python}}}}

\end{itemize}

\sphinxAtStartPar
The {\hyperref[\detokenize{python::doc}]{\sphinxcrossref{\DUrole{doc}{python}}}} command which can be used to define and
execute a Python function that you write the code for.  The Python
function can also be assigned to a LAMMPS python\sphinxhyphen{}style variable via
the {\hyperref[\detokenize{variable::doc}]{\sphinxcrossref{\DUrole{doc}{variable}}}} command.  Each time the variable is
evaluated, either in the LAMMPS input script itself, or by another
LAMMPS command that uses the variable, this will trigger the Python
function to be invoked.

\sphinxAtStartPar
The Python code for the function can be included directly in the input
script or in an auxiliary file.  The function can have arguments which
are mapped to LAMMPS variables (also defined in the input script) and
it can return a value to a LAMMPS variable.  This is thus a mechanism
for your input script to pass information to a piece of Python code,
ask Python to execute the code, and return information to your input
script.

\sphinxAtStartPar
Note that a Python function can be arbitrarily complex.  It can import
other Python modules, instantiate Python classes, call other Python
functions, etc.  The Python code that you provide can contain more
code than the single function.  It can contain other functions or
Python classes, as well as global variables or other mechanisms for
storing state between calls from LAMMPS to the function.

\sphinxAtStartPar
The Python function you provide can consist of “pure” Python code that
only performs operations provided by standard Python.  However, the
Python function can also “call back” to LAMMPS through its
Python\sphinxhyphen{}wrapped library interface, in the manner described in the
{\hyperref[\detokenize{Python_run::doc}]{\sphinxcrossref{\DUrole{doc}{Python run}}}} doc page.  This means it can issue LAMMPS
input script commands or query and set internal LAMMPS state.  As an
example, this can be useful in an input script to create a more
complex loop with branching logic, than can be created using the
simple looping and branching logic enabled by the {\hyperref[\detokenize{next::doc}]{\sphinxcrossref{\DUrole{doc}{next}}}} and
{\hyperref[\detokenize{if::doc}]{\sphinxcrossref{\DUrole{doc}{if}}}} commands.

\sphinxAtStartPar
See the {\hyperref[\detokenize{python::doc}]{\sphinxcrossref{\DUrole{doc}{python}}}} page and the {\hyperref[\detokenize{variable::doc}]{\sphinxcrossref{\DUrole{doc}{variable}}}}
doc page for its python\sphinxhyphen{}style variables for more info, including
examples of Python code you can write for both pure Python operations
and callbacks to LAMMPS.

\sphinxstepscope


\section{Output Readers}
\label{\detokenize{Python_formats:module-lammps.formats}}\label{\detokenize{Python_formats:output-readers}}\label{\detokenize{Python_formats::doc}}\index{module@\spxentry{module}!lammps.formats@\spxentry{lammps.formats}}\index{lammps.formats@\spxentry{lammps.formats}!module@\spxentry{module}}
\sphinxAtStartPar
The Python package contains the {\hyperref[\detokenize{Python_formats:module-lammps.formats}]{\sphinxcrossref{\sphinxcode{\sphinxupquote{lammps.formats}}}}} module, which
provides classes to post\sphinxhyphen{}process some of the output files generated by LAMMPS.


\begin{fulllineitems}

\pysigstartsignatures
\pysiglinewithargsret{\sphinxbfcode{\sphinxupquote{class\DUrole{w}{ }}}\sphinxcode{\sphinxupquote{lammps.formats.}}\sphinxbfcode{\sphinxupquote{LogFile}}}{\sphinxparam{\DUrole{n}{filename}}}{}
\pysigstopsignatures
\sphinxAtStartPar
Reads LAMMPS log files and extracts the thermo information

\sphinxAtStartPar
It supports the line, multi, and yaml thermo output styles.
\begin{quote}\begin{description}
\sphinxlineitem{Parameters}
\sphinxAtStartPar
\sphinxstyleliteralstrong{\sphinxupquote{filename}} (\sphinxstyleliteralemphasis{\sphinxupquote{str}}) \textendash{} path to log file

\sphinxlineitem{Variables}\begin{itemize}
\item {} 
\sphinxAtStartPar
\sphinxstyleliteralstrong{\sphinxupquote{runs}} \textendash{} List of LAMMPS runs in log file. Each run is a dictionary with
thermo fields as keys, storing the values over time

\item {} 
\sphinxAtStartPar
\sphinxstyleliteralstrong{\sphinxupquote{errors}} \textendash{} List of error lines in log file

\end{itemize}

\end{description}\end{quote}

\end{fulllineitems}



\begin{fulllineitems}

\pysigstartsignatures
\pysiglinewithargsret{\sphinxbfcode{\sphinxupquote{class\DUrole{w}{ }}}\sphinxcode{\sphinxupquote{lammps.formats.}}\sphinxbfcode{\sphinxupquote{AvgChunkFile}}}{\sphinxparam{\DUrole{n}{filename}}}{}
\pysigstopsignatures
\sphinxAtStartPar
Reads files generated by fix ave/chunk
\begin{quote}\begin{description}
\sphinxlineitem{Parameters}
\sphinxAtStartPar
\sphinxstyleliteralstrong{\sphinxupquote{filename}} (\sphinxstyleliteralemphasis{\sphinxupquote{str}}) \textendash{} path to ave/chunk file

\sphinxlineitem{Variables}\begin{itemize}
\item {} 
\sphinxAtStartPar
\sphinxstyleliteralstrong{\sphinxupquote{timesteps}} \textendash{} List of timesteps stored in file

\item {} 
\sphinxAtStartPar
\sphinxstyleliteralstrong{\sphinxupquote{total\_count}} \textendash{} total count over time

\item {} 
\sphinxAtStartPar
\sphinxstyleliteralstrong{\sphinxupquote{chunks}} \textendash{} List of chunks. Each chunk is a dictionary containing its ID, the coordinates, and the averaged quantities

\end{itemize}

\end{description}\end{quote}

\end{fulllineitems}


\sphinxstepscope


\section{Example Python scripts}
\label{\detokenize{Python_examples:example-python-scripts}}\label{\detokenize{Python_examples::doc}}
\sphinxAtStartPar
The \sphinxcode{\sphinxupquote{python/examples}} directory has Python scripts which show how Python
can run LAMMPS, grab data, change it, and put it back into LAMMPS.

\sphinxAtStartPar
These are the Python scripts included as demos in the \sphinxcode{\sphinxupquote{python/examples}}
directory of the LAMMPS distribution, to illustrate the kinds of
things that are possible when Python wraps LAMMPS.  If you create your
own scripts, send them to us and we can include them in the LAMMPS
distribution.


\begin{savenotes}\sphinxattablestart
\sphinxthistablewithglobalstyle
\centering
\begin{tabulary}{\linewidth}[t]{TT}
\sphinxtoprule
\sphinxtableatstartofbodyhook
\sphinxAtStartPar
\sphinxcode{\sphinxupquote{trivial.py}}
&
\sphinxAtStartPar
read/run a LAMMPS input script through Python
\\
\sphinxhline
\sphinxAtStartPar
\sphinxcode{\sphinxupquote{demo.py}}
&
\sphinxAtStartPar
invoke various LAMMPS library interface routines
\\
\sphinxhline
\sphinxAtStartPar
\sphinxcode{\sphinxupquote{simple.py}}
&
\sphinxAtStartPar
run in parallel, similar to \sphinxcode{\sphinxupquote{examples/COUPLE/simple/simple.cpp}}
\\
\sphinxhline
\sphinxAtStartPar
\sphinxcode{\sphinxupquote{split.py}}
&
\sphinxAtStartPar
same as \sphinxcode{\sphinxupquote{simple.py}} but running in parallel on a subset of procs
\\
\sphinxhline
\sphinxAtStartPar
\sphinxcode{\sphinxupquote{gui.py}}
&
\sphinxAtStartPar
GUI go/stop/temperature\sphinxhyphen{}slider to control LAMMPS
\\
\sphinxhline
\sphinxAtStartPar
\sphinxcode{\sphinxupquote{plot.py}}
&
\sphinxAtStartPar
real\sphinxhyphen{}time temperature plot with GnuPlot via Pizza.py
\\
\sphinxhline
\sphinxAtStartPar
\sphinxcode{\sphinxupquote{viz\_TOOL.py}}
&
\sphinxAtStartPar
real\sphinxhyphen{}time viz via some viz package
\\
\sphinxhline
\sphinxAtStartPar
\sphinxcode{\sphinxupquote{vizplotgui\_TOOL.py}}
&
\sphinxAtStartPar
combination of \sphinxcode{\sphinxupquote{viz\_TOOL.py}} and \sphinxcode{\sphinxupquote{plot.py}} and \sphinxcode{\sphinxupquote{gui.py}}
\\
\sphinxbottomrule
\end{tabulary}
\sphinxtableafterendhook\par
\sphinxattableend\end{savenotes}

\sphinxAtStartPar
For the \sphinxcode{\sphinxupquote{viz\_TOOL.py}} and \sphinxcode{\sphinxupquote{vizplotgui\_TOOL.py}} commands, replace \sphinxcode{\sphinxupquote{TOOL}}
with \sphinxcode{\sphinxupquote{gl}} or \sphinxcode{\sphinxupquote{atomeye}} or \sphinxcode{\sphinxupquote{pymol}} or \sphinxcode{\sphinxupquote{vmd}}, depending on what
visualization package you have installed.

\sphinxAtStartPar
Note that for GL, you need to be able to run the Pizza.py GL tool,
which is included in the pizza subdirectory.  See the Pizza.py doc pages for more info:
\begin{itemize}
\item {} 
\sphinxAtStartPar
\sphinxurl{https://lammps.github.io/pizza}

\end{itemize}

\sphinxAtStartPar
Note that for AtomEye, you need version 3, and there is a line in the
scripts that specifies the path and name of the executable.  See the
AtomEye web pages for more details:
\begin{itemize}
\item {} 
\sphinxAtStartPar
\sphinxurl{http://li.mit.edu/Archive/Graphics/A/}

\item {} 
\sphinxAtStartPar
\sphinxurl{http://li.mit.edu/Archive/Graphics/A3/A3.html}

\end{itemize}

\sphinxAtStartPar
The latter link is to AtomEye 3 which has the scripting capability
needed by these Python scripts.

\sphinxAtStartPar
Note that for PyMol, you need to have built and installed the
open\sphinxhyphen{}source version of PyMol in your Python, so that you can import it
from a Python script.  See the PyMol web pages for more details:
\begin{itemize}
\item {} 
\sphinxAtStartPar
\sphinxurl{https://www.pymol.org}

\item {} 
\sphinxAtStartPar
\sphinxurl{https://github.com/schrodinger/pymol-open-source}

\end{itemize}

\sphinxAtStartPar
The latter link is to the open\sphinxhyphen{}source version.

\sphinxAtStartPar
Note that for VMD, you need a fairly current version (1.8.7 works for
me) and there are some lines in the \sphinxcode{\sphinxupquote{pizza/vmd.py}} script for 4 PIZZA
variables that have to match the VMD installation on your system.


\bigskip\hrule\bigskip


\sphinxAtStartPar
See the \sphinxcode{\sphinxupquote{python/README}} file for instructions on how to run them and the
source code for individual scripts for comments about what they do.

\sphinxAtStartPar
Here are screenshots of the \sphinxcode{\sphinxupquote{vizplotgui\_tool.py}} script in action for
different visualization package options:

\sphinxAtStartPar
\sphinxincludegraphics[width=0.240\linewidth]{{screenshot_gl}.jpg}  \sphinxincludegraphics[width=0.240\linewidth]{{screenshot_atomeye}.jpg}  \sphinxincludegraphics[width=0.240\linewidth]{{screenshot_pymol}.jpg}  \sphinxincludegraphics[width=0.240\linewidth]{{screenshot_vmd}.jpg}



\sphinxstepscope


\section{Handling LAMMPS errors}
\label{\detokenize{Python_error:handling-lammps-errors}}\label{\detokenize{Python_error::doc}}
\sphinxAtStartPar
LAMMPS and the LAMMPS library are compiled with {\hyperref[\detokenize{Build_settings:exceptions}]{\sphinxcrossref{\DUrole{std,std-ref}{C++ exception support}}}} to provide a better error handling experience.  LAMMPS errors
trigger throwing a C++ exception. These exceptions allow capturing errors on
the C++ side and rethrowing them on the Python side.  This way LAMMPS errors
can be handled through the Python exception handling mechanism.

\begin{sphinxVerbatim}[commandchars=\\\{\}]
\PYG{k+kn}{from}\PYG{+w}{ }\PYG{n+nn}{lammps}\PYG{+w}{ }\PYG{k+kn}{import} \PYG{n}{lammps}\PYG{p}{,} \PYG{n}{MPIAbortException}

\PYG{n}{lmp} \PYG{o}{=} \PYG{n}{lammps}\PYG{p}{(}\PYG{p}{)}

\PYG{k}{try}\PYG{p}{:}
   \PYG{c+c1}{\PYGZsh{} LAMMPS will normally terminate itself and the running process if an error}
   \PYG{c+c1}{\PYGZsh{} occurs. This would kill the Python interpreter.  The library wrapper will}
   \PYG{c+c1}{\PYGZsh{} detect that an error has occurred and throw a Python exception}

   \PYG{n}{lmp}\PYG{o}{.}\PYG{n}{command}\PYG{p}{(}\PYG{l+s+s1}{\PYGZsq{}}\PYG{l+s+s1}{unknown}\PYG{l+s+s1}{\PYGZsq{}}\PYG{p}{)}
\PYG{k}{except} \PYG{n}{MPIAbortException} \PYG{k}{as} \PYG{n}{ae}\PYG{p}{:}
   \PYG{c+c1}{\PYGZsh{} Single MPI process got killed. This would normally be handled by an MPI abort}
   \PYG{k}{pass}
\PYG{k}{except} \PYG{n+ne}{Exception} \PYG{k}{as} \PYG{n}{e}\PYG{p}{:}
   \PYG{c+c1}{\PYGZsh{} All (MPI) processes have reached this error}
   \PYG{k}{pass}
\end{sphinxVerbatim}

\begin{sphinxadmonition}{warning}{Warning:}
\sphinxAtStartPar
Capturing a LAMMPS exception in Python can still mean that the
current LAMMPS process is in an illegal state and must be
terminated. It is advised to save your data and terminate the Python
instance as quickly as possible when running in parallel with MPI.
\end{sphinxadmonition}

\sphinxstepscope


\section{Troubleshooting}
\label{\detokenize{Python_trouble:troubleshooting}}\label{\detokenize{Python_trouble::doc}}

\subsection{Testing if Python can launch LAMMPS}
\label{\detokenize{Python_trouble:testing-if-python-can-launch-lammps}}
\sphinxAtStartPar
To test if LAMMPS is callable from Python, launch Python interactively
and type:

\begin{sphinxVerbatim}[commandchars=\\\{\}]
\PYG{g+gp}{\PYGZgt{}\PYGZgt{}\PYGZgt{} }\PYG{k+kn}{from}\PYG{+w}{ }\PYG{n+nn}{lammps}\PYG{+w}{ }\PYG{k+kn}{import} \PYG{n}{lammps}
\PYG{g+gp}{\PYGZgt{}\PYGZgt{}\PYGZgt{} }\PYG{n}{lmp} \PYG{o}{=} \PYG{n}{lammps}\PYG{p}{(}\PYG{p}{)}
\end{sphinxVerbatim}

\sphinxAtStartPar
If you get no errors, you’re ready to use LAMMPS from Python.  If the
second command fails, the most common error to see is

\begin{sphinxVerbatim}[commandchars=\\\{\}]
OSError:\PYG{+w}{ }Could\PYG{+w}{ }not\PYG{+w}{ }load\PYG{+w}{ }LAMMPS\PYG{+w}{ }dynamic\PYG{+w}{ }library
\end{sphinxVerbatim}

\sphinxAtStartPar
which means Python was unable to load the LAMMPS shared library.  This
typically occurs if the system can’t find the LAMMPS shared library or
one of the auxiliary shared libraries it depends on, or if something
about the library is incompatible with your Python.  The error message
should give you an indication of what went wrong.

\sphinxAtStartPar
If your shared library uses a suffix, such as \sphinxcode{\sphinxupquote{liblammps\_mpi.so}}, change
the constructor call as follows (see {\hyperref[\detokenize{Python_create:python-create-lammps}]{\sphinxcrossref{\DUrole{std,std-ref}{Creating or deleting a LAMMPS object}}}} for more details):

\begin{sphinxVerbatim}[commandchars=\\\{\}]
\PYG{g+gp}{\PYGZgt{}\PYGZgt{}\PYGZgt{} }\PYG{n}{lmp} \PYG{o}{=} \PYG{n}{lammps}\PYG{p}{(}\PYG{n}{name}\PYG{o}{=}\PYG{l+s+s1}{\PYGZsq{}}\PYG{l+s+s1}{mpi}\PYG{l+s+s1}{\PYGZsq{}}\PYG{p}{)}
\end{sphinxVerbatim}

\sphinxAtStartPar
You can also test the load directly in Python as follows, without
first importing from the \sphinxcode{\sphinxupquote{lammps}} module:

\begin{sphinxVerbatim}[commandchars=\\\{\}]
\PYG{g+gp}{\PYGZgt{}\PYGZgt{}\PYGZgt{} }\PYG{k+kn}{from}\PYG{+w}{ }\PYG{n+nn}{ctypes}\PYG{+w}{ }\PYG{k+kn}{import} \PYG{n}{CDLL}
\PYG{g+gp}{\PYGZgt{}\PYGZgt{}\PYGZgt{} }\PYG{n}{CDLL}\PYG{p}{(}\PYG{l+s+s2}{\PYGZdq{}}\PYG{l+s+s2}{liblammps.so}\PYG{l+s+s2}{\PYGZdq{}}\PYG{p}{)}
\end{sphinxVerbatim}

\sphinxAtStartPar
If an error occurs, carefully go through the steps in {\hyperref[\detokenize{Python_install:python-install-guides}]{\sphinxcrossref{\DUrole{std,std-ref}{Installing the LAMMPS Python Module and Shared Library}}}} and on the
{\hyperref[\detokenize{Build_basics::doc}]{\sphinxcrossref{\DUrole{doc}{Build\_basics}}}} page about building a shared library.

\sphinxAtStartPar
If you are not familiar with \sphinxhref{https://www.python.org}{Python}, it is a
powerful scripting and programming language which can do almost
everything that compiled languages like C, C++, or Fortran can do in
fewer lines of code. It also comes with a large collection of add\sphinxhyphen{}on
modules for many purposes (either bundled or easily installed from
Python code repositories).  The major drawback is slower execution speed
of the script code compared to compiled programming languages.  But when
the script code is interfaced to optimized compiled code, performance
can be on par with a standalone executable, for as long as the scripting
is restricted to high\sphinxhyphen{}level operations.  Thus Python is also convenient
to use as a “glue” language to “drive” a program through its library
interface, or to hook multiple pieces of software together, such as a
simulation code and a visualization tool, or to run a coupled
multi\sphinxhyphen{}scale or multi\sphinxhyphen{}physics model.

\sphinxAtStartPar
See the {\hyperref[\detokenize{Howto_couple::doc}]{\sphinxcrossref{\DUrole{doc}{Coupling LAMMPS to other codes}}}} page for more ideas about coupling LAMMPS
to other codes.  See the {\hyperref[\detokenize{Library::doc}]{\sphinxcrossref{\DUrole{doc}{LAMMPS Library Interfaces}}}} page for a description of the
LAMMPS library interfaces.  That interface is exposed to Python either
when calling LAMMPS from Python or when calling Python from a LAMMPS
input script and then calling back to LAMMPS from Python code.  The
C\sphinxhyphen{}library interface is designed to be easy to add functionality to,
thus the Python interface to LAMMPS is easy to extend as well.

\sphinxAtStartPar
If you create interesting Python scripts that run LAMMPS or
interesting Python functions that can be called from a LAMMPS input
script, that you think would be generally useful, please post them as
a pull request to our \sphinxhref{https://github.com/lammps/lammps}{GitHub site},
and they can be added to the LAMMPS distribution or web page.

\sphinxstepscope


\chapter{Modifying \& extending LAMMPS}
\label{\detokenize{Modify:modifying-extending-lammps}}\label{\detokenize{Modify::doc}}
\sphinxAtStartPar
LAMMPS has a modular design, so that it is easy to modify or extend with
new functionality.  In fact, about 95\% of its source code is optional.
The following pages give basic instructions on adding new features to
LAMMPS.  More in\sphinxhyphen{}depth explanations and documentation of individual
functions and classes are given in {\hyperref[\detokenize{Developer::doc}]{\sphinxcrossref{\DUrole{doc}{Information for Developers}}}}.

\sphinxAtStartPar
If you add a new feature to LAMMPS and think it will be of general
interest to other users, we encourage you to submit it for inclusion in
LAMMPS. This process is explained in the following three pages:
\begin{itemize}
\item {} 
\sphinxAtStartPar
{\hyperref[\detokenize{Modify_contribute::doc}]{\sphinxcrossref{\DUrole{doc}{how to prepare and submit your code}}}}

\item {} 
\sphinxAtStartPar
{\hyperref[\detokenize{Modify_requirements::doc}]{\sphinxcrossref{\DUrole{doc}{requirements for submissions}}}}

\item {} 
\sphinxAtStartPar
{\hyperref[\detokenize{Modify_style::doc}]{\sphinxcrossref{\DUrole{doc}{style guidelines}}}}

\end{itemize}

\sphinxAtStartPar
A summary description of various types of styles in LAMMPS follows.
A discussion of implementing specific styles from scratch is given
in {\hyperref[\detokenize{Developer_write::doc}]{\sphinxcrossref{\DUrole{doc}{writing new styles}}}}.

\sphinxstepscope


\section{Overview}
\label{\detokenize{Modify_overview:overview}}\label{\detokenize{Modify_overview::doc}}
\sphinxAtStartPar
The best way to add a new feature to LAMMPS is to find a similar
feature and look at the corresponding source and header files to
figure out what it does.  You will need some knowledge of C++ to
understand the high\sphinxhyphen{}level structure of LAMMPS and its class
organization.  Functions (class methods) that do actual computations
are mostly written in C\sphinxhyphen{}style code and operate on simple C\sphinxhyphen{}style data
structures (vectors and arrays).  A high\sphinxhyphen{}level overview of the
programming style choices in LAMMPS is {\hyperref[\detokenize{Developer_code_design::doc}]{\sphinxcrossref{\DUrole{doc}{given elsewhere}}}}.

\sphinxAtStartPar
Most of the new features described on the {\hyperref[\detokenize{Modify::doc}]{\sphinxcrossref{\DUrole{doc}{Modify}}}} doc
page require you to write a new C++ derived class (excluding exceptions
described below, this can often be done by making small edits to existing
files).  Creating a new class requires 2 files, a source code file
(*.cpp) and a header file (*.h).  The derived class must provide
certain methods to work as a new option.  Depending on how different
your new feature is compared to existing features, you can either
derive from the base class itself, or from a derived class that
already exists.  Enabling LAMMPS to invoke the new class is as simple
as putting the two source files in the src directory and re\sphinxhyphen{}building
LAMMPS.

\sphinxAtStartPar
The advantage of C++ and its object\sphinxhyphen{}orientation is that all the code
and variables needed to define the new feature are in the 2 files you
write.  Thus, it should not make the rest of LAMMPS more complex or
cause bugs through unwanted side effects.

\sphinxAtStartPar
Here is a concrete example.  Suppose you write 2 files
\sphinxcode{\sphinxupquote{pair\_foo.cpp}} and \sphinxcode{\sphinxupquote{pair\_foo.h}} that define a new class
\sphinxcode{\sphinxupquote{PairFoo}} which computes pairwise potentials described in the
classic 1997 {\hyperref[\detokenize{Modify_overview:foo}]{\sphinxcrossref{\DUrole{std,std-ref}{paper}}}} by Foo, \sphinxstyleemphasis{et al.} If you wish to invoke
those potentials in a LAMMPS input script with a command like:

\begin{sphinxVerbatim}[commandchars=\\\{\}]
\PYG{k}{pair\PYGZus{}style}\PYG{+w}{ }\PYG{n}{foo}\PYG{+w}{ }\PYG{l+m}{0.1}\PYG{+w}{ }\PYG{l+m}{3.5}
\end{sphinxVerbatim}

\sphinxAtStartPar
then your \sphinxcode{\sphinxupquote{pair\_foo.h}} file should be structured as follows:

\begin{sphinxVerbatim}[commandchars=\\\{\}]
\PYG{c+cp}{\PYGZsh{}}\PYG{c+cp}{ifdef PAIR\PYGZus{}CLASS}
\PYG{c+c1}{// clang\PYGZhy{}format off}
\PYG{n}{PairStyle}\PYG{p}{(}\PYG{n}{foo}\PYG{p}{,}\PYG{n}{PairFoo}\PYG{p}{)}\PYG{p}{;}
\PYG{c+cp}{\PYGZsh{}}\PYG{c+cp}{else}
\PYG{c+c1}{// clang\PYGZhy{}format on}
\PYG{p}{.}\PYG{p}{.}\PYG{p}{.}
\PYG{p}{(}\PYG{k}{class}\PYG{+w}{ }\PYG{n+nc}{definition}\PYG{+w}{ }\PYG{k}{for}\PYG{+w}{ }\PYG{n}{PairFoo}\PYG{p}{)}
\PYG{p}{.}\PYG{p}{.}\PYG{p}{.}
\PYG{c+cp}{\PYGZsh{}}\PYG{c+cp}{endif}
\end{sphinxVerbatim}

\sphinxAtStartPar
where “foo” is the style keyword in the pair\_style command, and
\sphinxcode{\sphinxupquote{PairFoo}} is the class name defined in your \sphinxcode{\sphinxupquote{pair\_foo.cpp}} and
\sphinxcode{\sphinxupquote{pair\_foo.h}} files.

\sphinxAtStartPar
When you re\sphinxhyphen{}build LAMMPS, your new pairwise potential becomes part of
the executable and can be invoked with a pair\_style command like the
example above.  Arguments like 0.1 and 3.5 can be defined and
processed by your new class.

\sphinxAtStartPar
As illustrated by this example, many features referred to in the
LAMMPS documentation are called a “style” of a particular command.

\sphinxAtStartPar
The {\hyperref[\detokenize{Modify::doc}]{\sphinxcrossref{\DUrole{doc}{Modify page}}}} lists all the common styles in LAMMPS,
and discusses the header file for the base class that these styles
derive from.  Public variables in that file can be used and set by the
derived classes, and may also be used by the base class.  Sometimes
they are also accessed by the rest of LAMMPS.  Pure functions, which
means functions declared as virtual in the base class header file and
which are also set to 0, are functions you \sphinxstylestrong{must} implement in your
new derived class to give it the functionality LAMMPS expects. Virtual
functions that are not set to 0 are functions you may override or not.
Those are usually defined with an empty function body.

\sphinxAtStartPar
Additionally, new output options can be added directly to the
thermo.cpp, dump\_custom.cpp, and variable.cpp files.  These are also
listed on the {\hyperref[\detokenize{Modify::doc}]{\sphinxcrossref{\DUrole{doc}{Modify page}}}}.

\sphinxAtStartPar
Here are additional guidelines for modifying LAMMPS and adding new
functionality:
\begin{itemize}
\item {} 
\sphinxAtStartPar
Think about whether what you want to do would be better as a pre\sphinxhyphen{} or
post\sphinxhyphen{}processing step.  Many computations are more easily and more
quickly done that way.

\item {} 
\sphinxAtStartPar
Do not try to do anything within the timestepping of a run that is not
parallel.  For example, do not accumulate a bunch of data on a single
processor and analyze it.  That would run the risk of seriously degrading
the parallel efficiency.

\item {} 
\sphinxAtStartPar
If your new feature reads arguments or writes output, make sure you
follow the unit conventions discussed by the {\hyperref[\detokenize{units::doc}]{\sphinxcrossref{\DUrole{doc}{units}}}}
command.

\end{itemize}


\bigskip\hrule\bigskip

\phantomsection\label{\detokenize{Modify_overview:foo}}
\sphinxAtStartPar
\sphinxstylestrong{(Foo)} Foo, Morefoo, and Maxfoo, J of Classic Potentials, 75, 345 (1997).

\sphinxstepscope


\section{Submitting new features for inclusion in LAMMPS}
\label{\detokenize{Modify_contribute:submitting-new-features-for-inclusion-in-lammps}}\label{\detokenize{Modify_contribute::doc}}
\sphinxAtStartPar
We encourage LAMMPS users to submit new features they write for LAMMPS
to be included in the LAMMPS distribution and thus become easily
accessible to all LAMMPS users.  The LAMMPS source code is managed
with git and public development is hosted on \sphinxhref{https://github.com/lammps/lammps}{GitHub}.  You can monitor the repository
to be notified of releases, follow the ongoing development, and
comment on topics of interest to you.

\sphinxAtStartPar
This section contains general information regarding the preparation
and submission of new features to LAMMPS. If you are new to
development in LAMMPS, we recommend you read one of the tutorials on
developing a new {\hyperref[\detokenize{Developer_write_pair::doc}]{\sphinxcrossref{\DUrole{doc}{pair style}}}} or {\hyperref[\detokenize{Developer_write_fix::doc}]{\sphinxcrossref{\DUrole{doc}{fix
style}}}} which provide a friendly introduction to
what LAMMPS development entails and common vocabulary used on this
section.


\subsection{Communication with the LAMMPS developers}
\label{\detokenize{Modify_contribute:communication-with-the-lammps-developers}}
\sphinxAtStartPar
For any larger modifications or programming project, you are
encouraged to contact the LAMMPS developers ahead of time to discuss
implementation strategies. That will make it easier to integrate your
contribution and typically results in less work for everyone involved.
You are also encouraged to search through the list of \sphinxhref{https://github.com/lammps/lammps/issues}{open issues on
GitHub} and submit a new
issue for a planned feature, to avoid duplicating work (and possibly
being scooped).

\sphinxAtStartPar
For informal communication with the LAMMPS developers, you may ask to
join the \sphinxhref{https://lammps.slack.com}{LAMMPS developers on Slack}.
This slack work space is by invitation only.  For access, please send
an e\sphinxhyphen{}mail to \sphinxcode{\sphinxupquote{slack@lammps.org}} explaining what part of LAMMPS you
are working on.  Only discussions related to LAMMPS development are
tolerated in that work space, so this is \sphinxstylestrong{NOT} for people looking
for help with compiling, installing, or using LAMMPS.  Please post a
message to the \sphinxhref{https://www.lammps.org/forum.html}{LAMMPS forum} for
those purposes.


\subsection{Time and effort required}
\label{\detokenize{Modify_contribute:time-and-effort-required}}
\sphinxAtStartPar
How quickly your contribution will be integrated can vary widely.  It
depends largely on how much effort is required by the LAMMPS
developers to integrate and test it, if any and what kind of changes
to the core code are required, how quickly you can address them, and
how much interest the contribution is to the larger LAMMPS
community. This process can be streamlined by following the
{\hyperref[\detokenize{Modify_requirements::doc}]{\sphinxcrossref{\DUrole{doc}{requirements}}}} and {\hyperref[\detokenize{Modify_style::doc}]{\sphinxcrossref{\DUrole{doc}{style
guidelines}}}}.  A small, modular, well written
contribution may be integrated within hours, but a complex change that
requires a re\sphinxhyphen{}design of a core functionality in LAMMPS can take months
before inclusion (though this is rare).


\subsection{Submission procedure}
\label{\detokenize{Modify_contribute:submission-procedure}}
\sphinxAtStartPar
All changes to LAMMPS (including those from LAMMPS developers) are
integrated via pull requests on GitHub and cannot be merged without
passing the automated testing and an approving review by a LAMMPS core
developer.  Before submitting your contribution, you should therefore
first ensure that your added or modified code compiles and works
correctly with the latest development version of LAMMPS and contains
all bug fixes from it.

\sphinxAtStartPar
Once you have prepared everything, see the {\hyperref[\detokenize{Howto_github::doc}]{\sphinxcrossref{\DUrole{doc}{LAMMPS GitHub
Tutorial}}}} page for instructions on how to submit your
changes or new files through a GitHub pull request.  If you are unable
or unwilling to submit via GitHub yourself, you may also send patch
files or full files to the \sphinxhref{https://www.lammps.org/authors.html}{LAMMPS developers} and ask them to submit a pull
request on GitHub on your behalf.  If this is the case, create a
gzipped tar file of all new or changed files or a corresponding patch
file using ‘diff \sphinxhyphen{}u’ or ‘diff \sphinxhyphen{}c’ format and compress it with gzip.
Please only use gzip compression, as this works well and is available
on all platforms.  This mode of submission may delay the integration
as it depends more on the LAMMPS developers.


\subsection{External contributions}
\label{\detokenize{Modify_contribute:external-contributions}}
\sphinxAtStartPar
If you prefer to do so, you can also develop and support your add\sphinxhyphen{}on
feature \sphinxstylestrong{without} having it included in the LAMMPS distribution, for
example as a download from a website of your own.  See the \sphinxhref{https://www.lammps.org/external.html}{External
LAMMPS packages and tools}
page of the LAMMPS website for examples of groups that do this.  We
are happy to advertise your package and website from that page.
Simply email the \sphinxhref{https://www.lammps.org/authors.html}{developers}
with info about your package, and we will post it there.  We recommend
naming external packages USER\sphinxhyphen{}\textless{}name\textgreater{} so they can be easily
distinguished from packages in the LAMMPS distribution which do not
have the USER\sphinxhyphen{} prefix.


\subsection{Location of files: individual files and packages}
\label{\detokenize{Modify_contribute:location-of-files-individual-files-and-packages}}
\sphinxAtStartPar
We rarely accept new styles in the core src folder.  Thus, please
review the list of {\hyperref[\detokenize{Packages_details::doc}]{\sphinxcrossref{\DUrole{doc}{available Packages}}}} to see
if your contribution should be added to one of them.  It should fit
into the general purpose of that package.  If it does not fit well, it
may be added to one of the EXTRA\sphinxhyphen{} packages or the MISC package.

\sphinxAtStartPar
However, if your project includes many related features that are not
covered by one of the existing packages or is dependent on a library
(bundled or external), it is best to create a new package with its own
directory (with a name like FOO).  In addition to your new files, the
directory should contain a README text file containing your name and
contact information and a brief description of what your new package
does.


\subsection{Changes to core LAMMPS files}
\label{\detokenize{Modify_contribute:changes-to-core-lammps-files}}
\sphinxAtStartPar
If designed correctly, most additions do not require any changes to
the core code of LAMMPS; they are simply add\sphinxhyphen{}on files that are
compiled with the rest of LAMMPS.  To make those styles work, you may
need some trivial changes to the core code.  An example of a trivial
change is making a parent\sphinxhyphen{}class method “virtual” when you derive a new
child class from it.  If your features involve more substantive
changes to the core LAMMPS files, it is particularly encouraged that
you communicate with the LAMMPS developers early in development.

\sphinxstepscope


\section{Requirements for contributions to LAMMPS}
\label{\detokenize{Modify_requirements:requirements-for-contributions-to-lammps}}\label{\detokenize{Modify_requirements::doc}}
\sphinxAtStartPar
The following is a summary of the current requirements and
recommendations for including contributed source code or documentation
into the LAMMPS software distribution.


\subsection{Motivation}
\label{\detokenize{Modify_requirements:motivation}}
\sphinxAtStartPar
The LAMMPS developers are committed to provide a software package that
is versatile, reliable, high\sphinxhyphen{}quality, efficient, portable, and easy to
maintain and modify.  Achieving all of these goals is challenging
since a large part of LAMMPS consists of contributed code from many
different authors who may not be professionally trained programmers or
familiar with the idiosyncrasies of maintaining a large software
package.  In addition, changes that interfere with the parallel
efficiency of the core code must be avoided.  As LAMMPS continues to
grow and more features and functionality are added, it is necessary to
follow established guidelines when accepting new contributions while
also working at the same time to improve the existing code.

\sphinxAtStartPar
The following requirements and recommendations are provided as a
guide.  They indicate which individual requirements are strict, and
which represent a preference and thus are negotiable or optional.
Please feel free to contact the LAMMPS core developers in case you
need additional explanations or clarifications, or you need assistance
in implementing the (strict) requirements for your contributions.
Requirements include:
\begin{itemize}
\item {} 
\sphinxAtStartPar
{\hyperref[\detokenize{Modify_requirements:reqlicense}]{\sphinxcrossref{\DUrole{std,std-ref}{Licensing requirements}}}} (strict)

\item {} 
\sphinxAtStartPar
{\hyperref[\detokenize{Modify_requirements:reqintegrationtesting}]{\sphinxcrossref{\DUrole{std,std-ref}{Integration testing}}}} (strict)

\item {} 
\sphinxAtStartPar
{\hyperref[\detokenize{Modify_requirements:reqdocumentation}]{\sphinxcrossref{\DUrole{std,std-ref}{Documentation}}}} (strict)

\item {} 
\sphinxAtStartPar
{\hyperref[\detokenize{Modify_requirements:reqprogrammingstandards}]{\sphinxcrossref{\DUrole{std,std-ref}{Programming language standards}}}} (strict)

\item {} 
\sphinxAtStartPar
{\hyperref[\detokenize{Modify_requirements:reqbuildsystem}]{\sphinxcrossref{\DUrole{std,std-ref}{Build system}}}} (strict)

\item {} 
\sphinxAtStartPar
{\hyperref[\detokenize{Modify_requirements:reqnaming}]{\sphinxcrossref{\DUrole{std,std-ref}{Command or style names}}}} (strict)

\item {} 
\sphinxAtStartPar
{\hyperref[\detokenize{Modify_requirements:reqprogrammingstyle}]{\sphinxcrossref{\DUrole{std,std-ref}{Programming style requirements}}}} (varied)

\item {} 
\sphinxAtStartPar
{\hyperref[\detokenize{Modify_requirements:reqexamples}]{\sphinxcrossref{\DUrole{std,std-ref}{Examples}}}} (preferred)

\item {} 
\sphinxAtStartPar
{\hyperref[\detokenize{Modify_requirements:reqerrormessages}]{\sphinxcrossref{\DUrole{std,std-ref}{Error or warning messages and explanations}}}} (preferred)

\item {} 
\sphinxAtStartPar
{\hyperref[\detokenize{Modify_requirements:reqcitation}]{\sphinxcrossref{\DUrole{std,std-ref}{Citation reminder}}}} (optional)

\item {} 
\sphinxAtStartPar
{\hyperref[\detokenize{Modify_requirements:requnittesting}]{\sphinxcrossref{\DUrole{std,std-ref}{Testing}}}} (optional)

\end{itemize}


\subsection{Licensing requirements (strict)}
\label{\detokenize{Modify_requirements:licensing-requirements-strict}}\label{\detokenize{Modify_requirements:reqlicense}}
\sphinxAtStartPar
Contributing authors agree when submitting a pull request that their
contributions can be distributed under the LAMMPS license conditions.
This is the GNU public license in version 2 (not 3 or later) for the
publicly distributed versions, e.g. on the LAMMPS homepage or on
GitHub.  We also have a version of LAMMPS under LGPL 2.1 terms which
is available on request; this will usually be the latest available or
a previous stable version with a few LGPL 2.1 incompatible files
removed.  More details are found on the {\hyperref[\detokenize{Intro_opensource::doc}]{\sphinxcrossref{\DUrole{doc}{LAMMPS open\sphinxhyphen{}source
license page}}}}.

\sphinxAtStartPar
Your new source files should have the LAMMPS copyright and GPL notice,
followed by your name and email address at the top, like other
user\sphinxhyphen{}contributed LAMMPS source files.

\sphinxAtStartPar
Contributions may be under a different license as long as that license
does not conflict with the aforementioned terms.  Contributions that
use code with a conflicting license can be split into two parts:
\begin{enumerate}
\sphinxsetlistlabels{\arabic}{enumi}{enumii}{}{.}%
\item {} 
\sphinxAtStartPar
the core parts (i.e. parts that must be in the \sphinxtitleref{src} tree) that are
licensed under compatible terms and bundled with the LAMMPS sources

\item {} 
\sphinxAtStartPar
an external library that must be downloaded and compiled (either
separately or as part of the LAMMPS compilation)

\end{enumerate}

\sphinxAtStartPar
Please note, that this split licensing mode may complicate including
the contribution in binary packages.


\subsection{Integration testing (strict)}
\label{\detokenize{Modify_requirements:integration-testing-strict}}\label{\detokenize{Modify_requirements:reqintegrationtesting}}
\sphinxAtStartPar
Where possible we use available continuous integration tools to search
for common programming mistakes, portability limitations, incompatible
formatting, and undesired side effects. Contributed code must pass the
automated tests on GitHub before it can be merged with the LAMMPS
distribution. These tests compile LAMMPS in a variety of environments
and settings and run the bundled unit tests.  At the discretion of the
LAMMPS developer managing the pull request, additional tests may be
activated that test for “side effects” on running a collection of
input decks and create consistent results.  The translation of the
documentation to HTML and PDF is also tested.

\sphinxAtStartPar
This means that contributed source code \sphinxstylestrong{must} compile with the most
current version of LAMMPS with \sphinxcode{\sphinxupquote{\sphinxhyphen{}DLAMMPS\_BIGBIG}} in addition to the
default setting of \sphinxcode{\sphinxupquote{\sphinxhyphen{}DLAMMPS\_SMALLBIG}}.  The code needs to work
correctly in both cases, and also in serial and parallel using MPI.

\sphinxAtStartPar
Some “disruptive” changes may break tests and require updates to the
testing tools or scripts or tests themselves.  This is rare.  If in
doubt, contact the LAMMPS developer that is assigned to the pull
request.


\subsection{Documentation (strict)}
\label{\detokenize{Modify_requirements:documentation-strict}}\label{\detokenize{Modify_requirements:reqdocumentation}}
\sphinxAtStartPar
Contributions that add new styles or commands or augment existing ones
must include the corresponding new or modified documentation in
\sphinxhref{https://www.sphinx-doc.org/en/master/usage/restructuredtext/index.html}{ReStructuredText format} (.rst files in the \sphinxcode{\sphinxupquote{doc/src/}}
folder). The documentation should be written in American English and the
.rst file must only use ASCII characters, so it can be cleanly
translated to PDF files (via \sphinxhref{https://www.sphinx-doc.org}{sphinx} and
PDFLaTeX).  Special characters may be included via embedded math
expression typeset in a LaTeX subset.

\sphinxAtStartPar
When adding new commands, they need to be integrated into the sphinx
documentation system, and the corresponding command tables and lists
updated. When translating the documentation into html files there
should be no warnings. When adding a new package, some lists
describing packages must also be updated as well as a package specific
description added.  Likewise, if necessary, some package specific
build instructions should be included.

\sphinxAtStartPar
As appropriate, the text files with the documentation can include
inline mathematical expressions or figures (see \sphinxcode{\sphinxupquote{doc/JPG}} for
examples).  Additional PDF files with further details may also be
included; see \sphinxcode{\sphinxupquote{doc/PDF}} for examples.  The page should also include
literature citations as appropriate; see the bottom of
\sphinxcode{\sphinxupquote{doc/fix\_nh.rst}} for examples and the earlier part of the same file
for how to format the cite itself.  Citation labels must be unique
across \sphinxstylestrong{all} .rst files.  The “Restrictions” section of the page
should indicate if your command is only available if LAMMPS is built
with the appropriate package.  See other command doc files for
examples of how to do this.

\sphinxAtStartPar
Please run at least “make html” and “make spelling” from within the
doc/src directory, and carefully inspect and proofread the resulting
HTML format doc page before submitting your code.  Upon submission of
a pull request, checks for error free completion of the HTML and PDF
build will be performed and also a spell check, a check for correct
anchors and labels, and a check for completeness of references to all
styles in their corresponding tables and lists is run.  In case the
spell check reports false positives, they can be added to the file
\sphinxcode{\sphinxupquote{doc/utils/sphinx\sphinxhyphen{}config/false\_positives.txt}}

\sphinxAtStartPar
Contributions that add or modify the library interface or “public”
APIs from the C++ code or the Fortran module must include suitable
doxygen comments in the source and corresponding changes to the
documentation sources for the “Programmer Guide” guide section of the
LAMMPS manual.

\sphinxAtStartPar
If your feature requires some more complex steps and explanations to
be used correctly or some external or bundled tools or scripts, we
recommend that you also contribute a {\hyperref[\detokenize{Howto::doc}]{\sphinxcrossref{\DUrole{doc}{Howto document}}}}
providing some more background information and some tutorial material.
This can also be used to provide more in\sphinxhyphen{}depth explanations of models
that require use of multiple commands.

\sphinxAtStartPar
As a rule\sphinxhyphen{}of\sphinxhyphen{}thumb, the more clear and self\sphinxhyphen{}explanatory you make your
documentation, README files and examples, and the easier you make it
for people to get started, the more likely it is that users will try
out your new feature.


\subsubsection{Programming language standards (strict)}
\label{\detokenize{Modify_requirements:programming-language-standards-strict}}\label{\detokenize{Modify_requirements:reqprogrammingstandards}}
\sphinxAtStartPar
The core of LAMMPS is written in C++11 in a style that can be mostly
described as “C with classes”.  Advanced C++ features like operator
overloading or excessive use of templates are avoided with the intent to
keep the code readable to programmers that have limited C++ programming
experience.  C++ constructs are acceptable when they help improve the
readability and reliability of the code, e.g. when using the
\sphinxtitleref{std::string} class instead of manipulating pointers and calling the
string functions of the C library.  In addition, a collection of
convenient {\hyperref[\detokenize{Developer_utils::doc}]{\sphinxcrossref{\DUrole{doc}{utility functions and classes}}}} for
recurring tasks and a collection of {\hyperref[\detokenize{Developer_platform::doc}]{\sphinxcrossref{\DUrole{doc}{platform neutral functions}}}} for improved portability are provided.
Contributions with code requiring more recent C++ standards are only
accepted as packages with the post C++11 standard code confined to the
package so that it is optional.

\sphinxAtStartPar
Included Fortran code has to be compatible with the Fortran 2003
standard.  Since not all platforms supported by LAMMPS provide good
support for compiling Fortran files, it should be considered to rewrite
these parts as C++ code, if possible, and thus allow for a wider adoption
of the contribution.  As of January 2023, all previously included
Fortran code for the LAMMPS executable has been replaced by equivalent
C++ code.

\sphinxAtStartPar
Python code must be compatible with Python 3.5 and later.  Large parts
of LAMMPS (including the {\hyperref[\detokenize{Packages_details:pkg-python}]{\sphinxcrossref{\DUrole{std,std-ref}{PYTHON package}}}}) are also
compatible with Python 2.7.  Compatibility with Python 2.7 is desirable,
but compatibility with Python 3.5 is \sphinxstylestrong{required}.

\sphinxAtStartPar
Compatibility with older programming language standards is very
important to maintain portability and availability of LAMMPS on many
platforms.  This applies especially to HPC cluster environments, which
tend to be running older software stacks and where LAMMPS users may be
required to use those older tools for access to advanced hardware
features or not have the option to install newer compilers or libraries.


\subsection{Build system (strict)}
\label{\detokenize{Modify_requirements:build-system-strict}}\label{\detokenize{Modify_requirements:reqbuildsystem}}
\sphinxAtStartPar
LAMMPS currently supports two build systems: one that is based on
{\hyperref[\detokenize{Build_make::doc}]{\sphinxcrossref{\DUrole{doc}{traditional Makefiles}}}} and one that is based on
{\hyperref[\detokenize{Build_cmake::doc}]{\sphinxcrossref{\DUrole{doc}{CMake}}}}.  Therefore, your contribution must be
compatible with and support both build systems.

\sphinxAtStartPar
For a single pair of header and implementation files that are an
independent feature, it is usually only required to add them to
\sphinxcode{\sphinxupquote{src/.gitignore}}.

\sphinxAtStartPar
For traditional make, if your contributed files or package depend on
other LAMMPS style files or packages also being installed
(e.g. because your file is a derived class from the other LAMMPS
class), then an \sphinxcode{\sphinxupquote{Install.sh}} file is also needed to check for those
dependencies and modifications to \sphinxcode{\sphinxupquote{src/Depend.sh}} to trigger the checks.
See other README and Install.sh files in other directories as
examples.

\sphinxAtStartPar
Similarly, for CMake support, changes may need to be made to
\sphinxcode{\sphinxupquote{cmake/CMakeLists.txt}}, some of the files in \sphinxcode{\sphinxupquote{cmake/presets}}, and
possibly a file with specific instructions needs to be added to
\sphinxcode{\sphinxupquote{cmake/Modules/Packages/}}.  Please check out how this is handled for
existing packages and ask the LAMMPS developers if you need assistance.


\subsection{Command or style names, file names, and keywords (strict)}
\label{\detokenize{Modify_requirements:command-or-style-names-file-names-and-keywords-strict}}\label{\detokenize{Modify_requirements:reqnaming}}
\sphinxAtStartPar
All user\sphinxhyphen{}visible command or style names should be all lower case and
should only use letters, numbers, or forward slashes.  They should be
descriptive and initialisms should be avoided unless they are well
established (e.g. lj for Lennard\sphinxhyphen{}Jones).  For a compute style
“some/name” the source files must be called \sphinxcode{\sphinxupquote{compute\_some\_name.h}} and
\sphinxcode{\sphinxupquote{compute\_some\_name.cpp}}. The “include guard” in the header file would
then be \sphinxcode{\sphinxupquote{LMP\_COMPUTE\_SOME\_NAME\_H}} and the class name
\sphinxcode{\sphinxupquote{ComputeSomeName}}.


\subsection{Programming style requirements (varied)}
\label{\detokenize{Modify_requirements:programming-style-requirements-varied}}\label{\detokenize{Modify_requirements:reqprogrammingstyle}}
\sphinxAtStartPar
To maintain source code consistency across contributions from many
people, there are various programming style requirements for
contributions to LAMMPS.  Some of these requirements are strict and
must be followed, while others are only preferred and thus may be
skipped.  An in\sphinxhyphen{}depth discussion of the style guidelines is provided
in the {\hyperref[\detokenize{Modify_style::doc}]{\sphinxcrossref{\DUrole{doc}{programming style doc page}}}}.


\subsection{Examples (preferred)}
\label{\detokenize{Modify_requirements:examples-preferred}}\label{\detokenize{Modify_requirements:reqexamples}}
\sphinxAtStartPar
For many new features, it is preferred that example scripts (simple,
small, fast to complete on 1 CPU) are included that demonstrate the
use of new or extended functionality. These are typically include
under the examples or examples/PACKAGES directory and are further
described on the {\hyperref[\detokenize{Examples::doc}]{\sphinxcrossref{\DUrole{doc}{examples page}}}}.  Guidelines for
input scripts include:
\begin{itemize}
\item {} 
\sphinxAtStartPar
commands that generate output should be commented out (except when the
output is the sole purpose or the feature, e.g. for a new compute)

\item {} 
\sphinxAtStartPar
commands like {\hyperref[\detokenize{log::doc}]{\sphinxcrossref{\DUrole{doc}{log}}}}, {\hyperref[\detokenize{echo::doc}]{\sphinxcrossref{\DUrole{doc}{echo}}}}, {\hyperref[\detokenize{package::doc}]{\sphinxcrossref{\DUrole{doc}{package}}}}, {\hyperref[\detokenize{processors::doc}]{\sphinxcrossref{\DUrole{doc}{processors}}}}, {\hyperref[\detokenize{suffix::doc}]{\sphinxcrossref{\DUrole{doc}{suffix}}}} may
\sphinxstylestrong{not} be used in the input file (exception: “processors * * 1” or
similar is acceptable when used to avoid unwanted domain decomposition
of empty volumes)

\item {} 
\sphinxAtStartPar
outside of the log files, no generated output should be included

\item {} 
\sphinxAtStartPar
custom thermo\_style settings may not include output measuring CPU or other
time as it complicates comparisons between different runs

\item {} 
\sphinxAtStartPar
input files should be named \sphinxcode{\sphinxupquote{in.name}}, data files should be named
\sphinxcode{\sphinxupquote{data.name}} and log files should be named \sphinxcode{\sphinxupquote{log.version.name.\textless{}compiler\textgreater{}.\textless{}ncpu\textgreater{}}}

\item {} 
\sphinxAtStartPar
the total file size of all the inputs and outputs should be small

\item {} 
\sphinxAtStartPar
where possible, potential files from the “potentials” folder or data
file from other folders should be re\sphinxhyphen{}used through symbolic links

\end{itemize}


\subsection{Error or warning messages and explanations (preferred)}
\label{\detokenize{Modify_requirements:error-or-warning-messages-and-explanations-preferred}}\label{\detokenize{Modify_requirements:reqerrormessages}}
\sphinxAtStartPar
\DUrole{versionmodified,changed}{Changed in version 4May2022.}

\sphinxAtStartPar
Starting with LAMMPS version 4 May 2022, the LAMMPS developers have
agreed on a new policy for error and warning messages.

\sphinxAtStartPar
Previously, all error and warning strings were supposed to be listed in
the class header files with an explanation.  Those would then be
regularly “harvested” and transferred to alphabetically sorted lists in
the manual.  To avoid excessively long lists and to reduce effort, this
came with a requirement to have rather generic error messages (e.g.
“Illegal … command”).  To identify the specific cause, the name of the
source file and the line number of the error location would be printed,
so that one could look up the cause by reading the source code.

\sphinxAtStartPar
The new policy encourages more specific error messages that ideally
indicate the cause directly, and requiring no further lookup. This is
aided by the \sphinxhref{https://fmt.dev}{\{fmt\} library} enabling Error class
methods that take a variable number of arguments and an error text that
will be treated like a \{fmt\} syntax format string. Error messages should
still preferably be kept to a single line or two lines at most.

\sphinxAtStartPar
For more complex explanations or errors that have multiple possible
reasons, a paragraph should be added to the \sphinxtitleref{Error\_details} page with an
error code reference (e.g. \sphinxcode{\sphinxupquote{.. \_err0001:}}) then the utility function
{\hyperref[\detokenize{Developer_utils:_CPPv4N9LAMMPS_NS5utils8errorurlEi}]{\sphinxcrossref{\sphinxcode{\sphinxupquote{utils::errorurl()}}}}} can be used
to generate a URL that will directly lead to that paragraph.  An error
for missing arguments can be easily generated using the
{\hyperref[\detokenize{Developer_utils:_CPPv4N9LAMMPS_NS5utils16missing_cmd_argsERKNSt6stringEiRKNSt6stringEP5Error}]{\sphinxcrossref{\sphinxcode{\sphinxupquote{utils::missing\_cmd\_args()}}}}} convenience function.
An example for this approach would be the
\sphinxcode{\sphinxupquote{src/read\_data.cpp}} and \sphinxcode{\sphinxupquote{src/atom.cpp}} files that implement the
{\hyperref[\detokenize{read_data::doc}]{\sphinxcrossref{\DUrole{doc}{read\_data}}}} and {\hyperref[\detokenize{atom_modify::doc}]{\sphinxcrossref{\DUrole{doc}{atom\_modify}}}}
commands and that may create {\hyperref[\detokenize{Errors_details:err0001}]{\sphinxcrossref{\DUrole{std,std-ref}{“Unknown identifier in data file”}}}}
errors that may have multiple possible reasons which complicates debugging,
and thus require some additional explanation.

\sphinxAtStartPar
The transformation of existing LAMMPS code to this new scheme is
ongoing.  Given the size of the LAMMPS code base, it will take a
significant amount of time to complete.  For new code, however,
following the new approach is strongly preferred.  The expectation is
that the new scheme will make understanding errors easier for LAMMPS
users, developers, and maintainers.


\subsection{Citation reminder (optional)}
\label{\detokenize{Modify_requirements:citation-reminder-optional}}\label{\detokenize{Modify_requirements:reqcitation}}
\sphinxAtStartPar
If there is a paper of yours describing your feature (either the
algorithm/science behind the feature itself, or its initial usage, or
its implementation in LAMMPS), you can add the citation to the *.cpp
source file.  See \sphinxcode{\sphinxupquote{src/DIFFRACTION/compute\_saed.cpp}} for an example.
A BibTeX format citation is stored in a string variable at the top of
the file, and a single line of code registering this variable is added
to the constructor of the class.  When your feature is used, then
LAMMPS (by default) will print the brief info and the DOI in the first
line to the screen and the full citation to the log file.

\sphinxAtStartPar
If there is additional functionality (which may have been added later)
described in a different publication, additional citation descriptions
may be added so long as they are only registered when the
corresponding keyword activating this functionality is used.

\sphinxAtStartPar
With these options, it is possible to have LAMMPS output a specific
citation reminder whenever a user invokes your feature from their
input script.  Please note that you should \sphinxstyleemphasis{only} use this for the
\sphinxstyleemphasis{most} relevant paper for a feature and a publication that you or your
group authored.  E.g. adding a citation in the source code for a paper
by Nose and Hoover if you write a fix that implements their integrator
is not the intended usage.  That kind of citation should just be
included in the documentation page you provide describing your
contribution.  If you are not sure what the best option would be,
please contact the LAMMPS developers for advice.


\subsection{Testing (optional)}
\label{\detokenize{Modify_requirements:testing-optional}}\label{\detokenize{Modify_requirements:requnittesting}}
\sphinxAtStartPar
If your contribution contains new utility functions or a supporting
class (i.e. anything that does not depend on a LAMMPS object), new
unit tests should be added to a suitable folder in the \sphinxcode{\sphinxupquote{unittest}}
tree.  When adding a new LAMMPS style computing forces or selected
fixes, a \sphinxcode{\sphinxupquote{.yaml}} file with a test configuration and reference data
should be added for the styles where a suitable tester program already
exists (e.g. pair styles, bond styles, etc.). Please see {\hyperref[\detokenize{Build_development:testing}]{\sphinxcrossref{\DUrole{std,std-ref}{this
section in the manual}}}} for more information on how to
enable, run, and expand testing.

\sphinxstepscope


\section{LAMMPS programming style}
\label{\detokenize{Modify_style:lammps-programming-style}}\label{\detokenize{Modify_style::doc}}
\sphinxAtStartPar
The aim of the LAMMPS developers is to use a consistent programming
style and naming conventions across the entire code base, as this
helps with maintenance, debugging, and understanding the code, both
for developers and users.  This page provides a list of standard style
choices used in LAMMPS.  Some of these standards are required, while
others are just preferred.  Following these conventions will make it
much easier to integrate your contribution.  If you are uncertain,
please ask.

\sphinxAtStartPar
The files \sphinxtitleref{pair\_lj\_cut.h}, \sphinxtitleref{pair\_lj\_cut.cpp}, \sphinxtitleref{utils.h}, and
\sphinxtitleref{utils.cpp} may serve as representative examples.


\subsection{Include files (varied)}
\label{\detokenize{Modify_style:include-files-varied}}\begin{itemize}
\item {} 
\sphinxAtStartPar
Header files that define a new LAMMPS style (i.e. that have a
\sphinxcode{\sphinxupquote{SomeStyle(some/name,SomeName);}} macro in them) should only use
the include file for the base class and otherwise use forward
declarations and pointers; when interfacing to a library use the
PIMPL (pointer to implementation) approach where you have a pointer
to a struct that contains all library specific data (and thus
requires the library header) but use a forward declaration and
define the struct only in the implementation file. This is a
\sphinxstylestrong{strict} requirement since this is where type clashes between
packages and hard\sphinxhyphen{}to\sphinxhyphen{}find bugs have regularly manifested in the
past.

\item {} 
\sphinxAtStartPar
Header files, especially those defining a “style”, should only use the
absolute minimum number of include files and \sphinxstylestrong{must not} contain any
\sphinxcode{\sphinxupquote{using}} statements. Typically, that would only be the header for the
base class.  Instead, any include statements should be put in the
corresponding implementation files and forward declarations be used.
For implementation files, the “include what you use” principle should
be employed.  However, there is the notable exception that when the
\sphinxcode{\sphinxupquote{pointers.h}} header is included (or the header of one of the classes
derived from it), certain headers will \sphinxstyleemphasis{always} be included and thus
do not need to be explicitly specified.  These are: \sphinxtitleref{mpi.h},
\sphinxtitleref{cstddef}, \sphinxtitleref{cstdio}, \sphinxtitleref{cstdlib}, \sphinxtitleref{string}, \sphinxtitleref{utils.h}, \sphinxtitleref{vector},
\sphinxtitleref{fmt/format.h}, \sphinxtitleref{climits}, \sphinxtitleref{cinttypes}.  This also means any such file
can assume that \sphinxtitleref{FILE}, \sphinxtitleref{NULL}, and \sphinxtitleref{INT\_MAX} are defined.

\item {} 
\sphinxAtStartPar
Class members variables should not be initialized in the header file,
but instead should be initialized either in the initializer list of
the constructor or explicitly assigned in the body of the constructor.
If the member variable is relevant to the functionality of a class
(for example when it stores a value from a command line argument), the
member variable declaration is followed by a brief comment explaining
its purpose and what its values can be.  Class members that are
pointers should always be initialized to \sphinxcode{\sphinxupquote{nullptr}} in the
initializer list of the constructor.  This reduces clutter in the
header and avoids accessing uninitialized pointers, which leads to
hard to debug issues, class members are often implicitly initialized
to \sphinxcode{\sphinxupquote{NULL}} on the first use (but \sphinxstyleemphasis{not} after a {\hyperref[\detokenize{clear::doc}]{\sphinxcrossref{\DUrole{doc}{clear command}}}}).  Please see the files \sphinxcode{\sphinxupquote{reset\_atoms\_mol.h}} and
\sphinxcode{\sphinxupquote{reset\_atoms\_mol.cpp}} as an example.

\item {} 
\sphinxAtStartPar
System headers or headers from installed libraries are included with
angular brackets (example: \sphinxcode{\sphinxupquote{\#include \textless{}vector\textgreater{}}}), while local
include files use double quotes (example: \sphinxcode{\sphinxupquote{\#include "atom.h"}})

\item {} 
\sphinxAtStartPar
When including system header files from the C library use the
C++\sphinxhyphen{}style names (\sphinxcode{\sphinxupquote{\textless{}cstdlib\textgreater{}}} or \sphinxcode{\sphinxupquote{\textless{}cstring\textgreater{}}}) instead of the
C\sphinxhyphen{}style names (\sphinxcode{\sphinxupquote{\textless{}stdlib.h\textgreater{}}} or \sphinxcode{\sphinxupquote{\textless{}string.h\textgreater{}}})

\item {} 
\sphinxAtStartPar
The order of \sphinxcode{\sphinxupquote{\#include}} statements in a file \sphinxcode{\sphinxupquote{some\_name.cpp}}
that implements a class \sphinxcode{\sphinxupquote{SomeName}} defined in a header file
\sphinxcode{\sphinxupquote{some\_name.h}} should be as follows:
\begin{itemize}
\item {} 
\sphinxAtStartPar
\sphinxcode{\sphinxupquote{\#include "some\_name.h"}} followed by an empty line

\item {} 
\sphinxAtStartPar
LAMMPS include files e.g. \sphinxcode{\sphinxupquote{\#include "comm.h"}} or \sphinxcode{\sphinxupquote{\#include
"modify.h"}} in alphabetical order followed by an empty line

\item {} 
\sphinxAtStartPar
System header files from the C++ or C standard library followed by
an empty line

\item {} 
\sphinxAtStartPar
\sphinxcode{\sphinxupquote{using namespace LAMMPS\_NS}} or other namespace imports.

\end{itemize}

\end{itemize}


\subsection{Whitespace (preferred)}
\label{\detokenize{Modify_style:whitespace-preferred}}
\sphinxAtStartPar
Source files should not contain TAB characters unless required by the
syntax (e.g. in makefiles) and no trailing whitespace.  Text files
should have Unix\sphinxhyphen{}style line endings (LF\sphinxhyphen{}only). Git will automatically
convert those in both directions when running on Windows; use dos2unix
on Linux machines to convert files to Unix\sphinxhyphen{}style line endings.  The
last line of text files include a line ending.

\sphinxAtStartPar
You can check for these issues with the python scripts in the
{\hyperref[\detokenize{Tools:coding-standard}]{\sphinxcrossref{\DUrole{std,std-ref}{“tools/coding\_standard”}}}} folder.  When run
normally with a source file or a source folder as argument, they will
list all non\sphinxhyphen{}conforming lines.  By adding the \sphinxtitleref{\sphinxhyphen{}f} flag to the command
line, they will modify the flagged files to try to remove the detected
issues.


\subsection{Constants (strongly preferred)}
\label{\detokenize{Modify_style:constants-strongly-preferred}}
\sphinxAtStartPar
Global or per\sphinxhyphen{}file constants should be declared as \sphinxtitleref{static constexpr}
variables rather than via the pre\sphinxhyphen{}processor with \sphinxtitleref{\#define}.  The name of
constants should be all uppercase.  This has multiple advantages:
\begin{itemize}
\item {} 
\sphinxAtStartPar
constants are easily identified as such by their all upper case name

\item {} 
\sphinxAtStartPar
rather than a pure text substitution during pre\sphinxhyphen{}processing, \sphinxtitleref{constexpr
variables} have a type associated with them and are processed later in
the parsing process where the syntax checks and type specific
processing (e.g. via overloads) can be applied to them.

\item {} 
\sphinxAtStartPar
compilers can emit a warning if the constant is not used and thus can
be removed (we regularly check for and remove dead code like this)

\item {} 
\sphinxAtStartPar
there are no unexpected substitutions and thus confusing syntax errors
when compiling leading to, for instance, conflicts so that LAMMPS
cannot be compiled with certain combinations of packages (this \sphinxstyleemphasis{has}
happened multiple times in the past).

\end{itemize}

\sphinxAtStartPar
Pre\sphinxhyphen{}processor defines should be limited to macros (but consider C++
templates) and conditional compilation.  If a per\sphinxhyphen{}processor define must
be used, it should be defined at the top of the .cpp file after the
include statements and at all cost it should be avoided to put them into
header files.

\sphinxAtStartPar
Some sets of commonly used constants are provided in the \sphinxcode{\sphinxupquote{MathConst}}
and \sphinxcode{\sphinxupquote{EwaldConst}} namespaces and implemented in the files
\sphinxcode{\sphinxupquote{math\_const.h}} and \sphinxcode{\sphinxupquote{ewald\_const.h}}, respectively.

\sphinxAtStartPar
There are always exceptions, special cases, and legacy code in LAMMPS,
so please contact the LAMMPS developers if you are not sure.


\subsection{Placement of braces (strongly preferred)}
\label{\detokenize{Modify_style:placement-of-braces-strongly-preferred}}
\sphinxAtStartPar
For new files added to the “src” tree, a \sphinxhref{https://clang.llvm.org/docs/ClangFormat.html}{clang\sphinxhyphen{}format} configuration file is
provided under the name \sphinxtitleref{.clang\sphinxhyphen{}format}.  This file is compatible with
clang\sphinxhyphen{}format version 8 and later. With that file present, files can be
reformatted according to the configuration with a command like:
\sphinxtitleref{clang\sphinxhyphen{}format \sphinxhyphen{}i new\sphinxhyphen{}file.cpp}.  Ideally, this is done while writing
the code or before a pull request is submitted.  Blocks of code where
the reformatting from clang\sphinxhyphen{}format yields hard\sphinxhyphen{}to\sphinxhyphen{}read or otherwise
undesirable output may be protected with placing a pair \sphinxtitleref{//
clang\sphinxhyphen{}format off} and \sphinxtitleref{// clang\sphinxhyphen{}format on} comments around that block.


\subsection{Miscellaneous standards (varied)}
\label{\detokenize{Modify_style:miscellaneous-standards-varied}}\begin{itemize}
\item {} 
\sphinxAtStartPar
I/O is done via the C\sphinxhyphen{}style stdio library and \sphinxstylestrong{not} iostreams.

\item {} 
\sphinxAtStartPar
Do not use so\sphinxhyphen{}called “alternative tokens” like \sphinxcode{\sphinxupquote{and}}, \sphinxcode{\sphinxupquote{or}},
\sphinxcode{\sphinxupquote{not}} and similar, but rather use the corresponding operators
\sphinxcode{\sphinxupquote{\&\&}}, \sphinxcode{\sphinxupquote{||}}, and \sphinxcode{\sphinxupquote{!}}.  The alternative tokens are not available
by default on all compilers.

\item {} 
\sphinxAtStartPar
Output to the screen and the logfile should use the corresponding
FILE pointers and only be done on MPI rank 0.  Use the
\sphinxcode{\sphinxupquote{utils::logmesg()}} convenience function where possible.

\item {} 
\sphinxAtStartPar
Usage of C++11 \sphinxtitleref{virtual}, \sphinxtitleref{override}, \sphinxtitleref{final} keywords: Please
follow the \sphinxhref{https://isocpp.github.io/CppCoreGuidelines/CppCoreGuidelines\#Rh-override}{C++ Core Guideline C.128}.
That means, you should only use \sphinxtitleref{virtual} to declare a new virtual
function, \sphinxtitleref{override} to indicate you are overriding an existing
virtual function, and \sphinxtitleref{final} to prevent any further overriding.

\item {} 
\sphinxAtStartPar
Trivial destructors: Do not write destructors when they are empty
and \sphinxtitleref{default}.

\begin{sphinxVerbatim}[commandchars=\\\{\}]
\PYG{c+c1}{// don\PYGZsq{}t write destructors for A or B like this}

\PYG{k}{class}\PYG{+w}{ }\PYG{n+nc}{A}\PYG{+w}{ }\PYG{o}{:}\PYG{+w}{ }\PYG{k}{protected}\PYG{+w}{ }\PYG{n}{Pointers}\PYG{+w}{ }\PYG{p}{\PYGZob{}}
\PYG{+w}{ }\PYG{k}{public}\PYG{o}{:}
\PYG{+w}{   }\PYG{n}{A}\PYG{p}{(}\PYG{p}{)}\PYG{p}{;}
\PYG{+w}{   }\PYG{o}{\PYGZti{}}\PYG{n}{A}\PYG{p}{(}\PYG{p}{)}\PYG{+w}{ }\PYG{k}{override}\PYG{+w}{ }\PYG{p}{\PYGZob{}}\PYG{p}{\PYGZcb{}}
\PYG{p}{\PYGZcb{}}\PYG{p}{;}

\PYG{k}{class}\PYG{+w}{ }\PYG{n+nc}{B}\PYG{+w}{ }\PYG{o}{:}\PYG{+w}{ }\PYG{k}{protected}\PYG{+w}{ }\PYG{n}{Pointers}\PYG{+w}{ }\PYG{p}{\PYGZob{}}
\PYG{+w}{ }\PYG{k}{public}\PYG{o}{:}
\PYG{+w}{   }\PYG{n}{B}\PYG{p}{(}\PYG{p}{)}\PYG{p}{;}
\PYG{+w}{   }\PYG{o}{\PYGZti{}}\PYG{n}{B}\PYG{p}{(}\PYG{p}{)}\PYG{+w}{ }\PYG{k}{override}\PYG{+w}{ }\PYG{o}{=}\PYG{+w}{ }\PYG{k}{default}\PYG{p}{;}
\PYG{p}{\PYGZcb{}}\PYG{p}{;}

\PYG{c+c1}{// instead, let the compiler create the implicit default destructor by not writing it}

\PYG{k}{class}\PYG{+w}{ }\PYG{n+nc}{A}\PYG{+w}{ }\PYG{o}{:}\PYG{+w}{ }\PYG{k}{protected}\PYG{+w}{ }\PYG{n}{Pointers}\PYG{+w}{ }\PYG{p}{\PYGZob{}}
\PYG{+w}{ }\PYG{k}{public}\PYG{o}{:}
\PYG{+w}{   }\PYG{n}{A}\PYG{p}{(}\PYG{p}{)}\PYG{p}{;}
\PYG{p}{\PYGZcb{}}\PYG{p}{;}

\PYG{k}{class}\PYG{+w}{ }\PYG{n+nc}{B}\PYG{+w}{ }\PYG{o}{:}\PYG{+w}{ }\PYG{k}{protected}\PYG{+w}{ }\PYG{n}{Pointers}\PYG{+w}{ }\PYG{p}{\PYGZob{}}
\PYG{+w}{ }\PYG{k}{public}\PYG{o}{:}
\PYG{+w}{   }\PYG{n}{B}\PYG{p}{(}\PYG{p}{)}\PYG{p}{;}
\PYG{p}{\PYGZcb{}}\PYG{p}{;}
\end{sphinxVerbatim}

\item {} 
\sphinxAtStartPar
Please use clang\sphinxhyphen{}format only to reformat files that you have
contributed.  For header files containing a \sphinxcode{\sphinxupquote{SomeStyle(keyword,
ClassName)}} macros it is required to have this macro embedded with
a pair of \sphinxcode{\sphinxupquote{// clang\sphinxhyphen{}format off}}, \sphinxcode{\sphinxupquote{// clang\sphinxhyphen{}format on}} comments
and the line must be terminated with a semicolon (;).  Example:

\begin{sphinxVerbatim}[commandchars=\\\{\}]
\PYG{c+cp}{\PYGZsh{}}\PYG{c+cp}{ifdef COMMAND\PYGZus{}CLASS}
\PYG{c+c1}{// clang\PYGZhy{}format off}
\PYG{n}{CommandStyle}\PYG{p}{(}\PYG{n}{run}\PYG{p}{,}\PYG{n}{Run}\PYG{p}{)}\PYG{p}{;}
\PYG{c+c1}{// clang\PYGZhy{}format on}
\PYG{c+cp}{\PYGZsh{}}\PYG{c+cp}{else}

\PYG{c+cp}{\PYGZsh{}}\PYG{c+cp}{ifndef LMP\PYGZus{}RUN\PYGZus{}H}
\PYG{p}{[}\PYG{p}{.}\PYG{p}{.}\PYG{p}{.}\PYG{p}{]}
\end{sphinxVerbatim}

\sphinxAtStartPar
You may also use \sphinxcode{\sphinxupquote{// clang\sphinxhyphen{}format on/off}} throughout your files to
protect individual sections from being reformatted.

\item {} 
\sphinxAtStartPar
All files should have 0644 permissions, i.e. writable by the user
only and readable by all and no executable permissions.  Executable
permissions (0755) should only be for shell scripts or python or
similar scripts for interpreted script languages.

\end{itemize}

\sphinxstepscope


\section{Atom styles}
\label{\detokenize{Modify_atom:atom-styles}}\label{\detokenize{Modify_atom::doc}}
\sphinxAtStartPar
Classes that define an {\hyperref[\detokenize{atom_style::doc}]{\sphinxcrossref{\DUrole{doc}{atom style}}}} are derived from
the AtomVec class and managed by the Atom class.  The atom style
determines what attributes are associated with an atom and
communicated when it is a ghost atom or migrates to a new processor.
A new atom style can be created if one of the existing atom styles
does not define all the attributes you need to store and communicate
with atoms.

\sphinxAtStartPar
The file \sphinxcode{\sphinxupquote{atom\_vec\_atomic.cpp}} is the simplest example of an atom style.
Examining the code for others will make these instructions more clear.

\sphinxAtStartPar
Note that the {\hyperref[\detokenize{atom_style::doc}]{\sphinxcrossref{\DUrole{doc}{atom style hybrid}}}} command can be used
to define atoms or particles which have the union of properties of
individual styles.  Also the {\hyperref[\detokenize{fix_property_atom::doc}]{\sphinxcrossref{\DUrole{doc}{fix property/atom}}}} command can be used to add a single property
(e.g. charge or a molecule ID) to a style that does not have it.  It can
also be used to add custom properties to an atom, with options to
communicate them with ghost atoms or read them from a data file.  Other
LAMMPS commands can access these custom properties, as can new pair,
fix, compute styles that are written to work with these properties.  For
example, the {\hyperref[\detokenize{set::doc}]{\sphinxcrossref{\DUrole{doc}{set}}}} command can be used to set the values of
custom per\sphinxhyphen{}atom properties from an input script.  All of these methods
are less work than writing and testing(!) code for a new atom style.

\sphinxAtStartPar
If you follow these directions your new style will automatically work
in tandem with others via the {\hyperref[\detokenize{atom_style::doc}]{\sphinxcrossref{\DUrole{doc}{atom\_style hybrid}}}}
command.

\sphinxAtStartPar
The first step is to define a set of string lists in the constructor of
the new derived class.  Each list will have zero or more comma\sphinxhyphen{}separated
strings that correspond to the variable names used in the \sphinxcode{\sphinxupquote{atom.h}}
header file for per\sphinxhyphen{}atom properties.  Note that some represent per\sphinxhyphen{}atom
vectors (q, molecule) while other are per\sphinxhyphen{}atom arrays (x,v).  For all
but the last two lists you do not need to specify any of
(id,type,x,v,f).  Those are included automatically as needed in the
other lists.


\begin{savenotes}\sphinxattablestart
\sphinxthistablewithglobalstyle
\centering
\begin{tabulary}{\linewidth}[t]{TT}
\sphinxtoprule
\sphinxtableatstartofbodyhook
\sphinxAtStartPar
fields\_grow
&
\sphinxAtStartPar
full list of properties which is allocated and stored
\\
\sphinxhline
\sphinxAtStartPar
fields\_copy
&
\sphinxAtStartPar
list of properties to copy atoms are rearranged on\sphinxhyphen{}processor
\\
\sphinxhline
\sphinxAtStartPar
fields\_comm
&
\sphinxAtStartPar
list of properties communicated to ghost atoms every step
\\
\sphinxhline
\sphinxAtStartPar
fields\_comm\_vel
&
\sphinxAtStartPar
additional properties communicated if {\hyperref[\detokenize{comm_modify::doc}]{\sphinxcrossref{\DUrole{doc}{comm\_modify vel}}}} is used
\\
\sphinxhline
\sphinxAtStartPar
fields\_reverse
&
\sphinxAtStartPar
list of properties summed from ghost atoms every step
\\
\sphinxhline
\sphinxAtStartPar
fields\_border
&
\sphinxAtStartPar
list of properties communicated with ghost atoms every reneighboring step
\\
\sphinxhline
\sphinxAtStartPar
fields\_border\_vel
&
\sphinxAtStartPar
additional properties communicated if {\hyperref[\detokenize{comm_modify::doc}]{\sphinxcrossref{\DUrole{doc}{comm\_modify vel}}}} is used
\\
\sphinxhline
\sphinxAtStartPar
fields\_exchange
&
\sphinxAtStartPar
list of properties communicated when an atom migrates to another processor
\\
\sphinxhline
\sphinxAtStartPar
fields\_restart
&
\sphinxAtStartPar
list of properties written/read to/from a restart file
\\
\sphinxhline
\sphinxAtStartPar
fields\_create
&
\sphinxAtStartPar
list of properties defined when an atom is created by {\hyperref[\detokenize{create_atoms::doc}]{\sphinxcrossref{\DUrole{doc}{create\_atoms}}}}
\\
\sphinxhline
\sphinxAtStartPar
fields\_data\_atom
&
\sphinxAtStartPar
list of properties (in order) in the Atoms section of a data file, as read by {\hyperref[\detokenize{read_data::doc}]{\sphinxcrossref{\DUrole{doc}{read\_data}}}}
\\
\sphinxhline
\sphinxAtStartPar
fields\_data\_vel
&
\sphinxAtStartPar
list of properties (in order) in the Velocities section of a data file, as read by {\hyperref[\detokenize{read_data::doc}]{\sphinxcrossref{\DUrole{doc}{read\_data}}}}
\\
\sphinxbottomrule
\end{tabulary}
\sphinxtableafterendhook\par
\sphinxattableend\end{savenotes}

\sphinxAtStartPar
In these lists you can list variable names which LAMMPS already defines
(in some other atom style), or you can create new variable names.  You
should not re\sphinxhyphen{}use a LAMMPS variable in your atom style that is used for
something with a different meaning in another atom style.  If the
meaning is related, but interpreted differently by your atom style, then
using the same variable name means a user must not use your style and
the other style together in a {\hyperref[\detokenize{atom_style::doc}]{\sphinxcrossref{\DUrole{doc}{atom\_style hybrid}}}}
command.  Because there will only be one value of the variable and
different parts of LAMMPS will then likely use it differently.  LAMMPS
has no way of checking for this.

\sphinxAtStartPar
If you are defining new variable names then make them descriptive and
unique to your new atom style.  For example choosing “e” for energy is
a bad choice; it is too generic.  A better choice would be “e\_foo”,
where “foo” is specific to your style.

\sphinxAtStartPar
If any of the variable names in your new atom style do not exist in
LAMMPS, you need to add them to the src/atom.h and atom.cpp files.

\sphinxAtStartPar
Search for the word “customize” or “customization” in these 2 files to
see where to add your variable.  Adding a flag to the 2nd customization
section in \sphinxcode{\sphinxupquote{atom.h}} is only necessary if your code (e.g. a pair style)
needs to check that a per\sphinxhyphen{}atom property is defined.  These flags should
also be set in the constructor of the atom style child class.

\sphinxAtStartPar
In \sphinxcode{\sphinxupquote{atom.cpp}}, aside from the constructor and destructor, there are 3
methods that a new variable name or flag needs to be added to.

\sphinxAtStartPar
In \sphinxcode{\sphinxupquote{Atom::peratom\_create()}} when using the \sphinxcode{\sphinxupquote{Atom::add\_peratom()}}
method, a cols argument of 0 is for per\sphinxhyphen{}atom vectors, a length \textgreater{}
1 is for per\sphinxhyphen{}atom arrays.  Note the use of the extra per\sphinxhyphen{}thread flag and
the add\_peratom\_vary() method when the last dimension of the array is
variable\sphinxhyphen{}length.

\sphinxAtStartPar
Adding the variable name to Atom::extract() enables the per\sphinxhyphen{}atom data
to be accessed through the {\hyperref[\detokenize{Howto_library::doc}]{\sphinxcrossref{\DUrole{doc}{LAMMPS library interface}}}} by a calling code, including from {\hyperref[\detokenize{Python_head::doc}]{\sphinxcrossref{\DUrole{doc}{Python}}}}.

\sphinxAtStartPar
The constructor of the new atom style will also typically set a few
flags which are defined at the top of \sphinxcode{\sphinxupquote{atom\_vec.h}}.  If these are
unclear, see how other atom styles use them.

\sphinxAtStartPar
The grow\_pointers() method is also required to make a copy of peratom
data pointers, as explained in the code.

\sphinxAtStartPar
There are a number of other optional methods which your atom style can
implement.  These are only needed if you need to do something
out\sphinxhyphen{}of\sphinxhyphen{}the\sphinxhyphen{}ordinary which the default operation of the AtomVec parent
class does not take care of.  The best way to figure out why they are
sometimes useful is to look at how other atom styles use them.
\begin{itemize}
\item {} 
\sphinxAtStartPar
process\_args = use if the atom style has arguments

\item {} 
\sphinxAtStartPar
init = called before each run

\item {} 
\sphinxAtStartPar
force\_clear = called before force computations each timestep

\end{itemize}

\sphinxAtStartPar
A few atom styles define “bonus” data associated with some or all of
their particles, such as {\hyperref[\detokenize{atom_style::doc}]{\sphinxcrossref{\DUrole{doc}{atom\_style ellipsoid or tri}}}}.  These methods work with that data:
\begin{itemize}
\item {} 
\sphinxAtStartPar
copy\_bonus

\item {} 
\sphinxAtStartPar
clear\_bonus

\item {} 
\sphinxAtStartPar
pack\_comm\_bonus

\item {} 
\sphinxAtStartPar
unpack\_comm\_bonus

\item {} 
\sphinxAtStartPar
pack\_border\_bonus

\item {} 
\sphinxAtStartPar
unpack\_border\_bonus

\item {} 
\sphinxAtStartPar
pack\_exchange\_bonus

\item {} 
\sphinxAtStartPar
unpack\_exchange\_bonus

\item {} 
\sphinxAtStartPar
size\_restart\_bonus

\item {} 
\sphinxAtStartPar
pack\_restart\_bonus

\item {} 
\sphinxAtStartPar
unpack\_restart\_bonus

\item {} 
\sphinxAtStartPar
data\_atom\_bonus

\item {} 
\sphinxAtStartPar
memory\_usage\_bonus

\end{itemize}

\sphinxAtStartPar
The {\hyperref[\detokenize{atom_style::doc}]{\sphinxcrossref{\DUrole{doc}{atom\_style body}}}} command can define a particle
geometry with an arbitrary number of values.  This method reads it
from a data file:
\begin{itemize}
\item {} 
\sphinxAtStartPar
data\_body

\end{itemize}

\sphinxAtStartPar
These methods are called before or after operations handled by the
parent AtomVec class.  They allow an atom style to do customized
operations on the per\sphinxhyphen{}atom values.  For example {\hyperref[\detokenize{atom_style::doc}]{\sphinxcrossref{\DUrole{doc}{atom\_style
sphere}}}} reads a diameter and density of each particle
from a data file.  But these need to be converted internally to a
radius and mass.  That operation is done in the data\_atom\_post()
method.
\begin{itemize}
\item {} 
\sphinxAtStartPar
pack\_restart\_pre

\item {} 
\sphinxAtStartPar
pack\_restart\_post

\item {} 
\sphinxAtStartPar
unpack\_restart\_init

\item {} 
\sphinxAtStartPar
create\_atom\_post

\item {} 
\sphinxAtStartPar
data\_atom\_post

\item {} 
\sphinxAtStartPar
pack\_data\_pre

\item {} 
\sphinxAtStartPar
pack\_data\_post

\end{itemize}

\sphinxAtStartPar
These methods enable the {\hyperref[\detokenize{compute_property_atom::doc}]{\sphinxcrossref{\DUrole{doc}{compute property/atom}}}}
command to access per\sphinxhyphen{}atom variables it does not
already define as arguments, so that they can be written to a dump
file or used by other LAMMPS commands.
\begin{itemize}
\item {} 
\sphinxAtStartPar
property\_atom

\item {} 
\sphinxAtStartPar
pack\_property\_atom

\end{itemize}

\sphinxstepscope


\section{Pair styles}
\label{\detokenize{Modify_pair:pair-styles}}\label{\detokenize{Modify_pair::doc}}
\sphinxAtStartPar
Classes that compute pairwise non\sphinxhyphen{}bonded interactions are derived from
the \sphinxcode{\sphinxupquote{Pair}} class.  In LAMMPS, pairwise force calculations include
many\sphinxhyphen{}body potentials such as EAM, Tersoff, or ReaxFF where particles
interact without an explicit bond topology but include interactions
beyond pairwise non\sphinxhyphen{}bonded contributions.  New styles can be created to
add support for additional pair potentials to LAMMPS.  When the
modifications are small, sometimes it is more effective to derive from
an existing pair style class.  This latter approach is also used by
{\hyperref[\detokenize{Speed_packages::doc}]{\sphinxcrossref{\DUrole{doc}{Accelerator packages}}}} where the accelerated style
names differ from their base classes by an appended suffix.

\sphinxAtStartPar
The file \sphinxcode{\sphinxupquote{src/pair\_lj\_cut.cpp}} is an example of a Pair class with a
very simple potential function.  It includes several optional methods to
enable its use with {\hyperref[\detokenize{run_style::doc}]{\sphinxcrossref{\DUrole{doc}{run\_style respa}}}} and {\hyperref[\detokenize{compute_group_group::doc}]{\sphinxcrossref{\DUrole{doc}{compute
group/group}}}}.  {\hyperref[\detokenize{Developer_write_pair::doc}]{\sphinxcrossref{\DUrole{doc}{Writing new pair styles}}}} contains
a detailed discussion of writing new pair styles from scratch, and how
simple and more complex pair styles can be implemented with examples
from existing pair styles.

\sphinxAtStartPar
Here is a brief list of some the class methods in the Pair class that
\sphinxstyleemphasis{must} be or \sphinxstyleemphasis{may} be overridden in a derived class for a new pair style.


\begin{savenotes}\sphinxattablestart
\sphinxthistablewithglobalstyle
\centering
\begin{tabulary}{\linewidth}[t]{TT}
\sphinxtoprule
\sphinxstyletheadfamily 
\sphinxAtStartPar
Required
&\sphinxstyletheadfamily 
\sphinxAtStartPar
“pure” methods that \sphinxstyleemphasis{must} be overridden in a derived class
\\
\sphinxmidrule
\sphinxtableatstartofbodyhook
\sphinxAtStartPar
compute
&
\sphinxAtStartPar
workhorse routine that computes pairwise interactions
\\
\sphinxhline
\sphinxAtStartPar
settings
&
\sphinxAtStartPar
processes the arguments to the pair\_style command
\\
\sphinxhline
\sphinxAtStartPar
coeff
&
\sphinxAtStartPar
set coefficients for one i,j type pair, called from pair\_coeff
\\
\sphinxbottomrule
\end{tabulary}
\sphinxtableafterendhook\par
\sphinxattableend\end{savenotes}


\begin{savenotes}\sphinxattablestart
\sphinxthistablewithglobalstyle
\centering
\begin{tabulary}{\linewidth}[t]{TT}
\sphinxtoprule
\sphinxstyletheadfamily 
\sphinxAtStartPar
Optional
&\sphinxstyletheadfamily 
\sphinxAtStartPar
methods that have a default or dummy implementation
\\
\sphinxmidrule
\sphinxtableatstartofbodyhook
\sphinxAtStartPar
init\_one
&
\sphinxAtStartPar
perform initialization for one i,j type pair
\\
\sphinxhline
\sphinxAtStartPar
init\_style
&
\sphinxAtStartPar
style initialization: request neighbor list(s), error checks
\\
\sphinxhline
\sphinxAtStartPar
init\_list
&
\sphinxAtStartPar
Neighbor class callback function to pass neighbor list to pair style
\\
\sphinxhline
\sphinxAtStartPar
single
&
\sphinxAtStartPar
force/r and energy of a single pairwise interaction between two atoms
\\
\sphinxhline
\sphinxAtStartPar
compute\_inner/middle/outer
&
\sphinxAtStartPar
versions of compute used by rRESPA
\\
\sphinxhline
\sphinxAtStartPar
memory\_usage
&
\sphinxAtStartPar
return estimated amount of memory used by the pair style
\\
\sphinxhline
\sphinxAtStartPar
modify\_params
&
\sphinxAtStartPar
process arguments to pair\_modify command
\\
\sphinxhline
\sphinxAtStartPar
extract
&
\sphinxAtStartPar
provide access to internal scalar or per\sphinxhyphen{}type data like cutoffs
\\
\sphinxhline
\sphinxAtStartPar
extract\_peratom
&
\sphinxAtStartPar
provide access to internal per\sphinxhyphen{}atom data
\\
\sphinxhline
\sphinxAtStartPar
setup
&
\sphinxAtStartPar
initialization at the beginning of a run
\\
\sphinxhline
\sphinxAtStartPar
finish
&
\sphinxAtStartPar
called at the end of a run, e.g. to print
\\
\sphinxhline
\sphinxAtStartPar
write \& read\_restart
&
\sphinxAtStartPar
write/read i,j pair coeffs to restart files
\\
\sphinxhline
\sphinxAtStartPar
write \& read\_restart\_settings
&
\sphinxAtStartPar
write/read global settings to restart files
\\
\sphinxhline
\sphinxAtStartPar
write\_data
&
\sphinxAtStartPar
write Pair Coeffs section to data file
\\
\sphinxhline
\sphinxAtStartPar
write\_data\_all
&
\sphinxAtStartPar
write PairIJ Coeffs section to data file
\\
\sphinxhline
\sphinxAtStartPar
pack \& unpack\_forward\_comm
&
\sphinxAtStartPar
copy data to and from buffer if style uses forward communication
\\
\sphinxhline
\sphinxAtStartPar
pack \& unpack\_reverse\_comm
&
\sphinxAtStartPar
copy data to and from buffer if style uses reverse communication
\\
\sphinxhline
\sphinxAtStartPar
reinit
&
\sphinxAtStartPar
reset all type\sphinxhyphen{}based parameters, called by fix adapt for example
\\
\sphinxhline
\sphinxAtStartPar
reset\_dt
&
\sphinxAtStartPar
called when the time step is changed by timestep or fix reset/dt
\\
\sphinxbottomrule
\end{tabulary}
\sphinxtableafterendhook\par
\sphinxattableend\end{savenotes}

\sphinxAtStartPar
Here is a list of flags or settings that should be set in the
constructor of the derived pair class when they differ from the default
setting.


\begin{savenotes}\sphinxattablestart
\sphinxthistablewithglobalstyle
\centering
\begin{tabulary}{\linewidth}[t]{TTT}
\sphinxtoprule
\sphinxstyletheadfamily 
\sphinxAtStartPar
Name of flag
&\sphinxstyletheadfamily 
\sphinxAtStartPar
Description
&\sphinxstyletheadfamily 
\sphinxAtStartPar
default
\\
\sphinxmidrule
\sphinxtableatstartofbodyhook
\sphinxAtStartPar
single\_enable
&
\sphinxAtStartPar
1 if single() method is implemented, 0 if missing
&
\sphinxAtStartPar
1
\\
\sphinxhline
\sphinxAtStartPar
respa\_enable
&
\sphinxAtStartPar
1 if pair style has compute\_inner/middle/outer()
&
\sphinxAtStartPar
0
\\
\sphinxhline
\sphinxAtStartPar
restartinfo
&
\sphinxAtStartPar
1 if pair style writes its settings to a restart
&
\sphinxAtStartPar
1
\\
\sphinxhline
\sphinxAtStartPar
one\_coeff
&
\sphinxAtStartPar
1 if only a pair\_coeff * * command is allowed
&
\sphinxAtStartPar
0
\\
\sphinxhline
\sphinxAtStartPar
manybody\_flag
&
\sphinxAtStartPar
1 if pair style is a manybody potential
&
\sphinxAtStartPar
0
\\
\sphinxhline
\sphinxAtStartPar
unit\_convert\_flag
&
\sphinxAtStartPar
value != 0 indicates support for unit conversion
&
\sphinxAtStartPar
0
\\
\sphinxhline
\sphinxAtStartPar
no\_virial\_fdotr\_compute
&
\sphinxAtStartPar
1 if pair style does not call virial\_fdotr\_compute()
&
\sphinxAtStartPar
0
\\
\sphinxhline
\sphinxAtStartPar
writedata
&
\sphinxAtStartPar
1 if write\_data() and write\_data\_all() are implemented
&
\sphinxAtStartPar
0
\\
\sphinxhline
\sphinxAtStartPar
comm\_forward
&
\sphinxAtStartPar
size of buffer (in doubles) for forward communication
&
\sphinxAtStartPar
0
\\
\sphinxhline
\sphinxAtStartPar
comm\_reverse
&
\sphinxAtStartPar
size of buffer (in doubles) for reverse communication
&
\sphinxAtStartPar
0
\\
\sphinxhline
\sphinxAtStartPar
ghostneigh
&
\sphinxAtStartPar
1 if cutghost is set and style uses neighbors of ghosts
&
\sphinxAtStartPar
0
\\
\sphinxhline
\sphinxAtStartPar
finitecutflag
&
\sphinxAtStartPar
1 if cutoff depends on diameter of atoms
&
\sphinxAtStartPar
0
\\
\sphinxhline
\sphinxAtStartPar
reinitflag
&
\sphinxAtStartPar
1 if style has reinit() and is compatible with fix adapt
&
\sphinxAtStartPar
0
\\
\sphinxhline
\sphinxAtStartPar
ewaldflag
&
\sphinxAtStartPar
1 if compatible with kspace\_style ewald
&
\sphinxAtStartPar
0
\\
\sphinxhline
\sphinxAtStartPar
pppmflag
&
\sphinxAtStartPar
1 if compatible with kspace\_style pppm
&
\sphinxAtStartPar
0
\\
\sphinxhline
\sphinxAtStartPar
msmflag
&
\sphinxAtStartPar
1 if compatible with kspace\_style msm
&
\sphinxAtStartPar
0
\\
\sphinxhline
\sphinxAtStartPar
dispersionflag
&
\sphinxAtStartPar
1 if compatible with ewald/disp or pppm/disp
&
\sphinxAtStartPar
0
\\
\sphinxhline
\sphinxAtStartPar
tip4pflag
&
\sphinxAtStartPar
1 if compatible with kspace\_style pppm/tip4p
&
\sphinxAtStartPar
0
\\
\sphinxhline
\sphinxAtStartPar
dipoleflag
&
\sphinxAtStartPar
1 if compatible with dipole kspace\_style
&
\sphinxAtStartPar
0
\\
\sphinxhline
\sphinxAtStartPar
spinflag
&
\sphinxAtStartPar
1 if compatible with spin kspace\_style
&
\sphinxAtStartPar
0
\\
\sphinxbottomrule
\end{tabulary}
\sphinxtableafterendhook\par
\sphinxattableend\end{savenotes}

\sphinxstepscope


\section{Bond, angle, dihedral, improper styles}
\label{\detokenize{Modify_bond:bond-angle-dihedral-improper-styles}}\label{\detokenize{Modify_bond::doc}}
\sphinxAtStartPar
Classes that compute molecular interactions are derived from the Bond,
Angle, Dihedral, and Improper classes.  New styles can be created to
add new potentials to LAMMPS.

\sphinxAtStartPar
Bond\_harmonic.cpp is the simplest example of a bond style.  Ditto for
the harmonic forms of the angle, dihedral, and improper style
commands.

\sphinxAtStartPar
Here is a brief description of common methods you define in your
new derived class.  See bond.h, angle.h, dihedral.h, and improper.h
for details and specific additional methods.


\begin{savenotes}\sphinxattablestart
\sphinxthistablewithglobalstyle
\centering
\begin{tabulary}{\linewidth}[t]{TT}
\sphinxtoprule
\sphinxstyletheadfamily 
\sphinxAtStartPar
Required
&\sphinxstyletheadfamily 
\sphinxAtStartPar
“pure” methods that \sphinxstyleemphasis{must} be overridden in a derived class
\\
\sphinxmidrule
\sphinxtableatstartofbodyhook
\sphinxAtStartPar
compute
&
\sphinxAtStartPar
compute the molecular interactions for all listed items
\\
\sphinxhline
\sphinxAtStartPar
coeff
&
\sphinxAtStartPar
set coefficients for one type
\\
\sphinxhline
\sphinxAtStartPar
equilibrium\_distance
&
\sphinxAtStartPar
length of bond, used by SHAKE (bond styles only)
\\
\sphinxhline
\sphinxAtStartPar
equilibrium\_angle
&
\sphinxAtStartPar
opening of angle, used by SHAKE (angle styles only)
\\
\sphinxhline
\sphinxAtStartPar
write \& read\_restart
&
\sphinxAtStartPar
writes/reads coeffs to restart files
\\
\sphinxhline
\sphinxAtStartPar
single
&
\sphinxAtStartPar
force/r (bond styles only) and energy of a single bond or angle
\\
\sphinxbottomrule
\end{tabulary}
\sphinxtableafterendhook\par
\sphinxattableend\end{savenotes}


\begin{savenotes}\sphinxattablestart
\sphinxthistablewithglobalstyle
\centering
\begin{tabulary}{\linewidth}[t]{TT}
\sphinxtoprule
\sphinxstyletheadfamily 
\sphinxAtStartPar
Optional
&\sphinxstyletheadfamily 
\sphinxAtStartPar
methods that have a default or dummy implementation
\\
\sphinxmidrule
\sphinxtableatstartofbodyhook
\sphinxAtStartPar
init
&
\sphinxAtStartPar
check if all coefficients are set, calls init\_style()
\\
\sphinxhline
\sphinxAtStartPar
init\_style
&
\sphinxAtStartPar
check if style specific conditions are met
\\
\sphinxhline
\sphinxAtStartPar
settings
&
\sphinxAtStartPar
apply global settings for all types
\\
\sphinxhline
\sphinxAtStartPar
write \& read\_restart\_settings
&
\sphinxAtStartPar
writes/reads global style settings to restart files
\\
\sphinxhline
\sphinxAtStartPar
write\_data
&
\sphinxAtStartPar
write corresponding Coeffs section(s) in data file
\\
\sphinxhline
\sphinxAtStartPar
memory\_usage
&
\sphinxAtStartPar
tally memory allocated by the style
\\
\sphinxhline
\sphinxAtStartPar
extract
&
\sphinxAtStartPar
provide access to internal data  (bond or angle styles only)
\\
\sphinxhline
\sphinxAtStartPar
reinit
&
\sphinxAtStartPar
reset all type\sphinxhyphen{}based parameters, called by fix adapt (bonds only)
\\
\sphinxhline
\sphinxAtStartPar
pack \& unpack\_forward\_comm
&
\sphinxAtStartPar
copy data to and from buffer in forward communication (bonds only)
\\
\sphinxhline
\sphinxAtStartPar
pack \& unpack\_reverse\_comm
&
\sphinxAtStartPar
copy data to and from buffer in reverse communication (bonds only)
\\
\sphinxbottomrule
\end{tabulary}
\sphinxtableafterendhook\par
\sphinxattableend\end{savenotes}

\sphinxAtStartPar
Here is a list of flags or settings that should be set in the
constructor of the derived class when they differ from the default
setting.


\begin{savenotes}\sphinxattablestart
\sphinxthistablewithglobalstyle
\centering
\begin{tabulary}{\linewidth}[t]{TTT}
\sphinxtoprule
\sphinxstyletheadfamily 
\sphinxAtStartPar
Name of flag
&\sphinxstyletheadfamily 
\sphinxAtStartPar
Description
&\sphinxstyletheadfamily 
\sphinxAtStartPar
default
\\
\sphinxmidrule
\sphinxtableatstartofbodyhook
\sphinxAtStartPar
writedata
&
\sphinxAtStartPar
1 if write\_data() is implemented
&
\sphinxAtStartPar
1
\\
\sphinxhline
\sphinxAtStartPar
single\_extra
&
\sphinxAtStartPar
number of extra single values calculated (bond styles only)
&
\sphinxAtStartPar
0
\\
\sphinxhline
\sphinxAtStartPar
partial\_flag
&
\sphinxAtStartPar
1 if bond type can be set to 0 and deleted (bond styles only)
&
\sphinxAtStartPar
0
\\
\sphinxhline
\sphinxAtStartPar
reinitflag
&
\sphinxAtStartPar
1 if style has reinit() and is compatible with fix adapt
&
\sphinxAtStartPar
1
\\
\sphinxhline
\sphinxAtStartPar
comm\_forward
&
\sphinxAtStartPar
size of buffer (in doubles) for forward communication  (bond styles only)
&
\sphinxAtStartPar
0
\\
\sphinxhline
\sphinxAtStartPar
comm\_reverse
&
\sphinxAtStartPar
size of buffer (in doubles) for reverse communication  (bond styles only)
&
\sphinxAtStartPar
0
\\
\sphinxhline
\sphinxAtStartPar
comm\_reverse\_off
&
\sphinxAtStartPar
size of buffer for reverse communication with newton off (bond styles only)
&
\sphinxAtStartPar
0
\\
\sphinxbottomrule
\end{tabulary}
\sphinxtableafterendhook\par
\sphinxattableend\end{savenotes}

\sphinxstepscope


\section{Compute styles}
\label{\detokenize{Modify_compute:compute-styles}}\label{\detokenize{Modify_compute::doc}}
\sphinxAtStartPar
Classes that compute scalar and vector quantities like temperature
and the pressure tensor, as well as classes that compute per\sphinxhyphen{}atom
quantities like kinetic energy and the centro\sphinxhyphen{}symmetry parameter
are derived from the Compute class.  New styles can be created
to add new calculations to LAMMPS.

\sphinxAtStartPar
Compute\_temp.cpp is a simple example of computing a scalar
temperature.  Compute\_ke\_atom.cpp is a simple example of computing
per\sphinxhyphen{}atom kinetic energy.

\sphinxAtStartPar
Here is a brief description of methods you define in your new derived
class.  See compute.h for details.


\begin{savenotes}\sphinxattablestart
\sphinxthistablewithglobalstyle
\centering
\begin{tabulary}{\linewidth}[t]{TT}
\sphinxtoprule
\sphinxtableatstartofbodyhook
\sphinxAtStartPar
init
&
\sphinxAtStartPar
perform one time setup (required)
\\
\sphinxhline
\sphinxAtStartPar
init\_list
&
\sphinxAtStartPar
neighbor list setup, if needed (optional)
\\
\sphinxhline
\sphinxAtStartPar
compute\_scalar
&
\sphinxAtStartPar
compute a scalar quantity (optional)
\\
\sphinxhline
\sphinxAtStartPar
compute\_vector
&
\sphinxAtStartPar
compute a vector of quantities (optional)
\\
\sphinxhline
\sphinxAtStartPar
compute\_peratom
&
\sphinxAtStartPar
compute one or more quantities per atom (optional)
\\
\sphinxhline
\sphinxAtStartPar
compute\_local
&
\sphinxAtStartPar
compute one or more quantities per processor (optional)
\\
\sphinxhline
\sphinxAtStartPar
pack\_comm
&
\sphinxAtStartPar
pack a buffer with items to communicate (optional)
\\
\sphinxhline
\sphinxAtStartPar
unpack\_comm
&
\sphinxAtStartPar
unpack the buffer (optional)
\\
\sphinxhline
\sphinxAtStartPar
pack\_reverse
&
\sphinxAtStartPar
pack a buffer with items to reverse communicate (optional)
\\
\sphinxhline
\sphinxAtStartPar
unpack\_reverse
&
\sphinxAtStartPar
unpack the buffer (optional)
\\
\sphinxhline
\sphinxAtStartPar
remove\_bias
&
\sphinxAtStartPar
remove velocity bias from one atom (optional)
\\
\sphinxhline
\sphinxAtStartPar
remove\_bias\_all
&
\sphinxAtStartPar
remove velocity bias from all atoms in group (optional)
\\
\sphinxhline
\sphinxAtStartPar
restore\_bias
&
\sphinxAtStartPar
restore velocity bias for one atom after remove\_bias (optional)
\\
\sphinxhline
\sphinxAtStartPar
restore\_bias\_all
&
\sphinxAtStartPar
same as before, but for all atoms in group (optional)
\\
\sphinxhline
\sphinxAtStartPar
pair\_tally\_callback
&
\sphinxAtStartPar
callback function for \sphinxstyleemphasis{tally}\sphinxhyphen{}style computes (optional).
\\
\sphinxhline
\sphinxAtStartPar
memory\_usage
&
\sphinxAtStartPar
tally memory usage (optional)
\\
\sphinxbottomrule
\end{tabulary}
\sphinxtableafterendhook\par
\sphinxattableend\end{savenotes}

\sphinxAtStartPar
Tally\sphinxhyphen{}style computes are a special case, as their computation is done
in two stages: the callback function is registered with the pair style
and then called from the Pair::ev\_tally() function, which is called for
each pair after force and energy has been computed for this pair. Then
the tallied values are retrieved with the standard compute\_scalar or
compute\_vector or compute\_peratom methods. The {\hyperref[\detokenize{compute_tally::doc}]{\sphinxcrossref{\DUrole{doc}{compute styles in the TALLY package}}}}
provide \sphinxstyleemphasis{examples} for utilizing this mechanism.

\sphinxstepscope


\section{Fix styles}
\label{\detokenize{Modify_fix:fix-styles}}\label{\detokenize{Modify_fix::doc}}
\sphinxAtStartPar
In LAMMPS, a “fix” is any operation that is computed during
timestepping that alters some property of the system.  Essentially
everything that happens during a simulation besides force computation,
neighbor list construction, and output, is a “fix”.  This includes
time integration (update of coordinates and velocities), force
constraints or boundary conditions (SHAKE or walls), and diagnostics
(compute a diffusion coefficient).  New styles can be created to add
new options to LAMMPS.

\sphinxAtStartPar
Fix\_setforce.cpp is a simple example of setting forces on atoms to
prescribed values.  There are dozens of fix options already in LAMMPS;
choose one as a template that is similar to what you want to
implement.

\sphinxAtStartPar
Here is a brief description of methods you can define in your new
derived class.  See fix.h for details.


\begin{savenotes}
\sphinxatlongtablestart
\sphinxthistablewithglobalstyle
\makeatletter
  \LTleft \@totalleftmargin plus1fill
  \LTright\dimexpr\columnwidth-\@totalleftmargin-\linewidth\relax plus1fill
\makeatother
\begin{longtable}{ll}
\sphinxtoprule
\endfirsthead

\multicolumn{2}{c}{\sphinxnorowcolor
    \makebox[0pt]{\sphinxtablecontinued{\tablename\ \thetable{} \textendash{} continued from previous page}}%
}\\
\sphinxtoprule
\endhead

\sphinxbottomrule
\multicolumn{2}{r}{\sphinxnorowcolor
    \makebox[0pt][r]{\sphinxtablecontinued{continues on next page}}%
}\\
\endfoot

\endlastfoot
\sphinxtableatstartofbodyhook

\sphinxAtStartPar
setmask
&
\sphinxAtStartPar
determines when the fix is called during the timestep (required)
\\
\sphinxhline
\sphinxAtStartPar
init
&
\sphinxAtStartPar
initialization before a run (optional)
\\
\sphinxhline
\sphinxAtStartPar
init\_list
&
\sphinxAtStartPar
store pointer to neighbor list; called by neighbor list code (optional)
\\
\sphinxhline
\sphinxAtStartPar
setup\_pre\_exchange
&
\sphinxAtStartPar
called before atom exchange in setup (optional)
\\
\sphinxhline
\sphinxAtStartPar
setup\_pre\_force
&
\sphinxAtStartPar
called before force computation in setup (optional)
\\
\sphinxhline
\sphinxAtStartPar
setup
&
\sphinxAtStartPar
called immediately before the first timestep and after forces are computed (optional)
\\
\sphinxhline
\sphinxAtStartPar
min\_setup\_pre\_force
&
\sphinxAtStartPar
like setup\_pre\_force, but for minimizations instead of MD runs (optional)
\\
\sphinxhline
\sphinxAtStartPar
min\_setup
&
\sphinxAtStartPar
like setup, but for minimizations instead of MD runs (optional)
\\
\sphinxhline
\sphinxAtStartPar
initial\_integrate
&
\sphinxAtStartPar
called at very beginning of each timestep (optional)
\\
\sphinxhline
\sphinxAtStartPar
pre\_exchange
&
\sphinxAtStartPar
called before atom exchange on re\sphinxhyphen{}neighboring steps (optional)
\\
\sphinxhline
\sphinxAtStartPar
pre\_neighbor
&
\sphinxAtStartPar
called before neighbor list build (optional)
\\
\sphinxhline
\sphinxAtStartPar
pre\_force
&
\sphinxAtStartPar
called before pair \& molecular forces are computed (optional)
\\
\sphinxhline
\sphinxAtStartPar
post\_force
&
\sphinxAtStartPar
called after pair \& molecular forces are computed and communicated (optional)
\\
\sphinxhline
\sphinxAtStartPar
final\_integrate
&
\sphinxAtStartPar
called at end of each timestep (optional)
\\
\sphinxhline
\sphinxAtStartPar
end\_of\_step
&
\sphinxAtStartPar
called at very end of timestep (optional)
\\
\sphinxhline
\sphinxAtStartPar
write\_restart
&
\sphinxAtStartPar
dumps fix info to restart file (optional)
\\
\sphinxhline
\sphinxAtStartPar
restart
&
\sphinxAtStartPar
uses info from restart file to re\sphinxhyphen{}initialize the fix (optional)
\\
\sphinxhline
\sphinxAtStartPar
grow\_arrays
&
\sphinxAtStartPar
allocate memory for atom\sphinxhyphen{}based arrays used by fix (optional)
\\
\sphinxhline
\sphinxAtStartPar
copy\_arrays
&
\sphinxAtStartPar
copy atom info when an atom migrates to a new processor (optional)
\\
\sphinxhline
\sphinxAtStartPar
pack\_exchange
&
\sphinxAtStartPar
store atom’s data in a buffer (optional)
\\
\sphinxhline
\sphinxAtStartPar
unpack\_exchange
&
\sphinxAtStartPar
retrieve atom’s data from a buffer (optional)
\\
\sphinxhline
\sphinxAtStartPar
pack\_restart
&
\sphinxAtStartPar
store atom’s data for writing to restart file (optional)
\\
\sphinxhline
\sphinxAtStartPar
unpack\_restart
&
\sphinxAtStartPar
retrieve atom’s data from a restart file buffer (optional)
\\
\sphinxhline
\sphinxAtStartPar
size\_restart
&
\sphinxAtStartPar
size of atom’s data (optional)
\\
\sphinxhline
\sphinxAtStartPar
maxsize\_restart
&
\sphinxAtStartPar
max size of atom’s data (optional)
\\
\sphinxhline
\sphinxAtStartPar
setup\_pre\_force\_respa
&
\sphinxAtStartPar
same as setup\_pre\_force, but for rRESPA (optional)
\\
\sphinxhline
\sphinxAtStartPar
initial\_integrate\_respa
&
\sphinxAtStartPar
same as initial\_integrate, but for rRESPA (optional)
\\
\sphinxhline
\sphinxAtStartPar
post\_integrate\_respa
&
\sphinxAtStartPar
called after the first half integration step is done in rRESPA (optional)
\\
\sphinxhline
\sphinxAtStartPar
pre\_force\_respa
&
\sphinxAtStartPar
same as pre\_force, but for rRESPA (optional)
\\
\sphinxhline
\sphinxAtStartPar
post\_force\_respa
&
\sphinxAtStartPar
same as post\_force, but for rRESPA (optional)
\\
\sphinxhline
\sphinxAtStartPar
final\_integrate\_respa
&
\sphinxAtStartPar
same as final\_integrate, but for rRESPA (optional)
\\
\sphinxhline
\sphinxAtStartPar
min\_pre\_force
&
\sphinxAtStartPar
called after pair \& molecular forces are computed in minimizer (optional)
\\
\sphinxhline
\sphinxAtStartPar
min\_post\_force
&
\sphinxAtStartPar
called after pair \& molecular forces are computed and communicated in minimizer (optional)
\\
\sphinxhline
\sphinxAtStartPar
min\_store
&
\sphinxAtStartPar
store extra data for linesearch based minimization on a LIFO stack (optional)
\\
\sphinxhline
\sphinxAtStartPar
min\_pushstore
&
\sphinxAtStartPar
push the minimization LIFO stack one element down (optional)
\\
\sphinxhline
\sphinxAtStartPar
min\_popstore
&
\sphinxAtStartPar
pop the minimization LIFO stack one element up (optional)
\\
\sphinxhline
\sphinxAtStartPar
min\_clearstore
&
\sphinxAtStartPar
clear minimization LIFO stack (optional)
\\
\sphinxhline
\sphinxAtStartPar
min\_step
&
\sphinxAtStartPar
reset or move forward on line search minimization (optional)
\\
\sphinxhline
\sphinxAtStartPar
min\_dof
&
\sphinxAtStartPar
report number of degrees of freedom \sphinxstyleemphasis{added} by this fix in minimization (optional)
\\
\sphinxhline
\sphinxAtStartPar
max\_alpha
&
\sphinxAtStartPar
report maximum allowed step size during linesearch minimization (optional)
\\
\sphinxhline
\sphinxAtStartPar
pack\_comm
&
\sphinxAtStartPar
pack a buffer to communicate a per\sphinxhyphen{}atom quantity (optional)
\\
\sphinxhline
\sphinxAtStartPar
unpack\_comm
&
\sphinxAtStartPar
unpack a buffer to communicate a per\sphinxhyphen{}atom quantity (optional)
\\
\sphinxhline
\sphinxAtStartPar
pack\_reverse\_comm
&
\sphinxAtStartPar
pack a buffer to reverse communicate a per\sphinxhyphen{}atom quantity (optional)
\\
\sphinxhline
\sphinxAtStartPar
unpack\_reverse\_comm
&
\sphinxAtStartPar
unpack a buffer to reverse communicate a per\sphinxhyphen{}atom quantity (optional)
\\
\sphinxhline
\sphinxAtStartPar
dof
&
\sphinxAtStartPar
report number of degrees of freedom \sphinxstyleemphasis{removed} by this fix during MD (optional)
\\
\sphinxhline
\sphinxAtStartPar
compute\_scalar
&
\sphinxAtStartPar
return a global scalar property that the fix computes (optional)
\\
\sphinxhline
\sphinxAtStartPar
compute\_vector
&
\sphinxAtStartPar
return a component of a vector property that the fix computes (optional)
\\
\sphinxhline
\sphinxAtStartPar
compute\_array
&
\sphinxAtStartPar
return a component of an array property that the fix computes (optional)
\\
\sphinxhline
\sphinxAtStartPar
deform
&
\sphinxAtStartPar
called when the box size is changed (optional)
\\
\sphinxhline
\sphinxAtStartPar
reset\_target
&
\sphinxAtStartPar
called when a change of the target temperature is requested during a run (optional)
\\
\sphinxhline
\sphinxAtStartPar
reset\_dt
&
\sphinxAtStartPar
is called when a change of the time step is requested during a run (optional)
\\
\sphinxhline
\sphinxAtStartPar
modify\_param
&
\sphinxAtStartPar
called when a fix\_modify request is executed (optional)
\\
\sphinxhline
\sphinxAtStartPar
memory\_usage
&
\sphinxAtStartPar
report memory used by fix (optional)
\\
\sphinxhline
\sphinxAtStartPar
thermo
&
\sphinxAtStartPar
compute quantities for thermodynamic output (optional)
\\
\sphinxbottomrule
\end{longtable}
\sphinxtableafterendhook
\sphinxatlongtableend
\end{savenotes}

\sphinxAtStartPar
Typically, only a small fraction of these methods are defined for a
particular fix.  Setmask is mandatory, as it determines when the fix
will be invoked during the timestep.  Fixes that perform time
integration (\sphinxstyleemphasis{nve}, \sphinxstyleemphasis{nvt}, \sphinxstyleemphasis{npt}) implement initial\_integrate() and
final\_integrate() to perform velocity Verlet updates.  Fixes that
constrain forces implement post\_force().

\sphinxAtStartPar
Fixes that perform diagnostics typically implement end\_of\_step().  For
an end\_of\_step fix, one of your fix arguments must be the variable
“nevery” which is used to determine when to call the fix and you must
set this variable in the constructor of your fix.  By convention, this
is the first argument the fix defines (after the ID, group\sphinxhyphen{}ID, style).

\sphinxAtStartPar
If the fix needs to store information for each atom that persists from
timestep to timestep, it can manage that memory and migrate the info
with the atoms as they move from processors to processor by
implementing the grow\_arrays, copy\_arrays, pack\_exchange, and
unpack\_exchange methods.  Similarly, the pack\_restart and
unpack\_restart methods can be implemented to store information about
the fix in restart files.  If you wish an integrator or force
constraint fix to work with rRESPA (see the {\hyperref[\detokenize{run_style::doc}]{\sphinxcrossref{\DUrole{doc}{run\_style}}}}
command), the initial\_integrate, post\_force\_integrate, and
final\_integrate\_respa methods can be implemented.  The thermo method
enables a fix to contribute values to thermodynamic output, as printed
quantities and/or to be summed to the potential energy of the system.

\sphinxstepscope


\section{Input script command style}
\label{\detokenize{Modify_command:input-script-command-style}}\label{\detokenize{Modify_command::doc}}
\sphinxAtStartPar
New commands can be added to LAMMPS input scripts by adding new classes
that are derived from the Command class and thus must have a “command”
method.  For example, the create\_atoms, read\_data, velocity, and run
commands are all implemented in this fashion.  When such a command is
encountered in the LAMMPS input script, LAMMPS simply creates a class
instance with the corresponding name, invokes the “command” method of
the class, and passes it the arguments from the input script.  The
command method can perform whatever operations it wishes on LAMMPS data
structures.

\sphinxAtStartPar
The single method your new class must define is as follows:


\begin{savenotes}\sphinxattablestart
\sphinxthistablewithglobalstyle
\centering
\begin{tabulary}{\linewidth}[t]{TT}
\sphinxtoprule
\sphinxtableatstartofbodyhook
\sphinxAtStartPar
command
&
\sphinxAtStartPar
operations performed by the new command
\\
\sphinxbottomrule
\end{tabulary}
\sphinxtableafterendhook\par
\sphinxattableend\end{savenotes}

\sphinxAtStartPar
Of course, the new class can define other methods and variables as
needed.

\sphinxstepscope


\section{Dump styles}
\label{\detokenize{Modify_dump:dump-styles}}\label{\detokenize{Modify_dump::doc}}
\sphinxAtStartPar
Classes that dump per\sphinxhyphen{}atom info to files are derived from the Dump
class.  To dump new quantities or in a new format, a new derived dump
class can be added, but it is typically simpler to modify the
DumpCustom class contained in the dump\_custom.cpp file.

\sphinxAtStartPar
Dump\_atom.cpp is a simple example of a derived dump class.

\sphinxAtStartPar
Here is a brief description of methods you define in your new derived
class.  See dump.h for details.


\begin{savenotes}\sphinxattablestart
\sphinxthistablewithglobalstyle
\centering
\begin{tabulary}{\linewidth}[t]{TT}
\sphinxtoprule
\sphinxtableatstartofbodyhook
\sphinxAtStartPar
write\_header
&
\sphinxAtStartPar
write the header section of a snapshot of atoms
\\
\sphinxhline
\sphinxAtStartPar
count
&
\sphinxAtStartPar
count the number of lines a processor will output
\\
\sphinxhline
\sphinxAtStartPar
pack
&
\sphinxAtStartPar
pack a proc’s output data into a buffer
\\
\sphinxhline
\sphinxAtStartPar
write\_data
&
\sphinxAtStartPar
write a proc’s data to a file
\\
\sphinxbottomrule
\end{tabulary}
\sphinxtableafterendhook\par
\sphinxattableend\end{savenotes}

\sphinxAtStartPar
See the {\hyperref[\detokenize{dump::doc}]{\sphinxcrossref{\DUrole{doc}{dump}}}} command and its \sphinxstyleemphasis{custom} style for a list of
keywords for atom information that can already be dumped by
DumpCustom.  It includes options to dump per\sphinxhyphen{}atom info from Compute
classes, so adding a new derived Compute class is one way to calculate
new quantities to dump.

\sphinxAtStartPar
Note that new keywords for atom properties are not typically
added to the {\hyperref[\detokenize{dump::doc}]{\sphinxcrossref{\DUrole{doc}{dump custom}}}} command.  Instead they are added
to the {\hyperref[\detokenize{compute_property_atom::doc}]{\sphinxcrossref{\DUrole{doc}{compute property/atom}}}} command.

\sphinxstepscope


\section{Kspace styles}
\label{\detokenize{Modify_kspace:kspace-styles}}\label{\detokenize{Modify_kspace::doc}}
\sphinxAtStartPar
Classes that compute long\sphinxhyphen{}range Coulombic interactions via K\sphinxhyphen{}space
representations (Ewald, PPPM) are derived from the KSpace class.  New
styles can be created to add new K\sphinxhyphen{}space options to LAMMPS.

\sphinxAtStartPar
Ewald.cpp is an example of computing K\sphinxhyphen{}space interactions.

\sphinxAtStartPar
Here is a brief description of methods you define in your new derived
class.  See kspace.h for details.


\begin{savenotes}\sphinxattablestart
\sphinxthistablewithglobalstyle
\centering
\begin{tabulary}{\linewidth}[t]{TT}
\sphinxtoprule
\sphinxtableatstartofbodyhook
\sphinxAtStartPar
init
&
\sphinxAtStartPar
initialize the calculation before a run
\\
\sphinxhline
\sphinxAtStartPar
setup
&
\sphinxAtStartPar
computation before the first timestep of a run
\\
\sphinxhline
\sphinxAtStartPar
compute
&
\sphinxAtStartPar
every\sphinxhyphen{}timestep computation
\\
\sphinxhline
\sphinxAtStartPar
memory\_usage
&
\sphinxAtStartPar
tally of memory usage
\\
\sphinxbottomrule
\end{tabulary}
\sphinxtableafterendhook\par
\sphinxattableend\end{savenotes}

\sphinxstepscope


\section{Minimization styles}
\label{\detokenize{Modify_min:minimization-styles}}\label{\detokenize{Modify_min::doc}}
\sphinxAtStartPar
Classes that perform energy minimization derived from the Min class.
New styles can be created to add new minimization algorithms to
LAMMPS.

\sphinxAtStartPar
Min\_cg.cpp is an example of conjugate gradient minimization.

\sphinxAtStartPar
Here is a brief description of methods you define in your new derived
class.  See min.h for details.


\begin{savenotes}\sphinxattablestart
\sphinxthistablewithglobalstyle
\centering
\begin{tabulary}{\linewidth}[t]{TT}
\sphinxtoprule
\sphinxtableatstartofbodyhook
\sphinxAtStartPar
init
&
\sphinxAtStartPar
initialize the minimization before a run
\\
\sphinxhline
\sphinxAtStartPar
run
&
\sphinxAtStartPar
perform the minimization
\\
\sphinxhline
\sphinxAtStartPar
memory\_usage
&
\sphinxAtStartPar
tally of memory usage
\\
\sphinxbottomrule
\end{tabulary}
\sphinxtableafterendhook\par
\sphinxattableend\end{savenotes}

\sphinxstepscope


\section{Region styles}
\label{\detokenize{Modify_region:region-styles}}\label{\detokenize{Modify_region::doc}}
\sphinxAtStartPar
Classes that define geometric regions are derived from the Region
class.  Regions are used elsewhere in LAMMPS to group atoms, delete
atoms to create a void, insert atoms in a specified region, etc.  New
styles can be created to add new region shapes to LAMMPS.

\sphinxAtStartPar
Region\_sphere.cpp is an example of a spherical region.

\sphinxAtStartPar
Here is a brief description of methods you define in your new derived
class.  See region.h for details.


\begin{savenotes}\sphinxattablestart
\sphinxthistablewithglobalstyle
\centering
\begin{tabulary}{\linewidth}[t]{TT}
\sphinxtoprule
\sphinxtableatstartofbodyhook
\sphinxAtStartPar
inside
&
\sphinxAtStartPar
determine whether a point is in the region
\\
\sphinxhline
\sphinxAtStartPar
surface\_interior
&
\sphinxAtStartPar
determine if a point is within a cutoff distance inside of surface
\\
\sphinxhline
\sphinxAtStartPar
surface\_exterior
&
\sphinxAtStartPar
determine if a point is within a cutoff distance outside of surface
\\
\sphinxhline
\sphinxAtStartPar
shape\_update
&
\sphinxAtStartPar
change region shape if set by time\sphinxhyphen{}dependent variable
\\
\sphinxbottomrule
\end{tabulary}
\sphinxtableafterendhook\par
\sphinxattableend\end{savenotes}

\sphinxstepscope


\section{Body styles}
\label{\detokenize{Modify_body:body-styles}}\label{\detokenize{Modify_body::doc}}
\sphinxAtStartPar
Classes that define body particles are derived from the Body class.
Body particles can represent complex entities, such as surface meshes
of discrete points, collections of sub\sphinxhyphen{}particles, deformable objects,
etc.

\sphinxAtStartPar
See the {\hyperref[\detokenize{Howto_body::doc}]{\sphinxcrossref{\DUrole{doc}{Howto body}}}} page for an overview of using
body particles and the various body styles LAMMPS supports.  New
styles can be created to add new kinds of body particles to LAMMPS.

\sphinxAtStartPar
Body\_nparticle.cpp is an example of a body particle that is treated as
a rigid body containing N sub\sphinxhyphen{}particles.

\sphinxAtStartPar
Here is a brief description of methods you define in your new derived
class.  See body.h for details.


\begin{savenotes}\sphinxattablestart
\sphinxthistablewithglobalstyle
\centering
\begin{tabulary}{\linewidth}[t]{TT}
\sphinxtoprule
\sphinxtableatstartofbodyhook
\sphinxAtStartPar
data\_body
&
\sphinxAtStartPar
process a line from the Bodies section of a data file
\\
\sphinxhline
\sphinxAtStartPar
noutrow
&
\sphinxAtStartPar
number of sub\sphinxhyphen{}particles output is generated for
\\
\sphinxhline
\sphinxAtStartPar
noutcol
&
\sphinxAtStartPar
number of values per\sphinxhyphen{}sub\sphinxhyphen{}particle output is generated for
\\
\sphinxhline
\sphinxAtStartPar
output
&
\sphinxAtStartPar
output values for the Mth sub\sphinxhyphen{}particle
\\
\sphinxhline
\sphinxAtStartPar
pack\_comm\_body
&
\sphinxAtStartPar
body attributes to communicate every timestep
\\
\sphinxhline
\sphinxAtStartPar
unpack\_comm\_body
&
\sphinxAtStartPar
unpacking of those attributes
\\
\sphinxhline
\sphinxAtStartPar
pack\_border\_body
&
\sphinxAtStartPar
body attributes to communicate when reneighboring is done
\\
\sphinxhline
\sphinxAtStartPar
unpack\_border\_body
&
\sphinxAtStartPar
unpacking of those attributes
\\
\sphinxbottomrule
\end{tabulary}
\sphinxtableafterendhook\par
\sphinxattableend\end{savenotes}

\sphinxstepscope


\section{Granular Sub\sphinxhyphen{}Model styles}
\label{\detokenize{Modify_gran_sub_mod:granular-sub-model-styles}}\label{\detokenize{Modify_gran_sub_mod::doc}}
\sphinxAtStartPar
In granular models, particles are spheres with a finite radius and rotational
degrees of freedom as further described in the
{\hyperref[\detokenize{Howto_granular::doc}]{\sphinxcrossref{\DUrole{doc}{Howto granular page}}}}. Interactions between pair of
particles or particles and walls may therefore depend on many different modes
of motion as described in {\hyperref[\detokenize{pair_granular::doc}]{\sphinxcrossref{\DUrole{doc}{pair granular}}}} and
{\hyperref[\detokenize{fix_wall_gran::doc}]{\sphinxcrossref{\DUrole{doc}{fix wall/gran}}}}. In both cases, the exchange of forces,
torques, and heat flow between two types of bodies is defined using a
GranularModel class. The GranularModel class organizes the details of an
interaction using a series of granular sub\sphinxhyphen{}models each of which describe a
particular interaction mode (e.g. normal forces or rolling friction). From a
parent GranSubMod class, several types of sub\sphinxhyphen{}model classes are derived:
\begin{itemize}
\item {} 
\sphinxAtStartPar
GranSubModNormal: normal force sub\sphinxhyphen{}model

\item {} 
\sphinxAtStartPar
GranSubModDamping: normal damping sub\sphinxhyphen{}model

\item {} 
\sphinxAtStartPar
GranSubModTangential: tangential forces and sliding friction sub\sphinxhyphen{}model

\item {} 
\sphinxAtStartPar
GranSubModRolling: rolling friction sub\sphinxhyphen{}model

\item {} 
\sphinxAtStartPar
GranSubModTwisting: twisting friction sub\sphinxhyphen{}model

\item {} 
\sphinxAtStartPar
GranSubModHeat: heat conduction sub\sphinxhyphen{}model

\end{itemize}

\sphinxAtStartPar
For each type of sub\sphinxhyphen{}model, more classes are further derived, each describing a
specific implementation. For instance, from the GranSubModNormal class the
GranSubModNormalHooke, GranSubModNormalHertz, and GranSubModNormalJKR classes
are derived which calculate Hookean, Hertzian, or JKR normal forces,
respectively.  This modular structure simplifies the addition of new granular
contact models as one only needs to create a new GranSubMod class without
having to modify the more complex PairGranular, FixGranWall, and GranularModel
classes. Most GranSubMod methods are also already defined by the parent classes
so new contact models typically only require edits to a few relevant methods
(e.g. methods that define coefficients and calculate forces).

\sphinxAtStartPar
Each GranSubMod class has a pointer to both the LAMMPS class and the GranularModel
class which owns it, \sphinxcode{\sphinxupquote{lmp}} and \sphinxcode{\sphinxupquote{gm}}, respectively. The GranularModel class
includes several public variables that describe the geometry/dynamics of the
contact such as


\begin{savenotes}\sphinxattablestart
\sphinxthistablewithglobalstyle
\centering
\begin{tabular}[t]{\X{25}{100}\X{75}{100}}
\sphinxtoprule
\sphinxtableatstartofbodyhook
\sphinxAtStartPar
\sphinxcode{\sphinxupquote{xi}} and \sphinxcode{\sphinxupquote{xj}}
&
\sphinxAtStartPar
Positions of the two contacting bodies
\\
\sphinxhline
\sphinxAtStartPar
\sphinxcode{\sphinxupquote{vi}} and \sphinxcode{\sphinxupquote{vj}}
&
\sphinxAtStartPar
Velocities of the two contacting bodies
\\
\sphinxhline
\sphinxAtStartPar
\sphinxcode{\sphinxupquote{omegai}} and \sphinxcode{\sphinxupquote{omegaj}}
&
\sphinxAtStartPar
Angular velocities of the two contacting bodies
\\
\sphinxhline
\sphinxAtStartPar
\sphinxcode{\sphinxupquote{dx}} and \sphinxcode{\sphinxupquote{nx}}
&
\sphinxAtStartPar
The displacement and normalized displacement vectors
\\
\sphinxhline
\sphinxAtStartPar
\sphinxcode{\sphinxupquote{r}}, \sphinxcode{\sphinxupquote{rsq}}, and \sphinxcode{\sphinxupquote{rinv}}
&
\sphinxAtStartPar
The distance, distance squared, and inverse distance
\\
\sphinxhline
\sphinxAtStartPar
\sphinxcode{\sphinxupquote{radsum}}
&
\sphinxAtStartPar
The sum of particle radii
\\
\sphinxhline
\sphinxAtStartPar
\sphinxcode{\sphinxupquote{vr}}, \sphinxcode{\sphinxupquote{vn}}, and \sphinxcode{\sphinxupquote{vt}}
&
\sphinxAtStartPar
The relative velocity vector and its normal and tangential components
\\
\sphinxhline
\sphinxAtStartPar
\sphinxcode{\sphinxupquote{wr}}
&
\sphinxAtStartPar
The relative rotational velocity
\\
\sphinxbottomrule
\end{tabular}
\sphinxtableafterendhook\par
\sphinxattableend\end{savenotes}

\sphinxAtStartPar
These quantities, among others, are calculated in the \sphinxcode{\sphinxupquote{GranularModel\sphinxhyphen{}\textgreater{}check\_contact()}}
and \sphinxcode{\sphinxupquote{GranularModel\sphinxhyphen{}\textgreater{}calculate\_forces()}} methods which can be referred to for more
details.

\sphinxAtStartPar
To create a new GranSubMod class, it is recommended that one first looks at similar
GranSubMod classes. All GranSubMod classes share several general methods which may
need to be defined


\begin{savenotes}\sphinxattablestart
\sphinxthistablewithglobalstyle
\centering
\begin{tabular}[t]{\X{25}{100}\X{75}{100}}
\sphinxtoprule
\sphinxtableatstartofbodyhook
\sphinxAtStartPar
\sphinxcode{\sphinxupquote{mix\_coeff()}}
&
\sphinxAtStartPar
Optional method to define how coefficients are mixed for different atom types. By default, coefficients are mixed using a geometric mean.
\\
\sphinxhline
\sphinxAtStartPar
\sphinxcode{\sphinxupquote{coeffs\_to\_local()}}
&
\sphinxAtStartPar
Parses coefficients to define local variables. Run once at model construction.
\\
\sphinxhline
\sphinxAtStartPar
\sphinxcode{\sphinxupquote{init()}}
&
\sphinxAtStartPar
Optional method to define local variables after other GranSubMod types were created. For instance, this method may be used by a tangential model that derives parameters from the normal model.
\\
\sphinxbottomrule
\end{tabular}
\sphinxtableafterendhook\par
\sphinxattableend\end{savenotes}

\sphinxAtStartPar
The Normal, Damping, Tangential, Twisting, and Rolling sub\sphinxhyphen{}models also have a
\sphinxcode{\sphinxupquote{calculate\_forces()}} method which calculate the respective forces/torques.
Correspondingly, the Heat sub\sphinxhyphen{}model has a \sphinxcode{\sphinxupquote{calculate\_heat()}} method. Lastly,
the Normal sub\sphinxhyphen{}model has a few extra optional methods:


\begin{savenotes}\sphinxattablestart
\sphinxthistablewithglobalstyle
\centering
\begin{tabular}[t]{\X{25}{100}\X{75}{100}}
\sphinxtoprule
\sphinxtableatstartofbodyhook
\sphinxAtStartPar
\sphinxcode{\sphinxupquote{touch()}}
&
\sphinxAtStartPar
Tests whether particles are in contact. By default, when particles overlap.
\\
\sphinxhline
\sphinxAtStartPar
\sphinxcode{\sphinxupquote{pulloff\_distance()}}
&
\sphinxAtStartPar
Returns the distance at which particles stop interacting. By default, when particles no longer overlap.
\\
\sphinxhline
\sphinxAtStartPar
\sphinxcode{\sphinxupquote{calculate\_radius()}}
&
\sphinxAtStartPar
Returns the radius of the contact. By default, the radius of the geometric cross section.
\\
\sphinxhline
\sphinxAtStartPar
\sphinxcode{\sphinxupquote{set\_fncrit()}}
&
\sphinxAtStartPar
Defines the critical force to break the contact used by some tangential, rolling, and twisting sub\sphinxhyphen{}models. By default, the current total normal force including damping.
\\
\sphinxbottomrule
\end{tabular}
\sphinxtableafterendhook\par
\sphinxattableend\end{savenotes}

\sphinxAtStartPar
As an example, say one wanted to create a new normal force option that consisted
of a Hookean force with a piecewise stiffness. This could be done by adding a new
set of files \sphinxcode{\sphinxupquote{gran\_sub\_mod\_custom.h}}:

\begin{sphinxVerbatim}[commandchars=\\\{\}]
\PYG{c+cp}{\PYGZsh{}}\PYG{c+cp}{ifdef GranSubMod\PYGZus{}CLASS}
\PYG{c+c1}{// clang\PYGZhy{}format off}
\PYG{n}{GranSubModStyle}\PYG{p}{(}\PYG{n}{hooke}\PYG{o}{/}\PYG{n}{piecewise}\PYG{p}{,}\PYG{n}{GranSubModNormalHookePiecewise}\PYG{p}{,}\PYG{n}{NORMAL}\PYG{p}{)}\PYG{p}{;}
\PYG{c+c1}{// clang\PYGZhy{}format on}
\PYG{c+cp}{\PYGZsh{}}\PYG{c+cp}{else}

\PYG{c+cp}{\PYGZsh{}}\PYG{c+cp}{ifndef GRAN\PYGZus{}SUB\PYGZus{}MOD\PYGZus{}CUSTOM\PYGZus{}H\PYGZus{}}
\PYG{c+cp}{\PYGZsh{}}\PYG{c+cp}{define GRAN\PYGZus{}SUB\PYGZus{}MOD\PYGZus{}CUSTOM\PYGZus{}H\PYGZus{}}

\PYG{c+cp}{\PYGZsh{}}\PYG{c+cp}{include}\PYG{+w}{ }\PYG{c+cpf}{\PYGZdq{}gran\PYGZus{}sub\PYGZus{}mod.h\PYGZdq{}}
\PYG{c+cp}{\PYGZsh{}}\PYG{c+cp}{include}\PYG{+w}{ }\PYG{c+cpf}{\PYGZdq{}gran\PYGZus{}sub\PYGZus{}mod\PYGZus{}normal.h\PYGZdq{}}

\PYG{k}{namespace}\PYG{+w}{ }\PYG{n+nn}{LAMMPS\PYGZus{}NS}\PYG{+w}{ }\PYG{p}{\PYGZob{}}
\PYG{k}{namespace}\PYG{+w}{ }\PYG{n+nn}{Granular\PYGZus{}NS}\PYG{+w}{ }\PYG{p}{\PYGZob{}}
\PYG{+w}{  }\PYG{k}{class}\PYG{+w}{ }\PYG{n+nc}{GranSubModNormalHookePiecewise}\PYG{+w}{ }\PYG{o}{:}\PYG{+w}{ }\PYG{k}{public}\PYG{+w}{ }\PYG{n}{GranSubModNormal}\PYG{+w}{ }\PYG{p}{\PYGZob{}}
\PYG{+w}{   }\PYG{k}{public}\PYG{o}{:}
\PYG{+w}{    }\PYG{n}{GranSubModNormalHookePiecewise}\PYG{p}{(}\PYG{k}{class}\PYG{+w}{ }\PYG{n+nc}{GranularModel}\PYG{+w}{ }\PYG{o}{*}\PYG{p}{,}\PYG{+w}{ }\PYG{k}{class}\PYG{+w}{ }\PYG{n+nc}{LAMMPS}\PYG{+w}{ }\PYG{o}{*}\PYG{p}{)}\PYG{p}{;}
\PYG{+w}{    }\PYG{k+kt}{void}\PYG{+w}{ }\PYG{n+nf}{coeffs\PYGZus{}to\PYGZus{}local}\PYG{p}{(}\PYG{p}{)}\PYG{+w}{ }\PYG{k}{override}\PYG{p}{;}
\PYG{+w}{    }\PYG{k+kt}{double}\PYG{+w}{ }\PYG{n+nf}{calculate\PYGZus{}forces}\PYG{p}{(}\PYG{p}{)}\PYG{+w}{ }\PYG{k}{override}\PYG{p}{;}
\PYG{+w}{   }\PYG{k}{protected}\PYG{o}{:}
\PYG{+w}{    }\PYG{k+kt}{double}\PYG{+w}{ }\PYG{n}{k1}\PYG{p}{,}\PYG{+w}{ }\PYG{n}{k2}\PYG{p}{,}\PYG{+w}{ }\PYG{n}{delta\PYGZus{}switch}\PYG{p}{;}
\PYG{+w}{  }\PYG{p}{\PYGZcb{}}\PYG{p}{;}
\PYG{p}{\PYGZcb{}}\PYG{+w}{    }\PYG{c+c1}{// namespace Granular\PYGZus{}NS}
\PYG{p}{\PYGZcb{}}\PYG{+w}{    }\PYG{c+c1}{// namespace LAMMPS\PYGZus{}NS}

\PYG{c+cp}{\PYGZsh{}}\PYG{c+cp}{endif }\PYG{c+cm}{/*GRAN\PYGZus{}SUB\PYGZus{}MOD\PYGZus{}CUSTOM\PYGZus{}H\PYGZus{} */}
\PYG{c+cp}{\PYGZsh{}}\PYG{c+cp}{endif }\PYG{c+cm}{/*GRAN\PYGZus{}SUB\PYGZus{}MOD\PYGZus{}CLASS\PYGZus{}H\PYGZus{} */}
\end{sphinxVerbatim}

\sphinxAtStartPar
and \sphinxcode{\sphinxupquote{gran\_sub\_mod\_custom.cpp}}

\begin{sphinxVerbatim}[commandchars=\\\{\}]
\PYG{c+cp}{\PYGZsh{}}\PYG{c+cp}{include}\PYG{+w}{ }\PYG{c+cpf}{\PYGZdq{}gran\PYGZus{}sub\PYGZus{}mod\PYGZus{}custom.h\PYGZdq{}}
\PYG{c+cp}{\PYGZsh{}}\PYG{c+cp}{include}\PYG{+w}{ }\PYG{c+cpf}{\PYGZdq{}gran\PYGZus{}sub\PYGZus{}mod\PYGZus{}normal.h\PYGZdq{}}
\PYG{c+cp}{\PYGZsh{}}\PYG{c+cp}{include}\PYG{+w}{ }\PYG{c+cpf}{\PYGZdq{}granular\PYGZus{}model.h\PYGZdq{}}

\PYG{k}{using}\PYG{+w}{ }\PYG{k}{namespace}\PYG{+w}{ }\PYG{n+nn}{LAMMPS\PYGZus{}NS}\PYG{p}{;}
\PYG{k}{using}\PYG{+w}{ }\PYG{k}{namespace}\PYG{+w}{ }\PYG{n+nn}{Granular\PYGZus{}NS}\PYG{p}{;}

\PYG{n}{GranSubModNormalHookePiecewise}\PYG{o}{:}\PYG{o}{:}\PYG{n}{GranSubModNormalHookePiecewise}\PYG{p}{(}\PYG{n}{GranularModel}\PYG{+w}{ }\PYG{o}{*}\PYG{n}{gm}\PYG{p}{,}\PYG{+w}{ }\PYG{n}{LAMMPS}\PYG{+w}{ }\PYG{o}{*}\PYG{n}{lmp}\PYG{p}{)}\PYG{+w}{ }\PYG{o}{:}
\PYG{+w}{    }\PYG{n}{GranSubModNormal}\PYG{p}{(}\PYG{n}{gm}\PYG{p}{,}\PYG{+w}{ }\PYG{n}{lmp}\PYG{p}{)}
\PYG{p}{\PYGZob{}}
\PYG{+w}{  }\PYG{n}{num\PYGZus{}coeffs}\PYG{+w}{ }\PYG{o}{=}\PYG{+w}{ }\PYG{l+m+mi}{4}\PYG{p}{;}
\PYG{p}{\PYGZcb{}}

\PYG{c+cm}{/* \PYGZhy{}\PYGZhy{}\PYGZhy{}\PYGZhy{}\PYGZhy{}\PYGZhy{}\PYGZhy{}\PYGZhy{}\PYGZhy{}\PYGZhy{}\PYGZhy{}\PYGZhy{}\PYGZhy{}\PYGZhy{}\PYGZhy{}\PYGZhy{}\PYGZhy{}\PYGZhy{}\PYGZhy{}\PYGZhy{}\PYGZhy{}\PYGZhy{}\PYGZhy{}\PYGZhy{}\PYGZhy{}\PYGZhy{}\PYGZhy{}\PYGZhy{}\PYGZhy{}\PYGZhy{}\PYGZhy{}\PYGZhy{}\PYGZhy{}\PYGZhy{}\PYGZhy{}\PYGZhy{}\PYGZhy{}\PYGZhy{}\PYGZhy{}\PYGZhy{}\PYGZhy{}\PYGZhy{}\PYGZhy{}\PYGZhy{}\PYGZhy{}\PYGZhy{}\PYGZhy{}\PYGZhy{}\PYGZhy{}\PYGZhy{}\PYGZhy{}\PYGZhy{}\PYGZhy{}\PYGZhy{}\PYGZhy{}\PYGZhy{}\PYGZhy{}\PYGZhy{}\PYGZhy{}\PYGZhy{}\PYGZhy{}\PYGZhy{}\PYGZhy{}\PYGZhy{}\PYGZhy{}\PYGZhy{}\PYGZhy{}\PYGZhy{}\PYGZhy{}\PYGZhy{} */}

\PYG{k+kt}{void}\PYG{+w}{ }\PYG{n}{GranSubModNormalHookePiecewise}\PYG{o}{:}\PYG{o}{:}\PYG{n}{coeffs\PYGZus{}to\PYGZus{}local}\PYG{p}{(}\PYG{p}{)}
\PYG{p}{\PYGZob{}}
\PYG{+w}{  }\PYG{n}{k1}\PYG{+w}{ }\PYG{o}{=}\PYG{+w}{ }\PYG{n}{coeffs}\PYG{p}{[}\PYG{l+m+mi}{0}\PYG{p}{]}\PYG{p}{;}
\PYG{+w}{  }\PYG{n}{k2}\PYG{+w}{ }\PYG{o}{=}\PYG{+w}{ }\PYG{n}{coeffs}\PYG{p}{[}\PYG{l+m+mi}{1}\PYG{p}{]}\PYG{p}{;}
\PYG{+w}{  }\PYG{n}{damp}\PYG{+w}{ }\PYG{o}{=}\PYG{+w}{ }\PYG{n}{coeffs}\PYG{p}{[}\PYG{l+m+mi}{2}\PYG{p}{]}\PYG{p}{;}
\PYG{+w}{  }\PYG{n}{delta\PYGZus{}switch}\PYG{+w}{ }\PYG{o}{=}\PYG{+w}{ }\PYG{n}{coeffs}\PYG{p}{[}\PYG{l+m+mi}{3}\PYG{p}{]}\PYG{p}{;}
\PYG{p}{\PYGZcb{}}

\PYG{c+cm}{/* \PYGZhy{}\PYGZhy{}\PYGZhy{}\PYGZhy{}\PYGZhy{}\PYGZhy{}\PYGZhy{}\PYGZhy{}\PYGZhy{}\PYGZhy{}\PYGZhy{}\PYGZhy{}\PYGZhy{}\PYGZhy{}\PYGZhy{}\PYGZhy{}\PYGZhy{}\PYGZhy{}\PYGZhy{}\PYGZhy{}\PYGZhy{}\PYGZhy{}\PYGZhy{}\PYGZhy{}\PYGZhy{}\PYGZhy{}\PYGZhy{}\PYGZhy{}\PYGZhy{}\PYGZhy{}\PYGZhy{}\PYGZhy{}\PYGZhy{}\PYGZhy{}\PYGZhy{}\PYGZhy{}\PYGZhy{}\PYGZhy{}\PYGZhy{}\PYGZhy{}\PYGZhy{}\PYGZhy{}\PYGZhy{}\PYGZhy{}\PYGZhy{}\PYGZhy{}\PYGZhy{}\PYGZhy{}\PYGZhy{}\PYGZhy{}\PYGZhy{}\PYGZhy{}\PYGZhy{}\PYGZhy{}\PYGZhy{}\PYGZhy{}\PYGZhy{}\PYGZhy{}\PYGZhy{}\PYGZhy{}\PYGZhy{}\PYGZhy{}\PYGZhy{}\PYGZhy{}\PYGZhy{}\PYGZhy{}\PYGZhy{}\PYGZhy{}\PYGZhy{}\PYGZhy{} */}

\PYG{k+kt}{double}\PYG{+w}{ }\PYG{n}{GranSubModNormalHookePiecewise}\PYG{o}{:}\PYG{o}{:}\PYG{n}{calculate\PYGZus{}forces}\PYG{p}{(}\PYG{p}{)}
\PYG{p}{\PYGZob{}}
\PYG{+w}{  }\PYG{k+kt}{double}\PYG{+w}{ }\PYG{n}{Fne}\PYG{p}{;}
\PYG{+w}{  }\PYG{k}{if}\PYG{+w}{ }\PYG{p}{(}\PYG{n}{gm}\PYG{o}{\PYGZhy{}}\PYG{o}{\PYGZgt{}}\PYG{n}{delta}\PYG{+w}{ }\PYG{o}{\PYGZgt{}}\PYG{o}{=}\PYG{+w}{ }\PYG{n}{delta\PYGZus{}switch}\PYG{p}{)}\PYG{+w}{ }\PYG{p}{\PYGZob{}}
\PYG{+w}{    }\PYG{n}{Fne}\PYG{+w}{ }\PYG{o}{=}\PYG{+w}{ }\PYG{n}{k1}\PYG{+w}{ }\PYG{o}{*}\PYG{+w}{ }\PYG{n}{delta\PYGZus{}switch}\PYG{+w}{ }\PYG{o}{+}\PYG{+w}{ }\PYG{n}{k2}\PYG{+w}{ }\PYG{o}{*}\PYG{+w}{ }\PYG{p}{(}\PYG{n}{gm}\PYG{o}{\PYGZhy{}}\PYG{o}{\PYGZgt{}}\PYG{n}{delta}\PYG{+w}{ }\PYG{o}{\PYGZhy{}}\PYG{+w}{ }\PYG{n}{delta\PYGZus{}switch}\PYG{p}{)}\PYG{p}{;}
\PYG{+w}{  }\PYG{p}{\PYGZcb{}}\PYG{+w}{ }\PYG{k}{else}\PYG{+w}{ }\PYG{p}{\PYGZob{}}
\PYG{+w}{    }\PYG{n}{Fne}\PYG{+w}{ }\PYG{o}{=}\PYG{+w}{ }\PYG{n}{k1}\PYG{+w}{ }\PYG{o}{*}\PYG{+w}{ }\PYG{n}{gm}\PYG{o}{\PYGZhy{}}\PYG{o}{\PYGZgt{}}\PYG{n}{delta}\PYG{p}{;}
\PYG{+w}{  }\PYG{p}{\PYGZcb{}}
\PYG{+w}{  }\PYG{k}{return}\PYG{+w}{ }\PYG{n}{Fne}\PYG{p}{;}
\PYG{p}{\PYGZcb{}}
\end{sphinxVerbatim}

\sphinxstepscope


\section{Thermodynamic output options}
\label{\detokenize{Modify_thermo:thermodynamic-output-options}}\label{\detokenize{Modify_thermo::doc}}
\sphinxAtStartPar
The \sphinxcode{\sphinxupquote{Thermo}} class computes and prints thermodynamic information to
the screen and log file; see the files \sphinxcode{\sphinxupquote{thermo.cpp}} and \sphinxcode{\sphinxupquote{thermo.h}}.

\sphinxAtStartPar
There are four styles defined in \sphinxcode{\sphinxupquote{thermo.cpp}}: “one”, “multi”, “yaml”,
and “custom”.  The “custom” style allows the user to explicitly list
keywords for individual quantities to print when thermodynamic output is
generated.  The others have a fixed list of keywords.  See the
{\hyperref[\detokenize{thermo_style::doc}]{\sphinxcrossref{\DUrole{doc}{thermo\_style}}}} command for a list of available
quantities.  The formatting of the “custom” style defaults to the “one”
style, but can be adapted using {\hyperref[\detokenize{thermo_modify::doc}]{\sphinxcrossref{\DUrole{doc}{thermo\_modify line}}}}.

\sphinxAtStartPar
The thermo styles (one, multi, etc) are defined by lists of keywords
with associated formats for integer and floating point numbers and
identified by an enumerator constant.  Adding a new style thus mostly
requires defining a new list of keywords and the associated formats and
then inserting the required output processing where the enumerators are
identified.  Search for the word “CUSTOMIZATION” with references to
“thermo style” in the \sphinxcode{\sphinxupquote{thermo.cpp}} file to see the locations where
code will need to be added.  The member function \sphinxcode{\sphinxupquote{Thermo::header()}}
prints output at the very beginning of a thermodynamic output block and
can be used to print column headers or other front matter.  The member
function \sphinxcode{\sphinxupquote{Thermo::footer()}} prints output at the end of a
thermodynamic output block.  The formatting of the output is done by
assembling a “line” (which may span multiple lines if the style inserts
newline characters (”n” as in the “multi” style).

\sphinxAtStartPar
New thermodynamic keywords can also be added to \sphinxcode{\sphinxupquote{thermo.cpp}} to
compute new quantities for output.  Search for the word “CUSTOMIZATION”
with references to “keyword” in \sphinxcode{\sphinxupquote{thermo.cpp}} to see the several
locations where code will need to be added.  Effectively, you need to
define a member function that computes the property, add an if statement
in \sphinxcode{\sphinxupquote{Thermo::parse\_fields()}} where the corresponding header string for
the keyword and the function pointer is registered by calling the
\sphinxcode{\sphinxupquote{Thermo::addfield()}} method, and add an if statement in
\sphinxcode{\sphinxupquote{Thermo::evaluate\_keyword()}} which is called from the \sphinxcode{\sphinxupquote{Variable}}
class when a thermo keyword is encountered.

\begin{sphinxadmonition}{note}{Note:}
\sphinxAtStartPar
The third argument to \sphinxcode{\sphinxupquote{Thermo::addfield()}} is a flag indicating
whether the function for the keyword computes a floating point
(FLOAT), regular integer (INT), or big integer (BIGINT) value.  This
information is used for formatting the thermodynamic output.  Inside
the function the result must then be stored either in the \sphinxcode{\sphinxupquote{dvalue}},
\sphinxcode{\sphinxupquote{ivalue}} or \sphinxcode{\sphinxupquote{bivalue}} member variable, respectively.
\end{sphinxadmonition}

\sphinxAtStartPar
Since the {\hyperref[\detokenize{thermo_style::doc}]{\sphinxcrossref{\DUrole{doc}{thermo\_style custom}}}} command allows to
use output of quantities calculated by {\hyperref[\detokenize{fix::doc}]{\sphinxcrossref{\DUrole{doc}{fixes}}}},
{\hyperref[\detokenize{compute::doc}]{\sphinxcrossref{\DUrole{doc}{computes}}}}, and {\hyperref[\detokenize{variable::doc}]{\sphinxcrossref{\DUrole{doc}{variables}}}}, it may often
be simpler to compute what you wish via one of those constructs, rather
than by adding a new keyword to the thermo\_style command.

\sphinxstepscope


\section{Variable options}
\label{\detokenize{Modify_variable:variable-options}}\label{\detokenize{Modify_variable::doc}}
\sphinxAtStartPar
The \sphinxcode{\sphinxupquote{Variable}} class computes and stores {\hyperref[\detokenize{variable::doc}]{\sphinxcrossref{\DUrole{doc}{variable}}}}
information in LAMMPS; see the file \sphinxcode{\sphinxupquote{variable.cpp}}.  The value
associated with a variable can be periodically printed to the screen
via the {\hyperref[\detokenize{print::doc}]{\sphinxcrossref{\DUrole{doc}{print}}}}, {\hyperref[\detokenize{fix_print::doc}]{\sphinxcrossref{\DUrole{doc}{fix print}}}}, or
{\hyperref[\detokenize{thermo_style::doc}]{\sphinxcrossref{\DUrole{doc}{thermo\_style custom}}}} commands.  Variables of style
“equal” can compute complex equations that involve the following types
of arguments:

\begin{sphinxVerbatim}[commandchars=\\\{\}]
thermo keywords = ke, vol, atoms, ...
other variables = v\PYGZus{}a, v\PYGZus{}myvar, ...
math functions = div(x,y), mult(x,y), add(x,y), ...
group functions = mass(group), xcm(group,x), ...
atom values = x[123], y[3], vx[34], ...
compute values = c\PYGZus{}mytemp[0], c\PYGZus{}thermo\PYGZus{}press[3], ...
\end{sphinxVerbatim}

\sphinxAtStartPar
Adding keywords for the {\hyperref[\detokenize{thermo_style::doc}]{\sphinxcrossref{\DUrole{doc}{thermo\_style custom}}}}
command (which can then be accessed by variables) is discussed in the
{\hyperref[\detokenize{Modify_thermo::doc}]{\sphinxcrossref{\DUrole{doc}{Modify thermo}}}} documentation.

\sphinxAtStartPar
Adding a new math function of one or two arguments can be done by
editing one section of the \sphinxcode{\sphinxupquote{Variable::evaluate()}} method.  Search for
the word “customize” to find the appropriate location.

\sphinxAtStartPar
Adding a new group function can be done by editing one section of the
\sphinxcode{\sphinxupquote{Variable::evaluate()}} method.  Search for the word “customize” to
find the appropriate location.  You may need to add a new method to the
Group class as well (see the \sphinxcode{\sphinxupquote{group.cpp}} file).

\sphinxAtStartPar
Accessing a new atom\sphinxhyphen{}based vector can be done by editing one section
of the Variable::evaluate() method.  Search for the word “customize”
to find the appropriate location.

\sphinxAtStartPar
Adding new {\hyperref[\detokenize{compute::doc}]{\sphinxcrossref{\DUrole{doc}{compute styles}}}} (whose calculated values can
then be accessed by variables) is discussed in the {\hyperref[\detokenize{Modify_compute::doc}]{\sphinxcrossref{\DUrole{doc}{Modify compute}}}} documentation.

\sphinxstepscope


\chapter{Information for Developers}
\label{\detokenize{Developer:information-for-developers}}\label{\detokenize{Developer::doc}}
\sphinxAtStartPar
This section describes the internal structure and basic algorithms
of the LAMMPS code. This is a work in progress and additional
information will be added incrementally depending on availability
of time and requests from the LAMMPS user community.

\sphinxstepscope


\section{Source files}
\label{\detokenize{Developer_org:source-files}}\label{\detokenize{Developer_org::doc}}
\sphinxAtStartPar
The source files of the LAMMPS code are found in two directories of the
distribution: \sphinxcode{\sphinxupquote{src}} and \sphinxcode{\sphinxupquote{lib}}.  Most of the code is written in C++
but there are small a number of files in several other languages like C,
Fortran, Shell script, or Python.

\sphinxAtStartPar
The core of the code is located in the \sphinxcode{\sphinxupquote{src}} folder and its
subdirectories.  A sizable number of these files are in the \sphinxcode{\sphinxupquote{src}}
directory itself, but there are plenty of {\hyperref[\detokenize{Packages::doc}]{\sphinxcrossref{\DUrole{doc}{packages}}}},
which can be included or excluded when LAMMPS is built.  See the
{\hyperref[\detokenize{Build_package::doc}]{\sphinxcrossref{\DUrole{doc}{Include packages in build}}}} section of the manual
for more information about that part of the build process.  LAMMPS
currently supports building with {\hyperref[\detokenize{Build_make::doc}]{\sphinxcrossref{\DUrole{doc}{conventional makefiles}}}} and through {\hyperref[\detokenize{Build_cmake::doc}]{\sphinxcrossref{\DUrole{doc}{CMake}}}}.  Those procedures
differ in how packages are enabled or disabled for inclusion into a
LAMMPS binary, so they cannot be mixed.  The source files for each
package are in all\sphinxhyphen{}uppercase subdirectories of the \sphinxcode{\sphinxupquote{src}} folder, for
example \sphinxcode{\sphinxupquote{src/MOLECULE}} or \sphinxcode{\sphinxupquote{src/EXTRA\sphinxhyphen{}MOLECULE}}.  The \sphinxcode{\sphinxupquote{src/STUBS}}
subdirectory is not a package but contains a dummy MPI library, that is
used when building a serial version of the code. The \sphinxcode{\sphinxupquote{src/MAKE}}
directory and its subdirectories contain makefiles with settings and
flags for a variety of configuration and machines for the build process
with traditional makefiles.

\sphinxAtStartPar
The \sphinxcode{\sphinxupquote{lib}} directory contains the source code for several supporting
libraries or files with configuration settings to use globally installed
libraries, that are required by some optional packages.  They may
include python scripts that can transparently download additional source
code on request.  Each subdirectory, like \sphinxcode{\sphinxupquote{lib/poems}} or \sphinxcode{\sphinxupquote{lib/gpu}},
contains the source files, some of which are in different languages such
as Fortran or CUDA. These libraries included in the LAMMPS build, if the
corresponding package is installed.

\sphinxAtStartPar
LAMMPS C++ source files almost always come in pairs, such as
\sphinxcode{\sphinxupquote{src/run.cpp}} (implementation file) and \sphinxcode{\sphinxupquote{src/run.h}} (header file).
Each pair of files defines a C++ class, for example the
\sphinxcode{\sphinxupquote{LAMMPS\_NS::Run}} class, which contains the code invoked by
the {\hyperref[\detokenize{run::doc}]{\sphinxcrossref{\DUrole{doc}{run}}}} command in a LAMMPS input script.  As this example
illustrates, source file and class names often have a one\sphinxhyphen{}to\sphinxhyphen{}one
correspondence with a command used in a LAMMPS input script.  Some
source files and classes do not have a corresponding input script
command, for example \sphinxcode{\sphinxupquote{src/force.cpp}} and the \sphinxcode{\sphinxupquote{LAMMPS\_NS::Force}}
class.  They are discussed in the next section.

\sphinxAtStartPar
The names of all source files are in lower case and may use the
underscore character ‘\_’ to separate words. Apart from bundled,
externally maintained libraries, which may have different conventions,
all C and C++ header files have a \sphinxcode{\sphinxupquote{.h}} extension, all C++ files have a
\sphinxcode{\sphinxupquote{.cpp}} extension, and C files a \sphinxcode{\sphinxupquote{.c}} extension.  A few C++ classes
and utility functions are implemented with only a \sphinxcode{\sphinxupquote{.h}} file. Examples
are the Pointers and Commands classes or the MathVec functions.


\section{Class topology}
\label{\detokenize{Developer_org:class-topology}}
\sphinxAtStartPar
Though LAMMPS has a lot of source files and classes, its class topology
is not very deep, which can be seen from the {\hyperref[\detokenize{Developer_org:id1}]{\sphinxcrossref{\DUrole{std,std-ref}{LAMMPS class topology}}}}
figure.  In that figure, each name refers to a class and has a pair of
associated source files in the \sphinxcode{\sphinxupquote{src}} folder, for example the class
\sphinxcode{\sphinxupquote{LAMMPS\_NS::Memory}} corresponds to the files \sphinxcode{\sphinxupquote{memory.cpp}}
and \sphinxcode{\sphinxupquote{memory.h}}, or the class \sphinxcode{\sphinxupquote{LAMMPS\_NS::AtomVec}}
corresponds to the files \sphinxcode{\sphinxupquote{atom\_vec.cpp}} and \sphinxcode{\sphinxupquote{atom\_vec.h}}.  Full
lines in the figure represent compositing: that is, the class at the
base of the arrow holds a pointer to an instance of the class at the
tip.  Dashed lines instead represent inheritance: the class at the tip
of the arrow is derived from the class at the base.  Classes with a red
boundary are not instantiated directly, but they represent the base
classes for “styles”.  Those “styles” make up the bulk of the LAMMPS
code and only a few representative examples are included in the figure,
so it remains readable.

\begin{figure}[htbp]
\centering
\capstart

\noindent\sphinxincludegraphics{{lammps-classes}.png}
\caption{LAMMPS class topology}\label{\detokenize{Developer_org:id2}}\label{\detokenize{Developer_org:id1}}
\begin{sphinxlegend}\begin{quote}

\sphinxAtStartPar
This figure shows relations of base classes of the LAMMPS
simulation package.  Full lines indicate that a class holds an
instance of the class it is pointing to; dashed lines point to
derived classes that are given as examples of what classes may be
instantiated during a LAMMPS run based on the input commands and
accessed through the API define by their respective base classes.
At the core is the {\hyperref[\detokenize{Classes_lammps:_CPPv4N9LAMMPS_NS6LAMMPSE}]{\sphinxcrossref{\sphinxcode{\sphinxupquote{LAMMPS}}}}} class,
which holds pointers to class instances with specific purposes.
Those may hold instances of other classes, sometimes directly, or
only temporarily, sometimes as derived classes or derived classes
of derived classes, which may also hold instances of other
classes.
\end{quote}
\end{sphinxlegend}
\end{figure}

\sphinxAtStartPar
The {\hyperref[\detokenize{Classes_lammps:_CPPv4N9LAMMPS_NS6LAMMPSE}]{\sphinxcrossref{\sphinxcode{\sphinxupquote{LAMMPS\_NS::LAMMPS}}}}} class is the topmost class and
represents what is generally referred to as an “instance of LAMMPS”.  It
is a composite holding pointers to instances of other core classes
providing the core functionality of the MD engine in LAMMPS and through
them abstractions of the required operations.  The constructor of the
LAMMPS class will instantiate those instances, process the command line
flags, initialize MPI (if not already done) and set up file pointers for
input and output.  The destructor will shut everything down and free all
associated memory.  Thus code for the standalone LAMMPS executable in
\sphinxcode{\sphinxupquote{main.cpp}} simply initializes MPI, instantiates a single instance of
LAMMPS while passing it the command line flags and input script. It
deletes the LAMMPS instance after the method reading the input returns
and shuts down the MPI environment before it exits the executable.

\sphinxAtStartPar
The {\hyperref[\detokenize{Classes_lammps:_CPPv4N9LAMMPS_NS8PointersE}]{\sphinxcrossref{\sphinxcode{\sphinxupquote{LAMMPS\_NS::Pointers}}}}} class is not shown in the
{\hyperref[\detokenize{Developer_org:id1}]{\sphinxcrossref{\DUrole{std,std-ref}{LAMMPS class topology}}}} figure for clarity.  It holds references to many
of the members of the \sphinxtitleref{LAMMPS\_NS::LAMMPS}, so that all classes derived
from {\hyperref[\detokenize{Classes_lammps:_CPPv4N9LAMMPS_NS8PointersE}]{\sphinxcrossref{\sphinxcode{\sphinxupquote{LAMMPS\_NS::Pointers}}}}} have direct access to those
references.  From the class topology all classes with blue boundary are
referenced in the Pointers class and all classes in the second and third
columns, that are not listed as derived classes, are instead derived
from {\hyperref[\detokenize{Classes_lammps:_CPPv4N9LAMMPS_NS8PointersE}]{\sphinxcrossref{\sphinxcode{\sphinxupquote{LAMMPS\_NS::Pointers}}}}}.  To initialize the pointer
references in Pointers, a pointer to the LAMMPS class instance needs to
be passed to the constructor. All constructors for classes derived from
it, must do so and thus pass that pointer to the constructor for
{\hyperref[\detokenize{Classes_lammps:_CPPv4N9LAMMPS_NS8PointersE}]{\sphinxcrossref{\sphinxcode{\sphinxupquote{LAMMPS\_NS::Pointers}}}}}.  The default constructor for
{\hyperref[\detokenize{Classes_lammps:_CPPv4N9LAMMPS_NS8PointersE}]{\sphinxcrossref{\sphinxcode{\sphinxupquote{LAMMPS\_NS::Pointers}}}}} is disabled to enforce this.

\sphinxAtStartPar
Since all storage is supposed to be encapsulated (there are a few
exceptions), the LAMMPS class can also be instantiated multiple times by
a calling code.  Outside the aforementioned exceptions, those LAMMPS
instances can be used alternately.  As of the time of this writing
(early 2023) LAMMPS is not yet sufficiently thread\sphinxhyphen{}safe for concurrent
execution.  When running in parallel with MPI, care has to be taken,
that suitable copies of communicators are used to not create conflicts
between different instances.

\sphinxAtStartPar
The LAMMPS class currently holds instances of 19 classes representing
the core functionality.  There are a handful of virtual parent classes
in LAMMPS that define what LAMMPS calls \sphinxcode{\sphinxupquote{styles}}.  These are shaded
red in the {\hyperref[\detokenize{Developer_org:id1}]{\sphinxcrossref{\DUrole{std,std-ref}{LAMMPS class topology}}}} figure.  Each of these are parents of a
number of child classes that implement the interface defined by the
parent class.  There are two main categories of these \sphinxcode{\sphinxupquote{styles}}: some
may only have one instance active at a time (e.g. atom, pair, bond,
angle, dihedral, improper, kspace, comm) and there is a dedicated
pointer variable for each of them in the corresponding composite class.
Setups that require a mix of different such styles have to use a
\sphinxstyleemphasis{hybrid} class instance that acts as a proxy, and manages and forwards
calls to the corresponding sub\sphinxhyphen{}style class instances for the designated
subset of atoms or data.  The composite class may also have lists of
class instances, e.g. \sphinxcode{\sphinxupquote{Modify}} handles lists of compute and fix
styles, while \sphinxcode{\sphinxupquote{Output}} handles a list of dump class instances.

\sphinxAtStartPar
The exception to this scheme are the \sphinxcode{\sphinxupquote{command}} style classes.  These
implement specific commands that can be invoked before, after, or in
between runs.  For these an instance of the class is created, its
command() method called and then, after completion, the class instance
deleted.  Examples for this are the create\_box, create\_atoms, minimize,
run, set, or velocity command styles.

\sphinxAtStartPar
For all those \sphinxcode{\sphinxupquote{styles}}, certain naming conventions are employed: for
the fix nve command the class is called FixNVE and the source files are
\sphinxcode{\sphinxupquote{fix\_nve.h}} and \sphinxcode{\sphinxupquote{fix\_nve.cpp}}. Similarly, for fix ave/time we have
FixAveTime and \sphinxcode{\sphinxupquote{fix\_ave\_time.h}} and \sphinxcode{\sphinxupquote{fix\_ave\_time.cpp}}.  Style names
are lower case and without spaces or special characters. A suffix or
words are appended with a forward slash ‘/’ which denotes a variant of
the corresponding class without the suffix.  To connect the style name
and the class name, LAMMPS uses macros like: \sphinxcode{\sphinxupquote{AtomStyle()}},
\sphinxcode{\sphinxupquote{PairStyle()}}, \sphinxcode{\sphinxupquote{BondStyle()}}, \sphinxcode{\sphinxupquote{RegionStyle()}}, and so on in the
corresponding header file.  During configuration or compilation, files
with the pattern \sphinxcode{\sphinxupquote{style\_\textless{}name\textgreater{}.h}} are created that consist of a list
of include statements including all headers of all styles of a given
type that are currently enabled (or “installed”).

\sphinxAtStartPar
More details on individual classes in the {\hyperref[\detokenize{Developer_org:id1}]{\sphinxcrossref{\DUrole{std,std-ref}{LAMMPS class topology}}}} are as
follows:
\begin{itemize}
\item {} 
\sphinxAtStartPar
The Memory class handles allocation of all large vectors and arrays.

\item {} 
\sphinxAtStartPar
The Error class prints all (terminal) error and warning messages.

\item {} 
\sphinxAtStartPar
The Universe class sets up one or more partitions of processors so
that one or multiple simulations can be run, on the processors
allocated for a run, e.g. by the mpirun command.

\item {} 
\sphinxAtStartPar
The Input class reads and processes input (strings and files), stores
variables, and invokes {\hyperref[\detokenize{Commands_all::doc}]{\sphinxcrossref{\DUrole{doc}{commands}}}}.

\item {} 
\sphinxAtStartPar
Command style classes are derived from the Command class. They provide
input script commands that perform one\sphinxhyphen{}time operations
before/after/between simulations or which invoke a simulation.  They
are usually instantiated from within the Input class, its \sphinxcode{\sphinxupquote{command}}
method invoked, and then immediately destructed.

\item {} 
\sphinxAtStartPar
The Finish class is instantiated to print statistics to the screen
after a simulation is performed, by commands like run and minimize.

\item {} 
\sphinxAtStartPar
The Special class walks the bond topology of a molecular system to
find first, second, third neighbors of each atom.  It is invoked by
several commands, like {\hyperref[\detokenize{read_data::doc}]{\sphinxcrossref{\DUrole{doc}{read\_data}}}},
{\hyperref[\detokenize{read_restart::doc}]{\sphinxcrossref{\DUrole{doc}{read\_restart}}}}, or {\hyperref[\detokenize{replicate::doc}]{\sphinxcrossref{\DUrole{doc}{replicate}}}}.

\item {} 
\sphinxAtStartPar
The Atom class stores per\sphinxhyphen{}atom properties associated with atom styles.
More precisely, they are allocated and managed by a class derived from
the AtomVec class, and the Atom class simply stores pointers to them.
The classes derived from AtomVec represent the different atom styles,
and they are instantiated through the {\hyperref[\detokenize{atom_style::doc}]{\sphinxcrossref{\DUrole{doc}{atom\_style}}}}
command.

\item {} 
\sphinxAtStartPar
The Update class holds instances of an integrator and a minimizer
class.  The Integrate class is a parent style for the Verlet and
r\sphinxhyphen{}RESPA time integrators, as defined by the {\hyperref[\detokenize{run_style::doc}]{\sphinxcrossref{\DUrole{doc}{run\_style}}}} command.  The Min class is a parent style for various
energy minimizers.

\item {} 
\sphinxAtStartPar
The Neighbor class builds and stores neighbor lists.  The NeighList
class stores a single list (for all atoms).  A NeighRequest class
instance is created by pair, fix, or compute styles when they need a
particular kind of neighbor list and use the NeighRequest properties
to select the neighbor list settings for the given request.  There can
be multiple instances of the NeighRequest class. The Neighbor class
will try to optimize how the requests are processed.  Depending on the
NeighRequest properties, neighbor lists are constructed from scratch,
aliased, or constructed by post\sphinxhyphen{}processing an existing list into
sub\sphinxhyphen{}lists.

\item {} 
\sphinxAtStartPar
The Comm class performs inter\sphinxhyphen{}processor communication, typically of
ghost atom information.  This usually involves MPI message exchanges
with 6 neighboring processors in the 3d logical grid of processors
mapped to the simulation box. There are two {\hyperref[\detokenize{comm_style::doc}]{\sphinxcrossref{\DUrole{doc}{communication styles}}}}, enabling different ways to perform the domain
decomposition.

\item {} 
\sphinxAtStartPar
The Irregular class is used, when atoms may migrate to arbitrary
processors.

\item {} 
\sphinxAtStartPar
The Domain class stores the simulation box geometry, as well as
geometric Regions and any user definition of a Lattice.  The latter
are defined by the {\hyperref[\detokenize{region::doc}]{\sphinxcrossref{\DUrole{doc}{region}}}} and {\hyperref[\detokenize{lattice::doc}]{\sphinxcrossref{\DUrole{doc}{lattice}}}}
commands in an input script.

\item {} 
\sphinxAtStartPar
The Force class computes various forces between atoms.  The Pair
parent class is for non\sphinxhyphen{}bonded or pairwise forces, which in LAMMPS
also includes many\sphinxhyphen{}body forces such as the Tersoff 3\sphinxhyphen{}body potential if
those are computed by walking pairwise neighbor lists.  The Bond,
Angle, Dihedral, Improper parent classes are styles for bonded
interactions within a static molecular topology.  The KSpace parent
class is for computing long\sphinxhyphen{}range Coulombic interactions.  One of its
child classes, PPPM, uses the FFT3D and Remap classes to redistribute
and communicate grid\sphinxhyphen{}based information across the parallel processors.

\item {} 
\sphinxAtStartPar
The Modify class stores lists of class instances derived from the
{\hyperref[\detokenize{fix::doc}]{\sphinxcrossref{\DUrole{doc}{Fix}}}} and {\hyperref[\detokenize{compute::doc}]{\sphinxcrossref{\DUrole{doc}{Compute}}}} base classes.

\item {} 
\sphinxAtStartPar
The Group class manipulates groups that atoms are assigned to via the
{\hyperref[\detokenize{group::doc}]{\sphinxcrossref{\DUrole{doc}{group}}}} command.  It also has functions to compute
various attributes of groups of atoms.

\item {} 
\sphinxAtStartPar
The Output class is used to generate 3 kinds of output from a LAMMPS
simulation: thermodynamic information printed to the screen and log
file, dump file snapshots, and restart files.  These correspond to the
{\hyperref[\detokenize{thermo_style::doc}]{\sphinxcrossref{\DUrole{doc}{Thermo}}}}, {\hyperref[\detokenize{dump::doc}]{\sphinxcrossref{\DUrole{doc}{Dump}}}}, and
{\hyperref[\detokenize{write_restart::doc}]{\sphinxcrossref{\DUrole{doc}{WriteRestart}}}} classes respectively.  The Dump
class is a base class, with several derived classes implementing
various dump style variants.

\item {} 
\sphinxAtStartPar
The Timer class logs timing information, output at the end
of a run.

\end{itemize}

\sphinxstepscope


\section{Code design}
\label{\detokenize{Developer_code_design:code-design}}\label{\detokenize{Developer_code_design::doc}}
\sphinxAtStartPar
This section explains some code design choices in LAMMPS with the goal
of helping developers write new code similar to the existing code.
Please see the section on {\hyperref[\detokenize{Modify_style::doc}]{\sphinxcrossref{\DUrole{doc}{Requirements for contributed code}}}} for more specific recommendations and guidelines.  While
that section is organized more in the form of a checklist for code
contributors, the focus here is on overall code design strategy, choices
made between possible alternatives, and discussing some relevant C++
programming language constructs.

\sphinxAtStartPar
Historically, the basic design philosophy of the LAMMPS C++ code was a
“C with classes” style.  The motivation was to make it easy to modify
LAMMPS for people without significant training in C++ programming.  Data
structures and code constructs were used that resemble the previous
implementation(s) in Fortran.  A contributing factor to this choice was
that at the time, C++ compilers were often not mature and some advanced
features contained bugs or did not function as the standard required.
There were also disagreements between compiler vendors as to how to
interpret the C++ standard documents.

\sphinxAtStartPar
However, C++ compilers and the C++ programming language have advanced
significantly.  In 2020, the LAMMPS developers decided to require the
C++11 standard as the minimum C++ language standard for LAMMPS.  Since
then, we have begun to replace C\sphinxhyphen{}style constructs with equivalent C++
functionality.  This was taken either from the C++ standard library or
implemented as custom classes or functions.  The goal is to improve
readability of the code and to increase code reuse through abstraction
of commonly used functionality.

\begin{sphinxadmonition}{note}{Note:}
\sphinxAtStartPar
Please note that as of spring 2023 there is still a sizable chunk of
legacy code in LAMMPS that has not yet been refactored to reflect
these style conventions in full.  LAMMPS has a large code base and
many contributors.  There is also a hierarchy of precedence in which
the code is adapted.  Highest priority has been the code in the
\sphinxcode{\sphinxupquote{src}} folder, followed by code in packages in order of their
popularity and complexity (simpler code gets adapted sooner), followed
by code in the \sphinxcode{\sphinxupquote{lib}} folder.  Source code that is downloaded from
external packages or libraries during compilation is not subject to
the conventions discussed here.
\end{sphinxadmonition}


\subsection{Object\sphinxhyphen{}oriented code}
\label{\detokenize{Developer_code_design:object-oriented-code}}
\sphinxAtStartPar
LAMMPS is designed to be an object\sphinxhyphen{}oriented code.  Each simulation is
represented by an instance of the LAMMPS class.  When running in
parallel, each MPI process creates such an instance.  This can be seen
in the \sphinxcode{\sphinxupquote{main.cpp}} file where the core steps of running a LAMMPS
simulation are the following 3 lines of code:

\begin{sphinxVerbatim}[commandchars=\\\{\}]
\PYG{n}{LAMMPS}\PYG{+w}{ }\PYG{o}{*}\PYG{n}{lammps}\PYG{+w}{ }\PYG{o}{=}\PYG{+w}{ }\PYG{k}{new}\PYG{+w}{ }\PYG{n}{LAMMPS}\PYG{p}{(}\PYG{n}{argc}\PYG{p}{,}\PYG{+w}{ }\PYG{n}{argv}\PYG{p}{,}\PYG{+w}{ }\PYG{n}{lammps\PYGZus{}comm}\PYG{p}{)}\PYG{p}{;}
\PYG{n}{lammps}\PYG{o}{\PYGZhy{}}\PYG{o}{\PYGZgt{}}\PYG{n}{input}\PYG{o}{\PYGZhy{}}\PYG{o}{\PYGZgt{}}\PYG{n}{file}\PYG{p}{(}\PYG{p}{)}\PYG{p}{;}
\PYG{k}{delete}\PYG{+w}{ }\PYG{n}{lammps}\PYG{p}{;}
\end{sphinxVerbatim}

\sphinxAtStartPar
The first line creates a LAMMPS class instance and passes the command
line arguments and the global communicator to its constructor.  The
second line triggers the LAMMPS instance to process the input (either
from standard input or a provided input file) until the simulation
ends.  The third line deletes the LAMMPS instance.  The remainder of
the main.cpp file has code for error handling, MPI configuration, and
other special features.

\sphinxAtStartPar
The basic LAMMPS class hierarchy which is created by the LAMMPS class
constructor is shown in {\hyperref[\detokenize{Developer_org:id1}]{\sphinxcrossref{\DUrole{std,std-ref}{LAMMPS class topology}}}}.  When input commands
are processed, additional class instances are created, or deleted, or
replaced.  Likewise, specific member functions of specific classes are
called to trigger actions such as creating atoms, computing forces,
computing properties, time\sphinxhyphen{}propagating the system, or writing output.


\subsubsection{Compositing and Inheritance}
\label{\detokenize{Developer_code_design:compositing-and-inheritance}}
\sphinxAtStartPar
LAMMPS makes extensive use of the object\sphinxhyphen{}oriented programming (OOP)
principles of \sphinxstyleemphasis{compositing} and \sphinxstyleemphasis{inheritance}. Classes like the
\sphinxcode{\sphinxupquote{LAMMPS}} class are a \sphinxstylestrong{composite} containing pointers to instances
of other classes like \sphinxcode{\sphinxupquote{Atom}}, \sphinxcode{\sphinxupquote{Comm}}, \sphinxcode{\sphinxupquote{Force}}, \sphinxcode{\sphinxupquote{Neighbor}},
\sphinxcode{\sphinxupquote{Modify}}, and so on.  Each of these classes implements certain
functionality by storing and manipulating data related to the
simulation and providing member functions that trigger certain
actions.  Some of those classes like \sphinxcode{\sphinxupquote{Force}} are themselves
composites, containing instances of classes describing different force
interactions.  Similarly, the \sphinxcode{\sphinxupquote{Modify}} class contains a list of
\sphinxcode{\sphinxupquote{Fix}} and \sphinxcode{\sphinxupquote{Compute}} classes.  If the input commands that
correspond to these classes include the word \sphinxstyleemphasis{style}, then LAMMPS
stores only a single instance of that class.  E.g. \sphinxstyleemphasis{atom\_style},
\sphinxstyleemphasis{comm\_style}, \sphinxstyleemphasis{pair\_style}, \sphinxstyleemphasis{bond\_style}.  If the input command does
\sphinxstylestrong{not} include the word \sphinxstyleemphasis{style}, then there may be many instances of
that class defined, for example \sphinxstyleemphasis{region}, \sphinxstyleemphasis{fix}, \sphinxstyleemphasis{compute}, \sphinxstyleemphasis{dump}.

\sphinxAtStartPar
\sphinxstylestrong{Inheritance} enables creation of \sphinxstyleemphasis{derived} classes that can share
common functionality in their base class while providing a consistent
interface.  The derived classes replace (dummy or pure) functions in
the base class.  The higher level classes can then call those methods
of the instantiated classes without having to know which specific
derived class variant was instantiated.  In LAMMPS these derived
classes are often referred to as “styles”, e.g.  pair styles, fix
styles, atom styles and so on.

\sphinxAtStartPar
This is the origin of the flexibility of LAMMPS.  For example, pair
styles implement a variety of different non\sphinxhyphen{}bonded interatomic
potentials functions.  All details for the implementation of a
potential are stored and executed in a single class.

\sphinxAtStartPar
As mentioned above, there can be multiple instances of classes derived
from the \sphinxcode{\sphinxupquote{Fix}} or \sphinxcode{\sphinxupquote{Compute}} base classes.  They represent a
different facet of LAMMPS’ flexibility, as they provide methods which
can be called at different points within a timestep, as explained in
the {\hyperref[\detokenize{Developer_flow::doc}]{\sphinxcrossref{\DUrole{doc}{How a timestep works}}}} doc page.  This allows
the input script to tailor how a specific simulation is run, what
diagnostic computations are performed, and how the output of those
computations is further processed or output.

\sphinxAtStartPar
Additional code sharing is possible by creating derived classes from the
derived classes (e.g., to implement an accelerated version of a pair
style) where only a subset of the derived class methods are replaced
with accelerated versions.


\subsubsection{Polymorphism}
\label{\detokenize{Developer_code_design:polymorphism}}
\sphinxAtStartPar
Polymorphism and dynamic dispatch are another OOP feature that play an
important role in how LAMMPS selects what code to execute.  In a
nutshell, this is a mechanism where the decision of which member
function to call from a class is determined at runtime and not when
the code is compiled.  To enable it, the function has to be declared
as \sphinxcode{\sphinxupquote{virtual}} and all corresponding functions in derived classes
should use the \sphinxcode{\sphinxupquote{override}} property. Below is a brief example.

\begin{sphinxVerbatim}[commandchars=\\\{\}]
\PYG{k}{class}\PYG{+w}{ }\PYG{n+nc}{Base}\PYG{+w}{ }\PYG{p}{\PYGZob{}}
\PYG{k}{public}\PYG{o}{:}
\PYG{+w}{ }\PYG{k}{virtual}\PYG{+w}{ }\PYG{o}{\PYGZti{}}\PYG{n}{Base}\PYG{p}{(}\PYG{p}{)}\PYG{+w}{ }\PYG{o}{=}\PYG{+w}{ }\PYG{k}{default}\PYG{p}{;}
\PYG{+w}{ }\PYG{k+kt}{void}\PYG{+w}{ }\PYG{n+nf}{call}\PYG{p}{(}\PYG{p}{)}\PYG{p}{;}
\PYG{+w}{ }\PYG{k+kt}{void}\PYG{+w}{ }\PYG{n+nf}{normal}\PYG{p}{(}\PYG{p}{)}\PYG{p}{;}
\PYG{+w}{ }\PYG{k}{virtual}\PYG{+w}{ }\PYG{k+kt}{void}\PYG{+w}{ }\PYG{n+nf}{poly}\PYG{p}{(}\PYG{p}{)}\PYG{p}{;}
\PYG{p}{\PYGZcb{}}\PYG{p}{;}

\PYG{k+kt}{void}\PYG{+w}{ }\PYG{n+nf}{Base::call}\PYG{p}{(}\PYG{p}{)}\PYG{+w}{ }\PYG{p}{\PYGZob{}}
\PYG{+w}{ }\PYG{n}{normal}\PYG{p}{(}\PYG{p}{)}\PYG{p}{;}
\PYG{+w}{ }\PYG{n}{poly}\PYG{p}{(}\PYG{p}{)}\PYG{p}{;}
\PYG{p}{\PYGZcb{}}

\PYG{k}{class}\PYG{+w}{ }\PYG{n+nc}{Derived}\PYG{+w}{ }\PYG{o}{:}\PYG{+w}{ }\PYG{k}{public}\PYG{+w}{ }\PYG{n}{Base}\PYG{+w}{ }\PYG{p}{\PYGZob{}}
\PYG{k}{public}\PYG{o}{:}
\PYG{+w}{ }\PYG{o}{\PYGZti{}}\PYG{n}{Derived}\PYG{p}{(}\PYG{p}{)}\PYG{+w}{ }\PYG{k}{override}\PYG{+w}{ }\PYG{o}{=}\PYG{+w}{ }\PYG{k}{default}\PYG{p}{;}
\PYG{+w}{ }\PYG{k+kt}{void}\PYG{+w}{ }\PYG{n+nf}{normal}\PYG{p}{(}\PYG{p}{)}\PYG{p}{;}
\PYG{+w}{ }\PYG{k+kt}{void}\PYG{+w}{ }\PYG{n+nf}{poly}\PYG{p}{(}\PYG{p}{)}\PYG{+w}{ }\PYG{k}{override}\PYG{p}{;}
\PYG{p}{\PYGZcb{}}\PYG{p}{;}

\PYG{c+c1}{// [....]}

\PYG{n}{Base}\PYG{+w}{ }\PYG{o}{*}\PYG{n}{base1}\PYG{+w}{ }\PYG{o}{=}\PYG{+w}{ }\PYG{k}{new}\PYG{+w}{ }\PYG{n}{Base}\PYG{p}{(}\PYG{p}{)}\PYG{p}{;}
\PYG{n}{Base}\PYG{+w}{ }\PYG{o}{*}\PYG{n}{base2}\PYG{+w}{ }\PYG{o}{=}\PYG{+w}{ }\PYG{k}{new}\PYG{+w}{ }\PYG{n}{Derived}\PYG{p}{(}\PYG{p}{)}\PYG{p}{;}

\PYG{n}{base1}\PYG{o}{\PYGZhy{}}\PYG{o}{\PYGZgt{}}\PYG{n}{call}\PYG{p}{(}\PYG{p}{)}\PYG{p}{;}
\PYG{n}{base2}\PYG{o}{\PYGZhy{}}\PYG{o}{\PYGZgt{}}\PYG{n}{call}\PYG{p}{(}\PYG{p}{)}\PYG{p}{;}
\end{sphinxVerbatim}

\sphinxAtStartPar
The difference in behavior of the \sphinxcode{\sphinxupquote{normal()}} and the \sphinxcode{\sphinxupquote{poly()}} member
functions is which of the two member functions is called when executing
\sphinxtitleref{base1\sphinxhyphen{}\textgreater{}call()} versus \sphinxtitleref{base2\sphinxhyphen{}\textgreater{}call()}.  Without polymorphism, a
function within the base class can only call member functions within the
same scope: that is, \sphinxcode{\sphinxupquote{Base::call()}} will always call
\sphinxcode{\sphinxupquote{Base::normal()}}.  But for the \sphinxtitleref{base2\sphinxhyphen{}\textgreater{}call()} case, the call of the
virtual member function will be dispatched to \sphinxcode{\sphinxupquote{Derived::poly()}}
instead.  This mechanism results in calling functions that are within
the scope of the class that was used to \sphinxstyleemphasis{create} the instance, even if
they are assigned to a pointer for their base class.  This is the
desired behavior, and this way LAMMPS can even use styles that are loaded
at runtime from a shared object file with the {\hyperref[\detokenize{plugin::doc}]{\sphinxcrossref{\DUrole{doc}{plugin command}}}}.

\sphinxAtStartPar
A special case of virtual functions are so\sphinxhyphen{}called pure functions.  These
are virtual functions that are initialized to 0 in the class declaration
(see example below).

\begin{sphinxVerbatim}[commandchars=\\\{\}]
\PYG{k}{class}\PYG{+w}{ }\PYG{n+nc}{Base}\PYG{+w}{ }\PYG{p}{\PYGZob{}}
\PYG{k}{public}\PYG{o}{:}
\PYG{+w}{ }\PYG{k}{virtual}\PYG{+w}{ }\PYG{k+kt}{void}\PYG{+w}{ }\PYG{n}{pure}\PYG{p}{(}\PYG{p}{)}\PYG{+w}{ }\PYG{o}{=}\PYG{+w}{ }\PYG{l+m+mi}{0}\PYG{p}{;}
\PYG{p}{\PYGZcb{}}\PYG{p}{;}
\end{sphinxVerbatim}

\sphinxAtStartPar
This has the effect that an instance of the base class cannot be
created and that derived classes \sphinxstylestrong{must} implement these functions.
Many of the functions listed with the various class styles in the
section {\hyperref[\detokenize{Modify::doc}]{\sphinxcrossref{\DUrole{doc}{Modifying \& extending LAMMPS}}}} are pure functions.  The motivation for this is
to define the interface or API of the functions, but defer their
implementation to the derived classes.

\sphinxAtStartPar
However, there are downsides to this. For example, calls to virtual
functions from within a constructor, will \sphinxstyleemphasis{not} be in the scope of the
derived class, and thus it is good practice to either avoid calling them
or to provide an explicit scope such as \sphinxcode{\sphinxupquote{Base::poly()}} or
\sphinxcode{\sphinxupquote{Derived::poly()}}.  Furthermore, any destructors in classes containing
virtual functions should be declared virtual too, so they will be
processed in the expected order before types are removed from dynamic
dispatch.

\begin{sphinxadmonition}{note}{Important Notes}

\sphinxAtStartPar
In order to be able to detect incompatibilities at compile time and
to avoid unexpected behavior, it is crucial that all member functions
that are intended to replace a virtual or pure function use the
\sphinxcode{\sphinxupquote{override}} property keyword.  For the same reason, the use of
overloads or default arguments for virtual functions should be
avoided, as they lead to confusion over which function is supposed to
override which, and which arguments need to be declared.
\end{sphinxadmonition}


\subsubsection{Style Factories}
\label{\detokenize{Developer_code_design:style-factories}}
\sphinxAtStartPar
In order to create class instances for different styles, LAMMPS often
uses a programming pattern called \sphinxtitleref{Factory}.  Those are functions that
create an instance of a specific derived class, say \sphinxcode{\sphinxupquote{PairLJCut}} and
return a pointer to the type of the common base class of that style,
\sphinxcode{\sphinxupquote{Pair}} in this case.  To associate the factory function with the
style keyword, a \sphinxcode{\sphinxupquote{std::map}} class is used with function pointers
indexed by their keyword (for example “lj/cut” for \sphinxcode{\sphinxupquote{PairLJCut}} and
“morse” for \sphinxcode{\sphinxupquote{PairMorse}}).  A couple of typedefs help keep the code
readable, and a template function is used to implement the actual
factory functions for the individual classes.  Below is an example
of such a factory function from the \sphinxcode{\sphinxupquote{Force}} class as declared in
\sphinxcode{\sphinxupquote{force.h}} and implemented in \sphinxcode{\sphinxupquote{force.cpp}}.  The file \sphinxcode{\sphinxupquote{style\_pair.h}}
is generated during compilation and includes all main header files
(i.e. those starting with \sphinxcode{\sphinxupquote{pair\_}}) of pair styles and then the
macro \sphinxcode{\sphinxupquote{PairStyle()}} will associate the style name “lj/cut”
with a factory function creating an instance of the \sphinxcode{\sphinxupquote{PairLJCut}}
class.

\begin{sphinxVerbatim}[commandchars=\\\{\}]
\PYG{c+c1}{// from force.h}
\PYG{k}{typedef}\PYG{+w}{ }\PYG{n}{Pair}\PYG{+w}{ }\PYG{o}{*}\PYG{p}{(}\PYG{o}{*}\PYG{n}{PairCreator}\PYG{p}{)}\PYG{p}{(}\PYG{n}{LAMMPS}\PYG{+w}{ }\PYG{o}{*}\PYG{p}{)}\PYG{p}{;}
\PYG{k}{typedef}\PYG{+w}{ }\PYG{n}{std}\PYG{o}{:}\PYG{o}{:}\PYG{n}{map}\PYG{o}{\PYGZlt{}}\PYG{n}{std}\PYG{o}{:}\PYG{o}{:}\PYG{n}{string}\PYG{p}{,}\PYG{+w}{ }\PYG{n}{PairCreator}\PYG{o}{\PYGZgt{}}\PYG{+w}{ }\PYG{n}{PairCreatorMap}\PYG{p}{;}
\PYG{n}{PairCreatorMap}\PYG{+w}{ }\PYG{o}{*}\PYG{n}{pair\PYGZus{}map}\PYG{p}{;}

\PYG{c+c1}{// from force.cpp}
\PYG{k}{template}\PYG{+w}{ }\PYG{o}{\PYGZlt{}}\PYG{k}{typename}\PYG{+w}{ }\PYG{n+nc}{S}\PYG{p}{,}\PYG{+w}{ }\PYG{k}{typename}\PYG{+w}{ }\PYG{n+nc}{T}\PYG{o}{\PYGZgt{}}\PYG{+w}{ }\PYG{k}{static}\PYG{+w}{ }\PYG{n}{S}\PYG{+w}{ }\PYG{o}{*}\PYG{n}{style\PYGZus{}creator}\PYG{p}{(}\PYG{n}{LAMMPS}\PYG{+w}{ }\PYG{o}{*}\PYG{n}{lmp}\PYG{p}{)}
\PYG{p}{\PYGZob{}}
\PYG{+w}{  }\PYG{k}{return}\PYG{+w}{ }\PYG{k}{new}\PYG{+w}{ }\PYG{n}{T}\PYG{p}{(}\PYG{n}{lmp}\PYG{p}{)}\PYG{p}{;}
\PYG{p}{\PYGZcb{}}

\PYG{c+c1}{// [...]}

\PYG{n}{pair\PYGZus{}map}\PYG{+w}{ }\PYG{o}{=}\PYG{+w}{ }\PYG{k}{new}\PYG{+w}{ }\PYG{n}{PairCreatorMap}\PYG{p}{(}\PYG{p}{)}\PYG{p}{;}

\PYG{c+cp}{\PYGZsh{}}\PYG{c+cp}{define PAIR\PYGZus{}CLASS}
\PYG{c+cp}{\PYGZsh{}}\PYG{c+cp}{define PairStyle(key, Class) (*pair\PYGZus{}map)[\PYGZsh{}key] = \PYGZam{}style\PYGZus{}creator\PYGZlt{}Pair, Class\PYGZgt{};}
\PYG{c+cp}{\PYGZsh{}}\PYG{c+cp}{include}\PYG{+w}{ }\PYG{c+cpf}{\PYGZdq{}style\PYGZus{}pair.h\PYGZdq{}}
\PYG{c+cp}{\PYGZsh{}}\PYG{c+cp}{undef PairStyle}
\PYG{c+cp}{\PYGZsh{}}\PYG{c+cp}{undef PAIR\PYGZus{}CLASS}

\PYG{c+c1}{// from pair\PYGZus{}lj\PYGZus{}cut.h}

\PYG{c+cp}{\PYGZsh{}}\PYG{c+cp}{ifdef PAIR\PYGZus{}CLASS}
\PYG{n}{PairStyle}\PYG{p}{(}\PYG{n}{lj}\PYG{o}{/}\PYG{n}{cut}\PYG{p}{,}\PYG{n}{PairLJCut}\PYG{p}{)}\PYG{p}{;}
\PYG{c+cp}{\PYGZsh{}}\PYG{c+cp}{else}
\PYG{c+c1}{// [...]}
\end{sphinxVerbatim}

\sphinxAtStartPar
Similar code constructs are present in other files like \sphinxcode{\sphinxupquote{modify.cpp}} and
\sphinxcode{\sphinxupquote{modify.h}} or \sphinxcode{\sphinxupquote{neighbor.cpp}} and \sphinxcode{\sphinxupquote{neighbor.h}}.  Those contain
similar macros and include \sphinxcode{\sphinxupquote{style\_*.h}} files for creating class instances
of styles they manage.


\subsection{I/O and output formatting}
\label{\detokenize{Developer_code_design:i-o-and-output-formatting}}

\subsubsection{C\sphinxhyphen{}style stdio versus C++ style iostreams}
\label{\detokenize{Developer_code_design:c-style-stdio-versus-c-style-iostreams}}
\sphinxAtStartPar
LAMMPS uses the “stdio” library of the standard C library for reading
from and writing to files and console instead of C++ “iostreams”.
This is mainly motivated by better performance, better control over
formatting, and less effort to achieve specific formatting.

\sphinxAtStartPar
Since mixing “stdio” and “iostreams” can lead to unexpected behavior,
use of the latter is strongly discouraged.  Output to the screen should
\sphinxstyleemphasis{not} use the predefined \sphinxcode{\sphinxupquote{stdout}} FILE pointer, but rather the
\sphinxcode{\sphinxupquote{screen}} and \sphinxcode{\sphinxupquote{logfile}} FILE pointers managed by the LAMMPS class.
Furthermore, output should generally only be done by MPI rank 0
(\sphinxcode{\sphinxupquote{comm\sphinxhyphen{}\textgreater{}me == 0}}).  Output that is sent to both \sphinxcode{\sphinxupquote{screen}} and
\sphinxcode{\sphinxupquote{logfile}} should use the {\hyperref[\detokenize{Developer_utils:_CPPv4IDpEN9LAMMPS_NS5utils7logmesgEvP6LAMMPSRKNSt6stringEDpRR4Args}]{\sphinxcrossref{\sphinxcode{\sphinxupquote{utils::logmesg() convenience
function}}}}}.

\sphinxAtStartPar
We discourage the use of stringstreams because the bundled \{fmt\} library
and the customized tokenizer classes provide the same functionality in a
cleaner way with better performance.  This also helps maintain a
consistent programming syntax with code from many different
contributors.


\subsubsection{Formatting with the \{fmt\} library}
\label{\detokenize{Developer_code_design:formatting-with-the-fmt-library}}
\sphinxAtStartPar
The LAMMPS source code includes a copy of the \sphinxhref{https://fmt.dev}{\{fmt\} library}, which is preferred over formatting with the
“printf()” family of functions.  The primary reason is that it allows
a typesafe default format for any type of supported data.  This is
particularly useful for formatting integers of a given size (32\sphinxhyphen{}bit or
64\sphinxhyphen{}bit) which may require different format strings depending on
compile time settings or compilers/operating systems.  Furthermore,
\{fmt\} gives better performance, has more functionality, a familiar
formatting syntax that has similarities to \sphinxcode{\sphinxupquote{format()}} in Python, and
provides a facility that can be used to integrate format strings and a
variable number of arguments into custom functions in a much simpler
way than the varargs mechanism of the C library.  Finally, \{fmt\} has
been included into the C++20 language standard, so changes to adopt it
are future\sphinxhyphen{}proof.

\sphinxAtStartPar
Formatted strings are frequently created by calling the
\sphinxcode{\sphinxupquote{fmt::format()}} function, which will return a string as a
\sphinxcode{\sphinxupquote{std::string}} class instance.  In contrast to the \sphinxcode{\sphinxupquote{\%}} placeholder in
\sphinxcode{\sphinxupquote{printf()}}, the \{fmt\} library uses \sphinxcode{\sphinxupquote{\{\}}} to embed format descriptors.
In the simplest case, no additional characters are needed, as \{fmt\} will
choose the default format based on the data type of the argument.
Otherwise, the \sphinxcode{\sphinxupquote{fmt::print()}} function may be used instead of
\sphinxcode{\sphinxupquote{printf()}} or \sphinxcode{\sphinxupquote{fprintf()}}.  In addition, several LAMMPS output
functions, that originally accepted a single string as argument have
been overloaded to accept a format string with optional arguments as
well (e.g., \sphinxcode{\sphinxupquote{Error::all()}}, \sphinxcode{\sphinxupquote{Error::one()}}, \sphinxcode{\sphinxupquote{utils::logmesg()}}).


\subsubsection{Summary of the \{fmt\} format syntax}
\label{\detokenize{Developer_code_design:summary-of-the-fmt-format-syntax}}
\sphinxAtStartPar
The syntax of the format string is “\{{[}\textless{}argument id\textgreater{}{]}{[}:\textless{}format spec\textgreater{}{]}\}”,
where either the argument id or the format spec (separated by a colon
‘:’) is optional.  The argument id is usually a number starting from 0
that is the index to the arguments following the format string.  By
default, these are assigned in order (i.e. 0, 1, 2, 3, 4 etc.).  The
most common case for using argument id would be to use the same argument
in multiple places in the format string without having to provide it as
an argument multiple times. The argument id is rarely used in the LAMMPS
source code.

\sphinxAtStartPar
More common is the use of a format specifier, which starts with a colon.
This may optionally be followed by a fill character (default is ‘ ‘). If
provided, the fill character \sphinxstylestrong{must} be followed by an alignment
character (‘\textless{}’, ‘\textasciicircum{}’, ‘\textgreater{}’ for left, centered, or right alignment
(default)).  The alignment character may be used without a fill
character.  The next important format parameter would be the minimum
width, which may be followed by a dot ‘.’ and a precision for floating
point numbers.  The final character in the format string would be an
indicator for the “presentation”, i.e. ‘d’ for decimal presentation of
integers, ‘x’ for hexadecimal, ‘o’ for octal, ‘c’ for character etc.
This mostly follows the “printf()” scheme, but without requiring an
additional length parameter to distinguish between different integer
widths.  The \{fmt\} library will detect those and adapt the formatting
accordingly.  For floating point numbers there are correspondingly, ‘g’
for generic presentation, ‘e’ for exponential presentation, and ‘f’ for
fixed point presentation.

\sphinxAtStartPar
The format string “\{:8\}” would thus represent \sphinxstyleemphasis{any} type argument and be
replaced by at least 8 characters; “\{:\textless{}8\}” would do this as left
aligned, “\{:\textasciicircum{}8\}” as centered, “\{:\textgreater{}8\}” as right aligned.  If a specific
presentation is selected, the argument type must be compatible or else
the \{fmt\} formatting code will throw an exception.  Some format string
examples are given below:

\begin{sphinxVerbatim}[commandchars=\\\{\}]
\PYG{k}{auto}\PYG{+w}{ }\PYG{n}{mesg}\PYG{+w}{ }\PYG{o}{=}\PYG{+w}{ }\PYG{n}{fmt}\PYG{o}{:}\PYG{o}{:}\PYG{n}{format}\PYG{p}{(}\PYG{l+s}{\PYGZdq{}}\PYG{l+s}{  CPU time: \PYGZob{}:4d\PYGZcb{}:\PYGZob{}:02d\PYGZcb{}:\PYGZob{}:02d\PYGZcb{}}\PYG{l+s+se}{\PYGZbs{}n}\PYG{l+s}{\PYGZdq{}}\PYG{p}{,}\PYG{+w}{ }\PYG{n}{cpuh}\PYG{p}{,}\PYG{+w}{ }\PYG{n}{cpum}\PYG{p}{,}\PYG{+w}{ }\PYG{n}{cpus}\PYG{p}{)}\PYG{p}{;}
\PYG{n}{mesg}\PYG{+w}{ }\PYG{o}{=}\PYG{+w}{ }\PYG{n}{fmt}\PYG{o}{:}\PYG{o}{:}\PYG{n}{format}\PYG{p}{(}\PYG{l+s}{\PYGZdq{}}\PYG{l+s}{\PYGZob{}:\PYGZlt{}8s\PYGZcb{}| \PYGZob{}:\PYGZlt{}10.5g\PYGZcb{} | \PYGZob{}:\PYGZlt{}10.5g\PYGZcb{} | \PYGZob{}:\PYGZlt{}10.5g\PYGZcb{} |\PYGZob{}:6.1f\PYGZcb{} |\PYGZob{}:6.2f\PYGZcb{}}\PYG{l+s+se}{\PYGZbs{}n}\PYG{l+s}{\PYGZdq{}}\PYG{p}{,}
\PYG{+w}{                   }\PYG{n}{label}\PYG{p}{,}\PYG{+w}{ }\PYG{n}{time\PYGZus{}min}\PYG{p}{,}\PYG{+w}{ }\PYG{n}{time}\PYG{p}{,}\PYG{+w}{ }\PYG{n}{time\PYGZus{}max}\PYG{p}{,}\PYG{+w}{ }\PYG{n}{time\PYGZus{}sq}\PYG{p}{,}\PYG{+w}{ }\PYG{n}{tmp}\PYG{p}{)}\PYG{p}{;}
\PYG{n}{utils}\PYG{o}{:}\PYG{o}{:}\PYG{n}{logmesg}\PYG{p}{(}\PYG{n}{lmp}\PYG{p}{,}\PYG{l+s}{\PYGZdq{}}\PYG{l+s}{\PYGZob{}:\PYGZgt{}6\PYGZcb{} = max \PYGZsh{} of 1\PYGZhy{}2 neighbors}\PYG{l+s+se}{\PYGZbs{}n}\PYG{l+s}{\PYGZdq{}}\PYG{p}{,}\PYG{n}{maxall}\PYG{p}{)}\PYG{p}{;}
\PYG{n}{utils}\PYG{o}{:}\PYG{o}{:}\PYG{n}{logmesg}\PYG{p}{(}\PYG{n}{lmp}\PYG{p}{,}\PYG{l+s}{\PYGZdq{}}\PYG{l+s}{Lattice spacing in x,y,z = \PYGZob{}:.8\PYGZcb{} \PYGZob{}:.8\PYGZcb{} \PYGZob{}:.8\PYGZcb{}}\PYG{l+s+se}{\PYGZbs{}n}\PYG{l+s}{\PYGZdq{}}\PYG{p}{,}
\PYG{+w}{               }\PYG{n}{xlattice}\PYG{p}{,}\PYG{n}{ylattice}\PYG{p}{,}\PYG{n}{zlattice}\PYG{p}{)}\PYG{p}{;}
\end{sphinxVerbatim}

\sphinxAtStartPar
which will create the following output lines:

\begin{sphinxVerbatim}[commandchars=\\\{\}]
CPU time:    0:02:16
Pair    | 2.0133     | 2.0133     | 2.0133     |   0.0 | 84.21
     4 = max \PYGZsh{} of 1\PYGZhy{}2 neighbors
Lattice spacing in x,y,z = 1.6795962 1.6795962 1.6795962
\end{sphinxVerbatim}

\sphinxAtStartPar
Finally, a special feature of the \{fmt\} library is that format
parameters like the width or the precision may be also provided as
arguments. In that case a nested format is used where a pair of curly
braces (with an optional argument id) “\{\}” are used instead of the
value, for example “\{:\{\}d\}” will consume two integer arguments, the
first will be the value shown and the second the minimum width.

\sphinxAtStartPar
For more details and examples, please consult the \sphinxhref{https://fmt.dev/latest/syntax.html}{\{fmt\} syntax
documentation} website.


\subsection{Memory management}
\label{\detokenize{Developer_code_design:memory-management}}
\sphinxAtStartPar
Dynamical allocation of small data and objects can be done with the C++
commands “new” and “delete/delete{[}{]}”.  Large data should use the member
functions of the \sphinxcode{\sphinxupquote{Memory}} class, most commonly, \sphinxcode{\sphinxupquote{Memory::create()}},
\sphinxcode{\sphinxupquote{Memory::grow()}}, and \sphinxcode{\sphinxupquote{Memory::destroy()}}, which provide variants
for vectors, 2d arrays, 3d arrays, etc.  These can also be used for
small data.

\sphinxAtStartPar
The use of \sphinxcode{\sphinxupquote{malloc()}}, \sphinxcode{\sphinxupquote{calloc()}}, \sphinxcode{\sphinxupquote{realloc()}} and \sphinxcode{\sphinxupquote{free()}}
directly is strongly discouraged.  To simplify adapting legacy code
into the LAMMPS code base the member functions \sphinxcode{\sphinxupquote{Memory::smalloc()}},
\sphinxcode{\sphinxupquote{Memory::srealloc()}}, and \sphinxcode{\sphinxupquote{Memory::sfree()}} are available, which
perform additional error checks for safety.

\sphinxAtStartPar
Use of these custom memory allocation functions is motivated by the
following considerations:
\begin{itemize}
\item {} 
\sphinxAtStartPar
Memory allocation failures on \sphinxstyleemphasis{any} MPI rank during a parallel run
will trigger an immediate abort of the entire parallel calculation.

\item {} 
\sphinxAtStartPar
A failing “new” will trigger an exception, which is also captured by
LAMMPS and triggers a global abort.

\item {} 
\sphinxAtStartPar
Allocation of multidimensional arrays will be done in a C compatible
fashion, but such that the storage of the actual data is stored in one
large contiguous block.  Thus, when MPI communication is needed,
the data can be communicated directly (similar to Fortran arrays).

\item {} 
\sphinxAtStartPar
The “destroy()” and “sfree()” functions may safely be called on NULL
pointers.

\item {} 
\sphinxAtStartPar
The “destroy()” functions will nullify the pointer variables, thus
making “use after free” errors easy to detect.

\item {} 
\sphinxAtStartPar
It is possible to use a larger than default memory alignment (not on
all operating systems, since the allocated storage pointers must be
compatible with \sphinxcode{\sphinxupquote{free()}} for technical reasons).

\end{itemize}

\sphinxAtStartPar
In the practical implementation of code this means, that any pointer
variables, that are class members should be initialized to a \sphinxcode{\sphinxupquote{nullptr}}
value in their respective constructors.  That way, it is safe to call
\sphinxcode{\sphinxupquote{Memory::destroy()}} or \sphinxcode{\sphinxupquote{delete{[}{]}}} on them before \sphinxstyleemphasis{any} allocation
outside the constructor.  This helps prevent memory leaks.

\sphinxstepscope


\section{Parallel algorithms}
\label{\detokenize{Developer_parallel:parallel-algorithms}}\label{\detokenize{Developer_parallel::doc}}
\sphinxAtStartPar
LAMMPS is designed to enable running simulations in parallel using the
MPI parallel communication standard with distributed data via domain
decomposition.  The parallelization aims to be efficient, and resulting
in good strong scaling (= good speedup for the same system) and good
weak scaling (= the computational cost of enlarging the system is
proportional to the system size).  Additional parallelization using GPUs
or OpenMP can also be applied within the subdomain assigned to an MPI
process.  For clarity, most of the following illustrations show the 2d
simulation case. The underlying algorithms in those cases, however,
apply to both 2d and 3d cases equally well.

\begin{sphinxadmonition}{note}{Note:}
\sphinxAtStartPar
The text and most of the figures in this chapter were adapted
for the manual from the section on parallel algorithms in the
{\hyperref[\detokenize{Intro_citing:lammps-paper}]{\sphinxcrossref{\DUrole{std,std-ref}{new LAMMPS paper}}}}.
\end{sphinxadmonition}

\sphinxstepscope


\subsection{Partitioning}
\label{\detokenize{Developer_par_part:partitioning}}\label{\detokenize{Developer_par_part::doc}}
\sphinxAtStartPar
The underlying spatial decomposition strategy used by LAMMPS for
distributed\sphinxhyphen{}memory parallelism is set with the {\hyperref[\detokenize{comm_style::doc}]{\sphinxcrossref{\DUrole{doc}{comm\_style command}}}} and can be either “brick” (a regular grid) or “tiled”.

\begin{figure}[htbp]
\centering
\capstart

\noindent\sphinxincludegraphics{{domain-decomp}.png}
\caption{Domain decomposition schemes}\label{\detokenize{Developer_par_part:id1}}\label{\detokenize{Developer_par_part:domain-decomposition}}
\begin{sphinxlegend}\begin{quote}

\sphinxAtStartPar
This figure shows the different kinds of domain decomposition used
for MPI parallelization: “brick” on the left with an orthogonal
(left) and a triclinic (middle) simulation domain, and a “tiled”
decomposition (right).  The black lines show the division into
subdomains, and the contained atoms are “owned” by the
corresponding MPI process. The green dashed lines indicate how
subdomains are extended with “ghost” atoms up to the communication
cutoff distance.
\end{quote}
\end{sphinxlegend}
\end{figure}

\sphinxAtStartPar
The LAMMPS simulation box is a 3d or 2d volume, which can be of
orthogonal or triclinic shape, as illustrated in the
{\hyperref[\detokenize{Developer_par_part:domain-decomposition}]{\sphinxcrossref{\DUrole{std,std-ref}{Domain decomposition schemes}}}} figure for the 2d case.  Orthogonal means
the box edges are aligned with the \sphinxstyleemphasis{x}, \sphinxstyleemphasis{y}, \sphinxstyleemphasis{z} Cartesian axes, and the
box faces are thus all rectangular.  Triclinic allows for a more general
parallelepiped shape in which edges are aligned with three arbitrary
vectors and the box faces are parallelograms.  In each dimension, box
faces can be periodic, or non\sphinxhyphen{}periodic with fixed or shrink\sphinxhyphen{}wrapped
boundaries.  In the fixed case, atoms which move outside the face are
deleted; shrink\sphinxhyphen{}wrapped means the position of the box face adjusts
continuously to enclose all the atoms.

\sphinxAtStartPar
For distributed\sphinxhyphen{}memory MPI parallelism, the simulation box is spatially
decomposed (partitioned) into non\sphinxhyphen{}overlapping subdomains which fill the
box. The default partitioning, “brick”, is most suitable when atom
density is roughly uniform, as shown in the left\sphinxhyphen{}side images of the
{\hyperref[\detokenize{Developer_par_part:domain-decomposition}]{\sphinxcrossref{\DUrole{std,std-ref}{Domain decomposition schemes}}}} figure.  The subdomains comprise a regular
grid, and all subdomains are identical in size and shape.  Both the
orthogonal and triclinic boxes can deform continuously during a
simulation, e.g. to compress a solid or shear a liquid, in which case
the processor subdomains likewise deform.

\sphinxAtStartPar
For models with non\sphinxhyphen{}uniform density, the number of particles per
processor can be load\sphinxhyphen{}imbalanced with the default partitioning.  This
reduces parallel efficiency, as the overall simulation rate is limited
by the slowest processor, i.e. the one with the largest computational
load.  For such models, LAMMPS supports multiple strategies to reduce
the load imbalance:
\begin{itemize}
\item {} 
\sphinxAtStartPar
The processor grid decomposition is by default based on the simulation
cell volume and tries to optimize the volume to surface ratio for the subdomains.
This can be changed with the {\hyperref[\detokenize{processors::doc}]{\sphinxcrossref{\DUrole{doc}{processors command}}}}.

\item {} 
\sphinxAtStartPar
The parallel planes defining the size of the subdomains can be shifted
with the {\hyperref[\detokenize{balance::doc}]{\sphinxcrossref{\DUrole{doc}{balance command}}}}. Which can be done in addition
to choosing a more optimal processor grid.

\item {} 
\sphinxAtStartPar
The recursive bisectioning algorithm in combination with the “tiled”
communication style can produce a partitioning with equal numbers of
particles in each subdomain.

\end{itemize}

\sphinxAtStartPar
\sphinxincludegraphics[width=0.240\linewidth]{{decomp-regular}.png}  \sphinxincludegraphics[width=0.240\linewidth]{{decomp-processors}.png}  \sphinxincludegraphics[width=0.240\linewidth]{{decomp-balance}.png}  \sphinxincludegraphics[width=0.240\linewidth]{{decomp-rcb}.png}

\sphinxAtStartPar
The pictures above demonstrate different decompositions for a 2d system
with 12 MPI ranks.  The atom colors indicate the load imbalance of each
subdomain, with green being optimal and red the least optimal.

\sphinxAtStartPar
Due to the vacuum in the system, the default decomposition is
unbalanced, with several MPI ranks without atoms (left). By forcing a
1x12x1 processor grid, every MPI rank does computations now, but the
number of atoms per subdomain is still uneven, and the thin slice shape
increases the amount of communication between subdomains (center
left). With a 2x6x1 processor grid and shifting the subdomain divisions,
the load imbalance is further reduced and the amount of communication
required between subdomains is less (center right).  And using the
recursive bisectioning leads to further improved decomposition (right).

\sphinxstepscope


\subsection{Communication}
\label{\detokenize{Developer_par_comm:communication}}\label{\detokenize{Developer_par_comm::doc}}
\sphinxAtStartPar
Following the selected partitioning scheme, all per\sphinxhyphen{}atom data is
distributed across the MPI processes, which allows LAMMPS to handle very
large systems provided it uses a correspondingly large number of MPI
processes.  To be able to compute the short\sphinxhyphen{}range interactions, MPI
processes need not only access to the data of atoms they “own” but also
information about atoms from neighboring subdomains, in LAMMPS referred
to as “ghost” atoms.  These are copies of atoms storing required
per\sphinxhyphen{}atom data for up to the communication cutoff distance. The green
dashed\sphinxhyphen{}line boxes in the {\hyperref[\detokenize{Developer_par_part:domain-decomposition}]{\sphinxcrossref{\DUrole{std,std-ref}{Domain decomposition schemes}}}} figure illustrate
the extended ghost\sphinxhyphen{}atom subdomain for one processor.

\sphinxAtStartPar
This approach is also used to implement periodic boundary
conditions: atoms that lie within the cutoff distance across a periodic
boundary are also stored as ghost atoms and taken from the periodic
replication of the subdomain, which may be the same subdomain, e.g. if
running in serial.  As a consequence of this, force computation in
LAMMPS is not subject to minimum image conventions and thus cutoffs may
be larger than half the simulation domain.

\begin{figure}[htbp]
\centering
\capstart

\noindent\sphinxincludegraphics{{ghost-comm}.png}
\caption{ghost atom communication}\label{\detokenize{Developer_par_comm:id1}}\label{\detokenize{Developer_par_comm:ghost-atom-comm}}
\begin{sphinxlegend}
\sphinxAtStartPar
This figure shows the ghost atom communication patterns between
subdomains for “brick” (left) and “tiled” communication styles for
2d simulations.  The numbers indicate MPI process ranks.  Here the
subdomains are drawn spatially separated for clarity.  The
dashed\sphinxhyphen{}line box is the extended subdomain of processor 0 which
includes its ghost atoms.  The red\sphinxhyphen{} and blue\sphinxhyphen{}shaded boxes are the
regions of communicated ghost atoms.
\end{sphinxlegend}
\end{figure}

\sphinxAtStartPar
Efficient communication patterns are needed to update the “ghost” atom
data, since that needs to be done at every MD time step or minimization
step.  The diagrams of the {\hyperref[\detokenize{Developer_par_comm:ghost-atom-comm}]{\sphinxcrossref{\DUrole{std,std-ref}{ghost atom communication}}}} figure illustrate how ghost
atom communication is performed in two stages for a 2d simulation (three
in 3d) for both a regular and irregular partitioning of the simulation
box.  For the regular case (left) atoms are exchanged first in the
\sphinxstyleemphasis{x}\sphinxhyphen{}direction, then in \sphinxstyleemphasis{y}, with four neighbors in the grid of processor
subdomains.

\sphinxAtStartPar
In the \sphinxstyleemphasis{x} stage, processor ranks 1 and 2 send owned atoms in their
red\sphinxhyphen{}shaded regions to rank 0 (and vice versa).  Then in the \sphinxstyleemphasis{y} stage,
ranks 3 and 4 send atoms in their blue\sphinxhyphen{}shaded regions to rank 0, which
includes ghost atoms they received in the \sphinxstyleemphasis{x} stage.  Rank 0 thus
acquires all its ghost atoms; atoms in the solid blue corner regions
are communicated twice before rank 0 receives them.

\sphinxAtStartPar
For the irregular case (right) the two stages are similar, but a
processor can have more than one neighbor in each direction.  In the
\sphinxstyleemphasis{x} stage, MPI ranks 1,2,3 send owned atoms in their red\sphinxhyphen{}shaded regions to
rank 0 (and vice versa).  These include only atoms between the lower
and upper \sphinxstyleemphasis{y}\sphinxhyphen{}boundary of rank 0’s subdomain.  In the \sphinxstyleemphasis{y} stage, ranks
4,5,6 send atoms in their blue\sphinxhyphen{}shaded regions to rank 0.  This may
include ghost atoms they received in the \sphinxstyleemphasis{x} stage, but only if they
are needed by rank 0 to fill its extended ghost atom regions in the
+/\sphinxhyphen{}\sphinxstyleemphasis{y} directions (blue rectangles).  Thus, in this case, ranks 5 and
6 do not include ghost atoms they received from each other (in the \sphinxstyleemphasis{x}
stage) in the atoms they send to rank 0.  The key point is that while
the pattern of communication is more complex in the irregular
partitioning case, it can still proceed in two stages (three in 3d)
via atom exchanges with only neighboring processors.

\sphinxAtStartPar
When attributes of owned atoms are sent to neighboring processors to
become attributes of their ghost atoms, LAMMPS calls this a “forward”
communication.  On timesteps when atoms migrate to new owning processors
and neighbor lists are rebuilt, each processor creates a list of its
owned atoms which are ghost atoms in each of its neighbor processors.
These lists are used to pack per\sphinxhyphen{}atom coordinates (for example) into
message buffers in subsequent steps until the next reneighboring.

\sphinxAtStartPar
A “reverse” communication is when computed ghost atom attributes are
sent back to the processor who owns the atom.  This is used (for
example) to sum partial forces on ghost atoms to the complete force on
owned atoms.  The order of the two stages described in the
{\hyperref[\detokenize{Developer_par_comm:ghost-atom-comm}]{\sphinxcrossref{\DUrole{std,std-ref}{ghost atom communication}}}} figure is inverted, and the same lists of atoms
are used to pack and unpack message buffers with per\sphinxhyphen{}atom forces.  When
a received buffer is unpacked, the ghost forces are summed to owned atom
forces.  As in forward communication, forces on atoms in the four blue
corners of the diagrams are sent, received, and summed twice (once at
each stage) before owning processors have the full force.

\sphinxAtStartPar
These two operations are used in many places within LAMMPS aside from
exchange of coordinates and forces, for example by manybody potentials
to share intermediate per\sphinxhyphen{}atom values, or by rigid\sphinxhyphen{}body integrators to
enable each atom in a body to access body properties.  Here are
additional details about how these communication operations are
performed in LAMMPS:
\begin{itemize}
\item {} 
\sphinxAtStartPar
When exchanging data with different processors, forward and reverse
communication is done using \sphinxcode{\sphinxupquote{MPI\_Send()}} and \sphinxcode{\sphinxupquote{MPI\_IRecv()}} calls.
If a processor is “exchanging” atoms with itself, only the pack and
unpack operations are performed, e.g. to create ghost atoms across
periodic boundaries when running on a single processor.

\item {} 
\sphinxAtStartPar
For forward communication of owned atom coordinates, periodic box
lengths are added and subtracted when the receiving processor is
across a periodic boundary from the sender.  There is then no need to
apply a minimum image convention when calculating distances between
atom pairs when building neighbor lists or computing forces.

\item {} 
\sphinxAtStartPar
The cutoff distance for exchanging ghost atoms is typically equal to
the neighbor cutoff.  But it can also set to a larger value if needed,
e.g. half the diameter of a rigid body composed of multiple atoms or
over 3x the length of a stretched bond for dihedral interactions.  It
can also exceed the periodic box size.  For the regular communication
pattern (left), if the cutoff distance extends beyond a neighbor
processor’s subdomain, then multiple exchanges are performed in the
same direction.  Each exchange is with the same neighbor processor,
but buffers are packed/unpacked using a different list of atoms. For
forward communication, in the first exchange, a processor sends only
owned atoms.  In subsequent exchanges, it sends ghost atoms received
in previous exchanges.  For the irregular pattern (right) overlaps of
a processor’s extended ghost\sphinxhyphen{}atom subdomain with all other processors
in each dimension are detected.

\end{itemize}

\sphinxstepscope


\subsection{Neighbor lists}
\label{\detokenize{Developer_par_neigh:neighbor-lists}}\label{\detokenize{Developer_par_neigh::doc}}
\sphinxAtStartPar
To compute forces efficiently, each processor creates a Verlet\sphinxhyphen{}style
neighbor list which enumerates all pairs of atoms \sphinxstyleemphasis{i,j} (\sphinxstyleemphasis{i} = owned,
\sphinxstyleemphasis{j} = owned or ghost) with separation less than the applicable neighbor
list cutoff distance.  In LAMMPS, the neighbor lists are stored in a
multiple\sphinxhyphen{}page data structure; each page is a contiguous chunk of memory
which stores vectors of neighbor atoms \sphinxstyleemphasis{j} for many \sphinxstyleemphasis{i} atoms.  This
allows pages to be incrementally allocated or deallocated in blocks as
needed.  Neighbor lists typically consume the most memory of any data
structure in LAMMPS.  The neighbor list is rebuilt (from scratch) once
every few timesteps, then used repeatedly each step for force or other
computations.  The neighbor cutoff distance is \(R_n = R_f +
\Delta_s\), where \(R_f\) is the (largest) force cutoff defined by
the interatomic potential for computing short\sphinxhyphen{}range pairwise or manybody
forces and \(\Delta_s\) is a “skin” distance that allows the list to
be used for multiple steps assuming that atoms do not move very far
between consecutive time steps.  Typically, the code triggers
reneighboring when any atom has moved half the skin distance since the
last reneighboring; this and other options of the neighbor list rebuild
can be adjusted with the {\hyperref[\detokenize{neigh_modify::doc}]{\sphinxcrossref{\DUrole{doc}{neigh\_modify}}}} command.

\sphinxAtStartPar
On steps when reneighboring is performed, atoms which have moved outside
their owning processor’s subdomain are first migrated to new processors
via communication.  Periodic boundary conditions are also (only)
enforced on these steps to ensure each atom is re\sphinxhyphen{}assigned to the
correct processor.  After migration, the atoms owned by each processor
are stored in a contiguous vector.  Periodically, each processor
spatially sorts owned atoms within its vector to reorder it for improved
cache efficiency in force computations and neighbor list building.  For
that, atoms are spatially binned and then reordered so that atoms in the
same bin are adjacent in the vector.  Atom sorting can be disabled or
its settings modified with the {\hyperref[\detokenize{atom_modify::doc}]{\sphinxcrossref{\DUrole{doc}{atom\_modify}}}} command.

\begin{figure}[htbp]
\centering
\capstart

\noindent\sphinxincludegraphics{{neigh-stencil}.png}
\caption{neighbor list stencils}\label{\detokenize{Developer_par_neigh:id1}}\label{\detokenize{Developer_par_neigh:neighbor-stencil}}
\begin{sphinxlegend}
\sphinxAtStartPar
A 2d simulation subdomain (thick black line) and the corresponding
ghost atom cutoff region (dashed blue line) for both orthogonal
(left) and triclinic (right) domains.  A regular grid of neighbor
bins (thin lines) overlays the entire simulation domain and need not
align with subdomain boundaries; only the portion overlapping the
augmented subdomain is shown.  In the triclinic case, it overlaps the
bounding box of the tilted rectangle.  The blue\sphinxhyphen{} and red\sphinxhyphen{}shaded bins
represent a stencil of bins searched to find neighbors of a particular
atom (black dot).
\end{sphinxlegend}
\end{figure}

\sphinxAtStartPar
To build a local neighbor list in linear time, the simulation domain is
overlaid (conceptually) with a regular 3d (or 2d) grid of neighbor bins,
as shown in the {\hyperref[\detokenize{Developer_par_neigh:neighbor-stencil}]{\sphinxcrossref{\DUrole{std,std-ref}{neighbor list stencils}}}} figure for 2d models and a
single MPI processor’s subdomain.  Each processor stores a set of
neighbor bins which overlap its subdomain, extended by the neighbor
cutoff distance \(R_n\).  As illustrated, the bins need not align
with processor boundaries; an integer number in each dimension is fit to
the size of the entire simulation box.

\sphinxAtStartPar
Most often, LAMMPS builds what is called a “half” neighbor list where
each \sphinxstyleemphasis{i,j} neighbor pair is stored only once, with either atom \sphinxstyleemphasis{i} or
\sphinxstyleemphasis{j} as the central atom.  The build can be done efficiently by using a
pre\sphinxhyphen{}computed “stencil” of bins around a central origin bin which
contains the atom whose neighbors are being searched for.  A stencil
is simply a list of integer offsets in \sphinxstyleemphasis{x,y,z} of nearby bins
surrounding the origin bin which are close enough to contain any
neighbor atom \sphinxstyleemphasis{j} within a distance \(R_n\) from any atom \sphinxstyleemphasis{i} in the
origin bin.  Note that for a half neighbor list, the stencil can be
asymmetric, since each atom only need store half its nearby neighbors.

\sphinxAtStartPar
These stencils are illustrated in the figure for a half list and a bin
size of \(\frac{1}{2} R_n\).  There are 13 red+blue stencil bins in
2d (for the orthogonal case, 15 for triclinic).  In 3d there would be
63, 13 in the plane of bins that contain the origin bin and 25 in each
of the two planes above it in the \sphinxstyleemphasis{z} direction (75 for triclinic).  The
triclinic stencil has extra bins because the bins tile the bounding box
of the entire triclinic domain, and thus are not periodic with respect
to the simulation box itself.  The stencil and logic for determining
which \sphinxstyleemphasis{i,j} pairs to include in the neighbor list are altered slightly
to account for this.

\sphinxAtStartPar
To build a neighbor list, a processor first loops over its “owned” plus
“ghost” atoms and assigns each to a neighbor bin.  This uses an integer
vector to create a linked list of atom indices within each bin.  It then
performs a triply\sphinxhyphen{}nested loop over its owned atoms \sphinxstyleemphasis{i}, the stencil of
bins surrounding atom \sphinxstyleemphasis{i}’s bin, and the \sphinxstyleemphasis{j} atoms in each stencil bin
(including ghost atoms).  If the distance \(r_{ij} < R_n\), then
atom \sphinxstyleemphasis{j} is added to the vector of atom \sphinxstyleemphasis{i}’s neighbors.

\sphinxAtStartPar
Here are additional details about neighbor list build options LAMMPS
supports:
\begin{itemize}
\item {} 
\sphinxAtStartPar
The choice of bin size is an option; a size half of \(R_n\) has
been found to be optimal for many typical cases.  Smaller bins incur
additional overhead to loop over; larger bins require more distance
calculations.  Note that for smaller bin sizes, the 2d stencil in the
figure would be of a more semicircular shape (hemispherical in 3d),
with bins near the corners of the square eliminated due to their
distance from the origin bin.

\item {} 
\sphinxAtStartPar
Depending on the interatomic potential(s) and other commands used in
an input script, multiple neighbor lists and stencils with different
attributes may be needed.  This includes lists with different cutoff
distances, e.g. for force computation versus occasional diagnostic
computations such as a radial distribution function, or for the
r\sphinxhyphen{}RESPA time integrator which can partition pairwise forces by
distance into subsets computed at different time intervals.  It
includes “full” lists (as opposed to half lists) where each \sphinxstyleemphasis{i,j} pair
appears twice, stored once with \sphinxstyleemphasis{i} and \sphinxstyleemphasis{j}, and which use a larger
symmetric stencil.  It also includes lists with partial enumeration of
ghost atom neighbors.  The full and ghost\sphinxhyphen{}atom lists are used by
various manybody interatomic potentials.  Lists may also use different
criteria for inclusion of a pairwise interaction.  Typically, this
simply depends only on the distance between two atoms and the cutoff
distance.  But for finite\sphinxhyphen{}size coarse\sphinxhyphen{}grained particles with
individual diameters (e.g. polydisperse granular particles), it can
also depend on the diameters of the two particles.

\item {} 
\sphinxAtStartPar
When using {\hyperref[\detokenize{pair_hybrid::doc}]{\sphinxcrossref{\DUrole{doc}{pair style hybrid}}}} multiple sub\sphinxhyphen{}lists
of the master neighbor list for the full system need to be generated,
one for each sub\sphinxhyphen{}style, which contains only the \sphinxstyleemphasis{i,j} pairs needed to
compute interactions between subsets of atoms for the corresponding
potential.  This means, not all \sphinxstyleemphasis{i} or \sphinxstyleemphasis{j} atoms owned by a processor
are included in a particular sub\sphinxhyphen{}list.

\item {} 
\sphinxAtStartPar
Some models use different cutoff lengths for pairwise interactions
between different kinds of particles, which are stored in a single
neighbor list.  One example is a solvated colloidal system with large
colloidal particles where colloid/colloid, colloid/solvent, and
solvent/solvent interaction cutoffs can be dramatically different.
Another is a model of polydisperse finite\sphinxhyphen{}size granular particles;
pairs of particles interact only when they are in contact with each
other.  Mixtures with particle size ratios as high as 10\sphinxhyphen{}100x may be
used to model realistic systems.  Efficient neighbor list building
algorithms for these kinds of systems are available in LAMMPS.  They
include a method which uses different stencils for different cutoff
lengths and trims the stencil to only include bins that straddle the
cutoff sphere surface.  More recently a method which uses both
multiple stencils and multiple bin sizes was developed; it builds
neighbor lists efficiently for systems with particles of any size
ratio, though other considerations (timestep size, force computations)
may limit the ability to model systems with huge polydispersity.

\item {} 
\sphinxAtStartPar
For small and sparse systems and as a fallback method, LAMMPS also
supports neighbor list construction without binning by using a full
\(O(N^2)\) loop over all \sphinxstyleemphasis{i,j} atom pairs in a subdomain when
using the {\hyperref[\detokenize{neighbor::doc}]{\sphinxcrossref{\DUrole{doc}{neighbor nsq}}}} command.

\item {} 
\sphinxAtStartPar
Dependent on the “pair” setting of the {\hyperref[\detokenize{newton::doc}]{\sphinxcrossref{\DUrole{doc}{newton}}}} command,
the “half” neighbor lists may contain \sphinxstylestrong{all} pairs of atoms where
atom \sphinxstyleemphasis{j} is a ghost atom (i.e. when the newton pair setting is \sphinxstyleemphasis{off}).
For the newton pair \sphinxstyleemphasis{on} setting the atom \sphinxstyleemphasis{j} is only added to the
list if its \sphinxstyleemphasis{z} coordinate is larger, or if equal the \sphinxstyleemphasis{y} coordinate
is larger, and that is equal, too, the \sphinxstyleemphasis{x} coordinate is larger.  For
homogeneously dense systems, that will result in picking neighbors from
a same size sector in always the same direction relative to the
“owned” atom, and thus it should lead to similar length neighbor lists
and reduce the chance of a load imbalance.

\end{itemize}

\sphinxstepscope


\subsection{Long\sphinxhyphen{}range interactions}
\label{\detokenize{Developer_par_long:long-range-interactions}}\label{\detokenize{Developer_par_long::doc}}
\sphinxAtStartPar
For charged systems, LAMMPS can compute long\sphinxhyphen{}range Coulombic
interactions via the FFT\sphinxhyphen{}based particle\sphinxhyphen{}particle/particle\sphinxhyphen{}mesh (PPPM)
method implemented in {\hyperref[\detokenize{kspace_style::doc}]{\sphinxcrossref{\DUrole{doc}{kspace style pppm and its variants}}}}.  For that Coulombic interactions are partitioned into
short\sphinxhyphen{} and long\sphinxhyphen{}range components.  The short\sphinxhyphen{}ranged portion is computed
in real space as a loop over pairs of charges within a cutoff distance,
using neighbor lists.  The long\sphinxhyphen{}range portion is computed in reciprocal
space using a kspace style.  For the PPPM implementation the simulation
cell is overlaid with a regular FFT grid in 3d. It proceeds in several stages:
\begin{enumerate}
\sphinxsetlistlabels{\alph}{enumi}{enumii}{}{)}%
\item {} 
\sphinxAtStartPar
each atom’s point charge is interpolated to nearby FFT grid points,

\item {} 
\sphinxAtStartPar
a forward 3d FFT is performed,

\item {} 
\sphinxAtStartPar
a convolution operation is performed in reciprocal space,

\item {} 
\sphinxAtStartPar
one or more inverse 3d FFTs are performed, and

\item {} 
\sphinxAtStartPar
electric field values from grid points near each atom are interpolated to compute
its forces.

\end{enumerate}

\sphinxAtStartPar
For any of the spatial\sphinxhyphen{}decomposition partitioning schemes each processor
owns the brick\sphinxhyphen{}shaped portion of FFT grid points contained within its
subdomain.  The two interpolation operations use a stencil of grid
points surrounding each atom.  To accommodate the stencil size, each
processor also stores a few layers of ghost grid points surrounding its
brick.  Forward and reverse communication of grid point values is
performed similar to the corresponding {\hyperref[\detokenize{Developer_par_comm::doc}]{\sphinxcrossref{\DUrole{doc}{atom data communication}}}}.  In this case, electric field values on owned
grid points are sent to neighboring processors to become ghost point
values.  Likewise charge values on ghost points are sent and summed to
values on owned points.

\sphinxAtStartPar
For triclinic simulation boxes, the FFT grid planes are parallel to
the box faces, but the mapping of charge and electric field values
to/from grid points is done in reduced coordinates where the tilted
box is conceptually a unit cube, so that the stencil and FFT
operations are unchanged.  However the FFT grid size required for a
given accuracy is larger for triclinic domains than it is for
orthogonal boxes.

\begin{figure}[htbp]
\centering
\capstart

\noindent\sphinxincludegraphics{{fft-decomp-parallel}.png}
\caption{Parallel FFT in PPPM}\label{\detokenize{Developer_par_long:id1}}\label{\detokenize{Developer_par_long:fft-parallel}}
\begin{sphinxlegend}\begin{quote}

\sphinxAtStartPar
Stages of a parallel FFT for a simulation domain overlaid with an
8x8x8 3d FFT grid, partitioned across 64 processors.  Within each
of the 4 diagrams, grid cells of the same color are owned by a
single processor; for simplicity, only cells owned by 4 or 8 of
the 64 processors are colored.  The two images on the left
illustrate brick\sphinxhyphen{}to\sphinxhyphen{}pencil communication.  The two images on the
right illustrate pencil\sphinxhyphen{}to\sphinxhyphen{}pencil communication, which in this
case transposes the \sphinxstyleemphasis{y} and \sphinxstyleemphasis{z} dimensions of the grid.
\end{quote}
\end{sphinxlegend}
\end{figure}

\sphinxAtStartPar
Parallel 3d FFTs require substantial communication relative to their
computational cost.  A 3d FFT is implemented by a series of 1d FFTs
along the \sphinxstyleemphasis{x\sphinxhyphen{}}, \sphinxstyleemphasis{y\sphinxhyphen{}}, and \sphinxstyleemphasis{z\sphinxhyphen{}}direction of the FFT grid.  Thus, the
FFT grid cannot be decomposed like atoms into 3 dimensions for parallel
processing of the FFTs but only in 1 (as planes) or 2 (as pencils)
dimensions and in between the steps the grid needs to be transposed to
have the FFT grid portion “owned” by each MPI process complete in the
direction of the 1d FFTs it has to perform. LAMMPS uses the
pencil\sphinxhyphen{}decomposition algorithm as shown in the {\hyperref[\detokenize{Developer_par_long:fft-parallel}]{\sphinxcrossref{\DUrole{std,std-ref}{Parallel FFT in PPPM}}}}
figure.

\sphinxAtStartPar
Initially (far left), each processor owns a brick of same\sphinxhyphen{}color grid
cells (actually grid points) contained within in its subdomain.  A
brick\sphinxhyphen{}to\sphinxhyphen{}pencil communication operation converts this layout to 1d
pencils in the \sphinxstyleemphasis{x}\sphinxhyphen{}dimension (center left).  Again, cells of the same
color are owned by the same processor.  Each processor can then compute
a 1d FFT on each pencil of data it wholly owns using a call to the
configured FFT library.  A pencil\sphinxhyphen{}to\sphinxhyphen{}pencil communication then converts
this layout to pencils in the \sphinxstyleemphasis{y} dimension (center right) which
effectively transposes the \sphinxstyleemphasis{x} and \sphinxstyleemphasis{y} dimensions of the grid, followed
by 1d FFTs in \sphinxstyleemphasis{y}.  A final transpose of pencils from \sphinxstyleemphasis{y} to \sphinxstyleemphasis{z} (far
right) followed by 1d FFTs in \sphinxstyleemphasis{z} completes the forward FFT.  The data
is left in a \sphinxstyleemphasis{z}\sphinxhyphen{}pencil layout for the convolution operation.  One or
more inverse FFTs then perform the sequence of 1d FFTs and communication
steps in reverse order; the final layout of resulting grid values is the
same as the initial brick layout.

\sphinxAtStartPar
Each communication operation within the FFT (brick\sphinxhyphen{}to\sphinxhyphen{}pencil or
pencil\sphinxhyphen{}to\sphinxhyphen{}pencil or pencil\sphinxhyphen{}to\sphinxhyphen{}brick) converts one tiling of the 3d grid
to another, where a tiling in this context means an assignment of a
small brick\sphinxhyphen{}shaped subset of grid points to each processor, the union of
which comprise the entire grid.  The parallel \sphinxhref{https://lammps.github.io/fftmpi/}{fftMPI library} written for LAMMPS allows arbitrary
definitions of the tiling so that an irregular partitioning of the
simulation domain can use it directly.  Transforming data from one
tiling to another is implemented in \sphinxtitleref{fftMPI} using point\sphinxhyphen{}to\sphinxhyphen{}point
communication, where each processor sends data to a few other
processors, since each tile in the initial tiling overlaps with a
handful of tiles in the final tiling.

\sphinxAtStartPar
The transformations could also be done using collective communication
across all \(P\) processors with a single call to \sphinxcode{\sphinxupquote{MPI\_Alltoall()}}, but
this is typically much slower.  However, for the specialized brick and
pencil tiling illustrated in {\hyperref[\detokenize{Developer_par_long:fft-parallel}]{\sphinxcrossref{\DUrole{std,std-ref}{Parallel FFT in PPPM}}}} figure, collective
communication across the entire MPI communicator is not required.  In
the example, an \(8^3\) grid with 512 grid cells is partitioned
across 64 processors; each processor owns a 2x2x2 3d brick of grid
cells.  The initial brick\sphinxhyphen{}to\sphinxhyphen{}pencil communication (upper left to upper
right) only requires collective communication within subgroups of 4
processors, as illustrated by the 4 colors.  More generally, a
brick\sphinxhyphen{}to\sphinxhyphen{}pencil communication can be performed by partitioning \sphinxstyleemphasis{P}
processors into \(P^{\frac{2}{3}}\) subgroups of
\(P^{\frac{1}{3}}\) processors each.  Each subgroup performs
collective communication only within its subgroup.  Similarly,
pencil\sphinxhyphen{}to\sphinxhyphen{}pencil communication can be performed by partitioning \sphinxstyleemphasis{P}
processors into \(P^{\frac{1}{2}}\) subgroups of
\(P^{\frac{1}{2}}\) processors each.  This is illustrated in the
figure for the \(y \Rightarrow z\) communication (center).  An
eight\sphinxhyphen{}processor subgroup owns the front \sphinxstyleemphasis{yz} plane of data and performs
collective communication within the subgroup to transpose from a
\sphinxstyleemphasis{y}\sphinxhyphen{}pencil to \sphinxstyleemphasis{z}\sphinxhyphen{}pencil layout.

\sphinxAtStartPar
LAMMPS invokes point\sphinxhyphen{}to\sphinxhyphen{}point communication by default, but also
provides the option of partitioned collective communication when using a
{\hyperref[\detokenize{kspace_modify::doc}]{\sphinxcrossref{\DUrole{doc}{kspace\_modify collective yes}}}} command to switch to
that mode.  In the latter case, the code detects the size of the
disjoint subgroups and partitions the single \sphinxstyleemphasis{P}\sphinxhyphen{}size communicator into
multiple smaller communicators, each of which invokes collective
communication.  Testing on a large IBM Blue Gene/Q machine at Argonne
National Labs showed a significant improvement in FFT performance for
large processor counts; partitioned collective communication was faster
than point\sphinxhyphen{}to\sphinxhyphen{}point communication or global collective communication
involving all \sphinxstyleemphasis{P} processors.

\sphinxAtStartPar
Here are some additional details about FFTs for long\sphinxhyphen{}range and related
grid/particle operations that LAMMPS supports:
\begin{itemize}
\item {} 
\sphinxAtStartPar
The fftMPI library allows each grid dimension to be a multiple of
small prime factors (2,3,5), and allows any number of processors to
perform the FFT.  The resulting brick and pencil decompositions are
thus not always as well\sphinxhyphen{}aligned, but the size of subgroups of
processors for the two modes of communication (brick/pencil and
pencil/pencil) still scale as \(O(P^{\frac{1}{3}})\) and
\(O(P^{\frac{1}{2}})\).

\item {} 
\sphinxAtStartPar
For efficiency in performing 1d FFTs, the grid transpose
operations illustrated in Figure {\hyperref[\detokenize{Developer_par_long:fft-parallel}]{\sphinxcrossref{\DUrole{std,std-ref}{Parallel FFT in PPPM}}}} also involve
reordering the 3d data so that a different dimension is contiguous
in memory.  This reordering can be done during the packing or
unpacking of buffers for MPI communication.

\item {} 
\sphinxAtStartPar
For large systems and particularly many MPI processes, the dominant
cost for parallel FFTs is often the communication, not the computation
of 1d FFTs, even though the latter scales as \(N \log(N)\) in the
number of grid points \sphinxstyleemphasis{N} per grid direction.  This is due to the fact
that only a 2d decomposition into pencils is possible while atom data
(and their corresponding short\sphinxhyphen{}range force and energy computations)
can be decomposed efficiently in 3d.

\sphinxAtStartPar
Reducing the number of MPI processes involved in the MPI communication
will reduce this kind of overhead.  By using a {\hyperref[\detokenize{Speed_omp::doc}]{\sphinxcrossref{\DUrole{doc}{hybrid MPI +
OpenMP parallelization}}}} it is still possible to use all
processes for parallel computation.  It will use OpenMP
parallelization inside the MPI domains. While that may have a lower
parallel efficiency for some part of the computation, that can be less
than the communication overhead in the 3d FFTs.

\sphinxAtStartPar
As an alternative, it is also possible to start a {\hyperref[\detokenize{Run_options:partition}]{\sphinxcrossref{\DUrole{std,std-ref}{multi\sphinxhyphen{}partition}}}} calculation and then use the {\hyperref[\detokenize{run_style::doc}]{\sphinxcrossref{\DUrole{doc}{verlet/split
integrator}}}} to perform the PPPM computation on a
dedicated, separate partition of MPI processes.  This uses an integer
“1:\sphinxstyleemphasis{p}” mapping of \sphinxstyleemphasis{p} subdomains of the atom decomposition to one
subdomain of the FFT grid decomposition and where pairwise non\sphinxhyphen{}bonded
and bonded forces and energies are computed on the larger partition
and the PPPM kspace computation concurrently on the smaller partition.

\item {} 
\sphinxAtStartPar
LAMMPS also implements PPPM\sphinxhyphen{}based solvers for other long\sphinxhyphen{}range
interactions, dipole and dispersion (Lennard\sphinxhyphen{}Jones), which can be used
in conjunction with long\sphinxhyphen{}range  Coulombics for point charges.

\item {} 
\sphinxAtStartPar
LAMMPS implements a \sphinxcode{\sphinxupquote{GridComm}} class which overlays the simulation
domain with a regular grid, partitions it across processors in a
manner consistent with processor subdomains, and provides methods for
forward and reverse communication of owned and ghost grid point
values.  It is used for PPPM as an FFT grid (as outlined above) and
also for the MSM algorithm, which uses a cascade of grid sizes from
fine to coarse to compute long\sphinxhyphen{}range Coulombic forces.  The GridComm
class is also useful for models where continuum fields interact with
particles.  For example, the two\sphinxhyphen{}temperature model (TTM) defines heat
transfer between atoms (particles) and electrons (continuum gas) where
spatial variations in the electron temperature are computed by finite
differences of a discretized heat equation on a regular grid.  The
{\hyperref[\detokenize{fix_ttm::doc}]{\sphinxcrossref{\DUrole{doc}{fix ttm/grid}}}} command uses the \sphinxcode{\sphinxupquote{GridComm}} class
internally to perform its grid operations on a distributed grid
instead of the original {\hyperref[\detokenize{fix_ttm::doc}]{\sphinxcrossref{\DUrole{doc}{fix ttm}}}} which uses a
replicated grid.

\end{itemize}

\sphinxstepscope


\subsection{OpenMP Parallelism}
\label{\detokenize{Developer_par_openmp:openmp-parallelism}}\label{\detokenize{Developer_par_openmp::doc}}
\sphinxAtStartPar
The styles in the INTEL, KOKKOS, and OPENMP packages offer to use OpenMP
thread parallelism to predominantly distribute loops over local data
and thus follow an orthogonal parallelization strategy to the
decomposition into spatial domains used by the {\hyperref[\detokenize{Developer_par_part::doc}]{\sphinxcrossref{\DUrole{doc}{MPI partitioning}}}}.  For clarity, this section discusses only the
implementation in the OPENMP package, as it is the simplest. The INTEL
and KOKKOS packages offer additional options and are more complex since
they support more features and different hardware like co\sphinxhyphen{}processors
or GPUs.

\sphinxAtStartPar
One of the key decisions when implementing the OPENMP package was to
keep the changes to the source code small, so that it would be easier to
maintain the code and keep it in sync with the non\sphinxhyphen{}threaded standard
implementation.  This is achieved by a) making the OPENMP version a
derived class from the regular version (e.g. \sphinxcode{\sphinxupquote{PairLJCutOMP}} from
\sphinxcode{\sphinxupquote{PairLJCut}}) and only overriding methods that are multi\sphinxhyphen{}threaded or
need to be modified to support multi\sphinxhyphen{}threading (similar to what was done
in the OPT package), b) keeping the structure in the modified code very
similar so that side\sphinxhyphen{}by\sphinxhyphen{}side comparisons are still useful, and c)
offloading additional functionality and multi\sphinxhyphen{}thread support functions
into three separate classes \sphinxcode{\sphinxupquote{ThrOMP}}, \sphinxcode{\sphinxupquote{ThrData}}, and \sphinxcode{\sphinxupquote{FixOMP}}.
\sphinxcode{\sphinxupquote{ThrOMP}} provides additional, multi\sphinxhyphen{}thread aware functionality not
available in the corresponding base class (e.g. \sphinxcode{\sphinxupquote{Pair}} for
\sphinxcode{\sphinxupquote{PairLJCutOMP}}) like multi\sphinxhyphen{}thread aware variants of the “tally”
functions. Those functions are made available through multiple
inheritance, so those new functions have to have unique names to avoid
ambiguities; typically \sphinxcode{\sphinxupquote{\_thr}} is appended to the name of the function.
\sphinxcode{\sphinxupquote{ThrData}} is a class that manages per\sphinxhyphen{}thread data structures.  It is
used instead of extending the corresponding storage to per\sphinxhyphen{}thread arrays
to avoid slowdowns due to “false sharing” when multiple threads update
adjacent elements in an array and thus force the CPU cache lines to be
reset and re\sphinxhyphen{}fetched.  \sphinxcode{\sphinxupquote{FixOMP}} finally manages the “multi\sphinxhyphen{}thread
state” like settings and access to per\sphinxhyphen{}thread storage, it is activated
by the {\hyperref[\detokenize{package::doc}]{\sphinxcrossref{\DUrole{doc}{package omp}}}} command.


\subsubsection{Avoiding data races}
\label{\detokenize{Developer_par_openmp:avoiding-data-races}}
\sphinxAtStartPar
A key problem when implementing thread parallelism in an MD code is
to avoid data races when updating accumulated properties like forces,
energies, and stresses.  When interactions are computed, they always
involve multiple atoms and thus there are race conditions when multiple
threads want to update per\sphinxhyphen{}atom data of the same atoms.  Five possible
strategies have been considered to avoid this:
\begin{enumerate}
\sphinxsetlistlabels{\arabic}{enumi}{enumii}{}{.}%
\item {} 
\sphinxAtStartPar
Restructure the code so that there is no overlapping access possible
when computing in parallel, e.g. by breaking lists into multiple
parts and synchronizing threads in between.

\item {} 
\sphinxAtStartPar
Have each thread be “responsible” for a specific group of atoms and
compute these interactions multiple times, once on each thread that
is responsible for a given atom, and then have each thread only update
the properties of this atom.

\item {} 
\sphinxAtStartPar
Use mutexes around functions and regions of code where the data race
could happen.

\item {} 
\sphinxAtStartPar
Use atomic operations when updating per\sphinxhyphen{}atom properties.

\item {} 
\sphinxAtStartPar
Use replicated per\sphinxhyphen{}thread data structures to accumulate data without
conflicts and then use a reduction to combine those results into the
data structures used by the regular style.

\end{enumerate}

\sphinxAtStartPar
Option 5 was chosen for the OPENMP package because it would retain the
performance for the case of a single thread and the code would be more
maintainable.  Option 1 would require extensive code changes,
particularly to the neighbor list code; option 2 would have incurred a
2x or more performance penalty for the serial case; option 3 causes
significant overhead and would enforce serialization of operations in
inner loops and thus defeat the purpose of multi\sphinxhyphen{}threading; option 4
slows down the serial case although not quite as bad as option 2.  The
downside of option 5 is that the overhead of the reduction operations
grows with the number of threads used, so there would be a crossover
point where options 2 or 4 would result in faster executing.  That is
why option 2 for example is used in the GPU package because a GPU is a
processor with a massive number of threads.  However, since the MPI
parallelization is generally more effective for typical MD systems, the
expectation is that thread parallelism is only used for a smaller number
of threads (2\sphinxhyphen{}8).  At the time of its implementation, that number was
equivalent to the number of CPU cores per CPU socket on high\sphinxhyphen{}end
supercomputers.

\sphinxAtStartPar
Thus arrays like the force array are dimensioned to the number of atoms
times the number of threads when enabling OpenMP support, and inside the
compute functions a pointer to a different chunk is obtained by each thread.
Similarly, accumulators like potential energy or virial are kept in
per\sphinxhyphen{}thread instances of the \sphinxcode{\sphinxupquote{ThrData}} class and then only reduced and
stored in their global counterparts at the end of the force computation.


\subsubsection{Loop scheduling}
\label{\detokenize{Developer_par_openmp:loop-scheduling}}
\sphinxAtStartPar
Multi\sphinxhyphen{}thread parallelization is applied by distributing (outer) loops
statically across threads.  Typically, this would be the loop over local
atoms \sphinxstyleemphasis{i} when processing \sphinxstyleemphasis{i,j} pairs of atoms from a neighbor list.
The design of the neighbor list code results in atoms having a similar
number of neighbors for homogeneous systems and thus load imbalances
across threads are not common and typically happen for systems where
also the MPI parallelization would be unbalanced, which would typically
have a more pronounced impact on the performance.  This same loop
scheduling scheme can also be applied to the reduction operations on
per\sphinxhyphen{}atom data to try and reduce the overhead of the reduction operation.


\subsubsection{Neighbor list parallelization}
\label{\detokenize{Developer_par_openmp:neighbor-list-parallelization}}
\sphinxAtStartPar
In addition to the parallelization of force computations, also the
generation of the neighbor lists is parallelized.  As explained
previously, neighbor lists are built by looping over “owned” atoms and
storing the neighbors in “pages”.  In the OPENMP variants of the
neighbor list code, each thread operates on a different chunk of “owned”
atoms and allocates and fills its own set of pages with neighbor list
data.  This is achieved by each thread keeping its own instance of the
{\hyperref[\detokenize{Developer_utils:_CPPv4I0EN9LAMMPS_NS6MyPageE}]{\sphinxcrossref{\sphinxcode{\sphinxupquote{MyPage}}}}} page allocator class.

\sphinxstepscope


\section{Accessing per\sphinxhyphen{}atom data}
\label{\detokenize{Developer_atom:accessing-per-atom-data}}\label{\detokenize{Developer_atom::doc}}
\sphinxAtStartPar
This page discusses how per\sphinxhyphen{}atom data is managed in LAMMPS, how it can
be accessed, what communication patterns apply, and some of the utility
functions that exist for a variety of purposes.


\subsection{Owned and ghost atoms}
\label{\detokenize{Developer_atom:owned-and-ghost-atoms}}
\sphinxAtStartPar
As described on the {\hyperref[\detokenize{Developer_par_part::doc}]{\sphinxcrossref{\DUrole{doc}{parallel partitioning algorithms}}}} page, LAMMPS uses a domain decomposition of the
simulation domain, either in a \sphinxstyleemphasis{brick} or \sphinxstyleemphasis{tiled} manner.  Each MPI
process \sphinxstyleemphasis{owns} exactly one subdomain and the atoms within it. To compute
forces for tuples of atoms that are spread across sub\sphinxhyphen{}domain boundaries,
also a “halo” of \sphinxstyleemphasis{ghost} atoms are maintained within the communication
cutoff distance of its subdomain.

\sphinxAtStartPar
The total number of atoms is stored in \sphinxtitleref{Atom::natoms} (within any
typical class this can be referred to at \sphinxtitleref{atom\sphinxhyphen{}\textgreater{}natoms}). The number of
\sphinxstyleemphasis{owned} (or “local” atoms) are stored in \sphinxtitleref{Atom::nlocal}; the number of
\sphinxstyleemphasis{ghost} atoms is stored in \sphinxtitleref{Atom::nghost}.  The sum of \sphinxtitleref{Atom::nlocal}
over all MPI processes should be \sphinxtitleref{Atom::natoms}. This is by default
regularly checked by the Thermo class, and if the sum does not match,
LAMMPS stops with a “lost atoms” error.  For convenience also the
property \sphinxtitleref{Atom::nmax} is available, this is the maximum of
\sphinxtitleref{Atom::nlocal + Atom::nghost} across all MPI processes.

\sphinxAtStartPar
Per\sphinxhyphen{}atom properties are either managed by the atom style, individual
classes,  or as custom arrays by the individual classes. If only access
to \sphinxstyleemphasis{owned} atoms is needed, they are usually allocated to be of size
\sphinxtitleref{Atom::nlocal}, otherwise of size \sphinxtitleref{Atom::nmax}. Please note that not all
per\sphinxhyphen{}atom properties are available or updated on \sphinxstyleemphasis{ghost} atoms. For
example, per\sphinxhyphen{}atom velocities are only updated with {\hyperref[\detokenize{comm_modify::doc}]{\sphinxcrossref{\DUrole{doc}{comm\_modify vel
yes}}}}.


\subsection{Atom indexing}
\label{\detokenize{Developer_atom:atom-indexing}}
\sphinxAtStartPar
When referring to individual atoms, they may be indexed by their local
\sphinxstyleemphasis{index}, their index in their \sphinxtitleref{Atom::x} array. This is densely populated
containing first all \sphinxstyleemphasis{owned} atoms (index \textless{} \sphinxtitleref{Atom::nlocal}) and then all
\sphinxstyleemphasis{ghost} atoms.  The order of atoms in these arrays can change due to
atoms migrating between between subdomains, atoms being added or
deleted, or atoms being sorted for better cache efficiency.  Atoms are
globally uniquely identified by their \sphinxstyleemphasis{atom ID}. There may be multiple
atoms with the same atom ID present, but only one of them may be an
\sphinxstyleemphasis{owned} atom.

\sphinxAtStartPar
To find the local \sphinxstyleemphasis{index} of an atom, when the \sphinxstyleemphasis{atom ID} is known, the
\sphinxtitleref{Atom::map()} function may be used. It will return the local atom index
or \sphinxhyphen{}1. If the returned value is between 0 (inclusive) and \sphinxtitleref{Atom::nlocal}
(exclusive) it is an \sphinxstyleemphasis{owned} or “local” atom; for larger values the atom
is present as a ghost atom; for a value of \sphinxhyphen{}1, the atom is not present
on the current subdomain at all.

\sphinxAtStartPar
If multiple atoms with the same tag exist in the same subdomain, they
can be found via the \sphinxtitleref{Atom::sametag} array. It points to the next atom
index with the same tag or \sphinxhyphen{}1 if there are no more atoms with the same
tag.  The list will be exhaustive when starting with an index of an
\sphinxstyleemphasis{owned} atom, since the atom IDs are unique, so there can only be one
such atom.  Example code to count atoms with same atom ID in a subdomain:

\begin{sphinxVerbatim}[commandchars=\\\{\}]
\PYG{k}{for}\PYG{+w}{ }\PYG{p}{(}\PYG{k+kt}{int}\PYG{+w}{ }\PYG{n}{i}\PYG{+w}{ }\PYG{o}{=}\PYG{+w}{ }\PYG{l+m+mi}{0}\PYG{p}{;}\PYG{+w}{ }\PYG{n}{i}\PYG{+w}{ }\PYG{o}{\PYGZlt{}}\PYG{+w}{ }\PYG{n}{atom}\PYG{o}{\PYGZhy{}}\PYG{o}{\PYGZgt{}}\PYG{n}{nlocal}\PYG{p}{;}\PYG{+w}{ }\PYG{o}{+}\PYG{o}{+}\PYG{n}{i}\PYG{p}{)}\PYG{+w}{ }\PYG{p}{\PYGZob{}}
\PYG{+w}{  }\PYG{k+kt}{int}\PYG{+w}{ }\PYG{n}{count}\PYG{+w}{ }\PYG{o}{=}\PYG{+w}{ }\PYG{l+m+mi}{0}\PYG{p}{;}
\PYG{+w}{  }\PYG{k}{while}\PYG{+w}{ }\PYG{p}{(}\PYG{n}{sametag}\PYG{p}{[}\PYG{n}{i}\PYG{p}{]}\PYG{+w}{ }\PYG{o}{\PYGZgt{}}\PYG{o}{=}\PYG{+w}{ }\PYG{l+m+mi}{0}\PYG{p}{)}\PYG{+w}{ }\PYG{p}{\PYGZob{}}
\PYG{+w}{    }\PYG{n}{i}\PYG{+w}{ }\PYG{o}{=}\PYG{+w}{ }\PYG{n}{sametag}\PYG{p}{[}\PYG{n}{i}\PYG{p}{]}\PYG{p}{;}
\PYG{+w}{    }\PYG{o}{+}\PYG{o}{+}\PYG{n}{count}\PYG{p}{;}
\PYG{+w}{  }\PYG{p}{\PYGZcb{}}
\PYG{+w}{  }\PYG{n}{printf}\PYG{p}{(}\PYG{l+s}{\PYGZdq{}}\PYG{l+s}{Atom ID: \PYGZpc{}ld is present \PYGZpc{}d times}\PYG{l+s+se}{\PYGZbs{}n}\PYG{l+s}{\PYGZdq{}}\PYG{p}{,}\PYG{+w}{ }\PYG{n}{atom}\PYG{o}{\PYGZhy{}}\PYG{o}{\PYGZgt{}}\PYG{n}{tag}\PYG{p}{[}\PYG{n}{i}\PYG{p}{]}\PYG{p}{,}\PYG{+w}{ }\PYG{n}{count}\PYG{p}{)}\PYG{p}{;}
\PYG{p}{\PYGZcb{}}
\end{sphinxVerbatim}


\subsection{Atom class versus AtomVec classes}
\label{\detokenize{Developer_atom:atom-class-versus-atomvec-classes}}
\sphinxAtStartPar
The \sphinxtitleref{Atom} class contains all kinds of flags and counters about atoms in
the system and that includes pointers to \sphinxstylestrong{all} per\sphinxhyphen{}atom properties
available for atoms.  However, only a subset of these pointers are
non\sphinxhyphen{}NULL and which those are depends on the atom style.  For each atom
style there is a corresponding \sphinxtitleref{AtomVecXXX} class derived from the
\sphinxtitleref{AtomVec} base class, where the XXX indicates the atom style.  This
\sphinxtitleref{AtomVecXXX} class will update the counters and per\sphinxhyphen{}atom pointers if
atoms are added or removed to the system or migrate between subdomains.

\sphinxstepscope


\section{Communication patterns}
\label{\detokenize{Developer_comm_ops:communication-patterns}}\label{\detokenize{Developer_comm_ops::doc}}
\sphinxAtStartPar
This page describes various inter\sphinxhyphen{}processor communication operations
provided by LAMMPS, mostly in the core \sphinxstyleemphasis{Comm} class.  These are operations
for common tasks implemented using MPI library calls.  They are used by
other classes to perform communication of different kinds.  These
operations are useful to know about when writing new code for LAMMPS
that needs to communicate data between processors.


\subsection{Owned and ghost atoms}
\label{\detokenize{Developer_comm_ops:owned-and-ghost-atoms}}
\sphinxAtStartPar
As described on the {\hyperref[\detokenize{Developer_par_part::doc}]{\sphinxcrossref{\DUrole{doc}{parallel partitioning algorithms}}}} page, LAMMPS spatially decomposes the simulation
domain, either in a \sphinxstyleemphasis{brick} or \sphinxstyleemphasis{tiled} manner.  Each processor (MPI
task) owns atoms within its subdomain and additionally stores ghost
atoms within a cutoff distance of its subdomain.


\subsubsection{Forward and reverse communication}
\label{\detokenize{Developer_comm_ops:forward-and-reverse-communication}}
\sphinxAtStartPar
As described on the {\hyperref[\detokenize{Developer_par_comm::doc}]{\sphinxcrossref{\DUrole{doc}{parallel communication algorithms}}}} page, the most common communication operations are
first, \sphinxstyleemphasis{forward communication} which sends owned atom information from
each processor to nearby processors to store with their ghost atoms.
The need to do this communication arises when data from the owned atoms
is updated (e.g. their positions) and this updated information needs to
be \sphinxstylestrong{copied} to the corresponding ghost atoms.

\sphinxAtStartPar
And second, \sphinxstyleemphasis{reverse communication}, which sends ghost atom information
from each processor to the owning processor to \sphinxstylestrong{accumulate} (sum)
the values with the corresponding owned atoms.  The need for this
arises when data is computed and also stored with ghost atoms
(e.g. forces when using a “half” neighbor list) and thus those terms
need to be added to their corresponding atoms on the process where
they are “owned” atoms.  Please note, that with the {\hyperref[\detokenize{newton::doc}]{\sphinxcrossref{\DUrole{doc}{newton off}}}} setting this does not happen and the neighbor lists are
constructed so that these interactions are computed on both MPI
processes containing one of the atoms and only the data pertaining to
the local atom is stored.

\sphinxAtStartPar
The time\sphinxhyphen{}integration classes in LAMMPS invoke these operations each
timestep via the \sphinxstyleemphasis{forward\_comm()} and \sphinxstyleemphasis{reverse\_comm()} methods in the
\sphinxstyleemphasis{Comm} class.  Which per\sphinxhyphen{}atom data is communicated depends on the
currently used {\hyperref[\detokenize{atom_style::doc}]{\sphinxcrossref{\DUrole{doc}{atom style}}}} and whether
{\hyperref[\detokenize{comm_modify::doc}]{\sphinxcrossref{\DUrole{doc}{comm\_modify vel}}}} setting is “no” (default) or
“yes”.

\sphinxAtStartPar
Similarly, \sphinxstyleemphasis{Pair} style classes can invoke the \sphinxstyleemphasis{forward\_comm(this)}
and \sphinxstyleemphasis{reverse\_comm(this)} methods in the \sphinxstyleemphasis{Comm} class to perform the
same operations on per\sphinxhyphen{}atom data that is generated and stored within
the pair style class. Note that this function requires passing the
\sphinxcode{\sphinxupquote{this}} pointer as the first argument to enable the \sphinxstyleemphasis{Comm} class to
call the “pack” and “unpack” functions discussed below.  An example of
the use of these functions are many\sphinxhyphen{}body pair styles like the
embedded\sphinxhyphen{}atom method (EAM) which compute intermediate values in the
first part of the compute() function that need to be stored by both
owned and ghost atoms for the second part of the force computation.
The \sphinxstyleemphasis{Comm} class methods perform the MPI communication for buffers of
per\sphinxhyphen{}atom data.  They “call back” to the \sphinxstyleemphasis{Pair} class, so it can \sphinxstyleemphasis{pack}
or \sphinxstyleemphasis{unpack} the buffer with data the \sphinxstyleemphasis{Pair} class owns.  There are 4
such methods that the \sphinxstyleemphasis{Pair} class must define, assuming it uses both
forward and reverse communication:
\begin{itemize}
\item {} 
\sphinxAtStartPar
pack\_forward\_comm()

\item {} 
\sphinxAtStartPar
unpack\_forward\_comm()

\item {} 
\sphinxAtStartPar
pack\_reverse\_comm()

\item {} 
\sphinxAtStartPar
unpack\_reverse\_comm()

\end{itemize}

\sphinxAtStartPar
The arguments to these methods include the buffer and a list of atoms
to pack or unpack.  The \sphinxstyleemphasis{Pair} class also must set the \sphinxstyleemphasis{comm\_forward}
and \sphinxstyleemphasis{comm\_reverse} variables, which store the number of values stored
in the communication buffers for each operation.  This means, if
desired, it can choose to store multiple per\sphinxhyphen{}atom values in the
buffer, and they will be communicated together to minimize
communication overhead.  The communication buffers are defined vectors
containing \sphinxcode{\sphinxupquote{double}} values.  To correctly store integers that may be
64\sphinxhyphen{}bit (bigint, tagint, imageint) in the buffer, you need to use the
{\hyperref[\detokenize{Developer_utils:communication-buffer-coding-with-ubuf}]{\sphinxcrossref{\DUrole{std,std-ref}{ubuf union}}}} construct.

\sphinxAtStartPar
The \sphinxstyleemphasis{Fix}, \sphinxstyleemphasis{Compute}, and \sphinxstyleemphasis{Dump} classes can also invoke the same kind
of forward and reverse communication operations using the same \sphinxstyleemphasis{Comm}
class methods.  Likewise, the same pack/unpack methods and
comm\_forward/comm\_reverse variables must be defined by the calling
\sphinxstyleemphasis{Fix}, \sphinxstyleemphasis{Compute}, or \sphinxstyleemphasis{Dump} class.

\sphinxAtStartPar
For \sphinxstyleemphasis{Fix} classes, there is an optional second argument to the
\sphinxstyleemphasis{forward\_comm()} and \sphinxstyleemphasis{reverse\_comm()} call which can be used when the
fix performs multiple modes of communication, with different numbers
of values per atom.  The fix should set the \sphinxstyleemphasis{comm\_forward} and
\sphinxstyleemphasis{comm\_reverse} variables to the maximum value, but can invoke the
communication for a particular mode with a smaller value.  For this
to work, the \sphinxstyleemphasis{pack\_forward\_comm()}, etc methods typically use a class
member variable to choose which values to pack/unpack into/from the
buffer.

\sphinxAtStartPar
Finally, for reverse communications in \sphinxstyleemphasis{Fix} classes there is also the
\sphinxstyleemphasis{reverse\_comm\_variable()} method that allows the communication to have
a different amount of data per\sphinxhyphen{}atom.  It invokes these corresponding
callback methods:
\begin{itemize}
\item {} 
\sphinxAtStartPar
pack\_reverse\_comm\_size()

\item {} 
\sphinxAtStartPar
unpack\_reverse\_comm\_size()

\end{itemize}

\sphinxAtStartPar
which have extra arguments to specify the amount of data stored
in the buffer for each atom.


\subsection{Higher level communication}
\label{\detokenize{Developer_comm_ops:higher-level-communication}}
\sphinxAtStartPar
There are also several higher\sphinxhyphen{}level communication operations provided
in LAMMPS which work for either \sphinxstyleemphasis{brick} or \sphinxstyleemphasis{tiled} decompositions.
They may be useful for a new class to invoke if it requires more
sophisticated communication than the \sphinxstyleemphasis{forward} and \sphinxstyleemphasis{reverse} methods
provide.  The 3 communication operations described here are
\begin{itemize}
\item {} 
\sphinxAtStartPar
ring

\item {} 
\sphinxAtStartPar
irregular

\item {} 
\sphinxAtStartPar
rendezvous

\end{itemize}

\sphinxAtStartPar
You can invoke these \sphinxstyleemphasis{grep} command in the LAMMPS src directory, to
see a list of classes that invoke the 3 operations.
\begin{itemize}
\item {} 
\sphinxAtStartPar
\sphinxcode{\sphinxupquote{grep "\textbackslash{}\sphinxhyphen{}\textgreater{}ring" *.cpp */*.cpp}}

\item {} 
\sphinxAtStartPar
\sphinxcode{\sphinxupquote{grep "irregular\textbackslash{}\sphinxhyphen{}\textgreater{}" *.cpp}}

\item {} 
\sphinxAtStartPar
\sphinxcode{\sphinxupquote{grep "\textbackslash{}\sphinxhyphen{}\textgreater{}rendezvous" *.cpp */*.cpp}}

\end{itemize}


\subsubsection{Ring operation}
\label{\detokenize{Developer_comm_ops:ring-operation}}
\sphinxAtStartPar
The \sphinxstyleemphasis{ring} operation is invoked via the \sphinxstyleemphasis{ring()} method in the \sphinxstyleemphasis{Comm}
class.

\sphinxAtStartPar
Each processor first creates a buffer with a list of values, typically
associated with a subset of the atoms it owns.  Now think of the \sphinxstyleemphasis{P}
processors as connected to each other in a \sphinxstyleemphasis{ring}.  Each processor \sphinxstyleemphasis{M}
sends data to the next \sphinxstyleemphasis{M+1} processor.  It receives data from the
preceding \sphinxstyleemphasis{M\sphinxhyphen{}1} processor.  The ring is periodic so that the last
processor sends to the first processor, and the first processor
receives from the last processor.

\sphinxAtStartPar
Invoking the \sphinxstyleemphasis{ring()} method passes each processor’s buffer in \sphinxstyleemphasis{P}
steps around the ring.  At each step a \sphinxstyleemphasis{callback} method, provided as
an argument to ring(), in the caller is invoked.  This allows each
processor to examine the data buffer provided by every other
processor.  It may extract values needed by its atoms from the
buffers, or it may alter placeholder values in the buffer.  In the
latter case, when the \sphinxstyleemphasis{ring} operation is complete, each processor can
examine its original buffer to extract modified values.

\sphinxAtStartPar
Note that the \sphinxstyleemphasis{ring} operation is similar to an MPI\_Alltoall()
operation, where every processor effectively sends and receives data to
every other processor.  The difference is that the \sphinxstyleemphasis{ring} operation
does it one step at a time, so the total volume of data does not need
to be stored by every processor.  However, the \sphinxstyleemphasis{ring} operation is
also less efficient than MPI\_Alltoall() because of the \sphinxstyleemphasis{P} stages
required.  So it is typically only suitable for small data buffers and
occasional operations that are not time\sphinxhyphen{}critical.


\subsubsection{Irregular operation}
\label{\detokenize{Developer_comm_ops:irregular-operation}}
\sphinxAtStartPar
The \sphinxstyleemphasis{irregular} operation is provided by the \sphinxstyleemphasis{Irregular} class.  What
LAMMPS terms irregular communication is when each processor knows what
data it needs to send to what processor, but does not know what
processors are sending it data.  An example is when load\sphinxhyphen{}balancing is
performed and each processor needs to send some of its atoms to new
processors.

\sphinxAtStartPar
The \sphinxstyleemphasis{Irregular} class provides 5 high\sphinxhyphen{}level methods useful in this
context:
\begin{itemize}
\item {} 
\sphinxAtStartPar
create\_data()

\item {} 
\sphinxAtStartPar
exchange\_data()

\item {} 
\sphinxAtStartPar
create\_atom()

\item {} 
\sphinxAtStartPar
exchange\_atom()

\item {} 
\sphinxAtStartPar
migrate\_atoms()

\end{itemize}

\sphinxAtStartPar
For the \sphinxstyleemphasis{create\_data()} method, each processor specifies a list of \sphinxstyleemphasis{N}
datums to send, each to a specified processor.  Internally, the method
creates efficient data structures for performing the communication.
The \sphinxstyleemphasis{exchange\_data()} method triggers the communication to be
performed.  Each processor provides the vector of \sphinxstyleemphasis{N} datums to send,
and the size of each datum.  All datums must be the same size.

\sphinxAtStartPar
The \sphinxstyleemphasis{create\_atom()} and \sphinxstyleemphasis{exchange\_atom()} methods are similar, except
that the size of each datum can be different.  Typically, this is used
to communicate atoms, each with a variable amount of per\sphinxhyphen{}atom data, to
other processors.

\sphinxAtStartPar
The \sphinxstyleemphasis{migrate\_atoms()} method is a convenience wrapper on the
\sphinxstyleemphasis{create\_atom()} and \sphinxstyleemphasis{exchange\_atom()} methods to simplify
communication of all the per\sphinxhyphen{}atom data associated with an atom so that
the atom can effectively migrate to a new owning processor.  It is
similar to the \sphinxstyleemphasis{exchange()} method in the \sphinxstyleemphasis{Comm} class invoked when
atoms move to neighboring processors (in the regular or tiled
decomposition) during timestepping, except that it allows atoms to
have moved arbitrarily long distances and still be properly
communicated to a new owning processor.


\subsubsection{Rendezvous operation}
\label{\detokenize{Developer_comm_ops:rendezvous-operation}}
\sphinxAtStartPar
Finally, the \sphinxstyleemphasis{rendezvous} operation is invoked via the \sphinxstyleemphasis{rendezvous()}
method in the \sphinxstyleemphasis{Comm} class.  Depending on how much communication is
needed and how many processors a LAMMPS simulation is running on, it
can be a much more efficient choice than the \sphinxstyleemphasis{ring()} method.  It uses
the \sphinxstyleemphasis{irregular} operation internally once or twice to do its
communication.  The rendezvous algorithm is described in detail in
{\hyperref[\detokenize{Developer_comm_ops:plimpton}]{\sphinxcrossref{\DUrole{std,std-ref}{(Plimpton)}}}}, including some LAMMPS use cases.

\sphinxAtStartPar
For the \sphinxstyleemphasis{rendezvous()} method, each processor specifies a list of \sphinxstyleemphasis{N}
datums to send and which processor to send each of them to.
Internally, this communication is performed as an irregular operation.
The received datums are returned to the caller via invocation of
\sphinxstyleemphasis{callback} function, provided as an argument to \sphinxstyleemphasis{rendezvous()}.  The
caller can then process the received datums and (optionally) assemble
a new list of datums to communicate to a new list of specific
processors.  When the callback function exits, the \sphinxstyleemphasis{rendezvous()}
method performs a second irregular communication on the new list of
datums.

\sphinxAtStartPar
Examples in LAMMPS of use of the \sphinxstyleemphasis{rendezvous} operation are the
{\hyperref[\detokenize{fix_rigid::doc}]{\sphinxcrossref{\DUrole{doc}{fix rigid/small}}}} and {\hyperref[\detokenize{fix_shake::doc}]{\sphinxcrossref{\DUrole{doc}{fix shake}}}} commands (for one\sphinxhyphen{}time identification of the rigid body
atom clusters) and the identification of special\_bond 1\sphinxhyphen{}2, 1\sphinxhyphen{}3 and 1\sphinxhyphen{}4
neighbors within molecules.  See the {\hyperref[\detokenize{special_bonds::doc}]{\sphinxcrossref{\DUrole{doc}{special\_bonds}}}}
command for context.


\bigskip\hrule\bigskip

\phantomsection\label{\detokenize{Developer_comm_ops:plimpton}}
\sphinxAtStartPar
\sphinxstylestrong{(Plimpton)} Plimpton and Knight, JPDC, 147, 184\sphinxhyphen{}195 (2021).

\sphinxstepscope


\section{How a timestep works}
\label{\detokenize{Developer_flow:how-a-timestep-works}}\label{\detokenize{Developer_flow::doc}}
\sphinxAtStartPar
The first and most fundamental operation within LAMMPS to understand is
how a timestep is structured.  Timestepping is performed by calling
methods of the Integrate class instance within the Update class.  Since
Integrate is a base class, it will point to an instance of a derived
class corresponding to what is selected by the {\hyperref[\detokenize{run_style::doc}]{\sphinxcrossref{\DUrole{doc}{run\_style}}}} input script command.

\sphinxAtStartPar
In this section, the timestep implemented by the Verlet class is
described.  A similar timestep protocol is implemented by the Respa
class, for the r\sphinxhyphen{}RESPA hierarchical timestepping method.

\sphinxAtStartPar
The Min base class performs energy minimization, so does not perform a
literal timestep.  But it has logic similar to what is described here,
to compute forces and invoke fixes at each iteration of a minimization.
Differences between time integration and minimization are highlighted at
the end of this section.

\sphinxAtStartPar
The Verlet class is encoded in the \sphinxcode{\sphinxupquote{src/verlet.cpp}} and \sphinxcode{\sphinxupquote{verlet.h}}
files.  It implements the velocity\sphinxhyphen{}Verlet timestepping algorithm.  The
workhorse method is \sphinxcode{\sphinxupquote{Verlet::run()}}, but first we highlight several
other methods in the class.
\begin{itemize}
\item {} 
\sphinxAtStartPar
The \sphinxcode{\sphinxupquote{init()}} method is called at the beginning of each dynamics
run.  It simply sets some internal flags, based on user settings in
other parts of the code.

\item {} 
\sphinxAtStartPar
The \sphinxcode{\sphinxupquote{setup()}} or \sphinxcode{\sphinxupquote{setup\_minimal()}} methods are also called before
each run.  The velocity\sphinxhyphen{}Verlet method requires current forces be
calculated before the first timestep, so these routines compute
forces due to all atomic interactions, using the same logic that
appears in the timestepping described next.  A few fixes are also
invoked, using the mechanism described in the next section.  Various
counters are also initialized before the run begins.  The
\sphinxcode{\sphinxupquote{setup\_minimal()}} method is a variant that has a flag for performing
less setup.  This is used when runs are continued and information
from the previous run is still valid.  For example, if repeated
short LAMMPS runs are being invoked, interleaved by other commands,
via the \sphinxstyleemphasis{pre no} and \sphinxstyleemphasis{every} options of the run command, the
\sphinxcode{\sphinxupquote{setup\_minimal()}} method is used.

\item {} 
\sphinxAtStartPar
The \sphinxcode{\sphinxupquote{force\_clear()}} method initializes force and other arrays to
zero before each timestep, so that forces (torques, etc) can be
accumulated.

\end{itemize}

\sphinxAtStartPar
Now for the \sphinxcode{\sphinxupquote{Verlet::run()}} method.  Its basic structure in hi\sphinxhyphen{}level
pseudocode is shown below.  In the actual code in \sphinxcode{\sphinxupquote{src/verlet.cpp}}
some of these operations are conditionally invoked.

\begin{sphinxVerbatim}[commandchars=\\\{\}]
\PYG{n}{loop} \PYG{n}{over} \PYG{n}{N} \PYG{n}{timesteps}\PYG{p}{:}
  \PYG{k}{if} \PYG{n}{timeout} \PYG{n}{condition}\PYG{p}{:} \PYG{k}{break}
  \PYG{n}{ev\PYGZus{}set}\PYG{p}{(}\PYG{p}{)}

  \PYG{n}{fix}\PYG{o}{\PYGZhy{}}\PYG{o}{\PYGZgt{}}\PYG{n}{initial\PYGZus{}integrate}\PYG{p}{(}\PYG{p}{)}
  \PYG{n}{fix}\PYG{o}{\PYGZhy{}}\PYG{o}{\PYGZgt{}}\PYG{n}{post\PYGZus{}integrate}\PYG{p}{(}\PYG{p}{)}

  \PYG{n}{nflag} \PYG{o}{=} \PYG{n}{neighbor}\PYG{o}{\PYGZhy{}}\PYG{o}{\PYGZgt{}}\PYG{n}{decide}\PYG{p}{(}\PYG{p}{)}
  \PYG{k}{if} \PYG{n}{nflag}\PYG{p}{:}
    \PYG{n}{fix}\PYG{o}{\PYGZhy{}}\PYG{o}{\PYGZgt{}}\PYG{n}{pre\PYGZus{}exchange}\PYG{p}{(}\PYG{p}{)}
    \PYG{n}{domain}\PYG{o}{\PYGZhy{}}\PYG{o}{\PYGZgt{}}\PYG{n}{pbc}\PYG{p}{(}\PYG{p}{)}
    \PYG{n}{domain}\PYG{o}{\PYGZhy{}}\PYG{o}{\PYGZgt{}}\PYG{n}{reset\PYGZus{}box}\PYG{p}{(}\PYG{p}{)}
    \PYG{n}{comm}\PYG{o}{\PYGZhy{}}\PYG{o}{\PYGZgt{}}\PYG{n}{setup}\PYG{p}{(}\PYG{p}{)}
    \PYG{n}{neighbor}\PYG{o}{\PYGZhy{}}\PYG{o}{\PYGZgt{}}\PYG{n}{setup\PYGZus{}bins}\PYG{p}{(}\PYG{p}{)}
    \PYG{n}{comm}\PYG{o}{\PYGZhy{}}\PYG{o}{\PYGZgt{}}\PYG{n}{exchange}\PYG{p}{(}\PYG{p}{)}
    \PYG{n}{comm}\PYG{o}{\PYGZhy{}}\PYG{o}{\PYGZgt{}}\PYG{n}{borders}\PYG{p}{(}\PYG{p}{)}
    \PYG{n}{fix}\PYG{o}{\PYGZhy{}}\PYG{o}{\PYGZgt{}}\PYG{n}{pre\PYGZus{}neighbor}\PYG{p}{(}\PYG{p}{)}
    \PYG{n}{neighbor}\PYG{o}{\PYGZhy{}}\PYG{o}{\PYGZgt{}}\PYG{n}{build}\PYG{p}{(}\PYG{p}{)}
    \PYG{n}{fix}\PYG{o}{\PYGZhy{}}\PYG{o}{\PYGZgt{}}\PYG{n}{post\PYGZus{}neighbor}\PYG{p}{(}\PYG{p}{)}
  \PYG{k}{else}\PYG{p}{:}
    \PYG{n}{comm}\PYG{o}{\PYGZhy{}}\PYG{o}{\PYGZgt{}}\PYG{n}{forward\PYGZus{}comm}\PYG{p}{(}\PYG{p}{)}

  \PYG{n}{force\PYGZus{}clear}\PYG{p}{(}\PYG{p}{)}
  \PYG{n}{fix}\PYG{o}{\PYGZhy{}}\PYG{o}{\PYGZgt{}}\PYG{n}{pre\PYGZus{}force}\PYG{p}{(}\PYG{p}{)}

  \PYG{n}{pair}\PYG{o}{\PYGZhy{}}\PYG{o}{\PYGZgt{}}\PYG{n}{compute}\PYG{p}{(}\PYG{p}{)}
  \PYG{n}{bond}\PYG{o}{\PYGZhy{}}\PYG{o}{\PYGZgt{}}\PYG{n}{compute}\PYG{p}{(}\PYG{p}{)}
  \PYG{n}{angle}\PYG{o}{\PYGZhy{}}\PYG{o}{\PYGZgt{}}\PYG{n}{compute}\PYG{p}{(}\PYG{p}{)}
  \PYG{n}{dihedral}\PYG{o}{\PYGZhy{}}\PYG{o}{\PYGZgt{}}\PYG{n}{compute}\PYG{p}{(}\PYG{p}{)}
  \PYG{n}{improper}\PYG{o}{\PYGZhy{}}\PYG{o}{\PYGZgt{}}\PYG{n}{compute}\PYG{p}{(}\PYG{p}{)}
  \PYG{n}{kspace}\PYG{o}{\PYGZhy{}}\PYG{o}{\PYGZgt{}}\PYG{n}{compute}\PYG{p}{(}\PYG{p}{)}

  \PYG{n}{fix}\PYG{o}{\PYGZhy{}}\PYG{o}{\PYGZgt{}}\PYG{n}{pre\PYGZus{}reverse}\PYG{p}{(}\PYG{p}{)}
  \PYG{n}{comm}\PYG{o}{\PYGZhy{}}\PYG{o}{\PYGZgt{}}\PYG{n}{reverse\PYGZus{}comm}\PYG{p}{(}\PYG{p}{)}

  \PYG{n}{fix}\PYG{o}{\PYGZhy{}}\PYG{o}{\PYGZgt{}}\PYG{n}{post\PYGZus{}force}\PYG{p}{(}\PYG{p}{)}
  \PYG{n}{fix}\PYG{o}{\PYGZhy{}}\PYG{o}{\PYGZgt{}}\PYG{n}{final\PYGZus{}integrate}\PYG{p}{(}\PYG{p}{)}
  \PYG{n}{fix}\PYG{o}{\PYGZhy{}}\PYG{o}{\PYGZgt{}}\PYG{n}{end\PYGZus{}of\PYGZus{}step}\PYG{p}{(}\PYG{p}{)}

  \PYG{k}{if} \PYG{n+nb}{any} \PYG{n}{output} \PYG{n}{on} \PYG{n}{this} \PYG{n}{step}\PYG{p}{:}
    \PYG{n}{output}\PYG{o}{\PYGZhy{}}\PYG{o}{\PYGZgt{}}\PYG{n}{write}\PYG{p}{(}\PYG{p}{)}

\PYG{c+c1}{\PYGZsh{} after loop}
\PYG{n}{fix}\PYG{o}{\PYGZhy{}}\PYG{o}{\PYGZgt{}}\PYG{n}{post\PYGZus{}run}\PYG{p}{(}\PYG{p}{)}
\end{sphinxVerbatim}

\sphinxAtStartPar
The \sphinxcode{\sphinxupquote{ev\_set()}} method (in the parent Integrate class), sets two flags
(\sphinxstyleemphasis{eflag} and \sphinxstyleemphasis{vflag}) for energy and virial computation.  Each flag
encodes whether global and/or per\sphinxhyphen{}atom energy and virial should be
calculated on this timestep, because some fix or variable or output will
need it.  These flags are passed to the various methods that compute
particle interactions, so that they either compute and tally the
corresponding data or can skip the extra calculations if the energy and
virial are not needed.  See the comments for the \sphinxcode{\sphinxupquote{Integrate::ev\_set()}}
method, which document the flag values.

\sphinxAtStartPar
At various points of the timestep, fixes are invoked,
e.g. \sphinxcode{\sphinxupquote{fix\sphinxhyphen{}\textgreater{}initial\_integrate()}}.  In the code, this is actually done
via the Modify class, which stores all the Fix objects and lists of which
should be invoked at what point in the timestep.  Fixes are the LAMMPS
mechanism for tailoring the operations of a timestep for a particular
simulation.  As described elsewhere, each fix has one or more methods,
each of which is invoked at a specific stage of the timestep, as show in
the timestep pseudocode.  All the active fixes defined in an input
script, that are flagged to have an \sphinxcode{\sphinxupquote{initial\_integrate()}} method, are
invoked at the beginning of each timestep.  Examples are {\hyperref[\detokenize{fix_nve::doc}]{\sphinxcrossref{\DUrole{doc}{fix nve}}}} or {\hyperref[\detokenize{fix_nh::doc}]{\sphinxcrossref{\DUrole{doc}{fix nvt or fix npt}}}} which perform the
start\sphinxhyphen{}of\sphinxhyphen{}timestep velocity\sphinxhyphen{}Verlet integration operations to update
velocities by a half\sphinxhyphen{}step, and coordinates by a full step.  The
\sphinxcode{\sphinxupquote{post\_integrate()}} method is next for operations that need to happen
immediately after those updates.  Only a few fixes use this, e.g. to
reflect particles off box boundaries in the {\hyperref[\detokenize{fix_wall_reflect::doc}]{\sphinxcrossref{\DUrole{doc}{FixWallReflect class}}}}.

\sphinxAtStartPar
The \sphinxcode{\sphinxupquote{decide()}} method in the Neighbor class determines whether
neighbor lists need to be rebuilt on the current timestep (conditions
can be changed using the {\hyperref[\detokenize{neigh_modify::doc}]{\sphinxcrossref{\DUrole{doc}{neigh\_modify every/delay/check}}}} command).  If not, coordinates of ghost atoms are
acquired by each processor via the \sphinxcode{\sphinxupquote{forward\_comm()}} method of the Comm
class.  If neighbor lists need to be built, several operations within
the inner if clause of the pseudocode are first invoked.  The
\sphinxcode{\sphinxupquote{pre\_exchange()}} method of any defined fixes is invoked first.
Typically, this inserts or deletes particles from the system.

\sphinxAtStartPar
Periodic boundary conditions are then applied by the Domain class via
its \sphinxcode{\sphinxupquote{pbc()}} method to remap particles that have moved outside the
simulation box back into the box.  Note that this is not done every
timestep, but only when neighbor lists are rebuilt.  This is so that
each processor’s subdomain will have consistent (nearby) atom
coordinates for its owned and ghost atoms.  It is also why dumped atom
coordinates may be slightly outside the simulation box if not dumped
on a step where the neighbor lists are rebuilt.

\sphinxAtStartPar
The box boundaries are then reset (if needed) via the \sphinxcode{\sphinxupquote{reset\_box()}}
method of the Domain class, e.g. if box boundaries are shrink\sphinxhyphen{}wrapped to
current particle coordinates.  A change in the box size or shape
requires internal information for communicating ghost atoms (Comm class)
and neighbor list bins (Neighbor class) to be updated.  The \sphinxcode{\sphinxupquote{setup()}}
method of the Comm class and \sphinxcode{\sphinxupquote{setup\_bins()}} method of the Neighbor
class perform the update.

\sphinxAtStartPar
The code is now ready to migrate atoms that have left a processor’s
geometric subdomain to new processors.  The \sphinxcode{\sphinxupquote{exchange()}} method of
the Comm class performs this operation.  The \sphinxcode{\sphinxupquote{borders()}} method of the
Comm class then identifies ghost atoms surrounding each processor’s
subdomain and communicates ghost atom information to neighboring
processors.  It does this by looping over all the atoms owned by a
processor to make lists of those to send to each neighbor processor.  On
subsequent timesteps, the lists are used by the \sphinxcode{\sphinxupquote{Comm::forward\_comm()}}
method.

\sphinxAtStartPar
Fixes with a \sphinxcode{\sphinxupquote{pre\_neighbor()}} method are then called.  These typically
re\sphinxhyphen{}build some data structure stored by the fix that depends on the
current atoms owned by each processor.

\sphinxAtStartPar
Now that each processor has a current list of its owned and ghost
atoms, LAMMPS is ready to rebuild neighbor lists via the \sphinxcode{\sphinxupquote{build()}}
method of the Neighbor class.  This is typically done by binning all
owned and ghost atoms, and scanning a stencil of bins around each
owned atom’s bin to make a Verlet list of neighboring atoms within the
force cutoff plus neighbor skin distance.

\sphinxAtStartPar
In the next portion of the timestep, all interaction forces between
particles are computed, after zeroing the per\sphinxhyphen{}atom force vector via the
\sphinxcode{\sphinxupquote{force\_clear()}} method.  If the newton flag is set to \sphinxstyleemphasis{on} by the
newton command, forces are added to both owned and ghost atoms, otherwise
only to owned (aka local) atoms.

\sphinxAtStartPar
Pairwise forces are calculated first, which enables the global virial
(if requested) to be calculated cheaply (at O(N) cost instead of O(N**2)
at the end of the \sphinxcode{\sphinxupquote{Pair::compute()}} method), by a dot product of atom
coordinates and forces.  By including owned and ghost atoms in the dot
product, the effect of periodic boundary conditions is correctly
accounted for.  Molecular topology interactions (bonds, angles,
dihedrals, impropers) are calculated next (if supported by the current
atom style).  The final contribution is from long\sphinxhyphen{}range Coulombic
interactions, invoked by the KSpace class.

\sphinxAtStartPar
The \sphinxcode{\sphinxupquote{pre\_reverse()}} method in fixes is used for operations that have to
be done \sphinxstyleemphasis{before} the upcoming reverse communication (e.g. to perform
additional data transfers or reductions for data computed during the
force computation and stored with ghost atoms).

\sphinxAtStartPar
If the newton flag is on, forces on ghost atoms are communicated and
summed back to their corresponding owned atoms.  The \sphinxcode{\sphinxupquote{reverse\_comm()}}
method of the Comm class performs this operation, which is essentially
the inverse operation of sending copies of owned atom coordinates to
other processor’s ghost atoms.

\sphinxAtStartPar
At this point in the timestep, the total force on each (local) atom is
known.  Additional force constraints (external forces, SHAKE, etc) are
applied by Fixes that have a \sphinxcode{\sphinxupquote{post\_force()}} method.  The second half
of the velocity\sphinxhyphen{}Verlet integration, \sphinxcode{\sphinxupquote{final\_integrate()}} is then
performed (another half\sphinxhyphen{}step update of the velocities) via fixes like
nve, nvt, npt.

\sphinxAtStartPar
At the end of the timestep, fixes that contain an \sphinxcode{\sphinxupquote{end\_of\_step()}}
method are invoked.  These typically perform a diagnostic calculation,
e.g. the ave/time and ave/spatial fixes.  The final operation of the
timestep is to perform any requested output, via the \sphinxcode{\sphinxupquote{write()}} method
of the Output class.  There are 3 kinds of LAMMPS output: thermodynamic
output to the screen and log file, snapshots of atom data to a dump
file, and restart files.  See the {\hyperref[\detokenize{thermo_style::doc}]{\sphinxcrossref{\DUrole{doc}{thermo\_style}}}},
{\hyperref[\detokenize{dump::doc}]{\sphinxcrossref{\DUrole{doc}{dump}}}}, and {\hyperref[\detokenize{restart::doc}]{\sphinxcrossref{\DUrole{doc}{restart}}}} commands for more
details.

\sphinxAtStartPar
The flow of control during energy minimization iterations is similar to
that of a molecular dynamics timestep.  Forces are computed, neighbor
lists are built as needed, atoms migrate to new processors, and atom
coordinates and forces are communicated to neighboring processors.  The
only difference is what Fix class operations are invoked when.  Only a
subset of LAMMPS fixes are useful during energy minimization, as
explained in their individual doc pages.  The relevant Fix class methods
are \sphinxcode{\sphinxupquote{min\_pre\_exchange()}}, \sphinxcode{\sphinxupquote{min\_pre\_force()}}, and
\sphinxcode{\sphinxupquote{min\_post\_force()}}.  Each fix is invoked at the appropriate place
within the minimization iteration.  For example, the
\sphinxcode{\sphinxupquote{min\_post\_force()}} method is analogous to the \sphinxcode{\sphinxupquote{post\_force()}} method
for dynamics; it is used to alter or constrain forces on each atom,
which affects the minimization procedure.

\sphinxAtStartPar
After all iterations are completed, there is a \sphinxcode{\sphinxupquote{cleanup}} step which
calls the \sphinxcode{\sphinxupquote{post\_run()}} method of fixes to perform operations only required
at the end of a calculation (like freeing temporary storage or creating
final outputs).

\sphinxstepscope


\section{Writing new styles}
\label{\detokenize{Developer_write:writing-new-styles}}\label{\detokenize{Developer_write::doc}}
\sphinxAtStartPar
The {\hyperref[\detokenize{Modify::doc}]{\sphinxcrossref{\DUrole{doc}{Modifying \& extending LAMMPS}}}} section of the manual gives an overview of how LAMMPS can
be extended by writing new classes that derive from existing
parent classes in LAMMPS.  Here, some specific coding
details are provided for writing code for LAMMPS.

\sphinxstepscope


\subsection{Writing new pair styles}
\label{\detokenize{Developer_write_pair:writing-new-pair-styles}}\label{\detokenize{Developer_write_pair::doc}}
\sphinxAtStartPar
Pair styles are at the core of most simulations with LAMMPS, since they
are used to compute the forces (plus energy and virial contributions, if
needed) on atoms for pairs or small clusters of atoms within a given
cutoff.  This is often the dominant computation in LAMMPS, and sometimes
even the only one.  Pair styles can be grouped into multiple categories:
\begin{enumerate}
\sphinxsetlistlabels{\arabic}{enumi}{enumii}{}{.}%
\item {} 
\sphinxAtStartPar
simple pairwise additive interactions of point particles
(e.g. {\hyperref[\detokenize{pair_lj::doc}]{\sphinxcrossref{\DUrole{doc}{Lennard\sphinxhyphen{}Jones}}}}, {\hyperref[\detokenize{pair_morse::doc}]{\sphinxcrossref{\DUrole{doc}{Morse}}}},
{\hyperref[\detokenize{pair_buck::doc}]{\sphinxcrossref{\DUrole{doc}{Buckingham}}}})

\item {} 
\sphinxAtStartPar
pairwise additive interactions of point particles with added
{\hyperref[\detokenize{pair_coul::doc}]{\sphinxcrossref{\DUrole{doc}{Coulomb}}}} interactions or \sphinxstyleemphasis{only} the Coulomb interactions

\item {} 
\sphinxAtStartPar
manybody interactions of point particles (e.g. {\hyperref[\detokenize{pair_eam::doc}]{\sphinxcrossref{\DUrole{doc}{EAM}}}},
{\hyperref[\detokenize{pair_tersoff::doc}]{\sphinxcrossref{\DUrole{doc}{Tersoff}}}})

\item {} 
\sphinxAtStartPar
complex interactions that include additional per\sphinxhyphen{}atom properties
(e.g. Discrete Element Models (DEM), Peridynamics, Ellipsoids)

\item {} 
\sphinxAtStartPar
special purpose pair styles that may not even compute forces like
{\hyperref[\detokenize{pair_zero::doc}]{\sphinxcrossref{\DUrole{doc}{pair\_style zero}}}} and {\hyperref[\detokenize{pair_tracker::doc}]{\sphinxcrossref{\DUrole{doc}{pair\_style tracker}}}}, or are a wrapper for multiple kinds of interactions
like {\hyperref[\detokenize{pair_hybrid::doc}]{\sphinxcrossref{\DUrole{doc}{pair\_style hybrid}}}}, {\hyperref[\detokenize{pair_list::doc}]{\sphinxcrossref{\DUrole{doc}{pair\_style list}}}},
and {\hyperref[\detokenize{pair_kim::doc}]{\sphinxcrossref{\DUrole{doc}{pair\_style kim}}}}

\end{enumerate}

\sphinxAtStartPar
In the text below, we will discuss aspects of implementing pair styles
in LAMMPS by looking at representative case studies.  The design of
LAMMPS allows developers to focus on the essentials, which is to compute
the forces (and energies or virial contributions), enter and manage the
global settings as well as the potential parameters, and the pair style
specific parts of reading and writing restart and data files.  Most of
the complex tasks like management of the atom positions, domain
decomposition and boundaries, or neighbor list creation are handled
transparently by other parts of the LAMMPS code.

\sphinxAtStartPar
As shown on the page for {\hyperref[\detokenize{Modify_pair::doc}]{\sphinxcrossref{\DUrole{doc}{writing or extending pair styles}}}}, in order to implement a new pair style, a new class must
be written that is either directly or indirectly derived from the
\sphinxcode{\sphinxupquote{Pair}} class.  If that class is directly derived from \sphinxcode{\sphinxupquote{Pair}}, there
are three methods that \sphinxstyleemphasis{must} be re\sphinxhyphen{}implemented, since they are “pure”
in the base class: \sphinxcode{\sphinxupquote{Pair::compute()}}, \sphinxcode{\sphinxupquote{Pair::settings()}},
\sphinxcode{\sphinxupquote{Pair::coeff()}}.  In addition, a custom constructor is needed.  All
other methods are optional and have default implementations in the base
class (most of which do nothing), but they may need to be overridden
depending on the requirements of the model.

\sphinxAtStartPar
We are looking at the following cases:
\begin{itemize}
\item {} 
\sphinxAtStartPar
{\hyperref[\detokenize{Developer_write_pair:case-1-a-pairwise-additive-model}]{\sphinxcrossref{Case 1: a pairwise additive model}}}

\item {} 
\sphinxAtStartPar
{\hyperref[\detokenize{Developer_write_pair:case-2-a-many-body-potential}]{\sphinxcrossref{Case 2: a many\sphinxhyphen{}body potential}}}

\item {} 
\sphinxAtStartPar
{\hyperref[\detokenize{Developer_write_pair:case-3-a-potential-requiring-communication}]{\sphinxcrossref{Case 3: a potential requiring communication}}}

\item {} 
\sphinxAtStartPar
{\hyperref[\detokenize{Developer_write_pair:case-4-potentials-without-a-compute-function}]{\sphinxcrossref{Case 4: potentials without a compute() function}}}

\end{itemize}


\subsection{Package and build system considerations}
\label{\detokenize{Developer_write_pair:package-and-build-system-considerations}}
\sphinxAtStartPar
In general, new pair styles should be added to the {\hyperref[\detokenize{Packages_details:pkg-extra-pair}]{\sphinxcrossref{\DUrole{std,std-ref}{EXTRA\sphinxhyphen{}PAIR
package}}}} unless they are an accelerated pair style and
then they should be added to the corresponding accelerator package
({\hyperref[\detokenize{Packages_details:pkg-gpu}]{\sphinxcrossref{\DUrole{std,std-ref}{GPU}}}}, {\hyperref[\detokenize{Packages_details:pkg-intel}]{\sphinxcrossref{\DUrole{std,std-ref}{INTEL}}}}, {\hyperref[\detokenize{Packages_details:pkg-kokkos}]{\sphinxcrossref{\DUrole{std,std-ref}{KOKKOS}}}}, {\hyperref[\detokenize{Packages_details:pkg-openmp}]{\sphinxcrossref{\DUrole{std,std-ref}{OPENMP}}}}, {\hyperref[\detokenize{Packages_details:pkg-opt}]{\sphinxcrossref{\DUrole{std,std-ref}{OPT}}}}).  If
you feel that your contribution should be added to a different package,
please consult with the LAMMPS developers first.

\sphinxAtStartPar
The contributed code needs to support the {\hyperref[\detokenize{Build_make::doc}]{\sphinxcrossref{\DUrole{doc}{traditional GNU make
build process}}}} \sphinxstylestrong{and} the {\hyperref[\detokenize{Build_cmake::doc}]{\sphinxcrossref{\DUrole{doc}{CMake build process}}}}.  For the GNU make process and if the package has an
\sphinxcode{\sphinxupquote{Install.sh}} file, most likely that file needs to be updated to
correctly copy the sources when installing the package and properly
delete them when uninstalling.  This is particularly important when
added a new pair style that is a derived class from an existing pair
style in a package, so that its installation depends on the the
installation status of the package of the derived class.  For the CMake
process, it is sometimes necessary to make changes to the package
specific CMake scripting in \sphinxcode{\sphinxupquote{cmake/Modules/Packages}}.


\bigskip\hrule\bigskip



\subsection{Case 1: a pairwise additive model}
\label{\detokenize{Developer_write_pair:case-1-a-pairwise-additive-model}}
\sphinxAtStartPar
In this section, we will describe the procedure of adding a simple pair
style to LAMMPS: an empirical model that can be used to model liquid
mercury.  The pair style shall be called {\hyperref[\detokenize{pair_born_gauss::doc}]{\sphinxcrossref{\DUrole{doc}{bond/gauss}}}} and the complete implementation can be found in the
files \sphinxcode{\sphinxupquote{src/EXTRA\sphinxhyphen{}PAIR/pair\_born\_gauss.cpp}} and
\sphinxcode{\sphinxupquote{src/EXTRA\sphinxhyphen{}PAIR/pair\_born\_gauss.h}} of the LAMMPS source code.


\subsubsection{Model and general considerations}
\label{\detokenize{Developer_write_pair:model-and-general-considerations}}
\sphinxAtStartPar
The functional form of the model according to {\hyperref[\detokenize{Developer_write_pair:bomont2}]{\sphinxcrossref{\DUrole{std,std-ref}{(Bomont)}}}}
consists of a repulsive Born\sphinxhyphen{}Mayer exponential term and a temperature
dependent, attractive Gaussian term.
\begin{equation*}
\begin{split}E = A_0 \exp \left( -\alpha r \right) - A_1 \exp\left[ -\beta \left(r - r_0 \right)^2 \right]\end{split}
\end{equation*}
\sphinxAtStartPar
For the application to mercury, the following parameters are listed:
\begin{itemize}
\item {} 
\sphinxAtStartPar
\(A_0 = 8.2464 \times 10^{13} \; \textrm{eV}\)

\item {} 
\sphinxAtStartPar
\(\alpha = 12.48 \; \AA^{-1}\)

\item {} 
\sphinxAtStartPar
\(\beta = 0.44 \; \AA^{-2}\)

\item {} 
\sphinxAtStartPar
\(r_0 = 3.56 \; \AA\)

\item {} 
\sphinxAtStartPar
\(A_1\) is temperature dependent and can be determined from
\(A_1 = a_0 + a_1 T + a_2 T^2\) with:
\begin{itemize}
\item {} 
\sphinxAtStartPar
\(a_0 = 1.97475 \times 10^{-2} \; \textrm{eV}\)

\item {} 
\sphinxAtStartPar
\(a_1 = 8.40841 \times 10^{-5} \; \textrm{eV/K}\)

\item {} 
\sphinxAtStartPar
\(a_2 = -2.58717 \times 10^{-8} \; \textrm{eV/K}^{-2}\)

\end{itemize}

\end{itemize}

\sphinxAtStartPar
With the optional cutoff, this means we have a total of 5 or 6
parameters for each pair of atom types. Additionally, we need to input a
default cutoff value as a global setting.

\sphinxAtStartPar
Because of the combination of Born\sphinxhyphen{}Mayer with a Gaussian, the pair style
shall be named “born/gauss” and thus the class name would be
\sphinxcode{\sphinxupquote{PairBornGauss}} and the source files \sphinxcode{\sphinxupquote{pair\_born\_gauss.h}} and
\sphinxcode{\sphinxupquote{pair\_born\_gauss.cpp}}.  Since this is a rather uncommon potential, it
shall be added to the {\hyperref[\detokenize{Packages_details:pkg-extra-pair}]{\sphinxcrossref{\DUrole{std,std-ref}{EXTRA\sphinxhyphen{}PAIR}}}} package.


\subsubsection{Header file}
\label{\detokenize{Developer_write_pair:header-file}}
\sphinxAtStartPar
The first segment of any LAMMPS source should be the copyright and
license statement.  Note the marker in the first line to indicate to
editors like emacs that this file is a C++ source, even though the .h
extension suggests a C source (this is a convention inherited from the
very beginning of the C++ version of LAMMPS).

\begin{sphinxVerbatim}[commandchars=\\\{\}]
\PYG{c+cm}{/* \PYGZhy{}*\PYGZhy{} c++ \PYGZhy{}*\PYGZhy{} \PYGZhy{}\PYGZhy{}\PYGZhy{}\PYGZhy{}\PYGZhy{}\PYGZhy{}\PYGZhy{}\PYGZhy{}\PYGZhy{}\PYGZhy{}\PYGZhy{}\PYGZhy{}\PYGZhy{}\PYGZhy{}\PYGZhy{}\PYGZhy{}\PYGZhy{}\PYGZhy{}\PYGZhy{}\PYGZhy{}\PYGZhy{}\PYGZhy{}\PYGZhy{}\PYGZhy{}\PYGZhy{}\PYGZhy{}\PYGZhy{}\PYGZhy{}\PYGZhy{}\PYGZhy{}\PYGZhy{}\PYGZhy{}\PYGZhy{}\PYGZhy{}\PYGZhy{}\PYGZhy{}\PYGZhy{}\PYGZhy{}\PYGZhy{}\PYGZhy{}\PYGZhy{}\PYGZhy{}\PYGZhy{}\PYGZhy{}\PYGZhy{}\PYGZhy{}\PYGZhy{}\PYGZhy{}\PYGZhy{}\PYGZhy{}\PYGZhy{}\PYGZhy{}\PYGZhy{}\PYGZhy{}\PYGZhy{}\PYGZhy{}\PYGZhy{}\PYGZhy{}}
\PYG{c+cm}{   LAMMPS \PYGZhy{} Large\PYGZhy{}scale Atomic/Molecular Massively Parallel Simulator}
\PYG{c+cm}{   https://www.lammps.org/, Sandia National Laboratories}
\PYG{c+cm}{   LAMMPS development team: developers@lammps.org}

\PYG{c+cm}{   Copyright (2003) Sandia Corporation.  Under the terms of Contract}
\PYG{c+cm}{   DE\PYGZhy{}AC04\PYGZhy{}94AL85000 with Sandia Corporation, the U.S. Government retains}
\PYG{c+cm}{   certain rights in this software.  This software is distributed under}
\PYG{c+cm}{   the GNU General Public License.}

\PYG{c+cm}{   See the README file in the top\PYGZhy{}level LAMMPS directory.}
\PYG{c+cm}{\PYGZhy{}\PYGZhy{}\PYGZhy{}\PYGZhy{}\PYGZhy{}\PYGZhy{}\PYGZhy{}\PYGZhy{}\PYGZhy{}\PYGZhy{}\PYGZhy{}\PYGZhy{}\PYGZhy{}\PYGZhy{}\PYGZhy{}\PYGZhy{}\PYGZhy{}\PYGZhy{}\PYGZhy{}\PYGZhy{}\PYGZhy{}\PYGZhy{}\PYGZhy{}\PYGZhy{}\PYGZhy{}\PYGZhy{}\PYGZhy{}\PYGZhy{}\PYGZhy{}\PYGZhy{}\PYGZhy{}\PYGZhy{}\PYGZhy{}\PYGZhy{}\PYGZhy{}\PYGZhy{}\PYGZhy{}\PYGZhy{}\PYGZhy{}\PYGZhy{}\PYGZhy{}\PYGZhy{}\PYGZhy{}\PYGZhy{}\PYGZhy{}\PYGZhy{}\PYGZhy{}\PYGZhy{}\PYGZhy{}\PYGZhy{}\PYGZhy{}\PYGZhy{}\PYGZhy{}\PYGZhy{}\PYGZhy{}\PYGZhy{}\PYGZhy{}\PYGZhy{}\PYGZhy{}\PYGZhy{}\PYGZhy{}\PYGZhy{}\PYGZhy{}\PYGZhy{}\PYGZhy{}\PYGZhy{}\PYGZhy{}\PYGZhy{}\PYGZhy{}\PYGZhy{}\PYGZhy{}\PYGZhy{}\PYGZhy{} */}
\end{sphinxVerbatim}

\sphinxAtStartPar
Every pair style must be registered in LAMMPS by including the following
lines of code in the second part of the header after the copyright
message and before the include guards for the class definition:

\begin{sphinxVerbatim}[commandchars=\\\{\}]
\PYG{c+cp}{\PYGZsh{}}\PYG{c+cp}{ifdef PAIR\PYGZus{}CLASS}
\PYG{c+c1}{// clang\PYGZhy{}format off}
\PYG{n}{PairStyle}\PYG{p}{(}\PYG{n}{born}\PYG{o}{/}\PYG{n}{gauss}\PYG{p}{,}\PYG{n}{PairBornGauss}\PYG{p}{)}\PYG{p}{;}
\PYG{c+c1}{// clang\PYGZhy{}format on}
\PYG{c+cp}{\PYGZsh{}}\PYG{c+cp}{else}

\PYG{c+cm}{/* the definition of the PairBornGauss class (see below) is inserted here */}

\PYG{c+cp}{\PYGZsh{}}\PYG{c+cp}{endif}
\end{sphinxVerbatim}

\sphinxAtStartPar
This block of between \sphinxcode{\sphinxupquote{\#ifdef PAIR\_CLASS}} and \sphinxcode{\sphinxupquote{\#else}} will be
included by the \sphinxcode{\sphinxupquote{Force}} class in \sphinxcode{\sphinxupquote{force.cpp}} to build a map of
“factory functions” that will create an instance of these classes and
return a pointer to it.  The map connects the name of the pair style,
“born/gauss”, to the name of the class, \sphinxcode{\sphinxupquote{PairBornGauss}}.  During
compilation, LAMMPS constructs a file \sphinxcode{\sphinxupquote{style\_pair.h}} that contains
\sphinxcode{\sphinxupquote{\#include}} statements for all “installed” pair styles.  Before
including \sphinxcode{\sphinxupquote{style\_pair.h}} into \sphinxcode{\sphinxupquote{force.cpp}}, the \sphinxcode{\sphinxupquote{PAIR\_CLASS}} define
is set and the \sphinxcode{\sphinxupquote{PairStyle(name,class)}} macro defined.  The code of the
macro adds the installed pair styles to the “factory map” which enables
the {\hyperref[\detokenize{pair_style::doc}]{\sphinxcrossref{\DUrole{doc}{pair\_style command}}}} to create the pair style
instance.

\sphinxAtStartPar
The list of header files to include is automatically updated by the
build system if there are new files, so the presence of the new header
file in the \sphinxcode{\sphinxupquote{src/EXTRA\sphinxhyphen{}PAIR}} folder and the enabling of the EXTRA\sphinxhyphen{}PAIR
package will trigger LAMMPS to include the new pair style when it is
(re\sphinxhyphen{})compiled.  The “// clang\sphinxhyphen{}format” format comments are needed so that
running {\hyperref[\detokenize{Build_development:clang-format}]{\sphinxcrossref{\DUrole{std,std-ref}{clang\sphinxhyphen{}format}}}} on the file will not insert
unwanted blanks between “born”, “/”, and “gauss” which would break the
\sphinxcode{\sphinxupquote{PairStyle}} macro.

\sphinxAtStartPar
The third part of the header file is the actual class definition of the
\sphinxcode{\sphinxupquote{PairBornGauss}} class.  This has the prototypes for all member
functions that will be implemented by this pair style.  This includes
{\hyperref[\detokenize{Modify_pair::doc}]{\sphinxcrossref{\DUrole{doc}{a few required and a number of optional functions}}}}.
All functions that were labeled in the base class as “virtual” must be
given the “override” property, as it is done in the code shown below.

\sphinxAtStartPar
The “override” property helps to detect unexpected mismatches because
compilation will stop with an error in case the signature of a function
is changed in the base class without also changing it in all derived
classes.  For example, if this change added an optional argument with a
default value, then all existing source code \sphinxstyleemphasis{calling} the function
would not need changes and still compile, but the function in the
derived class would no longer override the one in the base class due to
the different number of arguments and the behavior of the pair style is
thus changed in an unintended way.  Using the “override” keyword
prevents such issues.

\begin{sphinxVerbatim}[commandchars=\\\{\}]
\PYG{c+cp}{\PYGZsh{}}\PYG{c+cp}{ifndef LMP\PYGZus{}PAIR\PYGZus{}BORN\PYGZus{}GAUSS\PYGZus{}H}
\PYG{c+cp}{\PYGZsh{}}\PYG{c+cp}{define LMP\PYGZus{}PAIR\PYGZus{}BORN\PYGZus{}GAUSS\PYGZus{}H}

\PYG{c+cp}{\PYGZsh{}}\PYG{c+cp}{include}\PYG{+w}{ }\PYG{c+cpf}{\PYGZdq{}pair.h\PYGZdq{}}

\PYG{k}{namespace}\PYG{+w}{ }\PYG{n+nn}{LAMMPS\PYGZus{}NS}\PYG{+w}{ }\PYG{p}{\PYGZob{}}

\PYG{k}{class}\PYG{+w}{ }\PYG{n+nc}{PairBornGauss}\PYG{+w}{ }\PYG{o}{:}\PYG{+w}{ }\PYG{k}{public}\PYG{+w}{ }\PYG{n}{Pair}\PYG{+w}{ }\PYG{p}{\PYGZob{}}
\PYG{+w}{ }\PYG{k}{public}\PYG{o}{:}
\PYG{+w}{  }\PYG{n}{PairBornGauss}\PYG{p}{(}\PYG{k}{class}\PYG{+w}{ }\PYG{n+nc}{LAMMPS}\PYG{+w}{ }\PYG{o}{*}\PYG{p}{)}\PYG{p}{;}
\PYG{+w}{  }\PYG{o}{\PYGZti{}}\PYG{n}{PairBornGauss}\PYG{p}{(}\PYG{p}{)}\PYG{+w}{ }\PYG{k}{override}\PYG{p}{;}

\PYG{+w}{  }\PYG{k+kt}{void}\PYG{+w}{ }\PYG{n+nf}{compute}\PYG{p}{(}\PYG{k+kt}{int}\PYG{p}{,}\PYG{+w}{ }\PYG{k+kt}{int}\PYG{p}{)}\PYG{+w}{ }\PYG{k}{override}\PYG{p}{;}
\PYG{+w}{  }\PYG{k+kt}{void}\PYG{+w}{ }\PYG{n+nf}{settings}\PYG{p}{(}\PYG{k+kt}{int}\PYG{p}{,}\PYG{+w}{ }\PYG{k+kt}{char}\PYG{+w}{ }\PYG{o}{*}\PYG{o}{*}\PYG{p}{)}\PYG{+w}{ }\PYG{k}{override}\PYG{p}{;}
\PYG{+w}{  }\PYG{k+kt}{void}\PYG{+w}{ }\PYG{n+nf}{coeff}\PYG{p}{(}\PYG{k+kt}{int}\PYG{p}{,}\PYG{+w}{ }\PYG{k+kt}{char}\PYG{+w}{ }\PYG{o}{*}\PYG{o}{*}\PYG{p}{)}\PYG{+w}{ }\PYG{k}{override}\PYG{p}{;}
\PYG{+w}{  }\PYG{k+kt}{double}\PYG{+w}{ }\PYG{n+nf}{init\PYGZus{}one}\PYG{p}{(}\PYG{k+kt}{int}\PYG{p}{,}\PYG{+w}{ }\PYG{k+kt}{int}\PYG{p}{)}\PYG{+w}{ }\PYG{k}{override}\PYG{p}{;}

\PYG{+w}{  }\PYG{k+kt}{void}\PYG{+w}{ }\PYG{n+nf}{write\PYGZus{}restart}\PYG{p}{(}\PYG{k+kt}{FILE}\PYG{+w}{ }\PYG{o}{*}\PYG{p}{)}\PYG{+w}{ }\PYG{k}{override}\PYG{p}{;}
\PYG{+w}{  }\PYG{k+kt}{void}\PYG{+w}{ }\PYG{n+nf}{read\PYGZus{}restart}\PYG{p}{(}\PYG{k+kt}{FILE}\PYG{+w}{ }\PYG{o}{*}\PYG{p}{)}\PYG{+w}{ }\PYG{k}{override}\PYG{p}{;}
\PYG{+w}{  }\PYG{k+kt}{void}\PYG{+w}{ }\PYG{n+nf}{write\PYGZus{}restart\PYGZus{}settings}\PYG{p}{(}\PYG{k+kt}{FILE}\PYG{+w}{ }\PYG{o}{*}\PYG{p}{)}\PYG{+w}{ }\PYG{k}{override}\PYG{p}{;}
\PYG{+w}{  }\PYG{k+kt}{void}\PYG{+w}{ }\PYG{n+nf}{read\PYGZus{}restart\PYGZus{}settings}\PYG{p}{(}\PYG{k+kt}{FILE}\PYG{+w}{ }\PYG{o}{*}\PYG{p}{)}\PYG{+w}{ }\PYG{k}{override}\PYG{p}{;}
\PYG{+w}{  }\PYG{k+kt}{void}\PYG{+w}{ }\PYG{n+nf}{write\PYGZus{}data}\PYG{p}{(}\PYG{k+kt}{FILE}\PYG{+w}{ }\PYG{o}{*}\PYG{p}{)}\PYG{+w}{ }\PYG{k}{override}\PYG{p}{;}
\PYG{+w}{  }\PYG{k+kt}{void}\PYG{+w}{ }\PYG{n+nf}{write\PYGZus{}data\PYGZus{}all}\PYG{p}{(}\PYG{k+kt}{FILE}\PYG{+w}{ }\PYG{o}{*}\PYG{p}{)}\PYG{+w}{ }\PYG{k}{override}\PYG{p}{;}

\PYG{+w}{  }\PYG{k+kt}{double}\PYG{+w}{ }\PYG{n+nf}{single}\PYG{p}{(}\PYG{k+kt}{int}\PYG{p}{,}\PYG{+w}{ }\PYG{k+kt}{int}\PYG{p}{,}\PYG{+w}{ }\PYG{k+kt}{int}\PYG{p}{,}\PYG{+w}{ }\PYG{k+kt}{int}\PYG{p}{,}\PYG{+w}{ }\PYG{k+kt}{double}\PYG{p}{,}\PYG{+w}{ }\PYG{k+kt}{double}\PYG{p}{,}\PYG{+w}{ }\PYG{k+kt}{double}\PYG{p}{,}\PYG{+w}{ }\PYG{k+kt}{double}\PYG{+w}{ }\PYG{o}{\PYGZam{}}\PYG{p}{)}\PYG{+w}{ }\PYG{k}{override}\PYG{p}{;}
\PYG{+w}{  }\PYG{k+kt}{void}\PYG{+w}{ }\PYG{o}{*}\PYG{n+nf}{extract}\PYG{p}{(}\PYG{k}{const}\PYG{+w}{ }\PYG{k+kt}{char}\PYG{+w}{ }\PYG{o}{*}\PYG{p}{,}\PYG{+w}{ }\PYG{k+kt}{int}\PYG{+w}{ }\PYG{o}{\PYGZam{}}\PYG{p}{)}\PYG{+w}{ }\PYG{k}{override}\PYG{p}{;}
\end{sphinxVerbatim}

\sphinxAtStartPar
Also, variables and arrays for storing global settings and potential
parameters are defined.  Since these are internal to the class, they are
placed after a “protected:” label.

\begin{sphinxVerbatim}[commandchars=\\\{\}]
\PYG{+w}{ }\PYG{k}{protected}\PYG{o}{:}
\PYG{+w}{  }\PYG{k+kt}{double}\PYG{+w}{ }\PYG{n}{cut\PYGZus{}global}\PYG{p}{;}
\PYG{+w}{  }\PYG{k+kt}{double}\PYG{+w}{ }\PYG{o}{*}\PYG{o}{*}\PYG{n}{cut}\PYG{p}{;}
\PYG{+w}{  }\PYG{k+kt}{double}\PYG{+w}{ }\PYG{o}{*}\PYG{o}{*}\PYG{n}{biga0}\PYG{p}{,}\PYG{+w}{ }\PYG{o}{*}\PYG{o}{*}\PYG{n}{alpha}\PYG{p}{,}\PYG{+w}{ }\PYG{o}{*}\PYG{o}{*}\PYG{n}{biga1}\PYG{p}{,}\PYG{+w}{ }\PYG{o}{*}\PYG{o}{*}\PYG{n}{beta}\PYG{p}{,}\PYG{+w}{ }\PYG{o}{*}\PYG{o}{*}\PYG{n}{r0}\PYG{p}{;}
\PYG{+w}{  }\PYG{k+kt}{double}\PYG{+w}{ }\PYG{o}{*}\PYG{o}{*}\PYG{n}{a0}\PYG{p}{,}\PYG{+w}{ }\PYG{o}{*}\PYG{o}{*}\PYG{n}{a1}\PYG{p}{,}\PYG{+w}{ }\PYG{o}{*}\PYG{o}{*}\PYG{n}{a2}\PYG{p}{;}
\PYG{+w}{  }\PYG{k+kt}{double}\PYG{+w}{ }\PYG{o}{*}\PYG{o}{*}\PYG{n}{offset}\PYG{p}{;}

\PYG{+w}{  }\PYG{k}{virtual}\PYG{+w}{ }\PYG{k+kt}{void}\PYG{+w}{ }\PYG{n+nf}{allocate}\PYG{p}{(}\PYG{p}{)}\PYG{p}{;}
\PYG{p}{\PYGZcb{}}\PYG{p}{;}
\PYG{p}{\PYGZcb{}}\PYG{+w}{    }\PYG{c+c1}{// namespace LAMMPS\PYGZus{}NS}
\PYG{c+cp}{\PYGZsh{}}\PYG{c+cp}{endif}
\end{sphinxVerbatim}


\subsubsection{Implementation file}
\label{\detokenize{Developer_write_pair:implementation-file}}
\sphinxAtStartPar
We move on to the implementation of the \sphinxcode{\sphinxupquote{PairBornGauss}} class in the
\sphinxcode{\sphinxupquote{pair\_born\_gauss.cpp}} file.  This file also starts with a LAMMPS
copyright and license header.  Below that notice is typically the space
where comments may be added with additional information about this
specific file, the author(s), affiliation(s), and email address(es).
This way the contributing author(s) can be easily contacted, when
there are questions about the implementation later.  Since the file(s)
may be around for a long time, it is beneficial to use some kind of
“permanent” email address, if possible.

\begin{sphinxVerbatim}[commandchars=\\\{\}]
\PYG{c+cm}{/* \PYGZhy{}\PYGZhy{}\PYGZhy{}\PYGZhy{}\PYGZhy{}\PYGZhy{}\PYGZhy{}\PYGZhy{}\PYGZhy{}\PYGZhy{}\PYGZhy{}\PYGZhy{}\PYGZhy{}\PYGZhy{}\PYGZhy{}\PYGZhy{}\PYGZhy{}\PYGZhy{}\PYGZhy{}\PYGZhy{}\PYGZhy{}\PYGZhy{}\PYGZhy{}\PYGZhy{}\PYGZhy{}\PYGZhy{}\PYGZhy{}\PYGZhy{}\PYGZhy{}\PYGZhy{}\PYGZhy{}\PYGZhy{}\PYGZhy{}\PYGZhy{}\PYGZhy{}\PYGZhy{}\PYGZhy{}\PYGZhy{}\PYGZhy{}\PYGZhy{}\PYGZhy{}\PYGZhy{}\PYGZhy{}\PYGZhy{}\PYGZhy{}\PYGZhy{}\PYGZhy{}\PYGZhy{}\PYGZhy{}\PYGZhy{}\PYGZhy{}\PYGZhy{}\PYGZhy{}\PYGZhy{}\PYGZhy{}\PYGZhy{}\PYGZhy{}\PYGZhy{}\PYGZhy{}\PYGZhy{}\PYGZhy{}\PYGZhy{}\PYGZhy{}\PYGZhy{}\PYGZhy{}\PYGZhy{}\PYGZhy{}\PYGZhy{}\PYGZhy{}\PYGZhy{}}
\PYG{c+cm}{   LAMMPS \PYGZhy{} Large\PYGZhy{}scale Atomic/Molecular Massively Parallel Simulator}
\PYG{c+cm}{   https://www.lammps.org/, Sandia National Laboratories}
\PYG{c+cm}{   LAMMPS development team: developers@lammps.org}

\PYG{c+cm}{   Copyright (2003) Sandia Corporation.  Under the terms of Contract}
\PYG{c+cm}{   DE\PYGZhy{}AC04\PYGZhy{}94AL85000 with Sandia Corporation, the U.S. Government retains}
\PYG{c+cm}{   certain rights in this software.  This software is distributed under}
\PYG{c+cm}{   the GNU General Public License.}

\PYG{c+cm}{   See the README file in the top\PYGZhy{}level LAMMPS directory.}
\PYG{c+cm}{\PYGZhy{}\PYGZhy{}\PYGZhy{}\PYGZhy{}\PYGZhy{}\PYGZhy{}\PYGZhy{}\PYGZhy{}\PYGZhy{}\PYGZhy{}\PYGZhy{}\PYGZhy{}\PYGZhy{}\PYGZhy{}\PYGZhy{}\PYGZhy{}\PYGZhy{}\PYGZhy{}\PYGZhy{}\PYGZhy{}\PYGZhy{}\PYGZhy{}\PYGZhy{}\PYGZhy{}\PYGZhy{}\PYGZhy{}\PYGZhy{}\PYGZhy{}\PYGZhy{}\PYGZhy{}\PYGZhy{}\PYGZhy{}\PYGZhy{}\PYGZhy{}\PYGZhy{}\PYGZhy{}\PYGZhy{}\PYGZhy{}\PYGZhy{}\PYGZhy{}\PYGZhy{}\PYGZhy{}\PYGZhy{}\PYGZhy{}\PYGZhy{}\PYGZhy{}\PYGZhy{}\PYGZhy{}\PYGZhy{}\PYGZhy{}\PYGZhy{}\PYGZhy{}\PYGZhy{}\PYGZhy{}\PYGZhy{}\PYGZhy{}\PYGZhy{}\PYGZhy{}\PYGZhy{}\PYGZhy{}\PYGZhy{}\PYGZhy{}\PYGZhy{}\PYGZhy{}\PYGZhy{}\PYGZhy{}\PYGZhy{}\PYGZhy{}\PYGZhy{}\PYGZhy{}\PYGZhy{}\PYGZhy{}\PYGZhy{} */}

\PYG{c+c1}{// Contributing author: Axel Kohlmeyer, Temple University, akohlmey@gmail.com}

\PYG{c+cp}{\PYGZsh{}}\PYG{c+cp}{include}\PYG{+w}{ }\PYG{c+cpf}{\PYGZdq{}pair\PYGZus{}born\PYGZus{}gauss.h\PYGZdq{}}

\PYG{c+cp}{\PYGZsh{}}\PYG{c+cp}{include}\PYG{+w}{ }\PYG{c+cpf}{\PYGZdq{}atom.h\PYGZdq{}}
\PYG{c+cp}{\PYGZsh{}}\PYG{c+cp}{include}\PYG{+w}{ }\PYG{c+cpf}{\PYGZdq{}comm.h\PYGZdq{}}
\PYG{c+cp}{\PYGZsh{}}\PYG{c+cp}{include}\PYG{+w}{ }\PYG{c+cpf}{\PYGZdq{}error.h\PYGZdq{}}
\PYG{c+cp}{\PYGZsh{}}\PYG{c+cp}{include}\PYG{+w}{ }\PYG{c+cpf}{\PYGZdq{}fix.h\PYGZdq{}}
\PYG{c+cp}{\PYGZsh{}}\PYG{c+cp}{include}\PYG{+w}{ }\PYG{c+cpf}{\PYGZdq{}force.h\PYGZdq{}}
\PYG{c+cp}{\PYGZsh{}}\PYG{c+cp}{include}\PYG{+w}{ }\PYG{c+cpf}{\PYGZdq{}memory.h\PYGZdq{}}
\PYG{c+cp}{\PYGZsh{}}\PYG{c+cp}{include}\PYG{+w}{ }\PYG{c+cpf}{\PYGZdq{}neigh\PYGZus{}list.h\PYGZdq{}}

\PYG{c+cp}{\PYGZsh{}}\PYG{c+cp}{include}\PYG{+w}{ }\PYG{c+cpf}{\PYGZlt{}cmath\PYGZgt{}}
\PYG{c+cp}{\PYGZsh{}}\PYG{c+cp}{include}\PYG{+w}{ }\PYG{c+cpf}{\PYGZlt{}cstring\PYGZgt{}}

\PYG{k}{using}\PYG{+w}{ }\PYG{k}{namespace}\PYG{+w}{ }\PYG{n+nn}{LAMMPS\PYGZus{}NS}\PYG{p}{;}
\end{sphinxVerbatim}

\sphinxAtStartPar
The second section of the implementation file has various include
statements.  The include file for the class header has to come first,
then a block of LAMMPS classes (sorted alphabetically) followed by a
block of system headers and others, if needed.  Note the standardized
C++ notation for headers of C\sphinxhyphen{}library functions (\sphinxcode{\sphinxupquote{cmath}} instead of
\sphinxcode{\sphinxupquote{math.h}}).  The final statement of this segment imports the
\sphinxcode{\sphinxupquote{LAMMPS\_NS::}} namespace globally for this file.  This way, all LAMMPS
specific functions and classes do not have to be prefixed with
\sphinxcode{\sphinxupquote{LAMMPS\_NS::}}.


\subsubsection{Constructor and destructor (required)}
\label{\detokenize{Developer_write_pair:constructor-and-destructor-required}}
\sphinxAtStartPar
The first two functions in the implementation source file are typically
the constructor and the destructor.

\sphinxAtStartPar
Pair styles are different from most classes in LAMMPS that define a
“style”, as their constructor only uses the LAMMPS class instance
pointer as an argument, but \sphinxstylestrong{not} the command line arguments of the
{\hyperref[\detokenize{pair_style::doc}]{\sphinxcrossref{\DUrole{doc}{pair\_style command}}}}.  Instead, those arguments are
processed in the \sphinxcode{\sphinxupquote{Pair::settings()}} function (or rather the version in
the derived class).  The constructor is the place where global defaults
are set and specifically flags are set indicating which optional
features of a pair style are available.

\begin{sphinxVerbatim}[commandchars=\\\{\}]
\PYG{c+cm}{/* \PYGZhy{}\PYGZhy{}\PYGZhy{}\PYGZhy{}\PYGZhy{}\PYGZhy{}\PYGZhy{}\PYGZhy{}\PYGZhy{}\PYGZhy{}\PYGZhy{}\PYGZhy{}\PYGZhy{}\PYGZhy{}\PYGZhy{}\PYGZhy{}\PYGZhy{}\PYGZhy{}\PYGZhy{}\PYGZhy{}\PYGZhy{}\PYGZhy{}\PYGZhy{}\PYGZhy{}\PYGZhy{}\PYGZhy{}\PYGZhy{}\PYGZhy{}\PYGZhy{}\PYGZhy{}\PYGZhy{}\PYGZhy{}\PYGZhy{}\PYGZhy{}\PYGZhy{}\PYGZhy{}\PYGZhy{}\PYGZhy{}\PYGZhy{}\PYGZhy{}\PYGZhy{}\PYGZhy{}\PYGZhy{}\PYGZhy{}\PYGZhy{}\PYGZhy{}\PYGZhy{}\PYGZhy{}\PYGZhy{}\PYGZhy{}\PYGZhy{}\PYGZhy{}\PYGZhy{}\PYGZhy{}\PYGZhy{}\PYGZhy{}\PYGZhy{}\PYGZhy{}\PYGZhy{}\PYGZhy{}\PYGZhy{}\PYGZhy{}\PYGZhy{}\PYGZhy{}\PYGZhy{}\PYGZhy{}\PYGZhy{}\PYGZhy{}\PYGZhy{}\PYGZhy{} */}

\PYG{n}{PairBornGauss}\PYG{o}{:}\PYG{o}{:}\PYG{n}{PairBornGauss}\PYG{p}{(}\PYG{n}{LAMMPS}\PYG{+w}{ }\PYG{o}{*}\PYG{n}{lmp}\PYG{p}{)}\PYG{+w}{ }\PYG{o}{:}\PYG{+w}{ }\PYG{n}{Pair}\PYG{p}{(}\PYG{n}{lmp}\PYG{p}{)}
\PYG{p}{\PYGZob{}}
\PYG{+w}{  }\PYG{n}{writedata}\PYG{+w}{ }\PYG{o}{=}\PYG{+w}{ }\PYG{l+m+mi}{1}\PYG{p}{;}
\PYG{p}{\PYGZcb{}}
\end{sphinxVerbatim}

\sphinxAtStartPar
The \sphinxtitleref{writedata = 1;} statement indicates that the pair style is capable
of writing the current pair coefficient parameters to data files.  That
is, the class implements specific versions for \sphinxcode{\sphinxupquote{Pair::data()}} and
\sphinxcode{\sphinxupquote{Pair::data\_all()}}.  Other statements that could be added here would
be \sphinxtitleref{single\_enable = 1;} or \sphinxtitleref{respa\_enable = 0;} to indicate that the
\sphinxcode{\sphinxupquote{Pair::single()}} function is present and the
\sphinxcode{\sphinxupquote{Pair::compute\_(inner|middle|outer)}} functions are not, but those are
also the default settings and already set in the base class.

\sphinxAtStartPar
In the destructor, we need to delete all memory that was allocated by the
pair style, usually to hold force field parameters that were entered
with the {\hyperref[\detokenize{pair_coeff::doc}]{\sphinxcrossref{\DUrole{doc}{pair\_coeff command}}}}.  Most of those array
pointers will need to be declared in the derived class header, but some
(e.g. setflag, cutsq) are already declared in the base class.

\begin{sphinxVerbatim}[commandchars=\\\{\}]
\PYG{n}{PairBornGauss}\PYG{o}{:}\PYG{o}{:}\PYG{o}{\PYGZti{}}\PYG{n}{PairBornGauss}\PYG{p}{(}\PYG{p}{)}
\PYG{p}{\PYGZob{}}
\PYG{+w}{  }\PYG{k}{if}\PYG{+w}{ }\PYG{p}{(}\PYG{n}{allocated}\PYG{p}{)}\PYG{+w}{ }\PYG{p}{\PYGZob{}}
\PYG{+w}{    }\PYG{n}{memory}\PYG{o}{\PYGZhy{}}\PYG{o}{\PYGZgt{}}\PYG{n}{destroy}\PYG{p}{(}\PYG{n}{setflag}\PYG{p}{)}\PYG{p}{;}
\PYG{+w}{    }\PYG{n}{memory}\PYG{o}{\PYGZhy{}}\PYG{o}{\PYGZgt{}}\PYG{n}{destroy}\PYG{p}{(}\PYG{n}{cutsq}\PYG{p}{)}\PYG{p}{;}
\PYG{+w}{    }\PYG{n}{memory}\PYG{o}{\PYGZhy{}}\PYG{o}{\PYGZgt{}}\PYG{n}{destroy}\PYG{p}{(}\PYG{n}{cut}\PYG{p}{)}\PYG{p}{;}
\PYG{+w}{    }\PYG{n}{memory}\PYG{o}{\PYGZhy{}}\PYG{o}{\PYGZgt{}}\PYG{n}{destroy}\PYG{p}{(}\PYG{n}{biga0}\PYG{p}{)}\PYG{p}{;}
\PYG{+w}{    }\PYG{n}{memory}\PYG{o}{\PYGZhy{}}\PYG{o}{\PYGZgt{}}\PYG{n}{destroy}\PYG{p}{(}\PYG{n}{alpha}\PYG{p}{)}\PYG{p}{;}
\PYG{+w}{    }\PYG{n}{memory}\PYG{o}{\PYGZhy{}}\PYG{o}{\PYGZgt{}}\PYG{n}{destroy}\PYG{p}{(}\PYG{n}{biga1}\PYG{p}{)}\PYG{p}{;}
\PYG{+w}{    }\PYG{n}{memory}\PYG{o}{\PYGZhy{}}\PYG{o}{\PYGZgt{}}\PYG{n}{destroy}\PYG{p}{(}\PYG{n}{beta}\PYG{p}{)}\PYG{p}{;}
\PYG{+w}{    }\PYG{n}{memory}\PYG{o}{\PYGZhy{}}\PYG{o}{\PYGZgt{}}\PYG{n}{destroy}\PYG{p}{(}\PYG{n}{r0}\PYG{p}{)}\PYG{p}{;}
\PYG{+w}{    }\PYG{n}{memory}\PYG{o}{\PYGZhy{}}\PYG{o}{\PYGZgt{}}\PYG{n}{destroy}\PYG{p}{(}\PYG{n}{offset}\PYG{p}{)}\PYG{p}{;}
\PYG{+w}{  }\PYG{p}{\PYGZcb{}}
\PYG{p}{\PYGZcb{}}
\end{sphinxVerbatim}


\subsubsection{Settings and coefficients (required)}
\label{\detokenize{Developer_write_pair:settings-and-coefficients-required}}
\sphinxAtStartPar
To enter the global pair style settings and the pair style parameters,
the functions \sphinxcode{\sphinxupquote{Pair::settings()}} and \sphinxcode{\sphinxupquote{Pair::coeff()}} need to be
re\sphinxhyphen{}implemented.  The arguments to the \sphinxcode{\sphinxupquote{settings()}} function are the
arguments given to the {\hyperref[\detokenize{pair_style::doc}]{\sphinxcrossref{\DUrole{doc}{pair\_style command}}}}.
Normally, those would already be processed as part of the constructor,
but moving this to a separate function allows users to change global
settings like the default cutoff without having to reissue all
pair\_coeff commands or re\sphinxhyphen{}read the \sphinxcode{\sphinxupquote{Pair Coeffs}} sections from the
data file.  In the \sphinxcode{\sphinxupquote{settings()}} function, also the arrays for storing
parameters, to define cutoffs, track which pairs of parameters have been
explicitly set and allocated and, if needed, initialized.  In this case,
the memory allocation and initialization are moved to a function
\sphinxcode{\sphinxupquote{allocate()}}.

\begin{sphinxVerbatim}[commandchars=\\\{\}]
\PYG{c+cm}{/* \PYGZhy{}\PYGZhy{}\PYGZhy{}\PYGZhy{}\PYGZhy{}\PYGZhy{}\PYGZhy{}\PYGZhy{}\PYGZhy{}\PYGZhy{}\PYGZhy{}\PYGZhy{}\PYGZhy{}\PYGZhy{}\PYGZhy{}\PYGZhy{}\PYGZhy{}\PYGZhy{}\PYGZhy{}\PYGZhy{}\PYGZhy{}\PYGZhy{}\PYGZhy{}\PYGZhy{}\PYGZhy{}\PYGZhy{}\PYGZhy{}\PYGZhy{}\PYGZhy{}\PYGZhy{}\PYGZhy{}\PYGZhy{}\PYGZhy{}\PYGZhy{}\PYGZhy{}\PYGZhy{}\PYGZhy{}\PYGZhy{}\PYGZhy{}\PYGZhy{}\PYGZhy{}\PYGZhy{}\PYGZhy{}\PYGZhy{}\PYGZhy{}\PYGZhy{}\PYGZhy{}\PYGZhy{}\PYGZhy{}\PYGZhy{}\PYGZhy{}\PYGZhy{}\PYGZhy{}\PYGZhy{}\PYGZhy{}\PYGZhy{}\PYGZhy{}\PYGZhy{}\PYGZhy{}\PYGZhy{}\PYGZhy{}\PYGZhy{}\PYGZhy{}\PYGZhy{}\PYGZhy{}\PYGZhy{}\PYGZhy{}\PYGZhy{}\PYGZhy{}\PYGZhy{}}
\PYG{c+cm}{   allocate all arrays}
\PYG{c+cm}{\PYGZhy{}\PYGZhy{}\PYGZhy{}\PYGZhy{}\PYGZhy{}\PYGZhy{}\PYGZhy{}\PYGZhy{}\PYGZhy{}\PYGZhy{}\PYGZhy{}\PYGZhy{}\PYGZhy{}\PYGZhy{}\PYGZhy{}\PYGZhy{}\PYGZhy{}\PYGZhy{}\PYGZhy{}\PYGZhy{}\PYGZhy{}\PYGZhy{}\PYGZhy{}\PYGZhy{}\PYGZhy{}\PYGZhy{}\PYGZhy{}\PYGZhy{}\PYGZhy{}\PYGZhy{}\PYGZhy{}\PYGZhy{}\PYGZhy{}\PYGZhy{}\PYGZhy{}\PYGZhy{}\PYGZhy{}\PYGZhy{}\PYGZhy{}\PYGZhy{}\PYGZhy{}\PYGZhy{}\PYGZhy{}\PYGZhy{}\PYGZhy{}\PYGZhy{}\PYGZhy{}\PYGZhy{}\PYGZhy{}\PYGZhy{}\PYGZhy{}\PYGZhy{}\PYGZhy{}\PYGZhy{}\PYGZhy{}\PYGZhy{}\PYGZhy{}\PYGZhy{}\PYGZhy{}\PYGZhy{}\PYGZhy{}\PYGZhy{}\PYGZhy{}\PYGZhy{}\PYGZhy{}\PYGZhy{}\PYGZhy{}\PYGZhy{}\PYGZhy{}\PYGZhy{}\PYGZhy{}\PYGZhy{}\PYGZhy{} */}

\PYG{k+kt}{void}\PYG{+w}{ }\PYG{n+nf}{PairBornGauss::allocate}\PYG{p}{(}\PYG{p}{)}
\PYG{p}{\PYGZob{}}
\PYG{+w}{  }\PYG{n}{allocated}\PYG{+w}{ }\PYG{o}{=}\PYG{+w}{ }\PYG{l+m+mi}{1}\PYG{p}{;}
\PYG{+w}{  }\PYG{k+kt}{int}\PYG{+w}{ }\PYG{n}{np1}\PYG{+w}{ }\PYG{o}{=}\PYG{+w}{ }\PYG{n}{atom}\PYG{o}{\PYGZhy{}}\PYG{o}{\PYGZgt{}}\PYG{n}{ntypes}\PYG{+w}{ }\PYG{o}{+}\PYG{+w}{ }\PYG{l+m+mi}{1}\PYG{p}{;}

\PYG{+w}{  }\PYG{n}{memory}\PYG{o}{\PYGZhy{}}\PYG{o}{\PYGZgt{}}\PYG{n}{create}\PYG{p}{(}\PYG{n}{setflag}\PYG{p}{,}\PYG{+w}{ }\PYG{n}{np1}\PYG{p}{,}\PYG{+w}{ }\PYG{n}{np1}\PYG{p}{,}\PYG{+w}{ }\PYG{l+s}{\PYGZdq{}}\PYG{l+s}{pair:setflag}\PYG{l+s}{\PYGZdq{}}\PYG{p}{)}\PYG{p}{;}
\PYG{+w}{  }\PYG{k}{for}\PYG{+w}{ }\PYG{p}{(}\PYG{k+kt}{int}\PYG{+w}{ }\PYG{n}{i}\PYG{+w}{ }\PYG{o}{=}\PYG{+w}{ }\PYG{l+m+mi}{1}\PYG{p}{;}\PYG{+w}{ }\PYG{n}{i}\PYG{+w}{ }\PYG{o}{\PYGZlt{}}\PYG{+w}{ }\PYG{n}{np1}\PYG{p}{;}\PYG{+w}{ }\PYG{n}{i}\PYG{o}{+}\PYG{o}{+}\PYG{p}{)}
\PYG{+w}{    }\PYG{k}{for}\PYG{+w}{ }\PYG{p}{(}\PYG{k+kt}{int}\PYG{+w}{ }\PYG{n}{j}\PYG{+w}{ }\PYG{o}{=}\PYG{+w}{ }\PYG{n}{i}\PYG{p}{;}\PYG{+w}{ }\PYG{n}{j}\PYG{+w}{ }\PYG{o}{\PYGZlt{}}\PYG{+w}{ }\PYG{n}{np1}\PYG{p}{;}\PYG{+w}{ }\PYG{n}{j}\PYG{o}{+}\PYG{o}{+}\PYG{p}{)}\PYG{+w}{ }\PYG{n}{setflag}\PYG{p}{[}\PYG{n}{i}\PYG{p}{]}\PYG{p}{[}\PYG{n}{j}\PYG{p}{]}\PYG{+w}{ }\PYG{o}{=}\PYG{+w}{ }\PYG{l+m+mi}{0}\PYG{p}{;}

\PYG{+w}{  }\PYG{n}{memory}\PYG{o}{\PYGZhy{}}\PYG{o}{\PYGZgt{}}\PYG{n}{create}\PYG{p}{(}\PYG{n}{cutsq}\PYG{p}{,}\PYG{+w}{ }\PYG{n}{np1}\PYG{p}{,}\PYG{+w}{ }\PYG{n}{np1}\PYG{p}{,}\PYG{+w}{ }\PYG{l+s}{\PYGZdq{}}\PYG{l+s}{pair:cutsq}\PYG{l+s}{\PYGZdq{}}\PYG{p}{)}\PYG{p}{;}
\PYG{+w}{  }\PYG{n}{memory}\PYG{o}{\PYGZhy{}}\PYG{o}{\PYGZgt{}}\PYG{n}{create}\PYG{p}{(}\PYG{n}{cut}\PYG{p}{,}\PYG{+w}{ }\PYG{n}{np1}\PYG{p}{,}\PYG{+w}{ }\PYG{n}{np1}\PYG{p}{,}\PYG{+w}{ }\PYG{l+s}{\PYGZdq{}}\PYG{l+s}{pair:cut}\PYG{l+s}{\PYGZdq{}}\PYG{p}{)}\PYG{p}{;}
\PYG{+w}{  }\PYG{n}{memory}\PYG{o}{\PYGZhy{}}\PYG{o}{\PYGZgt{}}\PYG{n}{create}\PYG{p}{(}\PYG{n}{biga0}\PYG{p}{,}\PYG{+w}{ }\PYG{n}{np1}\PYG{p}{,}\PYG{+w}{ }\PYG{n}{np1}\PYG{p}{,}\PYG{+w}{ }\PYG{l+s}{\PYGZdq{}}\PYG{l+s}{pair:biga0}\PYG{l+s}{\PYGZdq{}}\PYG{p}{)}\PYG{p}{;}
\PYG{+w}{  }\PYG{n}{memory}\PYG{o}{\PYGZhy{}}\PYG{o}{\PYGZgt{}}\PYG{n}{create}\PYG{p}{(}\PYG{n}{alpha}\PYG{p}{,}\PYG{+w}{ }\PYG{n}{np1}\PYG{p}{,}\PYG{+w}{ }\PYG{n}{np1}\PYG{p}{,}\PYG{+w}{ }\PYG{l+s}{\PYGZdq{}}\PYG{l+s}{pair:alpha}\PYG{l+s}{\PYGZdq{}}\PYG{p}{)}\PYG{p}{;}
\PYG{+w}{  }\PYG{n}{memory}\PYG{o}{\PYGZhy{}}\PYG{o}{\PYGZgt{}}\PYG{n}{create}\PYG{p}{(}\PYG{n}{biga1}\PYG{p}{,}\PYG{+w}{ }\PYG{n}{np1}\PYG{p}{,}\PYG{+w}{ }\PYG{n}{np1}\PYG{p}{,}\PYG{+w}{ }\PYG{l+s}{\PYGZdq{}}\PYG{l+s}{pair:biga1}\PYG{l+s}{\PYGZdq{}}\PYG{p}{)}\PYG{p}{;}
\PYG{+w}{  }\PYG{n}{memory}\PYG{o}{\PYGZhy{}}\PYG{o}{\PYGZgt{}}\PYG{n}{create}\PYG{p}{(}\PYG{n}{beta}\PYG{p}{,}\PYG{+w}{ }\PYG{n}{np1}\PYG{p}{,}\PYG{+w}{ }\PYG{n}{np1}\PYG{p}{,}\PYG{+w}{ }\PYG{l+s}{\PYGZdq{}}\PYG{l+s}{pair:beta}\PYG{l+s}{\PYGZdq{}}\PYG{p}{)}\PYG{p}{;}
\PYG{+w}{  }\PYG{n}{memory}\PYG{o}{\PYGZhy{}}\PYG{o}{\PYGZgt{}}\PYG{n}{create}\PYG{p}{(}\PYG{n}{r0}\PYG{p}{,}\PYG{+w}{ }\PYG{n}{np1}\PYG{p}{,}\PYG{+w}{ }\PYG{n}{np1}\PYG{p}{,}\PYG{+w}{ }\PYG{l+s}{\PYGZdq{}}\PYG{l+s}{pair:r0}\PYG{l+s}{\PYGZdq{}}\PYG{p}{)}\PYG{p}{;}
\PYG{+w}{  }\PYG{n}{memory}\PYG{o}{\PYGZhy{}}\PYG{o}{\PYGZgt{}}\PYG{n}{create}\PYG{p}{(}\PYG{n}{offset}\PYG{p}{,}\PYG{+w}{ }\PYG{n}{np1}\PYG{p}{,}\PYG{+w}{ }\PYG{n}{np1}\PYG{p}{,}\PYG{+w}{ }\PYG{l+s}{\PYGZdq{}}\PYG{l+s}{pair:offset}\PYG{l+s}{\PYGZdq{}}\PYG{p}{)}\PYG{p}{;}
\PYG{p}{\PYGZcb{}}

\PYG{c+cm}{/* \PYGZhy{}\PYGZhy{}\PYGZhy{}\PYGZhy{}\PYGZhy{}\PYGZhy{}\PYGZhy{}\PYGZhy{}\PYGZhy{}\PYGZhy{}\PYGZhy{}\PYGZhy{}\PYGZhy{}\PYGZhy{}\PYGZhy{}\PYGZhy{}\PYGZhy{}\PYGZhy{}\PYGZhy{}\PYGZhy{}\PYGZhy{}\PYGZhy{}\PYGZhy{}\PYGZhy{}\PYGZhy{}\PYGZhy{}\PYGZhy{}\PYGZhy{}\PYGZhy{}\PYGZhy{}\PYGZhy{}\PYGZhy{}\PYGZhy{}\PYGZhy{}\PYGZhy{}\PYGZhy{}\PYGZhy{}\PYGZhy{}\PYGZhy{}\PYGZhy{}\PYGZhy{}\PYGZhy{}\PYGZhy{}\PYGZhy{}\PYGZhy{}\PYGZhy{}\PYGZhy{}\PYGZhy{}\PYGZhy{}\PYGZhy{}\PYGZhy{}\PYGZhy{}\PYGZhy{}\PYGZhy{}\PYGZhy{}\PYGZhy{}\PYGZhy{}\PYGZhy{}\PYGZhy{}\PYGZhy{}\PYGZhy{}\PYGZhy{}\PYGZhy{}\PYGZhy{}\PYGZhy{}\PYGZhy{}\PYGZhy{}\PYGZhy{}\PYGZhy{}\PYGZhy{}}
\PYG{c+cm}{   global settings}
\PYG{c+cm}{\PYGZhy{}\PYGZhy{}\PYGZhy{}\PYGZhy{}\PYGZhy{}\PYGZhy{}\PYGZhy{}\PYGZhy{}\PYGZhy{}\PYGZhy{}\PYGZhy{}\PYGZhy{}\PYGZhy{}\PYGZhy{}\PYGZhy{}\PYGZhy{}\PYGZhy{}\PYGZhy{}\PYGZhy{}\PYGZhy{}\PYGZhy{}\PYGZhy{}\PYGZhy{}\PYGZhy{}\PYGZhy{}\PYGZhy{}\PYGZhy{}\PYGZhy{}\PYGZhy{}\PYGZhy{}\PYGZhy{}\PYGZhy{}\PYGZhy{}\PYGZhy{}\PYGZhy{}\PYGZhy{}\PYGZhy{}\PYGZhy{}\PYGZhy{}\PYGZhy{}\PYGZhy{}\PYGZhy{}\PYGZhy{}\PYGZhy{}\PYGZhy{}\PYGZhy{}\PYGZhy{}\PYGZhy{}\PYGZhy{}\PYGZhy{}\PYGZhy{}\PYGZhy{}\PYGZhy{}\PYGZhy{}\PYGZhy{}\PYGZhy{}\PYGZhy{}\PYGZhy{}\PYGZhy{}\PYGZhy{}\PYGZhy{}\PYGZhy{}\PYGZhy{}\PYGZhy{}\PYGZhy{}\PYGZhy{}\PYGZhy{}\PYGZhy{}\PYGZhy{}\PYGZhy{}\PYGZhy{}\PYGZhy{}\PYGZhy{} */}

\PYG{k+kt}{void}\PYG{+w}{ }\PYG{n+nf}{PairBornGauss::settings}\PYG{p}{(}\PYG{k+kt}{int}\PYG{+w}{ }\PYG{n}{narg}\PYG{p}{,}\PYG{+w}{ }\PYG{k+kt}{char}\PYG{+w}{ }\PYG{o}{*}\PYG{o}{*}\PYG{n}{arg}\PYG{p}{)}
\PYG{p}{\PYGZob{}}
\PYG{+w}{  }\PYG{k}{if}\PYG{+w}{ }\PYG{p}{(}\PYG{n}{narg}\PYG{+w}{ }\PYG{o}{!}\PYG{o}{=}\PYG{+w}{ }\PYG{l+m+mi}{1}\PYG{p}{)}\PYG{+w}{ }\PYG{n}{error}\PYG{o}{\PYGZhy{}}\PYG{o}{\PYGZgt{}}\PYG{n}{all}\PYG{p}{(}\PYG{n}{FLERR}\PYG{p}{,}\PYG{+w}{ }\PYG{l+s}{\PYGZdq{}}\PYG{l+s}{Pair style bond/gauss must have exactly one argument}\PYG{l+s}{\PYGZdq{}}\PYG{p}{)}\PYG{p}{;}
\PYG{+w}{  }\PYG{n}{cut\PYGZus{}global}\PYG{+w}{ }\PYG{o}{=}\PYG{+w}{ }\PYG{n}{utils}\PYG{o}{:}\PYG{o}{:}\PYG{n}{numeric}\PYG{p}{(}\PYG{n}{FLERR}\PYG{p}{,}\PYG{+w}{ }\PYG{n}{arg}\PYG{p}{[}\PYG{l+m+mi}{0}\PYG{p}{]}\PYG{p}{,}\PYG{+w}{ }\PYG{n+nb}{false}\PYG{p}{,}\PYG{+w}{ }\PYG{n}{lmp}\PYG{p}{)}\PYG{p}{;}

\PYG{+w}{  }\PYG{c+c1}{// reset per\PYGZhy{}type pair cutoffs that have been explicitly set previously}

\PYG{+w}{  }\PYG{k}{if}\PYG{+w}{ }\PYG{p}{(}\PYG{n}{allocated}\PYG{p}{)}\PYG{+w}{ }\PYG{p}{\PYGZob{}}
\PYG{+w}{    }\PYG{k}{for}\PYG{+w}{ }\PYG{p}{(}\PYG{k+kt}{int}\PYG{+w}{ }\PYG{n}{i}\PYG{+w}{ }\PYG{o}{=}\PYG{+w}{ }\PYG{l+m+mi}{1}\PYG{p}{;}\PYG{+w}{ }\PYG{n}{i}\PYG{+w}{ }\PYG{o}{\PYGZlt{}}\PYG{o}{=}\PYG{+w}{ }\PYG{n}{atom}\PYG{o}{\PYGZhy{}}\PYG{o}{\PYGZgt{}}\PYG{n}{ntypes}\PYG{p}{;}\PYG{+w}{ }\PYG{n}{i}\PYG{o}{+}\PYG{o}{+}\PYG{p}{)}
\PYG{+w}{      }\PYG{k}{for}\PYG{+w}{ }\PYG{p}{(}\PYG{k+kt}{int}\PYG{+w}{ }\PYG{n}{j}\PYG{+w}{ }\PYG{o}{=}\PYG{+w}{ }\PYG{n}{i}\PYG{p}{;}\PYG{+w}{ }\PYG{n}{j}\PYG{+w}{ }\PYG{o}{\PYGZlt{}}\PYG{o}{=}\PYG{+w}{ }\PYG{n}{atom}\PYG{o}{\PYGZhy{}}\PYG{o}{\PYGZgt{}}\PYG{n}{ntypes}\PYG{p}{;}\PYG{+w}{ }\PYG{n}{j}\PYG{o}{+}\PYG{o}{+}\PYG{p}{)}
\PYG{+w}{        }\PYG{k}{if}\PYG{+w}{ }\PYG{p}{(}\PYG{n}{setflag}\PYG{p}{[}\PYG{n}{i}\PYG{p}{]}\PYG{p}{[}\PYG{n}{j}\PYG{p}{]}\PYG{p}{)}\PYG{+w}{ }\PYG{n}{cut}\PYG{p}{[}\PYG{n}{i}\PYG{p}{]}\PYG{p}{[}\PYG{n}{j}\PYG{p}{]}\PYG{+w}{ }\PYG{o}{=}\PYG{+w}{ }\PYG{n}{cut\PYGZus{}global}\PYG{p}{;}
\PYG{+w}{  }\PYG{p}{\PYGZcb{}}
\PYG{p}{\PYGZcb{}}
\end{sphinxVerbatim}

\sphinxAtStartPar
The arguments to the \sphinxcode{\sphinxupquote{coeff()}} function are the arguments to the
{\hyperref[\detokenize{pair_coeff::doc}]{\sphinxcrossref{\DUrole{doc}{pair\_coeff command}}}}.  The function is also called
when processing the \sphinxcode{\sphinxupquote{Pair Coeffs}} or \sphinxcode{\sphinxupquote{PairIJ Coeffs}} sections of
data files.  In the case of the \sphinxcode{\sphinxupquote{Pair Coeffs}} section, there is only
one atom type per line and thus the first argument is duplicated.  Since
the atom type arguments of the {\hyperref[\detokenize{pair_coeff::doc}]{\sphinxcrossref{\DUrole{doc}{pair\_coeff command}}}}
may be a range (e.g. *3 for atom types 1, 2, and 3), the
corresponding arguments are passed to the {\hyperref[\detokenize{Developer_utils:_CPPv4I0EN9LAMMPS_NS5utils6boundsEvPKciRKNSt6stringE6bigint6bigintR4TYPER4TYPEP5Error}]{\sphinxcrossref{\sphinxcode{\sphinxupquote{utils::bounds()}}}}} function which will then return the low
and high end of the range.  Note that the \sphinxcode{\sphinxupquote{setflag}} array is set to 1
for all pairs of atom types processed by this call.  This information is
later used in the \sphinxcode{\sphinxupquote{init\_one()}} function to determine if any coefficients
are missing and, if supported by the potential, generate those missing
coefficients from the selected mixing rule.

\begin{sphinxVerbatim}[commandchars=\\\{\}]
\PYG{c+cm}{/* \PYGZhy{}\PYGZhy{}\PYGZhy{}\PYGZhy{}\PYGZhy{}\PYGZhy{}\PYGZhy{}\PYGZhy{}\PYGZhy{}\PYGZhy{}\PYGZhy{}\PYGZhy{}\PYGZhy{}\PYGZhy{}\PYGZhy{}\PYGZhy{}\PYGZhy{}\PYGZhy{}\PYGZhy{}\PYGZhy{}\PYGZhy{}\PYGZhy{}\PYGZhy{}\PYGZhy{}\PYGZhy{}\PYGZhy{}\PYGZhy{}\PYGZhy{}\PYGZhy{}\PYGZhy{}\PYGZhy{}\PYGZhy{}\PYGZhy{}\PYGZhy{}\PYGZhy{}\PYGZhy{}\PYGZhy{}\PYGZhy{}\PYGZhy{}\PYGZhy{}\PYGZhy{}\PYGZhy{}\PYGZhy{}\PYGZhy{}\PYGZhy{}\PYGZhy{}\PYGZhy{}\PYGZhy{}\PYGZhy{}\PYGZhy{}\PYGZhy{}\PYGZhy{}\PYGZhy{}\PYGZhy{}\PYGZhy{}\PYGZhy{}\PYGZhy{}\PYGZhy{}\PYGZhy{}\PYGZhy{}\PYGZhy{}\PYGZhy{}\PYGZhy{}\PYGZhy{}\PYGZhy{}\PYGZhy{}\PYGZhy{}\PYGZhy{}\PYGZhy{}\PYGZhy{}}
\PYG{c+cm}{   set coeffs for one or more type pairs}
\PYG{c+cm}{\PYGZhy{}\PYGZhy{}\PYGZhy{}\PYGZhy{}\PYGZhy{}\PYGZhy{}\PYGZhy{}\PYGZhy{}\PYGZhy{}\PYGZhy{}\PYGZhy{}\PYGZhy{}\PYGZhy{}\PYGZhy{}\PYGZhy{}\PYGZhy{}\PYGZhy{}\PYGZhy{}\PYGZhy{}\PYGZhy{}\PYGZhy{}\PYGZhy{}\PYGZhy{}\PYGZhy{}\PYGZhy{}\PYGZhy{}\PYGZhy{}\PYGZhy{}\PYGZhy{}\PYGZhy{}\PYGZhy{}\PYGZhy{}\PYGZhy{}\PYGZhy{}\PYGZhy{}\PYGZhy{}\PYGZhy{}\PYGZhy{}\PYGZhy{}\PYGZhy{}\PYGZhy{}\PYGZhy{}\PYGZhy{}\PYGZhy{}\PYGZhy{}\PYGZhy{}\PYGZhy{}\PYGZhy{}\PYGZhy{}\PYGZhy{}\PYGZhy{}\PYGZhy{}\PYGZhy{}\PYGZhy{}\PYGZhy{}\PYGZhy{}\PYGZhy{}\PYGZhy{}\PYGZhy{}\PYGZhy{}\PYGZhy{}\PYGZhy{}\PYGZhy{}\PYGZhy{}\PYGZhy{}\PYGZhy{}\PYGZhy{}\PYGZhy{}\PYGZhy{}\PYGZhy{}\PYGZhy{}\PYGZhy{}\PYGZhy{} */}

\PYG{k+kt}{void}\PYG{+w}{ }\PYG{n+nf}{PairBornGauss::coeff}\PYG{p}{(}\PYG{k+kt}{int}\PYG{+w}{ }\PYG{n}{narg}\PYG{p}{,}\PYG{+w}{ }\PYG{k+kt}{char}\PYG{+w}{ }\PYG{o}{*}\PYG{o}{*}\PYG{n}{arg}\PYG{p}{)}
\PYG{p}{\PYGZob{}}
\PYG{+w}{  }\PYG{k}{if}\PYG{+w}{ }\PYG{p}{(}\PYG{n}{narg}\PYG{+w}{ }\PYG{o}{\PYGZlt{}}\PYG{+w}{ }\PYG{l+m+mi}{7}\PYG{+w}{ }\PYG{o}{|}\PYG{o}{|}\PYG{+w}{ }\PYG{n}{narg}\PYG{+w}{ }\PYG{o}{\PYGZgt{}}\PYG{+w}{ }\PYG{l+m+mi}{8}\PYG{p}{)}\PYG{+w}{ }\PYG{n}{error}\PYG{o}{\PYGZhy{}}\PYG{o}{\PYGZgt{}}\PYG{n}{all}\PYG{p}{(}\PYG{n}{FLERR}\PYG{p}{,}\PYG{+w}{ }\PYG{l+s}{\PYGZdq{}}\PYG{l+s}{Incorrect args for pair coefficients}\PYG{l+s}{\PYGZdq{}}\PYG{p}{)}\PYG{p}{;}
\PYG{+w}{  }\PYG{k}{if}\PYG{+w}{ }\PYG{p}{(}\PYG{o}{!}\PYG{n}{allocated}\PYG{p}{)}\PYG{+w}{ }\PYG{n}{allocate}\PYG{p}{(}\PYG{p}{)}\PYG{p}{;}

\PYG{+w}{  }\PYG{k+kt}{int}\PYG{+w}{ }\PYG{n}{ilo}\PYG{p}{,}\PYG{+w}{ }\PYG{n}{ihi}\PYG{p}{,}\PYG{+w}{ }\PYG{n}{jlo}\PYG{p}{,}\PYG{+w}{ }\PYG{n}{jhi}\PYG{p}{;}
\PYG{+w}{  }\PYG{n}{utils}\PYG{o}{:}\PYG{o}{:}\PYG{n}{bounds}\PYG{p}{(}\PYG{n}{FLERR}\PYG{p}{,}\PYG{+w}{ }\PYG{n}{arg}\PYG{p}{[}\PYG{l+m+mi}{0}\PYG{p}{]}\PYG{p}{,}\PYG{+w}{ }\PYG{l+m+mi}{1}\PYG{p}{,}\PYG{+w}{ }\PYG{n}{atom}\PYG{o}{\PYGZhy{}}\PYG{o}{\PYGZgt{}}\PYG{n}{ntypes}\PYG{p}{,}\PYG{+w}{ }\PYG{n}{ilo}\PYG{p}{,}\PYG{+w}{ }\PYG{n}{ihi}\PYG{p}{,}\PYG{+w}{ }\PYG{n}{error}\PYG{p}{)}\PYG{p}{;}
\PYG{+w}{  }\PYG{n}{utils}\PYG{o}{:}\PYG{o}{:}\PYG{n}{bounds}\PYG{p}{(}\PYG{n}{FLERR}\PYG{p}{,}\PYG{+w}{ }\PYG{n}{arg}\PYG{p}{[}\PYG{l+m+mi}{1}\PYG{p}{]}\PYG{p}{,}\PYG{+w}{ }\PYG{l+m+mi}{1}\PYG{p}{,}\PYG{+w}{ }\PYG{n}{atom}\PYG{o}{\PYGZhy{}}\PYG{o}{\PYGZgt{}}\PYG{n}{ntypes}\PYG{p}{,}\PYG{+w}{ }\PYG{n}{jlo}\PYG{p}{,}\PYG{+w}{ }\PYG{n}{jhi}\PYG{p}{,}\PYG{+w}{ }\PYG{n}{error}\PYG{p}{)}\PYG{p}{;}

\PYG{+w}{  }\PYG{k+kt}{double}\PYG{+w}{ }\PYG{n}{biga0\PYGZus{}one}\PYG{+w}{ }\PYG{o}{=}\PYG{+w}{ }\PYG{n}{utils}\PYG{o}{:}\PYG{o}{:}\PYG{n}{numeric}\PYG{p}{(}\PYG{n}{FLERR}\PYG{p}{,}\PYG{+w}{ }\PYG{n}{arg}\PYG{p}{[}\PYG{l+m+mi}{2}\PYG{p}{]}\PYG{p}{,}\PYG{+w}{ }\PYG{n+nb}{false}\PYG{p}{,}\PYG{+w}{ }\PYG{n}{lmp}\PYG{p}{)}\PYG{p}{;}
\PYG{+w}{  }\PYG{k+kt}{double}\PYG{+w}{ }\PYG{n}{alpha\PYGZus{}one}\PYG{+w}{ }\PYG{o}{=}\PYG{+w}{ }\PYG{n}{utils}\PYG{o}{:}\PYG{o}{:}\PYG{n}{numeric}\PYG{p}{(}\PYG{n}{FLERR}\PYG{p}{,}\PYG{+w}{ }\PYG{n}{arg}\PYG{p}{[}\PYG{l+m+mi}{3}\PYG{p}{]}\PYG{p}{,}\PYG{+w}{ }\PYG{n+nb}{false}\PYG{p}{,}\PYG{+w}{ }\PYG{n}{lmp}\PYG{p}{)}\PYG{p}{;}
\PYG{+w}{  }\PYG{k+kt}{double}\PYG{+w}{ }\PYG{n}{biga1\PYGZus{}one}\PYG{+w}{ }\PYG{o}{=}\PYG{+w}{ }\PYG{n}{utils}\PYG{o}{:}\PYG{o}{:}\PYG{n}{numeric}\PYG{p}{(}\PYG{n}{FLERR}\PYG{p}{,}\PYG{+w}{ }\PYG{n}{arg}\PYG{p}{[}\PYG{l+m+mi}{4}\PYG{p}{]}\PYG{p}{,}\PYG{+w}{ }\PYG{n+nb}{false}\PYG{p}{,}\PYG{+w}{ }\PYG{n}{lmp}\PYG{p}{)}\PYG{p}{;}
\PYG{+w}{  }\PYG{k+kt}{double}\PYG{+w}{ }\PYG{n}{beta\PYGZus{}one}\PYG{+w}{ }\PYG{o}{=}\PYG{+w}{ }\PYG{n}{utils}\PYG{o}{:}\PYG{o}{:}\PYG{n}{numeric}\PYG{p}{(}\PYG{n}{FLERR}\PYG{p}{,}\PYG{+w}{ }\PYG{n}{arg}\PYG{p}{[}\PYG{l+m+mi}{5}\PYG{p}{]}\PYG{p}{,}\PYG{+w}{ }\PYG{n+nb}{false}\PYG{p}{,}\PYG{+w}{ }\PYG{n}{lmp}\PYG{p}{)}\PYG{p}{;}
\PYG{+w}{  }\PYG{k+kt}{double}\PYG{+w}{ }\PYG{n}{r0\PYGZus{}one}\PYG{+w}{ }\PYG{o}{=}\PYG{+w}{ }\PYG{n}{utils}\PYG{o}{:}\PYG{o}{:}\PYG{n}{numeric}\PYG{p}{(}\PYG{n}{FLERR}\PYG{p}{,}\PYG{+w}{ }\PYG{n}{arg}\PYG{p}{[}\PYG{l+m+mi}{6}\PYG{p}{]}\PYG{p}{,}\PYG{+w}{ }\PYG{n+nb}{false}\PYG{p}{,}\PYG{+w}{ }\PYG{n}{lmp}\PYG{p}{)}\PYG{p}{;}
\PYG{+w}{  }\PYG{k+kt}{double}\PYG{+w}{ }\PYG{n}{cut\PYGZus{}one}\PYG{+w}{ }\PYG{o}{=}\PYG{+w}{ }\PYG{n}{cut\PYGZus{}global}\PYG{p}{;}
\PYG{+w}{  }\PYG{k}{if}\PYG{+w}{ }\PYG{p}{(}\PYG{n}{narg}\PYG{+w}{ }\PYG{o}{=}\PYG{o}{=}\PYG{+w}{ }\PYG{l+m+mi}{10}\PYG{p}{)}\PYG{+w}{ }\PYG{n}{cut\PYGZus{}one}\PYG{+w}{ }\PYG{o}{=}\PYG{+w}{ }\PYG{n}{utils}\PYG{o}{:}\PYG{o}{:}\PYG{n}{numeric}\PYG{p}{(}\PYG{n}{FLERR}\PYG{p}{,}\PYG{+w}{ }\PYG{n}{arg}\PYG{p}{[}\PYG{l+m+mi}{7}\PYG{p}{]}\PYG{p}{,}\PYG{+w}{ }\PYG{n+nb}{false}\PYG{p}{,}\PYG{+w}{ }\PYG{n}{lmp}\PYG{p}{)}\PYG{p}{;}

\PYG{+w}{  }\PYG{k+kt}{int}\PYG{+w}{ }\PYG{n}{count}\PYG{+w}{ }\PYG{o}{=}\PYG{+w}{ }\PYG{l+m+mi}{0}\PYG{p}{;}
\PYG{+w}{  }\PYG{k}{for}\PYG{+w}{ }\PYG{p}{(}\PYG{k+kt}{int}\PYG{+w}{ }\PYG{n}{i}\PYG{+w}{ }\PYG{o}{=}\PYG{+w}{ }\PYG{n}{ilo}\PYG{p}{;}\PYG{+w}{ }\PYG{n}{i}\PYG{+w}{ }\PYG{o}{\PYGZlt{}}\PYG{o}{=}\PYG{+w}{ }\PYG{n}{ihi}\PYG{p}{;}\PYG{+w}{ }\PYG{n}{i}\PYG{o}{+}\PYG{o}{+}\PYG{p}{)}\PYG{+w}{ }\PYG{p}{\PYGZob{}}
\PYG{+w}{    }\PYG{k}{for}\PYG{+w}{ }\PYG{p}{(}\PYG{k+kt}{int}\PYG{+w}{ }\PYG{n}{j}\PYG{+w}{ }\PYG{o}{=}\PYG{+w}{ }\PYG{n}{MAX}\PYG{p}{(}\PYG{n}{jlo}\PYG{p}{,}\PYG{+w}{ }\PYG{n}{i}\PYG{p}{)}\PYG{p}{;}\PYG{+w}{ }\PYG{n}{j}\PYG{+w}{ }\PYG{o}{\PYGZlt{}}\PYG{o}{=}\PYG{+w}{ }\PYG{n}{jhi}\PYG{p}{;}\PYG{+w}{ }\PYG{n}{j}\PYG{o}{+}\PYG{o}{+}\PYG{p}{)}\PYG{+w}{ }\PYG{p}{\PYGZob{}}
\PYG{+w}{      }\PYG{n}{biga0}\PYG{p}{[}\PYG{n}{i}\PYG{p}{]}\PYG{p}{[}\PYG{n}{j}\PYG{p}{]}\PYG{+w}{ }\PYG{o}{=}\PYG{+w}{ }\PYG{n}{biga0\PYGZus{}one}\PYG{p}{;}
\PYG{+w}{      }\PYG{n}{alpha}\PYG{p}{[}\PYG{n}{i}\PYG{p}{]}\PYG{p}{[}\PYG{n}{j}\PYG{p}{]}\PYG{+w}{ }\PYG{o}{=}\PYG{+w}{ }\PYG{n}{alpha\PYGZus{}one}\PYG{p}{;}
\PYG{+w}{      }\PYG{n}{biga1}\PYG{p}{[}\PYG{n}{i}\PYG{p}{]}\PYG{p}{[}\PYG{n}{j}\PYG{p}{]}\PYG{+w}{ }\PYG{o}{=}\PYG{+w}{ }\PYG{n}{biga1\PYGZus{}one}\PYG{p}{;}
\PYG{+w}{      }\PYG{n}{beta}\PYG{p}{[}\PYG{n}{i}\PYG{p}{]}\PYG{p}{[}\PYG{n}{j}\PYG{p}{]}\PYG{+w}{ }\PYG{o}{=}\PYG{+w}{ }\PYG{n}{beta\PYGZus{}one}\PYG{p}{;}
\PYG{+w}{      }\PYG{n}{r0}\PYG{p}{[}\PYG{n}{i}\PYG{p}{]}\PYG{p}{[}\PYG{n}{j}\PYG{p}{]}\PYG{+w}{ }\PYG{o}{=}\PYG{+w}{ }\PYG{n}{r0\PYGZus{}one}\PYG{p}{;}
\PYG{+w}{      }\PYG{n}{cut}\PYG{p}{[}\PYG{n}{i}\PYG{p}{]}\PYG{p}{[}\PYG{n}{j}\PYG{p}{]}\PYG{+w}{ }\PYG{o}{=}\PYG{+w}{ }\PYG{n}{cut\PYGZus{}one}\PYG{p}{;}
\PYG{+w}{      }\PYG{n}{setflag}\PYG{p}{[}\PYG{n}{i}\PYG{p}{]}\PYG{p}{[}\PYG{n}{j}\PYG{p}{]}\PYG{+w}{ }\PYG{o}{=}\PYG{+w}{ }\PYG{l+m+mi}{1}\PYG{p}{;}
\PYG{+w}{      }\PYG{n}{count}\PYG{o}{+}\PYG{o}{+}\PYG{p}{;}
\PYG{+w}{    }\PYG{p}{\PYGZcb{}}
\PYG{+w}{  }\PYG{p}{\PYGZcb{}}

\PYG{+w}{  }\PYG{k}{if}\PYG{+w}{ }\PYG{p}{(}\PYG{n}{count}\PYG{+w}{ }\PYG{o}{=}\PYG{o}{=}\PYG{+w}{ }\PYG{l+m+mi}{0}\PYG{p}{)}\PYG{+w}{ }\PYG{n}{error}\PYG{o}{\PYGZhy{}}\PYG{o}{\PYGZgt{}}\PYG{n}{all}\PYG{p}{(}\PYG{n}{FLERR}\PYG{p}{,}\PYG{+w}{ }\PYG{l+s}{\PYGZdq{}}\PYG{l+s}{Incorrect args for pair coefficients}\PYG{l+s}{\PYGZdq{}}\PYG{p}{)}\PYG{p}{;}
\PYG{p}{\PYGZcb{}}
\end{sphinxVerbatim}


\subsubsection{Initialization}
\label{\detokenize{Developer_write_pair:initialization}}
\sphinxAtStartPar
The \sphinxcode{\sphinxupquote{init\_one()}} function is called during the {\hyperref[\detokenize{Developer_flow::doc}]{\sphinxcrossref{\DUrole{doc}{“init” phase}}}} of a simulation.  This is where potential parameters
are checked for completeness, derived parameters computed (e.g. the
“offset” of the potential energy at the cutoff distance for use with the
{\hyperref[\detokenize{pair_modify::doc}]{\sphinxcrossref{\DUrole{doc}{pair\_modify shift yes}}}} command).  If a pair style
supports generating “mixed” parameters (i.e. where both atoms of a pair
have a different atom type) using a “mixing rule” from the parameters of
the type with itself, this is the place to compute and store those mixed
values.  The \sphinxstyleemphasis{born/gauss} pair style does not support mixing, so we only
check for completeness.  Another purpose of the \sphinxcode{\sphinxupquote{init\_one()}} function
is to symmetrize the potential parameter arrays.  The return value of
the function is the cutoff for the given pair of atom types.  This
information is used by the neighbor list code to determine the largest
cutoff and then build the neighbor lists accordingly.

\begin{sphinxVerbatim}[commandchars=\\\{\}]
\PYG{c+cm}{/* \PYGZhy{}\PYGZhy{}\PYGZhy{}\PYGZhy{}\PYGZhy{}\PYGZhy{}\PYGZhy{}\PYGZhy{}\PYGZhy{}\PYGZhy{}\PYGZhy{}\PYGZhy{}\PYGZhy{}\PYGZhy{}\PYGZhy{}\PYGZhy{}\PYGZhy{}\PYGZhy{}\PYGZhy{}\PYGZhy{}\PYGZhy{}\PYGZhy{}\PYGZhy{}\PYGZhy{}\PYGZhy{}\PYGZhy{}\PYGZhy{}\PYGZhy{}\PYGZhy{}\PYGZhy{}\PYGZhy{}\PYGZhy{}\PYGZhy{}\PYGZhy{}\PYGZhy{}\PYGZhy{}\PYGZhy{}\PYGZhy{}\PYGZhy{}\PYGZhy{}\PYGZhy{}\PYGZhy{}\PYGZhy{}\PYGZhy{}\PYGZhy{}\PYGZhy{}\PYGZhy{}\PYGZhy{}\PYGZhy{}\PYGZhy{}\PYGZhy{}\PYGZhy{}\PYGZhy{}\PYGZhy{}\PYGZhy{}\PYGZhy{}\PYGZhy{}\PYGZhy{}\PYGZhy{}\PYGZhy{}\PYGZhy{}\PYGZhy{}\PYGZhy{}\PYGZhy{}\PYGZhy{}\PYGZhy{}\PYGZhy{}\PYGZhy{}\PYGZhy{}\PYGZhy{}}
\PYG{c+cm}{   init for one type pair i,j and corresponding j,i}
\PYG{c+cm}{\PYGZhy{}\PYGZhy{}\PYGZhy{}\PYGZhy{}\PYGZhy{}\PYGZhy{}\PYGZhy{}\PYGZhy{}\PYGZhy{}\PYGZhy{}\PYGZhy{}\PYGZhy{}\PYGZhy{}\PYGZhy{}\PYGZhy{}\PYGZhy{}\PYGZhy{}\PYGZhy{}\PYGZhy{}\PYGZhy{}\PYGZhy{}\PYGZhy{}\PYGZhy{}\PYGZhy{}\PYGZhy{}\PYGZhy{}\PYGZhy{}\PYGZhy{}\PYGZhy{}\PYGZhy{}\PYGZhy{}\PYGZhy{}\PYGZhy{}\PYGZhy{}\PYGZhy{}\PYGZhy{}\PYGZhy{}\PYGZhy{}\PYGZhy{}\PYGZhy{}\PYGZhy{}\PYGZhy{}\PYGZhy{}\PYGZhy{}\PYGZhy{}\PYGZhy{}\PYGZhy{}\PYGZhy{}\PYGZhy{}\PYGZhy{}\PYGZhy{}\PYGZhy{}\PYGZhy{}\PYGZhy{}\PYGZhy{}\PYGZhy{}\PYGZhy{}\PYGZhy{}\PYGZhy{}\PYGZhy{}\PYGZhy{}\PYGZhy{}\PYGZhy{}\PYGZhy{}\PYGZhy{}\PYGZhy{}\PYGZhy{}\PYGZhy{}\PYGZhy{}\PYGZhy{}\PYGZhy{}\PYGZhy{}\PYGZhy{} */}

\PYG{k+kt}{double}\PYG{+w}{ }\PYG{n+nf}{PairBornGauss::init\PYGZus{}one}\PYG{p}{(}\PYG{k+kt}{int}\PYG{+w}{ }\PYG{n}{i}\PYG{p}{,}\PYG{+w}{ }\PYG{k+kt}{int}\PYG{+w}{ }\PYG{n}{j}\PYG{p}{)}
\PYG{p}{\PYGZob{}}
\PYG{+w}{  }\PYG{k}{if}\PYG{+w}{ }\PYG{p}{(}\PYG{n}{setflag}\PYG{p}{[}\PYG{n}{i}\PYG{p}{]}\PYG{p}{[}\PYG{n}{j}\PYG{p}{]}\PYG{+w}{ }\PYG{o}{=}\PYG{o}{=}\PYG{+w}{ }\PYG{l+m+mi}{0}\PYG{p}{)}\PYG{+w}{ }\PYG{n}{error}\PYG{o}{\PYGZhy{}}\PYG{o}{\PYGZgt{}}\PYG{n}{all}\PYG{p}{(}\PYG{n}{FLERR}\PYG{p}{,}\PYG{+w}{ }\PYG{l+s}{\PYGZdq{}}\PYG{l+s}{All pair coeffs are not set}\PYG{l+s}{\PYGZdq{}}\PYG{p}{)}\PYG{p}{;}

\PYG{+w}{  }\PYG{k}{if}\PYG{+w}{ }\PYG{p}{(}\PYG{n}{offset\PYGZus{}flag}\PYG{p}{)}\PYG{+w}{ }\PYG{p}{\PYGZob{}}
\PYG{+w}{    }\PYG{k+kt}{double}\PYG{+w}{ }\PYG{n}{dr}\PYG{+w}{ }\PYG{o}{=}\PYG{+w}{ }\PYG{n}{cut}\PYG{p}{[}\PYG{n}{i}\PYG{p}{]}\PYG{p}{[}\PYG{n}{j}\PYG{p}{]}\PYG{+w}{ }\PYG{o}{\PYGZhy{}}\PYG{+w}{ }\PYG{n}{r0}\PYG{p}{[}\PYG{n}{i}\PYG{p}{]}\PYG{p}{[}\PYG{n}{j}\PYG{p}{]}\PYG{p}{;}
\PYG{+w}{    }\PYG{n}{offset}\PYG{p}{[}\PYG{n}{i}\PYG{p}{]}\PYG{p}{[}\PYG{n}{j}\PYG{p}{]}\PYG{+w}{ }\PYG{o}{=}
\PYG{+w}{        }\PYG{n}{biga0}\PYG{p}{[}\PYG{n}{i}\PYG{p}{]}\PYG{p}{[}\PYG{n}{j}\PYG{p}{]}\PYG{+w}{ }\PYG{o}{*}\PYG{+w}{ }\PYG{n}{exp}\PYG{p}{(}\PYG{o}{\PYGZhy{}}\PYG{n}{alpha}\PYG{p}{[}\PYG{n}{i}\PYG{p}{]}\PYG{p}{[}\PYG{n}{j}\PYG{p}{]}\PYG{+w}{ }\PYG{o}{*}\PYG{+w}{ }\PYG{n}{cut}\PYG{p}{[}\PYG{n}{i}\PYG{p}{]}\PYG{p}{[}\PYG{n}{j}\PYG{p}{]}\PYG{p}{)}\PYG{+w}{ }\PYG{o}{\PYGZhy{}}\PYG{+w}{ }\PYG{n}{biga1}\PYG{p}{[}\PYG{n}{i}\PYG{p}{]}\PYG{p}{[}\PYG{n}{j}\PYG{p}{]}\PYG{+w}{ }\PYG{o}{*}\PYG{+w}{ }\PYG{n}{exp}\PYG{p}{(}\PYG{o}{\PYGZhy{}}\PYG{n}{beta}\PYG{p}{[}\PYG{n}{i}\PYG{p}{]}\PYG{p}{[}\PYG{n}{j}\PYG{p}{]}\PYG{+w}{ }\PYG{o}{*}\PYG{+w}{ }\PYG{n}{dr}\PYG{+w}{ }\PYG{o}{*}\PYG{+w}{ }\PYG{n}{dr}\PYG{p}{)}\PYG{p}{;}
\PYG{+w}{  }\PYG{p}{\PYGZcb{}}\PYG{+w}{ }\PYG{k}{else}
\PYG{+w}{    }\PYG{n}{offset}\PYG{p}{[}\PYG{n}{i}\PYG{p}{]}\PYG{p}{[}\PYG{n}{j}\PYG{p}{]}\PYG{+w}{ }\PYG{o}{=}\PYG{+w}{ }\PYG{l+m+mf}{0.0}\PYG{p}{;}

\PYG{+w}{  }\PYG{n}{biga0}\PYG{p}{[}\PYG{n}{j}\PYG{p}{]}\PYG{p}{[}\PYG{n}{i}\PYG{p}{]}\PYG{+w}{ }\PYG{o}{=}\PYG{+w}{ }\PYG{n}{biga0}\PYG{p}{[}\PYG{n}{i}\PYG{p}{]}\PYG{p}{[}\PYG{n}{j}\PYG{p}{]}\PYG{p}{;}
\PYG{+w}{  }\PYG{n}{alpha}\PYG{p}{[}\PYG{n}{j}\PYG{p}{]}\PYG{p}{[}\PYG{n}{i}\PYG{p}{]}\PYG{+w}{ }\PYG{o}{=}\PYG{+w}{ }\PYG{n}{alpha}\PYG{p}{[}\PYG{n}{i}\PYG{p}{]}\PYG{p}{[}\PYG{n}{j}\PYG{p}{]}\PYG{p}{;}
\PYG{+w}{  }\PYG{n}{biga1}\PYG{p}{[}\PYG{n}{j}\PYG{p}{]}\PYG{p}{[}\PYG{n}{i}\PYG{p}{]}\PYG{+w}{ }\PYG{o}{=}\PYG{+w}{ }\PYG{n}{biga1}\PYG{p}{[}\PYG{n}{i}\PYG{p}{]}\PYG{p}{[}\PYG{n}{j}\PYG{p}{]}\PYG{p}{;}
\PYG{+w}{  }\PYG{n}{beta}\PYG{p}{[}\PYG{n}{j}\PYG{p}{]}\PYG{p}{[}\PYG{n}{i}\PYG{p}{]}\PYG{+w}{ }\PYG{o}{=}\PYG{+w}{ }\PYG{n}{beta}\PYG{p}{[}\PYG{n}{i}\PYG{p}{]}\PYG{p}{[}\PYG{n}{j}\PYG{p}{]}\PYG{p}{;}
\PYG{+w}{  }\PYG{n}{r0}\PYG{p}{[}\PYG{n}{j}\PYG{p}{]}\PYG{p}{[}\PYG{n}{i}\PYG{p}{]}\PYG{+w}{ }\PYG{o}{=}\PYG{+w}{ }\PYG{n}{r0}\PYG{p}{[}\PYG{n}{i}\PYG{p}{]}\PYG{p}{[}\PYG{n}{j}\PYG{p}{]}\PYG{p}{;}
\PYG{+w}{  }\PYG{n}{offset}\PYG{p}{[}\PYG{n}{j}\PYG{p}{]}\PYG{p}{[}\PYG{n}{i}\PYG{p}{]}\PYG{+w}{ }\PYG{o}{=}\PYG{+w}{ }\PYG{n}{offset}\PYG{p}{[}\PYG{n}{i}\PYG{p}{]}\PYG{p}{[}\PYG{n}{j}\PYG{p}{]}\PYG{p}{;}

\PYG{+w}{  }\PYG{k}{return}\PYG{+w}{ }\PYG{n}{cut}\PYG{p}{[}\PYG{n}{i}\PYG{p}{]}\PYG{p}{[}\PYG{n}{j}\PYG{p}{]}\PYG{p}{;}
\PYG{p}{\PYGZcb{}}
\end{sphinxVerbatim}


\subsubsection{Computing forces from the neighbor list (required)}
\label{\detokenize{Developer_write_pair:computing-forces-from-the-neighbor-list-required}}
\sphinxAtStartPar
The \sphinxcode{\sphinxupquote{compute()}} function is the “workhorse” of a pair style.  This is
where we have the nested loops over all pairs of particles from the
neighbor list to compute forces and \sphinxhyphen{} if needed \sphinxhyphen{} energies and virials.

\sphinxAtStartPar
The first part is to define some variables for later use and store
cached copies of data or pointers that we need to access frequently.  Also,
this is a good place to call \sphinxcode{\sphinxupquote{Pair::ev\_init()}}, which initializes
several flags derived from the \sphinxtitleref{eflag} and \sphinxtitleref{vflag} parameters signaling
whether the energy and virial need to be tallied and whether only globally
or also per\sphinxhyphen{}atom.

\begin{sphinxVerbatim}[commandchars=\\\{\}]
\PYG{c+cm}{/* \PYGZhy{}\PYGZhy{}\PYGZhy{}\PYGZhy{}\PYGZhy{}\PYGZhy{}\PYGZhy{}\PYGZhy{}\PYGZhy{}\PYGZhy{}\PYGZhy{}\PYGZhy{}\PYGZhy{}\PYGZhy{}\PYGZhy{}\PYGZhy{}\PYGZhy{}\PYGZhy{}\PYGZhy{}\PYGZhy{}\PYGZhy{}\PYGZhy{}\PYGZhy{}\PYGZhy{}\PYGZhy{}\PYGZhy{}\PYGZhy{}\PYGZhy{}\PYGZhy{}\PYGZhy{}\PYGZhy{}\PYGZhy{}\PYGZhy{}\PYGZhy{}\PYGZhy{}\PYGZhy{}\PYGZhy{}\PYGZhy{}\PYGZhy{}\PYGZhy{}\PYGZhy{}\PYGZhy{}\PYGZhy{}\PYGZhy{}\PYGZhy{}\PYGZhy{}\PYGZhy{}\PYGZhy{}\PYGZhy{}\PYGZhy{}\PYGZhy{}\PYGZhy{}\PYGZhy{}\PYGZhy{}\PYGZhy{}\PYGZhy{}\PYGZhy{}\PYGZhy{}\PYGZhy{}\PYGZhy{}\PYGZhy{}\PYGZhy{}\PYGZhy{}\PYGZhy{}\PYGZhy{}\PYGZhy{}\PYGZhy{}\PYGZhy{}\PYGZhy{}\PYGZhy{} */}

\PYG{k+kt}{void}\PYG{+w}{ }\PYG{n+nf}{PairBornGauss::compute}\PYG{p}{(}\PYG{k+kt}{int}\PYG{+w}{ }\PYG{n}{eflag}\PYG{p}{,}\PYG{+w}{ }\PYG{k+kt}{int}\PYG{+w}{ }\PYG{n}{vflag}\PYG{p}{)}
\PYG{p}{\PYGZob{}}
\PYG{+w}{  }\PYG{k+kt}{int}\PYG{+w}{ }\PYG{n}{i}\PYG{p}{,}\PYG{+w}{ }\PYG{n}{j}\PYG{p}{,}\PYG{+w}{ }\PYG{n}{ii}\PYG{p}{,}\PYG{+w}{ }\PYG{n}{jj}\PYG{p}{,}\PYG{+w}{ }\PYG{n}{inum}\PYG{p}{,}\PYG{+w}{ }\PYG{n}{jnum}\PYG{p}{,}\PYG{+w}{ }\PYG{n}{itype}\PYG{p}{,}\PYG{+w}{ }\PYG{n}{jtype}\PYG{p}{;}
\PYG{+w}{  }\PYG{k+kt}{double}\PYG{+w}{ }\PYG{n}{xtmp}\PYG{p}{,}\PYG{+w}{ }\PYG{n}{ytmp}\PYG{p}{,}\PYG{+w}{ }\PYG{n}{ztmp}\PYG{p}{,}\PYG{+w}{ }\PYG{n}{delx}\PYG{p}{,}\PYG{+w}{ }\PYG{n}{dely}\PYG{p}{,}\PYG{+w}{ }\PYG{n}{delz}\PYG{p}{,}\PYG{+w}{ }\PYG{n}{evdwl}\PYG{p}{,}\PYG{+w}{ }\PYG{n}{fpair}\PYG{p}{;}
\PYG{+w}{  }\PYG{k+kt}{double}\PYG{+w}{ }\PYG{n}{rsq}\PYG{p}{,}\PYG{+w}{ }\PYG{n}{r}\PYG{p}{,}\PYG{+w}{ }\PYG{n}{dr}\PYG{p}{,}\PYG{+w}{ }\PYG{n}{aexp}\PYG{p}{,}\PYG{+w}{ }\PYG{n}{bexp}\PYG{p}{,}\PYG{+w}{ }\PYG{n}{factor\PYGZus{}lj}\PYG{p}{;}
\PYG{+w}{  }\PYG{k+kt}{int}\PYG{+w}{ }\PYG{o}{*}\PYG{n}{ilist}\PYG{p}{,}\PYG{+w}{ }\PYG{o}{*}\PYG{n}{jlist}\PYG{p}{,}\PYG{+w}{ }\PYG{o}{*}\PYG{n}{numneigh}\PYG{p}{,}\PYG{+w}{ }\PYG{o}{*}\PYG{o}{*}\PYG{n}{firstneigh}\PYG{p}{;}

\PYG{+w}{  }\PYG{n}{evdwl}\PYG{+w}{ }\PYG{o}{=}\PYG{+w}{ }\PYG{l+m+mf}{0.0}\PYG{p}{;}
\PYG{+w}{  }\PYG{n}{ev\PYGZus{}init}\PYG{p}{(}\PYG{n}{eflag}\PYG{p}{,}\PYG{+w}{ }\PYG{n}{vflag}\PYG{p}{)}\PYG{p}{;}

\PYG{+w}{  }\PYG{k+kt}{double}\PYG{+w}{ }\PYG{o}{*}\PYG{o}{*}\PYG{n}{x}\PYG{+w}{ }\PYG{o}{=}\PYG{+w}{ }\PYG{n}{atom}\PYG{o}{\PYGZhy{}}\PYG{o}{\PYGZgt{}}\PYG{n}{x}\PYG{p}{;}
\PYG{+w}{  }\PYG{k+kt}{double}\PYG{+w}{ }\PYG{o}{*}\PYG{o}{*}\PYG{n}{f}\PYG{+w}{ }\PYG{o}{=}\PYG{+w}{ }\PYG{n}{atom}\PYG{o}{\PYGZhy{}}\PYG{o}{\PYGZgt{}}\PYG{n}{f}\PYG{p}{;}
\PYG{+w}{  }\PYG{k+kt}{int}\PYG{+w}{ }\PYG{o}{*}\PYG{n}{type}\PYG{+w}{ }\PYG{o}{=}\PYG{+w}{ }\PYG{n}{atom}\PYG{o}{\PYGZhy{}}\PYG{o}{\PYGZgt{}}\PYG{n}{type}\PYG{p}{;}
\PYG{+w}{  }\PYG{k+kt}{int}\PYG{+w}{ }\PYG{n}{nlocal}\PYG{+w}{ }\PYG{o}{=}\PYG{+w}{ }\PYG{n}{atom}\PYG{o}{\PYGZhy{}}\PYG{o}{\PYGZgt{}}\PYG{n}{nlocal}\PYG{p}{;}
\PYG{+w}{  }\PYG{k+kt}{double}\PYG{+w}{ }\PYG{o}{*}\PYG{n}{special\PYGZus{}lj}\PYG{+w}{ }\PYG{o}{=}\PYG{+w}{ }\PYG{n}{force}\PYG{o}{\PYGZhy{}}\PYG{o}{\PYGZgt{}}\PYG{n}{special\PYGZus{}lj}\PYG{p}{;}
\PYG{+w}{  }\PYG{k+kt}{int}\PYG{+w}{ }\PYG{n}{newton\PYGZus{}pair}\PYG{+w}{ }\PYG{o}{=}\PYG{+w}{ }\PYG{n}{force}\PYG{o}{\PYGZhy{}}\PYG{o}{\PYGZgt{}}\PYG{n}{newton\PYGZus{}pair}\PYG{p}{;}

\PYG{+w}{  }\PYG{n}{inum}\PYG{+w}{ }\PYG{o}{=}\PYG{+w}{ }\PYG{n}{list}\PYG{o}{\PYGZhy{}}\PYG{o}{\PYGZgt{}}\PYG{n}{inum}\PYG{p}{;}
\PYG{+w}{  }\PYG{n}{ilist}\PYG{+w}{ }\PYG{o}{=}\PYG{+w}{ }\PYG{n}{list}\PYG{o}{\PYGZhy{}}\PYG{o}{\PYGZgt{}}\PYG{n}{ilist}\PYG{p}{;}
\PYG{+w}{  }\PYG{n}{numneigh}\PYG{+w}{ }\PYG{o}{=}\PYG{+w}{ }\PYG{n}{list}\PYG{o}{\PYGZhy{}}\PYG{o}{\PYGZgt{}}\PYG{n}{numneigh}\PYG{p}{;}
\PYG{+w}{  }\PYG{n}{firstneigh}\PYG{+w}{ }\PYG{o}{=}\PYG{+w}{ }\PYG{n}{list}\PYG{o}{\PYGZhy{}}\PYG{o}{\PYGZgt{}}\PYG{n}{firstneigh}\PYG{p}{;}
\end{sphinxVerbatim}

\sphinxAtStartPar
The outer loop (index \sphinxstyleemphasis{i}) is over local atoms of our sub\sphinxhyphen{}domain.
Typically, the value of \sphinxtitleref{inum} (the number of neighbor lists) is the
same as the number of local atoms (= atoms \sphinxstyleemphasis{owned} by this sub\sphinxhyphen{}domain).
But when the pair style is used as a sub\sphinxhyphen{}style of a {\hyperref[\detokenize{pair_hybrid::doc}]{\sphinxcrossref{\DUrole{doc}{hybrid pair
style}}}} or neighbor list entries are removed with
{\hyperref[\detokenize{neigh_modify::doc}]{\sphinxcrossref{\DUrole{doc}{neigh\_modify exclude}}}}, this number may be
smaller. The array \sphinxcode{\sphinxupquote{list\sphinxhyphen{}\textgreater{}ilist}} has the (local) indices of the atoms
for which neighbor lists have been created. Then \sphinxcode{\sphinxupquote{list\sphinxhyphen{}\textgreater{}numneigh}} is
an \sphinxtitleref{inum} sized array with the number of entries of each list of
neighbors, and \sphinxcode{\sphinxupquote{list\sphinxhyphen{}\textgreater{}firstneigh}} is a list of pointers to those lists.

\sphinxAtStartPar
For efficiency reasons, cached copies of some properties of the outer
loop atoms are also initialized.

\begin{sphinxVerbatim}[commandchars=\\\{\}]
\PYG{c+c1}{// loop over neighbors of my atoms}

\PYG{k}{for}\PYG{+w}{ }\PYG{p}{(}\PYG{n}{ii}\PYG{+w}{ }\PYG{o}{=}\PYG{+w}{ }\PYG{l+m+mi}{0}\PYG{p}{;}\PYG{+w}{ }\PYG{n}{ii}\PYG{+w}{ }\PYG{o}{\PYGZlt{}}\PYG{+w}{ }\PYG{n}{inum}\PYG{p}{;}\PYG{+w}{ }\PYG{n}{ii}\PYG{o}{+}\PYG{o}{+}\PYG{p}{)}\PYG{+w}{ }\PYG{p}{\PYGZob{}}
\PYG{+w}{  }\PYG{n}{i}\PYG{+w}{ }\PYG{o}{=}\PYG{+w}{ }\PYG{n}{ilist}\PYG{p}{[}\PYG{n}{ii}\PYG{p}{]}\PYG{p}{;}
\PYG{+w}{  }\PYG{n}{xtmp}\PYG{+w}{ }\PYG{o}{=}\PYG{+w}{ }\PYG{n}{x}\PYG{p}{[}\PYG{n}{i}\PYG{p}{]}\PYG{p}{[}\PYG{l+m+mi}{0}\PYG{p}{]}\PYG{p}{;}
\PYG{+w}{  }\PYG{n}{ytmp}\PYG{+w}{ }\PYG{o}{=}\PYG{+w}{ }\PYG{n}{x}\PYG{p}{[}\PYG{n}{i}\PYG{p}{]}\PYG{p}{[}\PYG{l+m+mi}{1}\PYG{p}{]}\PYG{p}{;}
\PYG{+w}{  }\PYG{n}{ztmp}\PYG{+w}{ }\PYG{o}{=}\PYG{+w}{ }\PYG{n}{x}\PYG{p}{[}\PYG{n}{i}\PYG{p}{]}\PYG{p}{[}\PYG{l+m+mi}{2}\PYG{p}{]}\PYG{p}{;}
\PYG{+w}{  }\PYG{n}{itype}\PYG{+w}{ }\PYG{o}{=}\PYG{+w}{ }\PYG{n}{type}\PYG{p}{[}\PYG{n}{i}\PYG{p}{]}\PYG{p}{;}
\PYG{+w}{  }\PYG{n}{jlist}\PYG{+w}{ }\PYG{o}{=}\PYG{+w}{ }\PYG{n}{firstneigh}\PYG{p}{[}\PYG{n}{i}\PYG{p}{]}\PYG{p}{;}
\PYG{+w}{  }\PYG{n}{jnum}\PYG{+w}{ }\PYG{o}{=}\PYG{+w}{ }\PYG{n}{numneigh}\PYG{p}{[}\PYG{n}{i}\PYG{p}{]}\PYG{p}{;}
\end{sphinxVerbatim}

\sphinxAtStartPar
The inner loop (index \sphinxstyleemphasis{j}) processes the neighbor lists.  The neighbor
list code encodes extra information using the upper 3 bits. The 2
highest bits encode whether a pair is a regular pair of neighbor (= 0)
or a pair of 1\sphinxhyphen{}2 (= 1), 1\sphinxhyphen{}3 (= 2), or 1\sphinxhyphen{}4 (= 3) {\hyperref[\detokenize{special_bonds::doc}]{\sphinxcrossref{\DUrole{doc}{“special” neighbor}}}}.  The next highest bit encodes whether the pair stores
data in a \sphinxcode{\sphinxupquote{fix neigh/history}} instance (an undocumented internal fix
style).  The \sphinxcode{\sphinxupquote{sbmask()}} inline function extracts those bits and
converts them into a number.  This number is used to look up the
corresponding scaling factor for the non\sphinxhyphen{}bonded interaction from the
\sphinxcode{\sphinxupquote{force\sphinxhyphen{}\textgreater{}special\_lj}} array and stores it in the \sphinxtitleref{factor\_lj} variable.
Due to the additional bits, the value of \sphinxstyleemphasis{j} would be out of range when
accessing data from per\sphinxhyphen{}atom arrays, so we apply the NEIGHMASK constant
with a bit\sphinxhyphen{}wise and operation to mask them out.  This step \sphinxstyleemphasis{must} be
done, even if a pair style does not use special bond scaling of forces
and energies to avoid segmentation faults.

\sphinxAtStartPar
With the corrected \sphinxstyleemphasis{j} index, it is now possible to compute the distance
of the pair.  For efficiency reasons, the square root is only taken
\sphinxstyleemphasis{after} the check for the cutoff (which has been stored as squared
cutoff by the \sphinxcode{\sphinxupquote{Pair}} base class).  For some pair styles, like the 12\sphinxhyphen{}6
Lennard\sphinxhyphen{}Jones potential, computing the square root can be avoided
entirely.

\begin{sphinxVerbatim}[commandchars=\\\{\}]
\PYG{k}{for}\PYG{+w}{ }\PYG{p}{(}\PYG{n}{jj}\PYG{+w}{ }\PYG{o}{=}\PYG{+w}{ }\PYG{l+m+mi}{0}\PYG{p}{;}\PYG{+w}{ }\PYG{n}{jj}\PYG{+w}{ }\PYG{o}{\PYGZlt{}}\PYG{+w}{ }\PYG{n}{jnum}\PYG{p}{;}\PYG{+w}{ }\PYG{n}{jj}\PYG{o}{+}\PYG{o}{+}\PYG{p}{)}\PYG{+w}{ }\PYG{p}{\PYGZob{}}
\PYG{+w}{  }\PYG{n}{j}\PYG{+w}{ }\PYG{o}{=}\PYG{+w}{ }\PYG{n}{jlist}\PYG{p}{[}\PYG{n}{jj}\PYG{p}{]}\PYG{p}{;}
\PYG{+w}{  }\PYG{n}{factor\PYGZus{}lj}\PYG{+w}{ }\PYG{o}{=}\PYG{+w}{ }\PYG{n}{special\PYGZus{}lj}\PYG{p}{[}\PYG{n}{sbmask}\PYG{p}{(}\PYG{n}{j}\PYG{p}{)}\PYG{p}{]}\PYG{p}{;}
\PYG{+w}{  }\PYG{n}{j}\PYG{+w}{ }\PYG{o}{\PYGZam{}}\PYG{o}{=}\PYG{+w}{ }\PYG{n}{NEIGHMASK}\PYG{p}{;}

\PYG{+w}{  }\PYG{n}{delx}\PYG{+w}{ }\PYG{o}{=}\PYG{+w}{ }\PYG{n}{xtmp}\PYG{+w}{ }\PYG{o}{\PYGZhy{}}\PYG{+w}{ }\PYG{n}{x}\PYG{p}{[}\PYG{n}{j}\PYG{p}{]}\PYG{p}{[}\PYG{l+m+mi}{0}\PYG{p}{]}\PYG{p}{;}
\PYG{+w}{  }\PYG{n}{dely}\PYG{+w}{ }\PYG{o}{=}\PYG{+w}{ }\PYG{n}{ytmp}\PYG{+w}{ }\PYG{o}{\PYGZhy{}}\PYG{+w}{ }\PYG{n}{x}\PYG{p}{[}\PYG{n}{j}\PYG{p}{]}\PYG{p}{[}\PYG{l+m+mi}{1}\PYG{p}{]}\PYG{p}{;}
\PYG{+w}{  }\PYG{n}{delz}\PYG{+w}{ }\PYG{o}{=}\PYG{+w}{ }\PYG{n}{ztmp}\PYG{+w}{ }\PYG{o}{\PYGZhy{}}\PYG{+w}{ }\PYG{n}{x}\PYG{p}{[}\PYG{n}{j}\PYG{p}{]}\PYG{p}{[}\PYG{l+m+mi}{2}\PYG{p}{]}\PYG{p}{;}
\PYG{+w}{  }\PYG{n}{rsq}\PYG{+w}{ }\PYG{o}{=}\PYG{+w}{ }\PYG{n}{delx}\PYG{+w}{ }\PYG{o}{*}\PYG{+w}{ }\PYG{n}{delx}\PYG{+w}{ }\PYG{o}{+}\PYG{+w}{ }\PYG{n}{dely}\PYG{+w}{ }\PYG{o}{*}\PYG{+w}{ }\PYG{n}{dely}\PYG{+w}{ }\PYG{o}{+}\PYG{+w}{ }\PYG{n}{delz}\PYG{+w}{ }\PYG{o}{*}\PYG{+w}{ }\PYG{n}{delz}\PYG{p}{;}
\PYG{+w}{  }\PYG{n}{jtype}\PYG{+w}{ }\PYG{o}{=}\PYG{+w}{ }\PYG{n}{type}\PYG{p}{[}\PYG{n}{j}\PYG{p}{]}\PYG{p}{;}
\end{sphinxVerbatim}

\sphinxAtStartPar
The following block of code is the actual application of the model
potential to compute the force.  Note, that \sphinxstyleemphasis{fpair} is the pair\sphinxhyphen{}wise
force divided by the distance, as this simplifies the projection of the
x\sphinxhyphen{}, y\sphinxhyphen{}, and z\sphinxhyphen{}components of the force vector by simply multiplying with
the respective distances in those directions.

\begin{sphinxVerbatim}[commandchars=\\\{\}]
\PYG{k}{if}\PYG{+w}{ }\PYG{p}{(}\PYG{n}{rsq}\PYG{+w}{ }\PYG{o}{\PYGZlt{}}\PYG{+w}{ }\PYG{n}{cutsq}\PYG{p}{[}\PYG{n}{itype}\PYG{p}{]}\PYG{p}{[}\PYG{n}{jtype}\PYG{p}{]}\PYG{p}{)}\PYG{+w}{ }\PYG{p}{\PYGZob{}}
\PYG{+w}{  }\PYG{n}{r}\PYG{+w}{ }\PYG{o}{=}\PYG{+w}{ }\PYG{n}{sqrt}\PYG{p}{(}\PYG{n}{rsq}\PYG{p}{)}\PYG{p}{;}
\PYG{+w}{  }\PYG{n}{dr}\PYG{+w}{ }\PYG{o}{=}\PYG{+w}{ }\PYG{n}{r}\PYG{+w}{ }\PYG{o}{\PYGZhy{}}\PYG{+w}{ }\PYG{n}{r0}\PYG{p}{[}\PYG{n}{itype}\PYG{p}{]}\PYG{p}{[}\PYG{n}{jtype}\PYG{p}{]}\PYG{p}{;}
\PYG{+w}{  }\PYG{n}{aexp}\PYG{+w}{ }\PYG{o}{=}\PYG{+w}{ }\PYG{n}{biga0}\PYG{p}{[}\PYG{n}{itype}\PYG{p}{]}\PYG{p}{[}\PYG{n}{jtype}\PYG{p}{]}\PYG{+w}{ }\PYG{o}{*}\PYG{+w}{ }\PYG{n}{exp}\PYG{p}{(}\PYG{o}{\PYGZhy{}}\PYG{n}{alpha}\PYG{p}{[}\PYG{n}{itype}\PYG{p}{]}\PYG{p}{[}\PYG{n}{jtype}\PYG{p}{]}\PYG{+w}{ }\PYG{o}{*}\PYG{+w}{ }\PYG{n}{r}\PYG{p}{)}\PYG{p}{;}
\PYG{+w}{  }\PYG{n}{bexp}\PYG{+w}{ }\PYG{o}{=}\PYG{+w}{ }\PYG{n}{biga1}\PYG{p}{[}\PYG{n}{itype}\PYG{p}{]}\PYG{p}{[}\PYG{n}{jtype}\PYG{p}{]}\PYG{+w}{ }\PYG{o}{*}\PYG{+w}{ }\PYG{n}{exp}\PYG{p}{(}\PYG{o}{\PYGZhy{}}\PYG{n}{beta}\PYG{p}{[}\PYG{n}{itype}\PYG{p}{]}\PYG{p}{[}\PYG{n}{jtype}\PYG{p}{]}\PYG{+w}{ }\PYG{o}{*}\PYG{+w}{ }\PYG{n}{dr}\PYG{+w}{ }\PYG{o}{*}\PYG{+w}{ }\PYG{n}{dr}\PYG{p}{)}\PYG{p}{;}
\PYG{+w}{  }\PYG{n}{fpair}\PYG{+w}{ }\PYG{o}{=}\PYG{+w}{ }\PYG{n}{alpha}\PYG{p}{[}\PYG{n}{itype}\PYG{p}{]}\PYG{p}{[}\PYG{n}{jtype}\PYG{p}{]}\PYG{+w}{ }\PYG{o}{*}\PYG{+w}{ }\PYG{n}{aexp}\PYG{p}{;}
\PYG{+w}{  }\PYG{n}{fpair}\PYG{+w}{ }\PYG{o}{\PYGZhy{}}\PYG{o}{=}\PYG{+w}{ }\PYG{l+m+mf}{2.0}\PYG{+w}{ }\PYG{o}{*}\PYG{+w}{ }\PYG{n}{beta}\PYG{p}{[}\PYG{n}{itype}\PYG{p}{]}\PYG{p}{[}\PYG{n}{jtype}\PYG{p}{]}\PYG{+w}{ }\PYG{o}{*}\PYG{+w}{ }\PYG{n}{dr}\PYG{+w}{ }\PYG{o}{*}\PYG{+w}{ }\PYG{n}{bexp}\PYG{p}{;}
\PYG{+w}{  }\PYG{n}{fpair}\PYG{+w}{ }\PYG{o}{*}\PYG{o}{=}\PYG{+w}{ }\PYG{n}{factor\PYGZus{}lj}\PYG{+w}{ }\PYG{o}{/}\PYG{+w}{ }\PYG{n}{r}\PYG{p}{;}
\end{sphinxVerbatim}

\sphinxAtStartPar
In the next block, the force is added to the per\sphinxhyphen{}atom force arrays.  This
pair style uses a “half” neighbor list (each pair is listed only once)
so we take advantage of the fact that \(\vec{F}_{ij} =
-\vec{F}_{ji}\), i.e.  apply Newton’s third law.  The force is \sphinxstyleemphasis{always}
stored when the atom is a “local” atom. Index \sphinxstyleemphasis{i} atoms are always “local”
(i.e. \sphinxstyleemphasis{i} \textless{} nlocal); index \sphinxstyleemphasis{j} atoms may be “ghost” atoms (\sphinxstyleemphasis{j} \textgreater{}= nlocal).

\sphinxAtStartPar
Depending on the settings used with the {\hyperref[\detokenize{newton::doc}]{\sphinxcrossref{\DUrole{doc}{newton command}}}},
those pairs are only listed once globally (newton\_pair == 1), then
forces must be stored even with ghost atoms and after all forces are
computed a “reverse communication” is performed to add those ghost atom
forces to their corresponding local atoms.  If the setting is disabled,
then the extra communication is skipped, since for pairs straddling
sub\sphinxhyphen{}domain boundaries, the forces are computed twice and only stored
with the local atoms in the domain that \sphinxstyleemphasis{owns} it.

\begin{sphinxVerbatim}[commandchars=\\\{\}]
\PYG{n}{f}\PYG{p}{[}\PYG{n}{i}\PYG{p}{]}\PYG{p}{[}\PYG{l+m+mi}{0}\PYG{p}{]}\PYG{+w}{ }\PYG{o}{+}\PYG{o}{=}\PYG{+w}{ }\PYG{n}{delx}\PYG{+w}{ }\PYG{o}{*}\PYG{+w}{ }\PYG{n}{fpair}\PYG{p}{;}
\PYG{n}{f}\PYG{p}{[}\PYG{n}{i}\PYG{p}{]}\PYG{p}{[}\PYG{l+m+mi}{1}\PYG{p}{]}\PYG{+w}{ }\PYG{o}{+}\PYG{o}{=}\PYG{+w}{ }\PYG{n}{dely}\PYG{+w}{ }\PYG{o}{*}\PYG{+w}{ }\PYG{n}{fpair}\PYG{p}{;}
\PYG{n}{f}\PYG{p}{[}\PYG{n}{i}\PYG{p}{]}\PYG{p}{[}\PYG{l+m+mi}{2}\PYG{p}{]}\PYG{+w}{ }\PYG{o}{+}\PYG{o}{=}\PYG{+w}{ }\PYG{n}{delz}\PYG{+w}{ }\PYG{o}{*}\PYG{+w}{ }\PYG{n}{fpair}\PYG{p}{;}
\PYG{k}{if}\PYG{+w}{ }\PYG{p}{(}\PYG{n}{newton\PYGZus{}pair}\PYG{+w}{ }\PYG{o}{|}\PYG{o}{|}\PYG{+w}{ }\PYG{n}{j}\PYG{+w}{ }\PYG{o}{\PYGZlt{}}\PYG{+w}{ }\PYG{n}{nlocal}\PYG{p}{)}\PYG{+w}{ }\PYG{p}{\PYGZob{}}
\PYG{+w}{  }\PYG{n}{f}\PYG{p}{[}\PYG{n}{j}\PYG{p}{]}\PYG{p}{[}\PYG{l+m+mi}{0}\PYG{p}{]}\PYG{+w}{ }\PYG{o}{\PYGZhy{}}\PYG{o}{=}\PYG{+w}{ }\PYG{n}{delx}\PYG{+w}{ }\PYG{o}{*}\PYG{+w}{ }\PYG{n}{fpair}\PYG{p}{;}
\PYG{+w}{  }\PYG{n}{f}\PYG{p}{[}\PYG{n}{j}\PYG{p}{]}\PYG{p}{[}\PYG{l+m+mi}{1}\PYG{p}{]}\PYG{+w}{ }\PYG{o}{\PYGZhy{}}\PYG{o}{=}\PYG{+w}{ }\PYG{n}{dely}\PYG{+w}{ }\PYG{o}{*}\PYG{+w}{ }\PYG{n}{fpair}\PYG{p}{;}
\PYG{+w}{  }\PYG{n}{f}\PYG{p}{[}\PYG{n}{j}\PYG{p}{]}\PYG{p}{[}\PYG{l+m+mi}{2}\PYG{p}{]}\PYG{+w}{ }\PYG{o}{\PYGZhy{}}\PYG{o}{=}\PYG{+w}{ }\PYG{n}{delz}\PYG{+w}{ }\PYG{o}{*}\PYG{+w}{ }\PYG{n}{fpair}\PYG{p}{;}
\PYG{p}{\PYGZcb{}}
\end{sphinxVerbatim}

\sphinxAtStartPar
The \sphinxcode{\sphinxupquote{ev\_tally()}} function tallies global or per\sphinxhyphen{}atom energy and
virial.  For typical MD simulations, the potential energy is merely a
diagnostic and only needed on output.  Similarly, the pressure may only
be computed for (infrequent) thermodynamic output.  For all timesteps
where this information is not needed either, \sphinxtitleref{eflag} or \sphinxtitleref{evflag} are
zero and the computation and call to the tally function skipped.  Note
that evdwl is initialized to zero at the beginning of the function, so
that it still is valid to access it, even if the energy is not computed
(e.g. when only the virial is needed).

\begin{sphinxVerbatim}[commandchars=\\\{\}]
\PYG{+w}{      }\PYG{k}{if}\PYG{+w}{ }\PYG{p}{(}\PYG{n}{eflag}\PYG{p}{)}\PYG{+w}{ }\PYG{n}{evdwl}\PYG{+w}{ }\PYG{o}{=}\PYG{+w}{ }\PYG{n}{factor\PYGZus{}lj}\PYG{+w}{ }\PYG{o}{*}\PYG{+w}{ }\PYG{p}{(}\PYG{n}{aexp}\PYG{+w}{ }\PYG{o}{\PYGZhy{}}\PYG{+w}{ }\PYG{n}{bexp}\PYG{+w}{ }\PYG{o}{\PYGZhy{}}\PYG{+w}{ }\PYG{n}{offset}\PYG{p}{[}\PYG{n}{itype}\PYG{p}{]}\PYG{p}{[}\PYG{n}{jtype}\PYG{p}{]}\PYG{p}{)}\PYG{p}{;}
\PYG{+w}{      }\PYG{k}{if}\PYG{+w}{ }\PYG{p}{(}\PYG{n}{evflag}\PYG{p}{)}\PYG{+w}{ }\PYG{n}{ev\PYGZus{}tally}\PYG{p}{(}\PYG{n}{i}\PYG{p}{,}\PYG{+w}{ }\PYG{n}{j}\PYG{p}{,}\PYG{+w}{ }\PYG{n}{nlocal}\PYG{p}{,}\PYG{+w}{ }\PYG{n}{newton\PYGZus{}pair}\PYG{p}{,}\PYG{+w}{ }\PYG{n}{evdwl}\PYG{p}{,}\PYG{+w}{ }\PYG{l+m+mf}{0.0}\PYG{p}{,}\PYG{+w}{ }\PYG{n}{fpair}\PYG{p}{,}\PYG{+w}{ }\PYG{n}{delx}\PYG{p}{,}\PYG{+w}{ }\PYG{n}{dely}\PYG{p}{,}\PYG{+w}{ }\PYG{n}{delz}\PYG{p}{)}\PYG{p}{;}
\PYG{+w}{    }\PYG{p}{\PYGZcb{}}
\PYG{+w}{  }\PYG{p}{\PYGZcb{}}
\PYG{p}{\PYGZcb{}}
\end{sphinxVerbatim}

\sphinxAtStartPar
If only the global virial is needed and no energy, then calls to
\sphinxcode{\sphinxupquote{ev\_tally()}} can be avoided altogether, and the global virial can be
computed more efficiently from the dot product of the total per\sphinxhyphen{}atom
force vector and the position vector of the corresponding atom,
\(\vec{F}\cdot\vec{r}\).  This has to be done \sphinxstyleemphasis{after} all pair\sphinxhyphen{}wise
forces are computed and \sphinxstyleemphasis{before} the reverse communication to collect
data from ghost atoms, since the position has to be the position that was
used to compute the force, i.e. \sphinxstyleemphasis{not} the “local” position if that ghost
atom is a periodic copy.

\begin{sphinxVerbatim}[commandchars=\\\{\}]
\PYG{+w}{  }\PYG{k}{if}\PYG{+w}{ }\PYG{p}{(}\PYG{n}{vflag\PYGZus{}fdotr}\PYG{p}{)}\PYG{+w}{ }\PYG{n}{virial\PYGZus{}fdotr\PYGZus{}compute}\PYG{p}{(}\PYG{p}{)}\PYG{p}{;}
\PYG{p}{\PYGZcb{}}
\end{sphinxVerbatim}


\subsubsection{Computing force and energy for a single pair}
\label{\detokenize{Developer_write_pair:computing-force-and-energy-for-a-single-pair}}
\sphinxAtStartPar
Certain features in LAMMPS only require computing interactions between
individual pairs of atoms and the (optional) \sphinxcode{\sphinxupquote{single()}} function is
needed to support those features (e.g. for tabulation of force and
energy with {\hyperref[\detokenize{pair_write::doc}]{\sphinxcrossref{\DUrole{doc}{pair\_write}}}}).  This is a repetition of
the force kernel in the \sphinxcode{\sphinxupquote{compute()}} function, but only for a single
pair of atoms, where the (squared) distance is provided as a parameter
(so it may not even be an existing distance between two specific atoms).
The energy is returned as the return value of the function and the force
as the \sphinxtitleref{fforce} reference.  Note, that this is, similar to how \sphinxstyleemphasis{fpair}
is used in the \sphinxcode{\sphinxupquote{compute()}} function, the magnitude of the force along
the vector between the two atoms \sphinxstyleemphasis{divided} by the distance.

\sphinxAtStartPar
The \sphinxcode{\sphinxupquote{single()}} function is optional, but it is expected to be
implemented for any true pair\sphinxhyphen{}wise additive potential. Many\sphinxhyphen{}body
potentials and special case potentials do not implement it. In a few
special cases (EAM, long\sphinxhyphen{}range Coulomb), the \sphinxcode{\sphinxupquote{single()}} function
implements the pairwise additive part of the complete force interaction
and depends on either pre\sphinxhyphen{}computed properties (derivative of embedding
term for EAM) or post\sphinxhyphen{}computed non\sphinxhyphen{}pair\sphinxhyphen{}wise force contributions (KSpace
style in case of long\sphinxhyphen{}range Coulomb).

\sphinxAtStartPar
The member variable \sphinxtitleref{single\_enable} should be set to 0 in the
constructor, if it is not implemented (its default value is 1).

\begin{sphinxVerbatim}[commandchars=\\\{\}]
\PYG{c+cm}{/* \PYGZhy{}\PYGZhy{}\PYGZhy{}\PYGZhy{}\PYGZhy{}\PYGZhy{}\PYGZhy{}\PYGZhy{}\PYGZhy{}\PYGZhy{}\PYGZhy{}\PYGZhy{}\PYGZhy{}\PYGZhy{}\PYGZhy{}\PYGZhy{}\PYGZhy{}\PYGZhy{}\PYGZhy{}\PYGZhy{}\PYGZhy{}\PYGZhy{}\PYGZhy{}\PYGZhy{}\PYGZhy{}\PYGZhy{}\PYGZhy{}\PYGZhy{}\PYGZhy{}\PYGZhy{}\PYGZhy{}\PYGZhy{}\PYGZhy{}\PYGZhy{}\PYGZhy{}\PYGZhy{}\PYGZhy{}\PYGZhy{}\PYGZhy{}\PYGZhy{}\PYGZhy{}\PYGZhy{}\PYGZhy{}\PYGZhy{}\PYGZhy{}\PYGZhy{}\PYGZhy{}\PYGZhy{}\PYGZhy{}\PYGZhy{}\PYGZhy{}\PYGZhy{}\PYGZhy{}\PYGZhy{}\PYGZhy{}\PYGZhy{}\PYGZhy{}\PYGZhy{}\PYGZhy{}\PYGZhy{}\PYGZhy{}\PYGZhy{}\PYGZhy{}\PYGZhy{}\PYGZhy{}\PYGZhy{}\PYGZhy{}\PYGZhy{}\PYGZhy{}\PYGZhy{} */}

\PYG{k+kt}{double}\PYG{+w}{ }\PYG{n+nf}{PairBornGauss::single}\PYG{p}{(}\PYG{k+kt}{int}\PYG{+w}{ }\PYG{c+cm}{/*i*/}\PYG{p}{,}\PYG{+w}{ }\PYG{k+kt}{int}\PYG{+w}{ }\PYG{c+cm}{/*j*/}\PYG{p}{,}\PYG{+w}{ }\PYG{k+kt}{int}\PYG{+w}{ }\PYG{n}{itype}\PYG{p}{,}\PYG{+w}{ }\PYG{k+kt}{int}\PYG{+w}{ }\PYG{n}{jtype}\PYG{p}{,}\PYG{+w}{ }\PYG{k+kt}{double}\PYG{+w}{ }\PYG{n}{rsq}\PYG{p}{,}
\PYG{+w}{                             }\PYG{k+kt}{double}\PYG{+w}{ }\PYG{c+cm}{/*factor\PYGZus{}coul*/}\PYG{p}{,}\PYG{+w}{ }\PYG{k+kt}{double}\PYG{+w}{ }\PYG{n}{factor\PYGZus{}lj}\PYG{p}{,}\PYG{+w}{ }\PYG{k+kt}{double}\PYG{+w}{ }\PYG{o}{\PYGZam{}}\PYG{n}{fforce}\PYG{p}{)}
\PYG{p}{\PYGZob{}}
\PYG{+w}{  }\PYG{k+kt}{double}\PYG{+w}{ }\PYG{n}{r}\PYG{p}{,}\PYG{+w}{ }\PYG{n}{dr}\PYG{p}{,}\PYG{+w}{ }\PYG{n}{aexp}\PYG{p}{,}\PYG{+w}{ }\PYG{n}{bexp}\PYG{p}{;}

\PYG{+w}{  }\PYG{n}{r}\PYG{+w}{ }\PYG{o}{=}\PYG{+w}{ }\PYG{n}{sqrt}\PYG{p}{(}\PYG{n}{rsq}\PYG{p}{)}\PYG{p}{;}
\PYG{+w}{  }\PYG{n}{dr}\PYG{+w}{ }\PYG{o}{=}\PYG{+w}{ }\PYG{n}{r}\PYG{+w}{ }\PYG{o}{\PYGZhy{}}\PYG{+w}{ }\PYG{n}{r0}\PYG{p}{[}\PYG{n}{itype}\PYG{p}{]}\PYG{p}{[}\PYG{n}{jtype}\PYG{p}{]}\PYG{p}{;}
\PYG{+w}{  }\PYG{n}{aexp}\PYG{+w}{ }\PYG{o}{=}\PYG{+w}{ }\PYG{n}{biga0}\PYG{p}{[}\PYG{n}{itype}\PYG{p}{]}\PYG{p}{[}\PYG{n}{jtype}\PYG{p}{]}\PYG{+w}{ }\PYG{o}{*}\PYG{+w}{ }\PYG{n}{exp}\PYG{p}{(}\PYG{o}{\PYGZhy{}}\PYG{n}{alpha}\PYG{p}{[}\PYG{n}{itype}\PYG{p}{]}\PYG{p}{[}\PYG{n}{jtype}\PYG{p}{]}\PYG{+w}{ }\PYG{o}{*}\PYG{+w}{ }\PYG{n}{r}\PYG{p}{)}\PYG{p}{;}
\PYG{+w}{  }\PYG{n}{bexp}\PYG{+w}{ }\PYG{o}{=}\PYG{+w}{ }\PYG{n}{biga1}\PYG{p}{[}\PYG{n}{itype}\PYG{p}{]}\PYG{p}{[}\PYG{n}{jtype}\PYG{p}{]}\PYG{+w}{ }\PYG{o}{*}\PYG{+w}{ }\PYG{n}{exp}\PYG{p}{(}\PYG{o}{\PYGZhy{}}\PYG{n}{beta}\PYG{p}{[}\PYG{n}{itype}\PYG{p}{]}\PYG{p}{[}\PYG{n}{jtype}\PYG{p}{]}\PYG{+w}{ }\PYG{o}{*}\PYG{+w}{ }\PYG{n}{dr}\PYG{+w}{ }\PYG{o}{*}\PYG{+w}{ }\PYG{n}{dr}\PYG{p}{)}\PYG{p}{;}

\PYG{+w}{  }\PYG{n}{fforce}\PYG{+w}{ }\PYG{o}{=}\PYG{+w}{ }\PYG{n}{factor\PYGZus{}lj}\PYG{+w}{ }\PYG{o}{*}\PYG{+w}{ }\PYG{p}{(}\PYG{n}{alpha}\PYG{p}{[}\PYG{n}{itype}\PYG{p}{]}\PYG{p}{[}\PYG{n}{jtype}\PYG{p}{]}\PYG{+w}{ }\PYG{o}{*}\PYG{+w}{ }\PYG{n}{aexp}\PYG{+w}{ }\PYG{o}{\PYGZhy{}}\PYG{+w}{ }\PYG{l+m+mf}{2.0}\PYG{+w}{ }\PYG{o}{*}\PYG{+w}{ }\PYG{n}{dr}\PYG{+w}{ }\PYG{o}{*}\PYG{+w}{ }\PYG{n}{beta}\PYG{p}{[}\PYG{n}{itype}\PYG{p}{]}\PYG{p}{[}\PYG{n}{jtype}\PYG{p}{]}\PYG{+w}{ }\PYG{o}{*}\PYG{+w}{ }\PYG{n}{bexp}\PYG{p}{)}\PYG{+w}{ }\PYG{o}{/}\PYG{+w}{ }\PYG{n}{r}\PYG{p}{;}
\PYG{+w}{  }\PYG{k}{return}\PYG{+w}{ }\PYG{n}{factor\PYGZus{}lj}\PYG{+w}{ }\PYG{o}{*}\PYG{+w}{ }\PYG{p}{(}\PYG{n}{aexp}\PYG{+w}{ }\PYG{o}{\PYGZhy{}}\PYG{+w}{ }\PYG{n}{bexp}\PYG{+w}{ }\PYG{o}{\PYGZhy{}}\PYG{+w}{ }\PYG{n}{offset}\PYG{p}{[}\PYG{n}{itype}\PYG{p}{]}\PYG{p}{[}\PYG{n}{jtype}\PYG{p}{]}\PYG{p}{)}\PYG{p}{;}
\PYG{p}{\PYGZcb{}}
\end{sphinxVerbatim}


\subsubsection{Reading and writing of restart files}
\label{\detokenize{Developer_write_pair:reading-and-writing-of-restart-files}}
\sphinxAtStartPar
Support for writing and reading binary restart files is provided by the
following four functions.  Writing is only done by MPI processor rank 0.
The output of global (not related to atom types) settings is usually
delegated to the \sphinxcode{\sphinxupquote{write\_restart\_settings()}} function.  This restart
facility is commonly only used, if there are small number of per\sphinxhyphen{}type
parameters.  For potentials that use per\sphinxhyphen{}element parameters or tabulated
data and read these from files, those parameters and the name of the
potential file are not written to restart files and the {\hyperref[\detokenize{pair_coeff::doc}]{\sphinxcrossref{\DUrole{doc}{pair\_coeff
command}}}} has to re\sphinxhyphen{}issued when restarting.  For pair styles
like “born/gauss” that do support writing to restart files, this is not
required.

\sphinxAtStartPar
Implementing the functions to read and write binary restart files is
optional.  The member variable \sphinxtitleref{restartinfo} should be set to 0 in the
constructor, if they are not implemented (its default value is 1).

\begin{sphinxVerbatim}[commandchars=\\\{\}]
\PYG{c+cm}{/* \PYGZhy{}\PYGZhy{}\PYGZhy{}\PYGZhy{}\PYGZhy{}\PYGZhy{}\PYGZhy{}\PYGZhy{}\PYGZhy{}\PYGZhy{}\PYGZhy{}\PYGZhy{}\PYGZhy{}\PYGZhy{}\PYGZhy{}\PYGZhy{}\PYGZhy{}\PYGZhy{}\PYGZhy{}\PYGZhy{}\PYGZhy{}\PYGZhy{}\PYGZhy{}\PYGZhy{}\PYGZhy{}\PYGZhy{}\PYGZhy{}\PYGZhy{}\PYGZhy{}\PYGZhy{}\PYGZhy{}\PYGZhy{}\PYGZhy{}\PYGZhy{}\PYGZhy{}\PYGZhy{}\PYGZhy{}\PYGZhy{}\PYGZhy{}\PYGZhy{}\PYGZhy{}\PYGZhy{}\PYGZhy{}\PYGZhy{}\PYGZhy{}\PYGZhy{}\PYGZhy{}\PYGZhy{}\PYGZhy{}\PYGZhy{}\PYGZhy{}\PYGZhy{}\PYGZhy{}\PYGZhy{}\PYGZhy{}\PYGZhy{}\PYGZhy{}\PYGZhy{}\PYGZhy{}\PYGZhy{}\PYGZhy{}\PYGZhy{}\PYGZhy{}\PYGZhy{}\PYGZhy{}\PYGZhy{}\PYGZhy{}\PYGZhy{}\PYGZhy{}\PYGZhy{}}
\PYG{c+cm}{   proc 0 writes to restart file}
\PYG{c+cm}{\PYGZhy{}\PYGZhy{}\PYGZhy{}\PYGZhy{}\PYGZhy{}\PYGZhy{}\PYGZhy{}\PYGZhy{}\PYGZhy{}\PYGZhy{}\PYGZhy{}\PYGZhy{}\PYGZhy{}\PYGZhy{}\PYGZhy{}\PYGZhy{}\PYGZhy{}\PYGZhy{}\PYGZhy{}\PYGZhy{}\PYGZhy{}\PYGZhy{}\PYGZhy{}\PYGZhy{}\PYGZhy{}\PYGZhy{}\PYGZhy{}\PYGZhy{}\PYGZhy{}\PYGZhy{}\PYGZhy{}\PYGZhy{}\PYGZhy{}\PYGZhy{}\PYGZhy{}\PYGZhy{}\PYGZhy{}\PYGZhy{}\PYGZhy{}\PYGZhy{}\PYGZhy{}\PYGZhy{}\PYGZhy{}\PYGZhy{}\PYGZhy{}\PYGZhy{}\PYGZhy{}\PYGZhy{}\PYGZhy{}\PYGZhy{}\PYGZhy{}\PYGZhy{}\PYGZhy{}\PYGZhy{}\PYGZhy{}\PYGZhy{}\PYGZhy{}\PYGZhy{}\PYGZhy{}\PYGZhy{}\PYGZhy{}\PYGZhy{}\PYGZhy{}\PYGZhy{}\PYGZhy{}\PYGZhy{}\PYGZhy{}\PYGZhy{}\PYGZhy{}\PYGZhy{}\PYGZhy{}\PYGZhy{}\PYGZhy{} */}

\PYG{k+kt}{void}\PYG{+w}{ }\PYG{n+nf}{PairBornGauss::write\PYGZus{}restart}\PYG{p}{(}\PYG{k+kt}{FILE}\PYG{+w}{ }\PYG{o}{*}\PYG{n}{fp}\PYG{p}{)}
\PYG{p}{\PYGZob{}}
\PYG{+w}{  }\PYG{n}{write\PYGZus{}restart\PYGZus{}settings}\PYG{p}{(}\PYG{n}{fp}\PYG{p}{)}\PYG{p}{;}

\PYG{+w}{  }\PYG{k+kt}{int}\PYG{+w}{ }\PYG{n}{i}\PYG{p}{,}\PYG{+w}{ }\PYG{n}{j}\PYG{p}{;}
\PYG{+w}{  }\PYG{k}{for}\PYG{+w}{ }\PYG{p}{(}\PYG{n}{i}\PYG{+w}{ }\PYG{o}{=}\PYG{+w}{ }\PYG{l+m+mi}{1}\PYG{p}{;}\PYG{+w}{ }\PYG{n}{i}\PYG{+w}{ }\PYG{o}{\PYGZlt{}}\PYG{o}{=}\PYG{+w}{ }\PYG{n}{atom}\PYG{o}{\PYGZhy{}}\PYG{o}{\PYGZgt{}}\PYG{n}{ntypes}\PYG{p}{;}\PYG{+w}{ }\PYG{n}{i}\PYG{o}{+}\PYG{o}{+}\PYG{p}{)}\PYG{+w}{ }\PYG{p}{\PYGZob{}}
\PYG{+w}{    }\PYG{k}{for}\PYG{+w}{ }\PYG{p}{(}\PYG{n}{j}\PYG{+w}{ }\PYG{o}{=}\PYG{+w}{ }\PYG{n}{i}\PYG{p}{;}\PYG{+w}{ }\PYG{n}{j}\PYG{+w}{ }\PYG{o}{\PYGZlt{}}\PYG{o}{=}\PYG{+w}{ }\PYG{n}{atom}\PYG{o}{\PYGZhy{}}\PYG{o}{\PYGZgt{}}\PYG{n}{ntypes}\PYG{p}{;}\PYG{+w}{ }\PYG{n}{j}\PYG{o}{+}\PYG{o}{+}\PYG{p}{)}\PYG{+w}{ }\PYG{p}{\PYGZob{}}
\PYG{+w}{      }\PYG{n}{fwrite}\PYG{p}{(}\PYG{o}{\PYGZam{}}\PYG{n}{setflag}\PYG{p}{[}\PYG{n}{i}\PYG{p}{]}\PYG{p}{[}\PYG{n}{j}\PYG{p}{]}\PYG{p}{,}\PYG{+w}{ }\PYG{k}{sizeof}\PYG{p}{(}\PYG{k+kt}{int}\PYG{p}{)}\PYG{p}{,}\PYG{+w}{ }\PYG{l+m+mi}{1}\PYG{p}{,}\PYG{+w}{ }\PYG{n}{fp}\PYG{p}{)}\PYG{p}{;}
\PYG{+w}{      }\PYG{k}{if}\PYG{+w}{ }\PYG{p}{(}\PYG{n}{setflag}\PYG{p}{[}\PYG{n}{i}\PYG{p}{]}\PYG{p}{[}\PYG{n}{j}\PYG{p}{]}\PYG{p}{)}\PYG{+w}{ }\PYG{p}{\PYGZob{}}
\PYG{+w}{        }\PYG{n}{fwrite}\PYG{p}{(}\PYG{o}{\PYGZam{}}\PYG{n}{biga0}\PYG{p}{[}\PYG{n}{i}\PYG{p}{]}\PYG{p}{[}\PYG{n}{j}\PYG{p}{]}\PYG{p}{,}\PYG{+w}{ }\PYG{k}{sizeof}\PYG{p}{(}\PYG{k+kt}{double}\PYG{p}{)}\PYG{p}{,}\PYG{+w}{ }\PYG{l+m+mi}{1}\PYG{p}{,}\PYG{+w}{ }\PYG{n}{fp}\PYG{p}{)}\PYG{p}{;}
\PYG{+w}{        }\PYG{n}{fwrite}\PYG{p}{(}\PYG{o}{\PYGZam{}}\PYG{n}{alpha}\PYG{p}{[}\PYG{n}{i}\PYG{p}{]}\PYG{p}{[}\PYG{n}{j}\PYG{p}{]}\PYG{p}{,}\PYG{+w}{ }\PYG{k}{sizeof}\PYG{p}{(}\PYG{k+kt}{double}\PYG{p}{)}\PYG{p}{,}\PYG{+w}{ }\PYG{l+m+mi}{1}\PYG{p}{,}\PYG{+w}{ }\PYG{n}{fp}\PYG{p}{)}\PYG{p}{;}
\PYG{+w}{        }\PYG{n}{fwrite}\PYG{p}{(}\PYG{o}{\PYGZam{}}\PYG{n}{biga1}\PYG{p}{[}\PYG{n}{i}\PYG{p}{]}\PYG{p}{[}\PYG{n}{j}\PYG{p}{]}\PYG{p}{,}\PYG{+w}{ }\PYG{k}{sizeof}\PYG{p}{(}\PYG{k+kt}{double}\PYG{p}{)}\PYG{p}{,}\PYG{+w}{ }\PYG{l+m+mi}{1}\PYG{p}{,}\PYG{+w}{ }\PYG{n}{fp}\PYG{p}{)}\PYG{p}{;}
\PYG{+w}{        }\PYG{n}{fwrite}\PYG{p}{(}\PYG{o}{\PYGZam{}}\PYG{n}{beta}\PYG{p}{[}\PYG{n}{i}\PYG{p}{]}\PYG{p}{[}\PYG{n}{j}\PYG{p}{]}\PYG{p}{,}\PYG{+w}{ }\PYG{k}{sizeof}\PYG{p}{(}\PYG{k+kt}{double}\PYG{p}{)}\PYG{p}{,}\PYG{+w}{ }\PYG{l+m+mi}{1}\PYG{p}{,}\PYG{+w}{ }\PYG{n}{fp}\PYG{p}{)}\PYG{p}{;}
\PYG{+w}{        }\PYG{n}{fwrite}\PYG{p}{(}\PYG{o}{\PYGZam{}}\PYG{n}{r0}\PYG{p}{[}\PYG{n}{i}\PYG{p}{]}\PYG{p}{[}\PYG{n}{j}\PYG{p}{]}\PYG{p}{,}\PYG{+w}{ }\PYG{k}{sizeof}\PYG{p}{(}\PYG{k+kt}{double}\PYG{p}{)}\PYG{p}{,}\PYG{+w}{ }\PYG{l+m+mi}{1}\PYG{p}{,}\PYG{+w}{ }\PYG{n}{fp}\PYG{p}{)}\PYG{p}{;}
\PYG{+w}{        }\PYG{n}{fwrite}\PYG{p}{(}\PYG{o}{\PYGZam{}}\PYG{n}{cut}\PYG{p}{[}\PYG{n}{i}\PYG{p}{]}\PYG{p}{[}\PYG{n}{j}\PYG{p}{]}\PYG{p}{,}\PYG{+w}{ }\PYG{k}{sizeof}\PYG{p}{(}\PYG{k+kt}{double}\PYG{p}{)}\PYG{p}{,}\PYG{+w}{ }\PYG{l+m+mi}{1}\PYG{p}{,}\PYG{+w}{ }\PYG{n}{fp}\PYG{p}{)}\PYG{p}{;}
\PYG{+w}{      }\PYG{p}{\PYGZcb{}}
\PYG{+w}{    }\PYG{p}{\PYGZcb{}}
\PYG{+w}{  }\PYG{p}{\PYGZcb{}}
\PYG{p}{\PYGZcb{}}

\PYG{c+cm}{/* \PYGZhy{}\PYGZhy{}\PYGZhy{}\PYGZhy{}\PYGZhy{}\PYGZhy{}\PYGZhy{}\PYGZhy{}\PYGZhy{}\PYGZhy{}\PYGZhy{}\PYGZhy{}\PYGZhy{}\PYGZhy{}\PYGZhy{}\PYGZhy{}\PYGZhy{}\PYGZhy{}\PYGZhy{}\PYGZhy{}\PYGZhy{}\PYGZhy{}\PYGZhy{}\PYGZhy{}\PYGZhy{}\PYGZhy{}\PYGZhy{}\PYGZhy{}\PYGZhy{}\PYGZhy{}\PYGZhy{}\PYGZhy{}\PYGZhy{}\PYGZhy{}\PYGZhy{}\PYGZhy{}\PYGZhy{}\PYGZhy{}\PYGZhy{}\PYGZhy{}\PYGZhy{}\PYGZhy{}\PYGZhy{}\PYGZhy{}\PYGZhy{}\PYGZhy{}\PYGZhy{}\PYGZhy{}\PYGZhy{}\PYGZhy{}\PYGZhy{}\PYGZhy{}\PYGZhy{}\PYGZhy{}\PYGZhy{}\PYGZhy{}\PYGZhy{}\PYGZhy{}\PYGZhy{}\PYGZhy{}\PYGZhy{}\PYGZhy{}\PYGZhy{}\PYGZhy{}\PYGZhy{}\PYGZhy{}\PYGZhy{}\PYGZhy{}\PYGZhy{}\PYGZhy{}}
\PYG{c+cm}{   proc 0 writes to restart file}
\PYG{c+cm}{\PYGZhy{}\PYGZhy{}\PYGZhy{}\PYGZhy{}\PYGZhy{}\PYGZhy{}\PYGZhy{}\PYGZhy{}\PYGZhy{}\PYGZhy{}\PYGZhy{}\PYGZhy{}\PYGZhy{}\PYGZhy{}\PYGZhy{}\PYGZhy{}\PYGZhy{}\PYGZhy{}\PYGZhy{}\PYGZhy{}\PYGZhy{}\PYGZhy{}\PYGZhy{}\PYGZhy{}\PYGZhy{}\PYGZhy{}\PYGZhy{}\PYGZhy{}\PYGZhy{}\PYGZhy{}\PYGZhy{}\PYGZhy{}\PYGZhy{}\PYGZhy{}\PYGZhy{}\PYGZhy{}\PYGZhy{}\PYGZhy{}\PYGZhy{}\PYGZhy{}\PYGZhy{}\PYGZhy{}\PYGZhy{}\PYGZhy{}\PYGZhy{}\PYGZhy{}\PYGZhy{}\PYGZhy{}\PYGZhy{}\PYGZhy{}\PYGZhy{}\PYGZhy{}\PYGZhy{}\PYGZhy{}\PYGZhy{}\PYGZhy{}\PYGZhy{}\PYGZhy{}\PYGZhy{}\PYGZhy{}\PYGZhy{}\PYGZhy{}\PYGZhy{}\PYGZhy{}\PYGZhy{}\PYGZhy{}\PYGZhy{}\PYGZhy{}\PYGZhy{}\PYGZhy{}\PYGZhy{}\PYGZhy{}\PYGZhy{} */}

\PYG{k+kt}{void}\PYG{+w}{ }\PYG{n+nf}{PairBornGauss::write\PYGZus{}restart\PYGZus{}settings}\PYG{p}{(}\PYG{k+kt}{FILE}\PYG{+w}{ }\PYG{o}{*}\PYG{n}{fp}\PYG{p}{)}
\PYG{p}{\PYGZob{}}
\PYG{+w}{  }\PYG{n}{fwrite}\PYG{p}{(}\PYG{o}{\PYGZam{}}\PYG{n}{cut\PYGZus{}global}\PYG{p}{,}\PYG{+w}{ }\PYG{k}{sizeof}\PYG{p}{(}\PYG{k+kt}{double}\PYG{p}{)}\PYG{p}{,}\PYG{+w}{ }\PYG{l+m+mi}{1}\PYG{p}{,}\PYG{+w}{ }\PYG{n}{fp}\PYG{p}{)}\PYG{p}{;}
\PYG{+w}{  }\PYG{n}{fwrite}\PYG{p}{(}\PYG{o}{\PYGZam{}}\PYG{n}{offset\PYGZus{}flag}\PYG{p}{,}\PYG{+w}{ }\PYG{k}{sizeof}\PYG{p}{(}\PYG{k+kt}{int}\PYG{p}{)}\PYG{p}{,}\PYG{+w}{ }\PYG{l+m+mi}{1}\PYG{p}{,}\PYG{+w}{ }\PYG{n}{fp}\PYG{p}{)}\PYG{p}{;}
\PYG{+w}{  }\PYG{n}{fwrite}\PYG{p}{(}\PYG{o}{\PYGZam{}}\PYG{n}{mix\PYGZus{}flag}\PYG{p}{,}\PYG{+w}{ }\PYG{k}{sizeof}\PYG{p}{(}\PYG{k+kt}{int}\PYG{p}{)}\PYG{p}{,}\PYG{+w}{ }\PYG{l+m+mi}{1}\PYG{p}{,}\PYG{+w}{ }\PYG{n}{fp}\PYG{p}{)}\PYG{p}{;}
\PYG{p}{\PYGZcb{}}
\end{sphinxVerbatim}

\sphinxAtStartPar
Similarly, on reading, only MPI processor rank 0 has opened the restart
file and will read the data.  The data is then distributed across all
parallel processes using calls to \sphinxcode{\sphinxupquote{MPI\_Bcast()}}.  Before reading atom
type specific data, the corresponding storage needs to be allocated.
Order and number or storage size of items read must be exactly the same
as when writing, or else the data will be read incorrectly.

\sphinxAtStartPar
Reading uses the {\hyperref[\detokenize{Developer_utils:_CPPv4N9LAMMPS_NS5utils6sfreadEPKciPv6size_t6size_tP4FILEPKcP5Error}]{\sphinxcrossref{\sphinxcode{\sphinxupquote{utils::sfread}}}}}
utility function to detect read errors and short reads, so that LAMMPS
can abort if that happens, e.g. when the restart file is corrupted.

\begin{sphinxVerbatim}[commandchars=\\\{\}]
\PYG{c+cm}{/* \PYGZhy{}\PYGZhy{}\PYGZhy{}\PYGZhy{}\PYGZhy{}\PYGZhy{}\PYGZhy{}\PYGZhy{}\PYGZhy{}\PYGZhy{}\PYGZhy{}\PYGZhy{}\PYGZhy{}\PYGZhy{}\PYGZhy{}\PYGZhy{}\PYGZhy{}\PYGZhy{}\PYGZhy{}\PYGZhy{}\PYGZhy{}\PYGZhy{}\PYGZhy{}\PYGZhy{}\PYGZhy{}\PYGZhy{}\PYGZhy{}\PYGZhy{}\PYGZhy{}\PYGZhy{}\PYGZhy{}\PYGZhy{}\PYGZhy{}\PYGZhy{}\PYGZhy{}\PYGZhy{}\PYGZhy{}\PYGZhy{}\PYGZhy{}\PYGZhy{}\PYGZhy{}\PYGZhy{}\PYGZhy{}\PYGZhy{}\PYGZhy{}\PYGZhy{}\PYGZhy{}\PYGZhy{}\PYGZhy{}\PYGZhy{}\PYGZhy{}\PYGZhy{}\PYGZhy{}\PYGZhy{}\PYGZhy{}\PYGZhy{}\PYGZhy{}\PYGZhy{}\PYGZhy{}\PYGZhy{}\PYGZhy{}\PYGZhy{}\PYGZhy{}\PYGZhy{}\PYGZhy{}\PYGZhy{}\PYGZhy{}\PYGZhy{}\PYGZhy{}\PYGZhy{}}
\PYG{c+cm}{   proc 0 reads from restart file, bcasts}
\PYG{c+cm}{\PYGZhy{}\PYGZhy{}\PYGZhy{}\PYGZhy{}\PYGZhy{}\PYGZhy{}\PYGZhy{}\PYGZhy{}\PYGZhy{}\PYGZhy{}\PYGZhy{}\PYGZhy{}\PYGZhy{}\PYGZhy{}\PYGZhy{}\PYGZhy{}\PYGZhy{}\PYGZhy{}\PYGZhy{}\PYGZhy{}\PYGZhy{}\PYGZhy{}\PYGZhy{}\PYGZhy{}\PYGZhy{}\PYGZhy{}\PYGZhy{}\PYGZhy{}\PYGZhy{}\PYGZhy{}\PYGZhy{}\PYGZhy{}\PYGZhy{}\PYGZhy{}\PYGZhy{}\PYGZhy{}\PYGZhy{}\PYGZhy{}\PYGZhy{}\PYGZhy{}\PYGZhy{}\PYGZhy{}\PYGZhy{}\PYGZhy{}\PYGZhy{}\PYGZhy{}\PYGZhy{}\PYGZhy{}\PYGZhy{}\PYGZhy{}\PYGZhy{}\PYGZhy{}\PYGZhy{}\PYGZhy{}\PYGZhy{}\PYGZhy{}\PYGZhy{}\PYGZhy{}\PYGZhy{}\PYGZhy{}\PYGZhy{}\PYGZhy{}\PYGZhy{}\PYGZhy{}\PYGZhy{}\PYGZhy{}\PYGZhy{}\PYGZhy{}\PYGZhy{}\PYGZhy{}\PYGZhy{}\PYGZhy{}\PYGZhy{} */}

\PYG{k+kt}{void}\PYG{+w}{ }\PYG{n+nf}{PairBornGauss::read\PYGZus{}restart}\PYG{p}{(}\PYG{k+kt}{FILE}\PYG{+w}{ }\PYG{o}{*}\PYG{n}{fp}\PYG{p}{)}
\PYG{p}{\PYGZob{}}
\PYG{+w}{  }\PYG{n}{read\PYGZus{}restart\PYGZus{}settings}\PYG{p}{(}\PYG{n}{fp}\PYG{p}{)}\PYG{p}{;}

\PYG{+w}{  }\PYG{n}{allocate}\PYG{p}{(}\PYG{p}{)}\PYG{p}{;}

\PYG{+w}{  }\PYG{k+kt}{int}\PYG{+w}{ }\PYG{n}{i}\PYG{p}{,}\PYG{+w}{ }\PYG{n}{j}\PYG{p}{;}
\PYG{+w}{  }\PYG{k+kt}{int}\PYG{+w}{ }\PYG{n}{me}\PYG{+w}{ }\PYG{o}{=}\PYG{+w}{ }\PYG{n}{comm}\PYG{o}{\PYGZhy{}}\PYG{o}{\PYGZgt{}}\PYG{n}{me}\PYG{p}{;}
\PYG{+w}{  }\PYG{k}{for}\PYG{+w}{ }\PYG{p}{(}\PYG{n}{i}\PYG{+w}{ }\PYG{o}{=}\PYG{+w}{ }\PYG{l+m+mi}{1}\PYG{p}{;}\PYG{+w}{ }\PYG{n}{i}\PYG{+w}{ }\PYG{o}{\PYGZlt{}}\PYG{o}{=}\PYG{+w}{ }\PYG{n}{atom}\PYG{o}{\PYGZhy{}}\PYG{o}{\PYGZgt{}}\PYG{n}{ntypes}\PYG{p}{;}\PYG{+w}{ }\PYG{n}{i}\PYG{o}{+}\PYG{o}{+}\PYG{p}{)}\PYG{+w}{ }\PYG{p}{\PYGZob{}}
\PYG{+w}{    }\PYG{k}{for}\PYG{+w}{ }\PYG{p}{(}\PYG{n}{j}\PYG{+w}{ }\PYG{o}{=}\PYG{+w}{ }\PYG{n}{i}\PYG{p}{;}\PYG{+w}{ }\PYG{n}{j}\PYG{+w}{ }\PYG{o}{\PYGZlt{}}\PYG{o}{=}\PYG{+w}{ }\PYG{n}{atom}\PYG{o}{\PYGZhy{}}\PYG{o}{\PYGZgt{}}\PYG{n}{ntypes}\PYG{p}{;}\PYG{+w}{ }\PYG{n}{j}\PYG{o}{+}\PYG{o}{+}\PYG{p}{)}\PYG{+w}{ }\PYG{p}{\PYGZob{}}
\PYG{+w}{      }\PYG{k}{if}\PYG{+w}{ }\PYG{p}{(}\PYG{n}{me}\PYG{+w}{ }\PYG{o}{=}\PYG{o}{=}\PYG{+w}{ }\PYG{l+m+mi}{0}\PYG{p}{)}\PYG{+w}{ }\PYG{n}{utils}\PYG{o}{:}\PYG{o}{:}\PYG{n}{sfread}\PYG{p}{(}\PYG{n}{FLERR}\PYG{p}{,}\PYG{+w}{ }\PYG{o}{\PYGZam{}}\PYG{n}{setflag}\PYG{p}{[}\PYG{n}{i}\PYG{p}{]}\PYG{p}{[}\PYG{n}{j}\PYG{p}{]}\PYG{p}{,}\PYG{+w}{ }\PYG{k}{sizeof}\PYG{p}{(}\PYG{k+kt}{int}\PYG{p}{)}\PYG{p}{,}\PYG{+w}{ }\PYG{l+m+mi}{1}\PYG{p}{,}\PYG{+w}{ }\PYG{n}{fp}\PYG{p}{,}\PYG{+w}{ }\PYG{n+nb}{nullptr}\PYG{p}{,}\PYG{+w}{ }\PYG{n}{error}\PYG{p}{)}\PYG{p}{;}
\PYG{+w}{      }\PYG{n}{MPI\PYGZus{}Bcast}\PYG{p}{(}\PYG{o}{\PYGZam{}}\PYG{n}{setflag}\PYG{p}{[}\PYG{n}{i}\PYG{p}{]}\PYG{p}{[}\PYG{n}{j}\PYG{p}{]}\PYG{p}{,}\PYG{+w}{ }\PYG{l+m+mi}{1}\PYG{p}{,}\PYG{+w}{ }\PYG{n}{MPI\PYGZus{}INT}\PYG{p}{,}\PYG{+w}{ }\PYG{l+m+mi}{0}\PYG{p}{,}\PYG{+w}{ }\PYG{n}{world}\PYG{p}{)}\PYG{p}{;}
\PYG{+w}{      }\PYG{k}{if}\PYG{+w}{ }\PYG{p}{(}\PYG{n}{setflag}\PYG{p}{[}\PYG{n}{i}\PYG{p}{]}\PYG{p}{[}\PYG{n}{j}\PYG{p}{]}\PYG{p}{)}\PYG{+w}{ }\PYG{p}{\PYGZob{}}
\PYG{+w}{        }\PYG{k}{if}\PYG{+w}{ }\PYG{p}{(}\PYG{n}{me}\PYG{+w}{ }\PYG{o}{=}\PYG{o}{=}\PYG{+w}{ }\PYG{l+m+mi}{0}\PYG{p}{)}\PYG{+w}{ }\PYG{p}{\PYGZob{}}
\PYG{+w}{          }\PYG{n}{utils}\PYG{o}{:}\PYG{o}{:}\PYG{n}{sfread}\PYG{p}{(}\PYG{n}{FLERR}\PYG{p}{,}\PYG{+w}{ }\PYG{o}{\PYGZam{}}\PYG{n}{biga0}\PYG{p}{[}\PYG{n}{i}\PYG{p}{]}\PYG{p}{[}\PYG{n}{j}\PYG{p}{]}\PYG{p}{,}\PYG{+w}{ }\PYG{k}{sizeof}\PYG{p}{(}\PYG{k+kt}{double}\PYG{p}{)}\PYG{p}{,}\PYG{+w}{ }\PYG{l+m+mi}{1}\PYG{p}{,}\PYG{+w}{ }\PYG{n}{fp}\PYG{p}{,}\PYG{+w}{ }\PYG{n+nb}{nullptr}\PYG{p}{,}\PYG{+w}{ }\PYG{n}{error}\PYG{p}{)}\PYG{p}{;}
\PYG{+w}{          }\PYG{n}{utils}\PYG{o}{:}\PYG{o}{:}\PYG{n}{sfread}\PYG{p}{(}\PYG{n}{FLERR}\PYG{p}{,}\PYG{+w}{ }\PYG{o}{\PYGZam{}}\PYG{n}{alpha}\PYG{p}{[}\PYG{n}{i}\PYG{p}{]}\PYG{p}{[}\PYG{n}{j}\PYG{p}{]}\PYG{p}{,}\PYG{+w}{ }\PYG{k}{sizeof}\PYG{p}{(}\PYG{k+kt}{double}\PYG{p}{)}\PYG{p}{,}\PYG{+w}{ }\PYG{l+m+mi}{1}\PYG{p}{,}\PYG{+w}{ }\PYG{n}{fp}\PYG{p}{,}\PYG{+w}{ }\PYG{n+nb}{nullptr}\PYG{p}{,}\PYG{+w}{ }\PYG{n}{error}\PYG{p}{)}\PYG{p}{;}
\PYG{+w}{          }\PYG{n}{utils}\PYG{o}{:}\PYG{o}{:}\PYG{n}{sfread}\PYG{p}{(}\PYG{n}{FLERR}\PYG{p}{,}\PYG{+w}{ }\PYG{o}{\PYGZam{}}\PYG{n}{biga1}\PYG{p}{[}\PYG{n}{i}\PYG{p}{]}\PYG{p}{[}\PYG{n}{j}\PYG{p}{]}\PYG{p}{,}\PYG{+w}{ }\PYG{k}{sizeof}\PYG{p}{(}\PYG{k+kt}{double}\PYG{p}{)}\PYG{p}{,}\PYG{+w}{ }\PYG{l+m+mi}{1}\PYG{p}{,}\PYG{+w}{ }\PYG{n}{fp}\PYG{p}{,}\PYG{+w}{ }\PYG{n+nb}{nullptr}\PYG{p}{,}\PYG{+w}{ }\PYG{n}{error}\PYG{p}{)}\PYG{p}{;}
\PYG{+w}{          }\PYG{n}{utils}\PYG{o}{:}\PYG{o}{:}\PYG{n}{sfread}\PYG{p}{(}\PYG{n}{FLERR}\PYG{p}{,}\PYG{+w}{ }\PYG{o}{\PYGZam{}}\PYG{n}{beta}\PYG{p}{[}\PYG{n}{i}\PYG{p}{]}\PYG{p}{[}\PYG{n}{j}\PYG{p}{]}\PYG{p}{,}\PYG{+w}{ }\PYG{k}{sizeof}\PYG{p}{(}\PYG{k+kt}{double}\PYG{p}{)}\PYG{p}{,}\PYG{+w}{ }\PYG{l+m+mi}{1}\PYG{p}{,}\PYG{+w}{ }\PYG{n}{fp}\PYG{p}{,}\PYG{+w}{ }\PYG{n+nb}{nullptr}\PYG{p}{,}\PYG{+w}{ }\PYG{n}{error}\PYG{p}{)}\PYG{p}{;}
\PYG{+w}{          }\PYG{n}{utils}\PYG{o}{:}\PYG{o}{:}\PYG{n}{sfread}\PYG{p}{(}\PYG{n}{FLERR}\PYG{p}{,}\PYG{+w}{ }\PYG{o}{\PYGZam{}}\PYG{n}{r0}\PYG{p}{[}\PYG{n}{i}\PYG{p}{]}\PYG{p}{[}\PYG{n}{j}\PYG{p}{]}\PYG{p}{,}\PYG{+w}{ }\PYG{k}{sizeof}\PYG{p}{(}\PYG{k+kt}{double}\PYG{p}{)}\PYG{p}{,}\PYG{+w}{ }\PYG{l+m+mi}{1}\PYG{p}{,}\PYG{+w}{ }\PYG{n}{fp}\PYG{p}{,}\PYG{+w}{ }\PYG{n+nb}{nullptr}\PYG{p}{,}\PYG{+w}{ }\PYG{n}{error}\PYG{p}{)}\PYG{p}{;}
\PYG{+w}{          }\PYG{n}{utils}\PYG{o}{:}\PYG{o}{:}\PYG{n}{sfread}\PYG{p}{(}\PYG{n}{FLERR}\PYG{p}{,}\PYG{+w}{ }\PYG{o}{\PYGZam{}}\PYG{n}{cut}\PYG{p}{[}\PYG{n}{i}\PYG{p}{]}\PYG{p}{[}\PYG{n}{j}\PYG{p}{]}\PYG{p}{,}\PYG{+w}{ }\PYG{k}{sizeof}\PYG{p}{(}\PYG{k+kt}{double}\PYG{p}{)}\PYG{p}{,}\PYG{+w}{ }\PYG{l+m+mi}{1}\PYG{p}{,}\PYG{+w}{ }\PYG{n}{fp}\PYG{p}{,}\PYG{+w}{ }\PYG{n+nb}{nullptr}\PYG{p}{,}\PYG{+w}{ }\PYG{n}{error}\PYG{p}{)}\PYG{p}{;}
\PYG{+w}{        }\PYG{p}{\PYGZcb{}}
\PYG{+w}{        }\PYG{n}{MPI\PYGZus{}Bcast}\PYG{p}{(}\PYG{o}{\PYGZam{}}\PYG{n}{biga0}\PYG{p}{[}\PYG{n}{i}\PYG{p}{]}\PYG{p}{[}\PYG{n}{j}\PYG{p}{]}\PYG{p}{,}\PYG{+w}{ }\PYG{l+m+mi}{1}\PYG{p}{,}\PYG{+w}{ }\PYG{n}{MPI\PYGZus{}DOUBLE}\PYG{p}{,}\PYG{+w}{ }\PYG{l+m+mi}{0}\PYG{p}{,}\PYG{+w}{ }\PYG{n}{world}\PYG{p}{)}\PYG{p}{;}
\PYG{+w}{        }\PYG{n}{MPI\PYGZus{}Bcast}\PYG{p}{(}\PYG{o}{\PYGZam{}}\PYG{n}{alpha}\PYG{p}{[}\PYG{n}{i}\PYG{p}{]}\PYG{p}{[}\PYG{n}{j}\PYG{p}{]}\PYG{p}{,}\PYG{+w}{ }\PYG{l+m+mi}{1}\PYG{p}{,}\PYG{+w}{ }\PYG{n}{MPI\PYGZus{}DOUBLE}\PYG{p}{,}\PYG{+w}{ }\PYG{l+m+mi}{0}\PYG{p}{,}\PYG{+w}{ }\PYG{n}{world}\PYG{p}{)}\PYG{p}{;}
\PYG{+w}{        }\PYG{n}{MPI\PYGZus{}Bcast}\PYG{p}{(}\PYG{o}{\PYGZam{}}\PYG{n}{biga1}\PYG{p}{[}\PYG{n}{i}\PYG{p}{]}\PYG{p}{[}\PYG{n}{j}\PYG{p}{]}\PYG{p}{,}\PYG{+w}{ }\PYG{l+m+mi}{1}\PYG{p}{,}\PYG{+w}{ }\PYG{n}{MPI\PYGZus{}DOUBLE}\PYG{p}{,}\PYG{+w}{ }\PYG{l+m+mi}{0}\PYG{p}{,}\PYG{+w}{ }\PYG{n}{world}\PYG{p}{)}\PYG{p}{;}
\PYG{+w}{        }\PYG{n}{MPI\PYGZus{}Bcast}\PYG{p}{(}\PYG{o}{\PYGZam{}}\PYG{n}{beta}\PYG{p}{[}\PYG{n}{i}\PYG{p}{]}\PYG{p}{[}\PYG{n}{j}\PYG{p}{]}\PYG{p}{,}\PYG{+w}{ }\PYG{l+m+mi}{1}\PYG{p}{,}\PYG{+w}{ }\PYG{n}{MPI\PYGZus{}DOUBLE}\PYG{p}{,}\PYG{+w}{ }\PYG{l+m+mi}{0}\PYG{p}{,}\PYG{+w}{ }\PYG{n}{world}\PYG{p}{)}\PYG{p}{;}
\PYG{+w}{        }\PYG{n}{MPI\PYGZus{}Bcast}\PYG{p}{(}\PYG{o}{\PYGZam{}}\PYG{n}{r0}\PYG{p}{[}\PYG{n}{i}\PYG{p}{]}\PYG{p}{[}\PYG{n}{j}\PYG{p}{]}\PYG{p}{,}\PYG{+w}{ }\PYG{l+m+mi}{1}\PYG{p}{,}\PYG{+w}{ }\PYG{n}{MPI\PYGZus{}DOUBLE}\PYG{p}{,}\PYG{+w}{ }\PYG{l+m+mi}{0}\PYG{p}{,}\PYG{+w}{ }\PYG{n}{world}\PYG{p}{)}\PYG{p}{;}
\PYG{+w}{        }\PYG{n}{MPI\PYGZus{}Bcast}\PYG{p}{(}\PYG{o}{\PYGZam{}}\PYG{n}{cut}\PYG{p}{[}\PYG{n}{i}\PYG{p}{]}\PYG{p}{[}\PYG{n}{j}\PYG{p}{]}\PYG{p}{,}\PYG{+w}{ }\PYG{l+m+mi}{1}\PYG{p}{,}\PYG{+w}{ }\PYG{n}{MPI\PYGZus{}DOUBLE}\PYG{p}{,}\PYG{+w}{ }\PYG{l+m+mi}{0}\PYG{p}{,}\PYG{+w}{ }\PYG{n}{world}\PYG{p}{)}\PYG{p}{;}
\PYG{+w}{      }\PYG{p}{\PYGZcb{}}
\PYG{+w}{    }\PYG{p}{\PYGZcb{}}
\PYG{+w}{  }\PYG{p}{\PYGZcb{}}
\PYG{p}{\PYGZcb{}}

\PYG{c+cm}{/* \PYGZhy{}\PYGZhy{}\PYGZhy{}\PYGZhy{}\PYGZhy{}\PYGZhy{}\PYGZhy{}\PYGZhy{}\PYGZhy{}\PYGZhy{}\PYGZhy{}\PYGZhy{}\PYGZhy{}\PYGZhy{}\PYGZhy{}\PYGZhy{}\PYGZhy{}\PYGZhy{}\PYGZhy{}\PYGZhy{}\PYGZhy{}\PYGZhy{}\PYGZhy{}\PYGZhy{}\PYGZhy{}\PYGZhy{}\PYGZhy{}\PYGZhy{}\PYGZhy{}\PYGZhy{}\PYGZhy{}\PYGZhy{}\PYGZhy{}\PYGZhy{}\PYGZhy{}\PYGZhy{}\PYGZhy{}\PYGZhy{}\PYGZhy{}\PYGZhy{}\PYGZhy{}\PYGZhy{}\PYGZhy{}\PYGZhy{}\PYGZhy{}\PYGZhy{}\PYGZhy{}\PYGZhy{}\PYGZhy{}\PYGZhy{}\PYGZhy{}\PYGZhy{}\PYGZhy{}\PYGZhy{}\PYGZhy{}\PYGZhy{}\PYGZhy{}\PYGZhy{}\PYGZhy{}\PYGZhy{}\PYGZhy{}\PYGZhy{}\PYGZhy{}\PYGZhy{}\PYGZhy{}\PYGZhy{}\PYGZhy{}\PYGZhy{}\PYGZhy{}\PYGZhy{}}
\PYG{c+cm}{   proc 0 reads from restart file, bcasts}
\PYG{c+cm}{\PYGZhy{}\PYGZhy{}\PYGZhy{}\PYGZhy{}\PYGZhy{}\PYGZhy{}\PYGZhy{}\PYGZhy{}\PYGZhy{}\PYGZhy{}\PYGZhy{}\PYGZhy{}\PYGZhy{}\PYGZhy{}\PYGZhy{}\PYGZhy{}\PYGZhy{}\PYGZhy{}\PYGZhy{}\PYGZhy{}\PYGZhy{}\PYGZhy{}\PYGZhy{}\PYGZhy{}\PYGZhy{}\PYGZhy{}\PYGZhy{}\PYGZhy{}\PYGZhy{}\PYGZhy{}\PYGZhy{}\PYGZhy{}\PYGZhy{}\PYGZhy{}\PYGZhy{}\PYGZhy{}\PYGZhy{}\PYGZhy{}\PYGZhy{}\PYGZhy{}\PYGZhy{}\PYGZhy{}\PYGZhy{}\PYGZhy{}\PYGZhy{}\PYGZhy{}\PYGZhy{}\PYGZhy{}\PYGZhy{}\PYGZhy{}\PYGZhy{}\PYGZhy{}\PYGZhy{}\PYGZhy{}\PYGZhy{}\PYGZhy{}\PYGZhy{}\PYGZhy{}\PYGZhy{}\PYGZhy{}\PYGZhy{}\PYGZhy{}\PYGZhy{}\PYGZhy{}\PYGZhy{}\PYGZhy{}\PYGZhy{}\PYGZhy{}\PYGZhy{}\PYGZhy{}\PYGZhy{}\PYGZhy{}\PYGZhy{} */}

\PYG{k+kt}{void}\PYG{+w}{ }\PYG{n+nf}{PairBornGauss::read\PYGZus{}restart\PYGZus{}settings}\PYG{p}{(}\PYG{k+kt}{FILE}\PYG{+w}{ }\PYG{o}{*}\PYG{n}{fp}\PYG{p}{)}
\PYG{p}{\PYGZob{}}
\PYG{+w}{  }\PYG{k}{if}\PYG{+w}{ }\PYG{p}{(}\PYG{n}{comm}\PYG{o}{\PYGZhy{}}\PYG{o}{\PYGZgt{}}\PYG{n}{me}\PYG{+w}{ }\PYG{o}{=}\PYG{o}{=}\PYG{+w}{ }\PYG{l+m+mi}{0}\PYG{p}{)}\PYG{+w}{ }\PYG{p}{\PYGZob{}}
\PYG{+w}{    }\PYG{n}{utils}\PYG{o}{:}\PYG{o}{:}\PYG{n}{sfread}\PYG{p}{(}\PYG{n}{FLERR}\PYG{p}{,}\PYG{+w}{ }\PYG{o}{\PYGZam{}}\PYG{n}{cut\PYGZus{}global}\PYG{p}{,}\PYG{+w}{ }\PYG{k}{sizeof}\PYG{p}{(}\PYG{k+kt}{double}\PYG{p}{)}\PYG{p}{,}\PYG{+w}{ }\PYG{l+m+mi}{1}\PYG{p}{,}\PYG{+w}{ }\PYG{n}{fp}\PYG{p}{,}\PYG{+w}{ }\PYG{n+nb}{nullptr}\PYG{p}{,}\PYG{+w}{ }\PYG{n}{error}\PYG{p}{)}\PYG{p}{;}
\PYG{+w}{    }\PYG{n}{utils}\PYG{o}{:}\PYG{o}{:}\PYG{n}{sfread}\PYG{p}{(}\PYG{n}{FLERR}\PYG{p}{,}\PYG{+w}{ }\PYG{o}{\PYGZam{}}\PYG{n}{offset\PYGZus{}flag}\PYG{p}{,}\PYG{+w}{ }\PYG{k}{sizeof}\PYG{p}{(}\PYG{k+kt}{int}\PYG{p}{)}\PYG{p}{,}\PYG{+w}{ }\PYG{l+m+mi}{1}\PYG{p}{,}\PYG{+w}{ }\PYG{n}{fp}\PYG{p}{,}\PYG{+w}{ }\PYG{n+nb}{nullptr}\PYG{p}{,}\PYG{+w}{ }\PYG{n}{error}\PYG{p}{)}\PYG{p}{;}
\PYG{+w}{    }\PYG{n}{utils}\PYG{o}{:}\PYG{o}{:}\PYG{n}{sfread}\PYG{p}{(}\PYG{n}{FLERR}\PYG{p}{,}\PYG{+w}{ }\PYG{o}{\PYGZam{}}\PYG{n}{mix\PYGZus{}flag}\PYG{p}{,}\PYG{+w}{ }\PYG{k}{sizeof}\PYG{p}{(}\PYG{k+kt}{int}\PYG{p}{)}\PYG{p}{,}\PYG{+w}{ }\PYG{l+m+mi}{1}\PYG{p}{,}\PYG{+w}{ }\PYG{n}{fp}\PYG{p}{,}\PYG{+w}{ }\PYG{n+nb}{nullptr}\PYG{p}{,}\PYG{+w}{ }\PYG{n}{error}\PYG{p}{)}\PYG{p}{;}
\PYG{+w}{  }\PYG{p}{\PYGZcb{}}
\PYG{+w}{  }\PYG{n}{MPI\PYGZus{}Bcast}\PYG{p}{(}\PYG{o}{\PYGZam{}}\PYG{n}{cut\PYGZus{}global}\PYG{p}{,}\PYG{+w}{ }\PYG{l+m+mi}{1}\PYG{p}{,}\PYG{+w}{ }\PYG{n}{MPI\PYGZus{}DOUBLE}\PYG{p}{,}\PYG{+w}{ }\PYG{l+m+mi}{0}\PYG{p}{,}\PYG{+w}{ }\PYG{n}{world}\PYG{p}{)}\PYG{p}{;}
\PYG{+w}{  }\PYG{n}{MPI\PYGZus{}Bcast}\PYG{p}{(}\PYG{o}{\PYGZam{}}\PYG{n}{offset\PYGZus{}flag}\PYG{p}{,}\PYG{+w}{ }\PYG{l+m+mi}{1}\PYG{p}{,}\PYG{+w}{ }\PYG{n}{MPI\PYGZus{}INT}\PYG{p}{,}\PYG{+w}{ }\PYG{l+m+mi}{0}\PYG{p}{,}\PYG{+w}{ }\PYG{n}{world}\PYG{p}{)}\PYG{p}{;}
\PYG{+w}{  }\PYG{n}{MPI\PYGZus{}Bcast}\PYG{p}{(}\PYG{o}{\PYGZam{}}\PYG{n}{mix\PYGZus{}flag}\PYG{p}{,}\PYG{+w}{ }\PYG{l+m+mi}{1}\PYG{p}{,}\PYG{+w}{ }\PYG{n}{MPI\PYGZus{}INT}\PYG{p}{,}\PYG{+w}{ }\PYG{l+m+mi}{0}\PYG{p}{,}\PYG{+w}{ }\PYG{n}{world}\PYG{p}{)}\PYG{p}{;}
\PYG{p}{\PYGZcb{}}
\end{sphinxVerbatim}


\subsubsection{Writing coefficients to data files}
\label{\detokenize{Developer_write_pair:writing-coefficients-to-data-files}}
\sphinxAtStartPar
The \sphinxcode{\sphinxupquote{write\_data()}} and \sphinxcode{\sphinxupquote{write\_data\_all()}} functions are optional and
write out the current state of the {\hyperref[\detokenize{pair_coeff::doc}]{\sphinxcrossref{\DUrole{doc}{pair\_coeff
settings}}}} as “Pair Coeffs” or “PairIJ Coeffs” sections to a
data file when using the {\hyperref[\detokenize{write_data::doc}]{\sphinxcrossref{\DUrole{doc}{write\_data command}}}}.  The
\sphinxcode{\sphinxupquote{write\_data()}} only writes out the diagonal elements of the pair
coefficient matrix, as that is required for the format of the “Pair
Coeffs” section.  It is called when the “pair” option of the
{\hyperref[\detokenize{write_data::doc}]{\sphinxcrossref{\DUrole{doc}{write\_data command}}}} is “ii” (the default).  This is
suitable for force fields where \sphinxstyleemphasis{all} off\sphinxhyphen{}diagonal terms of the pair
coefficient matrix are generated from mixing.  If explicit settings for
off\sphinxhyphen{}diagonal elements were made, LAMMPS will print a warning, as those
would be lost.  To avoid this, the “pair ij” option of {\hyperref[\detokenize{write_data::doc}]{\sphinxcrossref{\DUrole{doc}{write\_data}}}} can be used which will trigger calling the
\sphinxcode{\sphinxupquote{write\_data\_all()}} function instead, which will write out all settings
of the pair coefficient matrix (regardless of whether they were
originally created from mixing or not).

\sphinxAtStartPar
These data file output functions are only useful for true pair\sphinxhyphen{}wise
additive potentials, where the potential parameters can be entered
through \sphinxstyleemphasis{multiple} {\hyperref[\detokenize{pair_coeff::doc}]{\sphinxcrossref{\DUrole{doc}{pair\_coeff commands}}}}.  Pair styles
that require a single “pair\_coeff * *” command line are not compatible
with reading their parameters from data files.  For pair styles like
\sphinxstyleemphasis{born/gauss} that do support writing to data files, the potential
parameters will be read from the data file, if present, and
{\hyperref[\detokenize{pair_coeff::doc}]{\sphinxcrossref{\DUrole{doc}{pair\_coeff commands}}}} may not be needed.

\sphinxAtStartPar
The member variable \sphinxcode{\sphinxupquote{writedata}} should be set to 1 in the constructor,
if these functions are implemented (the default value is 0).

\begin{sphinxVerbatim}[commandchars=\\\{\}]
\PYG{c+cm}{/* \PYGZhy{}\PYGZhy{}\PYGZhy{}\PYGZhy{}\PYGZhy{}\PYGZhy{}\PYGZhy{}\PYGZhy{}\PYGZhy{}\PYGZhy{}\PYGZhy{}\PYGZhy{}\PYGZhy{}\PYGZhy{}\PYGZhy{}\PYGZhy{}\PYGZhy{}\PYGZhy{}\PYGZhy{}\PYGZhy{}\PYGZhy{}\PYGZhy{}\PYGZhy{}\PYGZhy{}\PYGZhy{}\PYGZhy{}\PYGZhy{}\PYGZhy{}\PYGZhy{}\PYGZhy{}\PYGZhy{}\PYGZhy{}\PYGZhy{}\PYGZhy{}\PYGZhy{}\PYGZhy{}\PYGZhy{}\PYGZhy{}\PYGZhy{}\PYGZhy{}\PYGZhy{}\PYGZhy{}\PYGZhy{}\PYGZhy{}\PYGZhy{}\PYGZhy{}\PYGZhy{}\PYGZhy{}\PYGZhy{}\PYGZhy{}\PYGZhy{}\PYGZhy{}\PYGZhy{}\PYGZhy{}\PYGZhy{}\PYGZhy{}\PYGZhy{}\PYGZhy{}\PYGZhy{}\PYGZhy{}\PYGZhy{}\PYGZhy{}\PYGZhy{}\PYGZhy{}\PYGZhy{}\PYGZhy{}\PYGZhy{}\PYGZhy{}\PYGZhy{}\PYGZhy{}}
\PYG{c+cm}{   proc 0 writes to data file}
\PYG{c+cm}{\PYGZhy{}\PYGZhy{}\PYGZhy{}\PYGZhy{}\PYGZhy{}\PYGZhy{}\PYGZhy{}\PYGZhy{}\PYGZhy{}\PYGZhy{}\PYGZhy{}\PYGZhy{}\PYGZhy{}\PYGZhy{}\PYGZhy{}\PYGZhy{}\PYGZhy{}\PYGZhy{}\PYGZhy{}\PYGZhy{}\PYGZhy{}\PYGZhy{}\PYGZhy{}\PYGZhy{}\PYGZhy{}\PYGZhy{}\PYGZhy{}\PYGZhy{}\PYGZhy{}\PYGZhy{}\PYGZhy{}\PYGZhy{}\PYGZhy{}\PYGZhy{}\PYGZhy{}\PYGZhy{}\PYGZhy{}\PYGZhy{}\PYGZhy{}\PYGZhy{}\PYGZhy{}\PYGZhy{}\PYGZhy{}\PYGZhy{}\PYGZhy{}\PYGZhy{}\PYGZhy{}\PYGZhy{}\PYGZhy{}\PYGZhy{}\PYGZhy{}\PYGZhy{}\PYGZhy{}\PYGZhy{}\PYGZhy{}\PYGZhy{}\PYGZhy{}\PYGZhy{}\PYGZhy{}\PYGZhy{}\PYGZhy{}\PYGZhy{}\PYGZhy{}\PYGZhy{}\PYGZhy{}\PYGZhy{}\PYGZhy{}\PYGZhy{}\PYGZhy{}\PYGZhy{}\PYGZhy{}\PYGZhy{}\PYGZhy{} */}

\PYG{k+kt}{void}\PYG{+w}{ }\PYG{n+nf}{PairBornGauss::write\PYGZus{}data}\PYG{p}{(}\PYG{k+kt}{FILE}\PYG{+w}{ }\PYG{o}{*}\PYG{n}{fp}\PYG{p}{)}
\PYG{p}{\PYGZob{}}
\PYG{+w}{  }\PYG{k}{for}\PYG{+w}{ }\PYG{p}{(}\PYG{k+kt}{int}\PYG{+w}{ }\PYG{n}{i}\PYG{+w}{ }\PYG{o}{=}\PYG{+w}{ }\PYG{l+m+mi}{1}\PYG{p}{;}\PYG{+w}{ }\PYG{n}{i}\PYG{+w}{ }\PYG{o}{\PYGZlt{}}\PYG{o}{=}\PYG{+w}{ }\PYG{n}{atom}\PYG{o}{\PYGZhy{}}\PYG{o}{\PYGZgt{}}\PYG{n}{ntypes}\PYG{p}{;}\PYG{+w}{ }\PYG{n}{i}\PYG{o}{+}\PYG{o}{+}\PYG{p}{)}
\PYG{+w}{    }\PYG{n}{fprintf}\PYG{p}{(}\PYG{n}{fp}\PYG{p}{,}\PYG{+w}{ }\PYG{l+s}{\PYGZdq{}}\PYG{l+s}{\PYGZpc{}d \PYGZpc{}g \PYGZpc{}g \PYGZpc{}g \PYGZpc{}g \PYGZpc{}g}\PYG{l+s+se}{\PYGZbs{}n}\PYG{l+s}{\PYGZdq{}}\PYG{p}{,}\PYG{+w}{ }\PYG{n}{i}\PYG{p}{,}\PYG{+w}{ }\PYG{n}{biga0}\PYG{p}{[}\PYG{n}{i}\PYG{p}{]}\PYG{p}{[}\PYG{n}{i}\PYG{p}{]}\PYG{p}{,}\PYG{+w}{ }\PYG{n}{alpha}\PYG{p}{[}\PYG{n}{i}\PYG{p}{]}\PYG{p}{[}\PYG{n}{i}\PYG{p}{]}\PYG{p}{,}\PYG{+w}{ }\PYG{n}{biga1}\PYG{p}{[}\PYG{n}{i}\PYG{p}{]}\PYG{p}{[}\PYG{n}{i}\PYG{p}{]}\PYG{p}{,}\PYG{+w}{ }\PYG{n}{beta}\PYG{p}{[}\PYG{n}{i}\PYG{p}{]}\PYG{p}{[}\PYG{n}{i}\PYG{p}{]}\PYG{p}{,}
\PYG{+w}{            }\PYG{n}{r0}\PYG{p}{[}\PYG{n}{i}\PYG{p}{]}\PYG{p}{[}\PYG{n}{i}\PYG{p}{]}\PYG{p}{)}\PYG{p}{;}
\PYG{p}{\PYGZcb{}}

\PYG{c+cm}{/* \PYGZhy{}\PYGZhy{}\PYGZhy{}\PYGZhy{}\PYGZhy{}\PYGZhy{}\PYGZhy{}\PYGZhy{}\PYGZhy{}\PYGZhy{}\PYGZhy{}\PYGZhy{}\PYGZhy{}\PYGZhy{}\PYGZhy{}\PYGZhy{}\PYGZhy{}\PYGZhy{}\PYGZhy{}\PYGZhy{}\PYGZhy{}\PYGZhy{}\PYGZhy{}\PYGZhy{}\PYGZhy{}\PYGZhy{}\PYGZhy{}\PYGZhy{}\PYGZhy{}\PYGZhy{}\PYGZhy{}\PYGZhy{}\PYGZhy{}\PYGZhy{}\PYGZhy{}\PYGZhy{}\PYGZhy{}\PYGZhy{}\PYGZhy{}\PYGZhy{}\PYGZhy{}\PYGZhy{}\PYGZhy{}\PYGZhy{}\PYGZhy{}\PYGZhy{}\PYGZhy{}\PYGZhy{}\PYGZhy{}\PYGZhy{}\PYGZhy{}\PYGZhy{}\PYGZhy{}\PYGZhy{}\PYGZhy{}\PYGZhy{}\PYGZhy{}\PYGZhy{}\PYGZhy{}\PYGZhy{}\PYGZhy{}\PYGZhy{}\PYGZhy{}\PYGZhy{}\PYGZhy{}\PYGZhy{}\PYGZhy{}\PYGZhy{}\PYGZhy{}\PYGZhy{}}
\PYG{c+cm}{   proc 0 writes all pairs to data file}
\PYG{c+cm}{\PYGZhy{}\PYGZhy{}\PYGZhy{}\PYGZhy{}\PYGZhy{}\PYGZhy{}\PYGZhy{}\PYGZhy{}\PYGZhy{}\PYGZhy{}\PYGZhy{}\PYGZhy{}\PYGZhy{}\PYGZhy{}\PYGZhy{}\PYGZhy{}\PYGZhy{}\PYGZhy{}\PYGZhy{}\PYGZhy{}\PYGZhy{}\PYGZhy{}\PYGZhy{}\PYGZhy{}\PYGZhy{}\PYGZhy{}\PYGZhy{}\PYGZhy{}\PYGZhy{}\PYGZhy{}\PYGZhy{}\PYGZhy{}\PYGZhy{}\PYGZhy{}\PYGZhy{}\PYGZhy{}\PYGZhy{}\PYGZhy{}\PYGZhy{}\PYGZhy{}\PYGZhy{}\PYGZhy{}\PYGZhy{}\PYGZhy{}\PYGZhy{}\PYGZhy{}\PYGZhy{}\PYGZhy{}\PYGZhy{}\PYGZhy{}\PYGZhy{}\PYGZhy{}\PYGZhy{}\PYGZhy{}\PYGZhy{}\PYGZhy{}\PYGZhy{}\PYGZhy{}\PYGZhy{}\PYGZhy{}\PYGZhy{}\PYGZhy{}\PYGZhy{}\PYGZhy{}\PYGZhy{}\PYGZhy{}\PYGZhy{}\PYGZhy{}\PYGZhy{}\PYGZhy{}\PYGZhy{}\PYGZhy{}\PYGZhy{} */}

\PYG{k+kt}{void}\PYG{+w}{ }\PYG{n+nf}{PairBornGauss::write\PYGZus{}data\PYGZus{}all}\PYG{p}{(}\PYG{k+kt}{FILE}\PYG{+w}{ }\PYG{o}{*}\PYG{n}{fp}\PYG{p}{)}
\PYG{p}{\PYGZob{}}
\PYG{+w}{  }\PYG{k}{for}\PYG{+w}{ }\PYG{p}{(}\PYG{k+kt}{int}\PYG{+w}{ }\PYG{n}{i}\PYG{+w}{ }\PYG{o}{=}\PYG{+w}{ }\PYG{l+m+mi}{1}\PYG{p}{;}\PYG{+w}{ }\PYG{n}{i}\PYG{+w}{ }\PYG{o}{\PYGZlt{}}\PYG{o}{=}\PYG{+w}{ }\PYG{n}{atom}\PYG{o}{\PYGZhy{}}\PYG{o}{\PYGZgt{}}\PYG{n}{ntypes}\PYG{p}{;}\PYG{+w}{ }\PYG{n}{i}\PYG{o}{+}\PYG{o}{+}\PYG{p}{)}
\PYG{+w}{    }\PYG{k}{for}\PYG{+w}{ }\PYG{p}{(}\PYG{k+kt}{int}\PYG{+w}{ }\PYG{n}{j}\PYG{+w}{ }\PYG{o}{=}\PYG{+w}{ }\PYG{n}{i}\PYG{p}{;}\PYG{+w}{ }\PYG{n}{j}\PYG{+w}{ }\PYG{o}{\PYGZlt{}}\PYG{o}{=}\PYG{+w}{ }\PYG{n}{atom}\PYG{o}{\PYGZhy{}}\PYG{o}{\PYGZgt{}}\PYG{n}{ntypes}\PYG{p}{;}\PYG{+w}{ }\PYG{n}{j}\PYG{o}{+}\PYG{o}{+}\PYG{p}{)}
\PYG{+w}{      }\PYG{n}{fprintf}\PYG{p}{(}\PYG{n}{fp}\PYG{p}{,}\PYG{+w}{ }\PYG{l+s}{\PYGZdq{}}\PYG{l+s}{\PYGZpc{}d \PYGZpc{}d \PYGZpc{}g \PYGZpc{}g \PYGZpc{}g \PYGZpc{}g \PYGZpc{}g \PYGZpc{}g}\PYG{l+s+se}{\PYGZbs{}n}\PYG{l+s}{\PYGZdq{}}\PYG{p}{,}\PYG{+w}{ }\PYG{n}{i}\PYG{p}{,}\PYG{+w}{ }\PYG{n}{j}\PYG{p}{,}\PYG{+w}{ }\PYG{n}{biga0}\PYG{p}{[}\PYG{n}{i}\PYG{p}{]}\PYG{p}{[}\PYG{n}{j}\PYG{p}{]}\PYG{p}{,}\PYG{+w}{ }\PYG{n}{alpha}\PYG{p}{[}\PYG{n}{i}\PYG{p}{]}\PYG{p}{[}\PYG{n}{j}\PYG{p}{]}\PYG{p}{,}\PYG{+w}{ }\PYG{n}{biga1}\PYG{p}{[}\PYG{n}{i}\PYG{p}{]}\PYG{p}{[}\PYG{n}{j}\PYG{p}{]}\PYG{p}{,}
\PYG{+w}{              }\PYG{n}{beta}\PYG{p}{[}\PYG{n}{i}\PYG{p}{]}\PYG{p}{[}\PYG{n}{j}\PYG{p}{]}\PYG{p}{,}\PYG{+w}{ }\PYG{n}{r0}\PYG{p}{[}\PYG{n}{i}\PYG{p}{]}\PYG{p}{[}\PYG{n}{j}\PYG{p}{]}\PYG{p}{,}\PYG{+w}{ }\PYG{n}{cut}\PYG{p}{[}\PYG{n}{i}\PYG{p}{]}\PYG{p}{[}\PYG{n}{j}\PYG{p}{]}\PYG{p}{)}\PYG{p}{;}
\PYG{p}{\PYGZcb{}}
\end{sphinxVerbatim}


\subsubsection{Give access to internal data}
\label{\detokenize{Developer_write_pair:give-access-to-internal-data}}
\sphinxAtStartPar
The purpose of the \sphinxcode{\sphinxupquote{extract()}} function is to facilitate access to
internal data of the pair style by other parts of LAMMPS.  One possible
application is to use {\hyperref[\detokenize{fix_adapt::doc}]{\sphinxcrossref{\DUrole{doc}{fix adapt}}}} to gradually change
potential parameters during a run.  Here, we implement access to the
pair coefficient matrix parameters.

\begin{sphinxVerbatim}[commandchars=\\\{\}]
\PYG{c+cm}{/* \PYGZhy{}\PYGZhy{}\PYGZhy{}\PYGZhy{}\PYGZhy{}\PYGZhy{}\PYGZhy{}\PYGZhy{}\PYGZhy{}\PYGZhy{}\PYGZhy{}\PYGZhy{}\PYGZhy{}\PYGZhy{}\PYGZhy{}\PYGZhy{}\PYGZhy{}\PYGZhy{}\PYGZhy{}\PYGZhy{}\PYGZhy{}\PYGZhy{}\PYGZhy{}\PYGZhy{}\PYGZhy{}\PYGZhy{}\PYGZhy{}\PYGZhy{}\PYGZhy{}\PYGZhy{}\PYGZhy{}\PYGZhy{}\PYGZhy{}\PYGZhy{}\PYGZhy{}\PYGZhy{}\PYGZhy{}\PYGZhy{}\PYGZhy{}\PYGZhy{}\PYGZhy{}\PYGZhy{}\PYGZhy{}\PYGZhy{}\PYGZhy{}\PYGZhy{}\PYGZhy{}\PYGZhy{}\PYGZhy{}\PYGZhy{}\PYGZhy{}\PYGZhy{}\PYGZhy{}\PYGZhy{}\PYGZhy{}\PYGZhy{}\PYGZhy{}\PYGZhy{}\PYGZhy{}\PYGZhy{}\PYGZhy{}\PYGZhy{}\PYGZhy{}\PYGZhy{}\PYGZhy{}\PYGZhy{}\PYGZhy{}\PYGZhy{}\PYGZhy{}\PYGZhy{} */}

\PYG{k+kt}{void}\PYG{+w}{ }\PYG{o}{*}\PYG{n+nf}{PairBornGauss::extract}\PYG{p}{(}\PYG{k}{const}\PYG{+w}{ }\PYG{k+kt}{char}\PYG{+w}{ }\PYG{o}{*}\PYG{n}{str}\PYG{p}{,}\PYG{+w}{ }\PYG{k+kt}{int}\PYG{+w}{ }\PYG{o}{\PYGZam{}}\PYG{n}{dim}\PYG{p}{)}
\PYG{p}{\PYGZob{}}
\PYG{+w}{  }\PYG{n}{dim}\PYG{+w}{ }\PYG{o}{=}\PYG{+w}{ }\PYG{l+m+mi}{2}\PYG{p}{;}
\PYG{+w}{  }\PYG{k}{if}\PYG{+w}{ }\PYG{p}{(}\PYG{n}{strcmp}\PYG{p}{(}\PYG{n}{str}\PYG{p}{,}\PYG{+w}{ }\PYG{l+s}{\PYGZdq{}}\PYG{l+s}{biga0}\PYG{l+s}{\PYGZdq{}}\PYG{p}{)}\PYG{+w}{ }\PYG{o}{=}\PYG{o}{=}\PYG{+w}{ }\PYG{l+m+mi}{0}\PYG{p}{)}\PYG{+w}{ }\PYG{k}{return}\PYG{+w}{ }\PYG{p}{(}\PYG{k+kt}{void}\PYG{+w}{ }\PYG{o}{*}\PYG{p}{)}\PYG{+w}{ }\PYG{n}{biga0}\PYG{p}{;}
\PYG{+w}{  }\PYG{k}{if}\PYG{+w}{ }\PYG{p}{(}\PYG{n}{strcmp}\PYG{p}{(}\PYG{n}{str}\PYG{p}{,}\PYG{+w}{ }\PYG{l+s}{\PYGZdq{}}\PYG{l+s}{biga1}\PYG{l+s}{\PYGZdq{}}\PYG{p}{)}\PYG{+w}{ }\PYG{o}{=}\PYG{o}{=}\PYG{+w}{ }\PYG{l+m+mi}{0}\PYG{p}{)}\PYG{+w}{ }\PYG{k}{return}\PYG{+w}{ }\PYG{p}{(}\PYG{k+kt}{void}\PYG{+w}{ }\PYG{o}{*}\PYG{p}{)}\PYG{+w}{ }\PYG{n}{biga1}\PYG{p}{;}
\PYG{+w}{  }\PYG{k}{if}\PYG{+w}{ }\PYG{p}{(}\PYG{n}{strcmp}\PYG{p}{(}\PYG{n}{str}\PYG{p}{,}\PYG{+w}{ }\PYG{l+s}{\PYGZdq{}}\PYG{l+s}{r0}\PYG{l+s}{\PYGZdq{}}\PYG{p}{)}\PYG{+w}{ }\PYG{o}{=}\PYG{o}{=}\PYG{+w}{ }\PYG{l+m+mi}{0}\PYG{p}{)}\PYG{+w}{ }\PYG{k}{return}\PYG{+w}{ }\PYG{p}{(}\PYG{k+kt}{void}\PYG{+w}{ }\PYG{o}{*}\PYG{p}{)}\PYG{+w}{ }\PYG{n}{r0}\PYG{p}{;}
\PYG{+w}{  }\PYG{k}{return}\PYG{+w}{ }\PYG{n+nb}{nullptr}\PYG{p}{;}
\PYG{p}{\PYGZcb{}}
\end{sphinxVerbatim}

\sphinxAtStartPar
Since the mercury potential, for which we have implemented the
born/gauss pair style, has a temperature dependent parameter “biga1”, we
can automatically adapt the potential based on the Taylor\sphinxhyphen{}MacLaurin
expansion for “biga1” when performing a simulation with a temperature
ramp.  LAMMPS commands for that application are given below:

\begin{sphinxVerbatim}[commandchars=\\\{\}]
\PYG{k}{variable }\PYG{n+nv+nvIdentifier}{tlo}\PYG{+w}{  }\PYG{n}{index}\PYG{+w}{ }\PYG{l+m}{300.0}
\PYG{k}{variable }\PYG{n+nv+nvIdentifier}{thi}\PYG{+w}{  }\PYG{n}{index}\PYG{+w}{ }\PYG{l+m}{600.0}
\PYG{k}{variable }\PYG{n+nv+nvIdentifier}{temp}\PYG{+w}{ }\PYG{n}{equal}\PYG{+w}{ }\PYG{n}{ramp}\PYG{n+nv}{(}\PYG{n+nv}{v\PYGZus{}tlo,v\PYGZus{}thi}\PYG{n+nv}{)}
\PYG{k}{variable }\PYG{n+nv+nvIdentifier}{biga1}\PYG{+w}{ }\PYG{n}{equal}\PYG{+w}{ }\PYG{n+nv}{(}\PYG{n+nv}{\PYGZhy{}2.58717e\PYGZhy{}8*v\PYGZus{}temp+8.40841e\PYGZhy{}5}\PYG{n+nv}{)}\PYG{o}{*}\PYG{n}{v\PYGZus{}temp}\PYG{o}{+}\PYG{l+m}{1.97475e\PYGZhy{}2}

\PYG{k}{fix             }\PYG{n+nv+nvIdentifier}{1}\PYG{+w}{ }\PYG{n+nv+nvIdentifier}{all}\PYG{+w}{ }\PYG{n}{nvt}\PYG{+w}{ }\PYG{n}{temp}\PYG{+w}{ }\PYG{n+nv}{\PYGZdl{}\PYGZob{}}\PYG{n+nv}{tlo}\PYG{n+nv}{\PYGZcb{}}\PYG{+w}{ }\PYG{n+nv}{\PYGZdl{}\PYGZob{}}\PYG{n+nv}{thi}\PYG{n+nv}{\PYGZcb{}}\PYG{+w}{ }\PYG{l+m}{0.1}
\PYG{k}{fix             }\PYG{n+nv+nvIdentifier}{2}\PYG{+w}{ }\PYG{n+nv+nvIdentifier}{all}\PYG{+w}{ }\PYG{n}{adapt}\PYG{+w}{ }\PYG{l+m}{1}\PYG{+w}{ }\PYG{n}{pair}\PYG{+w}{ }\PYG{n}{born}\PYG{o}{/}\PYG{n}{gauss}\PYG{+w}{ }\PYG{n}{biga1}\PYG{+w}{ }\PYG{o}{*}\PYG{+w}{ }\PYG{o}{*}\PYG{+w}{ }\PYG{n}{v\PYGZus{}biga1}
\end{sphinxVerbatim}


\subsection{Case 2: a many\sphinxhyphen{}body potential}
\label{\detokenize{Developer_write_pair:case-2-a-many-body-potential}}
\sphinxAtStartPar
Since there is a detailed description of the purpose and general layout
of a pair style in the previous case, we will focus on where the
implementation of a typical many\sphinxhyphen{}body potential \sphinxstyleemphasis{differs} from a
pair\sphinxhyphen{}wise additive potential.  We will use the implementation of the
Tersoff potential as {\hyperref[\detokenize{pair_tersoff::doc}]{\sphinxcrossref{\DUrole{doc}{pair\_style tersoff}}}} as an
example.  The complete implementation can be found in the files
\sphinxcode{\sphinxupquote{src/MANYBODY/pair\_tersoff.cpp}} and \sphinxcode{\sphinxupquote{src/MANYBODY/pair\_tersoff.h}} of
the LAMMPS source code.


\subsubsection{Constructor}
\label{\detokenize{Developer_write_pair:constructor}}
\sphinxAtStartPar
In the constructor, several {\hyperref[\detokenize{Modify_pair::doc}]{\sphinxcrossref{\DUrole{doc}{pair style flags}}}} must
be set differently for many\sphinxhyphen{}body potentials:
\begin{itemize}
\item {} 
\sphinxAtStartPar
the potential is not pair\sphinxhyphen{}wise additive, so the \sphinxcode{\sphinxupquote{single()}} function
cannot be used. This is indicated by setting the \sphinxtitleref{single\_enable}
member variable to 0 (default value is 1)

\item {} 
\sphinxAtStartPar
many\sphinxhyphen{}body potentials are usually not written to {\hyperref[\detokenize{write_restart::doc}]{\sphinxcrossref{\DUrole{doc}{binary
restart files}}}}.  This is indicated by setting the member
variable \sphinxtitleref{restartinfo} to 0 (default is 1)

\item {} 
\sphinxAtStartPar
many\sphinxhyphen{}body potentials typically read \sphinxstyleemphasis{all} parameters from a file which
stores parameters indexed with a string (e.g. the element).  For this,
only a single {\hyperref[\detokenize{pair_coeff::doc}]{\sphinxcrossref{\DUrole{doc}{pair\_coeff * *}}}} command is allowed.
This requirement is set and checked for, when the member variable
\sphinxtitleref{one\_coeff} is set to 1 (default value is 0)

\item {} 
\sphinxAtStartPar
many\sphinxhyphen{}body potentials can produce incorrect results if pairs of atoms
are excluded from the neighbor list, e.g. explicitly by
{\hyperref[\detokenize{neigh_modify::doc}]{\sphinxcrossref{\DUrole{doc}{neigh\_modify exclude}}}} or implicitly through
defining bonds, angles, etc. and having a {\hyperref[\detokenize{special_bonds::doc}]{\sphinxcrossref{\DUrole{doc}{special\_bonds setting}}}} that is not “special\_bonds lj/coul 1.0 1.0 1.0”.
LAMMPS will check for this and print a suitable warning, when the
member variable \sphinxtitleref{manybody\_flag} is set to 1 (default value is 0).

\end{itemize}

\begin{sphinxVerbatim}[commandchars=\\\{\}]
\PYG{n}{PairTersoff}\PYG{o}{:}\PYG{o}{:}\PYG{n}{PairTersoff}\PYG{p}{(}\PYG{n}{LAMMPS}\PYG{+w}{ }\PYG{o}{*}\PYG{n}{lmp}\PYG{p}{)}\PYG{+w}{ }\PYG{o}{:}\PYG{+w}{ }\PYG{n}{Pair}\PYG{p}{(}\PYG{n}{lmp}\PYG{p}{)}
\PYG{p}{\PYGZob{}}
\PYG{+w}{  }\PYG{n}{single\PYGZus{}enable}\PYG{+w}{ }\PYG{o}{=}\PYG{+w}{ }\PYG{l+m+mi}{0}\PYG{p}{;}
\PYG{+w}{  }\PYG{n}{restartinfo}\PYG{+w}{ }\PYG{o}{=}\PYG{+w}{ }\PYG{l+m+mi}{0}\PYG{p}{;}
\PYG{+w}{  }\PYG{n}{one\PYGZus{}coeff}\PYG{+w}{ }\PYG{o}{=}\PYG{+w}{ }\PYG{l+m+mi}{1}\PYG{p}{;}
\PYG{+w}{  }\PYG{n}{manybody\PYGZus{}flag}\PYG{+w}{ }\PYG{o}{=}\PYG{+w}{ }\PYG{l+m+mi}{1}\PYG{p}{;}
\end{sphinxVerbatim}


\subsubsection{Neighbor list request}
\label{\detokenize{Developer_write_pair:neighbor-list-request}}
\sphinxAtStartPar
For computing the three\sphinxhyphen{}body interactions of the Tersoff potential a
“full” neighbor list (both atoms of a pair are listed in each other’s
neighbor list) is required.  By default a “half” neighbor list is
requested (each pair is listed only once).  The request is made in
the \sphinxcode{\sphinxupquote{init\_style()}} function.  A more in\sphinxhyphen{}depth discussion of neighbor
lists in LAMMPS and how to request them is in {\hyperref[\detokenize{Developer_notes:request-neighbor-list}]{\sphinxcrossref{\DUrole{std,std-ref}{this section of the
documentation}}}}

\sphinxAtStartPar
Also, additional conditions must be met for some global settings which
are checked in the \sphinxcode{\sphinxupquote{init\_style()}} function.

\begin{sphinxVerbatim}[commandchars=\\\{\}]
\PYG{c+cm}{/* \PYGZhy{}\PYGZhy{}\PYGZhy{}\PYGZhy{}\PYGZhy{}\PYGZhy{}\PYGZhy{}\PYGZhy{}\PYGZhy{}\PYGZhy{}\PYGZhy{}\PYGZhy{}\PYGZhy{}\PYGZhy{}\PYGZhy{}\PYGZhy{}\PYGZhy{}\PYGZhy{}\PYGZhy{}\PYGZhy{}\PYGZhy{}\PYGZhy{}\PYGZhy{}\PYGZhy{}\PYGZhy{}\PYGZhy{}\PYGZhy{}\PYGZhy{}\PYGZhy{}\PYGZhy{}\PYGZhy{}\PYGZhy{}\PYGZhy{}\PYGZhy{}\PYGZhy{}\PYGZhy{}\PYGZhy{}\PYGZhy{}\PYGZhy{}\PYGZhy{}\PYGZhy{}\PYGZhy{}\PYGZhy{}\PYGZhy{}\PYGZhy{}\PYGZhy{}\PYGZhy{}\PYGZhy{}\PYGZhy{}\PYGZhy{}\PYGZhy{}\PYGZhy{}\PYGZhy{}\PYGZhy{}\PYGZhy{}\PYGZhy{}\PYGZhy{}\PYGZhy{}\PYGZhy{}\PYGZhy{}\PYGZhy{}\PYGZhy{}\PYGZhy{}\PYGZhy{}\PYGZhy{}\PYGZhy{}\PYGZhy{}\PYGZhy{}\PYGZhy{}\PYGZhy{}}
\PYG{c+cm}{   init specific to this pair style}
\PYG{c+cm}{\PYGZhy{}\PYGZhy{}\PYGZhy{}\PYGZhy{}\PYGZhy{}\PYGZhy{}\PYGZhy{}\PYGZhy{}\PYGZhy{}\PYGZhy{}\PYGZhy{}\PYGZhy{}\PYGZhy{}\PYGZhy{}\PYGZhy{}\PYGZhy{}\PYGZhy{}\PYGZhy{}\PYGZhy{}\PYGZhy{}\PYGZhy{}\PYGZhy{}\PYGZhy{}\PYGZhy{}\PYGZhy{}\PYGZhy{}\PYGZhy{}\PYGZhy{}\PYGZhy{}\PYGZhy{}\PYGZhy{}\PYGZhy{}\PYGZhy{}\PYGZhy{}\PYGZhy{}\PYGZhy{}\PYGZhy{}\PYGZhy{}\PYGZhy{}\PYGZhy{}\PYGZhy{}\PYGZhy{}\PYGZhy{}\PYGZhy{}\PYGZhy{}\PYGZhy{}\PYGZhy{}\PYGZhy{}\PYGZhy{}\PYGZhy{}\PYGZhy{}\PYGZhy{}\PYGZhy{}\PYGZhy{}\PYGZhy{}\PYGZhy{}\PYGZhy{}\PYGZhy{}\PYGZhy{}\PYGZhy{}\PYGZhy{}\PYGZhy{}\PYGZhy{}\PYGZhy{}\PYGZhy{}\PYGZhy{}\PYGZhy{}\PYGZhy{}\PYGZhy{}\PYGZhy{}\PYGZhy{}\PYGZhy{}\PYGZhy{} */}

\PYG{k+kt}{void}\PYG{+w}{ }\PYG{n+nf}{PairTersoff::init\PYGZus{}style}\PYG{p}{(}\PYG{p}{)}
\PYG{p}{\PYGZob{}}
\PYG{+w}{  }\PYG{k}{if}\PYG{+w}{ }\PYG{p}{(}\PYG{n}{atom}\PYG{o}{\PYGZhy{}}\PYG{o}{\PYGZgt{}}\PYG{n}{tag\PYGZus{}enable}\PYG{+w}{ }\PYG{o}{=}\PYG{o}{=}\PYG{+w}{ }\PYG{l+m+mi}{0}\PYG{p}{)}
\PYG{+w}{    }\PYG{n}{error}\PYG{o}{\PYGZhy{}}\PYG{o}{\PYGZgt{}}\PYG{n}{all}\PYG{p}{(}\PYG{n}{FLERR}\PYG{p}{,}\PYG{l+s}{\PYGZdq{}}\PYG{l+s}{Pair style Tersoff requires atom IDs}\PYG{l+s}{\PYGZdq{}}\PYG{p}{)}\PYG{p}{;}
\PYG{+w}{  }\PYG{k}{if}\PYG{+w}{ }\PYG{p}{(}\PYG{n}{force}\PYG{o}{\PYGZhy{}}\PYG{o}{\PYGZgt{}}\PYG{n}{newton\PYGZus{}pair}\PYG{+w}{ }\PYG{o}{=}\PYG{o}{=}\PYG{+w}{ }\PYG{l+m+mi}{0}\PYG{p}{)}
\PYG{+w}{    }\PYG{n}{error}\PYG{o}{\PYGZhy{}}\PYG{o}{\PYGZgt{}}\PYG{n}{all}\PYG{p}{(}\PYG{n}{FLERR}\PYG{p}{,}\PYG{l+s}{\PYGZdq{}}\PYG{l+s}{Pair style Tersoff requires newton pair on}\PYG{l+s}{\PYGZdq{}}\PYG{p}{)}\PYG{p}{;}

\PYG{+w}{  }\PYG{c+c1}{// need a full neighbor list}

\PYG{+w}{  }\PYG{n}{neighbor}\PYG{o}{\PYGZhy{}}\PYG{o}{\PYGZgt{}}\PYG{n}{add\PYGZus{}request}\PYG{p}{(}\PYG{k}{this}\PYG{p}{,}\PYG{n}{NeighConst}\PYG{o}{:}\PYG{o}{:}\PYG{n}{REQ\PYGZus{}FULL}\PYG{p}{)}\PYG{p}{;}
\PYG{p}{\PYGZcb{}}
\end{sphinxVerbatim}


\subsubsection{Computing forces from the neighbor list}
\label{\detokenize{Developer_write_pair:computing-forces-from-the-neighbor-list}}
\sphinxAtStartPar
Computing forces for a many\sphinxhyphen{}body potential is usually more complex than
for a pair\sphinxhyphen{}wise additive potential and there are multiple components.
For Tersoff, there is a pair\sphinxhyphen{}wise additive two\sphinxhyphen{}body term (two nested
loops over indices \sphinxstyleemphasis{i} and \sphinxstyleemphasis{j}) and a three\sphinxhyphen{}body term (three nested
loops over indices \sphinxstyleemphasis{i}, \sphinxstyleemphasis{j}, and \sphinxstyleemphasis{k}).  Since the neighbor list has
all neighbors up to the maximum cutoff (for the two\sphinxhyphen{}body term), but
the three\sphinxhyphen{}body interactions have a significantly shorter cutoff,
a “short neighbor list” is also constructed at the same time while computing
the two\sphinxhyphen{}body term and looping over the neighbor list for the first time.

\begin{sphinxVerbatim}[commandchars=\\\{\}]
\PYG{k}{if}\PYG{+w}{ }\PYG{p}{(}\PYG{n}{rsq}\PYG{+w}{ }\PYG{o}{\PYGZlt{}}\PYG{+w}{ }\PYG{n}{cutshortsq}\PYG{p}{)}\PYG{+w}{ }\PYG{p}{\PYGZob{}}
\PYG{+w}{  }\PYG{n}{neighshort}\PYG{p}{[}\PYG{n}{numshort}\PYG{o}{+}\PYG{o}{+}\PYG{p}{]}\PYG{+w}{ }\PYG{o}{=}\PYG{+w}{ }\PYG{n}{j}\PYG{p}{;}
\PYG{+w}{  }\PYG{k}{if}\PYG{+w}{ }\PYG{p}{(}\PYG{n}{numshort}\PYG{+w}{ }\PYG{o}{\PYGZgt{}}\PYG{o}{=}\PYG{+w}{ }\PYG{n}{maxshort}\PYG{p}{)}\PYG{+w}{ }\PYG{p}{\PYGZob{}}
\PYG{+w}{    }\PYG{n}{maxshort}\PYG{+w}{ }\PYG{o}{+}\PYG{o}{=}\PYG{+w}{ }\PYG{n}{maxshort}\PYG{o}{/}\PYG{l+m+mi}{2}\PYG{p}{;}
\PYG{+w}{    }\PYG{n}{memory}\PYG{o}{\PYGZhy{}}\PYG{o}{\PYGZgt{}}\PYG{n}{grow}\PYG{p}{(}\PYG{n}{neighshort}\PYG{p}{,}\PYG{n}{maxshort}\PYG{p}{,}\PYG{l+s}{\PYGZdq{}}\PYG{l+s}{pair:neighshort}\PYG{l+s}{\PYGZdq{}}\PYG{p}{)}\PYG{p}{;}
\PYG{+w}{  }\PYG{p}{\PYGZcb{}}
\PYG{p}{\PYGZcb{}}
\end{sphinxVerbatim}

\sphinxAtStartPar
For the two\sphinxhyphen{}body term, only a half neighbor list would be needed, even
though we have requested a full list (for the three\sphinxhyphen{}body loops).
Rather than computing all interactions twice, we skip over half of
the entries.  This is done in a slightly complex way to make certain
the same choice is made across all subdomains and so that there is
no load imbalance introduced.

\begin{sphinxVerbatim}[commandchars=\\\{\}]
\PYG{n}{jtag}\PYG{+w}{ }\PYG{o}{=}\PYG{+w}{ }\PYG{n}{tag}\PYG{p}{[}\PYG{n}{j}\PYG{p}{]}\PYG{p}{;}
\PYG{k}{if}\PYG{+w}{ }\PYG{p}{(}\PYG{n}{itag}\PYG{+w}{ }\PYG{o}{\PYGZgt{}}\PYG{+w}{ }\PYG{n}{jtag}\PYG{p}{)}\PYG{+w}{ }\PYG{p}{\PYGZob{}}
\PYG{+w}{  }\PYG{k}{if}\PYG{+w}{ }\PYG{p}{(}\PYG{p}{(}\PYG{n}{itag}\PYG{o}{+}\PYG{n}{jtag}\PYG{p}{)}\PYG{+w}{ }\PYG{o}{\PYGZpc{}}\PYG{+w}{ }\PYG{l+m+mi}{2}\PYG{+w}{ }\PYG{o}{=}\PYG{o}{=}\PYG{+w}{ }\PYG{l+m+mi}{0}\PYG{p}{)}\PYG{+w}{ }\PYG{k}{continue}\PYG{p}{;}
\PYG{p}{\PYGZcb{}}\PYG{+w}{ }\PYG{k}{else}\PYG{+w}{ }\PYG{k}{if}\PYG{+w}{ }\PYG{p}{(}\PYG{n}{itag}\PYG{+w}{ }\PYG{o}{\PYGZlt{}}\PYG{+w}{ }\PYG{n}{jtag}\PYG{p}{)}\PYG{+w}{ }\PYG{p}{\PYGZob{}}
\PYG{+w}{  }\PYG{k}{if}\PYG{+w}{ }\PYG{p}{(}\PYG{p}{(}\PYG{n}{itag}\PYG{o}{+}\PYG{n}{jtag}\PYG{p}{)}\PYG{+w}{ }\PYG{o}{\PYGZpc{}}\PYG{+w}{ }\PYG{l+m+mi}{2}\PYG{+w}{ }\PYG{o}{=}\PYG{o}{=}\PYG{+w}{ }\PYG{l+m+mi}{1}\PYG{p}{)}\PYG{+w}{ }\PYG{k}{continue}\PYG{p}{;}
\PYG{p}{\PYGZcb{}}\PYG{+w}{ }\PYG{k}{else}\PYG{+w}{ }\PYG{p}{\PYGZob{}}
\PYG{+w}{  }\PYG{k}{if}\PYG{+w}{ }\PYG{p}{(}\PYG{n}{x}\PYG{p}{[}\PYG{n}{j}\PYG{p}{]}\PYG{p}{[}\PYG{l+m+mi}{2}\PYG{p}{]}\PYG{+w}{ }\PYG{o}{\PYGZlt{}}\PYG{+w}{ }\PYG{n}{x}\PYG{p}{[}\PYG{n}{i}\PYG{p}{]}\PYG{p}{[}\PYG{l+m+mi}{2}\PYG{p}{]}\PYG{p}{)}\PYG{+w}{ }\PYG{k}{continue}\PYG{p}{;}
\PYG{+w}{  }\PYG{k}{if}\PYG{+w}{ }\PYG{p}{(}\PYG{n}{x}\PYG{p}{[}\PYG{n}{j}\PYG{p}{]}\PYG{p}{[}\PYG{l+m+mi}{2}\PYG{p}{]}\PYG{+w}{ }\PYG{o}{=}\PYG{o}{=}\PYG{+w}{ }\PYG{n}{ztmp}\PYG{+w}{ }\PYG{o}{\PYGZam{}}\PYG{o}{\PYGZam{}}\PYG{+w}{ }\PYG{n}{x}\PYG{p}{[}\PYG{n}{j}\PYG{p}{]}\PYG{p}{[}\PYG{l+m+mi}{1}\PYG{p}{]}\PYG{+w}{ }\PYG{o}{\PYGZlt{}}\PYG{+w}{ }\PYG{n}{ytmp}\PYG{p}{)}\PYG{+w}{ }\PYG{k}{continue}\PYG{p}{;}
\PYG{+w}{  }\PYG{k}{if}\PYG{+w}{ }\PYG{p}{(}\PYG{n}{x}\PYG{p}{[}\PYG{n}{j}\PYG{p}{]}\PYG{p}{[}\PYG{l+m+mi}{2}\PYG{p}{]}\PYG{+w}{ }\PYG{o}{=}\PYG{o}{=}\PYG{+w}{ }\PYG{n}{ztmp}\PYG{+w}{ }\PYG{o}{\PYGZam{}}\PYG{o}{\PYGZam{}}\PYG{+w}{ }\PYG{n}{x}\PYG{p}{[}\PYG{n}{j}\PYG{p}{]}\PYG{p}{[}\PYG{l+m+mi}{1}\PYG{p}{]}\PYG{+w}{ }\PYG{o}{=}\PYG{o}{=}\PYG{+w}{ }\PYG{n}{ytmp}\PYG{+w}{ }\PYG{o}{\PYGZam{}}\PYG{o}{\PYGZam{}}\PYG{+w}{ }\PYG{n}{x}\PYG{p}{[}\PYG{n}{j}\PYG{p}{]}\PYG{p}{[}\PYG{l+m+mi}{0}\PYG{p}{]}\PYG{+w}{ }\PYG{o}{\PYGZlt{}}\PYG{+w}{ }\PYG{n}{xtmp}\PYG{p}{)}\PYG{+w}{ }\PYG{k}{continue}\PYG{p}{;}
\PYG{p}{\PYGZcb{}}
\end{sphinxVerbatim}

\sphinxAtStartPar
For the three\sphinxhyphen{}body term, there is one additional nested loop and it uses
the “short” neighbor list, accumulated previously.

\begin{sphinxVerbatim}[commandchars=\\\{\}]
\PYG{c+c1}{// three\PYGZhy{}body interactions}
\PYG{c+c1}{// skip immediately if I\PYGZhy{}J is not within cutoff}
\PYG{k+kt}{double}\PYG{+w}{ }\PYG{n}{fjxtmp}\PYG{p}{,}\PYG{n}{fjytmp}\PYG{p}{,}\PYG{n}{fjztmp}\PYG{p}{;}

\PYG{k}{for}\PYG{+w}{ }\PYG{p}{(}\PYG{n}{jj}\PYG{+w}{ }\PYG{o}{=}\PYG{+w}{ }\PYG{l+m+mi}{0}\PYG{p}{;}\PYG{+w}{ }\PYG{n}{jj}\PYG{+w}{ }\PYG{o}{\PYGZlt{}}\PYG{+w}{ }\PYG{n}{numshort}\PYG{p}{;}\PYG{+w}{ }\PYG{n}{jj}\PYG{o}{+}\PYG{o}{+}\PYG{p}{)}\PYG{+w}{ }\PYG{p}{\PYGZob{}}
\PYG{+w}{  }\PYG{n}{j}\PYG{+w}{ }\PYG{o}{=}\PYG{+w}{ }\PYG{n}{neighshort}\PYG{p}{[}\PYG{n}{jj}\PYG{p}{]}\PYG{p}{;}
\PYG{+w}{  }\PYG{n}{jtype}\PYG{+w}{ }\PYG{o}{=}\PYG{+w}{ }\PYG{n}{map}\PYG{p}{[}\PYG{n}{type}\PYG{p}{[}\PYG{n}{j}\PYG{p}{]}\PYG{p}{]}\PYG{p}{;}

\PYG{+w}{  }\PYG{p}{[}\PYG{p}{.}\PYG{p}{.}\PYG{p}{.}\PYG{p}{]}

\PYG{+w}{  }\PYG{k}{for}\PYG{+w}{ }\PYG{p}{(}\PYG{n}{kk}\PYG{+w}{ }\PYG{o}{=}\PYG{+w}{ }\PYG{l+m+mi}{0}\PYG{p}{;}\PYG{+w}{ }\PYG{n}{kk}\PYG{+w}{ }\PYG{o}{\PYGZlt{}}\PYG{+w}{ }\PYG{n}{numshort}\PYG{p}{;}\PYG{+w}{ }\PYG{n}{kk}\PYG{o}{+}\PYG{o}{+}\PYG{p}{)}\PYG{+w}{ }\PYG{p}{\PYGZob{}}
\PYG{+w}{    }\PYG{k}{if}\PYG{+w}{ }\PYG{p}{(}\PYG{n}{jj}\PYG{+w}{ }\PYG{o}{=}\PYG{o}{=}\PYG{+w}{ }\PYG{n}{kk}\PYG{p}{)}\PYG{+w}{ }\PYG{k}{continue}\PYG{p}{;}
\PYG{+w}{    }\PYG{n}{k}\PYG{+w}{ }\PYG{o}{=}\PYG{+w}{ }\PYG{n}{neighshort}\PYG{p}{[}\PYG{n}{kk}\PYG{p}{]}\PYG{p}{;}
\PYG{+w}{    }\PYG{n}{ktype}\PYG{+w}{ }\PYG{o}{=}\PYG{+w}{ }\PYG{n}{map}\PYG{p}{[}\PYG{n}{type}\PYG{p}{[}\PYG{n}{k}\PYG{p}{]}\PYG{p}{]}\PYG{p}{;}

\PYG{+w}{    }\PYG{p}{[}\PYG{p}{.}\PYG{p}{.}\PYG{p}{.}\PYG{p}{]}
\PYG{+w}{  }\PYG{p}{\PYGZcb{}}
\PYG{p}{[}\PYG{p}{.}\PYG{p}{.}\PYG{p}{.}\PYG{p}{]}
\end{sphinxVerbatim}


\subsubsection{Reading potential parameters}
\label{\detokenize{Developer_write_pair:reading-potential-parameters}}
\sphinxAtStartPar
For the Tersoff potential, the parameters are listed in a file and
associated with triples of elements.  Because we have set the
\sphinxcode{\sphinxupquote{one\_coeff}} flag to 1 in the constructor, there may only be a single
{\hyperref[\detokenize{pair_coeff::doc}]{\sphinxcrossref{\DUrole{doc}{pair\_coeff * *}}}} line in the input for this pair
style, and as a consequence the \sphinxcode{\sphinxupquote{coeff()}} function will only be called
once.  Thus, the \sphinxcode{\sphinxupquote{coeff()}} function has to do three tasks, each of
which is delegated to a function in the \sphinxcode{\sphinxupquote{PairTersoff}} class:
\begin{enumerate}
\sphinxsetlistlabels{\arabic}{enumi}{enumii}{}{.}%
\item {} 
\sphinxAtStartPar
map elements to atom types.  Those follow the potential file name in the
command line arguments and are processed by the \sphinxcode{\sphinxupquote{map\_element2type()}} function.

\item {} 
\sphinxAtStartPar
read and parse the potential parameter file in the \sphinxcode{\sphinxupquote{read\_file()}} function.

\item {} 
\sphinxAtStartPar
Build data structures where the original and derived parameters are
indexed by all possible triples of atom types and thus can be looked
up quickly in the loops for the force computation

\end{enumerate}

\begin{sphinxVerbatim}[commandchars=\\\{\}]
\PYG{k+kt}{void}\PYG{+w}{ }\PYG{n+nf}{PairTersoff::coeff}\PYG{p}{(}\PYG{k+kt}{int}\PYG{+w}{ }\PYG{n}{narg}\PYG{p}{,}\PYG{+w}{ }\PYG{k+kt}{char}\PYG{+w}{ }\PYG{o}{*}\PYG{o}{*}\PYG{n}{arg}\PYG{p}{)}
\PYG{p}{\PYGZob{}}
\PYG{+w}{  }\PYG{k}{if}\PYG{+w}{ }\PYG{p}{(}\PYG{o}{!}\PYG{n}{allocated}\PYG{p}{)}\PYG{+w}{ }\PYG{n}{allocate}\PYG{p}{(}\PYG{p}{)}\PYG{p}{;}

\PYG{+w}{  }\PYG{n}{map\PYGZus{}element2type}\PYG{p}{(}\PYG{n}{narg}\PYG{l+m+mi}{\PYGZhy{}3}\PYG{p}{,}\PYG{n}{arg}\PYG{o}{+}\PYG{l+m+mi}{3}\PYG{p}{)}\PYG{p}{;}

\PYG{+w}{  }\PYG{c+c1}{// read potential file and initialize potential parameters}

\PYG{+w}{  }\PYG{n}{read\PYGZus{}file}\PYG{p}{(}\PYG{n}{arg}\PYG{p}{[}\PYG{l+m+mi}{2}\PYG{p}{]}\PYG{p}{)}\PYG{p}{;}
\PYG{+w}{  }\PYG{n}{setup\PYGZus{}params}\PYG{p}{(}\PYG{p}{)}\PYG{p}{;}
\PYG{p}{\PYGZcb{}}
\end{sphinxVerbatim}


\subsection{Case 3: a potential requiring communication}
\label{\detokenize{Developer_write_pair:case-3-a-potential-requiring-communication}}
\sphinxAtStartPar
For some models, the interactions between atoms depends on properties of
their environment which have to be computed \sphinxstyleemphasis{before} the the forces can
be computed.  Since LAMMPS is designed to run in parallel using a
{\hyperref[\detokenize{Developer_par_part::doc}]{\sphinxcrossref{\DUrole{doc}{domain decomposition strategy}}}}, not all
information of the atoms may be directly available and thus
communication steps may be need to collect data from ghost atoms of
neighboring subdomains or send data to ghost atoms for application
during the pairwise computation.

\sphinxAtStartPar
Specifically, two communication patterns are needed: a “reverse
communication” and a “forward communication”.  The reverse communication
collects data added to “ghost” atoms from neighboring sub\sphinxhyphen{}domains and
sums it to their corresponding “local” atoms.  This communication is
only required and thus executed when the \sphinxcode{\sphinxupquote{Force::newton\_pair}} setting
is 1 (i.e. {\hyperref[\detokenize{newton::doc}]{\sphinxcrossref{\DUrole{doc}{newton on}}}}, the default).  The forward
communication is used to copy computed per\sphinxhyphen{}atom data from “local” atoms
to their corresponding “ghost” atoms in neighboring sub\sphinxhyphen{}domains.

\sphinxAtStartPar
For this we will look at how the embedding term of the {\hyperref[\detokenize{pair_eam::doc}]{\sphinxcrossref{\DUrole{doc}{embedded
atom potential EAM}}}} is implemented in LAMMPS.  The complete
implementation of this pair style can be found in the files
\sphinxcode{\sphinxupquote{src/MANYBODY/pair\_eam.cpp}} and \sphinxcode{\sphinxupquote{src/MANYBODY/pair\_eam.h}} of the
LAMMPS source code.


\subsubsection{Allocating additional per\sphinxhyphen{}atom storage}
\label{\detokenize{Developer_write_pair:allocating-additional-per-atom-storage}}
\sphinxAtStartPar
First suitable (local) per\sphinxhyphen{}atom arrays (\sphinxtitleref{rho}, \sphinxtitleref{fp}, \sphinxtitleref{numforce}) are
allocated. These have to be large enough to include ghost atoms, are not
used outside the \sphinxcode{\sphinxupquote{compute()}} function and are re\sphinxhyphen{}initialized to zero
once per timestep.

\begin{sphinxVerbatim}[commandchars=\\\{\}]
\PYG{k}{if}\PYG{+w}{ }\PYG{p}{(}\PYG{n}{atom}\PYG{o}{\PYGZhy{}}\PYG{o}{\PYGZgt{}}\PYG{n}{nmax}\PYG{+w}{ }\PYG{o}{\PYGZgt{}}\PYG{+w}{ }\PYG{n}{nmax}\PYG{p}{)}\PYG{+w}{ }\PYG{p}{\PYGZob{}}
\PYG{+w}{  }\PYG{n}{memory}\PYG{o}{\PYGZhy{}}\PYG{o}{\PYGZgt{}}\PYG{n}{destroy}\PYG{p}{(}\PYG{n}{rho}\PYG{p}{)}\PYG{p}{;}
\PYG{+w}{  }\PYG{n}{memory}\PYG{o}{\PYGZhy{}}\PYG{o}{\PYGZgt{}}\PYG{n}{destroy}\PYG{p}{(}\PYG{n}{fp}\PYG{p}{)}\PYG{p}{;}
\PYG{+w}{  }\PYG{n}{memory}\PYG{o}{\PYGZhy{}}\PYG{o}{\PYGZgt{}}\PYG{n}{destroy}\PYG{p}{(}\PYG{n}{numforce}\PYG{p}{)}\PYG{p}{;}
\PYG{+w}{  }\PYG{n}{nmax}\PYG{+w}{ }\PYG{o}{=}\PYG{+w}{ }\PYG{n}{atom}\PYG{o}{\PYGZhy{}}\PYG{o}{\PYGZgt{}}\PYG{n}{nmax}\PYG{p}{;}
\PYG{+w}{  }\PYG{n}{memory}\PYG{o}{\PYGZhy{}}\PYG{o}{\PYGZgt{}}\PYG{n}{create}\PYG{p}{(}\PYG{n}{rho}\PYG{p}{,}\PYG{n}{nmax}\PYG{p}{,}\PYG{l+s}{\PYGZdq{}}\PYG{l+s}{pair:rho}\PYG{l+s}{\PYGZdq{}}\PYG{p}{)}\PYG{p}{;}
\PYG{+w}{  }\PYG{n}{memory}\PYG{o}{\PYGZhy{}}\PYG{o}{\PYGZgt{}}\PYG{n}{create}\PYG{p}{(}\PYG{n}{fp}\PYG{p}{,}\PYG{n}{nmax}\PYG{p}{,}\PYG{l+s}{\PYGZdq{}}\PYG{l+s}{pair:fp}\PYG{l+s}{\PYGZdq{}}\PYG{p}{)}\PYG{p}{;}
\PYG{+w}{  }\PYG{n}{memory}\PYG{o}{\PYGZhy{}}\PYG{o}{\PYGZgt{}}\PYG{n}{create}\PYG{p}{(}\PYG{n}{numforce}\PYG{p}{,}\PYG{n}{nmax}\PYG{p}{,}\PYG{l+s}{\PYGZdq{}}\PYG{l+s}{pair:numforce}\PYG{l+s}{\PYGZdq{}}\PYG{p}{)}\PYG{p}{;}
\PYG{p}{\PYGZcb{}}
\end{sphinxVerbatim}


\subsubsection{Reverse communication}
\label{\detokenize{Developer_write_pair:reverse-communication}}
\sphinxAtStartPar
Then a first loop over all pairs (\sphinxstyleemphasis{i} and \sphinxstyleemphasis{j}) is performed, where data
is stored in the \sphinxtitleref{rho} array representing the electron density at the site of
\sphinxstyleemphasis{i} contributed from all neighbors \sphinxstyleemphasis{j}.  Since the EAM pair style uses
a half neighbor list (for efficiency reasons), a reverse communication is
needed to collect the contributions to \sphinxtitleref{rho} from ghost atoms (only if
{\hyperref[\detokenize{newton::doc}]{\sphinxcrossref{\DUrole{doc}{newton on}}}} is set for pair styles).

\begin{sphinxVerbatim}[commandchars=\\\{\}]
\PYG{k}{if}\PYG{+w}{ }\PYG{p}{(}\PYG{n}{newton\PYGZus{}pair}\PYG{p}{)}\PYG{+w}{ }\PYG{n}{comm}\PYG{o}{\PYGZhy{}}\PYG{o}{\PYGZgt{}}\PYG{n}{reverse\PYGZus{}comm}\PYG{p}{(}\PYG{k}{this}\PYG{p}{)}\PYG{p}{;}
\end{sphinxVerbatim}

\sphinxAtStartPar
To support the reverse communication, two functions must be defined:
\sphinxcode{\sphinxupquote{pack\_reverse\_comm()}} that copies relevant data into a buffer for ghost
atoms and \sphinxcode{\sphinxupquote{unpack\_reverse\_comm()}} that takes the collected data and adds
it to the \sphinxtitleref{rho} array for the corresponding local atoms that match the
ghost atoms.  In order to allocate sufficiently sized buffers, a flag
must be set in the pair style constructor. Since in this case a single
double precision number is communicated per atom, the \sphinxtitleref{comm\_reverse}
member variable is set to 1 (default is 0 = no reverse communication).

\begin{sphinxVerbatim}[commandchars=\\\{\}]
\PYG{k+kt}{int}\PYG{+w}{ }\PYG{n+nf}{PairEAM::pack\PYGZus{}reverse\PYGZus{}comm}\PYG{p}{(}\PYG{k+kt}{int}\PYG{+w}{ }\PYG{n}{n}\PYG{p}{,}\PYG{+w}{ }\PYG{k+kt}{int}\PYG{+w}{ }\PYG{n}{first}\PYG{p}{,}\PYG{+w}{ }\PYG{k+kt}{double}\PYG{+w}{ }\PYG{o}{*}\PYG{n}{buf}\PYG{p}{)}
\PYG{p}{\PYGZob{}}
\PYG{+w}{  }\PYG{k+kt}{int}\PYG{+w}{ }\PYG{n}{i}\PYG{p}{,}\PYG{n}{m}\PYG{p}{,}\PYG{n}{last}\PYG{p}{;}

\PYG{+w}{  }\PYG{n}{m}\PYG{+w}{ }\PYG{o}{=}\PYG{+w}{ }\PYG{l+m+mi}{0}\PYG{p}{;}
\PYG{+w}{  }\PYG{n}{last}\PYG{+w}{ }\PYG{o}{=}\PYG{+w}{ }\PYG{n}{first}\PYG{+w}{ }\PYG{o}{+}\PYG{+w}{ }\PYG{n}{n}\PYG{p}{;}
\PYG{+w}{  }\PYG{k}{for}\PYG{+w}{ }\PYG{p}{(}\PYG{n}{i}\PYG{+w}{ }\PYG{o}{=}\PYG{+w}{ }\PYG{n}{first}\PYG{p}{;}\PYG{+w}{ }\PYG{n}{i}\PYG{+w}{ }\PYG{o}{\PYGZlt{}}\PYG{+w}{ }\PYG{n}{last}\PYG{p}{;}\PYG{+w}{ }\PYG{n}{i}\PYG{o}{+}\PYG{o}{+}\PYG{p}{)}\PYG{+w}{ }\PYG{n}{buf}\PYG{p}{[}\PYG{n}{m}\PYG{o}{+}\PYG{o}{+}\PYG{p}{]}\PYG{+w}{ }\PYG{o}{=}\PYG{+w}{ }\PYG{n}{rho}\PYG{p}{[}\PYG{n}{i}\PYG{p}{]}\PYG{p}{;}
\PYG{+w}{  }\PYG{k}{return}\PYG{+w}{ }\PYG{n}{m}\PYG{p}{;}
\PYG{p}{\PYGZcb{}}

\PYG{k+kt}{void}\PYG{+w}{ }\PYG{n+nf}{PairEAM::unpack\PYGZus{}reverse\PYGZus{}comm}\PYG{p}{(}\PYG{k+kt}{int}\PYG{+w}{ }\PYG{n}{n}\PYG{p}{,}\PYG{+w}{ }\PYG{k+kt}{int}\PYG{+w}{ }\PYG{o}{*}\PYG{n}{list}\PYG{p}{,}\PYG{+w}{ }\PYG{k+kt}{double}\PYG{+w}{ }\PYG{o}{*}\PYG{n}{buf}\PYG{p}{)}
\PYG{p}{\PYGZob{}}
\PYG{+w}{  }\PYG{k+kt}{int}\PYG{+w}{ }\PYG{n}{i}\PYG{p}{,}\PYG{n}{j}\PYG{p}{,}\PYG{n}{m}\PYG{p}{;}

\PYG{+w}{  }\PYG{n}{m}\PYG{+w}{ }\PYG{o}{=}\PYG{+w}{ }\PYG{l+m+mi}{0}\PYG{p}{;}
\PYG{+w}{  }\PYG{k}{for}\PYG{+w}{ }\PYG{p}{(}\PYG{n}{i}\PYG{+w}{ }\PYG{o}{=}\PYG{+w}{ }\PYG{l+m+mi}{0}\PYG{p}{;}\PYG{+w}{ }\PYG{n}{i}\PYG{+w}{ }\PYG{o}{\PYGZlt{}}\PYG{+w}{ }\PYG{n}{n}\PYG{p}{;}\PYG{+w}{ }\PYG{n}{i}\PYG{o}{+}\PYG{o}{+}\PYG{p}{)}\PYG{+w}{ }\PYG{p}{\PYGZob{}}
\PYG{+w}{    }\PYG{n}{j}\PYG{+w}{ }\PYG{o}{=}\PYG{+w}{ }\PYG{n}{list}\PYG{p}{[}\PYG{n}{i}\PYG{p}{]}\PYG{p}{;}
\PYG{+w}{    }\PYG{n}{rho}\PYG{p}{[}\PYG{n}{j}\PYG{p}{]}\PYG{+w}{ }\PYG{o}{+}\PYG{o}{=}\PYG{+w}{ }\PYG{n}{buf}\PYG{p}{[}\PYG{n}{m}\PYG{o}{+}\PYG{o}{+}\PYG{p}{]}\PYG{p}{;}
\PYG{+w}{  }\PYG{p}{\PYGZcb{}}
\PYG{p}{\PYGZcb{}}
\end{sphinxVerbatim}


\subsubsection{Forward communication}
\label{\detokenize{Developer_write_pair:forward-communication}}
\sphinxAtStartPar
From the density array \sphinxtitleref{rho}, the derivative of the embedding energy
\sphinxtitleref{fp} is computed. The computation is only done for “local” atoms, but
for the force computation, that property also is needed on ghost atoms.
For that a forward communication is needed.

\begin{sphinxVerbatim}[commandchars=\\\{\}]
\PYG{n}{comm}\PYG{o}{\PYGZhy{}}\PYG{o}{\PYGZgt{}}\PYG{n}{forward\PYGZus{}comm}\PYG{p}{(}\PYG{k}{this}\PYG{p}{)}\PYG{p}{;}
\end{sphinxVerbatim}

\sphinxAtStartPar
Similar to the reverse communication, this requires implementing a
\sphinxcode{\sphinxupquote{pack\_forward\_comm()}} and an \sphinxcode{\sphinxupquote{unpack\_forward\_comm()}} function.
Since there is one double precision number per atom that needs to be
communicated, we must set the \sphinxtitleref{comm\_forward} member variable to 1
(default is 0 = no forward communication).

\begin{sphinxVerbatim}[commandchars=\\\{\}]
\PYG{k+kt}{int}\PYG{+w}{ }\PYG{n+nf}{PairEAM::pack\PYGZus{}forward\PYGZus{}comm}\PYG{p}{(}\PYG{k+kt}{int}\PYG{+w}{ }\PYG{n}{n}\PYG{p}{,}\PYG{+w}{ }\PYG{k+kt}{int}\PYG{+w}{ }\PYG{o}{*}\PYG{n}{list}\PYG{p}{,}\PYG{+w}{ }\PYG{k+kt}{double}\PYG{+w}{ }\PYG{o}{*}\PYG{n}{buf}\PYG{p}{,}\PYG{+w}{ }\PYG{k+kt}{int}\PYG{+w}{ }\PYG{n}{pbc\PYGZus{}flag}\PYG{p}{,}\PYG{+w}{ }\PYG{k+kt}{int}\PYG{+w}{ }\PYG{o}{*}\PYG{n}{pbc}\PYG{p}{)}
\PYG{p}{\PYGZob{}}
\PYG{+w}{  }\PYG{k+kt}{int}\PYG{+w}{ }\PYG{n}{i}\PYG{p}{,}\PYG{n}{j}\PYG{p}{,}\PYG{n}{m}\PYG{p}{;}

\PYG{+w}{  }\PYG{n}{m}\PYG{+w}{ }\PYG{o}{=}\PYG{+w}{ }\PYG{l+m+mi}{0}\PYG{p}{;}
\PYG{+w}{  }\PYG{k}{for}\PYG{+w}{ }\PYG{p}{(}\PYG{n}{i}\PYG{+w}{ }\PYG{o}{=}\PYG{+w}{ }\PYG{l+m+mi}{0}\PYG{p}{;}\PYG{+w}{ }\PYG{n}{i}\PYG{+w}{ }\PYG{o}{\PYGZlt{}}\PYG{+w}{ }\PYG{n}{n}\PYG{p}{;}\PYG{+w}{ }\PYG{n}{i}\PYG{o}{+}\PYG{o}{+}\PYG{p}{)}\PYG{+w}{ }\PYG{p}{\PYGZob{}}
\PYG{+w}{    }\PYG{n}{j}\PYG{+w}{ }\PYG{o}{=}\PYG{+w}{ }\PYG{n}{list}\PYG{p}{[}\PYG{n}{i}\PYG{p}{]}\PYG{p}{;}
\PYG{+w}{    }\PYG{n}{buf}\PYG{p}{[}\PYG{n}{m}\PYG{o}{+}\PYG{o}{+}\PYG{p}{]}\PYG{+w}{ }\PYG{o}{=}\PYG{+w}{ }\PYG{n}{fp}\PYG{p}{[}\PYG{n}{j}\PYG{p}{]}\PYG{p}{;}
\PYG{+w}{  }\PYG{p}{\PYGZcb{}}
\PYG{+w}{  }\PYG{k}{return}\PYG{+w}{ }\PYG{n}{m}\PYG{p}{;}
\PYG{p}{\PYGZcb{}}

\PYG{k+kt}{void}\PYG{+w}{ }\PYG{n+nf}{PairEAM::unpack\PYGZus{}forward\PYGZus{}comm}\PYG{p}{(}\PYG{k+kt}{int}\PYG{+w}{ }\PYG{n}{n}\PYG{p}{,}\PYG{+w}{ }\PYG{k+kt}{int}\PYG{+w}{ }\PYG{n}{first}\PYG{p}{,}\PYG{+w}{ }\PYG{k+kt}{double}\PYG{+w}{ }\PYG{o}{*}\PYG{n}{buf}\PYG{p}{)}
\PYG{p}{\PYGZob{}}
\PYG{+w}{  }\PYG{k+kt}{int}\PYG{+w}{ }\PYG{n}{i}\PYG{p}{,}\PYG{n}{m}\PYG{p}{,}\PYG{n}{last}\PYG{p}{;}

\PYG{+w}{  }\PYG{n}{m}\PYG{+w}{ }\PYG{o}{=}\PYG{+w}{ }\PYG{l+m+mi}{0}\PYG{p}{;}
\PYG{+w}{  }\PYG{n}{last}\PYG{+w}{ }\PYG{o}{=}\PYG{+w}{ }\PYG{n}{first}\PYG{+w}{ }\PYG{o}{+}\PYG{+w}{ }\PYG{n}{n}\PYG{p}{;}
\PYG{+w}{  }\PYG{k}{for}\PYG{+w}{ }\PYG{p}{(}\PYG{n}{i}\PYG{+w}{ }\PYG{o}{=}\PYG{+w}{ }\PYG{n}{first}\PYG{p}{;}\PYG{+w}{ }\PYG{n}{i}\PYG{+w}{ }\PYG{o}{\PYGZlt{}}\PYG{+w}{ }\PYG{n}{last}\PYG{p}{;}\PYG{+w}{ }\PYG{n}{i}\PYG{o}{+}\PYG{o}{+}\PYG{p}{)}\PYG{+w}{ }\PYG{n}{fp}\PYG{p}{[}\PYG{n}{i}\PYG{p}{]}\PYG{+w}{ }\PYG{o}{=}\PYG{+w}{ }\PYG{n}{buf}\PYG{p}{[}\PYG{n}{m}\PYG{o}{+}\PYG{o}{+}\PYG{p}{]}\PYG{p}{;}
\PYG{p}{\PYGZcb{}}
\end{sphinxVerbatim}


\subsection{Case 4: potentials without a compute() function}
\label{\detokenize{Developer_write_pair:case-4-potentials-without-a-compute-function}}
\sphinxAtStartPar
A small number of pair style classes do not implement a \sphinxcode{\sphinxupquote{compute()}}
function, but instead use that of a different pair style.


\subsubsection{Embedded atom variants “eam/fs” and “eam/alloy”}
\label{\detokenize{Developer_write_pair:embedded-atom-variants-eam-fs-and-eam-alloy}}
\sphinxAtStartPar
The pair styles {\hyperref[\detokenize{pair_eam::doc}]{\sphinxcrossref{\DUrole{doc}{eam/fs and eam/alloy}}}} share the same
model and potential function as the {\hyperref[\detokenize{pair_eam::doc}]{\sphinxcrossref{\DUrole{doc}{eam pair style}}}}.
They differ in the format of the potential files.  Pair style {\hyperref[\detokenize{pair_eam::doc}]{\sphinxcrossref{\DUrole{doc}{eam}}}} supports only potential files for single elements.  For
multi\sphinxhyphen{}element systems, the mixed terms are computed from mixed
parameters.  The \sphinxstyleemphasis{eam/fs} and \sphinxstyleemphasis{eam/alloy} pair styles, however,
\sphinxstylestrong{require} the use of a single potential file for all elements where
the mixed element potential is included in the tabulation.  That enables
more accurate models for alloys, since the mixed terms can be adjusted
for a better representation of material properties compared to terms
created from mixing of per\sphinxhyphen{}element terms in the \sphinxcode{\sphinxupquote{PairEAM}} class.

\sphinxAtStartPar
We take a closer at the \sphinxstyleemphasis{eam/alloy} pair style.  The complete
implementation is in the files \sphinxcode{\sphinxupquote{src/MANYBODY/pair\_eam\_alloy.cpp}} and
\sphinxcode{\sphinxupquote{src/MANYBODY/pair\_eam\_alloy.h}}.

\sphinxAtStartPar
The \sphinxcode{\sphinxupquote{PairEAMAlloy}} class is derived from \sphinxcode{\sphinxupquote{PairEAM}} and not \sphinxcode{\sphinxupquote{Pair}}
and overrides only a small number of functions:

\begin{sphinxVerbatim}[commandchars=\\\{\}]
\PYG{k}{class}\PYG{+w}{ }\PYG{n+nc}{PairEAMAlloy}\PYG{+w}{ }\PYG{o}{:}\PYG{+w}{ }\PYG{k}{virtual}\PYG{+w}{ }\PYG{k}{public}\PYG{+w}{ }\PYG{n}{PairEAM}\PYG{+w}{ }\PYG{p}{\PYGZob{}}
\PYG{+w}{ }\PYG{k}{public}\PYG{o}{:}
\PYG{+w}{  }\PYG{n}{PairEAMAlloy}\PYG{p}{(}\PYG{k}{class}\PYG{+w}{ }\PYG{n+nc}{LAMMPS}\PYG{+w}{ }\PYG{o}{*}\PYG{p}{)}\PYG{p}{;}
\PYG{+w}{  }\PYG{k+kt}{void}\PYG{+w}{ }\PYG{n+nf}{coeff}\PYG{p}{(}\PYG{k+kt}{int}\PYG{p}{,}\PYG{+w}{ }\PYG{k+kt}{char}\PYG{+w}{ }\PYG{o}{*}\PYG{o}{*}\PYG{p}{)}\PYG{+w}{ }\PYG{k}{override}\PYG{p}{;}

\PYG{+w}{ }\PYG{k}{protected}\PYG{o}{:}
\PYG{+w}{  }\PYG{k+kt}{void}\PYG{+w}{ }\PYG{n}{read\PYGZus{}file}\PYG{p}{(}\PYG{k+kt}{char}\PYG{+w}{ }\PYG{o}{*}\PYG{p}{)}\PYG{+w}{ }\PYG{k}{override}\PYG{p}{;}
\PYG{+w}{  }\PYG{k+kt}{void}\PYG{+w}{ }\PYG{n+nf}{file2array}\PYG{p}{(}\PYG{p}{)}\PYG{+w}{ }\PYG{k}{override}\PYG{p}{;}
\PYG{p}{\PYGZcb{}}\PYG{p}{;}
\end{sphinxVerbatim}

\sphinxAtStartPar
All other functionality is inherited from the base classes.  In the
constructor we set the \sphinxcode{\sphinxupquote{one\_coeff}} flag and the \sphinxcode{\sphinxupquote{many\_body}} flag to
1 to indicate the different behavior.

\begin{sphinxVerbatim}[commandchars=\\\{\}]
\PYG{n}{PairEAMAlloy}\PYG{o}{:}\PYG{o}{:}\PYG{n}{PairEAMAlloy}\PYG{p}{(}\PYG{n}{LAMMPS}\PYG{+w}{ }\PYG{o}{*}\PYG{n}{lmp}\PYG{p}{)}\PYG{+w}{ }\PYG{o}{:}\PYG{+w}{ }\PYG{n}{PairEAM}\PYG{p}{(}\PYG{n}{lmp}\PYG{p}{)}
\PYG{p}{\PYGZob{}}
\PYG{+w}{  }\PYG{n}{one\PYGZus{}coeff}\PYG{+w}{ }\PYG{o}{=}\PYG{+w}{ }\PYG{l+m+mi}{1}\PYG{p}{;}
\PYG{+w}{  }\PYG{n}{manybody\PYGZus{}flag}\PYG{+w}{ }\PYG{o}{=}\PYG{+w}{ }\PYG{l+m+mi}{1}\PYG{p}{;}
\PYG{p}{\PYGZcb{}}
\end{sphinxVerbatim}

\sphinxAtStartPar
The \sphinxcode{\sphinxupquote{coeff()}} function (not shown here) implements the different
behavior when processing the {\hyperref[\detokenize{pair_coeff::doc}]{\sphinxcrossref{\DUrole{doc}{pair\_coeff command}}}}.
The \sphinxcode{\sphinxupquote{read\_file()}} and \sphinxcode{\sphinxupquote{file2array()}} replace the corresponding
\sphinxcode{\sphinxupquote{PairEAM}} class functions to accommodate the different data and
file format.


\subsubsection{AIREBO and AIREBO\sphinxhyphen{}M potentials}
\label{\detokenize{Developer_write_pair:airebo-and-airebo-m-potentials}}
\sphinxAtStartPar
The AIREBO\sphinxhyphen{}M potential differs from the better known AIREBO potential in
that it use a Morse potential instead of a Lennard\sphinxhyphen{}Jones potential for
non\sphinxhyphen{}bonded interactions.  Since this difference is very minimal compared
to the entire potential, both potentials are implemented in the
\sphinxcode{\sphinxupquote{PairAIREBO}} class and which non\sphinxhyphen{}bonded potential is used is
determined by the value of the \sphinxcode{\sphinxupquote{morseflag}} flag, which would be set to
either 0 or 1.

\begin{sphinxVerbatim}[commandchars=\\\{\}]
\PYG{k}{class}\PYG{+w}{ }\PYG{n+nc}{PairAIREBOMorse}\PYG{+w}{ }\PYG{o}{:}\PYG{+w}{ }\PYG{k}{public}\PYG{+w}{ }\PYG{n}{PairAIREBO}\PYG{+w}{ }\PYG{p}{\PYGZob{}}
\PYG{+w}{ }\PYG{k}{public}\PYG{o}{:}
\PYG{+w}{  }\PYG{n}{PairAIREBOMorse}\PYG{p}{(}\PYG{k}{class}\PYG{+w}{ }\PYG{n+nc}{LAMMPS}\PYG{+w}{ }\PYG{o}{*}\PYG{p}{)}\PYG{p}{;}
\PYG{+w}{  }\PYG{k+kt}{void}\PYG{+w}{ }\PYG{n+nf}{settings}\PYG{p}{(}\PYG{k+kt}{int}\PYG{p}{,}\PYG{+w}{ }\PYG{k+kt}{char}\PYG{+w}{ }\PYG{o}{*}\PYG{o}{*}\PYG{p}{)}\PYG{+w}{ }\PYG{k}{override}\PYG{p}{;}
\PYG{p}{\PYGZcb{}}\PYG{p}{;}
\end{sphinxVerbatim}

\sphinxAtStartPar
The \sphinxcode{\sphinxupquote{morseflag}} variable defaults to 0 and is set to 1 in the
\sphinxcode{\sphinxupquote{PairAIREBOMorse::settings()}} function which is called by the
{\hyperref[\detokenize{pair_style::doc}]{\sphinxcrossref{\DUrole{doc}{pair\_style}}}} command.  This function delegates
all command line processing and setting of other parameters to the
\sphinxcode{\sphinxupquote{PairAIREBO::settings()}} function of the base class.

\begin{sphinxVerbatim}[commandchars=\\\{\}]
\PYG{k+kt}{void}\PYG{+w}{ }\PYG{n+nf}{PairAIREBOMorse::settings}\PYG{p}{(}\PYG{k+kt}{int}\PYG{+w}{ }\PYG{n}{narg}\PYG{p}{,}\PYG{+w}{ }\PYG{k+kt}{char}\PYG{+w}{ }\PYG{o}{*}\PYG{o}{*}\PYG{n}{arg}\PYG{p}{)}
\PYG{p}{\PYGZob{}}
\PYG{+w}{  }\PYG{n}{PairAIREBO}\PYG{o}{:}\PYG{o}{:}\PYG{n}{settings}\PYG{p}{(}\PYG{n}{narg}\PYG{p}{,}\PYG{+w}{ }\PYG{n}{arg}\PYG{p}{)}\PYG{p}{;}

\PYG{+w}{  }\PYG{n}{morseflag}\PYG{+w}{ }\PYG{o}{=}\PYG{+w}{ }\PYG{l+m+mi}{1}\PYG{p}{;}
\PYG{p}{\PYGZcb{}}
\end{sphinxVerbatim}

\sphinxAtStartPar
The complete implementation is in the files
\sphinxcode{\sphinxupquote{src/MANYBODY/pair\_airebo.cpp}}, \sphinxcode{\sphinxupquote{src/MANYBODY/pair\_airebo.h}},
\sphinxcode{\sphinxupquote{src/MANYBODY/pair\_airebo\_morse.cpp}},
\sphinxcode{\sphinxupquote{src/MANYBODY/pair\_airebo\_morse.h}}.


\bigskip\hrule\bigskip

\phantomsection\label{\detokenize{Developer_write_pair:bomont2}}
\sphinxAtStartPar
\sphinxstylestrong{(Bomont)} Bomont, Bretonnet, J. Chem. Phys. 124, 054504 (2006)

\sphinxstepscope


\subsection{Writing a new fix style}
\label{\detokenize{Developer_write_fix:writing-a-new-fix-style}}\label{\detokenize{Developer_write_fix::doc}}
\sphinxAtStartPar
Writing fix styles is a flexible way of extending LAMMPS.  Users can
implement many things using fixes.  Some fix styles are only used
internally to support compute styles or pair styles:
\begin{itemize}
\item {} 
\sphinxAtStartPar
change particles attributes (positions, velocities, forces, etc.). Examples: \sphinxcode{\sphinxupquote{FixNVE}}, \sphinxcode{\sphinxupquote{FixFreeze}}.

\item {} 
\sphinxAtStartPar
read or write data.  Example: \sphinxcode{\sphinxupquote{FixRestart}}.

\item {} 
\sphinxAtStartPar
adding or modifying properties due to geometry. Example: \sphinxcode{\sphinxupquote{FixWall}}.

\item {} 
\sphinxAtStartPar
interacting with other subsystems or external code: Examples: \sphinxcode{\sphinxupquote{FixTTM}}, \sphinxcode{\sphinxupquote{FixExternal}}, \sphinxcode{\sphinxupquote{FixMDI}}

\item {} 
\sphinxAtStartPar
saving information for analysis or future use (previous positions,
for instance). Examples: \sphinxcode{\sphinxupquote{FixAveTime}}, \sphinxcode{\sphinxupquote{FixStoreState}}.

\end{itemize}

\sphinxAtStartPar
All fixes are derived from the \sphinxcode{\sphinxupquote{Fix}} base class and must have a
constructor with the signature: \sphinxcode{\sphinxupquote{FixPrintVel(class LAMMPS *, int, char **)}}.

\sphinxAtStartPar
Every fix must be registered in LAMMPS by writing the following lines
of code in the header before include guards:

\begin{sphinxVerbatim}[commandchars=\\\{\}]
\PYG{c+cp}{\PYGZsh{}}\PYG{c+cp}{ifdef FIX\PYGZus{}CLASS}
\PYG{c+c1}{// clang\PYGZhy{}format off}
\PYG{n}{FixStyle}\PYG{p}{(}\PYG{n}{print}\PYG{o}{/}\PYG{n}{vel}\PYG{p}{,}\PYG{n}{FixPrintVel}\PYG{p}{)}\PYG{p}{;}
\PYG{c+c1}{// clang\PYGZhy{}format on}
\PYG{c+cp}{\PYGZsh{}}\PYG{c+cp}{else}
\PYG{c+cm}{/* the definition of the FixPrintVel class comes here */}
\PYG{p}{.}\PYG{p}{.}\PYG{p}{.}
\PYG{c+cp}{\PYGZsh{}}\PYG{c+cp}{endif}
\end{sphinxVerbatim}

\sphinxAtStartPar
Where \sphinxcode{\sphinxupquote{print/vel}} is the style name of your fix in the input script and
\sphinxcode{\sphinxupquote{FixPrintVel}} is the name of the class. The header file would be called
\sphinxcode{\sphinxupquote{fix\_print\_vel.h}} and the implementation file \sphinxcode{\sphinxupquote{fix\_print\_vel.cpp}}.
These conventions allow LAMMPS to automatically integrate it into the
executable when compiling and associate your new fix class with the designated
keyword when it parses the input script.

\sphinxAtStartPar
Let’s write a simple fix which will print the average velocity at the end
of each timestep. First of all, implement a constructor:

\begin{sphinxVerbatim}[commandchars=\\\{\}]
\PYG{n}{FixPrintVel}\PYG{o}{:}\PYG{o}{:}\PYG{n}{FixPrintVel}\PYG{p}{(}\PYG{n}{LAMMPS}\PYG{+w}{ }\PYG{o}{*}\PYG{n}{lmp}\PYG{p}{,}\PYG{+w}{ }\PYG{k+kt}{int}\PYG{+w}{ }\PYG{n}{narg}\PYG{p}{,}\PYG{+w}{ }\PYG{k+kt}{char}\PYG{+w}{ }\PYG{o}{*}\PYG{o}{*}\PYG{n}{arg}\PYG{p}{)}
\PYG{o}{:}\PYG{+w}{ }\PYG{n}{Fix}\PYG{p}{(}\PYG{n}{lmp}\PYG{p}{,}\PYG{+w}{ }\PYG{n}{narg}\PYG{p}{,}\PYG{+w}{ }\PYG{n}{arg}\PYG{p}{)}
\PYG{p}{\PYGZob{}}
\PYG{+w}{  }\PYG{k}{if}\PYG{+w}{ }\PYG{p}{(}\PYG{n}{narg}\PYG{+w}{ }\PYG{o}{\PYGZlt{}}\PYG{+w}{ }\PYG{l+m+mi}{4}\PYG{p}{)}\PYG{+w}{ }\PYG{n}{utils}\PYG{o}{:}\PYG{o}{:}\PYG{n}{missing\PYGZus{}cmd\PYGZus{}args}\PYG{p}{(}\PYG{n}{FLERR}\PYG{p}{,}\PYG{+w}{ }\PYG{l+s}{\PYGZdq{}}\PYG{l+s}{fix print/vel}\PYG{l+s}{\PYGZdq{}}\PYG{p}{,}\PYG{+w}{ }\PYG{n}{error}\PYG{p}{)}\PYG{p}{;}

\PYG{+w}{  }\PYG{n}{nevery}\PYG{+w}{ }\PYG{o}{=}\PYG{+w}{ }\PYG{n}{utils}\PYG{o}{:}\PYG{o}{:}\PYG{n}{inumeric}\PYG{p}{(}\PYG{n}{FLERR}\PYG{p}{,}\PYG{n}{arg}\PYG{p}{[}\PYG{l+m+mi}{3}\PYG{p}{]}\PYG{p}{,}\PYG{n+nb}{false}\PYG{p}{,}\PYG{n}{lmp}\PYG{p}{)}\PYG{p}{;}
\PYG{+w}{  }\PYG{k}{if}\PYG{+w}{ }\PYG{p}{(}\PYG{n}{nevery}\PYG{+w}{ }\PYG{o}{\PYGZlt{}}\PYG{o}{=}\PYG{+w}{ }\PYG{l+m+mi}{0}\PYG{p}{)}
\PYG{+w}{    }\PYG{n}{error}\PYG{o}{\PYGZhy{}}\PYG{o}{\PYGZgt{}}\PYG{n}{all}\PYG{p}{(}\PYG{n}{FLERR}\PYG{p}{,}\PYG{l+s}{\PYGZdq{}}\PYG{l+s}{Illegal fix print/vel nevery value: \PYGZob{}\PYGZcb{}}\PYG{l+s}{\PYGZdq{}}\PYG{p}{,}\PYG{+w}{ }\PYG{n}{nevery}\PYG{p}{)}\PYG{p}{;}
\PYG{p}{\PYGZcb{}}
\end{sphinxVerbatim}

\sphinxAtStartPar
In the constructor you should parse the fix arguments which are
specified in the script.  All fixes have pretty much the same syntax:
\sphinxcode{\sphinxupquote{fix \textless{}fix\sphinxhyphen{}ID\textgreater{} \textless{}fix group\textgreater{} \textless{}fix name\textgreater{} \textless{}fix arguments ...\textgreater{}}}. The first 3
parameters are parsed by Fix base class constructor, while \sphinxcode{\sphinxupquote{\textless{}fix
arguments\textgreater{}}} should be parsed by you.  In our case, we need to specify
how often we want to print an average velocity. For instance, once in 50
timesteps: \sphinxcode{\sphinxupquote{fix 1 print/vel 50}}. There is a special variable in the
Fix class called \sphinxcode{\sphinxupquote{nevery}} which specifies how often the method
\sphinxcode{\sphinxupquote{end\_of\_step()}} is called.  Thus all we need to do is just set it up.

\sphinxAtStartPar
The next method we need to implement is \sphinxcode{\sphinxupquote{setmask()}}:

\begin{sphinxVerbatim}[commandchars=\\\{\}]
\PYG{k+kt}{int}\PYG{+w}{ }\PYG{n+nf}{FixPrintVel::setmask}\PYG{p}{(}\PYG{p}{)}
\PYG{p}{\PYGZob{}}
\PYG{+w}{  }\PYG{k+kt}{int}\PYG{+w}{ }\PYG{n}{mask}\PYG{+w}{ }\PYG{o}{=}\PYG{+w}{ }\PYG{l+m+mi}{0}\PYG{p}{;}
\PYG{+w}{  }\PYG{n}{mask}\PYG{+w}{ }\PYG{o}{|}\PYG{o}{=}\PYG{+w}{ }\PYG{n}{FixConst}\PYG{o}{:}\PYG{o}{:}\PYG{n}{END\PYGZus{}OF\PYGZus{}STEP}\PYG{p}{;}
\PYG{+w}{  }\PYG{k}{return}\PYG{+w}{ }\PYG{n}{mask}\PYG{p}{;}
\PYG{p}{\PYGZcb{}}
\end{sphinxVerbatim}

\sphinxAtStartPar
Here the we specify which methods of the fix should be called during
{\hyperref[\detokenize{Developer_flow::doc}]{\sphinxcrossref{\DUrole{doc}{execution of a timestep}}}}.  The constant
\sphinxcode{\sphinxupquote{END\_OF\_STEP}} corresponds to the \sphinxcode{\sphinxupquote{end\_of\_step()}} method.

\begin{sphinxVerbatim}[commandchars=\\\{\}]
\PYG{k+kt}{void}\PYG{+w}{ }\PYG{n+nf}{FixPrintVel::end\PYGZus{}of\PYGZus{}step}\PYG{p}{(}\PYG{p}{)}
\PYG{p}{\PYGZob{}}
\PYG{+w}{  }\PYG{c+c1}{// for add3, scale3}
\PYG{+w}{  }\PYG{k}{using}\PYG{+w}{ }\PYG{k}{namespace}\PYG{+w}{ }\PYG{n+nn}{MathExtra}\PYG{p}{;}

\PYG{+w}{  }\PYG{k+kt}{double}\PYG{o}{*}\PYG{o}{*}\PYG{+w}{ }\PYG{n}{v}\PYG{+w}{ }\PYG{o}{=}\PYG{+w}{ }\PYG{n}{atom}\PYG{o}{\PYGZhy{}}\PYG{o}{\PYGZgt{}}\PYG{n}{v}\PYG{p}{;}
\PYG{+w}{  }\PYG{k+kt}{int}\PYG{+w}{ }\PYG{n}{nlocal}\PYG{+w}{ }\PYG{o}{=}\PYG{+w}{ }\PYG{n}{atom}\PYG{o}{\PYGZhy{}}\PYG{o}{\PYGZgt{}}\PYG{n}{nlocal}\PYG{p}{;}
\PYG{+w}{  }\PYG{k+kt}{double}\PYG{+w}{ }\PYG{n}{localAvgVel}\PYG{p}{[}\PYG{l+m+mi}{4}\PYG{p}{]}\PYG{p}{;}\PYG{+w}{ }\PYG{c+c1}{// 4th element for particles count}
\PYG{+w}{  }\PYG{n}{memset}\PYG{p}{(}\PYG{n}{localAvgVel}\PYG{p}{,}\PYG{+w}{ }\PYG{l+m+mi}{0}\PYG{p}{,}\PYG{+w}{ }\PYG{l+m+mi}{4}\PYG{+w}{ }\PYG{o}{*}\PYG{+w}{ }\PYG{k}{sizeof}\PYG{p}{(}\PYG{k+kt}{double}\PYG{p}{)}\PYG{p}{)}\PYG{p}{;}
\PYG{+w}{  }\PYG{k}{for}\PYG{+w}{ }\PYG{p}{(}\PYG{k+kt}{int}\PYG{+w}{ }\PYG{n}{particleInd}\PYG{+w}{ }\PYG{o}{=}\PYG{+w}{ }\PYG{l+m+mi}{0}\PYG{p}{;}\PYG{+w}{ }\PYG{n}{particleInd}\PYG{+w}{ }\PYG{o}{\PYGZlt{}}\PYG{+w}{ }\PYG{n}{nlocal}\PYG{p}{;}\PYG{+w}{ }\PYG{o}{+}\PYG{o}{+}\PYG{n}{particleInd}\PYG{p}{)}\PYG{+w}{ }\PYG{p}{\PYGZob{}}
\PYG{+w}{    }\PYG{n}{add3}\PYG{p}{(}\PYG{n}{localAvgVel}\PYG{p}{,}\PYG{+w}{ }\PYG{n}{v}\PYG{p}{[}\PYG{n}{particleInd}\PYG{p}{]}\PYG{p}{,}\PYG{+w}{ }\PYG{n}{localAvgVel}\PYG{p}{)}\PYG{p}{;}
\PYG{+w}{  }\PYG{p}{\PYGZcb{}}
\PYG{+w}{  }\PYG{n}{localAvgVel}\PYG{p}{[}\PYG{l+m+mi}{3}\PYG{p}{]}\PYG{+w}{ }\PYG{o}{=}\PYG{+w}{ }\PYG{n}{nlocal}\PYG{p}{;}
\PYG{+w}{  }\PYG{k+kt}{double}\PYG{+w}{ }\PYG{n}{globalAvgVel}\PYG{p}{[}\PYG{l+m+mi}{4}\PYG{p}{]}\PYG{p}{;}
\PYG{+w}{  }\PYG{n}{memset}\PYG{p}{(}\PYG{n}{globalAvgVel}\PYG{p}{,}\PYG{+w}{ }\PYG{l+m+mi}{0}\PYG{p}{,}\PYG{+w}{ }\PYG{l+m+mi}{4}\PYG{+w}{ }\PYG{o}{*}\PYG{+w}{ }\PYG{k}{sizeof}\PYG{p}{(}\PYG{k+kt}{double}\PYG{p}{)}\PYG{p}{)}\PYG{p}{;}
\PYG{+w}{  }\PYG{n}{MPI\PYGZus{}Allreduce}\PYG{p}{(}\PYG{n}{localAvgVel}\PYG{p}{,}\PYG{+w}{ }\PYG{n}{globalAvgVel}\PYG{p}{,}\PYG{+w}{ }\PYG{l+m+mi}{4}\PYG{p}{,}\PYG{+w}{ }\PYG{n}{MPI\PYGZus{}DOUBLE}\PYG{p}{,}\PYG{+w}{ }\PYG{n}{MPI\PYGZus{}SUM}\PYG{p}{,}\PYG{+w}{ }\PYG{n}{world}\PYG{p}{)}\PYG{p}{;}
\PYG{+w}{  }\PYG{n}{scale3}\PYG{p}{(}\PYG{l+m+mf}{1.0}\PYG{+w}{ }\PYG{o}{/}\PYG{+w}{ }\PYG{n}{globalAvgVel}\PYG{p}{[}\PYG{l+m+mi}{3}\PYG{p}{]}\PYG{p}{,}\PYG{+w}{ }\PYG{n}{globalAvgVel}\PYG{p}{)}\PYG{p}{;}
\PYG{+w}{  }\PYG{k}{if}\PYG{+w}{ }\PYG{p}{(}\PYG{p}{(}\PYG{n}{comm}\PYG{o}{\PYGZhy{}}\PYG{o}{\PYGZgt{}}\PYG{n}{me}\PYG{+w}{ }\PYG{o}{=}\PYG{o}{=}\PYG{+w}{ }\PYG{l+m+mi}{0}\PYG{p}{)}\PYG{+w}{ }\PYG{o}{\PYGZam{}}\PYG{o}{\PYGZam{}}\PYG{+w}{ }\PYG{n}{screen}\PYG{p}{)}\PYG{+w}{ }\PYG{p}{\PYGZob{}}
\PYG{+w}{    }\PYG{n}{fmt}\PYG{o}{:}\PYG{o}{:}\PYG{n}{print}\PYG{p}{(}\PYG{n}{screen}\PYG{p}{,}\PYG{l+s}{\PYGZdq{}}\PYG{l+s}{\PYGZob{}\PYGZcb{}, \PYGZob{}\PYGZcb{}, \PYGZob{}\PYGZcb{}}\PYG{l+s+se}{\PYGZbs{}n}\PYG{l+s}{\PYGZdq{}}\PYG{p}{,}
\PYG{+w}{               }\PYG{n}{globalAvgVel}\PYG{p}{[}\PYG{l+m+mi}{0}\PYG{p}{]}\PYG{p}{,}\PYG{+w}{ }\PYG{n}{globalAvgVel}\PYG{p}{[}\PYG{l+m+mi}{1}\PYG{p}{]}\PYG{p}{,}\PYG{+w}{ }\PYG{n}{globalAvgVel}\PYG{p}{[}\PYG{l+m+mi}{2}\PYG{p}{]}\PYG{p}{)}\PYG{p}{;}
\PYG{+w}{  }\PYG{p}{\PYGZcb{}}
\PYG{p}{\PYGZcb{}}
\end{sphinxVerbatim}

\sphinxAtStartPar
In the code above, we use MathExtra routines defined in
\sphinxcode{\sphinxupquote{math\_extra.h}}.  There are bunch of math functions to work with
arrays of doubles as with math vectors.  It is also important to note
that LAMMPS code should always assume to be run in parallel and that
atom data is thus distributed across the MPI ranks.  Thus you can
only process data from local atoms directly and need to use MPI library
calls to combine or exchange data.  For serial execution, LAMMPS
comes bundled with the MPI STUBS library that contains the MPI library
function calls in dummy versions that only work for a single MPI rank.

\sphinxAtStartPar
In this code we use an instance of Atom class. This object is stored
in the Pointers class (see \sphinxcode{\sphinxupquote{pointers.h}}) which is the base class of
the Fix base class. This object contains references to various class
instances (the original instances are created and held by the LAMMPS
class) with all global information about the simulation system.
Data from the Pointers class is available to all classes inherited from
it using protected inheritance. Hence when you write you own class,
which is going to use LAMMPS data, don’t forget to inherit from Pointers
or pass a Pointer to it to all functions that need access. When writing
fixes we inherit from class Fix which is inherited from Pointers so
there is no need to inherit from it directly.

\sphinxAtStartPar
The code above computes average velocity for all particles in the
simulation.  Yet you have one unused parameter in fix call from the
script: \sphinxcode{\sphinxupquote{group\_name}}.  This parameter specifies the group of atoms
used in the fix. So we should compute the average for all particles in the
simulation only if \sphinxcode{\sphinxupquote{group\_name == "all"}}, but it can be any group.
The group membership information of an atom is contained in the \sphinxstyleemphasis{mask}
property of an atom and the bit corresponding to a given group is
stored in the groupbit variable which is defined in Fix base class:

\begin{sphinxVerbatim}[commandchars=\\\{\}]
\PYG{k}{for}\PYG{+w}{ }\PYG{p}{(}\PYG{k+kt}{int}\PYG{+w}{ }\PYG{n}{i}\PYG{+w}{ }\PYG{o}{=}\PYG{+w}{ }\PYG{l+m+mi}{0}\PYG{p}{;}\PYG{+w}{ }\PYG{n}{i}\PYG{+w}{ }\PYG{o}{\PYGZlt{}}\PYG{+w}{ }\PYG{n}{nlocal}\PYG{p}{;}\PYG{+w}{ }\PYG{o}{+}\PYG{o}{+}\PYG{n}{i}\PYG{p}{)}\PYG{+w}{ }\PYG{p}{\PYGZob{}}
\PYG{+w}{  }\PYG{k}{if}\PYG{+w}{ }\PYG{p}{(}\PYG{n}{atom}\PYG{o}{\PYGZhy{}}\PYG{o}{\PYGZgt{}}\PYG{n}{mask}\PYG{p}{[}\PYG{n}{i}\PYG{p}{]}\PYG{+w}{ }\PYG{o}{\PYGZam{}}\PYG{+w}{ }\PYG{n}{groupbit}\PYG{p}{)}\PYG{+w}{ }\PYG{p}{\PYGZob{}}
\PYG{+w}{  }\PYG{c+c1}{// Do all job here}
\PYG{+w}{  }\PYG{p}{\PYGZcb{}}
\PYG{p}{\PYGZcb{}}
\end{sphinxVerbatim}

\sphinxAtStartPar
Class Atom encapsulates atoms positions, velocities, forces, etc. Users
can access them using the particle index. Note, that particle indexes are
usually changed every few timesteps because of neighbor list rebuilds
and spatial sorting (to improve cache efficiency).

\sphinxAtStartPar
Let us consider another Fix example: We want to have a fix which stores
atoms position from the previous time step in your fix. The local atoms
indexes may not be valid on the next iteration. In order to handle
this situation there are several methods which should be implemented:
\begin{itemize}
\item {} 
\sphinxAtStartPar
\sphinxcode{\sphinxupquote{double memory\_usage()}}: return how much memory the fix uses (optional)

\item {} 
\sphinxAtStartPar
\sphinxcode{\sphinxupquote{void grow\_arrays(int)}}: do reallocation of the per\sphinxhyphen{}particle arrays in your fix

\item {} 
\sphinxAtStartPar
\sphinxcode{\sphinxupquote{void copy\_arrays(int i, int j, int delflag)}}: copy i\sphinxhyphen{}th per\sphinxhyphen{}particle
information to j\sphinxhyphen{}th particle position. Used when atom sorting is performed.
if delflag is set and atom j owns a body, move the body information to atom i.

\item {} 
\sphinxAtStartPar
\sphinxcode{\sphinxupquote{void set\_arrays(int i)}}: sets i\sphinxhyphen{}th particle related information to zero

\end{itemize}

\sphinxAtStartPar
Note, that if your class implements these methods, it must add calls of
add\_callback and delete\_callback to the constructor and destructor. Since we want
to store positions of atoms from the previous timestep, we need to add
\sphinxcode{\sphinxupquote{double** xold}} to the header file. Than add allocation code
to the constructor:

\begin{sphinxVerbatim}[commandchars=\\\{\}]
\PYG{n}{FixSavePos}\PYG{o}{:}\PYG{o}{:}\PYG{n}{FixSavePos}\PYG{p}{(}\PYG{n}{LAMMPS}\PYG{+w}{ }\PYG{o}{*}\PYG{n}{lmp}\PYG{p}{,}\PYG{+w}{ }\PYG{k+kt}{int}\PYG{+w}{ }\PYG{n}{narg}\PYG{p}{,}\PYG{+w}{ }\PYG{k+kt}{char}\PYG{+w}{ }\PYG{o}{*}\PYG{o}{*}\PYG{n}{arg}\PYG{p}{)}\PYG{p}{,}\PYG{+w}{ }\PYG{n}{xold}\PYG{p}{(}\PYG{n+nb}{nullptr}\PYG{p}{)}
\PYG{p}{\PYGZob{}}
\PYG{c+c1}{//...}
\PYG{+w}{  }\PYG{n}{memory}\PYG{o}{\PYGZhy{}}\PYG{o}{\PYGZgt{}}\PYG{n}{create}\PYG{p}{(}\PYG{n}{xold}\PYG{p}{,}\PYG{+w}{ }\PYG{n}{atom}\PYG{o}{\PYGZhy{}}\PYG{o}{\PYGZgt{}}\PYG{n}{nmax}\PYG{p}{,}\PYG{+w}{ }\PYG{l+m+mi}{3}\PYG{p}{,}\PYG{+w}{ }\PYG{l+s}{\PYGZdq{}}\PYG{l+s}{FixSavePos:x}\PYG{l+s}{\PYGZdq{}}\PYG{p}{)}\PYG{p}{;}
\PYG{+w}{  }\PYG{n}{atom}\PYG{o}{\PYGZhy{}}\PYG{o}{\PYGZgt{}}\PYG{n}{add\PYGZus{}callback}\PYG{p}{(}\PYG{l+m+mi}{0}\PYG{p}{)}\PYG{p}{;}
\PYG{p}{\PYGZcb{}}

\PYG{n}{FixSavePos}\PYG{o}{:}\PYG{o}{:}\PYG{o}{\PYGZti{}}\PYG{n}{FixSavePos}\PYG{p}{(}\PYG{p}{)}\PYG{+w}{ }\PYG{p}{\PYGZob{}}
\PYG{+w}{  }\PYG{n}{atom}\PYG{o}{\PYGZhy{}}\PYG{o}{\PYGZgt{}}\PYG{n}{delete\PYGZus{}callback}\PYG{p}{(}\PYG{n}{id}\PYG{p}{,}\PYG{+w}{ }\PYG{l+m+mi}{0}\PYG{p}{)}\PYG{p}{;}
\PYG{+w}{  }\PYG{n}{memory}\PYG{o}{\PYGZhy{}}\PYG{o}{\PYGZgt{}}\PYG{n}{destroy}\PYG{p}{(}\PYG{n}{xold}\PYG{p}{)}\PYG{p}{;}
\PYG{p}{\PYGZcb{}}
\end{sphinxVerbatim}

\sphinxAtStartPar
Implement the aforementioned methods:

\begin{sphinxVerbatim}[commandchars=\\\{\}]
\PYG{k+kt}{double}\PYG{+w}{ }\PYG{n+nf}{FixSavePos::memory\PYGZus{}usage}\PYG{p}{(}\PYG{p}{)}
\PYG{p}{\PYGZob{}}
\PYG{+w}{  }\PYG{k+kt}{int}\PYG{+w}{ }\PYG{n}{nmax}\PYG{+w}{ }\PYG{o}{=}\PYG{+w}{ }\PYG{n}{atom}\PYG{o}{\PYGZhy{}}\PYG{o}{\PYGZgt{}}\PYG{n}{nmax}\PYG{p}{;}
\PYG{+w}{  }\PYG{k+kt}{double}\PYG{+w}{ }\PYG{n}{bytes}\PYG{+w}{ }\PYG{o}{=}\PYG{+w}{ }\PYG{l+m+mf}{0.0}\PYG{p}{;}
\PYG{+w}{  }\PYG{n}{bytes}\PYG{+w}{ }\PYG{o}{+}\PYG{o}{=}\PYG{+w}{ }\PYG{n}{nmax}\PYG{+w}{ }\PYG{o}{*}\PYG{+w}{ }\PYG{l+m+mi}{3}\PYG{+w}{ }\PYG{o}{*}\PYG{+w}{ }\PYG{k}{sizeof}\PYG{p}{(}\PYG{k+kt}{double}\PYG{p}{)}\PYG{p}{;}
\PYG{+w}{  }\PYG{k}{return}\PYG{+w}{ }\PYG{n}{bytes}\PYG{p}{;}
\PYG{p}{\PYGZcb{}}

\PYG{k+kt}{void}\PYG{+w}{ }\PYG{n+nf}{FixSavePos::grow\PYGZus{}arrays}\PYG{p}{(}\PYG{k+kt}{int}\PYG{+w}{ }\PYG{n}{nmax}\PYG{p}{)}
\PYG{p}{\PYGZob{}}
\PYG{+w}{  }\PYG{n}{memory}\PYG{o}{\PYGZhy{}}\PYG{o}{\PYGZgt{}}\PYG{n}{grow}\PYG{p}{(}\PYG{n}{xold}\PYG{p}{,}\PYG{+w}{ }\PYG{n}{nmax}\PYG{p}{,}\PYG{+w}{ }\PYG{l+m+mi}{3}\PYG{p}{,}\PYG{+w}{ }\PYG{l+s}{\PYGZdq{}}\PYG{l+s}{FixSavePos:xold}\PYG{l+s}{\PYGZdq{}}\PYG{p}{)}\PYG{p}{;}
\PYG{p}{\PYGZcb{}}

\PYG{k+kt}{void}\PYG{+w}{ }\PYG{n+nf}{FixSavePos::copy\PYGZus{}arrays}\PYG{p}{(}\PYG{k+kt}{int}\PYG{+w}{ }\PYG{n}{i}\PYG{p}{,}\PYG{+w}{ }\PYG{k+kt}{int}\PYG{+w}{ }\PYG{n}{j}\PYG{p}{,}\PYG{+w}{ }\PYG{k+kt}{int}\PYG{+w}{ }\PYG{n}{delflag}\PYG{p}{)}
\PYG{p}{\PYGZob{}}
\PYG{+w}{  }\PYG{n}{memcpy}\PYG{p}{(}\PYG{n}{xold}\PYG{p}{[}\PYG{n}{j}\PYG{p}{]}\PYG{p}{,}\PYG{+w}{ }\PYG{n}{xold}\PYG{p}{[}\PYG{n}{i}\PYG{p}{]}\PYG{p}{,}\PYG{+w}{ }\PYG{k}{sizeof}\PYG{p}{(}\PYG{k+kt}{double}\PYG{p}{)}\PYG{+w}{ }\PYG{o}{*}\PYG{+w}{ }\PYG{l+m+mi}{3}\PYG{p}{)}\PYG{p}{;}
\PYG{p}{\PYGZcb{}}

\PYG{k+kt}{void}\PYG{+w}{ }\PYG{n+nf}{FixSavePos::set\PYGZus{}arrays}\PYG{p}{(}\PYG{k+kt}{int}\PYG{+w}{ }\PYG{n}{i}\PYG{p}{)}
\PYG{p}{\PYGZob{}}
\PYG{+w}{  }\PYG{n}{memset}\PYG{p}{(}\PYG{n}{xold}\PYG{p}{[}\PYG{n}{i}\PYG{p}{]}\PYG{p}{,}\PYG{+w}{ }\PYG{l+m+mi}{0}\PYG{p}{,}\PYG{+w}{ }\PYG{k}{sizeof}\PYG{p}{(}\PYG{k+kt}{double}\PYG{p}{)}\PYG{+w}{ }\PYG{o}{*}\PYG{+w}{ }\PYG{l+m+mi}{3}\PYG{p}{)}\PYG{p}{;}
\PYG{p}{\PYGZcb{}}

\PYG{k+kt}{int}\PYG{+w}{ }\PYG{n+nf}{FixSavePos::pack\PYGZus{}exchange}\PYG{p}{(}\PYG{k+kt}{int}\PYG{+w}{ }\PYG{n}{i}\PYG{p}{,}\PYG{+w}{ }\PYG{k+kt}{double}\PYG{+w}{ }\PYG{o}{*}\PYG{n}{buf}\PYG{p}{)}
\PYG{p}{\PYGZob{}}
\PYG{+w}{  }\PYG{k+kt}{int}\PYG{+w}{ }\PYG{n}{m}\PYG{+w}{ }\PYG{o}{=}\PYG{+w}{ }\PYG{l+m+mi}{0}\PYG{p}{;}
\PYG{+w}{  }\PYG{n}{buf}\PYG{p}{[}\PYG{n}{m}\PYG{o}{+}\PYG{o}{+}\PYG{p}{]}\PYG{+w}{ }\PYG{o}{=}\PYG{+w}{ }\PYG{n}{xold}\PYG{p}{[}\PYG{n}{i}\PYG{p}{]}\PYG{p}{[}\PYG{l+m+mi}{0}\PYG{p}{]}\PYG{p}{;}
\PYG{+w}{  }\PYG{n}{buf}\PYG{p}{[}\PYG{n}{m}\PYG{o}{+}\PYG{o}{+}\PYG{p}{]}\PYG{+w}{ }\PYG{o}{=}\PYG{+w}{ }\PYG{n}{xold}\PYG{p}{[}\PYG{n}{i}\PYG{p}{]}\PYG{p}{[}\PYG{l+m+mi}{1}\PYG{p}{]}\PYG{p}{;}
\PYG{+w}{  }\PYG{n}{buf}\PYG{p}{[}\PYG{n}{m}\PYG{o}{+}\PYG{o}{+}\PYG{p}{]}\PYG{+w}{ }\PYG{o}{=}\PYG{+w}{ }\PYG{n}{xold}\PYG{p}{[}\PYG{n}{i}\PYG{p}{]}\PYG{p}{[}\PYG{l+m+mi}{2}\PYG{p}{]}\PYG{p}{;}

\PYG{+w}{  }\PYG{k}{return}\PYG{+w}{ }\PYG{n}{m}\PYG{p}{;}
\PYG{p}{\PYGZcb{}}

\PYG{k+kt}{int}\PYG{+w}{ }\PYG{n+nf}{FixSavePos::unpack\PYGZus{}exchange}\PYG{p}{(}\PYG{k+kt}{int}\PYG{+w}{ }\PYG{n}{nlocal}\PYG{p}{,}\PYG{+w}{ }\PYG{k+kt}{double}\PYG{+w}{ }\PYG{o}{*}\PYG{n}{buf}\PYG{p}{)}
\PYG{p}{\PYGZob{}}
\PYG{+w}{  }\PYG{k+kt}{int}\PYG{+w}{ }\PYG{n}{m}\PYG{+w}{ }\PYG{o}{=}\PYG{+w}{ }\PYG{l+m+mi}{0}\PYG{p}{;}
\PYG{+w}{  }\PYG{n}{xold}\PYG{p}{[}\PYG{n}{nlocal}\PYG{p}{]}\PYG{p}{[}\PYG{l+m+mi}{0}\PYG{p}{]}\PYG{+w}{ }\PYG{o}{=}\PYG{+w}{ }\PYG{n}{buf}\PYG{p}{[}\PYG{n}{m}\PYG{o}{+}\PYG{o}{+}\PYG{p}{]}\PYG{p}{;}
\PYG{+w}{  }\PYG{n}{xold}\PYG{p}{[}\PYG{n}{nlocal}\PYG{p}{]}\PYG{p}{[}\PYG{l+m+mi}{1}\PYG{p}{]}\PYG{+w}{ }\PYG{o}{=}\PYG{+w}{ }\PYG{n}{buf}\PYG{p}{[}\PYG{n}{m}\PYG{o}{+}\PYG{o}{+}\PYG{p}{]}\PYG{p}{;}
\PYG{+w}{  }\PYG{n}{xold}\PYG{p}{[}\PYG{n}{nlocal}\PYG{p}{]}\PYG{p}{[}\PYG{l+m+mi}{2}\PYG{p}{]}\PYG{+w}{ }\PYG{o}{=}\PYG{+w}{ }\PYG{n}{buf}\PYG{p}{[}\PYG{n}{m}\PYG{o}{+}\PYG{o}{+}\PYG{p}{]}\PYG{p}{;}

\PYG{+w}{  }\PYG{k}{return}\PYG{+w}{ }\PYG{n}{m}\PYG{p}{;}
\PYG{p}{\PYGZcb{}}
\end{sphinxVerbatim}

\sphinxAtStartPar
Now, a little bit about memory allocation. We use the Memory class which
is just a bunch of template functions for allocating 1D and 2D
arrays. So you need to add include \sphinxcode{\sphinxupquote{memory.h}} to have access to them.

\sphinxAtStartPar
Finally, if you need to write/read some global information used in
your fix to the restart file, you might do it by setting the flag
\sphinxcode{\sphinxupquote{restart\_global = 1}} in the constructor and implementing methods
\sphinxcode{\sphinxupquote{void write\_restart(FILE *fp)}} and \sphinxcode{\sphinxupquote{void restart(char *buf)}}.
If, in addition, you want to write the per\sphinxhyphen{}atom property to restart
files then these additional settings and functions are needed:
\begin{itemize}
\item {} 
\sphinxAtStartPar
a fix flag indicating this needs to be set \sphinxcode{\sphinxupquote{restart\_peratom = 1;}}

\item {} 
\sphinxAtStartPar
\sphinxcode{\sphinxupquote{atom\sphinxhyphen{}\textgreater{}add\_callback()}} and \sphinxcode{\sphinxupquote{atom\sphinxhyphen{}\textgreater{}delete\_callback()}} must be called
a second time with the final argument set to 1 instead of 0 (indicating
restart processing instead of per\sphinxhyphen{}atom data memory management).

\item {} 
\sphinxAtStartPar
the functions \sphinxcode{\sphinxupquote{void pack\_restart(int i, double *buf)}} and
\sphinxcode{\sphinxupquote{void unpack\_restart(int nlocal, int nth)}} need to be implemented

\end{itemize}

\sphinxstepscope


\section{Notes for developers and code maintainers}
\label{\detokenize{Developer_notes:notes-for-developers-and-code-maintainers}}\label{\detokenize{Developer_notes::doc}}
\sphinxAtStartPar
This section documents how some of the code functionality within
LAMMPS works at a conceptual level.  Comments on code in source files
typically document what a variable stores, what a small section of
code does, or what a function does and its input/outputs.  The topics
on this page are intended to document code functionality at a higher level.

\sphinxAtStartPar
Available topics are:
\begin{itemize}
\item {} 
\sphinxAtStartPar
{\hyperref[\detokenize{Developer_notes:reading-and-parsing-of-text-and-text-files}]{\sphinxcrossref{Reading and parsing of text and text files}}}

\item {} 
\sphinxAtStartPar
{\hyperref[\detokenize{Developer_notes:requesting-and-accessing-neighbor-lists}]{\sphinxcrossref{Requesting and accessing neighbor lists}}}

\item {} 
\sphinxAtStartPar
{\hyperref[\detokenize{Developer_notes:choosing-between-a-custom-atom-style-fix-property-atom-and-fix-store-atom}]{\sphinxcrossref{Choosing between a custom atom style, fix property/atom, and fix STORE/ATOM}}}

\item {} 
\sphinxAtStartPar
{\hyperref[\detokenize{Developer_notes:fix-contributions-to-instantaneous-energy-virial-and-cumulative-energy}]{\sphinxcrossref{Fix contributions to instantaneous energy, virial, and cumulative energy}}}

\item {} 
\sphinxAtStartPar
{\hyperref[\detokenize{Developer_notes:kspace-pppm-fft-grids}]{\sphinxcrossref{KSpace PPPM FFT grids}}}

\end{itemize}


\bigskip\hrule\bigskip



\subsection{Reading and parsing of text and text files}
\label{\detokenize{Developer_notes:reading-and-parsing-of-text-and-text-files}}
\sphinxAtStartPar
Classes in LAMMPS frequently need to read in additional
data from a file, e.g. potential parameters from a potential file for
manybody potentials.  LAMMPS provides several custom classes and
convenience functions to simplify the process.  They offer the
following benefits:
\begin{itemize}
\item {} 
\sphinxAtStartPar
better code reuse and fewer lines of code needed to implement reading
and parsing data from a file

\item {} 
\sphinxAtStartPar
better detection of format errors, incompatible data, and better error messages

\item {} 
\sphinxAtStartPar
exit with an error message instead of silently converting only part of the
text to a number or returning a 0 on unrecognized text and thus reading incorrect values

\item {} 
\sphinxAtStartPar
re\sphinxhyphen{}entrant code through avoiding global static variables (as used by \sphinxcode{\sphinxupquote{strtok()}})

\item {} 
\sphinxAtStartPar
transparent support for translating unsupported UTF\sphinxhyphen{}8 characters to their ASCII equivalents
(the text\sphinxhyphen{}to\sphinxhyphen{}value conversion functions \sphinxstylestrong{only} accept ASCII characters)

\end{itemize}

\sphinxAtStartPar
In most cases (e.g. potential files) the same data is needed on all MPI
ranks.  Then it is best to do the reading and parsing only on MPI rank
0, and communicate the data later with one or more \sphinxcode{\sphinxupquote{MPI\_Bcast()}}
calls.  For reading generic text and potential parameter files the
custom classes {\hyperref[\detokenize{Developer_utils:_CPPv4N9LAMMPS_NS14TextFileReaderE}]{\sphinxcrossref{\sphinxcode{\sphinxupquote{TextFileReader}}}}}
and {\hyperref[\detokenize{Developer_utils:_CPPv4N9LAMMPS_NS19PotentialFileReaderE}]{\sphinxcrossref{\sphinxcode{\sphinxupquote{PotentialFileReader}}}}}
are available. They allow reading the file as individual lines for which
they can return a tokenizer class (see below) for parsing the line.  Or
they can return blocks of numbers as a vector directly.  The
documentation on {\hyperref[\detokenize{Developer_utils:file-reader-classes}]{\sphinxcrossref{\DUrole{std,std-ref}{File reader classes}}}}
contains an example for a typical case.

\sphinxAtStartPar
When reading per\sphinxhyphen{}atom data, the data on each line of the file usually
needs to include an atom ID so it can be associated with a particular
atom.  In that case the data can be read in multi\sphinxhyphen{}line chunks and
broadcast to all MPI ranks with
{\hyperref[\detokenize{Developer_utils:_CPPv4N9LAMMPS_NS5utils20read_lines_from_fileEP4FILEiiPci8MPI_Comm}]{\sphinxcrossref{\sphinxcode{\sphinxupquote{utils::read\_lines\_from\_file()}}}}}.  Those chunks are then
split into lines, parsed, and applied only to atoms the MPI rank
“owns”.

\sphinxAtStartPar
For splitting a string (incrementally) into words and optionally
converting those to numbers, the {\hyperref[\detokenize{Developer_utils:_CPPv4N9LAMMPS_NS9TokenizerE}]{\sphinxcrossref{\sphinxcode{\sphinxupquote{Tokenizer}}}}} and {\hyperref[\detokenize{Developer_utils:_CPPv4N9LAMMPS_NS14ValueTokenizerE}]{\sphinxcrossref{\sphinxcode{\sphinxupquote{ValueTokenizer}}}}} can be used.  Those provide a superset of
the functionality of \sphinxcode{\sphinxupquote{strtok()}} from the C\sphinxhyphen{}library and the latter
also includes conversion to different types.  Any errors while
processing the string in those classes will result in an exception,
which can be caught and the error processed as needed.  Unlike the
C\sphinxhyphen{}library functions \sphinxcode{\sphinxupquote{atoi()}}, \sphinxcode{\sphinxupquote{atof()}}, \sphinxcode{\sphinxupquote{strtol()}}, or
\sphinxcode{\sphinxupquote{strtod()}} the conversion will check if the converted text is a
valid integer or floating point number and will not silently return an
unexpected or incorrect value.  For example, \sphinxcode{\sphinxupquote{atoi()}} will return 12
when converting “12.5”, while the ValueTokenizer class will throw an
{\hyperref[\detokenize{Developer_utils:_CPPv4N9LAMMPS_NS23InvalidIntegerExceptionE}]{\sphinxcrossref{\sphinxcode{\sphinxupquote{InvalidIntegerException}}}}} if
{\hyperref[\detokenize{Developer_utils:_CPPv4N9LAMMPS_NS14ValueTokenizer8next_intEv}]{\sphinxcrossref{\sphinxcode{\sphinxupquote{ValueTokenizer::next\_int()}}}}} is called on the same string.


\subsection{Requesting and accessing neighbor lists}
\label{\detokenize{Developer_notes:requesting-and-accessing-neighbor-lists}}\label{\detokenize{Developer_notes:request-neighbor-list}}
\sphinxAtStartPar
LAMMPS uses Verlet\sphinxhyphen{}style neighbor lists to avoid having to loop over
\sphinxstyleemphasis{all} pairs of \sphinxstyleemphasis{all} atoms when computing pairwise properties with a
cutoff (e.g. pairwise forces or radial distribution functions).  There
are three main algorithms that can be selected by the {\hyperref[\detokenize{neighbor::doc}]{\sphinxcrossref{\DUrole{doc}{neighbor
command}}}}: \sphinxtitleref{bin} (the default, uses binning to achieve linear
scaling with system size), \sphinxtitleref{nsq} (without binning, quadratic scaling),
\sphinxtitleref{multi} (with binning, optimized for varying cutoffs or polydisperse
granular particles).  In addition to how the neighbor lists are
constructed a number of different variants of neighbor lists need to be
created (e.g. “full” or “half”) for different purposes and styles and
those may be required in every time step (“perpetual”) or on some steps
(“occasional”).

\sphinxAtStartPar
The neighbor list creation is managed by the \sphinxcode{\sphinxupquote{Neighbor}} class.
Individual classes can obtain a neighbor list by creating an instance of
a \sphinxcode{\sphinxupquote{NeighRequest}} class which is stored in a list inside the
\sphinxcode{\sphinxupquote{Neighbor}} class.  The \sphinxcode{\sphinxupquote{Neighbor}} class will then analyze the
various requests and apply optimizations where neighbor lists that have
the same settings will be created only once and then copied, or a list
may be constructed by processing a neighbor list from a different
request that is a superset of the requested list.  The neighbor list
build is then {\hyperref[\detokenize{Developer_par_neigh::doc}]{\sphinxcrossref{\DUrole{doc}{processed in parallel}}}}.

\sphinxAtStartPar
The most commonly required neighbor list is a so\sphinxhyphen{}called “half” neighbor
list, where each pair of atoms is listed only once (except when the
{\hyperref[\detokenize{newton::doc}]{\sphinxcrossref{\DUrole{doc}{newton command setting}}}} for pair is off; in that case
pairs straddling subdomains or periodic boundaries will be listed twice).
Thus these are the default settings when a neighbor list request is created in:

\begin{sphinxVerbatim}[commandchars=\\\{\}]
\PYG{k+kt}{void}\PYG{+w}{ }\PYG{n+nf}{Pair::init\PYGZus{}style}\PYG{p}{(}\PYG{p}{)}
\PYG{p}{\PYGZob{}}
\PYG{+w}{  }\PYG{n}{neighbor}\PYG{o}{\PYGZhy{}}\PYG{o}{\PYGZgt{}}\PYG{n}{add\PYGZus{}request}\PYG{p}{(}\PYG{k}{this}\PYG{p}{)}\PYG{p}{;}
\PYG{p}{\PYGZcb{}}

\PYG{k+kt}{void}\PYG{+w}{ }\PYG{n+nf}{Pair::init\PYGZus{}list}\PYG{p}{(}\PYG{k+kt}{int}\PYG{+w}{ }\PYG{c+cm}{/*id*/}\PYG{p}{,}\PYG{+w}{ }\PYG{n}{NeighList}\PYG{+w}{ }\PYG{o}{*}\PYG{n}{ptr}\PYG{p}{)}
\PYG{p}{\PYGZob{}}
\PYG{+w}{  }\PYG{n}{list}\PYG{+w}{ }\PYG{o}{=}\PYG{+w}{ }\PYG{n}{ptr}\PYG{p}{;}
\PYG{p}{\PYGZcb{}}
\end{sphinxVerbatim}

\sphinxAtStartPar
The \sphinxcode{\sphinxupquote{this}} pointer argument is required so the neighbor list code can
access the requesting class instance to store the assembled neighbor
list with that instance by calling its \sphinxcode{\sphinxupquote{init\_list()}} member function.
The optional second argument (omitted here) contains a bitmask of flags
that determines the kind of neighbor list requested.  The default value
used here asks for a perpetual “half” neighbor list.

\sphinxAtStartPar
Non\sphinxhyphen{}default values of the second argument need to be used to adjust a
neighbor list request to the specific needs of a style.  The {\hyperref[\detokenize{pair_tersoff::doc}]{\sphinxcrossref{\DUrole{doc}{tersoff}}}} pair style, for example, needs a “full” neighbor list:

\begin{sphinxVerbatim}[commandchars=\\\{\}]
\PYG{k+kt}{void}\PYG{+w}{ }\PYG{n+nf}{PairTersoff::init\PYGZus{}style}\PYG{p}{(}\PYG{p}{)}
\PYG{p}{\PYGZob{}}
\PYG{+w}{  }\PYG{c+c1}{// [...]}
\PYG{+w}{  }\PYG{n}{neighbor}\PYG{o}{\PYGZhy{}}\PYG{o}{\PYGZgt{}}\PYG{n}{add\PYGZus{}request}\PYG{p}{(}\PYG{k}{this}\PYG{p}{,}\PYG{+w}{ }\PYG{n}{NeighConst}\PYG{o}{:}\PYG{o}{:}\PYG{n}{REQ\PYGZus{}FULL}\PYG{p}{)}\PYG{p}{;}
\PYG{p}{\PYGZcb{}}
\end{sphinxVerbatim}

\sphinxAtStartPar
When a pair style supports r\sphinxhyphen{}RESPA time integration with different cutoff regions,
the request flag may depend on the corresponding r\sphinxhyphen{}RESPA settings. Here is an
example from pair style lj/cut:

\begin{sphinxVerbatim}[commandchars=\\\{\}]
\PYG{k+kt}{void}\PYG{+w}{ }\PYG{n+nf}{PairLJCut::init\PYGZus{}style}\PYG{p}{(}\PYG{p}{)}
\PYG{p}{\PYGZob{}}
\PYG{+w}{  }\PYG{k+kt}{int}\PYG{+w}{ }\PYG{n}{list\PYGZus{}style}\PYG{+w}{ }\PYG{o}{=}\PYG{+w}{ }\PYG{n}{NeighConst}\PYG{o}{:}\PYG{o}{:}\PYG{n}{REQ\PYGZus{}DEFAULT}\PYG{p}{;}

\PYG{+w}{  }\PYG{k}{if}\PYG{+w}{ }\PYG{p}{(}\PYG{n}{update}\PYG{o}{\PYGZhy{}}\PYG{o}{\PYGZgt{}}\PYG{n}{whichflag}\PYG{+w}{ }\PYG{o}{=}\PYG{o}{=}\PYG{+w}{ }\PYG{l+m+mi}{1}\PYG{+w}{ }\PYG{o}{\PYGZam{}}\PYG{o}{\PYGZam{}}\PYG{+w}{ }\PYG{n}{utils}\PYG{o}{:}\PYG{o}{:}\PYG{n}{strmatch}\PYG{p}{(}\PYG{n}{update}\PYG{o}{\PYGZhy{}}\PYG{o}{\PYGZgt{}}\PYG{n}{integrate\PYGZus{}style}\PYG{p}{,}\PYG{+w}{ }\PYG{l+s}{\PYGZdq{}}\PYG{l+s}{\PYGZca{}respa}\PYG{l+s}{\PYGZdq{}}\PYG{p}{)}\PYG{p}{)}\PYG{+w}{ }\PYG{p}{\PYGZob{}}
\PYG{+w}{    }\PYG{k}{auto}\PYG{+w}{ }\PYG{n}{respa}\PYG{+w}{ }\PYG{o}{=}\PYG{+w}{ }\PYG{p}{(}\PYG{n}{Respa}\PYG{+w}{ }\PYG{o}{*}\PYG{p}{)}\PYG{+w}{ }\PYG{n}{update}\PYG{o}{\PYGZhy{}}\PYG{o}{\PYGZgt{}}\PYG{n}{integrate}\PYG{p}{;}
\PYG{+w}{    }\PYG{k}{if}\PYG{+w}{ }\PYG{p}{(}\PYG{n}{respa}\PYG{o}{\PYGZhy{}}\PYG{o}{\PYGZgt{}}\PYG{n}{level\PYGZus{}inner}\PYG{+w}{ }\PYG{o}{\PYGZgt{}}\PYG{o}{=}\PYG{+w}{ }\PYG{l+m+mi}{0}\PYG{p}{)}\PYG{+w}{ }\PYG{n}{list\PYGZus{}style}\PYG{+w}{ }\PYG{o}{=}\PYG{+w}{ }\PYG{n}{NeighConst}\PYG{o}{:}\PYG{o}{:}\PYG{n}{REQ\PYGZus{}RESPA\PYGZus{}INOUT}\PYG{p}{;}
\PYG{+w}{    }\PYG{k}{if}\PYG{+w}{ }\PYG{p}{(}\PYG{n}{respa}\PYG{o}{\PYGZhy{}}\PYG{o}{\PYGZgt{}}\PYG{n}{level\PYGZus{}middle}\PYG{+w}{ }\PYG{o}{\PYGZgt{}}\PYG{o}{=}\PYG{+w}{ }\PYG{l+m+mi}{0}\PYG{p}{)}\PYG{+w}{ }\PYG{n}{list\PYGZus{}style}\PYG{+w}{ }\PYG{o}{=}\PYG{+w}{ }\PYG{n}{NeighConst}\PYG{o}{:}\PYG{o}{:}\PYG{n}{REQ\PYGZus{}RESPA\PYGZus{}ALL}\PYG{p}{;}
\PYG{+w}{  }\PYG{p}{\PYGZcb{}}
\PYG{+w}{  }\PYG{n}{neighbor}\PYG{o}{\PYGZhy{}}\PYG{o}{\PYGZgt{}}\PYG{n}{add\PYGZus{}request}\PYG{p}{(}\PYG{k}{this}\PYG{p}{,}\PYG{+w}{ }\PYG{n}{list\PYGZus{}style}\PYG{p}{)}\PYG{p}{;}
\PYG{+w}{  }\PYG{c+c1}{// [...]}
\PYG{p}{\PYGZcb{}}
\end{sphinxVerbatim}

\sphinxAtStartPar
Granular pair styles need neighbor lists based on particle sizes and not cutoff
and also may need to store data across timesteps (“history”).
For example with:

\begin{sphinxVerbatim}[commandchars=\\\{\}]
\PYG{k}{if}\PYG{+w}{ }\PYG{p}{(}\PYG{n}{use\PYGZus{}history}\PYG{p}{)}\PYG{+w}{ }\PYG{n}{neighbor}\PYG{o}{\PYGZhy{}}\PYG{o}{\PYGZgt{}}\PYG{n}{add\PYGZus{}request}\PYG{p}{(}\PYG{k}{this}\PYG{p}{,}\PYG{+w}{ }\PYG{n}{NeighConst}\PYG{o}{:}\PYG{o}{:}\PYG{n}{REQ\PYGZus{}SIZE}\PYG{+w}{ }\PYG{o}{|}\PYG{+w}{ }\PYG{n}{NeighConst}\PYG{o}{:}\PYG{o}{:}\PYG{n}{REQ\PYGZus{}HISTORY}\PYG{p}{)}\PYG{p}{;}
\PYG{k}{else}\PYG{+w}{ }\PYG{n}{neighbor}\PYG{o}{\PYGZhy{}}\PYG{o}{\PYGZgt{}}\PYG{n}{add\PYGZus{}request}\PYG{p}{(}\PYG{k}{this}\PYG{p}{,}\PYG{+w}{ }\PYG{n}{NeighConst}\PYG{o}{:}\PYG{o}{:}\PYG{n}{REQ\PYGZus{}SIZE}\PYG{p}{)}\PYG{p}{;}
\end{sphinxVerbatim}

\sphinxAtStartPar
In case a class would need to make multiple neighbor list requests with different
settings, each request can set an id which is then used in the corresponding
\sphinxcode{\sphinxupquote{init\_list()}} function to assign it to the suitable pointer variable. This is
done for example by the {\hyperref[\detokenize{pair_meam::doc}]{\sphinxcrossref{\DUrole{doc}{pair style meam}}}}:

\begin{sphinxVerbatim}[commandchars=\\\{\}]
\PYG{k+kt}{void}\PYG{+w}{ }\PYG{n+nf}{PairMEAM::init\PYGZus{}style}\PYG{p}{(}\PYG{p}{)}
\PYG{p}{\PYGZob{}}
\PYG{c+c1}{// [...]}
\PYG{+w}{  }\PYG{n}{neighbor}\PYG{o}{\PYGZhy{}}\PYG{o}{\PYGZgt{}}\PYG{n}{add\PYGZus{}request}\PYG{p}{(}\PYG{k}{this}\PYG{p}{,}\PYG{+w}{ }\PYG{n}{NeighConst}\PYG{o}{:}\PYG{o}{:}\PYG{n}{REQ\PYGZus{}FULL}\PYG{p}{)}\PYG{o}{\PYGZhy{}}\PYG{o}{\PYGZgt{}}\PYG{n}{set\PYGZus{}id}\PYG{p}{(}\PYG{l+m+mi}{1}\PYG{p}{)}\PYG{p}{;}
\PYG{+w}{  }\PYG{n}{neighbor}\PYG{o}{\PYGZhy{}}\PYG{o}{\PYGZgt{}}\PYG{n}{add\PYGZus{}request}\PYG{p}{(}\PYG{k}{this}\PYG{p}{)}\PYG{o}{\PYGZhy{}}\PYG{o}{\PYGZgt{}}\PYG{n}{set\PYGZus{}id}\PYG{p}{(}\PYG{l+m+mi}{2}\PYG{p}{)}\PYG{p}{;}
\PYG{p}{\PYGZcb{}}
\PYG{k+kt}{void}\PYG{+w}{ }\PYG{n+nf}{PairMEAM::init\PYGZus{}list}\PYG{p}{(}\PYG{k+kt}{int}\PYG{+w}{ }\PYG{n}{id}\PYG{p}{,}\PYG{+w}{ }\PYG{n}{NeighList}\PYG{+w}{ }\PYG{o}{*}\PYG{n}{ptr}\PYG{p}{)}
\PYG{p}{\PYGZob{}}
\PYG{+w}{  }\PYG{k}{if}\PYG{+w}{ }\PYG{p}{(}\PYG{n}{id}\PYG{+w}{ }\PYG{o}{=}\PYG{o}{=}\PYG{+w}{ }\PYG{l+m+mi}{1}\PYG{p}{)}\PYG{+w}{ }\PYG{n}{listfull}\PYG{+w}{ }\PYG{o}{=}\PYG{+w}{ }\PYG{n}{ptr}\PYG{p}{;}
\PYG{+w}{  }\PYG{k}{else}\PYG{+w}{ }\PYG{k}{if}\PYG{+w}{ }\PYG{p}{(}\PYG{n}{id}\PYG{+w}{ }\PYG{o}{=}\PYG{o}{=}\PYG{+w}{ }\PYG{l+m+mi}{2}\PYG{p}{)}\PYG{+w}{ }\PYG{n}{listhalf}\PYG{+w}{ }\PYG{o}{=}\PYG{+w}{ }\PYG{n}{ptr}\PYG{p}{;}
\PYG{p}{\PYGZcb{}}
\end{sphinxVerbatim}

\sphinxAtStartPar
Fixes may require a neighbor list that is only build occasionally (or
just once) and this can also be indicated by a flag.  As an example here
is the request from the \sphinxcode{\sphinxupquote{FixPeriNeigh}} class which is created
internally by {\hyperref[\detokenize{pair_peri::doc}]{\sphinxcrossref{\DUrole{doc}{Peridynamics pair styles}}}}:

\begin{sphinxVerbatim}[commandchars=\\\{\}]
\PYG{n}{neighbor}\PYG{o}{\PYGZhy{}}\PYG{o}{\PYGZgt{}}\PYG{n}{add\PYGZus{}request}\PYG{p}{(}\PYG{k}{this}\PYG{p}{,}\PYG{+w}{ }\PYG{n}{NeighConst}\PYG{o}{:}\PYG{o}{:}\PYG{n}{REQ\PYGZus{}FULL}\PYG{+w}{ }\PYG{o}{|}\PYG{+w}{ }\PYG{n}{NeighConst}\PYG{o}{:}\PYG{o}{:}\PYG{n}{REQ\PYGZus{}OCCASIONAL}\PYG{p}{)}\PYG{p}{;}
\end{sphinxVerbatim}

\sphinxAtStartPar
It is also possible to request a neighbor list that uses a different cutoff
than what is usually inferred from the pair style settings (largest cutoff of
all pair styles plus neighbor list skin).  The following is used in the
{\hyperref[\detokenize{compute_rdf::doc}]{\sphinxcrossref{\DUrole{doc}{compute rdf}}}} command implementation:

\begin{sphinxVerbatim}[commandchars=\\\{\}]
\PYG{k}{if}\PYG{+w}{ }\PYG{p}{(}\PYG{n}{cutflag}\PYG{p}{)}
\PYG{+w}{  }\PYG{n}{neighbor}\PYG{o}{\PYGZhy{}}\PYG{o}{\PYGZgt{}}\PYG{n}{add\PYGZus{}request}\PYG{p}{(}\PYG{k}{this}\PYG{p}{,}\PYG{+w}{ }\PYG{n}{NeighConst}\PYG{o}{:}\PYG{o}{:}\PYG{n}{REQ\PYGZus{}OCCASIONAL}\PYG{p}{)}\PYG{o}{\PYGZhy{}}\PYG{o}{\PYGZgt{}}\PYG{n}{set\PYGZus{}cutoff}\PYG{p}{(}\PYG{n}{mycutneigh}\PYG{p}{)}\PYG{p}{;}
\PYG{k}{else}
\PYG{+w}{  }\PYG{n}{neighbor}\PYG{o}{\PYGZhy{}}\PYG{o}{\PYGZgt{}}\PYG{n}{add\PYGZus{}request}\PYG{p}{(}\PYG{k}{this}\PYG{p}{,}\PYG{+w}{ }\PYG{n}{NeighConst}\PYG{o}{:}\PYG{o}{:}\PYG{n}{REQ\PYGZus{}OCCASIONAL}\PYG{p}{)}\PYG{p}{;}
\end{sphinxVerbatim}

\sphinxAtStartPar
The neighbor list request function has a slightly different set of arguments
when created by a command style.  In this case the neighbor list is
\sphinxstyleemphasis{always} an occasional neighbor list, so that flag is not needed. However
for printing the neighbor list summary the name of the requesting command
should be set.  Below is the request from the {\hyperref[\detokenize{delete_atoms::doc}]{\sphinxcrossref{\DUrole{doc}{delete atoms}}}}
command:

\begin{sphinxVerbatim}[commandchars=\\\{\}]
\PYG{n}{neighbor}\PYG{o}{\PYGZhy{}}\PYG{o}{\PYGZgt{}}\PYG{n}{add\PYGZus{}request}\PYG{p}{(}\PYG{k}{this}\PYG{p}{,}\PYG{+w}{ }\PYG{l+s}{\PYGZdq{}}\PYG{l+s}{delete\PYGZus{}atoms}\PYG{l+s}{\PYGZdq{}}\PYG{p}{,}\PYG{+w}{ }\PYG{n}{NeighConst}\PYG{o}{:}\PYG{o}{:}\PYG{n}{REQ\PYGZus{}FULL}\PYG{p}{)}\PYG{p}{;}
\end{sphinxVerbatim}


\subsection{Choosing between a custom atom style, fix property/atom, and fix STORE/ATOM}
\label{\detokenize{Developer_notes:choosing-between-a-custom-atom-style-fix-property-atom-and-fix-store-atom}}
\sphinxAtStartPar
There are multiple ways to manage per\sphinxhyphen{}atom data within LAMMPS.  Often
the per\sphinxhyphen{}atom storage is only used locally and managed by the class that
uses it.  If the data has to persist between multiple time steps and
migrate with atoms when they move from sub\sphinxhyphen{}domain to sub\sphinxhyphen{}domain or
across periodic boundaries, then using a custom atom style, or {\hyperref[\detokenize{fix_property_atom::doc}]{\sphinxcrossref{\DUrole{doc}{fix
property/atom}}}}, or the internal fix STORE/ATOM are
possible options.
\begin{itemize}
\item {} 
\sphinxAtStartPar
Using the atom style is usually the most programming effort and mostly
needed when the per\sphinxhyphen{}atom data is an integral part of the model like a
per\sphinxhyphen{}atom charge or diameter and thus should be part of the Atoms
section of a {\hyperref[\detokenize{read_data::doc}]{\sphinxcrossref{\DUrole{doc}{data file}}}}.

\item {} 
\sphinxAtStartPar
Fix property/atom is useful if the data is optional or should be
entered by the user, or accessed as a (named) custom property. In this
case the fix should be entered as part of the input (and not
internally) which allows to enter and store its content with data files.

\item {} 
\sphinxAtStartPar
Fix STORE/ATOM should be used when the data should be accessed internally
only and thus the fix can be created internally.

\end{itemize}


\subsection{Fix contributions to instantaneous energy, virial, and cumulative energy}
\label{\detokenize{Developer_notes:fix-contributions-to-instantaneous-energy-virial-and-cumulative-energy}}
\sphinxAtStartPar
Fixes can calculate contributions to the instantaneous energy and/or
virial of the system, both in a global and peratom sense.  Fixes that
perform thermostatting or barostatting can calculate the cumulative
energy they add to or subtract from the system, which is accessed by
the \sphinxstyleemphasis{ecouple} and \sphinxstyleemphasis{econserve} thermodynamic keywords.  This subsection
explains how both work and what flags to set in a new fix to enable
this functionality.

\sphinxAtStartPar
Let’s start with thermostatting and barostatting fixes.  Examples are
the {\hyperref[\detokenize{fix_langevin::doc}]{\sphinxcrossref{\DUrole{doc}{fix langevin}}}} and {\hyperref[\detokenize{fix_nh::doc}]{\sphinxcrossref{\DUrole{doc}{fix npt}}}}
commands.  Here is what the fix needs to do:
\begin{itemize}
\item {} 
\sphinxAtStartPar
Set the variable \sphinxstyleemphasis{ecouple\_flag} = 1 in the constructor.  Also set
\sphinxstyleemphasis{scalar\_flag} = 1, \sphinxstyleemphasis{extscalar} = 1, and \sphinxstyleemphasis{global\_freq} to a timestep
increment which matches how often the fix is invoked.

\item {} 
\sphinxAtStartPar
Implement a compute\_scalar() method that returns the cumulative
energy added or subtracted by the fix, e.g. by rescaling the
velocity of atoms.  The sign convention is that subtracted energy is
positive, added energy is negative.  This must be the total energy
added to the entire system, i.e. an “extensive” quantity, not a
per\sphinxhyphen{}atom energy.  Cumulative means the summed energy since the fix
was instantiated, even across multiple runs.  This is because the
energy is used by the \sphinxstyleemphasis{econserve} thermodynamic keyword to check
that the fix is conserving the total energy of the system,
i.e. potential energy + kinetic energy + coupling energy = a
constant.

\end{itemize}

\sphinxAtStartPar
And here is how the code operates:
\begin{itemize}
\item {} 
\sphinxAtStartPar
The Modify class makes a list of all fixes that set \sphinxstyleemphasis{ecouple\_flag} = 1.

\item {} 
\sphinxAtStartPar
The {\hyperref[\detokenize{thermo_style::doc}]{\sphinxcrossref{\DUrole{doc}{thermo\_style custom}}}} command defines
\sphinxstyleemphasis{ecouple} and \sphinxstyleemphasis{econserve} keywords.

\item {} 
\sphinxAtStartPar
These keywords sum the energy contributions from all the
\sphinxstyleemphasis{ecouple\_flag} = 1 fixes by invoking the \sphinxstyleemphasis{energy\_couple()} method in
the Modify class, which calls the \sphinxstyleemphasis{compute\_scalar()} method of each
fix in the list.

\end{itemize}


\bigskip\hrule\bigskip


\sphinxAtStartPar
Next, here is how a fix contributes to the instantaneous energy and
virial of the system.  First, it sets any or all of these flags to a
value of 1 in their constructor:
\begin{itemize}
\item {} 
\sphinxAtStartPar
\sphinxstyleemphasis{energy\_global\_flag} to contribute to global energy, example: {\hyperref[\detokenize{fix_indent::doc}]{\sphinxcrossref{\DUrole{doc}{fix indent}}}}

\item {} 
\sphinxAtStartPar
\sphinxstyleemphasis{energy\_peratom\_flag} to contribute to peratom energy, {\hyperref[\detokenize{fix_cmap::doc}]{\sphinxcrossref{\DUrole{doc}{fix cmap}}}}

\item {} 
\sphinxAtStartPar
\sphinxstyleemphasis{virial\_global\_flag} to contribute to global virial, example: {\hyperref[\detokenize{fix_wall::doc}]{\sphinxcrossref{\DUrole{doc}{fix wall}}}}

\item {} 
\sphinxAtStartPar
\sphinxstyleemphasis{virial\_peratom\_flag} to contribute to peratom virial, example: {\hyperref[\detokenize{fix_wall::doc}]{\sphinxcrossref{\DUrole{doc}{fix wall}}}}

\end{itemize}

\sphinxAtStartPar
The fix must also do the following:
\begin{itemize}
\item {} 
\sphinxAtStartPar
For global energy, implement a compute\_scalar() method that returns
the energy added or subtracted on this timestep.  Here the sign
convention is that added energy is positive, subtracted energy is
negative.

\item {} 
\sphinxAtStartPar
For peratom energy, invoke the ev\_init(eflag,vflag) function each
time the fix is invoked, which initializes per\sphinxhyphen{}atom energy storage.
The value of eflag may need to be stored from an earlier call to the
fix during the same timestep.  See how the {\hyperref[\detokenize{fix_cmap::doc}]{\sphinxcrossref{\DUrole{doc}{fix cmap}}}} command does this in src/MOLECULE/fix\_cmap.cpp.  When an
energy for one or more atoms is calculated, invoke the ev\_tally()
function to tally the contribution to each atom.  Both the ev\_init()
and ev\_tally() methods are in the parent Fix class.

\item {} 
\sphinxAtStartPar
For global and/or peratom virial, invoke the v\_init(vflag) function
each time the fix is invoked, which initializes virial storage.
When forces on one or more atoms are calculated, invoke the
v\_tally() function to tally the contribution.  Both the v\_init() and
v\_tally() methods are in the parent Fix class.  Note that there are
several variants of v\_tally(); choose the one appropriate to your
fix.

\end{itemize}

\begin{sphinxadmonition}{note}{Note:}
\sphinxAtStartPar
The ev\_init() and ev\_tally() methods also account for global and
peratom virial contributions.  Thus you do not need to invoke the
v\_init() and v\_tally() methods if the fix also calculates peratom
energies.
\end{sphinxadmonition}

\sphinxAtStartPar
The fix must also specify whether (by default) to include or exclude
these contributions to the global/peratom energy/virial of the system.
For the fix to include the contributions, set either or both of these
variables in the constructor:
\begin{itemize}
\item {} 
\sphinxAtStartPar
\sphinxstyleemphasis{thermo\_energy} = 1, for global and peratom energy

\item {} 
\sphinxAtStartPar
\sphinxstyleemphasis{thermo\_virial} = 1, for global and peratom virial

\end{itemize}

\sphinxAtStartPar
Note that these variables are zeroed in fix.cpp.  Thus if you don’t
set the variables, the contributions will be excluded (by default).

\sphinxAtStartPar
However, the user has ultimate control over whether to include or
exclude the contributions of the fix via the {\hyperref[\detokenize{fix_modify::doc}]{\sphinxcrossref{\DUrole{doc}{fix modify}}}} command:
\begin{itemize}
\item {} 
\sphinxAtStartPar
fix modify \sphinxstyleemphasis{energy yes} to include global and peratom energy contributions

\item {} 
\sphinxAtStartPar
fix modify \sphinxstyleemphasis{virial yes} to include global and peratom virial contributions

\end{itemize}

\sphinxAtStartPar
If the fix contributes to any of the global/peratom energy/virial
values for the system, it should be explained on the fix doc page,
along with the default values for the \sphinxstyleemphasis{energy yes/no} and \sphinxstyleemphasis{virial
yes/no} settings of the {\hyperref[\detokenize{fix_modify::doc}]{\sphinxcrossref{\DUrole{doc}{fix modify}}}} command.

\sphinxAtStartPar
Finally, these 4 contributions are included in the output of 4
computes:
\begin{itemize}
\item {} 
\sphinxAtStartPar
global energy in {\hyperref[\detokenize{compute_pe::doc}]{\sphinxcrossref{\DUrole{doc}{compute pe}}}}

\item {} 
\sphinxAtStartPar
peratom energy in {\hyperref[\detokenize{compute_pe_atom::doc}]{\sphinxcrossref{\DUrole{doc}{compute pe/atom}}}}

\item {} 
\sphinxAtStartPar
global virial in {\hyperref[\detokenize{compute_pressure::doc}]{\sphinxcrossref{\DUrole{doc}{compute pressure}}}}

\item {} 
\sphinxAtStartPar
peratom virial in {\hyperref[\detokenize{compute_stress_atom::doc}]{\sphinxcrossref{\DUrole{doc}{compute stress/atom}}}}

\end{itemize}

\sphinxAtStartPar
These computes invoke a method of the Modify class to include
contributions from fixes that have the corresponding flags set,
e.g. \sphinxstyleemphasis{energy\_peratom\_flag} and \sphinxstyleemphasis{thermo\_energy} for {\hyperref[\detokenize{compute_pe_atom::doc}]{\sphinxcrossref{\DUrole{doc}{compute
pe/atom}}}}.

\sphinxAtStartPar
Note that each compute has an optional keyword to either include or
exclude all contributions from fixes.  Also note that {\hyperref[\detokenize{compute_pe::doc}]{\sphinxcrossref{\DUrole{doc}{compute pe}}}} and {\hyperref[\detokenize{compute_pressure::doc}]{\sphinxcrossref{\DUrole{doc}{compute pressure}}}} are what
is used (by default) by {\hyperref[\detokenize{thermo_style::doc}]{\sphinxcrossref{\DUrole{doc}{thermodynamic output}}}} to
calculate values for its \sphinxstyleemphasis{pe} and \sphinxstyleemphasis{press} keywords.


\subsection{KSpace PPPM FFT grids}
\label{\detokenize{Developer_notes:kspace-pppm-fft-grids}}
\sphinxAtStartPar
The various {\hyperref[\detokenize{kspace_style::doc}]{\sphinxcrossref{\DUrole{doc}{KSpace PPPM}}}} styles in LAMMPS use
FFTs to solve Poisson’s equation.  This subsection describes:
\begin{itemize}
\item {} 
\sphinxAtStartPar
how FFT grids are defined

\item {} 
\sphinxAtStartPar
how they are decomposed across processors

\item {} 
\sphinxAtStartPar
how they are indexed by each processor

\item {} 
\sphinxAtStartPar
how particle charge and electric field values are mapped to/from
the grid

\end{itemize}

\sphinxAtStartPar
An FFT grid cell is a 3d volume; grid points are corners of a grid
cell and the code stores values assigned to grid points in vectors or
3d arrays.  A global 3d FFT grid has points indexed 0 to N\sphinxhyphen{}1 inclusive
in each dimension.

\sphinxAtStartPar
Each processor owns two subsets of the grid, each subset is
brick\sphinxhyphen{}shaped.  Depending on how it is used, these subsets are
allocated as a 1d vector or 3d array.  Either way, the ordering of
values within contiguous memory x fastest, then y, z slowest.

\sphinxAtStartPar
For the \sphinxcode{\sphinxupquote{3d decomposition}} of the grid, the global grid is
partitioned into bricks that correspond to the subdomains of the
simulation box that each processor owns.  Often, this is a regular 3d
array (Px by Py by Pz) of bricks, where P = number of processors =
Px * Py * Pz.  More generally it can be a tiled decomposition, where
each processor owns a brick and the union of all the bricks is the
global grid.  Tiled decompositions are produced by load balancing with
the RCB algorithm; see the {\hyperref[\detokenize{balance::doc}]{\sphinxcrossref{\DUrole{doc}{balance rcb}}}} command.

\sphinxAtStartPar
For the \sphinxcode{\sphinxupquote{FFT decompostion}} of the grid, each processor owns a brick
that spans the entire x dimension of the grid while the y and z
dimensions are partitioned as a regular 2d array (P1 by P2), where P =
P1 * P2.

\sphinxAtStartPar
The following indices store the inclusive bounds of the brick a
processor owns, within the global grid:

\begin{sphinxVerbatim}[commandchars=\\\{\}]
nFOO\PYGZus{}in = 3d decomposition brick
nFOO\PYGZus{}fft = FFT decomposition brick
nFOO\PYGZus{}out = 3d decomposition brick + ghost cells
\end{sphinxVerbatim}

\sphinxAtStartPar
where \sphinxcode{\sphinxupquote{FOO}} corresponds to \sphinxcode{\sphinxupquote{xlo, xhi, ylo, yhi, zlo,}} or \sphinxcode{\sphinxupquote{zhi}}.

\sphinxAtStartPar
The \sphinxcode{\sphinxupquote{in}} and \sphinxcode{\sphinxupquote{fft}} indices are from 0 to N\sphinxhyphen{}1 inclusive in each
dimension, where N is the grid size.

\sphinxAtStartPar
The \sphinxcode{\sphinxupquote{out}} indices index an array which stores the \sphinxcode{\sphinxupquote{in}} subset of
the grid plus ghost cells that surround it.  These indices can thus be
\textless{} 0 or \textgreater{}= N.

\sphinxAtStartPar
The number of ghost cells a processor owns in each of the 6 directions
is a function of:

\begin{sphinxVerbatim}[commandchars=\\\{\}]
neighbor skin distance (since atoms can move outside a proc subdomain)
qdist = offset or charge from atom due to TIP4P fictitious charge
order = mapping stencil size
shift = factor used when order is an even number (see below)
\end{sphinxVerbatim}

\sphinxAtStartPar
Here is an explanation of how the PPPM variables \sphinxcode{\sphinxupquote{order}},
\sphinxcode{\sphinxupquote{nlower}} / \sphinxcode{\sphinxupquote{nupper}}, \sphinxcode{\sphinxupquote{shift}}, and \sphinxcode{\sphinxupquote{OFFSET}} work. They are the
relevant variables that determine how atom charge is mapped to grid
points and how field values are mapped from grid points to atoms:

\begin{sphinxVerbatim}[commandchars=\\\{\}]
order = \PYGZsh{} of nearby grid points in each dim that atom charge/field are mapped to/from
nlower,nupper = extent of stencil around the grid point an atom is assigned to
OFFSET = large integer added/subtracted when mapping to avoid int(\PYGZhy{}0.75) = 0 when \PYGZhy{}1 is the desired result
\end{sphinxVerbatim}

\sphinxAtStartPar
The particle\_map() method assigns each atom to a grid point.

\sphinxAtStartPar
If order is even, say 4:

\begin{sphinxVerbatim}[commandchars=\\\{\}]
atom is assigned to grid point to its left (in each dim)
shift = OFFSET
nlower = \PYGZhy{}1, nupper = 2, which are offsets from assigned grid point
window of mapping grid pts is thus 2 grid points to left of atom, 2 to right
\end{sphinxVerbatim}

\sphinxAtStartPar
If order is odd, say 5:

\begin{sphinxVerbatim}[commandchars=\\\{\}]
atom is assigned to left/right grid pt it is closest to (in each dim)
shift = OFFSET + 0.5
nlower = 2, nupper = 2
if point is in left half of cell, then window of affected grid pts is 3 grid points to left of atom, 2 to right
if point is in right half of cell, then window of affected grid pts is 2 grid points to left of atom, 3 to right
\end{sphinxVerbatim}

\sphinxAtStartPar
These settings apply to each dimension, so that if order = 5, an
atom’s charge is mapped to 125 grid points that surround the atom.

\sphinxstepscope


\section{Notes for updating code written for older LAMMPS versions}
\label{\detokenize{Developer_updating:notes-for-updating-code-written-for-older-lammps-versions}}\label{\detokenize{Developer_updating::doc}}
\sphinxAtStartPar
This section documents how C++ source files that are available \sphinxstyleemphasis{outside
of the LAMMPS source distribution} (e.g. in external USER packages or as
source files provided as a supplement to a publication) that are written
for an older version of LAMMPS and thus need to be updated to be
compatible with the current version of LAMMPS.  Due to the active
development of LAMMPS it is likely to always be incomplete.  Please
contact \sphinxhref{mailto:developers@lammps.org}{developers@lammps.org} in case you run across an issue that is not
(yet) listed here.  Please also review the latest information about the
LAMMPS {\hyperref[\detokenize{Modify_style::doc}]{\sphinxcrossref{\DUrole{doc}{programming style conventions}}}}, especially
if you are considering to submit the updated version for inclusion into
the LAMMPS distribution.

\sphinxAtStartPar
Available topics in mostly chronological order are:
\begin{itemize}
\item {} 
\sphinxAtStartPar
{\hyperref[\detokenize{Developer_updating:setting-flags-in-the-constructor}]{\sphinxcrossref{Setting flags in the constructor}}}

\item {} 
\sphinxAtStartPar
{\hyperref[\detokenize{Developer_updating:rename-of-pack-unpack-comm-to-pack-unpack-forward-comm}]{\sphinxcrossref{Rename of pack/unpack\_comm() to pack/unpack\_forward\_comm()}}}

\item {} 
\sphinxAtStartPar
{\hyperref[\detokenize{Developer_updating:use-ev-init-to-initialize-variables-derived-from-eflag-and-vflag}]{\sphinxcrossref{Use ev\_init() to initialize variables derived from eflag and vflag}}}

\item {} 
\sphinxAtStartPar
{\hyperref[\detokenize{Developer_updating:use-utils-count-words-functions-instead-of-atom-count-words}]{\sphinxcrossref{Use utils::count\_words() functions instead of atom\sphinxhyphen{}\textgreater{}count\_words()}}}

\item {} 
\sphinxAtStartPar
{\hyperref[\detokenize{Developer_updating:use-utils-numeric-functions-instead-of-force-numeric}]{\sphinxcrossref{Use utils::numeric() functions instead of force\sphinxhyphen{}\textgreater{}numeric()}}}

\item {} 
\sphinxAtStartPar
{\hyperref[\detokenize{Developer_updating:use-utils-open-potential-function-to-open-potential-files}]{\sphinxcrossref{Use utils::open\_potential() function to open potential files}}}

\item {} 
\sphinxAtStartPar
{\hyperref[\detokenize{Developer_updating:use-symbolic-atom-and-atomvec-constants-instead-of-numerical-values}]{\sphinxcrossref{Use symbolic Atom and AtomVec constants instead of numerical values}}}

\item {} 
\sphinxAtStartPar
{\hyperref[\detokenize{Developer_updating:simplify-customized-error-messages}]{\sphinxcrossref{Simplify customized error messages}}}

\item {} 
\sphinxAtStartPar
{\hyperref[\detokenize{Developer_updating:use-of-override-instead-of-virtual}]{\sphinxcrossref{Use of “override” instead of “virtual”}}}

\item {} 
\sphinxAtStartPar
{\hyperref[\detokenize{Developer_updating:simplified-and-more-compact-neighbor-list-requests}]{\sphinxcrossref{Simplified and more compact neighbor list requests}}}

\item {} 
\sphinxAtStartPar
{\hyperref[\detokenize{Developer_updating:split-of-fix-store-into-fix-store-global-and-fix-store-peratom}]{\sphinxcrossref{Split of fix STORE into fix STORE/GLOBAL and fix STORE/PERATOM}}}

\item {} 
\sphinxAtStartPar
{\hyperref[\detokenize{Developer_updating:rename-of-fix-store-peratom-to-fix-store-atom-and-change-of-arguments}]{\sphinxcrossref{Rename of fix STORE/PERATOM to fix STORE/ATOM and change of arguments}}}

\item {} 
\sphinxAtStartPar
{\hyperref[\detokenize{Developer_updating:use-output-get-dump-by-id-instead-of-output-find-dump}]{\sphinxcrossref{Use Output::get\_dump\_by\_id() instead of Output::find\_dump()}}}

\item {} 
\sphinxAtStartPar
{\hyperref[\detokenize{Developer_updating:refactored-grid-communication-using-grid3d-grid2d-classes-instead-of-gridcomm}]{\sphinxcrossref{Refactored grid communication using Grid3d/Grid2d classes instead of GridComm}}}

\end{itemize}


\bigskip\hrule\bigskip



\subsection{Setting flags in the constructor}
\label{\detokenize{Developer_updating:setting-flags-in-the-constructor}}
\sphinxAtStartPar
As LAMMPS gains additional functionality, new flags may need to be set
in the constructor or a class to signal compatibility with such features.
Most of the time the defaults are chosen conservatively, but sometimes
the conservative choice is the uncommon choice, and then those settings
need to be made when updating code.

\sphinxAtStartPar
Pair styles:
\begin{itemize}
\item {} 
\sphinxAtStartPar
\sphinxcode{\sphinxupquote{manybody\_flag}}: set to 1 if your pair style is not pair\sphinxhyphen{}wise additive

\item {} 
\sphinxAtStartPar
\sphinxcode{\sphinxupquote{restartinfo}}: set to 0 if your pair style does not store data in restart files

\end{itemize}


\subsection{Rename of pack/unpack\_comm() to pack/unpack\_forward\_comm()}
\label{\detokenize{Developer_updating:rename-of-pack-unpack-comm-to-pack-unpack-forward-comm}}
\sphinxAtStartPar
\DUrole{versionmodified,changed}{Changed in version 8Aug2014.}

\sphinxAtStartPar
In this change set, the functions to pack/unpack data into communication buffers
for {\hyperref[\detokenize{Developer_comm_ops::doc}]{\sphinxcrossref{\DUrole{doc}{forward communications}}}} were renamed from
\sphinxcode{\sphinxupquote{pack\_comm()}} and \sphinxcode{\sphinxupquote{unpack\_comm()}} to \sphinxcode{\sphinxupquote{pack\_forward\_comm()}} and
\sphinxcode{\sphinxupquote{unpack\_forward\_comm()}}, respectively.  Also the meaning of the return
value of these functions was changed: rather than returning the number
of items per atom stored in the buffer, now the total number of items
added (or unpacked) needs to be returned.  Here is an example from the
\sphinxtitleref{PairEAM} class.  Of course the member function declaration in corresponding
header file needs to be updated accordingly.

\sphinxAtStartPar
Old:

\begin{sphinxVerbatim}[commandchars=\\\{\}]
\PYG{k+kt}{int}\PYG{+w}{ }\PYG{n+nf}{PairEAM::pack\PYGZus{}comm}\PYG{p}{(}\PYG{k+kt}{int}\PYG{+w}{ }\PYG{n}{n}\PYG{p}{,}\PYG{+w}{ }\PYG{k+kt}{int}\PYG{+w}{ }\PYG{o}{*}\PYG{n}{list}\PYG{p}{,}\PYG{+w}{ }\PYG{k+kt}{double}\PYG{+w}{ }\PYG{o}{*}\PYG{n}{buf}\PYG{p}{,}\PYG{+w}{ }\PYG{k+kt}{int}\PYG{+w}{ }\PYG{n}{pbc\PYGZus{}flag}\PYG{p}{,}\PYG{+w}{ }\PYG{k+kt}{int}\PYG{+w}{ }\PYG{o}{*}\PYG{n}{pbc}\PYG{p}{)}
\PYG{p}{\PYGZob{}}
\PYG{+w}{  }\PYG{k+kt}{int}\PYG{+w}{ }\PYG{n}{m}\PYG{+w}{ }\PYG{o}{=}\PYG{+w}{ }\PYG{l+m+mi}{0}\PYG{p}{;}
\PYG{+w}{  }\PYG{k}{for}\PYG{+w}{ }\PYG{p}{(}\PYG{k+kt}{int}\PYG{+w}{ }\PYG{n}{i}\PYG{+w}{ }\PYG{o}{=}\PYG{+w}{ }\PYG{l+m+mi}{0}\PYG{p}{;}\PYG{+w}{ }\PYG{n}{i}\PYG{+w}{ }\PYG{o}{\PYGZlt{}}\PYG{+w}{ }\PYG{n}{n}\PYG{p}{;}\PYG{+w}{ }\PYG{n}{i}\PYG{o}{+}\PYG{o}{+}\PYG{p}{)}\PYG{+w}{ }\PYG{p}{\PYGZob{}}
\PYG{+w}{    }\PYG{k+kt}{int}\PYG{+w}{ }\PYG{n}{j}\PYG{+w}{ }\PYG{o}{=}\PYG{+w}{ }\PYG{n}{list}\PYG{p}{[}\PYG{n}{i}\PYG{p}{]}\PYG{p}{;}
\PYG{+w}{    }\PYG{n}{buf}\PYG{p}{[}\PYG{n}{m}\PYG{o}{+}\PYG{o}{+}\PYG{p}{]}\PYG{+w}{ }\PYG{o}{=}\PYG{+w}{ }\PYG{n}{fp}\PYG{p}{[}\PYG{n}{j}\PYG{p}{]}\PYG{p}{;}
\PYG{+w}{  }\PYG{p}{\PYGZcb{}}
\PYG{+w}{  }\PYG{k}{return}\PYG{+w}{ }\PYG{l+m+mi}{1}\PYG{p}{;}
\PYG{p}{\PYGZcb{}}
\end{sphinxVerbatim}

\sphinxAtStartPar
New:

\begin{sphinxVerbatim}[commandchars=\\\{\}]
\PYG{k+kt}{int}\PYG{+w}{ }\PYG{n+nf}{PairEAM::pack\PYGZus{}forward\PYGZus{}comm}\PYG{p}{(}\PYG{k+kt}{int}\PYG{+w}{ }\PYG{n}{n}\PYG{p}{,}\PYG{+w}{ }\PYG{k+kt}{int}\PYG{+w}{ }\PYG{o}{*}\PYG{n}{list}\PYG{p}{,}\PYG{+w}{ }\PYG{k+kt}{double}\PYG{+w}{ }\PYG{o}{*}\PYG{n}{buf}\PYG{p}{,}\PYG{+w}{ }\PYG{k+kt}{int}\PYG{+w}{ }\PYG{n}{pbc\PYGZus{}flag}\PYG{p}{,}\PYG{+w}{ }\PYG{k+kt}{int}\PYG{+w}{ }\PYG{o}{*}\PYG{n}{pbc}\PYG{p}{)}
\PYG{p}{\PYGZob{}}
\PYG{+w}{  }\PYG{k+kt}{int}\PYG{+w}{ }\PYG{n}{m}\PYG{+w}{ }\PYG{o}{=}\PYG{+w}{ }\PYG{l+m+mi}{0}\PYG{p}{;}
\PYG{+w}{  }\PYG{k}{for}\PYG{+w}{ }\PYG{p}{(}\PYG{k+kt}{int}\PYG{+w}{ }\PYG{n}{i}\PYG{+w}{ }\PYG{o}{=}\PYG{+w}{ }\PYG{l+m+mi}{0}\PYG{p}{;}\PYG{+w}{ }\PYG{n}{i}\PYG{+w}{ }\PYG{o}{\PYGZlt{}}\PYG{+w}{ }\PYG{n}{n}\PYG{p}{;}\PYG{+w}{ }\PYG{n}{i}\PYG{o}{+}\PYG{o}{+}\PYG{p}{)}\PYG{+w}{ }\PYG{p}{\PYGZob{}}
\PYG{+w}{    }\PYG{k+kt}{int}\PYG{+w}{ }\PYG{n}{j}\PYG{+w}{ }\PYG{o}{=}\PYG{+w}{ }\PYG{n}{list}\PYG{p}{[}\PYG{n}{i}\PYG{p}{]}\PYG{p}{;}
\PYG{+w}{    }\PYG{n}{buf}\PYG{p}{[}\PYG{n}{m}\PYG{o}{+}\PYG{o}{+}\PYG{p}{]}\PYG{+w}{ }\PYG{o}{=}\PYG{+w}{ }\PYG{n}{fp}\PYG{p}{[}\PYG{n}{j}\PYG{p}{]}\PYG{p}{;}
\PYG{+w}{  }\PYG{p}{\PYGZcb{}}
\PYG{+w}{  }\PYG{k}{return}\PYG{+w}{ }\PYG{n}{m}\PYG{p}{;}
\PYG{p}{\PYGZcb{}}
\end{sphinxVerbatim}

\begin{sphinxadmonition}{note}{Note:}
\sphinxAtStartPar
Because the various “pack” and “unpack” functions are defined in the
respective base classes as dummy functions doing nothing, and because
of the the name mismatch the custom versions in the derived class
will no longer be called, there will be no compilation error when
this change is not applied.  Only calculations will suddenly produce
incorrect results because the required forward communication calls
will cease to function correctly.
\end{sphinxadmonition}


\subsection{Use ev\_init() to initialize variables derived from eflag and vflag}
\label{\detokenize{Developer_updating:use-ev-init-to-initialize-variables-derived-from-eflag-and-vflag}}
\sphinxAtStartPar
\DUrole{versionmodified,changed}{Changed in version 29Mar2019.}

\sphinxAtStartPar
There are several variables that need to be initialized based on
the values of the “eflag” and “vflag” variables and since sometimes
there are new bits added and new variables need to be set to 1 or 0.
To make this consistent across all styles, there is now an inline
function \sphinxcode{\sphinxupquote{ev\_init(eflag, vflag)}} that makes those settings
consistently and calls either \sphinxcode{\sphinxupquote{ev\_setup()}} or \sphinxcode{\sphinxupquote{ev\_unset()}}.
Example from a pair style:

\sphinxAtStartPar
Old:

\begin{sphinxVerbatim}[commandchars=\\\{\}]
\PYG{k}{if}\PYG{+w}{ }\PYG{p}{(}\PYG{n}{eflag}\PYG{+w}{ }\PYG{o}{|}\PYG{o}{|}\PYG{+w}{ }\PYG{n}{vflag}\PYG{p}{)}\PYG{+w}{ }\PYG{n}{ev\PYGZus{}setup}\PYG{p}{(}\PYG{n}{eflag}\PYG{p}{,}\PYG{+w}{ }\PYG{n}{vflag}\PYG{p}{)}\PYG{p}{;}
\PYG{k}{else}\PYG{+w}{ }\PYG{n}{evflag}\PYG{+w}{ }\PYG{o}{=}\PYG{+w}{ }\PYG{n}{vflag\PYGZus{}fdotr}\PYG{+w}{ }\PYG{o}{=}\PYG{+w}{ }\PYG{n}{eflag\PYGZus{}global}\PYG{+w}{ }\PYG{o}{=}\PYG{+w}{ }\PYG{n}{eflag\PYGZus{}atom}\PYG{+w}{ }\PYG{o}{=}\PYG{+w}{ }\PYG{l+m+mi}{0}\PYG{p}{;}
\end{sphinxVerbatim}

\sphinxAtStartPar
New:

\begin{sphinxVerbatim}[commandchars=\\\{\}]
\PYG{n}{ev\PYGZus{}init}\PYG{p}{(}\PYG{n}{eflag}\PYG{p}{,}\PYG{+w}{ }\PYG{n}{vflag}\PYG{p}{)}\PYG{p}{;}
\end{sphinxVerbatim}

\sphinxAtStartPar
Not applying this change will not cause a compilation error, but
can lead to inconsistent behavior and incorrect tallying of
energy or virial.


\subsection{Use utils::count\_words() functions instead of atom\sphinxhyphen{}\textgreater{}count\_words()}
\label{\detokenize{Developer_updating:use-utils-count-words-functions-instead-of-atom-count-words}}
\sphinxAtStartPar
\DUrole{versionmodified,changed}{Changed in version 2Jun2020.}

\sphinxAtStartPar
The “count\_words()” functions for parsing text have been moved from the
Atom class to the {\hyperref[\detokenize{Developer_utils::doc}]{\sphinxcrossref{\DUrole{doc}{utils namespace}}}}.  The
“count\_words()” function in “utils” uses the Tokenizer class internally
to split a line into words and count them, thus it will not modify the
argument string as the function in the Atoms class did and thus had a
variant using a copy buffer.  Unlike the old version, the new version
does not remove comments. For that you can use the
{\hyperref[\detokenize{Developer_utils:_CPPv4N9LAMMPS_NS5utils12trim_commentERKNSt6stringE}]{\sphinxcrossref{\sphinxcode{\sphinxupquote{utils::trim\_comment() function}}}}} as shown in the example below.

\sphinxAtStartPar
Old:

\begin{sphinxVerbatim}[commandchars=\\\{\}]
\PYG{n}{nwords}\PYG{+w}{ }\PYG{o}{=}\PYG{+w}{ }\PYG{n}{atom}\PYG{o}{\PYGZhy{}}\PYG{o}{\PYGZgt{}}\PYG{n}{count\PYGZus{}words}\PYG{p}{(}\PYG{n}{line}\PYG{p}{)}\PYG{p}{;}
\PYG{k+kt}{int}\PYG{+w}{ }\PYG{n}{nwords}\PYG{+w}{ }\PYG{o}{=}\PYG{+w}{ }\PYG{n}{atom}\PYG{o}{\PYGZhy{}}\PYG{o}{\PYGZgt{}}\PYG{n}{count\PYGZus{}words}\PYG{p}{(}\PYG{n}{buf}\PYG{p}{)}\PYG{p}{;}
\end{sphinxVerbatim}

\sphinxAtStartPar
New:

\begin{sphinxVerbatim}[commandchars=\\\{\}]
\PYG{n}{nwords}\PYG{+w}{ }\PYG{o}{=}\PYG{+w}{ }\PYG{n}{utils}\PYG{o}{:}\PYG{o}{:}\PYG{n}{count\PYGZus{}words}\PYG{p}{(}\PYG{n}{line}\PYG{p}{)}\PYG{p}{;}
\PYG{k+kt}{int}\PYG{+w}{ }\PYG{n}{nwords}\PYG{+w}{ }\PYG{o}{=}\PYG{+w}{ }\PYG{n}{utils}\PYG{o}{:}\PYG{o}{:}\PYG{n}{count\PYGZus{}words}\PYG{p}{(}\PYG{n}{utils}\PYG{o}{:}\PYG{o}{:}\PYG{n}{trim\PYGZus{}comment}\PYG{p}{(}\PYG{n}{buf}\PYG{p}{)}\PYG{p}{)}\PYG{p}{;}
\end{sphinxVerbatim}


\begin{sphinxseealso}{See also:}

\sphinxAtStartPar
{\hyperref[\detokenize{Developer_utils:_CPPv4N9LAMMPS_NS5utils11count_wordsERKNSt6stringERKNSt6stringE}]{\sphinxcrossref{\sphinxcode{\sphinxupquote{utils::count\_words()}}}}},
\sphinxcode{\sphinxupquote{utils::trim\_comments()}}


\end{sphinxseealso}



\subsection{Use utils::numeric() functions instead of force\sphinxhyphen{}\textgreater{}numeric()}
\label{\detokenize{Developer_updating:use-utils-numeric-functions-instead-of-force-numeric}}
\sphinxAtStartPar
\DUrole{versionmodified,changed}{Changed in version 18Sep2020.}

\sphinxAtStartPar
The “numeric()” conversion functions (including “inumeric()”,
“bnumeric()”, and “tnumeric()”) have been moved from the Force class to
the {\hyperref[\detokenize{Developer_utils::doc}]{\sphinxcrossref{\DUrole{doc}{utils namespace}}}}.  Also they take an
additional argument that selects whether the \sphinxcode{\sphinxupquote{Error::all()}} or
\sphinxcode{\sphinxupquote{Error::one()}} function should be called in case of an error.  The
former should be used when \sphinxstyleemphasis{all} MPI processes call the conversion
function and the latter \sphinxstyleemphasis{must} be used when they are called from only
one or a subset of the MPI processes.

\sphinxAtStartPar
Old:

\begin{sphinxVerbatim}[commandchars=\\\{\}]
\PYG{n}{val}\PYG{+w}{ }\PYG{o}{=}\PYG{+w}{ }\PYG{n}{force}\PYG{o}{\PYGZhy{}}\PYG{o}{\PYGZgt{}}\PYG{n}{numeric}\PYG{p}{(}\PYG{n}{FLERR}\PYG{p}{,}\PYG{+w}{ }\PYG{n}{arg}\PYG{p}{[}\PYG{l+m+mi}{1}\PYG{p}{]}\PYG{p}{)}\PYG{p}{;}
\PYG{n}{num}\PYG{+w}{ }\PYG{o}{=}\PYG{+w}{ }\PYG{n}{force}\PYG{o}{\PYGZhy{}}\PYG{o}{\PYGZgt{}}\PYG{n}{inumeric}\PYG{p}{(}\PYG{n}{FLERR}\PYG{p}{,}\PYG{+w}{ }\PYG{n}{arg}\PYG{p}{[}\PYG{l+m+mi}{2}\PYG{p}{]}\PYG{p}{)}\PYG{p}{;}
\end{sphinxVerbatim}

\sphinxAtStartPar
New:

\begin{sphinxVerbatim}[commandchars=\\\{\}]
\PYG{n}{val}\PYG{+w}{ }\PYG{o}{=}\PYG{+w}{ }\PYG{n}{utils}\PYG{o}{:}\PYG{o}{:}\PYG{n}{numeric}\PYG{p}{(}\PYG{n}{FLERR}\PYG{p}{,}\PYG{+w}{ }\PYG{n+nb}{true}\PYG{p}{,}\PYG{+w}{ }\PYG{n}{arg}\PYG{p}{[}\PYG{l+m+mi}{1}\PYG{p}{]}\PYG{p}{,}\PYG{+w}{ }\PYG{n}{lmp}\PYG{p}{)}\PYG{p}{;}
\PYG{n}{num}\PYG{+w}{ }\PYG{o}{=}\PYG{+w}{ }\PYG{n}{utils}\PYG{o}{:}\PYG{o}{:}\PYG{n}{inumeric}\PYG{p}{(}\PYG{n}{FLERR}\PYG{p}{,}\PYG{+w}{ }\PYG{n+nb}{false}\PYG{p}{,}\PYG{+w}{ }\PYG{n}{arg}\PYG{p}{[}\PYG{l+m+mi}{2}\PYG{p}{]}\PYG{p}{,}\PYG{+w}{ }\PYG{n}{lmp}\PYG{p}{)}\PYG{p}{;}
\end{sphinxVerbatim}


\begin{sphinxseealso}{See also:}

\sphinxAtStartPar
{\hyperref[\detokenize{Developer_utils:_CPPv4N9LAMMPS_NS5utils7numericEPKciPKcbP6LAMMPS}]{\sphinxcrossref{\sphinxcode{\sphinxupquote{utils::numeric()}}}}},
{\hyperref[\detokenize{Developer_utils:_CPPv4N9LAMMPS_NS5utils8inumericEPKciPKcbP6LAMMPS}]{\sphinxcrossref{\sphinxcode{\sphinxupquote{utils::inumeric()}}}}},
{\hyperref[\detokenize{Developer_utils:_CPPv4N9LAMMPS_NS5utils8bnumericEPKciPKcbP6LAMMPS}]{\sphinxcrossref{\sphinxcode{\sphinxupquote{utils::bnumeric()}}}}},
{\hyperref[\detokenize{Developer_utils:_CPPv4N9LAMMPS_NS5utils8tnumericEPKciPKcbP6LAMMPS}]{\sphinxcrossref{\sphinxcode{\sphinxupquote{utils::tnumeric()}}}}}


\end{sphinxseealso}



\subsection{Use utils::open\_potential() function to open potential files}
\label{\detokenize{Developer_updating:use-utils-open-potential-function-to-open-potential-files}}
\sphinxAtStartPar
\DUrole{versionmodified,changed}{Changed in version 18Sep2020.}

\sphinxAtStartPar
The {\hyperref[\detokenize{Developer_utils:_CPPv4N9LAMMPS_NS5utils14open_potentialERKNSt6stringEP6LAMMPSPi}]{\sphinxcrossref{\sphinxcode{\sphinxupquote{utils::open\_potential()}}}}} function must be used to replace
calls to \sphinxcode{\sphinxupquote{force\sphinxhyphen{}\textgreater{}open\_potential()}} and should be used to replace
\sphinxcode{\sphinxupquote{fopen()}} for opening potential files for reading.  The custom
function does three additional steps compared to \sphinxcode{\sphinxupquote{fopen()}}: 1) it will
try to parse the \sphinxcode{\sphinxupquote{UNITS:}} and \sphinxcode{\sphinxupquote{DATE:}} metadata and will stop with an
error on a units mismatch and will print the date info, if present, in
the log file; 2) for pair styles that support it, it will set up
possible automatic unit conversions based on the embedded unit
information and LAMMPS’ current units setting; 3) it will not only try
to open a potential file at the given path, but will also search in the
folders listed in the \sphinxcode{\sphinxupquote{LAMMPS\_POTENTIALS}} environment variable.  This
allows potential files to reside in a common location instead of having to
copy them around for simulations.

\sphinxAtStartPar
Old:

\begin{sphinxVerbatim}[commandchars=\\\{\}]
\PYG{n}{fp}\PYG{+w}{ }\PYG{o}{=}\PYG{+w}{ }\PYG{n}{force}\PYG{o}{\PYGZhy{}}\PYG{o}{\PYGZgt{}}\PYG{n}{open\PYGZus{}potential}\PYG{p}{(}\PYG{n}{filename}\PYG{p}{)}\PYG{p}{;}
\PYG{n}{fp}\PYG{+w}{ }\PYG{o}{=}\PYG{+w}{ }\PYG{n}{fopen}\PYG{p}{(}\PYG{n}{filename}\PYG{p}{,}\PYG{+w}{ }\PYG{l+s}{\PYGZdq{}}\PYG{l+s}{r}\PYG{l+s}{\PYGZdq{}}\PYG{p}{)}\PYG{p}{;}
\end{sphinxVerbatim}

\sphinxAtStartPar
New:

\begin{sphinxVerbatim}[commandchars=\\\{\}]
\PYG{n}{fp}\PYG{+w}{ }\PYG{o}{=}\PYG{+w}{ }\PYG{n}{utils}\PYG{o}{:}\PYG{o}{:}\PYG{n}{open\PYGZus{}potential}\PYG{p}{(}\PYG{n}{filename}\PYG{p}{,}\PYG{+w}{ }\PYG{n}{lmp}\PYG{p}{)}\PYG{p}{;}
\end{sphinxVerbatim}


\subsection{Use symbolic Atom and AtomVec constants instead of numerical values}
\label{\detokenize{Developer_updating:use-symbolic-atom-and-atomvec-constants-instead-of-numerical-values}}
\sphinxAtStartPar
\DUrole{versionmodified,changed}{Changed in version 18Sep2020.}

\sphinxAtStartPar
Properties in LAMMPS that were represented by integer values (0, 1,
2, 3) to indicate settings in the \sphinxcode{\sphinxupquote{Atom}} and \sphinxcode{\sphinxupquote{AtomVec}} classes (or
classes derived from it) (and its derived classes) have been converted
to use scoped enumerators instead.


\begin{savenotes}\sphinxattablestart
\sphinxthistablewithglobalstyle
\centering
\begin{tabular}[t]{\X{23}{99}\X{10}{99}\X{23}{99}\X{10}{99}\X{23}{99}\X{10}{99}}
\sphinxtoprule
\sphinxstyletheadfamily 
\sphinxAtStartPar
Symbolic Constant
&\sphinxstyletheadfamily 
\sphinxAtStartPar
Value
&\sphinxstyletheadfamily 
\sphinxAtStartPar
Symbolic Constant
&\sphinxstyletheadfamily 
\sphinxAtStartPar
Value
&\sphinxstyletheadfamily 
\sphinxAtStartPar
Symbolic Constant
&\sphinxstyletheadfamily 
\sphinxAtStartPar
Value
\\
\sphinxmidrule
\sphinxtableatstartofbodyhook
\sphinxAtStartPar
Atom::GROW
&
\sphinxAtStartPar
0
&
\sphinxAtStartPar
Atom::ATOMIC
&
\sphinxAtStartPar
0
&
\sphinxAtStartPar
Atom::MAP\_NONE
&
\sphinxAtStartPar
0
\\
\sphinxhline
\sphinxAtStartPar
Atom::RESTART
&
\sphinxAtStartPar
1
&
\sphinxAtStartPar
Atom::MOLECULAR
&
\sphinxAtStartPar
1
&
\sphinxAtStartPar
Atom::MAP\_ARRAY
&
\sphinxAtStartPar
1
\\
\sphinxhline
\sphinxAtStartPar
Atom::BORDER
&
\sphinxAtStartPar
2
&
\sphinxAtStartPar
Atom::TEMPLATE
&
\sphinxAtStartPar
2
&
\sphinxAtStartPar
Atom::MAP\_HASH
&
\sphinxAtStartPar
2
\\
\sphinxhline
\sphinxAtStartPar
AtomVec::PER\_ATOM
&
\sphinxAtStartPar
0
&
\sphinxAtStartPar
AtomVec::PER\_TYPE
&
\sphinxAtStartPar
1
&
\sphinxAtStartPar
Atom::MAP\_YES
&
\sphinxAtStartPar
3
\\
\sphinxbottomrule
\end{tabular}
\sphinxtableafterendhook\par
\sphinxattableend\end{savenotes}

\sphinxAtStartPar
Old:

\begin{sphinxVerbatim}[commandchars=\\\{\}]
\PYG{n}{molecular}\PYG{+w}{ }\PYG{o}{=}\PYG{+w}{ }\PYG{l+m+mi}{0}\PYG{p}{;}
\PYG{n}{mass\PYGZus{}type}\PYG{+w}{ }\PYG{o}{=}\PYG{+w}{ }\PYG{l+m+mi}{1}\PYG{p}{;}
\PYG{k}{if}\PYG{+w}{ }\PYG{p}{(}\PYG{n}{atom}\PYG{o}{\PYGZhy{}}\PYG{o}{\PYGZgt{}}\PYG{n}{molecular}\PYG{+w}{ }\PYG{o}{=}\PYG{o}{=}\PYG{+w}{ }\PYG{l+m+mi}{2}\PYG{p}{)}
\PYG{k}{if}\PYG{+w}{ }\PYG{p}{(}\PYG{n}{atom}\PYG{o}{\PYGZhy{}}\PYG{o}{\PYGZgt{}}\PYG{n}{map\PYGZus{}style}\PYG{+w}{ }\PYG{o}{=}\PYG{o}{=}\PYG{+w}{ }\PYG{l+m+mi}{2}\PYG{p}{)}
\PYG{n}{atom}\PYG{o}{\PYGZhy{}}\PYG{o}{\PYGZgt{}}\PYG{n}{add\PYGZus{}callback}\PYG{p}{(}\PYG{l+m+mi}{0}\PYG{p}{)}\PYG{p}{;}
\PYG{n}{atom}\PYG{o}{\PYGZhy{}}\PYG{o}{\PYGZgt{}}\PYG{n}{delete\PYGZus{}callback}\PYG{p}{(}\PYG{n}{id}\PYG{p}{,}\PYG{l+m+mi}{1}\PYG{p}{)}\PYG{p}{;}
\end{sphinxVerbatim}

\sphinxAtStartPar
New:

\begin{sphinxVerbatim}[commandchars=\\\{\}]
\PYG{n}{molecular}\PYG{+w}{ }\PYG{o}{=}\PYG{+w}{ }\PYG{n}{Atom}\PYG{o}{:}\PYG{o}{:}\PYG{n}{ATOMIC}\PYG{p}{;}
\PYG{n}{mass\PYGZus{}type}\PYG{+w}{ }\PYG{o}{=}\PYG{+w}{ }\PYG{n}{AtomVec}\PYG{o}{:}\PYG{o}{:}\PYG{n}{PER\PYGZus{}TYPE}\PYG{p}{;}
\PYG{k}{if}\PYG{+w}{ }\PYG{p}{(}\PYG{n}{atom}\PYG{o}{\PYGZhy{}}\PYG{o}{\PYGZgt{}}\PYG{n}{molecular}\PYG{+w}{ }\PYG{o}{=}\PYG{o}{=}\PYG{+w}{ }\PYG{n}{Atom}\PYG{o}{:}\PYG{o}{:}\PYG{n}{TEMPLATE}\PYG{p}{)}
\PYG{k}{if}\PYG{+w}{ }\PYG{p}{(}\PYG{n}{atom}\PYG{o}{\PYGZhy{}}\PYG{o}{\PYGZgt{}}\PYG{n}{map\PYGZus{}style}\PYG{+w}{ }\PYG{o}{=}\PYG{o}{=}\PYG{+w}{ }\PYG{n}{Atom}\PYG{o}{:}\PYG{o}{:}\PYG{n}{MAP\PYGZus{}HASH}\PYG{p}{)}
\PYG{n}{atom}\PYG{o}{\PYGZhy{}}\PYG{o}{\PYGZgt{}}\PYG{n}{add\PYGZus{}callback}\PYG{p}{(}\PYG{n}{Atom}\PYG{o}{:}\PYG{o}{:}\PYG{n}{GROW}\PYG{p}{)}\PYG{p}{;}
\PYG{n}{atom}\PYG{o}{\PYGZhy{}}\PYG{o}{\PYGZgt{}}\PYG{n}{delete\PYGZus{}callback}\PYG{p}{(}\PYG{n}{id}\PYG{p}{,}\PYG{n}{Atom}\PYG{o}{:}\PYG{o}{:}\PYG{n}{RESTART}\PYG{p}{)}\PYG{p}{;}
\end{sphinxVerbatim}


\subsection{Simplify customized error messages}
\label{\detokenize{Developer_updating:simplify-customized-error-messages}}
\sphinxAtStartPar
\DUrole{versionmodified,changed}{Changed in version 14May2021.}

\sphinxAtStartPar
Aided by features of the bundled \{fmt\} library, error messages now
can have a variable number of arguments and the string will be interpreted
as a \{fmt\} style format string so that error messages can be
easily customized without having to use temporary buffers and \sphinxcode{\sphinxupquote{sprintf()}}.
Example:

\sphinxAtStartPar
Old:

\begin{sphinxVerbatim}[commandchars=\\\{\}]
\PYG{k}{if}\PYG{+w}{ }\PYG{p}{(}\PYG{n}{fptr}\PYG{+w}{ }\PYG{o}{=}\PYG{o}{=}\PYG{+w}{ }\PYG{n+nb}{NULL}\PYG{p}{)}\PYG{+w}{ }\PYG{p}{\PYGZob{}}
\PYG{+w}{  }\PYG{k+kt}{char}\PYG{+w}{ }\PYG{n}{str}\PYG{p}{[}\PYG{l+m+mi}{128}\PYG{p}{]}\PYG{p}{;}
\PYG{+w}{  }\PYG{n}{sprintf}\PYG{p}{(}\PYG{n}{str}\PYG{p}{,}\PYG{l+s}{\PYGZdq{}}\PYG{l+s}{Cannot open AEAM potential file \PYGZpc{}s}\PYG{l+s}{\PYGZdq{}}\PYG{p}{,}\PYG{n}{filename}\PYG{p}{)}\PYG{p}{;}
\PYG{+w}{  }\PYG{n}{error}\PYG{o}{\PYGZhy{}}\PYG{o}{\PYGZgt{}}\PYG{n}{one}\PYG{p}{(}\PYG{n}{FLERR}\PYG{p}{,}\PYG{n}{str}\PYG{p}{)}\PYG{p}{;}
\PYG{p}{\PYGZcb{}}
\end{sphinxVerbatim}

\sphinxAtStartPar
New:

\begin{sphinxVerbatim}[commandchars=\\\{\}]
\PYG{k}{if}\PYG{+w}{ }\PYG{p}{(}\PYG{n}{fptr}\PYG{+w}{ }\PYG{o}{=}\PYG{o}{=}\PYG{+w}{ }\PYG{n+nb}{nullptr}\PYG{p}{)}
\PYG{+w}{  }\PYG{n}{error}\PYG{o}{\PYGZhy{}}\PYG{o}{\PYGZgt{}}\PYG{n}{one}\PYG{p}{(}\PYG{n}{FLERR}\PYG{p}{,}\PYG{+w}{ }\PYG{l+s}{\PYGZdq{}}\PYG{l+s}{Cannot open AEAM potential file \PYGZob{}\PYGZcb{}: \PYGZob{}\PYGZcb{}}\PYG{l+s}{\PYGZdq{}}\PYG{p}{,}\PYG{+w}{ }\PYG{n}{filename}\PYG{p}{,}\PYG{+w}{ }\PYG{n}{utils}\PYG{o}{:}\PYG{o}{:}\PYG{n}{getsyserror}\PYG{p}{(}\PYG{p}{)}\PYG{p}{)}\PYG{p}{;}
\end{sphinxVerbatim}


\subsection{Use of “override” instead of “virtual”}
\label{\detokenize{Developer_updating:use-of-override-instead-of-virtual}}
\sphinxAtStartPar
\DUrole{versionmodified,changed}{Changed in version 17Feb2022.}

\sphinxAtStartPar
Since LAMMPS requires C++11, we switched to use the “override” keyword
instead of “virtual” to indicate polymorphism in derived classes.  This
allows the C++ compiler to better detect inconsistencies when an
override is intended or not.  Please note that “override” has to be
added to \sphinxstylestrong{all} polymorph functions in derived classes and “virtual”
\sphinxstyleemphasis{only} to the function in the base class (or the destructor).  Here is
an example from the \sphinxcode{\sphinxupquote{FixWallReflect}} class:

\sphinxAtStartPar
Old:

\begin{sphinxVerbatim}[commandchars=\\\{\}]
\PYG{n}{FixWallReflect}\PYG{p}{(}\PYG{k}{class}\PYG{+w}{ }\PYG{n+nc}{LAMMPS}\PYG{+w}{ }\PYG{o}{*}\PYG{p}{,}\PYG{+w}{ }\PYG{k+kt}{int}\PYG{p}{,}\PYG{+w}{ }\PYG{k+kt}{char}\PYG{+w}{ }\PYG{o}{*}\PYG{o}{*}\PYG{p}{)}\PYG{p}{;}
\PYG{k}{virtual}\PYG{+w}{ }\PYG{o}{\PYGZti{}}\PYG{n}{FixWallReflect}\PYG{p}{(}\PYG{p}{)}\PYG{p}{;}
\PYG{k+kt}{int}\PYG{+w}{ }\PYG{n+nf}{setmask}\PYG{p}{(}\PYG{p}{)}\PYG{p}{;}
\PYG{k+kt}{void}\PYG{+w}{ }\PYG{n+nf}{init}\PYG{p}{(}\PYG{p}{)}\PYG{p}{;}
\PYG{k+kt}{void}\PYG{+w}{ }\PYG{n+nf}{post\PYGZus{}integrate}\PYG{p}{(}\PYG{p}{)}\PYG{p}{;}
\end{sphinxVerbatim}

\sphinxAtStartPar
New:

\begin{sphinxVerbatim}[commandchars=\\\{\}]
\PYG{n}{FixWallReflect}\PYG{p}{(}\PYG{k}{class}\PYG{+w}{ }\PYG{n+nc}{LAMMPS}\PYG{+w}{ }\PYG{o}{*}\PYG{p}{,}\PYG{+w}{ }\PYG{k+kt}{int}\PYG{p}{,}\PYG{+w}{ }\PYG{k+kt}{char}\PYG{+w}{ }\PYG{o}{*}\PYG{o}{*}\PYG{p}{)}\PYG{p}{;}
\PYG{o}{\PYGZti{}}\PYG{n}{FixWallReflect}\PYG{p}{(}\PYG{p}{)}\PYG{+w}{ }\PYG{k}{override}\PYG{p}{;}
\PYG{k+kt}{int}\PYG{+w}{ }\PYG{n+nf}{setmask}\PYG{p}{(}\PYG{p}{)}\PYG{+w}{ }\PYG{k}{override}\PYG{p}{;}
\PYG{k+kt}{void}\PYG{+w}{ }\PYG{n+nf}{init}\PYG{p}{(}\PYG{p}{)}\PYG{+w}{ }\PYG{k}{override}\PYG{p}{;}
\PYG{k+kt}{void}\PYG{+w}{ }\PYG{n+nf}{post\PYGZus{}integrate}\PYG{p}{(}\PYG{p}{)}\PYG{+w}{ }\PYG{k}{override}\PYG{p}{;}
\end{sphinxVerbatim}

\sphinxAtStartPar
This change set will neither cause a compilation failure, nor will it
change functionality, but if you plan to submit the updated code for
inclusion into the LAMMPS distribution, it will be requested for achieve
a consistent {\hyperref[\detokenize{Modify_style::doc}]{\sphinxcrossref{\DUrole{doc}{programming style}}}}.


\subsection{Simplified function names for forward and reverse communication}
\label{\detokenize{Developer_updating:simplified-function-names-for-forward-and-reverse-communication}}
\sphinxAtStartPar
\DUrole{versionmodified,changed}{Changed in version 24Mar2022.}

\sphinxAtStartPar
Rather than using the function name to distinguish between the different
forward and reverse communication functions for styles, LAMMPS now uses
the type of the “this” pointer argument.

\sphinxAtStartPar
Old:

\begin{sphinxVerbatim}[commandchars=\\\{\}]
\PYG{n}{comm}\PYG{o}{\PYGZhy{}}\PYG{o}{\PYGZgt{}}\PYG{n}{forward\PYGZus{}comm\PYGZus{}pair}\PYG{p}{(}\PYG{k}{this}\PYG{p}{)}\PYG{p}{;}
\PYG{n}{comm}\PYG{o}{\PYGZhy{}}\PYG{o}{\PYGZgt{}}\PYG{n}{forward\PYGZus{}comm\PYGZus{}fix}\PYG{p}{(}\PYG{k}{this}\PYG{p}{)}\PYG{p}{;}
\PYG{n}{comm}\PYG{o}{\PYGZhy{}}\PYG{o}{\PYGZgt{}}\PYG{n}{forward\PYGZus{}comm\PYGZus{}compute}\PYG{p}{(}\PYG{k}{this}\PYG{p}{)}\PYG{p}{;}
\PYG{n}{comm}\PYG{o}{\PYGZhy{}}\PYG{o}{\PYGZgt{}}\PYG{n}{forward\PYGZus{}comm\PYGZus{}dump}\PYG{p}{(}\PYG{k}{this}\PYG{p}{)}\PYG{p}{;}
\PYG{n}{comm}\PYG{o}{\PYGZhy{}}\PYG{o}{\PYGZgt{}}\PYG{n}{reverse\PYGZus{}comm\PYGZus{}pair}\PYG{p}{(}\PYG{k}{this}\PYG{p}{)}\PYG{p}{;}
\PYG{n}{comm}\PYG{o}{\PYGZhy{}}\PYG{o}{\PYGZgt{}}\PYG{n}{reverse\PYGZus{}comm\PYGZus{}fix}\PYG{p}{(}\PYG{k}{this}\PYG{p}{)}\PYG{p}{;}
\PYG{n}{comm}\PYG{o}{\PYGZhy{}}\PYG{o}{\PYGZgt{}}\PYG{n}{reverse\PYGZus{}comm\PYGZus{}compute}\PYG{p}{(}\PYG{k}{this}\PYG{p}{)}\PYG{p}{;}
\PYG{n}{comm}\PYG{o}{\PYGZhy{}}\PYG{o}{\PYGZgt{}}\PYG{n}{reverse\PYGZus{}comm\PYGZus{}dump}\PYG{p}{(}\PYG{k}{this}\PYG{p}{)}\PYG{p}{;}
\end{sphinxVerbatim}

\sphinxAtStartPar
New:

\begin{sphinxVerbatim}[commandchars=\\\{\}]
\PYG{n}{comm}\PYG{o}{\PYGZhy{}}\PYG{o}{\PYGZgt{}}\PYG{n}{forward\PYGZus{}comm}\PYG{p}{(}\PYG{k}{this}\PYG{p}{)}\PYG{p}{;}
\PYG{n}{comm}\PYG{o}{\PYGZhy{}}\PYG{o}{\PYGZgt{}}\PYG{n}{reverse\PYGZus{}comm}\PYG{p}{(}\PYG{k}{this}\PYG{p}{)}\PYG{p}{;}
\end{sphinxVerbatim}

\sphinxAtStartPar
This change is \sphinxstylestrong{required} or else the code will not compile.


\subsection{Simplified and more compact neighbor list requests}
\label{\detokenize{Developer_updating:simplified-and-more-compact-neighbor-list-requests}}
\sphinxAtStartPar
\DUrole{versionmodified,changed}{Changed in version 24Mar2022.}

\sphinxAtStartPar
This change set reduces the amount of code required to request a
neighbor list.  It enforces consistency and no longer requires to change
internal data of the request.  More information on neighbor list
requests can be {\hyperref[\detokenize{Developer_notes::doc}]{\sphinxcrossref{\DUrole{doc}{found here}}}}. Example from the
\sphinxcode{\sphinxupquote{ComputeRDF}} class:

\sphinxAtStartPar
Old:

\begin{sphinxVerbatim}[commandchars=\\\{\}]
\PYG{k+kt}{int}\PYG{+w}{ }\PYG{n}{irequest}\PYG{+w}{ }\PYG{o}{=}\PYG{+w}{ }\PYG{n}{neighbor}\PYG{o}{\PYGZhy{}}\PYG{o}{\PYGZgt{}}\PYG{n}{request}\PYG{p}{(}\PYG{k}{this}\PYG{p}{,}\PYG{n}{instance\PYGZus{}me}\PYG{p}{)}\PYG{p}{;}
\PYG{n}{neighbor}\PYG{o}{\PYGZhy{}}\PYG{o}{\PYGZgt{}}\PYG{n}{requests}\PYG{p}{[}\PYG{n}{irequest}\PYG{p}{]}\PYG{o}{\PYGZhy{}}\PYG{o}{\PYGZgt{}}\PYG{n}{pair}\PYG{+w}{ }\PYG{o}{=}\PYG{+w}{ }\PYG{l+m+mi}{0}\PYG{p}{;}
\PYG{n}{neighbor}\PYG{o}{\PYGZhy{}}\PYG{o}{\PYGZgt{}}\PYG{n}{requests}\PYG{p}{[}\PYG{n}{irequest}\PYG{p}{]}\PYG{o}{\PYGZhy{}}\PYG{o}{\PYGZgt{}}\PYG{n}{compute}\PYG{+w}{ }\PYG{o}{=}\PYG{+w}{ }\PYG{l+m+mi}{1}\PYG{p}{;}
\PYG{n}{neighbor}\PYG{o}{\PYGZhy{}}\PYG{o}{\PYGZgt{}}\PYG{n}{requests}\PYG{p}{[}\PYG{n}{irequest}\PYG{p}{]}\PYG{o}{\PYGZhy{}}\PYG{o}{\PYGZgt{}}\PYG{n}{occasional}\PYG{+w}{ }\PYG{o}{=}\PYG{+w}{ }\PYG{l+m+mi}{1}\PYG{p}{;}
\PYG{k}{if}\PYG{+w}{ }\PYG{p}{(}\PYG{n}{cutflag}\PYG{p}{)}\PYG{+w}{ }\PYG{p}{\PYGZob{}}
\PYG{+w}{  }\PYG{n}{neighbor}\PYG{o}{\PYGZhy{}}\PYG{o}{\PYGZgt{}}\PYG{n}{requests}\PYG{p}{[}\PYG{n}{irequest}\PYG{p}{]}\PYG{o}{\PYGZhy{}}\PYG{o}{\PYGZgt{}}\PYG{n}{cut}\PYG{+w}{ }\PYG{o}{=}\PYG{+w}{ }\PYG{l+m+mi}{1}\PYG{p}{;}
\PYG{+w}{  }\PYG{n}{neighbor}\PYG{o}{\PYGZhy{}}\PYG{o}{\PYGZgt{}}\PYG{n}{requests}\PYG{p}{[}\PYG{n}{irequest}\PYG{p}{]}\PYG{o}{\PYGZhy{}}\PYG{o}{\PYGZgt{}}\PYG{n}{cutoff}\PYG{+w}{ }\PYG{o}{=}\PYG{+w}{ }\PYG{n}{mycutneigh}\PYG{p}{;}
\PYG{p}{\PYGZcb{}}
\end{sphinxVerbatim}

\sphinxAtStartPar
New:

\begin{sphinxVerbatim}[commandchars=\\\{\}]
\PYG{k}{auto}\PYG{+w}{ }\PYG{n}{req}\PYG{+w}{ }\PYG{o}{=}\PYG{+w}{ }\PYG{n}{neighbor}\PYG{o}{\PYGZhy{}}\PYG{o}{\PYGZgt{}}\PYG{n}{add\PYGZus{}request}\PYG{p}{(}\PYG{k}{this}\PYG{p}{,}\PYG{+w}{ }\PYG{n}{NeighConst}\PYG{o}{:}\PYG{o}{:}\PYG{n}{REQ\PYGZus{}OCCASIONAL}\PYG{p}{)}\PYG{p}{;}
\PYG{k}{if}\PYG{+w}{ }\PYG{p}{(}\PYG{n}{cutflag}\PYG{p}{)}\PYG{+w}{ }\PYG{n}{req}\PYG{o}{\PYGZhy{}}\PYG{o}{\PYGZgt{}}\PYG{n}{set\PYGZus{}cutoff}\PYG{p}{(}\PYG{n}{mycutneigh}\PYG{p}{)}\PYG{p}{;}
\end{sphinxVerbatim}

\sphinxAtStartPar
Public access to the \sphinxcode{\sphinxupquote{NeighRequest}} class data members has been
removed so this update is \sphinxstylestrong{required} to avoid compilation failure.


\subsection{Split of fix STORE into fix STORE/GLOBAL and fix STORE/PERATOM}
\label{\detokenize{Developer_updating:split-of-fix-store-into-fix-store-global-and-fix-store-peratom}}
\sphinxAtStartPar
\DUrole{versionmodified,changed}{Changed in version 15Sep2022.}

\sphinxAtStartPar
This change splits the GLOBAL and PERATOM modes of fix STORE into two
separate fixes STORE/GLOBAL and STORE/PERATOM.  There was very little
shared code between the two fix STORE modes and the two different code
paths had to be prefixed with if statements.  Furthermore, some flags
were used differently in the two modes leading to confusion.  Splitting
the code into two fix styles, makes it more easily maintainable.  Since
these are internal fixes, there is no user visible change.

\sphinxAtStartPar
Old:

\begin{sphinxVerbatim}[commandchars=\\\{\}]
\PYG{c+cp}{\PYGZsh{}}\PYG{c+cp}{include}\PYG{+w}{ }\PYG{c+cpf}{\PYGZdq{}fix\PYGZus{}store.h\PYGZdq{}}

\PYG{n}{FixStore}\PYG{+w}{ }\PYG{o}{*}\PYG{n}{fix}\PYG{+w}{ }\PYG{o}{=}\PYG{+w}{ }\PYG{k}{dynamic\PYGZus{}cast}\PYG{o}{\PYGZlt{}}\PYG{n}{FixStore}\PYG{+w}{ }\PYG{o}{*}\PYG{o}{\PYGZgt{}}\PYG{p}{(}
\PYG{+w}{   }\PYG{n}{modify}\PYG{o}{\PYGZhy{}}\PYG{o}{\PYGZgt{}}\PYG{n}{add\PYGZus{}fix}\PYG{p}{(}\PYG{n}{fmt}\PYG{o}{:}\PYG{o}{:}\PYG{n}{format}\PYG{p}{(}\PYG{l+s}{\PYGZdq{}}\PYG{l+s}{\PYGZob{}\PYGZcb{} \PYGZob{}\PYGZcb{} STORE peratom 1 13}\PYG{l+s}{\PYGZdq{}}\PYG{p}{,}\PYG{n}{id\PYGZus{}pole}\PYG{p}{,}\PYG{n}{group}\PYG{o}{\PYGZhy{}}\PYG{o}{\PYGZgt{}}\PYG{n}{names}\PYG{p}{[}\PYG{l+m+mi}{0}\PYG{p}{]}\PYG{p}{)}\PYG{p}{)}\PYG{p}{;}

\PYG{n}{FixStore}\PYG{+w}{ }\PYG{o}{*}\PYG{n}{fix}\PYG{+w}{ }\PYG{o}{=}\PYG{+w}{ }\PYG{k}{dynamic\PYGZus{}cast}\PYG{o}{\PYGZlt{}}\PYG{n}{FixStore}\PYG{+w}{ }\PYG{o}{*}\PYG{o}{\PYGZgt{}}\PYG{p}{(}\PYG{n}{modify}\PYG{o}{\PYGZhy{}}\PYG{o}{\PYGZgt{}}\PYG{n}{get\PYGZus{}fix\PYGZus{}by\PYGZus{}id}\PYG{p}{(}\PYG{n}{id\PYGZus{}pole}\PYG{p}{)}\PYG{p}{)}\PYG{p}{;}
\end{sphinxVerbatim}

\sphinxAtStartPar
New:

\begin{sphinxVerbatim}[commandchars=\\\{\}]
\PYG{c+cp}{\PYGZsh{}}\PYG{c+cp}{include}\PYG{+w}{ }\PYG{c+cpf}{\PYGZdq{}fix\PYGZus{}store\PYGZus{}peratom.h\PYGZdq{}}

\PYG{n}{FixStorePeratom}\PYG{+w}{ }\PYG{o}{*}\PYG{n}{fix}\PYG{+w}{ }\PYG{o}{=}\PYG{+w}{ }\PYG{k}{dynamic\PYGZus{}cast}\PYG{o}{\PYGZlt{}}\PYG{n}{FixStorePeratom}\PYG{+w}{ }\PYG{o}{*}\PYG{o}{\PYGZgt{}}\PYG{p}{(}
\PYG{+w}{   }\PYG{n}{modify}\PYG{o}{\PYGZhy{}}\PYG{o}{\PYGZgt{}}\PYG{n}{add\PYGZus{}fix}\PYG{p}{(}\PYG{n}{fmt}\PYG{o}{:}\PYG{o}{:}\PYG{n}{format}\PYG{p}{(}\PYG{l+s}{\PYGZdq{}}\PYG{l+s}{\PYGZob{}\PYGZcb{} \PYGZob{}\PYGZcb{} STORE/PERATOM 1 13}\PYG{l+s}{\PYGZdq{}}\PYG{p}{,}\PYG{n}{id\PYGZus{}pole}\PYG{p}{,}\PYG{n}{group}\PYG{o}{\PYGZhy{}}\PYG{o}{\PYGZgt{}}\PYG{n}{names}\PYG{p}{[}\PYG{l+m+mi}{0}\PYG{p}{]}\PYG{p}{)}\PYG{p}{)}\PYG{p}{;}

\PYG{n}{FixStorePeratom}\PYG{+w}{ }\PYG{o}{*}\PYG{n}{fix}\PYG{+w}{ }\PYG{o}{=}\PYG{+w}{ }\PYG{k}{dynamic\PYGZus{}cast}\PYG{o}{\PYGZlt{}}\PYG{n}{FixStorePeratom}\PYG{+w}{ }\PYG{o}{*}\PYG{o}{\PYGZgt{}}\PYG{p}{(}\PYG{n}{modify}\PYG{o}{\PYGZhy{}}\PYG{o}{\PYGZgt{}}\PYG{n}{get\PYGZus{}fix\PYGZus{}by\PYGZus{}id}\PYG{p}{(}\PYG{n}{id\PYGZus{}pole}\PYG{p}{)}\PYG{p}{)}\PYG{p}{;}
\end{sphinxVerbatim}

\sphinxAtStartPar
Old:

\begin{sphinxVerbatim}[commandchars=\\\{\}]
\PYG{c+cp}{\PYGZsh{}}\PYG{c+cp}{include}\PYG{+w}{ }\PYG{c+cpf}{\PYGZdq{}fix\PYGZus{}store.h\PYGZdq{}}

\PYG{n}{FixStore}\PYG{+w}{ }\PYG{o}{*}\PYG{n}{fix}\PYG{+w}{ }\PYG{o}{=}\PYG{+w}{ }\PYG{k}{dynamic\PYGZus{}cast}\PYG{o}{\PYGZlt{}}\PYG{n}{FixStore}\PYG{+w}{ }\PYG{o}{*}\PYG{o}{\PYGZgt{}}\PYG{p}{(}
\PYG{+w}{   }\PYG{n}{modify}\PYG{o}{\PYGZhy{}}\PYG{o}{\PYGZgt{}}\PYG{n}{add\PYGZus{}fix}\PYG{p}{(}\PYG{n}{fmt}\PYG{o}{:}\PYG{o}{:}\PYG{n}{format}\PYG{p}{(}\PYG{l+s}{\PYGZdq{}}\PYG{l+s}{\PYGZob{}\PYGZcb{} \PYGZob{}\PYGZcb{} STORE global 1 1}\PYG{l+s}{\PYGZdq{}}\PYG{p}{,}\PYG{n}{id\PYGZus{}fix}\PYG{p}{,}\PYG{n}{group}\PYG{o}{\PYGZhy{}}\PYG{o}{\PYGZgt{}}\PYG{n}{names}\PYG{p}{[}\PYG{n}{igroup}\PYG{p}{]}\PYG{p}{)}\PYG{p}{)}\PYG{p}{;}

\PYG{n}{FixStore}\PYG{+w}{ }\PYG{o}{*}\PYG{n}{fix}\PYG{+w}{ }\PYG{o}{=}\PYG{+w}{ }\PYG{k}{dynamic\PYGZus{}cast}\PYG{o}{\PYGZlt{}}\PYG{n}{FixStore}\PYG{+w}{ }\PYG{o}{*}\PYG{o}{\PYGZgt{}}\PYG{p}{(}\PYG{n}{modify}\PYG{o}{\PYGZhy{}}\PYG{o}{\PYGZgt{}}\PYG{n}{get\PYGZus{}fix\PYGZus{}by\PYGZus{}id}\PYG{p}{(}\PYG{n}{id\PYGZus{}fix}\PYG{p}{)}\PYG{p}{)}\PYG{p}{;}
\end{sphinxVerbatim}

\sphinxAtStartPar
New:

\begin{sphinxVerbatim}[commandchars=\\\{\}]
\PYG{c+cp}{\PYGZsh{}}\PYG{c+cp}{include}\PYG{+w}{ }\PYG{c+cpf}{\PYGZdq{}fix\PYGZus{}store\PYGZus{}global.h\PYGZdq{}}

\PYG{n}{FixStoreGlobal}\PYG{+w}{ }\PYG{o}{*}\PYG{n}{fix}\PYG{+w}{ }\PYG{o}{=}\PYG{+w}{ }\PYG{k}{dynamic\PYGZus{}cast}\PYG{o}{\PYGZlt{}}\PYG{n}{FixStoreGlobal}\PYG{+w}{ }\PYG{o}{*}\PYG{o}{\PYGZgt{}}\PYG{p}{(}
\PYG{+w}{   }\PYG{n}{modify}\PYG{o}{\PYGZhy{}}\PYG{o}{\PYGZgt{}}\PYG{n}{add\PYGZus{}fix}\PYG{p}{(}\PYG{n}{fmt}\PYG{o}{:}\PYG{o}{:}\PYG{n}{format}\PYG{p}{(}\PYG{l+s}{\PYGZdq{}}\PYG{l+s}{\PYGZob{}\PYGZcb{} \PYGZob{}\PYGZcb{} STORE/GLOBAL 1 1}\PYG{l+s}{\PYGZdq{}}\PYG{p}{,}\PYG{n}{id\PYGZus{}fix}\PYG{p}{,}\PYG{n}{group}\PYG{o}{\PYGZhy{}}\PYG{o}{\PYGZgt{}}\PYG{n}{names}\PYG{p}{[}\PYG{n}{igroup}\PYG{p}{]}\PYG{p}{)}\PYG{p}{)}\PYG{p}{;}

\PYG{n}{FixStoreGlobal}\PYG{+w}{ }\PYG{o}{*}\PYG{n}{fix}\PYG{+w}{ }\PYG{o}{=}\PYG{+w}{ }\PYG{k}{dynamic\PYGZus{}cast}\PYG{o}{\PYGZlt{}}\PYG{n}{FixStoreGlobal}\PYG{+w}{ }\PYG{o}{*}\PYG{o}{\PYGZgt{}}\PYG{p}{(}\PYG{n}{modify}\PYG{o}{\PYGZhy{}}\PYG{o}{\PYGZgt{}}\PYG{n}{get\PYGZus{}fix\PYGZus{}by\PYGZus{}id}\PYG{p}{(}\PYG{n}{id\PYGZus{}fix}\PYG{p}{)}\PYG{p}{)}\PYG{p}{;}
\end{sphinxVerbatim}

\sphinxAtStartPar
This change is \sphinxstylestrong{required} or else the code will not compile.


\subsection{Rename of fix STORE/PERATOM to fix STORE/ATOM and change of arguments}
\label{\detokenize{Developer_updating:rename-of-fix-store-peratom-to-fix-store-atom-and-change-of-arguments}}
\sphinxAtStartPar
\DUrole{versionmodified,changed}{Changed in version 28Mar2023.}

\sphinxAtStartPar
The available functionality of the internal fix to store per\sphinxhyphen{}atom
properties was expanded to enable storing data with ghost atoms and to
support binary restart files.  With those changes, the fix was renamed
to fix STORE/ATOM and the number and order of (required) arguments has
changed.

\sphinxAtStartPar
Old syntax: \sphinxcode{\sphinxupquote{ID group\sphinxhyphen{}ID STORE/PERATOM rflag n1 n2 {[}n3{]}}}
\begin{itemize}
\item {} 
\sphinxAtStartPar
\sphinxstyleemphasis{rflag} = 0/1, \sphinxstyleemphasis{no}/\sphinxstyleemphasis{yes} store per\sphinxhyphen{}atom values in restart file

\item {} 
\sphinxAtStartPar
\(n1 = 1, n2 = 1, \mathrm{no}\;n3 \to\) per\sphinxhyphen{}atom vector, single value per atom

\item {} 
\sphinxAtStartPar
\(n1 = 1, n2 > 1, \mathrm{no}\;n3 \to\) per\sphinxhyphen{}atom array, \sphinxstyleemphasis{n2} values per atom

\item {} 
\sphinxAtStartPar
\(n1 = 1, n2 > 0, n3 > 0 \to\) per\sphinxhyphen{}atom tensor, \sphinxstyleemphasis{n2} x \sphinxstyleemphasis{n3} values per atom

\end{itemize}

\sphinxAtStartPar
New syntax:  \sphinxcode{\sphinxupquote{ID group\sphinxhyphen{}ID STORE/ATOM n1 n2 gflag rflag}}
\begin{itemize}
\item {} 
\sphinxAtStartPar
\(n1 = 1, n2 = 0 \to\) per\sphinxhyphen{}atom vector, single value per atom

\item {} 
\sphinxAtStartPar
\(n1 > 1, n2 = 0 \to\) per\sphinxhyphen{}atom array, \sphinxstyleemphasis{n1} values per atom

\item {} 
\sphinxAtStartPar
\(n1 > 0, n2 > 0 \to\) per\sphinxhyphen{}atom tensor, \sphinxstyleemphasis{n1} x \sphinxstyleemphasis{n2} values per atom

\item {} 
\sphinxAtStartPar
\sphinxstyleemphasis{gflag} = 0/1, \sphinxstyleemphasis{no}/\sphinxstyleemphasis{yes} communicate per\sphinxhyphen{}atom values with ghost atoms

\item {} 
\sphinxAtStartPar
\sphinxstyleemphasis{rflag} = 0/1, \sphinxstyleemphasis{no}/\sphinxstyleemphasis{yes} store per\sphinxhyphen{}atom values in restart file

\end{itemize}

\sphinxAtStartPar
Since this is an internal fix, there is no user visible change.


\subsection{Use Output::get\_dump\_by\_id() instead of Output::find\_dump()}
\label{\detokenize{Developer_updating:use-output-get-dump-by-id-instead-of-output-find-dump}}
\sphinxAtStartPar
\DUrole{versionmodified,changed}{Changed in version 15Sep2022.}

\sphinxAtStartPar
The accessor function to individual dump style instances has been changed
from \sphinxcode{\sphinxupquote{Output::find\_dump()}} returning the index of the dump instance in
the list of dumps to \sphinxcode{\sphinxupquote{Output::get\_dump\_by\_id()}} returning a pointer to
the dump directly.  Example:

\sphinxAtStartPar
Old:

\begin{sphinxVerbatim}[commandchars=\\\{\}]
\PYG{k+kt}{int}\PYG{+w}{ }\PYG{n}{idump}\PYG{+w}{ }\PYG{o}{=}\PYG{+w}{ }\PYG{n}{output}\PYG{o}{\PYGZhy{}}\PYG{o}{\PYGZgt{}}\PYG{n}{find\PYGZus{}dump}\PYG{p}{(}\PYG{n}{arg}\PYG{p}{[}\PYG{n}{iarg}\PYG{o}{+}\PYG{l+m+mi}{1}\PYG{p}{]}\PYG{p}{)}\PYG{p}{;}
\PYG{k}{if}\PYG{+w}{ }\PYG{p}{(}\PYG{n}{idump}\PYG{+w}{ }\PYG{o}{\PYGZlt{}}\PYG{+w}{ }\PYG{l+m+mi}{0}\PYG{p}{)}
\PYG{+w}{  }\PYG{n}{error}\PYG{o}{\PYGZhy{}}\PYG{o}{\PYGZgt{}}\PYG{n}{all}\PYG{p}{(}\PYG{n}{FLERR}\PYG{p}{,}\PYG{l+s}{\PYGZdq{}}\PYG{l+s}{Dump ID in hyper command does not exist}\PYG{l+s}{\PYGZdq{}}\PYG{p}{)}\PYG{p}{;}
\PYG{n}{memory}\PYG{o}{\PYGZhy{}}\PYG{o}{\PYGZgt{}}\PYG{n}{grow}\PYG{p}{(}\PYG{n}{dumplist}\PYG{p}{,}\PYG{n}{ndump}\PYG{o}{+}\PYG{l+m+mi}{1}\PYG{p}{,}\PYG{l+s}{\PYGZdq{}}\PYG{l+s}{hyper:dumplist}\PYG{l+s}{\PYGZdq{}}\PYG{p}{)}\PYG{p}{;}
\PYG{n}{dumplist}\PYG{p}{[}\PYG{n}{ndump}\PYG{o}{+}\PYG{o}{+}\PYG{p}{]}\PYG{+w}{ }\PYG{o}{=}\PYG{+w}{ }\PYG{n}{idump}\PYG{p}{;}

\PYG{p}{[}\PYG{p}{.}\PYG{p}{.}\PYG{p}{.}\PYG{p}{]}

\PYG{k}{if}\PYG{+w}{ }\PYG{p}{(}\PYG{n}{dumpflag}\PYG{p}{)}
\PYG{+w}{  }\PYG{k}{for}\PYG{+w}{ }\PYG{p}{(}\PYG{k+kt}{int}\PYG{+w}{ }\PYG{n}{idump}\PYG{+w}{ }\PYG{o}{=}\PYG{+w}{ }\PYG{l+m+mi}{0}\PYG{p}{;}\PYG{+w}{ }\PYG{n}{idump}\PYG{+w}{ }\PYG{o}{\PYGZlt{}}\PYG{+w}{ }\PYG{n}{ndump}\PYG{p}{;}\PYG{+w}{ }\PYG{n}{idump}\PYG{o}{+}\PYG{o}{+}\PYG{p}{)}
\PYG{+w}{    }\PYG{n}{output}\PYG{o}{\PYGZhy{}}\PYG{o}{\PYGZgt{}}\PYG{n}{dump}\PYG{p}{[}\PYG{n}{dumplist}\PYG{p}{[}\PYG{n}{idump}\PYG{p}{]}\PYG{p}{]}\PYG{o}{\PYGZhy{}}\PYG{o}{\PYGZgt{}}\PYG{n}{write}\PYG{p}{(}\PYG{p}{)}\PYG{p}{;}
\end{sphinxVerbatim}

\sphinxAtStartPar
New:

\begin{sphinxVerbatim}[commandchars=\\\{\}]
\PYG{k}{auto}\PYG{+w}{ }\PYG{n}{idump}\PYG{+w}{ }\PYG{o}{=}\PYG{+w}{ }\PYG{n}{output}\PYG{o}{\PYGZhy{}}\PYG{o}{\PYGZgt{}}\PYG{n}{get\PYGZus{}dump\PYGZus{}by\PYGZus{}id}\PYG{p}{(}\PYG{n}{arg}\PYG{p}{[}\PYG{n}{iarg}\PYG{o}{+}\PYG{l+m+mi}{1}\PYG{p}{]}\PYG{p}{)}\PYG{p}{;}
\PYG{k}{if}\PYG{+w}{ }\PYG{p}{(}\PYG{o}{!}\PYG{n}{idump}\PYG{p}{)}\PYG{+w}{ }\PYG{n}{error}\PYG{o}{\PYGZhy{}}\PYG{o}{\PYGZgt{}}\PYG{n}{all}\PYG{p}{(}\PYG{n}{FLERR}\PYG{p}{,}\PYG{l+s}{\PYGZdq{}}\PYG{l+s}{Dump ID \PYGZob{}\PYGZcb{} in hyper command does not exist}\PYG{l+s}{\PYGZdq{}}\PYG{p}{,}\PYG{+w}{ }\PYG{n}{arg}\PYG{p}{[}\PYG{n}{iarg}\PYG{o}{+}\PYG{l+m+mi}{1}\PYG{p}{]}\PYG{p}{)}\PYG{p}{;}
\PYG{n}{dumplist}\PYG{p}{.}\PYG{n}{emplace\PYGZus{}back}\PYG{p}{(}\PYG{n}{idump}\PYG{p}{)}\PYG{p}{;}

\PYG{p}{[}\PYG{p}{.}\PYG{p}{.}\PYG{p}{.}\PYG{p}{]}

\PYG{k}{if}\PYG{+w}{ }\PYG{p}{(}\PYG{n}{dumpflag}\PYG{p}{)}\PYG{+w}{ }\PYG{k}{for}\PYG{+w}{ }\PYG{p}{(}\PYG{k}{auto}\PYG{+w}{ }\PYG{n}{idump}\PYG{+w}{ }\PYG{o}{:}\PYG{+w}{ }\PYG{n}{dumplist}\PYG{p}{)}\PYG{+w}{ }\PYG{n}{idump}\PYG{o}{\PYGZhy{}}\PYG{o}{\PYGZgt{}}\PYG{n}{write}\PYG{p}{(}\PYG{p}{)}\PYG{p}{;}
\end{sphinxVerbatim}

\sphinxAtStartPar
This change is \sphinxstylestrong{required} or else the code will not compile.


\subsection{Refactored grid communication using Grid3d/Grid2d classes instead of GridComm}
\label{\detokenize{Developer_updating:refactored-grid-communication-using-grid3d-grid2d-classes-instead-of-gridcomm}}
\sphinxAtStartPar
\DUrole{versionmodified,changed}{Changed in version 22Dec2022.}

\sphinxAtStartPar
The \sphinxcode{\sphinxupquote{GridComm}} class was for creating and communicating distributed
grids was replaced by the \sphinxcode{\sphinxupquote{Grid3d}} class with added functionality.
A \sphinxcode{\sphinxupquote{Grid2d}} class was also added for additional flexibility.

\sphinxAtStartPar
The new functionality and commands using the two grid classes are
discussed on the following documentation pages:
\begin{itemize}
\item {} 
\sphinxAtStartPar
{\hyperref[\detokenize{Howto_grid::doc}]{\sphinxcrossref{\DUrole{doc}{Using distributed grids}}}}

\item {} 
\sphinxAtStartPar
{\hyperref[\detokenize{Developer_grid::doc}]{\sphinxcrossref{\DUrole{doc}{Use of distributed grids within style classes}}}}

\end{itemize}

\sphinxAtStartPar
If you have custom LAMMPS code, which uses the GridComm class, here are some notes
on how to adapt it for using the Grid3d class.
\begin{enumerate}
\sphinxsetlistlabels{\arabic}{enumi}{enumii}{(}{)}%
\item {} 
\sphinxAtStartPar
The constructor has changed to allow the \sphinxcode{\sphinxupquote{Grid3d}} / \sphinxcode{\sphinxupquote{Grid2d}}
classes to partition the global grid across processors, both for
owned and ghost grid cells.  Previously any class which called
\sphinxcode{\sphinxupquote{GridComm}} performed the partitioning itself and that information
was passed in the \sphinxcode{\sphinxupquote{GridComm::GridComm()}} constructor.  There are
several “set” functions which can be called to alter how \sphinxcode{\sphinxupquote{Grid3d}}
/ \sphinxcode{\sphinxupquote{Grid2d}} perform the partitioning.  They should be sufficient
for most use cases of the grid classes.

\item {} 
\sphinxAtStartPar
The partitioning is triggered by the \sphinxcode{\sphinxupquote{setup\_grid()}} method.

\item {} 
\sphinxAtStartPar
The \sphinxcode{\sphinxupquote{setup()}} method of the \sphinxcode{\sphinxupquote{GridComm}} class has been replaced
by the \sphinxcode{\sphinxupquote{setup\_comm()}} method in the new grid classes.  The syntax
for the \sphinxcode{\sphinxupquote{forward\_comm()}} and \sphinxcode{\sphinxupquote{reverse\_comm()}} methods is
slightly altered as is the syntax of the associated pack/unpack
callback methods.  But the functionality of these operations is the
same as before.

\item {} 
\sphinxAtStartPar
The new \sphinxcode{\sphinxupquote{Grid3d}} / \sphinxcode{\sphinxupquote{Grid2d}} classes have additional
functionality for dynamic load\sphinxhyphen{}balancing of grids and their
associated data across processors.  This did not exist in the
\sphinxcode{\sphinxupquote{GridComm}} class.

\end{enumerate}

\sphinxAtStartPar
This and more is explained in detail on the {\hyperref[\detokenize{Developer_grid::doc}]{\sphinxcrossref{\DUrole{doc}{Use of distributed grids within style classes}}}} page.
The following LAMMPS source files can be used as illustrative examples
for how the new grid classes are used by computes, fixes, and various
KSpace solvers which use distributed FFT grids:
\begin{itemize}
\item {} 
\sphinxAtStartPar
\sphinxcode{\sphinxupquote{src/fix\_ave\_grid.cpp}}

\item {} 
\sphinxAtStartPar
\sphinxcode{\sphinxupquote{src/compute\_property\_grid.cpp}}

\item {} 
\sphinxAtStartPar
\sphinxcode{\sphinxupquote{src/EXTRA\sphinxhyphen{}FIX/fix\_ttm\_grid.cpp}}

\item {} 
\sphinxAtStartPar
\sphinxcode{\sphinxupquote{src/KSPACE/pppm.cpp}}

\end{itemize}

\sphinxAtStartPar
This change is \sphinxstylestrong{required} or else the code will not compile.

\sphinxstepscope


\section{Writing plugins}
\label{\detokenize{Developer_plugins:writing-plugins}}\label{\detokenize{Developer_plugins::doc}}
\sphinxAtStartPar
Plugins provide a mechanism to add functionality to a LAMMPS executable
without recompiling LAMMPS.  The functionality for this and the
{\hyperref[\detokenize{plugin::doc}]{\sphinxcrossref{\DUrole{doc}{plugin command}}}} are implemented in the
{\hyperref[\detokenize{Packages_details:pkg-plugin}]{\sphinxcrossref{\DUrole{std,std-ref}{PLUGIN package}}}} which must be installed to use plugins.

\sphinxAtStartPar
Plugins use the operating system’s capability to load dynamic shared
object (DSO) files in a way similar shared libraries and then reference
specific functions in those DSOs.  Any DSO file with plugins has to
include an initialization function with a specific name,
“lammpsplugin\_init”, that has to follow specific rules described below.
When loading the DSO with the “plugin” command, this function is looked
up and called and will then register the contained plugin(s) with
LAMMPS.

\sphinxAtStartPar
When the environment variable \sphinxcode{\sphinxupquote{LAMMPS\_PLUGIN\_PATH}} is set, then LAMMPS
will search the directory (or directories) listed in this path for files
with names that end in \sphinxcode{\sphinxupquote{plugin.so}} (e.g. \sphinxcode{\sphinxupquote{helloplugin.so}}) and will
try to load the contained plugins automatically at start\sphinxhyphen{}up.  For
plugins that are loaded this way, the behavior of LAMMPS should be
identical to a binary where the corresponding code was compiled in
statically as a package.

\sphinxAtStartPar
From the programmer perspective this can work because of the object
oriented design of LAMMPS where all pair style commands are derived from
the class Pair, all fix style commands from the class Fix and so on and
usually only functions present in those base classes are called
directly.  When a {\hyperref[\detokenize{pair_style::doc}]{\sphinxcrossref{\DUrole{doc}{pair\_style command}}}} command or {\hyperref[\detokenize{fix::doc}]{\sphinxcrossref{\DUrole{doc}{fix command}}}} command is
issued a new instance of such a derived class is created.  This is done
by a so\sphinxhyphen{}called factory function which is mapped to the style name.  Thus
when, for example, the LAMMPS processes the command \sphinxcode{\sphinxupquote{pair\_style lj/cut
2.5}}, LAMMPS will look up the factory function for creating the
\sphinxcode{\sphinxupquote{PairLJCut}} class and then execute it.  The return value of that
function is a \sphinxcode{\sphinxupquote{Pair *}} pointer and the pointer will be assigned to the
location for the currently active pair style.

\sphinxAtStartPar
A DSO file with a plugin thus has to implement such a factory function
and register it with LAMMPS so that it gets added to the map of available
styles of the given category.  To register a plugin with LAMMPS an
initialization function has to be present in the DSO file called
\sphinxcode{\sphinxupquote{lammpsplugin\_init}} which is called with three \sphinxcode{\sphinxupquote{void *}} arguments:
a pointer to the current LAMMPS instance, a pointer to the opened DSO
handle, and a pointer to the registration function.  The registration
function takes two arguments: a pointer to a \sphinxcode{\sphinxupquote{lammpsplugin\_t}} struct
with information about the plugin and a pointer to the current LAMMPS
instance.  Please see below for an example of how the registration is
done.


\subsection{Members of \sphinxstyleliteralintitle{\sphinxupquote{lammpsplugin\_t}}}
\label{\detokenize{Developer_plugins:members-of-lammpsplugin-t}}

\begin{savenotes}\sphinxattablestart
\sphinxthistablewithglobalstyle
\centering
\begin{tabular}[t]{\X{15}{100}\X{85}{100}}
\sphinxtoprule
\sphinxstyletheadfamily 
\sphinxAtStartPar
Member
&\sphinxstyletheadfamily 
\sphinxAtStartPar
Description
\\
\sphinxmidrule
\sphinxtableatstartofbodyhook
\sphinxAtStartPar
version
&
\sphinxAtStartPar
LAMMPS Version string the plugin was compiled for
\\
\sphinxhline
\sphinxAtStartPar
style
&
\sphinxAtStartPar
Style of the plugin (pair, bond, fix, command, etc.)
\\
\sphinxhline
\sphinxAtStartPar
name
&
\sphinxAtStartPar
Name of the plugin style
\\
\sphinxhline
\sphinxAtStartPar
info
&
\sphinxAtStartPar
String with information about the plugin
\\
\sphinxhline
\sphinxAtStartPar
author
&
\sphinxAtStartPar
String with the name and email of the author
\\
\sphinxhline
\sphinxAtStartPar
creator.v1
&
\sphinxAtStartPar
Pointer to factory function for pair, bond, angle, dihedral, improper, kspace, or command styles
\\
\sphinxhline
\sphinxAtStartPar
creator.v2
&
\sphinxAtStartPar
Pointer to factory function for compute, fix, or region styles
\\
\sphinxhline
\sphinxAtStartPar
handle
&
\sphinxAtStartPar
Pointer to the open DSO file handle
\\
\sphinxbottomrule
\end{tabular}
\sphinxtableafterendhook\par
\sphinxattableend\end{savenotes}

\sphinxAtStartPar
Only one of the two alternate creator entries can be used at a time and
which of those is determined by the style of plugin. The “creator.v1”
element is for factory functions of supported styles computing forces
(i.e. pair, bond, angle, dihedral, or improper styles) or command styles
and the function takes as single argument the pointer to the LAMMPS
instance. The factory function is cast to the \sphinxcode{\sphinxupquote{lammpsplugin\_factory1}}
type before assignment.  The “creator.v2” element is for factory
functions creating an instance of a fix, compute, or region style and
takes three arguments: a pointer to the LAMMPS instance, an integer with
the length of the argument list and a \sphinxcode{\sphinxupquote{char **}} pointer to the list of
arguments. The factory function pointer needs to be cast to the
\sphinxcode{\sphinxupquote{lammpsplugin\_factory2}} type before assignment.


\subsection{Pair style example}
\label{\detokenize{Developer_plugins:pair-style-example}}
\sphinxAtStartPar
As an example, a hypothetical pair style plugin “morse2” implemented in
a class \sphinxcode{\sphinxupquote{PairMorse2}} in the files \sphinxcode{\sphinxupquote{pair\_morse2.h}} and
\sphinxcode{\sphinxupquote{pair\_morse2.cpp}} with the factory function and initialization
function would look like this:

\begin{sphinxVerbatim}[commandchars=\\\{\}]
\PYG{c+cp}{\PYGZsh{}}\PYG{c+cp}{include}\PYG{+w}{ }\PYG{c+cpf}{\PYGZdq{}lammpsplugin.h\PYGZdq{}}
\PYG{c+cp}{\PYGZsh{}}\PYG{c+cp}{include}\PYG{+w}{ }\PYG{c+cpf}{\PYGZdq{}version.h\PYGZdq{}}
\PYG{c+cp}{\PYGZsh{}}\PYG{c+cp}{include}\PYG{+w}{ }\PYG{c+cpf}{\PYGZdq{}pair\PYGZus{}morse2.h\PYGZdq{}}

\PYG{k}{using}\PYG{+w}{ }\PYG{k}{namespace}\PYG{+w}{ }\PYG{n+nn}{LAMMPS\PYGZus{}NS}\PYG{p}{;}

\PYG{k}{static}\PYG{+w}{ }\PYG{n}{Pair}\PYG{+w}{ }\PYG{o}{*}\PYG{n+nf}{morse2creator}\PYG{p}{(}\PYG{n}{LAMMPS}\PYG{+w}{ }\PYG{o}{*}\PYG{n}{lmp}\PYG{p}{)}
\PYG{p}{\PYGZob{}}
\PYG{+w}{  }\PYG{k}{return}\PYG{+w}{ }\PYG{k}{new}\PYG{+w}{ }\PYG{n}{PairMorse2}\PYG{p}{(}\PYG{n}{lmp}\PYG{p}{)}\PYG{p}{;}
\PYG{p}{\PYGZcb{}}

\PYG{k}{extern}\PYG{+w}{ }\PYG{l+s}{\PYGZdq{}}\PYG{l+s}{C}\PYG{l+s}{\PYGZdq{}}\PYG{+w}{ }\PYG{k+kt}{void}\PYG{+w}{ }\PYG{n}{lammpsplugin\PYGZus{}init}\PYG{p}{(}\PYG{k+kt}{void}\PYG{+w}{ }\PYG{o}{*}\PYG{n}{lmp}\PYG{p}{,}\PYG{+w}{ }\PYG{k+kt}{void}\PYG{+w}{ }\PYG{o}{*}\PYG{n}{handle}\PYG{p}{,}\PYG{+w}{ }\PYG{k+kt}{void}\PYG{+w}{ }\PYG{o}{*}\PYG{n}{regfunc}\PYG{p}{)}
\PYG{p}{\PYGZob{}}
\PYG{+w}{  }\PYG{n}{lammpsplugin\PYGZus{}regfunc}\PYG{+w}{ }\PYG{n}{register\PYGZus{}plugin}\PYG{+w}{ }\PYG{o}{=}\PYG{+w}{ }\PYG{p}{(}\PYG{n}{lammpsplugin\PYGZus{}regfunc}\PYG{p}{)}\PYG{+w}{ }\PYG{n}{regfunc}\PYG{p}{;}
\PYG{+w}{  }\PYG{n}{lammpsplugin\PYGZus{}t}\PYG{+w}{ }\PYG{n}{plugin}\PYG{p}{;}

\PYG{+w}{  }\PYG{n}{plugin}\PYG{p}{.}\PYG{n}{version}\PYG{+w}{ }\PYG{o}{=}\PYG{+w}{ }\PYG{n}{LAMMPS\PYGZus{}VERSION}\PYG{p}{;}
\PYG{+w}{  }\PYG{n}{plugin}\PYG{p}{.}\PYG{n}{style}\PYG{+w}{   }\PYG{o}{=}\PYG{+w}{ }\PYG{l+s}{\PYGZdq{}}\PYG{l+s}{pair}\PYG{l+s}{\PYGZdq{}}\PYG{p}{;}
\PYG{+w}{  }\PYG{n}{plugin}\PYG{p}{.}\PYG{n}{name}\PYG{+w}{    }\PYG{o}{=}\PYG{+w}{ }\PYG{l+s}{\PYGZdq{}}\PYG{l+s}{morse2}\PYG{l+s}{\PYGZdq{}}\PYG{p}{;}
\PYG{+w}{  }\PYG{n}{plugin}\PYG{p}{.}\PYG{n}{info}\PYG{+w}{    }\PYG{o}{=}\PYG{+w}{ }\PYG{l+s}{\PYGZdq{}}\PYG{l+s}{Morse2 variant pair style v1.0}\PYG{l+s}{\PYGZdq{}}\PYG{p}{;}
\PYG{+w}{  }\PYG{n}{plugin}\PYG{p}{.}\PYG{n}{author}\PYG{+w}{  }\PYG{o}{=}\PYG{+w}{ }\PYG{l+s}{\PYGZdq{}}\PYG{l+s}{Axel Kohlmeyer (akohlmey@gmail.com)}\PYG{l+s}{\PYGZdq{}}\PYG{p}{;}
\PYG{+w}{  }\PYG{n}{plugin}\PYG{p}{.}\PYG{n}{creator}\PYG{p}{.}\PYG{n}{v1}\PYG{+w}{ }\PYG{o}{=}\PYG{+w}{ }\PYG{p}{(}\PYG{n}{lammpsplugin\PYGZus{}factory1}\PYG{+w}{ }\PYG{o}{*}\PYG{p}{)}\PYG{+w}{ }\PYG{o}{\PYGZam{}}\PYG{n}{morse2creator}\PYG{p}{;}
\PYG{+w}{  }\PYG{n}{plugin}\PYG{p}{.}\PYG{n}{handle}\PYG{+w}{  }\PYG{o}{=}\PYG{+w}{ }\PYG{n}{handle}\PYG{p}{;}
\PYG{+w}{  }\PYG{p}{(}\PYG{o}{*}\PYG{n}{register\PYGZus{}plugin}\PYG{p}{)}\PYG{p}{(}\PYG{o}{\PYGZam{}}\PYG{n}{plugin}\PYG{p}{,}\PYG{n}{lmp}\PYG{p}{)}\PYG{p}{;}
\PYG{p}{\PYGZcb{}}
\end{sphinxVerbatim}

\sphinxAtStartPar
The factory function in this example is called \sphinxcode{\sphinxupquote{morse2creator()}}.  It
receives a pointer to the LAMMPS class as only argument and thus has to
be assigned to the \sphinxstyleemphasis{creator.v1} member of the plugin struct and cast to
the \sphinxcode{\sphinxupquote{lammpsplugin\_factory1}} function pointer type.  It returns a
pointer to the allocated class instance derived from the \sphinxcode{\sphinxupquote{Pair}} class.
This function may be declared static to avoid clashes with other
plugins.  The name of the derived class, \sphinxcode{\sphinxupquote{PairMorse2}}, however must be
unique inside the entire LAMMPS executable.


\subsection{Fix style example}
\label{\detokenize{Developer_plugins:fix-style-example}}
\sphinxAtStartPar
If the factory function is for a fix or compute, which take three
arguments (a pointer to the LAMMPS class, the number of arguments and the
list of argument strings), then the pointer type is \sphinxcode{\sphinxupquote{lammpsplugin\_factory2}}
and it must be assigned to the \sphinxstyleemphasis{creator.v2} member of the plugin struct.
Below is an example for that:

\begin{sphinxVerbatim}[commandchars=\\\{\}]
\PYG{c+cp}{\PYGZsh{}}\PYG{c+cp}{include}\PYG{+w}{ }\PYG{c+cpf}{\PYGZdq{}lammpsplugin.h\PYGZdq{}}
\PYG{c+cp}{\PYGZsh{}}\PYG{c+cp}{include}\PYG{+w}{ }\PYG{c+cpf}{\PYGZdq{}version.h\PYGZdq{}}
\PYG{c+cp}{\PYGZsh{}}\PYG{c+cp}{include}\PYG{+w}{ }\PYG{c+cpf}{\PYGZdq{}fix\PYGZus{}nve2.h\PYGZdq{}}

\PYG{k}{using}\PYG{+w}{ }\PYG{k}{namespace}\PYG{+w}{ }\PYG{n+nn}{LAMMPS\PYGZus{}NS}\PYG{p}{;}

\PYG{k}{static}\PYG{+w}{ }\PYG{n}{Fix}\PYG{+w}{ }\PYG{o}{*}\PYG{n+nf}{nve2creator}\PYG{p}{(}\PYG{n}{LAMMPS}\PYG{+w}{ }\PYG{o}{*}\PYG{n}{lmp}\PYG{p}{,}\PYG{+w}{ }\PYG{k+kt}{int}\PYG{+w}{ }\PYG{n}{argc}\PYG{p}{,}\PYG{+w}{ }\PYG{k+kt}{char}\PYG{+w}{ }\PYG{o}{*}\PYG{o}{*}\PYG{n}{argv}\PYG{p}{)}
\PYG{p}{\PYGZob{}}
\PYG{+w}{  }\PYG{k}{return}\PYG{+w}{ }\PYG{k}{new}\PYG{+w}{ }\PYG{n}{FixNVE2}\PYG{p}{(}\PYG{n}{lmp}\PYG{p}{,}\PYG{n}{argc}\PYG{p}{,}\PYG{n}{argv}\PYG{p}{)}\PYG{p}{;}
\PYG{p}{\PYGZcb{}}

\PYG{k}{extern}\PYG{+w}{ }\PYG{l+s}{\PYGZdq{}}\PYG{l+s}{C}\PYG{l+s}{\PYGZdq{}}\PYG{+w}{ }\PYG{k+kt}{void}\PYG{+w}{ }\PYG{n}{lammpsplugin\PYGZus{}init}\PYG{p}{(}\PYG{k+kt}{void}\PYG{+w}{ }\PYG{o}{*}\PYG{n}{lmp}\PYG{p}{,}\PYG{+w}{ }\PYG{k+kt}{void}\PYG{+w}{ }\PYG{o}{*}\PYG{n}{handle}\PYG{p}{,}\PYG{+w}{ }\PYG{k+kt}{void}\PYG{+w}{ }\PYG{o}{*}\PYG{n}{regfunc}\PYG{p}{)}
\PYG{p}{\PYGZob{}}
\PYG{+w}{  }\PYG{n}{lammpsplugin\PYGZus{}regfunc}\PYG{+w}{ }\PYG{n}{register\PYGZus{}plugin}\PYG{+w}{ }\PYG{o}{=}\PYG{+w}{ }\PYG{p}{(}\PYG{n}{lammpsplugin\PYGZus{}regfunc}\PYG{p}{)}\PYG{+w}{ }\PYG{n}{regfunc}\PYG{p}{;}
\PYG{+w}{  }\PYG{n}{lammpsplugin\PYGZus{}t}\PYG{+w}{ }\PYG{n}{plugin}\PYG{p}{;}

\PYG{+w}{  }\PYG{n}{plugin}\PYG{p}{.}\PYG{n}{version}\PYG{+w}{ }\PYG{o}{=}\PYG{+w}{ }\PYG{n}{LAMMPS\PYGZus{}VERSION}\PYG{p}{;}
\PYG{+w}{  }\PYG{n}{plugin}\PYG{p}{.}\PYG{n}{style}\PYG{+w}{   }\PYG{o}{=}\PYG{+w}{ }\PYG{l+s}{\PYGZdq{}}\PYG{l+s}{fix}\PYG{l+s}{\PYGZdq{}}\PYG{p}{;}
\PYG{+w}{  }\PYG{n}{plugin}\PYG{p}{.}\PYG{n}{name}\PYG{+w}{    }\PYG{o}{=}\PYG{+w}{ }\PYG{l+s}{\PYGZdq{}}\PYG{l+s}{nve2}\PYG{l+s}{\PYGZdq{}}\PYG{p}{;}
\PYG{+w}{  }\PYG{n}{plugin}\PYG{p}{.}\PYG{n}{info}\PYG{+w}{    }\PYG{o}{=}\PYG{+w}{ }\PYG{l+s}{\PYGZdq{}}\PYG{l+s}{NVE2 variant fix style v1.0}\PYG{l+s}{\PYGZdq{}}\PYG{p}{;}
\PYG{+w}{  }\PYG{n}{plugin}\PYG{p}{.}\PYG{n}{author}\PYG{+w}{  }\PYG{o}{=}\PYG{+w}{ }\PYG{l+s}{\PYGZdq{}}\PYG{l+s}{Axel Kohlmeyer (akohlmey@gmail.com)}\PYG{l+s}{\PYGZdq{}}\PYG{p}{;}
\PYG{+w}{  }\PYG{n}{plugin}\PYG{p}{.}\PYG{n}{creator}\PYG{p}{.}\PYG{n}{v2}\PYG{+w}{ }\PYG{o}{=}\PYG{+w}{ }\PYG{p}{(}\PYG{n}{lammpsplugin\PYGZus{}factory2}\PYG{+w}{ }\PYG{o}{*}\PYG{p}{)}\PYG{+w}{ }\PYG{o}{\PYGZam{}}\PYG{n}{nve2creator}\PYG{p}{;}
\PYG{+w}{  }\PYG{n}{plugin}\PYG{p}{.}\PYG{n}{handle}\PYG{+w}{  }\PYG{o}{=}\PYG{+w}{ }\PYG{n}{handle}\PYG{p}{;}
\PYG{+w}{  }\PYG{p}{(}\PYG{o}{*}\PYG{n}{register\PYGZus{}plugin}\PYG{p}{)}\PYG{p}{(}\PYG{o}{\PYGZam{}}\PYG{n}{plugin}\PYG{p}{,}\PYG{n}{lmp}\PYG{p}{)}\PYG{p}{;}
\PYG{p}{\PYGZcb{}}
\end{sphinxVerbatim}


\subsection{Command style example}
\label{\detokenize{Developer_plugins:command-style-example}}
\sphinxAtStartPar
Command styles also use the first variant of factory function as
demonstrated in the following example, which also shows that the
implementation of the plugin class may be within the same source
file as the plugin interface code:

\begin{sphinxVerbatim}[commandchars=\\\{\}]
\PYG{c+cp}{\PYGZsh{}}\PYG{c+cp}{include}\PYG{+w}{ }\PYG{c+cpf}{\PYGZdq{}lammpsplugin.h\PYGZdq{}}

\PYG{c+cp}{\PYGZsh{}}\PYG{c+cp}{include}\PYG{+w}{ }\PYG{c+cpf}{\PYGZdq{}comm.h\PYGZdq{}}
\PYG{c+cp}{\PYGZsh{}}\PYG{c+cp}{include}\PYG{+w}{ }\PYG{c+cpf}{\PYGZdq{}error.h\PYGZdq{}}
\PYG{c+cp}{\PYGZsh{}}\PYG{c+cp}{include}\PYG{+w}{ }\PYG{c+cpf}{\PYGZdq{}command.h\PYGZdq{}}
\PYG{c+cp}{\PYGZsh{}}\PYG{c+cp}{include}\PYG{+w}{ }\PYG{c+cpf}{\PYGZdq{}version.h\PYGZdq{}}

\PYG{c+cp}{\PYGZsh{}}\PYG{c+cp}{include}\PYG{+w}{ }\PYG{c+cpf}{\PYGZlt{}cstring\PYGZgt{}}

\PYG{k}{namespace}\PYG{+w}{ }\PYG{n+nn}{LAMMPS\PYGZus{}NS}\PYG{+w}{ }\PYG{p}{\PYGZob{}}
\PYG{+w}{  }\PYG{k}{class}\PYG{+w}{ }\PYG{n+nc}{Hello}\PYG{+w}{ }\PYG{o}{:}\PYG{+w}{ }\PYG{k}{public}\PYG{+w}{ }\PYG{n}{Command}\PYG{+w}{ }\PYG{p}{\PYGZob{}}
\PYG{+w}{   }\PYG{k}{public}\PYG{o}{:}
\PYG{+w}{    }\PYG{n}{Hello}\PYG{p}{(}\PYG{k}{class}\PYG{+w}{ }\PYG{n+nc}{LAMMPS}\PYG{+w}{ }\PYG{o}{*}\PYG{n}{lmp}\PYG{p}{)}\PYG{+w}{ }\PYG{o}{:}\PYG{+w}{ }\PYG{n}{Command}\PYG{p}{(}\PYG{n}{lmp}\PYG{p}{)}\PYG{+w}{ }\PYG{p}{\PYGZob{}}\PYG{p}{\PYGZcb{}}\PYG{p}{;}
\PYG{+w}{    }\PYG{k+kt}{void}\PYG{+w}{ }\PYG{n+nf}{command}\PYG{p}{(}\PYG{k+kt}{int}\PYG{p}{,}\PYG{+w}{ }\PYG{k+kt}{char}\PYG{+w}{ }\PYG{o}{*}\PYG{o}{*}\PYG{p}{)}\PYG{p}{;}
\PYG{+w}{  }\PYG{p}{\PYGZcb{}}\PYG{p}{;}
\PYG{p}{\PYGZcb{}}

\PYG{k}{using}\PYG{+w}{ }\PYG{k}{namespace}\PYG{+w}{ }\PYG{n+nn}{LAMMPS\PYGZus{}NS}\PYG{p}{;}

\PYG{k+kt}{void}\PYG{+w}{ }\PYG{n+nf}{Hello::command}\PYG{p}{(}\PYG{k+kt}{int}\PYG{+w}{ }\PYG{n}{argc}\PYG{p}{,}\PYG{+w}{ }\PYG{k+kt}{char}\PYG{+w}{ }\PYG{o}{*}\PYG{o}{*}\PYG{n}{argv}\PYG{p}{)}
\PYG{p}{\PYGZob{}}
\PYG{+w}{   }\PYG{k}{if}\PYG{+w}{ }\PYG{p}{(}\PYG{n}{argc}\PYG{+w}{ }\PYG{o}{!}\PYG{o}{=}\PYG{+w}{ }\PYG{l+m+mi}{1}\PYG{p}{)}\PYG{+w}{ }\PYG{n}{error}\PYG{o}{\PYGZhy{}}\PYG{o}{\PYGZgt{}}\PYG{n}{all}\PYG{p}{(}\PYG{n}{FLERR}\PYG{p}{,}\PYG{l+s}{\PYGZdq{}}\PYG{l+s}{Illegal hello command}\PYG{l+s}{\PYGZdq{}}\PYG{p}{)}\PYG{p}{;}
\PYG{+w}{   }\PYG{k}{if}\PYG{+w}{ }\PYG{p}{(}\PYG{n}{comm}\PYG{o}{\PYGZhy{}}\PYG{o}{\PYGZgt{}}\PYG{n}{me}\PYG{+w}{ }\PYG{o}{=}\PYG{o}{=}\PYG{+w}{ }\PYG{l+m+mi}{0}\PYG{p}{)}
\PYG{+w}{     }\PYG{n}{utils}\PYG{o}{:}\PYG{o}{:}\PYG{n}{logmesg}\PYG{p}{(}\PYG{n}{lmp}\PYG{p}{,}\PYG{n}{fmt}\PYG{o}{:}\PYG{o}{:}\PYG{n}{format}\PYG{p}{(}\PYG{l+s}{\PYGZdq{}}\PYG{l+s}{Hello, \PYGZob{}\PYGZcb{}!}\PYG{l+s+se}{\PYGZbs{}n}\PYG{l+s}{\PYGZdq{}}\PYG{p}{,}\PYG{n}{argv}\PYG{p}{[}\PYG{l+m+mi}{0}\PYG{p}{]}\PYG{p}{)}\PYG{p}{)}\PYG{p}{;}
\PYG{p}{\PYGZcb{}}

\PYG{k}{static}\PYG{+w}{ }\PYG{k+kt}{void}\PYG{+w}{ }\PYG{n+nf}{hellocreator}\PYG{p}{(}\PYG{n}{LAMMPS}\PYG{+w}{ }\PYG{o}{*}\PYG{n}{lmp}\PYG{p}{)}
\PYG{p}{\PYGZob{}}
\PYG{+w}{  }\PYG{k}{return}\PYG{+w}{ }\PYG{k}{new}\PYG{+w}{ }\PYG{n}{Hello}\PYG{p}{(}\PYG{n}{lmp}\PYG{p}{)}\PYG{p}{;}
\PYG{p}{\PYGZcb{}}

\PYG{k}{extern}\PYG{+w}{ }\PYG{l+s}{\PYGZdq{}}\PYG{l+s}{C}\PYG{l+s}{\PYGZdq{}}\PYG{+w}{ }\PYG{k+kt}{void}\PYG{+w}{ }\PYG{n}{lammpsplugin\PYGZus{}init}\PYG{p}{(}\PYG{k+kt}{void}\PYG{+w}{ }\PYG{o}{*}\PYG{n}{lmp}\PYG{p}{,}\PYG{+w}{ }\PYG{k+kt}{void}\PYG{+w}{ }\PYG{o}{*}\PYG{n}{handle}\PYG{p}{,}\PYG{+w}{ }\PYG{k+kt}{void}\PYG{+w}{ }\PYG{o}{*}\PYG{n}{regfunc}\PYG{p}{)}
\PYG{p}{\PYGZob{}}
\PYG{+w}{  }\PYG{n}{lammpsplugin\PYGZus{}t}\PYG{+w}{ }\PYG{n}{plugin}\PYG{p}{;}
\PYG{+w}{  }\PYG{n}{lammpsplugin\PYGZus{}regfunc}\PYG{+w}{ }\PYG{n}{register\PYGZus{}plugin}\PYG{+w}{ }\PYG{o}{=}\PYG{+w}{ }\PYG{p}{(}\PYG{n}{lammpsplugin\PYGZus{}regfunc}\PYG{p}{)}\PYG{+w}{ }\PYG{n}{regfunc}\PYG{p}{;}

\PYG{+w}{  }\PYG{n}{plugin}\PYG{p}{.}\PYG{n}{version}\PYG{+w}{ }\PYG{o}{=}\PYG{+w}{ }\PYG{n}{LAMMPS\PYGZus{}VERSION}\PYG{p}{;}
\PYG{+w}{  }\PYG{n}{plugin}\PYG{p}{.}\PYG{n}{style}\PYG{+w}{   }\PYG{o}{=}\PYG{+w}{ }\PYG{l+s}{\PYGZdq{}}\PYG{l+s}{command}\PYG{l+s}{\PYGZdq{}}\PYG{p}{;}
\PYG{+w}{  }\PYG{n}{plugin}\PYG{p}{.}\PYG{n}{name}\PYG{+w}{    }\PYG{o}{=}\PYG{+w}{ }\PYG{l+s}{\PYGZdq{}}\PYG{l+s}{hello}\PYG{l+s}{\PYGZdq{}}\PYG{p}{;}
\PYG{+w}{  }\PYG{n}{plugin}\PYG{p}{.}\PYG{n}{info}\PYG{+w}{    }\PYG{o}{=}\PYG{+w}{ }\PYG{l+s}{\PYGZdq{}}\PYG{l+s}{Hello world command v1.1}\PYG{l+s}{\PYGZdq{}}\PYG{p}{;}
\PYG{+w}{  }\PYG{n}{plugin}\PYG{p}{.}\PYG{n}{author}\PYG{+w}{  }\PYG{o}{=}\PYG{+w}{ }\PYG{l+s}{\PYGZdq{}}\PYG{l+s}{Axel Kohlmeyer (akohlmey@gmail.com)}\PYG{l+s}{\PYGZdq{}}\PYG{p}{;}
\PYG{+w}{  }\PYG{n}{plugin}\PYG{p}{.}\PYG{n}{creator}\PYG{p}{.}\PYG{n}{v1}\PYG{+w}{ }\PYG{o}{=}\PYG{+w}{ }\PYG{p}{(}\PYG{n}{lammpsplugin\PYGZus{}factory1}\PYG{+w}{ }\PYG{o}{*}\PYG{p}{)}\PYG{+w}{ }\PYG{o}{\PYGZam{}}\PYG{n}{hellocreator}\PYG{p}{;}
\PYG{+w}{  }\PYG{n}{plugin}\PYG{p}{.}\PYG{n}{handle}\PYG{+w}{  }\PYG{o}{=}\PYG{+w}{ }\PYG{n}{handle}\PYG{p}{;}
\PYG{+w}{  }\PYG{p}{(}\PYG{o}{*}\PYG{n}{register\PYGZus{}plugin}\PYG{p}{)}\PYG{p}{(}\PYG{o}{\PYGZam{}}\PYG{n}{plugin}\PYG{p}{,}\PYG{n}{lmp}\PYG{p}{)}\PYG{p}{;}
\PYG{p}{\PYGZcb{}}
\end{sphinxVerbatim}


\subsection{Additional Details}
\label{\detokenize{Developer_plugins:additional-details}}
\sphinxAtStartPar
The initialization function \sphinxstylestrong{must} be called \sphinxcode{\sphinxupquote{lammpsplugin\_init}}, it
\sphinxstylestrong{must} have C bindings and it takes three void pointers as arguments.
The first is a pointer to the LAMMPS class that calls it and it needs to
be passed to the registration function.  The second argument is a
pointer to the internal handle of the DSO file, this needs to be added
to the plugin info struct, so that the DSO can be closed and unloaded
when all its contained plugins are unloaded.  The third argument is a
function pointer to the registration function and needs to be stored
in a variable of \sphinxcode{\sphinxupquote{lammpsplugin\_regfunc}} type and then called with a
pointer to the \sphinxcode{\sphinxupquote{lammpsplugin\_t}} struct and the pointer to the LAMMPS
instance as arguments to register a single plugin.  There may be multiple
calls to multiple plugins in the same initialization function.

\sphinxAtStartPar
To register a plugin a struct of the \sphinxcode{\sphinxupquote{lammpsplugin\_t}} needs to be filled
with relevant info: current LAMMPS version string, kind of style, name of
style, info string, author string, pointer to factory function, and the
DSO handle.  The registration function is called with a pointer to the address
of this struct and the pointer of the LAMMPS class.  The registration function
will then add the factory function of the plugin style to the respective
style map under the provided name.  It will also make a copy of the struct
in a list of all loaded plugins and update the reference counter for loaded
plugins from this specific DSO file.

\sphinxAtStartPar
The pair style itself (i.e. the PairMorse2 class in this example) can be
written just like any other pair style that is included in LAMMPS.  For
a plugin, the use of the \sphinxcode{\sphinxupquote{PairStyle}} macro in the section encapsulated
by \sphinxcode{\sphinxupquote{\#ifdef PAIR\_CLASS}} is not needed, since the mapping of the class
name to the style name is done by the plugin registration function with
the information from the \sphinxcode{\sphinxupquote{lammpsplugin\_t}} struct.  It may be included
in case the new code is intended to be later included in LAMMPS directly.

\sphinxAtStartPar
A plugin may be registered under an existing style name.  In that case
the plugin will override the existing code.  This can be used to modify
the behavior of existing styles or to debug new versions of them without
having to re\sphinxhyphen{}compile or re\sphinxhyphen{}install all of LAMMPS.


\subsection{Compiling plugins}
\label{\detokenize{Developer_plugins:compiling-plugins}}
\sphinxAtStartPar
Plugins need to be compiled with the same compilers and libraries
(e.g. MPI) and compilation settings (MPI on/off, OpenMP, integer sizes)
as the LAMMPS executable and library.  Otherwise the plugin will likely
not load due to mismatches in the function signatures (LAMMPS is C++ so
scope, type, and number of arguments are encoded into the symbol names
and thus differences in them will lead to failed plugin load commands).
Compilation of the plugin can be managed via both, CMake or traditional
GNU makefiles.  Some examples that can be used as a template are in the
\sphinxcode{\sphinxupquote{examples/plugins}} folder.  The CMake script code has some small
adjustments to allow building the plugins for running unit tests with
them.

\sphinxAtStartPar
Another example that converts the KIM package into a plugin can be found
in the \sphinxcode{\sphinxupquote{examples/kim/plugin}} folder.  No changes to the sources of the
KIM package themselves are needed; only the plugin interface and loader
code needs to be added.  This example only supports building with CMake,
but is probably a more typical example. To compile you need to run CMake
with \sphinxcode{\sphinxupquote{\sphinxhyphen{}DLAMMPS\_SOURCE\_DIR=\textless{}path/to/lammps/src/folder\textgreater{}}}.  Other
configuration setting are identical to those for compiling LAMMPS.

\sphinxAtStartPar
A second example for a plugin from a package is in the
\sphinxcode{\sphinxupquote{examples/PACKAGES/pace/plugin}} folder that will create a plugin from
the ML\sphinxhyphen{}PACE package.  In this case the bulk of the code is in a static
external library that is being downloaded and compiled first and then
combined with the pair style wrapper and the plugin loader.  This
example also contains a NSIS script that can be used to create an
Installer package for Windows (the mutual licensing terms of the
external library and LAMMPS conflict when distributing binaries, so the
ML\sphinxhyphen{}PACE package cannot be linked statically, but the LAMMPS headers
required to build the plugin are also available under a less restrictive
license).  This will automatically set the required environment variable
and launching a (compatible) LAMMPS binary will load and register the
plugin and the ML\sphinxhyphen{}PACE package can then be used as it was linked into
LAMMPS.

\sphinxstepscope


\section{Adding tests for unit testing}
\label{\detokenize{Developer_unittest:adding-tests-for-unit-testing}}\label{\detokenize{Developer_unittest::doc}}
\sphinxAtStartPar
This section discusses adding or expanding tests for the unit test
infrastructure included into the LAMMPS source code distribution.
Unlike example inputs, unit tests focus on testing the “local” behavior
of individual features, tend to run fast, and should be set up to cover
as much of the added code as possible.  When contributing code to the
distribution, the LAMMPS developers will appreciate if additions to the
integrated unit test facility are included.

\sphinxAtStartPar
Given the complex nature of MD simulations where many operations can
only be performed when suitable “real” simulation environment has been
set up, not all tests will be unit tests in the strict definition of
the term.  They are rather executed on a more abstract level by issuing
LAMMPS script commands and then inspecting the changes to the internal
data.  For some classes of tests, generic test programs have been
written that can be applied to parts of LAMMPS that use the same
interface (via polymorphism) and those are driven by input files, so
tests can be added by simply adding more of those input files.  Those
tests should be seen more as a hybrid between unit and regression tests.

\sphinxAtStartPar
When adding tests it is recommended to also {\hyperref[\detokenize{Build_development:testing}]{\sphinxcrossref{\DUrole{std,std-ref}{enable support for
code coverage reporting}}}}, and study the coverage reports
so that it is possible to monitor which parts of the code of a given
file are executed during the tests and which tests would need to be
added to increase the coverage.

\sphinxAtStartPar
The tests are grouped into categories and corresponding folders.
The following sections describe how the tests are implemented and
executed in those categories with increasing complexity of tests
and implementation.


\subsection{Tests for utility functions}
\label{\detokenize{Developer_unittest:tests-for-utility-functions}}
\sphinxAtStartPar
These tests are driven by programs in the \sphinxcode{\sphinxupquote{unittest/utils}} folder
and most closely resemble conventional unit tests. There is one test
program for each namespace or group of classes or file. The naming
convention for the sources and executables is that they start with
with \sphinxcode{\sphinxupquote{test\_}}.  The following sources and groups of tests are currently
available:


\begin{savenotes}\sphinxattablestart
\sphinxthistablewithglobalstyle
\raggedright
\begin{tabular}[t]{\X{32}{100}\X{18}{100}\X{50}{100}}
\sphinxtoprule
\sphinxstyletheadfamily 
\sphinxAtStartPar
File name:
&\sphinxstyletheadfamily 
\sphinxAtStartPar
Test name:
&\sphinxstyletheadfamily 
\sphinxAtStartPar
Description:
\\
\sphinxmidrule
\sphinxtableatstartofbodyhook
\sphinxAtStartPar
\sphinxcode{\sphinxupquote{test\_argutils.cpp}}
&
\sphinxAtStartPar
ArgInfo
&
\sphinxAtStartPar
Tests for \sphinxcode{\sphinxupquote{ArgInfo}} class used by LAMMPS
\\
\sphinxhline
\sphinxAtStartPar
\sphinxcode{\sphinxupquote{test\_fmtlib.cpp}}
&
\sphinxAtStartPar
FmtLib
&
\sphinxAtStartPar
Tests for \sphinxcode{\sphinxupquote{fmtlib::}} functions used by LAMMPS
\\
\sphinxhline
\sphinxAtStartPar
\sphinxcode{\sphinxupquote{test\_math\_eigen\_impl.cpp}}
&
\sphinxAtStartPar
MathEigen
&
\sphinxAtStartPar
Tests for \sphinxcode{\sphinxupquote{MathEigen::}} classes and functions
\\
\sphinxhline
\sphinxAtStartPar
\sphinxcode{\sphinxupquote{test\_mempool.cpp}}
&
\sphinxAtStartPar
MemPool
&
\sphinxAtStartPar
Tests for {\hyperref[\detokenize{Developer_utils:_CPPv4I0EN9LAMMPS_NS6MyPageE}]{\sphinxcrossref{\sphinxcode{\sphinxupquote{MyPage}}}}} and {\hyperref[\detokenize{Developer_utils:_CPPv4I0EN9LAMMPS_NS11MyPoolChunkE}]{\sphinxcrossref{\sphinxcode{\sphinxupquote{MyPoolChunk}}}}}
\\
\sphinxhline
\sphinxAtStartPar
\sphinxcode{\sphinxupquote{test\_tokenizer.cpp}}
&
\sphinxAtStartPar
Tokenizer
&
\sphinxAtStartPar
Tests for {\hyperref[\detokenize{Developer_utils:_CPPv4N9LAMMPS_NS9TokenizerE}]{\sphinxcrossref{\sphinxcode{\sphinxupquote{Tokenizer}}}}} and {\hyperref[\detokenize{Developer_utils:_CPPv4N9LAMMPS_NS14ValueTokenizerE}]{\sphinxcrossref{\sphinxcode{\sphinxupquote{ValueTokenizer}}}}}
\\
\sphinxhline
\sphinxAtStartPar
\sphinxcode{\sphinxupquote{test\_utils.cpp}}
&
\sphinxAtStartPar
Utils
&
\sphinxAtStartPar
Tests for \sphinxcode{\sphinxupquote{utils::}} {\hyperref[\detokenize{Developer_utils::doc}]{\sphinxcrossref{\DUrole{doc}{functions}}}}
\\
\sphinxbottomrule
\end{tabular}
\sphinxtableafterendhook\par
\sphinxattableend\end{savenotes}

\sphinxAtStartPar
To add tests either an existing source file needs to be modified or a
new source file needs to be added to the distribution and enabled for
testing.  To add a new file suitable CMake script code needs to be added
to the \sphinxcode{\sphinxupquote{CMakeLists.txt}} file in the \sphinxcode{\sphinxupquote{unittest/utils}} folder.  Example:

\begin{sphinxVerbatim}[commandchars=\\\{\}]
\PYG{n+nb}{add\PYGZus{}executable}\PYG{p}{(}\PYG{l+s}{test\PYGZus{}tokenizer}\PYG{+w}{ }\PYG{l+s}{test\PYGZus{}tokenizer.cpp}\PYG{p}{)}
\PYG{n+nb}{target\PYGZus{}link\PYGZus{}libraries}\PYG{p}{(}\PYG{l+s}{test\PYGZus{}tokenizer}\PYG{+w}{ }\PYG{l+s}{PRIVATE}\PYG{+w}{ }\PYG{l+s}{lammps}\PYG{+w}{ }\PYG{l+s}{GTest::GMockMain}\PYG{+w}{ }\PYG{l+s}{GTest::GMock}\PYG{+w}{ }\PYG{l+s}{GTest::GTest}\PYG{p}{)}
\PYG{n+nb}{add\PYGZus{}test}\PYG{p}{(}\PYG{l+s}{Tokenizer}\PYG{+w}{ }\PYG{l+s}{test\PYGZus{}tokenizer}\PYG{p}{)}
\end{sphinxVerbatim}

\sphinxAtStartPar
This adds instructions to build the \sphinxcode{\sphinxupquote{test\_tokenizer}} executable from
\sphinxcode{\sphinxupquote{test\_tokenizer.cpp}} and links it with the GoogleTest libraries and the
LAMMPS library as well as it uses the \sphinxcode{\sphinxupquote{main()}} function from the
GoogleMock library of GoogleTest.  The third line registers the executable
as a test program to be run from \sphinxcode{\sphinxupquote{ctest}} under the name \sphinxcode{\sphinxupquote{Tokenizer}}.

\sphinxAtStartPar
The test executable itself will execute multiple individual tests
through the GoogleTest framework. In this case each test consists of
creating a tokenizer class instance with a given string and explicit or
default separator choice, and then executing member functions of the
class and comparing their results with expected values. A few examples:

\begin{sphinxVerbatim}[commandchars=\\\{\}]
\PYG{n}{TEST}\PYG{p}{(}\PYG{n}{Tokenizer}\PYG{p}{,}\PYG{+w}{ }\PYG{n}{empty\PYGZus{}string}\PYG{p}{)}
\PYG{p}{\PYGZob{}}
\PYG{+w}{    }\PYG{n}{Tokenizer}\PYG{+w}{ }\PYG{n}{t}\PYG{p}{(}\PYG{l+s}{\PYGZdq{}}\PYG{l+s}{\PYGZdq{}}\PYG{p}{,}\PYG{+w}{ }\PYG{l+s}{\PYGZdq{}}\PYG{l+s}{ }\PYG{l+s}{\PYGZdq{}}\PYG{p}{)}\PYG{p}{;}
\PYG{+w}{    }\PYG{n}{ASSERT\PYGZus{}EQ}\PYG{p}{(}\PYG{n}{t}\PYG{p}{.}\PYG{n}{count}\PYG{p}{(}\PYG{p}{)}\PYG{p}{,}\PYG{+w}{ }\PYG{l+m+mi}{0}\PYG{p}{)}\PYG{p}{;}
\PYG{p}{\PYGZcb{}}

\PYG{n}{TEST}\PYG{p}{(}\PYG{n}{Tokenizer}\PYG{p}{,}\PYG{+w}{ }\PYG{n}{two\PYGZus{}words}\PYG{p}{)}
\PYG{p}{\PYGZob{}}
\PYG{+w}{    }\PYG{n}{Tokenizer}\PYG{+w}{ }\PYG{n}{t}\PYG{p}{(}\PYG{l+s}{\PYGZdq{}}\PYG{l+s}{test word}\PYG{l+s}{\PYGZdq{}}\PYG{p}{,}\PYG{+w}{ }\PYG{l+s}{\PYGZdq{}}\PYG{l+s}{ }\PYG{l+s}{\PYGZdq{}}\PYG{p}{)}\PYG{p}{;}
\PYG{+w}{    }\PYG{n}{ASSERT\PYGZus{}EQ}\PYG{p}{(}\PYG{n}{t}\PYG{p}{.}\PYG{n}{count}\PYG{p}{(}\PYG{p}{)}\PYG{p}{,}\PYG{+w}{ }\PYG{l+m+mi}{2}\PYG{p}{)}\PYG{p}{;}
\PYG{p}{\PYGZcb{}}

\PYG{n}{TEST}\PYG{p}{(}\PYG{n}{Tokenizer}\PYG{p}{,}\PYG{+w}{ }\PYG{n}{default\PYGZus{}separators}\PYG{p}{)}
\PYG{p}{\PYGZob{}}
\PYG{+w}{    }\PYG{n}{Tokenizer}\PYG{+w}{ }\PYG{n}{t}\PYG{p}{(}\PYG{l+s}{\PYGZdq{}}\PYG{l+s}{ }\PYG{l+s+se}{\PYGZbs{}r}\PYG{l+s+se}{\PYGZbs{}n}\PYG{l+s}{ test }\PYG{l+s+se}{\PYGZbs{}t}\PYG{l+s}{ word }\PYG{l+s+se}{\PYGZbs{}f}\PYG{l+s}{\PYGZdq{}}\PYG{p}{)}\PYG{p}{;}
\PYG{+w}{    }\PYG{n}{ASSERT\PYGZus{}THAT}\PYG{p}{(}\PYG{n}{t}\PYG{p}{.}\PYG{n}{next}\PYG{p}{(}\PYG{p}{)}\PYG{p}{,}\PYG{+w}{ }\PYG{n}{Eq}\PYG{p}{(}\PYG{l+s}{\PYGZdq{}}\PYG{l+s}{test}\PYG{l+s}{\PYGZdq{}}\PYG{p}{)}\PYG{p}{)}\PYG{p}{;}
\PYG{+w}{    }\PYG{n}{ASSERT\PYGZus{}THAT}\PYG{p}{(}\PYG{n}{t}\PYG{p}{.}\PYG{n}{next}\PYG{p}{(}\PYG{p}{)}\PYG{p}{,}\PYG{+w}{ }\PYG{n}{Eq}\PYG{p}{(}\PYG{l+s}{\PYGZdq{}}\PYG{l+s}{word}\PYG{l+s}{\PYGZdq{}}\PYG{p}{)}\PYG{p}{)}\PYG{p}{;}
\PYG{+w}{    }\PYG{n}{ASSERT\PYGZus{}EQ}\PYG{p}{(}\PYG{n}{t}\PYG{p}{.}\PYG{n}{count}\PYG{p}{(}\PYG{p}{)}\PYG{p}{,}\PYG{+w}{ }\PYG{l+m+mi}{2}\PYG{p}{)}\PYG{p}{;}
\PYG{p}{\PYGZcb{}}
\end{sphinxVerbatim}

\sphinxAtStartPar
Each of these TEST functions will become an individual
test run by the test program. When using the \sphinxcode{\sphinxupquote{ctest}}
command as a front end to run the tests, their output
will be suppressed and only a summary printed, but adding
the ‘\sphinxhyphen{}V’ option will then produce output from the tests
above like the following:

\begin{sphinxVerbatim}[commandchars=\\\{\}]
\PYG{g+go}{[...]}
\PYG{g+go}{1: [ RUN      ] Tokenizer.empty\PYGZus{}string}
\PYG{g+go}{1: [       OK ] Tokenizer.empty\PYGZus{}string (0 ms)}
\PYG{g+go}{1: [ RUN      ] Tokenizer.two\PYGZus{}words}
\PYG{g+go}{1: [       OK ] Tokenizer.two\PYGZus{}words (0 ms)}
\PYG{g+go}{1: [ RUN      ] Tokenizer.default\PYGZus{}separators}
\PYG{g+go}{1: [       OK ] Tokenizer.default\PYGZus{}separators (0 ms)}
\PYG{g+go}{[...]}
\end{sphinxVerbatim}

\sphinxAtStartPar
The MathEigen test collection has been adapted from a standalone test
and does not use the GoogleTest framework and thus not representative.
The other test sources, however, can serve as guiding examples for
additional tests.


\subsection{Tests for individual LAMMPS commands}
\label{\detokenize{Developer_unittest:tests-for-individual-lammps-commands}}
\sphinxAtStartPar
The tests \sphinxcode{\sphinxupquote{unittest/commands}} are a bit more complex as they require
to first create a {\hyperref[\detokenize{Classes_lammps:_CPPv4N9LAMMPS_NS6LAMMPSE}]{\sphinxcrossref{\sphinxcode{\sphinxupquote{LAMMPS}}}}} class instance
and then use the {\hyperref[\detokenize{Cplusplus::doc}]{\sphinxcrossref{\DUrole{doc}{C++ API}}}} to pass individual commands
to that LAMMPS instance.  For that reason these tests use a GoogleTest
“test fixture”, i.e. a class derived from \sphinxcode{\sphinxupquote{testing::Test}} that will
create (and delete) the required LAMMPS class instance for each set of
tests in a \sphinxcode{\sphinxupquote{TEST\_F()}} function.  Please see the individual source files
for different examples of setting up suitable test fixtures.  Here is an
example for implementing a test using a fixture by first checking the
default value and then issuing LAMMPS commands and checking whether they
have the desired effect:

\begin{sphinxVerbatim}[commandchars=\\\{\}]
\PYG{n}{TEST\PYGZus{}F}\PYG{p}{(}\PYG{n}{SimpleCommandsTest}\PYG{p}{,}\PYG{+w}{ }\PYG{n}{ResetTimestep}\PYG{p}{)}
\PYG{p}{\PYGZob{}}
\PYG{+w}{    }\PYG{n}{ASSERT\PYGZus{}EQ}\PYG{p}{(}\PYG{n}{lmp}\PYG{o}{\PYGZhy{}}\PYG{o}{\PYGZgt{}}\PYG{n}{update}\PYG{o}{\PYGZhy{}}\PYG{o}{\PYGZgt{}}\PYG{n}{ntimestep}\PYG{p}{,}\PYG{+w}{ }\PYG{l+m+mi}{0}\PYG{p}{)}\PYG{p}{;}

\PYG{+w}{    }\PYG{n}{BEGIN\PYGZus{}HIDE\PYGZus{}OUTPUT}\PYG{p}{(}\PYG{p}{)}\PYG{p}{;}
\PYG{+w}{    }\PYG{n}{command}\PYG{p}{(}\PYG{l+s}{\PYGZdq{}}\PYG{l+s}{reset\PYGZus{}timestep 10}\PYG{l+s}{\PYGZdq{}}\PYG{p}{)}\PYG{p}{;}
\PYG{+w}{    }\PYG{n}{END\PYGZus{}HIDE\PYGZus{}OUTPUT}\PYG{p}{(}\PYG{p}{)}\PYG{p}{;}
\PYG{+w}{    }\PYG{n}{ASSERT\PYGZus{}EQ}\PYG{p}{(}\PYG{n}{lmp}\PYG{o}{\PYGZhy{}}\PYG{o}{\PYGZgt{}}\PYG{n}{update}\PYG{o}{\PYGZhy{}}\PYG{o}{\PYGZgt{}}\PYG{n}{ntimestep}\PYG{p}{,}\PYG{+w}{ }\PYG{l+m+mi}{10}\PYG{p}{)}\PYG{p}{;}

\PYG{+w}{    }\PYG{n}{BEGIN\PYGZus{}HIDE\PYGZus{}OUTPUT}\PYG{p}{(}\PYG{p}{)}\PYG{p}{;}
\PYG{+w}{    }\PYG{n}{command}\PYG{p}{(}\PYG{l+s}{\PYGZdq{}}\PYG{l+s}{reset\PYGZus{}timestep 0}\PYG{l+s}{\PYGZdq{}}\PYG{p}{)}\PYG{p}{;}
\PYG{+w}{    }\PYG{n}{END\PYGZus{}HIDE\PYGZus{}OUTPUT}\PYG{p}{(}\PYG{p}{)}\PYG{p}{;}
\PYG{+w}{    }\PYG{n}{ASSERT\PYGZus{}EQ}\PYG{p}{(}\PYG{n}{lmp}\PYG{o}{\PYGZhy{}}\PYG{o}{\PYGZgt{}}\PYG{n}{update}\PYG{o}{\PYGZhy{}}\PYG{o}{\PYGZgt{}}\PYG{n}{ntimestep}\PYG{p}{,}\PYG{+w}{ }\PYG{l+m+mi}{0}\PYG{p}{)}\PYG{p}{;}

\PYG{+w}{    }\PYG{n}{TEST\PYGZus{}FAILURE}\PYG{p}{(}\PYG{l+s}{\PYGZdq{}}\PYG{l+s}{.*ERROR: Timestep must be \PYGZgt{}= 0.*}\PYG{l+s}{\PYGZdq{}}\PYG{p}{,}\PYG{+w}{ }\PYG{n}{command}\PYG{p}{(}\PYG{l+s}{\PYGZdq{}}\PYG{l+s}{reset\PYGZus{}timestep \PYGZhy{}10}\PYG{l+s}{\PYGZdq{}}\PYG{p}{)}\PYG{p}{;}\PYG{p}{)}\PYG{p}{;}
\PYG{+w}{    }\PYG{n}{TEST\PYGZus{}FAILURE}\PYG{p}{(}\PYG{l+s}{\PYGZdq{}}\PYG{l+s}{.*ERROR: Illegal reset\PYGZus{}timestep .*}\PYG{l+s}{\PYGZdq{}}\PYG{p}{,}\PYG{+w}{ }\PYG{n}{command}\PYG{p}{(}\PYG{l+s}{\PYGZdq{}}\PYG{l+s}{reset\PYGZus{}timestep}\PYG{l+s}{\PYGZdq{}}\PYG{p}{)}\PYG{p}{;}\PYG{p}{)}\PYG{p}{;}
\PYG{+w}{    }\PYG{n}{TEST\PYGZus{}FAILURE}\PYG{p}{(}\PYG{l+s}{\PYGZdq{}}\PYG{l+s}{.*ERROR: Illegal reset\PYGZus{}timestep .*}\PYG{l+s}{\PYGZdq{}}\PYG{p}{,}\PYG{+w}{ }\PYG{n}{command}\PYG{p}{(}\PYG{l+s}{\PYGZdq{}}\PYG{l+s}{reset\PYGZus{}timestep 10 10}\PYG{l+s}{\PYGZdq{}}\PYG{p}{)}\PYG{p}{;}\PYG{p}{)}\PYG{p}{;}
\PYG{+w}{    }\PYG{n}{TEST\PYGZus{}FAILURE}\PYG{p}{(}\PYG{l+s}{\PYGZdq{}}\PYG{l+s}{.*ERROR: Expected integer .*}\PYG{l+s}{\PYGZdq{}}\PYG{p}{,}\PYG{+w}{ }\PYG{n}{command}\PYG{p}{(}\PYG{l+s}{\PYGZdq{}}\PYG{l+s}{reset\PYGZus{}timestep xxx}\PYG{l+s}{\PYGZdq{}}\PYG{p}{)}\PYG{p}{;}\PYG{p}{)}\PYG{p}{;}
\PYG{p}{\PYGZcb{}}
\end{sphinxVerbatim}

\sphinxAtStartPar
Please note the use of the \sphinxcode{\sphinxupquote{BEGIN\_HIDE\_OUTPUT}} and \sphinxcode{\sphinxupquote{END\_HIDE\_OUTPUT}}
functions that will capture output from running LAMMPS.  This is normally
discarded but by setting the verbose flag (via setting the \sphinxcode{\sphinxupquote{TEST\_ARGS}}
environment variable, \sphinxcode{\sphinxupquote{TEST\_ARGS=\sphinxhyphen{}v}}) it can be printed and used to
understand why tests fail unexpectedly.

\sphinxAtStartPar
The specifics of so\sphinxhyphen{}called “death tests”, i.e. conditions where LAMMPS
should fail and throw an exception, are implemented in the
\sphinxcode{\sphinxupquote{TEST\_FAILURE()}} macro. These tests operate by capturing the screen
output when executing the failing command and then comparing that with a
provided regular expression string pattern.  Example:

\begin{sphinxVerbatim}[commandchars=\\\{\}]
\PYG{n}{TEST\PYGZus{}F}\PYG{p}{(}\PYG{n}{SimpleCommandsTest}\PYG{p}{,}\PYG{+w}{ }\PYG{n}{UnknownCommand}\PYG{p}{)}
\PYG{p}{\PYGZob{}}
\PYG{+w}{    }\PYG{n}{TEST\PYGZus{}FAILURE}\PYG{p}{(}\PYG{l+s}{\PYGZdq{}}\PYG{l+s}{.*ERROR: Unknown command.*}\PYG{l+s}{\PYGZdq{}}\PYG{p}{,}\PYG{+w}{ }\PYG{n}{lmp}\PYG{o}{\PYGZhy{}}\PYG{o}{\PYGZgt{}}\PYG{n}{input}\PYG{o}{\PYGZhy{}}\PYG{o}{\PYGZgt{}}\PYG{n}{one}\PYG{p}{(}\PYG{l+s}{\PYGZdq{}}\PYG{l+s}{XXX one two}\PYG{l+s}{\PYGZdq{}}\PYG{p}{)}\PYG{p}{;}\PYG{p}{)}\PYG{p}{;}
\PYG{p}{\PYGZcb{}}
\end{sphinxVerbatim}

\sphinxAtStartPar
The following test programs are currently available:


\begin{savenotes}\sphinxattablestart
\sphinxthistablewithglobalstyle
\raggedright
\begin{tabulary}{\linewidth}[t]{TTT}
\sphinxtoprule
\sphinxstyletheadfamily 
\sphinxAtStartPar
File name:
&\sphinxstyletheadfamily 
\sphinxAtStartPar
Test name:
&\sphinxstyletheadfamily 
\sphinxAtStartPar
Description:
\\
\sphinxmidrule
\sphinxtableatstartofbodyhook
\sphinxAtStartPar
\sphinxcode{\sphinxupquote{test\_simple\_commands.cpp}}
&
\sphinxAtStartPar
SimpleCommands
&
\sphinxAtStartPar
Tests for LAMMPS commands that do not require a box
\\
\sphinxhline
\sphinxAtStartPar
\sphinxcode{\sphinxupquote{test\_lattice\_region.cpp}}
&
\sphinxAtStartPar
LatticeRegion
&
\sphinxAtStartPar
Tests to validate the {\hyperref[\detokenize{lattice::doc}]{\sphinxcrossref{\DUrole{doc}{lattice}}}} and {\hyperref[\detokenize{region::doc}]{\sphinxcrossref{\DUrole{doc}{region}}}} commands
\\
\sphinxhline
\sphinxAtStartPar
\sphinxcode{\sphinxupquote{test\_groups.cpp}}
&
\sphinxAtStartPar
GroupTest
&
\sphinxAtStartPar
Tests to validate the {\hyperref[\detokenize{group::doc}]{\sphinxcrossref{\DUrole{doc}{group}}}} command
\\
\sphinxhline
\sphinxAtStartPar
\sphinxcode{\sphinxupquote{test\_variables.cpp}}
&
\sphinxAtStartPar
VariableTest
&
\sphinxAtStartPar
Tests to validate the {\hyperref[\detokenize{variable::doc}]{\sphinxcrossref{\DUrole{doc}{variable}}}} command
\\
\sphinxhline
\sphinxAtStartPar
\sphinxcode{\sphinxupquote{test\_kim\_commands.cpp}}
&
\sphinxAtStartPar
KimCommands
&
\sphinxAtStartPar
Tests for several commands from the {\hyperref[\detokenize{Packages_details:pkg-kim}]{\sphinxcrossref{\DUrole{std,std-ref}{KIM package}}}}
\\
\sphinxhline
\sphinxAtStartPar
\sphinxcode{\sphinxupquote{test\_reset\_atoms.cpp}}
&
\sphinxAtStartPar
ResetAtoms
&
\sphinxAtStartPar
Tests to validate the {\hyperref[\detokenize{reset_atoms::doc}]{\sphinxcrossref{\DUrole{doc}{reset\_atoms}}}} sub\sphinxhyphen{}commands
\\
\sphinxbottomrule
\end{tabulary}
\sphinxtableafterendhook\par
\sphinxattableend\end{savenotes}


\subsection{Tests for the C\sphinxhyphen{}style library interface}
\label{\detokenize{Developer_unittest:tests-for-the-c-style-library-interface}}
\sphinxAtStartPar
Tests for validating the LAMMPS C\sphinxhyphen{}style library interface are in the
\sphinxcode{\sphinxupquote{unittest/c\sphinxhyphen{}library}} folder.  They are implemented either to be used
for utility functions or for LAMMPS commands, but use the functions
implemented in the \sphinxcode{\sphinxupquote{src/library.cpp}} file as much as possible.  There
may be some overlap with other tests, but only in as much as is required
to test the C\sphinxhyphen{}style library API.  The tests are distributed over
multiple test programs which try to match the grouping of the
functions in the source code and {\hyperref[\detokenize{Library:lammps-c-api}]{\sphinxcrossref{\DUrole{std,std-ref}{in the manual}}}}.

\sphinxAtStartPar
This group of tests also includes tests invoking LAMMPS in parallel
through the library interface, provided that LAMMPS was compiled with
MPI support.  These include tests where LAMMPS is run in multi\sphinxhyphen{}partition
mode or only on a subset of the MPI world communicator.  The CMake
script code for adding this kind of test looks like this:

\begin{sphinxVerbatim}[commandchars=\\\{\}]
\PYG{n+nb}{if}\PYG{+w}{ }\PYG{p}{(}\PYG{l+s}{BUILD\PYGZus{}MPI}\PYG{p}{)}
\PYG{+w}{  }\PYG{n+nb}{add\PYGZus{}executable}\PYG{p}{(}\PYG{l+s}{test\PYGZus{}library\PYGZus{}mpi}\PYG{+w}{ }\PYG{l+s}{test\PYGZus{}library\PYGZus{}mpi.cpp}\PYG{p}{)}
\PYG{+w}{  }\PYG{n+nb}{target\PYGZus{}link\PYGZus{}libraries}\PYG{p}{(}\PYG{l+s}{test\PYGZus{}library\PYGZus{}mpi}\PYG{+w}{ }\PYG{l+s}{PRIVATE}\PYG{+w}{ }\PYG{l+s}{lammps}\PYG{+w}{ }\PYG{l+s}{GTest::GTest}\PYG{+w}{ }\PYG{l+s}{GTest::GMock}\PYG{p}{)}
\PYG{+w}{  }\PYG{n+nb}{target\PYGZus{}compile\PYGZus{}definitions}\PYG{p}{(}\PYG{l+s}{test\PYGZus{}library\PYGZus{}mpi}\PYG{+w}{ }\PYG{l+s}{PRIVATE}\PYG{+w}{ }\PYG{o}{\PYGZdl{}\PYGZob{}}\PYG{n+nv}{TEST\PYGZus{}CONFIG\PYGZus{}DEFS}\PYG{o}{\PYGZcb{}}\PYG{p}{)}
\PYG{+w}{  }\PYG{n+nb}{add\PYGZus{}mpi\PYGZus{}test}\PYG{p}{(}\PYG{l+s}{NAME}\PYG{+w}{ }\PYG{l+s}{LibraryMPI}\PYG{+w}{ }\PYG{l+s}{NUM\PYGZus{}PROCS}\PYG{+w}{ }\PYG{l+s}{4}\PYG{+w}{ }\PYG{l+s}{COMMAND}\PYG{+w}{ }\PYG{o}{\PYGZdl{}\PYGZlt{}}\PYG{n+nv}{TARGET\PYGZus{}FILE:test\PYGZus{}library\PYGZus{}mpi}\PYG{o}{\PYGZgt{}}\PYG{p}{)}
\PYG{n+nb}{endif}\PYG{p}{(}\PYG{p}{)}
\end{sphinxVerbatim}

\sphinxAtStartPar
Note the custom function \sphinxcode{\sphinxupquote{add\_mpi\_test()}} which adapts how \sphinxcode{\sphinxupquote{ctest}}
will execute the test so it is launched in parallel (with 4 MPI ranks).


\subsection{Tests for the Python module and package}
\label{\detokenize{Developer_unittest:tests-for-the-python-module-and-package}}
\sphinxAtStartPar
The \sphinxcode{\sphinxupquote{unittest/python}} folder contains primarily tests for classes and
functions in the LAMMPS python module but also for commands in the
PYTHON package.  These tests are only enabled if the necessary
prerequisites are detected or enabled during configuration and
compilation of LAMMPS (shared library build enabled, Python interpreter
found, Python development files found).

\sphinxAtStartPar
The Python tests are implemented using the \sphinxcode{\sphinxupquote{unittest}} standard Python
module and split into multiple files with similar categories as the
tests for the C\sphinxhyphen{}style library interface.


\subsection{Tests for the Fortran interface}
\label{\detokenize{Developer_unittest:tests-for-the-fortran-interface}}
\sphinxAtStartPar
Tests for using the Fortran module are in the \sphinxcode{\sphinxupquote{unittest/fortran}}
folder.  Since they are also using the GoogleTest library, they require
implementing test wrappers in C++ that will call fortran functions
which provide a C function interface through ISO\_C\_BINDINGS that will in
turn call the functions in the LAMMPS Fortran module.


\subsection{Tests for the C++\sphinxhyphen{}style library interface}
\label{\detokenize{Developer_unittest:id1}}
\sphinxAtStartPar
The tests in the \sphinxcode{\sphinxupquote{unittest/cplusplus}} folder are somewhat similar to
the tests for the C\sphinxhyphen{}style library interface, but do not need to test the
several convenience and utility functions that are only available through
the C\sphinxhyphen{}style interface.  Instead it can focus on the more generic features
that are used internally.  This part of the unit tests is currently still
mostly in the planning stage.


\subsection{Tests for reading and writing file formats}
\label{\detokenize{Developer_unittest:tests-for-reading-and-writing-file-formats}}
\sphinxAtStartPar
The \sphinxcode{\sphinxupquote{unittest/formats}} folder contains test programs for reading and
writing files like data files, restart files, potential files or dump files.
This covers simple things like the file i/o convenience functions in the
\sphinxcode{\sphinxupquote{utils::}} namespace to complex tests of atom styles where creating and
deleting atoms with different properties is tested in different ways
and through script commands or reading and writing of data or restart files.


\subsection{Tests for styles computing or modifying forces}
\label{\detokenize{Developer_unittest:tests-for-styles-computing-or-modifying-forces}}
\sphinxAtStartPar
These are tests common configurations for pair styles, bond styles,
angle styles, kspace styles and certain fix styles.  Those are tests
driven by some test executables build from sources in the
\sphinxcode{\sphinxupquote{unittest/force\sphinxhyphen{}styles}} folder and use LAMMPS input template and data
files as well as input files in YAML format from the
\sphinxcode{\sphinxupquote{unittest/force\sphinxhyphen{}styles/tests}} folder. The YAML file names have to
follow some naming conventions so they get associated with the test
programs and categorized and listed with canonical names in the list
of tests as displayed by \sphinxcode{\sphinxupquote{ctest \sphinxhyphen{}N}}.  If you add a new YAML file,
you need to re\sphinxhyphen{}run CMake to update the corresponding list of tests.

\sphinxAtStartPar
A minimal YAML file for a (molecular) pair style test will looks
something like the following (see \sphinxcode{\sphinxupquote{mol\sphinxhyphen{}pair\sphinxhyphen{}zero.yaml}}):

\begin{sphinxVerbatim}[commandchars=\\\{\}]
\PYG{n+nn}{\PYGZhy{}\PYGZhy{}\PYGZhy{}}
\PYG{n+nt}{lammps\PYGZus{}version}\PYG{p}{:}\PYG{+w}{ }\PYG{l+lScalar+lScalarPlain}{24}\PYG{l+lScalar+lScalarPlain}{ }\PYG{l+lScalar+lScalarPlain}{Aug}\PYG{l+lScalar+lScalarPlain}{ }\PYG{l+lScalar+lScalarPlain}{2020}
\PYG{n+nt}{date\PYGZus{}generated}\PYG{p}{:}\PYG{+w}{ }\PYG{l+lScalar+lScalarPlain}{Tue}\PYG{l+lScalar+lScalarPlain}{ }\PYG{l+lScalar+lScalarPlain}{Sep}\PYG{l+lScalar+lScalarPlain}{ }\PYG{l+lScalar+lScalarPlain}{15}\PYG{l+lScalar+lScalarPlain}{ }\PYG{l+lScalar+lScalarPlain}{09:44:21}\PYG{l+lScalar+lScalarPlain}{ }\PYG{l+lScalar+lScalarPlain}{202}
\PYG{n+nt}{epsilon}\PYG{p}{:}\PYG{+w}{ }\PYG{l+lScalar+lScalarPlain}{1e\PYGZhy{}14}
\PYG{n+nt}{prerequisites}\PYG{p}{:}\PYG{+w}{ }\PYG{k+kt}{!}\PYG{+w}{ }\PYG{p+pIndicator}{|}
\PYG{+w}{  }\PYG{n+no}{atom full}
\PYG{+w}{  }\PYG{n+no}{pair zero}
\PYG{n+nt}{pre\PYGZus{}commands}\PYG{p}{:}\PYG{+w}{ }\PYG{k+kt}{!}\PYG{+w}{ }\PYG{l+s}{\PYGZdq{}}\PYG{l+s}{\PYGZdq{}}
\PYG{n+nt}{post\PYGZus{}commands}\PYG{p}{:}\PYG{+w}{ }\PYG{k+kt}{!}\PYG{+w}{ }\PYG{l+s}{\PYGZdq{}}\PYG{l+s}{\PYGZdq{}}
\PYG{n+nt}{input\PYGZus{}file}\PYG{p}{:}\PYG{+w}{ }\PYG{l+lScalar+lScalarPlain}{in.fourmol}
\PYG{n+nt}{pair\PYGZus{}style}\PYG{p}{:}\PYG{+w}{ }\PYG{l+lScalar+lScalarPlain}{zero}\PYG{l+lScalar+lScalarPlain}{ }\PYG{l+lScalar+lScalarPlain}{8.0}
\PYG{n+nt}{pair\PYGZus{}coeff}\PYG{p}{:}\PYG{+w}{ }\PYG{k+kt}{!}\PYG{+w}{ }\PYG{p+pIndicator}{|}
\PYG{+w}{  }\PYG{n+no}{* *}
\PYG{n+nt}{extract}\PYG{p}{:}\PYG{+w}{ }\PYG{k+kt}{!}\PYG{+w}{ }\PYG{l+s}{\PYGZdq{}}\PYG{l+s}{\PYGZdq{}}
\PYG{n+nt}{natoms}\PYG{p}{:}\PYG{+w}{ }\PYG{l+lScalar+lScalarPlain}{29}
\PYG{n+nt}{init\PYGZus{}vdwl}\PYG{p}{:}\PYG{+w}{ }\PYG{l+lScalar+lScalarPlain}{0}
\PYG{n+nt}{init\PYGZus{}coul}\PYG{p}{:}\PYG{+w}{ }\PYG{l+lScalar+lScalarPlain}{0}

\PYG{p+pIndicator}{[}\PYG{n+nv}{...}\PYG{p+pIndicator}{]}
\end{sphinxVerbatim}

\sphinxAtStartPar
The following table describes the available keys and their purpose for
testing pair styles:


\begin{savenotes}\sphinxattablestart
\sphinxthistablewithglobalstyle
\centering
\begin{tabulary}{\linewidth}[t]{TT}
\sphinxtoprule
\sphinxstyletheadfamily 
\sphinxAtStartPar
Key:
&\sphinxstyletheadfamily 
\sphinxAtStartPar
Description:
\\
\sphinxmidrule
\sphinxtableatstartofbodyhook
\sphinxAtStartPar
lammps\_version
&
\sphinxAtStartPar
LAMMPS version used to last update the reference data
\\
\sphinxhline
\sphinxAtStartPar
date\_generated
&
\sphinxAtStartPar
date when the file was last updated
\\
\sphinxhline
\sphinxAtStartPar
epsilon
&
\sphinxAtStartPar
base value for the relative precision required for tests to pass
\\
\sphinxhline
\sphinxAtStartPar
prerequisites
&
\sphinxAtStartPar
list of style kind / style name pairs required to run the test
\\
\sphinxhline
\sphinxAtStartPar
pre\_commands
&
\sphinxAtStartPar
LAMMPS commands to be executed before the input template file is read
\\
\sphinxhline
\sphinxAtStartPar
post\_commands
&
\sphinxAtStartPar
LAMMPS commands to be executed right before the actual tests
\\
\sphinxhline
\sphinxAtStartPar
input\_file
&
\sphinxAtStartPar
LAMMPS input file template based on pair style zero
\\
\sphinxhline
\sphinxAtStartPar
pair\_style
&
\sphinxAtStartPar
arguments to the pair\_style command to be tested
\\
\sphinxhline
\sphinxAtStartPar
pair\_coeff
&
\sphinxAtStartPar
list of pair\_coeff arguments to set parameters for the input template
\\
\sphinxhline
\sphinxAtStartPar
extract
&
\sphinxAtStartPar
list of keywords supported by \sphinxcode{\sphinxupquote{Pair::extract()}} and their dimension
\\
\sphinxhline
\sphinxAtStartPar
natoms
&
\sphinxAtStartPar
number of atoms in the input file template
\\
\sphinxhline
\sphinxAtStartPar
init\_vdwl
&
\sphinxAtStartPar
non\sphinxhyphen{}Coulomb pair energy after “run 0”
\\
\sphinxhline
\sphinxAtStartPar
init\_coul
&
\sphinxAtStartPar
Coulomb pair energy after “run 0”
\\
\sphinxhline
\sphinxAtStartPar
init\_stress
&
\sphinxAtStartPar
stress tensor after “run 0”
\\
\sphinxhline
\sphinxAtStartPar
init\_forces
&
\sphinxAtStartPar
forces on atoms after “run 0”
\\
\sphinxhline
\sphinxAtStartPar
run\_vdwl
&
\sphinxAtStartPar
non\sphinxhyphen{}Coulomb pair energy after “run 4”
\\
\sphinxhline
\sphinxAtStartPar
run\_coul
&
\sphinxAtStartPar
Coulomb pair energy after “run 4”
\\
\sphinxhline
\sphinxAtStartPar
run\_stress
&
\sphinxAtStartPar
stress tensor after “run 4”
\\
\sphinxhline
\sphinxAtStartPar
run\_forces
&
\sphinxAtStartPar
forces on atoms after “run 4”
\\
\sphinxbottomrule
\end{tabulary}
\sphinxtableafterendhook\par
\sphinxattableend\end{savenotes}

\sphinxAtStartPar
The test program will read all this data from the YAML file and then
create a LAMMPS instance, apply the settings/commands from the YAML file
as needed and then issue a “run 0” command, write out a restart file, a
data file and a coeff file. The actual test will then compare computed
energies, stresses, and forces with the reference data, issue a “run 4”
command and compare to the second set of reference data.  This will be
run with both the newton\_pair setting enabled and disabled and is
expected to generate the same results (allowing for some numerical
noise). Then it will restart from the previously generated restart and
compare with the reference and also start from the data file.  A final
check will use multi\sphinxhyphen{}cutoff r\sphinxhyphen{}RESPA (if supported by the pair style) at
a 1:1 split and compare to the Verlet results.  These sets of tests are
run with multiple test fixtures for accelerated styles (OPT, OPENMP,
INTEL, KOKKOS (OpenMP only)) and for the latter three with 4 OpenMP
threads enabled.  For these tests the relative error (epsilon) is lowered
by a common factor due to the additional numerical noise, but the tests
are still comparing to the same reference data.

\sphinxAtStartPar
Additional tests will check whether all listed extract keywords are
supported and have the correct dimensionality and the final set of tests
will set up a few pairs of atoms explicitly and in such a fashion that
the forces on the atoms computed from \sphinxcode{\sphinxupquote{Pair::compute()}} will match
individually with the results from \sphinxcode{\sphinxupquote{Pair::single()}}, if the pair style
does support that functionality.

\sphinxAtStartPar
With this scheme a large fraction of the code of any tested pair style
will be executed and consistent results are required for different
settings and between different accelerated pair style variants and the
base class, as well as for computing individual pairs through the
\sphinxcode{\sphinxupquote{Pair::single()}} method where supported.

\sphinxAtStartPar
The \sphinxcode{\sphinxupquote{test\_pair\_style}} tester is used with 4 categories of test inputs:
\begin{itemize}
\item {} 
\sphinxAtStartPar
pair styles compatible with molecular systems using bonded
interactions and exclusions.  For pair styles requiring a KSpace style
the KSpace computations are disabled.  The YAML files match the
pattern “mol\sphinxhyphen{}pair\sphinxhyphen{}*.yaml” and the tests are correspondingly labeled
with “MolPairStyle:*”

\item {} 
\sphinxAtStartPar
pair styles not compatible with the previous input template.
The YAML files match the pattern “atomic\sphinxhyphen{}pair\sphinxhyphen{}*.yaml” and the tests are
correspondingly labeled with “AtomicPairStyle:*”

\item {} 
\sphinxAtStartPar
manybody pair styles.
The YAML files match the pattern “atomic\sphinxhyphen{}pair\sphinxhyphen{}*.yaml” and the tests are
correspondingly labeled with “AtomicPairStyle:*”

\item {} 
\sphinxAtStartPar
kspace styles.
The YAML files match the pattern “kspace\sphinxhyphen{}*.yaml” and the tests are
correspondingly labeled with “KSpaceStyle:*”. In these cases a compatible
pair style is defined, but the computation of the pair style contributions
is disabled.

\end{itemize}

\sphinxAtStartPar
The \sphinxcode{\sphinxupquote{test\_bond\_style}}, \sphinxcode{\sphinxupquote{test\_angle\_style}}, \sphinxcode{\sphinxupquote{test\_dihedral\_style}}, and
\sphinxcode{\sphinxupquote{test\_improper\_style}} tester programs are set up in a similar fashion and
share support functions with the pair style tester.  The final group of
tests in this section is for fix styles that add/manipulate forces and
velocities, e.g. for time integration, thermostats and more.

\sphinxAtStartPar
Adding a new test is easiest done by copying and modifying an existing YAML
file for a style that is similar to one to be tested.  The file name should
follow the naming conventions described above and after copying the file,
the first step is to replace the style names where needed.  The coefficient
values do not have to be meaningful, just in a reasonable range for the
given system.  It does not matter if some forces are large, as long as
they do not diverge.

\sphinxAtStartPar
The template input files define a large number of index variables at the top
that can be modified inside the YAML file to control the behavior.  For example,
if a pair style requires a “newton on” setting, the following can be used in
as the “pre\_commands” section:

\begin{sphinxVerbatim}[commandchars=\\\{\}]
\PYG{n+nt}{pre\PYGZus{}commands}\PYG{p}{:}\PYG{+w}{ }\PYG{k+kt}{!}\PYG{+w}{ }\PYG{p+pIndicator}{|}
\PYG{+w}{  }\PYG{n+no}{variable newton\PYGZus{}pair delete}
\PYG{+w}{  }\PYG{n+no}{variable newton\PYGZus{}pair index on}
\end{sphinxVerbatim}

\sphinxAtStartPar
And for a pair style requiring a kspace solver the following would be used as
the “post\_commands” section:

\begin{sphinxVerbatim}[commandchars=\\\{\}]
\PYG{n+nt}{post\PYGZus{}commands}\PYG{p}{:}\PYG{+w}{ }\PYG{k+kt}{!}\PYG{+w}{ }\PYG{p+pIndicator}{|}
\PYG{+w}{  }\PYG{n+no}{pair\PYGZus{}modify table 0}
\PYG{+w}{  }\PYG{n+no}{kspace\PYGZus{}style pppm/tip4p 1.0e\PYGZhy{}6}
\PYG{+w}{  }\PYG{n+no}{kspace\PYGZus{}modify gewald 0.3}
\PYG{+w}{  }\PYG{n+no}{kspace\PYGZus{}modify compute no}
\end{sphinxVerbatim}

\sphinxAtStartPar
Note that this disables computing the kspace contribution, but still will run
the setup.  The “gewald” parameter should be set explicitly to speed up the run.
For styles with long\sphinxhyphen{}range electrostatics, typically two tests are added one using
the (slower) analytic approximation of the erfc() function and the other using
the tabulated coulomb, to test both code paths.  The reference results in the YAML
files then should be compared manually, if they agree well enough within the limits
of those two approximations.

\sphinxAtStartPar
The \sphinxcode{\sphinxupquote{test\_pair\_style}} and equivalent programs have special command line options
to update the YAML files. Running a command like

\begin{sphinxVerbatim}[commandchars=\\\{\}]
test\PYGZus{}pair\PYGZus{}style\PYG{+w}{ }mol\PYGZhy{}pair\PYGZhy{}lennard\PYGZus{}mdf.yaml\PYG{+w}{ }\PYGZhy{}g\PYG{+w}{ }new.yaml
\end{sphinxVerbatim}

\sphinxAtStartPar
will read the settings from the \sphinxcode{\sphinxupquote{mol\sphinxhyphen{}pair\sphinxhyphen{}lennard\_mdf.yaml}} file and then compute
the reference data and write a new file with to \sphinxcode{\sphinxupquote{new.yaml}}.  If this step fails,
there are likely some (LAMMPS or YAML) syntax issues in the YAML file that need to
be resolved and then one can compare the two files to see if the output is as expected.

\sphinxAtStartPar
It is also possible to do an update in place with:

\begin{sphinxVerbatim}[commandchars=\\\{\}]
test\PYGZus{}pair\PYGZus{}style\PYG{+w}{ }mol\PYGZhy{}pair\PYGZhy{}lennard\PYGZus{}mdf.yaml\PYG{+w}{ }\PYGZhy{}u
\end{sphinxVerbatim}

\sphinxAtStartPar
And one can finally run the full set of tests with:

\begin{sphinxVerbatim}[commandchars=\\\{\}]
test\PYGZus{}pair\PYGZus{}style\PYG{+w}{ }mol\PYGZhy{}pair\PYGZhy{}lennard\PYGZus{}mdf.yaml
\end{sphinxVerbatim}

\sphinxAtStartPar
This will just print a summary of the groups of tests.  When using the “\sphinxhyphen{}v” flag
the test will also keep any LAMMPS output and when using the “\sphinxhyphen{}s” flag, there
will be some statistics reported on the relative errors for the individual checks
which can help to figure out what would be a good choice of the epsilon parameter.
It should be as small as possible to catch any unintended side effects from changes
elsewhere, but large enough to accommodate the numerical noise due to the implementation
of the potentials and differences in compilers.

\begin{sphinxadmonition}{note}{Note:}
\sphinxAtStartPar
These kinds of tests can be very sensitive to compiler optimization and
thus the expectation is that they pass with compiler optimization turned
off. When compiler optimization is enabled, there may be some failures, but
one has to carefully check whether those are acceptable due to the enhanced
numerical noise from reordering floating\sphinxhyphen{}point math operations or due to
the compiler mis\sphinxhyphen{}compiling the code. That is not always obvious.
\end{sphinxadmonition}


\subsection{Tests for programs in the tools folder}
\label{\detokenize{Developer_unittest:tests-for-programs-in-the-tools-folder}}
\sphinxAtStartPar
The \sphinxcode{\sphinxupquote{unittest/tools}} folder contains tests for programs in the
\sphinxcode{\sphinxupquote{tools}} folder.  This currently only contains tests for the LAMMPS
shell, which are implemented as a python scripts using the \sphinxcode{\sphinxupquote{unittest}}
Python module and launching the tool commands through the \sphinxcode{\sphinxupquote{subprocess}}
Python module.


\subsection{Troubleshooting failed unit tests}
\label{\detokenize{Developer_unittest:troubleshooting-failed-unit-tests}}
\sphinxAtStartPar
The are by default no unit tests for newly added features (e.g. pair, fix,
or compute styles) unless your pull request also includes tests for the
added features.  If you are modifying some features, you may see failures
for existing tests, if your modifications have some unexpected side effects
or your changes render the existing test invalid.  If you are adding an
accelerated version of an existing style, then only tests for INTEL,
KOKKOS (with OpenMP only), OPENMP, and OPT will be run automatically.
Tests for the GPU package are time consuming and thus are only run
\sphinxstyleemphasis{after} a merge, or when a special label, \sphinxcode{\sphinxupquote{gpu\_unit\_tests}} is added
to the pull request.  After the test has started, it is often best to
remove the label since every PR activity will re\sphinxhyphen{}trigger the test (that
is a limitation of triggering a test with a label).  Support for unit
tests when using KOKKOS with GPU acceleration is currently not supported.

\sphinxAtStartPar
When you see a failed build on GitHub, click on \sphinxcode{\sphinxupquote{Details}} to be taken
to the corresponding LAMMPS Jenkins CI web page.  Click on the “Exit”
symbol near the \sphinxcode{\sphinxupquote{Logout}} button on the top right of that page to go to
the “classic view”.  In the classic view, there is a list of the
individual runs that make up this test run (they are shown but cannot be
inspected in the default view).  You can click on any of those.
Clicking on \sphinxcode{\sphinxupquote{Test Result}} will display the list of failed tests. Click
on the “Status” column to sort the tests based on their Failed or Passed
status.  Then click on the failed test to expand its output.

\sphinxAtStartPar
For example, the following output snippet shows the failed unit test

\begin{sphinxVerbatim}[commandchars=\\\{\}]
\PYG{g+go}{[ RUN      ] PairStyle.gpu}
\PYG{g+go}{/home/builder/workspace/dev/pull\PYGZus{}requests/ubuntu\PYGZus{}gpu/unit\PYGZus{}tests/cmake\PYGZus{}gpu\PYGZus{}opencl\PYGZus{}mixed\PYGZus{}smallbig\PYGZus{}clang\PYGZus{}static/unittest/force\PYGZhy{}styles/test\PYGZus{}main.cpp:63: Failure}
\PYG{g+go}{Expected: (err) \PYGZlt{}= (epsilon)}
\PYG{g+go}{Actual: 0.00018957912910606503 vs 0.0001}
\PYG{g+go}{Google Test trace:}
\PYG{g+go}{/home/builder/workspace/dev/pull\PYGZus{}requests/ubuntu\PYGZus{}gpu/unit\PYGZus{}tests/cmake\PYGZus{}gpu\PYGZus{}opencl\PYGZus{}mixed\PYGZus{}smallbig\PYGZus{}clang\PYGZus{}static/unittest/force\PYGZhy{}styles/test\PYGZus{}main.cpp:56: EXPECT\PYGZus{}FORCES: init\PYGZus{}forces (newton off)}
\PYG{g+go}{/home/builder/workspace/dev/pull\PYGZus{}requests/ubuntu\PYGZus{}gpu/unit\PYGZus{}tests/cmake\PYGZus{}gpu\PYGZus{}opencl\PYGZus{}mixed\PYGZus{}smallbig\PYGZus{}clang\PYGZus{}static/unittest/force\PYGZhy{}styles/test\PYGZus{}main.cpp:64: Failure}
\PYG{g+go}{Expected: (err) \PYGZlt{}= (epsilon)}
\PYG{g+go}{Actual: 0.00022892713393549854 vs 0.0001}
\end{sphinxVerbatim}

\sphinxAtStartPar
The failed assertions provide line numbers in the test source
(e.g. \sphinxcode{\sphinxupquote{test\_main.cpp:56}}), from which one can understand what
specific assertion failed.

\sphinxAtStartPar
Note that the force style engine runs one of a small number of systems
in a rather off\sphinxhyphen{}equilibrium configuration with a few atoms for a few
steps, writes data and restart files, uses {\hyperref[\detokenize{clear::doc}]{\sphinxcrossref{\DUrole{doc}{the clear command}}}} to reset LAMMPS, and then runs from those files with different
settings (e.g. newton on/off) and integrators (e.g. verlet vs. respa).
Beyond potential issues/bugs in the source code, the mismatch between
the expected and actual values could be that force arrays are not
properly cleared between multiple run commands or that class members are
not correctly initialized or written to or read from a data or restart
file.

\sphinxAtStartPar
While the epsilon (relative precision) for a single, \sphinxhref{https://en.wikipedia.org/wiki/IEEE\_754}{IEEE 754 compliant}, double precision floating
point operation is at about 2.2e\sphinxhyphen{}16, the achievable precision for the
tests is lower due to most numbers being sums over intermediate results
and the non\sphinxhyphen{}associativity of floating point math leading to larger
errors.  As a rule of thumb, the test epsilon can often be in the range
5.0e\sphinxhyphen{}14 to 1.0e\sphinxhyphen{}13.  But for “noisy” force kernels, e.g. those a larger
amount of arithmetic operations involving \sphinxtitleref{exp()}, \sphinxtitleref{log()} or \sphinxtitleref{sin()}
functions, and also due to the effect of compiler optimization or differences
between compilers or platforms, epsilon may need to be further relaxed,
sometimes epsilon can be relaxed to 1.0e\sphinxhyphen{}12. If interpolation or lookup
tables are used, epsilon may need to be set to 1.0e\sphinxhyphen{}10 or even higher.
For tests of accelerated styles, the per\sphinxhyphen{}test epsilon is multiplied
by empirical factors that take into account the differences in the order
of floating point operations or that some or most intermediate operations
may be done using approximations or with single precision floating point
math.

\sphinxAtStartPar
To rerun the failed unit test individually, change to the \sphinxcode{\sphinxupquote{build}} directory
and run the test with verbose output. For example,

\begin{sphinxVerbatim}[commandchars=\\\{\}]
env\PYG{+w}{ }\PYG{n+nv}{TEST\PYGZus{}ARGS}\PYG{o}{=}\PYGZhy{}v\PYG{+w}{ }ctest\PYG{+w}{ }\PYGZhy{}R\PYG{+w}{ }\PYGZca{}MolPairStyle:lj\PYGZus{}cut\PYGZus{}coul\PYGZus{}long\PYG{+w}{ }\PYGZhy{}V
\end{sphinxVerbatim}

\sphinxAtStartPar
\sphinxcode{\sphinxupquote{ctest}} with the \sphinxcode{\sphinxupquote{\sphinxhyphen{}V}} flag also shows the exact command line
of the test. One can then use \sphinxcode{\sphinxupquote{gdb \sphinxhyphen{}\sphinxhyphen{}args}} to further debug and
catch exceptions with the test command, for example,

\begin{sphinxVerbatim}[commandchars=\\\{\}]
gdb\PYG{+w}{ }\PYGZhy{}\PYGZhy{}args\PYG{+w}{ }/path/to/lammps/build/test\PYGZus{}pair\PYGZus{}style\PYG{+w}{ }/path/to/lammps/unittest/force\PYGZhy{}styles/tests/mol\PYGZhy{}pair\PYGZhy{}lj\PYGZus{}cut\PYGZus{}coul\PYGZus{}long.yaml
\end{sphinxVerbatim}

\sphinxAtStartPar
It is recommended to configure the build with \sphinxcode{\sphinxupquote{\sphinxhyphen{}D
BUILD\_SHARED\_LIBS=on}} and use a custom linker to shorten the build time
during recompilation.  Installing \sphinxtitleref{ccache} in your development
environment helps speed up recompilation by caching previous
compilations and detecting when the same compilation is being done
again.  Please see {\hyperref[\detokenize{Build_development::doc}]{\sphinxcrossref{\DUrole{doc}{Development build options}}}} for further details.

\sphinxstepscope


\section{C++ base classes}
\label{\detokenize{Classes:c-base-classes}}\label{\detokenize{Classes::doc}}
\sphinxAtStartPar
LAMMPS is designed to be used as a C++ class library where one can set
up and drive a simulation through creating a class instance and then
calling some abstract operations or commands on that class or its member
class instances.  These are interfaced to the {\hyperref[\detokenize{Library::doc}]{\sphinxcrossref{\DUrole{doc}{C library API}}}}, which providing an additional level of abstraction
simplification for common operations. The C API is also the basis for
calling LAMMPS from Python or Fortran.

\sphinxAtStartPar
When used from a C++ program, most of the symbols and functions in
LAMMPS are wrapped into the \sphinxcode{\sphinxupquote{LAMMPS\_NS}} namespace so they will not
collide with your own classes or other libraries. This, however, does
not extend to the additional libraries bundled with LAMMPS in the lib
folder and some of the low\sphinxhyphen{}level code of some packages.

\sphinxAtStartPar
Behind the scenes this is implemented through inheritance and
polymorphism where base classes define the abstract interface and
derived classes provide the specialized implementation for specific
models or optimizations or ports to accelerator platforms.  This
document will provide an outline of the fundamental class hierarchy and
some selected examples for derived classes of specific models.

\begin{sphinxadmonition}{note}{Note:}
\sphinxAtStartPar
Please see the {\hyperref[\detokenize{Library:thread-safety}]{\sphinxcrossref{\DUrole{std,std-ref}{note about thread\sphinxhyphen{}safety}}}}
in the library Howto doc page.
\end{sphinxadmonition}


\bigskip\hrule\bigskip


\sphinxstepscope


\subsection{LAMMPS Class}
\label{\detokenize{Classes_lammps:lammps-class}}\label{\detokenize{Classes_lammps::doc}}
\sphinxAtStartPar
The LAMMPS class is encapsulating an MD simulation state and thus it is
the class that needs to be created when starting a new simulation system
state.  The LAMMPS executable essentially creates one instance of this
class and passes the command line flags and tells it to process the
provided input (a file or \sphinxcode{\sphinxupquote{stdin}}).  It shuts the class down when
control is returned to it and then exits.  When using LAMMPS as a
library from another code it is required to create an instance of this
class, either directly from C++ with \sphinxcode{\sphinxupquote{new LAMMPS()}} or through one
of the library interface functions like {\hyperref[\detokenize{Library_create:_CPPv411lammps_openiPPc8MPI_CommPPv}]{\sphinxcrossref{\sphinxcode{\sphinxupquote{lammps\_open()}}}}} of the
C\sphinxhyphen{}library interface, or the {\hyperref[\detokenize{Python_module:lammps.lammps}]{\sphinxcrossref{\sphinxcode{\sphinxupquote{lammps.lammps}}}}} class constructor
of the Python module, or the \sphinxcode{\sphinxupquote{lammps()}} constructor of the Fortran
module.

\sphinxAtStartPar
In order to avoid clashes of function names, all of the core code in
LAMMPS is placed into the \sphinxcode{\sphinxupquote{LAMMPS\_NS}} namespace.  Functions or variables
outside of that namespace must be “static”, i.e.  visible only to the
scope of the file/object they are defined in.  Code in packages or the
libraries in the \sphinxcode{\sphinxupquote{lib}} folder may not adhere to this as some of them
are adapted from legacy code or consist of external libraries with their
own requirements and policies.


\bigskip\hrule\bigskip

\index{LAMMPS (C++ class)@\spxentry{LAMMPS}\spxextra{C++ class}}

\begin{fulllineitems}
\phantomsection\label{\detokenize{Classes_lammps:_CPPv4N9LAMMPS_NS6LAMMPSE}}\phantomsection\label{\detokenize{Classes_lammps:_CPPv3N9LAMMPS_NS6LAMMPSE}}\phantomsection\label{\detokenize{Classes_lammps:_CPPv2N9LAMMPS_NS6LAMMPSE}}\phantomsection\label{\detokenize{Classes_lammps:LAMMPS_NS::LAMMPS}}
\pysigstartsignatures
\pysigstartmultiline
\pysigline{\phantomsection\label{\detokenize{Classes_lammps:classLAMMPS__NS_1_1LAMMPS}}\DUrole{k}{class}\DUrole{w}{ }\sphinxbfcode{\sphinxupquote{\DUrole{n}{LAMMPS}}}}
\pysigstopmultiline
\pysigstopsignatures
\sphinxAtStartPar
{\hyperref[\detokenize{Classes_lammps:classLAMMPS__NS_1_1LAMMPS}]{\sphinxcrossref{\DUrole{std,std-ref}{LAMMPS}}}} simulation instance. 

\sphinxAtStartPar
The {\hyperref[\detokenize{Classes_lammps:classLAMMPS__NS_1_1LAMMPS}]{\sphinxcrossref{\DUrole{std,std-ref}{LAMMPS}}}} class contains pointers of all constituent class instances and global variables that are used by a {\hyperref[\detokenize{Classes_lammps:classLAMMPS__NS_1_1LAMMPS}]{\sphinxcrossref{\DUrole{std,std-ref}{LAMMPS}}}} simulation. Its contents represent the entire state of the simulation.

\sphinxAtStartPar
The {\hyperref[\detokenize{Classes_lammps:classLAMMPS__NS_1_1LAMMPS}]{\sphinxcrossref{\DUrole{std,std-ref}{LAMMPS}}}} class manages the components of an MD simulation by creating, deleting, and initializing instances of the classes it is composed of, processing command line flags, and providing access to some global properties. The specifics of setting up and running a simulation are handled by the individual component class instances. 

\begin{sphinxuseclass}{breathe-sectiondef}\subsubsection*{Public Functions}
\index{LAMMPS::non\_pair\_suffix (C++ function)@\spxentry{LAMMPS::non\_pair\_suffix}\spxextra{C++ function}}

\begin{fulllineitems}
\phantomsection\label{\detokenize{Classes_lammps:_CPPv4NK9LAMMPS_NS6LAMMPS15non_pair_suffixEv}}\phantomsection\label{\detokenize{Classes_lammps:_CPPv3NK9LAMMPS_NS6LAMMPS15non_pair_suffixEv}}\phantomsection\label{\detokenize{Classes_lammps:_CPPv2NK9LAMMPS_NS6LAMMPS15non_pair_suffixEv}}\phantomsection\label{\detokenize{Classes_lammps:LAMMPS_NS::LAMMPS::non_pair_suffixC}}
\pysigstartsignatures
\pysigstartmultiline
\pysiglinewithargsret{\phantomsection\label{\detokenize{Classes_lammps:classLAMMPS__NS_1_1LAMMPS_1aeadaf4767470104f970858ed7d7227a3}}\DUrole{k}{const}\DUrole{w}{ }\DUrole{kt}{char}\DUrole{w}{ }\DUrole{p}{*}\sphinxbfcode{\sphinxupquote{\DUrole{n}{non\_pair\_suffix}}}}{}{\DUrole{w}{ }\DUrole{k}{const}}
\pysigstopmultiline
\pysigstopsignatures
\sphinxAtStartPar
Return suffix for non\sphinxhyphen{}pair styles depending on pair\_only\_flag. 
\begin{quote}\begin{description}
\sphinxlineitem{Returns}
\sphinxAtStartPar
suffix or null pointer 

\end{description}\end{quote}

\end{fulllineitems}

\index{LAMMPS::match\_style (C++ function)@\spxentry{LAMMPS::match\_style}\spxextra{C++ function}}

\begin{fulllineitems}
\phantomsection\label{\detokenize{Classes_lammps:_CPPv4N9LAMMPS_NS6LAMMPS11match_styleEPKcPKc}}\phantomsection\label{\detokenize{Classes_lammps:_CPPv3N9LAMMPS_NS6LAMMPS11match_styleEPKcPKc}}\phantomsection\label{\detokenize{Classes_lammps:_CPPv2N9LAMMPS_NS6LAMMPS11match_styleEPKcPKc}}\phantomsection\label{\detokenize{Classes_lammps:LAMMPS_NS::LAMMPS::match_style__cCP.cCP}}
\pysigstartsignatures
\pysigstartmultiline
\pysiglinewithargsret{\phantomsection\label{\detokenize{Classes_lammps:classLAMMPS__NS_1_1LAMMPS_1a494c5512fbe010f7eddf168ee45878d9}}\DUrole{k}{const}\DUrole{w}{ }\DUrole{kt}{char}\DUrole{w}{ }\DUrole{p}{*}\sphinxbfcode{\sphinxupquote{\DUrole{n}{match\_style}}}}{\DUrole{k}{const}\DUrole{w}{ }\DUrole{kt}{char}\DUrole{w}{ }\DUrole{p}{*}\DUrole{n,sig-param}{style}\sphinxparamcomma \DUrole{k}{const}\DUrole{w}{ }\DUrole{kt}{char}\DUrole{w}{ }\DUrole{p}{*}\DUrole{n,sig-param}{name}}{}
\pysigstopmultiline
\pysigstopsignatures
\sphinxAtStartPar
Return name of package that a specific style belongs to. 

\sphinxAtStartPar
This function checks the given name against all list of styles for all types of styles and if the name and the style match, it returns which package this style belongs to.
\begin{quote}\begin{description}
\sphinxlineitem{Parameters}\begin{itemize}
\item {} 
\sphinxAtStartPar
\sphinxstyleliteralstrong{\sphinxupquote{style}} \textendash{} Type of style (e.g. atom, pair, fix, etc.) 

\item {} 
\sphinxAtStartPar
\sphinxstyleliteralstrong{\sphinxupquote{name}} \textendash{} Name of style 

\end{itemize}

\sphinxlineitem{Returns}
\sphinxAtStartPar
Name of the package this style is part of 

\end{description}\end{quote}

\end{fulllineitems}

\index{LAMMPS::LAMMPS (C++ function)@\spxentry{LAMMPS::LAMMPS}\spxextra{C++ function}}

\begin{fulllineitems}
\phantomsection\label{\detokenize{Classes_lammps:_CPPv4N9LAMMPS_NS6LAMMPS6LAMMPSER4argv8MPI_Comm}}\phantomsection\label{\detokenize{Classes_lammps:_CPPv3N9LAMMPS_NS6LAMMPS6LAMMPSER4argv8MPI_Comm}}\phantomsection\label{\detokenize{Classes_lammps:_CPPv2N9LAMMPS_NS6LAMMPS6LAMMPSER4argv8MPI_Comm}}\phantomsection\label{\detokenize{Classes_lammps:LAMMPS_NS::LAMMPS::LAMMPS__argvR.MPI_Comm}}
\pysigstartsignatures
\pysigstartmultiline
\pysiglinewithargsret{\phantomsection\label{\detokenize{Classes_lammps:classLAMMPS__NS_1_1LAMMPS_1a6b5422fe9b6d4523afe1bd28ccc4a34f}}\sphinxbfcode{\sphinxupquote{\DUrole{n}{LAMMPS}}}}{\DUrole{n}{argv}\DUrole{w}{ }\DUrole{p}{\&}\DUrole{n,sig-param}{args}\sphinxparamcomma \DUrole{n}{MPI\_Comm}}{}
\pysigstopmultiline
\pysigstopsignatures
\sphinxAtStartPar
Create a {\hyperref[\detokenize{Classes_lammps:classLAMMPS__NS_1_1LAMMPS}]{\sphinxcrossref{\DUrole{std,std-ref}{LAMMPS}}}} simulation instance
\begin{quote}\begin{description}
\sphinxlineitem{Parameters}\begin{itemize}
\item {} 
\sphinxAtStartPar
\sphinxstyleliteralstrong{\sphinxupquote{args}} \textendash{} list of arguments 

\item {} 
\sphinxAtStartPar
\sphinxstyleliteralstrong{\sphinxupquote{communicator}} \textendash{} MPI communicator used by this {\hyperref[\detokenize{Classes_lammps:classLAMMPS__NS_1_1LAMMPS}]{\sphinxcrossref{\DUrole{std,std-ref}{LAMMPS}}}} instance 

\end{itemize}

\end{description}\end{quote}

\end{fulllineitems}

\index{LAMMPS::LAMMPS (C++ function)@\spxentry{LAMMPS::LAMMPS}\spxextra{C++ function}}

\begin{fulllineitems}
\phantomsection\label{\detokenize{Classes_lammps:_CPPv4N9LAMMPS_NS6LAMMPS6LAMMPSEiPPc8MPI_Comm}}\phantomsection\label{\detokenize{Classes_lammps:_CPPv3N9LAMMPS_NS6LAMMPS6LAMMPSEiPPc8MPI_Comm}}\phantomsection\label{\detokenize{Classes_lammps:_CPPv2N9LAMMPS_NS6LAMMPS6LAMMPSEiPPc8MPI_Comm}}\phantomsection\label{\detokenize{Classes_lammps:LAMMPS_NS::LAMMPS::LAMMPS__i.cPP.MPI_Comm}}
\pysigstartsignatures
\pysigstartmultiline
\pysiglinewithargsret{\phantomsection\label{\detokenize{Classes_lammps:classLAMMPS__NS_1_1LAMMPS_1ad442476f0b5d6babaf85c9fdb3c3fbc7}}\sphinxbfcode{\sphinxupquote{\DUrole{n}{LAMMPS}}}}{\DUrole{kt}{int}\sphinxparamcomma \DUrole{kt}{char}\DUrole{p}{*}\DUrole{p}{*}\sphinxparamcomma \DUrole{n}{MPI\_Comm}}{}
\pysigstopmultiline
\pysigstopsignatures
\sphinxAtStartPar
Create a {\hyperref[\detokenize{Classes_lammps:classLAMMPS__NS_1_1LAMMPS}]{\sphinxcrossref{\DUrole{std,std-ref}{LAMMPS}}}} simulation instance

\sphinxAtStartPar
The {\hyperref[\detokenize{Classes_lammps:classLAMMPS__NS_1_1LAMMPS}]{\sphinxcrossref{\DUrole{std,std-ref}{LAMMPS}}}} constructor starts up a simulation by allocating all fundamental classes in the necessary order, parses input switches and their arguments, initializes communicators, screen and logfile output FILE pointers.
\begin{quote}\begin{description}
\sphinxlineitem{Parameters}\begin{itemize}
\item {} 
\sphinxAtStartPar
\sphinxstyleliteralstrong{\sphinxupquote{narg}} \textendash{} number of arguments 

\item {} 
\sphinxAtStartPar
\sphinxstyleliteralstrong{\sphinxupquote{arg}} \textendash{} list of arguments 

\item {} 
\sphinxAtStartPar
\sphinxstyleliteralstrong{\sphinxupquote{communicator}} \textendash{} MPI communicator used by this {\hyperref[\detokenize{Classes_lammps:classLAMMPS__NS_1_1LAMMPS}]{\sphinxcrossref{\DUrole{std,std-ref}{LAMMPS}}}} instance 

\end{itemize}

\end{description}\end{quote}

\end{fulllineitems}

\index{LAMMPS::\textasciitilde{}LAMMPS (C++ function)@\spxentry{LAMMPS::\textasciitilde{}LAMMPS}\spxextra{C++ function}}

\begin{fulllineitems}
\phantomsection\label{\detokenize{Classes_lammps:_CPPv4N9LAMMPS_NS6LAMMPSD0Ev}}\phantomsection\label{\detokenize{Classes_lammps:_CPPv3N9LAMMPS_NS6LAMMPSD0Ev}}\phantomsection\label{\detokenize{Classes_lammps:_CPPv2N9LAMMPS_NS6LAMMPSD0Ev}}\phantomsection\label{\detokenize{Classes_lammps:LAMMPS_NS::LAMMPS::_LAMMPS}}
\pysigstartsignatures
\pysigstartmultiline
\pysiglinewithargsret{\phantomsection\label{\detokenize{Classes_lammps:classLAMMPS__NS_1_1LAMMPS_1a92a3e5374e6a9e114052e6a5abcb4bec}}\sphinxbfcode{\sphinxupquote{\DUrole{n}{\textasciitilde{}LAMMPS}}}}{}{\DUrole{w}{ }\DUrole{k}{noexcept}\DUrole{p}{(}\DUrole{k}{false}\DUrole{p}{)}}
\pysigstopmultiline
\pysigstopsignatures
\sphinxAtStartPar
Shut down a {\hyperref[\detokenize{Classes_lammps:classLAMMPS__NS_1_1LAMMPS}]{\sphinxcrossref{\DUrole{std,std-ref}{LAMMPS}}}} simulation instance

\sphinxAtStartPar
The {\hyperref[\detokenize{Classes_lammps:classLAMMPS__NS_1_1LAMMPS}]{\sphinxcrossref{\DUrole{std,std-ref}{LAMMPS}}}} destructor shuts down the simulation by deleting top\sphinxhyphen{}level class instances, closing screen and log files for the global instance (aka “world”) and files and MPI communicators in sub\sphinxhyphen{}partitions (“universes”). Then it deletes the fundamental class instances and copies of data inside the class. 

\end{fulllineitems}


\end{sphinxuseclass}
\begin{sphinxuseclass}{breathe-sectiondef}\subsubsection*{Public Static Functions}
\index{LAMMPS::is\_installed\_pkg (C++ function)@\spxentry{LAMMPS::is\_installed\_pkg}\spxextra{C++ function}}

\begin{fulllineitems}
\phantomsection\label{\detokenize{Classes_lammps:_CPPv4N9LAMMPS_NS6LAMMPS16is_installed_pkgEPKc}}\phantomsection\label{\detokenize{Classes_lammps:_CPPv3N9LAMMPS_NS6LAMMPS16is_installed_pkgEPKc}}\phantomsection\label{\detokenize{Classes_lammps:_CPPv2N9LAMMPS_NS6LAMMPS16is_installed_pkgEPKc}}\phantomsection\label{\detokenize{Classes_lammps:LAMMPS_NS::LAMMPS::is_installed_pkg__cCP}}
\pysigstartsignatures
\pysigstartmultiline
\pysiglinewithargsret{\phantomsection\label{\detokenize{Classes_lammps:classLAMMPS__NS_1_1LAMMPS_1a0e0366f32d92ec48bd0c034826e9c101}}\DUrole{k}{static}\DUrole{w}{ }\DUrole{kt}{bool}\DUrole{w}{ }\sphinxbfcode{\sphinxupquote{\DUrole{n}{is\_installed\_pkg}}}}{\DUrole{k}{const}\DUrole{w}{ }\DUrole{kt}{char}\DUrole{w}{ }\DUrole{p}{*}\DUrole{n,sig-param}{pkg}}{}
\pysigstopmultiline
\pysigstopsignatures
\sphinxAtStartPar
Return true if a {\hyperref[\detokenize{Classes_lammps:classLAMMPS__NS_1_1LAMMPS}]{\sphinxcrossref{\DUrole{std,std-ref}{LAMMPS}}}} package is enabled in this binary
\begin{quote}\begin{description}
\sphinxlineitem{Parameters}
\sphinxAtStartPar
\sphinxstyleliteralstrong{\sphinxupquote{pkg}} \textendash{} name of package 

\sphinxlineitem{Returns}
\sphinxAtStartPar
true if yes, else false 

\end{description}\end{quote}

\end{fulllineitems}

\index{LAMMPS::argv\_pointers (C++ function)@\spxentry{LAMMPS::argv\_pointers}\spxextra{C++ function}}

\begin{fulllineitems}
\phantomsection\label{\detokenize{Classes_lammps:_CPPv4N9LAMMPS_NS6LAMMPS13argv_pointersER4argv}}\phantomsection\label{\detokenize{Classes_lammps:_CPPv3N9LAMMPS_NS6LAMMPS13argv_pointersER4argv}}\phantomsection\label{\detokenize{Classes_lammps:_CPPv2N9LAMMPS_NS6LAMMPS13argv_pointersER4argv}}\phantomsection\label{\detokenize{Classes_lammps:LAMMPS_NS::LAMMPS::argv_pointers__argvR}}
\pysigstartsignatures
\pysigstartmultiline
\pysiglinewithargsret{\phantomsection\label{\detokenize{Classes_lammps:classLAMMPS__NS_1_1LAMMPS_1a66bef5337d2a927cad7cd4235826defb}}\DUrole{k}{static}\DUrole{w}{ }\DUrole{n}{std}\DUrole{p}{::}\DUrole{n}{vector}\DUrole{p}{\textless{}}\DUrole{kt}{char}\DUrole{p}{*}\DUrole{p}{\textgreater{}}\DUrole{w}{ }\sphinxbfcode{\sphinxupquote{\DUrole{n}{argv\_pointers}}}}{\DUrole{n}{argv}\DUrole{w}{ }\DUrole{p}{\&}\DUrole{n,sig-param}{args}}{}
\pysigstopmultiline
\pysigstopsignatures
\sphinxAtStartPar
Create vector of argv string pointers including terminating nullptr element
\begin{quote}\begin{description}
\sphinxlineitem{Parameters}
\sphinxAtStartPar
\sphinxstyleliteralstrong{\sphinxupquote{args}} \textendash{} list of arguments 

\end{description}\end{quote}

\end{fulllineitems}


\end{sphinxuseclass}
\end{fulllineitems}

\index{Pointers (C++ class)@\spxentry{Pointers}\spxextra{C++ class}}

\begin{fulllineitems}
\phantomsection\label{\detokenize{Classes_lammps:_CPPv4N9LAMMPS_NS8PointersE}}\phantomsection\label{\detokenize{Classes_lammps:_CPPv3N9LAMMPS_NS8PointersE}}\phantomsection\label{\detokenize{Classes_lammps:_CPPv2N9LAMMPS_NS8PointersE}}\phantomsection\label{\detokenize{Classes_lammps:LAMMPS_NS::Pointers}}
\pysigstartsignatures
\pysigstartmultiline
\pysigline{\phantomsection\label{\detokenize{Classes_lammps:classLAMMPS__NS_1_1Pointers}}\DUrole{k}{class}\DUrole{w}{ }\sphinxbfcode{\sphinxupquote{\DUrole{n}{Pointers}}}}
\pysigstopmultiline
\pysigstopsignatures
\sphinxAtStartPar
Base class for {\hyperref[\detokenize{Classes_lammps:classLAMMPS__NS_1_1LAMMPS}]{\sphinxcrossref{\DUrole{std,std-ref}{LAMMPS}}}} features. 

\sphinxAtStartPar
The {\hyperref[\detokenize{Classes_lammps:classLAMMPS__NS_1_1Pointers}]{\sphinxcrossref{\DUrole{std,std-ref}{Pointers}}}} class contains references to many of the pointers and members of the {\hyperref[\detokenize{Classes_lammps:classLAMMPS__NS_1_1LAMMPS}]{\sphinxcrossref{\DUrole{std,std-ref}{LAMMPS\_NS::LAMMPS}}}} class. Derived classes thus gain access to the constituent class instances in the {\hyperref[\detokenize{Classes_lammps:classLAMMPS__NS_1_1LAMMPS}]{\sphinxcrossref{\DUrole{std,std-ref}{LAMMPS}}}} composite class and thus to the core functionality of {\hyperref[\detokenize{Classes_lammps:classLAMMPS__NS_1_1LAMMPS}]{\sphinxcrossref{\DUrole{std,std-ref}{LAMMPS}}}}.

\sphinxAtStartPar
This kind of construct is needed, since the {\hyperref[\detokenize{Classes_lammps:classLAMMPS__NS_1_1LAMMPS}]{\sphinxcrossref{\DUrole{std,std-ref}{LAMMPS}}}} constructor should only be run once per {\hyperref[\detokenize{Classes_lammps:classLAMMPS__NS_1_1LAMMPS}]{\sphinxcrossref{\DUrole{std,std-ref}{LAMMPS}}}} instance and thus classes cannot be derived from {\hyperref[\detokenize{Classes_lammps:classLAMMPS__NS_1_1LAMMPS}]{\sphinxcrossref{\DUrole{std,std-ref}{LAMMPS}}}} itself. The {\hyperref[\detokenize{Classes_lammps:classLAMMPS__NS_1_1Pointers}]{\sphinxcrossref{\DUrole{std,std-ref}{Pointers}}}} class constructor instead only initializes C++ references to component pointer in the {\hyperref[\detokenize{Classes_lammps:classLAMMPS__NS_1_1LAMMPS}]{\sphinxcrossref{\DUrole{std,std-ref}{LAMMPS}}}} class. 

\sphinxAtStartPar
Subclassed by {\hyperref[\detokenize{Classes_atom:classLAMMPS__NS_1_1Atom}]{\sphinxcrossref{\DUrole{std,std-ref}{Atom}}}}, {\hyperref[\detokenize{Classes_input:classLAMMPS__NS_1_1Input}]{\sphinxcrossref{\DUrole{std,std-ref}{Input}}}}, {\hyperref[\detokenize{Developer_utils:classLAMMPS__NS_1_1PotentialFileReader}]{\sphinxcrossref{\DUrole{std,std-ref}{PotentialFileReader}}}}

\end{fulllineitems}


\sphinxstepscope


\subsection{LAMMPS Atom and AtomVec Base Classes}
\label{\detokenize{Classes_atom:lammps-atom-and-atomvec-base-classes}}\label{\detokenize{Classes_atom::doc}}\index{Atom (C++ class)@\spxentry{Atom}\spxextra{C++ class}}

\begin{fulllineitems}
\phantomsection\label{\detokenize{Classes_atom:_CPPv4N9LAMMPS_NS4AtomE}}\phantomsection\label{\detokenize{Classes_atom:_CPPv3N9LAMMPS_NS4AtomE}}\phantomsection\label{\detokenize{Classes_atom:_CPPv2N9LAMMPS_NS4AtomE}}\phantomsection\label{\detokenize{Classes_atom:LAMMPS_NS::Atom}}
\pysigstartsignatures
\pysigstartmultiline
\pysigline{\phantomsection\label{\detokenize{Classes_atom:classLAMMPS__NS_1_1Atom}}\DUrole{k}{class}\DUrole{w}{ }\sphinxbfcode{\sphinxupquote{\DUrole{n}{Atom}}}\DUrole{w}{ }\DUrole{p}{:}\DUrole{w}{ }\DUrole{k}{protected}\DUrole{w}{ }{\hyperref[\detokenize{Classes_lammps:_CPPv4N9LAMMPS_NS8PointersE}]{\sphinxcrossref{\DUrole{n}{Pointers}}}}}
\pysigstopmultiline
\pysigstopsignatures
\sphinxAtStartPar
Class to provide access to atom data. 

\sphinxAtStartPar

\sphinxAtStartPar
The Atom class provides access to atom style related global settings and
per\sphinxhyphen{}atom data that is stored with atoms and migrates with them from
sub\sphinxhyphen{}domain to sub\sphinxhyphen{}domain as atoms move around.  This includes topology
data, which is stored with either one specific atom or all atoms involved
depending on the settings of the {\hyperref[\detokenize{newton::doc}]{\sphinxcrossref{\DUrole{doc}{newton command}}}}.

\sphinxAtStartPar
The actual per\sphinxhyphen{}atom data is allocated and managed by one of the various
classes derived from the AtomVec class as determined by
the {\hyperref[\detokenize{atom_style::doc}]{\sphinxcrossref{\DUrole{doc}{atom\_style command}}}}.  The pointers in the Atom class
are updated by the AtomVec class as needed.


\begin{sphinxuseclass}{breathe-sectiondef}\subsubsection*{Public Functions}
\index{Atom::Atom (C++ function)@\spxentry{Atom::Atom}\spxextra{C++ function}}

\begin{fulllineitems}
\phantomsection\label{\detokenize{Classes_atom:_CPPv4N9LAMMPS_NS4Atom4AtomEP6LAMMPS}}\phantomsection\label{\detokenize{Classes_atom:_CPPv3N9LAMMPS_NS4Atom4AtomEP6LAMMPS}}\phantomsection\label{\detokenize{Classes_atom:_CPPv2N9LAMMPS_NS4Atom4AtomEP6LAMMPS}}\phantomsection\label{\detokenize{Classes_atom:LAMMPS_NS::Atom::Atom__LAMMPSP}}
\pysigstartsignatures
\pysigstartmultiline
\pysiglinewithargsret{\phantomsection\label{\detokenize{Classes_atom:classLAMMPS__NS_1_1Atom_1a0d2e320788dddcb5ef150c187679f905}}\sphinxbfcode{\sphinxupquote{\DUrole{n}{Atom}}}}{\DUrole{k}{class}\DUrole{w}{ }{\hyperref[\detokenize{Classes_lammps:_CPPv4N9LAMMPS_NS6LAMMPSE}]{\sphinxcrossref{\DUrole{n}{LAMMPS}}}}\DUrole{p}{*}}{}
\pysigstopmultiline
\pysigstopsignatures
\sphinxAtStartPar
{\hyperref[\detokenize{Classes_atom:classLAMMPS__NS_1_1Atom}]{\sphinxcrossref{\DUrole{std,std-ref}{Atom}}}} class constructor

\sphinxAtStartPar
This resets and initializes all kinds of settings, parameters, and pointer variables for per\sphinxhyphen{}atom arrays. This also initializes the factory for creating instances of classes derived from the AtomVec base class, which correspond to the selected atom style.
\begin{quote}\begin{description}
\sphinxlineitem{Parameters}
\sphinxAtStartPar
\sphinxstyleliteralstrong{\sphinxupquote{\_lmp}} \textendash{} pointer to the base {\hyperref[\detokenize{Classes_lammps:classLAMMPS__NS_1_1LAMMPS}]{\sphinxcrossref{\DUrole{std,std-ref}{LAMMPS}}}} class 

\end{description}\end{quote}

\end{fulllineitems}

\index{Atom::find\_custom (C++ function)@\spxentry{Atom::find\_custom}\spxextra{C++ function}}

\begin{fulllineitems}
\phantomsection\label{\detokenize{Classes_atom:_CPPv4N9LAMMPS_NS4Atom11find_customEPKcRiRi}}\phantomsection\label{\detokenize{Classes_atom:_CPPv3N9LAMMPS_NS4Atom11find_customEPKcRiRi}}\phantomsection\label{\detokenize{Classes_atom:_CPPv2N9LAMMPS_NS4Atom11find_customEPKcRiRi}}\phantomsection\label{\detokenize{Classes_atom:LAMMPS_NS::Atom::find_custom__cCP.iR.iR}}
\pysigstartsignatures
\pysigstartmultiline
\pysiglinewithargsret{\phantomsection\label{\detokenize{Classes_atom:classLAMMPS__NS_1_1Atom_1a6f6bc746553cf0ac64e327aced14e958}}\DUrole{kt}{int}\DUrole{w}{ }\sphinxbfcode{\sphinxupquote{\DUrole{n}{find\_custom}}}}{\DUrole{k}{const}\DUrole{w}{ }\DUrole{kt}{char}\DUrole{p}{*}\sphinxparamcomma \DUrole{kt}{int}\DUrole{p}{\&}\sphinxparamcomma \DUrole{kt}{int}\DUrole{p}{\&}}{}
\pysigstopmultiline
\pysigstopsignatures
\sphinxAtStartPar
Find a custom per\sphinxhyphen{}atom property with given name. 

\sphinxAtStartPar

\sphinxAtStartPar
This function returns the list index of a custom per\sphinxhyphen{}atom property
with the name “name”, also returning by reference its data type and
number of values per atom.
 
\begin{quote}\begin{description}
\sphinxlineitem{Parameters}\begin{itemize}
\item {} 
\sphinxAtStartPar
\sphinxstyleliteralstrong{\sphinxupquote{name}} \textendash{} Name of the property (w/o a “i\_” or “d\_” or “i2\_” or “d2\_” prefix) 

\item {} 
\sphinxAtStartPar
\sphinxstyleliteralstrong{\sphinxupquote{\&flag}} \textendash{} Returns data type of property: 0 for int, 1 for double 

\item {} 
\sphinxAtStartPar
\sphinxstyleliteralstrong{\sphinxupquote{\&cols}} \textendash{} Returns number of values: 0 for a single value, 1 or more for a vector of values 

\end{itemize}

\sphinxlineitem{Returns}
\sphinxAtStartPar
index of property in the respective list of properties 

\end{description}\end{quote}

\end{fulllineitems}

\index{Atom::find\_custom\_ghost (C++ function)@\spxentry{Atom::find\_custom\_ghost}\spxextra{C++ function}}

\begin{fulllineitems}
\phantomsection\label{\detokenize{Classes_atom:_CPPv4N9LAMMPS_NS4Atom17find_custom_ghostEPKcRiRiRi}}\phantomsection\label{\detokenize{Classes_atom:_CPPv3N9LAMMPS_NS4Atom17find_custom_ghostEPKcRiRiRi}}\phantomsection\label{\detokenize{Classes_atom:_CPPv2N9LAMMPS_NS4Atom17find_custom_ghostEPKcRiRiRi}}\phantomsection\label{\detokenize{Classes_atom:LAMMPS_NS::Atom::find_custom_ghost__cCP.iR.iR.iR}}
\pysigstartsignatures
\pysigstartmultiline
\pysiglinewithargsret{\phantomsection\label{\detokenize{Classes_atom:classLAMMPS__NS_1_1Atom_1a3ce6ca4cdb300ea3ff044b800d1180ed}}\DUrole{kt}{int}\DUrole{w}{ }\sphinxbfcode{\sphinxupquote{\DUrole{n}{find\_custom\_ghost}}}}{\DUrole{k}{const}\DUrole{w}{ }\DUrole{kt}{char}\DUrole{p}{*}\sphinxparamcomma \DUrole{kt}{int}\DUrole{p}{\&}\sphinxparamcomma \DUrole{kt}{int}\DUrole{p}{\&}\sphinxparamcomma \DUrole{kt}{int}\DUrole{p}{\&}}{}
\pysigstopmultiline
\pysigstopsignatures
\sphinxAtStartPar
Find a custom per\sphinxhyphen{}atom property with given name and retrieve ghost property. 

\sphinxAtStartPar

\sphinxAtStartPar
This function returns the list index of a custom per\sphinxhyphen{}atom property
with the name “name”, also returning by reference its data type,
number of values per atom, and if it is communicated to ghost particles.
Classes rarely need to check on ghost communication and so \sphinxtitleref{find\_custom}
is typically preferred to this function. See {\hyperref[\detokenize{pair_amoeba::doc}]{\sphinxcrossref{\DUrole{doc}{pair amoeba}}}}
for an example where checking ghost communication is necessary.
 
\begin{quote}\begin{description}
\sphinxlineitem{Parameters}\begin{itemize}
\item {} 
\sphinxAtStartPar
\sphinxstyleliteralstrong{\sphinxupquote{name}} \textendash{} Name of the property (w/o a “i\_” or “d\_” or “i2\_” or “d2\_” prefix) 

\item {} 
\sphinxAtStartPar
\sphinxstyleliteralstrong{\sphinxupquote{\&flag}} \textendash{} Returns data type of property: 0 for int, 1 for double 

\item {} 
\sphinxAtStartPar
\sphinxstyleliteralstrong{\sphinxupquote{\&cols}} \textendash{} Returns number of values: 0 for a single value, 1 or more for a vector of values 

\item {} 
\sphinxAtStartPar
\sphinxstyleliteralstrong{\sphinxupquote{\&ghost}} \textendash{} Returns whether property is communicated to ghost atoms: 0 for no, 1 for yes 

\end{itemize}

\sphinxlineitem{Returns}
\sphinxAtStartPar
index of property in the respective list of properties 

\end{description}\end{quote}

\end{fulllineitems}

\index{Atom::add\_custom (C++ function)@\spxentry{Atom::add\_custom}\spxextra{C++ function}}

\begin{fulllineitems}
\phantomsection\label{\detokenize{Classes_atom:_CPPv4N9LAMMPS_NS4Atom10add_customEPKciii}}\phantomsection\label{\detokenize{Classes_atom:_CPPv3N9LAMMPS_NS4Atom10add_customEPKciii}}\phantomsection\label{\detokenize{Classes_atom:_CPPv2N9LAMMPS_NS4Atom10add_customEPKciii}}\phantomsection\label{\detokenize{Classes_atom:LAMMPS_NS::Atom::add_custom__cCP.i.i.i}}
\pysigstartsignatures
\pysigstartmultiline
\pysiglinewithargsret{\phantomsection\label{\detokenize{Classes_atom:classLAMMPS__NS_1_1Atom_1acd074d93188a3df2a43d08234b1691b9}}\DUrole{k}{virtual}\DUrole{w}{ }\DUrole{kt}{int}\DUrole{w}{ }\sphinxbfcode{\sphinxupquote{\DUrole{n}{add\_custom}}}}{\DUrole{k}{const}\DUrole{w}{ }\DUrole{kt}{char}\DUrole{p}{*}\sphinxparamcomma \DUrole{kt}{int}\sphinxparamcomma \DUrole{kt}{int}\sphinxparamcomma \DUrole{kt}{int}\DUrole{w}{ }\DUrole{n,sig-param}{ghost}\DUrole{w}{ }\DUrole{p}{=}\DUrole{w}{ }\DUrole{m}{0}}{}
\pysigstopmultiline
\pysigstopsignatures
\sphinxAtStartPar
Add a custom per\sphinxhyphen{}atom property with the given name and type and size. 

\sphinxAtStartPar

\sphinxAtStartPar
This function will add a custom per\sphinxhyphen{}atom property with one or more values
with the name “name” to the list of custom properties.
This function is called, e.g. from {\hyperref[\detokenize{fix_property_atom::doc}]{\sphinxcrossref{\DUrole{doc}{fix property/atom}}}}.
 
\begin{quote}\begin{description}
\sphinxlineitem{Parameters}\begin{itemize}
\item {} 
\sphinxAtStartPar
\sphinxstyleliteralstrong{\sphinxupquote{name}} \textendash{} Name of the property (w/o a “i\_” or “d\_” or “i2\_” or “d2\_” prefix) 

\item {} 
\sphinxAtStartPar
\sphinxstyleliteralstrong{\sphinxupquote{flag}} \textendash{} Data type of property: 0 for int, 1 for double 

\item {} 
\sphinxAtStartPar
\sphinxstyleliteralstrong{\sphinxupquote{cols}} \textendash{} Number of values: 0 for a single value, 1 or more for a vector of values 

\item {} 
\sphinxAtStartPar
\sphinxstyleliteralstrong{\sphinxupquote{ghost}} \textendash{} Whether property is communicated to ghost atoms: 0 for no, 1 for yes 

\end{itemize}

\sphinxlineitem{Returns}
\sphinxAtStartPar
index of property in the respective list of properties 

\end{description}\end{quote}

\end{fulllineitems}

\index{Atom::remove\_custom (C++ function)@\spxentry{Atom::remove\_custom}\spxextra{C++ function}}

\begin{fulllineitems}
\phantomsection\label{\detokenize{Classes_atom:_CPPv4N9LAMMPS_NS4Atom13remove_customEiii}}\phantomsection\label{\detokenize{Classes_atom:_CPPv3N9LAMMPS_NS4Atom13remove_customEiii}}\phantomsection\label{\detokenize{Classes_atom:_CPPv2N9LAMMPS_NS4Atom13remove_customEiii}}\phantomsection\label{\detokenize{Classes_atom:LAMMPS_NS::Atom::remove_custom__i.i.i}}
\pysigstartsignatures
\pysigstartmultiline
\pysiglinewithargsret{\phantomsection\label{\detokenize{Classes_atom:classLAMMPS__NS_1_1Atom_1aecf945aa87239b81163cc602aa7c5549}}\DUrole{k}{virtual}\DUrole{w}{ }\DUrole{kt}{void}\DUrole{w}{ }\sphinxbfcode{\sphinxupquote{\DUrole{n}{remove\_custom}}}}{\DUrole{kt}{int}\sphinxparamcomma \DUrole{kt}{int}\sphinxparamcomma \DUrole{kt}{int}}{}
\pysigstopmultiline
\pysigstopsignatures
\sphinxAtStartPar
Remove a custom per\sphinxhyphen{}atom property of a given type and size. 

\sphinxAtStartPar

\sphinxAtStartPar
This will remove a property that was requested, e.g. by the
{\hyperref[\detokenize{fix_property_atom::doc}]{\sphinxcrossref{\DUrole{doc}{fix property/atom}}}} command.  It frees the
allocated memory and sets the pointer to \sphinxcode{\sphinxupquote{nullptr}} for the entry in
the list so it can be reused. The lists of these pointers are never
compacted or shrunk, so that indices to name mappings remain valid.
 
\begin{quote}\begin{description}
\sphinxlineitem{Parameters}\begin{itemize}
\item {} 
\sphinxAtStartPar
\sphinxstyleliteralstrong{\sphinxupquote{index}} \textendash{} Index of property in the respective list of properties 

\item {} 
\sphinxAtStartPar
\sphinxstyleliteralstrong{\sphinxupquote{flag}} \textendash{} Data type of property: 0 for int, 1 for double 

\item {} 
\sphinxAtStartPar
\sphinxstyleliteralstrong{\sphinxupquote{cols}} \textendash{} Number of values: 0 for a single value, 1 or more for a vector of values 

\end{itemize}

\end{description}\end{quote}

\end{fulllineitems}

\index{Atom::extract (C++ function)@\spxentry{Atom::extract}\spxextra{C++ function}}

\begin{fulllineitems}
\phantomsection\label{\detokenize{Classes_atom:_CPPv4N9LAMMPS_NS4Atom7extractEPKc}}\phantomsection\label{\detokenize{Classes_atom:_CPPv3N9LAMMPS_NS4Atom7extractEPKc}}\phantomsection\label{\detokenize{Classes_atom:_CPPv2N9LAMMPS_NS4Atom7extractEPKc}}\phantomsection\label{\detokenize{Classes_atom:LAMMPS_NS::Atom::extract__cCP}}
\pysigstartsignatures
\pysigstartmultiline
\pysiglinewithargsret{\phantomsection\label{\detokenize{Classes_atom:classLAMMPS__NS_1_1Atom_1aa12f2b783db9b4523ddbf9f01d132d88}}\DUrole{kt}{void}\DUrole{w}{ }\DUrole{p}{*}\sphinxbfcode{\sphinxupquote{\DUrole{n}{extract}}}}{\DUrole{k}{const}\DUrole{w}{ }\DUrole{kt}{char}\DUrole{p}{*}}{}
\pysigstopmultiline
\pysigstopsignatures
\sphinxAtStartPar
Provide access to internal data of the {\hyperref[\detokenize{Classes_atom:classLAMMPS__NS_1_1Atom}]{\sphinxcrossref{\DUrole{std,std-ref}{Atom}}}} class by keyword

\sphinxAtStartPar

\sphinxAtStartPar
This function is a way to access internal per\sphinxhyphen{}atom data.  This data is
distributed across MPI ranks and thus only the data for “local” atoms
are expected to be available.  Whether also data for “ghost” atoms is
stored and up\sphinxhyphen{}to\sphinxhyphen{}date depends on various simulation settings.

\sphinxAtStartPar
This table lists a large part of the supported names, their data types,
length of the data area, and a short description.


\begin{savenotes}\sphinxattablestart
\sphinxthistablewithglobalstyle
\centering
\begin{tabulary}{\linewidth}[t]{TTTT}
\sphinxtoprule
\sphinxstyletheadfamily 
\sphinxAtStartPar
Name
&\sphinxstyletheadfamily 
\sphinxAtStartPar
Type
&\sphinxstyletheadfamily 
\sphinxAtStartPar
Items per atom
&\sphinxstyletheadfamily 
\sphinxAtStartPar
Description
\\
\sphinxmidrule
\sphinxtableatstartofbodyhook
\sphinxAtStartPar
mass
&
\sphinxAtStartPar
double
&
\sphinxAtStartPar
1
&
\sphinxAtStartPar
per\sphinxhyphen{}type mass. This array is \sphinxstylestrong{NOT} a per\sphinxhyphen{}atom array
but of length \sphinxcode{\sphinxupquote{ntypes+1}}, element 0 is ignored.
\\
\sphinxhline
\sphinxAtStartPar
id
&
\sphinxAtStartPar
tagint
&
\sphinxAtStartPar
1
&
\sphinxAtStartPar
atom ID of the particles
\\
\sphinxhline
\sphinxAtStartPar
type
&
\sphinxAtStartPar
int
&
\sphinxAtStartPar
1
&
\sphinxAtStartPar
atom type of the particles
\\
\sphinxhline
\sphinxAtStartPar
mask
&
\sphinxAtStartPar
int
&
\sphinxAtStartPar
1
&
\sphinxAtStartPar
bitmask for mapping to groups. Individual bits are set
to 0 or 1 for each group.
\\
\sphinxhline
\sphinxAtStartPar
image
&
\sphinxAtStartPar
imageint
&
\sphinxAtStartPar
1
&
\sphinxAtStartPar
3 image flags encoded into a single integer.
See {\hyperref[\detokenize{Library_utility:_CPPv425lammps_encode_image_flagsiii}]{\sphinxcrossref{\sphinxcode{\sphinxupquote{lammps\_encode\_image\_flags()}}}}}.
\\
\sphinxhline
\sphinxAtStartPar
x
&
\sphinxAtStartPar
double
&
\sphinxAtStartPar
3
&
\sphinxAtStartPar
x\sphinxhyphen{}, y\sphinxhyphen{}, and z\sphinxhyphen{}coordinate of the particles
\\
\sphinxhline
\sphinxAtStartPar
v
&
\sphinxAtStartPar
double
&
\sphinxAtStartPar
3
&
\sphinxAtStartPar
x\sphinxhyphen{}, y\sphinxhyphen{}, and z\sphinxhyphen{}component of the velocity of the particles
\\
\sphinxhline
\sphinxAtStartPar
f
&
\sphinxAtStartPar
double
&
\sphinxAtStartPar
3
&
\sphinxAtStartPar
x\sphinxhyphen{}, y\sphinxhyphen{}, and z\sphinxhyphen{}component of the force on the particles
\\
\sphinxhline
\sphinxAtStartPar
molecule
&
\sphinxAtStartPar
int
&
\sphinxAtStartPar
1
&
\sphinxAtStartPar
molecule ID of the particles
\\
\sphinxhline
\sphinxAtStartPar
q
&
\sphinxAtStartPar
double
&
\sphinxAtStartPar
1
&
\sphinxAtStartPar
charge of the particles
\\
\sphinxhline
\sphinxAtStartPar
mu
&
\sphinxAtStartPar
double
&
\sphinxAtStartPar
3
&
\sphinxAtStartPar
dipole moment of the particles
\\
\sphinxhline
\sphinxAtStartPar
omega
&
\sphinxAtStartPar
double
&
\sphinxAtStartPar
3
&
\sphinxAtStartPar
x\sphinxhyphen{}, y\sphinxhyphen{}, and z\sphinxhyphen{}component of rotational velocity of the particles
\\
\sphinxhline
\sphinxAtStartPar
angmom
&
\sphinxAtStartPar
double
&
\sphinxAtStartPar
3
&
\sphinxAtStartPar
x\sphinxhyphen{}, y\sphinxhyphen{}, and z\sphinxhyphen{}component of angular momentum of the particles
\\
\sphinxhline
\sphinxAtStartPar
torque
&
\sphinxAtStartPar
double
&
\sphinxAtStartPar
3
&
\sphinxAtStartPar
x\sphinxhyphen{}, y\sphinxhyphen{}, and z\sphinxhyphen{}component of the torque on the particles
\\
\sphinxhline
\sphinxAtStartPar
radius
&
\sphinxAtStartPar
double
&
\sphinxAtStartPar
1
&
\sphinxAtStartPar
radius of the (extended) particles
\\
\sphinxhline
\sphinxAtStartPar
rmass
&
\sphinxAtStartPar
double
&
\sphinxAtStartPar
1
&
\sphinxAtStartPar
per\sphinxhyphen{}atom mass of the particles. \sphinxcode{\sphinxupquote{nullptr}} if per\sphinxhyphen{}type masses are
used. See the {\hyperref[\detokenize{Library_properties:_CPPv422lammps_extract_settingPvPKc}]{\sphinxcrossref{\sphinxcode{\sphinxupquote{rmass\_flag}}}}} setting.
\\
\sphinxhline
\sphinxAtStartPar
ellipsoid
&
\sphinxAtStartPar
int
&
\sphinxAtStartPar
1
&
\sphinxAtStartPar
1 if the particle is an ellipsoidal particle, 0 if not
\\
\sphinxhline
\sphinxAtStartPar
line
&
\sphinxAtStartPar
int
&
\sphinxAtStartPar
1
&
\sphinxAtStartPar
1 if the particle is a line particle, 0 if not
\\
\sphinxhline
\sphinxAtStartPar
tri
&
\sphinxAtStartPar
int
&
\sphinxAtStartPar
1
&
\sphinxAtStartPar
1 if the particle is a triangulated particle, 0 if not
\\
\sphinxhline
\sphinxAtStartPar
body
&
\sphinxAtStartPar
int
&
\sphinxAtStartPar
1
&
\sphinxAtStartPar
1 if the particle is a body particle, 0 if not
\\
\sphinxhline
\sphinxAtStartPar
quat
&
\sphinxAtStartPar
double
&
\sphinxAtStartPar
4
&
\sphinxAtStartPar
four quaternion components of the particles
\\
\sphinxhline
\sphinxAtStartPar
temperature
&
\sphinxAtStartPar
double
&
\sphinxAtStartPar
1
&
\sphinxAtStartPar
temperature of the particles
\\
\sphinxhline
\sphinxAtStartPar
heatflow
&
\sphinxAtStartPar
double
&
\sphinxAtStartPar
1
&
\sphinxAtStartPar
heatflow of the particles
\\
\sphinxhline
\sphinxAtStartPar
i\_name
&
\sphinxAtStartPar
int
&
\sphinxAtStartPar
1
&
\sphinxAtStartPar
single integer value defined by fix property/atom vector name
\\
\sphinxhline
\sphinxAtStartPar
d\_name
&
\sphinxAtStartPar
double
&
\sphinxAtStartPar
1
&
\sphinxAtStartPar
single double value defined by fix property/atom vector name
\\
\sphinxhline
\sphinxAtStartPar
i2\_name
&
\sphinxAtStartPar
int
&
\sphinxAtStartPar
N
&
\sphinxAtStartPar
N integer values defined by fix property/atom array name
\\
\sphinxhline
\sphinxAtStartPar
d2\_name
&
\sphinxAtStartPar
double
&
\sphinxAtStartPar
N
&
\sphinxAtStartPar
N double values defined by fix property/atom array name
\\
\sphinxbottomrule
\end{tabulary}
\sphinxtableafterendhook\par
\sphinxattableend\end{savenotes}
\begin{description}
\sphinxlineitem{\sphinxstyleemphasis{See also}}
\sphinxAtStartPar
{\hyperref[\detokenize{Library_atoms:_CPPv419lammps_extract_atomPvPKc}]{\sphinxcrossref{\sphinxcode{\sphinxupquote{lammps\_extract\_atom()}}}}}

\end{description}


\sphinxAtStartPar


\begin{sphinxseealso}{See also:}

\sphinxAtStartPar
{\hyperref[\detokenize{Classes_atom:classLAMMPS__NS_1_1Atom_1ae25da0ff714bd58c9ebf8f50e4edb81b}]{\sphinxcrossref{\DUrole{std,std-ref}{extract\_datatype}}}}


\end{sphinxseealso}


\begin{quote}\begin{description}
\sphinxlineitem{Parameters}
\sphinxAtStartPar
\sphinxstyleliteralstrong{\sphinxupquote{name}} \textendash{} string with the keyword of the desired property. Typically the name of the pointer variable returned 

\sphinxlineitem{Returns}
\sphinxAtStartPar
pointer to the requested data cast to \sphinxcode{\sphinxupquote{void *}} or \sphinxcode{\sphinxupquote{nullptr}}

\end{description}\end{quote}

\end{fulllineitems}

\index{Atom::extract\_datatype (C++ function)@\spxentry{Atom::extract\_datatype}\spxextra{C++ function}}

\begin{fulllineitems}
\phantomsection\label{\detokenize{Classes_atom:_CPPv4N9LAMMPS_NS4Atom16extract_datatypeEPKc}}\phantomsection\label{\detokenize{Classes_atom:_CPPv3N9LAMMPS_NS4Atom16extract_datatypeEPKc}}\phantomsection\label{\detokenize{Classes_atom:_CPPv2N9LAMMPS_NS4Atom16extract_datatypeEPKc}}\phantomsection\label{\detokenize{Classes_atom:LAMMPS_NS::Atom::extract_datatype__cCP}}
\pysigstartsignatures
\pysigstartmultiline
\pysiglinewithargsret{\phantomsection\label{\detokenize{Classes_atom:classLAMMPS__NS_1_1Atom_1ae25da0ff714bd58c9ebf8f50e4edb81b}}\DUrole{kt}{int}\DUrole{w}{ }\sphinxbfcode{\sphinxupquote{\DUrole{n}{extract\_datatype}}}}{\DUrole{k}{const}\DUrole{w}{ }\DUrole{kt}{char}\DUrole{p}{*}}{}
\pysigstopmultiline
\pysigstopsignatures
\sphinxAtStartPar
Provide data type info about internal data of the {\hyperref[\detokenize{Classes_atom:classLAMMPS__NS_1_1Atom}]{\sphinxcrossref{\DUrole{std,std-ref}{Atom}}}} class

\sphinxAtStartPar

\sphinxAtStartPar
\DUrole{versionmodified,added}{Added in version 18Sep2020.}


\sphinxAtStartPar


\begin{sphinxseealso}{See also:}

\sphinxAtStartPar
{\hyperref[\detokenize{Classes_atom:classLAMMPS__NS_1_1Atom_1aa12f2b783db9b4523ddbf9f01d132d88}]{\sphinxcrossref{\DUrole{std,std-ref}{extract}}}}


\end{sphinxseealso}


\begin{quote}\begin{description}
\sphinxlineitem{Parameters}
\sphinxAtStartPar
\sphinxstyleliteralstrong{\sphinxupquote{name}} \textendash{} string with the keyword of the desired property. 

\sphinxlineitem{Returns}
\sphinxAtStartPar
data type constant for desired property or \sphinxhyphen{}1 

\end{description}\end{quote}

\end{fulllineitems}


\end{sphinxuseclass}\index{Atom::PerAtom (C++ struct)@\spxentry{Atom::PerAtom}\spxextra{C++ struct}}

\begin{fulllineitems}
\phantomsection\label{\detokenize{Classes_atom:_CPPv4N9LAMMPS_NS4Atom7PerAtomE}}\phantomsection\label{\detokenize{Classes_atom:_CPPv3N9LAMMPS_NS4Atom7PerAtomE}}\phantomsection\label{\detokenize{Classes_atom:_CPPv2N9LAMMPS_NS4Atom7PerAtomE}}\phantomsection\label{\detokenize{Classes_atom:LAMMPS_NS::Atom::PerAtom}}
\pysigstartsignatures
\pysigstartmultiline
\pysigline{\phantomsection\label{\detokenize{Classes_atom:structLAMMPS__NS_1_1Atom_1_1PerAtom}}\DUrole{k}{struct}\DUrole{w}{ }\sphinxbfcode{\sphinxupquote{\DUrole{n}{PerAtom}}}}
\pysigstopmultiline
\pysigstopsignatures
\end{fulllineitems}


\end{fulllineitems}


\sphinxstepscope


\subsection{LAMMPS Input Base Class}
\label{\detokenize{Classes_input:lammps-input-base-class}}\label{\detokenize{Classes_input::doc}}\index{Input (C++ class)@\spxentry{Input}\spxextra{C++ class}}

\begin{fulllineitems}
\phantomsection\label{\detokenize{Classes_input:_CPPv4N9LAMMPS_NS5InputE}}\phantomsection\label{\detokenize{Classes_input:_CPPv3N9LAMMPS_NS5InputE}}\phantomsection\label{\detokenize{Classes_input:_CPPv2N9LAMMPS_NS5InputE}}\phantomsection\label{\detokenize{Classes_input:LAMMPS_NS::Input}}
\pysigstartsignatures
\pysigstartmultiline
\pysigline{\phantomsection\label{\detokenize{Classes_input:classLAMMPS__NS_1_1Input}}\DUrole{k}{class}\DUrole{w}{ }\sphinxbfcode{\sphinxupquote{\DUrole{n}{Input}}}\DUrole{w}{ }\DUrole{p}{:}\DUrole{w}{ }\DUrole{k}{protected}\DUrole{w}{ }{\hyperref[\detokenize{Classes_lammps:_CPPv4N9LAMMPS_NS8PointersE}]{\sphinxcrossref{\DUrole{n}{Pointers}}}}}
\pysigstopmultiline
\pysigstopsignatures
\sphinxAtStartPar
Class for processing commands and input files. 

\sphinxAtStartPar
The Input class contains methods for reading, pre\sphinxhyphen{}processing and
parsing LAMMPS commands and input files and will dispatch commands
to the respective class instances or contain the code to execute
the commands directly.  It also contains the instance of the
Variable class which performs computations and text substitutions.

\begin{sphinxuseclass}{breathe-sectiondef}\subsubsection*{Public Functions}
\index{Input::Input (C++ function)@\spxentry{Input::Input}\spxextra{C++ function}}

\begin{fulllineitems}
\phantomsection\label{\detokenize{Classes_input:_CPPv4N9LAMMPS_NS5Input5InputEP6LAMMPSiPPc}}\phantomsection\label{\detokenize{Classes_input:_CPPv3N9LAMMPS_NS5Input5InputEP6LAMMPSiPPc}}\phantomsection\label{\detokenize{Classes_input:_CPPv2N9LAMMPS_NS5Input5InputEP6LAMMPSiPPc}}\phantomsection\label{\detokenize{Classes_input:LAMMPS_NS::Input::Input__LAMMPSP.i.cPP}}
\pysigstartsignatures
\pysigstartmultiline
\pysiglinewithargsret{\phantomsection\label{\detokenize{Classes_input:classLAMMPS__NS_1_1Input_1abd21b9a5d98f5ac7b05ef80bec6f706f}}\sphinxbfcode{\sphinxupquote{\DUrole{n}{Input}}}}{\DUrole{k}{class}\DUrole{w}{ }{\hyperref[\detokenize{Classes_lammps:_CPPv4N9LAMMPS_NS6LAMMPSE}]{\sphinxcrossref{\DUrole{n}{LAMMPS}}}}\DUrole{p}{*}\sphinxparamcomma \DUrole{kt}{int}\sphinxparamcomma \DUrole{kt}{char}\DUrole{p}{*}\DUrole{p}{*}}{}
\pysigstopmultiline
\pysigstopsignatures
\sphinxAtStartPar
{\hyperref[\detokenize{Classes_input:classLAMMPS__NS_1_1Input}]{\sphinxcrossref{\DUrole{std,std-ref}{Input}}}} class constructor

\sphinxAtStartPar

\sphinxAtStartPar
This sets up the input processing, processes the \sphinxstyleemphasis{\sphinxhyphen{}var} and \sphinxstyleemphasis{\sphinxhyphen{}echo}
command line flags, holds the factory of commands and creates and
initializes an instance of the Variable class.

\sphinxAtStartPar
To execute a command, a specific class instance, derived from
\sphinxcode{\sphinxupquote{Command}}, is created, then its \sphinxcode{\sphinxupquote{command()}} member
function executed, and finally the class instance is deleted.

\begin{quote}\begin{description}
\sphinxlineitem{Parameters}\begin{itemize}
\item {} 
\sphinxAtStartPar
\sphinxstyleliteralstrong{\sphinxupquote{lmp}} \textendash{} pointer to the base {\hyperref[\detokenize{Classes_lammps:classLAMMPS__NS_1_1LAMMPS}]{\sphinxcrossref{\DUrole{std,std-ref}{LAMMPS}}}} class 

\item {} 
\sphinxAtStartPar
\sphinxstyleliteralstrong{\sphinxupquote{argc}} \textendash{} number of entries in \sphinxstyleemphasis{argv}

\item {} 
\sphinxAtStartPar
\sphinxstyleliteralstrong{\sphinxupquote{argv}} \textendash{} argument vector 

\end{itemize}

\end{description}\end{quote}

\end{fulllineitems}

\index{Input::file (C++ function)@\spxentry{Input::file}\spxextra{C++ function}}

\begin{fulllineitems}
\phantomsection\label{\detokenize{Classes_input:_CPPv4N9LAMMPS_NS5Input4fileEv}}\phantomsection\label{\detokenize{Classes_input:_CPPv3N9LAMMPS_NS5Input4fileEv}}\phantomsection\label{\detokenize{Classes_input:_CPPv2N9LAMMPS_NS5Input4fileEv}}\phantomsection\label{\detokenize{Classes_input:LAMMPS_NS::Input::file}}
\pysigstartsignatures
\pysigstartmultiline
\pysiglinewithargsret{\phantomsection\label{\detokenize{Classes_input:classLAMMPS__NS_1_1Input_1a5dafc7f590ea20432d3a8ca63e7ee96a}}\DUrole{kt}{void}\DUrole{w}{ }\sphinxbfcode{\sphinxupquote{\DUrole{n}{file}}}}{}{}
\pysigstopmultiline
\pysigstopsignatures
\sphinxAtStartPar
Process all input from the \sphinxcode{\sphinxupquote{FILE *}} pointer \sphinxstyleemphasis{infile}

\sphinxAtStartPar
This will read lines from \sphinxstyleemphasis{infile}, parse and execute them until the end
of the file is reached.  The \sphinxstyleemphasis{infile} pointer will usually point to
\sphinxcode{\sphinxupquote{stdin}} or the input file given with the \sphinxcode{\sphinxupquote{\sphinxhyphen{}in}} command line flag.

\end{fulllineitems}

\index{Input::file (C++ function)@\spxentry{Input::file}\spxextra{C++ function}}

\begin{fulllineitems}
\phantomsection\label{\detokenize{Classes_input:_CPPv4N9LAMMPS_NS5Input4fileEPKc}}\phantomsection\label{\detokenize{Classes_input:_CPPv3N9LAMMPS_NS5Input4fileEPKc}}\phantomsection\label{\detokenize{Classes_input:_CPPv2N9LAMMPS_NS5Input4fileEPKc}}\phantomsection\label{\detokenize{Classes_input:LAMMPS_NS::Input::file__cCP}}
\pysigstartsignatures
\pysigstartmultiline
\pysiglinewithargsret{\phantomsection\label{\detokenize{Classes_input:classLAMMPS__NS_1_1Input_1afd15ad1408b89c5f6df242228147b370}}\DUrole{kt}{void}\DUrole{w}{ }\sphinxbfcode{\sphinxupquote{\DUrole{n}{file}}}}{\DUrole{k}{const}\DUrole{w}{ }\DUrole{kt}{char}\DUrole{p}{*}}{}
\pysigstopmultiline
\pysigstopsignatures
\sphinxAtStartPar
Process all input from the file \sphinxstyleemphasis{filename}

\sphinxAtStartPar
This function opens the file at the path \sphinxstyleemphasis{filename}, puts the current
file pointer stored in \sphinxstyleemphasis{infile} on a stack and instead assigns \sphinxstyleemphasis{infile}
with the newly opened file pointer.  Then it will call the
{\hyperref[\detokenize{Classes_input:_CPPv4N9LAMMPS_NS5Input4fileEPKc}]{\sphinxcrossref{\sphinxcode{\sphinxupquote{Input::file()}}}}} function to read,
parse and execute the contents of that file.  When the end of the file
is reached, it is closed and the previous file pointer from the infile
file pointer stack restored to \sphinxstyleemphasis{infile}.
\begin{quote}\begin{description}
\sphinxlineitem{Parameters}
\sphinxAtStartPar
\sphinxstyleliteralstrong{\sphinxupquote{filename}} \textendash{} name of file with {\hyperref[\detokenize{Classes_lammps:classLAMMPS__NS_1_1LAMMPS}]{\sphinxcrossref{\DUrole{std,std-ref}{LAMMPS}}}} commands 

\end{description}\end{quote}

\end{fulllineitems}

\index{Input::one (C++ function)@\spxentry{Input::one}\spxextra{C++ function}}

\begin{fulllineitems}
\phantomsection\label{\detokenize{Classes_input:_CPPv4N9LAMMPS_NS5Input3oneERKNSt6stringE}}\phantomsection\label{\detokenize{Classes_input:_CPPv3N9LAMMPS_NS5Input3oneERKNSt6stringE}}\phantomsection\label{\detokenize{Classes_input:_CPPv2N9LAMMPS_NS5Input3oneERKNSt6stringE}}\phantomsection\label{\detokenize{Classes_input:LAMMPS_NS::Input::one__ssCR}}
\pysigstartsignatures
\pysigstartmultiline
\pysiglinewithargsret{\phantomsection\label{\detokenize{Classes_input:classLAMMPS__NS_1_1Input_1ad2ba7d137bc5ca20182ef7cd050940ed}}\DUrole{kt}{char}\DUrole{w}{ }\DUrole{p}{*}\sphinxbfcode{\sphinxupquote{\DUrole{n}{one}}}}{\DUrole{k}{const}\DUrole{w}{ }\DUrole{n}{std}\DUrole{p}{::}\DUrole{n}{string}\DUrole{p}{\&}}{}
\pysigstopmultiline
\pysigstopsignatures
\sphinxAtStartPar
Process a single command from a string in \sphinxstyleemphasis{single}

\sphinxAtStartPar
This function takes the text in \sphinxstyleemphasis{single}, makes a copy, parses that,
executes the command and returns the name of the command (without the
arguments).  If there was no command in \sphinxstyleemphasis{single} it will return
\sphinxcode{\sphinxupquote{nullptr}}.
\begin{quote}\begin{description}
\sphinxlineitem{Parameters}
\sphinxAtStartPar
\sphinxstyleliteralstrong{\sphinxupquote{single}} \textendash{} string with {\hyperref[\detokenize{Classes_lammps:classLAMMPS__NS_1_1LAMMPS}]{\sphinxcrossref{\DUrole{std,std-ref}{LAMMPS}}}} command 

\sphinxlineitem{Returns}
\sphinxAtStartPar
string with name of the parsed command w/o arguments 

\end{description}\end{quote}

\end{fulllineitems}


\end{sphinxuseclass}
\end{fulllineitems}


\sphinxstepscope


\section{Platform abstraction functions}
\label{\detokenize{Developer_platform:platform-abstraction-functions}}\label{\detokenize{Developer_platform::doc}}
\sphinxAtStartPar
The \sphinxcode{\sphinxupquote{platform}} sub\sphinxhyphen{}namespace inside the \sphinxcode{\sphinxupquote{LAMMPS\_NS}} namespace
provides a collection of wrapper and convenience functions and utilities
that perform common tasks for which platform specific code would be
required or for which a more high\sphinxhyphen{}level abstraction would be convenient
and reduce duplicated code.  This reduces redundant implementations and
encourages consistent behavior and thus has some overlap with the
{\hyperref[\detokenize{Developer_utils::doc}]{\sphinxcrossref{\DUrole{doc}{“utils” sub\sphinxhyphen{}namespace}}}}.


\subsection{Time functions}
\label{\detokenize{Developer_platform:time-functions}}\index{platform::cputime (C++ function)@\spxentry{platform::cputime}\spxextra{C++ function}}

\begin{fulllineitems}
\phantomsection\label{\detokenize{Developer_platform:_CPPv4N9LAMMPS_NS8platform7cputimeEv}}\phantomsection\label{\detokenize{Developer_platform:_CPPv3N9LAMMPS_NS8platform7cputimeEv}}\phantomsection\label{\detokenize{Developer_platform:_CPPv2N9LAMMPS_NS8platform7cputimeEv}}\phantomsection\label{\detokenize{Developer_platform:LAMMPS_NS::platform::cputime}}
\pysigstartsignatures
\pysigstartmultiline
\pysiglinewithargsret{\phantomsection\label{\detokenize{Developer_platform:platform_8h_1a5f4b3548be5eb5c6ecb125ef68070fba}}\DUrole{kt}{double}\DUrole{w}{ }\sphinxcode{\sphinxupquote{\DUrole{n}{LAMMPS\_NS}\DUrole{p}{::}\DUrole{n}{platform}\DUrole{p}{::}}}\sphinxbfcode{\sphinxupquote{\DUrole{n}{cputime}}}}{}{}
\pysigstopmultiline
\pysigstopsignatures
\sphinxAtStartPar
Return the consumed CPU time for the current process in seconds

\sphinxAtStartPar
This is a wrapper around the POSIX function getrusage() and its Windows equivalent. It is to be used in a similar fashion as MPI\_Wtime(). Its resolution may be rather low so it can only be trusted when observing processes consuming CPU time of at least a few seconds.
\begin{quote}\begin{description}
\sphinxlineitem{Returns}
\sphinxAtStartPar
used CPU time in seconds 

\end{description}\end{quote}

\end{fulllineitems}

\index{platform::walltime (C++ function)@\spxentry{platform::walltime}\spxextra{C++ function}}

\begin{fulllineitems}
\phantomsection\label{\detokenize{Developer_platform:_CPPv4N9LAMMPS_NS8platform8walltimeEv}}\phantomsection\label{\detokenize{Developer_platform:_CPPv3N9LAMMPS_NS8platform8walltimeEv}}\phantomsection\label{\detokenize{Developer_platform:_CPPv2N9LAMMPS_NS8platform8walltimeEv}}\phantomsection\label{\detokenize{Developer_platform:LAMMPS_NS::platform::walltime}}
\pysigstartsignatures
\pysigstartmultiline
\pysiglinewithargsret{\phantomsection\label{\detokenize{Developer_platform:platform_8h_1aca5f89d1d77efd166713bf8dd6d7fcd4}}\DUrole{kt}{double}\DUrole{w}{ }\sphinxcode{\sphinxupquote{\DUrole{n}{LAMMPS\_NS}\DUrole{p}{::}\DUrole{n}{platform}\DUrole{p}{::}}}\sphinxbfcode{\sphinxupquote{\DUrole{n}{walltime}}}}{}{}
\pysigstopmultiline
\pysigstopsignatures
\sphinxAtStartPar
Return the wall clock state for the current process in seconds

\sphinxAtStartPar
This clock is counting continuous time and is initialized during load of the executable/library. Its absolute value must be considered arbitrary and thus elapsed wall times are measured in taking differences. It is therefore to be used in a similar fashion as MPI\_Wtime() but has a different offset, usually leading to better resolution.
\begin{quote}\begin{description}
\sphinxlineitem{Returns}
\sphinxAtStartPar
wall clock time in seconds 

\end{description}\end{quote}

\end{fulllineitems}

\index{platform::usleep (C++ function)@\spxentry{platform::usleep}\spxextra{C++ function}}

\begin{fulllineitems}
\phantomsection\label{\detokenize{Developer_platform:_CPPv4N9LAMMPS_NS8platform6usleepEi}}\phantomsection\label{\detokenize{Developer_platform:_CPPv3N9LAMMPS_NS8platform6usleepEi}}\phantomsection\label{\detokenize{Developer_platform:_CPPv2N9LAMMPS_NS8platform6usleepEi}}\phantomsection\label{\detokenize{Developer_platform:LAMMPS_NS::platform::usleep__i}}
\pysigstartsignatures
\pysigstartmultiline
\pysiglinewithargsret{\phantomsection\label{\detokenize{Developer_platform:platform_8h_1afb043f14811cda3d34904e2eec5f07b9}}\DUrole{kt}{void}\DUrole{w}{ }\sphinxcode{\sphinxupquote{\DUrole{n}{LAMMPS\_NS}\DUrole{p}{::}\DUrole{n}{platform}\DUrole{p}{::}}}\sphinxbfcode{\sphinxupquote{\DUrole{n}{usleep}}}}{\DUrole{kt}{int}\DUrole{w}{ }\DUrole{n,sig-param}{usec}}{}
\pysigstopmultiline
\pysigstopsignatures
\sphinxAtStartPar
Suspend execution for a microsecond interval

\sphinxAtStartPar
This emulates the usleep(3) BSD function call also mentioned in POSIX.1\sphinxhyphen{}2001. This is not a precise delay; it may be longer, but not shorter.
\begin{quote}\begin{description}
\sphinxlineitem{Parameters}
\sphinxAtStartPar
\sphinxstyleliteralstrong{\sphinxupquote{usec}} \textendash{} length of delay in microseconds 

\end{description}\end{quote}

\end{fulllineitems}



\subsection{Platform information functions}
\label{\detokenize{Developer_platform:platform-information-functions}}\index{platform::os\_info (C++ function)@\spxentry{platform::os\_info}\spxextra{C++ function}}

\begin{fulllineitems}
\phantomsection\label{\detokenize{Developer_platform:_CPPv4N9LAMMPS_NS8platform7os_infoEv}}\phantomsection\label{\detokenize{Developer_platform:_CPPv3N9LAMMPS_NS8platform7os_infoEv}}\phantomsection\label{\detokenize{Developer_platform:_CPPv2N9LAMMPS_NS8platform7os_infoEv}}\phantomsection\label{\detokenize{Developer_platform:LAMMPS_NS::platform::os_info}}
\pysigstartsignatures
\pysigstartmultiline
\pysiglinewithargsret{\phantomsection\label{\detokenize{Developer_platform:platform_8h_1aaf549959ee090489b6146d8918e07d0c}}\DUrole{n}{std}\DUrole{p}{::}\DUrole{n}{string}\DUrole{w}{ }\sphinxcode{\sphinxupquote{\DUrole{n}{LAMMPS\_NS}\DUrole{p}{::}\DUrole{n}{platform}\DUrole{p}{::}}}\sphinxbfcode{\sphinxupquote{\DUrole{n}{os\_info}}}}{}{}
\pysigstopmultiline
\pysigstopsignatures
\sphinxAtStartPar
Return string with the operating system version and architecture info
\begin{quote}\begin{description}
\sphinxlineitem{Returns}
\sphinxAtStartPar
string with info about the OS and the platform is is running on 

\end{description}\end{quote}

\end{fulllineitems}

\index{platform::compiler\_info (C++ function)@\spxentry{platform::compiler\_info}\spxextra{C++ function}}

\begin{fulllineitems}
\phantomsection\label{\detokenize{Developer_platform:_CPPv4N9LAMMPS_NS8platform13compiler_infoEv}}\phantomsection\label{\detokenize{Developer_platform:_CPPv3N9LAMMPS_NS8platform13compiler_infoEv}}\phantomsection\label{\detokenize{Developer_platform:_CPPv2N9LAMMPS_NS8platform13compiler_infoEv}}\phantomsection\label{\detokenize{Developer_platform:LAMMPS_NS::platform::compiler_info}}
\pysigstartsignatures
\pysigstartmultiline
\pysiglinewithargsret{\phantomsection\label{\detokenize{Developer_platform:platform_8h_1aff1bead546363da736aa6e7d5a356c8f}}\DUrole{n}{std}\DUrole{p}{::}\DUrole{n}{string}\DUrole{w}{ }\sphinxcode{\sphinxupquote{\DUrole{n}{LAMMPS\_NS}\DUrole{p}{::}\DUrole{n}{platform}\DUrole{p}{::}}}\sphinxbfcode{\sphinxupquote{\DUrole{n}{compiler\_info}}}}{}{}
\pysigstopmultiline
\pysigstopsignatures
\sphinxAtStartPar
Return string with compiler version info

\sphinxAtStartPar
This function uses predefined compiler macros to identify Compilers and their version and configuration info.
\begin{quote}\begin{description}
\sphinxlineitem{Returns}
\sphinxAtStartPar
string with the compiler information text 

\end{description}\end{quote}

\end{fulllineitems}

\index{platform::cxx\_standard (C++ function)@\spxentry{platform::cxx\_standard}\spxextra{C++ function}}

\begin{fulllineitems}
\phantomsection\label{\detokenize{Developer_platform:_CPPv4N9LAMMPS_NS8platform12cxx_standardEv}}\phantomsection\label{\detokenize{Developer_platform:_CPPv3N9LAMMPS_NS8platform12cxx_standardEv}}\phantomsection\label{\detokenize{Developer_platform:_CPPv2N9LAMMPS_NS8platform12cxx_standardEv}}\phantomsection\label{\detokenize{Developer_platform:LAMMPS_NS::platform::cxx_standard}}
\pysigstartsignatures
\pysigstartmultiline
\pysiglinewithargsret{\phantomsection\label{\detokenize{Developer_platform:platform_8h_1a6b0cd92425b2ba864c0827aa84634092}}\DUrole{n}{std}\DUrole{p}{::}\DUrole{n}{string}\DUrole{w}{ }\sphinxcode{\sphinxupquote{\DUrole{n}{LAMMPS\_NS}\DUrole{p}{::}\DUrole{n}{platform}\DUrole{p}{::}}}\sphinxbfcode{\sphinxupquote{\DUrole{n}{cxx\_standard}}}}{}{}
\pysigstopmultiline
\pysigstopsignatures
\sphinxAtStartPar
Return string with C++ standard version used to compile {\hyperref[\detokenize{Classes_lammps:classLAMMPS__NS_1_1LAMMPS}]{\sphinxcrossref{\DUrole{std,std-ref}{LAMMPS}}}}.

\sphinxAtStartPar
This function uses predefined compiler macros to identify the C++ standard version used to compile {\hyperref[\detokenize{Classes_lammps:classLAMMPS__NS_1_1LAMMPS}]{\sphinxcrossref{\DUrole{std,std-ref}{LAMMPS}}}} with.
\begin{quote}\begin{description}
\sphinxlineitem{Returns}
\sphinxAtStartPar
string with the C++ standard version or “unknown” 

\end{description}\end{quote}

\end{fulllineitems}

\index{platform::openmp\_standard (C++ function)@\spxentry{platform::openmp\_standard}\spxextra{C++ function}}

\begin{fulllineitems}
\phantomsection\label{\detokenize{Developer_platform:_CPPv4N9LAMMPS_NS8platform15openmp_standardEv}}\phantomsection\label{\detokenize{Developer_platform:_CPPv3N9LAMMPS_NS8platform15openmp_standardEv}}\phantomsection\label{\detokenize{Developer_platform:_CPPv2N9LAMMPS_NS8platform15openmp_standardEv}}\phantomsection\label{\detokenize{Developer_platform:LAMMPS_NS::platform::openmp_standard}}
\pysigstartsignatures
\pysigstartmultiline
\pysiglinewithargsret{\phantomsection\label{\detokenize{Developer_platform:platform_8h_1a9a1ea03475feccb71ba24c81ed47bcbe}}\DUrole{n}{std}\DUrole{p}{::}\DUrole{n}{string}\DUrole{w}{ }\sphinxcode{\sphinxupquote{\DUrole{n}{LAMMPS\_NS}\DUrole{p}{::}\DUrole{n}{platform}\DUrole{p}{::}}}\sphinxbfcode{\sphinxupquote{\DUrole{n}{openmp\_standard}}}}{}{}
\pysigstopmultiline
\pysigstopsignatures
\sphinxAtStartPar
Return string with OpenMP standard version info

\sphinxAtStartPar
This function uses predefined compiler macros to identify OpenMP support and the supported version of the standard.
\begin{quote}\begin{description}
\sphinxlineitem{Returns}
\sphinxAtStartPar
string with the openmp information text 

\end{description}\end{quote}

\end{fulllineitems}

\index{platform::mpi\_vendor (C++ function)@\spxentry{platform::mpi\_vendor}\spxextra{C++ function}}

\begin{fulllineitems}
\phantomsection\label{\detokenize{Developer_platform:_CPPv4N9LAMMPS_NS8platform10mpi_vendorEv}}\phantomsection\label{\detokenize{Developer_platform:_CPPv3N9LAMMPS_NS8platform10mpi_vendorEv}}\phantomsection\label{\detokenize{Developer_platform:_CPPv2N9LAMMPS_NS8platform10mpi_vendorEv}}\phantomsection\label{\detokenize{Developer_platform:LAMMPS_NS::platform::mpi_vendor}}
\pysigstartsignatures
\pysigstartmultiline
\pysiglinewithargsret{\phantomsection\label{\detokenize{Developer_platform:platform_8h_1a91f4f7bc79930f0ced868e82f5a0c85f}}\DUrole{n}{std}\DUrole{p}{::}\DUrole{n}{string}\DUrole{w}{ }\sphinxcode{\sphinxupquote{\DUrole{n}{LAMMPS\_NS}\DUrole{p}{::}\DUrole{n}{platform}\DUrole{p}{::}}}\sphinxbfcode{\sphinxupquote{\DUrole{n}{mpi\_vendor}}}}{}{}
\pysigstopmultiline
\pysigstopsignatures
\sphinxAtStartPar
Return string with MPI vendor info

\sphinxAtStartPar
This function uses predefined macros to identify the vendor of the MPI library used.
\begin{quote}\begin{description}
\sphinxlineitem{Returns}
\sphinxAtStartPar
string with the MPI vendor information text 

\end{description}\end{quote}

\end{fulllineitems}

\index{platform::mpi\_info (C++ function)@\spxentry{platform::mpi\_info}\spxextra{C++ function}}

\begin{fulllineitems}
\phantomsection\label{\detokenize{Developer_platform:_CPPv4N9LAMMPS_NS8platform8mpi_infoERiRi}}\phantomsection\label{\detokenize{Developer_platform:_CPPv3N9LAMMPS_NS8platform8mpi_infoERiRi}}\phantomsection\label{\detokenize{Developer_platform:_CPPv2N9LAMMPS_NS8platform8mpi_infoERiRi}}\phantomsection\label{\detokenize{Developer_platform:LAMMPS_NS::platform::mpi_info__iR.iR}}
\pysigstartsignatures
\pysigstartmultiline
\pysiglinewithargsret{\phantomsection\label{\detokenize{Developer_platform:platform_8h_1a138dd5b5a0f070a142ae99957f7ba161}}\DUrole{n}{std}\DUrole{p}{::}\DUrole{n}{string}\DUrole{w}{ }\sphinxcode{\sphinxupquote{\DUrole{n}{LAMMPS\_NS}\DUrole{p}{::}\DUrole{n}{platform}\DUrole{p}{::}}}\sphinxbfcode{\sphinxupquote{\DUrole{n}{mpi\_info}}}}{\DUrole{kt}{int}\DUrole{w}{ }\DUrole{p}{\&}\DUrole{n,sig-param}{major}\sphinxparamcomma \DUrole{kt}{int}\DUrole{w}{ }\DUrole{p}{\&}\DUrole{n,sig-param}{minor}}{}
\pysigstopmultiline
\pysigstopsignatures
\sphinxAtStartPar
Return string with MPI version info

\sphinxAtStartPar
This function uses predefined macros and MPI function calls to identify the version of the MPI library used.
\begin{quote}\begin{description}
\sphinxlineitem{Parameters}\begin{itemize}
\item {} 
\sphinxAtStartPar
\sphinxstyleliteralstrong{\sphinxupquote{major}} \textendash{} major version of the MPI standard (set on exit) 

\item {} 
\sphinxAtStartPar
\sphinxstyleliteralstrong{\sphinxupquote{minor}} \textendash{} minor version of the MPI standard (set on exit) 

\end{itemize}

\sphinxlineitem{Returns}
\sphinxAtStartPar
string with the MPI version information text 

\end{description}\end{quote}

\end{fulllineitems}

\index{platform::compress\_info (C++ function)@\spxentry{platform::compress\_info}\spxextra{C++ function}}

\begin{fulllineitems}
\phantomsection\label{\detokenize{Developer_platform:_CPPv4N9LAMMPS_NS8platform13compress_infoEv}}\phantomsection\label{\detokenize{Developer_platform:_CPPv3N9LAMMPS_NS8platform13compress_infoEv}}\phantomsection\label{\detokenize{Developer_platform:_CPPv2N9LAMMPS_NS8platform13compress_infoEv}}\phantomsection\label{\detokenize{Developer_platform:LAMMPS_NS::platform::compress_info}}
\pysigstartsignatures
\pysigstartmultiline
\pysiglinewithargsret{\phantomsection\label{\detokenize{Developer_platform:platform_8h_1aec1a0eb6ddfe528d22869aa0aa1b1004}}\DUrole{n}{std}\DUrole{p}{::}\DUrole{n}{string}\DUrole{w}{ }\sphinxcode{\sphinxupquote{\DUrole{n}{LAMMPS\_NS}\DUrole{p}{::}\DUrole{n}{platform}\DUrole{p}{::}}}\sphinxbfcode{\sphinxupquote{\DUrole{n}{compress\_info}}}}{}{}
\pysigstopmultiline
\pysigstopsignatures
\sphinxAtStartPar
Return string with list of available compression types and executables

\sphinxAtStartPar
This function tests which of the supported compression executables are available for reading or writing compressed files where supported.
\begin{quote}\begin{description}
\sphinxlineitem{Returns}
\sphinxAtStartPar
string with list of available compression tools 

\end{description}\end{quote}

\end{fulllineitems}



\subsection{File and path functions and global constants}
\label{\detokenize{Developer_platform:file-and-path-functions-and-global-constants}}\index{platform::filepathsep (C++ member)@\spxentry{platform::filepathsep}\spxextra{C++ member}}

\begin{fulllineitems}
\phantomsection\label{\detokenize{Developer_platform:_CPPv4N9LAMMPS_NS8platform11filepathsepE}}\phantomsection\label{\detokenize{Developer_platform:_CPPv3N9LAMMPS_NS8platform11filepathsepE}}\phantomsection\label{\detokenize{Developer_platform:_CPPv2N9LAMMPS_NS8platform11filepathsepE}}\phantomsection\label{\detokenize{Developer_platform:LAMMPS_NS::platform::filepathsep__cA}}
\pysigstartsignatures
\pysigstartmultiline
\pysigline{\phantomsection\label{\detokenize{Developer_platform:platform_8h_1a3174b82f12dfbffab20b2b96fbc442f3}}\DUrole{k}{constexpr}\DUrole{w}{ }\DUrole{kt}{char}\DUrole{w}{ }\sphinxcode{\sphinxupquote{\DUrole{n}{LAMMPS\_NS}\DUrole{p}{::}\DUrole{n}{platform}\DUrole{p}{::}}}\sphinxbfcode{\sphinxupquote{\DUrole{n}{filepathsep}}}\DUrole{p}{{[}}\DUrole{p}{{]}}\DUrole{w}{ }\DUrole{p}{=}\DUrole{w}{ }\DUrole{s}{"/"}}
\pysigstopmultiline
\pysigstopsignatures
\sphinxAtStartPar
Platform specific file path component separator

\sphinxAtStartPar
This is a string with the character that separates directories and filename in paths on a platform. If multiple are characters are provided, the first is the preferred one. 

\end{fulllineitems}

\index{platform::pathvarsep (C++ member)@\spxentry{platform::pathvarsep}\spxextra{C++ member}}

\begin{fulllineitems}
\phantomsection\label{\detokenize{Developer_platform:_CPPv4N9LAMMPS_NS8platform10pathvarsepE}}\phantomsection\label{\detokenize{Developer_platform:_CPPv3N9LAMMPS_NS8platform10pathvarsepE}}\phantomsection\label{\detokenize{Developer_platform:_CPPv2N9LAMMPS_NS8platform10pathvarsepE}}\phantomsection\label{\detokenize{Developer_platform:LAMMPS_NS::platform::pathvarsep__c}}
\pysigstartsignatures
\pysigstartmultiline
\pysigline{\phantomsection\label{\detokenize{Developer_platform:platform_8h_1a0427d2d3074fc8188d6731545bc33855}}\DUrole{k}{constexpr}\DUrole{w}{ }\DUrole{kt}{char}\DUrole{w}{ }\sphinxcode{\sphinxupquote{\DUrole{n}{LAMMPS\_NS}\DUrole{p}{::}\DUrole{n}{platform}\DUrole{p}{::}}}\sphinxbfcode{\sphinxupquote{\DUrole{n}{pathvarsep}}}\DUrole{w}{ }\DUrole{p}{=}\DUrole{w}{ }\DUrole{sc}{\textquotesingle{}:\textquotesingle{}}}
\pysigstopmultiline
\pysigstopsignatures
\sphinxAtStartPar
Platform specific path environment variable component separator

\sphinxAtStartPar
This is the character that separates entries in “PATH”\sphinxhyphen{}style environment variables. 

\end{fulllineitems}

\index{platform::guesspath (C++ function)@\spxentry{platform::guesspath}\spxextra{C++ function}}

\begin{fulllineitems}
\phantomsection\label{\detokenize{Developer_platform:_CPPv4N9LAMMPS_NS8platform9guesspathEP4FILEPci}}\phantomsection\label{\detokenize{Developer_platform:_CPPv3N9LAMMPS_NS8platform9guesspathEP4FILEPci}}\phantomsection\label{\detokenize{Developer_platform:_CPPv2N9LAMMPS_NS8platform9guesspathEP4FILEPci}}\phantomsection\label{\detokenize{Developer_platform:LAMMPS_NS::platform::guesspath__FILEP.cP.i}}
\pysigstartsignatures
\pysigstartmultiline
\pysiglinewithargsret{\phantomsection\label{\detokenize{Developer_platform:platform_8h_1aa344b662f8d3c394089253824eca591d}}\DUrole{k}{const}\DUrole{w}{ }\DUrole{kt}{char}\DUrole{w}{ }\DUrole{p}{*}\sphinxcode{\sphinxupquote{\DUrole{n}{LAMMPS\_NS}\DUrole{p}{::}\DUrole{n}{platform}\DUrole{p}{::}}}\sphinxbfcode{\sphinxupquote{\DUrole{n}{guesspath}}}}{\DUrole{n}{FILE}\DUrole{w}{ }\DUrole{p}{*}\DUrole{n,sig-param}{fp}\sphinxparamcomma \DUrole{kt}{char}\DUrole{w}{ }\DUrole{p}{*}\DUrole{n,sig-param}{buf}\sphinxparamcomma \DUrole{kt}{int}\DUrole{w}{ }\DUrole{n,sig-param}{len}}{}
\pysigstopmultiline
\pysigstopsignatures
\sphinxAtStartPar
Try to detect pathname from FILE pointer

\sphinxAtStartPar
Currently only supported on Linux, MacOS, and Windows. Otherwise will report “(unknown)”.

\sphinxAtStartPar

On Linux the folder /proc/self/fd holds symbolic links to the actual pathnames associated with each open file descriptor of the current process. On MacOS the same kind of information can be obtained using \sphinxcode{\sphinxupquote{fcntl(fd,F\_GETPATH,buf)}}. On Windows we use \sphinxcode{\sphinxupquote{GetFinalPathNameByHandleA()}} which is available with Windows Vista and later. If the buffer is too small (\textless{} 16 bytes) a null pointer is returned.

\sphinxAtStartPar
This function is used to provide a filename with error messages in functions where the filename is not passed as an argument, but the FILE * pointer. 
\begin{quote}\begin{description}
\sphinxlineitem{Parameters}\begin{itemize}
\item {} 
\sphinxAtStartPar
\sphinxstyleliteralstrong{\sphinxupquote{fp}} \textendash{} FILE pointer struct from STDIO library for which we want to detect the name 

\item {} 
\sphinxAtStartPar
\sphinxstyleliteralstrong{\sphinxupquote{buf}} \textendash{} storage buffer for pathname. output will be truncated if not large enough 

\item {} 
\sphinxAtStartPar
\sphinxstyleliteralstrong{\sphinxupquote{len}} \textendash{} size of storage buffer. output will be truncated to this length \sphinxhyphen{} 1 

\end{itemize}

\sphinxlineitem{Returns}
\sphinxAtStartPar
pointer to the storage buffer with path or a NULL pointer if buf is invalid or the buffer size is too small

\end{description}\end{quote}

\end{fulllineitems}

\index{platform::path\_basename (C++ function)@\spxentry{platform::path\_basename}\spxextra{C++ function}}

\begin{fulllineitems}
\phantomsection\label{\detokenize{Developer_platform:_CPPv4N9LAMMPS_NS8platform13path_basenameERKNSt6stringE}}\phantomsection\label{\detokenize{Developer_platform:_CPPv3N9LAMMPS_NS8platform13path_basenameERKNSt6stringE}}\phantomsection\label{\detokenize{Developer_platform:_CPPv2N9LAMMPS_NS8platform13path_basenameERKNSt6stringE}}\phantomsection\label{\detokenize{Developer_platform:LAMMPS_NS::platform::path_basename__ssCR}}
\pysigstartsignatures
\pysigstartmultiline
\pysiglinewithargsret{\phantomsection\label{\detokenize{Developer_platform:platform_8h_1af571f031125be18774c73496bf428f1e}}\DUrole{n}{std}\DUrole{p}{::}\DUrole{n}{string}\DUrole{w}{ }\sphinxcode{\sphinxupquote{\DUrole{n}{LAMMPS\_NS}\DUrole{p}{::}\DUrole{n}{platform}\DUrole{p}{::}}}\sphinxbfcode{\sphinxupquote{\DUrole{n}{path\_basename}}}}{\DUrole{k}{const}\DUrole{w}{ }\DUrole{n}{std}\DUrole{p}{::}\DUrole{n}{string}\DUrole{w}{ }\DUrole{p}{\&}\DUrole{n,sig-param}{path}}{}
\pysigstopmultiline
\pysigstopsignatures
\sphinxAtStartPar
Strip off leading part of path, return just the filename
\begin{quote}\begin{description}
\sphinxlineitem{Parameters}
\sphinxAtStartPar
\sphinxstyleliteralstrong{\sphinxupquote{path}} \textendash{} file path 

\sphinxlineitem{Returns}
\sphinxAtStartPar
file name 

\end{description}\end{quote}

\end{fulllineitems}

\index{platform::path\_join (C++ function)@\spxentry{platform::path\_join}\spxextra{C++ function}}

\begin{fulllineitems}
\phantomsection\label{\detokenize{Developer_platform:_CPPv4N9LAMMPS_NS8platform9path_joinERKNSt6stringERKNSt6stringE}}\phantomsection\label{\detokenize{Developer_platform:_CPPv3N9LAMMPS_NS8platform9path_joinERKNSt6stringERKNSt6stringE}}\phantomsection\label{\detokenize{Developer_platform:_CPPv2N9LAMMPS_NS8platform9path_joinERKNSt6stringERKNSt6stringE}}\phantomsection\label{\detokenize{Developer_platform:LAMMPS_NS::platform::path_join__ssCR.ssCR}}
\pysigstartsignatures
\pysigstartmultiline
\pysiglinewithargsret{\phantomsection\label{\detokenize{Developer_platform:platform_8h_1afdc0fd0cefb9e5241e4ba20852877f4e}}\DUrole{n}{std}\DUrole{p}{::}\DUrole{n}{string}\DUrole{w}{ }\sphinxcode{\sphinxupquote{\DUrole{n}{LAMMPS\_NS}\DUrole{p}{::}\DUrole{n}{platform}\DUrole{p}{::}}}\sphinxbfcode{\sphinxupquote{\DUrole{n}{path\_join}}}}{\DUrole{k}{const}\DUrole{w}{ }\DUrole{n}{std}\DUrole{p}{::}\DUrole{n}{string}\DUrole{w}{ }\DUrole{p}{\&}\DUrole{n,sig-param}{a}\sphinxparamcomma \DUrole{k}{const}\DUrole{w}{ }\DUrole{n}{std}\DUrole{p}{::}\DUrole{n}{string}\DUrole{w}{ }\DUrole{p}{\&}\DUrole{n,sig-param}{b}}{}
\pysigstopmultiline
\pysigstopsignatures
\sphinxAtStartPar
Join two pathname segments

\sphinxAtStartPar
This uses the forward slash ‘/’ character unless {\hyperref[\detokenize{Classes_lammps:classLAMMPS__NS_1_1LAMMPS}]{\sphinxcrossref{\DUrole{std,std-ref}{LAMMPS}}}} is compiled for Windows where it uses the backward slash ‘\textquotesingle{}
\begin{quote}\begin{description}
\sphinxlineitem{Parameters}\begin{itemize}
\item {} 
\sphinxAtStartPar
\sphinxstyleliteralstrong{\sphinxupquote{a}} \textendash{} first path 

\item {} 
\sphinxAtStartPar
\sphinxstyleliteralstrong{\sphinxupquote{b}} \textendash{} second path 

\end{itemize}

\sphinxlineitem{Returns}
\sphinxAtStartPar
combined path 

\end{description}\end{quote}

\end{fulllineitems}

\index{platform::file\_is\_readable (C++ function)@\spxentry{platform::file\_is\_readable}\spxextra{C++ function}}

\begin{fulllineitems}
\phantomsection\label{\detokenize{Developer_platform:_CPPv4N9LAMMPS_NS8platform16file_is_readableERKNSt6stringE}}\phantomsection\label{\detokenize{Developer_platform:_CPPv3N9LAMMPS_NS8platform16file_is_readableERKNSt6stringE}}\phantomsection\label{\detokenize{Developer_platform:_CPPv2N9LAMMPS_NS8platform16file_is_readableERKNSt6stringE}}\phantomsection\label{\detokenize{Developer_platform:LAMMPS_NS::platform::file_is_readable__ssCR}}
\pysigstartsignatures
\pysigstartmultiline
\pysiglinewithargsret{\phantomsection\label{\detokenize{Developer_platform:platform_8h_1a9b6a15dfd00fe9eebeb84f092402b6b1}}\DUrole{kt}{bool}\DUrole{w}{ }\sphinxcode{\sphinxupquote{\DUrole{n}{LAMMPS\_NS}\DUrole{p}{::}\DUrole{n}{platform}\DUrole{p}{::}}}\sphinxbfcode{\sphinxupquote{\DUrole{n}{file\_is\_readable}}}}{\DUrole{k}{const}\DUrole{w}{ }\DUrole{n}{std}\DUrole{p}{::}\DUrole{n}{string}\DUrole{w}{ }\DUrole{p}{\&}\DUrole{n,sig-param}{path}}{}
\pysigstopmultiline
\pysigstopsignatures
\sphinxAtStartPar
Check if file exists and is readable
\begin{quote}\begin{description}
\sphinxlineitem{Parameters}
\sphinxAtStartPar
\sphinxstyleliteralstrong{\sphinxupquote{path}} \textendash{} file path 

\sphinxlineitem{Returns}
\sphinxAtStartPar
true if file exists and is readable 

\end{description}\end{quote}

\end{fulllineitems}

\index{platform::is\_console (C++ function)@\spxentry{platform::is\_console}\spxextra{C++ function}}

\begin{fulllineitems}
\phantomsection\label{\detokenize{Developer_platform:_CPPv4N9LAMMPS_NS8platform10is_consoleEP4FILE}}\phantomsection\label{\detokenize{Developer_platform:_CPPv3N9LAMMPS_NS8platform10is_consoleEP4FILE}}\phantomsection\label{\detokenize{Developer_platform:_CPPv2N9LAMMPS_NS8platform10is_consoleEP4FILE}}\phantomsection\label{\detokenize{Developer_platform:LAMMPS_NS::platform::is_console__FILEP}}
\pysigstartsignatures
\pysigstartmultiline
\pysiglinewithargsret{\phantomsection\label{\detokenize{Developer_platform:platform_8h_1ad3dce4729b475039e1f3cd7ad5dc4cf5}}\DUrole{kt}{bool}\DUrole{w}{ }\sphinxcode{\sphinxupquote{\DUrole{n}{LAMMPS\_NS}\DUrole{p}{::}\DUrole{n}{platform}\DUrole{p}{::}}}\sphinxbfcode{\sphinxupquote{\DUrole{n}{is\_console}}}}{\DUrole{n}{FILE}\DUrole{w}{ }\DUrole{p}{*}\DUrole{n,sig-param}{fp}}{}
\pysigstopmultiline
\pysigstopsignatures
\sphinxAtStartPar
Check if a file pointer may be connected to a console
\begin{quote}\begin{description}
\sphinxlineitem{Parameters}
\sphinxAtStartPar
\sphinxstyleliteralstrong{\sphinxupquote{fp}} \textendash{} file pointer 

\sphinxlineitem{Returns}
\sphinxAtStartPar
true if the file pointer is flagged as a TTY 

\end{description}\end{quote}

\end{fulllineitems}

\index{platform::disk\_free (C++ function)@\spxentry{platform::disk\_free}\spxextra{C++ function}}

\begin{fulllineitems}
\phantomsection\label{\detokenize{Developer_platform:_CPPv4N9LAMMPS_NS8platform9disk_freeERKNSt6stringE}}\phantomsection\label{\detokenize{Developer_platform:_CPPv3N9LAMMPS_NS8platform9disk_freeERKNSt6stringE}}\phantomsection\label{\detokenize{Developer_platform:_CPPv2N9LAMMPS_NS8platform9disk_freeERKNSt6stringE}}\phantomsection\label{\detokenize{Developer_platform:LAMMPS_NS::platform::disk_free__ssCR}}
\pysigstartsignatures
\pysigstartmultiline
\pysiglinewithargsret{\phantomsection\label{\detokenize{Developer_platform:platform_8h_1a6ed051736e3bf0c472eacaa8d7ac86e6}}\DUrole{kt}{double}\DUrole{w}{ }\sphinxcode{\sphinxupquote{\DUrole{n}{LAMMPS\_NS}\DUrole{p}{::}\DUrole{n}{platform}\DUrole{p}{::}}}\sphinxbfcode{\sphinxupquote{\DUrole{n}{disk\_free}}}}{\DUrole{k}{const}\DUrole{w}{ }\DUrole{n}{std}\DUrole{p}{::}\DUrole{n}{string}\DUrole{w}{ }\DUrole{p}{\&}\DUrole{n,sig-param}{path}}{}
\pysigstopmultiline
\pysigstopsignatures
\sphinxAtStartPar
Return free disk space in bytes of file system pointed to by path

\sphinxAtStartPar
Returns \sphinxhyphen{}1.0 if the path is invalid or free space reporting not supported.
\begin{quote}\begin{description}
\sphinxlineitem{Parameters}
\sphinxAtStartPar
\sphinxstyleliteralstrong{\sphinxupquote{path}} \textendash{} file or folder path in file system 

\sphinxlineitem{Returns}
\sphinxAtStartPar


\end{description}\end{quote}

\end{fulllineitems}

\index{platform::path\_is\_directory (C++ function)@\spxentry{platform::path\_is\_directory}\spxextra{C++ function}}

\begin{fulllineitems}
\phantomsection\label{\detokenize{Developer_platform:_CPPv4N9LAMMPS_NS8platform17path_is_directoryERKNSt6stringE}}\phantomsection\label{\detokenize{Developer_platform:_CPPv3N9LAMMPS_NS8platform17path_is_directoryERKNSt6stringE}}\phantomsection\label{\detokenize{Developer_platform:_CPPv2N9LAMMPS_NS8platform17path_is_directoryERKNSt6stringE}}\phantomsection\label{\detokenize{Developer_platform:LAMMPS_NS::platform::path_is_directory__ssCR}}
\pysigstartsignatures
\pysigstartmultiline
\pysiglinewithargsret{\phantomsection\label{\detokenize{Developer_platform:platform_8h_1a9da6d417e9c497085bd74752bf541679}}\DUrole{kt}{bool}\DUrole{w}{ }\sphinxcode{\sphinxupquote{\DUrole{n}{LAMMPS\_NS}\DUrole{p}{::}\DUrole{n}{platform}\DUrole{p}{::}}}\sphinxbfcode{\sphinxupquote{\DUrole{n}{path\_is\_directory}}}}{\DUrole{k}{const}\DUrole{w}{ }\DUrole{n}{std}\DUrole{p}{::}\DUrole{n}{string}\DUrole{w}{ }\DUrole{p}{\&}\DUrole{n,sig-param}{path}}{}
\pysigstopmultiline
\pysigstopsignatures
\sphinxAtStartPar
Check if a path is a directory
\begin{quote}\begin{description}
\sphinxlineitem{Parameters}
\sphinxAtStartPar
\sphinxstyleliteralstrong{\sphinxupquote{path}} \textendash{} directory path 

\sphinxlineitem{Returns}
\sphinxAtStartPar
true if the directory exists 

\end{description}\end{quote}

\end{fulllineitems}

\index{platform::current\_directory (C++ function)@\spxentry{platform::current\_directory}\spxextra{C++ function}}

\begin{fulllineitems}
\phantomsection\label{\detokenize{Developer_platform:_CPPv4N9LAMMPS_NS8platform17current_directoryEv}}\phantomsection\label{\detokenize{Developer_platform:_CPPv3N9LAMMPS_NS8platform17current_directoryEv}}\phantomsection\label{\detokenize{Developer_platform:_CPPv2N9LAMMPS_NS8platform17current_directoryEv}}\phantomsection\label{\detokenize{Developer_platform:LAMMPS_NS::platform::current_directory}}
\pysigstartsignatures
\pysigstartmultiline
\pysiglinewithargsret{\phantomsection\label{\detokenize{Developer_platform:platform_8h_1a22659a82b5a27cf76e22d02759647f2a}}\DUrole{n}{std}\DUrole{p}{::}\DUrole{n}{string}\DUrole{w}{ }\sphinxcode{\sphinxupquote{\DUrole{n}{LAMMPS\_NS}\DUrole{p}{::}\DUrole{n}{platform}\DUrole{p}{::}}}\sphinxbfcode{\sphinxupquote{\DUrole{n}{current\_directory}}}}{}{}
\pysigstopmultiline
\pysigstopsignatures
\sphinxAtStartPar
Get string with path to the current directory
\begin{quote}\begin{description}
\sphinxlineitem{Returns}
\sphinxAtStartPar
path to the current directory or empty string 

\end{description}\end{quote}

\end{fulllineitems}

\index{platform::list\_directory (C++ function)@\spxentry{platform::list\_directory}\spxextra{C++ function}}

\begin{fulllineitems}
\phantomsection\label{\detokenize{Developer_platform:_CPPv4N9LAMMPS_NS8platform14list_directoryERKNSt6stringE}}\phantomsection\label{\detokenize{Developer_platform:_CPPv3N9LAMMPS_NS8platform14list_directoryERKNSt6stringE}}\phantomsection\label{\detokenize{Developer_platform:_CPPv2N9LAMMPS_NS8platform14list_directoryERKNSt6stringE}}\phantomsection\label{\detokenize{Developer_platform:LAMMPS_NS::platform::list_directory__ssCR}}
\pysigstartsignatures
\pysigstartmultiline
\pysiglinewithargsret{\phantomsection\label{\detokenize{Developer_platform:platform_8h_1ad9f70de7fed649daa7ef1957c9997e38}}\DUrole{n}{std}\DUrole{p}{::}\DUrole{n}{vector}\DUrole{p}{\textless{}}\DUrole{n}{std}\DUrole{p}{::}\DUrole{n}{string}\DUrole{p}{\textgreater{}}\DUrole{w}{ }\sphinxcode{\sphinxupquote{\DUrole{n}{LAMMPS\_NS}\DUrole{p}{::}\DUrole{n}{platform}\DUrole{p}{::}}}\sphinxbfcode{\sphinxupquote{\DUrole{n}{list\_directory}}}}{\DUrole{k}{const}\DUrole{w}{ }\DUrole{n}{std}\DUrole{p}{::}\DUrole{n}{string}\DUrole{w}{ }\DUrole{p}{\&}\DUrole{n,sig-param}{dir}}{}
\pysigstopmultiline
\pysigstopsignatures
\sphinxAtStartPar
Get list of entries in a directory

\sphinxAtStartPar
This provides a list of strings of the entries in the directory without the leading path name while also skipping over “..” and “.”.
\begin{quote}\begin{description}
\sphinxlineitem{Parameters}
\sphinxAtStartPar
\sphinxstyleliteralstrong{\sphinxupquote{dir}} \textendash{} path to directory 

\sphinxlineitem{Returns}
\sphinxAtStartPar
vector with strings of all directory entries 

\end{description}\end{quote}

\end{fulllineitems}

\index{platform::chdir (C++ function)@\spxentry{platform::chdir}\spxextra{C++ function}}

\begin{fulllineitems}
\phantomsection\label{\detokenize{Developer_platform:_CPPv4N9LAMMPS_NS8platform5chdirERKNSt6stringE}}\phantomsection\label{\detokenize{Developer_platform:_CPPv3N9LAMMPS_NS8platform5chdirERKNSt6stringE}}\phantomsection\label{\detokenize{Developer_platform:_CPPv2N9LAMMPS_NS8platform5chdirERKNSt6stringE}}\phantomsection\label{\detokenize{Developer_platform:LAMMPS_NS::platform::chdir__ssCR}}
\pysigstartsignatures
\pysigstartmultiline
\pysiglinewithargsret{\phantomsection\label{\detokenize{Developer_platform:platform_8h_1a8dc97878b620dd294214f33d426b25dd}}\DUrole{kt}{int}\DUrole{w}{ }\sphinxcode{\sphinxupquote{\DUrole{n}{LAMMPS\_NS}\DUrole{p}{::}\DUrole{n}{platform}\DUrole{p}{::}}}\sphinxbfcode{\sphinxupquote{\DUrole{n}{chdir}}}}{\DUrole{k}{const}\DUrole{w}{ }\DUrole{n}{std}\DUrole{p}{::}\DUrole{n}{string}\DUrole{w}{ }\DUrole{p}{\&}\DUrole{n,sig-param}{path}}{}
\pysigstopmultiline
\pysigstopsignatures
\sphinxAtStartPar
Change current directory
\begin{quote}\begin{description}
\sphinxlineitem{Parameters}
\sphinxAtStartPar
\sphinxstyleliteralstrong{\sphinxupquote{path}} \textendash{} new current working directory path 

\sphinxlineitem{Returns}
\sphinxAtStartPar
\sphinxhyphen{}1 if unsuccessful, otherwise \textgreater{}= 0 

\end{description}\end{quote}

\end{fulllineitems}

\index{platform::mkdir (C++ function)@\spxentry{platform::mkdir}\spxextra{C++ function}}

\begin{fulllineitems}
\phantomsection\label{\detokenize{Developer_platform:_CPPv4N9LAMMPS_NS8platform5mkdirERKNSt6stringE}}\phantomsection\label{\detokenize{Developer_platform:_CPPv3N9LAMMPS_NS8platform5mkdirERKNSt6stringE}}\phantomsection\label{\detokenize{Developer_platform:_CPPv2N9LAMMPS_NS8platform5mkdirERKNSt6stringE}}\phantomsection\label{\detokenize{Developer_platform:LAMMPS_NS::platform::mkdir__ssCR}}
\pysigstartsignatures
\pysigstartmultiline
\pysiglinewithargsret{\phantomsection\label{\detokenize{Developer_platform:platform_8h_1aeacc2f1b5e3104760c86fb4644cb684a}}\DUrole{kt}{int}\DUrole{w}{ }\sphinxcode{\sphinxupquote{\DUrole{n}{LAMMPS\_NS}\DUrole{p}{::}\DUrole{n}{platform}\DUrole{p}{::}}}\sphinxbfcode{\sphinxupquote{\DUrole{n}{mkdir}}}}{\DUrole{k}{const}\DUrole{w}{ }\DUrole{n}{std}\DUrole{p}{::}\DUrole{n}{string}\DUrole{w}{ }\DUrole{p}{\&}\DUrole{n,sig-param}{path}}{}
\pysigstopmultiline
\pysigstopsignatures
\sphinxAtStartPar
Create a directory

\sphinxAtStartPar
Unlike the the \sphinxcode{\sphinxupquote{mkdir()}} or \sphinxcode{\sphinxupquote{\_mkdir()}} functions of the C library, this function will also try to create non\sphinxhyphen{}existing sub\sphinxhyphen{}directories in case they don’t exist, and thus behaves like the \sphinxcode{\sphinxupquote{mkdir \sphinxhyphen{}p}} command rather than plain \sphinxcode{\sphinxupquote{mkdir}} or \textasciigrave{}\sphinxcode{\sphinxupquote{md}}.
\begin{quote}\begin{description}
\sphinxlineitem{Parameters}
\sphinxAtStartPar
\sphinxstyleliteralstrong{\sphinxupquote{path}} \textendash{} directory path 

\sphinxlineitem{Returns}
\sphinxAtStartPar
\sphinxhyphen{}1 if unsuccessful, otherwise \textgreater{}= 0 

\end{description}\end{quote}

\end{fulllineitems}

\index{platform::rmdir (C++ function)@\spxentry{platform::rmdir}\spxextra{C++ function}}

\begin{fulllineitems}
\phantomsection\label{\detokenize{Developer_platform:_CPPv4N9LAMMPS_NS8platform5rmdirERKNSt6stringE}}\phantomsection\label{\detokenize{Developer_platform:_CPPv3N9LAMMPS_NS8platform5rmdirERKNSt6stringE}}\phantomsection\label{\detokenize{Developer_platform:_CPPv2N9LAMMPS_NS8platform5rmdirERKNSt6stringE}}\phantomsection\label{\detokenize{Developer_platform:LAMMPS_NS::platform::rmdir__ssCR}}
\pysigstartsignatures
\pysigstartmultiline
\pysiglinewithargsret{\phantomsection\label{\detokenize{Developer_platform:platform_8h_1a3302c183fe617ad6bf01e5f1bfff8dd9}}\DUrole{kt}{int}\DUrole{w}{ }\sphinxcode{\sphinxupquote{\DUrole{n}{LAMMPS\_NS}\DUrole{p}{::}\DUrole{n}{platform}\DUrole{p}{::}}}\sphinxbfcode{\sphinxupquote{\DUrole{n}{rmdir}}}}{\DUrole{k}{const}\DUrole{w}{ }\DUrole{n}{std}\DUrole{p}{::}\DUrole{n}{string}\DUrole{w}{ }\DUrole{p}{\&}\DUrole{n,sig-param}{path}}{}
\pysigstopmultiline
\pysigstopsignatures
\sphinxAtStartPar
Delete a directory

\sphinxAtStartPar
Unlike the the \sphinxcode{\sphinxupquote{rmdir()}} or \sphinxcode{\sphinxupquote{\_rmdir()}} functions of the C library, this function will check for the contents of the folder and recurse into any sub\sphinxhyphen{}folders, if necessary, and delete all contained folders and their contents before deleting the folder \sphinxstyleemphasis{path}.
\begin{quote}\begin{description}
\sphinxlineitem{Parameters}
\sphinxAtStartPar
\sphinxstyleliteralstrong{\sphinxupquote{path}} \textendash{} directory path 

\sphinxlineitem{Returns}
\sphinxAtStartPar
\sphinxhyphen{}1 if unsuccessful, otherwise \textgreater{}= 0 

\end{description}\end{quote}

\end{fulllineitems}

\index{platform::unlink (C++ function)@\spxentry{platform::unlink}\spxextra{C++ function}}

\begin{fulllineitems}
\phantomsection\label{\detokenize{Developer_platform:_CPPv4N9LAMMPS_NS8platform6unlinkERKNSt6stringE}}\phantomsection\label{\detokenize{Developer_platform:_CPPv3N9LAMMPS_NS8platform6unlinkERKNSt6stringE}}\phantomsection\label{\detokenize{Developer_platform:_CPPv2N9LAMMPS_NS8platform6unlinkERKNSt6stringE}}\phantomsection\label{\detokenize{Developer_platform:LAMMPS_NS::platform::unlink__ssCR}}
\pysigstartsignatures
\pysigstartmultiline
\pysiglinewithargsret{\phantomsection\label{\detokenize{Developer_platform:platform_8h_1a6555cc4e8063ed94d56cd815278a9421}}\DUrole{kt}{int}\DUrole{w}{ }\sphinxcode{\sphinxupquote{\DUrole{n}{LAMMPS\_NS}\DUrole{p}{::}\DUrole{n}{platform}\DUrole{p}{::}}}\sphinxbfcode{\sphinxupquote{\DUrole{n}{unlink}}}}{\DUrole{k}{const}\DUrole{w}{ }\DUrole{n}{std}\DUrole{p}{::}\DUrole{n}{string}\DUrole{w}{ }\DUrole{p}{\&}\DUrole{n,sig-param}{path}}{}
\pysigstopmultiline
\pysigstopsignatures
\sphinxAtStartPar
Delete a file
\begin{quote}\begin{description}
\sphinxlineitem{Parameters}
\sphinxAtStartPar
\sphinxstyleliteralstrong{\sphinxupquote{path}} \textendash{} path to file to be deleted 

\sphinxlineitem{Returns}
\sphinxAtStartPar
0 on success, \sphinxhyphen{}1 on error 

\end{description}\end{quote}

\end{fulllineitems}



\subsection{Standard I/O function wrappers}
\label{\detokenize{Developer_platform:standard-i-o-function-wrappers}}\index{platform::END\_OF\_FILE (C++ member)@\spxentry{platform::END\_OF\_FILE}\spxextra{C++ member}}

\begin{fulllineitems}
\phantomsection\label{\detokenize{Developer_platform:_CPPv4N9LAMMPS_NS8platform11END_OF_FILEE}}\phantomsection\label{\detokenize{Developer_platform:_CPPv3N9LAMMPS_NS8platform11END_OF_FILEE}}\phantomsection\label{\detokenize{Developer_platform:_CPPv2N9LAMMPS_NS8platform11END_OF_FILEE}}\phantomsection\label{\detokenize{Developer_platform:LAMMPS_NS::platform::END_OF_FILE__bigint}}
\pysigstartsignatures
\pysigstartmultiline
\pysigline{\phantomsection\label{\detokenize{Developer_platform:platform_8h_1aff22a2f6cfc75229f3c6691a64679ab0}}\DUrole{k}{constexpr}\DUrole{w}{ }\DUrole{n}{bigint}\DUrole{w}{ }\sphinxcode{\sphinxupquote{\DUrole{n}{LAMMPS\_NS}\DUrole{p}{::}\DUrole{n}{platform}\DUrole{p}{::}}}\sphinxbfcode{\sphinxupquote{\DUrole{n}{END\_OF\_FILE}}}\DUrole{w}{ }\DUrole{p}{=}\DUrole{w}{ }\DUrole{o}{\sphinxhyphen{}}\DUrole{m}{1}}
\pysigstopmultiline
\pysigstopsignatures
\sphinxAtStartPar
constant to seek to the end of the file 

\end{fulllineitems}

\index{platform::ftell (C++ function)@\spxentry{platform::ftell}\spxextra{C++ function}}

\begin{fulllineitems}
\phantomsection\label{\detokenize{Developer_platform:_CPPv4N9LAMMPS_NS8platform5ftellEP4FILE}}\phantomsection\label{\detokenize{Developer_platform:_CPPv3N9LAMMPS_NS8platform5ftellEP4FILE}}\phantomsection\label{\detokenize{Developer_platform:_CPPv2N9LAMMPS_NS8platform5ftellEP4FILE}}\phantomsection\label{\detokenize{Developer_platform:LAMMPS_NS::platform::ftell__FILEP}}
\pysigstartsignatures
\pysigstartmultiline
\pysiglinewithargsret{\phantomsection\label{\detokenize{Developer_platform:platform_8h_1afd4bf7f50771420b35850e961c8fa487}}\DUrole{n}{bigint}\DUrole{w}{ }\sphinxcode{\sphinxupquote{\DUrole{n}{LAMMPS\_NS}\DUrole{p}{::}\DUrole{n}{platform}\DUrole{p}{::}}}\sphinxbfcode{\sphinxupquote{\DUrole{n}{ftell}}}}{\DUrole{n}{FILE}\DUrole{w}{ }\DUrole{p}{*}\DUrole{n,sig-param}{fp}}{}
\pysigstopmultiline
\pysigstopsignatures
\sphinxAtStartPar
Get current file position
\begin{quote}\begin{description}
\sphinxlineitem{Parameters}
\sphinxAtStartPar
\sphinxstyleliteralstrong{\sphinxupquote{fp}} \textendash{} FILE pointer of the given file 

\sphinxlineitem{Returns}
\sphinxAtStartPar
current FILE pointer position cast to a bigint 

\end{description}\end{quote}

\end{fulllineitems}

\index{platform::fseek (C++ function)@\spxentry{platform::fseek}\spxextra{C++ function}}

\begin{fulllineitems}
\phantomsection\label{\detokenize{Developer_platform:_CPPv4N9LAMMPS_NS8platform5fseekEP4FILE6bigint}}\phantomsection\label{\detokenize{Developer_platform:_CPPv3N9LAMMPS_NS8platform5fseekEP4FILE6bigint}}\phantomsection\label{\detokenize{Developer_platform:_CPPv2N9LAMMPS_NS8platform5fseekEP4FILE6bigint}}\phantomsection\label{\detokenize{Developer_platform:LAMMPS_NS::platform::fseek__FILEP.bigint}}
\pysigstartsignatures
\pysigstartmultiline
\pysiglinewithargsret{\phantomsection\label{\detokenize{Developer_platform:platform_8h_1af983f9763b57efcf3888b2a19f9babcd}}\DUrole{kt}{int}\DUrole{w}{ }\sphinxcode{\sphinxupquote{\DUrole{n}{LAMMPS\_NS}\DUrole{p}{::}\DUrole{n}{platform}\DUrole{p}{::}}}\sphinxbfcode{\sphinxupquote{\DUrole{n}{fseek}}}}{\DUrole{n}{FILE}\DUrole{w}{ }\DUrole{p}{*}\DUrole{n,sig-param}{fp}\sphinxparamcomma \DUrole{n}{bigint}\DUrole{w}{ }\DUrole{n,sig-param}{pos}}{}
\pysigstopmultiline
\pysigstopsignatures
\sphinxAtStartPar
Set absolute file position

\sphinxAtStartPar
If the absolute position is END\_OF\_FILE, then position at the end of the file.
\begin{quote}\begin{description}
\sphinxlineitem{Parameters}\begin{itemize}
\item {} 
\sphinxAtStartPar
\sphinxstyleliteralstrong{\sphinxupquote{fp}} \textendash{} FILE pointer of the given file 

\item {} 
\sphinxAtStartPar
\sphinxstyleliteralstrong{\sphinxupquote{pos}} \textendash{} new position of the FILE pointer 

\end{itemize}

\sphinxlineitem{Returns}
\sphinxAtStartPar
0 if successful, otherwise \sphinxhyphen{}1 

\end{description}\end{quote}

\end{fulllineitems}

\index{platform::ftruncate (C++ function)@\spxentry{platform::ftruncate}\spxextra{C++ function}}

\begin{fulllineitems}
\phantomsection\label{\detokenize{Developer_platform:_CPPv4N9LAMMPS_NS8platform9ftruncateEP4FILE6bigint}}\phantomsection\label{\detokenize{Developer_platform:_CPPv3N9LAMMPS_NS8platform9ftruncateEP4FILE6bigint}}\phantomsection\label{\detokenize{Developer_platform:_CPPv2N9LAMMPS_NS8platform9ftruncateEP4FILE6bigint}}\phantomsection\label{\detokenize{Developer_platform:LAMMPS_NS::platform::ftruncate__FILEP.bigint}}
\pysigstartsignatures
\pysigstartmultiline
\pysiglinewithargsret{\phantomsection\label{\detokenize{Developer_platform:platform_8h_1ad9bfc18919956b05289d2eff3643aff1}}\DUrole{kt}{int}\DUrole{w}{ }\sphinxcode{\sphinxupquote{\DUrole{n}{LAMMPS\_NS}\DUrole{p}{::}\DUrole{n}{platform}\DUrole{p}{::}}}\sphinxbfcode{\sphinxupquote{\DUrole{n}{ftruncate}}}}{\DUrole{n}{FILE}\DUrole{w}{ }\DUrole{p}{*}\DUrole{n,sig-param}{fp}\sphinxparamcomma \DUrole{n}{bigint}\DUrole{w}{ }\DUrole{n,sig-param}{length}}{}
\pysigstopmultiline
\pysigstopsignatures
\sphinxAtStartPar
Truncate file to a given length and re\sphinxhyphen{}position file pointer
\begin{quote}\begin{description}
\sphinxlineitem{Parameters}\begin{itemize}
\item {} 
\sphinxAtStartPar
\sphinxstyleliteralstrong{\sphinxupquote{fp}} \textendash{} FILE pointer of the given file 

\item {} 
\sphinxAtStartPar
\sphinxstyleliteralstrong{\sphinxupquote{length}} \textendash{} length to which the file is being truncated to 

\end{itemize}

\sphinxlineitem{Returns}
\sphinxAtStartPar
0 if successful, otherwise \sphinxhyphen{}1 

\end{description}\end{quote}

\end{fulllineitems}

\index{platform::popen (C++ function)@\spxentry{platform::popen}\spxextra{C++ function}}

\begin{fulllineitems}
\phantomsection\label{\detokenize{Developer_platform:_CPPv4N9LAMMPS_NS8platform5popenERKNSt6stringERKNSt6stringE}}\phantomsection\label{\detokenize{Developer_platform:_CPPv3N9LAMMPS_NS8platform5popenERKNSt6stringERKNSt6stringE}}\phantomsection\label{\detokenize{Developer_platform:_CPPv2N9LAMMPS_NS8platform5popenERKNSt6stringERKNSt6stringE}}\phantomsection\label{\detokenize{Developer_platform:LAMMPS_NS::platform::popen__ssCR.ssCR}}
\pysigstartsignatures
\pysigstartmultiline
\pysiglinewithargsret{\phantomsection\label{\detokenize{Developer_platform:platform_8h_1ac811c6a9953c81b80fd6b761701b7bfe}}\DUrole{n}{FILE}\DUrole{w}{ }\DUrole{p}{*}\sphinxcode{\sphinxupquote{\DUrole{n}{LAMMPS\_NS}\DUrole{p}{::}\DUrole{n}{platform}\DUrole{p}{::}}}\sphinxbfcode{\sphinxupquote{\DUrole{n}{popen}}}}{\DUrole{k}{const}\DUrole{w}{ }\DUrole{n}{std}\DUrole{p}{::}\DUrole{n}{string}\DUrole{w}{ }\DUrole{p}{\&}\DUrole{n,sig-param}{cmd}\sphinxparamcomma \DUrole{k}{const}\DUrole{w}{ }\DUrole{n}{std}\DUrole{p}{::}\DUrole{n}{string}\DUrole{w}{ }\DUrole{p}{\&}\DUrole{n,sig-param}{mode}}{}
\pysigstopmultiline
\pysigstopsignatures
\sphinxAtStartPar
Open a pipe to a command for reading or writing
\begin{quote}\begin{description}
\sphinxlineitem{Parameters}\begin{itemize}
\item {} 
\sphinxAtStartPar
\sphinxstyleliteralstrong{\sphinxupquote{cmd}} \textendash{} command for the pipe 

\item {} 
\sphinxAtStartPar
\sphinxstyleliteralstrong{\sphinxupquote{mode}} \textendash{} “r” for reading from \sphinxstyleemphasis{cmd} or “w” for writing to \sphinxstyleemphasis{cmd}

\end{itemize}

\sphinxlineitem{Returns}
\sphinxAtStartPar
file pointer to the pipe if successful or null 

\end{description}\end{quote}

\end{fulllineitems}

\index{platform::pclose (C++ function)@\spxentry{platform::pclose}\spxextra{C++ function}}

\begin{fulllineitems}
\phantomsection\label{\detokenize{Developer_platform:_CPPv4N9LAMMPS_NS8platform6pcloseEP4FILE}}\phantomsection\label{\detokenize{Developer_platform:_CPPv3N9LAMMPS_NS8platform6pcloseEP4FILE}}\phantomsection\label{\detokenize{Developer_platform:_CPPv2N9LAMMPS_NS8platform6pcloseEP4FILE}}\phantomsection\label{\detokenize{Developer_platform:LAMMPS_NS::platform::pclose__FILEP}}
\pysigstartsignatures
\pysigstartmultiline
\pysiglinewithargsret{\phantomsection\label{\detokenize{Developer_platform:platform_8h_1abca9eaf3eb5f2010062ee8f6a4017733}}\DUrole{kt}{int}\DUrole{w}{ }\sphinxcode{\sphinxupquote{\DUrole{n}{LAMMPS\_NS}\DUrole{p}{::}\DUrole{n}{platform}\DUrole{p}{::}}}\sphinxbfcode{\sphinxupquote{\DUrole{n}{pclose}}}}{\DUrole{n}{FILE}\DUrole{w}{ }\DUrole{p}{*}\DUrole{n,sig-param}{fp}}{}
\pysigstopmultiline
\pysigstopsignatures
\sphinxAtStartPar
Close a previously opened pipe
\begin{quote}\begin{description}
\sphinxlineitem{Parameters}
\sphinxAtStartPar
\sphinxstyleliteralstrong{\sphinxupquote{fp}} \textendash{} FILE pointer for the pipe 

\sphinxlineitem{Returns}
\sphinxAtStartPar
exit status of the pipe command or \sphinxhyphen{}1 in case of errors 

\end{description}\end{quote}

\end{fulllineitems}



\subsection{Environment variable functions}
\label{\detokenize{Developer_platform:environment-variable-functions}}\index{platform::putenv (C++ function)@\spxentry{platform::putenv}\spxextra{C++ function}}

\begin{fulllineitems}
\phantomsection\label{\detokenize{Developer_platform:_CPPv4N9LAMMPS_NS8platform6putenvERKNSt6stringE}}\phantomsection\label{\detokenize{Developer_platform:_CPPv3N9LAMMPS_NS8platform6putenvERKNSt6stringE}}\phantomsection\label{\detokenize{Developer_platform:_CPPv2N9LAMMPS_NS8platform6putenvERKNSt6stringE}}\phantomsection\label{\detokenize{Developer_platform:LAMMPS_NS::platform::putenv__ssCR}}
\pysigstartsignatures
\pysigstartmultiline
\pysiglinewithargsret{\phantomsection\label{\detokenize{Developer_platform:platform_8h_1aae6f2f54e54291876e6bbd8f952c9c5f}}\DUrole{kt}{int}\DUrole{w}{ }\sphinxcode{\sphinxupquote{\DUrole{n}{LAMMPS\_NS}\DUrole{p}{::}\DUrole{n}{platform}\DUrole{p}{::}}}\sphinxbfcode{\sphinxupquote{\DUrole{n}{putenv}}}}{\DUrole{k}{const}\DUrole{w}{ }\DUrole{n}{std}\DUrole{p}{::}\DUrole{n}{string}\DUrole{w}{ }\DUrole{p}{\&}\DUrole{n,sig-param}{vardef}}{}
\pysigstopmultiline
\pysigstopsignatures
\sphinxAtStartPar
Add variable to the environment
\begin{quote}\begin{description}
\sphinxlineitem{Parameters}
\sphinxAtStartPar
\sphinxstyleliteralstrong{\sphinxupquote{vardef}} \textendash{} variable name or variable definition (NAME=value) 

\sphinxlineitem{Returns}
\sphinxAtStartPar
\sphinxhyphen{}1 if failure otherwise 0 

\end{description}\end{quote}

\end{fulllineitems}

\index{platform::unsetenv (C++ function)@\spxentry{platform::unsetenv}\spxextra{C++ function}}

\begin{fulllineitems}
\phantomsection\label{\detokenize{Developer_platform:_CPPv4N9LAMMPS_NS8platform8unsetenvERKNSt6stringE}}\phantomsection\label{\detokenize{Developer_platform:_CPPv3N9LAMMPS_NS8platform8unsetenvERKNSt6stringE}}\phantomsection\label{\detokenize{Developer_platform:_CPPv2N9LAMMPS_NS8platform8unsetenvERKNSt6stringE}}\phantomsection\label{\detokenize{Developer_platform:LAMMPS_NS::platform::unsetenv__ssCR}}
\pysigstartsignatures
\pysigstartmultiline
\pysiglinewithargsret{\phantomsection\label{\detokenize{Developer_platform:platform_8h_1add0ec024cf89b7fe7a83b4c54fdf8f94}}\DUrole{kt}{int}\DUrole{w}{ }\sphinxcode{\sphinxupquote{\DUrole{n}{LAMMPS\_NS}\DUrole{p}{::}\DUrole{n}{platform}\DUrole{p}{::}}}\sphinxbfcode{\sphinxupquote{\DUrole{n}{unsetenv}}}}{\DUrole{k}{const}\DUrole{w}{ }\DUrole{n}{std}\DUrole{p}{::}\DUrole{n}{string}\DUrole{w}{ }\DUrole{p}{\&}\DUrole{n,sig-param}{variable}}{}
\pysigstopmultiline
\pysigstopsignatures
\sphinxAtStartPar
Delete variable from the environment
\begin{quote}\begin{description}
\sphinxlineitem{Parameters}
\sphinxAtStartPar
\sphinxstyleliteralstrong{\sphinxupquote{variable}} \textendash{} variable name 

\sphinxlineitem{Returns}
\sphinxAtStartPar
\sphinxhyphen{}1 if failure otherwise 0 

\end{description}\end{quote}

\end{fulllineitems}

\index{platform::list\_pathenv (C++ function)@\spxentry{platform::list\_pathenv}\spxextra{C++ function}}

\begin{fulllineitems}
\phantomsection\label{\detokenize{Developer_platform:_CPPv4N9LAMMPS_NS8platform12list_pathenvERKNSt6stringE}}\phantomsection\label{\detokenize{Developer_platform:_CPPv3N9LAMMPS_NS8platform12list_pathenvERKNSt6stringE}}\phantomsection\label{\detokenize{Developer_platform:_CPPv2N9LAMMPS_NS8platform12list_pathenvERKNSt6stringE}}\phantomsection\label{\detokenize{Developer_platform:LAMMPS_NS::platform::list_pathenv__ssCR}}
\pysigstartsignatures
\pysigstartmultiline
\pysiglinewithargsret{\phantomsection\label{\detokenize{Developer_platform:platform_8h_1ac52391dc7686d4a9e21b5d6ded324869}}\DUrole{n}{std}\DUrole{p}{::}\DUrole{n}{vector}\DUrole{p}{\textless{}}\DUrole{n}{std}\DUrole{p}{::}\DUrole{n}{string}\DUrole{p}{\textgreater{}}\DUrole{w}{ }\sphinxcode{\sphinxupquote{\DUrole{n}{LAMMPS\_NS}\DUrole{p}{::}\DUrole{n}{platform}\DUrole{p}{::}}}\sphinxbfcode{\sphinxupquote{\DUrole{n}{list\_pathenv}}}}{\DUrole{k}{const}\DUrole{w}{ }\DUrole{n}{std}\DUrole{p}{::}\DUrole{n}{string}\DUrole{w}{ }\DUrole{p}{\&}\DUrole{n,sig-param}{var}}{}
\pysigstopmultiline
\pysigstopsignatures
\sphinxAtStartPar
Get list of entries in a path environment variable

\sphinxAtStartPar
This provides a list of strings of the entries in an environment variable that is containing a “path” like “PATH” or “LD\_LIBRARY\_PATH”.
\begin{quote}\begin{description}
\sphinxlineitem{Parameters}
\sphinxAtStartPar
\sphinxstyleliteralstrong{\sphinxupquote{var}} \textendash{} name of the environment variable 

\sphinxlineitem{Returns}
\sphinxAtStartPar
vector with strings of all entries in that path variable 

\end{description}\end{quote}

\end{fulllineitems}

\index{platform::find\_exe\_path (C++ function)@\spxentry{platform::find\_exe\_path}\spxextra{C++ function}}

\begin{fulllineitems}
\phantomsection\label{\detokenize{Developer_platform:_CPPv4N9LAMMPS_NS8platform13find_exe_pathERKNSt6stringE}}\phantomsection\label{\detokenize{Developer_platform:_CPPv3N9LAMMPS_NS8platform13find_exe_pathERKNSt6stringE}}\phantomsection\label{\detokenize{Developer_platform:_CPPv2N9LAMMPS_NS8platform13find_exe_pathERKNSt6stringE}}\phantomsection\label{\detokenize{Developer_platform:LAMMPS_NS::platform::find_exe_path__ssCR}}
\pysigstartsignatures
\pysigstartmultiline
\pysiglinewithargsret{\phantomsection\label{\detokenize{Developer_platform:platform_8h_1ac7d2f135d11783248a70f4a66fddf671}}\DUrole{n}{std}\DUrole{p}{::}\DUrole{n}{string}\DUrole{w}{ }\sphinxcode{\sphinxupquote{\DUrole{n}{LAMMPS\_NS}\DUrole{p}{::}\DUrole{n}{platform}\DUrole{p}{::}}}\sphinxbfcode{\sphinxupquote{\DUrole{n}{find\_exe\_path}}}}{\DUrole{k}{const}\DUrole{w}{ }\DUrole{n}{std}\DUrole{p}{::}\DUrole{n}{string}\DUrole{w}{ }\DUrole{p}{\&}\DUrole{n,sig-param}{cmd}}{}
\pysigstopmultiline
\pysigstopsignatures
\sphinxAtStartPar
Find pathname of an executable in the standard search path

\sphinxAtStartPar
This function will traverse the list of directories in the PATH environment variable and look for the executable \sphinxstyleemphasis{cmd}. If the file exists and is executable the full path is returned as string, otherwise an empty string is returned.

\sphinxAtStartPar
On Windows the \sphinxstyleemphasis{cmd} string must not include an extension as this function will automatically append the extensions “.exe”, “.com” and “.bat” and look for those paths. On Windows also the current directory is checked (and first), but otherwise is not checked unless “.” exists in the PATH environment variable.

\sphinxAtStartPar
Because of the nature of the check, this will not detect shell functions built\sphinxhyphen{}in command or aliases.
\begin{quote}\begin{description}
\sphinxlineitem{Parameters}
\sphinxAtStartPar
\sphinxstyleliteralstrong{\sphinxupquote{cmd}} \textendash{} name of command 

\sphinxlineitem{Returns}
\sphinxAtStartPar
vector with strings of all directory entries 

\end{description}\end{quote}

\end{fulllineitems}



\subsection{Dynamically loaded object or library functions}
\label{\detokenize{Developer_platform:dynamically-loaded-object-or-library-functions}}\index{platform::dlopen (C++ function)@\spxentry{platform::dlopen}\spxextra{C++ function}}

\begin{fulllineitems}
\phantomsection\label{\detokenize{Developer_platform:_CPPv4N9LAMMPS_NS8platform6dlopenERKNSt6stringE}}\phantomsection\label{\detokenize{Developer_platform:_CPPv3N9LAMMPS_NS8platform6dlopenERKNSt6stringE}}\phantomsection\label{\detokenize{Developer_platform:_CPPv2N9LAMMPS_NS8platform6dlopenERKNSt6stringE}}\phantomsection\label{\detokenize{Developer_platform:LAMMPS_NS::platform::dlopen__ssCR}}
\pysigstartsignatures
\pysigstartmultiline
\pysiglinewithargsret{\phantomsection\label{\detokenize{Developer_platform:platform_8h_1acd2e0d5c0aa701ba30307f5dfee05b8c}}\DUrole{kt}{void}\DUrole{w}{ }\DUrole{p}{*}\sphinxcode{\sphinxupquote{\DUrole{n}{LAMMPS\_NS}\DUrole{p}{::}\DUrole{n}{platform}\DUrole{p}{::}}}\sphinxbfcode{\sphinxupquote{\DUrole{n}{dlopen}}}}{\DUrole{k}{const}\DUrole{w}{ }\DUrole{n}{std}\DUrole{p}{::}\DUrole{n}{string}\DUrole{w}{ }\DUrole{p}{\&}\DUrole{n,sig-param}{fname}}{}
\pysigstopmultiline
\pysigstopsignatures
\sphinxAtStartPar
Open a shared object file or library
\begin{quote}\begin{description}
\sphinxlineitem{Parameters}
\sphinxAtStartPar
\sphinxstyleliteralstrong{\sphinxupquote{fname}} \textendash{} name or path of the shared object 

\sphinxlineitem{Returns}
\sphinxAtStartPar
handle to the shared object or null 

\end{description}\end{quote}

\end{fulllineitems}

\index{platform::dlclose (C++ function)@\spxentry{platform::dlclose}\spxextra{C++ function}}

\begin{fulllineitems}
\phantomsection\label{\detokenize{Developer_platform:_CPPv4N9LAMMPS_NS8platform7dlcloseEPv}}\phantomsection\label{\detokenize{Developer_platform:_CPPv3N9LAMMPS_NS8platform7dlcloseEPv}}\phantomsection\label{\detokenize{Developer_platform:_CPPv2N9LAMMPS_NS8platform7dlcloseEPv}}\phantomsection\label{\detokenize{Developer_platform:LAMMPS_NS::platform::dlclose__voidP}}
\pysigstartsignatures
\pysigstartmultiline
\pysiglinewithargsret{\phantomsection\label{\detokenize{Developer_platform:platform_8h_1af483b6c4400965aa885f9f1a144138a5}}\DUrole{kt}{int}\DUrole{w}{ }\sphinxcode{\sphinxupquote{\DUrole{n}{LAMMPS\_NS}\DUrole{p}{::}\DUrole{n}{platform}\DUrole{p}{::}}}\sphinxbfcode{\sphinxupquote{\DUrole{n}{dlclose}}}}{\DUrole{kt}{void}\DUrole{w}{ }\DUrole{p}{*}\DUrole{n,sig-param}{handle}}{}
\pysigstopmultiline
\pysigstopsignatures
\sphinxAtStartPar
Close a shared object

\sphinxAtStartPar
This releases the object corresponding to the provided handle. Resolved symbols associated with this handle may not be used after this call
\begin{quote}\begin{description}
\sphinxlineitem{Parameters}
\sphinxAtStartPar
\sphinxstyleliteralstrong{\sphinxupquote{handle}} \textendash{} handle to an opened shared object 

\sphinxlineitem{Returns}
\sphinxAtStartPar
0 if successful, non\sphinxhyphen{}zero of not 

\end{description}\end{quote}

\end{fulllineitems}

\index{platform::dlsym (C++ function)@\spxentry{platform::dlsym}\spxextra{C++ function}}

\begin{fulllineitems}
\phantomsection\label{\detokenize{Developer_platform:_CPPv4N9LAMMPS_NS8platform5dlsymEPvRKNSt6stringE}}\phantomsection\label{\detokenize{Developer_platform:_CPPv3N9LAMMPS_NS8platform5dlsymEPvRKNSt6stringE}}\phantomsection\label{\detokenize{Developer_platform:_CPPv2N9LAMMPS_NS8platform5dlsymEPvRKNSt6stringE}}\phantomsection\label{\detokenize{Developer_platform:LAMMPS_NS::platform::dlsym__voidP.ssCR}}
\pysigstartsignatures
\pysigstartmultiline
\pysiglinewithargsret{\phantomsection\label{\detokenize{Developer_platform:platform_8h_1ab1b3f356af1d5cbe80f8188b1d63770b}}\DUrole{kt}{void}\DUrole{w}{ }\DUrole{p}{*}\sphinxcode{\sphinxupquote{\DUrole{n}{LAMMPS\_NS}\DUrole{p}{::}\DUrole{n}{platform}\DUrole{p}{::}}}\sphinxbfcode{\sphinxupquote{\DUrole{n}{dlsym}}}}{\DUrole{kt}{void}\DUrole{w}{ }\DUrole{p}{*}\DUrole{n,sig-param}{handle}\sphinxparamcomma \DUrole{k}{const}\DUrole{w}{ }\DUrole{n}{std}\DUrole{p}{::}\DUrole{n}{string}\DUrole{w}{ }\DUrole{p}{\&}\DUrole{n,sig-param}{symbol}}{}
\pysigstopmultiline
\pysigstopsignatures
\sphinxAtStartPar
Resolve a symbol in shared object
\begin{quote}\begin{description}
\sphinxlineitem{Parameters}\begin{itemize}
\item {} 
\sphinxAtStartPar
\sphinxstyleliteralstrong{\sphinxupquote{handle}} \textendash{} handle to an opened shared object 

\item {} 
\sphinxAtStartPar
\sphinxstyleliteralstrong{\sphinxupquote{symbol}} \textendash{} name of the symbol to extract 

\end{itemize}

\sphinxlineitem{Returns}
\sphinxAtStartPar
pointer to the resolved symbol or null 

\end{description}\end{quote}

\end{fulllineitems}

\index{platform::dlerror (C++ function)@\spxentry{platform::dlerror}\spxextra{C++ function}}

\begin{fulllineitems}
\phantomsection\label{\detokenize{Developer_platform:_CPPv4N9LAMMPS_NS8platform7dlerrorEv}}\phantomsection\label{\detokenize{Developer_platform:_CPPv3N9LAMMPS_NS8platform7dlerrorEv}}\phantomsection\label{\detokenize{Developer_platform:_CPPv2N9LAMMPS_NS8platform7dlerrorEv}}\phantomsection\label{\detokenize{Developer_platform:LAMMPS_NS::platform::dlerror}}
\pysigstartsignatures
\pysigstartmultiline
\pysiglinewithargsret{\phantomsection\label{\detokenize{Developer_platform:platform_8h_1a5cfc90f8002e180ad96c73a101c6aa23}}\DUrole{n}{std}\DUrole{p}{::}\DUrole{n}{string}\DUrole{w}{ }\sphinxcode{\sphinxupquote{\DUrole{n}{LAMMPS\_NS}\DUrole{p}{::}\DUrole{n}{platform}\DUrole{p}{::}}}\sphinxbfcode{\sphinxupquote{\DUrole{n}{dlerror}}}}{}{}
\pysigstopmultiline
\pysigstopsignatures
\sphinxAtStartPar
Obtain error diagnostic info after dynamic linking function calls

\sphinxAtStartPar
Return a human\sphinxhyphen{}readable string describing the most recent error that occurred when using one of the functions for dynamic loading objects the last call to this function. If there was no error, the string is empty.
\begin{quote}\begin{description}
\sphinxlineitem{Returns}
\sphinxAtStartPar
string with error message or empty 

\end{description}\end{quote}

\end{fulllineitems}



\subsection{Compressed file I/O functions}
\label{\detokenize{Developer_platform:compressed-file-i-o-functions}}\index{platform::has\_compress\_extension (C++ function)@\spxentry{platform::has\_compress\_extension}\spxextra{C++ function}}

\begin{fulllineitems}
\phantomsection\label{\detokenize{Developer_platform:_CPPv4N9LAMMPS_NS8platform22has_compress_extensionERKNSt6stringE}}\phantomsection\label{\detokenize{Developer_platform:_CPPv3N9LAMMPS_NS8platform22has_compress_extensionERKNSt6stringE}}\phantomsection\label{\detokenize{Developer_platform:_CPPv2N9LAMMPS_NS8platform22has_compress_extensionERKNSt6stringE}}\phantomsection\label{\detokenize{Developer_platform:LAMMPS_NS::platform::has_compress_extension__ssCR}}
\pysigstartsignatures
\pysigstartmultiline
\pysiglinewithargsret{\phantomsection\label{\detokenize{Developer_platform:platform_8h_1acb172eb1662e351bab844ccfafbd9427}}\DUrole{kt}{bool}\DUrole{w}{ }\sphinxcode{\sphinxupquote{\DUrole{n}{LAMMPS\_NS}\DUrole{p}{::}\DUrole{n}{platform}\DUrole{p}{::}}}\sphinxbfcode{\sphinxupquote{\DUrole{n}{has\_compress\_extension}}}}{\DUrole{k}{const}\DUrole{w}{ }\DUrole{n}{std}\DUrole{p}{::}\DUrole{n}{string}\DUrole{w}{ }\DUrole{p}{\&}\DUrole{n,sig-param}{file}}{}
\pysigstopmultiline
\pysigstopsignatures
\sphinxAtStartPar
Check if a file name ends in a known extension for a compressed file format

\sphinxAtStartPar
Currently supported file extensions are: .gz, .bz2, .zst, .xz, .lzma, lz4
\begin{quote}\begin{description}
\sphinxlineitem{Parameters}
\sphinxAtStartPar
\sphinxstyleliteralstrong{\sphinxupquote{file}} \textendash{} name of the file to check 

\sphinxlineitem{Returns}
\sphinxAtStartPar
true if the file has a known extension, otherwise false 

\end{description}\end{quote}

\end{fulllineitems}

\index{platform::compressed\_read (C++ function)@\spxentry{platform::compressed\_read}\spxextra{C++ function}}

\begin{fulllineitems}
\phantomsection\label{\detokenize{Developer_platform:_CPPv4N9LAMMPS_NS8platform15compressed_readERKNSt6stringE}}\phantomsection\label{\detokenize{Developer_platform:_CPPv3N9LAMMPS_NS8platform15compressed_readERKNSt6stringE}}\phantomsection\label{\detokenize{Developer_platform:_CPPv2N9LAMMPS_NS8platform15compressed_readERKNSt6stringE}}\phantomsection\label{\detokenize{Developer_platform:LAMMPS_NS::platform::compressed_read__ssCR}}
\pysigstartsignatures
\pysigstartmultiline
\pysiglinewithargsret{\phantomsection\label{\detokenize{Developer_platform:platform_8h_1a07977ea469d1405938f4225010fd0a3e}}\DUrole{n}{FILE}\DUrole{w}{ }\DUrole{p}{*}\sphinxcode{\sphinxupquote{\DUrole{n}{LAMMPS\_NS}\DUrole{p}{::}\DUrole{n}{platform}\DUrole{p}{::}}}\sphinxbfcode{\sphinxupquote{\DUrole{n}{compressed\_read}}}}{\DUrole{k}{const}\DUrole{w}{ }\DUrole{n}{std}\DUrole{p}{::}\DUrole{n}{string}\DUrole{w}{ }\DUrole{p}{\&}\DUrole{n,sig-param}{file}}{}
\pysigstopmultiline
\pysigstopsignatures
\sphinxAtStartPar
Open pipe to compressed text file for reading
\begin{quote}\begin{description}
\sphinxlineitem{Parameters}
\sphinxAtStartPar
\sphinxstyleliteralstrong{\sphinxupquote{file}} \textendash{} name of the file to open 

\sphinxlineitem{Returns}
\sphinxAtStartPar
FILE pointer to pipe using for reading the compressed file. 

\end{description}\end{quote}

\end{fulllineitems}

\index{platform::compressed\_write (C++ function)@\spxentry{platform::compressed\_write}\spxextra{C++ function}}

\begin{fulllineitems}
\phantomsection\label{\detokenize{Developer_platform:_CPPv4N9LAMMPS_NS8platform16compressed_writeERKNSt6stringE}}\phantomsection\label{\detokenize{Developer_platform:_CPPv3N9LAMMPS_NS8platform16compressed_writeERKNSt6stringE}}\phantomsection\label{\detokenize{Developer_platform:_CPPv2N9LAMMPS_NS8platform16compressed_writeERKNSt6stringE}}\phantomsection\label{\detokenize{Developer_platform:LAMMPS_NS::platform::compressed_write__ssCR}}
\pysigstartsignatures
\pysigstartmultiline
\pysiglinewithargsret{\phantomsection\label{\detokenize{Developer_platform:platform_8h_1a6789ec6955f5cbc9555c83bd0b3ef1c1}}\DUrole{n}{FILE}\DUrole{w}{ }\DUrole{p}{*}\sphinxcode{\sphinxupquote{\DUrole{n}{LAMMPS\_NS}\DUrole{p}{::}\DUrole{n}{platform}\DUrole{p}{::}}}\sphinxbfcode{\sphinxupquote{\DUrole{n}{compressed\_write}}}}{\DUrole{k}{const}\DUrole{w}{ }\DUrole{n}{std}\DUrole{p}{::}\DUrole{n}{string}\DUrole{w}{ }\DUrole{p}{\&}\DUrole{n,sig-param}{file}}{}
\pysigstopmultiline
\pysigstopsignatures
\sphinxAtStartPar
Open pipe to compressed text file for writing
\begin{quote}\begin{description}
\sphinxlineitem{Parameters}
\sphinxAtStartPar
\sphinxstyleliteralstrong{\sphinxupquote{file}} \textendash{} name of the file to open 

\sphinxlineitem{Returns}
\sphinxAtStartPar
FILE pointer to pipe using for reading the compressed file. 

\end{description}\end{quote}

\end{fulllineitems}


\sphinxstepscope


\section{Utility functions}
\label{\detokenize{Developer_utils:utility-functions}}\label{\detokenize{Developer_utils::doc}}
\sphinxAtStartPar
The \sphinxcode{\sphinxupquote{utils}} sub\sphinxhyphen{}namespace inside the \sphinxcode{\sphinxupquote{LAMMPS\_NS}} namespace provides
a collection of convenience functions and utilities that perform common
tasks that are required repeatedly throughout the LAMMPS code like
reading or writing to files with error checking or translation of
strings into specific types of numbers with checking for validity.  This
reduces redundant implementations and encourages consistent behavior and
thus has some overlap with the {\hyperref[\detokenize{Developer_platform::doc}]{\sphinxcrossref{\DUrole{doc}{“platform” sub\sphinxhyphen{}namespace}}}}.


\subsection{I/O with status check and similar functions}
\label{\detokenize{Developer_utils:i-o-with-status-check-and-similar-functions}}
\sphinxAtStartPar
The the first two functions are wrappers around the corresponding C
library calls \sphinxcode{\sphinxupquote{fgets()}} or \sphinxcode{\sphinxupquote{fread()}}.  They will check if there
were errors on reading or an unexpected end\sphinxhyphen{}of\sphinxhyphen{}file state was reached.
In that case, the functions will stop with an error message, indicating
the name of the problematic file, if possible unless the \sphinxstyleemphasis{error} argument
is a NULL pointer.

\sphinxAtStartPar
The {\hyperref[\detokenize{Developer_utils:_CPPv4N9LAMMPS_NS5utils11fgets_truncEPciP4FILE}]{\sphinxcrossref{\sphinxcode{\sphinxupquote{utils::fgets\_trunc()}}}}}
function will work similar for \sphinxcode{\sphinxupquote{fgets()}} but it will read in a whole
line (i.e. until the end of line or end of file), but store only as many
characters as will fit into the buffer including a final newline
character and the terminating NULL byte.  If the line in the file is
longer it will thus be truncated in the buffer.  This function is used
by {\hyperref[\detokenize{Developer_utils:_CPPv4N9LAMMPS_NS5utils20read_lines_from_fileEP4FILEiiPci8MPI_Comm}]{\sphinxcrossref{\sphinxcode{\sphinxupquote{utils::read\_lines\_from\_file()}}}}} to read individual lines but
make certain they follow the size constraints.

\sphinxAtStartPar
The {\hyperref[\detokenize{Developer_utils:_CPPv4N9LAMMPS_NS5utils20read_lines_from_fileEP4FILEiiPci8MPI_Comm}]{\sphinxcrossref{\sphinxcode{\sphinxupquote{utils::read\_lines\_from\_file()}}}}} function will read the
requested number of lines of a maximum length into a buffer and will
return 0 if successful or 1 if not. It also guarantees that all lines
are terminated with a newline character and the entire buffer with a
NULL character.


\bigskip\hrule\bigskip

\index{utils::sfgets (C++ function)@\spxentry{utils::sfgets}\spxextra{C++ function}}

\begin{fulllineitems}
\phantomsection\label{\detokenize{Developer_utils:_CPPv4N9LAMMPS_NS5utils6sfgetsEPKciPciP4FILEPKcP5Error}}\phantomsection\label{\detokenize{Developer_utils:_CPPv3N9LAMMPS_NS5utils6sfgetsEPKciPciP4FILEPKcP5Error}}\phantomsection\label{\detokenize{Developer_utils:_CPPv2N9LAMMPS_NS5utils6sfgetsEPKciPciP4FILEPKcP5Error}}\phantomsection\label{\detokenize{Developer_utils:LAMMPS_NS::utils::sfgets__cCP.i.cP.i.FILEP.cCP.ErrorP}}
\pysigstartsignatures
\pysigstartmultiline
\pysiglinewithargsret{\phantomsection\label{\detokenize{Developer_utils:utils_8h_1a151f0249d5fa35b2f79f9fd411f0c61b}}\DUrole{kt}{void}\DUrole{w}{ }\sphinxcode{\sphinxupquote{\DUrole{n}{LAMMPS\_NS}\DUrole{p}{::}\DUrole{n}{utils}\DUrole{p}{::}}}\sphinxbfcode{\sphinxupquote{\DUrole{n}{sfgets}}}}{\DUrole{k}{const}\DUrole{w}{ }\DUrole{kt}{char}\DUrole{w}{ }\DUrole{p}{*}\DUrole{n,sig-param}{srcname}\sphinxparamcomma \DUrole{kt}{int}\DUrole{w}{ }\DUrole{n,sig-param}{srcline}\sphinxparamcomma \DUrole{kt}{char}\DUrole{w}{ }\DUrole{p}{*}\DUrole{n,sig-param}{s}\sphinxparamcomma \DUrole{kt}{int}\DUrole{w}{ }\DUrole{n,sig-param}{size}\sphinxparamcomma \DUrole{n}{FILE}\DUrole{w}{ }\DUrole{p}{*}\DUrole{n,sig-param}{fp}\sphinxparamcomma \DUrole{k}{const}\DUrole{w}{ }\DUrole{kt}{char}\DUrole{w}{ }\DUrole{p}{*}\DUrole{n,sig-param}{filename}\sphinxparamcomma \DUrole{n}{Error}\DUrole{w}{ }\DUrole{p}{*}\DUrole{n,sig-param}{error}}{}
\pysigstopmultiline
\pysigstopsignatures
\sphinxAtStartPar
Safe wrapper around fgets() which aborts on errors or EOF and prints a suitable error message to help debugging.

\sphinxAtStartPar
Use nullptr as the error parameter to avoid the abort on EOF or error.
\begin{quote}\begin{description}
\sphinxlineitem{Parameters}\begin{itemize}
\item {} 
\sphinxAtStartPar
\sphinxstyleliteralstrong{\sphinxupquote{srcname}} \textendash{} name of the calling source file (from FLERR macro) 

\item {} 
\sphinxAtStartPar
\sphinxstyleliteralstrong{\sphinxupquote{srcline}} \textendash{} line in the calling source file (from FLERR macro) 

\item {} 
\sphinxAtStartPar
\sphinxstyleliteralstrong{\sphinxupquote{s}} \textendash{} buffer for storing the result of fgets() 

\item {} 
\sphinxAtStartPar
\sphinxstyleliteralstrong{\sphinxupquote{size}} \textendash{} size of buffer s (max number of bytes read by fgets()) 

\item {} 
\sphinxAtStartPar
\sphinxstyleliteralstrong{\sphinxupquote{fp}} \textendash{} file pointer used by fgets() 

\item {} 
\sphinxAtStartPar
\sphinxstyleliteralstrong{\sphinxupquote{filename}} \textendash{} file name associated with fp (may be a null pointer; then {\hyperref[\detokenize{Classes_lammps:classLAMMPS__NS_1_1LAMMPS}]{\sphinxcrossref{\DUrole{std,std-ref}{LAMMPS}}}} will try to detect) 

\item {} 
\sphinxAtStartPar
\sphinxstyleliteralstrong{\sphinxupquote{error}} \textendash{} pointer to Error class instance (for abort) or nullptr 

\end{itemize}

\end{description}\end{quote}

\end{fulllineitems}

\index{utils::sfread (C++ function)@\spxentry{utils::sfread}\spxextra{C++ function}}

\begin{fulllineitems}
\phantomsection\label{\detokenize{Developer_utils:_CPPv4N9LAMMPS_NS5utils6sfreadEPKciPv6size_t6size_tP4FILEPKcP5Error}}\phantomsection\label{\detokenize{Developer_utils:_CPPv3N9LAMMPS_NS5utils6sfreadEPKciPv6size_t6size_tP4FILEPKcP5Error}}\phantomsection\label{\detokenize{Developer_utils:_CPPv2N9LAMMPS_NS5utils6sfreadEPKciPv6size_t6size_tP4FILEPKcP5Error}}\phantomsection\label{\detokenize{Developer_utils:LAMMPS_NS::utils::sfread__cCP.i.voidP.s.s.FILEP.cCP.ErrorP}}
\pysigstartsignatures
\pysigstartmultiline
\pysiglinewithargsret{\phantomsection\label{\detokenize{Developer_utils:utils_8h_1a5fc62aaa4658e16ec52523996d92af83}}\DUrole{kt}{void}\DUrole{w}{ }\sphinxcode{\sphinxupquote{\DUrole{n}{LAMMPS\_NS}\DUrole{p}{::}\DUrole{n}{utils}\DUrole{p}{::}}}\sphinxbfcode{\sphinxupquote{\DUrole{n}{sfread}}}}{\DUrole{k}{const}\DUrole{w}{ }\DUrole{kt}{char}\DUrole{w}{ }\DUrole{p}{*}\DUrole{n,sig-param}{srcname}\sphinxparamcomma \DUrole{kt}{int}\DUrole{w}{ }\DUrole{n,sig-param}{srcline}\sphinxparamcomma \DUrole{kt}{void}\DUrole{w}{ }\DUrole{p}{*}\DUrole{n,sig-param}{s}\sphinxparamcomma \DUrole{n}{size\_t}\DUrole{w}{ }\DUrole{n,sig-param}{size}\sphinxparamcomma \DUrole{n}{size\_t}\DUrole{w}{ }\DUrole{n,sig-param}{num}\sphinxparamcomma \DUrole{n}{FILE}\DUrole{w}{ }\DUrole{p}{*}\DUrole{n,sig-param}{fp}\sphinxparamcomma \DUrole{k}{const}\DUrole{w}{ }\DUrole{kt}{char}\DUrole{w}{ }\DUrole{p}{*}\DUrole{n,sig-param}{filename}\sphinxparamcomma \DUrole{n}{Error}\DUrole{w}{ }\DUrole{p}{*}\DUrole{n,sig-param}{error}}{}
\pysigstopmultiline
\pysigstopsignatures
\sphinxAtStartPar
Safe wrapper around fread() which aborts on errors or EOF and prints a suitable error message to help debugging.

\sphinxAtStartPar
Use nullptr as the error parameter to avoid the abort on EOF or error.
\begin{quote}\begin{description}
\sphinxlineitem{Parameters}\begin{itemize}
\item {} 
\sphinxAtStartPar
\sphinxstyleliteralstrong{\sphinxupquote{srcname}} \textendash{} name of the calling source file (from FLERR macro) 

\item {} 
\sphinxAtStartPar
\sphinxstyleliteralstrong{\sphinxupquote{srcline}} \textendash{} line in the calling source file (from FLERR macro) 

\item {} 
\sphinxAtStartPar
\sphinxstyleliteralstrong{\sphinxupquote{s}} \textendash{} buffer for storing the result of fread() 

\item {} 
\sphinxAtStartPar
\sphinxstyleliteralstrong{\sphinxupquote{size}} \textendash{} size of data elements read by fread() 

\item {} 
\sphinxAtStartPar
\sphinxstyleliteralstrong{\sphinxupquote{num}} \textendash{} number of data elements read by fread() 

\item {} 
\sphinxAtStartPar
\sphinxstyleliteralstrong{\sphinxupquote{fp}} \textendash{} file pointer used by fread() 

\item {} 
\sphinxAtStartPar
\sphinxstyleliteralstrong{\sphinxupquote{filename}} \textendash{} file name associated with fp (may be a null pointer; then {\hyperref[\detokenize{Classes_lammps:classLAMMPS__NS_1_1LAMMPS}]{\sphinxcrossref{\DUrole{std,std-ref}{LAMMPS}}}} will try to detect) 

\item {} 
\sphinxAtStartPar
\sphinxstyleliteralstrong{\sphinxupquote{error}} \textendash{} pointer to Error class instance (for abort) or nullptr 

\end{itemize}

\end{description}\end{quote}

\end{fulllineitems}

\index{utils::fgets\_trunc (C++ function)@\spxentry{utils::fgets\_trunc}\spxextra{C++ function}}

\begin{fulllineitems}
\phantomsection\label{\detokenize{Developer_utils:_CPPv4N9LAMMPS_NS5utils11fgets_truncEPciP4FILE}}\phantomsection\label{\detokenize{Developer_utils:_CPPv3N9LAMMPS_NS5utils11fgets_truncEPciP4FILE}}\phantomsection\label{\detokenize{Developer_utils:_CPPv2N9LAMMPS_NS5utils11fgets_truncEPciP4FILE}}\phantomsection\label{\detokenize{Developer_utils:LAMMPS_NS::utils::fgets_trunc__cP.i.FILEP}}
\pysigstartsignatures
\pysigstartmultiline
\pysiglinewithargsret{\phantomsection\label{\detokenize{Developer_utils:utils_8h_1ae2af7b66ca9a1980b73c1941779e2047}}\DUrole{kt}{char}\DUrole{w}{ }\DUrole{p}{*}\sphinxcode{\sphinxupquote{\DUrole{n}{LAMMPS\_NS}\DUrole{p}{::}\DUrole{n}{utils}\DUrole{p}{::}}}\sphinxbfcode{\sphinxupquote{\DUrole{n}{fgets\_trunc}}}}{\DUrole{kt}{char}\DUrole{w}{ }\DUrole{p}{*}\DUrole{n,sig-param}{s}\sphinxparamcomma \DUrole{kt}{int}\DUrole{w}{ }\DUrole{n,sig-param}{size}\sphinxparamcomma \DUrole{n}{FILE}\DUrole{w}{ }\DUrole{p}{*}\DUrole{n,sig-param}{fp}}{}
\pysigstopmultiline
\pysigstopsignatures
\sphinxAtStartPar
Wrapper around fgets() which reads whole lines but truncates the data to the buffer size and ensures a newline char at the end.

\sphinxAtStartPar
This function is useful for reading line based text files with possible comments that should be parsed later. This applies to data files, potential files, atomfile variable files and so on. It is used instead of fgets() by utils::read\_lines\_from\_file().
\begin{quote}\begin{description}
\sphinxlineitem{Parameters}\begin{itemize}
\item {} 
\sphinxAtStartPar
\sphinxstyleliteralstrong{\sphinxupquote{s}} \textendash{} buffer for storing the result of fgets() 

\item {} 
\sphinxAtStartPar
\sphinxstyleliteralstrong{\sphinxupquote{size}} \textendash{} size of buffer s (max number of bytes returned) 

\item {} 
\sphinxAtStartPar
\sphinxstyleliteralstrong{\sphinxupquote{fp}} \textendash{} file pointer used by fgets() 

\end{itemize}

\end{description}\end{quote}

\end{fulllineitems}

\index{utils::read\_lines\_from\_file (C++ function)@\spxentry{utils::read\_lines\_from\_file}\spxextra{C++ function}}

\begin{fulllineitems}
\phantomsection\label{\detokenize{Developer_utils:_CPPv4N9LAMMPS_NS5utils20read_lines_from_fileEP4FILEiiPci8MPI_Comm}}\phantomsection\label{\detokenize{Developer_utils:_CPPv3N9LAMMPS_NS5utils20read_lines_from_fileEP4FILEiiPci8MPI_Comm}}\phantomsection\label{\detokenize{Developer_utils:_CPPv2N9LAMMPS_NS5utils20read_lines_from_fileEP4FILEiiPci8MPI_Comm}}\phantomsection\label{\detokenize{Developer_utils:LAMMPS_NS::utils::read_lines_from_file__FILEP.i.i.cP.i.MPI_Comm}}
\pysigstartsignatures
\pysigstartmultiline
\pysiglinewithargsret{\phantomsection\label{\detokenize{Developer_utils:utils_8h_1a7de7c579867fe5cac7f68affc3c9d416}}\DUrole{kt}{int}\DUrole{w}{ }\sphinxcode{\sphinxupquote{\DUrole{n}{LAMMPS\_NS}\DUrole{p}{::}\DUrole{n}{utils}\DUrole{p}{::}}}\sphinxbfcode{\sphinxupquote{\DUrole{n}{read\_lines\_from\_file}}}}{\DUrole{n}{FILE}\DUrole{w}{ }\DUrole{p}{*}\DUrole{n,sig-param}{fp}\sphinxparamcomma \DUrole{kt}{int}\DUrole{w}{ }\DUrole{n,sig-param}{nlines}\sphinxparamcomma \DUrole{kt}{int}\DUrole{w}{ }\DUrole{n,sig-param}{nmax}\sphinxparamcomma \DUrole{kt}{char}\DUrole{w}{ }\DUrole{p}{*}\DUrole{n,sig-param}{buffer}\sphinxparamcomma \DUrole{kt}{int}\DUrole{w}{ }\DUrole{n,sig-param}{me}\sphinxparamcomma \DUrole{n}{MPI\_Comm}\DUrole{w}{ }\DUrole{n,sig-param}{comm}}{}
\pysigstopmultiline
\pysigstopsignatures
\sphinxAtStartPar
Read N lines of text from file into buffer and broadcast them

\sphinxAtStartPar
This function uses repeated calls to fread() to fill a buffer with newline terminated text. If a line does not end in a newline (e.g. at the end of a file), it is added. The caller has to allocate an nlines by nmax sized buffer for storing the text data. Reading is done by MPI rank 0 of the given communicator only, and thus only MPI rank 0 needs to provide a valid file pointer.
\begin{quote}\begin{description}
\sphinxlineitem{Parameters}\begin{itemize}
\item {} 
\sphinxAtStartPar
\sphinxstyleliteralstrong{\sphinxupquote{fp}} \textendash{} file pointer used by fread 

\item {} 
\sphinxAtStartPar
\sphinxstyleliteralstrong{\sphinxupquote{nlines}} \textendash{} number of lines to be read 

\item {} 
\sphinxAtStartPar
\sphinxstyleliteralstrong{\sphinxupquote{nmax}} \textendash{} maximum length of a single line 

\item {} 
\sphinxAtStartPar
\sphinxstyleliteralstrong{\sphinxupquote{buffer}} \textendash{} buffer for storing the data. 

\item {} 
\sphinxAtStartPar
\sphinxstyleliteralstrong{\sphinxupquote{me}} \textendash{} MPI rank of calling process in MPI communicator 

\item {} 
\sphinxAtStartPar
\sphinxstyleliteralstrong{\sphinxupquote{comm}} \textendash{} MPI communicator for broadcast 

\end{itemize}

\sphinxlineitem{Returns}
\sphinxAtStartPar
1 if the read was short, 0 if read was successful 

\end{description}\end{quote}

\end{fulllineitems}



\bigskip\hrule\bigskip



\subsection{String to number conversions with validity check}
\label{\detokenize{Developer_utils:string-to-number-conversions-with-validity-check}}
\sphinxAtStartPar
These functions should be used to convert strings to numbers. They are
are strongly preferred over C library calls like \sphinxcode{\sphinxupquote{atoi()}} or
\sphinxcode{\sphinxupquote{atof()}} since they check if the \sphinxstylestrong{entire} string is a valid
(floating\sphinxhyphen{}point or integer) number, and will error out instead of
silently returning the result of a partial conversion or zero in cases
where the string is not a valid number.  This behavior improves
detecting typos or issues when processing input files.

\sphinxAtStartPar
Similarly the {\hyperref[\detokenize{Developer_utils:_CPPv4N9LAMMPS_NS5utils7logicalEPKciPKcbP6LAMMPS}]{\sphinxcrossref{\sphinxcode{\sphinxupquote{utils::logical()}}}}} function
will convert a string into a boolean and will only accept certain words.

\sphinxAtStartPar
The \sphinxstyleemphasis{do\_abort} flag should be set to \sphinxcode{\sphinxupquote{true}} in case  this function
is called only on a single MPI rank, as that will then trigger the
a call to \sphinxcode{\sphinxupquote{Error::one()}} for errors instead of \sphinxcode{\sphinxupquote{Error::all()}}
and avoids a “hanging” calculation when run in parallel.

\sphinxAtStartPar
Please also see {\hyperref[\detokenize{Developer_utils:_CPPv4N9LAMMPS_NS5utils10is_integerERKNSt6stringE}]{\sphinxcrossref{\sphinxcode{\sphinxupquote{utils::is\_integer()}}}}}
and {\hyperref[\detokenize{Developer_utils:_CPPv4N9LAMMPS_NS5utils9is_doubleERKNSt6stringE}]{\sphinxcrossref{\sphinxcode{\sphinxupquote{utils::is\_double()}}}}} for testing
strings for compliance without conversion.


\bigskip\hrule\bigskip

\index{utils::numeric (C++ function)@\spxentry{utils::numeric}\spxextra{C++ function}}

\begin{fulllineitems}
\phantomsection\label{\detokenize{Developer_utils:_CPPv4N9LAMMPS_NS5utils7numericEPKciRKNSt6stringEbP6LAMMPS}}\phantomsection\label{\detokenize{Developer_utils:_CPPv3N9LAMMPS_NS5utils7numericEPKciRKNSt6stringEbP6LAMMPS}}\phantomsection\label{\detokenize{Developer_utils:_CPPv2N9LAMMPS_NS5utils7numericEPKciRKNSt6stringEbP6LAMMPS}}\phantomsection\label{\detokenize{Developer_utils:LAMMPS_NS::utils::numeric__cCP.i.ssCR.b.LAMMPSP}}
\pysigstartsignatures
\pysigstartmultiline
\pysiglinewithargsret{\phantomsection\label{\detokenize{Developer_utils:utils_8h_1a68ccfc10e677337b215ccd51363a03aa}}\DUrole{kt}{double}\DUrole{w}{ }\sphinxcode{\sphinxupquote{\DUrole{n}{LAMMPS\_NS}\DUrole{p}{::}\DUrole{n}{utils}\DUrole{p}{::}}}\sphinxbfcode{\sphinxupquote{\DUrole{n}{numeric}}}}{\DUrole{k}{const}\DUrole{w}{ }\DUrole{kt}{char}\DUrole{w}{ }\DUrole{p}{*}\DUrole{n,sig-param}{file}\sphinxparamcomma \DUrole{kt}{int}\DUrole{w}{ }\DUrole{n,sig-param}{line}\sphinxparamcomma \DUrole{k}{const}\DUrole{w}{ }\DUrole{n}{std}\DUrole{p}{::}\DUrole{n}{string}\DUrole{w}{ }\DUrole{p}{\&}\DUrole{n,sig-param}{str}\sphinxparamcomma \DUrole{kt}{bool}\DUrole{w}{ }\DUrole{n,sig-param}{do\_abort}\sphinxparamcomma {\hyperref[\detokenize{Classes_lammps:_CPPv4N9LAMMPS_NS6LAMMPSE}]{\sphinxcrossref{\DUrole{n}{LAMMPS}}}}\DUrole{w}{ }\DUrole{p}{*}\DUrole{n,sig-param}{lmp}}{}
\pysigstopmultiline
\pysigstopsignatures
\sphinxAtStartPar
Convert a string to a floating point number while checking if it is a valid floating point or integer number
\begin{quote}\begin{description}
\sphinxlineitem{Parameters}\begin{itemize}
\item {} 
\sphinxAtStartPar
\sphinxstyleliteralstrong{\sphinxupquote{file}} \textendash{} name of source file for error message 

\item {} 
\sphinxAtStartPar
\sphinxstyleliteralstrong{\sphinxupquote{line}} \textendash{} line number in source file for error message 

\item {} 
\sphinxAtStartPar
\sphinxstyleliteralstrong{\sphinxupquote{str}} \textendash{} string to be converted to number 

\item {} 
\sphinxAtStartPar
\sphinxstyleliteralstrong{\sphinxupquote{do\_abort}} \textendash{} determines whether to call Error::one() or Error::all() 

\item {} 
\sphinxAtStartPar
\sphinxstyleliteralstrong{\sphinxupquote{lmp}} \textendash{} pointer to top\sphinxhyphen{}level {\hyperref[\detokenize{Classes_lammps:classLAMMPS__NS_1_1LAMMPS}]{\sphinxcrossref{\DUrole{std,std-ref}{LAMMPS}}}} class instance 

\end{itemize}

\sphinxlineitem{Returns}
\sphinxAtStartPar
double precision floating point number 

\end{description}\end{quote}

\end{fulllineitems}

\index{utils::numeric (C++ function)@\spxentry{utils::numeric}\spxextra{C++ function}}

\begin{fulllineitems}
\phantomsection\label{\detokenize{Developer_utils:_CPPv4N9LAMMPS_NS5utils7numericEPKciPKcbP6LAMMPS}}\phantomsection\label{\detokenize{Developer_utils:_CPPv3N9LAMMPS_NS5utils7numericEPKciPKcbP6LAMMPS}}\phantomsection\label{\detokenize{Developer_utils:_CPPv2N9LAMMPS_NS5utils7numericEPKciPKcbP6LAMMPS}}\phantomsection\label{\detokenize{Developer_utils:LAMMPS_NS::utils::numeric__cCP.i.cCP.b.LAMMPSP}}
\pysigstartsignatures
\pysigstartmultiline
\pysiglinewithargsret{\phantomsection\label{\detokenize{Developer_utils:utils_8h_1a503fd271ff7cf07725ad0b03b6d7f762}}\DUrole{kt}{double}\DUrole{w}{ }\sphinxcode{\sphinxupquote{\DUrole{n}{LAMMPS\_NS}\DUrole{p}{::}\DUrole{n}{utils}\DUrole{p}{::}}}\sphinxbfcode{\sphinxupquote{\DUrole{n}{numeric}}}}{\DUrole{k}{const}\DUrole{w}{ }\DUrole{kt}{char}\DUrole{w}{ }\DUrole{p}{*}\DUrole{n,sig-param}{file}\sphinxparamcomma \DUrole{kt}{int}\DUrole{w}{ }\DUrole{n,sig-param}{line}\sphinxparamcomma \DUrole{k}{const}\DUrole{w}{ }\DUrole{kt}{char}\DUrole{w}{ }\DUrole{p}{*}\DUrole{n,sig-param}{str}\sphinxparamcomma \DUrole{kt}{bool}\DUrole{w}{ }\DUrole{n,sig-param}{do\_abort}\sphinxparamcomma {\hyperref[\detokenize{Classes_lammps:_CPPv4N9LAMMPS_NS6LAMMPSE}]{\sphinxcrossref{\DUrole{n}{LAMMPS}}}}\DUrole{w}{ }\DUrole{p}{*}\DUrole{n,sig-param}{lmp}}{}
\pysigstopmultiline
\pysigstopsignatures
\sphinxAtStartPar
This is an overloaded member function, provided for convenience. It differs from the above function only in what argument(s) it accepts.
\begin{quote}\begin{description}
\sphinxlineitem{Parameters}\begin{itemize}
\item {} 
\sphinxAtStartPar
\sphinxstyleliteralstrong{\sphinxupquote{file}} \textendash{} name of source file for error message 

\item {} 
\sphinxAtStartPar
\sphinxstyleliteralstrong{\sphinxupquote{line}} \textendash{} line number in source file for error message 

\item {} 
\sphinxAtStartPar
\sphinxstyleliteralstrong{\sphinxupquote{str}} \textendash{} string to be converted to number 

\item {} 
\sphinxAtStartPar
\sphinxstyleliteralstrong{\sphinxupquote{do\_abort}} \textendash{} determines whether to call Error::one() or Error::all() 

\item {} 
\sphinxAtStartPar
\sphinxstyleliteralstrong{\sphinxupquote{lmp}} \textendash{} pointer to top\sphinxhyphen{}level {\hyperref[\detokenize{Classes_lammps:classLAMMPS__NS_1_1LAMMPS}]{\sphinxcrossref{\DUrole{std,std-ref}{LAMMPS}}}} class instance 

\end{itemize}

\sphinxlineitem{Returns}
\sphinxAtStartPar
double precision floating point number 

\end{description}\end{quote}

\end{fulllineitems}

\index{utils::inumeric (C++ function)@\spxentry{utils::inumeric}\spxextra{C++ function}}

\begin{fulllineitems}
\phantomsection\label{\detokenize{Developer_utils:_CPPv4N9LAMMPS_NS5utils8inumericEPKciRKNSt6stringEbP6LAMMPS}}\phantomsection\label{\detokenize{Developer_utils:_CPPv3N9LAMMPS_NS5utils8inumericEPKciRKNSt6stringEbP6LAMMPS}}\phantomsection\label{\detokenize{Developer_utils:_CPPv2N9LAMMPS_NS5utils8inumericEPKciRKNSt6stringEbP6LAMMPS}}\phantomsection\label{\detokenize{Developer_utils:LAMMPS_NS::utils::inumeric__cCP.i.ssCR.b.LAMMPSP}}
\pysigstartsignatures
\pysigstartmultiline
\pysiglinewithargsret{\phantomsection\label{\detokenize{Developer_utils:utils_8h_1a72a1fc411c4fea8f116157d2deb9ecb0}}\DUrole{kt}{int}\DUrole{w}{ }\sphinxcode{\sphinxupquote{\DUrole{n}{LAMMPS\_NS}\DUrole{p}{::}\DUrole{n}{utils}\DUrole{p}{::}}}\sphinxbfcode{\sphinxupquote{\DUrole{n}{inumeric}}}}{\DUrole{k}{const}\DUrole{w}{ }\DUrole{kt}{char}\DUrole{w}{ }\DUrole{p}{*}\DUrole{n,sig-param}{file}\sphinxparamcomma \DUrole{kt}{int}\DUrole{w}{ }\DUrole{n,sig-param}{line}\sphinxparamcomma \DUrole{k}{const}\DUrole{w}{ }\DUrole{n}{std}\DUrole{p}{::}\DUrole{n}{string}\DUrole{w}{ }\DUrole{p}{\&}\DUrole{n,sig-param}{str}\sphinxparamcomma \DUrole{kt}{bool}\DUrole{w}{ }\DUrole{n,sig-param}{do\_abort}\sphinxparamcomma {\hyperref[\detokenize{Classes_lammps:_CPPv4N9LAMMPS_NS6LAMMPSE}]{\sphinxcrossref{\DUrole{n}{LAMMPS}}}}\DUrole{w}{ }\DUrole{p}{*}\DUrole{n,sig-param}{lmp}}{}
\pysigstopmultiline
\pysigstopsignatures
\sphinxAtStartPar
Convert a string to an integer number while checking if it is a valid integer number (regular int)
\begin{quote}\begin{description}
\sphinxlineitem{Parameters}\begin{itemize}
\item {} 
\sphinxAtStartPar
\sphinxstyleliteralstrong{\sphinxupquote{file}} \textendash{} name of source file for error message 

\item {} 
\sphinxAtStartPar
\sphinxstyleliteralstrong{\sphinxupquote{line}} \textendash{} line number in source file for error message 

\item {} 
\sphinxAtStartPar
\sphinxstyleliteralstrong{\sphinxupquote{str}} \textendash{} string to be converted to number 

\item {} 
\sphinxAtStartPar
\sphinxstyleliteralstrong{\sphinxupquote{do\_abort}} \textendash{} determines whether to call Error::one() or Error::all() 

\item {} 
\sphinxAtStartPar
\sphinxstyleliteralstrong{\sphinxupquote{lmp}} \textendash{} pointer to top\sphinxhyphen{}level {\hyperref[\detokenize{Classes_lammps:classLAMMPS__NS_1_1LAMMPS}]{\sphinxcrossref{\DUrole{std,std-ref}{LAMMPS}}}} class instance 

\end{itemize}

\sphinxlineitem{Returns}
\sphinxAtStartPar
integer number (regular int) 

\end{description}\end{quote}

\end{fulllineitems}

\index{utils::inumeric (C++ function)@\spxentry{utils::inumeric}\spxextra{C++ function}}

\begin{fulllineitems}
\phantomsection\label{\detokenize{Developer_utils:_CPPv4N9LAMMPS_NS5utils8inumericEPKciPKcbP6LAMMPS}}\phantomsection\label{\detokenize{Developer_utils:_CPPv3N9LAMMPS_NS5utils8inumericEPKciPKcbP6LAMMPS}}\phantomsection\label{\detokenize{Developer_utils:_CPPv2N9LAMMPS_NS5utils8inumericEPKciPKcbP6LAMMPS}}\phantomsection\label{\detokenize{Developer_utils:LAMMPS_NS::utils::inumeric__cCP.i.cCP.b.LAMMPSP}}
\pysigstartsignatures
\pysigstartmultiline
\pysiglinewithargsret{\phantomsection\label{\detokenize{Developer_utils:utils_8h_1a0fa1caaa9b4c196fb960712e6f2dfe5d}}\DUrole{kt}{int}\DUrole{w}{ }\sphinxcode{\sphinxupquote{\DUrole{n}{LAMMPS\_NS}\DUrole{p}{::}\DUrole{n}{utils}\DUrole{p}{::}}}\sphinxbfcode{\sphinxupquote{\DUrole{n}{inumeric}}}}{\DUrole{k}{const}\DUrole{w}{ }\DUrole{kt}{char}\DUrole{w}{ }\DUrole{p}{*}\DUrole{n,sig-param}{file}\sphinxparamcomma \DUrole{kt}{int}\DUrole{w}{ }\DUrole{n,sig-param}{line}\sphinxparamcomma \DUrole{k}{const}\DUrole{w}{ }\DUrole{kt}{char}\DUrole{w}{ }\DUrole{p}{*}\DUrole{n,sig-param}{str}\sphinxparamcomma \DUrole{kt}{bool}\DUrole{w}{ }\DUrole{n,sig-param}{do\_abort}\sphinxparamcomma {\hyperref[\detokenize{Classes_lammps:_CPPv4N9LAMMPS_NS6LAMMPSE}]{\sphinxcrossref{\DUrole{n}{LAMMPS}}}}\DUrole{w}{ }\DUrole{p}{*}\DUrole{n,sig-param}{lmp}}{}
\pysigstopmultiline
\pysigstopsignatures
\sphinxAtStartPar
This is an overloaded member function, provided for convenience. It differs from the above function only in what argument(s) it accepts.
\begin{quote}\begin{description}
\sphinxlineitem{Parameters}\begin{itemize}
\item {} 
\sphinxAtStartPar
\sphinxstyleliteralstrong{\sphinxupquote{file}} \textendash{} name of source file for error message 

\item {} 
\sphinxAtStartPar
\sphinxstyleliteralstrong{\sphinxupquote{line}} \textendash{} line number in source file for error message 

\item {} 
\sphinxAtStartPar
\sphinxstyleliteralstrong{\sphinxupquote{str}} \textendash{} string to be converted to number 

\item {} 
\sphinxAtStartPar
\sphinxstyleliteralstrong{\sphinxupquote{do\_abort}} \textendash{} determines whether to call Error::one() or Error::all() 

\item {} 
\sphinxAtStartPar
\sphinxstyleliteralstrong{\sphinxupquote{lmp}} \textendash{} pointer to top\sphinxhyphen{}level {\hyperref[\detokenize{Classes_lammps:classLAMMPS__NS_1_1LAMMPS}]{\sphinxcrossref{\DUrole{std,std-ref}{LAMMPS}}}} class instance 

\end{itemize}

\sphinxlineitem{Returns}
\sphinxAtStartPar
double precision floating point number 

\end{description}\end{quote}

\end{fulllineitems}

\index{utils::bnumeric (C++ function)@\spxentry{utils::bnumeric}\spxextra{C++ function}}

\begin{fulllineitems}
\phantomsection\label{\detokenize{Developer_utils:_CPPv4N9LAMMPS_NS5utils8bnumericEPKciRKNSt6stringEbP6LAMMPS}}\phantomsection\label{\detokenize{Developer_utils:_CPPv3N9LAMMPS_NS5utils8bnumericEPKciRKNSt6stringEbP6LAMMPS}}\phantomsection\label{\detokenize{Developer_utils:_CPPv2N9LAMMPS_NS5utils8bnumericEPKciRKNSt6stringEbP6LAMMPS}}\phantomsection\label{\detokenize{Developer_utils:LAMMPS_NS::utils::bnumeric__cCP.i.ssCR.b.LAMMPSP}}
\pysigstartsignatures
\pysigstartmultiline
\pysiglinewithargsret{\phantomsection\label{\detokenize{Developer_utils:utils_8h_1aa2fc5058efa5708a27f9f755f3b1e405}}\DUrole{n}{bigint}\DUrole{w}{ }\sphinxcode{\sphinxupquote{\DUrole{n}{LAMMPS\_NS}\DUrole{p}{::}\DUrole{n}{utils}\DUrole{p}{::}}}\sphinxbfcode{\sphinxupquote{\DUrole{n}{bnumeric}}}}{\DUrole{k}{const}\DUrole{w}{ }\DUrole{kt}{char}\DUrole{w}{ }\DUrole{p}{*}\DUrole{n,sig-param}{file}\sphinxparamcomma \DUrole{kt}{int}\DUrole{w}{ }\DUrole{n,sig-param}{line}\sphinxparamcomma \DUrole{k}{const}\DUrole{w}{ }\DUrole{n}{std}\DUrole{p}{::}\DUrole{n}{string}\DUrole{w}{ }\DUrole{p}{\&}\DUrole{n,sig-param}{str}\sphinxparamcomma \DUrole{kt}{bool}\DUrole{w}{ }\DUrole{n,sig-param}{do\_abort}\sphinxparamcomma {\hyperref[\detokenize{Classes_lammps:_CPPv4N9LAMMPS_NS6LAMMPSE}]{\sphinxcrossref{\DUrole{n}{LAMMPS}}}}\DUrole{w}{ }\DUrole{p}{*}\DUrole{n,sig-param}{lmp}}{}
\pysigstopmultiline
\pysigstopsignatures
\sphinxAtStartPar
Convert a string to an integer number while checking if it is a valid integer number (bigint)
\begin{quote}\begin{description}
\sphinxlineitem{Parameters}\begin{itemize}
\item {} 
\sphinxAtStartPar
\sphinxstyleliteralstrong{\sphinxupquote{file}} \textendash{} name of source file for error message 

\item {} 
\sphinxAtStartPar
\sphinxstyleliteralstrong{\sphinxupquote{line}} \textendash{} line number in source file for error message 

\item {} 
\sphinxAtStartPar
\sphinxstyleliteralstrong{\sphinxupquote{str}} \textendash{} string to be converted to number 

\item {} 
\sphinxAtStartPar
\sphinxstyleliteralstrong{\sphinxupquote{do\_abort}} \textendash{} determines whether to call Error::one() or Error::all() 

\item {} 
\sphinxAtStartPar
\sphinxstyleliteralstrong{\sphinxupquote{lmp}} \textendash{} pointer to top\sphinxhyphen{}level {\hyperref[\detokenize{Classes_lammps:classLAMMPS__NS_1_1LAMMPS}]{\sphinxcrossref{\DUrole{std,std-ref}{LAMMPS}}}} class instance 

\end{itemize}

\sphinxlineitem{Returns}
\sphinxAtStartPar
integer number (bigint) 

\end{description}\end{quote}

\end{fulllineitems}

\index{utils::bnumeric (C++ function)@\spxentry{utils::bnumeric}\spxextra{C++ function}}

\begin{fulllineitems}
\phantomsection\label{\detokenize{Developer_utils:_CPPv4N9LAMMPS_NS5utils8bnumericEPKciPKcbP6LAMMPS}}\phantomsection\label{\detokenize{Developer_utils:_CPPv3N9LAMMPS_NS5utils8bnumericEPKciPKcbP6LAMMPS}}\phantomsection\label{\detokenize{Developer_utils:_CPPv2N9LAMMPS_NS5utils8bnumericEPKciPKcbP6LAMMPS}}\phantomsection\label{\detokenize{Developer_utils:LAMMPS_NS::utils::bnumeric__cCP.i.cCP.b.LAMMPSP}}
\pysigstartsignatures
\pysigstartmultiline
\pysiglinewithargsret{\phantomsection\label{\detokenize{Developer_utils:utils_8h_1af773a99c24a68d2055c9116ae1835f0a}}\DUrole{n}{bigint}\DUrole{w}{ }\sphinxcode{\sphinxupquote{\DUrole{n}{LAMMPS\_NS}\DUrole{p}{::}\DUrole{n}{utils}\DUrole{p}{::}}}\sphinxbfcode{\sphinxupquote{\DUrole{n}{bnumeric}}}}{\DUrole{k}{const}\DUrole{w}{ }\DUrole{kt}{char}\DUrole{w}{ }\DUrole{p}{*}\DUrole{n,sig-param}{file}\sphinxparamcomma \DUrole{kt}{int}\DUrole{w}{ }\DUrole{n,sig-param}{line}\sphinxparamcomma \DUrole{k}{const}\DUrole{w}{ }\DUrole{kt}{char}\DUrole{w}{ }\DUrole{p}{*}\DUrole{n,sig-param}{str}\sphinxparamcomma \DUrole{kt}{bool}\DUrole{w}{ }\DUrole{n,sig-param}{do\_abort}\sphinxparamcomma {\hyperref[\detokenize{Classes_lammps:_CPPv4N9LAMMPS_NS6LAMMPSE}]{\sphinxcrossref{\DUrole{n}{LAMMPS}}}}\DUrole{w}{ }\DUrole{p}{*}\DUrole{n,sig-param}{lmp}}{}
\pysigstopmultiline
\pysigstopsignatures
\sphinxAtStartPar
This is an overloaded member function, provided for convenience. It differs from the above function only in what argument(s) it accepts.
\begin{quote}\begin{description}
\sphinxlineitem{Parameters}\begin{itemize}
\item {} 
\sphinxAtStartPar
\sphinxstyleliteralstrong{\sphinxupquote{file}} \textendash{} name of source file for error message 

\item {} 
\sphinxAtStartPar
\sphinxstyleliteralstrong{\sphinxupquote{line}} \textendash{} line number in source file for error message 

\item {} 
\sphinxAtStartPar
\sphinxstyleliteralstrong{\sphinxupquote{str}} \textendash{} string to be converted to number 

\item {} 
\sphinxAtStartPar
\sphinxstyleliteralstrong{\sphinxupquote{do\_abort}} \textendash{} determines whether to call Error::one() or Error::all() 

\item {} 
\sphinxAtStartPar
\sphinxstyleliteralstrong{\sphinxupquote{lmp}} \textendash{} pointer to top\sphinxhyphen{}level {\hyperref[\detokenize{Classes_lammps:classLAMMPS__NS_1_1LAMMPS}]{\sphinxcrossref{\DUrole{std,std-ref}{LAMMPS}}}} class instance 

\end{itemize}

\sphinxlineitem{Returns}
\sphinxAtStartPar
double precision floating point number 

\end{description}\end{quote}

\end{fulllineitems}

\index{utils::tnumeric (C++ function)@\spxentry{utils::tnumeric}\spxextra{C++ function}}

\begin{fulllineitems}
\phantomsection\label{\detokenize{Developer_utils:_CPPv4N9LAMMPS_NS5utils8tnumericEPKciRKNSt6stringEbP6LAMMPS}}\phantomsection\label{\detokenize{Developer_utils:_CPPv3N9LAMMPS_NS5utils8tnumericEPKciRKNSt6stringEbP6LAMMPS}}\phantomsection\label{\detokenize{Developer_utils:_CPPv2N9LAMMPS_NS5utils8tnumericEPKciRKNSt6stringEbP6LAMMPS}}\phantomsection\label{\detokenize{Developer_utils:LAMMPS_NS::utils::tnumeric__cCP.i.ssCR.b.LAMMPSP}}
\pysigstartsignatures
\pysigstartmultiline
\pysiglinewithargsret{\phantomsection\label{\detokenize{Developer_utils:utils_8h_1a64a163aaea54730784f7c1797369f106}}\DUrole{n}{tagint}\DUrole{w}{ }\sphinxcode{\sphinxupquote{\DUrole{n}{LAMMPS\_NS}\DUrole{p}{::}\DUrole{n}{utils}\DUrole{p}{::}}}\sphinxbfcode{\sphinxupquote{\DUrole{n}{tnumeric}}}}{\DUrole{k}{const}\DUrole{w}{ }\DUrole{kt}{char}\DUrole{w}{ }\DUrole{p}{*}\DUrole{n,sig-param}{file}\sphinxparamcomma \DUrole{kt}{int}\DUrole{w}{ }\DUrole{n,sig-param}{line}\sphinxparamcomma \DUrole{k}{const}\DUrole{w}{ }\DUrole{n}{std}\DUrole{p}{::}\DUrole{n}{string}\DUrole{w}{ }\DUrole{p}{\&}\DUrole{n,sig-param}{str}\sphinxparamcomma \DUrole{kt}{bool}\DUrole{w}{ }\DUrole{n,sig-param}{do\_abort}\sphinxparamcomma {\hyperref[\detokenize{Classes_lammps:_CPPv4N9LAMMPS_NS6LAMMPSE}]{\sphinxcrossref{\DUrole{n}{LAMMPS}}}}\DUrole{w}{ }\DUrole{p}{*}\DUrole{n,sig-param}{lmp}}{}
\pysigstopmultiline
\pysigstopsignatures
\sphinxAtStartPar
Convert a string to an integer number while checking if it is a valid integer number (tagint)
\begin{quote}\begin{description}
\sphinxlineitem{Parameters}\begin{itemize}
\item {} 
\sphinxAtStartPar
\sphinxstyleliteralstrong{\sphinxupquote{file}} \textendash{} name of source file for error message 

\item {} 
\sphinxAtStartPar
\sphinxstyleliteralstrong{\sphinxupquote{line}} \textendash{} line number in source file for error message 

\item {} 
\sphinxAtStartPar
\sphinxstyleliteralstrong{\sphinxupquote{str}} \textendash{} string to be converted to number 

\item {} 
\sphinxAtStartPar
\sphinxstyleliteralstrong{\sphinxupquote{do\_abort}} \textendash{} determines whether to call Error::one() or Error::all() 

\item {} 
\sphinxAtStartPar
\sphinxstyleliteralstrong{\sphinxupquote{lmp}} \textendash{} pointer to top\sphinxhyphen{}level {\hyperref[\detokenize{Classes_lammps:classLAMMPS__NS_1_1LAMMPS}]{\sphinxcrossref{\DUrole{std,std-ref}{LAMMPS}}}} class instance 

\end{itemize}

\sphinxlineitem{Returns}
\sphinxAtStartPar
integer number (tagint) 

\end{description}\end{quote}

\end{fulllineitems}

\index{utils::tnumeric (C++ function)@\spxentry{utils::tnumeric}\spxextra{C++ function}}

\begin{fulllineitems}
\phantomsection\label{\detokenize{Developer_utils:_CPPv4N9LAMMPS_NS5utils8tnumericEPKciPKcbP6LAMMPS}}\phantomsection\label{\detokenize{Developer_utils:_CPPv3N9LAMMPS_NS5utils8tnumericEPKciPKcbP6LAMMPS}}\phantomsection\label{\detokenize{Developer_utils:_CPPv2N9LAMMPS_NS5utils8tnumericEPKciPKcbP6LAMMPS}}\phantomsection\label{\detokenize{Developer_utils:LAMMPS_NS::utils::tnumeric__cCP.i.cCP.b.LAMMPSP}}
\pysigstartsignatures
\pysigstartmultiline
\pysiglinewithargsret{\phantomsection\label{\detokenize{Developer_utils:utils_8h_1a4e8b8a6816b717518c30fd561e6e42fa}}\DUrole{n}{tagint}\DUrole{w}{ }\sphinxcode{\sphinxupquote{\DUrole{n}{LAMMPS\_NS}\DUrole{p}{::}\DUrole{n}{utils}\DUrole{p}{::}}}\sphinxbfcode{\sphinxupquote{\DUrole{n}{tnumeric}}}}{\DUrole{k}{const}\DUrole{w}{ }\DUrole{kt}{char}\DUrole{w}{ }\DUrole{p}{*}\DUrole{n,sig-param}{file}\sphinxparamcomma \DUrole{kt}{int}\DUrole{w}{ }\DUrole{n,sig-param}{line}\sphinxparamcomma \DUrole{k}{const}\DUrole{w}{ }\DUrole{kt}{char}\DUrole{w}{ }\DUrole{p}{*}\DUrole{n,sig-param}{str}\sphinxparamcomma \DUrole{kt}{bool}\DUrole{w}{ }\DUrole{n,sig-param}{do\_abort}\sphinxparamcomma {\hyperref[\detokenize{Classes_lammps:_CPPv4N9LAMMPS_NS6LAMMPSE}]{\sphinxcrossref{\DUrole{n}{LAMMPS}}}}\DUrole{w}{ }\DUrole{p}{*}\DUrole{n,sig-param}{lmp}}{}
\pysigstopmultiline
\pysigstopsignatures
\sphinxAtStartPar
This is an overloaded member function, provided for convenience. It differs from the above function only in what argument(s) it accepts.
\begin{quote}\begin{description}
\sphinxlineitem{Parameters}\begin{itemize}
\item {} 
\sphinxAtStartPar
\sphinxstyleliteralstrong{\sphinxupquote{file}} \textendash{} name of source file for error message 

\item {} 
\sphinxAtStartPar
\sphinxstyleliteralstrong{\sphinxupquote{line}} \textendash{} line number in source file for error message 

\item {} 
\sphinxAtStartPar
\sphinxstyleliteralstrong{\sphinxupquote{str}} \textendash{} string to be converted to number 

\item {} 
\sphinxAtStartPar
\sphinxstyleliteralstrong{\sphinxupquote{do\_abort}} \textendash{} determines whether to call Error::one() or Error::all() 

\item {} 
\sphinxAtStartPar
\sphinxstyleliteralstrong{\sphinxupquote{lmp}} \textendash{} pointer to top\sphinxhyphen{}level {\hyperref[\detokenize{Classes_lammps:classLAMMPS__NS_1_1LAMMPS}]{\sphinxcrossref{\DUrole{std,std-ref}{LAMMPS}}}} class instance 

\end{itemize}

\sphinxlineitem{Returns}
\sphinxAtStartPar
double precision floating point number 

\end{description}\end{quote}

\end{fulllineitems}

\index{utils::logical (C++ function)@\spxentry{utils::logical}\spxextra{C++ function}}

\begin{fulllineitems}
\phantomsection\label{\detokenize{Developer_utils:_CPPv4N9LAMMPS_NS5utils7logicalEPKciRKNSt6stringEbP6LAMMPS}}\phantomsection\label{\detokenize{Developer_utils:_CPPv3N9LAMMPS_NS5utils7logicalEPKciRKNSt6stringEbP6LAMMPS}}\phantomsection\label{\detokenize{Developer_utils:_CPPv2N9LAMMPS_NS5utils7logicalEPKciRKNSt6stringEbP6LAMMPS}}\phantomsection\label{\detokenize{Developer_utils:LAMMPS_NS::utils::logical__cCP.i.ssCR.b.LAMMPSP}}
\pysigstartsignatures
\pysigstartmultiline
\pysiglinewithargsret{\phantomsection\label{\detokenize{Developer_utils:utils_8h_1a57df85d91e1593d0893945852f26eb12}}\DUrole{kt}{int}\DUrole{w}{ }\sphinxcode{\sphinxupquote{\DUrole{n}{LAMMPS\_NS}\DUrole{p}{::}\DUrole{n}{utils}\DUrole{p}{::}}}\sphinxbfcode{\sphinxupquote{\DUrole{n}{logical}}}}{\DUrole{k}{const}\DUrole{w}{ }\DUrole{kt}{char}\DUrole{w}{ }\DUrole{p}{*}\DUrole{n,sig-param}{file}\sphinxparamcomma \DUrole{kt}{int}\DUrole{w}{ }\DUrole{n,sig-param}{line}\sphinxparamcomma \DUrole{k}{const}\DUrole{w}{ }\DUrole{n}{std}\DUrole{p}{::}\DUrole{n}{string}\DUrole{w}{ }\DUrole{p}{\&}\DUrole{n,sig-param}{str}\sphinxparamcomma \DUrole{kt}{bool}\DUrole{w}{ }\DUrole{n,sig-param}{do\_abort}\sphinxparamcomma {\hyperref[\detokenize{Classes_lammps:_CPPv4N9LAMMPS_NS6LAMMPSE}]{\sphinxcrossref{\DUrole{n}{LAMMPS}}}}\DUrole{w}{ }\DUrole{p}{*}\DUrole{n,sig-param}{lmp}}{}
\pysigstopmultiline
\pysigstopsignatures
\sphinxAtStartPar
Convert a string to a boolean while checking whether it is a valid boolean term. Valid terms are ‘yes’, ‘no’, ‘true’, ‘false’, ‘on’, ‘off’, and ‘1’, ‘0’. Only lower case is accepted.
\begin{quote}\begin{description}
\sphinxlineitem{Parameters}\begin{itemize}
\item {} 
\sphinxAtStartPar
\sphinxstyleliteralstrong{\sphinxupquote{file}} \textendash{} name of source file for error message 

\item {} 
\sphinxAtStartPar
\sphinxstyleliteralstrong{\sphinxupquote{line}} \textendash{} line number in source file for error message 

\item {} 
\sphinxAtStartPar
\sphinxstyleliteralstrong{\sphinxupquote{str}} \textendash{} string to be converted to logical 

\item {} 
\sphinxAtStartPar
\sphinxstyleliteralstrong{\sphinxupquote{do\_abort}} \textendash{} determines whether to call Error::one() or Error::all() 

\item {} 
\sphinxAtStartPar
\sphinxstyleliteralstrong{\sphinxupquote{lmp}} \textendash{} pointer to top\sphinxhyphen{}level {\hyperref[\detokenize{Classes_lammps:classLAMMPS__NS_1_1LAMMPS}]{\sphinxcrossref{\DUrole{std,std-ref}{LAMMPS}}}} class instance 

\end{itemize}

\sphinxlineitem{Returns}
\sphinxAtStartPar
1 if string resolves to “true”, otherwise 0 

\end{description}\end{quote}

\end{fulllineitems}

\index{utils::logical (C++ function)@\spxentry{utils::logical}\spxextra{C++ function}}

\begin{fulllineitems}
\phantomsection\label{\detokenize{Developer_utils:_CPPv4N9LAMMPS_NS5utils7logicalEPKciPKcbP6LAMMPS}}\phantomsection\label{\detokenize{Developer_utils:_CPPv3N9LAMMPS_NS5utils7logicalEPKciPKcbP6LAMMPS}}\phantomsection\label{\detokenize{Developer_utils:_CPPv2N9LAMMPS_NS5utils7logicalEPKciPKcbP6LAMMPS}}\phantomsection\label{\detokenize{Developer_utils:LAMMPS_NS::utils::logical__cCP.i.cCP.b.LAMMPSP}}
\pysigstartsignatures
\pysigstartmultiline
\pysiglinewithargsret{\phantomsection\label{\detokenize{Developer_utils:utils_8h_1ae01f31ea9af04975dcdb062b3b340c2e}}\DUrole{kt}{int}\DUrole{w}{ }\sphinxcode{\sphinxupquote{\DUrole{n}{LAMMPS\_NS}\DUrole{p}{::}\DUrole{n}{utils}\DUrole{p}{::}}}\sphinxbfcode{\sphinxupquote{\DUrole{n}{logical}}}}{\DUrole{k}{const}\DUrole{w}{ }\DUrole{kt}{char}\DUrole{w}{ }\DUrole{p}{*}\DUrole{n,sig-param}{file}\sphinxparamcomma \DUrole{kt}{int}\DUrole{w}{ }\DUrole{n,sig-param}{line}\sphinxparamcomma \DUrole{k}{const}\DUrole{w}{ }\DUrole{kt}{char}\DUrole{w}{ }\DUrole{p}{*}\DUrole{n,sig-param}{str}\sphinxparamcomma \DUrole{kt}{bool}\DUrole{w}{ }\DUrole{n,sig-param}{do\_abort}\sphinxparamcomma {\hyperref[\detokenize{Classes_lammps:_CPPv4N9LAMMPS_NS6LAMMPSE}]{\sphinxcrossref{\DUrole{n}{LAMMPS}}}}\DUrole{w}{ }\DUrole{p}{*}\DUrole{n,sig-param}{lmp}}{}
\pysigstopmultiline
\pysigstopsignatures
\sphinxAtStartPar
This is an overloaded member function, provided for convenience. It differs from the above function only in what argument(s) it accepts.
\begin{quote}\begin{description}
\sphinxlineitem{Parameters}\begin{itemize}
\item {} 
\sphinxAtStartPar
\sphinxstyleliteralstrong{\sphinxupquote{file}} \textendash{} name of source file for error message 

\item {} 
\sphinxAtStartPar
\sphinxstyleliteralstrong{\sphinxupquote{line}} \textendash{} line number in source file for error message 

\item {} 
\sphinxAtStartPar
\sphinxstyleliteralstrong{\sphinxupquote{str}} \textendash{} string to be converted to logical 

\item {} 
\sphinxAtStartPar
\sphinxstyleliteralstrong{\sphinxupquote{do\_abort}} \textendash{} determines whether to call Error::one() or Error::all() 

\item {} 
\sphinxAtStartPar
\sphinxstyleliteralstrong{\sphinxupquote{lmp}} \textendash{} pointer to top\sphinxhyphen{}level {\hyperref[\detokenize{Classes_lammps:classLAMMPS__NS_1_1LAMMPS}]{\sphinxcrossref{\DUrole{std,std-ref}{LAMMPS}}}} class instance 

\end{itemize}

\sphinxlineitem{Returns}
\sphinxAtStartPar
1 if string resolves to “true”, otherwise 0 

\end{description}\end{quote}

\end{fulllineitems}



\subsection{String processing}
\label{\detokenize{Developer_utils:string-processing}}
\sphinxAtStartPar
The following are functions to help with processing strings
and parsing files or arguments.


\bigskip\hrule\bigskip

\index{utils::strdup (C++ function)@\spxentry{utils::strdup}\spxextra{C++ function}}

\begin{fulllineitems}
\phantomsection\label{\detokenize{Developer_utils:_CPPv4N9LAMMPS_NS5utils6strdupERKNSt6stringE}}\phantomsection\label{\detokenize{Developer_utils:_CPPv3N9LAMMPS_NS5utils6strdupERKNSt6stringE}}\phantomsection\label{\detokenize{Developer_utils:_CPPv2N9LAMMPS_NS5utils6strdupERKNSt6stringE}}\phantomsection\label{\detokenize{Developer_utils:LAMMPS_NS::utils::strdup__ssCR}}
\pysigstartsignatures
\pysigstartmultiline
\pysiglinewithargsret{\phantomsection\label{\detokenize{Developer_utils:utils_8h_1a2a9f5a378d6155ae527b04b419fbd3e5}}\DUrole{kt}{char}\DUrole{w}{ }\DUrole{p}{*}\sphinxcode{\sphinxupquote{\DUrole{n}{LAMMPS\_NS}\DUrole{p}{::}\DUrole{n}{utils}\DUrole{p}{::}}}\sphinxbfcode{\sphinxupquote{\DUrole{n}{strdup}}}}{\DUrole{k}{const}\DUrole{w}{ }\DUrole{n}{std}\DUrole{p}{::}\DUrole{n}{string}\DUrole{w}{ }\DUrole{p}{\&}\DUrole{n,sig-param}{text}}{}
\pysigstopmultiline
\pysigstopsignatures
\sphinxAtStartPar
Make C\sphinxhyphen{}style copy of string in new storage

\sphinxAtStartPar
This allocates a storage buffer and copies the C\sphinxhyphen{}style or C++ style string into it. The buffer is allocated with “new” and thus needs to be deallocated with “delete{[}{]}”.
\begin{quote}\begin{description}
\sphinxlineitem{Parameters}
\sphinxAtStartPar
\sphinxstyleliteralstrong{\sphinxupquote{text}} \textendash{} string that should be copied 

\sphinxlineitem{Returns}
\sphinxAtStartPar
new buffer with copy of string 

\end{description}\end{quote}

\end{fulllineitems}

\index{utils::lowercase (C++ function)@\spxentry{utils::lowercase}\spxextra{C++ function}}

\begin{fulllineitems}
\phantomsection\label{\detokenize{Developer_utils:_CPPv4N9LAMMPS_NS5utils9lowercaseERKNSt6stringE}}\phantomsection\label{\detokenize{Developer_utils:_CPPv3N9LAMMPS_NS5utils9lowercaseERKNSt6stringE}}\phantomsection\label{\detokenize{Developer_utils:_CPPv2N9LAMMPS_NS5utils9lowercaseERKNSt6stringE}}\phantomsection\label{\detokenize{Developer_utils:LAMMPS_NS::utils::lowercase__ssCR}}
\pysigstartsignatures
\pysigstartmultiline
\pysiglinewithargsret{\phantomsection\label{\detokenize{Developer_utils:utils_8h_1aa27f09f6c641f8b4c2de6903006d10a2}}\DUrole{n}{std}\DUrole{p}{::}\DUrole{n}{string}\DUrole{w}{ }\sphinxcode{\sphinxupquote{\DUrole{n}{LAMMPS\_NS}\DUrole{p}{::}\DUrole{n}{utils}\DUrole{p}{::}}}\sphinxbfcode{\sphinxupquote{\DUrole{n}{lowercase}}}}{\DUrole{k}{const}\DUrole{w}{ }\DUrole{n}{std}\DUrole{p}{::}\DUrole{n}{string}\DUrole{w}{ }\DUrole{p}{\&}\DUrole{n,sig-param}{line}}{}
\pysigstopmultiline
\pysigstopsignatures
\sphinxAtStartPar
Convert string to lowercase
\begin{quote}\begin{description}
\sphinxlineitem{Parameters}
\sphinxAtStartPar
\sphinxstyleliteralstrong{\sphinxupquote{line}} \textendash{} string that should be converted 

\sphinxlineitem{Returns}
\sphinxAtStartPar
new string with all lowercase characters 

\end{description}\end{quote}

\end{fulllineitems}

\index{utils::uppercase (C++ function)@\spxentry{utils::uppercase}\spxextra{C++ function}}

\begin{fulllineitems}
\phantomsection\label{\detokenize{Developer_utils:_CPPv4N9LAMMPS_NS5utils9uppercaseERKNSt6stringE}}\phantomsection\label{\detokenize{Developer_utils:_CPPv3N9LAMMPS_NS5utils9uppercaseERKNSt6stringE}}\phantomsection\label{\detokenize{Developer_utils:_CPPv2N9LAMMPS_NS5utils9uppercaseERKNSt6stringE}}\phantomsection\label{\detokenize{Developer_utils:LAMMPS_NS::utils::uppercase__ssCR}}
\pysigstartsignatures
\pysigstartmultiline
\pysiglinewithargsret{\phantomsection\label{\detokenize{Developer_utils:utils_8h_1a985e66e0bdcd8f73d1ee6a477f9ddc2a}}\DUrole{n}{std}\DUrole{p}{::}\DUrole{n}{string}\DUrole{w}{ }\sphinxcode{\sphinxupquote{\DUrole{n}{LAMMPS\_NS}\DUrole{p}{::}\DUrole{n}{utils}\DUrole{p}{::}}}\sphinxbfcode{\sphinxupquote{\DUrole{n}{uppercase}}}}{\DUrole{k}{const}\DUrole{w}{ }\DUrole{n}{std}\DUrole{p}{::}\DUrole{n}{string}\DUrole{w}{ }\DUrole{p}{\&}\DUrole{n,sig-param}{line}}{}
\pysigstopmultiline
\pysigstopsignatures
\sphinxAtStartPar
Convert string to uppercase
\begin{quote}\begin{description}
\sphinxlineitem{Parameters}
\sphinxAtStartPar
\sphinxstyleliteralstrong{\sphinxupquote{line}} \textendash{} string that should be converted 

\sphinxlineitem{Returns}
\sphinxAtStartPar
new string with all uppercase characters 

\end{description}\end{quote}

\end{fulllineitems}

\index{utils::trim (C++ function)@\spxentry{utils::trim}\spxextra{C++ function}}

\begin{fulllineitems}
\phantomsection\label{\detokenize{Developer_utils:_CPPv4N9LAMMPS_NS5utils4trimERKNSt6stringE}}\phantomsection\label{\detokenize{Developer_utils:_CPPv3N9LAMMPS_NS5utils4trimERKNSt6stringE}}\phantomsection\label{\detokenize{Developer_utils:_CPPv2N9LAMMPS_NS5utils4trimERKNSt6stringE}}\phantomsection\label{\detokenize{Developer_utils:LAMMPS_NS::utils::trim__ssCR}}
\pysigstartsignatures
\pysigstartmultiline
\pysiglinewithargsret{\phantomsection\label{\detokenize{Developer_utils:utils_8h_1a2eaed9f7998793904b0953fe0d9bbf6a}}\DUrole{n}{std}\DUrole{p}{::}\DUrole{n}{string}\DUrole{w}{ }\sphinxcode{\sphinxupquote{\DUrole{n}{LAMMPS\_NS}\DUrole{p}{::}\DUrole{n}{utils}\DUrole{p}{::}}}\sphinxbfcode{\sphinxupquote{\DUrole{n}{trim}}}}{\DUrole{k}{const}\DUrole{w}{ }\DUrole{n}{std}\DUrole{p}{::}\DUrole{n}{string}\DUrole{w}{ }\DUrole{p}{\&}\DUrole{n,sig-param}{line}}{}
\pysigstopmultiline
\pysigstopsignatures
\sphinxAtStartPar
Trim leading and trailing whitespace. Like TRIM() in Fortran.
\begin{quote}\begin{description}
\sphinxlineitem{Parameters}
\sphinxAtStartPar
\sphinxstyleliteralstrong{\sphinxupquote{line}} \textendash{} string that should be trimmed 

\sphinxlineitem{Returns}
\sphinxAtStartPar
new string without whitespace (string) 

\end{description}\end{quote}

\end{fulllineitems}

\index{utils::trim\_comment (C++ function)@\spxentry{utils::trim\_comment}\spxextra{C++ function}}

\begin{fulllineitems}
\phantomsection\label{\detokenize{Developer_utils:_CPPv4N9LAMMPS_NS5utils12trim_commentERKNSt6stringE}}\phantomsection\label{\detokenize{Developer_utils:_CPPv3N9LAMMPS_NS5utils12trim_commentERKNSt6stringE}}\phantomsection\label{\detokenize{Developer_utils:_CPPv2N9LAMMPS_NS5utils12trim_commentERKNSt6stringE}}\phantomsection\label{\detokenize{Developer_utils:LAMMPS_NS::utils::trim_comment__ssCR}}
\pysigstartsignatures
\pysigstartmultiline
\pysiglinewithargsret{\phantomsection\label{\detokenize{Developer_utils:utils_8h_1ae197fb8181a875095b3dbc0f1e2cdf2d}}\DUrole{n}{std}\DUrole{p}{::}\DUrole{n}{string}\DUrole{w}{ }\sphinxcode{\sphinxupquote{\DUrole{n}{LAMMPS\_NS}\DUrole{p}{::}\DUrole{n}{utils}\DUrole{p}{::}}}\sphinxbfcode{\sphinxupquote{\DUrole{n}{trim\_comment}}}}{\DUrole{k}{const}\DUrole{w}{ }\DUrole{n}{std}\DUrole{p}{::}\DUrole{n}{string}\DUrole{w}{ }\DUrole{p}{\&}\DUrole{n,sig-param}{line}}{}
\pysigstopmultiline
\pysigstopsignatures
\sphinxAtStartPar
Return string with anything from the first ‘\#’ character onward removed
\begin{quote}\begin{description}
\sphinxlineitem{Parameters}
\sphinxAtStartPar
\sphinxstyleliteralstrong{\sphinxupquote{line}} \textendash{} string that should be trimmed 

\sphinxlineitem{Returns}
\sphinxAtStartPar
new string without comment (string) 

\end{description}\end{quote}

\end{fulllineitems}

\index{utils::strip\_style\_suffix (C++ function)@\spxentry{utils::strip\_style\_suffix}\spxextra{C++ function}}

\begin{fulllineitems}
\phantomsection\label{\detokenize{Developer_utils:_CPPv4N9LAMMPS_NS5utils18strip_style_suffixERKNSt6stringEP6LAMMPS}}\phantomsection\label{\detokenize{Developer_utils:_CPPv3N9LAMMPS_NS5utils18strip_style_suffixERKNSt6stringEP6LAMMPS}}\phantomsection\label{\detokenize{Developer_utils:_CPPv2N9LAMMPS_NS5utils18strip_style_suffixERKNSt6stringEP6LAMMPS}}\phantomsection\label{\detokenize{Developer_utils:LAMMPS_NS::utils::strip_style_suffix__ssCR.LAMMPSP}}
\pysigstartsignatures
\pysigstartmultiline
\pysiglinewithargsret{\phantomsection\label{\detokenize{Developer_utils:utils_8h_1ab59c3628bad8784f2df865818395686a}}\DUrole{n}{std}\DUrole{p}{::}\DUrole{n}{string}\DUrole{w}{ }\sphinxcode{\sphinxupquote{\DUrole{n}{LAMMPS\_NS}\DUrole{p}{::}\DUrole{n}{utils}\DUrole{p}{::}}}\sphinxbfcode{\sphinxupquote{\DUrole{n}{strip\_style\_suffix}}}}{\DUrole{k}{const}\DUrole{w}{ }\DUrole{n}{std}\DUrole{p}{::}\DUrole{n}{string}\DUrole{w}{ }\DUrole{p}{\&}\DUrole{n,sig-param}{style}\sphinxparamcomma {\hyperref[\detokenize{Classes_lammps:_CPPv4N9LAMMPS_NS6LAMMPSE}]{\sphinxcrossref{\DUrole{n}{LAMMPS}}}}\DUrole{w}{ }\DUrole{p}{*}\DUrole{n,sig-param}{lmp}}{}
\pysigstopmultiline
\pysigstopsignatures
\sphinxAtStartPar
Remove style suffix from string if suffix flag is active

\sphinxAtStartPar

\sphinxAtStartPar
This will try to undo the effect from using the {\hyperref[\detokenize{suffix::doc}]{\sphinxcrossref{\DUrole{doc}{suffix command}}}}
or the \sphinxstyleemphasis{\sphinxhyphen{}suffix/\sphinxhyphen{}sf} command line flag and return correspondingly modified string.
 
\begin{sphinxVerbatim}[commandchars=\\\{\}]
 \PYGZbs{}param style  string of style name
 \PYGZbs{}param lmp    pointer to the LAMMPS class (has suffix\PYGZus{}flag and suffix strings)
 \PYGZbs{}return  processed string
\end{sphinxVerbatim}


\end{fulllineitems}

\index{utils::star\_subst (C++ function)@\spxentry{utils::star\_subst}\spxextra{C++ function}}

\begin{fulllineitems}
\phantomsection\label{\detokenize{Developer_utils:_CPPv4N9LAMMPS_NS5utils10star_substERKNSt6stringE6biginti}}\phantomsection\label{\detokenize{Developer_utils:_CPPv3N9LAMMPS_NS5utils10star_substERKNSt6stringE6biginti}}\phantomsection\label{\detokenize{Developer_utils:_CPPv2N9LAMMPS_NS5utils10star_substERKNSt6stringE6biginti}}\phantomsection\label{\detokenize{Developer_utils:LAMMPS_NS::utils::star_subst__ssCR.bigint.i}}
\pysigstartsignatures
\pysigstartmultiline
\pysiglinewithargsret{\phantomsection\label{\detokenize{Developer_utils:utils_8h_1a36e52b4e7c0e6b89dded52d86413c11d}}\DUrole{n}{std}\DUrole{p}{::}\DUrole{n}{string}\DUrole{w}{ }\sphinxcode{\sphinxupquote{\DUrole{n}{LAMMPS\_NS}\DUrole{p}{::}\DUrole{n}{utils}\DUrole{p}{::}}}\sphinxbfcode{\sphinxupquote{\DUrole{n}{star\_subst}}}}{\DUrole{k}{const}\DUrole{w}{ }\DUrole{n}{std}\DUrole{p}{::}\DUrole{n}{string}\DUrole{w}{ }\DUrole{p}{\&}\DUrole{n,sig-param}{name}\sphinxparamcomma \DUrole{n}{bigint}\DUrole{w}{ }\DUrole{n,sig-param}{step}\sphinxparamcomma \DUrole{kt}{int}\DUrole{w}{ }\DUrole{n,sig-param}{pad}}{}
\pysigstopmultiline
\pysigstopsignatures
\sphinxAtStartPar
Replace first ‘*’ character in a string with a number, optionally zero\sphinxhyphen{}padded

\sphinxAtStartPar
If there is no ‘*’ character in the string, return the original string. If the number requires more characters than the value of the \sphinxstyleemphasis{pad} argument, do not add zeros; otherwise add as many zeroes as needed to the left to make the the number representation \sphinxstyleemphasis{pad} characters wide.
\begin{quote}\begin{description}
\sphinxlineitem{Parameters}\begin{itemize}
\item {} 
\sphinxAtStartPar
\sphinxstyleliteralstrong{\sphinxupquote{name}} \textendash{} string with file containing a ‘*’ (or not) 

\item {} 
\sphinxAtStartPar
\sphinxstyleliteralstrong{\sphinxupquote{step}} \textendash{} step number to replace the (first) ‘*’ 

\item {} 
\sphinxAtStartPar
\sphinxstyleliteralstrong{\sphinxupquote{pad}} \textendash{} zero\sphinxhyphen{}padding (may be zero) 

\end{itemize}

\sphinxlineitem{Returns}
\sphinxAtStartPar
processed string 

\end{description}\end{quote}

\end{fulllineitems}

\index{utils::has\_utf8 (C++ function)@\spxentry{utils::has\_utf8}\spxextra{C++ function}}

\begin{fulllineitems}
\phantomsection\label{\detokenize{Developer_utils:_CPPv4N9LAMMPS_NS5utils8has_utf8ERKNSt6stringE}}\phantomsection\label{\detokenize{Developer_utils:_CPPv3N9LAMMPS_NS5utils8has_utf8ERKNSt6stringE}}\phantomsection\label{\detokenize{Developer_utils:_CPPv2N9LAMMPS_NS5utils8has_utf8ERKNSt6stringE}}\phantomsection\label{\detokenize{Developer_utils:LAMMPS_NS::utils::has_utf8__ssCR}}
\pysigstartsignatures
\pysigstartmultiline
\pysiglinewithargsret{\phantomsection\label{\detokenize{Developer_utils:utils_8h_1a5cad02698eb9150f24cef34c476681cd}}\DUrole{k}{inline}\DUrole{w}{ }\DUrole{kt}{bool}\DUrole{w}{ }\sphinxcode{\sphinxupquote{\DUrole{n}{LAMMPS\_NS}\DUrole{p}{::}\DUrole{n}{utils}\DUrole{p}{::}}}\sphinxbfcode{\sphinxupquote{\DUrole{n}{has\_utf8}}}}{\DUrole{k}{const}\DUrole{w}{ }\DUrole{n}{std}\DUrole{p}{::}\DUrole{n}{string}\DUrole{w}{ }\DUrole{p}{\&}\DUrole{n,sig-param}{line}}{}
\pysigstopmultiline
\pysigstopsignatures
\sphinxAtStartPar
Check if a string will likely have UTF\sphinxhyphen{}8 encoded characters

\sphinxAtStartPar
UTF\sphinxhyphen{}8 uses the 7\sphinxhyphen{}bit standard ASCII table for the first 127 characters and all other characters are encoded as multiple bytes. For the multi\sphinxhyphen{}byte characters the first byte has either the highest two, three, or four bits set followed by a zero bit and followed by one, two, or three more bytes, respectively, where the highest bit is set and the second highest bit set to 0. The remaining bits combined are the character code, which is thus limited to 21\sphinxhyphen{}bits.

\sphinxAtStartPar
For the sake of efficiency this test only checks if a character in the string has the highest bit set and thus is very likely an UTF\sphinxhyphen{}8 character. It will not be able to tell this this is a valid UTF\sphinxhyphen{}8 character or whether it is a 2\sphinxhyphen{}byte, 3\sphinxhyphen{}byte, or 4\sphinxhyphen{}byte character.

\sphinxAtStartPar

\sphinxAtStartPar
\begin{description}
\sphinxlineitem{\sphinxstyleemphasis{See also}}
\sphinxAtStartPar
{\hyperref[\detokenize{Developer_utils:_CPPv4N9LAMMPS_NS5utils10utf8_substERKNSt6stringE}]{\sphinxcrossref{\sphinxcode{\sphinxupquote{utils::utf8\_subst()}}}}}

\end{description}

 
\begin{quote}\begin{description}
\sphinxlineitem{Parameters}
\sphinxAtStartPar
\sphinxstyleliteralstrong{\sphinxupquote{line}} \textendash{} string that should be checked 

\sphinxlineitem{Returns}
\sphinxAtStartPar
true if string contains UTF\sphinxhyphen{}8 encoded characters (bool) 

\end{description}\end{quote}

\end{fulllineitems}

\index{utils::utf8\_subst (C++ function)@\spxentry{utils::utf8\_subst}\spxextra{C++ function}}

\begin{fulllineitems}
\phantomsection\label{\detokenize{Developer_utils:_CPPv4N9LAMMPS_NS5utils10utf8_substERKNSt6stringE}}\phantomsection\label{\detokenize{Developer_utils:_CPPv3N9LAMMPS_NS5utils10utf8_substERKNSt6stringE}}\phantomsection\label{\detokenize{Developer_utils:_CPPv2N9LAMMPS_NS5utils10utf8_substERKNSt6stringE}}\phantomsection\label{\detokenize{Developer_utils:LAMMPS_NS::utils::utf8_subst__ssCR}}
\pysigstartsignatures
\pysigstartmultiline
\pysiglinewithargsret{\phantomsection\label{\detokenize{Developer_utils:utils_8h_1a6ab67644ed865096c60e1fd456a52bd4}}\DUrole{n}{std}\DUrole{p}{::}\DUrole{n}{string}\DUrole{w}{ }\sphinxcode{\sphinxupquote{\DUrole{n}{LAMMPS\_NS}\DUrole{p}{::}\DUrole{n}{utils}\DUrole{p}{::}}}\sphinxbfcode{\sphinxupquote{\DUrole{n}{utf8\_subst}}}}{\DUrole{k}{const}\DUrole{w}{ }\DUrole{n}{std}\DUrole{p}{::}\DUrole{n}{string}\DUrole{w}{ }\DUrole{p}{\&}\DUrole{n,sig-param}{line}}{}
\pysigstopmultiline
\pysigstopsignatures
\sphinxAtStartPar
Replace known UTF\sphinxhyphen{}8 characters with ASCII equivalents

\sphinxAtStartPar

\sphinxAtStartPar
\begin{description}
\sphinxlineitem{\sphinxstyleemphasis{See also}}
\sphinxAtStartPar
{\hyperref[\detokenize{Developer_utils:_CPPv4N9LAMMPS_NS5utils8has_utf8ERKNSt6stringE}]{\sphinxcrossref{\sphinxcode{\sphinxupquote{utils::has\_utf8()}}}}}

\end{description}

 
\begin{quote}\begin{description}
\sphinxlineitem{Parameters}
\sphinxAtStartPar
\sphinxstyleliteralstrong{\sphinxupquote{line}} \textendash{} string that should be converted 

\sphinxlineitem{Returns}
\sphinxAtStartPar
new string with ascii replacements (string) 

\end{description}\end{quote}

\end{fulllineitems}

\index{utils::count\_words (C++ function)@\spxentry{utils::count\_words}\spxextra{C++ function}}

\begin{fulllineitems}
\phantomsection\label{\detokenize{Developer_utils:_CPPv4N9LAMMPS_NS5utils11count_wordsEPKc}}\phantomsection\label{\detokenize{Developer_utils:_CPPv3N9LAMMPS_NS5utils11count_wordsEPKc}}\phantomsection\label{\detokenize{Developer_utils:_CPPv2N9LAMMPS_NS5utils11count_wordsEPKc}}\phantomsection\label{\detokenize{Developer_utils:LAMMPS_NS::utils::count_words__cCP}}
\pysigstartsignatures
\pysigstartmultiline
\pysiglinewithargsret{\phantomsection\label{\detokenize{Developer_utils:utils_8h_1a8cda2fde7eb5a65ed3f071ed17f42054}}\DUrole{n}{size\_t}\DUrole{w}{ }\sphinxcode{\sphinxupquote{\DUrole{n}{LAMMPS\_NS}\DUrole{p}{::}\DUrole{n}{utils}\DUrole{p}{::}}}\sphinxbfcode{\sphinxupquote{\DUrole{n}{count\_words}}}}{\DUrole{k}{const}\DUrole{w}{ }\DUrole{kt}{char}\DUrole{w}{ }\DUrole{p}{*}\DUrole{n,sig-param}{text}}{}
\pysigstopmultiline
\pysigstopsignatures
\sphinxAtStartPar
Count words in C\sphinxhyphen{}string, ignore any whitespace matching “ \textbackslash{}t\textbackslash{}r\textbackslash{}n\textbackslash{}f”
\begin{quote}\begin{description}
\sphinxlineitem{Parameters}
\sphinxAtStartPar
\sphinxstyleliteralstrong{\sphinxupquote{text}} \textendash{} string that should be searched 

\sphinxlineitem{Returns}
\sphinxAtStartPar
number of words found 

\end{description}\end{quote}

\end{fulllineitems}

\index{utils::count\_words (C++ function)@\spxentry{utils::count\_words}\spxextra{C++ function}}

\begin{fulllineitems}
\phantomsection\label{\detokenize{Developer_utils:_CPPv4N9LAMMPS_NS5utils11count_wordsERKNSt6stringE}}\phantomsection\label{\detokenize{Developer_utils:_CPPv3N9LAMMPS_NS5utils11count_wordsERKNSt6stringE}}\phantomsection\label{\detokenize{Developer_utils:_CPPv2N9LAMMPS_NS5utils11count_wordsERKNSt6stringE}}\phantomsection\label{\detokenize{Developer_utils:LAMMPS_NS::utils::count_words__ssCR}}
\pysigstartsignatures
\pysigstartmultiline
\pysiglinewithargsret{\phantomsection\label{\detokenize{Developer_utils:utils_8h_1ab469ccad212a812ee2f67727307509ee}}\DUrole{n}{size\_t}\DUrole{w}{ }\sphinxcode{\sphinxupquote{\DUrole{n}{LAMMPS\_NS}\DUrole{p}{::}\DUrole{n}{utils}\DUrole{p}{::}}}\sphinxbfcode{\sphinxupquote{\DUrole{n}{count\_words}}}}{\DUrole{k}{const}\DUrole{w}{ }\DUrole{n}{std}\DUrole{p}{::}\DUrole{n}{string}\DUrole{w}{ }\DUrole{p}{\&}\DUrole{n,sig-param}{text}}{}
\pysigstopmultiline
\pysigstopsignatures
\sphinxAtStartPar
Count words in string, ignore any whitespace matching “ \textbackslash{}t\textbackslash{}r\textbackslash{}n\textbackslash{}f”
\begin{quote}\begin{description}
\sphinxlineitem{Parameters}
\sphinxAtStartPar
\sphinxstyleliteralstrong{\sphinxupquote{text}} \textendash{} string that should be searched 

\sphinxlineitem{Returns}
\sphinxAtStartPar
number of words found 

\end{description}\end{quote}

\end{fulllineitems}

\index{utils::count\_words (C++ function)@\spxentry{utils::count\_words}\spxextra{C++ function}}

\begin{fulllineitems}
\phantomsection\label{\detokenize{Developer_utils:_CPPv4N9LAMMPS_NS5utils11count_wordsERKNSt6stringERKNSt6stringE}}\phantomsection\label{\detokenize{Developer_utils:_CPPv3N9LAMMPS_NS5utils11count_wordsERKNSt6stringERKNSt6stringE}}\phantomsection\label{\detokenize{Developer_utils:_CPPv2N9LAMMPS_NS5utils11count_wordsERKNSt6stringERKNSt6stringE}}\phantomsection\label{\detokenize{Developer_utils:LAMMPS_NS::utils::count_words__ssCR.ssCR}}
\pysigstartsignatures
\pysigstartmultiline
\pysiglinewithargsret{\phantomsection\label{\detokenize{Developer_utils:utils_8h_1a31175af2f2305fea6e26a2f780e0a491}}\DUrole{n}{size\_t}\DUrole{w}{ }\sphinxcode{\sphinxupquote{\DUrole{n}{LAMMPS\_NS}\DUrole{p}{::}\DUrole{n}{utils}\DUrole{p}{::}}}\sphinxbfcode{\sphinxupquote{\DUrole{n}{count\_words}}}}{\DUrole{k}{const}\DUrole{w}{ }\DUrole{n}{std}\DUrole{p}{::}\DUrole{n}{string}\DUrole{w}{ }\DUrole{p}{\&}\DUrole{n,sig-param}{text}\sphinxparamcomma \DUrole{k}{const}\DUrole{w}{ }\DUrole{n}{std}\DUrole{p}{::}\DUrole{n}{string}\DUrole{w}{ }\DUrole{p}{\&}\DUrole{n,sig-param}{separators}}{}
\pysigstopmultiline
\pysigstopsignatures
\sphinxAtStartPar
Count words in string with custom choice of separating characters
\begin{quote}\begin{description}
\sphinxlineitem{Parameters}\begin{itemize}
\item {} 
\sphinxAtStartPar
\sphinxstyleliteralstrong{\sphinxupquote{text}} \textendash{} string that should be searched 

\item {} 
\sphinxAtStartPar
\sphinxstyleliteralstrong{\sphinxupquote{separators}} \textendash{} string containing characters that will be treated as whitespace 

\end{itemize}

\sphinxlineitem{Returns}
\sphinxAtStartPar
number of words found 

\end{description}\end{quote}

\end{fulllineitems}

\index{utils::trim\_and\_count\_words (C++ function)@\spxentry{utils::trim\_and\_count\_words}\spxextra{C++ function}}

\begin{fulllineitems}
\phantomsection\label{\detokenize{Developer_utils:_CPPv4N9LAMMPS_NS5utils20trim_and_count_wordsERKNSt6stringERKNSt6stringE}}\phantomsection\label{\detokenize{Developer_utils:_CPPv3N9LAMMPS_NS5utils20trim_and_count_wordsERKNSt6stringERKNSt6stringE}}\phantomsection\label{\detokenize{Developer_utils:_CPPv2N9LAMMPS_NS5utils20trim_and_count_wordsERKNSt6stringERKNSt6stringE}}\phantomsection\label{\detokenize{Developer_utils:LAMMPS_NS::utils::trim_and_count_words__ssCR.ssCR}}
\pysigstartsignatures
\pysigstartmultiline
\pysiglinewithargsret{\phantomsection\label{\detokenize{Developer_utils:utils_8h_1af92e41728a60e9d949fa37e68f132178}}\DUrole{n}{size\_t}\DUrole{w}{ }\sphinxcode{\sphinxupquote{\DUrole{n}{LAMMPS\_NS}\DUrole{p}{::}\DUrole{n}{utils}\DUrole{p}{::}}}\sphinxbfcode{\sphinxupquote{\DUrole{n}{trim\_and\_count\_words}}}}{\DUrole{k}{const}\DUrole{w}{ }\DUrole{n}{std}\DUrole{p}{::}\DUrole{n}{string}\DUrole{w}{ }\DUrole{p}{\&}\DUrole{n,sig-param}{text}\sphinxparamcomma \DUrole{k}{const}\DUrole{w}{ }\DUrole{n}{std}\DUrole{p}{::}\DUrole{n}{string}\DUrole{w}{ }\DUrole{p}{\&}\DUrole{n,sig-param}{separators}\DUrole{w}{ }\DUrole{p}{=}\DUrole{w}{ }\DUrole{s}{" \textbackslash{}t\textbackslash{}r\textbackslash{}n\textbackslash{}f"}}{}
\pysigstopmultiline
\pysigstopsignatures
\sphinxAtStartPar
Count words in a single line, trim anything from ‘\#’ onward
\begin{quote}\begin{description}
\sphinxlineitem{Parameters}\begin{itemize}
\item {} 
\sphinxAtStartPar
\sphinxstyleliteralstrong{\sphinxupquote{text}} \textendash{} string that should be trimmed and searched 

\item {} 
\sphinxAtStartPar
\sphinxstyleliteralstrong{\sphinxupquote{separators}} \textendash{} string containing characters that will be treated as whitespace 

\end{itemize}

\sphinxlineitem{Returns}
\sphinxAtStartPar
number of words found 

\end{description}\end{quote}

\end{fulllineitems}

\index{utils::join\_words (C++ function)@\spxentry{utils::join\_words}\spxextra{C++ function}}

\begin{fulllineitems}
\phantomsection\label{\detokenize{Developer_utils:_CPPv4N9LAMMPS_NS5utils10join_wordsERKNSt6vectorINSt6stringEEERKNSt6stringE}}\phantomsection\label{\detokenize{Developer_utils:_CPPv3N9LAMMPS_NS5utils10join_wordsERKNSt6vectorINSt6stringEEERKNSt6stringE}}\phantomsection\label{\detokenize{Developer_utils:_CPPv2N9LAMMPS_NS5utils10join_wordsERKNSt6vectorINSt6stringEEERKNSt6stringE}}\phantomsection\label{\detokenize{Developer_utils:LAMMPS_NS::utils::join_words__std::vector:ss:CR.ssCR}}
\pysigstartsignatures
\pysigstartmultiline
\pysiglinewithargsret{\phantomsection\label{\detokenize{Developer_utils:utils_8h_1a52b303c232166990f27299fb49b19ccb}}\DUrole{n}{std}\DUrole{p}{::}\DUrole{n}{string}\DUrole{w}{ }\sphinxcode{\sphinxupquote{\DUrole{n}{LAMMPS\_NS}\DUrole{p}{::}\DUrole{n}{utils}\DUrole{p}{::}}}\sphinxbfcode{\sphinxupquote{\DUrole{n}{join\_words}}}}{\DUrole{k}{const}\DUrole{w}{ }\DUrole{n}{std}\DUrole{p}{::}\DUrole{n}{vector}\DUrole{p}{\textless{}}\DUrole{n}{std}\DUrole{p}{::}\DUrole{n}{string}\DUrole{p}{\textgreater{}}\DUrole{w}{ }\DUrole{p}{\&}\DUrole{n,sig-param}{words}\sphinxparamcomma \DUrole{k}{const}\DUrole{w}{ }\DUrole{n}{std}\DUrole{p}{::}\DUrole{n}{string}\DUrole{w}{ }\DUrole{p}{\&}\DUrole{n,sig-param}{sep}}{}
\pysigstopmultiline
\pysigstopsignatures
\sphinxAtStartPar
Take list of words and join them with a given separator text.

\sphinxAtStartPar
This is the inverse operation of what the split\_words() function Tokenizer classes do.
\begin{quote}\begin{description}
\sphinxlineitem{Parameters}\begin{itemize}
\item {} 
\sphinxAtStartPar
\sphinxstyleliteralstrong{\sphinxupquote{words}} \textendash{} STL vector with strings 

\item {} 
\sphinxAtStartPar
\sphinxstyleliteralstrong{\sphinxupquote{sep}} \textendash{} separator string (may be empty) 

\end{itemize}

\sphinxlineitem{Returns}
\sphinxAtStartPar
string with the concatenated words and separators 

\end{description}\end{quote}

\end{fulllineitems}

\index{utils::split\_words (C++ function)@\spxentry{utils::split\_words}\spxextra{C++ function}}

\begin{fulllineitems}
\phantomsection\label{\detokenize{Developer_utils:_CPPv4N9LAMMPS_NS5utils11split_wordsERKNSt6stringE}}\phantomsection\label{\detokenize{Developer_utils:_CPPv3N9LAMMPS_NS5utils11split_wordsERKNSt6stringE}}\phantomsection\label{\detokenize{Developer_utils:_CPPv2N9LAMMPS_NS5utils11split_wordsERKNSt6stringE}}\phantomsection\label{\detokenize{Developer_utils:LAMMPS_NS::utils::split_words__ssCR}}
\pysigstartsignatures
\pysigstartmultiline
\pysiglinewithargsret{\phantomsection\label{\detokenize{Developer_utils:utils_8h_1ac6b258d9949ba630b8383fa857f9bbb9}}\DUrole{n}{std}\DUrole{p}{::}\DUrole{n}{vector}\DUrole{p}{\textless{}}\DUrole{n}{std}\DUrole{p}{::}\DUrole{n}{string}\DUrole{p}{\textgreater{}}\DUrole{w}{ }\sphinxcode{\sphinxupquote{\DUrole{n}{LAMMPS\_NS}\DUrole{p}{::}\DUrole{n}{utils}\DUrole{p}{::}}}\sphinxbfcode{\sphinxupquote{\DUrole{n}{split\_words}}}}{\DUrole{k}{const}\DUrole{w}{ }\DUrole{n}{std}\DUrole{p}{::}\DUrole{n}{string}\DUrole{w}{ }\DUrole{p}{\&}\DUrole{n,sig-param}{text}}{}
\pysigstopmultiline
\pysigstopsignatures
\sphinxAtStartPar
Take text and split into non\sphinxhyphen{}whitespace words.

\sphinxAtStartPar
This can handle strings with single and double quotes, escaped quotes, and escaped codes within quotes, but due to using an STL container and STL strings is rather slow because of making copies. Designed for parsing command lines and similar text and not for time critical processing. Use a tokenizer class if performance matters.

\sphinxAtStartPar

\sphinxAtStartPar
\begin{description}
\sphinxlineitem{\sphinxstyleemphasis{See also}}
\sphinxAtStartPar
{\hyperref[\detokenize{Developer_utils:_CPPv4N9LAMMPS_NS9TokenizerE}]{\sphinxcrossref{\sphinxcode{\sphinxupquote{Tokenizer}}}}}, {\hyperref[\detokenize{Developer_utils:_CPPv4N9LAMMPS_NS14ValueTokenizerE}]{\sphinxcrossref{\sphinxcode{\sphinxupquote{ValueTokenizer}}}}}

\end{description}

 
\begin{quote}\begin{description}
\sphinxlineitem{Parameters}
\sphinxAtStartPar
\sphinxstyleliteralstrong{\sphinxupquote{text}} \textendash{} string that should be split 

\sphinxlineitem{Returns}
\sphinxAtStartPar
STL vector with the words 

\end{description}\end{quote}

\end{fulllineitems}

\index{utils::split\_lines (C++ function)@\spxentry{utils::split\_lines}\spxextra{C++ function}}

\begin{fulllineitems}
\phantomsection\label{\detokenize{Developer_utils:_CPPv4N9LAMMPS_NS5utils11split_linesERKNSt6stringE}}\phantomsection\label{\detokenize{Developer_utils:_CPPv3N9LAMMPS_NS5utils11split_linesERKNSt6stringE}}\phantomsection\label{\detokenize{Developer_utils:_CPPv2N9LAMMPS_NS5utils11split_linesERKNSt6stringE}}\phantomsection\label{\detokenize{Developer_utils:LAMMPS_NS::utils::split_lines__ssCR}}
\pysigstartsignatures
\pysigstartmultiline
\pysiglinewithargsret{\phantomsection\label{\detokenize{Developer_utils:utils_8h_1aa1a7c146c8809600ba053e45c914e84e}}\DUrole{n}{std}\DUrole{p}{::}\DUrole{n}{vector}\DUrole{p}{\textless{}}\DUrole{n}{std}\DUrole{p}{::}\DUrole{n}{string}\DUrole{p}{\textgreater{}}\DUrole{w}{ }\sphinxcode{\sphinxupquote{\DUrole{n}{LAMMPS\_NS}\DUrole{p}{::}\DUrole{n}{utils}\DUrole{p}{::}}}\sphinxbfcode{\sphinxupquote{\DUrole{n}{split\_lines}}}}{\DUrole{k}{const}\DUrole{w}{ }\DUrole{n}{std}\DUrole{p}{::}\DUrole{n}{string}\DUrole{w}{ }\DUrole{p}{\&}\DUrole{n,sig-param}{text}}{}
\pysigstopmultiline
\pysigstopsignatures
\sphinxAtStartPar
Take multi\sphinxhyphen{}line text and split into lines
\begin{quote}\begin{description}
\sphinxlineitem{Parameters}
\sphinxAtStartPar
\sphinxstyleliteralstrong{\sphinxupquote{text}} \textendash{} string that should be split 

\sphinxlineitem{Returns}
\sphinxAtStartPar
STL vector with the lines 

\end{description}\end{quote}

\end{fulllineitems}

\index{utils::strmatch (C++ function)@\spxentry{utils::strmatch}\spxextra{C++ function}}

\begin{fulllineitems}
\phantomsection\label{\detokenize{Developer_utils:_CPPv4N9LAMMPS_NS5utils8strmatchERKNSt6stringERKNSt6stringE}}\phantomsection\label{\detokenize{Developer_utils:_CPPv3N9LAMMPS_NS5utils8strmatchERKNSt6stringERKNSt6stringE}}\phantomsection\label{\detokenize{Developer_utils:_CPPv2N9LAMMPS_NS5utils8strmatchERKNSt6stringERKNSt6stringE}}\phantomsection\label{\detokenize{Developer_utils:LAMMPS_NS::utils::strmatch__ssCR.ssCR}}
\pysigstartsignatures
\pysigstartmultiline
\pysiglinewithargsret{\phantomsection\label{\detokenize{Developer_utils:utils_8h_1a3a2003eb8f17075d799ecf645e4f2c30}}\DUrole{kt}{bool}\DUrole{w}{ }\sphinxcode{\sphinxupquote{\DUrole{n}{LAMMPS\_NS}\DUrole{p}{::}\DUrole{n}{utils}\DUrole{p}{::}}}\sphinxbfcode{\sphinxupquote{\DUrole{n}{strmatch}}}}{\DUrole{k}{const}\DUrole{w}{ }\DUrole{n}{std}\DUrole{p}{::}\DUrole{n}{string}\DUrole{w}{ }\DUrole{p}{\&}\DUrole{n,sig-param}{text}\sphinxparamcomma \DUrole{k}{const}\DUrole{w}{ }\DUrole{n}{std}\DUrole{p}{::}\DUrole{n}{string}\DUrole{w}{ }\DUrole{p}{\&}\DUrole{n,sig-param}{pattern}}{}
\pysigstopmultiline
\pysigstopsignatures
\sphinxAtStartPar
Match text against a simplified regex pattern

\sphinxAtStartPar

More flexible and specific matching of a string against a pattern. This function is supposed to be a more safe, more specific and simple to use API to find pattern matches. The purpose is to replace uses of either strncmp() or strstr() in the code base to find sub\sphinxhyphen{}strings safely. With strncmp() finding prefixes, the number of characters to match must be counted, which can lead to errors, while using “\textasciicircum{}pattern” will do the same with less problems. Matching for suffixes using strstr() is not as specific as ‘pattern\$’, and complex matches, e.g. “\textasciicircum{}rigid.*\textbackslash{}/small.*”, to match all small body optimized rigid fixes require only one test.

\sphinxAtStartPar
The use of std::string arguments allows for simple concatenation even with char * type variables. Example: utils::strmatch(text, std::string(“\textasciicircum{}”) + charptr) 
\begin{quote}\begin{description}
\sphinxlineitem{Parameters}\begin{itemize}
\item {} 
\sphinxAtStartPar
\sphinxstyleliteralstrong{\sphinxupquote{text}} \textendash{} the text to be matched against the pattern 

\item {} 
\sphinxAtStartPar
\sphinxstyleliteralstrong{\sphinxupquote{pattern}} \textendash{} the search pattern, which may contain regexp markers 

\end{itemize}

\sphinxlineitem{Returns}
\sphinxAtStartPar
true if the pattern matches, false if not

\end{description}\end{quote}

\end{fulllineitems}

\index{utils::strfind (C++ function)@\spxentry{utils::strfind}\spxextra{C++ function}}

\begin{fulllineitems}
\phantomsection\label{\detokenize{Developer_utils:_CPPv4N9LAMMPS_NS5utils7strfindERKNSt6stringERKNSt6stringE}}\phantomsection\label{\detokenize{Developer_utils:_CPPv3N9LAMMPS_NS5utils7strfindERKNSt6stringERKNSt6stringE}}\phantomsection\label{\detokenize{Developer_utils:_CPPv2N9LAMMPS_NS5utils7strfindERKNSt6stringERKNSt6stringE}}\phantomsection\label{\detokenize{Developer_utils:LAMMPS_NS::utils::strfind__ssCR.ssCR}}
\pysigstartsignatures
\pysigstartmultiline
\pysiglinewithargsret{\phantomsection\label{\detokenize{Developer_utils:utils_8h_1a0fc913e925ac5ad63addfae419332e31}}\DUrole{n}{std}\DUrole{p}{::}\DUrole{n}{string}\DUrole{w}{ }\sphinxcode{\sphinxupquote{\DUrole{n}{LAMMPS\_NS}\DUrole{p}{::}\DUrole{n}{utils}\DUrole{p}{::}}}\sphinxbfcode{\sphinxupquote{\DUrole{n}{strfind}}}}{\DUrole{k}{const}\DUrole{w}{ }\DUrole{n}{std}\DUrole{p}{::}\DUrole{n}{string}\DUrole{w}{ }\DUrole{p}{\&}\DUrole{n,sig-param}{text}\sphinxparamcomma \DUrole{k}{const}\DUrole{w}{ }\DUrole{n}{std}\DUrole{p}{::}\DUrole{n}{string}\DUrole{w}{ }\DUrole{p}{\&}\DUrole{n,sig-param}{pattern}}{}
\pysigstopmultiline
\pysigstopsignatures
\sphinxAtStartPar
Find sub\sphinxhyphen{}string that matches a simplified regex pattern

\sphinxAtStartPar

This function is a companion function to utils::strmatch(). Arguments and logic is the same, but instead of a boolean, it returns the sub\sphinxhyphen{}string that matches the regex pattern. There can be only one match. This can be used as a more flexible alternative to strstr(). 
\begin{quote}\begin{description}
\sphinxlineitem{Parameters}\begin{itemize}
\item {} 
\sphinxAtStartPar
\sphinxstyleliteralstrong{\sphinxupquote{text}} \textendash{} the text to be matched against the pattern 

\item {} 
\sphinxAtStartPar
\sphinxstyleliteralstrong{\sphinxupquote{pattern}} \textendash{} the search pattern, which may contain regexp markers 

\end{itemize}

\sphinxlineitem{Returns}
\sphinxAtStartPar
the string that matches the pattern or an empty one

\end{description}\end{quote}

\end{fulllineitems}

\index{utils::is\_integer (C++ function)@\spxentry{utils::is\_integer}\spxextra{C++ function}}

\begin{fulllineitems}
\phantomsection\label{\detokenize{Developer_utils:_CPPv4N9LAMMPS_NS5utils10is_integerERKNSt6stringE}}\phantomsection\label{\detokenize{Developer_utils:_CPPv3N9LAMMPS_NS5utils10is_integerERKNSt6stringE}}\phantomsection\label{\detokenize{Developer_utils:_CPPv2N9LAMMPS_NS5utils10is_integerERKNSt6stringE}}\phantomsection\label{\detokenize{Developer_utils:LAMMPS_NS::utils::is_integer__ssCR}}
\pysigstartsignatures
\pysigstartmultiline
\pysiglinewithargsret{\phantomsection\label{\detokenize{Developer_utils:utils_8h_1a66b3c703e53c70f60ba745a557cace67}}\DUrole{kt}{bool}\DUrole{w}{ }\sphinxcode{\sphinxupquote{\DUrole{n}{LAMMPS\_NS}\DUrole{p}{::}\DUrole{n}{utils}\DUrole{p}{::}}}\sphinxbfcode{\sphinxupquote{\DUrole{n}{is\_integer}}}}{\DUrole{k}{const}\DUrole{w}{ }\DUrole{n}{std}\DUrole{p}{::}\DUrole{n}{string}\DUrole{w}{ }\DUrole{p}{\&}\DUrole{n,sig-param}{str}}{}
\pysigstopmultiline
\pysigstopsignatures
\sphinxAtStartPar
Check if string can be converted to valid integer
\begin{quote}\begin{description}
\sphinxlineitem{Parameters}
\sphinxAtStartPar
\sphinxstyleliteralstrong{\sphinxupquote{str}} \textendash{} string that should be checked 

\sphinxlineitem{Returns}
\sphinxAtStartPar
true, if string contains valid a integer, false otherwise 

\end{description}\end{quote}

\end{fulllineitems}

\index{utils::is\_double (C++ function)@\spxentry{utils::is\_double}\spxextra{C++ function}}

\begin{fulllineitems}
\phantomsection\label{\detokenize{Developer_utils:_CPPv4N9LAMMPS_NS5utils9is_doubleERKNSt6stringE}}\phantomsection\label{\detokenize{Developer_utils:_CPPv3N9LAMMPS_NS5utils9is_doubleERKNSt6stringE}}\phantomsection\label{\detokenize{Developer_utils:_CPPv2N9LAMMPS_NS5utils9is_doubleERKNSt6stringE}}\phantomsection\label{\detokenize{Developer_utils:LAMMPS_NS::utils::is_double__ssCR}}
\pysigstartsignatures
\pysigstartmultiline
\pysiglinewithargsret{\phantomsection\label{\detokenize{Developer_utils:utils_8h_1a3f130082bbb981e0253f32bc46c1b3b3}}\DUrole{kt}{bool}\DUrole{w}{ }\sphinxcode{\sphinxupquote{\DUrole{n}{LAMMPS\_NS}\DUrole{p}{::}\DUrole{n}{utils}\DUrole{p}{::}}}\sphinxbfcode{\sphinxupquote{\DUrole{n}{is\_double}}}}{\DUrole{k}{const}\DUrole{w}{ }\DUrole{n}{std}\DUrole{p}{::}\DUrole{n}{string}\DUrole{w}{ }\DUrole{p}{\&}\DUrole{n,sig-param}{str}}{}
\pysigstopmultiline
\pysigstopsignatures
\sphinxAtStartPar
Check if string can be converted to valid floating\sphinxhyphen{}point number
\begin{quote}\begin{description}
\sphinxlineitem{Parameters}
\sphinxAtStartPar
\sphinxstyleliteralstrong{\sphinxupquote{str}} \textendash{} string that should be checked 

\sphinxlineitem{Returns}
\sphinxAtStartPar
true, if string contains valid number, false otherwise 

\end{description}\end{quote}

\end{fulllineitems}

\index{utils::is\_id (C++ function)@\spxentry{utils::is\_id}\spxextra{C++ function}}

\begin{fulllineitems}
\phantomsection\label{\detokenize{Developer_utils:_CPPv4N9LAMMPS_NS5utils5is_idERKNSt6stringE}}\phantomsection\label{\detokenize{Developer_utils:_CPPv3N9LAMMPS_NS5utils5is_idERKNSt6stringE}}\phantomsection\label{\detokenize{Developer_utils:_CPPv2N9LAMMPS_NS5utils5is_idERKNSt6stringE}}\phantomsection\label{\detokenize{Developer_utils:LAMMPS_NS::utils::is_id__ssCR}}
\pysigstartsignatures
\pysigstartmultiline
\pysiglinewithargsret{\phantomsection\label{\detokenize{Developer_utils:utils_8h_1a8b19cdd3cde5fe93a7b5f87c3271e117}}\DUrole{kt}{bool}\DUrole{w}{ }\sphinxcode{\sphinxupquote{\DUrole{n}{LAMMPS\_NS}\DUrole{p}{::}\DUrole{n}{utils}\DUrole{p}{::}}}\sphinxbfcode{\sphinxupquote{\DUrole{n}{is\_id}}}}{\DUrole{k}{const}\DUrole{w}{ }\DUrole{n}{std}\DUrole{p}{::}\DUrole{n}{string}\DUrole{w}{ }\DUrole{p}{\&}\DUrole{n,sig-param}{str}}{}
\pysigstopmultiline
\pysigstopsignatures
\sphinxAtStartPar
Check if string is a valid ID ID strings may contain only letters, numbers, and underscores.
\begin{quote}\begin{description}
\sphinxlineitem{Parameters}
\sphinxAtStartPar
\sphinxstyleliteralstrong{\sphinxupquote{str}} \textendash{} string that should be checked 

\sphinxlineitem{Returns}
\sphinxAtStartPar
true, if string contains valid id, false otherwise 

\end{description}\end{quote}

\end{fulllineitems}

\index{utils::is\_type (C++ function)@\spxentry{utils::is\_type}\spxextra{C++ function}}

\begin{fulllineitems}
\phantomsection\label{\detokenize{Developer_utils:_CPPv4N9LAMMPS_NS5utils7is_typeERKNSt6stringE}}\phantomsection\label{\detokenize{Developer_utils:_CPPv3N9LAMMPS_NS5utils7is_typeERKNSt6stringE}}\phantomsection\label{\detokenize{Developer_utils:_CPPv2N9LAMMPS_NS5utils7is_typeERKNSt6stringE}}\phantomsection\label{\detokenize{Developer_utils:LAMMPS_NS::utils::is_type__ssCR}}
\pysigstartsignatures
\pysigstartmultiline
\pysiglinewithargsret{\phantomsection\label{\detokenize{Developer_utils:utils_8h_1adf45c3878dd1d69dc6a14887244b5f36}}\DUrole{kt}{int}\DUrole{w}{ }\sphinxcode{\sphinxupquote{\DUrole{n}{LAMMPS\_NS}\DUrole{p}{::}\DUrole{n}{utils}\DUrole{p}{::}}}\sphinxbfcode{\sphinxupquote{\DUrole{n}{is\_type}}}}{\DUrole{k}{const}\DUrole{w}{ }\DUrole{n}{std}\DUrole{p}{::}\DUrole{n}{string}\DUrole{w}{ }\DUrole{p}{\&}\DUrole{n,sig-param}{str}}{}
\pysigstopmultiline
\pysigstopsignatures
\sphinxAtStartPar
Check if string is a valid type label, or numeric type, or numeric type range. Numeric type or type range may only contain digits or the ‘*’ character. Type label strings may not contain a digit, or a ‘*’, or a ‘\#’ character as the first character to distinguish them from comments and numeric types or type ranges. They also may not contain any whitespace. If the string is a valid numeric type or type range the function returns 0, if it is a valid type label the function returns 1, otherwise it returns \sphinxhyphen{}1.
\begin{quote}\begin{description}
\sphinxlineitem{Parameters}
\sphinxAtStartPar
\sphinxstyleliteralstrong{\sphinxupquote{str}} \textendash{} string that should be checked 

\sphinxlineitem{Returns}
\sphinxAtStartPar
0, 1, or \sphinxhyphen{}1, depending on whether the string is valid numeric type, valid type label or neither, respectively 

\end{description}\end{quote}

\end{fulllineitems}



\subsection{Potential file functions}
\label{\detokenize{Developer_utils:potential-file-functions}}\index{utils::get\_potential\_file\_path (C++ function)@\spxentry{utils::get\_potential\_file\_path}\spxextra{C++ function}}

\begin{fulllineitems}
\phantomsection\label{\detokenize{Developer_utils:_CPPv4N9LAMMPS_NS5utils23get_potential_file_pathERKNSt6stringE}}\phantomsection\label{\detokenize{Developer_utils:_CPPv3N9LAMMPS_NS5utils23get_potential_file_pathERKNSt6stringE}}\phantomsection\label{\detokenize{Developer_utils:_CPPv2N9LAMMPS_NS5utils23get_potential_file_pathERKNSt6stringE}}\phantomsection\label{\detokenize{Developer_utils:LAMMPS_NS::utils::get_potential_file_path__ssCR}}
\pysigstartsignatures
\pysigstartmultiline
\pysiglinewithargsret{\phantomsection\label{\detokenize{Developer_utils:utils_8h_1a54601c14224556c45ed863cb95e8d49c}}\DUrole{n}{std}\DUrole{p}{::}\DUrole{n}{string}\DUrole{w}{ }\sphinxcode{\sphinxupquote{\DUrole{n}{LAMMPS\_NS}\DUrole{p}{::}\DUrole{n}{utils}\DUrole{p}{::}}}\sphinxbfcode{\sphinxupquote{\DUrole{n}{get\_potential\_file\_path}}}}{\DUrole{k}{const}\DUrole{w}{ }\DUrole{n}{std}\DUrole{p}{::}\DUrole{n}{string}\DUrole{w}{ }\DUrole{p}{\&}\DUrole{n,sig-param}{path}}{}
\pysigstopmultiline
\pysigstopsignatures
\sphinxAtStartPar
Determine full path of potential file. If file is not found in current directory, search directories listed in LAMMPS\_POTENTIALS environment variable
\begin{quote}\begin{description}
\sphinxlineitem{Parameters}
\sphinxAtStartPar
\sphinxstyleliteralstrong{\sphinxupquote{path}} \textendash{} file path 

\sphinxlineitem{Returns}
\sphinxAtStartPar
full path to potential file 

\end{description}\end{quote}

\end{fulllineitems}

\index{utils::get\_potential\_date (C++ function)@\spxentry{utils::get\_potential\_date}\spxextra{C++ function}}

\begin{fulllineitems}
\phantomsection\label{\detokenize{Developer_utils:_CPPv4N9LAMMPS_NS5utils18get_potential_dateERKNSt6stringERKNSt6stringE}}\phantomsection\label{\detokenize{Developer_utils:_CPPv3N9LAMMPS_NS5utils18get_potential_dateERKNSt6stringERKNSt6stringE}}\phantomsection\label{\detokenize{Developer_utils:_CPPv2N9LAMMPS_NS5utils18get_potential_dateERKNSt6stringERKNSt6stringE}}\phantomsection\label{\detokenize{Developer_utils:LAMMPS_NS::utils::get_potential_date__ssCR.ssCR}}
\pysigstartsignatures
\pysigstartmultiline
\pysiglinewithargsret{\phantomsection\label{\detokenize{Developer_utils:utils_8h_1a86fd64c9c2c2183daada07c001e91392}}\DUrole{n}{std}\DUrole{p}{::}\DUrole{n}{string}\DUrole{w}{ }\sphinxcode{\sphinxupquote{\DUrole{n}{LAMMPS\_NS}\DUrole{p}{::}\DUrole{n}{utils}\DUrole{p}{::}}}\sphinxbfcode{\sphinxupquote{\DUrole{n}{get\_potential\_date}}}}{\DUrole{k}{const}\DUrole{w}{ }\DUrole{n}{std}\DUrole{p}{::}\DUrole{n}{string}\DUrole{w}{ }\DUrole{p}{\&}\DUrole{n,sig-param}{path}\sphinxparamcomma \DUrole{k}{const}\DUrole{w}{ }\DUrole{n}{std}\DUrole{p}{::}\DUrole{n}{string}\DUrole{w}{ }\DUrole{p}{\&}\DUrole{n,sig-param}{potential\_name}}{}
\pysigstopmultiline
\pysigstopsignatures
\sphinxAtStartPar
Read potential file and return DATE field if it is present
\begin{quote}\begin{description}
\sphinxlineitem{Parameters}\begin{itemize}
\item {} 
\sphinxAtStartPar
\sphinxstyleliteralstrong{\sphinxupquote{path}} \textendash{} file path 

\item {} 
\sphinxAtStartPar
\sphinxstyleliteralstrong{\sphinxupquote{potential\_name}} \textendash{} name of potential that is being read 

\end{itemize}

\sphinxlineitem{Returns}
\sphinxAtStartPar
DATE field if present 

\end{description}\end{quote}

\end{fulllineitems}

\index{utils::get\_potential\_units (C++ function)@\spxentry{utils::get\_potential\_units}\spxextra{C++ function}}

\begin{fulllineitems}
\phantomsection\label{\detokenize{Developer_utils:_CPPv4N9LAMMPS_NS5utils19get_potential_unitsERKNSt6stringERKNSt6stringE}}\phantomsection\label{\detokenize{Developer_utils:_CPPv3N9LAMMPS_NS5utils19get_potential_unitsERKNSt6stringERKNSt6stringE}}\phantomsection\label{\detokenize{Developer_utils:_CPPv2N9LAMMPS_NS5utils19get_potential_unitsERKNSt6stringERKNSt6stringE}}\phantomsection\label{\detokenize{Developer_utils:LAMMPS_NS::utils::get_potential_units__ssCR.ssCR}}
\pysigstartsignatures
\pysigstartmultiline
\pysiglinewithargsret{\phantomsection\label{\detokenize{Developer_utils:utils_8h_1ab8b096aefd75e4cef0162f113b800c2e}}\DUrole{n}{std}\DUrole{p}{::}\DUrole{n}{string}\DUrole{w}{ }\sphinxcode{\sphinxupquote{\DUrole{n}{LAMMPS\_NS}\DUrole{p}{::}\DUrole{n}{utils}\DUrole{p}{::}}}\sphinxbfcode{\sphinxupquote{\DUrole{n}{get\_potential\_units}}}}{\DUrole{k}{const}\DUrole{w}{ }\DUrole{n}{std}\DUrole{p}{::}\DUrole{n}{string}\DUrole{w}{ }\DUrole{p}{\&}\DUrole{n,sig-param}{path}\sphinxparamcomma \DUrole{k}{const}\DUrole{w}{ }\DUrole{n}{std}\DUrole{p}{::}\DUrole{n}{string}\DUrole{w}{ }\DUrole{p}{\&}\DUrole{n,sig-param}{potential\_name}}{}
\pysigstopmultiline
\pysigstopsignatures
\sphinxAtStartPar
Read potential file and return UNITS field if it is present
\begin{quote}\begin{description}
\sphinxlineitem{Parameters}\begin{itemize}
\item {} 
\sphinxAtStartPar
\sphinxstyleliteralstrong{\sphinxupquote{path}} \textendash{} file path 

\item {} 
\sphinxAtStartPar
\sphinxstyleliteralstrong{\sphinxupquote{potential\_name}} \textendash{} name of potential that is being read 

\end{itemize}

\sphinxlineitem{Returns}
\sphinxAtStartPar
UNITS field if present 

\end{description}\end{quote}

\end{fulllineitems}

\index{utils::get\_supported\_conversions (C++ function)@\spxentry{utils::get\_supported\_conversions}\spxextra{C++ function}}

\begin{fulllineitems}
\phantomsection\label{\detokenize{Developer_utils:_CPPv4N9LAMMPS_NS5utils25get_supported_conversionsEKi}}\phantomsection\label{\detokenize{Developer_utils:_CPPv3N9LAMMPS_NS5utils25get_supported_conversionsEKi}}\phantomsection\label{\detokenize{Developer_utils:_CPPv2N9LAMMPS_NS5utils25get_supported_conversionsEKi}}\phantomsection\label{\detokenize{Developer_utils:LAMMPS_NS::utils::get_supported_conversions__iC}}
\pysigstartsignatures
\pysigstartmultiline
\pysiglinewithargsret{\phantomsection\label{\detokenize{Developer_utils:utils_8h_1a4d70c1d4e44dad6d82bb7fb643adb095}}\DUrole{kt}{int}\DUrole{w}{ }\sphinxcode{\sphinxupquote{\DUrole{n}{LAMMPS\_NS}\DUrole{p}{::}\DUrole{n}{utils}\DUrole{p}{::}}}\sphinxbfcode{\sphinxupquote{\DUrole{n}{get\_supported\_conversions}}}}{\DUrole{k}{const}\DUrole{w}{ }\DUrole{kt}{int}\DUrole{w}{ }\DUrole{n,sig-param}{property}}{}
\pysigstopmultiline
\pysigstopsignatures
\sphinxAtStartPar
Return bitmask of available conversion factors for a given property
\begin{quote}\begin{description}
\sphinxlineitem{Parameters}
\sphinxAtStartPar
\sphinxstyleliteralstrong{\sphinxupquote{property}} \textendash{} property to be converted 

\sphinxlineitem{Returns}
\sphinxAtStartPar
bitmask indicating available conversions 

\end{description}\end{quote}

\end{fulllineitems}

\index{utils::get\_conversion\_factor (C++ function)@\spxentry{utils::get\_conversion\_factor}\spxextra{C++ function}}

\begin{fulllineitems}
\phantomsection\label{\detokenize{Developer_utils:_CPPv4N9LAMMPS_NS5utils21get_conversion_factorEKiKi}}\phantomsection\label{\detokenize{Developer_utils:_CPPv3N9LAMMPS_NS5utils21get_conversion_factorEKiKi}}\phantomsection\label{\detokenize{Developer_utils:_CPPv2N9LAMMPS_NS5utils21get_conversion_factorEKiKi}}\phantomsection\label{\detokenize{Developer_utils:LAMMPS_NS::utils::get_conversion_factor__iC.iC}}
\pysigstartsignatures
\pysigstartmultiline
\pysiglinewithargsret{\phantomsection\label{\detokenize{Developer_utils:utils_8h_1ae069d34c47dc4eeb168de963e592f5ed}}\DUrole{kt}{double}\DUrole{w}{ }\sphinxcode{\sphinxupquote{\DUrole{n}{LAMMPS\_NS}\DUrole{p}{::}\DUrole{n}{utils}\DUrole{p}{::}}}\sphinxbfcode{\sphinxupquote{\DUrole{n}{get\_conversion\_factor}}}}{\DUrole{k}{const}\DUrole{w}{ }\DUrole{kt}{int}\DUrole{w}{ }\DUrole{n,sig-param}{property}\sphinxparamcomma \DUrole{k}{const}\DUrole{w}{ }\DUrole{kt}{int}\DUrole{w}{ }\DUrole{n,sig-param}{conversion}}{}
\pysigstopmultiline
\pysigstopsignatures
\sphinxAtStartPar
Return unit conversion factor for given property and selected from/to units
\begin{quote}\begin{description}
\sphinxlineitem{Parameters}\begin{itemize}
\item {} 
\sphinxAtStartPar
\sphinxstyleliteralstrong{\sphinxupquote{property}} \textendash{} property to be converted 

\item {} 
\sphinxAtStartPar
\sphinxstyleliteralstrong{\sphinxupquote{conversion}} \textendash{} constant indicating the conversion 

\end{itemize}

\sphinxlineitem{Returns}
\sphinxAtStartPar
conversion factor 

\end{description}\end{quote}

\end{fulllineitems}

\index{utils::open\_potential (C++ function)@\spxentry{utils::open\_potential}\spxextra{C++ function}}

\begin{fulllineitems}
\phantomsection\label{\detokenize{Developer_utils:_CPPv4N9LAMMPS_NS5utils14open_potentialERKNSt6stringEP6LAMMPSPi}}\phantomsection\label{\detokenize{Developer_utils:_CPPv3N9LAMMPS_NS5utils14open_potentialERKNSt6stringEP6LAMMPSPi}}\phantomsection\label{\detokenize{Developer_utils:_CPPv2N9LAMMPS_NS5utils14open_potentialERKNSt6stringEP6LAMMPSPi}}\phantomsection\label{\detokenize{Developer_utils:LAMMPS_NS::utils::open_potential__ssCR.LAMMPSP.iP}}
\pysigstartsignatures
\pysigstartmultiline
\pysiglinewithargsret{\phantomsection\label{\detokenize{Developer_utils:utils_8h_1a1b92528e9bbdff202464fded75539b84}}\DUrole{n}{FILE}\DUrole{w}{ }\DUrole{p}{*}\sphinxcode{\sphinxupquote{\DUrole{n}{LAMMPS\_NS}\DUrole{p}{::}\DUrole{n}{utils}\DUrole{p}{::}}}\sphinxbfcode{\sphinxupquote{\DUrole{n}{open\_potential}}}}{\DUrole{k}{const}\DUrole{w}{ }\DUrole{n}{std}\DUrole{p}{::}\DUrole{n}{string}\DUrole{w}{ }\DUrole{p}{\&}\DUrole{n,sig-param}{name}\sphinxparamcomma {\hyperref[\detokenize{Classes_lammps:_CPPv4N9LAMMPS_NS6LAMMPSE}]{\sphinxcrossref{\DUrole{n}{LAMMPS}}}}\DUrole{w}{ }\DUrole{p}{*}\DUrole{n,sig-param}{lmp}\sphinxparamcomma \DUrole{kt}{int}\DUrole{w}{ }\DUrole{p}{*}\DUrole{n,sig-param}{auto\_convert}}{}
\pysigstopmultiline
\pysigstopsignatures
\sphinxAtStartPar
Open a potential file as specified by \sphinxstyleemphasis{name}

\sphinxAtStartPar
If opening the file directly fails, the function will search for it in the list of folder pointed to by the environment variable \sphinxcode{\sphinxupquote{LAMMPS\_POTENTIALS}} (if it is set).

\sphinxAtStartPar
If the potential file has a \sphinxcode{\sphinxupquote{UNITS}} tag in the first line, the tag’s value is compared to the current unit style setting. The behavior of the function then depends on the value of the \sphinxstyleemphasis{auto\_convert} parameter. If it is a null pointer, then the unit values must match or else the open will fail with an error. Otherwise the bitmask that \sphinxstyleemphasis{auto\_convert} points to is used check for compatibility with possible automatic conversions by the calling function. If compatible, the bitmask is set to the required conversion or \sphinxcode{\sphinxupquote{utils::NOCONVERT}}.
\begin{quote}\begin{description}
\sphinxlineitem{Parameters}\begin{itemize}
\item {} 
\sphinxAtStartPar
\sphinxstyleliteralstrong{\sphinxupquote{name}} \textendash{} file\sphinxhyphen{} or pathname of the potential file 

\item {} 
\sphinxAtStartPar
\sphinxstyleliteralstrong{\sphinxupquote{lmp}} \textendash{} pointer to top\sphinxhyphen{}level {\hyperref[\detokenize{Classes_lammps:classLAMMPS__NS_1_1LAMMPS}]{\sphinxcrossref{\DUrole{std,std-ref}{LAMMPS}}}} class instance 

\item {} 
\sphinxAtStartPar
\sphinxstyleliteralstrong{\sphinxupquote{auto\_convert}} \textendash{} pointer to unit conversion bitmask or \sphinxcode{\sphinxupquote{nullptr}}

\end{itemize}

\sphinxlineitem{Returns}
\sphinxAtStartPar
FILE pointer of the opened potential file or \sphinxcode{\sphinxupquote{nullptr}}

\end{description}\end{quote}

\end{fulllineitems}



\subsection{Argument processing}
\label{\detokenize{Developer_utils:argument-processing}}\index{utils::bounds (C++ function)@\spxentry{utils::bounds}\spxextra{C++ function}}

\begin{fulllineitems}
\phantomsection\label{\detokenize{Developer_utils:_CPPv4I0EN9LAMMPS_NS5utils6boundsEvPKciRKNSt6stringE6bigint6bigintR4TYPER4TYPEP5Error}}\phantomsection\label{\detokenize{Developer_utils:_CPPv3I0EN9LAMMPS_NS5utils6boundsEPKciRKNSt6stringE6bigint6bigintR4TYPER4TYPEP5Error}}\phantomsection\label{\detokenize{Developer_utils:_CPPv2I0EN9LAMMPS_NS5utils6boundsEPKciRKNSt6stringE6bigint6bigintR4TYPER4TYPEP5Error}}
\pysigstartsignatures
\pysigstartmultiline
\pysigline{\DUrole{k}{template}\DUrole{p}{\textless{}}\DUrole{k}{typename}\DUrole{w}{ }\sphinxbfcode{\sphinxupquote{\DUrole{n}{TYPE}}}\DUrole{p}{\textgreater{}}}
\pysiglinewithargsret{\phantomsection\label{\detokenize{Developer_utils:utils_8h_1aec32e4cae8ced7029e6f162a7442bcf9}}\DUrole{kt}{void}\DUrole{w}{ }\sphinxcode{\sphinxupquote{\DUrole{n}{LAMMPS\_NS}\DUrole{p}{::}\DUrole{n}{utils}\DUrole{p}{::}}}\sphinxbfcode{\sphinxupquote{\DUrole{n}{bounds}}}}{\DUrole{k}{const}\DUrole{w}{ }\DUrole{kt}{char}\DUrole{w}{ }\DUrole{p}{*}\DUrole{n,sig-param}{file}\sphinxparamcomma \DUrole{kt}{int}\DUrole{w}{ }\DUrole{n,sig-param}{line}\sphinxparamcomma \DUrole{k}{const}\DUrole{w}{ }\DUrole{n}{std}\DUrole{p}{::}\DUrole{n}{string}\DUrole{w}{ }\DUrole{p}{\&}\DUrole{n,sig-param}{str}\sphinxparamcomma \DUrole{n}{bigint}\DUrole{w}{ }\DUrole{n,sig-param}{nmin}\sphinxparamcomma \DUrole{n}{bigint}\DUrole{w}{ }\DUrole{n,sig-param}{nmax}\sphinxparamcomma {\hyperref[\detokenize{Developer_utils:_CPPv4I0EN9LAMMPS_NS5utils6boundsEvPKciRKNSt6stringE6bigint6bigintR4TYPER4TYPEP5Error}]{\sphinxcrossref{\DUrole{n}{TYPE}}}}\DUrole{w}{ }\DUrole{p}{\&}\DUrole{n,sig-param}{nlo}\sphinxparamcomma {\hyperref[\detokenize{Developer_utils:_CPPv4I0EN9LAMMPS_NS5utils6boundsEvPKciRKNSt6stringE6bigint6bigintR4TYPER4TYPEP5Error}]{\sphinxcrossref{\DUrole{n}{TYPE}}}}\DUrole{w}{ }\DUrole{p}{\&}\DUrole{n,sig-param}{nhi}\sphinxparamcomma \DUrole{n}{Error}\DUrole{w}{ }\DUrole{p}{*}\DUrole{n,sig-param}{error}}{}
\pysigstopmultiline
\pysigstopsignatures
\sphinxAtStartPar
Compute index bounds derived from a string with a possible wildcard

\sphinxAtStartPar
This functions processes the string in \sphinxstyleemphasis{str} and set the values of \sphinxstyleemphasis{nlo} and \sphinxstyleemphasis{nhi} according to the following five cases:

\sphinxAtStartPar
\begin{itemize}
\item {} 
\sphinxAtStartPar
a single number, i: nlo = i; nhi = i;

\item {} 
\sphinxAtStartPar
a single asterisk, *: nlo = nmin; nhi = nmax;

\item {} 
\sphinxAtStartPar
a single number followed by an asterisk, i*: nlo = i; nhi = nmax;

\item {} 
\sphinxAtStartPar
a single asterisk followed by a number, *i: nlo = nmin; nhi = i;

\item {} 
\sphinxAtStartPar
two numbers with an asterisk in between. i*j: nlo = i; nhi = j;

\end{itemize}

\begin{quote}\begin{description}
\sphinxlineitem{Parameters}\begin{itemize}
\item {} 
\sphinxAtStartPar
\sphinxstyleliteralstrong{\sphinxupquote{file}} \textendash{} name of source file for error message 

\item {} 
\sphinxAtStartPar
\sphinxstyleliteralstrong{\sphinxupquote{line}} \textendash{} line number in source file for error message 

\item {} 
\sphinxAtStartPar
\sphinxstyleliteralstrong{\sphinxupquote{str}} \textendash{} string to be processed 

\item {} 
\sphinxAtStartPar
\sphinxstyleliteralstrong{\sphinxupquote{nmin}} \textendash{} smallest possible lower bound 

\item {} 
\sphinxAtStartPar
\sphinxstyleliteralstrong{\sphinxupquote{nmax}} \textendash{} largest allowed upper bound 

\item {} 
\sphinxAtStartPar
\sphinxstyleliteralstrong{\sphinxupquote{nlo}} \textendash{} lower bound 

\item {} 
\sphinxAtStartPar
\sphinxstyleliteralstrong{\sphinxupquote{nhi}} \textendash{} upper bound 

\item {} 
\sphinxAtStartPar
\sphinxstyleliteralstrong{\sphinxupquote{error}} \textendash{} pointer to Error class for out\sphinxhyphen{}of\sphinxhyphen{}bounds messages 

\end{itemize}

\end{description}\end{quote}

\end{fulllineitems}

\index{utils::bounds\_typelabel (C++ function)@\spxentry{utils::bounds\_typelabel}\spxextra{C++ function}}

\begin{fulllineitems}
\phantomsection\label{\detokenize{Developer_utils:_CPPv4I0EN9LAMMPS_NS5utils16bounds_typelabelEvPKciRKNSt6stringE6bigint6bigintR4TYPER4TYPEP6LAMMPSi}}\phantomsection\label{\detokenize{Developer_utils:_CPPv3I0EN9LAMMPS_NS5utils16bounds_typelabelEPKciRKNSt6stringE6bigint6bigintR4TYPER4TYPEP6LAMMPSi}}\phantomsection\label{\detokenize{Developer_utils:_CPPv2I0EN9LAMMPS_NS5utils16bounds_typelabelEPKciRKNSt6stringE6bigint6bigintR4TYPER4TYPEP6LAMMPSi}}
\pysigstartsignatures
\pysigstartmultiline
\pysigline{\DUrole{k}{template}\DUrole{p}{\textless{}}\DUrole{k}{typename}\DUrole{w}{ }\sphinxbfcode{\sphinxupquote{\DUrole{n}{TYPE}}}\DUrole{p}{\textgreater{}}}
\pysiglinewithargsret{\phantomsection\label{\detokenize{Developer_utils:utils_8h_1a0591b9a07cdbb8d72183972157e7ff02}}\DUrole{kt}{void}\DUrole{w}{ }\sphinxcode{\sphinxupquote{\DUrole{n}{LAMMPS\_NS}\DUrole{p}{::}\DUrole{n}{utils}\DUrole{p}{::}}}\sphinxbfcode{\sphinxupquote{\DUrole{n}{bounds\_typelabel}}}}{\DUrole{k}{const}\DUrole{w}{ }\DUrole{kt}{char}\DUrole{w}{ }\DUrole{p}{*}\DUrole{n,sig-param}{file}\sphinxparamcomma \DUrole{kt}{int}\DUrole{w}{ }\DUrole{n,sig-param}{line}\sphinxparamcomma \DUrole{k}{const}\DUrole{w}{ }\DUrole{n}{std}\DUrole{p}{::}\DUrole{n}{string}\DUrole{w}{ }\DUrole{p}{\&}\DUrole{n,sig-param}{str}\sphinxparamcomma \DUrole{n}{bigint}\DUrole{w}{ }\DUrole{n,sig-param}{nmin}\sphinxparamcomma \DUrole{n}{bigint}\DUrole{w}{ }\DUrole{n,sig-param}{nmax}\sphinxparamcomma {\hyperref[\detokenize{Developer_utils:_CPPv4I0EN9LAMMPS_NS5utils16bounds_typelabelEvPKciRKNSt6stringE6bigint6bigintR4TYPER4TYPEP6LAMMPSi}]{\sphinxcrossref{\DUrole{n}{TYPE}}}}\DUrole{w}{ }\DUrole{p}{\&}\DUrole{n,sig-param}{nlo}\sphinxparamcomma {\hyperref[\detokenize{Developer_utils:_CPPv4I0EN9LAMMPS_NS5utils16bounds_typelabelEvPKciRKNSt6stringE6bigint6bigintR4TYPER4TYPEP6LAMMPSi}]{\sphinxcrossref{\DUrole{n}{TYPE}}}}\DUrole{w}{ }\DUrole{p}{\&}\DUrole{n,sig-param}{nhi}\sphinxparamcomma {\hyperref[\detokenize{Classes_lammps:_CPPv4N9LAMMPS_NS6LAMMPSE}]{\sphinxcrossref{\DUrole{n}{LAMMPS}}}}\DUrole{w}{ }\DUrole{p}{*}\DUrole{n,sig-param}{lmp}\sphinxparamcomma \DUrole{kt}{int}\DUrole{w}{ }\DUrole{n,sig-param}{mode}}{}
\pysigstopmultiline
\pysigstopsignatures
\sphinxAtStartPar
Same as utils::bounds(), but string may be a typelabel

\sphinxAtStartPar

\sphinxAtStartPar

\sphinxAtStartPar
\DUrole{versionmodified,added}{Added in version 27June2024.}

\sphinxAtStartPar
This functions adds the following case to {\hyperref[\detokenize{Developer_utils:_CPPv4I0EN9LAMMPS_NS5utils6boundsEvPKciRKNSt6stringE6bigint6bigintR4TYPER4TYPEP5Error}]{\sphinxcrossref{\sphinxcode{\sphinxupquote{utils::bounds()}}}}}:
\begin{itemize}
\item {} 
\sphinxAtStartPar
a single type label, typestr: nlo = nhi = label2type(typestr)

\end{itemize}

 
\begin{sphinxVerbatim}[commandchars=\\\{\}]
 \PYGZbs{}param file     name of source file for error message
 \PYGZbs{}param line     line number in source file for error message
 \PYGZbs{}param str      string to be processed
 \PYGZbs{}param nmin     smallest possible lower bound
 \PYGZbs{}param nmax     largest allowed upper bound
 \PYGZbs{}param nlo      lower bound
 \PYGZbs{}param nhi      upper bound
 \PYGZbs{}param lmp      pointer to top\PYGZhy{}level LAMMPS class instance
 \PYGZbs{}param mode     select labelmap using constants from Atom class
\end{sphinxVerbatim}


\end{fulllineitems}

\index{utils::expand\_args (C++ function)@\spxentry{utils::expand\_args}\spxextra{C++ function}}

\begin{fulllineitems}
\phantomsection\label{\detokenize{Developer_utils:_CPPv4N9LAMMPS_NS5utils11expand_argsEPKciiPPciRPPcP6LAMMPS}}\phantomsection\label{\detokenize{Developer_utils:_CPPv3N9LAMMPS_NS5utils11expand_argsEPKciiPPciRPPcP6LAMMPS}}\phantomsection\label{\detokenize{Developer_utils:_CPPv2N9LAMMPS_NS5utils11expand_argsEPKciiPPciRPPcP6LAMMPS}}\phantomsection\label{\detokenize{Developer_utils:LAMMPS_NS::utils::expand_args__cCP.i.i.cPP.i.cPPR.LAMMPSP}}
\pysigstartsignatures
\pysigstartmultiline
\pysiglinewithargsret{\phantomsection\label{\detokenize{Developer_utils:utils_8h_1a4672127156cfeac7cfb3376dfb7f6ca9}}\DUrole{kt}{int}\DUrole{w}{ }\sphinxcode{\sphinxupquote{\DUrole{n}{LAMMPS\_NS}\DUrole{p}{::}\DUrole{n}{utils}\DUrole{p}{::}}}\sphinxbfcode{\sphinxupquote{\DUrole{n}{expand\_args}}}}{\DUrole{k}{const}\DUrole{w}{ }\DUrole{kt}{char}\DUrole{w}{ }\DUrole{p}{*}\DUrole{n,sig-param}{file}\sphinxparamcomma \DUrole{kt}{int}\DUrole{w}{ }\DUrole{n,sig-param}{line}\sphinxparamcomma \DUrole{kt}{int}\DUrole{w}{ }\DUrole{n,sig-param}{narg}\sphinxparamcomma \DUrole{kt}{char}\DUrole{w}{ }\DUrole{p}{*}\DUrole{p}{*}\DUrole{n,sig-param}{arg}\sphinxparamcomma \DUrole{kt}{int}\DUrole{w}{ }\DUrole{n,sig-param}{mode}\sphinxparamcomma \DUrole{kt}{char}\DUrole{w}{ }\DUrole{p}{*}\DUrole{p}{*}\DUrole{p}{\&}\DUrole{n,sig-param}{earg}\sphinxparamcomma {\hyperref[\detokenize{Classes_lammps:_CPPv4N9LAMMPS_NS6LAMMPSE}]{\sphinxcrossref{\DUrole{n}{LAMMPS}}}}\DUrole{w}{ }\DUrole{p}{*}\DUrole{n,sig-param}{lmp}}{}
\pysigstopmultiline
\pysigstopsignatures
\sphinxAtStartPar
Expand list of arguments when containing fix/compute wildcards

\sphinxAtStartPar
This function searches the list of arguments in \sphinxstyleemphasis{arg} for strings of the kind c\_ID{[}*{]}, f\_ID{[}*{]}, v\_ID{[}*{]}, i2\_ID{[}*{]}, d2\_ID{[}*{]}, or c\_ID:gname:dname{[}*{]} referring to computes, fixes, vector style variables, custom per\sphinxhyphen{}atom arrays, or grids, respectively. Any such strings are replaced by one or more strings with the ‘*’ character replaced by the corresponding possible numbers as determined from the fix, compute, variable, property, or grid instance. Unrecognized strings are just copied. If the \sphinxstyleemphasis{mode} parameter is set to 0, expand global vectors, but not global arrays; if it is set to 1, expand global arrays (by column) but not global vectors.

\sphinxAtStartPar
If any expansion happens, the earg list and all its strings are new allocations and must be freed explicitly by the caller. Otherwise arg and earg will point to the same address and no explicit de\sphinxhyphen{}allocation is needed by the caller.
\begin{quote}\begin{description}
\sphinxlineitem{Parameters}\begin{itemize}
\item {} 
\sphinxAtStartPar
\sphinxstyleliteralstrong{\sphinxupquote{file}} \textendash{} name of source file for error message 

\item {} 
\sphinxAtStartPar
\sphinxstyleliteralstrong{\sphinxupquote{line}} \textendash{} line number in source file for error message 

\item {} 
\sphinxAtStartPar
\sphinxstyleliteralstrong{\sphinxupquote{narg}} \textendash{} number of arguments in current list 

\item {} 
\sphinxAtStartPar
\sphinxstyleliteralstrong{\sphinxupquote{arg}} \textendash{} argument list, possibly containing wildcards 

\item {} 
\sphinxAtStartPar
\sphinxstyleliteralstrong{\sphinxupquote{mode}} \textendash{} select between global vectors(=0) and arrays (=1) 

\item {} 
\sphinxAtStartPar
\sphinxstyleliteralstrong{\sphinxupquote{earg}} \textendash{} new argument list with wildcards expanded 

\item {} 
\sphinxAtStartPar
\sphinxstyleliteralstrong{\sphinxupquote{lmp}} \textendash{} pointer to top\sphinxhyphen{}level {\hyperref[\detokenize{Classes_lammps:classLAMMPS__NS_1_1LAMMPS}]{\sphinxcrossref{\DUrole{std,std-ref}{LAMMPS}}}} class instance 

\end{itemize}

\sphinxlineitem{Returns}
\sphinxAtStartPar
number of arguments in expanded list 

\end{description}\end{quote}

\end{fulllineitems}

\index{utils::parse\_grid\_id (C++ function)@\spxentry{utils::parse\_grid\_id}\spxextra{C++ function}}

\begin{fulllineitems}
\phantomsection\label{\detokenize{Developer_utils:_CPPv4N9LAMMPS_NS5utils13parse_grid_idEPKciRKNSt6stringEP5Error}}\phantomsection\label{\detokenize{Developer_utils:_CPPv3N9LAMMPS_NS5utils13parse_grid_idEPKciRKNSt6stringEP5Error}}\phantomsection\label{\detokenize{Developer_utils:_CPPv2N9LAMMPS_NS5utils13parse_grid_idEPKciRKNSt6stringEP5Error}}\phantomsection\label{\detokenize{Developer_utils:LAMMPS_NS::utils::parse_grid_id__cCP.i.ssCR.ErrorP}}
\pysigstartsignatures
\pysigstartmultiline
\pysiglinewithargsret{\phantomsection\label{\detokenize{Developer_utils:utils_8h_1a91407dd860f7f2a4fb3fb5c91911f6ca}}\DUrole{n}{std}\DUrole{p}{::}\DUrole{n}{vector}\DUrole{p}{\textless{}}\DUrole{n}{std}\DUrole{p}{::}\DUrole{n}{string}\DUrole{p}{\textgreater{}}\DUrole{w}{ }\sphinxcode{\sphinxupquote{\DUrole{n}{LAMMPS\_NS}\DUrole{p}{::}\DUrole{n}{utils}\DUrole{p}{::}}}\sphinxbfcode{\sphinxupquote{\DUrole{n}{parse\_grid\_id}}}}{\DUrole{k}{const}\DUrole{w}{ }\DUrole{kt}{char}\DUrole{w}{ }\DUrole{p}{*}\DUrole{n,sig-param}{file}\sphinxparamcomma \DUrole{kt}{int}\DUrole{w}{ }\DUrole{n,sig-param}{line}\sphinxparamcomma \DUrole{k}{const}\DUrole{w}{ }\DUrole{n}{std}\DUrole{p}{::}\DUrole{n}{string}\DUrole{w}{ }\DUrole{p}{\&}\DUrole{n,sig-param}{name}\sphinxparamcomma \DUrole{n}{Error}\DUrole{w}{ }\DUrole{p}{*}\DUrole{n,sig-param}{error}}{}
\pysigstopmultiline
\pysigstopsignatures
\sphinxAtStartPar
Parse grid reference into 3 sub\sphinxhyphen{}strings

\sphinxAtStartPar
Format of grid ID reference = id:gname:dname. Return vector with the 3 sub\sphinxhyphen{}strings.
\begin{quote}\begin{description}
\sphinxlineitem{Parameters}
\sphinxAtStartPar
\sphinxstyleliteralstrong{\sphinxupquote{name}} \textendash{} = complete grid ID 

\sphinxlineitem{Returns}
\sphinxAtStartPar
std::vector\textless{}std::string\textgreater{} containing the 3 sub\sphinxhyphen{}strings 

\end{description}\end{quote}

\end{fulllineitems}

\index{utils::expand\_type (C++ function)@\spxentry{utils::expand\_type}\spxextra{C++ function}}

\begin{fulllineitems}
\phantomsection\label{\detokenize{Developer_utils:_CPPv4N9LAMMPS_NS5utils11expand_typeEPKciRKNSt6stringEiP6LAMMPS}}\phantomsection\label{\detokenize{Developer_utils:_CPPv3N9LAMMPS_NS5utils11expand_typeEPKciRKNSt6stringEiP6LAMMPS}}\phantomsection\label{\detokenize{Developer_utils:_CPPv2N9LAMMPS_NS5utils11expand_typeEPKciRKNSt6stringEiP6LAMMPS}}\phantomsection\label{\detokenize{Developer_utils:LAMMPS_NS::utils::expand_type__cCP.i.ssCR.i.LAMMPSP}}
\pysigstartsignatures
\pysigstartmultiline
\pysiglinewithargsret{\phantomsection\label{\detokenize{Developer_utils:utils_8h_1ac69ac30ade492d1a68c14ce5dd92d5bc}}\DUrole{kt}{char}\DUrole{w}{ }\DUrole{p}{*}\sphinxcode{\sphinxupquote{\DUrole{n}{LAMMPS\_NS}\DUrole{p}{::}\DUrole{n}{utils}\DUrole{p}{::}}}\sphinxbfcode{\sphinxupquote{\DUrole{n}{expand\_type}}}}{\DUrole{k}{const}\DUrole{w}{ }\DUrole{kt}{char}\DUrole{w}{ }\DUrole{p}{*}\DUrole{n,sig-param}{file}\sphinxparamcomma \DUrole{kt}{int}\DUrole{w}{ }\DUrole{n,sig-param}{line}\sphinxparamcomma \DUrole{k}{const}\DUrole{w}{ }\DUrole{n}{std}\DUrole{p}{::}\DUrole{n}{string}\DUrole{w}{ }\DUrole{p}{\&}\DUrole{n,sig-param}{str}\sphinxparamcomma \DUrole{kt}{int}\DUrole{w}{ }\DUrole{n,sig-param}{mode}\sphinxparamcomma {\hyperref[\detokenize{Classes_lammps:_CPPv4N9LAMMPS_NS6LAMMPSE}]{\sphinxcrossref{\DUrole{n}{LAMMPS}}}}\DUrole{w}{ }\DUrole{p}{*}\DUrole{n,sig-param}{lmp}}{}
\pysigstopmultiline
\pysigstopsignatures
\sphinxAtStartPar
Expand type label string into its equivalent numeric type

\sphinxAtStartPar
This function checks if a given string may be a type label and then searches the labelmap type indicated by the \sphinxstyleemphasis{mode} argument for the corresponding numeric type. If this is found, a copy of the numeric type string is made and returned. Otherwise a null pointer is returned. If a string is returned, the calling code must free it with delete{[}{]}.
\begin{quote}\begin{description}
\sphinxlineitem{Parameters}\begin{itemize}
\item {} 
\sphinxAtStartPar
\sphinxstyleliteralstrong{\sphinxupquote{file}} \textendash{} name of source file for error message 

\item {} 
\sphinxAtStartPar
\sphinxstyleliteralstrong{\sphinxupquote{line}} \textendash{} line number in source file for error message 

\item {} 
\sphinxAtStartPar
\sphinxstyleliteralstrong{\sphinxupquote{str}} \textendash{} type string to be expanded 

\item {} 
\sphinxAtStartPar
\sphinxstyleliteralstrong{\sphinxupquote{mode}} \textendash{} select labelmap using constants from {\hyperref[\detokenize{Classes_atom:classLAMMPS__NS_1_1Atom}]{\sphinxcrossref{\DUrole{std,std-ref}{Atom}}}} class 

\item {} 
\sphinxAtStartPar
\sphinxstyleliteralstrong{\sphinxupquote{lmp}} \textendash{} pointer to top\sphinxhyphen{}level {\hyperref[\detokenize{Classes_lammps:classLAMMPS__NS_1_1LAMMPS}]{\sphinxcrossref{\DUrole{std,std-ref}{LAMMPS}}}} class instance 

\end{itemize}

\sphinxlineitem{Returns}
\sphinxAtStartPar
pointer to expanded string or null pointer 

\end{description}\end{quote}

\end{fulllineitems}



\subsection{Convenience functions}
\label{\detokenize{Developer_utils:convenience-functions}}\index{utils::logmesg (C++ function)@\spxentry{utils::logmesg}\spxextra{C++ function}}

\begin{fulllineitems}
\phantomsection\label{\detokenize{Developer_utils:_CPPv4IDpEN9LAMMPS_NS5utils7logmesgEvP6LAMMPSRKNSt6stringEDpRR4Args}}\phantomsection\label{\detokenize{Developer_utils:_CPPv3IDpEN9LAMMPS_NS5utils7logmesgEP6LAMMPSRKNSt6stringEDpRR4Args}}\phantomsection\label{\detokenize{Developer_utils:_CPPv2IDpEN9LAMMPS_NS5utils7logmesgEP6LAMMPSRKNSt6stringEDpRR4Args}}
\pysigstartsignatures
\pysigstartmultiline
\pysigline{\DUrole{k}{template}\DUrole{p}{\textless{}}\DUrole{k}{typename}\DUrole{w}{ }\DUrole{p}{...}\sphinxbfcode{\sphinxupquote{\DUrole{n}{Args}}}\DUrole{p}{\textgreater{}}}
\pysiglinewithargsret{\phantomsection\label{\detokenize{Developer_utils:utils_8h_1a293e69d2ef72383ce1e6373c23aab9d4}}\DUrole{kt}{void}\DUrole{w}{ }\sphinxcode{\sphinxupquote{\DUrole{n}{LAMMPS\_NS}\DUrole{p}{::}\DUrole{n}{utils}\DUrole{p}{::}}}\sphinxbfcode{\sphinxupquote{\DUrole{n}{logmesg}}}}{{\hyperref[\detokenize{Classes_lammps:_CPPv4N9LAMMPS_NS6LAMMPSE}]{\sphinxcrossref{\DUrole{n}{LAMMPS}}}}\DUrole{w}{ }\DUrole{p}{*}\DUrole{n,sig-param}{lmp}\sphinxparamcomma \DUrole{k}{const}\DUrole{w}{ }\DUrole{n}{std}\DUrole{p}{::}\DUrole{n}{string}\DUrole{w}{ }\DUrole{p}{\&}\DUrole{n,sig-param}{format}\sphinxparamcomma {\hyperref[\detokenize{Developer_utils:_CPPv4IDpEN9LAMMPS_NS5utils7logmesgEvP6LAMMPSRKNSt6stringEDpRR4Args}]{\sphinxcrossref{\DUrole{n}{Args}}}}\DUrole{p}{\&}\DUrole{p}{\&}\DUrole{p}{...}\DUrole{w}{ }\DUrole{n,sig-param}{args}}{}
\pysigstopmultiline
\pysigstopsignatures
\sphinxAtStartPar
Send formatted message to screen and logfile, if available

\sphinxAtStartPar
This function simplifies the repetitive task of outputting some message to both the screen and/or the log file. The template wrapper with fmtlib format and argument processing allows this function to work similar to \sphinxcode{\sphinxupquote{fmt::print()}}.
\begin{quote}\begin{description}
\sphinxlineitem{Parameters}\begin{itemize}
\item {} 
\sphinxAtStartPar
\sphinxstyleliteralstrong{\sphinxupquote{lmp}} \textendash{} pointer to {\hyperref[\detokenize{Classes_lammps:classLAMMPS__NS_1_1LAMMPS}]{\sphinxcrossref{\DUrole{std,std-ref}{LAMMPS}}}} class instance 

\item {} 
\sphinxAtStartPar
\sphinxstyleliteralstrong{\sphinxupquote{format}} \textendash{} format string of message to be printed 

\item {} 
\sphinxAtStartPar
\sphinxstyleliteralstrong{\sphinxupquote{args}} \textendash{} arguments to format string 

\end{itemize}

\end{description}\end{quote}

\end{fulllineitems}

\index{utils::logmesg (C++ function)@\spxentry{utils::logmesg}\spxextra{C++ function}}

\begin{fulllineitems}
\phantomsection\label{\detokenize{Developer_utils:_CPPv4N9LAMMPS_NS5utils7logmesgEP6LAMMPSRKNSt6stringE}}\phantomsection\label{\detokenize{Developer_utils:_CPPv3N9LAMMPS_NS5utils7logmesgEP6LAMMPSRKNSt6stringE}}\phantomsection\label{\detokenize{Developer_utils:_CPPv2N9LAMMPS_NS5utils7logmesgEP6LAMMPSRKNSt6stringE}}\phantomsection\label{\detokenize{Developer_utils:LAMMPS_NS::utils::logmesg__LAMMPSP.ssCR}}
\pysigstartsignatures
\pysigstartmultiline
\pysiglinewithargsret{\phantomsection\label{\detokenize{Developer_utils:utils_8h_1a0b5a7272e51011310cc877ba444f9ac7}}\DUrole{kt}{void}\DUrole{w}{ }\sphinxcode{\sphinxupquote{\DUrole{n}{LAMMPS\_NS}\DUrole{p}{::}\DUrole{n}{utils}\DUrole{p}{::}}}\sphinxbfcode{\sphinxupquote{\DUrole{n}{logmesg}}}}{{\hyperref[\detokenize{Classes_lammps:_CPPv4N9LAMMPS_NS6LAMMPSE}]{\sphinxcrossref{\DUrole{n}{LAMMPS}}}}\DUrole{w}{ }\DUrole{p}{*}\DUrole{n,sig-param}{lmp}\sphinxparamcomma \DUrole{k}{const}\DUrole{w}{ }\DUrole{n}{std}\DUrole{p}{::}\DUrole{n}{string}\DUrole{w}{ }\DUrole{p}{\&}\DUrole{n,sig-param}{mesg}}{}
\pysigstopmultiline
\pysigstopsignatures
\sphinxAtStartPar
This is an overloaded member function, provided for convenience. It differs from the above function only in what argument(s) it accepts.
\begin{quote}\begin{description}
\sphinxlineitem{Parameters}\begin{itemize}
\item {} 
\sphinxAtStartPar
\sphinxstyleliteralstrong{\sphinxupquote{lmp}} \textendash{} pointer to {\hyperref[\detokenize{Classes_lammps:classLAMMPS__NS_1_1LAMMPS}]{\sphinxcrossref{\DUrole{std,std-ref}{LAMMPS}}}} class instance 

\item {} 
\sphinxAtStartPar
\sphinxstyleliteralstrong{\sphinxupquote{mesg}} \textendash{} string with message to be printed 

\end{itemize}

\end{description}\end{quote}

\end{fulllineitems}

\index{utils::errorurl (C++ function)@\spxentry{utils::errorurl}\spxextra{C++ function}}

\begin{fulllineitems}
\phantomsection\label{\detokenize{Developer_utils:_CPPv4N9LAMMPS_NS5utils8errorurlEi}}\phantomsection\label{\detokenize{Developer_utils:_CPPv3N9LAMMPS_NS5utils8errorurlEi}}\phantomsection\label{\detokenize{Developer_utils:_CPPv2N9LAMMPS_NS5utils8errorurlEi}}\phantomsection\label{\detokenize{Developer_utils:LAMMPS_NS::utils::errorurl__i}}
\pysigstartsignatures
\pysigstartmultiline
\pysiglinewithargsret{\phantomsection\label{\detokenize{Developer_utils:utils_8h_1a0042d23b4761c1937f4cfebee627cada}}\DUrole{n}{std}\DUrole{p}{::}\DUrole{n}{string}\DUrole{w}{ }\sphinxcode{\sphinxupquote{\DUrole{n}{LAMMPS\_NS}\DUrole{p}{::}\DUrole{n}{utils}\DUrole{p}{::}}}\sphinxbfcode{\sphinxupquote{\DUrole{n}{errorurl}}}}{\DUrole{kt}{int}\DUrole{w}{ }\DUrole{n,sig-param}{errorcode}}{}
\pysigstopmultiline
\pysigstopsignatures
\sphinxAtStartPar
Return text redirecting the user to a specific paragraph in the manual

\sphinxAtStartPar
The {\hyperref[\detokenize{Classes_lammps:classLAMMPS__NS_1_1LAMMPS}]{\sphinxcrossref{\DUrole{std,std-ref}{LAMMPS}}}} manual contains detailed explanations for errors and warnings where a simple error message may not be sufficient. These can be reached through URLs with a numeric code. This function creates the corresponding text to be included into the error message that redirects the user to that URL.
\begin{quote}\begin{description}
\sphinxlineitem{Parameters}
\sphinxAtStartPar
\sphinxstyleliteralstrong{\sphinxupquote{errorcode}} \textendash{} number pointing to a paragraph in the manual 

\end{description}\end{quote}

\end{fulllineitems}

\index{utils::missing\_cmd\_args (C++ function)@\spxentry{utils::missing\_cmd\_args}\spxextra{C++ function}}

\begin{fulllineitems}
\phantomsection\label{\detokenize{Developer_utils:_CPPv4N9LAMMPS_NS5utils16missing_cmd_argsERKNSt6stringEiRKNSt6stringEP5Error}}\phantomsection\label{\detokenize{Developer_utils:_CPPv3N9LAMMPS_NS5utils16missing_cmd_argsERKNSt6stringEiRKNSt6stringEP5Error}}\phantomsection\label{\detokenize{Developer_utils:_CPPv2N9LAMMPS_NS5utils16missing_cmd_argsERKNSt6stringEiRKNSt6stringEP5Error}}\phantomsection\label{\detokenize{Developer_utils:LAMMPS_NS::utils::missing_cmd_args__ssCR.i.ssCR.ErrorP}}
\pysigstartsignatures
\pysigstartmultiline
\pysiglinewithargsret{\phantomsection\label{\detokenize{Developer_utils:utils_8h_1a9ebdf3470338d3fddf4d0462b440bb8b}}\DUrole{kt}{void}\DUrole{w}{ }\sphinxcode{\sphinxupquote{\DUrole{n}{LAMMPS\_NS}\DUrole{p}{::}\DUrole{n}{utils}\DUrole{p}{::}}}\sphinxbfcode{\sphinxupquote{\DUrole{n}{missing\_cmd\_args}}}}{\DUrole{k}{const}\DUrole{w}{ }\DUrole{n}{std}\DUrole{p}{::}\DUrole{n}{string}\DUrole{w}{ }\DUrole{p}{\&}\DUrole{n,sig-param}{file}\sphinxparamcomma \DUrole{kt}{int}\DUrole{w}{ }\DUrole{n,sig-param}{line}\sphinxparamcomma \DUrole{k}{const}\DUrole{w}{ }\DUrole{n}{std}\DUrole{p}{::}\DUrole{n}{string}\DUrole{w}{ }\DUrole{p}{\&}\DUrole{n,sig-param}{cmd}\sphinxparamcomma \DUrole{n}{Error}\DUrole{w}{ }\DUrole{p}{*}\DUrole{n,sig-param}{error}}{}
\pysigstopmultiline
\pysigstopsignatures
\sphinxAtStartPar
Print error message about missing arguments for command

\sphinxAtStartPar
This function simplifies the repetitive reporting missing arguments to a command.
\begin{quote}\begin{description}
\sphinxlineitem{Parameters}\begin{itemize}
\item {} 
\sphinxAtStartPar
\sphinxstyleliteralstrong{\sphinxupquote{file}} \textendash{} name of source file for error message 

\item {} 
\sphinxAtStartPar
\sphinxstyleliteralstrong{\sphinxupquote{line}} \textendash{} line number in source file for error message 

\item {} 
\sphinxAtStartPar
\sphinxstyleliteralstrong{\sphinxupquote{cmd}} \textendash{} name of the failing command 

\item {} 
\sphinxAtStartPar
\sphinxstyleliteralstrong{\sphinxupquote{error}} \textendash{} pointer to Error class instance (for abort) or nullptr 

\end{itemize}

\end{description}\end{quote}

\end{fulllineitems}

\index{utils::flush\_buffers (C++ function)@\spxentry{utils::flush\_buffers}\spxextra{C++ function}}

\begin{fulllineitems}
\phantomsection\label{\detokenize{Developer_utils:_CPPv4N9LAMMPS_NS5utils13flush_buffersEP6LAMMPS}}\phantomsection\label{\detokenize{Developer_utils:_CPPv3N9LAMMPS_NS5utils13flush_buffersEP6LAMMPS}}\phantomsection\label{\detokenize{Developer_utils:_CPPv2N9LAMMPS_NS5utils13flush_buffersEP6LAMMPS}}\phantomsection\label{\detokenize{Developer_utils:LAMMPS_NS::utils::flush_buffers__LAMMPSP}}
\pysigstartsignatures
\pysigstartmultiline
\pysiglinewithargsret{\phantomsection\label{\detokenize{Developer_utils:utils_8h_1a1721a8d265f9cde71f2c20ac1ba581e5}}\DUrole{kt}{void}\DUrole{w}{ }\sphinxcode{\sphinxupquote{\DUrole{n}{LAMMPS\_NS}\DUrole{p}{::}\DUrole{n}{utils}\DUrole{p}{::}}}\sphinxbfcode{\sphinxupquote{\DUrole{n}{flush\_buffers}}}}{{\hyperref[\detokenize{Classes_lammps:_CPPv4N9LAMMPS_NS6LAMMPSE}]{\sphinxcrossref{\DUrole{n}{LAMMPS}}}}\DUrole{w}{ }\DUrole{p}{*}\DUrole{n,sig-param}{lmp}}{}
\pysigstopmultiline
\pysigstopsignatures
\sphinxAtStartPar
Flush output buffers

\sphinxAtStartPar
This function calls fflush() on screen and logfile FILE pointers if available and thus tells the operating system to output all currently buffered data. This is local operation and independent from buffering by a file system or an MPI library. 

\end{fulllineitems}

\index{utils::getsyserror (C++ function)@\spxentry{utils::getsyserror}\spxextra{C++ function}}

\begin{fulllineitems}
\phantomsection\label{\detokenize{Developer_utils:_CPPv4N9LAMMPS_NS5utils11getsyserrorEv}}\phantomsection\label{\detokenize{Developer_utils:_CPPv3N9LAMMPS_NS5utils11getsyserrorEv}}\phantomsection\label{\detokenize{Developer_utils:_CPPv2N9LAMMPS_NS5utils11getsyserrorEv}}\phantomsection\label{\detokenize{Developer_utils:LAMMPS_NS::utils::getsyserror}}
\pysigstartsignatures
\pysigstartmultiline
\pysiglinewithargsret{\phantomsection\label{\detokenize{Developer_utils:utils_8h_1a59872d3df515856d461cc81f8de6777e}}\DUrole{n}{std}\DUrole{p}{::}\DUrole{n}{string}\DUrole{w}{ }\sphinxcode{\sphinxupquote{\DUrole{n}{LAMMPS\_NS}\DUrole{p}{::}\DUrole{n}{utils}\DUrole{p}{::}}}\sphinxbfcode{\sphinxupquote{\DUrole{n}{getsyserror}}}}{}{}
\pysigstopmultiline
\pysigstopsignatures
\sphinxAtStartPar
Return a string representing the current system error status

\sphinxAtStartPar
This is a wrapper around calling strerror(errno).
\begin{quote}\begin{description}
\sphinxlineitem{Returns}
\sphinxAtStartPar
error string 

\end{description}\end{quote}

\end{fulllineitems}

\index{utils::check\_packages\_for\_style (C++ function)@\spxentry{utils::check\_packages\_for\_style}\spxextra{C++ function}}

\begin{fulllineitems}
\phantomsection\label{\detokenize{Developer_utils:_CPPv4N9LAMMPS_NS5utils24check_packages_for_styleERKNSt6stringERKNSt6stringEP6LAMMPS}}\phantomsection\label{\detokenize{Developer_utils:_CPPv3N9LAMMPS_NS5utils24check_packages_for_styleERKNSt6stringERKNSt6stringEP6LAMMPS}}\phantomsection\label{\detokenize{Developer_utils:_CPPv2N9LAMMPS_NS5utils24check_packages_for_styleERKNSt6stringERKNSt6stringEP6LAMMPS}}\phantomsection\label{\detokenize{Developer_utils:LAMMPS_NS::utils::check_packages_for_style__ssCR.ssCR.LAMMPSP}}
\pysigstartsignatures
\pysigstartmultiline
\pysiglinewithargsret{\phantomsection\label{\detokenize{Developer_utils:utils_8h_1a7958600997ce9aecc9c589d6a340b93b}}\DUrole{n}{std}\DUrole{p}{::}\DUrole{n}{string}\DUrole{w}{ }\sphinxcode{\sphinxupquote{\DUrole{n}{LAMMPS\_NS}\DUrole{p}{::}\DUrole{n}{utils}\DUrole{p}{::}}}\sphinxbfcode{\sphinxupquote{\DUrole{n}{check\_packages\_for\_style}}}}{\DUrole{k}{const}\DUrole{w}{ }\DUrole{n}{std}\DUrole{p}{::}\DUrole{n}{string}\DUrole{w}{ }\DUrole{p}{\&}\DUrole{n,sig-param}{style}\sphinxparamcomma \DUrole{k}{const}\DUrole{w}{ }\DUrole{n}{std}\DUrole{p}{::}\DUrole{n}{string}\DUrole{w}{ }\DUrole{p}{\&}\DUrole{n,sig-param}{name}\sphinxparamcomma {\hyperref[\detokenize{Classes_lammps:_CPPv4N9LAMMPS_NS6LAMMPSE}]{\sphinxcrossref{\DUrole{n}{LAMMPS}}}}\DUrole{w}{ }\DUrole{p}{*}\DUrole{n,sig-param}{lmp}}{}
\pysigstopmultiline
\pysigstopsignatures
\sphinxAtStartPar
Report if a requested style is in a package or may have a typo
\begin{quote}\begin{description}
\sphinxlineitem{Parameters}\begin{itemize}
\item {} 
\sphinxAtStartPar
\sphinxstyleliteralstrong{\sphinxupquote{style}} \textendash{} type of style that is to be checked for 

\item {} 
\sphinxAtStartPar
\sphinxstyleliteralstrong{\sphinxupquote{name}} \textendash{} name of style that was not found 

\item {} 
\sphinxAtStartPar
\sphinxstyleliteralstrong{\sphinxupquote{lmp}} \textendash{} pointer to top\sphinxhyphen{}level {\hyperref[\detokenize{Classes_lammps:classLAMMPS__NS_1_1LAMMPS}]{\sphinxcrossref{\DUrole{std,std-ref}{LAMMPS}}}} class instance 

\end{itemize}

\sphinxlineitem{Returns}
\sphinxAtStartPar
string usable for error messages 

\end{description}\end{quote}

\end{fulllineitems}

\index{utils::timespec2seconds (C++ function)@\spxentry{utils::timespec2seconds}\spxextra{C++ function}}

\begin{fulllineitems}
\phantomsection\label{\detokenize{Developer_utils:_CPPv4N9LAMMPS_NS5utils16timespec2secondsERKNSt6stringE}}\phantomsection\label{\detokenize{Developer_utils:_CPPv3N9LAMMPS_NS5utils16timespec2secondsERKNSt6stringE}}\phantomsection\label{\detokenize{Developer_utils:_CPPv2N9LAMMPS_NS5utils16timespec2secondsERKNSt6stringE}}\phantomsection\label{\detokenize{Developer_utils:LAMMPS_NS::utils::timespec2seconds__ssCR}}
\pysigstartsignatures
\pysigstartmultiline
\pysiglinewithargsret{\phantomsection\label{\detokenize{Developer_utils:utils_8h_1afd22548fe9c6d4570ffcde59947a1764}}\DUrole{kt}{double}\DUrole{w}{ }\sphinxcode{\sphinxupquote{\DUrole{n}{LAMMPS\_NS}\DUrole{p}{::}\DUrole{n}{utils}\DUrole{p}{::}}}\sphinxbfcode{\sphinxupquote{\DUrole{n}{timespec2seconds}}}}{\DUrole{k}{const}\DUrole{w}{ }\DUrole{n}{std}\DUrole{p}{::}\DUrole{n}{string}\DUrole{w}{ }\DUrole{p}{\&}\DUrole{n,sig-param}{timespec}}{}
\pysigstopmultiline
\pysigstopsignatures
\sphinxAtStartPar
Convert a time string to seconds

\sphinxAtStartPar
The strings “off” and “unlimited” result in \sphinxhyphen{}1
\begin{quote}\begin{description}
\sphinxlineitem{Parameters}
\sphinxAtStartPar
\sphinxstyleliteralstrong{\sphinxupquote{timespec}} \textendash{} a string in the following format: ({[}{[}HH:{]}MM:{]}SS) 

\sphinxlineitem{Returns}
\sphinxAtStartPar
total in seconds 

\end{description}\end{quote}

\end{fulllineitems}

\index{utils::date2num (C++ function)@\spxentry{utils::date2num}\spxextra{C++ function}}

\begin{fulllineitems}
\phantomsection\label{\detokenize{Developer_utils:_CPPv4N9LAMMPS_NS5utils8date2numERKNSt6stringE}}\phantomsection\label{\detokenize{Developer_utils:_CPPv3N9LAMMPS_NS5utils8date2numERKNSt6stringE}}\phantomsection\label{\detokenize{Developer_utils:_CPPv2N9LAMMPS_NS5utils8date2numERKNSt6stringE}}\phantomsection\label{\detokenize{Developer_utils:LAMMPS_NS::utils::date2num__ssCR}}
\pysigstartsignatures
\pysigstartmultiline
\pysiglinewithargsret{\phantomsection\label{\detokenize{Developer_utils:utils_8h_1af07ce9a5537cbc035845686e894f3491}}\DUrole{kt}{int}\DUrole{w}{ }\sphinxcode{\sphinxupquote{\DUrole{n}{LAMMPS\_NS}\DUrole{p}{::}\DUrole{n}{utils}\DUrole{p}{::}}}\sphinxbfcode{\sphinxupquote{\DUrole{n}{date2num}}}}{\DUrole{k}{const}\DUrole{w}{ }\DUrole{n}{std}\DUrole{p}{::}\DUrole{n}{string}\DUrole{w}{ }\DUrole{p}{\&}\DUrole{n,sig-param}{date}}{}
\pysigstopmultiline
\pysigstopsignatures
\sphinxAtStartPar
Convert a {\hyperref[\detokenize{Classes_lammps:classLAMMPS__NS_1_1LAMMPS}]{\sphinxcrossref{\DUrole{std,std-ref}{LAMMPS}}}} version date to a number

\sphinxAtStartPar
This will generate a number YYYYMMDD from a date string (with or without blanks) that is suitable for numerical comparisons, i.e. later dates will generate a larger number.

\sphinxAtStartPar
The day may or may not have a leading zero, the month is identified by the first 3 letters (so there may be more) and the year may be 2 or 4 digits (the missing 2 digits will be assumed as 20. That is 04 corresponds to 2004).

\sphinxAtStartPar
No check is made whether the date is valid.
\begin{quote}\begin{description}
\sphinxlineitem{Parameters}
\sphinxAtStartPar
\sphinxstyleliteralstrong{\sphinxupquote{date}} \textendash{} string in the format (Day Month Year) 

\sphinxlineitem{Returns}
\sphinxAtStartPar
date code 

\end{description}\end{quote}

\end{fulllineitems}

\index{utils::current\_date (C++ function)@\spxentry{utils::current\_date}\spxextra{C++ function}}

\begin{fulllineitems}
\phantomsection\label{\detokenize{Developer_utils:_CPPv4N9LAMMPS_NS5utils12current_dateEv}}\phantomsection\label{\detokenize{Developer_utils:_CPPv3N9LAMMPS_NS5utils12current_dateEv}}\phantomsection\label{\detokenize{Developer_utils:_CPPv2N9LAMMPS_NS5utils12current_dateEv}}\phantomsection\label{\detokenize{Developer_utils:LAMMPS_NS::utils::current_date}}
\pysigstartsignatures
\pysigstartmultiline
\pysiglinewithargsret{\phantomsection\label{\detokenize{Developer_utils:utils_8h_1a7a7a3c2717c15c16a93621e9e7d0f288}}\DUrole{n}{std}\DUrole{p}{::}\DUrole{n}{string}\DUrole{w}{ }\sphinxcode{\sphinxupquote{\DUrole{n}{LAMMPS\_NS}\DUrole{p}{::}\DUrole{n}{utils}\DUrole{p}{::}}}\sphinxbfcode{\sphinxupquote{\DUrole{n}{current\_date}}}}{}{}
\pysigstopmultiline
\pysigstopsignatures
\sphinxAtStartPar
Return current date as string

\sphinxAtStartPar
This will generate a string containing the current date in YYYY\sphinxhyphen{}MM\sphinxhyphen{}DD format.
\begin{quote}\begin{description}
\sphinxlineitem{Returns}
\sphinxAtStartPar
string with current date 

\end{description}\end{quote}

\end{fulllineitems}



\subsection{Customized standard functions}
\label{\detokenize{Developer_utils:customized-standard-functions}}\index{utils::binary\_search (C++ function)@\spxentry{utils::binary\_search}\spxextra{C++ function}}

\begin{fulllineitems}
\phantomsection\label{\detokenize{Developer_utils:_CPPv4N9LAMMPS_NS5utils13binary_searchEKdKiPKd}}\phantomsection\label{\detokenize{Developer_utils:_CPPv3N9LAMMPS_NS5utils13binary_searchEKdKiPKd}}\phantomsection\label{\detokenize{Developer_utils:_CPPv2N9LAMMPS_NS5utils13binary_searchEKdKiPKd}}\phantomsection\label{\detokenize{Developer_utils:LAMMPS_NS::utils::binary_search__doubleC.iC.doubleCP}}
\pysigstartsignatures
\pysigstartmultiline
\pysiglinewithargsret{\phantomsection\label{\detokenize{Developer_utils:utils_8h_1a57ddd765821bb497c07a157b66725a01}}\DUrole{kt}{int}\DUrole{w}{ }\sphinxcode{\sphinxupquote{\DUrole{n}{LAMMPS\_NS}\DUrole{p}{::}\DUrole{n}{utils}\DUrole{p}{::}}}\sphinxbfcode{\sphinxupquote{\DUrole{n}{binary\_search}}}}{\DUrole{k}{const}\DUrole{w}{ }\DUrole{kt}{double}\DUrole{w}{ }\DUrole{n,sig-param}{needle}\sphinxparamcomma \DUrole{k}{const}\DUrole{w}{ }\DUrole{kt}{int}\DUrole{w}{ }\DUrole{n,sig-param}{n}\sphinxparamcomma \DUrole{k}{const}\DUrole{w}{ }\DUrole{kt}{double}\DUrole{w}{ }\DUrole{p}{*}\DUrole{n,sig-param}{haystack}}{}
\pysigstopmultiline
\pysigstopsignatures
\sphinxAtStartPar
Binary search in a vector of ascending doubles of length N

\sphinxAtStartPar
If the value is smaller than the smallest value in the vector, 0 is returned. If the value is larger or equal than the largest value in the vector, N\sphinxhyphen{}1 is returned. Otherwise the index that satisfies the condition

\sphinxAtStartPar
haystack{[}index{]} \textless{}= value \textless{} haystack{[}index+1{]}

\sphinxAtStartPar
is returned, i.e. a value from 1 to N\sphinxhyphen{}2. Note that if there are tied values in the haystack, always the larger index is returned as only that satisfied the condition.
\begin{quote}\begin{description}
\sphinxlineitem{Parameters}\begin{itemize}
\item {} 
\sphinxAtStartPar
\sphinxstyleliteralstrong{\sphinxupquote{needle}} \textendash{} search value for which are are looking for the closest index 

\item {} 
\sphinxAtStartPar
\sphinxstyleliteralstrong{\sphinxupquote{n}} \textendash{} size of the haystack array 

\item {} 
\sphinxAtStartPar
\sphinxstyleliteralstrong{\sphinxupquote{haystack}} \textendash{} array with data in ascending order. 

\end{itemize}

\sphinxlineitem{Returns}
\sphinxAtStartPar
index of value in the haystack array smaller or equal to needle 

\end{description}\end{quote}

\end{fulllineitems}

\index{utils::merge\_sort (C++ function)@\spxentry{utils::merge\_sort}\spxextra{C++ function}}

\begin{fulllineitems}
\phantomsection\label{\detokenize{Developer_utils:_CPPv4N9LAMMPS_NS5utils10merge_sortEPiiPvPFiiiPvE}}\phantomsection\label{\detokenize{Developer_utils:_CPPv3N9LAMMPS_NS5utils10merge_sortEPiiPvPFiiiPvE}}\phantomsection\label{\detokenize{Developer_utils:_CPPv2N9LAMMPS_NS5utils10merge_sortEPiiPvPFiiiPvE}}
\pysigstartsignatures
\pysigstartmultiline
\pysiglinewithargsret{\phantomsection\label{\detokenize{Developer_utils:utils_8h_1af677fa53e7e5713dfb0fa6d6c3c6f81e}}\DUrole{kt}{void}\DUrole{w}{ }\sphinxcode{\sphinxupquote{\DUrole{n}{LAMMPS\_NS}\DUrole{p}{::}\DUrole{n}{utils}\DUrole{p}{::}}}\sphinxbfcode{\sphinxupquote{\DUrole{n}{merge\_sort}}}}{\DUrole{kt}{int}\DUrole{w}{ }\DUrole{p}{*}\DUrole{n,sig-param}{index}\sphinxparamcomma \DUrole{kt}{int}\DUrole{w}{ }\DUrole{n,sig-param}{num}\sphinxparamcomma \DUrole{kt}{void}\DUrole{w}{ }\DUrole{p}{*}\DUrole{n,sig-param}{ptr}\sphinxparamcomma \DUrole{kt}{int}\DUrole{w}{ }\DUrole{p}{(}\DUrole{p}{*}\DUrole{n,sig-param}{comp}\DUrole{p}{)}\DUrole{p}{(}\DUrole{kt}{int}\DUrole{p}{,}\DUrole{w}{ }\DUrole{kt}{int}\DUrole{p}{,}\DUrole{w}{ }\DUrole{kt}{void}\DUrole{p}{*}\DUrole{p}{)}}{}
\pysigstopmultiline
\pysigstopsignatures
\sphinxAtStartPar
Custom merge sort implementation

\sphinxAtStartPar
This function provides a custom upward hybrid merge sort implementation with support to pass an opaque pointer to the comparison function, e.g. for access to class members. This avoids having to use global variables. For improved performance, it uses an in\sphinxhyphen{}place insertion sort on initial chunks of up to 64 elements and switches to merge sort from then on.
\begin{quote}\begin{description}
\sphinxlineitem{Parameters}\begin{itemize}
\item {} 
\sphinxAtStartPar
\sphinxstyleliteralstrong{\sphinxupquote{index}} \textendash{} Array with indices to be sorted 

\item {} 
\sphinxAtStartPar
\sphinxstyleliteralstrong{\sphinxupquote{num}} \textendash{} Length of the index array 

\item {} 
\sphinxAtStartPar
\sphinxstyleliteralstrong{\sphinxupquote{ptr}} \textendash{} Pointer to opaque object passed to comparison function 

\item {} 
\sphinxAtStartPar
\sphinxstyleliteralstrong{\sphinxupquote{comp}} \textendash{} Pointer to comparison function 

\end{itemize}

\end{description}\end{quote}

\end{fulllineitems}



\bigskip\hrule\bigskip



\section{Special Math functions}
\label{\detokenize{Developer_utils:special-math-functions}}
\sphinxAtStartPar
The \sphinxcode{\sphinxupquote{MathSpecial}} namespace implements a selection of custom and optimized
mathematical functions for a variety of applications.
\index{MathSpecial::factorial (C++ function)@\spxentry{MathSpecial::factorial}\spxextra{C++ function}}

\begin{fulllineitems}
\phantomsection\label{\detokenize{Developer_utils:_CPPv4N9LAMMPS_NS11MathSpecial9factorialEKi}}\phantomsection\label{\detokenize{Developer_utils:_CPPv3N9LAMMPS_NS11MathSpecial9factorialEKi}}\phantomsection\label{\detokenize{Developer_utils:_CPPv2N9LAMMPS_NS11MathSpecial9factorialEKi}}\phantomsection\label{\detokenize{Developer_utils:LAMMPS_NS::MathSpecial::factorial__iC}}
\pysigstartsignatures
\pysigstartmultiline
\pysiglinewithargsret{\phantomsection\label{\detokenize{Developer_utils:math__special_8h_1a9a6be409a4fa9407a40b40ab6d226a8f}}\DUrole{kt}{double}\DUrole{w}{ }\sphinxcode{\sphinxupquote{\DUrole{n}{LAMMPS\_NS}\DUrole{p}{::}\DUrole{n}{MathSpecial}\DUrole{p}{::}}}\sphinxbfcode{\sphinxupquote{\DUrole{n}{factorial}}}}{\DUrole{k}{const}\DUrole{w}{ }\DUrole{kt}{int}\DUrole{w}{ }\DUrole{n,sig-param}{n}}{}
\pysigstopmultiline
\pysigstopsignatures
\sphinxAtStartPar
Fast tabulated factorial function

\sphinxAtStartPar
This function looks up pre\sphinxhyphen{}computed factorial values for arguments of n = 0 to a maximum of 167, which is the maximal value representable by a double precision floating point number. For other values of n a NaN value is returned.
\begin{quote}\begin{description}
\sphinxlineitem{Parameters}
\sphinxAtStartPar
\sphinxstyleliteralstrong{\sphinxupquote{n}} \textendash{} argument (valid: 0 \textless{}= n \textless{}= 167) 

\sphinxlineitem{Returns}
\sphinxAtStartPar
value of n! as double precision number or NaN 

\end{description}\end{quote}

\end{fulllineitems}

\index{MathSpecial::exp2\_x86 (C++ function)@\spxentry{MathSpecial::exp2\_x86}\spxextra{C++ function}}

\begin{fulllineitems}
\phantomsection\label{\detokenize{Developer_utils:_CPPv4N9LAMMPS_NS11MathSpecial8exp2_x86Ed}}\phantomsection\label{\detokenize{Developer_utils:_CPPv3N9LAMMPS_NS11MathSpecial8exp2_x86Ed}}\phantomsection\label{\detokenize{Developer_utils:_CPPv2N9LAMMPS_NS11MathSpecial8exp2_x86Ed}}\phantomsection\label{\detokenize{Developer_utils:LAMMPS_NS::MathSpecial::exp2_x86__double}}
\pysigstartsignatures
\pysigstartmultiline
\pysiglinewithargsret{\phantomsection\label{\detokenize{Developer_utils:math__special_8h_1a46561242143bc684f0307a2260686505}}\DUrole{kt}{double}\DUrole{w}{ }\sphinxcode{\sphinxupquote{\DUrole{n}{LAMMPS\_NS}\DUrole{p}{::}\DUrole{n}{MathSpecial}\DUrole{p}{::}}}\sphinxbfcode{\sphinxupquote{\DUrole{n}{exp2\_x86}}}}{\DUrole{kt}{double}\DUrole{w}{ }\DUrole{n,sig-param}{x}}{}
\pysigstopmultiline
\pysigstopsignatures
\sphinxAtStartPar
Fast implementation of 2\textasciicircum{}x without argument checks for little endian CPUs

\sphinxAtStartPar
This function implements an optimized version of pow(2.0, x) that does not check for valid arguments and thus may only be used where arguments are well behaved. The implementation makes assumptions about the layout of double precision floating point numbers in memory and thus will only work on little endian CPUs. If little endian cannot be safely detected, the result of calling pow(2.0, x) will be returned. This function also is the basis for the fast exponential fm\_exp(x).
\begin{quote}\begin{description}
\sphinxlineitem{Parameters}
\sphinxAtStartPar
\sphinxstyleliteralstrong{\sphinxupquote{x}} \textendash{} argument 

\sphinxlineitem{Returns}
\sphinxAtStartPar
value of 2\textasciicircum{}x as double precision number 

\end{description}\end{quote}

\end{fulllineitems}

\index{MathSpecial::fm\_exp (C++ function)@\spxentry{MathSpecial::fm\_exp}\spxextra{C++ function}}

\begin{fulllineitems}
\phantomsection\label{\detokenize{Developer_utils:_CPPv4N9LAMMPS_NS11MathSpecial6fm_expEd}}\phantomsection\label{\detokenize{Developer_utils:_CPPv3N9LAMMPS_NS11MathSpecial6fm_expEd}}\phantomsection\label{\detokenize{Developer_utils:_CPPv2N9LAMMPS_NS11MathSpecial6fm_expEd}}\phantomsection\label{\detokenize{Developer_utils:LAMMPS_NS::MathSpecial::fm_exp__double}}
\pysigstartsignatures
\pysigstartmultiline
\pysiglinewithargsret{\phantomsection\label{\detokenize{Developer_utils:math__special_8h_1a918af8d4a779a4da394bf9cc4a92d628}}\DUrole{kt}{double}\DUrole{w}{ }\sphinxcode{\sphinxupquote{\DUrole{n}{LAMMPS\_NS}\DUrole{p}{::}\DUrole{n}{MathSpecial}\DUrole{p}{::}}}\sphinxbfcode{\sphinxupquote{\DUrole{n}{fm\_exp}}}}{\DUrole{kt}{double}\DUrole{w}{ }\DUrole{n,sig-param}{x}}{}
\pysigstopmultiline
\pysigstopsignatures
\sphinxAtStartPar
Fast implementation of exp(x) for little endian CPUs

\sphinxAtStartPar
This function implements an optimized version of exp(x) for little endian CPUs. It calls the exp2\_x86(x) function with a suitable prefactor to x to return exp(x). The implementation makes assumptions about the layout of double precision floating point numbers in memory and thus will only work on little endian CPUs. If little endian cannot be safely detected, the result of calling the exp(x) implementation in the standard math library will be returned.
\begin{quote}\begin{description}
\sphinxlineitem{Parameters}
\sphinxAtStartPar
\sphinxstyleliteralstrong{\sphinxupquote{x}} \textendash{} argument 

\sphinxlineitem{Returns}
\sphinxAtStartPar
value of e\textasciicircum{}x as double precision number 

\end{description}\end{quote}

\end{fulllineitems}

\index{MathSpecial::my\_erfcx (C++ function)@\spxentry{MathSpecial::my\_erfcx}\spxextra{C++ function}}

\begin{fulllineitems}
\phantomsection\label{\detokenize{Developer_utils:_CPPv4N9LAMMPS_NS11MathSpecial8my_erfcxEKd}}\phantomsection\label{\detokenize{Developer_utils:_CPPv3N9LAMMPS_NS11MathSpecial8my_erfcxEKd}}\phantomsection\label{\detokenize{Developer_utils:_CPPv2N9LAMMPS_NS11MathSpecial8my_erfcxEKd}}\phantomsection\label{\detokenize{Developer_utils:LAMMPS_NS::MathSpecial::my_erfcx__doubleC}}
\pysigstartsignatures
\pysigstartmultiline
\pysiglinewithargsret{\phantomsection\label{\detokenize{Developer_utils:math__special_8h_1a4f48d97dbd22584d59e064996cbd0694}}\DUrole{k}{static}\DUrole{w}{ }\DUrole{k}{inline}\DUrole{w}{ }\DUrole{kt}{double}\DUrole{w}{ }\sphinxcode{\sphinxupquote{\DUrole{n}{LAMMPS\_NS}\DUrole{p}{::}\DUrole{n}{MathSpecial}\DUrole{p}{::}}}\sphinxbfcode{\sphinxupquote{\DUrole{n}{my\_erfcx}}}}{\DUrole{k}{const}\DUrole{w}{ }\DUrole{kt}{double}\DUrole{w}{ }\DUrole{n,sig-param}{x}}{}
\pysigstopmultiline
\pysigstopsignatures
\sphinxAtStartPar
Fast scaled error function complement exp(x*x)*erfc(x) for coul/long styles

\sphinxAtStartPar
This is a portable fast implementation of exp(x*x)*erfc(x) that can be used in coul/long pair styles as a replacement for the polynomial expansion that is/was widely used. Unlike the polynomial expansion, that is only accurate at the level of single precision floating point it provides full double precision accuracy, but at comparable speed (unlike the erfc() implementation shipped with GNU standard math library).
\begin{quote}\begin{description}
\sphinxlineitem{Parameters}
\sphinxAtStartPar
\sphinxstyleliteralstrong{\sphinxupquote{x}} \textendash{} argument 

\sphinxlineitem{Returns}
\sphinxAtStartPar
value of e\textasciicircum{}(x*x)*erfc(x) 

\end{description}\end{quote}

\end{fulllineitems}

\index{MathSpecial::expmsq (C++ function)@\spxentry{MathSpecial::expmsq}\spxextra{C++ function}}

\begin{fulllineitems}
\phantomsection\label{\detokenize{Developer_utils:_CPPv4N9LAMMPS_NS11MathSpecial6expmsqEd}}\phantomsection\label{\detokenize{Developer_utils:_CPPv3N9LAMMPS_NS11MathSpecial6expmsqEd}}\phantomsection\label{\detokenize{Developer_utils:_CPPv2N9LAMMPS_NS11MathSpecial6expmsqEd}}\phantomsection\label{\detokenize{Developer_utils:LAMMPS_NS::MathSpecial::expmsq__double}}
\pysigstartsignatures
\pysigstartmultiline
\pysiglinewithargsret{\phantomsection\label{\detokenize{Developer_utils:math__special_8h_1a7ccf069f22f9a4e9dfeed7e04f3ca447}}\DUrole{k}{static}\DUrole{w}{ }\DUrole{k}{inline}\DUrole{w}{ }\DUrole{kt}{double}\DUrole{w}{ }\sphinxcode{\sphinxupquote{\DUrole{n}{LAMMPS\_NS}\DUrole{p}{::}\DUrole{n}{MathSpecial}\DUrole{p}{::}}}\sphinxbfcode{\sphinxupquote{\DUrole{n}{expmsq}}}}{\DUrole{kt}{double}\DUrole{w}{ }\DUrole{n,sig-param}{x}}{}
\pysigstopmultiline
\pysigstopsignatures
\sphinxAtStartPar
Fast implementation of exp(\sphinxhyphen{}x*x) for little endian CPUs for coul/long styles

\sphinxAtStartPar
This function implements an optimized version of exp(\sphinxhyphen{}x*x) based on exp2\_x86() for use with little endian CPUs. If little endian cannot be safely detected, the result of calling the exp(\sphinxhyphen{}x*x) implementation in the standard math library will be returned.
\begin{quote}\begin{description}
\sphinxlineitem{Parameters}
\sphinxAtStartPar
\sphinxstyleliteralstrong{\sphinxupquote{x}} \textendash{} argument 

\sphinxlineitem{Returns}
\sphinxAtStartPar
value of e\textasciicircum{}(\sphinxhyphen{}x*x) as double precision number 

\end{description}\end{quote}

\end{fulllineitems}

\index{MathSpecial::square (C++ function)@\spxentry{MathSpecial::square}\spxextra{C++ function}}

\begin{fulllineitems}
\phantomsection\label{\detokenize{Developer_utils:_CPPv4N9LAMMPS_NS11MathSpecial6squareERKd}}\phantomsection\label{\detokenize{Developer_utils:_CPPv3N9LAMMPS_NS11MathSpecial6squareERKd}}\phantomsection\label{\detokenize{Developer_utils:_CPPv2N9LAMMPS_NS11MathSpecial6squareERKd}}\phantomsection\label{\detokenize{Developer_utils:LAMMPS_NS::MathSpecial::square__doubleCR}}
\pysigstartsignatures
\pysigstartmultiline
\pysiglinewithargsret{\phantomsection\label{\detokenize{Developer_utils:math__special_8h_1aa7dd862a7b7b244516c8ea1de509ee09}}\DUrole{k}{static}\DUrole{w}{ }\DUrole{k}{inline}\DUrole{w}{ }\DUrole{kt}{double}\DUrole{w}{ }\sphinxcode{\sphinxupquote{\DUrole{n}{LAMMPS\_NS}\DUrole{p}{::}\DUrole{n}{MathSpecial}\DUrole{p}{::}}}\sphinxbfcode{\sphinxupquote{\DUrole{n}{square}}}}{\DUrole{k}{const}\DUrole{w}{ }\DUrole{kt}{double}\DUrole{w}{ }\DUrole{p}{\&}\DUrole{n,sig-param}{x}}{}
\pysigstopmultiline
\pysigstopsignatures
\sphinxAtStartPar
Fast inline version of pow(x, 2.0)
\begin{quote}\begin{description}
\sphinxlineitem{Parameters}
\sphinxAtStartPar
\sphinxstyleliteralstrong{\sphinxupquote{x}} \textendash{} argument 

\sphinxlineitem{Returns}
\sphinxAtStartPar
x*x 

\end{description}\end{quote}

\end{fulllineitems}

\index{MathSpecial::cube (C++ function)@\spxentry{MathSpecial::cube}\spxextra{C++ function}}

\begin{fulllineitems}
\phantomsection\label{\detokenize{Developer_utils:_CPPv4N9LAMMPS_NS11MathSpecial4cubeERKd}}\phantomsection\label{\detokenize{Developer_utils:_CPPv3N9LAMMPS_NS11MathSpecial4cubeERKd}}\phantomsection\label{\detokenize{Developer_utils:_CPPv2N9LAMMPS_NS11MathSpecial4cubeERKd}}\phantomsection\label{\detokenize{Developer_utils:LAMMPS_NS::MathSpecial::cube__doubleCR}}
\pysigstartsignatures
\pysigstartmultiline
\pysiglinewithargsret{\phantomsection\label{\detokenize{Developer_utils:math__special_8h_1ac0baf9eac8e92c1431410002ab7e6619}}\DUrole{k}{static}\DUrole{w}{ }\DUrole{k}{inline}\DUrole{w}{ }\DUrole{kt}{double}\DUrole{w}{ }\sphinxcode{\sphinxupquote{\DUrole{n}{LAMMPS\_NS}\DUrole{p}{::}\DUrole{n}{MathSpecial}\DUrole{p}{::}}}\sphinxbfcode{\sphinxupquote{\DUrole{n}{cube}}}}{\DUrole{k}{const}\DUrole{w}{ }\DUrole{kt}{double}\DUrole{w}{ }\DUrole{p}{\&}\DUrole{n,sig-param}{x}}{}
\pysigstopmultiline
\pysigstopsignatures
\sphinxAtStartPar
Fast inline version of pow(x, 3.0)
\begin{quote}\begin{description}
\sphinxlineitem{Parameters}
\sphinxAtStartPar
\sphinxstyleliteralstrong{\sphinxupquote{x}} \textendash{} argument 

\sphinxlineitem{Returns}
\sphinxAtStartPar
x*x 

\end{description}\end{quote}

\end{fulllineitems}

\index{MathSpecial::powsign (C++ function)@\spxentry{MathSpecial::powsign}\spxextra{C++ function}}

\begin{fulllineitems}
\phantomsection\label{\detokenize{Developer_utils:_CPPv4N9LAMMPS_NS11MathSpecial7powsignEKi}}\phantomsection\label{\detokenize{Developer_utils:_CPPv3N9LAMMPS_NS11MathSpecial7powsignEKi}}\phantomsection\label{\detokenize{Developer_utils:_CPPv2N9LAMMPS_NS11MathSpecial7powsignEKi}}\phantomsection\label{\detokenize{Developer_utils:LAMMPS_NS::MathSpecial::powsign__iC}}
\pysigstartsignatures
\pysigstartmultiline
\pysiglinewithargsret{\phantomsection\label{\detokenize{Developer_utils:math__special_8h_1acb0afe1baf4875b3a3fcaba95c88b0a3}}\DUrole{k}{static}\DUrole{w}{ }\DUrole{k}{inline}\DUrole{w}{ }\DUrole{kt}{double}\DUrole{w}{ }\sphinxcode{\sphinxupquote{\DUrole{n}{LAMMPS\_NS}\DUrole{p}{::}\DUrole{n}{MathSpecial}\DUrole{p}{::}}}\sphinxbfcode{\sphinxupquote{\DUrole{n}{powsign}}}}{\DUrole{k}{const}\DUrole{w}{ }\DUrole{kt}{int}\DUrole{w}{ }\DUrole{n,sig-param}{n}}{}
\pysigstopmultiline
\pysigstopsignatures
\end{fulllineitems}

\index{MathSpecial::powint (C++ function)@\spxentry{MathSpecial::powint}\spxextra{C++ function}}

\begin{fulllineitems}
\phantomsection\label{\detokenize{Developer_utils:_CPPv4N9LAMMPS_NS11MathSpecial6powintERKdKi}}\phantomsection\label{\detokenize{Developer_utils:_CPPv3N9LAMMPS_NS11MathSpecial6powintERKdKi}}\phantomsection\label{\detokenize{Developer_utils:_CPPv2N9LAMMPS_NS11MathSpecial6powintERKdKi}}\phantomsection\label{\detokenize{Developer_utils:LAMMPS_NS::MathSpecial::powint__doubleCR.iC}}
\pysigstartsignatures
\pysigstartmultiline
\pysiglinewithargsret{\phantomsection\label{\detokenize{Developer_utils:math__special_8h_1a1a7436c36200b55ffbfff5159a747964}}\DUrole{k}{static}\DUrole{w}{ }\DUrole{k}{inline}\DUrole{w}{ }\DUrole{kt}{double}\DUrole{w}{ }\sphinxcode{\sphinxupquote{\DUrole{n}{LAMMPS\_NS}\DUrole{p}{::}\DUrole{n}{MathSpecial}\DUrole{p}{::}}}\sphinxbfcode{\sphinxupquote{\DUrole{n}{powint}}}}{\DUrole{k}{const}\DUrole{w}{ }\DUrole{kt}{double}\DUrole{w}{ }\DUrole{p}{\&}\DUrole{n,sig-param}{x}\sphinxparamcomma \DUrole{k}{const}\DUrole{w}{ }\DUrole{kt}{int}\DUrole{w}{ }\DUrole{n,sig-param}{n}}{}
\pysigstopmultiline
\pysigstopsignatures
\end{fulllineitems}

\index{MathSpecial::powsinxx (C++ function)@\spxentry{MathSpecial::powsinxx}\spxextra{C++ function}}

\begin{fulllineitems}
\phantomsection\label{\detokenize{Developer_utils:_CPPv4N9LAMMPS_NS11MathSpecial8powsinxxERKdi}}\phantomsection\label{\detokenize{Developer_utils:_CPPv3N9LAMMPS_NS11MathSpecial8powsinxxERKdi}}\phantomsection\label{\detokenize{Developer_utils:_CPPv2N9LAMMPS_NS11MathSpecial8powsinxxERKdi}}\phantomsection\label{\detokenize{Developer_utils:LAMMPS_NS::MathSpecial::powsinxx__doubleCR.i}}
\pysigstartsignatures
\pysigstartmultiline
\pysiglinewithargsret{\phantomsection\label{\detokenize{Developer_utils:math__special_8h_1a151acd0bdb5ac0eca20cf6de0a9c1a0c}}\DUrole{k}{static}\DUrole{w}{ }\DUrole{k}{inline}\DUrole{w}{ }\DUrole{kt}{double}\DUrole{w}{ }\sphinxcode{\sphinxupquote{\DUrole{n}{LAMMPS\_NS}\DUrole{p}{::}\DUrole{n}{MathSpecial}\DUrole{p}{::}}}\sphinxbfcode{\sphinxupquote{\DUrole{n}{powsinxx}}}}{\DUrole{k}{const}\DUrole{w}{ }\DUrole{kt}{double}\DUrole{w}{ }\DUrole{p}{\&}\DUrole{n,sig-param}{x}\sphinxparamcomma \DUrole{kt}{int}\DUrole{w}{ }\DUrole{n,sig-param}{n}}{}
\pysigstopmultiline
\pysigstopsignatures
\end{fulllineitems}



\bigskip\hrule\bigskip



\section{Tokenizer classes}
\label{\detokenize{Developer_utils:tokenizer-classes}}
\sphinxAtStartPar
The purpose of the tokenizer classes is to simplify the recurring task
of breaking lines of text down into words and/or numbers.
Traditionally, LAMMPS code would be using the \sphinxcode{\sphinxupquote{strtok()}} function from
the C library for that purpose, but that function has two significant
disadvantages: 1) it cannot be used concurrently from different LAMMPS
instances since it stores its status in a global variable and 2) it
modifies the string that it is processing.  These classes were
implemented to avoid both of these issues and also to reduce the amount
of code that needs to be written.

\sphinxAtStartPar
The basic procedure is to create an instance of the tokenizer class with
the string to be processed as an argument and then do a loop until all
available tokens are read.  The constructor has a default set of
separator characters, but that can be overridden. The default separators
are all “whitespace” characters, i.e. the space character, the tabulator
character, the carriage return character, the linefeed character, and
the form feed character.
\sphinxSetupCaptionForVerbatim{Tokenizer class example listing entries of the PATH environment variable}
\def\sphinxLiteralBlockLabel{\label{\detokenize{Developer_utils:id3}}}
\begin{sphinxVerbatim}[commandchars=\\\{\}]
\PYG{c+cp}{\PYGZsh{}}\PYG{c+cp}{include}\PYG{+w}{ }\PYG{c+cpf}{\PYGZdq{}tokenizer.h\PYGZdq{}}
\PYG{c+cp}{\PYGZsh{}}\PYG{c+cp}{include}\PYG{+w}{ }\PYG{c+cpf}{\PYGZlt{}cstdlib\PYGZgt{}}
\PYG{c+cp}{\PYGZsh{}}\PYG{c+cp}{include}\PYG{+w}{ }\PYG{c+cpf}{\PYGZlt{}string\PYGZgt{}}
\PYG{c+cp}{\PYGZsh{}}\PYG{c+cp}{include}\PYG{+w}{ }\PYG{c+cpf}{\PYGZlt{}iostream\PYGZgt{}}

\PYG{k}{using}\PYG{+w}{ }\PYG{k}{namespace}\PYG{+w}{ }\PYG{n+nn}{LAMMPS\PYGZus{}NS}\PYG{p}{;}

\PYG{k+kt}{int}\PYG{+w}{ }\PYG{n+nf}{main}\PYG{p}{(}\PYG{k+kt}{int}\PYG{p}{,}\PYG{+w}{ }\PYG{k+kt}{char}\PYG{+w}{ }\PYG{o}{*}\PYG{o}{*}\PYG{p}{)}
\PYG{p}{\PYGZob{}}
\PYG{+w}{    }\PYG{k}{const}\PYG{+w}{ }\PYG{k+kt}{char}\PYG{+w}{ }\PYG{o}{*}\PYG{n}{path}\PYG{+w}{ }\PYG{o}{=}\PYG{+w}{ }\PYG{n}{getenv}\PYG{p}{(}\PYG{l+s}{\PYGZdq{}}\PYG{l+s}{PATH}\PYG{l+s}{\PYGZdq{}}\PYG{p}{)}\PYG{p}{;}

\PYG{+w}{    }\PYG{k}{if}\PYG{+w}{ }\PYG{p}{(}\PYG{n}{path}\PYG{+w}{ }\PYG{o}{!}\PYG{o}{=}\PYG{+w}{ }\PYG{n+nb}{nullptr}\PYG{p}{)}\PYG{+w}{ }\PYG{p}{\PYGZob{}}
\PYG{+w}{        }\PYG{n}{Tokenizer}\PYG{+w}{ }\PYG{n}{p}\PYG{p}{(}\PYG{n}{path}\PYG{p}{,}\PYG{l+s}{\PYGZdq{}}\PYG{l+s}{:}\PYG{l+s}{\PYGZdq{}}\PYG{p}{)}\PYG{p}{;}
\PYG{+w}{        }\PYG{k}{while}\PYG{+w}{ }\PYG{p}{(}\PYG{n}{p}\PYG{p}{.}\PYG{n}{has\PYGZus{}next}\PYG{p}{(}\PYG{p}{)}\PYG{p}{)}
\PYG{+w}{            }\PYG{n}{std}\PYG{o}{:}\PYG{o}{:}\PYG{n}{cout}\PYG{+w}{ }\PYG{o}{\PYGZlt{}}\PYG{o}{\PYGZlt{}}\PYG{+w}{ }\PYG{l+s}{\PYGZdq{}}\PYG{l+s}{Entry: }\PYG{l+s}{\PYGZdq{}}\PYG{+w}{ }\PYG{o}{\PYGZlt{}}\PYG{o}{\PYGZlt{}}\PYG{+w}{ }\PYG{n}{p}\PYG{p}{.}\PYG{n}{next}\PYG{p}{(}\PYG{p}{)}\PYG{+w}{ }\PYG{o}{\PYGZlt{}}\PYG{o}{\PYGZlt{}}\PYG{+w}{ }\PYG{l+s}{\PYGZdq{}}\PYG{l+s+se}{\PYGZbs{}n}\PYG{l+s}{\PYGZdq{}}\PYG{p}{;}
\PYG{+w}{    }\PYG{p}{\PYGZcb{}}
\PYG{+w}{    }\PYG{k}{return}\PYG{+w}{ }\PYG{l+m+mi}{0}\PYG{p}{;}
\PYG{p}{\PYGZcb{}}
\end{sphinxVerbatim}

\sphinxAtStartPar
Most tokenizer operations cannot fail except for
{\hyperref[\detokenize{Developer_utils:_CPPv4N9LAMMPS_NS9Tokenizer4nextEv}]{\sphinxcrossref{\sphinxcode{\sphinxupquote{LAMMPS\_NS::Tokenizer::next()}}}}} (when used without first
checking with {\hyperref[\detokenize{Developer_utils:_CPPv4NK9LAMMPS_NS9Tokenizer8has_nextEv}]{\sphinxcrossref{\sphinxcode{\sphinxupquote{LAMMPS\_NS::Tokenizer::has\_next()}}}}}) and
{\hyperref[\detokenize{Developer_utils:_CPPv4N9LAMMPS_NS9Tokenizer4skipEi}]{\sphinxcrossref{\sphinxcode{\sphinxupquote{LAMMPS\_NS::Tokenizer::skip()}}}}}.  In case of failure, the class
will throw an exception, so you may need to wrap the code using the
tokenizer into a \sphinxcode{\sphinxupquote{try}} / \sphinxcode{\sphinxupquote{catch}} block to handle errors.  The
{\hyperref[\detokenize{Developer_utils:_CPPv4N9LAMMPS_NS14ValueTokenizerE}]{\sphinxcrossref{\sphinxcode{\sphinxupquote{LAMMPS\_NS::ValueTokenizer}}}}} class may also throw an exception
when a (type of) number is requested as next token that is not
compatible with the string representing the next word.
\sphinxSetupCaptionForVerbatim{ValueTokenizer class example with exception handling}
\def\sphinxLiteralBlockLabel{\label{\detokenize{Developer_utils:id4}}}
\begin{sphinxVerbatim}[commandchars=\\\{\}]
\PYG{c+cp}{\PYGZsh{}}\PYG{c+cp}{include}\PYG{+w}{ }\PYG{c+cpf}{\PYGZdq{}tokenizer.h\PYGZdq{}}
\PYG{c+cp}{\PYGZsh{}}\PYG{c+cp}{include}\PYG{+w}{ }\PYG{c+cpf}{\PYGZlt{}cstdlib\PYGZgt{}}
\PYG{c+cp}{\PYGZsh{}}\PYG{c+cp}{include}\PYG{+w}{ }\PYG{c+cpf}{\PYGZlt{}string\PYGZgt{}}
\PYG{c+cp}{\PYGZsh{}}\PYG{c+cp}{include}\PYG{+w}{ }\PYG{c+cpf}{\PYGZlt{}iostream\PYGZgt{}}

\PYG{k}{using}\PYG{+w}{ }\PYG{k}{namespace}\PYG{+w}{ }\PYG{n+nn}{LAMMPS\PYGZus{}NS}\PYG{p}{;}

\PYG{k+kt}{int}\PYG{+w}{ }\PYG{n+nf}{main}\PYG{p}{(}\PYG{k+kt}{int}\PYG{p}{,}\PYG{+w}{ }\PYG{k+kt}{char}\PYG{+w}{ }\PYG{o}{*}\PYG{o}{*}\PYG{p}{)}
\PYG{p}{\PYGZob{}}
\PYG{+w}{    }\PYG{k}{const}\PYG{+w}{ }\PYG{k+kt}{char}\PYG{+w}{ }\PYG{o}{*}\PYG{n}{text}\PYG{+w}{ }\PYG{o}{=}\PYG{+w}{ }\PYG{l+s}{\PYGZdq{}}\PYG{l+s}{1 2 3 4 5 20.0 21 twentytwo 2.3}\PYG{l+s}{\PYGZdq{}}\PYG{p}{;}
\PYG{+w}{    }\PYG{k+kt}{double}\PYG{+w}{ }\PYG{n}{num1}\PYG{p}{(}\PYG{l+m+mi}{0}\PYG{p}{)}\PYG{p}{,}\PYG{n}{num2}\PYG{p}{(}\PYG{l+m+mi}{0}\PYG{p}{)}\PYG{p}{,}\PYG{n}{num3}\PYG{p}{(}\PYG{l+m+mi}{0}\PYG{p}{)}\PYG{p}{,}\PYG{n}{num4}\PYG{p}{(}\PYG{l+m+mi}{0}\PYG{p}{)}\PYG{p}{;}

\PYG{+w}{    }\PYG{n}{ValueTokenizer}\PYG{+w}{ }\PYG{n}{t}\PYG{p}{(}\PYG{n}{text}\PYG{p}{)}\PYG{p}{;}
\PYG{+w}{    }\PYG{c+c1}{// read 4 doubles after skipping over 5 numbers}
\PYG{+w}{    }\PYG{k}{try}\PYG{+w}{ }\PYG{p}{\PYGZob{}}
\PYG{+w}{        }\PYG{n}{t}\PYG{p}{.}\PYG{n}{skip}\PYG{p}{(}\PYG{l+m+mi}{5}\PYG{p}{)}\PYG{p}{;}
\PYG{+w}{        }\PYG{n}{num1}\PYG{+w}{ }\PYG{o}{=}\PYG{+w}{ }\PYG{n}{t}\PYG{p}{.}\PYG{n}{next\PYGZus{}double}\PYG{p}{(}\PYG{p}{)}\PYG{p}{;}
\PYG{+w}{        }\PYG{n}{num2}\PYG{+w}{ }\PYG{o}{=}\PYG{+w}{ }\PYG{n}{t}\PYG{p}{.}\PYG{n}{next\PYGZus{}double}\PYG{p}{(}\PYG{p}{)}\PYG{p}{;}
\PYG{+w}{        }\PYG{n}{num3}\PYG{+w}{ }\PYG{o}{=}\PYG{+w}{ }\PYG{n}{t}\PYG{p}{.}\PYG{n}{next\PYGZus{}double}\PYG{p}{(}\PYG{p}{)}\PYG{p}{;}
\PYG{+w}{        }\PYG{n}{num4}\PYG{+w}{ }\PYG{o}{=}\PYG{+w}{ }\PYG{n}{t}\PYG{p}{.}\PYG{n}{next\PYGZus{}double}\PYG{p}{(}\PYG{p}{)}\PYG{p}{;}
\PYG{+w}{    }\PYG{p}{\PYGZcb{}}\PYG{+w}{ }\PYG{k}{catch}\PYG{+w}{ }\PYG{p}{(}\PYG{n}{TokenizerException}\PYG{+w}{ }\PYG{o}{\PYGZam{}}\PYG{n}{e}\PYG{p}{)}\PYG{+w}{ }\PYG{p}{\PYGZob{}}
\PYG{+w}{        }\PYG{n}{std}\PYG{o}{:}\PYG{o}{:}\PYG{n}{cout}\PYG{+w}{ }\PYG{o}{\PYGZlt{}}\PYG{o}{\PYGZlt{}}\PYG{+w}{ }\PYG{l+s}{\PYGZdq{}}\PYG{l+s}{Reading numbers failed: }\PYG{l+s}{\PYGZdq{}}\PYG{+w}{ }\PYG{o}{\PYGZlt{}}\PYG{o}{\PYGZlt{}}\PYG{+w}{ }\PYG{n}{e}\PYG{p}{.}\PYG{n}{what}\PYG{p}{(}\PYG{p}{)}\PYG{+w}{ }\PYG{o}{\PYGZlt{}}\PYG{o}{\PYGZlt{}}\PYG{+w}{ }\PYG{l+s}{\PYGZdq{}}\PYG{l+s+se}{\PYGZbs{}n}\PYG{l+s}{\PYGZdq{}}\PYG{p}{;}
\PYG{+w}{    }\PYG{p}{\PYGZcb{}}
\PYG{+w}{    }\PYG{n}{std}\PYG{o}{:}\PYG{o}{:}\PYG{n}{cout}\PYG{+w}{ }\PYG{o}{\PYGZlt{}}\PYG{o}{\PYGZlt{}}\PYG{+w}{ }\PYG{l+s}{\PYGZdq{}}\PYG{l+s}{Values: }\PYG{l+s}{\PYGZdq{}}\PYG{+w}{ }\PYG{o}{\PYGZlt{}}\PYG{o}{\PYGZlt{}}\PYG{+w}{ }\PYG{n}{num1}\PYG{+w}{ }\PYG{o}{\PYGZlt{}}\PYG{o}{\PYGZlt{}}\PYG{+w}{ }\PYG{l+s}{\PYGZdq{}}\PYG{l+s}{ }\PYG{l+s}{\PYGZdq{}}\PYG{+w}{ }\PYG{o}{\PYGZlt{}}\PYG{o}{\PYGZlt{}}\PYG{+w}{ }\PYG{n}{num2}\PYG{+w}{ }\PYG{o}{\PYGZlt{}}\PYG{o}{\PYGZlt{}}\PYG{+w}{ }\PYG{l+s}{\PYGZdq{}}\PYG{l+s}{ }\PYG{l+s}{\PYGZdq{}}\PYG{+w}{ }\PYG{o}{\PYGZlt{}}\PYG{o}{\PYGZlt{}}\PYG{+w}{ }\PYG{n}{num3}\PYG{+w}{ }\PYG{o}{\PYGZlt{}}\PYG{o}{\PYGZlt{}}\PYG{+w}{ }\PYG{l+s}{\PYGZdq{}}\PYG{l+s}{ }\PYG{l+s}{\PYGZdq{}}\PYG{+w}{ }\PYG{o}{\PYGZlt{}}\PYG{o}{\PYGZlt{}}\PYG{+w}{ }\PYG{n}{num4}\PYG{+w}{ }\PYG{o}{\PYGZlt{}}\PYG{o}{\PYGZlt{}}\PYG{+w}{ }\PYG{l+s}{\PYGZdq{}}\PYG{l+s+se}{\PYGZbs{}n}\PYG{l+s}{\PYGZdq{}}\PYG{p}{;}
\PYG{+w}{    }\PYG{k}{return}\PYG{+w}{ }\PYG{l+m+mi}{0}\PYG{p}{;}
\PYG{p}{\PYGZcb{}}
\end{sphinxVerbatim}

\sphinxAtStartPar
This code example should produce the following output:

\begin{sphinxVerbatim}[commandchars=\\\{\}]
Reading numbers failed: Not a valid floating\PYGZhy{}point number: \PYGZsq{}twentytwo\PYGZsq{}
Values: 20 21 0 0
\end{sphinxVerbatim}


\bigskip\hrule\bigskip

\index{Tokenizer (C++ class)@\spxentry{Tokenizer}\spxextra{C++ class}}

\begin{fulllineitems}
\phantomsection\label{\detokenize{Developer_utils:_CPPv4N9LAMMPS_NS9TokenizerE}}\phantomsection\label{\detokenize{Developer_utils:_CPPv3N9LAMMPS_NS9TokenizerE}}\phantomsection\label{\detokenize{Developer_utils:_CPPv2N9LAMMPS_NS9TokenizerE}}\phantomsection\label{\detokenize{Developer_utils:LAMMPS_NS::Tokenizer}}
\pysigstartsignatures
\pysigstartmultiline
\pysigline{\phantomsection\label{\detokenize{Developer_utils:classLAMMPS__NS_1_1Tokenizer}}\DUrole{k}{class}\DUrole{w}{ }\sphinxbfcode{\sphinxupquote{\DUrole{n}{Tokenizer}}}}
\pysigstopmultiline
\pysigstopsignatures
\begin{sphinxuseclass}{breathe-sectiondef}\subsubsection*{Public Functions}
\index{Tokenizer::Tokenizer (C++ function)@\spxentry{Tokenizer::Tokenizer}\spxextra{C++ function}}

\begin{fulllineitems}
\phantomsection\label{\detokenize{Developer_utils:_CPPv4N9LAMMPS_NS9Tokenizer9TokenizerENSt6stringENSt6stringE}}\phantomsection\label{\detokenize{Developer_utils:_CPPv3N9LAMMPS_NS9Tokenizer9TokenizerENSt6stringENSt6stringE}}\phantomsection\label{\detokenize{Developer_utils:_CPPv2N9LAMMPS_NS9Tokenizer9TokenizerENSt6stringENSt6stringE}}\phantomsection\label{\detokenize{Developer_utils:LAMMPS_NS::Tokenizer::Tokenizer__ss.ss}}
\pysigstartsignatures
\pysigstartmultiline
\pysiglinewithargsret{\phantomsection\label{\detokenize{Developer_utils:classLAMMPS__NS_1_1Tokenizer_1a1a572462a4d7e941633c746e21bb7586}}\sphinxbfcode{\sphinxupquote{\DUrole{n}{Tokenizer}}}}{\DUrole{n}{std}\DUrole{p}{::}\DUrole{n}{string}\DUrole{w}{ }\DUrole{n,sig-param}{str}\sphinxparamcomma \DUrole{n}{std}\DUrole{p}{::}\DUrole{n}{string}\DUrole{w}{ }\DUrole{n,sig-param}{separators}\DUrole{w}{ }\DUrole{p}{=}\DUrole{w}{ }\DUrole{n}{TOKENIZER\_DEFAULT\_SEPARATORS}}{}
\pysigstopmultiline
\pysigstopsignatures
\sphinxAtStartPar
Class for splitting text into words

\sphinxAtStartPar
This tokenizer will break down a string into sub\sphinxhyphen{}strings (i.e words) separated by the given separator characters. If the string contains certain known UTF\sphinxhyphen{}8 characters they will be replaced by their ASCII equivalents processing the string.

\sphinxAtStartPar
\begin{description}
\sphinxlineitem{\sphinxstyleemphasis{See also}}
\sphinxAtStartPar
{\hyperref[\detokenize{Developer_utils:_CPPv4N9LAMMPS_NS14ValueTokenizerE}]{\sphinxcrossref{\sphinxcode{\sphinxupquote{ValueTokenizer}}}}}, {\hyperref[\detokenize{Developer_utils:_CPPv4N9LAMMPS_NS5utils11split_wordsERKNSt6stringE}]{\sphinxcrossref{\sphinxcode{\sphinxupquote{utils::split\_words()}}}}}, {\hyperref[\detokenize{Developer_utils:_CPPv4N9LAMMPS_NS5utils10utf8_substERKNSt6stringE}]{\sphinxcrossref{\sphinxcode{\sphinxupquote{utils::utf8\_subst()}}}}}

\end{description}

\begin{quote}\begin{description}
\sphinxlineitem{Parameters}\begin{itemize}
\item {} 
\sphinxAtStartPar
\sphinxstyleliteralstrong{\sphinxupquote{str}} \textendash{} string to be processed 

\item {} 
\sphinxAtStartPar
\sphinxstyleliteralstrong{\sphinxupquote{\_separators}} \textendash{} string with separator characters (default: “ \textbackslash{}t\textbackslash{}r\textbackslash{}n\textbackslash{}f”) 

\end{itemize}

\end{description}\end{quote}

\end{fulllineitems}

\index{Tokenizer::reset (C++ function)@\spxentry{Tokenizer::reset}\spxextra{C++ function}}

\begin{fulllineitems}
\phantomsection\label{\detokenize{Developer_utils:_CPPv4N9LAMMPS_NS9Tokenizer5resetEv}}\phantomsection\label{\detokenize{Developer_utils:_CPPv3N9LAMMPS_NS9Tokenizer5resetEv}}\phantomsection\label{\detokenize{Developer_utils:_CPPv2N9LAMMPS_NS9Tokenizer5resetEv}}\phantomsection\label{\detokenize{Developer_utils:LAMMPS_NS::Tokenizer::reset}}
\pysigstartsignatures
\pysigstartmultiline
\pysiglinewithargsret{\phantomsection\label{\detokenize{Developer_utils:classLAMMPS__NS_1_1Tokenizer_1ad20897c5c8bd47f5d4005989bead0e55}}\DUrole{kt}{void}\DUrole{w}{ }\sphinxbfcode{\sphinxupquote{\DUrole{n}{reset}}}}{}{}
\pysigstopmultiline
\pysigstopsignatures
\sphinxAtStartPar
Re\sphinxhyphen{}position the tokenizer state to the first word, i.e. the first non\sphinxhyphen{}separator character 

\end{fulllineitems}

\index{Tokenizer::skip (C++ function)@\spxentry{Tokenizer::skip}\spxextra{C++ function}}

\begin{fulllineitems}
\phantomsection\label{\detokenize{Developer_utils:_CPPv4N9LAMMPS_NS9Tokenizer4skipEi}}\phantomsection\label{\detokenize{Developer_utils:_CPPv3N9LAMMPS_NS9Tokenizer4skipEi}}\phantomsection\label{\detokenize{Developer_utils:_CPPv2N9LAMMPS_NS9Tokenizer4skipEi}}\phantomsection\label{\detokenize{Developer_utils:LAMMPS_NS::Tokenizer::skip__i}}
\pysigstartsignatures
\pysigstartmultiline
\pysiglinewithargsret{\phantomsection\label{\detokenize{Developer_utils:classLAMMPS__NS_1_1Tokenizer_1abda0e6255008b1441f06a62f197da9a7}}\DUrole{kt}{void}\DUrole{w}{ }\sphinxbfcode{\sphinxupquote{\DUrole{n}{skip}}}}{\DUrole{kt}{int}\DUrole{w}{ }\DUrole{n,sig-param}{n}\DUrole{w}{ }\DUrole{p}{=}\DUrole{w}{ }\DUrole{m}{1}}{}
\pysigstopmultiline
\pysigstopsignatures
\sphinxAtStartPar
Skip over a given number of tokens
\begin{quote}\begin{description}
\sphinxlineitem{Parameters}
\sphinxAtStartPar
\sphinxstyleliteralstrong{\sphinxupquote{n}} \textendash{} number of tokens to skip over 

\end{description}\end{quote}

\end{fulllineitems}

\index{Tokenizer::has\_next (C++ function)@\spxentry{Tokenizer::has\_next}\spxextra{C++ function}}

\begin{fulllineitems}
\phantomsection\label{\detokenize{Developer_utils:_CPPv4NK9LAMMPS_NS9Tokenizer8has_nextEv}}\phantomsection\label{\detokenize{Developer_utils:_CPPv3NK9LAMMPS_NS9Tokenizer8has_nextEv}}\phantomsection\label{\detokenize{Developer_utils:_CPPv2NK9LAMMPS_NS9Tokenizer8has_nextEv}}\phantomsection\label{\detokenize{Developer_utils:LAMMPS_NS::Tokenizer::has_nextC}}
\pysigstartsignatures
\pysigstartmultiline
\pysiglinewithargsret{\phantomsection\label{\detokenize{Developer_utils:classLAMMPS__NS_1_1Tokenizer_1a5ac70fe9462200eec1fe01b48ea992f3}}\DUrole{kt}{bool}\DUrole{w}{ }\sphinxbfcode{\sphinxupquote{\DUrole{n}{has\_next}}}}{}{\DUrole{w}{ }\DUrole{k}{const}}
\pysigstopmultiline
\pysigstopsignatures
\sphinxAtStartPar
Indicate whether more tokens are available
\begin{quote}\begin{description}
\sphinxlineitem{Returns}
\sphinxAtStartPar
true if there are more tokens, false if not 

\end{description}\end{quote}

\end{fulllineitems}

\index{Tokenizer::contains (C++ function)@\spxentry{Tokenizer::contains}\spxextra{C++ function}}

\begin{fulllineitems}
\phantomsection\label{\detokenize{Developer_utils:_CPPv4NK9LAMMPS_NS9Tokenizer8containsERKNSt6stringE}}\phantomsection\label{\detokenize{Developer_utils:_CPPv3NK9LAMMPS_NS9Tokenizer8containsERKNSt6stringE}}\phantomsection\label{\detokenize{Developer_utils:_CPPv2NK9LAMMPS_NS9Tokenizer8containsERKNSt6stringE}}\phantomsection\label{\detokenize{Developer_utils:LAMMPS_NS::Tokenizer::contains__ssCRC}}
\pysigstartsignatures
\pysigstartmultiline
\pysiglinewithargsret{\phantomsection\label{\detokenize{Developer_utils:classLAMMPS__NS_1_1Tokenizer_1abf5acdb9f865006dd3d41f084403f74f}}\DUrole{kt}{bool}\DUrole{w}{ }\sphinxbfcode{\sphinxupquote{\DUrole{n}{contains}}}}{\DUrole{k}{const}\DUrole{w}{ }\DUrole{n}{std}\DUrole{p}{::}\DUrole{n}{string}\DUrole{w}{ }\DUrole{p}{\&}\DUrole{n,sig-param}{str}}{\DUrole{w}{ }\DUrole{k}{const}}
\pysigstopmultiline
\pysigstopsignatures
\sphinxAtStartPar
Search the text to be processed for a sub\sphinxhyphen{}string.

\sphinxAtStartPar
This method does a generic sub\sphinxhyphen{}string match.
\begin{quote}\begin{description}
\sphinxlineitem{Parameters}
\sphinxAtStartPar
\sphinxstyleliteralstrong{\sphinxupquote{str}} \textendash{} string to be searched for 

\sphinxlineitem{Returns}
\sphinxAtStartPar
true if string was found, false if not 

\end{description}\end{quote}

\end{fulllineitems}

\index{Tokenizer::matches (C++ function)@\spxentry{Tokenizer::matches}\spxextra{C++ function}}

\begin{fulllineitems}
\phantomsection\label{\detokenize{Developer_utils:_CPPv4NK9LAMMPS_NS9Tokenizer7matchesERKNSt6stringE}}\phantomsection\label{\detokenize{Developer_utils:_CPPv3NK9LAMMPS_NS9Tokenizer7matchesERKNSt6stringE}}\phantomsection\label{\detokenize{Developer_utils:_CPPv2NK9LAMMPS_NS9Tokenizer7matchesERKNSt6stringE}}\phantomsection\label{\detokenize{Developer_utils:LAMMPS_NS::Tokenizer::matches__ssCRC}}
\pysigstartsignatures
\pysigstartmultiline
\pysiglinewithargsret{\phantomsection\label{\detokenize{Developer_utils:classLAMMPS__NS_1_1Tokenizer_1a07a6c962ca23635328ffde2dbac74533}}\DUrole{kt}{bool}\DUrole{w}{ }\sphinxbfcode{\sphinxupquote{\DUrole{n}{matches}}}}{\DUrole{k}{const}\DUrole{w}{ }\DUrole{n}{std}\DUrole{p}{::}\DUrole{n}{string}\DUrole{w}{ }\DUrole{p}{\&}\DUrole{n,sig-param}{str}}{\DUrole{w}{ }\DUrole{k}{const}}
\pysigstopmultiline
\pysigstopsignatures
\sphinxAtStartPar
Search the text to be processed for regular expression match.

\sphinxAtStartPar
This method matches the current string against a regular expression using
the {\hyperref[\detokenize{Developer_utils:_CPPv4N9LAMMPS_NS5utils8strmatchERKNSt6stringERKNSt6stringE}]{\sphinxcrossref{\sphinxcode{\sphinxupquote{utils::strmatch()}}}}} function.
\begin{quote}\begin{description}
\sphinxlineitem{Parameters}
\sphinxAtStartPar
\sphinxstyleliteralstrong{\sphinxupquote{str}} \textendash{} regular expression to be matched against 

\sphinxlineitem{Returns}
\sphinxAtStartPar
true if string was found, false if not 

\end{description}\end{quote}

\end{fulllineitems}

\index{Tokenizer::next (C++ function)@\spxentry{Tokenizer::next}\spxextra{C++ function}}

\begin{fulllineitems}
\phantomsection\label{\detokenize{Developer_utils:_CPPv4N9LAMMPS_NS9Tokenizer4nextEv}}\phantomsection\label{\detokenize{Developer_utils:_CPPv3N9LAMMPS_NS9Tokenizer4nextEv}}\phantomsection\label{\detokenize{Developer_utils:_CPPv2N9LAMMPS_NS9Tokenizer4nextEv}}\phantomsection\label{\detokenize{Developer_utils:LAMMPS_NS::Tokenizer::next}}
\pysigstartsignatures
\pysigstartmultiline
\pysiglinewithargsret{\phantomsection\label{\detokenize{Developer_utils:classLAMMPS__NS_1_1Tokenizer_1afba811f5984f17a43875432cb2845bc3}}\DUrole{n}{std}\DUrole{p}{::}\DUrole{n}{string}\DUrole{w}{ }\sphinxbfcode{\sphinxupquote{\DUrole{n}{next}}}}{}{}
\pysigstopmultiline
\pysigstopsignatures
\sphinxAtStartPar
Retrieve next token.
\begin{quote}\begin{description}
\sphinxlineitem{Returns}
\sphinxAtStartPar
string with the next token 

\end{description}\end{quote}

\end{fulllineitems}

\index{Tokenizer::count (C++ function)@\spxentry{Tokenizer::count}\spxextra{C++ function}}

\begin{fulllineitems}
\phantomsection\label{\detokenize{Developer_utils:_CPPv4N9LAMMPS_NS9Tokenizer5countEv}}\phantomsection\label{\detokenize{Developer_utils:_CPPv3N9LAMMPS_NS9Tokenizer5countEv}}\phantomsection\label{\detokenize{Developer_utils:_CPPv2N9LAMMPS_NS9Tokenizer5countEv}}\phantomsection\label{\detokenize{Developer_utils:LAMMPS_NS::Tokenizer::count}}
\pysigstartsignatures
\pysigstartmultiline
\pysiglinewithargsret{\phantomsection\label{\detokenize{Developer_utils:classLAMMPS__NS_1_1Tokenizer_1ad873dc786f051a2c04e32d1541d3702e}}\DUrole{n}{size\_t}\DUrole{w}{ }\sphinxbfcode{\sphinxupquote{\DUrole{n}{count}}}}{}{}
\pysigstopmultiline
\pysigstopsignatures
\sphinxAtStartPar
Count number of tokens in text.
\begin{quote}\begin{description}
\sphinxlineitem{Returns}
\sphinxAtStartPar
number of counted tokens 

\end{description}\end{quote}

\end{fulllineitems}

\index{Tokenizer::as\_vector (C++ function)@\spxentry{Tokenizer::as\_vector}\spxextra{C++ function}}

\begin{fulllineitems}
\phantomsection\label{\detokenize{Developer_utils:_CPPv4N9LAMMPS_NS9Tokenizer9as_vectorEv}}\phantomsection\label{\detokenize{Developer_utils:_CPPv3N9LAMMPS_NS9Tokenizer9as_vectorEv}}\phantomsection\label{\detokenize{Developer_utils:_CPPv2N9LAMMPS_NS9Tokenizer9as_vectorEv}}\phantomsection\label{\detokenize{Developer_utils:LAMMPS_NS::Tokenizer::as_vector}}
\pysigstartsignatures
\pysigstartmultiline
\pysiglinewithargsret{\phantomsection\label{\detokenize{Developer_utils:classLAMMPS__NS_1_1Tokenizer_1abc12e1da3f00c1a120a02b136f0730dd}}\DUrole{n}{std}\DUrole{p}{::}\DUrole{n}{vector}\DUrole{p}{\textless{}}\DUrole{n}{std}\DUrole{p}{::}\DUrole{n}{string}\DUrole{p}{\textgreater{}}\DUrole{w}{ }\sphinxbfcode{\sphinxupquote{\DUrole{n}{as\_vector}}}}{}{}
\pysigstopmultiline
\pysigstopsignatures
\sphinxAtStartPar
Retrieve the entire text converted to an STL vector of tokens.
\begin{quote}\begin{description}
\sphinxlineitem{Returns}
\sphinxAtStartPar
The STL vector 

\end{description}\end{quote}

\end{fulllineitems}


\end{sphinxuseclass}
\end{fulllineitems}

\index{TokenizerException (C++ class)@\spxentry{TokenizerException}\spxextra{C++ class}}

\begin{fulllineitems}
\phantomsection\label{\detokenize{Developer_utils:_CPPv4N9LAMMPS_NS18TokenizerExceptionE}}\phantomsection\label{\detokenize{Developer_utils:_CPPv3N9LAMMPS_NS18TokenizerExceptionE}}\phantomsection\label{\detokenize{Developer_utils:_CPPv2N9LAMMPS_NS18TokenizerExceptionE}}\phantomsection\label{\detokenize{Developer_utils:LAMMPS_NS::TokenizerException}}
\pysigstartsignatures
\pysigstartmultiline
\pysigline{\phantomsection\label{\detokenize{Developer_utils:classLAMMPS__NS_1_1TokenizerException}}\DUrole{k}{class}\DUrole{w}{ }\sphinxbfcode{\sphinxupquote{\DUrole{n}{TokenizerException}}}\DUrole{w}{ }\DUrole{p}{:}\DUrole{w}{ }\DUrole{k}{public}\DUrole{w}{ }\DUrole{n}{exception}}
\pysigstopmultiline
\pysigstopsignatures
\sphinxAtStartPar
General Tokenizer exception class 

\sphinxAtStartPar
Subclassed by {\hyperref[\detokenize{Developer_utils:classLAMMPS__NS_1_1InvalidFloatException}]{\sphinxcrossref{\DUrole{std,std-ref}{InvalidFloatException}}}}, {\hyperref[\detokenize{Developer_utils:classLAMMPS__NS_1_1InvalidIntegerException}]{\sphinxcrossref{\DUrole{std,std-ref}{InvalidIntegerException}}}}

\begin{sphinxuseclass}{breathe-sectiondef}\subsubsection*{Public Functions}
\index{TokenizerException::TokenizerException (C++ function)@\spxentry{TokenizerException::TokenizerException}\spxextra{C++ function}}

\begin{fulllineitems}
\phantomsection\label{\detokenize{Developer_utils:_CPPv4N9LAMMPS_NS18TokenizerException18TokenizerExceptionERKNSt6stringERKNSt6stringE}}\phantomsection\label{\detokenize{Developer_utils:_CPPv3N9LAMMPS_NS18TokenizerException18TokenizerExceptionERKNSt6stringERKNSt6stringE}}\phantomsection\label{\detokenize{Developer_utils:_CPPv2N9LAMMPS_NS18TokenizerException18TokenizerExceptionERKNSt6stringERKNSt6stringE}}\phantomsection\label{\detokenize{Developer_utils:LAMMPS_NS::TokenizerException::TokenizerException__ssCR.ssCR}}
\pysigstartsignatures
\pysigstartmultiline
\pysiglinewithargsret{\phantomsection\label{\detokenize{Developer_utils:classLAMMPS__NS_1_1TokenizerException_1afcd938a3385e0b950738c2e24cdbc484}}\DUrole{k}{explicit}\DUrole{w}{ }\sphinxbfcode{\sphinxupquote{\DUrole{n}{TokenizerException}}}}{\DUrole{k}{const}\DUrole{w}{ }\DUrole{n}{std}\DUrole{p}{::}\DUrole{n}{string}\DUrole{w}{ }\DUrole{p}{\&}\DUrole{n,sig-param}{msg}\sphinxparamcomma \DUrole{k}{const}\DUrole{w}{ }\DUrole{n}{std}\DUrole{p}{::}\DUrole{n}{string}\DUrole{w}{ }\DUrole{p}{\&}\DUrole{n,sig-param}{token}}{}
\pysigstopmultiline
\pysigstopsignatures
\sphinxAtStartPar
Thrown during retrieving or skipping tokens
\begin{quote}\begin{description}
\sphinxlineitem{Parameters}\begin{itemize}
\item {} 
\sphinxAtStartPar
\sphinxstyleliteralstrong{\sphinxupquote{msg}} \textendash{} String with error message 

\item {} 
\sphinxAtStartPar
\sphinxstyleliteralstrong{\sphinxupquote{token}} \textendash{} String of the token/word that caused the error 

\end{itemize}

\end{description}\end{quote}

\end{fulllineitems}

\index{TokenizerException::what (C++ function)@\spxentry{TokenizerException::what}\spxextra{C++ function}}

\begin{fulllineitems}
\phantomsection\label{\detokenize{Developer_utils:_CPPv4NK9LAMMPS_NS18TokenizerException4whatEv}}\phantomsection\label{\detokenize{Developer_utils:_CPPv3NK9LAMMPS_NS18TokenizerException4whatEv}}\phantomsection\label{\detokenize{Developer_utils:_CPPv2NK9LAMMPS_NS18TokenizerException4whatEv}}\phantomsection\label{\detokenize{Developer_utils:LAMMPS_NS::TokenizerException::whatC}}
\pysigstartsignatures
\pysigstartmultiline
\pysiglinewithargsret{\phantomsection\label{\detokenize{Developer_utils:classLAMMPS__NS_1_1TokenizerException_1aac1c33909e01e1dfc85617cf2477d3eb}}\DUrole{k}{inline}\DUrole{w}{ }\DUrole{k}{const}\DUrole{w}{ }\DUrole{kt}{char}\DUrole{w}{ }\DUrole{p}{*}\sphinxbfcode{\sphinxupquote{\DUrole{n}{what}}}}{}{\DUrole{w}{ }\DUrole{k}{const}\DUrole{w}{ }\DUrole{k}{noexcept}\DUrole{w}{ }\DUrole{k}{override}}
\pysigstopmultiline
\pysigstopsignatures
\sphinxAtStartPar
Retrieve message describing the thrown exception 
\begin{quote}\begin{description}
\sphinxlineitem{Returns}
\sphinxAtStartPar
string with error message 

\end{description}\end{quote}

\end{fulllineitems}


\end{sphinxuseclass}
\end{fulllineitems}

\index{ValueTokenizer (C++ class)@\spxentry{ValueTokenizer}\spxextra{C++ class}}

\begin{fulllineitems}
\phantomsection\label{\detokenize{Developer_utils:_CPPv4N9LAMMPS_NS14ValueTokenizerE}}\phantomsection\label{\detokenize{Developer_utils:_CPPv3N9LAMMPS_NS14ValueTokenizerE}}\phantomsection\label{\detokenize{Developer_utils:_CPPv2N9LAMMPS_NS14ValueTokenizerE}}\phantomsection\label{\detokenize{Developer_utils:LAMMPS_NS::ValueTokenizer}}
\pysigstartsignatures
\pysigstartmultiline
\pysigline{\phantomsection\label{\detokenize{Developer_utils:classLAMMPS__NS_1_1ValueTokenizer}}\DUrole{k}{class}\DUrole{w}{ }\sphinxbfcode{\sphinxupquote{\DUrole{n}{ValueTokenizer}}}}
\pysigstopmultiline
\pysigstopsignatures
\begin{sphinxuseclass}{breathe-sectiondef}\subsubsection*{Public Functions}
\index{ValueTokenizer::ValueTokenizer (C++ function)@\spxentry{ValueTokenizer::ValueTokenizer}\spxextra{C++ function}}

\begin{fulllineitems}
\phantomsection\label{\detokenize{Developer_utils:_CPPv4N9LAMMPS_NS14ValueTokenizer14ValueTokenizerERKNSt6stringERKNSt6stringE}}\phantomsection\label{\detokenize{Developer_utils:_CPPv3N9LAMMPS_NS14ValueTokenizer14ValueTokenizerERKNSt6stringERKNSt6stringE}}\phantomsection\label{\detokenize{Developer_utils:_CPPv2N9LAMMPS_NS14ValueTokenizer14ValueTokenizerERKNSt6stringERKNSt6stringE}}\phantomsection\label{\detokenize{Developer_utils:LAMMPS_NS::ValueTokenizer::ValueTokenizer__ssCR.ssCR}}
\pysigstartsignatures
\pysigstartmultiline
\pysiglinewithargsret{\phantomsection\label{\detokenize{Developer_utils:classLAMMPS__NS_1_1ValueTokenizer_1a9fbf461f611028a7d95b7df0dc8dbec4}}\sphinxbfcode{\sphinxupquote{\DUrole{n}{ValueTokenizer}}}}{\DUrole{k}{const}\DUrole{w}{ }\DUrole{n}{std}\DUrole{p}{::}\DUrole{n}{string}\DUrole{w}{ }\DUrole{p}{\&}\DUrole{n,sig-param}{str}\sphinxparamcomma \DUrole{k}{const}\DUrole{w}{ }\DUrole{n}{std}\DUrole{p}{::}\DUrole{n}{string}\DUrole{w}{ }\DUrole{p}{\&}\DUrole{n,sig-param}{separators}\DUrole{w}{ }\DUrole{p}{=}\DUrole{w}{ }\DUrole{n}{TOKENIZER\_DEFAULT\_SEPARATORS}}{}
\pysigstopmultiline
\pysigstopsignatures
\sphinxAtStartPar
Class for reading text with numbers

\sphinxAtStartPar
\begin{description}
\sphinxlineitem{\sphinxstyleemphasis{See also}}
\sphinxAtStartPar
{\hyperref[\detokenize{Developer_utils:_CPPv4N9LAMMPS_NS9TokenizerE}]{\sphinxcrossref{\sphinxcode{\sphinxupquote{Tokenizer}}}}}

\end{description}


\sphinxAtStartPar


\begin{sphinxseealso}{See also:}

\sphinxAtStartPar
Tokenizer {\hyperref[\detokenize{Developer_utils:classLAMMPS__NS_1_1InvalidIntegerException}]{\sphinxcrossref{\DUrole{std,std-ref}{InvalidIntegerException}}}} {\hyperref[\detokenize{Developer_utils:classLAMMPS__NS_1_1InvalidFloatException}]{\sphinxcrossref{\DUrole{std,std-ref}{InvalidFloatException}}}}


\end{sphinxseealso}


\begin{quote}\begin{description}
\sphinxlineitem{Parameters}\begin{itemize}
\item {} 
\sphinxAtStartPar
\sphinxstyleliteralstrong{\sphinxupquote{str}} \textendash{} String to be processed 

\item {} 
\sphinxAtStartPar
\sphinxstyleliteralstrong{\sphinxupquote{separators}} \textendash{} String with separator characters (default: “ \textbackslash{}t\textbackslash{}r\textbackslash{}n\textbackslash{}f”)

\end{itemize}

\end{description}\end{quote}

\end{fulllineitems}

\index{ValueTokenizer::next\_string (C++ function)@\spxentry{ValueTokenizer::next\_string}\spxextra{C++ function}}

\begin{fulllineitems}
\phantomsection\label{\detokenize{Developer_utils:_CPPv4N9LAMMPS_NS14ValueTokenizer11next_stringEv}}\phantomsection\label{\detokenize{Developer_utils:_CPPv3N9LAMMPS_NS14ValueTokenizer11next_stringEv}}\phantomsection\label{\detokenize{Developer_utils:_CPPv2N9LAMMPS_NS14ValueTokenizer11next_stringEv}}\phantomsection\label{\detokenize{Developer_utils:LAMMPS_NS::ValueTokenizer::next_string}}
\pysigstartsignatures
\pysigstartmultiline
\pysiglinewithargsret{\phantomsection\label{\detokenize{Developer_utils:classLAMMPS__NS_1_1ValueTokenizer_1ad19df6455b8dfbfd2e5b1d93fc72d257}}\DUrole{n}{std}\DUrole{p}{::}\DUrole{n}{string}\DUrole{w}{ }\sphinxbfcode{\sphinxupquote{\DUrole{n}{next\_string}}}}{}{}
\pysigstopmultiline
\pysigstopsignatures
\sphinxAtStartPar
Retrieve next token
\begin{quote}\begin{description}
\sphinxlineitem{Returns}
\sphinxAtStartPar
string with next token 

\end{description}\end{quote}

\end{fulllineitems}

\index{ValueTokenizer::next\_tagint (C++ function)@\spxentry{ValueTokenizer::next\_tagint}\spxextra{C++ function}}

\begin{fulllineitems}
\phantomsection\label{\detokenize{Developer_utils:_CPPv4N9LAMMPS_NS14ValueTokenizer11next_tagintEv}}\phantomsection\label{\detokenize{Developer_utils:_CPPv3N9LAMMPS_NS14ValueTokenizer11next_tagintEv}}\phantomsection\label{\detokenize{Developer_utils:_CPPv2N9LAMMPS_NS14ValueTokenizer11next_tagintEv}}\phantomsection\label{\detokenize{Developer_utils:LAMMPS_NS::ValueTokenizer::next_tagint}}
\pysigstartsignatures
\pysigstartmultiline
\pysiglinewithargsret{\phantomsection\label{\detokenize{Developer_utils:classLAMMPS__NS_1_1ValueTokenizer_1a388a19376fd8647693fa10077e29a68f}}\DUrole{n}{tagint}\DUrole{w}{ }\sphinxbfcode{\sphinxupquote{\DUrole{n}{next\_tagint}}}}{}{}
\pysigstopmultiline
\pysigstopsignatures
\sphinxAtStartPar
Retrieve next token and convert to tagint
\begin{quote}\begin{description}
\sphinxlineitem{Returns}
\sphinxAtStartPar
value of next token 

\end{description}\end{quote}

\end{fulllineitems}

\index{ValueTokenizer::next\_bigint (C++ function)@\spxentry{ValueTokenizer::next\_bigint}\spxextra{C++ function}}

\begin{fulllineitems}
\phantomsection\label{\detokenize{Developer_utils:_CPPv4N9LAMMPS_NS14ValueTokenizer11next_bigintEv}}\phantomsection\label{\detokenize{Developer_utils:_CPPv3N9LAMMPS_NS14ValueTokenizer11next_bigintEv}}\phantomsection\label{\detokenize{Developer_utils:_CPPv2N9LAMMPS_NS14ValueTokenizer11next_bigintEv}}\phantomsection\label{\detokenize{Developer_utils:LAMMPS_NS::ValueTokenizer::next_bigint}}
\pysigstartsignatures
\pysigstartmultiline
\pysiglinewithargsret{\phantomsection\label{\detokenize{Developer_utils:classLAMMPS__NS_1_1ValueTokenizer_1a960e60b9cac96e1abfc9beb6a436e414}}\DUrole{n}{bigint}\DUrole{w}{ }\sphinxbfcode{\sphinxupquote{\DUrole{n}{next\_bigint}}}}{}{}
\pysigstopmultiline
\pysigstopsignatures
\sphinxAtStartPar
Retrieve next token and convert to bigint
\begin{quote}\begin{description}
\sphinxlineitem{Returns}
\sphinxAtStartPar
value of next token 

\end{description}\end{quote}

\end{fulllineitems}

\index{ValueTokenizer::next\_int (C++ function)@\spxentry{ValueTokenizer::next\_int}\spxextra{C++ function}}

\begin{fulllineitems}
\phantomsection\label{\detokenize{Developer_utils:_CPPv4N9LAMMPS_NS14ValueTokenizer8next_intEv}}\phantomsection\label{\detokenize{Developer_utils:_CPPv3N9LAMMPS_NS14ValueTokenizer8next_intEv}}\phantomsection\label{\detokenize{Developer_utils:_CPPv2N9LAMMPS_NS14ValueTokenizer8next_intEv}}\phantomsection\label{\detokenize{Developer_utils:LAMMPS_NS::ValueTokenizer::next_int}}
\pysigstartsignatures
\pysigstartmultiline
\pysiglinewithargsret{\phantomsection\label{\detokenize{Developer_utils:classLAMMPS__NS_1_1ValueTokenizer_1a659db78ef3cd8cdbfa35654fbea28929}}\DUrole{kt}{int}\DUrole{w}{ }\sphinxbfcode{\sphinxupquote{\DUrole{n}{next\_int}}}}{}{}
\pysigstopmultiline
\pysigstopsignatures
\sphinxAtStartPar
Retrieve next token and convert to int
\begin{quote}\begin{description}
\sphinxlineitem{Returns}
\sphinxAtStartPar
value of next token 

\end{description}\end{quote}

\end{fulllineitems}

\index{ValueTokenizer::next\_double (C++ function)@\spxentry{ValueTokenizer::next\_double}\spxextra{C++ function}}

\begin{fulllineitems}
\phantomsection\label{\detokenize{Developer_utils:_CPPv4N9LAMMPS_NS14ValueTokenizer11next_doubleEv}}\phantomsection\label{\detokenize{Developer_utils:_CPPv3N9LAMMPS_NS14ValueTokenizer11next_doubleEv}}\phantomsection\label{\detokenize{Developer_utils:_CPPv2N9LAMMPS_NS14ValueTokenizer11next_doubleEv}}\phantomsection\label{\detokenize{Developer_utils:LAMMPS_NS::ValueTokenizer::next_double}}
\pysigstartsignatures
\pysigstartmultiline
\pysiglinewithargsret{\phantomsection\label{\detokenize{Developer_utils:classLAMMPS__NS_1_1ValueTokenizer_1aa6b2325cb771c1d1aa11cc5bd3c3f2e9}}\DUrole{kt}{double}\DUrole{w}{ }\sphinxbfcode{\sphinxupquote{\DUrole{n}{next\_double}}}}{}{}
\pysigstopmultiline
\pysigstopsignatures
\sphinxAtStartPar
Retrieve next token and convert to double
\begin{quote}\begin{description}
\sphinxlineitem{Returns}
\sphinxAtStartPar
value of next token 

\end{description}\end{quote}

\end{fulllineitems}

\index{ValueTokenizer::has\_next (C++ function)@\spxentry{ValueTokenizer::has\_next}\spxextra{C++ function}}

\begin{fulllineitems}
\phantomsection\label{\detokenize{Developer_utils:_CPPv4NK9LAMMPS_NS14ValueTokenizer8has_nextEv}}\phantomsection\label{\detokenize{Developer_utils:_CPPv3NK9LAMMPS_NS14ValueTokenizer8has_nextEv}}\phantomsection\label{\detokenize{Developer_utils:_CPPv2NK9LAMMPS_NS14ValueTokenizer8has_nextEv}}\phantomsection\label{\detokenize{Developer_utils:LAMMPS_NS::ValueTokenizer::has_nextC}}
\pysigstartsignatures
\pysigstartmultiline
\pysiglinewithargsret{\phantomsection\label{\detokenize{Developer_utils:classLAMMPS__NS_1_1ValueTokenizer_1a5ac70fe9462200eec1fe01b48ea992f3}}\DUrole{kt}{bool}\DUrole{w}{ }\sphinxbfcode{\sphinxupquote{\DUrole{n}{has\_next}}}}{}{\DUrole{w}{ }\DUrole{k}{const}}
\pysigstopmultiline
\pysigstopsignatures
\sphinxAtStartPar
Indicate whether more tokens are available
\begin{quote}\begin{description}
\sphinxlineitem{Returns}
\sphinxAtStartPar
true if there are more tokens, false if not 

\end{description}\end{quote}

\end{fulllineitems}

\index{ValueTokenizer::contains (C++ function)@\spxentry{ValueTokenizer::contains}\spxextra{C++ function}}

\begin{fulllineitems}
\phantomsection\label{\detokenize{Developer_utils:_CPPv4NK9LAMMPS_NS14ValueTokenizer8containsERKNSt6stringE}}\phantomsection\label{\detokenize{Developer_utils:_CPPv3NK9LAMMPS_NS14ValueTokenizer8containsERKNSt6stringE}}\phantomsection\label{\detokenize{Developer_utils:_CPPv2NK9LAMMPS_NS14ValueTokenizer8containsERKNSt6stringE}}\phantomsection\label{\detokenize{Developer_utils:LAMMPS_NS::ValueTokenizer::contains__ssCRC}}
\pysigstartsignatures
\pysigstartmultiline
\pysiglinewithargsret{\phantomsection\label{\detokenize{Developer_utils:classLAMMPS__NS_1_1ValueTokenizer_1a48b6af4c43173e0a07e80d29f125ae62}}\DUrole{kt}{bool}\DUrole{w}{ }\sphinxbfcode{\sphinxupquote{\DUrole{n}{contains}}}}{\DUrole{k}{const}\DUrole{w}{ }\DUrole{n}{std}\DUrole{p}{::}\DUrole{n}{string}\DUrole{w}{ }\DUrole{p}{\&}\DUrole{n,sig-param}{value}}{\DUrole{w}{ }\DUrole{k}{const}}
\pysigstopmultiline
\pysigstopsignatures
\sphinxAtStartPar
Search the text to be processed for a sub\sphinxhyphen{}string.

\sphinxAtStartPar
This method does a generic sub\sphinxhyphen{}string match.
\begin{quote}\begin{description}
\sphinxlineitem{Parameters}
\sphinxAtStartPar
\sphinxstyleliteralstrong{\sphinxupquote{value}} \textendash{} string with value to be searched for 

\sphinxlineitem{Returns}
\sphinxAtStartPar
true if string was found, false if not 

\end{description}\end{quote}

\end{fulllineitems}

\index{ValueTokenizer::matches (C++ function)@\spxentry{ValueTokenizer::matches}\spxextra{C++ function}}

\begin{fulllineitems}
\phantomsection\label{\detokenize{Developer_utils:_CPPv4NK9LAMMPS_NS14ValueTokenizer7matchesERKNSt6stringE}}\phantomsection\label{\detokenize{Developer_utils:_CPPv3NK9LAMMPS_NS14ValueTokenizer7matchesERKNSt6stringE}}\phantomsection\label{\detokenize{Developer_utils:_CPPv2NK9LAMMPS_NS14ValueTokenizer7matchesERKNSt6stringE}}\phantomsection\label{\detokenize{Developer_utils:LAMMPS_NS::ValueTokenizer::matches__ssCRC}}
\pysigstartsignatures
\pysigstartmultiline
\pysiglinewithargsret{\phantomsection\label{\detokenize{Developer_utils:classLAMMPS__NS_1_1ValueTokenizer_1a07a6c962ca23635328ffde2dbac74533}}\DUrole{kt}{bool}\DUrole{w}{ }\sphinxbfcode{\sphinxupquote{\DUrole{n}{matches}}}}{\DUrole{k}{const}\DUrole{w}{ }\DUrole{n}{std}\DUrole{p}{::}\DUrole{n}{string}\DUrole{w}{ }\DUrole{p}{\&}\DUrole{n,sig-param}{str}}{\DUrole{w}{ }\DUrole{k}{const}}
\pysigstopmultiline
\pysigstopsignatures
\sphinxAtStartPar
Search the text to be processed for regular expression match.

\sphinxAtStartPar
This method matches the current string against a regular expression using
the {\hyperref[\detokenize{Developer_utils:_CPPv4N9LAMMPS_NS5utils8strmatchERKNSt6stringERKNSt6stringE}]{\sphinxcrossref{\sphinxcode{\sphinxupquote{utils::strmatch()}}}}} function.
\begin{quote}\begin{description}
\sphinxlineitem{Parameters}
\sphinxAtStartPar
\sphinxstyleliteralstrong{\sphinxupquote{str}} \textendash{} regular expression to be matched against 

\sphinxlineitem{Returns}
\sphinxAtStartPar
true if string was found, false if not 

\end{description}\end{quote}

\end{fulllineitems}

\index{ValueTokenizer::skip (C++ function)@\spxentry{ValueTokenizer::skip}\spxextra{C++ function}}

\begin{fulllineitems}
\phantomsection\label{\detokenize{Developer_utils:_CPPv4N9LAMMPS_NS14ValueTokenizer4skipEi}}\phantomsection\label{\detokenize{Developer_utils:_CPPv3N9LAMMPS_NS14ValueTokenizer4skipEi}}\phantomsection\label{\detokenize{Developer_utils:_CPPv2N9LAMMPS_NS14ValueTokenizer4skipEi}}\phantomsection\label{\detokenize{Developer_utils:LAMMPS_NS::ValueTokenizer::skip__i}}
\pysigstartsignatures
\pysigstartmultiline
\pysiglinewithargsret{\phantomsection\label{\detokenize{Developer_utils:classLAMMPS__NS_1_1ValueTokenizer_1a8031283fcb7534c047cd51957d3bfb80}}\DUrole{kt}{void}\DUrole{w}{ }\sphinxbfcode{\sphinxupquote{\DUrole{n}{skip}}}}{\DUrole{kt}{int}\DUrole{w}{ }\DUrole{n,sig-param}{ntokens}\DUrole{w}{ }\DUrole{p}{=}\DUrole{w}{ }\DUrole{m}{1}}{}
\pysigstopmultiline
\pysigstopsignatures
\sphinxAtStartPar
Skip over a given number of tokens
\begin{quote}\begin{description}
\sphinxlineitem{Parameters}
\sphinxAtStartPar
\sphinxstyleliteralstrong{\sphinxupquote{n}} \textendash{} number of tokens to skip over 

\end{description}\end{quote}

\end{fulllineitems}

\index{ValueTokenizer::count (C++ function)@\spxentry{ValueTokenizer::count}\spxextra{C++ function}}

\begin{fulllineitems}
\phantomsection\label{\detokenize{Developer_utils:_CPPv4N9LAMMPS_NS14ValueTokenizer5countEv}}\phantomsection\label{\detokenize{Developer_utils:_CPPv3N9LAMMPS_NS14ValueTokenizer5countEv}}\phantomsection\label{\detokenize{Developer_utils:_CPPv2N9LAMMPS_NS14ValueTokenizer5countEv}}\phantomsection\label{\detokenize{Developer_utils:LAMMPS_NS::ValueTokenizer::count}}
\pysigstartsignatures
\pysigstartmultiline
\pysiglinewithargsret{\phantomsection\label{\detokenize{Developer_utils:classLAMMPS__NS_1_1ValueTokenizer_1ad873dc786f051a2c04e32d1541d3702e}}\DUrole{n}{size\_t}\DUrole{w}{ }\sphinxbfcode{\sphinxupquote{\DUrole{n}{count}}}}{}{}
\pysigstopmultiline
\pysigstopsignatures
\sphinxAtStartPar
Count number of tokens in text.
\begin{quote}\begin{description}
\sphinxlineitem{Returns}
\sphinxAtStartPar
number of counted tokens 

\end{description}\end{quote}

\end{fulllineitems}


\end{sphinxuseclass}
\end{fulllineitems}

\index{InvalidIntegerException (C++ class)@\spxentry{InvalidIntegerException}\spxextra{C++ class}}

\begin{fulllineitems}
\phantomsection\label{\detokenize{Developer_utils:_CPPv4N9LAMMPS_NS23InvalidIntegerExceptionE}}\phantomsection\label{\detokenize{Developer_utils:_CPPv3N9LAMMPS_NS23InvalidIntegerExceptionE}}\phantomsection\label{\detokenize{Developer_utils:_CPPv2N9LAMMPS_NS23InvalidIntegerExceptionE}}\phantomsection\label{\detokenize{Developer_utils:LAMMPS_NS::InvalidIntegerException}}
\pysigstartsignatures
\pysigstartmultiline
\pysigline{\phantomsection\label{\detokenize{Developer_utils:classLAMMPS__NS_1_1InvalidIntegerException}}\DUrole{k}{class}\DUrole{w}{ }\sphinxbfcode{\sphinxupquote{\DUrole{n}{InvalidIntegerException}}}\DUrole{w}{ }\DUrole{p}{:}\DUrole{w}{ }\DUrole{k}{public}\DUrole{w}{ }{\hyperref[\detokenize{Developer_utils:_CPPv4N9LAMMPS_NS18TokenizerExceptionE}]{\sphinxcrossref{\DUrole{n}{TokenizerException}}}}}
\pysigstopmultiline
\pysigstopsignatures
\sphinxAtStartPar
Exception thrown by ValueTokenizer when trying to convert an invalid integer string 

\begin{sphinxuseclass}{breathe-sectiondef}\subsubsection*{Public Functions}
\index{InvalidIntegerException::InvalidIntegerException (C++ function)@\spxentry{InvalidIntegerException::InvalidIntegerException}\spxextra{C++ function}}

\begin{fulllineitems}
\phantomsection\label{\detokenize{Developer_utils:_CPPv4N9LAMMPS_NS23InvalidIntegerException23InvalidIntegerExceptionERKNSt6stringE}}\phantomsection\label{\detokenize{Developer_utils:_CPPv3N9LAMMPS_NS23InvalidIntegerException23InvalidIntegerExceptionERKNSt6stringE}}\phantomsection\label{\detokenize{Developer_utils:_CPPv2N9LAMMPS_NS23InvalidIntegerException23InvalidIntegerExceptionERKNSt6stringE}}\phantomsection\label{\detokenize{Developer_utils:LAMMPS_NS::InvalidIntegerException::InvalidIntegerException__ssCR}}
\pysigstartsignatures
\pysigstartmultiline
\pysiglinewithargsret{\phantomsection\label{\detokenize{Developer_utils:classLAMMPS__NS_1_1InvalidIntegerException_1abbff68ef6488f500630d129910eb9136}}\DUrole{k}{inline}\DUrole{w}{ }\DUrole{k}{explicit}\DUrole{w}{ }\sphinxbfcode{\sphinxupquote{\DUrole{n}{InvalidIntegerException}}}}{\DUrole{k}{const}\DUrole{w}{ }\DUrole{n}{std}\DUrole{p}{::}\DUrole{n}{string}\DUrole{w}{ }\DUrole{p}{\&}\DUrole{n,sig-param}{token}}{}
\pysigstopmultiline
\pysigstopsignatures
\sphinxAtStartPar
Thrown during converting string to integer number
\begin{quote}\begin{description}
\sphinxlineitem{Parameters}
\sphinxAtStartPar
\sphinxstyleliteralstrong{\sphinxupquote{token}} \textendash{} String of the token/word that caused the error 

\end{description}\end{quote}

\end{fulllineitems}


\end{sphinxuseclass}
\end{fulllineitems}

\index{InvalidFloatException (C++ class)@\spxentry{InvalidFloatException}\spxextra{C++ class}}

\begin{fulllineitems}
\phantomsection\label{\detokenize{Developer_utils:_CPPv4N9LAMMPS_NS21InvalidFloatExceptionE}}\phantomsection\label{\detokenize{Developer_utils:_CPPv3N9LAMMPS_NS21InvalidFloatExceptionE}}\phantomsection\label{\detokenize{Developer_utils:_CPPv2N9LAMMPS_NS21InvalidFloatExceptionE}}\phantomsection\label{\detokenize{Developer_utils:LAMMPS_NS::InvalidFloatException}}
\pysigstartsignatures
\pysigstartmultiline
\pysigline{\phantomsection\label{\detokenize{Developer_utils:classLAMMPS__NS_1_1InvalidFloatException}}\DUrole{k}{class}\DUrole{w}{ }\sphinxbfcode{\sphinxupquote{\DUrole{n}{InvalidFloatException}}}\DUrole{w}{ }\DUrole{p}{:}\DUrole{w}{ }\DUrole{k}{public}\DUrole{w}{ }{\hyperref[\detokenize{Developer_utils:_CPPv4N9LAMMPS_NS18TokenizerExceptionE}]{\sphinxcrossref{\DUrole{n}{TokenizerException}}}}}
\pysigstopmultiline
\pysigstopsignatures
\sphinxAtStartPar
Exception thrown by ValueTokenizer when trying to convert an floating point string 

\begin{sphinxuseclass}{breathe-sectiondef}\subsubsection*{Public Functions}
\index{InvalidFloatException::InvalidFloatException (C++ function)@\spxentry{InvalidFloatException::InvalidFloatException}\spxextra{C++ function}}

\begin{fulllineitems}
\phantomsection\label{\detokenize{Developer_utils:_CPPv4N9LAMMPS_NS21InvalidFloatException21InvalidFloatExceptionERKNSt6stringE}}\phantomsection\label{\detokenize{Developer_utils:_CPPv3N9LAMMPS_NS21InvalidFloatException21InvalidFloatExceptionERKNSt6stringE}}\phantomsection\label{\detokenize{Developer_utils:_CPPv2N9LAMMPS_NS21InvalidFloatException21InvalidFloatExceptionERKNSt6stringE}}\phantomsection\label{\detokenize{Developer_utils:LAMMPS_NS::InvalidFloatException::InvalidFloatException__ssCR}}
\pysigstartsignatures
\pysigstartmultiline
\pysiglinewithargsret{\phantomsection\label{\detokenize{Developer_utils:classLAMMPS__NS_1_1InvalidFloatException_1af72f7f4352493364c8f734f17f416ded}}\DUrole{k}{inline}\DUrole{w}{ }\DUrole{k}{explicit}\DUrole{w}{ }\sphinxbfcode{\sphinxupquote{\DUrole{n}{InvalidFloatException}}}}{\DUrole{k}{const}\DUrole{w}{ }\DUrole{n}{std}\DUrole{p}{::}\DUrole{n}{string}\DUrole{w}{ }\DUrole{p}{\&}\DUrole{n,sig-param}{token}}{}
\pysigstopmultiline
\pysigstopsignatures
\sphinxAtStartPar
Thrown during converting string to floating point number
\begin{quote}\begin{description}
\sphinxlineitem{Parameters}
\sphinxAtStartPar
\sphinxstyleliteralstrong{\sphinxupquote{token}} \textendash{} String of the token/word that caused the error 

\end{description}\end{quote}

\end{fulllineitems}


\end{sphinxuseclass}
\end{fulllineitems}



\bigskip\hrule\bigskip



\section{Argument parsing classes}
\label{\detokenize{Developer_utils:argument-parsing-classes}}
\sphinxAtStartPar
The purpose of argument parsing classes it to simplify and unify how
arguments of commands in LAMMPS are parsed and to make abstractions of
repetitive tasks.

\sphinxAtStartPar
The {\hyperref[\detokenize{Developer_utils:_CPPv4N9LAMMPS_NS7ArgInfoE}]{\sphinxcrossref{\sphinxcode{\sphinxupquote{LAMMPS\_NS::ArgInfo}}}}} class provides an abstraction
for parsing references to compute or fix styles, variables or custom
integer or double properties handled by {\hyperref[\detokenize{fix_property_atom::doc}]{\sphinxcrossref{\DUrole{doc}{fix property/atom}}}}.
These would start with a “c\_”, “f\_”, “v\_”, “d\_”, “d2\_”, “i\_”, or “i2\_”
followed by the ID or name of than instance and may be postfixed with
one or two array indices “{[}\textless{}number\textgreater{}{]}” with numbers \textgreater{} 0.

\sphinxAtStartPar
A typical code segment would look like this:
\sphinxSetupCaptionForVerbatim{Usage example for ArgInfo class}
\def\sphinxLiteralBlockLabel{\label{\detokenize{Developer_utils:id5}}}
\begin{sphinxVerbatim}[commandchars=\\\{\}]
\PYG{k+kt}{int}\PYG{+w}{ }\PYG{n}{nvalues}\PYG{+w}{ }\PYG{o}{=}\PYG{+w}{ }\PYG{l+m+mi}{0}\PYG{p}{;}
\PYG{k}{for}\PYG{+w}{ }\PYG{p}{(}\PYG{n}{iarg}\PYG{+w}{ }\PYG{o}{=}\PYG{+w}{ }\PYG{l+m+mi}{0}\PYG{p}{;}\PYG{+w}{ }\PYG{n}{iarg}\PYG{+w}{ }\PYG{o}{\PYGZlt{}}\PYG{+w}{ }\PYG{n}{nargnew}\PYG{p}{;}\PYG{+w}{ }\PYG{n}{iarg}\PYG{o}{+}\PYG{o}{+}\PYG{p}{)}\PYG{+w}{ }\PYG{p}{\PYGZob{}}
\PYG{+w}{  }\PYG{n}{ArgInfo}\PYG{+w}{ }\PYG{n+nf}{argi}\PYG{p}{(}\PYG{n}{arg}\PYG{p}{[}\PYG{n}{iarg}\PYG{p}{]}\PYG{p}{)}\PYG{p}{;}

\PYG{+w}{  }\PYG{n}{which}\PYG{p}{[}\PYG{n}{nvalues}\PYG{p}{]}\PYG{+w}{ }\PYG{o}{=}\PYG{+w}{ }\PYG{n}{argi}\PYG{p}{.}\PYG{n}{get\PYGZus{}type}\PYG{p}{(}\PYG{p}{)}\PYG{p}{;}
\PYG{+w}{  }\PYG{n}{argindex}\PYG{p}{[}\PYG{n}{nvalues}\PYG{p}{]}\PYG{+w}{ }\PYG{o}{=}\PYG{+w}{ }\PYG{n}{argi}\PYG{p}{.}\PYG{n}{get\PYGZus{}index1}\PYG{p}{(}\PYG{p}{)}\PYG{p}{;}
\PYG{+w}{  }\PYG{n}{ids}\PYG{p}{[}\PYG{n}{nvalues}\PYG{p}{]}\PYG{+w}{ }\PYG{o}{=}\PYG{+w}{ }\PYG{n}{argi}\PYG{p}{.}\PYG{n}{copy\PYGZus{}name}\PYG{p}{(}\PYG{p}{)}\PYG{p}{;}

\PYG{+w}{  }\PYG{k}{if}\PYG{+w}{ }\PYG{p}{(}\PYG{p}{(}\PYG{n}{which}\PYG{p}{[}\PYG{n}{nvalues}\PYG{p}{]}\PYG{+w}{ }\PYG{o}{=}\PYG{o}{=}\PYG{+w}{ }\PYG{n}{ArgInfo}\PYG{o}{:}\PYG{o}{:}\PYG{n}{UNKNOWN}\PYG{p}{)}
\PYG{+w}{       }\PYG{o}{|}\PYG{o}{|}\PYG{+w}{ }\PYG{p}{(}\PYG{n}{which}\PYG{p}{[}\PYG{n}{nvalues}\PYG{p}{]}\PYG{+w}{ }\PYG{o}{=}\PYG{o}{=}\PYG{+w}{ }\PYG{n}{ArgInfo}\PYG{o}{:}\PYG{o}{:}\PYG{n}{NONE}\PYG{p}{)}
\PYG{+w}{       }\PYG{o}{|}\PYG{o}{|}\PYG{+w}{ }\PYG{p}{(}\PYG{n}{argi}\PYG{p}{.}\PYG{n}{get\PYGZus{}dim}\PYG{p}{(}\PYG{p}{)}\PYG{+w}{ }\PYG{o}{\PYGZgt{}}\PYG{+w}{ }\PYG{l+m+mi}{1}\PYG{p}{)}\PYG{p}{)}
\PYG{+w}{    }\PYG{n}{error}\PYG{o}{\PYGZhy{}}\PYG{o}{\PYGZgt{}}\PYG{n}{all}\PYG{p}{(}\PYG{n}{FLERR}\PYG{p}{,}\PYG{l+s}{\PYGZdq{}}\PYG{l+s}{Illegal compute XXX command}\PYG{l+s}{\PYGZdq{}}\PYG{p}{)}\PYG{p}{;}

\PYG{+w}{  }\PYG{n}{nvalues}\PYG{o}{+}\PYG{o}{+}\PYG{p}{;}
\PYG{p}{\PYGZcb{}}
\end{sphinxVerbatim}


\bigskip\hrule\bigskip

\index{ArgInfo (C++ class)@\spxentry{ArgInfo}\spxextra{C++ class}}

\begin{fulllineitems}
\phantomsection\label{\detokenize{Developer_utils:_CPPv4N9LAMMPS_NS7ArgInfoE}}\phantomsection\label{\detokenize{Developer_utils:_CPPv3N9LAMMPS_NS7ArgInfoE}}\phantomsection\label{\detokenize{Developer_utils:_CPPv2N9LAMMPS_NS7ArgInfoE}}\phantomsection\label{\detokenize{Developer_utils:LAMMPS_NS::ArgInfo}}
\pysigstartsignatures
\pysigstartmultiline
\pysigline{\phantomsection\label{\detokenize{Developer_utils:classLAMMPS__NS_1_1ArgInfo}}\DUrole{k}{class}\DUrole{w}{ }\sphinxbfcode{\sphinxupquote{\DUrole{n}{ArgInfo}}}}
\pysigstopmultiline
\pysigstopsignatures
\begin{sphinxuseclass}{breathe-sectiondef}\subsubsection*{Public Types}
\index{ArgInfo::ArgTypes (C++ enum)@\spxentry{ArgInfo::ArgTypes}\spxextra{C++ enum}}

\begin{fulllineitems}
\phantomsection\label{\detokenize{Developer_utils:_CPPv4N9LAMMPS_NS7ArgInfo8ArgTypesE}}\phantomsection\label{\detokenize{Developer_utils:_CPPv3N9LAMMPS_NS7ArgInfo8ArgTypesE}}\phantomsection\label{\detokenize{Developer_utils:_CPPv2N9LAMMPS_NS7ArgInfo8ArgTypesE}}
\pysigstartsignatures
\pysigstartmultiline
\pysigline{\phantomsection\label{\detokenize{Developer_utils:classLAMMPS__NS_1_1ArgInfo_1a6df44aeb82f88a507de3bccbb74a6b13}}\DUrole{k}{enum}\DUrole{w}{ }\sphinxbfcode{\sphinxupquote{\DUrole{n}{ArgTypes}}}}
\pysigstopmultiline
\pysigstopsignatures
\sphinxAtStartPar
constants for argument types 

\sphinxAtStartPar
\sphinxstyleemphasis{Values:}
\index{ArgInfo::ArgTypes::ERROR (C++ enumerator)@\spxentry{ArgInfo::ArgTypes::ERROR}\spxextra{C++ enumerator}}

\begin{fulllineitems}
\phantomsection\label{\detokenize{Developer_utils:_CPPv4N9LAMMPS_NS7ArgInfo8ArgTypes5ERRORE}}\phantomsection\label{\detokenize{Developer_utils:_CPPv3N9LAMMPS_NS7ArgInfo8ArgTypes5ERRORE}}\phantomsection\label{\detokenize{Developer_utils:_CPPv2N9LAMMPS_NS7ArgInfo8ArgTypes5ERRORE}}
\pysigstartsignatures
\pysigstartmultiline
\pysigline{\phantomsection\label{\detokenize{Developer_utils:classLAMMPS__NS_1_1ArgInfo_1a6df44aeb82f88a507de3bccbb74a6b13a2fd6f336d08340583bd620a7f5694c90}}\DUrole{k}{enumerator}\DUrole{w}{ }\sphinxbfcode{\sphinxupquote{\DUrole{n}{ERROR}}}}
\pysigstopmultiline
\pysigstopsignatures
\end{fulllineitems}

\index{ArgInfo::ArgTypes::UNKNOWN (C++ enumerator)@\spxentry{ArgInfo::ArgTypes::UNKNOWN}\spxextra{C++ enumerator}}

\begin{fulllineitems}
\phantomsection\label{\detokenize{Developer_utils:_CPPv4N9LAMMPS_NS7ArgInfo8ArgTypes7UNKNOWNE}}\phantomsection\label{\detokenize{Developer_utils:_CPPv3N9LAMMPS_NS7ArgInfo8ArgTypes7UNKNOWNE}}\phantomsection\label{\detokenize{Developer_utils:_CPPv2N9LAMMPS_NS7ArgInfo8ArgTypes7UNKNOWNE}}
\pysigstartsignatures
\pysigstartmultiline
\pysigline{\phantomsection\label{\detokenize{Developer_utils:classLAMMPS__NS_1_1ArgInfo_1a6df44aeb82f88a507de3bccbb74a6b13a6ce26a62afab55d7606ad4e92428b30c}}\DUrole{k}{enumerator}\DUrole{w}{ }\sphinxbfcode{\sphinxupquote{\DUrole{n}{UNKNOWN}}}}
\pysigstopmultiline
\pysigstopsignatures
\end{fulllineitems}

\index{ArgInfo::ArgTypes::NONE (C++ enumerator)@\spxentry{ArgInfo::ArgTypes::NONE}\spxextra{C++ enumerator}}

\begin{fulllineitems}
\phantomsection\label{\detokenize{Developer_utils:_CPPv4N9LAMMPS_NS7ArgInfo8ArgTypes4NONEE}}\phantomsection\label{\detokenize{Developer_utils:_CPPv3N9LAMMPS_NS7ArgInfo8ArgTypes4NONEE}}\phantomsection\label{\detokenize{Developer_utils:_CPPv2N9LAMMPS_NS7ArgInfo8ArgTypes4NONEE}}
\pysigstartsignatures
\pysigstartmultiline
\pysigline{\phantomsection\label{\detokenize{Developer_utils:classLAMMPS__NS_1_1ArgInfo_1a6df44aeb82f88a507de3bccbb74a6b13ac157bdf0b85a40d2619cbc8bc1ae5fe2}}\DUrole{k}{enumerator}\DUrole{w}{ }\sphinxbfcode{\sphinxupquote{\DUrole{n}{NONE}}}}
\pysigstopmultiline
\pysigstopsignatures
\end{fulllineitems}

\index{ArgInfo::ArgTypes::X (C++ enumerator)@\spxentry{ArgInfo::ArgTypes::X}\spxextra{C++ enumerator}}

\begin{fulllineitems}
\phantomsection\label{\detokenize{Developer_utils:_CPPv4N9LAMMPS_NS7ArgInfo8ArgTypes1XE}}\phantomsection\label{\detokenize{Developer_utils:_CPPv3N9LAMMPS_NS7ArgInfo8ArgTypes1XE}}\phantomsection\label{\detokenize{Developer_utils:_CPPv2N9LAMMPS_NS7ArgInfo8ArgTypes1XE}}
\pysigstartsignatures
\pysigstartmultiline
\pysigline{\phantomsection\label{\detokenize{Developer_utils:classLAMMPS__NS_1_1ArgInfo_1a6df44aeb82f88a507de3bccbb74a6b13a58833a3110c570fb05130d40c365d1e4}}\DUrole{k}{enumerator}\DUrole{w}{ }\sphinxbfcode{\sphinxupquote{\DUrole{n}{X}}}}
\pysigstopmultiline
\pysigstopsignatures
\end{fulllineitems}

\index{ArgInfo::ArgTypes::V (C++ enumerator)@\spxentry{ArgInfo::ArgTypes::V}\spxextra{C++ enumerator}}

\begin{fulllineitems}
\phantomsection\label{\detokenize{Developer_utils:_CPPv4N9LAMMPS_NS7ArgInfo8ArgTypes1VE}}\phantomsection\label{\detokenize{Developer_utils:_CPPv3N9LAMMPS_NS7ArgInfo8ArgTypes1VE}}\phantomsection\label{\detokenize{Developer_utils:_CPPv2N9LAMMPS_NS7ArgInfo8ArgTypes1VE}}
\pysigstartsignatures
\pysigstartmultiline
\pysigline{\phantomsection\label{\detokenize{Developer_utils:classLAMMPS__NS_1_1ArgInfo_1a6df44aeb82f88a507de3bccbb74a6b13a28f41f1144eee94834387e9a6a088bc1}}\DUrole{k}{enumerator}\DUrole{w}{ }\sphinxbfcode{\sphinxupquote{\DUrole{n}{V}}}}
\pysigstopmultiline
\pysigstopsignatures
\end{fulllineitems}

\index{ArgInfo::ArgTypes::F (C++ enumerator)@\spxentry{ArgInfo::ArgTypes::F}\spxextra{C++ enumerator}}

\begin{fulllineitems}
\phantomsection\label{\detokenize{Developer_utils:_CPPv4N9LAMMPS_NS7ArgInfo8ArgTypes1FE}}\phantomsection\label{\detokenize{Developer_utils:_CPPv3N9LAMMPS_NS7ArgInfo8ArgTypes1FE}}\phantomsection\label{\detokenize{Developer_utils:_CPPv2N9LAMMPS_NS7ArgInfo8ArgTypes1FE}}
\pysigstartsignatures
\pysigstartmultiline
\pysigline{\phantomsection\label{\detokenize{Developer_utils:classLAMMPS__NS_1_1ArgInfo_1a6df44aeb82f88a507de3bccbb74a6b13af382a63cc3d6491bf26b59e66f46826d}}\DUrole{k}{enumerator}\DUrole{w}{ }\sphinxbfcode{\sphinxupquote{\DUrole{n}{F}}}}
\pysigstopmultiline
\pysigstopsignatures
\end{fulllineitems}

\index{ArgInfo::ArgTypes::COMPUTE (C++ enumerator)@\spxentry{ArgInfo::ArgTypes::COMPUTE}\spxextra{C++ enumerator}}

\begin{fulllineitems}
\phantomsection\label{\detokenize{Developer_utils:_CPPv4N9LAMMPS_NS7ArgInfo8ArgTypes7COMPUTEE}}\phantomsection\label{\detokenize{Developer_utils:_CPPv3N9LAMMPS_NS7ArgInfo8ArgTypes7COMPUTEE}}\phantomsection\label{\detokenize{Developer_utils:_CPPv2N9LAMMPS_NS7ArgInfo8ArgTypes7COMPUTEE}}
\pysigstartsignatures
\pysigstartmultiline
\pysigline{\phantomsection\label{\detokenize{Developer_utils:classLAMMPS__NS_1_1ArgInfo_1a6df44aeb82f88a507de3bccbb74a6b13a391a0f3463a51c90bdf19aa7f35e0b46}}\DUrole{k}{enumerator}\DUrole{w}{ }\sphinxbfcode{\sphinxupquote{\DUrole{n}{COMPUTE}}}}
\pysigstopmultiline
\pysigstopsignatures
\end{fulllineitems}

\index{ArgInfo::ArgTypes::FIX (C++ enumerator)@\spxentry{ArgInfo::ArgTypes::FIX}\spxextra{C++ enumerator}}

\begin{fulllineitems}
\phantomsection\label{\detokenize{Developer_utils:_CPPv4N9LAMMPS_NS7ArgInfo8ArgTypes3FIXE}}\phantomsection\label{\detokenize{Developer_utils:_CPPv3N9LAMMPS_NS7ArgInfo8ArgTypes3FIXE}}\phantomsection\label{\detokenize{Developer_utils:_CPPv2N9LAMMPS_NS7ArgInfo8ArgTypes3FIXE}}
\pysigstartsignatures
\pysigstartmultiline
\pysigline{\phantomsection\label{\detokenize{Developer_utils:classLAMMPS__NS_1_1ArgInfo_1a6df44aeb82f88a507de3bccbb74a6b13a4c71825a48c0068b0761fd432c96d451}}\DUrole{k}{enumerator}\DUrole{w}{ }\sphinxbfcode{\sphinxupquote{\DUrole{n}{FIX}}}}
\pysigstopmultiline
\pysigstopsignatures
\end{fulllineitems}

\index{ArgInfo::ArgTypes::VARIABLE (C++ enumerator)@\spxentry{ArgInfo::ArgTypes::VARIABLE}\spxextra{C++ enumerator}}

\begin{fulllineitems}
\phantomsection\label{\detokenize{Developer_utils:_CPPv4N9LAMMPS_NS7ArgInfo8ArgTypes8VARIABLEE}}\phantomsection\label{\detokenize{Developer_utils:_CPPv3N9LAMMPS_NS7ArgInfo8ArgTypes8VARIABLEE}}\phantomsection\label{\detokenize{Developer_utils:_CPPv2N9LAMMPS_NS7ArgInfo8ArgTypes8VARIABLEE}}
\pysigstartsignatures
\pysigstartmultiline
\pysigline{\phantomsection\label{\detokenize{Developer_utils:classLAMMPS__NS_1_1ArgInfo_1a6df44aeb82f88a507de3bccbb74a6b13a39031ce5df6f91d3778590d6d644b9ea}}\DUrole{k}{enumerator}\DUrole{w}{ }\sphinxbfcode{\sphinxupquote{\DUrole{n}{VARIABLE}}}}
\pysigstopmultiline
\pysigstopsignatures
\end{fulllineitems}

\index{ArgInfo::ArgTypes::KEYWORD (C++ enumerator)@\spxentry{ArgInfo::ArgTypes::KEYWORD}\spxextra{C++ enumerator}}

\begin{fulllineitems}
\phantomsection\label{\detokenize{Developer_utils:_CPPv4N9LAMMPS_NS7ArgInfo8ArgTypes7KEYWORDE}}\phantomsection\label{\detokenize{Developer_utils:_CPPv3N9LAMMPS_NS7ArgInfo8ArgTypes7KEYWORDE}}\phantomsection\label{\detokenize{Developer_utils:_CPPv2N9LAMMPS_NS7ArgInfo8ArgTypes7KEYWORDE}}
\pysigstartsignatures
\pysigstartmultiline
\pysigline{\phantomsection\label{\detokenize{Developer_utils:classLAMMPS__NS_1_1ArgInfo_1a6df44aeb82f88a507de3bccbb74a6b13a129281444e94f5f509cba213d51a814d}}\DUrole{k}{enumerator}\DUrole{w}{ }\sphinxbfcode{\sphinxupquote{\DUrole{n}{KEYWORD}}}}
\pysigstopmultiline
\pysigstopsignatures
\end{fulllineitems}

\index{ArgInfo::ArgTypes::TYPE (C++ enumerator)@\spxentry{ArgInfo::ArgTypes::TYPE}\spxextra{C++ enumerator}}

\begin{fulllineitems}
\phantomsection\label{\detokenize{Developer_utils:_CPPv4N9LAMMPS_NS7ArgInfo8ArgTypes4TYPEE}}\phantomsection\label{\detokenize{Developer_utils:_CPPv3N9LAMMPS_NS7ArgInfo8ArgTypes4TYPEE}}\phantomsection\label{\detokenize{Developer_utils:_CPPv2N9LAMMPS_NS7ArgInfo8ArgTypes4TYPEE}}
\pysigstartsignatures
\pysigstartmultiline
\pysigline{\phantomsection\label{\detokenize{Developer_utils:classLAMMPS__NS_1_1ArgInfo_1a6df44aeb82f88a507de3bccbb74a6b13ab47ea8bb955afd0adc0ef98517dd6084}}\DUrole{k}{enumerator}\DUrole{w}{ }\sphinxbfcode{\sphinxupquote{\DUrole{n}{TYPE}}}}
\pysigstopmultiline
\pysigstopsignatures
\end{fulllineitems}

\index{ArgInfo::ArgTypes::MOLECULE (C++ enumerator)@\spxentry{ArgInfo::ArgTypes::MOLECULE}\spxextra{C++ enumerator}}

\begin{fulllineitems}
\phantomsection\label{\detokenize{Developer_utils:_CPPv4N9LAMMPS_NS7ArgInfo8ArgTypes8MOLECULEE}}\phantomsection\label{\detokenize{Developer_utils:_CPPv3N9LAMMPS_NS7ArgInfo8ArgTypes8MOLECULEE}}\phantomsection\label{\detokenize{Developer_utils:_CPPv2N9LAMMPS_NS7ArgInfo8ArgTypes8MOLECULEE}}
\pysigstartsignatures
\pysigstartmultiline
\pysigline{\phantomsection\label{\detokenize{Developer_utils:classLAMMPS__NS_1_1ArgInfo_1a6df44aeb82f88a507de3bccbb74a6b13a52a658fe0a308541e9fd6a5e1248d500}}\DUrole{k}{enumerator}\DUrole{w}{ }\sphinxbfcode{\sphinxupquote{\DUrole{n}{MOLECULE}}}}
\pysigstopmultiline
\pysigstopsignatures
\end{fulllineitems}

\index{ArgInfo::ArgTypes::DNAME (C++ enumerator)@\spxentry{ArgInfo::ArgTypes::DNAME}\spxextra{C++ enumerator}}

\begin{fulllineitems}
\phantomsection\label{\detokenize{Developer_utils:_CPPv4N9LAMMPS_NS7ArgInfo8ArgTypes5DNAMEE}}\phantomsection\label{\detokenize{Developer_utils:_CPPv3N9LAMMPS_NS7ArgInfo8ArgTypes5DNAMEE}}\phantomsection\label{\detokenize{Developer_utils:_CPPv2N9LAMMPS_NS7ArgInfo8ArgTypes5DNAMEE}}
\pysigstartsignatures
\pysigstartmultiline
\pysigline{\phantomsection\label{\detokenize{Developer_utils:classLAMMPS__NS_1_1ArgInfo_1a6df44aeb82f88a507de3bccbb74a6b13acc5ed32d1085edc116f9e5b127775f6d}}\DUrole{k}{enumerator}\DUrole{w}{ }\sphinxbfcode{\sphinxupquote{\DUrole{n}{DNAME}}}}
\pysigstopmultiline
\pysigstopsignatures
\end{fulllineitems}

\index{ArgInfo::ArgTypes::INAME (C++ enumerator)@\spxentry{ArgInfo::ArgTypes::INAME}\spxextra{C++ enumerator}}

\begin{fulllineitems}
\phantomsection\label{\detokenize{Developer_utils:_CPPv4N9LAMMPS_NS7ArgInfo8ArgTypes5INAMEE}}\phantomsection\label{\detokenize{Developer_utils:_CPPv3N9LAMMPS_NS7ArgInfo8ArgTypes5INAMEE}}\phantomsection\label{\detokenize{Developer_utils:_CPPv2N9LAMMPS_NS7ArgInfo8ArgTypes5INAMEE}}
\pysigstartsignatures
\pysigstartmultiline
\pysigline{\phantomsection\label{\detokenize{Developer_utils:classLAMMPS__NS_1_1ArgInfo_1a6df44aeb82f88a507de3bccbb74a6b13a42992693b11698580c1c11967cbca39f}}\DUrole{k}{enumerator}\DUrole{w}{ }\sphinxbfcode{\sphinxupquote{\DUrole{n}{INAME}}}}
\pysigstopmultiline
\pysigstopsignatures
\end{fulllineitems}

\index{ArgInfo::ArgTypes::DENSITY\_NUMBER (C++ enumerator)@\spxentry{ArgInfo::ArgTypes::DENSITY\_NUMBER}\spxextra{C++ enumerator}}

\begin{fulllineitems}
\phantomsection\label{\detokenize{Developer_utils:_CPPv4N9LAMMPS_NS7ArgInfo8ArgTypes14DENSITY_NUMBERE}}\phantomsection\label{\detokenize{Developer_utils:_CPPv3N9LAMMPS_NS7ArgInfo8ArgTypes14DENSITY_NUMBERE}}\phantomsection\label{\detokenize{Developer_utils:_CPPv2N9LAMMPS_NS7ArgInfo8ArgTypes14DENSITY_NUMBERE}}
\pysigstartsignatures
\pysigstartmultiline
\pysigline{\phantomsection\label{\detokenize{Developer_utils:classLAMMPS__NS_1_1ArgInfo_1a6df44aeb82f88a507de3bccbb74a6b13a36aba5a6b73317be76adbe7e40189e16}}\DUrole{k}{enumerator}\DUrole{w}{ }\sphinxbfcode{\sphinxupquote{\DUrole{n}{DENSITY\_NUMBER}}}}
\pysigstopmultiline
\pysigstopsignatures
\end{fulllineitems}

\index{ArgInfo::ArgTypes::DENSITY\_MASS (C++ enumerator)@\spxentry{ArgInfo::ArgTypes::DENSITY\_MASS}\spxextra{C++ enumerator}}

\begin{fulllineitems}
\phantomsection\label{\detokenize{Developer_utils:_CPPv4N9LAMMPS_NS7ArgInfo8ArgTypes12DENSITY_MASSE}}\phantomsection\label{\detokenize{Developer_utils:_CPPv3N9LAMMPS_NS7ArgInfo8ArgTypes12DENSITY_MASSE}}\phantomsection\label{\detokenize{Developer_utils:_CPPv2N9LAMMPS_NS7ArgInfo8ArgTypes12DENSITY_MASSE}}
\pysigstartsignatures
\pysigstartmultiline
\pysigline{\phantomsection\label{\detokenize{Developer_utils:classLAMMPS__NS_1_1ArgInfo_1a6df44aeb82f88a507de3bccbb74a6b13a4763598e36da157b838fe8f74e8063ba}}\DUrole{k}{enumerator}\DUrole{w}{ }\sphinxbfcode{\sphinxupquote{\DUrole{n}{DENSITY\_MASS}}}}
\pysigstopmultiline
\pysigstopsignatures
\end{fulllineitems}

\index{ArgInfo::ArgTypes::MASS (C++ enumerator)@\spxentry{ArgInfo::ArgTypes::MASS}\spxextra{C++ enumerator}}

\begin{fulllineitems}
\phantomsection\label{\detokenize{Developer_utils:_CPPv4N9LAMMPS_NS7ArgInfo8ArgTypes4MASSE}}\phantomsection\label{\detokenize{Developer_utils:_CPPv3N9LAMMPS_NS7ArgInfo8ArgTypes4MASSE}}\phantomsection\label{\detokenize{Developer_utils:_CPPv2N9LAMMPS_NS7ArgInfo8ArgTypes4MASSE}}
\pysigstartsignatures
\pysigstartmultiline
\pysigline{\phantomsection\label{\detokenize{Developer_utils:classLAMMPS__NS_1_1ArgInfo_1a6df44aeb82f88a507de3bccbb74a6b13a1c910486b66cf1ccce2ef81d4318227f}}\DUrole{k}{enumerator}\DUrole{w}{ }\sphinxbfcode{\sphinxupquote{\DUrole{n}{MASS}}}}
\pysigstopmultiline
\pysigstopsignatures
\end{fulllineitems}

\index{ArgInfo::ArgTypes::TEMPERATURE (C++ enumerator)@\spxentry{ArgInfo::ArgTypes::TEMPERATURE}\spxextra{C++ enumerator}}

\begin{fulllineitems}
\phantomsection\label{\detokenize{Developer_utils:_CPPv4N9LAMMPS_NS7ArgInfo8ArgTypes11TEMPERATUREE}}\phantomsection\label{\detokenize{Developer_utils:_CPPv3N9LAMMPS_NS7ArgInfo8ArgTypes11TEMPERATUREE}}\phantomsection\label{\detokenize{Developer_utils:_CPPv2N9LAMMPS_NS7ArgInfo8ArgTypes11TEMPERATUREE}}
\pysigstartsignatures
\pysigstartmultiline
\pysigline{\phantomsection\label{\detokenize{Developer_utils:classLAMMPS__NS_1_1ArgInfo_1a6df44aeb82f88a507de3bccbb74a6b13ac4ae6787ff1d8b2d1cf0ae9aa696e56c}}\DUrole{k}{enumerator}\DUrole{w}{ }\sphinxbfcode{\sphinxupquote{\DUrole{n}{TEMPERATURE}}}}
\pysigstopmultiline
\pysigstopsignatures
\end{fulllineitems}

\index{ArgInfo::ArgTypes::BIN1D (C++ enumerator)@\spxentry{ArgInfo::ArgTypes::BIN1D}\spxextra{C++ enumerator}}

\begin{fulllineitems}
\phantomsection\label{\detokenize{Developer_utils:_CPPv4N9LAMMPS_NS7ArgInfo8ArgTypes5BIN1DE}}\phantomsection\label{\detokenize{Developer_utils:_CPPv3N9LAMMPS_NS7ArgInfo8ArgTypes5BIN1DE}}\phantomsection\label{\detokenize{Developer_utils:_CPPv2N9LAMMPS_NS7ArgInfo8ArgTypes5BIN1DE}}
\pysigstartsignatures
\pysigstartmultiline
\pysigline{\phantomsection\label{\detokenize{Developer_utils:classLAMMPS__NS_1_1ArgInfo_1a6df44aeb82f88a507de3bccbb74a6b13af113047abcb2980258dbb091edc74919}}\DUrole{k}{enumerator}\DUrole{w}{ }\sphinxbfcode{\sphinxupquote{\DUrole{n}{BIN1D}}}}
\pysigstopmultiline
\pysigstopsignatures
\end{fulllineitems}

\index{ArgInfo::ArgTypes::BIN2D (C++ enumerator)@\spxentry{ArgInfo::ArgTypes::BIN2D}\spxextra{C++ enumerator}}

\begin{fulllineitems}
\phantomsection\label{\detokenize{Developer_utils:_CPPv4N9LAMMPS_NS7ArgInfo8ArgTypes5BIN2DE}}\phantomsection\label{\detokenize{Developer_utils:_CPPv3N9LAMMPS_NS7ArgInfo8ArgTypes5BIN2DE}}\phantomsection\label{\detokenize{Developer_utils:_CPPv2N9LAMMPS_NS7ArgInfo8ArgTypes5BIN2DE}}
\pysigstartsignatures
\pysigstartmultiline
\pysigline{\phantomsection\label{\detokenize{Developer_utils:classLAMMPS__NS_1_1ArgInfo_1a6df44aeb82f88a507de3bccbb74a6b13a172ec7ac0684d66a619471d3d6cb1520}}\DUrole{k}{enumerator}\DUrole{w}{ }\sphinxbfcode{\sphinxupquote{\DUrole{n}{BIN2D}}}}
\pysigstopmultiline
\pysigstopsignatures
\end{fulllineitems}

\index{ArgInfo::ArgTypes::BIN3D (C++ enumerator)@\spxentry{ArgInfo::ArgTypes::BIN3D}\spxextra{C++ enumerator}}

\begin{fulllineitems}
\phantomsection\label{\detokenize{Developer_utils:_CPPv4N9LAMMPS_NS7ArgInfo8ArgTypes5BIN3DE}}\phantomsection\label{\detokenize{Developer_utils:_CPPv3N9LAMMPS_NS7ArgInfo8ArgTypes5BIN3DE}}\phantomsection\label{\detokenize{Developer_utils:_CPPv2N9LAMMPS_NS7ArgInfo8ArgTypes5BIN3DE}}
\pysigstartsignatures
\pysigstartmultiline
\pysigline{\phantomsection\label{\detokenize{Developer_utils:classLAMMPS__NS_1_1ArgInfo_1a6df44aeb82f88a507de3bccbb74a6b13a32f02f6ff0b9d0b294a92bd63c5658b6}}\DUrole{k}{enumerator}\DUrole{w}{ }\sphinxbfcode{\sphinxupquote{\DUrole{n}{BIN3D}}}}
\pysigstopmultiline
\pysigstopsignatures
\end{fulllineitems}

\index{ArgInfo::ArgTypes::BINSPHERE (C++ enumerator)@\spxentry{ArgInfo::ArgTypes::BINSPHERE}\spxextra{C++ enumerator}}

\begin{fulllineitems}
\phantomsection\label{\detokenize{Developer_utils:_CPPv4N9LAMMPS_NS7ArgInfo8ArgTypes9BINSPHEREE}}\phantomsection\label{\detokenize{Developer_utils:_CPPv3N9LAMMPS_NS7ArgInfo8ArgTypes9BINSPHEREE}}\phantomsection\label{\detokenize{Developer_utils:_CPPv2N9LAMMPS_NS7ArgInfo8ArgTypes9BINSPHEREE}}
\pysigstartsignatures
\pysigstartmultiline
\pysigline{\phantomsection\label{\detokenize{Developer_utils:classLAMMPS__NS_1_1ArgInfo_1a6df44aeb82f88a507de3bccbb74a6b13aea2305105c66b4fd68a8acd3b62b548f}}\DUrole{k}{enumerator}\DUrole{w}{ }\sphinxbfcode{\sphinxupquote{\DUrole{n}{BINSPHERE}}}}
\pysigstopmultiline
\pysigstopsignatures
\end{fulllineitems}

\index{ArgInfo::ArgTypes::BINCYLINDER (C++ enumerator)@\spxentry{ArgInfo::ArgTypes::BINCYLINDER}\spxextra{C++ enumerator}}

\begin{fulllineitems}
\phantomsection\label{\detokenize{Developer_utils:_CPPv4N9LAMMPS_NS7ArgInfo8ArgTypes11BINCYLINDERE}}\phantomsection\label{\detokenize{Developer_utils:_CPPv3N9LAMMPS_NS7ArgInfo8ArgTypes11BINCYLINDERE}}\phantomsection\label{\detokenize{Developer_utils:_CPPv2N9LAMMPS_NS7ArgInfo8ArgTypes11BINCYLINDERE}}
\pysigstartsignatures
\pysigstartmultiline
\pysigline{\phantomsection\label{\detokenize{Developer_utils:classLAMMPS__NS_1_1ArgInfo_1a6df44aeb82f88a507de3bccbb74a6b13a751fb391ace04b5f6bd4bc9b8209bbc0}}\DUrole{k}{enumerator}\DUrole{w}{ }\sphinxbfcode{\sphinxupquote{\DUrole{n}{BINCYLINDER}}}}
\pysigstopmultiline
\pysigstopsignatures
\end{fulllineitems}


\end{fulllineitems}


\end{sphinxuseclass}
\begin{sphinxuseclass}{breathe-sectiondef}\subsubsection*{Public Functions}
\index{ArgInfo::ArgInfo (C++ function)@\spxentry{ArgInfo::ArgInfo}\spxextra{C++ function}}

\begin{fulllineitems}
\phantomsection\label{\detokenize{Developer_utils:_CPPv4N9LAMMPS_NS7ArgInfo7ArgInfoERKNSt6stringEi}}\phantomsection\label{\detokenize{Developer_utils:_CPPv3N9LAMMPS_NS7ArgInfo7ArgInfoERKNSt6stringEi}}\phantomsection\label{\detokenize{Developer_utils:_CPPv2N9LAMMPS_NS7ArgInfo7ArgInfoERKNSt6stringEi}}\phantomsection\label{\detokenize{Developer_utils:LAMMPS_NS::ArgInfo::ArgInfo__ssCR.i}}
\pysigstartsignatures
\pysigstartmultiline
\pysiglinewithargsret{\phantomsection\label{\detokenize{Developer_utils:classLAMMPS__NS_1_1ArgInfo_1a81a8a9e941f241bc7bd8df5048edd72e}}\sphinxbfcode{\sphinxupquote{\DUrole{n}{ArgInfo}}}}{\DUrole{k}{const}\DUrole{w}{ }\DUrole{n}{std}\DUrole{p}{::}\DUrole{n}{string}\DUrole{w}{ }\DUrole{p}{\&}\DUrole{n,sig-param}{arg}\sphinxparamcomma \DUrole{kt}{int}\DUrole{w}{ }\DUrole{n,sig-param}{allowed}\DUrole{w}{ }\DUrole{p}{=}\DUrole{w}{ }\DUrole{n}{COMPUTE}\DUrole{w}{ }\DUrole{o}{|}\DUrole{w}{ }\DUrole{n}{FIX}\DUrole{w}{ }\DUrole{o}{|}\DUrole{w}{ }\DUrole{n}{VARIABLE}}{}
\pysigstopmultiline
\pysigstopsignatures
\sphinxAtStartPar
Class for processing references to fixes, computes and variables

\sphinxAtStartPar
This class provides an abstraction for the repetitive task of parsing arguments that may contain references to fixes, computes, variables, or custom per\sphinxhyphen{}atom properties. It will identify the name and the index value in the first and second dimension, if present.
\begin{quote}\begin{description}
\sphinxlineitem{Parameters}\begin{itemize}
\item {} 
\sphinxAtStartPar
\sphinxstyleliteralstrong{\sphinxupquote{arg}} \textendash{} string with possible reference 

\item {} 
\sphinxAtStartPar
\sphinxstyleliteralstrong{\sphinxupquote{allowed}} \textendash{} integer with bitmap of allowed types of references 

\end{itemize}

\end{description}\end{quote}

\end{fulllineitems}

\index{ArgInfo::get\_type (C++ function)@\spxentry{ArgInfo::get\_type}\spxextra{C++ function}}

\begin{fulllineitems}
\phantomsection\label{\detokenize{Developer_utils:_CPPv4NK9LAMMPS_NS7ArgInfo8get_typeEv}}\phantomsection\label{\detokenize{Developer_utils:_CPPv3NK9LAMMPS_NS7ArgInfo8get_typeEv}}\phantomsection\label{\detokenize{Developer_utils:_CPPv2NK9LAMMPS_NS7ArgInfo8get_typeEv}}\phantomsection\label{\detokenize{Developer_utils:LAMMPS_NS::ArgInfo::get_typeC}}
\pysigstartsignatures
\pysigstartmultiline
\pysiglinewithargsret{\phantomsection\label{\detokenize{Developer_utils:classLAMMPS__NS_1_1ArgInfo_1a328fa3a7797b2464bfb1079d365d0edb}}\DUrole{k}{inline}\DUrole{w}{ }\DUrole{kt}{int}\DUrole{w}{ }\sphinxbfcode{\sphinxupquote{\DUrole{n}{get\_type}}}}{}{\DUrole{w}{ }\DUrole{k}{const}}
\pysigstopmultiline
\pysigstopsignatures
\sphinxAtStartPar
get type of reference

\sphinxAtStartPar
Return a type constant for the reference. This may be either COMPUTE, FIX, VARIABLE (if not restricted to a subset of those by the “allowed” argument of the constructor) or NONE, if it if not a recognized or allowed reference, or UNKNOWN, in case some error happened identifying or parsing the values of the indices
\begin{quote}\begin{description}
\sphinxlineitem{Returns}
\sphinxAtStartPar
integer with a constant from ArgTypes enumerator 

\end{description}\end{quote}

\end{fulllineitems}

\index{ArgInfo::get\_dim (C++ function)@\spxentry{ArgInfo::get\_dim}\spxextra{C++ function}}

\begin{fulllineitems}
\phantomsection\label{\detokenize{Developer_utils:_CPPv4NK9LAMMPS_NS7ArgInfo7get_dimEv}}\phantomsection\label{\detokenize{Developer_utils:_CPPv3NK9LAMMPS_NS7ArgInfo7get_dimEv}}\phantomsection\label{\detokenize{Developer_utils:_CPPv2NK9LAMMPS_NS7ArgInfo7get_dimEv}}\phantomsection\label{\detokenize{Developer_utils:LAMMPS_NS::ArgInfo::get_dimC}}
\pysigstartsignatures
\pysigstartmultiline
\pysiglinewithargsret{\phantomsection\label{\detokenize{Developer_utils:classLAMMPS__NS_1_1ArgInfo_1a812caf9be8bbe92a2fd8273dd2dda065}}\DUrole{k}{inline}\DUrole{w}{ }\DUrole{kt}{int}\DUrole{w}{ }\sphinxbfcode{\sphinxupquote{\DUrole{n}{get\_dim}}}}{}{\DUrole{w}{ }\DUrole{k}{const}}
\pysigstopmultiline
\pysigstopsignatures
\sphinxAtStartPar
get dimension of reference

\sphinxAtStartPar
This will return either 0, 1, 2 depending on whether the reference has no, one or two “{[}\{number\}{]}” postfixes.
\begin{quote}\begin{description}
\sphinxlineitem{Returns}
\sphinxAtStartPar
integer with the dimensionality of the reference 

\end{description}\end{quote}

\end{fulllineitems}

\index{ArgInfo::get\_index1 (C++ function)@\spxentry{ArgInfo::get\_index1}\spxextra{C++ function}}

\begin{fulllineitems}
\phantomsection\label{\detokenize{Developer_utils:_CPPv4NK9LAMMPS_NS7ArgInfo10get_index1Ev}}\phantomsection\label{\detokenize{Developer_utils:_CPPv3NK9LAMMPS_NS7ArgInfo10get_index1Ev}}\phantomsection\label{\detokenize{Developer_utils:_CPPv2NK9LAMMPS_NS7ArgInfo10get_index1Ev}}\phantomsection\label{\detokenize{Developer_utils:LAMMPS_NS::ArgInfo::get_index1C}}
\pysigstartsignatures
\pysigstartmultiline
\pysiglinewithargsret{\phantomsection\label{\detokenize{Developer_utils:classLAMMPS__NS_1_1ArgInfo_1a1ef3adbaad24c6a95b2cbc7475585ffd}}\DUrole{k}{inline}\DUrole{w}{ }\DUrole{kt}{int}\DUrole{w}{ }\sphinxbfcode{\sphinxupquote{\DUrole{n}{get\_index1}}}}{}{\DUrole{w}{ }\DUrole{k}{const}}
\pysigstopmultiline
\pysigstopsignatures
\sphinxAtStartPar
get index of first dimension

\sphinxAtStartPar
This will return the number in the first “{[}\{number\}{]}” postfix or 0 if there is no postfix.
\begin{quote}\begin{description}
\sphinxlineitem{Returns}
\sphinxAtStartPar
integer with index or the postfix or 0 

\end{description}\end{quote}

\end{fulllineitems}

\index{ArgInfo::get\_index2 (C++ function)@\spxentry{ArgInfo::get\_index2}\spxextra{C++ function}}

\begin{fulllineitems}
\phantomsection\label{\detokenize{Developer_utils:_CPPv4NK9LAMMPS_NS7ArgInfo10get_index2Ev}}\phantomsection\label{\detokenize{Developer_utils:_CPPv3NK9LAMMPS_NS7ArgInfo10get_index2Ev}}\phantomsection\label{\detokenize{Developer_utils:_CPPv2NK9LAMMPS_NS7ArgInfo10get_index2Ev}}\phantomsection\label{\detokenize{Developer_utils:LAMMPS_NS::ArgInfo::get_index2C}}
\pysigstartsignatures
\pysigstartmultiline
\pysiglinewithargsret{\phantomsection\label{\detokenize{Developer_utils:classLAMMPS__NS_1_1ArgInfo_1adcbe39c7eaf369281347891bd0818678}}\DUrole{k}{inline}\DUrole{w}{ }\DUrole{kt}{int}\DUrole{w}{ }\sphinxbfcode{\sphinxupquote{\DUrole{n}{get\_index2}}}}{}{\DUrole{w}{ }\DUrole{k}{const}}
\pysigstopmultiline
\pysigstopsignatures
\sphinxAtStartPar
get index of second dimension

\sphinxAtStartPar
This will return the number in the second “{[}\{number\}{]}” postfix or \sphinxhyphen{}1 if there is no second postfix.
\begin{quote}\begin{description}
\sphinxlineitem{Returns}
\sphinxAtStartPar
integer with index of the postfix or \sphinxhyphen{}1 

\end{description}\end{quote}

\end{fulllineitems}

\index{ArgInfo::get\_name (C++ function)@\spxentry{ArgInfo::get\_name}\spxextra{C++ function}}

\begin{fulllineitems}
\phantomsection\label{\detokenize{Developer_utils:_CPPv4NK9LAMMPS_NS7ArgInfo8get_nameEv}}\phantomsection\label{\detokenize{Developer_utils:_CPPv3NK9LAMMPS_NS7ArgInfo8get_nameEv}}\phantomsection\label{\detokenize{Developer_utils:_CPPv2NK9LAMMPS_NS7ArgInfo8get_nameEv}}\phantomsection\label{\detokenize{Developer_utils:LAMMPS_NS::ArgInfo::get_nameC}}
\pysigstartsignatures
\pysigstartmultiline
\pysiglinewithargsret{\phantomsection\label{\detokenize{Developer_utils:classLAMMPS__NS_1_1ArgInfo_1a272dda5e08548168f78061e194511fcf}}\DUrole{k}{inline}\DUrole{w}{ }\DUrole{k}{const}\DUrole{w}{ }\DUrole{kt}{char}\DUrole{w}{ }\DUrole{p}{*}\sphinxbfcode{\sphinxupquote{\DUrole{n}{get\_name}}}}{}{\DUrole{w}{ }\DUrole{k}{const}}
\pysigstopmultiline
\pysigstopsignatures
\sphinxAtStartPar
return reference to the ID or name of the reference

\sphinxAtStartPar
This string is pointing to an internal storage element and is only valid to use while the ArgInfo class instance is in scope. If you need a long\sphinxhyphen{}lived string make a copy with copy\_name().
\begin{quote}\begin{description}
\sphinxlineitem{Returns}
\sphinxAtStartPar
C\sphinxhyphen{}style char * string 

\end{description}\end{quote}

\end{fulllineitems}

\index{ArgInfo::copy\_name (C++ function)@\spxentry{ArgInfo::copy\_name}\spxextra{C++ function}}

\begin{fulllineitems}
\phantomsection\label{\detokenize{Developer_utils:_CPPv4N9LAMMPS_NS7ArgInfo9copy_nameEv}}\phantomsection\label{\detokenize{Developer_utils:_CPPv3N9LAMMPS_NS7ArgInfo9copy_nameEv}}\phantomsection\label{\detokenize{Developer_utils:_CPPv2N9LAMMPS_NS7ArgInfo9copy_nameEv}}\phantomsection\label{\detokenize{Developer_utils:LAMMPS_NS::ArgInfo::copy_name}}
\pysigstartsignatures
\pysigstartmultiline
\pysiglinewithargsret{\phantomsection\label{\detokenize{Developer_utils:classLAMMPS__NS_1_1ArgInfo_1aec1f62f5dcbdb968f21542406ecafb1c}}\DUrole{kt}{char}\DUrole{w}{ }\DUrole{p}{*}\sphinxbfcode{\sphinxupquote{\DUrole{n}{copy\_name}}}}{}{}
\pysigstopmultiline
\pysigstopsignatures
\sphinxAtStartPar
make copy of the ID of the reference as C\sphinxhyphen{}style string

\sphinxAtStartPar
The ID is copied into a buffer allocated with “new” and thus must be later deleted with “delete {[}{]}” to avoid a memory leak. Because it is a full copy in a newly allocated buffer, the lifetime of this string extends beyond the the time the ArgInfo class is in scope.
\begin{quote}\begin{description}
\sphinxlineitem{Returns}
\sphinxAtStartPar
copy of string as char * 

\end{description}\end{quote}

\end{fulllineitems}


\end{sphinxuseclass}
\end{fulllineitems}



\bigskip\hrule\bigskip



\section{File reader classes}
\label{\detokenize{Developer_utils:file-reader-classes}}\label{\detokenize{Developer_utils:id1}}
\sphinxAtStartPar
The purpose of the file reader classes is to simplify the recurring task
of reading and parsing files. They can use the
{\hyperref[\detokenize{Developer_utils:_CPPv4N9LAMMPS_NS14ValueTokenizerE}]{\sphinxcrossref{\sphinxcode{\sphinxupquote{ValueTokenizer}}}}} class to process
the read in text.  The {\hyperref[\detokenize{Developer_utils:_CPPv4N9LAMMPS_NS14TextFileReaderE}]{\sphinxcrossref{\sphinxcode{\sphinxupquote{TextFileReader}}}}} is a more general version while
{\hyperref[\detokenize{Developer_utils:_CPPv4N9LAMMPS_NS19PotentialFileReaderE}]{\sphinxcrossref{\sphinxcode{\sphinxupquote{PotentialFileReader}}}}} is
specialized to implement the behavior expected for looking up and
reading/parsing files with potential parameters in LAMMPS.  The
potential file reader class requires a LAMMPS instance, requires to be
run on MPI rank 0 only, will use the
{\hyperref[\detokenize{Developer_utils:_CPPv4N9LAMMPS_NS5utils23get_potential_file_pathERKNSt6stringE}]{\sphinxcrossref{\sphinxcode{\sphinxupquote{utils::get\_potential\_file\_path}}}}} function to look up and
open the file, and will call the \sphinxcode{\sphinxupquote{LAMMPS\_NS::Error}} class in
case of failures to read or to convert numbers, so that LAMMPS will be
aborted.
\sphinxSetupCaptionForVerbatim{Use of PotentialFileReader class in pair style coul/streitz}
\def\sphinxLiteralBlockLabel{\label{\detokenize{Developer_utils:id6}}}
\begin{sphinxVerbatim}[commandchars=\\\{\}]
\PYG{+w}{ }\PYG{n}{PotentialFileReader}\PYG{+w}{ }\PYG{n}{reader}\PYG{p}{(}\PYG{n}{lmp}\PYG{p}{,}\PYG{+w}{ }\PYG{n}{file}\PYG{p}{,}\PYG{+w}{ }\PYG{l+s}{\PYGZdq{}}\PYG{l+s}{coul/streitz}\PYG{l+s}{\PYGZdq{}}\PYG{p}{)}\PYG{p}{;}
\PYG{+w}{ }\PYG{k+kt}{char}\PYG{+w}{ }\PYG{o}{*}\PYG{+w}{ }\PYG{n}{line}\PYG{p}{;}

\PYG{+w}{ }\PYG{k}{while}\PYG{p}{(}\PYG{p}{(}\PYG{n}{line}\PYG{+w}{ }\PYG{o}{=}\PYG{+w}{ }\PYG{n}{reader}\PYG{p}{.}\PYG{n}{next\PYGZus{}line}\PYG{p}{(}\PYG{n}{NPARAMS\PYGZus{}PER\PYGZus{}LINE}\PYG{p}{)}\PYG{p}{)}\PYG{p}{)}\PYG{+w}{ }\PYG{p}{\PYGZob{}}
\PYG{+w}{   }\PYG{k}{try}\PYG{+w}{ }\PYG{p}{\PYGZob{}}
\PYG{+w}{     }\PYG{n}{ValueTokenizer}\PYG{+w}{ }\PYG{n+nf}{values}\PYG{p}{(}\PYG{n}{line}\PYG{p}{)}\PYG{p}{;}
\PYG{+w}{     }\PYG{n}{std}\PYG{o}{:}\PYG{o}{:}\PYG{n}{string}\PYG{+w}{ }\PYG{n}{iname}\PYG{+w}{ }\PYG{o}{=}\PYG{+w}{ }\PYG{n}{values}\PYG{p}{.}\PYG{n}{next\PYGZus{}string}\PYG{p}{(}\PYG{p}{)}\PYG{p}{;}

\PYG{+w}{     }\PYG{k+kt}{int}\PYG{+w}{ }\PYG{n}{ielement}\PYG{p}{;}
\PYG{+w}{     }\PYG{k}{for}\PYG{+w}{ }\PYG{p}{(}\PYG{n}{ielement}\PYG{+w}{ }\PYG{o}{=}\PYG{+w}{ }\PYG{l+m+mi}{0}\PYG{p}{;}\PYG{+w}{ }\PYG{n}{ielement}\PYG{+w}{ }\PYG{o}{\PYGZlt{}}\PYG{+w}{ }\PYG{n}{nelements}\PYG{p}{;}\PYG{+w}{ }\PYG{n}{ielement}\PYG{o}{+}\PYG{o}{+}\PYG{p}{)}
\PYG{+w}{       }\PYG{k}{if}\PYG{+w}{ }\PYG{p}{(}\PYG{n}{iname}\PYG{+w}{ }\PYG{o}{=}\PYG{o}{=}\PYG{+w}{ }\PYG{n}{elements}\PYG{p}{[}\PYG{n}{ielement}\PYG{p}{]}\PYG{p}{)}\PYG{+w}{ }\PYG{k}{break}\PYG{p}{;}

\PYG{+w}{     }\PYG{k}{if}\PYG{+w}{ }\PYG{p}{(}\PYG{n}{nparams}\PYG{+w}{ }\PYG{o}{=}\PYG{o}{=}\PYG{+w}{ }\PYG{n}{maxparam}\PYG{p}{)}\PYG{+w}{ }\PYG{p}{\PYGZob{}}
\PYG{+w}{       }\PYG{n}{maxparam}\PYG{+w}{ }\PYG{o}{+}\PYG{o}{=}\PYG{+w}{ }\PYG{n}{DELTA}\PYG{p}{;}
\PYG{+w}{       }\PYG{n}{params}\PYG{+w}{ }\PYG{o}{=}\PYG{+w}{ }\PYG{p}{(}\PYG{n}{Param}\PYG{+w}{ }\PYG{o}{*}\PYG{p}{)}\PYG{+w}{ }\PYG{n}{memory}\PYG{o}{\PYGZhy{}}\PYG{o}{\PYGZgt{}}\PYG{n}{srealloc}\PYG{p}{(}\PYG{n}{params}\PYG{p}{,}\PYG{n}{maxparam}\PYG{o}{*}\PYG{k}{sizeof}\PYG{p}{(}\PYG{n}{Param}\PYG{p}{)}\PYG{p}{,}
\PYG{+w}{                                           }\PYG{l+s}{\PYGZdq{}}\PYG{l+s}{pair:params}\PYG{l+s}{\PYGZdq{}}\PYG{p}{)}\PYG{p}{;}
\PYG{+w}{     }\PYG{p}{\PYGZcb{}}

\PYG{+w}{     }\PYG{n}{params}\PYG{p}{[}\PYG{n}{nparams}\PYG{p}{]}\PYG{p}{.}\PYG{n}{ielement}\PYG{+w}{ }\PYG{o}{=}\PYG{+w}{ }\PYG{n}{ielement}\PYG{p}{;}
\PYG{+w}{     }\PYG{n}{params}\PYG{p}{[}\PYG{n}{nparams}\PYG{p}{]}\PYG{p}{.}\PYG{n}{chi}\PYG{+w}{ }\PYG{o}{=}\PYG{+w}{ }\PYG{n}{values}\PYG{p}{.}\PYG{n}{next\PYGZus{}double}\PYG{p}{(}\PYG{p}{)}\PYG{p}{;}
\PYG{+w}{     }\PYG{n}{params}\PYG{p}{[}\PYG{n}{nparams}\PYG{p}{]}\PYG{p}{.}\PYG{n}{eta}\PYG{+w}{ }\PYG{o}{=}\PYG{+w}{ }\PYG{n}{values}\PYG{p}{.}\PYG{n}{next\PYGZus{}double}\PYG{p}{(}\PYG{p}{)}\PYG{p}{;}
\PYG{+w}{     }\PYG{n}{params}\PYG{p}{[}\PYG{n}{nparams}\PYG{p}{]}\PYG{p}{.}\PYG{n}{gamma}\PYG{+w}{ }\PYG{o}{=}\PYG{+w}{ }\PYG{n}{values}\PYG{p}{.}\PYG{n}{next\PYGZus{}double}\PYG{p}{(}\PYG{p}{)}\PYG{p}{;}
\PYG{+w}{     }\PYG{n}{params}\PYG{p}{[}\PYG{n}{nparams}\PYG{p}{]}\PYG{p}{.}\PYG{n}{zeta}\PYG{+w}{ }\PYG{o}{=}\PYG{+w}{ }\PYG{n}{values}\PYG{p}{.}\PYG{n}{next\PYGZus{}double}\PYG{p}{(}\PYG{p}{)}\PYG{p}{;}
\PYG{+w}{     }\PYG{n}{params}\PYG{p}{[}\PYG{n}{nparams}\PYG{p}{]}\PYG{p}{.}\PYG{n}{zcore}\PYG{+w}{ }\PYG{o}{=}\PYG{+w}{ }\PYG{n}{values}\PYG{p}{.}\PYG{n}{next\PYGZus{}double}\PYG{p}{(}\PYG{p}{)}\PYG{p}{;}

\PYG{+w}{   }\PYG{p}{\PYGZcb{}}\PYG{+w}{ }\PYG{k}{catch}\PYG{+w}{ }\PYG{p}{(}\PYG{n}{TokenizerException}\PYG{+w}{ }\PYG{o}{\PYGZam{}}\PYG{+w}{ }\PYG{n}{e}\PYG{p}{)}\PYG{+w}{ }\PYG{p}{\PYGZob{}}
\PYG{+w}{     }\PYG{n}{error}\PYG{o}{\PYGZhy{}}\PYG{o}{\PYGZgt{}}\PYG{n}{one}\PYG{p}{(}\PYG{n}{FLERR}\PYG{p}{,}\PYG{+w}{ }\PYG{n}{e}\PYG{p}{.}\PYG{n}{what}\PYG{p}{(}\PYG{p}{)}\PYG{p}{)}\PYG{p}{;}
\PYG{+w}{   }\PYG{p}{\PYGZcb{}}
\PYG{+w}{   }\PYG{n}{nparams}\PYG{o}{+}\PYG{o}{+}\PYG{p}{;}
\PYG{+w}{ }\PYG{p}{\PYGZcb{}}
\end{sphinxVerbatim}

\sphinxAtStartPar
A file that would be parsed by the reader code fragment looks like this:
\begin{sphinxalltt}
\# DATE: 2015\sphinxhyphen{}02\sphinxhyphen{}19 UNITS: metal CONTRIBUTOR: Ray Shan CITATION: Streitz and Mintmire, Phys Rev B, 50, 11996\sphinxhyphen{}12003 (1994)
\#
\# X (eV)                J (eV)          gamma (1/AA)   zeta (1/AA)    Z (e)

Al      0.000000        10.328655       0.000000        0.968438        0.763905
O       5.484763        14.035715       0.000000        2.143957        0.000000
\end{sphinxalltt}


\bigskip\hrule\bigskip

\index{TextFileReader (C++ class)@\spxentry{TextFileReader}\spxextra{C++ class}}

\begin{fulllineitems}
\phantomsection\label{\detokenize{Developer_utils:_CPPv4N9LAMMPS_NS14TextFileReaderE}}\phantomsection\label{\detokenize{Developer_utils:_CPPv3N9LAMMPS_NS14TextFileReaderE}}\phantomsection\label{\detokenize{Developer_utils:_CPPv2N9LAMMPS_NS14TextFileReaderE}}\phantomsection\label{\detokenize{Developer_utils:LAMMPS_NS::TextFileReader}}
\pysigstartsignatures
\pysigstartmultiline
\pysigline{\phantomsection\label{\detokenize{Developer_utils:classLAMMPS__NS_1_1TextFileReader}}\DUrole{k}{class}\DUrole{w}{ }\sphinxbfcode{\sphinxupquote{\DUrole{n}{TextFileReader}}}}
\pysigstopmultiline
\pysigstopsignatures
\begin{sphinxuseclass}{breathe-sectiondef}\subsubsection*{Public Functions}
\index{TextFileReader::TextFileReader (C++ function)@\spxentry{TextFileReader::TextFileReader}\spxextra{C++ function}}

\begin{fulllineitems}
\phantomsection\label{\detokenize{Developer_utils:_CPPv4N9LAMMPS_NS14TextFileReader14TextFileReaderERKNSt6stringERKNSt6stringE}}\phantomsection\label{\detokenize{Developer_utils:_CPPv3N9LAMMPS_NS14TextFileReader14TextFileReaderERKNSt6stringERKNSt6stringE}}\phantomsection\label{\detokenize{Developer_utils:_CPPv2N9LAMMPS_NS14TextFileReader14TextFileReaderERKNSt6stringERKNSt6stringE}}\phantomsection\label{\detokenize{Developer_utils:LAMMPS_NS::TextFileReader::TextFileReader__ssCR.ssCR}}
\pysigstartsignatures
\pysigstartmultiline
\pysiglinewithargsret{\phantomsection\label{\detokenize{Developer_utils:classLAMMPS__NS_1_1TextFileReader_1a2043b5f5a6222c33ed29f356e09fa50a}}\sphinxbfcode{\sphinxupquote{\DUrole{n}{TextFileReader}}}}{\DUrole{k}{const}\DUrole{w}{ }\DUrole{n}{std}\DUrole{p}{::}\DUrole{n}{string}\DUrole{w}{ }\DUrole{p}{\&}\DUrole{n,sig-param}{filename}\sphinxparamcomma \DUrole{k}{const}\DUrole{w}{ }\DUrole{n}{std}\DUrole{p}{::}\DUrole{n}{string}\DUrole{w}{ }\DUrole{p}{\&}\DUrole{n,sig-param}{filetype}}{}
\pysigstopmultiline
\pysigstopsignatures
\sphinxAtStartPar
Class for reading and parsing text files

\sphinxAtStartPar
The value of the class member variable \sphinxstyleemphasis{ignore\_comments} controls whether any text following the pound sign (\#) should be ignored (true) or not (false). Default: true, i.e. ignore. 
\sphinxAtStartPar
\begin{description}
\sphinxlineitem{\sphinxstyleemphasis{See also}}
\sphinxAtStartPar
{\hyperref[\detokenize{Developer_utils:_CPPv4N9LAMMPS_NS14TextFileReaderE}]{\sphinxcrossref{\sphinxcode{\sphinxupquote{TextFileReader}}}}}

\end{description}


\begin{quote}\begin{description}
\sphinxlineitem{Parameters}\begin{itemize}
\item {} 
\sphinxAtStartPar
\sphinxstyleliteralstrong{\sphinxupquote{filename}} \textendash{} Name of file to be read 

\item {} 
\sphinxAtStartPar
\sphinxstyleliteralstrong{\sphinxupquote{filetype}} \textendash{} Description of file type for error messages 

\end{itemize}

\end{description}\end{quote}

\end{fulllineitems}

\index{TextFileReader::TextFileReader (C++ function)@\spxentry{TextFileReader::TextFileReader}\spxextra{C++ function}}

\begin{fulllineitems}
\phantomsection\label{\detokenize{Developer_utils:_CPPv4N9LAMMPS_NS14TextFileReader14TextFileReaderEP4FILENSt6stringE}}\phantomsection\label{\detokenize{Developer_utils:_CPPv3N9LAMMPS_NS14TextFileReader14TextFileReaderEP4FILENSt6stringE}}\phantomsection\label{\detokenize{Developer_utils:_CPPv2N9LAMMPS_NS14TextFileReader14TextFileReaderEP4FILENSt6stringE}}\phantomsection\label{\detokenize{Developer_utils:LAMMPS_NS::TextFileReader::TextFileReader__FILEP.ss}}
\pysigstartsignatures
\pysigstartmultiline
\pysiglinewithargsret{\phantomsection\label{\detokenize{Developer_utils:classLAMMPS__NS_1_1TextFileReader_1ae85518b9022db4c1075338ccd8fe5de8}}\sphinxbfcode{\sphinxupquote{\DUrole{n}{TextFileReader}}}}{\DUrole{n}{FILE}\DUrole{w}{ }\DUrole{p}{*}\DUrole{n,sig-param}{fp}\sphinxparamcomma \DUrole{n}{std}\DUrole{p}{::}\DUrole{n}{string}\DUrole{w}{ }\DUrole{n,sig-param}{filetype}}{}
\pysigstopmultiline
\pysigstopsignatures
\sphinxAtStartPar
This is an overloaded member function, provided for convenience. It differs from the above function only in what argument(s) it accepts.

\sphinxAtStartPar
This function is useful in combination with {\hyperref[\detokenize{Developer_utils:_CPPv4N9LAMMPS_NS5utils14open_potentialERKNSt6stringEP6LAMMPSPi}]{\sphinxcrossref{\sphinxcode{\sphinxupquote{utils::open\_potential()}}}}}.

\begin{sphinxadmonition}{note}{Note:}
\sphinxAtStartPar
The FILE pointer is not closed in the destructor, but will be advanced
when reading from it.
\end{sphinxadmonition}
\begin{quote}\begin{description}
\sphinxlineitem{Parameters}\begin{itemize}
\item {} 
\sphinxAtStartPar
\sphinxstyleliteralstrong{\sphinxupquote{fp}} \textendash{} File descriptor of the already opened file 

\item {} 
\sphinxAtStartPar
\sphinxstyleliteralstrong{\sphinxupquote{filetype}} \textendash{} Description of file type for error messages 

\end{itemize}

\end{description}\end{quote}

\end{fulllineitems}

\index{TextFileReader::\textasciitilde{}TextFileReader (C++ function)@\spxentry{TextFileReader::\textasciitilde{}TextFileReader}\spxextra{C++ function}}

\begin{fulllineitems}
\phantomsection\label{\detokenize{Developer_utils:_CPPv4N9LAMMPS_NS14TextFileReaderD0Ev}}\phantomsection\label{\detokenize{Developer_utils:_CPPv3N9LAMMPS_NS14TextFileReaderD0Ev}}\phantomsection\label{\detokenize{Developer_utils:_CPPv2N9LAMMPS_NS14TextFileReaderD0Ev}}\phantomsection\label{\detokenize{Developer_utils:LAMMPS_NS::TextFileReader::_TextFileReader}}
\pysigstartsignatures
\pysigstartmultiline
\pysiglinewithargsret{\phantomsection\label{\detokenize{Developer_utils:classLAMMPS__NS_1_1TextFileReader_1a60ce1f5ac0e90a8deedc0f0359eacc49}}\DUrole{k}{virtual}\DUrole{w}{ }\sphinxbfcode{\sphinxupquote{\DUrole{n}{\textasciitilde{}TextFileReader}}}}{}{}
\pysigstopmultiline
\pysigstopsignatures
\sphinxAtStartPar
Closes the file 

\end{fulllineitems}

\index{TextFileReader::set\_bufsize (C++ function)@\spxentry{TextFileReader::set\_bufsize}\spxextra{C++ function}}

\begin{fulllineitems}
\phantomsection\label{\detokenize{Developer_utils:_CPPv4N9LAMMPS_NS14TextFileReader11set_bufsizeEi}}\phantomsection\label{\detokenize{Developer_utils:_CPPv3N9LAMMPS_NS14TextFileReader11set_bufsizeEi}}\phantomsection\label{\detokenize{Developer_utils:_CPPv2N9LAMMPS_NS14TextFileReader11set_bufsizeEi}}\phantomsection\label{\detokenize{Developer_utils:LAMMPS_NS::TextFileReader::set_bufsize__i}}
\pysigstartsignatures
\pysigstartmultiline
\pysiglinewithargsret{\phantomsection\label{\detokenize{Developer_utils:classLAMMPS__NS_1_1TextFileReader_1af2f99f85d2cb5e54a4190e4acdeeb4aa}}\DUrole{kt}{void}\DUrole{w}{ }\sphinxbfcode{\sphinxupquote{\DUrole{n}{set\_bufsize}}}}{\DUrole{kt}{int}}{}
\pysigstopmultiline
\pysigstopsignatures
\sphinxAtStartPar
adjust line buffer size 

\end{fulllineitems}

\index{TextFileReader::rewind (C++ function)@\spxentry{TextFileReader::rewind}\spxextra{C++ function}}

\begin{fulllineitems}
\phantomsection\label{\detokenize{Developer_utils:_CPPv4N9LAMMPS_NS14TextFileReader6rewindEv}}\phantomsection\label{\detokenize{Developer_utils:_CPPv3N9LAMMPS_NS14TextFileReader6rewindEv}}\phantomsection\label{\detokenize{Developer_utils:_CPPv2N9LAMMPS_NS14TextFileReader6rewindEv}}\phantomsection\label{\detokenize{Developer_utils:LAMMPS_NS::TextFileReader::rewind}}
\pysigstartsignatures
\pysigstartmultiline
\pysiglinewithargsret{\phantomsection\label{\detokenize{Developer_utils:classLAMMPS__NS_1_1TextFileReader_1ab8734e666421c9fe3b6380a818c6c727}}\DUrole{kt}{void}\DUrole{w}{ }\sphinxbfcode{\sphinxupquote{\DUrole{n}{rewind}}}}{}{}
\pysigstopmultiline
\pysigstopsignatures
\sphinxAtStartPar
Reset file to the beginning 

\end{fulllineitems}

\index{TextFileReader::skip\_line (C++ function)@\spxentry{TextFileReader::skip\_line}\spxextra{C++ function}}

\begin{fulllineitems}
\phantomsection\label{\detokenize{Developer_utils:_CPPv4N9LAMMPS_NS14TextFileReader9skip_lineEv}}\phantomsection\label{\detokenize{Developer_utils:_CPPv3N9LAMMPS_NS14TextFileReader9skip_lineEv}}\phantomsection\label{\detokenize{Developer_utils:_CPPv2N9LAMMPS_NS14TextFileReader9skip_lineEv}}\phantomsection\label{\detokenize{Developer_utils:LAMMPS_NS::TextFileReader::skip_line}}
\pysigstartsignatures
\pysigstartmultiline
\pysiglinewithargsret{\phantomsection\label{\detokenize{Developer_utils:classLAMMPS__NS_1_1TextFileReader_1a13171cc94df93641fa41710249bd2828}}\DUrole{kt}{void}\DUrole{w}{ }\sphinxbfcode{\sphinxupquote{\DUrole{n}{skip\_line}}}}{}{}
\pysigstopmultiline
\pysigstopsignatures
\sphinxAtStartPar
Read the next line and ignore it 

\end{fulllineitems}

\index{TextFileReader::next\_line (C++ function)@\spxentry{TextFileReader::next\_line}\spxextra{C++ function}}

\begin{fulllineitems}
\phantomsection\label{\detokenize{Developer_utils:_CPPv4N9LAMMPS_NS14TextFileReader9next_lineEi}}\phantomsection\label{\detokenize{Developer_utils:_CPPv3N9LAMMPS_NS14TextFileReader9next_lineEi}}\phantomsection\label{\detokenize{Developer_utils:_CPPv2N9LAMMPS_NS14TextFileReader9next_lineEi}}\phantomsection\label{\detokenize{Developer_utils:LAMMPS_NS::TextFileReader::next_line__i}}
\pysigstartsignatures
\pysigstartmultiline
\pysiglinewithargsret{\phantomsection\label{\detokenize{Developer_utils:classLAMMPS__NS_1_1TextFileReader_1ac8028df52c806625ec32fe25eb9e87c0}}\DUrole{kt}{char}\DUrole{w}{ }\DUrole{p}{*}\sphinxbfcode{\sphinxupquote{\DUrole{n}{next\_line}}}}{\DUrole{kt}{int}\DUrole{w}{ }\DUrole{n,sig-param}{nparams}\DUrole{w}{ }\DUrole{p}{=}\DUrole{w}{ }\DUrole{m}{0}}{}
\pysigstopmultiline
\pysigstopsignatures
\sphinxAtStartPar
Read the next line(s) until \sphinxstyleemphasis{nparams} words have been read.

\sphinxAtStartPar
This reads a line and counts the words in it, if the number is less than the requested number, it will read the next line, as well. Output will be a string with all read lines combined. The purpose is to somewhat replicate the reading behavior of formatted files in Fortran.

\sphinxAtStartPar
If the \sphinxstyleemphasis{ignore\_comments} class member has the value \sphinxstyleemphasis{true}, then any text read in is truncated at the first ‘\#’ character.
\begin{quote}\begin{description}
\sphinxlineitem{Parameters}
\sphinxAtStartPar
\sphinxstyleliteralstrong{\sphinxupquote{nparams}} \textendash{} Number of words that must be read. Default: 0 

\sphinxlineitem{Returns}
\sphinxAtStartPar
String with the concatenated text 

\end{description}\end{quote}

\end{fulllineitems}

\index{TextFileReader::next\_dvector (C++ function)@\spxentry{TextFileReader::next\_dvector}\spxextra{C++ function}}

\begin{fulllineitems}
\phantomsection\label{\detokenize{Developer_utils:_CPPv4N9LAMMPS_NS14TextFileReader12next_dvectorEPdi}}\phantomsection\label{\detokenize{Developer_utils:_CPPv3N9LAMMPS_NS14TextFileReader12next_dvectorEPdi}}\phantomsection\label{\detokenize{Developer_utils:_CPPv2N9LAMMPS_NS14TextFileReader12next_dvectorEPdi}}\phantomsection\label{\detokenize{Developer_utils:LAMMPS_NS::TextFileReader::next_dvector__doubleP.i}}
\pysigstartsignatures
\pysigstartmultiline
\pysiglinewithargsret{\phantomsection\label{\detokenize{Developer_utils:classLAMMPS__NS_1_1TextFileReader_1a4c74ed481f64285779f60633e6301a47}}\DUrole{kt}{void}\DUrole{w}{ }\sphinxbfcode{\sphinxupquote{\DUrole{n}{next\_dvector}}}}{\DUrole{kt}{double}\DUrole{w}{ }\DUrole{p}{*}\DUrole{n,sig-param}{list}\sphinxparamcomma \DUrole{kt}{int}\DUrole{w}{ }\DUrole{n,sig-param}{n}}{}
\pysigstopmultiline
\pysigstopsignatures
\sphinxAtStartPar
Read lines until \sphinxstyleemphasis{n} doubles have been read and stored in array \sphinxstyleemphasis{list}

\sphinxAtStartPar
This reads lines from the file using the next\_line() function, and splits them into floating\sphinxhyphen{}point numbers using the ValueTokenizer class and stores the number in the provided list.
\begin{quote}\begin{description}
\sphinxlineitem{Parameters}\begin{itemize}
\item {} 
\sphinxAtStartPar
\sphinxstyleliteralstrong{\sphinxupquote{list}} \textendash{} Pointer to array with suitable storage for \sphinxstyleemphasis{n} doubles 

\item {} 
\sphinxAtStartPar
\sphinxstyleliteralstrong{\sphinxupquote{n}} \textendash{} Number of doubles to be read 

\end{itemize}

\end{description}\end{quote}

\end{fulllineitems}

\index{TextFileReader::next\_values (C++ function)@\spxentry{TextFileReader::next\_values}\spxextra{C++ function}}

\begin{fulllineitems}
\phantomsection\label{\detokenize{Developer_utils:_CPPv4N9LAMMPS_NS14TextFileReader11next_valuesEiRKNSt6stringE}}\phantomsection\label{\detokenize{Developer_utils:_CPPv3N9LAMMPS_NS14TextFileReader11next_valuesEiRKNSt6stringE}}\phantomsection\label{\detokenize{Developer_utils:_CPPv2N9LAMMPS_NS14TextFileReader11next_valuesEiRKNSt6stringE}}\phantomsection\label{\detokenize{Developer_utils:LAMMPS_NS::TextFileReader::next_values__i.ssCR}}
\pysigstartsignatures
\pysigstartmultiline
\pysiglinewithargsret{\phantomsection\label{\detokenize{Developer_utils:classLAMMPS__NS_1_1TextFileReader_1a8a0ce533407ad4d77f92c807de0131b0}}{\hyperref[\detokenize{Developer_utils:_CPPv4N9LAMMPS_NS14ValueTokenizerE}]{\sphinxcrossref{\DUrole{n}{ValueTokenizer}}}}\DUrole{w}{ }\sphinxbfcode{\sphinxupquote{\DUrole{n}{next\_values}}}}{\DUrole{kt}{int}\DUrole{w}{ }\DUrole{n,sig-param}{nparams}\sphinxparamcomma \DUrole{k}{const}\DUrole{w}{ }\DUrole{n}{std}\DUrole{p}{::}\DUrole{n}{string}\DUrole{w}{ }\DUrole{p}{\&}\DUrole{n,sig-param}{separators}\DUrole{w}{ }\DUrole{p}{=}\DUrole{w}{ }\DUrole{n}{TOKENIZER\_DEFAULT\_SEPARATORS}}{}
\pysigstopmultiline
\pysigstopsignatures
\sphinxAtStartPar
Read text until \sphinxstyleemphasis{nparams} words are read and passed to a tokenizer object for custom parsing.

\sphinxAtStartPar
This reads lines from the file using the next\_line() function, and splits them into floating\sphinxhyphen{}point numbers using the ValueTokenizer class and stores the number in the provided list.
\begin{quote}\begin{description}
\sphinxlineitem{Parameters}\begin{itemize}
\item {} 
\sphinxAtStartPar
\sphinxstyleliteralstrong{\sphinxupquote{nparams}} \textendash{} Number of words to be read 

\item {} 
\sphinxAtStartPar
\sphinxstyleliteralstrong{\sphinxupquote{separators}} \textendash{} String with list of separators. 

\end{itemize}

\sphinxlineitem{Returns}
\sphinxAtStartPar
ValueTokenizer object for read in text 

\end{description}\end{quote}

\end{fulllineitems}


\end{sphinxuseclass}
\begin{sphinxuseclass}{breathe-sectiondef}\subsubsection*{Public Members}
\index{TextFileReader::ignore\_comments (C++ member)@\spxentry{TextFileReader::ignore\_comments}\spxextra{C++ member}}

\begin{fulllineitems}
\phantomsection\label{\detokenize{Developer_utils:_CPPv4N9LAMMPS_NS14TextFileReader15ignore_commentsE}}\phantomsection\label{\detokenize{Developer_utils:_CPPv3N9LAMMPS_NS14TextFileReader15ignore_commentsE}}\phantomsection\label{\detokenize{Developer_utils:_CPPv2N9LAMMPS_NS14TextFileReader15ignore_commentsE}}\phantomsection\label{\detokenize{Developer_utils:LAMMPS_NS::TextFileReader::ignore_comments__b}}
\pysigstartsignatures
\pysigstartmultiline
\pysigline{\phantomsection\label{\detokenize{Developer_utils:classLAMMPS__NS_1_1TextFileReader_1a87625d6c23cbea82833fceab91c31bc4}}\DUrole{kt}{bool}\DUrole{w}{ }\sphinxbfcode{\sphinxupquote{\DUrole{n}{ignore\_comments}}}}
\pysigstopmultiline
\pysigstopsignatures
\sphinxAtStartPar
Controls whether comments are ignored. 

\end{fulllineitems}


\end{sphinxuseclass}
\end{fulllineitems}

\index{PotentialFileReader (C++ class)@\spxentry{PotentialFileReader}\spxextra{C++ class}}

\begin{fulllineitems}
\phantomsection\label{\detokenize{Developer_utils:_CPPv4N9LAMMPS_NS19PotentialFileReaderE}}\phantomsection\label{\detokenize{Developer_utils:_CPPv3N9LAMMPS_NS19PotentialFileReaderE}}\phantomsection\label{\detokenize{Developer_utils:_CPPv2N9LAMMPS_NS19PotentialFileReaderE}}\phantomsection\label{\detokenize{Developer_utils:LAMMPS_NS::PotentialFileReader}}
\pysigstartsignatures
\pysigstartmultiline
\pysigline{\phantomsection\label{\detokenize{Developer_utils:classLAMMPS__NS_1_1PotentialFileReader}}\DUrole{k}{class}\DUrole{w}{ }\sphinxbfcode{\sphinxupquote{\DUrole{n}{PotentialFileReader}}}\DUrole{w}{ }\DUrole{p}{:}\DUrole{w}{ }\DUrole{k}{protected}\DUrole{w}{ }{\hyperref[\detokenize{Classes_lammps:_CPPv4N9LAMMPS_NS8PointersE}]{\sphinxcrossref{\DUrole{n}{Pointers}}}}}
\pysigstopmultiline
\pysigstopsignatures
\begin{sphinxuseclass}{breathe-sectiondef}\subsubsection*{Public Functions}
\index{PotentialFileReader::PotentialFileReader (C++ function)@\spxentry{PotentialFileReader::PotentialFileReader}\spxextra{C++ function}}

\begin{fulllineitems}
\phantomsection\label{\detokenize{Developer_utils:_CPPv4N9LAMMPS_NS19PotentialFileReader19PotentialFileReaderEP6LAMMPSRKNSt6stringERKNSt6stringERKNSt6stringEKi}}\phantomsection\label{\detokenize{Developer_utils:_CPPv3N9LAMMPS_NS19PotentialFileReader19PotentialFileReaderEP6LAMMPSRKNSt6stringERKNSt6stringERKNSt6stringEKi}}\phantomsection\label{\detokenize{Developer_utils:_CPPv2N9LAMMPS_NS19PotentialFileReader19PotentialFileReaderEP6LAMMPSRKNSt6stringERKNSt6stringERKNSt6stringEKi}}\phantomsection\label{\detokenize{Developer_utils:LAMMPS_NS::PotentialFileReader::PotentialFileReader__LAMMPSP.ssCR.ssCR.ssCR.iC}}
\pysigstartsignatures
\pysigstartmultiline
\pysiglinewithargsret{\phantomsection\label{\detokenize{Developer_utils:classLAMMPS__NS_1_1PotentialFileReader_1a149f46f04cd28f4d9d293ff333380d54}}\sphinxbfcode{\sphinxupquote{\DUrole{n}{PotentialFileReader}}}}{\DUrole{k}{class}\DUrole{w}{ }{\hyperref[\detokenize{Classes_lammps:_CPPv4N9LAMMPS_NS6LAMMPSE}]{\sphinxcrossref{\DUrole{n}{LAMMPS}}}}\DUrole{w}{ }\DUrole{p}{*}\DUrole{n,sig-param}{lmp}\sphinxparamcomma \DUrole{k}{const}\DUrole{w}{ }\DUrole{n}{std}\DUrole{p}{::}\DUrole{n}{string}\DUrole{w}{ }\DUrole{p}{\&}\DUrole{n,sig-param}{filename}\sphinxparamcomma \DUrole{k}{const}\DUrole{w}{ }\DUrole{n}{std}\DUrole{p}{::}\DUrole{n}{string}\DUrole{w}{ }\DUrole{p}{\&}\DUrole{n,sig-param}{potential\_name}\sphinxparamcomma \DUrole{k}{const}\DUrole{w}{ }\DUrole{n}{std}\DUrole{p}{::}\DUrole{n}{string}\DUrole{w}{ }\DUrole{p}{\&}\DUrole{n,sig-param}{name\_suffix}\sphinxparamcomma \DUrole{k}{const}\DUrole{w}{ }\DUrole{kt}{int}\DUrole{w}{ }\DUrole{n,sig-param}{auto\_convert}\DUrole{w}{ }\DUrole{p}{=}\DUrole{w}{ }\DUrole{m}{0}}{}
\pysigstopmultiline
\pysigstopsignatures
\sphinxAtStartPar
Class for reading and parsing {\hyperref[\detokenize{Classes_lammps:classLAMMPS__NS_1_1LAMMPS}]{\sphinxcrossref{\DUrole{std,std-ref}{LAMMPS}}}} potential files

\sphinxAtStartPar
The value of the class member variable \sphinxstyleemphasis{ignore\_comments} controls whether any text following the pound sign (\#) should be ignored (true) or not (false). Default: true, i.e. ignore. 
\sphinxAtStartPar
\begin{description}
\sphinxlineitem{\sphinxstyleemphasis{See also}}
\sphinxAtStartPar
{\hyperref[\detokenize{Developer_utils:_CPPv4N9LAMMPS_NS14TextFileReaderE}]{\sphinxcrossref{\sphinxcode{\sphinxupquote{TextFileReader}}}}}

\end{description}


\begin{quote}\begin{description}
\sphinxlineitem{Parameters}\begin{itemize}
\item {} 
\sphinxAtStartPar
\sphinxstyleliteralstrong{\sphinxupquote{lmp}} \textendash{} Pointer to {\hyperref[\detokenize{Classes_lammps:classLAMMPS__NS_1_1LAMMPS}]{\sphinxcrossref{\DUrole{std,std-ref}{LAMMPS}}}} instance 

\item {} 
\sphinxAtStartPar
\sphinxstyleliteralstrong{\sphinxupquote{filename}} \textendash{} Name of file to be read 

\item {} 
\sphinxAtStartPar
\sphinxstyleliteralstrong{\sphinxupquote{potential\_name}} \textendash{} Name of potential style for error messages 

\item {} 
\sphinxAtStartPar
\sphinxstyleliteralstrong{\sphinxupquote{name\_suffix}} \textendash{} Suffix added to potential name in error messages 

\item {} 
\sphinxAtStartPar
\sphinxstyleliteralstrong{\sphinxupquote{auto\_convert}} \textendash{} Bitmask of supported unit conversions 

\end{itemize}

\end{description}\end{quote}

\end{fulllineitems}

\index{PotentialFileReader::\textasciitilde{}PotentialFileReader (C++ function)@\spxentry{PotentialFileReader::\textasciitilde{}PotentialFileReader}\spxextra{C++ function}}

\begin{fulllineitems}
\phantomsection\label{\detokenize{Developer_utils:_CPPv4N9LAMMPS_NS19PotentialFileReaderD0Ev}}\phantomsection\label{\detokenize{Developer_utils:_CPPv3N9LAMMPS_NS19PotentialFileReaderD0Ev}}\phantomsection\label{\detokenize{Developer_utils:_CPPv2N9LAMMPS_NS19PotentialFileReaderD0Ev}}\phantomsection\label{\detokenize{Developer_utils:LAMMPS_NS::PotentialFileReader::_PotentialFileReader}}
\pysigstartsignatures
\pysigstartmultiline
\pysiglinewithargsret{\phantomsection\label{\detokenize{Developer_utils:classLAMMPS__NS_1_1PotentialFileReader_1a2fc140e8f6be2389ddffb5f7edce870e}}\sphinxbfcode{\sphinxupquote{\DUrole{n}{\textasciitilde{}PotentialFileReader}}}}{}{\DUrole{w}{ }\DUrole{k}{override}}
\pysigstopmultiline
\pysigstopsignatures
\sphinxAtStartPar
Closes the file 

\end{fulllineitems}

\index{PotentialFileReader::ignore\_comments (C++ function)@\spxentry{PotentialFileReader::ignore\_comments}\spxextra{C++ function}}

\begin{fulllineitems}
\phantomsection\label{\detokenize{Developer_utils:_CPPv4N9LAMMPS_NS19PotentialFileReader15ignore_commentsEb}}\phantomsection\label{\detokenize{Developer_utils:_CPPv3N9LAMMPS_NS19PotentialFileReader15ignore_commentsEb}}\phantomsection\label{\detokenize{Developer_utils:_CPPv2N9LAMMPS_NS19PotentialFileReader15ignore_commentsEb}}\phantomsection\label{\detokenize{Developer_utils:LAMMPS_NS::PotentialFileReader::ignore_comments__b}}
\pysigstartsignatures
\pysigstartmultiline
\pysiglinewithargsret{\phantomsection\label{\detokenize{Developer_utils:classLAMMPS__NS_1_1PotentialFileReader_1a84753973f8cd425807695e3c889f97f9}}\DUrole{kt}{void}\DUrole{w}{ }\sphinxbfcode{\sphinxupquote{\DUrole{n}{ignore\_comments}}}}{\DUrole{kt}{bool}\DUrole{w}{ }\DUrole{n,sig-param}{value}}{}
\pysigstopmultiline
\pysigstopsignatures
\sphinxAtStartPar
Set comment (= text after ‘\#’) handling preference for the file to be read
\begin{quote}\begin{description}
\sphinxlineitem{Parameters}
\sphinxAtStartPar
\sphinxstyleliteralstrong{\sphinxupquote{value}} \textendash{} Comment text is ignored if true, or not if false 

\end{description}\end{quote}

\end{fulllineitems}

\index{PotentialFileReader::rewind (C++ function)@\spxentry{PotentialFileReader::rewind}\spxextra{C++ function}}

\begin{fulllineitems}
\phantomsection\label{\detokenize{Developer_utils:_CPPv4N9LAMMPS_NS19PotentialFileReader6rewindEv}}\phantomsection\label{\detokenize{Developer_utils:_CPPv3N9LAMMPS_NS19PotentialFileReader6rewindEv}}\phantomsection\label{\detokenize{Developer_utils:_CPPv2N9LAMMPS_NS19PotentialFileReader6rewindEv}}\phantomsection\label{\detokenize{Developer_utils:LAMMPS_NS::PotentialFileReader::rewind}}
\pysigstartsignatures
\pysigstartmultiline
\pysiglinewithargsret{\phantomsection\label{\detokenize{Developer_utils:classLAMMPS__NS_1_1PotentialFileReader_1ab8734e666421c9fe3b6380a818c6c727}}\DUrole{kt}{void}\DUrole{w}{ }\sphinxbfcode{\sphinxupquote{\DUrole{n}{rewind}}}}{}{}
\pysigstopmultiline
\pysigstopsignatures
\sphinxAtStartPar
Reset file to the beginning 

\end{fulllineitems}

\index{PotentialFileReader::skip\_line (C++ function)@\spxentry{PotentialFileReader::skip\_line}\spxextra{C++ function}}

\begin{fulllineitems}
\phantomsection\label{\detokenize{Developer_utils:_CPPv4N9LAMMPS_NS19PotentialFileReader9skip_lineEv}}\phantomsection\label{\detokenize{Developer_utils:_CPPv3N9LAMMPS_NS19PotentialFileReader9skip_lineEv}}\phantomsection\label{\detokenize{Developer_utils:_CPPv2N9LAMMPS_NS19PotentialFileReader9skip_lineEv}}\phantomsection\label{\detokenize{Developer_utils:LAMMPS_NS::PotentialFileReader::skip_line}}
\pysigstartsignatures
\pysigstartmultiline
\pysiglinewithargsret{\phantomsection\label{\detokenize{Developer_utils:classLAMMPS__NS_1_1PotentialFileReader_1a13171cc94df93641fa41710249bd2828}}\DUrole{kt}{void}\DUrole{w}{ }\sphinxbfcode{\sphinxupquote{\DUrole{n}{skip\_line}}}}{}{}
\pysigstopmultiline
\pysigstopsignatures
\sphinxAtStartPar
Read a line but ignore its content 

\end{fulllineitems}

\index{PotentialFileReader::next\_line (C++ function)@\spxentry{PotentialFileReader::next\_line}\spxextra{C++ function}}

\begin{fulllineitems}
\phantomsection\label{\detokenize{Developer_utils:_CPPv4N9LAMMPS_NS19PotentialFileReader9next_lineEi}}\phantomsection\label{\detokenize{Developer_utils:_CPPv3N9LAMMPS_NS19PotentialFileReader9next_lineEi}}\phantomsection\label{\detokenize{Developer_utils:_CPPv2N9LAMMPS_NS19PotentialFileReader9next_lineEi}}\phantomsection\label{\detokenize{Developer_utils:LAMMPS_NS::PotentialFileReader::next_line__i}}
\pysigstartsignatures
\pysigstartmultiline
\pysiglinewithargsret{\phantomsection\label{\detokenize{Developer_utils:classLAMMPS__NS_1_1PotentialFileReader_1ac8028df52c806625ec32fe25eb9e87c0}}\DUrole{kt}{char}\DUrole{w}{ }\DUrole{p}{*}\sphinxbfcode{\sphinxupquote{\DUrole{n}{next\_line}}}}{\DUrole{kt}{int}\DUrole{w}{ }\DUrole{n,sig-param}{nparams}\DUrole{w}{ }\DUrole{p}{=}\DUrole{w}{ }\DUrole{m}{0}}{}
\pysigstopmultiline
\pysigstopsignatures
\sphinxAtStartPar
Read the next line(s) until \sphinxstyleemphasis{nparams} words have been read.

\sphinxAtStartPar
This reads a line and counts the words in it, if the number is less than the requested number, it will read the next line, as well. Output will be a string with all read lines combined. The purpose is to somewhat replicate the reading behavior of formatted files in Fortran.
\begin{quote}\begin{description}
\sphinxlineitem{Parameters}
\sphinxAtStartPar
\sphinxstyleliteralstrong{\sphinxupquote{nparams}} \textendash{} Number of words that must be read. Default: 0 

\sphinxlineitem{Returns}
\sphinxAtStartPar
String with the concatenated text 

\end{description}\end{quote}

\end{fulllineitems}

\index{PotentialFileReader::next\_dvector (C++ function)@\spxentry{PotentialFileReader::next\_dvector}\spxextra{C++ function}}

\begin{fulllineitems}
\phantomsection\label{\detokenize{Developer_utils:_CPPv4N9LAMMPS_NS19PotentialFileReader12next_dvectorEPdi}}\phantomsection\label{\detokenize{Developer_utils:_CPPv3N9LAMMPS_NS19PotentialFileReader12next_dvectorEPdi}}\phantomsection\label{\detokenize{Developer_utils:_CPPv2N9LAMMPS_NS19PotentialFileReader12next_dvectorEPdi}}\phantomsection\label{\detokenize{Developer_utils:LAMMPS_NS::PotentialFileReader::next_dvector__doubleP.i}}
\pysigstartsignatures
\pysigstartmultiline
\pysiglinewithargsret{\phantomsection\label{\detokenize{Developer_utils:classLAMMPS__NS_1_1PotentialFileReader_1a4c74ed481f64285779f60633e6301a47}}\DUrole{kt}{void}\DUrole{w}{ }\sphinxbfcode{\sphinxupquote{\DUrole{n}{next\_dvector}}}}{\DUrole{kt}{double}\DUrole{w}{ }\DUrole{p}{*}\DUrole{n,sig-param}{list}\sphinxparamcomma \DUrole{kt}{int}\DUrole{w}{ }\DUrole{n,sig-param}{n}}{}
\pysigstopmultiline
\pysigstopsignatures
\sphinxAtStartPar
Read lines until \sphinxstyleemphasis{n} doubles have been read and stored in array \sphinxstyleemphasis{list}

\sphinxAtStartPar
This reads lines from the file using the next\_line() function, and splits them into floating\sphinxhyphen{}point numbers using the ValueTokenizer class and stores the number in the provided list.
\begin{quote}\begin{description}
\sphinxlineitem{Parameters}\begin{itemize}
\item {} 
\sphinxAtStartPar
\sphinxstyleliteralstrong{\sphinxupquote{list}} \textendash{} Pointer to array with suitable storage for \sphinxstyleemphasis{n} doubles 

\item {} 
\sphinxAtStartPar
\sphinxstyleliteralstrong{\sphinxupquote{n}} \textendash{} Number of doubles to be read 

\end{itemize}

\end{description}\end{quote}

\end{fulllineitems}

\index{PotentialFileReader::next\_values (C++ function)@\spxentry{PotentialFileReader::next\_values}\spxextra{C++ function}}

\begin{fulllineitems}
\phantomsection\label{\detokenize{Developer_utils:_CPPv4N9LAMMPS_NS19PotentialFileReader11next_valuesEiRKNSt6stringE}}\phantomsection\label{\detokenize{Developer_utils:_CPPv3N9LAMMPS_NS19PotentialFileReader11next_valuesEiRKNSt6stringE}}\phantomsection\label{\detokenize{Developer_utils:_CPPv2N9LAMMPS_NS19PotentialFileReader11next_valuesEiRKNSt6stringE}}\phantomsection\label{\detokenize{Developer_utils:LAMMPS_NS::PotentialFileReader::next_values__i.ssCR}}
\pysigstartsignatures
\pysigstartmultiline
\pysiglinewithargsret{\phantomsection\label{\detokenize{Developer_utils:classLAMMPS__NS_1_1PotentialFileReader_1a8a0ce533407ad4d77f92c807de0131b0}}{\hyperref[\detokenize{Developer_utils:_CPPv4N9LAMMPS_NS14ValueTokenizerE}]{\sphinxcrossref{\DUrole{n}{ValueTokenizer}}}}\DUrole{w}{ }\sphinxbfcode{\sphinxupquote{\DUrole{n}{next\_values}}}}{\DUrole{kt}{int}\DUrole{w}{ }\DUrole{n,sig-param}{nparams}\sphinxparamcomma \DUrole{k}{const}\DUrole{w}{ }\DUrole{n}{std}\DUrole{p}{::}\DUrole{n}{string}\DUrole{w}{ }\DUrole{p}{\&}\DUrole{n,sig-param}{separators}\DUrole{w}{ }\DUrole{p}{=}\DUrole{w}{ }\DUrole{n}{TOKENIZER\_DEFAULT\_SEPARATORS}}{}
\pysigstopmultiline
\pysigstopsignatures
\sphinxAtStartPar
Read text until \sphinxstyleemphasis{nparams} words are read and passed to a tokenizer object for custom parsing.

\sphinxAtStartPar
This reads lines from the file using the next\_line() function, and splits them into floating\sphinxhyphen{}point numbers using the ValueTokenizer class and stores the number in the provided list.
\begin{quote}\begin{description}
\sphinxlineitem{Parameters}\begin{itemize}
\item {} 
\sphinxAtStartPar
\sphinxstyleliteralstrong{\sphinxupquote{nparams}} \textendash{} Number of words to be read 

\item {} 
\sphinxAtStartPar
\sphinxstyleliteralstrong{\sphinxupquote{separators}} \textendash{} String with list of separators. 

\end{itemize}

\sphinxlineitem{Returns}
\sphinxAtStartPar
ValueTokenizer object for read in text 

\end{description}\end{quote}

\end{fulllineitems}

\index{PotentialFileReader::next\_double (C++ function)@\spxentry{PotentialFileReader::next\_double}\spxextra{C++ function}}

\begin{fulllineitems}
\phantomsection\label{\detokenize{Developer_utils:_CPPv4N9LAMMPS_NS19PotentialFileReader11next_doubleEv}}\phantomsection\label{\detokenize{Developer_utils:_CPPv3N9LAMMPS_NS19PotentialFileReader11next_doubleEv}}\phantomsection\label{\detokenize{Developer_utils:_CPPv2N9LAMMPS_NS19PotentialFileReader11next_doubleEv}}\phantomsection\label{\detokenize{Developer_utils:LAMMPS_NS::PotentialFileReader::next_double}}
\pysigstartsignatures
\pysigstartmultiline
\pysiglinewithargsret{\phantomsection\label{\detokenize{Developer_utils:classLAMMPS__NS_1_1PotentialFileReader_1aa6b2325cb771c1d1aa11cc5bd3c3f2e9}}\DUrole{kt}{double}\DUrole{w}{ }\sphinxbfcode{\sphinxupquote{\DUrole{n}{next\_double}}}}{}{}
\pysigstopmultiline
\pysigstopsignatures
\sphinxAtStartPar
Read next line and convert first word to a double
\begin{quote}\begin{description}
\sphinxlineitem{Returns}
\sphinxAtStartPar
Value of first word in line as double 

\end{description}\end{quote}

\end{fulllineitems}

\index{PotentialFileReader::next\_int (C++ function)@\spxentry{PotentialFileReader::next\_int}\spxextra{C++ function}}

\begin{fulllineitems}
\phantomsection\label{\detokenize{Developer_utils:_CPPv4N9LAMMPS_NS19PotentialFileReader8next_intEv}}\phantomsection\label{\detokenize{Developer_utils:_CPPv3N9LAMMPS_NS19PotentialFileReader8next_intEv}}\phantomsection\label{\detokenize{Developer_utils:_CPPv2N9LAMMPS_NS19PotentialFileReader8next_intEv}}\phantomsection\label{\detokenize{Developer_utils:LAMMPS_NS::PotentialFileReader::next_int}}
\pysigstartsignatures
\pysigstartmultiline
\pysiglinewithargsret{\phantomsection\label{\detokenize{Developer_utils:classLAMMPS__NS_1_1PotentialFileReader_1a659db78ef3cd8cdbfa35654fbea28929}}\DUrole{kt}{int}\DUrole{w}{ }\sphinxbfcode{\sphinxupquote{\DUrole{n}{next\_int}}}}{}{}
\pysigstopmultiline
\pysigstopsignatures
\sphinxAtStartPar
Read next line and convert first word to an int
\begin{quote}\begin{description}
\sphinxlineitem{Returns}
\sphinxAtStartPar
Value of first word in line as int 

\end{description}\end{quote}

\end{fulllineitems}

\index{PotentialFileReader::next\_tagint (C++ function)@\spxentry{PotentialFileReader::next\_tagint}\spxextra{C++ function}}

\begin{fulllineitems}
\phantomsection\label{\detokenize{Developer_utils:_CPPv4N9LAMMPS_NS19PotentialFileReader11next_tagintEv}}\phantomsection\label{\detokenize{Developer_utils:_CPPv3N9LAMMPS_NS19PotentialFileReader11next_tagintEv}}\phantomsection\label{\detokenize{Developer_utils:_CPPv2N9LAMMPS_NS19PotentialFileReader11next_tagintEv}}\phantomsection\label{\detokenize{Developer_utils:LAMMPS_NS::PotentialFileReader::next_tagint}}
\pysigstartsignatures
\pysigstartmultiline
\pysiglinewithargsret{\phantomsection\label{\detokenize{Developer_utils:classLAMMPS__NS_1_1PotentialFileReader_1a388a19376fd8647693fa10077e29a68f}}\DUrole{n}{tagint}\DUrole{w}{ }\sphinxbfcode{\sphinxupquote{\DUrole{n}{next\_tagint}}}}{}{}
\pysigstopmultiline
\pysigstopsignatures
\sphinxAtStartPar
Read next line and convert first word to a tagint
\begin{quote}\begin{description}
\sphinxlineitem{Returns}
\sphinxAtStartPar
Value of first word in line as tagint 

\end{description}\end{quote}

\end{fulllineitems}

\index{PotentialFileReader::next\_bigint (C++ function)@\spxentry{PotentialFileReader::next\_bigint}\spxextra{C++ function}}

\begin{fulllineitems}
\phantomsection\label{\detokenize{Developer_utils:_CPPv4N9LAMMPS_NS19PotentialFileReader11next_bigintEv}}\phantomsection\label{\detokenize{Developer_utils:_CPPv3N9LAMMPS_NS19PotentialFileReader11next_bigintEv}}\phantomsection\label{\detokenize{Developer_utils:_CPPv2N9LAMMPS_NS19PotentialFileReader11next_bigintEv}}\phantomsection\label{\detokenize{Developer_utils:LAMMPS_NS::PotentialFileReader::next_bigint}}
\pysigstartsignatures
\pysigstartmultiline
\pysiglinewithargsret{\phantomsection\label{\detokenize{Developer_utils:classLAMMPS__NS_1_1PotentialFileReader_1a960e60b9cac96e1abfc9beb6a436e414}}\DUrole{n}{bigint}\DUrole{w}{ }\sphinxbfcode{\sphinxupquote{\DUrole{n}{next\_bigint}}}}{}{}
\pysigstopmultiline
\pysigstopsignatures
\sphinxAtStartPar
Read next line and convert first word to a bigint
\begin{quote}\begin{description}
\sphinxlineitem{Returns}
\sphinxAtStartPar
Value of first word in line as bigint 

\end{description}\end{quote}

\end{fulllineitems}

\index{PotentialFileReader::next\_string (C++ function)@\spxentry{PotentialFileReader::next\_string}\spxextra{C++ function}}

\begin{fulllineitems}
\phantomsection\label{\detokenize{Developer_utils:_CPPv4N9LAMMPS_NS19PotentialFileReader11next_stringEv}}\phantomsection\label{\detokenize{Developer_utils:_CPPv3N9LAMMPS_NS19PotentialFileReader11next_stringEv}}\phantomsection\label{\detokenize{Developer_utils:_CPPv2N9LAMMPS_NS19PotentialFileReader11next_stringEv}}\phantomsection\label{\detokenize{Developer_utils:LAMMPS_NS::PotentialFileReader::next_string}}
\pysigstartsignatures
\pysigstartmultiline
\pysiglinewithargsret{\phantomsection\label{\detokenize{Developer_utils:classLAMMPS__NS_1_1PotentialFileReader_1ad19df6455b8dfbfd2e5b1d93fc72d257}}\DUrole{n}{std}\DUrole{p}{::}\DUrole{n}{string}\DUrole{w}{ }\sphinxbfcode{\sphinxupquote{\DUrole{n}{next\_string}}}}{}{}
\pysigstopmultiline
\pysigstopsignatures
\sphinxAtStartPar
Read next line and return first word
\begin{quote}\begin{description}
\sphinxlineitem{Returns}
\sphinxAtStartPar
First word of read in line 

\end{description}\end{quote}

\end{fulllineitems}


\end{sphinxuseclass}
\end{fulllineitems}



\bigskip\hrule\bigskip



\section{Memory pool classes}
\label{\detokenize{Developer_utils:memory-pool-classes}}
\sphinxAtStartPar
The memory pool classes are used for cases where otherwise many
small memory allocations would be needed and where the data would
be either all used or all freed.  One example for that is the
storage of neighbor lists.  The memory management strategy is
based on the assumption that allocations will be in chunks of similar
sizes.  The allocation is then not done per individual call for a
reserved chunk of memory, but for a “page” that can hold multiple
chunks of data.  A parameter for the maximum chunk size must be
provided, as that is used to determine whether a new page of memory
must be used.

\sphinxAtStartPar
The {\hyperref[\detokenize{Developer_utils:_CPPv4I0EN9LAMMPS_NS6MyPageE}]{\sphinxcrossref{\sphinxcode{\sphinxupquote{MyPage}}}}} class offers two ways to
reserve a chunk: 1) with {\hyperref[\detokenize{Developer_utils:_CPPv4N9LAMMPS_NS6MyPage3getEi}]{\sphinxcrossref{\sphinxcode{\sphinxupquote{MyPage::get()}}}}} the
chunk size needs to be known in advance, 2) with {\hyperref[\detokenize{Developer_utils:_CPPv4N9LAMMPS_NS6MyPage4vgetEv}]{\sphinxcrossref{\sphinxcode{\sphinxupquote{MyPage::vget()}}}}} a pointer to the next chunk is returned, but
its size is registered later with {\hyperref[\detokenize{Developer_utils:_CPPv4N9LAMMPS_NS6MyPage4vgotEi}]{\sphinxcrossref{\sphinxcode{\sphinxupquote{MyPage::vgot()}}}}}.
\sphinxSetupCaptionForVerbatim{Example of using {\hyperref[\detokenize{Developer_utils:_CPPv4I0EN9LAMMPS_NS6MyPageE}]{\sphinxcrossref{\sphinxcode{\sphinxupquote{MyPage}}}}}}
\def\sphinxLiteralBlockLabel{\label{\detokenize{Developer_utils:id7}}}
\begin{sphinxVerbatim}[commandchars=\\\{\}]
\PYG{+w}{   }\PYG{c+cp}{\PYGZsh{}}\PYG{c+cp}{include}\PYG{+w}{ }\PYG{c+cpf}{\PYGZdq{}my\PYGZus{}page.h\PYGZdq{}}
\PYG{+w}{   }\PYG{k}{using}\PYG{+w}{ }\PYG{k}{namespace}\PYG{+w}{ }\PYG{n+nn}{LAMMPS\PYGZus{}NS}\PYG{p}{;}

\PYG{+w}{   }\PYG{n}{MyPage}\PYG{o}{\PYGZlt{}}\PYG{k+kt}{double}\PYG{o}{\PYGZgt{}}\PYG{+w}{ }\PYG{o}{*}\PYG{n}{dpage}\PYG{+w}{ }\PYG{o}{=}\PYG{+w}{ }\PYG{k}{new}\PYG{+w}{ }\PYG{n}{MyPage}\PYG{o}{\PYGZlt{}}\PYG{k+kt}{double}\PYG{o}{\PYGZgt{}}\PYG{p}{;}
\PYG{+w}{   }\PYG{c+c1}{// max size of chunk: 256, size of page: 10240 doubles (=81920 bytes)}
\PYG{+w}{   }\PYG{n}{dpage}\PYG{o}{\PYGZhy{}}\PYG{o}{\PYGZgt{}}\PYG{n}{init}\PYG{p}{(}\PYG{l+m+mi}{256}\PYG{p}{,}\PYG{l+m+mi}{10240}\PYG{p}{)}\PYG{p}{;}

\PYG{+w}{   }\PYG{k+kt}{double}\PYG{+w}{ }\PYG{o}{*}\PYG{o}{*}\PYG{n+nf}{build\PYGZus{}some\PYGZus{}lists}\PYG{p}{(}\PYG{k+kt}{int}\PYG{+w}{ }\PYG{n}{num}\PYG{p}{)}
\PYG{+w}{   }\PYG{p}{\PYGZob{}}
\PYG{+w}{       }\PYG{n}{dpage}\PYG{o}{\PYGZhy{}}\PYG{o}{\PYGZgt{}}\PYG{n}{reset}\PYG{p}{(}\PYG{p}{)}\PYG{p}{;}
\PYG{+w}{       }\PYG{k+kt}{double}\PYG{+w}{ }\PYG{o}{*}\PYG{o}{*}\PYG{n}{dlist}\PYG{+w}{ }\PYG{o}{=}\PYG{+w}{ }\PYG{k}{new}\PYG{+w}{ }\PYG{k+kt}{double}\PYG{o}{*}\PYG{p}{[}\PYG{n}{num}\PYG{p}{]}\PYG{p}{;}
\PYG{+w}{       }\PYG{k}{for}\PYG{+w}{ }\PYG{p}{(}\PYG{k+kt}{int}\PYG{+w}{ }\PYG{n}{i}\PYG{o}{=}\PYG{l+m+mi}{0}\PYG{p}{;}\PYG{+w}{ }\PYG{n}{i}\PYG{+w}{ }\PYG{o}{\PYGZlt{}}\PYG{+w}{ }\PYG{n}{num}\PYG{p}{;}\PYG{+w}{ }\PYG{o}{+}\PYG{o}{+}\PYG{n}{i}\PYG{p}{)}\PYG{+w}{ }\PYG{p}{\PYGZob{}}
\PYG{+w}{           }\PYG{k+kt}{double}\PYG{+w}{ }\PYG{o}{*}\PYG{n}{dptr}\PYG{+w}{ }\PYG{o}{=}\PYG{+w}{ }\PYG{n}{dpage}\PYG{p}{.}\PYG{n}{vget}\PYG{p}{(}\PYG{p}{)}\PYG{p}{;}
\PYG{+w}{           }\PYG{k+kt}{int}\PYG{+w}{ }\PYG{n}{jnum}\PYG{+w}{ }\PYG{o}{=}\PYG{+w}{ }\PYG{l+m+mi}{0}\PYG{p}{;}
\PYG{+w}{           }\PYG{k}{for}\PYG{+w}{ }\PYG{p}{(}\PYG{k+kt}{int}\PYG{+w}{ }\PYG{n}{j}\PYG{o}{=}\PYG{l+m+mi}{0}\PYG{p}{;}\PYG{+w}{ }\PYG{n}{j}\PYG{+w}{ }\PYG{o}{\PYGZlt{}}\PYG{+w}{ }\PYG{n}{jmax}\PYG{p}{;}\PYG{+w}{ }\PYG{o}{+}\PYG{o}{+}\PYG{n}{j}\PYG{p}{)}\PYG{+w}{ }\PYG{p}{\PYGZob{}}
\PYG{+w}{               }\PYG{c+c1}{// compute some dvalue for eligible loop index j}
\PYG{+w}{               }\PYG{n}{dptr}\PYG{p}{[}\PYG{n}{j}\PYG{p}{]}\PYG{+w}{ }\PYG{o}{=}\PYG{+w}{ }\PYG{n}{dvalue}\PYG{p}{;}
\PYG{+w}{               }\PYG{o}{+}\PYG{o}{+}\PYG{n}{jnum}\PYG{p}{;}
\PYG{+w}{           }\PYG{p}{\PYGZcb{}}
\PYG{+w}{           }\PYG{k}{if}\PYG{+w}{ }\PYG{p}{(}\PYG{n}{dpage}\PYG{p}{.}\PYG{n}{status}\PYG{p}{(}\PYG{p}{)}\PYG{+w}{ }\PYG{o}{!}\PYG{o}{=}\PYG{+w}{ }\PYG{l+m+mi}{0}\PYG{p}{)}\PYG{+w}{ }\PYG{p}{\PYGZob{}}
\PYG{+w}{               }\PYG{c+c1}{// handle out of memory or jnum too large errors}
\PYG{+w}{           }\PYG{p}{\PYGZcb{}}
\PYG{+w}{           }\PYG{n}{dpage}\PYG{p}{.}\PYG{n}{vgot}\PYG{p}{(}\PYG{n}{jnum}\PYG{p}{)}\PYG{p}{;}
\PYG{+w}{           }\PYG{n}{dlist}\PYG{p}{[}\PYG{n}{i}\PYG{p}{]}\PYG{+w}{ }\PYG{o}{=}\PYG{+w}{ }\PYG{n}{dptr}\PYG{p}{;}
\PYG{+w}{       }\PYG{p}{\PYGZcb{}}
\PYG{+w}{       }\PYG{k}{return}\PYG{+w}{ }\PYG{n}{dlist}\PYG{p}{;}
\PYG{+w}{   }\PYG{p}{\PYGZcb{}}
\end{sphinxVerbatim}


\bigskip\hrule\bigskip

\index{MyPage (C++ class)@\spxentry{MyPage}\spxextra{C++ class}}

\begin{fulllineitems}
\phantomsection\label{\detokenize{Developer_utils:_CPPv4I0EN9LAMMPS_NS6MyPageE}}\phantomsection\label{\detokenize{Developer_utils:_CPPv3I0EN9LAMMPS_NS6MyPageE}}\phantomsection\label{\detokenize{Developer_utils:_CPPv2I0EN9LAMMPS_NS6MyPageE}}
\pysigstartsignatures
\pysigstartmultiline
\pysigline{\DUrole{k}{template}\DUrole{p}{\textless{}}\DUrole{k}{class}\DUrole{w}{ }\sphinxbfcode{\sphinxupquote{\DUrole{n}{T}}}\DUrole{p}{\textgreater{}}}
\pysigline{\phantomsection\label{\detokenize{Developer_utils:classLAMMPS__NS_1_1MyPage}}\DUrole{k}{class}\DUrole{w}{ }\sphinxbfcode{\sphinxupquote{\DUrole{n}{MyPage}}}}
\pysigstopmultiline
\pysigstopsignatures
\sphinxAtStartPar
Templated class for storing chunks of datums in pages. 

\sphinxAtStartPar
The size of the chunk may vary from call to call, but must be less or equal than the \sphinxstyleemphasis{maxchunk} setting. The chunks are not returnable like with malloc() (i.e. you cannot call free() on them individually). One can only reset and start over. The purpose of this class is to replace many small memory allocations via malloc() with a few large ones. Since the pages are never freed until the class is re\sphinxhyphen{}initialized, they can be re\sphinxhyphen{}used without having to re\sphinxhyphen{}allocate them by calling the {\hyperref[\detokenize{Developer_utils:classLAMMPS__NS_1_1MyPage_1ad20897c5c8bd47f5d4005989bead0e55}]{\sphinxcrossref{\DUrole{std,std-ref}{reset()}}}} method.

\sphinxAtStartPar
The settings \sphinxstyleemphasis{maxchunk}, \sphinxstyleemphasis{pagesize}, and \sphinxstyleemphasis{pagedelta} control the memory allocation strategy. The \sphinxstyleemphasis{maxchunk} value represents the expected largest number of items per chunk. If there is less space left on the current page, a new page is allocated for the next chunk. The \sphinxstyleemphasis{pagesize} value represents how many items can fit on a single page. It should have space for multiple chunks of size \sphinxstyleemphasis{maxchunk}. The combination of these two parameters determines how much memory is wasted by either switching to the next page too soon or allocating too large pages that never get properly used. An error is issued if a requested chunk is larger than \sphinxstyleemphasis{maxchunk}. The \sphinxstyleemphasis{pagedelta} parameter determines how many pages are allocated in one go. In combination with the \sphinxstyleemphasis{pagesize} setting, this determines how often blocks of memory get allocated (fewer allocations will result in faster execution).

\begin{sphinxadmonition}{note}{Note:}
\sphinxAtStartPar
This is a template class with explicit instantiation. If the class is used with a new data type, a new explicit instantiation may need to be added at the end of the file \sphinxcode{\sphinxupquote{src/my\_page.cpp}} to avoid symbol lookup errors. 
\end{sphinxadmonition}

\begin{sphinxuseclass}{breathe-sectiondef}\subsubsection*{Public Functions}
\index{MyPage::MyPage (C++ function)@\spxentry{MyPage::MyPage}\spxextra{C++ function}}

\begin{fulllineitems}
\phantomsection\label{\detokenize{Developer_utils:_CPPv4N9LAMMPS_NS6MyPage6MyPageEv}}\phantomsection\label{\detokenize{Developer_utils:_CPPv3N9LAMMPS_NS6MyPage6MyPageEv}}\phantomsection\label{\detokenize{Developer_utils:_CPPv2N9LAMMPS_NS6MyPage6MyPageEv}}\phantomsection\label{\detokenize{Developer_utils:LAMMPS_NS::MyPage::MyPage}}
\pysigstartsignatures
\pysigstartmultiline
\pysiglinewithargsret{\phantomsection\label{\detokenize{Developer_utils:classLAMMPS__NS_1_1MyPage_1aaea419ba80ccc8fc23bcfd0870c1e875}}\sphinxbfcode{\sphinxupquote{\DUrole{n}{MyPage}}}}{}{}
\pysigstopmultiline
\pysigstopsignatures
\sphinxAtStartPar
Create a class instance

\sphinxAtStartPar
Need to call {\hyperref[\detokenize{Developer_utils:classLAMMPS__NS_1_1MyPage_1a381c748265a923272bd332b3905a9891}]{\sphinxcrossref{\DUrole{std,std-ref}{init()}}}} before use to define allocation settings 

\end{fulllineitems}

\index{MyPage::init (C++ function)@\spxentry{MyPage::init}\spxextra{C++ function}}

\begin{fulllineitems}
\phantomsection\label{\detokenize{Developer_utils:_CPPv4N9LAMMPS_NS6MyPage4initEiii}}\phantomsection\label{\detokenize{Developer_utils:_CPPv3N9LAMMPS_NS6MyPage4initEiii}}\phantomsection\label{\detokenize{Developer_utils:_CPPv2N9LAMMPS_NS6MyPage4initEiii}}\phantomsection\label{\detokenize{Developer_utils:LAMMPS_NS::MyPage::init__i.i.i}}
\pysigstartsignatures
\pysigstartmultiline
\pysiglinewithargsret{\phantomsection\label{\detokenize{Developer_utils:classLAMMPS__NS_1_1MyPage_1a381c748265a923272bd332b3905a9891}}\DUrole{kt}{int}\DUrole{w}{ }\sphinxbfcode{\sphinxupquote{\DUrole{n}{init}}}}{\DUrole{kt}{int}\DUrole{w}{ }\DUrole{n,sig-param}{user\_maxchunk}\DUrole{w}{ }\DUrole{p}{=}\DUrole{w}{ }\DUrole{m}{1}\sphinxparamcomma \DUrole{kt}{int}\DUrole{w}{ }\DUrole{n,sig-param}{user\_pagesize}\DUrole{w}{ }\DUrole{p}{=}\DUrole{w}{ }\DUrole{m}{1024}\sphinxparamcomma \DUrole{kt}{int}\DUrole{w}{ }\DUrole{n,sig-param}{user\_pagedelta}\DUrole{w}{ }\DUrole{p}{=}\DUrole{w}{ }\DUrole{m}{1}}{}
\pysigstopmultiline
\pysigstopsignatures
\sphinxAtStartPar
(Re\sphinxhyphen{})initialize the set of pages and allocation parameters.

\sphinxAtStartPar
This also frees all previously allocated storage and allocates the first page(s).
\begin{quote}\begin{description}
\sphinxlineitem{Parameters}\begin{itemize}
\item {} 
\sphinxAtStartPar
\sphinxstyleliteralstrong{\sphinxupquote{user\_maxchunk}} \textendash{} Expected maximum number of items for one chunk 

\item {} 
\sphinxAtStartPar
\sphinxstyleliteralstrong{\sphinxupquote{user\_pagesize}} \textendash{} Number of items on a single memory page 

\item {} 
\sphinxAtStartPar
\sphinxstyleliteralstrong{\sphinxupquote{user\_pagedelta}} \textendash{} Number of pages to allocate with one malloc 

\end{itemize}

\sphinxlineitem{Returns}
\sphinxAtStartPar
1 if there were invalid parameters, 2 if there was an allocation error or 0 if successful 

\end{description}\end{quote}

\end{fulllineitems}

\index{MyPage::get (C++ function)@\spxentry{MyPage::get}\spxextra{C++ function}}

\begin{fulllineitems}
\phantomsection\label{\detokenize{Developer_utils:_CPPv4N9LAMMPS_NS6MyPage3getEi}}\phantomsection\label{\detokenize{Developer_utils:_CPPv3N9LAMMPS_NS6MyPage3getEi}}\phantomsection\label{\detokenize{Developer_utils:_CPPv2N9LAMMPS_NS6MyPage3getEi}}\phantomsection\label{\detokenize{Developer_utils:LAMMPS_NS::MyPage::get__i}}
\pysigstartsignatures
\pysigstartmultiline
\pysiglinewithargsret{\phantomsection\label{\detokenize{Developer_utils:classLAMMPS__NS_1_1MyPage_1ac0fe0e7127a43512b203e64202414cab}}{\hyperref[\detokenize{Developer_utils:_CPPv4I0EN9LAMMPS_NS6MyPageE}]{\sphinxcrossref{\DUrole{n}{T}}}}\DUrole{w}{ }\DUrole{p}{*}\sphinxbfcode{\sphinxupquote{\DUrole{n}{get}}}}{\DUrole{kt}{int}\DUrole{w}{ }\DUrole{n,sig-param}{n}\DUrole{w}{ }\DUrole{p}{=}\DUrole{w}{ }\DUrole{m}{1}}{}
\pysigstopmultiline
\pysigstopsignatures
\sphinxAtStartPar
Pointer to location that can store N items.

\sphinxAtStartPar
This will allocate more pages as needed. If the parameter \sphinxstyleemphasis{N} is larger than the \sphinxstyleemphasis{maxchunk} setting, an error is flagged.
\begin{quote}\begin{description}
\sphinxlineitem{Parameters}
\sphinxAtStartPar
\sphinxstyleliteralstrong{\sphinxupquote{n}} \textendash{} number of items for which storage is requested 

\sphinxlineitem{Returns}
\sphinxAtStartPar
memory location or null pointer, if error or allocation failed 

\end{description}\end{quote}

\end{fulllineitems}

\index{MyPage::vget (C++ function)@\spxentry{MyPage::vget}\spxextra{C++ function}}

\begin{fulllineitems}
\phantomsection\label{\detokenize{Developer_utils:_CPPv4N9LAMMPS_NS6MyPage4vgetEv}}\phantomsection\label{\detokenize{Developer_utils:_CPPv3N9LAMMPS_NS6MyPage4vgetEv}}\phantomsection\label{\detokenize{Developer_utils:_CPPv2N9LAMMPS_NS6MyPage4vgetEv}}\phantomsection\label{\detokenize{Developer_utils:LAMMPS_NS::MyPage::vget}}
\pysigstartsignatures
\pysigstartmultiline
\pysiglinewithargsret{\phantomsection\label{\detokenize{Developer_utils:classLAMMPS__NS_1_1MyPage_1a9d22e2eb992cdc9902cf841fce95564c}}\DUrole{k}{inline}\DUrole{w}{ }{\hyperref[\detokenize{Developer_utils:_CPPv4I0EN9LAMMPS_NS6MyPageE}]{\sphinxcrossref{\DUrole{n}{T}}}}\DUrole{w}{ }\DUrole{p}{*}\sphinxbfcode{\sphinxupquote{\DUrole{n}{vget}}}}{}{}
\pysigstopmultiline
\pysigstopsignatures
\sphinxAtStartPar
Get pointer to location that can store \sphinxstyleemphasis{maxchunk} items.

\sphinxAtStartPar
This will return the same pointer as the previous call to this function unless {\hyperref[\detokenize{Developer_utils:classLAMMPS__NS_1_1MyPage_1a8b3b7dc16eae4b0445e534a37eb85806}]{\sphinxcrossref{\DUrole{std,std-ref}{vgot()}}}} is called afterwards to record how many items of the chunk were actually used.
\begin{quote}\begin{description}
\sphinxlineitem{Returns}
\sphinxAtStartPar
pointer to chunk of memory or null pointer if run out of memory 

\end{description}\end{quote}

\end{fulllineitems}

\index{MyPage::vgot (C++ function)@\spxentry{MyPage::vgot}\spxextra{C++ function}}

\begin{fulllineitems}
\phantomsection\label{\detokenize{Developer_utils:_CPPv4N9LAMMPS_NS6MyPage4vgotEi}}\phantomsection\label{\detokenize{Developer_utils:_CPPv3N9LAMMPS_NS6MyPage4vgotEi}}\phantomsection\label{\detokenize{Developer_utils:_CPPv2N9LAMMPS_NS6MyPage4vgotEi}}\phantomsection\label{\detokenize{Developer_utils:LAMMPS_NS::MyPage::vgot__i}}
\pysigstartsignatures
\pysigstartmultiline
\pysiglinewithargsret{\phantomsection\label{\detokenize{Developer_utils:classLAMMPS__NS_1_1MyPage_1a8b3b7dc16eae4b0445e534a37eb85806}}\DUrole{k}{inline}\DUrole{w}{ }\DUrole{kt}{void}\DUrole{w}{ }\sphinxbfcode{\sphinxupquote{\DUrole{n}{vgot}}}}{\DUrole{kt}{int}\DUrole{w}{ }\DUrole{n,sig-param}{n}}{}
\pysigstopmultiline
\pysigstopsignatures
\sphinxAtStartPar
Mark \sphinxstyleemphasis{N} items as used of the chunk reserved with a preceding call to {\hyperref[\detokenize{Developer_utils:classLAMMPS__NS_1_1MyPage_1a9d22e2eb992cdc9902cf841fce95564c}]{\sphinxcrossref{\DUrole{std,std-ref}{vget()}}}}.

\sphinxAtStartPar
This will advance the internal pointer inside the current memory page. It is not necessary to call this function for \sphinxstyleemphasis{N} = 0, implying the reserved storage was not used. A following call to {\hyperref[\detokenize{Developer_utils:classLAMMPS__NS_1_1MyPage_1a9d22e2eb992cdc9902cf841fce95564c}]{\sphinxcrossref{\DUrole{std,std-ref}{vget()}}}} will then reserve the same location again. It is an error if \sphinxstyleemphasis{N} \textgreater{} \sphinxstyleemphasis{maxchunk}.
\begin{quote}\begin{description}
\sphinxlineitem{Parameters}
\sphinxAtStartPar
\sphinxstyleliteralstrong{\sphinxupquote{n}} \textendash{} Number of items used in previously reserved chunk 

\end{description}\end{quote}

\end{fulllineitems}

\index{MyPage::reset (C++ function)@\spxentry{MyPage::reset}\spxextra{C++ function}}

\begin{fulllineitems}
\phantomsection\label{\detokenize{Developer_utils:_CPPv4N9LAMMPS_NS6MyPage5resetEv}}\phantomsection\label{\detokenize{Developer_utils:_CPPv3N9LAMMPS_NS6MyPage5resetEv}}\phantomsection\label{\detokenize{Developer_utils:_CPPv2N9LAMMPS_NS6MyPage5resetEv}}\phantomsection\label{\detokenize{Developer_utils:LAMMPS_NS::MyPage::reset}}
\pysigstartsignatures
\pysigstartmultiline
\pysiglinewithargsret{\phantomsection\label{\detokenize{Developer_utils:classLAMMPS__NS_1_1MyPage_1ad20897c5c8bd47f5d4005989bead0e55}}\DUrole{kt}{void}\DUrole{w}{ }\sphinxbfcode{\sphinxupquote{\DUrole{n}{reset}}}}{}{}
\pysigstopmultiline
\pysigstopsignatures
\sphinxAtStartPar
Reset state of memory pool without freeing any memory 

\end{fulllineitems}

\index{MyPage::size (C++ function)@\spxentry{MyPage::size}\spxextra{C++ function}}

\begin{fulllineitems}
\phantomsection\label{\detokenize{Developer_utils:_CPPv4NK9LAMMPS_NS6MyPage4sizeEv}}\phantomsection\label{\detokenize{Developer_utils:_CPPv3NK9LAMMPS_NS6MyPage4sizeEv}}\phantomsection\label{\detokenize{Developer_utils:_CPPv2NK9LAMMPS_NS6MyPage4sizeEv}}\phantomsection\label{\detokenize{Developer_utils:LAMMPS_NS::MyPage::sizeC}}
\pysigstartsignatures
\pysigstartmultiline
\pysiglinewithargsret{\phantomsection\label{\detokenize{Developer_utils:classLAMMPS__NS_1_1MyPage_1ac3fe9e82d5f4f1d55a5c104a2a5d37a6}}\DUrole{k}{inline}\DUrole{w}{ }\DUrole{kt}{double}\DUrole{w}{ }\sphinxbfcode{\sphinxupquote{\DUrole{n}{size}}}}{}{\DUrole{w}{ }\DUrole{k}{const}}
\pysigstopmultiline
\pysigstopsignatures
\sphinxAtStartPar
Return total size of allocated pages
\begin{quote}\begin{description}
\sphinxlineitem{Returns}
\sphinxAtStartPar
total storage used in bytes 

\end{description}\end{quote}

\end{fulllineitems}

\index{MyPage::status (C++ function)@\spxentry{MyPage::status}\spxextra{C++ function}}

\begin{fulllineitems}
\phantomsection\label{\detokenize{Developer_utils:_CPPv4NK9LAMMPS_NS6MyPage6statusEv}}\phantomsection\label{\detokenize{Developer_utils:_CPPv3NK9LAMMPS_NS6MyPage6statusEv}}\phantomsection\label{\detokenize{Developer_utils:_CPPv2NK9LAMMPS_NS6MyPage6statusEv}}\phantomsection\label{\detokenize{Developer_utils:LAMMPS_NS::MyPage::statusC}}
\pysigstartsignatures
\pysigstartmultiline
\pysiglinewithargsret{\phantomsection\label{\detokenize{Developer_utils:classLAMMPS__NS_1_1MyPage_1a92d2adf8350ab6a66718eab4a990dffb}}\DUrole{k}{inline}\DUrole{w}{ }\DUrole{kt}{int}\DUrole{w}{ }\sphinxbfcode{\sphinxupquote{\DUrole{n}{status}}}}{}{\DUrole{w}{ }\DUrole{k}{const}}
\pysigstopmultiline
\pysigstopsignatures
\sphinxAtStartPar
Return error status
\begin{quote}\begin{description}
\sphinxlineitem{Returns}
\sphinxAtStartPar
0 if no error, 1 requested chunk size \textgreater{} maxchunk, 2 if malloc failed 

\end{description}\end{quote}

\end{fulllineitems}


\end{sphinxuseclass}
\end{fulllineitems}

\index{MyPoolChunk (C++ class)@\spxentry{MyPoolChunk}\spxextra{C++ class}}

\begin{fulllineitems}
\phantomsection\label{\detokenize{Developer_utils:_CPPv4I0EN9LAMMPS_NS11MyPoolChunkE}}\phantomsection\label{\detokenize{Developer_utils:_CPPv3I0EN9LAMMPS_NS11MyPoolChunkE}}\phantomsection\label{\detokenize{Developer_utils:_CPPv2I0EN9LAMMPS_NS11MyPoolChunkE}}
\pysigstartsignatures
\pysigstartmultiline
\pysigline{\DUrole{k}{template}\DUrole{p}{\textless{}}\DUrole{k}{class}\DUrole{w}{ }\sphinxbfcode{\sphinxupquote{\DUrole{n}{T}}}\DUrole{p}{\textgreater{}}}
\pysigline{\phantomsection\label{\detokenize{Developer_utils:classLAMMPS__NS_1_1MyPoolChunk}}\DUrole{k}{class}\DUrole{w}{ }\sphinxbfcode{\sphinxupquote{\DUrole{n}{MyPoolChunk}}}}
\pysigstopmultiline
\pysigstopsignatures
\sphinxAtStartPar
Templated class for storing chunks of datums in pages. 

\sphinxAtStartPar
The size of the chunk may vary from call to call between the \sphinxstyleemphasis{minchunk} and \sphinxstyleemphasis{maxchunk} setting. Chunks may be returned to the pool for re\sphinxhyphen{}use. Chunks can be reserved in \sphinxstyleemphasis{nbin} different sizes between \sphinxstyleemphasis{minchunk} and \sphinxstyleemphasis{maxchunk}. The \sphinxstyleemphasis{chunksperpage} setting specifies how many chunks are stored on any page and the \sphinxstyleemphasis{pagedelta} setting determines how many pages are allocated in one go. Pages are never freed, so they can be re\sphinxhyphen{}used without re\sphinxhyphen{}allocation.

\begin{sphinxadmonition}{note}{Note:}
\sphinxAtStartPar
This is a template class with explicit instantiation. If the class is used with a new data type, a new explicit instantiation may need to be added at the end of the file \sphinxcode{\sphinxupquote{src/my\_pool\_chunk.cpp}} to avoid symbol lookup errors. 
\end{sphinxadmonition}

\begin{sphinxuseclass}{breathe-sectiondef}\subsubsection*{Public Functions}
\index{MyPoolChunk::MyPoolChunk (C++ function)@\spxentry{MyPoolChunk::MyPoolChunk}\spxextra{C++ function}}

\begin{fulllineitems}
\phantomsection\label{\detokenize{Developer_utils:_CPPv4N9LAMMPS_NS11MyPoolChunk11MyPoolChunkEiiiii}}\phantomsection\label{\detokenize{Developer_utils:_CPPv3N9LAMMPS_NS11MyPoolChunk11MyPoolChunkEiiiii}}\phantomsection\label{\detokenize{Developer_utils:_CPPv2N9LAMMPS_NS11MyPoolChunk11MyPoolChunkEiiiii}}\phantomsection\label{\detokenize{Developer_utils:LAMMPS_NS::MyPoolChunk::MyPoolChunk__i.i.i.i.i}}
\pysigstartsignatures
\pysigstartmultiline
\pysiglinewithargsret{\phantomsection\label{\detokenize{Developer_utils:classLAMMPS__NS_1_1MyPoolChunk_1a60daa98fd41e5ff09753fe55f044b453}}\sphinxbfcode{\sphinxupquote{\DUrole{n}{MyPoolChunk}}}}{\DUrole{kt}{int}\DUrole{w}{ }\DUrole{n,sig-param}{user\_minchunk}\DUrole{w}{ }\DUrole{p}{=}\DUrole{w}{ }\DUrole{m}{1}\sphinxparamcomma \DUrole{kt}{int}\DUrole{w}{ }\DUrole{n,sig-param}{user\_maxchunk}\DUrole{w}{ }\DUrole{p}{=}\DUrole{w}{ }\DUrole{m}{1}\sphinxparamcomma \DUrole{kt}{int}\DUrole{w}{ }\DUrole{n,sig-param}{user\_nbin}\DUrole{w}{ }\DUrole{p}{=}\DUrole{w}{ }\DUrole{m}{1}\sphinxparamcomma \DUrole{kt}{int}\DUrole{w}{ }\DUrole{n,sig-param}{user\_chunkperpage}\DUrole{w}{ }\DUrole{p}{=}\DUrole{w}{ }\DUrole{m}{1024}\sphinxparamcomma \DUrole{kt}{int}\DUrole{w}{ }\DUrole{n,sig-param}{user\_pagedelta}\DUrole{w}{ }\DUrole{p}{=}\DUrole{w}{ }\DUrole{m}{1}}{}
\pysigstopmultiline
\pysigstopsignatures
\sphinxAtStartPar
Create a class instance and set memory pool parameters
\begin{quote}\begin{description}
\sphinxlineitem{Parameters}\begin{itemize}
\item {} 
\sphinxAtStartPar
\sphinxstyleliteralstrong{\sphinxupquote{user\_minchunk}} \textendash{} Minimal chunk size 

\item {} 
\sphinxAtStartPar
\sphinxstyleliteralstrong{\sphinxupquote{user\_maxchunk}} \textendash{} Maximal chunk size 

\item {} 
\sphinxAtStartPar
\sphinxstyleliteralstrong{\sphinxupquote{user\_nbin}} \textendash{} Number of bins of different chunk sizes 

\item {} 
\sphinxAtStartPar
\sphinxstyleliteralstrong{\sphinxupquote{user\_chunkperpage}} \textendash{} Number of chunks per page 

\item {} 
\sphinxAtStartPar
\sphinxstyleliteralstrong{\sphinxupquote{user\_pagedelta}} \textendash{} Number of pages to allocate in one go 

\end{itemize}

\end{description}\end{quote}

\end{fulllineitems}

\index{MyPoolChunk::\textasciitilde{}MyPoolChunk (C++ function)@\spxentry{MyPoolChunk::\textasciitilde{}MyPoolChunk}\spxextra{C++ function}}

\begin{fulllineitems}
\phantomsection\label{\detokenize{Developer_utils:_CPPv4N9LAMMPS_NS11MyPoolChunkD0Ev}}\phantomsection\label{\detokenize{Developer_utils:_CPPv3N9LAMMPS_NS11MyPoolChunkD0Ev}}\phantomsection\label{\detokenize{Developer_utils:_CPPv2N9LAMMPS_NS11MyPoolChunkD0Ev}}\phantomsection\label{\detokenize{Developer_utils:LAMMPS_NS::MyPoolChunk::_MyPoolChunk}}
\pysigstartsignatures
\pysigstartmultiline
\pysiglinewithargsret{\phantomsection\label{\detokenize{Developer_utils:classLAMMPS__NS_1_1MyPoolChunk_1aa185474081f1f923a35914ed67077e11}}\sphinxbfcode{\sphinxupquote{\DUrole{n}{\textasciitilde{}MyPoolChunk}}}}{}{}
\pysigstopmultiline
\pysigstopsignatures
\sphinxAtStartPar
Destroy class instance and free all allocated memory 

\end{fulllineitems}

\index{MyPoolChunk::get (C++ function)@\spxentry{MyPoolChunk::get}\spxextra{C++ function}}

\begin{fulllineitems}
\phantomsection\label{\detokenize{Developer_utils:_CPPv4N9LAMMPS_NS11MyPoolChunk3getERi}}\phantomsection\label{\detokenize{Developer_utils:_CPPv3N9LAMMPS_NS11MyPoolChunk3getERi}}\phantomsection\label{\detokenize{Developer_utils:_CPPv2N9LAMMPS_NS11MyPoolChunk3getERi}}\phantomsection\label{\detokenize{Developer_utils:LAMMPS_NS::MyPoolChunk::get__iR}}
\pysigstartsignatures
\pysigstartmultiline
\pysiglinewithargsret{\phantomsection\label{\detokenize{Developer_utils:classLAMMPS__NS_1_1MyPoolChunk_1a591380d4a3a168110139f35e1d56a3c9}}{\hyperref[\detokenize{Developer_utils:_CPPv4I0EN9LAMMPS_NS11MyPoolChunkE}]{\sphinxcrossref{\DUrole{n}{T}}}}\DUrole{w}{ }\DUrole{p}{*}\sphinxbfcode{\sphinxupquote{\DUrole{n}{get}}}}{\DUrole{kt}{int}\DUrole{w}{ }\DUrole{p}{\&}\DUrole{n,sig-param}{index}}{}
\pysigstopmultiline
\pysigstopsignatures
\sphinxAtStartPar
Return pointer/index of unused chunk of size maxchunk
\begin{quote}\begin{description}
\sphinxlineitem{Parameters}
\sphinxAtStartPar
\sphinxstyleliteralstrong{\sphinxupquote{index}} \textendash{} Index of chunk in memory pool 

\sphinxlineitem{Returns}
\sphinxAtStartPar
Pointer to requested chunk of storage 

\end{description}\end{quote}

\end{fulllineitems}

\index{MyPoolChunk::get (C++ function)@\spxentry{MyPoolChunk::get}\spxextra{C++ function}}

\begin{fulllineitems}
\phantomsection\label{\detokenize{Developer_utils:_CPPv4N9LAMMPS_NS11MyPoolChunk3getEiRi}}\phantomsection\label{\detokenize{Developer_utils:_CPPv3N9LAMMPS_NS11MyPoolChunk3getEiRi}}\phantomsection\label{\detokenize{Developer_utils:_CPPv2N9LAMMPS_NS11MyPoolChunk3getEiRi}}\phantomsection\label{\detokenize{Developer_utils:LAMMPS_NS::MyPoolChunk::get__i.iR}}
\pysigstartsignatures
\pysigstartmultiline
\pysiglinewithargsret{\phantomsection\label{\detokenize{Developer_utils:classLAMMPS__NS_1_1MyPoolChunk_1a781c43c8f6e18dcb99c613292a38cc4e}}{\hyperref[\detokenize{Developer_utils:_CPPv4I0EN9LAMMPS_NS11MyPoolChunkE}]{\sphinxcrossref{\DUrole{n}{T}}}}\DUrole{w}{ }\DUrole{p}{*}\sphinxbfcode{\sphinxupquote{\DUrole{n}{get}}}}{\DUrole{kt}{int}\DUrole{w}{ }\DUrole{n,sig-param}{n}\sphinxparamcomma \DUrole{kt}{int}\DUrole{w}{ }\DUrole{p}{\&}\DUrole{n,sig-param}{index}}{}
\pysigstopmultiline
\pysigstopsignatures
\sphinxAtStartPar
Return pointer/index of unused chunk of size N
\begin{quote}\begin{description}
\sphinxlineitem{Parameters}\begin{itemize}
\item {} 
\sphinxAtStartPar
\sphinxstyleliteralstrong{\sphinxupquote{n}} \textendash{} Size of chunk 

\item {} 
\sphinxAtStartPar
\sphinxstyleliteralstrong{\sphinxupquote{index}} \textendash{} Index of chunk in memory pool 

\end{itemize}

\sphinxlineitem{Returns}
\sphinxAtStartPar
Pointer to requested chunk of storage 

\end{description}\end{quote}

\end{fulllineitems}

\index{MyPoolChunk::put (C++ function)@\spxentry{MyPoolChunk::put}\spxextra{C++ function}}

\begin{fulllineitems}
\phantomsection\label{\detokenize{Developer_utils:_CPPv4N9LAMMPS_NS11MyPoolChunk3putEi}}\phantomsection\label{\detokenize{Developer_utils:_CPPv3N9LAMMPS_NS11MyPoolChunk3putEi}}\phantomsection\label{\detokenize{Developer_utils:_CPPv2N9LAMMPS_NS11MyPoolChunk3putEi}}\phantomsection\label{\detokenize{Developer_utils:LAMMPS_NS::MyPoolChunk::put__i}}
\pysigstartsignatures
\pysigstartmultiline
\pysiglinewithargsret{\phantomsection\label{\detokenize{Developer_utils:classLAMMPS__NS_1_1MyPoolChunk_1acf93b81dd253b533195c20ee85621412}}\DUrole{kt}{void}\DUrole{w}{ }\sphinxbfcode{\sphinxupquote{\DUrole{n}{put}}}}{\DUrole{kt}{int}\DUrole{w}{ }\DUrole{n,sig-param}{index}}{}
\pysigstopmultiline
\pysigstopsignatures
\sphinxAtStartPar
Put indexed chunk back into memory pool via free list
\begin{quote}\begin{description}
\sphinxlineitem{Parameters}
\sphinxAtStartPar
\sphinxstyleliteralstrong{\sphinxupquote{index}} \textendash{} Memory chunk index returned by call to {\hyperref[\detokenize{Developer_utils:classLAMMPS__NS_1_1MyPoolChunk_1a591380d4a3a168110139f35e1d56a3c9}]{\sphinxcrossref{\DUrole{std,std-ref}{get()}}}}

\end{description}\end{quote}

\end{fulllineitems}

\index{MyPoolChunk::size (C++ function)@\spxentry{MyPoolChunk::size}\spxextra{C++ function}}

\begin{fulllineitems}
\phantomsection\label{\detokenize{Developer_utils:_CPPv4NK9LAMMPS_NS11MyPoolChunk4sizeEv}}\phantomsection\label{\detokenize{Developer_utils:_CPPv3NK9LAMMPS_NS11MyPoolChunk4sizeEv}}\phantomsection\label{\detokenize{Developer_utils:_CPPv2NK9LAMMPS_NS11MyPoolChunk4sizeEv}}\phantomsection\label{\detokenize{Developer_utils:LAMMPS_NS::MyPoolChunk::sizeC}}
\pysigstartsignatures
\pysigstartmultiline
\pysiglinewithargsret{\phantomsection\label{\detokenize{Developer_utils:classLAMMPS__NS_1_1MyPoolChunk_1ac3fe9e82d5f4f1d55a5c104a2a5d37a6}}\DUrole{kt}{double}\DUrole{w}{ }\sphinxbfcode{\sphinxupquote{\DUrole{n}{size}}}}{}{\DUrole{w}{ }\DUrole{k}{const}}
\pysigstopmultiline
\pysigstopsignatures
\sphinxAtStartPar
Return total size of allocated pages
\begin{quote}\begin{description}
\sphinxlineitem{Returns}
\sphinxAtStartPar
total storage used in bytes 

\end{description}\end{quote}

\end{fulllineitems}

\index{MyPoolChunk::status (C++ function)@\spxentry{MyPoolChunk::status}\spxextra{C++ function}}

\begin{fulllineitems}
\phantomsection\label{\detokenize{Developer_utils:_CPPv4NK9LAMMPS_NS11MyPoolChunk6statusEv}}\phantomsection\label{\detokenize{Developer_utils:_CPPv3NK9LAMMPS_NS11MyPoolChunk6statusEv}}\phantomsection\label{\detokenize{Developer_utils:_CPPv2NK9LAMMPS_NS11MyPoolChunk6statusEv}}\phantomsection\label{\detokenize{Developer_utils:LAMMPS_NS::MyPoolChunk::statusC}}
\pysigstartsignatures
\pysigstartmultiline
\pysiglinewithargsret{\phantomsection\label{\detokenize{Developer_utils:classLAMMPS__NS_1_1MyPoolChunk_1a92d2adf8350ab6a66718eab4a990dffb}}\DUrole{k}{inline}\DUrole{w}{ }\DUrole{kt}{int}\DUrole{w}{ }\sphinxbfcode{\sphinxupquote{\DUrole{n}{status}}}}{}{\DUrole{w}{ }\DUrole{k}{const}}
\pysigstopmultiline
\pysigstopsignatures
\sphinxAtStartPar
Return error status
\begin{quote}\begin{description}
\sphinxlineitem{Returns}
\sphinxAtStartPar
0 if no error, 1 if invalid input, 2 if malloc() failed, 3 if chunk \textgreater{} maxchunk 

\end{description}\end{quote}

\end{fulllineitems}


\end{sphinxuseclass}
\end{fulllineitems}



\bigskip\hrule\bigskip



\section{Eigensolver functions}
\label{\detokenize{Developer_utils:eigensolver-functions}}
\sphinxAtStartPar
The \sphinxcode{\sphinxupquote{MathEigen}} sub\sphinxhyphen{}namespace of the \sphinxcode{\sphinxupquote{LAMMPS\_NS}} namespace contains
functions and classes for eigensolvers. Currently only the
{\hyperref[\detokenize{Developer_utils:_CPPv4N9MathEigen7jacobi3EAL3E_AL3E_KdPdAL3E_AL3E_di}]{\sphinxcrossref{\sphinxcode{\sphinxupquote{jacobi3 function}}}}} is used in various
places in LAMMPS.  That function is built on top of a group of more
generic eigensolvers that are maintained in the \sphinxcode{\sphinxupquote{math\_eigen\_impl.h}}
header file.  This header contains the implementation of three template
classes:
\begin{enumerate}
\sphinxsetlistlabels{\arabic}{enumi}{enumii}{}{.}%
\item {} 
\sphinxAtStartPar
“Jacobi” calculates all of the eigenvalues and eigenvectors
of a dense, symmetric, real matrix.

\item {} 
\sphinxAtStartPar
The “PEigenDense” class only calculates the principal eigenvalue
(i.e. the largest or smallest eigenvalue), and its corresponding
eigenvector.  However it is much more efficient than “Jacobi” when
applied to large matrices (larger than 13x13).  PEigenDense also can
understand complex\sphinxhyphen{}valued Hermitian matrices.

\item {} 
\sphinxAtStartPar
The “LambdaLanczos” class is a generalization of “PEigenDense” which can be
applied to arbitrary sparse matrices.

\end{enumerate}

\sphinxAtStartPar
The “math\_eigen\_impl.h” code is an amalgamation of \sphinxhref{https://github.com/jewettaij/jacobi\_pd}{jacobi\_pd} by Andrew Jewett at Scripps
Research (under CC0\sphinxhyphen{}1.0 license) and \sphinxhref{https://github.com/mrcdr/lambda-lanczos}{Lambda Lanczos} by Yuya Kurebayashi at
Tohoku University (under MIT license)


\bigskip\hrule\bigskip

\index{MathEigen::jacobi3 (C++ function)@\spxentry{MathEigen::jacobi3}\spxextra{C++ function}}

\begin{fulllineitems}
\phantomsection\label{\detokenize{Developer_utils:_CPPv4N9MathEigen7jacobi3EPPCKdPdPPdi}}\phantomsection\label{\detokenize{Developer_utils:_CPPv3N9MathEigen7jacobi3EPPCKdPdPPdi}}\phantomsection\label{\detokenize{Developer_utils:_CPPv2N9MathEigen7jacobi3EPPCKdPdPPdi}}\phantomsection\label{\detokenize{Developer_utils:MathEigen::jacobi3__doubleCPCP.doubleP.doublePP.i}}
\pysigstartsignatures
\pysigstartmultiline
\pysiglinewithargsret{\phantomsection\label{\detokenize{Developer_utils:math__eigen_8h_1ab5f679ef0f37f9d96beb719cd8594dfe}}\DUrole{kt}{int}\DUrole{w}{ }\sphinxcode{\sphinxupquote{\DUrole{n}{MathEigen}\DUrole{p}{::}}}\sphinxbfcode{\sphinxupquote{\DUrole{n}{jacobi3}}}}{\DUrole{kt}{double}\DUrole{w}{ }\DUrole{k}{const}\DUrole{w}{ }\DUrole{p}{*}\DUrole{k}{const}\DUrole{w}{ }\DUrole{p}{*}\DUrole{n,sig-param}{mat}\sphinxparamcomma \DUrole{kt}{double}\DUrole{w}{ }\DUrole{p}{*}\DUrole{n,sig-param}{eval}\sphinxparamcomma \DUrole{kt}{double}\DUrole{w}{ }\DUrole{p}{*}\DUrole{p}{*}\DUrole{n,sig-param}{evec}\sphinxparamcomma \DUrole{kt}{int}\DUrole{w}{ }\DUrole{n,sig-param}{sort}\DUrole{w}{ }\DUrole{p}{=}\DUrole{w}{ }\DUrole{o}{\sphinxhyphen{}}\DUrole{m}{1}}{}
\pysigstopmultiline
\pysigstopsignatures
\sphinxAtStartPar
A specialized function which finds the eigenvalues and eigenvectors of a 3x3 matrix (in double ** format).
\begin{quote}\begin{description}
\sphinxlineitem{Parameters}\begin{itemize}
\item {} 
\sphinxAtStartPar
\sphinxstyleliteralstrong{\sphinxupquote{mat}} \textendash{} the 3x3 matrix you wish to diagonalize 

\item {} 
\sphinxAtStartPar
\sphinxstyleliteralstrong{\sphinxupquote{eval}} \textendash{} store the eigenvalues here 

\item {} 
\sphinxAtStartPar
\sphinxstyleliteralstrong{\sphinxupquote{evec}} \textendash{} store the eigenvectors here… 

\item {} 
\sphinxAtStartPar
\sphinxstyleliteralstrong{\sphinxupquote{sort}} \textendash{} order eigenvalues and \sphinxhyphen{}vectors (\sphinxhyphen{}1 decreasing (default), 1 increasing, 0 unsorted) 

\end{itemize}

\sphinxlineitem{Returns}
\sphinxAtStartPar
0 if eigenvalue calculation converged, 1 if it failed 

\end{description}\end{quote}

\end{fulllineitems}

\index{MathEigen::jacobi3 (C++ function)@\spxentry{MathEigen::jacobi3}\spxextra{C++ function}}

\begin{fulllineitems}
\phantomsection\label{\detokenize{Developer_utils:_CPPv4N9MathEigen7jacobi3EAL3E_AL3E_KdPdAL3E_AL3E_di}}\phantomsection\label{\detokenize{Developer_utils:_CPPv3N9MathEigen7jacobi3EAL3E_AL3E_KdPdAL3E_AL3E_di}}\phantomsection\label{\detokenize{Developer_utils:_CPPv2N9MathEigen7jacobi3EA3_A3_KdPdA3_A3_di}}\phantomsection\label{\detokenize{Developer_utils:MathEigen::jacobi3__doubleCAA.doubleP.doubleAA.i}}
\pysigstartsignatures
\pysigstartmultiline
\pysiglinewithargsret{\phantomsection\label{\detokenize{Developer_utils:math__eigen_8h_1a00305ba13433b680f4c3e6bf2459a5b0}}\DUrole{kt}{int}\DUrole{w}{ }\sphinxcode{\sphinxupquote{\DUrole{n}{MathEigen}\DUrole{p}{::}}}\sphinxbfcode{\sphinxupquote{\DUrole{n}{jacobi3}}}}{\DUrole{kt}{double}\DUrole{w}{ }\DUrole{k}{const}\DUrole{w}{ }\DUrole{n,sig-param}{mat}\DUrole{p}{{[}}\DUrole{m}{3}\DUrole{p}{{]}}\DUrole{p}{{[}}\DUrole{m}{3}\DUrole{p}{{]}}\sphinxparamcomma \DUrole{kt}{double}\DUrole{w}{ }\DUrole{p}{*}\DUrole{n,sig-param}{eval}\sphinxparamcomma \DUrole{kt}{double}\DUrole{w}{ }\DUrole{n,sig-param}{evec}\DUrole{p}{{[}}\DUrole{m}{3}\DUrole{p}{{]}}\DUrole{p}{{[}}\DUrole{m}{3}\DUrole{p}{{]}}\sphinxparamcomma \DUrole{kt}{int}\DUrole{w}{ }\DUrole{n,sig-param}{sort}\DUrole{w}{ }\DUrole{p}{=}\DUrole{w}{ }\DUrole{o}{\sphinxhyphen{}}\DUrole{m}{1}}{}
\pysigstopmultiline
\pysigstopsignatures
\sphinxAtStartPar
This is an overloaded member function, provided for convenience. It differs from the above function only in what argument(s) it accepts. 

\end{fulllineitems}



\bigskip\hrule\bigskip



\section{Communication buffer coding with \sphinxstyleemphasis{ubuf}}
\label{\detokenize{Developer_utils:communication-buffer-coding-with-ubuf}}\label{\detokenize{Developer_utils:id2}}
\sphinxAtStartPar
LAMMPS uses communication buffers where it collects data from various
class instances and then exchanges the data with neighboring subdomains.
For simplicity those buffers are defined as \sphinxcode{\sphinxupquote{double}} buffers and
used for doubles and integer numbers. This presents a unique problem
when 64\sphinxhyphen{}bit integers are used.  While the storage needed for a \sphinxcode{\sphinxupquote{double}}
is also 64\sphinxhyphen{}bit, it cannot be used by a simple assignment.  To get around
that limitation, LAMMPS uses the {\hyperref[\detokenize{Developer_utils:_CPPv4N9LAMMPS_NS4ubufE}]{\sphinxcrossref{\sphinxcode{\sphinxupquote{ubuf}}}}}
union.  It is used in the various “pack” and “unpack” functions in the
LAMMPS classes to store and retrieve integers that may be 64\sphinxhyphen{}bit from
the communication buffers.


\bigskip\hrule\bigskip

\index{ubuf (C++ union)@\spxentry{ubuf}\spxextra{C++ union}}

\begin{fulllineitems}
\phantomsection\label{\detokenize{Developer_utils:_CPPv4N9LAMMPS_NS4ubufE}}\phantomsection\label{\detokenize{Developer_utils:_CPPv3N9LAMMPS_NS4ubufE}}\phantomsection\label{\detokenize{Developer_utils:_CPPv2N9LAMMPS_NS4ubufE}}
\pysigstartsignatures
\pysigstartmultiline
\pysigline{\phantomsection\label{\detokenize{Developer_utils:unionLAMMPS__NS_1_1ubuf}}\DUrole{k}{union}\DUrole{w}{ }\sphinxbfcode{\sphinxupquote{\DUrole{n}{ubuf}}}}
\pysigstopmultiline
\pysigstopsignatures\sphinxstyleemphasis{\#include \textless{}lmptype.h\textgreater{}}
\sphinxAtStartPar
Data structure for packing 32\sphinxhyphen{}bit and 64\sphinxhyphen{}bit integers into double (communication) buffers

\sphinxAtStartPar
Using this union avoids aliasing issues by having member types (double, int) referencing the same buffer memory location.

\sphinxAtStartPar
The explicit constructor for 32\sphinxhyphen{}bit integers prevents compilers from (incorrectly) calling the double constructor when storing an int into a double buffer. 
\sphinxAtStartPar

\sphinxAtStartPar
\sphinxstylestrong{Usage:}
\sphinxSetupCaptionForVerbatim{To copy an integer into a double buffer:}
\def\sphinxLiteralBlockLabel{\label{\detokenize{Developer_utils:id8}}}
\begin{sphinxVerbatim}[commandchars=\\\{\}]
\PYG{k+kt}{double}\PYG{+w}{ }\PYG{n}{buf}\PYG{p}{[}\PYG{l+m+mi}{2}\PYG{p}{]}\PYG{p}{;}
\PYG{k+kt}{int}\PYG{+w}{    }\PYG{n}{foo}\PYG{+w}{ }\PYG{o}{=}\PYG{+w}{   }\PYG{l+m+mi}{1}\PYG{p}{;}
\PYG{n}{tagint}\PYG{+w}{ }\PYG{n}{bar}\PYG{+w}{ }\PYG{o}{=}\PYG{+w}{ }\PYG{l+m+mi}{2}\PYG{o}{\PYGZlt{}}\PYG{o}{\PYGZlt{}}\PYG{l+m+mi}{40}\PYG{p}{;}
\PYG{n}{buf}\PYG{p}{[}\PYG{l+m+mi}{1}\PYG{p}{]}\PYG{+w}{ }\PYG{o}{=}\PYG{+w}{ }\PYG{n}{ubuf}\PYG{p}{(}\PYG{n}{foo}\PYG{p}{)}\PYG{p}{.}\PYG{n}{d}\PYG{p}{;}
\PYG{n}{buf}\PYG{p}{[}\PYG{l+m+mi}{2}\PYG{p}{]}\PYG{+w}{ }\PYG{o}{=}\PYG{+w}{ }\PYG{n}{ubuf}\PYG{p}{(}\PYG{n}{bar}\PYG{p}{)}\PYG{p}{.}\PYG{n}{d}\PYG{p}{;}
\end{sphinxVerbatim}
\sphinxSetupCaptionForVerbatim{To copy from a double buffer back to an int:}
\def\sphinxLiteralBlockLabel{\label{\detokenize{Developer_utils:id9}}}
\begin{sphinxVerbatim}[commandchars=\\\{\}]
\PYG{n}{foo}\PYG{+w}{ }\PYG{o}{=}\PYG{+w}{ }\PYG{p}{(}\PYG{k+kt}{int}\PYG{p}{)}\PYG{+w}{    }\PYG{n}{ubuf}\PYG{p}{(}\PYG{n}{buf}\PYG{p}{[}\PYG{l+m+mi}{1}\PYG{p}{]}\PYG{p}{)}\PYG{p}{.}\PYG{n}{i}\PYG{p}{;}
\PYG{n}{bar}\PYG{+w}{ }\PYG{o}{=}\PYG{+w}{ }\PYG{p}{(}\PYG{n}{tagint}\PYG{p}{)}\PYG{+w}{ }\PYG{n}{ubuf}\PYG{p}{(}\PYG{n}{buf}\PYG{p}{[}\PYG{l+m+mi}{2}\PYG{p}{]}\PYG{p}{)}\PYG{p}{.}\PYG{n}{i}\PYG{p}{;}
\end{sphinxVerbatim}

\sphinxAtStartPar
The typecasts prevent compiler warnings about possible truncation issues.



\begin{sphinxuseclass}{breathe-sectiondef}\subsubsection*{Public Functions}
\index{ubuf::ubuf (C++ function)@\spxentry{ubuf::ubuf}\spxextra{C++ function}}

\begin{fulllineitems}
\phantomsection\label{\detokenize{Developer_utils:_CPPv4N9LAMMPS_NS4ubuf4ubufERKd}}\phantomsection\label{\detokenize{Developer_utils:_CPPv3N9LAMMPS_NS4ubuf4ubufERKd}}\phantomsection\label{\detokenize{Developer_utils:_CPPv2N9LAMMPS_NS4ubuf4ubufERKd}}\phantomsection\label{\detokenize{Developer_utils:LAMMPS_NS::ubuf::ubuf__doubleCR}}
\pysigstartsignatures
\pysigstartmultiline
\pysiglinewithargsret{\phantomsection\label{\detokenize{Developer_utils:unionLAMMPS__NS_1_1ubuf_1a5cdd4245af25171a133da99e9808170b}}\DUrole{k}{inline}\DUrole{w}{ }\sphinxbfcode{\sphinxupquote{\DUrole{n}{ubuf}}}}{\DUrole{k}{const}\DUrole{w}{ }\DUrole{kt}{double}\DUrole{w}{ }\DUrole{p}{\&}\DUrole{n,sig-param}{arg}}{}
\pysigstopmultiline
\pysigstopsignatures
\end{fulllineitems}

\index{ubuf::ubuf (C++ function)@\spxentry{ubuf::ubuf}\spxextra{C++ function}}

\begin{fulllineitems}
\phantomsection\label{\detokenize{Developer_utils:_CPPv4N9LAMMPS_NS4ubuf4ubufERK7int64_t}}\phantomsection\label{\detokenize{Developer_utils:_CPPv3N9LAMMPS_NS4ubuf4ubufERK7int64_t}}\phantomsection\label{\detokenize{Developer_utils:_CPPv2N9LAMMPS_NS4ubuf4ubufERK7int64_t}}\phantomsection\label{\detokenize{Developer_utils:LAMMPS_NS::ubuf::ubuf__int64_tCR}}
\pysigstartsignatures
\pysigstartmultiline
\pysiglinewithargsret{\phantomsection\label{\detokenize{Developer_utils:unionLAMMPS__NS_1_1ubuf_1a86bb79057870b1bd9b4f9275e5d929fa}}\DUrole{k}{inline}\DUrole{w}{ }\sphinxbfcode{\sphinxupquote{\DUrole{n}{ubuf}}}}{\DUrole{k}{const}\DUrole{w}{ }\DUrole{n}{int64\_t}\DUrole{w}{ }\DUrole{p}{\&}\DUrole{n,sig-param}{arg}}{}
\pysigstopmultiline
\pysigstopsignatures
\end{fulllineitems}

\index{ubuf::ubuf (C++ function)@\spxentry{ubuf::ubuf}\spxextra{C++ function}}

\begin{fulllineitems}
\phantomsection\label{\detokenize{Developer_utils:_CPPv4N9LAMMPS_NS4ubuf4ubufERKi}}\phantomsection\label{\detokenize{Developer_utils:_CPPv3N9LAMMPS_NS4ubuf4ubufERKi}}\phantomsection\label{\detokenize{Developer_utils:_CPPv2N9LAMMPS_NS4ubuf4ubufERKi}}\phantomsection\label{\detokenize{Developer_utils:LAMMPS_NS::ubuf::ubuf__iCR}}
\pysigstartsignatures
\pysigstartmultiline
\pysiglinewithargsret{\phantomsection\label{\detokenize{Developer_utils:unionLAMMPS__NS_1_1ubuf_1a76e1c4df9fb2687d7f2055f90b08360d}}\DUrole{k}{inline}\DUrole{w}{ }\sphinxbfcode{\sphinxupquote{\DUrole{n}{ubuf}}}}{\DUrole{k}{const}\DUrole{w}{ }\DUrole{kt}{int}\DUrole{w}{ }\DUrole{p}{\&}\DUrole{n,sig-param}{arg}}{}
\pysigstopmultiline
\pysigstopsignatures
\end{fulllineitems}


\end{sphinxuseclass}
\begin{sphinxuseclass}{breathe-sectiondef}\subsubsection*{Public Members}
\index{ubuf::d (C++ member)@\spxentry{ubuf::d}\spxextra{C++ member}}

\begin{fulllineitems}
\phantomsection\label{\detokenize{Developer_utils:_CPPv4N9LAMMPS_NS4ubuf1dE}}\phantomsection\label{\detokenize{Developer_utils:_CPPv3N9LAMMPS_NS4ubuf1dE}}\phantomsection\label{\detokenize{Developer_utils:_CPPv2N9LAMMPS_NS4ubuf1dE}}\phantomsection\label{\detokenize{Developer_utils:LAMMPS_NS::ubuf::d__double}}
\pysigstartsignatures
\pysigstartmultiline
\pysigline{\phantomsection\label{\detokenize{Developer_utils:unionLAMMPS__NS_1_1ubuf_1a873684cefeb665f3d5e6b495de57fc0d}}\DUrole{kt}{double}\DUrole{w}{ }\sphinxbfcode{\sphinxupquote{\DUrole{n}{d}}}}
\pysigstopmultiline
\pysigstopsignatures
\end{fulllineitems}

\index{ubuf::i (C++ member)@\spxentry{ubuf::i}\spxextra{C++ member}}

\begin{fulllineitems}
\phantomsection\label{\detokenize{Developer_utils:_CPPv4N9LAMMPS_NS4ubuf1iE}}\phantomsection\label{\detokenize{Developer_utils:_CPPv3N9LAMMPS_NS4ubuf1iE}}\phantomsection\label{\detokenize{Developer_utils:_CPPv2N9LAMMPS_NS4ubuf1iE}}\phantomsection\label{\detokenize{Developer_utils:LAMMPS_NS::ubuf::i__int64_t}}
\pysigstartsignatures
\pysigstartmultiline
\pysigline{\phantomsection\label{\detokenize{Developer_utils:unionLAMMPS__NS_1_1ubuf_1a193e7d4afbc8ac1c220aa7130e279451}}\DUrole{n}{int64\_t}\DUrole{w}{ }\sphinxbfcode{\sphinxupquote{\DUrole{n}{i}}}}
\pysigstopmultiline
\pysigstopsignatures
\end{fulllineitems}


\end{sphinxuseclass}
\end{fulllineitems}


\sphinxstepscope


\section{Use of distributed grids within style classes}
\label{\detokenize{Developer_grid:use-of-distributed-grids-within-style-classes}}\label{\detokenize{Developer_grid::doc}}
\sphinxAtStartPar
\DUrole{versionmodified,added}{Added in version 22Dec2022.}

\sphinxAtStartPar
The LAMMPS source code includes two classes which facilitate the
creation and use of distributed grids.  These are the Grid2d and
Grid3d classes in the src/grid2d.cpp.h and src/grid3d.cpp.h files
respectively.  As the names imply, they are used for 2d or 3d
simulations, as defined by the {\hyperref[\detokenize{dimension::doc}]{\sphinxcrossref{\DUrole{doc}{dimension}}}} command.

\sphinxAtStartPar
The {\hyperref[\detokenize{Howto_grid::doc}]{\sphinxcrossref{\DUrole{doc}{Howto\_grid}}}} page gives an overview of how
distributed grids are defined from a user perspective, lists LAMMPS
commands which use them, and explains how grid cell data is referenced
from an input script.  Please read that page first as it motivates the
coding details discussed here.

\sphinxAtStartPar
This doc page is for users who wish to write new styles (input script
commands) which use distributed grids.  There are a variety of
material models and analysis methods which use atoms (or
coarse\sphinxhyphen{}grained particles) and grids in tandem.

\sphinxAtStartPar
A \sphinxstyleemphasis{distributed} grid means each processor owns a subset of the grid
cells.  In LAMMPS, the subset for each processor will be a sub\sphinxhyphen{}block
of grid cells with low and high index bounds in each dimension of the
grid.  The union of the sub\sphinxhyphen{}blocks across all processors is the global
grid.

\sphinxAtStartPar
More specifically, a grid point is defined for each cell (by default
the center point), and a processor owns a grid cell if its point is
within the processor’s spatial subdomain.  The union of processor
subdomains is the global simulation box.  If a grid point is on the
boundary of two subdomains, the lower processor owns the grid cell.  A
processor may also store copies of ghost cells which surround its
owned cells.


\bigskip\hrule\bigskip



\subsection{Style commands}
\label{\detokenize{Developer_grid:style-commands}}
\sphinxAtStartPar
Style commands which can define and use distributed grids include the
{\hyperref[\detokenize{compute::doc}]{\sphinxcrossref{\DUrole{doc}{compute}}}}, {\hyperref[\detokenize{fix::doc}]{\sphinxcrossref{\DUrole{doc}{fix}}}}, {\hyperref[\detokenize{pair_style::doc}]{\sphinxcrossref{\DUrole{doc}{pair}}}},
and {\hyperref[\detokenize{kspace_style::doc}]{\sphinxcrossref{\DUrole{doc}{kspace}}}} styles.  If you wish grid cell data
to persist across timesteps, then use a fix.  If you wish grid cell
data to be accessible by other commands, then use a fix or compute.
Currently in LAMMPS, the {\hyperref[\detokenize{pair_amoeba::doc}]{\sphinxcrossref{\DUrole{doc}{pair\_style amoeba}}}},
{\hyperref[\detokenize{kspace_style::doc}]{\sphinxcrossref{\DUrole{doc}{kspace\_style pppm}}}}, and {\hyperref[\detokenize{kspace_style::doc}]{\sphinxcrossref{\DUrole{doc}{kspace\_style msm}}}} commands use distributed grids but do not require
either of these capabilities; they thus create and use distributed
grids internally.  Note that a pair style which needs grid cell data
to persist could be coded to work in tandem with a fix style which
provides that capability.

\sphinxAtStartPar
The \sphinxstyleemphasis{size} of a grid is specified by the number of grid cells in each
dimension of the simulation domain.  In any dimension the size can be
any value \textgreater{}= 1.  Thus a 10x10x1 grid for a 3d simulation is
effectively a 2d grid, where each grid cell spans the entire
z\sphinxhyphen{}dimension.  A 1x100x1 grid for a 3d simulation is effectively a 1d
grid, where grid cells are a series of thin xz slabs in the
y\sphinxhyphen{}dimension.  It is even possible to define a 1x1x1 3d grid, though it
may be inefficient to use it in a computational sense.

\sphinxAtStartPar
Note that the choice of grid size is independent of the number of
processors or their layout in a grid of processor subdomains which
overlays the simulations domain.  Depending on the distributed grid
size, a single processor may own many 1000s or no grid cells.

\sphinxAtStartPar
A command can define multiple grids, each of a different size.  Each
grid is an instantiation of the Grid2d or Grid3d class.

\sphinxAtStartPar
The command also defines what data it will store for each grid it
creates and it allocates the multidimensional array(s) needed to
store the data.  No grid cell data is stored within the Grid2d or
Grid3d classes.

\sphinxAtStartPar
If a single value per grid cell is needed, the data array will have
the same dimension as the grid, i.e. a 2d array for a 2d grid,
likewise for 3d.  If multiple values per grid cell are needed, the
data array will have one more dimension than the grid, i.e. a 3d array
for a 2d grid, or 4d array for a 3d grid.  A command can choose to
define multiple data arrays for each grid it defines.


\bigskip\hrule\bigskip



\subsection{Grid data allocation and access}
\label{\detokenize{Developer_grid:grid-data-allocation-and-access}}
\sphinxAtStartPar
The simplest way for a command to allocate and access grid cell data
is to use the \sphinxstyleemphasis{create\_offset()} methods provided by the Memory class.
Arguments for these methods can be values returned by the
\sphinxstyleemphasis{setup\_grid()} method (described below), which define the extent of
the grid cells (owned+ghost) the processor owns.  These 4 methods
allocate memory for 2d (first two) and 3d (second two) grid data.  The
two methods that end in “\_offset” allocate an array which stores a single
value per grid cell.  The two that end in “\_last” allocate an array
which stores \sphinxstyleemphasis{Nvalues} per grid cell.

\begin{sphinxVerbatim}[commandchars=\\\{\}]
\PYG{c+c1}{// single value per cell for a 2d grid = 2d array}
\PYG{n}{memory}\PYG{o}{\PYGZhy{}}\PYG{o}{\PYGZgt{}}\PYG{n}{create2d\PYGZus{}offset}\PYG{p}{(}\PYG{n}{data2d\PYGZus{}one}\PYG{p}{,}\PYG{+w}{ }\PYG{n}{nylo\PYGZus{}out}\PYG{p}{,}\PYG{+w}{ }\PYG{n}{nyhi\PYGZus{}out}\PYG{p}{,}
\PYG{+w}{                        }\PYG{n}{nxlo\PYGZus{}out}\PYG{p}{,}\PYG{+w}{ }\PYG{n}{nxhi\PYGZus{}out}\PYG{p}{,}\PYG{+w}{ }\PYG{l+s}{\PYGZdq{}}\PYG{l+s}{data2d\PYGZus{}one}\PYG{l+s}{\PYGZdq{}}\PYG{p}{)}\PYG{p}{;}

\PYG{c+c1}{// nvalues per cell for a 2d grid = 3d array}
\PYG{n}{memory}\PYG{o}{\PYGZhy{}}\PYG{o}{\PYGZgt{}}\PYG{n}{create3d\PYGZus{}offset\PYGZus{}last}\PYG{p}{(}\PYG{n}{data2d\PYGZus{}multi}\PYG{p}{,}\PYG{+w}{ }\PYG{n}{nylo\PYGZus{}out}\PYG{p}{,}\PYG{+w}{ }\PYG{n}{nyhi\PYGZus{}out}\PYG{p}{,}
\PYG{+w}{                             }\PYG{n}{nxlo\PYGZus{}out}\PYG{p}{,}\PYG{+w}{ }\PYG{n}{nxhi\PYGZus{}out}\PYG{p}{,}\PYG{+w}{ }\PYG{n}{nvalues}\PYG{p}{,}\PYG{+w}{ }\PYG{l+s}{\PYGZdq{}}\PYG{l+s}{data2d\PYGZus{}multi}\PYG{l+s}{\PYGZdq{}}\PYG{p}{)}\PYG{p}{;}

\PYG{c+c1}{// single value per cell for a 3d grid = 3d array}
\PYG{n}{memory}\PYG{o}{\PYGZhy{}}\PYG{o}{\PYGZgt{}}\PYG{n}{create3d\PYGZus{}offset}\PYG{p}{(}\PYG{n}{data3d\PYGZus{}one}\PYG{p}{,}\PYG{+w}{ }\PYG{n}{nzlo\PYGZus{}out}\PYG{p}{,}\PYG{+w}{ }\PYG{n}{nzhi\PYGZus{}out}\PYG{p}{,}\PYG{+w}{ }\PYG{n}{nylo\PYGZus{}out}\PYG{p}{,}
\PYG{+w}{                        }\PYG{n}{nyhi\PYGZus{}out}\PYG{p}{,}\PYG{+w}{ }\PYG{n}{nxlo\PYGZus{}out}\PYG{p}{,}\PYG{+w}{ }\PYG{n}{nxhi\PYGZus{}out}\PYG{p}{,}\PYG{+w}{ }\PYG{l+s}{\PYGZdq{}}\PYG{l+s}{data3d\PYGZus{}one}\PYG{l+s}{\PYGZdq{}}\PYG{p}{)}\PYG{p}{;}

\PYG{c+c1}{// nvalues per cell for a 3d grid = 4d array}
\PYG{n}{memory}\PYG{o}{\PYGZhy{}}\PYG{o}{\PYGZgt{}}\PYG{n}{create4d\PYGZus{}offset\PYGZus{}last}\PYG{p}{(}\PYG{n}{data3d\PYGZus{}multi}\PYG{p}{,}\PYG{+w}{ }\PYG{n}{nzlo\PYGZus{}out}\PYG{p}{,}\PYG{+w}{ }\PYG{n}{nzhi\PYGZus{}out}\PYG{p}{,}\PYG{+w}{ }\PYG{n}{nylo\PYGZus{}out}\PYG{p}{,}
\PYG{+w}{                             }\PYG{n}{nyhi\PYGZus{}out}\PYG{p}{,}\PYG{+w}{ }\PYG{n}{nxlo\PYGZus{}out}\PYG{p}{,}\PYG{+w}{ }\PYG{n}{nxhi\PYGZus{}out}\PYG{p}{,}\PYG{+w}{ }\PYG{n}{nvalues}\PYG{p}{,}
\PYG{+w}{                             }\PYG{l+s}{\PYGZdq{}}\PYG{l+s}{data3d\PYGZus{}multi}\PYG{l+s}{\PYGZdq{}}\PYG{p}{)}\PYG{p}{;}
\end{sphinxVerbatim}

\sphinxAtStartPar
Note that these multidimensional arrays are allocated as contiguous
chunks of memory where the x\sphinxhyphen{}index of the grid varies fastest, then y,
and the z\sphinxhyphen{}index slowest.  For multiple values per grid cell, the
Nvalues are contiguous, so their index varies even faster than the
x\sphinxhyphen{}index.

\sphinxAtStartPar
The key point is that the “offset” methods create arrays which are
indexed by the range of indices which are the bounds of the sub\sphinxhyphen{}block
of the global grid owned by this processor.  This means loops like
these can be written in the caller code to loop over owned grid cells,
where the “i” loop bounds are the range of owned grid cells for the
processor.  These are the bounds returned by the \sphinxstyleemphasis{setup\_grid()}
method:

\begin{sphinxVerbatim}[commandchars=\\\{\}]
\PYG{k}{for}\PYG{+w}{ }\PYG{p}{(}\PYG{k+kt}{int}\PYG{+w}{ }\PYG{n}{iy}\PYG{+w}{ }\PYG{o}{=}\PYG{+w}{ }\PYG{n}{iylo}\PYG{p}{;}\PYG{+w}{ }\PYG{n}{iy}\PYG{+w}{ }\PYG{o}{\PYGZlt{}}\PYG{o}{=}\PYG{+w}{ }\PYG{n}{iyhi}\PYG{p}{;}\PYG{+w}{ }\PYG{n}{iy}\PYG{o}{+}\PYG{o}{+}\PYG{p}{)}
\PYG{+w}{  }\PYG{k}{for}\PYG{+w}{ }\PYG{p}{(}\PYG{k+kt}{int}\PYG{+w}{ }\PYG{n}{ix}\PYG{+w}{ }\PYG{o}{=}\PYG{+w}{ }\PYG{n}{ixlo}\PYG{p}{;}\PYG{+w}{ }\PYG{n}{ix}\PYG{+w}{ }\PYG{o}{\PYGZlt{}}\PYG{o}{=}\PYG{+w}{ }\PYG{n}{ixhi}\PYG{p}{;}\PYG{+w}{ }\PYG{n}{ix}\PYG{o}{+}\PYG{o}{+}\PYG{p}{)}
\PYG{+w}{    }\PYG{n}{data2d\PYGZus{}one}\PYG{p}{[}\PYG{n}{iy}\PYG{p}{]}\PYG{p}{[}\PYG{n}{ix}\PYG{p}{]}\PYG{+w}{ }\PYG{o}{=}\PYG{+w}{ }\PYG{l+m+mf}{0.0}\PYG{p}{;}

\PYG{k}{for}\PYG{+w}{ }\PYG{p}{(}\PYG{k+kt}{int}\PYG{+w}{ }\PYG{n}{iy}\PYG{+w}{ }\PYG{o}{=}\PYG{+w}{ }\PYG{n}{iylo}\PYG{p}{;}\PYG{+w}{ }\PYG{n}{iy}\PYG{+w}{ }\PYG{o}{\PYGZlt{}}\PYG{o}{=}\PYG{+w}{ }\PYG{n}{iyhi}\PYG{p}{;}\PYG{+w}{ }\PYG{n}{iy}\PYG{o}{+}\PYG{o}{+}\PYG{p}{)}
\PYG{+w}{  }\PYG{k}{for}\PYG{+w}{ }\PYG{p}{(}\PYG{k+kt}{int}\PYG{+w}{ }\PYG{n}{ix}\PYG{+w}{ }\PYG{o}{=}\PYG{+w}{ }\PYG{n}{ixlo}\PYG{p}{;}\PYG{+w}{ }\PYG{n}{ix}\PYG{+w}{ }\PYG{o}{\PYGZlt{}}\PYG{o}{=}\PYG{+w}{ }\PYG{n}{ixhi}\PYG{p}{;}\PYG{+w}{ }\PYG{n}{ix}\PYG{o}{+}\PYG{o}{+}\PYG{p}{)}
\PYG{+w}{    }\PYG{k}{for}\PYG{+w}{ }\PYG{p}{(}\PYG{k+kt}{int}\PYG{+w}{ }\PYG{n}{m}\PYG{+w}{ }\PYG{o}{=}\PYG{+w}{ }\PYG{l+m+mi}{0}\PYG{p}{;}\PYG{+w}{ }\PYG{n}{m}\PYG{+w}{ }\PYG{o}{\PYGZlt{}}\PYG{+w}{ }\PYG{n}{nvalues}\PYG{p}{;}\PYG{+w}{ }\PYG{n}{m}\PYG{o}{+}\PYG{o}{+}\PYG{p}{)}
\PYG{+w}{      }\PYG{n}{data2d\PYGZus{}multi}\PYG{p}{[}\PYG{n}{iy}\PYG{p}{]}\PYG{p}{[}\PYG{n}{ix}\PYG{p}{]}\PYG{p}{[}\PYG{n}{m}\PYG{p}{]}\PYG{+w}{ }\PYG{o}{=}\PYG{+w}{ }\PYG{l+m+mf}{0.0}\PYG{p}{;}

\PYG{k}{for}\PYG{+w}{ }\PYG{p}{(}\PYG{k+kt}{int}\PYG{+w}{ }\PYG{n}{iz}\PYG{+w}{ }\PYG{o}{=}\PYG{+w}{ }\PYG{n}{izlo}\PYG{p}{;}\PYG{+w}{ }\PYG{n}{iz}\PYG{+w}{ }\PYG{o}{\PYGZlt{}}\PYG{o}{=}\PYG{+w}{ }\PYG{n}{izhi}\PYG{p}{;}\PYG{+w}{ }\PYG{n}{iz}\PYG{o}{+}\PYG{o}{+}\PYG{p}{)}
\PYG{+w}{  }\PYG{k}{for}\PYG{+w}{ }\PYG{p}{(}\PYG{k+kt}{int}\PYG{+w}{ }\PYG{n}{iy}\PYG{+w}{ }\PYG{o}{=}\PYG{+w}{ }\PYG{n}{iylo}\PYG{p}{;}\PYG{+w}{ }\PYG{n}{iy}\PYG{+w}{ }\PYG{o}{\PYGZlt{}}\PYG{o}{=}\PYG{+w}{ }\PYG{n}{iyhi}\PYG{p}{;}\PYG{+w}{ }\PYG{n}{iy}\PYG{o}{+}\PYG{o}{+}\PYG{p}{)}
\PYG{+w}{    }\PYG{k}{for}\PYG{+w}{ }\PYG{p}{(}\PYG{k+kt}{int}\PYG{+w}{ }\PYG{n}{ix}\PYG{+w}{ }\PYG{o}{=}\PYG{+w}{ }\PYG{n}{ixlo}\PYG{p}{;}\PYG{+w}{ }\PYG{n}{ix}\PYG{+w}{ }\PYG{o}{\PYGZlt{}}\PYG{o}{=}\PYG{+w}{ }\PYG{n}{ixhi}\PYG{p}{;}\PYG{+w}{ }\PYG{n}{ix}\PYG{o}{+}\PYG{o}{+}\PYG{p}{)}
\PYG{+w}{      }\PYG{n}{data3d\PYGZus{}one}\PYG{p}{[}\PYG{n}{iz}\PYG{p}{]}\PYG{p}{[}\PYG{n}{iy}\PYG{p}{]}\PYG{p}{[}\PYG{n}{ix}\PYG{p}{]}\PYG{+w}{ }\PYG{o}{=}\PYG{+w}{ }\PYG{l+m+mf}{0.0}\PYG{p}{;}

\PYG{k}{for}\PYG{+w}{ }\PYG{p}{(}\PYG{k+kt}{int}\PYG{+w}{ }\PYG{n}{iz}\PYG{+w}{ }\PYG{o}{=}\PYG{+w}{ }\PYG{n}{izlo}\PYG{p}{;}\PYG{+w}{ }\PYG{n}{iz}\PYG{+w}{ }\PYG{o}{\PYGZlt{}}\PYG{o}{=}\PYG{+w}{ }\PYG{n}{izhi}\PYG{p}{;}\PYG{+w}{ }\PYG{n}{iz}\PYG{o}{+}\PYG{o}{+}\PYG{p}{)}
\PYG{+w}{  }\PYG{k}{for}\PYG{+w}{ }\PYG{p}{(}\PYG{k+kt}{int}\PYG{+w}{ }\PYG{n}{iy}\PYG{+w}{ }\PYG{o}{=}\PYG{+w}{ }\PYG{n}{iylo}\PYG{p}{;}\PYG{+w}{ }\PYG{n}{iy}\PYG{+w}{ }\PYG{o}{\PYGZlt{}}\PYG{o}{=}\PYG{+w}{ }\PYG{n}{iyhi}\PYG{p}{;}\PYG{+w}{ }\PYG{n}{iy}\PYG{o}{+}\PYG{o}{+}\PYG{p}{)}
\PYG{+w}{    }\PYG{k}{for}\PYG{+w}{ }\PYG{p}{(}\PYG{k+kt}{int}\PYG{+w}{ }\PYG{n}{ix}\PYG{+w}{ }\PYG{o}{=}\PYG{+w}{ }\PYG{n}{ixlo}\PYG{p}{;}\PYG{+w}{ }\PYG{n}{ix}\PYG{+w}{ }\PYG{o}{\PYGZlt{}}\PYG{o}{=}\PYG{+w}{ }\PYG{n}{ixhi}\PYG{p}{;}\PYG{+w}{ }\PYG{n}{ix}\PYG{o}{+}\PYG{o}{+}\PYG{p}{)}
\PYG{+w}{       }\PYG{k}{for}\PYG{+w}{ }\PYG{p}{(}\PYG{k+kt}{int}\PYG{+w}{ }\PYG{n}{m}\PYG{+w}{ }\PYG{o}{=}\PYG{+w}{ }\PYG{l+m+mi}{0}\PYG{p}{;}\PYG{+w}{ }\PYG{n}{m}\PYG{+w}{ }\PYG{o}{\PYGZlt{}}\PYG{+w}{ }\PYG{n}{nvalues}\PYG{p}{;}\PYG{+w}{ }\PYG{n}{m}\PYG{o}{+}\PYG{o}{+}\PYG{p}{)}
\PYG{+w}{          }\PYG{n}{data3d\PYGZus{}multi}\PYG{p}{[}\PYG{n}{iz}\PYG{p}{]}\PYG{p}{[}\PYG{n}{iy}\PYG{p}{]}\PYG{p}{[}\PYG{n}{ix}\PYG{p}{]}\PYG{p}{[}\PYG{n}{m}\PYG{p}{]}\PYG{+w}{ }\PYG{o}{=}\PYG{+w}{ }\PYG{l+m+mf}{0.0}\PYG{p}{;}
\end{sphinxVerbatim}

\sphinxAtStartPar
Simply replacing the “i” bounds with “o” bounds, also returned by the
\sphinxstyleemphasis{setup\_grid()} method, would alter this code to loop over owned+ghost
cells (the entire allocated grid).


\bigskip\hrule\bigskip



\subsection{Grid class constructors}
\label{\detokenize{Developer_grid:grid-class-constructors}}
\sphinxAtStartPar
The following subsections describe the public methods of the Grid3d
class which a style command can invoke.  The Grid2d methods are
similar; simply remove arguments which refer to the z\sphinxhyphen{}dimension.

\sphinxAtStartPar
There are 2 constructors which can be used.  They differ in the extra
i/o xyz lo/hi arguments:

\begin{sphinxVerbatim}[commandchars=\\\{\}]
\PYG{n}{Grid3d}\PYG{p}{(}\PYG{k}{class}\PYG{+w}{ }\PYG{n+nc}{LAMMPS}\PYG{+w}{ }\PYG{o}{*}\PYG{n}{lmp}\PYG{p}{,}\PYG{+w}{ }\PYG{n}{MPI\PYGZus{}Comm}\PYG{+w}{ }\PYG{n}{gcomm}\PYG{p}{,}\PYG{+w}{ }\PYG{k+kt}{int}\PYG{+w}{ }\PYG{n}{gnx}\PYG{p}{,}\PYG{+w}{ }\PYG{k+kt}{int}\PYG{+w}{ }\PYG{n}{gny}\PYG{p}{,}\PYG{+w}{ }\PYG{k+kt}{int}\PYG{+w}{ }\PYG{n}{gnz}\PYG{p}{)}
\PYG{n}{Grid3d}\PYG{p}{(}\PYG{k}{class}\PYG{+w}{ }\PYG{n+nc}{LAMMPS}\PYG{+w}{ }\PYG{o}{*}\PYG{n}{lmp}\PYG{p}{,}\PYG{+w}{ }\PYG{n}{MPI\PYGZus{}Comm}\PYG{+w}{ }\PYG{n}{gcomm}\PYG{p}{,}\PYG{+w}{ }\PYG{k+kt}{int}\PYG{+w}{ }\PYG{n}{gnx}\PYG{p}{,}\PYG{+w}{ }\PYG{k+kt}{int}\PYG{+w}{ }\PYG{n}{gny}\PYG{p}{,}\PYG{+w}{ }\PYG{k+kt}{int}\PYG{+w}{ }\PYG{n}{gnz}\PYG{p}{,}
\PYG{+w}{       }\PYG{k+kt}{int}\PYG{+w}{ }\PYG{n}{ixlo}\PYG{p}{,}\PYG{+w}{ }\PYG{k+kt}{int}\PYG{+w}{ }\PYG{n}{ixhi}\PYG{p}{,}\PYG{+w}{ }\PYG{k+kt}{int}\PYG{+w}{ }\PYG{n}{iylo}\PYG{p}{,}\PYG{+w}{ }\PYG{k+kt}{int}\PYG{+w}{ }\PYG{n}{iyhi}\PYG{p}{,}\PYG{+w}{ }\PYG{k+kt}{int}\PYG{+w}{ }\PYG{n}{izlo}\PYG{p}{,}\PYG{+w}{ }\PYG{k+kt}{int}\PYG{+w}{ }\PYG{n}{izhi}\PYG{p}{,}
\PYG{+w}{       }\PYG{k+kt}{int}\PYG{+w}{ }\PYG{n}{oxlo}\PYG{p}{,}\PYG{+w}{ }\PYG{k+kt}{int}\PYG{+w}{ }\PYG{n}{oxhi}\PYG{p}{,}\PYG{+w}{ }\PYG{k+kt}{int}\PYG{+w}{ }\PYG{n}{oylo}\PYG{p}{,}\PYG{+w}{ }\PYG{k+kt}{int}\PYG{+w}{ }\PYG{n}{oyhi}\PYG{p}{,}\PYG{+w}{ }\PYG{k+kt}{int}\PYG{+w}{ }\PYG{n}{ozlo}\PYG{p}{,}\PYG{+w}{ }\PYG{k+kt}{int}\PYG{+w}{ }\PYG{n}{ozhi}\PYG{p}{)}
\end{sphinxVerbatim}

\sphinxAtStartPar
Both constructors take the LAMMPS instance pointer and a communicator
over which the grid will be distributed.  Typically this is the
\sphinxstyleemphasis{world} communicator the LAMMPS instance is using.  The
{\hyperref[\detokenize{kspace_style::doc}]{\sphinxcrossref{\DUrole{doc}{kspace\_style msm}}}} command creates a series of
grids, each of different size, which are partitioned across different
sub\sphinxhyphen{}communicators of processors.  Both constructors are also passed
the global grid size: \sphinxstyleemphasis{gnx} by \sphinxstyleemphasis{gny} by \sphinxstyleemphasis{gnz}.

\sphinxAtStartPar
The first constructor is used when the caller wants the Grid class to
partition the global grid across processors; the Grid class defines
which grid cells each processor owns and also which it stores as ghost
cells.  A subsequent call to \sphinxstyleemphasis{setup\_grid()}, discussed below, returns
this info to the caller.

\sphinxAtStartPar
The second constructor allows the caller to define the extent of owned
and ghost cells, and pass them to the Grid class.  The 6 arguments
which start with “i” are the inclusive lower and upper index bounds of
the owned (inner) grid cells this processor owns in each of the 3
dimensions within the global grid.  Owned grid cells are indexed from
0 to N\sphinxhyphen{}1 in each dimension.

\sphinxAtStartPar
The 6 arguments which start with “o” are the inclusive bounds of the
owned+ghost (outer) grid cells it stores.  If the ghost cells are on
the other side of a periodic boundary, then these indices may be \textless{} 0
or \textgreater{}= N in any dimension, so that oxlo \textless{}= ixlo and ixhi \textgreater{}= ixhi is
always the case.

\sphinxAtStartPar
For example, if Nx = 100, then a processor might pass ixlo=50,
ixhi=60, oxlo=48, oxhi=62 to the Grid class.  Or ixlo=0, ixhi=10,
oxlo=\sphinxhyphen{}2, oxhi=13.  If a processor owns no grid cells in a dimension,
then the ihi value should be specified as one less than the ilo value.

\sphinxAtStartPar
Note that the only reason to use the second constructor is if the
logic for assigning ghost cells is too complex for the Grid class to
compute, using the various set() methods described next.  Currently
only the kspace\_style pppm/electrode and kspace\_style msm commands use
the second constructor.


\bigskip\hrule\bigskip



\subsection{Grid class set methods}
\label{\detokenize{Developer_grid:grid-class-set-methods}}
\sphinxAtStartPar
The following methods affect how the Grid class computes which owned
and ghost cells are assigned to each processor.  \sphinxstyleemphasis{Set\_shift\_grid()} is
the only method which influences owned cell assignment; all the rest
influence ghost cell assignment.  These methods are only used with the
first constructor; they are ignored if the second constructor is used.
These methods must be called before the \sphinxstyleemphasis{setup\_grid()} method is
invoked, because they influence its operation.

\begin{sphinxVerbatim}[commandchars=\\\{\}]
\PYG{k+kt}{void}\PYG{+w}{ }\PYG{n+nf}{set\PYGZus{}shift\PYGZus{}grid}\PYG{p}{(}\PYG{k+kt}{double}\PYG{+w}{ }\PYG{n}{shift}\PYG{p}{)}\PYG{p}{;}
\PYG{k+kt}{void}\PYG{+w}{ }\PYG{n+nf}{set\PYGZus{}distance}\PYG{p}{(}\PYG{k+kt}{double}\PYG{+w}{ }\PYG{n}{distance}\PYG{p}{)}\PYG{p}{;}
\PYG{k+kt}{void}\PYG{+w}{ }\PYG{n+nf}{set\PYGZus{}stencil\PYGZus{}atom}\PYG{p}{(}\PYG{k+kt}{int}\PYG{+w}{ }\PYG{n}{lo}\PYG{p}{,}\PYG{+w}{ }\PYG{k+kt}{int}\PYG{+w}{ }\PYG{n}{hi}\PYG{p}{)}\PYG{p}{;}
\PYG{k+kt}{void}\PYG{+w}{ }\PYG{n+nf}{set\PYGZus{}shift\PYGZus{}atom}\PYG{p}{(}\PYG{k+kt}{double}\PYG{+w}{ }\PYG{n}{shift\PYGZus{}lo}\PYG{p}{,}\PYG{+w}{ }\PYG{k+kt}{double}\PYG{+w}{ }\PYG{n}{shift\PYGZus{}hi}\PYG{p}{)}\PYG{p}{;}
\PYG{k+kt}{void}\PYG{+w}{ }\PYG{n+nf}{set\PYGZus{}stencil\PYGZus{}grid}\PYG{p}{(}\PYG{k+kt}{int}\PYG{+w}{ }\PYG{n}{lo}\PYG{p}{,}\PYG{+w}{ }\PYG{k+kt}{int}\PYG{+w}{ }\PYG{n}{hi}\PYG{p}{)}\PYG{p}{;}
\PYG{k+kt}{void}\PYG{+w}{ }\PYG{n+nf}{set\PYGZus{}zfactor}\PYG{p}{(}\PYG{k+kt}{double}\PYG{+w}{ }\PYG{n}{factor}\PYG{p}{)}\PYG{p}{;}
\end{sphinxVerbatim}

\sphinxAtStartPar
Processors own a grid cell if a point within the grid cell is inside
the processor’s subdomain.  By default this is the center point of the
grid cell.  The \sphinxstyleemphasis{set\_shift\_grid()} method can change this.  The \sphinxstyleemphasis{shift}
argument is a value from 0.0 to 1.0 (inclusive) which is the offset of
the point within the grid cell in each dimension.  The default is 0.5
for the center of the cell.  A value of 0.0 is the lower left corner
point; a value of 1.0 is the upper right corner point.  There is
typically no need to change the default as it is optimal for
minimizing the number of ghost cells needed.

\sphinxAtStartPar
If a processor maps its particles to grid cells, it needs to allow for
its particles being outside its subdomain between reneighboring.  The
\sphinxstyleemphasis{distance} argument of the \sphinxstyleemphasis{set\_distance()} method sets the furthest
distance outside a processor’s subdomain which a particle can move.
Typically this is half the neighbor skin distance, assuming
reneighboring is done appropriately.  This distance is used in
determining how many ghost cells a processor needs to store to enable
its particles to be mapped to grid cells.  The default value is 0.0.

\sphinxAtStartPar
Some commands, like the {\hyperref[\detokenize{kspace_style::doc}]{\sphinxcrossref{\DUrole{doc}{kspace\_style pppm}}}}
command, map values (charge in the case of PPPM) to a stencil of grid
cells beyond the grid cell the particle is in.  The stencil extent may
be different in the low and high directions.  The \sphinxstyleemphasis{set\_stencil\_atom()}
method defines the maximum values of those 2 extents, assumed to be
the same in each of the 3 dimensions.  Both the lo and hi values are
specified as positive integers.  The default values are both 0.

\sphinxAtStartPar
Some commands, like the {\hyperref[\detokenize{kspace_style::doc}]{\sphinxcrossref{\DUrole{doc}{kspace\_style pppm}}}}
command, shift the position of an atom when mapping it to a grid cell,
based on the size of the stencil used to map values to the grid
(charge in the case of PPPM).  The lo and hi arguments of the
\sphinxstyleemphasis{set\_shift\_atom()} method are the minimum shift in the low direction
and the maximum shift in the high direction, assumed to be the same in
each of the 3 dimensions.  The shifts should be fractions of a grid
cell size with values between 0.0 and 1.0 inclusive.  The default
values are both 0.0.  See the src/pppm.cpp file for examples of these
lo/hi values for regular and staggered grids.

\sphinxAtStartPar
Some methods like the {\hyperref[\detokenize{fix_ttm::doc}]{\sphinxcrossref{\DUrole{doc}{fix ttm/grid}}}} command, perform
finite difference kinds of operations on the grid, to diffuse electron
heat in the case of the two\sphinxhyphen{}temperature model (TTM).  This operation
uses ghost grid values beyond the owned grid values the processor
updates.  The \sphinxstyleemphasis{set\_stencil\_grid()} method defines the extent of this
stencil in both directions, assumed to be the same in each of the 3
dimensions.  Both the lo and hi values are specified as positive
integers.  The default values are both 0.

\sphinxAtStartPar
The kspace\_style pppm commands allow a grid to be defined which
overlays a volume which extends beyond the simulation box in the z
dimension.  This is for the purpose of modeling a 2d\sphinxhyphen{}periodic slab
(non\sphinxhyphen{}periodic in z) as if it were a larger 3d periodic system,
extended (with empty space) in the z dimension.  The
{\hyperref[\detokenize{kspace_modify::doc}]{\sphinxcrossref{\DUrole{doc}{kspace\_modify slab}}}} command is used to specify
the ratio of the larger volume to the simulation volume; a volume
ratio of \textasciitilde{}3 is typical.  For this kind of model, the PPPM caller sets
the global grid size \sphinxstyleemphasis{gnz} \textasciitilde{}3x larger than it would be otherwise.
This same ratio is passed by the PPPM caller as the \sphinxstyleemphasis{factor} argument
to the Grid class via the \sphinxstyleemphasis{set\_zfactor()} method (\sphinxstyleemphasis{set\_yfactor()} for
2d grids).  The Grid class will then assign ownership of the 1/3 of
grid cells that overlay the simulation box to the processors which
also overlay the simulation box.  The remaining 2/3 of the grid cells
are assigned to processors whose subdomains are adjacent to the upper
z boundary of the simulation box.


\bigskip\hrule\bigskip



\subsection{Grid class setup\_grid method}
\label{\detokenize{Developer_grid:grid-class-setup-grid-method}}
\sphinxAtStartPar
The \sphinxstyleemphasis{setup\_grid()} method is called after the first constructor
(above) to partition the grid across processors, which determines
which grid cells each processor owns.  It also calculates how many
ghost grid cells in each dimension and each direction each processor
needs to store.

\sphinxAtStartPar
Note that this method is NOT called if the second constructor above is
used.  In that case, the caller assigns owned and ghost cells to each
processor.

\sphinxAtStartPar
Also note that this method must be invoked after any \sphinxstyleemphasis{set\_*()} methods have
been used, since they can influence the assignment of owned and ghost
cells.

\begin{sphinxVerbatim}[commandchars=\\\{\}]
\PYG{k+kt}{void}\PYG{+w}{ }\PYG{n}{setup\PYGZus{}grid}\PYG{p}{(}\PYG{k+kt}{int}\PYG{+w}{ }\PYG{o}{\PYGZam{}}\PYG{n}{ixlo}\PYG{p}{,}\PYG{+w}{ }\PYG{k+kt}{int}\PYG{+w}{ }\PYG{o}{\PYGZam{}}\PYG{n}{ixhi}\PYG{p}{,}\PYG{+w}{ }\PYG{k+kt}{int}\PYG{+w}{ }\PYG{o}{\PYGZam{}}\PYG{n}{iylo}\PYG{p}{,}\PYG{+w}{ }\PYG{k+kt}{int}\PYG{+w}{ }\PYG{o}{\PYGZam{}}\PYG{n}{iyhi}\PYG{p}{,}\PYG{+w}{ }\PYG{k+kt}{int}\PYG{+w}{ }\PYG{o}{\PYGZam{}}\PYG{n}{izlo}\PYG{p}{,}\PYG{+w}{ }\PYG{k+kt}{int}\PYG{+w}{ }\PYG{o}{\PYGZam{}}\PYG{n}{izhi}\PYG{p}{,}
\PYG{+w}{                }\PYG{k+kt}{int}\PYG{+w}{ }\PYG{o}{\PYGZam{}}\PYG{n}{oxlo}\PYG{p}{,}\PYG{+w}{ }\PYG{k+kt}{int}\PYG{+w}{ }\PYG{o}{\PYGZam{}}\PYG{n}{oxhi}\PYG{p}{,}\PYG{+w}{ }\PYG{k+kt}{int}\PYG{+w}{ }\PYG{o}{\PYGZam{}}\PYG{n}{oylo}\PYG{p}{,}\PYG{+w}{ }\PYG{k+kt}{int}\PYG{+w}{ }\PYG{o}{\PYGZam{}}\PYG{n}{oyhi}\PYG{p}{,}\PYG{+w}{ }\PYG{k+kt}{int}\PYG{+w}{ }\PYG{o}{\PYGZam{}}\PYG{n}{ozlo}\PYG{p}{,}\PYG{+w}{ }\PYG{k+kt}{int}\PYG{+w}{ }\PYG{o}{\PYGZam{}}\PYG{n}{ozhi}\PYG{p}{)}
\end{sphinxVerbatim}

\sphinxAtStartPar
The 6 return arguments which start with “i” are the inclusive lower
and upper index bounds of the owned (inner) grid cells this processor
owns in each of the 3 dimensions within the global grid.  Owned grid
cells are indexed from 0 to N\sphinxhyphen{}1 in each dimension.

\sphinxAtStartPar
The 6 return arguments which start with “o” are the inclusive bounds of
the owned+ghost cells it owns.  If the ghost cells are on the other
side of a periodic boundary, then these indices may be \textless{} 0 or \textgreater{}= N in
any dimension, so that oxlo \textless{}= ixlo and ixhi \textgreater{}= ixhi is always the
case.


\bigskip\hrule\bigskip



\subsection{More grid class set methods}
\label{\detokenize{Developer_grid:more-grid-class-set-methods}}
\sphinxAtStartPar
The following 2 methods can be used to override settings made by the
constructors above.  If used, they must be called called before the
\sphinxstyleemphasis{setup\_comm()} method is invoked, since it uses the settings that
these methods override.  In LAMMPS these methods are called by by the
{\hyperref[\detokenize{kspace_style::doc}]{\sphinxcrossref{\DUrole{doc}{kspace\_style msm}}}} command for the grids it
instantiates using the 2nd constructor above.

\begin{sphinxVerbatim}[commandchars=\\\{\}]
\PYG{k+kt}{void}\PYG{+w}{ }\PYG{n}{set\PYGZus{}proc\PYGZus{}neighs}\PYG{p}{(}\PYG{k+kt}{int}\PYG{+w}{ }\PYG{n}{pxlo}\PYG{p}{,}\PYG{+w}{ }\PYG{k+kt}{int}\PYG{+w}{ }\PYG{n}{pxhi}\PYG{p}{,}\PYG{+w}{ }\PYG{k+kt}{int}\PYG{+w}{ }\PYG{n}{pylo}\PYG{p}{,}\PYG{+w}{ }\PYG{k+kt}{int}\PYG{+w}{ }\PYG{n}{pyhi}\PYG{p}{,}\PYG{+w}{ }\PYG{k+kt}{int}\PYG{+w}{ }\PYG{n}{pzlo}\PYG{p}{,}\PYG{+w}{ }\PYG{k+kt}{int}\PYG{+w}{ }\PYG{n}{pzhi}\PYG{p}{)}
\PYG{k+kt}{void}\PYG{+w}{ }\PYG{n}{set\PYGZus{}caller\PYGZus{}grid}\PYG{p}{(}\PYG{k+kt}{int}\PYG{+w}{ }\PYG{n}{fxlo}\PYG{p}{,}\PYG{+w}{ }\PYG{k+kt}{int}\PYG{+w}{ }\PYG{n}{fxhi}\PYG{p}{,}\PYG{+w}{ }\PYG{k+kt}{int}\PYG{+w}{ }\PYG{n}{fylo}\PYG{p}{,}\PYG{+w}{ }\PYG{k+kt}{int}\PYG{+w}{ }\PYG{n}{fyhi}\PYG{p}{,}\PYG{+w}{ }\PYG{k+kt}{int}\PYG{+w}{ }\PYG{n}{fzlo}\PYG{p}{,}\PYG{+w}{ }\PYG{k+kt}{int}\PYG{+w}{ }\PYG{n}{fzhi}\PYG{p}{)}
\end{sphinxVerbatim}

\sphinxAtStartPar
The \sphinxstyleemphasis{set\_proc\_neighs()} method sets the processor IDs of the 6
neighboring processors for each processor.  Normally these would match
the processor grid neighbors which LAMMPS creates to overlay the
simulation box (the default).  However, MSM excludes non\sphinxhyphen{}participating
processors from coarse grid communication when less processors are
used.  This method allows MSM to override the default values.

\sphinxAtStartPar
The \sphinxstyleemphasis{set\_caller\_grid()} method species the size of the data arrays the
caller allocates.  Normally these would match the extent of the ghost
grid cells (the default).  However the MSM caller allocates a larger
data array (more ghost cells) for its finest\sphinxhyphen{}level grid, for use in
other operations besides owned/ghost cell communication.  This method
allows MSM to override the default values.


\bigskip\hrule\bigskip



\subsection{Grid class get methods}
\label{\detokenize{Developer_grid:grid-class-get-methods}}
\sphinxAtStartPar
The following methods allow the caller to query the settings for a
specific grid, whether it created the grid or another command created
it.

\begin{sphinxVerbatim}[commandchars=\\\{\}]
\PYG{k+kt}{void}\PYG{+w}{ }\PYG{n+nf}{get\PYGZus{}size}\PYG{p}{(}\PYG{k+kt}{int}\PYG{+w}{ }\PYG{o}{\PYGZam{}}\PYG{n}{nxgrid}\PYG{p}{,}\PYG{+w}{ }\PYG{k+kt}{int}\PYG{+w}{ }\PYG{o}{\PYGZam{}}\PYG{n}{nygrid}\PYG{p}{,}\PYG{+w}{ }\PYG{k+kt}{int}\PYG{+w}{ }\PYG{o}{\PYGZam{}}\PYG{n}{nzgrid}\PYG{p}{)}\PYG{p}{;}
\PYG{k+kt}{void}\PYG{+w}{ }\PYG{n}{get\PYGZus{}bounds\PYGZus{}owned}\PYG{p}{(}\PYG{k+kt}{int}\PYG{+w}{ }\PYG{o}{\PYGZam{}}\PYG{n}{xlo}\PYG{p}{,}\PYG{+w}{ }\PYG{k+kt}{int}\PYG{+w}{ }\PYG{o}{\PYGZam{}}\PYG{n}{xhi}\PYG{p}{,}\PYG{+w}{ }\PYG{k+kt}{int}\PYG{+w}{ }\PYG{o}{\PYGZam{}}\PYG{n}{ylo}\PYG{p}{,}\PYG{+w}{ }\PYG{k+kt}{int}\PYG{+w}{ }\PYG{o}{\PYGZam{}}\PYG{n}{yhi}\PYG{p}{,}\PYG{+w}{ }\PYG{k+kt}{int}\PYG{+w}{ }\PYG{o}{\PYGZam{}}\PYG{n}{zlo}\PYG{p}{,}\PYG{+w}{ }\PYG{k+kt}{int}\PYG{+w}{ }\PYG{o}{\PYGZam{}}\PYG{n}{zhi}\PYG{p}{)}
\PYG{k+kt}{void}\PYG{+w}{ }\PYG{n}{get\PYGZus{}bounds\PYGZus{}ghost}\PYG{p}{(}\PYG{k+kt}{int}\PYG{+w}{ }\PYG{o}{\PYGZam{}}\PYG{n}{xlo}\PYG{p}{,}\PYG{+w}{ }\PYG{k+kt}{int}\PYG{+w}{ }\PYG{o}{\PYGZam{}}\PYG{n}{xhi}\PYG{p}{,}\PYG{+w}{ }\PYG{k+kt}{int}\PYG{+w}{ }\PYG{o}{\PYGZam{}}\PYG{n}{ylo}\PYG{p}{,}\PYG{+w}{ }\PYG{k+kt}{int}\PYG{+w}{ }\PYG{o}{\PYGZam{}}\PYG{n}{yhi}\PYG{p}{,}\PYG{+w}{ }\PYG{k+kt}{int}\PYG{+w}{ }\PYG{o}{\PYGZam{}}\PYG{n}{zlo}\PYG{p}{,}\PYG{+w}{ }\PYG{k+kt}{int}\PYG{+w}{ }\PYG{o}{\PYGZam{}}\PYG{n}{zhi}\PYG{p}{)}
\end{sphinxVerbatim}

\sphinxAtStartPar
The \sphinxstyleemphasis{get\_size()} method returns the size of the global grid in each dimension.

\sphinxAtStartPar
The \sphinxstyleemphasis{get\_bounds\_owned()} method return the inclusive index bounds of
the grid cells this processor owns.  The values range from 0 to N\sphinxhyphen{}1 in
each dimension.  These values are the same as the “i” values returned
by \sphinxstyleemphasis{setup\_grid()}.

\sphinxAtStartPar
The \sphinxstyleemphasis{get\_bounds\_ghost()} method return the inclusive index bounds of
the owned+ghost grid cells this processor stores.  The owned cell
indices range from 0 to N\sphinxhyphen{}1, so these indices may be less than 0 or
greater than or equal to N in each dimension.  These values are the
same as the “o” values returned by \sphinxstyleemphasis{setup\_grid()}.


\bigskip\hrule\bigskip



\subsection{Grid class owned/ghost communication}
\label{\detokenize{Developer_grid:grid-class-owned-ghost-communication}}
\sphinxAtStartPar
If needed by the command, the following methods setup and perform
communication of grid data to/from neighboring processors.  The
\sphinxstyleemphasis{forward\_comm()} method sends owned grid cell data to the
corresponding ghost grid cells on other processors.  The
\sphinxstyleemphasis{reverse\_comm()} method sends ghost grid cell data to the
corresponding owned grid cells on another processor.  The caller can
choose to sum ghost grid cell data to the owned grid cell or simply
copy it.

\begin{sphinxVerbatim}[commandchars=\\\{\}]
\PYG{k+kt}{void}\PYG{+w}{ }\PYG{n+nf}{setup\PYGZus{}comm}\PYG{p}{(}\PYG{k+kt}{int}\PYG{+w}{ }\PYG{o}{\PYGZam{}}\PYG{n}{nbuf1}\PYG{p}{,}\PYG{+w}{ }\PYG{k+kt}{int}\PYG{+w}{ }\PYG{o}{\PYGZam{}}\PYG{n}{nbuf2}\PYG{p}{)}
\PYG{k+kt}{void}\PYG{+w}{ }\PYG{n}{forward\PYGZus{}comm}\PYG{p}{(}\PYG{k+kt}{int}\PYG{+w}{ }\PYG{n}{caller}\PYG{p}{,}\PYG{+w}{ }\PYG{k+kt}{void}\PYG{+w}{ }\PYG{o}{*}\PYG{n}{ptr}\PYG{p}{,}\PYG{+w}{ }\PYG{k+kt}{int}\PYG{+w}{ }\PYG{n}{which}\PYG{p}{,}\PYG{+w}{ }\PYG{k+kt}{int}\PYG{+w}{ }\PYG{n}{nper}\PYG{p}{,}\PYG{+w}{ }\PYG{k+kt}{int}\PYG{+w}{ }\PYG{n}{nbyte}\PYG{p}{,}
\PYG{+w}{                  }\PYG{k+kt}{void}\PYG{+w}{ }\PYG{o}{*}\PYG{n}{buf1}\PYG{p}{,}\PYG{+w}{ }\PYG{k+kt}{void}\PYG{+w}{ }\PYG{o}{*}\PYG{n}{buf2}\PYG{p}{,}\PYG{+w}{ }\PYG{n}{MPI\PYGZus{}Datatype}\PYG{+w}{ }\PYG{n}{datatype}\PYG{p}{)}\PYG{p}{;}
\PYG{k+kt}{void}\PYG{+w}{ }\PYG{n+nf}{reverse\PYGZus{}comm}\PYG{p}{(}\PYG{k+kt}{int}\PYG{+w}{ }\PYG{n}{caller}\PYG{p}{,}\PYG{+w}{ }\PYG{k+kt}{void}\PYG{+w}{ }\PYG{o}{*}\PYG{n}{ptr}\PYG{p}{,}\PYG{+w}{ }\PYG{k+kt}{int}\PYG{+w}{ }\PYG{n}{which}\PYG{p}{,}\PYG{+w}{ }\PYG{k+kt}{int}\PYG{+w}{ }\PYG{n}{nper}\PYG{p}{,}\PYG{+w}{ }\PYG{k+kt}{int}\PYG{+w}{ }\PYG{n}{nbyte}\PYG{p}{,}
\PYG{+w}{                  }\PYG{k+kt}{void}\PYG{+w}{ }\PYG{o}{*}\PYG{n}{buf1}\PYG{p}{,}\PYG{+w}{ }\PYG{k+kt}{void}\PYG{+w}{ }\PYG{o}{*}\PYG{n}{buf2}\PYG{p}{,}\PYG{+w}{ }\PYG{n}{MPI\PYGZus{}Datatype}\PYG{+w}{ }\PYG{n}{datatype}\PYG{p}{)}
\PYG{k+kt}{int}\PYG{+w}{ }\PYG{n}{ghost\PYGZus{}adjacent}\PYG{p}{(}\PYG{p}{)}\PYG{p}{;}
\end{sphinxVerbatim}

\sphinxAtStartPar
The \sphinxstyleemphasis{setup\_comm()} method must be called one time before performing
\sphinxstyleemphasis{forward} or \sphinxstyleemphasis{reverse} communication (multiple times if needed).  It
returns two integers, which should be used to allocate two buffers.
The \sphinxstyleemphasis{nbuf1} and \sphinxstyleemphasis{nbuf2} values are the number of grid cells whose data
will be stored in two buffers by the Grid class when \sphinxstyleemphasis{forward} or
\sphinxstyleemphasis{reverse} communication is performed.  The caller should thus allocate
them to a size large enough to hold all the data used in any single
forward or reverse communication operation it performs.  Note that the
caller may allocate and communicate multiple data arrays for a grid it
instantiates.  This size includes the bytes needed for the data type
of the grid data it stores, e.g. double precision values.

\sphinxAtStartPar
The \sphinxstyleemphasis{forward\_comm()} and \sphinxstyleemphasis{reverse\_comm()} methods send grid cell data
from owned to ghost cells, or ghost to owned cells, respectively, as
described above.  The \sphinxstyleemphasis{caller} argument should be one of these values
\textendash{} Grid3d::COMPUTE, Grid3d::FIX, Grid3d::KSPACE, Grid3d::PAIR \textendash{}
depending on the style of the caller class.  The \sphinxstyleemphasis{ptr} argument is the
“this” pointer to the caller class.  These two arguments are used to
call back to pack()/unpack() functions in the caller class, as
explained below.

\sphinxAtStartPar
The \sphinxstyleemphasis{which} argument is a flag the caller can set which is passed to
the caller’s pack()/unpack() methods.  This allows a single callback
method to pack/unpack data for several different flavors of
forward/reverse communication, e.g. operating on different grids or
grid data.

\sphinxAtStartPar
The \sphinxstyleemphasis{nper} argument is the number of values per grid cell to be
communicated.  The \sphinxstyleemphasis{nbyte} argument is the number of bytes per value,
e.g. 8 for double\sphinxhyphen{}precision values.  The \sphinxstyleemphasis{buf1} and \sphinxstyleemphasis{buf2} arguments
are the two allocated buffers described above.  So long as they are
allocated for the maximum size communication, they can be re\sphinxhyphen{}used for
any \sphinxstyleemphasis{forward\_comm()/reverse\_comm()} call.  The \sphinxstyleemphasis{datatype} argument is
the MPI\_Datatype setting, which should match the buffer allocation and
the \sphinxstyleemphasis{nbyte} argument.  E.g. MPI\_DOUBLE for buffers storing double
precision values.

\sphinxAtStartPar
To use the \sphinxstyleemphasis{forward\_grid()} method, the caller must provide two
callback functions; likewise for use of the \sphinxstyleemphasis{reverse\_grid()} methods.
These are the 4 functions, their arguments are all the same.

\begin{sphinxVerbatim}[commandchars=\\\{\}]
\PYG{k+kt}{void}\PYG{+w}{ }\PYG{n+nf}{pack\PYGZus{}forward\PYGZus{}grid}\PYG{p}{(}\PYG{k+kt}{int}\PYG{+w}{ }\PYG{n}{which}\PYG{p}{,}\PYG{+w}{ }\PYG{k+kt}{void}\PYG{+w}{ }\PYG{o}{*}\PYG{n}{vbuf}\PYG{p}{,}\PYG{+w}{ }\PYG{k+kt}{int}\PYG{+w}{ }\PYG{n}{nlist}\PYG{p}{,}\PYG{+w}{ }\PYG{k+kt}{int}\PYG{+w}{ }\PYG{o}{*}\PYG{n}{list}\PYG{p}{)}\PYG{p}{;}
\PYG{k+kt}{void}\PYG{+w}{ }\PYG{n+nf}{unpack\PYGZus{}forward\PYGZus{}grid}\PYG{p}{(}\PYG{k+kt}{int}\PYG{+w}{ }\PYG{n}{which}\PYG{p}{,}\PYG{+w}{ }\PYG{k+kt}{void}\PYG{+w}{ }\PYG{o}{*}\PYG{n}{vbuf}\PYG{p}{,}\PYG{+w}{ }\PYG{k+kt}{int}\PYG{+w}{ }\PYG{n}{nlist}\PYG{p}{,}\PYG{+w}{ }\PYG{k+kt}{int}\PYG{+w}{ }\PYG{o}{*}\PYG{n}{list}\PYG{p}{)}\PYG{p}{;}
\PYG{k+kt}{void}\PYG{+w}{ }\PYG{n+nf}{pack\PYGZus{}reverse\PYGZus{}grid}\PYG{p}{(}\PYG{k+kt}{int}\PYG{+w}{ }\PYG{n}{which}\PYG{p}{,}\PYG{+w}{ }\PYG{k+kt}{void}\PYG{+w}{ }\PYG{o}{*}\PYG{n}{vbuf}\PYG{p}{,}\PYG{+w}{ }\PYG{k+kt}{int}\PYG{+w}{ }\PYG{n}{nlist}\PYG{p}{,}\PYG{+w}{ }\PYG{k+kt}{int}\PYG{+w}{ }\PYG{o}{*}\PYG{n}{list}\PYG{p}{)}\PYG{p}{;}
\PYG{k+kt}{void}\PYG{+w}{ }\PYG{n+nf}{unpack\PYGZus{}reverse\PYGZus{}grid}\PYG{p}{(}\PYG{k+kt}{int}\PYG{+w}{ }\PYG{n}{which}\PYG{p}{,}\PYG{+w}{ }\PYG{k+kt}{void}\PYG{+w}{ }\PYG{o}{*}\PYG{n}{vbuf}\PYG{p}{,}\PYG{+w}{ }\PYG{k+kt}{int}\PYG{+w}{ }\PYG{n}{nlist}\PYG{p}{,}\PYG{+w}{ }\PYG{k+kt}{int}\PYG{+w}{ }\PYG{o}{*}\PYG{n}{list}\PYG{p}{)}\PYG{p}{;}
\end{sphinxVerbatim}

\sphinxAtStartPar
The \sphinxstyleemphasis{which} argument is set to the \sphinxstyleemphasis{which} value of the
\sphinxstyleemphasis{forward\_comm()} or \sphinxstyleemphasis{reverse\_comm()} calls.  It allows the pack/unpack
function to select what data values to pack/unpack.  \sphinxstyleemphasis{Vbuf} is the
buffer to pack/unpack the data to/from.  It is a void pointer so that
the caller can cast it to whatever data type it chooses, e.g. double
precision values.  \sphinxstyleemphasis{Nlist} is the number of grid cells to pack/unpack
and \sphinxstyleemphasis{list} is a vector (nlist in length) of offsets to where the data
for each grid cell resides in the caller’s data arrays, which is best
illustrated with an example from the src/EXTRA\sphinxhyphen{}FIX/fix\_ttm\_grid.cpp
class which stores the scalar electron temperature for 3d system in a
3d grid (one value per grid cell):

\begin{sphinxVerbatim}[commandchars=\\\{\}]
\PYG{k+kt}{void}\PYG{+w}{ }\PYG{n+nf}{FixTTMGrid::pack\PYGZus{}forward\PYGZus{}grid}\PYG{p}{(}\PYG{k+kt}{int}\PYG{+w}{ }\PYG{c+cm}{/*which*/}\PYG{p}{,}\PYG{+w}{ }\PYG{k+kt}{void}\PYG{+w}{ }\PYG{o}{*}\PYG{n}{vbuf}\PYG{p}{,}\PYG{+w}{ }\PYG{k+kt}{int}\PYG{+w}{ }\PYG{n}{nlist}\PYG{p}{,}\PYG{+w}{ }\PYG{k+kt}{int}\PYG{+w}{ }\PYG{o}{*}\PYG{n}{list}\PYG{p}{)}
\PYG{p}{\PYGZob{}}
\PYG{+w}{  }\PYG{k}{auto}\PYG{+w}{ }\PYG{n}{buf}\PYG{+w}{ }\PYG{o}{=}\PYG{+w}{ }\PYG{p}{(}\PYG{k+kt}{double}\PYG{+w}{ }\PYG{o}{*}\PYG{p}{)}\PYG{+w}{ }\PYG{n}{vbuf}\PYG{p}{;}
\PYG{+w}{  }\PYG{k+kt}{double}\PYG{+w}{ }\PYG{o}{*}\PYG{n}{src}\PYG{+w}{ }\PYG{o}{=}\PYG{+w}{ }\PYG{o}{\PYGZam{}}\PYG{n}{T\PYGZus{}electron}\PYG{p}{[}\PYG{n}{nzlo\PYGZus{}out}\PYG{p}{]}\PYG{p}{[}\PYG{n}{nylo\PYGZus{}out}\PYG{p}{]}\PYG{p}{[}\PYG{n}{nxlo\PYGZus{}out}\PYG{p}{]}\PYG{p}{;}
\PYG{+w}{  }\PYG{k}{for}\PYG{+w}{ }\PYG{p}{(}\PYG{k+kt}{int}\PYG{+w}{ }\PYG{n}{i}\PYG{+w}{ }\PYG{o}{=}\PYG{+w}{ }\PYG{l+m+mi}{0}\PYG{p}{;}\PYG{+w}{ }\PYG{n}{i}\PYG{+w}{ }\PYG{o}{\PYGZlt{}}\PYG{+w}{ }\PYG{n}{nlist}\PYG{p}{;}\PYG{+w}{ }\PYG{n}{i}\PYG{o}{+}\PYG{o}{+}\PYG{p}{)}\PYG{+w}{ }\PYG{n}{buf}\PYG{p}{[}\PYG{n}{i}\PYG{p}{]}\PYG{+w}{ }\PYG{o}{=}\PYG{+w}{ }\PYG{n}{src}\PYG{p}{[}\PYG{n}{list}\PYG{p}{[}\PYG{n}{i}\PYG{p}{]}\PYG{p}{]}\PYG{p}{;}
\PYG{p}{\PYGZcb{}}
\end{sphinxVerbatim}

\sphinxAtStartPar
In this case, the \sphinxstyleemphasis{which} argument is not used, \sphinxstyleemphasis{vbuf} points to a
buffer of doubles, and the electron temperature is stored by the
FixTTMGrid class in a 3d array of owned+ghost cells called T\_electron.
That array is allocated by the \sphinxstyleemphasis{memory\sphinxhyphen{}\textgreater{}create\_3d\_offset()} method
described above so that the first grid cell it stores is indexed as
T\_electron{[}nzlo\_out{]}{[}nylo\_out{]}{[}nxlo\_out{]}.  The \sphinxstyleemphasis{nlist} values in
\sphinxstyleemphasis{list} are integer offsets from that first grid cell.  Setting \sphinxstyleemphasis{src}
to the address of the first cell allows those offsets to be used to
access the temperatures to pack into the buffer.

\sphinxAtStartPar
Here is a similar portion of code from the src/fix\_ave\_grid.cpp class
which can store two kinds of data, a scalar count of atoms in a grid
cell, and one or more grid\sphinxhyphen{}cell\sphinxhyphen{}averaged atom properties.  The code
from its \sphinxstyleemphasis{unpack\_reverse\_grid()} function for 2d grids and multiple
per\sphinxhyphen{}atom properties per grid cell (\sphinxstyleemphasis{nvalues}) is shown here:

\begin{sphinxVerbatim}[commandchars=\\\{\}]
\PYG{k+kt}{void}\PYG{+w}{ }\PYG{n+nf}{FixAveGrid::unpack\PYGZus{}reverse\PYGZus{}grid}\PYG{p}{(}\PYG{k+kt}{int}\PYG{+w}{ }\PYG{c+cm}{/*which*/}\PYG{p}{,}\PYG{+w}{ }\PYG{k+kt}{void}\PYG{+w}{ }\PYG{o}{*}\PYG{n}{vbuf}\PYG{p}{,}\PYG{+w}{ }\PYG{k+kt}{int}\PYG{+w}{ }\PYG{n}{nlist}\PYG{p}{,}\PYG{+w}{ }\PYG{k+kt}{int}\PYG{+w}{ }\PYG{o}{*}\PYG{n}{list}\PYG{p}{)}
\PYG{p}{\PYGZob{}}
\PYG{+w}{  }\PYG{k}{auto}\PYG{+w}{ }\PYG{n}{buf}\PYG{+w}{ }\PYG{o}{=}\PYG{+w}{ }\PYG{p}{(}\PYG{k+kt}{double}\PYG{+w}{ }\PYG{o}{*}\PYG{p}{)}\PYG{+w}{ }\PYG{n}{vbuf}\PYG{p}{;}
\PYG{+w}{  }\PYG{k+kt}{double}\PYG{+w}{ }\PYG{o}{*}\PYG{n}{count}\PYG{p}{,}\PYG{o}{*}\PYG{n}{data}\PYG{p}{,}\PYG{o}{*}\PYG{n}{values}\PYG{p}{;}
\PYG{+w}{  }\PYG{n}{count}\PYG{+w}{ }\PYG{o}{=}\PYG{+w}{ }\PYG{o}{\PYGZam{}}\PYG{n}{count2d}\PYG{p}{[}\PYG{n}{nylo\PYGZus{}out}\PYG{p}{]}\PYG{p}{[}\PYG{n}{nxlo\PYGZus{}out}\PYG{p}{]}\PYG{p}{;}
\PYG{+w}{  }\PYG{n}{data}\PYG{+w}{ }\PYG{o}{=}\PYG{+w}{ }\PYG{o}{\PYGZam{}}\PYG{n}{array2d}\PYG{p}{[}\PYG{n}{nylo\PYGZus{}out}\PYG{p}{]}\PYG{p}{[}\PYG{n}{nxlo\PYGZus{}out}\PYG{p}{]}\PYG{p}{[}\PYG{l+m+mi}{0}\PYG{p}{]}\PYG{p}{;}
\PYG{+w}{  }\PYG{n}{m}\PYG{+w}{ }\PYG{o}{=}\PYG{+w}{ }\PYG{l+m+mi}{0}\PYG{p}{;}
\PYG{+w}{  }\PYG{k}{for}\PYG{+w}{ }\PYG{p}{(}\PYG{n}{i}\PYG{+w}{ }\PYG{o}{=}\PYG{+w}{ }\PYG{l+m+mi}{0}\PYG{p}{;}\PYG{+w}{ }\PYG{n}{i}\PYG{+w}{ }\PYG{o}{\PYGZlt{}}\PYG{+w}{ }\PYG{n}{nlist}\PYG{p}{;}\PYG{+w}{ }\PYG{n}{i}\PYG{o}{+}\PYG{o}{+}\PYG{p}{)}\PYG{+w}{ }\PYG{p}{\PYGZob{}}
\PYG{+w}{    }\PYG{n}{count}\PYG{p}{[}\PYG{n}{list}\PYG{p}{[}\PYG{n}{i}\PYG{p}{]}\PYG{p}{]}\PYG{+w}{ }\PYG{o}{+}\PYG{o}{=}\PYG{+w}{ }\PYG{n}{buf}\PYG{p}{[}\PYG{n}{m}\PYG{o}{+}\PYG{o}{+}\PYG{p}{]}\PYG{p}{;}
\PYG{+w}{    }\PYG{n}{values}\PYG{+w}{ }\PYG{o}{=}\PYG{+w}{ }\PYG{o}{\PYGZam{}}\PYG{n}{data}\PYG{p}{[}\PYG{n}{nvalues}\PYG{o}{*}\PYG{n}{list}\PYG{p}{[}\PYG{n}{i}\PYG{p}{]}\PYG{p}{]}\PYG{p}{;}
\PYG{+w}{    }\PYG{k}{for}\PYG{+w}{ }\PYG{p}{(}\PYG{n}{j}\PYG{+w}{ }\PYG{o}{=}\PYG{+w}{ }\PYG{l+m+mi}{0}\PYG{p}{;}\PYG{+w}{ }\PYG{n}{j}\PYG{+w}{ }\PYG{o}{\PYGZlt{}}\PYG{+w}{ }\PYG{n}{nvalues}\PYG{p}{;}\PYG{+w}{ }\PYG{n}{j}\PYG{o}{+}\PYG{o}{+}\PYG{p}{)}
\PYG{+w}{     }\PYG{n}{values}\PYG{p}{[}\PYG{n}{j}\PYG{p}{]}\PYG{+w}{ }\PYG{o}{+}\PYG{o}{=}\PYG{+w}{ }\PYG{n}{buf}\PYG{p}{[}\PYG{n}{m}\PYG{o}{+}\PYG{o}{+}\PYG{p}{]}\PYG{p}{;}
\PYG{+w}{  }\PYG{p}{\PYGZcb{}}
\PYG{p}{\PYGZcb{}}
\end{sphinxVerbatim}

\sphinxAtStartPar
Both the count and the multiple values per grid cell are communicated
in \sphinxstyleemphasis{vbuf}.  Note that \sphinxstyleemphasis{data} is now a pointer to the first value in
the first grid cell.  And \sphinxstyleemphasis{values} points to where the first value in
\sphinxstyleemphasis{data} is for an offset of grid cells, calculated by multiplying
\sphinxstyleemphasis{nvalues} by \sphinxstyleemphasis{list{[}i{]}}.  Finally, because this is reverse
communication, the communicated buffer values are summed to the caller
values.

\sphinxAtStartPar
The \sphinxstyleemphasis{ghost\_adjacent()} method returns a 1 if every processor can
perform the necessary owned/ghost communication with only its nearest
neighbor processors (4 in 2d, 6 in 3d).  It returns a 0 if any
processor’s ghost cells extend further than nearest neighbor
processors.

\sphinxAtStartPar
This can be checked by callers who have the option to change the
global grid size to ensure more efficient nearest\sphinxhyphen{}neighbor\sphinxhyphen{}only
communication if they wish.  In this case, they instantiate a grid of
a given size (resolution), then invoke \sphinxstyleemphasis{setup\_comm()} followed by
\sphinxstyleemphasis{ghost\_adjacent()}.  If the ghost cells are not adjacent, they destroy
the grid instance and start over with a higher\sphinxhyphen{}resolution grid.
Several of the {\hyperref[\detokenize{kspace_style::doc}]{\sphinxcrossref{\DUrole{doc}{kspace\_style pppm}}}} command
variants have this option.


\bigskip\hrule\bigskip



\subsection{Grid class remap methods for load balancing}
\label{\detokenize{Developer_grid:grid-class-remap-methods-for-load-balancing}}
\sphinxAtStartPar
The following methods are used when a load\sphinxhyphen{}balancing operation,
triggered by the {\hyperref[\detokenize{balance::doc}]{\sphinxcrossref{\DUrole{doc}{balance}}}} or {\hyperref[\detokenize{fix_balance::doc}]{\sphinxcrossref{\DUrole{doc}{fix balance}}}} commands, changes the partitioning of the simulation
domain into processor subdomains.

\sphinxAtStartPar
In order to work with load\sphinxhyphen{}balancing, any style command (compute, fix,
pair, or kspace style) which allocates a grid and stores per\sphinxhyphen{}grid data
should define a \sphinxstyleemphasis{reset\_grid()} method; it takes no arguments.  It will
be called by the two balance commands after they have reset processor
subdomains and migrated atoms (particles) to new owning processors.
The \sphinxstyleemphasis{reset\_grid()} method will typically perform some or all of the
following operations.  See the src/fix\_ave\_grid.cpp and
src/EXTRA\_FIX/fix\_ttm\_grid.cpp files for examples of \sphinxstyleemphasis{reset\_grid()}
methods, as well as the \sphinxstyleemphasis{pack\_remap\_grid()} and \sphinxstyleemphasis{unpack\_remap\_grid()}
functions.

\sphinxAtStartPar
First, the \sphinxstyleemphasis{reset\_grid()} method can instantiate new grid(s) of the
same global size, then call \sphinxstyleemphasis{setup\_grid()} to partition them via the
new processor subdomains.  At this point, it can invoke the
\sphinxstyleemphasis{identical()} method which compares the owned and ghost grid cell
index bounds between two grids, the old grid passed as a pointer
argument, and the new grid whose \sphinxstyleemphasis{identical()} method is being called.
It returns 1 if the indices match on all processors, otherwise 0.  If
they all match, then the new grids can be deleted; the command can
continue to use the old grids.

\sphinxAtStartPar
If not, then the command should allocate new grid data array(s) which
depend on the new partitioning.  If the command does not need to
persist its grid data from the old partitioning to the new one, then
the command can simply delete the old data array(s) and grid
instance(s).  It can then return.

\sphinxAtStartPar
If the grid data does need to persist, then the data for each grid
needs to be “remapped” from the old grid partitioning to the new grid
partitioning.  The \sphinxstyleemphasis{setup\_remap()} and \sphinxstyleemphasis{remap()} methods are used for
that purpose.

\begin{sphinxVerbatim}[commandchars=\\\{\}]
\PYG{k+kt}{int}\PYG{+w}{ }\PYG{n+nf}{identical}\PYG{p}{(}\PYG{n}{Grid3d}\PYG{+w}{ }\PYG{o}{*}\PYG{n}{old}\PYG{p}{)}\PYG{p}{;}
\PYG{k+kt}{void}\PYG{+w}{ }\PYG{n}{setup\PYGZus{}remap}\PYG{p}{(}\PYG{n}{Grid3d}\PYG{+w}{ }\PYG{o}{*}\PYG{n}{old}\PYG{p}{,}\PYG{+w}{ }\PYG{k+kt}{int}\PYG{+w}{ }\PYG{o}{\PYGZam{}}\PYG{n}{nremap\PYGZus{}buf1}\PYG{p}{,}\PYG{+w}{ }\PYG{k+kt}{int}\PYG{+w}{ }\PYG{o}{\PYGZam{}}\PYG{n}{nremap\PYGZus{}buf2}\PYG{p}{)}
\PYG{k+kt}{void}\PYG{+w}{ }\PYG{n}{remap}\PYG{p}{(}\PYG{k+kt}{int}\PYG{+w}{ }\PYG{n}{caller}\PYG{p}{,}\PYG{+w}{ }\PYG{k+kt}{void}\PYG{+w}{ }\PYG{o}{*}\PYG{n}{ptr}\PYG{p}{,}\PYG{+w}{ }\PYG{k+kt}{int}\PYG{+w}{ }\PYG{n}{which}\PYG{p}{,}\PYG{+w}{ }\PYG{k+kt}{int}\PYG{+w}{ }\PYG{n}{nper}\PYG{p}{,}\PYG{+w}{ }\PYG{k+kt}{int}\PYG{+w}{ }\PYG{n}{nbyte}\PYG{p}{,}
\PYG{+w}{           }\PYG{k+kt}{void}\PYG{+w}{ }\PYG{o}{*}\PYG{n}{buf1}\PYG{p}{,}\PYG{+w}{ }\PYG{k+kt}{void}\PYG{+w}{ }\PYG{o}{*}\PYG{n}{buf2}\PYG{p}{,}\PYG{+w}{ }\PYG{n}{MPI\PYGZus{}Datatype}\PYG{+w}{ }\PYG{n}{datatype}\PYG{p}{)}
\end{sphinxVerbatim}

\sphinxAtStartPar
The arguments to these methods are identical to those for
the \sphinxstyleemphasis{setup\_comm()} and \sphinxstyleemphasis{forward\_comm()} or \sphinxstyleemphasis{reverse\_comm()} methods.
However the returned \sphinxstyleemphasis{nremap\_buf1} and \sphinxstyleemphasis{nremap2\_buf} values will be
different than the \sphinxstyleemphasis{nbuf1} and \sphinxstyleemphasis{nbuf2} values.  They should be used to
allocate two different remap buffers, separate from the owned/ghost
communication buffers.

\sphinxAtStartPar
To use the \sphinxstyleemphasis{remap()} method, the caller must provide two
callback functions:

\begin{sphinxVerbatim}[commandchars=\\\{\}]
\PYG{k+kt}{void}\PYG{+w}{ }\PYG{n+nf}{pack\PYGZus{}remap\PYGZus{}grid}\PYG{p}{(}\PYG{k+kt}{int}\PYG{+w}{ }\PYG{n}{which}\PYG{p}{,}\PYG{+w}{ }\PYG{k+kt}{void}\PYG{+w}{ }\PYG{o}{*}\PYG{n}{vbuf}\PYG{p}{,}\PYG{+w}{ }\PYG{k+kt}{int}\PYG{+w}{ }\PYG{n}{nlist}\PYG{p}{,}\PYG{+w}{ }\PYG{k+kt}{int}\PYG{+w}{ }\PYG{o}{*}\PYG{n}{list}\PYG{p}{)}\PYG{p}{;}
\PYG{k+kt}{void}\PYG{+w}{ }\PYG{n+nf}{unpack\PYGZus{}remap\PYGZus{}grid}\PYG{p}{(}\PYG{k+kt}{int}\PYG{+w}{ }\PYG{n}{which}\PYG{p}{,}\PYG{+w}{ }\PYG{k+kt}{void}\PYG{+w}{ }\PYG{o}{*}\PYG{n}{vbuf}\PYG{p}{,}\PYG{+w}{ }\PYG{k+kt}{int}\PYG{+w}{ }\PYG{n}{list}\PYG{p}{,}\PYG{+w}{ }\PYG{k+kt}{int}\PYG{+w}{ }\PYG{o}{*}\PYG{n}{list}\PYG{p}{)}\PYG{p}{;}
\end{sphinxVerbatim}

\sphinxAtStartPar
Their arguments are identical to those for the \sphinxstyleemphasis{pack\_forward\_grid()}
and \sphinxstyleemphasis{unpack\_forward\_grid()} callback functions (or the reverse
variants) discussed above.  Normally, both these methods pack/unpack
all the data arrays for a given grid.  The \sphinxstyleemphasis{which} argument of the
\sphinxstyleemphasis{remap()} method sets the \sphinxstyleemphasis{which} value for the pack/unpack functions.
If the command instantiates multiple grids (of different sizes), it
can be used within the pack/unpack methods to select which grid’s data
is being remapped.

\sphinxAtStartPar
Note that the \sphinxstyleemphasis{pack\_remap\_grid()} function must copy values from the
OLD grid data arrays into the \sphinxstyleemphasis{vbuf} buffer. The \sphinxstyleemphasis{unpack\_remap\_grid()}
function must copy values from the \sphinxstyleemphasis{vbuf} buffer into the NEW grid
data arrays.

\sphinxAtStartPar
After the remap operation for grid cell data has been performed, the
\sphinxstyleemphasis{reset\_grid()} method can deallocate the two remap buffers it created,
and can then exit.


\bigskip\hrule\bigskip



\subsection{Grid class I/O methods}
\label{\detokenize{Developer_grid:grid-class-i-o-methods}}
\sphinxAtStartPar
There are two I/O methods in the Grid classes which can be used to
read and write grid cell data to files.  The caller can decide on the
precise format of each file, e.g. whether header lines are prepended
or comment lines are allowed.  Fundamentally, the file should contain
one line per grid cell for the entire global grid.  Each line should
contain identifying info as to which grid cell it is, e.g. a unique
grid cell ID or the ix,iy,iz indices of the cell within a 3d grid.
The line should also contain one or more data values which are stored
within the grid data arrays created by the command

\sphinxAtStartPar
For grid cell IDs, the LAMMPS convention is that the IDs run from 1 to
N, where N = Nx * Ny for 2d grids and N = Nx * Ny * Nz for 3d grids.
The x\sphinxhyphen{}index of the grid cell varies fastest, then y, and the z\sphinxhyphen{}index
varies slowest.  So for a 10x10x10 grid the cell IDs from 901\sphinxhyphen{}1000
would be in the top xy layer of the z dimension.

\sphinxAtStartPar
The \sphinxstyleemphasis{read\_file()} method does something simple.  It reads a chunk of
consecutive lines from the file and passes them back to the caller to
process.  The caller provides a \sphinxstyleemphasis{unpack\_read\_grid()} function for this
purpose.  The function checks the grid cell ID or indices and only
stores grid cell data for the grid cells it owns.

\sphinxAtStartPar
The \sphinxstyleemphasis{write\_file()} method does something slightly more complex.  Each
processor packs the data for its owned grid cells into a buffer.  The
caller provides a \sphinxstyleemphasis{pack\_write\_grid()} function for this purpose.  The
\sphinxstyleemphasis{write\_file()} method then loops over all processors and each sends
its buffer one at a time to processor 0, along with the 3d (or 2d)
index bounds of its grid cell data within the global grid.  Processor
0 calls back to the \sphinxstyleemphasis{unpack\_write\_grid()} function provided by the
caller with the buffer.  The function writes one line per grid cell to
the file.

\sphinxAtStartPar
See the src/EXTRA\_FIX/fix\_ttm\_grid.cpp file for examples of now both
these methods are used to read/write electron temperature values
from/to a file, as well as for implementations of the the pack/unpack
functions described below.

\sphinxAtStartPar
Here are the details of the two I/O methods and the 3 callback
functions.  See the src/fix\_ave\_grid.cpp file for examples of all of
them.

\begin{sphinxVerbatim}[commandchars=\\\{\}]
\PYG{k+kt}{void}\PYG{+w}{ }\PYG{n}{read\PYGZus{}file}\PYG{p}{(}\PYG{k+kt}{int}\PYG{+w}{ }\PYG{n}{caller}\PYG{p}{,}\PYG{+w}{ }\PYG{k+kt}{void}\PYG{+w}{ }\PYG{o}{*}\PYG{n}{ptr}\PYG{p}{,}\PYG{+w}{ }\PYG{k+kt}{FILE}\PYG{+w}{ }\PYG{o}{*}\PYG{n}{fp}\PYG{p}{,}\PYG{+w}{ }\PYG{k+kt}{int}\PYG{+w}{ }\PYG{n}{nchunk}\PYG{p}{,}\PYG{+w}{ }\PYG{k+kt}{int}\PYG{+w}{ }\PYG{n}{maxline}\PYG{p}{)}
\PYG{k+kt}{void}\PYG{+w}{ }\PYG{n}{write\PYGZus{}file}\PYG{p}{(}\PYG{k+kt}{int}\PYG{+w}{ }\PYG{n}{caller}\PYG{p}{,}\PYG{+w}{ }\PYG{k+kt}{void}\PYG{+w}{ }\PYG{o}{*}\PYG{n}{ptr}\PYG{p}{,}\PYG{+w}{ }\PYG{k+kt}{int}\PYG{+w}{ }\PYG{n}{which}\PYG{p}{,}
\PYG{+w}{                }\PYG{k+kt}{int}\PYG{+w}{ }\PYG{n}{nper}\PYG{p}{,}\PYG{+w}{ }\PYG{k+kt}{int}\PYG{+w}{ }\PYG{n}{nbyte}\PYG{p}{,}\PYG{+w}{ }\PYG{n}{MPI\PYGZus{}Datatype}\PYG{+w}{ }\PYG{n}{datatype}
\end{sphinxVerbatim}

\sphinxAtStartPar
The \sphinxstyleemphasis{caller} argument in both methods should be one of these values \textendash{}
Grid3d::COMPUTE, Grid3d::FIX, Grid3d::KSPACE, Grid3d::PAIR \textendash{}
depending on the style of the caller class.  The \sphinxstyleemphasis{ptr} argument in
both methods is the “this” pointer to the caller class.  These 2
arguments are used to call back to pack()/unpack() functions in the
caller class, as explained below.

\sphinxAtStartPar
For the \sphinxstyleemphasis{read\_file()} method, the \sphinxstyleemphasis{fp} argument is a file pointer to
the file to be read from, opened on processor 0 by the caller.
\sphinxstyleemphasis{Nchunk} is the number of lines to read per chunk, and \sphinxstyleemphasis{maxline} is
the maximum number of characters per line.  The Grid class will
allocate a buffer for storing chunks of lines based on these values.

\sphinxAtStartPar
For the \sphinxstyleemphasis{write\_file()} method, the \sphinxstyleemphasis{which} argument is a flag the
caller can set which is passed back to the caller’s pack()/unpack()
methods.  If the command instantiates multiple grids (of different
sizes), this flag can be used within the pack/unpack methods to select
which grid’s data is being written out (presumably to different
files).  the \sphinxstyleemphasis{nper} argument is the number of values per grid cell to
be written out.  The \sphinxstyleemphasis{nbyte} argument is the number of bytes per
value, e.g. 8 for double\sphinxhyphen{}precision values.  The \sphinxstyleemphasis{datatype} argument is
the MPI\_Datatype setting, which should match the \sphinxstyleemphasis{nbyte} argument.
E.g. MPI\_DOUBLE for double precision values.

\sphinxAtStartPar
To use the \sphinxstyleemphasis{read\_grid()} method, the caller must provide one callback
function.  To use the \sphinxstyleemphasis{write\_grid()} method, it provides two callback
functions:

\begin{sphinxVerbatim}[commandchars=\\\{\}]
\PYG{k+kt}{int}\PYG{+w}{ }\PYG{n}{unpack\PYGZus{}read\PYGZus{}grid}\PYG{p}{(}\PYG{k+kt}{int}\PYG{+w}{ }\PYG{n}{nlines}\PYG{p}{,}\PYG{+w}{ }\PYG{k+kt}{char}\PYG{+w}{ }\PYG{o}{*}\PYG{n}{buffer}\PYG{p}{)}
\PYG{k+kt}{void}\PYG{+w}{ }\PYG{n}{pack\PYGZus{}write\PYGZus{}grid}\PYG{p}{(}\PYG{k+kt}{int}\PYG{+w}{ }\PYG{n}{which}\PYG{p}{,}\PYG{+w}{ }\PYG{k+kt}{void}\PYG{+w}{ }\PYG{o}{*}\PYG{n}{vbuf}\PYG{p}{)}
\PYG{k+kt}{void}\PYG{+w}{ }\PYG{n}{unpack\PYGZus{}write\PYGZus{}grid}\PYG{p}{(}\PYG{k+kt}{int}\PYG{+w}{ }\PYG{n}{which}\PYG{p}{,}\PYG{+w}{ }\PYG{k+kt}{void}\PYG{+w}{ }\PYG{o}{*}\PYG{n}{vbuf}\PYG{p}{,}\PYG{+w}{ }\PYG{k+kt}{int}\PYG{+w}{ }\PYG{o}{*}\PYG{n}{bounds}\PYG{p}{)}
\end{sphinxVerbatim}

\sphinxAtStartPar
For \sphinxstyleemphasis{unpack\_read\_grid()} the \sphinxstyleemphasis{nlines} argument is the number of lines
of character data read from the file and contained in \sphinxstyleemphasis{buffer}.  The
lines each include a newline character at the end.  When the function
processes the lines, it may choose to skip some of them (header or
comment lines).  It returns an integer count of the number of grid
cell lines it processed.  This enables the Grid class \sphinxstyleemphasis{read\_file()}
method to know when it has read the correct number of lines.

\sphinxAtStartPar
For \sphinxstyleemphasis{pack\_write\_grid()} and \sphinxstyleemphasis{unpack\_write\_grid()}, the \sphinxstyleemphasis{vbuf} argument
is the buffer to pack/unpack data to/from.  It is a void pointer so
that the caller can cast it to whatever data type it chooses,
e.g. double precision values.  the \sphinxstyleemphasis{which} argument is set to the
\sphinxstyleemphasis{which} value of the \sphinxstyleemphasis{write\_file()} method.  It allows the caller to
choose which grid data to operate on.

\sphinxAtStartPar
For \sphinxstyleemphasis{unpack\_write\_grid()}, the \sphinxstyleemphasis{bounds} argument is a vector of 4 or 6
integer grid indices (4 for 2d, 6 for 3d).  They are the
xlo,xhi,ylo,yhi,zlo,zhi index bounds of the portion of the global grid
which the \sphinxstyleemphasis{vbuf} holds owned grid cell data values for.  The caller
should loop over the values in \sphinxstyleemphasis{vbuf} with a double loop (2d) or
triple loop (3d), similar to the code snippets listed above.  The
x\sphinxhyphen{}index varies fastest, then y, and the z\sphinxhyphen{}index slowest.  If there are
multiple values per grid cell, the index for those values varies
fastest of all.  The caller can add the x,y,z indices of the grid cell
(or the corresponding grid cell ID) to the data value(s) written as
one line to the output file.


\bigskip\hrule\bigskip



\subsection{Style class grid access methods}
\label{\detokenize{Developer_grid:style-class-grid-access-methods}}
\sphinxAtStartPar
A style command can enable its grid cell data to be accessible from
other commands.  For example {\hyperref[\detokenize{fix_ave_grid::doc}]{\sphinxcrossref{\DUrole{doc}{fix ave/grid}}}} or
{\hyperref[\detokenize{dump::doc}]{\sphinxcrossref{\DUrole{doc}{dump grid}}}} or {\hyperref[\detokenize{dump::doc}]{\sphinxcrossref{\DUrole{doc}{dump grid/vtk}}}}.  Those
commands access the grid cell data by using a \sphinxstyleemphasis{grid reference} in
their input script syntax, as described on the {\hyperref[\detokenize{Howto_grid::doc}]{\sphinxcrossref{\DUrole{doc}{Howto\_grid}}}} doc page.  They look like this:
\begin{itemize}
\item {} 
\sphinxAtStartPar
c\_ID:gname:dname

\item {} 
\sphinxAtStartPar
c\_ID:gname:dname{[}I{]}

\item {} 
\sphinxAtStartPar
f\_ID:gname:dname

\item {} 
\sphinxAtStartPar
f\_ID:gname:dname{[}I{]}

\end{itemize}

\sphinxAtStartPar
Each grid command instantiates has a unique \sphinxstyleemphasis{gname}, defined by the
command.  Likewise each grid cell data structure (scalar or vector)
associated with a grid has a unique \sphinxstyleemphasis{dname}, also defined by the
command.

\sphinxAtStartPar
To provide access to its grid cell data, a style command needs to
implement the following 4 methods:

\begin{sphinxVerbatim}[commandchars=\\\{\}]
\PYG{k+kt}{int}\PYG{+w}{ }\PYG{n+nf}{get\PYGZus{}grid\PYGZus{}by\PYGZus{}name}\PYG{p}{(}\PYG{k}{const}\PYG{+w}{ }\PYG{n}{std}\PYG{o}{:}\PYG{o}{:}\PYG{n}{string}\PYG{+w}{ }\PYG{o}{\PYGZam{}}\PYG{n}{name}\PYG{p}{,}\PYG{+w}{ }\PYG{k+kt}{int}\PYG{+w}{ }\PYG{o}{\PYGZam{}}\PYG{n}{dim}\PYG{p}{)}\PYG{p}{;}
\PYG{k+kt}{void}\PYG{+w}{ }\PYG{o}{*}\PYG{n+nf}{get\PYGZus{}grid\PYGZus{}by\PYGZus{}index}\PYG{p}{(}\PYG{k+kt}{int}\PYG{+w}{ }\PYG{n}{index}\PYG{p}{)}\PYG{p}{;}
\PYG{k+kt}{int}\PYG{+w}{ }\PYG{n+nf}{get\PYGZus{}griddata\PYGZus{}by\PYGZus{}name}\PYG{p}{(}\PYG{k+kt}{int}\PYG{+w}{ }\PYG{n}{igrid}\PYG{p}{,}\PYG{+w}{ }\PYG{k}{const}\PYG{+w}{ }\PYG{n}{std}\PYG{o}{:}\PYG{o}{:}\PYG{n}{string}\PYG{+w}{ }\PYG{o}{\PYGZam{}}\PYG{n}{name}\PYG{p}{,}\PYG{+w}{ }\PYG{k+kt}{int}\PYG{+w}{ }\PYG{o}{\PYGZam{}}\PYG{n}{ncol}\PYG{p}{)}\PYG{p}{;}
\PYG{k+kt}{void}\PYG{+w}{ }\PYG{o}{*}\PYG{n+nf}{get\PYGZus{}griddata\PYGZus{}by\PYGZus{}index}\PYG{p}{(}\PYG{k+kt}{int}\PYG{+w}{ }\PYG{n}{index}\PYG{p}{)}\PYG{p}{;}
\end{sphinxVerbatim}

\sphinxAtStartPar
Currently only computes and fixes can implement these methods.  If it
does so, the compute of fix should also set the variable
\sphinxstyleemphasis{pergrid\_flag} to 1.  See any of the compute or fix commands which set
“pergrid\_flag = 1” for examples of how these 4 functions can be
implemented.

\sphinxAtStartPar
The \sphinxstyleemphasis{get\_grid\_by\_name()} method takes a grid name as input and returns
two values.  The \sphinxstyleemphasis{dim} argument is returned as 2 or 3 for the
dimensionality of the grid.  The function return is a grid index from
0 to G\sphinxhyphen{}1 where \sphinxstyleemphasis{G} is the number of grids the command instantiates.  A
value of \sphinxhyphen{}1 is returned if the grid name is not recognized.

\sphinxAtStartPar
The \sphinxstyleemphasis{get\_grid\_by\_index()} method is called after the
\sphinxstyleemphasis{get\_grid\_by\_name()} method, using the grid index it returned as its
argument.  This method will return a pointer to the Grid2d or Grid3d
class.  The caller can use this to query grid attributes, such as the
global size of the grid, to ensure it is of the expected size.

\sphinxAtStartPar
The \sphinxstyleemphasis{get\_griddata\_by\_name()} method takes a grid index \sphinxstyleemphasis{igrid} and a
data name as input.  It returns two values.  The \sphinxstyleemphasis{ncol} argument is
returned as a 0 if the grid data is a single value (scalar) per grid
cell, or an integer M \textgreater{} 0 if there are M values (vector) per grid
cell.  Note that even if M = 1, it is still a 1\sphinxhyphen{}length vector, not a
scalar.  The function return is a data index from 0 to D\sphinxhyphen{}1 where \sphinxstyleemphasis{D}
is the number of data sets associated with that grid by the command.
A value of \sphinxhyphen{}1 is returned if the data name is not recognized.

\sphinxAtStartPar
The \sphinxstyleemphasis{get\_griddata\_by\_index()} method is called after the
\sphinxstyleemphasis{get\_griddata\_by\_name()} method, using the data index it returned as
its argument.  This method will return a pointer to the
multidimensional array which stores the requested data.

\sphinxAtStartPar
As in the discussion above of the Memory class \sphinxstyleemphasis{create\_offset()}
methods, the dimensionality of the array associated with the returned
pointer depends on whether it is a 2d or 3d grid and whether there is
a single or multiple values stored for each grid cell:
\begin{itemize}
\item {} 
\sphinxAtStartPar
single value per cell for a 2d grid = 2d array pointer

\item {} 
\sphinxAtStartPar
multiple values per cell for a 2d grid = 3d array pointer

\item {} 
\sphinxAtStartPar
single value per cell for a 3d grid = 3d array pointer

\item {} 
\sphinxAtStartPar
multiple values per cell for a 3d grid = 4d array pointer

\end{itemize}

\sphinxAtStartPar
The caller will typically access the data by casting the void pointer
to the corresponding array pointer and using nested loops in x,y,z
between owned or ghost index bounds returned by the
\sphinxstyleemphasis{get\_bounds\_owned()} or \sphinxstyleemphasis{get\_bounds\_ghost()} methods to index into the
array.  Example code snippets with this logic were listed above,


\bigskip\hrule\bigskip



\subsection{Final notes}
\label{\detokenize{Developer_grid:final-notes}}
\sphinxAtStartPar
Finally, here are some additional issues to pay attention to for
writing any style command which uses distributed grids via the Grid2d
or Grid3d class.

\sphinxAtStartPar
The command destructor should delete all instances of the Grid class,
any buffers it allocated for forward/reverse or remap communication,
and any data arrays it allocated to store grid cell data.

\sphinxAtStartPar
If a command is intended to work for either 2d or 3d simulations, then
it should have logic to instantiate either 2d or 3d grids and their
associated data arrays, depending on the dimension of the simulation
box.  The {\hyperref[\detokenize{fix_ave_grid::doc}]{\sphinxcrossref{\DUrole{doc}{fix ave/grid}}}} command is an example of
such a command.

\sphinxAtStartPar
When a command maps its particles to the grid and updates grid cell
values, it should check that it is not updating or accessing a grid
cell value outside the range of its owned+ghost cells, and generate an
error message if that is the case.  This could happen, for example, if
a particle has moved further than half the neighbor skin distance,
because the neighbor list update criterion are not adequate to prevent
it from happening.  See the src/KSPACE/pppm.cpp file and its
\sphinxstyleemphasis{particle\_map()} method for an example of this kind of error check.


\part{Command Reference}
\label{\detokenize{Manual:command-reference}}\phantomsection\label{\detokenize{Manual:id1}}
\sphinxstepscope


\chapter{Commands}
\label{\detokenize{commands_list:commands}}\label{\detokenize{commands_list::doc}}
\sphinxstepscope

\index{angle\_coeff@\spxentry{angle\_coeff}}\ignorespaces 

\section{angle\_coeff command}
\label{\detokenize{angle_coeff:angle-coeff-command}}\label{\detokenize{angle_coeff:index-0}}\label{\detokenize{angle_coeff::doc}}

\subsection{Syntax}
\label{\detokenize{angle_coeff:syntax}}
\begin{sphinxVerbatim}[commandchars=\\\{\}]
\PYG{k}{angle\PYGZus{}coeff}\PYG{+w}{ }\PYG{n}{N}\PYG{+w}{ }\PYG{n}{args}
\end{sphinxVerbatim}
\begin{itemize}
\item {} 
\sphinxAtStartPar
N = numeric angle type (see asterisk form below), or type label

\item {} 
\sphinxAtStartPar
args = coefficients for one or more angle types

\end{itemize}


\subsection{Examples}
\label{\detokenize{angle_coeff:examples}}
\begin{sphinxVerbatim}[commandchars=\\\{\}]
\PYG{k}{angle\PYGZus{}coeff}\PYG{+w}{ }\PYG{l+m}{1}\PYG{+w}{ }\PYG{l+m}{300.0}\PYG{+w}{ }\PYG{l+m}{107.0}
\PYG{k}{angle\PYGZus{}coeff}\PYG{+w}{ }\PYG{o}{*}\PYG{+w}{ }\PYG{l+m}{5.0}
\PYG{k}{angle\PYGZus{}coeff}\PYG{+w}{ }\PYG{l+m}{2}\PYG{o}{*}\PYG{l+m}{10}\PYG{+w}{ }\PYG{l+m}{5.0}

\PYG{k}{labelmap}\PYG{+w}{ }\PYG{n}{angle}\PYG{+w}{ }\PYG{l+m}{1}\PYG{+w}{ }\PYG{n}{hydroxyl}
\PYG{k}{angle\PYGZus{}coeff}\PYG{+w}{ }\PYG{n}{hydroxyl}\PYG{+w}{ }\PYG{l+m}{300.0}\PYG{+w}{ }\PYG{l+m}{107.0}
\end{sphinxVerbatim}


\subsection{Description}
\label{\detokenize{angle_coeff:description}}
\sphinxAtStartPar
Specify the angle force field coefficients for one or more angle types.
The number and meaning of the coefficients depends on the angle style.
Angle coefficients can also be set in the data file read by the
{\hyperref[\detokenize{read_data::doc}]{\sphinxcrossref{\DUrole{doc}{read\_data}}}} command or in a restart file.

\sphinxAtStartPar
\(N\) can be specified in one of two ways.  An explicit numeric
value can be used, as in the first example above.  Or \(N\) can be a
type label, which is an alphanumeric string defined by the
{\hyperref[\detokenize{labelmap::doc}]{\sphinxcrossref{\DUrole{doc}{labelmap}}}} command or in a section of a data file read
by the {\hyperref[\detokenize{read_data::doc}]{\sphinxcrossref{\DUrole{doc}{read\_data}}}} command.

\sphinxAtStartPar
For numeric values only, a wild\sphinxhyphen{}card asterisk can be used to set the
coefficients for multiple angle types.  This takes the form “*” or
“*n” or “n*” or “m*n”.  If \(N\) is the number of angle types,
then an asterisk with no numeric values means all types from 1 to
\(N\).  A leading asterisk means all types from 1 to n (inclusive).
A trailing asterisk means all types from n to \(N\) (inclusive).  A
middle asterisk means all types from m to n (inclusive).

\sphinxAtStartPar
Note that using an {\hyperref[\detokenize{angle_coeff::doc}]{\sphinxcrossref{\DUrole{doc}{angle\_coeff}}}} command can
override a previous setting for the same angle type.  For example,
these commands set the coeffs for all angle types, then overwrite the
coeffs for just angle type 2:

\begin{sphinxVerbatim}[commandchars=\\\{\}]
\PYG{k}{angle\PYGZus{}coeff}\PYG{+w}{ }\PYG{o}{*}\PYG{+w}{ }\PYG{l+m}{200.0}\PYG{+w}{ }\PYG{l+m}{107.0}\PYG{+w}{ }\PYG{l+m}{1.2}
\PYG{k}{angle\PYGZus{}coeff}\PYG{+w}{ }\PYG{l+m}{2}\PYG{+w}{ }\PYG{l+m}{50.0}\PYG{+w}{ }\PYG{l+m}{107.0}
\end{sphinxVerbatim}

\sphinxAtStartPar
A line in a data file that specifies angle coefficients uses the exact
same format as the arguments of the {\hyperref[\detokenize{angle_coeff::doc}]{\sphinxcrossref{\DUrole{doc}{angle\_coeff}}}}
command in an input script, except that wild\sphinxhyphen{}card asterisks should not
be used since coefficients for all \(N\) types must be listed in the
file.  For example, under the “Angle Coeffs” section of a data file, the
line that corresponds to the first example above would be listed as

\begin{sphinxVerbatim}[commandchars=\\\{\}]
1 300.0 107.0
\end{sphinxVerbatim}

\sphinxAtStartPar
The {\hyperref[\detokenize{angle_class2::doc}]{\sphinxcrossref{\DUrole{doc}{angle\_style class2}}}} is an exception to this
rule, in that an additional argument is used in the input script to
allow specification of the cross\sphinxhyphen{}term coefficients.  See its doc page
for details.


\bigskip\hrule\bigskip


\sphinxAtStartPar
The list of all angle styles defined in LAMMPS is given on the
{\hyperref[\detokenize{angle_style::doc}]{\sphinxcrossref{\DUrole{doc}{angle\_style}}}} doc page.  They are also listed in more
compact form on the {\hyperref[\detokenize{Commands_bond:angle}]{\sphinxcrossref{\DUrole{std,std-ref}{Commands angle}}}} doc page.

\sphinxAtStartPar
On either of those pages, click on the style to display the formula it
computes and its coefficients as specified by the associated
{\hyperref[\detokenize{angle_coeff::doc}]{\sphinxcrossref{\DUrole{doc}{angle\_coeff}}}} command.


\bigskip\hrule\bigskip



\subsection{Restrictions}
\label{\detokenize{angle_coeff:restrictions}}
\sphinxAtStartPar
This command must come after the simulation box is defined by a
{\hyperref[\detokenize{read_data::doc}]{\sphinxcrossref{\DUrole{doc}{read\_data}}}}, {\hyperref[\detokenize{read_restart::doc}]{\sphinxcrossref{\DUrole{doc}{read\_restart}}}}, or
{\hyperref[\detokenize{create_box::doc}]{\sphinxcrossref{\DUrole{doc}{create\_box}}}} command.

\sphinxAtStartPar
An angle style must be defined before any angle coefficients are
set, either in the input script or in a data file.


\subsection{Related commands}
\label{\detokenize{angle_coeff:related-commands}}
\sphinxAtStartPar
{\hyperref[\detokenize{angle_style::doc}]{\sphinxcrossref{\DUrole{doc}{angle\_style}}}}


\subsection{Default}
\label{\detokenize{angle_coeff:default}}
\sphinxAtStartPar
none

\sphinxstepscope

\index{angle\_style@\spxentry{angle\_style}}\ignorespaces 

\section{angle\_style command}
\label{\detokenize{angle_style:angle-style-command}}\label{\detokenize{angle_style:index-0}}\label{\detokenize{angle_style::doc}}

\subsection{Syntax}
\label{\detokenize{angle_style:syntax}}
\begin{sphinxVerbatim}[commandchars=\\\{\}]
\PYG{k}{angle\PYGZus{}style}\PYG{+w}{ }\PYG{n}{style}
\end{sphinxVerbatim}
\begin{itemize}
\item {} 
\sphinxAtStartPar
style = \sphinxstyleemphasis{none} or \sphinxstyleemphasis{zero} or \sphinxstyleemphasis{hybrid} or \sphinxstyleemphasis{amoeba} or \sphinxstyleemphasis{charmm} or \sphinxstyleemphasis{class2} or \sphinxstyleemphasis{class2/p6} or \sphinxstyleemphasis{cosine} or \sphinxstyleemphasis{cosine/buck6d} or \sphinxstyleemphasis{cosine/delta} or \sphinxstyleemphasis{cosine/periodic} or \sphinxstyleemphasis{cosine/shift} or \sphinxstyleemphasis{cosine/shift/exp} or \sphinxstyleemphasis{cosine/squared} or \sphinxstyleemphasis{cosine/squared/restricted} or \sphinxstyleemphasis{cross} or \sphinxstyleemphasis{dipole} or \sphinxstyleemphasis{fourier} or \sphinxstyleemphasis{fourier/simple} or \sphinxstyleemphasis{gaussian} or \sphinxstyleemphasis{harmonic} or \sphinxstyleemphasis{lepton} or \sphinxstyleemphasis{mm3} or \sphinxstyleemphasis{quartic} or \sphinxstyleemphasis{spica} or \sphinxstyleemphasis{table}

\end{itemize}


\subsection{Examples}
\label{\detokenize{angle_style:examples}}
\begin{sphinxVerbatim}[commandchars=\\\{\}]
\PYG{k}{angle\PYGZus{}style}\PYG{+w}{ }\PYG{n}{harmonic}
\PYG{k}{angle\PYGZus{}style}\PYG{+w}{ }\PYG{n}{charmm}
\PYG{k}{angle\PYGZus{}style}\PYG{+w}{ }\PYG{n}{hybrid}\PYG{+w}{ }\PYG{n}{harmonic}\PYG{+w}{ }\PYG{n}{cosine}
\end{sphinxVerbatim}


\subsection{Description}
\label{\detokenize{angle_style:description}}
\sphinxAtStartPar
Set the formula(s) LAMMPS uses to compute angle interactions between
triplets of atoms, which remain in force for the duration of the
simulation.  The list of angle triplets is read in by a
{\hyperref[\detokenize{read_data::doc}]{\sphinxcrossref{\DUrole{doc}{read\_data}}}} or {\hyperref[\detokenize{read_restart::doc}]{\sphinxcrossref{\DUrole{doc}{read\_restart}}}} command
from a data or restart file.

\sphinxAtStartPar
Hybrid models where angles are computed using different angle
potentials can be setup using the \sphinxstyleemphasis{hybrid} angle style.

\sphinxAtStartPar
The coefficients associated with a angle style can be specified in a
data or restart file or via the {\hyperref[\detokenize{angle_coeff::doc}]{\sphinxcrossref{\DUrole{doc}{angle\_coeff}}}} command.

\sphinxAtStartPar
All angle potentials store their coefficient data in binary restart
files which means angle\_style and {\hyperref[\detokenize{angle_coeff::doc}]{\sphinxcrossref{\DUrole{doc}{angle\_coeff}}}}
commands do not need to be re\sphinxhyphen{}specified in an input script that
restarts a simulation.  See the {\hyperref[\detokenize{read_restart::doc}]{\sphinxcrossref{\DUrole{doc}{read\_restart}}}}
command for details on how to do this.  The one exception is that
angle\_style \sphinxstyleemphasis{hybrid} only stores the list of sub\sphinxhyphen{}styles in the restart
file; angle coefficients need to be re\sphinxhyphen{}specified.

\begin{sphinxadmonition}{note}{Note:}
\sphinxAtStartPar
When both an angle and pair style is defined, the
{\hyperref[\detokenize{special_bonds::doc}]{\sphinxcrossref{\DUrole{doc}{special\_bonds}}}} command often needs to be used to
turn off (or weight) the pairwise interaction that would otherwise
exist between 3 bonded atoms.
\end{sphinxadmonition}

\sphinxAtStartPar
In the formulas listed for each angle style, \sphinxstyleemphasis{theta} is the angle
between the three atoms in the angle.


\bigskip\hrule\bigskip


\sphinxAtStartPar
Here is an alphabetic list of angle styles defined in LAMMPS.  Click on
the style to display the formula it computes and coefficients
specified by the associated {\hyperref[\detokenize{angle_coeff::doc}]{\sphinxcrossref{\DUrole{doc}{angle\_coeff}}}} command.

\sphinxAtStartPar
Click on the style to display the formula it computes, any additional
arguments specified in the angle\_style command, and coefficients
specified by the associated {\hyperref[\detokenize{angle_coeff::doc}]{\sphinxcrossref{\DUrole{doc}{angle\_coeff}}}} command.

\sphinxAtStartPar
There are also additional accelerated pair styles included in the
LAMMPS distribution for faster performance on CPUs, GPUs, and KNLs.
The individual style names on the {\hyperref[\detokenize{Commands_bond:angle}]{\sphinxcrossref{\DUrole{std,std-ref}{Commands angle}}}} page are followed by one or more
of (g,i,k,o,t) to indicate which accelerated styles exist.
\begin{itemize}
\item {} 
\sphinxAtStartPar
{\hyperref[\detokenize{angle_none::doc}]{\sphinxcrossref{\DUrole{doc}{none}}}} \sphinxhyphen{} turn off angle interactions

\item {} 
\sphinxAtStartPar
{\hyperref[\detokenize{angle_zero::doc}]{\sphinxcrossref{\DUrole{doc}{zero}}}} \sphinxhyphen{} topology but no interactions

\item {} 
\sphinxAtStartPar
{\hyperref[\detokenize{angle_hybrid::doc}]{\sphinxcrossref{\DUrole{doc}{hybrid}}}} \sphinxhyphen{} define multiple styles of angle interactions

\item {} 
\sphinxAtStartPar
{\hyperref[\detokenize{angle_amoeba::doc}]{\sphinxcrossref{\DUrole{doc}{amoeba}}}} \sphinxhyphen{} AMOEBA angle

\item {} 
\sphinxAtStartPar
{\hyperref[\detokenize{angle_charmm::doc}]{\sphinxcrossref{\DUrole{doc}{charmm}}}} \sphinxhyphen{} CHARMM angle

\item {} 
\sphinxAtStartPar
{\hyperref[\detokenize{angle_class2::doc}]{\sphinxcrossref{\DUrole{doc}{class2}}}} \sphinxhyphen{} COMPASS (class 2) angle

\item {} 
\sphinxAtStartPar
{\hyperref[\detokenize{angle_class2::doc}]{\sphinxcrossref{\DUrole{doc}{class2/p6}}}} \sphinxhyphen{} COMPASS (class 2) angle expanded to 6th order

\item {} 
\sphinxAtStartPar
{\hyperref[\detokenize{angle_cosine::doc}]{\sphinxcrossref{\DUrole{doc}{cosine}}}} \sphinxhyphen{} angle with cosine term

\item {} 
\sphinxAtStartPar
{\hyperref[\detokenize{angle_cosine_buck6d::doc}]{\sphinxcrossref{\DUrole{doc}{cosine/buck6d}}}} \sphinxhyphen{} same as cosine with Buckingham term between 1\sphinxhyphen{}3 atoms

\item {} 
\sphinxAtStartPar
{\hyperref[\detokenize{angle_cosine_delta::doc}]{\sphinxcrossref{\DUrole{doc}{cosine/delta}}}} \sphinxhyphen{} angle with difference of cosines

\item {} 
\sphinxAtStartPar
{\hyperref[\detokenize{angle_cosine_periodic::doc}]{\sphinxcrossref{\DUrole{doc}{cosine/periodic}}}} \sphinxhyphen{} DREIDING angle

\item {} 
\sphinxAtStartPar
{\hyperref[\detokenize{angle_cosine_shift::doc}]{\sphinxcrossref{\DUrole{doc}{cosine/shift}}}} \sphinxhyphen{} angle cosine with a shift

\item {} 
\sphinxAtStartPar
{\hyperref[\detokenize{angle_cosine_shift_exp::doc}]{\sphinxcrossref{\DUrole{doc}{cosine/shift/exp}}}} \sphinxhyphen{} cosine with shift and exponential term in spring constant

\item {} 
\sphinxAtStartPar
{\hyperref[\detokenize{angle_cosine_squared::doc}]{\sphinxcrossref{\DUrole{doc}{cosine/squared}}}} \sphinxhyphen{} angle with cosine squared term

\item {} 
\sphinxAtStartPar
{\hyperref[\detokenize{angle_cosine_squared_restricted::doc}]{\sphinxcrossref{\DUrole{doc}{cosine/squared/restricted}}}} \sphinxhyphen{} angle with restricted cosine squared term

\item {} 
\sphinxAtStartPar
{\hyperref[\detokenize{angle_cross::doc}]{\sphinxcrossref{\DUrole{doc}{cross}}}} \sphinxhyphen{} cross term coupling angle and bond lengths

\item {} 
\sphinxAtStartPar
{\hyperref[\detokenize{angle_dipole::doc}]{\sphinxcrossref{\DUrole{doc}{dipole}}}} \sphinxhyphen{} angle that controls orientation of a point dipole

\item {} 
\sphinxAtStartPar
{\hyperref[\detokenize{angle_fourier::doc}]{\sphinxcrossref{\DUrole{doc}{fourier}}}} \sphinxhyphen{} angle with multiple cosine terms

\item {} 
\sphinxAtStartPar
{\hyperref[\detokenize{angle_fourier_simple::doc}]{\sphinxcrossref{\DUrole{doc}{fourier/simple}}}} \sphinxhyphen{} angle with a single cosine term

\item {} 
\sphinxAtStartPar
{\hyperref[\detokenize{angle_gaussian::doc}]{\sphinxcrossref{\DUrole{doc}{gaussian}}}} \sphinxhyphen{} multi\sphinxhyphen{}centered Gaussian\sphinxhyphen{}based angle potential

\item {} 
\sphinxAtStartPar
{\hyperref[\detokenize{angle_harmonic::doc}]{\sphinxcrossref{\DUrole{doc}{harmonic}}}} \sphinxhyphen{} harmonic angle

\item {} 
\sphinxAtStartPar
{\hyperref[\detokenize{angle_lepton::doc}]{\sphinxcrossref{\DUrole{doc}{lepton}}}} \sphinxhyphen{} angle potential from evaluating a string

\item {} 
\sphinxAtStartPar
{\hyperref[\detokenize{angle_mesocnt::doc}]{\sphinxcrossref{\DUrole{doc}{mesocnt}}}} \sphinxhyphen{} piecewise harmonic and linear angle for bending\sphinxhyphen{}buckling of nanotubes

\item {} 
\sphinxAtStartPar
{\hyperref[\detokenize{angle_mm3::doc}]{\sphinxcrossref{\DUrole{doc}{mm3}}}} \sphinxhyphen{} anharmonic angle

\item {} 
\sphinxAtStartPar
{\hyperref[\detokenize{angle_quartic::doc}]{\sphinxcrossref{\DUrole{doc}{quartic}}}} \sphinxhyphen{} angle with cubic and quartic terms

\item {} 
\sphinxAtStartPar
{\hyperref[\detokenize{angle_spica::doc}]{\sphinxcrossref{\DUrole{doc}{spica}}}} \sphinxhyphen{} harmonic angle with repulsive SPICA pair style between 1\sphinxhyphen{}3 atoms

\item {} 
\sphinxAtStartPar
{\hyperref[\detokenize{angle_table::doc}]{\sphinxcrossref{\DUrole{doc}{table}}}} \sphinxhyphen{} tabulated by angle

\end{itemize}


\bigskip\hrule\bigskip



\subsection{Restrictions}
\label{\detokenize{angle_style:restrictions}}
\sphinxAtStartPar
Angle styles can only be set for atom\_styles that allow angles to be
defined.

\sphinxAtStartPar
Most angle styles are part of the MOLECULE package.  They are only
enabled if LAMMPS was built with that package.  See the {\hyperref[\detokenize{Build_package::doc}]{\sphinxcrossref{\DUrole{doc}{Build package}}}} page for more info.  The doc pages for
individual bond potentials tell if it is part of a package.


\subsection{Related commands}
\label{\detokenize{angle_style:related-commands}}
\sphinxAtStartPar
{\hyperref[\detokenize{angle_coeff::doc}]{\sphinxcrossref{\DUrole{doc}{angle\_coeff}}}}


\subsection{Default}
\label{\detokenize{angle_style:default}}
\begin{sphinxVerbatim}[commandchars=\\\{\}]
\PYG{k}{angle\PYGZus{}style}\PYG{+w}{ }\PYG{n}{none}
\end{sphinxVerbatim}

\sphinxstepscope

\index{angle\_write@\spxentry{angle\_write}}\ignorespaces 

\section{angle\_write command}
\label{\detokenize{angle_write:angle-write-command}}\label{\detokenize{angle_write:index-0}}\label{\detokenize{angle_write::doc}}

\subsection{Syntax}
\label{\detokenize{angle_write:syntax}}
\begin{sphinxVerbatim}[commandchars=\\\{\}]
\PYG{k}{angle\PYGZus{}write}\PYG{+w}{ }\PYG{n}{atype}\PYG{+w}{ }\PYG{n}{N}\PYG{+w}{ }\PYG{n}{file}\PYG{+w}{ }\PYG{n}{keyword}
\end{sphinxVerbatim}
\begin{itemize}
\item {} 
\sphinxAtStartPar
atype = angle type

\item {} 
\sphinxAtStartPar
N = \# of values

\item {} 
\sphinxAtStartPar
file = name of file to write values to

\item {} 
\sphinxAtStartPar
keyword = section name in file for this set of tabulated values

\end{itemize}


\subsection{Examples}
\label{\detokenize{angle_write:examples}}
\begin{sphinxVerbatim}[commandchars=\\\{\}]
\PYG{k}{angle\PYGZus{}write}\PYG{+w}{ }\PYG{l+m}{1}\PYG{+w}{  }\PYG{l+m}{500}\PYG{+w}{ }\PYG{n}{table.txt}\PYG{+w}{ }\PYG{n}{Harmonic\PYGZus{}1}
\PYG{k}{angle\PYGZus{}write}\PYG{+w}{ }\PYG{l+m}{3}\PYG{+w}{ }\PYG{l+m}{1000}\PYG{+w}{ }\PYG{n}{table.txt}\PYG{+w}{ }\PYG{n}{Harmonic\PYGZus{}3}
\end{sphinxVerbatim}


\subsection{Description}
\label{\detokenize{angle_write:description}}
\sphinxAtStartPar
\DUrole{versionmodified,added}{Added in version 8Feb2023.}

\sphinxAtStartPar
Write energy and force values to a file as a function of angle for the
currently defined angle potential.  Force in this context means the
force with respect to the angle, not the force on individual atoms.
This is useful for plotting the potential function or otherwise
debugging its values.  The resulting file can also be used as input for
use with {\hyperref[\detokenize{angle_table::doc}]{\sphinxcrossref{\DUrole{doc}{angle style table}}}}.

\sphinxAtStartPar
If the file already exists, the table of values is appended to the end
of the file to allow multiple tables of energy and force to be included
in one file.  The individual sections may be identified by the \sphinxstyleemphasis{keyword}.

\sphinxAtStartPar
The energy and force values are computed for angles ranging from 0
degrees to 180 degrees for 3 interacting atoms forming an angle type
atype, using the appropriate {\hyperref[\detokenize{angle_coeff::doc}]{\sphinxcrossref{\DUrole{doc}{angle\_coeff}}}}
coefficients. N evenly spaced angles are used.

\sphinxAtStartPar
For example, for N = 6, values are computed at \(\theta = 0, 36,
72, 108, 144, 180\).

\sphinxAtStartPar
The file is written in the format used as input for the
{\hyperref[\detokenize{angle_table::doc}]{\sphinxcrossref{\DUrole{doc}{angle\_style table}}}} option with \sphinxstyleemphasis{keyword} as the
section name.  Each line written to the file lists an index number
(1\sphinxhyphen{}N), an angle (in degrees), an energy (in energy units), and a force
(in force units per radians\textasciicircum{}2).  In case a new file is created, the
first line will be a comment with a “DATE:” and “UNITS:” tag with the
current date and {\hyperref[\detokenize{units::doc}]{\sphinxcrossref{\DUrole{doc}{units}}}} settings.  For subsequent
invocations of the \sphinxstyleemphasis{angle\_write} command for the same file, data will be
appended and the current units settings will be compared to the data
from the header, if present. The \sphinxstyleemphasis{angle\_write} will refuse to add a
table to an existing file if the units are not the same.


\subsection{Restrictions}
\label{\detokenize{angle_write:restrictions}}
\sphinxAtStartPar
All force field coefficients for angle and other kinds of interactions
must be set before this command can be invoked.

\sphinxAtStartPar
The table of the angle energy and force data data is created by using a
separate, internally created, new LAMMPS instance with a dummy system of
3 atoms for which the angle potential energy is computed after
transferring the angle style and coefficients and arranging the three atoms
into the corresponding geometries.  The angle force is then determined
from the potential energies through numerical differentiation.  As a
consequence of this approach, not all angle styles are compatible. The
following conditions must be met:
\begin{itemize}
\item {} 
\sphinxAtStartPar
The angle style must be able to write its coefficients to a data file.
This condition excludes for example {\hyperref[\detokenize{angle_hybrid::doc}]{\sphinxcrossref{\DUrole{doc}{angle style hybrid}}}} and
{\hyperref[\detokenize{angle_table::doc}]{\sphinxcrossref{\DUrole{doc}{angle style table}}}}.

\item {} 
\sphinxAtStartPar
The potential function must not have any terms that depend on geometry
properties other than the angle. This condition excludes for example
{\hyperref[\detokenize{angle_class2::doc}]{\sphinxcrossref{\DUrole{doc}{angle style class2}}}} all angle types for
{\hyperref[\detokenize{angle_charmm::doc}]{\sphinxcrossref{\DUrole{doc}{angle style charmm}}}} that have non\sphinxhyphen{}zero
Urey\sphinxhyphen{}Bradley terms.  Please note that the \sphinxstyleemphasis{write\_angle} command has no
way of checking for this condition, so the resulting tables may be
bogus if the requirement is not met.  It is thus recommended to make
careful tests for any created tables.

\end{itemize}


\subsection{Related commands}
\label{\detokenize{angle_write:related-commands}}
\sphinxAtStartPar
{\hyperref[\detokenize{angle_table::doc}]{\sphinxcrossref{\DUrole{doc}{angle\_style table}}}}, {\hyperref[\detokenize{bond_write::doc}]{\sphinxcrossref{\DUrole{doc}{bond\_write}}}},
{\hyperref[\detokenize{dihedral_write::doc}]{\sphinxcrossref{\DUrole{doc}{dihedral\_write}}}}, {\hyperref[\detokenize{angle_style::doc}]{\sphinxcrossref{\DUrole{doc}{angle\_style}}}},
{\hyperref[\detokenize{angle_coeff::doc}]{\sphinxcrossref{\DUrole{doc}{angle\_coeff}}}}


\subsection{Default}
\label{\detokenize{angle_write:default}}
\sphinxAtStartPar
none

\sphinxstepscope

\index{atom\_modify@\spxentry{atom\_modify}}\ignorespaces 

\section{atom\_modify command}
\label{\detokenize{atom_modify:atom-modify-command}}\label{\detokenize{atom_modify:index-0}}\label{\detokenize{atom_modify::doc}}

\subsection{Syntax}
\label{\detokenize{atom_modify:syntax}}
\begin{sphinxVerbatim}[commandchars=\\\{\}]
\PYG{k}{atom\PYGZus{}modify}\PYG{+w}{ }\PYG{n}{keyword}\PYG{+w}{ }\PYG{n}{values}\PYG{+w}{ }\PYG{n}{...}
\end{sphinxVerbatim}
\begin{itemize}
\item {} 
\sphinxAtStartPar
one or more keyword/value pairs may be appended

\item {} 
\sphinxAtStartPar
keyword = \sphinxstyleemphasis{id} or \sphinxstyleemphasis{map} or \sphinxstyleemphasis{first} or \sphinxstyleemphasis{sort}
\begin{sphinxalltt}
\sphinxstyleemphasis{id} value = \sphinxstyleemphasis{yes} or \sphinxstyleemphasis{no}
\sphinxstyleemphasis{map} value = \sphinxstyleemphasis{yes} or \sphinxstyleemphasis{array} or \sphinxstyleemphasis{hash}
\sphinxstyleemphasis{first} value = group\sphinxhyphen{}ID = group whose atoms will appear first in internal atom lists
\sphinxstyleemphasis{sort} values = Nfreq binsize
  Nfreq = sort atoms spatially every this many time steps
  binsize = bin size for spatial sorting (distance units)
\end{sphinxalltt}

\end{itemize}


\subsection{Examples}
\label{\detokenize{atom_modify:examples}}
\begin{sphinxVerbatim}[commandchars=\\\{\}]
\PYG{k}{atom\PYGZus{}modify}\PYG{+w}{ }\PYG{n}{map}\PYG{+w}{ }\PYG{n}{yes}
\PYG{k}{atom\PYGZus{}modify}\PYG{+w}{ }\PYG{n}{map}\PYG{+w}{ }\PYG{n}{hash}\PYG{+w}{ }\PYG{n}{sort}\PYG{+w}{ }\PYG{l+m}{10000}\PYG{+w}{ }\PYG{l+m}{2.0}
\PYG{k}{atom\PYGZus{}modify}\PYG{+w}{ }\PYG{n}{first}\PYG{+w}{ }\PYG{n}{colloid}
\end{sphinxVerbatim}


\subsection{Description}
\label{\detokenize{atom_modify:description}}
\sphinxAtStartPar
Modify certain attributes of atoms defined and stored within LAMMPS,
in addition to what is specified by the {\hyperref[\detokenize{atom_style::doc}]{\sphinxcrossref{\DUrole{doc}{atom\_style}}}}
command.  The \sphinxstyleemphasis{id} and \sphinxstyleemphasis{map} keywords must be specified before a
simulation box is defined; other keywords can be specified any time.

\sphinxAtStartPar
The \sphinxstyleemphasis{id} keyword determines whether non\sphinxhyphen{}zero atom IDs can be assigned
to each atom.  If the value is \sphinxstyleemphasis{yes}, which is the default, IDs are
assigned, whether you use the {\hyperref[\detokenize{create_atoms::doc}]{\sphinxcrossref{\DUrole{doc}{create atoms}}}} or
{\hyperref[\detokenize{read_data::doc}]{\sphinxcrossref{\DUrole{doc}{read\_data}}}} or {\hyperref[\detokenize{read_restart::doc}]{\sphinxcrossref{\DUrole{doc}{read\_restart}}}}
commands to initialize atoms.  If the value is \sphinxstyleemphasis{no} the IDs for all
atoms are assumed to be 0.

\sphinxAtStartPar
If atom IDs are used, they must all be positive integers.  They should
also be unique, though LAMMPS does not check for this.  Typically they
should also be consecutively numbered (from 1 to Natoms), though this
is not required.  Molecular {\hyperref[\detokenize{atom_style::doc}]{\sphinxcrossref{\DUrole{doc}{atom styles}}}} are those
that store bond topology information (styles bond, angle, molecular,
full).  These styles require atom IDs since the IDs are used to encode
the topology.  Some other LAMMPS commands also require the use of atom
IDs.  E.g. some many\sphinxhyphen{}body pair styles use them to avoid double
computation of the I\sphinxhyphen{}J interaction between two atoms.

\sphinxAtStartPar
The only reason not to use atom IDs is if you are running an atomic
simulation so large that IDs cannot be uniquely assigned.  For a
default LAMMPS build this limit is 2\textasciicircum{}31 or about 2 billion atoms.
However, even in this case, you can use 64\sphinxhyphen{}bit atom IDs, allowing 2\textasciicircum{}63
or about 9e18 atoms, if you build LAMMPS with the \sphinxhyphen{} DLAMMPS\_BIGBIG
switch.  This is described on the {\hyperref[\detokenize{Build_settings::doc}]{\sphinxcrossref{\DUrole{doc}{Build\_settings}}}}
doc page.  If atom IDs are not used, they must be specified as 0 for
all atoms, e.g. in a data or restart file.

\begin{sphinxadmonition}{note}{Note:}
\sphinxAtStartPar
If a {\hyperref[\detokenize{Howto_triclinic::doc}]{\sphinxcrossref{\DUrole{doc}{triclinic simulation box}}}} is used,
atom IDs are required, due to how neighbor lists are built.
\end{sphinxadmonition}

\sphinxAtStartPar
The \sphinxstyleemphasis{map} keyword determines how atoms with specific IDs are found
when required.  For example, the bond (angle, etc) methods need to
find the local index of an atom with a specific global ID which is a
bond (angle, etc) partner.  LAMMPS performs this operation efficiently
by creating a “map”, which is either an \sphinxstyleemphasis{array} or \sphinxstyleemphasis{hash} table, as
described below.

\sphinxAtStartPar
When the \sphinxstyleemphasis{map} keyword is not specified in your input script, LAMMPS
only creates a map for {\hyperref[\detokenize{atom_style::doc}]{\sphinxcrossref{\DUrole{doc}{atom\_styles}}}} for molecular
systems which have permanent bonds (angles, etc).  No map is created
for atomic systems, since it is normally not needed.  However some
LAMMPS commands require a map, even for atomic systems, and will
generate an error if one does not exist.  The \sphinxstyleemphasis{map} keyword thus
allows you to force the creation of a map.

\sphinxAtStartPar
Specifying a value of \sphinxstyleemphasis{yes} will create either an array\sphinxhyphen{}style or
hash\sphinxhyphen{}style map, depending on the size of the system.  If no atom ID is
larger than 1 million, then an array\sphinxhyphen{}style map is used, otherwise a
hash\sphinxhyphen{}style map is used.  Specifying a value of \sphinxstyleemphasis{array} or \sphinxstyleemphasis{hash}
creates an array\sphinxhyphen{}style or hash\sphinxhyphen{}style map respectively, regardless of
the size of the system.

\sphinxAtStartPar
For an array\sphinxhyphen{}style map, each processor stores a lookup table of length
N, where N is the largest atom ID in the system.  This is a fast,
simple method for many simulations, but requires too much memory for
large simulations.  For a hash\sphinxhyphen{}style map, a hash table is created on
each processor, which finds an atom ID in constant time (independent
of the global number of atom IDs).  It can be slightly slower than the
\sphinxstyleemphasis{array} map, but its memory cost is proportional to the number of
atoms owned by a processor, i.e. N/P when N is the total number of
atoms in the system and P is the number of processors.

\sphinxAtStartPar
The \sphinxstyleemphasis{first} keyword allows a {\hyperref[\detokenize{group::doc}]{\sphinxcrossref{\DUrole{doc}{group}}}} to be specified
whose atoms will be maintained as the first atoms in each processor’s
list of owned atoms.  This in only useful when the specified group is
a small fraction of all the atoms, and there are other operations
LAMMPS is performing that will be sped\sphinxhyphen{}up significantly by being able
to loop over the smaller set of atoms.  Otherwise the reordering
required by this option will be a net slow\sphinxhyphen{}down.  The
{\hyperref[\detokenize{neigh_modify::doc}]{\sphinxcrossref{\DUrole{doc}{neigh\_modify include}}}} and {\hyperref[\detokenize{comm_modify::doc}]{\sphinxcrossref{\DUrole{doc}{comm\_modify group}}}} commands are two examples of commands that require this
setting to work efficiently.  Several {\hyperref[\detokenize{fix::doc}]{\sphinxcrossref{\DUrole{doc}{fixes}}}}, most notably
time integration fixes like {\hyperref[\detokenize{fix_nve::doc}]{\sphinxcrossref{\DUrole{doc}{fix nve}}}}, also take
advantage of this setting if the group they operate on is the group
specified by this command.  Note that specifying “all” as the group\sphinxhyphen{}ID
effectively turns off the \sphinxstyleemphasis{first} option.

\sphinxAtStartPar
It is OK to use the \sphinxstyleemphasis{first} keyword with a group that has not yet been
defined, e.g. to use the atom\_modify first command at the beginning of
your input script.  LAMMPS does not use the group until a simulation
is run.

\sphinxAtStartPar
The \sphinxstyleemphasis{sort} keyword turns on a spatial sorting or reordering of atoms
within each processor’s subdomain every \sphinxstyleemphasis{Nfreq} timesteps.  If
\sphinxstyleemphasis{Nfreq} is set to 0, then sorting is turned off.  Sorting can improve
cache performance and thus speed\sphinxhyphen{}up a LAMMPS simulation, as discussed
in a paper by {\hyperref[\detokenize{atom_modify:meloni}]{\sphinxcrossref{\DUrole{std,std-ref}{(Meloni)}}}}.  Its efficacy depends on the problem
size (atoms/processor), how quickly the system becomes disordered, and
various other factors.  As a general rule, sorting is typically more
effective at speeding up simulations of liquids as opposed to solids.
In tests we have done, the speed\sphinxhyphen{}up can range from zero to 3\sphinxhyphen{}4x.

\sphinxAtStartPar
Reordering is performed every \sphinxstyleemphasis{Nfreq} timesteps during a dynamics run
or iterations during a minimization.  More precisely, reordering
occurs at the first reneighboring that occurs after the target
timestep.  The reordering is performed locally by each processor,
using bins of the specified \sphinxstyleemphasis{binsize}.  If \sphinxstyleemphasis{binsize} is set to 0.0,
then a binsize equal to half the {\hyperref[\detokenize{neighbor::doc}]{\sphinxcrossref{\DUrole{doc}{neighbor}}}} cutoff
distance (force cutoff plus skin distance) is used, which is a
reasonable value.  After the atoms have been binned, they are
reordered so that atoms in the same bin are adjacent to each other in
the processor’s 1d list of atoms.

\sphinxAtStartPar
The goal of this procedure is for atoms to put atoms close to each
other in the processor’s one\sphinxhyphen{}dimensional list of atoms that are also
near to each other spatially.  This can improve cache performance when
pairwise interactions and neighbor lists are computed.  Note that if
bins are too small, there will be few atoms/bin.  Likewise if bins are
too large, there will be many atoms/bin.  In both cases, the goal of
cache locality will be undermined.

\begin{sphinxadmonition}{note}{Note:}
\sphinxAtStartPar
Running a simulation with sorting on versus off should not change
the simulation results in a statistical sense.  However, a
different ordering will induce round\sphinxhyphen{}off differences, which will
lead to diverging trajectories over time when comparing two
simulations.  Various commands, particularly those which use random
numbers (e.g. {\hyperref[\detokenize{velocity::doc}]{\sphinxcrossref{\DUrole{doc}{velocity create}}}}, and {\hyperref[\detokenize{fix_langevin::doc}]{\sphinxcrossref{\DUrole{doc}{fix
langevin}}}}), may generate (statistically identical)
results which depend on the order in which atoms are processed.
The order of atoms in a {\hyperref[\detokenize{dump::doc}]{\sphinxcrossref{\DUrole{doc}{dump}}}} file will also typically
change if sorting is enabled.
\end{sphinxadmonition}

\begin{sphinxadmonition}{note}{Note:}
\sphinxAtStartPar
When running simple pair\sphinxhyphen{}wise potentials like Lennard Jones on GPUs
with the KOKKOS package, using a larger binsize (e.g. 2x larger than
default) and a more frequent reordering than default (e.g. every 100
time steps) may improve performance.
\end{sphinxadmonition}


\subsection{Restrictions}
\label{\detokenize{atom_modify:restrictions}}
\sphinxAtStartPar
The \sphinxstyleemphasis{first} and \sphinxstyleemphasis{sort} options cannot be used together.  Since sorting
is on by default, it will be turned off if the \sphinxstyleemphasis{first} keyword is
used with a group\sphinxhyphen{}ID that is not “all”.


\subsection{Related commands}
\label{\detokenize{atom_modify:related-commands}}
\sphinxAtStartPar
none


\subsection{Default}
\label{\detokenize{atom_modify:default}}
\sphinxAtStartPar
By default, \sphinxstyleemphasis{id} is yes.  By default, atomic systems (no bond topology
info) do not use a map.  For molecular systems (with bond topology
info), the default is to use a map of either \sphinxstyleemphasis{array} or \sphinxstyleemphasis{hash} style
depending on the size of the system, as explained above for the \sphinxstyleemphasis{map
yes} keyword/value option.  By default, a \sphinxstyleemphasis{first} group is not
defined.  By default, sorting is enabled with a frequency of 1000 and
a binsize of 0.0, which means the neighbor cutoff will be used to set
the bin size. If no neighbor cutoff is defined, sorting will be turned
off.


\bigskip\hrule\bigskip

\phantomsection\label{\detokenize{atom_modify:meloni}}
\sphinxAtStartPar
\sphinxstylestrong{(Meloni)} Meloni, Rosati and Colombo, J Chem Phys, 126, 121102 (2007).

\sphinxstepscope

\index{atom\_style@\spxentry{atom\_style}}\ignorespaces 

\section{atom\_style command}
\label{\detokenize{atom_style:atom-style-command}}\label{\detokenize{atom_style:index-0}}\label{\detokenize{atom_style::doc}}

\subsection{Syntax}
\label{\detokenize{atom_style:syntax}}
\begin{sphinxVerbatim}[commandchars=\\\{\}]
\PYG{k}{atom\PYGZus{}style}\PYG{+w}{ }\PYG{n}{style}\PYG{+w}{ }\PYG{n}{args}
\end{sphinxVerbatim}
\begin{itemize}
\item {} 
\sphinxAtStartPar
style = \sphinxstyleemphasis{amoeba} or \sphinxstyleemphasis{angle} or \sphinxstyleemphasis{atomic} or \sphinxstyleemphasis{body} or \sphinxstyleemphasis{bond} or \sphinxstyleemphasis{charge} or \sphinxstyleemphasis{dielectric} or \sphinxstyleemphasis{dipole} or  \sphinxstyleemphasis{dpd} or \sphinxstyleemphasis{edpd} or \sphinxstyleemphasis{electron} or \sphinxstyleemphasis{ellipsoid} or \sphinxstyleemphasis{full} or \sphinxstyleemphasis{line} or \sphinxstyleemphasis{mdpd} or \sphinxstyleemphasis{molecular} or \sphinxstyleemphasis{oxdna} or \sphinxstyleemphasis{peri} or \sphinxstyleemphasis{smd} or \sphinxstyleemphasis{sph} or \sphinxstyleemphasis{sphere} or \sphinxstyleemphasis{bpm/sphere} or \sphinxstyleemphasis{spin} or \sphinxstyleemphasis{tdpd} or \sphinxstyleemphasis{tri} or \sphinxstyleemphasis{template} or \sphinxstyleemphasis{wavepacket} or \sphinxstyleemphasis{hybrid}
\begin{sphinxalltt}
args = none for any style except the following
  \sphinxstyleemphasis{body} args = bstyle bstyle\sphinxhyphen{}args
    bstyle = style of body particles
    bstyle\sphinxhyphen{}args = additional arguments specific to the bstyle
                  see the {\hyperref[\detokenize{Howto_body::doc}]{\sphinxcrossref{\DUrole{doc}{Howto body}}}} doc
                  page for details
  \sphinxstyleemphasis{sphere} arg = 0/1 (optional) for static/dynamic particle radii
  \sphinxstyleemphasis{bpm/sphere} arg = 0/1 (optional) for static/dynamic particle radii
  \sphinxstyleemphasis{tdpd} arg = Nspecies
    Nspecies = \# of chemical species
  \sphinxstyleemphasis{template} arg = template\sphinxhyphen{}ID
    template\sphinxhyphen{}ID = ID of molecule template specified in a separate {\hyperref[\detokenize{molecule::doc}]{\sphinxcrossref{\DUrole{doc}{molecule}}}} command
  \sphinxstyleemphasis{hybrid} args = list of one or more sub\sphinxhyphen{}styles, each with their args
\end{sphinxalltt}

\item {} 
\sphinxAtStartPar
accelerated styles (with same args) = \sphinxstyleemphasis{angle/kk} or \sphinxstyleemphasis{atomic/kk} or \sphinxstyleemphasis{bond/kk} or \sphinxstyleemphasis{charge/kk} or \sphinxstyleemphasis{full/kk} or \sphinxstyleemphasis{molecular/kk} or \sphinxstyleemphasis{spin/kk}

\end{itemize}


\subsection{Examples}
\label{\detokenize{atom_style:examples}}
\begin{sphinxVerbatim}[commandchars=\\\{\}]
\PYG{k}{atom\PYGZus{}style}\PYG{+w}{ }\PYG{n}{atomic}
\PYG{k}{atom\PYGZus{}style}\PYG{+w}{ }\PYG{n}{bond}
\PYG{k}{atom\PYGZus{}style}\PYG{+w}{ }\PYG{n}{full}
\PYG{k}{atom\PYGZus{}style}\PYG{+w}{ }\PYG{n}{body}\PYG{+w}{ }\PYG{n}{nparticle}\PYG{+w}{ }\PYG{l+m}{2}\PYG{+w}{ }\PYG{l+m}{10}
\PYG{k}{atom\PYGZus{}style}\PYG{+w}{ }\PYG{n}{hybrid}\PYG{+w}{ }\PYG{n}{charge}\PYG{+w}{ }\PYG{n}{bond}
\PYG{k}{atom\PYGZus{}style}\PYG{+w}{ }\PYG{n}{hybrid}\PYG{+w}{ }\PYG{n}{charge}\PYG{+w}{ }\PYG{n}{body}\PYG{+w}{ }\PYG{n}{nparticle}\PYG{+w}{ }\PYG{l+m}{2}\PYG{+w}{ }\PYG{l+m}{5}
\PYG{k}{atom\PYGZus{}style}\PYG{+w}{ }\PYG{n}{spin}
\PYG{k}{atom\PYGZus{}style}\PYG{+w}{ }\PYG{n}{template}\PYG{+w}{ }\PYG{n}{myMols}
\PYG{k}{atom\PYGZus{}style}\PYG{+w}{ }\PYG{n}{hybrid}\PYG{+w}{ }\PYG{n}{template}\PYG{+w}{ }\PYG{n}{twomols}\PYG{+w}{ }\PYG{n}{charge}
\PYG{k}{atom\PYGZus{}style}\PYG{+w}{ }\PYG{n}{tdpd}\PYG{+w}{ }\PYG{l+m}{2}
\end{sphinxVerbatim}


\subsection{Description}
\label{\detokenize{atom_style:description}}
\sphinxAtStartPar
The \sphinxstyleemphasis{atom\_style} command selects which per\sphinxhyphen{}atom attributes are
associated with atoms in a LAMMPS simulation and thus stored and
communicated with those atoms as well as read from and stored in data
and restart files.  Different models (e.g. {\hyperref[\detokenize{pair_style::doc}]{\sphinxcrossref{\DUrole{doc}{pair styles}}}}) require access to specific per\sphinxhyphen{}atom attributes and thus
require a specific atom style.  For example, to compute Coulomb
interactions, the atom must have a “charge” (aka “q”) attribute.

\sphinxAtStartPar
A number of distinct atom styles exist that combine attributes.  Some
atom styles are a superset of other atom styles.  Further attributes
may be added to atoms either via using a hybrid style which provides a
union of the attributes of the sub\sphinxhyphen{}styles, or via the {\hyperref[\detokenize{fix_property_atom::doc}]{\sphinxcrossref{\DUrole{doc}{fix
property/atom}}}} command.  The \sphinxstyleemphasis{atom\_style} command
must be used before a simulation is setup via a {\hyperref[\detokenize{read_data::doc}]{\sphinxcrossref{\DUrole{doc}{read\_data}}}}, {\hyperref[\detokenize{read_restart::doc}]{\sphinxcrossref{\DUrole{doc}{read\_restart}}}}, or {\hyperref[\detokenize{create_box::doc}]{\sphinxcrossref{\DUrole{doc}{create\_box}}}} command.

\begin{sphinxadmonition}{note}{Note:}
\sphinxAtStartPar
Many of the atom styles discussed here are only enabled if LAMMPS was
built with a specific package, as listed below in the Restrictions
section.
\end{sphinxadmonition}

\sphinxAtStartPar
Once a style is selected and the simulation box defined, it cannot be
changed but only augmented with the {\hyperref[\detokenize{fix_property_atom::doc}]{\sphinxcrossref{\DUrole{doc}{fix property/atom}}}} command.  So one should select an atom style
general enough to encompass all attributes required.  E.g. with atom
style \sphinxstyleemphasis{bond}, it is not possible to define angles and use angle styles.

\sphinxAtStartPar
It is OK to use a style more general than needed, though it may be
slightly inefficient because it will allocate and communicate
additional unused data.


\subsection{Atom style attributes}
\label{\detokenize{atom_style:atom-style-attributes}}
\sphinxAtStartPar
The atom style \sphinxstyleemphasis{atomic} has the minimum subset of per\sphinxhyphen{}atom attributes
and is also the default setting.  It encompasses the following per\sphinxhyphen{}atom
attributes (name of the vector or array in the {\hyperref[\detokenize{Classes_atom::doc}]{\sphinxcrossref{\DUrole{doc}{Atom class}}}} is given in parenthesis): atom\sphinxhyphen{}ID (tag), type (type),
position (x), velocities (v), forces (f), image flags (image), group
membership (mask).  Since all atom styles are a superset of atom style
\sphinxstyleemphasis{atomic}, they all include these attributes.

\sphinxAtStartPar
This table lists all the available atom styles, which attributes they
provide, which {\hyperref[\detokenize{Packages::doc}]{\sphinxcrossref{\DUrole{doc}{package}}}} is required to use them, and
what the typical applications are that use them.  See the
{\hyperref[\detokenize{read_data::doc}]{\sphinxcrossref{\DUrole{doc}{read\_data}}}}, {\hyperref[\detokenize{create_atoms::doc}]{\sphinxcrossref{\DUrole{doc}{create\_atoms}}}}, and
{\hyperref[\detokenize{set::doc}]{\sphinxcrossref{\DUrole{doc}{set}}}} commands for details on how to set these various
quantities.  More information about many of the styles is provided in
the Additional Information section below.


\begin{savenotes}\sphinxattablestart
\sphinxthistablewithglobalstyle
\centering
\begin{tabular}[t]{\X{12}{90}\X{28}{90}\X{18}{90}\X{32}{90}}
\sphinxtoprule
\sphinxstyletheadfamily 
\sphinxAtStartPar
Atom style
&\sphinxstyletheadfamily 
\sphinxAtStartPar
Attributes
&\sphinxstyletheadfamily 
\sphinxAtStartPar
Required package
&\sphinxstyletheadfamily 
\sphinxAtStartPar
Applications
\\
\sphinxmidrule
\sphinxtableatstartofbodyhook
\sphinxAtStartPar
\sphinxstyleemphasis{amoeba}
&
\sphinxAtStartPar
\sphinxstyleemphasis{full} + “1\sphinxhyphen{}5 special neighbor data”
&
\sphinxAtStartPar
{\hyperref[\detokenize{Packages_details:pkg-amoeba}]{\sphinxcrossref{\DUrole{std,std-ref}{AMOEBA}}}}
&
\sphinxAtStartPar
AMOEBA/HIPPO force fields
\\
\sphinxhline
\sphinxAtStartPar
\sphinxstyleemphasis{angle}
&
\sphinxAtStartPar
\sphinxstyleemphasis{bond} + “angle data”
&
\sphinxAtStartPar
{\hyperref[\detokenize{Packages_details:pkg-molecule}]{\sphinxcrossref{\DUrole{std,std-ref}{MOLECULE}}}}
&
\sphinxAtStartPar
bead\sphinxhyphen{}spring polymers with stiffness
\\
\sphinxhline
\sphinxAtStartPar
\sphinxstyleemphasis{atomic}
&
\sphinxAtStartPar
tag, type, x, v, f, image, mask
&&
\sphinxAtStartPar
atomic liquids, solids, metals
\\
\sphinxhline
\sphinxAtStartPar
\sphinxstyleemphasis{body}
&
\sphinxAtStartPar
\sphinxstyleemphasis{atomic} + radius, rmass, angmom, torque, body
&
\sphinxAtStartPar
{\hyperref[\detokenize{Packages_details:pkg-body}]{\sphinxcrossref{\DUrole{std,std-ref}{BODY}}}}
&
\sphinxAtStartPar
arbitrary bodies, see {\hyperref[\detokenize{Howto_body::doc}]{\sphinxcrossref{\DUrole{doc}{body howto}}}}
\\
\sphinxhline
\sphinxAtStartPar
\sphinxstyleemphasis{bond}
&
\sphinxAtStartPar
\sphinxstyleemphasis{atomic} + molecule, nspecial, special + “bond data”
&
\sphinxAtStartPar
{\hyperref[\detokenize{Packages_details:pkg-molecule}]{\sphinxcrossref{\DUrole{std,std-ref}{MOLECULE}}}}
&
\sphinxAtStartPar
bead\sphinxhyphen{}spring polymers
\\
\sphinxhline
\sphinxAtStartPar
\sphinxstyleemphasis{bpm/sphere}
&
\sphinxAtStartPar
\sphinxstyleemphasis{bond} + radius, rmass, omega, torque, quat
&
\sphinxAtStartPar
{\hyperref[\detokenize{Packages_details:pkg-bpm}]{\sphinxcrossref{\DUrole{std,std-ref}{BPM}}}}
&
\sphinxAtStartPar
granular bonded particle models, see {\hyperref[\detokenize{Howto_bpm::doc}]{\sphinxcrossref{\DUrole{doc}{BPM howto}}}}
\\
\sphinxhline
\sphinxAtStartPar
\sphinxstyleemphasis{charge}
&
\sphinxAtStartPar
\sphinxstyleemphasis{atomic} + q
&&
\sphinxAtStartPar
atomic systems with charges
\\
\sphinxhline
\sphinxAtStartPar
\sphinxstyleemphasis{dielectric}
&
\sphinxAtStartPar
\sphinxstyleemphasis{full} + mu, area, ed, em, epsilon, curvature, q\_scaled
&
\sphinxAtStartPar
{\hyperref[\detokenize{Packages_details:pkg-dielectric}]{\sphinxcrossref{\DUrole{std,std-ref}{DIELECTRIC}}}}
&
\sphinxAtStartPar
systems with surface polarization
\\
\sphinxhline
\sphinxAtStartPar
\sphinxstyleemphasis{dipole}
&
\sphinxAtStartPar
\sphinxstyleemphasis{charge} + mu
&
\sphinxAtStartPar
{\hyperref[\detokenize{Packages_details:pkg-dipole}]{\sphinxcrossref{\DUrole{std,std-ref}{DIPOLE}}}}
&
\sphinxAtStartPar
atomic systems with charges and point dipoles
\\
\sphinxhline
\sphinxAtStartPar
\sphinxstyleemphasis{dpd}
&
\sphinxAtStartPar
\sphinxstyleemphasis{atomic} + rho + “reactive DPD data”
&
\sphinxAtStartPar
{\hyperref[\detokenize{Packages_details:pkg-dpd-react}]{\sphinxcrossref{\DUrole{std,std-ref}{DPD\sphinxhyphen{}REACT}}}}
&
\sphinxAtStartPar
reactive DPD
\\
\sphinxhline
\sphinxAtStartPar
\sphinxstyleemphasis{edpd}
&
\sphinxAtStartPar
\sphinxstyleemphasis{atomic} + “eDPD data”
&
\sphinxAtStartPar
{\hyperref[\detokenize{Packages_details:pkg-dpd-meso}]{\sphinxcrossref{\DUrole{std,std-ref}{DPD\sphinxhyphen{}MESO}}}}
&
\sphinxAtStartPar
Energy conservative DPD (eDPD)
\\
\sphinxhline
\sphinxAtStartPar
\sphinxstyleemphasis{electron}
&
\sphinxAtStartPar
\sphinxstyleemphasis{charge} + espin, eradius, ervel, erforce
&
\sphinxAtStartPar
{\hyperref[\detokenize{Packages_details:pkg-eff}]{\sphinxcrossref{\DUrole{std,std-ref}{EFF}}}}
&
\sphinxAtStartPar
Electron force field systems
\\
\sphinxhline
\sphinxAtStartPar
\sphinxstyleemphasis{ellipsoid}
&
\sphinxAtStartPar
\sphinxstyleemphasis{atomic} + rmass, angmom, torque, ellipsoid
&&
\sphinxAtStartPar
aspherical particles
\\
\sphinxhline
\sphinxAtStartPar
\sphinxstyleemphasis{full}
&
\sphinxAtStartPar
\sphinxstyleemphasis{molecular} + q
&
\sphinxAtStartPar
{\hyperref[\detokenize{Packages_details:pkg-molecule}]{\sphinxcrossref{\DUrole{std,std-ref}{MOLECULE}}}}
&
\sphinxAtStartPar
molecular force fields
\\
\sphinxhline
\sphinxAtStartPar
\sphinxstyleemphasis{line}
&
\sphinxAtStartPar
\sphinxstyleemphasis{atomic} + molecule, radius, rmass, omega, torque, line
&&
\sphinxAtStartPar
2\sphinxhyphen{}d rigid body particles
\\
\sphinxhline
\sphinxAtStartPar
\sphinxstyleemphasis{mdpd}
&
\sphinxAtStartPar
\sphinxstyleemphasis{atomic} + rho, drho, vest
&
\sphinxAtStartPar
{\hyperref[\detokenize{Packages_details:pkg-dpd-meso}]{\sphinxcrossref{\DUrole{std,std-ref}{DPD\sphinxhyphen{}MESO}}}}
&
\sphinxAtStartPar
Many\sphinxhyphen{}body DPD (mDPD)
\\
\sphinxhline
\sphinxAtStartPar
\sphinxstyleemphasis{molecular}
&
\sphinxAtStartPar
\sphinxstyleemphasis{angle} + “dihedral and improper data”
&
\sphinxAtStartPar
{\hyperref[\detokenize{Packages_details:pkg-molecule}]{\sphinxcrossref{\DUrole{std,std-ref}{MOLECULE}}}}
&
\sphinxAtStartPar
apolar and uncharged molecules
\\
\sphinxhline
\sphinxAtStartPar
\sphinxstyleemphasis{oxdna}
&
\sphinxAtStartPar
\sphinxstyleemphasis{atomic} + id5p
&
\sphinxAtStartPar
{\hyperref[\detokenize{Packages_details:pkg-cg-dna}]{\sphinxcrossref{\DUrole{std,std-ref}{CG\sphinxhyphen{}DNA}}}}
&
\sphinxAtStartPar
coarse\sphinxhyphen{}grained DNA and RNA models
\\
\sphinxhline
\sphinxAtStartPar
\sphinxstyleemphasis{peri}
&
\sphinxAtStartPar
\sphinxstyleemphasis{atomic} + rmass, vfrac, s0, x0
&
\sphinxAtStartPar
{\hyperref[\detokenize{Packages_details:pkg-peri}]{\sphinxcrossref{\DUrole{std,std-ref}{PERI}}}}
&
\sphinxAtStartPar
mesoscopic Peridynamics models
\\
\sphinxhline
\sphinxAtStartPar
\sphinxstyleemphasis{smd}
&
\sphinxAtStartPar
\sphinxstyleemphasis{atomic} + molecule, radius, rmass + “smd data”
&
\sphinxAtStartPar
{\hyperref[\detokenize{Packages_details:pkg-machdyn}]{\sphinxcrossref{\DUrole{std,std-ref}{MACHDYN}}}}
&
\sphinxAtStartPar
Smooth Mach Dynamics models
\\
\sphinxhline
\sphinxAtStartPar
\sphinxstyleemphasis{rheo}
&
\sphinxAtStartPar
\sphinxstyleemphasis{atomic} + rho, status
&
\sphinxAtStartPar
{\hyperref[\detokenize{Packages_details:pkg-rheo}]{\sphinxcrossref{\DUrole{std,std-ref}{RHEO}}}}
&
\sphinxAtStartPar
solid and fluid RHEO particles
\\
\sphinxhline
\sphinxAtStartPar
\sphinxstyleemphasis{rheo/thermal}
&
\sphinxAtStartPar
\sphinxstyleemphasis{atomic} + rho, status, energy, temperature
&
\sphinxAtStartPar
{\hyperref[\detokenize{Packages_details:pkg-rheo}]{\sphinxcrossref{\DUrole{std,std-ref}{RHEO}}}}
&
\sphinxAtStartPar
RHEO particles with temperature
\\
\sphinxhline
\sphinxAtStartPar
\sphinxstyleemphasis{sph}
&
\sphinxAtStartPar
\sphinxstyleemphasis{atomic} + “sph data”
&
\sphinxAtStartPar
{\hyperref[\detokenize{Packages_details:pkg-sph}]{\sphinxcrossref{\DUrole{std,std-ref}{SPH}}}}
&
\sphinxAtStartPar
Smoothed particle hydrodynamics models
\\
\sphinxhline
\sphinxAtStartPar
\sphinxstyleemphasis{sphere}
&
\sphinxAtStartPar
\sphinxstyleemphasis{atomic} + radius, rmass, omega, torque
&&
\sphinxAtStartPar
finite size spherical particles, e.g. granular models
\\
\sphinxhline
\sphinxAtStartPar
\sphinxstyleemphasis{spin}
&
\sphinxAtStartPar
\sphinxstyleemphasis{atomic} + “magnetic moment data”
&
\sphinxAtStartPar
{\hyperref[\detokenize{Packages_details:pkg-spin}]{\sphinxcrossref{\DUrole{std,std-ref}{SPIN}}}}
&
\sphinxAtStartPar
magnetic particles
\\
\sphinxhline
\sphinxAtStartPar
\sphinxstyleemphasis{tdpd}
&
\sphinxAtStartPar
\sphinxstyleemphasis{atomic} + cc, cc\_flux, vest
&
\sphinxAtStartPar
{\hyperref[\detokenize{Packages_details:pkg-dpd-meso}]{\sphinxcrossref{\DUrole{std,std-ref}{DPD\sphinxhyphen{}MESO}}}}
&
\sphinxAtStartPar
Transport DPD (tDPD)
\\
\sphinxhline
\sphinxAtStartPar
\sphinxstyleemphasis{template}
&
\sphinxAtStartPar
\sphinxstyleemphasis{atomic} + molecule, molindex, molatom
&
\sphinxAtStartPar
{\hyperref[\detokenize{Packages_details:pkg-molecule}]{\sphinxcrossref{\DUrole{std,std-ref}{MOLECULE}}}}
&
\sphinxAtStartPar
molecular systems where attributes are taken from {\hyperref[\detokenize{molecule::doc}]{\sphinxcrossref{\DUrole{doc}{molecule files}}}}
\\
\sphinxhline
\sphinxAtStartPar
\sphinxstyleemphasis{tri}
&
\sphinxAtStartPar
\sphinxstyleemphasis{sphere} + molecule, angmom, tri
&&
\sphinxAtStartPar
3\sphinxhyphen{}d triangulated rigid body LJ particles
\\
\sphinxhline
\sphinxAtStartPar
\sphinxstyleemphasis{wavepacket}
&
\sphinxAtStartPar
\sphinxstyleemphasis{charge} + “wavepacket data”
&
\sphinxAtStartPar
{\hyperref[\detokenize{Packages_details:pkg-awpmd}]{\sphinxcrossref{\DUrole{std,std-ref}{AWPMD}}}}
&
\sphinxAtStartPar
Antisymmetrized wave packet MD
\\
\sphinxbottomrule
\end{tabular}
\sphinxtableafterendhook\par
\sphinxattableend\end{savenotes}

\begin{sphinxadmonition}{note}{Note:}
\sphinxAtStartPar
It is possible to add some attributes, such as a molecule ID and
charge, to atom styles that do not have them built in using the
{\hyperref[\detokenize{fix_property_atom::doc}]{\sphinxcrossref{\DUrole{doc}{fix property/atom}}}} command.  This command
also allows new custom\sphinxhyphen{}named attributes consisting of extra integer
or floating\sphinxhyphen{}point values or vectors to be added to atoms.  See the
{\hyperref[\detokenize{fix_property_atom::doc}]{\sphinxcrossref{\DUrole{doc}{fix property/atom}}}} page for examples of
cases where this is useful and details on how to initialize,
access, and output these custom values.
\end{sphinxadmonition}


\bigskip\hrule\bigskip



\subsection{Particle size and mass}
\label{\detokenize{atom_style:particle-size-and-mass}}
\sphinxAtStartPar
All of the atom styles define point particles unless they (1) define
finite\sphinxhyphen{}size spherical particles via the \sphinxstyleemphasis{radius} attribute, or (2)
define finite\sphinxhyphen{}size aspherical particles (e.g. the \sphinxstyleemphasis{body}, \sphinxstyleemphasis{ellipsoid},
\sphinxstyleemphasis{line}, and \sphinxstyleemphasis{tri} styles).  Most of these styles can also be used with
mixtures of point and finite\sphinxhyphen{}size particles.

\sphinxAtStartPar
Note that the \sphinxstyleemphasis{radius} property may need to be provided as a
\sphinxstyleemphasis{diameter} (e.g. in {\hyperref[\detokenize{molecule::doc}]{\sphinxcrossref{\DUrole{doc}{molecule files}}}} or {\hyperref[\detokenize{read_data::doc}]{\sphinxcrossref{\DUrole{doc}{data
files}}}}).  See the {\hyperref[\detokenize{Howto_spherical::doc}]{\sphinxcrossref{\DUrole{doc}{Howto spherical}}}}
page for an overview of using finite\sphinxhyphen{}size spherical and aspherical
particle models with LAMMPS.

\sphinxAtStartPar
Unless an atom style defines the per\sphinxhyphen{}atom \sphinxstyleemphasis{rmass} attribute, particle
masses are defined on a per\sphinxhyphen{}type basis, using the {\hyperref[\detokenize{mass::doc}]{\sphinxcrossref{\DUrole{doc}{mass}}}}
command.  This means each particle’s mass is indexed by its atom
\sphinxstyleemphasis{type}.

\sphinxAtStartPar
A few styles define the per\sphinxhyphen{}atom \sphinxstyleemphasis{rmass} attribute which can also be
added using the {\hyperref[\detokenize{fix_property_atom::doc}]{\sphinxcrossref{\DUrole{doc}{fix property/atom}}}} command.
In this case each particle stores its own mass.  Atom styles that have
a per\sphinxhyphen{}atom rmass may define it indirectly through setting particle
diameter and density on a per\sphinxhyphen{}particle basis.  If both per\sphinxhyphen{}type mass
and per\sphinxhyphen{}atom \sphinxstyleemphasis{rmass} are defined (e.g. in a hybrid style), the
per\sphinxhyphen{}atom mass will take precedence in any operation which which works
with both flavors of mass.


\bigskip\hrule\bigskip



\subsection{Additional information about specific atom styles}
\label{\detokenize{atom_style:additional-information-about-specific-atom-styles}}
\sphinxAtStartPar
For the \sphinxstyleemphasis{body} style, the particles are arbitrary bodies with internal
attributes defined by the “style” of the bodies, which is specified by
the \sphinxstyleemphasis{bstyle} argument.  Body particles can represent complex entities,
such as surface meshes of discrete points, collections of
sub\sphinxhyphen{}particles, deformable objects, etc.

\sphinxAtStartPar
The {\hyperref[\detokenize{Howto_body::doc}]{\sphinxcrossref{\DUrole{doc}{Howto body}}}} page describes the body styles
LAMMPS currently supports, and provides more details as to the kind of
body particles they represent.  For all styles, each body particle
stores moments of inertia and a quaternion 4\sphinxhyphen{}vector, so that its
orientation and position can be time integrated due to forces and
torques.

\sphinxAtStartPar
Note that there may be additional arguments required along with the
\sphinxstyleemphasis{bstyle} specification, in the atom\_style body command.  These
arguments are described on the {\hyperref[\detokenize{Howto_body::doc}]{\sphinxcrossref{\DUrole{doc}{Howto body}}}} doc page.

\sphinxAtStartPar
For the \sphinxstyleemphasis{dielectric} style, each particle can be either a physical
particle (e.g. an ion), or an interface particle representing a
boundary element between two regions of different dielectric
constant. For interface particles, in addition to the properties
associated with atom\_style full, each particle also should be assigned
a unit dipole vector (mu) representing the direction of the induced
dipole moment at each interface particle, an area (area/patch), the
difference and mean of the dielectric constants of two sides of the
interface along the direction of the normal vector (ed and em), the
local dielectric constant at the boundary element (epsilon), and a
mean local curvature (curv).  Physical particles must be assigned
these values, as well, but only their local dielectric constants will
be used; see documentation for associated {\hyperref[\detokenize{pair_dielectric::doc}]{\sphinxcrossref{\DUrole{doc}{pair styles}}}} and {\hyperref[\detokenize{fix_polarize::doc}]{\sphinxcrossref{\DUrole{doc}{fixes}}}}.  The distinction
between the physical and interface particles is only meaningful when
{\hyperref[\detokenize{fix_polarize::doc}]{\sphinxcrossref{\DUrole{doc}{fix polarize}}}} commands are applied to the
interface particles. This style is part of the DIELECTRIC package.

\sphinxAtStartPar
For the \sphinxstyleemphasis{dipole} style, a point dipole vector mu is defined for each
point particle.  Note that if you wish the particles to be finite\sphinxhyphen{}size
spheres as in a Stockmayer potential for a dipolar fluid, so that the
particles can rotate due to dipole\sphinxhyphen{}dipole interactions, then you need
to use the command \sphinxtitleref{atom\_style hybrid sphere dipole}, which will
assign both a diameter and dipole moment to each particle.  This also
requires using an integrator with a “/sphere” suffix like {\hyperref[\detokenize{fix_nve_sphere::doc}]{\sphinxcrossref{\DUrole{doc}{fix
nve/sphere}}}} or {\hyperref[\detokenize{fix_nvt_sphere::doc}]{\sphinxcrossref{\DUrole{doc}{fix nvt/sphere}}}}
and the “update dipole” or “update dlm” parameters to the fix
commands.

\sphinxAtStartPar
The \sphinxstyleemphasis{dpd} style is for reactive dissipative particle dynamics (DPD)
particles.  Note that it is part of the DPD\sphinxhyphen{}REACT package, and is not
required for use with the {\hyperref[\detokenize{pair_dpd::doc}]{\sphinxcrossref{\DUrole{doc}{pair\_style dpd or dpd/stat}}}}
commands, which only require the attributes from atom\_style \sphinxstyleemphasis{atomic}.
Atom\_style \sphinxstyleemphasis{dpd} extends DPD particle properties with internal
temperature (dpdTheta), internal conductive energy (uCond), internal
mechanical energy (uMech), and internal chemical energy (uChem).

\sphinxAtStartPar
The \sphinxstyleemphasis{edpd} style is for energy\sphinxhyphen{}conserving dissipative particle
dynamics (eDPD) particles which store a temperature (edpd\_temp), and
heat capacity (edpd\_cv).

\sphinxAtStartPar
For the \sphinxstyleemphasis{electron} style, the particles representing electrons are 3d
Gaussians with a specified position and bandwidth or uncertainty in
position, which is represented by the eradius = electron size.

\sphinxAtStartPar
For the \sphinxstyleemphasis{ellipsoid} style, particles can be ellipsoids which each
stores a shape vector with the 3 diameters of the ellipsoid and a
quaternion 4\sphinxhyphen{}vector with its orientation.  Each particle stores a flag
in the ellipsoid vector which indicates whether it is an ellipsoid (1)
or a point particle (0).

\sphinxAtStartPar
For the \sphinxstyleemphasis{line} style, particles can be are idealized line segments
which store a per\sphinxhyphen{}particle mass and length and orientation (i.e. the
end points of the line segment).  Each particle stores a flag in the
line vector which indicates whether it is a line segment (1) or a
point particle (0).

\sphinxAtStartPar
The \sphinxstyleemphasis{mdpd} style is for many\sphinxhyphen{}body dissipative particle dynamics (mDPD)
particles which store a density (rho) for considering density\sphinxhyphen{}dependent
many\sphinxhyphen{}body interactions.

\sphinxAtStartPar
The \sphinxstyleemphasis{oxdna} style is for coarse\sphinxhyphen{}grained nucleotides and stores the
3’\sphinxhyphen{}to\sphinxhyphen{}5’ polarity of the nucleotide strand, which is set through
the bond topology in the data file. The first (second) atom in a
bond definition is understood to point towards the 3’\sphinxhyphen{}end (5’\sphinxhyphen{}end)
of the strand.

\sphinxAtStartPar
For the \sphinxstyleemphasis{peri} style, the particles are spherical and each stores a
per\sphinxhyphen{}particle mass and volume.

\sphinxAtStartPar
The \sphinxstyleemphasis{smd} style is for Smooth Particle Mach dynamics.  Both fluids and
solids can be modeled.  Particles store the mass and volume of an
integration point, a kernel diameter used for calculating the field
variables (e.g. stress and deformation) and a contact radius for
calculating repulsive forces which prevent individual physical bodies
from penetrating each other.

\sphinxAtStartPar
The \sphinxstyleemphasis{sph} style is for smoothed particle hydrodynamics (SPH) particles
which store a density (rho), energy (esph), and heat capacity (cv).

\sphinxAtStartPar
For the \sphinxstyleemphasis{spin} style, a magnetic spin is associated with each atom.
Those spins have a norm (their magnetic moment) and a direction.

\sphinxAtStartPar
The \sphinxstyleemphasis{tdpd} style is for transport dissipative particle dynamics (tDPD)
particles which store a set of chemical concentration. An integer
“cc\_species” is required to specify the number of chemical species
involved in a tDPD system.

\sphinxAtStartPar
The \sphinxstyleemphasis{wavepacket} style is similar to the \sphinxstyleemphasis{electron} style, but the
electrons may consist of several Gaussian wave packets, summed up with
coefficients cs= (cs\_re,cs\_im).  Each of the wave packets is treated
as a separate particle in LAMMPS, wave packets belonging to the same
electron must have identical \sphinxstyleemphasis{etag} values.

\sphinxAtStartPar
The \sphinxstyleemphasis{sphere} and \sphinxstyleemphasis{bpm/sphere} styles allow particles to be either point
particles or finite\sphinxhyphen{}size particles.  If the \sphinxstyleemphasis{radius} attribute is \textgreater{}
0.0, the particle is a finite\sphinxhyphen{}size sphere.  If the diameter = 0.0, it
is a point particle.  Note that by using the \sphinxstyleemphasis{disc} keyword with the
{\hyperref[\detokenize{fix_nve_sphere::doc}]{\sphinxcrossref{\DUrole{doc}{fix nve/sphere}}}}, {\hyperref[\detokenize{fix_nvt_sphere::doc}]{\sphinxcrossref{\DUrole{doc}{fix nvt/sphere}}}}, {\hyperref[\detokenize{fix_nph_sphere::doc}]{\sphinxcrossref{\DUrole{doc}{fix nph/sphere}}}}, {\hyperref[\detokenize{fix_npt_sphere::doc}]{\sphinxcrossref{\DUrole{doc}{fix
npt/sphere}}}} commands for the \sphinxstyleemphasis{sphere} style, spheres
can be effectively treated as 2d discs for a 2d simulation if desired.
See also the {\hyperref[\detokenize{set::doc}]{\sphinxcrossref{\DUrole{doc}{set density/disc}}}} command.  These styles also
take an optional 0 or 1 argument.  A value of 0 means the radius of
each sphere is constant for the duration of the simulation (this is
the default).  A value of 1 means the radii may vary dynamically
during the simulation, e.g. due to use of the {\hyperref[\detokenize{fix_adapt::doc}]{\sphinxcrossref{\DUrole{doc}{fix adapt}}}} command.

\sphinxAtStartPar
The \sphinxstyleemphasis{template} style allows molecular topology (bonds,angles,etc) to be
defined via a molecule template using the {\hyperref[\detokenize{molecule::doc}]{\sphinxcrossref{\DUrole{doc}{molecule}}}}
command.  The template stores one or more molecules with a single copy
of the topology info (bonds,angles,etc) of each.  Individual atoms only
store a template index and template atom to identify which molecule and
which atom\sphinxhyphen{}within\sphinxhyphen{}the\sphinxhyphen{}molecule they represent.  Using the \sphinxstyleemphasis{template}
style instead of the \sphinxstyleemphasis{bond}, \sphinxstyleemphasis{angle}, \sphinxstyleemphasis{molecular} styles can save memory
for systems comprised of a large number of small molecules, all of a
single type (or small number of types).  See the paper by Grime and
Voth, in {\hyperref[\detokenize{atom_style:grime}]{\sphinxcrossref{\DUrole{std,std-ref}{(Grime)}}}}, for examples of how this can be
advantageous for large\sphinxhyphen{}scale coarse\sphinxhyphen{}grained systems.  The
\sphinxcode{\sphinxupquote{examples/template}} directory has a few demo inputs and examples
showing the use of the \sphinxstyleemphasis{template} atom style versus \sphinxstyleemphasis{molecular}.

\begin{sphinxadmonition}{note}{Note:}
\sphinxAtStartPar
When using the \sphinxstyleemphasis{template} style with a {\hyperref[\detokenize{molecule::doc}]{\sphinxcrossref{\DUrole{doc}{molecule template}}}} that contains multiple molecules, you should ensure the
atom types, bond types, angle\_types, etc in all the molecules are
consistent.  E.g. if one molecule represents H2O and another CO2,
then you probably do not want each molecule file to define two atom
types and a single bond type, because they will conflict with each
other when a mixture system of H2O and CO2 molecules is defined,
e.g. by the {\hyperref[\detokenize{read_data::doc}]{\sphinxcrossref{\DUrole{doc}{read\_data}}}} command.  Rather the H2O
molecule should define atom types 1 and 2, and bond type 1.  And
the CO2 molecule should define atom types 3 and 4 (or atom types 3
and 2 if a single oxygen type is desired), and bond type 2.
\end{sphinxadmonition}

\sphinxAtStartPar
For the \sphinxstyleemphasis{tri} style, particles can be planar triangles which each
stores a per\sphinxhyphen{}particle mass and size and orientation (i.e. the corner
points of the triangle).  Each particle stores a flag in the tri
vector which indicates whether it is a triangle (1) or a point
particle (0).


\bigskip\hrule\bigskip


\sphinxAtStartPar
Typically, simulations require only a single (non\sphinxhyphen{}hybrid) atom style.
If some atoms in the simulation do not have all the properties defined
by a particular style, use the simplest style that defines all the
needed properties by any atom.  For example, if some atoms in a
simulation are charged, but others are not, use the \sphinxstyleemphasis{charge} style.
If some atoms have bonds, but others do not, use the \sphinxstyleemphasis{bond} style.

\sphinxAtStartPar
The only scenario where the \sphinxstyleemphasis{hybrid} style is needed is if there is no
single style which defines all needed properties of all atoms.  For
example, as mentioned above, if you want dipolar particles which will
rotate due to torque, you need to use “atom\_style hybrid sphere
dipole”.  When a hybrid style is used, atoms store and communicate the
union of all quantities implied by the individual styles.

\sphinxAtStartPar
When using the \sphinxstyleemphasis{hybrid} style, you cannot combine the \sphinxstyleemphasis{template} style
with another molecular style that stores bond, angle, etc info on a
per\sphinxhyphen{}atom basis.

\sphinxAtStartPar
LAMMPS can be extended with new atom styles as well as new body styles;
see the corresponding manual page on {\hyperref[\detokenize{Modify_atom::doc}]{\sphinxcrossref{\DUrole{doc}{modifying \& extending LAMMPS}}}}.


\bigskip\hrule\bigskip


\sphinxAtStartPar
Styles with a \sphinxstyleemphasis{gpu}, \sphinxstyleemphasis{intel}, \sphinxstyleemphasis{kk}, \sphinxstyleemphasis{omp}, or \sphinxstyleemphasis{opt} suffix are
functionally the same as the corresponding style without the suffix.
They have been optimized to run faster, depending on your available
hardware, as discussed on the {\hyperref[\detokenize{Speed_packages::doc}]{\sphinxcrossref{\DUrole{doc}{Accelerator packages}}}}
page.  The accelerated styles take the same arguments and should
produce the same results, except for round\sphinxhyphen{}off and precision issues.

\sphinxAtStartPar
These accelerated styles are part of the GPU, INTEL, KOKKOS,
OPENMP, and OPT packages, respectively.  They are only enabled if
LAMMPS was built with those packages.  See the {\hyperref[\detokenize{Build_package::doc}]{\sphinxcrossref{\DUrole{doc}{Build package}}}} page for more info.

\sphinxAtStartPar
You can specify the accelerated styles explicitly in your input script
by including their suffix, or you can use the {\hyperref[\detokenize{Run_options::doc}]{\sphinxcrossref{\DUrole{doc}{\sphinxhyphen{}suffix command\sphinxhyphen{}line switch}}}} when you invoke LAMMPS, or you can use the
{\hyperref[\detokenize{suffix::doc}]{\sphinxcrossref{\DUrole{doc}{suffix}}}} command in your input script.

\sphinxAtStartPar
See the {\hyperref[\detokenize{Speed_packages::doc}]{\sphinxcrossref{\DUrole{doc}{Accelerator packages}}}} page for more
instructions on how to use the accelerated styles effectively.


\bigskip\hrule\bigskip



\subsection{Restrictions}
\label{\detokenize{atom_style:restrictions}}
\sphinxAtStartPar
This command cannot be used after the simulation box is defined by a
{\hyperref[\detokenize{read_data::doc}]{\sphinxcrossref{\DUrole{doc}{read\_data}}}} or {\hyperref[\detokenize{create_box::doc}]{\sphinxcrossref{\DUrole{doc}{create\_box}}}} command.

\sphinxAtStartPar
Many of the styles listed above are only enabled if LAMMPS was built
with a specific package, as listed below.  See the {\hyperref[\detokenize{Build_package::doc}]{\sphinxcrossref{\DUrole{doc}{Build package}}}} page for more info.  The table above lists which package
is required for individual atom styles.


\subsection{Related commands}
\label{\detokenize{atom_style:related-commands}}
\sphinxAtStartPar
{\hyperref[\detokenize{read_data::doc}]{\sphinxcrossref{\DUrole{doc}{read\_data}}}}, {\hyperref[\detokenize{pair_style::doc}]{\sphinxcrossref{\DUrole{doc}{pair\_style}}}},
{\hyperref[\detokenize{fix_property_atom::doc}]{\sphinxcrossref{\DUrole{doc}{fix property/atom}}}}, {\hyperref[\detokenize{set::doc}]{\sphinxcrossref{\DUrole{doc}{set}}}}


\subsection{Default}
\label{\detokenize{atom_style:default}}
\sphinxAtStartPar
The default atom style is \sphinxstyleemphasis{atomic}.  If atom\_style \sphinxstyleemphasis{sphere} or
\sphinxstyleemphasis{bpm/sphere} is used, its default argument is 0.


\bigskip\hrule\bigskip

\phantomsection\label{\detokenize{atom_style:grime}}
\sphinxAtStartPar
\sphinxstylestrong{(Grime)} Grime and Voth, to appear in J Chem Theory \& Computation
(2014).

\sphinxstepscope

\index{balance@\spxentry{balance}}\ignorespaces 

\section{balance command}
\label{\detokenize{balance:balance-command}}\label{\detokenize{balance:index-0}}\label{\detokenize{balance::doc}}

\subsection{Syntax}
\label{\detokenize{balance:syntax}}
\begin{sphinxVerbatim}[commandchars=\\\{\}]
\PYG{k}{balance}\PYG{+w}{ }\PYG{n}{thresh}\PYG{+w}{ }\PYG{n}{style}\PYG{+w}{ }\PYG{n}{args}\PYG{+w}{ }\PYG{n}{...}\PYG{+w}{ }\PYG{n}{keyword}\PYG{+w}{ }\PYG{n}{args}\PYG{+w}{ }\PYG{n}{...}
\end{sphinxVerbatim}
\begin{itemize}
\item {} 
\sphinxAtStartPar
thresh = imbalance threshold that must be exceeded to perform a re\sphinxhyphen{}balance

\item {} 
\sphinxAtStartPar
one style/arg pair can be used (or multiple for \sphinxstyleemphasis{x},\sphinxstyleemphasis{y},\sphinxstyleemphasis{z})

\item {} 
\sphinxAtStartPar
style = \sphinxstyleemphasis{x} or \sphinxstyleemphasis{y} or \sphinxstyleemphasis{z} or \sphinxstyleemphasis{shift} or \sphinxstyleemphasis{rcb}
\begin{sphinxalltt}
\sphinxstyleemphasis{x} args = \sphinxstyleemphasis{uniform} or Px\sphinxhyphen{}1 numbers between 0 and 1
  \sphinxstyleemphasis{uniform} = evenly spaced cuts between processors in x dimension
  numbers = Px\sphinxhyphen{}1 ascending values between 0 and 1, Px \sphinxhyphen{} \# of processors in x dimension
  \sphinxstyleemphasis{x} can be specified together with \sphinxstyleemphasis{y} or \sphinxstyleemphasis{z}
\sphinxstyleemphasis{y} args = \sphinxstyleemphasis{uniform} or Py\sphinxhyphen{}1 numbers between 0 and 1
  \sphinxstyleemphasis{uniform} = evenly spaced cuts between processors in y dimension
  numbers = Py\sphinxhyphen{}1 ascending values between 0 and 1, Py \sphinxhyphen{} \# of processors in y dimension
  \sphinxstyleemphasis{y} can be specified together with \sphinxstyleemphasis{x} or \sphinxstyleemphasis{z}
\sphinxstyleemphasis{z} args = \sphinxstyleemphasis{uniform} or Pz\sphinxhyphen{}1 numbers between 0 and 1
  \sphinxstyleemphasis{uniform} = evenly spaced cuts between processors in z dimension
  numbers = Pz\sphinxhyphen{}1 ascending values between 0 and 1, Pz \sphinxhyphen{} \# of processors in z dimension
  \sphinxstyleemphasis{z} can be specified together with \sphinxstyleemphasis{x} or \sphinxstyleemphasis{y}
\sphinxstyleemphasis{shift} args = dimstr Niter stopthresh
  dimstr = sequence of letters containing "x" or "y" or "z", each not more than once
  Niter = \# of times to iterate within each dimension of dimstr sequence
  stopthresh = stop balancing when this imbalance threshold is reached
\sphinxstyleemphasis{rcb} args = none
\end{sphinxalltt}

\item {} 
\sphinxAtStartPar
zero or more keyword/arg pairs may be appended

\item {} 
\sphinxAtStartPar
keyword = \sphinxstyleemphasis{weight} or \sphinxstyleemphasis{out}
\begin{sphinxalltt}
\sphinxstyleemphasis{weight} style args = use weighted particle counts for the balancing
  \sphinxstyleemphasis{style} = \sphinxstyleemphasis{group} or \sphinxstyleemphasis{neigh} or \sphinxstyleemphasis{time} or \sphinxstyleemphasis{var} or \sphinxstyleemphasis{store}
    \sphinxstyleemphasis{group} args = Ngroup group1 weight1 group2 weight2 ...
      Ngroup = number of groups with assigned weights
      group1, group2, ... = group IDs
      weight1, weight2, ...   = corresponding weight factors
    \sphinxstyleemphasis{neigh} factor = compute weight based on number of neighbors
      factor = scaling factor (\textgreater{} 0)
    \sphinxstyleemphasis{time} factor = compute weight based on time spend computing
      factor = scaling factor (\textgreater{} 0)
    \sphinxstyleemphasis{var} name = take weight from atom\sphinxhyphen{}style variable
      name = name of the atom\sphinxhyphen{}style variable
    \sphinxstyleemphasis{store} name = store weight in custom atom property defined by {\hyperref[\detokenize{fix_property_atom::doc}]{\sphinxcrossref{\DUrole{doc}{fix property/atom}}}} command
      name = atom property name (without d\_ prefix)
\sphinxstyleemphasis{sort} arg = \sphinxstyleemphasis{no} or \sphinxstyleemphasis{yes}
\sphinxstyleemphasis{out} arg = filename
  filename = write each processor\textquotesingle{}s subdomain to a file
\end{sphinxalltt}

\end{itemize}


\subsection{Examples}
\label{\detokenize{balance:examples}}
\begin{sphinxVerbatim}[commandchars=\\\{\}]
\PYG{k}{balance}\PYG{+w}{ }\PYG{l+m}{0.9}\PYG{+w}{ }\PYG{n}{x}\PYG{+w}{ }\PYG{n}{uniform}\PYG{+w}{ }\PYG{n}{y}\PYG{+w}{ }\PYG{l+m}{0.4}\PYG{+w}{ }\PYG{l+m}{0.5}\PYG{+w}{ }\PYG{l+m}{0.6}
\PYG{k}{balance}\PYG{+w}{ }\PYG{l+m}{1.2}\PYG{+w}{ }\PYG{n}{shift}\PYG{+w}{ }\PYG{n}{xz}\PYG{+w}{ }\PYG{l+m}{5}\PYG{+w}{ }\PYG{l+m}{1.1}
\PYG{k}{balance}\PYG{+w}{ }\PYG{l+m}{1.0}\PYG{+w}{ }\PYG{n}{shift}\PYG{+w}{ }\PYG{n}{xz}\PYG{+w}{ }\PYG{l+m}{5}\PYG{+w}{ }\PYG{l+m}{1.1}
\PYG{k}{balance}\PYG{+w}{ }\PYG{l+m}{1.1}\PYG{+w}{ }\PYG{n}{rcb}
\PYG{k}{balance}\PYG{+w}{ }\PYG{l+m}{1.0}\PYG{+w}{ }\PYG{n}{shift}\PYG{+w}{ }\PYG{n}{x}\PYG{+w}{ }\PYG{l+m}{10}\PYG{+w}{ }\PYG{l+m}{1.1}\PYG{+w}{ }\PYG{n}{weight}\PYG{+w}{ }\PYG{k}{group }\PYG{n+nv+nvIdentifier}{2}\PYG{+w}{ }\PYG{n}{fast}\PYG{+w}{ }\PYG{l+m}{0.5}\PYG{+w}{ }\PYG{n}{slow}\PYG{+w}{ }\PYG{l+m}{2.0}
\PYG{k}{balance}\PYG{+w}{ }\PYG{l+m}{1.0}\PYG{+w}{ }\PYG{n}{shift}\PYG{+w}{ }\PYG{n}{x}\PYG{+w}{ }\PYG{l+m}{10}\PYG{+w}{ }\PYG{l+m}{1.1}\PYG{+w}{ }\PYG{n}{weight}\PYG{+w}{ }\PYG{n}{time}\PYG{+w}{ }\PYG{l+m}{0.8}\PYG{+w}{ }\PYG{n}{weight}\PYG{+w}{ }\PYG{n}{neigh}\PYG{+w}{ }\PYG{l+m}{0.5}\PYG{+w}{ }\PYG{n}{weight}\PYG{+w}{ }\PYG{n}{store}\PYG{+w}{ }\PYG{n}{balance}
\PYG{k}{balance}\PYG{+w}{ }\PYG{l+m}{1.0}\PYG{+w}{ }\PYG{n}{shift}\PYG{+w}{ }\PYG{n}{x}\PYG{+w}{ }\PYG{l+m}{20}\PYG{+w}{ }\PYG{l+m}{1.0}\PYG{+w}{ }\PYG{n}{out}\PYG{+w}{ }\PYG{n}{tmp.balance}
\end{sphinxVerbatim}


\subsection{Description}
\label{\detokenize{balance:description}}
\sphinxAtStartPar
This command adjusts the size and shape of processor subdomains
within the simulation box, to attempt to balance the number of atoms
or particles and thus indirectly the computational cost (load) more
evenly across processors.  The load balancing is “static” in the sense
that this command performs the balancing once, before or between
simulations.  The processor subdomains will then remain static during
the subsequent run.  To perform “dynamic” balancing, see the {\hyperref[\detokenize{fix_balance::doc}]{\sphinxcrossref{\DUrole{doc}{fix
balance}}}} command, which can adjust processor subdomain
sizes and shapes on\sphinxhyphen{}the\sphinxhyphen{}fly during a {\hyperref[\detokenize{run::doc}]{\sphinxcrossref{\DUrole{doc}{run}}}}.

\sphinxAtStartPar
Load\sphinxhyphen{}balancing is typically most useful if the particles in the
simulation box have a spatially\sphinxhyphen{}varying density distribution or when
the computational cost varies significantly between different
particles.  E.g. a model of a vapor/liquid interface, or a solid with
an irregular\sphinxhyphen{}shaped geometry containing void regions, or {\hyperref[\detokenize{pair_hybrid::doc}]{\sphinxcrossref{\DUrole{doc}{hybrid
pair style simulations}}}} which combine pair styles with
different computational cost.  In these cases, the LAMMPS default of
dividing the simulation box volume into a regular\sphinxhyphen{}spaced grid of 3d
bricks, with one equal\sphinxhyphen{}volume subdomain per processor, may assign
numbers of particles per processor in a way that the computational
effort varies significantly.  This can lead to poor performance when
the simulation is run in parallel.

\sphinxAtStartPar
The balancing can be performed with or without per\sphinxhyphen{}particle weighting.
With no weighting, the balancing attempts to assign an equal number of
particles to each processor.  With weighting, the balancing attempts
to assign an equal aggregate computational weight to each processor,
which typically induces a different number of atoms assigned to each
processor.  Details on the various weighting options and examples for
how they can be used are {\hyperref[\detokenize{balance:weighted-balance}]{\sphinxcrossref{\DUrole{std,std-ref}{given below}}}}.

\sphinxAtStartPar
Note that the {\hyperref[\detokenize{processors::doc}]{\sphinxcrossref{\DUrole{doc}{processors}}}} command allows some
control over how the box volume is split across processors.
Specifically, for a Px by Py by Pz grid of processors, it allows
choice of Px, Py, and Pz, subject to the constraint that Px * Py *
Pz = P, the total number of processors.  This is sufficient to achieve
good load\sphinxhyphen{}balance for some problems on some processor counts.
However, all the processor subdomains will still have the same shape
and same volume.

\sphinxAtStartPar
The requested load\sphinxhyphen{}balancing operation is only performed if the
current “imbalance factor” in particles owned by each processor
exceeds the specified \sphinxstyleemphasis{thresh} parameter.  The imbalance factor is
defined as the maximum number of particles (or weight) owned by any
processor, divided by the average number of particles (or weight) per
processor.  Thus an imbalance factor of 1.0 is perfect balance.

\sphinxAtStartPar
As an example, for 10000 particles running on 10 processors, if the
most heavily loaded processor has 1200 particles, then the factor is
1.2, meaning there is a 20\% imbalance.  Note that a re\sphinxhyphen{}balance can be
forced even if the current balance is perfect (1.0) be specifying a
\sphinxstyleemphasis{thresh} \textless{} 1.0.

\begin{sphinxadmonition}{note}{Note:}
\sphinxAtStartPar
Balancing is performed even if the imbalance factor does not
exceed the \sphinxstyleemphasis{thresh} parameter if a “grid” style is specified when the
current partitioning is “tiled”.  The meaning of “grid” vs “tiled” is
explained below.  This is to allow forcing of the partitioning to
“grid” so that the {\hyperref[\detokenize{comm_style::doc}]{\sphinxcrossref{\DUrole{doc}{comm\_style brick}}}} command can then
be used to replace a current {\hyperref[\detokenize{comm_style::doc}]{\sphinxcrossref{\DUrole{doc}{comm\_style tiled}}}}
setting.
\end{sphinxadmonition}

\sphinxAtStartPar
When the balance command completes, it prints statistics about the
result, including the change in the imbalance factor and the change in
the maximum number of particles on any processor.  For “grid” methods
(defined below) that create a logical 3d grid of processors, the
positions of all cutting planes in each of the 3 dimensions (as
fractions of the box length) are also printed.

\begin{sphinxadmonition}{note}{Note:}
\sphinxAtStartPar
This command attempts to minimize the imbalance factor, as
defined above.  But depending on the method a perfect balance (1.0)
may not be achieved.  For example, “grid” methods (defined below) that
create a logical 3d grid cannot achieve perfect balance for many
irregular distributions of particles.  Likewise, if a portion of the
system is a perfect lattice, e.g. the initial system is generated by
the {\hyperref[\detokenize{create_atoms::doc}]{\sphinxcrossref{\DUrole{doc}{create\_atoms}}}} command, then “grid” methods may
be unable to achieve exact balance.  This is because entire lattice
planes will be owned or not owned by a single processor.
\end{sphinxadmonition}

\begin{sphinxadmonition}{note}{Note:}
\sphinxAtStartPar
The imbalance factor is also an estimate of the maximum speed\sphinxhyphen{}up
you can hope to achieve by running a perfectly balanced simulation
versus an imbalanced one.  In the example above, the 10000 particle
simulation could run up to 20\% faster if it were perfectly balanced,
versus when imbalanced.  However, computational cost is not strictly
proportional to particle count, and changing the relative size and
shape of processor subdomains may lead to additional computational
and communication overheads, e.g. in the PPPM solver used via the
{\hyperref[\detokenize{kspace_style::doc}]{\sphinxcrossref{\DUrole{doc}{kspace\_style}}}} command.  Thus you should benchmark
the run times of a simulation before and after balancing.
\end{sphinxadmonition}


\bigskip\hrule\bigskip


\sphinxAtStartPar
The method used to perform a load balance is specified by one of the
listed styles (or more in the case of \sphinxstyleemphasis{x},\sphinxstyleemphasis{y},\sphinxstyleemphasis{z}), which are
described in detail below.  There are 2 kinds of styles.

\sphinxAtStartPar
The \sphinxstyleemphasis{x}, \sphinxstyleemphasis{y}, \sphinxstyleemphasis{z}, and \sphinxstyleemphasis{shift} styles are “grid” methods which
produce a logical 3d grid of processors.  They operate by changing the
cutting planes (or lines) between processors in 3d (or 2d), to adjust
the volume (area in 2d) assigned to each processor, as in the
following 2d diagram where processor subdomains are shown and
particles are colored by the processor that owns them.

\sphinxAtStartPar
\sphinxincludegraphics[width=0.320\linewidth]{{balance_uniform}.jpg}  \sphinxincludegraphics[width=0.320\linewidth]{{balance_nonuniform}.jpg}  \sphinxincludegraphics[width=0.320\linewidth]{{balance_rcb}.jpg}

\sphinxAtStartPar
The leftmost diagram is the default partitioning of the simulation box
across processors (one sub\sphinxhyphen{}box for each of 16 processors); the middle
diagram is after a “grid” method has been applied.  The \sphinxstyleemphasis{rcb} style is
a “tiling” method which does not produce a logical 3d grid of
processors.  Rather it tiles the simulation domain with rectangular
sub\sphinxhyphen{}boxes of varying size and shape in an irregular fashion so as to
have equal numbers of particles (or weight) in each sub\sphinxhyphen{}box, as in the
rightmost diagram above.

\sphinxAtStartPar
The “grid” methods can be used with either of the {\hyperref[\detokenize{comm_style::doc}]{\sphinxcrossref{\DUrole{doc}{comm\_style}}}} command options, \sphinxstyleemphasis{brick} or \sphinxstyleemphasis{tiled}.  The “tiling”
methods can only be used with {\hyperref[\detokenize{comm_style::doc}]{\sphinxcrossref{\DUrole{doc}{comm\_style tiled}}}}.
Note that it can be useful to use a “grid” method with
{\hyperref[\detokenize{comm_style::doc}]{\sphinxcrossref{\DUrole{doc}{comm\_style tiled}}}} to return the domain partitioning
to a logical 3d grid of processors so that “comm\_style brick” can
afterwords be specified for subsequent {\hyperref[\detokenize{run::doc}]{\sphinxcrossref{\DUrole{doc}{run}}}} commands.

\sphinxAtStartPar
When a “grid” method is specified, the current domain partitioning can
be either a logical 3d grid or a tiled partitioning.  In the former
case, the current logical 3d grid is used as a starting point and
changes are made to improve the imbalance factor.  In the latter case,
the tiled partitioning is discarded and a logical 3d grid is created
with uniform spacing in all dimensions.  This becomes the starting
point for the balancing operation.

\sphinxAtStartPar
When a “tiling” method is specified, the current domain partitioning
(“grid” or “tiled”) is ignored, and a new partitioning is computed
from scratch.


\bigskip\hrule\bigskip


\sphinxAtStartPar
The \sphinxstyleemphasis{x}, \sphinxstyleemphasis{y}, and \sphinxstyleemphasis{z} styles invoke a “grid” method for balancing, as
described above.  Note that any or all of these 3 styles can be
specified together, one after the other, but they cannot be used with
any other style.  This style adjusts the position of cutting planes
between processor subdomains in specific dimensions.  Only the
specified dimensions are altered.

\sphinxAtStartPar
The \sphinxstyleemphasis{uniform} argument spaces the planes evenly, as in the left
diagrams above.  The \sphinxstyleemphasis{numeric} argument requires listing Ps\sphinxhyphen{}1 numbers
that specify the position of the cutting planes.  This requires
knowing Ps = Px or Py or Pz = the number of processors assigned by
LAMMPS to the relevant dimension.  This assignment is made (and the
Px, Py, Pz values printed out) when the simulation box is created by
the “create\_box” or “read\_data” or “read\_restart” command and is
influenced by the settings of the {\hyperref[\detokenize{processors::doc}]{\sphinxcrossref{\DUrole{doc}{processors}}}}
command.

\sphinxAtStartPar
Each of the numeric values must be between 0 and 1, and they must be
listed in ascending order.  They represent the fractional position of
the cutting place.  The left (or lower) edge of the box is 0.0, and
the right (or upper) edge is 1.0.  Neither of these values is
specified.  Only the interior Ps\sphinxhyphen{}1 positions are specified.  Thus is
there are 2 processors in the x dimension, you specify a single value
such as 0.75, which would make the left processor’s subdomain 3x
larger than the right processor’s subdomain.


\bigskip\hrule\bigskip


\sphinxAtStartPar
The \sphinxstyleemphasis{shift} style invokes a “grid” method for balancing, as
described above.  It changes the positions of cutting planes between
processors in an iterative fashion, seeking to reduce the imbalance
factor, similar to how the {\hyperref[\detokenize{fix_balance::doc}]{\sphinxcrossref{\DUrole{doc}{fix balance shift}}}}
command operates.

\sphinxAtStartPar
The \sphinxstyleemphasis{dimstr} argument is a string of characters, each of which must be
an “x” or “y” or “z”. Each character can appear zero or one time,
since there is no advantage to balancing on a dimension more than
once.  You should normally only list dimensions where you expect there
to be a density variation in the particles.

\sphinxAtStartPar
Balancing proceeds by adjusting the cutting planes in each of the
dimensions listed in \sphinxstyleemphasis{dimstr}, one dimension at a time.  For a single
dimension, the balancing operation (described below) is iterated on up
to \sphinxstyleemphasis{Niter} times.  After each dimension finishes, the imbalance factor
is re\sphinxhyphen{}computed, and the balancing operation halts if the \sphinxstyleemphasis{stopthresh}
criterion is met.

\sphinxAtStartPar
A re\sphinxhyphen{}balance operation in a single dimension is performed using a
recursive multisectioning algorithm, where the position of each
cutting plane (line in 2d) in the dimension is adjusted independently.
This is similar to a recursive bisectioning for a single value, except
that the bounds used for each bisectioning take advantage of
information from neighboring cuts if possible.  At each iteration, the
count of particles on either side of each plane is tallied.  If the
counts do not match the target value for the plane, the position of
the cut is adjusted to be halfway between a low and high bound.  The
low and high bounds are adjusted on each iteration, using new count
information, so that they become closer together over time.  Thus as
the recursion progresses, the count of particles on either side of the
plane gets closer to the target value.

\sphinxAtStartPar
After the balanced plane positions are determined, if any pair of
adjacent planes are closer together than the neighbor skin distance
(as specified by the {\hyperref[\detokenize{neigh_modify::doc}]{\sphinxcrossref{\DUrole{doc}{neigh\_modify}}}} command), then
the plane positions are shifted to separate them by at least this
amount.  This is to prevent particles being lost when dynamics are run
with processor subdomains that are too narrow in one or more
dimensions.

\sphinxAtStartPar
Once the re\sphinxhyphen{}balancing is complete and final processor subdomains
assigned, particles are migrated to their new owning processor, and
the balance procedure ends.

\begin{sphinxadmonition}{note}{Note:}
\sphinxAtStartPar
At each re\sphinxhyphen{}balance operation, the bisectioning for each cutting
plane (line in 2d) typically starts with low and high bounds separated
by the extent of a processor’s subdomain in one dimension.  The size
of this bracketing region shrinks by 1/2 every iteration.  Thus if
\sphinxstyleemphasis{Niter} is specified as 10, the cutting plane will typically be
positioned to 1 part in 1000 accuracy (relative to the perfect target
position).  For \sphinxstyleemphasis{Niter} = 20, it will be accurate to 1 part in a
million.  Thus there is no need to set \sphinxstyleemphasis{Niter} to a large value.
LAMMPS will check if the threshold accuracy is reached (in a
dimension) is less iterations than \sphinxstyleemphasis{Niter} and exit early.  However,
\sphinxstyleemphasis{Niter} should also not be set too small, since it will take roughly
the same number of iterations to converge even if the cutting plane is
initially close to the target value.
\end{sphinxadmonition}


\bigskip\hrule\bigskip


\sphinxAtStartPar
The \sphinxstyleemphasis{rcb} style invokes a “tiled” method for balancing, as described
above.  It performs a recursive coordinate bisectioning (RCB) of the
simulation domain. The basic idea is as follows.

\sphinxAtStartPar
The simulation domain is cut into 2 boxes by an axis\sphinxhyphen{}aligned cut in
one of the dimensions, leaving one new sub\sphinxhyphen{}box on either side of the
cut.  Which dimension is chosen for the cut depends on the particle
(weight) distribution within the parent box.  Normally the longest
dimension of the box is cut, but if all (or most) of the particles are
at one end of the box, a cut may be performed in another dimension to
induce sub\sphinxhyphen{}boxes that are more cube\sphinxhyphen{}ish (3d) or square\sphinxhyphen{}ish (2d) in
shape.

\sphinxAtStartPar
After the cut is made, all the processors are also partitioned into 2
groups, half assigned to the box on the lower side of the cut, and
half to the box on the upper side.  (If the processor count is odd,
one side gets an extra processor.)  The cut is positioned so that the
number of (weighted) particles in the lower box is exactly the number
that the processors assigned to that box should own for load balance
to be perfect.  This also makes load balance for the upper box
perfect.  The positioning of the cut is done iteratively, by a
bisectioning method (median search).  Note that counting particles on
either side of the cut requires communication between all processors
at each iteration.

\sphinxAtStartPar
That is the procedure for the first cut.  Subsequent cuts are made
recursively, in exactly the same manner.  The subset of processors
assigned to each box make a new cut in one dimension of that box,
splitting the box, the subset of processors, and the particles in the
box in two.  The recursion continues until every processor is assigned
a sub\sphinxhyphen{}box of the entire simulation domain, and owns the (weighted)
particles in that sub\sphinxhyphen{}box.


\bigskip\hrule\bigskip

\phantomsection\label{\detokenize{balance:weighted-balance}}
\sphinxAtStartPar
This subsection describes how to perform weighted load balancing
using the \sphinxstyleemphasis{weight} keyword.

\sphinxAtStartPar
By default, all particles have a weight of 1.0, which means each
particle is assumed to require the same amount of computation during a
timestep.  There are, however, scenarios where this is not a good
assumption.  Measuring the computational cost for each particle
accurately would be impractical and slow down the computation.
Instead the \sphinxstyleemphasis{weight} keyword implements several ways to influence the
per\sphinxhyphen{}particle weights empirically by properties readily available or
using the user’s knowledge of the system.  Note that the absolute
value of the weights are not important; only their relative ratios
affect which particle is assigned to which processor.  A particle with
a weight of 2.5 is assumed to require 5x more computational than a
particle with a weight of 0.5.  For all the options below the weight
assigned to a particle must be a positive value; an error will be be
generated if a weight is \textless{}= 0.0.

\sphinxAtStartPar
Below is a list of possible weight options with a short description of
their usage and some example scenarios where they might be applicable.
It is possible to apply multiple weight flags and the weightings they
induce will be combined through multiplication.  Most of the time,
however, it is sufficient to use just one method.

\sphinxAtStartPar
The \sphinxstyleemphasis{group} weight style assigns weight factors to specified
{\hyperref[\detokenize{group::doc}]{\sphinxcrossref{\DUrole{doc}{groups}}}} of particles.  The \sphinxstyleemphasis{group} style keyword is
followed by the number of groups, then pairs of group IDs and the
corresponding weight factor.  If a particle belongs to none of the
specified groups, its weight is not changed.  If it belongs to
multiple groups, its weight is the product of the weight factors.

\sphinxAtStartPar
This weight style is useful in combination with pair style
{\hyperref[\detokenize{pair_hybrid::doc}]{\sphinxcrossref{\DUrole{doc}{hybrid}}}}, e.g. when combining a more costly many\sphinxhyphen{}body
potential with a fast pairwise potential.  It is also useful when
using {\hyperref[\detokenize{run_style::doc}]{\sphinxcrossref{\DUrole{doc}{run\_style respa}}}} where some portions of the
system have many bonded interactions and others none.  It assumes that
the computational cost for each group remains constant over time.
This is a purely empirical weighting, so a series test runs to tune
the assigned weight factors for optimal performance is recommended.

\sphinxAtStartPar
The \sphinxstyleemphasis{neigh} weight style assigns the same weight to each particle
owned by a processor based on the total count of neighbors in the
neighbor list owned by that processor.  The motivation is that more
neighbors means a higher computational cost.  The style does not use
neighbors per atom to assign a unique weight to each atom, because
that value can vary depending on how the neighbor list is built.

\sphinxAtStartPar
The \sphinxstyleemphasis{factor} setting is applied as an overall scale factor to the
\sphinxstyleemphasis{neigh} weights which allows adjustment of their impact on the
balancing operation.  The specified \sphinxstyleemphasis{factor} value must be positive.
A value \textgreater{} 1.0 will increase the weights so that the ratio of max
weight to min weight increases by \sphinxstyleemphasis{factor}.  A value \textless{} 1.0 will
decrease the weights so that the ratio of max weight to min weight
decreases by \sphinxstyleemphasis{factor}.  In both cases the intermediate weight values
increase/decrease proportionally as well.  A value = 1.0 has no effect
on the \sphinxstyleemphasis{neigh} weights.  As a rule of thumb, we have found a \sphinxstyleemphasis{factor}
of about 0.8 often results in the best performance, since the number
of neighbors is likely to overestimate the ideal weight.

\sphinxAtStartPar
This weight style is useful for systems where there are different
cutoffs used for different pairs of interactions, or the density
fluctuates, or a large number of particles are in the vicinity of a
wall, or a combination of these effects.  If a simulation uses
multiple neighbor lists, this weight style will use the first suitable
neighbor list it finds.  It will not request or compute a new list.  A
warning will be issued if there is no suitable neighbor list available
or if it is not current, e.g. if the balance command is used before a
{\hyperref[\detokenize{run::doc}]{\sphinxcrossref{\DUrole{doc}{run}}}} or {\hyperref[\detokenize{minimize::doc}]{\sphinxcrossref{\DUrole{doc}{minimize}}}} command is used, in which
case the neighbor list may not yet have been built.  In this case no
weights are computed.  Inserting a {\hyperref[\detokenize{run::doc}]{\sphinxcrossref{\DUrole{doc}{run 0 post no}}}} command
before issuing the \sphinxstyleemphasis{balance} command, may be a workaround for this
case, as it will induce the neighbor list to be built.

\sphinxAtStartPar
The \sphinxstyleemphasis{time} weight style uses {\hyperref[\detokenize{timer::doc}]{\sphinxcrossref{\DUrole{doc}{timer data}}}} to estimate
weights.  It assigns the same weight to each particle owned by a
processor based on the total computational time spent by that
processor.  See details below on what time window is used.  It uses
the same timing information as is used for the {\hyperref[\detokenize{Run_output::doc}]{\sphinxcrossref{\DUrole{doc}{MPI task timing
breakdown}}}}, namely, for sections \sphinxstyleemphasis{Pair}, \sphinxstyleemphasis{Bond},
\sphinxstyleemphasis{Kspace}, and \sphinxstyleemphasis{Neigh}.  The time spent in those portions of the
timestep are measured for each MPI rank, summed, then divided by the
number of particles owned by that processor.  I.e. the weight is an
effective CPU time/particle averaged over the particles on that
processor.

\sphinxAtStartPar
The \sphinxstyleemphasis{factor} setting is applied as an overall scale factor to the
\sphinxstyleemphasis{time} weights which allows adjustment of their impact on the
balancing operation.  The specified \sphinxstyleemphasis{factor} value must be positive.
A value \textgreater{} 1.0 will increase the weights so that the ratio of max
weight to min weight increases by \sphinxstyleemphasis{factor}.  A value \textless{} 1.0 will
decrease the weights so that the ratio of max weight to min weight
decreases by \sphinxstyleemphasis{factor}.  In both cases the intermediate weight values
increase/decrease proportionally as well.  A value = 1.0 has no effect
on the \sphinxstyleemphasis{time} weights.  As a rule of thumb, effective values to use
are typically between 0.5 and 1.2.  Note that the timer quantities
mentioned above can be affected by communication which occurs in the
middle of the operations, e.g. pair styles with intermediate exchange
of data witin the force computation, and likewise for KSpace solves.

\sphinxAtStartPar
When using the \sphinxstyleemphasis{time} weight style with the \sphinxstyleemphasis{balance} command, the
timing data is taken from the preceding run command, i.e. the timings
are for the entire previous run.  For the \sphinxstyleemphasis{fix balance} command the
timing data is for only the timesteps since the last balancing
operation was performed.  If timing information for the required
sections is not available, e.g. at the beginning of a run, or when the
{\hyperref[\detokenize{timer::doc}]{\sphinxcrossref{\DUrole{doc}{timer}}}} command is set to either \sphinxstyleemphasis{loop} or \sphinxstyleemphasis{off}, a warning
is issued.  In this case no weights are computed.

\begin{sphinxadmonition}{note}{Note:}
\sphinxAtStartPar
The \sphinxstyleemphasis{time} weight style is the most generic option, and should
be tried first, unless the \sphinxstyleemphasis{group} style is easily applicable.
However, since the computed cost function is averaged over all
particles on a processor, the weights may not be highly accurate.
This style can also be effective as a secondary weight in combination
with either \sphinxstyleemphasis{group} or \sphinxstyleemphasis{neigh} to offset some of inaccuracies in
either of those heuristics.
\end{sphinxadmonition}

\sphinxAtStartPar
The \sphinxstyleemphasis{var} weight style assigns per\sphinxhyphen{}particle weights by evaluating an
{\hyperref[\detokenize{variable::doc}]{\sphinxcrossref{\DUrole{doc}{atom\sphinxhyphen{}style variable}}}} specified by \sphinxstyleemphasis{name}.  This is
provided as a more flexible alternative to the \sphinxstyleemphasis{group} weight style,
allowing definition of a more complex heuristics based on information
(global and per atom) available inside of LAMMPS.  For example,
atom\sphinxhyphen{}style variables can reference the position of a particle, its
velocity, the volume of its Voronoi cell, etc.

\sphinxAtStartPar
The \sphinxstyleemphasis{store} weight style does not compute a weight factor.  Instead it
stores the current accumulated weights in a custom per\sphinxhyphen{}atom vector
specified by \sphinxstyleemphasis{name}.  This must be a vector defined as \sphinxstyleemphasis{d\_name} via
the {\hyperref[\detokenize{fix_property_atom::doc}]{\sphinxcrossref{\DUrole{doc}{fix property/atom}}}} command.  This means
the values in the vector can be read as part of a data file with the
{\hyperref[\detokenize{read_data::doc}]{\sphinxcrossref{\DUrole{doc}{read\_data}}}} command or specified with the {\hyperref[\detokenize{set::doc}]{\sphinxcrossref{\DUrole{doc}{set}}}} command.  These weights can also be output in a {\hyperref[\detokenize{dump::doc}]{\sphinxcrossref{\DUrole{doc}{dump}}}} file, so this is a way to examine, debug, or visualize the
per\sphinxhyphen{}particle weights used during the load\sphinxhyphen{}balancing operation.

\sphinxAtStartPar
Note that the name of the custom per\sphinxhyphen{}atom vector is specified just
as \sphinxstyleemphasis{name}, not as \sphinxstyleemphasis{d\_name} as it is for other commands that use
different kinds of custom atom vectors or arrays as arguments.


\bigskip\hrule\bigskip


\sphinxAtStartPar
The \sphinxstyleemphasis{sort} keyword determines whether the communication of per\sphinxhyphen{}atom
data to other processors during load\sphinxhyphen{}balancing will be random or
deterministic.  Random is generally faster; deterministic will ensure
the new ordering of atoms on each processor is the same each time the
same simulation is run.  This can be useful for debugging purposes.
Since the balance command is a one\sphinxhyphen{}time operation, the default is
\sphinxstyleemphasis{yes} to perform sorting.

\sphinxAtStartPar
The \sphinxstyleemphasis{out} keyword writes a text file to the specified \sphinxstyleemphasis{filename} with
the results of the balancing operation.  The file contains the bounds
of the subdomain for each processor after the balancing operation
completes.  The format of the file is compatible with the
\sphinxhref{https://lammps.github.io/pizza}{Pizza.py} \sphinxstyleemphasis{mdump} tool which has support for manipulating and
visualizing mesh files.  An example is shown here for a balancing by 4
processors for a 2d problem:

\begin{sphinxVerbatim}[commandchars=\\\{\}]
ITEM: TIMESTEP
0
ITEM: NUMBER OF NODES
16
ITEM: BOX BOUNDS
0 10
0 10
0 10
ITEM: NODES
1 1 0 0 0
2 1 5 0 0
3 1 5 5 0
4 1 0 5 0
5 1 5 0 0
6 1 10 0 0
7 1 10 5 0
8 1 5 5 0
9 1 0 5 0
10 1 5 5 0
11 1 5 10 0
12 1 10 5 0
13 1 5 5 0
14 1 10 5 0
15 1 10 10 0
16 1 5 10 0
ITEM: TIMESTEP
0
ITEM: NUMBER OF SQUARES
4
ITEM: SQUARES
1 1 1 2 3 4
2 1 5 6 7 8
3 1 9 10 11 12
4 1 13 14 15 16
\end{sphinxVerbatim}

\sphinxAtStartPar
The coordinates of all the vertices are listed in the NODES section, 5
per processor.  Note that the 4 subdomains share vertices, so there
will be duplicate nodes in the list.

\sphinxAtStartPar
The “SQUARES” section lists the node IDs of the 4 vertices in a
rectangle for each processor (1 to 4).

\sphinxAtStartPar
For a 3d problem, the syntax is similar with 8 vertices listed for
each processor, instead of 4, and “SQUARES” replaced by “CUBES”.


\bigskip\hrule\bigskip



\subsection{Restrictions}
\label{\detokenize{balance:restrictions}}
\sphinxAtStartPar
For 2d simulations, the \sphinxstyleemphasis{z} style cannot be used.  Nor can a “z”
appear in \sphinxstyleemphasis{dimstr} for the \sphinxstyleemphasis{shift} style.

\sphinxAtStartPar
Balancing through recursive bisectioning (\sphinxstyleemphasis{rcb} style) requires
{\hyperref[\detokenize{comm_style::doc}]{\sphinxcrossref{\DUrole{doc}{comm\_style tiled}}}}


\subsection{Related commands}
\label{\detokenize{balance:related-commands}}
\sphinxAtStartPar
{\hyperref[\detokenize{group::doc}]{\sphinxcrossref{\DUrole{doc}{group}}}}, {\hyperref[\detokenize{processors::doc}]{\sphinxcrossref{\DUrole{doc}{processors}}}},
{\hyperref[\detokenize{fix_balance::doc}]{\sphinxcrossref{\DUrole{doc}{fix balance}}}}, {\hyperref[\detokenize{comm_style::doc}]{\sphinxcrossref{\DUrole{doc}{comm\_style}}}}


\subsection{Default}
\label{\detokenize{balance:default}}
\sphinxAtStartPar
The default setting is sort = yes.

\sphinxstepscope

\index{bond\_coeff@\spxentry{bond\_coeff}}\ignorespaces 

\section{bond\_coeff command}
\label{\detokenize{bond_coeff:bond-coeff-command}}\label{\detokenize{bond_coeff:index-0}}\label{\detokenize{bond_coeff::doc}}

\subsection{Syntax}
\label{\detokenize{bond_coeff:syntax}}
\begin{sphinxVerbatim}[commandchars=\\\{\}]
\PYG{k}{bond\PYGZus{}coeff}\PYG{+w}{ }\PYG{n}{N}\PYG{+w}{ }\PYG{n}{args}
\end{sphinxVerbatim}
\begin{itemize}
\item {} 
\sphinxAtStartPar
N = numeric bond type (see asterisk form below), or type label

\item {} 
\sphinxAtStartPar
args = coefficients for one or more bond types

\end{itemize}


\subsection{Examples}
\label{\detokenize{bond_coeff:examples}}
\begin{sphinxVerbatim}[commandchars=\\\{\}]
\PYG{k}{bond\PYGZus{}coeff}\PYG{+w}{ }\PYG{l+m}{5}\PYG{+w}{ }\PYG{l+m}{80.0}\PYG{+w}{ }\PYG{l+m}{1.2}
\PYG{k}{bond\PYGZus{}coeff}\PYG{+w}{ }\PYG{o}{*}\PYG{+w}{ }\PYG{l+m}{30.0}\PYG{+w}{ }\PYG{l+m}{1.5}\PYG{+w}{ }\PYG{l+m}{1.0}\PYG{+w}{ }\PYG{l+m}{1.0}
\PYG{k}{bond\PYGZus{}coeff}\PYG{+w}{ }\PYG{l+m}{1}\PYG{o}{*}\PYG{l+m}{4}\PYG{+w}{ }\PYG{l+m}{30.0}\PYG{+w}{ }\PYG{l+m}{1.5}\PYG{+w}{ }\PYG{l+m}{1.0}\PYG{+w}{ }\PYG{l+m}{1.0}
\PYG{k}{bond\PYGZus{}coeff}\PYG{+w}{ }\PYG{l+m}{1}\PYG{+w}{ }\PYG{n}{harmonic}\PYG{+w}{ }\PYG{l+m}{200.0}\PYG{+w}{ }\PYG{l+m}{1.0}\PYG{+w}{ }\PYG{c}{\PYGZsh{} (for bond\PYGZus{}style hybrid)}

\PYG{k}{labelmap}\PYG{+w}{ }\PYG{n}{bond}\PYG{+w}{ }\PYG{l+m}{5}\PYG{+w}{ }\PYG{n}{carbonyl}
\PYG{k}{bond\PYGZus{}coeff}\PYG{+w}{ }\PYG{n}{carbonyl}\PYG{+w}{ }\PYG{l+m}{80.0}\PYG{+w}{ }\PYG{l+m}{1.2}
\end{sphinxVerbatim}


\subsection{Description}
\label{\detokenize{bond_coeff:description}}
\sphinxAtStartPar
Specify the bond force field coefficients for one or more bond types.
The number and meaning of the coefficients depends on the bond style.
Bond coefficients can also be set in the data file read by the
{\hyperref[\detokenize{read_data::doc}]{\sphinxcrossref{\DUrole{doc}{read\_data}}}} command or in a restart file.

\sphinxAtStartPar
\(N\) can be specified in one of several ways.  An explicit numeric
value can be used, as in the first example above.  Or \(N\) can be a
type label, which is an alphanumeric string defined by the
{\hyperref[\detokenize{labelmap::doc}]{\sphinxcrossref{\DUrole{doc}{labelmap}}}} command or in a section of a data file read
by the {\hyperref[\detokenize{read_data::doc}]{\sphinxcrossref{\DUrole{doc}{read\_data}}}} command.

\sphinxAtStartPar
For numeric values only, a wild\sphinxhyphen{}card asterisk can be used to set the
coefficients for multiple bond types.  This takes the form “*” or “*n”
or “n*” or “m*n”.  If \(N\) is the number of bond types, then an
asterisk with no numeric values means all types from 1 to \(N\).  A
leading asterisk means all types from 1 to n (inclusive).  A trailing
asterisk means all types from n to \(N\) (inclusive).  A middle
asterisk means all types from m to n (inclusive).

\sphinxAtStartPar
Note that using a bond\_coeff command can override a previous setting
for the same bond type.  For example, these commands set the coeffs
for all bond types, then overwrite the coeffs for just bond type 2:

\begin{sphinxVerbatim}[commandchars=\\\{\}]
\PYG{k}{bond\PYGZus{}coeff}\PYG{+w}{ }\PYG{o}{*}\PYG{+w}{ }\PYG{l+m}{100.0}\PYG{+w}{ }\PYG{l+m}{1.2}
\PYG{k}{bond\PYGZus{}coeff}\PYG{+w}{ }\PYG{l+m}{2}\PYG{+w}{ }\PYG{l+m}{200.0}\PYG{+w}{ }\PYG{l+m}{1.2}
\end{sphinxVerbatim}

\sphinxAtStartPar
A line in a data file that specifies bond coefficients uses the exact
same format as the arguments of the bond\_coeff command in an input
script, except that wild\sphinxhyphen{}card asterisks should not be used since
coefficients for all \(N\) types must be listed in the file.  For
example, under the “Bond Coeffs” section of a data file, the line that
corresponds to the first example above would be listed as

\begin{sphinxVerbatim}[commandchars=\\\{\}]
5 80.0 1.2
\end{sphinxVerbatim}


\bigskip\hrule\bigskip


\sphinxAtStartPar
The list of all bond styles defined in LAMMPS is given on the
{\hyperref[\detokenize{bond_style::doc}]{\sphinxcrossref{\DUrole{doc}{bond\_style}}}} doc page.  They are also listed in more
compact form on the {\hyperref[\detokenize{Commands_bond::doc}]{\sphinxcrossref{\DUrole{doc}{Commands bond}}}} doc page.

\sphinxAtStartPar
On either of those pages, click on the style to display the formula it
computes and its coefficients as specified by the associated
bond\_coeff command.


\bigskip\hrule\bigskip



\subsection{Restrictions}
\label{\detokenize{bond_coeff:restrictions}}
\sphinxAtStartPar
This command must come after the simulation box is defined by a
{\hyperref[\detokenize{read_data::doc}]{\sphinxcrossref{\DUrole{doc}{read\_data}}}}, {\hyperref[\detokenize{read_restart::doc}]{\sphinxcrossref{\DUrole{doc}{read\_restart}}}}, or
{\hyperref[\detokenize{create_box::doc}]{\sphinxcrossref{\DUrole{doc}{create\_box}}}} command.

\sphinxAtStartPar
A bond style must be defined before any bond coefficients are set,
either in the input script or in a data file.


\subsection{Related commands}
\label{\detokenize{bond_coeff:related-commands}}
\sphinxAtStartPar
{\hyperref[\detokenize{bond_style::doc}]{\sphinxcrossref{\DUrole{doc}{bond\_style}}}}


\subsection{Default}
\label{\detokenize{bond_coeff:default}}
\sphinxAtStartPar
none

\sphinxstepscope

\index{bond\_style@\spxentry{bond\_style}}\ignorespaces 

\section{bond\_style command}
\label{\detokenize{bond_style:bond-style-command}}\label{\detokenize{bond_style:index-0}}\label{\detokenize{bond_style::doc}}

\subsection{Syntax}
\label{\detokenize{bond_style:syntax}}
\begin{sphinxVerbatim}[commandchars=\\\{\}]
\PYG{k}{bond\PYGZus{}style}\PYG{+w}{ }\PYG{n}{style}\PYG{+w}{ }\PYG{n}{args}
\end{sphinxVerbatim}
\begin{itemize}
\item {} 
\sphinxAtStartPar
style = \sphinxstyleemphasis{none} or \sphinxstyleemphasis{zero} or \sphinxstyleemphasis{hybrid} or \sphinxstyleemphasis{bpm/rotational} or \sphinxstyleemphasis{bpm/spring} or \sphinxstyleemphasis{class2} or \sphinxstyleemphasis{fene} or \sphinxstyleemphasis{fene/expand} or \sphinxstyleemphasis{fene/nm} or \sphinxstyleemphasis{gaussian} or \sphinxstyleemphasis{gromos} or \sphinxstyleemphasis{harmonic} or \sphinxstyleemphasis{harmonic/restrain} \sphinxstyleemphasis{harmonic/shift} or \sphinxstyleemphasis{harmonic/shift/cut} or \sphinxstyleemphasis{lepton} or \sphinxstyleemphasis{morse} or \sphinxstyleemphasis{nonlinear} or \sphinxstyleemphasis{oxdna/fene} or \sphinxstyleemphasis{oxdena2/fene} or \sphinxstyleemphasis{oxrna2/fene} or \sphinxstyleemphasis{quartic} or \sphinxstyleemphasis{special} or \sphinxstyleemphasis{table}

\item {} 
\sphinxAtStartPar
args = none for any style except \sphinxstyleemphasis{hybrid}
\begin{itemize}
\item {} 
\sphinxAtStartPar
\sphinxstyleemphasis{hybrid} args = list of one or more styles

\end{itemize}

\end{itemize}


\subsection{Examples}
\label{\detokenize{bond_style:examples}}
\begin{sphinxVerbatim}[commandchars=\\\{\}]
\PYG{k}{bond\PYGZus{}style}\PYG{+w}{ }\PYG{n}{harmonic}
\PYG{k}{bond\PYGZus{}style}\PYG{+w}{ }\PYG{n}{fene}
\PYG{k}{bond\PYGZus{}style}\PYG{+w}{ }\PYG{n}{hybrid}\PYG{+w}{ }\PYG{n}{harmonic}\PYG{+w}{ }\PYG{n}{fene}
\end{sphinxVerbatim}


\subsection{Description}
\label{\detokenize{bond_style:description}}
\sphinxAtStartPar
Set the formula(s) LAMMPS uses to compute bond interactions between
pairs of atoms.  In LAMMPS, a bond differs from a pairwise
interaction, which are set via the {\hyperref[\detokenize{pair_style::doc}]{\sphinxcrossref{\DUrole{doc}{pair\_style}}}}
command.  Bonds are defined between specified pairs of atoms and
remain in force for the duration of the simulation (unless new bonds
are created or existing bonds break, which is possible in some fixes
and bond potentials).  The list of bonded atoms is read in by a
{\hyperref[\detokenize{read_data::doc}]{\sphinxcrossref{\DUrole{doc}{read\_data}}}} or {\hyperref[\detokenize{read_restart::doc}]{\sphinxcrossref{\DUrole{doc}{read\_restart}}}}
command from a data or restart file.  By contrast, pair potentials are
typically defined between all pairs of atoms within a cutoff distance
and the set of active interactions changes over time.

\sphinxAtStartPar
Hybrid models where bonds are computed using different bond potentials
can be setup using the \sphinxstyleemphasis{hybrid} bond style.

\sphinxAtStartPar
The coefficients associated with a bond style can be specified in a
data or restart file or via the {\hyperref[\detokenize{bond_coeff::doc}]{\sphinxcrossref{\DUrole{doc}{bond\_coeff}}}} command.

\sphinxAtStartPar
All bond potentials store their coefficient data in binary restart
files which means bond\_style and {\hyperref[\detokenize{bond_coeff::doc}]{\sphinxcrossref{\DUrole{doc}{bond\_coeff}}}} commands
do not need to be re\sphinxhyphen{}specified in an input script that restarts a
simulation.  See the {\hyperref[\detokenize{read_restart::doc}]{\sphinxcrossref{\DUrole{doc}{read\_restart}}}} command for
details on how to do this.  The one exception is that bond\_style
\sphinxstyleemphasis{hybrid} only stores the list of sub\sphinxhyphen{}styles in the restart file; bond
coefficients need to be re\sphinxhyphen{}specified.

\begin{sphinxadmonition}{note}{Note:}
\sphinxAtStartPar
When both a bond and pair style is defined, the
{\hyperref[\detokenize{special_bonds::doc}]{\sphinxcrossref{\DUrole{doc}{special\_bonds}}}} command often needs to be used to
turn off (or weight) the pairwise interaction that would otherwise
exist between two bonded atoms.
\end{sphinxadmonition}

\sphinxAtStartPar
In the formulas listed for each bond style, \sphinxstyleemphasis{r} is the distance
between the two atoms in the bond.


\bigskip\hrule\bigskip


\sphinxAtStartPar
Here is an alphabetic list of bond styles defined in LAMMPS.  Click on
the style to display the formula it computes and coefficients
specified by the associated {\hyperref[\detokenize{bond_coeff::doc}]{\sphinxcrossref{\DUrole{doc}{bond\_coeff}}}} command.

\sphinxAtStartPar
Click on the style to display the formula it computes, any additional
arguments specified in the bond\_style command, and coefficients
specified by the associated {\hyperref[\detokenize{bond_coeff::doc}]{\sphinxcrossref{\DUrole{doc}{bond\_coeff}}}} command.

\sphinxAtStartPar
There are also additional accelerated pair styles included in the
LAMMPS distribution for faster performance on CPUs, GPUs, and KNLs.
The individual style names on the {\hyperref[\detokenize{Commands_bond::doc}]{\sphinxcrossref{\DUrole{doc}{Commands bond}}}}
doc page are followed by one or more of (g,i,k,o,t) to indicate which
accelerated styles exist.
\begin{itemize}
\item {} 
\sphinxAtStartPar
{\hyperref[\detokenize{bond_none::doc}]{\sphinxcrossref{\DUrole{doc}{none}}}} \sphinxhyphen{} turn off bonded interactions

\item {} 
\sphinxAtStartPar
{\hyperref[\detokenize{bond_zero::doc}]{\sphinxcrossref{\DUrole{doc}{zero}}}} \sphinxhyphen{} topology but no interactions

\item {} 
\sphinxAtStartPar
{\hyperref[\detokenize{bond_hybrid::doc}]{\sphinxcrossref{\DUrole{doc}{hybrid}}}} \sphinxhyphen{} define multiple styles of bond interactions

\item {} 
\sphinxAtStartPar
{\hyperref[\detokenize{bond_bpm_rotational::doc}]{\sphinxcrossref{\DUrole{doc}{bpm/rotational}}}} \sphinxhyphen{} breakable bond with forces and torques based on deviation from reference state

\item {} 
\sphinxAtStartPar
{\hyperref[\detokenize{bond_bpm_spring::doc}]{\sphinxcrossref{\DUrole{doc}{bpm/spring}}}} \sphinxhyphen{} breakable bond with forces based on deviation from reference length

\item {} 
\sphinxAtStartPar
{\hyperref[\detokenize{bond_class2::doc}]{\sphinxcrossref{\DUrole{doc}{class2}}}} \sphinxhyphen{} COMPASS (class 2) bond

\item {} 
\sphinxAtStartPar
{\hyperref[\detokenize{bond_fene::doc}]{\sphinxcrossref{\DUrole{doc}{fene}}}} \sphinxhyphen{} FENE (finite\sphinxhyphen{}extensible non\sphinxhyphen{}linear elastic) bond

\item {} 
\sphinxAtStartPar
{\hyperref[\detokenize{bond_fene_expand::doc}]{\sphinxcrossref{\DUrole{doc}{fene/expand}}}} \sphinxhyphen{} FENE bonds with variable size particles

\item {} 
\sphinxAtStartPar
{\hyperref[\detokenize{bond_fene::doc}]{\sphinxcrossref{\DUrole{doc}{fene/nm}}}} \sphinxhyphen{} FENE bonds with a generalized Lennard\sphinxhyphen{}Jones potential

\item {} 
\sphinxAtStartPar
{\hyperref[\detokenize{bond_gaussian::doc}]{\sphinxcrossref{\DUrole{doc}{gaussian}}}} \sphinxhyphen{} multicentered Gaussian\sphinxhyphen{}based bond potential

\item {} 
\sphinxAtStartPar
{\hyperref[\detokenize{bond_gromos::doc}]{\sphinxcrossref{\DUrole{doc}{gromos}}}} \sphinxhyphen{} GROMOS force field bond

\item {} 
\sphinxAtStartPar
{\hyperref[\detokenize{bond_harmonic::doc}]{\sphinxcrossref{\DUrole{doc}{harmonic}}}} \sphinxhyphen{} harmonic bond

\item {} 
\sphinxAtStartPar
{\hyperref[\detokenize{bond_harmonic_restrain::doc}]{\sphinxcrossref{\DUrole{doc}{harmonic/restrain}}}} \sphinxhyphen{} harmonic bond to restrain to original bond distance

\item {} 
\sphinxAtStartPar
{\hyperref[\detokenize{bond_harmonic_shift::doc}]{\sphinxcrossref{\DUrole{doc}{harmonic/shift}}}} \sphinxhyphen{} shifted harmonic bond

\item {} 
\sphinxAtStartPar
{\hyperref[\detokenize{bond_harmonic_shift_cut::doc}]{\sphinxcrossref{\DUrole{doc}{harmonic/shift/cut}}}} \sphinxhyphen{} shifted harmonic bond with a cutoff

\item {} 
\sphinxAtStartPar
{\hyperref[\detokenize{bond_lepton::doc}]{\sphinxcrossref{\DUrole{doc}{lepton}}}} \sphinxhyphen{} bond potential from evaluating a string

\item {} 
\sphinxAtStartPar
{\hyperref[\detokenize{bond_mesocnt::doc}]{\sphinxcrossref{\DUrole{doc}{mesocnt}}}} \sphinxhyphen{} Harmonic bond wrapper with parameterization presets for nanotubes

\item {} 
\sphinxAtStartPar
{\hyperref[\detokenize{bond_mm3::doc}]{\sphinxcrossref{\DUrole{doc}{mm3}}}} \sphinxhyphen{} MM3 anharmonic bond

\item {} 
\sphinxAtStartPar
{\hyperref[\detokenize{bond_morse::doc}]{\sphinxcrossref{\DUrole{doc}{morse}}}} \sphinxhyphen{} Morse bond

\item {} 
\sphinxAtStartPar
{\hyperref[\detokenize{bond_nonlinear::doc}]{\sphinxcrossref{\DUrole{doc}{nonlinear}}}} \sphinxhyphen{} nonlinear bond

\item {} 
\sphinxAtStartPar
{\hyperref[\detokenize{bond_oxdna::doc}]{\sphinxcrossref{\DUrole{doc}{oxdna/fene}}}} \sphinxhyphen{} modified FENE bond suitable for DNA modeling

\item {} 
\sphinxAtStartPar
{\hyperref[\detokenize{bond_oxdna::doc}]{\sphinxcrossref{\DUrole{doc}{oxdna2/fene}}}} \sphinxhyphen{} same as oxdna but used with different pair styles

\item {} 
\sphinxAtStartPar
{\hyperref[\detokenize{bond_oxdna::doc}]{\sphinxcrossref{\DUrole{doc}{oxrna2/fene}}}} \sphinxhyphen{} modified FENE bond suitable for RNA modeling

\item {} 
\sphinxAtStartPar
{\hyperref[\detokenize{bond_quartic::doc}]{\sphinxcrossref{\DUrole{doc}{quartic}}}} \sphinxhyphen{} breakable quartic bond

\item {} 
\sphinxAtStartPar
{\hyperref[\detokenize{bond_rheo_shell::doc}]{\sphinxcrossref{\DUrole{doc}{rheo/shell}}}} \sphinxhyphen{} shell bond for oxidation modeling in RHEO

\item {} 
\sphinxAtStartPar
{\hyperref[\detokenize{bond_special::doc}]{\sphinxcrossref{\DUrole{doc}{special}}}} \sphinxhyphen{} enable special bond exclusions for 1\sphinxhyphen{}5 pairs and beyond

\item {} 
\sphinxAtStartPar
{\hyperref[\detokenize{bond_table::doc}]{\sphinxcrossref{\DUrole{doc}{table}}}} \sphinxhyphen{} tabulated by bond length

\end{itemize}


\bigskip\hrule\bigskip



\subsection{Restrictions}
\label{\detokenize{bond_style:restrictions}}
\sphinxAtStartPar
Bond styles can only be set for atom styles that allow bonds to be
defined.

\sphinxAtStartPar
Most bond styles are part of the MOLECULE package.  They are only
enabled if LAMMPS was built with that package.  See the {\hyperref[\detokenize{Build_package::doc}]{\sphinxcrossref{\DUrole{doc}{Build package}}}} page for more info.  The doc pages for
individual bond potentials tell if it is part of a package.


\subsection{Related commands}
\label{\detokenize{bond_style:related-commands}}
\sphinxAtStartPar
{\hyperref[\detokenize{bond_coeff::doc}]{\sphinxcrossref{\DUrole{doc}{bond\_coeff}}}}, {\hyperref[\detokenize{delete_bonds::doc}]{\sphinxcrossref{\DUrole{doc}{delete\_bonds}}}}


\subsection{Default}
\label{\detokenize{bond_style:default}}
\begin{sphinxVerbatim}[commandchars=\\\{\}]
\PYG{k}{bond\PYGZus{}style}\PYG{+w}{ }\PYG{n}{none}
\end{sphinxVerbatim}

\sphinxstepscope

\index{bond\_write@\spxentry{bond\_write}}\ignorespaces 

\section{bond\_write command}
\label{\detokenize{bond_write:bond-write-command}}\label{\detokenize{bond_write:index-0}}\label{\detokenize{bond_write::doc}}

\subsection{Syntax}
\label{\detokenize{bond_write:syntax}}
\begin{sphinxVerbatim}[commandchars=\\\{\}]
\PYG{k}{bond\PYGZus{}write}\PYG{+w}{ }\PYG{n}{btype}\PYG{+w}{ }\PYG{n}{N}\PYG{+w}{ }\PYG{n}{inner}\PYG{+w}{ }\PYG{n}{outer}\PYG{+w}{ }\PYG{n}{file}\PYG{+w}{ }\PYG{n}{keyword}\PYG{+w}{ }\PYG{n}{itype}\PYG{+w}{ }\PYG{n}{jtype}
\end{sphinxVerbatim}
\begin{itemize}
\item {} 
\sphinxAtStartPar
btype = bond type

\item {} 
\sphinxAtStartPar
N = \# of values

\item {} 
\sphinxAtStartPar
inner,outer = inner and outer bond length (distance units)

\item {} 
\sphinxAtStartPar
file = name of file to write values to

\item {} 
\sphinxAtStartPar
keyword = section name in file for this set of tabulated values

\item {} 
\sphinxAtStartPar
itype,jtype = two atom types (optional)

\end{itemize}


\subsection{Examples}
\label{\detokenize{bond_write:examples}}
\begin{sphinxVerbatim}[commandchars=\\\{\}]
\PYG{k}{bond\PYGZus{}write}\PYG{+w}{ }\PYG{l+m}{1}\PYG{+w}{ }\PYG{l+m}{500}\PYG{+w}{ }\PYG{l+m}{0.5}\PYG{+w}{ }\PYG{l+m}{3.5}\PYG{+w}{ }\PYG{n}{table.txt}\PYG{+w}{ }\PYG{n}{Harmonic\PYGZus{}1}
\PYG{k}{bond\PYGZus{}write}\PYG{+w}{ }\PYG{l+m}{3}\PYG{+w}{ }\PYG{l+m}{1000}\PYG{+w}{ }\PYG{l+m}{0.1}\PYG{+w}{ }\PYG{l+m}{6.0}\PYG{+w}{ }\PYG{n}{table.txt}\PYG{+w}{ }\PYG{n}{Morse}
\end{sphinxVerbatim}


\subsection{Description}
\label{\detokenize{bond_write:description}}
\sphinxAtStartPar
Write energy and force values to a file as a function of distance for
the currently defined {\hyperref[\detokenize{bond_style::doc}]{\sphinxcrossref{\DUrole{doc}{bond style}}}} for a selected bond
type.  This is useful for plotting the potential function or otherwise
debugging its values.  The resulting file can also be used as input for
use with {\hyperref[\detokenize{bond_table::doc}]{\sphinxcrossref{\DUrole{doc}{bond style table}}}}.

\sphinxAtStartPar
If the file already exists, the table of values is appended to the end
of the file to allow multiple tables of energy and force to be included
in one file.  The individual sections may be identified by the \sphinxstyleemphasis{keyword}.

\sphinxAtStartPar
The energy and force values are computed at distances from \sphinxstyleemphasis{inner} to
\sphinxstyleemphasis{outer} for two interacting atoms forming a bond of type \sphinxstyleemphasis{btype}, using
the appropriate {\hyperref[\detokenize{bond_coeff::doc}]{\sphinxcrossref{\DUrole{doc}{bond\_coeff}}}} coefficients. N evenly
spaced distances are used.

\sphinxAtStartPar
For example, for N = 7, inner = 1.0, and outer = 4.0,
values are computed at r = 1.0, 1.5, 2.0, 2.5, 3.0, 3.5, 4.0.

\sphinxAtStartPar
The file is written in the format used as input for the {\hyperref[\detokenize{bond_table::doc}]{\sphinxcrossref{\DUrole{doc}{bond\_style
table}}}} option with \sphinxstyleemphasis{keyword} as the section name.  Each
line written to the file lists an index number (1\sphinxhyphen{}N), a distance (in
distance units), an energy (in energy units), and a force (in force
units).  In case a new file is created, the first line will be a comment
with a “DATE:” and “UNITS:” tag with the current date and {\hyperref[\detokenize{units::doc}]{\sphinxcrossref{\DUrole{doc}{units}}}} settings.  For subsequent invocations of the \sphinxstyleemphasis{bond\_write}
command for the same file, data will be appended and the current units
settings will be compared to the data from the header, if present.  The
\sphinxstyleemphasis{bond\_write} command will refuse to add a table to an existing file if
the units are not the same.


\subsection{Restrictions}
\label{\detokenize{bond_write:restrictions}}
\sphinxAtStartPar
All force field coefficients for bond and other kinds of interactions
must be set before this command can be invoked.

\sphinxAtStartPar
Due to how the bond force is computed, an inner value \textgreater{} 0.0 must
be specified even if the potential has a finite value at r = 0.0.


\subsection{Related commands}
\label{\detokenize{bond_write:related-commands}}
\sphinxAtStartPar
{\hyperref[\detokenize{bond_table::doc}]{\sphinxcrossref{\DUrole{doc}{bond\_style table}}}}, {\hyperref[\detokenize{angle_write::doc}]{\sphinxcrossref{\DUrole{doc}{angle\_write}}}},
{\hyperref[\detokenize{bond_style::doc}]{\sphinxcrossref{\DUrole{doc}{bond\_style}}}}, {\hyperref[\detokenize{bond_coeff::doc}]{\sphinxcrossref{\DUrole{doc}{bond\_coeff}}}}


\subsection{Default}
\label{\detokenize{bond_write:default}}
\sphinxAtStartPar
none

\sphinxstepscope

\index{boundary@\spxentry{boundary}}\ignorespaces 

\section{boundary command}
\label{\detokenize{boundary:boundary-command}}\label{\detokenize{boundary:index-0}}\label{\detokenize{boundary::doc}}

\subsection{Syntax}
\label{\detokenize{boundary:syntax}}
\begin{sphinxVerbatim}[commandchars=\\\{\}]
\PYG{k}{boundary}\PYG{+w}{ }\PYG{n}{x}\PYG{+w}{ }\PYG{n}{y}\PYG{+w}{ }\PYG{n}{z}
\end{sphinxVerbatim}
\begin{itemize}
\item {} 
\sphinxAtStartPar
x,y,z = \sphinxstyleemphasis{p} or \sphinxstyleemphasis{s} or \sphinxstyleemphasis{f} or \sphinxstyleemphasis{m}, one or two letters
\begin{sphinxalltt}
\sphinxstyleemphasis{p} is periodic
\sphinxstyleemphasis{f} is non\sphinxhyphen{}periodic and fixed
\sphinxstyleemphasis{s} is non\sphinxhyphen{}periodic and shrink\sphinxhyphen{}wrapped
\sphinxstyleemphasis{m} is non\sphinxhyphen{}periodic and shrink\sphinxhyphen{}wrapped with a minimum value
\end{sphinxalltt}

\end{itemize}


\subsection{Examples}
\label{\detokenize{boundary:examples}}
\begin{sphinxVerbatim}[commandchars=\\\{\}]
\PYG{k}{boundary}\PYG{+w}{ }\PYG{n}{p}\PYG{+w}{ }\PYG{n}{p}\PYG{+w}{ }\PYG{n}{f}
\PYG{k}{boundary}\PYG{+w}{ }\PYG{n}{p}\PYG{+w}{ }\PYG{n}{fs}\PYG{+w}{ }\PYG{n}{p}
\PYG{k}{boundary}\PYG{+w}{ }\PYG{n}{s}\PYG{+w}{ }\PYG{n}{f}\PYG{+w}{ }\PYG{n}{fm}
\end{sphinxVerbatim}


\subsection{Description}
\label{\detokenize{boundary:description}}
\sphinxAtStartPar
Set the style of boundaries for the global simulation box in each
dimension.  A single letter assigns the same style to both the lower
and upper face of the box.  Two letters assigns the first style to the
lower face and the second style to the upper face.  The initial size
of the simulation box is set by the {\hyperref[\detokenize{read_data::doc}]{\sphinxcrossref{\DUrole{doc}{read\_data}}}},
{\hyperref[\detokenize{read_restart::doc}]{\sphinxcrossref{\DUrole{doc}{read\_restart}}}}, or {\hyperref[\detokenize{create_box::doc}]{\sphinxcrossref{\DUrole{doc}{create\_box}}}}
commands.

\sphinxAtStartPar
The style \sphinxstyleemphasis{p} means the box is periodic, so that particles interact
across the boundary, and they can exit one end of the box and re\sphinxhyphen{}enter
the other end.  A periodic dimension can change in size due to
constant pressure boundary conditions or box deformation (see the
{\hyperref[\detokenize{fix_nh::doc}]{\sphinxcrossref{\DUrole{doc}{fix npt}}}} and {\hyperref[\detokenize{fix_deform::doc}]{\sphinxcrossref{\DUrole{doc}{fix deform}}}} commands).
The \sphinxstyleemphasis{p} style must be applied to both faces of a dimension.  For 2d
simulations the z dimension must be periodic (which is the default).

\sphinxAtStartPar
The styles \sphinxstyleemphasis{f}, \sphinxstyleemphasis{s}, and \sphinxstyleemphasis{m} mean the box is non\sphinxhyphen{}periodic, so that
particles do not interact across the boundary and do not move from one
side of the box to the other.

\sphinxAtStartPar
For style \sphinxstyleemphasis{f}, the position of the face is fixed.  If an atom moves
outside the face it will be deleted on the next timestep that
reneighboring occurs.  This will typically generate an error unless
you have set the {\hyperref[\detokenize{thermo_modify::doc}]{\sphinxcrossref{\DUrole{doc}{thermo\_modify lost}}}} option to
allow for lost atoms.

\sphinxAtStartPar
For style \sphinxstyleemphasis{s}, the position of the face is set so as to encompass the
atoms in that dimension (shrink\sphinxhyphen{}wrapping), no matter how far they
move. Note that when the difference between the current box dimensions
and the shrink\sphinxhyphen{}wrap box dimensions is large, this can lead to lost
atoms at the beginning of a run when running in parallel. This is due
to the large change in the (global) box dimensions also causing
significant changes in the individual subdomain sizes. If these
changes are farther than the communication cutoff, atoms will be lost.
This is best addressed by setting initial box dimensions to match the
shrink\sphinxhyphen{}wrapped dimensions more closely, by using \sphinxstyleemphasis{m} style boundaries
(see below).

\sphinxAtStartPar
For style \sphinxstyleemphasis{m}, shrink\sphinxhyphen{}wrapping occurs, but is bounded by the value
specified in the data or restart file or set by the
{\hyperref[\detokenize{create_box::doc}]{\sphinxcrossref{\DUrole{doc}{create\_box}}}} command.  For example, if the upper z
face has a value of 50.0 in the data file, the face will always be
positioned at 50.0 or above, even if the maximum z\sphinxhyphen{}extent of all the
atoms becomes less than 50.0.  This can be useful if you start a
simulation with an empty box or if you wish to leave room on one side
of the box, e.g. for atoms to evaporate from a surface.

\sphinxAtStartPar
LAMMPS also allows use of triclinic (non\sphinxhyphen{}orthogonal) simulation boxes.
See the {\hyperref[\detokenize{Howto_triclinic::doc}]{\sphinxcrossref{\DUrole{doc}{Howto triclinic}}}} page for a
description of both general and restricted triclinic boxes and how to
define them.  General triclinic boxes (arbitrary edge vectors \sphinxstylestrong{A},
\sphinxstylestrong{B}, and \sphinxstylestrong{C}) are converted internally to restricted triclinic
boxes with tilt factors (xy,xz,yz) which skew an otherwise orthogonal
box.

\sphinxAtStartPar
The boundary \textless{}boundary\textgreater{} command settings explained above for the 6
faces of an orthogonal box also apply in similar manner to the 6 faces
of a restricted triclinic box (and thus to the corresponding 6 faces
of a general triclinic box), with the following context.

\sphinxAtStartPar
if the second dimension of a tilt factor (e.g. y for xy) is periodic,
then the periodicity is enforced with the tilt factor offset.  This
means that for y periodicity a particle which exits the lower y
boundary is displaced in the x\sphinxhyphen{}direction by xy before it re\sphinxhyphen{}enters the
upper y boundary.  And vice versa if a particle exits the upper y
boundary.  Likewise the ghost atoms surrounding a particle near the
lower y boundary include images of particles near the upper y\sphinxhyphen{}boundary
which are displaced in the x\sphinxhyphen{}direction by xy.  Similar rules apply for
z\sphinxhyphen{}periodicity and the xz and/or yz tilt factors.

\sphinxAtStartPar
If the first dimension of a tilt factor is shrink\sphinxhyphen{}wrapped, then the
shrink wrapping is applied to the tilted box face, to encompass the
atoms.  E.g. for a positive xy tilt, the xlo and xhi faces of the box
are planes tilting in the +y direction as y increases.  The position
of these tilted planes are adjusted dynamically to shrink\sphinxhyphen{}wrap around
the atoms to determine the xlo and xhi extents of the box.


\subsection{Restrictions}
\label{\detokenize{boundary:restrictions}}
\sphinxAtStartPar
This command cannot be used after the simulation box is defined by a
{\hyperref[\detokenize{read_data::doc}]{\sphinxcrossref{\DUrole{doc}{read\_data}}}} or {\hyperref[\detokenize{create_box::doc}]{\sphinxcrossref{\DUrole{doc}{create\_box}}}} command
or {\hyperref[\detokenize{read_restart::doc}]{\sphinxcrossref{\DUrole{doc}{read\_restart}}}} command.  See the
{\hyperref[\detokenize{change_box::doc}]{\sphinxcrossref{\DUrole{doc}{change\_box}}}} command for how to change the
simulation box boundaries after it has been defined.

\sphinxAtStartPar
For 2d simulations, the z dimension must be periodic.


\subsection{Related commands}
\label{\detokenize{boundary:related-commands}}
\sphinxAtStartPar
See the {\hyperref[\detokenize{thermo_modify::doc}]{\sphinxcrossref{\DUrole{doc}{thermo\_modify}}}} command for a discussion
of lost atoms.


\subsection{Default}
\label{\detokenize{boundary:default}}
\begin{sphinxVerbatim}[commandchars=\\\{\}]
\PYG{k}{boundary}\PYG{+w}{ }\PYG{n}{p}\PYG{+w}{ }\PYG{n}{p}\PYG{+w}{ }\PYG{n}{p}
\end{sphinxVerbatim}

\sphinxstepscope

\index{change\_box@\spxentry{change\_box}}\ignorespaces 

\section{change\_box command}
\label{\detokenize{change_box:change-box-command}}\label{\detokenize{change_box:index-0}}\label{\detokenize{change_box::doc}}

\subsection{Syntax}
\label{\detokenize{change_box:syntax}}
\begin{sphinxVerbatim}[commandchars=\\\{\}]
\PYG{k}{change\PYGZus{}box }\PYG{n+nv+nvIdentifier}{group\PYGZhy{}ID}\PYG{+w}{ }\PYG{n}{parameter}\PYG{+w}{ }\PYG{n}{args}\PYG{+w}{ }\PYG{n}{...}\PYG{+w}{ }\PYG{n}{keyword}\PYG{+w}{ }\PYG{n}{args}\PYG{+w}{ }\PYG{n}{...}
\end{sphinxVerbatim}
\begin{itemize}
\item {} 
\sphinxAtStartPar
group\sphinxhyphen{}ID = ID of group of atoms to (optionally) displace

\item {} 
\sphinxAtStartPar
one or more parameter/arg pairs may be appended
\begin{sphinxalltt}
parameter = \sphinxstyleemphasis{x} or \sphinxstyleemphasis{y} or \sphinxstyleemphasis{z} or \sphinxstyleemphasis{xy} or \sphinxstyleemphasis{xz} or \sphinxstyleemphasis{yz} or \sphinxstyleemphasis{boundary} or \sphinxstyleemphasis{ortho} or \sphinxstyleemphasis{triclinic} or \sphinxstyleemphasis{set} or \sphinxstyleemphasis{remap}
  \sphinxstyleemphasis{x}, \sphinxstyleemphasis{y}, \sphinxstyleemphasis{z} args = style value(s)
    style = \sphinxstyleemphasis{final} or \sphinxstyleemphasis{delta} or \sphinxstyleemphasis{scale} or \sphinxstyleemphasis{volume}
      \sphinxstyleemphasis{final} values = lo hi
        lo hi = box boundaries after displacement (distance units)
      \sphinxstyleemphasis{delta} values = dlo dhi
        dlo dhi = change in box boundaries after displacement (distance units)
      \sphinxstyleemphasis{scale} values = factor
        factor = multiplicative factor for change in box length after displacement
      \sphinxstyleemphasis{volume} value = none = adjust this dim to preserve volume of system
  \sphinxstyleemphasis{xy}, \sphinxstyleemphasis{xz}, \sphinxstyleemphasis{yz} args = style value
    style = \sphinxstyleemphasis{final} or \sphinxstyleemphasis{delta}
      \sphinxstyleemphasis{final} value = tilt
        tilt = tilt factor after displacement (distance units)
      \sphinxstyleemphasis{delta} value = dtilt
        dtilt = change in tilt factor after displacement (distance units)
  \sphinxstyleemphasis{boundary} args = x y z
    x,y,z = \sphinxstyleemphasis{p} or \sphinxstyleemphasis{s} or \sphinxstyleemphasis{f} or \sphinxstyleemphasis{m}, one or two letters
    \sphinxstyleemphasis{p} is periodic
    \sphinxstyleemphasis{f} is non\sphinxhyphen{}periodic and fixed
    \sphinxstyleemphasis{s} is non\sphinxhyphen{}periodic and shrink\sphinxhyphen{}wrapped
    \sphinxstyleemphasis{m} is non\sphinxhyphen{}periodic and shrink\sphinxhyphen{}wrapped with a minimum value
  \sphinxstyleemphasis{ortho} args = none = change box to orthogonal
  \sphinxstyleemphasis{triclinic} args = none = change box to triclinic
  \sphinxstyleemphasis{set} args = none = store state of current box
  \sphinxstyleemphasis{remap} args = none = remap atom coords from last saved state to current box
\end{sphinxalltt}

\item {} 
\sphinxAtStartPar
zero or more keyword/value pairs may be appended

\item {} 
\sphinxAtStartPar
keyword = \sphinxstyleemphasis{units}
\begin{sphinxalltt}
\sphinxstyleemphasis{units} value = \sphinxstyleemphasis{lattice} or \sphinxstyleemphasis{box}
  lattice = distances are defined in lattice units
  box = distances are defined in simulation box units
\end{sphinxalltt}

\end{itemize}


\subsection{Examples}
\label{\detokenize{change_box:examples}}
\begin{sphinxVerbatim}[commandchars=\\\{\}]
\PYG{k}{change\PYGZus{}box }\PYG{n+nv+nvIdentifier}{all}\PYG{+w}{ }\PYG{n}{xy}\PYG{+w}{ }\PYG{n}{final}\PYG{+w}{ }\PYG{o}{\PYGZhy{}}\PYG{l+m}{2.0}\PYG{+w}{ }\PYG{n}{z}\PYG{+w}{ }\PYG{n}{final}\PYG{+w}{ }\PYG{l+m}{0.0}\PYG{+w}{ }\PYG{l+m}{5.0}\PYG{+w}{ }\PYG{n}{boundary}\PYG{+w}{ }\PYG{n}{p}\PYG{+w}{ }\PYG{n}{p}\PYG{+w}{ }\PYG{n}{f}\PYG{+w}{ }\PYG{n}{remap}\PYG{+w}{ }\PYG{n}{units}\PYG{+w}{ }\PYG{n}{box}
\PYG{k}{change\PYGZus{}box }\PYG{n+nv+nvIdentifier}{all}\PYG{+w}{ }\PYG{n}{x}\PYG{+w}{ }\PYG{n}{scale}\PYG{+w}{ }\PYG{l+m}{1.1}\PYG{+w}{ }\PYG{n}{y}\PYG{+w}{ }\PYG{n}{volume}\PYG{+w}{ }\PYG{n}{z}\PYG{+w}{ }\PYG{n}{volume}\PYG{+w}{ }\PYG{n}{remap}
\end{sphinxVerbatim}


\subsection{Description}
\label{\detokenize{change_box:description}}
\sphinxAtStartPar
Change the volume and/or shape and/or boundary conditions for the
simulation box.  Orthogonal simulation boxes have 3 adjustable size
parameters (x,y,z).  Triclinic (non\sphinxhyphen{}orthogonal) simulation boxes have
6 adjustable size/shape parameters (x,y,z,xy,xz,yz).  Any or all of
them can be adjusted independently by this command.  Thus it can be
used to expand or contract a box, or to apply a shear strain to a
non\sphinxhyphen{}orthogonal box.  It can also be used to change the boundary
conditions for the simulation box, similar to the
{\hyperref[\detokenize{boundary::doc}]{\sphinxcrossref{\DUrole{doc}{boundary}}}} command.

\sphinxAtStartPar
The size and shape of the initial simulation box are specified by the
{\hyperref[\detokenize{create_box::doc}]{\sphinxcrossref{\DUrole{doc}{create\_box}}}} or {\hyperref[\detokenize{read_data::doc}]{\sphinxcrossref{\DUrole{doc}{read\_data}}}} or
{\hyperref[\detokenize{read_restart::doc}]{\sphinxcrossref{\DUrole{doc}{read\_restart}}}} command used to setup the simulation.
The size and shape may be altered by subsequent runs, e.g. by use of
the {\hyperref[\detokenize{fix_nh::doc}]{\sphinxcrossref{\DUrole{doc}{fix npt}}}} or {\hyperref[\detokenize{fix_deform::doc}]{\sphinxcrossref{\DUrole{doc}{fix deform}}}} commands.
The {\hyperref[\detokenize{create_box::doc}]{\sphinxcrossref{\DUrole{doc}{create\_box}}}}, {\hyperref[\detokenize{read_data::doc}]{\sphinxcrossref{\DUrole{doc}{read data}}}}, and
{\hyperref[\detokenize{read_restart::doc}]{\sphinxcrossref{\DUrole{doc}{read\_restart}}}} commands also determine whether the
simulation box is orthogonal or triclinic and their doc pages explain
the meaning of the xy,xz,yz tilt factors.

\sphinxAtStartPar
See the {\hyperref[\detokenize{Howto_triclinic::doc}]{\sphinxcrossref{\DUrole{doc}{Howto triclinic}}}} page for a
geometric description of triclinic boxes, as defined by LAMMPS, and
how to transform these parameters to and from other commonly used
triclinic representations.

\sphinxAtStartPar
The keywords used in this command are applied sequentially to the
simulation box and the atoms in it, in the order specified.

\sphinxAtStartPar
Before the sequence of keywords are invoked, the current box
size/shape is stored, in case a \sphinxstyleemphasis{remap} keyword is used to map the
atom coordinates from a previously stored box size/shape to the
current one.

\sphinxAtStartPar
After all the keywords have been processed, any shrink\sphinxhyphen{}wrap boundary
conditions are invoked (see the {\hyperref[\detokenize{boundary::doc}]{\sphinxcrossref{\DUrole{doc}{boundary}}}} command)
which may change simulation box boundaries, and atoms are migrated to
new owning processors.

\begin{sphinxadmonition}{note}{Note:}
\sphinxAtStartPar
This means that you cannot use the change\_box command to enlarge
a shrink\sphinxhyphen{}wrapped box, e.g. to make room to insert more atoms via the
{\hyperref[\detokenize{create_atoms::doc}]{\sphinxcrossref{\DUrole{doc}{create\_atoms}}}} command, because the simulation box
will be re\sphinxhyphen{}shrink\sphinxhyphen{}wrapped before the change\_box command completes.
Instead you could do something like this, assuming the simulation box
is non\sphinxhyphen{}periodic and atoms extend from 0 to 20 in all dimensions:
\end{sphinxadmonition}

\begin{sphinxVerbatim}[commandchars=\\\{\}]
\PYG{k}{change\PYGZus{}box }\PYG{n+nv+nvIdentifier}{all}\PYG{+w}{ }\PYG{n}{x}\PYG{+w}{ }\PYG{n}{final}\PYG{+w}{ }\PYG{o}{\PYGZhy{}}\PYG{l+m}{10}\PYG{+w}{ }\PYG{l+m}{20}
\PYG{k}{create\PYGZus{}atoms}\PYG{+w}{ }\PYG{l+m}{1}\PYG{+w}{ }\PYG{n}{single}\PYG{+w}{ }\PYG{o}{\PYGZhy{}}\PYG{l+m}{5}\PYG{+w}{ }\PYG{l+m}{5}\PYG{+w}{ }\PYG{l+m}{5}\PYG{+w}{       }\PYG{c}{\PYGZsh{} this will fail to insert an atom}

\PYG{k}{change\PYGZus{}box }\PYG{n+nv+nvIdentifier}{all}\PYG{+w}{ }\PYG{n}{x}\PYG{+w}{ }\PYG{n}{final}\PYG{+w}{ }\PYG{o}{\PYGZhy{}}\PYG{l+m}{10}\PYG{+w}{ }\PYG{l+m}{20}\PYG{+w}{ }\PYG{n}{boundary}\PYG{+w}{ }\PYG{n}{f}\PYG{+w}{ }\PYG{n}{s}\PYG{+w}{ }\PYG{n}{s}
\PYG{k}{create\PYGZus{}atoms}\PYG{+w}{ }\PYG{l+m}{1}\PYG{+w}{ }\PYG{n}{single}\PYG{+w}{ }\PYG{o}{\PYGZhy{}}\PYG{l+m}{5}\PYG{+w}{ }\PYG{l+m}{5}\PYG{+w}{ }\PYG{l+m}{5}
\PYG{k}{change\PYGZus{}box }\PYG{n+nv+nvIdentifier}{all}\PYG{+w}{ }\PYG{n}{boundary}\PYG{+w}{ }\PYG{n}{s}\PYG{+w}{ }\PYG{n}{s}\PYG{+w}{ }\PYG{n}{s}\PYG{+w}{      }\PYG{c}{\PYGZsh{} this will work}
\end{sphinxVerbatim}

\begin{sphinxadmonition}{note}{Note:}
\sphinxAtStartPar
Unlike the earlier “displace\_box” version of this command, atom
remapping is NOT performed by default.  This command allows remapping
to be done in a more general way, exactly when you specify it (zero or
more times) in the sequence of transformations.  Thus if you do not
use the \sphinxstyleemphasis{remap} keyword, atom coordinates will not be changed even if
the box size/shape changes.  If a uniformly strained state is desired,
the \sphinxstyleemphasis{remap} keyword should be specified.
\end{sphinxadmonition}

\begin{sphinxadmonition}{note}{Note:}
\sphinxAtStartPar
It is possible to lose atoms with this command.  E.g. by
changing the box without remapping the atoms, and having atoms end up
outside of non\sphinxhyphen{}periodic boundaries.  It is also possible to alter
bonds between atoms straddling a boundary in bad ways.  E.g. by
converting a boundary from periodic to non\sphinxhyphen{}periodic.  It is also
possible when remapping atoms to put them (nearly) on top of each
other.  E.g. by converting a boundary from non\sphinxhyphen{}periodic to periodic.
All of these will typically lead to bad dynamics and/or generate error
messages.
\end{sphinxadmonition}

\begin{sphinxadmonition}{note}{Note:}
\sphinxAtStartPar
The simulation box size/shape can be changed by arbitrarily large
amounts by this command.  This is not a problem, except that the
mapping of processors to the simulation box is not changed from its
initial 3d configuration; see the {\hyperref[\detokenize{processors::doc}]{\sphinxcrossref{\DUrole{doc}{processors}}}}
command.  Thus, if the box size/shape changes dramatically, the
mapping of processors to the simulation box may not end up as
optimal as the initial mapping attempted to be.  You may wish to
re\sphinxhyphen{}balance the atoms by using the {\hyperref[\detokenize{balance::doc}]{\sphinxcrossref{\DUrole{doc}{balance}}}} command
if that is the case.
\end{sphinxadmonition}

\begin{sphinxadmonition}{note}{Note:}
\sphinxAtStartPar
You cannot use this command after reading a restart file (and
before a run is performed) if the restart file stored per\sphinxhyphen{}atom
information from a fix and any of the specified keywords change the
box size or shape or boundary conditions.  This is because atoms
may be moved to new processors and the restart info will not
migrate with them.  LAMMPS will generate an error if this could
happen.  Only the \sphinxstyleemphasis{ortho} and \sphinxstyleemphasis{triclinic} keywords do not trigger
this error.  One solution is to perform a “run 0” command before
using the change\_box command.  This clears the per\sphinxhyphen{}atom restart
data, whether it has been re\sphinxhyphen{}assigned to a new fix or not.
\end{sphinxadmonition}

\begin{sphinxadmonition}{note}{Note:}
\sphinxAtStartPar
Because the keywords used in this command are applied one at a time
to the simulation box and the atoms in it, care must be taken with
triclinic cells to avoid exceeding the limits on skew after each
transformation in the sequence.  If skew is exceeded before the
final transformation this can be avoided by changing the order of
the sequence, or breaking the transformation into two or more
smaller transformations.  For more information on the allowed
limits for box skew see the discussion on triclinic boxes on
{\hyperref[\detokenize{Howto_triclinic::doc}]{\sphinxcrossref{\DUrole{doc}{Howto triclinic}}}} doc page.
\end{sphinxadmonition}


\bigskip\hrule\bigskip


\sphinxAtStartPar
For the \sphinxstyleemphasis{x}, \sphinxstyleemphasis{y}, and \sphinxstyleemphasis{z} parameters, this is the meaning of their
styles and values.

\sphinxAtStartPar
For style \sphinxstyleemphasis{final}, the final lo and hi box boundaries of a dimension
are specified.  The values can be in lattice or box distance units.
See the discussion of the units keyword below.

\sphinxAtStartPar
For style \sphinxstyleemphasis{delta}, plus or minus changes in the lo/hi box boundaries
of a dimension are specified.  The values can be in lattice or box
distance units.  See the discussion of the units keyword below.

\sphinxAtStartPar
For style \sphinxstyleemphasis{scale}, a multiplicative factor to apply to the box length
of a dimension is specified.  For example, if the initial box length
is 10, and the factor is 1.1, then the final box length will be 11.  A
factor less than 1.0 means compression.

\sphinxAtStartPar
The \sphinxstyleemphasis{volume} style changes the specified dimension in such a way that
the overall box volume remains constant with respect to the operation
performed by the preceding keyword.  The \sphinxstyleemphasis{volume} style can only be
used following a keyword that changed the volume, which is any of the
\sphinxstyleemphasis{x}, \sphinxstyleemphasis{y}, \sphinxstyleemphasis{z} keywords.  If the preceding keyword “key” had a \sphinxstyleemphasis{volume}
style, then both it and the current keyword apply to the keyword
preceding “key”.  I.e. this sequence of keywords is allowed:

\begin{sphinxVerbatim}[commandchars=\\\{\}]
\PYG{k}{change\PYGZus{}box }\PYG{n+nv+nvIdentifier}{all}\PYG{+w}{ }\PYG{n}{x}\PYG{+w}{ }\PYG{n}{scale}\PYG{+w}{ }\PYG{l+m}{1.1}\PYG{+w}{ }\PYG{n}{y}\PYG{+w}{ }\PYG{n}{volume}\PYG{+w}{ }\PYG{n}{z}\PYG{+w}{ }\PYG{n}{volume}
\end{sphinxVerbatim}

\sphinxAtStartPar
The \sphinxstyleemphasis{volume} style changes the associated dimension so that the
overall box volume is unchanged relative to its value before the
preceding keyword was invoked.

\sphinxAtStartPar
If the following command is used, then the z box length will shrink by
the same 1.1 factor the x box length was increased by:

\begin{sphinxVerbatim}[commandchars=\\\{\}]
\PYG{k}{change\PYGZus{}box }\PYG{n+nv+nvIdentifier}{all}\PYG{+w}{ }\PYG{n}{x}\PYG{+w}{ }\PYG{n}{scale}\PYG{+w}{ }\PYG{l+m}{1.1}\PYG{+w}{ }\PYG{n}{z}\PYG{+w}{ }\PYG{n}{volume}
\end{sphinxVerbatim}

\sphinxAtStartPar
If the following command is used, then the y,z box lengths will each
shrink by sqrt(1.1) to keep the volume constant.  In this case, the
y,z box lengths shrink so as to keep their relative aspect ratio
constant:

\begin{sphinxVerbatim}[commandchars=\\\{\}]
\PYG{k}{change\PYGZus{}box }\PYG{n+nv+nvIdentifier}{all}\PYG{+w}{ }\PYG{n}{x}\PYG{+w}{ }\PYG{n}{scale}\PYG{+w}{ }\PYG{l+m}{1.1}\PYG{+w}{ }\PYG{n}{y}\PYG{+w}{ }\PYG{n}{volume}\PYG{+w}{ }\PYG{n}{z}\PYG{+w}{ }\PYG{n}{volume}
\end{sphinxVerbatim}

\sphinxAtStartPar
If the following command is used, then the final box will be a factor
of 10\% larger in x and y, and a factor of 21\% smaller in z, so as to
keep the volume constant:

\begin{sphinxVerbatim}[commandchars=\\\{\}]
\PYG{k}{change\PYGZus{}box }\PYG{n+nv+nvIdentifier}{all}\PYG{+w}{ }\PYG{n}{x}\PYG{+w}{ }\PYG{n}{scale}\PYG{+w}{ }\PYG{l+m}{1.1}\PYG{+w}{ }\PYG{n}{z}\PYG{+w}{ }\PYG{n}{volume}\PYG{+w}{ }\PYG{n}{y}\PYG{+w}{ }\PYG{n}{scale}\PYG{+w}{ }\PYG{l+m}{1.1}\PYG{+w}{ }\PYG{n}{z}\PYG{+w}{ }\PYG{n}{volume}
\end{sphinxVerbatim}

\begin{sphinxadmonition}{note}{Note:}
\sphinxAtStartPar
For solids or liquids, when one dimension of the box is
expanded, it may be physically undesirable to hold the other 2 box
lengths constant since that implies a density change.  For solids,
adjusting the other dimensions via the \sphinxstyleemphasis{volume} style may make
physical sense (just as for a liquid), but may not be correct for
materials and potentials whose Poisson ratio is not 0.5.
\end{sphinxadmonition}

\sphinxAtStartPar
For the \sphinxstyleemphasis{scale} and \sphinxstyleemphasis{volume} styles, the box length is expanded or
compressed around its mid point.


\bigskip\hrule\bigskip


\sphinxAtStartPar
For the \sphinxstyleemphasis{xy}, \sphinxstyleemphasis{xz}, and \sphinxstyleemphasis{yz} parameters, this is the meaning of their
styles and values.  Note that changing the tilt factors of a triclinic
box does not change its volume.

\sphinxAtStartPar
For style \sphinxstyleemphasis{final}, the final tilt factor is specified.  The value
can be in lattice or box distance units.  See the discussion of the
units keyword below.

\sphinxAtStartPar
For style \sphinxstyleemphasis{delta}, a plus or minus change in the tilt factor is
specified.  The value can be in lattice or box distance units.  See
the discussion of the units keyword below.

\sphinxAtStartPar
All of these styles change the xy, xz, yz tilt factors.  In LAMMPS,
tilt factors (xy,xz,yz) for triclinic boxes are required to be no more
than half the distance of the parallel box length.  For example, if
xlo = 2 and xhi = 12, then the x box length is 10 and the xy tilt
factor must be between \sphinxhyphen{}5 and 5.  Similarly, both xz and yz must be
between \sphinxhyphen{}(xhi\sphinxhyphen{}xlo)/2 and +(yhi\sphinxhyphen{}ylo)/2.  Note that this is not a
limitation, since if the maximum tilt factor is 5 (as in this
example), then configurations with tilt = …, \sphinxhyphen{}15, \sphinxhyphen{}5, 5, 15, 25,
… are all equivalent.  Any tilt factor specified by this command
must be within these limits.


\bigskip\hrule\bigskip


\sphinxAtStartPar
The \sphinxstyleemphasis{boundary} keyword takes arguments that have exactly the same
meaning as they do for the {\hyperref[\detokenize{boundary::doc}]{\sphinxcrossref{\DUrole{doc}{boundary}}}} command.  In each
dimension, a single letter assigns the same style to both the lower
and upper face of the box.  Two letters assigns the first style to the
lower face and the second style to the upper face.

\sphinxAtStartPar
The style \sphinxstyleemphasis{p} means the box is periodic; the other styles mean
non\sphinxhyphen{}periodic. For style \sphinxstyleemphasis{f}, the position of the face is fixed.  For
style \sphinxstyleemphasis{s}, the position of the face is set so as to encompass the
atoms in that dimension (shrink\sphinxhyphen{}wrapping), no matter how far they
move.  For style \sphinxstyleemphasis{m}, shrink\sphinxhyphen{}wrapping occurs, but is bounded by the
current box edge in that dimension, so that the box will become no
smaller.  See the {\hyperref[\detokenize{boundary::doc}]{\sphinxcrossref{\DUrole{doc}{boundary}}}} command for more
explanation of these style options.

\sphinxAtStartPar
Note that the “boundary” command itself can only be used before the
simulation box is defined via a {\hyperref[\detokenize{read_data::doc}]{\sphinxcrossref{\DUrole{doc}{read\_data}}}} or
{\hyperref[\detokenize{create_box::doc}]{\sphinxcrossref{\DUrole{doc}{create\_box}}}} or {\hyperref[\detokenize{read_restart::doc}]{\sphinxcrossref{\DUrole{doc}{read\_restart}}}}
command.  This command allows the boundary conditions to be changed
later in your input script.  Also note that the
{\hyperref[\detokenize{read_restart::doc}]{\sphinxcrossref{\DUrole{doc}{read\_restart}}}} will change boundary conditions to
match what is stored in the restart file.  So if you wish to change
them, you should use the change\_box command after the read\_restart
command.

\begin{sphinxadmonition}{note}{Note:}
\sphinxAtStartPar
Changing a periodic boundary to a non\sphinxhyphen{}periodic one will also cause
the image flag for that dimension of all atoms to be reset to 0.
LAMMPS will print a warning message, if that happens.
Please note that this reset can lead to undesired changes when
atoms are involved in bonded interactions that straddle periodic
boundaries and thus the values of the image flag differs for those
atoms.
\end{sphinxadmonition}


\bigskip\hrule\bigskip


\sphinxAtStartPar
The \sphinxstyleemphasis{ortho} and \sphinxstyleemphasis{triclinic} keywords convert the simulation box to be
orthogonal or triclinic (non\sphinxhyphen{}orthogonal).

\sphinxAtStartPar
The simulation box is defined as either orthogonal or triclinic when
it is created via the {\hyperref[\detokenize{create_box::doc}]{\sphinxcrossref{\DUrole{doc}{create\_box}}}},
{\hyperref[\detokenize{read_data::doc}]{\sphinxcrossref{\DUrole{doc}{read\_data}}}}, or {\hyperref[\detokenize{read_restart::doc}]{\sphinxcrossref{\DUrole{doc}{read\_restart}}}}
commands.

\sphinxAtStartPar
These keywords allow you to toggle the existing simulation box from
orthogonal to triclinic and vice versa.  For example, an initial
equilibration simulation can be run in an orthogonal box, the box can
be toggled to triclinic, and then a {\hyperref[\detokenize{Howto_nemd::doc}]{\sphinxcrossref{\DUrole{doc}{non\sphinxhyphen{}equilibrium MD (NEMD) simulation}}}} can be run with deformation via the {\hyperref[\detokenize{fix_deform::doc}]{\sphinxcrossref{\DUrole{doc}{fix deform}}}} command.

\sphinxAtStartPar
If the simulation box is currently triclinic and has non\sphinxhyphen{}zero tilt in
xy, yz, or xz, then it cannot be converted to an orthogonal box.


\bigskip\hrule\bigskip


\sphinxAtStartPar
The \sphinxstyleemphasis{set} keyword saves the current box size/shape.  This can be
useful if you wish to use the \sphinxstyleemphasis{remap} keyword more than once or if you
wish it to be applied to an intermediate box size/shape in a sequence
of keyword operations.  Note that the box size/shape is saved before
any of the keywords are processed, i.e. the box size/shape at the time
the create\_box command is encountered in the input script.

\sphinxAtStartPar
The \sphinxstyleemphasis{remap} keyword remaps atom coordinates from the last saved box
size/shape to the current box state.  For example, if you stretch the
box in the x dimension or tilt it in the xy plane via the \sphinxstyleemphasis{x} and \sphinxstyleemphasis{xy}
keywords, then the \sphinxstyleemphasis{remap} command will dilate or tilt the atoms to
conform to the new box size/shape, as if the atoms moved with the box
as it deformed.

\sphinxAtStartPar
Note that this operation is performed without regard to periodic
boundaries.  Also, any shrink\sphinxhyphen{}wrapping of non\sphinxhyphen{}periodic boundaries (see
the {\hyperref[\detokenize{boundary::doc}]{\sphinxcrossref{\DUrole{doc}{boundary}}}} command) occurs after all keywords,
including this one, have been processed.

\sphinxAtStartPar
Only atoms in the specified group are remapped.


\bigskip\hrule\bigskip


\sphinxAtStartPar
The \sphinxstyleemphasis{units} keyword determines the meaning of the distance units used
to define various arguments.  A \sphinxstyleemphasis{box} value selects standard distance
units as defined by the {\hyperref[\detokenize{units::doc}]{\sphinxcrossref{\DUrole{doc}{units}}}} command, e.g. Angstroms for
units = real or metal.  A \sphinxstyleemphasis{lattice} value means the distance units are
in lattice spacings.  The {\hyperref[\detokenize{lattice::doc}]{\sphinxcrossref{\DUrole{doc}{lattice}}}} command must have
been previously used to define the lattice spacing.


\bigskip\hrule\bigskip



\subsection{Restrictions}
\label{\detokenize{change_box:restrictions}}
\sphinxAtStartPar
If you use the \sphinxstyleemphasis{ortho} or \sphinxstyleemphasis{triclinic} keywords, then at the point in
the input script when this command is issued, no {\hyperref[\detokenize{dump::doc}]{\sphinxcrossref{\DUrole{doc}{dumps}}}} can
be active, nor can a {\hyperref[\detokenize{fix_deform::doc}]{\sphinxcrossref{\DUrole{doc}{fix deform}}}} be active.  This is
because these commands test whether the simulation box is orthogonal
when they are first issued.  Note that these commands can be used in
your script before a change\_box command is issued, so long as an
{\hyperref[\detokenize{undump::doc}]{\sphinxcrossref{\DUrole{doc}{undump}}}} or {\hyperref[\detokenize{unfix::doc}]{\sphinxcrossref{\DUrole{doc}{unfix}}}} command is also used to
turn them off.


\subsection{Related commands}
\label{\detokenize{change_box:related-commands}}
\sphinxAtStartPar
{\hyperref[\detokenize{fix_deform::doc}]{\sphinxcrossref{\DUrole{doc}{fix deform}}}}, {\hyperref[\detokenize{boundary::doc}]{\sphinxcrossref{\DUrole{doc}{boundary}}}}


\subsection{Default}
\label{\detokenize{change_box:default}}
\sphinxAtStartPar
The option default is units = lattice.

\sphinxstepscope

\index{clear@\spxentry{clear}}\ignorespaces 

\section{clear command}
\label{\detokenize{clear:clear-command}}\label{\detokenize{clear:index-0}}\label{\detokenize{clear::doc}}

\subsection{Syntax}
\label{\detokenize{clear:syntax}}
\begin{sphinxVerbatim}[commandchars=\\\{\}]
\PYG{k}{clear}
\end{sphinxVerbatim}


\subsection{Examples}
\label{\detokenize{clear:examples}}
\begin{sphinxVerbatim}[commandchars=\\\{\}]
\PYG{c}{\PYGZsh{} (commands for 1st simulation)}
\PYG{k}{clear}
\PYG{c}{\PYGZsh{} (commands for 2nd simulation)}
\end{sphinxVerbatim}


\subsection{Description}
\label{\detokenize{clear:description}}
\sphinxAtStartPar
This command deletes all atoms, restores all settings to their default
values, and frees all memory allocated by LAMMPS.  Once a clear command
has been executed, it is almost as if LAMMPS were starting over, with
only the exceptions noted below.  This command enables multiple jobs to
be run sequentially from one input script.

\sphinxAtStartPar
These settings are not affected by a clear command: the working
directory ({\hyperref[\detokenize{shell::doc}]{\sphinxcrossref{\DUrole{doc}{shell}}}} command), log file status ({\hyperref[\detokenize{log::doc}]{\sphinxcrossref{\DUrole{doc}{log}}}} command), echo status ({\hyperref[\detokenize{echo::doc}]{\sphinxcrossref{\DUrole{doc}{echo}}}} command), and input
script variables except for \sphinxstyleemphasis{atomfile} style variables ({\hyperref[\detokenize{variable::doc}]{\sphinxcrossref{\DUrole{doc}{variable}}}} command).


\subsection{Restrictions}
\label{\detokenize{clear:restrictions}}
\sphinxAtStartPar
none


\subsection{Related commands}
\label{\detokenize{clear:related-commands}}
\sphinxAtStartPar
none


\subsection{Default}
\label{\detokenize{clear:default}}
\sphinxAtStartPar
none

\sphinxstepscope

\index{comm\_modify@\spxentry{comm\_modify}}\ignorespaces 

\section{comm\_modify command}
\label{\detokenize{comm_modify:comm-modify-command}}\label{\detokenize{comm_modify:index-0}}\label{\detokenize{comm_modify::doc}}

\subsection{Syntax}
\label{\detokenize{comm_modify:syntax}}
\begin{sphinxVerbatim}[commandchars=\\\{\}]
\PYG{k}{comm\PYGZus{}modify}\PYG{+w}{ }\PYG{n}{keyword}\PYG{+w}{ }\PYG{n}{value}\PYG{+w}{ }\PYG{n}{...}
\end{sphinxVerbatim}
\begin{itemize}
\item {} 
\sphinxAtStartPar
one or more keyword/value pairs may be appended

\item {} 
\sphinxAtStartPar
keyword = \sphinxstyleemphasis{mode} or \sphinxstyleemphasis{cutoff} or \sphinxstyleemphasis{cutoff/multi} or \sphinxstyleemphasis{group} or \sphinxstyleemphasis{reduce/multi} or \sphinxstyleemphasis{vel}
\begin{sphinxalltt}
\sphinxstyleemphasis{mode} value = \sphinxstyleemphasis{single}, \sphinxstyleemphasis{multi}, or \sphinxstyleemphasis{multi/old} = communicate atoms within a single or multiple distances
\sphinxstyleemphasis{cutoff} value = Rcut (distance units) = communicate atoms from this far away
\sphinxstyleemphasis{cutoff/multi} collection value
   collection = atom collection or collection range (supports asterisk notation)
   value = Rcut (distance units) = communicate atoms for selected types from this far away
\sphinxstyleemphasis{reduce/multi} arg = none = reduce number of communicated ghost atoms for multi style
\sphinxstyleemphasis{cutoff/multi/old} type value
   type = atom type or type range (supports asterisk notation)
   value = Rcut (distance units) = communicate atoms for selected types from this far away
\sphinxstyleemphasis{group} value = group\sphinxhyphen{}ID = only communicate atoms in the group
\sphinxstyleemphasis{vel} value = \sphinxstyleemphasis{yes} or \sphinxstyleemphasis{no} = do or do not communicate velocity info with ghost atoms
\end{sphinxalltt}

\end{itemize}


\subsection{Examples}
\label{\detokenize{comm_modify:examples}}
\begin{sphinxVerbatim}[commandchars=\\\{\}]
\PYG{k}{comm\PYGZus{}modify}\PYG{+w}{ }\PYG{n}{mode}\PYG{+w}{ }\PYG{n}{multi}\PYG{+w}{ }\PYG{n}{reduce}\PYG{o}{/}\PYG{n}{multi}
\PYG{k}{comm\PYGZus{}modify}\PYG{+w}{ }\PYG{n}{mode}\PYG{+w}{ }\PYG{n}{multi}\PYG{+w}{ }\PYG{k}{group }\PYG{n+nv+nvIdentifier}{solvent}
\PYG{k}{comm\PYGZus{}modify}\PYG{+w}{ }\PYG{n}{mode}\PYG{+w}{ }\PYG{n}{multi}\PYG{+w}{ }\PYG{n}{cutoff}\PYG{o}{/}\PYG{n}{multi}\PYG{+w}{ }\PYG{l+m}{1}\PYG{+w}{ }\PYG{l+m}{10.0}\PYG{+w}{ }\PYG{n}{cutoff}\PYG{o}{/}\PYG{n}{multi}\PYG{+w}{ }\PYG{l+m}{2}\PYG{o}{*}\PYG{l+m}{4}\PYG{+w}{ }\PYG{l+m}{15.0}
\PYG{k}{comm\PYGZus{}modify}\PYG{+w}{ }\PYG{n}{vel}\PYG{+w}{ }\PYG{n}{yes}
\PYG{k}{comm\PYGZus{}modify}\PYG{+w}{ }\PYG{n}{mode}\PYG{+w}{ }\PYG{n}{single}\PYG{+w}{ }\PYG{n}{cutoff}\PYG{+w}{ }\PYG{l+m}{5.0}\PYG{+w}{ }\PYG{n}{vel}\PYG{+w}{ }\PYG{n}{yes}
\PYG{k}{comm\PYGZus{}modify}\PYG{+w}{ }\PYG{n}{cutoff}\PYG{o}{/}\PYG{n}{multi}\PYG{+w}{ }\PYG{o}{*}\PYG{+w}{ }\PYG{l+m}{0.0}
\end{sphinxVerbatim}


\subsection{Description}
\label{\detokenize{comm_modify:description}}
\sphinxAtStartPar
This command sets parameters that affect the inter\sphinxhyphen{}processor
communication of atom information that occurs each timestep as
coordinates and other properties are exchanged between neighboring
processors and stored as properties of ghost atoms.

\begin{sphinxadmonition}{note}{Note:}
\sphinxAtStartPar
These options apply to the currently defined comm style.  When
you specify a {\hyperref[\detokenize{comm_style::doc}]{\sphinxcrossref{\DUrole{doc}{comm\_style}}}} or
{\hyperref[\detokenize{read_restart::doc}]{\sphinxcrossref{\DUrole{doc}{read\_restart}}}} command, all communication settings
are restored to their default or stored values, including those
previously reset by a comm\_modify command.  Thus if your input script
specifies a comm\_style or read\_restart command, you should use the
comm\_modify command after it.
\end{sphinxadmonition}

\sphinxAtStartPar
The \sphinxstyleemphasis{mode} keyword determines whether a single or multiple cutoff
distances are used to determine which atoms to communicate.

\sphinxAtStartPar
The default mode is \sphinxstyleemphasis{single} which means each processor acquires
information for ghost atoms that are within a single distance from its
subdomain.  The distance is by default the maximum of the neighbor
cutoff across all atom type pairs.

\sphinxAtStartPar
For many systems this is an efficient algorithm, but for systems with
widely varying cutoffs for different type pairs, the \sphinxstyleemphasis{multi} or \sphinxstyleemphasis{multi/old} mode can
be faster.  In \sphinxstyleemphasis{multi}, each atom is assigned to a collection which should
correspond to a set of atoms with similar interaction cutoffs.
See the {\hyperref[\detokenize{neighbor::doc}]{\sphinxcrossref{\DUrole{doc}{neighbor}}}} command for a detailed description of collections.
In this case, each atom collection is assigned its own distance
cutoff for communication purposes, and fewer atoms will be
communicated. in \sphinxstyleemphasis{multi/old}, a similar technique is used but atoms
are grouped by atom type. See the {\hyperref[\detokenize{neighbor::doc}]{\sphinxcrossref{\DUrole{doc}{neighbor multi}}}}  and
{\hyperref[\detokenize{neighbor::doc}]{\sphinxcrossref{\DUrole{doc}{neighbor multi/old}}}} commands for
neighbor list construction options that may also be beneficial for
simulations of this kind. The \sphinxstyleemphasis{multi} communication mode is only compatible
with the \sphinxstyleemphasis{multi} neighbor style. The \sphinxstyleemphasis{multi/old} communication mode is comparable
with both the \sphinxstyleemphasis{multi} and \sphinxstyleemphasis{multi/old} neighbor styles.

\sphinxAtStartPar
The \sphinxstyleemphasis{cutoff} keyword allows you to extend the ghost cutoff distance
for communication mode \sphinxstyleemphasis{single}, which is the distance from the borders
of a processor’s subdomain at which ghost atoms are acquired from other
processors.  By default the ghost cutoff = neighbor cutoff = pairwise
force cutoff + neighbor skin.  See the {\hyperref[\detokenize{neighbor::doc}]{\sphinxcrossref{\DUrole{doc}{neighbor}}}} command
for more information about the skin distance.  If the specified Rcut is
greater than the neighbor cutoff, then extra ghost atoms will be acquired.
If the provided cutoff is smaller, the provided value will be ignored,
the ghost cutoff is set to the neighbor cutoff and a warning will be
printed. Specifying a cutoff value of 0.0 will reset any previous value
to the default. If bonded interactions exist and equilibrium bond length
information is available, then also a heuristic based on that bond length
is computed. It is used as communication cutoff, if there is no pair
style present and no \sphinxstyleemphasis{comm\_modify cutoff} command used. Otherwise a
warning is printed, if this bond based estimate is larger than the
communication cutoff used.

\sphinxAtStartPar
The \sphinxstyleemphasis{cutoff/multi} option is equivalent to \sphinxstyleemphasis{cutoff}, but applies to
communication mode \sphinxstyleemphasis{multi} instead. Since the communication cutoffs are
determined per atom collections, a collection specifier is needed and
cutoff for one or multiple collections can be extended. Also ranges of
collections using the usual asterisk notation can be given.  Collections
are indexed from 1 to N where N is the total number of collections.
Note that the arguments for \sphinxstyleemphasis{cutoff/multi} are parsed right before each
simulation to account for potential changes in the number of
collections.  Custom cutoffs are preserved between runs but if
collections are redefined, one may want to re\sphinxhyphen{}specify the communication
cutoffs.  For granular pair styles,the default cutoff is set to the sum
of the current maximum atomic radii for each collection.  The
\sphinxstyleemphasis{cutoff/multi/old} option is similar to \sphinxstyleemphasis{cutoff/multi} except it
operates on atom types as opposed to collections.

\sphinxAtStartPar
The \sphinxstyleemphasis{reduce/multi} option applies to \sphinxstyleemphasis{multi} and sets the communication
cutoff for a particle equal to the maximum interaction distance between particles
in the same collection. This reduces the number of
ghost atoms that need to be communicated. This method is only compatible with the
\sphinxstyleemphasis{multi} neighbor style and requires a half neighbor list and Newton on.
See the {\hyperref[\detokenize{neighbor::doc}]{\sphinxcrossref{\DUrole{doc}{neighbor multi}}}} command for more information.

\sphinxAtStartPar
These are simulation scenarios in which it may be useful or even
necessary to set a ghost cutoff \textgreater{} neighbor cutoff:
\begin{itemize}
\item {} 
\sphinxAtStartPar
a single polymer chain with bond interactions, but no pairwise interactions

\item {} 
\sphinxAtStartPar
bonded interactions (e.g. dihedrals) extend further than the pairwise cutoff

\item {} 
\sphinxAtStartPar
ghost atoms beyond the pairwise cutoff are needed for some computation

\end{itemize}

\sphinxAtStartPar
In the first scenario, a pairwise potential is not defined.  Thus the
pairwise neighbor cutoff will be 0.0.  But ghost atoms are still
needed for computing bond, angle, etc interactions between atoms on
different processors, or when the interaction straddles a periodic
boundary.

\sphinxAtStartPar
The appropriate ghost cutoff depends on the {\hyperref[\detokenize{newton::doc}]{\sphinxcrossref{\DUrole{doc}{newton bond}}}}
setting.  For newton bond \sphinxstyleemphasis{off}, the distance needs to be the furthest
distance between any two atoms in the bond, angle, etc.  E.g. the
distance between 1\sphinxhyphen{}4 atoms in a dihedral.  For newton bond \sphinxstyleemphasis{on}, the
distance between the central atom in the bond, angle, etc and any
other atom is sufficient.  E.g. the distance between 2\sphinxhyphen{}4 atoms in a
dihedral.

\sphinxAtStartPar
In the second scenario, a pairwise potential is defined, but its
neighbor cutoff is not sufficiently long enough to enable bond, angle,
etc terms to be computed.  As in the previous scenario, an appropriate
ghost cutoff should be set.

\sphinxAtStartPar
In the last scenario, a {\hyperref[\detokenize{fix::doc}]{\sphinxcrossref{\DUrole{doc}{fix}}}} or {\hyperref[\detokenize{compute::doc}]{\sphinxcrossref{\DUrole{doc}{compute}}}} or
{\hyperref[\detokenize{pair_style::doc}]{\sphinxcrossref{\DUrole{doc}{pairwise potential}}}} needs to calculate with ghost
atoms beyond the normal pairwise cutoff for some computation it
performs (e.g. locate neighbors of ghost atoms in a manybody pair
potential).  Setting the ghost cutoff appropriately can ensure it will
find the needed atoms.

\begin{sphinxadmonition}{note}{Note:}
\sphinxAtStartPar
In these scenarios, if you do not set the ghost cutoff long
enough, and if there is only one processor in a periodic dimension
(e.g. you are running in serial), then LAMMPS may “find” the atom it
is looking for (e.g. the partner atom in a bond), that is on the far
side of the simulation box, across a periodic boundary.  This will
typically lead to bad dynamics (i.e. the bond length is now the
simulation box length).  To detect if this is happening, see the
{\hyperref[\detokenize{neigh_modify::doc}]{\sphinxcrossref{\DUrole{doc}{neigh\_modify cluster}}}} command.
\end{sphinxadmonition}

\sphinxAtStartPar
The \sphinxstyleemphasis{group} keyword will limit communication to atoms in the specified
group.  This can be useful for models where no ghost atoms are needed
for some kinds of particles.  All atoms (not just those in the
specified group) will still migrate to new processors as they move.
The group specified with this option must also be specified via the
{\hyperref[\detokenize{atom_modify::doc}]{\sphinxcrossref{\DUrole{doc}{atom\_modify first}}}} command.

\sphinxAtStartPar
The \sphinxstyleemphasis{vel} keyword enables velocity information to be communicated with
ghost particles.  Depending on the {\hyperref[\detokenize{atom_style::doc}]{\sphinxcrossref{\DUrole{doc}{atom\_style}}}},
velocity info includes the translational velocity, angular velocity,
and angular momentum of a particle.  If the \sphinxstyleemphasis{vel} option is set to
\sphinxstyleemphasis{yes}, then ghost atoms store these quantities; if \sphinxstyleemphasis{no} then they do
not.  The \sphinxstyleemphasis{yes} setting is needed by some pair styles which require
the velocity state of both the I and J particles to compute a pairwise
I,J interaction, as well as by some compute and fix commands.

\sphinxAtStartPar
Note that if the {\hyperref[\detokenize{fix_deform::doc}]{\sphinxcrossref{\DUrole{doc}{fix deform}}}} command is being used
with its “remap v” option enabled, then the velocities for ghost atoms
(in the fix deform group) mirrored across a periodic boundary will
also include components due to any velocity shift that occurs across
that boundary (e.g. due to dilation or shear).


\subsection{Restrictions}
\label{\detokenize{comm_modify:restrictions}}
\sphinxAtStartPar
Communication mode \sphinxstyleemphasis{multi} is currently only available for
{\hyperref[\detokenize{comm_style::doc}]{\sphinxcrossref{\DUrole{doc}{comm\_style}}}} \sphinxstyleemphasis{brick}.


\subsection{Related commands}
\label{\detokenize{comm_modify:related-commands}}
\sphinxAtStartPar
{\hyperref[\detokenize{comm_style::doc}]{\sphinxcrossref{\DUrole{doc}{comm\_style}}}}, {\hyperref[\detokenize{neighbor::doc}]{\sphinxcrossref{\DUrole{doc}{neighbor}}}}


\subsection{Default}
\label{\detokenize{comm_modify:default}}
\sphinxAtStartPar
The option defaults are mode = single, group = all, cutoff = 0.0, vel =
no.  The cutoff default of 0.0 means that ghost cutoff = neighbor
cutoff = pairwise force cutoff + neighbor skin.

\sphinxstepscope

\index{comm\_style@\spxentry{comm\_style}}\ignorespaces 

\section{comm\_style command}
\label{\detokenize{comm_style:comm-style-command}}\label{\detokenize{comm_style:index-0}}\label{\detokenize{comm_style::doc}}

\subsection{Syntax}
\label{\detokenize{comm_style:syntax}}
\begin{sphinxVerbatim}[commandchars=\\\{\}]
\PYG{k}{comm\PYGZus{}style}\PYG{+w}{ }\PYG{n}{style}
\end{sphinxVerbatim}
\begin{itemize}
\item {} 
\sphinxAtStartPar
style = \sphinxstyleemphasis{brick} or \sphinxstyleemphasis{tiled}

\end{itemize}


\subsection{Examples}
\label{\detokenize{comm_style:examples}}
\begin{sphinxVerbatim}[commandchars=\\\{\}]
\PYG{k}{comm\PYGZus{}style}\PYG{+w}{ }\PYG{n}{brick}
\PYG{k}{comm\PYGZus{}style}\PYG{+w}{ }\PYG{n}{tiled}
\end{sphinxVerbatim}


\subsection{Description}
\label{\detokenize{comm_style:description}}
\sphinxAtStartPar
This command sets the style of inter\sphinxhyphen{}processor communication of atom
information that occurs each timestep as coordinates and other
properties are exchanged between neighboring processors and stored as
properties of ghost atoms.

\sphinxAtStartPar
For the default \sphinxstyleemphasis{brick} style, the domain decomposition used by LAMMPS
to partition the simulation box must be a regular 3d grid of bricks,
one per processor.  Each processor communicates with its 6 Cartesian
neighbors in the grid to acquire information for nearby atoms.

\sphinxAtStartPar
For the \sphinxstyleemphasis{tiled} style, a more general domain decomposition can be
used, as triggered by the {\hyperref[\detokenize{balance::doc}]{\sphinxcrossref{\DUrole{doc}{balance}}}} or {\hyperref[\detokenize{fix_balance::doc}]{\sphinxcrossref{\DUrole{doc}{fix balance}}}} commands.  The simulation box can be
partitioned into non\sphinxhyphen{}overlapping rectangular\sphinxhyphen{}shaped “tiles” or varying
sizes and shapes.  Again there is one tile per processor.  To acquire
information for nearby atoms, communication must now be done with a
more complex pattern of neighboring processors.

\sphinxAtStartPar
Note that this command does not actually define a partitioning of the
simulation box (a domain decomposition), rather it determines what
kinds of decompositions are allowed and the pattern of communication
used to enable the decomposition.  A decomposition is created when the
simulation box is first created, via the {\hyperref[\detokenize{create_box::doc}]{\sphinxcrossref{\DUrole{doc}{create\_box}}}}
or {\hyperref[\detokenize{read_data::doc}]{\sphinxcrossref{\DUrole{doc}{read\_data}}}} or {\hyperref[\detokenize{read_restart::doc}]{\sphinxcrossref{\DUrole{doc}{read\_restart}}}}
commands.  For both the \sphinxstyleemphasis{brick} and \sphinxstyleemphasis{tiled} styles, the initial
decomposition will be the same, as described by
{\hyperref[\detokenize{create_box::doc}]{\sphinxcrossref{\DUrole{doc}{create\_box}}}} and {\hyperref[\detokenize{processors::doc}]{\sphinxcrossref{\DUrole{doc}{processors}}}}
commands.  The decomposition can be changed via the
{\hyperref[\detokenize{balance::doc}]{\sphinxcrossref{\DUrole{doc}{balance}}}} or {\hyperref[\detokenize{fix_balance::doc}]{\sphinxcrossref{\DUrole{doc}{fix balance}}}} commands.


\subsection{Restrictions}
\label{\detokenize{comm_style:restrictions}}
\sphinxAtStartPar
none


\subsection{Related commands}
\label{\detokenize{comm_style:related-commands}}
\sphinxAtStartPar
{\hyperref[\detokenize{comm_modify::doc}]{\sphinxcrossref{\DUrole{doc}{comm\_modify}}}}, {\hyperref[\detokenize{processors::doc}]{\sphinxcrossref{\DUrole{doc}{processors}}}},
{\hyperref[\detokenize{balance::doc}]{\sphinxcrossref{\DUrole{doc}{balance}}}}, {\hyperref[\detokenize{fix_balance::doc}]{\sphinxcrossref{\DUrole{doc}{fix balance}}}}


\subsection{Default}
\label{\detokenize{comm_style:default}}
\sphinxAtStartPar
The default style is brick.

\sphinxstepscope

\index{compute@\spxentry{compute}}\ignorespaces 

\section{compute command}
\label{\detokenize{compute:compute-command}}\label{\detokenize{compute:index-0}}\label{\detokenize{compute::doc}}

\subsection{Syntax}
\label{\detokenize{compute:syntax}}
\begin{sphinxVerbatim}[commandchars=\\\{\}]
\PYG{k}{compute }\PYG{n+nv+nvIdentifier}{ID}\PYG{+w}{ }\PYG{n+nv+nvIdentifier}{group\PYGZhy{}ID}\PYG{+w}{ }\PYG{n}{style}\PYG{+w}{ }\PYG{n}{args}
\end{sphinxVerbatim}
\begin{itemize}
\item {} 
\sphinxAtStartPar
ID = user\sphinxhyphen{}assigned name for the computation

\item {} 
\sphinxAtStartPar
group\sphinxhyphen{}ID = ID of the group of atoms to perform the computation on

\item {} 
\sphinxAtStartPar
style = one of a list of possible style names (see below)

\item {} 
\sphinxAtStartPar
args = arguments used by a particular style

\end{itemize}


\subsection{Examples}
\label{\detokenize{compute:examples}}
\begin{sphinxVerbatim}[commandchars=\\\{\}]
\PYG{k}{compute }\PYG{n+nv+nvIdentifier}{1}\PYG{+w}{ }\PYG{n+nv+nvIdentifier}{all}\PYG{+w}{ }\PYG{n}{temp}
\PYG{k}{compute }\PYG{n+nv+nvIdentifier}{newtemp}\PYG{+w}{ }\PYG{n+nv+nvIdentifier}{flow}\PYG{+w}{ }\PYG{n}{temp}\PYG{o}{/}\PYG{n}{partial}\PYG{+w}{ }\PYG{l+m}{1}\PYG{+w}{ }\PYG{l+m}{1}\PYG{+w}{ }\PYG{l+m}{0}
\PYG{k}{compute }\PYG{n+nv+nvIdentifier}{3}\PYG{+w}{ }\PYG{n+nv+nvIdentifier}{all}\PYG{+w}{ }\PYG{n}{ke}\PYG{o}{/}\PYG{n}{atom}
\end{sphinxVerbatim}


\subsection{Description}
\label{\detokenize{compute:description}}
\sphinxAtStartPar
Define a diagnostic computation that will be performed on a group of
atoms.  Quantities calculated by a compute are instantaneous values,
meaning they are calculated from information about atoms on the
current timestep or iteration, though internally a compute may store
some information about a previous state of the system.  Defining a
compute does not perform the computation.  Instead computes are
invoked by other LAMMPS commands as needed (e.g., to calculate a
temperature needed for a thermostat fix or to generate thermodynamic
or dump file output).  See the {\hyperref[\detokenize{Howto_output::doc}]{\sphinxcrossref{\DUrole{doc}{Howto output}}}} page
for a summary of various LAMMPS output options, many of which involve
computes.

\sphinxAtStartPar
The ID of a compute can only contain alphanumeric characters and
underscores.


\bigskip\hrule\bigskip


\sphinxAtStartPar
Computes calculate and store any of four \sphinxstyleemphasis{styles} of quantities:
global, per\sphinxhyphen{}atom, local, or per\sphinxhyphen{}grid.

\sphinxAtStartPar
A global quantity is one or more system\sphinxhyphen{}wide values, e.g. the
temperature of the system.  A per\sphinxhyphen{}atom quantity is one or more values
per atom, e.g. the kinetic energy of each atom.  Per\sphinxhyphen{}atom values are
set to 0.0 for atoms not in the specified compute group.  Local
quantities are calculated by each processor based on the atoms it
owns, but there may be zero or more per atom, e.g. a list of bond
distances.  Per\sphinxhyphen{}grid quantities are calculated on a regular 2d or 3d
grid which overlays a 2d or 3d simulation domain.  The grid points and
the data they store are distributed across processors; each processor
owns the grid points which fall within its subdomain.

\sphinxAtStartPar
As a general rule of thumb, computes that produce per\sphinxhyphen{}atom quantities
have the word “atom” at the end of their style, e.g. \sphinxstyleemphasis{ke/atom}.
Computes that produce local quantities have the word “local” at the
end of their style, e.g. \sphinxstyleemphasis{bond/local}.  Computes that produce
per\sphinxhyphen{}grid quantities have the word “grid” at the end of their style,
e.g. \sphinxstyleemphasis{property/grid}.  And styles with neither “atom” or “local” or
“grid” at the end of their style name produce global quantities.

\sphinxAtStartPar
Global, per\sphinxhyphen{}atom, local, and per\sphinxhyphen{}grid quantities can also be of three
\sphinxstyleemphasis{kinds}: a single scalar value (global only), a vector of values, or a
2d array of values.  For per\sphinxhyphen{}atom, local, and per\sphinxhyphen{}grid quantities, a
“vector” means a single value for each atom, each local entity
(e.g. bond), or grid cell.  Likewise an “array”, means multiple values
for each atom, each local entity, or each grid cell.

\sphinxAtStartPar
Note that a single compute can produce any combination of global,
per\sphinxhyphen{}atom, local, or per\sphinxhyphen{}grid values.  Likewise it can produce any
combination of scalar, vector, or array output for each style.  The
exception is that for per\sphinxhyphen{}atom, local, and per\sphinxhyphen{}grid output, either a
vector or array can be produced, but not both.  The doc page for each
compute explains the values it produces.

\sphinxAtStartPar
When a compute output is accessed by another input script command it
is referenced via the following bracket notation, where ID is the ID
of the compute:


\begin{savenotes}\sphinxattablestart
\sphinxthistablewithglobalstyle
\centering
\begin{tabulary}{\linewidth}[t]{TT}
\sphinxtoprule
\sphinxtableatstartofbodyhook
\sphinxAtStartPar
c\_ID
&
\sphinxAtStartPar
entire scalar, vector, or array
\\
\sphinxhline
\sphinxAtStartPar
c\_ID{[}I{]}
&
\sphinxAtStartPar
one element of vector, one column of array
\\
\sphinxhline
\sphinxAtStartPar
c\_ID{[}I{]}{[}J{]}
&
\sphinxAtStartPar
one element of array
\\
\sphinxbottomrule
\end{tabulary}
\sphinxtableafterendhook\par
\sphinxattableend\end{savenotes}

\sphinxAtStartPar
In other words, using one bracket reduces the dimension of the
quantity once (vector \(\to\) scalar, array \(\to\) vector).
Using two brackets reduces the dimension twice (array \(\to\)
scalar).  Thus, for example, a command that uses global scalar compute
values as input can also process elements of a vector or array.
Depending on the command, this can either be done directly using the
syntax in the table, or by first defining a {\hyperref[\detokenize{variable::doc}]{\sphinxcrossref{\DUrole{doc}{variable}}}}
of the appropriate style to store the quantity, then using the
variable as an input to the command.

\sphinxAtStartPar
Note that commands and {\hyperref[\detokenize{variable::doc}]{\sphinxcrossref{\DUrole{doc}{variables}}}} which take compute
outputs as input typically do not allow for all styles and kinds of
data (e.g., a command may require global but not per\sphinxhyphen{}atom values, or
it may require a vector of values, not a scalar).  This means there is
typically no ambiguity about referring to a compute output as c\_ID
even if it produces, for example, both a scalar and vector.  The doc
pages for various commands explain the details, including how any
ambiguities are resolved.


\bigskip\hrule\bigskip


\sphinxAtStartPar
In LAMMPS, the values generated by a compute can be used in several
ways:
\begin{itemize}
\item {} 
\sphinxAtStartPar
The results of computes that calculate a global temperature or
pressure can be used by fixes that do thermostatting or barostatting
or when atom velocities are created.

\item {} 
\sphinxAtStartPar
Global values can be output via the {\hyperref[\detokenize{thermo_style::doc}]{\sphinxcrossref{\DUrole{doc}{thermo\_style custom}}}} or {\hyperref[\detokenize{fix_ave_time::doc}]{\sphinxcrossref{\DUrole{doc}{fix ave/time}}}} command.
Or the values can be referenced in a {\hyperref[\detokenize{variable::doc}]{\sphinxcrossref{\DUrole{doc}{variable equal}}}} or
{\hyperref[\detokenize{variable::doc}]{\sphinxcrossref{\DUrole{doc}{variable atom}}}} command.

\item {} 
\sphinxAtStartPar
Per\sphinxhyphen{}atom values can be output via the {\hyperref[\detokenize{dump::doc}]{\sphinxcrossref{\DUrole{doc}{dump custom}}}} command.
Or they can be time\sphinxhyphen{}averaged via the {\hyperref[\detokenize{fix_ave_atom::doc}]{\sphinxcrossref{\DUrole{doc}{fix ave/atom}}}}
command or reduced by the {\hyperref[\detokenize{compute_reduce::doc}]{\sphinxcrossref{\DUrole{doc}{compute reduce}}}}
command.  Or the per\sphinxhyphen{}atom values can be referenced in an {\hyperref[\detokenize{variable::doc}]{\sphinxcrossref{\DUrole{doc}{atom\sphinxhyphen{}style variable}}}}.

\item {} 
\sphinxAtStartPar
Local values can be reduced by the {\hyperref[\detokenize{compute_reduce::doc}]{\sphinxcrossref{\DUrole{doc}{compute reduce}}}} command, or histogrammed by the {\hyperref[\detokenize{fix_ave_histo::doc}]{\sphinxcrossref{\DUrole{doc}{fix ave/histo}}}} command, or output by the {\hyperref[\detokenize{dump::doc}]{\sphinxcrossref{\DUrole{doc}{dump local}}}} command.

\end{itemize}

\sphinxAtStartPar
The results of computes that calculate global quantities can be either
“intensive” or “extensive” values.  Intensive means the value is
independent of the number of atoms in the simulation
(e.g., temperature).  Extensive means the value scales with the number of
atoms in the simulation (e.g., total rotational kinetic energy).
{\hyperref[\detokenize{thermo_style::doc}]{\sphinxcrossref{\DUrole{doc}{Thermodynamic output}}}} will normalize extensive
values by the number of atoms in the system, depending on the
“thermo\_modify norm” setting.  It will not normalize intensive values.
If a compute value is accessed in another way (e.g., by a
{\hyperref[\detokenize{variable::doc}]{\sphinxcrossref{\DUrole{doc}{variable}}}}), you may want to know whether it is an
intensive or extensive value.  See the page for individual
computes for further info.


\bigskip\hrule\bigskip


\sphinxAtStartPar
LAMMPS creates its own computes internally for thermodynamic output.
Three computes are always created, named “thermo\_temp”,
“thermo\_press”, and “thermo\_pe”, as if these commands had been invoked
in the input script:

\begin{sphinxVerbatim}[commandchars=\\\{\}]
\PYG{k}{compute }\PYG{n+nv+nvIdentifier}{thermo\PYGZus{}temp}\PYG{+w}{ }\PYG{n+nv+nvIdentifier}{all}\PYG{+w}{ }\PYG{n}{temp}
\PYG{k}{compute }\PYG{n+nv+nvIdentifier}{thermo\PYGZus{}press}\PYG{+w}{ }\PYG{n+nv+nvIdentifier}{all}\PYG{+w}{ }\PYG{n}{pressure}\PYG{+w}{ }\PYG{n}{thermo\PYGZus{}temp}
\PYG{k}{compute }\PYG{n+nv+nvIdentifier}{thermo\PYGZus{}pe}\PYG{+w}{ }\PYG{n+nv+nvIdentifier}{all}\PYG{+w}{ }\PYG{n}{pe}
\end{sphinxVerbatim}

\sphinxAtStartPar
Additional computes for other quantities are created if the thermo
style requires it.  See the documentation for the
{\hyperref[\detokenize{thermo_style::doc}]{\sphinxcrossref{\DUrole{doc}{thermo\_style}}}} command.

\sphinxAtStartPar
Fixes that calculate temperature or pressure, i.e. for thermostatting
or barostatting, may also create computes.  These are discussed in the
documentation for specific {\hyperref[\detokenize{fix::doc}]{\sphinxcrossref{\DUrole{doc}{fix}}}} commands.

\sphinxAtStartPar
In all these cases, the default computes LAMMPS creates can be
replaced by computes defined by the user in the input script, as
described by the {\hyperref[\detokenize{thermo_modify::doc}]{\sphinxcrossref{\DUrole{doc}{thermo\_modify}}}} and {\hyperref[\detokenize{fix_modify::doc}]{\sphinxcrossref{\DUrole{doc}{fix modify}}}} commands.

\sphinxAtStartPar
Properties of either a default or user\sphinxhyphen{}defined compute can be modified
via the {\hyperref[\detokenize{compute_modify::doc}]{\sphinxcrossref{\DUrole{doc}{compute\_modify}}}} command.

\sphinxAtStartPar
Computes can be deleted with the {\hyperref[\detokenize{uncompute::doc}]{\sphinxcrossref{\DUrole{doc}{uncompute}}}} command.

\sphinxAtStartPar
Code for new computes can be added to LAMMPS; see the
{\hyperref[\detokenize{Modify::doc}]{\sphinxcrossref{\DUrole{doc}{Modify}}}} page for details.  The results of their
calculations accessed in the various ways described above.


\bigskip\hrule\bigskip


\sphinxAtStartPar
Each compute style has its own page which describes its arguments
and what it does.  Here is an alphabetic list of compute styles
available in LAMMPS.  They are also listed in more compact form on the
{\hyperref[\detokenize{Commands_compute::doc}]{\sphinxcrossref{\DUrole{doc}{Commands compute}}}} doc page.

\sphinxAtStartPar
There are also additional accelerated compute styles included in the
LAMMPS distribution for faster performance on CPUs, GPUs, and KNLs.
The individual style names on the {\hyperref[\detokenize{Commands_compute::doc}]{\sphinxcrossref{\DUrole{doc}{Commands compute}}}} page are followed by one or more of
(g,i,k,o,t) to indicate which accelerated styles exist.
\begin{itemize}
\item {} 
\sphinxAtStartPar
{\hyperref[\detokenize{compute_ackland_atom::doc}]{\sphinxcrossref{\DUrole{doc}{ackland/atom}}}} \sphinxhyphen{} determines the local lattice structure based on the Ackland formulation

\item {} 
\sphinxAtStartPar
{\hyperref[\detokenize{compute_adf::doc}]{\sphinxcrossref{\DUrole{doc}{adf}}}} \sphinxhyphen{} angular distribution function of triples of atoms

\item {} 
\sphinxAtStartPar
{\hyperref[\detokenize{compute_cluster_atom::doc}]{\sphinxcrossref{\DUrole{doc}{aggregate/atom}}}} \sphinxhyphen{} aggregate ID for each atom

\item {} 
\sphinxAtStartPar
{\hyperref[\detokenize{compute_angle::doc}]{\sphinxcrossref{\DUrole{doc}{angle}}}} \sphinxhyphen{} energy of each angle sub\sphinxhyphen{}style

\item {} 
\sphinxAtStartPar
{\hyperref[\detokenize{compute_angle_local::doc}]{\sphinxcrossref{\DUrole{doc}{angle/local}}}} \sphinxhyphen{} theta and energy of each angle

\item {} 
\sphinxAtStartPar
{\hyperref[\detokenize{compute_angmom_chunk::doc}]{\sphinxcrossref{\DUrole{doc}{angmom/chunk}}}} \sphinxhyphen{} angular momentum for each chunk

\item {} 
\sphinxAtStartPar
{\hyperref[\detokenize{compute_ave_sphere_atom::doc}]{\sphinxcrossref{\DUrole{doc}{ave/sphere/atom}}}} \sphinxhyphen{} compute local density and temperature around each atom

\item {} 
\sphinxAtStartPar
{\hyperref[\detokenize{compute_basal_atom::doc}]{\sphinxcrossref{\DUrole{doc}{basal/atom}}}} \sphinxhyphen{} calculates the hexagonal close\sphinxhyphen{}packed “c” lattice vector of each atom

\item {} 
\sphinxAtStartPar
{\hyperref[\detokenize{compute_body_local::doc}]{\sphinxcrossref{\DUrole{doc}{body/local}}}} \sphinxhyphen{} attributes of body sub\sphinxhyphen{}particles

\item {} 
\sphinxAtStartPar
{\hyperref[\detokenize{compute_bond::doc}]{\sphinxcrossref{\DUrole{doc}{bond}}}} \sphinxhyphen{} energy of each bond sub\sphinxhyphen{}style

\item {} 
\sphinxAtStartPar
{\hyperref[\detokenize{compute_bond_local::doc}]{\sphinxcrossref{\DUrole{doc}{bond/local}}}} \sphinxhyphen{} distance and energy of each bond

\item {} 
\sphinxAtStartPar
{\hyperref[\detokenize{compute_born_matrix::doc}]{\sphinxcrossref{\DUrole{doc}{born/matrix}}}} \sphinxhyphen{} second derivative or potential with respect to strain

\item {} 
\sphinxAtStartPar
{\hyperref[\detokenize{compute_centro_atom::doc}]{\sphinxcrossref{\DUrole{doc}{centro/atom}}}} \sphinxhyphen{} centro\sphinxhyphen{}symmetry parameter for each atom

\item {} 
\sphinxAtStartPar
{\hyperref[\detokenize{compute_stress_atom::doc}]{\sphinxcrossref{\DUrole{doc}{centroid/stress/atom}}}} \sphinxhyphen{} centroid based stress tensor for each atom

\item {} 
\sphinxAtStartPar
{\hyperref[\detokenize{compute_chunk_atom::doc}]{\sphinxcrossref{\DUrole{doc}{chunk/atom}}}} \sphinxhyphen{} assign chunk IDs to each atom

\item {} 
\sphinxAtStartPar
{\hyperref[\detokenize{compute_chunk_spread_atom::doc}]{\sphinxcrossref{\DUrole{doc}{chunk/spread/atom}}}} \sphinxhyphen{} spreads chunk values to each atom in chunk

\item {} 
\sphinxAtStartPar
{\hyperref[\detokenize{compute_cluster_atom::doc}]{\sphinxcrossref{\DUrole{doc}{cluster/atom}}}} \sphinxhyphen{} cluster ID for each atom

\item {} 
\sphinxAtStartPar
{\hyperref[\detokenize{compute_cna_atom::doc}]{\sphinxcrossref{\DUrole{doc}{cna/atom}}}} \sphinxhyphen{} common neighbor analysis (CNA) for each atom

\item {} 
\sphinxAtStartPar
{\hyperref[\detokenize{compute_cnp_atom::doc}]{\sphinxcrossref{\DUrole{doc}{cnp/atom}}}} \sphinxhyphen{} common neighborhood parameter (CNP) for each atom

\item {} 
\sphinxAtStartPar
{\hyperref[\detokenize{compute_com::doc}]{\sphinxcrossref{\DUrole{doc}{com}}}} \sphinxhyphen{} center of mass of group of atoms

\item {} 
\sphinxAtStartPar
{\hyperref[\detokenize{compute_com_chunk::doc}]{\sphinxcrossref{\DUrole{doc}{com/chunk}}}} \sphinxhyphen{} center of mass for each chunk

\item {} 
\sphinxAtStartPar
{\hyperref[\detokenize{compute_contact_atom::doc}]{\sphinxcrossref{\DUrole{doc}{contact/atom}}}} \sphinxhyphen{} contact count for each spherical particle

\item {} 
\sphinxAtStartPar
{\hyperref[\detokenize{compute_coord_atom::doc}]{\sphinxcrossref{\DUrole{doc}{coord/atom}}}} \sphinxhyphen{} coordination number for each atom

\item {} 
\sphinxAtStartPar
{\hyperref[\detokenize{compute_count_type::doc}]{\sphinxcrossref{\DUrole{doc}{count/type}}}} \sphinxhyphen{} count of atoms or bonds by type

\item {} 
\sphinxAtStartPar
{\hyperref[\detokenize{compute_damage_atom::doc}]{\sphinxcrossref{\DUrole{doc}{damage/atom}}}} \sphinxhyphen{} Peridynamic damage for each atom

\item {} 
\sphinxAtStartPar
{\hyperref[\detokenize{compute_dihedral::doc}]{\sphinxcrossref{\DUrole{doc}{dihedral}}}} \sphinxhyphen{} energy of each dihedral sub\sphinxhyphen{}style

\item {} 
\sphinxAtStartPar
{\hyperref[\detokenize{compute_dihedral_local::doc}]{\sphinxcrossref{\DUrole{doc}{dihedral/local}}}} \sphinxhyphen{} angle of each dihedral

\item {} 
\sphinxAtStartPar
{\hyperref[\detokenize{compute_dilatation_atom::doc}]{\sphinxcrossref{\DUrole{doc}{dilatation/atom}}}} \sphinxhyphen{} Peridynamic dilatation for each atom

\item {} 
\sphinxAtStartPar
{\hyperref[\detokenize{compute_dipole::doc}]{\sphinxcrossref{\DUrole{doc}{dipole}}}} \sphinxhyphen{} dipole vector and total dipole

\item {} 
\sphinxAtStartPar
{\hyperref[\detokenize{compute_dipole_chunk::doc}]{\sphinxcrossref{\DUrole{doc}{dipole/chunk}}}} \sphinxhyphen{} dipole vector and total dipole for each chunk

\item {} 
\sphinxAtStartPar
{\hyperref[\detokenize{compute_dipole::doc}]{\sphinxcrossref{\DUrole{doc}{dipole/tip4p}}}} \sphinxhyphen{} dipole vector and total dipole with TIP4P pair style

\item {} 
\sphinxAtStartPar
{\hyperref[\detokenize{compute_dipole_chunk::doc}]{\sphinxcrossref{\DUrole{doc}{dipole/tip4p/chunk}}}} \sphinxhyphen{} dipole vector and total dipole for each chunk with TIP4P pair style

\item {} 
\sphinxAtStartPar
{\hyperref[\detokenize{compute_displace_atom::doc}]{\sphinxcrossref{\DUrole{doc}{displace/atom}}}} \sphinxhyphen{} displacement of each atom

\item {} 
\sphinxAtStartPar
{\hyperref[\detokenize{compute_dpd::doc}]{\sphinxcrossref{\DUrole{doc}{dpd}}}} \sphinxhyphen{} total values of internal conductive energy, internal mechanical energy, chemical energy, and harmonic average of internal temperature

\item {} 
\sphinxAtStartPar
{\hyperref[\detokenize{compute_dpd_atom::doc}]{\sphinxcrossref{\DUrole{doc}{dpd/atom}}}} \sphinxhyphen{} per\sphinxhyphen{}particle values of internal conductive energy, internal mechanical energy, chemical energy, and internal temperature

\item {} 
\sphinxAtStartPar
{\hyperref[\detokenize{compute_edpd_temp_atom::doc}]{\sphinxcrossref{\DUrole{doc}{edpd/temp/atom}}}} \sphinxhyphen{} per\sphinxhyphen{}atom temperature for each eDPD particle in a group

\item {} 
\sphinxAtStartPar
{\hyperref[\detokenize{compute_efield_atom::doc}]{\sphinxcrossref{\DUrole{doc}{efield/atom}}}} \sphinxhyphen{} electric field at each atom

\item {} 
\sphinxAtStartPar
{\hyperref[\detokenize{compute_efield_wolf_atom::doc}]{\sphinxcrossref{\DUrole{doc}{efield/wolf/atom}}}} \sphinxhyphen{} electric field at each atom

\item {} 
\sphinxAtStartPar
{\hyperref[\detokenize{compute_entropy_atom::doc}]{\sphinxcrossref{\DUrole{doc}{entropy/atom}}}} \sphinxhyphen{} pair entropy fingerprint of each atom

\item {} 
\sphinxAtStartPar
{\hyperref[\detokenize{compute_erotate_asphere::doc}]{\sphinxcrossref{\DUrole{doc}{erotate/asphere}}}} \sphinxhyphen{} rotational energy of aspherical particles

\item {} 
\sphinxAtStartPar
{\hyperref[\detokenize{compute_erotate_rigid::doc}]{\sphinxcrossref{\DUrole{doc}{erotate/rigid}}}} \sphinxhyphen{} rotational energy of rigid bodies

\item {} 
\sphinxAtStartPar
{\hyperref[\detokenize{compute_erotate_sphere::doc}]{\sphinxcrossref{\DUrole{doc}{erotate/sphere}}}} \sphinxhyphen{} rotational energy of spherical particles

\item {} 
\sphinxAtStartPar
{\hyperref[\detokenize{compute_erotate_sphere_atom::doc}]{\sphinxcrossref{\DUrole{doc}{erotate/sphere/atom}}}} \sphinxhyphen{} rotational energy for each spherical particle

\item {} 
\sphinxAtStartPar
{\hyperref[\detokenize{compute_event_displace::doc}]{\sphinxcrossref{\DUrole{doc}{event/displace}}}} \sphinxhyphen{} detect event on atom displacement

\item {} 
\sphinxAtStartPar
{\hyperref[\detokenize{compute_fabric::doc}]{\sphinxcrossref{\DUrole{doc}{fabric}}}} \sphinxhyphen{} calculates fabric tensors from pair interactions

\item {} 
\sphinxAtStartPar
{\hyperref[\detokenize{compute_fep::doc}]{\sphinxcrossref{\DUrole{doc}{fep}}}} \sphinxhyphen{} compute free energies for alchemical transformation from perturbation theory

\item {} 
\sphinxAtStartPar
{\hyperref[\detokenize{compute_fep_ta::doc}]{\sphinxcrossref{\DUrole{doc}{fep/ta}}}} \sphinxhyphen{} compute free energies for a test area perturbation

\item {} 
\sphinxAtStartPar
{\hyperref[\detokenize{compute_tally::doc}]{\sphinxcrossref{\DUrole{doc}{force/tally}}}} \sphinxhyphen{} force between two groups of atoms via the tally callback mechanism

\item {} 
\sphinxAtStartPar
{\hyperref[\detokenize{compute_cluster_atom::doc}]{\sphinxcrossref{\DUrole{doc}{fragment/atom}}}} \sphinxhyphen{} fragment ID for each atom

\item {} 
\sphinxAtStartPar
{\hyperref[\detokenize{compute_global_atom::doc}]{\sphinxcrossref{\DUrole{doc}{global/atom}}}} \sphinxhyphen{} assign global values to each atom from arrays of global values

\item {} 
\sphinxAtStartPar
{\hyperref[\detokenize{compute_group_group::doc}]{\sphinxcrossref{\DUrole{doc}{group/group}}}} \sphinxhyphen{} energy/force between two groups of atoms

\item {} 
\sphinxAtStartPar
{\hyperref[\detokenize{compute_gyration::doc}]{\sphinxcrossref{\DUrole{doc}{gyration}}}} \sphinxhyphen{} radius of gyration of group of atoms

\item {} 
\sphinxAtStartPar
{\hyperref[\detokenize{compute_gyration_chunk::doc}]{\sphinxcrossref{\DUrole{doc}{gyration/chunk}}}} \sphinxhyphen{} radius of gyration for each chunk

\item {} 
\sphinxAtStartPar
{\hyperref[\detokenize{compute_gyration_shape::doc}]{\sphinxcrossref{\DUrole{doc}{gyration/shape}}}} \sphinxhyphen{} shape parameters from gyration tensor

\item {} 
\sphinxAtStartPar
{\hyperref[\detokenize{compute_gyration_shape_chunk::doc}]{\sphinxcrossref{\DUrole{doc}{gyration/shape/chunk}}}} \sphinxhyphen{} shape parameters from gyration tensor for each chunk

\item {} 
\sphinxAtStartPar
{\hyperref[\detokenize{compute_heat_flux::doc}]{\sphinxcrossref{\DUrole{doc}{heat/flux}}}} \sphinxhyphen{} heat flux through a group of atoms

\item {} 
\sphinxAtStartPar
{\hyperref[\detokenize{compute_tally::doc}]{\sphinxcrossref{\DUrole{doc}{heat/flux/tally}}}} \sphinxhyphen{} heat flux through a group of atoms via the tally callback mechanism

\item {} 
\sphinxAtStartPar
{\hyperref[\detokenize{compute_tally::doc}]{\sphinxcrossref{\DUrole{doc}{heat/flux/virial/tally}}}} \sphinxhyphen{} virial heat flux between two groups via the tally callback mechanism

\item {} 
\sphinxAtStartPar
{\hyperref[\detokenize{compute_hexorder_atom::doc}]{\sphinxcrossref{\DUrole{doc}{hexorder/atom}}}} \sphinxhyphen{} bond orientational order parameter q6

\item {} 
\sphinxAtStartPar
{\hyperref[\detokenize{compute_hma::doc}]{\sphinxcrossref{\DUrole{doc}{hma}}}} \sphinxhyphen{} harmonically mapped averaging for atomic crystals

\item {} 
\sphinxAtStartPar
{\hyperref[\detokenize{compute_improper::doc}]{\sphinxcrossref{\DUrole{doc}{improper}}}} \sphinxhyphen{} energy of each improper sub\sphinxhyphen{}style

\item {} 
\sphinxAtStartPar
{\hyperref[\detokenize{compute_improper_local::doc}]{\sphinxcrossref{\DUrole{doc}{improper/local}}}} \sphinxhyphen{} angle of each improper

\item {} 
\sphinxAtStartPar
{\hyperref[\detokenize{compute_inertia_chunk::doc}]{\sphinxcrossref{\DUrole{doc}{inertia/chunk}}}} \sphinxhyphen{} inertia tensor for each chunk

\item {} 
\sphinxAtStartPar
{\hyperref[\detokenize{compute_ke::doc}]{\sphinxcrossref{\DUrole{doc}{ke}}}} \sphinxhyphen{} translational kinetic energy

\item {} 
\sphinxAtStartPar
{\hyperref[\detokenize{compute_ke_atom::doc}]{\sphinxcrossref{\DUrole{doc}{ke/atom}}}} \sphinxhyphen{} kinetic energy for each atom

\item {} 
\sphinxAtStartPar
{\hyperref[\detokenize{compute_ke_atom_eff::doc}]{\sphinxcrossref{\DUrole{doc}{ke/atom/eff}}}} \sphinxhyphen{} per\sphinxhyphen{}atom translational and radial kinetic energy in the electron force field model

\item {} 
\sphinxAtStartPar
{\hyperref[\detokenize{compute_ke_eff::doc}]{\sphinxcrossref{\DUrole{doc}{ke/eff}}}} \sphinxhyphen{} kinetic energy of a group of nuclei and electrons in the electron force field model

\item {} 
\sphinxAtStartPar
{\hyperref[\detokenize{compute_ke_rigid::doc}]{\sphinxcrossref{\DUrole{doc}{ke/rigid}}}} \sphinxhyphen{} translational kinetic energy of rigid bodies

\item {} 
\sphinxAtStartPar
{\hyperref[\detokenize{compute_composition_atom::doc}]{\sphinxcrossref{\DUrole{doc}{composition/atom}}}} \sphinxhyphen{} local composition for each atom

\item {} 
\sphinxAtStartPar
{\hyperref[\detokenize{compute_mliap::doc}]{\sphinxcrossref{\DUrole{doc}{mliap}}}} \sphinxhyphen{} gradients of energy and forces with respect to model parameters and related quantities for training machine learning interatomic potentials

\item {} 
\sphinxAtStartPar
{\hyperref[\detokenize{compute_momentum::doc}]{\sphinxcrossref{\DUrole{doc}{momentum}}}} \sphinxhyphen{} translational momentum

\item {} 
\sphinxAtStartPar
{\hyperref[\detokenize{compute_msd::doc}]{\sphinxcrossref{\DUrole{doc}{msd}}}} \sphinxhyphen{} mean\sphinxhyphen{}squared displacement of group of atoms

\item {} 
\sphinxAtStartPar
{\hyperref[\detokenize{compute_msd_chunk::doc}]{\sphinxcrossref{\DUrole{doc}{msd/chunk}}}} \sphinxhyphen{} mean\sphinxhyphen{}squared displacement for each chunk

\item {} 
\sphinxAtStartPar
{\hyperref[\detokenize{compute_msd_nongauss::doc}]{\sphinxcrossref{\DUrole{doc}{msd/nongauss}}}} \sphinxhyphen{} MSD and non\sphinxhyphen{}Gaussian parameter of group of atoms

\item {} 
\sphinxAtStartPar
{\hyperref[\detokenize{compute_nbond_atom::doc}]{\sphinxcrossref{\DUrole{doc}{nbond/atom}}}} \sphinxhyphen{} calculates number of bonds per atom

\item {} 
\sphinxAtStartPar
{\hyperref[\detokenize{compute_omega_chunk::doc}]{\sphinxcrossref{\DUrole{doc}{omega/chunk}}}} \sphinxhyphen{} angular velocity for each chunk

\item {} 
\sphinxAtStartPar
{\hyperref[\detokenize{compute_orientorder_atom::doc}]{\sphinxcrossref{\DUrole{doc}{orientorder/atom}}}} \sphinxhyphen{} Steinhardt bond orientational order parameters Ql

\item {} 
\sphinxAtStartPar
{\hyperref[\detokenize{compute_pace::doc}]{\sphinxcrossref{\DUrole{doc}{pace}}}} \sphinxhyphen{} atomic cluster expansion descriptors and related quantities

\item {} 
\sphinxAtStartPar
{\hyperref[\detokenize{compute_pair::doc}]{\sphinxcrossref{\DUrole{doc}{pair}}}} \sphinxhyphen{} values computed by a pair style

\item {} 
\sphinxAtStartPar
{\hyperref[\detokenize{compute_pair_local::doc}]{\sphinxcrossref{\DUrole{doc}{pair/local}}}} \sphinxhyphen{} distance/energy/force of each pairwise interaction

\item {} 
\sphinxAtStartPar
{\hyperref[\detokenize{compute_pe::doc}]{\sphinxcrossref{\DUrole{doc}{pe}}}} \sphinxhyphen{} potential energy

\item {} 
\sphinxAtStartPar
{\hyperref[\detokenize{compute_pe_atom::doc}]{\sphinxcrossref{\DUrole{doc}{pe/atom}}}} \sphinxhyphen{} potential energy for each atom

\item {} 
\sphinxAtStartPar
{\hyperref[\detokenize{compute_tally::doc}]{\sphinxcrossref{\DUrole{doc}{pe/mol/tally}}}} \sphinxhyphen{} potential energy between two groups of atoms separated into intermolecular and intramolecular components via the tally callback mechanism

\item {} 
\sphinxAtStartPar
{\hyperref[\detokenize{compute_tally::doc}]{\sphinxcrossref{\DUrole{doc}{pe/tally}}}} \sphinxhyphen{} potential energy between two groups of atoms via the tally callback mechanism

\item {} 
\sphinxAtStartPar
{\hyperref[\detokenize{compute_plasticity_atom::doc}]{\sphinxcrossref{\DUrole{doc}{plasticity/atom}}}} \sphinxhyphen{} Peridynamic plasticity for each atom

\item {} 
\sphinxAtStartPar
{\hyperref[\detokenize{compute_pod_atom::doc}]{\sphinxcrossref{\DUrole{doc}{pod/atom}}}} \sphinxhyphen{} POD descriptors for each atom

\item {} 
\sphinxAtStartPar
{\hyperref[\detokenize{compute_pod_atom::doc}]{\sphinxcrossref{\DUrole{doc}{podd/atom}}}} \sphinxhyphen{} derivative of POD descriptors for each atom

\item {} 
\sphinxAtStartPar
{\hyperref[\detokenize{compute_pod_atom::doc}]{\sphinxcrossref{\DUrole{doc}{pod/local}}}} \sphinxhyphen{} local POD descriptors and their derivatives

\item {} 
\sphinxAtStartPar
{\hyperref[\detokenize{compute_pod_atom::doc}]{\sphinxcrossref{\DUrole{doc}{pod/global}}}} \sphinxhyphen{} global POD descriptors and their derivatives

\item {} 
\sphinxAtStartPar
{\hyperref[\detokenize{compute_pressure::doc}]{\sphinxcrossref{\DUrole{doc}{pressure}}}} \sphinxhyphen{} total pressure and pressure tensor

\item {} 
\sphinxAtStartPar
{\hyperref[\detokenize{compute_pressure_alchemy::doc}]{\sphinxcrossref{\DUrole{doc}{pressure/alchemy}}}} \sphinxhyphen{} mixed system total pressure and pressure tensor for {\hyperref[\detokenize{fix_alchemy::doc}]{\sphinxcrossref{\DUrole{doc}{fix alchemy}}}} runs

\item {} 
\sphinxAtStartPar
{\hyperref[\detokenize{compute_pressure_uef::doc}]{\sphinxcrossref{\DUrole{doc}{pressure/uef}}}} \sphinxhyphen{} pressure tensor in the reference frame of an applied flow field

\item {} 
\sphinxAtStartPar
{\hyperref[\detokenize{compute_property_atom::doc}]{\sphinxcrossref{\DUrole{doc}{property/atom}}}} \sphinxhyphen{} convert atom attributes to per\sphinxhyphen{}atom vectors/arrays

\item {} 
\sphinxAtStartPar
{\hyperref[\detokenize{compute_property_chunk::doc}]{\sphinxcrossref{\DUrole{doc}{property/chunk}}}} \sphinxhyphen{} extract various per\sphinxhyphen{}chunk attributes

\item {} 
\sphinxAtStartPar
{\hyperref[\detokenize{compute_property_grid::doc}]{\sphinxcrossref{\DUrole{doc}{property/grid}}}} \sphinxhyphen{} convert per\sphinxhyphen{}grid attributes to per\sphinxhyphen{}grid vectors/arrays

\item {} 
\sphinxAtStartPar
{\hyperref[\detokenize{compute_property_local::doc}]{\sphinxcrossref{\DUrole{doc}{property/local}}}} \sphinxhyphen{} convert local attributes to local vectors/arrays

\item {} 
\sphinxAtStartPar
{\hyperref[\detokenize{compute_ptm_atom::doc}]{\sphinxcrossref{\DUrole{doc}{ptm/atom}}}} \sphinxhyphen{} determines the local lattice structure based on the Polyhedral Template Matching method

\item {} 
\sphinxAtStartPar
{\hyperref[\detokenize{compute_rattlers_atom::doc}]{\sphinxcrossref{\DUrole{doc}{rattlers/atom}}}} \sphinxhyphen{} identify under\sphinxhyphen{}coordinated rattler atoms

\item {} 
\sphinxAtStartPar
{\hyperref[\detokenize{compute_rdf::doc}]{\sphinxcrossref{\DUrole{doc}{rdf}}}} \sphinxhyphen{} radial distribution function \(g(r)\) histogram of group of atoms

\item {} 
\sphinxAtStartPar
{\hyperref[\detokenize{compute_reaxff_atom::doc}]{\sphinxcrossref{\DUrole{doc}{reaxff/atom}}}} \sphinxhyphen{} extract ReaxFF bond information

\item {} 
\sphinxAtStartPar
{\hyperref[\detokenize{compute_reduce::doc}]{\sphinxcrossref{\DUrole{doc}{reduce}}}} \sphinxhyphen{} combine per\sphinxhyphen{}atom quantities into a single global value

\item {} 
\sphinxAtStartPar
{\hyperref[\detokenize{compute_reduce_chunk::doc}]{\sphinxcrossref{\DUrole{doc}{reduce/chunk}}}} \sphinxhyphen{} reduce per\sphinxhyphen{}atom quantities within each chunk

\item {} 
\sphinxAtStartPar
{\hyperref[\detokenize{compute_reduce::doc}]{\sphinxcrossref{\DUrole{doc}{reduce/region}}}} \sphinxhyphen{} same as compute reduce, within a region

\item {} 
\sphinxAtStartPar
{\hyperref[\detokenize{compute_rheo_property_atom::doc}]{\sphinxcrossref{\DUrole{doc}{rheo/property/atom}}}} \sphinxhyphen{} convert atom attributes in RHEO package to per\sphinxhyphen{}atom vectors/arrays

\item {} 
\sphinxAtStartPar
{\hyperref[\detokenize{compute_rigid_local::doc}]{\sphinxcrossref{\DUrole{doc}{rigid/local}}}} \sphinxhyphen{} extract rigid body attributes

\item {} 
\sphinxAtStartPar
{\hyperref[\detokenize{compute_saed::doc}]{\sphinxcrossref{\DUrole{doc}{saed}}}} \sphinxhyphen{} electron diffraction intensity on a mesh of reciprocal lattice nodes

\item {} 
\sphinxAtStartPar
{\hyperref[\detokenize{compute_slcsa_atom::doc}]{\sphinxcrossref{\DUrole{doc}{slcsa/atom}}}} \sphinxhyphen{} perform Supervised Learning Crystal Structure Analysis (SL\sphinxhyphen{}CSA)

\item {} 
\sphinxAtStartPar
{\hyperref[\detokenize{compute_slice::doc}]{\sphinxcrossref{\DUrole{doc}{slice}}}} \sphinxhyphen{} extract values from global vector or array

\item {} 
\sphinxAtStartPar
{\hyperref[\detokenize{compute_smd_contact_radius::doc}]{\sphinxcrossref{\DUrole{doc}{smd/contact/radius}}}} \sphinxhyphen{} contact radius for Smooth Mach Dynamics

\item {} 
\sphinxAtStartPar
{\hyperref[\detokenize{compute_smd_damage::doc}]{\sphinxcrossref{\DUrole{doc}{smd/damage}}}} \sphinxhyphen{} damage status of SPH particles in Smooth Mach Dynamics

\item {} 
\sphinxAtStartPar
{\hyperref[\detokenize{compute_smd_hourglass_error::doc}]{\sphinxcrossref{\DUrole{doc}{smd/hourglass/error}}}} \sphinxhyphen{} error associated with approximated relative separation in Smooth Mach Dynamics

\item {} 
\sphinxAtStartPar
{\hyperref[\detokenize{compute_smd_internal_energy::doc}]{\sphinxcrossref{\DUrole{doc}{smd/internal/energy}}}} \sphinxhyphen{} per\sphinxhyphen{}particle enthalpy in Smooth Mach Dynamics

\item {} 
\sphinxAtStartPar
{\hyperref[\detokenize{compute_smd_plastic_strain::doc}]{\sphinxcrossref{\DUrole{doc}{smd/plastic/strain}}}} \sphinxhyphen{} equivalent plastic strain per particle in Smooth Mach Dynamics

\item {} 
\sphinxAtStartPar
{\hyperref[\detokenize{compute_smd_plastic_strain_rate::doc}]{\sphinxcrossref{\DUrole{doc}{smd/plastic/strain/rate}}}} \sphinxhyphen{} time rate of the equivalent plastic strain in Smooth Mach Dynamics

\item {} 
\sphinxAtStartPar
{\hyperref[\detokenize{compute_smd_rho::doc}]{\sphinxcrossref{\DUrole{doc}{smd/rho}}}} \sphinxhyphen{} per\sphinxhyphen{}particle mass density in Smooth Mach Dynamics

\item {} 
\sphinxAtStartPar
{\hyperref[\detokenize{compute_smd_tlsph_defgrad::doc}]{\sphinxcrossref{\DUrole{doc}{smd/tlsph/defgrad}}}} \sphinxhyphen{} deformation gradient in Smooth Mach Dynamics

\item {} 
\sphinxAtStartPar
{\hyperref[\detokenize{compute_smd_tlsph_dt::doc}]{\sphinxcrossref{\DUrole{doc}{smd/tlsph/dt}}}} \sphinxhyphen{} CFL\sphinxhyphen{}stable time increment per particle in Smooth Mach Dynamics

\item {} 
\sphinxAtStartPar
{\hyperref[\detokenize{compute_smd_tlsph_num_neighs::doc}]{\sphinxcrossref{\DUrole{doc}{smd/tlsph/num/neighs}}}} \sphinxhyphen{} number of particles inside the smoothing kernel radius for Smooth Mach Dynamics

\item {} 
\sphinxAtStartPar
{\hyperref[\detokenize{compute_smd_tlsph_shape::doc}]{\sphinxcrossref{\DUrole{doc}{smd/tlsph/shape}}}} \sphinxhyphen{} current shape of the volume of a particle for Smooth Mach Dynamics

\item {} 
\sphinxAtStartPar
{\hyperref[\detokenize{compute_smd_tlsph_strain::doc}]{\sphinxcrossref{\DUrole{doc}{smd/tlsph/strain}}}} \sphinxhyphen{} Green\textendash{}Lagrange strain tensor for Smooth Mach Dynamics

\item {} 
\sphinxAtStartPar
{\hyperref[\detokenize{compute_smd_tlsph_strain_rate::doc}]{\sphinxcrossref{\DUrole{doc}{smd/tlsph/strain/rate}}}} \sphinxhyphen{} rate of strain for Smooth Mach Dynamics

\item {} 
\sphinxAtStartPar
{\hyperref[\detokenize{compute_smd_tlsph_stress::doc}]{\sphinxcrossref{\DUrole{doc}{smd/tlsph/stress}}}} \sphinxhyphen{} per\sphinxhyphen{}particle Cauchy stress tensor for SPH particles

\item {} 
\sphinxAtStartPar
{\hyperref[\detokenize{compute_smd_triangle_vertices::doc}]{\sphinxcrossref{\DUrole{doc}{smd/triangle/vertices}}}} \sphinxhyphen{} coordinates of vertices corresponding to the triangle elements of a mesh for Smooth Mach Dynamics

\item {} 
\sphinxAtStartPar
{\hyperref[\detokenize{compute_smd_ulsph_effm::doc}]{\sphinxcrossref{\DUrole{doc}{smd/ulsph/effm}}}} \sphinxhyphen{} per\sphinxhyphen{}particle effective shear modulus

\item {} 
\sphinxAtStartPar
{\hyperref[\detokenize{compute_smd_ulsph_num_neighs::doc}]{\sphinxcrossref{\DUrole{doc}{smd/ulsph/num/neighs}}}} \sphinxhyphen{} number of neighbor particles inside the smoothing kernel radius for Smooth Mach Dynamics

\item {} 
\sphinxAtStartPar
{\hyperref[\detokenize{compute_smd_ulsph_strain::doc}]{\sphinxcrossref{\DUrole{doc}{smd/ulsph/strain}}}} \sphinxhyphen{} logarithmic strain tensor for Smooth Mach Dynamics

\item {} 
\sphinxAtStartPar
{\hyperref[\detokenize{compute_smd_ulsph_strain_rate::doc}]{\sphinxcrossref{\DUrole{doc}{smd/ulsph/strain/rate}}}} \sphinxhyphen{} logarithmic strain rate for Smooth Mach Dynamics

\item {} 
\sphinxAtStartPar
{\hyperref[\detokenize{compute_smd_ulsph_stress::doc}]{\sphinxcrossref{\DUrole{doc}{smd/ulsph/stress}}}} \sphinxhyphen{} per\sphinxhyphen{}particle Cauchy stress tensor and von Mises equivalent stress in Smooth Mach Dynamics

\item {} 
\sphinxAtStartPar
{\hyperref[\detokenize{compute_smd_vol::doc}]{\sphinxcrossref{\DUrole{doc}{smd/vol}}}} \sphinxhyphen{} per\sphinxhyphen{}particle volumes and their sum in Smooth Mach Dynamics

\item {} 
\sphinxAtStartPar
{\hyperref[\detokenize{compute_sna_atom::doc}]{\sphinxcrossref{\DUrole{doc}{snap}}}} \sphinxhyphen{} gradients of SNAP energy and forces with respect to linear coefficients and related quantities for fitting SNAP potentials

\item {} 
\sphinxAtStartPar
{\hyperref[\detokenize{compute_sna_atom::doc}]{\sphinxcrossref{\DUrole{doc}{sna/atom}}}} \sphinxhyphen{} bispectrum components for each atom

\item {} 
\sphinxAtStartPar
{\hyperref[\detokenize{compute_sna_atom::doc}]{\sphinxcrossref{\DUrole{doc}{sna/grid}}}} \sphinxhyphen{} global array of bispectrum components on a regular grid

\item {} 
\sphinxAtStartPar
{\hyperref[\detokenize{compute_sna_atom::doc}]{\sphinxcrossref{\DUrole{doc}{sna/grid/local}}}} \sphinxhyphen{} local array of bispectrum components on a regular grid

\item {} 
\sphinxAtStartPar
{\hyperref[\detokenize{compute_sna_atom::doc}]{\sphinxcrossref{\DUrole{doc}{snad/atom}}}} \sphinxhyphen{} derivative of bispectrum components for each atom

\item {} 
\sphinxAtStartPar
{\hyperref[\detokenize{compute_sna_atom::doc}]{\sphinxcrossref{\DUrole{doc}{snav/atom}}}} \sphinxhyphen{} virial contribution from bispectrum components for each atom

\item {} 
\sphinxAtStartPar
{\hyperref[\detokenize{compute_sph_e_atom::doc}]{\sphinxcrossref{\DUrole{doc}{sph/e/atom}}}} \sphinxhyphen{} per\sphinxhyphen{}atom internal energy of Smooth\sphinxhyphen{}Particle Hydrodynamics atoms

\item {} 
\sphinxAtStartPar
{\hyperref[\detokenize{compute_sph_rho_atom::doc}]{\sphinxcrossref{\DUrole{doc}{sph/rho/atom}}}} \sphinxhyphen{} per\sphinxhyphen{}atom density of Smooth\sphinxhyphen{}Particle Hydrodynamics atoms

\item {} 
\sphinxAtStartPar
{\hyperref[\detokenize{compute_sph_t_atom::doc}]{\sphinxcrossref{\DUrole{doc}{sph/t/atom}}}} \sphinxhyphen{} per\sphinxhyphen{}atom internal temperature of Smooth\sphinxhyphen{}Particle Hydrodynamics atoms

\item {} 
\sphinxAtStartPar
{\hyperref[\detokenize{compute_spin::doc}]{\sphinxcrossref{\DUrole{doc}{spin}}}} \sphinxhyphen{} magnetic quantities for a system of atoms having spins

\item {} 
\sphinxAtStartPar
{\hyperref[\detokenize{compute_stress_atom::doc}]{\sphinxcrossref{\DUrole{doc}{stress/atom}}}} \sphinxhyphen{} stress tensor for each atom

\item {} 
\sphinxAtStartPar
{\hyperref[\detokenize{compute_stress_cartesian::doc}]{\sphinxcrossref{\DUrole{doc}{stress/cartesian}}}} \sphinxhyphen{} stress tensor in cartesian coordinates

\item {} 
\sphinxAtStartPar
{\hyperref[\detokenize{compute_stress_curvilinear::doc}]{\sphinxcrossref{\DUrole{doc}{stress/cylinder}}}} \sphinxhyphen{} stress tensor in cylindrical coordinates

\item {} 
\sphinxAtStartPar
{\hyperref[\detokenize{compute_stress_mop::doc}]{\sphinxcrossref{\DUrole{doc}{stress/mop}}}} \sphinxhyphen{} normal components of the local stress tensor using the method of planes

\item {} 
\sphinxAtStartPar
{\hyperref[\detokenize{compute_stress_mop::doc}]{\sphinxcrossref{\DUrole{doc}{stress/mop/profile}}}} \sphinxhyphen{} profile of the normal components of the local stress tensor using the method of planes

\item {} 
\sphinxAtStartPar
{\hyperref[\detokenize{compute_stress_curvilinear::doc}]{\sphinxcrossref{\DUrole{doc}{stress/spherical}}}} \sphinxhyphen{} stress tensor in spherical coordinates

\item {} 
\sphinxAtStartPar
{\hyperref[\detokenize{compute_tally::doc}]{\sphinxcrossref{\DUrole{doc}{stress/tally}}}} \sphinxhyphen{} stress between two groups of atoms via the tally callback mechanism

\item {} 
\sphinxAtStartPar
{\hyperref[\detokenize{compute_tdpd_cc_atom::doc}]{\sphinxcrossref{\DUrole{doc}{tdpd/cc/atom}}}} \sphinxhyphen{} per\sphinxhyphen{}atom chemical concentration of a specified species for each tDPD particle

\item {} 
\sphinxAtStartPar
{\hyperref[\detokenize{compute_temp::doc}]{\sphinxcrossref{\DUrole{doc}{temp}}}} \sphinxhyphen{} temperature of group of atoms

\item {} 
\sphinxAtStartPar
{\hyperref[\detokenize{compute_temp_asphere::doc}]{\sphinxcrossref{\DUrole{doc}{temp/asphere}}}} \sphinxhyphen{} temperature of aspherical particles

\item {} 
\sphinxAtStartPar
{\hyperref[\detokenize{compute_temp_body::doc}]{\sphinxcrossref{\DUrole{doc}{temp/body}}}} \sphinxhyphen{} temperature of body particles

\item {} 
\sphinxAtStartPar
{\hyperref[\detokenize{compute_temp_chunk::doc}]{\sphinxcrossref{\DUrole{doc}{temp/chunk}}}} \sphinxhyphen{} temperature of each chunk

\item {} 
\sphinxAtStartPar
{\hyperref[\detokenize{compute_temp_com::doc}]{\sphinxcrossref{\DUrole{doc}{temp/com}}}} \sphinxhyphen{} temperature after subtracting center\sphinxhyphen{}of\sphinxhyphen{}mass velocity

\item {} 
\sphinxAtStartPar
{\hyperref[\detokenize{compute_temp_cs::doc}]{\sphinxcrossref{\DUrole{doc}{temp/cs}}}} \sphinxhyphen{} temperature based on the center\sphinxhyphen{}of\sphinxhyphen{}mass velocity of atom pairs that are bonded to each other

\item {} 
\sphinxAtStartPar
{\hyperref[\detokenize{compute_temp_deform::doc}]{\sphinxcrossref{\DUrole{doc}{temp/deform}}}} \sphinxhyphen{} temperature excluding box deformation velocity

\item {} 
\sphinxAtStartPar
{\hyperref[\detokenize{compute_temp_deform_eff::doc}]{\sphinxcrossref{\DUrole{doc}{temp/deform/eff}}}} \sphinxhyphen{} temperature excluding box deformation velocity in the electron force field model

\item {} 
\sphinxAtStartPar
{\hyperref[\detokenize{compute_temp_drude::doc}]{\sphinxcrossref{\DUrole{doc}{temp/drude}}}} \sphinxhyphen{} temperature of Core\textendash{}Drude pairs

\item {} 
\sphinxAtStartPar
{\hyperref[\detokenize{compute_temp_eff::doc}]{\sphinxcrossref{\DUrole{doc}{temp/eff}}}} \sphinxhyphen{} temperature of a group of nuclei and electrons in the electron force field model

\item {} 
\sphinxAtStartPar
{\hyperref[\detokenize{compute_temp_partial::doc}]{\sphinxcrossref{\DUrole{doc}{temp/partial}}}} \sphinxhyphen{} temperature excluding one or more dimensions of velocity

\item {} 
\sphinxAtStartPar
{\hyperref[\detokenize{compute_temp_profile::doc}]{\sphinxcrossref{\DUrole{doc}{temp/profile}}}} \sphinxhyphen{} temperature excluding a binned velocity profile

\item {} 
\sphinxAtStartPar
{\hyperref[\detokenize{compute_temp_ramp::doc}]{\sphinxcrossref{\DUrole{doc}{temp/ramp}}}} \sphinxhyphen{} temperature excluding ramped velocity component

\item {} 
\sphinxAtStartPar
{\hyperref[\detokenize{compute_temp_region::doc}]{\sphinxcrossref{\DUrole{doc}{temp/region}}}} \sphinxhyphen{} temperature of a region of atoms

\item {} 
\sphinxAtStartPar
{\hyperref[\detokenize{compute_temp_region_eff::doc}]{\sphinxcrossref{\DUrole{doc}{temp/region/eff}}}} \sphinxhyphen{} temperature of a region of nuclei and electrons in the electron force field model

\item {} 
\sphinxAtStartPar
{\hyperref[\detokenize{compute_temp_rotate::doc}]{\sphinxcrossref{\DUrole{doc}{temp/rotate}}}} \sphinxhyphen{} temperature of a group of atoms after subtracting out their center\sphinxhyphen{}of\sphinxhyphen{}mass and angular velocities

\item {} 
\sphinxAtStartPar
{\hyperref[\detokenize{compute_temp_sphere::doc}]{\sphinxcrossref{\DUrole{doc}{temp/sphere}}}} \sphinxhyphen{} temperature of spherical particles

\item {} 
\sphinxAtStartPar
{\hyperref[\detokenize{compute_temp_uef::doc}]{\sphinxcrossref{\DUrole{doc}{temp/uef}}}} \sphinxhyphen{} kinetic energy tensor in the reference frame of an applied flow field

\item {} 
\sphinxAtStartPar
{\hyperref[\detokenize{compute_ti::doc}]{\sphinxcrossref{\DUrole{doc}{ti}}}} \sphinxhyphen{} thermodynamic integration free energy values

\item {} 
\sphinxAtStartPar
{\hyperref[\detokenize{compute_torque_chunk::doc}]{\sphinxcrossref{\DUrole{doc}{torque/chunk}}}} \sphinxhyphen{} torque applied on each chunk

\item {} 
\sphinxAtStartPar
{\hyperref[\detokenize{compute_vacf::doc}]{\sphinxcrossref{\DUrole{doc}{vacf}}}} \sphinxhyphen{} velocity auto\sphinxhyphen{}correlation function of group of atoms

\item {} 
\sphinxAtStartPar
{\hyperref[\detokenize{compute_vcm_chunk::doc}]{\sphinxcrossref{\DUrole{doc}{vcm/chunk}}}} \sphinxhyphen{} velocity of center\sphinxhyphen{}of\sphinxhyphen{}mass for each chunk

\item {} 
\sphinxAtStartPar
{\hyperref[\detokenize{compute_viscosity_cos::doc}]{\sphinxcrossref{\DUrole{doc}{viscosity/cos}}}} \sphinxhyphen{} velocity profile under cosine\sphinxhyphen{}shaped acceleration

\item {} 
\sphinxAtStartPar
{\hyperref[\detokenize{compute_voronoi_atom::doc}]{\sphinxcrossref{\DUrole{doc}{voronoi/atom}}}} \sphinxhyphen{} Voronoi volume and neighbors for each atom

\item {} 
\sphinxAtStartPar
{\hyperref[\detokenize{compute_xrd::doc}]{\sphinxcrossref{\DUrole{doc}{xrd}}}} \sphinxhyphen{} X\sphinxhyphen{}ray diffraction intensity on a mesh of reciprocal lattice nodes

\end{itemize}


\subsection{Restrictions}
\label{\detokenize{compute:restrictions}}
\sphinxAtStartPar
none


\subsection{Related commands}
\label{\detokenize{compute:related-commands}}
\sphinxAtStartPar
{\hyperref[\detokenize{uncompute::doc}]{\sphinxcrossref{\DUrole{doc}{uncompute}}}}, {\hyperref[\detokenize{compute_modify::doc}]{\sphinxcrossref{\DUrole{doc}{compute\_modify}}}},
{\hyperref[\detokenize{fix_ave_atom::doc}]{\sphinxcrossref{\DUrole{doc}{fix ave/atom}}}}, {\hyperref[\detokenize{fix_ave_time::doc}]{\sphinxcrossref{\DUrole{doc}{fix ave/time}}}},
{\hyperref[\detokenize{fix_ave_histo::doc}]{\sphinxcrossref{\DUrole{doc}{fix ave/histo}}}}


\subsection{Default}
\label{\detokenize{compute:default}}
\sphinxAtStartPar
none

\sphinxstepscope

\index{compute\_modify@\spxentry{compute\_modify}}\ignorespaces 

\section{compute\_modify command}
\label{\detokenize{compute_modify:compute-modify-command}}\label{\detokenize{compute_modify:index-0}}\label{\detokenize{compute_modify::doc}}

\subsection{Syntax}
\label{\detokenize{compute_modify:syntax}}
\begin{sphinxVerbatim}[commandchars=\\\{\}]
\PYG{k}{compute\PYGZus{}modify }\PYG{n+nv+nvIdentifier}{compute\PYGZhy{}ID}\PYG{+w}{ }\PYG{n}{keyword}\PYG{+w}{ }\PYG{n}{value}\PYG{+w}{ }\PYG{n}{...}
\end{sphinxVerbatim}
\begin{itemize}
\item {} 
\sphinxAtStartPar
compute\sphinxhyphen{}ID = ID of the compute to modify

\item {} 
\sphinxAtStartPar
one or more keyword/value pairs may be listed

\item {} 
\sphinxAtStartPar
keyword = \sphinxstyleemphasis{extra/dof} or \sphinxstyleemphasis{dynamic/dof}
\begin{sphinxalltt}
\sphinxstyleemphasis{extra/dof} value = N
  N = \# of extra degrees of freedom to subtract
\sphinxstyleemphasis{dynamic/dof} value = \sphinxstyleemphasis{yes} or \sphinxstyleemphasis{no}
  yes/no = do or do not re\sphinxhyphen{}compute the number of degrees of freedom (DOF) contributing to the temperature
\end{sphinxalltt}

\end{itemize}


\subsection{Examples}
\label{\detokenize{compute_modify:examples}}
\begin{sphinxVerbatim}[commandchars=\\\{\}]
\PYG{k}{compute\PYGZus{}modify }\PYG{n+nv+nvIdentifier}{myTemp}\PYG{+w}{ }\PYG{n}{extra}\PYG{o}{/}\PYG{n}{dof}\PYG{+w}{ }\PYG{l+m}{0}
\PYG{k}{compute\PYGZus{}modify }\PYG{n+nv+nvIdentifier}{newtemp}\PYG{+w}{ }\PYG{n}{dynamic}\PYG{o}{/}\PYG{n}{dof}\PYG{+w}{ }\PYG{n}{yes}\PYG{+w}{ }\PYG{n}{extra}\PYG{o}{/}\PYG{n}{dof}\PYG{+w}{ }\PYG{l+m}{600}
\end{sphinxVerbatim}


\subsection{Description}
\label{\detokenize{compute_modify:description}}
\sphinxAtStartPar
Modify one or more parameters of a previously defined compute.  Not
all compute styles support all parameters.

\sphinxAtStartPar
The \sphinxstyleemphasis{extra/dof} keyword refers to how many degrees of freedom are
subtracted (typically from \(3N\)) as a normalizing factor in a
temperature computation.  Only computes that compute a temperature use
this option.  The default is 2 or 3 for {\hyperref[\detokenize{dimension::doc}]{\sphinxcrossref{\DUrole{doc}{2d or 3d systems}}}} which is a correction factor for an ensemble of velocities
with zero total linear momentum. For compute temp/partial, if one or
more velocity components are excluded, the value used for \sphinxstyleemphasis{extra/dof} is
scaled accordingly. You can use a negative number for the \sphinxstyleemphasis{extra/dof}
parameter if you need to add degrees\sphinxhyphen{}of\sphinxhyphen{}freedom.  See the {\hyperref[\detokenize{compute_temp_asphere::doc}]{\sphinxcrossref{\DUrole{doc}{compute
temp/asphere}}}} command for an example.

\sphinxAtStartPar
The \sphinxstyleemphasis{dynamic/dof} keyword determines whether the number
of atoms \(N\) in the compute group and their associated degrees
of freedom (DOF) are re\sphinxhyphen{}computed each time a temperature is computed.
Only compute styles that calculate a temperature use this option.  By
default, \(N\) and their DOF are assumed to be constant.  If you
are adding atoms or molecules to the system (see the {\hyperref[\detokenize{fix_pour::doc}]{\sphinxcrossref{\DUrole{doc}{fix pour}}}}, {\hyperref[\detokenize{fix_deposit::doc}]{\sphinxcrossref{\DUrole{doc}{fix deposit}}}}, and {\hyperref[\detokenize{fix_gcmc::doc}]{\sphinxcrossref{\DUrole{doc}{fix gcmc}}}} commands) or expect atoms or molecules to be lost
(e.g. due to exiting the simulation box or via {\hyperref[\detokenize{fix_evaporate::doc}]{\sphinxcrossref{\DUrole{doc}{fix evaporate}}}}), then this option should be used to ensure the
temperature is correctly normalized.


\subsection{Restrictions}
\label{\detokenize{compute_modify:restrictions}}\begin{quote}

\sphinxAtStartPar
none
\end{quote}


\subsection{Related commands}
\label{\detokenize{compute_modify:related-commands}}
\sphinxAtStartPar
{\hyperref[\detokenize{compute::doc}]{\sphinxcrossref{\DUrole{doc}{compute}}}}


\subsection{Default}
\label{\detokenize{compute_modify:default}}
\sphinxAtStartPar
The option defaults are extra/dof = 2 or 3 for 2d or 3d systems,
respectively, and dynamic/dof = \sphinxstyleemphasis{no}.

\sphinxstepscope

\index{create\_atoms@\spxentry{create\_atoms}}\ignorespaces 

\section{create\_atoms command}
\label{\detokenize{create_atoms:create-atoms-command}}\label{\detokenize{create_atoms:index-0}}\label{\detokenize{create_atoms::doc}}

\subsection{Syntax}
\label{\detokenize{create_atoms:syntax}}
\begin{sphinxVerbatim}[commandchars=\\\{\}]
\PYG{k}{create\PYGZus{}atoms}\PYG{+w}{ }\PYG{n}{type}\PYG{+w}{ }\PYG{n}{style}\PYG{+w}{ }\PYG{n}{args}\PYG{+w}{ }\PYG{n}{keyword}\PYG{+w}{ }\PYG{n}{values}\PYG{+w}{ }\PYG{n}{...}
\end{sphinxVerbatim}
\begin{itemize}
\item {} 
\sphinxAtStartPar
type = atom type (1\sphinxhyphen{}Ntypes or type label) of atoms to create (offset for molecule creation)

\item {} 
\sphinxAtStartPar
style = \sphinxstyleemphasis{box} or \sphinxstyleemphasis{region} or \sphinxstyleemphasis{single} or \sphinxstyleemphasis{mesh} or \sphinxstyleemphasis{random}
\begin{sphinxalltt}
\sphinxstyleemphasis{box} args = none
\sphinxstyleemphasis{region} args = region\sphinxhyphen{}ID
  region\sphinxhyphen{}ID = particles will only be created if contained in the region
\sphinxstyleemphasis{single} args = x y z
  x,y,z = coordinates of a single particle (distance units)
\sphinxstyleemphasis{mesh} args = STL\sphinxhyphen{}file
  STL\sphinxhyphen{}file = file with triangle mesh in STL format
\sphinxstyleemphasis{random} args = N seed region\sphinxhyphen{}ID
  N = number of particles to create
  seed = random \# seed (positive integer)
  region\sphinxhyphen{}ID = create atoms within this region, use NULL for entire simulation box
\end{sphinxalltt}

\item {} 
\sphinxAtStartPar
zero or more keyword/value pairs may be appended

\item {} 
\sphinxAtStartPar
keyword = \sphinxstyleemphasis{mol} or \sphinxstyleemphasis{basis} or \sphinxstyleemphasis{ratio} or \sphinxstyleemphasis{subset} or \sphinxstyleemphasis{remap} or \sphinxstyleemphasis{var} or \sphinxstyleemphasis{set} or \sphinxstyleemphasis{radscale} or \sphinxstyleemphasis{meshmode} or \sphinxstyleemphasis{rotate} or \sphinxstyleemphasis{overlap} or \sphinxstyleemphasis{maxtry} or \sphinxstyleemphasis{units}
\begin{sphinxalltt}
\sphinxstyleemphasis{mol} values = template\sphinxhyphen{}ID seed
  template\sphinxhyphen{}ID = ID of molecule template specified in a separate {\hyperref[\detokenize{molecule::doc}]{\sphinxcrossref{\DUrole{doc}{molecule}}}} command
  seed = random \# seed (positive integer)
\sphinxstyleemphasis{basis} values = M itype
  M = which basis atom
  itype = atom type (1\sphinxhyphen{}Ntypes or type label) to assign to this basis atom
\sphinxstyleemphasis{ratio} values = frac seed
  frac = fraction of lattice sites (0 to 1) to populate randomly
  seed = random \# seed (positive integer)
\sphinxstyleemphasis{subset} values = Nsubset seed
  Nsubset = \# of lattice sites to populate randomly
  seed = random \# seed (positive integer)
\sphinxstyleemphasis{remap} value = \sphinxstyleemphasis{yes} or \sphinxstyleemphasis{no}
\sphinxstyleemphasis{var} value = name = variable name to evaluate for test of atom creation
\sphinxstyleemphasis{set} values = dim name
  dim = \sphinxstyleemphasis{x} or \sphinxstyleemphasis{y} or \sphinxstyleemphasis{z}
  name = name of variable to set with x, y, or z atom position
\sphinxstyleemphasis{radscale} value = factor
  factor = scale factor for setting atom radius
\sphinxstyleemphasis{meshmode} values = mode arg
  mode = \sphinxstyleemphasis{bisect} or \sphinxstyleemphasis{qrand}
  \sphinxstyleemphasis{bisect} arg = radthresh
    radthresh = threshold value for \sphinxstyleemphasis{mesh} to determine when to split triangles (distance units)
  \sphinxstyleemphasis{qrand} arg = density
    density = minimum number density for atoms place on \sphinxstyleemphasis{mesh} triangles (inverse distance squared units)
\sphinxstyleemphasis{rotate} values = theta Rx Ry Rz
  theta = rotation angle for single molecule (degrees)
  Rx,Ry,Rz = rotation vector for single molecule
\sphinxstyleemphasis{overlap} value = Doverlap
  Doverlap = only insert if at least this distance from all existing atoms
\sphinxstyleemphasis{maxtry} value = Ntry
  Ntry = number of attempts to insert a particle before failure
\sphinxstyleemphasis{units} value = \sphinxstyleemphasis{lattice} or \sphinxstyleemphasis{box}
  \sphinxstyleemphasis{lattice} = the geometry is defined in lattice units
  \sphinxstyleemphasis{box} = the geometry is defined in simulation box units
\end{sphinxalltt}

\end{itemize}


\subsection{Examples}
\label{\detokenize{create_atoms:examples}}
\begin{sphinxVerbatim}[commandchars=\\\{\}]
\PYG{k}{create\PYGZus{}atoms}\PYG{+w}{ }\PYG{l+m}{1}\PYG{+w}{ }\PYG{n}{box}

\PYG{k}{labelmap}\PYG{+w}{ }\PYG{n}{atom}\PYG{+w}{ }\PYG{l+m}{1}\PYG{+w}{ }\PYG{n}{Pt}
\PYG{k}{create\PYGZus{}atoms}\PYG{+w}{ }\PYG{n}{Pt}\PYG{+w}{ }\PYG{n}{box}

\PYG{k}{labelmap}\PYG{+w}{ }\PYG{n}{atom}\PYG{+w}{ }\PYG{l+m}{1}\PYG{+w}{ }\PYG{n}{C}\PYG{+w}{ }\PYG{l+m}{2}\PYG{+w}{ }\PYG{n}{Si}
\PYG{k}{create\PYGZus{}atoms}\PYG{+w}{ }\PYG{n}{C}\PYG{+w}{ }\PYG{k}{region }\PYG{n+nv+nvIdentifier}{regsphere}\PYG{+w}{ }\PYG{n}{basis}\PYG{+w}{ }\PYG{n}{Si}\PYG{+w}{ }\PYG{n}{C}

\PYG{k}{create\PYGZus{}atoms}\PYG{+w}{ }\PYG{l+m}{3}\PYG{+w}{ }\PYG{k}{region }\PYG{n+nv+nvIdentifier}{regsphere}\PYG{+w}{ }\PYG{n}{basis}\PYG{+w}{ }\PYG{l+m}{2}\PYG{+w}{ }\PYG{l+m}{3}
\PYG{k}{create\PYGZus{}atoms}\PYG{+w}{ }\PYG{l+m}{3}\PYG{+w}{ }\PYG{k}{region }\PYG{n+nv+nvIdentifier}{regsphere}\PYG{+w}{ }\PYG{n}{basis}\PYG{+w}{ }\PYG{l+m}{2}\PYG{+w}{ }\PYG{l+m}{3}\PYG{+w}{ }\PYG{n}{ratio}\PYG{+w}{ }\PYG{l+m}{0.5}\PYG{+w}{ }\PYG{l+m}{74637}
\PYG{k}{create\PYGZus{}atoms}\PYG{+w}{ }\PYG{l+m}{3}\PYG{+w}{ }\PYG{n}{single}\PYG{+w}{ }\PYG{l+m}{0}\PYG{+w}{ }\PYG{l+m}{0}\PYG{+w}{ }\PYG{l+m}{5}
\PYG{k}{create\PYGZus{}atoms}\PYG{+w}{ }\PYG{l+m}{1}\PYG{+w}{ }\PYG{n}{box}\PYG{+w}{ }\PYG{n}{var}\PYG{+w}{ }\PYG{n}{v}\PYG{+w}{ }\PYG{n}{set}\PYG{+w}{ }\PYG{n}{x}\PYG{+w}{ }\PYG{n}{xpos}\PYG{+w}{ }\PYG{n}{set}\PYG{+w}{ }\PYG{n}{y}\PYG{+w}{ }\PYG{n}{ypos}
\PYG{k}{create\PYGZus{}atoms}\PYG{+w}{ }\PYG{l+m}{2}\PYG{+w}{ }\PYG{n}{random}\PYG{+w}{ }\PYG{l+m}{50}\PYG{+w}{ }\PYG{l+m}{12345}\PYG{+w}{ }\PYG{n}{NULL}\PYG{+w}{ }\PYG{n}{overlap}\PYG{+w}{ }\PYG{l+m}{2.0}\PYG{+w}{ }\PYG{n}{maxtry}\PYG{+w}{ }\PYG{l+m}{50}
\PYG{k}{create\PYGZus{}atoms}\PYG{+w}{ }\PYG{l+m}{1}\PYG{+w}{ }\PYG{n}{mesh}\PYG{+w}{ }\PYG{n}{open\PYGZus{}box.stl}\PYG{+w}{ }\PYG{n}{meshmode}\PYG{+w}{ }\PYG{n}{qrand}\PYG{+w}{ }\PYG{l+m}{0.1}\PYG{+w}{ }\PYG{n}{units}\PYG{+w}{ }\PYG{n}{box}
\PYG{k}{create\PYGZus{}atoms}\PYG{+w}{ }\PYG{l+m}{1}\PYG{+w}{ }\PYG{n}{mesh}\PYG{+w}{ }\PYG{n}{funnel.stl}\PYG{+w}{ }\PYG{n}{meshmode}\PYG{+w}{ }\PYG{n}{bisect}\PYG{+w}{ }\PYG{l+m}{4.0}\PYG{+w}{ }\PYG{n}{units}\PYG{+w}{ }\PYG{n}{box}\PYG{+w}{ }\PYG{n}{radscale}\PYG{+w}{ }\PYG{l+m}{0.9}
\end{sphinxVerbatim}


\subsection{Description}
\label{\detokenize{create_atoms:description}}
\sphinxAtStartPar
This command creates atoms (or molecules) within the simulation box,
either on a lattice, or at random points, or on a surface defined by a
triangulated mesh.  Or it creates a single atom (or molecule) at a
specified point.  It is an alternative to reading in atom coordinates
explicitly via a {\hyperref[\detokenize{read_data::doc}]{\sphinxcrossref{\DUrole{doc}{read\_data}}}} or {\hyperref[\detokenize{read_restart::doc}]{\sphinxcrossref{\DUrole{doc}{read\_restart}}}} command.

\sphinxAtStartPar
To use this command a simulation box must already exist, which is
typically created via the {\hyperref[\detokenize{create_box::doc}]{\sphinxcrossref{\DUrole{doc}{create\_box}}}} command.
Before using this command, a lattice must typically also be defined
using the {\hyperref[\detokenize{lattice::doc}]{\sphinxcrossref{\DUrole{doc}{lattice}}}} command, unless you specify the
\sphinxstyleemphasis{single} or \sphinxstyleemphasis{mesh} style with units = box or the \sphinxstyleemphasis{random} style.  To
create atoms on a lattice for general triclinic boxes, see the
discussion below.

\sphinxAtStartPar
For the remainder of this doc page, a created atom or molecule is
referred to as a “particle”.

\sphinxAtStartPar
If created particles are individual atoms, they are assigned the
specified atom \sphinxstyleemphasis{type}, though this can be altered via the \sphinxstyleemphasis{basis}
keyword as discussed below.  If molecules are being created, the type
of each atom in the created molecule is specified in a specified file
read by the {\hyperref[\detokenize{molecule::doc}]{\sphinxcrossref{\DUrole{doc}{molecule}}}} command, and those values are
added to the specified atom \sphinxstyleemphasis{type} (e.g., if \sphinxstyleemphasis{type} = 2 and the file
specifies atom types 1, 2, and 3, then each created molecule will have
atom types 3, 4, and 5).

\begin{sphinxadmonition}{note}{Note:}
\sphinxAtStartPar
You cannot use this command to create atoms that are outside the
simulation box; they will just be ignored by LAMMPS.  This is true
even if you are using shrink\sphinxhyphen{}wrapped box boundaries, as specified
by the {\hyperref[\detokenize{boundary::doc}]{\sphinxcrossref{\DUrole{doc}{boundary}}}} command.  However, you can first
use the {\hyperref[\detokenize{change_box::doc}]{\sphinxcrossref{\DUrole{doc}{change\_box}}}} command to temporarily
expand the box, then add atoms via create\_atoms, then finally use
change\_box command again if needed to re\sphinxhyphen{}shrink\sphinxhyphen{}wrap the new atoms.
See the {\hyperref[\detokenize{change_box::doc}]{\sphinxcrossref{\DUrole{doc}{change\_box}}}} doc page for an example of
how to do this, using the create\_atoms \sphinxstyleemphasis{single} style to insert a
new atom outside the current simulation box.
\end{sphinxadmonition}


\bigskip\hrule\bigskip


\sphinxAtStartPar
For the \sphinxstyleemphasis{box} style, the create\_atoms command fills the entire
simulation box with particles on the lattice.  If your simulation box
is periodic, you should ensure its size is a multiple of the lattice
spacings, to avoid unwanted atom overlaps at the box boundaries.  If
your box is periodic and a multiple of the lattice spacing in a
particular dimension, LAMMPS is careful to put exactly one particle at
the boundary (on either side of the box), not zero or two.

\sphinxAtStartPar
For the \sphinxstyleemphasis{region} style, a geometric volume is filled with particles on
the lattice.  This volume is what is both inside the simulation box
and also consistent with the region volume.  See the {\hyperref[\detokenize{region::doc}]{\sphinxcrossref{\DUrole{doc}{region}}}} command for details.  Note that a region can be specified so
that its “volume” is either inside or outside its geometric boundary.
Also note that if a region is the same size as a periodic simulation
box (in some dimension), LAMMPS does NOT implement the same logic
described above for the \sphinxstyleemphasis{box} style, to ensure exactly one particle at
periodic boundaries.  If this is desired, you should either use the
\sphinxstyleemphasis{box} style, or tweak the region size to get precisely the particles
you want.


\bigskip\hrule\bigskip


\sphinxAtStartPar
If the simulation box is formulated as a general triclinic box defined
by arbitrary edge vectors \sphinxstylestrong{A}, \sphinxstylestrong{B}, \sphinxstylestrong{C}, then the \sphinxstyleemphasis{box} and
\sphinxstyleemphasis{region} styles will create atoms on a lattice commensurate with those
edge vectors.  See the {\hyperref[\detokenize{Howto_triclinic::doc}]{\sphinxcrossref{\DUrole{doc}{Howto\_triclinic}}}} doc
page for a detailed explanation of orthogonal, restricted triclinic,
and general triclinic simulation boxes.  As with the {\hyperref[\detokenize{create_box::doc}]{\sphinxcrossref{\DUrole{doc}{create\_box}}}} command, the {\hyperref[\detokenize{lattice::doc}]{\sphinxcrossref{\DUrole{doc}{lattice}}}} command used by
this command must be of style \sphinxstyleemphasis{custom} and use its \sphinxstyleemphasis{triclinic/general}
option.  The \sphinxstyleemphasis{a1, *a2}, \sphinxstyleemphasis{a3} settings of the {\hyperref[\detokenize{lattice::doc}]{\sphinxcrossref{\DUrole{doc}{lattice}}}}
command define the edge vectors of a unit cell of the general
triclinic lattice. The {\hyperref[\detokenize{create_box::doc}]{\sphinxcrossref{\DUrole{doc}{create\_box}}}} command creates
a simulation box which replicates that unit cell along each of the
\sphinxstylestrong{A}, \sphinxstylestrong{B}, \sphinxstylestrong{C} edge vectors.

\begin{sphinxadmonition}{note}{Note:}
\sphinxAtStartPar
LAMMPS allows specification of general triclinic simulation boxes
as a convenience for users who may be converting data from
solid\sphinxhyphen{}state crystallographic representations or from DFT codes for
input to LAMMPS.  However, as explained on the
{\hyperref[\detokenize{Howto_triclinic::doc}]{\sphinxcrossref{\DUrole{doc}{Howto\_triclinic}}}} doc page, internally,
LAMMPS only uses restricted triclinic simulation boxes.  This means
the box created by the {\hyperref[\detokenize{create_box::doc}]{\sphinxcrossref{\DUrole{doc}{create\_box}}}} command as
well as the atoms created by this command with their per\sphinxhyphen{}atom
information (e.g. coordinates, velocities) are converted (rotated)
from general to restricted triclinic form when the two commands are
invoked.  The {\hyperref[\detokenize{Howto_triclinic::doc}]{\sphinxcrossref{\DUrole{doc}{Howto\_triclinic}}}} doc page also
discusses other LAMMPS commands which can input/output general
triclinic representations of the simulation box and per\sphinxhyphen{}atom data.
\end{sphinxadmonition}

\sphinxAtStartPar
The \sphinxstyleemphasis{box} style will fill the entire general triclinic box with
particles on the lattice, as explained above.

\begin{sphinxadmonition}{note}{Note:}
\sphinxAtStartPar
The \sphinxstyleemphasis{region} style also operates as explained above, but the check
for particles inside the region is performed \sphinxstyleemphasis{after} the particle
coordinates have been converted to the restricted triclinic box.
This means the region must also be defined with respect to the
restricted triclinic box, not the general triclinic box.
\end{sphinxadmonition}

\sphinxAtStartPar
If the simulation box is general triclinic, the \sphinxstyleemphasis{single}, \sphinxstyleemphasis{random},
and \sphinxstyleemphasis{mesh} styles described next operate on the box \sphinxstyleemphasis{after} it has
been converted to restricted triclinic.  So all the settings for those
styles should be made in that context.


\bigskip\hrule\bigskip


\sphinxAtStartPar
For the \sphinxstyleemphasis{single} style, a single particle is added to the system at
the specified coordinates.  This can be useful for debugging purposes
or to create a tiny system with a handful of particles at specified
positions.  For a 2d simulation the specified z coordinate must be
0.0.

\sphinxAtStartPar
\DUrole{versionmodified,changed}{Changed in version 2Jun2022.}

\sphinxAtStartPar
The \sphinxstyleemphasis{porosity} style has been renamed to \sphinxstyleemphasis{random} with added functionality.

\sphinxAtStartPar
For the \sphinxstyleemphasis{random} style, \sphinxstyleemphasis{N} particles are added to the system at
randomly generated coordinates, which can be useful for generating an
amorphous system.  For 2d simulations, the z coordinates of all added
atoms will be 0.0.

\sphinxAtStartPar
The particles are created one by one using the specified random number
\sphinxstyleemphasis{seed}, resulting in the same set of particle coordinates, independent
of how many processors are being used in the simulation.  Unless the
\sphinxstyleemphasis{overlap} keyword is specified, particles created by the \sphinxstyleemphasis{random}
style will typically be highly overlapped.  Various additional
criteria can be used to accept or reject a random particle insertion;
see the keyword discussion below.  Multiple attempts per particle are
made (see the \sphinxstyleemphasis{maxtry} keyword) until the insertion is either
successful or fails.  If this command fails to add all requested \sphinxstyleemphasis{N}
particles, a warning will be output.

\sphinxAtStartPar
If the \sphinxstyleemphasis{region\sphinxhyphen{}ID} argument is specified as NULL, then the randomly
created particles will be anywhere in the simulation box.  If a
\sphinxstyleemphasis{region\sphinxhyphen{}ID} is specified, a geometric volume is filled that is both
inside the simulation box and is also consistent with the region
volume.  See the {\hyperref[\detokenize{region::doc}]{\sphinxcrossref{\DUrole{doc}{region}}}} command for details.  Note
that a region can be specified so that its “volume” is either inside
or outside its geometric boundary.

\sphinxAtStartPar
Note that the create\_atoms command adds particles to those that
already exist.  This means it can be used to add particles to a system
previously read in from a data or restart file.  Or the create\_atoms
command can be used multiple times, to add multiple sets of particles
to the simulation.  For example, grain boundaries can be created, by
interleaving the create\_atoms command with {\hyperref[\detokenize{lattice::doc}]{\sphinxcrossref{\DUrole{doc}{lattice}}}}
commands specifying different orientations.

\sphinxAtStartPar
When this command is used, care should be taken to ensure the
resulting system does not contain particles that are highly
overlapped.  Such overlaps will cause many interatomic potentials to
compute huge energies and forces, leading to bad dynamics.  There are
several strategies to avoid this problem:
\begin{itemize}
\item {} 
\sphinxAtStartPar
Use the {\hyperref[\detokenize{delete_atoms::doc}]{\sphinxcrossref{\DUrole{doc}{delete\_atoms overlap}}}} command after
create\_atoms.  For example, this strategy can be used to overlay and
surround a large protein molecule with a volume of water molecules,
then delete water molecules that overlap with the protein atoms.

\item {} 
\sphinxAtStartPar
For the \sphinxstyleemphasis{random} style, use the optional \sphinxstyleemphasis{overlap} keyword to avoid
overlaps when each new particle is created.

\item {} 
\sphinxAtStartPar
Before running dynamics on an overlapped system, perform an
{\hyperref[\detokenize{minimize::doc}]{\sphinxcrossref{\DUrole{doc}{energy minimization}}}}.  Or run initial dynamics with
{\hyperref[\detokenize{pair_soft::doc}]{\sphinxcrossref{\DUrole{doc}{pair\_style soft}}}} or with {\hyperref[\detokenize{fix_nve_limit::doc}]{\sphinxcrossref{\DUrole{doc}{fix nve/limit}}}} to un\sphinxhyphen{}overlap the particles, before running normal
dynamics.

\end{itemize}


\begin{wrapfigure}{r}{0.330\linewidth}
\centering
\sphinxhref{\_images/marble\_race.jpg}{\sphinxincludegraphics{{marble_race}.jpg}}
\end{wrapfigure}

\mbox{}\par\vskip-\dimexpr\baselineskip+\parskip\relax

\sphinxAtStartPar
\DUrole{versionmodified,added}{Added in version 2Jun2022.}

\sphinxAtStartPar
For the \sphinxstyleemphasis{mesh} style, a file with a triangle mesh in \sphinxhref{https://en.wikipedia.org/wiki/STL\_(file\_format)}{STL format} is read and one or
more particles are placed into the area of each triangle.  The reader
supports both ASCII and binary files conforming to the format on the
Wikipedia page.  Binary STL files (e.g. as frequently offered for
3d\sphinxhyphen{}printing) can also be first converted to ASCII for editing with the
{\hyperref[\detokenize{Tools:stlconvert}]{\sphinxcrossref{\DUrole{std,std-ref}{stl\_bin2txt tool}}}}.  The use of the \sphinxstyleemphasis{units box} option
is required. There are two algorithms available for placing atoms:
\sphinxstyleemphasis{bisect} and \sphinxstyleemphasis{qrand}. They can be selected via the \sphinxstyleemphasis{meshmode} option;
\sphinxstyleemphasis{bisect} is the default.  If the atom style allows it, the radius will
be set to a value depending on the algorithm and the value of the
\sphinxstyleemphasis{radscale} parameter (see below), and the atoms created from the mesh
are assigned a new molecule ID.

\sphinxAtStartPar
In \sphinxstyleemphasis{bisect} mode a particle is created at the center of each triangle
unless the average distance of the triangle vertices from its center is
larger than the \sphinxstyleemphasis{radthresh} value (default is lattice spacing in
x\sphinxhyphen{}direction).  In case the average distance is over the threshold, the
triangle is recursively split into two halves along the the longest side
until the threshold is reached. There will be at least one sphere per
triangle. The value of \sphinxstyleemphasis{radthresh} is set as an argument to \sphinxstyleemphasis{meshmode
bisect}.  The average distance of the vertices from the center is also
used to set the radius.

\sphinxAtStartPar
In \sphinxstyleemphasis{qrand} mode a quasi\sphinxhyphen{}random sequence is used to distribute particles
on mesh triangles using an approach by {\hyperref[\detokenize{create_atoms:roberts2019}]{\sphinxcrossref{\DUrole{std,std-ref}{(Roberts)}}}}.
Particles are added to the triangle until the minimum number density is
met or exceeded such that every triangle will have at least one
particle.  The minimum number density is set as an argument to the
\sphinxstyleemphasis{qrand} option.  The radius will be set so that the sum of the area of
the radius of the particles created in place of a triangle will be equal
to the area of that triangle.

\begin{sphinxadmonition}{note}{Note:}
\sphinxAtStartPar
The atom placement algorithms in the \sphinxstyleemphasis{mesh} style benefit from meshes
where triangles are close to equilateral.  It is therefore
recommended to pre\sphinxhyphen{}process STL files to optimize the mesh
accordingly.  There are multiple open source and commercial software
tools available with the capability to generate optimized meshes.
\end{sphinxadmonition}

\begin{sphinxadmonition}{note}{Note:}
\sphinxAtStartPar
In most cases the atoms created in \sphinxstyleemphasis{mesh} style will become an
immobile or rigid object that would not be time integrated or moved
by {\hyperref[\detokenize{fix_move::doc}]{\sphinxcrossref{\DUrole{doc}{fix move}}}} or {\hyperref[\detokenize{fix_rigid::doc}]{\sphinxcrossref{\DUrole{doc}{fix rigid}}}}.  For
computational efficiency \sphinxstyleemphasis{and} to avoid undesired contributions to
pressure and potential energy due to close contacts, it is usually
beneficial to exclude computing interactions between the created
particles using {\hyperref[\detokenize{neigh_modify::doc}]{\sphinxcrossref{\DUrole{doc}{neigh\_modify exclude}}}}.
\end{sphinxadmonition}


\bigskip\hrule\bigskip


\sphinxAtStartPar
Individual atoms are inserted by this command, unless the \sphinxstyleemphasis{mol}
keyword is used.  It specifies a \sphinxstyleemphasis{template\sphinxhyphen{}ID} previously defined
using the {\hyperref[\detokenize{molecule::doc}]{\sphinxcrossref{\DUrole{doc}{molecule}}}} command, which reads a file that
defines the molecule.  The coordinates, atom types, charges, etc, as
well as any bond/angle/etc and special neighbor information for the
molecule can be specified in the molecule file.  See the
{\hyperref[\detokenize{molecule::doc}]{\sphinxcrossref{\DUrole{doc}{molecule}}}} command for details.  The only settings
required to be in this file are the coordinates and types of atoms in
the molecule.

\begin{sphinxadmonition}{note}{Note:}
\sphinxAtStartPar
If you are using the \sphinxstyleemphasis{mol} keyword in combination with the
{\hyperref[\detokenize{atom_style::doc}]{\sphinxcrossref{\DUrole{doc}{atom style template}}}} command, they must use
the same molecule template\sphinxhyphen{}ID.
\end{sphinxadmonition}

\sphinxAtStartPar
Using a lattice to add molecules, e.g. via the \sphinxstyleemphasis{box} or \sphinxstyleemphasis{region} or
\sphinxstyleemphasis{single} styles, is exactly the same as adding atoms on lattice
points, except that entire molecules are added at each point, i.e. on
the point defined by each basis atom in the unit cell as it tiles the
simulation box or region.  This is done by placing the geometric
center of the molecule at the lattice point, and (by default) giving
the molecule a random orientation about the point.  The random \sphinxstyleemphasis{seed}
specified with the \sphinxstyleemphasis{mol} keyword is used for this operation, and the
random numbers generated by each processor are different.  This means
the coordinates of individual atoms (in the molecules) will be
different when running on different numbers of processors, unlike when
atoms are being created in parallel.

\sphinxAtStartPar
Note that with random rotations, it may be important to use a lattice
with a large enough spacing that adjacent molecules will not overlap,
regardless of their relative orientations.  See the description of the
\sphinxstyleemphasis{rotate} keyword below, which overrides the default random orientation
and inserts all molecules at a specified orientation.

\begin{sphinxadmonition}{note}{Note:}
\sphinxAtStartPar
If the {\hyperref[\detokenize{create_box::doc}]{\sphinxcrossref{\DUrole{doc}{create\_box}}}} command is used to create
the simulation box, followed by the create\_atoms command with its
\sphinxstyleemphasis{mol} option for adding molecules, then you typically need to use the
optional keywords allowed by the {\hyperref[\detokenize{create_box::doc}]{\sphinxcrossref{\DUrole{doc}{create\_box}}}} command
for extra bonds (angles,etc) or extra special neighbors.  This is
because by default, the {\hyperref[\detokenize{create_box::doc}]{\sphinxcrossref{\DUrole{doc}{create\_box}}}} command sets up a
non\sphinxhyphen{}molecular system that does not allow molecules to be added.
\end{sphinxadmonition}


\bigskip\hrule\bigskip


\sphinxAtStartPar
This is the meaning of the other optional keywords.

\sphinxAtStartPar
The \sphinxstyleemphasis{basis} keyword is only used when atoms (not molecules) are being
created.  It specifies an atom type that will be assigned to specific
basis atoms as they are created.  See the {\hyperref[\detokenize{lattice::doc}]{\sphinxcrossref{\DUrole{doc}{lattice}}}}
command for specifics on how basis atoms are defined for the unit cell
of the lattice.  By default, all created atoms are assigned the
argument \sphinxstyleemphasis{type} as their atom type.

\sphinxAtStartPar
The \sphinxstyleemphasis{ratio} and \sphinxstyleemphasis{subset} keywords can be used in conjunction with the
\sphinxstyleemphasis{box} or \sphinxstyleemphasis{region} styles to limit the total number of particles
inserted.  The lattice defines a set of \sphinxstyleemphasis{Nlatt} eligible sites for
inserting particles, which may be limited by the \sphinxstyleemphasis{region} style or the
\sphinxstyleemphasis{var} and \sphinxstyleemphasis{set} keywords.  For the \sphinxstyleemphasis{ratio} keyword, only the specified
fraction of them (\(0 \le f \le 1\)) will be assigned particles.
For the \sphinxstyleemphasis{subset} keyword only the specified \sphinxstyleemphasis{Nsubset} of them will be
assigned particles.  In both cases the assigned lattice sites are
chosen randomly.  An iterative algorithm is used that ensures the
correct number of particles are inserted, in a perfectly random
fashion.  Which lattice sites are selected will change with the number
of processors used.

\sphinxAtStartPar
The \sphinxstyleemphasis{remap} keyword only applies to the \sphinxstyleemphasis{single} style.  If it is set
to \sphinxstyleemphasis{yes}, then if the specified position is outside the simulation
box, it will mapped back into the box, assuming the relevant
dimensions are periodic.  If it is set to \sphinxstyleemphasis{no}, no remapping is done
and no particle is created if its position is outside the box.

\sphinxAtStartPar
The \sphinxstyleemphasis{var} and \sphinxstyleemphasis{set} keywords can be used together to provide a
criterion for accepting or rejecting the addition of an individual
atom, based on its coordinates.  They apply to all styles except
\sphinxstyleemphasis{single}.  The \sphinxstyleemphasis{name} specified for the \sphinxstyleemphasis{var} keyword is the name of
an {\hyperref[\detokenize{variable::doc}]{\sphinxcrossref{\DUrole{doc}{equal\sphinxhyphen{}style variable}}}} that should evaluate to a
zero or non\sphinxhyphen{}zero value based on one or two or three variables that
will store the \sphinxstyleemphasis{x}, \sphinxstyleemphasis{y}, or \sphinxstyleemphasis{z} coordinates of an atom (one variable per
coordinate).  If used, these other variables must be
{\hyperref[\detokenize{variable::doc}]{\sphinxcrossref{\DUrole{doc}{internal\sphinxhyphen{}style variables}}}} defined in the input
script; their initial numeric value can be anything.  They must be
internal\sphinxhyphen{}style variables, because this command resets their values
directly.  The \sphinxstyleemphasis{set} keyword is used to identify the names of these
other variables, one variable for the \sphinxstyleemphasis{x}\sphinxhyphen{}coordinate of a created atom,
one for \sphinxstyleemphasis{y}, and one for \sphinxstyleemphasis{z}.


\begin{wrapfigure}{r}{0.500\linewidth}
\centering
\sphinxhref{\_images/sinusoid.jpg}{\sphinxincludegraphics{{sinusoid}.jpg}}
\end{wrapfigure}

\mbox{}\par\vskip-\dimexpr\baselineskip+\parskip\relax

\sphinxAtStartPar
When an atom is created, its \((x,y,z)\) coordinates become the values for
any \sphinxstyleemphasis{set} variable that is defined.  The \sphinxstyleemphasis{var} variable is then
evaluated.  If the returned value is 0.0, the atom is not created.  If
it is non\sphinxhyphen{}zero, the atom is created.

\sphinxAtStartPar
As an example, these commands can be used in a 2d simulation, to
create a sinusoidal surface.  Note that the surface is “rough” due to
individual lattice points being “above” or “below” the mathematical
expression for the sinusoidal curve.  If a finer lattice were used,
the sinusoid would appear to be “smoother”.  Also note the use of the
“xlat” and “ylat” {\hyperref[\detokenize{thermo_style::doc}]{\sphinxcrossref{\DUrole{doc}{thermo\_style}}}} keywords, which
converts lattice spacings to distance.

\begin{sphinxVerbatim}[commandchars=\\\{\}]
\PYG{k}{dimension}\PYG{+w}{   }\PYG{l+m}{2}
\PYG{k}{variable    }\PYG{n+nv+nvIdentifier}{x}\PYG{+w}{ }\PYG{n}{equal}\PYG{+w}{ }\PYG{l+m}{100}
\PYG{k}{variable    }\PYG{n+nv+nvIdentifier}{y}\PYG{+w}{ }\PYG{n}{equal}\PYG{+w}{ }\PYG{l+m}{25}
\PYG{k}{lattice}\PYG{+w}{     }\PYG{n}{hex}\PYG{+w}{ }\PYG{l+m}{0.8442}
\PYG{k}{region      }\PYG{n+nv+nvIdentifier}{box}\PYG{+w}{ }\PYG{n}{block}\PYG{+w}{ }\PYG{l+m}{0}\PYG{+w}{ }\PYG{n+nv}{\PYGZdl{}x}\PYG{+w}{ }\PYG{l+m}{0}\PYG{+w}{ }\PYG{n+nv}{\PYGZdl{}y}\PYG{+w}{ }\PYG{o}{\PYGZhy{}}\PYG{l+m}{0.5}\PYG{+w}{ }\PYG{l+m}{0.5}
\PYG{k}{create\PYGZus{}box  }\PYG{n+nv+nvIdentifier}{1}\PYG{+w}{ }\PYG{n+nv+nvIdentifier}{box}

\PYG{k}{variable    }\PYG{n+nv+nvIdentifier}{xx}\PYG{+w}{ }\PYG{n}{internal}\PYG{+w}{ }\PYG{l+m}{0.0}
\PYG{k}{variable    }\PYG{n+nv+nvIdentifier}{yy}\PYG{+w}{ }\PYG{n}{internal}\PYG{+w}{ }\PYG{l+m}{0.0}
\PYG{k}{variable    }\PYG{n+nv+nvIdentifier}{v}\PYG{+w}{ }\PYG{n}{equal}\PYG{+w}{ }\PYG{l+s}{\PYGZdq{}}\PYG{l+s}{(0.2*v\PYGZus{}y*ylat * cos(v\PYGZus{}xx/xlat * 2.0*PI*4.0/v\PYGZus{}x) + 0.5*v\PYGZus{}y*ylat \PYGZhy{} v\PYGZus{}yy) \PYGZgt{} 0.0}\PYG{l+s}{\PYGZdq{}}
\PYG{k}{create\PYGZus{}atoms}\PYG{+w}{  }\PYG{l+m}{1}\PYG{+w}{ }\PYG{n}{box}\PYG{+w}{ }\PYG{n}{var}\PYG{+w}{ }\PYG{n}{v}\PYG{+w}{ }\PYG{n}{set}\PYG{+w}{ }\PYG{n}{x}\PYG{+w}{ }\PYG{n}{xx}\PYG{+w}{ }\PYG{n}{set}\PYG{+w}{ }\PYG{n}{y}\PYG{+w}{ }\PYG{n}{yy}
\PYG{k}{write\PYGZus{}dump  }\PYG{n+nv+nvIdentifier}{all}\PYG{+w}{ }\PYG{n}{atom}\PYG{+w}{ }\PYG{n}{sinusoid.lammpstrj}
\end{sphinxVerbatim}


\bigskip\hrule\bigskip


\sphinxAtStartPar
The \sphinxstyleemphasis{rotate} keyword allows specification of the orientation
at which molecules are inserted.  The axis of rotation is
determined by the rotation vector \((R_x,R_y,R_z)\) that goes through the
insertion point.  The specified \sphinxstyleemphasis{theta} determines the angle of
rotation around that axis.  Note that the direction of rotation for
the atoms around the rotation axis is consistent with the right\sphinxhyphen{}hand
rule: if your right\sphinxhyphen{}hand’s thumb points along \sphinxstyleemphasis{R}, then your fingers
wrap around the axis in the direction of rotation.

\sphinxAtStartPar
The \sphinxstyleemphasis{radscale} keyword only applies to the \sphinxstyleemphasis{mesh} style and adjusts the
radius of created particles (see above), provided this is supported by
the atom style.  Its value is a prefactor (must be \(>\) 0.0, default is
1.0) that is applied to the atom radius inferred from the size of the
individual triangles in the triangle mesh that the particle corresponds
to.

\sphinxAtStartPar
\DUrole{versionmodified,added}{Added in version 2Jun2022.}

\sphinxAtStartPar
The \sphinxstyleemphasis{overlap} keyword only applies to the \sphinxstyleemphasis{random} style.  It prevents
newly created particles from being created closer than the specified
\sphinxstyleemphasis{Doverlap} distance from any other particle.  If particles have finite
size (see {\hyperref[\detokenize{atom_style::doc}]{\sphinxcrossref{\DUrole{doc}{atom\_style sphere}}}} for example) \sphinxstyleemphasis{Doverlap}
should be specified large enough to include the particle size in the
non\sphinxhyphen{}overlapping criterion.  If molecules are being randomly inserted, then
an insertion is only accepted if each particle in the molecule meets the
overlap criterion with respect to other particles (not including particles
in the molecule itself).

\begin{sphinxadmonition}{note}{Note:}
\sphinxAtStartPar
Checking for overlaps is a costly \(\mathcal{O}(N(N+M))\) operation for
inserting \sphinxstyleemphasis{N} new particles into a system with \sphinxstyleemphasis{M} existing particles.
This is because distances to all \sphinxstyleemphasis{M} existing particles are computed for
each new particle that is added.  Thus the intended use of this
keyword is to add relatively small numbers of particles to systems
that remain at a relatively low density even after the new
particles are created.  Careful use of the \sphinxstyleemphasis{maxtry} keyword in
combination with \sphinxstyleemphasis{overlap} is recommended.  See the discussion
above about systems with overlapped particles for alternate
strategies that allow for overlapped insertions.
\end{sphinxadmonition}

\sphinxAtStartPar
\DUrole{versionmodified,added}{Added in version 2Jun2022.}

\sphinxAtStartPar
The \sphinxstyleemphasis{maxtry} keyword only applies to the \sphinxstyleemphasis{random} style.  It limits
the number of attempts to generate valid coordinates for a single new
particle that satisfy all requirements imposed by the \sphinxstyleemphasis{region}, \sphinxstyleemphasis{var},
and \sphinxstyleemphasis{overlap} keywords.  The default is 10 attempts per particle
before the loop over the requested \sphinxstyleemphasis{N} particles advances to the next
particle.  Note that if insertion success is unlikely (e.g., inserting
new particles into a dense system using the \sphinxstyleemphasis{overlap} keyword),
setting the \sphinxstyleemphasis{maxtry} keyword to a large value may result in this
command running for a long time.


\begin{wrapfigure}{r}{0.300\linewidth}
\centering
\sphinxhref{\_images/overlap.png}{\sphinxincludegraphics{{overlap}.png}}
\end{wrapfigure}

\mbox{}\par\vskip-\dimexpr\baselineskip+\parskip\relax

\sphinxAtStartPar
Here is an example for the \sphinxstyleemphasis{random} style using these commands

\begin{sphinxVerbatim}[commandchars=\\\{\}]
\PYG{k}{units}\PYG{+w}{         }\PYG{n}{lj}
\PYG{k}{dimension}\PYG{+w}{     }\PYG{l+m}{2}
\PYG{k}{region        }\PYG{n+nv+nvIdentifier}{box}\PYG{+w}{ }\PYG{n}{block}\PYG{+w}{ }\PYG{l+m}{0}\PYG{+w}{ }\PYG{l+m}{50}\PYG{+w}{ }\PYG{l+m}{0}\PYG{+w}{ }\PYG{l+m}{50}\PYG{+w}{ }\PYG{o}{\PYGZhy{}}\PYG{l+m}{0.5}\PYG{+w}{ }\PYG{l+m}{0.5}
\PYG{k}{create\PYGZus{}box    }\PYG{n+nv+nvIdentifier}{1}\PYG{+w}{ }\PYG{n+nv+nvIdentifier}{box}
\PYG{k}{create\PYGZus{}atoms}\PYG{+w}{  }\PYG{l+m}{1}\PYG{+w}{ }\PYG{n}{random}\PYG{+w}{ }\PYG{l+m}{2000}\PYG{+w}{ }\PYG{l+m}{13487}\PYG{+w}{ }\PYG{n}{NULL}\PYG{+w}{ }\PYG{n}{overlap}\PYG{+w}{ }\PYG{l+m}{1.0}\PYG{+w}{ }\PYG{n}{maxtry}\PYG{+w}{ }\PYG{l+m}{50}
\PYG{k}{pair\PYGZus{}style}\PYG{+w}{    }\PYG{n}{lj}\PYG{o}{/}\PYG{n}{cut}\PYG{+w}{ }\PYG{l+m}{2.5}
\PYG{k}{pair\PYGZus{}coeff}\PYG{+w}{    }\PYG{l+m}{1}\PYG{+w}{ }\PYG{l+m}{1}\PYG{+w}{ }\PYG{l+m}{1.0}\PYG{+w}{ }\PYG{l+m}{1.0}\PYG{+w}{ }\PYG{l+m}{2.5}
\end{sphinxVerbatim}

\sphinxAtStartPar
to produce a system as shown in the image with 1520 particles (out of
2000 requested) that are moderately dense and which have no overlaps
sufficient to prevent the LJ pair\_style from running properly (because
the overlap criterion is 1.0).  The create\_atoms command ran for 0.3 s
on a single CPU core.


\bigskip\hrule\bigskip


\sphinxAtStartPar
The \sphinxstyleemphasis{units} keyword determines the meaning of the distance units used
by parameters for various styles.  A \sphinxstyleemphasis{box} value selects standard
distance units as defined by the {\hyperref[\detokenize{units::doc}]{\sphinxcrossref{\DUrole{doc}{units}}}} command (e.g.,
\(\AA\) for units = \sphinxstyleemphasis{real} or \sphinxstyleemphasis{metal}.  A \sphinxstyleemphasis{lattice} value means
the distance units are in lattice spacings.  These are affected settings:
\begin{itemize}
\item {} 
\sphinxAtStartPar
for \sphinxstyleemphasis{single} style: coordinates of the particle created

\item {} 
\sphinxAtStartPar
for \sphinxstyleemphasis{random} style: overlap distance \sphinxstyleemphasis{Doverlap} by the \sphinxstyleemphasis{overlap} keyword

\item {} 
\sphinxAtStartPar
for \sphinxstyleemphasis{mesh} style: \sphinxstyleemphasis{bisect} threshold value for \sphinxstyleemphasis{meshmode} = \sphinxstyleemphasis{bisect}

\item {} 
\sphinxAtStartPar
for \sphinxstyleemphasis{mesh} style: \sphinxstyleemphasis{radthresh} value for \sphinxstyleemphasis{meshmode} = \sphinxstyleemphasis{bisect}

\item {} 
\sphinxAtStartPar
for \sphinxstyleemphasis{mesh} style: \sphinxstyleemphasis{density} value for \sphinxstyleemphasis{meshmode} = \sphinxstyleemphasis{qrand}

\end{itemize}

\sphinxAtStartPar
Since \sphinxstyleemphasis{density} represents an area (distance \textasciicircum{}2), the lattice spacing
factor is also squared.


\bigskip\hrule\bigskip


\sphinxAtStartPar
Atom IDs are assigned to created atoms in the following way.  The
collection of created atoms are assigned consecutive IDs that start
immediately following the largest atom ID existing before the
create\_atoms command was invoked.  This is done by the processor’s
communicating the number of atoms they each own, the first processor
numbering its atoms from \(1\) to \(N_1\), the second processor from
\(N_1+1\) to \(N_2\), and so on, where \(N_1\) is the number of
atoms owned by the first processor, \(N_2\) is the number owned by the
second processor, and so forth.  Thus, when the same simulation is performed on
different numbers of processors, there is no guarantee a particular created
atom will be assigned the same ID in both simulations.  If molecules are being
created, molecule IDs are assigned to created molecules in a similar fashion.

\sphinxAtStartPar
Aside from their ID, atom type, and \(xyz\) position, other properties of
created atoms are set to default values, depending on which quantities
are defined by the chosen {\hyperref[\detokenize{atom_style::doc}]{\sphinxcrossref{\DUrole{doc}{atom style}}}}.  See the
{\hyperref[\detokenize{atom_style::doc}]{\sphinxcrossref{\DUrole{doc}{atom style}}}} command for more details.  See the
{\hyperref[\detokenize{set::doc}]{\sphinxcrossref{\DUrole{doc}{set}}}} and {\hyperref[\detokenize{velocity::doc}]{\sphinxcrossref{\DUrole{doc}{velocity}}}} commands for info on
how to change these values.
\begin{itemize}
\item {} 
\sphinxAtStartPar
charge = 0.0

\item {} 
\sphinxAtStartPar
dipole moment magnitude = 0.0

\item {} 
\sphinxAtStartPar
diameter = 1.0

\item {} 
\sphinxAtStartPar
shape = 0.0 0.0 0.0

\item {} 
\sphinxAtStartPar
density = 1.0

\item {} 
\sphinxAtStartPar
volume = 1.0

\item {} 
\sphinxAtStartPar
velocity = 0.0 0.0 0.0

\item {} 
\sphinxAtStartPar
angular velocity = 0.0 0.0 0.0

\item {} 
\sphinxAtStartPar
angular momentum = 0.0 0.0 0.0

\item {} 
\sphinxAtStartPar
quaternion = (1,0,0,0)

\item {} 
\sphinxAtStartPar
bonds, angles, dihedrals, impropers = none

\end{itemize}

\sphinxAtStartPar
If molecules are being created, these defaults can be overridden by
values specified in the file read by the {\hyperref[\detokenize{molecule::doc}]{\sphinxcrossref{\DUrole{doc}{molecule}}}}
command. That is, the file typically defines bonds (angles, etc.) between
atoms in the molecule, and can optionally define charges on each atom.

\sphinxAtStartPar
Note that the \sphinxstyleemphasis{sphere} atom style sets the default particle diameter to
1.0 as well as the density.  This means the mass for the particle is not
1.0, but is \(\frac{\pi}{6} d^3 = 0.5236\), where \(d\) is the
diameter.  When using the \sphinxstyleemphasis{mesh} style, the particle diameter is adjusted from
the size of the individual triangles in the triangle mesh.

\sphinxAtStartPar
Note that the \sphinxstyleemphasis{ellipsoid} atom style sets the default particle shape
to (0.0 0.0 0.0) and the density to 1.0, which means it is a point
particle, not an ellipsoid, and has a mass of 1.0.

\sphinxAtStartPar
Note that the \sphinxstyleemphasis{peri} style sets the default volume and density to 1.0
and thus also set the mass for the particle to 1.0.

\sphinxAtStartPar
The {\hyperref[\detokenize{set::doc}]{\sphinxcrossref{\DUrole{doc}{set}}}} command can be used to override many of these
default settings.


\bigskip\hrule\bigskip



\subsection{Restrictions}
\label{\detokenize{create_atoms:restrictions}}
\sphinxAtStartPar
An {\hyperref[\detokenize{atom_style::doc}]{\sphinxcrossref{\DUrole{doc}{atom\_style}}}} must be previously defined to use this
command.

\sphinxAtStartPar
A rotation vector specified for a single molecule must be in
the z\sphinxhyphen{}direction for a 2d model.

\sphinxAtStartPar
For {\hyperref[\detokenize{molecule::doc}]{\sphinxcrossref{\DUrole{doc}{molecule templates}}}} that are created from multiple
files, i.e. contain multiple molecule \sphinxstyleemphasis{sets}, only the first set is
used.  To create multiple molecules the files currently need to be
merged and different molecule IDs assigned with a Molecules section.


\subsection{Related commands}
\label{\detokenize{create_atoms:related-commands}}
\sphinxAtStartPar
{\hyperref[\detokenize{lattice::doc}]{\sphinxcrossref{\DUrole{doc}{lattice}}}}, {\hyperref[\detokenize{region::doc}]{\sphinxcrossref{\DUrole{doc}{region}}}},
{\hyperref[\detokenize{create_box::doc}]{\sphinxcrossref{\DUrole{doc}{create\_box}}}}, {\hyperref[\detokenize{read_data::doc}]{\sphinxcrossref{\DUrole{doc}{read\_data}}}},
{\hyperref[\detokenize{read_restart::doc}]{\sphinxcrossref{\DUrole{doc}{read\_restart}}}}


\subsection{Default}
\label{\detokenize{create_atoms:default}}
\sphinxAtStartPar
The default for the \sphinxstyleemphasis{basis} keyword is that all created atoms are
assigned the argument \sphinxstyleemphasis{type} as their atom type (when single atoms are
being created).  The other defaults are \sphinxstyleemphasis{remap} = no, \sphinxstyleemphasis{rotate} = random,
\sphinxstyleemphasis{radscale} = 1.0, \sphinxstyleemphasis{radthresh} = x\sphinxhyphen{}lattice spacing, \sphinxstyleemphasis{overlap} not
checked, \sphinxstyleemphasis{maxtry} = 10, and \sphinxstyleemphasis{units} = lattice.


\bigskip\hrule\bigskip

\phantomsection\label{\detokenize{create_atoms:roberts2019}}
\sphinxAtStartPar
\sphinxstylestrong{(Roberts)} R. Roberts (2019) “Evenly Distributing Points in a Triangle.” Extreme Learning.  \sphinxurl{http://extremelearning.com.au/evenly-distributing-points-in-a-triangle/}

\sphinxstepscope

\index{create\_bonds@\spxentry{create\_bonds}}\ignorespaces 

\section{create\_bonds command}
\label{\detokenize{create_bonds:create-bonds-command}}\label{\detokenize{create_bonds:index-0}}\label{\detokenize{create_bonds::doc}}

\subsection{Syntax}
\label{\detokenize{create_bonds:syntax}}
\begin{sphinxVerbatim}[commandchars=\\\{\}]
\PYG{k}{create\PYGZus{}bonds}\PYG{+w}{ }\PYG{n}{style}\PYG{+w}{ }\PYG{n}{args}\PYG{+w}{ }\PYG{n}{...}\PYG{+w}{ }\PYG{n}{keyword}\PYG{+w}{ }\PYG{n}{value}\PYG{+w}{ }\PYG{n}{...}
\end{sphinxVerbatim}
\begin{itemize}
\item {} 
\sphinxAtStartPar
style = \sphinxstyleemphasis{many} or \sphinxstyleemphasis{single/bond} or \sphinxstyleemphasis{single/angle} or \sphinxstyleemphasis{single/dihedral} or \sphinxstyleemphasis{single/improper}

\end{itemize}
\begin{sphinxalltt}
\sphinxstyleemphasis{many} args = group\sphinxhyphen{}ID group2\sphinxhyphen{}ID btype rmin rmax
  group\sphinxhyphen{}ID = ID of first group
  group2\sphinxhyphen{}ID = ID of second group, bonds will be between atoms in the 2 groups
  btype = bond type of created bonds
  rmin = minimum distance between pair of atoms to bond together
  rmax = maximum distance between pair of atoms to bond together
\sphinxstyleemphasis{single/bond} args = btype batom1 batom2
  btype = bond type of new bond
  batom1,batom2 = atom IDs for two atoms in bond
\sphinxstyleemphasis{single/angle} args = atype aatom1 aatom2 aatom3
  atype = angle type of new angle
  aatom1,aatom2,aatom3 = atom IDs for three atoms in angle
\sphinxstyleemphasis{single/dihedral} args = dtype datom1 datom2 datom3 datom4
  dtype = dihedral type of new dihedral
  datom1,datom2,datom3,datom4 = atom IDs for four atoms in dihedral
\sphinxstyleemphasis{single/improper} args = itype iatom1 iatom2 iatom3 iatom4
  itype = improper type of new improper
  iatom1,iatom2,iatom3,iatom4 = atom IDs for four atoms in improper
\end{sphinxalltt}
\begin{itemize}
\item {} 
\sphinxAtStartPar
zero or more keyword/value pairs may be appended

\item {} 
\sphinxAtStartPar
keyword = \sphinxstyleemphasis{special}

\end{itemize}
\begin{sphinxalltt}
\sphinxstyleemphasis{special} value = \sphinxstyleemphasis{yes} or \sphinxstyleemphasis{no}
\end{sphinxalltt}


\subsection{Examples}
\label{\detokenize{create_bonds:examples}}
\begin{sphinxVerbatim}[commandchars=\\\{\}]
\PYG{k}{create\PYGZus{}bonds}\PYG{+w}{ }\PYG{n}{many}\PYG{+w}{ }\PYG{n}{all}\PYG{+w}{ }\PYG{n}{all}\PYG{+w}{ }\PYG{l+m}{1}\PYG{+w}{ }\PYG{l+m}{1.0}\PYG{+w}{ }\PYG{l+m}{1.2}
\PYG{k}{create\PYGZus{}bonds}\PYG{+w}{ }\PYG{n}{many}\PYG{+w}{ }\PYG{n}{surf}\PYG{+w}{ }\PYG{n}{solvent}\PYG{+w}{ }\PYG{l+m}{3}\PYG{+w}{ }\PYG{l+m}{2.0}\PYG{+w}{ }\PYG{l+m}{2.4}
\PYG{k}{create\PYGZus{}bonds}\PYG{+w}{ }\PYG{n}{single}\PYG{o}{/}\PYG{n}{bond}\PYG{+w}{ }\PYG{l+m}{1}\PYG{+w}{ }\PYG{l+m}{1}\PYG{+w}{ }\PYG{l+m}{2}
\PYG{k}{create\PYGZus{}bonds}\PYG{+w}{ }\PYG{n}{single}\PYG{o}{/}\PYG{n}{angle}\PYG{+w}{ }\PYG{l+m}{5}\PYG{+w}{ }\PYG{l+m}{52}\PYG{+w}{ }\PYG{l+m}{98}\PYG{+w}{ }\PYG{l+m}{107}\PYG{+w}{ }\PYG{n}{special}\PYG{+w}{ }\PYG{n}{no}
\PYG{k}{create\PYGZus{}bonds}\PYG{+w}{ }\PYG{n}{single}\PYG{o}{/}\PYG{n}{dihedral}\PYG{+w}{ }\PYG{l+m}{2}\PYG{+w}{ }\PYG{l+m}{4}\PYG{+w}{ }\PYG{l+m}{19}\PYG{+w}{ }\PYG{l+m}{27}\PYG{+w}{ }\PYG{l+m}{101}
\PYG{k}{create\PYGZus{}bonds}\PYG{+w}{ }\PYG{n}{single}\PYG{o}{/}\PYG{n}{improper}\PYG{+w}{ }\PYG{l+m}{3}\PYG{+w}{ }\PYG{l+m}{23}\PYG{+w}{ }\PYG{l+m}{26}\PYG{+w}{ }\PYG{l+m}{31}\PYG{+w}{ }\PYG{l+m}{57}
\end{sphinxVerbatim}


\subsection{Description}
\label{\detokenize{create_bonds:description}}
\sphinxAtStartPar
Create bonds between pairs of atoms that meet a specified distance
criteria.  Or create a single bond, angle, dihedral or improper between 2, 3,
or 4 specified atoms.

\sphinxAtStartPar
The new bond (angle, dihedral, improper) interactions will then be computed
during a simulation by the bond (angle, dihedral, improper) potential defined by
the {\hyperref[\detokenize{bond_style::doc}]{\sphinxcrossref{\DUrole{doc}{bond\_style}}}}, {\hyperref[\detokenize{bond_coeff::doc}]{\sphinxcrossref{\DUrole{doc}{bond\_coeff}}}},
{\hyperref[\detokenize{angle_style::doc}]{\sphinxcrossref{\DUrole{doc}{angle\_style}}}}, {\hyperref[\detokenize{angle_coeff::doc}]{\sphinxcrossref{\DUrole{doc}{angle\_coeff}}}},
{\hyperref[\detokenize{dihedral_style::doc}]{\sphinxcrossref{\DUrole{doc}{dihedral\_style}}}},
{\hyperref[\detokenize{dihedral_coeff::doc}]{\sphinxcrossref{\DUrole{doc}{dihedral\_coeff}}}}, {\hyperref[\detokenize{improper_style::doc}]{\sphinxcrossref{\DUrole{doc}{improper\_style}}}},
{\hyperref[\detokenize{improper_coeff::doc}]{\sphinxcrossref{\DUrole{doc}{improper\_coeff}}}} commands.

\sphinxAtStartPar
The \sphinxstyleemphasis{many} style is useful for adding bonds to a system (e.g., between
nearest neighbors in a lattice of atoms) without having to enumerate
all the bonds in the data file read by the {\hyperref[\detokenize{read_data::doc}]{\sphinxcrossref{\DUrole{doc}{read\_data}}}}
command.

\sphinxAtStartPar
The \sphinxstyleemphasis{single} styles are useful for adding bonds, angles, dihedrals, and
impropers to a system incrementally, then continuing a simulation.

\sphinxAtStartPar
Note that this command does not auto\sphinxhyphen{}create any angle, dihedral, or improper
interactions when a bond is added, nor does it auto\sphinxhyphen{}create any bonds
when an angle, dihedral, or improper is added.  It also will not auto\sphinxhyphen{}create
any angles when a dihedral or improper is added.  Thus, the flexibility of this
command is limited.  It can be used several times to create different types of
bond at different distances, but it cannot typically auto\sphinxhyphen{}create all the
bonds or angles or dihedrals or impropers that would normally be defined in a
data file for a complex system of molecules.

\begin{sphinxadmonition}{note}{Note:}
\sphinxAtStartPar
If the system has no bonds (angles, dihedrals, impropers) to begin
with, or if more bonds per atom are being added than currently exist,
then you must ensure that the number of bond types and the maximum
number of bonds per atom are set to large enough values, and
similarly for angles, dihedrals, impropers, and special neighbors,
otherwise an error may occur when too many bonds (angles, dihedrals,
impropers) are added to an atom.  If the {\hyperref[\detokenize{read_data::doc}]{\sphinxcrossref{\DUrole{doc}{read\_data}}}}
command is used to define the system, these parameters can be set via
its optional \sphinxstyleemphasis{extra/bond/types}, \sphinxstyleemphasis{extra/bond/per/atom}, and similar
keywords to the command.  If the {\hyperref[\detokenize{create_box::doc}]{\sphinxcrossref{\DUrole{doc}{create\_box}}}}
command is used to define the system, these two parameters can be set
via its optional \sphinxstyleemphasis{bond/types} and \sphinxstyleemphasis{extra/bond/per/atom} arguments,
and similarly for angles, dihedrals, and impropers.  See the
corresponding documentation pages for these two commands for details.
\end{sphinxadmonition}


\bigskip\hrule\bigskip


\sphinxAtStartPar
The \sphinxstyleemphasis{many} style will create bonds between pairs of atoms \(I,J\),
where \(I\) is in one of the two specified groups and \(J\) is in the
other.  The two groups can be the same (e.g., group “all”).  The created bonds
will be of bond type \sphinxstyleemphasis{btype}, where \sphinxstyleemphasis{btype} must be a value between 1 and the
number of bond types defined.

\sphinxAtStartPar
For a bond to be created, an \(I,J\) pair of atoms must be a distance
\(D\) apart such that \(r_\text{min} \le D \le r_\text{max}\).

\sphinxAtStartPar
The following settings must have been made in an input script before
the \sphinxstyleemphasis{many} style is used:
\begin{itemize}
\item {} 
\sphinxAtStartPar
special\_bonds weight for 1\textendash{}2 interactions must be 0.0

\item {} 
\sphinxAtStartPar
a {\hyperref[\detokenize{pair_style::doc}]{\sphinxcrossref{\DUrole{doc}{pair\_style}}}} must be defined

\item {} 
\sphinxAtStartPar
no {\hyperref[\detokenize{kspace_style::doc}]{\sphinxcrossref{\DUrole{doc}{kspace\_style}}}} defined

\item {} 
\sphinxAtStartPar
minimum {\hyperref[\detokenize{pair_style::doc}]{\sphinxcrossref{\DUrole{doc}{pair\_style}}}} cutoff + {\hyperref[\detokenize{neighbor::doc}]{\sphinxcrossref{\DUrole{doc}{neighbor}}}}
skin \(\ge r_\text{max}\)

\end{itemize}

\sphinxAtStartPar
These settings are required so that a neighbor list can be created to
search for nearby atoms.  Pairs of atoms that are already bonded cannot
appear in the neighbor list, to avoid creation of duplicate bonds.  The
neighbor list for all atom type pairs must also extend to a distance
that encompasses the \sphinxstyleemphasis{rmax} for new bonds to create.  When using
periodic boundary conditions, the box length in each periodic dimension
must be larger than \sphinxstyleemphasis{rmax}, so that no bonds are created between the
system and its own periodic image.

\begin{sphinxadmonition}{note}{Note:}
\sphinxAtStartPar
If you want to create bonds between pairs of 1\textendash{}3 or 1\textendash{}4 atoms in
the current bond topology, then you need to use {\hyperref[\detokenize{special_bonds::doc}]{\sphinxcrossref{\DUrole{doc}{special\_bonds
lj 0 1 1}}}} to ensure those pairs appear in the
neighbor list.  They will not appear with the default special\_bonds
settings, which are zero for 1\textendash{}2, 1\textendash{}3, and 1\textendash{}4 atoms.  1\textendash{}3 or 1\textendash{}4
atoms are those which are two hops or three hops apart in the bond
topology.
\end{sphinxadmonition}

\sphinxAtStartPar
An additional requirement for this style is that your system must be
ready to perform a simulation.  This means, for example, that all
{\hyperref[\detokenize{pair_style::doc}]{\sphinxcrossref{\DUrole{doc}{pair\_style}}}} coefficients be set via the
{\hyperref[\detokenize{pair_coeff::doc}]{\sphinxcrossref{\DUrole{doc}{pair\_coeff}}}} command.  A {\hyperref[\detokenize{bond_style::doc}]{\sphinxcrossref{\DUrole{doc}{bond\_style}}}}
command and all bond coefficients must also be set, even if no bonds
exist before this command is invoked.  This is because the building of
neighbor list requires initialization and setup of a simulation,
similar to what a {\hyperref[\detokenize{run::doc}]{\sphinxcrossref{\DUrole{doc}{run}}}} command would require.

\sphinxAtStartPar
Note that you can change any of these settings after this command
executes (e.g., if you wish to use long\sphinxhyphen{}range Coulombic interactions)
via the {\hyperref[\detokenize{kspace_style::doc}]{\sphinxcrossref{\DUrole{doc}{kspace\_style}}}} command for your subsequent
simulation.


\bigskip\hrule\bigskip


\sphinxAtStartPar
The \sphinxstyleemphasis{single/bond} style creates a single bond of type \sphinxstyleemphasis{btype} between
two atoms with IDs \sphinxstyleemphasis{batom1} and \sphinxstyleemphasis{batom2}.  \sphinxstyleemphasis{Btype} must be a value
between 1 and the number of bond types defined.

\sphinxAtStartPar
The \sphinxstyleemphasis{single/angle} style creates a single angle of type \sphinxstyleemphasis{atype}
between three atoms with IDs \sphinxstyleemphasis{aatom1}, \sphinxstyleemphasis{aatom2}, and \sphinxstyleemphasis{aatom3}.  The
ordering of the atoms is the same as in the \sphinxstyleemphasis{Angles} section of a data
file read by the {\hyperref[\detokenize{read_data::doc}]{\sphinxcrossref{\DUrole{doc}{read\_data}}}} command (i.e., the three atoms
are ordered linearly within the angle; the central atom is \sphinxstyleemphasis{aatom2}).
\sphinxstyleemphasis{Atype} must be a value between 1 and the number of angle types
defined.

\sphinxAtStartPar
The \sphinxstyleemphasis{single/dihedral} style creates a single dihedral of type \sphinxstyleemphasis{dtype}
between four atoms with IDs \sphinxstyleemphasis{datom1}, \sphinxstyleemphasis{datom2}, \sphinxstyleemphasis{datom3}, and \sphinxstyleemphasis{datom4}.  The
ordering of the atoms is the same as in the \sphinxstyleemphasis{Dihedrals} section of a data file
read by the {\hyperref[\detokenize{read_data::doc}]{\sphinxcrossref{\DUrole{doc}{read\_data}}}} command.  I.e. the 4 atoms are ordered
linearly within the dihedral.  \sphinxstyleemphasis{dtype} must be a value between 1 and
the number of dihedral types defined.

\sphinxAtStartPar
The \sphinxstyleemphasis{single/improper} style creates a single improper of type \sphinxstyleemphasis{itype}
between four atoms with IDs \sphinxstyleemphasis{iatom1}, \sphinxstyleemphasis{iatom2}, \sphinxstyleemphasis{iatom3}, and \sphinxstyleemphasis{iatom4}.  The
ordering of the atoms is the same as in the \sphinxstyleemphasis{Impropers} section of a data file
read by the {\hyperref[\detokenize{read_data::doc}]{\sphinxcrossref{\DUrole{doc}{read\_data}}}} command.  I.e. the 4 atoms are ordered
linearly within the improper.  \sphinxstyleemphasis{itype} must be a value between 1 and
the number of improper types defined.


\bigskip\hrule\bigskip


\sphinxAtStartPar
The keyword \sphinxstyleemphasis{special} controls whether an internal list of special
bonds is created after one or more bonds, or a single angle, dihedral, or
improper is added to the system.

\sphinxAtStartPar
The default value is \sphinxstyleemphasis{yes}.  A value of \sphinxstyleemphasis{no} cannot be used
with the \sphinxstyleemphasis{many} style.

\sphinxAtStartPar
This is an expensive operation since the bond topology for the system
must be walked to find all 1\textendash{}2, 1\textendash{}3, and 1\textendash{}4 interactions to store in an
internal list, which is used when pairwise interactions are weighted;
see the {\hyperref[\detokenize{special_bonds::doc}]{\sphinxcrossref{\DUrole{doc}{special\_bonds}}}} command for details.

\sphinxAtStartPar
Thus if you are adding a few bonds or a large list of angles all at
the same time, by using this command repeatedly, it is more efficient
to only trigger the internal list to be created once, after the last
bond (or angle, or dihedral, or improper) is added:

\begin{sphinxVerbatim}[commandchars=\\\{\}]
\PYG{k}{create\PYGZus{}bonds}\PYG{+w}{ }\PYG{n}{single}\PYG{o}{/}\PYG{n}{bond}\PYG{+w}{ }\PYG{l+m}{5}\PYG{+w}{ }\PYG{l+m}{52}\PYG{+w}{ }\PYG{l+m}{98}\PYG{+w}{ }\PYG{n}{special}\PYG{+w}{ }\PYG{n}{no}
\PYG{k}{create\PYGZus{}bonds}\PYG{+w}{ }\PYG{n}{single}\PYG{o}{/}\PYG{n}{bond}\PYG{+w}{ }\PYG{l+m}{5}\PYG{+w}{ }\PYG{l+m}{73}\PYG{+w}{ }\PYG{l+m}{74}\PYG{+w}{ }\PYG{n}{special}\PYG{+w}{ }\PYG{n}{no}
\PYG{n}{...}
\PYG{k}{create\PYGZus{}bonds}\PYG{+w}{ }\PYG{n}{single}\PYG{o}{/}\PYG{n}{bond}\PYG{+w}{ }\PYG{l+m}{5}\PYG{+w}{ }\PYG{l+m}{17}\PYG{+w}{ }\PYG{l+m}{386}\PYG{+w}{ }\PYG{n}{special}\PYG{+w}{ }\PYG{n}{no}
\PYG{k}{create\PYGZus{}bonds}\PYG{+w}{ }\PYG{n}{single}\PYG{o}{/}\PYG{n}{bond}\PYG{+w}{ }\PYG{l+m}{4}\PYG{+w}{ }\PYG{l+m}{112}\PYG{+w}{ }\PYG{l+m}{183}\PYG{+w}{ }\PYG{n}{special}\PYG{+w}{ }\PYG{n}{yes}
\end{sphinxVerbatim}

\sphinxAtStartPar
Note that you \sphinxstylestrong{must} ensure the internal list is rebuilt after the last
bond (angle, dihedral, improper) is added, \sphinxstyleemphasis{before} performing a simulation.
Otherwise, pairwise interactions will not be properly excluded or
weighted.  LAMMPS does \sphinxstylestrong{not} check that you have done this correctly.


\bigskip\hrule\bigskip



\subsection{Restrictions}
\label{\detokenize{create_bonds:restrictions}}
\sphinxAtStartPar
This command cannot be used with molecular systems defined using
molecule template files via the {\hyperref[\detokenize{molecule::doc}]{\sphinxcrossref{\DUrole{doc}{molecule}}}} and
{\hyperref[\detokenize{atom_style::doc}]{\sphinxcrossref{\DUrole{doc}{atom\_style template}}}} commands.

\sphinxAtStartPar
For style \sphinxstyleemphasis{many}, no {\hyperref[\detokenize{kspace_style::doc}]{\sphinxcrossref{\DUrole{doc}{kspace style}}}} must be
defined. Also, the \sphinxstyleemphasis{rmax} value must be smaller than any periodic box
length and the neighbor list cutoff (largest pair cutoff plus neighbor
skin).


\subsection{Related commands}
\label{\detokenize{create_bonds:related-commands}}
\sphinxAtStartPar
{\hyperref[\detokenize{create_atoms::doc}]{\sphinxcrossref{\DUrole{doc}{create\_atoms}}}}, {\hyperref[\detokenize{delete_bonds::doc}]{\sphinxcrossref{\DUrole{doc}{delete\_bonds}}}}


\subsection{Default}
\label{\detokenize{create_bonds:default}}
\sphinxAtStartPar
The keyword default is special = yes.

\sphinxstepscope

\index{create\_box@\spxentry{create\_box}}\ignorespaces 

\section{create\_box command}
\label{\detokenize{create_box:create-box-command}}\label{\detokenize{create_box:index-0}}\label{\detokenize{create_box::doc}}

\subsection{Syntax}
\label{\detokenize{create_box:syntax}}
\begin{sphinxVerbatim}[commandchars=\\\{\}]
\PYG{k}{create\PYGZus{}box }\PYG{n+nv+nvIdentifier}{N}\PYG{+w}{ }\PYG{n+nv+nvIdentifier}{region\PYGZhy{}ID}\PYG{+w}{ }\PYG{n}{keyword}\PYG{+w}{ }\PYG{n}{value}\PYG{+w}{ }\PYG{n}{...}
\PYG{k}{create\PYGZus{}box }\PYG{n+nv+nvIdentifier}{N}\PYG{+w}{ }\PYG{n+nv+nvIdentifier}{NULL}\PYG{+w}{ }\PYG{n}{alo}\PYG{+w}{ }\PYG{n}{ahi}\PYG{+w}{ }\PYG{n}{blo}\PYG{+w}{ }\PYG{n}{bhi}\PYG{+w}{ }\PYG{n}{clo}\PYG{+w}{ }\PYG{n}{chi}\PYG{+w}{ }\PYG{n}{keyword}\PYG{+w}{ }\PYG{n}{value}\PYG{+w}{ }\PYG{n}{...}
\end{sphinxVerbatim}
\begin{itemize}
\item {} 
\sphinxAtStartPar
N = \# of atom types to use in this simulation

\item {} 
\sphinxAtStartPar
region\sphinxhyphen{}ID = ID of region to use as simulation domain or NULL for general triclinic box

\item {} 
\sphinxAtStartPar
alo,ahi,blo,bhi,clo,chi = multipliers on a1,a2,a3 vectors defined by {\hyperref[\detokenize{lattice::doc}]{\sphinxcrossref{\DUrole{doc}{lattice}}}} command (only when region\sphinxhyphen{}ID = NULL)

\item {} 
\sphinxAtStartPar
zero or more keyword/value pairs may be appended

\item {} 
\sphinxAtStartPar
keyword = \sphinxstyleemphasis{bond/types} or \sphinxstyleemphasis{angle/types} or \sphinxstyleemphasis{dihedral/types} or \sphinxstyleemphasis{improper/types} or \sphinxstyleemphasis{extra/bond/per/atom} or \sphinxstyleemphasis{extra/angle/per/atom} or \sphinxstyleemphasis{extra/dihedral/per/atom} or \sphinxstyleemphasis{extra/improper/per/atom} or \sphinxstyleemphasis{extra/special/per/atom}
\begin{sphinxalltt}
\sphinxstyleemphasis{bond/types} value = \# of bond types
\sphinxstyleemphasis{angle/types} value = \# of angle types
\sphinxstyleemphasis{dihedral/types} value = \# of dihedral types
\sphinxstyleemphasis{improper/types} value = \# of improper types
\sphinxstyleemphasis{extra/bond/per/atom} value = \# of bonds per atom
\sphinxstyleemphasis{extra/angle/per/atom} value = \# of angles per atom
\sphinxstyleemphasis{extra/dihedral/per/atom} value = \# of dihedrals per atom
\sphinxstyleemphasis{extra/improper/per/atom} value = \# of impropers per atom
\sphinxstyleemphasis{extra/special/per/atom} value = \# of special neighbors per atom
\end{sphinxalltt}

\end{itemize}


\subsection{Examples}
\label{\detokenize{create_box:examples}}
\begin{sphinxVerbatim}[commandchars=\\\{\}]
\PYG{c}{\PYGZsh{} orthogonal or restricted triclinic box using regionID = mybox}
\PYG{k}{create\PYGZus{}box }\PYG{n+nv+nvIdentifier}{2}\PYG{+w}{ }\PYG{n+nv+nvIdentifier}{mybox}
\PYG{k}{create\PYGZus{}box }\PYG{n+nv+nvIdentifier}{2}\PYG{+w}{ }\PYG{n+nv+nvIdentifier}{mybox}\PYG{+w}{ }\PYG{n}{bond}\PYG{o}{/}\PYG{n}{types}\PYG{+w}{ }\PYG{l+m}{2}\PYG{+w}{ }\PYG{n}{extra}\PYG{o}{/}\PYG{n}{bond}\PYG{o}{/}\PYG{n}{per}\PYG{o}{/}\PYG{n}{atom}\PYG{+w}{ }\PYG{l+m}{1}
\end{sphinxVerbatim}

\begin{sphinxVerbatim}[commandchars=\\\{\}]
\PYG{c}{\PYGZsh{} 2d general triclinic box using primitive cell for 2d hex lattice}
\PYG{k}{lattice}\PYG{+w}{       }\PYG{n}{custom}\PYG{+w}{ }\PYG{l+m}{1.0}\PYG{+w}{ }\PYG{n}{a1}\PYG{+w}{ }\PYG{l+m}{1.0}\PYG{+w}{ }\PYG{l+m}{0.0}\PYG{+w}{ }\PYG{l+m}{0.0}\PYG{+w}{ }\PYG{n}{a2}\PYG{+w}{ }\PYG{l+m}{0.5}\PYG{+w}{ }\PYG{l+m}{0.86602540378}\PYG{+w}{ }\PYG{l+m}{0.0}\PYG{+w}{ }\PYG{o}{\PYGZam{}}
\PYG{+w}{              }\PYG{n}{a3}\PYG{+w}{ }\PYG{l+m}{0.0}\PYG{+w}{ }\PYG{l+m}{0.0}\PYG{+w}{ }\PYG{l+m}{1.0}\PYG{+w}{ }\PYG{n}{basis}\PYG{+w}{ }\PYG{l+m}{0.0}\PYG{+w}{ }\PYG{l+m}{0.0}\PYG{+w}{ }\PYG{l+m}{0.0}\PYG{+w}{ }\PYG{n}{triclinic}\PYG{o}{/}\PYG{n}{general}
\PYG{k}{create\PYGZus{}box    }\PYG{n+nv+nvIdentifier}{1}\PYG{+w}{ }\PYG{n+nv+nvIdentifier}{NULL}\PYG{+w}{ }\PYG{l+m}{0}\PYG{+w}{ }\PYG{l+m}{5}\PYG{+w}{ }\PYG{l+m}{0}\PYG{+w}{ }\PYG{l+m}{5}\PYG{+w}{ }\PYG{o}{\PYGZhy{}}\PYG{l+m}{0.5}\PYG{+w}{ }\PYG{l+m}{0.5}
\end{sphinxVerbatim}

\begin{sphinxVerbatim}[commandchars=\\\{\}]
\PYG{c}{\PYGZsh{} 3d general triclinic box using primitive cell for 3d fcc lattice}
\PYG{k}{lattice}\PYG{+w}{ }\PYG{n}{custom}\PYG{+w}{ }\PYG{l+m}{1.0}\PYG{+w}{ }\PYG{n}{a2}\PYG{+w}{ }\PYG{l+m}{0.0}\PYG{+w}{ }\PYG{l+m}{0.5}\PYG{+w}{ }\PYG{l+m}{0.5}\PYG{+w}{ }\PYG{n}{a1}\PYG{+w}{ }\PYG{l+m}{0.5}\PYG{+w}{ }\PYG{l+m}{0.0}\PYG{+w}{ }\PYG{l+m}{0.5}\PYG{+w}{ }\PYG{n}{a3}\PYG{+w}{ }\PYG{l+m}{0.5}\PYG{+w}{ }\PYG{l+m}{0.5}\PYG{+w}{ }\PYG{l+m}{0.0}\PYG{+w}{ }\PYG{n}{basis}\PYG{+w}{ }\PYG{l+m}{0.0}\PYG{+w}{ }\PYG{l+m}{0.0}\PYG{+w}{ }\PYG{l+m}{0.0}\PYG{+w}{ }\PYG{n}{triclinic}\PYG{o}{/}\PYG{n}{general}
\PYG{n}{create}\PYG{+w}{ }\PYG{n}{box}\PYG{+w}{ }\PYG{l+m}{1}\PYG{+w}{ }\PYG{n}{NULL}\PYG{+w}{ }\PYG{o}{\PYGZhy{}}\PYG{l+m}{5}\PYG{+w}{ }\PYG{l+m}{5}\PYG{+w}{ }\PYG{o}{\PYGZhy{}}\PYG{l+m}{10}\PYG{+w}{ }\PYG{l+m}{10}\PYG{+w}{ }\PYG{l+m}{0}\PYG{+w}{ }\PYG{l+m}{20}
\end{sphinxVerbatim}


\subsection{Description}
\label{\detokenize{create_box:description}}
\sphinxAtStartPar
This command creates a simulation box. It also partitions the box into
a regular 3d grid of smaller sub\sphinxhyphen{}boxes, one per processor (MPI task).
The geometry of the partitioning is based on the size and shape of the
simulation box, the number of processors being used and the settings
of the {\hyperref[\detokenize{processors::doc}]{\sphinxcrossref{\DUrole{doc}{processors}}}} command.  The partitioning can
later be changed by the {\hyperref[\detokenize{balance::doc}]{\sphinxcrossref{\DUrole{doc}{balance}}}} or {\hyperref[\detokenize{fix_balance::doc}]{\sphinxcrossref{\DUrole{doc}{fix balance}}}} commands.

\sphinxAtStartPar
Simulation boxes in LAMMPS can be either orthogonal or triclinic in
shape.  Orthogonal boxes are a brick in 3d (rectangle in 2d) with 6
faces that are each perpendicular to one of the standard xyz
coordinate axes.  Triclinic boxes are a parallelepiped in 3d
(parallelogram in 2d) with opposite pairs of faces parallel to each
other.  LAMMPS supports two forms of triclinic boxes, restricted and
general, which differ in how the box is oriented with respect to the
xyz coordinate axes.  See the {\hyperref[\detokenize{Howto_triclinic::doc}]{\sphinxcrossref{\DUrole{doc}{Howto triclinic}}}}
for a detailed description of all 3 kinds of simulation boxes.

\sphinxAtStartPar
The argument \sphinxstyleemphasis{N} is the number of atom types that will be used in the
simulation.

\sphinxAtStartPar
Orthogonal and restricted triclinic boxes are created by specifying a
region ID previously defined by the {\hyperref[\detokenize{region::doc}]{\sphinxcrossref{\DUrole{doc}{region}}}} command.
General triclinic boxes are discussed below.

\sphinxAtStartPar
If the region is not of style \sphinxstyleemphasis{prism}, then LAMMPS encloses the region
(block, sphere, etc.) with an axis\sphinxhyphen{}aligned orthogonal bounding box
which becomes the simulation domain.  For a 2d simulation, the zlo and
zhi values of the simulation box must straddle zero.

\sphinxAtStartPar
If the region is of style \sphinxstyleemphasis{prism}, LAMMPS creates a non\sphinxhyphen{}orthogonal
simulation domain shaped as a parallelepiped with triclinic symmetry.
As defined by the {\hyperref[\detokenize{region::doc}]{\sphinxcrossref{\DUrole{doc}{region prism}}}} command, the
parallelepiped has an “origin” at (xlo,ylo,zlo) and three edge vectors
starting from the origin given by \(\vec a =
(x_\text{hi}-x_\text{lo},0,0)\); \(\vec b =
(xy,y_\text{hi}-y_\text{lo},0)\); and \(\vec c =
(xz,yz,z_\text{hi}-z_\text{lo})\).  In LAMMPS lingo, this is a
restricted triclinic box because the three edge vectors cannot be
defined in arbitrary (general) directions.  The parameters \sphinxstyleemphasis{xy},
\sphinxstyleemphasis{xz}, and \sphinxstyleemphasis{yz} can be 0.0 or positive or negative values and are
called “tilt factors” because they are the amount of displacement
applied to faces of an originally orthogonal box to transform it into
the parallelepiped.  For a 2d simulation, the zlo and zhi values of
the simulation box must straddle zero.

\sphinxAtStartPar
Typically a \sphinxstyleemphasis{prism} region used with the create\_box command should
have tilt factors \((xy,xz,yz)\) that do not skew the box more
than half the distance of the parallel box length.  For example, if
\(x_\text{lo} = 2\) and \(x_\text{hi} = 12\), then the
\(x\) box length is 10 and the \(xy\) tilt factor must be
between \(-5\) and \(5\).  Similarly, both \(xz\) and
\(yz\) must be between \(-(x_\text{hi}-x_\text{lo})/2\) and
\(+(y_\text{hi}-y_\text{lo})/2\).  Note that this is not a
limitation, since if the maximum tilt factor is 5 (as in this
example), then configurations with tilt \(= \dots, -15\),
\(-5\), \(5\), \(15\), \(25, \dots\) are all
geometrically equivalent.

\sphinxAtStartPar
LAMMPS will issue a warning if the tilt factors of the created box do
not meet this criterion.  This is because simulations with large tilt
factors may run inefficiently, since they require more ghost atoms and
thus more communication.  With very large tilt factors, LAMMPS may
eventually produce incorrect trajectories and stop with errors due to
lost atoms or similar issues.

\sphinxAtStartPar
See the {\hyperref[\detokenize{Howto_triclinic::doc}]{\sphinxcrossref{\DUrole{doc}{Howto triclinic}}}} page for geometric
descriptions of triclinic boxes and tilt factors, as well as how to
transform the restricted triclinic parameters to and from other
commonly used triclinic representations.

\sphinxAtStartPar
When a prism region is used, the simulation domain should normally be
periodic in the dimension that the tilt is applied to, which is given
by the second dimension of the tilt factor (e.g., \(y\) for
\(xy\) tilt).  This is so that pairs of atoms interacting across
that boundary will have one of them shifted by the tilt factor.
Periodicity is set by the {\hyperref[\detokenize{boundary::doc}]{\sphinxcrossref{\DUrole{doc}{boundary}}}} command.  For
example, if the \(xy\) tilt factor is non\sphinxhyphen{}zero, then the \(y\)
dimension should be periodic.  Similarly, the \(z\) dimension
should be periodic if \(xz\) or \(yz\) is non\sphinxhyphen{}zero.  LAMMPS
does not require this periodicity, but you may lose atoms if this is
not the case.

\sphinxAtStartPar
Note that if your simulation will tilt the box (e.g., via the
{\hyperref[\detokenize{fix_deform::doc}]{\sphinxcrossref{\DUrole{doc}{fix deform}}}} command), the simulation box must be
created as triclinic, even if the tilt factors are initially 0.0.  You
can also change an orthogonal box to a triclinic box or vice versa by
using the {\hyperref[\detokenize{change_box::doc}]{\sphinxcrossref{\DUrole{doc}{change box}}}} command with its \sphinxstyleemphasis{ortho} and
\sphinxstyleemphasis{triclinic} options.

\begin{sphinxadmonition}{note}{Note:}
\sphinxAtStartPar
If the system is non\sphinxhyphen{}periodic (in a dimension), then you should not
make the lo/hi box dimensions (as defined in your {\hyperref[\detokenize{region::doc}]{\sphinxcrossref{\DUrole{doc}{region}}}} command) radically smaller/larger than the extent of the
atoms you eventually plan to create (e.g., via the
{\hyperref[\detokenize{create_atoms::doc}]{\sphinxcrossref{\DUrole{doc}{create\_atoms}}}} command).  For example, if your
atoms extend from 0 to 50, you should not specify the box bounds as
\(-10000\) and \(10000\). This is because as described
above, LAMMPS uses the specified box size to lay out the 3d grid of
processors.  A huge (mostly empty) box will be sub\sphinxhyphen{}optimal for
performance when using “fixed” boundary conditions (see the
{\hyperref[\detokenize{boundary::doc}]{\sphinxcrossref{\DUrole{doc}{boundary}}}} command).  When using “shrink\sphinxhyphen{}wrap”
boundary conditions (see the {\hyperref[\detokenize{boundary::doc}]{\sphinxcrossref{\DUrole{doc}{boundary}}}} command), a
huge (mostly empty) box may cause a parallel simulation to lose
atoms the first time that LAMMPS shrink\sphinxhyphen{}wraps the box around the
atoms.
\end{sphinxadmonition}


\bigskip\hrule\bigskip


\sphinxAtStartPar
As noted above, general triclinic boxes in LAMMPS allow the box to
have arbitrary edge vectors \sphinxstylestrong{A}, \sphinxstylestrong{B}, \sphinxstylestrong{C}.  The only
restrictions are that the three vectors be distinct, non\sphinxhyphen{}zero, and not
co\sphinxhyphen{}planar.  They must also define a right\sphinxhyphen{}handed system such that
(\sphinxstylestrong{A} x \sphinxstylestrong{B}) points in the direction of \sphinxstylestrong{C}.  Note that a
left\sphinxhyphen{}handed system can be converted to a right\sphinxhyphen{}handed system by simply
swapping the order of any pair of the \sphinxstylestrong{A}, \sphinxstylestrong{B}, \sphinxstylestrong{C} vectors.

\sphinxAtStartPar
To create a general triclinic boxes, the region is specified as NULL
and the next 6 parameters (alo,ahi,blo,bhi,clo,chi) define the three
edge vectors \sphinxstylestrong{A}, \sphinxstylestrong{B}, \sphinxstylestrong{C} using additional information
previously defined by the {\hyperref[\detokenize{lattice::doc}]{\sphinxcrossref{\DUrole{doc}{lattice}}}} command.

\sphinxAtStartPar
The lattice must be of style \sphinxstyleemphasis{custom} and use its \sphinxstyleemphasis{triclinic/general}
option.  This insures the lattice satisfies the restrictions listed
above.  The \sphinxstyleemphasis{a1, *a2}, \sphinxstyleemphasis{a3} settings of the {\hyperref[\detokenize{lattice::doc}]{\sphinxcrossref{\DUrole{doc}{lattice}}}}
command define the edge vectors of a unit cell of the general
triclinic lattice.  This command uses them to define the three edge
vectors and origin of the general triclinic box as:
\begin{itemize}
\item {} 
\sphinxAtStartPar
\sphinxstylestrong{A} = (ahi\sphinxhyphen{}alo) * \sphinxstyleemphasis{a1}

\item {} 
\sphinxAtStartPar
\sphinxstylestrong{B} = (bhi\sphinxhyphen{}blo) * \sphinxstyleemphasis{a2}

\item {} 
\sphinxAtStartPar
\sphinxstylestrong{C} = (chi\sphinxhyphen{}clo) * \sphinxstyleemphasis{a3}

\item {} 
\sphinxAtStartPar
origin = (alo*a1 + blo*a2 + clo*a3)

\end{itemize}

\sphinxAtStartPar
For 2d general triclinic boxes, clo = \sphinxhyphen{}0.5 and chi = 0.5 is required.

\begin{sphinxadmonition}{note}{Note:}
\sphinxAtStartPar
LAMMPS allows specification of general triclinic simulation boxes
as a convenience for users who may be converting data from
solid\sphinxhyphen{}state crystallographic representations or from DFT codes for
input to LAMMPS.  However, as explained on the
{\hyperref[\detokenize{Howto_triclinic::doc}]{\sphinxcrossref{\DUrole{doc}{Howto\_triclinic}}}} doc page, internally,
LAMMPS only uses restricted triclinic simulation boxes.  This means
the box defined by this command and per\sphinxhyphen{}atom information
(e.g. coordinates, velocities) defined by the {\hyperref[\detokenize{create_atoms::doc}]{\sphinxcrossref{\DUrole{doc}{create\_atoms}}}} command are converted (rotated) from general to
restricted triclinic form when the two commands are invoked.  The
\textless{}Howto\_triclinic\textgreater{}\textasciigrave{} doc page also discusses other LAMMPS commands
which can input/output general triclinic representations of the
simulation box and per\sphinxhyphen{}atom data.
\end{sphinxadmonition}


\bigskip\hrule\bigskip


\sphinxAtStartPar
The optional keywords can be used to create a system that allows for
bond (angle, dihedral, improper) interactions, or for molecules with
special 1\textendash{}2, 1\textendash{}3, or 1\textendash{}4 neighbors to be added later.  These
optional keywords serve the same purpose as the analogous keywords
that can be used in a data file which are recognized by the
{\hyperref[\detokenize{read_data::doc}]{\sphinxcrossref{\DUrole{doc}{read\_data}}}} command when it sets up a system.

\sphinxAtStartPar
Note that if these keywords are not used, then the create\_box command
creates an atomic (non\sphinxhyphen{}molecular) simulation that does not allow bonds
between pairs of atoms to be defined, or a {\hyperref[\detokenize{bond_style::doc}]{\sphinxcrossref{\DUrole{doc}{bond potential}}}} to be specified, or for molecules with special neighbors
to be added to the system by commands such as {\hyperref[\detokenize{create_atoms::doc}]{\sphinxcrossref{\DUrole{doc}{create\_atoms mol}}}}, {\hyperref[\detokenize{fix_deposit::doc}]{\sphinxcrossref{\DUrole{doc}{fix deposit}}}} or {\hyperref[\detokenize{fix_pour::doc}]{\sphinxcrossref{\DUrole{doc}{fix pour}}}}.

\sphinxAtStartPar
As an example, see the examples/deposit/in.deposit.molecule script,
which deposits molecules onto a substrate.  Initially there are no
molecules in the system, but they are added later by the {\hyperref[\detokenize{fix_deposit::doc}]{\sphinxcrossref{\DUrole{doc}{fix
deposit}}}} command.  The create\_box command in the script
uses the bond/types and extra/bond/per/atom keywords to allow this.
If the added molecule contained more than one special bond (allowed by
default), an extra/special/per/atom keyword would also need to be
specified.


\bigskip\hrule\bigskip



\subsection{Restrictions}
\label{\detokenize{create_box:restrictions}}
\sphinxAtStartPar
An {\hyperref[\detokenize{atom_style::doc}]{\sphinxcrossref{\DUrole{doc}{atom\_style}}}} and {\hyperref[\detokenize{region::doc}]{\sphinxcrossref{\DUrole{doc}{region}}}} must have
been previously defined to use this command.


\subsection{Related commands}
\label{\detokenize{create_box:related-commands}}
\sphinxAtStartPar
{\hyperref[\detokenize{read_data::doc}]{\sphinxcrossref{\DUrole{doc}{read\_data}}}}, {\hyperref[\detokenize{create_atoms::doc}]{\sphinxcrossref{\DUrole{doc}{create\_atoms}}}},
{\hyperref[\detokenize{region::doc}]{\sphinxcrossref{\DUrole{doc}{region}}}}


\subsection{Default}
\label{\detokenize{create_box:default}}
\sphinxAtStartPar
none

\sphinxstepscope

\index{delete\_atoms@\spxentry{delete\_atoms}}\ignorespaces 

\section{delete\_atoms command}
\label{\detokenize{delete_atoms:delete-atoms-command}}\label{\detokenize{delete_atoms:index-0}}\label{\detokenize{delete_atoms::doc}}

\subsection{Syntax}
\label{\detokenize{delete_atoms:syntax}}
\begin{sphinxVerbatim}[commandchars=\\\{\}]
\PYG{k}{delete\PYGZus{}atoms}\PYG{+w}{ }\PYG{n}{style}\PYG{+w}{ }\PYG{n}{args}\PYG{+w}{ }\PYG{n}{keyword}\PYG{+w}{ }\PYG{n}{value}\PYG{+w}{ }\PYG{n}{...}
\end{sphinxVerbatim}
\begin{itemize}
\item {} 
\sphinxAtStartPar
style = \sphinxstyleemphasis{group} or \sphinxstyleemphasis{region} or \sphinxstyleemphasis{overlap} or \sphinxstyleemphasis{random} or \sphinxstyleemphasis{variable}
\begin{sphinxalltt}
\sphinxstyleemphasis{group} args = group\sphinxhyphen{}ID
\sphinxstyleemphasis{region} args = region\sphinxhyphen{}ID
\sphinxstyleemphasis{overlap} args = cutoff group1\sphinxhyphen{}ID group2\sphinxhyphen{}ID
  cutoff = delete one atom from pairs of atoms within the cutoff (distance units)
  group1\sphinxhyphen{}ID = one atom in pair must be in this group
  group2\sphinxhyphen{}ID = other atom in pair must be in this group
\sphinxstyleemphasis{random} args = ranstyle value eflag group\sphinxhyphen{}ID region\sphinxhyphen{}ID seed
  ranstyle = \sphinxstyleemphasis{fraction} or \sphinxstyleemphasis{count}
    for \sphinxstyleemphasis{fraction}:
      value = fraction (0.0 to 1.0) of eligible atoms to delete
      eflag =  \sphinxstyleemphasis{no} for fast approximate deletion, \sphinxstyleemphasis{yes} for exact deletion
    for \sphinxstyleemphasis{count}:
      value = number of atoms to delete
      eflag = \sphinxstyleemphasis{no} for warning if count \textgreater{} eligible atoms, \sphinxstyleemphasis{yes} for error
  group\sphinxhyphen{}ID = group within which to perform deletions
  region\sphinxhyphen{}ID = region within which to perform deletions
              or NULL to only impose the group criterion
  seed = random number seed (positive integer)
\sphinxstyleemphasis{variable} args = variable\sphinxhyphen{}name
\end{sphinxalltt}

\item {} 
\sphinxAtStartPar
zero or more keyword/value pairs may be appended

\item {} 
\sphinxAtStartPar
keyword = \sphinxstyleemphasis{compress} or \sphinxstyleemphasis{bond} or \sphinxstyleemphasis{mol}
\begin{sphinxalltt}
\sphinxstyleemphasis{compress} value = \sphinxstyleemphasis{no} or \sphinxstyleemphasis{yes}
\sphinxstyleemphasis{bond} value = \sphinxstyleemphasis{no} or \sphinxstyleemphasis{yes}
\sphinxstyleemphasis{mol} value = \sphinxstyleemphasis{no} or \sphinxstyleemphasis{yes}
\end{sphinxalltt}

\end{itemize}


\subsection{Examples}
\label{\detokenize{delete_atoms:examples}}
\begin{sphinxVerbatim}[commandchars=\\\{\}]
\PYG{k}{delete\PYGZus{}atoms}\PYG{+w}{ }\PYG{k}{group }\PYG{n+nv+nvIdentifier}{edge}
\PYG{k}{delete\PYGZus{}atoms}\PYG{+w}{ }\PYG{k}{region }\PYG{n+nv+nvIdentifier}{sphere}\PYG{+w}{ }\PYG{n}{compress}\PYG{+w}{ }\PYG{n}{no}
\PYG{k}{delete\PYGZus{}atoms}\PYG{+w}{ }\PYG{n}{overlap}\PYG{+w}{ }\PYG{l+m}{0.3}\PYG{+w}{ }\PYG{n}{all}\PYG{+w}{ }\PYG{n}{all}
\PYG{k}{delete\PYGZus{}atoms}\PYG{+w}{ }\PYG{n}{overlap}\PYG{+w}{ }\PYG{l+m}{0.5}\PYG{+w}{ }\PYG{n}{solvent}\PYG{+w}{ }\PYG{n}{colloid}
\PYG{k}{delete\PYGZus{}atoms}\PYG{+w}{ }\PYG{n}{random}\PYG{+w}{ }\PYG{n}{fraction}\PYG{+w}{ }\PYG{l+m}{0.1}\PYG{+w}{ }\PYG{n}{yes}\PYG{+w}{ }\PYG{n}{all}\PYG{+w}{ }\PYG{n}{cube}\PYG{+w}{ }\PYG{l+m}{482793}\PYG{+w}{ }\PYG{n}{bond}\PYG{+w}{ }\PYG{n}{yes}
\PYG{k}{delete\PYGZus{}atoms}\PYG{+w}{ }\PYG{n}{random}\PYG{+w}{ }\PYG{n}{fraction}\PYG{+w}{ }\PYG{l+m}{0.3}\PYG{+w}{ }\PYG{n}{no}\PYG{+w}{ }\PYG{n}{polymer}\PYG{+w}{ }\PYG{n}{NULL}\PYG{+w}{ }\PYG{l+m}{482793}\PYG{+w}{ }\PYG{n}{bond}\PYG{+w}{ }\PYG{n}{yes}
\PYG{k}{delete\PYGZus{}atoms}\PYG{+w}{ }\PYG{n}{random}\PYG{+w}{ }\PYG{n}{count}\PYG{+w}{ }\PYG{l+m}{500}\PYG{+w}{ }\PYG{n}{no}\PYG{+w}{ }\PYG{n}{ions}\PYG{+w}{ }\PYG{n}{NULL}\PYG{+w}{ }\PYG{l+m}{482793}
\PYG{k}{delete\PYGZus{}atoms}\PYG{+w}{ }\PYG{k}{variable }\PYG{n+nv+nvIdentifier}{checkers}
\end{sphinxVerbatim}


\subsection{Description}
\label{\detokenize{delete_atoms:description}}
\sphinxAtStartPar
Delete the specified atoms.  This command can be used, for example, to
carve out voids from a block of material or to delete created atoms
that are too close to each other (e.g., at a grain boundary).

\sphinxAtStartPar
For style \sphinxstyleemphasis{group}, all atoms belonging to the group are deleted.

\sphinxAtStartPar
For style \sphinxstyleemphasis{region}, all atoms in the region volume are deleted.
Additional atoms can be deleted if they are in a molecule for which
one or more atoms were deleted within the region; see the \sphinxstyleemphasis{mol}
keyword discussion below.

\sphinxAtStartPar
For style \sphinxstyleemphasis{overlap} pairs of atoms whose distance of separation is
within the specified cutoff distance are searched for, and one of the
two atoms is deleted.  Only pairs where one of the two atoms is in the
first group specified and the other atom is in the second group are
considered.  The atom that is in the first group is the one that is
deleted.

\sphinxAtStartPar
Note that it is OK for the two group IDs to be the same (e.g., group
\sphinxstyleemphasis{all}), or for some atoms to be members of both groups.  In these
cases, either atom in the pair may be deleted.  Also note that if
there are atoms which are members of both groups, the only guarantee
is that at the end of the deletion operation, enough deletions will
have occurred that no atom pairs within the cutoff will remain
(subject to the group restriction).  There is no guarantee that the
minimum number of atoms will be deleted, or that the same atoms will
be deleted when running on different numbers of processors.

\sphinxAtStartPar
For style \sphinxstyleemphasis{random} a subset of eligible atoms are deleted.  Which
atoms to delete are chosen randomly using the specified random number
\sphinxstyleemphasis{seed}.  Which atoms are deleted may vary when running on different
numbers of processors.

\sphinxAtStartPar
For \sphinxstyleemphasis{ranstyle} = \sphinxstyleemphasis{fraction}, the specified fractional \sphinxstyleemphasis{value} (0.0 to
1.0) of eligible atoms are deleted.  If \sphinxstyleemphasis{eflag} is set to \sphinxstyleemphasis{no}, then
the number of deleted atoms will be approximate, but the operation
will be fast.  If \sphinxstyleemphasis{eflag} is set to \sphinxstyleemphasis{yes}, then the number deleted
will match the requested fraction, but for large systems the selection
of deleted atoms may take additional time to determine.

\sphinxAtStartPar
For \sphinxstyleemphasis{ranstyle} = \sphinxstyleemphasis{count}, the specified integer \sphinxstyleemphasis{value} is the number
of eligible atoms are deleted.  If \sphinxstyleemphasis{eflag} is set to \sphinxstyleemphasis{no}, then if the
requested number is larger then the number of eligible atoms, a
warning is issued and only the eligible atoms are deleted instead of
the requested \sphinxstyleemphasis{value}.  If \sphinxstyleemphasis{eflag} is set to \sphinxstyleemphasis{yes}, an error is
triggered instead and LAMMPS will exit.  For large systems the
selection of atoms to delete may take additional time to determine,
the same as for requesting an exact fraction with \sphinxstyleemphasis{pstyle} =
\sphinxstyleemphasis{fraction}.

\sphinxAtStartPar
Which atoms are eligible for deletion for style \sphinxstyleemphasis{random} is determined
by the specified \sphinxstyleemphasis{group\sphinxhyphen{}ID} and \sphinxstyleemphasis{region\sphinxhyphen{}ID}.  To be eligible, an atom
must be in both the specified group and region.  If \sphinxstyleemphasis{group\sphinxhyphen{}ID} = all,
there is effectively no group criterion.  If \sphinxstyleemphasis{region\sphinxhyphen{}ID} is specified
as NULL, no region criterion is imposed.

\sphinxAtStartPar
\DUrole{versionmodified,added}{Added in version 4May2022.}

\sphinxAtStartPar
For style \sphinxstyleemphasis{variable}, all atoms for which the atom\sphinxhyphen{}style variable with
the given name evaluates to non\sphinxhyphen{}zero will be deleted. Additional atoms
can be deleted if they are in a molecule for which one or more atoms
were deleted within the region; see the \sphinxstyleemphasis{mol} keyword discussion below.
This option allows complex selections of atoms not covered by the
other options listed above.


\bigskip\hrule\bigskip


\sphinxAtStartPar
Here is the meaning of the optional keywords.

\sphinxAtStartPar
If the \sphinxstyleemphasis{compress} keyword is set to \sphinxstyleemphasis{yes}, then after atoms are
deleted, then atom IDs are re\sphinxhyphen{}assigned so that they run from 1 to the
number of atoms in the system.  Note that this is not done for
molecular systems (see the {\hyperref[\detokenize{atom_style::doc}]{\sphinxcrossref{\DUrole{doc}{atom\_style}}}} command),
regardless of the \sphinxstyleemphasis{compress} setting, since it would foul up the bond
connectivity that has already been assigned.  However, the
{\hyperref[\detokenize{reset_atoms::doc}]{\sphinxcrossref{\DUrole{doc}{reset\_atoms id}}}} command can be used after this
command to accomplish the same thing.

\sphinxAtStartPar
Note that the re\sphinxhyphen{}assignment of IDs is not really a compression, where
gaps in atom IDs are removed by decrementing atom IDs that are larger.
Instead the IDs for all atoms are erased, and new IDs are assigned so
that the atoms owned by individual processors have consecutive IDs, as
the {\hyperref[\detokenize{create_atoms::doc}]{\sphinxcrossref{\DUrole{doc}{create\_atoms}}}} command explains.

\sphinxAtStartPar
A molecular system with fixed bonds, angles, dihedrals, or improper
interactions, is one where the topology of the interactions is
typically defined in the data file read by the {\hyperref[\detokenize{read_data::doc}]{\sphinxcrossref{\DUrole{doc}{read\_data}}}} command, and where the interactions themselves are
defined with the {\hyperref[\detokenize{bond_style::doc}]{\sphinxcrossref{\DUrole{doc}{bond\_style}}}}, {\hyperref[\detokenize{angle_style::doc}]{\sphinxcrossref{\DUrole{doc}{angle\_style}}}}, etc. commands.  If you delete atoms from such a system,
you must be careful not to end up with bonded interactions that are
stored by remaining atoms but which include deleted atoms.  This will
cause LAMMPS to generate a “missing atoms” error when the bonded
interaction is computed.  The \sphinxstyleemphasis{bond} and \sphinxstyleemphasis{mol} keywords offer two ways
to do that.

\sphinxAtStartPar
It the \sphinxstyleemphasis{bond} keyword is set to \sphinxstyleemphasis{yes} then any bond or angle or
dihedral or improper interaction that includes a deleted atom is also
removed from the lists of such interactions stored by non\sphinxhyphen{}deleted
atoms.  Note that simply deleting interactions due to dangling bonds
(e.g., at a surface) may result in a inaccurate or invalid model for
the remaining atoms.

\sphinxAtStartPar
It the \sphinxstyleemphasis{mol} keyword is set to \sphinxstyleemphasis{yes}, then for every atom that is
deleted, all other atoms in the same molecule (with the same molecule
ID) will also be deleted.  This is not done for atoms with molecule ID
= 0, since such an ID is assumed to flag isolated atoms that are not
part of molecules.

\begin{sphinxadmonition}{note}{Note:}
\sphinxAtStartPar
The molecule deletion operation is invoked after all individual
atoms have been deleted using the rules described above for each
style.  This means additional atoms may be deleted that are not in the
group or region, that are not required by the overlap cutoff
criterion, or that will create a higher fraction of porosity than was
requested.
\end{sphinxadmonition}


\subsection{Restrictions}
\label{\detokenize{delete_atoms:restrictions}}
\sphinxAtStartPar
The \sphinxstyleemphasis{overlap} styles requires inter\sphinxhyphen{}processor communication to acquire
ghost atoms and build a neighbor list.  This means that your system
must be ready to perform a simulation before using this command (force
fields setup, atom masses set, etc.).  Since a neighbor list is used to
find overlapping atom pairs, it also means that you must define a
{\hyperref[\detokenize{pair_style::doc}]{\sphinxcrossref{\DUrole{doc}{pair style}}}} with the minimum force cutoff distance
between any pair of atoms types (plus the {\hyperref[\detokenize{neighbor::doc}]{\sphinxcrossref{\DUrole{doc}{neighbor}}}}
skin) \(\ge\) the specified overlap cutoff.

\sphinxAtStartPar
If the {\hyperref[\detokenize{special_bonds::doc}]{\sphinxcrossref{\DUrole{doc}{special\_bonds}}}} command is used with a
setting of 0, then a pair of bonded atoms (1\textendash{}2, 1\textendash{}3, or 1\textendash{}4) will not
appear in the neighbor list, and thus will not be considered for
deletion by the \sphinxstyleemphasis{overlap} styles.  You probably do not want to
delete one atom in a bonded pair anyway.

\sphinxAtStartPar
The \sphinxstyleemphasis{bond yes} option cannot be used with molecular systems defined
using molecule template files via the {\hyperref[\detokenize{molecule::doc}]{\sphinxcrossref{\DUrole{doc}{molecule}}}} and
{\hyperref[\detokenize{atom_style::doc}]{\sphinxcrossref{\DUrole{doc}{atom\_style template}}}} commands.


\subsection{Related commands}
\label{\detokenize{delete_atoms:related-commands}}
\sphinxAtStartPar
{\hyperref[\detokenize{create_atoms::doc}]{\sphinxcrossref{\DUrole{doc}{create\_atoms}}}}, {\hyperref[\detokenize{reset_atoms::doc}]{\sphinxcrossref{\DUrole{doc}{reset\_atoms id}}}}


\subsection{Default}
\label{\detokenize{delete_atoms:default}}
\sphinxAtStartPar
The option defaults are compress = yes, bond = no, mol = no.

\sphinxstepscope

\index{delete\_bonds@\spxentry{delete\_bonds}}\ignorespaces 

\section{delete\_bonds command}
\label{\detokenize{delete_bonds:delete-bonds-command}}\label{\detokenize{delete_bonds:index-0}}\label{\detokenize{delete_bonds::doc}}

\subsection{Syntax}
\label{\detokenize{delete_bonds:syntax}}
\begin{sphinxVerbatim}[commandchars=\\\{\}]
\PYG{k}{delete\PYGZus{}bonds }\PYG{n+nv+nvIdentifier}{group\PYGZhy{}ID}\PYG{+w}{ }\PYG{n}{style}\PYG{+w}{ }\PYG{n}{arg}\PYG{+w}{ }\PYG{n}{keyword}\PYG{+w}{ }\PYG{n}{...}
\end{sphinxVerbatim}
\begin{itemize}
\item {} 
\sphinxAtStartPar
group\sphinxhyphen{}ID = group ID

\item {} 
\sphinxAtStartPar
style = \sphinxstyleemphasis{multi} or \sphinxstyleemphasis{atom} or \sphinxstyleemphasis{bond} or \sphinxstyleemphasis{angle} or \sphinxstyleemphasis{dihedral} or \sphinxstyleemphasis{improper} or \sphinxstyleemphasis{stats}
\begin{sphinxalltt}
\sphinxstyleemphasis{multi} arg = none
\sphinxstyleemphasis{atom} arg = an atom type or range of types (see below)
\sphinxstyleemphasis{bond} arg = a bond type or range of types (see below)
\sphinxstyleemphasis{angle} arg = an angle type or range of types (see below)
\sphinxstyleemphasis{dihedral} arg = a dihedral type or range of types (see below)
\sphinxstyleemphasis{improper} arg = an improper type or range of types (see below)
\sphinxstyleemphasis{stats} arg = none
\end{sphinxalltt}

\item {} 
\sphinxAtStartPar
zero or more keywords may be appended

\item {} 
\sphinxAtStartPar
keyword = \sphinxstyleemphasis{any} or \sphinxstyleemphasis{undo} or \sphinxstyleemphasis{remove} or \sphinxstyleemphasis{special}
\begin{sphinxalltt}
\sphinxstyleemphasis{any} arg = none = turn off interactions if any atoms are in the group (or on if \sphinxstyleemphasis{undo} is also used)
\sphinxstyleemphasis{undo} arg = none = turn specified bonds on instead of off
\sphinxstyleemphasis{remove} arg = permanently remove bonds that have been turned off
\sphinxstyleemphasis{special} arg = recompute pairwise 1\sphinxhyphen{}2, 1\sphinxhyphen{}3, and 1\sphinxhyphen{}4 lists
\end{sphinxalltt}

\end{itemize}


\subsection{Examples}
\label{\detokenize{delete_bonds:examples}}
\begin{sphinxVerbatim}[commandchars=\\\{\}]
\PYG{k}{delete\PYGZus{}bonds }\PYG{n+nv+nvIdentifier}{frozen}\PYG{+w}{ }\PYG{n}{multi}\PYG{+w}{ }\PYG{n}{remove}
\PYG{k}{delete\PYGZus{}bonds }\PYG{n+nv+nvIdentifier}{all}\PYG{+w}{ }\PYG{n}{atom}\PYG{+w}{ }\PYG{l+m}{4}\PYG{+w}{ }\PYG{n}{special}
\PYG{k}{delete\PYGZus{}bonds }\PYG{n+nv+nvIdentifier}{all}\PYG{+w}{ }\PYG{n}{bond}\PYG{+w}{ }\PYG{l+m}{0}\PYG{o}{*}\PYG{l+m}{3}\PYG{+w}{ }\PYG{n}{special}
\PYG{k}{delete\PYGZus{}bonds }\PYG{n+nv+nvIdentifier}{all}\PYG{+w}{ }\PYG{n}{stats}

\PYG{k}{labelmap}\PYG{+w}{ }\PYG{n}{atom}\PYG{+w}{ }\PYG{l+m}{4}\PYG{+w}{ }\PYG{n}{hc}
\PYG{k}{delete\PYGZus{}bonds }\PYG{n+nv+nvIdentifier}{all}\PYG{+w}{ }\PYG{n}{atom}\PYG{+w}{ }\PYG{n}{hc}\PYG{+w}{ }\PYG{n}{special}
\end{sphinxVerbatim}


\subsection{Description}
\label{\detokenize{delete_bonds:description}}
\sphinxAtStartPar
Turn off (or on) molecular topology interactions (i.e., bonds, angles,
dihedrals, and/or impropers).  This command is useful for deleting
interactions that have been previously turned off by bond\sphinxhyphen{}breaking
potentials.  It is also useful for turning off topology interactions
between frozen or rigid atoms.  Pairwise interactions can be turned
off via the {\hyperref[\detokenize{neigh_modify::doc}]{\sphinxcrossref{\DUrole{doc}{neigh\_modify exclude}}}} command.  The
{\hyperref[\detokenize{fix_shake::doc}]{\sphinxcrossref{\DUrole{doc}{fix shake}}}} command also effectively turns off certain
bond and angle interactions.

\sphinxAtStartPar
For all styles, by default, an interaction is only turned off (or on)
if all the atoms involved are in the specified group.  See the \sphinxstyleemphasis{any}
keyword to change the behavior.

\sphinxAtStartPar
Several of the styles (\sphinxstyleemphasis{atom}, \sphinxstyleemphasis{bond}, \sphinxstyleemphasis{angle}, \sphinxstyleemphasis{dihedral}, \sphinxstyleemphasis{improper})
take a \sphinxstyleemphasis{type} as an argument.  The specified \sphinxstyleemphasis{type} can be a
{\hyperref[\detokenize{Howto_type_labels::doc}]{\sphinxcrossref{\DUrole{doc}{type label}}}}.  Otherwise, the type should be an
integer from 0 to \(N\), where \(N\) is the number of relevant
types (atom types, bond types, etc.).  A value of 0 is only relevant for
style \sphinxstyleemphasis{bond}; see details below.  For numeric types, a wildcard asterisk
can be used in place of or in conjunction with the \sphinxstyleemphasis{type} argument to
specify a range of types.  This takes the form “*” or “*n” or “m*” or
“m*n”.  If \(N\) is the number of types, then an asterisk with no
numeric values means all types from 0 to \(N\).  A leading asterisk
means all types from 0 to n (inclusive).  A trailing asterisk means all
types from m to N (inclusive).  A middle asterisk means all types from m to
n (inclusive).  Note that it is fine to include a type of 0 for non\sphinxhyphen{}bond
styles; it will simply be ignored.

\sphinxAtStartPar
For style \sphinxstyleemphasis{multi} all bond, angle, dihedral, and improper interactions
of any type, involving atoms in the group, are turned off.

\sphinxAtStartPar
Style \sphinxstyleemphasis{atom} is the same as style \sphinxstyleemphasis{multi} except that in addition, one
or more of the atoms involved in the bond, angle, dihedral, or
improper interaction must also be of the specified atom type.

\sphinxAtStartPar
For style \sphinxstyleemphasis{bond}, only bonds are candidates for turn\sphinxhyphen{}off, and the bond
must also be of the specified type.  Styles \sphinxstyleemphasis{angle}, \sphinxstyleemphasis{dihedral}, and
\sphinxstyleemphasis{improper} are treated similarly.

\sphinxAtStartPar
For style \sphinxstyleemphasis{bond}, you can set the type to 0 to delete bonds that have
been previously broken by a bond\sphinxhyphen{}breaking potential (which sets the
bond type to 0 when a bond is broken); for example, see the
{\hyperref[\detokenize{bond_style::doc}]{\sphinxcrossref{\DUrole{doc}{bond\_style quartic}}}} command.

\sphinxAtStartPar
For style \sphinxstyleemphasis{stats} no interactions are turned off (or on); the status
of all interactions in the specified group is simply reported.  This
is useful for diagnostic purposes if bonds have been turned off by a
bond\sphinxhyphen{}breaking potential during a previous run.

\sphinxAtStartPar
The default behavior of the delete\_bonds command is to turn off
interactions by toggling their type to a negative value, but not to
permanently remove the interaction.  For example, a bond\_type of 2 is set to
\(-2.\)  The neighbor list creation routines will not include such an
interaction in their interaction lists.  The default is also to not
alter the list of 1\textendash{}2, 1\textendash{}3, or 1\textendash{}4 neighbors computed by the
{\hyperref[\detokenize{special_bonds::doc}]{\sphinxcrossref{\DUrole{doc}{special\_bonds}}}} command and used to weight pairwise
force and energy calculations.  This means that pairwise computations
will proceed as if the bond (or angle, etc.) were still turned on.

\sphinxAtStartPar
Several keywords can be appended to the argument list to alter the
default behaviors.

\sphinxAtStartPar
The \sphinxstyleemphasis{any} keyword changes the requirement that all atoms in the bond
(angle, etc.) must be in the specified group in order to turn off the
interaction.  Instead, if any of the atoms in the interaction are in
the specified group, it will be turned off (or on if the \sphinxstyleemphasis{undo}
keyword is used).

\sphinxAtStartPar
The \sphinxstyleemphasis{undo} keyword inverts the delete\_bonds command so that the
specified bonds, angles, etc. are turned on if they are currently
turned off.  This means a negative value is toggled to positive.  For
example, for style \sphinxstyleemphasis{angle}, if \sphinxstyleemphasis{type} is specified as 2, then all
angles with current type = \(-2\) are reset to type = \(2\).
Note that the {\hyperref[\detokenize{fix_shake::doc}]{\sphinxcrossref{\DUrole{doc}{fix shake}}}} command also sets bond and angle
types negative, so this option should not be used on those interactions.

\sphinxAtStartPar
The \sphinxstyleemphasis{remove} keyword is invoked at the end of the delete\_bonds
operation.  It causes turned\sphinxhyphen{}off bonds (angles, etc.) to be removed
from each atom’s data structure and then adjusts the global bond
(angle, etc.) counts accordingly.  Removal is a permanent change;
removed bonds cannot be turned back on via the \sphinxstyleemphasis{undo} keyword.
Removal does not alter the pairwise 1\textendash{}2, 1\textendash{}3, or 1\textendash{}4 weighting list.

\sphinxAtStartPar
The \sphinxstyleemphasis{special} keyword is invoked at the end of the delete\_bonds
operation, after (optional) removal.  It re\sphinxhyphen{}computes the pairwise 1\textendash{}2,
1\textendash{}3, 1\textendash{}4 weighting list.  The weighting list computation treats
turned\sphinxhyphen{}off bonds the same as turned\sphinxhyphen{}on.  Thus, turned\sphinxhyphen{}off bonds must
be removed if you wish to change the weighting list.

\sphinxAtStartPar
Note that the choice of \sphinxstyleemphasis{remove} and \sphinxstyleemphasis{special} options affects how
1\textendash{}2, 1\textendash{}3, 1\textendash{}4 pairwise interactions will be computed across bonds that
have been modified by the delete\_bonds command.


\subsection{Restrictions}
\label{\detokenize{delete_bonds:restrictions}}
\sphinxAtStartPar
This command requires inter\sphinxhyphen{}processor communication to acquire ghost
atoms, to coordinate the deleting of bonds, angles, etc. between atoms
shared by multiple processors.  This means that your system must be
ready to perform a simulation before using this command (force fields
setup, atom masses set, etc.).  Just as would be needed to run
dynamics, the force field you define should define a cutoff
(e.g., through a {\hyperref[\detokenize{pair_style::doc}]{\sphinxcrossref{\DUrole{doc}{pair\_style}}}} command) which is long
enough for a processor to acquire the ghost atoms its needs to compute
bond, angle, etc. interactions.

\sphinxAtStartPar
If deleted bonds (or angles, etc.) are removed but the 1\textendash{}2, 1\textendash{}3, and 1\textendash{}4
weighting list is not recomputed, this can cause a later
{\hyperref[\detokenize{fix_shake::doc}]{\sphinxcrossref{\DUrole{doc}{fix shake}}}} command to fail due to an atom’s bonds being
inconsistent with the weighting list.  This should only happen if the
group used in the fix command includes both atoms in the bond, in
which case you probably should be recomputing the weighting list.


\subsection{Related commands}
\label{\detokenize{delete_bonds:related-commands}}
\sphinxAtStartPar
{\hyperref[\detokenize{neigh_modify::doc}]{\sphinxcrossref{\DUrole{doc}{neigh\_modify}}}} exclude,
{\hyperref[\detokenize{special_bonds::doc}]{\sphinxcrossref{\DUrole{doc}{special\_bonds}}}}, {\hyperref[\detokenize{fix_shake::doc}]{\sphinxcrossref{\DUrole{doc}{fix shake}}}}


\subsection{Default}
\label{\detokenize{delete_bonds:default}}
\sphinxAtStartPar
none

\sphinxstepscope

\index{dielectric@\spxentry{dielectric}}\ignorespaces 

\section{dielectric command}
\label{\detokenize{dielectric:dielectric-command}}\label{\detokenize{dielectric:index-0}}\label{\detokenize{dielectric::doc}}

\subsection{Syntax}
\label{\detokenize{dielectric:syntax}}
\begin{sphinxVerbatim}[commandchars=\\\{\}]
\PYG{k}{dielectric}\PYG{+w}{ }\PYG{n}{value}
\end{sphinxVerbatim}
\begin{itemize}
\item {} 
\sphinxAtStartPar
value = dielectric constant

\end{itemize}


\subsection{Examples}
\label{\detokenize{dielectric:examples}}
\begin{sphinxVerbatim}[commandchars=\\\{\}]
\PYG{k}{dielectric}\PYG{+w}{ }\PYG{l+m}{2.0}
\end{sphinxVerbatim}


\subsection{Description}
\label{\detokenize{dielectric:description}}
\sphinxAtStartPar
Set the dielectric constant for Coulombic interactions (pairwise and
long\sphinxhyphen{}range) to this value.  The constant is unitless, since it is used
to reduce the strength of the interactions.  The value is used in the
denominator of the formulas for Coulombic interactions (e.g., a value
of 4.0 reduces the Coulombic interactions to 25\% of their default
strength).  See the {\hyperref[\detokenize{pair_style::doc}]{\sphinxcrossref{\DUrole{doc}{pair\_style}}}} command for more
details.


\subsection{Restrictions}
\label{\detokenize{dielectric:restrictions}}\begin{quote}

\sphinxAtStartPar
none
\end{quote}


\subsection{Related commands}
\label{\detokenize{dielectric:related-commands}}
\sphinxAtStartPar
{\hyperref[\detokenize{pair_style::doc}]{\sphinxcrossref{\DUrole{doc}{pair\_style}}}}


\subsection{Default}
\label{\detokenize{dielectric:default}}
\begin{sphinxVerbatim}[commandchars=\\\{\}]
\PYG{k}{dielectric}\PYG{+w}{ }\PYG{l+m}{1.0}
\end{sphinxVerbatim}

\sphinxstepscope

\index{dihedral\_coeff@\spxentry{dihedral\_coeff}}\ignorespaces 

\section{dihedral\_coeff command}
\label{\detokenize{dihedral_coeff:dihedral-coeff-command}}\label{\detokenize{dihedral_coeff:index-0}}\label{\detokenize{dihedral_coeff::doc}}

\subsection{Syntax}
\label{\detokenize{dihedral_coeff:syntax}}
\begin{sphinxVerbatim}[commandchars=\\\{\}]
\PYG{k}{dihedral\PYGZus{}coeff}\PYG{+w}{ }\PYG{n}{N}\PYG{+w}{ }\PYG{n}{args}
\end{sphinxVerbatim}
\begin{itemize}
\item {} 
\sphinxAtStartPar
N = numeric dihedral type (see asterisk form below) or alphanumeric type label

\item {} 
\sphinxAtStartPar
args = coefficients for one or more dihedral types

\end{itemize}


\subsection{Examples}
\label{\detokenize{dihedral_coeff:examples}}
\begin{sphinxVerbatim}[commandchars=\\\{\}]
\PYG{k}{dihedral\PYGZus{}coeff}\PYG{+w}{ }\PYG{l+m}{1}\PYG{+w}{ }\PYG{l+m}{80.0}\PYG{+w}{ }\PYG{l+m}{1}\PYG{+w}{ }\PYG{l+m}{3}
\PYG{k}{dihedral\PYGZus{}coeff}\PYG{+w}{ }\PYG{o}{*}\PYG{+w}{ }\PYG{l+m}{80.0}\PYG{+w}{ }\PYG{l+m}{1}\PYG{+w}{ }\PYG{l+m}{3}\PYG{+w}{ }\PYG{l+m}{0.5}
\PYG{k}{dihedral\PYGZus{}coeff}\PYG{+w}{ }\PYG{l+m}{2}\PYG{o}{*}\PYG{+w}{ }\PYG{l+m}{80.0}\PYG{+w}{ }\PYG{l+m}{1}\PYG{+w}{ }\PYG{l+m}{3}\PYG{+w}{ }\PYG{l+m}{0.5}

\PYG{k}{labelmap}\PYG{+w}{ }\PYG{n}{dihedral}\PYG{+w}{ }\PYG{l+m}{1}\PYG{+w}{ }\PYG{n}{backbone}
\PYG{k}{dihedral\PYGZus{}coeff}\PYG{+w}{ }\PYG{n}{backbone}\PYG{+w}{ }\PYG{l+m}{80.0}\PYG{+w}{ }\PYG{l+m}{1}\PYG{+w}{ }\PYG{l+m}{3}
\end{sphinxVerbatim}


\subsection{Description}
\label{\detokenize{dihedral_coeff:description}}
\sphinxAtStartPar
Specify the dihedral force field coefficients for one or more dihedral
types.  The number and meaning of the coefficients depends on the
dihedral style.  Dihedral coefficients can also be set in the data file
read by the {\hyperref[\detokenize{read_data::doc}]{\sphinxcrossref{\DUrole{doc}{read\_data}}}} command or in a restart file.

\sphinxAtStartPar
\(N\) can be specified in one of two ways.  An explicit numeric
value can be used, as in the first example above.  Or \(N\) can be
an alphanumeric type label, which is a string defined by the
{\hyperref[\detokenize{labelmap::doc}]{\sphinxcrossref{\DUrole{doc}{labelmap}}}} command or in a corresponding section of a
data file read by the {\hyperref[\detokenize{read_data::doc}]{\sphinxcrossref{\DUrole{doc}{read\_data}}}} command.

\sphinxAtStartPar
For numeric values only, a wild\sphinxhyphen{}card asterisk can be used to set the
coefficients for multiple dihedral types.  This takes the form “*” or
“*n” or “n*” or “m*n”.  If \(N\) is the number of dihedral types,
then an asterisk with no numeric values means all types from 1 to
\(N\).  A leading asterisk means all types from 1 to n (inclusive).
A trailing asterisk means all types from n to \(N\) (inclusive).  A
middle asterisk means all types from m to n (inclusive).

\sphinxAtStartPar
Note that using a dihedral\_coeff command can override a previous setting
for the same dihedral type.  For example, these commands set the coeffs
for all dihedral types, then overwrite the coeffs for just dihedral type
2:

\begin{sphinxVerbatim}[commandchars=\\\{\}]
\PYG{k}{dihedral\PYGZus{}coeff}\PYG{+w}{ }\PYG{o}{*}\PYG{+w}{ }\PYG{l+m}{80.0}\PYG{+w}{ }\PYG{l+m}{1}\PYG{+w}{ }\PYG{l+m}{3}
\PYG{k}{dihedral\PYGZus{}coeff}\PYG{+w}{ }\PYG{l+m}{2}\PYG{+w}{ }\PYG{l+m}{200.0}\PYG{+w}{ }\PYG{l+m}{1}\PYG{+w}{ }\PYG{l+m}{3}
\end{sphinxVerbatim}

\sphinxAtStartPar
A line in a data file that specifies dihedral coefficients uses the exact
same format as the arguments of the dihedral\_coeff command in an input
script, except that wild\sphinxhyphen{}card asterisks should not be used since
coefficients for all \(N\) types must be listed in the file.  For example,
under the “Dihedral Coeffs” section of a data file, the line that
corresponds to the first example above would be listed as

\begin{sphinxVerbatim}[commandchars=\\\{\}]
1 80.0 1 3
\end{sphinxVerbatim}

\sphinxAtStartPar
The {\hyperref[\detokenize{dihedral_class2::doc}]{\sphinxcrossref{\DUrole{doc}{dihedral\_style class2}}}} is an exception to
this rule, in that an additional argument is used in the input script
to allow specification of the cross\sphinxhyphen{}term coefficients.  See its doc
page for details.

\begin{sphinxadmonition}{note}{Note:}
\sphinxAtStartPar
When comparing the formulas and coefficients for various LAMMPS
dihedral styles with dihedral equations defined by other force fields,
note that some force field implementations divide/multiply the energy
prefactor \sphinxstyleemphasis{K} by the multiple number of torsions that contain the
\sphinxstyleemphasis{J}\textendash{}\sphinxstyleemphasis{K} bond in an \sphinxstyleemphasis{I}\sphinxhyphen{}\sphinxstyleemphasis{J}\sphinxhyphen{}\sphinxstyleemphasis{K}\sphinxhyphen{}\sphinxstyleemphasis{L} torsion.  LAMMPS does
not do this (i.e., the listed dihedral equation applies to each individual
dihedral).  Thus, you need to define \sphinxstyleemphasis{K} appropriately to account for this
difference, if necessary.
\end{sphinxadmonition}


\bigskip\hrule\bigskip


\sphinxAtStartPar
The list of all dihedral styles defined in LAMMPS is given on the
{\hyperref[\detokenize{dihedral_style::doc}]{\sphinxcrossref{\DUrole{doc}{dihedral\_style}}}} doc page.  They are also listed
in more compact form on the {\hyperref[\detokenize{Commands_bond:dihedral}]{\sphinxcrossref{\DUrole{std,std-ref}{Commands dihedral}}}} doc page.

\sphinxAtStartPar
On either of those pages, click on the style to display the formula it
computes and its coefficients as specified by the associated
dihedral\_coeff command.


\bigskip\hrule\bigskip



\subsection{Restrictions}
\label{\detokenize{dihedral_coeff:restrictions}}
\sphinxAtStartPar
This command must come after the simulation box is defined by a
{\hyperref[\detokenize{read_data::doc}]{\sphinxcrossref{\DUrole{doc}{read\_data}}}}, {\hyperref[\detokenize{read_restart::doc}]{\sphinxcrossref{\DUrole{doc}{read\_restart}}}}, or
{\hyperref[\detokenize{create_box::doc}]{\sphinxcrossref{\DUrole{doc}{create\_box}}}} command.

\sphinxAtStartPar
A dihedral style must be defined before any dihedral coefficients are
set, either in the input script or in a data file.


\subsection{Related commands}
\label{\detokenize{dihedral_coeff:related-commands}}
\sphinxAtStartPar
{\hyperref[\detokenize{dihedral_style::doc}]{\sphinxcrossref{\DUrole{doc}{dihedral\_style}}}}


\subsection{Default}
\label{\detokenize{dihedral_coeff:default}}
\sphinxAtStartPar
none

\sphinxstepscope

\index{dihedral\_style@\spxentry{dihedral\_style}}\ignorespaces 

\section{dihedral\_style command}
\label{\detokenize{dihedral_style:dihedral-style-command}}\label{\detokenize{dihedral_style:index-0}}\label{\detokenize{dihedral_style::doc}}

\subsection{Syntax}
\label{\detokenize{dihedral_style:syntax}}
\begin{sphinxVerbatim}[commandchars=\\\{\}]
\PYG{k}{dihedral\PYGZus{}style}\PYG{+w}{ }\PYG{n}{style}
\end{sphinxVerbatim}
\begin{itemize}
\item {} 
\sphinxAtStartPar
style = \sphinxstyleemphasis{none} or \sphinxstyleemphasis{zero} or \sphinxstyleemphasis{hybrid} or \sphinxstyleemphasis{charmm} or \sphinxstyleemphasis{charmmfsw} or \sphinxstyleemphasis{class2} or \sphinxstyleemphasis{cosine/shift/exp} or \sphinxstyleemphasis{cosine/squared/restricted} or \sphinxstyleemphasis{fourier} or \sphinxstyleemphasis{harmonic} or \sphinxstyleemphasis{helix} or \sphinxstyleemphasis{lepton} or \sphinxstyleemphasis{multi/harmonic} or \sphinxstyleemphasis{nharmonic} or \sphinxstyleemphasis{opls} or \sphinxstyleemphasis{spherical} or \sphinxstyleemphasis{table} or \sphinxstyleemphasis{table/cut}

\end{itemize}


\subsection{Examples}
\label{\detokenize{dihedral_style:examples}}
\begin{sphinxVerbatim}[commandchars=\\\{\}]
\PYG{k}{dihedral\PYGZus{}style}\PYG{+w}{ }\PYG{n}{harmonic}
\PYG{k}{dihedral\PYGZus{}style}\PYG{+w}{ }\PYG{n}{multi}\PYG{o}{/}\PYG{n}{harmonic}
\PYG{k}{dihedral\PYGZus{}style}\PYG{+w}{ }\PYG{n}{hybrid}\PYG{+w}{ }\PYG{n}{harmonic}\PYG{+w}{ }\PYG{n}{charmm}
\end{sphinxVerbatim}


\subsection{Description}
\label{\detokenize{dihedral_style:description}}
\sphinxAtStartPar
Set the formula(s) LAMMPS uses to compute dihedral interactions
between quadruplets of atoms, which remain in force for the duration
of the simulation.  The list of dihedral quadruplets is read in by a
{\hyperref[\detokenize{read_data::doc}]{\sphinxcrossref{\DUrole{doc}{read\_data}}}} or {\hyperref[\detokenize{read_restart::doc}]{\sphinxcrossref{\DUrole{doc}{read\_restart}}}} command
from a data or restart file.

\sphinxAtStartPar
Hybrid models where dihedrals are computed using different dihedral
potentials can be setup using the \sphinxstyleemphasis{hybrid} dihedral style.

\sphinxAtStartPar
The coefficients associated with a dihedral style can be specified in
a data or restart file or via the {\hyperref[\detokenize{dihedral_coeff::doc}]{\sphinxcrossref{\DUrole{doc}{dihedral\_coeff}}}}
command.

\sphinxAtStartPar
All dihedral potentials store their coefficient data in binary restart
files which means dihedral\_style and
{\hyperref[\detokenize{dihedral_coeff::doc}]{\sphinxcrossref{\DUrole{doc}{dihedral\_coeff}}}} commands do not need to be
re\sphinxhyphen{}specified in an input script that restarts a simulation.  See the
{\hyperref[\detokenize{read_restart::doc}]{\sphinxcrossref{\DUrole{doc}{read\_restart}}}} command for details on how to do
this.  The one exception is that dihedral\_style \sphinxstyleemphasis{hybrid} only stores
the list of sub\sphinxhyphen{}styles in the restart file; dihedral coefficients need
to be re\sphinxhyphen{}specified.

\begin{sphinxadmonition}{note}{Note:}
\sphinxAtStartPar
When both a dihedral and pair style is defined, the
{\hyperref[\detokenize{special_bonds::doc}]{\sphinxcrossref{\DUrole{doc}{special\_bonds}}}} command often needs to be used to
turn off (or weight) the pairwise interaction that would otherwise
exist between four bonded atoms.
\end{sphinxadmonition}

\sphinxAtStartPar
In the formulas listed for each dihedral style, \sphinxstyleemphasis{phi} is the torsional
angle defined by the quadruplet of atoms.  This angle has a sign
convention as shown in this diagram:

\noindent{\hspace*{\fill}\sphinxincludegraphics{{dihedral_sign}.jpg}\hspace*{\fill}}

\sphinxAtStartPar
where the \(I,J,K,L\) ordering of the four atoms that define the dihedral
is from left to right.

\sphinxAtStartPar
This sign convention effects several of the dihedral styles listed
below (e.g., charmm, helix) in the sense that the energy formula
depends on the sign of phi, which may be reflected in the value of the
coefficients you specify.

\begin{sphinxadmonition}{note}{Note:}
\sphinxAtStartPar
When comparing the formulas and coefficients for various LAMMPS
dihedral styles with dihedral equations defined by other force fields,
note that some force field implementations divide/multiply the energy
prefactor \sphinxstyleemphasis{K} by the multiple number of torsions that contain the
\sphinxstyleemphasis{J}\textendash{}\sphinxstyleemphasis{K} bond in an \sphinxstyleemphasis{I}\textendash{}\sphinxstyleemphasis{J}\textendash{}\sphinxstyleemphasis{K}\textendash{}\sphinxstyleemphasis{L} torsion.  LAMMPS does
not do this (i.e., the listed dihedral equation applies to each individual
dihedral).  Thus, you need to define \sphinxstyleemphasis{K} appropriately via the
{\hyperref[\detokenize{dihedral_coeff::doc}]{\sphinxcrossref{\DUrole{doc}{dihedral\_coeff}}}} command to account for this
difference if necessary.
\end{sphinxadmonition}


\bigskip\hrule\bigskip


\sphinxAtStartPar
Here is an alphabetic list of dihedral styles defined in LAMMPS.  Click on
the style to display the formula it computes and coefficients
specified by the associated {\hyperref[\detokenize{dihedral_coeff::doc}]{\sphinxcrossref{\DUrole{doc}{dihedral\_coeff}}}} command.

\sphinxAtStartPar
Click on the style to display the formula it computes, any additional
arguments specified in the dihedral\_style command, and coefficients
specified by the associated {\hyperref[\detokenize{dihedral_coeff::doc}]{\sphinxcrossref{\DUrole{doc}{dihedral\_coeff}}}}
command.

\sphinxAtStartPar
There are also additional accelerated pair styles included in the
LAMMPS distribution for faster performance on CPUs, GPUs, and KNLs.
The individual style names on the {\hyperref[\detokenize{Commands_bond:dihedral}]{\sphinxcrossref{\DUrole{std,std-ref}{Commands dihedral}}}} page are
followed by one or more of (g,i,k,o,t) to indicate which accelerated styles
exist.
\begin{itemize}
\item {} 
\sphinxAtStartPar
{\hyperref[\detokenize{dihedral_none::doc}]{\sphinxcrossref{\DUrole{doc}{none}}}} \sphinxhyphen{} turn off dihedral interactions

\item {} 
\sphinxAtStartPar
{\hyperref[\detokenize{dihedral_zero::doc}]{\sphinxcrossref{\DUrole{doc}{zero}}}} \sphinxhyphen{} topology but no interactions

\item {} 
\sphinxAtStartPar
{\hyperref[\detokenize{dihedral_hybrid::doc}]{\sphinxcrossref{\DUrole{doc}{hybrid}}}} \sphinxhyphen{} define multiple styles of dihedral interactions

\item {} 
\sphinxAtStartPar
{\hyperref[\detokenize{dihedral_charmm::doc}]{\sphinxcrossref{\DUrole{doc}{charmm}}}} \sphinxhyphen{} CHARMM dihedral

\item {} 
\sphinxAtStartPar
{\hyperref[\detokenize{dihedral_charmm::doc}]{\sphinxcrossref{\DUrole{doc}{charmmfsw}}}} \sphinxhyphen{} CHARMM dihedral with force switching

\item {} 
\sphinxAtStartPar
{\hyperref[\detokenize{dihedral_class2::doc}]{\sphinxcrossref{\DUrole{doc}{class2}}}} \sphinxhyphen{} COMPASS (class 2) dihedral

\item {} 
\sphinxAtStartPar
{\hyperref[\detokenize{dihedral_cosine_shift_exp::doc}]{\sphinxcrossref{\DUrole{doc}{cosine/shift/exp}}}} \sphinxhyphen{} dihedral with exponential in spring constant

\item {} 
\sphinxAtStartPar
{\hyperref[\detokenize{dihedral_cosine_squared_restricted::doc}]{\sphinxcrossref{\DUrole{doc}{cosine/squared/restricted}}}} \sphinxhyphen{} squared cosine dihedral with restricted term

\item {} 
\sphinxAtStartPar
{\hyperref[\detokenize{dihedral_fourier::doc}]{\sphinxcrossref{\DUrole{doc}{fourier}}}} \sphinxhyphen{} dihedral with multiple cosine terms

\item {} 
\sphinxAtStartPar
{\hyperref[\detokenize{dihedral_harmonic::doc}]{\sphinxcrossref{\DUrole{doc}{harmonic}}}} \sphinxhyphen{} harmonic dihedral

\item {} 
\sphinxAtStartPar
{\hyperref[\detokenize{dihedral_helix::doc}]{\sphinxcrossref{\DUrole{doc}{helix}}}} \sphinxhyphen{} helix dihedral

\item {} 
\sphinxAtStartPar
{\hyperref[\detokenize{dihedral_lepton::doc}]{\sphinxcrossref{\DUrole{doc}{lepton}}}} \sphinxhyphen{} dihedral potential from evaluating a string

\item {} 
\sphinxAtStartPar
{\hyperref[\detokenize{dihedral_multi_harmonic::doc}]{\sphinxcrossref{\DUrole{doc}{multi/harmonic}}}} \sphinxhyphen{} dihedral with 5 harmonic terms

\item {} 
\sphinxAtStartPar
{\hyperref[\detokenize{dihedral_nharmonic::doc}]{\sphinxcrossref{\DUrole{doc}{nharmonic}}}} \sphinxhyphen{} same as multi\sphinxhyphen{}harmonic with N terms

\item {} 
\sphinxAtStartPar
{\hyperref[\detokenize{dihedral_opls::doc}]{\sphinxcrossref{\DUrole{doc}{opls}}}} \sphinxhyphen{} OPLS dihedral

\item {} 
\sphinxAtStartPar
{\hyperref[\detokenize{dihedral_quadratic::doc}]{\sphinxcrossref{\DUrole{doc}{quadratic}}}} \sphinxhyphen{} dihedral with quadratic term in angle

\item {} 
\sphinxAtStartPar
{\hyperref[\detokenize{dihedral_spherical::doc}]{\sphinxcrossref{\DUrole{doc}{spherical}}}} \sphinxhyphen{} dihedral which includes angle terms to avoid singularities

\item {} 
\sphinxAtStartPar
{\hyperref[\detokenize{dihedral_table::doc}]{\sphinxcrossref{\DUrole{doc}{table}}}} \sphinxhyphen{} tabulated dihedral

\item {} 
\sphinxAtStartPar
{\hyperref[\detokenize{dihedral_table::doc}]{\sphinxcrossref{\DUrole{doc}{table/cut}}}} \sphinxhyphen{} tabulated dihedral with analytic cutoff

\end{itemize}


\bigskip\hrule\bigskip



\subsection{Restrictions}
\label{\detokenize{dihedral_style:restrictions}}
\sphinxAtStartPar
Dihedral styles can only be set for atom styles that allow dihedrals
to be defined.

\sphinxAtStartPar
Most dihedral styles are part of the MOLECULE package.  They are only
enabled if LAMMPS was built with that package.  See the {\hyperref[\detokenize{Build_package::doc}]{\sphinxcrossref{\DUrole{doc}{Build package}}}} page for more info.  The doc pages for
individual dihedral potentials tell if it is part of a package.


\subsection{Related commands}
\label{\detokenize{dihedral_style:related-commands}}
\sphinxAtStartPar
{\hyperref[\detokenize{dihedral_coeff::doc}]{\sphinxcrossref{\DUrole{doc}{dihedral\_coeff}}}}


\subsection{Default}
\label{\detokenize{dihedral_style:default}}
\sphinxAtStartPar
dihedral\_style none

\sphinxstepscope

\index{dihedral\_write@\spxentry{dihedral\_write}}\ignorespaces 

\section{dihedral\_write command}
\label{\detokenize{dihedral_write:dihedral-write-command}}\label{\detokenize{dihedral_write:index-0}}\label{\detokenize{dihedral_write::doc}}

\subsection{Syntax}
\label{\detokenize{dihedral_write:syntax}}
\begin{sphinxVerbatim}[commandchars=\\\{\}]
\PYG{k}{dihedral\PYGZus{}write}\PYG{+w}{ }\PYG{n}{dtype}\PYG{+w}{ }\PYG{n}{N}\PYG{+w}{ }\PYG{n}{file}\PYG{+w}{ }\PYG{n}{keyword}
\end{sphinxVerbatim}
\begin{itemize}
\item {} 
\sphinxAtStartPar
dtype = dihedral type

\item {} 
\sphinxAtStartPar
N = \# of values

\item {} 
\sphinxAtStartPar
file = name of file to write values to

\item {} 
\sphinxAtStartPar
keyword = section name in file for this set of tabulated values

\end{itemize}


\subsection{Examples}
\label{\detokenize{dihedral_write:examples}}
\begin{sphinxVerbatim}[commandchars=\\\{\}]
\PYG{k}{dihedral\PYGZus{}write}\PYG{+w}{ }\PYG{l+m}{1}\PYG{+w}{  }\PYG{l+m}{500}\PYG{+w}{ }\PYG{n}{table.txt}\PYG{+w}{ }\PYG{n}{Harmonic\PYGZus{}1}
\PYG{k}{dihedral\PYGZus{}write}\PYG{+w}{ }\PYG{l+m}{3}\PYG{+w}{ }\PYG{l+m}{1000}\PYG{+w}{ }\PYG{n}{table.txt}\PYG{+w}{ }\PYG{n}{Harmonic\PYGZus{}3}
\end{sphinxVerbatim}


\subsection{Description}
\label{\detokenize{dihedral_write:description}}
\sphinxAtStartPar
\DUrole{versionmodified,added}{Added in version 8Feb2023.}

\sphinxAtStartPar
Write energy and force values to a file as a function of the dihedral
angle for the currently defined dihedral potential.  Force in this
context means the force with respect to the dihedral angle, not the
force on individual atoms.  This is useful for plotting the potential
function or otherwise debugging its values.  The resulting file can also
be used as input for use with {\hyperref[\detokenize{dihedral_table::doc}]{\sphinxcrossref{\DUrole{doc}{dihedral style table}}}}.

\sphinxAtStartPar
If the file already exists, the table of values is appended to the end
of the file to allow multiple tables of energy and force to be included
in one file.  The individual sections may be identified by the \sphinxstyleemphasis{keyword}.

\sphinxAtStartPar
The energy and force values are computed for dihedrals ranging from 0
degrees to 360 degrees for 4 interacting atoms forming an dihedral type
dtype, using the appropriate {\hyperref[\detokenize{dihedral_coeff::doc}]{\sphinxcrossref{\DUrole{doc}{dihedral\_coeff}}}}
coefficients. N evenly spaced dihedrals are used.  Since 0 and 360
degrees are the same dihedral angle, the latter entry is skipped.

\sphinxAtStartPar
For example, for N = 6, values would be computed at
\(\phi = 0, 60, 120, 180, 240, 300\).

\sphinxAtStartPar
The file is written in the format used as input for the
{\hyperref[\detokenize{dihedral_table::doc}]{\sphinxcrossref{\DUrole{doc}{dihedral\_style table}}}} option with \sphinxstyleemphasis{keyword} as
the section name.  Each line written to the file lists an index number
(1\sphinxhyphen{}N), an dihedral angle (in degrees), an energy (in energy units), and
a force (in force units per radians\textasciicircum{}2).  In case a new file is created,
the first line will be a comment with a “DATE:” and “UNITS:” tag with
the current date and {\hyperref[\detokenize{units::doc}]{\sphinxcrossref{\DUrole{doc}{units}}}} settings.  For subsequent
invocations of the \sphinxstyleemphasis{dihedral\_write} command for the same file, data will
be appended and the current units settings will be compared to the data
from the header, if present. The \sphinxstyleemphasis{dihedral\_write} will refuse to add a
table to an existing file if the units are not the same.


\subsection{Restrictions}
\label{\detokenize{dihedral_write:restrictions}}
\sphinxAtStartPar
All force field coefficients for dihedrals and other kinds of interactions
must be set before this command can be invoked.

\sphinxAtStartPar
The table of the dihedral energy and force data data is created by using a
separate, internally created, new LAMMPS instance with a dummy system of
4 atoms for which the dihedral potential energy is computed after
transferring the dihedral style and coefficients and arranging the 4 atoms
into the corresponding geometries.  The dihedral force is then determined
from the potential energies through numerical differentiation.  As a
consequence of this approach, not all dihedral styles are compatible. The
following conditions must be met:
\begin{itemize}
\item {} 
\sphinxAtStartPar
The dihedral style must be able to write its coefficients to a data file.
This condition excludes for example {\hyperref[\detokenize{dihedral_hybrid::doc}]{\sphinxcrossref{\DUrole{doc}{dihedral style hybrid}}}} and
{\hyperref[\detokenize{dihedral_table::doc}]{\sphinxcrossref{\DUrole{doc}{dihedral style table}}}}.

\item {} 
\sphinxAtStartPar
The potential function must not have any terms that depend on geometry
properties other than the dihedral.  This condition excludes for
example {\hyperref[\detokenize{dihedral_class2::doc}]{\sphinxcrossref{\DUrole{doc}{dihedral style class2}}}}.  Please note
that the \sphinxstyleemphasis{write\_dihedral} command has no way of checking for this
condition.  It will check the style name against an internal list of
known to be incompatible styles.  The resulting tables may be bogus
for unlisted dihedral styles if the requirement is not met.  It is
thus recommended to make careful tests for any created tables.

\end{itemize}


\subsection{Related commands}
\label{\detokenize{dihedral_write:related-commands}}
\sphinxAtStartPar
{\hyperref[\detokenize{dihedral_table::doc}]{\sphinxcrossref{\DUrole{doc}{dihedral\_style table}}}}, {\hyperref[\detokenize{bond_write::doc}]{\sphinxcrossref{\DUrole{doc}{bond\_write}}}},
{\hyperref[\detokenize{angle_write::doc}]{\sphinxcrossref{\DUrole{doc}{angle\_write}}}}, {\hyperref[\detokenize{dihedral_style::doc}]{\sphinxcrossref{\DUrole{doc}{dihedral\_style}}}},
{\hyperref[\detokenize{dihedral_coeff::doc}]{\sphinxcrossref{\DUrole{doc}{dihedral\_coeff}}}}


\subsection{Default}
\label{\detokenize{dihedral_write:default}}
\sphinxAtStartPar
none

\sphinxstepscope

\index{dimension@\spxentry{dimension}}\ignorespaces 

\section{dimension command}
\label{\detokenize{dimension:dimension-command}}\label{\detokenize{dimension:index-0}}\label{\detokenize{dimension::doc}}

\subsection{Syntax}
\label{\detokenize{dimension:syntax}}
\begin{sphinxVerbatim}[commandchars=\\\{\}]
\PYG{k}{dimension}\PYG{+w}{ }\PYG{n}{N}
\end{sphinxVerbatim}
\begin{itemize}
\item {} 
\sphinxAtStartPar
N = 2 or 3

\end{itemize}


\subsection{Examples}
\label{\detokenize{dimension:examples}}
\begin{sphinxVerbatim}[commandchars=\\\{\}]
\PYG{k}{dimension}\PYG{+w}{ }\PYG{l+m}{2}
\end{sphinxVerbatim}


\subsection{Description}
\label{\detokenize{dimension:description}}
\sphinxAtStartPar
Set the dimensionality of the simulation.  By default LAMMPS runs 3d
simulations.  To run a 2d simulation, this command should be used
prior to setting up a simulation box via the
{\hyperref[\detokenize{create_box::doc}]{\sphinxcrossref{\DUrole{doc}{create\_box}}}} or {\hyperref[\detokenize{read_data::doc}]{\sphinxcrossref{\DUrole{doc}{read\_data}}}} commands.
Restart files also store this setting.

\sphinxAtStartPar
See the discussion on the {\hyperref[\detokenize{Howto_2d::doc}]{\sphinxcrossref{\DUrole{doc}{Howto 2d}}}} page for
additional instructions on how to run 2d simulations.

\begin{sphinxadmonition}{note}{Note:}
\sphinxAtStartPar
Some models in LAMMPS treat particles as finite\sphinxhyphen{}size spheres or
ellipsoids, as opposed to point particles.  In 2d, the particles will
still be spheres or ellipsoids, not circular disks or ellipses,
meaning their moment of inertia will be the same as in 3d.
\end{sphinxadmonition}


\subsection{Restrictions}
\label{\detokenize{dimension:restrictions}}
\sphinxAtStartPar
This command must be used before the simulation box is defined by a
{\hyperref[\detokenize{read_data::doc}]{\sphinxcrossref{\DUrole{doc}{read\_data}}}} or {\hyperref[\detokenize{create_box::doc}]{\sphinxcrossref{\DUrole{doc}{create\_box}}}} command.


\subsection{Related commands}
\label{\detokenize{dimension:related-commands}}
\sphinxAtStartPar
{\hyperref[\detokenize{fix_enforce2d::doc}]{\sphinxcrossref{\DUrole{doc}{fix enforce2d}}}}


\subsection{Default}
\label{\detokenize{dimension:default}}
\begin{sphinxVerbatim}[commandchars=\\\{\}]
\PYG{k}{dimension}\PYG{+w}{ }\PYG{l+m}{3}
\end{sphinxVerbatim}

\sphinxstepscope

\index{displace\_atoms@\spxentry{displace\_atoms}}\ignorespaces 

\section{displace\_atoms command}
\label{\detokenize{displace_atoms:displace-atoms-command}}\label{\detokenize{displace_atoms:index-0}}\label{\detokenize{displace_atoms::doc}}

\subsection{Syntax}
\label{\detokenize{displace_atoms:syntax}}
\begin{sphinxVerbatim}[commandchars=\\\{\}]
\PYG{k}{displace\PYGZus{}atoms }\PYG{n+nv+nvIdentifier}{group\PYGZhy{}ID}\PYG{+w}{ }\PYG{n}{style}\PYG{+w}{ }\PYG{n}{args}\PYG{+w}{ }\PYG{n}{keyword}\PYG{+w}{ }\PYG{n}{value}\PYG{+w}{ }\PYG{n}{...}
\end{sphinxVerbatim}
\begin{itemize}
\item {} 
\sphinxAtStartPar
group\sphinxhyphen{}ID = ID of group of atoms to displace

\item {} 
\sphinxAtStartPar
style = \sphinxstyleemphasis{move} or \sphinxstyleemphasis{ramp} or \sphinxstyleemphasis{random} or \sphinxstyleemphasis{rotate}
\begin{sphinxalltt}
\sphinxstyleemphasis{move} args = delx dely delz
  delx,dely,delz = distance to displace in each dimension (distance units)
  any of delx,dely,delz can be a variable (see below)
\sphinxstyleemphasis{ramp} args = ddim dlo dhi dim clo chi
  ddim = \sphinxstyleemphasis{x} or \sphinxstyleemphasis{y} or \sphinxstyleemphasis{z}
  dlo,dhi = displacement distance between dlo and dhi (distance units)
  dim = \sphinxstyleemphasis{x} or \sphinxstyleemphasis{y} or \sphinxstyleemphasis{z}
  clo,chi = lower and upper bound of domain to displace (distance units)
\sphinxstyleemphasis{random} args = dx dy dz seed
  dx,dy,dz = random displacement magnitude in each dimension (distance units)
  seed = random \# seed (positive integer)
\sphinxstyleemphasis{rotate} args = Px Py Pz Rx Ry Rz theta
  Px,Py,Pz = origin point of axis of rotation (distance units)
  Rx,Ry,Rz = axis of rotation vector
  theta = angle of rotation (degrees)
\end{sphinxalltt}

\item {} 
\sphinxAtStartPar
zero or more keyword/value pairs may be appended
\begin{sphinxalltt}
keyword = \sphinxstyleemphasis{units}
  \sphinxstyleemphasis{units} value = \sphinxstyleemphasis{box} or \sphinxstyleemphasis{lattice}
\end{sphinxalltt}

\end{itemize}


\subsection{Examples}
\label{\detokenize{displace_atoms:examples}}
\begin{sphinxVerbatim}[commandchars=\\\{\}]
\PYG{k}{displace\PYGZus{}atoms }\PYG{n+nv+nvIdentifier}{top}\PYG{+w}{ }\PYG{n}{move}\PYG{+w}{ }\PYG{l+m}{0}\PYG{+w}{ }\PYG{o}{\PYGZhy{}}\PYG{l+m}{5}\PYG{+w}{ }\PYG{l+m}{0}\PYG{+w}{ }\PYG{n}{units}\PYG{+w}{ }\PYG{n}{box}
\PYG{k}{displace\PYGZus{}atoms }\PYG{n+nv+nvIdentifier}{flow}\PYG{+w}{ }\PYG{n}{ramp}\PYG{+w}{ }\PYG{n}{x}\PYG{+w}{ }\PYG{l+m}{0.0}\PYG{+w}{ }\PYG{l+m}{5.0}\PYG{+w}{ }\PYG{n}{y}\PYG{+w}{ }\PYG{l+m}{2.0}\PYG{+w}{ }\PYG{l+m}{20.5}
\end{sphinxVerbatim}


\subsection{Description}
\label{\detokenize{displace_atoms:description}}
\sphinxAtStartPar
Displace a group of atoms.  This can be used to move atoms a large
distance before beginning a simulation or to randomize atoms initially
on a lattice.  For example, in a shear simulation, an initial strain
can be imposed on the system.  Or two groups of atoms can be brought
into closer proximity.

\sphinxAtStartPar
The \sphinxstyleemphasis{move} style displaces the group of atoms by the specified 3d
displacement vector.  Any of the three quantities defining the vector
components can be specified as an equal\sphinxhyphen{}style or atom\sphinxhyphen{}style
{\hyperref[\detokenize{variable::doc}]{\sphinxcrossref{\DUrole{doc}{variable}}}}.  If the value is a variable, it should be
specified as v\_name, where name is the variable name.  In this case,
the variable will be evaluated, and its value(s) used for the
displacement(s).  The scale factor implied by the \sphinxstyleemphasis{units} keyword will
also be applied to the variable result.

\sphinxAtStartPar
Equal\sphinxhyphen{}style variables can specify formulas with various mathematical
functions, and include {\hyperref[\detokenize{thermo_style::doc}]{\sphinxcrossref{\DUrole{doc}{thermo\_style}}}} command
keywords for the simulation box parameters and timestep and elapsed
time.  Atom\sphinxhyphen{}style variables can specify the same formulas as
equal\sphinxhyphen{}style variables but can also include per\sphinxhyphen{}atom values, such as
atom coordinates or per\sphinxhyphen{}atom values read from a file.  Note that if
the variable references other {\hyperref[\detokenize{compute::doc}]{\sphinxcrossref{\DUrole{doc}{compute}}}} or {\hyperref[\detokenize{fix::doc}]{\sphinxcrossref{\DUrole{doc}{fix}}}}
commands, those values must be up\sphinxhyphen{}to\sphinxhyphen{}date for the current timestep.
See the “Variable Accuracy” section of the {\hyperref[\detokenize{variable::doc}]{\sphinxcrossref{\DUrole{doc}{variable}}}}
doc page for more details.

\sphinxAtStartPar
The \sphinxstyleemphasis{ramp} style displaces atoms a variable amount in one dimension
depending on the atom’s coordinate in a (possibly) different
dimension.  For example, the second example command displaces atoms in
the \sphinxstyleemphasis{x}\sphinxhyphen{}direction an amount between 0.0 and 5.0 distance units.  Each
atom’s displacement depends on the fractional distance its \sphinxstyleemphasis{y}
coordinate is between 2.0 and 20.5.  Atoms with \sphinxstyleemphasis{y}\sphinxhyphen{}coordinates outside
those bounds will be moved the minimum (0.0) or maximum (5.0) amount.

\sphinxAtStartPar
The \sphinxstyleemphasis{random} style independently moves each atom in the group by a
random displacement, uniformly sampled from a value between \(-dx\) and
\(+dx\) in the \sphinxstyleemphasis{x} dimension, and similarly for \sphinxstyleemphasis{y} and \sphinxstyleemphasis{z}.
Random numbers are used in such a way that the displacement of a particular
atom is the same, regardless of how many processors are being used.

\sphinxAtStartPar
The \sphinxstyleemphasis{rotate} style rotates each atom in the group by the angle \sphinxstyleemphasis{theta}
around a rotation axis \(R = (R_x,R_y,R_z)\) that goes through a point
\(P = (P_x,P_y,P_z)\).  The direction of rotation for the atoms around the
rotation axis is consistent with the right\sphinxhyphen{}hand rule: if your
right\sphinxhyphen{}hand thumb points along \sphinxstyleemphasis{R}, then your fingers wrap around the
axis in the direction of positive theta.

\sphinxAtStartPar
If the defined {\hyperref[\detokenize{atom_style::doc}]{\sphinxcrossref{\DUrole{doc}{atom\_style}}}} assigns an orientation to
each atom ({\hyperref[\detokenize{atom_style::doc}]{\sphinxcrossref{\DUrole{doc}{atom styles}}}} ellipsoid, line, tri, body),
then that property is also updated appropriately to correspond to the
atom’s rotation.

\sphinxAtStartPar
Distance units for displacements and the origin point of the \sphinxstyleemphasis{rotate}
style are determined by the setting of \sphinxstyleemphasis{box} or \sphinxstyleemphasis{lattice} for the
\sphinxstyleemphasis{units} keyword.  \sphinxstyleemphasis{Box} means distance units as defined by the
{\hyperref[\detokenize{units::doc}]{\sphinxcrossref{\DUrole{doc}{units}}}} command (e.g., \(\AA\) for
\sphinxstyleemphasis{real} or \sphinxstyleemphasis{metal} units). \sphinxstyleemphasis{Lattice} means distance units are in lattice
spacings.  The {\hyperref[\detokenize{lattice::doc}]{\sphinxcrossref{\DUrole{doc}{lattice}}}} command must have been previously used
to define the lattice spacing.


\bigskip\hrule\bigskip


\begin{sphinxadmonition}{note}{Note:}
\sphinxAtStartPar
Care should be taken not to move atoms on top of other atoms.
After the move, atoms are remapped into the periodic simulation box if
needed, and any shrink\sphinxhyphen{}wrap boundary conditions (see the
{\hyperref[\detokenize{boundary::doc}]{\sphinxcrossref{\DUrole{doc}{boundary}}}} command) are enforced which may change the
box size.  Other than this effect, this command does not change the
size or shape of the simulation box.  See the
{\hyperref[\detokenize{change_box::doc}]{\sphinxcrossref{\DUrole{doc}{change\_box}}}} command if that effect is desired.
\end{sphinxadmonition}

\begin{sphinxadmonition}{note}{Note:}
\sphinxAtStartPar
Atoms can be moved arbitrarily long distances by this command.
If the simulation box is non\sphinxhyphen{}periodic and shrink\sphinxhyphen{}wrapped (see the
{\hyperref[\detokenize{boundary::doc}]{\sphinxcrossref{\DUrole{doc}{boundary}}}} command), this can change its size or shape.
This is not a problem, except that the mapping of processors to the
simulation box is not changed by this command from its initial 3d
configuration; see the {\hyperref[\detokenize{processors::doc}]{\sphinxcrossref{\DUrole{doc}{processors}}}} command.  Thus, if
the box size/shape changes dramatically, the mapping of processors to
the simulation box may not end up as optimal as the initial mapping
attempted to be.
\end{sphinxadmonition}


\bigskip\hrule\bigskip



\subsection{Restrictions}
\label{\detokenize{displace_atoms:restrictions}}
\sphinxAtStartPar
For a 2d simulation, only rotations around the a vector parallel to
the \(z\)\sphinxhyphen{}axis are allowed.


\subsection{Related commands}
\label{\detokenize{displace_atoms:related-commands}}
\sphinxAtStartPar
{\hyperref[\detokenize{lattice::doc}]{\sphinxcrossref{\DUrole{doc}{lattice}}}}, {\hyperref[\detokenize{change_box::doc}]{\sphinxcrossref{\DUrole{doc}{change\_box}}}},
{\hyperref[\detokenize{fix_move::doc}]{\sphinxcrossref{\DUrole{doc}{fix move}}}}


\subsection{Default}
\label{\detokenize{displace_atoms:default}}
\sphinxAtStartPar
The option defaults are units = lattice.

\sphinxstepscope

\index{dynamical\_matrix@\spxentry{dynamical\_matrix}}\ignorespaces 
\index{dynamical\_matrix/kk@\spxentry{dynamical\_matrix/kk}}\ignorespaces 

\section{dynamical\_matrix command}
\label{\detokenize{dynamical_matrix:dynamical-matrix-command}}\label{\detokenize{dynamical_matrix:index-1}}\label{\detokenize{dynamical_matrix:index-0}}\label{\detokenize{dynamical_matrix::doc}}
\sphinxAtStartPar
Accelerator Variant: dynamical\_matrix/kk


\subsection{Syntax}
\label{\detokenize{dynamical_matrix:syntax}}
\begin{sphinxVerbatim}[commandchars=\\\{\}]
\PYG{k}{dynamical\PYGZus{}matrix }\PYG{n+nv+nvIdentifier}{group\PYGZhy{}ID}\PYG{+w}{ }\PYG{n}{style}\PYG{+w}{ }\PYG{n}{gamma}\PYG{+w}{ }\PYG{n}{args}\PYG{+w}{ }\PYG{n}{keyword}\PYG{+w}{ }\PYG{n}{value}\PYG{+w}{ }\PYG{n}{...}
\end{sphinxVerbatim}
\begin{itemize}
\item {} 
\sphinxAtStartPar
group\sphinxhyphen{}ID = ID of group of atoms to displace

\item {} 
\sphinxAtStartPar
style = \sphinxstyleemphasis{regular} or \sphinxstyleemphasis{eskm}

\item {} 
\sphinxAtStartPar
gamma = finite different displacement length (distance units)

\item {} 
\sphinxAtStartPar
one or more keyword/arg pairs may be appended
\begin{sphinxalltt}
keyword = \sphinxstyleemphasis{file} or \sphinxstyleemphasis{binary}
  \sphinxstyleemphasis{file} name = name of output file for the dynamical matrix
  \sphinxstyleemphasis{binary} arg = \sphinxstyleemphasis{yes} or \sphinxstyleemphasis{no} or \sphinxstyleemphasis{gzip}
\end{sphinxalltt}

\end{itemize}


\subsection{Examples}
\label{\detokenize{dynamical_matrix:examples}}
\begin{sphinxVerbatim}[commandchars=\\\{\}]
\PYG{k}{dynamical\PYGZus{}matrix }\PYG{n+nv+nvIdentifier}{1}\PYG{+w}{ }\PYG{n}{regular}\PYG{+w}{ }\PYG{l+m}{0.000001}
\PYG{k}{dynamical\PYGZus{}matrix }\PYG{n+nv+nvIdentifier}{1}\PYG{+w}{ }\PYG{n}{eskm}\PYG{+w}{ }\PYG{l+m}{0.000001}
\PYG{k}{dynamical\PYGZus{}matrix }\PYG{n+nv+nvIdentifier}{3}\PYG{+w}{ }\PYG{n}{regular}\PYG{+w}{ }\PYG{l+m}{0.00004}\PYG{+w}{ }\PYG{n}{file}\PYG{+w}{ }\PYG{n}{dynmat.dat}
\PYG{k}{dynamical\PYGZus{}matrix }\PYG{n+nv+nvIdentifier}{5}\PYG{+w}{ }\PYG{n}{eskm}\PYG{+w}{ }\PYG{l+m}{0.00000001}\PYG{+w}{ }\PYG{n}{file}\PYG{+w}{ }\PYG{n}{dynamical.dat}\PYG{+w}{ }\PYG{n}{binary}\PYG{+w}{ }\PYG{n}{yes}
\end{sphinxVerbatim}


\subsection{Description}
\label{\detokenize{dynamical_matrix:description}}
\sphinxAtStartPar
Calculate the dynamical matrix by finite difference of the selected group,
\begin{equation*}
\begin{split}D = \frac{\Phi_{ij}^{\alpha\beta}}{\sqrt{M_i M_j}}\end{split}
\end{equation*}
\sphinxAtStartPar
where D is the dynamical matrix and \(\Phi\) is the force constant
matrix defined by
\begin{equation*}
\begin{split}\Phi_{ij}^{\alpha\beta} = \frac{\partial^2 U}{\partial x_{i,\alpha} \partial x_{j,\beta}}\end{split}
\end{equation*}
\sphinxAtStartPar
The output for the dynamical matrix is printed three elements at a time.
The three elements are the three \(\beta\) elements for a respective
i/\(\alpha\)/j combination.  Each line is printed in order of j
increasing first, \(\alpha\) second, and i last.

\sphinxAtStartPar
If the style eskm is selected, the dynamical matrix will be in units of
inverse squared femtoseconds. These units will then conveniently leave
frequencies in THz.


\bigskip\hrule\bigskip


\sphinxAtStartPar
Styles with a \sphinxstyleemphasis{gpu}, \sphinxstyleemphasis{intel}, \sphinxstyleemphasis{kk}, \sphinxstyleemphasis{omp}, or \sphinxstyleemphasis{opt} suffix are
functionally the same as the corresponding style without the suffix.
They have been optimized to run faster, depending on your available
hardware, as discussed on the {\hyperref[\detokenize{Speed_packages::doc}]{\sphinxcrossref{\DUrole{doc}{Accelerator packages}}}}
page.  The accelerated styles take the same arguments and should
produce the same results, except for round\sphinxhyphen{}off and precision issues.

\sphinxAtStartPar
These accelerated styles are part of the GPU, INTEL, KOKKOS,
OPENMP, and OPT packages, respectively.  They are only enabled if
LAMMPS was built with those packages.  See the {\hyperref[\detokenize{Build_package::doc}]{\sphinxcrossref{\DUrole{doc}{Build package}}}} page for more info.

\sphinxAtStartPar
You can specify the accelerated styles explicitly in your input script
by including their suffix, or you can use the {\hyperref[\detokenize{Run_options::doc}]{\sphinxcrossref{\DUrole{doc}{\sphinxhyphen{}suffix command\sphinxhyphen{}line switch}}}} when you invoke LAMMPS, or you can use the
{\hyperref[\detokenize{suffix::doc}]{\sphinxcrossref{\DUrole{doc}{suffix}}}} command in your input script.

\sphinxAtStartPar
See the {\hyperref[\detokenize{Speed_packages::doc}]{\sphinxcrossref{\DUrole{doc}{Accelerator packages}}}} page for more
instructions on how to use the accelerated styles effectively.


\bigskip\hrule\bigskip



\subsection{Restrictions}
\label{\detokenize{dynamical_matrix:restrictions}}
\sphinxAtStartPar
The command collects an array of nine times the number of atoms in a group
on every single MPI rank, so the memory requirements can be very significant
for large systems.

\sphinxAtStartPar
This command is part of the PHONON package.  It is only enabled if
LAMMPS was built with that package.
See the {\hyperref[\detokenize{Build_package::doc}]{\sphinxcrossref{\DUrole{doc}{Build package}}}} page for more info.


\subsection{Related commands}
\label{\detokenize{dynamical_matrix:related-commands}}
\sphinxAtStartPar
{\hyperref[\detokenize{fix_phonon::doc}]{\sphinxcrossref{\DUrole{doc}{fix phonon}}}}, {\hyperref[\detokenize{fix_numdiff::doc}]{\sphinxcrossref{\DUrole{doc}{fix numdiff}}}},

\sphinxAtStartPar
{\hyperref[\detokenize{compute_hma::doc}]{\sphinxcrossref{\DUrole{doc}{compute hma}}}} uses an analytic formulation of the
Hessian provided by a pair\_style’s Pair::single\_hessian() function,
if implemented.


\subsection{Default}
\label{\detokenize{dynamical_matrix:default}}
\sphinxAtStartPar
The default settings are file = “dynmat.dyn”, binary = no

\sphinxstepscope

\index{echo@\spxentry{echo}}\ignorespaces 

\section{echo command}
\label{\detokenize{echo:echo-command}}\label{\detokenize{echo:index-0}}\label{\detokenize{echo::doc}}

\subsection{Syntax}
\label{\detokenize{echo:syntax}}
\begin{sphinxVerbatim}[commandchars=\\\{\}]
\PYG{k}{echo}\PYG{+w}{ }\PYG{n}{style}
\end{sphinxVerbatim}
\begin{itemize}
\item {} 
\sphinxAtStartPar
style = \sphinxstyleemphasis{none} or \sphinxstyleemphasis{screen} or \sphinxstyleemphasis{log} or \sphinxstyleemphasis{both}

\end{itemize}


\subsection{Examples}
\label{\detokenize{echo:examples}}
\begin{sphinxVerbatim}[commandchars=\\\{\}]
\PYG{k}{echo}\PYG{+w}{ }\PYG{n}{both}
\PYG{k}{echo}\PYG{+w}{ }\PYG{n}{log}
\end{sphinxVerbatim}


\subsection{Description}
\label{\detokenize{echo:description}}
\sphinxAtStartPar
This command determines whether LAMMPS echoes each input script
command to the screen and/or log file as it is read and processed.  If
an input script has errors, it can be useful to look at echoed output
to see the last command processed.

\sphinxAtStartPar
The {\hyperref[\detokenize{Run_options::doc}]{\sphinxcrossref{\DUrole{doc}{command\sphinxhyphen{}line switch}}}} \sphinxhyphen{}echo can be used in place
of this command.


\subsection{Restrictions}
\label{\detokenize{echo:restrictions}}\begin{quote}

\sphinxAtStartPar
none
\end{quote}


\subsection{Related commands}
\label{\detokenize{echo:related-commands}}
\sphinxAtStartPar
none


\subsection{Default}
\label{\detokenize{echo:default}}
\begin{sphinxVerbatim}[commandchars=\\\{\}]
\PYG{k}{echo}\PYG{+w}{ }\PYG{n}{log}
\end{sphinxVerbatim}

\sphinxstepscope

\index{fix@\spxentry{fix}}\ignorespaces 

\section{fix command}
\label{\detokenize{fix:fix-command}}\label{\detokenize{fix:index-0}}\label{\detokenize{fix::doc}}

\subsection{Syntax}
\label{\detokenize{fix:syntax}}
\begin{sphinxVerbatim}[commandchars=\\\{\}]
\PYG{k}{fix }\PYG{n+nv+nvIdentifier}{ID}\PYG{+w}{ }\PYG{n+nv+nvIdentifier}{group\PYGZhy{}ID}\PYG{+w}{ }\PYG{n}{style}\PYG{+w}{ }\PYG{n}{args}
\end{sphinxVerbatim}
\begin{itemize}
\item {} 
\sphinxAtStartPar
ID = user\sphinxhyphen{}assigned name for the fix

\item {} 
\sphinxAtStartPar
group\sphinxhyphen{}ID = ID of the group of atoms to apply the fix to

\item {} 
\sphinxAtStartPar
style = one of a long list of possible style names (see below)

\item {} 
\sphinxAtStartPar
args = arguments used by a particular style

\end{itemize}


\subsection{Examples}
\label{\detokenize{fix:examples}}
\begin{sphinxVerbatim}[commandchars=\\\{\}]
\PYG{k}{fix }\PYG{n+nv+nvIdentifier}{1}\PYG{+w}{ }\PYG{n+nv+nvIdentifier}{all}\PYG{+w}{ }\PYG{n}{nve}
\PYG{k}{fix }\PYG{n+nv+nvIdentifier}{3}\PYG{+w}{ }\PYG{n+nv+nvIdentifier}{all}\PYG{+w}{ }\PYG{n}{nvt}\PYG{+w}{ }\PYG{n}{temp}\PYG{+w}{ }\PYG{l+m}{300.0}\PYG{+w}{ }\PYG{l+m}{300.0}\PYG{+w}{ }\PYG{l+m}{0.01}
\PYG{k}{fix }\PYG{n+nv+nvIdentifier}{mine}\PYG{+w}{ }\PYG{n+nv+nvIdentifier}{top}\PYG{+w}{ }\PYG{n}{setforce}\PYG{+w}{ }\PYG{l+m}{0.0}\PYG{+w}{ }\PYG{n}{NULL}\PYG{+w}{ }\PYG{l+m}{0.0}
\end{sphinxVerbatim}


\subsection{Description}
\label{\detokenize{fix:description}}
\sphinxAtStartPar
Set a fix that will be applied to a group of atoms.  In LAMMPS, a
“fix” is any operation that is applied to the system during
timestepping or minimization.  Examples include updating of atom
positions and velocities due to time integration, controlling
temperature, applying constraint forces to atoms, enforcing boundary
conditions, computing diagnostics, etc.  There are hundreds of fixes
defined in LAMMPS and new ones can be added; see the
{\hyperref[\detokenize{Modify::doc}]{\sphinxcrossref{\DUrole{doc}{Modify}}}} page for details.

\sphinxAtStartPar
Fixes perform their operations at different stages of the timestep.
If two or more fixes operate at the same stage of the timestep, they are
invoked in the order they were specified in the input script.

\sphinxAtStartPar
The ID of a fix can only contain alphanumeric characters and
underscores.

\sphinxAtStartPar
Fixes can be deleted with the {\hyperref[\detokenize{unfix::doc}]{\sphinxcrossref{\DUrole{doc}{unfix}}}} command.

\begin{sphinxadmonition}{note}{Note:}
\sphinxAtStartPar
The {\hyperref[\detokenize{unfix::doc}]{\sphinxcrossref{\DUrole{doc}{unfix}}}} command is the only way to turn off a
fix; simply specifying a new fix with a similar style will not turn
off the first one.  This is especially important to realize for
integration fixes.  For example, using a {\hyperref[\detokenize{fix_nve::doc}]{\sphinxcrossref{\DUrole{doc}{fix nve}}}}
command for a second run after using a {\hyperref[\detokenize{fix_nh::doc}]{\sphinxcrossref{\DUrole{doc}{fix nvt}}}} command
for the first run will not cancel out the NVT time integration
invoked by the “fix nvt” command.  Thus, two time integrators would be
in place!
\end{sphinxadmonition}

\sphinxAtStartPar
If you specify a new fix with the same ID and style as an existing
fix, the old fix is deleted and the new one is created (presumably
with new settings).  This is the same as if an “unfix” command were
first performed on the old fix, except that the new fix is kept in the
same order relative to the existing fixes as the old one originally
was.  Note that this operation also wipes out any additional changes
made to the old fix via the {\hyperref[\detokenize{fix_modify::doc}]{\sphinxcrossref{\DUrole{doc}{fix\_modify}}}} command.

\sphinxAtStartPar
The {\hyperref[\detokenize{fix_modify::doc}]{\sphinxcrossref{\DUrole{doc}{fix modify}}}} command allows settings for some
fixes to be reset.  See the page for individual fixes for details.

\sphinxAtStartPar
Some fixes store an internal “state” which is written to binary
restart files via the {\hyperref[\detokenize{restart::doc}]{\sphinxcrossref{\DUrole{doc}{restart}}}} or
{\hyperref[\detokenize{write_restart::doc}]{\sphinxcrossref{\DUrole{doc}{write\_restart}}}} commands.  This allows the fix to
continue on with its calculations in a restarted simulation.  See the
{\hyperref[\detokenize{read_restart::doc}]{\sphinxcrossref{\DUrole{doc}{read\_restart}}}} command for info on how to re\sphinxhyphen{}specify
a fix in an input script that reads a restart file.  See the doc pages
for individual fixes for info on which ones can be restarted.


\bigskip\hrule\bigskip


\sphinxAtStartPar
Some fixes calculate and store any of four \sphinxstyleemphasis{styles} of quantities:
global, per\sphinxhyphen{}atom, local, or per\sphinxhyphen{}grid.

\sphinxAtStartPar
A global quantity is one or more system\sphinxhyphen{}wide values, e.g. the energy
of a wall interacting with particles.  A per\sphinxhyphen{}atom quantity is one or
more values per atom, e.g. the original coordinates of each atom at
time 0.  Per\sphinxhyphen{}atom values are set to 0.0 for atoms not in the specified
fix group.  Local quantities are calculated by each processor based on
the atoms it owns, but there may be zero or more per atom, e.g. values
for each bond.  Per\sphinxhyphen{}grid quantities are calculated on a regular 2d or
3d grid which overlays a 2d or 3d simulation domain.  The grid points
and the data they store are distributed across processors; each
processor owns the grid points which fall within its subdomain.

\sphinxAtStartPar
As a general rule of thumb, fixes that produce per\sphinxhyphen{}atom quantities
have the word “atom” at the end of their style, e.g. \sphinxstyleemphasis{ave/atom}.
Fixes that produce local quantities have the word “local” at the end
of their style, e.g. \sphinxstyleemphasis{store/local}.  Fixes that produce per\sphinxhyphen{}grid
quantities have the word “grid” at the end of their style,
e.g. \sphinxstyleemphasis{ave/grid}.

\sphinxAtStartPar
Global, per\sphinxhyphen{}atom, local, and per\sphinxhyphen{}grid quantities can also be of three
\sphinxstyleemphasis{kinds}: a single scalar value (global only), a vector of values, or a
2d array of values.  For per\sphinxhyphen{}atom, local, and per\sphinxhyphen{}grid quantities, a
“vector” means a single value for each atom, each local entity
(e.g. bond), or grid cell.  Likewise an “array”, means multiple values
for each atom, each local entity, or each grid cell.

\sphinxAtStartPar
Note that a single fix can produce any combination of global,
per\sphinxhyphen{}atom, local, or per\sphinxhyphen{}grid values.  Likewise it can produce any
combination of scalar, vector, or array output for each style.  The
exception is that for per\sphinxhyphen{}atom, local, and per\sphinxhyphen{}grid output, either a
vector or array can be produced, but not both.  The doc page for each
fix explains the values it produces, if any.

\sphinxAtStartPar
When a fix output is accessed by another input script command it is
referenced via the following bracket notation, where ID is the ID of
the fix:


\begin{savenotes}\sphinxattablestart
\sphinxthistablewithglobalstyle
\centering
\begin{tabulary}{\linewidth}[t]{TT}
\sphinxtoprule
\sphinxtableatstartofbodyhook
\sphinxAtStartPar
f\_ID
&
\sphinxAtStartPar
entire scalar, vector, or array
\\
\sphinxhline
\sphinxAtStartPar
f\_ID{[}I{]}
&
\sphinxAtStartPar
one element of vector, one column of array
\\
\sphinxhline
\sphinxAtStartPar
f\_ID{[}I{]}{[}J{]}
&
\sphinxAtStartPar
one element of array
\\
\sphinxbottomrule
\end{tabulary}
\sphinxtableafterendhook\par
\sphinxattableend\end{savenotes}

\sphinxAtStartPar
In other words, using one bracket reduces the dimension of the
quantity once (vector \(\to\) scalar, array \(\to\) vector).
Using two brackets reduces the dimension twice (array \(\to\)
scalar).  Thus, for example, a command that uses global scalar fix
values as input can also process elements of a vector or array.
Depending on the command, this can either be done directly using the
syntax in the table, or by first defining a {\hyperref[\detokenize{variable::doc}]{\sphinxcrossref{\DUrole{doc}{variable}}}}
of the appropriate style to store the quantity, then using the
variable as an input to the command.

\sphinxAtStartPar
Note that commands and {\hyperref[\detokenize{variable::doc}]{\sphinxcrossref{\DUrole{doc}{variables}}}} which take fix
outputs as input typically do not allow for all styles and kinds of
data (e.g., a command may require global but not per\sphinxhyphen{}atom values, or
it may require a vector of values, not a scalar).  This means there is
typically no ambiguity about referring to a fix output as c\_ID even if
it produces, for example, both a scalar and vector.  The doc pages for
various commands explain the details, including how any ambiguities
are resolved.


\bigskip\hrule\bigskip


\sphinxAtStartPar
In LAMMPS, the values generated by a fix can be used in several ways:
\begin{itemize}
\item {} 
\sphinxAtStartPar
Global values can be output via the {\hyperref[\detokenize{thermo_style::doc}]{\sphinxcrossref{\DUrole{doc}{thermo\_style custom}}}} or {\hyperref[\detokenize{fix_ave_time::doc}]{\sphinxcrossref{\DUrole{doc}{fix ave/time}}}} command.
Alternatively, the values can be referenced in an
{\hyperref[\detokenize{variable::doc}]{\sphinxcrossref{\DUrole{doc}{equal\sphinxhyphen{}style variable}}}} command.

\item {} 
\sphinxAtStartPar
Per\sphinxhyphen{}atom values can be output via the {\hyperref[\detokenize{dump::doc}]{\sphinxcrossref{\DUrole{doc}{dump custom}}}} command,
or they can be time\sphinxhyphen{}averaged via the {\hyperref[\detokenize{fix_ave_atom::doc}]{\sphinxcrossref{\DUrole{doc}{fix ave/atom}}}}
command or reduced by the {\hyperref[\detokenize{compute_reduce::doc}]{\sphinxcrossref{\DUrole{doc}{compute reduce}}}}
command.  Alternatively, per\sphinxhyphen{}atom values can be referenced in an
{\hyperref[\detokenize{variable::doc}]{\sphinxcrossref{\DUrole{doc}{atom\sphinxhyphen{}style variable}}}}.

\item {} 
\sphinxAtStartPar
Local values can be reduced by the {\hyperref[\detokenize{compute_reduce::doc}]{\sphinxcrossref{\DUrole{doc}{compute reduce}}}}
command or histogrammed by the {\hyperref[\detokenize{fix_ave_histo::doc}]{\sphinxcrossref{\DUrole{doc}{fix ave/histo}}}} command.
They can also be output by the {\hyperref[\detokenize{dump::doc}]{\sphinxcrossref{\DUrole{doc}{dump local}}}} command.

\end{itemize}

\sphinxAtStartPar
See the {\hyperref[\detokenize{Howto_output::doc}]{\sphinxcrossref{\DUrole{doc}{Howto output}}}} page for a summary of
various LAMMPS output options, many of which involve fixes.

\sphinxAtStartPar
The results of fixes that calculate global quantities can be either
“intensive” or “extensive” values.  Intensive means the value is
independent of the number of atoms in the simulation
(e.g., temperature).  Extensive means the value scales with the number of
atoms in the simulation (e.g., total rotational kinetic energy).
{\hyperref[\detokenize{thermo_style::doc}]{\sphinxcrossref{\DUrole{doc}{Thermodynamic output}}}} will normalize extensive
values by the number of atoms in the system, depending on the
“thermo\_modify norm” setting.  It will not normalize intensive values.
If a fix value is accessed in another way (e.g., by a
{\hyperref[\detokenize{variable::doc}]{\sphinxcrossref{\DUrole{doc}{variable}}}}), you may want to know whether it is an
intensive or extensive value.  See the page for individual fix styles
for further info.


\bigskip\hrule\bigskip


\sphinxAtStartPar
Each fix style has its own page that describes its arguments and
what it does, as listed below.  Here is an alphabetical list of fix
styles available in LAMMPS.  They are also listed in more compact form
on the {\hyperref[\detokenize{Commands_fix::doc}]{\sphinxcrossref{\DUrole{doc}{Commands fix}}}} doc page.

\sphinxAtStartPar
There are also additional accelerated fix styles included in the
LAMMPS distribution for faster performance on CPUs, GPUs, and KNLs.
The individual style names on the {\hyperref[\detokenize{Commands_fix::doc}]{\sphinxcrossref{\DUrole{doc}{Commands fix}}}} doc
page are followed by one or more of (g,i,k,o,t) to indicate which
accelerated styles exist.
\begin{itemize}
\item {} 
\sphinxAtStartPar
{\hyperref[\detokenize{fix_accelerate_cos::doc}]{\sphinxcrossref{\DUrole{doc}{accelerate/cos}}}} \sphinxhyphen{} apply cosine\sphinxhyphen{}shaped acceleration to atoms

\item {} 
\sphinxAtStartPar
{\hyperref[\detokenize{fix_acks2_reaxff::doc}]{\sphinxcrossref{\DUrole{doc}{acks2/reaxff}}}} \sphinxhyphen{} apply ACKS2 charge equilibration

\item {} 
\sphinxAtStartPar
{\hyperref[\detokenize{fix_adapt::doc}]{\sphinxcrossref{\DUrole{doc}{adapt}}}} \sphinxhyphen{} change a simulation parameter over time

\item {} 
\sphinxAtStartPar
{\hyperref[\detokenize{fix_adapt_fep::doc}]{\sphinxcrossref{\DUrole{doc}{adapt/fep}}}} \sphinxhyphen{} enhanced version of fix adapt

\item {} 
\sphinxAtStartPar
{\hyperref[\detokenize{fix_addforce::doc}]{\sphinxcrossref{\DUrole{doc}{addforce}}}} \sphinxhyphen{} add a force to each atom

\item {} 
\sphinxAtStartPar
{\hyperref[\detokenize{fix_add_heat::doc}]{\sphinxcrossref{\DUrole{doc}{add/heat}}}} \sphinxhyphen{} add a heat flux to each atom

\item {} 
\sphinxAtStartPar
{\hyperref[\detokenize{fix_addtorque::doc}]{\sphinxcrossref{\DUrole{doc}{addtorque}}}} \sphinxhyphen{} add a torque to a group of atoms

\item {} 
\sphinxAtStartPar
{\hyperref[\detokenize{fix_alchemy::doc}]{\sphinxcrossref{\DUrole{doc}{alchemy}}}} \sphinxhyphen{} perform an “alchemical transformation” between two partitions

\item {} 
\sphinxAtStartPar
{\hyperref[\detokenize{fix_amoeba_bitorsion::doc}]{\sphinxcrossref{\DUrole{doc}{amoeba/bitorsion}}}} \sphinxhyphen{} torsion/torsion terms in AMOEBA force field

\item {} 
\sphinxAtStartPar
{\hyperref[\detokenize{fix_amoeba_pitorsion::doc}]{\sphinxcrossref{\DUrole{doc}{amoeba/pitorsion}}}} \sphinxhyphen{} 6\sphinxhyphen{}body terms in AMOEBA force field

\item {} 
\sphinxAtStartPar
{\hyperref[\detokenize{fix_append_atoms::doc}]{\sphinxcrossref{\DUrole{doc}{append/atoms}}}} \sphinxhyphen{} append atoms to a running simulation

\item {} 
\sphinxAtStartPar
{\hyperref[\detokenize{fix_atc::doc}]{\sphinxcrossref{\DUrole{doc}{atc}}}} \sphinxhyphen{} initiates a coupled MD/FE simulation

\item {} 
\sphinxAtStartPar
{\hyperref[\detokenize{fix_atom_swap::doc}]{\sphinxcrossref{\DUrole{doc}{atom/swap}}}} \sphinxhyphen{} Monte Carlo atom type swapping

\item {} 
\sphinxAtStartPar
{\hyperref[\detokenize{fix_ave_atom::doc}]{\sphinxcrossref{\DUrole{doc}{ave/atom}}}} \sphinxhyphen{} compute per\sphinxhyphen{}atom time\sphinxhyphen{}averaged quantities

\item {} 
\sphinxAtStartPar
{\hyperref[\detokenize{fix_ave_chunk::doc}]{\sphinxcrossref{\DUrole{doc}{ave/chunk}}}} \sphinxhyphen{} compute per\sphinxhyphen{}chunk time\sphinxhyphen{}averaged quantities

\item {} 
\sphinxAtStartPar
{\hyperref[\detokenize{fix_ave_correlate::doc}]{\sphinxcrossref{\DUrole{doc}{ave/correlate}}}} \sphinxhyphen{} compute/output time correlations

\item {} 
\sphinxAtStartPar
{\hyperref[\detokenize{fix_ave_correlate_long::doc}]{\sphinxcrossref{\DUrole{doc}{ave/correlate/long}}}} \sphinxhyphen{} alternative to {\hyperref[\detokenize{fix_ave_correlate::doc}]{\sphinxcrossref{\DUrole{doc}{ave/correlate}}}} that allows efficient calculation over long time windows

\item {} 
\sphinxAtStartPar
{\hyperref[\detokenize{fix_ave_grid::doc}]{\sphinxcrossref{\DUrole{doc}{ave/grid}}}} \sphinxhyphen{} compute per\sphinxhyphen{}grid time\sphinxhyphen{}averaged quantities

\item {} 
\sphinxAtStartPar
{\hyperref[\detokenize{fix_ave_histo::doc}]{\sphinxcrossref{\DUrole{doc}{ave/histo}}}} \sphinxhyphen{} compute/output time\sphinxhyphen{}averaged histograms

\item {} 
\sphinxAtStartPar
{\hyperref[\detokenize{fix_ave_histo::doc}]{\sphinxcrossref{\DUrole{doc}{ave/histo/weight}}}} \sphinxhyphen{} weighted version of fix ave/histo

\item {} 
\sphinxAtStartPar
{\hyperref[\detokenize{fix_ave_time::doc}]{\sphinxcrossref{\DUrole{doc}{ave/time}}}} \sphinxhyphen{} compute/output global time\sphinxhyphen{}averaged quantities

\item {} 
\sphinxAtStartPar
{\hyperref[\detokenize{fix_aveforce::doc}]{\sphinxcrossref{\DUrole{doc}{aveforce}}}} \sphinxhyphen{} add an averaged force to each atom

\item {} 
\sphinxAtStartPar
{\hyperref[\detokenize{fix_balance::doc}]{\sphinxcrossref{\DUrole{doc}{balance}}}} \sphinxhyphen{} perform dynamic load\sphinxhyphen{}balancing

\item {} 
\sphinxAtStartPar
{\hyperref[\detokenize{fix_brownian::doc}]{\sphinxcrossref{\DUrole{doc}{brownian}}}} \sphinxhyphen{} overdamped translational brownian motion

\item {} 
\sphinxAtStartPar
{\hyperref[\detokenize{fix_brownian::doc}]{\sphinxcrossref{\DUrole{doc}{brownian/asphere}}}} \sphinxhyphen{} overdamped translational and rotational brownian motion for ellipsoids

\item {} 
\sphinxAtStartPar
{\hyperref[\detokenize{fix_brownian::doc}]{\sphinxcrossref{\DUrole{doc}{brownian/sphere}}}} \sphinxhyphen{} overdamped translational and rotational brownian motion for spheres

\item {} 
\sphinxAtStartPar
{\hyperref[\detokenize{fix_bocs::doc}]{\sphinxcrossref{\DUrole{doc}{bocs}}}} \sphinxhyphen{} NPT style time integration with pressure correction

\item {} 
\sphinxAtStartPar
{\hyperref[\detokenize{fix_bond_break::doc}]{\sphinxcrossref{\DUrole{doc}{bond/break}}}} \sphinxhyphen{} break bonds on the fly

\item {} 
\sphinxAtStartPar
{\hyperref[\detokenize{fix_bond_create::doc}]{\sphinxcrossref{\DUrole{doc}{bond/create}}}} \sphinxhyphen{} create bonds on the fly

\item {} 
\sphinxAtStartPar
{\hyperref[\detokenize{fix_bond_create::doc}]{\sphinxcrossref{\DUrole{doc}{bond/create/angle}}}} \sphinxhyphen{} create bonds on the fly with angle constraints

\item {} 
\sphinxAtStartPar
{\hyperref[\detokenize{fix_bond_react::doc}]{\sphinxcrossref{\DUrole{doc}{bond/react}}}} \sphinxhyphen{} apply topology changes to model reactions

\item {} 
\sphinxAtStartPar
{\hyperref[\detokenize{fix_bond_swap::doc}]{\sphinxcrossref{\DUrole{doc}{bond/swap}}}} \sphinxhyphen{} Monte Carlo bond swapping

\item {} 
\sphinxAtStartPar
{\hyperref[\detokenize{fix_box_relax::doc}]{\sphinxcrossref{\DUrole{doc}{box/relax}}}} \sphinxhyphen{} relax box size during energy minimization

\item {} 
\sphinxAtStartPar
{\hyperref[\detokenize{fix_charge_regulation::doc}]{\sphinxcrossref{\DUrole{doc}{charge/regulation}}}} \sphinxhyphen{} Monte Carlo sampling of charge regulation

\item {} 
\sphinxAtStartPar
{\hyperref[\detokenize{fix_cmap::doc}]{\sphinxcrossref{\DUrole{doc}{cmap}}}} \sphinxhyphen{} CMAP torsion/torsion terms in CHARMM force field

\item {} 
\sphinxAtStartPar
{\hyperref[\detokenize{fix_colvars::doc}]{\sphinxcrossref{\DUrole{doc}{colvars}}}} \sphinxhyphen{} interface to the collective variables “Colvars” library

\item {} 
\sphinxAtStartPar
{\hyperref[\detokenize{fix_controller::doc}]{\sphinxcrossref{\DUrole{doc}{controller}}}} \sphinxhyphen{} apply control loop feedback mechanism

\item {} 
\sphinxAtStartPar
{\hyperref[\detokenize{fix_damping_cundall::doc}]{\sphinxcrossref{\DUrole{doc}{damping/cundall}}}} \sphinxhyphen{} Cundall non\sphinxhyphen{}viscous damping for granular simulations

\item {} 
\sphinxAtStartPar
{\hyperref[\detokenize{fix_deform::doc}]{\sphinxcrossref{\DUrole{doc}{deform}}}} \sphinxhyphen{} change the simulation box size/shape

\item {} 
\sphinxAtStartPar
{\hyperref[\detokenize{fix_deform_pressure::doc}]{\sphinxcrossref{\DUrole{doc}{deform/pressure}}}} \sphinxhyphen{} change the simulation box size/shape with additional loading conditions

\item {} 
\sphinxAtStartPar
{\hyperref[\detokenize{fix_deposit::doc}]{\sphinxcrossref{\DUrole{doc}{deposit}}}} \sphinxhyphen{} add new atoms above a surface

\item {} 
\sphinxAtStartPar
{\hyperref[\detokenize{fix_dpd_energy::doc}]{\sphinxcrossref{\DUrole{doc}{dpd/energy}}}} \sphinxhyphen{} constant energy dissipative particle dynamics

\item {} 
\sphinxAtStartPar
{\hyperref[\detokenize{fix_drag::doc}]{\sphinxcrossref{\DUrole{doc}{drag}}}} \sphinxhyphen{} drag atoms towards a defined coordinate

\item {} 
\sphinxAtStartPar
{\hyperref[\detokenize{fix_drude::doc}]{\sphinxcrossref{\DUrole{doc}{drude}}}} \sphinxhyphen{} part of Drude oscillator polarization model

\item {} 
\sphinxAtStartPar
{\hyperref[\detokenize{fix_drude_transform::doc}]{\sphinxcrossref{\DUrole{doc}{drude/transform/direct}}}} \sphinxhyphen{}  part of Drude oscillator polarization model

\item {} 
\sphinxAtStartPar
{\hyperref[\detokenize{fix_drude_transform::doc}]{\sphinxcrossref{\DUrole{doc}{drude/transform/inverse}}}} \sphinxhyphen{}  part of Drude oscillator polarization model

\item {} 
\sphinxAtStartPar
{\hyperref[\detokenize{fix_dt_reset::doc}]{\sphinxcrossref{\DUrole{doc}{dt/reset}}}} \sphinxhyphen{} reset the timestep based on velocity, forces

\item {} 
\sphinxAtStartPar
{\hyperref[\detokenize{fix_dpd_source::doc}]{\sphinxcrossref{\DUrole{doc}{edpd/source}}}} \sphinxhyphen{} add heat source to eDPD simulations

\item {} 
\sphinxAtStartPar
{\hyperref[\detokenize{fix_efield::doc}]{\sphinxcrossref{\DUrole{doc}{efield}}}} \sphinxhyphen{} impose electric field on system

\item {} 
\sphinxAtStartPar
{\hyperref[\detokenize{fix_efield::doc}]{\sphinxcrossref{\DUrole{doc}{efield/tip4p}}}} \sphinxhyphen{} impose electric field on system with TIP4P molecules

\item {} 
\sphinxAtStartPar
{\hyperref[\detokenize{fix_ehex::doc}]{\sphinxcrossref{\DUrole{doc}{ehex}}}} \sphinxhyphen{} enhanced heat exchange algorithm

\item {} 
\sphinxAtStartPar
{\hyperref[\detokenize{fix_electrode::doc}]{\sphinxcrossref{\DUrole{doc}{electrode/conp}}}} \sphinxhyphen{} impose electric potential

\item {} 
\sphinxAtStartPar
{\hyperref[\detokenize{fix_electrode::doc}]{\sphinxcrossref{\DUrole{doc}{electrode/conq}}}} \sphinxhyphen{} impose total electric charge

\item {} 
\sphinxAtStartPar
{\hyperref[\detokenize{fix_electrode::doc}]{\sphinxcrossref{\DUrole{doc}{electrode/thermo}}}} \sphinxhyphen{} apply thermo\sphinxhyphen{}potentiostat

\item {} 
\sphinxAtStartPar
{\hyperref[\detokenize{fix_electron_stopping::doc}]{\sphinxcrossref{\DUrole{doc}{electron/stopping}}}} \sphinxhyphen{} electronic stopping power as a friction force

\item {} 
\sphinxAtStartPar
{\hyperref[\detokenize{fix_electron_stopping::doc}]{\sphinxcrossref{\DUrole{doc}{electron/stopping/fit}}}} \sphinxhyphen{} electronic stopping power as a friction force

\item {} 
\sphinxAtStartPar
{\hyperref[\detokenize{fix_enforce2d::doc}]{\sphinxcrossref{\DUrole{doc}{enforce2d}}}} \sphinxhyphen{} zero out \sphinxstyleemphasis{z}\sphinxhyphen{}dimension velocity and force

\item {} 
\sphinxAtStartPar
{\hyperref[\detokenize{fix_eos_cv::doc}]{\sphinxcrossref{\DUrole{doc}{eos/cv}}}} \sphinxhyphen{} applies a mesoparticle equation of state to relate the particle internal energy to the particle internal temperature

\item {} 
\sphinxAtStartPar
{\hyperref[\detokenize{fix_eos_table::doc}]{\sphinxcrossref{\DUrole{doc}{eos/table}}}} \sphinxhyphen{} applies a tabulated mesoparticle equation of state to relate the particle internal energy to the particle internal temperature

\item {} 
\sphinxAtStartPar
{\hyperref[\detokenize{fix_eos_table_rx::doc}]{\sphinxcrossref{\DUrole{doc}{eos/table/rx}}}} \sphinxhyphen{} applies a tabulated mesoparticle equation of state to relate the concentration\sphinxhyphen{}dependent particle internal energy to the particle internal temperature

\item {} 
\sphinxAtStartPar
{\hyperref[\detokenize{fix_evaporate::doc}]{\sphinxcrossref{\DUrole{doc}{evaporate}}}} \sphinxhyphen{} remove atoms from simulation periodically

\item {} 
\sphinxAtStartPar
{\hyperref[\detokenize{fix_external::doc}]{\sphinxcrossref{\DUrole{doc}{external}}}} \sphinxhyphen{} callback to an external driver program

\item {} 
\sphinxAtStartPar
{\hyperref[\detokenize{fix_ffl::doc}]{\sphinxcrossref{\DUrole{doc}{ffl}}}} \sphinxhyphen{} apply a Fast\sphinxhyphen{}Forward Langevin equation thermostat

\item {} 
\sphinxAtStartPar
{\hyperref[\detokenize{fix_filter_corotate::doc}]{\sphinxcrossref{\DUrole{doc}{filter/corotate}}}} \sphinxhyphen{} implement corotation filter to allow larger timesteps with r\sphinxhyphen{}RESPA

\item {} 
\sphinxAtStartPar
{\hyperref[\detokenize{fix_flow_gauss::doc}]{\sphinxcrossref{\DUrole{doc}{flow/gauss}}}} \sphinxhyphen{} Gaussian dynamics for constant mass flux

\item {} 
\sphinxAtStartPar
{\hyperref[\detokenize{fix_freeze::doc}]{\sphinxcrossref{\DUrole{doc}{freeze}}}} \sphinxhyphen{} freeze atoms in a granular simulation

\item {} 
\sphinxAtStartPar
{\hyperref[\detokenize{fix_gcmc::doc}]{\sphinxcrossref{\DUrole{doc}{gcmc}}}} \sphinxhyphen{} grand canonical insertions/deletions

\item {} 
\sphinxAtStartPar
{\hyperref[\detokenize{fix_gld::doc}]{\sphinxcrossref{\DUrole{doc}{gld}}}} \sphinxhyphen{} generalized Langevin dynamics integrator

\item {} 
\sphinxAtStartPar
{\hyperref[\detokenize{fix_gle::doc}]{\sphinxcrossref{\DUrole{doc}{gle}}}} \sphinxhyphen{} generalized Langevin equation thermostat

\item {} 
\sphinxAtStartPar
{\hyperref[\detokenize{fix_gravity::doc}]{\sphinxcrossref{\DUrole{doc}{gravity}}}} \sphinxhyphen{} add gravity to atoms in a granular simulation

\item {} 
\sphinxAtStartPar
{\hyperref[\detokenize{fix_grem::doc}]{\sphinxcrossref{\DUrole{doc}{grem}}}} \sphinxhyphen{} implements the generalized replica exchange method

\item {} 
\sphinxAtStartPar
{\hyperref[\detokenize{fix_halt::doc}]{\sphinxcrossref{\DUrole{doc}{halt}}}} \sphinxhyphen{} terminate a dynamics run or minimization

\item {} 
\sphinxAtStartPar
{\hyperref[\detokenize{fix_heat::doc}]{\sphinxcrossref{\DUrole{doc}{heat}}}} \sphinxhyphen{} add/subtract momentum\sphinxhyphen{}conserving heat

\item {} 
\sphinxAtStartPar
{\hyperref[\detokenize{fix_heat_flow::doc}]{\sphinxcrossref{\DUrole{doc}{heat/flow}}}} \sphinxhyphen{} plain time integration of heat flow with per\sphinxhyphen{}atom temperature updates

\item {} 
\sphinxAtStartPar
{\hyperref[\detokenize{fix_hyper_global::doc}]{\sphinxcrossref{\DUrole{doc}{hyper/global}}}} \sphinxhyphen{} global hyperdynamics

\item {} 
\sphinxAtStartPar
{\hyperref[\detokenize{fix_hyper_local::doc}]{\sphinxcrossref{\DUrole{doc}{hyper/local}}}} \sphinxhyphen{} local hyperdynamics

\item {} 
\sphinxAtStartPar
{\hyperref[\detokenize{fix_imd::doc}]{\sphinxcrossref{\DUrole{doc}{imd}}}} \sphinxhyphen{} implements the “Interactive MD” (IMD) protocol

\item {} 
\sphinxAtStartPar
{\hyperref[\detokenize{fix_indent::doc}]{\sphinxcrossref{\DUrole{doc}{indent}}}} \sphinxhyphen{} impose force due to an indenter

\item {} 
\sphinxAtStartPar
{\hyperref[\detokenize{fix_ipi::doc}]{\sphinxcrossref{\DUrole{doc}{ipi}}}} \sphinxhyphen{} enable LAMMPS to run as a client for i\sphinxhyphen{}PI path\sphinxhyphen{}integral simulations

\item {} 
\sphinxAtStartPar
{\hyperref[\detokenize{fix_langevin::doc}]{\sphinxcrossref{\DUrole{doc}{langevin}}}} \sphinxhyphen{} Langevin temperature control

\item {} 
\sphinxAtStartPar
{\hyperref[\detokenize{fix_langevin_drude::doc}]{\sphinxcrossref{\DUrole{doc}{langevin/drude}}}} \sphinxhyphen{} Langevin temperature control of Drude oscillators

\item {} 
\sphinxAtStartPar
{\hyperref[\detokenize{fix_langevin_eff::doc}]{\sphinxcrossref{\DUrole{doc}{langevin/eff}}}} \sphinxhyphen{} Langevin temperature control for the electron force field model

\item {} 
\sphinxAtStartPar
{\hyperref[\detokenize{fix_langevin_spin::doc}]{\sphinxcrossref{\DUrole{doc}{langevin/spin}}}} \sphinxhyphen{} Langevin temperature control for a spin or spin\sphinxhyphen{}lattice system

\item {} 
\sphinxAtStartPar
{\hyperref[\detokenize{fix_lb_fluid::doc}]{\sphinxcrossref{\DUrole{doc}{lb/fluid}}}} \sphinxhyphen{} lattice\sphinxhyphen{}Boltzmann fluid on a uniform mesh

\item {} 
\sphinxAtStartPar
{\hyperref[\detokenize{fix_lb_momentum::doc}]{\sphinxcrossref{\DUrole{doc}{lb/momentum}}}} \sphinxhyphen{} {\hyperref[\detokenize{fix_momentum::doc}]{\sphinxcrossref{\DUrole{doc}{fix momentum}}}} replacement for use with a lattice\sphinxhyphen{}Boltzmann fluid

\item {} 
\sphinxAtStartPar
{\hyperref[\detokenize{fix_lb_viscous::doc}]{\sphinxcrossref{\DUrole{doc}{lb/viscous}}}} \sphinxhyphen{} {\hyperref[\detokenize{fix_viscous::doc}]{\sphinxcrossref{\DUrole{doc}{fix viscous}}}} replacement for use with a lattice\sphinxhyphen{}Boltzmann fluid

\item {} 
\sphinxAtStartPar
{\hyperref[\detokenize{fix_lineforce::doc}]{\sphinxcrossref{\DUrole{doc}{lineforce}}}} \sphinxhyphen{} constrain atoms to move in a line

\item {} 
\sphinxAtStartPar
{\hyperref[\detokenize{fix_manifoldforce::doc}]{\sphinxcrossref{\DUrole{doc}{manifoldforce}}}} \sphinxhyphen{} restrain atoms to a manifold during minimization

\item {} 
\sphinxAtStartPar
{\hyperref[\detokenize{fix_mdi_qm::doc}]{\sphinxcrossref{\DUrole{doc}{mdi/qm}}}} \sphinxhyphen{} LAMMPS operates as a client for a quantum code via the MolSSI Driver Interface (MDI)

\item {} 
\sphinxAtStartPar
{\hyperref[\detokenize{fix_mdi_qmmm::doc}]{\sphinxcrossref{\DUrole{doc}{mdi/qmmm}}}} \sphinxhyphen{} LAMMPS operates as client for QM/MM simulation with a quantum code via the MolSSI Driver Interface (MDI)

\item {} 
\sphinxAtStartPar
{\hyperref[\detokenize{fix_meso_move::doc}]{\sphinxcrossref{\DUrole{doc}{meso/move}}}} \sphinxhyphen{} move mesoscopic SPH/SDPD particles in a prescribed fashion

\item {} 
\sphinxAtStartPar
{\hyperref[\detokenize{fix_mol_swap::doc}]{\sphinxcrossref{\DUrole{doc}{mol/swap}}}} \sphinxhyphen{} Monte Carlo atom type swapping with a molecule

\item {} 
\sphinxAtStartPar
{\hyperref[\detokenize{fix_momentum::doc}]{\sphinxcrossref{\DUrole{doc}{momentum}}}} \sphinxhyphen{} zero the linear and/or angular momentum of a group of atoms

\item {} 
\sphinxAtStartPar
{\hyperref[\detokenize{fix_momentum::doc}]{\sphinxcrossref{\DUrole{doc}{momentum/chunk}}}} \sphinxhyphen{} zero the linear and/or angular momentum of a chunk of atoms

\item {} 
\sphinxAtStartPar
{\hyperref[\detokenize{fix_move::doc}]{\sphinxcrossref{\DUrole{doc}{move}}}} \sphinxhyphen{} move atoms in a prescribed fashion

\item {} 
\sphinxAtStartPar
{\hyperref[\detokenize{fix_msst::doc}]{\sphinxcrossref{\DUrole{doc}{msst}}}} \sphinxhyphen{} multi\sphinxhyphen{}scale shock technique (MSST) integration

\item {} 
\sphinxAtStartPar
{\hyperref[\detokenize{fix_mvv_dpd::doc}]{\sphinxcrossref{\DUrole{doc}{mvv/dpd}}}} \sphinxhyphen{} DPD using the modified velocity\sphinxhyphen{}Verlet integration algorithm

\item {} 
\sphinxAtStartPar
{\hyperref[\detokenize{fix_mvv_dpd::doc}]{\sphinxcrossref{\DUrole{doc}{mvv/edpd}}}} \sphinxhyphen{} constant energy DPD using the modified velocity\sphinxhyphen{}Verlet algorithm

\item {} 
\sphinxAtStartPar
{\hyperref[\detokenize{fix_mvv_dpd::doc}]{\sphinxcrossref{\DUrole{doc}{mvv/tdpd}}}} \sphinxhyphen{} constant temperature DPD using the modified velocity\sphinxhyphen{}Verlet algorithm

\item {} 
\sphinxAtStartPar
{\hyperref[\detokenize{fix_neb::doc}]{\sphinxcrossref{\DUrole{doc}{neb}}}} \sphinxhyphen{} nudged elastic band (NEB) spring forces

\item {} 
\sphinxAtStartPar
{\hyperref[\detokenize{fix_neb_spin::doc}]{\sphinxcrossref{\DUrole{doc}{neb/spin}}}} \sphinxhyphen{} nudged elastic band (NEB) spring forces for spins

\item {} 
\sphinxAtStartPar
{\hyperref[\detokenize{fix_nonaffine_displacement::doc}]{\sphinxcrossref{\DUrole{doc}{nonaffine/displacement}}}} \sphinxhyphen{} calculate nonaffine displacement of atoms

\item {} 
\sphinxAtStartPar
{\hyperref[\detokenize{fix_nh::doc}]{\sphinxcrossref{\DUrole{doc}{nph}}}} \sphinxhyphen{} constant NPH time integration via Nose/Hoover

\item {} 
\sphinxAtStartPar
{\hyperref[\detokenize{fix_nph_asphere::doc}]{\sphinxcrossref{\DUrole{doc}{nph/asphere}}}} \sphinxhyphen{} NPH for aspherical particles

\item {} 
\sphinxAtStartPar
{\hyperref[\detokenize{fix_nph_body::doc}]{\sphinxcrossref{\DUrole{doc}{nph/body}}}} \sphinxhyphen{} NPH for body particles

\item {} 
\sphinxAtStartPar
{\hyperref[\detokenize{fix_nh_eff::doc}]{\sphinxcrossref{\DUrole{doc}{nph/eff}}}} \sphinxhyphen{} NPH for  nuclei and electrons in the electron force field model

\item {} 
\sphinxAtStartPar
{\hyperref[\detokenize{fix_nph_sphere::doc}]{\sphinxcrossref{\DUrole{doc}{nph/sphere}}}} \sphinxhyphen{} NPH for spherical particles

\item {} 
\sphinxAtStartPar
{\hyperref[\detokenize{fix_nphug::doc}]{\sphinxcrossref{\DUrole{doc}{nphug}}}} \sphinxhyphen{} constant\sphinxhyphen{}stress Hugoniostat integration

\item {} 
\sphinxAtStartPar
{\hyperref[\detokenize{fix_nh::doc}]{\sphinxcrossref{\DUrole{doc}{npt}}}} \sphinxhyphen{} constant NPT time integration via Nose/Hoover

\item {} 
\sphinxAtStartPar
{\hyperref[\detokenize{fix_npt_asphere::doc}]{\sphinxcrossref{\DUrole{doc}{npt/asphere}}}} \sphinxhyphen{} NPT for aspherical particles

\item {} 
\sphinxAtStartPar
{\hyperref[\detokenize{fix_npt_body::doc}]{\sphinxcrossref{\DUrole{doc}{npt/body}}}} \sphinxhyphen{} NPT for body particles

\item {} 
\sphinxAtStartPar
{\hyperref[\detokenize{fix_npt_cauchy::doc}]{\sphinxcrossref{\DUrole{doc}{npt/cauchy}}}} \sphinxhyphen{} NPT with Cauchy stress

\item {} 
\sphinxAtStartPar
{\hyperref[\detokenize{fix_nh_eff::doc}]{\sphinxcrossref{\DUrole{doc}{npt/eff}}}} \sphinxhyphen{} NPT for  nuclei and electrons in the electron force field model

\item {} 
\sphinxAtStartPar
{\hyperref[\detokenize{fix_npt_sphere::doc}]{\sphinxcrossref{\DUrole{doc}{npt/sphere}}}} \sphinxhyphen{} NPT for spherical particles

\item {} 
\sphinxAtStartPar
{\hyperref[\detokenize{fix_nh_uef::doc}]{\sphinxcrossref{\DUrole{doc}{npt/uef}}}} \sphinxhyphen{} NPT style time integration with diagonal flow

\item {} 
\sphinxAtStartPar
{\hyperref[\detokenize{fix_numdiff::doc}]{\sphinxcrossref{\DUrole{doc}{numdiff}}}} \sphinxhyphen{} numerically approximate atomic forces using finite energy differences

\item {} 
\sphinxAtStartPar
{\hyperref[\detokenize{fix_numdiff_virial::doc}]{\sphinxcrossref{\DUrole{doc}{numdiff/virial}}}} \sphinxhyphen{} numerically approximate virial stress tensor using finite energy differences

\item {} 
\sphinxAtStartPar
{\hyperref[\detokenize{fix_nve::doc}]{\sphinxcrossref{\DUrole{doc}{nve}}}} \sphinxhyphen{} constant NVE time integration

\item {} 
\sphinxAtStartPar
{\hyperref[\detokenize{fix_nve_asphere::doc}]{\sphinxcrossref{\DUrole{doc}{nve/asphere}}}} \sphinxhyphen{} NVE for aspherical particles

\item {} 
\sphinxAtStartPar
{\hyperref[\detokenize{fix_nve_asphere_noforce::doc}]{\sphinxcrossref{\DUrole{doc}{nve/asphere/noforce}}}} \sphinxhyphen{} NVE for aspherical particles without forces

\item {} 
\sphinxAtStartPar
{\hyperref[\detokenize{fix_nve_awpmd::doc}]{\sphinxcrossref{\DUrole{doc}{nve/awpmd}}}} \sphinxhyphen{} NVE for the Antisymmetrized Wave Packet Molecular Dynamics model

\item {} 
\sphinxAtStartPar
{\hyperref[\detokenize{fix_nve_body::doc}]{\sphinxcrossref{\DUrole{doc}{nve/body}}}} \sphinxhyphen{} NVE for body particles

\item {} 
\sphinxAtStartPar
{\hyperref[\detokenize{fix_nve_dot::doc}]{\sphinxcrossref{\DUrole{doc}{nve/dot}}}} \sphinxhyphen{} rigid body constant energy time integrator for coarse grain models

\item {} 
\sphinxAtStartPar
{\hyperref[\detokenize{fix_nve_dotc_langevin::doc}]{\sphinxcrossref{\DUrole{doc}{nve/dotc/langevin}}}} \sphinxhyphen{} Langevin style rigid body time integrator for coarse grain models

\item {} 
\sphinxAtStartPar
{\hyperref[\detokenize{fix_nve_eff::doc}]{\sphinxcrossref{\DUrole{doc}{nve/eff}}}} \sphinxhyphen{} NVE for  nuclei and electrons in the electron force field model

\item {} 
\sphinxAtStartPar
{\hyperref[\detokenize{fix_nve_limit::doc}]{\sphinxcrossref{\DUrole{doc}{nve/limit}}}} \sphinxhyphen{} NVE with limited step length

\item {} 
\sphinxAtStartPar
{\hyperref[\detokenize{fix_nve_line::doc}]{\sphinxcrossref{\DUrole{doc}{nve/line}}}} \sphinxhyphen{} NVE for line segments

\item {} 
\sphinxAtStartPar
{\hyperref[\detokenize{fix_nve_manifold_rattle::doc}]{\sphinxcrossref{\DUrole{doc}{nve/manifold/rattle}}}} \sphinxhyphen{} NVE time integration for atoms constrained to a curved surface (manifold)

\item {} 
\sphinxAtStartPar
{\hyperref[\detokenize{fix_nve_noforce::doc}]{\sphinxcrossref{\DUrole{doc}{nve/noforce}}}} \sphinxhyphen{} NVE without forces (update positions only)

\item {} 
\sphinxAtStartPar
{\hyperref[\detokenize{fix_nve_sphere::doc}]{\sphinxcrossref{\DUrole{doc}{nve/sphere}}}} \sphinxhyphen{} NVE for spherical particles

\item {} 
\sphinxAtStartPar
{\hyperref[\detokenize{fix_nve_bpm_sphere::doc}]{\sphinxcrossref{\DUrole{doc}{nve/bpm/sphere}}}} \sphinxhyphen{} NVE for spherical particles used in the BPM package

\item {} 
\sphinxAtStartPar
{\hyperref[\detokenize{fix_nve_spin::doc}]{\sphinxcrossref{\DUrole{doc}{nve/spin}}}} \sphinxhyphen{} NVE for a spin or spin\sphinxhyphen{}lattice system

\item {} 
\sphinxAtStartPar
{\hyperref[\detokenize{fix_nve_tri::doc}]{\sphinxcrossref{\DUrole{doc}{nve/tri}}}} \sphinxhyphen{} NVE for triangles

\item {} 
\sphinxAtStartPar
{\hyperref[\detokenize{fix_nvk::doc}]{\sphinxcrossref{\DUrole{doc}{nvk}}}} \sphinxhyphen{} constant kinetic energy time integration

\item {} 
\sphinxAtStartPar
{\hyperref[\detokenize{fix_nh::doc}]{\sphinxcrossref{\DUrole{doc}{nvt}}}} \sphinxhyphen{} NVT time integration via Nose/Hoover

\item {} 
\sphinxAtStartPar
{\hyperref[\detokenize{fix_nvt_asphere::doc}]{\sphinxcrossref{\DUrole{doc}{nvt/asphere}}}} \sphinxhyphen{} NVT for aspherical particles

\item {} 
\sphinxAtStartPar
{\hyperref[\detokenize{fix_nvt_body::doc}]{\sphinxcrossref{\DUrole{doc}{nvt/body}}}} \sphinxhyphen{} NVT for body particles

\item {} 
\sphinxAtStartPar
{\hyperref[\detokenize{fix_nh_eff::doc}]{\sphinxcrossref{\DUrole{doc}{nvt/eff}}}} \sphinxhyphen{} NVE for  nuclei and electrons in the electron force field model

\item {} 
\sphinxAtStartPar
{\hyperref[\detokenize{fix_nvt_manifold_rattle::doc}]{\sphinxcrossref{\DUrole{doc}{nvt/manifold/rattle}}}} \sphinxhyphen{} NVT time integration for atoms constrained to a curved surface (manifold)

\item {} 
\sphinxAtStartPar
{\hyperref[\detokenize{fix_nvt_sllod::doc}]{\sphinxcrossref{\DUrole{doc}{nvt/sllod}}}} \sphinxhyphen{} NVT for NEMD with SLLOD equations

\item {} 
\sphinxAtStartPar
{\hyperref[\detokenize{fix_nvt_sllod_eff::doc}]{\sphinxcrossref{\DUrole{doc}{nvt/sllod/eff}}}} \sphinxhyphen{} NVT for NEMD with SLLOD equations for the electron force field model

\item {} 
\sphinxAtStartPar
{\hyperref[\detokenize{fix_nvt_sphere::doc}]{\sphinxcrossref{\DUrole{doc}{nvt/sphere}}}} \sphinxhyphen{} NVT for spherical particles

\item {} 
\sphinxAtStartPar
{\hyperref[\detokenize{fix_nh_uef::doc}]{\sphinxcrossref{\DUrole{doc}{nvt/uef}}}} \sphinxhyphen{} NVT style time integration with diagonal flow

\item {} 
\sphinxAtStartPar
{\hyperref[\detokenize{fix_oneway::doc}]{\sphinxcrossref{\DUrole{doc}{oneway}}}} \sphinxhyphen{} constrain particles on move in one direction

\item {} 
\sphinxAtStartPar
{\hyperref[\detokenize{fix_orient::doc}]{\sphinxcrossref{\DUrole{doc}{orient/bcc}}}} \sphinxhyphen{} add grain boundary migration force for BCC

\item {} 
\sphinxAtStartPar
{\hyperref[\detokenize{fix_orient::doc}]{\sphinxcrossref{\DUrole{doc}{orient/fcc}}}} \sphinxhyphen{} add grain boundary migration force for FCC

\item {} 
\sphinxAtStartPar
{\hyperref[\detokenize{fix_orient_eco::doc}]{\sphinxcrossref{\DUrole{doc}{orient/eco}}}} \sphinxhyphen{} add generalized grain boundary migration force

\item {} 
\sphinxAtStartPar
{\hyperref[\detokenize{fix_pafi::doc}]{\sphinxcrossref{\DUrole{doc}{pafi}}}} \sphinxhyphen{} constrained force averages on hyper\sphinxhyphen{}planes to compute free energies (PAFI)

\item {} 
\sphinxAtStartPar
{\hyperref[\detokenize{fix_pair::doc}]{\sphinxcrossref{\DUrole{doc}{pair}}}} \sphinxhyphen{} access per\sphinxhyphen{}atom info from pair styles

\item {} 
\sphinxAtStartPar
{\hyperref[\detokenize{fix_phonon::doc}]{\sphinxcrossref{\DUrole{doc}{phonon}}}} \sphinxhyphen{} calculate dynamical matrix from MD simulations

\item {} 
\sphinxAtStartPar
{\hyperref[\detokenize{fix_pimd::doc}]{\sphinxcrossref{\DUrole{doc}{pimd/langevin}}}} \sphinxhyphen{} Feynman path\sphinxhyphen{}integral molecular dynamics with stochastic thermostat

\item {} 
\sphinxAtStartPar
{\hyperref[\detokenize{fix_pimd::doc}]{\sphinxcrossref{\DUrole{doc}{pimd/nvt}}}} \sphinxhyphen{} Feynman path\sphinxhyphen{}integral molecular dynamics with Nose\sphinxhyphen{}Hoover thermostat

\item {} 
\sphinxAtStartPar
{\hyperref[\detokenize{fix_planeforce::doc}]{\sphinxcrossref{\DUrole{doc}{planeforce}}}} \sphinxhyphen{} constrain atoms to move in a plane

\item {} 
\sphinxAtStartPar
{\hyperref[\detokenize{fix_plumed::doc}]{\sphinxcrossref{\DUrole{doc}{plumed}}}} \sphinxhyphen{} wrapper on PLUMED free energy library

\item {} 
\sphinxAtStartPar
{\hyperref[\detokenize{fix_poems::doc}]{\sphinxcrossref{\DUrole{doc}{poems}}}} \sphinxhyphen{} constrain clusters of atoms to move as coupled rigid bodies

\item {} 
\sphinxAtStartPar
{\hyperref[\detokenize{fix_polarize::doc}]{\sphinxcrossref{\DUrole{doc}{polarize/bem/gmres}}}} \sphinxhyphen{} compute induced charges at the interface between impermeable media with different dielectric constants with generalized minimum residual (GMRES)

\item {} 
\sphinxAtStartPar
{\hyperref[\detokenize{fix_polarize::doc}]{\sphinxcrossref{\DUrole{doc}{polarize/bem/icc}}}} \sphinxhyphen{} compute induced charges at the interface between impermeable media with different dielectric constants with the successive over\sphinxhyphen{}relaxation algorithm

\item {} 
\sphinxAtStartPar
{\hyperref[\detokenize{fix_polarize::doc}]{\sphinxcrossref{\DUrole{doc}{polarize/functional}}}} \sphinxhyphen{} compute induced charges at the interface between impermeable media with different dielectric constants with the energy variational approach

\item {} 
\sphinxAtStartPar
{\hyperref[\detokenize{fix_pour::doc}]{\sphinxcrossref{\DUrole{doc}{pour}}}} \sphinxhyphen{} pour new atoms/molecules into a granular simulation domain

\item {} 
\sphinxAtStartPar
{\hyperref[\detokenize{fix_precession_spin::doc}]{\sphinxcrossref{\DUrole{doc}{precession/spin}}}} \sphinxhyphen{} apply a precession torque to each magnetic spin

\item {} 
\sphinxAtStartPar
{\hyperref[\detokenize{fix_press_berendsen::doc}]{\sphinxcrossref{\DUrole{doc}{press/berendsen}}}} \sphinxhyphen{} pressure control by Berendsen barostat

\item {} 
\sphinxAtStartPar
{\hyperref[\detokenize{fix_press_langevin::doc}]{\sphinxcrossref{\DUrole{doc}{press/langevin}}}} \sphinxhyphen{} pressure control by Langevin barostat

\item {} 
\sphinxAtStartPar
{\hyperref[\detokenize{fix_print::doc}]{\sphinxcrossref{\DUrole{doc}{print}}}} \sphinxhyphen{} print text and variables during a simulation

\item {} 
\sphinxAtStartPar
{\hyperref[\detokenize{fix_propel_self::doc}]{\sphinxcrossref{\DUrole{doc}{propel/self}}}} \sphinxhyphen{} model self\sphinxhyphen{}propelled particles

\item {} 
\sphinxAtStartPar
{\hyperref[\detokenize{fix_property_atom::doc}]{\sphinxcrossref{\DUrole{doc}{property/atom}}}} \sphinxhyphen{} add customized per\sphinxhyphen{}atom values

\item {} 
\sphinxAtStartPar
{\hyperref[\detokenize{fix_python_invoke::doc}]{\sphinxcrossref{\DUrole{doc}{python/invoke}}}} \sphinxhyphen{} call a Python function during a simulation

\item {} 
\sphinxAtStartPar
{\hyperref[\detokenize{fix_python_move::doc}]{\sphinxcrossref{\DUrole{doc}{python/move}}}} \sphinxhyphen{} move particles using a Python function during a simulation run

\item {} 
\sphinxAtStartPar
{\hyperref[\detokenize{fix_qbmsst::doc}]{\sphinxcrossref{\DUrole{doc}{qbmsst}}}} \sphinxhyphen{} quantum bath multi\sphinxhyphen{}scale shock technique time integrator

\item {} 
\sphinxAtStartPar
{\hyperref[\detokenize{fix_qeq_comb::doc}]{\sphinxcrossref{\DUrole{doc}{qeq/comb}}}} \sphinxhyphen{} charge equilibration for COMB potential

\item {} 
\sphinxAtStartPar
{\hyperref[\detokenize{fix_qeq::doc}]{\sphinxcrossref{\DUrole{doc}{qeq/dynamic}}}} \sphinxhyphen{} charge equilibration via dynamic method

\item {} 
\sphinxAtStartPar
{\hyperref[\detokenize{fix_qeq::doc}]{\sphinxcrossref{\DUrole{doc}{qeq/fire}}}} \sphinxhyphen{} charge equilibration via FIRE minimizer

\item {} 
\sphinxAtStartPar
{\hyperref[\detokenize{fix_qeq::doc}]{\sphinxcrossref{\DUrole{doc}{qeq/point}}}} \sphinxhyphen{} charge equilibration via point method

\item {} 
\sphinxAtStartPar
{\hyperref[\detokenize{fix_qeq_reaxff::doc}]{\sphinxcrossref{\DUrole{doc}{qeq/reaxff}}}} \sphinxhyphen{} charge equilibration for ReaxFF potential

\item {} 
\sphinxAtStartPar
{\hyperref[\detokenize{fix_qeq::doc}]{\sphinxcrossref{\DUrole{doc}{qeq/shielded}}}} \sphinxhyphen{} charge equilibration via shielded method

\item {} 
\sphinxAtStartPar
{\hyperref[\detokenize{fix_qeq::doc}]{\sphinxcrossref{\DUrole{doc}{qeq/slater}}}} \sphinxhyphen{} charge equilibration via Slater method

\item {} 
\sphinxAtStartPar
{\hyperref[\detokenize{fix_qmmm::doc}]{\sphinxcrossref{\DUrole{doc}{qmmm}}}} \sphinxhyphen{} functionality to enable a quantum mechanics/molecular mechanics coupling

\item {} 
\sphinxAtStartPar
{\hyperref[\detokenize{fix_qtb::doc}]{\sphinxcrossref{\DUrole{doc}{qtb}}}} \sphinxhyphen{} implement quantum thermal bath scheme

\item {} 
\sphinxAtStartPar
{\hyperref[\detokenize{fix_shake::doc}]{\sphinxcrossref{\DUrole{doc}{rattle}}}} \sphinxhyphen{} RATTLE constraints on bonds and/or angles

\item {} 
\sphinxAtStartPar
{\hyperref[\detokenize{fix_reaxff_bonds::doc}]{\sphinxcrossref{\DUrole{doc}{reaxff/bonds}}}} \sphinxhyphen{} write out ReaxFF bond information

\item {} 
\sphinxAtStartPar
{\hyperref[\detokenize{fix_reaxff_species::doc}]{\sphinxcrossref{\DUrole{doc}{reaxff/species}}}} \sphinxhyphen{} write out ReaxFF molecule information

\item {} 
\sphinxAtStartPar
{\hyperref[\detokenize{fix_recenter::doc}]{\sphinxcrossref{\DUrole{doc}{recenter}}}} \sphinxhyphen{} constrain the center\sphinxhyphen{}of\sphinxhyphen{}mass position of a group of atoms

\item {} 
\sphinxAtStartPar
{\hyperref[\detokenize{fix_restrain::doc}]{\sphinxcrossref{\DUrole{doc}{restrain}}}} \sphinxhyphen{} constrain a bond, angle, dihedral

\item {} 
\sphinxAtStartPar
{\hyperref[\detokenize{fix_rheo::doc}]{\sphinxcrossref{\DUrole{doc}{rheo}}}} \sphinxhyphen{} integrator for the RHEO package

\item {} 
\sphinxAtStartPar
{\hyperref[\detokenize{fix_rheo_thermal::doc}]{\sphinxcrossref{\DUrole{doc}{rheo/thermal}}}} \sphinxhyphen{} thermal integrator for the RHEO package

\item {} 
\sphinxAtStartPar
{\hyperref[\detokenize{fix_rheo_oxidation::doc}]{\sphinxcrossref{\DUrole{doc}{rheo/oxidation}}}} \sphinxhyphen{} create oxidation bonds for the RHEO package

\item {} 
\sphinxAtStartPar
{\hyperref[\detokenize{fix_rheo_pressure::doc}]{\sphinxcrossref{\DUrole{doc}{rheo/pressure}}}} \sphinxhyphen{} pressure calculation for the RHEO package

\item {} 
\sphinxAtStartPar
{\hyperref[\detokenize{fix_rheo_pressure::doc}]{\sphinxcrossref{\DUrole{doc}{rheo/viscosity}}}} \sphinxhyphen{} viscosity calculation for the RHEO package

\item {} 
\sphinxAtStartPar
{\hyperref[\detokenize{fix_rhok::doc}]{\sphinxcrossref{\DUrole{doc}{rhok}}}} \sphinxhyphen{} add bias potential for long\sphinxhyphen{}range ordered systems

\item {} 
\sphinxAtStartPar
{\hyperref[\detokenize{fix_rigid::doc}]{\sphinxcrossref{\DUrole{doc}{rigid}}}} \sphinxhyphen{} constrain one or more clusters of atoms to move as a rigid body with NVE integration

\item {} 
\sphinxAtStartPar
{\hyperref[\detokenize{fix_rigid_meso::doc}]{\sphinxcrossref{\DUrole{doc}{rigid/meso}}}} \sphinxhyphen{} constrain clusters of mesoscopic SPH/SDPD particles to move as a rigid body

\item {} 
\sphinxAtStartPar
{\hyperref[\detokenize{fix_rigid::doc}]{\sphinxcrossref{\DUrole{doc}{rigid/nph}}}} \sphinxhyphen{} constrain one or more clusters of atoms to move as a rigid body with NPH integration

\item {} 
\sphinxAtStartPar
{\hyperref[\detokenize{fix_rigid::doc}]{\sphinxcrossref{\DUrole{doc}{rigid/nph/small}}}} \sphinxhyphen{} constrain many small clusters of atoms to move as a rigid body with NPH integration

\item {} 
\sphinxAtStartPar
{\hyperref[\detokenize{fix_rigid::doc}]{\sphinxcrossref{\DUrole{doc}{rigid/npt}}}} \sphinxhyphen{} constrain one or more clusters of atoms to move as a rigid body with NPT integration

\item {} 
\sphinxAtStartPar
{\hyperref[\detokenize{fix_rigid::doc}]{\sphinxcrossref{\DUrole{doc}{rigid/npt/small}}}} \sphinxhyphen{} constrain many small clusters of atoms to move as a rigid body with NPT integration

\item {} 
\sphinxAtStartPar
{\hyperref[\detokenize{fix_rigid::doc}]{\sphinxcrossref{\DUrole{doc}{rigid/nve}}}} \sphinxhyphen{} constrain one or more clusters of atoms to move as a rigid body with alternate NVE integration

\item {} 
\sphinxAtStartPar
{\hyperref[\detokenize{fix_rigid::doc}]{\sphinxcrossref{\DUrole{doc}{rigid/nve/small}}}} \sphinxhyphen{} constrain many small clusters of atoms to move as a rigid body with alternate NVE integration

\item {} 
\sphinxAtStartPar
{\hyperref[\detokenize{fix_rigid::doc}]{\sphinxcrossref{\DUrole{doc}{rigid/nvt}}}} \sphinxhyphen{} constrain one or more clusters of atoms to move as a rigid body with NVT integration

\item {} 
\sphinxAtStartPar
{\hyperref[\detokenize{fix_rigid::doc}]{\sphinxcrossref{\DUrole{doc}{rigid/nvt/small}}}} \sphinxhyphen{} constrain many small clusters of atoms to move as a rigid body with NVT integration

\item {} 
\sphinxAtStartPar
{\hyperref[\detokenize{fix_rigid::doc}]{\sphinxcrossref{\DUrole{doc}{rigid/small}}}} \sphinxhyphen{} constrain many small clusters of atoms to move as a rigid body with NVE integration

\item {} 
\sphinxAtStartPar
{\hyperref[\detokenize{fix_rx::doc}]{\sphinxcrossref{\DUrole{doc}{rx}}}} \sphinxhyphen{} solve reaction kinetic ODEs for a defined reaction set

\item {} 
\sphinxAtStartPar
{\hyperref[\detokenize{fix_saed_vtk::doc}]{\sphinxcrossref{\DUrole{doc}{saed/vtk}}}} \sphinxhyphen{} time\sphinxhyphen{}average the intensities from {\hyperref[\detokenize{compute_saed::doc}]{\sphinxcrossref{\DUrole{doc}{compute saed}}}}

\item {} 
\sphinxAtStartPar
{\hyperref[\detokenize{fix_setforce::doc}]{\sphinxcrossref{\DUrole{doc}{setforce}}}} \sphinxhyphen{} set the force on each atom

\item {} 
\sphinxAtStartPar
{\hyperref[\detokenize{fix_setforce::doc}]{\sphinxcrossref{\DUrole{doc}{setforce/spin}}}} \sphinxhyphen{} set magnetic precession vectors on each atom

\item {} 
\sphinxAtStartPar
{\hyperref[\detokenize{fix_sgcmc::doc}]{\sphinxcrossref{\DUrole{doc}{sgcmc}}}} \sphinxhyphen{} fix for hybrid semi\sphinxhyphen{}grand canonical MD/MC simulations

\item {} 
\sphinxAtStartPar
{\hyperref[\detokenize{fix_shake::doc}]{\sphinxcrossref{\DUrole{doc}{shake}}}} \sphinxhyphen{} SHAKE constraints on bonds and/or angles

\item {} 
\sphinxAtStartPar
{\hyperref[\detokenize{fix_shardlow::doc}]{\sphinxcrossref{\DUrole{doc}{shardlow}}}} \sphinxhyphen{} integration of DPD equations of motion using the Shardlow splitting

\item {} 
\sphinxAtStartPar
{\hyperref[\detokenize{fix_smd::doc}]{\sphinxcrossref{\DUrole{doc}{smd}}}} \sphinxhyphen{} applied a steered MD force to a group

\item {} 
\sphinxAtStartPar
{\hyperref[\detokenize{fix_smd_adjust_dt::doc}]{\sphinxcrossref{\DUrole{doc}{smd/adjust\_dt}}}} \sphinxhyphen{} calculate a new stable time increment for use with SMD integrators

\item {} 
\sphinxAtStartPar
{\hyperref[\detokenize{fix_smd_integrate_tlsph::doc}]{\sphinxcrossref{\DUrole{doc}{smd/integrate\_tlsph}}}} \sphinxhyphen{} explicit time integration with total Lagrangian SPH pair style

\item {} 
\sphinxAtStartPar
{\hyperref[\detokenize{fix_smd_integrate_ulsph::doc}]{\sphinxcrossref{\DUrole{doc}{smd/integrate\_ulsph}}}} \sphinxhyphen{} explicit time integration with updated Lagrangian SPH pair style

\item {} 
\sphinxAtStartPar
{\hyperref[\detokenize{fix_smd_move_triangulated_surface::doc}]{\sphinxcrossref{\DUrole{doc}{smd/move\_tri\_surf}}}} \sphinxhyphen{} update position and velocity near rigid surfaces using SPH integrators

\item {} 
\sphinxAtStartPar
{\hyperref[\detokenize{fix_smd_setvel::doc}]{\sphinxcrossref{\DUrole{doc}{smd/setvel}}}} \sphinxhyphen{} sets each velocity component, ignoring forces, for Smooth Mach Dynamics

\item {} 
\sphinxAtStartPar
{\hyperref[\detokenize{fix_smd_wall_surface::doc}]{\sphinxcrossref{\DUrole{doc}{smd/wall\_surface}}}} \sphinxhyphen{} create a rigid wall with a triangulated surface for use in Smooth Mach Dynamics

\item {} 
\sphinxAtStartPar
{\hyperref[\detokenize{fix_sph::doc}]{\sphinxcrossref{\DUrole{doc}{sph}}}} \sphinxhyphen{} time integration for SPH/DPDE particles

\item {} 
\sphinxAtStartPar
{\hyperref[\detokenize{fix_sph_stationary::doc}]{\sphinxcrossref{\DUrole{doc}{sph/stationary}}}} \sphinxhyphen{} update energy and density but not position or velocity in Smooth Particle Hydrodynamics

\item {} 
\sphinxAtStartPar
{\hyperref[\detokenize{fix_spring::doc}]{\sphinxcrossref{\DUrole{doc}{spring}}}} \sphinxhyphen{} apply harmonic spring force to group of atoms

\item {} 
\sphinxAtStartPar
{\hyperref[\detokenize{fix_spring_chunk::doc}]{\sphinxcrossref{\DUrole{doc}{spring/chunk}}}} \sphinxhyphen{} apply harmonic spring force to each chunk of atoms

\item {} 
\sphinxAtStartPar
{\hyperref[\detokenize{fix_spring_rg::doc}]{\sphinxcrossref{\DUrole{doc}{spring/rg}}}} \sphinxhyphen{} spring on radius of gyration of group of atoms

\item {} 
\sphinxAtStartPar
{\hyperref[\detokenize{fix_spring_self::doc}]{\sphinxcrossref{\DUrole{doc}{spring/self}}}} \sphinxhyphen{} spring from each atom to its origin

\item {} 
\sphinxAtStartPar
{\hyperref[\detokenize{fix_srd::doc}]{\sphinxcrossref{\DUrole{doc}{srd}}}} \sphinxhyphen{} stochastic rotation dynamics (SRD)

\item {} 
\sphinxAtStartPar
{\hyperref[\detokenize{fix_store_force::doc}]{\sphinxcrossref{\DUrole{doc}{store/force}}}} \sphinxhyphen{} store force on each atom

\item {} 
\sphinxAtStartPar
{\hyperref[\detokenize{fix_store_state::doc}]{\sphinxcrossref{\DUrole{doc}{store/state}}}} \sphinxhyphen{} store attributes for each atom

\item {} 
\sphinxAtStartPar
{\hyperref[\detokenize{fix_dpd_source::doc}]{\sphinxcrossref{\DUrole{doc}{tdpd/source}}}} \sphinxhyphen{} add external concentration source

\item {} 
\sphinxAtStartPar
{\hyperref[\detokenize{fix_temp_berendsen::doc}]{\sphinxcrossref{\DUrole{doc}{temp/berendsen}}}} \sphinxhyphen{} temperature control by Berendsen thermostat

\item {} 
\sphinxAtStartPar
{\hyperref[\detokenize{fix_temp_csvr::doc}]{\sphinxcrossref{\DUrole{doc}{temp/csld}}}} \sphinxhyphen{} canonical sampling thermostat with Langevin dynamics

\item {} 
\sphinxAtStartPar
{\hyperref[\detokenize{fix_temp_csvr::doc}]{\sphinxcrossref{\DUrole{doc}{temp/csvr}}}} \sphinxhyphen{} canonical sampling thermostat with Hamiltonian dynamics

\item {} 
\sphinxAtStartPar
{\hyperref[\detokenize{fix_temp_rescale::doc}]{\sphinxcrossref{\DUrole{doc}{temp/rescale}}}} \sphinxhyphen{} temperature control by velocity rescaling

\item {} 
\sphinxAtStartPar
{\hyperref[\detokenize{fix_temp_rescale_eff::doc}]{\sphinxcrossref{\DUrole{doc}{temp/rescale/eff}}}} \sphinxhyphen{} temperature control by velocity rescaling in the electron force field model

\item {} 
\sphinxAtStartPar
{\hyperref[\detokenize{fix_tfmc::doc}]{\sphinxcrossref{\DUrole{doc}{tfmc}}}} \sphinxhyphen{} perform force\sphinxhyphen{}bias Monte Carlo with time\sphinxhyphen{}stamped method

\item {} 
\sphinxAtStartPar
{\hyperref[\detokenize{fix_tgnh_drude::doc}]{\sphinxcrossref{\DUrole{doc}{tgnvt/drude}}}} \sphinxhyphen{} NVT time integration for Drude polarizable model via temperature\sphinxhyphen{}grouped Nose\sphinxhyphen{}Hoover

\item {} 
\sphinxAtStartPar
{\hyperref[\detokenize{fix_tgnh_drude::doc}]{\sphinxcrossref{\DUrole{doc}{tgnpt/drude}}}} \sphinxhyphen{} NPT time integration for Drude polarizable model via temperature\sphinxhyphen{}grouped Nose\sphinxhyphen{}Hoover

\item {} 
\sphinxAtStartPar
{\hyperref[\detokenize{fix_thermal_conductivity::doc}]{\sphinxcrossref{\DUrole{doc}{thermal/conductivity}}}} \sphinxhyphen{} Mueller\sphinxhyphen{}Plathe kinetic energy exchange for thermal conductivity calculation

\item {} 
\sphinxAtStartPar
{\hyperref[\detokenize{fix_ti_spring::doc}]{\sphinxcrossref{\DUrole{doc}{ti/spring}}}} \sphinxhyphen{} perform thermodynamic integration between a solid and an Einstein crystal

\item {} 
\sphinxAtStartPar
{\hyperref[\detokenize{fix_tmd::doc}]{\sphinxcrossref{\DUrole{doc}{tmd}}}} \sphinxhyphen{} guide a group of atoms to a new configuration

\item {} 
\sphinxAtStartPar
{\hyperref[\detokenize{fix_ttm::doc}]{\sphinxcrossref{\DUrole{doc}{ttm}}}} \sphinxhyphen{} two\sphinxhyphen{}temperature model for electronic/atomic coupling (replicated grid)

\item {} 
\sphinxAtStartPar
{\hyperref[\detokenize{fix_ttm::doc}]{\sphinxcrossref{\DUrole{doc}{ttm/grid}}}} \sphinxhyphen{} two\sphinxhyphen{}temperature model for electronic/atomic coupling (distributed grid)

\item {} 
\sphinxAtStartPar
{\hyperref[\detokenize{fix_ttm::doc}]{\sphinxcrossref{\DUrole{doc}{ttm/mod}}}} \sphinxhyphen{} enhanced two\sphinxhyphen{}temperature model with additional options

\item {} 
\sphinxAtStartPar
{\hyperref[\detokenize{fix_tune_kspace::doc}]{\sphinxcrossref{\DUrole{doc}{tune/kspace}}}} \sphinxhyphen{} auto\sphinxhyphen{}tune \(k\)\sphinxhyphen{}space parameters

\item {} 
\sphinxAtStartPar
{\hyperref[\detokenize{fix_vector::doc}]{\sphinxcrossref{\DUrole{doc}{vector}}}} \sphinxhyphen{} accumulate a global vector every \sphinxstyleemphasis{N} timesteps

\item {} 
\sphinxAtStartPar
{\hyperref[\detokenize{fix_viscosity::doc}]{\sphinxcrossref{\DUrole{doc}{viscosity}}}} \sphinxhyphen{} Mueller\sphinxhyphen{}Plathe momentum exchange for viscosity calculation

\item {} 
\sphinxAtStartPar
{\hyperref[\detokenize{fix_viscous::doc}]{\sphinxcrossref{\DUrole{doc}{viscous}}}} \sphinxhyphen{} viscous damping for granular simulations

\item {} 
\sphinxAtStartPar
{\hyperref[\detokenize{fix_viscous_sphere::doc}]{\sphinxcrossref{\DUrole{doc}{viscous/sphere}}}} \sphinxhyphen{} viscous damping on angular velocity for granular simulations

\item {} 
\sphinxAtStartPar
{\hyperref[\detokenize{fix_wall_body_polygon::doc}]{\sphinxcrossref{\DUrole{doc}{wall/body/polygon}}}} \sphinxhyphen{} time integration for body particles of style {\hyperref[\detokenize{Howto_body::doc}]{\sphinxcrossref{\DUrole{doc}{rounded/polygon}}}}

\item {} 
\sphinxAtStartPar
{\hyperref[\detokenize{fix_wall_body_polyhedron::doc}]{\sphinxcrossref{\DUrole{doc}{wall/body/polyhedron}}}} \sphinxhyphen{} time integration for body particles of style {\hyperref[\detokenize{Howto_body::doc}]{\sphinxcrossref{\DUrole{doc}{rounded/polyhedron}}}}

\item {} 
\sphinxAtStartPar
{\hyperref[\detokenize{fix_wall::doc}]{\sphinxcrossref{\DUrole{doc}{wall/colloid}}}} \sphinxhyphen{} Lennard\sphinxhyphen{}Jones wall interacting with finite\sphinxhyphen{}size particles

\item {} 
\sphinxAtStartPar
{\hyperref[\detokenize{fix_wall_ees::doc}]{\sphinxcrossref{\DUrole{doc}{wall/ees}}}} \sphinxhyphen{} wall for ellipsoidal particles

\item {} 
\sphinxAtStartPar
{\hyperref[\detokenize{fix_wall_flow::doc}]{\sphinxcrossref{\DUrole{doc}{wall/flow}}}} \sphinxhyphen{} flow boundary conditions

\item {} 
\sphinxAtStartPar
{\hyperref[\detokenize{fix_wall_gran::doc}]{\sphinxcrossref{\DUrole{doc}{wall/gran}}}} \sphinxhyphen{} frictional wall(s) for granular simulations

\item {} 
\sphinxAtStartPar
{\hyperref[\detokenize{fix_wall_gran_region::doc}]{\sphinxcrossref{\DUrole{doc}{wall/gran/region}}}} \sphinxhyphen{} {\hyperref[\detokenize{fix_wall_region::doc}]{\sphinxcrossref{\DUrole{doc}{fix wall/region}}}} equivalent for use with granular particles

\item {} 
\sphinxAtStartPar
{\hyperref[\detokenize{fix_wall::doc}]{\sphinxcrossref{\DUrole{doc}{wall/harmonic}}}} \sphinxhyphen{} harmonic spring wall

\item {} 
\sphinxAtStartPar
{\hyperref[\detokenize{fix_wall::doc}]{\sphinxcrossref{\DUrole{doc}{wall/lj1043}}}} \sphinxhyphen{} Lennard\sphinxhyphen{}Jones 10\textendash{}4\textendash{}3 wall

\item {} 
\sphinxAtStartPar
{\hyperref[\detokenize{fix_wall::doc}]{\sphinxcrossref{\DUrole{doc}{wall/lj126}}}} \sphinxhyphen{} Lennard\sphinxhyphen{}Jones 12\textendash{}6 wall

\item {} 
\sphinxAtStartPar
{\hyperref[\detokenize{fix_wall::doc}]{\sphinxcrossref{\DUrole{doc}{wall/lj93}}}} \sphinxhyphen{} Lennard\sphinxhyphen{}Jones 9\textendash{}3 wall

\item {} 
\sphinxAtStartPar
{\hyperref[\detokenize{fix_wall::doc}]{\sphinxcrossref{\DUrole{doc}{wall/lepton}}}} \sphinxhyphen{} Custom Lepton expression wall

\item {} 
\sphinxAtStartPar
{\hyperref[\detokenize{fix_wall::doc}]{\sphinxcrossref{\DUrole{doc}{wall/morse}}}} \sphinxhyphen{} Morse potential wall

\item {} 
\sphinxAtStartPar
{\hyperref[\detokenize{fix_wall_piston::doc}]{\sphinxcrossref{\DUrole{doc}{wall/piston}}}} \sphinxhyphen{} moving reflective piston wall

\item {} 
\sphinxAtStartPar
{\hyperref[\detokenize{fix_wall_reflect::doc}]{\sphinxcrossref{\DUrole{doc}{wall/reflect}}}} \sphinxhyphen{} reflecting wall(s)

\item {} 
\sphinxAtStartPar
{\hyperref[\detokenize{fix_wall_reflect_stochastic::doc}]{\sphinxcrossref{\DUrole{doc}{wall/reflect/stochastic}}}} \sphinxhyphen{} reflecting wall(s) with finite temperature

\item {} 
\sphinxAtStartPar
{\hyperref[\detokenize{fix_wall_region::doc}]{\sphinxcrossref{\DUrole{doc}{wall/region}}}} \sphinxhyphen{} use region surface as wall

\item {} 
\sphinxAtStartPar
{\hyperref[\detokenize{fix_wall_ees::doc}]{\sphinxcrossref{\DUrole{doc}{wall/region/ees}}}} \sphinxhyphen{} use region surface as wall for ellipsoidal particles

\item {} 
\sphinxAtStartPar
{\hyperref[\detokenize{fix_wall_srd::doc}]{\sphinxcrossref{\DUrole{doc}{wall/srd}}}} \sphinxhyphen{} slip/no\sphinxhyphen{}slip wall for SRD particles

\item {} 
\sphinxAtStartPar
{\hyperref[\detokenize{fix_wall::doc}]{\sphinxcrossref{\DUrole{doc}{wall/table}}}} \sphinxhyphen{} Tabulated potential wall wall

\item {} 
\sphinxAtStartPar
{\hyperref[\detokenize{fix_widom::doc}]{\sphinxcrossref{\DUrole{doc}{widom}}}} \sphinxhyphen{} Widom insertions of atoms or molecules

\end{itemize}


\subsection{Restrictions}
\label{\detokenize{fix:restrictions}}
\sphinxAtStartPar
Some fix styles are part of specific packages.  They are only enabled
if LAMMPS was built with that package.  See the
{\hyperref[\detokenize{Build_package::doc}]{\sphinxcrossref{\DUrole{doc}{Build package}}}} page for more info.  The doc pages for
individual fixes tell if it is part of a package.


\subsection{Related commands}
\label{\detokenize{fix:related-commands}}
\sphinxAtStartPar
{\hyperref[\detokenize{unfix::doc}]{\sphinxcrossref{\DUrole{doc}{unfix}}}}, {\hyperref[\detokenize{fix_modify::doc}]{\sphinxcrossref{\DUrole{doc}{fix\_modify}}}}


\subsection{Default}
\label{\detokenize{fix:default}}
\sphinxAtStartPar
none

\sphinxstepscope

\index{fix\_modify@\spxentry{fix\_modify}}\ignorespaces 

\section{fix\_modify command}
\label{\detokenize{fix_modify:fix-modify-command}}\label{\detokenize{fix_modify:index-0}}\label{\detokenize{fix_modify::doc}}

\subsection{Syntax}
\label{\detokenize{fix_modify:syntax}}
\begin{sphinxVerbatim}[commandchars=\\\{\}]
\PYG{k}{fix\PYGZus{}modify }\PYG{n+nv+nvIdentifier}{fix\PYGZhy{}ID}\PYG{+w}{ }\PYG{n}{keyword}\PYG{+w}{ }\PYG{n}{value}\PYG{+w}{ }\PYG{n}{...}
\end{sphinxVerbatim}
\begin{itemize}
\item {} 
\sphinxAtStartPar
fix\sphinxhyphen{}ID = ID of the fix to modify

\item {} 
\sphinxAtStartPar
one or more keyword/value pairs may be appended

\item {} 
\sphinxAtStartPar
keyword = \sphinxstyleemphasis{bodyforces} or \sphinxstyleemphasis{colname} or \sphinxstyleemphasis{dynamic/dof} or \sphinxstyleemphasis{energy} or \sphinxstyleemphasis{press} or \sphinxstyleemphasis{respa} or \sphinxstyleemphasis{temp} or \sphinxstyleemphasis{virial}
\begin{sphinxalltt}
\sphinxstyleemphasis{bodyforces} value = \sphinxstyleemphasis{early} or \sphinxstyleemphasis{late}
  early/late = compute rigid\sphinxhyphen{}body forces/torques early or late in the timestep
\sphinxstyleemphasis{colname} values =  ID string
  string = new column header name
  ID = integer from 1 to N, or integer from \sphinxhyphen{}1 to \sphinxhyphen{}N, where N = \# of quantities being output
       \sphinxstyleemphasis{or} a fix output property keyword or reference to compute, fix, property or variable.
\sphinxstyleemphasis{dynamic/dof} value = \sphinxstyleemphasis{yes} or \sphinxstyleemphasis{no}
  yes/no = do or do not re\sphinxhyphen{}compute the number of degrees of freedom (DOF) contributing to the temperature
\sphinxstyleemphasis{energy} value = \sphinxstyleemphasis{yes} or \sphinxstyleemphasis{no}
\sphinxstyleemphasis{press} value = compute ID that calculates a pressure
\sphinxstyleemphasis{respa} value = \sphinxstyleemphasis{1} to \sphinxstyleemphasis{max respa level} or \sphinxstyleemphasis{0} (for outermost level)
\sphinxstyleemphasis{temp} value = compute ID that calculates a temperature
\sphinxstyleemphasis{virial} value = \sphinxstyleemphasis{yes} or \sphinxstyleemphasis{no}
\end{sphinxalltt}

\end{itemize}


\subsection{Examples}
\label{\detokenize{fix_modify:examples}}
\begin{sphinxVerbatim}[commandchars=\\\{\}]
\PYG{k}{fix\PYGZus{}modify }\PYG{n+nv+nvIdentifier}{3}\PYG{+w}{ }\PYG{n}{temp}\PYG{+w}{ }\PYG{n}{myTemp}\PYG{+w}{ }\PYG{n}{press}\PYG{+w}{ }\PYG{n}{myPress}
\PYG{k}{fix\PYGZus{}modify }\PYG{n+nv+nvIdentifier}{1}\PYG{+w}{ }\PYG{n}{energy}\PYG{+w}{ }\PYG{n}{yes}
\PYG{k}{fix\PYGZus{}modify }\PYG{n+nv+nvIdentifier}{tether}\PYG{+w}{ }\PYG{n}{respa}\PYG{+w}{ }\PYG{l+m}{2}
\PYG{k}{fix\PYGZus{}modify }\PYG{n+nv+nvIdentifier}{ave}\PYG{+w}{ }\PYG{n}{colname}\PYG{+w}{ }\PYG{n}{c\PYGZus{}thermo\PYGZus{}press}\PYG{+w}{ }\PYG{n}{Pressure}\PYG{+w}{ }\PYG{n}{colname}\PYG{+w}{ }\PYG{l+m}{1}\PYG{+w}{ }\PYG{n}{Temperature}
\end{sphinxVerbatim}


\subsection{Description}
\label{\detokenize{fix_modify:description}}
\sphinxAtStartPar
Modify one or more parameters of a previously defined fix.  Only
specific fix styles support specific parameters.  See the doc pages
for individual fix commands for info on which ones support which
fix\_modify parameters.

\sphinxAtStartPar
The \sphinxstyleemphasis{temp} keyword is used to determine how a fix computes
temperature.  The specified compute ID must have been previously
defined by the user via the {\hyperref[\detokenize{compute::doc}]{\sphinxcrossref{\DUrole{doc}{compute}}}} command and it must
be a style of compute that calculates a temperature.  All fixes that
compute temperatures define their own compute by default, as described
in their documentation.  Thus this option allows the user to override
the default method for computing T.

\sphinxAtStartPar
The \sphinxstyleemphasis{press} keyword is used to determine how a fix computes pressure.
The specified compute ID must have been previously defined by the user
via the {\hyperref[\detokenize{compute::doc}]{\sphinxcrossref{\DUrole{doc}{compute}}}} command and it must be a style of
compute that calculates a pressure.  All fixes that compute pressures
define their own compute by default, as described in their
documentation.  Thus this option allows the user to override the
default method for computing P.

\sphinxAtStartPar
The \sphinxstyleemphasis{energy} keyword can be used with fixes that support it, which is
explained at the bottom of their doc page.  \sphinxstyleemphasis{Energy yes} will add a
contribution to the potential energy of the system.  More
specifically, the fix’s global or per\sphinxhyphen{}atom energy is included in the
calculation performed by the {\hyperref[\detokenize{compute_pe::doc}]{\sphinxcrossref{\DUrole{doc}{compute pe}}}} or
{\hyperref[\detokenize{compute_pe_atom::doc}]{\sphinxcrossref{\DUrole{doc}{compute pe/atom}}}} commands.  The former is what
is used the {\hyperref[\detokenize{thermo_style::doc}]{\sphinxcrossref{\DUrole{doc}{thermo\_style}}}} command for output of
any quantity that includes the global potential energy of the system.
Note that the {\hyperref[\detokenize{compute_pe::doc}]{\sphinxcrossref{\DUrole{doc}{compute pe}}}} and {\hyperref[\detokenize{compute_pe_atom::doc}]{\sphinxcrossref{\DUrole{doc}{compute pe/atom}}}} commands also have an option to include or exclude
the contribution from fixes.  For fixes that tally a global energy, it
can also be printed with thermodynamic output by using the keyword
f\_ID in the thermo\_style custom command, where ID is the fix\sphinxhyphen{}ID of the
appropriate fix.

\begin{sphinxadmonition}{note}{Note:}
\sphinxAtStartPar
If you are performing an {\hyperref[\detokenize{minimize::doc}]{\sphinxcrossref{\DUrole{doc}{energy minimization}}}} with
one of these fixes and want the energy and forces it produces to be
part of the optimization criteria, you must specify the \sphinxstyleemphasis{energy
yes} setting.
\end{sphinxadmonition}

\begin{sphinxadmonition}{note}{Note:}
\sphinxAtStartPar
For most fixes that support the \sphinxstyleemphasis{energy} keyword, the default
setting is \sphinxstyleemphasis{no}.  For a few it is \sphinxstyleemphasis{yes}, when a user would expect
that to be the case.  The page of each fix gives the default.
\end{sphinxadmonition}

\sphinxAtStartPar
The \sphinxstyleemphasis{virial} keyword can be used with fixes that support it, which is
explained at the bottom of their doc page.  \sphinxstyleemphasis{Virial yes} will add a
contribution to the virial of the system.  More specifically, the
fix’s global or per\sphinxhyphen{}atom virial is included in the calculation
performed by the {\hyperref[\detokenize{compute_pressure::doc}]{\sphinxcrossref{\DUrole{doc}{compute pressure}}}} or
{\hyperref[\detokenize{compute_stress_atom::doc}]{\sphinxcrossref{\DUrole{doc}{compute stress/atom}}}} commands.  The former
is what is used the {\hyperref[\detokenize{thermo_style::doc}]{\sphinxcrossref{\DUrole{doc}{thermo\_style}}}} command for
output of any quantity that includes the global pressure of the
system.  Note that the {\hyperref[\detokenize{compute_pressure::doc}]{\sphinxcrossref{\DUrole{doc}{compute pressure}}}} and
{\hyperref[\detokenize{compute_stress_atom::doc}]{\sphinxcrossref{\DUrole{doc}{compute stress/atom}}}} commands also have an
option to include or exclude the contribution from fixes.

\begin{sphinxadmonition}{note}{Note:}
\sphinxAtStartPar
If you are performing an {\hyperref[\detokenize{minimize::doc}]{\sphinxcrossref{\DUrole{doc}{energy minimization}}}} with
{\hyperref[\detokenize{fix_box_relax::doc}]{\sphinxcrossref{\DUrole{doc}{box relaxation}}}} and one of these fixes and
want the virial contribution of the fix to be part of the
optimization criteria, you must specify the \sphinxstyleemphasis{virial yes} setting.
\end{sphinxadmonition}

\begin{sphinxadmonition}{note}{Note:}
\sphinxAtStartPar
For most fixes that support the \sphinxstyleemphasis{virial} keyword, the default
setting is \sphinxstyleemphasis{no}.  For a few it is \sphinxstyleemphasis{yes}, when a user would expect
that to be the case.  The page of each fix gives the default.
\end{sphinxadmonition}

\sphinxAtStartPar
For fixes that set or modify forces, it may be possible to select at
which {\hyperref[\detokenize{run_style::doc}]{\sphinxcrossref{\DUrole{doc}{r\sphinxhyphen{}RESPA}}}} level the fix operates via the \sphinxstyleemphasis{respa}
keyword. The RESPA level at which the fix is active can be selected.
This is a number ranging from 1 to the number of levels. If the RESPA
level is larger than the current maximum, the outermost level will be
used, which is also the default setting. This default can be restored
using a value of \sphinxstyleemphasis{0} for the RESPA level. The affected fix has to be
enabled to support this feature; if not, \sphinxstyleemphasis{fix\_modify} will report an
error. Active fixes with a custom RESPA level setting are reported
with their specified level at the beginning of a r\sphinxhyphen{}RESPA run.

\sphinxAtStartPar
The \sphinxstyleemphasis{dynamic/dof} keyword determines whether the number of atoms N in
the fix group and their associated degrees of freedom are re\sphinxhyphen{}computed
each time a temperature is computed.  Only fix styles that calculate
their own internal temperature use this option.  Currently this is only
the {\hyperref[\detokenize{fix_rigid::doc}]{\sphinxcrossref{\DUrole{doc}{fix rigid/nvt/small}}}} and {\hyperref[\detokenize{fix_rigid::doc}]{\sphinxcrossref{\DUrole{doc}{fix rigid/npt/small}}}} commands for the purpose of thermostatting rigid body
translation and rotation.  By default, N and their DOF are assumed to be
constant.  If you are adding atoms or molecules to the system (see the
{\hyperref[\detokenize{fix_pour::doc}]{\sphinxcrossref{\DUrole{doc}{fix pour}}}}, {\hyperref[\detokenize{fix_deposit::doc}]{\sphinxcrossref{\DUrole{doc}{fix deposit}}}}, and
{\hyperref[\detokenize{fix_gcmc::doc}]{\sphinxcrossref{\DUrole{doc}{fix gcmc}}}} commands) or expect atoms or molecules to be
lost (e.g. due to exiting the simulation box or via {\hyperref[\detokenize{fix_evaporate::doc}]{\sphinxcrossref{\DUrole{doc}{fix evaporate}}}}), then this option should be used to ensure the
temperature is correctly normalized.

\begin{sphinxadmonition}{note}{Note:}
\sphinxAtStartPar
Other thermostatting fixes, such as {\hyperref[\detokenize{fix_nh::doc}]{\sphinxcrossref{\DUrole{doc}{fix nvt}}}}, do not
use the \sphinxstyleemphasis{dynamic/dof} keyword because they use a temperature compute
to calculate temperature.  See the {\hyperref[\detokenize{compute_modify::doc}]{\sphinxcrossref{\DUrole{doc}{compute\_modify dynamic/dof}}}} command for a similar way to ensure correct
temperature normalization for those thermostats.
\end{sphinxadmonition}

\sphinxAtStartPar
The \sphinxstyleemphasis{bodyforces} keyword determines whether the forces and torques
acting on rigid bodies are computed \sphinxstyleemphasis{early} at the post\sphinxhyphen{}force stage of
each timestep (right after per\sphinxhyphen{}atom forces have been computed and
communicated among processors), or \sphinxstyleemphasis{late} at the final\sphinxhyphen{}integrate stage
of each timestep (after any other fixes have finished their post\sphinxhyphen{}force
tasks).  Only the rigid\sphinxhyphen{}body integration fixes use this option, which
includes {\hyperref[\detokenize{fix_rigid::doc}]{\sphinxcrossref{\DUrole{doc}{fix rigid}}}} and {\hyperref[\detokenize{fix_rigid::doc}]{\sphinxcrossref{\DUrole{doc}{fix rigid/small}}}}, and their variants, and also {\hyperref[\detokenize{fix_poems::doc}]{\sphinxcrossref{\DUrole{doc}{fix poems}}}}.

\sphinxAtStartPar
The default is \sphinxstyleemphasis{late}.  If there are other fixes that add forces to
individual atoms, then the rigid\sphinxhyphen{}body constraints will include these
forces when time\sphinxhyphen{}integrating the rigid bodies.  If \sphinxstyleemphasis{early} is
specified, then new fixes can be written that use or modify the
per\sphinxhyphen{}body force and torque, before time\sphinxhyphen{}integration of the rigid bodies
occurs.  Note however this has the side effect, that fixes such as
{\hyperref[\detokenize{fix_addforce::doc}]{\sphinxcrossref{\DUrole{doc}{fix addforce}}}}, {\hyperref[\detokenize{fix_setforce::doc}]{\sphinxcrossref{\DUrole{doc}{fix setforce}}}},
{\hyperref[\detokenize{fix_spring::doc}]{\sphinxcrossref{\DUrole{doc}{fix spring}}}}, which add forces to individual atoms
will have no effect on the motion of the rigid bodies if they are
specified in the input script after the fix rigid command.  LAMMPS
will give a warning if that is the case.

\sphinxAtStartPar
The \sphinxstyleemphasis{colname} keyword can be used to change the default header keywords
in output files of fix styles that support it: currently only {\hyperref[\detokenize{fix_ave_time::doc}]{\sphinxcrossref{\DUrole{doc}{fix
ave/time}}}} is supported.  The setting for \sphinxstyleemphasis{ID string}
replaces the default text with the provided string.  \sphinxstyleemphasis{ID} can be a
positive integer when it represents the column number counting from the
left, a negative integer when it represents the column number from the
right (i.e. \sphinxhyphen{}1 is the last column/keyword), or a custom fix output
keyword (or compute, fix, property, or variable reference) and then it
replaces the string for that specific keyword. The \sphinxstyleemphasis{colname} keyword can
be used multiple times. If multiple \sphinxstyleemphasis{colname} settings refer to the same
keyword, the last setting has precedence.


\subsection{Restrictions}
\label{\detokenize{fix_modify:restrictions}}
\sphinxAtStartPar
none


\subsection{Related commands}
\label{\detokenize{fix_modify:related-commands}}
\sphinxAtStartPar
{\hyperref[\detokenize{fix::doc}]{\sphinxcrossref{\DUrole{doc}{fix}}}}, {\hyperref[\detokenize{compute_temp::doc}]{\sphinxcrossref{\DUrole{doc}{compute temp}}}},
{\hyperref[\detokenize{compute_pressure::doc}]{\sphinxcrossref{\DUrole{doc}{compute pressure}}}}, {\hyperref[\detokenize{thermo_style::doc}]{\sphinxcrossref{\DUrole{doc}{thermo\_style}}}}


\subsection{Default}
\label{\detokenize{fix_modify:default}}
\sphinxAtStartPar
The option defaults are temp = ID defined by fix, press = ID defined
by fix, energy = no, virial = different for each fix style, respa = 0,
bodyforce = late.

\sphinxstepscope

\index{fitpod@\spxentry{fitpod}}\ignorespaces 

\section{fitpod command}
\label{\detokenize{fitpod_command:fitpod-command}}\label{\detokenize{fitpod_command:index-0}}\label{\detokenize{fitpod_command::doc}}

\subsection{Syntax}
\label{\detokenize{fitpod_command:syntax}}
\begin{sphinxVerbatim}[commandchars=\\\{\}]
\PYG{k}{fitpod}\PYG{+w}{ }\PYG{n}{Ta\PYGZus{}param.pod}\PYG{+w}{ }\PYG{n}{Ta\PYGZus{}data.pod}\PYG{+w}{ }\PYG{n}{Ta\PYGZus{}coefficients.pod}
\end{sphinxVerbatim}
\begin{itemize}
\item {} 
\sphinxAtStartPar
fitpod = style name of this command

\item {} 
\sphinxAtStartPar
Ta\_param.pod = an input file that describes proper orthogonal descriptors (PODs)

\item {} 
\sphinxAtStartPar
Ta\_data.pod = an input file that specifies DFT data used to fit a POD potential

\item {} 
\sphinxAtStartPar
Ta\_coefficients.pod (optional) = an input file that specifies trainable coefficients of a POD potential

\end{itemize}


\subsection{Examples}
\label{\detokenize{fitpod_command:examples}}
\begin{sphinxVerbatim}[commandchars=\\\{\}]
\PYG{k}{fitpod}\PYG{+w}{ }\PYG{n}{Ta\PYGZus{}param.pod}\PYG{+w}{ }\PYG{n}{Ta\PYGZus{}data.pod}
\PYG{k}{fitpod}\PYG{+w}{ }\PYG{n}{Ta\PYGZus{}param.pod}\PYG{+w}{ }\PYG{n}{Ta\PYGZus{}data.pod}\PYG{+w}{ }\PYG{n}{Ta\PYGZus{}coefficients.pod}
\end{sphinxVerbatim}


\subsection{Description}
\label{\detokenize{fitpod_command:description}}
\sphinxAtStartPar
\DUrole{versionmodified,added}{Added in version 22Dec2022.}

\sphinxAtStartPar
Fit a machine\sphinxhyphen{}learning interatomic potential (ML\sphinxhyphen{}IAP) based on proper
orthogonal descriptors (POD); please see {\hyperref[\detokenize{fitpod_command:nguyen20222a}]{\sphinxcrossref{\DUrole{std,std-ref}{(Nguyen and Rohskopf)}}}}, {\hyperref[\detokenize{fitpod_command:nguyen20232a}]{\sphinxcrossref{\DUrole{std,std-ref}{(Nguyen2023)}}}}, {\hyperref[\detokenize{fitpod_command:nguyen20242a}]{\sphinxcrossref{\DUrole{std,std-ref}{(Nguyen2024)}}}}, and {\hyperref[\detokenize{fitpod_command:nguyen20243a}]{\sphinxcrossref{\DUrole{std,std-ref}{(Nguyen and Sema)}}}} for details.
The fitted POD potential can be used to run MD simulations via
{\hyperref[\detokenize{pair_pod::doc}]{\sphinxcrossref{\DUrole{doc}{pair\_style pod}}}}.

\sphinxAtStartPar
Two input files are required for this command. The first input file
describes a POD potential parameter settings, while the second input
file specifies the DFT data used for the fitting procedure. All keywords
except \sphinxstyleemphasis{species} have default values. If a keyword is not set in the
input file, its default value is used. The table below has one\sphinxhyphen{}line
descriptions of all the keywords that can be used in the first input
file (i.e. \sphinxcode{\sphinxupquote{Ta\_param.pod}})


\begin{savenotes}\sphinxattablestart
\sphinxthistablewithglobalstyle
\centering
\begin{tabular}[t]{\X{40}{100}\X{9}{100}\X{10}{100}\X{41}{100}}
\sphinxtoprule
\sphinxstyletheadfamily 
\sphinxAtStartPar
Keyword
&\sphinxstyletheadfamily 
\sphinxAtStartPar
Default
&\sphinxstyletheadfamily 
\sphinxAtStartPar
Type
&\sphinxstyletheadfamily 
\sphinxAtStartPar
Description
\\
\sphinxmidrule
\sphinxtableatstartofbodyhook
\sphinxAtStartPar
species
&
\sphinxAtStartPar
(none)
&
\sphinxAtStartPar
STRING
&
\sphinxAtStartPar
Chemical symbols for all elements in the system and have to match XYZ training files.
\\
\sphinxhline
\sphinxAtStartPar
pbc
&
\sphinxAtStartPar
1 1 1
&
\sphinxAtStartPar
INT
&
\sphinxAtStartPar
three integer constants specify boundary conditions
\\
\sphinxhline
\sphinxAtStartPar
rin
&
\sphinxAtStartPar
0.5
&
\sphinxAtStartPar
REAL
&
\sphinxAtStartPar
a real number specifies the inner cut\sphinxhyphen{}off radius
\\
\sphinxhline
\sphinxAtStartPar
rcut
&
\sphinxAtStartPar
5.0
&
\sphinxAtStartPar
REAL
&
\sphinxAtStartPar
a real number specifies the outer cut\sphinxhyphen{}off radius
\\
\sphinxhline
\sphinxAtStartPar
bessel\_polynomial\_degree
&
\sphinxAtStartPar
4
&
\sphinxAtStartPar
INT
&
\sphinxAtStartPar
the maximum degree of Bessel polynomials
\\
\sphinxhline
\sphinxAtStartPar
inverse\_polynomial\_degree
&
\sphinxAtStartPar
8
&
\sphinxAtStartPar
INT
&
\sphinxAtStartPar
the maximum degree of inverse radial basis functions
\\
\sphinxhline
\sphinxAtStartPar
number\_of\_environment\_clusters
&
\sphinxAtStartPar
1
&
\sphinxAtStartPar
INT
&
\sphinxAtStartPar
the number of clusters for environment\sphinxhyphen{}adaptive potentials
\\
\sphinxhline
\sphinxAtStartPar
number\_of\_principal\_components
&
\sphinxAtStartPar
2
&
\sphinxAtStartPar
INT
&
\sphinxAtStartPar
the number of principal components for dimensionality reduction
\\
\sphinxhline
\sphinxAtStartPar
onebody
&
\sphinxAtStartPar
1
&
\sphinxAtStartPar
BOOL
&
\sphinxAtStartPar
turns on/off one\sphinxhyphen{}body potential
\\
\sphinxhline
\sphinxAtStartPar
twobody\_number\_radial\_basis\_functions
&
\sphinxAtStartPar
8
&
\sphinxAtStartPar
INT
&
\sphinxAtStartPar
number of radial basis functions for two\sphinxhyphen{}body potential
\\
\sphinxhline
\sphinxAtStartPar
threebody\_number\_radial\_basis\_functions
&
\sphinxAtStartPar
6
&
\sphinxAtStartPar
INT
&
\sphinxAtStartPar
number of radial basis functions for three\sphinxhyphen{}body potential
\\
\sphinxhline
\sphinxAtStartPar
threebody\_angular\_degree
&
\sphinxAtStartPar
5
&
\sphinxAtStartPar
INT
&
\sphinxAtStartPar
angular degree for three\sphinxhyphen{}body potential
\\
\sphinxhline
\sphinxAtStartPar
fourbody\_number\_radial\_basis\_functions
&
\sphinxAtStartPar
4
&
\sphinxAtStartPar
INT
&
\sphinxAtStartPar
number of radial basis functions for four\sphinxhyphen{}body potential
\\
\sphinxhline
\sphinxAtStartPar
fourbody\_angular\_degree
&
\sphinxAtStartPar
3
&
\sphinxAtStartPar
INT
&
\sphinxAtStartPar
angular degree for four\sphinxhyphen{}body potential
\\
\sphinxhline
\sphinxAtStartPar
fivebody\_number\_radial\_basis\_functions
&
\sphinxAtStartPar
0
&
\sphinxAtStartPar
INT
&
\sphinxAtStartPar
number of radial basis functions for five\sphinxhyphen{}body potential
\\
\sphinxhline
\sphinxAtStartPar
fivebody\_angular\_degree
&
\sphinxAtStartPar
0
&
\sphinxAtStartPar
INT
&
\sphinxAtStartPar
angular degree for five\sphinxhyphen{}body potential
\\
\sphinxhline
\sphinxAtStartPar
sixbody\_number\_radial\_basis\_functions
&
\sphinxAtStartPar
0
&
\sphinxAtStartPar
INT
&
\sphinxAtStartPar
number of radial basis functions for six\sphinxhyphen{}body potential
\\
\sphinxhline
\sphinxAtStartPar
sixbody\_angular\_degree
&
\sphinxAtStartPar
0
&
\sphinxAtStartPar
INT
&
\sphinxAtStartPar
angular degree for six\sphinxhyphen{}body potential
\\
\sphinxhline
\sphinxAtStartPar
sevenbody\_number\_radial\_basis\_functions
&
\sphinxAtStartPar
0
&
\sphinxAtStartPar
INT
&
\sphinxAtStartPar
number of radial basis functions for seven\sphinxhyphen{}body potential
\\
\sphinxhline
\sphinxAtStartPar
sevenbody\_angular\_degree
&
\sphinxAtStartPar
0
&
\sphinxAtStartPar
INT
&
\sphinxAtStartPar
angular degree for seven\sphinxhyphen{}body potential
\\
\sphinxbottomrule
\end{tabular}
\sphinxtableafterendhook\par
\sphinxattableend\end{savenotes}

\sphinxAtStartPar
Note that both the number of radial basis functions and angular degree
must decrease as the body order increases. The next table describes all
keywords that can be used in the second input file (i.e. \sphinxcode{\sphinxupquote{Ta\_data.pod}}
in the example above):


\begin{savenotes}\sphinxattablestart
\sphinxthistablewithglobalstyle
\centering
\begin{tabular}[t]{\X{38}{100}\X{9}{100}\X{10}{100}\X{43}{100}}
\sphinxtoprule
\sphinxstyletheadfamily 
\sphinxAtStartPar
Keyword
&\sphinxstyletheadfamily 
\sphinxAtStartPar
Default
&\sphinxstyletheadfamily 
\sphinxAtStartPar
Type
&\sphinxstyletheadfamily 
\sphinxAtStartPar
Description
\\
\sphinxmidrule
\sphinxtableatstartofbodyhook
\sphinxAtStartPar
file\_format
&
\sphinxAtStartPar
extxyz
&
\sphinxAtStartPar
STRING
&
\sphinxAtStartPar
only the extended xyz format (extxyz) is currently supported
\\
\sphinxhline
\sphinxAtStartPar
file\_extension
&
\sphinxAtStartPar
xyz
&
\sphinxAtStartPar
STRING
&
\sphinxAtStartPar
extension of the data files
\\
\sphinxhline
\sphinxAtStartPar
path\_to\_training\_data\_set
&
\sphinxAtStartPar
(none)
&
\sphinxAtStartPar
STRING
&
\sphinxAtStartPar
specifies the path to training data files in double quotes
\\
\sphinxhline
\sphinxAtStartPar
path\_to\_test\_data\_set
&
\sphinxAtStartPar
“”
&
\sphinxAtStartPar
STRING
&
\sphinxAtStartPar
specifies the path to test data files in double quotes
\\
\sphinxhline
\sphinxAtStartPar
path\_to\_environment\_configuration\_set
&
\sphinxAtStartPar
“”
&
\sphinxAtStartPar
STRING
&
\sphinxAtStartPar
specifies the path to environment configuration files in double quotes
\\
\sphinxhline
\sphinxAtStartPar
fraction\_training\_data\_set
&
\sphinxAtStartPar
1.0
&
\sphinxAtStartPar
REAL
&
\sphinxAtStartPar
a real number (\textless{}= 1.0) specifies the fraction of the training set used to fit POD
\\
\sphinxhline
\sphinxAtStartPar
randomize\_training\_data\_set
&
\sphinxAtStartPar
0
&
\sphinxAtStartPar
BOOL
&
\sphinxAtStartPar
turns on/off randomization of the training set
\\
\sphinxhline
\sphinxAtStartPar
fraction\_test\_data\_set
&
\sphinxAtStartPar
1.0
&
\sphinxAtStartPar
REAL
&
\sphinxAtStartPar
a real number (\textless{}= 1.0) specifies the fraction of the test set used to validate POD
\\
\sphinxhline
\sphinxAtStartPar
randomize\_test\_data\_set
&
\sphinxAtStartPar
0
&
\sphinxAtStartPar
BOOL
&
\sphinxAtStartPar
turns on/off randomization of the test set
\\
\sphinxhline
\sphinxAtStartPar
fitting\_weight\_energy
&
\sphinxAtStartPar
100.0
&
\sphinxAtStartPar
REAL
&
\sphinxAtStartPar
a real constant specifies the weight for energy in the least\sphinxhyphen{}squares fit
\\
\sphinxhline
\sphinxAtStartPar
fitting\_weight\_force
&
\sphinxAtStartPar
1.0
&
\sphinxAtStartPar
REAL
&
\sphinxAtStartPar
a real constant specifies the weight for force in the least\sphinxhyphen{}squares fit
\\
\sphinxhline
\sphinxAtStartPar
fitting\_regularization\_parameter
&
\sphinxAtStartPar
1.0e\sphinxhyphen{}10
&
\sphinxAtStartPar
REAL
&
\sphinxAtStartPar
a real constant specifies the regularization parameter in the least\sphinxhyphen{}squares fit
\\
\sphinxhline
\sphinxAtStartPar
error\_analysis\_for\_training\_data\_set
&
\sphinxAtStartPar
0
&
\sphinxAtStartPar
BOOL
&
\sphinxAtStartPar
turns on/off error analysis for the training data set
\\
\sphinxhline
\sphinxAtStartPar
error\_analysis\_for\_test\_data\_set
&
\sphinxAtStartPar
0
&
\sphinxAtStartPar
BOOL
&
\sphinxAtStartPar
turns on/off error analysis for the test data set
\\
\sphinxhline
\sphinxAtStartPar
basename\_for\_output\_files
&
\sphinxAtStartPar
pod
&
\sphinxAtStartPar
STRING
&
\sphinxAtStartPar
a basename string added to the output files
\\
\sphinxhline
\sphinxAtStartPar
precision\_for\_pod\_coefficients
&
\sphinxAtStartPar
8
&
\sphinxAtStartPar
INT
&
\sphinxAtStartPar
number of digits after the decimal points for numbers in the coefficient file
\\
\sphinxhline
\sphinxAtStartPar
group\_weights
&
\sphinxAtStartPar
global
&
\sphinxAtStartPar
STRING
&
\sphinxAtStartPar
\sphinxcode{\sphinxupquote{table}} uses group weights defined for each group named by filename
\\
\sphinxbottomrule
\end{tabular}
\sphinxtableafterendhook\par
\sphinxattableend\end{savenotes}

\sphinxAtStartPar
All keywords except \sphinxstyleemphasis{path\_to\_training\_data\_set} have default values. If
a keyword is not set in the input file, its default value is used.  After
successful training, a number of output files are produced, if enabled:
\begin{itemize}
\item {} 
\sphinxAtStartPar
\sphinxcode{\sphinxupquote{\textless{}basename\textgreater{}\_training\_errors.pod}}  reports the errors in energy and forces for the training data set

\item {} 
\sphinxAtStartPar
\sphinxcode{\sphinxupquote{\textless{}basename\textgreater{}\_training\_analysis.pod}} reports detailed errors for all training configurations

\item {} 
\sphinxAtStartPar
\sphinxcode{\sphinxupquote{\textless{}basename\textgreater{}\_test\_errors.pod}} reports errors for the test data set

\item {} 
\sphinxAtStartPar
\sphinxcode{\sphinxupquote{\textless{}basename\textgreater{}\_test\_analysis.pod}} reports detailed errors for all test configurations

\item {} 
\sphinxAtStartPar
\sphinxcode{\sphinxupquote{\textless{}basename\textgreater{}\_coefficients.pod}} contains the coefficients of the POD potential

\end{itemize}

\sphinxAtStartPar
After training the POD potential, \sphinxcode{\sphinxupquote{Ta\_param.pod}} and
\sphinxcode{\sphinxupquote{\textless{}basename\textgreater{}\_coefficients.pod}} are the two files needed to use the POD
potential in LAMMPS.  See {\hyperref[\detokenize{pair_pod::doc}]{\sphinxcrossref{\DUrole{doc}{pair\_style pod}}}} for using the
POD potential. Examples about training and using POD potentials are
found in the directory lammps/examples/PACKAGES/pod and the Github repo
\sphinxurl{https://github.com/cesmix-mit/pod-examples}.


\subsubsection{Loss Function Group Weights}
\label{\detokenize{fitpod_command:loss-function-group-weights}}
\sphinxAtStartPar
The \sphinxstyleemphasis{group\_weights} keyword in the \sphinxcode{\sphinxupquote{data.pod}} file is responsible for
weighting certain groups of configurations in the loss function. For
example:

\begin{sphinxVerbatim}[commandchars=\\\{\}]
\PYG{n}{group\PYGZus{}weights}\PYG{+w}{ }\PYG{n}{table}
\PYG{n}{Displaced\PYGZus{}A15}\PYG{+w}{ }\PYG{l+m}{100.0}\PYG{+w}{ }\PYG{l+m}{1.0}
\PYG{n}{Displaced\PYGZus{}BCC}\PYG{+w}{ }\PYG{l+m}{100.0}\PYG{+w}{ }\PYG{l+m}{1.0}
\PYG{n}{Displaced\PYGZus{}FCC}\PYG{+w}{ }\PYG{l+m}{100.0}\PYG{+w}{ }\PYG{l+m}{1.0}
\PYG{n}{Elastic\PYGZus{}BCC}\PYG{+w}{   }\PYG{l+m}{100.0}\PYG{+w}{ }\PYG{l+m}{1.0}
\PYG{n}{Elastic\PYGZus{}FCC}\PYG{+w}{   }\PYG{l+m}{100.0}\PYG{+w}{ }\PYG{l+m}{1.0}
\PYG{n}{GSF\PYGZus{}110}\PYG{+w}{       }\PYG{l+m}{100.0}\PYG{+w}{ }\PYG{l+m}{1.0}
\PYG{n}{GSF\PYGZus{}112}\PYG{+w}{       }\PYG{l+m}{100.0}\PYG{+w}{ }\PYG{l+m}{1.0}
\PYG{n}{Liquid}\PYG{+w}{        }\PYG{l+m}{100.0}\PYG{+w}{ }\PYG{l+m}{1.0}
\PYG{n}{Surface}\PYG{+w}{       }\PYG{l+m}{100.0}\PYG{+w}{ }\PYG{l+m}{1.0}
\PYG{n}{Volume\PYGZus{}A15}\PYG{+w}{    }\PYG{l+m}{100.0}\PYG{+w}{ }\PYG{l+m}{1.0}
\PYG{n}{Volume\PYGZus{}BCC}\PYG{+w}{    }\PYG{l+m}{100.0}\PYG{+w}{ }\PYG{l+m}{1.0}
\PYG{n}{Volume\PYGZus{}FCC}\PYG{+w}{    }\PYG{l+m}{100.0}\PYG{+w}{ }\PYG{l+m}{1.0}
\end{sphinxVerbatim}

\sphinxAtStartPar
This will apply an energy weight of \sphinxcode{\sphinxupquote{100.0}} and a force weight of
\sphinxcode{\sphinxupquote{1.0}} for all groups in the \sphinxcode{\sphinxupquote{Ta}} example. The groups are named by
their respective filename. If certain groups are left out of this table,
then the globally defined weights from the \sphinxcode{\sphinxupquote{fitting\_weight\_energy}} and
\sphinxcode{\sphinxupquote{fitting\_weight\_force}} keywords will be used.


\subsection{POD Potential}
\label{\detokenize{fitpod_command:pod-potential}}
\sphinxAtStartPar
We consider a multi\sphinxhyphen{}element system of \sphinxstyleemphasis{N} atoms with \(N_{\rm e}\)
unique elements.  We denote by \(\boldsymbol r_n\) and \(Z_n\)
position vector and type of an atom \sphinxstyleemphasis{n} in the system,
respectively. Note that we have \(Z_n \in \{1, \ldots, N_{\rm e}
\}\), \(\boldsymbol R = (\boldsymbol r_1, \boldsymbol r_2, \ldots,
\boldsymbol r_N) \in \mathbb{R}^{3N}\), and \(\boldsymbol Z = (Z_1,
Z_2, \ldots, Z_N) \in \mathbb{N}^{N}\). The total energy of the
POD potential is expressed as \(E(\boldsymbol R, \boldsymbol Z) =
\sum_{i=1}^N E_i(\boldsymbol R_i, \boldsymbol Z_i)\), where
\begin{equation*}
\begin{split}E_i(\boldsymbol R_i, \boldsymbol Z_i) \ = \ \sum_{m=1}^M c_m \mathcal{D}_{im}(\boldsymbol R_i, \boldsymbol Z_i)\end{split}
\end{equation*}
\sphinxAtStartPar
Here \(c_m\) are trainable coefficients and
\(\mathcal{D}_{im}(\boldsymbol R_i, \boldsymbol Z_i)\) are per\sphinxhyphen{}atom
POD descriptors. Summing the per\sphinxhyphen{}atom descriptors over \(i\) yields
the global descriptors \(d_m(\boldsymbol R, \boldsymbol Z) =
\sum_{i=1}^N \mathcal{D}_{im}(\boldsymbol R_i, \boldsymbol Z_i)\).  It
thus follows that \(E(\boldsymbol R, \boldsymbol Z) = \sum_{m=1}^M
c_m d_m(\boldsymbol R, \boldsymbol Z)\).

\sphinxAtStartPar
The per\sphinxhyphen{}atom POD descriptors include one, two, three, four, five, six,
and seven\sphinxhyphen{}body descriptors, which can be specified in the first input
file. Furthermore, the per\sphinxhyphen{}atom POD descriptors also depend on the
number of environment clusters specified in the first input file.
Please see {\hyperref[\detokenize{fitpod_command:nguyen20242a}]{\sphinxcrossref{\DUrole{std,std-ref}{(Nguyen2024)}}}} and {\hyperref[\detokenize{fitpod_command:nguyen20243a}]{\sphinxcrossref{\DUrole{std,std-ref}{(Nguyen and Sema)}}}} for the detailed description of the per\sphinxhyphen{}atom POD
descriptors.


\subsection{Training}
\label{\detokenize{fitpod_command:training}}
\sphinxAtStartPar
A POD potential is trained using the least\sphinxhyphen{}squares regression against
density functional theory (DFT) data.  Let \(J\) be the number of
training configurations, with \(N_j\) being the number of atoms in
the j\sphinxhyphen{}th configuration. The training configurations are extracted from
the extended XYZ files located in a directory (i.e.,
path\_to\_training\_data\_set in the second input file).  Let
\(\{E^{\star}_j\}_{j=1}^{J}\) and \(\{\boldsymbol
F^{\star}_j\}_{j=1}^{J}\) be the DFT energies and forces for \(J\)
configurations. Next, we calculate the global descriptors and their
derivatives for all training configurations. Let \(d_{jm}, 1 \le m
\le M\), be the global descriptors associated with the j\sphinxhyphen{}th
configuration, where \(M\) is the number of global descriptors. We
then form a matrix \(\boldsymbol A \in \mathbb{R}^{J \times M}\)
with entries \(A_{jm} = d_{jm}/ N_j\) for \(j=1,\ldots,J\) and
\(m=1,\ldots,M\).  Moreover, we form a matrix \(\boldsymbol B
\in \mathbb{R}^{\mathcal{N} \times M}\) by stacking the derivatives of
the global descriptors for all training configurations from top to
bottom, where \(\mathcal{N} = 3\sum_{j=1}^{J} N_j\).

\sphinxAtStartPar
The coefficient vector \(\boldsymbol c\) of the POD potential is
found by solving the following least\sphinxhyphen{}squares problem
\begin{equation*}
\begin{split}{\min}_{\boldsymbol c \in \mathbb{R}^{M}} \ w_E \|\boldsymbol A \boldsymbol c - \bar{\boldsymbol E}^{\star} \|^2 + w_F \|\boldsymbol B \boldsymbol c + \boldsymbol F^{\star} \|^2 + w_R \|\boldsymbol c \|^2,\end{split}
\end{equation*}
\sphinxAtStartPar
where \(w_E\) and \(w_F\) are weights for the energy
(\sphinxstyleemphasis{fitting\_weight\_energy}) and force (\sphinxstyleemphasis{fitting\_weight\_force}),
respectively; and \(w_R\) is the regularization parameter
(\sphinxstyleemphasis{fitting\_regularization\_parameter}).  Here \(\bar{\boldsymbol
E}^{\star} \in \mathbb{R}^{J}\) is a vector of with entries
\(\bar{E}^{\star}_j = E^{\star}_j/N_j\) and \(\boldsymbol
F^{\star}\) is a vector of \(\mathcal{N}\) entries obtained by
stacking \(\{\boldsymbol F^{\star}_j\}_{j=1}^{J}\) from top to
bottom.


\subsection{Validation}
\label{\detokenize{fitpod_command:validation}}
\sphinxAtStartPar
POD potential can be validated on a test dataset in a directory
specified by setting path\_to\_test\_data\_set in the second input file.  It
is possible to validate the POD potential after the training is
complete.  This is done by providing the coefficient file as an input to
{\hyperref[\detokenize{fitpod_command::doc}]{\sphinxcrossref{\DUrole{doc}{fitpod}}}}, for example,

\begin{sphinxVerbatim}[commandchars=\\\{\}]
\PYG{k}{fitpod}\PYG{+w}{ }\PYG{n}{Ta\PYGZus{}param.pod}\PYG{+w}{ }\PYG{n}{Ta\PYGZus{}data.pod}\PYG{+w}{ }\PYG{n}{Ta\PYGZus{}coefficients.pod}
\end{sphinxVerbatim}


\subsection{Restrictions}
\label{\detokenize{fitpod_command:restrictions}}
\sphinxAtStartPar
This command is part of the ML\sphinxhyphen{}POD package.  It is only enabled if
LAMMPS was built with that package. See the {\hyperref[\detokenize{Build_package::doc}]{\sphinxcrossref{\DUrole{doc}{Build package}}}} page for more info.


\subsection{Related commands}
\label{\detokenize{fitpod_command:related-commands}}
\sphinxAtStartPar
{\hyperref[\detokenize{pair_pod::doc}]{\sphinxcrossref{\DUrole{doc}{pair\_style pod}}}},
{\hyperref[\detokenize{compute_pod_atom::doc}]{\sphinxcrossref{\DUrole{doc}{compute pod/atom}}}},
{\hyperref[\detokenize{compute_pod_atom::doc}]{\sphinxcrossref{\DUrole{doc}{compute podd/atom}}}},
{\hyperref[\detokenize{compute_pod_atom::doc}]{\sphinxcrossref{\DUrole{doc}{compute pod/local}}}},
{\hyperref[\detokenize{compute_pod_atom::doc}]{\sphinxcrossref{\DUrole{doc}{compute pod/global}}}}


\subsection{Default}
\label{\detokenize{fitpod_command:default}}
\sphinxAtStartPar
The keyword defaults are also given in the description of the input files.


\bigskip\hrule\bigskip

\phantomsection\label{\detokenize{fitpod_command:nguyen20222a}}
\sphinxAtStartPar
\sphinxstylestrong{(Nguyen and Rohskopf)} Nguyen and Rohskopf,  Journal of Computational Physics, 480, 112030, (2023).

\phantomsection\label{\detokenize{fitpod_command:nguyen20232a}}
\sphinxAtStartPar
\sphinxstylestrong{(Nguyen2023)} Nguyen, Physical Review B, 107(14), 144103, (2023).

\phantomsection\label{\detokenize{fitpod_command:nguyen20242a}}
\sphinxAtStartPar
\sphinxstylestrong{(Nguyen2024)} Nguyen, Journal of Computational Physics, 113102, (2024).

\phantomsection\label{\detokenize{fitpod_command:nguyen20243a}}
\sphinxAtStartPar
\sphinxstylestrong{(Nguyen and Sema)} Nguyen and Sema, \sphinxurl{https://arxiv.org/abs/2405.00306}, (2024).

\sphinxstepscope

\index{geturl@\spxentry{geturl}}\ignorespaces 

\section{geturl command}
\label{\detokenize{geturl:geturl-command}}\label{\detokenize{geturl:index-0}}\label{\detokenize{geturl::doc}}

\subsection{Syntax}
\label{\detokenize{geturl:syntax}}
\begin{sphinxVerbatim}[commandchars=\\\{\}]
\PYG{k}{geturl}\PYG{+w}{ }\PYG{n}{url}\PYG{+w}{ }\PYG{n}{keyword}\PYG{+w}{ }\PYG{n}{args}\PYG{+w}{ }\PYG{n}{...}
\end{sphinxVerbatim}
\begin{itemize}
\item {} 
\sphinxAtStartPar
url = URL of the file to download

\item {} 
\sphinxAtStartPar
zero or more keyword argument pairs may be provided

\item {} 
\sphinxAtStartPar
keyword = \sphinxstyleemphasis{output} or \sphinxstyleemphasis{verify} or \sphinxstyleemphasis{overwrite} or \sphinxstyleemphasis{verbose}
\begin{sphinxalltt}
\sphinxstyleemphasis{output} filename = write to \sphinxstyleemphasis{filename} instead of inferring the name from the URL
\sphinxstyleemphasis{verify} yes/no = verify SSL certificate and hostname if \sphinxstyleemphasis{yes}, do not if \sphinxstyleemphasis{no}
\sphinxstyleemphasis{overwrite} yes/no = if \sphinxstyleemphasis{yes} overwrite the output file in case it exists, do not if \sphinxstyleemphasis{no}
\sphinxstyleemphasis{verbose} yes/no = if \sphinxstyleemphasis{yes} write verbose debug output from libcurl to screen, do not if \sphinxstyleemphasis{no}
\end{sphinxalltt}

\end{itemize}


\subsection{Examples}
\label{\detokenize{geturl:examples}}
\begin{sphinxVerbatim}[commandchars=\\\{\}]
\PYG{k}{geturl}\PYG{+w}{ }\PYG{n}{https}://www.ctcms.nist.gov/potentials/Download/1990\PYGZhy{}\PYGZhy{}Ackland\PYGZhy{}G\PYGZhy{}J\PYGZhy{}Vitek\PYGZhy{}V\PYGZhy{}\PYGZhy{}Cu/2/Cu2.eam.fs
\PYG{k}{geturl}\PYG{+w}{ }\PYG{n}{https}://github.com/lammps/lammps/blob/develop/bench/in.lj\PYG{+w}{ }\PYG{n}{output}\PYG{+w}{ }\PYG{n}{in.bench}\PYG{o}{\PYGZhy{}}\PYG{n}{lj}
\end{sphinxVerbatim}


\subsection{Description}
\label{\detokenize{geturl:description}}
\sphinxAtStartPar
\DUrole{versionmodified,added}{Added in version 29Aug2024.}

\sphinxAtStartPar
Download a file from an URL to the local disk. This is implemented with
the \sphinxhref{https:://curl.se/libcurl/}{libcurl library} which supports a
large variety of protocols including “http”, “https”, “ftp”, “scp”,
“sftp”, “file”.  The transfer will only be performed on MPI rank 0.

\sphinxAtStartPar
The \sphinxstyleemphasis{output} keyword can be used to set the filename. By default, the last part
of the URL is used.

\sphinxAtStartPar
The \sphinxstyleemphasis{verify} keyword determines whether \sphinxcode{\sphinxupquote{libcurl}} will validate the
SSL certificate and hostname for encrypted connections.  Turning this
off may be required when using a proxy or connecting to a server with a
self\sphinxhyphen{}signed SSL certificate.

\sphinxAtStartPar
The \sphinxstyleemphasis{overwrite} keyword determines whether a file should be overwritten if it
already exists.  If the argument is \sphinxstyleemphasis{no}, then the download will be skipped
if the file exists.

\sphinxAtStartPar
The \sphinxstyleemphasis{verbose} keyword determines whether a detailed protocol of the steps
performed by libcurl is written to the screen.  Using the argument \sphinxstyleemphasis{yes}
can be used to debug connection issues when the \sphinxstyleemphasis{geturl} command does not
behave as expected.  If the argument is \sphinxstyleemphasis{no}, geturl will operate silently
and only report the error status number provided by libcurl, in case of a
failure.


\bigskip\hrule\bigskip



\subsection{Restrictions}
\label{\detokenize{geturl:restrictions}}
\sphinxAtStartPar
This command is part of the EXTRA\sphinxhyphen{}COMMAND package.  It is only enabled
if LAMMPS was built with that package.  See the {\hyperref[\detokenize{Build_package::doc}]{\sphinxcrossref{\DUrole{doc}{Build package}}}} page for more info.  It also requires that LAMMPS was
built with support for \sphinxhref{https://curl.se/libcurl/}{the libcurl library}.  See the page about {\hyperref[\detokenize{Build_settings:libcurl}]{\sphinxcrossref{\DUrole{std,std-ref}{Compiling LAMMPS
with libcurl support}}}} for further info.  If support for
libcurl is not included, using \sphinxstyleemphasis{geturl} will trigger an error.


\subsection{Related commands}
\label{\detokenize{geturl:related-commands}}
\sphinxAtStartPar
{\hyperref[\detokenize{shell::doc}]{\sphinxcrossref{\DUrole{doc}{shell}}}}


\subsection{Default}
\label{\detokenize{geturl:default}}
\sphinxAtStartPar
\sphinxstyleemphasis{verify} = yes, \sphinxstyleemphasis{overwrite} = yes

\sphinxstepscope

\index{group@\spxentry{group}}\ignorespaces 

\section{group command}
\label{\detokenize{group:group-command}}\label{\detokenize{group:index-0}}\label{\detokenize{group::doc}}

\subsection{Syntax}
\label{\detokenize{group:syntax}}
\begin{sphinxVerbatim}[commandchars=\\\{\}]
\PYG{k}{group }\PYG{n+nv+nvIdentifier}{ID}\PYG{+w}{ }\PYG{n}{style}\PYG{+w}{ }\PYG{n}{args}
\end{sphinxVerbatim}
\begin{itemize}
\item {} 
\sphinxAtStartPar
ID = user\sphinxhyphen{}defined name of the group

\item {} 
\sphinxAtStartPar
style = \sphinxstyleemphasis{delete} or \sphinxstyleemphasis{clear} or \sphinxstyleemphasis{empty} or \sphinxstyleemphasis{region} or \sphinxstyleemphasis{type} or \sphinxstyleemphasis{id} or \sphinxstyleemphasis{molecule} or \sphinxstyleemphasis{variable} or \sphinxstyleemphasis{include} or \sphinxstyleemphasis{subtract} or \sphinxstyleemphasis{union} or \sphinxstyleemphasis{intersect} or \sphinxstyleemphasis{dynamic} or \sphinxstyleemphasis{static}
\begin{sphinxalltt}
\sphinxstyleemphasis{delete} = no args
\sphinxstyleemphasis{clear} = no args
\sphinxstyleemphasis{empty} = no args
\sphinxstyleemphasis{region} args = region\sphinxhyphen{}ID
\sphinxstyleemphasis{type} or \sphinxstyleemphasis{id} or \sphinxstyleemphasis{molecule}
  args = list of one or more atom types (1\sphinxhyphen{}Ntypes or type label), atom IDs, or molecule IDs
    any numeric entry in list can be a sequence formatted as A:B or A:B:C where
    A = starting index, B = ending index,
    C = increment between indices, 1 if not specified
  args = logical value
    logical = "\textless{}" or "\textless{}=" or "\textgreater{}" or "\textgreater{}=" or "==" or "!="
    value = an atom type (1\sphinxhyphen{}Ntypes or type label) or atom ID or molecule ID (depending on \sphinxstyleemphasis{style})
  args = logical value1 value2
    logical = "\textless{}\textgreater{}"
    value1,value2 = atom types or atom IDs or molecule IDs (depending on \sphinxstyleemphasis{style})
\sphinxstyleemphasis{variable} args = variable\sphinxhyphen{}name
\sphinxstyleemphasis{include} args = molecule
  molecule = add atoms to group with same molecule ID as atoms already in group
\sphinxstyleemphasis{subtract} args = two or more group IDs
\sphinxstyleemphasis{union} args = one or more group IDs
\sphinxstyleemphasis{intersect} args = two or more group IDs
\sphinxstyleemphasis{dynamic} args = parent\sphinxhyphen{}ID keyword value ...
  one or more keyword/value pairs may be appended
  keyword = \sphinxstyleemphasis{region} or \sphinxstyleemphasis{var} or \sphinxstyleemphasis{property} or \sphinxstyleemphasis{every}
    \sphinxstyleemphasis{region} value = region\sphinxhyphen{}ID
    \sphinxstyleemphasis{var} value = name of variable
    \sphinxstyleemphasis{property} value = name of custom integer or floating point vector
    \sphinxstyleemphasis{every} value = N = update group every this many timesteps
\sphinxstyleemphasis{static} = no args
\end{sphinxalltt}

\end{itemize}


\subsection{Examples}
\label{\detokenize{group:examples}}
\begin{sphinxVerbatim}[commandchars=\\\{\}]
\PYG{k}{group }\PYG{n+nv+nvIdentifier}{edge}\PYG{+w}{ }\PYG{k}{region }\PYG{n+nv+nvIdentifier}{regstrip}
\PYG{k}{group }\PYG{n+nv+nvIdentifier}{water}\PYG{+w}{ }\PYG{n}{type}\PYG{+w}{ }\PYG{l+m}{3}\PYG{+w}{ }\PYG{l+m}{4}
\PYG{k}{group }\PYG{n+nv+nvIdentifier}{water}\PYG{+w}{ }\PYG{n}{type}\PYG{+w}{ }\PYG{n}{OW}\PYG{+w}{ }\PYG{n}{HT}
\PYG{k}{group }\PYG{n+nv+nvIdentifier}{sub}\PYG{+w}{ }\PYG{n}{id}\PYG{+w}{ }\PYG{l+m}{10}\PYG{+w}{ }\PYG{l+m}{25}\PYG{+w}{ }\PYG{l+m}{50}
\PYG{k}{group }\PYG{n+nv+nvIdentifier}{sub}\PYG{+w}{ }\PYG{n}{id}\PYG{+w}{ }\PYG{l+m}{10}\PYG{+w}{ }\PYG{l+m}{25}\PYG{+w}{ }\PYG{l+m}{50}\PYG{+w}{ }\PYG{l+m}{500:1000}
\PYG{k}{group }\PYG{n+nv+nvIdentifier}{sub}\PYG{+w}{ }\PYG{n}{id}\PYG{+w}{ }\PYG{l+m}{100:10000:10}
\PYG{k}{group }\PYG{n+nv+nvIdentifier}{sub}\PYG{+w}{ }\PYG{n}{id}\PYG{+w}{ }\PYG{o}{\PYGZlt{}}\PYG{o}{=}\PYG{+w}{ }\PYG{l+m}{150}
\PYG{k}{group }\PYG{n+nv+nvIdentifier}{polyA}\PYG{+w}{ }\PYG{n}{molecule}\PYG{+w}{ }\PYG{o}{\PYGZlt{}}\PYG{o}{\PYGZgt{}}\PYG{+w}{ }\PYG{l+m}{50}\PYG{+w}{ }\PYG{l+m}{250}
\PYG{k}{group }\PYG{n+nv+nvIdentifier}{hienergy}\PYG{+w}{ }\PYG{k}{variable }\PYG{n+nv+nvIdentifier}{eng}
\PYG{k}{group }\PYG{n+nv+nvIdentifier}{hienergy}\PYG{+w}{ }\PYG{n}{include}\PYG{+w}{ }\PYG{n}{molecule}
\PYG{k}{group }\PYG{n+nv+nvIdentifier}{boundary}\PYG{+w}{ }\PYG{n}{subtract}\PYG{+w}{ }\PYG{n}{all}\PYG{+w}{ }\PYG{n}{a2}\PYG{+w}{ }\PYG{n}{a3}
\PYG{k}{group }\PYG{n+nv+nvIdentifier}{boundary}\PYG{+w}{ }\PYG{n}{union}\PYG{+w}{ }\PYG{n}{lower}\PYG{+w}{ }\PYG{n}{upper}
\PYG{k}{group }\PYG{n+nv+nvIdentifier}{boundary}\PYG{+w}{ }\PYG{n}{intersect}\PYG{+w}{ }\PYG{n}{upper}\PYG{+w}{ }\PYG{n}{flow}
\PYG{k}{group }\PYG{n+nv+nvIdentifier}{boundary}\PYG{+w}{ }\PYG{n}{delete}
\PYG{k}{group }\PYG{n+nv+nvIdentifier}{mine}\PYG{+w}{ }\PYG{n}{dynamic}\PYG{+w}{ }\PYG{n}{all}\PYG{+w}{ }\PYG{k}{region }\PYG{n+nv+nvIdentifier}{myRegion}\PYG{+w}{ }\PYG{n}{every}\PYG{+w}{ }\PYG{l+m}{100}
\end{sphinxVerbatim}


\subsection{Description}
\label{\detokenize{group:description}}
\sphinxAtStartPar
Identify a collection of atoms as belonging to a group.  The group ID
can then be used in other commands such as {\hyperref[\detokenize{fix::doc}]{\sphinxcrossref{\DUrole{doc}{fix}}}},
{\hyperref[\detokenize{compute::doc}]{\sphinxcrossref{\DUrole{doc}{compute}}}}, {\hyperref[\detokenize{dump::doc}]{\sphinxcrossref{\DUrole{doc}{dump}}}}, or {\hyperref[\detokenize{velocity::doc}]{\sphinxcrossref{\DUrole{doc}{velocity}}}}
to act on those atoms together.

\sphinxAtStartPar
If the group ID already exists, the group command adds the specified
atoms to the group.

\begin{sphinxadmonition}{note}{Note:}
\sphinxAtStartPar
By default groups are static, meaning the atoms are permanently
assigned to the group.  For example, if the \sphinxstyleemphasis{region} style is used to
assign atoms to a group, the atoms will remain in the group even if
they later move out of the region.  As explained below, the \sphinxstyleemphasis{dynamic}
style can be used to make a group dynamic so that a periodic
determination is made as to which atoms are in the group.  Since many
LAMMPS commands operate on groups of atoms, you should think carefully
about whether making a group dynamic makes sense for your model.
\end{sphinxadmonition}

\sphinxAtStartPar
A group with the ID \sphinxstyleemphasis{all} is predefined.  All atoms belong to this group.
This group cannot be deleted, or made dynamic.

\sphinxAtStartPar
The \sphinxstyleemphasis{delete} style removes the named group and un\sphinxhyphen{}assigns all atoms
that were assigned to that group.  Since there is a restriction (see
below) that no more than 32 groups can be defined at any time, the
\sphinxstyleemphasis{delete} style allows you to remove groups that are no longer needed,
so that more can be specified.  You cannot delete a group if it has
been used to define a current {\hyperref[\detokenize{fix::doc}]{\sphinxcrossref{\DUrole{doc}{fix}}}} or {\hyperref[\detokenize{compute::doc}]{\sphinxcrossref{\DUrole{doc}{compute}}}}
or {\hyperref[\detokenize{dump::doc}]{\sphinxcrossref{\DUrole{doc}{dump}}}}.

\sphinxAtStartPar
The \sphinxstyleemphasis{clear} style un\sphinxhyphen{}assigns all atoms that were assigned to that
group.  This may be dangerous to do during a simulation run
(e.g., using the {\hyperref[\detokenize{run::doc}]{\sphinxcrossref{\DUrole{doc}{run every}}}} command if a fix or compute or
other operation expects the atoms in the group to remain constant), but
LAMMPS does not check for this.

\sphinxAtStartPar
The \sphinxstyleemphasis{empty} style creates an empty group, which is useful for commands
like {\hyperref[\detokenize{fix_gcmc::doc}]{\sphinxcrossref{\DUrole{doc}{fix gcmc}}}} or with complex scripts that add atoms
to a group.

\sphinxAtStartPar
The \sphinxstyleemphasis{region} style puts all atoms in the region volume into the group.
Note that this is a static one\sphinxhyphen{}time assignment.  The atoms remain
assigned (or not assigned) to the group even in they later move out of
the region volume.

\sphinxAtStartPar
The \sphinxstyleemphasis{type}, \sphinxstyleemphasis{id}, and \sphinxstyleemphasis{molecule} styles put all atoms with the
specified atom types, atom IDs, or molecule IDs into the group.  These
three styles can use arguments specified in one of two formats.

\sphinxAtStartPar
The first format is a list of values (types or IDs).  For example, the
second command in the examples above puts all atoms of type 3 or 4 into
the group named \sphinxstyleemphasis{water}.  Each numeric entry in the list can be a
colon\sphinxhyphen{}separated sequence \sphinxcode{\sphinxupquote{A:B}} or \sphinxcode{\sphinxupquote{A:B:C}}, as in two of the examples
above.  A “sequence” generates a sequence of values (types or IDs),
with an optional increment.  The first example with \sphinxcode{\sphinxupquote{500:1000}} has the
default increment of 1 and would add all atom IDs from 500 to 1000
(inclusive) to the group sub, along with 10, 25, and 50 since they also
appear in the list of values.  The second example with \sphinxcode{\sphinxupquote{100:10000:10}}
uses an increment of 10 and would thus would add atoms IDs
\(100, 110, 120, \dots, 9990, 10000\) to the group sub.

\sphinxAtStartPar
The second format is a \sphinxstyleemphasis{logical} followed by one or two values (type
or ID).  The 7 valid logicals are listed above.  All the logicals
except \sphinxcode{\sphinxupquote{\textless{}\textgreater{}}} take a single argument.  The third example above adds all
atoms with IDs from 1 to 150 to the group named \sphinxstyleemphasis{sub}.  The logical \sphinxcode{\sphinxupquote{\textless{}\textgreater{}}}
means “between” and takes 2 arguments.  The fourth example above adds all
atoms belonging to molecules with IDs from 50 to 250 (inclusive) to
the group named polyA.  For the \sphinxstyleemphasis{type} style, type labels are converted into
numeric types before being evaluated.

\sphinxAtStartPar
The \sphinxstyleemphasis{variable} style evaluates a variable to determine which atoms to
add to the group.  It must be an {\hyperref[\detokenize{variable::doc}]{\sphinxcrossref{\DUrole{doc}{atom\sphinxhyphen{}style variable}}}}
previously defined in the input script.  If the variable evaluates
to a non\sphinxhyphen{}zero value for a particular atom, then that atom is added
to the specified group.

\sphinxAtStartPar
Atom\sphinxhyphen{}style variables can specify formulas that include thermodynamic
quantities, per\sphinxhyphen{}atom values such as atom coordinates, or per\sphinxhyphen{}atom
quantities calculated by computes, fixes, or other variables.  They
can also include Boolean logic where two numeric values are compared to
yield a 1 or 0 (effectively a true or false).  Thus, using the
\sphinxstyleemphasis{variable} style is a general way to flag specific atoms to include
or exclude from a group.

\sphinxAtStartPar
For example, these lines define a variable “eatom” that calculates the
potential energy of each atom and includes it in the group if its
potential energy is above the threshold value \(-3.0\).

\begin{sphinxVerbatim}[commandchars=\\\{\}]
\PYG{k}{compute         }\PYG{n+nv+nvIdentifier}{1}\PYG{+w}{ }\PYG{n+nv+nvIdentifier}{all}\PYG{+w}{ }\PYG{n}{pe}\PYG{o}{/}\PYG{n}{atom}
\PYG{k}{compute         }\PYG{n+nv+nvIdentifier}{2}\PYG{+w}{ }\PYG{n+nv+nvIdentifier}{all}\PYG{+w}{ }\PYG{n}{reduce}\PYG{+w}{ }\PYG{n}{sum}\PYG{+w}{ }\PYG{n}{c\PYGZus{}1}
\PYG{k}{thermo\PYGZus{}style}\PYG{+w}{    }\PYG{n}{custom}\PYG{+w}{ }\PYG{n}{step}\PYG{+w}{ }\PYG{n}{temp}\PYG{+w}{ }\PYG{n}{pe}\PYG{+w}{ }\PYG{n}{c\PYGZus{}2}
\PYG{k}{run}\PYG{+w}{             }\PYG{l+m}{0}

\PYG{k}{variable        }\PYG{n+nv+nvIdentifier}{eatom}\PYG{+w}{ }\PYG{n}{atom}\PYG{+w}{ }\PYG{l+s}{\PYGZdq{}}\PYG{l+s}{c\PYGZus{}1 \PYGZgt{} \PYGZhy{}3.0}\PYG{l+s}{\PYGZdq{}}
\PYG{k}{group           }\PYG{n+nv+nvIdentifier}{hienergy}\PYG{+w}{ }\PYG{k}{variable }\PYG{n+nv+nvIdentifier}{eatom}
\end{sphinxVerbatim}

\sphinxAtStartPar
Note that these lines

\begin{sphinxVerbatim}[commandchars=\\\{\}]
\PYG{k}{compute         }\PYG{n+nv+nvIdentifier}{2}\PYG{+w}{ }\PYG{n+nv+nvIdentifier}{all}\PYG{+w}{ }\PYG{n}{reduce}\PYG{+w}{ }\PYG{n}{sum}\PYG{+w}{ }\PYG{n}{c\PYGZus{}1}
\PYG{k}{thermo\PYGZus{}style}\PYG{+w}{    }\PYG{n}{custom}\PYG{+w}{ }\PYG{n}{step}\PYG{+w}{ }\PYG{n}{temp}\PYG{+w}{ }\PYG{n}{pe}\PYG{+w}{ }\PYG{n}{c\PYGZus{}2}
\PYG{k}{run}\PYG{+w}{             }\PYG{l+m}{0}
\end{sphinxVerbatim}

\sphinxAtStartPar
are necessary to ensure that the “eatom” variable is current when the
group command invokes it.  Because the eatom variable computes the
per\sphinxhyphen{}atom energy via the pe/atom compute, it will only be current if a
run has been performed which evaluated pairwise energies, and the
pe/atom compute was actually invoked during the run.  Printing the
thermodynamic info for compute 2 ensures that this is the case, since
it sums the pe/atom compute values (in the reduce compute) to output
them to the screen.  See the “Variable Accuracy” section of the
{\hyperref[\detokenize{variable::doc}]{\sphinxcrossref{\DUrole{doc}{variable}}}} page for more details on ensuring that
variables are current when they are evaluated between runs.

\sphinxAtStartPar
The \sphinxstyleemphasis{include} style with its arg \sphinxstyleemphasis{molecule} adds atoms to a group that
have the same molecule ID as atoms already in the group.  The molecule
ID = 0 is ignored in this operation, since it is assumed to flag
isolated atoms that are not part of molecules.  An example of where
this operation is useful is if the \sphinxstyleemphasis{region} style has been used
previously to add atoms to a group that are within a geometric region.
If molecules straddle the region boundary, then atoms outside the
region that are part of molecules with atoms inside the region will
not be in the group.  Using the group command a second time with \sphinxstyleemphasis{include
molecule} will add those atoms that are outside the region to the
group.

\begin{sphinxadmonition}{note}{Note:}
\sphinxAtStartPar
The \sphinxstyleemphasis{include molecule} operation is relatively expensive in a
parallel sense.  This is because it requires communication of relevant
molecule IDs between all the processors and each processor to loop
over its atoms once per processor, to compare its atoms to the list of
molecule IDs from every other processor.  Hence it scales as N, rather
than N/P as most of the group operations do, where N is the number of
atoms, and P is the number of processors.
\end{sphinxadmonition}

\sphinxAtStartPar
The \sphinxstyleemphasis{subtract} style takes a list of two or more existing group names
as arguments.  All atoms that belong to the first group, but not to any
of the other groups are added to the specified group.

\sphinxAtStartPar
The \sphinxstyleemphasis{union} style takes a list of one or more existing group names as
arguments.  All atoms that belong to any of the listed groups are
added to the specified group.

\sphinxAtStartPar
The \sphinxstyleemphasis{intersect} style takes a list of two or more existing group names
as arguments.  Atoms that belong to every one of the listed groups are
added to the specified group.


\bigskip\hrule\bigskip


\sphinxAtStartPar
The \sphinxstyleemphasis{dynamic} style flags an existing or new group as dynamic.  This
means atoms will be (re)assigned to the group periodically as a
simulation runs.  This is in contrast to static groups where atoms are
permanently assigned to the group.  The way the assignment occurs is
as follows.  Only atoms in the group specified as the parent group via
the parent\sphinxhyphen{}ID are assigned to the dynamic group before the following
conditions are applied.

\sphinxAtStartPar
If the \sphinxstyleemphasis{region} keyword is used, atoms not in the specified region are
removed from the dynamic group.

\sphinxAtStartPar
If the \sphinxstyleemphasis{var} keyword is used, the variable name must be an atom\sphinxhyphen{}style
or atomfile\sphinxhyphen{}style variable.  The variable is evaluated and atoms whose
per\sphinxhyphen{}atom values are 0.0, are removed from the dynamic group.

\sphinxAtStartPar
If the \sphinxstyleemphasis{property} keyword is used, the name refers to a custom integer
or floating point per\sphinxhyphen{}atom vector defined via the {\hyperref[\detokenize{fix_property_atom::doc}]{\sphinxcrossref{\DUrole{doc}{fix
property/atom}}}} command.  This means the values in
the vector can be read as part of a data file with the {\hyperref[\detokenize{read_data::doc}]{\sphinxcrossref{\DUrole{doc}{read\_data}}}} command or specified with the {\hyperref[\detokenize{set::doc}]{\sphinxcrossref{\DUrole{doc}{set}}}} command.
Or accessed and changed via the {\hyperref[\detokenize{Howto_library::doc}]{\sphinxcrossref{\DUrole{doc}{library interface to LAMMPS}}}}, or by styles you add to LAMMPS (pair, fix, compute,
etc) which access the custom vector and modify its values.  Which
means the values can be modified between or during simulations.  Atoms
whose values in the custom vector are zero are removed from the
dynamic group.  Note that the name of the custom per\sphinxhyphen{}atom vector is
specified just as \sphinxstyleemphasis{name}, not as \sphinxstyleemphasis{i\_name} or \sphinxstyleemphasis{d\_name} as it is for
other commands that use different kinds of custom atom vectors or
arrays as arguments.

\sphinxAtStartPar
The assignment of atoms to a dynamic group is done at the beginning of
each run and on every timestep that is a multiple of \sphinxstyleemphasis{N}, which is
the argument for the \sphinxstyleemphasis{every} keyword (\(N = 1\) is the default).  For an
energy minimization, via the {\hyperref[\detokenize{minimize::doc}]{\sphinxcrossref{\DUrole{doc}{minimize}}}} command, an
assignment is made at the beginning of the minimization, but not
during the iterations of the minimizer.

\sphinxAtStartPar
The point in the timestep at which atoms are assigned to a dynamic
group is after interatomic forces have been computed, but before any
fixes which alter forces or otherwise update the system have been
invoked.  This means that atom positions have been updated, neighbor
lists and ghost atoms are current, and both intermolecular and
intramolecular forces have been calculated based on the new
coordinates.  Thus the region criterion, if applied, should be
accurate.  Also, any computes invoked by an atom\sphinxhyphen{}style variable should
use updated information for that timestep (e.g., potential energy/atom
or coordination number/atom).  Similarly, fixes or computes which are
invoked after that point in the timestep, should operate on the new
group of atoms.

\begin{sphinxadmonition}{note}{Note:}
\sphinxAtStartPar
If the \sphinxstyleemphasis{region} keyword is used to determine what atoms are in
the dynamic group, atoms can move outside of the simulation box
between reneighboring events.  Thus if you want to include all atoms
on the left side of the simulation box, you probably want to set the
left boundary of the region to be outside the simulation box by some
reasonable amount (e.g., up to the cutoff of the potential), else they
may be excluded from the dynamic region.
\end{sphinxadmonition}

\sphinxAtStartPar
Here is an example of using a dynamic group to shrink the set of atoms
being integrated by using a spherical region with a variable radius
(shrinking from 18 to 5 over the course of the run).  This could be
used to model a quench of the system, freezing atoms outside the
shrinking sphere, then converting the remaining atoms to a static
group and running further.

\begin{sphinxVerbatim}[commandchars=\\\{\}]
\PYG{k}{variable        }\PYG{n+nv+nvIdentifier}{nsteps}\PYG{+w}{ }\PYG{n}{equal}\PYG{+w}{ }\PYG{l+m}{5000}
\PYG{k}{variable        }\PYG{n+nv+nvIdentifier}{rad}\PYG{+w}{ }\PYG{n}{equal}\PYG{+w}{ }\PYG{l+m}{18}\PYG{o}{\PYGZhy{}}\PYG{n+nv}{(}\PYG{n+nv}{step/v\PYGZus{}nsteps}\PYG{n+nv}{)}\PYG{o}{*}\PYG{n+nv}{(}\PYG{n+nv}{18\PYGZhy{}5}\PYG{n+nv}{)}
\PYG{k}{region          }\PYG{n+nv+nvIdentifier}{ss}\PYG{+w}{ }\PYG{n}{sphere}\PYG{+w}{ }\PYG{l+m}{20}\PYG{+w}{ }\PYG{l+m}{20}\PYG{+w}{ }\PYG{l+m}{0}\PYG{+w}{ }\PYG{n}{v\PYGZus{}rad}
\PYG{k}{group           }\PYG{n+nv+nvIdentifier}{mobile}\PYG{+w}{ }\PYG{n}{dynamic}\PYG{+w}{ }\PYG{n}{all}\PYG{+w}{ }\PYG{k}{region }\PYG{n+nv+nvIdentifier}{ss}
\PYG{k}{fix             }\PYG{n+nv+nvIdentifier}{1}\PYG{+w}{ }\PYG{n+nv+nvIdentifier}{mobile}\PYG{+w}{ }\PYG{n}{nve}
\PYG{k}{run}\PYG{+w}{             }\PYG{n+nv}{\PYGZdl{}\PYGZob{}}\PYG{n+nv}{nsteps}\PYG{n+nv}{\PYGZcb{}}
\PYG{k}{group           }\PYG{n+nv+nvIdentifier}{mobile}\PYG{+w}{ }\PYG{n}{static}
\PYG{k}{run}\PYG{+w}{             }\PYG{n+nv}{\PYGZdl{}\PYGZob{}}\PYG{n+nv}{nsteps}\PYG{n+nv}{\PYGZcb{}}
\end{sphinxVerbatim}

\begin{sphinxadmonition}{note}{Note:}
\sphinxAtStartPar
All fixes and computes take a group ID as an argument, but they do
not all allow for use of a dynamic group.  If you get an error
message that this is not allowed, but feel that it should be for the
fix or compute in question, then please post your reasoning to the
\sphinxhref{https://matsci.org/c/lammps-development/42}{LAMMPS forum at MatSci}
and we can look into changing it.  The same applies if you come
across inconsistent behavior when dynamic groups are allowed.
\end{sphinxadmonition}

\sphinxAtStartPar
The \sphinxstyleemphasis{static} style removes the setting for a dynamic group, converting
it to a static group (the default).  The atoms in the static group are
those currently in the dynamic group.


\bigskip\hrule\bigskip



\subsection{Restrictions}
\label{\detokenize{group:restrictions}}
\sphinxAtStartPar
There can be no more than 32 groups defined at one time, including
“all”.

\sphinxAtStartPar
The parent group of a dynamic group cannot itself be a dynamic group.


\subsection{Related commands}
\label{\detokenize{group:related-commands}}
\sphinxAtStartPar
{\hyperref[\detokenize{dump::doc}]{\sphinxcrossref{\DUrole{doc}{dump}}}}, {\hyperref[\detokenize{fix::doc}]{\sphinxcrossref{\DUrole{doc}{fix}}}}, {\hyperref[\detokenize{region::doc}]{\sphinxcrossref{\DUrole{doc}{region}}}},
{\hyperref[\detokenize{velocity::doc}]{\sphinxcrossref{\DUrole{doc}{velocity}}}}


\subsection{Default}
\label{\detokenize{group:default}}
\sphinxAtStartPar
All atoms belong to the “all” group.

\sphinxstepscope

\index{group2ndx@\spxentry{group2ndx}}\ignorespaces 
\index{ndx2group@\spxentry{ndx2group}}\ignorespaces 

\section{group2ndx command}
\label{\detokenize{group2ndx:group2ndx-command}}\label{\detokenize{group2ndx:index-1}}\label{\detokenize{group2ndx:index-0}}\label{\detokenize{group2ndx::doc}}

\section{ndx2group command}
\label{\detokenize{group2ndx:ndx2group-command}}

\subsection{Syntax}
\label{\detokenize{group2ndx:syntax}}
\begin{sphinxVerbatim}[commandchars=\\\{\}]
\PYG{k}{group2ndx }\PYG{n+nv+nvIdentifier}{file}\PYG{+w}{ }\PYG{n}{args}
\PYG{k}{ndx2group }\PYG{n+nv+nvIdentifier}{file}\PYG{+w}{ }\PYG{n}{args}
\end{sphinxVerbatim}
\begin{itemize}
\item {} 
\sphinxAtStartPar
file = name of index file to write out or read in

\item {} 
\sphinxAtStartPar
args = zero or more group IDs may be appended

\end{itemize}


\subsection{Examples}
\label{\detokenize{group2ndx:examples}}
\begin{sphinxVerbatim}[commandchars=\\\{\}]
\PYG{k}{group2ndx }\PYG{n+nv+nvIdentifier}{allindex.ndx}
\PYG{k}{group2ndx }\PYG{n+nv+nvIdentifier}{someindex.ndx}\PYG{+w}{ }\PYG{n}{upper}\PYG{+w}{ }\PYG{n}{lower}\PYG{+w}{ }\PYG{n}{mobile}
\PYG{k}{ndx2group }\PYG{n+nv+nvIdentifier}{someindex.ndx}
\PYG{k}{ndx2group }\PYG{n+nv+nvIdentifier}{someindex.ndx}\PYG{+w}{ }\PYG{n}{mobile}
\end{sphinxVerbatim}


\subsection{Description}
\label{\detokenize{group2ndx:description}}
\sphinxAtStartPar
Write or read a Gromacs style index file in text format that associates
atom IDs with the corresponding group definitions. This index file can be
used with in combination with Gromacs analysis tools or to import group
definitions into the {\hyperref[\detokenize{fix_colvars::doc}]{\sphinxcrossref{\DUrole{doc}{fix colvars}}}} input file.

\sphinxAtStartPar
It can also be used to save and restore group definitions for static groups
using the individual atom IDs. This may be important if the original
group definition depends on a region or otherwise on the geometry and thus
cannot be easily recreated.

\sphinxAtStartPar
Another application would be to import atom groups defined for Gromacs
simulation into LAMMPS.  When translating Gromacs topology and geometry
data to LAMMPS.

\sphinxAtStartPar
The \sphinxstyleemphasis{group2ndx} command will write group definitions to an index file.
Without specifying any group IDs, all groups will be written to the
index file.  When specifying group IDs, only those groups will be
written to the index file.  In order to follow the Gromacs conventions,
the group \sphinxstyleemphasis{all} will be renamed to \sphinxstyleemphasis{System} in the index file.

\sphinxAtStartPar
The \sphinxstyleemphasis{ndx2group} command will create of update group definitions from
those stored in an index file.  Without specifying any group IDs, all
groups except \sphinxstyleemphasis{System} will be read from the index file and the
corresponding groups recreated.  If a group of the same name already
exists, it will be completely reset.  When specifying group IDs, those
groups, if present, will be read from the index file and restored.


\subsection{File Format}
\label{\detokenize{group2ndx:file-format}}
\sphinxAtStartPar
The file format is equivalent and compatible with what is produced by
the \sphinxhref{https://manual.gromacs.org/current/onlinehelp/gmx-make\_ndx.html}{Gromacs make\_ndx command}.
and follows the \sphinxhref{https://manual.gromacs.org/current/reference-manual/file-formats.html\#ndx}{Gromacs definition of an ndx file}

\sphinxAtStartPar
Each group definition begins with the group name in square brackets with
blanks, e.g. {[} water {]} and is then followed by the list of atom
indices, which may be spread over multiple lines.  Here is a small
example file:

\begin{sphinxVerbatim}[commandchars=\\\{\}]
\PYG{k}{[ Oxygen ]}
\PYG{n+na}{1  4  7}
\PYG{k}{[ Hydrogen ]}
\PYG{n+na}{2  3  5  6}
\PYG{n+na}{8  9}
\PYG{k}{[ Water ]}
\PYG{n+na}{1 2 3 4 5 6 7 8 9}
\end{sphinxVerbatim}

\sphinxAtStartPar
The index file defines 3 groups: Oxygen, Hydrogen, and Water and the
latter happens to be the union of the first two.


\bigskip\hrule\bigskip



\subsection{Restrictions}
\label{\detokenize{group2ndx:restrictions}}
\sphinxAtStartPar
These commands require that atoms have atom IDs, since this is the
information that is written to the index file.

\sphinxAtStartPar
These commands are part of the EXTRA\sphinxhyphen{}COMMAND package.  They are only
enabled if LAMMPS was built with that package.  See the
{\hyperref[\detokenize{Build_package::doc}]{\sphinxcrossref{\DUrole{doc}{Build package}}}} page for more info.


\subsection{Related commands}
\label{\detokenize{group2ndx:related-commands}}
\sphinxAtStartPar
{\hyperref[\detokenize{group::doc}]{\sphinxcrossref{\DUrole{doc}{group}}}}, {\hyperref[\detokenize{dump::doc}]{\sphinxcrossref{\DUrole{doc}{dump}}}}, {\hyperref[\detokenize{fix_colvars::doc}]{\sphinxcrossref{\DUrole{doc}{fix colvars}}}}


\subsection{Default}
\label{\detokenize{group2ndx:default}}
\sphinxAtStartPar
none

\sphinxstepscope

\index{hyper@\spxentry{hyper}}\ignorespaces 

\section{hyper command}
\label{\detokenize{hyper:hyper-command}}\label{\detokenize{hyper:index-0}}\label{\detokenize{hyper::doc}}

\subsection{Syntax}
\label{\detokenize{hyper:syntax}}
\begin{sphinxVerbatim}[commandchars=\\\{\}]
\PYG{k}{hyper}\PYG{+w}{ }\PYG{n}{N}\PYG{+w}{ }\PYG{n}{Nevent}\PYG{+w}{ }\PYG{n}{fix}\PYG{o}{\PYGZhy{}}\PYG{n}{ID}\PYG{+w}{ }\PYG{n}{compute}\PYG{o}{\PYGZhy{}}\PYG{n}{ID}\PYG{+w}{ }\PYG{n}{keyword}\PYG{+w}{ }\PYG{n}{values}\PYG{+w}{ }\PYG{n}{...}
\end{sphinxVerbatim}
\begin{itemize}
\item {} 
\sphinxAtStartPar
N = \# of timesteps to run

\item {} 
\sphinxAtStartPar
Nevent = check for events every this many steps

\item {} 
\sphinxAtStartPar
fix\sphinxhyphen{}ID = ID of a fix that applies a global or local bias potential, can be NULL

\item {} 
\sphinxAtStartPar
compute\sphinxhyphen{}ID = ID of a compute that identifies when an event has occurred

\item {} 
\sphinxAtStartPar
zero or more keyword/value pairs may be appended

\item {} 
\sphinxAtStartPar
keyword = \sphinxstyleemphasis{min} or \sphinxstyleemphasis{dump} or \sphinxstyleemphasis{rebond}
\begin{sphinxalltt}
\sphinxstyleemphasis{min} values = etol ftol maxiter maxeval
  etol = stopping tolerance for energy, used in quenching
  ftol = stopping tolerance for force, used in quenching
  maxiter = max iterations of minimize, used in quenching
  maxeval = max number of force/energy evaluations, used in quenching
\sphinxstyleemphasis{dump} value = dump\sphinxhyphen{}ID
  dump\sphinxhyphen{}ID = ID of dump to trigger whenever an event takes place
\sphinxstyleemphasis{rebond} value = Nrebond
  Nrebond = frequency at which to reset bonds, even if no event has occurred
\end{sphinxalltt}

\end{itemize}


\subsection{Examples}
\label{\detokenize{hyper:examples}}
\begin{sphinxVerbatim}[commandchars=\\\{\}]
\PYG{k}{compute }\PYG{n+nv+nvIdentifier}{event}\PYG{+w}{ }\PYG{n+nv+nvIdentifier}{all}\PYG{+w}{ }\PYG{n}{event}\PYG{o}{/}\PYG{n}{displace}\PYG{+w}{ }\PYG{l+m}{1.0}
\PYG{k}{fix }\PYG{n+nv+nvIdentifier}{HG}\PYG{+w}{ }\PYG{n+nv+nvIdentifier}{mobile}\PYG{+w}{ }\PYG{n}{hyper}\PYG{o}{/}\PYG{n}{global}\PYG{+w}{ }\PYG{l+m}{3.0}\PYG{+w}{ }\PYG{l+m}{0.3}\PYG{+w}{ }\PYG{l+m}{0.4}\PYG{+w}{ }\PYG{l+m}{800.0}
\PYG{k}{hyper}\PYG{+w}{ }\PYG{l+m}{5000}\PYG{+w}{ }\PYG{l+m}{100}\PYG{+w}{ }\PYG{n}{HG}\PYG{+w}{ }\PYG{n}{event}\PYG{+w}{ }\PYG{n}{min}\PYG{+w}{ }\PYG{l+m}{1.0e\PYGZhy{}6}\PYG{+w}{ }\PYG{l+m}{1.0e\PYGZhy{}6}\PYG{+w}{ }\PYG{l+m}{100}\PYG{+w}{ }\PYG{l+m}{100}\PYG{+w}{ }\PYG{k}{dump }\PYG{n+nv+nvIdentifier}{1}\PYG{+w}{ }\PYG{n+nv+nvIdentifier}{dump}\PYG{+w}{ }\PYG{l+m}{5}
\end{sphinxVerbatim}


\subsection{Description}
\label{\detokenize{hyper:description}}
\sphinxAtStartPar
Run a bond\sphinxhyphen{}boost hyperdynamics (HD) simulation where time is
accelerated by application of a bias potential to one or more pairs of
nearby atoms in the system.  This command can be used to run both
global and local hyperdynamics.  In global HD a single bond within the
system is biased on each timestep.  In local HD multiple bonds
(separated by a sufficient distance) can be biased simultaneously at
each timestep.  In the bond\sphinxhyphen{}boost hyperdynamics context, a “bond” is
not a covalent bond between a pair of atoms in a molecule.  Rather it
is simply a pair of nearby atoms as discussed below.

\sphinxAtStartPar
Both global and local HD are described in {\hyperref[\detokenize{hyper:voter2013}]{\sphinxcrossref{\DUrole{std,std-ref}{(Voter2013)}}}} by Art Voter and collaborators.  Similar to parallel
replica dynamics (PRD), global and local HD are methods for performing
accelerated dynamics that are suitable for infrequent\sphinxhyphen{}event systems
that obey first\sphinxhyphen{}order kinetics.  A good overview of accelerated
dynamics methods (AMD) for such systems in given in {\hyperref[\detokenize{hyper:voter2002hd}]{\sphinxcrossref{\DUrole{std,std-ref}{(Voter2002)}}}} from the same group.  To quote from the review paper:
“The dynamical evolution is characterized by vibrational excursions
within a potential basin, punctuated by occasional transitions between
basins.  The transition probability is characterized by p(t) =
k*exp(\sphinxhyphen{}kt) where k is the rate constant.”

\sphinxAtStartPar
Both HD and PRD produce a time\sphinxhyphen{}accurate trajectory that effectively
extends the timescale over which a system can be simulated, but they
do it differently.  HD uses a single replica of the system and
accelerates time by biasing the interaction potential in a manner such
that each timestep is effectively longer.  PRD creates Nr replicas of
the system and runs dynamics on each independently with a normal
unbiased potential until an event occurs in one of the replicas.  The
time between events is reduced by a factor of Nr replicas.  For both
methods, per CPU second, more physical time elapses and more events
occur.  See the {\hyperref[\detokenize{prd::doc}]{\sphinxcrossref{\DUrole{doc}{prd}}}} page for more info about PRD.

\sphinxAtStartPar
An HD run has several stages, which are repeated each time an event
occurs, as explained below.  The logic for an HD run is as follows:

\begin{sphinxVerbatim}[commandchars=\\\{\}]
quench
create initial list of bonds

while (time remains):
  run dynamics for Nevent steps
  quench
  check for an event
  if event occurred: reset list of bonds
  restore pre\PYGZhy{}quench state
\end{sphinxVerbatim}

\sphinxAtStartPar
The list of bonds is the list of atom pairs of atoms that are within a
short cutoff distance of each other after the system energy is
minimized (quenched).  This list is created and reset by a {\hyperref[\detokenize{fix_hyper_global::doc}]{\sphinxcrossref{\DUrole{doc}{fix hyper/global}}}} or {\hyperref[\detokenize{fix_hyper_local::doc}]{\sphinxcrossref{\DUrole{doc}{fix hyper/local}}}} command specified as \sphinxstyleemphasis{fix\sphinxhyphen{}ID}.  At
every dynamics timestep, the same fix selects one of more bonds to
apply a bias potential to.

\begin{sphinxadmonition}{note}{Note:}
\sphinxAtStartPar
The style of fix associated with the specified \sphinxstyleemphasis{fix\sphinxhyphen{}ID}
determines whether you are running the global versus local
hyperdynamics algorithm.
\end{sphinxadmonition}

\sphinxAtStartPar
Dynamics (with the bias potential) is run continuously, stopping every
\sphinxstyleemphasis{Nevent} steps to check if a transition event has occurred.  The
specified \sphinxstyleemphasis{N} for total steps must be a multiple of \sphinxstyleemphasis{Nevent}.  check
is performed by quenching the system and comparing the resulting atom
coordinates to the coordinates from the previous basin.

\sphinxAtStartPar
A quench is an energy minimization and is performed by whichever
algorithm has been defined by the {\hyperref[\detokenize{min_style::doc}]{\sphinxcrossref{\DUrole{doc}{min\_style}}}} command.
Minimization parameters may be set via the
{\hyperref[\detokenize{min_modify::doc}]{\sphinxcrossref{\DUrole{doc}{min\_modify}}}} command and by the \sphinxstyleemphasis{min} keyword of the
hyper command.  The latter are the settings that would be used with
the {\hyperref[\detokenize{minimize::doc}]{\sphinxcrossref{\DUrole{doc}{minimize}}}} command.  Note that typically, you do not
need to perform a highly\sphinxhyphen{}converged minimization to detect a transition
event, though you may need to in order to prevent a set of atoms in
the system from relaxing to a saddle point.

\sphinxAtStartPar
The event check is performed by a compute with the specified
\sphinxstyleemphasis{compute\sphinxhyphen{}ID}.  Currently there is only one compute that works with the
hyper command, which is the {\hyperref[\detokenize{compute_event_displace::doc}]{\sphinxcrossref{\DUrole{doc}{compute event/displace}}}} command.  Other
event\sphinxhyphen{}checking computes may be added.  {\hyperref[\detokenize{compute_event_displace::doc}]{\sphinxcrossref{\DUrole{doc}{Compute event/displace}}}} checks whether any atom in
the compute group has moved further than a specified threshold
distance.  If so, an event has occurred.

\sphinxAtStartPar
If this happens, the list of bonds is reset, since some bond pairs
are likely now too far apart, and new pairs are likely close enough
to be considered a bond.  The pre\sphinxhyphen{}quenched state of the
system (coordinates and velocities) is restored, and dynamics continue.

\sphinxAtStartPar
At the end of the hyper run, a variety of statistics are output to the
screen and logfile.  These include info relevant to both global and
local hyperdynamics, such as the number of events and the elapsed
hyper time (accelerated time), And it includes info specific to one or
the other, depending on which style of fix was specified by \sphinxstyleemphasis{fix\sphinxhyphen{}ID}.


\bigskip\hrule\bigskip


\sphinxAtStartPar
The optional keywords operate as follows.

\sphinxAtStartPar
As explained above, the \sphinxstyleemphasis{min} keyword can be used to specify
parameters for the quench.  Their meaning is the same
as for the {\hyperref[\detokenize{minimize::doc}]{\sphinxcrossref{\DUrole{doc}{minimize}}}} command

\sphinxAtStartPar
The \sphinxstyleemphasis{dump} keyword can be used to trigger a specific dump command with
the specified \sphinxstyleemphasis{dump\sphinxhyphen{}ID} to output a snapshot each time an event is
detected.  It can be specified multiple times with different \sphinxstyleemphasis{dump\sphinxhyphen{}ID}
values, as in the example above.  These snapshots will be for the
quenched state of the system on a timestep that is a multiple of
\sphinxstyleemphasis{Nevent}, i.e. a timestep after the event has occurred.  Note that any
dump command in the input script will also output snapshots at
whatever timestep interval it defines via its \sphinxstyleemphasis{N} argument; see the
{\hyperref[\detokenize{dump::doc}]{\sphinxcrossref{\DUrole{doc}{dump}}}} command for details.  This means if you only want a
particular dump to output snapshots when events are detected, you
should specify its \sphinxstyleemphasis{N} as a value larger than the length of the
hyperdynamics run.

\sphinxAtStartPar
As in the code logic above, the bond list is normally only reset when
an event occurs.  The \sphinxstyleemphasis{rebond} keyword will force a reset of the bond
list every \sphinxstyleemphasis{Nrebond} steps, even if an event has not occurred.
\sphinxstyleemphasis{Nrebond} must be a multiple of \sphinxstyleemphasis{Nevent}.  This can be useful to check
if more frequent resets alter event statistics, perhaps because the
parameters chosen for defining what is a bond and what is an event are
producing bad dynamics in the presence of the bias potential.


\bigskip\hrule\bigskip



\subsection{Restrictions}
\label{\detokenize{hyper:restrictions}}
\sphinxAtStartPar
This command can only be used if LAMMPS was built with the REPLICA
package.  See the {\hyperref[\detokenize{Build_package::doc}]{\sphinxcrossref{\DUrole{doc}{Build package}}}} doc
page for more info.


\subsection{Related commands}
\label{\detokenize{hyper:related-commands}}
\sphinxAtStartPar
{\hyperref[\detokenize{fix_hyper_global::doc}]{\sphinxcrossref{\DUrole{doc}{fix hyper/global}}}}, {\hyperref[\detokenize{fix_hyper_local::doc}]{\sphinxcrossref{\DUrole{doc}{fix hyper/local}}}}, {\hyperref[\detokenize{compute_event_displace::doc}]{\sphinxcrossref{\DUrole{doc}{compute event/displace}}}}, {\hyperref[\detokenize{prd::doc}]{\sphinxcrossref{\DUrole{doc}{prd}}}}


\subsection{Default}
\label{\detokenize{hyper:default}}
\sphinxAtStartPar
The option defaults are min = 0.1 0.1 40 50 and time = steps.


\bigskip\hrule\bigskip

\phantomsection\label{\detokenize{hyper:voter2013}}
\sphinxAtStartPar
\sphinxstylestrong{(Voter2013)} S. Y. Kim, D. Perez, A. F. Voter, J Chem Phys, 139,
144110 (2013).

\phantomsection\label{\detokenize{hyper:voter2002hd}}
\sphinxAtStartPar
\sphinxstylestrong{(Voter2002)} Voter, Montalenti, Germann, Annual Review of Materials
Research 32, 321 (2002).

\sphinxstepscope

\index{if@\spxentry{if}}\ignorespaces 

\section{if command}
\label{\detokenize{if:if-command}}\label{\detokenize{if:index-0}}\label{\detokenize{if::doc}}

\subsection{Syntax}
\label{\detokenize{if:syntax}}
\begin{sphinxVerbatim}[commandchars=\\\{\}]
\PYG{k}{if}\PYG{+w}{ }\PYG{n}{boolean}\PYG{+w}{ }\PYG{k}{then}\PYG{+w}{ }\PYG{n}{t1}\PYG{+w}{ }\PYG{n}{t2}\PYG{+w}{ }\PYG{n}{...}\PYG{+w}{ }\PYG{k}{elif}\PYG{+w}{ }\PYG{n}{boolean}\PYG{+w}{ }\PYG{n}{f1}\PYG{+w}{ }\PYG{n}{f2}\PYG{+w}{ }\PYG{n}{...}\PYG{+w}{ }\PYG{k}{elif}\PYG{+w}{ }\PYG{n}{boolean}\PYG{+w}{ }\PYG{n}{f1}\PYG{+w}{ }\PYG{n}{f2}\PYG{+w}{ }\PYG{n}{...}\PYG{+w}{ }\PYG{k}{else}\PYG{+w}{ }\PYG{n}{e1}\PYG{+w}{ }\PYG{n}{e2}\PYG{+w}{ }\PYG{n}{...}
\end{sphinxVerbatim}
\begin{itemize}
\item {} 
\sphinxAtStartPar
boolean = a Boolean expression evaluated as TRUE or FALSE (see below)

\item {} 
\sphinxAtStartPar
then = required word

\item {} 
\sphinxAtStartPar
t1,t2,…,tN = one or more LAMMPS commands to execute if condition is met, each enclosed in quotes

\item {} 
\sphinxAtStartPar
elif = optional word, can appear multiple times

\item {} 
\sphinxAtStartPar
f1,f2,…,fN = one or more LAMMPS commands to execute if elif condition is met, each enclosed in quotes (optional arguments)

\item {} 
\sphinxAtStartPar
else = optional argument

\item {} 
\sphinxAtStartPar
e1,e2,…,eN = one or more LAMMPS commands to execute if no condition is met, each enclosed in quotes (optional arguments)

\end{itemize}


\subsection{Examples}
\label{\detokenize{if:examples}}
\begin{sphinxVerbatim}[commandchars=\\\{\}]
\PYG{k}{if}\PYG{+w}{ }\PYG{l+s}{\PYGZdq{}}\PYG{l+s}{\PYGZdl{}\PYGZob{}steps\PYGZcb{} \PYGZgt{} 1000}\PYG{l+s}{\PYGZdq{}}\PYG{+w}{ }\PYG{k}{then}\PYG{+w}{ }\PYG{n}{quit}
\PYG{k}{if}\PYG{+w}{ }\PYG{l+s}{\PYGZdq{}}\PYG{l+s}{\PYGZdl{}\PYGZob{}myString\PYGZcb{} == a10}\PYG{l+s}{\PYGZdq{}}\PYG{+w}{ }\PYG{k}{then}\PYG{+w}{ }\PYG{n}{quit}
\PYG{k}{if}\PYG{+w}{ }\PYG{l+s}{\PYGZdq{}}\PYG{l+s}{\PYGZdl{}x \PYGZlt{}= \PYGZdl{}y}\PYG{l+s}{\PYGZdq{}}\PYG{+w}{ }\PYG{k}{then}\PYG{+w}{ }\PYG{l+s}{\PYGZdq{}}\PYG{l+s}{print \PYGZsq{}X is smaller = \PYGZdl{}x\PYGZsq{}}\PYG{l+s}{\PYGZdq{}}\PYG{+w}{ }\PYG{k}{else}\PYG{+w}{ }\PYG{l+s}{\PYGZdq{}}\PYG{l+s}{print \PYGZsq{}Y is smaller = \PYGZdl{}y\PYGZsq{}}\PYG{l+s}{\PYGZdq{}}
\PYG{k}{if}\PYG{+w}{ }\PYG{l+s}{\PYGZdq{}}\PYG{l+s}{(\PYGZdl{}\PYGZob{}eng\PYGZcb{} \PYGZgt{} 0.0) || (\PYGZdl{}n \PYGZlt{} 1000)}\PYG{l+s}{\PYGZdq{}}\PYG{+w}{ }\PYG{k}{then}\PYG{+w}{ }\PYG{o}{\PYGZam{}}
\PYG{+w}{  }\PYG{l+s}{\PYGZdq{}}\PYG{l+s}{timestep 0.005}\PYG{l+s}{\PYGZdq{}}\PYG{+w}{ }\PYG{o}{\PYGZam{}}
\PYG{k}{elif}\PYG{+w}{ }\PYG{n+nv}{\PYGZdl{}n}\PYG{o}{\PYGZlt{}}\PYG{l+m}{10000}\PYG{+w}{ }\PYG{o}{\PYGZam{}}
\PYG{+w}{  }\PYG{l+s}{\PYGZdq{}}\PYG{l+s}{timestep 0.01}\PYG{l+s}{\PYGZdq{}}\PYG{+w}{ }\PYG{o}{\PYGZam{}}
\PYG{k}{else}\PYG{+w}{ }\PYG{o}{\PYGZam{}}
\PYG{+w}{  }\PYG{l+s}{\PYGZdq{}}\PYG{l+s}{timestep 0.02}\PYG{l+s}{\PYGZdq{}}\PYG{+w}{ }\PYG{o}{\PYGZam{}}
\PYG{+w}{  }\PYG{l+s}{\PYGZdq{}}\PYG{l+s}{print \PYGZsq{}Max step reached\PYGZsq{}}\PYG{l+s}{\PYGZdq{}}
\PYG{k}{if}\PYG{+w}{ }\PYG{l+s}{\PYGZdq{}}\PYG{l+s}{\PYGZdl{}\PYGZob{}eng\PYGZcb{} \PYGZgt{} \PYGZdl{}\PYGZob{}eng\PYGZus{}previous\PYGZcb{}}\PYG{l+s}{\PYGZdq{}}\PYG{+w}{ }\PYG{k}{then}\PYG{+w}{ }\PYG{l+s}{\PYGZdq{}}\PYG{l+s}{jump file1}\PYG{l+s}{\PYGZdq{}}\PYG{+w}{ }\PYG{k}{else}\PYG{+w}{ }\PYG{l+s}{\PYGZdq{}}\PYG{l+s}{jump file2}\PYG{l+s}{\PYGZdq{}}
\end{sphinxVerbatim}


\subsection{Description}
\label{\detokenize{if:description}}
\sphinxAtStartPar
This command provides an if\sphinxhyphen{}then\sphinxhyphen{}else capability within an input
script.  A Boolean expression is evaluated and the result is TRUE or
FALSE.  Note that as in the examples above, the expression can contain
variables, as defined by the {\hyperref[\detokenize{variable::doc}]{\sphinxcrossref{\DUrole{doc}{variable}}}} command, which
will be evaluated as part of the expression.  Thus a user\sphinxhyphen{}defined
formula that reflects the current state of the simulation can be used
to issue one or more new commands.

\sphinxAtStartPar
If the result of the Boolean expression is TRUE, then one or more
commands (t1, t2, …, tN) are executed.  If it is FALSE, then Boolean
expressions associated with successive elif keywords are evaluated
until one is found to be true, in which case its commands (f1, f2,
…, fN) are executed.  If no Boolean expression is TRUE, then the
commands associated with the else keyword, namely (e1, e2, …, eN),
are executed.  The elif and else keywords and their associated
commands are optional.  If they are not specified and the initial
Boolean expression is FALSE, then no commands are executed.

\sphinxAtStartPar
The syntax for Boolean expressions is described below.

\sphinxAtStartPar
Each command (t1, f1, e1, etc.) can be any valid LAMMPS input script
command.  If the command is more than one word, it must enclosed in
quotes, so it will be treated as a single argument, as in the examples
above.

\begin{sphinxadmonition}{note}{Note:}
\sphinxAtStartPar
If a command itself requires a quoted argument (e.g., a
{\hyperref[\detokenize{print::doc}]{\sphinxcrossref{\DUrole{doc}{print}}}} command), then double and single quotes can be used
and nested in the usual manner, as in the examples above and below.
The {\hyperref[\detokenize{Commands_parse::doc}]{\sphinxcrossref{\DUrole{doc}{Commands parse}}}} page has more details on
using quotes in arguments.  Only one of level of nesting is allowed,
but that should be sufficient for most use cases.
\end{sphinxadmonition}

\sphinxAtStartPar
Note that by using the line continuation character “\&”, the if command
can be spread across many lines, though it is still a single command:

\begin{sphinxVerbatim}[commandchars=\\\{\}]
\PYG{k}{if}\PYG{+w}{ }\PYG{l+s}{\PYGZdq{}}\PYG{l+s}{\PYGZdl{}a \PYGZlt{} \PYGZdl{}b}\PYG{l+s}{\PYGZdq{}}\PYG{+w}{ }\PYG{k}{then}\PYG{+w}{ }\PYG{o}{\PYGZam{}}
\PYG{+w}{  }\PYG{l+s}{\PYGZdq{}}\PYG{l+s}{print \PYGZsq{}Minimum value = \PYGZdl{}a\PYGZsq{}}\PYG{l+s}{\PYGZdq{}}\PYG{+w}{ }\PYG{o}{\PYGZam{}}
\PYG{+w}{  }\PYG{l+s}{\PYGZdq{}}\PYG{l+s}{run 1000}\PYG{l+s}{\PYGZdq{}}\PYG{+w}{ }\PYG{o}{\PYGZam{}}
\PYG{k}{else}\PYG{+w}{ }\PYG{o}{\PYGZam{}}
\PYG{+w}{  }\PYG{l+s}{\PYGZsq{}}\PYG{l+s}{print \PYGZdq{}Minimum value = \PYGZdl{}b\PYGZdq{}}\PYG{l+s}{\PYGZsq{}}\PYG{+w}{ }\PYG{o}{\PYGZam{}}
\PYG{+w}{  }\PYG{l+s}{\PYGZdq{}}\PYG{l+s}{minimize 0.001 0.001 1000 10000}\PYG{l+s}{\PYGZdq{}}
\end{sphinxVerbatim}

\sphinxAtStartPar
Note that if one of the commands to execute is {\hyperref[\detokenize{quit::doc}]{\sphinxcrossref{\DUrole{doc}{quit}}}}, as in
the first example above, then executing the command will cause LAMMPS
to halt.

\sphinxAtStartPar
Note that by jumping to a label in the same input script, the if
command can be used to break out of a loop.  See the {\hyperref[\detokenize{variable::doc}]{\sphinxcrossref{\DUrole{doc}{variable delete}}}} command for info on how to delete the associated
loop variable, so that it can be re\sphinxhyphen{}used later in the input script.

\sphinxAtStartPar
Here is an example of a loop which checks every 1000 steps if the
system temperature has reached a certain value, and if so, breaks out
of the loop to finish the run.  Note that any variable could be
checked, so long as it is current on the timestep when the run
completes.  As explained on the {\hyperref[\detokenize{variable::doc}]{\sphinxcrossref{\DUrole{doc}{variable}}}} doc page,
this can be ensured by including the variable in thermodynamic output.

\begin{sphinxVerbatim}[commandchars=\\\{\}]
\PYG{k}{variable }\PYG{n+nv+nvIdentifier}{myTemp}\PYG{+w}{ }\PYG{n}{equal}\PYG{+w}{ }\PYG{n}{temp}
\PYG{k}{label }\PYG{l+s+sc}{loop}
\PYG{k}{variable }\PYG{n+nv+nvIdentifier}{a}\PYG{+w}{ }\PYG{n}{loop}\PYG{+w}{ }\PYG{l+m}{1000}
\PYG{k}{run}\PYG{+w}{ }\PYG{l+m}{1000}
\PYG{k}{if}\PYG{+w}{ }\PYG{l+s}{\PYGZdq{}}\PYG{l+s}{\PYGZdl{}\PYGZob{}myTemp\PYGZcb{} \PYGZlt{} 300.0}\PYG{l+s}{\PYGZdq{}}\PYG{+w}{ }\PYG{k}{then}\PYG{+w}{ }\PYG{l+s}{\PYGZdq{}}\PYG{l+s}{jump SELF break}\PYG{l+s}{\PYGZdq{}}
\PYG{k}{next }\PYG{n+nv+nvIdentifier}{a}
\PYG{k}{jump }\PYG{l+s+sc}{SELF}\PYG{+w}{ }\PYG{n}{loop}
\PYG{k}{label }\PYG{l+s+sc}{break}
\PYG{k}{print}\PYG{+w}{ }\PYG{l+s}{\PYGZdq{}}\PYG{l+s}{ALL DONE}\PYG{l+s}{\PYGZdq{}}
\end{sphinxVerbatim}

\sphinxAtStartPar
Here is an example of a double loop which uses the if and
{\hyperref[\detokenize{jump::doc}]{\sphinxcrossref{\DUrole{doc}{jump}}}} commands to break out of the inner loop when a
condition is met, then continues iterating through the outer loop.

\begin{sphinxVerbatim}[commandchars=\\\{\}]
\PYG{k}{label       }\PYG{l+s+sc}{loopa}
\PYG{k}{variable    }\PYG{n+nv+nvIdentifier}{a}\PYG{+w}{ }\PYG{n}{loop}\PYG{+w}{ }\PYG{l+m}{5}
\PYG{+w}{  }\PYG{k}{label     }\PYG{l+s+sc}{loopb}
\PYG{+w}{  }\PYG{k}{variable  }\PYG{n+nv+nvIdentifier}{b}\PYG{+w}{ }\PYG{n}{loop}\PYG{+w}{ }\PYG{l+m}{5}
\PYG{+w}{    }\PYG{n}{print}\PYG{+w}{     }\PYG{l+s}{\PYGZdq{}}\PYG{l+s}{A,B = \PYGZdl{}a,\PYGZdl{}b}\PYG{l+s}{\PYGZdq{}}
\PYG{+w}{    }\PYG{n}{run}\PYG{+w}{       }\PYG{l+m}{10000}
\PYG{+w}{    }\PYG{k}{if}\PYG{+w}{        }\PYG{l+s}{\PYGZdq{}}\PYG{l+s}{\PYGZdl{}b \PYGZgt{} 2}\PYG{l+s}{\PYGZdq{}}\PYG{+w}{ }\PYG{k}{then}\PYG{+w}{ }\PYG{l+s}{\PYGZdq{}}\PYG{l+s}{jump SELF break}\PYG{l+s}{\PYGZdq{}}
\PYG{+w}{  }\PYG{k}{next      }\PYG{n+nv+nvIdentifier}{b}
\PYG{+w}{  }\PYG{k}{jump      }\PYG{l+s+sc}{in.script}\PYG{+w}{ }\PYG{n}{loopb}
\PYG{+w}{  }\PYG{k}{label       }\PYG{l+s+sc}{break}
\PYG{+w}{  }\PYG{k}{variable    }\PYG{n+nv+nvIdentifier}{b}\PYG{+w}{ }\PYG{n}{delete}
\PYG{k}{next        }\PYG{n+nv+nvIdentifier}{a}
\PYG{k}{jump        }\PYG{l+s+sc}{SELF}\PYG{+w}{ }\PYG{n}{loopa}
\end{sphinxVerbatim}


\bigskip\hrule\bigskip


\sphinxAtStartPar
The Boolean expressions for the if and elif keywords have a C\sphinxhyphen{}like
syntax.  Note that each expression is a single argument within the if
command.  Thus if you want to include spaces in the expression for
clarity, you must enclose the entire expression in quotes.

\sphinxAtStartPar
An expression is built out of numbers (which start with a digit or
period or minus sign) or strings (which start with a letter and can
contain alphanumeric characters, underscores, or forward slashes):

\begin{sphinxVerbatim}[commandchars=\\\{\}]
0.2, 100, 1.0e20, \PYGZhy{}15.4, ...
InP, myString, a123, ab\PYGZus{}23\PYGZus{}cd, lj/cut, ...
\end{sphinxVerbatim}

\sphinxAtStartPar
and Boolean operators:
\begin{sphinxalltt}
A == B, A != B, A \textless{} B, A \textless{}= B, A \textgreater{} B, A \textgreater{}= B, A \&\& B, A || B, A |\textasciicircum{} B, !A
\end{sphinxalltt}

\sphinxAtStartPar
Each A and B is a number or string or a variable reference like \sphinxcode{\sphinxupquote{\$a}} or
\sphinxcode{\sphinxupquote{\$\{abc\}}}, or A or B can be another Boolean expression.

\sphinxAtStartPar
Note that all variables used will be substituted for before the
Boolean expression in evaluated.  A variable can produce a number,
like an {\hyperref[\detokenize{variable::doc}]{\sphinxcrossref{\DUrole{doc}{equal\sphinxhyphen{}style variable}}}}, or it can produce a
string, like an {\hyperref[\detokenize{variable::doc}]{\sphinxcrossref{\DUrole{doc}{index\sphinxhyphen{}style variable}}}}.

\sphinxAtStartPar
The Boolean operators \sphinxcode{\sphinxupquote{==}} and \sphinxcode{\sphinxupquote{!=}} can operate on a pair or strings
or numbers.  They cannot compare a number to a string.  All the other
Boolean operations can only operate on numbers.

\sphinxAtStartPar
Expressions are evaluated left to right and have the usual C\sphinxhyphen{}style
precedence: the unary logical NOT operator \sphinxcode{\sphinxupquote{!}} has the highest
precedence, the 4 relational operators \sphinxcode{\sphinxupquote{\textless{}}}, \sphinxcode{\sphinxupquote{\textless{}=}}, \sphinxcode{\sphinxupquote{\textgreater{}}}, and \sphinxcode{\sphinxupquote{\textgreater{}=}} are
next; the two remaining relational operators \sphinxcode{\sphinxupquote{==}} and \sphinxcode{\sphinxupquote{!=}} are next;
then the logical AND operator \sphinxcode{\sphinxupquote{\&\&}}; and finally the logical OR
operator \sphinxcode{\sphinxupquote{||}} and logical XOR (exclusive or) operator \sphinxcode{\sphinxupquote{|\textasciicircum{}}} have
the lowest precedence.  Parenthesis can be used to group one or more
portions of an expression and/or enforce a different order of
evaluation than what would occur with the default precedence.

\sphinxAtStartPar
When the six relational operators (first six in list above) compare two
numbers, they return either a 1.0 or 0.0 depending on whether the
relationship between A and B is TRUE or FALSE.

\sphinxAtStartPar
When the three logical operators (last three in list above) compare two numbers,
they also return either a 1.0 or 0.0 depending on whether the
relationship between A and B is TRUE or FALSE (or just A).  The
logical AND operator will return 1.0 if both its arguments are
non\sphinxhyphen{}zero, else it returns 0.0.  The logical OR operator will return
1.0 if either of its arguments is non\sphinxhyphen{}zero, else it returns 0.0.  The
logical XOR operator will return 1.0 if one of its arguments is zero
and the other non\sphinxhyphen{}zero, else it returns 0.0.  The logical NOT operator
returns 1.0 if its argument is 0.0, else it returns 0.0.  The 3
logical operators can only be used to operate on numbers, not on
strings.

\sphinxAtStartPar
The overall Boolean expression produces a TRUE result if the numeric
result is non\sphinxhyphen{}zero.  If the result is zero, the expression result is
FALSE.

\begin{sphinxadmonition}{note}{Note:}
\sphinxAtStartPar
If the Boolean expression is a single numeric value with no Boolean
operators, it will be FALSE if the value = 0.0, otherwise TRUE.  If
the Boolean expression is a single string, an error message will be
issued.
\end{sphinxadmonition}


\bigskip\hrule\bigskip



\subsection{Restrictions}
\label{\detokenize{if:restrictions}}
\sphinxAtStartPar
none


\subsection{Related commands}
\label{\detokenize{if:related-commands}}
\sphinxAtStartPar
{\hyperref[\detokenize{variable::doc}]{\sphinxcrossref{\DUrole{doc}{variable}}}}, {\hyperref[\detokenize{print::doc}]{\sphinxcrossref{\DUrole{doc}{print}}}}


\subsection{Default}
\label{\detokenize{if:default}}
\sphinxAtStartPar
none

\sphinxstepscope

\index{improper\_coeff@\spxentry{improper\_coeff}}\ignorespaces 

\section{improper\_coeff command}
\label{\detokenize{improper_coeff:improper-coeff-command}}\label{\detokenize{improper_coeff:index-0}}\label{\detokenize{improper_coeff::doc}}

\subsection{Syntax}
\label{\detokenize{improper_coeff:syntax}}
\begin{sphinxVerbatim}[commandchars=\\\{\}]
\PYG{k}{improper\PYGZus{}coeff}\PYG{+w}{ }\PYG{n}{N}\PYG{+w}{ }\PYG{n}{args}
\end{sphinxVerbatim}
\begin{itemize}
\item {} 
\sphinxAtStartPar
N = numeric improper type (see asterisk form below), or type label

\item {} 
\sphinxAtStartPar
args = coefficients for one or more improper types

\end{itemize}


\subsection{Examples}
\label{\detokenize{improper_coeff:examples}}
\begin{sphinxVerbatim}[commandchars=\\\{\}]
\PYG{k}{improper\PYGZus{}coeff}\PYG{+w}{ }\PYG{l+m}{1}\PYG{+w}{ }\PYG{l+m}{300.0}\PYG{+w}{ }\PYG{l+m}{0.0}
\PYG{k}{improper\PYGZus{}coeff}\PYG{+w}{ }\PYG{o}{*}\PYG{+w}{ }\PYG{l+m}{80.2}\PYG{+w}{ }\PYG{o}{\PYGZhy{}}\PYG{l+m}{1}\PYG{+w}{ }\PYG{l+m}{2}
\PYG{k}{improper\PYGZus{}coeff}\PYG{+w}{ }\PYG{o}{*}\PYG{l+m}{4}\PYG{+w}{ }\PYG{l+m}{80.2}\PYG{+w}{ }\PYG{o}{\PYGZhy{}}\PYG{l+m}{1}\PYG{+w}{ }\PYG{l+m}{2}

\PYG{k}{labelmap}\PYG{+w}{ }\PYG{n}{improper}\PYG{+w}{ }\PYG{l+m}{1}\PYG{+w}{ }\PYG{n}{benzene}
\PYG{k}{improper\PYGZus{}coeff}\PYG{+w}{ }\PYG{n}{benzene}\PYG{+w}{ }\PYG{l+m}{300.0}\PYG{+w}{ }\PYG{l+m}{0.0}
\end{sphinxVerbatim}


\subsection{Description}
\label{\detokenize{improper_coeff:description}}
\sphinxAtStartPar
Specify the improper force field coefficients for one or more improper
types.  The number and meaning of the coefficients depends on the
improper style.  Improper coefficients can also be set in the data file
read by the {\hyperref[\detokenize{read_data::doc}]{\sphinxcrossref{\DUrole{doc}{read\_data}}}} command or in a restart file.

\sphinxAtStartPar
\(N\) can be specified in one of two ways.  An explicit numeric
value can be used, as in the first example above.  Or \(N\) can be a
type label, which is an alphanumeric string defined by the
{\hyperref[\detokenize{labelmap::doc}]{\sphinxcrossref{\DUrole{doc}{labelmap}}}} command or in a section of a data file read
by the {\hyperref[\detokenize{read_data::doc}]{\sphinxcrossref{\DUrole{doc}{read\_data}}}} command.

\sphinxAtStartPar
For numeric values only, a wild\sphinxhyphen{}card asterisk can be used to set the
coefficients for multiple improper types.  This takes the form “*” or
“*n” or “n*” or “m*n”.  If \(N\) = the number of improper types,
then an asterisk with no numeric values means all types from 1 to
\(N\).  A leading asterisk means all types from 1 to n (inclusive).
A trailing asterisk means all types from n to \(N\) (inclusive).  A
middle asterisk means all types from m to n (inclusive).

\sphinxAtStartPar
Note that using an improper\_coeff command can override a previous
setting for the same improper type.  For example, these commands set the
coeffs for all improper types, then overwrite the coeffs for just
improper type 2:

\begin{sphinxVerbatim}[commandchars=\\\{\}]
\PYG{k}{improper\PYGZus{}coeff}\PYG{+w}{ }\PYG{o}{*}\PYG{+w}{ }\PYG{l+m}{300.0}\PYG{+w}{ }\PYG{l+m}{0.0}
\PYG{k}{improper\PYGZus{}coeff}\PYG{+w}{ }\PYG{l+m}{2}\PYG{+w}{ }\PYG{l+m}{50.0}\PYG{+w}{ }\PYG{l+m}{0.0}
\end{sphinxVerbatim}

\sphinxAtStartPar
A line in a data file that specifies improper coefficients uses the
exact same format as the arguments of the improper\_coeff command in an
input script, except that wild\sphinxhyphen{}card asterisks should not be used since
coefficients for all \(N\) types must be listed in the file.  For
example, under the “Improper Coeffs” section of a data file, the line
that corresponds to the first example above would be listed as

\begin{sphinxVerbatim}[commandchars=\\\{\}]
1 300.0 0.0
\end{sphinxVerbatim}

\sphinxAtStartPar
The {\hyperref[\detokenize{improper_class2::doc}]{\sphinxcrossref{\DUrole{doc}{improper\_style class2}}}} is an exception to
this rule, in that an additional argument is used in the input script
to allow specification of the cross\sphinxhyphen{}term coefficients.  See its doc
page for details.


\bigskip\hrule\bigskip


\sphinxAtStartPar
The list of all improper styles defined in LAMMPS is given on the
{\hyperref[\detokenize{improper_style::doc}]{\sphinxcrossref{\DUrole{doc}{improper\_style}}}} doc page.  They are also listed
in more compact form on the {\hyperref[\detokenize{Commands_bond:improper}]{\sphinxcrossref{\DUrole{std,std-ref}{Commands improper}}}} doc page.

\sphinxAtStartPar
On either of those pages, click on the style to display the formula it
computes and its coefficients as specified by the associated
improper\_coeff command.


\bigskip\hrule\bigskip



\subsection{Restrictions}
\label{\detokenize{improper_coeff:restrictions}}
\sphinxAtStartPar
This command must come after the simulation box is defined by a
{\hyperref[\detokenize{read_data::doc}]{\sphinxcrossref{\DUrole{doc}{read\_data}}}}, {\hyperref[\detokenize{read_restart::doc}]{\sphinxcrossref{\DUrole{doc}{read\_restart}}}}, or
{\hyperref[\detokenize{create_box::doc}]{\sphinxcrossref{\DUrole{doc}{create\_box}}}} command.

\sphinxAtStartPar
An improper style must be defined before any improper coefficients are
set, either in the input script or in a data file.


\subsection{Related commands}
\label{\detokenize{improper_coeff:related-commands}}
\sphinxAtStartPar
{\hyperref[\detokenize{improper_style::doc}]{\sphinxcrossref{\DUrole{doc}{improper\_style}}}}


\subsection{Default}
\label{\detokenize{improper_coeff:default}}
\sphinxAtStartPar
none

\sphinxstepscope

\index{improper\_style@\spxentry{improper\_style}}\ignorespaces 

\section{improper\_style command}
\label{\detokenize{improper_style:improper-style-command}}\label{\detokenize{improper_style:index-0}}\label{\detokenize{improper_style::doc}}

\subsection{Syntax}
\label{\detokenize{improper_style:syntax}}
\begin{sphinxVerbatim}[commandchars=\\\{\}]
\PYG{k}{improper\PYGZus{}style}\PYG{+w}{ }\PYG{n}{style}
\end{sphinxVerbatim}
\begin{itemize}
\item {} 
\sphinxAtStartPar
style = \sphinxstyleemphasis{none} or \sphinxstyleemphasis{hybrid} or \sphinxstyleemphasis{class2} or \sphinxstyleemphasis{cvff} or \sphinxstyleemphasis{harmonic}

\end{itemize}


\subsection{Examples}
\label{\detokenize{improper_style:examples}}
\begin{sphinxVerbatim}[commandchars=\\\{\}]
\PYG{k}{improper\PYGZus{}style}\PYG{+w}{ }\PYG{n}{harmonic}
\PYG{k}{improper\PYGZus{}style}\PYG{+w}{ }\PYG{n}{cvff}
\PYG{k}{improper\PYGZus{}style}\PYG{+w}{ }\PYG{n}{hybrid}\PYG{+w}{ }\PYG{n}{cvff}\PYG{+w}{ }\PYG{n}{harmonic}
\end{sphinxVerbatim}


\subsection{Description}
\label{\detokenize{improper_style:description}}
\sphinxAtStartPar
Set the formula(s) LAMMPS uses to compute improper interactions
between quadruplets of atoms, which remain in force for the duration
of the simulation.  The list of improper quadruplets is read in by a
{\hyperref[\detokenize{read_data::doc}]{\sphinxcrossref{\DUrole{doc}{read\_data}}}} or {\hyperref[\detokenize{read_restart::doc}]{\sphinxcrossref{\DUrole{doc}{read\_restart}}}} command
from a data or restart file.  Note that the ordering of the 4 atoms in
an improper quadruplet determines the definition of the improper
angle used in the formula for each style.  See the doc pages of
individual styles for details.

\sphinxAtStartPar
Hybrid models where impropers are computed using different improper
potentials can be setup using the \sphinxstyleemphasis{hybrid} improper style.

\sphinxAtStartPar
The coefficients associated with an improper style can be specified in
a data or restart file or via the {\hyperref[\detokenize{improper_coeff::doc}]{\sphinxcrossref{\DUrole{doc}{improper\_coeff}}}}
command.

\sphinxAtStartPar
All improper potentials store their coefficient data in binary restart
files which means improper\_style and
{\hyperref[\detokenize{improper_coeff::doc}]{\sphinxcrossref{\DUrole{doc}{improper\_coeff}}}} commands do not need to be
re\sphinxhyphen{}specified in an input script that restarts a simulation.  See the
{\hyperref[\detokenize{read_restart::doc}]{\sphinxcrossref{\DUrole{doc}{read\_restart}}}} command for details on how to do
this.  The one exception is that improper\_style \sphinxstyleemphasis{hybrid} only stores
the list of sub\sphinxhyphen{}styles in the restart file; improper coefficients need
to be re\sphinxhyphen{}specified.

\begin{sphinxadmonition}{note}{Note:}
\sphinxAtStartPar
When both an improper and pair style is defined, the
{\hyperref[\detokenize{special_bonds::doc}]{\sphinxcrossref{\DUrole{doc}{special\_bonds}}}} command often needs to be used to
turn off (or weight) the pairwise interaction that would otherwise
exist between a group of 4 bonded atoms.
\end{sphinxadmonition}


\bigskip\hrule\bigskip


\sphinxAtStartPar
Here is an alphabetic list of improper styles defined in LAMMPS.
Click on the style to display the formula it computes and coefficients
specified by the associated {\hyperref[\detokenize{improper_coeff::doc}]{\sphinxcrossref{\DUrole{doc}{improper\_coeff}}}}
command.

\sphinxAtStartPar
Click on the style to display the formula it computes, any additional
arguments specified in the improper\_style command, and coefficients
specified by the associated {\hyperref[\detokenize{improper_coeff::doc}]{\sphinxcrossref{\DUrole{doc}{improper\_coeff}}}}
command.

\sphinxAtStartPar
There are also additional accelerated pair styles included in the
LAMMPS distribution for faster performance on CPUs, GPUs, and KNLs.
The individual style names on the {\hyperref[\detokenize{Commands_bond:improper}]{\sphinxcrossref{\DUrole{std,std-ref}{Commands improper}}}} page are followed by one or
more of (g,i,k,o,t) to indicate which accelerated styles exist.
\begin{itemize}
\item {} 
\sphinxAtStartPar
{\hyperref[\detokenize{improper_none::doc}]{\sphinxcrossref{\DUrole{doc}{none}}}} \sphinxhyphen{} turn off improper interactions

\item {} 
\sphinxAtStartPar
{\hyperref[\detokenize{improper_zero::doc}]{\sphinxcrossref{\DUrole{doc}{zero}}}} \sphinxhyphen{} topology but no interactions

\item {} 
\sphinxAtStartPar
{\hyperref[\detokenize{improper_hybrid::doc}]{\sphinxcrossref{\DUrole{doc}{hybrid}}}} \sphinxhyphen{} define multiple styles of improper interactions

\item {} 
\sphinxAtStartPar
{\hyperref[\detokenize{improper_amoeba::doc}]{\sphinxcrossref{\DUrole{doc}{amoeba}}}} \sphinxhyphen{} AMOEBA out\sphinxhyphen{}of\sphinxhyphen{}plane improper

\item {} 
\sphinxAtStartPar
{\hyperref[\detokenize{improper_class2::doc}]{\sphinxcrossref{\DUrole{doc}{class2}}}} \sphinxhyphen{} COMPASS (class 2) improper

\item {} 
\sphinxAtStartPar
{\hyperref[\detokenize{improper_cossq::doc}]{\sphinxcrossref{\DUrole{doc}{cossq}}}} \sphinxhyphen{} improper with a cosine squared term

\item {} 
\sphinxAtStartPar
{\hyperref[\detokenize{improper_cvff::doc}]{\sphinxcrossref{\DUrole{doc}{cvff}}}} \sphinxhyphen{} CVFF improper

\item {} 
\sphinxAtStartPar
{\hyperref[\detokenize{improper_distance::doc}]{\sphinxcrossref{\DUrole{doc}{distance}}}} \sphinxhyphen{} improper based on distance between atom planes

\item {} 
\sphinxAtStartPar
{\hyperref[\detokenize{improper_distharm::doc}]{\sphinxcrossref{\DUrole{doc}{distharm}}}} \sphinxhyphen{} improper that is harmonic in the out\sphinxhyphen{}of\sphinxhyphen{}plane distance

\item {} 
\sphinxAtStartPar
{\hyperref[\detokenize{improper_fourier::doc}]{\sphinxcrossref{\DUrole{doc}{fourier}}}} \sphinxhyphen{} improper with multiple cosine terms

\item {} 
\sphinxAtStartPar
{\hyperref[\detokenize{improper_harmonic::doc}]{\sphinxcrossref{\DUrole{doc}{harmonic}}}} \sphinxhyphen{} harmonic improper

\item {} 
\sphinxAtStartPar
{\hyperref[\detokenize{improper_inversion_harmonic::doc}]{\sphinxcrossref{\DUrole{doc}{inversion/harmonic}}}} \sphinxhyphen{} harmonic improper with Wilson\sphinxhyphen{}Decius out\sphinxhyphen{}of\sphinxhyphen{}plane definition

\item {} 
\sphinxAtStartPar
{\hyperref[\detokenize{improper_ring::doc}]{\sphinxcrossref{\DUrole{doc}{ring}}}} \sphinxhyphen{} improper which prevents planar conformations

\item {} 
\sphinxAtStartPar
{\hyperref[\detokenize{improper_umbrella::doc}]{\sphinxcrossref{\DUrole{doc}{umbrella}}}} \sphinxhyphen{} DREIDING improper

\item {} 
\sphinxAtStartPar
{\hyperref[\detokenize{improper_sqdistharm::doc}]{\sphinxcrossref{\DUrole{doc}{sqdistharm}}}} \sphinxhyphen{} improper that is harmonic in the square of the out\sphinxhyphen{}of\sphinxhyphen{}plane distance

\end{itemize}


\bigskip\hrule\bigskip



\subsection{Restrictions}
\label{\detokenize{improper_style:restrictions}}
\sphinxAtStartPar
Improper styles can only be set for atom\_style choices that allow
impropers to be defined.

\sphinxAtStartPar
Most improper styles are part of the MOLECULE package.  They are only
enabled if LAMMPS was built with that package.  See the {\hyperref[\detokenize{Build_package::doc}]{\sphinxcrossref{\DUrole{doc}{Build package}}}} page for more info.  The doc pages for
individual improper potentials tell if it is part of a package.


\subsection{Related commands}
\label{\detokenize{improper_style:related-commands}}
\sphinxAtStartPar
{\hyperref[\detokenize{improper_coeff::doc}]{\sphinxcrossref{\DUrole{doc}{improper\_coeff}}}}


\subsection{Default}
\label{\detokenize{improper_style:default}}
\begin{sphinxVerbatim}[commandchars=\\\{\}]
\PYG{k}{improper\PYGZus{}style}\PYG{+w}{ }\PYG{n}{none}
\end{sphinxVerbatim}

\sphinxstepscope

\index{include@\spxentry{include}}\ignorespaces 

\section{include command}
\label{\detokenize{include:include-command}}\label{\detokenize{include:index-0}}\label{\detokenize{include::doc}}

\subsection{Syntax}
\label{\detokenize{include:syntax}}
\begin{sphinxVerbatim}[commandchars=\\\{\}]
\PYG{k}{include}\PYG{+w}{ }\PYG{n}{file}
\end{sphinxVerbatim}
\begin{itemize}
\item {} 
\sphinxAtStartPar
file = filename of new input script to switch to

\end{itemize}


\subsection{Examples}
\label{\detokenize{include:examples}}
\begin{sphinxVerbatim}[commandchars=\\\{\}]
\PYG{k}{include}\PYG{+w}{ }\PYG{n}{newfile}
\PYG{k}{include}\PYG{+w}{ }\PYG{n}{in.run2}
\end{sphinxVerbatim}


\subsection{Description}
\label{\detokenize{include:description}}
\sphinxAtStartPar
This command opens a new input script file and begins reading LAMMPS
commands from that file.  When the new file is finished, the original
file is returned to.  Include files can be nested as deeply as
desired.  If input script A includes script B, and B includes A, then
LAMMPS could run for a long time.

\sphinxAtStartPar
If the filename is a variable (see the {\hyperref[\detokenize{variable::doc}]{\sphinxcrossref{\DUrole{doc}{variable}}}}
command), different processor partitions can run different input
scripts.


\subsection{Restrictions}
\label{\detokenize{include:restrictions}}
\sphinxAtStartPar
none


\subsection{Related commands}
\label{\detokenize{include:related-commands}}
\sphinxAtStartPar
{\hyperref[\detokenize{variable::doc}]{\sphinxcrossref{\DUrole{doc}{variable}}}}, {\hyperref[\detokenize{jump::doc}]{\sphinxcrossref{\DUrole{doc}{jump}}}}


\subsection{Default}
\label{\detokenize{include:default}}
\sphinxAtStartPar
none

\sphinxstepscope

\index{info@\spxentry{info}}\ignorespaces 

\section{info command}
\label{\detokenize{info:info-command}}\label{\detokenize{info:index-0}}\label{\detokenize{info::doc}}

\subsection{Syntax}
\label{\detokenize{info:syntax}}
\begin{sphinxVerbatim}[commandchars=\\\{\}]
\PYG{k}{info}\PYG{+w}{ }\PYG{n}{args}
\end{sphinxVerbatim}
\begin{itemize}
\item {} 
\sphinxAtStartPar
args = one or more of the following keywords: \sphinxstyleemphasis{out}, \sphinxstyleemphasis{all}, \sphinxstyleemphasis{system}, \sphinxstyleemphasis{memory}, \sphinxstyleemphasis{communication}, \sphinxstyleemphasis{computes}, \sphinxstyleemphasis{dumps}, \sphinxstyleemphasis{fixes}, \sphinxstyleemphasis{groups}, \sphinxstyleemphasis{regions}, \sphinxstyleemphasis{variables}, \sphinxstyleemphasis{coeffs}, \sphinxstyleemphasis{styles}, \sphinxstyleemphasis{time}, \sphinxstyleemphasis{accelerator}, \sphinxstyleemphasis{fft} or \sphinxstyleemphasis{configuration}

\item {} 
\sphinxAtStartPar
\sphinxstyleemphasis{out} values = \sphinxstyleemphasis{screen}, \sphinxstyleemphasis{log}, \sphinxstyleemphasis{append} filename, \sphinxstyleemphasis{overwrite} filename

\item {} 
\sphinxAtStartPar
\sphinxstyleemphasis{styles} values = \sphinxstyleemphasis{all}, \sphinxstyleemphasis{angle}, \sphinxstyleemphasis{atom}, \sphinxstyleemphasis{bond}, \sphinxstyleemphasis{compute}, \sphinxstyleemphasis{command}, \sphinxstyleemphasis{dump}, \sphinxstyleemphasis{dihedral}, \sphinxstyleemphasis{fix}, \sphinxstyleemphasis{improper}, \sphinxstyleemphasis{integrate}, \sphinxstyleemphasis{kspace}, \sphinxstyleemphasis{minimize}, \sphinxstyleemphasis{pair}, \sphinxstyleemphasis{region}

\end{itemize}


\subsection{Examples}
\label{\detokenize{info:examples}}
\begin{sphinxVerbatim}[commandchars=\\\{\}]
\PYG{k}{info}\PYG{+w}{ }\PYG{n}{system}
\PYG{k}{info}\PYG{+w}{ }\PYG{n}{groups}\PYG{+w}{ }\PYG{n}{computes}\PYG{+w}{ }\PYG{n}{variables}
\PYG{k}{info}\PYG{+w}{ }\PYG{n}{all}\PYG{+w}{ }\PYG{n}{out}\PYG{+w}{ }\PYG{n}{log}
\PYG{k}{info}\PYG{+w}{ }\PYG{n}{all}\PYG{+w}{ }\PYG{n}{out}\PYG{+w}{ }\PYG{n}{append}\PYG{+w}{ }\PYG{n}{info.txt}
\PYG{k}{info}\PYG{+w}{ }\PYG{n}{styles}\PYG{+w}{ }\PYG{n}{all}
\PYG{k}{info}\PYG{+w}{ }\PYG{n}{styles}\PYG{+w}{ }\PYG{n}{atom}\PYG{+w}{ }\PYG{n}{styles}\PYG{+w}{ }\PYG{n}{command}
\end{sphinxVerbatim}


\subsection{Description}
\label{\detokenize{info:description}}
\sphinxAtStartPar
Print out information about the current internal state of the running
LAMMPS process. This can be helpful when debugging or validating complex
input scripts.  Several output categories are available and one or more
output categories may be requested.  All category keywords take no
arguments, only \sphinxstyleemphasis{out} and \sphinxstyleemphasis{styles} take arguments as shown below.  The
keywords are cumulative, may be abbreviated, and unknown keywords are
ignored.

\sphinxAtStartPar
The \sphinxstyleemphasis{out} flag controls where the output is sent. It can only be sent
to one target. By default this is the screen, if it is active. The
\sphinxstyleemphasis{log} argument selects the log file instead. With the \sphinxstyleemphasis{append} and
\sphinxstyleemphasis{overwrite} option, followed by a filename, the output is written
to that file, which is either appended to or overwritten, respectively.

\sphinxAtStartPar
The \sphinxstyleemphasis{all} flag activates printing all categories listed below.

\sphinxAtStartPar
The \sphinxstyleemphasis{configuration} category prints some information about the
LAMMPS version as well as architecture and OS it is run on.

\sphinxAtStartPar
The \sphinxstyleemphasis{memory} category prints some information about the current
memory allocation of MPI rank 0 (this the amount of dynamically
allocated memory reported by LAMMPS classes). Where supported,
also some OS specific information about the size of the reserved
memory pool size (this is where malloc() and the new operator
request memory from) and the maximum resident set size is reported
(this is the maximum amount of physical memory occupied so far).

\sphinxAtStartPar
The \sphinxstyleemphasis{system} category prints a general system overview listing.  This
includes the unit style, atom style, number of atoms, bonds, angles,
dihedrals, and impropers and the number of the respective types, box
dimensions and properties, force computing styles and more.

\sphinxAtStartPar
The \sphinxstyleemphasis{communication} category prints a variety of information about
communication and parallelization: the MPI library version level,
the number of MPI ranks and OpenMP threads, the communication style
and layout, the processor grid dimensions, ghost atom communication
mode, cutoff, and related settings.

\sphinxAtStartPar
The \sphinxstyleemphasis{computes} category prints a list of all currently defined
computes, their IDs and styles and groups they operate on.

\sphinxAtStartPar
The \sphinxstyleemphasis{dumps} category prints a list of all currently active dumps,
their IDs, styles, filenames, groups, and dump frequencies.

\sphinxAtStartPar
The \sphinxstyleemphasis{fixes} category prints a list of all currently defined fixes,
their IDs and styles and groups they operate on.

\sphinxAtStartPar
The \sphinxstyleemphasis{groups} category prints a list of all currently defined groups.

\sphinxAtStartPar
The \sphinxstyleemphasis{regions} category prints a list of all currently defined regions,
their IDs and styles and whether “inside” or “outside” atoms are
selected.

\sphinxAtStartPar
The \sphinxstyleemphasis{variables} category prints a list of all currently defined
variables, their names, styles, definition and last computed value, if
available.

\sphinxAtStartPar
The \sphinxstyleemphasis{coeffs} category prints a list for each defined force style
(pair, bond, angle, dihedral, improper) indicating which of the
corresponding coefficients have been set. This can be very helpful
to debug error messages like “All pair coeffs are not set”.

\sphinxAtStartPar
The \sphinxstyleemphasis{accelerator} category prints out information about compile time
settings of included accelerator support for the GPU, KOKKOS, INTEL,
and OPENMP packages.

\sphinxAtStartPar
\DUrole{versionmodified,added}{Added in version 7Feb2024.}

\sphinxAtStartPar
The \sphinxstyleemphasis{fft} category prints out information about the included 3d\sphinxhyphen{}FFT
support.  This lists the 3d\sphinxhyphen{}FFT engine, FFT precision, FFT library
used by the FFT engine. If the KOKKOS package is included, the settings
used for the KOKKOS package are displayed as well.

\sphinxAtStartPar
The \sphinxstyleemphasis{styles} category prints the list of styles available in the current
LAMMPS binary. The \sphinxstyleemphasis{styles} keyword without option is the same as using
the “all” option.  One of the following options may be used to control
which category of styles is printed out.  To select multiple categories,
the styles keyword needs to be used multiple times with the desired
categories:
\begin{itemize}
\item {} 
\sphinxAtStartPar
all

\item {} 
\sphinxAtStartPar
angle

\item {} 
\sphinxAtStartPar
atom

\item {} 
\sphinxAtStartPar
bond

\item {} 
\sphinxAtStartPar
compute

\item {} 
\sphinxAtStartPar
command

\item {} 
\sphinxAtStartPar
dump

\item {} 
\sphinxAtStartPar
dihedral

\item {} 
\sphinxAtStartPar
fix

\item {} 
\sphinxAtStartPar
improper

\item {} 
\sphinxAtStartPar
integrate

\item {} 
\sphinxAtStartPar
kspace

\item {} 
\sphinxAtStartPar
minimize

\item {} 
\sphinxAtStartPar
pair

\item {} 
\sphinxAtStartPar
region

\end{itemize}

\sphinxAtStartPar
The \sphinxstyleemphasis{time} category prints the accumulated CPU and wall time for the
process that writes output (usually MPI rank 0).


\subsection{Restrictions}
\label{\detokenize{info:restrictions}}
\sphinxAtStartPar
none


\subsection{Related commands}
\label{\detokenize{info:related-commands}}
\sphinxAtStartPar
{\hyperref[\detokenize{print::doc}]{\sphinxcrossref{\DUrole{doc}{print}}}}


\subsection{Default}
\label{\detokenize{info:default}}
\sphinxAtStartPar
The \sphinxstyleemphasis{out} option has the default \sphinxstyleemphasis{screen}.

\sphinxAtStartPar
The \sphinxstyleemphasis{styles} option has the default \sphinxstyleemphasis{all}.

\sphinxstepscope

\index{jump@\spxentry{jump}}\ignorespaces 

\section{jump command}
\label{\detokenize{jump:jump-command}}\label{\detokenize{jump:index-0}}\label{\detokenize{jump::doc}}

\subsection{Syntax}
\label{\detokenize{jump:syntax}}
\begin{sphinxVerbatim}[commandchars=\\\{\}]
\PYG{k}{jump }\PYG{l+s+sc}{file}\PYG{+w}{ }\PYG{k}{label}
\end{sphinxVerbatim}
\begin{itemize}
\item {} 
\sphinxAtStartPar
file = filename of new input script to switch to

\item {} 
\sphinxAtStartPar
label = optional label within file to jump to

\end{itemize}


\subsection{Examples}
\label{\detokenize{jump:examples}}
\begin{sphinxVerbatim}[commandchars=\\\{\}]
\PYG{k}{jump }\PYG{l+s+sc}{newfile}
\PYG{k}{jump }\PYG{l+s+sc}{in.run2}\PYG{+w}{ }\PYG{n}{runloop}
\PYG{k}{jump }\PYG{l+s+sc}{SELF}\PYG{+w}{ }\PYG{n}{runloop}
\end{sphinxVerbatim}


\subsection{Description}
\label{\detokenize{jump:description}}
\sphinxAtStartPar
This command closes the current input script file, opens the file with
the specified name, and begins reading LAMMPS commands from that file.
Unlike the {\hyperref[\detokenize{include::doc}]{\sphinxcrossref{\DUrole{doc}{include}}}} command, the original file is not
returned to, although by using multiple jump commands it is possible
to chain from file to file or back to the original file.

\sphinxAtStartPar
If the word “SELF” is used for the filename, then the current input
script is re\sphinxhyphen{}opened and read again.

\begin{sphinxadmonition}{note}{Note:}
\sphinxAtStartPar
The SELF option is not guaranteed to work when the current input
script is being read through stdin (standard input), e.g.
\end{sphinxadmonition}

\begin{sphinxVerbatim}[commandchars=\\\{\}]
lmp\PYGZus{}g++\PYG{+w}{ }\PYGZlt{}\PYG{+w}{ }\PYG{k}{in}.script
\end{sphinxVerbatim}

\sphinxAtStartPar
since the SELF option invokes the C\sphinxhyphen{}library rewind() call, which may
not be supported for stdin on some systems or by some MPI
implementations.  This can be worked around by using the {\hyperref[\detokenize{Run_options::doc}]{\sphinxcrossref{\DUrole{doc}{\sphinxhyphen{}in command\sphinxhyphen{}line switch}}}}, e.g.

\begin{sphinxVerbatim}[commandchars=\\\{\}]
lmp\PYGZus{}g++\PYG{+w}{ }\PYGZhy{}in\PYG{+w}{ }\PYG{k}{in}.script
\end{sphinxVerbatim}

\sphinxAtStartPar
or by using the {\hyperref[\detokenize{Run_options::doc}]{\sphinxcrossref{\DUrole{doc}{\sphinxhyphen{}var command\sphinxhyphen{}line switch}}}} to pass
the script name as a variable to the input script.  In the latter
case, a {\hyperref[\detokenize{variable::doc}]{\sphinxcrossref{\DUrole{doc}{variable}}}} called “fname” could be used in place
of SELF, e.g.

\begin{sphinxVerbatim}[commandchars=\\\{\}]
lmp\PYGZus{}g++\PYG{+w}{ }\PYGZhy{}var\PYG{+w}{ }fname\PYG{+w}{ }\PYG{k}{in}.script\PYG{+w}{ }\PYGZlt{}\PYG{+w}{ }\PYG{k}{in}.script
\end{sphinxVerbatim}

\sphinxAtStartPar
The second argument to the jump command is optional.  If specified, it is
treated as a label and the new file is scanned (without executing
commands) until the label is found, and commands are executed from
that point forward.  This can be used to loop over a portion of the
input script, as in this example.  These commands perform 10 runs,
each of 10000 steps, and create 10 dump files named file.1, file.2,
etc.  The {\hyperref[\detokenize{next::doc}]{\sphinxcrossref{\DUrole{doc}{next}}}} command is used to exit the loop after 10
iterations.  When the “a” variable has been incremented for the tenth
time, it will cause the next jump command to be skipped.

\begin{sphinxVerbatim}[commandchars=\\\{\}]
\PYG{k}{variable }\PYG{n+nv+nvIdentifier}{a}\PYG{+w}{ }\PYG{n}{loop}\PYG{+w}{ }\PYG{l+m}{10}
\PYG{k}{label }\PYG{l+s+sc}{loop}
\PYG{k}{dump }\PYG{n+nv+nvIdentifier}{1}\PYG{+w}{ }\PYG{n+nv+nvIdentifier}{all}\PYG{+w}{ }\PYG{n}{atom}\PYG{+w}{ }\PYG{l+m}{100}\PYG{+w}{ }\PYG{n}{file.}\PYG{n+nv}{\PYGZdl{}a}
\PYG{k}{run}\PYG{+w}{ }\PYG{l+m}{10000}
\PYG{k}{undump }\PYG{n+nv+nvIdentifier}{1}
\PYG{k}{next }\PYG{n+nv+nvIdentifier}{a}
\PYG{k}{jump }\PYG{l+s+sc}{in.lj}\PYG{+w}{ }\PYG{n}{loop}
\end{sphinxVerbatim}

\sphinxAtStartPar
If the jump \sphinxstyleemphasis{file} argument is a variable, the jump command can be
used to cause different processor partitions to run different input
scripts.  In this example, LAMMPS is run on 40 processors, with 4
partitions of 10 procs each.  An in.file containing the example
variable and jump command will cause each partition to run a different
simulation.

\begin{sphinxVerbatim}[commandchars=\\\{\}]
\PYG{n}{mpirun}\PYG{+w}{ }\PYG{o}{\PYGZhy{}}\PYG{n}{np}\PYG{+w}{ }\PYG{l+m}{40}\PYG{+w}{ }\PYG{n}{lmp\PYGZus{}ibm}\PYG{+w}{ }\PYG{o}{\PYGZhy{}}\PYG{n}{partition}\PYG{+w}{ }\PYG{l+m}{4}\PYG{n}{x10}\PYG{+w}{ }\PYG{o}{\PYGZhy{}}\PYG{n}{in}\PYG{+w}{ }\PYG{n}{in.file}

\PYG{k}{variable }\PYG{n+nv+nvIdentifier}{f}\PYG{+w}{ }\PYG{n}{world}\PYG{+w}{ }\PYG{n}{script.1}\PYG{+w}{ }\PYG{n}{script.2}\PYG{+w}{ }\PYG{n}{script.3}\PYG{+w}{ }\PYG{n}{script.4}
\PYG{k}{jump }\PYG{n+nv}{\PYGZdl{}f}
\end{sphinxVerbatim}

\sphinxAtStartPar
Here is an example of a loop which checks every 1000 steps if the
system temperature has reached a certain value, and if so, breaks out
of the loop to finish the run.  Note that any variable could be
checked, so long as it is current on the timestep when the run
completes.  As explained on the {\hyperref[\detokenize{variable::doc}]{\sphinxcrossref{\DUrole{doc}{variable}}}} doc page,
this can be ensured by including the variable in thermodynamic output.

\begin{sphinxVerbatim}[commandchars=\\\{\}]
\PYG{k}{variable }\PYG{n+nv+nvIdentifier}{myTemp}\PYG{+w}{ }\PYG{n}{equal}\PYG{+w}{ }\PYG{n}{temp}
\PYG{k}{label }\PYG{l+s+sc}{loop}
\PYG{k}{variable }\PYG{n+nv+nvIdentifier}{a}\PYG{+w}{ }\PYG{n}{loop}\PYG{+w}{ }\PYG{l+m}{1000}
\PYG{k}{run}\PYG{+w}{ }\PYG{l+m}{1000}
\PYG{k}{if}\PYG{+w}{ }\PYG{l+s}{\PYGZdq{}}\PYG{l+s}{\PYGZdl{}\PYGZob{}myTemp\PYGZcb{} \PYGZlt{} 300.0}\PYG{l+s}{\PYGZdq{}}\PYG{+w}{ }\PYG{k}{then}\PYG{+w}{ }\PYG{l+s}{\PYGZdq{}}\PYG{l+s}{jump SELF break}\PYG{l+s}{\PYGZdq{}}
\PYG{k}{next }\PYG{n+nv+nvIdentifier}{a}
\PYG{k}{jump }\PYG{l+s+sc}{SELF}\PYG{+w}{ }\PYG{n}{loop}
\PYG{k}{label }\PYG{l+s+sc}{break}
\PYG{k}{print}\PYG{+w}{ }\PYG{l+s}{\PYGZdq{}}\PYG{l+s}{ALL DONE}\PYG{l+s}{\PYGZdq{}}
\end{sphinxVerbatim}

\sphinxAtStartPar
Here is an example of a double loop which uses the if and
{\hyperref[\detokenize{jump::doc}]{\sphinxcrossref{\DUrole{doc}{jump}}}} commands to break out of the inner loop when a
condition is met, then continues iterating through the outer loop.

\begin{sphinxVerbatim}[commandchars=\\\{\}]
\PYG{k}{label       }\PYG{l+s+sc}{loopa}
\PYG{k}{variable    }\PYG{n+nv+nvIdentifier}{a}\PYG{+w}{ }\PYG{n}{loop}\PYG{+w}{ }\PYG{l+m}{5}
\PYG{+w}{  }\PYG{k}{label     }\PYG{l+s+sc}{loopb}
\PYG{+w}{  }\PYG{k}{variable  }\PYG{n+nv+nvIdentifier}{b}\PYG{+w}{ }\PYG{n}{loop}\PYG{+w}{ }\PYG{l+m}{5}
\PYG{+w}{    }\PYG{n}{print}\PYG{+w}{     }\PYG{l+s}{\PYGZdq{}}\PYG{l+s}{A,B = \PYGZdl{}a,\PYGZdl{}b}\PYG{l+s}{\PYGZdq{}}
\PYG{+w}{    }\PYG{n}{run}\PYG{+w}{       }\PYG{l+m}{10000}
\PYG{+w}{    }\PYG{k}{if}\PYG{+w}{        }\PYG{l+s}{\PYGZdq{}}\PYG{l+s}{\PYGZdl{}b \PYGZgt{} 2}\PYG{l+s}{\PYGZdq{}}\PYG{+w}{ }\PYG{k}{then}\PYG{+w}{ }\PYG{l+s}{\PYGZdq{}}\PYG{l+s}{jump SELF break}\PYG{l+s}{\PYGZdq{}}
\PYG{+w}{  }\PYG{k}{next      }\PYG{n+nv+nvIdentifier}{b}
\PYG{+w}{  }\PYG{k}{jump      }\PYG{l+s+sc}{in.script}\PYG{+w}{ }\PYG{n}{loopb}
\PYG{+w}{  }\PYG{k}{label       }\PYG{l+s+sc}{break}
\PYG{+w}{  }\PYG{k}{variable    }\PYG{n+nv+nvIdentifier}{b}\PYG{+w}{ }\PYG{n}{delete}
\PYG{k}{next        }\PYG{n+nv+nvIdentifier}{a}
\PYG{k}{jump        }\PYG{l+s+sc}{SELF}\PYG{+w}{ }\PYG{n}{loopa}
\end{sphinxVerbatim}


\subsection{Restrictions}
\label{\detokenize{jump:restrictions}}
\sphinxAtStartPar
If you jump to a file and it does not contain the specified label,
LAMMPS will come to the end of the file and exit.


\subsection{Related commands}
\label{\detokenize{jump:related-commands}}
\sphinxAtStartPar
{\hyperref[\detokenize{variable::doc}]{\sphinxcrossref{\DUrole{doc}{variable}}}}, {\hyperref[\detokenize{include::doc}]{\sphinxcrossref{\DUrole{doc}{include}}}}, {\hyperref[\detokenize{label::doc}]{\sphinxcrossref{\DUrole{doc}{label}}}},
{\hyperref[\detokenize{next::doc}]{\sphinxcrossref{\DUrole{doc}{next}}}}


\subsection{Default}
\label{\detokenize{jump:default}}
\sphinxAtStartPar
none

\sphinxstepscope

\index{kim\_commands@\spxentry{kim\_commands}}\ignorespaces 

\section{kim command}
\label{\detokenize{kim_commands:kim-command}}\label{\detokenize{kim_commands:index-0}}\label{\detokenize{kim_commands::doc}}

\subsection{Syntax}
\label{\detokenize{kim_commands:syntax}}
\begin{sphinxVerbatim}[commandchars=\\\{\}]
\PYG{k}{kim }\PYG{l+s+s1}{sub\PYGZhy{}command}\PYG{+w}{ }\PYG{n}{args}
\end{sphinxVerbatim}
\begin{itemize}
\item {} 
\sphinxAtStartPar
sub\sphinxhyphen{}command = {\hyperref[\detokenize{kim_commands:kim-init}]{\sphinxcrossref{\DUrole{std,std-ref}{init}}}} or {\hyperref[\detokenize{kim_commands:interactions}]{\sphinxcrossref{\DUrole{std,std-ref}{interactions}}}} or
{\hyperref[\detokenize{kim_commands:query}]{\sphinxcrossref{\DUrole{std,std-ref}{query}}}} or {\hyperref[\detokenize{kim_commands:param}]{\sphinxcrossref{\DUrole{std,std-ref}{param}}}} or {\hyperref[\detokenize{kim_commands:property}]{\sphinxcrossref{\DUrole{std,std-ref}{property}}}}

\item {} 
\sphinxAtStartPar
args = arguments used by a particular sub\sphinxhyphen{}command

\end{itemize}


\subsection{Examples}
\label{\detokenize{kim_commands:examples}}
\begin{sphinxVerbatim}[commandchars=\\\{\}]
\PYG{k}{kim }\PYG{l+s+s1}{init}\PYG{+w}{ }\PYG{n}{args}
\PYG{k}{kim }\PYG{l+s+s1}{interactions}\PYG{+w}{ }\PYG{n}{args}
\PYG{k}{kim }\PYG{l+s+s1}{query}\PYG{+w}{ }\PYG{n}{args}
\PYG{k}{kim }\PYG{l+s+s1}{param}\PYG{+w}{ }\PYG{n}{args}
\PYG{k}{kim }\PYG{l+s+s1}{property}\PYG{+w}{ }\PYG{n}{args}
\end{sphinxVerbatim}


\subsection{Description}
\label{\detokenize{kim_commands:description}}
\sphinxAtStartPar
The \sphinxstyleemphasis{kim command} includes a set of sub\sphinxhyphen{}commands that allow LAMMPS
users to use interatomic models (IM) (potentials and force fields) and
their predictions for various physical properties archived in the
\sphinxhref{https://openkim.org}{Open Knowledgebase of Interatomic Models (OpenKIM)} repository.

\sphinxAtStartPar
Using OpenKIM provides LAMMPS users with immediate access to a large
number of verified IMs and their predictions. OpenKIM IMs have
multiple benefits including \sphinxhref{https://openkim.org/doc/overview/kim-features/}{reliability, reproducibility and
convenience}.

\phantomsection\label{\detokenize{kim_commands:im-types}}
\sphinxAtStartPar
There are two types of IMs archived in OpenKIM:
\begin{enumerate}
\sphinxsetlistlabels{\arabic}{enumi}{enumii}{}{.}%
\item {} 
\sphinxAtStartPar
The first type is called a \sphinxstyleemphasis{KIM Portable Model} (PM).  A KIM PM is an
independent computer implementation of an IM written in one of the languages
supported by KIM (C, C++, Fortran) that conforms to the KIM Application
Programming Interface (\sphinxhref{https://openkim.org/kim-api/}{KIM API}) Portable
Model Interface (PMI) standard.  A KIM PM will work seamlessly with any
simulation code that supports the KIM API/PMI standard (including LAMMPS; see
\sphinxhref{https://openkim.org/projects-using-kim/}{complete list of supported codes}).

\item {} 
\sphinxAtStartPar
The second type is called a \sphinxstyleemphasis{KIM Simulator Model} (SM).  A KIM SM is an IM
that is implemented natively within a simulation code (\sphinxstyleemphasis{simulator}) that
supports the KIM API Simulator Model Interface (SMI); in this case LAMMPS.  A
separate SM package is archived in OpenKIM for each parameterization of the
IM, which includes all of the necessary parameter files, LAMMPS commands, and
metadata (supported species, units, etc.) needed to run the IM in LAMMPS.

\end{enumerate}

\sphinxAtStartPar
With these two IM types, OpenKIM can archive and test almost all IMs that can be
used by LAMMPS. (It is easy to contribute new IMs to OpenKIM, see the
\sphinxhref{https://openkim.org/doc/repository/adding-content/}{upload instructions}.)

\sphinxAtStartPar
OpenKIM IMs are uniquely identified by a
\sphinxhref{https://openkim.org/doc/schema/kim-ids/}{KIM ID}.
The extended KIM ID consists of
a human\sphinxhyphen{}readable prefix identifying the type of IM, authors, publication year,
and supported species, separated by two underscores from the KIM ID itself,
which begins with an IM code
(\sphinxstyleemphasis{MO} for a KIM Portable Model, and \sphinxstyleemphasis{SM} for a KIM Simulator Model)
followed by a unique 12\sphinxhyphen{}digit code and a 3\sphinxhyphen{}digit version identifier.
By convention SM prefixes begin with \sphinxstyleemphasis{Sim\_} to readily identify them.

\begin{sphinxVerbatim}[commandchars=\\\{\}]
SW\PYGZus{}StillingerWeber\PYGZus{}1985\PYGZus{}Si\PYGZus{}\PYGZus{}MO\PYGZus{}405512056662\PYGZus{}005
Sim\PYGZus{}LAMMPS\PYGZus{}ReaxFF\PYGZus{}StrachanVanDuinChakraborty\PYGZus{}2003\PYGZus{}CHNO\PYGZus{}\PYGZus{}SM\PYGZus{}107643900657\PYGZus{}001
\end{sphinxVerbatim}

\sphinxAtStartPar
Each OpenKIM IM has a dedicated “Model Page” on \sphinxhref{https://openkim.org}{OpenKIM}
providing all the information on the IM including a title, description,
authorship and citation information, test and verification check results,
visualizations of results, a wiki with documentation and user comments, and
access to raw files, and other information.
The URL for the Model Page is constructed from the
\sphinxhref{https://openkim.org/doc/schema/kim-ids/}{extended KIM ID} of the IM:

\begin{sphinxVerbatim}[commandchars=\\\{\}]
https://openkim.org/id/extended\PYGZus{}KIM\PYGZus{}ID
\end{sphinxVerbatim}

\sphinxAtStartPar
For example, for the Stillinger\sphinxhyphen{}Weber potential listed above the Model Page is
located at:
\begin{sphinxalltt}
\sphinxurl{https://openkim.org/id/SW\_StillingerWeber\_1985\_Si\_\_MO\_405512056662\_005}
\end{sphinxalltt}

\sphinxAtStartPar
See the
\sphinxhref{https://openkim.org/browse/models/by-species}{current list of KIM PMs and SMs archived in OpenKIM}.
This list is sorted by species and can be filtered to display only IMs for
certain species combinations.

\sphinxAtStartPar
See \sphinxhref{https://openkim.org/doc/usage/obtaining-models}{Obtaining KIM Models} to
learn how to install a pre\sphinxhyphen{}built binary of the OpenKIM Repository of Models.

\begin{sphinxadmonition}{note}{Note:}
\sphinxAtStartPar
It is also possible to locally install IMs not archived in OpenKIM,
in which case their names do not have to conform to the KIM ID format.
\end{sphinxadmonition}


\subsection{Using OpenKIM IMs with LAMMPS (\sphinxstyleemphasis{kim init}, \sphinxstyleemphasis{kim interactions})}
\label{\detokenize{kim_commands:using-openkim-ims-with-lammps-kim-init-kim-interactions}}
\sphinxAtStartPar
Two sub\sphinxhyphen{}commands are employed when using OpenKIM IMs in LAMMPS, one to select
the IM and perform necessary initialization (\sphinxstyleemphasis{kim init}), and the second to
set up the IM for use by executing any necessary LAMMPS commands
(\sphinxstyleemphasis{kim interactions}).  Both are required.


\subsubsection{Syntax}
\label{\detokenize{kim_commands:id1}}
\begin{sphinxVerbatim}[commandchars=\\\{\}]
\PYG{k}{kim }\PYG{l+s+s1}{init}\PYG{+w}{ }\PYG{n}{model}\PYG{+w}{ }\PYG{n}{user\PYGZus{}units}\PYG{+w}{ }\PYG{n}{unitarg}
\PYG{k}{kim }\PYG{l+s+s1}{interactions}\PYG{+w}{ }\PYG{n}{typeargs}
\end{sphinxVerbatim}
\phantomsection\label{\detokenize{kim_commands:typeargs-options}}\begin{itemize}
\item {} 
\sphinxAtStartPar
model = name of the KIM interatomic model (the KIM ID for models archived in
OpenKIM)

\item {} 
\sphinxAtStartPar
user\_units = the LAMMPS {\hyperref[\detokenize{units::doc}]{\sphinxcrossref{\DUrole{doc}{units}}}} style assumed in the LAMMPS
input script

\item {} 
\sphinxAtStartPar
unitarg = \sphinxstyleemphasis{unit\_conversion\_mode} (optional)

\item {} 
\sphinxAtStartPar
typeargs = atom type to species mapping (one entry per atom type) or
\sphinxstyleemphasis{fixed\_types} for models with a preset fixed mapping

\end{itemize}


\subsubsection{Examples}
\label{\detokenize{kim_commands:id2}}
\begin{sphinxVerbatim}[commandchars=\\\{\}]
\PYG{k}{kim }\PYG{l+s+s1}{init}\PYG{+w}{ }\PYG{n}{SW\PYGZus{}StillingerWeber\PYGZus{}1985\PYGZus{}Si\PYGZus{}\PYGZus{}MO\PYGZus{}405512056662\PYGZus{}005}\PYG{+w}{ }\PYG{n}{metal}
\PYG{k}{kim }\PYG{l+s+s1}{interactions}\PYG{+w}{ }\PYG{n}{Si}

\PYG{k}{kim }\PYG{l+s+s1}{init}\PYG{+w}{ }\PYG{n}{Sim\PYGZus{}LAMMPS\PYGZus{}ReaxFF\PYGZus{}StrachanVanDuinChakraborty\PYGZus{}2003\PYGZus{}CHNO\PYGZus{}\PYGZus{}SM\PYGZus{}107643900657\PYGZus{}001}\PYG{+w}{ }\PYG{n}{real}
\PYG{k}{kim }\PYG{l+s+s1}{init}\PYG{+w}{ }\PYG{n}{Sim\PYGZus{}LAMMPS\PYGZus{}ReaxFF\PYGZus{}StrachanVanDuinChakraborty\PYGZus{}2003\PYGZus{}CHNO\PYGZus{}\PYGZus{}SM\PYGZus{}107643900657\PYGZus{}001}\PYG{+w}{ }\PYG{n}{metal}\PYG{+w}{ }\PYG{n}{unit\PYGZus{}conversion\PYGZus{}mode}
\PYG{k}{kim }\PYG{l+s+s1}{interactions}\PYG{+w}{ }\PYG{n}{C}\PYG{+w}{ }\PYG{n}{H}\PYG{+w}{ }\PYG{n}{O}

\PYG{k}{kim }\PYG{l+s+s1}{init}\PYG{+w}{ }\PYG{n}{Sim\PYGZus{}LAMMPS\PYGZus{}IFF\PYGZus{}PCFF\PYGZus{}HeinzMishraLinEmami\PYGZus{}2015Ver1v5\PYGZus{}FccmetalsMineralsSolventsPolymers\PYGZus{}\PYGZus{}SM\PYGZus{}039297821658\PYGZus{}000}\PYG{+w}{ }\PYG{n}{real}
\PYG{k}{kim }\PYG{l+s+s1}{interactions}\PYG{+w}{ }\PYG{n}{fixed\PYGZus{}types}
\end{sphinxVerbatim}

\sphinxAtStartPar
See the \sphinxstyleemphasis{examples/kim} directory for example input scripts that use KIM PMs
and KIM SMs.


\subsubsection{OpenKIM IM Initialization (\sphinxstyleemphasis{kim init})}
\label{\detokenize{kim_commands:openkim-im-initialization-kim-init}}\label{\detokenize{kim_commands:kim-init}}
\sphinxAtStartPar
The \sphinxstyleemphasis{kim} command followed by \sphinxstyleemphasis{init} sub\sphinxhyphen{}command must be issued \sphinxstylestrong{before}
the simulation box is created (normally at the top of the file).  This command
sets the OpenKIM IM that will be used and may issue additional commands changing
LAMMPS default settings that are required for using the selected IM (such as
{\hyperref[\detokenize{units::doc}]{\sphinxcrossref{\DUrole{doc}{units}}}} or {\hyperref[\detokenize{atom_style::doc}]{\sphinxcrossref{\DUrole{doc}{atom\_style}}}}).  If needed, those
settings can be overridden, however, typically a script containing a \sphinxstyleemphasis{kim init}
command would not include \sphinxstyleemphasis{units} and \sphinxstyleemphasis{atom\_style} commands.

\sphinxAtStartPar
The required arguments of \sphinxstyleemphasis{kim init} are the \sphinxstyleemphasis{model} name of the IM to be used
in the simulation (for an IM archived in OpenKIM this is its
\sphinxhref{https://openkim.org/doc/schema/kim-ids/}{extended KIM ID}, and the
\sphinxstyleemphasis{user\_units}, which are the LAMMPS {\hyperref[\detokenize{units::doc}]{\sphinxcrossref{\DUrole{doc}{units style}}}} used in the input
script. (Any dimensioned numerical values in the input script and values read in
from files are expected to be in the \sphinxstyleemphasis{user\_units} system.)

\sphinxAtStartPar
The selected IM can be either a {\hyperref[\detokenize{kim_commands:im-types}]{\sphinxcrossref{\DUrole{std,std-ref}{KIM PM or a KIM SM}}}}.  For a KIM
SM, the \sphinxstyleemphasis{kim init} command verifies that the SM is designed to work with LAMMPS
(and not another simulation code).  In addition, the LAMMPS version used for
defining the SM and the LAMMPS version being currently run are printed to help
diagnose any incompatible changes to input script or command syntax between the
two LAMMPS versions.

\sphinxAtStartPar
Based on the selected model \sphinxstyleemphasis{kim init} may modify the
{\hyperref[\detokenize{atom_style::doc}]{\sphinxcrossref{\DUrole{doc}{atom\_style}}}}.  Some SMs have requirements for this setting.
If this is the case, then \sphinxstyleemphasis{atom\_style} will be set to the required style.
Otherwise, the value is left unchanged (which in the absence of an \sphinxstyleemphasis{atom\_style}
command in the input script is the
{\hyperref[\detokenize{atom_style::doc}]{\sphinxcrossref{\DUrole{doc}{default atom\_style value}}}}).

\sphinxAtStartPar
Regarding units, the \sphinxstyleemphasis{kim init} behaves in different ways depending on whether
or not \sphinxstyleemphasis{unit conversion mode} is activated as indicated by the optional
\sphinxstyleemphasis{unitarg} argument.
If unit conversion mode is \sphinxstylestrong{not} active, then \sphinxstyleemphasis{user\_units} must either match
the required units of the IM or the IM must be able to adjust its units to
match. (The latter is only possible with some KIM PMs; SMs can never adjust
their units.) If a match is possible, the LAMMPS {\hyperref[\detokenize{units::doc}]{\sphinxcrossref{\DUrole{doc}{units}}}} command is
called to set the units to \sphinxstyleemphasis{user\_units}.  If the match fails, the simulation
is terminated with an error.  The \sphinxstyleemphasis{kim init} command also sets the
default value for the {\hyperref[\detokenize{neighbor::doc}]{\sphinxcrossref{\DUrole{doc}{skin}}}} (extra distance beyond force
cutoff) as 2.0 Angstroms and sets the default value for the
{\hyperref[\detokenize{timestep::doc}]{\sphinxcrossref{\DUrole{doc}{timestep}}}} size as 1.0 femtosecond.

\sphinxAtStartPar
Here is an example of a LAMMPS script to compute the cohesive energy of a
face\sphinxhyphen{}centered cubic (fcc) lattice for the MEAM potential by Pascuet and
Fernandez (2015) for Al.

\begin{sphinxVerbatim}[commandchars=\\\{\}]
\PYG{k}{kim          }\PYG{l+s+s1}{init}\PYG{+w}{ }\PYG{n}{Sim\PYGZus{}LAMMPS\PYGZus{}MEAM\PYGZus{}PascuetFernandez\PYGZus{}2015\PYGZus{}Al\PYGZus{}\PYGZus{}SM\PYGZus{}811588957187\PYGZus{}000}\PYG{+w}{ }\PYG{n}{metal}
\PYG{k}{boundary}\PYG{+w}{     }\PYG{n}{p}\PYG{+w}{ }\PYG{n}{p}\PYG{+w}{ }\PYG{n}{p}
\PYG{k}{lattice}\PYG{+w}{      }\PYG{n}{fcc}\PYG{+w}{ }\PYG{l+m}{4.049}
\PYG{k}{region       }\PYG{n+nv+nvIdentifier}{simbox}\PYG{+w}{ }\PYG{n}{block}\PYG{+w}{ }\PYG{l+m}{0}\PYG{+w}{ }\PYG{l+m}{1}\PYG{+w}{ }\PYG{l+m}{0}\PYG{+w}{ }\PYG{l+m}{1}\PYG{+w}{ }\PYG{l+m}{0}\PYG{+w}{ }\PYG{l+m}{1}\PYG{+w}{ }\PYG{n}{units}\PYG{+w}{ }\PYG{n}{lattice}
\PYG{k}{create\PYGZus{}box   }\PYG{n+nv+nvIdentifier}{1}\PYG{+w}{ }\PYG{n+nv+nvIdentifier}{simbox}
\PYG{k}{create\PYGZus{}atoms}\PYG{+w}{ }\PYG{l+m}{1}\PYG{+w}{ }\PYG{n}{box}
\PYG{k}{mass}\PYG{+w}{         }\PYG{l+m}{1}\PYG{+w}{ }\PYG{l+m}{26.981539}
\PYG{k}{kim          }\PYG{l+s+s1}{interactions}\PYG{+w}{ }\PYG{n}{Al}
\PYG{k}{run}\PYG{+w}{          }\PYG{l+m}{0}
\PYG{k}{variable     }\PYG{n+nv+nvIdentifier}{Ec}\PYG{+w}{ }\PYG{n}{equal}\PYG{+w}{ }\PYG{n+nv}{(}\PYG{n+nv}{pe/count}\PYG{n+nv}{(}\PYG{n+nv}{all}\PYG{n+nv}{)}\PYG{n+nv}{)}
\PYG{k}{print}\PYG{+w}{        }\PYG{l+s}{\PYGZdq{}}\PYG{l+s}{Cohesive Energy = \PYGZdl{}\PYGZob{}Ec\PYGZcb{} eV}\PYG{l+s}{\PYGZdq{}}
\end{sphinxVerbatim}

\sphinxAtStartPar
The above script will end with an error in the \sphinxstyleemphasis{kim init} line if the IM is
changed to another potential for Al that does not work with \sphinxstyleemphasis{metal} units.  To
address this, \sphinxstyleemphasis{kim init} offers the \sphinxstyleemphasis{unit\_conversion\_mode} as shown below.

\sphinxAtStartPar
If unit conversion mode \sphinxstyleemphasis{is} active, then \sphinxstyleemphasis{kim init} calls the LAMMPS
{\hyperref[\detokenize{units::doc}]{\sphinxcrossref{\DUrole{doc}{units}}}} command to set the units to the IM’s required or preferred
units.  Conversion factors between the IM’s units and the \sphinxstyleemphasis{user\_units} are
defined for all {\hyperref[\detokenize{units::doc}]{\sphinxcrossref{\DUrole{doc}{physical quantities}}}} (mass, distance, etc.).
(Note that converting to or from the “lj” unit style is not supported.)
These factors are stored as {\hyperref[\detokenize{variable::doc}]{\sphinxcrossref{\DUrole{doc}{internal style variables}}}} with the
following standard names:

\begin{sphinxVerbatim}[commandchars=\\\{\}]
\PYGZus{}u\PYGZus{}mass
\PYGZus{}u\PYGZus{}distance
\PYGZus{}u\PYGZus{}time
\PYGZus{}u\PYGZus{}energy
\PYGZus{}u\PYGZus{}velocity
\PYGZus{}u\PYGZus{}force
\PYGZus{}u\PYGZus{}torque
\PYGZus{}u\PYGZus{}temperature
\PYGZus{}u\PYGZus{}pressure
\PYGZus{}u\PYGZus{}viscosity
\PYGZus{}u\PYGZus{}charge
\PYGZus{}u\PYGZus{}dipole
\PYGZus{}u\PYGZus{}efield
\PYGZus{}u\PYGZus{}density
\end{sphinxVerbatim}

\sphinxAtStartPar
If desired, the input script can be designed to work with these conversion
factors so that the script will work without change with any OpenKIM IM.
(This approach is used in the
\sphinxhref{https://openkim.org/doc/evaluation/kim-tests/}{OpenKIM Testing Framework}.)

\sphinxAtStartPar
For example, the script given above for the cohesive energy of fcc Al can be
rewritten to work with any IM regardless of units.  The following script
constructs an fcc lattice with a lattice parameter defined in meters, computes
the total energy, and prints the cohesive energy in Joules regardless of the
units of the IM.

\begin{sphinxVerbatim}[commandchars=\\\{\}]
\PYG{k}{kim          }\PYG{l+s+s1}{init}\PYG{+w}{ }\PYG{n}{Sim\PYGZus{}LAMMPS\PYGZus{}MEAM\PYGZus{}PascuetFernandez\PYGZus{}2015\PYGZus{}Al\PYGZus{}\PYGZus{}SM\PYGZus{}811588957187\PYGZus{}000}\PYG{+w}{ }\PYG{n}{si}\PYG{+w}{ }\PYG{n}{unit\PYGZus{}conversion\PYGZus{}mode}
\PYG{k}{boundary}\PYG{+w}{     }\PYG{n}{p}\PYG{+w}{ }\PYG{n}{p}\PYG{+w}{ }\PYG{n}{p}
\PYG{k}{lattice}\PYG{+w}{      }\PYG{n}{fcc}\PYG{+w}{ }\PYG{n+nv}{\PYGZdl{}(}\PYG{n+nv}{4.049e\PYGZhy{}10*v\PYGZus{}\PYGZus{}u\PYGZus{}distance}\PYG{n+nv}{)}
\PYG{k}{region       }\PYG{n+nv+nvIdentifier}{simbox}\PYG{+w}{ }\PYG{n}{block}\PYG{+w}{ }\PYG{l+m}{0}\PYG{+w}{ }\PYG{l+m}{1}\PYG{+w}{ }\PYG{l+m}{0}\PYG{+w}{ }\PYG{l+m}{1}\PYG{+w}{ }\PYG{l+m}{0}\PYG{+w}{ }\PYG{l+m}{1}\PYG{+w}{ }\PYG{n}{units}\PYG{+w}{ }\PYG{n}{lattice}
\PYG{k}{create\PYGZus{}box   }\PYG{n+nv+nvIdentifier}{1}\PYG{+w}{ }\PYG{n+nv+nvIdentifier}{simbox}
\PYG{k}{create\PYGZus{}atoms}\PYG{+w}{ }\PYG{l+m}{1}\PYG{+w}{ }\PYG{n}{box}
\PYG{k}{mass}\PYG{+w}{         }\PYG{l+m}{1}\PYG{+w}{ }\PYG{n+nv}{\PYGZdl{}(}\PYG{n+nv}{4.480134e\PYGZhy{}26*v\PYGZus{}\PYGZus{}u\PYGZus{}mass}\PYG{n+nv}{)}
\PYG{k}{kim          }\PYG{l+s+s1}{interactions}\PYG{+w}{ }\PYG{n}{Al}
\PYG{k}{neighbor}\PYG{+w}{     }\PYG{n+nv}{\PYGZdl{}(}\PYG{n+nv}{2e\PYGZhy{}10*v\PYGZus{}\PYGZus{}u\PYGZus{}distance}\PYG{n+nv}{)}\PYG{+w}{ }\PYG{n}{bin}
\PYG{k}{run}\PYG{+w}{          }\PYG{l+m}{0}
\PYG{k}{variable     }\PYG{n+nv+nvIdentifier}{Ec\PYGZus{}in\PYGZus{}J}\PYG{+w}{ }\PYG{n}{equal}\PYG{+w}{ }\PYG{n+nv}{(}\PYG{n+nv}{pe/count}\PYG{n+nv}{(}\PYG{n+nv}{all}\PYG{n+nv}{)}\PYG{n+nv}{)}\PYG{o}{/}\PYG{n}{v\PYGZus{}\PYGZus{}u\PYGZus{}energy}
\PYG{k}{print}\PYG{+w}{        }\PYG{l+s}{\PYGZdq{}}\PYG{l+s}{Cohesive Energy = \PYGZdl{}\PYGZob{}Ec\PYGZus{}in\PYGZus{}J\PYGZcb{} J}\PYG{l+s}{\PYGZdq{}}
\end{sphinxVerbatim}

\sphinxAtStartPar
Note the multiplication by \sphinxtitleref{v\_\_u\_distance} and \sphinxtitleref{v\_\_u\_mass} to convert from SI
units (specified in the \sphinxstyleemphasis{kim init} command) to whatever units the IM uses (metal
in this case), and the division by \sphinxtitleref{v\_\_u\_energy} to convert from the IM’s energy
units to SI units (Joule).  This script will work correctly for any IM for Al
(KIM PM or SM) selected by the \sphinxstyleemphasis{kim init} command.

\sphinxAtStartPar
Care must be taken to apply unit conversion to dimensional variables read in
from a file.  For example, if a configuration of atoms is read in from a dump
file using the {\hyperref[\detokenize{read_dump::doc}]{\sphinxcrossref{\DUrole{doc}{read\_dump}}}} command, the following can be done
to convert the box and all atomic positions to the correct units:

\begin{sphinxVerbatim}[commandchars=\\\{\}]
\PYG{k}{change\PYGZus{}box }\PYG{n+nv+nvIdentifier}{all}\PYG{+w}{ }\PYG{n}{x}\PYG{+w}{ }\PYG{n}{scale}\PYG{+w}{ }\PYG{n+nv}{\PYGZdl{}\PYGZob{}}\PYG{n+nv}{\PYGZus{}u\PYGZus{}distance}\PYG{n+nv}{\PYGZcb{}}\PYG{+w}{ }\PYG{o}{\PYGZam{}}
\PYG{+w}{               }\PYG{n}{y}\PYG{+w}{ }\PYG{n}{scale}\PYG{+w}{ }\PYG{n+nv}{\PYGZdl{}\PYGZob{}}\PYG{n+nv}{\PYGZus{}u\PYGZus{}distance}\PYG{n+nv}{\PYGZcb{}}\PYG{+w}{ }\PYG{o}{\PYGZam{}}
\PYG{+w}{               }\PYG{n}{z}\PYG{+w}{ }\PYG{n}{scale}\PYG{+w}{ }\PYG{n+nv}{\PYGZdl{}\PYGZob{}}\PYG{n+nv}{\PYGZus{}u\PYGZus{}distance}\PYG{n+nv}{\PYGZcb{}}\PYG{+w}{ }\PYG{o}{\PYGZam{}}
\PYG{+w}{               }\PYG{n}{xy}\PYG{+w}{ }\PYG{n}{final}\PYG{+w}{ }\PYG{n+nv}{\PYGZdl{}(}\PYG{n+nv}{xy*v\PYGZus{}\PYGZus{}u\PYGZus{}distance}\PYG{n+nv}{)}\PYG{+w}{ }\PYG{o}{\PYGZam{}}
\PYG{+w}{               }\PYG{n}{xz}\PYG{+w}{ }\PYG{n}{final}\PYG{+w}{ }\PYG{n+nv}{\PYGZdl{}(}\PYG{n+nv}{xz*v\PYGZus{}\PYGZus{}u\PYGZus{}distance}\PYG{n+nv}{)}\PYG{+w}{ }\PYG{o}{\PYGZam{}}
\PYG{+w}{               }\PYG{n}{yz}\PYG{+w}{ }\PYG{n}{final}\PYG{+w}{ }\PYG{n+nv}{\PYGZdl{}(}\PYG{n+nv}{yz*v\PYGZus{}\PYGZus{}u\PYGZus{}distance}\PYG{n+nv}{)}\PYG{+w}{ }\PYG{o}{\PYGZam{}}
\PYG{+w}{               }\PYG{n}{remap}
\end{sphinxVerbatim}

\begin{sphinxadmonition}{note}{Note:}
\sphinxAtStartPar
Unit conversion will only work if the conversion factors are placed in all
appropriate places in the input script.  It is up to the user to do this
correctly.
\end{sphinxadmonition}


\subsubsection{OpenKIM IM Execution (\sphinxstyleemphasis{kim interactions})}
\label{\detokenize{kim_commands:openkim-im-execution-kim-interactions}}\label{\detokenize{kim_commands:interactions}}
\sphinxAtStartPar
The second and final step in using an OpenKIM IM is to execute the
\sphinxstyleemphasis{kim interactions} command.  This command must be preceded by a \sphinxstyleemphasis{kim init}
command and a command that defines the number of atom types \sphinxstyleemphasis{N} (such as
{\hyperref[\detokenize{create_box::doc}]{\sphinxcrossref{\DUrole{doc}{create\_box}}}}).
The \sphinxstyleemphasis{kim interactions} command has one argument \sphinxstyleemphasis{typeargs}.  This argument
contains either a list of \sphinxstyleemphasis{N} chemical species, which defines a mapping between
atom types in LAMMPS to the available species in the OpenKIM IM, or the keyword
\sphinxstyleemphasis{fixed\_types} for models that have a preset fixed mapping (i.e.  the mapping
between LAMMPS atom types and chemical species is defined by the model and
cannot be changed).  In the latter case, the user must consult the model
documentation to see how many atom types there are and how they map to the
chemical species.

\sphinxAtStartPar
For example, consider an OpenKIM IM that supports Si and C species.  If the
LAMMPS simulation has four atom types, where the first three are Si, and the
fourth is C, the following \sphinxstyleemphasis{kim interactions} command would be used:

\begin{sphinxVerbatim}[commandchars=\\\{\}]
\PYG{k}{kim }\PYG{l+s+s1}{interactions}\PYG{+w}{ }\PYG{n}{Si}\PYG{+w}{ }\PYG{n}{Si}\PYG{+w}{ }\PYG{n}{Si}\PYG{+w}{ }\PYG{n}{C}
\end{sphinxVerbatim}

\sphinxAtStartPar
Alternatively, for a model with a fixed mapping the command would be:

\begin{sphinxVerbatim}[commandchars=\\\{\}]
\PYG{k}{kim }\PYG{l+s+s1}{interactions}\PYG{+w}{ }\PYG{n}{fixed\PYGZus{}types}
\end{sphinxVerbatim}

\sphinxAtStartPar
The \sphinxstyleemphasis{kim interactions} command performs all the necessary steps to set up the
OpenKIM IM selected in the \sphinxstyleemphasis{kim init} command.  The specific actions depend on
whether the IM is a KIM PM or a KIM SM.  For a KIM PM, a
{\hyperref[\detokenize{pair_kim::doc}]{\sphinxcrossref{\DUrole{doc}{pair\_style kim}}}} command is executed followed by the appropriate
\sphinxstyleemphasis{pair\_coeff} command.  For example, for the Ercolessi and Adams (1994) KIM PM
for Al set by the following commands:

\begin{sphinxVerbatim}[commandchars=\\\{\}]
\PYG{k}{kim }\PYG{l+s+s1}{init}\PYG{+w}{ }\PYG{n}{EAM\PYGZus{}Dynamo\PYGZus{}ErcolessiAdams\PYGZus{}1994\PYGZus{}Al\PYGZus{}\PYGZus{}MO\PYGZus{}123629422045\PYGZus{}005}\PYG{+w}{ }\PYG{n}{metal}
\PYG{n}{...}
\PYG{n}{...}\PYG{+w}{  }\PYG{n}{box}\PYG{+w}{ }\PYG{n}{specification}\PYG{+w}{ }\PYG{n}{lines}\PYG{+w}{ }\PYG{n}{skipped}
\PYG{n}{...}
\PYG{k}{kim }\PYG{l+s+s1}{interactions}\PYG{+w}{ }\PYG{n}{Al}
\end{sphinxVerbatim}

\sphinxAtStartPar
the \sphinxstyleemphasis{kim interactions} command executes the following LAMMPS input commands:

\begin{sphinxVerbatim}[commandchars=\\\{\}]
\PYG{k}{pair\PYGZus{}style}\PYG{+w}{ }\PYG{k}{kim }\PYG{l+s+s1}{EAM\PYGZus{}Dynamo\PYGZus{}ErcolessiAdams\PYGZus{}1994\PYGZus{}Al\PYGZus{}\PYGZus{}MO\PYGZus{}123629422045\PYGZus{}005}
\PYG{k}{pair\PYGZus{}coeff}\PYG{+w}{ }\PYG{o}{*}\PYG{+w}{ }\PYG{o}{*}\PYG{+w}{ }\PYG{n}{Al}
\end{sphinxVerbatim}

\sphinxAtStartPar
For a KIM SM, the generated input commands may be more complex and require that
LAMMPS is built with the required packages included for the type of potential
being used.  The set of commands to be executed is defined in the SM
specification file, which is part of the SM package.  For example, for the
Strachan et al.  (2003) ReaxFF SM set by the following commands:

\begin{sphinxVerbatim}[commandchars=\\\{\}]
\PYG{k}{kim }\PYG{l+s+s1}{init}\PYG{+w}{ }\PYG{n}{Sim\PYGZus{}LAMMPS\PYGZus{}ReaxFF\PYGZus{}StrachanVanDuinChakraborty\PYGZus{}2003\PYGZus{}CHNO\PYGZus{}\PYGZus{}SM\PYGZus{}107643900657\PYGZus{}000}\PYG{+w}{ }\PYG{n}{real}
\PYG{n}{...}
\PYG{n}{...}\PYG{+w}{  }\PYG{n}{box}\PYG{+w}{ }\PYG{n}{specification}\PYG{+w}{ }\PYG{n}{lines}\PYG{+w}{ }\PYG{n}{skipped}
\PYG{n}{...}
\PYG{k}{kim }\PYG{l+s+s1}{interactions}\PYG{+w}{ }\PYG{n}{C}\PYG{+w}{ }\PYG{n}{H}\PYG{+w}{ }\PYG{n}{N}\PYG{+w}{ }\PYG{n}{O}
\end{sphinxVerbatim}

\sphinxAtStartPar
the \sphinxstyleemphasis{kim interactions} command executes the following LAMMPS input commands:

\begin{sphinxVerbatim}[commandchars=\\\{\}]
\PYG{k}{pair\PYGZus{}style}\PYG{+w}{ }\PYG{n}{reaxff}\PYG{+w}{ }\PYG{n}{lmp\PYGZus{}control}\PYG{+w}{ }\PYG{n}{safezone}\PYG{+w}{ }\PYG{l+m}{2.0}\PYG{+w}{ }\PYG{n}{mincap}\PYG{+w}{ }\PYG{l+m}{100}
\PYG{k}{pair\PYGZus{}coeff}\PYG{+w}{ }\PYG{o}{*}\PYG{+w}{ }\PYG{o}{*}\PYG{+w}{ }\PYG{n}{ffield.reax.rdx}\PYG{+w}{ }\PYG{n}{C}\PYG{+w}{ }\PYG{n}{H}\PYG{+w}{ }\PYG{n}{N}\PYG{+w}{ }\PYG{n}{O}
\PYG{k}{fix }\PYG{n+nv+nvIdentifier}{reaxqeq}\PYG{+w}{ }\PYG{n+nv+nvIdentifier}{all}\PYG{+w}{ }\PYG{n}{qeq}\PYG{o}{/}\PYG{n}{reaxff}\PYG{+w}{ }\PYG{l+m}{1}\PYG{+w}{ }\PYG{l+m}{0.0}\PYG{+w}{ }\PYG{l+m}{10.0}\PYG{+w}{ }\PYG{l+m}{1.0e\PYGZhy{}6}\PYG{+w}{ }\PYG{n}{param.qeq}
\end{sphinxVerbatim}

\begin{sphinxadmonition}{note}{Note:}
\sphinxAtStartPar
The files \sphinxstyleemphasis{lmp\_control}, \sphinxstyleemphasis{ffield.reax.rdx} and \sphinxstyleemphasis{param.qeq} are specific to
the Strachan et al.  (2003) ReaxFF parameterization and are archived as part
of the SM package in OpenKIM.
\end{sphinxadmonition}

\begin{sphinxadmonition}{note}{Note:}
\sphinxAtStartPar
Parameters like cutoff radii and charge tolerances, which have an effect on
IM predictions, are also included in the SM definition ensuring
reproducibility.
\end{sphinxadmonition}

\begin{sphinxadmonition}{note}{Note:}
\sphinxAtStartPar
When using \sphinxstyleemphasis{kim init} and \sphinxstyleemphasis{kim interactions} to select and set up an OpenKIM
IM, other LAMMPS commands for the same functions (such as pair\_style,
pair\_coeff, bond\_style, bond\_coeff, fixes related to charge equilibration,
etc.) should normally not appear in the input script.
\end{sphinxadmonition}

\begin{sphinxadmonition}{note}{Note:}
\sphinxAtStartPar
\sphinxstyleemphasis{kim interactions} must be called each time after the
{\hyperref[\detokenize{change_box::doc}]{\sphinxcrossref{\DUrole{doc}{change\_box}}}} command to provide the correct settings (it
should be called with the same {\hyperref[\detokenize{kim_commands:typeargs-options}]{\sphinxcrossref{\DUrole{std,std-ref}{typeargs}}}} as the
first call.)  The reason is that changing a periodic boundary to a
non\sphinxhyphen{}periodic one, or in general, using the \sphinxstyleemphasis{change\_box} command after the
interactions are set via \sphinxstyleemphasis{kim interactions} or \sphinxstyleemphasis{pair\_coeff} commands might
affect some of the settings.  For example, SM models containing Coulombic
terms in the interactions require different settings if a periodic boundary
changes to a non\sphinxhyphen{}periodic one.  In other cases, the second call to
\sphinxstyleemphasis{kim interactions} does not affect any other settings.
\end{sphinxadmonition}


\subsection{Using OpenKIM Web Queries in LAMMPS (\sphinxstyleemphasis{kim query})}
\label{\detokenize{kim_commands:using-openkim-web-queries-in-lammps-kim-query}}\label{\detokenize{kim_commands:query}}
\sphinxAtStartPar
The \sphinxstyleemphasis{kim query} command performs a web query to retrieve the predictions of an
IM set by \sphinxstyleemphasis{kim init} for material properties archived in
\sphinxhref{https://openkim.org}{OpenKIM}.


\subsubsection{Syntax}
\label{\detokenize{kim_commands:id5}}
\begin{sphinxVerbatim}[commandchars=\\\{\}]
\PYG{k}{kim }\PYG{l+s+s1}{query}\PYG{+w}{ }\PYG{k}{variable }\PYG{n+nv+nvIdentifier}{formatarg}\PYG{+w}{ }\PYG{n}{query\PYGZus{}function}\PYG{+w}{ }\PYG{n}{queryargs}
\end{sphinxVerbatim}
\begin{itemize}
\item {} 
\sphinxAtStartPar
variable(s) = single name or list of names of (string style) LAMMPS
variable(s) where a query result or parameter get result is stored.  Variables
that do not exist will be created by the command

\item {} 
\sphinxAtStartPar
formatarg = \sphinxstyleemphasis{list} or \sphinxstyleemphasis{split} or \sphinxstyleemphasis{index} (optional)
\begin{sphinxalltt}
\sphinxstyleemphasis{list} = returns a single string with a list of space separated values
       (e.g. "1.0 2.0 3.0"), which is placed in a LAMMPS variable as
       defined by the \sphinxstyleemphasis{variable} argument. {[}default{]}
\sphinxstyleemphasis{split} = returns the values separately in new variables with names based
       on the prefix specified in \sphinxstyleemphasis{variable} and a number appended to
       indicate which element in the list of values is in the variable
\sphinxstyleemphasis{index} = returns a variable style index that can be incremented via the
       next command.  This enables the construction of simple loops
\end{sphinxalltt}

\item {} 
\sphinxAtStartPar
query\_function = name of the OpenKIM web API query function to be used

\item {} 
\sphinxAtStartPar
queryargs = a series of \sphinxstyleemphasis{keyword=value} pairs that represent the web query;
supported keywords depend on the query function

\end{itemize}


\subsubsection{Examples}
\label{\detokenize{kim_commands:id6}}
\begin{sphinxVerbatim}[commandchars=\\\{\}]
\PYG{k}{kim }\PYG{l+s+s1}{query}\PYG{+w}{ }\PYG{n}{a0}\PYG{+w}{ }\PYG{n}{get\PYGZus{}lattice\PYGZus{}constant\PYGZus{}cubic}\PYG{+w}{ }\PYG{n}{crystal}\PYG{o}{=}\PYG{n}{[fcc]}\PYG{+w}{ }\PYG{n}{species}\PYG{o}{=}\PYG{n}{[Al]}\PYG{+w}{ }\PYG{n}{units}\PYG{o}{=}\PYG{n}{[angstrom]}
\PYG{k}{kim }\PYG{l+s+s1}{query}\PYG{+w}{ }\PYG{n}{model}\PYG{+w}{ }\PYG{n}{index}\PYG{+w}{ }\PYG{n}{get\PYGZus{}available\PYGZus{}models}\PYG{+w}{ }\PYG{n}{species}\PYG{o}{=}\PYG{n}{[Al]}\PYG{+w}{ }\PYG{n}{potential\PYGZus{}type}\PYG{o}{=}\PYG{n}{[eam]}
\end{sphinxVerbatim}

\sphinxAtStartPar
The result of the query is stored in one or more
{\hyperref[\detokenize{variable::doc}]{\sphinxcrossref{\DUrole{doc}{string style variables}}}} as determined by the optional
\sphinxstyleemphasis{formatarg} argument.  For the “list”
setting of \sphinxstyleemphasis{formatarg} (or if \sphinxstyleemphasis{formatarg} is not specified), the result is
returned as a space\sphinxhyphen{}separated list of values in \sphinxstyleemphasis{variable}.  The \sphinxstyleemphasis{formatarg}
keyword “split” separates the result values into individual variables of the
form \sphinxstyleemphasis{prefix\_I}, where \sphinxstyleemphasis{prefix} is set to the \sphinxstyleemphasis{kim query} \sphinxstyleemphasis{variable} argument
and \sphinxstyleemphasis{I} ranges from 1 to the number of returned values.  The number and order of
the returned values is determined by the type of query performed.  The
\sphinxstyleemphasis{formatarg} keyword “index” returns a {\hyperref[\detokenize{variable::doc}]{\sphinxcrossref{\DUrole{doc}{variable style index}}}}
that can be incremented via the {\hyperref[\detokenize{next::doc}]{\sphinxcrossref{\DUrole{doc}{next}}}} command.  This enables the
construction of simple loops over the returned values by the type of query
performed.

\begin{sphinxadmonition}{note}{Note:}
\sphinxAtStartPar
\sphinxstyleemphasis{kim query} only supports queries that return a single result or an array of
values.  More complex queries that return a JSON structure are not currently
supported.  An attempt to use \sphinxstyleemphasis{kim query} in such cases will generate an
error.
\end{sphinxadmonition}

\sphinxAtStartPar
The second required argument \sphinxstyleemphasis{query\_function} is the name of the query function
to be called (e.g.  \sphinxstyleemphasis{get\_lattice\_constant\_cubic}).  All following
{\hyperref[\detokenize{Commands_parse::doc}]{\sphinxcrossref{\DUrole{doc}{arguments}}}} are parameters handed over to the web query in
the format \sphinxstyleemphasis{keyword=value}, where \sphinxstyleemphasis{value} is always an array of one or more
comma\sphinxhyphen{}separated items in brackets.  The list of supported keywords and the type
and format of their values depend on the query function used.  The current list
of query functions is available on the OpenKIM webpage at
\sphinxurl{https://openkim.org/doc/usage/kim-query}.

\begin{sphinxadmonition}{note}{Note:}
\sphinxAtStartPar
All query functions, except \sphinxstyleemphasis{get\_available\_models}, require the
\sphinxstyleemphasis{model} keyword, which identifies the IM whose predictions are being queried.
\sphinxstyleemphasis{kim query} automatically generates the \sphinxstyleemphasis{model} keyword based on the IM set
in by \sphinxstyleemphasis{kim init}, and it can be overwritten if specified as an argument to
the \sphinxstyleemphasis{kim query}.  Where \sphinxstyleemphasis{kim init} is not specified, the \sphinxstyleemphasis{model} keyword
must be provided as an argument to the \sphinxstyleemphasis{kim query}.
\end{sphinxadmonition}

\begin{sphinxadmonition}{note}{Note:}
\sphinxAtStartPar
Each \sphinxstyleemphasis{query\_function} is associated with a default method (implemented as a
\sphinxhref{https://openkim.org/doc/evaluation/kim-tests/}{KIM Test}) used to compute
this property.  In cases where there are multiple methods in OpenKIM for
computing a property, a \sphinxstyleemphasis{method} keyword can be provided to select the method
of choice.  See the
\sphinxhref{https://openkim.org/doc/usage/kim-query}{query documentation} to see which
methods are available for a given \sphinxstyleemphasis{query\_function}.
\end{sphinxadmonition}


\subsubsection{\sphinxstyleemphasis{kim query} Usage Examples and Further Clarifications}
\label{\detokenize{kim_commands:kim-query-usage-examples-and-further-clarifications}}
\sphinxAtStartPar
The data obtained by \sphinxstyleemphasis{kim query} commands can be used as part of the setup or
analysis phases of LAMMPS simulations.  Some examples are given below.

\sphinxAtStartPar
\sphinxstylestrong{Define an equilibrium fcc crystal}

\begin{sphinxVerbatim}[commandchars=\\\{\}]
\PYG{k}{kim      }\PYG{l+s+s1}{init}\PYG{+w}{ }\PYG{n}{EAM\PYGZus{}Dynamo\PYGZus{}ErcolessiAdams\PYGZus{}1994\PYGZus{}Al\PYGZus{}\PYGZus{}MO\PYGZus{}123629422045\PYGZus{}005}\PYG{+w}{ }\PYG{n}{metal}
\PYG{k}{boundary}\PYG{+w}{ }\PYG{n}{p}\PYG{+w}{ }\PYG{n}{p}\PYG{+w}{ }\PYG{n}{p}
\PYG{k}{kim      }\PYG{l+s+s1}{query}\PYG{+w}{ }\PYG{n}{a0}\PYG{+w}{ }\PYG{n}{get\PYGZus{}lattice\PYGZus{}constant\PYGZus{}cubic}\PYG{+w}{ }\PYG{n}{crystal}\PYG{o}{=}\PYG{n}{[fcc]}\PYG{+w}{ }\PYG{n}{species}\PYG{o}{=}\PYG{n}{[Al]}\PYG{+w}{ }\PYG{n}{units}\PYG{o}{=}\PYG{n}{[angstrom]}
\PYG{k}{lattice}\PYG{+w}{  }\PYG{n}{fcc}\PYG{+w}{ }\PYG{n+nv}{\PYGZdl{}\PYGZob{}}\PYG{n+nv}{a0}\PYG{n+nv}{\PYGZcb{}}
\PYG{n}{...}
\end{sphinxVerbatim}

\begin{sphinxVerbatim}[commandchars=\\\{\}]
\PYG{k}{units}\PYG{+w}{   }\PYG{n}{metal}
\PYG{k}{kim     }\PYG{l+s+s1}{query}\PYG{+w}{ }\PYG{n}{a0}\PYG{+w}{ }\PYG{n}{get\PYGZus{}lattice\PYGZus{}constant\PYGZus{}cubic}\PYG{+w}{ }\PYG{n}{crystal}\PYG{o}{=}\PYG{n}{[fcc]}\PYG{+w}{ }\PYG{n}{species}\PYG{o}{=}\PYG{n}{[Al]}\PYG{+w}{ }\PYG{n}{units}\PYG{o}{=}\PYG{n}{[angstrom]}\PYG{+w}{ }\PYG{n}{model}\PYG{o}{=}\PYG{n}{[EAM\PYGZus{}Dynamo\PYGZus{}ErcolessiAdams\PYGZus{}1994\PYGZus{}Al\PYGZus{}\PYGZus{}MO\PYGZus{}123629422045\PYGZus{}005]}
\PYG{k}{lattice}\PYG{+w}{ }\PYG{n}{fcc}\PYG{+w}{ }\PYG{n+nv}{\PYGZdl{}\PYGZob{}}\PYG{n+nv}{a0}\PYG{n+nv}{\PYGZcb{}}
\PYG{n}{...}
\end{sphinxVerbatim}

\sphinxAtStartPar
The \sphinxstyleemphasis{kim query} command retrieves from \sphinxhref{https://openkim.org}{OpenKIM} the
equilibrium lattice constant predicted by the Ercolessi and Adams (1994)
potential for the fcc structure and places it in variable \sphinxstyleemphasis{a0}.  This variable
is then used on the next line to set up the crystal.  By using \sphinxstyleemphasis{kim query}, the
user is saved the trouble and possible error of tracking this value down, or of
having to perform an energy minimization to find the equilibrium lattice
constant.

\begin{sphinxadmonition}{note}{Note:}
\sphinxAtStartPar
In \sphinxstyleemphasis{unit\_conversion\_mode} the results obtained from a \sphinxstyleemphasis{kim query} would need
to be converted to the appropriate units system.  For example, in the above
script, the lattice command would need to be changed to:
“lattice fcc \$(v\_a0*v\_\_u\_distance)”.
\end{sphinxadmonition}

\sphinxAtStartPar
\sphinxstylestrong{Define an equilibrium hcp crystal}

\begin{sphinxVerbatim}[commandchars=\\\{\}]
\PYG{k}{kim      }\PYG{l+s+s1}{init}\PYG{+w}{ }\PYG{n}{EAM\PYGZus{}Dynamo\PYGZus{}MendelevAckland\PYGZus{}2007v3\PYGZus{}Zr\PYGZus{}\PYGZus{}MO\PYGZus{}004835508849\PYGZus{}000}\PYG{+w}{ }\PYG{n}{metal}
\PYG{k}{boundary}\PYG{+w}{ }\PYG{n}{p}\PYG{+w}{ }\PYG{n}{p}\PYG{+w}{ }\PYG{n}{p}
\PYG{k}{kim      }\PYG{l+s+s1}{query}\PYG{+w}{ }\PYG{n}{latconst}\PYG{+w}{ }\PYG{n}{split}\PYG{+w}{ }\PYG{n}{get\PYGZus{}lattice\PYGZus{}constant\PYGZus{}hexagonal}\PYG{+w}{ }\PYG{n}{crystal}\PYG{o}{=}\PYG{n}{[hcp]}\PYG{+w}{ }\PYG{n}{species}\PYG{o}{=}\PYG{n}{[Zr]}\PYG{+w}{ }\PYG{n}{units}\PYG{o}{=}\PYG{n}{[angstrom]}
\PYG{k}{lattice}\PYG{+w}{  }\PYG{n}{custom}\PYG{+w}{ }\PYG{n+nv}{\PYGZdl{}\PYGZob{}}\PYG{n+nv}{latconst\PYGZus{}1}\PYG{n+nv}{\PYGZcb{}}\PYG{+w}{ }\PYG{n}{a1}\PYG{+w}{ }\PYG{l+m}{0.5}\PYG{+w}{ }\PYG{o}{\PYGZhy{}}\PYG{l+m}{0.866025}\PYG{+w}{ }\PYG{l+m}{0}\PYG{+w}{ }\PYG{n}{a2}\PYG{+w}{ }\PYG{l+m}{0.5}\PYG{+w}{ }\PYG{l+m}{0.866025}\PYG{+w}{ }\PYG{l+m}{0}\PYG{+w}{ }\PYG{n}{a3}\PYG{+w}{ }\PYG{l+m}{0}\PYG{+w}{ }\PYG{l+m}{0}\PYG{+w}{ }\PYG{n+nv}{\PYGZdl{}(}\PYG{n+nv}{latconst\PYGZus{}2/latconst\PYGZus{}1}\PYG{n+nv}{)}\PYG{+w}{ }\PYG{o}{\PYGZam{}}
\PYG{+w}{         }\PYG{n}{basis}\PYG{+w}{ }\PYG{l+m}{0.333333}\PYG{+w}{ }\PYG{l+m}{0.666666}\PYG{+w}{ }\PYG{l+m}{0.25}\PYG{+w}{ }\PYG{n}{basis}\PYG{+w}{ }\PYG{l+m}{0.666666}\PYG{+w}{ }\PYG{l+m}{0.333333}\PYG{+w}{ }\PYG{l+m}{0.75}
\PYG{n}{...}
\end{sphinxVerbatim}

\sphinxAtStartPar
In this case the \sphinxstyleemphasis{kim query} returns two arguments (since the hexagonal close
packed (hcp) structure has two independent lattice constants).  The \sphinxstyleemphasis{formatarg}
keyword “split” places the two values into the variables \sphinxstyleemphasis{latconst\_1} and
\sphinxstyleemphasis{latconst\_2}.  (These variables are created if they do not already exist.)

\sphinxAtStartPar
\sphinxstylestrong{Define a crystal at finite temperature accounting for thermal expansion}

\begin{sphinxVerbatim}[commandchars=\\\{\}]
\PYG{k}{kim      }\PYG{l+s+s1}{init}\PYG{+w}{ }\PYG{n}{EAM\PYGZus{}Dynamo\PYGZus{}ErcolessiAdams\PYGZus{}1994\PYGZus{}Al\PYGZus{}\PYGZus{}MO\PYGZus{}123629422045\PYGZus{}005}\PYG{+w}{ }\PYG{n}{metal}
\PYG{k}{boundary}\PYG{+w}{ }\PYG{n}{p}\PYG{+w}{ }\PYG{n}{p}\PYG{+w}{ }\PYG{n}{p}
\PYG{k}{kim      }\PYG{l+s+s1}{query}\PYG{+w}{ }\PYG{n}{a0}\PYG{+w}{ }\PYG{n}{get\PYGZus{}lattice\PYGZus{}constant\PYGZus{}cubic}\PYG{+w}{ }\PYG{n}{crystal}\PYG{o}{=}\PYG{n}{[fcc]}\PYG{+w}{ }\PYG{n}{species}\PYG{o}{=}\PYG{n}{[Al]}\PYG{+w}{ }\PYG{n}{units}\PYG{o}{=}\PYG{n}{[angstrom]}
\PYG{k}{kim      }\PYG{l+s+s1}{query}\PYG{+w}{ }\PYG{n}{alpha}\PYG{+w}{ }\PYG{n}{get\PYGZus{}linear\PYGZus{}thermal\PYGZus{}expansion\PYGZus{}coefficient\PYGZus{}cubic}\PYG{+w}{ }\PYG{n}{crystal}\PYG{o}{=}\PYG{n}{[fcc]}\PYG{+w}{ }\PYG{n}{species}\PYG{o}{=}\PYG{n}{[Al]}\PYG{+w}{ }\PYG{n}{units}\PYG{o}{=}\PYG{n}{[1}\PYG{o}{/}\PYG{n}{K]}\PYG{+w}{ }\PYG{n}{temperature}\PYG{o}{=}\PYG{n}{[293.15]}\PYG{+w}{ }\PYG{n}{temperature\PYGZus{}units}\PYG{o}{=}\PYG{n}{[K]}
\PYG{k}{variable }\PYG{n+nv+nvIdentifier}{DeltaT}\PYG{+w}{ }\PYG{n}{equal}\PYG{+w}{ }\PYG{l+m}{300}
\PYG{k}{lattice}\PYG{+w}{  }\PYG{n}{fcc}\PYG{+w}{ }\PYG{n+nv}{\PYGZdl{}(}\PYG{n+nv}{v\PYGZus{}a0*v\PYGZus{}alpha*v\PYGZus{}DeltaT}\PYG{n+nv}{)}
\PYG{n}{...}
\end{sphinxVerbatim}

\sphinxAtStartPar
As in the previous example, the equilibrium lattice constant is obtained for the
Ercolessi and Adams (1994) potential.  However, in this case the crystal is
scaled to the appropriate lattice constant at room temperature (293.15 K) by
using the linear thermal expansion constant predicted by the potential.

\begin{sphinxadmonition}{note}{Note:}
\sphinxAtStartPar
When passing numerical values as arguments (as in the case of the temperature
in the above example) it is also possible to pass a tolerance indicating how
close to the value is considered a match.  If no tolerance is passed a default
value is used.  If multiple results are returned (indicating that the
tolerance is too large), \sphinxstyleemphasis{kim query} will return an error.  See the
\sphinxhref{https://openkim.org/doc/usage/kim-query}{query documentation} to see which
numerical arguments and tolerances are available for a given
\sphinxstyleemphasis{query\_function}.
\end{sphinxadmonition}

\sphinxAtStartPar
\sphinxstylestrong{Compute defect formation energy}

\begin{sphinxVerbatim}[commandchars=\\\{\}]
\PYG{k}{kim      }\PYG{l+s+s1}{init}\PYG{+w}{ }\PYG{n}{EAM\PYGZus{}Dynamo\PYGZus{}ErcolessiAdams\PYGZus{}1994\PYGZus{}Al\PYGZus{}\PYGZus{}MO\PYGZus{}123629422045\PYGZus{}005}\PYG{+w}{ }\PYG{n}{metal}
\PYG{n}{...}
\PYG{n}{...}\PYG{+w}{ }\PYG{n}{Build}\PYG{+w}{ }\PYG{n}{fcc}\PYG{+w}{ }\PYG{n}{crystal}\PYG{+w}{ }\PYG{n}{containing}\PYG{+w}{ }\PYG{n}{some}\PYG{+w}{ }\PYG{n}{defect}\PYG{+w}{ }\PYG{n}{and}\PYG{+w}{ }\PYG{k}{compute }\PYG{n+nv+nvIdentifier}{the}\PYG{+w}{ }\PYG{n+nv+nvIdentifier}{total}\PYG{+w}{ }\PYG{n}{energy}
\PYG{n}{...}\PYG{+w}{ }\PYG{n}{which}\PYG{+w}{ }\PYG{n}{is}\PYG{+w}{ }\PYG{n}{stored}\PYG{+w}{ }\PYG{n}{in}\PYG{+w}{ }\PYG{n}{the}\PYG{+w}{ }\PYG{k}{variable }\PYG{o}{*}\PYG{n}{Etot}\PYG{o}{*}
\PYG{n}{...}
\PYG{k}{kim      }\PYG{l+s+s1}{query}\PYG{+w}{ }\PYG{n}{Ec}\PYG{+w}{ }\PYG{n}{get\PYGZus{}cohesive\PYGZus{}energy\PYGZus{}cubic}\PYG{+w}{ }\PYG{n}{crystal}\PYG{o}{=}\PYG{n}{[fcc]}\PYG{+w}{ }\PYG{n}{species}\PYG{o}{=}\PYG{n}{[Al]}\PYG{+w}{ }\PYG{n}{units}\PYG{o}{=}\PYG{n}{[eV]}
\PYG{k}{variable }\PYG{n+nv+nvIdentifier}{Eform}\PYG{+w}{ }\PYG{n}{equal}\PYG{+w}{ }\PYG{n+nv}{\PYGZdl{}\PYGZob{}}\PYG{n+nv}{Etot}\PYG{n+nv}{\PYGZcb{}}\PYG{+w}{ }\PYG{o}{\PYGZhy{}}\PYG{+w}{ }\PYG{n}{count}\PYG{n+nv}{(}\PYG{n+nv}{all}\PYG{n+nv}{)}\PYG{o}{*}\PYG{n+nv}{\PYGZdl{}\PYGZob{}}\PYG{n+nv}{Ec}\PYG{n+nv}{\PYGZcb{}}
\PYG{n}{...}
\end{sphinxVerbatim}

\sphinxAtStartPar
The defect formation energy \sphinxstyleemphasis{Eform} is computed by subtracting the ideal fcc
cohesive energy of the atoms in the system from \sphinxstyleemphasis{Etot}.  The ideal fcc
cohesive energy of the atoms is obtained from
\sphinxhref{https://openkim.org}{OpenKIM} for the Ercolessi and Adams (1994) potential.

\sphinxAtStartPar
\sphinxstylestrong{Retrieve equilibrium fcc crystal of all EAM potentials that support a specific species}

\begin{sphinxVerbatim}[commandchars=\\\{\}]
\PYG{k}{kim   }\PYG{l+s+s1}{query}\PYG{+w}{ }\PYG{n}{model}\PYG{+w}{ }\PYG{n}{index}\PYG{+w}{ }\PYG{n}{get\PYGZus{}available\PYGZus{}models}\PYG{+w}{ }\PYG{n}{species}\PYG{o}{=}\PYG{n}{[Al]}\PYG{+w}{ }\PYG{n}{potential\PYGZus{}type}\PYG{o}{=}\PYG{n}{[eam]}
\PYG{k}{label }\PYG{l+s+sc}{model\PYGZus{}loop}
\PYG{k}{kim   }\PYG{l+s+s1}{query}\PYG{+w}{ }\PYG{n}{latconst}\PYG{+w}{ }\PYG{n}{get\PYGZus{}lattice\PYGZus{}constant\PYGZus{}cubic}\PYG{+w}{ }\PYG{n}{crystal}\PYG{o}{=}\PYG{n}{[fcc]}\PYG{+w}{ }\PYG{n}{species}\PYG{o}{=}\PYG{n}{[Al]}\PYG{+w}{ }\PYG{n}{units}\PYG{o}{=}\PYG{n}{[angstrom]}\PYG{+w}{ }\PYG{n}{model}\PYG{o}{=}\PYG{n}{[}\PYG{n+nv}{\PYGZdl{}\PYGZob{}}\PYG{n+nv}{model}\PYG{n+nv}{\PYGZcb{}}\PYG{n}{]}
\PYG{k}{print}\PYG{+w}{ }\PYG{l+s}{\PYGZdq{}}\PYG{l+s}{FCC lattice constant (\PYGZdl{}\PYGZob{}model\PYGZcb{} potential) = \PYGZdl{}\PYGZob{}latconst\PYGZcb{}}\PYG{l+s}{\PYGZdq{}}
\PYG{n}{...}
\PYG{n}{...}\PYG{+w}{   }\PYG{n}{do}\PYG{+w}{ }\PYG{n}{something}\PYG{+w}{ }\PYG{n}{with}\PYG{+w}{ }\PYG{n}{current}\PYG{+w}{ }\PYG{n}{value}\PYG{+w}{ }\PYG{n}{of}\PYG{+w}{ }\PYG{n}{latconst}
\PYG{n}{...}
\PYG{k}{next  }\PYG{n+nv+nvIdentifier}{model}
\PYG{k}{jump  }\PYG{l+s+sc}{SELF}\PYG{+w}{ }\PYG{n}{model\PYGZus{}loop}
\end{sphinxVerbatim}

\sphinxAtStartPar
In this example, the \sphinxstyleemphasis{index} mode of \sphinxstyleemphasis{formatarg} is used.  The first \sphinxstyleemphasis{kim query}
returns the list of all available EAM potentials that support the \sphinxstyleemphasis{Al} species
and archived in \sphinxhref{https://openkim.org}{OpenKIM}.  The result of the query
operation is stored in the LAMMPS variable \sphinxstyleemphasis{model} as an index \sphinxstyleemphasis{variable}.
This variable is used later to access the values one at a time within a loop as
shown in the example.  The second \sphinxstyleemphasis{kim query} command retrieves from
\sphinxhref{https://openkim.org}{OpenKIM} the equilibrium lattice constant predicted by
each potential for the fcc structure and places it in variable \sphinxstyleemphasis{latconst}.

\begin{sphinxadmonition}{note}{Note:}
\sphinxAtStartPar
\sphinxstyleemphasis{kim query} commands return results archived in
\sphinxhref{https://openkim.org}{OpenKIM}.  These results are obtained using programs
for computing material properties (KIM Tests and KIM Test Drivers) that were
contributed to OpenKIM.  In order to give credit to Test developers, the
number of times results from these programs are queried is tracked.  No other
information about the nature of the query or its source is recorded.
\end{sphinxadmonition}


\subsection{Accessing KIM Model Parameters from LAMMPS (\sphinxstyleemphasis{kim param})}
\label{\detokenize{kim_commands:accessing-kim-model-parameters-from-lammps-kim-param}}\label{\detokenize{kim_commands:param}}
\sphinxAtStartPar
All IMs are functional forms containing a set of parameters.  These parameters’
values are typically selected to best reproduce a training set of quantum
mechanical calculations or available experimental data.  For example, a
Lennard\sphinxhyphen{}Jones potential intended to model argon might have the values of its two
parameters, epsilon, and sigma, fit to the dimer dissociation energy or
thermodynamic properties at a critical point of the phase diagram.

\sphinxAtStartPar
Normally a user employing an IM should not modify its parameters since, as noted
above, these are selected to reproduce material properties.  However, there are
cases where accessing and modifying IM parameters is desired, such as for
assessing uncertainty, fitting an IM, or working with an ensemble of IMs.  As
explained {\hyperref[\detokenize{kim_commands:im-types}]{\sphinxcrossref{\DUrole{std,std-ref}{above}}}}, IMs archived in OpenKIM are either Portable
Models (PMs) or Simulator Models (SMs).  KIM PMs are complete independent
implementations of an IM, whereas KIM SMs are wrappers to an IM implemented
within LAMMPS.  Two different mechanisms are provided for accessing IM
parameters in these two cases:
\begin{itemize}
\item {} 
\sphinxAtStartPar
For a KIM PM, the \sphinxstyleemphasis{kim param} command can be used to \sphinxstyleemphasis{get} and \sphinxstyleemphasis{set} the
values of the PM’s parameters as explained below.

\item {} 
\sphinxAtStartPar
For a KIM SM, the user should consult the documentation page for the specific
IM and follow instructions there for how to modify its parameters (if
possible).

\end{itemize}

\sphinxAtStartPar
The \sphinxstyleemphasis{kim param get} and \sphinxstyleemphasis{kim param set} commands provide an interface to access
and change the parameters of a KIM PM that “publishes” its parameters and makes
them publicly available (see the
\sphinxhref{https://kim-api.readthedocs.io/en/devel/features.html}{KIM API documentation}
for details).

\begin{sphinxadmonition}{note}{Note:}
\sphinxAtStartPar
The \sphinxstyleemphasis{kim param set/get} command must be preceded by a \sphinxstyleemphasis{kim interactions}
command (or alternatively by a \sphinxstyleemphasis{pair\_style kim} and \sphinxstyleemphasis{pair\_coeff} commands).
The \sphinxstyleemphasis{kim param set} command may be used wherever a \sphinxstyleemphasis{pair\_coeff} command may
occur.
\end{sphinxadmonition}


\subsubsection{Syntax}
\label{\detokenize{kim_commands:id13}}
\begin{sphinxVerbatim}[commandchars=\\\{\}]
\PYG{k}{kim }\PYG{l+s+s1}{param}\PYG{+w}{ }\PYG{n}{get}\PYG{+w}{ }\PYG{n}{param\PYGZus{}name}\PYG{+w}{ }\PYG{n}{index\PYGZus{}range}\PYG{+w}{ }\PYG{k}{variable }\PYG{n+nv+nvIdentifier}{formatarg}
\PYG{k}{kim }\PYG{l+s+s1}{param}\PYG{+w}{ }\PYG{n}{set}\PYG{+w}{ }\PYG{n}{param\PYGZus{}name}\PYG{+w}{ }\PYG{n}{index\PYGZus{}range}\PYG{+w}{ }\PYG{n}{values}
\end{sphinxVerbatim}
\phantomsection\label{\detokenize{kim_commands:formatarg-options}}\begin{itemize}
\item {} 
\sphinxAtStartPar
param\_name = name of a KIM portable model parameter (which is published by the
PM and available for access). The specific string used to identify a parameter
is defined by the PM.  For example, for the
\sphinxhref{https://openkim.org/id/SW\_StillingerWeber\_1985\_Si\_\_MO\_405512056662\_005}{Stillinger\sphinxhyphen{}Weber (SW) potential in OpenKIM},
the parameter names are \sphinxstyleemphasis{A, B, p, q, sigma, gamma, cutoff, lambda, costheta0}

\item {} 
\sphinxAtStartPar
index\_range = KIM portable model parameter index range (an integer for a
single element, or pair of integers separated by a colon for a range of
elements)

\item {} 
\sphinxAtStartPar
variable(s) = single name or list of names of (string style) LAMMPS
variable(s) where a query result or parameter get result is stored.  Variables
that do not exist will be created by the command

\item {} 
\sphinxAtStartPar
formatarg = \sphinxstyleemphasis{list} or \sphinxstyleemphasis{split} or \sphinxstyleemphasis{explicit} (optional)
\begin{sphinxalltt}
\sphinxstyleemphasis{list} = returns a single string with a list of space separated values
       (e.g. "1.0 2.0 3.0"), which is placed in a LAMMPS variable as
       defined by the \sphinxstyleemphasis{variable} argument
\sphinxstyleemphasis{split} = returns the values separately in new variables with names based
       on the prefix specified in \sphinxstyleemphasis{variable} and a number appended to
       indicate which element in the list of values is in the variable
\sphinxstyleemphasis{explicit} = returns the values separately in one more more variable names
       provided as arguments that precede \sphinxstyleemphasis{formatarg} (default)
\end{sphinxalltt}

\item {} 
\sphinxAtStartPar
values = new value(s) to replace the current value(s) of a KIM portable model
parameter

\end{itemize}

\begin{sphinxadmonition}{note}{Note:}
\sphinxAtStartPar
The list of all the parameters that a PM exposes for access/mutation are
automatically written to the lammps log file when \sphinxstyleemphasis{kim init} is called.
\end{sphinxadmonition}

\sphinxAtStartPar
Each published parameter of a KIM PM takes the form of an array of numerical
values.  The array can contain one element for a single\sphinxhyphen{}valued parameter, or a
set of values.  For example, the
\sphinxhref{https://openkim.org/id/SW\_ZhouWardMartin\_2013\_CdTeZnSeHgS\_\_MO\_503261197030\_002}{multispecies SW potential for the Zn\sphinxhyphen{}Cd\sphinxhyphen{}Hg\sphinxhyphen{}S\sphinxhyphen{}Se\sphinxhyphen{}Te system}
has the same parameter names as the
\sphinxhref{https://openkim.org/id/SW\_StillingerWeber\_1985\_Si\_\_MO\_405512056662\_005}{single\sphinxhyphen{}species SW potential},
but each parameter array contains 21 entries that correspond to the parameter
values used for each pairwise combination of the model’s six supported species
(this model does not have parameters specific to individual ternary combinations
of its supported species).

\sphinxAtStartPar
The \sphinxstyleemphasis{index\_range} argument may either be an integer referring to a specific
element within the array associated with the parameter specified by
\sphinxstyleemphasis{param\_name}, or a pair of integers separated by a colon that refer to a slice
of this array.  In both cases, one\sphinxhyphen{}based indexing is used to refer to the
entries of the array.

\sphinxAtStartPar
The result of a \sphinxstyleemphasis{get} operation for a specific \sphinxstyleemphasis{index\_range} is stored in
one or more {\hyperref[\detokenize{variable::doc}]{\sphinxcrossref{\DUrole{doc}{LAMMPS string style variables}}}} as determined by the
optional \sphinxstyleemphasis{formatarg} argument {\hyperref[\detokenize{kim_commands:formatarg-options}]{\sphinxcrossref{\DUrole{std,std-ref}{documented above.}}}} If
not specified, the default for \sphinxstyleemphasis{formatarg} is “explicit” for the \sphinxstyleemphasis{kim param}
command.

\sphinxAtStartPar
For the case where the result is an array with multiple values (i.e.
\sphinxstyleemphasis{index\_range} contains a range), the optional “split” or “explicit” \sphinxstyleemphasis{formatarg}
keywords can be used to separate the results into multiple variables; see the
examples below.  Multiple parameters can be retrieved with a single call to
\sphinxstyleemphasis{kim param get} by repeating the argument list following \sphinxstyleemphasis{get}.

\sphinxAtStartPar
For a \sphinxstyleemphasis{set} operation, the \sphinxstyleemphasis{values} argument contains the new value(s) for the
element(s) of the parameter specified by \sphinxstyleemphasis{index\_range}.  For the case where
multiple values are being set, \sphinxstyleemphasis{values} contains a set of values separated by
spaces.  Multiple parameters can be set with a single call to \sphinxstyleemphasis{kim param set} by
repeating the argument list following \sphinxstyleemphasis{set}.


\subsubsection{\sphinxstyleemphasis{kim param} Usage Examples and Further Clarifications}
\label{\detokenize{kim_commands:kim-param-usage-examples-and-further-clarifications}}
\sphinxAtStartPar
Examples of getting and setting KIM PM parameters with further clarifications
are provided below.

\sphinxAtStartPar
\sphinxstylestrong{Getting a scalar parameter}

\begin{sphinxVerbatim}[commandchars=\\\{\}]
\PYG{k}{kim }\PYG{l+s+s1}{init}\PYG{+w}{ }\PYG{n}{SW\PYGZus{}StillingerWeber\PYGZus{}1985\PYGZus{}Si\PYGZus{}\PYGZus{}MO\PYGZus{}405512056662\PYGZus{}005}\PYG{+w}{ }\PYG{n}{metal}
\PYG{n}{...}
\PYG{k}{kim }\PYG{l+s+s1}{interactions}\PYG{+w}{ }\PYG{n}{Si}
\PYG{k}{kim }\PYG{l+s+s1}{param}\PYG{+w}{ }\PYG{n}{get}\PYG{+w}{ }\PYG{n}{A}\PYG{+w}{ }\PYG{l+m}{1}\PYG{+w}{ }\PYG{n}{VARA}
\end{sphinxVerbatim}

\sphinxAtStartPar
or

\begin{sphinxVerbatim}[commandchars=\\\{\}]
\PYG{n}{...}
\PYG{k}{pair\PYGZus{}style}\PYG{+w}{ }\PYG{k}{kim }\PYG{l+s+s1}{SW\PYGZus{}StillingerWeber\PYGZus{}1985\PYGZus{}Si\PYGZus{}\PYGZus{}MO\PYGZus{}405512056662\PYGZus{}005}
\PYG{k}{pair\PYGZus{}coeff}\PYG{+w}{ }\PYG{o}{*}\PYG{+w}{ }\PYG{o}{*}\PYG{+w}{ }\PYG{n}{Si}
\PYG{k}{kim }\PYG{l+s+s1}{param}\PYG{+w}{ }\PYG{n}{get}\PYG{+w}{ }\PYG{n}{A}\PYG{+w}{ }\PYG{l+m}{1}\PYG{+w}{ }\PYG{n}{VARA}
\end{sphinxVerbatim}

\sphinxAtStartPar
In these cases, the value of the SW \sphinxstyleemphasis{A} parameter is retrieved and placed in the
LAMMPS variable \sphinxstyleemphasis{VARA}.  The variable \sphinxstyleemphasis{VARA} can be used in the remainder of
the input script in the same manner as any other LAMMPS variable.

\sphinxAtStartPar
\sphinxstylestrong{Getting multiple scalar parameters with a single call}

\begin{sphinxVerbatim}[commandchars=\\\{\}]
\PYG{n}{...}
\PYG{k}{kim }\PYG{l+s+s1}{interactions}\PYG{+w}{ }\PYG{n}{Si}
\PYG{k}{kim }\PYG{l+s+s1}{param}\PYG{+w}{ }\PYG{n}{get}\PYG{+w}{ }\PYG{n}{A}\PYG{+w}{ }\PYG{l+m}{1}\PYG{+w}{ }\PYG{n}{VARA}\PYG{+w}{ }\PYG{n}{B}\PYG{+w}{ }\PYG{l+m}{1}\PYG{+w}{ }\PYG{n}{VARB}
\end{sphinxVerbatim}

\sphinxAtStartPar
In this example, it is shown how to retrieve the \sphinxstyleemphasis{A} and \sphinxstyleemphasis{B} parameters of the
SW potential and store them in the LAMMPS variables \sphinxstyleemphasis{VARA} and \sphinxstyleemphasis{VARB}.

\sphinxAtStartPar
\sphinxstylestrong{Getting a range of values from a parameter}

\sphinxAtStartPar
There are several options when getting a range of values from a parameter
determined by the \sphinxstyleemphasis{formatarg} argument.

\begin{sphinxVerbatim}[commandchars=\\\{\}]
\PYG{k}{kim }\PYG{l+s+s1}{init}\PYG{+w}{ }\PYG{n}{SW\PYGZus{}ZhouWardMartin\PYGZus{}2013\PYGZus{}CdTeZnSeHgS\PYGZus{}\PYGZus{}MO\PYGZus{}503261197030\PYGZus{}002}\PYG{+w}{ }\PYG{n}{metal}
\PYG{n}{...}
\PYG{k}{kim }\PYG{l+s+s1}{interactions}\PYG{+w}{ }\PYG{n}{Te}\PYG{+w}{ }\PYG{n}{Zn}\PYG{+w}{ }\PYG{n}{Se}
\PYG{k}{kim }\PYG{l+s+s1}{param}\PYG{+w}{ }\PYG{n}{get}\PYG{+w}{ }\PYG{n}{lambda}\PYG{+w}{ }\PYG{l+m}{7:9}\PYG{+w}{ }\PYG{n}{LAM\PYGZus{}TeTe}\PYG{+w}{ }\PYG{n}{LAM\PYGZus{}TeZn}\PYG{+w}{ }\PYG{n}{LAM\PYGZus{}TeSe}
\end{sphinxVerbatim}

\sphinxAtStartPar
In this case, \sphinxstyleemphasis{formatarg} is not specified and therefore the default
“explicit” mode is used. (The behavior would be the same if the word
\sphinxstyleemphasis{explicit} were added after \sphinxstyleemphasis{LAM\_TeSe}.) Elements 7, 8 and 9 of parameter
lambda retrieved by the \sphinxstyleemphasis{get} operation are placed in the LAMMPS variables
\sphinxstyleemphasis{LAM\_TeTe}, \sphinxstyleemphasis{LAM\_TeZn} and \sphinxstyleemphasis{LAM\_TeSe}, respectively.

\begin{sphinxadmonition}{note}{Note:}
\sphinxAtStartPar
In the above example, elements 7\sphinxhyphen{}9 of the lambda parameter correspond to
Te\sphinxhyphen{}Te, Te\sphinxhyphen{}Zm and Te\sphinxhyphen{}Se interactions.  This can be determined by visiting the
\sphinxhref{https://openkim.org/id/SW\_ZhouWardMartin\_2013\_CdTeZnSeHgS\_\_MO\_503261197030\_002}{model page for the specified potential}
and looking at its parameter file linked to at the bottom of the page (file
with .param ending) and consulting the README documentation provided with the
driver for the PM being used.  A link to the driver is provided at the top of
the model page.
\end{sphinxadmonition}

\begin{sphinxVerbatim}[commandchars=\\\{\}]
\PYG{n}{...}
\PYG{k}{kim      }\PYG{l+s+s1}{interactions}\PYG{+w}{ }\PYG{n}{Te}\PYG{+w}{ }\PYG{n}{Zn}\PYG{+w}{ }\PYG{n}{Se}
\PYG{k}{kim      }\PYG{l+s+s1}{param}\PYG{+w}{ }\PYG{n}{get}\PYG{+w}{ }\PYG{n}{lambda}\PYG{+w}{ }\PYG{l+m}{15:17}\PYG{+w}{ }\PYG{n}{LAMS}\PYG{+w}{ }\PYG{n}{list}
\PYG{k}{variable }\PYG{n+nv+nvIdentifier}{LAM\PYGZus{}VALUE}\PYG{+w}{ }\PYG{n}{index}\PYG{+w}{ }\PYG{n+nv}{\PYGZdl{}\PYGZob{}}\PYG{n+nv}{LAMS}\PYG{n+nv}{\PYGZcb{}}
\PYG{k}{label    }\PYG{l+s+sc}{loop\PYGZus{}on\PYGZus{}lambda}
\PYG{n}{...}
\PYG{n}{...}\PYG{+w}{      }\PYG{n}{do}\PYG{+w}{ }\PYG{n}{something}\PYG{+w}{ }\PYG{n}{with}\PYG{+w}{ }\PYG{n}{the}\PYG{+w}{ }\PYG{n}{current}\PYG{+w}{ }\PYG{n}{value}\PYG{+w}{ }\PYG{n}{of}\PYG{+w}{ }\PYG{n}{lambda}
\PYG{n}{...}
\PYG{k}{next     }\PYG{n+nv+nvIdentifier}{LAM\PYGZus{}VALUE}
\PYG{k}{jump     }\PYG{l+s+sc}{SELF}\PYG{+w}{ }\PYG{n}{loop\PYGZus{}on\PYGZus{}lambda}
\end{sphinxVerbatim}

\sphinxAtStartPar
In this case, the “list” mode of \sphinxstyleemphasis{formatarg} is used.  The result of the \sphinxstyleemphasis{get}
operation is stored in the LAMMPS variable \sphinxstyleemphasis{LAMS} as a string containing the
three retrieved values separated by spaces, e.g “1.0 2.0 3.0”.  This can be used
in LAMMPS with an \sphinxstyleemphasis{index} variable to access the values one at a time within a
loop as shown in the example.  At each iteration of the loop \sphinxstyleemphasis{LAM\_VALUE}
contains the current value of lambda.

\begin{sphinxVerbatim}[commandchars=\\\{\}]
\PYG{n}{...}
\PYG{k}{kim }\PYG{l+s+s1}{interactions}\PYG{+w}{ }\PYG{n}{Te}\PYG{+w}{ }\PYG{n}{Zn}\PYG{+w}{ }\PYG{n}{Se}
\PYG{k}{kim }\PYG{l+s+s1}{param}\PYG{+w}{ }\PYG{n}{get}\PYG{+w}{ }\PYG{n}{lambda}\PYG{+w}{ }\PYG{l+m}{15:17}\PYG{+w}{ }\PYG{n}{LAM}\PYG{+w}{ }\PYG{n}{split}
\end{sphinxVerbatim}

\sphinxAtStartPar
In this case, the “split” mode of \sphinxstyleemphasis{formatarg} is used.  The three values
retrieved by the \sphinxstyleemphasis{get} operation are stored in the three LAMMPS variables
\sphinxstyleemphasis{LAM\_15}, \sphinxstyleemphasis{LAM\_16} and \sphinxstyleemphasis{LAM\_17}.  The provided name “LAM” is used as prefix
and the location in the lambda array is appended to create the variable names.

\sphinxAtStartPar
\sphinxstylestrong{Setting a scalar parameter}

\begin{sphinxVerbatim}[commandchars=\\\{\}]
\PYG{k}{kim }\PYG{l+s+s1}{init}\PYG{+w}{ }\PYG{n}{SW\PYGZus{}StillingerWeber\PYGZus{}1985\PYGZus{}Si\PYGZus{}\PYGZus{}MO\PYGZus{}405512056662\PYGZus{}005}\PYG{+w}{ }\PYG{n}{metal}
\PYG{n}{...}
\PYG{k}{kim }\PYG{l+s+s1}{interactions}\PYG{+w}{ }\PYG{n}{Si}
\PYG{k}{kim }\PYG{l+s+s1}{param}\PYG{+w}{ }\PYG{n}{set}\PYG{+w}{ }\PYG{n}{gamma}\PYG{+w}{ }\PYG{l+m}{1}\PYG{+w}{ }\PYG{l+m}{2.6}
\end{sphinxVerbatim}

\sphinxAtStartPar
Here, the SW potential’s gamma parameter is set to 2.6.  Note that the \sphinxstyleemphasis{get}
and \sphinxstyleemphasis{set} commands work together, so that a \sphinxstyleemphasis{get} following a \sphinxstyleemphasis{set} operation
will return the new value that was set.  For example,

\begin{sphinxVerbatim}[commandchars=\\\{\}]
\PYG{n}{...}
\PYG{k}{kim }\PYG{l+s+s1}{interactions}\PYG{+w}{ }\PYG{n}{Si}
\PYG{k}{kim }\PYG{l+s+s1}{param}\PYG{+w}{ }\PYG{n}{get}\PYG{+w}{ }\PYG{n}{gamma}\PYG{+w}{ }\PYG{l+m}{1}\PYG{+w}{ }\PYG{n}{ORIG\PYGZus{}GAMMA}
\PYG{k}{kim }\PYG{l+s+s1}{param}\PYG{+w}{ }\PYG{n}{set}\PYG{+w}{ }\PYG{n}{gamma}\PYG{+w}{ }\PYG{l+m}{1}\PYG{+w}{ }\PYG{l+m}{2.6}
\PYG{k}{kim }\PYG{l+s+s1}{param}\PYG{+w}{ }\PYG{n}{get}\PYG{+w}{ }\PYG{n}{gamma}\PYG{+w}{ }\PYG{l+m}{1}\PYG{+w}{ }\PYG{n}{NEW\PYGZus{}GAMMA}
\PYG{n}{...}
\PYG{k}{print}\PYG{+w}{ }\PYG{l+s}{\PYGZdq{}}\PYG{l+s}{original gamma = \PYGZdl{}\PYGZob{}ORIG\PYGZus{}GAMMA\PYGZcb{}, new gamma = \PYGZdl{}\PYGZob{}NEW\PYGZus{}GAMMA\PYGZcb{}}\PYG{l+s}{\PYGZdq{}}
\end{sphinxVerbatim}

\sphinxAtStartPar
Here, \sphinxstyleemphasis{ORIG\_GAMMA} will contain the original gamma value for the SW potential,
while \sphinxstyleemphasis{NEW\_GAMMA} will contain the value 2.6.

\sphinxAtStartPar
\sphinxstylestrong{Setting multiple scalar parameters with a single call}

\begin{sphinxVerbatim}[commandchars=\\\{\}]
\PYG{k}{kim      }\PYG{l+s+s1}{init}\PYG{+w}{ }\PYG{n}{SW\PYGZus{}ZhouWardMartin\PYGZus{}2013\PYGZus{}CdTeZnSeHgS\PYGZus{}\PYGZus{}MO\PYGZus{}503261197030\PYGZus{}002}\PYG{+w}{ }\PYG{n}{metal}
\PYG{n}{...}
\PYG{k}{kim      }\PYG{l+s+s1}{interactions}\PYG{+w}{ }\PYG{n}{Cd}\PYG{+w}{ }\PYG{n}{Te}
\PYG{k}{variable }\PYG{n+nv+nvIdentifier}{VARG}\PYG{+w}{ }\PYG{n}{equal}\PYG{+w}{ }\PYG{l+m}{2.6}
\PYG{k}{variable }\PYG{n+nv+nvIdentifier}{VARS}\PYG{+w}{ }\PYG{n}{equal}\PYG{+w}{ }\PYG{l+m}{2.0951}
\PYG{k}{kim      }\PYG{l+s+s1}{param}\PYG{+w}{ }\PYG{n}{set}\PYG{+w}{ }\PYG{n}{gamma}\PYG{+w}{ }\PYG{l+m}{1}\PYG{+w}{ }\PYG{n+nv}{\PYGZdl{}\PYGZob{}}\PYG{n+nv}{VARG}\PYG{n+nv}{\PYGZcb{}}\PYG{+w}{ }\PYG{n}{sigma}\PYG{+w}{ }\PYG{l+m}{3}\PYG{+w}{ }\PYG{n+nv}{\PYGZdl{}\PYGZob{}}\PYG{n+nv}{VARS}\PYG{n+nv}{\PYGZcb{}}
\end{sphinxVerbatim}

\sphinxAtStartPar
In this case, the first element of the \sphinxstyleemphasis{gamma} parameter and third element of
the \sphinxstyleemphasis{sigma} parameter are set to 2.6 and 2.0951, respectively.  This example
also shows how LAMMPS variables can be used when setting parameters.

\sphinxAtStartPar
\sphinxstylestrong{Setting a range of values of a parameter}

\begin{sphinxVerbatim}[commandchars=\\\{\}]
\PYG{k}{kim }\PYG{l+s+s1}{init}\PYG{+w}{ }\PYG{n}{SW\PYGZus{}ZhouWardMartin\PYGZus{}2013\PYGZus{}CdTeZnSeHgS\PYGZus{}\PYGZus{}MO\PYGZus{}503261197030\PYGZus{}002}\PYG{+w}{ }\PYG{n}{metal}
\PYG{n}{...}
\PYG{k}{kim }\PYG{l+s+s1}{interactions}\PYG{+w}{ }\PYG{n}{Cd}\PYG{+w}{ }\PYG{n}{Te}\PYG{+w}{ }\PYG{n}{Zn}\PYG{+w}{ }\PYG{n}{Se}\PYG{+w}{ }\PYG{n}{Hg}\PYG{+w}{ }\PYG{n}{S}
\PYG{k}{kim }\PYG{l+s+s1}{param}\PYG{+w}{ }\PYG{n}{set}\PYG{+w}{ }\PYG{n}{sigma}\PYG{+w}{ }\PYG{l+m}{2:6}\PYG{+w}{ }\PYG{l+m}{2.35214}\PYG{+w}{ }\PYG{l+m}{2.23869}\PYG{+w}{ }\PYG{l+m}{2.04516}\PYG{+w}{ }\PYG{l+m}{2.43269}\PYG{+w}{ }\PYG{l+m}{1.80415}
\end{sphinxVerbatim}

\sphinxAtStartPar
In this case, elements 2 through 6 of the parameter \sphinxstyleemphasis{sigma} are set to the
values 2.35214, 2.23869, 2.04516, 2.43269 and 1.80415 in order.


\subsection{Writing material properties in standard KIM Property Instance format (\sphinxstyleemphasis{kim property})}
\label{\detokenize{kim_commands:writing-material-properties-in-standard-kim-property-instance-format-kim-property}}\label{\detokenize{kim_commands:property}}
\sphinxAtStartPar
The OpenKIM system includes a
collection of Tests (material property calculation codes), Models (interatomic
potentials), Predictions, and Reference Data (DFT or experiments).
Specifically, a KIM Test is a computation that when coupled with a KIM Model
generates the prediction of that model for a specific material property
rigorously defined by a KIM Property Definition (see the
\sphinxhref{https://openkim.org/doc/schema/properties-framework/}{KIM Properties Framework}
for further details).  A prediction of a material property for a given model is
a specific numerical realization of a property definition, referred to as a
“Property Instance.” The objective of the \sphinxstyleemphasis{kim property} command is to make it
easy to output material properties in a standardized, machine readable, format
that can be easily ingested by other programs.  Additionally, it aims to make it
as easy as possible to convert a LAMMPS script that computes a material property
into a KIM Test that can then be uploaded to
\sphinxhref{https://openkim.org}{openkim.org}

\sphinxAtStartPar
A developer interested in creating a KIM Test using a LAMMPS script should first
determine whether a property definition that applies to their calculation
already exists in OpenKIM by searching the \sphinxhref{https://openkim.org/properties}{properties page}.  If none exists, it is possible to use a
locally defined property definition contained in a file until it can be
uploaded to the official repository (see below).  Once one or more applicable
property definitions have been identified, the \sphinxstyleemphasis{kim property create},
\sphinxstyleemphasis{kim property modify}, \sphinxstyleemphasis{kim property remove}, and \sphinxstyleemphasis{kim property destroy},
commands provide an interface to create, set, modify, remove, and destroy
instances of them within a LAMMPS script.


\subsubsection{Syntax}
\label{\detokenize{kim_commands:id14}}
\begin{sphinxVerbatim}[commandchars=\\\{\}]
\PYG{k}{kim }\PYG{l+s+s1}{property}\PYG{+w}{ }\PYG{n}{create}\PYG{+w}{  }\PYG{n}{instance\PYGZus{}id}\PYG{+w}{ }\PYG{n}{property\PYGZus{}id}
\PYG{k}{kim }\PYG{l+s+s1}{property}\PYG{+w}{ }\PYG{n}{modify}\PYG{+w}{  }\PYG{n}{instance\PYGZus{}id}\PYG{+w}{ }\PYG{n}{key}\PYG{+w}{ }\PYG{n}{key\PYGZus{}name}\PYG{+w}{ }\PYG{n}{key\PYGZus{}name\PYGZus{}key}\PYG{+w}{ }\PYG{n}{key\PYGZus{}name\PYGZus{}value}
\PYG{k}{kim }\PYG{l+s+s1}{property}\PYG{+w}{ }\PYG{n}{remove}\PYG{+w}{  }\PYG{n}{instance\PYGZus{}id}\PYG{+w}{ }\PYG{n}{key}\PYG{+w}{ }\PYG{n}{key\PYGZus{}name}
\PYG{k}{kim }\PYG{l+s+s1}{property}\PYG{+w}{ }\PYG{n}{destroy}\PYG{+w}{ }\PYG{n}{instance\PYGZus{}id}
\PYG{k}{kim }\PYG{l+s+s1}{property}\PYG{+w}{ }\PYG{k}{dump    }\PYG{n+nv+nvIdentifier}{file}
\end{sphinxVerbatim}
\begin{itemize}
\item {} 
\sphinxAtStartPar
instance\_id = a positive integer identifying the KIM property instance;
(note that the results file can contain multiple property instances)

\item {} 
\sphinxAtStartPar
property\_id = identifier of a
\sphinxhref{https://openkim.org/properties}{KIM Property Definition}, which can be (1)
a property short name, (2) the full unique ID of the property (including the
contributor and date), (3) a file name corresponding to a local property
definition file

\item {} 
\sphinxAtStartPar
key\_name = one of the keys belonging to the specified KIM property definition

\item {} 
\sphinxAtStartPar
key\_name\_key = a key belonging to a key\sphinxhyphen{}value pair (standardized in the
\sphinxhref{https://openkim.org/doc/schema/properties-framework}{KIM Properties Framework})

\item {} 
\sphinxAtStartPar
key\_name\_value = value to be associated with a key\_name\_key in a key\sphinxhyphen{}value
pair

\item {} 
\sphinxAtStartPar
file = name of a file to write the currently defined set of KIM property
instances to

\end{itemize}

\sphinxAtStartPar
Examples of each of the three \sphinxstyleemphasis{property\_id} cases are shown below,

\begin{sphinxVerbatim}[commandchars=\\\{\}]
\PYG{k}{kim }\PYG{l+s+s1}{property}\PYG{+w}{ }\PYG{n}{create}\PYG{+w}{ }\PYG{l+m}{1}\PYG{+w}{ }\PYG{n}{atomic}\PYG{o}{\PYGZhy{}}\PYG{n}{mass}
\PYG{k}{kim }\PYG{l+s+s1}{property}\PYG{+w}{ }\PYG{n}{create}\PYG{+w}{ }\PYG{l+m}{2}\PYG{+w}{ }\PYG{n}{cohesive}\PYG{o}{\PYGZhy{}}\PYG{n}{energy}\PYG{o}{\PYGZhy{}}\PYG{n}{relation}\PYG{o}{\PYGZhy{}}\PYG{n}{cubic}\PYG{o}{\PYGZhy{}}\PYG{n}{crystal}
\end{sphinxVerbatim}

\begin{sphinxVerbatim}[commandchars=\\\{\}]
\PYG{k}{kim }\PYG{l+s+s1}{property}\PYG{+w}{ }\PYG{n}{create}\PYG{+w}{ }\PYG{l+m}{1}\PYG{+w}{ }\PYG{n}{tag}:brunnels@noreply.openkim.org,2016\PYGZhy{}05\PYGZhy{}11:property/atomic\PYGZhy{}mass
\PYG{k}{kim }\PYG{l+s+s1}{property}\PYG{+w}{ }\PYG{n}{create}\PYG{+w}{ }\PYG{l+m}{2}\PYG{+w}{ }\PYG{n}{tag}:staff@noreply.openkim.org,2014\PYGZhy{}04\PYGZhy{}15:property/cohesive\PYGZhy{}energy\PYGZhy{}relation\PYGZhy{}cubic\PYGZhy{}crystal
\end{sphinxVerbatim}

\begin{sphinxVerbatim}[commandchars=\\\{\}]
\PYG{k}{kim }\PYG{l+s+s1}{property}\PYG{+w}{ }\PYG{n}{create}\PYG{+w}{ }\PYG{l+m}{1}\PYG{+w}{ }\PYG{n}{new}\PYG{o}{\PYGZhy{}}\PYG{n}{property.edn}
\PYG{k}{kim }\PYG{l+s+s1}{property}\PYG{+w}{ }\PYG{n}{create}\PYG{+w}{ }\PYG{l+m}{2}\PYG{+w}{ }\PYG{o}{/}\PYG{n}{home}\PYG{o}{/}\PYG{n}{mary}\PYG{o}{/}\PYG{n}{marys}\PYG{o}{\PYGZhy{}}\PYG{n}{kim}\PYG{o}{\PYGZhy{}}\PYG{n}{properties}\PYG{o}{/}\PYG{n}{dissociation}\PYG{o}{\PYGZhy{}}\PYG{n}{energy.edn}
\end{sphinxVerbatim}

\sphinxAtStartPar
In the last example, “new\sphinxhyphen{}property.edn” and
“/home/mary/marys\sphinxhyphen{}kim\sphinxhyphen{}properties/dissociation\sphinxhyphen{}energy.edn” are the names of files
that contain user\sphinxhyphen{}defined (local) property definitions.

\sphinxAtStartPar
A KIM property instance takes the form of a “map”, i.e. a set of key\sphinxhyphen{}value
pairs akin to Perl’s hash, Python’s dictionary, or Java’s Hashtable.  It
consists of a set of property key names, each of which is referred to here by
the \sphinxstyleemphasis{key\_name} argument, that are defined as part of the relevant KIM Property
Definition and include only lowercase alphanumeric characters and dashes.  The
value paired with each property key is itself a map whose possible keys are
defined as part of the \sphinxhref{https://openkim.org/doc/schema/properties-framework}{KIM Properties Framework}; these keys are
referred to by the \sphinxstyleemphasis{key\_name\_key} argument and their associated values by the
\sphinxstyleemphasis{key\_name\_value} argument.  These values may either be scalars or arrays,
as stipulated in the property definition.

\begin{sphinxadmonition}{note}{Note:}
\sphinxAtStartPar
Each map assigned to a \sphinxstyleemphasis{key\_name} must contain the \sphinxstyleemphasis{key\_name\_key}
“source\sphinxhyphen{}value” and an associated \sphinxstyleemphasis{key\_name\_value} of the appropriate
type (as defined in the relevant KIM Property Definition).  For keys that
are defined as having physical units, the “source\sphinxhyphen{}unit” \sphinxstyleemphasis{key\_name\_key} must
also be given a string value recognized by
\sphinxhref{https://www.gnu.org/software/units/}{GNU units}.
\end{sphinxadmonition}

\sphinxAtStartPar
Once a \sphinxstyleemphasis{kim property create} command has been given to instantiate a property
instance, maps associated with the property’s keys can be edited using the
\sphinxstyleemphasis{kim property modify} command.  In using this command, the special keyword
“key” should be given, followed by the property key name and the key\sphinxhyphen{}value pair
in the map associated with the key that is to be set.  For example, the
\sphinxhref{https://openkim.org/properties/show/2016-05-11/brunnels@noreply.openkim.org/atomic-mass}{atomic\sphinxhyphen{}mass}
property definition consists of two property keys named “mass” and “species.”
An instance of this property could be created like so:

\begin{sphinxVerbatim}[commandchars=\\\{\}]
\PYG{k}{kim }\PYG{l+s+s1}{property}\PYG{+w}{ }\PYG{n}{create}\PYG{+w}{ }\PYG{l+m}{1}\PYG{+w}{ }\PYG{n}{atomic}\PYG{o}{\PYGZhy{}}\PYG{n}{mass}
\PYG{k}{kim }\PYG{l+s+s1}{property}\PYG{+w}{ }\PYG{n}{modify}\PYG{+w}{ }\PYG{l+m}{1}\PYG{+w}{ }\PYG{n}{key}\PYG{+w}{ }\PYG{n}{species}\PYG{+w}{ }\PYG{n}{source}\PYG{o}{\PYGZhy{}}\PYG{n}{value}\PYG{+w}{ }\PYG{n}{Al}
\PYG{k}{kim }\PYG{l+s+s1}{property}\PYG{+w}{ }\PYG{n}{modify}\PYG{+w}{ }\PYG{l+m}{1}\PYG{+w}{ }\PYG{n}{key}\PYG{+w}{ }\PYG{n}{mass}\PYG{+w}{    }\PYG{n}{source}\PYG{o}{\PYGZhy{}}\PYG{n}{value}\PYG{+w}{ }\PYG{l+m}{26.98154}
\PYG{k}{kim }\PYG{l+s+s1}{property}\PYG{+w}{ }\PYG{n}{modify}\PYG{+w}{ }\PYG{l+m}{1}\PYG{+w}{ }\PYG{n}{key}\PYG{+w}{ }\PYG{n}{mass}\PYG{+w}{    }\PYG{n}{source}\PYG{o}{\PYGZhy{}}\PYG{n}{unit}\PYG{+w}{ }\PYG{n}{amu}
\end{sphinxVerbatim}

\sphinxAtStartPar
or, equivalently,

\begin{sphinxVerbatim}[commandchars=\\\{\}]
\PYG{k}{kim }\PYG{l+s+s1}{property}\PYG{+w}{ }\PYG{n}{create}\PYG{+w}{ }\PYG{l+m}{1}\PYG{+w}{ }\PYG{n}{atomic}\PYG{o}{\PYGZhy{}}\PYG{n}{mass}
\PYG{k}{kim }\PYG{l+s+s1}{property}\PYG{+w}{ }\PYG{n}{modify}\PYG{+w}{ }\PYG{l+m}{1}\PYG{+w}{ }\PYG{n}{key}\PYG{+w}{ }\PYG{n}{species}\PYG{+w}{ }\PYG{n}{source}\PYG{o}{\PYGZhy{}}\PYG{n}{value}\PYG{+w}{ }\PYG{n}{Al}\PYG{+w}{       }\PYG{o}{\PYGZam{}}
\PYG{+w}{                      }\PYG{n}{key}\PYG{+w}{ }\PYG{n}{mass}\PYG{+w}{    }\PYG{n}{source}\PYG{o}{\PYGZhy{}}\PYG{n}{value}\PYG{+w}{ }\PYG{l+m}{26.98154}\PYG{+w}{ }\PYG{o}{\PYGZam{}}
\PYG{+w}{                                  }\PYG{n}{source}\PYG{o}{\PYGZhy{}}\PYG{n}{unit}\PYG{+w}{  }\PYG{n}{amu}
\end{sphinxVerbatim}


\subsubsection{\sphinxstyleemphasis{kim property} Usage Examples and Further Clarifications}
\label{\detokenize{kim_commands:kim-property-usage-examples-and-further-clarifications}}
\sphinxAtStartPar
\sphinxstylestrong{Create}

\begin{sphinxVerbatim}[commandchars=\\\{\}]
\PYG{k}{kim }\PYG{l+s+s1}{property}\PYG{+w}{ }\PYG{n}{create}\PYG{+w}{ }\PYG{n}{instance\PYGZus{}id}\PYG{+w}{ }\PYG{n}{property\PYGZus{}id}
\end{sphinxVerbatim}

\sphinxAtStartPar
The \sphinxstyleemphasis{kim property create} command takes as input a property instance ID and the
property definition name, and creates an initial empty property instance data
structure.  For example,

\begin{sphinxVerbatim}[commandchars=\\\{\}]
\PYG{k}{kim }\PYG{l+s+s1}{property}\PYG{+w}{ }\PYG{n}{create}\PYG{+w}{ }\PYG{l+m}{1}\PYG{+w}{ }\PYG{n}{atomic}\PYG{o}{\PYGZhy{}}\PYG{n}{mass}
\PYG{k}{kim }\PYG{l+s+s1}{property}\PYG{+w}{ }\PYG{n}{create}\PYG{+w}{ }\PYG{l+m}{2}\PYG{+w}{ }\PYG{n}{cohesive}\PYG{o}{\PYGZhy{}}\PYG{n}{energy}\PYG{o}{\PYGZhy{}}\PYG{n}{relation}\PYG{o}{\PYGZhy{}}\PYG{n}{cubic}\PYG{o}{\PYGZhy{}}\PYG{n}{crystal}
\end{sphinxVerbatim}

\sphinxAtStartPar
creates an empty property instance of the “atomic\sphinxhyphen{}mass” property definition
with instance ID 1 and an empty instance of the
“cohesive\sphinxhyphen{}energy\sphinxhyphen{}relation\sphinxhyphen{}cubic\sphinxhyphen{}crystal” property with ID 2.  A list of
published property definitions in OpenKIM can be found on the \sphinxhref{https://openkim.org/properties}{properties page}.

\sphinxAtStartPar
One can also provide the name of a file in the current working directory or the
path of a file containing a valid property definition.  For example,

\begin{sphinxVerbatim}[commandchars=\\\{\}]
\PYG{k}{kim }\PYG{l+s+s1}{property}\PYG{+w}{ }\PYG{n}{create}\PYG{+w}{ }\PYG{l+m}{1}\PYG{+w}{ }\PYG{n}{new}\PYG{o}{\PYGZhy{}}\PYG{n}{property.edn}
\end{sphinxVerbatim}

\sphinxAtStartPar
where “new\sphinxhyphen{}property.edn” refers to a file name containing a new property
definition that does not exist in OpenKIM.

\sphinxAtStartPar
If the \sphinxstyleemphasis{property\_id} given cannot be found in OpenKIM and no file of this name
containing a valid property definition can be found, this command will produce
an error with an appropriate message.  Calling \sphinxstyleemphasis{kim property create} with the
same instance ID multiple times will also produce an error.

\sphinxAtStartPar
\sphinxstylestrong{Modify}

\begin{sphinxVerbatim}[commandchars=\\\{\}]
\PYG{k}{kim }\PYG{l+s+s1}{property}\PYG{+w}{ }\PYG{n}{modify}\PYG{+w}{ }\PYG{n}{instance\PYGZus{}id}\PYG{+w}{ }\PYG{n}{key}\PYG{+w}{ }\PYG{n}{key\PYGZus{}name}\PYG{+w}{ }\PYG{n}{key\PYGZus{}name\PYGZus{}key}\PYG{+w}{ }\PYG{n}{key\PYGZus{}name\PYGZus{}value}
\end{sphinxVerbatim}

\sphinxAtStartPar
The \sphinxstyleemphasis{kim property modify} command incrementally builds the property instance
by receiving property definition keys along with associated arguments.  Each
\sphinxstyleemphasis{key\_name} is associated with a map containing one or more key\sphinxhyphen{}value pairs (in
the form of \sphinxstyleemphasis{key\_name\_key}\sphinxhyphen{}\sphinxstyleemphasis{key\_name\_value} pairs).  For example,

\begin{sphinxVerbatim}[commandchars=\\\{\}]
\PYG{k}{kim }\PYG{l+s+s1}{property}\PYG{+w}{ }\PYG{n}{modify}\PYG{+w}{ }\PYG{l+m}{1}\PYG{+w}{ }\PYG{n}{key}\PYG{+w}{ }\PYG{n}{species}\PYG{+w}{ }\PYG{n}{source}\PYG{o}{\PYGZhy{}}\PYG{n}{value}\PYG{+w}{ }\PYG{n}{Al}
\PYG{k}{kim }\PYG{l+s+s1}{property}\PYG{+w}{ }\PYG{n}{modify}\PYG{+w}{ }\PYG{l+m}{1}\PYG{+w}{ }\PYG{n}{key}\PYG{+w}{ }\PYG{n}{mass}\PYG{+w}{    }\PYG{n}{source}\PYG{o}{\PYGZhy{}}\PYG{n}{value}\PYG{+w}{ }\PYG{l+m}{26.98154}
\PYG{k}{kim }\PYG{l+s+s1}{property}\PYG{+w}{ }\PYG{n}{modify}\PYG{+w}{ }\PYG{l+m}{1}\PYG{+w}{ }\PYG{n}{key}\PYG{+w}{ }\PYG{n}{mass}\PYG{+w}{    }\PYG{n}{source}\PYG{o}{\PYGZhy{}}\PYG{n}{unit}\PYG{+w}{  }\PYG{n}{amu}
\end{sphinxVerbatim}

\sphinxAtStartPar
where the special keyword “key” is followed by a \sphinxstyleemphasis{key\_name} (“species” or
“mass” in the above) and one or more key\sphinxhyphen{}value pairs.  These key\sphinxhyphen{}value pairs
may continue until either another “key” keyword is given or the end of the
command line is reached.  Thus, the above could equivalently be written as

\begin{sphinxVerbatim}[commandchars=\\\{\}]
\PYG{k}{kim }\PYG{l+s+s1}{property}\PYG{+w}{ }\PYG{n}{modify}\PYG{+w}{ }\PYG{l+m}{1}\PYG{+w}{ }\PYG{n}{key}\PYG{+w}{ }\PYG{n}{species}\PYG{+w}{ }\PYG{n}{source}\PYG{o}{\PYGZhy{}}\PYG{n}{value}\PYG{+w}{ }\PYG{n}{Al}\PYG{+w}{       }\PYG{o}{\PYGZam{}}
\PYG{+w}{                      }\PYG{n}{key}\PYG{+w}{ }\PYG{n}{mass}\PYG{+w}{    }\PYG{n}{source}\PYG{o}{\PYGZhy{}}\PYG{n}{value}\PYG{+w}{ }\PYG{l+m}{26.98154}\PYG{+w}{ }\PYG{o}{\PYGZam{}}
\PYG{+w}{                      }\PYG{n}{key}\PYG{+w}{ }\PYG{n}{mass}\PYG{+w}{    }\PYG{n}{source}\PYG{o}{\PYGZhy{}}\PYG{n}{unit}\PYG{+w}{  }\PYG{n}{amu}
\end{sphinxVerbatim}

\sphinxAtStartPar
As an example of modifying multiple key\sphinxhyphen{}value pairs belonging to the map of a
single property key, the following command modifies the map of the
“cohesive\sphinxhyphen{}potential\sphinxhyphen{}energy” property key to contain the key “source\sphinxhyphen{}unit” which
is assigned a value of “eV” and the key “digits” which is assigned a value of
5,

\begin{sphinxVerbatim}[commandchars=\\\{\}]
\PYG{k}{kim }\PYG{l+s+s1}{property}\PYG{+w}{ }\PYG{n}{modify}\PYG{+w}{ }\PYG{l+m}{2}\PYG{+w}{ }\PYG{n}{key}\PYG{+w}{ }\PYG{n}{cohesive}\PYG{o}{\PYGZhy{}}\PYG{n}{potential}\PYG{o}{\PYGZhy{}}\PYG{n}{energy}\PYG{+w}{ }\PYG{n}{source}\PYG{o}{\PYGZhy{}}\PYG{n}{unit}\PYG{+w}{ }\PYG{n}{eV}\PYG{+w}{ }\PYG{n}{digits}\PYG{+w}{ }\PYG{l+m}{5}
\end{sphinxVerbatim}

\begin{sphinxadmonition}{note}{Note:}
\sphinxAtStartPar
The relevant data types of the values in the map are handled automatically
based on the specification of the key in the KIM Property Definition.  In
the example above, this means that the value “eV” will automatically be
interpreted as a string while the value 5 will be interpreted as an integer.
\end{sphinxadmonition}

\sphinxAtStartPar
The values contained in maps can either be scalars, as in all of the examples
above, or arrays depending on which is stipulated in the corresponding Property
Definition.  For one\sphinxhyphen{}dimensional arrays, a single one\sphinxhyphen{}based index must be
supplied that indicates which element of the array is to be modified.  For
multidimensional arrays, multiple indices must be given depending on the
dimensionality of the array.

\begin{sphinxadmonition}{note}{Note:}
\sphinxAtStartPar
All array indexing used by \sphinxstyleemphasis{kim property modify} is one\sphinxhyphen{}based, i.e. the
indices are enumerated 1, 2, 3, …
\end{sphinxadmonition}

\begin{sphinxadmonition}{note}{Note:}
\sphinxAtStartPar
The dimensionality of arrays are defined in the the corresponding Property
Definition.  The extent of each dimension of an array can either be a
specific finite number or indefinite and determined at run time.  If
an array has a fixed extent, attempting to modify an out\sphinxhyphen{}of\sphinxhyphen{}range index will
fail with an error message.
\end{sphinxadmonition}

\sphinxAtStartPar
For example, the “species” property key of the
\sphinxhref{https://openkim.org/properties/show/2014-04-15/staff@noreply.openkim.org/cohesive-energy-relation-cubic-crystal}{cohesive\sphinxhyphen{}energy\sphinxhyphen{}relation\sphinxhyphen{}cubic\sphinxhyphen{}crystal}
property is a one\sphinxhyphen{}dimensional array that can contain any number of entries
based on the number of atoms in the unit cell of a given cubic crystal.  To
assign an array containing the string “Al” four times to the “source\sphinxhyphen{}value” key
of the “species” property key, we can do so by issuing:

\begin{sphinxVerbatim}[commandchars=\\\{\}]
\PYG{k}{kim }\PYG{l+s+s1}{property}\PYG{+w}{ }\PYG{n}{modify}\PYG{+w}{ }\PYG{l+m}{2}\PYG{+w}{ }\PYG{n}{key}\PYG{+w}{ }\PYG{n}{species}\PYG{+w}{ }\PYG{n}{source}\PYG{o}{\PYGZhy{}}\PYG{n}{value}\PYG{+w}{ }\PYG{l+m}{1}\PYG{+w}{ }\PYG{n}{Al}
\PYG{k}{kim }\PYG{l+s+s1}{property}\PYG{+w}{ }\PYG{n}{modify}\PYG{+w}{ }\PYG{l+m}{2}\PYG{+w}{ }\PYG{n}{key}\PYG{+w}{ }\PYG{n}{species}\PYG{+w}{ }\PYG{n}{source}\PYG{o}{\PYGZhy{}}\PYG{n}{value}\PYG{+w}{ }\PYG{l+m}{2}\PYG{+w}{ }\PYG{n}{Al}
\PYG{k}{kim }\PYG{l+s+s1}{property}\PYG{+w}{ }\PYG{n}{modify}\PYG{+w}{ }\PYG{l+m}{2}\PYG{+w}{ }\PYG{n}{key}\PYG{+w}{ }\PYG{n}{species}\PYG{+w}{ }\PYG{n}{source}\PYG{o}{\PYGZhy{}}\PYG{n}{value}\PYG{+w}{ }\PYG{l+m}{3}\PYG{+w}{ }\PYG{n}{Al}
\PYG{k}{kim }\PYG{l+s+s1}{property}\PYG{+w}{ }\PYG{n}{modify}\PYG{+w}{ }\PYG{l+m}{2}\PYG{+w}{ }\PYG{n}{key}\PYG{+w}{ }\PYG{n}{species}\PYG{+w}{ }\PYG{n}{source}\PYG{o}{\PYGZhy{}}\PYG{n}{value}\PYG{+w}{ }\PYG{l+m}{4}\PYG{+w}{ }\PYG{n}{Al}
\end{sphinxVerbatim}

\begin{sphinxadmonition}{note}{Note:}
\sphinxAtStartPar
No declaration of the number of elements in this array was given;
\sphinxstyleemphasis{kim property modify} will automatically handle memory management to allow
an arbitrary number of elements to be added to the array.
\end{sphinxadmonition}

\begin{sphinxadmonition}{note}{Note:}
\sphinxAtStartPar
In the event that \sphinxstyleemphasis{kim property modify} is used to set the value of an
array index without having set the values of all lesser indices, they will
be assigned default values based on the data type associated with the key in
the map:


\begin{savenotes}\sphinxattablestart
\sphinxthistablewithglobalstyle
\centering
\begin{tabulary}{\linewidth}[t]{TT}
\sphinxtoprule
\sphinxtableatstartofbodyhook
\sphinxAtStartPar
Data type
&
\sphinxAtStartPar
Default value
\\
\sphinxhline
\sphinxAtStartPar
int
&
\sphinxAtStartPar
0
\\
\sphinxhline
\sphinxAtStartPar
float
&
\sphinxAtStartPar
0.0
\\
\sphinxhline
\sphinxAtStartPar
string
&
\sphinxAtStartPar
""
\\
\sphinxhline
\sphinxAtStartPar
file
&
\sphinxAtStartPar
""
\\
\sphinxbottomrule
\end{tabulary}
\sphinxtableafterendhook\par
\sphinxattableend\end{savenotes}

\sphinxAtStartPar
For example, doing the following:

\begin{sphinxVerbatim}[commandchars=\\\{\}]
\PYG{k}{kim }\PYG{l+s+s1}{property}\PYG{+w}{ }\PYG{n}{create}\PYG{+w}{ }\PYG{l+m}{2}\PYG{+w}{ }\PYG{n}{cohesive}\PYG{o}{\PYGZhy{}}\PYG{n}{energy}\PYG{o}{\PYGZhy{}}\PYG{n}{relation}\PYG{o}{\PYGZhy{}}\PYG{n}{cubic}\PYG{o}{\PYGZhy{}}\PYG{n}{crystal}
\PYG{k}{kim }\PYG{l+s+s1}{property}\PYG{+w}{ }\PYG{n}{modify}\PYG{+w}{ }\PYG{l+m}{2}\PYG{+w}{ }\PYG{n}{key}\PYG{+w}{ }\PYG{n}{species}\PYG{+w}{ }\PYG{n}{source}\PYG{o}{\PYGZhy{}}\PYG{n}{value}\PYG{+w}{ }\PYG{l+m}{4}\PYG{+w}{ }\PYG{n}{Al}
\end{sphinxVerbatim}

\sphinxAtStartPar
will result in the “source\sphinxhyphen{}value” key in the map for the property key
“species” being assigned the array {[}“”, “”, “”, “Al”{]}.
\end{sphinxadmonition}

\sphinxAtStartPar
For convenience, the index argument provided may refer to an inclusive range of
indices by specifying two integers separated by a colon (the first integer must
be less than or equal to the second integer, and no whitespace should be
included).  Thus, the snippet above could equivalently be written:

\begin{sphinxVerbatim}[commandchars=\\\{\}]
\PYG{k}{kim }\PYG{l+s+s1}{property}\PYG{+w}{ }\PYG{n}{modify}\PYG{+w}{ }\PYG{l+m}{2}\PYG{+w}{ }\PYG{n}{key}\PYG{+w}{ }\PYG{n}{species}\PYG{+w}{ }\PYG{n}{source}\PYG{o}{\PYGZhy{}}\PYG{n}{value}\PYG{+w}{ }\PYG{l+m}{1:4}\PYG{+w}{ }\PYG{n}{Al}\PYG{+w}{ }\PYG{n}{Al}\PYG{+w}{ }\PYG{n}{Al}\PYG{+w}{ }\PYG{n}{Al}
\end{sphinxVerbatim}

\sphinxAtStartPar
Calling this command with a non\sphinxhyphen{}positive index, e.g.
\sphinxcode{\sphinxupquote{kim property modify 2 key species source\sphinxhyphen{}value 0 Al}}, or an incorrect
number of input arguments, e.g.
\sphinxcode{\sphinxupquote{kim property modify 2 key species source\sphinxhyphen{}value 1:4 Al Al}}, will result in an
error.

\sphinxAtStartPar
As an example of modifying multidimensional arrays, consider the “basis\sphinxhyphen{}atoms”
key in the \sphinxhref{https://openkim.org/properties/show/2014-04-15/staff@noreply.openkim.org/cohesive-energy-relation-cubic-crystal}{cohesive\sphinxhyphen{}energy\sphinxhyphen{}relation\sphinxhyphen{}cubic\sphinxhyphen{}crystal}
property definition.  This is a two\sphinxhyphen{}dimensional array containing the fractional
coordinates of atoms in the unit cell of the cubic crystal.  In the case of,
e.g. a conventional fcc unit cell, the “source\sphinxhyphen{}value” key in the map associated
with this key should be assigned the following value:

\begin{sphinxVerbatim}[commandchars=\\\{\}]
[[0.0, 0.0, 0.0],
 [0.5, 0.5, 0.0],
 [0.5, 0.0, 0.5],
 [0.0, 0.5, 0.5]]
\end{sphinxVerbatim}

\sphinxAtStartPar
While each of the twelve components could be set individually, we can instead
set each row at a time using colon notation:

\begin{sphinxVerbatim}[commandchars=\\\{\}]
\PYG{k}{kim }\PYG{l+s+s1}{property}\PYG{+w}{ }\PYG{n}{modify}\PYG{+w}{ }\PYG{l+m}{2}\PYG{+w}{ }\PYG{n}{key}\PYG{+w}{ }\PYG{n}{basis}\PYG{o}{\PYGZhy{}}\PYG{n}{atom}\PYG{o}{\PYGZhy{}}\PYG{n}{coordinates}\PYG{+w}{ }\PYG{n}{source}\PYG{o}{\PYGZhy{}}\PYG{n}{value}\PYG{+w}{ }\PYG{l+m}{1}\PYG{+w}{ }\PYG{l+m}{1:3}\PYG{+w}{ }\PYG{l+m}{0.0}\PYG{+w}{ }\PYG{l+m}{0.0}\PYG{+w}{ }\PYG{l+m}{0.0}
\PYG{k}{kim }\PYG{l+s+s1}{property}\PYG{+w}{ }\PYG{n}{modify}\PYG{+w}{ }\PYG{l+m}{2}\PYG{+w}{ }\PYG{n}{key}\PYG{+w}{ }\PYG{n}{basis}\PYG{o}{\PYGZhy{}}\PYG{n}{atom}\PYG{o}{\PYGZhy{}}\PYG{n}{coordinates}\PYG{+w}{ }\PYG{n}{source}\PYG{o}{\PYGZhy{}}\PYG{n}{value}\PYG{+w}{ }\PYG{l+m}{2}\PYG{+w}{ }\PYG{l+m}{1:3}\PYG{+w}{ }\PYG{l+m}{0.5}\PYG{+w}{ }\PYG{l+m}{0.5}\PYG{+w}{ }\PYG{l+m}{0.0}
\PYG{k}{kim }\PYG{l+s+s1}{property}\PYG{+w}{ }\PYG{n}{modify}\PYG{+w}{ }\PYG{l+m}{2}\PYG{+w}{ }\PYG{n}{key}\PYG{+w}{ }\PYG{n}{basis}\PYG{o}{\PYGZhy{}}\PYG{n}{atom}\PYG{o}{\PYGZhy{}}\PYG{n}{coordinates}\PYG{+w}{ }\PYG{n}{source}\PYG{o}{\PYGZhy{}}\PYG{n}{value}\PYG{+w}{ }\PYG{l+m}{3}\PYG{+w}{ }\PYG{l+m}{1:3}\PYG{+w}{ }\PYG{l+m}{0.5}\PYG{+w}{ }\PYG{l+m}{0.0}\PYG{+w}{ }\PYG{l+m}{0.5}
\PYG{k}{kim }\PYG{l+s+s1}{property}\PYG{+w}{ }\PYG{n}{modify}\PYG{+w}{ }\PYG{l+m}{2}\PYG{+w}{ }\PYG{n}{key}\PYG{+w}{ }\PYG{n}{basis}\PYG{o}{\PYGZhy{}}\PYG{n}{atom}\PYG{o}{\PYGZhy{}}\PYG{n}{coordinates}\PYG{+w}{ }\PYG{n}{source}\PYG{o}{\PYGZhy{}}\PYG{n}{value}\PYG{+w}{ }\PYG{l+m}{4}\PYG{+w}{ }\PYG{l+m}{1:3}\PYG{+w}{ }\PYG{l+m}{0.0}\PYG{+w}{ }\PYG{l+m}{0.5}\PYG{+w}{ }\PYG{l+m}{0.5}
\end{sphinxVerbatim}

\sphinxAtStartPar
Where the first index given refers to a row and the second index refers to a
column.  We could, instead, choose to set each column at a time like so:

\begin{sphinxVerbatim}[commandchars=\\\{\}]
\PYG{k}{kim }\PYG{l+s+s1}{property}\PYG{+w}{ }\PYG{n}{modify}\PYG{+w}{ }\PYG{l+m}{2}\PYG{+w}{ }\PYG{n}{key}\PYG{+w}{ }\PYG{n}{basis}\PYG{o}{\PYGZhy{}}\PYG{n}{atom}\PYG{o}{\PYGZhy{}}\PYG{n}{coordinates}\PYG{+w}{ }\PYG{n}{source}\PYG{o}{\PYGZhy{}}\PYG{n}{value}\PYG{+w}{ }\PYG{l+m}{1:4}\PYG{+w}{ }\PYG{l+m}{1}\PYG{+w}{ }\PYG{l+m}{0.0}\PYG{+w}{ }\PYG{l+m}{0.5}\PYG{+w}{ }\PYG{l+m}{0.5}\PYG{+w}{ }\PYG{l+m}{0.0}\PYG{+w}{ }\PYG{o}{\PYGZam{}}
\PYG{+w}{                      }\PYG{n}{key}\PYG{+w}{ }\PYG{n}{basis}\PYG{o}{\PYGZhy{}}\PYG{n}{atom}\PYG{o}{\PYGZhy{}}\PYG{n}{coordinates}\PYG{+w}{ }\PYG{n}{source}\PYG{o}{\PYGZhy{}}\PYG{n}{value}\PYG{+w}{ }\PYG{l+m}{1:4}\PYG{+w}{ }\PYG{l+m}{2}\PYG{+w}{ }\PYG{l+m}{0.0}\PYG{+w}{ }\PYG{l+m}{0.5}\PYG{+w}{ }\PYG{l+m}{0.0}\PYG{+w}{ }\PYG{l+m}{0.5}\PYG{+w}{ }\PYG{o}{\PYGZam{}}
\PYG{+w}{                      }\PYG{n}{key}\PYG{+w}{ }\PYG{n}{basis}\PYG{o}{\PYGZhy{}}\PYG{n}{atom}\PYG{o}{\PYGZhy{}}\PYG{n}{coordinates}\PYG{+w}{ }\PYG{n}{source}\PYG{o}{\PYGZhy{}}\PYG{n}{value}\PYG{+w}{ }\PYG{l+m}{1:4}\PYG{+w}{ }\PYG{l+m}{3}\PYG{+w}{ }\PYG{l+m}{0.0}\PYG{+w}{ }\PYG{l+m}{0.0}\PYG{+w}{ }\PYG{l+m}{0.5}\PYG{+w}{ }\PYG{l+m}{0.5}
\end{sphinxVerbatim}

\begin{sphinxadmonition}{note}{Note:}
\sphinxAtStartPar
Multiple calls of \sphinxstyleemphasis{kim property modify} made for the same instance ID
can be combined into a single invocation, meaning the following are
both valid:

\begin{sphinxVerbatim}[commandchars=\\\{\}]
\PYG{k}{kim }\PYG{l+s+s1}{property}\PYG{+w}{ }\PYG{n}{modify}\PYG{+w}{ }\PYG{l+m}{2}\PYG{+w}{ }\PYG{n}{key}\PYG{+w}{ }\PYG{n}{basis}\PYG{o}{\PYGZhy{}}\PYG{n}{atom}\PYG{o}{\PYGZhy{}}\PYG{n}{coordinates}\PYG{+w}{ }\PYG{n}{source}\PYG{o}{\PYGZhy{}}\PYG{n}{value}\PYG{+w}{ }\PYG{l+m}{1}\PYG{+w}{ }\PYG{l+m}{1:3}\PYG{+w}{ }\PYG{l+m}{0.0}\PYG{+w}{ }\PYG{l+m}{0.0}\PYG{+w}{ }\PYG{l+m}{0.0}\PYG{+w}{ }\PYG{o}{\PYGZam{}}
\PYG{+w}{                      }\PYG{n}{key}\PYG{+w}{ }\PYG{n}{basis}\PYG{o}{\PYGZhy{}}\PYG{n}{atom}\PYG{o}{\PYGZhy{}}\PYG{n}{coordinates}\PYG{+w}{ }\PYG{n}{source}\PYG{o}{\PYGZhy{}}\PYG{n}{value}\PYG{+w}{ }\PYG{l+m}{2}\PYG{+w}{ }\PYG{l+m}{1:3}\PYG{+w}{ }\PYG{l+m}{0.5}\PYG{+w}{ }\PYG{l+m}{0.5}\PYG{+w}{ }\PYG{l+m}{0.0}\PYG{+w}{ }\PYG{o}{\PYGZam{}}
\PYG{+w}{                      }\PYG{n}{key}\PYG{+w}{ }\PYG{n}{basis}\PYG{o}{\PYGZhy{}}\PYG{n}{atom}\PYG{o}{\PYGZhy{}}\PYG{n}{coordinates}\PYG{+w}{ }\PYG{n}{source}\PYG{o}{\PYGZhy{}}\PYG{n}{value}\PYG{+w}{ }\PYG{l+m}{3}\PYG{+w}{ }\PYG{l+m}{1:3}\PYG{+w}{ }\PYG{l+m}{0.5}\PYG{+w}{ }\PYG{l+m}{0.0}\PYG{+w}{ }\PYG{l+m}{0.5}\PYG{+w}{ }\PYG{o}{\PYGZam{}}
\PYG{+w}{                      }\PYG{n}{key}\PYG{+w}{ }\PYG{n}{basis}\PYG{o}{\PYGZhy{}}\PYG{n}{atom}\PYG{o}{\PYGZhy{}}\PYG{n}{coordinates}\PYG{+w}{ }\PYG{n}{source}\PYG{o}{\PYGZhy{}}\PYG{n}{value}\PYG{+w}{ }\PYG{l+m}{4}\PYG{+w}{ }\PYG{l+m}{1:3}\PYG{+w}{ }\PYG{l+m}{0.0}\PYG{+w}{ }\PYG{l+m}{0.5}\PYG{+w}{ }\PYG{l+m}{0.5}
\end{sphinxVerbatim}

\begin{sphinxVerbatim}[commandchars=\\\{\}]
\PYG{k}{kim }\PYG{l+s+s1}{property}\PYG{+w}{ }\PYG{n}{modify}\PYG{+w}{ }\PYG{l+m}{2}\PYG{+w}{ }\PYG{n}{key}\PYG{+w}{ }\PYG{n}{short}\PYG{o}{\PYGZhy{}}\PYG{n}{name}\PYG{+w}{ }\PYG{n}{source}\PYG{o}{\PYGZhy{}}\PYG{n}{value}\PYG{+w}{ }\PYG{l+m}{1}\PYG{+w}{ }\PYG{n}{fcc}\PYG{+w}{                         }\PYG{o}{\PYGZam{}}
\PYG{+w}{                      }\PYG{n}{key}\PYG{+w}{ }\PYG{n}{species}\PYG{+w}{ }\PYG{n}{source}\PYG{o}{\PYGZhy{}}\PYG{n}{value}\PYG{+w}{ }\PYG{l+m}{1:4}\PYG{+w}{ }\PYG{n}{Al}\PYG{+w}{ }\PYG{n}{Al}\PYG{+w}{ }\PYG{n}{Al}\PYG{+w}{ }\PYG{n}{Al}\PYG{+w}{                  }\PYG{o}{\PYGZam{}}
\PYG{+w}{                      }\PYG{n}{key}\PYG{+w}{ }\PYG{n}{a}\PYG{+w}{ }\PYG{n}{source}\PYG{o}{\PYGZhy{}}\PYG{n}{value}\PYG{+w}{ }\PYG{l+m}{1:5}\PYG{+w}{ }\PYG{l+m}{3.9149}\PYG{+w}{ }\PYG{l+m}{4.0000}\PYG{+w}{ }\PYG{l+m}{4.032}\PYG{+w}{ }\PYG{l+m}{4.0817}\PYG{+w}{ }\PYG{l+m}{4.1602}\PYG{+w}{  }\PYG{o}{\PYGZam{}}
\PYG{+w}{                            }\PYG{n}{source}\PYG{o}{\PYGZhy{}}\PYG{n}{unit}\PYG{+w}{ }\PYG{n}{angstrom}\PYG{+w}{                                }\PYG{o}{\PYGZam{}}
\PYG{+w}{                            }\PYG{n}{digits}\PYG{+w}{ }\PYG{l+m}{5}\PYG{+w}{                                            }\PYG{o}{\PYGZam{}}
\PYG{+w}{                      }\PYG{n}{key}\PYG{+w}{ }\PYG{n}{basis}\PYG{o}{\PYGZhy{}}\PYG{n}{atom}\PYG{o}{\PYGZhy{}}\PYG{n}{coordinates}\PYG{+w}{ }\PYG{n}{source}\PYG{o}{\PYGZhy{}}\PYG{n}{value}\PYG{+w}{ }\PYG{l+m}{1}\PYG{+w}{ }\PYG{l+m}{1:3}\PYG{+w}{ }\PYG{l+m}{0.0}\PYG{+w}{ }\PYG{l+m}{0.0}\PYG{+w}{ }\PYG{l+m}{0.0}\PYG{+w}{ }\PYG{o}{\PYGZam{}}
\PYG{+w}{                      }\PYG{n}{key}\PYG{+w}{ }\PYG{n}{basis}\PYG{o}{\PYGZhy{}}\PYG{n}{atom}\PYG{o}{\PYGZhy{}}\PYG{n}{coordinates}\PYG{+w}{ }\PYG{n}{source}\PYG{o}{\PYGZhy{}}\PYG{n}{value}\PYG{+w}{ }\PYG{l+m}{2}\PYG{+w}{ }\PYG{l+m}{1:3}\PYG{+w}{ }\PYG{l+m}{0.5}\PYG{+w}{ }\PYG{l+m}{0.5}\PYG{+w}{ }\PYG{l+m}{0.0}\PYG{+w}{ }\PYG{o}{\PYGZam{}}
\PYG{+w}{                      }\PYG{n}{key}\PYG{+w}{ }\PYG{n}{basis}\PYG{o}{\PYGZhy{}}\PYG{n}{atom}\PYG{o}{\PYGZhy{}}\PYG{n}{coordinates}\PYG{+w}{ }\PYG{n}{source}\PYG{o}{\PYGZhy{}}\PYG{n}{value}\PYG{+w}{ }\PYG{l+m}{3}\PYG{+w}{ }\PYG{l+m}{1:3}\PYG{+w}{ }\PYG{l+m}{0.5}\PYG{+w}{ }\PYG{l+m}{0.0}\PYG{+w}{ }\PYG{l+m}{0.5}\PYG{+w}{ }\PYG{o}{\PYGZam{}}
\PYG{+w}{                      }\PYG{n}{key}\PYG{+w}{ }\PYG{n}{basis}\PYG{o}{\PYGZhy{}}\PYG{n}{atom}\PYG{o}{\PYGZhy{}}\PYG{n}{coordinates}\PYG{+w}{ }\PYG{n}{source}\PYG{o}{\PYGZhy{}}\PYG{n}{value}\PYG{+w}{ }\PYG{l+m}{4}\PYG{+w}{ }\PYG{l+m}{1:3}\PYG{+w}{ }\PYG{l+m}{0.0}\PYG{+w}{ }\PYG{l+m}{0.5}\PYG{+w}{ }\PYG{l+m}{0.5}
\end{sphinxVerbatim}
\end{sphinxadmonition}

\begin{sphinxadmonition}{note}{Note:}
\sphinxAtStartPar
For multidimensional arrays, only one colon\sphinxhyphen{}separated range is allowed
in the index listing.  Therefore,

\begin{sphinxVerbatim}[commandchars=\\\{\}]
\PYG{k}{kim }\PYG{l+s+s1}{property}\PYG{+w}{ }\PYG{n}{modify}\PYG{+w}{ }\PYG{l+m}{2}\PYG{+w}{ }\PYG{n}{key}\PYG{+w}{ }\PYG{n}{basis}\PYG{o}{\PYGZhy{}}\PYG{n}{atom}\PYG{o}{\PYGZhy{}}\PYG{n}{coordinates}\PYG{+w}{ }\PYG{l+m}{1}\PYG{+w}{ }\PYG{l+m}{1:3}\PYG{+w}{ }\PYG{l+m}{0.0}\PYG{+w}{ }\PYG{l+m}{0.0}\PYG{+w}{ }\PYG{l+m}{0.0}
\end{sphinxVerbatim}

\sphinxAtStartPar
is valid but

\begin{sphinxVerbatim}[commandchars=\\\{\}]
\PYG{k}{kim }\PYG{l+s+s1}{property}\PYG{+w}{ }\PYG{n}{modify}\PYG{+w}{ }\PYG{l+m}{2}\PYG{+w}{ }\PYG{n}{key}\PYG{+w}{ }\PYG{n}{basis}\PYG{o}{\PYGZhy{}}\PYG{n}{atom}\PYG{o}{\PYGZhy{}}\PYG{n}{coordinates}\PYG{+w}{ }\PYG{l+m}{1:2}\PYG{+w}{ }\PYG{l+m}{1:3}\PYG{+w}{ }\PYG{l+m}{0.0}\PYG{+w}{ }\PYG{l+m}{0.0}\PYG{+w}{ }\PYG{l+m}{0.0}\PYG{+w}{ }\PYG{l+m}{0.0}\PYG{+w}{ }\PYG{l+m}{0.0}\PYG{+w}{ }\PYG{l+m}{0.0}
\end{sphinxVerbatim}

\sphinxAtStartPar
is not.
\end{sphinxadmonition}

\begin{sphinxadmonition}{note}{Note:}
\sphinxAtStartPar
After one sets a value in a map with the \sphinxstyleemphasis{kim property modify} command,
additional calls will overwrite the previous value.
\end{sphinxadmonition}

\sphinxAtStartPar
\sphinxstylestrong{Remove}

\begin{sphinxVerbatim}[commandchars=\\\{\}]
\PYG{k}{kim }\PYG{l+s+s1}{property}\PYG{+w}{ }\PYG{n}{remove}\PYG{+w}{ }\PYG{n}{instance\PYGZus{}id}\PYG{+w}{ }\PYG{n}{key}\PYG{+w}{ }\PYG{n}{key\PYGZus{}name}
\end{sphinxVerbatim}

\sphinxAtStartPar
The \sphinxstyleemphasis{kim property remove} command can be used to remove a property key from a
property instance.  For example,

\begin{sphinxVerbatim}[commandchars=\\\{\}]
\PYG{k}{kim }\PYG{l+s+s1}{property}\PYG{+w}{ }\PYG{n}{remove}\PYG{+w}{ }\PYG{l+m}{2}\PYG{+w}{ }\PYG{n}{key}\PYG{+w}{ }\PYG{n}{basis}\PYG{o}{\PYGZhy{}}\PYG{n}{atom}\PYG{o}{\PYGZhy{}}\PYG{n}{coordinates}
\end{sphinxVerbatim}

\sphinxAtStartPar
\sphinxstylestrong{Destroy}

\begin{sphinxVerbatim}[commandchars=\\\{\}]
\PYG{k}{kim }\PYG{l+s+s1}{property}\PYG{+w}{ }\PYG{n}{destroy}\PYG{+w}{ }\PYG{n}{instance\PYGZus{}id}
\end{sphinxVerbatim}

\sphinxAtStartPar
The \sphinxstyleemphasis{kim property destroy} command deletes a previously created property
instance ID.  For example,

\begin{sphinxVerbatim}[commandchars=\\\{\}]
\PYG{k}{kim }\PYG{l+s+s1}{property}\PYG{+w}{ }\PYG{n}{destroy}\PYG{+w}{ }\PYG{l+m}{2}
\end{sphinxVerbatim}

\begin{sphinxadmonition}{note}{Note:}
\sphinxAtStartPar
If this command is called with an instance ID that does not exist, no
error is raised.
\end{sphinxadmonition}

\sphinxAtStartPar
\sphinxstylestrong{Dump}

\sphinxAtStartPar
The \sphinxstyleemphasis{kim property dump}  command can be used to write the content of all
currently defined property instances to a file:

\begin{sphinxVerbatim}[commandchars=\\\{\}]
\PYG{k}{kim }\PYG{l+s+s1}{property}\PYG{+w}{ }\PYG{k}{dump }\PYG{n+nv+nvIdentifier}{file}
\end{sphinxVerbatim}

\sphinxAtStartPar
For example,

\begin{sphinxVerbatim}[commandchars=\\\{\}]
\PYG{k}{kim }\PYG{l+s+s1}{property}\PYG{+w}{ }\PYG{k}{dump }\PYG{n+nv+nvIdentifier}{results.edn}
\end{sphinxVerbatim}

\begin{sphinxadmonition}{note}{Note:}
\sphinxAtStartPar
Issuing the \sphinxstyleemphasis{kim property dump} command clears all existing property
instances from memory.
\end{sphinxadmonition}


\subsection{Citation of OpenKIM IMs}
\label{\detokenize{kim_commands:citation-of-openkim-ims}}
\sphinxAtStartPar
When publishing results obtained using OpenKIM IMs researchers are
requested to cite the OpenKIM project {\hyperref[\detokenize{kim_commands:kim-mainpaper}]{\sphinxcrossref{\DUrole{std,std-ref}{(Tadmor)}}}},
KIM API {\hyperref[\detokenize{kim_commands:id18}]{\sphinxcrossref{\DUrole{std,std-ref}{(Elliott)}}}}, and the specific IM codes used in
the simulations, in addition to the relevant scientific references for
the IM. The citation format for an IM is displayed on its page on
\sphinxhref{https://openkim.org}{OpenKIM} along with the corresponding BibTex
file, and is automatically added to the LAMMPS citation reminder.

\sphinxAtStartPar
Citing the IM software (KIM infrastructure and specific PM or SM codes) used in
the simulation gives credit to the researchers who developed them and enables
open source efforts like OpenKIM to function.


\subsection{Restrictions}
\label{\detokenize{kim_commands:restrictions}}
\sphinxAtStartPar
The \sphinxstyleemphasis{kim} command is part of the KIM package.  It is only enabled if
LAMMPS is built with that package.  A requirement for the KIM package,
is the KIM API library that must be downloaded from the \sphinxhref{https://openkim.org/kim-api/}{OpenKIM website} and installed before LAMMPS is
compiled.  When installing LAMMPS from binary, the kim\sphinxhyphen{}api package is a
dependency that is automatically downloaded and installed.  The \sphinxstyleemphasis{kim
query} command requires the \sphinxstyleemphasis{libcurl} library to be installed.  The \sphinxstyleemphasis{kim
property} command requires \sphinxstyleemphasis{Python} 3.6 or later and the \sphinxstyleemphasis{kim\sphinxhyphen{}property}
python package to be installed.  See the KIM section of the
{\hyperref[\detokenize{Packages_details::doc}]{\sphinxcrossref{\DUrole{doc}{Packages details}}}} for details.

\sphinxAtStartPar
Furthermore, when using \sphinxstyleemphasis{kim} command to run KIM SMs, any packages required by
the native potential being used or other commands or fixes that it invokes must
be installed.


\subsection{Related commands}
\label{\detokenize{kim_commands:related-commands}}
\sphinxAtStartPar
{\hyperref[\detokenize{pair_kim::doc}]{\sphinxcrossref{\DUrole{doc}{pair\_style kim}}}}


\bigskip\hrule\bigskip

\phantomsection\label{\detokenize{kim_commands:kim-mainpaper}}
\sphinxAtStartPar
\sphinxstylestrong{(Tadmor)} Tadmor, Elliott, Sethna, Miller and Becker, JOM, 63, 17 (2011).
doi: \sphinxurl{https://doi.org/10.1007/s11837-011-0102-6}

\phantomsection\label{\detokenize{kim_commands:id18}}
\sphinxAtStartPar
\sphinxstylestrong{(Elliott)} Elliott, Tadmor and Bernstein, \sphinxurl{https://openkim.org/kim-api} (2011)
doi: \sphinxurl{https://doi.org/10.25950/FF8F563A}

\sphinxstepscope

\index{kspace\_modify@\spxentry{kspace\_modify}}\ignorespaces 

\section{kspace\_modify command}
\label{\detokenize{kspace_modify:kspace-modify-command}}\label{\detokenize{kspace_modify:index-0}}\label{\detokenize{kspace_modify::doc}}

\subsection{Syntax}
\label{\detokenize{kspace_modify:syntax}}
\begin{sphinxVerbatim}[commandchars=\\\{\}]
\PYG{k}{kspace\PYGZus{}modify}\PYG{+w}{ }\PYG{n}{keyword}\PYG{+w}{ }\PYG{n}{value}\PYG{+w}{ }\PYG{n}{...}
\end{sphinxVerbatim}
\begin{itemize}
\item {} 
\sphinxAtStartPar
one or more keyword/value pairs may be listed

\item {} 
\sphinxAtStartPar
keyword = \sphinxstyleemphasis{collective} or \sphinxstyleemphasis{compute} or \sphinxstyleemphasis{cutoff/adjust} or \sphinxstyleemphasis{diff} or \sphinxstyleemphasis{disp/auto} or \sphinxstyleemphasis{fftbench} or \sphinxstyleemphasis{force/disp/kspace} or \sphinxstyleemphasis{force/disp/real} or \sphinxstyleemphasis{force} or \sphinxstyleemphasis{gewald/disp} or \sphinxstyleemphasis{gewald} or \sphinxstyleemphasis{kmax/ewald} or \sphinxstyleemphasis{mesh} or \sphinxstyleemphasis{minorder} or \sphinxstyleemphasis{mix/disp} or \sphinxstyleemphasis{order/disp} or \sphinxstyleemphasis{order} or \sphinxstyleemphasis{overlap} or \sphinxstyleemphasis{scafacos} or \sphinxstyleemphasis{slab} or \sphinxstyleemphasis{splittol} or \sphinxstyleemphasis{wire}
\begin{sphinxalltt}
\sphinxstyleemphasis{collective} value = \sphinxstyleemphasis{yes} or \sphinxstyleemphasis{no}
\sphinxstyleemphasis{compute} value = \sphinxstyleemphasis{yes} or \sphinxstyleemphasis{no}
\sphinxstyleemphasis{cutoff/adjust} value = \sphinxstyleemphasis{yes} or \sphinxstyleemphasis{no}
\sphinxstyleemphasis{diff} value = \sphinxstyleemphasis{ad} or \sphinxstyleemphasis{ik} = 2 or 4 FFTs for PPPM in smoothed or non\sphinxhyphen{}smoothed mode
\sphinxstyleemphasis{disp/auto} value = yes or no
\sphinxstyleemphasis{fftbench} value = \sphinxstyleemphasis{yes} or \sphinxstyleemphasis{no}
\sphinxstyleemphasis{force/disp/real} value = accuracy (force units)
\sphinxstyleemphasis{force/disp/kspace} value = accuracy (force units)
\sphinxstyleemphasis{force} value = accuracy (force units)
\sphinxstyleemphasis{gewald} value = rinv (1/distance units)
  rinv = G\sphinxhyphen{}ewald parameter for Coulombics
\sphinxstyleemphasis{gewald/disp} value = rinv (1/distance units)
  rinv = G\sphinxhyphen{}ewald parameter for dispersion
\sphinxstyleemphasis{kmax/ewald} value = kx ky kz
  kx,ky,kz = number of Ewald sum kspace vectors in each dimension
\sphinxstyleemphasis{mesh} value = x y z
  x,y,z = grid size in each dimension for long\sphinxhyphen{}range Coulombics
\sphinxstyleemphasis{mesh/disp} value = x y z
  x,y,z = grid size in each dimension for 1/r\textasciicircum{}6 dispersion
\sphinxstyleemphasis{minorder} value = M
  M = min allowed extent of Gaussian when auto\sphinxhyphen{}adjusting to minimize grid communication
\sphinxstyleemphasis{mix/disp} value = \sphinxstyleemphasis{pair} or \sphinxstyleemphasis{geom} or \sphinxstyleemphasis{none}
\sphinxstyleemphasis{order} value = N
  N = extent of Gaussian for PPPM or MSM mapping of charge to grid
\sphinxstyleemphasis{order/disp} value = N
  N = extent of Gaussian for PPPM mapping of dispersion term to grid
\sphinxstyleemphasis{overlap} = \sphinxstyleemphasis{yes} or \sphinxstyleemphasis{no} = whether the grid stencil for PPPM is allowed to overlap into more than the nearest\sphinxhyphen{}neighbor processor
\sphinxstyleemphasis{pressure/scalar} value = \sphinxstyleemphasis{yes} or \sphinxstyleemphasis{no}
\sphinxstyleemphasis{scafacos} values = option value1 value2 ...
  option = \sphinxstyleemphasis{tolerance}
    value = \sphinxstyleemphasis{energy} or \sphinxstyleemphasis{energy\_rel} or \sphinxstyleemphasis{field} or \sphinxstyleemphasis{field\_rel} or \sphinxstyleemphasis{potential} or \sphinxstyleemphasis{potential\_rel}
  option = \sphinxstyleemphasis{fmm\_tuning}
    value = \sphinxstyleemphasis{0} or \sphinxstyleemphasis{1}
\sphinxstyleemphasis{slab} value = volfactor or \sphinxstyleemphasis{nozforce}
  volfactor = ratio of the total extended volume used in the
    2d approximation compared with the volume of the simulation domain
  \sphinxstyleemphasis{nozforce} turns off kspace forces in the z direction
\sphinxstyleemphasis{splittol} value = tol
  tol = relative size of two eigenvalues (see discussion below)
\sphinxstyleemphasis{wire} value = volfactor (available with ELECTRODE package)
  volfactor = ratio of the total extended dimension used in the 1d
    approximation compared with the dimension of the simulation domain
\end{sphinxalltt}

\end{itemize}


\subsection{Examples}
\label{\detokenize{kspace_modify:examples}}
\begin{sphinxVerbatim}[commandchars=\\\{\}]
\PYG{k}{kspace\PYGZus{}modify}\PYG{+w}{ }\PYG{n}{mesh}\PYG{+w}{ }\PYG{l+m}{24}\PYG{+w}{ }\PYG{l+m}{24}\PYG{+w}{ }\PYG{l+m}{30}\PYG{+w}{ }\PYG{n}{order}\PYG{+w}{ }\PYG{l+m}{6}
\PYG{k}{kspace\PYGZus{}modify}\PYG{+w}{ }\PYG{n}{slab}\PYG{+w}{ }\PYG{l+m}{3.0}
\PYG{k}{kspace\PYGZus{}modify}\PYG{+w}{ }\PYG{n}{scafacos}\PYG{+w}{ }\PYG{n}{tolerance}\PYG{+w}{ }\PYG{n}{energy}
\end{sphinxVerbatim}


\subsection{Description}
\label{\detokenize{kspace_modify:description}}
\sphinxAtStartPar
Set parameters used by the kspace solvers defined by the
{\hyperref[\detokenize{kspace_style::doc}]{\sphinxcrossref{\DUrole{doc}{kspace\_style}}}} command.  Not all parameters are
relevant to all kspace styles.


\bigskip\hrule\bigskip


\sphinxAtStartPar
The \sphinxstyleemphasis{collective} keyword applies only to PPPM.  It is set to \sphinxstyleemphasis{no} by
default, except on IBM BlueGene machines.  If this option is set to
\sphinxstyleemphasis{yes}, LAMMPS will use MPI collective operations to remap data for
3d\sphinxhyphen{}FFT operations instead of the default point\sphinxhyphen{}to\sphinxhyphen{}point communication.
This is faster on IBM BlueGene machines, and may also be faster on
other machines if they have an efficient implementation of MPI
collective operations and adequate hardware.


\bigskip\hrule\bigskip


\sphinxAtStartPar
The \sphinxstyleemphasis{compute} keyword allows Kspace computations to be turned off,
even though a {\hyperref[\detokenize{kspace_style::doc}]{\sphinxcrossref{\DUrole{doc}{kspace\_style}}}} is defined.  This is
not useful for running a real simulation, but can be useful for
debugging purposes or for computing only partial forces that do not
include the Kspace contribution.  You can also do this by simply not
defining a {\hyperref[\detokenize{kspace_style::doc}]{\sphinxcrossref{\DUrole{doc}{kspace\_style}}}}, but a Kspace\sphinxhyphen{}compatible
{\hyperref[\detokenize{pair_style::doc}]{\sphinxcrossref{\DUrole{doc}{pair\_style}}}} requires a kspace style to be defined.
This keyword gives you that option.


\bigskip\hrule\bigskip


\sphinxAtStartPar
The \sphinxstyleemphasis{cutoff/adjust} keyword applies only to MSM. If this option is
turned on, the Coulombic cutoff will be automatically adjusted at the
beginning of the run to give the desired estimated error. Other
cutoffs such as LJ will not be affected. If the grid is not set using
the \sphinxstyleemphasis{mesh} command, this command will also attempt to use the optimal
grid that minimizes cost using an estimate given by
{\hyperref[\detokenize{kspace_modify:hardy1}]{\sphinxcrossref{\DUrole{std,std-ref}{(Hardy)}}}}. Note that this cost estimate is not exact, somewhat
experimental, and still may not yield the optimal parameters.


\bigskip\hrule\bigskip


\sphinxAtStartPar
The \sphinxstyleemphasis{diff} keyword specifies the differentiation scheme used by the
PPPM method to compute forces on particles given electrostatic
potentials on the PPPM mesh.  The \sphinxstyleemphasis{ik} approach is the default for
PPPM and is the original formulation used in {\hyperref[\detokenize{kspace_modify:hockney1}]{\sphinxcrossref{\DUrole{std,std-ref}{(Hockney)}}}}.  It
performs differentiation in Kspace, and uses 3 FFTs to transfer each
component of the computed fields back to real space for total of 4
FFTs per timestep.

\sphinxAtStartPar
The analytic differentiation \sphinxstyleemphasis{ad} approach uses only 1 FFT to transfer
information back to real space for a total of 2 FFTs per timestep.  It
then performs analytic differentiation on the single quantity to
generate the 3 components of the electric field at each grid point.
This is sometimes referred to as “smoothed” PPPM.  This approach
requires a somewhat larger PPPM mesh to achieve the same accuracy as
the \sphinxstyleemphasis{ik} method. Currently, only the \sphinxstyleemphasis{ik} method (default) can be
used for a triclinic simulation cell with PPPM. The \sphinxstyleemphasis{ad} method is
always used for MSM.

\begin{sphinxadmonition}{note}{Note:}
\sphinxAtStartPar
Currently, not all PPPM styles support the \sphinxstyleemphasis{ad} option.  Support
for those PPPM variants will be added later.
\end{sphinxadmonition}


\bigskip\hrule\bigskip


\sphinxAtStartPar
The \sphinxstyleemphasis{disp/auto} option controls whether the pppm/disp is allowed to
generate PPPM parameters automatically. If set to \sphinxstyleemphasis{no}, parameters have
to be specified using the \sphinxstyleemphasis{gewald/disp}, \sphinxstyleemphasis{mesh/disp},
\sphinxstyleemphasis{force/disp/real} or \sphinxstyleemphasis{force/disp/kspace} keywords, or
the code will stop with an error message. When this option is set to
\sphinxstyleemphasis{yes}, the error message will not appear and the simulation will start.
For a typical application, using the automatic parameter generation
will provide simulations that are either inaccurate or slow. Using this
option is thus not recommended. For guidelines on how to obtain good
parameters, see the {\hyperref[\detokenize{Howto_dispersion::doc}]{\sphinxcrossref{\DUrole{doc}{long\sphinxhyphen{}range dispersion howto}}}}
discussion.


\bigskip\hrule\bigskip


\sphinxAtStartPar
The \sphinxstyleemphasis{fftbench} keyword applies only to PPPM. It is off by default. If
this option is turned on, LAMMPS will perform a short FFT benchmark
computation and report its timings, and will thus finish some seconds
later than it would if this option were off.


\bigskip\hrule\bigskip


\sphinxAtStartPar
The \sphinxstyleemphasis{force/disp/real} and \sphinxstyleemphasis{force/disp/kspace} keywords set the force
accuracy for the real and reciprocal space computations for the dispersion
part of pppm/disp. As shown in {\hyperref[\detokenize{kspace_modify:isele-holder1}]{\sphinxcrossref{\DUrole{std,std-ref}{(Isele\sphinxhyphen{}Holder)}}}},
optimal performance and accuracy in the results is obtained when these
values are different.


\bigskip\hrule\bigskip


\sphinxAtStartPar
The \sphinxstyleemphasis{force} keyword overrides the relative accuracy parameter set by
the {\hyperref[\detokenize{kspace_style::doc}]{\sphinxcrossref{\DUrole{doc}{kspace\_style}}}} command with an absolute
accuracy.  The accuracy determines the RMS error in per\sphinxhyphen{}atom forces
calculated by the long\sphinxhyphen{}range solver and is thus specified in force
units.  A negative value for the accuracy setting means to use the
relative accuracy parameter.  The accuracy setting is used in
conjunction with the pairwise cutoff to determine the number of
K\sphinxhyphen{}space vectors for style \sphinxstyleemphasis{ewald}, the FFT grid size for style
\sphinxstyleemphasis{pppm}, or the real space grid size for style \sphinxstyleemphasis{msm}.


\bigskip\hrule\bigskip


\sphinxAtStartPar
The \sphinxstyleemphasis{gewald} keyword sets the value of the Ewald or PPPM G\sphinxhyphen{}ewald
parameter for charge as \sphinxstyleemphasis{rinv} in reciprocal distance units.  Without
this setting, LAMMPS chooses the parameter automatically as a function
of cutoff, precision, grid spacing, etc.  This means it can vary from
one simulation to the next which may not be desirable for matching a
KSpace solver to a pre\sphinxhyphen{}tabulated pairwise potential.  This setting can
also be useful if Ewald or PPPM fails to choose a good grid spacing
and G\sphinxhyphen{}ewald parameter automatically.  If the value is set to 0.0,
LAMMPS will choose the G\sphinxhyphen{}ewald parameter automatically.  MSM does not
use the \sphinxstyleemphasis{gewald} parameter.


\bigskip\hrule\bigskip


\sphinxAtStartPar
The \sphinxstyleemphasis{gewald/disp} keyword sets the value of the Ewald or PPPM G\sphinxhyphen{}ewald
parameter for dispersion as \sphinxstyleemphasis{rinv} in reciprocal distance units.  It
has the same meaning as the \sphinxstyleemphasis{gewald} setting for Coulombics.


\bigskip\hrule\bigskip


\sphinxAtStartPar
The \sphinxstyleemphasis{kmax/ewald} keyword sets the number of kspace vectors in each
dimension for kspace style \sphinxstyleemphasis{ewald}.  The three values must be positive
integers, or else (0,0,0), which unsets the option.  When this option
is not set, the Ewald sum scheme chooses its own kspace vectors,
consistent with the user\sphinxhyphen{}specified accuracy and pairwise cutoff. In
any case, if kspace style \sphinxstyleemphasis{ewald} is invoked, the values used are
printed to the screen and the log file at the start of the run.


\bigskip\hrule\bigskip


\sphinxAtStartPar
The \sphinxstyleemphasis{mesh} keyword sets the grid size for kspace style \sphinxstyleemphasis{pppm} or
\sphinxstyleemphasis{msm}.  In the case of PPPM, this is the FFT mesh, and each dimension
must be factorizable into powers of 2, 3, and 5.  In the case of MSM,
this is the finest scale real\sphinxhyphen{}space mesh, and each dimension must be
factorizable into powers of 2.  When this option is not set, the PPPM
or MSM solver chooses its own grid size, consistent with the
user\sphinxhyphen{}specified accuracy and pairwise cutoff.  Values for x,y,z of
0,0,0 unset the option.


\bigskip\hrule\bigskip


\sphinxAtStartPar
The \sphinxstyleemphasis{mesh/disp} keyword sets the grid size for kspace style
\sphinxstyleemphasis{pppm/disp}.  This is the FFT mesh for long\sphinxhyphen{}range dispersion and ach
dimension must be factorizable into powers of 2, 3, and 5.  When this
option is not set, the PPPM solver chooses its own grid size,
consistent with the user\sphinxhyphen{}specified accuracy and pairwise cutoff.
Values for x,y,z of 0,0,0 unset the option.


\bigskip\hrule\bigskip


\sphinxAtStartPar
The \sphinxstyleemphasis{minorder} keyword allows LAMMPS to reduce the \sphinxstyleemphasis{order} setting if
necessary to keep the communication of ghost grid point limited to
exchanges between nearest\sphinxhyphen{}neighbor processors.  See the discussion of
the \sphinxstyleemphasis{overlap} keyword for details.  If the \sphinxstyleemphasis{overlap} keyword is set to
\sphinxstyleemphasis{yes}, which is the default, this is never needed.  If it set to \sphinxstyleemphasis{no}
and overlap occurs, then LAMMPS will reduce the order setting, one
step at a time, until the ghost grid overlap only extends to nearest
neighbor processors.  The \sphinxstyleemphasis{minorder} keyword limits how small the
\sphinxstyleemphasis{order} setting can become.  The minimum allowed value for PPPM is 2,
which is the default.  If \sphinxstyleemphasis{minorder} is set to the same value as
\sphinxstyleemphasis{order} then no reduction is allowed, and LAMMPS will generate an
error if the grid communication is non\sphinxhyphen{}nearest\sphinxhyphen{}neighbor and \sphinxstyleemphasis{overlap}
is set to \sphinxstyleemphasis{no}. The \sphinxstyleemphasis{minorder} keyword is not currently supported in
MSM.


\bigskip\hrule\bigskip


\sphinxAtStartPar
The \sphinxstyleemphasis{mix/disp} keyword selects the mixing rule for the dispersion
coefficients.  With \sphinxstyleemphasis{pair}, the dispersion coefficients of unlike
types are computed as indicated with {\hyperref[\detokenize{pair_modify::doc}]{\sphinxcrossref{\DUrole{doc}{pair\_modify}}}}.
With \sphinxstyleemphasis{geom}, geometric mixing is enforced on the dispersion
coefficients in the kspace coefficients. When using the arithmetic
mixing rule, this will speed\sphinxhyphen{}up the simulations but introduces some
error in the force computations, as shown in {\hyperref[\detokenize{kspace_modify:wennberg}]{\sphinxcrossref{\DUrole{std,std-ref}{(Wennberg)}}}}.
With \sphinxstyleemphasis{none}, it is assumed that no mixing rule is
applicable. Splitting of the dispersion coefficients will be performed
as described in {\hyperref[\detokenize{kspace_modify:isele-holder1}]{\sphinxcrossref{\DUrole{std,std-ref}{(Isele\sphinxhyphen{}Holder)}}}}.

\sphinxAtStartPar
This splitting can be influenced with the \sphinxstyleemphasis{splittol} keywords.  Only
the eigenvalues that are larger than tol compared to the largest
eigenvalues are included. Using this keywords the original matrix of
dispersion coefficients is approximated. This leads to faster
computations, but the accuracy in the reciprocal space computations of
the dispersion part is decreased.


\bigskip\hrule\bigskip


\sphinxAtStartPar
The \sphinxstyleemphasis{order} keyword determines how many grid spacings an atom’s charge
extends when it is mapped to the grid in kspace style \sphinxstyleemphasis{pppm} or \sphinxstyleemphasis{msm}.
The default for this parameter is 5 for PPPM and 8 for MSM, which
means each charge spans 5 or 8 grid cells in each dimension,
respectively.  For the LAMMPS implementation of MSM, the order can
range from 4 to 10 and must be even. For PPPM, the minimum allowed
setting is 2 and the maximum allowed setting is 7.  The larger the
value of this parameter, the smaller that LAMMPS will set the grid
size, to achieve the requested accuracy.  Conversely, the smaller the
order value, the larger the grid size will be.  Note that there is an
inherent trade\sphinxhyphen{}off involved: a small grid will lower the cost of FFTs
or MSM direct sum, but a larger order parameter will increase the cost
of interpolating charge/fields to/from the grid.

\sphinxAtStartPar
The PPPM order parameter may be reset by LAMMPS when it sets up the
FFT grid if the implied grid stencil extends beyond the grid cells
owned by neighboring processors.  Typically this will only occur when
small problems are run on large numbers of processors.  A warning will
be generated indicating the order parameter is being reduced to allow
LAMMPS to run the problem. Automatic adjustment of the order parameter
is not supported in MSM.


\bigskip\hrule\bigskip


\sphinxAtStartPar
The \sphinxstyleemphasis{order/disp} keyword determines how many grid spacings an atom’s
dispersion term extends when it is mapped to the grid in kspace style
\sphinxstyleemphasis{pppm/disp}.  It has the same meaning as the \sphinxstyleemphasis{order} setting for
Coulombics.


\bigskip\hrule\bigskip


\sphinxAtStartPar
The \sphinxstyleemphasis{overlap} keyword can be used in conjunction with the \sphinxstyleemphasis{minorder}
keyword with the PPPM styles to adjust the amount of communication
that occurs when values on the FFT grid are exchanged between
processors.  This communication is distinct from the communication
inherent in the parallel FFTs themselves, and is required because
processors interpolate charge and field values using grid point values
owned by neighboring processors (i.e. ghost point communication).  If
the \sphinxstyleemphasis{overlap} keyword is set to \sphinxstyleemphasis{yes} then this communication is
allowed to extend beyond nearest\sphinxhyphen{}neighbor processors, e.g. when using
lots of processors on a small problem.  If it is set to \sphinxstyleemphasis{no} then the
communication will be limited to nearest\sphinxhyphen{}neighbor processors and the
\sphinxstyleemphasis{order} setting will be reduced if necessary, as explained by the
\sphinxstyleemphasis{minorder} keyword discussion. The \sphinxstyleemphasis{overlap} keyword is always set to
\sphinxstyleemphasis{yes} in MSM.


\bigskip\hrule\bigskip


\sphinxAtStartPar
The \sphinxstyleemphasis{pressure/scalar} keyword applies only to MSM. If this option is
turned on, only the scalar pressure (i.e. (Pxx + Pyy + Pzz)/3.0) will
be computed, which can be used, for example, to run an isotropic barostat.
Computing the full pressure tensor with MSM is expensive, and this option
provides a faster alternative. The scalar pressure is computed using a
relationship between the Coulombic energy and pressure {\hyperref[\detokenize{kspace_modify:hummer}]{\sphinxcrossref{\DUrole{std,std-ref}{(Hummer)}}}}
instead of using the virial equation. This option cannot be used to access
individual components of the pressure tensor, to compute per\sphinxhyphen{}atom virial,
or with suffix kspace/pair styles of MSM, like OMP or GPU.


\bigskip\hrule\bigskip


\sphinxAtStartPar
The \sphinxstyleemphasis{scafacos} keyword is used for settings that are passed to the
ScaFaCoS library when using {\hyperref[\detokenize{kspace_style::doc}]{\sphinxcrossref{\DUrole{doc}{kspace\_style scafacos}}}}.

\sphinxAtStartPar
The \sphinxstyleemphasis{tolerance} option affects how the \sphinxstyleemphasis{accuracy} specified with the
{\hyperref[\detokenize{kspace_style::doc}]{\sphinxcrossref{\DUrole{doc}{kspace\_style}}}} command is interpreted by ScaFaCoS.
The following values may be used:
\begin{itemize}
\item {} 
\sphinxAtStartPar
energy = absolute accuracy in total Coulombic energy

\item {} 
\sphinxAtStartPar
energy\_rel = relative accuracy in total Coulombic energy

\item {} 
\sphinxAtStartPar
potential = absolute accuracy in total Coulombic potential

\item {} 
\sphinxAtStartPar
potential\_rel = relative accuracy in total Coulombic potential

\item {} 
\sphinxAtStartPar
field = absolute accuracy in electric field

\item {} 
\sphinxAtStartPar
field\_rel = relative accuracy in electric field

\end{itemize}

\sphinxAtStartPar
The values with suffix \_rel indicate the tolerance is a relative
tolerance; the other values impose an absolute tolerance on the given
quantity. Absolute tolerance in this case means, that for a given
quantity q and a given absolute tolerance of t\_a the result should
be between q\sphinxhyphen{}t\_a and q+t\_a. For a relative tolerance t\_r the relative
error should not be greater than t\_r, i.e. abs(1 \sphinxhyphen{} (result/q)) \textless{} t\_r.
As a consequence of this, the tolerance type should be checked, when
performing computations with a high absolute field / energy. E.g.
if the total energy in the system is 1000000.0 an absolute tolerance
of 1e\sphinxhyphen{}3 would mean that the result has to be between 999999.999 and
1000000.001, which would be equivalent to a relative tolerance of
1e\sphinxhyphen{}9.

\sphinxAtStartPar
The energy and energy\_rel values, set a tolerance based on the total
Coulombic energy of the system.  The potential and potential\_rel set a
tolerance based on the per\sphinxhyphen{}atom Coulombic energy.  The field and
field\_rel tolerance types set a tolerance based on the electric field
values computed by ScaFaCoS.  Since per\sphinxhyphen{}atom forces are derived from
the per\sphinxhyphen{}atom electric field, this effectively sets a tolerance on the
forces, similar to other LAMMPS KSpace styles, as explained on the
{\hyperref[\detokenize{kspace_style::doc}]{\sphinxcrossref{\DUrole{doc}{kspace\_style}}}} doc page.

\sphinxAtStartPar
Note that not all ScaFaCoS solvers support all tolerance types.
These are the allowed values for each method:
\begin{itemize}
\item {} 
\sphinxAtStartPar
fmm = energy and energy\_rel

\item {} 
\sphinxAtStartPar
p2nfft = field (1d\sphinxhyphen{},2d\sphinxhyphen{},3d\sphinxhyphen{}periodic systems) or potential (0d\sphinxhyphen{}periodic)

\item {} 
\sphinxAtStartPar
p3m = field

\item {} 
\sphinxAtStartPar
ewald = field

\item {} 
\sphinxAtStartPar
direct = has no tolerance tuning

\end{itemize}

\sphinxAtStartPar
If the tolerance type is not changed, the default values for the
tolerance type are the first values in the above list, e.g. energy
is the default tolerance type for the fmm solver.

\sphinxAtStartPar
The \sphinxstyleemphasis{fmm\_tuning} option is only relevant when using the FMM method.
It activates (value=1) or deactivates (value=0) an internal tuning
mechanism for the FMM solver.  The tuning operation runs sequentially
and can be very time\sphinxhyphen{}consuming.  Usually it is not needed for systems
with a homogeneous charge distribution. The default for this option is
therefore \sphinxstyleemphasis{0}. The FMM internal tuning is performed once, when the
solver is set up.


\bigskip\hrule\bigskip


\sphinxAtStartPar
The \sphinxstyleemphasis{slab} keyword allows an Ewald or PPPM solver to be used for a
systems that are periodic in x,y but non\sphinxhyphen{}periodic in z \sphinxhyphen{} a
{\hyperref[\detokenize{boundary::doc}]{\sphinxcrossref{\DUrole{doc}{boundary}}}} setting of “boundary p p f”.  This is done
by treating the system as if it were periodic in z, but inserting
empty volume between atom slabs and removing dipole inter\sphinxhyphen{}slab
interactions so that slab\sphinxhyphen{}slab interactions are effectively turned
off.  The volfactor value sets the ratio of the extended dimension in
z divided by the actual dimension in z.  It must be a value \textgreater{}= 1.0.  A
value of 1.0 (the default) means the slab approximation is not used.

\sphinxAtStartPar
The recommended value for volfactor is 3.0.  A larger value is
inefficient; a smaller value introduces unwanted slab\sphinxhyphen{}slab
interactions.  The use of fixed boundaries in z means that the user
must prevent particle migration beyond the initial z\sphinxhyphen{}bounds, typically
by providing a wall\sphinxhyphen{}style fix.  The methodology behind the \sphinxstyleemphasis{slab}
option is explained in the paper by {\hyperref[\detokenize{kspace_modify:yeh}]{\sphinxcrossref{\DUrole{std,std-ref}{(Yeh)}}}}.  The \sphinxstyleemphasis{slab}
option is also extended to non\sphinxhyphen{}neutral systems {\hyperref[\detokenize{kspace_modify:ballenegger}]{\sphinxcrossref{\DUrole{std,std-ref}{(Ballenegger)}}}}.

\sphinxAtStartPar
An alternative slab option can be invoked with the \sphinxstyleemphasis{nozforce} keyword
in lieu of the volfactor.  This turns off all kspace forces in the z
direction.  The \sphinxstyleemphasis{nozforce} option is not supported by MSM. For MSM,
any combination of periodic, non\sphinxhyphen{}periodic, or shrink\sphinxhyphen{}wrapped
boundaries can be set using {\hyperref[\detokenize{boundary::doc}]{\sphinxcrossref{\DUrole{doc}{boundary}}}} (the slab
approximation in not needed).  The \sphinxstyleemphasis{slab} keyword is not currently
supported by Ewald or PPPM when using a triclinic simulation cell. The
slab correction has also been extended to point dipole interactions
{\hyperref[\detokenize{kspace_modify:klapp}]{\sphinxcrossref{\DUrole{std,std-ref}{(Klapp)}}}} in {\hyperref[\detokenize{kspace_style::doc}]{\sphinxcrossref{\DUrole{doc}{kspace\_style}}}}
\sphinxstyleemphasis{ewald/disp}, \sphinxstyleemphasis{ewald/dipole}, and \sphinxstyleemphasis{pppm/dipole}.

\begin{sphinxadmonition}{note}{Note:}
\sphinxAtStartPar
If you wish to apply an electric field in the Z\sphinxhyphen{}direction, in
conjunction with the \sphinxstyleemphasis{slab} keyword, you should do it by adding
explicit charged particles to the +/\sphinxhyphen{} Z surfaces.  If you do it via
the {\hyperref[\detokenize{fix_efield::doc}]{\sphinxcrossref{\DUrole{doc}{fix efield}}}} command, it will not give the correct
dielectric constant due to the Yeh/Berkowitz {\hyperref[\detokenize{kspace_modify:yeh}]{\sphinxcrossref{\DUrole{std,std-ref}{(Yeh)}}}} correction
not being compatible with how {\hyperref[\detokenize{fix_efield::doc}]{\sphinxcrossref{\DUrole{doc}{fix efield}}}} works.
\end{sphinxadmonition}


\bigskip\hrule\bigskip


\sphinxAtStartPar
The \sphinxstyleemphasis{force/disp/real} and \sphinxstyleemphasis{force/disp/kspace} keywords set the force
accuracy for the real and reciprocal space computations for the dispersion
part of pppm/disp. As shown in {\hyperref[\detokenize{kspace_modify:isele-holder1}]{\sphinxcrossref{\DUrole{std,std-ref}{(Isele\sphinxhyphen{}Holder)}}}},
optimal performance and accuracy in the results is obtained when these
values are different.

\sphinxAtStartPar
The \sphinxstyleemphasis{disp/auto} option controls whether the pppm/disp is allowed to
generate PPPM parameters automatically. If set to \sphinxstyleemphasis{no}, parameters
have to be specified using the \sphinxstyleemphasis{gewald/disp}, \sphinxstyleemphasis{mesh/disp},
\sphinxstyleemphasis{force/disp/real} or \sphinxstyleemphasis{force/disp/kspace} keywords, or the code will
stop with an error message. When this option is set to \sphinxstyleemphasis{yes}, the
error message will not appear and the simulation will start.  For a
typical application, using the automatic parameter generation will
provide simulations that are either inaccurate or slow. Using this
option is thus not recommended.  For guidelines on how to obtain good
parameters, see the {\hyperref[\detokenize{Howto_dispersion::doc}]{\sphinxcrossref{\DUrole{doc}{Howto dispersion}}}} doc page.


\bigskip\hrule\bigskip



\subsection{Restrictions}
\label{\detokenize{kspace_modify:restrictions}}
\sphinxAtStartPar
none


\subsection{Related commands}
\label{\detokenize{kspace_modify:related-commands}}
\sphinxAtStartPar
{\hyperref[\detokenize{kspace_style::doc}]{\sphinxcrossref{\DUrole{doc}{kspace\_style}}}}, {\hyperref[\detokenize{boundary::doc}]{\sphinxcrossref{\DUrole{doc}{boundary}}}}


\subsection{Default}
\label{\detokenize{kspace_modify:default}}
\sphinxAtStartPar
The option defaults are as follows:
\begin{itemize}
\item {} 
\sphinxAtStartPar
compute = yes

\item {} 
\sphinxAtStartPar
cutoff/adjust = yes (MSM)

\item {} 
\sphinxAtStartPar
diff = ik (PPPM)

\item {} 
\sphinxAtStartPar
disp/auto = no

\item {} 
\sphinxAtStartPar
fftbench = no (PPPM)

\item {} 
\sphinxAtStartPar
force = \sphinxhyphen{}1.0

\item {} 
\sphinxAtStartPar
force/disp/kspace = \sphinxhyphen{}1.0

\item {} 
\sphinxAtStartPar
force/disp/real = \sphinxhyphen{}1.0

\item {} 
\sphinxAtStartPar
gewald = gewald/disp = 0.0

\item {} 
\sphinxAtStartPar
mesh = mesh/disp = 0 0 0

\item {} 
\sphinxAtStartPar
minorder = 2

\item {} 
\sphinxAtStartPar
mix/disp = pair

\item {} 
\sphinxAtStartPar
order = 10 (MSM)

\item {} 
\sphinxAtStartPar
order = order/disp = 5 (PPPM)

\item {} 
\sphinxAtStartPar
order = order/disp = 7 (PPPM/intel)

\item {} 
\sphinxAtStartPar
overlap = yes

\item {} 
\sphinxAtStartPar
pressure/scalar = yes (MSM)

\item {} 
\sphinxAtStartPar
slab = 1.0

\item {} 
\sphinxAtStartPar
split = 0

\item {} 
\sphinxAtStartPar
tol = 1.0e\sphinxhyphen{}6

\end{itemize}

\sphinxAtStartPar
For scafacos settings, the scafacos tolerance option depends on the
method chosen, as documented above.  The scafacos fmm\_tuning default
= 0.


\bigskip\hrule\bigskip

\phantomsection\label{\detokenize{kspace_modify:hockney1}}
\sphinxAtStartPar
\sphinxstylestrong{(Hockney)} Hockney and Eastwood, Computer Simulation Using Particles,
Adam Hilger, NY (1989).

\phantomsection\label{\detokenize{kspace_modify:yeh}}
\sphinxAtStartPar
\sphinxstylestrong{(Yeh)} Yeh and Berkowitz, J Chem Phys, 111, 3155 (1999).

\phantomsection\label{\detokenize{kspace_modify:ballenegger}}
\sphinxAtStartPar
\sphinxstylestrong{(Ballenegger)} Ballenegger, Arnold, Cerda, J Chem Phys, 131, 094107
(2009).

\phantomsection\label{\detokenize{kspace_modify:klapp}}
\sphinxAtStartPar
\sphinxstylestrong{(Klapp)} Klapp, Schoen, J Chem Phys, 117, 8050 (2002).

\phantomsection\label{\detokenize{kspace_modify:hardy1}}
\sphinxAtStartPar
\sphinxstylestrong{(Hardy)} David Hardy thesis: Multilevel Summation for the Fast
Evaluation of Forces for the Simulation of Biomolecules, University of
Illinois at Urbana\sphinxhyphen{}Champaign, (2006).

\phantomsection\label{\detokenize{kspace_modify:hummer}}
\sphinxAtStartPar
\sphinxstylestrong{(Hummer)} Hummer, Gronbech\sphinxhyphen{}Jensen, Neumann, J Chem Phys, 109, 2791 (1998)

\phantomsection\label{\detokenize{kspace_modify:isele-holder1}}
\sphinxAtStartPar
\sphinxstylestrong{(Isele\sphinxhyphen{}Holder)} Isele\sphinxhyphen{}Holder, Mitchell, Hammond, Kohlmeyer, Ismail, J
Chem Theory Comput, 9, 5412 (2013).

\phantomsection\label{\detokenize{kspace_modify:wennberg}}
\sphinxAtStartPar
\sphinxstylestrong{(Wennberg)} Wennberg, Murtola, Hess, Lindahl, J Chem Theory Comput,
9, 3527 (2013).

\sphinxstepscope

\index{kspace\_style ewald@\spxentry{kspace\_style ewald}}\ignorespaces 
\index{kspace\_style ewald/dipole@\spxentry{kspace\_style ewald/dipole}}\ignorespaces 
\index{kspace\_style ewald/dipole/spin@\spxentry{kspace\_style ewald/dipole/spin}}\ignorespaces 
\index{kspace\_style ewald/disp@\spxentry{kspace\_style ewald/disp}}\ignorespaces 
\index{kspace\_style ewald/disp/dipole@\spxentry{kspace\_style ewald/disp/dipole}}\ignorespaces 
\index{kspace\_style ewald/omp@\spxentry{kspace\_style ewald/omp}}\ignorespaces 
\index{kspace\_style ewald/electrode@\spxentry{kspace\_style ewald/electrode}}\ignorespaces 
\index{kspace\_style pppm@\spxentry{kspace\_style pppm}}\ignorespaces 
\index{kspace\_style pppm/kk@\spxentry{kspace\_style pppm/kk}}\ignorespaces 
\index{kspace\_style pppm/omp@\spxentry{kspace\_style pppm/omp}}\ignorespaces 
\index{kspace\_style pppm/gpu@\spxentry{kspace\_style pppm/gpu}}\ignorespaces 
\index{kspace\_style pppm/intel@\spxentry{kspace\_style pppm/intel}}\ignorespaces 
\index{kspace\_style pppm/cg@\spxentry{kspace\_style pppm/cg}}\ignorespaces 
\index{kspace\_style pppm/dielectric@\spxentry{kspace\_style pppm/dielectric}}\ignorespaces 
\index{kspace\_style pppm/dipole@\spxentry{kspace\_style pppm/dipole}}\ignorespaces 
\index{kspace\_style pppm/dipole/spin@\spxentry{kspace\_style pppm/dipole/spin}}\ignorespaces 
\index{kspace\_style pppm/disp@\spxentry{kspace\_style pppm/disp}}\ignorespaces 
\index{kspace\_style pppm/disp/omp@\spxentry{kspace\_style pppm/disp/omp}}\ignorespaces 
\index{kspace\_style pppm/disp/tip4p@\spxentry{kspace\_style pppm/disp/tip4p}}\ignorespaces 
\index{kspace\_style pppm/disp/tip4p/omp@\spxentry{kspace\_style pppm/disp/tip4p/omp}}\ignorespaces 
\index{kspace\_style pppm/disp/intel@\spxentry{kspace\_style pppm/disp/intel}}\ignorespaces 
\index{kspace\_style pppm/disp/dielectric@\spxentry{kspace\_style pppm/disp/dielectric}}\ignorespaces 
\index{kspace\_style pppm/cg/omp@\spxentry{kspace\_style pppm/cg/omp}}\ignorespaces 
\index{kspace\_style pppm/stagger@\spxentry{kspace\_style pppm/stagger}}\ignorespaces 
\index{kspace\_style pppm/tip4p@\spxentry{kspace\_style pppm/tip4p}}\ignorespaces 
\index{kspace\_style pppm/tip4p/omp@\spxentry{kspace\_style pppm/tip4p/omp}}\ignorespaces 
\index{kspace\_style pppm/electrode@\spxentry{kspace\_style pppm/electrode}}\ignorespaces 
\index{kspace\_style pppm/electrode/intel@\spxentry{kspace\_style pppm/electrode/intel}}\ignorespaces 
\index{kspace\_style msm@\spxentry{kspace\_style msm}}\ignorespaces 
\index{kspace\_style msm/omp@\spxentry{kspace\_style msm/omp}}\ignorespaces 
\index{kspace\_style msm/cg@\spxentry{kspace\_style msm/cg}}\ignorespaces 
\index{kspace\_style msm/cg/omp@\spxentry{kspace\_style msm/cg/omp}}\ignorespaces 
\index{kspace\_style msm/dielectric@\spxentry{kspace\_style msm/dielectric}}\ignorespaces 
\index{kspace\_style scafacos@\spxentry{kspace\_style scafacos}}\ignorespaces 

\section{kspace\_style command}
\label{\detokenize{kspace_style:kspace-style-command}}\label{\detokenize{kspace_style:index-33}}\label{\detokenize{kspace_style:index-32}}\label{\detokenize{kspace_style:index-31}}\label{\detokenize{kspace_style:index-30}}\label{\detokenize{kspace_style:index-29}}\label{\detokenize{kspace_style:index-28}}\label{\detokenize{kspace_style:index-27}}\label{\detokenize{kspace_style:index-26}}\label{\detokenize{kspace_style:index-25}}\label{\detokenize{kspace_style:index-24}}\label{\detokenize{kspace_style:index-23}}\label{\detokenize{kspace_style:index-22}}\label{\detokenize{kspace_style:index-21}}\label{\detokenize{kspace_style:index-20}}\label{\detokenize{kspace_style:index-19}}\label{\detokenize{kspace_style:index-18}}\label{\detokenize{kspace_style:index-17}}\label{\detokenize{kspace_style:index-16}}\label{\detokenize{kspace_style:index-15}}\label{\detokenize{kspace_style:index-14}}\label{\detokenize{kspace_style:index-13}}\label{\detokenize{kspace_style:index-12}}\label{\detokenize{kspace_style:index-11}}\label{\detokenize{kspace_style:index-10}}\label{\detokenize{kspace_style:index-9}}\label{\detokenize{kspace_style:index-8}}\label{\detokenize{kspace_style:index-7}}\label{\detokenize{kspace_style:index-6}}\label{\detokenize{kspace_style:index-5}}\label{\detokenize{kspace_style:index-4}}\label{\detokenize{kspace_style:index-3}}\label{\detokenize{kspace_style:index-2}}\label{\detokenize{kspace_style:index-1}}\label{\detokenize{kspace_style:index-0}}\label{\detokenize{kspace_style::doc}}

\subsection{Syntax}
\label{\detokenize{kspace_style:syntax}}
\begin{sphinxVerbatim}[commandchars=\\\{\}]
\PYG{k}{kspace\PYGZus{}style}\PYG{+w}{ }\PYG{n}{style}\PYG{+w}{ }\PYG{n}{value}
\end{sphinxVerbatim}
\begin{itemize}
\item {} 
\sphinxAtStartPar
style = \sphinxstyleemphasis{none} or \sphinxstyleemphasis{ewald} or \sphinxstyleemphasis{ewald/dipole} or \sphinxstyleemphasis{ewald/dipole/spin} or \sphinxstyleemphasis{ewald/disp} or \sphinxstyleemphasis{ewald/disp/dipole} or \sphinxstyleemphasis{ewald/omp} or \sphinxstyleemphasis{ewald/electrode} or \sphinxstyleemphasis{pppm} or \sphinxstyleemphasis{pppm/cg} or \sphinxstyleemphasis{pppm/disp} or \sphinxstyleemphasis{pppm/tip4p} or \sphinxstyleemphasis{pppm/stagger} or \sphinxstyleemphasis{pppm/disp/tip4p} or \sphinxstyleemphasis{pppm/gpu} or \sphinxstyleemphasis{pppm/intel} or \sphinxstyleemphasis{pppm/disp/intel} or \sphinxstyleemphasis{pppm/kk} or \sphinxstyleemphasis{pppm/omp} or \sphinxstyleemphasis{pppm/cg/omp} or \sphinxstyleemphasis{pppm/disp/tip4p/omp} or \sphinxstyleemphasis{pppm/tip4p/omp} or \sphinxstyleemphasis{pppm/dielectic} or \sphinxstyleemphasis{pppm/disp/dielectric} or \sphinxstyleemphasis{pppm/electrode} or \sphinxstyleemphasis{pppm/electrode/intel} or \sphinxstyleemphasis{msm} or \sphinxstyleemphasis{msm/cg} or \sphinxstyleemphasis{msm/omp} or \sphinxstyleemphasis{msm/cg/omp} or \sphinxstyleemphasis{msm/dielectric} or \sphinxstyleemphasis{scafacos}
\begin{sphinxalltt}
\sphinxstyleemphasis{none} value = none
\sphinxstyleemphasis{ewald} value = accuracy
  accuracy = desired relative error in forces
\sphinxstyleemphasis{ewald/dipole} value = accuracy
  accuracy = desired relative error in forces
\sphinxstyleemphasis{ewald/dipole/spin} value = accuracy
  accuracy = desired relative error in forces
\sphinxstyleemphasis{ewald/disp} value = accuracy
  accuracy = desired relative error in forces
\sphinxstyleemphasis{ewald/disp/dipole} value = accuracy
  accuracy = desired relative error in forces
\sphinxstyleemphasis{ewald/omp} value = accuracy
  accuracy = desired relative error in forces
\sphinxstyleemphasis{ewald/electrode} value = accuracy
  accuracy = desired relative error in forces
\sphinxstyleemphasis{pppm} value = accuracy
  accuracy = desired relative error in forces
\sphinxstyleemphasis{pppm/cg} values = accuracy (smallq)
  accuracy = desired relative error in forces
  smallq = cutoff for charges to be considered (optional) (charge units)
\sphinxstyleemphasis{pppm/dipole} value = accuracy
  accuracy = desired relative error in forces
\sphinxstyleemphasis{pppm/dipole/spin} value = accuracy
  accuracy = desired relative error in forces
\sphinxstyleemphasis{pppm/disp} value = accuracy
  accuracy = desired relative error in forces
\sphinxstyleemphasis{pppm/tip4p} value = accuracy
  accuracy = desired relative error in forces
\sphinxstyleemphasis{pppm/disp/tip4p} value = accuracy
  accuracy = desired relative error in forces
\sphinxstyleemphasis{pppm/gpu} value = accuracy
  accuracy = desired relative error in forces
\sphinxstyleemphasis{pppm/intel} value = accuracy
  accuracy = desired relative error in forces
\sphinxstyleemphasis{pppm/disp/intel} value = accuracy
  accuracy = desired relative error in forces
\sphinxstyleemphasis{pppm/kk} value = accuracy
  accuracy = desired relative error in forces
\sphinxstyleemphasis{pppm/omp} value = accuracy
  accuracy = desired relative error in forces
\sphinxstyleemphasis{pppm/cg/omp} values = accuracy (smallq)
  accuracy = desired relative error in forces
  smallq = cutoff for charges to be considered (optional) (charge units)
\sphinxstyleemphasis{pppm/disp/omp} value = accuracy
  accuracy = desired relative error in forces
\sphinxstyleemphasis{pppm/tip4p/omp} value = accuracy
  accuracy = desired relative error in forces
\sphinxstyleemphasis{pppm/disp/tip4p/omp} value = accuracy
  accuracy = desired relative error in forces
\sphinxstyleemphasis{pppm/stagger} value = accuracy
  accuracy = desired relative error in forces
\sphinxstyleemphasis{pppm/dielectric} value = accuracy
  accuracy = desired relative error in forces
\sphinxstyleemphasis{pppm/disp/dielectric} value = accuracy
  accuracy = desired relative error in forces
\sphinxstyleemphasis{pppm/electrode} value = accuracy
  accuracy = desired relative error in forces
\sphinxstyleemphasis{pppm/electrode/intel} value = accuracy
  accuracy = desired relative error in forces
\sphinxstyleemphasis{msm} value = accuracy
  accuracy = desired relative error in forces
\sphinxstyleemphasis{msm/cg} value = accuracy (smallq)
  accuracy = desired relative error in forces
  smallq = cutoff for charges to be considered (optional) (charge units)
\sphinxstyleemphasis{msm/omp} value = accuracy
  accuracy = desired relative error in forces
\sphinxstyleemphasis{msm/cg/omp} value = accuracy (smallq)
  accuracy = desired relative error in forces
  smallq = cutoff for charges to be considered (optional) (charge units)
\sphinxstyleemphasis{msm/dielectric} value = accuracy
  accuracy = desired relative error in forces
\sphinxstyleemphasis{scafacos} values = method accuracy
  method = fmm or p2nfft or p3m or ewald or direct
  accuracy = desired relative error in forces
\end{sphinxalltt}

\end{itemize}


\subsection{Examples}
\label{\detokenize{kspace_style:examples}}
\begin{sphinxVerbatim}[commandchars=\\\{\}]
\PYG{k}{kspace\PYGZus{}style}\PYG{+w}{ }\PYG{n}{pppm}\PYG{+w}{ }\PYG{l+m}{1.0e\PYGZhy{}4}
\PYG{k}{kspace\PYGZus{}style}\PYG{+w}{ }\PYG{n}{pppm}\PYG{o}{/}\PYG{n}{cg}\PYG{+w}{ }\PYG{l+m}{1.0e\PYGZhy{}5}\PYG{+w}{ }\PYG{l+m}{1.0e\PYGZhy{}6}
\PYG{k}{kspace\PYGZus{}style}\PYG{+w}{ }\PYG{n}{msm}\PYG{+w}{ }\PYG{l+m}{1.0e\PYGZhy{}4}
\PYG{k}{kspace\PYGZus{}style}\PYG{+w}{ }\PYG{n}{scafacos}\PYG{+w}{ }\PYG{n}{fmm}\PYG{+w}{ }\PYG{l+m}{1.0e\PYGZhy{}4}
\PYG{k}{kspace\PYGZus{}style}\PYG{+w}{ }\PYG{n}{none}
\end{sphinxVerbatim}

\sphinxAtStartPar
Used in input scripts:
\begin{quote}

\begin{sphinxVerbatim}[commandchars=\\\{\}]
examples/peptide/in.peptide
\end{sphinxVerbatim}
\end{quote}


\subsection{Description}
\label{\detokenize{kspace_style:description}}
\sphinxAtStartPar
Define a long\sphinxhyphen{}range solver for LAMMPS to use each timestep to compute
long\sphinxhyphen{}range Coulombic interactions or long\sphinxhyphen{}range \(1/r^6\) interactions.
Most of the long\sphinxhyphen{}range solvers perform their computation in K\sphinxhyphen{}space,
hence the name of this command.

\sphinxAtStartPar
When such a solver is used in conjunction with an appropriate pair
style, the cutoff for Coulombic or \(1/r^N\) interactions is effectively
infinite.  If the Coulombic case, this means each charge in the system
interacts with charges in an infinite array of periodic images of the
simulation domain.

\sphinxAtStartPar
Note that using a long\sphinxhyphen{}range solver requires use of a matching {\hyperref[\detokenize{pair_style::doc}]{\sphinxcrossref{\DUrole{doc}{pair style}}}} to perform consistent short\sphinxhyphen{}range pairwise
calculations.  This means that the name of the pair style contains a
matching keyword to the name of the KSpace style, as in this table:


\begin{savenotes}\sphinxattablestart
\sphinxthistablewithglobalstyle
\centering
\begin{tabulary}{\linewidth}[t]{TT}
\sphinxtoprule
\sphinxtableatstartofbodyhook
\sphinxAtStartPar
Pair style
&
\sphinxAtStartPar
KSpace style
\\
\sphinxhline
\sphinxAtStartPar
coul/long
&
\sphinxAtStartPar
ewald or pppm
\\
\sphinxhline
\sphinxAtStartPar
coul/msm
&
\sphinxAtStartPar
msm
\\
\sphinxhline
\sphinxAtStartPar
lj/long or buck/long
&
\sphinxAtStartPar
disp (for dispersion)
\\
\sphinxhline
\sphinxAtStartPar
tip4p/long
&
\sphinxAtStartPar
tip4p
\\
\sphinxhline
\sphinxAtStartPar
dipole/long
&
\sphinxAtStartPar
dipole
\\
\sphinxbottomrule
\end{tabulary}
\sphinxtableafterendhook\par
\sphinxattableend\end{savenotes}


\bigskip\hrule\bigskip


\sphinxAtStartPar
The \sphinxstyleemphasis{ewald} style performs a standard Ewald summation as described in
any solid\sphinxhyphen{}state physics text.

\sphinxAtStartPar
The \sphinxstyleemphasis{ewald/disp} style adds a long\sphinxhyphen{}range dispersion sum option for
\(1/r^6\) potentials and is useful for simulation of interfaces
{\hyperref[\detokenize{kspace_style:veld}]{\sphinxcrossref{\DUrole{std,std-ref}{(Veld)}}}}.  It also performs standard Coulombic Ewald summations,
but in a more efficient manner than the \sphinxstyleemphasis{ewald} style.  The \(1/r^6\)
capability means that Lennard\sphinxhyphen{}Jones or Buckingham potentials can be
used without a cutoff, i.e. they become full long\sphinxhyphen{}range potentials.

\sphinxAtStartPar
The \sphinxstyleemphasis{ewald/disp/dipole} style can also be used with point\sphinxhyphen{}dipoles, see
{\hyperref[\detokenize{kspace_style:toukmaji}]{\sphinxcrossref{\DUrole{std,std-ref}{(Toukmaji)}}}}.

\sphinxAtStartPar
The \sphinxstyleemphasis{ewald/dipole} style adds long\sphinxhyphen{}range standard Ewald summations
for dipole\sphinxhyphen{}dipole interactions, see {\hyperref[\detokenize{kspace_style:toukmaji}]{\sphinxcrossref{\DUrole{std,std-ref}{(Toukmaji)}}}}.

\sphinxAtStartPar
The \sphinxstyleemphasis{ewald/dipole/spin} style adds long\sphinxhyphen{}range standard Ewald
summations for magnetic dipole\sphinxhyphen{}dipole interactions between
magnetic spins.


\bigskip\hrule\bigskip


\sphinxAtStartPar
The \sphinxstyleemphasis{pppm} style invokes a particle\sphinxhyphen{}particle particle\sphinxhyphen{}mesh solver
{\hyperref[\detokenize{kspace_style:hockney}]{\sphinxcrossref{\DUrole{std,std-ref}{(Hockney)}}}} which maps atom charge to a 3d mesh, uses 3d FFTs
to solve Poisson’s equation on the mesh, then interpolates electric
fields on the mesh points back to the atoms.  It is closely related to
the particle\sphinxhyphen{}mesh Ewald technique (PME) {\hyperref[\detokenize{kspace_style:darden}]{\sphinxcrossref{\DUrole{std,std-ref}{(Darden)}}}} used in
AMBER and CHARMM.  The cost of traditional Ewald summation scales as
\(N^{\frac{3}{2}}\) where \(N\) is the number of atoms in the system.  The PPPM solver
scales as \(N \log{N}\) due to the FFTs, so it is almost always a faster
choice {\hyperref[\detokenize{kspace_style:pollock}]{\sphinxcrossref{\DUrole{std,std-ref}{(Pollock)}}}}.

\sphinxAtStartPar
The \sphinxstyleemphasis{pppm/cg} style is identical to the \sphinxstyleemphasis{pppm} style except that it
has an optimization for systems where most particles are uncharged.
Similarly the \sphinxstyleemphasis{msm/cg} style implements the same optimization for \sphinxstyleemphasis{msm}.
The optional \sphinxstyleemphasis{smallq} argument defines the cutoff for the absolute
charge value which determines whether a particle is considered charged
or not.  Its default value is 1.0e\sphinxhyphen{}5.

\sphinxAtStartPar
The \sphinxstyleemphasis{pppm/dipole} style invokes a particle\sphinxhyphen{}particle particle\sphinxhyphen{}mesh solver
for dipole\sphinxhyphen{}dipole interactions, following the method of {\hyperref[\detokenize{kspace_style:cerda2008}]{\sphinxcrossref{\DUrole{std,std-ref}{(Cerda)}}}}.

\sphinxAtStartPar
The \sphinxstyleemphasis{pppm/dipole/spin} style invokes a particle\sphinxhyphen{}particle particle\sphinxhyphen{}mesh solver
for magnetic dipole\sphinxhyphen{}dipole interactions between magnetic spins.

\sphinxAtStartPar
The \sphinxstyleemphasis{pppm/tip4p} style is identical to the \sphinxstyleemphasis{pppm} style except that it
adds a charge at the massless fourth site in each TIP4P water molecule.
It should be used with {\hyperref[\detokenize{pair_style::doc}]{\sphinxcrossref{\DUrole{doc}{pair styles}}}} with a
\sphinxstyleemphasis{tip4p/long} in their style name.

\sphinxAtStartPar
The \sphinxstyleemphasis{pppm/stagger} style performs calculations using two different
meshes, one shifted slightly with respect to the other.  This can
reduce force aliasing errors and increase the accuracy of the method
for a given mesh size.  Or a coarser mesh can be used for the same
target accuracy, which saves CPU time.  However, there is a trade\sphinxhyphen{}off
since FFTs on two meshes are now performed which increases the
computation required.  See {\hyperref[\detokenize{kspace_style:cerutti}]{\sphinxcrossref{\DUrole{std,std-ref}{(Cerutti)}}}}, {\hyperref[\detokenize{kspace_style:neelov}]{\sphinxcrossref{\DUrole{std,std-ref}{(Neelov)}}}},
and {\hyperref[\detokenize{kspace_style:hockney}]{\sphinxcrossref{\DUrole{std,std-ref}{(Hockney)}}}} for details of the method.

\sphinxAtStartPar
For high relative accuracy, using staggered PPPM allows the mesh size
to be reduced by a factor of 2 in each dimension as compared to
regular PPPM (for the same target accuracy).  This can give up to a 4x
speedup in the KSpace time (8x less mesh points, 2x more expensive).
However, for low relative accuracy, the staggered PPPM mesh size may
be essentially the same as for regular PPPM, which means the method
will be up to 2x slower in the KSpace time (simply 2x more expensive).
For more details and timings, see the {\hyperref[\detokenize{Speed_tips::doc}]{\sphinxcrossref{\DUrole{doc}{Speed tips}}}} doc
page.

\begin{sphinxadmonition}{note}{Note:}
\sphinxAtStartPar
Using \sphinxstyleemphasis{pppm/stagger} may not give the same increase in the
accuracy of energy and pressure as it does in forces, so some caution
must be used if energy and/or pressure are quantities of interest,
such as when using a barostat.
\end{sphinxadmonition}


\bigskip\hrule\bigskip


\sphinxAtStartPar
The \sphinxstyleemphasis{pppm/disp} and \sphinxstyleemphasis{pppm/disp/tip4p} styles add a mesh\sphinxhyphen{}based long\sphinxhyphen{}range
dispersion sum option for 1/r\textasciicircum{}6 potentials {\hyperref[\detokenize{kspace_style:isele-holder2012}]{\sphinxcrossref{\DUrole{std,std-ref}{(Isele\sphinxhyphen{}Holder)}}}},
similar to the \sphinxstyleemphasis{ewald/disp} style. The 1/r\textasciicircum{}6 capability means
that Lennard\sphinxhyphen{}Jones or Buckingham potentials can be used without a cutoff,
i.e. they become full long\sphinxhyphen{}range potentials.

\sphinxAtStartPar
For these styles, you will possibly want to adjust the default choice
of parameters by using the {\hyperref[\detokenize{kspace_modify::doc}]{\sphinxcrossref{\DUrole{doc}{kspace\_modify}}}} command.
This can be done by either choosing the Ewald and grid parameters, or
by specifying separate accuracies for the real and kspace
calculations. When not making any settings, the simulation will stop
with an error message. Further information on the influence of the
parameters and how to choose them is described in
{\hyperref[\detokenize{kspace_style:isele-holder2012}]{\sphinxcrossref{\DUrole{std,std-ref}{(Isele\sphinxhyphen{}Holder)}}}},
{\hyperref[\detokenize{kspace_style:isele-holder2013}]{\sphinxcrossref{\DUrole{std,std-ref}{(Isele\sphinxhyphen{}Holder2)}}}} and the {\hyperref[\detokenize{Howto_dispersion::doc}]{\sphinxcrossref{\DUrole{doc}{Howto dispersion}}}} doc page.


\bigskip\hrule\bigskip


\begin{sphinxadmonition}{note}{Note:}
\sphinxAtStartPar
All of the PPPM styles can be used with single\sphinxhyphen{}precision FFTs by
using the compiler switch \sphinxhyphen{}DFFT\_SINGLE for the FFT\_INC setting in your
low\sphinxhyphen{}level Makefile.  This setting also changes some of the PPPM
operations (e.g. mapping charge to mesh and interpolating electric
fields to particles) to be performed in single precision.  This option
can speed\sphinxhyphen{}up long\sphinxhyphen{}range calculations, particularly in parallel or on
GPUs.  The use of the \sphinxhyphen{}DFFT\_SINGLE flag is discussed on the {\hyperref[\detokenize{Build_settings::doc}]{\sphinxcrossref{\DUrole{doc}{Build settings}}}} doc page. MSM does not currently support
the \sphinxhyphen{}DFFT\_SINGLE compiler switch.
\end{sphinxadmonition}


\bigskip\hrule\bigskip


\sphinxAtStartPar
The \sphinxstyleemphasis{electrode} styles add methods that are required for the constant potential
method implemented in {\hyperref[\detokenize{fix_electrode::doc}]{\sphinxcrossref{\DUrole{doc}{fix electrode/*}}}}.  The styles
\sphinxstyleemphasis{ewald/electrode}, \sphinxstyleemphasis{pppm/electrode} and \sphinxstyleemphasis{pppm/electrode/intel} are available.
These styles do not support the \sphinxtitleref{kspace\_modify slab nozforce} command.


\bigskip\hrule\bigskip


\sphinxAtStartPar
The \sphinxstyleemphasis{msm} style invokes a multi\sphinxhyphen{}level summation method MSM solver,
{\hyperref[\detokenize{kspace_style:hardy2006}]{\sphinxcrossref{\DUrole{std,std-ref}{(Hardy)}}}} or {\hyperref[\detokenize{kspace_style:hardy2009}]{\sphinxcrossref{\DUrole{std,std-ref}{(Hardy2)}}}}, which maps atom charge
to a 3d mesh, and uses a multi\sphinxhyphen{}level hierarchy of coarser and coarser
meshes on which direct Coulomb solvers are done.  This method does not
use FFTs and scales as \(N\). It may therefore be faster than the other
K\sphinxhyphen{}space solvers for relatively large problems when running on large
core counts. MSM can also be used for non\sphinxhyphen{}periodic boundary conditions
and for mixed periodic and non\sphinxhyphen{}periodic boundaries.

\sphinxAtStartPar
MSM is most competitive versus Ewald and PPPM when only relatively
low accuracy forces, about 1e\sphinxhyphen{}4 relative error or less accurate,
are needed. Note that use of a larger Coulombic cutoff (i.e. 15
Angstroms instead of 10 Angstroms) provides better MSM accuracy for
both the real space and grid computed forces.

\sphinxAtStartPar
Currently calculation of the full pressure tensor in MSM is expensive.
Using the {\hyperref[\detokenize{kspace_modify::doc}]{\sphinxcrossref{\DUrole{doc}{kspace\_modify}}}} \sphinxstyleemphasis{pressure/scalar yes}
command provides a less expensive way to compute the scalar pressure
(Pxx + Pyy + Pzz)/3.0. The scalar pressure can be used, for example,
to run an isotropic barostat. If the full pressure tensor is needed,
then calculating the pressure at every timestep or using a fixed
pressure simulation with MSM will cause the code to run slower.


\bigskip\hrule\bigskip


\sphinxAtStartPar
The \sphinxstyleemphasis{scafacos} style is a wrapper on the \sphinxhref{http://www.scafacos.de}{ScaFaCoS Coulomb solver
library} which provides a variety of solver
methods which can be used with LAMMPS.  The paper by {\hyperref[\detokenize{kspace_style:sutmann2014}]{\sphinxcrossref{\DUrole{std,std-ref}{(Sutman)}}}} gives an overview of ScaFaCoS.

\sphinxAtStartPar
ScaFaCoS was developed by a consortium of German research facilities
with a BMBF (German Ministry of Science and Education) funded project
in 2009\sphinxhyphen{}2012. Participants of the consortium were the Universities of
Bonn, Chemnitz, Stuttgart, and Wuppertal as well as the
Forschungszentrum Juelich.

\sphinxAtStartPar
The library is available for download at “\sphinxurl{http://scafacos.de}” or can
be cloned from the git\sphinxhyphen{}repository
“\sphinxurl{https://github.com/scafacos/scafacos.git}”.

\sphinxAtStartPar
In order to use this KSpace style, you must download and build the
ScaFaCoS library, then build LAMMPS with the SCAFACOS package
installed package which links LAMMPS to the ScaFaCoS library.
See details on {\hyperref[\detokenize{Build_extras:scafacos}]{\sphinxcrossref{\DUrole{std,std-ref}{this page}}}}.

\begin{sphinxadmonition}{note}{Note:}
\sphinxAtStartPar
Unlike other KSpace solvers in LAMMPS, ScaFaCoS computes all
Coulombic interactions, both short\sphinxhyphen{} and long\sphinxhyphen{}range.  Thus you should
NOT use a Coulombic pair style when using kspace\_style scafacos.  This
also means the total Coulombic energy (short\sphinxhyphen{} and long\sphinxhyphen{}range) will be
tallied for {\hyperref[\detokenize{thermo_style::doc}]{\sphinxcrossref{\DUrole{doc}{thermodynamic output}}}} command as part
of the \sphinxstyleemphasis{elong} keyword; the \sphinxstyleemphasis{ecoul} keyword will be zero.
\end{sphinxadmonition}

\begin{sphinxadmonition}{note}{Note:}
\sphinxAtStartPar
See the current restriction below about use of ScaFaCoS in
LAMMPS with molecular charged systems or the TIP4P water model.
\end{sphinxadmonition}

\sphinxAtStartPar
The specified \sphinxstyleemphasis{method} determines which ScaFaCoS algorithm is used.
These are the ScaFaCoS methods currently available from LAMMPS:
\begin{itemize}
\item {} 
\sphinxAtStartPar
\sphinxstyleemphasis{fmm} = Fast Multi\sphinxhyphen{}Pole method

\item {} 
\sphinxAtStartPar
\sphinxstyleemphasis{p2nfft} = FFT\sphinxhyphen{}based Coulomb solver

\item {} 
\sphinxAtStartPar
\sphinxstyleemphasis{ewald} = Ewald summation

\item {} 
\sphinxAtStartPar
\sphinxstyleemphasis{direct} = direct O(N\textasciicircum{}2) summation

\item {} 
\sphinxAtStartPar
\sphinxstyleemphasis{p3m} = PPPM

\end{itemize}

\sphinxAtStartPar
We plan to support additional ScaFaCoS solvers from LAMMPS in the
future.  For an overview of the included solvers, refer to
{\hyperref[\detokenize{kspace_style:sutmann2013}]{\sphinxcrossref{\DUrole{std,std-ref}{(Sutmann)}}}}

\sphinxAtStartPar
The specified \sphinxstyleemphasis{accuracy} is similar to the accuracy setting for other
LAMMPS KSpace styles, but is passed to ScaFaCoS, which can interpret
it in different ways for different methods it supports.  Within the
ScaFaCoS library the \sphinxstyleemphasis{accuracy} is treated as a tolerance level
(either absolute or relative) for the chosen quantity, where the
quantity can be either the Columic field values, the per\sphinxhyphen{}atom Columic
energy or the total Columic energy.  To select from these options, see
the {\hyperref[\detokenize{kspace_modify::doc}]{\sphinxcrossref{\DUrole{doc}{kspace\_modify scafacos accuracy}}}} doc page.

\sphinxAtStartPar
The {\hyperref[\detokenize{kspace_modify::doc}]{\sphinxcrossref{\DUrole{doc}{kspace\_modify scafacos}}}} command also explains
other ScaFaCoS options currently exposed to LAMMPS.


\bigskip\hrule\bigskip


\sphinxAtStartPar
The specified \sphinxstyleemphasis{accuracy} determines the relative RMS error in per\sphinxhyphen{}atom
forces calculated by the long\sphinxhyphen{}range solver.  It is set as a
dimensionless number, relative to the force that two unit point
charges (e.g. 2 monovalent ions) exert on each other at a distance of
1 Angstrom.  This reference value was chosen as representative of the
magnitude of electrostatic forces in atomic systems.  Thus an accuracy
value of 1.0e\sphinxhyphen{}4 means that the RMS error will be a factor of 10000
smaller than the reference force.

\sphinxAtStartPar
The accuracy setting is used in conjunction with the pairwise cutoff
to determine the number of K\sphinxhyphen{}space vectors for style \sphinxstyleemphasis{ewald} or the
grid size for style \sphinxstyleemphasis{pppm} or \sphinxstyleemphasis{msm}.

\sphinxAtStartPar
Note that style \sphinxstyleemphasis{pppm} only computes the grid size at the beginning of
a simulation, so if the length or triclinic tilt of the simulation
cell increases dramatically during the course of the simulation, the
accuracy of the simulation may degrade.  Likewise, if the
{\hyperref[\detokenize{kspace_modify::doc}]{\sphinxcrossref{\DUrole{doc}{kspace\_modify slab}}}} option is used with
shrink\sphinxhyphen{}wrap boundaries in the z\sphinxhyphen{}dimension, and the box size changes
dramatically in z.  For example, for a triclinic system with all three
tilt factors set to the maximum limit, the PPPM grid should be
increased roughly by a factor of 1.5 in the y direction and 2.0 in the
z direction as compared to the same system using a cubic orthogonal
simulation cell.  One way to handle this issue if you have a long
simulation where the box size changes dramatically, is to break it
into shorter simulations (multiple {\hyperref[\detokenize{run::doc}]{\sphinxcrossref{\DUrole{doc}{run}}}} commands).  This
works because the grid size is re\sphinxhyphen{}computed at the beginning of each
run.  Another way to ensure the described accuracy requirement is met
is to run a short simulation at the maximum expected tilt or length,
note the required grid size, and then use the
{\hyperref[\detokenize{kspace_modify::doc}]{\sphinxcrossref{\DUrole{doc}{kspace\_modify}}}} \sphinxstyleemphasis{mesh} command to manually set the
PPPM grid size to this value for the long run.  The simulation then
will be “too accurate” for some portion of the run.

\sphinxAtStartPar
RMS force errors in real space for \sphinxstyleemphasis{ewald} and \sphinxstyleemphasis{pppm} are estimated
using equation 18 of {\hyperref[\detokenize{kspace_style:kolafa}]{\sphinxcrossref{\DUrole{std,std-ref}{(Kolafa)}}}}, which is also referenced as
equation 9 of {\hyperref[\detokenize{kspace_style:petersen}]{\sphinxcrossref{\DUrole{std,std-ref}{(Petersen)}}}}. RMS force errors in K\sphinxhyphen{}space for
\sphinxstyleemphasis{ewald} are estimated using equation 11 of {\hyperref[\detokenize{kspace_style:petersen}]{\sphinxcrossref{\DUrole{std,std-ref}{(Petersen)}}}},
which is similar to equation 32 of {\hyperref[\detokenize{kspace_style:kolafa}]{\sphinxcrossref{\DUrole{std,std-ref}{(Kolafa)}}}}. RMS force
errors in K\sphinxhyphen{}space for \sphinxstyleemphasis{pppm} are estimated using equation 38 of
{\hyperref[\detokenize{kspace_style:deserno}]{\sphinxcrossref{\DUrole{std,std-ref}{(Deserno)}}}}. RMS force errors for \sphinxstyleemphasis{msm} are estimated
using ideas from chapter 3 of {\hyperref[\detokenize{kspace_style:hardy2006}]{\sphinxcrossref{\DUrole{std,std-ref}{(Hardy)}}}}, with equation 3.197
of particular note. When using \sphinxstyleemphasis{msm} with non\sphinxhyphen{}periodic boundary
conditions, it is expected that the error estimation will be too
pessimistic. RMS force errors for dipoles when using \sphinxstyleemphasis{ewald/disp}
or \sphinxstyleemphasis{ewald/dipole} are estimated using equations 33 and 46 of
{\hyperref[\detokenize{kspace_style:wang}]{\sphinxcrossref{\DUrole{std,std-ref}{(Wang)}}}}. The RMS force errors for \sphinxstyleemphasis{pppm/dipole} are estimated
using the equations in {\hyperref[\detokenize{kspace_style:cerda2008}]{\sphinxcrossref{\DUrole{std,std-ref}{(Cerda)}}}}.

\sphinxAtStartPar
See the {\hyperref[\detokenize{kspace_modify::doc}]{\sphinxcrossref{\DUrole{doc}{kspace\_modify}}}} command for additional
options of the K\sphinxhyphen{}space solvers that can be set, including a \sphinxstyleemphasis{force}
option for setting an absolute RMS error in forces, as opposed to a
relative RMS error.


\bigskip\hrule\bigskip


\sphinxAtStartPar
Styles with a \sphinxstyleemphasis{gpu}, \sphinxstyleemphasis{intel}, \sphinxstyleemphasis{kk}, \sphinxstyleemphasis{omp}, or \sphinxstyleemphasis{opt} suffix are
functionally the same as the corresponding style without the suffix.
They have been optimized to run faster, depending on your available
hardware, as discussed on the {\hyperref[\detokenize{Speed_packages::doc}]{\sphinxcrossref{\DUrole{doc}{Accelerator packages}}}}
page.  The accelerated styles take the same arguments and should
produce the same results, except for round\sphinxhyphen{}off and precision issues.

\sphinxAtStartPar
These accelerated styles are part of the GPU, INTEL, KOKKOS,
OPENMP, and OPT packages, respectively.  They are only enabled if
LAMMPS was built with those packages.  See the {\hyperref[\detokenize{Build_package::doc}]{\sphinxcrossref{\DUrole{doc}{Build package}}}} page for more info.

\sphinxAtStartPar
You can specify the accelerated styles explicitly in your input script
by including their suffix, or you can use the {\hyperref[\detokenize{Run_options::doc}]{\sphinxcrossref{\DUrole{doc}{\sphinxhyphen{}suffix command\sphinxhyphen{}line switch}}}} when you invoke LAMMPS, or you can use the
{\hyperref[\detokenize{suffix::doc}]{\sphinxcrossref{\DUrole{doc}{suffix}}}} command in your input script.

\sphinxAtStartPar
See the {\hyperref[\detokenize{Speed_packages::doc}]{\sphinxcrossref{\DUrole{doc}{Accelerator packages}}}} page for more
instructions on how to use the accelerated styles effectively.

\begin{sphinxadmonition}{note}{Note:}
\sphinxAtStartPar
For the GPU package, the \sphinxstyleemphasis{pppm/gpu} style performs charge assignment
and force interpolation calculations on the GPU.  These processes
are performed either in single or double precision, depending on
whether the \sphinxhyphen{}DFFT\_SINGLE setting was specified in your low\sphinxhyphen{}level
Makefile, as discussed above.  The FFTs themselves are still
calculated on the CPU.  If \sphinxstyleemphasis{pppm/gpu} is used with a GPU\sphinxhyphen{}enabled
pair style, part of the PPPM calculation can be performed
concurrently on the GPU while other calculations for non\sphinxhyphen{}bonded and
bonded force calculation are performed on the CPU.
\end{sphinxadmonition}

\begin{sphinxadmonition}{note}{Note:}
\sphinxAtStartPar
For the KOKKOS package, the \sphinxstyleemphasis{pppm/kk} style performs charge
assignment and force interpolation calculations, along with the FFTs
themselves, on the GPU or (optionally) threaded on the CPU when
using OpenMP and FFTW3. The specific FFT library is selected using
the FFT\_KOKKOS CMake parameter. See the
{\hyperref[\detokenize{Build_settings::doc}]{\sphinxcrossref{\DUrole{doc}{Build settings}}}} doc page for how to select a
3rd\sphinxhyphen{}party FFT library.
\end{sphinxadmonition}


\bigskip\hrule\bigskip



\subsection{Restrictions}
\label{\detokenize{kspace_style:restrictions}}
\sphinxAtStartPar
Note that the long\sphinxhyphen{}range electrostatic solvers in LAMMPS assume conducting
metal (tinfoil) boundary conditions for both charge and dipole
interactions. Vacuum boundary conditions are not currently supported.

\sphinxAtStartPar
The \sphinxstyleemphasis{ewald/disp}, \sphinxstyleemphasis{ewald}, \sphinxstyleemphasis{pppm}, and \sphinxstyleemphasis{msm} styles support
non\sphinxhyphen{}orthogonal (triclinic symmetry) simulation boxes. However,
triclinic simulation cells may not yet be supported by all suffix
versions of these styles.

\sphinxAtStartPar
Most of the base kspace styles are part of the KSPACE package.  They are
only enabled if LAMMPS was built with that package.  See the {\hyperref[\detokenize{Build_package::doc}]{\sphinxcrossref{\DUrole{doc}{Build
package}}}} page for more info.

\sphinxAtStartPar
The \sphinxstyleemphasis{msm/dielectric} and \sphinxstyleemphasis{pppm/dielectric} kspace styles are part of the
DIELECTRIC package. They are only enabled if LAMMPS was built with
that package \sphinxstylestrong{and} the KSPACE package.  See the {\hyperref[\detokenize{Build_package::doc}]{\sphinxcrossref{\DUrole{doc}{Build package}}}} page for more info.

\sphinxAtStartPar
For MSM, a simulation must be 3d and one can use any combination of
periodic, non\sphinxhyphen{}periodic, but not shrink\sphinxhyphen{}wrapped boundaries (specified
using the {\hyperref[\detokenize{boundary::doc}]{\sphinxcrossref{\DUrole{doc}{boundary}}}} command).

\sphinxAtStartPar
For Ewald and PPPM, a simulation must be 3d and periodic in all
dimensions.  The only exception is if the slab option is set with
{\hyperref[\detokenize{kspace_modify::doc}]{\sphinxcrossref{\DUrole{doc}{kspace\_modify}}}}, in which case the xy dimensions
must be periodic and the z dimension must be non\sphinxhyphen{}periodic.

\sphinxAtStartPar
The scafacos KSpace style will only be enabled if LAMMPS is built with
the SCAFACOS package.  See the {\hyperref[\detokenize{Build_package::doc}]{\sphinxcrossref{\DUrole{doc}{Build package}}}}
doc page for more info.

\sphinxAtStartPar
The use of ScaFaCos in LAMMPS does not yet support molecular charged
systems where the short\sphinxhyphen{}range Coulombic interactions between atoms in
the same bond/angle/dihedral are weighted by the
{\hyperref[\detokenize{special_bonds::doc}]{\sphinxcrossref{\DUrole{doc}{special\_bonds}}}} command.  Likewise it does not
support the “TIP4P water style” where a fictitious charge site is
introduced in each water molecule.
Finally, the methods \sphinxstyleemphasis{p3m} and \sphinxstyleemphasis{ewald} do not support computing the
virial, so this contribution is not included.


\subsection{Related commands}
\label{\detokenize{kspace_style:related-commands}}
\sphinxAtStartPar
{\hyperref[\detokenize{kspace_modify::doc}]{\sphinxcrossref{\DUrole{doc}{kspace\_modify}}}}, {\hyperref[\detokenize{pair_lj_cut_coul::doc}]{\sphinxcrossref{\DUrole{doc}{pair\_style lj/cut/coul/long}}}}, {\hyperref[\detokenize{pair_charmm::doc}]{\sphinxcrossref{\DUrole{doc}{pair\_style lj/charmm/coul/long}}}}, {\hyperref[\detokenize{pair_lj_long::doc}]{\sphinxcrossref{\DUrole{doc}{pair\_style lj/long/coul/long}}}}, {\hyperref[\detokenize{pair_buck::doc}]{\sphinxcrossref{\DUrole{doc}{pair\_style buck/coul/long}}}}


\subsection{Default}
\label{\detokenize{kspace_style:default}}
\begin{sphinxVerbatim}[commandchars=\\\{\}]
\PYG{k}{kspace\PYGZus{}style}\PYG{+w}{ }\PYG{n}{none}
\end{sphinxVerbatim}


\bigskip\hrule\bigskip

\phantomsection\label{\detokenize{kspace_style:darden}}
\sphinxAtStartPar
\sphinxstylestrong{(Darden)} Darden, York, Pedersen, J Chem Phys, 98, 10089 (1993).

\phantomsection\label{\detokenize{kspace_style:deserno}}
\sphinxAtStartPar
\sphinxstylestrong{(Deserno)} Deserno and Holm, J Chem Phys, 109, 7694 (1998).

\phantomsection\label{\detokenize{kspace_style:hockney}}
\sphinxAtStartPar
\sphinxstylestrong{(Hockney)} Hockney and Eastwood, Computer Simulation Using Particles,
Adam Hilger, NY (1989).

\phantomsection\label{\detokenize{kspace_style:kolafa}}
\sphinxAtStartPar
\sphinxstylestrong{(Kolafa)} Kolafa and Perram, Molecular Simulation, 9, 351 (1992).

\phantomsection\label{\detokenize{kspace_style:petersen}}
\sphinxAtStartPar
\sphinxstylestrong{(Petersen)} Petersen, J Chem Phys, 103, 3668 (1995).

\phantomsection\label{\detokenize{kspace_style:wang}}
\sphinxAtStartPar
\sphinxstylestrong{(Wang)} Wang and Holm, J Chem Phys, 115, 6277 (2001).

\phantomsection\label{\detokenize{kspace_style:pollock}}
\sphinxAtStartPar
\sphinxstylestrong{(Pollock)} Pollock and Glosli, Comp Phys Comm, 95, 93 (1996).

\phantomsection\label{\detokenize{kspace_style:cerutti}}
\sphinxAtStartPar
\sphinxstylestrong{(Cerutti)} Cerutti, Duke, Darden, Lybrand, Journal of Chemical Theory
and Computation 5, 2322 (2009)

\phantomsection\label{\detokenize{kspace_style:neelov}}
\sphinxAtStartPar
\sphinxstylestrong{(Neelov)} Neelov, Holm, J Chem Phys 132, 234103 (2010)

\phantomsection\label{\detokenize{kspace_style:veld}}
\sphinxAtStartPar
\sphinxstylestrong{(Veld)} In ‘t Veld, Ismail, Grest, J Chem Phys, 127, 144711 (2007).

\phantomsection\label{\detokenize{kspace_style:toukmaji}}
\sphinxAtStartPar
\sphinxstylestrong{(Toukmaji)} Toukmaji, Sagui, Board, and Darden, J Chem Phys, 113,
10913 (2000).

\phantomsection\label{\detokenize{kspace_style:isele-holder2012}}
\sphinxAtStartPar
\sphinxstylestrong{(Isele\sphinxhyphen{}Holder)} Isele\sphinxhyphen{}Holder, Mitchell, Ismail, J Chem Phys, 137,
174107 (2012).

\phantomsection\label{\detokenize{kspace_style:isele-holder2013}}
\sphinxAtStartPar
\sphinxstylestrong{(Isele\sphinxhyphen{}Holder2)} Isele\sphinxhyphen{}Holder, Mitchell, Hammond, Kohlmeyer, Ismail,
J Chem Theory Comput 9, 5412 (2013).

\phantomsection\label{\detokenize{kspace_style:hardy2006}}
\sphinxAtStartPar
\sphinxstylestrong{(Hardy)} David Hardy thesis: Multilevel Summation for the Fast
Evaluation of Forces for the Simulation of Biomolecules, University of
Illinois at Urbana\sphinxhyphen{}Champaign, (2006).

\phantomsection\label{\detokenize{kspace_style:hardy2009}}
\sphinxAtStartPar
\sphinxstylestrong{(Hardy2)} Hardy, Stone, Schulten, Parallel Computing, 35, 164\sphinxhyphen{}177
(2009).

\phantomsection\label{\detokenize{kspace_style:sutmann2013}}
\sphinxAtStartPar
\sphinxstylestrong{(Sutmann)} Sutmann, Arnold, Fahrenberger, et. al., Physical review / E 88(6), 063308 (2013)

\phantomsection\label{\detokenize{kspace_style:cerda2008}}
\sphinxAtStartPar
\sphinxstylestrong{(Cerda)} Cerda, Ballenegger, Lenz, Holm, J Chem Phys 129, 234104 (2008)

\phantomsection\label{\detokenize{kspace_style:sutmann2014}}\begin{description}
\sphinxlineitem{\sphinxstylestrong{(Sutmann)} G. Sutmann. ScaFaCoS \sphinxhyphen{} a Scalable library of Fast Coulomb Solvers for particle Systems.}
\sphinxAtStartPar
In Bajaj, Zavattieri, Koslowski, Siegmund, Proceedings of the Society of Engineering Science 51st Annual Technical Meeting. 2014.

\end{description}

\sphinxstepscope

\index{label@\spxentry{label}}\ignorespaces 

\section{label command}
\label{\detokenize{label:label-command}}\label{\detokenize{label:index-0}}\label{\detokenize{label::doc}}

\subsection{Syntax}
\label{\detokenize{label:syntax}}
\begin{sphinxVerbatim}[commandchars=\\\{\}]
\PYG{k}{label }\PYG{l+s+sc}{ID}
\end{sphinxVerbatim}
\begin{itemize}
\item {} 
\sphinxAtStartPar
ID = string used as label name

\end{itemize}


\subsection{Examples}
\label{\detokenize{label:examples}}
\begin{sphinxVerbatim}[commandchars=\\\{\}]
\PYG{k}{label }\PYG{l+s+sc}{xyz}
\PYG{k}{label }\PYG{l+s+sc}{loop}
\end{sphinxVerbatim}


\subsection{Description}
\label{\detokenize{label:description}}
\sphinxAtStartPar
Label this line of the input script with the chosen ID.  Unless a jump
command was used previously, this does nothing.  But if a
{\hyperref[\detokenize{jump::doc}]{\sphinxcrossref{\DUrole{doc}{jump}}}} command was used with a label argument to begin
invoking this script file, then all command lines in the script prior
to this line will be ignored.  I.e. execution of the script will begin
at this line.  This is useful for looping over a section of the input
script as discussed in the {\hyperref[\detokenize{jump::doc}]{\sphinxcrossref{\DUrole{doc}{jump}}}} command.


\subsection{Restrictions}
\label{\detokenize{label:restrictions}}\begin{quote}

\sphinxAtStartPar
none
\end{quote}


\subsection{Related commands}
\label{\detokenize{label:related-commands}}
\sphinxAtStartPar
{\hyperref[\detokenize{jump::doc}]{\sphinxcrossref{\DUrole{doc}{jump}}}}, {\hyperref[\detokenize{next::doc}]{\sphinxcrossref{\DUrole{doc}{next}}}}


\subsection{Default}
\label{\detokenize{label:default}}
\sphinxAtStartPar
none

\sphinxstepscope

\index{labelmap@\spxentry{labelmap}}\ignorespaces 

\section{labelmap command}
\label{\detokenize{labelmap:labelmap-command}}\label{\detokenize{labelmap:index-0}}\label{\detokenize{labelmap::doc}}

\subsection{Syntax}
\label{\detokenize{labelmap:syntax}}
\begin{sphinxVerbatim}[commandchars=\\\{\}]
\PYG{k}{labelmap}\PYG{+w}{ }\PYG{n}{option}\PYG{+w}{ }\PYG{n}{args}
\end{sphinxVerbatim}
\begin{itemize}
\item {} 
\sphinxAtStartPar
\sphinxstyleemphasis{option} = \sphinxstyleemphasis{atom} or \sphinxstyleemphasis{bond} or \sphinxstyleemphasis{angle} or \sphinxstyleemphasis{dihedral} or \sphinxstyleemphasis{improper} or \sphinxstyleemphasis{clear} or \sphinxstyleemphasis{write}
\begin{sphinxalltt}
\sphinxstyleemphasis{clear} = no args
\sphinxstyleemphasis{write} arg = filename
\sphinxstyleemphasis{atom} or \sphinxstyleemphasis{bond} or \sphinxstyleemphasis{angle} or \sphinxstyleemphasis{dihedral} or \sphinxstyleemphasis{improper}
  args = list of one or more numeric\sphinxhyphen{}type/type\sphinxhyphen{}label pairs
\end{sphinxalltt}

\end{itemize}


\subsection{Examples}
\label{\detokenize{labelmap:examples}}
\begin{sphinxVerbatim}[commandchars=\\\{\}]
\PYG{k}{labelmap}\PYG{+w}{ }\PYG{n}{atom}\PYG{+w}{ }\PYG{l+m}{1}\PYG{+w}{ }\PYG{n}{c1}\PYG{+w}{ }\PYG{l+m}{2}\PYG{+w}{ }\PYG{n}{hc}\PYG{+w}{ }\PYG{l+m}{3}\PYG{+w}{ }\PYG{n}{cp}\PYG{+w}{ }\PYG{l+m}{4}\PYG{+w}{ }\PYG{n}{nt}
\PYG{k}{labelmap}\PYG{+w}{ }\PYG{n}{atom}\PYG{+w}{ }\PYG{l+m}{3}\PYG{+w}{ }\PYG{n}{carbon}\PYG{+w}{ }\PYG{l+m}{4}\PYG{+w}{ }\PYG{l+s}{\PYGZsq{}}\PYG{l+s}{c3\PYGZdq{}}\PYG{l+s}{\PYGZsq{}}\PYG{+w}{ }\PYG{l+m}{5}\PYG{+w}{ }\PYG{l+s}{\PYGZdq{}}\PYG{l+s}{c1\PYGZsq{}}\PYG{l+s}{\PYGZdq{}}\PYG{+w}{ }\PYG{l+m}{6}\PYG{+w}{ }\PYG{l+s}{\PYGZdq{}}\PYG{l+s}{c\PYGZsh{}}\PYG{l+s}{\PYGZdq{}}
\PYG{k}{labelmap}\PYG{+w}{ }\PYG{n}{atom}\PYG{+w}{ }\PYG{n+nv}{\PYGZdl{}(}\PYG{n+nv}{label2type}\PYG{n+nv}{(}\PYG{n+nv}{atom,carbon}\PYG{n+nv}{)}\PYG{n+nv}{)}\PYG{+w}{ }\PYG{n}{C}\PYG{+w}{  }\PYG{c}{\PYGZsh{} change type label from \PYGZsq{}carbon\PYGZsq{} to \PYGZsq{}C\PYGZsq{}}
\PYG{k}{labelmap}\PYG{+w}{ }\PYG{n}{clear}
\PYG{k}{labelmap}\PYG{+w}{ }\PYG{n}{write}\PYG{+w}{ }\PYG{n}{mymap.include}
\PYG{k}{labelmap}\PYG{+w}{ }\PYG{n}{bond}\PYG{+w}{ }\PYG{l+m}{1}\PYG{+w}{ }\PYG{n}{carbonyl}\PYG{+w}{ }\PYG{l+m}{2}\PYG{+w}{ }\PYG{n}{nitrile}\PYG{+w}{ }\PYG{l+m}{3}\PYG{+w}{ }\PYG{l+s}{\PYGZdq{}}\PYG{l+s}{\PYGZdq{}}\PYG{l+s}{\PYGZdq{}}\PYG{l+s}{ c1\PYGZsq{}\PYGZhy{}c2}\PYG{l+s}{\PYGZdq{}}\PYG{+w}{ }\PYG{l+s}{\PYGZdq{}}\PYG{l+s}{\PYGZdq{}}\PYG{l+s}{\PYGZdq{}}
\end{sphinxVerbatim}


\subsection{Description}
\label{\detokenize{labelmap:description}}
\sphinxAtStartPar
\DUrole{versionmodified,added}{Added in version 15Sep2022.}

\sphinxAtStartPar
Define alphanumeric type labels to associate with one or more numeric
atom, bond, angle, dihedral or improper types.  A collection of type
labels for all atom types, bond types, etc. is stored as a label map.

\sphinxAtStartPar
The label map can also be defined by the {\hyperref[\detokenize{read_data::doc}]{\sphinxcrossref{\DUrole{doc}{read\_data}}}}
command when it reads these sections in a data file: Atom Type Labels,
Bond Type Labels, etc.  See the {\hyperref[\detokenize{Howto_type_labels::doc}]{\sphinxcrossref{\DUrole{doc}{Howto type labels}}}} doc page for a general discussion of how type
labels can be used.  See {\hyperref[\detokenize{labelmap:typelabel1}]{\sphinxcrossref{\DUrole{std,std-ref}{(Gissinger)}}}} for a
discussion of the type label implementation in LAMMPS and its uses.

\sphinxAtStartPar
Valid type labels can contain any alphanumeric character, but must not
start with a number, a ‘\#’, or a ‘*’ character.  They can contain other
standard ASCII characters such as angular or square brackets ‘\textless{}’ and ‘\textgreater{}’
or ‘{[}’ and ‘{]}’, parenthesis ‘(’ and ‘)’, dash ‘\sphinxhyphen{}’, underscore ‘\_’, plus
‘+’ and equals ‘=’ signs and more.  They must not contain blanks or any
other whitespace.  Note that type labels must be put in single or double
quotation marks if they contain the ‘\#’ character or if they contain a
double (”) or single quotation mark (‘).  If the label contains both
a single and a double quotation mark, then triple quotation (“””) must
be used.  When enclosing a type label with quotation marks, the
LAMMPS input parser may require adding leading or trailing blanks
around the type label so it can identify the enclosing quotation
marks.  Those blanks will be removed when defining the label.

\sphinxAtStartPar
A \sphinxstyleemphasis{labelmap} command can only modify the label map for one type\sphinxhyphen{}kind
(atom types, bond types, etc).  Any number of numeric\sphinxhyphen{}type/type\sphinxhyphen{}label
pairs may follow.  If a type label already exists for the same numeric
type, it will be overwritten.  Type labels must be unique; assigning the
same type label to multiple numeric types within the same type\sphinxhyphen{}kind is
not allowed.  When reading and writing data files, it is required that
there is a label defined for \sphinxstyleemphasis{every} numeric type within a given
type\sphinxhyphen{}kind in order to write out the type label section for that
type\sphinxhyphen{}kind.

\sphinxAtStartPar
The \sphinxstyleemphasis{clear} option resets the label map and thus discards all previous
settings.

\sphinxAtStartPar
The \sphinxstyleemphasis{write} option takes a filename as argument and writes the current
label mappings to a file as a sequence of \sphinxstyleemphasis{labelmap} commands, so the
file can be copied into a new LAMMPS input file or read in using the
{\hyperref[\detokenize{include::doc}]{\sphinxcrossref{\DUrole{doc}{include}}}} command.


\bigskip\hrule\bigskip



\subsection{Restrictions}
\label{\detokenize{labelmap:restrictions}}
\sphinxAtStartPar
This command must come after the simulation box is defined by a
{\hyperref[\detokenize{read_data::doc}]{\sphinxcrossref{\DUrole{doc}{read\_data}}}}, {\hyperref[\detokenize{read_restart::doc}]{\sphinxcrossref{\DUrole{doc}{read\_restart}}}}, or
{\hyperref[\detokenize{create_box::doc}]{\sphinxcrossref{\DUrole{doc}{create\_box}}}} command.

\sphinxAtStartPar
Label maps are currently not supported when using the KOKKOS package.


\subsection{Related commands}
\label{\detokenize{labelmap:related-commands}}
\sphinxAtStartPar
{\hyperref[\detokenize{read_data::doc}]{\sphinxcrossref{\DUrole{doc}{read\_data}}}}, {\hyperref[\detokenize{write_data::doc}]{\sphinxcrossref{\DUrole{doc}{write\_data}}}},
{\hyperref[\detokenize{molecule::doc}]{\sphinxcrossref{\DUrole{doc}{molecule}}}}, {\hyperref[\detokenize{fix_bond_react::doc}]{\sphinxcrossref{\DUrole{doc}{fix bond/react}}}}


\subsection{Default}
\label{\detokenize{labelmap:default}}
\sphinxAtStartPar
none


\bigskip\hrule\bigskip

\phantomsection\label{\detokenize{labelmap:typelabel1}}
\sphinxAtStartPar
\sphinxstylestrong{(Gissinger)} J. R. Gissinger, I. Nikiforov, Y. Afshar, B. Waters, M. Choi, D. S. Karls, A. Stukowski, W. Im, H. Heinz, A. Kohlmeyer, and E. B. Tadmor, J Phys Chem B, 128, 3282\sphinxhyphen{}3297 (2024).

\sphinxstepscope

\index{lattice@\spxentry{lattice}}\ignorespaces 

\section{lattice command}
\label{\detokenize{lattice:lattice-command}}\label{\detokenize{lattice:index-0}}\label{\detokenize{lattice::doc}}

\subsection{Syntax}
\label{\detokenize{lattice:syntax}}
\begin{sphinxVerbatim}[commandchars=\\\{\}]
\PYG{k}{lattice}\PYG{+w}{ }\PYG{n}{style}\PYG{+w}{ }\PYG{n}{scale}\PYG{+w}{ }\PYG{n}{keyword}\PYG{+w}{ }\PYG{n}{values}\PYG{+w}{ }\PYG{n}{...}
\end{sphinxVerbatim}
\begin{itemize}
\item {} 
\sphinxAtStartPar
style = \sphinxstyleemphasis{none} or \sphinxstyleemphasis{sc} or \sphinxstyleemphasis{bcc} or \sphinxstyleemphasis{fcc} or \sphinxstyleemphasis{hcp} or \sphinxstyleemphasis{diamond} or         \sphinxstyleemphasis{sq} or \sphinxstyleemphasis{sq2} or \sphinxstyleemphasis{hex} or \sphinxstyleemphasis{custom}

\item {} 
\sphinxAtStartPar
scale = scale factor between lattice and simulation box
\begin{sphinxalltt}
scale = reduced density rho* (for LJ units)
scale = lattice constant in distance units (for all other units)
\end{sphinxalltt}

\item {} 
\sphinxAtStartPar
zero or more keyword/value pairs may be appended

\item {} 
\sphinxAtStartPar
keyword = \sphinxstyleemphasis{origin} or \sphinxstyleemphasis{orient} or \sphinxstyleemphasis{spacing} or \sphinxstyleemphasis{a1} or \sphinxstyleemphasis{a2} or \sphinxstyleemphasis{a3} or \sphinxstyleemphasis{basis} or \sphinxstyleemphasis{triclinic/general}
\begin{sphinxalltt}
\sphinxstyleemphasis{origin} values = x y z
  x,y,z = fractions of a unit cell (0 \textless{}= x,y,z \textless{} 1)
\sphinxstyleemphasis{orient} values = dim i j k
  dim = \sphinxstyleemphasis{x} or \sphinxstyleemphasis{y} or \sphinxstyleemphasis{z}
  i,j,k = integer lattice directions
\sphinxstyleemphasis{spacing} values = dx dy dz
  dx,dy,dz = lattice spacings in the x,y,z box directions
\sphinxstyleemphasis{a1},\sphinxstyleemphasis{a2},\sphinxstyleemphasis{a3} values = x y z
  x,y,z = primitive vector components that define unit cell
\sphinxstyleemphasis{basis} values = x y z
  x,y,z = fractional coords of a basis atom (0 \textless{}= x,y,z \textless{} 1)
\sphinxstyleemphasis{triclinic/general} values = no values
\end{sphinxalltt}

\end{itemize}


\subsection{Examples}
\label{\detokenize{lattice:examples}}
\begin{sphinxVerbatim}[commandchars=\\\{\}]
\PYG{k}{lattice}\PYG{+w}{ }\PYG{n}{fcc}\PYG{+w}{ }\PYG{l+m}{3.52}
\PYG{k}{lattice}\PYG{+w}{ }\PYG{n}{hex}\PYG{+w}{ }\PYG{l+m}{0.85}
\PYG{k}{lattice}\PYG{+w}{ }\PYG{n}{sq}\PYG{+w}{ }\PYG{l+m}{0.8}\PYG{+w}{ }\PYG{n}{origin}\PYG{+w}{ }\PYG{l+m}{0.0}\PYG{+w}{ }\PYG{l+m}{0.5}\PYG{+w}{ }\PYG{l+m}{0.0}\PYG{+w}{ }\PYG{n}{orient}\PYG{+w}{ }\PYG{n}{x}\PYG{+w}{ }\PYG{l+m}{1}\PYG{+w}{ }\PYG{l+m}{1}\PYG{+w}{ }\PYG{l+m}{0}\PYG{+w}{ }\PYG{n}{orient}\PYG{+w}{ }\PYG{n}{y}\PYG{+w}{ }\PYG{o}{\PYGZhy{}}\PYG{l+m}{1}\PYG{+w}{ }\PYG{l+m}{1}\PYG{+w}{ }\PYG{l+m}{0}
\PYG{k}{lattice}\PYG{+w}{ }\PYG{n}{custom}\PYG{+w}{ }\PYG{l+m}{3.52}\PYG{+w}{ }\PYG{n}{a1}\PYG{+w}{ }\PYG{l+m}{1.0}\PYG{+w}{ }\PYG{l+m}{0.0}\PYG{+w}{ }\PYG{l+m}{0.0}\PYG{+w}{ }\PYG{n}{a2}\PYG{+w}{ }\PYG{l+m}{0.5}\PYG{+w}{ }\PYG{l+m}{1.0}\PYG{+w}{ }\PYG{l+m}{0.0}\PYG{+w}{ }\PYG{n}{a3}\PYG{+w}{ }\PYG{l+m}{0.0}\PYG{+w}{ }\PYG{l+m}{0.0}\PYG{+w}{ }\PYG{l+m}{0.5}\PYG{+w}{ }\PYG{o}{\PYGZam{}}
\PYG{+w}{                    }\PYG{n}{basis}\PYG{+w}{ }\PYG{l+m}{0.0}\PYG{+w}{ }\PYG{l+m}{0.0}\PYG{+w}{ }\PYG{l+m}{0.0}\PYG{+w}{ }\PYG{n}{basis}\PYG{+w}{ }\PYG{l+m}{0.5}\PYG{+w}{ }\PYG{l+m}{0.5}\PYG{+w}{ }\PYG{l+m}{0.5}\PYG{+w}{ }\PYG{n}{triclinic}\PYG{o}{/}\PYG{n}{general}
\PYG{k}{lattice}\PYG{+w}{ }\PYG{n}{none}\PYG{+w}{ }\PYG{l+m}{2.0}
\end{sphinxVerbatim}


\subsection{Description}
\label{\detokenize{lattice:description}}
\sphinxAtStartPar
Define a lattice for use by other commands.  In LAMMPS, a lattice is
simply a set of points in space, determined by a unit cell with basis
atoms, that is replicated infinitely in all dimensions.  The arguments
of the lattice command can be used to define a wide variety of
crystallographic lattices.

\sphinxAtStartPar
A lattice is used by LAMMPS in two ways.  First, the
{\hyperref[\detokenize{create_atoms::doc}]{\sphinxcrossref{\DUrole{doc}{create\_atoms}}}} command creates atoms on the lattice
points inside the simulation box.  Note that the
{\hyperref[\detokenize{create_atoms::doc}]{\sphinxcrossref{\DUrole{doc}{create\_atoms}}}} command allows different atom types
to be assigned to different basis atoms of the lattice.  Second, the
lattice spacing in the x,y,z dimensions implied by the lattice, can be
used by other commands as distance units
(e.g. {\hyperref[\detokenize{create_box::doc}]{\sphinxcrossref{\DUrole{doc}{create\_box}}}}, {\hyperref[\detokenize{region::doc}]{\sphinxcrossref{\DUrole{doc}{region}}}} and
{\hyperref[\detokenize{velocity::doc}]{\sphinxcrossref{\DUrole{doc}{velocity}}}}), which are often convenient to use when the
underlying problem geometry is atoms on a lattice.

\sphinxAtStartPar
The lattice style must be consistent with the dimension of the
simulation \sphinxhyphen{} see the {\hyperref[\detokenize{dimension::doc}]{\sphinxcrossref{\DUrole{doc}{dimension}}}} command.  Styles \sphinxstyleemphasis{sc}
or \sphinxstyleemphasis{bcc} or \sphinxstyleemphasis{fcc} or \sphinxstyleemphasis{hcp} or \sphinxstyleemphasis{diamond} are for 3d problems.  Styles
\sphinxstyleemphasis{sq} or \sphinxstyleemphasis{sq2} or \sphinxstyleemphasis{hex} are for 2d problems.  Style \sphinxstyleemphasis{custom} can be
used for either 2d or 3d problems.

\sphinxAtStartPar
A lattice consists of a unit cell, a set of basis atoms within that
cell, and a set of transformation parameters (scale, origin, orient)
that map the unit cell into the simulation box.  The vectors a1,a2,a3
are the edge vectors of the unit cell.  This is the nomenclature for
“primitive” vectors in solid\sphinxhyphen{}state crystallography, but in LAMMPS the
unit cell they determine does not have to be a “primitive cell” of
minimum volume.

\sphinxAtStartPar
Note that the lattice command can be used multiple times in an input
script.  Each time it is invoked, the lattice attributes are
re\sphinxhyphen{}defined and are used for all subsequent commands (that use lattice
attributes).  For example, a sequence of lattice,
{\hyperref[\detokenize{region::doc}]{\sphinxcrossref{\DUrole{doc}{region}}}}, and {\hyperref[\detokenize{create_atoms::doc}]{\sphinxcrossref{\DUrole{doc}{create\_atoms}}}} commands
can be repeated multiple times to build a poly\sphinxhyphen{}crystalline model with
different geometric regions populated with atoms in different lattice
orientations.


\bigskip\hrule\bigskip


\sphinxAtStartPar
A lattice of style \sphinxstyleemphasis{none} does not define a unit cell and basis set,
so it cannot be used with the {\hyperref[\detokenize{create_atoms::doc}]{\sphinxcrossref{\DUrole{doc}{create\_atoms}}}}
command.  However it does define a lattice spacing via the specified
scale parameter.  As explained above the lattice spacings in x,y,z can
be used by other commands as distance units.  No additional
keyword/value pairs can be specified for the \sphinxstyleemphasis{none} style.  By
default, a “lattice none 1.0” is defined, which means the lattice
spacing is the same as one distance unit, as defined by the
{\hyperref[\detokenize{units::doc}]{\sphinxcrossref{\DUrole{doc}{units}}}} command.

\sphinxAtStartPar
Lattices of style \sphinxstyleemphasis{sc}, \sphinxstyleemphasis{fcc}, \sphinxstyleemphasis{bcc}, and \sphinxstyleemphasis{diamond} are 3d lattices
that define a cubic unit cell with edge length = 1.0.  This means a1 =
1 0 0, a2 = 0 1 0, and a3 = 0 0 1.  Style \sphinxstyleemphasis{hcp} has a1 = 1 0 0, a2 = 0
sqrt(3) 0, and a3 = 0 0 sqrt(8/3).  The placement of the basis atoms
within the unit cell are described in any solid\sphinxhyphen{}state physics text.  A
\sphinxstyleemphasis{sc} lattice has 1 basis atom at the lower\sphinxhyphen{}left\sphinxhyphen{}bottom corner of the
cube.  A \sphinxstyleemphasis{bcc} lattice has 2 basis atoms, one at the corner and one at
the center of the cube.  A \sphinxstyleemphasis{fcc} lattice has 4 basis atoms, one at the
corner and 3 at the cube face centers.  A \sphinxstyleemphasis{hcp} lattice has 4 basis
atoms, two in the z = 0 plane and 2 in the z = 0.5 plane.  A \sphinxstyleemphasis{diamond}
lattice has 8 basis atoms.

\sphinxAtStartPar
Lattices of style \sphinxstyleemphasis{sq} and \sphinxstyleemphasis{sq2} are 2d lattices that define a square
unit cell with edge length = 1.0.  This means a1 = 1 0 0 and a2 = 0 1
0.  A \sphinxstyleemphasis{sq} lattice has 1 basis atom at the lower\sphinxhyphen{}left corner of the
square.  A \sphinxstyleemphasis{sq2} lattice has 2 basis atoms, one at the corner and one
at the center of the square.  A \sphinxstyleemphasis{hex} style is also a 2d lattice, but
the unit cell is rectangular, with a1 = 1 0 0 and a2 = 0 sqrt(3) 0.
It has 2 basis atoms, one at the corner and one at the center of the
rectangle.

\sphinxAtStartPar
A lattice of style \sphinxstyleemphasis{custom} allows you to specify a1, a2, a3, and a
list of basis atoms to put in the unit cell.  By default, a1 and a2
and a3 are 3 orthogonal unit vectors (edges of a unit cube).  But you
can specify them to be of any length and non\sphinxhyphen{}orthogonal to each other,
so that they describe a tilted parallelepiped.  Via the \sphinxstyleemphasis{basis}
keyword you add atoms, one at a time, to the unit cell.  Its arguments
are fractional coordinates (0.0 \textless{}= x,y,z \textless{} 1.0).  For 2d simulations,
the fractional z coordinate for any basis atom must be 0.0.

\sphinxAtStartPar
The position vector x of a basis atom within the unit cell is a linear
combination of the unit cell’s 3 edge vectors, i.e. x = bx a1 + by
a2 + bz a3, where bx,by,bz are the 3 values specified for the \sphinxstyleemphasis{basis}
keyword.


\bigskip\hrule\bigskip


\sphinxAtStartPar
This subsection discusses the arguments that determine how the
idealized unit cell is transformed into a lattice of points within the
simulation box.

\sphinxAtStartPar
The \sphinxstyleemphasis{scale} argument determines how the size of the unit cell will be
scaled when mapping it into the simulation box.  I.e. it determines a
multiplicative factor to apply to the unit cell, to convert it to a
lattice of the desired size and distance units in the simulation box.
The meaning of the \sphinxstyleemphasis{scale} argument depends on the {\hyperref[\detokenize{units::doc}]{\sphinxcrossref{\DUrole{doc}{units}}}}
being used in your simulation.

\sphinxAtStartPar
For all unit styles except \sphinxstyleemphasis{lj}, the scale argument is specified in
the distance units defined by the unit style.  For example, in \sphinxstyleemphasis{real}
or \sphinxstyleemphasis{metal} units, if the unit cell is a unit cube with edge length
1.0, specifying scale = 3.52 would create a cubic lattice with a
spacing of 3.52 Angstroms.  In \sphinxstyleemphasis{cgs} units, the spacing would be 3.52
cm.

\sphinxAtStartPar
For unit style \sphinxstyleemphasis{lj}, the scale argument is the Lennard\sphinxhyphen{}Jones reduced
density, typically written as rho*.  LAMMPS converts this value into
the multiplicative factor via the formula “factor\textasciicircum{}dim = rho/rho*”,
where rho = N/V with V = the volume of the lattice unit cell and N =
the number of basis atoms in the unit cell (described below), and dim
= 2 or 3 for the dimensionality of the simulation.  Effectively, this
means that if LJ particles of size sigma = 1.0 are used in the
simulation, the lattice of particles will be at the desired reduced
density.

\sphinxAtStartPar
The \sphinxstyleemphasis{origin} option specifies how the unit cell will be shifted or
translated when mapping it into the simulation box.  The x,y,z values
are fractional values (0.0 \textless{}= x,y,z \textless{} 1.0) meaning shift the lattice
by a fraction of the lattice spacing in each dimension.  The meaning
of “lattice spacing” is discussed below.  For 2d simulations, the
\sphinxstyleemphasis{origin} z value must be 0.0.

\sphinxAtStartPar
The \sphinxstyleemphasis{orient} option specifies how the unit cell will be rotated when
mapping it into the simulation box.  The \sphinxstyleemphasis{dim} argument is one of the
3 coordinate axes in the simulation box.  The other 3 arguments are
the crystallographic direction in the lattice that you want to orient
along that axis, specified as integers.  E.g. “orient x 2 1 0” means
the x\sphinxhyphen{}axis in the simulation box will be the {[}210{]} lattice direction,
and similarly for y and z.  The 3 lattice directions you specify do
not have to be unit vectors, but they must be mutually orthogonal and
obey the right\sphinxhyphen{}hand rule, i.e. (X cross Y) points in the Z direction.
For 2d simulations, the \sphinxstyleemphasis{orient} x and y vectors must define 0 for
their 3rd component.  Similarly the \sphinxstyleemphasis{orient} z vector must define 0
for its 1st and 2nd components.

\begin{sphinxadmonition}{note}{Note:}
\sphinxAtStartPar
The preceding paragraph describing lattice directions is only
valid for orthogonal cubic unit cells (or square in 2d).  If you are
using a \sphinxstyleemphasis{hcp} or \sphinxstyleemphasis{hex} lattice or the more general lattice style
\sphinxstyleemphasis{custom} with non\sphinxhyphen{}orthogonal a1,a2,a3 vectors, then you should think
of the 3 \sphinxstyleemphasis{orient} vectors as creating a 3x3 rotation matrix which is
applied to a1,a2,a3 to rotate the original unit cell to a new
orientation in the simulation box.
\end{sphinxadmonition}


\bigskip\hrule\bigskip


\sphinxAtStartPar
The \sphinxstyleemphasis{triclinic/general} option specifies that the defined lattice is
for use with a general triclinic simulation box, as opposed to an
orthogonal or restricted triclinic box.  The {\hyperref[\detokenize{Howto_triclinic::doc}]{\sphinxcrossref{\DUrole{doc}{Howto triclinic}}}} doc page explains all 3 kinds of simulation boxes
LAMMPS supports.

\sphinxAtStartPar
If this option is specified, a \sphinxstyleemphasis{custom} lattice style must be used.
The \sphinxstyleemphasis{a1}, \sphinxstyleemphasis{a2}, \sphinxstyleemphasis{a3} vectors should define the edge vectors of a
single unit cell of the lattice with one or more basis atoms.  They
edge vectors can be arbitrary so long as they are non\sphinxhyphen{}zero, distinct,
and not co\sphinxhyphen{}planar.  In addition, they must define a right\sphinxhyphen{}handed
system, such that (\sphinxstyleemphasis{a1} cross \sphinxstyleemphasis{a2}) points in the direction of \sphinxstyleemphasis{a3}.
Note that a left\sphinxhyphen{}handed system can be converted to a right\sphinxhyphen{}handed
system by simply swapping the order of any pair of the \sphinxstyleemphasis{a1}, \sphinxstyleemphasis{a2},
\sphinxstyleemphasis{a3} vectors.  For 2d simulations, the \sphinxstyleemphasis{a3} vector must be specified
as (0.0,0.0,1.0), which is its default value.

\sphinxAtStartPar
If this option is used, the \sphinxstyleemphasis{origin} and \sphinxstyleemphasis{orient} settings must have
their default values.  Namely (0.0,0.0,0.0) for the \sphinxstyleemphasis{origin} and
(100), (010), (001) for the \sphinxstyleemphasis{orient} vectors.

\sphinxAtStartPar
The {\hyperref[\detokenize{create_box::doc}]{\sphinxcrossref{\DUrole{doc}{create\_box}}}} command can be used to create a
general triclinic box that replicates the \sphinxstyleemphasis{a1}, \sphinxstyleemphasis{a2}, \sphinxstyleemphasis{a3} unit cell
vectors in each direction to create the 3 arbitrary edge vectors of
the overall simulation box.  It requires a lattice with the
\sphinxstyleemphasis{triclinic/general} option.

\sphinxAtStartPar
Likewise, the {\hyperref[\detokenize{create_atoms::doc}]{\sphinxcrossref{\DUrole{doc}{create\_atoms}}}} command can be used
to add atoms (or molecules) to a general triclinic box which lie on
the lattice points defined by \sphinxstyleemphasis{a1}, \sphinxstyleemphasis{a2}, \sphinxstyleemphasis{a3} and the unit cell basis
atoms.  To do this, it also requires a lattice with the
\sphinxstyleemphasis{triclinic/general} option.

\begin{sphinxadmonition}{note}{Note:}
\sphinxAtStartPar
LAMMPS allows specification of general triclinic lattices and
simulation boxes as a convenience for users who may be converting
data from solid\sphinxhyphen{}state crystallographic representations or from DFT
codes for input to LAMMPS.  However, as explained on the
{\hyperref[\detokenize{Howto_triclinic::doc}]{\sphinxcrossref{\DUrole{doc}{Howto\_triclinic}}}} doc page, internally,
LAMMPS only uses restricted triclinic simulation boxes.  This means
the box and per\sphinxhyphen{}atom information (e.g. coordinates, velocities)
defined by the {\hyperref[\detokenize{create_box::doc}]{\sphinxcrossref{\DUrole{doc}{create\_box}}}} and
{\hyperref[\detokenize{create_atoms::doc}]{\sphinxcrossref{\DUrole{doc}{create\_atoms}}}} commands are converted from
general to restricted triclinic form when the two commands are
invoked.  It also means that any other commands which use lattice
spacings from this command (e.g. the region command), will be
operating on a restricted triclinic simulation box, even if the
\sphinxstyleemphasis{triclinic/general} option was used to define the lattice.  See the
next section for details.
\end{sphinxadmonition}


\bigskip\hrule\bigskip


\sphinxAtStartPar
Several LAMMPS commands have the option to use distance units that are
inferred from “lattice spacings” in the x,y,z box directions.
E.g. the {\hyperref[\detokenize{region::doc}]{\sphinxcrossref{\DUrole{doc}{region}}}} command can create a block of size
10x20x20, where 10 means 10 lattice spacings in the x direction.

\begin{sphinxadmonition}{note}{Note:}
\sphinxAtStartPar
Though they are called lattice spacings, all the commands that
have a “units lattice” option, simply use the 3 values as scale
factors on the distance units defined by the {\hyperref[\detokenize{units::doc}]{\sphinxcrossref{\DUrole{doc}{units}}}}
command.  Thus if you do not like the lattice spacings computed by
LAMMPS (e.g. for a non\sphinxhyphen{}orthogonal or rotated unit cell), you can
define the 3 values to be whatever you wish, via the \sphinxstyleemphasis{spacing} option.
\end{sphinxadmonition}

\sphinxAtStartPar
If the \sphinxstyleemphasis{spacing} option is not specified, the lattice spacings are
computed by LAMMPS in the following way.  A unit cell of the lattice
is mapped into the simulation box (scaled and rotated), so that it now
has (perhaps) a modified size and orientation.  The lattice spacing in
X is defined as the difference between the min/max extent of the x
coordinates of the 8 corner points of the modified unit cell (4 in
2d).  Similarly, the Y and Z lattice spacings are defined as the
difference in the min/max of the y and z coordinates.

\begin{sphinxadmonition}{note}{Note:}
\sphinxAtStartPar
If the \sphinxstyleemphasis{triclinic/general} option is specified, the unit cell
defined by \sphinxstyleemphasis{a1}, \sphinxstyleemphasis{a2}, \sphinxstyleemphasis{a3} edge vectors is first converted to a
restricted triclinic orientation, which is a rotation operation.
The min/max extent of the 8 corner points is then determined, as
described in the preceding paragraph, to set the lattice
spacings. As explained for the \sphinxstyleemphasis{triclinic/general} option above,
this is because any use of the lattice spacings by other commands
will be for a restricted triclinic simulation box, not a general
triclinic box.
\end{sphinxadmonition}

\sphinxAtStartPar
Note that if the unit cell is orthogonal with axis\sphinxhyphen{}aligned edges (no
rotation via the \sphinxstyleemphasis{orient} keyword), then the lattice spacings in each
dimension are simply the scale factor (described above) multiplied by
the length of a1,a2,a3.  Thus a \sphinxstyleemphasis{hex} style lattice with a scale
factor of 3.0 Angstroms, would have a lattice spacing of 3.0 in x and
3*sqrt(3.0) in y.

\begin{sphinxadmonition}{note}{Note:}
\sphinxAtStartPar
For non\sphinxhyphen{}orthogonal unit cells and/or when a rotation is applied
via the \sphinxstyleemphasis{orient} keyword, then the lattice spacings computed by LAMMPS
are typically less intuitive.  In particular, in these cases, there is
no guarantee that a particular lattice spacing is an integer multiple
of the periodicity of the lattice in that direction.  Thus, if you
create an orthogonal periodic simulation box whose size in a dimension
is a multiple of the lattice spacing, and then fill it with atoms via
the {\hyperref[\detokenize{create_atoms::doc}]{\sphinxcrossref{\DUrole{doc}{create\_atoms}}}} command, you will NOT necessarily
create a periodic system.  I.e. atoms may overlap incorrectly at the
faces of the simulation box.
\end{sphinxadmonition}

\sphinxAtStartPar
The \sphinxstyleemphasis{spacing} option sets the 3 lattice spacings directly.  All must
be non\sphinxhyphen{}zero (use 1.0 for dz in a 2d simulation).  The specified values
are multiplied by the multiplicative factor described above that is
associated with the scale factor.  Thus a spacing of 1.0 means one
unit cell edge length independent of the scale factor.  As mentioned
above, this option can be useful if the spacings LAMMPS computes are
inconvenient to use in subsequent commands, which can be the case for
non\sphinxhyphen{}orthogonal or rotated lattices.

\sphinxAtStartPar
Note that whenever the lattice command is used, the values of the
lattice spacings LAMMPS calculates are printed out.  Thus their effect
in commands that use the spacings should be decipherable.


\bigskip\hrule\bigskip


\sphinxAtStartPar
Example commands for generating a Wurtzite crystal.
The lattice constants approximate those of CdSe.
The \(\sqrt{3}\times 1\) orthorhombic supercell is used
with the x, y, and z directions oriented
along \([\bar{1}\bar{2}30]\),
\([10\bar{1}0]\), and \([0001]\), respectively.

\begin{sphinxVerbatim}[commandchars=\\\{\}]
\PYG{k}{variable }\PYG{n+nv+nvIdentifier}{a}\PYG{+w}{ }\PYG{n}{equal}\PYG{+w}{  }\PYG{l+m}{4.34}
\PYG{k}{variable }\PYG{n+nv+nvIdentifier}{b}\PYG{+w}{ }\PYG{n}{equal}\PYG{+w}{  }\PYG{n+nv}{\PYGZdl{}a}\PYG{o}{*}\PYG{n}{sqrt}\PYG{n+nv}{(}\PYG{n+nv}{3.0}\PYG{n+nv}{)}
\PYG{k}{variable }\PYG{n+nv+nvIdentifier}{c}\PYG{+w}{ }\PYG{n}{equal}\PYG{+w}{  }\PYG{n+nv}{\PYGZdl{}a}\PYG{o}{*}\PYG{n}{sqrt}\PYG{n+nv}{(}\PYG{n+nv}{8.0/3.0}\PYG{n+nv}{)}

\PYG{k}{variable }\PYG{n+nv+nvIdentifier}{third}\PYG{+w}{ }\PYG{n}{equal}\PYG{+w}{ }\PYG{l+m}{1.0}\PYG{o}{/}\PYG{l+m}{3.0}
\PYG{k}{variable }\PYG{n+nv+nvIdentifier}{five6}\PYG{+w}{ }\PYG{n}{equal}\PYG{+w}{ }\PYG{l+m}{5.0}\PYG{o}{/}\PYG{l+m}{6.0}

\PYG{k}{lattice}\PYG{+w}{ }\PYG{n}{custom}\PYG{+w}{    }\PYG{l+m}{1.0}\PYG{+w}{     }\PYG{o}{\PYGZam{}}
\PYG{+w}{        }\PYG{n}{a1}\PYG{+w}{      }\PYG{n+nv}{\PYGZdl{}b}\PYG{+w}{       }\PYG{l+m}{0.0}\PYG{+w}{     }\PYG{l+m}{0.0}\PYG{+w}{     }\PYG{o}{\PYGZam{}}
\PYG{+w}{        }\PYG{n}{a2}\PYG{+w}{      }\PYG{l+m}{0.0}\PYG{+w}{      }\PYG{n+nv}{\PYGZdl{}a}\PYG{+w}{      }\PYG{l+m}{0.0}\PYG{+w}{     }\PYG{o}{\PYGZam{}}
\PYG{+w}{        }\PYG{n}{a3}\PYG{+w}{      }\PYG{l+m}{0.0}\PYG{+w}{      }\PYG{l+m}{0.0}\PYG{+w}{     }\PYG{n+nv}{\PYGZdl{}c}\PYG{+w}{      }\PYG{o}{\PYGZam{}}
\PYG{+w}{        }\PYG{n}{basis}\PYG{+w}{   }\PYG{l+m}{0.0}\PYG{+w}{      }\PYG{l+m}{0.0}\PYG{+w}{     }\PYG{l+m}{0.0}\PYG{+w}{     }\PYG{o}{\PYGZam{}}
\PYG{+w}{        }\PYG{n}{basis}\PYG{+w}{   }\PYG{l+m}{0.5}\PYG{+w}{      }\PYG{l+m}{0.5}\PYG{+w}{     }\PYG{l+m}{0.0}\PYG{+w}{     }\PYG{o}{\PYGZam{}}
\PYG{+w}{        }\PYG{n}{basis}\PYG{+w}{   }\PYG{n+nv}{\PYGZdl{}\PYGZob{}}\PYG{n+nv}{third}\PYG{n+nv}{\PYGZcb{}}\PYG{+w}{ }\PYG{l+m}{0.0}\PYG{+w}{     }\PYG{l+m}{0.5}\PYG{+w}{     }\PYG{o}{\PYGZam{}}
\PYG{+w}{        }\PYG{n}{basis}\PYG{+w}{   }\PYG{n+nv}{\PYGZdl{}\PYGZob{}}\PYG{n+nv}{five6}\PYG{n+nv}{\PYGZcb{}}\PYG{+w}{ }\PYG{l+m}{0.5}\PYG{+w}{     }\PYG{l+m}{0.5}\PYG{+w}{     }\PYG{o}{\PYGZam{}}
\PYG{+w}{        }\PYG{n}{basis}\PYG{+w}{   }\PYG{l+m}{0.0}\PYG{+w}{      }\PYG{l+m}{0.0}\PYG{+w}{     }\PYG{l+m}{0.625}\PYG{+w}{   }\PYG{o}{\PYGZam{}}
\PYG{+w}{        }\PYG{n}{basis}\PYG{+w}{   }\PYG{l+m}{0.5}\PYG{+w}{      }\PYG{l+m}{0.5}\PYG{+w}{     }\PYG{l+m}{0.625}\PYG{+w}{   }\PYG{o}{\PYGZam{}}
\PYG{+w}{        }\PYG{n}{basis}\PYG{+w}{   }\PYG{n+nv}{\PYGZdl{}\PYGZob{}}\PYG{n+nv}{third}\PYG{n+nv}{\PYGZcb{}}\PYG{+w}{ }\PYG{l+m}{0.0}\PYG{+w}{     }\PYG{l+m}{0.125}\PYG{+w}{   }\PYG{o}{\PYGZam{}}
\PYG{+w}{        }\PYG{n}{basis}\PYG{+w}{   }\PYG{n+nv}{\PYGZdl{}\PYGZob{}}\PYG{n+nv}{five6}\PYG{n+nv}{\PYGZcb{}}\PYG{+w}{ }\PYG{l+m}{0.5}\PYG{+w}{     }\PYG{l+m}{0.125}

\PYG{k}{region }\PYG{n+nv+nvIdentifier}{myreg}\PYG{+w}{ }\PYG{n}{block}\PYG{+w}{ }\PYG{l+m}{0}\PYG{+w}{ }\PYG{l+m}{1}\PYG{+w}{ }\PYG{l+m}{0}\PYG{+w}{ }\PYG{l+m}{1}\PYG{+w}{ }\PYG{l+m}{0}\PYG{+w}{ }\PYG{l+m}{1}
\PYG{k}{create\PYGZus{}box      }\PYG{n+nv+nvIdentifier}{2}\PYG{+w}{ }\PYG{n+nv+nvIdentifier}{myreg}
\PYG{k}{create\PYGZus{}atoms}\PYG{+w}{    }\PYG{l+m}{1}\PYG{+w}{ }\PYG{n}{box}\PYG{+w}{      }\PYG{o}{\PYGZam{}}
\PYG{+w}{        }\PYG{n}{basis}\PYG{+w}{   }\PYG{l+m}{5}\PYG{+w}{       }\PYG{l+m}{2}\PYG{+w}{       }\PYG{o}{\PYGZam{}}
\PYG{+w}{        }\PYG{n}{basis}\PYG{+w}{   }\PYG{l+m}{6}\PYG{+w}{       }\PYG{l+m}{2}\PYG{+w}{       }\PYG{o}{\PYGZam{}}
\PYG{+w}{        }\PYG{n}{basis}\PYG{+w}{   }\PYG{l+m}{7}\PYG{+w}{       }\PYG{l+m}{2}\PYG{+w}{       }\PYG{o}{\PYGZam{}}
\PYG{+w}{        }\PYG{n}{basis}\PYG{+w}{   }\PYG{l+m}{8}\PYG{+w}{       }\PYG{l+m}{2}
\end{sphinxVerbatim}


\bigskip\hrule\bigskip



\subsection{Restrictions}
\label{\detokenize{lattice:restrictions}}
\sphinxAtStartPar
The \sphinxstyleemphasis{a1,a2,a3,basis} keywords can only be used with style \sphinxstyleemphasis{custom}.


\subsection{Related commands}
\label{\detokenize{lattice:related-commands}}
\sphinxAtStartPar
{\hyperref[\detokenize{dimension::doc}]{\sphinxcrossref{\DUrole{doc}{dimension}}}}, {\hyperref[\detokenize{create_atoms::doc}]{\sphinxcrossref{\DUrole{doc}{create\_atoms}}}},
{\hyperref[\detokenize{region::doc}]{\sphinxcrossref{\DUrole{doc}{region}}}}


\subsection{Default}
\label{\detokenize{lattice:default}}
\begin{sphinxVerbatim}[commandchars=\\\{\}]
\PYG{k}{lattice}\PYG{+w}{ }\PYG{n}{none}\PYG{+w}{ }\PYG{l+m}{1.0}
\end{sphinxVerbatim}

\sphinxAtStartPar
For other lattice styles, the option defaults are origin = 0.0 0.0
0.0, orient = x 1 0 0, orient = y 0 1 0, orient = z 0 0 1, a1 = 1 0 0,
a2 = 0 1 0, and a3 = 0 0 1.

\sphinxstepscope

\index{log@\spxentry{log}}\ignorespaces 

\section{log command}
\label{\detokenize{log:log-command}}\label{\detokenize{log:index-0}}\label{\detokenize{log::doc}}

\subsection{Syntax}
\label{\detokenize{log:syntax}}
\begin{sphinxVerbatim}[commandchars=\\\{\}]
\PYG{k}{log}\PYG{+w}{ }\PYG{n}{file}\PYG{+w}{ }\PYG{n}{keyword}
\end{sphinxVerbatim}
\begin{itemize}
\item {} 
\sphinxAtStartPar
file = name of new logfile

\item {} 
\sphinxAtStartPar
keyword = \sphinxstyleemphasis{append} if output should be appended to logfile (optional)

\end{itemize}


\subsection{Examples}
\label{\detokenize{log:examples}}
\begin{sphinxVerbatim}[commandchars=\\\{\}]
\PYG{k}{log}\PYG{+w}{ }\PYG{n}{log.equil}
\PYG{k}{log}\PYG{+w}{ }\PYG{n}{log.equil}\PYG{+w}{ }\PYG{n}{append}
\end{sphinxVerbatim}


\subsection{Description}
\label{\detokenize{log:description}}
\sphinxAtStartPar
This command closes the current LAMMPS log file, opens a new file with
the specified name, and begins logging information to it.  If the
specified file name is \sphinxstyleemphasis{none}, then no new log file is opened.  If the
optional keyword \sphinxstyleemphasis{append} is specified, then output will be appended
to an existing log file, instead of overwriting it.

\sphinxAtStartPar
If multiple processor partitions are being used, the file name should
be a variable, so that different processors do not attempt to write to
the same log file.

\sphinxAtStartPar
The file “log.lammps” is the default log file for a LAMMPS run.  The
name of the initial log file can also be set by the {\hyperref[\detokenize{Run_options::doc}]{\sphinxcrossref{\DUrole{doc}{\sphinxhyphen{}log command\sphinxhyphen{}line switch}}}}.


\subsection{Restrictions}
\label{\detokenize{log:restrictions}}
\sphinxAtStartPar
none


\subsection{Related commands}
\label{\detokenize{log:related-commands}}
\sphinxAtStartPar
none


\subsection{Default}
\label{\detokenize{log:default}}
\sphinxAtStartPar
The default LAMMPS log file is named log.lammps

\sphinxstepscope

\index{mass@\spxentry{mass}}\ignorespaces 

\section{mass command}
\label{\detokenize{mass:mass-command}}\label{\detokenize{mass:index-0}}\label{\detokenize{mass::doc}}

\subsection{Syntax}
\label{\detokenize{mass:syntax}}
\begin{sphinxVerbatim}[commandchars=\\\{\}]
\PYG{k}{mass}\PYG{+w}{ }\PYG{n}{I}\PYG{+w}{ }\PYG{n}{value}
\end{sphinxVerbatim}
\begin{itemize}
\item {} 
\sphinxAtStartPar
I = atom type (see asterisk form below), or type label

\item {} 
\sphinxAtStartPar
value = mass

\end{itemize}


\subsection{Examples}
\label{\detokenize{mass:examples}}
\begin{sphinxVerbatim}[commandchars=\\\{\}]
\PYG{k}{mass}\PYG{+w}{ }\PYG{l+m}{1}\PYG{+w}{ }\PYG{l+m}{1.0}
\PYG{k}{mass}\PYG{+w}{ }\PYG{o}{*}\PYG{+w}{ }\PYG{l+m}{62.5}
\PYG{k}{mass}\PYG{+w}{ }\PYG{l+m}{2}\PYG{o}{*}\PYG{+w}{ }\PYG{l+m}{62.5}

\PYG{k}{labelmap}\PYG{+w}{ }\PYG{n}{atom}\PYG{+w}{ }\PYG{l+m}{1}\PYG{+w}{ }\PYG{n}{C}
\PYG{k}{mass}\PYG{+w}{ }\PYG{n}{C}\PYG{+w}{ }\PYG{l+m}{12.01}
\end{sphinxVerbatim}


\subsection{Description}
\label{\detokenize{mass:description}}
\sphinxAtStartPar
Set the mass for all atoms of one or more atom types.  Per\sphinxhyphen{}type mass
values can also be set in the {\hyperref[\detokenize{read_data::doc}]{\sphinxcrossref{\DUrole{doc}{read\_data}}}} data file
using the “Masses” keyword.  See the {\hyperref[\detokenize{units::doc}]{\sphinxcrossref{\DUrole{doc}{units}}}} command for
what mass units to use.

\sphinxAtStartPar
The I index can be specified in one of several ways.  An explicit
numeric value can be used, as in the first example above.  Or I can be
a type label, which is an alphanumeric string defined by the
{\hyperref[\detokenize{labelmap::doc}]{\sphinxcrossref{\DUrole{doc}{labelmap}}}} command or in a section of a data file read
by the {\hyperref[\detokenize{read_data::doc}]{\sphinxcrossref{\DUrole{doc}{read\_data}}}} command, and which converts
internally to a numeric type. Or a wild\sphinxhyphen{}card asterisk can be used to
set the mass for multiple atom types.  This takes the form “*” or
“*n” or “n*” or “m*n”, where m and n are numbers.  If N = the
number of atom types, then an asterisk with no numeric values means
all types from 1 to N.  A leading asterisk means all types from 1 to n
(inclusive).  A trailing asterisk means all types from n to N
(inclusive).  A middle asterisk means all types from m to n
(inclusive).

\sphinxAtStartPar
A line in a {\hyperref[\detokenize{read_data::doc}]{\sphinxcrossref{\DUrole{doc}{data file}}}} that follows the “Masses”
keyword specifies mass using the same format as the arguments of the
mass command in an input script, except that no wild\sphinxhyphen{}card asterisk can
be used.  For example, under the “Masses” section of a data file, the
line that corresponds to the first example above would be listed as

\begin{sphinxVerbatim}[commandchars=\\\{\}]
1 1.0
\end{sphinxVerbatim}

\sphinxAtStartPar
Note that the mass command can only be used if the {\hyperref[\detokenize{atom_style::doc}]{\sphinxcrossref{\DUrole{doc}{atom style}}}} requires per\sphinxhyphen{}type atom mass to be set.
Currently, all but the \sphinxstyleemphasis{sphere} and \sphinxstyleemphasis{ellipsoid} and \sphinxstyleemphasis{peri} styles do.
They require mass to be set for individual particles, not types.
Per\sphinxhyphen{}atom masses are defined in the data file read by the
{\hyperref[\detokenize{read_data::doc}]{\sphinxcrossref{\DUrole{doc}{read\_data}}}} command, or set to default values by the
{\hyperref[\detokenize{create_atoms::doc}]{\sphinxcrossref{\DUrole{doc}{create\_atoms}}}} command.  Per\sphinxhyphen{}atom masses can also be
set to new values by the {\hyperref[\detokenize{set::doc}]{\sphinxcrossref{\DUrole{doc}{set mass}}}} or {\hyperref[\detokenize{set::doc}]{\sphinxcrossref{\DUrole{doc}{set density}}}}
commands.

\sphinxAtStartPar
Also note that {\hyperref[\detokenize{pair_eam::doc}]{\sphinxcrossref{\DUrole{doc}{pair\_style eam}}}} and {\hyperref[\detokenize{pair_bop::doc}]{\sphinxcrossref{\DUrole{doc}{pair\_style bop}}}} commands define the masses of atom types in their
respective potential files, in which case the mass command is normally
not used.

\sphinxAtStartPar
If you define a {\hyperref[\detokenize{atom_style::doc}]{\sphinxcrossref{\DUrole{doc}{hybrid atom style}}}} which includes one
(or more) sub\sphinxhyphen{}styles which require per\sphinxhyphen{}type mass and one (or more)
sub\sphinxhyphen{}styles which require per\sphinxhyphen{}atom mass, then you must define both.
However, in this case the per\sphinxhyphen{}type mass will be ignored; only the
per\sphinxhyphen{}atom mass will be used by LAMMPS.


\subsection{Restrictions}
\label{\detokenize{mass:restrictions}}
\sphinxAtStartPar
This command must come after the simulation box is defined by a
{\hyperref[\detokenize{read_data::doc}]{\sphinxcrossref{\DUrole{doc}{read\_data}}}}, {\hyperref[\detokenize{read_restart::doc}]{\sphinxcrossref{\DUrole{doc}{read\_restart}}}}, or
{\hyperref[\detokenize{create_box::doc}]{\sphinxcrossref{\DUrole{doc}{create\_box}}}} command.

\sphinxAtStartPar
All masses must be defined before a simulation is run.  They must also
all be defined before a {\hyperref[\detokenize{velocity::doc}]{\sphinxcrossref{\DUrole{doc}{velocity}}}} or {\hyperref[\detokenize{fix_shake::doc}]{\sphinxcrossref{\DUrole{doc}{fix shake}}}} command is used.

\sphinxAtStartPar
The mass assigned to any type or atom must be \textgreater{} 0.0.


\subsection{Related commands}
\label{\detokenize{mass:related-commands}}
\sphinxAtStartPar
none


\subsection{Default}
\label{\detokenize{mass:default}}
\sphinxAtStartPar
none

\sphinxstepscope

\index{mdi@\spxentry{mdi}}\ignorespaces 

\section{mdi command}
\label{\detokenize{mdi:mdi-command}}\label{\detokenize{mdi:index-0}}\label{\detokenize{mdi::doc}}

\subsection{Syntax}
\label{\detokenize{mdi:syntax}}
\begin{sphinxVerbatim}[commandchars=\\\{\}]
\PYG{k}{mdi}\PYG{+w}{ }\PYG{n}{option}\PYG{+w}{ }\PYG{n}{args}
\end{sphinxVerbatim}
\begin{itemize}
\item {} 
\sphinxAtStartPar
option = \sphinxstyleemphasis{engine} or \sphinxstyleemphasis{plugin} or \sphinxstyleemphasis{connect} or \sphinxstyleemphasis{exit}
\begin{sphinxalltt}
\sphinxstyleemphasis{engine} args = zero or more keyword/args pairs
  keywords = \sphinxstyleemphasis{elements}
    \sphinxstyleemphasis{elements} args = N\_1 N\_2 ... N\_ntypes
      N\_1,N\_2,...N\_ntypes = chemical symbol for each of ntypes LAMMPS atom types
\sphinxstyleemphasis{plugin} args = name keyword value keyword value ...
  name = name of plugin library (e.g., \sphinxstyleemphasis{lammps} means a liblammps.so library will be loaded)
  keyword/value pairs in any order, some are required, some are optional
  keywords = \sphinxstyleemphasis{mdi} or \sphinxstyleemphasis{infile} or \sphinxstyleemphasis{extra} or \sphinxstyleemphasis{command}
    \sphinxstyleemphasis{mdi} value = args passed to MDI for driver to operate with plugins (required)
    \sphinxstyleemphasis{infile} value = filename the engine will read at start\sphinxhyphen{}up (optional)
    \sphinxstyleemphasis{extra} value = aditional command\sphinxhyphen{}line args to pass to engine library when loaded (optional)
    \sphinxstyleemphasis{command} value = a LAMMPS input script command to execute (required)
\sphinxstyleemphasis{connect} args = none
\sphinxstyleemphasis{exit} args = none
\end{sphinxalltt}

\end{itemize}


\subsection{Examples}
\label{\detokenize{mdi:examples}}
\begin{sphinxVerbatim}[commandchars=\\\{\}]
\PYG{k}{mdi}\PYG{+w}{ }\PYG{n}{engine}
\PYG{k}{mdi}\PYG{+w}{ }\PYG{n}{engine}\PYG{+w}{ }\PYG{n}{elements}\PYG{+w}{ }\PYG{n}{Al}\PYG{+w}{ }\PYG{n}{Cu}
\PYG{k}{mdi}\PYG{+w}{ }\PYG{n}{plugin}\PYG{+w}{ }\PYG{n}{lammps}\PYG{+w}{ }\PYG{n}{mdi}\PYG{+w}{ }\PYG{l+s}{\PYGZdq{}}\PYG{l+s}{\PYGZhy{}role ENGINE \PYGZhy{}name lammps \PYGZhy{}method LINK}\PYG{l+s}{\PYGZdq{}}\PYG{+w}{ }\PYG{o}{\PYGZam{}}
\PYG{+w}{           }\PYG{n}{infile}\PYG{+w}{ }\PYG{n}{in.aimd.engine}\PYG{+w}{ }\PYG{n}{extra}\PYG{+w}{ }\PYG{l+s}{\PYGZdq{}}\PYG{l+s}{\PYGZhy{}log log.aimd.engine.plugin}\PYG{l+s}{\PYGZdq{}}\PYG{+w}{ }\PYG{o}{\PYGZam{}}
\PYG{+w}{           }\PYG{n}{command}\PYG{+w}{ }\PYG{l+s}{\PYGZdq{}}\PYG{l+s}{run 5}\PYG{l+s}{\PYGZdq{}}
\PYG{k}{mdi}\PYG{+w}{ }\PYG{n}{connect}
\PYG{k}{mdi}\PYG{+w}{ }\PYG{n}{exit}
\end{sphinxVerbatim}


\subsection{Description}
\label{\detokenize{mdi:description}}
\sphinxAtStartPar
This command implements operations within LAMMPS to use the \sphinxtitleref{MDI
Library \textless{}https://molssi\sphinxhyphen{}mdi.github.io/MDI\_Library/html/index.html\textgreater{}}
for coupling to other codes in a client/server protocol.

\sphinxAtStartPar
See the Howto MDI doc page for a discussion of all the different ways
2 or more codes can interact via MDI.

\sphinxAtStartPar
The examples/mdi directory has examples which use LAMMPS in 4
different modes: as a driver using an engine as either a stand\sphinxhyphen{}alone
code or as a plugin, and as an engine operating as either a
stand\sphinxhyphen{}alone code or as a plugin.  The README file in that directory
shows how to launch and couple codes for all the 4 usage modes, and so
they communicate via the MDI library using either MPI or sockets.

\sphinxAtStartPar
The scripts in that directory illustrate the use of all the options
for this command.

\sphinxAtStartPar
The \sphinxstyleemphasis{engine} option enables LAMMPS to act as an MDI engine (server),
responding to requests from an MDI driver (client) code.

\sphinxAtStartPar
The \sphinxstyleemphasis{plugin} option enables LAMMPS to act as an MDI driver (client),
and load the MDI engine (server) code as a library plugin.  In this
case the MDI engine is a library plugin.  An MDI engine can also be a
stand\sphinxhyphen{}alone code, launched separately from LAMMPS, in which case the
mdi plugin command is not used.

\sphinxAtStartPar
The \sphinxstyleemphasis{connect} and \sphinxstyleemphasis{exit} options are only used when LAMMPS is acting
as an MDI driver.  As explained below, these options are normally not
needed, except for a specific kind of use case.


\bigskip\hrule\bigskip


\sphinxAtStartPar
The \sphinxstyleemphasis{mdi engine} command is used to make LAMMPS operate as an MDI
engine.  It is typically used in an input script after LAMMPS has
setup the system it is going to model consistent with what the driver
code expects.  Depending on when the driver code tells the LAMMPS
engine to exit, other commands can be executed after this command, but
typically it is used at the end of a LAMMPS input script.

\sphinxAtStartPar
To act as an MDI engine operating as an MD code (or surrogate QM
code), this is the list of standard MDI commands issued by a driver
code which LAMMPS currently recognizes.  Using standard commands
defined by the MDI library means that a driver code can work
interchangeably with LAMMPS or other MD codes or with QM codes which
support the MDI standard.  See more details about these commands in
the \sphinxhref{https://molssi-mdi.github.io/MDI\_Library/html/mdi\_standard.html}{MDI library documentation}

\sphinxAtStartPar
These commands are valid at the @DEFAULT node defined by MDI.
Commands that start with “\textgreater{}” mean the driver is sending information to
LAMMPS.  Commands that start with “\textless{}” are requests by the driver for
LAMMPS to send it information.  Commands that start with an alphabetic
letter perform actions.  Commands that start with “@” are MDI “node”
commands, which are described further below.


\begin{savenotes}\sphinxattablestart
\sphinxthistablewithglobalstyle
\centering
\begin{tabular}[t]{\X{33}{100}\X{67}{100}}
\sphinxtoprule
\sphinxstyletheadfamily 
\sphinxAtStartPar
Command name
&\sphinxstyletheadfamily 
\sphinxAtStartPar
Action
\\
\sphinxmidrule
\sphinxtableatstartofbodyhook
\sphinxAtStartPar
\textgreater{}CELL or \textless{}CELL
&
\sphinxAtStartPar
Send/request 3 simulation box edge vectors (9 values)
\\
\sphinxhline
\sphinxAtStartPar
\textgreater{}CELL\_DISPL or \textless{}CELL\_DISPL
&
\sphinxAtStartPar
Send/request displacement of the simulation box from the origin (3 values)
\\
\sphinxhline
\sphinxAtStartPar
\textgreater{}CHARGES or \textless{}CHARGES
&
\sphinxAtStartPar
Send/request charge on each atom (N values)
\\
\sphinxhline
\sphinxAtStartPar
\textgreater{}COORDS or \textless{}COORDS
&
\sphinxAtStartPar
Send/request coordinates of each atom (3N values)
\\
\sphinxhline
\sphinxAtStartPar
\textgreater{}ELEMENTS
&
\sphinxAtStartPar
Send elements (atomic numbers) for each atom (N values)
\\
\sphinxhline
\sphinxAtStartPar
\textless{}ENERGY
&
\sphinxAtStartPar
Request total energy (potential + kinetic) of the system (1 value)
\\
\sphinxhline
\sphinxAtStartPar
\textgreater{}FORCES or \textless{}FORCES
&
\sphinxAtStartPar
Send/request forces on each atom (3N values)
\\
\sphinxhline
\sphinxAtStartPar
\textgreater{}+FORCES
&
\sphinxAtStartPar
Send forces to add to each atom (3N values)
\\
\sphinxhline
\sphinxAtStartPar
\textless{}LABELS
&
\sphinxAtStartPar
Request string label of each atom (N values)
\\
\sphinxhline
\sphinxAtStartPar
\textless{}MASSES
&
\sphinxAtStartPar
Request mass of each atom (N values)
\\
\sphinxhline
\sphinxAtStartPar
MD
&
\sphinxAtStartPar
Perform an MD simulation for N timesteps (most recent \textgreater{}NSTEPS value)
\\
\sphinxhline
\sphinxAtStartPar
OPTG
&
\sphinxAtStartPar
Perform an energy minimization to convergence (most recent \textgreater{}TOLERANCE values)
\\
\sphinxhline
\sphinxAtStartPar
\textgreater{}NATOMS or \textless{}NATOMS
&
\sphinxAtStartPar
Sends/request number of atoms in the system (1 value)
\\
\sphinxhline
\sphinxAtStartPar
\textgreater{}NSTEPS
&
\sphinxAtStartPar
Send number of timesteps for next MD dynamics run via MD command
\\
\sphinxhline
\sphinxAtStartPar
\textless{}PE
&
\sphinxAtStartPar
Request potential energy of the system (1 value)
\\
\sphinxhline
\sphinxAtStartPar
\textless{}STRESS
&
\sphinxAtStartPar
Request symmetric stress tensor (virial) of the system (9 values)
\\
\sphinxhline
\sphinxAtStartPar
\textgreater{}TOLERANCE
&
\sphinxAtStartPar
Send 4 tolerance parameters for next MD minimization via OPTG command
\\
\sphinxhline
\sphinxAtStartPar
\textgreater{}TYPES or \textless{}TYPES
&
\sphinxAtStartPar
Send/request the LAMMPS atom type for each atom (N values)
\\
\sphinxhline
\sphinxAtStartPar
\textgreater{}VELOCITIES or \textless{}VELOCITIES
&
\sphinxAtStartPar
Send/request the velocity of each atom (3N values)
\\
\sphinxhline
\sphinxAtStartPar
@INIT\_MD or @INIT\_OPTG
&
\sphinxAtStartPar
Driver tells LAMMPS to start single\sphinxhyphen{}step dynamics or minimization (see below)
\\
\sphinxhline
\sphinxAtStartPar
EXIT
&
\sphinxAtStartPar
Driver tells LAMMPS to exit engine mode
\\
\sphinxbottomrule
\end{tabular}
\sphinxtableafterendhook\par
\sphinxattableend\end{savenotes}

\begin{sphinxadmonition}{note}{Note:}
\sphinxAtStartPar
The \textless{}ENERGY, \textless{}FORCES, \textless{}PE, and \textless{}STRESS commands trigger LAMMPS to
compute atomic interactions for the current configuration of atoms
and size/shape of the simulation box.  I.e. LAMMPS invokes its
pair, bond, angle, …, kspace styles.  If the driver is updating
the atom coordinates and/or box incrementally (as in an MD
simulation which the driver is managing), then the LAMMPS engine
will do the same, and only occasionally trigger neighbor list
builds.  If the change in atom positions is large (since the
previous \textgreater{}COORDS command), then LAMMPS will do a more expensive
operation to migrate atoms to new processors as needed and
re\sphinxhyphen{}neighbor.  If the \textgreater{}NATOMS or \textgreater{}TYPES or \textgreater{}ELEMENTS commands have
been sent (since the previous \textgreater{}COORDS command), then LAMMPS assumes
the system is new and re\sphinxhyphen{}initializes an entirely new simulation.
\end{sphinxadmonition}

\begin{sphinxadmonition}{note}{Note:}
\sphinxAtStartPar
The \textgreater{}TYPES or \textgreater{}ELEMENTS commands are how the MDI driver tells the
LAMMPS engine which LAMMPS atom type to assign to each atom.  If
both the MDI driver and the LAMMPS engine are initialized so that
atom type values are consistent in both codes, then the \textgreater{}TYPES
command can be used.  If not, the optional \sphinxstyleemphasis{elements} keyword can
be used to specify what element each LAMMPS atom type corresponds
to.  This is specified by the chemical symbol of the element,
e.g. C or Al or Si.  A symbol must be specified for each of the
ntypes LAMMPS atom types.  Each LAMMPS type must map to a unique
element; two or more types cannot map to the same element.  Ntypes
is typically specified via the {\hyperref[\detokenize{create_box::doc}]{\sphinxcrossref{\DUrole{doc}{create\_box}}}}
command or in the data file read by the {\hyperref[\detokenize{read_data::doc}]{\sphinxcrossref{\DUrole{doc}{read\_data}}}} command.  Once this has been done, the MDI driver can
send an \textgreater{}ELEMENTS command to the LAMMPS driver with the atomic
number of each atom and the LAMMPS engine will be able to map it to
a LAMMPS atom type.
\end{sphinxadmonition}

\sphinxAtStartPar
The MD and OPTG commands perform an entire MD simulation or energy
minimization (to convergence) with no communication from the driver
until the simulation is complete.  By contrast, the @INIT\_MD and
@INIT\_OPTG commands allow the driver to communicate with the engine at
each timestep of a dynamics run or iteration of a minimization; see
more info below.

\sphinxAtStartPar
The MD command performs a simulation using the most recent \textgreater{}NSTEPS
value.  The OPTG command performs a minimization using the 4
convergence parameters from the most recent \textgreater{}TOLERANCE command.  The 4
parameters sent are those used by the {\hyperref[\detokenize{minimize::doc}]{\sphinxcrossref{\DUrole{doc}{minimize}}}}
command in LAMMPS: etol, ftol, maxiter, and maxeval.

\sphinxAtStartPar
The mdi engine command also implements the following custom MDI
commands which are LAMMPS\sphinxhyphen{}specific.  These commands are also valid at
the @DEFAULT node defined by MDI:
\begin{itemize}
\item {} \begin{itemize}
\item {} 
\sphinxAtStartPar
Command name

\item {} 
\sphinxAtStartPar
Action

\end{itemize}

\item {} \begin{itemize}
\item {} 
\sphinxAtStartPar
\textgreater{}NBYTES

\item {} 
\sphinxAtStartPar
Send \# of datums in a subsequent command (1 value)

\end{itemize}

\item {} \begin{itemize}
\item {} 
\sphinxAtStartPar
\textgreater{}COMMAND

\item {} 
\sphinxAtStartPar
Send a LAMMPS input script command as a string (Nbytes in length)

\end{itemize}

\item {} \begin{itemize}
\item {} 
\sphinxAtStartPar
\textgreater{}COMMANDS

\item {} 
\sphinxAtStartPar
Send multiple LAMMPS input script commands as a newline\sphinxhyphen{}separated string (Nbytes in length)

\end{itemize}

\item {} \begin{itemize}
\item {} 
\sphinxAtStartPar
\textgreater{}INFILE

\item {} 
\sphinxAtStartPar
Send filename of an input script to execute (filename Nbytes in length)

\end{itemize}

\item {} \begin{itemize}
\item {} 
\sphinxAtStartPar
\textless{}KE

\item {} 
\sphinxAtStartPar
Request kinetic energy of the system (1 value)

\end{itemize}

\end{itemize}

\sphinxAtStartPar
Note that other custom commands can easily be added if these are not
sufficient to support what a user\sphinxhyphen{}written driver code needs.  Code to
support new commands can be added to the MDI package within LAMMPS,
specifically to the src/MDI/mdi\_engine.cpp file.

\sphinxAtStartPar
MDI also defines a standard mechanism for the driver to request that
an MD engine (LAMMPS) perform a dynamics simulation one step at a time
or an energy minimization one iteration at a time.  This is so that
the driver can (optionally) communicate with LAMMPS at intermediate
points of the timestep or iteration by issuing MDI node commands which
start with “@”.

\sphinxAtStartPar
To tell LAMMPS to run dynamics in single\sphinxhyphen{}step mode, the driver sends
as @INIT\_MD command followed by the these commands.  The driver
can interact with LAMMPS at 3 node locations within each
timestep: @COORDS, @FORCES, @ENDSTEP:
\begin{itemize}
\item {} \begin{itemize}
\item {} 
\sphinxAtStartPar
Command name

\item {} 
\sphinxAtStartPar
Action

\end{itemize}

\item {} \begin{itemize}
\item {} 
\sphinxAtStartPar
@COORDS

\item {} 
\sphinxAtStartPar
Proceed to next @COORDS node = post\sphinxhyphen{}integrate location in LAMMPS timestep

\end{itemize}

\item {} \begin{itemize}
\item {} 
\sphinxAtStartPar
@FORCES

\item {} 
\sphinxAtStartPar
Proceed to next @FORCES node = post\sphinxhyphen{}force location in LAMMPS timestep

\end{itemize}

\item {} \begin{itemize}
\item {} 
\sphinxAtStartPar
@ENDSTEP

\item {} 
\sphinxAtStartPar
Proceed to next @ENDSTEP node = end\sphinxhyphen{}of\sphinxhyphen{}step location in LAMMPS timestep

\end{itemize}

\item {} \begin{itemize}
\item {} 
\sphinxAtStartPar
@DEFAULT

\item {} 
\sphinxAtStartPar
Exit MD simulation, return to @DEFAULT node

\end{itemize}

\item {} \begin{itemize}
\item {} 
\sphinxAtStartPar
EXIT

\item {} 
\sphinxAtStartPar
Driver tells LAMMPS to exit the MD simulation and engine mode

\end{itemize}

\end{itemize}

\sphinxAtStartPar
To tell LAMMPS to run an energy minimization in single\sphinxhyphen{}iteration mode.
The driver can interact with LAMMPS at 2 node locations within each
iteration of the minimizer: @COORDS, @FORCES:
\begin{itemize}
\item {} \begin{itemize}
\item {} 
\sphinxAtStartPar
Command name

\item {} 
\sphinxAtStartPar
Action

\end{itemize}

\item {} \begin{itemize}
\item {} 
\sphinxAtStartPar
@COORDS

\item {} 
\sphinxAtStartPar
Proceed to next @COORDS node = min\sphinxhyphen{}pre\sphinxhyphen{}force location in LAMMPS min iteration

\end{itemize}

\item {} \begin{itemize}
\item {} 
\sphinxAtStartPar
@FORCES

\item {} 
\sphinxAtStartPar
Proceed to next @FORCES node = min\sphinxhyphen{}post\sphinxhyphen{}force location in LAMMPS min iteration

\end{itemize}

\item {} \begin{itemize}
\item {} 
\sphinxAtStartPar
@DEFAULT

\item {} 
\sphinxAtStartPar
Exit minimization, return to @DEFAULT node

\end{itemize}

\item {} \begin{itemize}
\item {} 
\sphinxAtStartPar
EXIT

\item {} 
\sphinxAtStartPar
Driver tells LAMMPS to exit the minimization and engine mode

\end{itemize}

\end{itemize}

\sphinxAtStartPar
While LAMMPS is at its @COORDS node, the following standard MDI
commands are supported, as documented above: \textgreater{}COORDS or \textless{}COORDS,
@COORDS, @FORCES, @ENDSTEP, @DEFAULT, EXIT.

\sphinxAtStartPar
While LAMMPS is at its @FORCES node, the following standard MDI
commands are supported, as documented above: \textless{}COORDS, \textless{}ENERGY, \textgreater{}FORCES
or \textgreater{}+FORCES or \textless{}FORCES, \textless{}KE, \textless{}PE, \textless{}STRESS, @COORDS, @FORCES, @ENDSTEP,
@DEFAULT, EXIT.

\sphinxAtStartPar
While LAMMPS is at its @ENDSTEP node, the following standard MDI
commands are supported, as documented above: \textless{}ENERGY, \textless{}FORCES, \textless{}KE,
\textless{}PE, \textless{}STRESS, @COORDS, @FORCES, @ENDSTEP, @DEFAULT, EXIT.


\bigskip\hrule\bigskip


\sphinxAtStartPar
The \sphinxstyleemphasis{mdi plugin} command is used to make LAMMPS operate as an MDI
driver which loads an MDI engine as a plugin library.  It is typically
used in an input script after LAMMPS has setup the system it is going
to model consistent with the engine code.

\sphinxAtStartPar
The \sphinxstyleemphasis{name} argument specifies which plugin library to load.  A name
like “lammps” is converted to a filename liblammps.so.  The path for
where this file is located is specified by the \sphinxhyphen{}plugin\_path switch
within the \sphinxhyphen{}mdi command\sphinxhyphen{}line switch, which is specified when LAMMPS is
launched.  See the examples/mdi/README files for examples of how this
is done.

\sphinxAtStartPar
The \sphinxstyleemphasis{mdi} keyword is required and is used as the \sphinxhyphen{}mdi argument passed
to the library when it is launched.  The \sphinxhyphen{}role and \sphinxhyphen{}method settings
are required.  The \sphinxhyphen{}name setting can be anything you choose.  MDI
drivers and engines can query their names to verify they are values
they expect.

\sphinxAtStartPar
The \sphinxstyleemphasis{infile} keyword is optional.  It sets the name of an input script
which the engine will open and process.  MDI will pass it as a
command\sphinxhyphen{}line argument to the library when it is launched.  The file
typically contains settings that an MD or QM code will use for its
calculations.

\sphinxAtStartPar
The \sphinxstyleemphasis{extra} keyword is optional.  It contains additional command\sphinxhyphen{}line
arguments which MDI will pass to the library when it is launched.

\sphinxAtStartPar
The \sphinxstyleemphasis{command} keyword is required.  It specifies a LAMMPS input script
command (as a single argument in quotes if it is multiple words).
Once the plugin library is launched, LAMMPS will execute this command.
Other previously\sphinxhyphen{}defined commands in the input script, such as the
{\hyperref[\detokenize{fix_mdi_qm::doc}]{\sphinxcrossref{\DUrole{doc}{fix mdi/qm}}}} command, should perform MDI
communication with the engine, while the specified \sphinxstyleemphasis{command} executes.
Note that if \sphinxstyleemphasis{command} is an {\hyperref[\detokenize{include::doc}]{\sphinxcrossref{\DUrole{doc}{include}}}} command, then it
could specify a filename with multiple LAMMPS commands.

\begin{sphinxadmonition}{note}{Note:}
\sphinxAtStartPar
When the \sphinxstyleemphasis{command} is complete, LAMMPS will send an MDI EXIT
command to the plugin engine and the plugin will be removed.  The
“mdi plugin” command will then exit and the next command (if any)
in the LAMMPS input script will be processed.  A subsequent “mdi
plugin” command could then load the same or a different MDI
plugin if desired.
\end{sphinxadmonition}


\bigskip\hrule\bigskip


\sphinxAtStartPar
The \sphinxstyleemphasis{mdi connect} and \sphinxstyleemphasis{mdi exit} commands are only used when LAMMPS is
operating as an MDI driver.  And when other LAMMPS command(s) which
send MDI commands and associated data to/from the MDI engine are not
able to initiate and terminate the connection to the engine code.

\sphinxAtStartPar
The only current MDI driver command in LAMMPS is the {\hyperref[\detokenize{fix_mdi_qm::doc}]{\sphinxcrossref{\DUrole{doc}{fix mdi/qm}}}} command.  If it is only used once in an input script
then it can initiate and terminate the connection, but if it is being
issued multiple times (e.g., in a loop that issues a {\hyperref[\detokenize{clear::doc}]{\sphinxcrossref{\DUrole{doc}{clear}}}} command), then it cannot initiate or terminate the connection
multiple times.  Instead, the \sphinxstyleemphasis{mdi connect} and \sphinxstyleemphasis{mdi exit} commands
should be used outside the loop to initiate or terminate the connection.

\sphinxAtStartPar
See the examples/mdi/in.series.driver script for an example of how
this is done.  The LOOP in that script is reading a series of data
file configurations and passing them to an MDI engine (e.g., quantum
code) for energy and force evaluation.  A \sphinxstyleemphasis{clear} command inside the
loop wipes out the current system so a new one can be defined.  This
operation also destroys all fixes.  So the {\hyperref[\detokenize{fix_mdi_qm::doc}]{\sphinxcrossref{\DUrole{doc}{fix mdi/qm}}}} command is issued once per loop iteration.  Note that it
includes a “connect no” option which disables the initiate/terminate
logic within that fix.


\subsection{Restrictions}
\label{\detokenize{mdi:restrictions}}
\sphinxAtStartPar
This command is part of the MDI package.  It is only enabled if LAMMPS
was built with that package.  See the {\hyperref[\detokenize{Build_package::doc}]{\sphinxcrossref{\DUrole{doc}{Build package}}}} page for more info.

\sphinxAtStartPar
To use LAMMPS in conjunction with other MDI\sphinxhyphen{}enabled atomistic codes,
the {\hyperref[\detokenize{units::doc}]{\sphinxcrossref{\DUrole{doc}{units}}}} command should be used to specify \sphinxstyleemphasis{real} or
\sphinxstyleemphasis{metal} units.  This will ensure the correct unit conversions between
LAMMPS and MDI units, which the other codes will also perform in their
preferred units.

\sphinxAtStartPar
LAMMPS can also be used as an MDI engine in other unit choices it
supports (e.g., \sphinxstyleemphasis{lj}), but then no unit conversion is performed.


\subsection{Related commands}
\label{\detokenize{mdi:related-commands}}
\sphinxAtStartPar
{\hyperref[\detokenize{fix_mdi_qm::doc}]{\sphinxcrossref{\DUrole{doc}{fix mdi/qm}}}}


\subsection{Default}
\label{\detokenize{mdi:default}}
\sphinxAtStartPar
None

\sphinxstepscope

\index{min\_modify@\spxentry{min\_modify}}\ignorespaces 

\section{min\_modify command}
\label{\detokenize{min_modify:min-modify-command}}\label{\detokenize{min_modify:index-0}}\label{\detokenize{min_modify::doc}}

\subsection{Syntax}
\label{\detokenize{min_modify:syntax}}
\begin{sphinxVerbatim}[commandchars=\\\{\}]
\PYG{k}{min\PYGZus{}modify}\PYG{+w}{ }\PYG{n}{keyword}\PYG{+w}{ }\PYG{n}{values}\PYG{+w}{ }\PYG{n}{...}
\end{sphinxVerbatim}
\begin{itemize}
\item {} 
\sphinxAtStartPar
one or more keyword/value pairs may be listed
\begin{sphinxalltt}
keyword = \sphinxstyleemphasis{dmax} or \sphinxstyleemphasis{line} or \sphinxstyleemphasis{norm} or \sphinxstyleemphasis{alpha\_damp} or \sphinxstyleemphasis{discrete\_factor} or \sphinxstyleemphasis{integrator} or \sphinxstyleemphasis{abcfire} or \sphinxstyleemphasis{tmax}
  \sphinxstyleemphasis{dmax} value = max
    max = maximum distance for line search to move (distance units)
  \sphinxstyleemphasis{line} value = \sphinxstyleemphasis{backtrack} or \sphinxstyleemphasis{quadratic} or \sphinxstyleemphasis{forcezero} or \sphinxstyleemphasis{spin\_cubic} or \sphinxstyleemphasis{spin\_none}
    backtrack,quadratic,forcezero,spin\_cubic,spin\_none = style of linesearch to use
  \sphinxstyleemphasis{norm} value = \sphinxstyleemphasis{two} or \sphinxstyleemphasis{inf} or \sphinxstyleemphasis{max}
    two = Euclidean two\sphinxhyphen{}norm (length of 3N vector)
    inf = max force component across all 3\sphinxhyphen{}vectors
    max = max force norm across all 3\sphinxhyphen{}vectors
  \sphinxstyleemphasis{alpha\_damp} value = damping
    damping = fictitious magnetic damping for spin minimization (adim)
  \sphinxstyleemphasis{discrete\_factor} value = factor
    factor = discretization factor for adaptive spin timestep (adim)
  \sphinxstyleemphasis{integrator} value = \sphinxstyleemphasis{eulerimplicit} or \sphinxstyleemphasis{verlet} or \sphinxstyleemphasis{leapfrog} or \sphinxstyleemphasis{eulerexplicit}
    time integration scheme for fire minimization
  \sphinxstyleemphasis{abcfire} value = yes or no (default no)
    yes = use ABC\sphinxhyphen{}FIRE variant of fire minimization style
    no  = use default FIRE variant of fire minimization style
  \sphinxstyleemphasis{tmax} value = factor
    factor = maximum adaptive timestep for fire minimization (adim)
\end{sphinxalltt}

\end{itemize}


\subsection{Examples}
\label{\detokenize{min_modify:examples}}
\begin{sphinxVerbatim}[commandchars=\\\{\}]
\PYG{k}{min\PYGZus{}modify}\PYG{+w}{ }\PYG{n}{dmax}\PYG{+w}{ }\PYG{l+m}{0.2}
\PYG{k}{min\PYGZus{}modify}\PYG{+w}{ }\PYG{n}{integrator}\PYG{+w}{ }\PYG{n}{verlet}\PYG{+w}{ }\PYG{n}{tmax}\PYG{+w}{ }\PYG{l+m}{4}
\end{sphinxVerbatim}


\subsection{Description}
\label{\detokenize{min_modify:description}}
\sphinxAtStartPar
This command sets parameters that affect the energy minimization
algorithms selected by the {\hyperref[\detokenize{min_style::doc}]{\sphinxcrossref{\DUrole{doc}{min\_style}}}} command.  The
various settings may affect the convergence rate and overall number of
force evaluations required by a minimization, so users can experiment
with these parameters to tune their minimizations.

\sphinxAtStartPar
The \sphinxstyleemphasis{cg} and \sphinxstyleemphasis{sd} minimization styles have an outer iteration and an
inner iteration which is steps along a one\sphinxhyphen{}dimensional line search in
a particular search direction.  The \sphinxstyleemphasis{dmax} parameter is how far any
atom can move in a single line search in any dimension (x, y, or z).
For the \sphinxstyleemphasis{quickmin} and \sphinxstyleemphasis{fire} minimization styles, the \sphinxstyleemphasis{dmax} setting
is how far any atom can move in a single iteration (timestep).  Thus a
value of 0.1 in real {\hyperref[\detokenize{units::doc}]{\sphinxcrossref{\DUrole{doc}{units}}}} means no atom will move
further than 0.1 Angstroms in a single outer iteration.  This prevents
highly overlapped atoms from being moved long distances (e.g. through
another atom) due to large forces.

\sphinxAtStartPar
The choice of line search algorithm for the \sphinxstyleemphasis{cg} and \sphinxstyleemphasis{sd} minimization
styles can be selected via the \sphinxstyleemphasis{line} keyword.  The default
\sphinxstyleemphasis{quadratic} line search algorithm starts out using the robust
backtracking method described below. However, once the system gets
close to a local minimum and the linesearch steps get small, so that
the energy is approximately quadratic in the step length, it uses the
estimated location of zero gradient as the linesearch step, provided
the energy change is downhill.  This becomes more efficient than
backtracking for highly\sphinxhyphen{}converged relaxations. The \sphinxstyleemphasis{forcezero} line
search algorithm is similar to \sphinxstyleemphasis{quadratic}.  It may be more
efficient than \sphinxstyleemphasis{quadratic} on some systems.

\sphinxAtStartPar
The backtracking search is robust and should always find a local
energy minimum.  However, it will “converge” when it can no longer
reduce the energy of the system.  Individual atom forces may still be
larger than desired at this point, because the energy change is
measured as the difference of two large values (energy before and
energy after) and that difference may be smaller than machine epsilon
even if atoms could move in the gradient direction to reduce forces
further.

\sphinxAtStartPar
The choice of a norm can be modified for the min styles \sphinxstyleemphasis{cg}, \sphinxstyleemphasis{sd}, \sphinxstyleemphasis{quickmin}, \sphinxstyleemphasis{fire}, \sphinxstyleemphasis{fire/old}, \sphinxstyleemphasis{spin}, \sphinxstyleemphasis{spin/cg} and
\sphinxstyleemphasis{spin/lbfgs} using the \sphinxstyleemphasis{norm} keyword.  The default \sphinxstyleemphasis{two} norm computes
the 2\sphinxhyphen{}norm (Euclidean length) of the global force vector:
\begin{equation*}
\begin{split}|| \vec{F} ||_{2} = \sqrt{\vec{F}_1^2+ \cdots + \vec{F}_N^2}\end{split}
\end{equation*}
\sphinxAtStartPar
The \sphinxstyleemphasis{max} norm computes the length of the 3\sphinxhyphen{}vector force
for each atom  (2\sphinxhyphen{}norm), and takes the maximum value of those across
all atoms
\begin{equation*}
\begin{split}|| \vec{F} ||_{max} = {\rm max}\left(||\vec{F}_1||, \cdots, ||\vec{F}_N||\right)\end{split}
\end{equation*}
\sphinxAtStartPar
The \sphinxstyleemphasis{inf} norm takes the maximum component across the forces of
all atoms in the system:
\begin{equation*}
\begin{split}|| \vec{F} ||_{inf} = {\rm max}\left(|F_1^1|, |F_1^2|, |F_1^3| \cdots, |F_N^1|, |F_N^2|, |F_N^3|\right)\end{split}
\end{equation*}
\sphinxAtStartPar
For the min styles \sphinxstyleemphasis{spin}, \sphinxstyleemphasis{spin/cg} and \sphinxstyleemphasis{spin/lbfgs}, the force
norm is replaced by the spin\sphinxhyphen{}torque norm.

\sphinxAtStartPar
Keywords \sphinxstyleemphasis{alpha\_damp} and \sphinxstyleemphasis{discrete\_factor} only make sense when
a {\hyperref[\detokenize{min_spin::doc}]{\sphinxcrossref{\DUrole{doc}{min\_spin}}}} command is declared.
Keyword \sphinxstyleemphasis{alpha\_damp} defines an analog of a magnetic damping.
It defines a relaxation rate toward an equilibrium for a given
magnetic system.
Keyword \sphinxstyleemphasis{discrete\_factor} defines a discretization factor for the
adaptive timestep used in the \sphinxstyleemphasis{spin} minimization.
See {\hyperref[\detokenize{min_spin::doc}]{\sphinxcrossref{\DUrole{doc}{min\_spin}}}} for more information about those
quantities.

\sphinxAtStartPar
The choice of a line search algorithm for the \sphinxstyleemphasis{spin/cg} and
\sphinxstyleemphasis{spin/lbfgs} styles can be specified via the \sphinxstyleemphasis{line} keyword.  The
\sphinxstyleemphasis{spin\_cubic} and \sphinxstyleemphasis{spin\_none} keywords only make sense when one of those two
minimization styles is declared.  The \sphinxstyleemphasis{spin\_cubic} performs the line
search based on a cubic interpolation of the energy along the search
direction. The \sphinxstyleemphasis{spin\_none} keyword deactivates the line search
procedure.  The \sphinxstyleemphasis{spin\_none} is a default value for \sphinxstyleemphasis{line} keyword for
both \sphinxstyleemphasis{spin/lbfgs} and \sphinxstyleemphasis{spin/cg}. Convergence of \sphinxstyleemphasis{spin/lbfgs} can be
more robust if \sphinxstyleemphasis{spin\_cubic} line search is used.

\sphinxAtStartPar
The Newton \sphinxstyleemphasis{integrator} used for \sphinxstyleemphasis{fire} minimization can be selected to
be either the symplectic Euler (\sphinxstyleemphasis{eulerimplicit}), velocity Verlet (\sphinxstyleemphasis{verlet}), Leapfrog (\sphinxstyleemphasis{leapfrog}) or non\sphinxhyphen{}symplectic forward Euler
(\sphinxstyleemphasis{eulerexplicit} ). The keyword \sphinxstyleemphasis{tmax} defines the maximum value for
the adaptive timestep during a \sphinxstyleemphasis{fire} minimization. It is a
multiplication factor applied to the current {\hyperref[\detokenize{timestep::doc}]{\sphinxcrossref{\DUrole{doc}{timestep}}}}
(not in time unit). For example, \sphinxstyleemphasis{tmax} = 4.0 with a {\hyperref[\detokenize{timestep::doc}]{\sphinxcrossref{\DUrole{doc}{timestep}}}} of 2fs, means that the maximum value the timestep can reach
during a \sphinxstyleemphasis{fire} minimization is 4fs.  Note that parameter defaults has
been chosen to be reliable in most cases, but one should consider
adjusting {\hyperref[\detokenize{timestep::doc}]{\sphinxcrossref{\DUrole{doc}{timestep}}}} and \sphinxstyleemphasis{tmax} to optimize the
minimization for large or complex systems.  Other parameters of the
\sphinxstyleemphasis{fire} minimization can be tuned (\sphinxstyleemphasis{tmin}, \sphinxstyleemphasis{delaystep}, \sphinxstyleemphasis{dtgrow},
\sphinxstyleemphasis{dtshrink}, \sphinxstyleemphasis{alpha0}, and \sphinxstyleemphasis{alphashrink}). Please refer to the
references describing the {\hyperref[\detokenize{min_style::doc}]{\sphinxcrossref{\DUrole{doc}{min\_style}}}} \sphinxstyleemphasis{fire}.  An
additional stopping criteria \sphinxstyleemphasis{vdfmax} is used by \sphinxstyleemphasis{fire} in order to
avoid unnecessary looping when it is reasonable to think the system will
not be relaxed further.  Note that in this case the system will NOT have
reached your minimization criteria. This could happen when the system
comes to be stuck in a local basin of the phase space.  \sphinxstyleemphasis{vdfmax} is the
maximum number of consecutive iterations with P(t) \textless{} 0.

\sphinxAtStartPar
\DUrole{versionmodified,added}{Added in version 8Feb2023.}

\sphinxAtStartPar
The \sphinxstyleemphasis{abcfire} keyword allows to activate the ABC\sphinxhyphen{}FIRE variant of the
\sphinxstyleemphasis{fire} minimization algorithm. ABC\sphinxhyphen{}FIRE introduces an additional factor
that modifies the bias and scaling of the velocities of the atoms during
the mixing step {\hyperref[\detokenize{min_modify:echeverrirestrepo}]{\sphinxcrossref{\DUrole{std,std-ref}{(Echeverri Restrepo)}}}}.  This
can lead to faster convergence of the minimizer.

\sphinxAtStartPar
The {\hyperref[\detokenize{min_style::doc}]{\sphinxcrossref{\DUrole{doc}{min\_style}}}} \sphinxstyleemphasis{fire} is an optimized implementation of
{\hyperref[\detokenize{min_style::doc}]{\sphinxcrossref{\DUrole{doc}{min\_style}}}} \sphinxstyleemphasis{fire/old}. It can however behave similarly
to the \sphinxstyleemphasis{fire/old} style by using the following set of parameters:

\begin{sphinxVerbatim}[commandchars=\\\{\}]
\PYG{k}{min\PYGZus{}modify}\PYG{+w}{ }\PYG{n}{integrator}\PYG{+w}{ }\PYG{n}{eulerexplicit}\PYG{+w}{ }\PYG{n}{tmax}\PYG{+w}{ }\PYG{l+m}{10.0}\PYG{+w}{ }\PYG{n}{tmin}\PYG{+w}{ }\PYG{l+m}{0.0}\PYG{+w}{ }\PYG{n}{delaystep}\PYG{+w}{ }\PYG{l+m}{5}\PYG{+w}{ }\PYG{o}{\PYGZam{}}
\PYG{+w}{          }\PYG{n}{dtgrow}\PYG{+w}{ }\PYG{l+m}{1.1}\PYG{+w}{ }\PYG{n}{dtshrink}\PYG{+w}{ }\PYG{l+m}{0.5}\PYG{+w}{ }\PYG{n}{alpha0}\PYG{+w}{ }\PYG{l+m}{0.1}\PYG{+w}{ }\PYG{n}{alphashrink}\PYG{+w}{ }\PYG{l+m}{0.99}\PYG{+w}{ }\PYG{o}{\PYGZam{}}
\PYG{+w}{          }\PYG{n}{vdfmax}\PYG{+w}{ }\PYG{l+m}{100000}\PYG{+w}{ }\PYG{n}{halfstepback}\PYG{+w}{ }\PYG{n}{no}\PYG{+w}{ }\PYG{n}{initialdelay}\PYG{+w}{ }\PYG{n}{no}
\end{sphinxVerbatim}


\subsection{Restrictions}
\label{\detokenize{min_modify:restrictions}}
\sphinxAtStartPar
For magnetic GNEB calculations, only \sphinxstyleemphasis{spin\_none} value for \sphinxstyleemphasis{line}
keyword can be used when minimization styles \sphinxstyleemphasis{spin/cg} and \sphinxstyleemphasis{spin/lbfgs} are
employed.  See {\hyperref[\detokenize{neb_spin::doc}]{\sphinxcrossref{\DUrole{doc}{neb/spin}}}} for more explanation.


\subsection{Related commands}
\label{\detokenize{min_modify:related-commands}}
\sphinxAtStartPar
{\hyperref[\detokenize{min_style::doc}]{\sphinxcrossref{\DUrole{doc}{min\_style}}}}, {\hyperref[\detokenize{minimize::doc}]{\sphinxcrossref{\DUrole{doc}{minimize}}}}


\subsection{Default}
\label{\detokenize{min_modify:default}}
\sphinxAtStartPar
The option defaults are dmax = 0.1, line = quadratic and norm = two.

\sphinxAtStartPar
For the \sphinxstyleemphasis{spin}, \sphinxstyleemphasis{spin/cg} and \sphinxstyleemphasis{spin/lbfgs} styles, the option
defaults are alpha\_damp = 1.0, discrete\_factor = 10.0, line =
spin\_none, and norm = euclidean.

\sphinxAtStartPar
For the \sphinxstyleemphasis{fire} style, the option defaults are integrator =
eulerimplicit, tmax = 10.0, tmin = 0.02, delaystep = 20, dtgrow = 1.1,
dtshrink = 0.5, alpha0 = 0.25, alphashrink = 0.99, vdfmax = 2000,
halfstepback = yes and initialdelay = yes.

\phantomsection\label{\detokenize{min_modify:echeverrirestrepo}}
\sphinxAtStartPar
\sphinxstylestrong{(EcheverriRestrepo)} Echeverri Restrepo, Andric, Comput Mater Sci, 218, 111978 (2023).

\sphinxstepscope

\index{min\_style spin@\spxentry{min\_style spin}}\ignorespaces 

\section{min\_style spin command}
\label{\detokenize{min_spin:min-style-spin-command}}\label{\detokenize{min_spin:index-0}}\label{\detokenize{min_spin::doc}}

\section{min\_style spin/cg command}
\label{\detokenize{min_spin:min-style-spin-cg-command}}

\section{min\_style spin/lbfgs command}
\label{\detokenize{min_spin:min-style-spin-lbfgs-command}}

\subsection{Syntax}
\label{\detokenize{min_spin:syntax}}
\begin{sphinxVerbatim}[commandchars=\\\{\}]
\PYG{k}{min\PYGZus{}style}\PYG{+w}{ }\PYG{n}{spin}
\PYG{k}{min\PYGZus{}style}\PYG{+w}{ }\PYG{n}{spin}\PYG{o}{/}\PYG{n}{cg}
\PYG{k}{min\PYGZus{}style}\PYG{+w}{ }\PYG{n}{spin}\PYG{o}{/}\PYG{n}{lbfgs}
\end{sphinxVerbatim}


\subsection{Examples}
\label{\detokenize{min_spin:examples}}
\begin{sphinxVerbatim}[commandchars=\\\{\}]
\PYG{k}{min\PYGZus{}style}\PYG{+w}{  }\PYG{n}{spin}\PYG{o}{/}\PYG{n}{lbfgs}
\PYG{k}{min\PYGZus{}modify}\PYG{+w}{ }\PYG{n}{line}\PYG{+w}{ }\PYG{n}{spin\PYGZus{}cubic}\PYG{+w}{ }\PYG{n}{discrete\PYGZus{}factor}\PYG{+w}{ }\PYG{l+m}{10.0}
\end{sphinxVerbatim}


\subsection{Description}
\label{\detokenize{min_spin:description}}
\sphinxAtStartPar
Apply a minimization algorithm to use when a {\hyperref[\detokenize{minimize::doc}]{\sphinxcrossref{\DUrole{doc}{minimize}}}}
command is performed.

\sphinxAtStartPar
Style \sphinxstyleemphasis{spin} defines a damped spin dynamics with an adaptive
timestep, according to:
\begin{equation*}
\begin{split}\frac{d \vec{s}_{i}}{dt} = \lambda\, \vec{s}_{i} \times\left( \vec{\omega}_{i} \times\vec{s}_{i} \right)\end{split}
\end{equation*}
\sphinxAtStartPar
with \(\lambda\) a damping coefficient (similar to a magnetic
damping). \(\lambda\) can be defined by setting the
\sphinxstyleemphasis{alpha\_damp} keyword with the {\hyperref[\detokenize{min_modify::doc}]{\sphinxcrossref{\DUrole{doc}{min\_modify}}}} command.

\sphinxAtStartPar
The minimization procedure solves this equation using an
adaptive timestep. The value of this timestep is defined
by the largest precession frequency that has to be solved in the
system:
\begin{equation*}
\begin{split}{\Delta t}_{\rm max} = \frac{2\pi}{\kappa \left|\vec{\omega}_{\rm max} \right|}\end{split}
\end{equation*}
\sphinxAtStartPar
with \(\left|\vec{\omega}_{\rm max}\right|\) the norm of the largest precession
frequency in the system (across all processes, and across all replicas if a
spin/neb calculation is performed).

\sphinxAtStartPar
\(\kappa\) defines a discretization factor \sphinxstyleemphasis{discrete\_factor} for
the definition of this timestep.  \sphinxstyleemphasis{discrete\_factor} can be defined with
the {\hyperref[\detokenize{min_modify::doc}]{\sphinxcrossref{\DUrole{doc}{min\_modify}}}} command.

\sphinxAtStartPar
Style \sphinxstyleemphasis{spin/cg} defines an orthogonal spin optimization
(OSO) combined to a conjugate gradient (CG) algorithm.
The {\hyperref[\detokenize{min_modify::doc}]{\sphinxcrossref{\DUrole{doc}{min\_modify}}}} command can be used to
couple the \sphinxstyleemphasis{spin/cg} to a line search procedure, and to modify the
discretization factor \sphinxstyleemphasis{discrete\_factor}.
By default, style \sphinxstyleemphasis{spin/cg} does not employ the line search procedure
and uses the adaptive time\sphinxhyphen{}step technique in the same way as style \sphinxstyleemphasis{spin}.

\sphinxAtStartPar
Style \sphinxstyleemphasis{spin/lbfgs} defines an orthogonal spin optimization (OSO)
combined to a limited\sphinxhyphen{}memory Broyden\sphinxhyphen{}Fletcher\sphinxhyphen{}Goldfarb\sphinxhyphen{}Shanno (L\sphinxhyphen{}BFGS)
algorithm.  By default, style \sphinxstyleemphasis{spin/lbfgs} does not employ line search
procedure.  If the line search procedure is not used then the discrete
factor defines the maximum root mean squared rotation angle of spins by
equation \sphinxstyleemphasis{pi/(5*Kappa)}.  The default value for Kappa is 10.  The
\sphinxstyleemphasis{spin\_cubic} line search option can improve the convergence of the
\sphinxstyleemphasis{spin/lbfgs} algorithm.

\sphinxAtStartPar
The {\hyperref[\detokenize{min_modify::doc}]{\sphinxcrossref{\DUrole{doc}{min\_modify}}}} command can be used to
activate the line search procedure, and to modify the
discretization factor \sphinxstyleemphasis{discrete\_factor}.

\sphinxAtStartPar
For more information about styles \sphinxstyleemphasis{spin/cg} and \sphinxstyleemphasis{spin/lbfgs},
see their implementation reported in {\hyperref[\detokenize{min_spin:ivanov1}]{\sphinxcrossref{\DUrole{std,std-ref}{(Ivanov)}}}}.

\begin{sphinxadmonition}{note}{Note:}
\sphinxAtStartPar
All the \sphinxstyleemphasis{spin} styles replace the force tolerance by a torque
tolerance. See {\hyperref[\detokenize{minimize::doc}]{\sphinxcrossref{\DUrole{doc}{minimize}}}} for more explanation.
\end{sphinxadmonition}

\begin{sphinxadmonition}{note}{Note:}
\sphinxAtStartPar
The \sphinxstyleemphasis{spin/cg} and \sphinxstyleemphasis{spin/lbfgs} styles can be used
for magnetic NEB calculations only if the line search procedure
is deactivated. See {\hyperref[\detokenize{neb_spin::doc}]{\sphinxcrossref{\DUrole{doc}{neb/spin}}}} for more explanation.
\end{sphinxadmonition}


\subsection{Restrictions}
\label{\detokenize{min_spin:restrictions}}
\sphinxAtStartPar
The \sphinxstyleemphasis{spin}, \sphinxstyleemphasis{spin/cg}, and \sphinxstyleemphasis{spin/lbfgps} styles are part of the SPIN
package.  They are only enabled if LAMMPS was built with that package.
See the {\hyperref[\detokenize{Build_package::doc}]{\sphinxcrossref{\DUrole{doc}{Build package}}}} page for more info.

\sphinxAtStartPar
This minimization procedure is only applied to spin degrees of
freedom for a frozen lattice configuration.


\subsection{Related commands}
\label{\detokenize{min_spin:related-commands}}
\sphinxAtStartPar
{\hyperref[\detokenize{min_style::doc}]{\sphinxcrossref{\DUrole{doc}{min\_style}}}}, {\hyperref[\detokenize{minimize::doc}]{\sphinxcrossref{\DUrole{doc}{minimize}}}},
{\hyperref[\detokenize{min_modify::doc}]{\sphinxcrossref{\DUrole{doc}{min\_modify}}}}


\subsection{Default}
\label{\detokenize{min_spin:default}}
\sphinxAtStartPar
The option defaults are \sphinxstyleemphasis{alpha\_damp} = 1.0, \sphinxstyleemphasis{discrete\_factor} =
10.0, \sphinxstyleemphasis{line} = spin\_none and \sphinxstyleemphasis{norm} = euclidean.


\bigskip\hrule\bigskip

\phantomsection\label{\detokenize{min_spin:ivanov1}}
\sphinxAtStartPar
\sphinxstylestrong{(Ivanov)} Ivanov, Uzdin, Jonsson. arXiv preprint arXiv:1904.02669, (2019).

\sphinxstepscope

\index{min\_style@\spxentry{min\_style}}\ignorespaces 

\section{min\_style cg command}
\label{\detokenize{min_style:min-style-cg-command}}\label{\detokenize{min_style:index-0}}\label{\detokenize{min_style::doc}}

\section{min\_style hftn command}
\label{\detokenize{min_style:min-style-hftn-command}}

\section{min\_style sd command}
\label{\detokenize{min_style:min-style-sd-command}}

\section{min\_style quickmin command}
\label{\detokenize{min_style:min-style-quickmin-command}}

\section{min\_style fire command}
\label{\detokenize{min_style:min-style-fire-command}}

\section{min\_style spin command}
\label{\detokenize{min_style:min-style-spin-command}}

\section{min\_style spin/cg command}
\label{\detokenize{min_style:min-style-spin-cg-command}}

\section{min\_style spin/lbfgs command}
\label{\detokenize{min_style:min-style-spin-lbfgs-command}}

\subsection{Syntax}
\label{\detokenize{min_style:syntax}}
\begin{sphinxVerbatim}[commandchars=\\\{\}]
\PYG{k}{min\PYGZus{}style}\PYG{+w}{ }\PYG{n}{style}
\end{sphinxVerbatim}
\begin{itemize}
\item {} 
\sphinxAtStartPar
style = \sphinxstyleemphasis{cg} or \sphinxstyleemphasis{hftn} or \sphinxstyleemphasis{sd} or \sphinxstyleemphasis{quickmin} or \sphinxstyleemphasis{fire} or \sphinxstyleemphasis{spin} or \sphinxstyleemphasis{spin/cg} or \sphinxstyleemphasis{spin/lbfgs}
\begin{sphinxalltt}
\sphinxstyleemphasis{spin} is discussed briefly here and fully on {\hyperref[\detokenize{min_spin::doc}]{\sphinxcrossref{\DUrole{doc}{min\_style spin}}}} doc page
\sphinxstyleemphasis{spin/cg} is discussed briefly here and fully on {\hyperref[\detokenize{min_spin::doc}]{\sphinxcrossref{\DUrole{doc}{min\_style spin}}}} doc page
\sphinxstyleemphasis{spin/lbfgs} is discussed briefly here and fully on {\hyperref[\detokenize{min_spin::doc}]{\sphinxcrossref{\DUrole{doc}{min\_style spin}}}} doc page
\end{sphinxalltt}

\end{itemize}


\subsection{Examples}
\label{\detokenize{min_style:examples}}
\begin{sphinxVerbatim}[commandchars=\\\{\}]
\PYG{k}{min\PYGZus{}style}\PYG{+w}{ }\PYG{n}{cg}
\PYG{k}{min\PYGZus{}style}\PYG{+w}{ }\PYG{n}{fire}
\PYG{k}{min\PYGZus{}style}\PYG{+w}{ }\PYG{n}{spin}
\end{sphinxVerbatim}


\subsection{Description}
\label{\detokenize{min_style:description}}
\sphinxAtStartPar
Choose a minimization algorithm to use when a {\hyperref[\detokenize{minimize::doc}]{\sphinxcrossref{\DUrole{doc}{minimize}}}} command is performed.

\sphinxAtStartPar
Style \sphinxstyleemphasis{cg} is the Polak\sphinxhyphen{}Ribiere version of the conjugate gradient (CG)
algorithm.  At each iteration the force gradient is combined with the
previous iteration information to compute a new search direction
perpendicular (conjugate) to the previous search direction.  The PR
variant affects how the direction is chosen and how the CG method is
restarted when it ceases to make progress.  The PR variant is thought
to be the most effective CG choice for most problems.

\sphinxAtStartPar
Style \sphinxstyleemphasis{hftn} is a Hessian\sphinxhyphen{}free truncated Newton algorithm.  At each
iteration a quadratic model of the energy potential is solved by a
conjugate gradient inner iteration.  The Hessian (second derivatives)
of the energy is not formed directly, but approximated in each
conjugate search direction by a finite difference directional
derivative.  When close to an energy minimum, the algorithm behaves
like a Newton method and exhibits a quadratic convergence rate to high
accuracy.  In most cases the behavior of \sphinxstyleemphasis{hftn} is similar to \sphinxstyleemphasis{cg},
but it offers an alternative if \sphinxstyleemphasis{cg} seems to perform poorly.  This
style is not affected by the {\hyperref[\detokenize{min_modify::doc}]{\sphinxcrossref{\DUrole{doc}{min\_modify}}}} command.

\sphinxAtStartPar
Style \sphinxstyleemphasis{sd} is a steepest descent algorithm.  At each iteration, the
search direction is set to the downhill direction corresponding to the
force vector (negative gradient of energy).  Typically, steepest
descent will not converge as quickly as CG, but may be more robust in
some situations.

\sphinxAtStartPar
Style \sphinxstyleemphasis{quickmin} is a damped dynamics method described in
{\hyperref[\detokenize{min_style:sheppard}]{\sphinxcrossref{\DUrole{std,std-ref}{(Sheppard)}}}}, where the damping parameter is related
to the projection of the velocity vector along the current force
vector for each atom.  The velocity of each atom is initialized to 0.0
by this style, at the beginning of a minimization.

\sphinxAtStartPar
Style \sphinxstyleemphasis{fire} is a damped dynamics method described in {\hyperref[\detokenize{min_style:bitzek}]{\sphinxcrossref{\DUrole{std,std-ref}{(Bitzek)}}}}, which is similar to \sphinxstyleemphasis{quickmin} but adds a variable timestep
and alters the projection operation to maintain components of the
velocity non\sphinxhyphen{}parallel to the current force vector.  The velocity of each
atom is initialized to 0.0 by this style, at the beginning of a
minimization.  This style correspond to an optimized version described
in {\hyperref[\detokenize{min_style:guenole}]{\sphinxcrossref{\DUrole{std,std-ref}{(Guenole)}}}} that include different time integration
schemes and default parameters.  The default parameters can be modified
with the command {\hyperref[\detokenize{min_modify::doc}]{\sphinxcrossref{\DUrole{doc}{min\_modify}}}}.

\sphinxAtStartPar
Style \sphinxstyleemphasis{spin} is a damped spin dynamics with an adaptive timestep.

\sphinxAtStartPar
Style \sphinxstyleemphasis{spin/cg} uses an orthogonal spin optimization (OSO) combined to
a conjugate gradient (CG) approach to minimize spin configurations.

\sphinxAtStartPar
Style \sphinxstyleemphasis{spin/lbfgs} uses an orthogonal spin optimization (OSO) combined
to a limited\sphinxhyphen{}memory Broyden\sphinxhyphen{}Fletcher\sphinxhyphen{}Goldfarb\sphinxhyphen{}Shanno (LBFGS) approach
to minimize spin configurations.

\sphinxAtStartPar
See the {\hyperref[\detokenize{min_spin::doc}]{\sphinxcrossref{\DUrole{doc}{min/spin}}}} page for more information about
the \sphinxstyleemphasis{spin}, \sphinxstyleemphasis{spin/cg} and \sphinxstyleemphasis{spin/lbfgs} styles.

\sphinxAtStartPar
Either the \sphinxstyleemphasis{quickmin} or the \sphinxstyleemphasis{fire} styles are useful in the context of
nudged elastic band (NEB) calculations via the {\hyperref[\detokenize{neb::doc}]{\sphinxcrossref{\DUrole{doc}{neb}}}} command.

\sphinxAtStartPar
Either the \sphinxstyleemphasis{spin}, \sphinxstyleemphasis{spin/cg}, or \sphinxstyleemphasis{spin/lbfgs} styles are useful in the
context of magnetic geodesic nudged elastic band (GNEB) calculations via
the {\hyperref[\detokenize{neb_spin::doc}]{\sphinxcrossref{\DUrole{doc}{neb/spin}}}} command.

\begin{sphinxadmonition}{note}{Note:}
\sphinxAtStartPar
The damped dynamic minimizers use whatever timestep you have
defined via the {\hyperref[\detokenize{timestep::doc}]{\sphinxcrossref{\DUrole{doc}{timestep}}}} command.  Often they
will converge more quickly if you use a timestep about 10x larger
than you would normally use for dynamics simulations.
For \sphinxstyleemphasis{fire}, the default timestep is recommended to be equal to
the one you would normally use for dynamics simulations.
\end{sphinxadmonition}

\begin{sphinxadmonition}{note}{Note:}
\sphinxAtStartPar
The \sphinxstyleemphasis{quickmin}, \sphinxstyleemphasis{fire}, \sphinxstyleemphasis{hftn}, and \sphinxstyleemphasis{cg/kk} styles do not yet
support the use of the {\hyperref[\detokenize{fix_box_relax::doc}]{\sphinxcrossref{\DUrole{doc}{fix box/relax}}}} command
or minimizations involving the electron radius in {\hyperref[\detokenize{pair_eff::doc}]{\sphinxcrossref{\DUrole{doc}{eFF}}}} models.
\end{sphinxadmonition}


\bigskip\hrule\bigskip


\sphinxAtStartPar
Styles with a \sphinxstyleemphasis{gpu}, \sphinxstyleemphasis{intel}, \sphinxstyleemphasis{kk}, \sphinxstyleemphasis{omp}, or \sphinxstyleemphasis{opt} suffix are
functionally the same as the corresponding style without the suffix.
They have been optimized to run faster, depending on your available
hardware, as discussed on the {\hyperref[\detokenize{Speed_packages::doc}]{\sphinxcrossref{\DUrole{doc}{Accelerator packages}}}}
page.  The accelerated styles take the same arguments and should
produce the same results, except for round\sphinxhyphen{}off and precision issues.

\sphinxAtStartPar
These accelerated styles are part of the GPU, INTEL, KOKKOS,
OPENMP, and OPT packages, respectively.  They are only enabled if
LAMMPS was built with those packages.  See the {\hyperref[\detokenize{Build_package::doc}]{\sphinxcrossref{\DUrole{doc}{Build package}}}} page for more info.

\sphinxAtStartPar
You can specify the accelerated styles explicitly in your input script
by including their suffix, or you can use the {\hyperref[\detokenize{Run_options::doc}]{\sphinxcrossref{\DUrole{doc}{\sphinxhyphen{}suffix command\sphinxhyphen{}line switch}}}} when you invoke LAMMPS, or you can use the
{\hyperref[\detokenize{suffix::doc}]{\sphinxcrossref{\DUrole{doc}{suffix}}}} command in your input script.

\sphinxAtStartPar
See the {\hyperref[\detokenize{Speed_packages::doc}]{\sphinxcrossref{\DUrole{doc}{Accelerator packages}}}} page for more
instructions on how to use the accelerated styles effectively.


\bigskip\hrule\bigskip



\subsection{Restrictions}
\label{\detokenize{min_style:restrictions}}
\sphinxAtStartPar
The \sphinxstyleemphasis{spin}, \sphinxstyleemphasis{spin/cg}, and \sphinxstyleemphasis{spin/lbfgps} styles are part of the SPIN
package.  They are only enabled if LAMMPS was built with that package.
See the {\hyperref[\detokenize{Build_package::doc}]{\sphinxcrossref{\DUrole{doc}{Build package}}}} page for more info.


\subsection{Related commands}
\label{\detokenize{min_style:related-commands}}
\sphinxAtStartPar
{\hyperref[\detokenize{min_modify::doc}]{\sphinxcrossref{\DUrole{doc}{min\_modify}}}}, {\hyperref[\detokenize{minimize::doc}]{\sphinxcrossref{\DUrole{doc}{minimize}}}}, {\hyperref[\detokenize{neb::doc}]{\sphinxcrossref{\DUrole{doc}{neb}}}}


\subsection{Default}
\label{\detokenize{min_style:default}}
\begin{sphinxVerbatim}[commandchars=\\\{\}]
\PYG{k}{min\PYGZus{}style}\PYG{+w}{ }\PYG{n}{cg}
\end{sphinxVerbatim}


\bigskip\hrule\bigskip

\phantomsection\label{\detokenize{min_style:sheppard}}
\sphinxAtStartPar
\sphinxstylestrong{(Sheppard)} Sheppard, Terrell, Henkelman, J Chem Phys, 128, 134106
(2008).  See ref 1 in this paper for original reference to Qmin in
Jonsson, Mills, Jacobsen.

\phantomsection\label{\detokenize{min_style:bitzek}}
\sphinxAtStartPar
\sphinxstylestrong{(Bitzek)} Bitzek, Koskinen, Gahler, Moseler, Gumbsch, Phys Rev Lett,
97, 170201 (2006).

\phantomsection\label{\detokenize{min_style:guenole}}
\sphinxAtStartPar
\sphinxstylestrong{(Guenole)} Guenole, Noehring, Vaid, Houlle, Xie, Prakash, Bitzek,
Comput Mater Sci, 175, 109584 (2020).

\sphinxstepscope

\index{minimize@\spxentry{minimize}}\ignorespaces 
\index{minimize/kk@\spxentry{minimize/kk}}\ignorespaces 

\section{minimize command}
\label{\detokenize{minimize:minimize-command}}\label{\detokenize{minimize:index-1}}\label{\detokenize{minimize:index-0}}\label{\detokenize{minimize::doc}}
\sphinxAtStartPar
Accelerator Variant: minimize/kk


\subsection{Syntax}
\label{\detokenize{minimize:syntax}}
\begin{sphinxVerbatim}[commandchars=\\\{\}]
\PYG{k}{minimize}\PYG{+w}{ }\PYG{n}{etol}\PYG{+w}{ }\PYG{n}{ftol}\PYG{+w}{ }\PYG{n}{maxiter}\PYG{+w}{ }\PYG{n}{maxeval}
\end{sphinxVerbatim}
\begin{itemize}
\item {} 
\sphinxAtStartPar
etol = stopping tolerance for energy (unitless)

\item {} 
\sphinxAtStartPar
ftol = stopping tolerance for force (force units)

\item {} 
\sphinxAtStartPar
maxiter = max iterations of minimizer

\item {} 
\sphinxAtStartPar
maxeval = max number of force/energy evaluations

\end{itemize}


\subsection{Examples}
\label{\detokenize{minimize:examples}}
\begin{sphinxVerbatim}[commandchars=\\\{\}]
\PYG{k}{minimize}\PYG{+w}{ }\PYG{l+m}{1.0e\PYGZhy{}4}\PYG{+w}{ }\PYG{l+m}{1.0e\PYGZhy{}6}\PYG{+w}{ }\PYG{l+m}{100}\PYG{+w}{ }\PYG{l+m}{1000}
\PYG{k}{minimize}\PYG{+w}{ }\PYG{l+m}{0.0}\PYG{+w}{ }\PYG{l+m}{1.0e\PYGZhy{}8}\PYG{+w}{ }\PYG{l+m}{1000}\PYG{+w}{ }\PYG{l+m}{100000}
\end{sphinxVerbatim}


\subsection{Description}
\label{\detokenize{minimize:description}}
\sphinxAtStartPar
Perform an energy minimization of the system, by iteratively adjusting
atom coordinates.  Iterations are terminated when one of the stopping
criteria is satisfied.  At that point the configuration will hopefully
be in a local potential energy minimum.  More precisely, the
configuration should approximate a critical point for the objective
function (see below), which may or may not be a local minimum.

\sphinxAtStartPar
The minimization algorithm used is set by the {\hyperref[\detokenize{min_style::doc}]{\sphinxcrossref{\DUrole{doc}{min\_style}}}} command.  Other options are set by the {\hyperref[\detokenize{min_modify::doc}]{\sphinxcrossref{\DUrole{doc}{min\_modify}}}} command.  Minimize commands can be interspersed with
{\hyperref[\detokenize{run::doc}]{\sphinxcrossref{\DUrole{doc}{run}}}} commands to alternate between relaxation and dynamics.
The minimizers bound the distance atoms may move in one iteration, so
that you can relax systems with highly overlapped atoms (large energies
and forces) by pushing the atoms off of each other.

\begin{sphinxadmonition}{note}{Neighbor list update settings}

\sphinxAtStartPar
The distance that atoms can move during individual minimization steps
can be quite large, especially at the beginning of a minimization.
Thus {\hyperref[\detokenize{neigh_modify::doc}]{\sphinxcrossref{\DUrole{doc}{neighbor list settings}}}} of \sphinxstyleemphasis{every = 1} and
\sphinxstyleemphasis{delay = 0} are \sphinxstylestrong{required}.  This may be combined with either
\sphinxstyleemphasis{check = no} (always update the neighbor list) or \sphinxstyleemphasis{check = yes} (only
update the neighbor list if at least one atom has moved more than
half the {\hyperref[\detokenize{neighbor::doc}]{\sphinxcrossref{\DUrole{doc}{neighbor list skin}}}} distance since the last
reneighboring).  Using \sphinxstyleemphasis{check = yes} is recommended since it avoids
unneeded reneighboring steps when the system is closer to the minimum
and thus atoms move only small distances.  Using \sphinxstyleemphasis{check = no} may be
required for debugging or when coupling LAMMPS with external codes
that require a predictable sequence of neighbor list updates.

\sphinxAtStartPar
If the settings are \sphinxstylestrong{not} \sphinxstyleemphasis{every = 1} and \sphinxstyleemphasis{delay = 0}, LAMMPS will
temporarily apply a {\hyperref[\detokenize{neigh_modify::doc}]{\sphinxcrossref{\DUrole{doc}{neigh\_modify every 1 delay 0 check yes}}}} setting during the minimization and restore the
original setting at the end of the minimization.  A corresponding
message will be printed to the screen and log file, if this happens.
\end{sphinxadmonition}

\sphinxAtStartPar
Alternate means of relaxing a system are to run dynamics with a small or
{\hyperref[\detokenize{fix_nve_limit::doc}]{\sphinxcrossref{\DUrole{doc}{limited timestep}}}}.  Or dynamics can be run using
{\hyperref[\detokenize{fix_viscous::doc}]{\sphinxcrossref{\DUrole{doc}{fix viscous}}}} to impose a damping force that slowly
drains all kinetic energy from the system.  The {\hyperref[\detokenize{pair_soft::doc}]{\sphinxcrossref{\DUrole{doc}{pair\_style soft}}}} potential can be used to un\sphinxhyphen{}overlap atoms while running
dynamics.

\sphinxAtStartPar
Note that you can minimize some atoms in the system while holding the
coordinates of other atoms fixed by applying {\hyperref[\detokenize{fix_setforce::doc}]{\sphinxcrossref{\DUrole{doc}{fix setforce 0.0 0.0
0.0}}}} to the other atoms.  See a more detailed discussion
of {\hyperref[\detokenize{minimize:fix-minimize}]{\sphinxcrossref{\DUrole{std,std-ref}{using fixes while minimizing below}}}}.

\sphinxAtStartPar
The {\hyperref[\detokenize{min_style::doc}]{\sphinxcrossref{\DUrole{doc}{minimization styles}}}} \sphinxstyleemphasis{cg}, \sphinxstyleemphasis{sd}, and \sphinxstyleemphasis{hftn}
involves an outer iteration loop which sets the search direction along
which atom coordinates are changed.  An inner iteration is then
performed using a line search algorithm.  The line search typically
evaluates forces and energies several times to set new coordinates.
Currently, a backtracking algorithm is used which may not be optimal
in terms of the number of force evaluations performed, but appears to
be more robust than previous line searches we have tried.  The
backtracking method is described in Nocedal and Wright’s Numerical
Optimization (Procedure 3.1 on p 41).

\sphinxAtStartPar
The {\hyperref[\detokenize{min_style::doc}]{\sphinxcrossref{\DUrole{doc}{minimization styles}}}} \sphinxstyleemphasis{quickmin}, \sphinxstyleemphasis{fire} and
\sphinxstyleemphasis{fire/old} perform damped dynamics using an Euler integration step.  Thus
they require a {\hyperref[\detokenize{timestep::doc}]{\sphinxcrossref{\DUrole{doc}{timestep}}}} be defined.

\begin{sphinxadmonition}{note}{Note:}
\sphinxAtStartPar
The damped dynamic minimizer algorithms will use the timestep you
have defined via the {\hyperref[\detokenize{timestep::doc}]{\sphinxcrossref{\DUrole{doc}{timestep}}}} command or its
default value.  Often they will converge more quickly if you use a
timestep about 10x larger than you would normally use for regular
molecular dynamics simulations.
\end{sphinxadmonition}


\bigskip\hrule\bigskip


\sphinxAtStartPar
In all cases, the objective function being minimized is the total
potential energy of the system as a function of the N atom
coordinates:
\begin{equation*}
\begin{split}E(r_1,r_2, \ldots ,r_N)  = & \sum_{i,j} E_{\it pair}(r_i,r_j) +
                             \sum_{ij} E_{\it bond}(r_i,r_j) +
                             \sum_{ijk} E_{\it angle}(r_i,r_j,r_k) + \\
                           & \sum_{ijkl} E_{\it dihedral}(r_i,r_j,r_k,r_l) +
                             \sum_{ijkl} E_{\it improper}(r_i,r_j,r_k,r_l) +
                             \sum_i E_{\it fix}(r_i)\end{split}
\end{equation*}
\sphinxAtStartPar
where the first term is the sum of all non\sphinxhyphen{}bonded {\hyperref[\detokenize{pair_style::doc}]{\sphinxcrossref{\DUrole{doc}{pairwise
interactions}}}} including {\hyperref[\detokenize{kspace_style::doc}]{\sphinxcrossref{\DUrole{doc}{long\sphinxhyphen{}range Coulombic
interactions}}}}, the second through fifth terms are {\hyperref[\detokenize{bond_style::doc}]{\sphinxcrossref{\DUrole{doc}{bond}}}}, {\hyperref[\detokenize{angle_style::doc}]{\sphinxcrossref{\DUrole{doc}{angle}}}}, {\hyperref[\detokenize{dihedral_style::doc}]{\sphinxcrossref{\DUrole{doc}{dihedral}}}}, and {\hyperref[\detokenize{improper_style::doc}]{\sphinxcrossref{\DUrole{doc}{improper}}}} interactions
respectively, and the last term is energy due to {\hyperref[\detokenize{fix::doc}]{\sphinxcrossref{\DUrole{doc}{fixes}}}}
which can act as constraints or apply force to atoms, such as through
interaction with a wall.  See the discussion below about how fix
commands affect minimization.

\sphinxAtStartPar
The starting point for the minimization is the current configuration
of the atoms.


\bigskip\hrule\bigskip


\sphinxAtStartPar
The minimization procedure stops if any of several criteria are met:
\begin{itemize}
\item {} 
\sphinxAtStartPar
the change in energy between outer iterations is less than \sphinxstyleemphasis{etol}

\item {} 
\sphinxAtStartPar
the 2\sphinxhyphen{}norm (length) of the global force vector is less than the \sphinxstyleemphasis{ftol}

\item {} 
\sphinxAtStartPar
the line search fails because the step distance backtracks to 0.0

\item {} 
\sphinxAtStartPar
the number of outer iterations or timesteps exceeds \sphinxstyleemphasis{maxiter}

\item {} 
\sphinxAtStartPar
the number of total force evaluations exceeds \sphinxstyleemphasis{maxeval}

\end{itemize}

\begin{sphinxadmonition}{note}{Note:}
\sphinxAtStartPar
the {\hyperref[\detokenize{min_style::doc}]{\sphinxcrossref{\DUrole{doc}{minimization style}}}} \sphinxstyleemphasis{spin},
\sphinxstyleemphasis{spin/cg}, and \sphinxstyleemphasis{spin/lbfgs} replace
the force tolerance \sphinxstyleemphasis{ftol} by a torque tolerance.
The minimization procedure stops if the 2\sphinxhyphen{}norm (length) of the torque vector on atom
(defined as the cross product between the
atomic spin and its precession vectors omega) is less than \sphinxstyleemphasis{ftol},
or if any of the other criteria are met. Torque have the same units as the energy.
\end{sphinxadmonition}

\begin{sphinxadmonition}{note}{Note:}
\sphinxAtStartPar
You can also use the {\hyperref[\detokenize{fix_halt::doc}]{\sphinxcrossref{\DUrole{doc}{fix halt}}}} command to specify
a general criterion for exiting a minimization, that is a
calculation performed on the state of the current system, as
defined by an {\hyperref[\detokenize{variable::doc}]{\sphinxcrossref{\DUrole{doc}{equal\sphinxhyphen{}style variable}}}}.
\end{sphinxadmonition}

\sphinxAtStartPar
For the first criterion, the specified energy tolerance \sphinxstyleemphasis{etol} is
unitless; it is met when the energy change between successive
iterations divided by the energy magnitude is less than or equal to
the tolerance.  For example, a setting of 1.0e\sphinxhyphen{}4 for \sphinxstyleemphasis{etol} means an
energy tolerance of one part in 10\textasciicircum{}4.  For the damped dynamics
minimizers this check is not performed for a few steps after
velocities are reset to 0, otherwise the minimizer would prematurely
converge.

\sphinxAtStartPar
For the second criterion, the specified force tolerance \sphinxstyleemphasis{ftol} is in
force units, since it is the length of the global force vector for all
atoms, e.g. a vector of size 3N for N atoms.  Since many of the
components will be near zero after minimization, you can think of
\sphinxstyleemphasis{ftol} as an upper bound on the final force on any component of any
atom.  For example, a setting of 1.0e\sphinxhyphen{}4 for \sphinxstyleemphasis{ftol} means no x, y, or z
component of force on any atom will be larger than 1.0e\sphinxhyphen{}4 (in force
units) after minimization.

\sphinxAtStartPar
Either or both of the \sphinxstyleemphasis{etol} and \sphinxstyleemphasis{ftol} values can be set to 0.0, in
which case some other criterion will terminate the minimization.

\sphinxAtStartPar
During a minimization, the outer iteration count is treated as a
timestep.  Output is triggered by this timestep, e.g. thermodynamic
output or dump and restart files.

\sphinxAtStartPar
Using the {\hyperref[\detokenize{thermo_style::doc}]{\sphinxcrossref{\DUrole{doc}{thermo\_style custom}}}} command with the
\sphinxstyleemphasis{fmax} or \sphinxstyleemphasis{fnorm} keywords can be useful for monitoring the progress
of the minimization.  Note that these outputs will be calculated only
from forces on the atoms, and will not include any extra degrees of
freedom, such as from the {\hyperref[\detokenize{fix_box_relax::doc}]{\sphinxcrossref{\DUrole{doc}{fix box/relax}}}} command.

\sphinxAtStartPar
Following minimization, a statistical summary is printed that lists
which convergence criterion caused the minimizer to stop, as well as
information about the energy, force, final line search, and iteration
counts.  An example is as follows:

\begin{sphinxVerbatim}[commandchars=\\\{\}]
Minimization stats:
  Stopping criterion = max iterations
  Energy initial, next\PYGZhy{}to\PYGZhy{}last, final =
       \PYGZhy{}0.626828169302     \PYGZhy{}2.82642039062     \PYGZhy{}2.82643549739
  Force two\PYGZhy{}norm initial, final = 2052.1 91.9642
  Force max component initial, final = 346.048 9.78056
  Final line search alpha, max atom move = 2.23899e\PYGZhy{}06 2.18986e\PYGZhy{}05
  Iterations, force evaluations = 2000 12724
\end{sphinxVerbatim}

\sphinxAtStartPar
The 3 energy values are for before and after the minimization and on
the next\sphinxhyphen{}to\sphinxhyphen{}last iteration.  This is what the \sphinxstyleemphasis{etol} parameter checks.

\sphinxAtStartPar
The two\sphinxhyphen{}norm force values are the length of the global force vector
before and after minimization.  This is what the \sphinxstyleemphasis{ftol} parameter
checks.

\sphinxAtStartPar
The max\sphinxhyphen{}component force values are the absolute value of the largest
component (x,y,z) in the global force vector, i.e. the infinity\sphinxhyphen{}norm
of the force vector.

\sphinxAtStartPar
The alpha parameter for the line\sphinxhyphen{}search, when multiplied by the max
force component (on the last iteration), gives the max distance any
atom moved during the last iteration.  Alpha will be 0.0 if the line
search could not reduce the energy.  Even if alpha is non\sphinxhyphen{}zero, if the
“max atom move” distance is tiny compared to typical atom coordinates,
then it is possible the last iteration effectively caused no atom
movement and thus the evaluated energy did not change and the
minimizer terminated.  Said another way, even with non\sphinxhyphen{}zero forces,
it’s possible the effect of those forces is to move atoms a distance
less than machine precision, so that the energy cannot be further
reduced.

\sphinxAtStartPar
The iterations and force evaluation values are what is checked by the
\sphinxstyleemphasis{maxiter} and \sphinxstyleemphasis{maxeval} parameters.


\bigskip\hrule\bigskip


\begin{sphinxadmonition}{note}{Note:}
\sphinxAtStartPar
There are several force fields in LAMMPS which have discontinuities
or other approximations which may prevent you from performing an
energy minimization to tight tolerances.  For example, you should use
a {\hyperref[\detokenize{pair_style::doc}]{\sphinxcrossref{\DUrole{doc}{pair style}}}} that goes to 0.0 at the cutoff
distance when performing minimization (even if you later change it
when running dynamics).  If you do not do this, the total energy of
the system will have discontinuities when the relative distance
between any pair of atoms changes from cutoff \sphinxstyleemphasis{plus} epsilon to
cutoff \sphinxstyleemphasis{minus} epsilon and the minimizer may thus behave poorly.
Some of the many\sphinxhyphen{}body potentials use splines and other internal
cutoffs that inherently have this problem.  The {\hyperref[\detokenize{kspace_style::doc}]{\sphinxcrossref{\DUrole{doc}{long\sphinxhyphen{}range
Coulombic styles}}}} (PPPM, Ewald) are approximate to
within the user\sphinxhyphen{}specified tolerance, which means their energy and
forces may not agree to a higher precision than the Kspace\sphinxhyphen{}specified
tolerance.  This agreement is further reduced when using tabulation
to speed up the computation of the real\sphinxhyphen{}space part of the Coulomb
interactions, which is enabled by default.  In all these cases, the
minimizer may give up and stop before finding a minimum to the
specified energy or force tolerance.
\end{sphinxadmonition}

\sphinxAtStartPar
Note that a cutoff Lennard\sphinxhyphen{}Jones potential (and others) can be shifted
so that its energy is 0.0 at the cutoff via the {\hyperref[\detokenize{pair_modify::doc}]{\sphinxcrossref{\DUrole{doc}{pair\_modify}}}} command.  See the doc pages for individual {\hyperref[\detokenize{pair_style::doc}]{\sphinxcrossref{\DUrole{doc}{pair
styles}}}} for details.  Note that most Coulombic potentials
have a cutoff, unless versions with a long\sphinxhyphen{}range component are used
(e.g. {\hyperref[\detokenize{pair_lj_cut_coul::doc}]{\sphinxcrossref{\DUrole{doc}{pair\_style lj/cut/coul/long}}}}) or some
other damping/smoothing schemes are used.  The CHARMM potentials go to
0.0 at the cutoff (e.g. {\hyperref[\detokenize{pair_charmm::doc}]{\sphinxcrossref{\DUrole{doc}{pair\_style lj/charmm/coul/charmm}}}}), as do the GROMACS potentials (e.g. {\hyperref[\detokenize{pair_gromacs::doc}]{\sphinxcrossref{\DUrole{doc}{pair\_style
lj/gromacs}}}}).

\sphinxAtStartPar
If a soft potential ({\hyperref[\detokenize{pair_soft::doc}]{\sphinxcrossref{\DUrole{doc}{pair\_style soft}}}}) is used the
Astop value is used for the prefactor (no time dependence).

\phantomsection\label{\detokenize{minimize:fix-minimize}}
\sphinxAtStartPar
The {\hyperref[\detokenize{fix_box_relax::doc}]{\sphinxcrossref{\DUrole{doc}{fix box/relax}}}} command can be used to apply an
external pressure to the simulation box and allow it to shrink/expand
during the minimization.

\sphinxAtStartPar
Only a few other fixes (typically those that add forces) are invoked
during minimization.  See the doc pages for individual {\hyperref[\detokenize{fix::doc}]{\sphinxcrossref{\DUrole{doc}{fix}}}}
commands to see which ones are relevant.  Current examples of fixes
that can be used include:
\begin{itemize}
\item {} 
\sphinxAtStartPar
{\hyperref[\detokenize{fix_addforce::doc}]{\sphinxcrossref{\DUrole{doc}{fix addforce}}}}

\item {} 
\sphinxAtStartPar
{\hyperref[\detokenize{fix_addtorque::doc}]{\sphinxcrossref{\DUrole{doc}{fix addtorque}}}}

\item {} 
\sphinxAtStartPar
{\hyperref[\detokenize{fix_efield::doc}]{\sphinxcrossref{\DUrole{doc}{fix efield}}}}

\item {} 
\sphinxAtStartPar
{\hyperref[\detokenize{fix_enforce2d::doc}]{\sphinxcrossref{\DUrole{doc}{fix enforce2d}}}}

\item {} 
\sphinxAtStartPar
{\hyperref[\detokenize{fix_indent::doc}]{\sphinxcrossref{\DUrole{doc}{fix indent}}}}

\item {} 
\sphinxAtStartPar
{\hyperref[\detokenize{fix_lineforce::doc}]{\sphinxcrossref{\DUrole{doc}{fix lineforce}}}}

\item {} 
\sphinxAtStartPar
{\hyperref[\detokenize{fix_planeforce::doc}]{\sphinxcrossref{\DUrole{doc}{fix planeforce}}}}

\item {} 
\sphinxAtStartPar
{\hyperref[\detokenize{fix_setforce::doc}]{\sphinxcrossref{\DUrole{doc}{fix setforce}}}}

\item {} 
\sphinxAtStartPar
{\hyperref[\detokenize{fix_spring::doc}]{\sphinxcrossref{\DUrole{doc}{fix spring}}}}

\item {} 
\sphinxAtStartPar
{\hyperref[\detokenize{fix_spring_self::doc}]{\sphinxcrossref{\DUrole{doc}{fix spring/self}}}}

\item {} 
\sphinxAtStartPar
{\hyperref[\detokenize{fix_viscous::doc}]{\sphinxcrossref{\DUrole{doc}{fix viscous}}}}

\item {} 
\sphinxAtStartPar
{\hyperref[\detokenize{fix_wall::doc}]{\sphinxcrossref{\DUrole{doc}{fix wall}}}}

\item {} 
\sphinxAtStartPar
{\hyperref[\detokenize{fix_wall_region::doc}]{\sphinxcrossref{\DUrole{doc}{fix wall/region}}}}

\end{itemize}

\begin{sphinxadmonition}{note}{Note:}
\sphinxAtStartPar
Some fixes which are invoked during minimization have an
associated potential energy.  For that energy to be included in the
total potential energy of the system (the quantity being minimized),
you MUST enable the {\hyperref[\detokenize{fix_modify::doc}]{\sphinxcrossref{\DUrole{doc}{fix\_modify}}}} \sphinxstyleemphasis{energy} option for
that fix.  The doc pages for individual {\hyperref[\detokenize{fix::doc}]{\sphinxcrossref{\DUrole{doc}{fix}}}} commands
specify if this should be done.
\end{sphinxadmonition}

\begin{sphinxadmonition}{note}{Note:}
\sphinxAtStartPar
The minimizers in LAMMPS do not allow for bonds (or angles, etc)
to be held fixed while atom coordinates are being relaxed, e.g. via
{\hyperref[\detokenize{fix_shake::doc}]{\sphinxcrossref{\DUrole{doc}{fix shake}}}} or {\hyperref[\detokenize{fix_rigid::doc}]{\sphinxcrossref{\DUrole{doc}{fix rigid}}}}.  See more
info in the Restrictions section below.
\end{sphinxadmonition}


\bigskip\hrule\bigskip


\sphinxAtStartPar
Styles with a \sphinxstyleemphasis{gpu}, \sphinxstyleemphasis{intel}, \sphinxstyleemphasis{kk}, \sphinxstyleemphasis{omp}, or \sphinxstyleemphasis{opt} suffix are
functionally the same as the corresponding style without the suffix.
They have been optimized to run faster, depending on your available
hardware, as discussed on the {\hyperref[\detokenize{Speed_packages::doc}]{\sphinxcrossref{\DUrole{doc}{Accelerator packages}}}}
page.  The accelerated styles take the same arguments and should
produce the same results, except for round\sphinxhyphen{}off and precision issues.

\sphinxAtStartPar
These accelerated styles are part of the GPU, INTEL, KOKKOS,
OPENMP, and OPT packages, respectively.  They are only enabled if
LAMMPS was built with those packages.  See the {\hyperref[\detokenize{Build_package::doc}]{\sphinxcrossref{\DUrole{doc}{Build package}}}} page for more info.

\sphinxAtStartPar
You can specify the accelerated styles explicitly in your input script
by including their suffix, or you can use the {\hyperref[\detokenize{Run_options::doc}]{\sphinxcrossref{\DUrole{doc}{\sphinxhyphen{}suffix command\sphinxhyphen{}line switch}}}} when you invoke LAMMPS, or you can use the
{\hyperref[\detokenize{suffix::doc}]{\sphinxcrossref{\DUrole{doc}{suffix}}}} command in your input script.

\sphinxAtStartPar
See the {\hyperref[\detokenize{Speed_packages::doc}]{\sphinxcrossref{\DUrole{doc}{Accelerator packages}}}} page for more
instructions on how to use the accelerated styles effectively.


\bigskip\hrule\bigskip



\subsection{Restrictions}
\label{\detokenize{minimize:restrictions}}
\sphinxAtStartPar
Features that are not yet implemented are listed here, in case someone
knows how they could be coded:

\sphinxAtStartPar
It is an error to use {\hyperref[\detokenize{fix_shake::doc}]{\sphinxcrossref{\DUrole{doc}{fix shake}}}} with minimization
because it turns off bonds that should be included in the potential
energy of the system.  The effect of a fix shake can be approximated
during a minimization by using stiff spring constants for the bonds
and/or angles that would normally be constrained by the SHAKE
algorithm.

\sphinxAtStartPar
{\hyperref[\detokenize{fix_rigid::doc}]{\sphinxcrossref{\DUrole{doc}{Fix rigid}}}} is also not supported by minimization.  It
is not an error to have it defined, but the energy minimization will
not keep the defined body(s) rigid during the minimization.  Note that
if bonds, angles, etc internal to a rigid body have been turned off
(e.g. via {\hyperref[\detokenize{neigh_modify::doc}]{\sphinxcrossref{\DUrole{doc}{neigh\_modify exclude}}}}), they will not
contribute to the potential energy which is probably not what is
desired.

\sphinxAtStartPar
Pair potentials that produce torque on a particle (e.g. {\hyperref[\detokenize{pair_gran::doc}]{\sphinxcrossref{\DUrole{doc}{granular potentials}}}} or the {\hyperref[\detokenize{pair_gayberne::doc}]{\sphinxcrossref{\DUrole{doc}{GayBerne potential}}}} for ellipsoidal particles) are not
relaxed by a minimization.  More specifically, radial relaxations are
induced, but no rotations are induced by a minimization, so such a
system will not fully relax.


\subsection{Related commands}
\label{\detokenize{minimize:related-commands}}
\sphinxAtStartPar
{\hyperref[\detokenize{min_modify::doc}]{\sphinxcrossref{\DUrole{doc}{min\_modify}}}}, {\hyperref[\detokenize{min_style::doc}]{\sphinxcrossref{\DUrole{doc}{min\_style}}}},
{\hyperref[\detokenize{run_style::doc}]{\sphinxcrossref{\DUrole{doc}{run\_style}}}}


\subsection{Default}
\label{\detokenize{minimize:default}}
\sphinxAtStartPar
none

\sphinxstepscope

\index{molecule@\spxentry{molecule}}\ignorespaces 

\section{molecule command}
\label{\detokenize{molecule:molecule-command}}\label{\detokenize{molecule:index-0}}\label{\detokenize{molecule::doc}}

\subsection{Syntax}
\label{\detokenize{molecule:syntax}}
\begin{sphinxVerbatim}[commandchars=\\\{\}]
\PYG{k}{molecule}\PYG{+w}{ }\PYG{n}{ID}\PYG{+w}{ }\PYG{n}{file1}\PYG{+w}{ }\PYG{n}{keyword}\PYG{+w}{ }\PYG{n}{values}\PYG{+w}{ }\PYG{n}{...}\PYG{+w}{ }\PYG{n}{file2}\PYG{+w}{ }\PYG{n}{keyword}\PYG{+w}{ }\PYG{n}{values}\PYG{+w}{ }\PYG{n}{...}\PYG{+w}{ }\PYG{n}{fileN}\PYG{+w}{ }\PYG{n}{...}
\end{sphinxVerbatim}
\begin{itemize}
\item {} 
\sphinxAtStartPar
ID = user\sphinxhyphen{}assigned name for the molecule template

\item {} 
\sphinxAtStartPar
file1,file2,… = names of files containing molecule descriptions

\item {} 
\sphinxAtStartPar
zero or more keyword/value pairs may be appended after each file

\item {} 
\sphinxAtStartPar
keyword = \sphinxstyleemphasis{offset} or \sphinxstyleemphasis{toff} or \sphinxstyleemphasis{boff} or \sphinxstyleemphasis{aoff} or \sphinxstyleemphasis{doff} or \sphinxstyleemphasis{ioff} or \sphinxstyleemphasis{scale}
\begin{sphinxalltt}
\sphinxstyleemphasis{offset} values = Toff Boff Aoff Doff Ioff
  Toff = offset to add to atom types
  Boff = offset to add to bond types
  Aoff = offset to add to angle types
  Doff = offset to add to dihedral types
  Ioff = offset to add to improper types
\sphinxstyleemphasis{toff} value = Toff
  Toff = offset to add to atom types
\sphinxstyleemphasis{boff} value = Boff
  Boff = offset to add to bond types
\sphinxstyleemphasis{aoff} value = Aoff
  Aoff = offset to add to angle types
\sphinxstyleemphasis{doff} value = Doff
  Doff = offset to add to dihedral types
\sphinxstyleemphasis{ioff} value = Ioff
  Ioff = offset to add to improper types
\sphinxstyleemphasis{scale} value = sfactor
  sfactor = scale factor to apply to the size and mass of the molecule
\end{sphinxalltt}

\end{itemize}


\subsection{Examples}
\label{\detokenize{molecule:examples}}
\begin{sphinxVerbatim}[commandchars=\\\{\}]
\PYG{k}{molecule}\PYG{+w}{ }\PYG{l+m}{1}\PYG{+w}{ }\PYG{n}{mymol.txt}
\PYG{k}{molecule}\PYG{+w}{ }\PYG{l+m}{1}\PYG{+w}{ }\PYG{n}{co2.txt}\PYG{+w}{ }\PYG{n}{h2o.txt}
\PYG{k}{molecule}\PYG{+w}{ }\PYG{n}{CO2}\PYG{+w}{ }\PYG{n}{co2.txt}\PYG{+w}{ }\PYG{n}{boff}\PYG{+w}{ }\PYG{l+m}{3}\PYG{+w}{ }\PYG{n}{aoff}\PYG{+w}{ }\PYG{l+m}{2}
\PYG{k}{molecule}\PYG{+w}{ }\PYG{l+m}{1}\PYG{+w}{ }\PYG{n}{mymol.txt}\PYG{+w}{ }\PYG{n}{offset}\PYG{+w}{ }\PYG{l+m}{6}\PYG{+w}{ }\PYG{l+m}{9}\PYG{+w}{ }\PYG{l+m}{18}\PYG{+w}{ }\PYG{l+m}{23}\PYG{+w}{ }\PYG{l+m}{14}
\PYG{k}{molecule}\PYG{+w}{ }\PYG{n}{objects}\PYG{+w}{ }\PYG{n}{file.1}\PYG{+w}{ }\PYG{n}{scale}\PYG{+w}{ }\PYG{l+m}{1.5}\PYG{+w}{ }\PYG{n}{file.1}\PYG{+w}{ }\PYG{n}{scale}\PYG{+w}{ }\PYG{l+m}{2.0}\PYG{+w}{ }\PYG{n}{file.2}\PYG{+w}{ }\PYG{n}{scale}\PYG{+w}{ }\PYG{l+m}{1.3}
\end{sphinxVerbatim}


\subsection{Description}
\label{\detokenize{molecule:description}}
\sphinxAtStartPar
Define a molecule template that can be used as part of other LAMMPS
commands, typically to define a collection of particles as a bonded
molecule or a rigid body.  Commands that currently use molecule
templates include:
\begin{itemize}
\item {} 
\sphinxAtStartPar
{\hyperref[\detokenize{fix_deposit::doc}]{\sphinxcrossref{\DUrole{doc}{fix deposit}}}}

\item {} 
\sphinxAtStartPar
{\hyperref[\detokenize{fix_pour::doc}]{\sphinxcrossref{\DUrole{doc}{fix pour}}}}

\item {} 
\sphinxAtStartPar
{\hyperref[\detokenize{fix_rigid::doc}]{\sphinxcrossref{\DUrole{doc}{fix rigid/small}}}}

\item {} 
\sphinxAtStartPar
{\hyperref[\detokenize{fix_shake::doc}]{\sphinxcrossref{\DUrole{doc}{fix shake}}}}

\item {} 
\sphinxAtStartPar
{\hyperref[\detokenize{fix_gcmc::doc}]{\sphinxcrossref{\DUrole{doc}{fix gcmc}}}}

\item {} 
\sphinxAtStartPar
{\hyperref[\detokenize{fix_bond_react::doc}]{\sphinxcrossref{\DUrole{doc}{fix bond/react}}}}

\item {} 
\sphinxAtStartPar
{\hyperref[\detokenize{create_atoms::doc}]{\sphinxcrossref{\DUrole{doc}{create\_atoms}}}}

\item {} 
\sphinxAtStartPar
{\hyperref[\detokenize{atom_style::doc}]{\sphinxcrossref{\DUrole{doc}{atom\_style template}}}}

\end{itemize}

\sphinxAtStartPar
The ID of a molecule template can only contain alphanumeric characters
and underscores.

\sphinxAtStartPar
A single template can contain multiple molecules, listed one per file.
Some of the commands listed above currently use only the first
molecule in the template, and will issue a warning if the template
contains multiple molecules.  The {\hyperref[\detokenize{atom_style::doc}]{\sphinxcrossref{\DUrole{doc}{atom\_style template}}}} command allows multiple\sphinxhyphen{}molecule templates to define a
system with more than one templated molecule.

\sphinxAtStartPar
Each filename can be followed by optional keywords which are applied
only to the molecule in the file as used in this template.  This is to
make it easy to use the same molecule file in different molecule
templates or in different simulations.  You can specify the same file
multiple times with different optional keywords.

\sphinxAtStartPar
The \sphinxstyleemphasis{offset}, \sphinxstyleemphasis{toff}, \sphinxstyleemphasis{boff}, \sphinxstyleemphasis{aoff}, \sphinxstyleemphasis{doff}, \sphinxstyleemphasis{ioff} keywords
add the specified offset values to the atom types, bond types, angle
types, dihedral types, and/or improper types as they are read from the
molecule file.  E.g. if \sphinxstyleemphasis{toff} = 2, and the file uses atom types
1,2,3, then each created molecule will have atom types 3,4,5.  For the
\sphinxstyleemphasis{offset} keyword, all five offset values must be specified, but
individual values will be ignored if the molecule template does not
use that attribute (e.g. no bonds).

\begin{sphinxadmonition}{note}{Note:}
\sphinxAtStartPar
Offsets are \sphinxstylestrong{ignored} on lines using type labels, as the type
labels will determine the actual types directly depending on the
current {\hyperref[\detokenize{labelmap::doc}]{\sphinxcrossref{\DUrole{doc}{labelmap}}}} settings.
\end{sphinxadmonition}

\sphinxAtStartPar
The \sphinxstyleemphasis{scale} keyword scales the size of the molecule.  This can be
useful for modeling polydisperse granular rigid bodies.  The scale
factor is applied to each of these properties in the molecule file, if
they are defined: the individual particle coordinates (Coords
section), the individual mass of each particle (Masses section), the
individual diameters of each particle (Diameters section), the total
mass of the molecule (header keyword = mass), the center\sphinxhyphen{}of\sphinxhyphen{}mass of
the molecule (header keyword = com), and the moments of inertia of the
molecule (header keyword = inertia).

\begin{sphinxadmonition}{note}{Note:}
\sphinxAtStartPar
The molecule command can be used to define molecules with bonds,
angles, dihedrals, impropers, or special bond lists of neighbors
within a molecular topology, so that you can later add the molecules
to your simulation, via one or more of the commands listed above.
Since this topology\sphinxhyphen{}related information requires that suitable storage
is reserved when LAMMPS creates the simulation box (e.g. when using
the {\hyperref[\detokenize{create_box::doc}]{\sphinxcrossref{\DUrole{doc}{create\_box}}}} command or the
{\hyperref[\detokenize{read_data::doc}]{\sphinxcrossref{\DUrole{doc}{read\_data}}}} command) suitable space has to be reserved
so you do not overflow those pre\sphinxhyphen{}allocated data structures when adding
molecules later.  Both the {\hyperref[\detokenize{create_box::doc}]{\sphinxcrossref{\DUrole{doc}{create\_box}}}} command and
the {\hyperref[\detokenize{read_data::doc}]{\sphinxcrossref{\DUrole{doc}{read\_data}}}} command have “extra” options which
ensure space is allocated for storing topology info for molecules that
are added later.
\end{sphinxadmonition}


\bigskip\hrule\bigskip



\subsection{Format of a molecule file}
\label{\detokenize{molecule:format-of-a-molecule-file}}
\sphinxAtStartPar
The format of an individual molecule file looks similar but is
different than that of a data file read by the {\hyperref[\detokenize{read_data::doc}]{\sphinxcrossref{\DUrole{doc}{read\_data}}}}
commands.  Here is a simple example for a TIP3P water molecule:

\begin{sphinxVerbatim}[commandchars=\\\{\}]
\PYGZsh{} Water molecule. TIP3P geometry
\PYGZsh{} header section:
3 atoms
2 bonds
1 angles

\PYGZsh{} body section:
Coords

1    0.00000  \PYGZhy{}0.06556   0.00000
2    0.75695   0.52032   0.00000
3   \PYGZhy{}0.75695   0.52032   0.00000

Types

1        1   \PYGZsh{} O
2        2   \PYGZsh{} H
3        2   \PYGZsh{} H

Charges

1       \PYGZhy{}0.834
2        0.417
3        0.417

Bonds

1   1      1      2
2   1      1      3

Angles

1   1      2      1      3
\end{sphinxVerbatim}

\sphinxAtStartPar
A molecule file has a header and a body.  The header appears first.  The
first line of the header and thus of the molecule file is \sphinxstyleemphasis{always}
skipped; it typically contains a description of the file or a comment
from the software that created the file.

\sphinxAtStartPar
Then lines are read one line at a time.  Lines can have a trailing
comment starting with ‘\#’ that is ignored.  There \sphinxstyleemphasis{must} be at least one
blank between any valid content and the comment.  If the line is blank
(i.e. contains only white\sphinxhyphen{}space after comments are deleted), it is
skipped.  If the line contains a header keyword, the corresponding
value(s) is/are read from the line.  A line that is \sphinxstyleemphasis{not} blank and does
\sphinxstyleemphasis{not} contains a header keyword begins the body of the file.

\sphinxAtStartPar
The body of the file contains zero or more sections.  The first line
of a section has only a keyword.  The next line is skipped.  The
remaining lines of the section contain values.  The number of lines
depends on the section keyword as described below.  Zero or more blank
lines can be used between sections.  Sections can appear in any order,
with a few exceptions as noted below.

\sphinxAtStartPar
These are the recognized header keywords.  Header lines can come in
any order.  The numeric value(s) are read from the beginning of the
line.  The keyword should appear at the end of the line.  All these
settings have default values, as explained below.  A line need only
appear if the value(s) are different than the default, except when
defining a \sphinxstyleemphasis{body} particle, which requires setting the number of
\sphinxstyleemphasis{atoms} to 1, and setting the \sphinxstyleemphasis{inertia} in a specific section (see below).


\begin{savenotes}\sphinxattablestart
\sphinxthistablewithglobalstyle
\centering
\begin{tabular}[t]{\X{20}{90}\X{13}{90}\X{42}{90}\X{15}{90}}
\sphinxtoprule
\sphinxstyletheadfamily 
\sphinxAtStartPar
Number(s)
&\sphinxstyletheadfamily 
\sphinxAtStartPar
Keyword
&\sphinxstyletheadfamily 
\sphinxAtStartPar
Meaning
&\sphinxstyletheadfamily 
\sphinxAtStartPar
Default Value
\\
\sphinxmidrule
\sphinxtableatstartofbodyhook
\sphinxAtStartPar
N
&
\sphinxAtStartPar
atoms
&
\sphinxAtStartPar
\# of atoms N in molecule
&
\sphinxAtStartPar
0
\\
\sphinxhline
\sphinxAtStartPar
Nb
&
\sphinxAtStartPar
bonds
&
\sphinxAtStartPar
\# of bonds Nb in molecule
&
\sphinxAtStartPar
0
\\
\sphinxhline
\sphinxAtStartPar
Na
&
\sphinxAtStartPar
angles
&
\sphinxAtStartPar
\# of angles Na in molecule
&
\sphinxAtStartPar
0
\\
\sphinxhline
\sphinxAtStartPar
Nd
&
\sphinxAtStartPar
dihedrals
&
\sphinxAtStartPar
\# of dihedrals Nd in molecule
&
\sphinxAtStartPar
0
\\
\sphinxhline
\sphinxAtStartPar
Ni
&
\sphinxAtStartPar
impropers
&
\sphinxAtStartPar
\# of impropers Ni in molecule
&
\sphinxAtStartPar
0
\\
\sphinxhline
\sphinxAtStartPar
Nf
&
\sphinxAtStartPar
fragments
&
\sphinxAtStartPar
\# of fragments Nf in molecule
&
\sphinxAtStartPar
0
\\
\sphinxhline
\sphinxAtStartPar
Ninteger Ndouble
&
\sphinxAtStartPar
body
&
\sphinxAtStartPar
\# of integer and floating\sphinxhyphen{}point values in body particle
&
\sphinxAtStartPar
0
\\
\sphinxhline
\sphinxAtStartPar
Mtotal
&
\sphinxAtStartPar
mass
&
\sphinxAtStartPar
total mass of molecule
&
\sphinxAtStartPar
computed
\\
\sphinxhline
\sphinxAtStartPar
Xc Yc Zc
&
\sphinxAtStartPar
com
&
\sphinxAtStartPar
coordinates of center\sphinxhyphen{}of\sphinxhyphen{}mass of molecule
&
\sphinxAtStartPar
computed
\\
\sphinxhline
\sphinxAtStartPar
Ixx Iyy Izz Ixy Ixz Iyz
&
\sphinxAtStartPar
inertia
&
\sphinxAtStartPar
6 components of inertia tensor of molecule
&
\sphinxAtStartPar
computed
\\
\sphinxbottomrule
\end{tabular}
\sphinxtableafterendhook\par
\sphinxattableend\end{savenotes}

\sphinxAtStartPar
For \sphinxstyleemphasis{mass}, \sphinxstyleemphasis{com}, and \sphinxstyleemphasis{inertia}, the default is for LAMMPS to calculate
this quantity itself if needed, assuming the molecules consist of a set
of point particles or finite\sphinxhyphen{}size particles (with a non\sphinxhyphen{}zero diameter)
that do \sphinxstylestrong{not} overlap.  If finite\sphinxhyphen{}size particles in the molecule
\sphinxstylestrong{do} overlap, LAMMPS will not account for the overlap effects when
calculating any of these 3 quantities, so you should pre\sphinxhyphen{}compute them
yourself and list the values in the file.

\sphinxAtStartPar
The mass and center\sphinxhyphen{}of\sphinxhyphen{}mass coordinates (Xc,Yc,Zc) are
self\sphinxhyphen{}explanatory.  The 6 moments of inertia (ixx,iyy,izz,ixy,ixz,iyz)
should be the values consistent with the current orientation of the
rigid body around its center of mass.  The values are with respect to
the simulation box XYZ axes, not with respect to the principal axes of
the rigid body itself.  LAMMPS performs the latter calculation
internally.

\sphinxAtStartPar
These are the allowed section keywords for the body of the file.
\begin{itemize}
\item {} 
\sphinxAtStartPar
\sphinxstyleemphasis{Coords, Types, Molecules, Fragments, Charges, Diameters, Dipoles, Masses} = atom\sphinxhyphen{}property sections

\item {} 
\sphinxAtStartPar
\sphinxstyleemphasis{Bonds, Angles, Dihedrals, Impropers} = molecular topology sections

\item {} 
\sphinxAtStartPar
\sphinxstyleemphasis{Special Bond Counts, Special Bonds} = special neighbor info

\item {} 
\sphinxAtStartPar
\sphinxstyleemphasis{Shake Flags, Shake Atoms, Shake Bond Types} = SHAKE info

\item {} 
\sphinxAtStartPar
\sphinxstyleemphasis{Body Integers, Body Doubles} = body\sphinxhyphen{}property sections

\end{itemize}

\sphinxAtStartPar
For the Types, Bonds, Angles, Dihedrals, and Impropers sections, each
atom/bond/angle/etc type can be specified either as a number (numeric
type) or as an alphanumeric type label.  The latter is only allowed if
type labels have been defined, either by the {\hyperref[\detokenize{labelmap::doc}]{\sphinxcrossref{\DUrole{doc}{labelmap}}}} command or in data files read by the {\hyperref[\detokenize{read_data::doc}]{\sphinxcrossref{\DUrole{doc}{read\_data}}}} command which have sections for Atom Type Labels, Bond
Type Labels, Angle Type Labels, etc.  See the {\hyperref[\detokenize{Howto_type_labels::doc}]{\sphinxcrossref{\DUrole{doc}{Howto type labels}}}} doc page for the allowed syntax of type labels
and a general discussion of how type labels can be used.
When using type labels, any values specified as \sphinxstyleemphasis{offset} are ignored.

\sphinxAtStartPar
If a Bonds section is specified then the Special Bond Counts and
Special Bonds sections can also be used, if desired, to explicitly
list the 1\sphinxhyphen{}2, 1\sphinxhyphen{}3, 1\sphinxhyphen{}4 neighbors within the molecule topology (see
details below).  This is optional since if these sections are not
included, LAMMPS will auto\sphinxhyphen{}generate this information.  Note that
LAMMPS uses this info to properly exclude or weight bonded pairwise
interactions between bonded atoms.  See the {\hyperref[\detokenize{special_bonds::doc}]{\sphinxcrossref{\DUrole{doc}{special\_bonds}}}} command for more details.  One reason to list the
special bond info explicitly is for the {\hyperref[\detokenize{Howto_drude::doc}]{\sphinxcrossref{\DUrole{doc}{thermalized Drude
oscillator model}}}} which treats the bonds between nuclear
cores and Drude electrons in a different manner.

\begin{sphinxadmonition}{note}{Note:}
\sphinxAtStartPar
Whether a section is required depends on how the molecule template
is used by other LAMMPS commands.  For example, to add a molecule
via the {\hyperref[\detokenize{fix_deposit::doc}]{\sphinxcrossref{\DUrole{doc}{fix deposit}}}} command, the Coords and
Types sections are required.  To add a rigid body via the {\hyperref[\detokenize{fix_pour::doc}]{\sphinxcrossref{\DUrole{doc}{fix
pour}}}} command, the Bonds (Angles, etc) sections are not
required, since the molecule will be treated as a rigid body.  Some
sections are optional.  For example, the {\hyperref[\detokenize{fix_pour::doc}]{\sphinxcrossref{\DUrole{doc}{fix pour}}}}
command can be used to add “molecules” which are clusters of
finite\sphinxhyphen{}size granular particles.  If the Diameters section is not
specified, each particle in the molecule will have a default
diameter of 1.0.  See the doc pages for LAMMPS commands that use
molecule templates for more details.
\end{sphinxadmonition}

\sphinxAtStartPar
Each section is listed below in alphabetic order.  The format of each
section is described including the number of lines it must contain and
rules (if any) for whether it can appear in the data file.  For per\sphinxhyphen{}
atom sections, entries should be numbered from 1 to Natoms (where
Natoms is the number of atoms in the template), indicating which atom
(or bond, etc) the entry applies to.  Per\sphinxhyphen{}atom sections need to
include a setting for every atom, but the atoms can be listed in any
order.


\bigskip\hrule\bigskip


\sphinxAtStartPar
\sphinxstyleemphasis{Coords} section:
\begin{itemize}
\item {} 
\sphinxAtStartPar
one line per atom

\item {} 
\sphinxAtStartPar
line syntax: ID x y z

\item {} 
\sphinxAtStartPar
x,y,z = coordinate of atom

\end{itemize}


\bigskip\hrule\bigskip


\sphinxAtStartPar
\sphinxstyleemphasis{Types} section:
\begin{itemize}
\item {} 
\sphinxAtStartPar
one line per atom

\item {} 
\sphinxAtStartPar
line syntax: ID type

\item {} 
\sphinxAtStartPar
type = atom type of atom (1\sphinxhyphen{}Natomtype, or type label)

\end{itemize}


\bigskip\hrule\bigskip


\sphinxAtStartPar
\sphinxstyleemphasis{Molecules} section:
\begin{itemize}
\item {} 
\sphinxAtStartPar
one line per atom

\item {} 
\sphinxAtStartPar
line syntax: ID molecule\sphinxhyphen{}ID

\item {} 
\sphinxAtStartPar
molecule\sphinxhyphen{}ID = molecule ID of atom

\end{itemize}


\bigskip\hrule\bigskip


\sphinxAtStartPar
\sphinxstyleemphasis{Fragments} section:
\begin{itemize}
\item {} 
\sphinxAtStartPar
one line per fragment

\item {} 
\sphinxAtStartPar
line syntax: ID a b c d …

\item {} 
\sphinxAtStartPar
a,b,c,d,… = IDs of atoms in fragment

\end{itemize}

\sphinxAtStartPar
The ID of a fragment can only contain alphanumeric characters and
underscores.  The atom IDs should be values from 1 to Natoms, where
Natoms = \# of atoms in the molecule.


\bigskip\hrule\bigskip


\sphinxAtStartPar
\sphinxstyleemphasis{Charges} section:
\begin{itemize}
\item {} 
\sphinxAtStartPar
one line per atom

\item {} 
\sphinxAtStartPar
line syntax: ID q

\item {} 
\sphinxAtStartPar
q = charge on atom

\end{itemize}

\sphinxAtStartPar
This section is only allowed for {\hyperref[\detokenize{atom_style::doc}]{\sphinxcrossref{\DUrole{doc}{atom styles}}}} that
support charge.  If this section is not included, the default charge
on each atom in the molecule is 0.0.


\bigskip\hrule\bigskip


\sphinxAtStartPar
\sphinxstyleemphasis{Diameters} section:
\begin{itemize}
\item {} 
\sphinxAtStartPar
one line per atom

\item {} 
\sphinxAtStartPar
line syntax: ID diam

\item {} 
\sphinxAtStartPar
diam = diameter of atom

\end{itemize}

\sphinxAtStartPar
This section is only allowed for {\hyperref[\detokenize{atom_style::doc}]{\sphinxcrossref{\DUrole{doc}{atom styles}}}} that
support finite\sphinxhyphen{}size spherical particles, e.g. atom\_style sphere.  If
not listed, the default diameter of each atom in the molecule is 1.0.


\bigskip\hrule\bigskip


\sphinxAtStartPar
\DUrole{versionmodified,added}{Added in version 7Feb2024.}

\sphinxAtStartPar
\sphinxstyleemphasis{Dipoles} section:
\begin{itemize}
\item {} 
\sphinxAtStartPar
one line per atom

\item {} 
\sphinxAtStartPar
line syntax: ID mux muy muz

\item {} 
\sphinxAtStartPar
mux,muy,muz = x\sphinxhyphen{}, y\sphinxhyphen{}, and z\sphinxhyphen{}component of point dipole vector of atom

\end{itemize}

\sphinxAtStartPar
This section is only allowed for {\hyperref[\detokenize{atom_style::doc}]{\sphinxcrossref{\DUrole{doc}{atom styles}}}} that
support particles with point dipoles, e.g. atom\_style dipole.  If not
listed, the default dipole component of each atom in the molecule is set
to 0.0.


\bigskip\hrule\bigskip


\sphinxAtStartPar
\sphinxstyleemphasis{Masses} section:
\begin{itemize}
\item {} 
\sphinxAtStartPar
one line per atom

\item {} 
\sphinxAtStartPar
line syntax: ID mass

\item {} 
\sphinxAtStartPar
mass = mass of atom

\end{itemize}

\sphinxAtStartPar
This section is only allowed for {\hyperref[\detokenize{atom_style::doc}]{\sphinxcrossref{\DUrole{doc}{atom styles}}}} that
support per\sphinxhyphen{}atom mass, as opposed to per\sphinxhyphen{}type mass.  See the
{\hyperref[\detokenize{mass::doc}]{\sphinxcrossref{\DUrole{doc}{mass}}}} command for details.  If this section is not
included, the default mass for each atom is derived from its volume
(see Diameters section) and a default density of 1.0, in
{\hyperref[\detokenize{units::doc}]{\sphinxcrossref{\DUrole{doc}{units}}}} of mass/volume.


\bigskip\hrule\bigskip


\sphinxAtStartPar
\sphinxstyleemphasis{Bonds} section:
\begin{itemize}
\item {} 
\sphinxAtStartPar
one line per bond

\item {} 
\sphinxAtStartPar
line syntax: ID type atom1 atom2

\item {} 
\sphinxAtStartPar
type = bond type (1\sphinxhyphen{}Nbondtype, or type label)

\item {} 
\sphinxAtStartPar
atom1,atom2 = IDs of atoms in bond

\end{itemize}

\sphinxAtStartPar
The IDs for the two atoms in each bond should be values
from 1 to Natoms, where Natoms = \# of atoms in the molecule.


\bigskip\hrule\bigskip


\sphinxAtStartPar
\sphinxstyleemphasis{Angles} section:
\begin{itemize}
\item {} 
\sphinxAtStartPar
one line per angle

\item {} 
\sphinxAtStartPar
line syntax: ID type atom1 atom2 atom3

\item {} 
\sphinxAtStartPar
type = angle type (1\sphinxhyphen{}Nangletype, or type label)

\item {} 
\sphinxAtStartPar
atom1,atom2,atom3 = IDs of atoms in angle

\end{itemize}

\sphinxAtStartPar
The IDs for the three atoms in each angle should be values from 1 to
Natoms, where Natoms = \# of atoms in the molecule.  The three atoms are
ordered linearly within the angle.  Thus the central atom (around
which the angle is computed) is the atom2 in the list.


\bigskip\hrule\bigskip


\sphinxAtStartPar
\sphinxstyleemphasis{Dihedrals} section:
\begin{itemize}
\item {} 
\sphinxAtStartPar
one line per dihedral

\item {} 
\sphinxAtStartPar
line syntax: ID type atom1 atom2 atom3 atom4

\item {} 
\sphinxAtStartPar
type = dihedral type (1\sphinxhyphen{}Ndihedraltype, or type label)

\item {} 
\sphinxAtStartPar
atom1,atom2,atom3,atom4 = IDs of atoms in dihedral

\end{itemize}

\sphinxAtStartPar
The IDs for the four atoms in each dihedral should be values from 1 to
Natoms, where Natoms = \# of atoms in the molecule.  The 4 atoms are
ordered linearly within the dihedral.


\bigskip\hrule\bigskip


\sphinxAtStartPar
\sphinxstyleemphasis{Impropers} section:
\begin{itemize}
\item {} 
\sphinxAtStartPar
one line per improper

\item {} 
\sphinxAtStartPar
line syntax: ID type atom1 atom2 atom3 atom4

\item {} 
\sphinxAtStartPar
type = improper type (1\sphinxhyphen{}Nimpropertype, or type label)

\item {} 
\sphinxAtStartPar
atom1,atom2,atom3,atom4 = IDs of atoms in improper

\end{itemize}

\sphinxAtStartPar
The IDs for the four atoms in each improper should be values from 1 to
Natoms, where Natoms = \# of atoms in the molecule.  The ordering of
the 4 atoms determines the definition of the improper angle used in
the formula for the defined {\hyperref[\detokenize{improper_style::doc}]{\sphinxcrossref{\DUrole{doc}{improper style}}}}.  See
the doc pages for individual styles for details.


\bigskip\hrule\bigskip


\sphinxAtStartPar
\sphinxstyleemphasis{Special Bond Counts} section:
\begin{itemize}
\item {} 
\sphinxAtStartPar
one line per atom

\item {} 
\sphinxAtStartPar
line syntax: ID N1 N2 N3

\item {} 
\sphinxAtStartPar
N1 = \# of 1\sphinxhyphen{}2 bonds

\item {} 
\sphinxAtStartPar
N2 = \# of 1\sphinxhyphen{}3 bonds

\item {} 
\sphinxAtStartPar
N3 = \# of 1\sphinxhyphen{}4 bonds

\end{itemize}

\sphinxAtStartPar
N1, N2, N3 are the number of 1\sphinxhyphen{}2, 1\sphinxhyphen{}3, 1\sphinxhyphen{}4 neighbors respectively of
this atom within the topology of the molecule.  See the
{\hyperref[\detokenize{special_bonds::doc}]{\sphinxcrossref{\DUrole{doc}{special\_bonds}}}} page for more discussion of
1\sphinxhyphen{}2, 1\sphinxhyphen{}3, 1\sphinxhyphen{}4 neighbors.  If this section appears, the Special Bonds
section must also appear.

\sphinxAtStartPar
As explained above, LAMMPS will auto\sphinxhyphen{}generate this information if this
section is not specified.  If specified, this section will
override what would be auto\sphinxhyphen{}generated.


\bigskip\hrule\bigskip


\sphinxAtStartPar
\sphinxstyleemphasis{Special Bonds} section:
\begin{itemize}
\item {} 
\sphinxAtStartPar
one line per atom

\item {} 
\sphinxAtStartPar
line syntax: ID a b c d …

\item {} 
\sphinxAtStartPar
a,b,c,d,… = IDs of atoms in N1+N2+N3 special bonds

\end{itemize}

\sphinxAtStartPar
A, b, c, d, etc are the IDs of the n1+n2+n3 atoms that are 1\sphinxhyphen{}2, 1\sphinxhyphen{}3,
1\sphinxhyphen{}4 neighbors of this atom.  The IDs should be values from 1 to
Natoms, where Natoms = \# of atoms in the molecule.  The first N1
values should be the 1\sphinxhyphen{}2 neighbors, the next N2 should be the 1\sphinxhyphen{}3
neighbors, the last N3 should be the 1\sphinxhyphen{}4 neighbors.  No atom ID should
appear more than once.  See the {\hyperref[\detokenize{special_bonds::doc}]{\sphinxcrossref{\DUrole{doc}{special\_bonds}}}} doc
page for more discussion of 1\sphinxhyphen{}2, 1\sphinxhyphen{}3, 1\sphinxhyphen{}4 neighbors.  If this section
appears, the Special Bond Counts section must also appear.

\sphinxAtStartPar
As explained above, LAMMPS will auto\sphinxhyphen{}generate this information if this
section is not specified.  If specified, this section will override
what would be auto\sphinxhyphen{}generated.


\bigskip\hrule\bigskip


\sphinxAtStartPar
\sphinxstyleemphasis{Shake Flags} section:
\begin{itemize}
\item {} 
\sphinxAtStartPar
one line per atom

\item {} 
\sphinxAtStartPar
line syntax: ID flag

\item {} 
\sphinxAtStartPar
flag = 0,1,2,3,4

\end{itemize}

\sphinxAtStartPar
This section is only needed when molecules created using the template
will be constrained by SHAKE via the “fix shake” command.  The other
two Shake sections must also appear in the file, following this one.

\sphinxAtStartPar
The meaning of the flag for each atom is as follows.  See the {\hyperref[\detokenize{fix_shake::doc}]{\sphinxcrossref{\DUrole{doc}{fix shake}}}} page for a further description of SHAKE
clusters.
\begin{itemize}
\item {} 
\sphinxAtStartPar
0 = not part of a SHAKE cluster

\item {} 
\sphinxAtStartPar
1 = part of a SHAKE angle cluster (two bonds and the angle they form)

\item {} 
\sphinxAtStartPar
2 = part of a 2\sphinxhyphen{}atom SHAKE cluster with a single bond

\item {} 
\sphinxAtStartPar
3 = part of a 3\sphinxhyphen{}atom SHAKE cluster with two bonds

\item {} 
\sphinxAtStartPar
4 = part of a 4\sphinxhyphen{}atom SHAKE cluster with three bonds

\end{itemize}


\bigskip\hrule\bigskip


\sphinxAtStartPar
\sphinxstyleemphasis{Shake Atoms} section:
\begin{itemize}
\item {} 
\sphinxAtStartPar
one line per atom

\item {} 
\sphinxAtStartPar
line syntax: ID a b c d

\item {} 
\sphinxAtStartPar
a,b,c,d = IDs of atoms in cluster

\end{itemize}

\sphinxAtStartPar
This section is only needed when molecules created using the template
will be constrained by SHAKE via the “fix shake” command.  The other
two Shake sections must also appear in the file.

\sphinxAtStartPar
The a,b,c,d values are atom IDs (from 1 to Natoms) for all the atoms
in the SHAKE cluster that this atom belongs to.  The number of values
that must appear is determined by the shake flag for the atom (see the
Shake Flags section above).  All atoms in a particular cluster should
list their a,b,c,d values identically.

\sphinxAtStartPar
If flag = 0, no a,b,c,d values are listed on the line, just the
(ignored) ID.

\sphinxAtStartPar
If flag = 1, a,b,c are listed, where a = ID of central atom in the
angle, and b,c the other two atoms in the angle.

\sphinxAtStartPar
If flag = 2, a,b are listed, where a = ID of atom in bond with the
lowest ID, and b = ID of atom in bond with the highest ID.

\sphinxAtStartPar
If flag = 3, a,b,c are listed, where a = ID of central atom,
and b,c = IDs of other two atoms bonded to the central atom.

\sphinxAtStartPar
If flag = 4, a,b,c,d are listed, where a = ID of central atom,
and b,c,d = IDs of other three atoms bonded to the central atom.

\sphinxAtStartPar
See the {\hyperref[\detokenize{fix_shake::doc}]{\sphinxcrossref{\DUrole{doc}{fix shake}}}} page for a further description
of SHAKE clusters.


\bigskip\hrule\bigskip


\sphinxAtStartPar
\sphinxstyleemphasis{Shake Bond Types} section:
\begin{itemize}
\item {} 
\sphinxAtStartPar
one line per atom

\item {} 
\sphinxAtStartPar
line syntax: ID a b c

\item {} 
\sphinxAtStartPar
a,b,c = bond types (or angle type) of bonds (or angle) in cluster

\end{itemize}

\sphinxAtStartPar
This section is only needed when molecules created using the template
will be constrained by SHAKE via the “fix shake” command.  The other
two Shake sections must also appear in the file.

\sphinxAtStartPar
The a,b,c values are bond types for all bonds in the SHAKE cluster that
this atom belongs to.  Bond types may be either numbers (from 1 to Nbondtypes)
or bond type labels as defined by the {\hyperref[\detokenize{labelmap::doc}]{\sphinxcrossref{\DUrole{doc}{labelmap}}}} command
or a “Bond Type Labels” section of a data file.

\sphinxAtStartPar
The number of values that must appear is determined by the shake flag
for the atom (see the Shake Flags section above).  All atoms in a
particular cluster should list their a,b,c values identically.

\sphinxAtStartPar
If flag = 0, no a,b,c values are listed on the line, just the
(ignored) ID.

\sphinxAtStartPar
If flag = 1, a,b,c are listed, where a = bondtype of the bond between
the central atom and the first non\sphinxhyphen{}central atom (value b in the Shake
Atoms section), b = bondtype of the bond between the central atom and
the second non\sphinxhyphen{}central atom (value c in the Shake Atoms section), and c
= the angle type (1 to Nangletypes, or angle type label) of the angle
between the three atoms.

\sphinxAtStartPar
If flag = 2, only a is listed, where a = bondtype of the bond between
the two atoms in the cluster.

\sphinxAtStartPar
If flag = 3, a,b are listed, where a = bondtype of the bond between
the central atom and the first non\sphinxhyphen{}central atom (value b in the Shake
Atoms section), and b = bondtype of the bond between the central atom
and the second non\sphinxhyphen{}central atom (value c in the Shake Atoms section).

\sphinxAtStartPar
If flag = 4, a,b,c are listed, where a = bondtype of the bond between
the central atom and the first non\sphinxhyphen{}central atom (value b in the Shake
Atoms section), b = bondtype of the bond between the central atom and
the second non\sphinxhyphen{}central atom (value c in the Shake Atoms section), and c
= bondtype of the bond between the central atom and the third
non\sphinxhyphen{}central atom (value d in the Shake Atoms section).

\sphinxAtStartPar
See the {\hyperref[\detokenize{fix_shake::doc}]{\sphinxcrossref{\DUrole{doc}{fix shake}}}} page for a further description
of SHAKE clusters.


\bigskip\hrule\bigskip


\sphinxAtStartPar
\sphinxstyleemphasis{Body Integers} section:
\begin{itemize}
\item {} 
\sphinxAtStartPar
one line

\item {} 
\sphinxAtStartPar
line syntax: N E F

\item {} 
\sphinxAtStartPar
N = number of sub\sphinxhyphen{}particles or number or vertices

\item {} 
\sphinxAtStartPar
E,F = number of edges and faces

\end{itemize}

\sphinxAtStartPar
This section is only needed when the molecule is a body particle. the other
Body section must also appear in the file.

\sphinxAtStartPar
The total number of values that must appear is determined by the body style, and
must be equal to the Ninteger value given in the \sphinxstyleemphasis{body} header.

\sphinxAtStartPar
For \sphinxstyleemphasis{nparticle} and \sphinxstyleemphasis{rounded/polygon}, only the number of sub\sphinxhyphen{}particles or
vertices N is required, and Ninteger should have a value of 1.

\sphinxAtStartPar
For \sphinxstyleemphasis{rounded/polyhedron}, the number of edges E and faces F is required, and
Ninteger should have a value of 3.

\sphinxAtStartPar
See the {\hyperref[\detokenize{Howto_body::doc}]{\sphinxcrossref{\DUrole{doc}{Howto body}}}} page for a further description of
the file format.


\bigskip\hrule\bigskip


\sphinxAtStartPar
\sphinxstyleemphasis{Body Doubles} section:
\begin{itemize}
\item {} 
\sphinxAtStartPar
first line

\item {} 
\sphinxAtStartPar
line syntax: Ixx Iyy Izz Ixy Ixz Iyz

\item {} 
\sphinxAtStartPar
Ixx Iyy Izz Ixy Ixz Iyz = 6 components of inertia tensor of body particle

\item {} 
\sphinxAtStartPar
one line per sub\sphinxhyphen{}particle or vertex

\item {} 
\sphinxAtStartPar
line syntax: x y z

\item {} 
\sphinxAtStartPar
x, y, z = coordinates of sub\sphinxhyphen{}particle or vertex

\item {} 
\sphinxAtStartPar
one line per edge

\item {} 
\sphinxAtStartPar
line syntax: N1 N2

\item {} 
\sphinxAtStartPar
N1, N2 = vertex indices

\item {} 
\sphinxAtStartPar
one line per face

\item {} 
\sphinxAtStartPar
line syntax: N1 N2 N3 N4

\item {} 
\sphinxAtStartPar
N1, N2, N3, N4 = vertex indices

\item {} 
\sphinxAtStartPar
last line

\item {} 
\sphinxAtStartPar
line syntax: diam

\item {} 
\sphinxAtStartPar
diam = rounded diameter that surrounds each vertex

\end{itemize}

\sphinxAtStartPar
This section is only needed when the molecule is a body particle. the other
Body section must also appear in the file.

\sphinxAtStartPar
The total number of values that must appear is determined by the body style, and
must be equal to the Ndouble value given in the \sphinxstyleemphasis{body} header. The 6 moments of
inertia and the 3N coordinates of the sub\sphinxhyphen{}particles or vertices are required
for all body styles.

\sphinxAtStartPar
For \sphinxstyleemphasis{rounded/polygon}, the E = 6 + 3*N + 1 edges are automatically determined
from the vertices.

\sphinxAtStartPar
For \sphinxstyleemphasis{rounded/polyhedron}, the 2E vertex indices for the end points of the edges
and 4F vertex indices defining the faces are required.

\sphinxAtStartPar
See the {\hyperref[\detokenize{Howto_body::doc}]{\sphinxcrossref{\DUrole{doc}{Howto body}}}} page for a further description of
the file format.


\bigskip\hrule\bigskip



\subsection{Restrictions}
\label{\detokenize{molecule:restrictions}}
\sphinxAtStartPar
None


\subsection{Related commands}
\label{\detokenize{molecule:related-commands}}
\sphinxAtStartPar
{\hyperref[\detokenize{fix_deposit::doc}]{\sphinxcrossref{\DUrole{doc}{fix deposit}}}}, {\hyperref[\detokenize{fix_pour::doc}]{\sphinxcrossref{\DUrole{doc}{fix pour}}}},
{\hyperref[\detokenize{fix_gcmc::doc}]{\sphinxcrossref{\DUrole{doc}{fix gcmc}}}}


\subsection{Default}
\label{\detokenize{molecule:default}}
\sphinxAtStartPar
The default keywords values are offset 0 0 0 0 0 and scale = 1.0.

\sphinxstepscope

\index{neb@\spxentry{neb}}\ignorespaces 

\section{neb command}
\label{\detokenize{neb:neb-command}}\label{\detokenize{neb:index-0}}\label{\detokenize{neb::doc}}

\subsection{Syntax}
\label{\detokenize{neb:syntax}}
\begin{sphinxVerbatim}[commandchars=\\\{\}]
\PYG{k}{neb}\PYG{+w}{ }\PYG{n}{etol}\PYG{+w}{ }\PYG{n}{ftol}\PYG{+w}{ }\PYG{n}{N1}\PYG{+w}{ }\PYG{n}{N2}\PYG{+w}{ }\PYG{n}{Nevery}\PYG{+w}{ }\PYG{n}{file}\PYG{o}{\PYGZhy{}}\PYG{n}{style}\PYG{+w}{ }\PYG{n}{arg}\PYG{+w}{ }\PYG{n}{keyword}\PYG{+w}{ }\PYG{n}{values}
\end{sphinxVerbatim}
\begin{itemize}
\item {} 
\sphinxAtStartPar
etol = stopping tolerance for energy (dimensionless)

\item {} 
\sphinxAtStartPar
ftol = stopping tolerance for force (force units)

\item {} 
\sphinxAtStartPar
N1 = max \# of iterations (timesteps) to run initial NEB

\item {} 
\sphinxAtStartPar
N2 = max \# of iterations (timesteps) to run barrier\sphinxhyphen{}climbing NEB

\item {} 
\sphinxAtStartPar
Nevery = print replica energies and reaction coordinates every this many timesteps

\item {} 
\sphinxAtStartPar
file\sphinxhyphen{}style = \sphinxstyleemphasis{final} or \sphinxstyleemphasis{each} or \sphinxstyleemphasis{none}
\begin{sphinxalltt}
\sphinxstyleemphasis{final} arg = filename
  filename = file with initial coords for final replica
    coords for intermediate replicas are linearly interpolated
    between first and last replica
\sphinxstyleemphasis{each} arg = filename
  filename = unique filename for each replica (except first)
    with its initial coords
\sphinxstyleemphasis{none} arg = no argument all replicas assumed to already have
    their initial coords
\end{sphinxalltt}

\item {} 
\sphinxAtStartPar
zero or more keyword/value pairs may be appended

\item {} 
\sphinxAtStartPar
keyword = \sphinxstyleemphasis{verbosity}
\begin{sphinxalltt}
\sphinxstyleemphasis{verbosity} value = \sphinxstyleemphasis{verbose} or \sphinxstyleemphasis{default} or \sphinxstyleemphasis{terse}
  \sphinxstyleemphasis{verbose} = very detailed per\sphinxhyphen{}replica output
  \sphinxstyleemphasis{default} = some per\sphinxhyphen{}replica output
  \sphinxstyleemphasis{terse} = only global state output
\end{sphinxalltt}

\end{itemize}


\subsection{Examples}
\label{\detokenize{neb:examples}}
\begin{sphinxVerbatim}[commandchars=\\\{\}]
\PYG{k}{neb}\PYG{+w}{ }\PYG{l+m}{0.1}\PYG{+w}{ }\PYG{l+m}{0.0}\PYG{+w}{ }\PYG{l+m}{1000}\PYG{+w}{ }\PYG{l+m}{500}\PYG{+w}{ }\PYG{l+m}{50}\PYG{+w}{ }\PYG{n}{final}\PYG{+w}{ }\PYG{n}{coords.final}
\PYG{k}{neb}\PYG{+w}{ }\PYG{l+m}{0.0}\PYG{+w}{ }\PYG{l+m}{0.001}\PYG{+w}{ }\PYG{l+m}{1000}\PYG{+w}{ }\PYG{l+m}{500}\PYG{+w}{ }\PYG{l+m}{50}\PYG{+w}{ }\PYG{n}{each}\PYG{+w}{ }\PYG{n}{coords.initial.}\PYG{n+nv}{\PYGZdl{}i}
\PYG{k}{neb}\PYG{+w}{ }\PYG{l+m}{0.0}\PYG{+w}{ }\PYG{l+m}{0.001}\PYG{+w}{ }\PYG{l+m}{1000}\PYG{+w}{ }\PYG{l+m}{500}\PYG{+w}{ }\PYG{l+m}{50}\PYG{+w}{ }\PYG{n}{none}\PYG{+w}{ }\PYG{n}{verbose}
\end{sphinxVerbatim}


\subsection{Description}
\label{\detokenize{neb:description}}
\sphinxAtStartPar
Perform a nudged elastic band (NEB) calculation using multiple
replicas of a system.  Two or more replicas must be used; the first
and last are the end points of the transition path.

\sphinxAtStartPar
NEB is a method for finding both the atomic configurations and height of
the energy barrier associated with a transition state, e.g. for an atom
to perform a diffusive hop from one energy basin to another in a
coordinated fashion with its neighbors.  The implementation in LAMMPS
follows the discussion in these 4 papers: {\hyperref[\detokenize{neb:henkelmana}]{\sphinxcrossref{\DUrole{std,std-ref}{(HenkelmanA)}}}}, {\hyperref[\detokenize{neb:henkelmanb}]{\sphinxcrossref{\DUrole{std,std-ref}{(HenkelmanB)}}}}, {\hyperref[\detokenize{neb:nakano3}]{\sphinxcrossref{\DUrole{std,std-ref}{(Nakano)}}}} and {\hyperref[\detokenize{neb:maras2}]{\sphinxcrossref{\DUrole{std,std-ref}{(Maras)}}}}.

\sphinxAtStartPar
Each replica runs on a partition of one or more processors.  Processor
partitions are defined at run\sphinxhyphen{}time using the {\hyperref[\detokenize{Run_options::doc}]{\sphinxcrossref{\DUrole{doc}{\sphinxhyphen{}partition
command\sphinxhyphen{}line switch}}}}.  Note that if you have MPI
installed, you can run a multi\sphinxhyphen{}replica simulation with more replicas
(partitions) than you have physical processors, e.g you can run a
10\sphinxhyphen{}replica simulation on just one or two processors.  You will simply
not get the performance speed\sphinxhyphen{}up you would see with one or more physical
processors per replica.  See the {\hyperref[\detokenize{Howto_replica::doc}]{\sphinxcrossref{\DUrole{doc}{Howto replica}}}}
doc page for further discussion.

\begin{sphinxadmonition}{note}{Note:}
\sphinxAtStartPar
As explained below, a NEB calculation performs a damped dynamics
minimization across all the replicas.  The minimizer uses whatever
timestep you have defined in your input script, via the
{\hyperref[\detokenize{timestep::doc}]{\sphinxcrossref{\DUrole{doc}{timestep}}}} command.  Often NEB will converge more
quickly if you use a timestep about 10x larger than you would normally
use for dynamics simulations.
\end{sphinxadmonition}

\sphinxAtStartPar
When a NEB calculation is performed, it is assumed that each replica
is running the same system, though LAMMPS does not check for this.
I.e. the simulation domain, the number of atoms, the interaction
potentials, and the starting configuration when the neb command is
issued should be the same for every replica.

\sphinxAtStartPar
In a NEB calculation each replica is connected to other replicas by
inter\sphinxhyphen{}replica nudging forces.  These forces are imposed by the
{\hyperref[\detokenize{fix_neb::doc}]{\sphinxcrossref{\DUrole{doc}{fix neb}}}} command, which must be used in conjunction
with the neb command.  The group used to define the fix neb command
defines the NEB atoms which are the only ones that inter\sphinxhyphen{}replica springs
are applied to.  If the group does not include all atoms, then non\sphinxhyphen{}NEB
atoms have no inter\sphinxhyphen{}replica springs and the forces they feel and their
motion is computed in the usual way due only to other atoms within
their replica.  Conceptually, the non\sphinxhyphen{}NEB atoms provide a background
force field for the NEB atoms.  They can be allowed to move during the
NEB minimization procedure (which will typically induce different
coordinates for non\sphinxhyphen{}NEB atoms in different replicas), or held fixed
using other LAMMPS commands such as {\hyperref[\detokenize{fix_setforce::doc}]{\sphinxcrossref{\DUrole{doc}{fix setforce}}}}.
Note that the {\hyperref[\detokenize{partition::doc}]{\sphinxcrossref{\DUrole{doc}{partition}}}} command can be used to invoke
a command on a subset of the replicas, e.g. if you wish to hold NEB or
non\sphinxhyphen{}NEB atoms fixed in only the end\sphinxhyphen{}point replicas.

\sphinxAtStartPar
The initial atomic configuration for each of the replicas can be
specified in different manners via the \sphinxstyleemphasis{file\sphinxhyphen{}style} setting, as
discussed below.  Only atoms whose initial coordinates should differ
from the current configuration need be specified.

\sphinxAtStartPar
Conceptually, the initial and final configurations for the first
replica should be states on either side of an energy barrier.

\sphinxAtStartPar
As explained below, the initial configurations of intermediate
replicas can be atomic coordinates interpolated in a linear fashion
between the first and last replicas.  This is often adequate for
simple transitions.  For more complex transitions, it may lead to slow
convergence or even bad results if the minimum energy path (MEP, see
below) of states over the barrier cannot be correctly converged to
from such an initial path.  In this case, you will want to generate
initial states for the intermediate replicas that are geometrically
closer to the MEP and read them in.


\bigskip\hrule\bigskip


\sphinxAtStartPar
For a \sphinxstyleemphasis{file\sphinxhyphen{}style} setting of \sphinxstyleemphasis{final}, a filename is specified which
contains atomic coordinates for zero or more atoms, in the format
described below.  For each atom that appears in the file, the new
coordinates are assigned to that atom in the final replica.  Each
intermediate replica also assigns a new position to that atom in an
interpolated manner.  This is done by using the current position of
the atom as the starting point and the read\sphinxhyphen{}in position as the final
point.  The distance between them is calculated, and the new position
is assigned to be a fraction of the distance.  E.g. if there are 10
replicas, the second replica will assign a position that is 10\% of the
distance along a line between the starting and final point, and the
9th replica will assign a position that is 90\% of the distance along
the line.  Note that for this procedure to produce consistent
coordinates across all the replicas, the current coordinates need to
be the same in all replicas.  LAMMPS does not check for this, but
invalid initial configurations will likely result if it is not the
case.

\begin{sphinxadmonition}{note}{Note:}
\sphinxAtStartPar
The “distance” between the starting and final point is
calculated in a minimum\sphinxhyphen{}image sense for a periodic simulation box.
This means that if the two positions are on opposite sides of a box
(periodic in that dimension), the distance between them will be small,
because the periodic image of one of the atoms is close to the other.
Similarly, even if the assigned position resulting from the
interpolation is outside the periodic box, the atom will be wrapped
back into the box when the NEB calculation begins.
\end{sphinxadmonition}

\sphinxAtStartPar
For a \sphinxstyleemphasis{file\sphinxhyphen{}style} setting of \sphinxstyleemphasis{each}, a filename is specified which is
assumed to be unique to each replica.  This can be done by using a
variable in the filename, e.g.

\begin{sphinxVerbatim}[commandchars=\\\{\}]
\PYG{k}{variable }\PYG{n+nv+nvIdentifier}{i}\PYG{+w}{ }\PYG{n}{equal}\PYG{+w}{ }\PYG{n}{part}
\PYG{k}{neb}\PYG{+w}{ }\PYG{l+m}{0.0}\PYG{+w}{ }\PYG{l+m}{0.001}\PYG{+w}{ }\PYG{l+m}{1000}\PYG{+w}{ }\PYG{l+m}{500}\PYG{+w}{ }\PYG{l+m}{50}\PYG{+w}{ }\PYG{n}{each}\PYG{+w}{ }\PYG{n}{coords.initial.}\PYG{n+nv}{\PYGZdl{}i}
\end{sphinxVerbatim}

\sphinxAtStartPar
which in this case will substitute the partition ID (0 to N\sphinxhyphen{}1) for the
variable I, which is also effectively the replica ID.  See the
{\hyperref[\detokenize{variable::doc}]{\sphinxcrossref{\DUrole{doc}{variable}}}} command for other options, such as using
world\sphinxhyphen{}, universe\sphinxhyphen{}, or uloop\sphinxhyphen{}style variables.

\sphinxAtStartPar
Each replica (except the first replica) will read its file, formatted
as described below, and for any atom that appears in the file, assign
the specified coordinates to its atom.  The various files do not need
to contain the same set of atoms.

\sphinxAtStartPar
For a \sphinxstyleemphasis{file\sphinxhyphen{}style} setting of \sphinxstyleemphasis{none}, no filename is specified.  Each
replica is assumed to already be in its initial configuration at the
time the neb command is issued.  This allows each replica to define
its own configuration by reading a replica\sphinxhyphen{}specific data or restart or
dump file, via the {\hyperref[\detokenize{read_data::doc}]{\sphinxcrossref{\DUrole{doc}{read\_data}}}},
{\hyperref[\detokenize{read_restart::doc}]{\sphinxcrossref{\DUrole{doc}{read\_restart}}}}, or {\hyperref[\detokenize{read_dump::doc}]{\sphinxcrossref{\DUrole{doc}{read\_dump}}}}
commands.  The replica\sphinxhyphen{}specific names of these files can be specified
as in the discussion above for the \sphinxstyleemphasis{each} file\sphinxhyphen{}style.  Also see the
section below for how a NEB calculation can produce restart files, so
that a long calculation can be restarted if needed.

\begin{sphinxadmonition}{note}{Note:}
\sphinxAtStartPar
None of the \sphinxstyleemphasis{file\sphinxhyphen{}style} settings change the initial
configuration of any atom in the first replica.  The first replica
must thus be in the correct initial configuration at the time the neb
command is issued.
\end{sphinxadmonition}


\bigskip\hrule\bigskip


\sphinxAtStartPar
A NEB calculation proceeds in two stages, each of which is a
minimization procedure, performed via damped dynamics.  To enable
this, you must first define a damped dynamics
{\hyperref[\detokenize{min_style::doc}]{\sphinxcrossref{\DUrole{doc}{min\_style}}}}, such as \sphinxstyleemphasis{quickmin} or \sphinxstyleemphasis{fire}.  The \sphinxstyleemphasis{cg},
\sphinxstyleemphasis{sd}, and \sphinxstyleemphasis{hftn} styles cannot be used, since they perform iterative
line searches in their inner loop, which cannot be easily synchronized
across multiple replicas.

\sphinxAtStartPar
The minimizer tolerances for energy and force are set by \sphinxstyleemphasis{etol} and
\sphinxstyleemphasis{ftol}, the same as for the {\hyperref[\detokenize{minimize::doc}]{\sphinxcrossref{\DUrole{doc}{minimize}}}} command.

\sphinxAtStartPar
A non\sphinxhyphen{}zero \sphinxstyleemphasis{etol} means that the NEB calculation will terminate if the
energy criterion is met by every replica.  The energies being compared
to \sphinxstyleemphasis{etol} do not include any contribution from the inter\sphinxhyphen{}replica
nudging forces, since these are non\sphinxhyphen{}conservative.  A non\sphinxhyphen{}zero \sphinxstyleemphasis{ftol}
means that the NEB calculation will terminate if the force criterion
is met by every replica.  The forces being compared to \sphinxstyleemphasis{ftol} include
the inter\sphinxhyphen{}replica nudging forces.

\sphinxAtStartPar
The maximum number of iterations in each stage is set by \sphinxstyleemphasis{N1} and
\sphinxstyleemphasis{N2}.  These are effectively timestep counts since each iteration of
damped dynamics is like a single timestep in a dynamics
{\hyperref[\detokenize{run::doc}]{\sphinxcrossref{\DUrole{doc}{run}}}}.  During both stages, the potential energy of each
replica and its normalized distance along the reaction path (reaction
coordinate RD) will be printed to the screen and log file every
\sphinxstyleemphasis{Nevery} timesteps.  The RD is 0 and 1 for the first and last replica.
For intermediate replicas, it is the cumulative distance (normalized
by the total cumulative distance) between adjacent replicas, where
“distance” is defined as the length of the 3N\sphinxhyphen{}vector of differences in
atomic coordinates, where N is the number of NEB atoms involved in the
transition.  These outputs allow you to monitor NEB’s progress in
finding a good energy barrier.  \sphinxstyleemphasis{N1} and \sphinxstyleemphasis{N2} must both be multiples
of \sphinxstyleemphasis{Nevery}.

\sphinxAtStartPar
In the first stage of NEB, the set of replicas should converge toward
a minimum energy path (MEP) of conformational states that transition
over a barrier.  The MEP for a transition is defined as a sequence of
3N\sphinxhyphen{}dimensional states, each of which has a potential energy gradient
parallel to the MEP itself.  The configuration of highest energy along
a MEP corresponds to a saddle point.  The replica states will also be
roughly equally spaced along the MEP due to the inter\sphinxhyphen{}replica nudging
force added by the {\hyperref[\detokenize{fix_neb::doc}]{\sphinxcrossref{\DUrole{doc}{fix neb}}}} command.

\sphinxAtStartPar
In the second stage of NEB, the replica with the highest energy is
selected and the inter\sphinxhyphen{}replica forces on it are converted to a force
that drives its atom coordinates to the top or saddle point of the
barrier, via the barrier\sphinxhyphen{}climbing calculation described in
{\hyperref[\detokenize{neb:henkelmanb}]{\sphinxcrossref{\DUrole{std,std-ref}{(HenkelmanB)}}}}.  As before, the other replicas rearrange
themselves along the MEP so as to be roughly equally spaced.

\sphinxAtStartPar
When both stages are complete, if the NEB calculation was successful,
the configurations of the replicas should be along (close to) the MEP
and the replica with the highest energy should be an atomic
configuration at (close to) the saddle point of the transition. The
potential energies for the set of replicas represents the energy
profile of the transition along the MEP.


\bigskip\hrule\bigskip


\sphinxAtStartPar
A few other settings in your input script are required or advised to
perform a NEB calculation.  See the NOTE about the choice of timestep
at the beginning of this doc page.

\sphinxAtStartPar
An atom map must be defined which it is not by default for {\hyperref[\detokenize{atom_style::doc}]{\sphinxcrossref{\DUrole{doc}{atom\_style atomic}}}} problems.  The {\hyperref[\detokenize{atom_modify::doc}]{\sphinxcrossref{\DUrole{doc}{atom\_modify map}}}} command can be used to do this.

\sphinxAtStartPar
The minimizers in LAMMPS operate on all atoms in your system, even
non\sphinxhyphen{}NEB atoms, as defined above.  To prevent non\sphinxhyphen{}NEB atoms from moving
during the minimization, you should use the {\hyperref[\detokenize{fix_setforce::doc}]{\sphinxcrossref{\DUrole{doc}{fix setforce}}}} command to set the force on each of those
atoms to 0.0.  This is not required, and may not even be desired in
some cases, but if those atoms move too far (e.g. because the initial
state of your system was not well\sphinxhyphen{}minimized), it can cause problems
for the NEB procedure.

\sphinxAtStartPar
The damped dynamics {\hyperref[\detokenize{min_style::doc}]{\sphinxcrossref{\DUrole{doc}{minimizers}}}}, such as \sphinxstyleemphasis{quickmin}
and \sphinxstyleemphasis{fire}), adjust the position and velocity of the atoms via an
Euler integration step.  Thus you must define an appropriate
{\hyperref[\detokenize{timestep::doc}]{\sphinxcrossref{\DUrole{doc}{timestep}}}} to use with NEB.  As mentioned above, NEB
will often converge more quickly if you use a timestep about 10x
larger than you would normally use for dynamics simulations.


\bigskip\hrule\bigskip


\sphinxAtStartPar
Each file read by the neb command containing atomic coordinates used
to initialize one or more replicas must be formatted as follows.

\sphinxAtStartPar
The file can be ASCII text or a gzipped text file (detected by a .gz
suffix).  The file can contain initial blank lines or comment lines
starting with “\#” which are ignored.  The first non\sphinxhyphen{}blank, non\sphinxhyphen{}comment
line should list N = the number of lines to follow.  The N successive
lines contain the following information:

\begin{sphinxVerbatim}[commandchars=\\\{\}]
ID1 x1 y1 z1
ID2 x2 y2 z2
...
IDN xN yN zN
\end{sphinxVerbatim}

\sphinxAtStartPar
The fields are the atom ID, followed by the x,y,z coordinates.  The
lines can be listed in any order.  Additional trailing information on
the line is OK, such as a comment.

\sphinxAtStartPar
Note that for a typical NEB calculation you do not need to specify
initial coordinates for very many atoms to produce differing starting
and final replicas whose intermediate replicas will converge to the
energy barrier.  Typically only new coordinates for atoms
geometrically near the barrier need be specified.

\sphinxAtStartPar
Also note there is no requirement that the atoms in the file
correspond to the NEB atoms in the group defined by the {\hyperref[\detokenize{fix_neb::doc}]{\sphinxcrossref{\DUrole{doc}{fix neb}}}} command.  Not every NEB atom need be in the file,
and non\sphinxhyphen{}NEB atoms can be listed in the file.


\bigskip\hrule\bigskip


\sphinxAtStartPar
Four kinds of output can be generated during a NEB calculation: energy
barrier statistics, thermodynamic output by each replica, dump files,
and restart files.

\sphinxAtStartPar
When running with multiple partitions (each of which is a replica in
this case), the print\sphinxhyphen{}out to the screen and master log.lammps file
contains a line of output, printed once every \sphinxstyleemphasis{Nevery} timesteps.  The
amount of information printed in this line can be selected with the
\sphinxstyleemphasis{verbosity} keyword.  Available options are \sphinxstyleemphasis{terse}, \sphinxstyleemphasis{default}, and
\sphinxstyleemphasis{verbose}.

\sphinxAtStartPar
With the \sphinxstyleemphasis{terse} setting, it contains the timestep, the maximum force of
a replica, the maximum force per atom (in any replica), potential
gradients in the initial, final, and climbing replicas, the forward and
backward energy barriers, the total reaction coordinate (RDT).

\sphinxAtStartPar
With the \sphinxstyleemphasis{default} setting, additionally the normalized
reaction coordinate and potential energy of each replica are printed.

\sphinxAtStartPar
With the \sphinxstyleemphasis{verbose} setting, additional per\sphinxhyphen{}replica properties are
printed: the “path angle” (pathangle), the angle between the 3N\sphinxhyphen{}length
tangent vector and the 3N\sphinxhyphen{}length force vector at image \sphinxstyleemphasis{i}
(angletangrad), the angle between the 3N\sphinxhyphen{}length energy gradient vector
of replica \sphinxstyleemphasis{i} and that of replica \sphinxstyleemphasis{i}+1 (anglegrad), the norm of the
energy gradient (gradV), the the two\sphinxhyphen{}norm of the 3N\sphinxhyphen{}length force vector
(RepForce), and the maximum force component of any atom (MaxAtomForce).

\sphinxAtStartPar
The “maximum force per replica” is the two\sphinxhyphen{}norm of the 3N\sphinxhyphen{}length force
vector for the atoms in each replica, maximized across replicas, which
is what the \sphinxstyleemphasis{ftol} setting is checking against.  In this case, N is
all the atoms in each replica.  The “maximum force per atom” is the
maximum force component of any atom in any replica.  The potential
gradients are the two\sphinxhyphen{}norm of the 3N\sphinxhyphen{}length force vector solely due to
the interaction potential i.e.  without adding in inter\sphinxhyphen{}replica
forces.

\sphinxAtStartPar
The “reaction coordinate” (RD) for each replica is the two\sphinxhyphen{}norm of the
3N\sphinxhyphen{}length vector of distances between its atoms and the preceding
replica’s atoms, added to the RD of the preceding replica. The RD of
the first replica RD1 = 0.0; the RD of the final replica RDN = RDT,
the total reaction coordinate.  The normalized RDs are divided by RDT,
so that they form a monotonically increasing sequence from zero to
one. When computing RD, N only includes the atoms being operated on by
the fix neb command.

\sphinxAtStartPar
The forward (reverse) energy barrier is the potential energy of the
highest replica minus the energy of the first (last) replica.

\sphinxAtStartPar
The “path angle” (pathangle) for the replica i which is the angle
between the 3N\sphinxhyphen{}length vectors \((R_{i-1} - R_i)\) and
\((R_{i+1} - R_i)\) (where \(R_i\) is the atomic coordinates of
replica \sphinxstyleemphasis{i}). A “path angle” of 180 indicates that replicas \sphinxstyleemphasis{i}\sphinxhyphen{}1, \sphinxstyleemphasis{i}
and \sphinxstyleemphasis{i}+1 are aligned.  “angletangrad” is the angle between the
3N\sphinxhyphen{}length tangent vector and the 3N\sphinxhyphen{}length force vector at image
\sphinxstyleemphasis{i}. The tangent vector is calculated as in {\hyperref[\detokenize{neb:henkelmana}]{\sphinxcrossref{\DUrole{std,std-ref}{(HenkelmanA)}}}} for all intermediate replicas and at R2 \sphinxhyphen{} R1 and RM \sphinxhyphen{} RM\sphinxhyphen{}1
for the first and last replica, respectively.  “anglegrad” is the angle
between the 3N\sphinxhyphen{}length energy gradient vector of replica \sphinxstyleemphasis{i} and that of
replica \sphinxstyleemphasis{i}+1. It is not defined for the final replica and reads nan.
gradV is the norm of the energy gradient of image \sphinxstyleemphasis{i} (\(\nabla
V\)).  ReplicaForce is the two\sphinxhyphen{}norm of the 3N\sphinxhyphen{}length force vector
(including nudging forces) for replica \sphinxstyleemphasis{i}.  MaxAtomForce is the maximum
force component of any atom in replica \sphinxstyleemphasis{i}.

\sphinxAtStartPar
When a NEB calculation does not converge properly, the supplementary
information can help understanding what is going wrong. For instance
when the path angle becomes acute, the definition of tangent used in
the NEB calculation is questionable and the NEB cannot may diverge
{\hyperref[\detokenize{neb:maras2}]{\sphinxcrossref{\DUrole{std,std-ref}{(Maras)}}}}.

\sphinxAtStartPar
When running on multiple partitions, LAMMPS produces additional log
files for each partition, e.g. log.lammps.0, log.lammps.1, etc.  For a
NEB calculation, these contain the thermodynamic output for each
replica.

\sphinxAtStartPar
If {\hyperref[\detokenize{dump::doc}]{\sphinxcrossref{\DUrole{doc}{dump}}}} commands in the input script define a filename
that includes a \sphinxstyleemphasis{universe} or \sphinxstyleemphasis{uloop} style {\hyperref[\detokenize{variable::doc}]{\sphinxcrossref{\DUrole{doc}{variable}}}},
then one dump file (per dump command) will be created for each
replica.  At the end of the NEB calculation, the final snapshot in
each file will contain the sequence of snapshots that transition the
system over the energy barrier.  Earlier snapshots will show the
convergence of the replicas to the MEP.

\sphinxAtStartPar
Likewise, {\hyperref[\detokenize{restart::doc}]{\sphinxcrossref{\DUrole{doc}{restart}}}} filenames can be specified with a
\sphinxstyleemphasis{universe} or \sphinxstyleemphasis{uloop} style {\hyperref[\detokenize{variable::doc}]{\sphinxcrossref{\DUrole{doc}{variable}}}}, to generate
restart files for each replica.  These may be useful if the NEB
calculation fails to converge properly to the MEP, and you wish to
restart the calculation from an intermediate point with altered
parameters.

\sphinxAtStartPar
There are 2 Python scripts provided in the tools/python directory,
neb\_combine.py and neb\_final.py, which are useful in analyzing output
from a NEB calculation.  Assume a NEB simulation with M replicas, and
the NEB atoms labeled with a specific atom type.

\sphinxAtStartPar
The neb\_combine.py script extracts atom coords for the NEB atoms from
all M dump files and creates a single dump file where each snapshot
contains the NEB atoms from all the replicas and one copy of non\sphinxhyphen{}NEB
atoms from the first replica (presumed to be identical in other
replicas).  This can be visualized/animated to see how the NEB atoms
relax as the NEB calculation proceeds.

\sphinxAtStartPar
The neb\_final.py script extracts the final snapshot from each of the M
dump files to create a single dump file with M snapshots.  This can be
visualized to watch the system make its transition over the energy
barrier.

\sphinxAtStartPar
To illustrate, here are images from the final snapshot produced by the
neb\_combine.py script run on the dump files produced by the two
example input scripts in examples/neb.

\sphinxAtStartPar
\sphinxincludegraphics[width=0.480\linewidth]{{hop1}.jpg}  \sphinxincludegraphics[width=0.480\linewidth]{{hop2}.jpg}




\bigskip\hrule\bigskip



\subsection{Restrictions}
\label{\detokenize{neb:restrictions}}
\sphinxAtStartPar
This command can only be used if LAMMPS was built with the REPLICA
package.  See the {\hyperref[\detokenize{Build_package::doc}]{\sphinxcrossref{\DUrole{doc}{Build package}}}} doc
page for more info.


\bigskip\hrule\bigskip



\subsection{Related commands}
\label{\detokenize{neb:related-commands}}
\sphinxAtStartPar
{\hyperref[\detokenize{prd::doc}]{\sphinxcrossref{\DUrole{doc}{prd}}}}, {\hyperref[\detokenize{temper::doc}]{\sphinxcrossref{\DUrole{doc}{temper}}}}, {\hyperref[\detokenize{fix_langevin::doc}]{\sphinxcrossref{\DUrole{doc}{fix langevin}}}},
{\hyperref[\detokenize{fix_viscous::doc}]{\sphinxcrossref{\DUrole{doc}{fix viscous}}}}, {\hyperref[\detokenize{fix_neb::doc}]{\sphinxcrossref{\DUrole{doc}{fix neb}}}}


\subsection{Default}
\label{\detokenize{neb:default}}
\sphinxAtStartPar
\sphinxstyleemphasis{verbosity} = \sphinxstyleemphasis{default}


\bigskip\hrule\bigskip

\phantomsection\label{\detokenize{neb:henkelmana}}
\sphinxAtStartPar
\sphinxstylestrong{(HenkelmanA)} Henkelman and Jonsson, J Chem Phys, 113, 9978\sphinxhyphen{}9985 (2000).

\phantomsection\label{\detokenize{neb:henkelmanb}}
\sphinxAtStartPar
\sphinxstylestrong{(HenkelmanB)} Henkelman, Uberuaga, Jonsson, J Chem Phys, 113,
9901\sphinxhyphen{}9904 (2000).

\phantomsection\label{\detokenize{neb:nakano3}}
\sphinxAtStartPar
\sphinxstylestrong{(Nakano)} Nakano, Comp Phys Comm, 178, 280\sphinxhyphen{}289 (2008).

\phantomsection\label{\detokenize{neb:maras2}}
\sphinxAtStartPar
\sphinxstylestrong{(Maras)} Maras, Trushin, Stukowski, Ala\sphinxhyphen{}Nissila, Jonsson,
Comp Phys Comm, 205, 13\sphinxhyphen{}21 (2016)

\sphinxstepscope

\index{neb/spin@\spxentry{neb/spin}}\ignorespaces 

\section{neb/spin command}
\label{\detokenize{neb_spin:neb-spin-command}}\label{\detokenize{neb_spin:index-0}}\label{\detokenize{neb_spin::doc}}

\subsection{Syntax}
\label{\detokenize{neb_spin:syntax}}
\begin{sphinxVerbatim}[commandchars=\\\{\}]
\PYG{k}{neb/spin}\PYG{+w}{ }\PYG{n}{etol}\PYG{+w}{ }\PYG{n}{ttol}\PYG{+w}{ }\PYG{n}{N1}\PYG{+w}{ }\PYG{n}{N2}\PYG{+w}{ }\PYG{n}{Nevery}\PYG{+w}{ }\PYG{n}{file}\PYG{o}{\PYGZhy{}}\PYG{n}{style}\PYG{+w}{ }\PYG{n}{arg}\PYG{+w}{ }\PYG{n}{keyword}
\end{sphinxVerbatim}
\begin{itemize}
\item {} 
\sphinxAtStartPar
etol = stopping tolerance for energy (energy units)

\item {} 
\sphinxAtStartPar
ttol = stopping tolerance for torque ( units)

\item {} 
\sphinxAtStartPar
N1 = max \# of iterations (timesteps) to run initial NEB

\item {} 
\sphinxAtStartPar
N2 = max \# of iterations (timesteps) to run barrier\sphinxhyphen{}climbing NEB

\item {} 
\sphinxAtStartPar
Nevery = print replica energies and reaction coordinates every this many timesteps

\item {} 
\sphinxAtStartPar
file\sphinxhyphen{}style = \sphinxstyleemphasis{final} or \sphinxstyleemphasis{each} or \sphinxstyleemphasis{none}
\begin{sphinxalltt}
\sphinxstyleemphasis{final} arg = filename
  filename = file with initial coords for final replica
    coords for intermediate replicas are linearly interpolated
    between first and last replica
\sphinxstyleemphasis{each} arg = filename
  filename = unique filename for each replica (except first)
    with its initial coords
\sphinxstyleemphasis{none} arg = no argument all replicas assumed to already have
    their initial coords
\end{sphinxalltt}

\item {} 
\sphinxAtStartPar
keyword = \sphinxstyleemphasis{verbose}
\begin{sphinxalltt}
\sphinxstyleemphasis{verbose} = print supplemental information
\end{sphinxalltt}

\end{itemize}


\subsection{Examples}
\label{\detokenize{neb_spin:examples}}
\begin{sphinxVerbatim}[commandchars=\\\{\}]
\PYG{k}{neb/spin}\PYG{+w}{ }\PYG{l+m}{0.1}\PYG{+w}{ }\PYG{l+m}{0.0}\PYG{+w}{ }\PYG{l+m}{1000}\PYG{+w}{ }\PYG{l+m}{500}\PYG{+w}{ }\PYG{l+m}{50}\PYG{+w}{ }\PYG{n}{final}\PYG{+w}{ }\PYG{n}{coords.final}
\PYG{k}{neb/spin}\PYG{+w}{ }\PYG{l+m}{0.0}\PYG{+w}{ }\PYG{l+m}{0.001}\PYG{+w}{ }\PYG{l+m}{1000}\PYG{+w}{ }\PYG{l+m}{500}\PYG{+w}{ }\PYG{l+m}{50}\PYG{+w}{ }\PYG{n}{each}\PYG{+w}{ }\PYG{n}{coords.initial.}\PYG{n+nv}{\PYGZdl{}i}
\PYG{k}{neb/spin}\PYG{+w}{ }\PYG{l+m}{0.0}\PYG{+w}{ }\PYG{l+m}{0.001}\PYG{+w}{ }\PYG{l+m}{1000}\PYG{+w}{ }\PYG{l+m}{500}\PYG{+w}{ }\PYG{l+m}{50}\PYG{+w}{ }\PYG{n}{none}\PYG{+w}{ }\PYG{n}{verbose}
\end{sphinxVerbatim}


\subsection{Description}
\label{\detokenize{neb_spin:description}}
\sphinxAtStartPar
Perform a geodesic nudged elastic band (GNEB) calculation using multiple
replicas of a system.  Two or more replicas must be used; the first
and last are the end points of the transition path.

\sphinxAtStartPar
GNEB is a method for finding both the spin configurations and height
of the energy barrier associated with a transition state, e.g.
spins to perform a collective rotation from one energy basin to
another.
The implementation in LAMMPS follows the discussion in the
following paper: {\hyperref[\detokenize{neb_spin:bessaraba}]{\sphinxcrossref{\DUrole{std,std-ref}{(Bessarab)}}}}.

\sphinxAtStartPar
Each replica runs on a partition of one or more processors.  Processor
partitions are defined at run\sphinxhyphen{}time using the {\hyperref[\detokenize{Run_options::doc}]{\sphinxcrossref{\DUrole{doc}{\sphinxhyphen{}partition command\sphinxhyphen{}line switch}}}}.  Note that if you have MPI installed, you
can run a multi\sphinxhyphen{}replica simulation with more replicas (partitions)
than you have physical processors, e.g you can run a 10\sphinxhyphen{}replica
simulation on just one or two processors.  You will simply not get the
performance speed\sphinxhyphen{}up you would see with one or more physical
processors per replica.  See the {\hyperref[\detokenize{Howto_replica::doc}]{\sphinxcrossref{\DUrole{doc}{Howto replica}}}}
doc page for further discussion.

\begin{sphinxadmonition}{note}{Note:}
\sphinxAtStartPar
As explained below, a GNEB calculation performs a
minimization across all the replicas. One of the {\hyperref[\detokenize{min_spin::doc}]{\sphinxcrossref{\DUrole{doc}{spin}}}}
style minimizers has to be defined in your input script.
\end{sphinxadmonition}

\sphinxAtStartPar
When a GNEB calculation is performed, it is assumed that each replica
is running the same system, though LAMMPS does not check for this.
I.e. the simulation domain, the number of magnetic atoms, the
interaction potentials, and the starting configuration when the neb
command is issued should be the same for every replica.

\sphinxAtStartPar
In a GNEB calculation each replica is connected to other replicas by
inter\sphinxhyphen{}replica nudging forces.  These forces are imposed by the {\hyperref[\detokenize{fix_neb_spin::doc}]{\sphinxcrossref{\DUrole{doc}{fix neb/spin}}}} command, which must be used in conjunction
with the neb command.
The group used to define the fix neb/spin command defines the
GNEB magnetic atoms which are the only ones that inter\sphinxhyphen{}replica springs
are applied to.
If the group does not include all magnetic atoms, then non\sphinxhyphen{}GNEB
magnetic atoms have no inter\sphinxhyphen{}replica springs and the torques they feel
and their precession motion is computed in the usual way due only
to other magnetic atoms within their replica.
Conceptually, the non\sphinxhyphen{}GNEB atoms provide a background force field for
the GNEB atoms.
Their magnetic spins can be allowed to evolve during the GNEB
minimization procedure.

\sphinxAtStartPar
The initial spin configuration for each of the replicas can be
specified in different manners via the \sphinxstyleemphasis{file\sphinxhyphen{}style} setting, as
discussed below.  Only atomic spins whose initial coordinates should
differ from the current configuration need to be specified.

\sphinxAtStartPar
Conceptually, the initial and final configurations for the first
replica should be states on either side of an energy barrier.

\sphinxAtStartPar
As explained below, the initial configurations of intermediate
replicas can be spin coordinates interpolated in a linear fashion
between the first and last replicas.  This is often adequate for
simple transitions.  For more complex transitions, it may lead to slow
convergence or even bad results if the minimum energy path (MEP, see
below) of states over the barrier cannot be correctly converged to
from such an initial path.  In this case, you will want to generate
initial states for the intermediate replicas that are geometrically
closer to the MEP and read them in.


\bigskip\hrule\bigskip


\sphinxAtStartPar
For a \sphinxstyleemphasis{file\sphinxhyphen{}style} setting of \sphinxstyleemphasis{final}, a filename is specified which
contains atomic and spin coordinates for zero or more atoms, in the
format described below.
For each atom that appears in the file, the new coordinates are
assigned to that atom in the final replica.  Each intermediate replica
also assigns a new spin to that atom in an interpolated manner.
This is done by using the current direction of the spin at the starting
point and the read\sphinxhyphen{}in direction as the final point.
The “angular distance” between them is calculated, and the new direction
is assigned to be a fraction of the angular distance.

\begin{sphinxadmonition}{note}{Note:}
\sphinxAtStartPar
The “angular distance” between the starting and final point is
evaluated in the geodesic sense, as described in
{\hyperref[\detokenize{neb_spin:bessaraba}]{\sphinxcrossref{\DUrole{std,std-ref}{(Bessarab)}}}}.
\end{sphinxadmonition}

\begin{sphinxadmonition}{note}{Note:}
\sphinxAtStartPar
The angular interpolation between the starting and final point
is achieved using Rodrigues formula:
\end{sphinxadmonition}
\begin{equation*}
\begin{split}\vec{m}_i^{\nu} = \vec{m}_i^{I} \cos(\omega_i^{\nu})
+ (\vec{k}_i \times \vec{m}_i^{I}) \sin(\omega_i^{\nu})
+ (1.0-\cos(\omega_i^{\nu})) \vec{k}_i (\vec{k}_i\cdot
\vec{m}_i^{I})\end{split}
\end{equation*}
\sphinxAtStartPar
where \(\vec{m}_i^I\) is the initial spin configuration for
spin i, \(\omega_i^{\nu}\) is a rotation angle defined as:
\begin{equation*}
\begin{split}\omega_i^{\nu} = (\nu - 1) \Delta \omega_i {\rm ~~and~~} \Delta \omega_i = \frac{\omega_i}{Q-1}\end{split}
\end{equation*}
\sphinxAtStartPar
with \(\nu\) the image number, Q the total number of images, and
\(\omega_i\) the total rotation between the initial and final spins.
\(\vec{k}_i\) defines a rotation axis such as:
\begin{equation*}
\begin{split}\vec{k}_i =  \frac{\vec{m}_i^I \times \vec{m}_i^F}{\left|\vec{m}_i^I \times \vec{m}_i^F\right|}\end{split}
\end{equation*}
\sphinxAtStartPar
if the initial and final spins are not aligned.
If the initial and final spins are aligned, then their cross
product is null, and the expression above does not apply.
If they point toward the same direction, the intermediate images
conserve the same orientation.
If the initial and final spins are aligned, but point toward
opposite directions, an arbitrary rotation vector belonging to
the plane perpendicular to initial and final spins is chosen.
In this case, a warning message is displayed.

\sphinxAtStartPar
For a \sphinxstyleemphasis{file\sphinxhyphen{}style} setting of \sphinxstyleemphasis{each}, a filename is specified which is
assumed to be unique to each replica.
See the {\hyperref[\detokenize{neb::doc}]{\sphinxcrossref{\DUrole{doc}{neb}}}} documentation page for more information about this
option.

\sphinxAtStartPar
For a \sphinxstyleemphasis{file\sphinxhyphen{}style} setting of \sphinxstyleemphasis{none}, no filename is specified.  Each
replica is assumed to already be in its initial configuration at the
time the neb command is issued.  This allows each replica to define
its own configuration by reading a replica\sphinxhyphen{}specific data or restart or
dump file, via the {\hyperref[\detokenize{read_data::doc}]{\sphinxcrossref{\DUrole{doc}{read\_data}}}},
{\hyperref[\detokenize{read_restart::doc}]{\sphinxcrossref{\DUrole{doc}{read\_restart}}}}, or {\hyperref[\detokenize{read_dump::doc}]{\sphinxcrossref{\DUrole{doc}{read\_dump}}}}
commands.  The replica\sphinxhyphen{}specific names of these files can be specified
as in the discussion above for the \sphinxstyleemphasis{each} file\sphinxhyphen{}style.  Also see the
section below for how a NEB calculation can produce restart files, so
that a long calculation can be restarted if needed.

\begin{sphinxadmonition}{note}{Note:}
\sphinxAtStartPar
None of the \sphinxstyleemphasis{file\sphinxhyphen{}style} settings change the initial
configuration of any atom in the first replica.  The first replica
must thus be in the correct initial configuration at the time the neb
command is issued.
\end{sphinxadmonition}


\bigskip\hrule\bigskip


\sphinxAtStartPar
A NEB calculation proceeds in two stages, each of which is a
minimization procedure.  To enable
this, you must first define a
{\hyperref[\detokenize{min_style::doc}]{\sphinxcrossref{\DUrole{doc}{min\_style}}}}, using either the \sphinxstyleemphasis{spin},
\sphinxstyleemphasis{spin/cg}, or \sphinxstyleemphasis{spin/lbfgs} style (see
{\hyperref[\detokenize{min_spin::doc}]{\sphinxcrossref{\DUrole{doc}{min\_spin}}}} for more information).
The other styles cannot be used, since they relax the lattice
degrees of freedom instead of the spins.

\sphinxAtStartPar
The minimizer tolerances for energy and force are set by \sphinxstyleemphasis{etol} and
\sphinxstyleemphasis{ttol}, the same as for the {\hyperref[\detokenize{minimize::doc}]{\sphinxcrossref{\DUrole{doc}{minimize}}}} command.

\sphinxAtStartPar
A non\sphinxhyphen{}zero \sphinxstyleemphasis{etol} means that the GNEB calculation will terminate if the
energy criterion is met by every replica.  The energies being compared
to \sphinxstyleemphasis{etol} do not include any contribution from the inter\sphinxhyphen{}replica
nudging forces, since these are non\sphinxhyphen{}conservative.  A non\sphinxhyphen{}zero \sphinxstyleemphasis{ttol}
means that the GNEB calculation will terminate if the torque criterion
is met by every replica.  The torques being compared to \sphinxstyleemphasis{ttol} include
the inter\sphinxhyphen{}replica nudging forces.

\sphinxAtStartPar
The maximum number of iterations in each stage is set by \sphinxstyleemphasis{N1} and
\sphinxstyleemphasis{N2}.  These are effectively timestep counts since each iteration of
damped dynamics is like a single timestep in a dynamics
{\hyperref[\detokenize{run::doc}]{\sphinxcrossref{\DUrole{doc}{run}}}}.  During both stages, the potential energy of each
replica and its normalized distance along the reaction path (reaction
coordinate RD) will be printed to the screen and log file every
\sphinxstyleemphasis{Nevery} timesteps.  The RD is 0 and 1 for the first and last replica.
For intermediate replicas, it is the cumulative angular distance
(normalized by the total cumulative angular distance) between adjacent
replicas, where “distance” is defined as the length of the 3N\sphinxhyphen{}vector of
the geodesic distances in spin coordinates, with N the number of
GNEB spins involved (see equation (13) in {\hyperref[\detokenize{neb_spin:bessaraba}]{\sphinxcrossref{\DUrole{std,std-ref}{(Bessarab)}}}}).
These outputs allow you to monitor NEB’s progress in
finding a good energy barrier.  \sphinxstyleemphasis{N1} and \sphinxstyleemphasis{N2} must both be multiples
of \sphinxstyleemphasis{Nevery}.

\sphinxAtStartPar
In the first stage of GNEB, the set of replicas should converge toward
a minimum energy path (MEP) of conformational states that transition
over a barrier.  The MEP for a transition is defined as a sequence of
3N\sphinxhyphen{}dimensional spin states, each of which has a potential energy
gradient parallel to the MEP itself.
The configuration of highest energy along a MEP corresponds to a saddle
point.  The replica states will also be roughly equally spaced along
the MEP due to the inter\sphinxhyphen{}replica nudging force added by the
{\hyperref[\detokenize{fix_neb::doc}]{\sphinxcrossref{\DUrole{doc}{fix neb}}}} command.

\sphinxAtStartPar
In the second stage of GNEB, the replica with the highest energy is
selected and the inter\sphinxhyphen{}replica forces on it are converted to a force
that drives its spin coordinates to the top or saddle point of the
barrier, via the barrier\sphinxhyphen{}climbing calculation described in
{\hyperref[\detokenize{neb_spin:bessaraba}]{\sphinxcrossref{\DUrole{std,std-ref}{(Bessarab)}}}}.  As before, the other replicas rearrange
themselves along the MEP so as to be roughly equally spaced.

\sphinxAtStartPar
When both stages are complete, if the GNEB calculation was successful,
the configurations of the replicas should be along (close to) the MEP
and the replica with the highest energy should be a spin
configuration at (close to) the saddle point of the transition. The
potential energies for the set of replicas represents the energy
profile of the transition along the MEP.


\bigskip\hrule\bigskip


\sphinxAtStartPar
An atom map must be defined which it is not by default for {\hyperref[\detokenize{atom_style::doc}]{\sphinxcrossref{\DUrole{doc}{atom\_style atomic}}}} problems.  The {\hyperref[\detokenize{atom_modify::doc}]{\sphinxcrossref{\DUrole{doc}{atom\_modify map}}}} command can be used to do this.

\sphinxAtStartPar
An initial value can be defined for the timestep. Although, the \sphinxstyleemphasis{spin}
minimization algorithm is an adaptive timestep methodology, so that
this timestep is likely to evolve during the calculation.

\sphinxAtStartPar
The minimizers in LAMMPS operate on all spins in your system, even
non\sphinxhyphen{}GNEB atoms, as defined above.


\bigskip\hrule\bigskip


\sphinxAtStartPar
Each file read by the neb/spin command containing spin coordinates used
to initialize one or more replicas must be formatted as follows.

\sphinxAtStartPar
The file can be ASCII text or a gzipped text file (detected by a .gz
suffix).  The file can contain initial blank lines or comment lines
starting with “\#” which are ignored.  The first non\sphinxhyphen{}blank, non\sphinxhyphen{}comment
line should list N = the number of lines to follow.  The N successive
lines contain the following information:

\begin{sphinxVerbatim}[commandchars=\\\{\}]
ID1 g1 x1 y1 z1 sx1 sy1 sz1
ID2 g2 x2 y2 z2 sx2 sy2 sz2
...
IDN gN yN zN sxN syN szN
\end{sphinxVerbatim}

\sphinxAtStartPar
The fields are the atom ID, the norm of the associated magnetic spin,
followed by the \sphinxstyleemphasis{x,y,z} coordinates and the \sphinxstyleemphasis{sx,sy,sz} spin coordinates.
The lines can be listed in any order.  Additional trailing information on
the line is OK, such as a comment.

\sphinxAtStartPar
Note that for a typical GNEB calculation you do not need to specify
initial spin coordinates for very many atoms to produce differing starting
and final replicas whose intermediate replicas will converge to the
energy barrier.  Typically only new spin coordinates for atoms
geometrically near the barrier need be specified.

\sphinxAtStartPar
Also note there is no requirement that the atoms in the file
correspond to the GNEB atoms in the group defined by the {\hyperref[\detokenize{fix_neb::doc}]{\sphinxcrossref{\DUrole{doc}{fix neb}}}} command.  Not every GNEB atom need be in the file,
and non\sphinxhyphen{}GNEB atoms can be listed in the file.


\bigskip\hrule\bigskip


\sphinxAtStartPar
Four kinds of output can be generated during a GNEB calculation: energy
barrier statistics, thermodynamic output by each replica, dump files,
and restart files.

\sphinxAtStartPar
When running with multiple partitions (each of which is a replica in
this case), the print\sphinxhyphen{}out to the screen and master log.lammps file
contains a line of output, printed once every \sphinxstyleemphasis{Nevery} timesteps.  It
contains the timestep, the maximum torque per replica, the maximum
torque per atom (in any replica), potential gradients in the initial,
final, and climbing replicas, the forward and backward energy
barriers, the total reaction coordinate (RDT), and the normalized
reaction coordinate and potential energy of each replica.

\sphinxAtStartPar
The “maximum torque per replica” is the two\sphinxhyphen{}norm of the
3N\sphinxhyphen{}length vector given by the cross product of a spin by its
precession vector omega, in each replica, maximized across replicas,
which is what the \sphinxstyleemphasis{ttol} setting is checking against.  In this case, N is
all the atoms in each replica.  The “maximum torque per atom” is the
maximum torque component of any atom in any replica.  The potential
gradients are the two\sphinxhyphen{}norm of the 3N\sphinxhyphen{}length magnetic precession vector
solely due to the interaction potential i.e. without adding in
inter\sphinxhyphen{}replica forces, and projected along the path tangent (as detailed
in Appendix D of {\hyperref[\detokenize{neb_spin:bessaraba}]{\sphinxcrossref{\DUrole{std,std-ref}{(Bessarab)}}}}).

\sphinxAtStartPar
The “reaction coordinate” (RD) for each replica is the two\sphinxhyphen{}norm of the
3N\sphinxhyphen{}length vector of geodesic distances between its spins and the preceding
replica’s spins (see equation (13) of {\hyperref[\detokenize{neb_spin:bessaraba}]{\sphinxcrossref{\DUrole{std,std-ref}{(Bessarab)}}}}), added to
the RD of the preceding replica. The RD of the first replica RD1 = 0.0;
the RD of the final replica RDN = RDT, the total reaction coordinate.
The normalized RDs are divided by RDT, so that they form a monotonically
increasing sequence from zero to one. When computing RD, N only includes
the spins being operated on by the fix neb/spin command.

\sphinxAtStartPar
The forward (reverse) energy barrier is the potential energy of the
highest replica minus the energy of the first (last) replica.

\sphinxAtStartPar
Supplementary information for all replicas can be printed out to the
screen and master log.lammps file by adding the \sphinxstyleemphasis{verbose} keyword. This
information include the following.
The “GradVidottan” are the projections of the potential gradient for
the replica i on its tangent vector (as detailed in Appendix D of
{\hyperref[\detokenize{neb_spin:bessaraba}]{\sphinxcrossref{\DUrole{std,std-ref}{(Bessarab)}}}}).
The “DNi” are the non normalized geodesic distances (see equation (13)
of {\hyperref[\detokenize{neb_spin:bessaraba}]{\sphinxcrossref{\DUrole{std,std-ref}{(Bessarab)}}}}), between a replica i and the next replica
i+1. For the last replica, this distance is not defined and a “NAN”
value is the corresponding output.

\sphinxAtStartPar
When a NEB calculation does not converge properly, the supplementary
information can help understanding what is going wrong.

\sphinxAtStartPar
When running on multiple partitions, LAMMPS produces additional log
files for each partition, e.g. log.lammps.0, log.lammps.1, etc.  For a
GNEB calculation, these contain the thermodynamic output for each
replica.

\sphinxAtStartPar
If {\hyperref[\detokenize{dump::doc}]{\sphinxcrossref{\DUrole{doc}{dump}}}} commands in the input script define a filename
that includes a \sphinxstyleemphasis{universe} or \sphinxstyleemphasis{uloop} style {\hyperref[\detokenize{variable::doc}]{\sphinxcrossref{\DUrole{doc}{variable}}}},
then one dump file (per dump command) will be created for each
replica.  At the end of the GNEB calculation, the final snapshot in
each file will contain the sequence of snapshots that transition the
system over the energy barrier.  Earlier snapshots will show the
convergence of the replicas to the MEP.

\sphinxAtStartPar
Likewise, {\hyperref[\detokenize{restart::doc}]{\sphinxcrossref{\DUrole{doc}{restart}}}} filenames can be specified with a
\sphinxstyleemphasis{universe} or \sphinxstyleemphasis{uloop} style {\hyperref[\detokenize{variable::doc}]{\sphinxcrossref{\DUrole{doc}{variable}}}}, to generate
restart files for each replica.  These may be useful if the GNEB
calculation fails to converge properly to the MEP, and you wish to
restart the calculation from an intermediate point with altered
parameters.

\sphinxAtStartPar
A c file script in provided in the tool/spin/interpolate\_gneb
directory, that interpolates the MEP given the information provided
by the \sphinxstyleemphasis{verbose} output option (as detailed in Appendix D of
{\hyperref[\detokenize{neb_spin:bessaraba}]{\sphinxcrossref{\DUrole{std,std-ref}{(Bessarab)}}}}).


\bigskip\hrule\bigskip



\subsection{Restrictions}
\label{\detokenize{neb_spin:restrictions}}
\sphinxAtStartPar
This command can only be used if LAMMPS was built with the SPIN
package.  See the {\hyperref[\detokenize{Build_package::doc}]{\sphinxcrossref{\DUrole{doc}{Build package}}}} doc
page for more info.

\sphinxAtStartPar
For magnetic GNEB calculations, only the \sphinxstyleemphasis{spin\_none} value for the
\sphinxstyleemphasis{line} keyword can be used when minimization styles \sphinxstyleemphasis{spin/cg} and
\sphinxstyleemphasis{spin/lbfgs} are employed.


\bigskip\hrule\bigskip



\subsection{Related commands}
\label{\detokenize{neb_spin:related-commands}}
\sphinxAtStartPar
{\hyperref[\detokenize{min_spin::doc}]{\sphinxcrossref{\DUrole{doc}{min/spin}}}}, {\hyperref[\detokenize{fix_neb_spin::doc}]{\sphinxcrossref{\DUrole{doc}{fix neb/spin}}}}


\subsection{Default}
\label{\detokenize{neb_spin:default}}
\sphinxAtStartPar
none


\bigskip\hrule\bigskip

\phantomsection\label{\detokenize{neb_spin:bessaraba}}
\sphinxAtStartPar
\sphinxstylestrong{(Bessarab)} Bessarab, Uzdin, Jonsson, Comp Phys Comm, 196,
335\sphinxhyphen{}347 (2015).

\sphinxstepscope

\index{neigh\_modify@\spxentry{neigh\_modify}}\ignorespaces 

\section{neigh\_modify command}
\label{\detokenize{neigh_modify:neigh-modify-command}}\label{\detokenize{neigh_modify:index-0}}\label{\detokenize{neigh_modify::doc}}

\subsection{Syntax}
\label{\detokenize{neigh_modify:syntax}}
\begin{sphinxVerbatim}[commandchars=\\\{\}]
\PYG{k}{neigh\PYGZus{}modify}\PYG{+w}{ }\PYG{n}{keyword}\PYG{+w}{ }\PYG{n}{values}\PYG{+w}{ }\PYG{n}{...}
\end{sphinxVerbatim}
\begin{itemize}
\item {} 
\sphinxAtStartPar
one or more keyword/value pairs may be listed
\begin{sphinxalltt}
keyword = \sphinxstyleemphasis{delay} or \sphinxstyleemphasis{every} or \sphinxstyleemphasis{check} or \sphinxstyleemphasis{once} or \sphinxstyleemphasis{cluster} or \sphinxstyleemphasis{include} or \sphinxstyleemphasis{exclude} or \sphinxstyleemphasis{page} or \sphinxstyleemphasis{one} or \sphinxstyleemphasis{binsize} or \sphinxstyleemphasis{collection/type} or \sphinxstyleemphasis{collection/interval}
  \sphinxstyleemphasis{delay} value = N
    N = delay building neighbor lists until this many steps since last build
  \sphinxstyleemphasis{every} value = M
    M = consider building neighbor lists every this many steps
  \sphinxstyleemphasis{check} value = \sphinxstyleemphasis{yes} or \sphinxstyleemphasis{no}
    \sphinxstyleemphasis{yes} = only build if at least one atom has moved half the skin distance or more
    \sphinxstyleemphasis{no} = always build on 1st step where \sphinxstyleemphasis{every} and \sphinxstyleemphasis{delay} are conditions are satisfied
  \sphinxstyleemphasis{once} value = \sphinxstyleemphasis{yes} or \sphinxstyleemphasis{no}
    \sphinxstyleemphasis{yes} = only build neighbor list once at start of run and never rebuild
    \sphinxstyleemphasis{no} = rebuild neighbor list according to other settings
  \sphinxstyleemphasis{cluster} value = \sphinxstyleemphasis{yes} or \sphinxstyleemphasis{no}
    \sphinxstyleemphasis{yes} = check bond,angle,etc neighbor list for nearby clusters
    \sphinxstyleemphasis{no} = do not check bond,angle,etc neighbor list for nearby clusters
  \sphinxstyleemphasis{include} value = group\sphinxhyphen{}ID
    group\sphinxhyphen{}ID = only build pair neighbor lists for atoms in this group
  \sphinxstyleemphasis{exclude} values:
    \sphinxstyleemphasis{type} M N
      M,N = exclude if one atom in pair is type M, other is type N (M and N may be type labels)
    \sphinxstyleemphasis{group} group1\sphinxhyphen{}ID group2\sphinxhyphen{}ID
      group1\sphinxhyphen{}ID,group2\sphinxhyphen{}ID = exclude if one atom is in 1st group, other in 2nd
    \sphinxstyleemphasis{molecule/intra} group\sphinxhyphen{}ID
      group\sphinxhyphen{}ID = exclude if both atoms are in the same molecule and in group
    \sphinxstyleemphasis{molecule/inter} group\sphinxhyphen{}ID
      group\sphinxhyphen{}ID = exclude if both atoms are in different molecules and in group
    \sphinxstyleemphasis{none}
      delete all exclude settings
  \sphinxstyleemphasis{page} value = N
    N = number of pairs stored in a single neighbor page
  \sphinxstyleemphasis{one} value = N
    N = max number of neighbors of one atom
  \sphinxstyleemphasis{binsize} value = size
    size = bin size for neighbor list construction (distance units)
  \sphinxstyleemphasis{collection/type} values = N arg1 ... argN
    N = number of custom collections
    arg = N separate lists of types (see below)
  \sphinxstyleemphasis{collection/interval} values = N arg1 ... argN
    N = number of custom collections
    arg = N separate cutoffs for intervals (see below)
\end{sphinxalltt}

\end{itemize}


\subsection{Examples}
\label{\detokenize{neigh_modify:examples}}
\begin{sphinxVerbatim}[commandchars=\\\{\}]
\PYG{k}{neigh\PYGZus{}modify}\PYG{+w}{ }\PYG{n}{every}\PYG{+w}{ }\PYG{l+m}{2}\PYG{+w}{ }\PYG{n}{delay}\PYG{+w}{ }\PYG{l+m}{10}\PYG{+w}{ }\PYG{n}{check}\PYG{+w}{ }\PYG{n}{yes}\PYG{+w}{ }\PYG{n}{page}\PYG{+w}{ }\PYG{l+m}{100000}
\PYG{k}{neigh\PYGZus{}modify}\PYG{+w}{ }\PYG{n}{exclude}\PYG{+w}{ }\PYG{n}{type}\PYG{+w}{ }\PYG{l+m}{2}\PYG{+w}{ }\PYG{l+m}{3}
\PYG{k}{neigh\PYGZus{}modify}\PYG{+w}{ }\PYG{n}{exclude}\PYG{+w}{ }\PYG{k}{group }\PYG{n+nv+nvIdentifier}{frozen}\PYG{+w}{ }\PYG{n}{frozen}\PYG{+w}{ }\PYG{n}{check}\PYG{+w}{ }\PYG{n}{no}
\PYG{k}{neigh\PYGZus{}modify}\PYG{+w}{ }\PYG{n}{exclude}\PYG{+w}{ }\PYG{k}{group }\PYG{n+nv+nvIdentifier}{residue1}\PYG{+w}{ }\PYG{n}{chain3}
\PYG{k}{neigh\PYGZus{}modify}\PYG{+w}{ }\PYG{n}{exclude}\PYG{+w}{ }\PYG{n}{molecule}\PYG{o}{/}\PYG{n}{intra}\PYG{+w}{ }\PYG{n}{rigid}
\PYG{k}{neigh\PYGZus{}modify}\PYG{+w}{ }\PYG{n}{collection}\PYG{o}{/}\PYG{n}{type}\PYG{+w}{ }\PYG{l+m}{2}\PYG{+w}{ }\PYG{l+m}{1}\PYG{o}{*}\PYG{l+m}{2},5\PYG{+w}{ }\PYG{l+m}{3}\PYG{o}{*}\PYG{l+m}{4}
\PYG{k}{neigh\PYGZus{}modify}\PYG{+w}{ }\PYG{n}{collection}\PYG{o}{/}\PYG{n}{interval}\PYG{+w}{ }\PYG{l+m}{2}\PYG{+w}{ }\PYG{l+m}{1.0}\PYG{+w}{ }\PYG{l+m}{10.0}
\end{sphinxVerbatim}


\subsection{Description}
\label{\detokenize{neigh_modify:description}}
\sphinxAtStartPar
This command sets parameters that affect the building and use of
pairwise neighbor lists.  Depending on what pair interactions and other
commands are defined, a simulation may require one or more neighbor
lists.

\sphinxAtStartPar
The \sphinxstyleemphasis{every}, \sphinxstyleemphasis{delay}, \sphinxstyleemphasis{check}, and \sphinxstyleemphasis{once} options affect how often lists
are built as a simulation runs.  The \sphinxstyleemphasis{delay} setting means never build
new lists until at least N steps after the previous build.  The \sphinxstyleemphasis{every}
setting means attempt to build lists every M steps (after the delay has
passed).  If the \sphinxstyleemphasis{check} setting is \sphinxstyleemphasis{no}, the lists are built on the
first step that satisfies the \sphinxstyleemphasis{delay} and \sphinxstyleemphasis{every} settings.  If the
\sphinxstyleemphasis{check} setting is \sphinxstyleemphasis{yes}, then the \sphinxstyleemphasis{every} and \sphinxstyleemphasis{delay} settings
determine when a build may possibly be performed, but an actual build
only occurs if at least one atom has moved more than half the neighbor
skin distance (specified in the {\hyperref[\detokenize{neighbor::doc}]{\sphinxcrossref{\DUrole{doc}{neighbor}}}} command)
since the last neighbor list build.

\begin{sphinxadmonition}{note}{Impact of neighbor list settings}

\sphinxAtStartPar
The choice of neighbor list settings can have a significant impact on
the (parallel) performance of LAMMPS and the correctness of the
simulation results.  Since building the neighbor lists is time
consuming, doing it less frequently can speed up a calculation.  If
the lists are rebuilt too infrequently, however, interacting pairs
may be missing and thus the resulting pairwise interactions
incorrect.  The optimal settings depend on many factors like the
properties of the simulated system (density, geometry, topology,
temperature, pressure), the force field parameters and settings, the
size of the timestep, neighbor list skin distance and more.  The
default settings are chosen to be very conservative to guarantee
correctness of the simulation.  They depend on the \sphinxstyleemphasis{check} flag
heuristics to reduce the number of neighbor list rebuilds at a minor
expense for executing the check.  Determining the correctness of a
specific choice of neighbor list settings is complicated by the fact
that a neighbor list rebuild changes the order in which pairwise
interactions are computed and thus
\sphinxhyphen{} due to the limitations of floating\sphinxhyphen{}point math \sphinxhyphen{} the trajectory.
\end{sphinxadmonition}

\sphinxAtStartPar
If the \sphinxstyleemphasis{once} setting is yes, then the neighbor list is only built once
at the beginning of each run, and never rebuilt, except on steps when a
restart file is written, or steps when a fix forces a rebuild to occur
(e.g. fixes that create or delete atoms, such as {\hyperref[\detokenize{fix_deposit::doc}]{\sphinxcrossref{\DUrole{doc}{fix deposit}}}} or {\hyperref[\detokenize{fix_evaporate::doc}]{\sphinxcrossref{\DUrole{doc}{fix evaporate}}}}).  This setting
should only be made if you are certain atoms will not move far enough
that the neighbor list should be rebuilt, e.g. running a simulation of a
cold crystal.  Note that it is not that expensive to check if neighbor
lists should be rebuilt.

\sphinxAtStartPar
When the rRESPA integrator is used (see the {\hyperref[\detokenize{run_style::doc}]{\sphinxcrossref{\DUrole{doc}{run\_style}}}}
command), the \sphinxstyleemphasis{every} and \sphinxstyleemphasis{delay} parameters refer to the longest
(outermost) timestep.

\sphinxAtStartPar
The \sphinxstyleemphasis{cluster} option does a sanity test every time neighbor lists are
built for bond, angle, dihedral, and improper interactions, to check
that each set of 2, 3, or 4 atoms is a cluster of nearby atoms.  It
does this by computing the distance between pairs of atoms in the
interaction and ensuring they are not further apart than half the
periodic box length.  If they are, an error is generated, since the
interaction would be computed between far\sphinxhyphen{}away atoms instead of their
nearby periodic images.  The only way this should happen is if the
pairwise cutoff is so short that atoms that are part of the same
interaction are not communicated as ghost atoms.  This is an unusual
model (e.g. no pair interactions at all) and the problem can be fixed
by use of the {\hyperref[\detokenize{comm_modify::doc}]{\sphinxcrossref{\DUrole{doc}{comm\_modify cutoff}}}} command.  Note
that to save time, the default \sphinxstyleemphasis{cluster} setting is \sphinxstyleemphasis{no}, so that this
check is not performed.

\sphinxAtStartPar
The \sphinxstyleemphasis{include} option limits the building of pairwise neighbor lists to
atoms in the specified group.  This can be useful for models where a
large portion of the simulation is particles that do not interact with
other particles or with each other via pairwise interactions.  The
group specified with this option must also be specified via the
{\hyperref[\detokenize{atom_modify::doc}]{\sphinxcrossref{\DUrole{doc}{atom\_modify first}}}} command.  Note that specifying
“all” as the group\sphinxhyphen{}ID effectively turns off the \sphinxstyleemphasis{include} option.

\sphinxAtStartPar
The \sphinxstyleemphasis{exclude} option turns off pairwise interactions between certain
pairs of atoms, by not including them in the neighbor list.  These are
sample scenarios where this is useful:
\begin{itemize}
\item {} 
\sphinxAtStartPar
In crack simulations, pairwise interactions can be shut off between 2
slabs of atoms to effectively create a crack.

\item {} 
\sphinxAtStartPar
When a large collection of atoms is treated as frozen, interactions
between those atoms can be turned off to save needless
computation. E.g. Using the {\hyperref[\detokenize{fix_setforce::doc}]{\sphinxcrossref{\DUrole{doc}{fix setforce}}}} command
to freeze a wall or portion of a bio\sphinxhyphen{}molecule.

\item {} 
\sphinxAtStartPar
When one or more rigid bodies are specified, interactions within each
body can be turned off to save needless computation.  See the {\hyperref[\detokenize{fix_rigid::doc}]{\sphinxcrossref{\DUrole{doc}{fix rigid}}}} command for more details.

\end{itemize}

\sphinxAtStartPar
\DUrole{versionmodified,changed}{Changed in version 29Aug2024: }Support for type labels was added.

\sphinxAtStartPar
The \sphinxstyleemphasis{exclude type} option turns off the pairwise interaction if one atom
is of type M and the other of type N.  M can equal N.  The \sphinxstyleemphasis{exclude
group} option turns off the interaction if one atom is in the first
group and the other is the second.  Group1\sphinxhyphen{}ID can equal group2\sphinxhyphen{}ID.  The
\sphinxstyleemphasis{exclude molecule/intra} option turns off the interaction if both atoms
are in the specified group and in the same molecule, as determined by
their molecule ID.  The \sphinxstyleemphasis{exclude molecule/inter} turns off the
interaction between pairs of atoms that have different molecule IDs and
are both in the specified group.

\sphinxAtStartPar
Each of the exclude options can be specified multiple times.  The
\sphinxstyleemphasis{exclude type} option is the most efficient option to use; it requires
only a single check, no matter how many times it has been specified.
The other exclude options are more expensive if specified multiple
times; they require one check for each time they have been specified.

\sphinxAtStartPar
Note that the exclude options only affect pairwise interactions; see
the {\hyperref[\detokenize{delete_bonds::doc}]{\sphinxcrossref{\DUrole{doc}{delete\_bonds}}}} command for information on
turning off bond interactions.

\begin{sphinxadmonition}{note}{Note:}
\sphinxAtStartPar
Excluding pairwise interactions will not work correctly when
also using a long\sphinxhyphen{}range solver via the
{\hyperref[\detokenize{kspace_style::doc}]{\sphinxcrossref{\DUrole{doc}{kspace\_style}}}} command.  LAMMPS will give a warning
to this effect.  This is because the short\sphinxhyphen{}range pairwise interaction
needs to subtract off a term from the total energy for pairs whose
short\sphinxhyphen{}range interaction is excluded, to compensate for how the
long\sphinxhyphen{}range solver treats the interaction.  This is done correctly for
pairwise interactions that are excluded (or weighted) via the
{\hyperref[\detokenize{special_bonds::doc}]{\sphinxcrossref{\DUrole{doc}{special\_bonds}}}} command.  But it is not done for
interactions that are excluded via these neigh\_modify exclude options.
\end{sphinxadmonition}

\sphinxAtStartPar
The \sphinxstyleemphasis{page} and \sphinxstyleemphasis{one} options affect how memory is allocated for the
neighbor lists.  For most simulations the default settings for these
options are fine, but if a very large problem is being run or a very
long cutoff is being used, these parameters can be tuned.  The indices
of neighboring atoms are stored in “pages”, which are allocated one
after another as they fill up.  The size of each page is set by the
\sphinxstyleemphasis{page} value.  A new page is allocated when the next atom’s neighbors
could potentially overflow the list.  This threshold is set by the
\sphinxstyleemphasis{one} value which tells LAMMPS the maximum number of neighbor’s one
atom can have.

\begin{sphinxadmonition}{note}{Note:}
\sphinxAtStartPar
LAMMPS can crash without an error message if the number of
neighbors for a single particle is larger than the \sphinxstyleemphasis{page} setting,
which means it is much, much larger than the \sphinxstyleemphasis{one} setting.  This is
because LAMMPS does not error check these limits for every pairwise
interaction (too costly), but only after all the particle’s neighbors
have been found.  This problem usually means something is very wrong
with the way you have setup your problem (particle spacing, cutoff
length, neighbor skin distance, etc).  If you really expect that many
neighbors per particle, then boost the \sphinxstyleemphasis{one} and \sphinxstyleemphasis{page} settings
accordingly.
\end{sphinxadmonition}

\sphinxAtStartPar
The \sphinxstyleemphasis{binsize} option allows you to specify what size of bins will be
used in neighbor list construction to sort and find neighboring atoms.
By default, for {\hyperref[\detokenize{neighbor::doc}]{\sphinxcrossref{\DUrole{doc}{neighbor style bin}}}}, LAMMPS uses bins
that are 1/2 the size of the maximum pair cutoff.  For {\hyperref[\detokenize{neighbor::doc}]{\sphinxcrossref{\DUrole{doc}{neighbor
style multi}}}}, the bins are 1/2 the size of the collection
interaction cutoff.  Typically these are good values for minimizing the
time for neighbor list construction.  This setting overrides the
default.  If you make it too big, there is little overhead due to
looping over bins, but more atoms are checked.  If you make it too
small, the optimal number of atoms is checked, but bin overhead goes up.
If you set the binsize to 0.0, LAMMPS will use the default binsize of
1/2 the cutoff.

\sphinxAtStartPar
The \sphinxstyleemphasis{collection/type} option allows you to define collections of atom
types, used by the \sphinxstyleemphasis{multi} neighbor mode.  By grouping atom types with
similar physical size or interaction cutoff lengths, one may be able to
improve performance by reducing overhead.  You must first specify the
number of collections N to be defined followed by N lists of types.
Each list consists of a series of type ranges separated by commas. The
range can be specified as a single numeric value, or a wildcard asterisk
can be used to specify a range of values.  This takes the form “*” or
“*n” or “n*” or “m*n”.  For example, if M = the number of atom types,
then an asterisk with no numeric values means all types from 1 to M.  A
leading asterisk means all types from 1 to n (inclusive).  A trailing
asterisk means all types from n to M (inclusive).  A middle asterisk
means all types from m to n (inclusive).  Note that all atom types must
be included in exactly one of the N collections.

\sphinxAtStartPar
The \sphinxstyleemphasis{collection/interval} option provides a similar capability.  This
command allows a user to define collections by specifying a series of
cutoff intervals.  LAMMPS will automatically sort atoms into these
intervals based on their type\sphinxhyphen{}dependent cutoffs or their finite size.
You must first specify the number of collections N to be defined
followed by N values representing the upper cutoff of each interval.
This command is particularly useful for granular pair styles where the
interaction distance of particles depends on their radius and may not
depend on their atom type.


\subsection{Restrictions}
\label{\detokenize{neigh_modify:restrictions}}
\sphinxAtStartPar
If the \sphinxstyleemphasis{delay} setting is non\sphinxhyphen{}zero, then it must be a multiple of the
\sphinxstyleemphasis{every} setting.

\sphinxAtStartPar
The \sphinxstyleemphasis{molecule/intra} and \sphinxstyleemphasis{molecule/inter} exclusion options can only
be used with atom styles that define molecule IDs.

\sphinxAtStartPar
The value of the \sphinxstyleemphasis{page} setting must be at least 10x larger than the
\sphinxstyleemphasis{one} setting.  This ensures neighbor pages are not mostly empty
space.

\sphinxAtStartPar
The \sphinxstyleemphasis{exclude group} setting is currently not compatible with dynamic groups.


\subsection{Related commands}
\label{\detokenize{neigh_modify:related-commands}}
\sphinxAtStartPar
{\hyperref[\detokenize{neighbor::doc}]{\sphinxcrossref{\DUrole{doc}{neighbor}}}}, {\hyperref[\detokenize{delete_bonds::doc}]{\sphinxcrossref{\DUrole{doc}{delete\_bonds}}}}


\subsection{Default}
\label{\detokenize{neigh_modify:default}}
\sphinxAtStartPar
The option defaults are delay = 0, every = 1, check = yes, once = no,
cluster = no, include = all (same as no include option defined),
exclude = none, page = 100000, one = 2000, and binsize = 0.0.

\sphinxstepscope

\index{neighbor@\spxentry{neighbor}}\ignorespaces 

\section{neighbor command}
\label{\detokenize{neighbor:neighbor-command}}\label{\detokenize{neighbor:index-0}}\label{\detokenize{neighbor::doc}}

\subsection{Syntax}
\label{\detokenize{neighbor:syntax}}
\begin{sphinxVerbatim}[commandchars=\\\{\}]
\PYG{k}{neighbor}\PYG{+w}{ }\PYG{n}{skin}\PYG{+w}{ }\PYG{n}{style}
\end{sphinxVerbatim}
\begin{itemize}
\item {} 
\sphinxAtStartPar
skin = extra distance beyond force cutoff (distance units)

\item {} 
\sphinxAtStartPar
style = \sphinxstyleemphasis{bin} or \sphinxstyleemphasis{nsq} or \sphinxstyleemphasis{multi} or \sphinxstyleemphasis{multi/old}

\end{itemize}


\subsection{Examples}
\label{\detokenize{neighbor:examples}}
\begin{sphinxVerbatim}[commandchars=\\\{\}]
\PYG{k}{neighbor}\PYG{+w}{ }\PYG{l+m}{0.3}\PYG{+w}{ }\PYG{n}{bin}
\PYG{k}{neighbor}\PYG{+w}{ }\PYG{l+m}{2.0}\PYG{+w}{ }\PYG{n}{nsq}
\end{sphinxVerbatim}


\subsection{Description}
\label{\detokenize{neighbor:description}}
\sphinxAtStartPar
This command sets parameters that affect the building of pairwise
neighbor lists.  All atom pairs within a neighbor cutoff distance
equal to the their force cutoff plus the \sphinxstyleemphasis{skin} distance are stored in
the list.  Typically, the larger the skin distance, the less often
neighbor lists need to be built, but more pairs must be checked for
possible force interactions every timestep.  The default value for
\sphinxstyleemphasis{skin} depends on the choice of units for the simulation; see the
default values below.

\sphinxAtStartPar
The \sphinxstyleemphasis{skin} distance is also used to determine how often atoms migrate
to new processors if the \sphinxstyleemphasis{check} option of the
{\hyperref[\detokenize{neigh_modify::doc}]{\sphinxcrossref{\DUrole{doc}{neigh\_modify}}}} command is set to \sphinxstyleemphasis{yes}.  Atoms are
migrated (communicated) to new processors on the same timestep that
neighbor lists are re\sphinxhyphen{}built.

\sphinxAtStartPar
The \sphinxstyleemphasis{style} value selects what algorithm is used to build the list.
The \sphinxstyleemphasis{bin} style creates the list by binning which is an operation that
scales linearly with N/P, the number of atoms per processor where N =
total number of atoms and P = number of processors.  It is almost
always faster than the \sphinxstyleemphasis{nsq} style which scales as (N/P)\textasciicircum{}2.  For
unsolvated small molecules in a non\sphinxhyphen{}periodic box, the \sphinxstyleemphasis{nsq} choice can
sometimes be faster.  Either style should give the same answers.

\sphinxAtStartPar
The \sphinxstyleemphasis{multi} style is a modified binning algorithm that is useful for
systems with a wide range of cutoff distances, e.g. due to different
size particles. For granular pair styles, cutoffs are set to the sum of
the maximum atomic radii for each atom type.  For the \sphinxstyleemphasis{bin} style, the
bin size is set to 1/2 of the largest cutoff distance between any pair
of atom types and a single set of bins is defined to search over for all
atom types.  This can be inefficient if one pair of types has a very
long cutoff, but other type pairs have a much shorter cutoff. The
\sphinxstyleemphasis{multi} style uses different sized bins for collections of different
sized particles, where “size” may mean the physical size of the particle
or its cutoff distance for interacting with other particles. Different
sets of bins are then used to construct the neighbor lists as as further
described by Shire, Hanley, and Stratford {\hyperref[\detokenize{neighbor:multi-shire}]{\sphinxcrossref{\DUrole{std,std-ref}{(Shire)}}}}
and Monti et al. {\hyperref[\detokenize{neighbor:multi-monti}]{\sphinxcrossref{\DUrole{std,std-ref}{(Monti)}}}}. This imposes some extra
setup overhead, but the searches themselves may be much faster.

\sphinxAtStartPar
For instance in a dense binary system in d\sphinxhyphen{}dimensions with a ratio of the size
of the largest to smallest collection bin \(\lambda\), the computational
costs of building a default neighbor list grows as \(\lambda^{2d}\) while
the costs for \sphinxstyleemphasis{multi} grows as \(\lambda^d\), equivalent to the cost
of force evaluations, as argued in Monti et al. {\hyperref[\detokenize{neighbor:multi-monti}]{\sphinxcrossref{\DUrole{std,std-ref}{(Monti)}}}}.
In other words, the neighboring costs of \sphinxstyleemphasis{multi} are expected to scale the
same as force calculations, such that its relative cost is independent of
the particle size ratio. This is not the case for the default style which
becomes substantially more expensive with increasing size ratios.

\sphinxAtStartPar
By default in \sphinxstyleemphasis{multi}, each atom type defines a separate collection
of particles. For systems where two or more atom types have the same
size (either physical size or cutoff distance), the definition of
collections can be customized, which can result in less overhead and
faster performance. See the {\hyperref[\detokenize{neigh_modify::doc}]{\sphinxcrossref{\DUrole{doc}{neigh\_modify}}}} command
for how to define custom collections. Whether the collection definition
is customized or not, also see the {\hyperref[\detokenize{comm_modify::doc}]{\sphinxcrossref{\DUrole{doc}{comm\_modify mode multi}}}} command for communication options that further improve
performance in a manner consistent with neighbor style multi.

\sphinxAtStartPar
An alternate style, \sphinxstyleemphasis{multi/old}, sets the bin size to 1/2 of the shortest
cutoff distance and multiple sets of bins are defined to search over for
different atom types. This algorithm used to be the default \sphinxstyleemphasis{multi}
algorithm in LAMMPS but was found to be significantly slower than the new
approach. For the dense binary system, computational costs still grew as
\(\lambda^{2d}\) at large enough \(\lambda\). This is equivalent
to the default style, albeit with a smaller prefactor. For now we are
keeping the old option in case there are use cases where multi/old
outperforms the new multi style.

\begin{sphinxadmonition}{note}{Note:}
\sphinxAtStartPar
If there are multiple sub\sphinxhyphen{}styles in a {\hyperref[\detokenize{pair_hybrid::doc}]{\sphinxcrossref{\DUrole{doc}{hybrid/overlay pair style}}}} that cover the same atom types, but have significantly
different cutoffs, the \sphinxstyleemphasis{multi} style does not apply.  Instead, the
{\hyperref[\detokenize{pair_modify::doc}]{\sphinxcrossref{\DUrole{doc}{pair\_modify neigh/trim}}}} setting applies (which is
\sphinxstyleemphasis{yes} by default).  Please check the neighbor list summary printed at
the beginning of a calculation to verify that the desired set of
neighbor list builds is performed.
\end{sphinxadmonition}

\sphinxAtStartPar
The {\hyperref[\detokenize{neigh_modify::doc}]{\sphinxcrossref{\DUrole{doc}{neigh\_modify}}}} command has additional options
that control how often neighbor lists are built and which pairs are
stored in the list.

\sphinxAtStartPar
When a run is finished, counts of the number of neighbors stored in
the pairwise list and the number of times neighbor lists were built
are printed to the screen and log file.  See the {\hyperref[\detokenize{Run_output::doc}]{\sphinxcrossref{\DUrole{doc}{Run output}}}} page for details.


\subsection{Restrictions}
\label{\detokenize{neighbor:restrictions}}\begin{quote}

\sphinxAtStartPar
none
\end{quote}


\subsection{Related commands}
\label{\detokenize{neighbor:related-commands}}
\sphinxAtStartPar
{\hyperref[\detokenize{neigh_modify::doc}]{\sphinxcrossref{\DUrole{doc}{neigh\_modify}}}}, {\hyperref[\detokenize{units::doc}]{\sphinxcrossref{\DUrole{doc}{units}}}},
{\hyperref[\detokenize{comm_modify::doc}]{\sphinxcrossref{\DUrole{doc}{comm\_modify}}}}


\subsection{Default}
\label{\detokenize{neighbor:default}}
\begin{DUlineblock}{0em}
\item[] 0.3 bin for units = lj, skin = 0.3 sigma
\item[] 2.0 bin for units = real or metal, skin = 2.0 Angstroms
\item[] 0.001 bin for units = si, skin = 0.001 meters = 1.0 mm
\item[] 0.1 bin for units = cgs, skin = 0.1 cm = 1.0 mm
\item[] 
\end{DUlineblock}


\bigskip\hrule\bigskip

\phantomsection\label{\detokenize{neighbor:multi-shire}}
\sphinxAtStartPar
\sphinxstylestrong{(Shire)} Shire, Hanley and Stratford, Comp. Part. Mech., (2020).

\phantomsection\label{\detokenize{neighbor:multi-monti}}
\sphinxAtStartPar
\sphinxstylestrong{(Monti)} Monti, Clemmer, Srivastava, Silbert, Grest, and Lechman, Phys. Rev. E, (2022).

\sphinxstepscope

\index{newton@\spxentry{newton}}\ignorespaces 

\section{newton command}
\label{\detokenize{newton:newton-command}}\label{\detokenize{newton:index-0}}\label{\detokenize{newton::doc}}

\subsection{Syntax}
\label{\detokenize{newton:syntax}}
\begin{sphinxVerbatim}[commandchars=\\\{\}]
\PYG{k}{newton}\PYG{+w}{ }\PYG{n}{flag}
\PYG{k}{newton}\PYG{+w}{ }\PYG{n}{flag1}\PYG{+w}{ }\PYG{n}{flag2}
\end{sphinxVerbatim}
\begin{itemize}
\item {} 
\sphinxAtStartPar
flag = \sphinxstyleemphasis{on} or \sphinxstyleemphasis{off} for both pairwise and bonded interactions

\item {} 
\sphinxAtStartPar
flag1 = \sphinxstyleemphasis{on} or \sphinxstyleemphasis{off} for pairwise interactions

\item {} 
\sphinxAtStartPar
flag2 = \sphinxstyleemphasis{on} or \sphinxstyleemphasis{off} for bonded interactions

\end{itemize}


\subsection{Examples}
\label{\detokenize{newton:examples}}
\begin{sphinxVerbatim}[commandchars=\\\{\}]
\PYG{k}{newton}\PYG{+w}{ }\PYG{n}{off}
\PYG{k}{newton}\PYG{+w}{ }\PYG{n}{on}\PYG{+w}{ }\PYG{n}{off}
\end{sphinxVerbatim}


\subsection{Description}
\label{\detokenize{newton:description}}
\sphinxAtStartPar
This command turns Newton’s third law \sphinxstyleemphasis{on} or \sphinxstyleemphasis{off} for pairwise and
bonded interactions.  For most problems, setting Newton’s third law to
\sphinxstyleemphasis{on} means a modest savings in computation at the cost of two times
more communication.  Whether this is faster depends on problem size,
force cutoff lengths, a machine’s compute/communication ratio, and how
many processors are being used.

\sphinxAtStartPar
Setting the pairwise newton flag to \sphinxstyleemphasis{off} means that if two
interacting atoms are on different processors, both processors compute
their interaction and the resulting force information is not
communicated.  Similarly, for bonded interactions, newton \sphinxstyleemphasis{off} means
that if a bond, angle, dihedral, or improper interaction contains
atoms on 2 or more processors, the interaction is computed by each
processor.

\sphinxAtStartPar
LAMMPS should produce the same answers for any newton flag settings,
except for round\sphinxhyphen{}off issues.

\sphinxAtStartPar
With {\hyperref[\detokenize{run_style::doc}]{\sphinxcrossref{\DUrole{doc}{run\_style}}}} \sphinxstyleemphasis{respa} and only bonded interactions
(bond, angle, etc) computed in the innermost timestep, it may be
faster to turn newton \sphinxstyleemphasis{off} for bonded interactions, to avoid extra
communication in the innermost loop.


\subsection{Restrictions}
\label{\detokenize{newton:restrictions}}
\sphinxAtStartPar
The newton bond setting cannot be changed after the simulation box is
defined by a {\hyperref[\detokenize{read_data::doc}]{\sphinxcrossref{\DUrole{doc}{read\_data}}}} or
{\hyperref[\detokenize{create_box::doc}]{\sphinxcrossref{\DUrole{doc}{create\_box}}}} command.


\subsection{Related commands}
\label{\detokenize{newton:related-commands}}
\sphinxAtStartPar
{\hyperref[\detokenize{run_style::doc}]{\sphinxcrossref{\DUrole{doc}{run\_style}}}} respa


\subsection{Default}
\label{\detokenize{newton:default}}
\begin{sphinxVerbatim}[commandchars=\\\{\}]
\PYG{k}{newton}\PYG{+w}{ }\PYG{n}{on}
\end{sphinxVerbatim}

\sphinxstepscope

\index{next@\spxentry{next}}\ignorespaces 

\section{next command}
\label{\detokenize{next:next-command}}\label{\detokenize{next:index-0}}\label{\detokenize{next::doc}}

\subsection{Syntax}
\label{\detokenize{next:syntax}}
\begin{sphinxVerbatim}[commandchars=\\\{\}]
\PYG{k}{next }\PYG{n+nv+nvIdentifier}{variables}
\end{sphinxVerbatim}
\begin{itemize}
\item {} 
\sphinxAtStartPar
variables = one or more variable names

\end{itemize}


\subsection{Examples}
\label{\detokenize{next:examples}}
\begin{sphinxVerbatim}[commandchars=\\\{\}]
\PYG{k}{next }\PYG{n+nv+nvIdentifier}{x}
\PYG{k}{next }\PYG{n+nv+nvIdentifier}{a}\PYG{+w}{ }\PYG{n}{t}\PYG{+w}{ }\PYG{n}{x}\PYG{+w}{ }\PYG{n}{myTemp}
\end{sphinxVerbatim}


\subsection{Description}
\label{\detokenize{next:description}}
\sphinxAtStartPar
This command is used with variables defined by the
{\hyperref[\detokenize{variable::doc}]{\sphinxcrossref{\DUrole{doc}{variable}}}} command.  It assigns the next value to the
variable from the list of values defined for that variable by the
{\hyperref[\detokenize{variable::doc}]{\sphinxcrossref{\DUrole{doc}{variable}}}} command.  Thus when that variable is
subsequently substituted for in an input script command, the new value
is used.

\sphinxAtStartPar
See the {\hyperref[\detokenize{variable::doc}]{\sphinxcrossref{\DUrole{doc}{variable}}}} command for info on how to define and
use different kinds of variables in LAMMPS input scripts.  If a
variable name is a single lower\sphinxhyphen{}case character from “a” to “z”, it can
be used in an input script command as \$a or \$z.  If it is multiple
letters, it can be used as \$\{myTemp\}.

\sphinxAtStartPar
If multiple variables are used as arguments to the \sphinxstyleemphasis{next} command,
then all must be of the same variable style: \sphinxstyleemphasis{index}, \sphinxstyleemphasis{loop}, \sphinxstyleemphasis{file},
\sphinxstyleemphasis{universe}, or \sphinxstyleemphasis{uloop}.  An exception is that \sphinxstyleemphasis{universe}\sphinxhyphen{} and
\sphinxstyleemphasis{uloop}\sphinxhyphen{}style variables can be mixed in the same \sphinxstyleemphasis{next} command.

\sphinxAtStartPar
All the variables specified with the next command are incremented by
one value from their respective list of values.  A \sphinxstyleemphasis{file}\sphinxhyphen{}style
variable reads the next line from its associated file.  An
\sphinxstyleemphasis{atomfile}\sphinxhyphen{}style variable reads the next set of lines (one per atom)
from its associated file.  \sphinxstyleemphasis{String\sphinxhyphen{}} or \sphinxstyleemphasis{atom}\sphinxhyphen{} or \sphinxstyleemphasis{equal}\sphinxhyphen{} or
\sphinxstyleemphasis{world}\sphinxhyphen{}style variables cannot be used with the next command,
since they only store a single value.

\sphinxAtStartPar
When any of the variables in the next command has no more values, a
flag is set that causes the input script to skip the next
{\hyperref[\detokenize{jump::doc}]{\sphinxcrossref{\DUrole{doc}{jump}}}} command encountered.  This enables a loop containing
a next command to exit.  As explained in the {\hyperref[\detokenize{variable::doc}]{\sphinxcrossref{\DUrole{doc}{variable}}}}
command, the variable that has exhausted its values is also deleted.
This allows it to be used and re\sphinxhyphen{}defined later in the input script.
\sphinxstyleemphasis{File}\sphinxhyphen{}style and \sphinxstyleemphasis{atomfile}\sphinxhyphen{}style variables are exhausted when the
end\sphinxhyphen{}of\sphinxhyphen{}file is reached.

\sphinxAtStartPar
When the next command is used with \sphinxstyleemphasis{index}\sphinxhyphen{} or \sphinxstyleemphasis{loop}\sphinxhyphen{}style variables,
the next value is assigned to the variable for all processors.  When
the next command is used with \sphinxstyleemphasis{file}\sphinxhyphen{}style variables, the next line is
read from its file and the string assigned to the variable.  When the
next command is used with \sphinxstyleemphasis{atomfile}\sphinxhyphen{}style variables, the next set of
per\sphinxhyphen{}atom values is read from its file and assigned to the variable.

\sphinxAtStartPar
When the next command is used with \sphinxstyleemphasis{universe}\sphinxhyphen{} or \sphinxstyleemphasis{uloop}\sphinxhyphen{}style
variables, all \sphinxstyleemphasis{universe}\sphinxhyphen{} or \sphinxstyleemphasis{uloop}\sphinxhyphen{}style variables must be listed
in the next command.  This is because of the manner in which the
incrementing is done, using a single lock file for all variables.  The
next value (for each variable) is assigned to whichever processor
partition executes the command first.  All processors in the partition
are assigned the same value(s).  Running LAMMPS on multiple partitions
of processors via the {\hyperref[\detokenize{Run_options::doc}]{\sphinxcrossref{\DUrole{doc}{\sphinxhyphen{}partition command\sphinxhyphen{}line switch}}}}.  \sphinxstyleemphasis{Universe}\sphinxhyphen{} and \sphinxstyleemphasis{uloop}\sphinxhyphen{}style variables are
incremented using the files “tmp.lammps.variable” and
“tmp.lammps.variable.lock” which you will see in your directory during
and after such a LAMMPS run.

\sphinxAtStartPar
Here is an example of running a series of simulations using the next
command with an \sphinxstyleemphasis{index}\sphinxhyphen{}style variable.  If this input script is named
in.polymer, 8 simulations would be run using data files from
directories run1 through run8.

\begin{sphinxVerbatim}[commandchars=\\\{\}]
\PYG{k}{variable }\PYG{n+nv+nvIdentifier}{d}\PYG{+w}{ }\PYG{n}{index}\PYG{+w}{ }\PYG{n}{run1}\PYG{+w}{ }\PYG{n}{run2}\PYG{+w}{ }\PYG{n}{run3}\PYG{+w}{ }\PYG{n}{run4}\PYG{+w}{ }\PYG{n}{run5}\PYG{+w}{ }\PYG{n}{run6}\PYG{+w}{ }\PYG{n}{run7}\PYG{+w}{ }\PYG{n}{run8}
\PYG{k}{shell}\PYG{+w}{ }\PYG{n}{cd}\PYG{+w}{ }\PYG{n+nv}{\PYGZdl{}d}
\PYG{k}{read\PYGZus{}data}\PYG{+w}{ }\PYG{n}{data.polymer}
\PYG{k}{run}\PYG{+w}{ }\PYG{l+m}{10000}
\PYG{k}{shell}\PYG{+w}{ }\PYG{n}{cd}\PYG{+w}{ }\PYG{n}{..}
\PYG{k}{clear}
\PYG{k}{next }\PYG{n+nv+nvIdentifier}{d}
\PYG{k}{jump }\PYG{l+s+sc}{in.polymer}
\end{sphinxVerbatim}

\sphinxAtStartPar
If the variable “d” were of style \sphinxstyleemphasis{universe}, and the same in.polymer
input script were run on 3 partitions of processors, then the first 3
simulations would begin, one on each set of processors.  Whichever
partition finished first, it would assign variable “d” the fourth value
and run another simulation, and so forth until all 8 simulations were
finished.

\sphinxAtStartPar
Jump and next commands can also be nested to enable multi\sphinxhyphen{}level loops.
For example, this script will run 15 simulations in a double loop.

\begin{sphinxVerbatim}[commandchars=\\\{\}]
\PYG{k}{variable }\PYG{n+nv+nvIdentifier}{i}\PYG{+w}{ }\PYG{n}{loop}\PYG{+w}{ }\PYG{l+m}{3}
\PYG{+w}{  }\PYG{k}{variable }\PYG{n+nv+nvIdentifier}{j}\PYG{+w}{ }\PYG{n}{loop}\PYG{+w}{ }\PYG{l+m}{5}
\PYG{+w}{  }\PYG{n}{clear}
\PYG{+w}{  }\PYG{n}{...}
\PYG{+w}{  }\PYG{n}{read\PYGZus{}data}\PYG{+w}{ }\PYG{n}{data.polymer.}\PYG{n+nv}{\PYGZdl{}i}\PYG{n+nv}{\PYGZdl{}j}
\PYG{+w}{  }\PYG{n}{print}\PYG{+w}{ }\PYG{n}{Running}\PYG{+w}{ }\PYG{n}{simulation}\PYG{+w}{ }\PYG{n+nv}{\PYGZdl{}i}\PYG{n}{.}\PYG{n+nv}{\PYGZdl{}j}
\PYG{+w}{  }\PYG{n}{run}\PYG{+w}{ }\PYG{l+m}{10000}
\PYG{+w}{  }\PYG{k}{next }\PYG{n+nv+nvIdentifier}{j}
\PYG{+w}{  }\PYG{k}{jump }\PYG{l+s+sc}{in.script}
\PYG{k}{next }\PYG{n+nv+nvIdentifier}{i}
\PYG{k}{jump }\PYG{l+s+sc}{in.script}
\end{sphinxVerbatim}

\sphinxAtStartPar
Here is an example of a double loop which uses the {\hyperref[\detokenize{if::doc}]{\sphinxcrossref{\DUrole{doc}{if}}}} and
{\hyperref[\detokenize{jump::doc}]{\sphinxcrossref{\DUrole{doc}{jump}}}} commands to break out of the inner loop when a
condition is met, then continues iterating through the outer loop.

\begin{sphinxVerbatim}[commandchars=\\\{\}]
\PYG{k}{label       }\PYG{l+s+sc}{loopa}
\PYG{k}{variable    }\PYG{n+nv+nvIdentifier}{a}\PYG{+w}{ }\PYG{n}{loop}\PYG{+w}{ }\PYG{l+m}{5}
\PYG{+w}{  }\PYG{k}{label     }\PYG{l+s+sc}{loopb}
\PYG{+w}{  }\PYG{k}{variable  }\PYG{n+nv+nvIdentifier}{b}\PYG{+w}{ }\PYG{n}{loop}\PYG{+w}{ }\PYG{l+m}{5}
\PYG{+w}{  }\PYG{n}{print}\PYG{+w}{     }\PYG{l+s}{\PYGZdq{}}\PYG{l+s}{A,B = \PYGZdl{}a,\PYGZdl{}b}\PYG{l+s}{\PYGZdq{}}
\PYG{+w}{  }\PYG{n}{run}\PYG{+w}{       }\PYG{l+m}{10000}
\PYG{+w}{  }\PYG{k}{if}\PYG{+w}{        }\PYG{n+nv}{\PYGZdl{}b}\PYG{+w}{ }\PYG{o}{\PYGZgt{}}\PYG{+w}{ }\PYG{l+m}{2}\PYG{+w}{ }\PYG{k}{then}\PYG{+w}{ }\PYG{l+s}{\PYGZdq{}}\PYG{l+s}{jump in.script break}\PYG{l+s}{\PYGZdq{}}
\PYG{+w}{  }\PYG{k}{next      }\PYG{n+nv+nvIdentifier}{b}
\PYG{+w}{  }\PYG{k}{jump      }\PYG{l+s+sc}{in.script}\PYG{+w}{ }\PYG{n}{loopb}
\PYG{k}{label       }\PYG{l+s+sc}{break}
\PYG{k}{variable    }\PYG{n+nv+nvIdentifier}{b}\PYG{+w}{ }\PYG{n}{delete}

\PYG{k}{next        }\PYG{n+nv+nvIdentifier}{a}
\PYG{k}{jump        }\PYG{l+s+sc}{in.script}\PYG{+w}{ }\PYG{n}{loopa}
\end{sphinxVerbatim}


\subsection{Restrictions}
\label{\detokenize{next:restrictions}}
\sphinxAtStartPar
As described above.


\subsection{Related commands}
\label{\detokenize{next:related-commands}}
\sphinxAtStartPar
{\hyperref[\detokenize{jump::doc}]{\sphinxcrossref{\DUrole{doc}{jump}}}}, {\hyperref[\detokenize{include::doc}]{\sphinxcrossref{\DUrole{doc}{include}}}}, {\hyperref[\detokenize{shell::doc}]{\sphinxcrossref{\DUrole{doc}{shell}}}},
{\hyperref[\detokenize{variable::doc}]{\sphinxcrossref{\DUrole{doc}{variable}}}},


\subsection{Default}
\label{\detokenize{next:default}}
\sphinxAtStartPar
none

\sphinxstepscope

\index{package@\spxentry{package}}\ignorespaces 

\section{package command}
\label{\detokenize{package:package-command}}\label{\detokenize{package:index-0}}\label{\detokenize{package::doc}}

\subsection{Syntax}
\label{\detokenize{package:syntax}}
\begin{sphinxVerbatim}[commandchars=\\\{\}]
\PYG{k}{package}\PYG{+w}{ }\PYG{n}{style}\PYG{+w}{ }\PYG{n}{args}
\end{sphinxVerbatim}
\begin{itemize}
\item {} 
\sphinxAtStartPar
style = \sphinxstyleemphasis{gpu} or \sphinxstyleemphasis{intel} or \sphinxstyleemphasis{kokkos} or \sphinxstyleemphasis{omp}

\item {} 
\sphinxAtStartPar
args = arguments specific to the style
\begin{sphinxalltt}
\sphinxstyleemphasis{gpu} args = Ngpu keyword value ...
  Ngpu = \# of GPUs per node
  zero or more keyword/value pairs may be appended
  keywords = \sphinxstyleemphasis{neigh} or \sphinxstyleemphasis{newton} or \sphinxstyleemphasis{pair/only} or \sphinxstyleemphasis{binsize} or \sphinxstyleemphasis{split} or \sphinxstyleemphasis{gpuID} or \sphinxstyleemphasis{tpa} or \sphinxstyleemphasis{blocksize} or \sphinxstyleemphasis{omp} or \sphinxstyleemphasis{platform} or \sphinxstyleemphasis{device\_type} or \sphinxstyleemphasis{ocl\_args}
    \sphinxstyleemphasis{neigh} value = \sphinxstyleemphasis{yes} or \sphinxstyleemphasis{no}
      \sphinxstyleemphasis{yes} = neighbor list build on GPU (default)
      \sphinxstyleemphasis{no} = neighbor list build on CPU
    \sphinxstyleemphasis{newton} = \sphinxstyleemphasis{off} or \sphinxstyleemphasis{on}
      \sphinxstyleemphasis{off} = set Newton pairwise flag off (default and required)
      \sphinxstyleemphasis{on} = set Newton pairwise flag on (currently not allowed)
    \sphinxstyleemphasis{pair/only} = \sphinxstyleemphasis{off} or \sphinxstyleemphasis{on}
      \sphinxstyleemphasis{off} = apply "gpu" suffix to all available styles in the GPU package (default)
      \sphinxstyleemphasis{on} = apply "gpu" suffix only pair styles
    \sphinxstyleemphasis{binsize} value = size
      size = bin size for neighbor list construction (distance units)
    \sphinxstyleemphasis{split} = fraction
      fraction = fraction of atoms assigned to GPU (default = 1.0)
    \sphinxstyleemphasis{tpa} value = Nlanes
      Nlanes = \# of GPU vector lanes (CUDA threads) used per atom
    \sphinxstyleemphasis{blocksize} value = size
      size = thread block size for pair force computation
    \sphinxstyleemphasis{omp} value = Nthreads
      Nthreads = number of OpenMP threads to use on CPU (default = 0)
    \sphinxstyleemphasis{platform} value = id
      id = For OpenCL, platform ID for the GPU or accelerator
    \sphinxstyleemphasis{gpuID} values = id
      id = ID of first GPU to be used on each node
    \sphinxstyleemphasis{device\_type} value = \sphinxstyleemphasis{intelgpu} or \sphinxstyleemphasis{nvidiagpu} or \sphinxstyleemphasis{amdgpu} or \sphinxstyleemphasis{applegpu} or \sphinxstyleemphasis{generic} or \sphinxstyleemphasis{custom},val1,val2,...
      val1,val2,... = custom OpenCL accelerator configuration parameters (see below for details)
    \sphinxstyleemphasis{ocl\_args} value = args
      args = List of additional OpenCL compiler arguments delimited by colons
\sphinxstyleemphasis{intel} args = NPhi keyword value ...
  Nphi = \# of co\sphinxhyphen{}processors per node
  zero or more keyword/value pairs may be appended
  keywords = \sphinxstyleemphasis{mode} or \sphinxstyleemphasis{omp} or \sphinxstyleemphasis{lrt} or \sphinxstyleemphasis{balance} or \sphinxstyleemphasis{ghost} or \sphinxstyleemphasis{tpc} or \sphinxstyleemphasis{tptask} or \sphinxstyleemphasis{pppm\_table} or \sphinxstyleemphasis{no\_affinity}
    \sphinxstyleemphasis{mode} value = \sphinxstyleemphasis{single} or \sphinxstyleemphasis{mixed} or \sphinxstyleemphasis{double}
      single = perform force calculations in single precision
      mixed = perform force calculations in mixed precision
      double = perform force calculations in double precision
    \sphinxstyleemphasis{omp} value = Nthreads
      Nthreads = number of OpenMP threads to use on CPU (default = 0)
    \sphinxstyleemphasis{lrt} value = \sphinxstyleemphasis{yes} or \sphinxstyleemphasis{no}
      \sphinxstyleemphasis{yes} = use additional thread dedicated for some PPPM calculations
      \sphinxstyleemphasis{no} = do not dedicate an extra thread for some PPPM calculations
    \sphinxstyleemphasis{balance} value = split
      split = fraction of work to offload to co\sphinxhyphen{}processor, \sphinxhyphen{}1 for dynamic
    \sphinxstyleemphasis{ghost} value = \sphinxstyleemphasis{yes} or \sphinxstyleemphasis{no}
      \sphinxstyleemphasis{yes} = include ghost atoms for offload
      \sphinxstyleemphasis{no} = do not include ghost atoms for offload
    \sphinxstyleemphasis{tpc} value = Ntpc
      Ntpc = max number of co\sphinxhyphen{}processor threads per co\sphinxhyphen{}processor core (default = 4)
    \sphinxstyleemphasis{tptask} value = Ntptask
      Ntptask = max number of co\sphinxhyphen{}processor threads per MPI task (default = 240)
    \sphinxstyleemphasis{pppm\_table} value = \sphinxstyleemphasis{yes} or \sphinxstyleemphasis{no}
      \sphinxstyleemphasis{yes} = Precompute pppm values in table (doesn\textquotesingle{}t change accuracy)
      \sphinxstyleemphasis{no} = Compute pppm values on the fly
    \sphinxstyleemphasis{no\_affinity} values = none
\sphinxstyleemphasis{kokkos} args = keyword value ...
  zero or more keyword/value pairs may be appended
  keywords = \sphinxstyleemphasis{neigh} or \sphinxstyleemphasis{neigh/qeq} or \sphinxstyleemphasis{neigh/thread} or \sphinxstyleemphasis{neigh/transpose} or \sphinxstyleemphasis{newton} or \sphinxstyleemphasis{binsize} or \sphinxstyleemphasis{comm} or \sphinxstyleemphasis{comm/exchange} or \sphinxstyleemphasis{comm/forward} or \sphinxstyleemphasis{comm/pair/forward} or \sphinxstyleemphasis{comm/fix/forward} or \sphinxstyleemphasis{comm/reverse} or \sphinxstyleemphasis{comm/pair/reverse} or \sphinxstyleemphasis{sort} or \sphinxstyleemphasis{atom/map} or \sphinxstyleemphasis{gpu/aware} or \sphinxstyleemphasis{pair/only}
    \sphinxstyleemphasis{neigh} value = \sphinxstyleemphasis{full} or \sphinxstyleemphasis{half}
      full = full neighbor list
      half = half neighbor list built in thread\sphinxhyphen{}safe manner
    \sphinxstyleemphasis{neigh/qeq} value = \sphinxstyleemphasis{full} or \sphinxstyleemphasis{half}
      full = full neighbor list
      half = half neighbor list built in thread\sphinxhyphen{}safe manner
    \sphinxstyleemphasis{neigh/thread} value = \sphinxstyleemphasis{off} or \sphinxstyleemphasis{on}
      \sphinxstyleemphasis{off} = thread only over atoms
      \sphinxstyleemphasis{on} = thread over both atoms and neighbors
    \sphinxstyleemphasis{neigh/transpose} value = \sphinxstyleemphasis{off} or \sphinxstyleemphasis{on}
      \sphinxstyleemphasis{off} = use same memory layout for GPU neigh list build as pair style
      \sphinxstyleemphasis{on} = use transposed memory layout for GPU neigh list build
    \sphinxstyleemphasis{newton} = \sphinxstyleemphasis{off} or \sphinxstyleemphasis{on}
      \sphinxstyleemphasis{off} = set Newton pairwise and bonded flags off
      \sphinxstyleemphasis{on} = set Newton pairwise and bonded flags on
    \sphinxstyleemphasis{binsize} value = size
      size = bin size for neighbor list construction (distance units)
    \sphinxstyleemphasis{comm} value = \sphinxstyleemphasis{no} or \sphinxstyleemphasis{host} or \sphinxstyleemphasis{device}
      use value for comm/exchange and comm/forward and comm/pair/forward and comm/fix/forward and comm/reverse
    \sphinxstyleemphasis{comm/exchange} value = \sphinxstyleemphasis{no} or \sphinxstyleemphasis{host} or \sphinxstyleemphasis{device}
    \sphinxstyleemphasis{comm/forward} value = \sphinxstyleemphasis{no} or \sphinxstyleemphasis{host} or \sphinxstyleemphasis{device}
    \sphinxstyleemphasis{comm/pair/forward} value = \sphinxstyleemphasis{no} or \sphinxstyleemphasis{device}
    \sphinxstyleemphasis{comm/fix/forward} value = \sphinxstyleemphasis{no} or \sphinxstyleemphasis{device}
    \sphinxstyleemphasis{comm/reverse} value = \sphinxstyleemphasis{no} or \sphinxstyleemphasis{host} or \sphinxstyleemphasis{device}
      \sphinxstyleemphasis{no} = perform communication pack/unpack in non\sphinxhyphen{}KOKKOS mode
      \sphinxstyleemphasis{host} = perform pack/unpack on host (e.g. with OpenMP threading)
      \sphinxstyleemphasis{device} = perform pack/unpack on device (e.g. on GPU)
    \sphinxstyleemphasis{comm/pair/reverse} value = \sphinxstyleemphasis{no} or \sphinxstyleemphasis{device}
      \sphinxstyleemphasis{no} = perform communication pack/unpack in non\sphinxhyphen{}KOKKOS mode
      \sphinxstyleemphasis{device} = perform pack/unpack on device (e.g. on GPU)
    \sphinxstyleemphasis{sort} value = \sphinxstyleemphasis{no} or \sphinxstyleemphasis{device}
      \sphinxstyleemphasis{no} = perform atom sorting in non\sphinxhyphen{}KOKKOS mode
      \sphinxstyleemphasis{device} = perform atom sorting on device (e.g. on GPU)
    \sphinxstyleemphasis{atom/map} value = \sphinxstyleemphasis{no} or \sphinxstyleemphasis{device}
      \sphinxstyleemphasis{no} = build atom map in non\sphinxhyphen{}KOKKOS mode
      \sphinxstyleemphasis{device} = build atom map on device (e.g. on GPU)
    \sphinxstyleemphasis{gpu/aware} = \sphinxstyleemphasis{off} or \sphinxstyleemphasis{on}
      \sphinxstyleemphasis{off} = do not use GPU\sphinxhyphen{}aware MPI
      \sphinxstyleemphasis{on} = use GPU\sphinxhyphen{}aware MPI (default)
    \sphinxstyleemphasis{pair/only} = \sphinxstyleemphasis{off} or \sphinxstyleemphasis{on}
      \sphinxstyleemphasis{off} = use device acceleration (e.g. GPU) for all available styles in the KOKKOS package (default)
      \sphinxstyleemphasis{on}  = use device acceleration only for pair styles (and host acceleration for others)
\sphinxstyleemphasis{omp} args = Nthreads keyword value ...
  Nthreads = \# of OpenMP threads to associate with each MPI process
  zero or more keyword/value pairs may be appended
  keywords = \sphinxstyleemphasis{neigh}
    \sphinxstyleemphasis{neigh} value = \sphinxstyleemphasis{yes} or \sphinxstyleemphasis{no}
      \sphinxstyleemphasis{yes} = threaded neighbor list build (default)
      \sphinxstyleemphasis{no} = non\sphinxhyphen{}threaded neighbor list build
\end{sphinxalltt}

\end{itemize}


\subsection{Examples}
\label{\detokenize{package:examples}}
\begin{sphinxVerbatim}[commandchars=\\\{\}]
\PYG{k}{package}\PYG{+w}{ }\PYG{n}{gpu}\PYG{+w}{ }\PYG{l+m}{0}
\PYG{k}{package}\PYG{+w}{ }\PYG{n}{gpu}\PYG{+w}{ }\PYG{l+m}{1}\PYG{+w}{ }\PYG{n}{split}\PYG{+w}{ }\PYG{l+m}{0.75}
\PYG{k}{package}\PYG{+w}{ }\PYG{n}{gpu}\PYG{+w}{ }\PYG{l+m}{2}\PYG{+w}{ }\PYG{n}{split}\PYG{+w}{ }\PYG{o}{\PYGZhy{}}\PYG{l+m}{1.0}
\PYG{k}{package}\PYG{+w}{ }\PYG{n}{gpu}\PYG{+w}{ }\PYG{l+m}{0}\PYG{+w}{ }\PYG{n}{omp}\PYG{+w}{ }\PYG{l+m}{2}\PYG{+w}{ }\PYG{n}{device\PYGZus{}type}\PYG{+w}{ }\PYG{n}{intelgpu}
\PYG{k}{package}\PYG{+w}{ }\PYG{n}{kokkos}\PYG{+w}{ }\PYG{n}{neigh}\PYG{+w}{ }\PYG{n}{half}\PYG{+w}{ }\PYG{n}{comm}\PYG{+w}{ }\PYG{n}{device}
\PYG{k}{package}\PYG{+w}{ }\PYG{n}{omp}\PYG{+w}{ }\PYG{l+m}{0}\PYG{+w}{ }\PYG{n}{neigh}\PYG{+w}{ }\PYG{n}{no}
\PYG{k}{package}\PYG{+w}{ }\PYG{n}{omp}\PYG{+w}{ }\PYG{l+m}{4}
\PYG{k}{package}\PYG{+w}{ }\PYG{n}{intel}\PYG{+w}{ }\PYG{l+m}{1}
\PYG{k}{package}\PYG{+w}{ }\PYG{n}{intel}\PYG{+w}{ }\PYG{l+m}{2}\PYG{+w}{ }\PYG{n}{omp}\PYG{+w}{ }\PYG{l+m}{4}\PYG{+w}{ }\PYG{n}{mode}\PYG{+w}{ }\PYG{n}{mixed}\PYG{+w}{ }\PYG{n}{balance}\PYG{+w}{ }\PYG{l+m}{0.5}
\end{sphinxVerbatim}


\subsection{Description}
\label{\detokenize{package:description}}
\sphinxAtStartPar
This command invokes package\sphinxhyphen{}specific settings for the various
accelerator packages available in LAMMPS.  Currently the following
packages use settings from this command: GPU, INTEL, KOKKOS, and
OPENMP.

\sphinxAtStartPar
If this command is specified in an input script, it must be near the
top of the script, before the simulation box has been defined.  This
is because it specifies settings that the accelerator packages use in
their initialization, before a simulation is defined.

\sphinxAtStartPar
This command can also be specified from the command\sphinxhyphen{}line when
launching LAMMPS, using the “\sphinxhyphen{}pk” {\hyperref[\detokenize{Run_options::doc}]{\sphinxcrossref{\DUrole{doc}{command\sphinxhyphen{}line switch}}}}.  The syntax is exactly the same as when used
in an input script.

\sphinxAtStartPar
Note that all of the accelerator packages require the package command
to be specified (except the OPT package), if the package is to be used
in a simulation (LAMMPS can be built with an accelerator package
without using it in a particular simulation).  However, in all cases,
a default version of the command is typically invoked by other
accelerator settings.

\sphinxAtStartPar
The KOKKOS package requires a “\sphinxhyphen{}k on” {\hyperref[\detokenize{Run_options::doc}]{\sphinxcrossref{\DUrole{doc}{command\sphinxhyphen{}line switch}}}} respectively, which invokes a “package
kokkos” command with default settings.

\sphinxAtStartPar
For the GPU, INTEL, and OPENMP packages, if a “\sphinxhyphen{}sf gpu” or “\sphinxhyphen{}sf
intel” or “\sphinxhyphen{}sf omp” {\hyperref[\detokenize{Run_options::doc}]{\sphinxcrossref{\DUrole{doc}{command\sphinxhyphen{}line switch}}}} is used to
auto\sphinxhyphen{}append accelerator suffixes to various styles in the input
script, then those switches also invoke a “package gpu”, “package
intel”, or “package omp” command with default settings.

\begin{sphinxadmonition}{note}{Note:}
\sphinxAtStartPar
A package command for a particular style can be invoked multiple
times when a simulation is setup, e.g. by the {\hyperref[\detokenize{Run_options::doc}]{\sphinxcrossref{\DUrole{doc}{\sphinxhyphen{}c on, \sphinxhyphen{}k on, \sphinxhyphen{}sf, and \sphinxhyphen{}pk command\sphinxhyphen{}line switches}}}}, and by using this command
in an input script.  Each time it is used all of the style options are
set, either to default values or to specified settings.  I.e. settings
from previous invocations do not persist across multiple invocations.
\end{sphinxadmonition}

\sphinxAtStartPar
See the {\hyperref[\detokenize{Speed_packages::doc}]{\sphinxcrossref{\DUrole{doc}{Accelerator packages}}}} page for more details
about using the various accelerator packages for speeding up LAMMPS
simulations.


\bigskip\hrule\bigskip


\sphinxAtStartPar
The \sphinxstyleemphasis{gpu} style invokes settings associated with the use of the GPU
package.

\sphinxAtStartPar
The \sphinxstyleemphasis{Ngpu} argument sets the number of GPUs per node. If \sphinxstyleemphasis{Ngpu} is 0
and no other keywords are specified, GPU or accelerator devices are
auto\sphinxhyphen{}selected. In this process, all platforms are searched for
accelerator devices and GPUs are chosen if available. The device with
the highest number of compute cores is selected. The number of devices
is increased to be the number of matching accelerators with the same
number of compute cores. If there are more devices than MPI tasks,
the additional devices will be unused. The auto\sphinxhyphen{}selection of GPUs/
accelerator devices and platforms can be restricted by specifying
a non\sphinxhyphen{}zero value for \sphinxstyleemphasis{Ngpu} and / or using the \sphinxstyleemphasis{gpuID}, \sphinxstyleemphasis{platform},
and \sphinxstyleemphasis{device\_type} keywords as described below. If there are more MPI
tasks (per node) than GPUs, multiple MPI tasks will share each GPU.

\sphinxAtStartPar
Optional keyword/value pairs can also be specified.  Each has a
default value as listed below.

\sphinxAtStartPar
The \sphinxstyleemphasis{neigh} keyword specifies where neighbor lists for pair style
computation will be built.  If \sphinxstyleemphasis{neigh} is \sphinxstyleemphasis{yes}, which is the default,
neighbor list building is performed on the GPU.  If \sphinxstyleemphasis{neigh} is \sphinxstyleemphasis{no},
neighbor list building is performed on the CPU.  GPU neighbor list
building currently cannot be used with a triclinic box.  GPU neighbor
lists are not compatible with commands that are not GPU\sphinxhyphen{}enabled.  When
a non\sphinxhyphen{}GPU enabled command requires a neighbor list, it will also be
built on the CPU.  In these cases, it will typically be more efficient
to only use CPU neighbor list builds.

\sphinxAtStartPar
The \sphinxstyleemphasis{newton} keyword sets the Newton flags for pairwise (not bonded)
interactions to \sphinxstyleemphasis{off} or \sphinxstyleemphasis{on}, the same as the {\hyperref[\detokenize{newton::doc}]{\sphinxcrossref{\DUrole{doc}{newton}}}}
command allows.  Currently, only an \sphinxstyleemphasis{off} value is allowed, since all
the GPU package pair styles require this setting.  This means more
computation is done, but less communication.  In the future a value of
\sphinxstyleemphasis{on} may be allowed, so the \sphinxstyleemphasis{newton} keyword is included as an option
for compatibility with the package command for other accelerator
styles.  Note that the newton setting for bonded interactions is not
affected by this keyword.

\sphinxAtStartPar
The \sphinxstyleemphasis{pair/only} keyword can change how any “gpu” suffix is applied.
By default a suffix is applied to all styles for which an accelerated
variant is available.  However, that is not always the most effective
way to use an accelerator.  With \sphinxstyleemphasis{pair/only} set to \sphinxstyleemphasis{on} the suffix
will only by applied to supported pair styles, which tend to be the
most effective in using an accelerator and their operation can be
overlapped with all other computations on the CPU.

\sphinxAtStartPar
The \sphinxstyleemphasis{binsize} keyword sets the size of bins used to bin atoms in
neighbor list builds performed on the GPU, if \sphinxstyleemphasis{neigh} = \sphinxstyleemphasis{yes} is set.
If \sphinxstyleemphasis{binsize} is set to 0.0 (the default), then the binsize is set
automatically using heuristics in the GPU package.

\sphinxAtStartPar
The \sphinxstyleemphasis{split} keyword can be used for load balancing force calculations
between CPU and GPU cores in GPU\sphinxhyphen{}enabled pair styles. If 0 \textless{} \sphinxstyleemphasis{split} \textless{}
1.0, a fixed fraction of particles is offloaded to the GPU while force
calculation for the other particles occurs simultaneously on the CPU.
If \sphinxstyleemphasis{split} \textless{} 0.0, the optimal fraction (based on CPU and GPU timings)
is calculated every 25 timesteps, i.e. dynamic load\sphinxhyphen{}balancing across
the CPU and GPU is performed.  If \sphinxstyleemphasis{split} = 1.0, all force
calculations for GPU accelerated pair styles are performed on the GPU.
In this case, other {\hyperref[\detokenize{pair_hybrid::doc}]{\sphinxcrossref{\DUrole{doc}{hybrid}}}} pair interactions,
{\hyperref[\detokenize{bond_style::doc}]{\sphinxcrossref{\DUrole{doc}{bond}}}}, {\hyperref[\detokenize{angle_style::doc}]{\sphinxcrossref{\DUrole{doc}{angle}}}},
{\hyperref[\detokenize{dihedral_style::doc}]{\sphinxcrossref{\DUrole{doc}{dihedral}}}}, {\hyperref[\detokenize{improper_style::doc}]{\sphinxcrossref{\DUrole{doc}{improper}}}}, and
{\hyperref[\detokenize{kspace_style::doc}]{\sphinxcrossref{\DUrole{doc}{long\sphinxhyphen{}range}}}} calculations can be performed on the
CPU while the GPU is performing force calculations for the GPU\sphinxhyphen{}enabled
pair style.  If all CPU force computations complete before the GPU
completes, LAMMPS will block until the GPU has finished before
continuing the timestep.

\sphinxAtStartPar
As an example, if you have two GPUs per node and 8 CPU cores per node,
and would like to run on 4 nodes (32 cores) with dynamic balancing of
force calculation across CPU and GPU cores, you could specify

\begin{sphinxVerbatim}[commandchars=\\\{\}]
mpirun\PYG{+w}{ }\PYGZhy{}np\PYG{+w}{ }\PYG{l+m}{32}\PYG{+w}{ }\PYGZhy{}sf\PYG{+w}{ }gpu\PYG{+w}{ }\PYGZhy{}in\PYG{+w}{ }\PYG{k}{in}.script\PYG{+w}{    }\PYG{c+c1}{\PYGZsh{} launch command}
package\PYG{+w}{ }gpu\PYG{+w}{ }\PYG{l+m}{2}\PYG{+w}{ }split\PYG{+w}{ }\PYGZhy{}1\PYG{+w}{                 }\PYG{c+c1}{\PYGZsh{} input script command}
\end{sphinxVerbatim}

\sphinxAtStartPar
In this case, all CPU cores and GPU devices on the nodes would be
utilized.  Each GPU device would be shared by 4 CPU cores. The CPU
cores would perform force calculations for some fraction of the
particles at the same time the GPUs performed force calculation for
the other particles.

\sphinxAtStartPar
The \sphinxstyleemphasis{gpuID} keyword is used to specify the first ID for the GPU or
other accelerator that LAMMPS will use. For example, if the ID is
1 and \sphinxstyleemphasis{Ngpu} is 3, GPUs 1\sphinxhyphen{}3 will be used. Device IDs should be
determined from the output of nvc\_get\_devices, ocl\_get\_devices,
or hip\_get\_devices
as provided in the lib/gpu directory. When using OpenCL with
accelerators that have main memory NUMA, the accelerators can be
split into smaller virtual accelerators for more efficient use
with MPI.

\sphinxAtStartPar
The \sphinxstyleemphasis{tpa} keyword sets the number of GPU vector lanes per atom used to
perform force calculations.  With a default value of 1, the number of
lanes will be chosen based on the pair style, however, the value can
be set explicitly with this keyword to fine\sphinxhyphen{}tune performance.  For
large cutoffs or with a small number of particles per GPU, increasing
the value can improve performance. The number of lanes per atom must
be a power of 2 and currently cannot be greater than the SIMD width
for the GPU / accelerator. In the case it exceeds the SIMD width, it
will automatically be decreased to meet the restriction.

\sphinxAtStartPar
The \sphinxstyleemphasis{blocksize} keyword allows you to tweak the number of threads used
per thread block. This number should be a multiple of 32 (for GPUs)
and its maximum depends on the specific GPU hardware. Typical choices
are 64, 128, or 256. A larger block size increases occupancy of
individual GPU cores, but reduces the total number of thread blocks,
thus may lead to load imbalance. On modern hardware, the sensitivity
to the blocksize is typically low.

\sphinxAtStartPar
The \sphinxstyleemphasis{Nthreads} value for the \sphinxstyleemphasis{omp} keyword sets the number of OpenMP
threads allocated for each MPI task. This setting controls OpenMP
parallelism only for routines run on the CPUs. For more details on
setting the number of OpenMP threads, see the discussion of the
\sphinxstyleemphasis{Nthreads} setting on this page for the “package omp” command.
The meaning of \sphinxstyleemphasis{Nthreads} is exactly the same for the GPU, INTEL,
and GPU packages.

\sphinxAtStartPar
The \sphinxstyleemphasis{platform} keyword is only used with OpenCL to specify the ID for
an OpenCL platform. See the output from ocl\_get\_devices in the lib/gpu
directory. In LAMMPS only one platform can be active at a time and by
default (id=\sphinxhyphen{}1) the platform is auto\sphinxhyphen{}selected to find the GPU with the
most compute cores. When \sphinxstyleemphasis{Ngpu} or other keywords are specified, the
auto\sphinxhyphen{}selection is appropriately restricted. For example, if \sphinxstyleemphasis{Ngpu} is
3, only platforms with at least 3 accelerators are considered. Similar
restrictions can be enforced by the \sphinxstyleemphasis{gpuID} and \sphinxstyleemphasis{device\_type} keywords.

\sphinxAtStartPar
The \sphinxstyleemphasis{device\_type} keyword can be used for OpenCL to specify the type of
GPU to use or specify a custom configuration for an accelerator. In most
cases this selection will be automatic and there is no need to use the
keyword. The \sphinxstyleemphasis{applegpu} type is not specific to a particular GPU vendor,
but is separate due to the more restrictive Apple OpenCL implementation.
For expert users, to specify a custom configuration, the \sphinxstyleemphasis{custom} keyword
followed by the next parameters can be specified:

\sphinxAtStartPar
CONFIG\_ID, SIMD\_SIZE, MEM\_THREADS, SHUFFLE\_AVAIL, FAST\_MATH,
THREADS\_PER\_ATOM, THREADS\_PER\_CHARGE, THREADS\_PER\_THREE, BLOCK\_PAIR,
BLOCK\_BIO\_PAIR, BLOCK\_ELLIPSE, PPPM\_BLOCK\_1D, BLOCK\_NBOR\_BUILD,
BLOCK\_CELL\_2D, BLOCK\_CELL\_ID, MAX\_SHARED\_TYPES, MAX\_BIO\_SHARED\_TYPES,
PPPM\_MAX\_SPLINE, NBOR\_PREFETCH.

\sphinxAtStartPar
CONFIG\_ID can be 0. SHUFFLE\_AVAIL in \{0,1\} indicates that inline\sphinxhyphen{}PTX
(NVIDIA) or OpenCL extensions (Intel) should be used for horizontal
vector operations. FAST\_MATH in \{0,1\} indicates that OpenCL fast math
optimizations are used during the build and hardware\sphinxhyphen{}accelerated
transcendental functions are used when available. THREADS\_PER\_* give the
default \sphinxstyleemphasis{tpa} values for ellipsoidal models, styles using charge, and
any other styles. The BLOCK\_* parameters specify the block sizes for
various kernel calls and the MAX\_*SHARED*\_ parameters are used to
determine the amount of local shared memory to use for storing model
parameters.

\sphinxAtStartPar
For OpenCL, the routines are compiled at runtime for the specified GPU
or accelerator architecture. The \sphinxstyleemphasis{ocl\_args} keyword can be used to
specify additional flags for the runtime build.


\bigskip\hrule\bigskip


\sphinxAtStartPar
The \sphinxstyleemphasis{intel} style invokes settings associated with the use of the INTEL
package.  The keywords \sphinxstyleemphasis{balance}, \sphinxstyleemphasis{ghost}, \sphinxstyleemphasis{tpc}, and \sphinxstyleemphasis{tptask} are
\sphinxstylestrong{only} applicable if LAMMPS was built with Xeon Phi co\sphinxhyphen{}processor
support and are otherwise ignored.

\sphinxAtStartPar
The \sphinxstyleemphasis{Nphi} argument sets the number of co\sphinxhyphen{}processors per node.
This can be set to any value, including 0, if LAMMPS was not
built with co\sphinxhyphen{}processor support.

\sphinxAtStartPar
Optional keyword/value pairs can also be specified.  Each has a
default value as listed below.

\sphinxAtStartPar
The \sphinxstyleemphasis{Nthreads} value for the \sphinxstyleemphasis{omp} keyword sets the number of OpenMP
threads allocated for each MPI task. This setting controls OpenMP
parallelism only for routines run on the CPUs. For more details on
setting the number of OpenMP threads, see the discussion of the
\sphinxstyleemphasis{Nthreads} setting on this page for the “package omp” command.
The meaning of \sphinxstyleemphasis{Nthreads} is exactly the same for the GPU, INTEL,
and GPU packages.

\sphinxAtStartPar
The \sphinxstyleemphasis{mode} keyword determines the precision mode to use for
computing pair style forces, either on the CPU or on the co\sphinxhyphen{}processor,
when using a INTEL supported {\hyperref[\detokenize{pair_style::doc}]{\sphinxcrossref{\DUrole{doc}{pair style}}}}.  It
can take a value of \sphinxstyleemphasis{single}, \sphinxstyleemphasis{mixed} which is the default, or
\sphinxstyleemphasis{double}.  \sphinxstyleemphasis{Single} means single precision is used for the entire
force calculation.  \sphinxstyleemphasis{Mixed} means forces between a pair of atoms are
computed in single precision, but accumulated and stored in double
precision, including storage of forces, torques, energies, and virial
quantities.  \sphinxstyleemphasis{Double} means double precision is used for the entire
force calculation.

\sphinxAtStartPar
The \sphinxstyleemphasis{lrt} keyword can be used to enable “Long Range Thread (LRT)”
mode. It can take a value of \sphinxstyleemphasis{yes} to enable and \sphinxstyleemphasis{no} to disable.
LRT mode generates an extra thread (in addition to any OpenMP threads
specified with the OMP\_NUM\_THREADS environment variable or the \sphinxstyleemphasis{omp}
keyword). The extra thread is dedicated for performing part of the
{\hyperref[\detokenize{kspace_style::doc}]{\sphinxcrossref{\DUrole{doc}{PPPM solver}}}} computations and communications. This
can improve parallel performance on processors supporting
Simultaneous Multithreading (SMT) such as Hyper\sphinxhyphen{}Threading (HT) on Intel
processors. In this mode, one additional thread is generated per MPI
process. LAMMPS will generate a warning in the case that more threads
are used than available in SMT hardware on a node. If the PPPM solver
from the INTEL package is not used, then the LRT setting is
ignored and no extra threads are generated. Enabling LRT will replace
the {\hyperref[\detokenize{run_style::doc}]{\sphinxcrossref{\DUrole{doc}{run\_style}}}} with the \sphinxstyleemphasis{verlet/lrt/intel} style that
is identical to the default \sphinxstyleemphasis{verlet} style aside from supporting the
LRT feature. This feature requires setting the pre\sphinxhyphen{}processor flag
\sphinxhyphen{}DLMP\_INTEL\_USELRT in the makefile when compiling LAMMPS.

\sphinxAtStartPar
The \sphinxstyleemphasis{balance} keyword sets the fraction of {\hyperref[\detokenize{pair_style::doc}]{\sphinxcrossref{\DUrole{doc}{pair style}}}} work
offloaded to the co\sphinxhyphen{}processor for split values between 0.0 and 1.0 inclusive.
While this fraction of work is running on the co\sphinxhyphen{}processor, other calculations
will run on the host, including neighbor and pair calculations that are not
offloaded, as well as angle, bond, dihedral, kspace, and some MPI
communications.  If \sphinxstyleemphasis{split} is set to \sphinxhyphen{}1, the fraction of work is dynamically
adjusted automatically throughout the run.  This typically give performance
within 5 to 10 percent of the optimal fixed fraction.

\sphinxAtStartPar
The \sphinxstyleemphasis{ghost} keyword determines whether or not ghost atoms, i.e. atoms
at the boundaries of processor subdomains, are offloaded for neighbor
and force calculations.  When the value = “no”, ghost atoms are not
offloaded.  This option can reduce the amount of data transfer with
the co\sphinxhyphen{}processor and can also overlap MPI communication of forces with
computation on the co\sphinxhyphen{}processor when the {\hyperref[\detokenize{newton::doc}]{\sphinxcrossref{\DUrole{doc}{newton pair}}}}
setting is “on”.  When the value = “yes”, ghost atoms are offloaded.
In some cases this can provide better performance, especially if the
\sphinxstyleemphasis{balance} fraction is high.

\sphinxAtStartPar
The \sphinxstyleemphasis{tpc} keyword sets the max \# of co\sphinxhyphen{}processor threads \sphinxstyleemphasis{Ntpc} that
will run on each core of the co\sphinxhyphen{}processor.  The default value = 4,
which is the number of hardware threads per core supported by the
current generation Xeon Phi chips.

\sphinxAtStartPar
The \sphinxstyleemphasis{tptask} keyword sets the max \# of co\sphinxhyphen{}processor threads (Ntptask*
assigned to each MPI task.  The default value = 240, which is the
total \# of threads an entire current generation Xeon Phi chip can run
(240 = 60 cores * 4 threads/core).  This means each MPI task assigned
to the Phi will enough threads for the chip to run the max allowed,
even if only 1 MPI task is assigned.  If 8 MPI tasks are assigned to
the Phi, each will run with 30 threads.  If you wish to limit the
number of threads per MPI task, set \sphinxstyleemphasis{tptask} to a smaller value.
E.g. for \sphinxstyleemphasis{tptask} = 16, if 8 MPI tasks are assigned, each will run
with 16 threads, for a total of 128.

\sphinxAtStartPar
Note that the default settings for \sphinxstyleemphasis{tpc} and \sphinxstyleemphasis{tptask} are fine for
most problems, regardless of how many MPI tasks you assign to a Phi.

\sphinxAtStartPar
\DUrole{versionmodified,added}{Added in version 15Jun2023.}

\sphinxAtStartPar
The \sphinxstyleemphasis{pppm\_table} keyword with the argument yes allows to use a
pre\sphinxhyphen{}computed table to efficiently spread the charge to the PPPM grid.
This feature is enabled by default but can be turned off using the
keyword with the argument \sphinxstyleemphasis{no}.

\sphinxAtStartPar
The \sphinxstyleemphasis{no\_affinity} keyword will turn off automatic setting of core
affinity for MPI tasks and OpenMP threads on the host when using
offload to a co\sphinxhyphen{}processor. Affinity settings are used when possible
to prevent MPI tasks and OpenMP threads from being on separate NUMA
domains and to prevent offload threads from interfering with other
processes/threads used for LAMMPS.


\bigskip\hrule\bigskip


\sphinxAtStartPar
The \sphinxstyleemphasis{kokkos} style invokes settings associated with the use of the
KOKKOS package.

\sphinxAtStartPar
All of the settings are optional keyword/value pairs. Each has a default
value as listed below.

\sphinxAtStartPar
The \sphinxstyleemphasis{neigh} keyword determines how neighbor lists are built. A value of
\sphinxstyleemphasis{half} uses a thread\sphinxhyphen{}safe variant of half\sphinxhyphen{}neighbor lists, the same as
used by most pair styles in LAMMPS, which is the default when running on
CPUs (i.e. the Kokkos CUDA back end is not enabled).

\sphinxAtStartPar
A value of \sphinxstyleemphasis{full} uses a full neighbor lists and is the default when
running on GPUs. This performs twice as much computation as the \sphinxstyleemphasis{half}
option, however that is often a win because it is thread\sphinxhyphen{}safe and
does not require atomic operations in the calculation of pair forces. For
that reason, \sphinxstyleemphasis{full} is the default setting for GPUs. However, when
running on CPUs, a \sphinxstyleemphasis{half} neighbor list is the default because it are
often faster, just as it is for non\sphinxhyphen{}accelerated pair styles. Similarly,
the \sphinxstyleemphasis{neigh/qeq} keyword determines how neighbor lists are built for
{\hyperref[\detokenize{fix_qeq_reaxff::doc}]{\sphinxcrossref{\DUrole{doc}{fix qeq/reaxff/kk}}}}.

\sphinxAtStartPar
If the \sphinxstyleemphasis{neigh/thread} keyword is set to \sphinxstyleemphasis{off}, then the KOKKOS package
threads only over atoms. However, for small systems, this may not expose
enough parallelism to keep a GPU busy. When this keyword is set to \sphinxstyleemphasis{on},
the KOKKOS package threads over both atoms and neighbors of atoms. When
using \sphinxstyleemphasis{neigh/thread} \sphinxstyleemphasis{on}, the {\hyperref[\detokenize{newton::doc}]{\sphinxcrossref{\DUrole{doc}{newton pair}}}} setting must
be “off”. Using \sphinxstyleemphasis{neigh/thread} \sphinxstyleemphasis{on} may be slower for large systems, so
this this option is turned on by default only when running on one or
more GPUs and there are 16k atoms or less owned by an MPI rank. Not all
KOKKOS\sphinxhyphen{}enabled potentials support this keyword yet, and only thread over
atoms. Many simple pairwise potentials such as Lennard\sphinxhyphen{}Jones do support
threading over both atoms and neighbors.

\sphinxAtStartPar
If the \sphinxstyleemphasis{neigh/transpose} keyword is set to \sphinxstyleemphasis{off}, then the KOKKOS
package will use the same memory layout for building the neighbor list on
GPUs as used for the pair style. When this keyword is set to \sphinxstyleemphasis{on} it
will use a different (transposed) memory layout to build the neighbor
list on GPUs. This can be faster in some cases (e.g. ReaxFF HNS
benchmark) but slower in others (e.g. Lennard Jones benchmark). The
copy between different memory layouts is done out of place and
therefore doubles the memory overhead of the neighbor list, which can
be significant.

\sphinxAtStartPar
The \sphinxstyleemphasis{newton} keyword sets the Newton flags for pairwise and bonded
interactions to \sphinxstyleemphasis{off} or \sphinxstyleemphasis{on}, the same as the {\hyperref[\detokenize{newton::doc}]{\sphinxcrossref{\DUrole{doc}{newton}}}}
command allows. The default for GPUs is \sphinxstyleemphasis{off} because this will almost
always give better performance for the KOKKOS package. This means more
computation is done, but less communication. However, when running on
CPUs a value of \sphinxstyleemphasis{on} is the default since it can often be faster, just
as it is for non\sphinxhyphen{}accelerated pair styles

\sphinxAtStartPar
The \sphinxstyleemphasis{binsize} keyword sets the size of bins used to bin atoms during
neighbor list builds. The same value can be set by the
{\hyperref[\detokenize{neigh_modify::doc}]{\sphinxcrossref{\DUrole{doc}{neigh\_modify binsize}}}} command. Making it an option
in the package kokkos command allows it to be set from the command line.
The default value for CPUs is 0.0, which means the LAMMPS default will be
used, which is bins = 1/2 the size of the pairwise cutoff + neighbor skin
distance. This is fine when neighbor lists are built on the CPU. For GPU
builds, a 2x larger binsize equal to the pairwise cutoff + neighbor skin
is often faster, which is the default. Note that if you use a
longer\sphinxhyphen{}than\sphinxhyphen{}usual pairwise cutoff, e.g. to allow for a smaller fraction
of KSpace work with a {\hyperref[\detokenize{kspace_style::doc}]{\sphinxcrossref{\DUrole{doc}{long\sphinxhyphen{}range Coulombic solver}}}}
because the GPU is faster at performing pairwise interactions, then this
rule of thumb may give too large a binsize and the default should be
overridden with a smaller value.

\sphinxAtStartPar
The \sphinxstyleemphasis{comm} and \sphinxstyleemphasis{comm/exchange} and \sphinxstyleemphasis{comm/forward} and \sphinxstyleemphasis{comm/pair/forward}
and \sphinxstyleemphasis{comm/fix/forward} and \sphinxstyleemphasis{comm/reverse} and \sphinxstyleemphasis{comm/pair/reverse}
keywords determine whether the host or device performs the packing and
unpacking of data when communicating per\sphinxhyphen{}atom data between processors.
“Exchange” communication happens only on timesteps that neighbor lists
are rebuilt. The data is only for atoms that migrate to new processors.
“Forward” communication happens every timestep. “Reverse” communication
happens every timestep if the \sphinxstyleemphasis{newton} option is on. The data is for
atom coordinates and any other atom properties that needs to be updated
for ghost atoms owned by each processor. “Pair/comm” controls additional
communication in pair styles, such as pair\_style EAM. “Fix/comm” controls
additional communication in fixes, such as fix SHAKE.

\sphinxAtStartPar
The \sphinxstyleemphasis{comm} keyword is simply a short\sphinxhyphen{}cut to set the same value for all
the comm keywords.

\sphinxAtStartPar
The value options for the keywords are \sphinxstyleemphasis{no} or \sphinxstyleemphasis{host} or \sphinxstyleemphasis{device}. A
value of \sphinxstyleemphasis{no} means to use the standard non\sphinxhyphen{}KOKKOS method of
packing/unpacking data for the communication. A value of \sphinxstyleemphasis{host} means to
use the host, typically a multicore CPU, and perform the
packing/unpacking in parallel with threads. A value of \sphinxstyleemphasis{device} means to
use the device, typically a GPU, to perform the packing/unpacking
operation.

\sphinxAtStartPar
For the \sphinxstyleemphasis{comm/pair/forward} or \sphinxstyleemphasis{comm/fix/forward} or \sphinxstyleemphasis{comm/pair/reverse}
keywords, if a value of \sphinxstyleemphasis{host} is used it will be automatically
be changed to \sphinxstyleemphasis{no} since these keywords don’t support \sphinxstyleemphasis{host} mode. The
value of \sphinxstyleemphasis{no} will also always be used when running on the CPU, i.e. setting
the value to \sphinxstyleemphasis{device} will have no effect if the pair/fix style is
running on the CPU. For the \sphinxstyleemphasis{comm/fix/forward} or \sphinxstyleemphasis{comm/pair/reverse}
keywords, not all styles support \sphinxstyleemphasis{device} mode and in that case will run
in \sphinxstyleemphasis{no} mode instead.

\sphinxAtStartPar
The optimal choice for these keywords depends on the input script and
the hardware used. The \sphinxstyleemphasis{no} value is useful for verifying that the
Kokkos\sphinxhyphen{}based \sphinxstyleemphasis{host} and \sphinxstyleemphasis{device} values are working correctly. It is the
default when running on CPUs since it is usually the fastest.

\sphinxAtStartPar
When running on CPUs or Xeon Phi, the \sphinxstyleemphasis{host} and \sphinxstyleemphasis{device} values work
identically. When using GPUs, the \sphinxstyleemphasis{device} value is the default since it
will typically be optimal if all of your styles used in your input
script are supported by the KOKKOS package. In this case data can stay
on the GPU for many timesteps without being moved between the host and
GPU, if you use the \sphinxstyleemphasis{device} value. If your script uses styles (e.g.
fixes) which are not yet supported by the KOKKOS package, then data has
to be moved between the host and device anyway, so it is typically faster
to let the host handle communication, by using the \sphinxstyleemphasis{host} value. Using
\sphinxstyleemphasis{host} instead of \sphinxstyleemphasis{no} will enable use of multiple threads to
pack/unpack communicated data. When running small systems on a GPU,
performing the exchange pack/unpack on the host CPU can give speedup
since it reduces the number of CUDA kernel launches.

\sphinxAtStartPar
The \sphinxstyleemphasis{sort} keyword determines whether the host or device performs atom
sorting, see the {\hyperref[\detokenize{atom_modify::doc}]{\sphinxcrossref{\DUrole{doc}{atom\_modify sort}}}} command.  The value
options for the \sphinxstyleemphasis{sort} keyword are \sphinxstyleemphasis{no} or \sphinxstyleemphasis{device} similar to the \sphinxstyleemphasis{comm}
keywords above. If a value of \sphinxstyleemphasis{host} is used it will be automatically be
changed to \sphinxstyleemphasis{no} since the \sphinxstyleemphasis{sort} keyword does not support \sphinxstyleemphasis{host} mode. Not
all fix styles with extra atom data support \sphinxstyleemphasis{device} mode and in that case
a warning will be given and atom sorting will run in \sphinxstyleemphasis{no} mode instead.

\sphinxAtStartPar
\DUrole{versionmodified,added}{Added in version 17Apr2024.}

\sphinxAtStartPar
The \sphinxstyleemphasis{atom/map} keyword determines whether the host or device builds the
atom\_map, see the {\hyperref[\detokenize{atom_modify::doc}]{\sphinxcrossref{\DUrole{doc}{atom\_modify map}}}} command.  The
value options for the \sphinxstyleemphasis{atom/map} keyword are identical to the \sphinxstyleemphasis{sort}
keyword above.

\sphinxAtStartPar
The \sphinxstyleemphasis{gpu/aware} keyword chooses whether GPU\sphinxhyphen{}aware MPI will be used. When
this keyword is set to \sphinxstyleemphasis{on}, buffers in GPU memory are passed directly
through MPI send/receive calls. This reduces overhead of first copying
the data to the host CPU. However GPU\sphinxhyphen{}aware MPI is not supported on all
systems, which can lead to segmentation faults and would require using a
value of \sphinxstyleemphasis{off}. If LAMMPS can safely detect that GPU\sphinxhyphen{}aware MPI is not
available (currently only possible with OpenMPI v2.0.0 or later), then
the \sphinxstyleemphasis{gpu/aware} keyword is automatically set to \sphinxstyleemphasis{off} by default. When
the \sphinxstyleemphasis{gpu/aware} keyword is set to \sphinxstyleemphasis{off} while any of the \sphinxstyleemphasis{comm}
keywords are set to \sphinxstyleemphasis{device}, the value for these \sphinxstyleemphasis{comm} keywords will
be automatically changed to \sphinxstyleemphasis{no}. This setting has no effect if not
running on GPUs or if using only one MPI rank. GPU\sphinxhyphen{}aware MPI is available
for OpenMPI 1.8 (or later versions), Mvapich2 1.9 (or later) when the
“MV2\_USE\_CUDA” environment variable is set to “1”, CrayMPI, and IBM
Spectrum MPI when the “\sphinxhyphen{}gpu” flag is used.

\sphinxAtStartPar
The \sphinxstyleemphasis{pair/only} keyword can change how the KOKKOS suffix “kk” is applied
when using an accelerator device.  By default device acceleration is always
used for all available styles.  With \sphinxstyleemphasis{pair/only} set to \sphinxstyleemphasis{on} the suffix
setting will choose device acceleration only for pair styles and run all
other force computations on the host CPU.  The \sphinxstyleemphasis{comm} flags, along with the
\sphinxstyleemphasis{sort} and \sphinxstyleemphasis{atom/map} keywords will also automatically be changed to \sphinxstyleemphasis{no}.
This can result in better performance for certain configurations and
system sizes.


\bigskip\hrule\bigskip


\sphinxAtStartPar
The \sphinxstyleemphasis{omp} style invokes settings associated with the use of the
OPENMP package.

\sphinxAtStartPar
The \sphinxstyleemphasis{Nthreads} argument sets the number of OpenMP threads allocated for
each MPI task.  For example, if your system has nodes with dual
quad\sphinxhyphen{}core processors, it has a total of 8 cores per node.  You could
use two MPI tasks per node (e.g. using the \sphinxhyphen{}ppn option of the mpirun
command in MPICH or \sphinxhyphen{}npernode in OpenMPI), and set \sphinxstyleemphasis{Nthreads} = 4.
This would use all 8 cores on each node.  Note that the product of MPI
tasks * threads/task should not exceed the physical number of cores
(on a node), otherwise performance will suffer.

\sphinxAtStartPar
Setting \sphinxstyleemphasis{Nthreads} = 0 instructs LAMMPS to use whatever value is the
default for the given OpenMP environment. This is usually determined
via the \sphinxstyleemphasis{OMP\_NUM\_THREADS} environment variable or the compiler
runtime.  Note that in most cases the default for OpenMP capable
compilers is to use one thread for each available CPU core when
\sphinxstyleemphasis{OMP\_NUM\_THREADS} is not explicitly set, which can lead to poor
performance.

\sphinxAtStartPar
Here are examples of how to set the environment variable when
launching LAMMPS:

\begin{sphinxVerbatim}[commandchars=\\\{\}]
env\PYG{+w}{ }\PYG{n+nv}{OMP\PYGZus{}NUM\PYGZus{}THREADS}\PYG{o}{=}\PYG{l+m}{4}\PYG{+w}{ }lmp\PYGZus{}machine\PYG{+w}{ }\PYGZhy{}sf\PYG{+w}{ }omp\PYG{+w}{ }\PYGZhy{}in\PYG{+w}{ }\PYG{k}{in}.script
env\PYG{+w}{ }\PYG{n+nv}{OMP\PYGZus{}NUM\PYGZus{}THREADS}\PYG{o}{=}\PYG{l+m}{2}\PYG{+w}{ }mpirun\PYG{+w}{ }\PYGZhy{}np\PYG{+w}{ }\PYG{l+m}{2}\PYG{+w}{ }lmp\PYGZus{}machine\PYG{+w}{ }\PYGZhy{}sf\PYG{+w}{ }omp\PYG{+w}{ }\PYGZhy{}in\PYG{+w}{ }\PYG{k}{in}.script
mpirun\PYG{+w}{ }\PYGZhy{}x\PYG{+w}{ }\PYG{n+nv}{OMP\PYGZus{}NUM\PYGZus{}THREADS}\PYG{o}{=}\PYG{l+m}{2}\PYG{+w}{ }\PYGZhy{}np\PYG{+w}{ }\PYG{l+m}{2}\PYG{+w}{ }lmp\PYGZus{}machine\PYG{+w}{ }\PYGZhy{}sf\PYG{+w}{ }omp\PYG{+w}{ }\PYGZhy{}in\PYG{+w}{ }\PYG{k}{in}.script
\end{sphinxVerbatim}

\sphinxAtStartPar
or you can set it permanently in your shell’s start\sphinxhyphen{}up script.
All three of these examples use a total of 4 CPU cores.

\sphinxAtStartPar
Note that different MPI implementations have different ways of passing
the OMP\_NUM\_THREADS environment variable to all MPI processes.  The
second example line above is for MPICH; the third example line with \sphinxhyphen{}x is
for OpenMPI.  Check your MPI documentation for additional details.

\sphinxAtStartPar
What combination of threads and MPI tasks gives the best performance
is difficult to predict and can depend on many components of your
input.  Not all features of LAMMPS support OpenMP threading via the
OPENMP package and the parallel efficiency can be very different,
too.

\begin{sphinxadmonition}{note}{Note:}
\sphinxAtStartPar
If you build LAMMPS with the GPU, INTEL, and / or OPENMP
packages, be aware these packages all allow setting of the \sphinxstyleemphasis{Nthreads}
value via their package commands, but there is only a single global
\sphinxstyleemphasis{Nthreads} value used by OpenMP.  Thus if multiple package commands are
invoked, you should ensure the values are consistent.  If they are
not, the last one invoked will take precedence, for all packages.
Also note that if the {\hyperref[\detokenize{Run_options::doc}]{\sphinxcrossref{\DUrole{doc}{\sphinxhyphen{}sf hybrid intel omp command\sphinxhyphen{}line switch}}}} is used, it invokes a “package intel” command, followed by a
“package omp” command, both with a setting of \sphinxstyleemphasis{Nthreads} = 0. Likewise
for a hybrid suffix for gpu and omp. Note that KOKKOS also supports
setting the number of OpenMP threads from the command line using the
“\sphinxhyphen{}k on” {\hyperref[\detokenize{Run_options::doc}]{\sphinxcrossref{\DUrole{doc}{command\sphinxhyphen{}line switch}}}}. The default for
KOKKOS is 1 thread per MPI task, so any other number of threads should
be explicitly set using the “\sphinxhyphen{}k on” command\sphinxhyphen{}line switch (and this
setting should be consistent with settings from any other packages
used).
\end{sphinxadmonition}

\sphinxAtStartPar
Optional keyword/value pairs can also be specified.  Each has a
default value as listed below.

\sphinxAtStartPar
The \sphinxstyleemphasis{neigh} keyword specifies whether neighbor list building will be
multi\sphinxhyphen{}threaded in addition to force calculations.  If \sphinxstyleemphasis{neigh} is set
to \sphinxstyleemphasis{no} then neighbor list calculation is performed only by MPI tasks
with no OpenMP threading.  If \sphinxstyleemphasis{mode} is \sphinxstyleemphasis{yes} (the default), a
multi\sphinxhyphen{}threaded neighbor list build is used.  Using \sphinxstyleemphasis{neigh} = \sphinxstyleemphasis{yes} is
almost always faster and should produce identical neighbor lists at the
expense of using more memory.  Specifically, neighbor list pages are
allocated for all threads at the same time and each thread works
within its own pages.


\bigskip\hrule\bigskip



\subsection{Restrictions}
\label{\detokenize{package:restrictions}}
\sphinxAtStartPar
This command cannot be used after the simulation box is defined by a
{\hyperref[\detokenize{read_data::doc}]{\sphinxcrossref{\DUrole{doc}{read\_data}}}} or {\hyperref[\detokenize{create_box::doc}]{\sphinxcrossref{\DUrole{doc}{create\_box}}}} command.

\sphinxAtStartPar
The \sphinxstyleemphasis{gpu} style of this command can only be invoked if LAMMPS was built
with the GPU package.  See the {\hyperref[\detokenize{Build_package::doc}]{\sphinxcrossref{\DUrole{doc}{Build package}}}} doc
page for more info.

\sphinxAtStartPar
The \sphinxstyleemphasis{intel} style of this command can only be invoked if LAMMPS was
built with the INTEL package.  See the {\hyperref[\detokenize{Build_package::doc}]{\sphinxcrossref{\DUrole{doc}{Build package}}}} page for more info.

\sphinxAtStartPar
The \sphinxstyleemphasis{kokkos} style of this command can only be invoked if LAMMPS was built
with the KOKKOS package.  See the {\hyperref[\detokenize{Build_package::doc}]{\sphinxcrossref{\DUrole{doc}{Build package}}}}
doc page for more info.

\sphinxAtStartPar
The \sphinxstyleemphasis{omp} style of this command can only be invoked if LAMMPS was built
with the OPENMP package.  See the {\hyperref[\detokenize{Build_package::doc}]{\sphinxcrossref{\DUrole{doc}{Build package}}}}
doc page for more info.


\subsection{Related commands}
\label{\detokenize{package:related-commands}}
\sphinxAtStartPar
{\hyperref[\detokenize{suffix::doc}]{\sphinxcrossref{\DUrole{doc}{suffix}}}}, {\hyperref[\detokenize{Run_options::doc}]{\sphinxcrossref{\DUrole{doc}{\sphinxhyphen{}pk command\sphinxhyphen{}line switch}}}}


\subsection{Defaults}
\label{\detokenize{package:defaults}}
\sphinxAtStartPar
For the GPU package, the default parameters and settings are:

\begin{sphinxVerbatim}[commandchars=\\\{\}]
Ngpu = 0, neigh = yes, newton = off, binsize = 0.0, split = 1.0, gpuID = 0 to Ngpu\PYGZhy{}1, tpa = 1, omp = 0, platform=\PYGZhy{}1.
\end{sphinxVerbatim}

\sphinxAtStartPar
These settings are made automatically if the “\sphinxhyphen{}sf gpu”
{\hyperref[\detokenize{Run_options::doc}]{\sphinxcrossref{\DUrole{doc}{command\sphinxhyphen{}line switch}}}} is used.  If it is not used,
you must invoke the package gpu command in your input script or via the
“\sphinxhyphen{}pk gpu” {\hyperref[\detokenize{Run_options::doc}]{\sphinxcrossref{\DUrole{doc}{command\sphinxhyphen{}line switch}}}}.

\sphinxAtStartPar
For the INTEL package, the default parameters and settings are:

\begin{sphinxVerbatim}[commandchars=\\\{\}]
Nphi = 1, omp = 0, mode = mixed, lrt = no, balance = \PYGZhy{}1, tpc = 4, tptask = 240, pppm\PYGZus{}table = yes
\end{sphinxVerbatim}

\sphinxAtStartPar
The default ghost option is determined by the pair style being used.
This value is output to the screen in the offload report at the end of each
run.  Note that all of these settings, except “omp” and “mode”, are ignored if
LAMMPS was not built with Xeon Phi co\sphinxhyphen{}processor support.  These settings are
made automatically if the “\sphinxhyphen{}sf intel” {\hyperref[\detokenize{Run_options::doc}]{\sphinxcrossref{\DUrole{doc}{command\sphinxhyphen{}line switch}}}}
is used.  If it is not used, you must invoke the package intel command in your
input script or via the “\sphinxhyphen{}pk intel” {\hyperref[\detokenize{Run_options::doc}]{\sphinxcrossref{\DUrole{doc}{command\sphinxhyphen{}line switch}}}}.

\sphinxAtStartPar
For the KOKKOS package when using GPUs, the option defaults are:

\begin{sphinxVerbatim}[commandchars=\\\{\}]
neigh = full, neigh/qeq = full, newton = off, binsize = 2x LAMMPS default value, comm = device, sort = device, atom/map = device, neigh/transpose = off, gpu/aware = on
\end{sphinxVerbatim}

\sphinxAtStartPar
For GPUs, option neigh/thread = on when there are 16k atoms or less on
an MPI rank, otherwise it is “off”. When LAMMPS can safely detect that
GPU\sphinxhyphen{}aware MPI is not available, the default value of gpu/aware becomes
“off”.

\sphinxAtStartPar
For the KOKKOS package when using CPUs or Xeon Phis, the option defaults are:

\begin{sphinxVerbatim}[commandchars=\\\{\}]
neigh = half, neigh/qeq = half, newton = on, binsize = 0.0, comm = no, sort = no, atom/map = no
\end{sphinxVerbatim}

\sphinxAtStartPar
These settings are made automatically by
the required “\sphinxhyphen{}k on” {\hyperref[\detokenize{Run_options::doc}]{\sphinxcrossref{\DUrole{doc}{command\sphinxhyphen{}line switch}}}}.  You can
change them by using the package kokkos command in your input script or
via the {\hyperref[\detokenize{Run_options::doc}]{\sphinxcrossref{\DUrole{doc}{\sphinxhyphen{}pk kokkos command\sphinxhyphen{}line switch}}}}.

\sphinxAtStartPar
For the OMP package, the defaults are

\begin{sphinxVerbatim}[commandchars=\\\{\}]
Nthreads = 0, neigh = yes
\end{sphinxVerbatim}

\sphinxAtStartPar
These settings are made automatically if the “\sphinxhyphen{}sf omp”
{\hyperref[\detokenize{Run_options::doc}]{\sphinxcrossref{\DUrole{doc}{command\sphinxhyphen{}line switch}}}} is used.  If it is not used,
you must invoke the package omp command in your input script or via the
“\sphinxhyphen{}pk omp” {\hyperref[\detokenize{Run_options::doc}]{\sphinxcrossref{\DUrole{doc}{command\sphinxhyphen{}line switch}}}}.

\sphinxstepscope

\index{pair\_coeff@\spxentry{pair\_coeff}}\ignorespaces 

\section{pair\_coeff command}
\label{\detokenize{pair_coeff:pair-coeff-command}}\label{\detokenize{pair_coeff:index-0}}\label{\detokenize{pair_coeff::doc}}

\subsection{Syntax}
\label{\detokenize{pair_coeff:syntax}}
\begin{sphinxVerbatim}[commandchars=\\\{\}]
\PYG{k}{pair\PYGZus{}coeff}\PYG{+w}{ }\PYG{n}{I}\PYG{+w}{ }\PYG{n}{J}\PYG{+w}{ }\PYG{n}{args}
\end{sphinxVerbatim}
\begin{itemize}
\item {} 
\sphinxAtStartPar
I,J = numeric atom types (see asterisk form below), or type labels

\item {} 
\sphinxAtStartPar
args = coefficients for one or more pairs of atom types

\end{itemize}


\subsection{Examples}
\label{\detokenize{pair_coeff:examples}}
\begin{sphinxVerbatim}[commandchars=\\\{\}]
\PYG{k}{pair\PYGZus{}coeff}\PYG{+w}{ }\PYG{l+m}{1}\PYG{+w}{ }\PYG{l+m}{2}\PYG{+w}{ }\PYG{l+m}{1.0}\PYG{+w}{ }\PYG{l+m}{1.0}\PYG{+w}{ }\PYG{l+m}{2.5}
\PYG{k}{pair\PYGZus{}coeff}\PYG{+w}{ }\PYG{l+m}{2}\PYG{+w}{ }\PYG{o}{*}\PYG{+w}{ }\PYG{l+m}{1.0}\PYG{+w}{ }\PYG{l+m}{1.0}
\PYG{k}{pair\PYGZus{}coeff}\PYG{+w}{ }\PYG{l+m}{3}\PYG{o}{*}\PYG{+w}{ }\PYG{l+m}{1}\PYG{o}{*}\PYG{l+m}{2}\PYG{+w}{ }\PYG{l+m}{1.0}\PYG{+w}{ }\PYG{l+m}{1.0}\PYG{+w}{ }\PYG{l+m}{2.5}
\PYG{k}{pair\PYGZus{}coeff}\PYG{+w}{ }\PYG{o}{*}\PYG{+w}{ }\PYG{o}{*}\PYG{+w}{ }\PYG{l+m}{1.0}\PYG{+w}{ }\PYG{l+m}{1.0}
\PYG{k}{pair\PYGZus{}coeff}\PYG{+w}{ }\PYG{o}{*}\PYG{+w}{ }\PYG{o}{*}\PYG{+w}{ }\PYG{n}{nialhjea}\PYG{+w}{ }\PYG{l+m}{1}\PYG{+w}{ }\PYG{l+m}{1}\PYG{+w}{ }\PYG{l+m}{2}
\PYG{k}{pair\PYGZus{}coeff}\PYG{+w}{ }\PYG{o}{*}\PYG{+w}{ }\PYG{l+m}{3}\PYG{+w}{ }\PYG{n}{morse.table}\PYG{+w}{ }\PYG{n}{ENTRY1}
\PYG{k}{pair\PYGZus{}coeff}\PYG{+w}{ }\PYG{l+m}{1}\PYG{+w}{ }\PYG{l+m}{2}\PYG{+w}{ }\PYG{n}{lj}\PYG{o}{/}\PYG{n}{cut}\PYG{+w}{ }\PYG{l+m}{1.0}\PYG{+w}{ }\PYG{l+m}{1.0}\PYG{+w}{ }\PYG{l+m}{2.5}\PYG{+w}{ }\PYG{c}{\PYGZsh{} (for pair\PYGZus{}style hybrid)}

\PYG{k}{labelmap}\PYG{+w}{ }\PYG{n}{atom}\PYG{+w}{ }\PYG{l+m}{1}\PYG{+w}{ }\PYG{n}{C}
\PYG{k}{labelmap}\PYG{+w}{ }\PYG{n}{atom}\PYG{+w}{ }\PYG{l+m}{2}\PYG{+w}{ }\PYG{n}{H}
\PYG{k}{pair\PYGZus{}coeff}\PYG{+w}{ }\PYG{n}{C}\PYG{+w}{ }\PYG{n}{H}\PYG{+w}{ }\PYG{l+m}{1.0}\PYG{+w}{ }\PYG{l+m}{1.0}\PYG{+w}{ }\PYG{l+m}{2.5}
\end{sphinxVerbatim}


\subsection{Description}
\label{\detokenize{pair_coeff:description}}
\sphinxAtStartPar
Specify the pairwise force field coefficients for one or more pairs of
atom types.  The number and meaning of the coefficients depends on the
pair style.  Pair coefficients can also be set in the data file read
by the {\hyperref[\detokenize{read_data::doc}]{\sphinxcrossref{\DUrole{doc}{read\_data}}}} command or in a restart file.

\sphinxAtStartPar
I and J can be specified in one of several ways.  Explicit numeric
values can be used for each, as in the first example above.  Or, one
or both of the types in the I,J pair can be a type label, which is an
alphanumeric string defined by the {\hyperref[\detokenize{labelmap::doc}]{\sphinxcrossref{\DUrole{doc}{labelmap}}}} command
or in a section of a data file read by the {\hyperref[\detokenize{read_data::doc}]{\sphinxcrossref{\DUrole{doc}{read\_data}}}} command, and which converts internally to a numeric type.
Internally, LAMMPS will set coefficients for the symmetric J,I
interaction to the same values as the I,J interaction.

\sphinxAtStartPar
For numeric values only, a wildcard asterisk can be used in place of or
in conjunction with the I,J arguments to set the coefficients for
multiple pairs of atom types.  This takes the form “*” or “*n” or
“n*” or “m*n”.  If \(N\) is the number of atom types, then an
asterisk with no numeric values means all types from 1 to \(N\).  A
leading asterisk means all types from 1 to n (inclusive).  A trailing
asterisk means all types from n to \(N\) (inclusive).  A middle
asterisk means all types from m to n (inclusive).  For the asterisk
syntax, only type pairs with I \textless{}= J are considered; if asterisks imply
type pairs where J \textless{} I, they are ignored. Again internally, LAMMPS will
set the coefficients for the symmetric J,I interactions to the same
values as the I \textless{}= J interactions.

\sphinxAtStartPar
Note that a pair\_coeff command can override a previous setting for the
same I,J pair.  For example, these commands set the coeffs for all I,J
pairs, then overwrite the coeffs for just the I,J = 2,3 pair:

\begin{sphinxVerbatim}[commandchars=\\\{\}]
\PYG{k}{pair\PYGZus{}coeff}\PYG{+w}{ }\PYG{o}{*}\PYG{+w}{ }\PYG{o}{*}\PYG{+w}{ }\PYG{l+m}{1.0}\PYG{+w}{ }\PYG{l+m}{1.0}\PYG{+w}{ }\PYG{l+m}{2.5}
\PYG{k}{pair\PYGZus{}coeff}\PYG{+w}{ }\PYG{l+m}{2}\PYG{+w}{ }\PYG{l+m}{3}\PYG{+w}{ }\PYG{l+m}{2.0}\PYG{+w}{ }\PYG{l+m}{1.0}\PYG{+w}{ }\PYG{l+m}{1.12}
\end{sphinxVerbatim}

\sphinxAtStartPar
A line in a data file that specifies pair coefficients uses the exact
same format as the arguments of the pair\_coeff command in an input
script, with the exception of the I,J type arguments.  In each line of
the “Pair Coeffs” section of a data file, only a single type I is
specified, which sets the coefficients for type I interacting with
type I.  This is because the section has exactly \(N\) lines, where
\(N\) is the number of atom types.  For this reason, the wild\sphinxhyphen{}card
asterisk should also not be used as part of the I argument.  Thus in a
data file, the line corresponding to the first example above would be
listed as

\begin{sphinxVerbatim}[commandchars=\\\{\}]
2 1.0 1.0 2.5
\end{sphinxVerbatim}

\sphinxAtStartPar
For many potentials, if coefficients for type pairs with I != J are
not set explicitly by a pair\_coeff command, the values are inferred
from the I,I and J,J settings by mixing rules; see the
{\hyperref[\detokenize{pair_modify::doc}]{\sphinxcrossref{\DUrole{doc}{pair\_modify}}}} command for a discussion.  Details on
this option as it pertains to individual potentials are described on
the page for the potential.

\sphinxAtStartPar
Many pair styles, typically for many\sphinxhyphen{}body potentials, use tabulated
potential files as input, when specifying the pair\_coeff command.
Potential files provided with LAMMPS are in the potentials directory
of the distribution.  For some potentials, such as EAM, other archives
of suitable files can be found on the Web.  They can be used with
LAMMPS so long as they are in the format LAMMPS expects, as discussed
on the individual doc pages.  The first line of potential files may
contain metadata with upper case tags followed their value. These may
be parsed and used by LAMMPS.  Currently supported are the “DATE:”
tag and the \sphinxcode{\sphinxupquote{UNITS:}} tag.  For pair styles that have been programmed
to support the metadata, the value of the “DATE:” tag is printed to
the screen and logfile so that the version of a potential file can be
later identified.  The \sphinxcode{\sphinxupquote{UNITS:}} tag indicates the {\hyperref[\detokenize{units::doc}]{\sphinxcrossref{\DUrole{doc}{units}}}}
setting required for this particular potential file.  If the potential
file was created for a different sets of units, LAMMPS will terminate
with an error.  If the potential file does not contain the tag, no
check will be made and it is the responsibility of the user to determine
that the unit style is correct.

\sphinxAtStartPar
In some select cases and for specific combinations of unit styles,
LAMMPS is capable of automatically converting potential parameters
from a file. In those cases, a warning message signaling that an
automatic conversion has happened is printed to the screen.

\sphinxAtStartPar
When a pair\_coeff command using a potential file is specified, LAMMPS
looks for the potential file in 2 places.  First it looks in the
location specified.  E.g. if the file is specified as “niu3.eam”, it
is looked for in the current working directory.  If it is specified as
“../potentials/niu3.eam”, then it is looked for in the potentials
directory, assuming it is a sister directory of the current working
directory.  If the file is not found, it is then looked for in one of
the directories specified by the \sphinxcode{\sphinxupquote{LAMMPS\_POTENTIALS}} environment variable.
Thus if this is set to the potentials directory in the LAMMPS distribution,
then you can use those files from anywhere on your system, without
copying them into your working directory.  Environment variables are
set in different ways for different shells.  Here are example settings
for

\sphinxAtStartPar
csh, tcsh:

\begin{sphinxVerbatim}[commandchars=\\\{\}]
\PYG{n+nb}{setenv }LAMMPS\PYGZus{}POTENTIALS /path/to/lammps/potentials
\end{sphinxVerbatim}

\sphinxAtStartPar
bash:

\begin{sphinxVerbatim}[commandchars=\\\{\}]
\PYG{n+nb}{export}\PYG{+w}{ }\PYG{n+nv}{LAMMPS\PYGZus{}POTENTIALS}\PYG{o}{=}/path/to/lammps/potentials
\end{sphinxVerbatim}

\sphinxAtStartPar
Windows:

\begin{sphinxVerbatim}[commandchars=\\\{\}]
\PYG{g+go}{set LAMMPS\PYGZus{}POTENTIALS=\PYGZdq{}C:\PYGZbs{}\PYGZbs{}Path to LAMMPS\PYGZbs{}\PYGZbs{}Potentials\PYGZdq{}}
\end{sphinxVerbatim}

\sphinxAtStartPar
The \sphinxcode{\sphinxupquote{LAMMPS\_POTENTIALS}} environment variable may contain paths
to multiple folders, if they are separated by “;” on Windows and
“:” on all other operating systems, just like the \sphinxcode{\sphinxupquote{PATH}} and
similar environment variables.


\bigskip\hrule\bigskip


\sphinxAtStartPar
The alphabetic list of pair styles defined in LAMMPS is given on the
{\hyperref[\detokenize{pair_style::doc}]{\sphinxcrossref{\DUrole{doc}{pair\_style}}}} doc page.  They are also listed in more
compact form on the {\hyperref[\detokenize{Commands_pair::doc}]{\sphinxcrossref{\DUrole{doc}{Commands pair}}}} doc page.

\sphinxAtStartPar
Click on the style to display the formula it computes and its
coefficients as specified by the associated pair\_coeff command.


\bigskip\hrule\bigskip



\subsection{Restrictions}
\label{\detokenize{pair_coeff:restrictions}}
\sphinxAtStartPar
This command must come after the simulation box is defined by a
{\hyperref[\detokenize{read_data::doc}]{\sphinxcrossref{\DUrole{doc}{read\_data}}}}, {\hyperref[\detokenize{read_restart::doc}]{\sphinxcrossref{\DUrole{doc}{read\_restart}}}}, or
{\hyperref[\detokenize{create_box::doc}]{\sphinxcrossref{\DUrole{doc}{create\_box}}}} command.


\subsection{Related commands}
\label{\detokenize{pair_coeff:related-commands}}
\sphinxAtStartPar
{\hyperref[\detokenize{pair_style::doc}]{\sphinxcrossref{\DUrole{doc}{pair\_style}}}}, {\hyperref[\detokenize{pair_modify::doc}]{\sphinxcrossref{\DUrole{doc}{pair\_modify}}}},
{\hyperref[\detokenize{read_data::doc}]{\sphinxcrossref{\DUrole{doc}{read\_data}}}}, {\hyperref[\detokenize{read_restart::doc}]{\sphinxcrossref{\DUrole{doc}{read\_restart}}}},
{\hyperref[\detokenize{pair_write::doc}]{\sphinxcrossref{\DUrole{doc}{pair\_write}}}}


\subsection{Default}
\label{\detokenize{pair_coeff:default}}
\sphinxAtStartPar
none

\sphinxstepscope

\index{pair\_modify@\spxentry{pair\_modify}}\ignorespaces 

\section{pair\_modify command}
\label{\detokenize{pair_modify:pair-modify-command}}\label{\detokenize{pair_modify:index-0}}\label{\detokenize{pair_modify::doc}}

\subsection{Syntax}
\label{\detokenize{pair_modify:syntax}}
\begin{sphinxVerbatim}[commandchars=\\\{\}]
\PYG{k}{pair\PYGZus{}modify}\PYG{+w}{ }\PYG{n}{keyword}\PYG{+w}{ }\PYG{n}{values}\PYG{+w}{ }\PYG{n}{...}
\end{sphinxVerbatim}
\begin{itemize}
\item {} 
\sphinxAtStartPar
one or more keyword/value pairs may be listed

\item {} 
\sphinxAtStartPar
keyword = \sphinxstyleemphasis{pair} or \sphinxstyleemphasis{shift} or \sphinxstyleemphasis{mix} or \sphinxstyleemphasis{table} or \sphinxstyleemphasis{table/disp} or \sphinxstyleemphasis{tabinner}
or \sphinxstyleemphasis{tabinner/disp} or \sphinxstyleemphasis{tail} or \sphinxstyleemphasis{compute} or \sphinxstyleemphasis{nofdotr} or \sphinxstyleemphasis{special} or
\sphinxstyleemphasis{compute/tally} or \sphinxstyleemphasis{neigh/trim}
\begin{sphinxalltt}
\sphinxstyleemphasis{pair} value = sub\sphinxhyphen{}style N
  sub\sphinxhyphen{}style = sub\sphinxhyphen{}style of {\hyperref[\detokenize{pair_hybrid::doc}]{\sphinxcrossref{\DUrole{doc}{pair hybrid}}}}
  N = which instance of sub\sphinxhyphen{}style (1 to M), only specify if sub\sphinxhyphen{}style is used multiple times
\sphinxstyleemphasis{mix} value = \sphinxstyleemphasis{geometric} or \sphinxstyleemphasis{arithmetic} or \sphinxstyleemphasis{sixthpower}
\sphinxstyleemphasis{shift} value = \sphinxstyleemphasis{yes} or \sphinxstyleemphasis{no}
\sphinxstyleemphasis{table} value = N
  2\textasciicircum{}N = \# of values in table
\sphinxstyleemphasis{table/disp} value = N
  2\textasciicircum{}N = \# of values in table
\sphinxstyleemphasis{tabinner} value = cutoff
  cutoff = inner cutoff at which to begin table (distance units)
\sphinxstyleemphasis{tabinner/disp} value = cutoff
  cutoff = inner cutoff at which to begin table (distance units)
\sphinxstyleemphasis{tail} value = \sphinxstyleemphasis{yes} or \sphinxstyleemphasis{no}
\sphinxstyleemphasis{compute} value = \sphinxstyleemphasis{yes} or \sphinxstyleemphasis{no}
\sphinxstyleemphasis{nofdotr} value = none
\sphinxstyleemphasis{special} values = which wt1 wt2 wt3
   which = \sphinxstyleemphasis{lj/coul} or \sphinxstyleemphasis{lj} or \sphinxstyleemphasis{coul}
   w1,w2,w3 = 1\sphinxhyphen{}2, 1\sphinxhyphen{}3, 1\sphinxhyphen{}4 weights from 0.0 to 1.0 inclusive
\sphinxstyleemphasis{compute/tally} value = \sphinxstyleemphasis{yes} or \sphinxstyleemphasis{no}
\sphinxstyleemphasis{neigh/trim} value = \sphinxstyleemphasis{yes} or \sphinxstyleemphasis{no}
\end{sphinxalltt}

\end{itemize}


\subsection{Examples}
\label{\detokenize{pair_modify:examples}}
\begin{sphinxVerbatim}[commandchars=\\\{\}]
\PYG{k}{pair\PYGZus{}modify}\PYG{+w}{ }\PYG{n}{shift}\PYG{+w}{ }\PYG{n}{yes}\PYG{+w}{ }\PYG{n}{mix}\PYG{+w}{ }\PYG{n}{geometric}
\PYG{k}{pair\PYGZus{}modify}\PYG{+w}{ }\PYG{n}{tail}\PYG{+w}{ }\PYG{n}{yes}
\PYG{k}{pair\PYGZus{}modify}\PYG{+w}{ }\PYG{n}{table}\PYG{+w}{ }\PYG{l+m}{12}
\PYG{k}{pair\PYGZus{}modify}\PYG{+w}{ }\PYG{n}{pair}\PYG{+w}{ }\PYG{n}{lj}\PYG{o}{/}\PYG{n}{cut}\PYG{+w}{ }\PYG{k}{compute }\PYG{n+nv+nvIdentifier}{no}
\PYG{n+nv+nvIdentifier}{pair\PYGZus{}modify}\PYG{+w}{ }\PYG{n}{pair}\PYG{+w}{ }\PYG{n}{tersoff}\PYG{+w}{ }\PYG{n}{compute}\PYG{o}{/}\PYG{n}{tally}\PYG{+w}{ }\PYG{n}{no}
\PYG{k}{pair\PYGZus{}modify}\PYG{+w}{ }\PYG{n}{pair}\PYG{+w}{ }\PYG{n}{lj}\PYG{o}{/}\PYG{n}{cut}\PYG{o}{/}\PYG{n}{coul}\PYG{o}{/}\PYG{n}{long}\PYG{+w}{ }\PYG{l+m}{1}\PYG{+w}{ }\PYG{n}{special}\PYG{+w}{ }\PYG{n}{lj}\PYG{o}{/}\PYG{n}{coul}\PYG{+w}{ }\PYG{l+m}{0.0}\PYG{+w}{ }\PYG{l+m}{0.0}\PYG{+w}{ }\PYG{l+m}{0.0}
\PYG{k}{pair\PYGZus{}modify}\PYG{+w}{ }\PYG{n}{pair}\PYG{+w}{ }\PYG{n}{lj}\PYG{o}{/}\PYG{n}{cut}\PYG{o}{/}\PYG{n}{coul}\PYG{o}{/}\PYG{n}{long}\PYG{+w}{ }\PYG{n}{special}\PYG{+w}{ }\PYG{n}{lj}\PYG{+w}{ }\PYG{l+m}{0.0}\PYG{+w}{ }\PYG{l+m}{0.0}\PYG{+w}{ }\PYG{l+m}{0.5}\PYG{+w}{ }\PYG{n}{special}\PYG{+w}{ }\PYG{n}{coul}\PYG{+w}{ }\PYG{l+m}{0.0}\PYG{+w}{ }\PYG{l+m}{0.0}\PYG{+w}{ }\PYG{l+m}{0.8333333}
\end{sphinxVerbatim}


\subsection{Description}
\label{\detokenize{pair_modify:description}}
\sphinxAtStartPar
Modify the parameters of the currently defined pair style.  If the
pair style is {\hyperref[\detokenize{pair_hybrid::doc}]{\sphinxcrossref{\DUrole{doc}{hybrid or hybrid/overlay}}}}, then the
specified parameters are by default modified for all the hybrid sub\sphinxhyphen{}styles.

\begin{sphinxadmonition}{note}{Note:}
\sphinxAtStartPar
The behavior for hybrid pair styles can be changed by using the \sphinxstyleemphasis{pair}
keyword, which allows selection of a specific sub\sphinxhyphen{}style to apply all
remaining keywords to.
The \sphinxstyleemphasis{special} and \sphinxstyleemphasis{compute/tally} keywords can \sphinxstylestrong{only} be
used in conjunction with the \sphinxstyleemphasis{pair} keyword.  See further details about
these 3 keywords below.
\end{sphinxadmonition}

\sphinxAtStartPar
The \sphinxstyleemphasis{mix} keyword affects pair coefficients for interactions between
atoms of type I and J, when I != J and the coefficients are not
explicitly set in the input script.  Note that coefficients for I = J
must be set explicitly, either in the input script via the
{\hyperref[\detokenize{pair_coeff::doc}]{\sphinxcrossref{\DUrole{doc}{pair\_coeff}}}} command or in the “Pair Coeffs” or “PairIJ Coeffs”
sections of the {\hyperref[\detokenize{read_data::doc}]{\sphinxcrossref{\DUrole{doc}{data file}}}}.  For some pair styles it is not
necessary to specify coefficients when I != J, since a “mixing” rule
will create them from the I,I and J,J settings.  The pair\_modify
\sphinxstyleemphasis{mix} value determines what formulas are used to compute the mixed
coefficients.  In each case, the cutoff distance is mixed the same way
as sigma.

\sphinxAtStartPar
Note that not all pair styles support mixing and some mix options are
not available for certain pair styles. Also, there are additional
restrictions when using {\hyperref[\detokenize{pair_hybrid::doc}]{\sphinxcrossref{\DUrole{doc}{pair style hybrid or hybrid/overlay}}}}.  See the page for individual pair styles for those
restrictions.  Note also that the {\hyperref[\detokenize{pair_coeff::doc}]{\sphinxcrossref{\DUrole{doc}{pair\_coeff}}}} command
also can be used to directly set coefficients for a specific I != J
pairing, in which case no mixing is performed.  If possible, LAMMPS will
print an informational message about how many of the mixed pair
coefficients were generated and which mixing rule was applied.
\begin{itemize}
\item {} 
\sphinxAtStartPar
mix \sphinxstyleemphasis{geometric}
\begin{equation*}
\begin{split}\epsilon_{ij} = & \sqrt{\epsilon_i  \epsilon_j} \\
\sigma_{ij}   = & \sqrt{\sigma_i  \sigma_j}\end{split}
\end{equation*}
\item {} 
\sphinxAtStartPar
mix \sphinxstyleemphasis{arithmetic}
\begin{equation*}
\begin{split}\epsilon_{ij} = & \sqrt{\epsilon_i  \epsilon_j} \\
\sigma_{ij}   = & \frac{1}{2} (\sigma_i + \sigma_j)\end{split}
\end{equation*}
\item {} 
\sphinxAtStartPar
mix \sphinxstyleemphasis{sixthpower}
\begin{equation*}
\begin{split}\epsilon_{ij} = & \frac{2 \sqrt{\epsilon_i \epsilon_j} \sigma_i^3 \sigma_j^3}{\sigma_i^6 + \sigma_j^6} \\
\sigma_{ij} =   & \left(\frac{1}{2} (\sigma_i^6 + \sigma_j^6) \right)^{\frac{1}{6}}\end{split}
\end{equation*}
\end{itemize}

\sphinxAtStartPar
The \sphinxstyleemphasis{shift} keyword determines whether a Lennard\sphinxhyphen{}Jones potential is
shifted at its cutoff to 0.0.  If so, this adds an energy term to each
pairwise interaction which will be included in the thermodynamic
output, but does not affect pair forces or atom trajectories.  See the
doc page for individual pair styles to see which ones support this
option.

\sphinxAtStartPar
The \sphinxstyleemphasis{table} and \sphinxstyleemphasis{table/disp} keywords apply to pair styles with a
long\sphinxhyphen{}range Coulombic term or long\sphinxhyphen{}range dispersion term respectively;
see the page for individual styles to see which potentials support
these options.  If N is non\sphinxhyphen{}zero, a table of length 2\textasciicircum{}N is
pre\sphinxhyphen{}computed for forces and energies, which can shrink their
computational cost by up to a factor of 2.  The table is indexed via a
bit\sphinxhyphen{}mapping technique {\hyperref[\detokenize{pair_modify:wolff1}]{\sphinxcrossref{\DUrole{std,std-ref}{(Wolff)}}}} and a linear
interpolation is performed between adjacent table values.  In our
experiments with different table styles (lookup, linear, spline), this
method typically gave the best performance in terms of speed and
accuracy.

\sphinxAtStartPar
The choice of table length is a tradeoff in accuracy versus speed.  A
larger N yields more accurate force computations, but requires more
memory which can slow down the computation due to cache misses.  A
reasonable value of N is between 8 and 16.  The default value of 12
(table of length 4096) gives approximately the same accuracy as the
no\sphinxhyphen{}table (N = 0) option.  For N = 0, forces and energies are computed
directly, using a polynomial fit for the needed erfc() function
evaluation, which is what earlier versions of LAMMPS did.  Values
greater than 16 typically slow down the simulation and will not
improve accuracy; values from 1 to 8 give unreliable results.

\sphinxAtStartPar
The \sphinxstyleemphasis{tabinner} and \sphinxstyleemphasis{tabinner/disp} keywords set an inner cutoff above
which the pairwise computation is done by table lookup (if tables are
invoked), for the corresponding Coulombic and dispersion tables
discussed with the \sphinxstyleemphasis{table} and \sphinxstyleemphasis{table/disp} keywords.  The smaller the
cutoff is set, the less accurate the table becomes (for a given number
of table values), which can require use of larger tables.  The default
cutoff value is sqrt(2.0) distance units which means nearly all
pairwise interactions are computed via table lookup for simulations
with “real” units, but some close pairs may be computed directly
(non\sphinxhyphen{}table) for simulations with “lj” units.

\sphinxAtStartPar
When the \sphinxstyleemphasis{tail} keyword is set to \sphinxstyleemphasis{yes}, certain pair styles will
add a long\sphinxhyphen{}range VanderWaals tail “correction” to the energy and
pressure.  These corrections are bookkeeping terms which do not affect
dynamics, unless a constant\sphinxhyphen{}pressure simulation is being performed.
See the page for individual styles to see which support this
option.  These corrections are included in the calculation and
printing of thermodynamic quantities (see the {\hyperref[\detokenize{thermo_style::doc}]{\sphinxcrossref{\DUrole{doc}{thermo\_style}}}} command).  Their effect will also be included in
constant NPT or NPH simulations where the pressure influences the
simulation box dimensions (e.g. the {\hyperref[\detokenize{fix_nh::doc}]{\sphinxcrossref{\DUrole{doc}{fix npt}}}} and
{\hyperref[\detokenize{fix_nh::doc}]{\sphinxcrossref{\DUrole{doc}{fix nph}}}} commands).  The formulas used for the
long\sphinxhyphen{}range corrections come from equation 5 of {\hyperref[\detokenize{pair_modify:sun}]{\sphinxcrossref{\DUrole{std,std-ref}{(Sun)}}}}.

\begin{sphinxadmonition}{note}{Note:}
\sphinxAtStartPar
The tail correction terms are computed at the beginning of each
run, using the current atom counts of each atom type.  If atoms are
deleted (or lost) or created during a simulation, e.g. via the
{\hyperref[\detokenize{fix_gcmc::doc}]{\sphinxcrossref{\DUrole{doc}{fix gcmc}}}} command, the correction factors are not
re\sphinxhyphen{}computed.  If you expect the counts to change dramatically, you
can break a run into a series of shorter runs so that the
correction factors are re\sphinxhyphen{}computed more frequently.
\end{sphinxadmonition}

\sphinxAtStartPar
Several additional assumptions are inherent in using tail corrections,
including the following:
\begin{itemize}
\item {} 
\sphinxAtStartPar
The simulated system is a 3d bulk homogeneous liquid. This option
should not be used for systems that are non\sphinxhyphen{}liquid, 2d, have a slab
geometry (only 2d periodic), or inhomogeneous.

\item {} 
\sphinxAtStartPar
G(r), the radial distribution function (rdf), is unity beyond the
cutoff, so a fairly large cutoff should be used (i.e. 2.5 sigma for
an LJ fluid), and it is probably a good idea to verify this
assumption by checking the rdf.  The rdf is not exactly unity beyond
the cutoff for each pair of interaction types, so the tail
correction is necessarily an approximation.

\sphinxAtStartPar
The tail corrections are computed at the beginning of each
simulation run.  If the number of atoms changes during the run,
e.g. due to atoms leaving the simulation domain, or use of the
{\hyperref[\detokenize{fix_gcmc::doc}]{\sphinxcrossref{\DUrole{doc}{fix gcmc}}}} command, then the corrections are not
updated to reflect the changed atom count.  If this is a large
effect in your simulation, you should break the long run into
several short runs, so that the correction factors are re\sphinxhyphen{}computed
multiple times.

\item {} 
\sphinxAtStartPar
Thermophysical properties obtained from calculations with this
option enabled will not be thermodynamically consistent with the
truncated force\sphinxhyphen{}field that was used.  In other words, atoms do not
feel any LJ pair interactions beyond the cutoff, but the energy and
pressure reported by the simulation include an estimated
contribution from those interactions.

\end{itemize}

\sphinxAtStartPar
The \sphinxstyleemphasis{compute} keyword allows pairwise computations to be turned off,
even though a {\hyperref[\detokenize{pair_style::doc}]{\sphinxcrossref{\DUrole{doc}{pair\_style}}}} is defined.  This is not
useful for running a real simulation, but can be useful for debugging
purposes or for performing a {\hyperref[\detokenize{rerun::doc}]{\sphinxcrossref{\DUrole{doc}{rerun}}}} simulation, when you
only wish to compute partial forces that do not include the pairwise
contribution.

\sphinxAtStartPar
Two examples are as follows.  First, this option allows you to perform
a simulation with {\hyperref[\detokenize{pair_hybrid::doc}]{\sphinxcrossref{\DUrole{doc}{pair\_style hybrid}}}} with only a
subset of the hybrid sub\sphinxhyphen{}styles enabled.  Second, this option allows
you to perform a simulation with only long\sphinxhyphen{}range interactions but no
short\sphinxhyphen{}range pairwise interactions.  Doing this by simply not defining
a pair style will not work, because the {\hyperref[\detokenize{kspace_style::doc}]{\sphinxcrossref{\DUrole{doc}{kspace\_style}}}} command requires a Kspace\sphinxhyphen{}compatible pair style be
defined.

\sphinxAtStartPar
The \sphinxstyleemphasis{nofdotr} keyword allows to disable an optimization that computes
the global stress tensor from the total forces and atom positions
rather than from summing forces between individual pairs of atoms.


\bigskip\hrule\bigskip


\sphinxAtStartPar
The \sphinxstyleemphasis{pair} keyword can only be used with the {\hyperref[\detokenize{pair_hybrid::doc}]{\sphinxcrossref{\DUrole{doc}{hybrid and
hybrid/overlay}}}} pair styles.  If used, it must appear
first in the list of keywords.

\sphinxAtStartPar
Its meaning is that all the following parameters will only be modified
for the specified sub\sphinxhyphen{}style.  If the sub\sphinxhyphen{}style is defined multiple
times, then an additional numeric argument \sphinxstyleemphasis{N} must also be specified,
which is a number from 1 to M where M is the number of times the
sub\sphinxhyphen{}style was listed in the {\hyperref[\detokenize{pair_hybrid::doc}]{\sphinxcrossref{\DUrole{doc}{pair\_style hybrid}}}}
command.  The extra number indicates which instance of the sub\sphinxhyphen{}style
the remaining keywords will be applied to.

\sphinxAtStartPar
The \sphinxstyleemphasis{special} and \sphinxstyleemphasis{compute/tally} keywords can \sphinxstylestrong{only} be used in
conjunction with the \sphinxstyleemphasis{pair} keyword and they must directly follow it.
I.e. any other keyword, must appear after \sphinxstyleemphasis{pair}, \sphinxstyleemphasis{special}, and
\sphinxstyleemphasis{compute/tally}.

\sphinxAtStartPar
The \sphinxstyleemphasis{special} keyword overrides the global {\hyperref[\detokenize{special_bonds::doc}]{\sphinxcrossref{\DUrole{doc}{special\_bonds}}}}
1\sphinxhyphen{}2, 1\sphinxhyphen{}3, 1\sphinxhyphen{}4 exclusion settings (weights) for the sub\sphinxhyphen{}style selected
by the \sphinxstyleemphasis{pair} keyword.

\sphinxAtStartPar
Similar to the {\hyperref[\detokenize{special_bonds::doc}]{\sphinxcrossref{\DUrole{doc}{special\_bonds}}}} command, it takes
4 arguments.  The \sphinxstyleemphasis{which} argument can be \sphinxstyleemphasis{lj} to change only the
non\sphinxhyphen{}Coulomb weights (e.g. Lennard\sphinxhyphen{}Jones or Buckingham), \sphinxstyleemphasis{coul} to change
only the Coulombic settings, or \sphinxstyleemphasis{lj/coul} to change both to the same
values.  The \sphinxstyleemphasis{wt1,wt2,wt3} values are numeric weights from 0.0 to 1.0
inclusive, for the 1\sphinxhyphen{}2, 1\sphinxhyphen{}3, and 1\sphinxhyphen{}4 bond topology neighbors, respectively.
The \sphinxstyleemphasis{special} keyword can be used multiple times, e.g. to set the \sphinxstyleemphasis{lj}
and \sphinxstyleemphasis{coul} settings to different values.

\begin{sphinxadmonition}{note}{Note:}
\sphinxAtStartPar
The \sphinxstyleemphasis{special} keyword is not compatible with pair styles from the
GPU or the INTEL package and attempting to use it will cause
an error.
\end{sphinxadmonition}

\begin{sphinxadmonition}{note}{Note:}
\sphinxAtStartPar
Weights of exactly 0.0 or 1.0 in the {\hyperref[\detokenize{special_bonds::doc}]{\sphinxcrossref{\DUrole{doc}{special\_bonds}}}}
command have implications on the neighbor list construction, which
means that they cannot be overridden by using the \sphinxstyleemphasis{special} keyword.
One workaround for this restriction is to use the {\hyperref[\detokenize{special_bonds::doc}]{\sphinxcrossref{\DUrole{doc}{special\_bonds}}}}
command with weights like 1.0e\sphinxhyphen{}10 or 0.999999999 instead of 0.0 or 1.0,
respectively, which enables to reset each them to any value between 0.0
and 1.0 inclusively.  Otherwise you can set \sphinxstylestrong{all} global weights to
an arbitrary number between 0.0 or 1.0, like 0.5, and then you have
to override \sphinxstylestrong{all} \sphinxstyleemphasis{special} settings for \sphinxstylestrong{all} sub\sphinxhyphen{}styles which use
the 1\sphinxhyphen{}2, 1\sphinxhyphen{}3, and 1\sphinxhyphen{}4 exclusion weights in their force/energy computation.
\end{sphinxadmonition}

\sphinxAtStartPar
The \sphinxstyleemphasis{compute/tally} keyword disables or enables registering {\hyperref[\detokenize{compute_tally::doc}]{\sphinxcrossref{\DUrole{doc}{compute
*/tally}}}} computes for the sub\sphinxhyphen{}style specified by
the \sphinxstyleemphasis{pair} keyword.  Use \sphinxstyleemphasis{no} to disable, or \sphinxstyleemphasis{yes} to enable.

\begin{sphinxadmonition}{note}{Note:}
\sphinxAtStartPar
The “pair\_modify pair compute/tally” command must be issued
\sphinxstylestrong{before} the corresponding compute style is defined.
\end{sphinxadmonition}

\sphinxAtStartPar
\DUrole{versionmodified,added}{Added in version 3Aug2022.}

\sphinxAtStartPar
The \sphinxstyleemphasis{neigh/trim} keyword controls whether an explicit cutoff is set for
each neighbor list request issued by individual pair sub\sphinxhyphen{}styles when
using {\hyperref[\detokenize{pair_hybrid::doc}]{\sphinxcrossref{\DUrole{doc}{pair hybrid/overlay}}}}.  When this keyword is
set to \sphinxstyleemphasis{no}, then the cutoff of each pair sub\sphinxhyphen{}style neighbor list will
be set equal to the largest cutoff, even if a shorter cutoff is
specified for a particular sub\sphinxhyphen{}style.  If possible the neighbor list
will be copied directly from another list.  When this keyword is set to
\sphinxstyleemphasis{yes} then the cutoff of the neighbor list will be explicitly set to the
value requested by the pair sub\sphinxhyphen{}style, and if possible the list will be
created by trimming neighbors from another list with a longer cutoff,
otherwise a new neighbor list will be created with the specified cutoff.
The \sphinxstyleemphasis{yes} option can be faster when there are multiple pair styles with
different cutoffs since the number of pair\sphinxhyphen{}wise distance checks between
neighbors is reduced (but the time required to build the neighbor lists
is increased). The \sphinxstyleemphasis{no} option could be faster when two or more neighbor
lists have similar (but not exactly the same) cutoffs.

\begin{sphinxadmonition}{note}{Note:}
\sphinxAtStartPar
The “pair\_modify neigh/trim” command \sphinxstyleemphasis{only} applies when there are
multiple pair sub\sphinxhyphen{}styles for the same atoms with different cutoffs,
i.e. when using pair style hybrid/overlay.  If you have different
cutoffs for different pairs for atoms type, the {\hyperref[\detokenize{neighbor::doc}]{\sphinxcrossref{\DUrole{doc}{neighbor style
multi}}}} should be used to create optimized neighbor lists.
\end{sphinxadmonition}


\bigskip\hrule\bigskip



\subsection{Restrictions}
\label{\detokenize{pair_modify:restrictions}}
\sphinxAtStartPar
You cannot use \sphinxstyleemphasis{shift} yes with \sphinxstyleemphasis{tail} yes, since those are
conflicting options.  You cannot use \sphinxstyleemphasis{tail} yes with 2d simulations.
You cannot use \sphinxstyleemphasis{special} with pair styles from the GPU or
INTEL package.


\subsection{Related commands}
\label{\detokenize{pair_modify:related-commands}}
\sphinxAtStartPar
{\hyperref[\detokenize{pair_style::doc}]{\sphinxcrossref{\DUrole{doc}{pair\_style}}}}, {\hyperref[\detokenize{pair_hybrid::doc}]{\sphinxcrossref{\DUrole{doc}{pair\_style hybrid}}}},
{\hyperref[\detokenize{pair_coeff::doc}]{\sphinxcrossref{\DUrole{doc}{pair\_coeff}}}}, {\hyperref[\detokenize{thermo_style::doc}]{\sphinxcrossref{\DUrole{doc}{thermo\_style}}}},
{\hyperref[\detokenize{compute_tally::doc}]{\sphinxcrossref{\DUrole{doc}{compute */tally}}}}, {\hyperref[\detokenize{neighbor::doc}]{\sphinxcrossref{\DUrole{doc}{neighbor multi}}}}


\subsection{Default}
\label{\detokenize{pair_modify:default}}
\sphinxAtStartPar
The option defaults are mix = geometric, shift = no, table = 12,
tabinner = sqrt(2.0), tail = no, compute = yes, and neigh/trim yes.

\sphinxAtStartPar
Note that some pair styles perform mixing, but only a certain style of
mixing.  See the doc pages for individual pair styles for details.


\bigskip\hrule\bigskip

\phantomsection\label{\detokenize{pair_modify:wolff1}}
\sphinxAtStartPar
\sphinxstylestrong{(Wolff)} Wolff and Rudd, Comp Phys Comm, 120, 200\sphinxhyphen{}32 (1999).

\phantomsection\label{\detokenize{pair_modify:sun}}
\sphinxAtStartPar
\sphinxstylestrong{(Sun)} Sun, J Phys Chem B, 102, 7338\sphinxhyphen{}7364 (1998).

\sphinxstepscope

\index{pair\_style@\spxentry{pair\_style}}\ignorespaces 

\section{pair\_style command}
\label{\detokenize{pair_style:pair-style-command}}\label{\detokenize{pair_style:index-0}}\label{\detokenize{pair_style::doc}}

\subsection{Syntax}
\label{\detokenize{pair_style:syntax}}
\begin{sphinxVerbatim}[commandchars=\\\{\}]
\PYG{k}{pair\PYGZus{}style}\PYG{+w}{ }\PYG{n}{style}\PYG{+w}{ }\PYG{n}{args}
\end{sphinxVerbatim}
\begin{itemize}
\item {} 
\sphinxAtStartPar
style = one of the styles from the list below

\item {} 
\sphinxAtStartPar
args = arguments used by a particular style

\end{itemize}


\subsection{Examples}
\label{\detokenize{pair_style:examples}}
\begin{sphinxVerbatim}[commandchars=\\\{\}]
\PYG{k}{pair\PYGZus{}style}\PYG{+w}{ }\PYG{n}{lj}\PYG{o}{/}\PYG{n}{cut}\PYG{+w}{ }\PYG{l+m}{2.5}
\PYG{k}{pair\PYGZus{}style}\PYG{+w}{ }\PYG{n}{eam}\PYG{o}{/}\PYG{n}{alloy}
\PYG{k}{pair\PYGZus{}style}\PYG{+w}{ }\PYG{n}{hybrid}\PYG{+w}{ }\PYG{n}{lj}\PYG{o}{/}\PYG{n}{charmm}\PYG{o}{/}\PYG{n}{coul}\PYG{o}{/}\PYG{n}{long}\PYG{+w}{ }\PYG{l+m}{10.0}\PYG{+w}{ }\PYG{n}{eam}
\PYG{k}{pair\PYGZus{}style}\PYG{+w}{ }\PYG{n}{table}\PYG{+w}{ }\PYG{n}{linear}\PYG{+w}{ }\PYG{l+m}{1000}
\PYG{k}{pair\PYGZus{}style}\PYG{+w}{ }\PYG{n}{none}
\end{sphinxVerbatim}


\subsection{Description}
\label{\detokenize{pair_style:description}}
\sphinxAtStartPar
Set the formula(s) LAMMPS uses to compute pairwise interactions.  In
LAMMPS, pair potentials are defined between pairs of atoms that are
within a cutoff distance and the set of active interactions typically
changes over time.  See the {\hyperref[\detokenize{bond_style::doc}]{\sphinxcrossref{\DUrole{doc}{bond\_style}}}} command to
define potentials between pairs of bonded atoms, which typically
remain in place for the duration of a simulation.

\sphinxAtStartPar
In LAMMPS, pairwise force fields encompass a variety of interactions,
some of which include many\sphinxhyphen{}body effects, e.g. EAM, Stillinger\sphinxhyphen{}Weber,
Tersoff, REBO potentials.  They are still classified as “pairwise”
potentials because the set of interacting atoms changes with time
(unlike molecular bonds) and thus a neighbor list is used to find
nearby interacting atoms.

\sphinxAtStartPar
Hybrid models where specified pairs of atom types interact via
different pair potentials can be setup using the \sphinxstyleemphasis{hybrid} pair style.

\sphinxAtStartPar
The coefficients associated with a pair style are typically set for
each pair of atom types, and are specified by the
{\hyperref[\detokenize{pair_coeff::doc}]{\sphinxcrossref{\DUrole{doc}{pair\_coeff}}}} command or read from a file by the
{\hyperref[\detokenize{read_data::doc}]{\sphinxcrossref{\DUrole{doc}{read\_data}}}} or {\hyperref[\detokenize{read_restart::doc}]{\sphinxcrossref{\DUrole{doc}{read\_restart}}}}
commands.

\sphinxAtStartPar
The {\hyperref[\detokenize{pair_modify::doc}]{\sphinxcrossref{\DUrole{doc}{pair\_modify}}}} command sets options for mixing of
type I\sphinxhyphen{}J interaction coefficients and adding energy offsets or tail
corrections to Lennard\sphinxhyphen{}Jones potentials.  Details on these options as
they pertain to individual potentials are described on the doc page
for the potential.  Likewise, info on whether the potential
information is stored in a {\hyperref[\detokenize{write_restart::doc}]{\sphinxcrossref{\DUrole{doc}{restart file}}}} is listed
on the potential doc page.

\sphinxAtStartPar
In the formulas listed for each pair style, \sphinxstyleemphasis{E} is the energy of a
pairwise interaction between two atoms separated by a distance \sphinxstyleemphasis{r}.
The force between the atoms is the negative derivative of this
expression.

\sphinxAtStartPar
If the pair\_style command has a cutoff argument, it sets global
cutoffs for all pairs of atom types.  The distance(s) can be smaller
or larger than the dimensions of the simulation box.

\sphinxAtStartPar
In many cases, the global cutoff value can be overridden for a
specific pair of atom types by the {\hyperref[\detokenize{pair_coeff::doc}]{\sphinxcrossref{\DUrole{doc}{pair\_coeff}}}}
command.

\sphinxAtStartPar
If a new pair\_style command is specified with a new style, all
previous {\hyperref[\detokenize{pair_coeff::doc}]{\sphinxcrossref{\DUrole{doc}{pair\_coeff}}}} and {\hyperref[\detokenize{pair_modify::doc}]{\sphinxcrossref{\DUrole{doc}{pair\_modify}}}} command settings are erased; those commands must be
re\sphinxhyphen{}specified if necessary.

\sphinxAtStartPar
If a new pair\_style command is specified with the same style, then
only the global settings in that command are reset.  Any previous
doc:\sphinxtitleref{pair\_coeff \textless{}pair\_coeff\textgreater{}} and {\hyperref[\detokenize{pair_modify::doc}]{\sphinxcrossref{\DUrole{doc}{pair\_modify}}}}
command settings are preserved.  The only exception is that if the
global cutoff in the pair\_style command is changed, it will override
the corresponding cutoff in any of the previous {\hyperref[\detokenize{pair_coeff::doc}]{\sphinxcrossref{\DUrole{doc}{pair\_modify}}}} commands.

\sphinxAtStartPar
Two pair styles which do not follow this rule are the pair\_style
\sphinxstyleemphasis{table} and \sphinxstyleemphasis{hybrid} commands.  A new pair\_style command for these
styles will wipe out all previously specified {\hyperref[\detokenize{pair_coeff::doc}]{\sphinxcrossref{\DUrole{doc}{pair\_coeff}}}} and {\hyperref[\detokenize{pair_modify::doc}]{\sphinxcrossref{\DUrole{doc}{pair\_modify}}}} settings, including
for the sub\sphinxhyphen{}styles of the \sphinxstyleemphasis{hybrid} command.


\bigskip\hrule\bigskip


\sphinxAtStartPar
Here is an alphabetic list of pair styles defined in LAMMPS.  They are
also listed in more compact form on the {\hyperref[\detokenize{Commands_pair::doc}]{\sphinxcrossref{\DUrole{doc}{Commands pair}}}} doc page.

\sphinxAtStartPar
Click on the style to display the formula it computes, any additional
arguments specified in the pair\_style command, and coefficients
specified by the associated {\hyperref[\detokenize{pair_coeff::doc}]{\sphinxcrossref{\DUrole{doc}{pair\_coeff}}}} command.

\sphinxAtStartPar
There are also additional accelerated pair styles included in the
LAMMPS distribution for faster performance on CPUs, GPUs, and KNLs.
The individual style names on the {\hyperref[\detokenize{Commands_pair::doc}]{\sphinxcrossref{\DUrole{doc}{Commands pair}}}}
doc page are followed by one or more of (g,i,k,o,t) to indicate which
accelerated styles exist.
\begin{itemize}
\item {} 
\sphinxAtStartPar
{\hyperref[\detokenize{pair_none::doc}]{\sphinxcrossref{\DUrole{doc}{none}}}} \sphinxhyphen{} turn off pairwise interactions

\item {} 
\sphinxAtStartPar
{\hyperref[\detokenize{pair_hybrid::doc}]{\sphinxcrossref{\DUrole{doc}{hybrid}}}} \sphinxhyphen{} multiple styles of pairwise interactions

\item {} 
\sphinxAtStartPar
{\hyperref[\detokenize{pair_hybrid::doc}]{\sphinxcrossref{\DUrole{doc}{hybrid/molecular}}}} \sphinxhyphen{} different pair styles for intra\sphinxhyphen{} and inter\sphinxhyphen{}molecular interactions

\item {} 
\sphinxAtStartPar
{\hyperref[\detokenize{pair_hybrid::doc}]{\sphinxcrossref{\DUrole{doc}{hybrid/overlay}}}} \sphinxhyphen{} multiple styles of superposed pairwise interactions

\item {} 
\sphinxAtStartPar
{\hyperref[\detokenize{pair_hybrid::doc}]{\sphinxcrossref{\DUrole{doc}{hybrid/scaled}}}} \sphinxhyphen{} multiple styles of scaled superposed pairwise interactions

\item {} 
\sphinxAtStartPar
{\hyperref[\detokenize{pair_zero::doc}]{\sphinxcrossref{\DUrole{doc}{zero}}}} \sphinxhyphen{} neighbor list but no interactions

\item {} 
\sphinxAtStartPar
{\hyperref[\detokenize{pair_adp::doc}]{\sphinxcrossref{\DUrole{doc}{adp}}}} \sphinxhyphen{} angular dependent potential (ADP) of Mishin

\item {} 
\sphinxAtStartPar
{\hyperref[\detokenize{pair_agni::doc}]{\sphinxcrossref{\DUrole{doc}{agni}}}} \sphinxhyphen{} AGNI machine\sphinxhyphen{}learning potential

\item {} 
\sphinxAtStartPar
{\hyperref[\detokenize{pair_aip_water_2dm::doc}]{\sphinxcrossref{\DUrole{doc}{aip/water/2dm}}}} \sphinxhyphen{} anisotropic interfacial potential for water in 2d geometries

\item {} 
\sphinxAtStartPar
{\hyperref[\detokenize{pair_airebo::doc}]{\sphinxcrossref{\DUrole{doc}{airebo}}}} \sphinxhyphen{} AIREBO potential of Stuart

\item {} 
\sphinxAtStartPar
{\hyperref[\detokenize{pair_airebo::doc}]{\sphinxcrossref{\DUrole{doc}{airebo/morse}}}} \sphinxhyphen{} AIREBO with Morse instead of LJ

\item {} 
\sphinxAtStartPar
{\hyperref[\detokenize{pair_amoeba::doc}]{\sphinxcrossref{\DUrole{doc}{amoeba}}}} \sphinxhyphen{}

\item {} 
\sphinxAtStartPar
{\hyperref[\detokenize{pair_atm::doc}]{\sphinxcrossref{\DUrole{doc}{atm}}}} \sphinxhyphen{} Axilrod\sphinxhyphen{}Teller\sphinxhyphen{}Muto potential

\item {} 
\sphinxAtStartPar
{\hyperref[\detokenize{pair_awpmd::doc}]{\sphinxcrossref{\DUrole{doc}{awpmd/cut}}}} \sphinxhyphen{} Antisymmetrized Wave Packet MD potential for atoms and electrons

\item {} 
\sphinxAtStartPar
{\hyperref[\detokenize{pair_beck::doc}]{\sphinxcrossref{\DUrole{doc}{beck}}}} \sphinxhyphen{} Beck potential

\item {} 
\sphinxAtStartPar
{\hyperref[\detokenize{pair_body_nparticle::doc}]{\sphinxcrossref{\DUrole{doc}{body/nparticle}}}} \sphinxhyphen{} interactions between body particles

\item {} 
\sphinxAtStartPar
{\hyperref[\detokenize{pair_body_rounded_polygon::doc}]{\sphinxcrossref{\DUrole{doc}{body/rounded/polygon}}}} \sphinxhyphen{} granular\sphinxhyphen{}style 2d polygon potential

\item {} 
\sphinxAtStartPar
{\hyperref[\detokenize{pair_body_rounded_polyhedron::doc}]{\sphinxcrossref{\DUrole{doc}{body/rounded/polyhedron}}}} \sphinxhyphen{} granular\sphinxhyphen{}style 3d polyhedron potential

\item {} 
\sphinxAtStartPar
{\hyperref[\detokenize{pair_bop::doc}]{\sphinxcrossref{\DUrole{doc}{bop}}}} \sphinxhyphen{} BOP potential of Pettifor

\item {} 
\sphinxAtStartPar
{\hyperref[\detokenize{pair_born::doc}]{\sphinxcrossref{\DUrole{doc}{born}}}} \sphinxhyphen{} Born\sphinxhyphen{}Mayer\sphinxhyphen{}Huggins potential

\item {} 
\sphinxAtStartPar
{\hyperref[\detokenize{pair_born::doc}]{\sphinxcrossref{\DUrole{doc}{born/coul/dsf}}}} \sphinxhyphen{} Born with damped\sphinxhyphen{}shifted\sphinxhyphen{}force model

\item {} 
\sphinxAtStartPar
{\hyperref[\detokenize{pair_cs::doc}]{\sphinxcrossref{\DUrole{doc}{born/coul/dsf/cs}}}} \sphinxhyphen{} Born with damped\sphinxhyphen{}shifted\sphinxhyphen{}force and core/shell model

\item {} 
\sphinxAtStartPar
{\hyperref[\detokenize{pair_born::doc}]{\sphinxcrossref{\DUrole{doc}{born/coul/long}}}} \sphinxhyphen{} Born with long\sphinxhyphen{}range Coulomb

\item {} 
\sphinxAtStartPar
{\hyperref[\detokenize{pair_cs::doc}]{\sphinxcrossref{\DUrole{doc}{born/coul/long/cs}}}} \sphinxhyphen{} Born with long\sphinxhyphen{}range Coulomb and core/shell

\item {} 
\sphinxAtStartPar
{\hyperref[\detokenize{pair_born::doc}]{\sphinxcrossref{\DUrole{doc}{born/coul/msm}}}} \sphinxhyphen{} Born with long\sphinxhyphen{}range MSM Coulomb

\item {} 
\sphinxAtStartPar
{\hyperref[\detokenize{pair_born::doc}]{\sphinxcrossref{\DUrole{doc}{born/coul/wolf}}}} \sphinxhyphen{} Born with Wolf potential for Coulomb

\item {} 
\sphinxAtStartPar
{\hyperref[\detokenize{pair_cs::doc}]{\sphinxcrossref{\DUrole{doc}{born/coul/wolf/cs}}}} \sphinxhyphen{} Born with Wolf potential for Coulomb and core/shell model

\item {} 
\sphinxAtStartPar
{\hyperref[\detokenize{pair_born_gauss::doc}]{\sphinxcrossref{\DUrole{doc}{born/gauss}}}} \sphinxhyphen{} Born\sphinxhyphen{}Mayer / Gaussian potential

\item {} 
\sphinxAtStartPar
{\hyperref[\detokenize{pair_bpm_spring::doc}]{\sphinxcrossref{\DUrole{doc}{bpm/spring}}}} \sphinxhyphen{} repulsive harmonic force with damping

\item {} 
\sphinxAtStartPar
{\hyperref[\detokenize{pair_brownian::doc}]{\sphinxcrossref{\DUrole{doc}{brownian}}}} \sphinxhyphen{} Brownian potential for Fast Lubrication Dynamics

\item {} 
\sphinxAtStartPar
{\hyperref[\detokenize{pair_brownian::doc}]{\sphinxcrossref{\DUrole{doc}{brownian/poly}}}} \sphinxhyphen{} Brownian potential for Fast Lubrication Dynamics with polydispersity

\item {} 
\sphinxAtStartPar
{\hyperref[\detokenize{pair_buck::doc}]{\sphinxcrossref{\DUrole{doc}{buck}}}} \sphinxhyphen{} Buckingham potential

\item {} 
\sphinxAtStartPar
{\hyperref[\detokenize{pair_buck::doc}]{\sphinxcrossref{\DUrole{doc}{buck/coul/cut}}}} \sphinxhyphen{} Buckingham with cutoff Coulomb

\item {} 
\sphinxAtStartPar
{\hyperref[\detokenize{pair_buck::doc}]{\sphinxcrossref{\DUrole{doc}{buck/coul/long}}}} \sphinxhyphen{} Buckingham with long\sphinxhyphen{}range Coulomb

\item {} 
\sphinxAtStartPar
{\hyperref[\detokenize{pair_cs::doc}]{\sphinxcrossref{\DUrole{doc}{buck/coul/long/cs}}}} \sphinxhyphen{} Buckingham with long\sphinxhyphen{}range Coulomb and core/shell

\item {} 
\sphinxAtStartPar
{\hyperref[\detokenize{pair_buck::doc}]{\sphinxcrossref{\DUrole{doc}{buck/coul/msm}}}} \sphinxhyphen{} Buckingham with long\sphinxhyphen{}range MSM Coulomb

\item {} 
\sphinxAtStartPar
{\hyperref[\detokenize{pair_buck_long::doc}]{\sphinxcrossref{\DUrole{doc}{buck/long/coul/long}}}} \sphinxhyphen{} long\sphinxhyphen{}range Buckingham with long\sphinxhyphen{}range Coulomb

\item {} 
\sphinxAtStartPar
{\hyperref[\detokenize{pair_mdf::doc}]{\sphinxcrossref{\DUrole{doc}{buck/mdf}}}} \sphinxhyphen{} Buckingham with a taper function

\item {} 
\sphinxAtStartPar
{\hyperref[\detokenize{pair_buck6d_coul_gauss::doc}]{\sphinxcrossref{\DUrole{doc}{buck6d/coul/gauss/dsf}}}} \sphinxhyphen{} dispersion\sphinxhyphen{}damped Buckingham with damped\sphinxhyphen{}shift\sphinxhyphen{}force model

\item {} 
\sphinxAtStartPar
{\hyperref[\detokenize{pair_buck6d_coul_gauss::doc}]{\sphinxcrossref{\DUrole{doc}{buck6d/coul/gauss/long}}}} \sphinxhyphen{} dispersion\sphinxhyphen{}damped Buckingham with long\sphinxhyphen{}range Coulomb

\item {} 
\sphinxAtStartPar
{\hyperref[\detokenize{pair_colloid::doc}]{\sphinxcrossref{\DUrole{doc}{colloid}}}} \sphinxhyphen{} integrated colloidal potential

\item {} 
\sphinxAtStartPar
{\hyperref[\detokenize{pair_comb::doc}]{\sphinxcrossref{\DUrole{doc}{comb}}}} \sphinxhyphen{} charge\sphinxhyphen{}optimized many\sphinxhyphen{}body (COMB) potential

\item {} 
\sphinxAtStartPar
{\hyperref[\detokenize{pair_comb::doc}]{\sphinxcrossref{\DUrole{doc}{comb3}}}} \sphinxhyphen{} charge\sphinxhyphen{}optimized many\sphinxhyphen{}body (COMB3) potential

\item {} 
\sphinxAtStartPar
{\hyperref[\detokenize{pair_cosine_squared::doc}]{\sphinxcrossref{\DUrole{doc}{cosine/squared}}}} \sphinxhyphen{} Cooke\sphinxhyphen{}Kremer\sphinxhyphen{}Deserno membrane model potential

\item {} 
\sphinxAtStartPar
{\hyperref[\detokenize{pair_coul::doc}]{\sphinxcrossref{\DUrole{doc}{coul/cut}}}} \sphinxhyphen{} cutoff Coulomb potential

\item {} 
\sphinxAtStartPar
{\hyperref[\detokenize{pair_dielectric::doc}]{\sphinxcrossref{\DUrole{doc}{coul/cut/dielectric}}}} \sphinxhyphen{}

\item {} 
\sphinxAtStartPar
{\hyperref[\detokenize{pair_coul::doc}]{\sphinxcrossref{\DUrole{doc}{coul/cut/global}}}} \sphinxhyphen{} cutoff Coulomb potential

\item {} 
\sphinxAtStartPar
{\hyperref[\detokenize{pair_fep_soft::doc}]{\sphinxcrossref{\DUrole{doc}{coul/cut/soft}}}} \sphinxhyphen{} Coulomb potential with a soft core

\item {} 
\sphinxAtStartPar
{\hyperref[\detokenize{pair_coul::doc}]{\sphinxcrossref{\DUrole{doc}{coul/debye}}}} \sphinxhyphen{} cutoff Coulomb potential with Debye screening

\item {} 
\sphinxAtStartPar
{\hyperref[\detokenize{pair_coul_diel::doc}]{\sphinxcrossref{\DUrole{doc}{coul/diel}}}} \sphinxhyphen{} Coulomb potential with dielectric permittivity

\item {} 
\sphinxAtStartPar
{\hyperref[\detokenize{pair_coul::doc}]{\sphinxcrossref{\DUrole{doc}{coul/dsf}}}} \sphinxhyphen{} Coulomb with damped\sphinxhyphen{}shifted\sphinxhyphen{}force model

\item {} 
\sphinxAtStartPar
{\hyperref[\detokenize{pair_coul::doc}]{\sphinxcrossref{\DUrole{doc}{coul/exclude}}}} \sphinxhyphen{} subtract Coulomb potential for excluded pairs

\item {} 
\sphinxAtStartPar
{\hyperref[\detokenize{pair_coul::doc}]{\sphinxcrossref{\DUrole{doc}{coul/long}}}} \sphinxhyphen{} long\sphinxhyphen{}range Coulomb potential

\item {} 
\sphinxAtStartPar
{\hyperref[\detokenize{pair_cs::doc}]{\sphinxcrossref{\DUrole{doc}{coul/long/cs}}}} \sphinxhyphen{} long\sphinxhyphen{}range Coulomb potential and core/shell

\item {} 
\sphinxAtStartPar
{\hyperref[\detokenize{pair_dielectric::doc}]{\sphinxcrossref{\DUrole{doc}{coul/long/dielectric}}}} \sphinxhyphen{}

\item {} 
\sphinxAtStartPar
{\hyperref[\detokenize{pair_fep_soft::doc}]{\sphinxcrossref{\DUrole{doc}{coul/long/soft}}}} \sphinxhyphen{} long\sphinxhyphen{}range Coulomb potential with a soft core

\item {} 
\sphinxAtStartPar
{\hyperref[\detokenize{pair_coul::doc}]{\sphinxcrossref{\DUrole{doc}{coul/msm}}}} \sphinxhyphen{} long\sphinxhyphen{}range MSM Coulomb

\item {} 
\sphinxAtStartPar
{\hyperref[\detokenize{pair_coul::doc}]{\sphinxcrossref{\DUrole{doc}{coul/slater/cut}}}} \sphinxhyphen{} smeared out Coulomb

\item {} 
\sphinxAtStartPar
{\hyperref[\detokenize{pair_coul::doc}]{\sphinxcrossref{\DUrole{doc}{coul/slater/long}}}} \sphinxhyphen{} long\sphinxhyphen{}range smeared out Coulomb

\item {} 
\sphinxAtStartPar
{\hyperref[\detokenize{pair_coul_shield::doc}]{\sphinxcrossref{\DUrole{doc}{coul/shield}}}} \sphinxhyphen{} Coulomb for boron nitride for use with {\hyperref[\detokenize{pair_ilp_graphene_hbn::doc}]{\sphinxcrossref{\DUrole{doc}{ilp/graphene/hbn}}}} potential

\item {} 
\sphinxAtStartPar
{\hyperref[\detokenize{pair_coul::doc}]{\sphinxcrossref{\DUrole{doc}{coul/streitz}}}} \sphinxhyphen{} Coulomb via Streitz/Mintmire Slater orbitals

\item {} 
\sphinxAtStartPar
{\hyperref[\detokenize{pair_coul_tt::doc}]{\sphinxcrossref{\DUrole{doc}{coul/tt}}}} \sphinxhyphen{} damped charge\sphinxhyphen{}dipole Coulomb for Drude dipoles

\item {} 
\sphinxAtStartPar
{\hyperref[\detokenize{pair_coul::doc}]{\sphinxcrossref{\DUrole{doc}{coul/wolf}}}} \sphinxhyphen{} Coulomb via Wolf potential

\item {} 
\sphinxAtStartPar
{\hyperref[\detokenize{pair_cs::doc}]{\sphinxcrossref{\DUrole{doc}{coul/wolf/cs}}}} \sphinxhyphen{} Coulomb via Wolf potential with core/shell adjustments

\item {} 
\sphinxAtStartPar
{\hyperref[\detokenize{pair_dpd::doc}]{\sphinxcrossref{\DUrole{doc}{dpd}}}} \sphinxhyphen{} dissipative particle dynamics (DPD)

\item {} 
\sphinxAtStartPar
{\hyperref[\detokenize{pair_dpd_coul_slater_long::doc}]{\sphinxcrossref{\DUrole{doc}{dpd/coul/slater/long}}}} \sphinxhyphen{} dissipative particle dynamics (DPD) with electrostatic interactions

\item {} 
\sphinxAtStartPar
{\hyperref[\detokenize{pair_dpd_ext::doc}]{\sphinxcrossref{\DUrole{doc}{dpd/ext}}}} \sphinxhyphen{} generalized force field for DPD

\item {} 
\sphinxAtStartPar
{\hyperref[\detokenize{pair_dpd_ext::doc}]{\sphinxcrossref{\DUrole{doc}{dpd/ext/tstat}}}} \sphinxhyphen{} pairwise DPD thermostatting  with generalized force field

\item {} 
\sphinxAtStartPar
{\hyperref[\detokenize{pair_dpd_fdt::doc}]{\sphinxcrossref{\DUrole{doc}{dpd/fdt}}}} \sphinxhyphen{} DPD for constant temperature and pressure

\item {} 
\sphinxAtStartPar
{\hyperref[\detokenize{pair_dpd_fdt::doc}]{\sphinxcrossref{\DUrole{doc}{dpd/fdt/energy}}}} \sphinxhyphen{} DPD for constant energy and enthalpy

\item {} 
\sphinxAtStartPar
{\hyperref[\detokenize{pair_dpd::doc}]{\sphinxcrossref{\DUrole{doc}{dpd/tstat}}}} \sphinxhyphen{} pairwise DPD thermostatting

\item {} 
\sphinxAtStartPar
{\hyperref[\detokenize{pair_dsmc::doc}]{\sphinxcrossref{\DUrole{doc}{dsmc}}}} \sphinxhyphen{} Direct Simulation Monte Carlo (DSMC)

\item {} 
\sphinxAtStartPar
{\hyperref[\detokenize{pair_e3b::doc}]{\sphinxcrossref{\DUrole{doc}{e3b}}}} \sphinxhyphen{} Explicit\sphinxhyphen{}three body (E3B) water model

\item {} 
\sphinxAtStartPar
{\hyperref[\detokenize{pair_drip::doc}]{\sphinxcrossref{\DUrole{doc}{drip}}}} \sphinxhyphen{} Dihedral\sphinxhyphen{}angle\sphinxhyphen{}corrected registry\sphinxhyphen{}dependent interlayer potential (DRIP)

\item {} 
\sphinxAtStartPar
{\hyperref[\detokenize{pair_eam::doc}]{\sphinxcrossref{\DUrole{doc}{eam}}}} \sphinxhyphen{} embedded atom method (EAM)

\item {} 
\sphinxAtStartPar
{\hyperref[\detokenize{pair_eam::doc}]{\sphinxcrossref{\DUrole{doc}{eam/alloy}}}} \sphinxhyphen{} alloy EAM

\item {} 
\sphinxAtStartPar
{\hyperref[\detokenize{pair_eam::doc}]{\sphinxcrossref{\DUrole{doc}{eam/cd}}}} \sphinxhyphen{} concentration\sphinxhyphen{}dependent EAM

\item {} 
\sphinxAtStartPar
{\hyperref[\detokenize{pair_eam::doc}]{\sphinxcrossref{\DUrole{doc}{eam/cd/old}}}} \sphinxhyphen{} older two\sphinxhyphen{}site model for concentration\sphinxhyphen{}dependent EAM

\item {} 
\sphinxAtStartPar
{\hyperref[\detokenize{pair_eam::doc}]{\sphinxcrossref{\DUrole{doc}{eam/fs}}}} \sphinxhyphen{} Finnis\sphinxhyphen{}Sinclair EAM

\item {} 
\sphinxAtStartPar
{\hyperref[\detokenize{pair_eam::doc}]{\sphinxcrossref{\DUrole{doc}{eam/he}}}} \sphinxhyphen{} Finnis\sphinxhyphen{}Sinclair EAM modified for Helium in metals

\item {} 
\sphinxAtStartPar
{\hyperref[\detokenize{pair_edip::doc}]{\sphinxcrossref{\DUrole{doc}{edip}}}} \sphinxhyphen{} three\sphinxhyphen{}body EDIP potential

\item {} 
\sphinxAtStartPar
{\hyperref[\detokenize{pair_edip::doc}]{\sphinxcrossref{\DUrole{doc}{edip/multi}}}} \sphinxhyphen{} multi\sphinxhyphen{}element EDIP potential

\item {} 
\sphinxAtStartPar
{\hyperref[\detokenize{pair_mesodpd::doc}]{\sphinxcrossref{\DUrole{doc}{edpd}}}} \sphinxhyphen{} eDPD particle interactions

\item {} 
\sphinxAtStartPar
{\hyperref[\detokenize{pair_eff::doc}]{\sphinxcrossref{\DUrole{doc}{eff/cut}}}} \sphinxhyphen{} electron force field with a cutoff

\item {} 
\sphinxAtStartPar
{\hyperref[\detokenize{pair_eim::doc}]{\sphinxcrossref{\DUrole{doc}{eim}}}} \sphinxhyphen{} embedded ion method (EIM)

\item {} 
\sphinxAtStartPar
{\hyperref[\detokenize{pair_exp6_rx::doc}]{\sphinxcrossref{\DUrole{doc}{exp6/rx}}}} \sphinxhyphen{} reactive DPD potential

\item {} 
\sphinxAtStartPar
{\hyperref[\detokenize{pair_extep::doc}]{\sphinxcrossref{\DUrole{doc}{extep}}}} \sphinxhyphen{} extended Tersoff potential

\item {} 
\sphinxAtStartPar
{\hyperref[\detokenize{pair_gauss::doc}]{\sphinxcrossref{\DUrole{doc}{gauss}}}} \sphinxhyphen{} Gaussian potential

\item {} 
\sphinxAtStartPar
{\hyperref[\detokenize{pair_gauss::doc}]{\sphinxcrossref{\DUrole{doc}{gauss/cut}}}} \sphinxhyphen{} generalized Gaussian potential

\item {} 
\sphinxAtStartPar
{\hyperref[\detokenize{pair_gayberne::doc}]{\sphinxcrossref{\DUrole{doc}{gayberne}}}} \sphinxhyphen{} Gay\sphinxhyphen{}Berne ellipsoidal potential

\item {} 
\sphinxAtStartPar
{\hyperref[\detokenize{pair_granular::doc}]{\sphinxcrossref{\DUrole{doc}{granular}}}} \sphinxhyphen{} Generalized granular potential

\item {} 
\sphinxAtStartPar
{\hyperref[\detokenize{pair_gran::doc}]{\sphinxcrossref{\DUrole{doc}{gran/hertz/history}}}} \sphinxhyphen{} granular potential with Hertzian interactions

\item {} 
\sphinxAtStartPar
{\hyperref[\detokenize{pair_gran::doc}]{\sphinxcrossref{\DUrole{doc}{gran/hooke}}}} \sphinxhyphen{} granular potential with history effects

\item {} 
\sphinxAtStartPar
{\hyperref[\detokenize{pair_gran::doc}]{\sphinxcrossref{\DUrole{doc}{gran/hooke/history}}}} \sphinxhyphen{} granular potential without history effects

\item {} 
\sphinxAtStartPar
{\hyperref[\detokenize{pair_gw::doc}]{\sphinxcrossref{\DUrole{doc}{gw}}}} \sphinxhyphen{} Gao\sphinxhyphen{}Weber potential

\item {} 
\sphinxAtStartPar
{\hyperref[\detokenize{pair_gw::doc}]{\sphinxcrossref{\DUrole{doc}{gw/zbl}}}} \sphinxhyphen{} Gao\sphinxhyphen{}Weber potential with a repulsive ZBL core

\item {} 
\sphinxAtStartPar
{\hyperref[\detokenize{pair_harmonic_cut::doc}]{\sphinxcrossref{\DUrole{doc}{harmonic/cut}}}} \sphinxhyphen{} repulsive\sphinxhyphen{}only harmonic potential

\item {} 
\sphinxAtStartPar
{\hyperref[\detokenize{pair_hbond_dreiding::doc}]{\sphinxcrossref{\DUrole{doc}{hbond/dreiding/lj}}}} \sphinxhyphen{} DREIDING hydrogen bonding LJ potential

\item {} 
\sphinxAtStartPar
{\hyperref[\detokenize{pair_hbond_dreiding::doc}]{\sphinxcrossref{\DUrole{doc}{hbond/dreiding/morse}}}} \sphinxhyphen{} DREIDING hydrogen bonding Morse potential

\item {} 
\sphinxAtStartPar
{\hyperref[\detokenize{pair_hdnnp::doc}]{\sphinxcrossref{\DUrole{doc}{hdnnp}}}} \sphinxhyphen{} High\sphinxhyphen{}dimensional neural network potential

\item {} 
\sphinxAtStartPar
{\hyperref[\detokenize{pair_amoeba::doc}]{\sphinxcrossref{\DUrole{doc}{hippo}}}} \sphinxhyphen{}

\item {} 
\sphinxAtStartPar
{\hyperref[\detokenize{pair_ilp_graphene_hbn::doc}]{\sphinxcrossref{\DUrole{doc}{ilp/graphene/hbn}}}} \sphinxhyphen{} registry\sphinxhyphen{}dependent interlayer potential (ILP)

\item {} 
\sphinxAtStartPar
{\hyperref[\detokenize{pair_ilp_tmd::doc}]{\sphinxcrossref{\DUrole{doc}{ilp/tmd}}}} \sphinxhyphen{} interlayer potential (ILP) potential for transition metal dichalcogenides (TMD)

\item {} 
\sphinxAtStartPar
{\hyperref[\detokenize{pair_kim::doc}]{\sphinxcrossref{\DUrole{doc}{kim}}}} \sphinxhyphen{} interface to potentials provided by KIM project

\item {} 
\sphinxAtStartPar
{\hyperref[\detokenize{pair_kolmogorov_crespi_full::doc}]{\sphinxcrossref{\DUrole{doc}{kolmogorov/crespi/full}}}} \sphinxhyphen{} Kolmogorov\sphinxhyphen{}Crespi (KC) potential with no simplifications

\item {} 
\sphinxAtStartPar
{\hyperref[\detokenize{pair_kolmogorov_crespi_z::doc}]{\sphinxcrossref{\DUrole{doc}{kolmogorov/crespi/z}}}} \sphinxhyphen{} Kolmogorov\sphinxhyphen{}Crespi (KC) potential with normals along z\sphinxhyphen{}axis

\item {} 
\sphinxAtStartPar
{\hyperref[\detokenize{pair_lcbop::doc}]{\sphinxcrossref{\DUrole{doc}{lcbop}}}} \sphinxhyphen{} long\sphinxhyphen{}range bond\sphinxhyphen{}order potential (LCBOP)

\item {} 
\sphinxAtStartPar
{\hyperref[\detokenize{pair_lebedeva_z::doc}]{\sphinxcrossref{\DUrole{doc}{lebedeva/z}}}} \sphinxhyphen{} Lebedeva interlayer potential for graphene with normals along z\sphinxhyphen{}axis

\item {} 
\sphinxAtStartPar
{\hyperref[\detokenize{pair_mdf::doc}]{\sphinxcrossref{\DUrole{doc}{lennard/mdf}}}} \sphinxhyphen{} LJ potential in A/B form with a taper function

\item {} 
\sphinxAtStartPar
{\hyperref[\detokenize{pair_lepton::doc}]{\sphinxcrossref{\DUrole{doc}{lepton}}}} \sphinxhyphen{} pair potential from evaluating a string

\item {} 
\sphinxAtStartPar
{\hyperref[\detokenize{pair_lepton::doc}]{\sphinxcrossref{\DUrole{doc}{lepton/coul}}}} \sphinxhyphen{} pair potential from evaluating a string with support for charges

\item {} 
\sphinxAtStartPar
{\hyperref[\detokenize{pair_lepton::doc}]{\sphinxcrossref{\DUrole{doc}{lepton/sphere}}}} \sphinxhyphen{} pair potential from evaluating a string with support for radii

\item {} 
\sphinxAtStartPar
{\hyperref[\detokenize{pair_line_lj::doc}]{\sphinxcrossref{\DUrole{doc}{line/lj}}}} \sphinxhyphen{} LJ potential between line segments

\item {} 
\sphinxAtStartPar
{\hyperref[\detokenize{pair_list::doc}]{\sphinxcrossref{\DUrole{doc}{list}}}} \sphinxhyphen{} potential between pairs of atoms explicitly listed in an input file

\item {} 
\sphinxAtStartPar
{\hyperref[\detokenize{pair_charmm::doc}]{\sphinxcrossref{\DUrole{doc}{lj/charmm/coul/charmm}}}} \sphinxhyphen{} CHARMM potential with cutoff Coulomb

\item {} 
\sphinxAtStartPar
{\hyperref[\detokenize{pair_charmm::doc}]{\sphinxcrossref{\DUrole{doc}{lj/charmm/coul/charmm/implicit}}}} \sphinxhyphen{} CHARMM for implicit solvent

\item {} 
\sphinxAtStartPar
{\hyperref[\detokenize{pair_charmm::doc}]{\sphinxcrossref{\DUrole{doc}{lj/charmm/coul/long}}}} \sphinxhyphen{} CHARMM with long\sphinxhyphen{}range Coulomb

\item {} 
\sphinxAtStartPar
{\hyperref[\detokenize{pair_fep_soft::doc}]{\sphinxcrossref{\DUrole{doc}{lj/charmm/coul/long/soft}}}} \sphinxhyphen{} CHARMM with long\sphinxhyphen{}range Coulomb and a soft core

\item {} 
\sphinxAtStartPar
{\hyperref[\detokenize{pair_charmm::doc}]{\sphinxcrossref{\DUrole{doc}{lj/charmm/coul/msm}}}} \sphinxhyphen{} CHARMM with long\sphinxhyphen{}range MSM Coulomb

\item {} 
\sphinxAtStartPar
{\hyperref[\detokenize{pair_charmm::doc}]{\sphinxcrossref{\DUrole{doc}{lj/charmmfsw/coul/charmmfsh}}}} \sphinxhyphen{} CHARMM with force switching and shifting

\item {} 
\sphinxAtStartPar
{\hyperref[\detokenize{pair_charmm::doc}]{\sphinxcrossref{\DUrole{doc}{lj/charmmfsw/coul/long}}}} \sphinxhyphen{} CHARMM with force switching and long\sphinxhyphen{}rnage Coulomb

\item {} 
\sphinxAtStartPar
{\hyperref[\detokenize{pair_class2::doc}]{\sphinxcrossref{\DUrole{doc}{lj/class2}}}} \sphinxhyphen{} COMPASS (class 2) force field without Coulomb

\item {} 
\sphinxAtStartPar
{\hyperref[\detokenize{pair_class2::doc}]{\sphinxcrossref{\DUrole{doc}{lj/class2/coul/cut}}}} \sphinxhyphen{} COMPASS with cutoff Coulomb

\item {} 
\sphinxAtStartPar
{\hyperref[\detokenize{pair_fep_soft::doc}]{\sphinxcrossref{\DUrole{doc}{lj/class2/coul/cut/soft}}}} \sphinxhyphen{} COMPASS with cutoff Coulomb with a soft core

\item {} 
\sphinxAtStartPar
{\hyperref[\detokenize{pair_class2::doc}]{\sphinxcrossref{\DUrole{doc}{lj/class2/coul/long}}}} \sphinxhyphen{} COMPASS with long\sphinxhyphen{}range Coulomb

\item {} 
\sphinxAtStartPar
{\hyperref[\detokenize{pair_cs::doc}]{\sphinxcrossref{\DUrole{doc}{lj/class2/coul/long/cs}}}} \sphinxhyphen{} COMPASS with long\sphinxhyphen{}range Coulomb with core/shell adjustments

\item {} 
\sphinxAtStartPar
{\hyperref[\detokenize{pair_fep_soft::doc}]{\sphinxcrossref{\DUrole{doc}{lj/class2/coul/long/soft}}}} \sphinxhyphen{} COMPASS with long\sphinxhyphen{}range Coulomb with a soft core

\item {} 
\sphinxAtStartPar
{\hyperref[\detokenize{pair_fep_soft::doc}]{\sphinxcrossref{\DUrole{doc}{lj/class2/soft}}}} \sphinxhyphen{} COMPASS (class 2) force field with no Coulomb with a soft core

\item {} 
\sphinxAtStartPar
{\hyperref[\detokenize{pair_lj_cubic::doc}]{\sphinxcrossref{\DUrole{doc}{lj/cubic}}}} \sphinxhyphen{} LJ with cubic after inflection point

\item {} 
\sphinxAtStartPar
{\hyperref[\detokenize{pair_lj::doc}]{\sphinxcrossref{\DUrole{doc}{lj/cut}}}} \sphinxhyphen{} cutoff Lennard\sphinxhyphen{}Jones potential without Coulomb

\item {} 
\sphinxAtStartPar
{\hyperref[\detokenize{pair_lj_cut_coul::doc}]{\sphinxcrossref{\DUrole{doc}{lj/cut/coul/cut}}}} \sphinxhyphen{} LJ with cutoff Coulomb

\item {} 
\sphinxAtStartPar
{\hyperref[\detokenize{pair_dielectric::doc}]{\sphinxcrossref{\DUrole{doc}{lj/cut/coul/cut/dielectric}}}} \sphinxhyphen{}

\item {} 
\sphinxAtStartPar
{\hyperref[\detokenize{pair_fep_soft::doc}]{\sphinxcrossref{\DUrole{doc}{lj/cut/coul/cut/soft}}}} \sphinxhyphen{} LJ with cutoff Coulomb with a soft core

\item {} 
\sphinxAtStartPar
{\hyperref[\detokenize{pair_lj_cut_coul::doc}]{\sphinxcrossref{\DUrole{doc}{lj/cut/coul/debye}}}} \sphinxhyphen{} LJ with Debye screening added to Coulomb

\item {} 
\sphinxAtStartPar
{\hyperref[\detokenize{pair_dielectric::doc}]{\sphinxcrossref{\DUrole{doc}{lj/cut/coul/debye/dielectric}}}} \sphinxhyphen{}

\item {} 
\sphinxAtStartPar
{\hyperref[\detokenize{pair_lj_cut_coul::doc}]{\sphinxcrossref{\DUrole{doc}{lj/cut/coul/dsf}}}} \sphinxhyphen{} LJ with Coulomb via damped shifted forces

\item {} 
\sphinxAtStartPar
{\hyperref[\detokenize{pair_lj_cut_coul::doc}]{\sphinxcrossref{\DUrole{doc}{lj/cut/coul/long}}}} \sphinxhyphen{} LJ with long\sphinxhyphen{}range Coulomb

\item {} 
\sphinxAtStartPar
{\hyperref[\detokenize{pair_cs::doc}]{\sphinxcrossref{\DUrole{doc}{lj/cut/coul/long/cs}}}} \sphinxhyphen{} LJ with long\sphinxhyphen{}range Coulomb with core/shell adjustments

\item {} 
\sphinxAtStartPar
{\hyperref[\detokenize{pair_dielectric::doc}]{\sphinxcrossref{\DUrole{doc}{lj/cut/coul/long/dielectric}}}} \sphinxhyphen{}

\item {} 
\sphinxAtStartPar
{\hyperref[\detokenize{pair_fep_soft::doc}]{\sphinxcrossref{\DUrole{doc}{lj/cut/coul/long/soft}}}} \sphinxhyphen{} LJ with long\sphinxhyphen{}range Coulomb with a soft core

\item {} 
\sphinxAtStartPar
{\hyperref[\detokenize{pair_lj_cut_coul::doc}]{\sphinxcrossref{\DUrole{doc}{lj/cut/coul/msm}}}} \sphinxhyphen{} LJ with long\sphinxhyphen{}range MSM Coulomb

\item {} 
\sphinxAtStartPar
{\hyperref[\detokenize{pair_dielectric::doc}]{\sphinxcrossref{\DUrole{doc}{lj/cut/coul/msm/dielectric}}}} \sphinxhyphen{}

\item {} 
\sphinxAtStartPar
{\hyperref[\detokenize{pair_lj_cut_coul::doc}]{\sphinxcrossref{\DUrole{doc}{lj/cut/coul/wolf}}}} \sphinxhyphen{} LJ with Coulomb via Wolf potential

\item {} 
\sphinxAtStartPar
{\hyperref[\detokenize{pair_dipole::doc}]{\sphinxcrossref{\DUrole{doc}{lj/cut/dipole/cut}}}} \sphinxhyphen{} point dipoles with cutoff

\item {} 
\sphinxAtStartPar
{\hyperref[\detokenize{pair_dipole::doc}]{\sphinxcrossref{\DUrole{doc}{lj/cut/dipole/long}}}} \sphinxhyphen{} point dipoles with long\sphinxhyphen{}range Ewald

\item {} 
\sphinxAtStartPar
{\hyperref[\detokenize{pair_fep_soft::doc}]{\sphinxcrossref{\DUrole{doc}{lj/cut/soft}}}} \sphinxhyphen{} LJ with a soft core

\item {} 
\sphinxAtStartPar
{\hyperref[\detokenize{pair_lj_cut_sphere::doc}]{\sphinxcrossref{\DUrole{doc}{lj/cut/sphere}}}} \sphinxhyphen{} LJ where per\sphinxhyphen{}atom radius is used as LJ sigma

\item {} 
\sphinxAtStartPar
{\hyperref[\detokenize{pair_thole::doc}]{\sphinxcrossref{\DUrole{doc}{lj/cut/thole/long}}}} \sphinxhyphen{} LJ with Coulomb with thole damping

\item {} 
\sphinxAtStartPar
{\hyperref[\detokenize{pair_lj_cut_tip4p::doc}]{\sphinxcrossref{\DUrole{doc}{lj/cut/tip4p/cut}}}} \sphinxhyphen{} LJ with cutoff Coulomb for TIP4P water

\item {} 
\sphinxAtStartPar
{\hyperref[\detokenize{pair_lj_cut_tip4p::doc}]{\sphinxcrossref{\DUrole{doc}{lj/cut/tip4p/long}}}} \sphinxhyphen{} LJ with long\sphinxhyphen{}range Coulomb for TIP4P water

\item {} 
\sphinxAtStartPar
{\hyperref[\detokenize{pair_fep_soft::doc}]{\sphinxcrossref{\DUrole{doc}{lj/cut/tip4p/long/soft}}}} \sphinxhyphen{} LJ with cutoff Coulomb for TIP4P water with a soft core

\item {} 
\sphinxAtStartPar
{\hyperref[\detokenize{pair_lj_expand::doc}]{\sphinxcrossref{\DUrole{doc}{lj/expand}}}} \sphinxhyphen{} Lennard\sphinxhyphen{}Jones for variable size particles

\item {} 
\sphinxAtStartPar
{\hyperref[\detokenize{pair_lj_expand::doc}]{\sphinxcrossref{\DUrole{doc}{lj/expand/coul/long}}}} \sphinxhyphen{} Lennard\sphinxhyphen{}Jones for variable size particles with long\sphinxhyphen{}range Coulomb

\item {} 
\sphinxAtStartPar
{\hyperref[\detokenize{pair_lj_expand_sphere::doc}]{\sphinxcrossref{\DUrole{doc}{lj/expand/sphere}}}} \sphinxhyphen{} Variable size LJ where per\sphinxhyphen{}atom radius is used as delta (size)

\item {} 
\sphinxAtStartPar
{\hyperref[\detokenize{pair_gromacs::doc}]{\sphinxcrossref{\DUrole{doc}{lj/gromacs}}}} \sphinxhyphen{} GROMACS\sphinxhyphen{}style Lennard\sphinxhyphen{}Jones potential

\item {} 
\sphinxAtStartPar
{\hyperref[\detokenize{pair_gromacs::doc}]{\sphinxcrossref{\DUrole{doc}{lj/gromacs/coul/gromacs}}}} \sphinxhyphen{} GROMACS\sphinxhyphen{}style LJ and Coulomb potential

\item {} 
\sphinxAtStartPar
{\hyperref[\detokenize{pair_lj_long::doc}]{\sphinxcrossref{\DUrole{doc}{lj/long/coul/long}}}} \sphinxhyphen{} long\sphinxhyphen{}range LJ and long\sphinxhyphen{}range Coulomb

\item {} 
\sphinxAtStartPar
{\hyperref[\detokenize{pair_dielectric::doc}]{\sphinxcrossref{\DUrole{doc}{lj/long/coul/long/dielectric}}}} \sphinxhyphen{}

\item {} 
\sphinxAtStartPar
{\hyperref[\detokenize{pair_dipole::doc}]{\sphinxcrossref{\DUrole{doc}{lj/long/dipole/long}}}} \sphinxhyphen{} long\sphinxhyphen{}range LJ and long\sphinxhyphen{}range point dipoles

\item {} 
\sphinxAtStartPar
{\hyperref[\detokenize{pair_lj_long::doc}]{\sphinxcrossref{\DUrole{doc}{lj/long/tip4p/long}}}} \sphinxhyphen{} long\sphinxhyphen{}range LJ and long\sphinxhyphen{}range Coulomb for TIP4P water

\item {} 
\sphinxAtStartPar
{\hyperref[\detokenize{pair_mdf::doc}]{\sphinxcrossref{\DUrole{doc}{lj/mdf}}}} \sphinxhyphen{} LJ potential with a taper function

\item {} 
\sphinxAtStartPar
{\hyperref[\detokenize{pair_lj_relres::doc}]{\sphinxcrossref{\DUrole{doc}{lj/relres}}}} \sphinxhyphen{} LJ using multiscale Relative Resolution (RelRes) methodology {\hyperref[\detokenize{pair_lj_relres:chaimovich2}]{\sphinxcrossref{\DUrole{std,std-ref}{(Chaimovich)}}}}.

\item {} 
\sphinxAtStartPar
{\hyperref[\detokenize{pair_spica::doc}]{\sphinxcrossref{\DUrole{doc}{lj/spica}}}} \sphinxhyphen{} LJ for SPICA coarse\sphinxhyphen{}graining

\item {} 
\sphinxAtStartPar
{\hyperref[\detokenize{pair_spica::doc}]{\sphinxcrossref{\DUrole{doc}{lj/spica/coul/long}}}} \sphinxhyphen{} LJ for SPICA coarse\sphinxhyphen{}graining with long\sphinxhyphen{}range Coulomb

\item {} 
\sphinxAtStartPar
{\hyperref[\detokenize{pair_spica::doc}]{\sphinxcrossref{\DUrole{doc}{lj/spica/coul/msm}}}} \sphinxhyphen{} LJ for SPICA coarse\sphinxhyphen{}graining with long\sphinxhyphen{}range Coulomb via MSM

\item {} 
\sphinxAtStartPar
{\hyperref[\detokenize{pair_dipole::doc}]{\sphinxcrossref{\DUrole{doc}{lj/sf/dipole/sf}}}} \sphinxhyphen{} LJ with dipole interaction with shifted forces

\item {} 
\sphinxAtStartPar
{\hyperref[\detokenize{pair_lj_smooth::doc}]{\sphinxcrossref{\DUrole{doc}{lj/smooth}}}} \sphinxhyphen{} smoothed Lennard\sphinxhyphen{}Jones potential

\item {} 
\sphinxAtStartPar
{\hyperref[\detokenize{pair_lj_smooth_linear::doc}]{\sphinxcrossref{\DUrole{doc}{lj/smooth/linear}}}} \sphinxhyphen{} linear smoothed LJ potential

\item {} 
\sphinxAtStartPar
{\hyperref[\detokenize{pair_lj_switch3_coulgauss_long::doc}]{\sphinxcrossref{\DUrole{doc}{lj/switch3/coulgauss/long}}}} \sphinxhyphen{} smoothed LJ vdW potential with Gaussian electrostatics

\item {} 
\sphinxAtStartPar
{\hyperref[\detokenize{pair_lj96::doc}]{\sphinxcrossref{\DUrole{doc}{lj96/cut}}}} \sphinxhyphen{} Lennard\sphinxhyphen{}Jones 9/6 potential

\item {} 
\sphinxAtStartPar
{\hyperref[\detokenize{pair_local_density::doc}]{\sphinxcrossref{\DUrole{doc}{local/density}}}} \sphinxhyphen{} Generalized basic local density potential

\item {} 
\sphinxAtStartPar
{\hyperref[\detokenize{pair_lubricate::doc}]{\sphinxcrossref{\DUrole{doc}{lubricate}}}} \sphinxhyphen{} Hydrodynamic lubrication forces

\item {} 
\sphinxAtStartPar
{\hyperref[\detokenize{pair_lubricate::doc}]{\sphinxcrossref{\DUrole{doc}{lubricate/poly}}}} \sphinxhyphen{} Hydrodynamic lubrication forces with polydispersity

\item {} 
\sphinxAtStartPar
{\hyperref[\detokenize{pair_lubricateU::doc}]{\sphinxcrossref{\DUrole{doc}{lubricateU}}}} \sphinxhyphen{} Hydrodynamic lubrication forces for Fast Lubrication Dynamics

\item {} 
\sphinxAtStartPar
{\hyperref[\detokenize{pair_lubricateU::doc}]{\sphinxcrossref{\DUrole{doc}{lubricateU/poly}}}} \sphinxhyphen{} Hydrodynamic lubrication forces for Fast Lubrication with polydispersity

\item {} 
\sphinxAtStartPar
{\hyperref[\detokenize{pair_mesodpd::doc}]{\sphinxcrossref{\DUrole{doc}{mdpd}}}} \sphinxhyphen{} mDPD particle interactions

\item {} 
\sphinxAtStartPar
{\hyperref[\detokenize{pair_mesodpd::doc}]{\sphinxcrossref{\DUrole{doc}{mdpd/rhosum}}}} \sphinxhyphen{} mDPD particle interactions for mass density

\item {} 
\sphinxAtStartPar
{\hyperref[\detokenize{pair_meam::doc}]{\sphinxcrossref{\DUrole{doc}{meam}}}} \sphinxhyphen{} Modified embedded atom method (MEAM)

\item {} 
\sphinxAtStartPar
{\hyperref[\detokenize{pair_meam::doc}]{\sphinxcrossref{\DUrole{doc}{meam/ms}}}} \sphinxhyphen{} Multi\sphinxhyphen{}state modified embedded atom method (MS\sphinxhyphen{}MEAM)

\item {} 
\sphinxAtStartPar
{\hyperref[\detokenize{pair_meam_spline::doc}]{\sphinxcrossref{\DUrole{doc}{meam/spline}}}} \sphinxhyphen{} Splined version of MEAM

\item {} 
\sphinxAtStartPar
{\hyperref[\detokenize{pair_meam_sw_spline::doc}]{\sphinxcrossref{\DUrole{doc}{meam/sw/spline}}}} \sphinxhyphen{} Splined version of MEAM with a Stillinger\sphinxhyphen{}Weber term

\item {} 
\sphinxAtStartPar
{\hyperref[\detokenize{pair_mesocnt::doc}]{\sphinxcrossref{\DUrole{doc}{mesocnt}}}} \sphinxhyphen{} Mesoscopic vdW potential for (carbon) nanotubes

\item {} 
\sphinxAtStartPar
{\hyperref[\detokenize{pair_mesocnt::doc}]{\sphinxcrossref{\DUrole{doc}{mesocnt/viscous}}}} \sphinxhyphen{} Mesoscopic vdW potential for (carbon) nanotubes with friction

\item {} 
\sphinxAtStartPar
{\hyperref[\detokenize{pair_mgpt::doc}]{\sphinxcrossref{\DUrole{doc}{mgpt}}}} \sphinxhyphen{} Simplified model generalized pseudopotential theory (MGPT) potential

\item {} 
\sphinxAtStartPar
{\hyperref[\detokenize{pair_mie::doc}]{\sphinxcrossref{\DUrole{doc}{mie/cut}}}} \sphinxhyphen{} Mie potential

\item {} 
\sphinxAtStartPar
{\hyperref[\detokenize{pair_mliap::doc}]{\sphinxcrossref{\DUrole{doc}{mliap}}}} \sphinxhyphen{} Multiple styles of machine\sphinxhyphen{}learning potential

\item {} 
\sphinxAtStartPar
{\hyperref[\detokenize{pair_lj_switch3_coulgauss_long::doc}]{\sphinxcrossref{\DUrole{doc}{mm3/switch3/coulgauss/long}}}} \sphinxhyphen{} Smoothed MM3 vdW potential with Gaussian electrostatics

\item {} 
\sphinxAtStartPar
{\hyperref[\detokenize{pair_momb::doc}]{\sphinxcrossref{\DUrole{doc}{momb}}}} \sphinxhyphen{} Many\sphinxhyphen{}Body Metal\sphinxhyphen{}Organic (MOMB) force field

\item {} 
\sphinxAtStartPar
{\hyperref[\detokenize{pair_morse::doc}]{\sphinxcrossref{\DUrole{doc}{morse}}}} \sphinxhyphen{} Morse potential

\item {} 
\sphinxAtStartPar
{\hyperref[\detokenize{pair_morse::doc}]{\sphinxcrossref{\DUrole{doc}{morse/smooth/linear}}}} \sphinxhyphen{} Linear smoothed Morse potential

\item {} 
\sphinxAtStartPar
{\hyperref[\detokenize{pair_morse::doc}]{\sphinxcrossref{\DUrole{doc}{morse/soft}}}} \sphinxhyphen{} Morse potential with a soft core

\item {} 
\sphinxAtStartPar
{\hyperref[\detokenize{pair_multi_lucy::doc}]{\sphinxcrossref{\DUrole{doc}{multi/lucy}}}} \sphinxhyphen{} DPD potential with density\sphinxhyphen{}dependent force

\item {} 
\sphinxAtStartPar
{\hyperref[\detokenize{pair_multi_lucy_rx::doc}]{\sphinxcrossref{\DUrole{doc}{multi/lucy/rx}}}} \sphinxhyphen{} reactive DPD potential with density\sphinxhyphen{}dependent force

\item {} 
\sphinxAtStartPar
{\hyperref[\detokenize{pair_nb3b::doc}]{\sphinxcrossref{\DUrole{doc}{nb3b/harmonic}}}} \sphinxhyphen{} Non\sphinxhyphen{}bonded 3\sphinxhyphen{}body harmonic potential

\item {} 
\sphinxAtStartPar
{\hyperref[\detokenize{pair_nb3b::doc}]{\sphinxcrossref{\DUrole{doc}{nb3b/screened}}}} \sphinxhyphen{} Non\sphinxhyphen{}bonded 3\sphinxhyphen{}body screened harmonic potential

\item {} 
\sphinxAtStartPar
{\hyperref[\detokenize{pair_nm::doc}]{\sphinxcrossref{\DUrole{doc}{nm/cut}}}} \sphinxhyphen{} N\sphinxhyphen{}M potential

\item {} 
\sphinxAtStartPar
{\hyperref[\detokenize{pair_nm::doc}]{\sphinxcrossref{\DUrole{doc}{nm/cut/coul/cut}}}} \sphinxhyphen{} N\sphinxhyphen{}M potential with cutoff Coulomb

\item {} 
\sphinxAtStartPar
{\hyperref[\detokenize{pair_nm::doc}]{\sphinxcrossref{\DUrole{doc}{nm/cut/coul/long}}}} \sphinxhyphen{} N\sphinxhyphen{}M potential with long\sphinxhyphen{}range Coulomb

\item {} 
\sphinxAtStartPar
{\hyperref[\detokenize{pair_nm::doc}]{\sphinxcrossref{\DUrole{doc}{nm/cut/split}}}} \sphinxhyphen{} Split 12\sphinxhyphen{}6 Lennard\sphinxhyphen{}Jones and N\sphinxhyphen{}M potential

\item {} 
\sphinxAtStartPar
{\hyperref[\detokenize{pair_oxdna::doc}]{\sphinxcrossref{\DUrole{doc}{oxdna/coaxstk}}}} \sphinxhyphen{}

\item {} 
\sphinxAtStartPar
{\hyperref[\detokenize{pair_oxdna::doc}]{\sphinxcrossref{\DUrole{doc}{oxdna/excv}}}} \sphinxhyphen{}

\item {} 
\sphinxAtStartPar
{\hyperref[\detokenize{pair_oxdna::doc}]{\sphinxcrossref{\DUrole{doc}{oxdna/hbond}}}} \sphinxhyphen{}

\item {} 
\sphinxAtStartPar
{\hyperref[\detokenize{pair_oxdna::doc}]{\sphinxcrossref{\DUrole{doc}{oxdna/stk}}}} \sphinxhyphen{}

\item {} 
\sphinxAtStartPar
{\hyperref[\detokenize{pair_oxdna::doc}]{\sphinxcrossref{\DUrole{doc}{oxdna/xstk}}}} \sphinxhyphen{}

\item {} 
\sphinxAtStartPar
{\hyperref[\detokenize{pair_oxdna2::doc}]{\sphinxcrossref{\DUrole{doc}{oxdna2/coaxstk}}}} \sphinxhyphen{}

\item {} 
\sphinxAtStartPar
{\hyperref[\detokenize{pair_oxdna2::doc}]{\sphinxcrossref{\DUrole{doc}{oxdna2/dh}}}} \sphinxhyphen{}

\item {} 
\sphinxAtStartPar
{\hyperref[\detokenize{pair_oxdna2::doc}]{\sphinxcrossref{\DUrole{doc}{oxdna2/excv}}}} \sphinxhyphen{}

\item {} 
\sphinxAtStartPar
{\hyperref[\detokenize{pair_oxdna2::doc}]{\sphinxcrossref{\DUrole{doc}{oxdna2/hbond}}}} \sphinxhyphen{}

\item {} 
\sphinxAtStartPar
{\hyperref[\detokenize{pair_oxdna2::doc}]{\sphinxcrossref{\DUrole{doc}{oxdna2/stk}}}} \sphinxhyphen{}

\item {} 
\sphinxAtStartPar
{\hyperref[\detokenize{pair_oxdna2::doc}]{\sphinxcrossref{\DUrole{doc}{oxdna2/xstk}}}} \sphinxhyphen{}

\item {} 
\sphinxAtStartPar
{\hyperref[\detokenize{pair_oxrna2::doc}]{\sphinxcrossref{\DUrole{doc}{oxrna2/coaxstk}}}} \sphinxhyphen{}

\item {} 
\sphinxAtStartPar
{\hyperref[\detokenize{pair_oxrna2::doc}]{\sphinxcrossref{\DUrole{doc}{oxrna2/dh}}}} \sphinxhyphen{}

\item {} 
\sphinxAtStartPar
{\hyperref[\detokenize{pair_oxrna2::doc}]{\sphinxcrossref{\DUrole{doc}{oxrna2/excv}}}} \sphinxhyphen{}

\item {} 
\sphinxAtStartPar
{\hyperref[\detokenize{pair_oxrna2::doc}]{\sphinxcrossref{\DUrole{doc}{oxrna2/hbond}}}} \sphinxhyphen{}

\item {} 
\sphinxAtStartPar
{\hyperref[\detokenize{pair_oxrna2::doc}]{\sphinxcrossref{\DUrole{doc}{oxrna2/stk}}}} \sphinxhyphen{}

\item {} 
\sphinxAtStartPar
{\hyperref[\detokenize{pair_oxrna2::doc}]{\sphinxcrossref{\DUrole{doc}{oxrna2/xstk}}}} \sphinxhyphen{}

\item {} 
\sphinxAtStartPar
{\hyperref[\detokenize{pair_pace::doc}]{\sphinxcrossref{\DUrole{doc}{pace}}}} \sphinxhyphen{} Atomic Cluster Expansion (ACE) machine\sphinxhyphen{}learning potential

\item {} 
\sphinxAtStartPar
{\hyperref[\detokenize{pair_pace::doc}]{\sphinxcrossref{\DUrole{doc}{pace/extrapolation}}}} \sphinxhyphen{} Atomic Cluster Expansion (ACE) machine\sphinxhyphen{}learning potential with extrapolation grades

\item {} 
\sphinxAtStartPar
{\hyperref[\detokenize{pair_pedone::doc}]{\sphinxcrossref{\DUrole{doc}{pedone}}}} \sphinxhyphen{} Pedone (PMMCS) potential (non\sphinxhyphen{}Coulomb part)

\item {} 
\sphinxAtStartPar
{\hyperref[\detokenize{pair_pod::doc}]{\sphinxcrossref{\DUrole{doc}{pod}}}} \sphinxhyphen{} Proper orthogonal decomposition (POD) machine\sphinxhyphen{}learning potential

\item {} 
\sphinxAtStartPar
{\hyperref[\detokenize{pair_peri::doc}]{\sphinxcrossref{\DUrole{doc}{peri/eps}}}} \sphinxhyphen{} Peridynamic EPS potential

\item {} 
\sphinxAtStartPar
{\hyperref[\detokenize{pair_peri::doc}]{\sphinxcrossref{\DUrole{doc}{peri/lps}}}} \sphinxhyphen{} Peridynamic LPS potential

\item {} 
\sphinxAtStartPar
{\hyperref[\detokenize{pair_peri::doc}]{\sphinxcrossref{\DUrole{doc}{peri/pmb}}}} \sphinxhyphen{} Peridynamic PMB potential

\item {} 
\sphinxAtStartPar
{\hyperref[\detokenize{pair_peri::doc}]{\sphinxcrossref{\DUrole{doc}{peri/ves}}}} \sphinxhyphen{} Peridynamic VES potential

\item {} 
\sphinxAtStartPar
{\hyperref[\detokenize{pair_polymorphic::doc}]{\sphinxcrossref{\DUrole{doc}{polymorphic}}}} \sphinxhyphen{} Polymorphic 3\sphinxhyphen{}body potential

\item {} 
\sphinxAtStartPar
{\hyperref[\detokenize{pair_python::doc}]{\sphinxcrossref{\DUrole{doc}{python}}}} \sphinxhyphen{}

\item {} 
\sphinxAtStartPar
{\hyperref[\detokenize{pair_quip::doc}]{\sphinxcrossref{\DUrole{doc}{quip}}}} \sphinxhyphen{}

\item {} 
\sphinxAtStartPar
{\hyperref[\detokenize{pair_rann::doc}]{\sphinxcrossref{\DUrole{doc}{rann}}}} \sphinxhyphen{}

\item {} 
\sphinxAtStartPar
{\hyperref[\detokenize{pair_reaxff::doc}]{\sphinxcrossref{\DUrole{doc}{reaxff}}}} \sphinxhyphen{} ReaxFF potential

\item {} 
\sphinxAtStartPar
{\hyperref[\detokenize{pair_airebo::doc}]{\sphinxcrossref{\DUrole{doc}{rebo}}}} \sphinxhyphen{} Second generation REBO potential of Brenner

\item {} 
\sphinxAtStartPar
{\hyperref[\detokenize{pair_rebomos::doc}]{\sphinxcrossref{\DUrole{doc}{rebomos}}}} \sphinxhyphen{} REBOMoS potential for MoS2

\item {} 
\sphinxAtStartPar
{\hyperref[\detokenize{pair_rheo::doc}]{\sphinxcrossref{\DUrole{doc}{rheo}}}} \sphinxhyphen{} fluid interactions in RHEO package

\item {} 
\sphinxAtStartPar
{\hyperref[\detokenize{pair_rheo_solid::doc}]{\sphinxcrossref{\DUrole{doc}{rheo/solid}}}} \sphinxhyphen{} solid interactions in RHEO package

\item {} 
\sphinxAtStartPar
{\hyperref[\detokenize{pair_resquared::doc}]{\sphinxcrossref{\DUrole{doc}{resquared}}}} \sphinxhyphen{} Everaers RE\sphinxhyphen{}Squared ellipsoidal potential

\item {} 
\sphinxAtStartPar
{\hyperref[\detokenize{pair_saip_metal::doc}]{\sphinxcrossref{\DUrole{doc}{saip/metal}}}} \sphinxhyphen{} Interlayer potential for hetero\sphinxhyphen{}junctions formed with hexagonal 2D materials and metal surfaces

\item {} 
\sphinxAtStartPar
{\hyperref[\detokenize{pair_sdpd_taitwater_isothermal::doc}]{\sphinxcrossref{\DUrole{doc}{sdpd/taitwater/isothermal}}}} \sphinxhyphen{} Smoothed dissipative particle dynamics for water at isothermal conditions

\item {} 
\sphinxAtStartPar
{\hyperref[\detokenize{pair_smatb::doc}]{\sphinxcrossref{\DUrole{doc}{smatb}}}} \sphinxhyphen{} Second Moment Approximation to the Tight Binding

\item {} 
\sphinxAtStartPar
{\hyperref[\detokenize{pair_smatb::doc}]{\sphinxcrossref{\DUrole{doc}{smatb/single}}}} \sphinxhyphen{} Second Moment Approximation to the Tight Binding for single\sphinxhyphen{}element systems

\item {} 
\sphinxAtStartPar
{\hyperref[\detokenize{pair_smd_hertz::doc}]{\sphinxcrossref{\DUrole{doc}{smd/hertz}}}} \sphinxhyphen{}

\item {} 
\sphinxAtStartPar
{\hyperref[\detokenize{pair_smd_tlsph::doc}]{\sphinxcrossref{\DUrole{doc}{smd/tlsph}}}} \sphinxhyphen{}

\item {} 
\sphinxAtStartPar
{\hyperref[\detokenize{pair_smd_triangulated_surface::doc}]{\sphinxcrossref{\DUrole{doc}{smd/tri\_surface}}}} \sphinxhyphen{}

\item {} 
\sphinxAtStartPar
{\hyperref[\detokenize{pair_smd_ulsph::doc}]{\sphinxcrossref{\DUrole{doc}{smd/ulsph}}}} \sphinxhyphen{}

\item {} 
\sphinxAtStartPar
{\hyperref[\detokenize{pair_smtbq::doc}]{\sphinxcrossref{\DUrole{doc}{smtbq}}}} \sphinxhyphen{}

\item {} 
\sphinxAtStartPar
{\hyperref[\detokenize{pair_snap::doc}]{\sphinxcrossref{\DUrole{doc}{snap}}}} \sphinxhyphen{} SNAP machine\sphinxhyphen{}learning potential

\item {} 
\sphinxAtStartPar
{\hyperref[\detokenize{pair_soft::doc}]{\sphinxcrossref{\DUrole{doc}{soft}}}} \sphinxhyphen{} Soft (cosine) potential

\item {} 
\sphinxAtStartPar
{\hyperref[\detokenize{pair_sph_heatconduction::doc}]{\sphinxcrossref{\DUrole{doc}{sph/heatconduction}}}} \sphinxhyphen{}

\item {} 
\sphinxAtStartPar
{\hyperref[\detokenize{pair_sph_idealgas::doc}]{\sphinxcrossref{\DUrole{doc}{sph/idealgas}}}} \sphinxhyphen{}

\item {} 
\sphinxAtStartPar
{\hyperref[\detokenize{pair_sph_lj::doc}]{\sphinxcrossref{\DUrole{doc}{sph/lj}}}} \sphinxhyphen{}

\item {} 
\sphinxAtStartPar
{\hyperref[\detokenize{pair_sph_rhosum::doc}]{\sphinxcrossref{\DUrole{doc}{sph/rhosum}}}} \sphinxhyphen{}

\item {} 
\sphinxAtStartPar
{\hyperref[\detokenize{pair_sph_taitwater::doc}]{\sphinxcrossref{\DUrole{doc}{sph/taitwater}}}} \sphinxhyphen{}

\item {} 
\sphinxAtStartPar
{\hyperref[\detokenize{pair_sph_taitwater_morris::doc}]{\sphinxcrossref{\DUrole{doc}{sph/taitwater/morris}}}} \sphinxhyphen{}

\item {} 
\sphinxAtStartPar
{\hyperref[\detokenize{pair_spin_dipole::doc}]{\sphinxcrossref{\DUrole{doc}{spin/dipole/cut}}}} \sphinxhyphen{}

\item {} 
\sphinxAtStartPar
{\hyperref[\detokenize{pair_spin_dipole::doc}]{\sphinxcrossref{\DUrole{doc}{spin/dipole/long}}}} \sphinxhyphen{}

\item {} 
\sphinxAtStartPar
{\hyperref[\detokenize{pair_spin_dmi::doc}]{\sphinxcrossref{\DUrole{doc}{spin/dmi}}}} \sphinxhyphen{}

\item {} 
\sphinxAtStartPar
{\hyperref[\detokenize{pair_spin_exchange::doc}]{\sphinxcrossref{\DUrole{doc}{spin/exchange}}}} \sphinxhyphen{}

\item {} 
\sphinxAtStartPar
{\hyperref[\detokenize{pair_spin_exchange::doc}]{\sphinxcrossref{\DUrole{doc}{spin/exchange/biquadratic}}}} \sphinxhyphen{}

\item {} 
\sphinxAtStartPar
{\hyperref[\detokenize{pair_spin_magelec::doc}]{\sphinxcrossref{\DUrole{doc}{spin/magelec}}}} \sphinxhyphen{}

\item {} 
\sphinxAtStartPar
{\hyperref[\detokenize{pair_spin_neel::doc}]{\sphinxcrossref{\DUrole{doc}{spin/neel}}}} \sphinxhyphen{}

\item {} 
\sphinxAtStartPar
{\hyperref[\detokenize{pair_srp::doc}]{\sphinxcrossref{\DUrole{doc}{srp}}}} \sphinxhyphen{}

\item {} 
\sphinxAtStartPar
{\hyperref[\detokenize{pair_srp::doc}]{\sphinxcrossref{\DUrole{doc}{srp/react}}}} \sphinxhyphen{}

\item {} 
\sphinxAtStartPar
{\hyperref[\detokenize{pair_sw::doc}]{\sphinxcrossref{\DUrole{doc}{sw}}}} \sphinxhyphen{} Stillinger\sphinxhyphen{}Weber 3\sphinxhyphen{}body potential

\item {} 
\sphinxAtStartPar
{\hyperref[\detokenize{pair_sw_angle_table::doc}]{\sphinxcrossref{\DUrole{doc}{sw/angle/table}}}} \sphinxhyphen{} Stillinger\sphinxhyphen{}Weber potential with tabulated angular term

\item {} 
\sphinxAtStartPar
{\hyperref[\detokenize{pair_sw::doc}]{\sphinxcrossref{\DUrole{doc}{sw/mod}}}} \sphinxhyphen{} modified Stillinger\sphinxhyphen{}Weber 3\sphinxhyphen{}body potential

\item {} 
\sphinxAtStartPar
{\hyperref[\detokenize{pair_table::doc}]{\sphinxcrossref{\DUrole{doc}{table}}}} \sphinxhyphen{} tabulated pair potential

\item {} 
\sphinxAtStartPar
{\hyperref[\detokenize{pair_table_rx::doc}]{\sphinxcrossref{\DUrole{doc}{table/rx}}}} \sphinxhyphen{}

\item {} 
\sphinxAtStartPar
{\hyperref[\detokenize{pair_mesodpd::doc}]{\sphinxcrossref{\DUrole{doc}{tdpd}}}} \sphinxhyphen{} tDPD particle interactions

\item {} 
\sphinxAtStartPar
{\hyperref[\detokenize{pair_tersoff::doc}]{\sphinxcrossref{\DUrole{doc}{tersoff}}}} \sphinxhyphen{} Tersoff 3\sphinxhyphen{}body potential

\item {} 
\sphinxAtStartPar
{\hyperref[\detokenize{pair_tersoff_mod::doc}]{\sphinxcrossref{\DUrole{doc}{tersoff/mod}}}} \sphinxhyphen{} modified Tersoff 3\sphinxhyphen{}body potential

\item {} 
\sphinxAtStartPar
{\hyperref[\detokenize{pair_tersoff_mod::doc}]{\sphinxcrossref{\DUrole{doc}{tersoff/mod/c}}}} \sphinxhyphen{}

\item {} 
\sphinxAtStartPar
{\hyperref[\detokenize{pair_tersoff::doc}]{\sphinxcrossref{\DUrole{doc}{tersoff/table}}}} \sphinxhyphen{}

\item {} 
\sphinxAtStartPar
{\hyperref[\detokenize{pair_tersoff_zbl::doc}]{\sphinxcrossref{\DUrole{doc}{tersoff/zbl}}}} \sphinxhyphen{} Tersoff/ZBL 3\sphinxhyphen{}body potential

\item {} 
\sphinxAtStartPar
{\hyperref[\detokenize{pair_thole::doc}]{\sphinxcrossref{\DUrole{doc}{thole}}}} \sphinxhyphen{} Coulomb interactions with thole damping

\item {} 
\sphinxAtStartPar
{\hyperref[\detokenize{pair_threebody_table::doc}]{\sphinxcrossref{\DUrole{doc}{threebody/table}}}} \sphinxhyphen{} generic tabulated three\sphinxhyphen{}body potential

\item {} 
\sphinxAtStartPar
{\hyperref[\detokenize{pair_coul::doc}]{\sphinxcrossref{\DUrole{doc}{tip4p/cut}}}} \sphinxhyphen{} Coulomb for TIP4P water w/out LJ

\item {} 
\sphinxAtStartPar
{\hyperref[\detokenize{pair_coul::doc}]{\sphinxcrossref{\DUrole{doc}{tip4p/long}}}} \sphinxhyphen{} long\sphinxhyphen{}range Coulomb for TIP4P water w/out LJ

\item {} 
\sphinxAtStartPar
{\hyperref[\detokenize{pair_fep_soft::doc}]{\sphinxcrossref{\DUrole{doc}{tip4p/long/soft}}}} \sphinxhyphen{}

\item {} 
\sphinxAtStartPar
{\hyperref[\detokenize{pair_tracker::doc}]{\sphinxcrossref{\DUrole{doc}{tracker}}}} \sphinxhyphen{} monitor information about pairwise interactions

\item {} 
\sphinxAtStartPar
{\hyperref[\detokenize{pair_tri_lj::doc}]{\sphinxcrossref{\DUrole{doc}{tri/lj}}}} \sphinxhyphen{} LJ potential between triangles

\item {} 
\sphinxAtStartPar
{\hyperref[\detokenize{pair_ufm::doc}]{\sphinxcrossref{\DUrole{doc}{ufm}}}} \sphinxhyphen{}

\item {} 
\sphinxAtStartPar
{\hyperref[\detokenize{pair_uf3::doc}]{\sphinxcrossref{\DUrole{doc}{uf3}}}} \sphinxhyphen{} UF3 machine\sphinxhyphen{}learning potential

\item {} 
\sphinxAtStartPar
{\hyperref[\detokenize{pair_vashishta::doc}]{\sphinxcrossref{\DUrole{doc}{vashishta}}}} \sphinxhyphen{} Vashishta 2\sphinxhyphen{}body and 3\sphinxhyphen{}body potential

\item {} 
\sphinxAtStartPar
{\hyperref[\detokenize{pair_vashishta::doc}]{\sphinxcrossref{\DUrole{doc}{vashishta/table}}}} \sphinxhyphen{}

\item {} 
\sphinxAtStartPar
{\hyperref[\detokenize{pair_wf_cut::doc}]{\sphinxcrossref{\DUrole{doc}{wf/cut}}}} \sphinxhyphen{} Wang\sphinxhyphen{}Frenkel Potential for short\sphinxhyphen{}ranged interactions

\item {} 
\sphinxAtStartPar
{\hyperref[\detokenize{pair_ylz::doc}]{\sphinxcrossref{\DUrole{doc}{ylz}}}} \sphinxhyphen{} Yuan\sphinxhyphen{}Li\sphinxhyphen{}Zhang Potential for anisotropic interactions

\item {} 
\sphinxAtStartPar
{\hyperref[\detokenize{pair_yukawa::doc}]{\sphinxcrossref{\DUrole{doc}{yukawa}}}} \sphinxhyphen{} Yukawa potential

\item {} 
\sphinxAtStartPar
{\hyperref[\detokenize{pair_yukawa_colloid::doc}]{\sphinxcrossref{\DUrole{doc}{yukawa/colloid}}}} \sphinxhyphen{} screened Yukawa potential for finite\sphinxhyphen{}size particles

\item {} 
\sphinxAtStartPar
{\hyperref[\detokenize{pair_zbl::doc}]{\sphinxcrossref{\DUrole{doc}{zbl}}}} \sphinxhyphen{} Ziegler\sphinxhyphen{}Biersack\sphinxhyphen{}Littmark potential

\end{itemize}


\bigskip\hrule\bigskip



\subsection{Restrictions}
\label{\detokenize{pair_style:restrictions}}
\sphinxAtStartPar
This command must be used before any coefficients are set by the
{\hyperref[\detokenize{pair_coeff::doc}]{\sphinxcrossref{\DUrole{doc}{pair\_coeff}}}}, {\hyperref[\detokenize{read_data::doc}]{\sphinxcrossref{\DUrole{doc}{read\_data}}}}, or
{\hyperref[\detokenize{read_restart::doc}]{\sphinxcrossref{\DUrole{doc}{read\_restart}}}} commands.

\sphinxAtStartPar
Some pair styles are part of specific packages.  They are only enabled
if LAMMPS was built with that package.  See the {\hyperref[\detokenize{Build_package::doc}]{\sphinxcrossref{\DUrole{doc}{Build package}}}} page for more info.  The doc pages for
individual pair potentials tell if it is part of a package.


\subsection{Related commands}
\label{\detokenize{pair_style:related-commands}}
\sphinxAtStartPar
{\hyperref[\detokenize{pair_coeff::doc}]{\sphinxcrossref{\DUrole{doc}{pair\_coeff}}}}, {\hyperref[\detokenize{read_data::doc}]{\sphinxcrossref{\DUrole{doc}{read\_data}}}},
{\hyperref[\detokenize{pair_modify::doc}]{\sphinxcrossref{\DUrole{doc}{pair\_modify}}}}, {\hyperref[\detokenize{kspace_style::doc}]{\sphinxcrossref{\DUrole{doc}{kspace\_style}}}},
{\hyperref[\detokenize{dielectric::doc}]{\sphinxcrossref{\DUrole{doc}{dielectric}}}}, {\hyperref[\detokenize{pair_write::doc}]{\sphinxcrossref{\DUrole{doc}{pair\_write}}}}


\subsection{Default}
\label{\detokenize{pair_style:default}}
\begin{sphinxVerbatim}[commandchars=\\\{\}]
\PYG{k}{pair\PYGZus{}style}\PYG{+w}{ }\PYG{n}{none}
\end{sphinxVerbatim}

\sphinxstepscope

\index{pair\_write@\spxentry{pair\_write}}\ignorespaces 

\section{pair\_write command}
\label{\detokenize{pair_write:pair-write-command}}\label{\detokenize{pair_write:index-0}}\label{\detokenize{pair_write::doc}}

\subsection{Syntax}
\label{\detokenize{pair_write:syntax}}
\begin{sphinxVerbatim}[commandchars=\\\{\}]
\PYG{k}{pair\PYGZus{}write}\PYG{+w}{ }\PYG{n}{itype}\PYG{+w}{ }\PYG{n}{jtype}\PYG{+w}{ }\PYG{n}{N}\PYG{+w}{ }\PYG{n}{style}\PYG{+w}{ }\PYG{n}{inner}\PYG{+w}{ }\PYG{n}{outer}\PYG{+w}{ }\PYG{n}{file}\PYG{+w}{ }\PYG{n}{keyword}\PYG{+w}{ }\PYG{n}{Qi}\PYG{+w}{ }\PYG{n}{Qj}
\end{sphinxVerbatim}
\begin{itemize}
\item {} 
\sphinxAtStartPar
itype,jtype = 2 atom types (numeric or type label)

\item {} 
\sphinxAtStartPar
N = \# of values

\item {} 
\sphinxAtStartPar
style = \sphinxstyleemphasis{r} or \sphinxstyleemphasis{rsq} or \sphinxstyleemphasis{bitmap}

\item {} 
\sphinxAtStartPar
inner,outer = inner and outer cutoff (distance units)

\item {} 
\sphinxAtStartPar
file = name of file to write values to

\item {} 
\sphinxAtStartPar
keyword = section name in file for this set of tabulated values

\item {} 
\sphinxAtStartPar
Qi,Qj = 2 atom charges (charge units) (optional)

\end{itemize}


\subsection{Examples}
\label{\detokenize{pair_write:examples}}
\begin{sphinxVerbatim}[commandchars=\\\{\}]
\PYG{k}{pair\PYGZus{}write}\PYG{+w}{ }\PYG{l+m}{1}\PYG{+w}{ }\PYG{l+m}{3}\PYG{+w}{ }\PYG{l+m}{500}\PYG{+w}{ }\PYG{n}{r}\PYG{+w}{ }\PYG{l+m}{1.0}\PYG{+w}{ }\PYG{l+m}{10.0}\PYG{+w}{ }\PYG{n}{table.txt}\PYG{+w}{ }\PYG{n}{LJ}
\PYG{k}{pair\PYGZus{}write}\PYG{+w}{ }\PYG{l+m}{1}\PYG{+w}{ }\PYG{l+m}{1}\PYG{+w}{ }\PYG{l+m}{1000}\PYG{+w}{ }\PYG{n}{rsq}\PYG{+w}{ }\PYG{l+m}{2.0}\PYG{+w}{ }\PYG{l+m}{8.0}\PYG{+w}{ }\PYG{n}{table.txt}\PYG{+w}{ }\PYG{n}{Yukawa\PYGZus{}1\PYGZus{}1}\PYG{+w}{ }\PYG{o}{\PYGZhy{}}\PYG{l+m}{0.5}\PYG{+w}{ }\PYG{l+m}{0.5}

\PYG{k}{labelmap}\PYG{+w}{ }\PYG{n}{atom}\PYG{+w}{ }\PYG{l+m}{1}\PYG{+w}{ }\PYG{n}{C}\PYG{+w}{ }\PYG{l+m}{2}\PYG{+w}{ }\PYG{n}{H}
\PYG{k}{pair\PYGZus{}write}\PYG{+w}{ }\PYG{n}{C}\PYG{+w}{ }\PYG{n}{H}\PYG{+w}{ }\PYG{l+m}{500}\PYG{+w}{ }\PYG{n}{r}\PYG{+w}{ }\PYG{l+m}{1.0}\PYG{+w}{ }\PYG{l+m}{10.0}\PYG{+w}{ }\PYG{n}{table.txt}\PYG{+w}{ }\PYG{n}{LJ}
\end{sphinxVerbatim}


\subsection{Description}
\label{\detokenize{pair_write:description}}
\sphinxAtStartPar
Write energy and force values to a file as a function of distance for
the currently defined pair potential.  This is useful for plotting the
potential function or otherwise debugging its values.  If the file
already exists, the table of values is appended to the end of the file
to allow multiple tables of energy and force to be included in one
file.  In case a new file is created, the first line will be a comment
containing a “DATE:” and “UNITS:” tag with the current date and the
current {\hyperref[\detokenize{units::doc}]{\sphinxcrossref{\DUrole{doc}{units}}}} setting as argument.  For subsequent
invocations of the pair\_write command, the current units setting is
compared against the entry in the file, if present, and pair\_write
will refuse to add a table if the units are not the same.

\sphinxAtStartPar
The energy and force values are computed at distances from inner to
outer for 2 interacting atoms of type \sphinxstyleemphasis{itype} and \sphinxstyleemphasis{jtype}, using the
appropriate {\hyperref[\detokenize{pair_coeff::doc}]{\sphinxcrossref{\DUrole{doc}{pair\_coeff}}}} coefficients.  If the style
is \sphinxstyleemphasis{r}, then N distances are used, evenly spaced in r; if the style is
\sphinxstyleemphasis{rsq}, N distances are used, evenly spaced in r\textasciicircum{}2.

\sphinxAtStartPar
For example, for N = 7, style = \sphinxstyleemphasis{r}, inner = 1.0, and outer = 4.0,
values are computed at r = 1.0, 1.5, 2.0, 2.5, 3.0, 3.5, 4.0.

\sphinxAtStartPar
If the style is \sphinxstyleemphasis{bitmap}, then 2\textasciicircum{}N values are written to the file in a
format and order consistent with how they are read in by the
{\hyperref[\detokenize{pair_coeff::doc}]{\sphinxcrossref{\DUrole{doc}{pair\_coeff}}}} command for pair style \sphinxstyleemphasis{table}.  For
reasonable accuracy in a bitmapped table, choose N \textgreater{}= 12, an \sphinxstyleemphasis{inner}
value that is smaller than the distance of closest approach of 2
atoms, and an \sphinxstyleemphasis{outer} value \textless{}= cutoff of the potential.

\sphinxAtStartPar
If the pair potential is computed between charged atoms, the charges
of the pair of interacting atoms can optionally be specified.  If not
specified, values of Qi = Qj = 1.0 are used.

\sphinxAtStartPar
The file is written in the format used as input for the
{\hyperref[\detokenize{pair_style::doc}]{\sphinxcrossref{\DUrole{doc}{pair\_style}}}} \sphinxstyleemphasis{table} option with \sphinxstyleemphasis{keyword} as the
section name.  Each line written to the file lists an index number
(1\sphinxhyphen{}N), a distance (in distance units), an energy (in energy units),
and a force (in force units).


\subsection{Restrictions}
\label{\detokenize{pair_write:restrictions}}
\sphinxAtStartPar
All force field coefficients for pair and other kinds of interactions
must be set before this command can be invoked.

\sphinxAtStartPar
Due to how the pairwise force is computed, an inner value \textgreater{} 0.0 must
be specified even if the potential has a finite value at r = 0.0.

\sphinxAtStartPar
The \sphinxstyleemphasis{pair\_write} command can only be used for pairwise additive
interactions for which a \sphinxtitleref{Pair::single()} function can be and has
been implemented.  This excludes for example manybody potentials
or TIP4P coulomb styles.


\subsection{Related commands}
\label{\detokenize{pair_write:related-commands}}
\sphinxAtStartPar
{\hyperref[\detokenize{pair_table::doc}]{\sphinxcrossref{\DUrole{doc}{pair\_style table}}}},
{\hyperref[\detokenize{pair_style::doc}]{\sphinxcrossref{\DUrole{doc}{pair\_style}}}}, {\hyperref[\detokenize{pair_coeff::doc}]{\sphinxcrossref{\DUrole{doc}{pair\_coeff}}}}


\subsection{Default}
\label{\detokenize{pair_write:default}}
\sphinxAtStartPar
none

\sphinxstepscope

\index{partition@\spxentry{partition}}\ignorespaces 

\section{partition command}
\label{\detokenize{partition:partition-command}}\label{\detokenize{partition:index-0}}\label{\detokenize{partition::doc}}

\subsection{Syntax}
\label{\detokenize{partition:syntax}}
\begin{sphinxVerbatim}[commandchars=\\\{\}]
\PYG{k}{partition}\PYG{+w}{ }\PYG{n}{style}\PYG{+w}{ }\PYG{n}{N}\PYG{+w}{ }\PYG{n}{command}\PYG{+w}{ }\PYG{n}{...}
\end{sphinxVerbatim}
\begin{itemize}
\item {} 
\sphinxAtStartPar
style = \sphinxstyleemphasis{yes} or \sphinxstyleemphasis{no}

\item {} 
\sphinxAtStartPar
N = partition number (see asterisk form below)

\item {} 
\sphinxAtStartPar
command = any LAMMPS command

\end{itemize}


\subsection{Examples}
\label{\detokenize{partition:examples}}
\begin{sphinxVerbatim}[commandchars=\\\{\}]
\PYG{k}{partition}\PYG{+w}{ }\PYG{n}{yes}\PYG{+w}{ }\PYG{l+m}{1}\PYG{+w}{ }\PYG{n}{processors}\PYG{+w}{ }\PYG{l+m}{4}\PYG{+w}{ }\PYG{l+m}{10}\PYG{+w}{ }\PYG{l+m}{6}
\PYG{k}{partition}\PYG{+w}{ }\PYG{n}{no}\PYG{+w}{ }\PYG{l+m}{5}\PYG{+w}{ }\PYG{n}{print}\PYG{+w}{ }\PYG{l+s}{\PYGZdq{}}\PYG{l+s}{Active partition}\PYG{l+s}{\PYGZdq{}}
\PYG{k}{partition}\PYG{+w}{ }\PYG{n}{yes}\PYG{+w}{ }\PYG{o}{*}\PYG{l+m}{5}\PYG{+w}{ }\PYG{k}{fix }\PYG{n+nv+nvIdentifier}{all}\PYG{+w}{ }\PYG{n+nv+nvIdentifier}{nve}
\PYG{k}{partition}\PYG{+w}{ }\PYG{n}{yes}\PYG{+w}{ }\PYG{l+m}{6}\PYG{o}{*}\PYG{+w}{ }\PYG{k}{fix }\PYG{n+nv+nvIdentifier}{all}\PYG{+w}{ }\PYG{n+nv+nvIdentifier}{nvt}\PYG{+w}{ }\PYG{n}{temp}\PYG{+w}{ }\PYG{l+m}{1.0}\PYG{+w}{ }\PYG{l+m}{1.0}\PYG{+w}{ }\PYG{l+m}{0.1}
\end{sphinxVerbatim}


\subsection{Description}
\label{\detokenize{partition:description}}
\sphinxAtStartPar
This command invokes the specified command on a subset of the
partitions of processors you have defined via the {\hyperref[\detokenize{Run_options::doc}]{\sphinxcrossref{\DUrole{doc}{\sphinxhyphen{}partition command\sphinxhyphen{}line switch}}}}.

\sphinxAtStartPar
Normally, every input script command in your script is invoked by
every partition.  This behavior can be modified by defining world\sphinxhyphen{} or
universe\sphinxhyphen{}style {\hyperref[\detokenize{variable::doc}]{\sphinxcrossref{\DUrole{doc}{variables}}}} that have different values
for each partition.  This mechanism can be used to cause your script
to jump to different input script files on different partitions, if
such a variable is used in a {\hyperref[\detokenize{jump::doc}]{\sphinxcrossref{\DUrole{doc}{jump}}}} command.

\sphinxAtStartPar
The “partition” command is another mechanism for having as input
script operate differently on different partitions.  It is basically a
prefix on any LAMMPS command.  The command will only be invoked on
the partition(s) specified by the \sphinxstyleemphasis{style} and \sphinxstyleemphasis{N} arguments.

\sphinxAtStartPar
If the \sphinxstyleemphasis{style} is \sphinxstyleemphasis{yes}, the command will be invoked on any partition
which matches the \sphinxstyleemphasis{N} argument.  If the \sphinxstyleemphasis{style} is \sphinxstyleemphasis{no} the command
will be invoked on all the partitions which do not match the Np
argument.

\sphinxAtStartPar
Partitions are numbered from 1 to Np, where Np is the number of
partitions specified by the {\hyperref[\detokenize{Run_options::doc}]{\sphinxcrossref{\DUrole{doc}{\sphinxhyphen{}partition command\sphinxhyphen{}line switch}}}}.

\sphinxAtStartPar
\sphinxstyleemphasis{N} can be specified in one of two ways.  An explicit numeric value
can be used, as in the first example above.  Or a wild\sphinxhyphen{}card asterisk can
be used to span a range of partition numbers.  This takes the form “*”
or “*n” or “n*” or “m*n”.  An asterisk with no numeric values means
all partitions from 1 to Np.  A leading asterisk means all partitions
from 1 to n (inclusive).  A trailing asterisk means all partitions
from n to Np (inclusive).  A middle asterisk means all partitions from
m to n (inclusive).

\sphinxAtStartPar
This command can be useful for the “run\_style verlet/split” command
which imposed requirements on how the {\hyperref[\detokenize{processors::doc}]{\sphinxcrossref{\DUrole{doc}{processors}}}}
command lays out a 3d grid of processors in each of 2 partitions.


\subsection{Restrictions}
\label{\detokenize{partition:restrictions}}\begin{quote}

\sphinxAtStartPar
none
\end{quote}


\subsection{Related commands}
\label{\detokenize{partition:related-commands}}
\sphinxAtStartPar
{\hyperref[\detokenize{run_style::doc}]{\sphinxcrossref{\DUrole{doc}{run\_style verlet/split}}}}


\subsection{Default}
\label{\detokenize{partition:default}}
\sphinxAtStartPar
none

\sphinxstepscope

\index{plugin@\spxentry{plugin}}\ignorespaces 

\section{plugin command}
\label{\detokenize{plugin:plugin-command}}\label{\detokenize{plugin:index-0}}\label{\detokenize{plugin::doc}}

\subsection{Syntax}
\label{\detokenize{plugin:syntax}}
\begin{sphinxVerbatim}[commandchars=\\\{\}]
\PYG{k}{plugin}\PYG{+w}{ }\PYG{n}{command}\PYG{+w}{ }\PYG{n}{args}
\end{sphinxVerbatim}
\begin{itemize}
\item {} 
\sphinxAtStartPar
command = \sphinxstyleemphasis{load} or \sphinxstyleemphasis{unload} or \sphinxstyleemphasis{list} or \sphinxstyleemphasis{clear}

\item {} 
\sphinxAtStartPar
args = list of arguments for a particular plugin command
\begin{sphinxalltt}
\sphinxstyleemphasis{load} file = load plugin(s) from shared object in \sphinxstyleemphasis{file}
\sphinxstyleemphasis{unload} style name = unload plugin \sphinxstyleemphasis{name} of style \sphinxstyleemphasis{style}
    \sphinxstyleemphasis{style} = \sphinxstyleemphasis{pair} or \sphinxstyleemphasis{bond} or \sphinxstyleemphasis{angle} or \sphinxstyleemphasis{dihedral} or \sphinxstyleemphasis{improper} or \sphinxstyleemphasis{kspace} or \sphinxstyleemphasis{compute} or \sphinxstyleemphasis{fix} or \sphinxstyleemphasis{region} or \sphinxstyleemphasis{command}
\sphinxstyleemphasis{list} = print a list of currently loaded plugins
\sphinxstyleemphasis{clear} = unload all currently loaded plugins
\end{sphinxalltt}

\end{itemize}


\subsection{Examples}
\label{\detokenize{plugin:examples}}
\begin{sphinxVerbatim}[commandchars=\\\{\}]
\PYG{k}{plugin}\PYG{+w}{ }\PYG{n}{load}\PYG{+w}{ }\PYG{n}{morse2plugin.so}
\PYG{k}{plugin}\PYG{+w}{ }\PYG{n}{unload}\PYG{+w}{ }\PYG{n}{pair}\PYG{+w}{ }\PYG{n}{morse2}\PYG{o}{/}\PYG{n}{omp}
\PYG{k}{plugin}\PYG{+w}{ }\PYG{n}{unload}\PYG{+w}{ }\PYG{n}{command}\PYG{+w}{ }\PYG{n}{hello}
\PYG{k}{plugin}\PYG{+w}{ }\PYG{n}{list}
\PYG{k}{plugin}\PYG{+w}{ }\PYG{n}{clear}
\end{sphinxVerbatim}


\subsection{Description}
\label{\detokenize{plugin:description}}
\sphinxAtStartPar
The plugin command allows to load (and unload) additional styles and
commands into a LAMMPS binary from so\sphinxhyphen{}called dynamic shared object (DSO)
files.  This enables to add new functionality to an existing LAMMPS
binary without having to recompile and link the entire executable.

\sphinxAtStartPar
The \sphinxstyleemphasis{load} command will load and initialize all plugins contained in the
plugin DSO with the given filename.  A message with information the
plugin style and name and more will be printed.  Individual DSO files
may contain multiple plugins.  More details about how to write and
compile the plugin DSO is given in programmer’s guide part of the manual
under {\hyperref[\detokenize{Developer_plugins::doc}]{\sphinxcrossref{\DUrole{doc}{Writing plugins}}}}.

\sphinxAtStartPar
The \sphinxstyleemphasis{unload} command will remove the given style or the given name from
the list of available styles.  If the plugin style is currently in use,
that style instance will be deleted.

\sphinxAtStartPar
The \sphinxstyleemphasis{list} command will print a list of the loaded plugins and their
styles and names.

\sphinxAtStartPar
The \sphinxstyleemphasis{clear} command will unload all currently loaded plugins.

\begin{sphinxadmonition}{note}{Automatic loading of plugins}

\sphinxAtStartPar
\DUrole{versionmodified,added}{Added in version 4May2022.}

\sphinxAtStartPar
When the environment variable \sphinxcode{\sphinxupquote{LAMMPS\_PLUGIN\_PATH}} is set, then
LAMMPS will search the directory (or directories) listed in this path
for files with names that end in \sphinxcode{\sphinxupquote{plugin.so}}
(e.g. \sphinxcode{\sphinxupquote{helloplugin.so}}) and will try to load the contained plugins
automatically at start\sphinxhyphen{}up.
\end{sphinxadmonition}


\subsection{Restrictions}
\label{\detokenize{plugin:restrictions}}
\sphinxAtStartPar
The \sphinxstyleemphasis{plugin} command is part of the PLUGIN package.  It is
only enabled if LAMMPS was built with that package.  See
the {\hyperref[\detokenize{Build_package::doc}]{\sphinxcrossref{\DUrole{doc}{Build package}}}} page for more info.

\sphinxAtStartPar
If plugins access functions or classes from a package,
LAMMPS must have been compiled with that package included.

\sphinxAtStartPar
Plugins are dependent on the LAMMPS binary interface (ABI)
and particularly the MPI library used. So they are not guaranteed
to work when the plugin was compiled with a different MPI library
or different compilation settings or a different LAMMPS version.
There are no checks, so if there is a mismatch the plugin object
will either not load or data corruption and crashes may happen.


\subsection{Related commands}
\label{\detokenize{plugin:related-commands}}
\sphinxAtStartPar
none


\subsection{Default}
\label{\detokenize{plugin:default}}
\sphinxAtStartPar
none

\sphinxstepscope

\index{prd@\spxentry{prd}}\ignorespaces 

\section{prd command}
\label{\detokenize{prd:prd-command}}\label{\detokenize{prd:index-0}}\label{\detokenize{prd::doc}}

\subsection{Syntax}
\label{\detokenize{prd:syntax}}
\begin{sphinxVerbatim}[commandchars=\\\{\}]
\PYG{k}{prd}\PYG{+w}{ }\PYG{n}{N}\PYG{+w}{ }\PYG{n}{t\PYGZus{}event}\PYG{+w}{ }\PYG{n}{n\PYGZus{}dephase}\PYG{+w}{ }\PYG{n}{t\PYGZus{}dephase}\PYG{+w}{ }\PYG{n}{t\PYGZus{}correlate}\PYG{+w}{ }\PYG{n}{compute}\PYG{o}{\PYGZhy{}}\PYG{n}{ID}\PYG{+w}{ }\PYG{n}{seed}\PYG{+w}{ }\PYG{n}{keyword}\PYG{+w}{ }\PYG{n}{value}\PYG{+w}{ }\PYG{n}{...}
\end{sphinxVerbatim}
\begin{itemize}
\item {} 
\sphinxAtStartPar
N = \# of timesteps to run (not including dephasing/quenching)

\item {} 
\sphinxAtStartPar
t\_event = timestep interval between event checks

\item {} 
\sphinxAtStartPar
n\_dephase = number of velocity randomizations to perform in each dephase run

\item {} 
\sphinxAtStartPar
t\_dephase = number of timesteps to run dynamics after each velocity randomization during dephase

\item {} 
\sphinxAtStartPar
t\_correlate = number of timesteps within which 2 consecutive events are considered to be correlated

\item {} 
\sphinxAtStartPar
compute\sphinxhyphen{}ID = ID of the compute used for event detection

\item {} 
\sphinxAtStartPar
random\_seed = random \# seed (positive integer)

\item {} 
\sphinxAtStartPar
zero or more keyword/value pairs may be appended

\item {} 
\sphinxAtStartPar
keyword = \sphinxstyleemphasis{min} or \sphinxstyleemphasis{temp} or \sphinxstyleemphasis{vel} or \sphinxstyleemphasis{time}
\begin{sphinxalltt}
\sphinxstyleemphasis{min} values = etol ftol maxiter maxeval
  etol = stopping tolerance for energy, used in quenching
  ftol = stopping tolerance for force, used in quenching
  maxiter = max iterations of minimize, used in quenching
  maxeval = max number of force/energy evaluations, used in quenching
\sphinxstyleemphasis{temp} value = Tdephase
  Tdephase = target temperature for velocity randomization, used in dephasing
\sphinxstyleemphasis{vel} values = loop dist
  loop = \sphinxstyleemphasis{all} or \sphinxstyleemphasis{local} or \sphinxstyleemphasis{geom}, used in dephasing
  dist = \sphinxstyleemphasis{uniform} or \sphinxstyleemphasis{gaussian}, used in dephasing
\sphinxstyleemphasis{time} value = \sphinxstyleemphasis{steps} or \sphinxstyleemphasis{clock}
  \sphinxstyleemphasis{steps} = simulation runs for N timesteps on each replica (default)
  \sphinxstyleemphasis{clock} = simulation runs for N timesteps across all replicas
\end{sphinxalltt}

\end{itemize}


\subsection{Examples}
\label{\detokenize{prd:examples}}
\begin{sphinxVerbatim}[commandchars=\\\{\}]
\PYG{k}{prd}\PYG{+w}{ }\PYG{l+m}{5000}\PYG{+w}{ }\PYG{l+m}{100}\PYG{+w}{ }\PYG{l+m}{10}\PYG{+w}{ }\PYG{l+m}{10}\PYG{+w}{ }\PYG{l+m}{100}\PYG{+w}{ }\PYG{l+m}{1}\PYG{+w}{ }\PYG{l+m}{54982}
\PYG{k}{prd}\PYG{+w}{ }\PYG{l+m}{5000}\PYG{+w}{ }\PYG{l+m}{100}\PYG{+w}{ }\PYG{l+m}{10}\PYG{+w}{ }\PYG{l+m}{10}\PYG{+w}{ }\PYG{l+m}{100}\PYG{+w}{ }\PYG{l+m}{1}\PYG{+w}{ }\PYG{l+m}{54982}\PYG{+w}{ }\PYG{n}{min}\PYG{+w}{ }\PYG{l+m}{0.1}\PYG{+w}{ }\PYG{l+m}{0.1}\PYG{+w}{ }\PYG{l+m}{100}\PYG{+w}{ }\PYG{l+m}{200}
\end{sphinxVerbatim}


\subsection{Description}
\label{\detokenize{prd:description}}
\sphinxAtStartPar
Run a parallel replica dynamics (PRD) simulation using multiple
replicas of a system.  One or more replicas can be used.  The total
number of steps \sphinxstyleemphasis{N} to run can be interpreted in one of two ways; see
discussion of the \sphinxstyleemphasis{time} keyword below.

\sphinxAtStartPar
PRD is described in {\hyperref[\detokenize{prd:voter1998}]{\sphinxcrossref{\DUrole{std,std-ref}{(Voter1998)}}}} by Art Voter.
Similar to global or local hyperdynamics (HD), PRD is a method for
performing accelerated dynamics that is suitable for infrequent\sphinxhyphen{}event
systems that obey first\sphinxhyphen{}order kinetics.  A good overview of
accelerated dynamics methods (AMD) for such systems in given in this
review paper {\hyperref[\detokenize{prd:voter2002prd}]{\sphinxcrossref{\DUrole{std,std-ref}{(Voter2002)}}}} from Art’s group.  To
quote from the paper: “The dynamical evolution is characterized by
vibrational excursions within a potential basin, punctuated by
occasional transitions between basins.  The transition probability is
characterized by p(t) = k*exp(\sphinxhyphen{}kt) where k is the rate constant.”

\sphinxAtStartPar
Both PRD and HD produce a time\sphinxhyphen{}accurate trajectory that effectively
extends the timescale over which a system can be simulated, but they
do it differently.  PRD creates Nr replicas of the system and runs
dynamics on each independently with a normal unbiased potential until
an event occurs in one of the replicas.  The time between events is
reduced by a factor of Nr replicas.  HD uses a single replica of the
system and accelerates time by biasing the interaction potential in a
manner such that each timestep is effectively longer.  For both
methods, per CPU second, more physical time elapses and more events
occur.  See the {\hyperref[\detokenize{hyper::doc}]{\sphinxcrossref{\DUrole{doc}{hyper}}}} page for more info about HD.

\sphinxAtStartPar
In PRD, each replica runs on a partition of one or more processors.
Processor partitions are defined at run\sphinxhyphen{}time using the {\hyperref[\detokenize{Run_options::doc}]{\sphinxcrossref{\DUrole{doc}{\sphinxhyphen{}partition command\sphinxhyphen{}line switch}}}}.  Note that if you have MPI
installed, you can run a multi\sphinxhyphen{}replica simulation with more replicas
(partitions) than you have physical processors, e.g you can run a
10\sphinxhyphen{}replica simulation on one or two processors.  However for PRD, this
makes little sense, since running a replica on virtual instead of
physical processors,offers no effective parallel speed\sphinxhyphen{}up in searching
for infrequent events.  See the {\hyperref[\detokenize{Howto_replica::doc}]{\sphinxcrossref{\DUrole{doc}{Howto replica}}}} doc
page for further discussion.

\sphinxAtStartPar
When a PRD simulation is performed, it is assumed that each replica is
running the same model, though LAMMPS does not check for this.
I.e. the simulation domain, the number of atoms, the interaction
potentials, etc should be the same for every replica.

\sphinxAtStartPar
A PRD run has several stages, which are repeated each time an “event”
occurs in one of the replicas, as explained below.  The logic for a
PRD run is as follows:
\begin{sphinxalltt}
while (time remains):
  dephase for n\_dephase*t\_dephase steps
  until (event occurs on some replica):
    run dynamics for t\_event steps
    quench
    check for uncorrelated event on any replica
  until (no correlated event occurs):
    run dynamics for t\_correlate steps
    quench
    check for correlated event on this replica
  event replica shares state with all replicas
\end{sphinxalltt}

\sphinxAtStartPar
Before this loop begins, the state of the system on replica 0 is
shared with all replicas, so that all replicas begin from the same
initial state. The first potential energy basin is identified by
quenching (an energy minimization, see below) the initial state and
storing the resulting coordinates for reference.

\sphinxAtStartPar
In the first stage, dephasing is performed by each replica
independently to eliminate correlations between replicas.  This is
done by choosing a random set of velocities, based on the
\sphinxstyleemphasis{random\_seed} that is specified, and running \sphinxstyleemphasis{t\_dephase} timesteps of
dynamics.  This is repeated \sphinxstyleemphasis{n\_dephase} times.  At each of the
\sphinxstyleemphasis{n\_dephase} stages, if an event occurs during the \sphinxstyleemphasis{t\_dephase} steps of
dynamics for a particular replica, the replica repeats the stage until
no event occurs.

\sphinxAtStartPar
If the \sphinxstyleemphasis{temp} keyword is not specified, the target temperature for
velocity randomization for each replica is the current temperature of
that replica.  Otherwise, it is the specified \sphinxstyleemphasis{Tdephase} temperature.
The style of velocity randomization is controlled using the keyword
\sphinxstyleemphasis{vel} with arguments that have the same meaning as their counterparts
in the {\hyperref[\detokenize{velocity::doc}]{\sphinxcrossref{\DUrole{doc}{velocity}}}} command.

\sphinxAtStartPar
In the second stage, each replica runs dynamics continuously, stopping
every \sphinxstyleemphasis{t\_event} steps to check if a transition event has occurred.
This check is performed by quenching the system and comparing the
resulting atom coordinates to the coordinates from the previous basin.
The first time through the PRD loop, the “previous basin” is the set
of quenched coordinates from the initial state of the system.

\sphinxAtStartPar
A quench is an energy minimization and is performed by whichever
algorithm has been defined by the {\hyperref[\detokenize{min_style::doc}]{\sphinxcrossref{\DUrole{doc}{min\_style}}}} command.
Minimization parameters may be set via the
{\hyperref[\detokenize{min_modify::doc}]{\sphinxcrossref{\DUrole{doc}{min\_modify}}}} command and by the \sphinxstyleemphasis{min} keyword of the
PRD command.  The latter are the settings that would be used with the
{\hyperref[\detokenize{minimize::doc}]{\sphinxcrossref{\DUrole{doc}{minimize}}}} command.  Note that typically, you do not
need to perform a highly\sphinxhyphen{}converged minimization to detect a transition
event, though you may need to in order to prevent a set of atoms in
the system from relaxing to a saddle point.

\sphinxAtStartPar
The event check is performed by a compute with the specified
\sphinxstyleemphasis{compute\sphinxhyphen{}ID}.  Currently there is only one compute that works with the
PRD command, which is the {\hyperref[\detokenize{compute_event_displace::doc}]{\sphinxcrossref{\DUrole{doc}{compute event/displace}}}} command.  Other
event\sphinxhyphen{}checking computes may be added.  {\hyperref[\detokenize{compute_event_displace::doc}]{\sphinxcrossref{\DUrole{doc}{Compute event/displace}}}} checks whether any atom in
the compute group has moved further than a specified threshold
distance.  If so, an “event” has occurred.

\sphinxAtStartPar
In the third stage, the replica on which the event occurred (event
replica) continues to run dynamics to search for correlated events.
This is done by running dynamics for \sphinxstyleemphasis{t\_correlate} steps, quenching
every \sphinxstyleemphasis{t\_event} steps, and checking if another event has occurred.

\sphinxAtStartPar
The first time no correlated event occurs, the final state of the
event replica is shared with all replicas, the new basin reference
coordinates are updated with the quenched state, and the outer loop
begins again. While the replica event is searching for correlated
events, all the other replicas also run dynamics and event checking
with the same schedule, but the final states are always overwritten by
the state of the event replica.

\sphinxAtStartPar
The outer loop of the pseudocode above continues until \sphinxstyleemphasis{N} steps of
dynamics have been performed.  Note that \sphinxstyleemphasis{N} only includes the
dynamics of stages 2 and 3, not the steps taken during dephasing or
the minimization iterations of quenching.  The specified \sphinxstyleemphasis{N} is
interpreted in one of two ways, depending on the \sphinxstyleemphasis{time} keyword.  If
the \sphinxstyleemphasis{time} value is \sphinxstyleemphasis{steps}, which is the default, then each replica
runs for \sphinxstyleemphasis{N} timesteps.  If the \sphinxstyleemphasis{time} value is \sphinxstyleemphasis{clock}, then the
simulation runs until \sphinxstyleemphasis{N} aggregate timesteps across all replicas have
elapsed.  This aggregate time is the “clock” time defined below, which
typically advances nearly M times faster than the timestepping on a
single replica, where M is the number of replicas.


\bigskip\hrule\bigskip


\sphinxAtStartPar
Four kinds of output can be generated during a PRD run: event
statistics, thermodynamic output by each replica, dump files, and
restart files.

\sphinxAtStartPar
When running with multiple partitions (each of which is a replica in
this case), the print\sphinxhyphen{}out to the screen and master log.lammps file is
limited to event statistics.  Note that if a PRD run is performed on
only a single replica then the event statistics will be intermixed
with the usual thermodynamic output discussed below.

\sphinxAtStartPar
The quantities printed each time an event occurs are the timestep, CPU
time, clock, event number, a correlation flag, the number of
coincident events, and the replica number of the chosen event.

\sphinxAtStartPar
The timestep is the usual LAMMPS timestep, except that time does not
advance during dephasing or quenches, but only during dynamics.  Note
that are two kinds of dynamics in the PRD loop listed above that
contribute to this timestepping.  The first is when all replicas are
performing independent dynamics, waiting for an event to occur.  The
second is when correlated events are being searched for, but only one
replica is running dynamics.

\sphinxAtStartPar
The CPU time is the total elapsed time on each processor, since the
start of the PRD run.

\sphinxAtStartPar
The clock is the same as the timestep except that it advances by M
steps per timestep during the first kind of dynamics when the M
replicas are running independently.  The clock advances by only 1 step
per timestep during the second kind of dynamics, when only a single
replica is checking for a correlated event.  Thus “clock” time
represents the aggregate time (in steps) that has effectively elapsed
during a PRD simulation on M replicas.  If most of the PRD run is
spent in the second stage of the loop above, searching for infrequent
events, then the clock will advance nearly M times faster than it
would if a single replica was running.  Note the clock time between
successive events should be drawn from p(t).

\sphinxAtStartPar
The event number is a counter that increments with each event, whether
it is uncorrelated or correlated.

\sphinxAtStartPar
The correlation flag will be 0 when an uncorrelated event occurs
during the second stage of the loop listed above, i.e. when all
replicas are running independently.  The correlation flag will be 1
when a correlated event occurs during the third stage of the loop
listed above, i.e. when only one replica is running dynamics.

\sphinxAtStartPar
When more than one replica detects an event at the end of the same
event check (every \sphinxstyleemphasis{t\_event} steps) during the second stage, then
one of them is chosen at random.  The number of coincident events is
the number of replicas that detected an event.  Normally, this value
should be 1.  If it is often greater than 1, then either the number of
replicas is too large, or \sphinxstyleemphasis{t\_event} is too large.

\sphinxAtStartPar
The replica number is the ID of the replica (from 0 to M\sphinxhyphen{}1) in which
the event occurred.


\bigskip\hrule\bigskip


\sphinxAtStartPar
When running on multiple partitions, LAMMPS produces additional log
files for each partition, e.g. log.lammps.0, log.lammps.1, etc.  For
the PRD command, these contain the thermodynamic output for each
replica.  You will see short runs and minimizations corresponding to
the dynamics and quench operations of the loop listed above.  The
timestep will be reset appropriately depending on whether the
operation advances time or not.

\sphinxAtStartPar
After the PRD command completes, timing statistics for the PRD run are
printed in each replica’s log file, giving a breakdown of how much CPU
time was spent in each stage (dephasing, dynamics, quenching, etc).


\bigskip\hrule\bigskip


\sphinxAtStartPar
Any {\hyperref[\detokenize{dump::doc}]{\sphinxcrossref{\DUrole{doc}{dump files}}}} defined in the input script, will be
written to during a PRD run at timesteps corresponding to both
uncorrelated and correlated events.  This means the requested dump
frequency in the {\hyperref[\detokenize{dump::doc}]{\sphinxcrossref{\DUrole{doc}{dump}}}} command is ignored.  There will be
one dump file (per dump command) created for all partitions.

\sphinxAtStartPar
The atom coordinates of the dump snapshot are those of the minimum
energy configuration resulting from quenching following a transition
event.  The timesteps written into the dump files correspond to the
timestep at which the event occurred and NOT the clock.  A dump
snapshot corresponding to the initial minimum state used for event
detection is written to the dump file at the beginning of each PRD
run.


\bigskip\hrule\bigskip


\sphinxAtStartPar
If the {\hyperref[\detokenize{restart::doc}]{\sphinxcrossref{\DUrole{doc}{restart}}}} command is used, a single restart file
for all the partitions is generated, which allows a PRD run to be
continued by a new input script in the usual manner.

\sphinxAtStartPar
The restart file is generated at the end of the loop listed above.  If
no correlated events are found, this means it contains a snapshot of
the system at time T + \sphinxstyleemphasis{t\_correlate}, where T is the time at which the
uncorrelated event occurred.  If correlated events were found, then it
contains a snapshot of the system at time T + \sphinxstyleemphasis{t\_correlate}, where T
is the time of the last correlated event.

\sphinxAtStartPar
The restart frequency specified in the {\hyperref[\detokenize{restart::doc}]{\sphinxcrossref{\DUrole{doc}{restart}}}} command
is interpreted differently when performing a PRD run.  It does not
mean the timestep interval between restart files.  Instead it means an
event interval for uncorrelated events.  Thus a frequency of 1 means
write a restart file every time an uncorrelated event occurs.  A
frequency of 10 means write a restart file every 10th uncorrelated
event.

\sphinxAtStartPar
When an input script reads a restart file from a previous PRD run, the
new script can be run on a different number of replicas or processors.
However, it is assumed that \sphinxstyleemphasis{t\_correlate} in the new PRD command is
the same as it was previously.  If not, the calculation of the “clock”
value for the first event in the new run will be slightly off.


\bigskip\hrule\bigskip



\subsection{Restrictions}
\label{\detokenize{prd:restrictions}}
\sphinxAtStartPar
This command can only be used if LAMMPS was built with the REPLICA
package.  See the {\hyperref[\detokenize{Build_package::doc}]{\sphinxcrossref{\DUrole{doc}{Build package}}}} doc
page for more info.

\sphinxAtStartPar
The \sphinxstyleemphasis{N} and \sphinxstyleemphasis{t\_correlate} settings must be integer multiples of
\sphinxstyleemphasis{t\_event}.

\sphinxAtStartPar
Runs restarted from restart file written during a PRD run will not
produce identical results due to changes in the random numbers used
for dephasing.

\sphinxAtStartPar
This command cannot be used when any fixes are defined that keep track
of elapsed time to perform time\sphinxhyphen{}dependent operations.  Examples
include the “ave” fixes such as {\hyperref[\detokenize{fix_ave_chunk::doc}]{\sphinxcrossref{\DUrole{doc}{fix ave/chunk}}}}.
Also {\hyperref[\detokenize{fix_dt_reset::doc}]{\sphinxcrossref{\DUrole{doc}{fix dt/reset}}}} and {\hyperref[\detokenize{fix_deposit::doc}]{\sphinxcrossref{\DUrole{doc}{fix deposit}}}}.


\subsection{Related commands}
\label{\detokenize{prd:related-commands}}
\sphinxAtStartPar
{\hyperref[\detokenize{compute_event_displace::doc}]{\sphinxcrossref{\DUrole{doc}{compute event/displace}}}},
{\hyperref[\detokenize{min_modify::doc}]{\sphinxcrossref{\DUrole{doc}{min\_modify}}}}, {\hyperref[\detokenize{min_style::doc}]{\sphinxcrossref{\DUrole{doc}{min\_style}}}},
{\hyperref[\detokenize{run_style::doc}]{\sphinxcrossref{\DUrole{doc}{run\_style}}}}, {\hyperref[\detokenize{minimize::doc}]{\sphinxcrossref{\DUrole{doc}{minimize}}}},
{\hyperref[\detokenize{velocity::doc}]{\sphinxcrossref{\DUrole{doc}{velocity}}}}, {\hyperref[\detokenize{temper::doc}]{\sphinxcrossref{\DUrole{doc}{temper}}}}, {\hyperref[\detokenize{neb::doc}]{\sphinxcrossref{\DUrole{doc}{neb}}}},
{\hyperref[\detokenize{tad::doc}]{\sphinxcrossref{\DUrole{doc}{tad}}}}, {\hyperref[\detokenize{hyper::doc}]{\sphinxcrossref{\DUrole{doc}{hyper}}}}


\subsection{Default}
\label{\detokenize{prd:default}}
\sphinxAtStartPar
The option defaults are min = 0.1 0.1 40 50, no temp setting, vel =
geom gaussian, and time = steps.


\bigskip\hrule\bigskip

\phantomsection\label{\detokenize{prd:voter1998}}
\sphinxAtStartPar
\sphinxstylestrong{(Voter1998)} Voter, Phys Rev B, 57, 13985 (1998).

\phantomsection\label{\detokenize{prd:voter2002prd}}
\sphinxAtStartPar
\sphinxstylestrong{(Voter2002)} Voter, Montalenti, Germann, Annual Review of Materials
Research 32, 321 (2002).

\sphinxstepscope

\index{print@\spxentry{print}}\ignorespaces 

\section{print command}
\label{\detokenize{print:print-command}}\label{\detokenize{print:index-0}}\label{\detokenize{print::doc}}

\subsection{Syntax}
\label{\detokenize{print:syntax}}
\begin{sphinxVerbatim}[commandchars=\\\{\}]
\PYG{k}{print}\PYG{+w}{ }\PYG{n}{string}\PYG{+w}{ }\PYG{n}{keyword}\PYG{+w}{ }\PYG{n}{value}
\end{sphinxVerbatim}
\begin{itemize}
\item {} 
\sphinxAtStartPar
string = text string to print, which may contain variables

\item {} 
\sphinxAtStartPar
zero or more keyword/value pairs may be appended

\item {} 
\sphinxAtStartPar
keyword = \sphinxstyleemphasis{file} or \sphinxstyleemphasis{append} or \sphinxstyleemphasis{screen} or \sphinxstyleemphasis{universe}
\begin{sphinxalltt}
\sphinxstyleemphasis{file} value = filename
\sphinxstyleemphasis{append} value = filename
\sphinxstyleemphasis{screen} value = \sphinxstyleemphasis{yes} or \sphinxstyleemphasis{no}
\sphinxstyleemphasis{universe} value = \sphinxstyleemphasis{yes} or \sphinxstyleemphasis{no}
\end{sphinxalltt}

\end{itemize}


\subsection{Examples}
\label{\detokenize{print:examples}}
\begin{sphinxVerbatim}[commandchars=\\\{\}]
\PYG{k}{print}\PYG{+w}{ }\PYG{l+s}{\PYGZdq{}}\PYG{l+s}{Done with equilibration}\PYG{l+s}{\PYGZdq{}}\PYG{+w}{ }\PYG{n}{file}\PYG{+w}{ }\PYG{n}{info.dat}
\PYG{k}{print}\PYG{+w}{ }\PYG{n}{Vol}\PYG{o}{=}\PYG{n+nv}{\PYGZdl{}v}\PYG{+w}{ }\PYG{n}{append}\PYG{+w}{ }\PYG{n}{info.dat}\PYG{+w}{ }\PYG{n}{screen}\PYG{+w}{ }\PYG{n}{no}
\PYG{k}{print}\PYG{+w}{ }\PYG{l+s}{\PYGZdq{}}\PYG{l+s}{The system volume is now \PYGZdl{}v}\PYG{l+s}{\PYGZdq{}}
\PYG{k}{print}\PYG{+w}{ }\PYG{l+s}{\PYGZsq{}}\PYG{l+s}{The system volume is now \PYGZdl{}v}\PYG{l+s}{\PYGZsq{}}
\PYG{k}{print}\PYG{+w}{ }\PYG{l+s}{\PYGZdq{}}\PYG{l+s}{NEB calculation 1 complete}\PYG{l+s}{\PYGZdq{}}\PYG{+w}{ }\PYG{n}{screen}\PYG{+w}{ }\PYG{n}{no}\PYG{+w}{ }\PYG{n}{universe}\PYG{+w}{ }\PYG{n}{yes}
\PYG{k}{print}\PYG{+w}{ }\PYG{l+s}{\PYGZdq{}}\PYG{l+s}{\PYGZdq{}}\PYG{l+s}{\PYGZdq{}}
\PYG{l+s}{System volume = \PYGZdl{}v}
\PYG{l+s}{System temperature = \PYGZdl{}t}
\PYG{l+s}{\PYGZdq{}}\PYG{l+s}{\PYGZdq{}}\PYG{l+s}{\PYGZdq{}}
\end{sphinxVerbatim}


\subsection{Description}
\label{\detokenize{print:description}}
\sphinxAtStartPar
Print a text string to the screen and logfile.  The text string must
be a single argument, so if it is one line but more than one word, it
should be enclosed in single or double quotes.  To generate multiple
lines of output, the string can be enclosed in triple quotes, as in
the last example above.  If the text string contains variables, they
will be evaluated and their current values printed.

\sphinxAtStartPar
\DUrole{versionmodified,added}{Added in version 15Jun2023: }support for vector style variables

\sphinxAtStartPar
See the {\hyperref[\detokenize{variable::doc}]{\sphinxcrossref{\DUrole{doc}{variable}}}} command for a description of
\sphinxstyleemphasis{equal} and \sphinxstyleemphasis{vector} style variables which are typically the most
useful ones to use with the print command.  Equal\sphinxhyphen{} and vector\sphinxhyphen{}style
variables can calculate formulas involving mathematical operations,
atom properties, group properties, thermodynamic properties, global
values calculated by a {\hyperref[\detokenize{compute::doc}]{\sphinxcrossref{\DUrole{doc}{compute}}}} or {\hyperref[\detokenize{fix::doc}]{\sphinxcrossref{\DUrole{doc}{fix}}}},
or references to other {\hyperref[\detokenize{variable::doc}]{\sphinxcrossref{\DUrole{doc}{variables}}}}.  Vector\sphinxhyphen{}style
variables are printed in a bracketed, comma\sphinxhyphen{}separated format,
e.g. {[}1,2,3,4{]} or {[}12.5,2,4.6,10.1{]}.

\begin{sphinxadmonition}{note}{Note:}
\sphinxAtStartPar
As discussed on the {\hyperref[\detokenize{Commands_parse::doc}]{\sphinxcrossref{\DUrole{doc}{Commands parse}}}} doc
page, the text string can use “immediate” variables, specified as
\$(formula) with parenthesis, where the numeric formula has the same
syntax as equal\sphinxhyphen{}style variables described on the {\hyperref[\detokenize{variable::doc}]{\sphinxcrossref{\DUrole{doc}{variable}}}} doc page.  This is a convenient way to evaluate a
formula immediately without using the variable command to define a
named variable and then use that variable in the text string.  The
formula can include a trailing colon and format string which
determines the precision with which the numeric value is output.
This is also explained on the {\hyperref[\detokenize{Commands_parse::doc}]{\sphinxcrossref{\DUrole{doc}{Commands parse}}}} doc page.
\end{sphinxadmonition}

\sphinxAtStartPar
If you want the print command to be executed multiple times (with
changing variable values), there are 3 options.  First, consider using
the {\hyperref[\detokenize{fix_print::doc}]{\sphinxcrossref{\DUrole{doc}{fix print}}}} command, which will print a string
periodically during a simulation.  Second, the print command can be
used as an argument to the \sphinxstyleemphasis{every} option of the {\hyperref[\detokenize{run::doc}]{\sphinxcrossref{\DUrole{doc}{run}}}}
command.  Third, the print command could appear in a section of the
input script that is looped over (see the {\hyperref[\detokenize{jump::doc}]{\sphinxcrossref{\DUrole{doc}{jump}}}} and
{\hyperref[\detokenize{next::doc}]{\sphinxcrossref{\DUrole{doc}{next}}}} commands).

\sphinxAtStartPar
If the \sphinxstyleemphasis{file} or \sphinxstyleemphasis{append} keyword is used, a filename is specified to
which the output will be written.  If \sphinxstyleemphasis{file} is used, then the
filename is overwritten if it already exists.  If \sphinxstyleemphasis{append} is used,
then the filename is appended to if it already exists, or created if
it does not exist.

\sphinxAtStartPar
If the \sphinxstyleemphasis{screen} keyword is used, output to the screen and logfile can
be turned on or off as desired.

\sphinxAtStartPar
If the \sphinxstyleemphasis{universe} keyword is used, output to the global screen and
logfile can be turned on or off as desired. In multi\sphinxhyphen{}partition
calculations, the \sphinxstyleemphasis{screen} option and the corresponding output only
apply to the screen and logfile of the individual partition.


\subsection{Restrictions}
\label{\detokenize{print:restrictions}}\begin{quote}

\sphinxAtStartPar
none
\end{quote}


\subsection{Related commands}
\label{\detokenize{print:related-commands}}
\sphinxAtStartPar
{\hyperref[\detokenize{fix_print::doc}]{\sphinxcrossref{\DUrole{doc}{fix print}}}}, {\hyperref[\detokenize{variable::doc}]{\sphinxcrossref{\DUrole{doc}{variable}}}}


\subsection{Default}
\label{\detokenize{print:default}}
\sphinxAtStartPar
The option defaults are no file output, screen = yes, and universe = no.

\sphinxstepscope

\index{processors@\spxentry{processors}}\ignorespaces 

\section{processors command}
\label{\detokenize{processors:processors-command}}\label{\detokenize{processors:index-0}}\label{\detokenize{processors::doc}}

\subsection{Syntax}
\label{\detokenize{processors:syntax}}
\begin{sphinxVerbatim}[commandchars=\\\{\}]
\PYG{k}{processors}\PYG{+w}{ }\PYG{n}{Px}\PYG{+w}{ }\PYG{n}{Py}\PYG{+w}{ }\PYG{n}{Pz}\PYG{+w}{ }\PYG{n}{keyword}\PYG{+w}{ }\PYG{n}{args}\PYG{+w}{ }\PYG{n}{...}
\end{sphinxVerbatim}
\begin{itemize}
\item {} 
\sphinxAtStartPar
Px,Py,Pz = \# of processors in each dimension of 3d grid overlaying the simulation domain

\item {} 
\sphinxAtStartPar
zero or more keyword/arg pairs may be appended

\item {} 
\sphinxAtStartPar
keyword = \sphinxstyleemphasis{grid} or \sphinxstyleemphasis{map} or \sphinxstyleemphasis{part} or \sphinxstyleemphasis{file}
\begin{sphinxalltt}
\sphinxstyleemphasis{grid} arg = gstyle params ...
  gstyle = \sphinxstyleemphasis{onelevel} or \sphinxstyleemphasis{twolevel} or \sphinxstyleemphasis{numa} or \sphinxstyleemphasis{custom}
    \sphinxstyleemphasis{onelevel} params = none
    \sphinxstyleemphasis{twolevel} params = Nc Cx Cy Cz
      Nc = number of cores per node
      Cx,Cy,Cz = \# of cores in each dimension of 3d sub\sphinxhyphen{}grid assigned to each node
    \sphinxstyleemphasis{numa} params = none
    \sphinxstyleemphasis{custom} params = infile
      infile = file containing grid layout
\sphinxstyleemphasis{numa\_nodes} arg = Nn
      Nn = number of numa domains per node
\sphinxstyleemphasis{map} arg = \sphinxstyleemphasis{cart} or \sphinxstyleemphasis{cart/reorder} or \sphinxstyleemphasis{xyz} or \sphinxstyleemphasis{xzy} or \sphinxstyleemphasis{yxz} or \sphinxstyleemphasis{yzx} or \sphinxstyleemphasis{zxy} or \sphinxstyleemphasis{zyx}
   cart = use MPI\_Cart() methods to map processors to 3d grid with reorder = 0
   cart/reorder = use MPI\_Cart() methods to map processors to 3d grid with reorder = 1
   xyz,xzy,yxz,yzx,zxy,zyx = map processors to 3d grid in IJK ordering
\sphinxstyleemphasis{part} args = Psend Precv cstyle
  Psend = partition \# (1 to Np) which will send its processor layout
  Precv = partition \# (1 to Np) which will recv the processor layout
  cstyle = \sphinxstyleemphasis{multiple}
    \sphinxstyleemphasis{multiple} = Psend grid will be multiple of Precv grid in each dimension
\sphinxstyleemphasis{file} arg = outfile
  outfile = name of file to write 3d grid of processors to
\end{sphinxalltt}

\end{itemize}


\subsection{Examples}
\label{\detokenize{processors:examples}}
\begin{sphinxVerbatim}[commandchars=\\\{\}]
\PYG{k}{processors}\PYG{+w}{ }\PYG{o}{*}\PYG{+w}{ }\PYG{o}{*}\PYG{+w}{ }\PYG{l+m}{5}
\PYG{k}{processors}\PYG{+w}{ }\PYG{l+m}{2}\PYG{+w}{ }\PYG{l+m}{4}\PYG{+w}{ }\PYG{l+m}{4}
\PYG{k}{processors}\PYG{+w}{ }\PYG{o}{*}\PYG{+w}{ }\PYG{o}{*}\PYG{+w}{ }\PYG{l+m}{8}\PYG{+w}{ }\PYG{n}{map}\PYG{+w}{ }\PYG{n}{xyz}
\PYG{k}{processors}\PYG{+w}{ }\PYG{o}{*}\PYG{+w}{ }\PYG{o}{*}\PYG{+w}{ }\PYG{o}{*}\PYG{+w}{ }\PYG{n}{grid}\PYG{+w}{ }\PYG{n}{numa}
\PYG{k}{processors}\PYG{+w}{ }\PYG{o}{*}\PYG{+w}{ }\PYG{o}{*}\PYG{+w}{ }\PYG{o}{*}\PYG{+w}{ }\PYG{n}{grid}\PYG{+w}{ }\PYG{n}{twolevel}\PYG{+w}{ }\PYG{l+m}{4}\PYG{+w}{ }\PYG{o}{*}\PYG{+w}{ }\PYG{o}{*}\PYG{+w}{ }\PYG{l+m}{1}
\PYG{k}{processors}\PYG{+w}{ }\PYG{l+m}{4}\PYG{+w}{ }\PYG{l+m}{8}\PYG{+w}{ }\PYG{l+m}{16}\PYG{+w}{ }\PYG{n}{grid}\PYG{+w}{ }\PYG{n}{custom}\PYG{+w}{ }\PYG{n}{myfile}
\PYG{k}{processors}\PYG{+w}{ }\PYG{o}{*}\PYG{+w}{ }\PYG{o}{*}\PYG{+w}{ }\PYG{o}{*}\PYG{+w}{ }\PYG{n}{part}\PYG{+w}{ }\PYG{l+m}{1}\PYG{+w}{ }\PYG{l+m}{2}\PYG{+w}{ }\PYG{n}{multiple}
\end{sphinxVerbatim}


\subsection{Description}
\label{\detokenize{processors:description}}
\sphinxAtStartPar
Specify how processors are mapped as a regular 3d grid to the global
simulation box.  The mapping involves 2 steps.  First if there are P
processors it means choosing a factorization P = Px by Py by Pz so
that there are Px processors in the x dimension, and similarly for the
y and z dimensions.  Second, the P processors are mapped to the
regular 3d grid.  The arguments to this command control each of these
2 steps.

\sphinxAtStartPar
The Px, Py, Pz parameters affect the factorization.  Any of the 3
parameters can be specified with an asterisk “*”, which means LAMMPS
will choose the number of processors in that dimension of the grid.
It will do this based on the size and shape of the global simulation
box so as to minimize the surface\sphinxhyphen{}to\sphinxhyphen{}volume ratio of each processor’s
subdomain.

\sphinxAtStartPar
Choosing explicit values for Px or Py or Pz can be used to override
the default manner in which LAMMPS will create the regular 3d grid of
processors, if it is known to be sub\sphinxhyphen{}optimal for a particular problem.
E.g. a problem where the extent of atoms will change dramatically in a
particular dimension over the course of the simulation.

\sphinxAtStartPar
The product of Px, Py, Pz must equal P, the total \# of processors
LAMMPS is running on.  For a {\hyperref[\detokenize{dimension::doc}]{\sphinxcrossref{\DUrole{doc}{2d simulation}}}}, Pz must
equal 1.

\sphinxAtStartPar
Note that if you run on a prime number of processors P, then a grid
such as 1 x P x 1 will be required, which may incur extra
communication costs due to the high surface area of each processor’s
subdomain.

\sphinxAtStartPar
Also note that if multiple partitions are being used then P is the
number of processors in this partition; see the {\hyperref[\detokenize{Run_options::doc}]{\sphinxcrossref{\DUrole{doc}{\sphinxhyphen{}partition command\sphinxhyphen{}line switch}}}} page for details.  Also note
that you can prefix the processors command with the
{\hyperref[\detokenize{partition::doc}]{\sphinxcrossref{\DUrole{doc}{partition}}}} command to easily specify different
Px,Py,Pz values for different partitions.

\sphinxAtStartPar
You can use the {\hyperref[\detokenize{partition::doc}]{\sphinxcrossref{\DUrole{doc}{partition}}}} command to specify
different processor grids for different partitions, e.g.

\begin{sphinxVerbatim}[commandchars=\\\{\}]
\PYG{k}{partition}\PYG{+w}{ }\PYG{n}{yes}\PYG{+w}{ }\PYG{l+m}{1}\PYG{+w}{ }\PYG{n}{processors}\PYG{+w}{ }\PYG{l+m}{4}\PYG{+w}{ }\PYG{l+m}{4}\PYG{+w}{ }\PYG{l+m}{4}
\PYG{k}{partition}\PYG{+w}{ }\PYG{n}{yes}\PYG{+w}{ }\PYG{l+m}{2}\PYG{+w}{ }\PYG{n}{processors}\PYG{+w}{ }\PYG{l+m}{2}\PYG{+w}{ }\PYG{l+m}{3}\PYG{+w}{ }\PYG{l+m}{2}
\end{sphinxVerbatim}

\begin{sphinxadmonition}{note}{Note:}
\sphinxAtStartPar
This command only affects the initial regular 3d grid created
when the simulation box is first specified via a
{\hyperref[\detokenize{create_box::doc}]{\sphinxcrossref{\DUrole{doc}{create\_box}}}} or {\hyperref[\detokenize{read_data::doc}]{\sphinxcrossref{\DUrole{doc}{read\_data}}}} or
{\hyperref[\detokenize{read_restart::doc}]{\sphinxcrossref{\DUrole{doc}{read\_restart}}}} command.  Or if the simulation box is
re\sphinxhyphen{}created via the {\hyperref[\detokenize{replicate::doc}]{\sphinxcrossref{\DUrole{doc}{replicate}}}} command.  The same
regular grid is initially created, regardless of which
{\hyperref[\detokenize{comm_style::doc}]{\sphinxcrossref{\DUrole{doc}{comm\_style}}}} command is in effect.
\end{sphinxadmonition}

\sphinxAtStartPar
If load\sphinxhyphen{}balancing is never invoked via the {\hyperref[\detokenize{balance::doc}]{\sphinxcrossref{\DUrole{doc}{balance}}}} or
{\hyperref[\detokenize{fix_balance::doc}]{\sphinxcrossref{\DUrole{doc}{fix balance}}}} commands, then the initial regular grid
will persist for all simulations.  If balancing is performed, some of
the methods invoked by those commands retain the logical topology of
the initial 3d grid, and the mapping of processors to the grid
specified by the processors command.  However the grid spacings in
different dimensions may change, so that processors own subdomains of
different sizes.  If the {\hyperref[\detokenize{comm_style::doc}]{\sphinxcrossref{\DUrole{doc}{comm\_style tiled}}}} command is
used, methods invoked by the balancing commands may discard the 3d
grid of processors and tile the simulation domain with subdomains of
different sizes and shapes which no longer have a logical 3d
connectivity.  If that occurs, all the information specified by the
processors command is ignored.


\bigskip\hrule\bigskip


\sphinxAtStartPar
The \sphinxstyleemphasis{grid} keyword affects the factorization of P into Px,Py,Pz and it
can also affect how the P processor IDs are mapped to the 3d grid of
processors.

\sphinxAtStartPar
The \sphinxstyleemphasis{onelevel} style creates a 3d grid that is compatible with the
Px,Py,Pz settings, and which minimizes the surface\sphinxhyphen{}to\sphinxhyphen{}volume ratio of
each processor’s subdomain, as described above.  The mapping of
processors to the grid is determined by the \sphinxstyleemphasis{map} keyword setting.

\sphinxAtStartPar
The \sphinxstyleemphasis{twolevel} style can be used on machines with multicore nodes to
minimize off\sphinxhyphen{}node communication.  It ensures that contiguous
subsections of the 3d grid are assigned to all the cores of a node.
For example if \sphinxstyleemphasis{Nc} is 4, then 2x2x1 or 2x1x2 or 1x2x2 subsections of
the 3d grid will correspond to the cores of each node.  This affects
both the factorization and mapping steps.

\sphinxAtStartPar
The \sphinxstyleemphasis{Cx}, \sphinxstyleemphasis{Cy}, \sphinxstyleemphasis{Cz} settings are similar to the \sphinxstyleemphasis{Px}, \sphinxstyleemphasis{Py}, \sphinxstyleemphasis{Pz}
settings, only their product should equal \sphinxstyleemphasis{Nc}.  Any of the 3
parameters can be specified with an asterisk “*”, which means LAMMPS
will choose the number of cores in that dimension of the node’s
sub\sphinxhyphen{}grid.  As with Px,Py,Pz, it will do this based on the size and
shape of the global simulation box so as to minimize the
surface\sphinxhyphen{}to\sphinxhyphen{}volume ratio of each processor’s subdomain.

\begin{sphinxadmonition}{note}{Note:}
\sphinxAtStartPar
For the \sphinxstyleemphasis{twolevel} style to work correctly, it assumes the MPI
ranks of processors LAMMPS is running on are ordered by core and then
by node.  E.g. if you are running on 2 quad\sphinxhyphen{}core nodes, for a total of
8 processors, then it assumes processors 0,1,2,3 are on node 1, and
processors 4,5,6,7 are on node 2.  This is the default rank ordering
for most MPI implementations, but some MPIs provide options for this
ordering, e.g. via environment variable settings.
\end{sphinxadmonition}

\sphinxAtStartPar
The \sphinxstyleemphasis{numa} style operates similar to the \sphinxstyleemphasis{twolevel} keyword except that
it auto\sphinxhyphen{}detects which cores are running on which nodes.  It will also
subdivide the cores into numa domains. Currently, the number of numa
domains is not auto\sphinxhyphen{}detected and must be specified using the
\sphinxstyleemphasis{numa\_nodes} keyword; otherwise, the default value is used. The \sphinxstyleemphasis{numa}
style uses a different algorithm than the \sphinxstyleemphasis{twolevel} keyword for doing
the two\sphinxhyphen{}level factorization of the simulation box into a 3d processor
grid to minimize off\sphinxhyphen{}node communication and communication across numa
domains. It does its own MPI\sphinxhyphen{}based mapping of nodes and cores to the
regular 3d grid.  Thus it may produce a different layout of the
processors than the \sphinxstyleemphasis{twolevel} options.

\sphinxAtStartPar
The \sphinxstyleemphasis{numa} style will give an error if the number of MPI processes is
not divisible by the number of cores used per node, or any of the Px
or Py or Pz values is greater than 1.

\begin{sphinxadmonition}{note}{Note:}
\sphinxAtStartPar
Unlike the \sphinxstyleemphasis{twolevel} style, the \sphinxstyleemphasis{numa} style does not require
any particular ordering of MPI ranks in order to work correctly. This
is because it auto\sphinxhyphen{}detects which processes are running on which nodes.
However, it assumes that the lowest ranks are in the first numa
domain, and so forth. MPI rank orderings that do not preserve this
property might result in more intra\sphinxhyphen{}node communication between CPUs.
\end{sphinxadmonition}

\sphinxAtStartPar
The \sphinxstyleemphasis{custom} style uses the file \sphinxstyleemphasis{infile} to define both the 3d
factorization and the mapping of processors to the grid.

\sphinxAtStartPar
The file should have the following format.  Any number of initial
blank or comment lines (starting with a “\#” character) can be present.
The first non\sphinxhyphen{}blank, non\sphinxhyphen{}comment line should have
3 values:

\begin{sphinxVerbatim}[commandchars=\\\{\}]
Px Py Py
\end{sphinxVerbatim}

\sphinxAtStartPar
These must be compatible with the total number of processors
and the Px, Py, Pz settings of the processors command.

\sphinxAtStartPar
This line should be immediately followed by
P = Px*Py*Pz lines of the form:

\begin{sphinxVerbatim}[commandchars=\\\{\}]
ID I J K
\end{sphinxVerbatim}

\sphinxAtStartPar
where ID is a processor ID (from 0 to P\sphinxhyphen{}1) and I,J,K are the
processors location in the 3d grid.  I must be a number from 1 to Px
(inclusive) and similarly for J and K.  The P lines can be listed in
any order, but no processor ID should appear more than once.


\bigskip\hrule\bigskip


\sphinxAtStartPar
The \sphinxstyleemphasis{numa\_nodes} keyword is used to specify the number of numa domains
per node. It is currently only used by the \sphinxstyleemphasis{numa} style for two\sphinxhyphen{}level
factorization to reduce the amount of MPI communications between CPUs.
A good setting for this will typically be equal to the number of CPU
sockets per node.


\bigskip\hrule\bigskip


\sphinxAtStartPar
The \sphinxstyleemphasis{map} keyword affects how the P processor IDs (from 0 to P\sphinxhyphen{}1) are
mapped to the 3d grid of processors.  It is only used by the
\sphinxstyleemphasis{onelevel} and \sphinxstyleemphasis{twolevel} grid settings.

\sphinxAtStartPar
The \sphinxstyleemphasis{cart} style uses the family of MPI Cartesian functions to perform
the mapping, namely MPI\_Cart\_create(), MPI\_Cart\_get(),
MPI\_Cart\_shift(), and MPI\_Cart\_rank().  It invokes the
MPI\_Cart\_create() function with its reorder flag = 0, so that MPI is
not free to reorder the processors.

\sphinxAtStartPar
The \sphinxstyleemphasis{cart/reorder} style does the same thing as the \sphinxstyleemphasis{cart} style
except it sets the reorder flag to 1, so that MPI can reorder
processors if it desires.

\sphinxAtStartPar
The \sphinxstyleemphasis{xyz}, \sphinxstyleemphasis{xzy}, \sphinxstyleemphasis{yxz}, \sphinxstyleemphasis{yzx}, \sphinxstyleemphasis{zxy}, and \sphinxstyleemphasis{zyx} styles are all
similar.  If the style is IJK, then it maps the P processors to the
grid so that the processor ID in the I direction varies fastest, the
processor ID in the J direction varies next fastest, and the processor
ID in the K direction varies slowest.  For example, if you select
style \sphinxstyleemphasis{xyz} and you have a 2x2x2 grid of 8 processors, the assignments
of the 8 octants of the simulation domain will be:

\begin{sphinxVerbatim}[commandchars=\\\{\}]
proc 0 = lo x, lo y, lo z octant
proc 1 = hi x, lo y, lo z octant
proc 2 = lo x, hi y, lo z octant
proc 3 = hi x, hi y, lo z octant
proc 4 = lo x, lo y, hi z octant
proc 5 = hi x, lo y, hi z octant
proc 6 = lo x, hi y, hi z octant
proc 7 = hi x, hi y, hi z octant
\end{sphinxVerbatim}

\sphinxAtStartPar
Note that, in principle, an MPI implementation on a particular machine
should be aware of both the machine’s network topology and the
specific subset of processors and nodes that were assigned to your
simulation.  Thus its MPI\_Cart calls can optimize the assignment of
MPI processes to the 3d grid to minimize communication costs.  In
practice, however, few if any MPI implementations actually do this.
So it is likely that the \sphinxstyleemphasis{cart} and \sphinxstyleemphasis{cart/reorder} styles simply give
the same result as one of the IJK styles.

\sphinxAtStartPar
Also note, that for the \sphinxstyleemphasis{twolevel} grid style, the \sphinxstyleemphasis{map} setting is
used to first map the nodes to the 3d grid, then again to the cores
within each node.  For the latter step, the \sphinxstyleemphasis{cart} and \sphinxstyleemphasis{cart/reorder}
styles are not supported, so an \sphinxstyleemphasis{xyz} style is used in their place.


\bigskip\hrule\bigskip


\sphinxAtStartPar
The \sphinxstyleemphasis{part} keyword affects the factorization of P into Px,Py,Pz.

\sphinxAtStartPar
It can be useful when running in multi\sphinxhyphen{}partition mode, e.g. with the
{\hyperref[\detokenize{run_style::doc}]{\sphinxcrossref{\DUrole{doc}{run\_style verlet/split}}}} command.  It specifies a
dependency between a sending partition \sphinxstyleemphasis{Psend} and a receiving
partition \sphinxstyleemphasis{Precv} which is enforced when each is setting up their own
mapping of their processors to the simulation box.  Each of \sphinxstyleemphasis{Psend}
and \sphinxstyleemphasis{Precv} must be integers from 1 to Np, where Np is the number of
partitions you have defined via the {\hyperref[\detokenize{Run_options::doc}]{\sphinxcrossref{\DUrole{doc}{\sphinxhyphen{}partition command\sphinxhyphen{}line switch}}}}.

\sphinxAtStartPar
A “dependency” means that the sending partition will create its
regular 3d grid as Px by Py by Pz and after it has done this, it will
send the Px,Py,Pz values to the receiving partition.  The receiving
partition will wait to receive these values before creating its own
regular 3d grid and will use the sender’s Px,Py,Pz values as a
constraint.  The nature of the constraint is determined by the
\sphinxstyleemphasis{cstyle} argument.

\sphinxAtStartPar
For a \sphinxstyleemphasis{cstyle} of \sphinxstyleemphasis{multiple}, each dimension of the sender’s processor
grid is required to be an integer multiple of the corresponding
dimension in the receiver’s processor grid.  This is a requirement of
the {\hyperref[\detokenize{run_style::doc}]{\sphinxcrossref{\DUrole{doc}{run\_style verlet/split}}}} command.

\sphinxAtStartPar
For example, assume the sending partition creates a 4x6x10 grid = 240
processor grid.  If the receiving partition is running on 80
processors, it could create a 4x2x10 grid, but it will not create a
2x4x10 grid, since in the y\sphinxhyphen{}dimension, 6 is not an integer multiple of
4.

\begin{sphinxadmonition}{note}{Note:}
\sphinxAtStartPar
If you use the {\hyperref[\detokenize{partition::doc}]{\sphinxcrossref{\DUrole{doc}{partition}}}} command to invoke
different “processors” commands on different partitions, and you also
use the \sphinxstyleemphasis{part} keyword, then you must ensure that both the sending and
receiving partitions invoke the “processors” command that connects the
2 partitions via the \sphinxstyleemphasis{part} keyword.  LAMMPS cannot easily check for
this, but your simulation will likely hang in its setup phase if this
error has been made.
\end{sphinxadmonition}


\bigskip\hrule\bigskip


\sphinxAtStartPar
The \sphinxstyleemphasis{file} keyword writes the mapping of the factorization of P
processors and their mapping to the 3d grid to the specified file
\sphinxstyleemphasis{outfile}.  This is useful to check that you assigned physical
processors in the manner you desired, which can be tricky to figure
out, especially when running on multiple partitions or on, a multicore
machine or when the processor ranks were reordered by use of the
{\hyperref[\detokenize{Run_options::doc}]{\sphinxcrossref{\DUrole{doc}{\sphinxhyphen{}reorder command\sphinxhyphen{}line switch}}}} or due to use of
MPI\sphinxhyphen{}specific launch options such as a config file.

\sphinxAtStartPar
If you have multiple partitions you should ensure that each one writes
to a different file, e.g. using a {\hyperref[\detokenize{variable::doc}]{\sphinxcrossref{\DUrole{doc}{world\sphinxhyphen{}style variable}}}}
for the filename.  The file has a self\sphinxhyphen{}explanatory header, followed by
one\sphinxhyphen{}line per processor in this format:

\sphinxAtStartPar
world\sphinxhyphen{}ID universe\sphinxhyphen{}ID original\sphinxhyphen{}ID: I J K: name

\sphinxAtStartPar
The IDs are the processor’s rank in this simulation (the world), the
universe (of multiple simulations), and the original MPI communicator
used to instantiate LAMMPS, respectively.  The world and universe IDs
will only be different if you are running on more than one partition;
see the {\hyperref[\detokenize{Run_options::doc}]{\sphinxcrossref{\DUrole{doc}{\sphinxhyphen{}partition command\sphinxhyphen{}line switch}}}}.  The
universe and original IDs will only be different if you used the
{\hyperref[\detokenize{Run_options::doc}]{\sphinxcrossref{\DUrole{doc}{\sphinxhyphen{}reorder command\sphinxhyphen{}line switch}}}} to reorder the
processors differently than their rank in the original communicator
LAMMPS was instantiated with.

\sphinxAtStartPar
I,J,K are the indices of the processor in the regular 3d grid, each
from 1 to Nd, where Nd is the number of processors in that dimension
of the grid.

\sphinxAtStartPar
The \sphinxstyleemphasis{name} is what is returned by a call to MPI\_Get\_processor\_name()
and should represent an identifier relevant to the physical processors
in your machine.  Note that depending on the MPI implementation,
multiple cores can have the same \sphinxstyleemphasis{name}.


\bigskip\hrule\bigskip



\subsection{Restrictions}
\label{\detokenize{processors:restrictions}}
\sphinxAtStartPar
This command cannot be used after the simulation box is defined by a
{\hyperref[\detokenize{read_data::doc}]{\sphinxcrossref{\DUrole{doc}{read\_data}}}} or {\hyperref[\detokenize{create_box::doc}]{\sphinxcrossref{\DUrole{doc}{create\_box}}}} command.
It can be used before a restart file is read to change the 3d
processor grid from what is specified in the restart file.

\sphinxAtStartPar
The \sphinxstyleemphasis{grid numa} keyword only currently works with the \sphinxstyleemphasis{map cart}
option.

\sphinxAtStartPar
The \sphinxstyleemphasis{part} keyword (for the receiving partition) only works with the
\sphinxstyleemphasis{grid onelevel} or \sphinxstyleemphasis{grid twolevel} options.


\subsection{Related commands}
\label{\detokenize{processors:related-commands}}
\sphinxAtStartPar
{\hyperref[\detokenize{partition::doc}]{\sphinxcrossref{\DUrole{doc}{partition}}}}, {\hyperref[\detokenize{Run_options::doc}]{\sphinxcrossref{\DUrole{doc}{\sphinxhyphen{}reorder command\sphinxhyphen{}line switch}}}}


\subsection{Default}
\label{\detokenize{processors:default}}
\sphinxAtStartPar
The option defaults are Px Py Pz = * * *, grid = onelevel, map =
cart, and numa\_nodes = 2.

\sphinxstepscope

\index{python@\spxentry{python}}\ignorespaces 

\section{python command}
\label{\detokenize{python:python-command}}\label{\detokenize{python:index-0}}\label{\detokenize{python::doc}}

\subsection{Syntax}
\label{\detokenize{python:syntax}}
\begin{sphinxVerbatim}[commandchars=\\\{\}]
\PYG{k}{python}\PYG{+w}{ }\PYG{n}{mode}\PYG{+w}{ }\PYG{n}{keyword}\PYG{+w}{ }\PYG{n}{args}\PYG{+w}{ }\PYG{n}{...}
\end{sphinxVerbatim}
\begin{itemize}
\item {} 
\sphinxAtStartPar
mode = \sphinxstyleemphasis{source} or name of Python function

\sphinxAtStartPar
if mode is \sphinxstyleemphasis{source}:
\begin{sphinxalltt}
keyword = \sphinxstyleemphasis{here} or name of a \sphinxstyleemphasis{Python file}
  \sphinxstyleemphasis{here} arg = inline
    inline = one or more lines of Python code which defines func
             must be a single argument, typically enclosed between triple quotes
  \sphinxstyleemphasis{Python file} = name of a file with Python code which will be executed immediately
\end{sphinxalltt}

\item {} 
\sphinxAtStartPar
if \sphinxstyleemphasis{mode} is the name of a Python function, one or more keywords with/without arguments must be appended
\begin{sphinxalltt}
keyword = \sphinxstyleemphasis{invoke} or \sphinxstyleemphasis{input} or \sphinxstyleemphasis{return} or \sphinxstyleemphasis{format} or \sphinxstyleemphasis{length} or \sphinxstyleemphasis{file} or \sphinxstyleemphasis{here} or \sphinxstyleemphasis{exists}
  \sphinxstyleemphasis{invoke} arg = none = invoke the previously defined Python function
  \sphinxstyleemphasis{input} args = N i1 i2 ... iN
    N = \# of inputs to function
    i1,...,iN = value, SELF, or LAMMPS variable name
      value = integer number, floating point number, or string
      SELF = reference to LAMMPS itself which can be accessed by Python function
      variable = v\_name, where name = name of LAMMPS variable, e.g. v\_abc
  \sphinxstyleemphasis{return} arg = varReturn
    varReturn = v\_name  = LAMMPS variable name which the return value of the Python function will be assigned to
  \sphinxstyleemphasis{format} arg = fstring with M characters
    M = N if no return value, where N = \# of inputs
    M = N+1 if there is a return value
    fstring = each character (i,f,s,p) corresponds in order to an input or return value
    \textquotesingle{}i\textquotesingle{} = integer, \textquotesingle{}f\textquotesingle{} = floating point, \textquotesingle{}s\textquotesingle{} = string, \textquotesingle{}p\textquotesingle{} = SELF
  \sphinxstyleemphasis{length} arg = Nlen
    Nlen = max length of string returned from Python function
  \sphinxstyleemphasis{file} arg = filename
    filename = file of Python code, which defines func
  \sphinxstyleemphasis{here} arg = inline
    inline = one or more lines of Python code which defines func
             must be a single argument, typically enclosed between triple quotes
  \sphinxstyleemphasis{exists} arg = none = Python code has been loaded by previous python command
\end{sphinxalltt}

\end{itemize}


\subsection{Examples}
\label{\detokenize{python:examples}}
\begin{sphinxVerbatim}[commandchars=\\\{\}]
\PYG{k}{python}\PYG{+w}{ }\PYG{n}{pForce}\PYG{+w}{ }\PYG{n}{input}\PYG{+w}{ }\PYG{l+m}{2}\PYG{+w}{ }\PYG{n}{v\PYGZus{}x}\PYG{+w}{ }\PYG{l+m}{20.0}\PYG{+w}{ }\PYG{n}{return}\PYG{+w}{ }\PYG{n}{v\PYGZus{}f}\PYG{+w}{ }\PYG{n}{format}\PYG{+w}{ }\PYG{n}{fff}\PYG{+w}{ }\PYG{n}{file}\PYG{+w}{ }\PYG{n}{force.py}
\PYG{k}{python}\PYG{+w}{ }\PYG{n}{pForce}\PYG{+w}{ }\PYG{n}{invoke}

\PYG{k}{python}\PYG{+w}{ }\PYG{n}{factorial}\PYG{+w}{ }\PYG{n}{input}\PYG{+w}{ }\PYG{l+m}{1}\PYG{+w}{ }\PYG{n}{myN}\PYG{+w}{ }\PYG{n}{return}\PYG{+w}{ }\PYG{n}{v\PYGZus{}fac}\PYG{+w}{ }\PYG{n}{format}\PYG{+w}{ }\PYG{n}{ii}\PYG{+w}{ }\PYG{n}{here}\PYG{+w}{ }\PYG{l+s}{\PYGZdq{}}\PYG{l+s}{\PYGZdq{}}\PYG{l+s}{\PYGZdq{}}
\PYG{l+s}{def factorial(n):}
\PYG{l+s}{  if n == 1: return n}
\PYG{l+s}{  return n * factorial(n\PYGZhy{}1)}
\PYG{l+s}{ }\PYG{l+s}{\PYGZdq{}}\PYG{l+s}{\PYGZdq{}}\PYG{l+s}{\PYGZdq{}}

\PYG{k}{python}\PYG{+w}{ }\PYG{n}{loop}\PYG{+w}{ }\PYG{n}{input}\PYG{+w}{ }\PYG{l+m}{1}\PYG{+w}{ }\PYG{n}{SELF}\PYG{+w}{ }\PYG{n}{return}\PYG{+w}{ }\PYG{n}{v\PYGZus{}value}\PYG{+w}{ }\PYG{n}{format}\PYG{+w}{ }\PYG{n}{pf}\PYG{+w}{ }\PYG{n}{here}\PYG{+w}{ }\PYG{l+s}{\PYGZdq{}}\PYG{l+s}{\PYGZdq{}}\PYG{l+s}{\PYGZdq{}}
\PYG{l+s}{def loop(lmpptr,N,cut0):}
\PYG{l+s}{  from lammps import lammps}
\PYG{l+s}{  lmp = lammps(ptr=lmpptr)}

\PYG{l+s}{  \PYGZsh{} loop N times, increasing cutoff each time}

\PYG{l+s}{  for i in range(N):}
\PYG{l+s}{    cut = cut0 + i*0.1}
\PYG{l+s}{    lmp.set\PYGZus{}variable(}\PYG{l+s}{\PYGZdq{}}\PYG{n}{cut}\PYG{l+s}{\PYGZdq{}}\PYG{l+s}{,cut)               \PYGZsh{} set a variable in LAMMPS}
\PYG{l+s}{    lmp.command(}\PYG{l+s}{\PYGZdq{}}\PYG{n}{pair\PYGZus{}style}\PYG{+w}{ }\PYG{n}{lj}\PYG{o}{/}\PYG{n}{cut}\PYG{+w}{ }\PYG{n+nv}{\PYGZdl{}\PYGZob{}}\PYG{n+nv}{cut}\PYG{n+nv}{\PYGZcb{}}\PYG{l+s}{\PYGZdq{}}\PYG{l+s}{)   \PYGZsh{} LAMMPS commands}
\PYG{l+s}{    lmp.command(}\PYG{l+s}{\PYGZdq{}}\PYG{n}{pair\PYGZus{}coeff}\PYG{+w}{ }\PYG{o}{*}\PYG{+w}{ }\PYG{o}{*}\PYG{+w}{ }\PYG{l+m}{1.0}\PYG{+w}{ }\PYG{l+m}{1.0}\PYG{l+s}{\PYGZdq{}}\PYG{l+s}{)}
\PYG{l+s}{    lmp.command(}\PYG{l+s}{\PYGZdq{}}\PYG{n}{run}\PYG{+w}{ }\PYG{l+m}{100}\PYG{l+s}{\PYGZdq{}}\PYG{l+s}{)}
\PYG{l+s}{\PYGZdq{}}\PYG{l+s}{\PYGZdq{}}\PYG{l+s}{\PYGZdq{}}

\PYG{k}{python}\PYG{+w}{ }\PYG{n}{source}\PYG{+w}{ }\PYG{n}{funcdef.py}

\PYG{k}{python}\PYG{+w}{ }\PYG{n}{source}\PYG{+w}{ }\PYG{n}{here}\PYG{+w}{ }\PYG{l+s}{\PYGZdq{}}\PYG{l+s}{from lammps import lammps}\PYG{l+s}{\PYGZdq{}}
\end{sphinxVerbatim}


\subsection{Description}
\label{\detokenize{python:description}}
\sphinxAtStartPar
The \sphinxstyleemphasis{python} command allows interfacing LAMMPS with an embedded Python
interpreter and enables either executing arbitrary python code in that
interpreter, registering a Python function for future execution (as a
python style variable, from a fix interfaced with python, or for direct
invocation), or invoking such a previously registered function.

\sphinxAtStartPar
Arguments, including LAMMPS variables, can be passed to the function
from the LAMMPS input script and a value returned by the Python function
assigned to a LAMMPS variable.  The Python code for the function can be included
directly in the input script or in a separate Python file.  The function
can be standard Python code or it can make “callbacks” to LAMMPS through
its library interface to query or set internal values within LAMMPS.
This is a powerful mechanism for performing complex operations in a
LAMMPS input script that are not possible with the simple input script
and variable syntax which LAMMPS defines.  Thus your input script can
operate more like a true programming language.

\sphinxAtStartPar
Use of this command requires building LAMMPS with the PYTHON package
which links to the Python library so that the Python interpreter is
embedded in LAMMPS.  More details about this process are given below.

\sphinxAtStartPar
There are two ways to invoke a Python function once it has been
registered.  One is using the \sphinxstyleemphasis{invoke} keyword.  The other is to assign
the function to a {\hyperref[\detokenize{variable::doc}]{\sphinxcrossref{\DUrole{doc}{python\sphinxhyphen{}style variable}}}} defined in
your input script.  Whenever the variable is evaluated, it will execute
the Python function to assign a value to the variable.  Note that
variables can be evaluated in many different ways within LAMMPS.  They
can be substituted with their result directly in an input script, or
they can be passed to various commands as arguments, so that the
variable is evaluated during a simulation run.

\sphinxAtStartPar
A broader overview of how Python can be used with LAMMPS is given in the
{\hyperref[\detokenize{Python_head::doc}]{\sphinxcrossref{\DUrole{doc}{Use Python with LAMMPS}}}} section of the
documentation.  There also is an \sphinxcode{\sphinxupquote{examples/python}} directory which
illustrates use of the python command.


\bigskip\hrule\bigskip


\sphinxAtStartPar
The first argument of the \sphinxstyleemphasis{python} command is either the \sphinxstyleemphasis{source}
keyword or the name of a Python function.  This defines the mode
of the python command.

\sphinxAtStartPar
\DUrole{versionmodified,changed}{Changed in version 22Dec2022.}

\sphinxAtStartPar
If the \sphinxstyleemphasis{source} keyword is used, it is followed by either a file name or
the \sphinxstyleemphasis{here} keyword.  No other keywords can be used.  The \sphinxstyleemphasis{here} keyword
is followed by a string with python commands, either on a single line
enclosed in quotes, or as multiple lines enclosed in triple quotes.
These Python commands will be passed to the python interpreter and
executed immediately without registering a Python function for future
execution.  The code will be loaded into and run in the “main” module of
the Python interpreter.  This allows running arbitrary Python code at
any time while processing the LAMMPS input file.  This can be used to
pre\sphinxhyphen{}load Python modules, initialize global variables, define functions
or classes, or perform operations using the python programming language.
The Python code will be executed in parallel on all MPI processes.  No
arguments can be passed.

\sphinxAtStartPar
In all other cases, the first argument is the name of a Python function
that will be registered with LAMMPS for future execution.  The function
may already be defined (see \sphinxstyleemphasis{exists} keyword) or must be defined using
the \sphinxstyleemphasis{file} or \sphinxstyleemphasis{here} keywords as explained below.

\sphinxAtStartPar
If the \sphinxstyleemphasis{invoke} keyword is used, no other keywords can be used, and a
previous \sphinxstyleemphasis{python} command must have registered the Python function
referenced by this command.  This invokes the Python function with the
previously defined arguments and the return value is processed as
explained below.  You can invoke the function as many times as you wish
in your input script.

\sphinxAtStartPar
The \sphinxstyleemphasis{input} keyword defines how many arguments \sphinxstyleemphasis{N} the Python function
expects.  If it takes no arguments, then the \sphinxstyleemphasis{input} keyword should not
be used.  Each argument can be specified directly as a value, e.g. ‘6’
or ‘3.14159’ or ‘abc’ (a string of characters).  The type of each
argument is specified by the \sphinxstyleemphasis{format} keyword as explained below, so
that Python will know how to interpret the value.  If the word SELF is
used for an argument it has a special meaning.  A pointer is passed to
the Python function which it can convert into a reference to LAMMPS
itself using the {\hyperref[\detokenize{Python_module::doc}]{\sphinxcrossref{\DUrole{doc}{LAMMPS Python module}}}}.  This
enables the function to call back to LAMMPS through its library
interface as explained below.  This allows the Python function to query
or set values internal to LAMMPS which can affect the subsequent
execution of the input script.  A LAMMPS variable can also be used as an
argument, specified as v\_name, where “name” is the name of the variable.
Any style of LAMMPS variable returning a scalar or a string can be used,
as defined by the {\hyperref[\detokenize{variable::doc}]{\sphinxcrossref{\DUrole{doc}{variable}}}} command.  The \sphinxstyleemphasis{format}
keyword must be used to set the type of data that is passed to Python.
Each time the Python function is invoked, the LAMMPS variable is
evaluated and its value is passed to the Python function.

\sphinxAtStartPar
The \sphinxstyleemphasis{return} keyword is only needed if the Python function returns a
value.  The specified \sphinxstyleemphasis{varReturn} must be of the form v\_name, where
“name” is the name of a python\sphinxhyphen{}style LAMMPS variable, defined by the
{\hyperref[\detokenize{variable::doc}]{\sphinxcrossref{\DUrole{doc}{variable}}}} command.  The Python function can return a
numeric or string value, as specified by the \sphinxstyleemphasis{format} keyword.

\sphinxAtStartPar
As explained on the {\hyperref[\detokenize{variable::doc}]{\sphinxcrossref{\DUrole{doc}{variable}}}} doc page, the definition
of a python\sphinxhyphen{}style variable associates a Python function name with the
variable.  This must match the \sphinxstyleemphasis{Python function name} first argument of
the \sphinxstyleemphasis{python} command.  For example these two commands would be
consistent:

\begin{sphinxVerbatim}[commandchars=\\\{\}]
\PYG{k}{variable }\PYG{n+nv+nvIdentifier}{foo}\PYG{+w}{ }\PYG{n}{python}\PYG{+w}{ }\PYG{n}{myMultiply}
\PYG{k}{python}\PYG{+w}{ }\PYG{n}{myMultiply}\PYG{+w}{ }\PYG{n}{return}\PYG{+w}{ }\PYG{n}{v\PYGZus{}foo}\PYG{+w}{ }\PYG{n}{format}\PYG{+w}{ }\PYG{n}{f}\PYG{+w}{ }\PYG{n}{file}\PYG{+w}{ }\PYG{n}{funcs.py}
\end{sphinxVerbatim}

\sphinxAtStartPar
The two commands can appear in either order in the input script so
long as both are specified before the Python function is invoked for
the first time.  Afterwards, the variable ‘foo’ is associated with
the Python function ‘myMultiply’.

\sphinxAtStartPar
The \sphinxstyleemphasis{format} keyword must be used if the \sphinxstyleemphasis{input} or \sphinxstyleemphasis{return} keywords
are used.  It defines an \sphinxstyleemphasis{fstring} with M characters, where M = sum of
number of inputs and outputs.  The order of characters corresponds to
the N inputs, followed by the return value (if it exists).  Each
character must be one of the following: “i” for integer, “f” for
floating point, “s” for string, or “p” for SELF.  Each character defines
the type of the corresponding input or output value of the Python
function and affects the type conversion that is performed internally as
data is passed back and forth between LAMMPS and Python.  Note that it
is permissible to use a {\hyperref[\detokenize{variable::doc}]{\sphinxcrossref{\DUrole{doc}{python\sphinxhyphen{}style variable}}}} in a
LAMMPS command that allows for an equal\sphinxhyphen{}style variable as an argument,
but only if the output of the Python function is flagged as a numeric
value (“i” or “f”) via the \sphinxstyleemphasis{format} keyword.

\sphinxAtStartPar
If the \sphinxstyleemphasis{return} keyword is used and the \sphinxstyleemphasis{format} keyword specifies the
output as a string, then the default maximum length of that string is
63 characters (64\sphinxhyphen{}1 for the string terminator).  If you want to return
a longer string, the \sphinxstyleemphasis{length} keyword can be specified with its \sphinxstyleemphasis{Nlen}
value set to a larger number (the code allocates space for Nlen+1 to
include the string terminator).  If the Python function generates a
string longer than the default 63 or the specified \sphinxstyleemphasis{Nlen}, it will be
truncated.


\bigskip\hrule\bigskip


\sphinxAtStartPar
Either the \sphinxstyleemphasis{file}, \sphinxstyleemphasis{here}, or \sphinxstyleemphasis{exists} keyword must be used, but only
one of them.  These keywords specify what Python code to load into the
Python interpreter.  The \sphinxstyleemphasis{file} keyword gives the name of a file
containing Python code, which should end with a “.py” suffix.  The code
will be immediately loaded into and run in the “main” module of the
Python interpreter.  The Python code will be executed in parallel on all
MPI processes.  Note that Python code which contains a function
definition does not “execute” the function when it is run; it simply
defines the function so that it can be invoked later.

\sphinxAtStartPar
The \sphinxstyleemphasis{here} keyword does the same thing, except that the Python code
follows as a single argument to the \sphinxstyleemphasis{here} keyword.  This can be done
using triple quotes as delimiters, as in the examples above.  This
allows Python code to be listed verbatim in your input script, with
proper indentation, blank lines, and comments, as desired.  See the
{\hyperref[\detokenize{Commands_parse::doc}]{\sphinxcrossref{\DUrole{doc}{Commands parse}}}} doc page, for an explanation of
how triple quotes can be used as part of input script syntax.

\sphinxAtStartPar
The \sphinxstyleemphasis{exists} keyword takes no argument.  It means that Python code
containing the required Python function with the given name has already
been executed, for example by a \sphinxstyleemphasis{python source} command or in the same
file that was used previously with the \sphinxstyleemphasis{file} keyword.

\sphinxAtStartPar
Note that the Python code that is loaded and run must contain a function
with the specified function name.  To operate properly when later
invoked, the function code must match the \sphinxstyleemphasis{input} and \sphinxstyleemphasis{return} and
\sphinxstyleemphasis{format} keywords specified by the python command.  Otherwise Python
will generate an error.


\bigskip\hrule\bigskip


\sphinxAtStartPar
This section describes how Python code can be written to work with
LAMMPS.

\sphinxAtStartPar
Whether you load Python code from a file or directly from your input
script, via the \sphinxstyleemphasis{file} and \sphinxstyleemphasis{here} keywords, the code can be identical.
It must be indented properly as Python requires.  It can contain
comments or blank lines.  If the code is in your input script, it cannot
however contain triple\sphinxhyphen{}quoted Python strings, since that will conflict
with the triple\sphinxhyphen{}quote parsing that the LAMMPS input script performs.

\sphinxAtStartPar
All the Python code you specify via one or more python commands is
loaded into the Python “main” module, i.e. \sphinxcode{\sphinxupquote{\_\_name\_\_ == \textquotesingle{}\_\_main\_\_\textquotesingle{}}}.
The code can define global variables, define global functions, define
classes or execute statements that are outside of function definitions.
It can contain multiple functions, only one of which matches the \sphinxstyleemphasis{func}
setting in the python command.  This means you can use the \sphinxstyleemphasis{file}
keyword once to load several functions, and the \sphinxstyleemphasis{exists} keyword
thereafter in subsequent python commands to register the other functions
that were previously loaded with LAMMPS.

\sphinxAtStartPar
A Python function you define (or more generally, the code you load)
can import other Python modules or classes, it can make calls to other
system functions or functions you define, and it can access or modify
global variables (in the “main” module) which will persist between
successive function calls.  The latter can be useful, for example, to
prevent a function from being invoke multiple times per timestep by
different commands in a LAMMPS input script that access the returned
python\sphinxhyphen{}style variable associated with the function.  For example,
consider this function loaded with two global variables defined
outside the function:

\begin{sphinxVerbatim}[commandchars=\\\{\}]
\PYG{n}{nsteplast} \PYG{o}{=} \PYG{o}{\PYGZhy{}}\PYG{l+m+mi}{1}
\PYG{n}{nvaluelast} \PYG{o}{=} \PYG{l+m+mi}{0}

\PYG{k}{def}\PYG{+w}{ }\PYG{n+nf}{expensive}\PYG{p}{(}\PYG{n}{nstep}\PYG{p}{)}\PYG{p}{:}
  \PYG{k}{global} \PYG{n}{nsteplast}\PYG{p}{,}\PYG{n}{nvaluelast}
  \PYG{k}{if} \PYG{n}{nstep} \PYG{o}{==} \PYG{n}{nsteplast}\PYG{p}{:} \PYG{k}{return} \PYG{n}{nvaluelast}
  \PYG{n}{nsteplast} \PYG{o}{=} \PYG{n}{nstep}
  \PYG{c+c1}{\PYGZsh{} perform complicated calculation}
  \PYG{n}{nvalue} \PYG{o}{=} \PYG{o}{.}\PYG{o}{.}\PYG{o}{.}
  \PYG{n}{nvaluelast} \PYG{o}{=} \PYG{n}{nvalue}
  \PYG{k}{return} \PYG{n}{nvalue}
\end{sphinxVerbatim}

\sphinxAtStartPar
The variable ‘nsteplast’ stores the previous timestep the function was
invoked (passed as an argument to the function).  The variable
‘nvaluelast’ stores the return value computed on the last function
invocation.  If the function is invoked again on the same timestep, the
previous value is simply returned, without re\sphinxhyphen{}computing it.  The
“global” statement inside the Python function allows it to overwrite the
global variables from within the local context of the function.

\sphinxAtStartPar
Note that if you load Python code multiple times (via multiple python
commands), you can overwrite previously loaded variables and functions
if you are not careful.  E.g. if the code above were loaded twice, the
global variables would be re\sphinxhyphen{}initialized, which might not be what you
want.  Likewise, if a function with the same name exists in two chunks
of Python code you load, the function loaded second will override the
function loaded first.

\sphinxAtStartPar
It’s important to realize that if you are running LAMMPS in parallel,
each MPI task will load the Python interpreter and execute a local
copy of the Python function(s) you define.  There is no connection
between the Python interpreters running on different processors.
This implies three important things.

\sphinxAtStartPar
First, if you put a print or other statement creating output to the
screen in your Python function, you will see P copies of the output,
when running on P processors.  If the prints occur at (nearly) the same
time, the P copies of the output may be mixed together.  When loading
the LAMMPS Python module into the embedded Python interpreter, it is
possible to pass the pointer to the current LAMMPS class instance and
via the Python interface to the LAMMPS library interface, it is possible
to determine the MPI rank of the current process and thus adapt the
Python code so that output will only appear on MPI rank 0.  The
following LAMMPS input demonstrates how this could be done. The text
‘Hello, LAMMPS!’ should be printed only once, even when running LAMMPS
in parallel.

\begin{sphinxVerbatim}[commandchars=\\\{\}]
\PYG{k}{python}\PYG{+w}{ }\PYG{n}{python\PYGZus{}hello}\PYG{+w}{ }\PYG{n}{input}\PYG{+w}{ }\PYG{l+m}{1}\PYG{+w}{ }\PYG{n}{SELF}\PYG{+w}{ }\PYG{n}{format}\PYG{+w}{ }\PYG{n}{p}\PYG{+w}{ }\PYG{n}{here}\PYG{+w}{ }\PYG{l+s}{\PYGZdq{}}\PYG{l+s}{\PYGZdq{}}\PYG{l+s}{\PYGZdq{}}
\PYG{l+s}{def python\PYGZus{}hello(handle):}
\PYG{l+s}{    from lammps import lammps}
\PYG{l+s}{    lmp = lammps(ptr=handle)}
\PYG{l+s}{    me = lmp.extract\PYGZus{}setting(\PYGZsq{}world\PYGZus{}rank\PYGZsq{})}
\PYG{l+s}{    if me == 0:}
\PYG{l+s}{        print(\PYGZsq{}Hello, LAMMPS!\PYGZsq{})}
\PYG{l+s}{\PYGZdq{}}\PYG{l+s}{\PYGZdq{}}\PYG{l+s}{\PYGZdq{}}

\PYG{k}{python}\PYG{+w}{ }\PYG{n}{python\PYGZus{}hello}\PYG{+w}{ }\PYG{n}{invoke}
\end{sphinxVerbatim}

\sphinxAtStartPar
If your Python code loads Python modules that are not pre\sphinxhyphen{}loaded by the
Python library, then it will load the module from disk.  This may be a
bottleneck if 1000s of processors try to load a module at the same time.
On some large supercomputers, loading of modules from disk by Python may
be disabled.  In this case you would need to pre\sphinxhyphen{}build a Python library
that has the required modules pre\sphinxhyphen{}loaded and link LAMMPS with that
library.

\sphinxAtStartPar
Third, if your Python code calls back to LAMMPS (discussed in the
next section) and causes LAMMPS to perform an MPI operation requires
global communication (e.g. via MPI\_Allreduce), such as computing the
global temperature of the system, then you must ensure all your Python
functions (running independently on different processors) call back to
LAMMPS.  Otherwise the code may hang.


\bigskip\hrule\bigskip


\sphinxAtStartPar
Your Python function can “call back” to LAMMPS through its
library interface, if you use the SELF input to pass Python
a pointer to LAMMPS.  The mechanism for doing this in your
Python function is as follows:

\begin{sphinxVerbatim}[commandchars=\\\{\}]
\PYG{k}{def}\PYG{+w}{ }\PYG{n+nf}{foo}\PYG{p}{(}\PYG{n}{handle}\PYG{p}{,}\PYG{o}{.}\PYG{o}{.}\PYG{o}{.}\PYG{p}{)}\PYG{p}{:}
  \PYG{k+kn}{from}\PYG{+w}{ }\PYG{n+nn}{lammps}\PYG{+w}{ }\PYG{k+kn}{import} \PYG{n}{lammps}
  \PYG{n}{lmp} \PYG{o}{=} \PYG{n}{lammps}\PYG{p}{(}\PYG{n}{ptr}\PYG{o}{=}\PYG{n}{handle}\PYG{p}{)}
  \PYG{n}{lmp}\PYG{o}{.}\PYG{n}{command}\PYG{p}{(}\PYG{l+s+s1}{\PYGZsq{}}\PYG{l+s+s1}{print }\PYG{l+s+s1}{\PYGZdq{}}\PYG{l+s+s1}{Hello from inside Python}\PYG{l+s+s1}{\PYGZdq{}}\PYG{l+s+s1}{\PYGZsq{}}\PYG{p}{)}
  \PYG{o}{.}\PYG{o}{.}\PYG{o}{.}
\end{sphinxVerbatim}

\sphinxAtStartPar
The function definition must include a variable (‘handle’ in this case)
which corresponds to SELF in the \sphinxstyleemphasis{python} command.  The first line of
the function imports the {\hyperref[\detokenize{Python_module::doc}]{\sphinxcrossref{\DUrole{doc}{“lammps” Python module}}}}.
The second line creates a Python object \sphinxcode{\sphinxupquote{lmp}} which wraps the instance
of LAMMPS that called the function.  The ‘ptr=handle’ argument is what
makes that happen.  The third line invokes the command() function in the
LAMMPS library interface.  It takes a single string argument which is a
LAMMPS input script command for LAMMPS to execute, the same as if it
appeared in your input script.  In this case, LAMMPS should output

\begin{sphinxVerbatim}[commandchars=\\\{\}]
Hello from inside Python
\end{sphinxVerbatim}

\sphinxAtStartPar
to the screen and log file.  Note that since the LAMMPS print command
itself takes a string in quotes as its argument, the Python string
must be delimited with a different style of quotes.

\sphinxAtStartPar
The {\hyperref[\detokenize{Python_head::doc}]{\sphinxcrossref{\DUrole{doc}{Use Python with LAMMPS}}}} page describes the syntax
for how Python wraps the various functions included in the LAMMPS
library interface.

\sphinxAtStartPar
A more interesting example is in the \sphinxcode{\sphinxupquote{examples/python/in.python}} script
which loads and runs the following function from \sphinxcode{\sphinxupquote{examples/python/funcs.py}}:

\begin{sphinxVerbatim}[commandchars=\\\{\}]
\PYG{k}{def}\PYG{+w}{ }\PYG{n+nf}{loop}\PYG{p}{(}\PYG{n}{N}\PYG{p}{,}\PYG{n}{cut0}\PYG{p}{,}\PYG{n}{thresh}\PYG{p}{,}\PYG{n}{lmpptr}\PYG{p}{)}\PYG{p}{:}
  \PYG{n+nb}{print}\PYG{p}{(}\PYG{l+s+s2}{\PYGZdq{}}\PYG{l+s+s2}{LOOP ARGS}\PYG{l+s+s2}{\PYGZdq{}}\PYG{p}{,} \PYG{n}{N}\PYG{p}{,} \PYG{n}{cut0}\PYG{p}{,} \PYG{n}{thresh}\PYG{p}{,} \PYG{n}{lmpptr}\PYG{p}{)}
  \PYG{k+kn}{from}\PYG{+w}{ }\PYG{n+nn}{lammps}\PYG{+w}{ }\PYG{k+kn}{import} \PYG{n}{lammps}
  \PYG{n}{lmp} \PYG{o}{=} \PYG{n}{lammps}\PYG{p}{(}\PYG{n}{ptr}\PYG{o}{=}\PYG{n}{lmpptr}\PYG{p}{)}
  \PYG{n}{natoms} \PYG{o}{=} \PYG{n}{lmp}\PYG{o}{.}\PYG{n}{get\PYGZus{}natoms}\PYG{p}{(}\PYG{p}{)}

  \PYG{k}{for} \PYG{n}{i} \PYG{o+ow}{in} \PYG{n+nb}{range}\PYG{p}{(}\PYG{n}{N}\PYG{p}{)}\PYG{p}{:}
    \PYG{n}{cut} \PYG{o}{=} \PYG{n}{cut0} \PYG{o}{+} \PYG{n}{i}\PYG{o}{*}\PYG{l+m+mf}{0.1}

    \PYG{n}{lmp}\PYG{o}{.}\PYG{n}{set\PYGZus{}variable}\PYG{p}{(}\PYG{l+s+s2}{\PYGZdq{}}\PYG{l+s+s2}{cut}\PYG{l+s+s2}{\PYGZdq{}}\PYG{p}{,}\PYG{n}{cut}\PYG{p}{)}                 \PYG{c+c1}{\PYGZsh{} set a variable in LAMMPS}
    \PYG{n}{lmp}\PYG{o}{.}\PYG{n}{command}\PYG{p}{(}\PYG{l+s+s2}{\PYGZdq{}}\PYG{l+s+s2}{pair\PYGZus{}style lj/cut \PYGZdl{}}\PYG{l+s+si}{\PYGZob{}cut\PYGZcb{}}\PYG{l+s+s2}{\PYGZdq{}}\PYG{p}{)}     \PYG{c+c1}{\PYGZsh{} LAMMPS command}
    \PYG{c+c1}{\PYGZsh{}lmp.command(\PYGZdq{}pair\PYGZus{}style lj/cut \PYGZpc{}d\PYGZdq{} \PYGZpc{} cut)  \PYGZsh{} LAMMPS command option}

    \PYG{n}{lmp}\PYG{o}{.}\PYG{n}{command}\PYG{p}{(}\PYG{l+s+s2}{\PYGZdq{}}\PYG{l+s+s2}{pair\PYGZus{}coeff * * 1.0 1.0}\PYG{l+s+s2}{\PYGZdq{}}\PYG{p}{)}       \PYG{c+c1}{\PYGZsh{} ditto}
    \PYG{n}{lmp}\PYG{o}{.}\PYG{n}{command}\PYG{p}{(}\PYG{l+s+s2}{\PYGZdq{}}\PYG{l+s+s2}{run 10}\PYG{l+s+s2}{\PYGZdq{}}\PYG{p}{)}                       \PYG{c+c1}{\PYGZsh{} ditto}
    \PYG{n}{pe} \PYG{o}{=} \PYG{n}{lmp}\PYG{o}{.}\PYG{n}{extract\PYGZus{}compute}\PYG{p}{(}\PYG{l+s+s2}{\PYGZdq{}}\PYG{l+s+s2}{thermo\PYGZus{}pe}\PYG{l+s+s2}{\PYGZdq{}}\PYG{p}{,}\PYG{l+m+mi}{0}\PYG{p}{,}\PYG{l+m+mi}{0}\PYG{p}{)}   \PYG{c+c1}{\PYGZsh{} extract total PE from LAMMPS}
    \PYG{n+nb}{print}\PYG{p}{(}\PYG{l+s+s2}{\PYGZdq{}}\PYG{l+s+s2}{PE}\PYG{l+s+s2}{\PYGZdq{}}\PYG{p}{,} \PYG{n}{pe}\PYG{o}{/}\PYG{n}{natoms}\PYG{p}{,} \PYG{n}{thresh}\PYG{p}{)}
    \PYG{k}{if} \PYG{n}{pe}\PYG{o}{/}\PYG{n}{natoms} \PYG{o}{\PYGZlt{}} \PYG{n}{thresh}\PYG{p}{:} \PYG{k}{return}
\end{sphinxVerbatim}

\sphinxAtStartPar
with these input script commands:

\begin{sphinxVerbatim}[commandchars=\\\{\}]
\PYG{k}{python}\PYG{+w}{          }\PYG{n}{loop}\PYG{+w}{ }\PYG{n}{input}\PYG{+w}{ }\PYG{l+m}{4}\PYG{+w}{ }\PYG{l+m}{10}\PYG{+w}{ }\PYG{l+m}{1.0}\PYG{+w}{ }\PYG{o}{\PYGZhy{}}\PYG{l+m}{4.0}\PYG{+w}{ }\PYG{n}{SELF}\PYG{+w}{ }\PYG{n}{format}\PYG{+w}{ }\PYG{n}{iffp}\PYG{+w}{ }\PYG{n}{file}\PYG{+w}{ }\PYG{n}{funcs.py}
\PYG{k}{python}\PYG{+w}{          }\PYG{n}{loop}\PYG{+w}{ }\PYG{n}{invoke}
\end{sphinxVerbatim}

\sphinxAtStartPar
This has the effect of looping over a series of 10 short runs (10
timesteps each) where the pair style cutoff is increased from a value
of 1.0 in distance units, in increments of 0.1.  The looping stops
when the per\sphinxhyphen{}atom potential energy falls below a threshold of \sphinxhyphen{}4.0 in
energy units.  More generally, Python can be used to implement a loop
with complex logic, much more so than can be created using the LAMMPS
{\hyperref[\detokenize{jump::doc}]{\sphinxcrossref{\DUrole{doc}{jump}}}} and {\hyperref[\detokenize{if::doc}]{\sphinxcrossref{\DUrole{doc}{if}}}} commands.

\sphinxAtStartPar
Several LAMMPS library functions are called from the loop function.
Get\_natoms() returns the number of atoms in the simulation, so that it
can be used to normalize the potential energy that is returned by
extract\_compute() for the “thermo\_pe” compute that is defined by
default for LAMMPS thermodynamic output.  Set\_variable() sets the
value of a string variable defined in LAMMPS.  This library function
is a useful way for a Python function to return multiple values to
LAMMPS, more than the single value that can be passed back via a
return statement.  This cutoff value in the “cut” variable is then
substituted (by LAMMPS) in the pair\_style command that is executed
next.  Alternatively, the “LAMMPS command option” line could be used
in place of the 2 preceding lines, to have Python insert the value
into the LAMMPS command string.

\begin{sphinxadmonition}{note}{Note:}
\sphinxAtStartPar
When using the callback mechanism just described, recognize that
there are some operations you should not attempt because LAMMPS cannot
execute them correctly.  If the Python function is invoked between
runs in the LAMMPS input script, then it should be OK to invoke any
LAMMPS input script command via the library interface command() or
file() functions, so long as the command would work if it were
executed in the LAMMPS input script directly at the same point.
\end{sphinxadmonition}

\sphinxAtStartPar
However, a Python function can also be invoked during a run, whenever
an associated LAMMPS variable it is assigned to is evaluated.  If the
variable is an input argument to another LAMMPS command (e.g. {\hyperref[\detokenize{fix_setforce::doc}]{\sphinxcrossref{\DUrole{doc}{fix setforce}}}}), then the Python function will be invoked
inside the class for that command, in one of its methods that is
invoked in the middle of a timestep.  You cannot execute arbitrary
input script commands from the Python function (again, via the
command() or file() functions) at that point in the run and expect it
to work.  Other library functions such as those that invoke computes
or other variables may have hidden side effects as well.  In these
cases, LAMMPS has no simple way to check that something illogical is
being attempted.

\sphinxAtStartPar
The same applies to Python functions called during a simulation run at
each time step using {\hyperref[\detokenize{fix_python_invoke::doc}]{\sphinxcrossref{\DUrole{doc}{fix python/invoke}}}}.


\bigskip\hrule\bigskip


\sphinxAtStartPar
If you run Python code directly on your workstation, either
interactively or by using Python to launch a Python script stored in a
file, and your code has an error, you will typically see informative
error messages.  That is not the case when you run Python code from
LAMMPS using an embedded Python interpreter.  The code will typically
fail silently.  LAMMPS will catch some errors but cannot tell you
where in the Python code the problem occurred.  For example, if the
Python code cannot be loaded and run because it has syntax or other
logic errors, you may get an error from Python pointing to the
offending line, or you may get one of these generic errors from
LAMMPS:

\begin{sphinxVerbatim}[commandchars=\\\{\}]
Could not process Python file
Could not process Python string
\end{sphinxVerbatim}

\sphinxAtStartPar
When the Python function is invoked, if it does not return properly,
you will typically get this generic error from LAMMPS:

\begin{sphinxVerbatim}[commandchars=\\\{\}]
Python function evaluation failed
\end{sphinxVerbatim}

\sphinxAtStartPar
Here are three suggestions for debugging your Python code while
running it under LAMMPS.

\sphinxAtStartPar
First, don’t run it under LAMMPS, at least to start with!  Debug it
using plain Python.  Load and invoke your function, pass it arguments,
check return values, etc.

\sphinxAtStartPar
Second, add Python print statements to the function to check how far
it gets and intermediate values it calculates.  See the discussion
above about printing from Python when running in parallel.

\sphinxAtStartPar
Third, use Python exception handling.  For example, say this statement
in your Python function is failing, because you have not initialized the
variable foo:

\begin{sphinxVerbatim}[commandchars=\\\{\}]
\PYG{n}{foo} \PYG{o}{+}\PYG{o}{=} \PYG{l+m+mi}{1}
\end{sphinxVerbatim}

\sphinxAtStartPar
If you put one (or more) statements inside a “try” statement,
like this:

\begin{sphinxVerbatim}[commandchars=\\\{\}]
\PYG{k+kn}{import}\PYG{+w}{ }\PYG{n+nn}{exceptions}
\PYG{n+nb}{print}\PYG{p}{(}\PYG{l+s+s2}{\PYGZdq{}}\PYG{l+s+s2}{Inside simple function}\PYG{l+s+s2}{\PYGZdq{}}\PYG{p}{)}
\PYG{k}{try}\PYG{p}{:}
  \PYG{n}{foo} \PYG{o}{+}\PYG{o}{=} \PYG{l+m+mi}{1}      \PYG{c+c1}{\PYGZsh{} one or more statements here}
\PYG{k}{except} \PYG{n+ne}{Exception} \PYG{k}{as} \PYG{n}{e}\PYG{p}{:}
  \PYG{n+nb}{print}\PYG{p}{(}\PYG{l+s+s2}{\PYGZdq{}}\PYG{l+s+s2}{FOO error:}\PYG{l+s+s2}{\PYGZdq{}}\PYG{p}{,} \PYG{n}{e}\PYG{p}{)}
\end{sphinxVerbatim}

\sphinxAtStartPar
then you will get this message printed to the screen:

\begin{sphinxVerbatim}[commandchars=\\\{\}]
FOO error: local variable \PYGZsq{}foo\PYGZsq{} referenced before assignment
\end{sphinxVerbatim}

\sphinxAtStartPar
If there is no error in the try statements, then nothing is printed.
Either way the function continues on (unless you put a return or
sys.exit() in the except clause).


\bigskip\hrule\bigskip



\subsection{Restrictions}
\label{\detokenize{python:restrictions}}
\sphinxAtStartPar
This command is part of the PYTHON package.  It is only enabled if
LAMMPS was built with that package.  See the {\hyperref[\detokenize{Build_package::doc}]{\sphinxcrossref{\DUrole{doc}{Build package}}}} page for more info.

\sphinxAtStartPar
Building LAMMPS with the PYTHON package will link LAMMPS with the Python
library on your system.  Settings to enable this are in the
lib/python/Makefile.lammps file.  See the lib/python/README file for
information on those settings.

\sphinxAtStartPar
If you use Python code which calls back to LAMMPS, via the SELF input
argument explained above, there is an extra step required when building
LAMMPS.  LAMMPS must also be built as a shared library and your Python
function must be able to load the {\hyperref[\detokenize{Python_module::doc}]{\sphinxcrossref{\DUrole{doc}{“lammps” Python module}}}} that wraps the LAMMPS library interface.  These are the
same steps required to use Python by itself to wrap LAMMPS.  Details on
these steps are explained on the {\hyperref[\detokenize{Python_head::doc}]{\sphinxcrossref{\DUrole{doc}{Python}}}} doc page.
Note that it is important that the stand\sphinxhyphen{}alone LAMMPS executable and the
LAMMPS shared library be consistent (built from the same source code
files) in order for this to work.  If the two have been built at
different times using different source files, problems may occur.

\sphinxAtStartPar
Another limitation of calling back to Python from the LAMMPS module
using the \sphinxstyleemphasis{python} command in a LAMMPS input is that both, the Python
interpreter and LAMMPS, must be linked to the same Python runtime as a
shared library.  If the Python interpreter is linked to Python
statically (which seems to happen with Conda) then loading the shared
LAMMPS library will create a second python “main” module that hides the
one from the Python interpreter and all previous defined function and
global variables will become invisible.


\subsection{Related commands}
\label{\detokenize{python:related-commands}}
\sphinxAtStartPar
{\hyperref[\detokenize{shell::doc}]{\sphinxcrossref{\DUrole{doc}{shell}}}}, {\hyperref[\detokenize{variable::doc}]{\sphinxcrossref{\DUrole{doc}{variable}}}}, {\hyperref[\detokenize{fix_python_invoke::doc}]{\sphinxcrossref{\DUrole{doc}{fix python/invoke}}}}


\subsection{Default}
\label{\detokenize{python:default}}
\sphinxAtStartPar
none

\sphinxstepscope

\index{quit@\spxentry{quit}}\ignorespaces 

\section{quit command}
\label{\detokenize{quit:quit-command}}\label{\detokenize{quit:index-0}}\label{\detokenize{quit::doc}}

\subsection{Syntax}
\label{\detokenize{quit:syntax}}
\begin{sphinxVerbatim}[commandchars=\\\{\}]
\PYG{k}{quit}\PYG{+w}{ }\PYG{n}{status}
\end{sphinxVerbatim}

\sphinxAtStartPar
status = numerical exit status (optional)


\subsection{Examples}
\label{\detokenize{quit:examples}}
\begin{sphinxVerbatim}[commandchars=\\\{\}]
\PYG{k}{quit}
\PYG{k}{if}\PYG{+w}{ }\PYG{l+s}{\PYGZdq{}}\PYG{l+s}{\PYGZdl{}n \PYGZgt{} 10000}\PYG{l+s}{\PYGZdq{}}\PYG{+w}{ }\PYG{k}{then}\PYG{+w}{ }\PYG{l+s}{\PYGZdq{}}\PYG{l+s}{quit 1}\PYG{l+s}{\PYGZdq{}}
\end{sphinxVerbatim}


\subsection{Description}
\label{\detokenize{quit:description}}
\sphinxAtStartPar
This command causes LAMMPS to exit, after shutting down all output
cleanly.

\sphinxAtStartPar
It can be used as a debug statement in an input script, to terminate
the script at some intermediate point.

\sphinxAtStartPar
It can also be used as an invoked command inside the “then” or “else”
portion of an {\hyperref[\detokenize{if::doc}]{\sphinxcrossref{\DUrole{doc}{if}}}} command.

\sphinxAtStartPar
The optional status argument is an integer which signals the return
status to a program calling LAMMPS.  A return status of 0 usually
indicates success.  A status != 0 is failure, where the specified
value can be used to distinguish the kind of error, e.g. where in the
input script the quit was invoked.  If not specified, a status of 0 is
returned.


\subsection{Restrictions}
\label{\detokenize{quit:restrictions}}\begin{quote}

\sphinxAtStartPar
none
\end{quote}


\subsection{Related commands}
\label{\detokenize{quit:related-commands}}
\sphinxAtStartPar
{\hyperref[\detokenize{if::doc}]{\sphinxcrossref{\DUrole{doc}{if}}}}


\subsection{Default}
\label{\detokenize{quit:default}}
\sphinxAtStartPar
none

\sphinxstepscope

\index{read\_data@\spxentry{read\_data}}\ignorespaces 

\section{read\_data command}
\label{\detokenize{read_data:read-data-command}}\label{\detokenize{read_data:index-0}}\label{\detokenize{read_data::doc}}

\subsection{Syntax}
\label{\detokenize{read_data:syntax}}
\begin{sphinxVerbatim}[commandchars=\\\{\}]
\PYG{k}{read\PYGZus{}data}\PYG{+w}{ }\PYG{n}{file}\PYG{+w}{ }\PYG{n}{keyword}\PYG{+w}{ }\PYG{n}{args}\PYG{+w}{ }\PYG{n}{...}
\end{sphinxVerbatim}
\begin{itemize}
\item {} 
\sphinxAtStartPar
file = name of data file to read in

\item {} 
\sphinxAtStartPar
zero or more keyword/arg pairs may be appended

\item {} 
\sphinxAtStartPar
keyword = \sphinxstyleemphasis{add} or \sphinxstyleemphasis{offset} or \sphinxstyleemphasis{shift} or \sphinxstyleemphasis{extra/atom/types} or \sphinxstyleemphasis{extra/bond/types} or \sphinxstyleemphasis{extra/angle/types} or \sphinxstyleemphasis{extra/dihedral/types} or \sphinxstyleemphasis{extra/improper/types} or \sphinxstyleemphasis{extra/bond/per/atom} or \sphinxstyleemphasis{extra/angle/per/atom} or \sphinxstyleemphasis{extra/dihedral/per/atom} or \sphinxstyleemphasis{extra/improper/per/atom} or \sphinxstyleemphasis{extra/special/per/atom} or \sphinxstyleemphasis{group} or \sphinxstyleemphasis{nocoeff} or \sphinxstyleemphasis{fix}
\begin{sphinxalltt}
\sphinxstyleemphasis{add} arg = \sphinxstyleemphasis{append} or \sphinxstyleemphasis{IDoffset} or \sphinxstyleemphasis{IDoffset MOLoffset} or \sphinxstyleemphasis{merge}
  append = add new atoms with atom IDs appended to current IDs
  IDoffset = add new atoms with atom IDs having IDoffset added
  MOLoffset = add new atoms with molecule IDs having MOLoffset added (only when molecule IDs are enabled)
  merge = add new atoms with their atom IDs (and molecule IDs) unchanged
\sphinxstyleemphasis{offset} args = toff boff aoff doff ioff
  toff = offset to add to atom types
  boff = offset to add to bond types
  aoff = offset to add to angle types
  doff = offset to add to dihedral types
  ioff = offset to add to improper types
\sphinxstyleemphasis{shift} args = Sx Sy Sz
  Sx,Sy,Sz = distance to shift atoms when adding to system (distance units)
\sphinxstyleemphasis{extra/atom/types} arg = \# of extra atom types
\sphinxstyleemphasis{extra/bond/types} arg = \# of extra bond types
\sphinxstyleemphasis{extra/angle/types} arg = \# of extra angle types
\sphinxstyleemphasis{extra/dihedral/types} arg = \# of extra dihedral types
\sphinxstyleemphasis{extra/improper/types} arg = \# of extra improper types
\sphinxstyleemphasis{extra/bond/per/atom} arg = leave space for this many new bonds per atom
\sphinxstyleemphasis{extra/angle/per/atom} arg = leave space for this many new angles per atom
\sphinxstyleemphasis{extra/dihedral/per/atom} arg = leave space for this many new dihedrals per atom
\sphinxstyleemphasis{extra/improper/per/atom} arg = leave space for this many new impropers per atom
\sphinxstyleemphasis{extra/special/per/atom} arg = leave space for extra 1\sphinxhyphen{}2,1\sphinxhyphen{}3,1\sphinxhyphen{}4 interactions per atom
\sphinxstyleemphasis{group} args = groupID
  groupID = add atoms in data file to this group
\sphinxstyleemphasis{nocoeff} = ignore force field parameters
\sphinxstyleemphasis{fix} args = fix\sphinxhyphen{}ID header\sphinxhyphen{}string section\sphinxhyphen{}string
  fix\sphinxhyphen{}ID = ID of fix to process header lines and sections of data file
  header\sphinxhyphen{}string = header lines containing this string will be passed to fix
  section\sphinxhyphen{}string = section names with this string will be passed to fix
\end{sphinxalltt}

\end{itemize}


\subsection{Examples}
\label{\detokenize{read_data:examples}}
\begin{sphinxVerbatim}[commandchars=\\\{\}]
\PYG{k}{read\PYGZus{}data}\PYG{+w}{ }\PYG{n}{data.lj}
\PYG{k}{read\PYGZus{}data}\PYG{+w}{ }\PYG{n}{..}\PYG{o}{/}\PYG{n}{run7}\PYG{o}{/}\PYG{n}{data.polymer.gz}
\PYG{k}{read\PYGZus{}data}\PYG{+w}{ }\PYG{n}{data.protein}\PYG{+w}{ }\PYG{k}{fix }\PYG{n+nv+nvIdentifier}{mycmap}\PYG{+w}{ }\PYG{n+nv+nvIdentifier}{crossterm}\PYG{+w}{ }\PYG{n}{CMAP}
\PYG{k}{read\PYGZus{}data}\PYG{+w}{ }\PYG{n}{data.water}\PYG{+w}{ }\PYG{n}{add}\PYG{+w}{ }\PYG{n}{append}\PYG{+w}{ }\PYG{n}{offset}\PYG{+w}{ }\PYG{l+m}{3}\PYG{+w}{ }\PYG{l+m}{1}\PYG{+w}{ }\PYG{l+m}{1}\PYG{+w}{ }\PYG{l+m}{1}\PYG{+w}{ }\PYG{l+m}{1}\PYG{+w}{ }\PYG{n}{shift}\PYG{+w}{ }\PYG{l+m}{0.0}\PYG{+w}{ }\PYG{l+m}{0.0}\PYG{+w}{ }\PYG{l+m}{50.0}
\PYG{k}{read\PYGZus{}data}\PYG{+w}{ }\PYG{n}{data.water}\PYG{+w}{ }\PYG{n}{add}\PYG{+w}{ }\PYG{n}{merge}\PYG{+w}{ }\PYG{k}{group }\PYG{n+nv+nvIdentifier}{solvent}
\end{sphinxVerbatim}


\subsection{Description}
\label{\detokenize{read_data:description}}
\sphinxAtStartPar
Read in a data file containing information LAMMPS needs to run a
simulation.  The file can be ASCII text or a gzipped text file
(detected by a .gz suffix).

\sphinxAtStartPar
This is one of 3 ways to specify the simulation box: see the
{\hyperref[\detokenize{create_box::doc}]{\sphinxcrossref{\DUrole{doc}{create\_box}}}} and {\hyperref[\detokenize{read_restart::doc}]{\sphinxcrossref{\DUrole{doc}{read\_restart}}}}
and commands for alternative methods.  It is also one of 3 ways to
specify initial atom coordinates: see the {\hyperref[\detokenize{create_atoms::doc}]{\sphinxcrossref{\DUrole{doc}{create\_atoms}}}} and {\hyperref[\detokenize{read_restart::doc}]{\sphinxcrossref{\DUrole{doc}{read\_restart}}}} and commands
for alternative methods.  Also see the explanation of the
{\hyperref[\detokenize{Run_options::doc}]{\sphinxcrossref{\DUrole{doc}{\sphinxhyphen{}restart command\sphinxhyphen{}line switch}}}} which can convert a
restart file to a data file.

\sphinxAtStartPar
This command can be used multiple times to add new atoms and their
properties to an existing system by using the \sphinxstyleemphasis{add}, \sphinxstyleemphasis{offset}, and
\sphinxstyleemphasis{shift} keywords.  However, it is important to understand that several
system parameters, like the number of types of different kinds and per
atom settings are \sphinxstylestrong{locked in} after the first \sphinxstyleemphasis{read\_data} command,
which means that no type ID (including its offset) may have a larger
value when processing additional data files than what is set by the
first data file and the corresponding \sphinxstyleemphasis{read\_data} command options.  See
more details on this situation below, which includes the use case for
the \sphinxstyleemphasis{extra} keywords.

\sphinxAtStartPar
The \sphinxstyleemphasis{group} keyword adds all the atoms in the data file to the
specified group\sphinxhyphen{}ID.  The group will be created if it does not already
exist.  This is useful if you are reading multiple data files and wish
to put sets of atoms into different groups so they can be operated on
later.  E.g. a group of added atoms can be moved to new positions via
the {\hyperref[\detokenize{displace_atoms::doc}]{\sphinxcrossref{\DUrole{doc}{displace\_atoms}}}} command.  Note that atoms
read from the data file are also always added to the “all” group.  The
{\hyperref[\detokenize{group::doc}]{\sphinxcrossref{\DUrole{doc}{group}}}} command discusses atom groups, as used in LAMMPS.

\sphinxAtStartPar
The \sphinxstyleemphasis{nocoeff} keyword tells read\_data to ignore force field parameters.
The various Coeff sections are still read and have to have the correct
number of lines, but they are not applied. This also allows to read a
data file without having any pair, bond, angle, dihedral or improper
styles defined, or to read a data file for a different force field.

\sphinxAtStartPar
The use of the \sphinxstyleemphasis{fix} keyword is discussed below.


\bigskip\hrule\bigskip



\subsection{Reading multiple data files}
\label{\detokenize{read_data:reading-multiple-data-files}}
\sphinxAtStartPar
The read\_data command can be used multiple times with the same or
different data files to build up a complex system from components
contained in individual data files.  For example one data file could
contain fluid in a confined domain; a second could contain wall atoms,
and the second file could be read a third time to create a wall on the
other side of the fluid.  The third set of atoms could be rotated to
an opposing direction using the {\hyperref[\detokenize{displace_atoms::doc}]{\sphinxcrossref{\DUrole{doc}{displace\_atoms}}}}
command, after the third read\_data command is used.

\sphinxAtStartPar
The \sphinxstyleemphasis{add}, \sphinxstyleemphasis{offset}, \sphinxstyleemphasis{shift}, \sphinxstyleemphasis{extra}, and \sphinxstyleemphasis{group} keywords are
useful in this context.

\sphinxAtStartPar
If a simulation box does not yet exist, the \sphinxstyleemphasis{add} keyword cannot be
used; the read\_data command is being used for the first time.  If a
simulation box does exist, due to using the {\hyperref[\detokenize{create_box::doc}]{\sphinxcrossref{\DUrole{doc}{create\_box}}}} command, or a previous read\_data command, then the \sphinxstyleemphasis{add}
keyword must be used.

\begin{sphinxadmonition}{note}{Note:}
\sphinxAtStartPar
If the first read\_data command defined an orthogonal or restricted
triclinic or general triclinic simulation box (see the sub\sphinxhyphen{}section
below on header keywords), then subsequent data files must define
the same kind of simulation box.  For orthogonal boxes, the new box
can be a different size; see the next Note.  For a restricted
triclinic box, the 3 new tilt factors (“xy xz yz” keyword) must
have the same values as in the original data file.  For a general
triclinic box, the new avec, bvec, cvec, and “abc origin” keywords
must have the same values in the original data file.  files.  Also
the \sphinxstyleemphasis{shift} keyword cannot be used in subsequent read\_data commands
for a general triclinic box.
\end{sphinxadmonition}

\begin{sphinxadmonition}{note}{Note:}
\sphinxAtStartPar
For orthogonal boxes, the simulation box size in the new data file
will be merged with the existing simulation box to create a large
enough box in each dimension to contain both the existing and new
atoms.  Each box dimension never shrinks due to this merge
operation, it only stays the same or grows.  Care must be used if
you are growing the existing simulation box in a periodic
dimension.  If there are existing atoms with bonds that straddle
that periodic boundary, then the atoms may become far apart if the
box size grows.  This will separate the atoms in the bond, which
can lead to “lost” bond atoms or bad dynamics.
\end{sphinxadmonition}

\sphinxAtStartPar
The three choices for the \sphinxstyleemphasis{add} argument affect how the atom IDs and
molecule IDs of atoms in the data file are treated.

\sphinxAtStartPar
If \sphinxstyleemphasis{append} is specified, atoms in the data file are added to the
current system, with their atom IDs reset so that an atom\sphinxhyphen{}ID = M in
the data file becomes atom\sphinxhyphen{}ID = N+M, where N is the largest atom ID in
the current system.  This rule is applied to all occurrences of atom
IDs in the data file, e.g. in the Velocity or Bonds section. This is
also done for molecule IDs, if the atom style does support molecule
IDs or they are enabled via fix property/atom.

\sphinxAtStartPar
If \sphinxstyleemphasis{IDoffset} is specified, then \sphinxstyleemphasis{IDoffset} is a numeric value is
given, e.g. 1000, so that an atom\sphinxhyphen{}ID = M in the data file becomes
atom\sphinxhyphen{}ID = 1000+M. For systems with enabled molecule IDs, another
numerical argument \sphinxstyleemphasis{MOLoffset} is required representing the equivalent
offset for molecule IDs.

\sphinxAtStartPar
If \sphinxstyleemphasis{merge} is specified, the data file atoms are added to the current
system without changing their IDs.  They are assumed to merge (without
duplication) with the currently defined atoms.  It is up to you to
ensure there are no multiply defined atom IDs, as LAMMPS only performs
an incomplete check that this is the case by ensuring the resulting
max atom\sphinxhyphen{}ID \textgreater{}= the number of atoms. For molecule IDs, there is no
check done at all.

\sphinxAtStartPar
The \sphinxstyleemphasis{offset} and \sphinxstyleemphasis{shift} keywords can only be used if the \sphinxstyleemphasis{add}
keyword is also specified.

\sphinxAtStartPar
The \sphinxstyleemphasis{offset} keyword adds the specified offset values to the atom
types, bond types, angle types, dihedral types, and improper types as
they are read from the data file.  E.g. if \sphinxstyleemphasis{toff} = 2, and the file
uses atom types 1,2,3, then the added atoms will have atom types
3,4,5.  These offsets apply to all occurrences of types in the data
file, e.g. for the Atoms or Masses or Pair Coeffs or Bond Coeffs
sections.  This makes it easy to use atoms and molecules and their
attributes from a data file in different simulations, where you want
their types (atom, bond, angle, etc) to be different depending on what
other types already exist.  All five offset values must be specified,
but individual values will be ignored if the data file does not use
that attribute (e.g. no bonds).

\begin{sphinxadmonition}{note}{Note:}
\sphinxAtStartPar
Offsets are \sphinxstylestrong{ignored} on lines using type labels, as the type
labels will determine the actual types directly depending on the
current {\hyperref[\detokenize{labelmap::doc}]{\sphinxcrossref{\DUrole{doc}{labelmap}}}} settings.
\end{sphinxadmonition}

\sphinxAtStartPar
The \sphinxstyleemphasis{shift} keyword can be used to specify an (Sx, Sy, Sz)
displacement applied to the coordinates of each atom.  Sz must be 0.0
for a 2d simulation.  This is a mechanism for adding structured
collections of atoms at different locations within the simulation box,
to build up a complex geometry.  It is up to you to ensure atoms do
not end up overlapping unphysically which would lead to bad dynamics.
Note that the {\hyperref[\detokenize{displace_atoms::doc}]{\sphinxcrossref{\DUrole{doc}{displace\_atoms}}}} command can be used
to move a subset of atoms after they have been read from a data file.
Likewise, the {\hyperref[\detokenize{delete_atoms::doc}]{\sphinxcrossref{\DUrole{doc}{delete\_atoms}}}} command can be used to
remove overlapping atoms.  Note that the shift values (Sx, Sy, Sz) are
also added to the simulation box information (xlo, xhi, ylo, yhi, zlo,
zhi) in the data file to shift its boundaries.  E.g. xlo\_new = xlo +
Sx, xhi\_new = xhi + Sx.

\sphinxAtStartPar
The \sphinxstyleemphasis{extra} keywords can only be used the first time the read\_data
command is used.  They are useful if you intend to add new atom, bond,
angle, etc types later with additional read\_data commands.  This is
because the maximum number of allowed atom, bond, angle, etc types is
set by LAMMPS when the system is first initialized.  If you do not use
the \sphinxstyleemphasis{extra} keywords, then the number of these types will be limited
to what appears in the first data file you read.  For example, if the
first data file is a solid substrate of Si, it will likely specify a
single atom type.  If you read a second data file with a different
material (water molecules) that sit on top of the substrate, you will
want to use different atom types for those atoms.  You can only do
this if you set the \sphinxstyleemphasis{extra/atom/types} keyword to a sufficiently large
value when reading the substrate data file.  Note that use of the
\sphinxstyleemphasis{extra} keywords also allows each data file to contain sections like
Masses or Pair Coeffs or Bond Coeffs which are sized appropriately for
the number of types in that data file.  If the \sphinxstyleemphasis{offset} keyword is
used appropriately when each data file is read, the values in those
sections will be stored correctly in the larger data structures
allocated by the use of the \sphinxstyleemphasis{extra} keywords.  E.g. the substrate file
can list mass and pair coefficients for type 1 silicon atoms.  The
water file can list mass and pair coefficients for type 1 and type 2
hydrogen and oxygen atoms.  Use of the \sphinxstyleemphasis{extra} and \sphinxstyleemphasis{offset} keywords
will store those mass and pair coefficient values appropriately in
data structures that allow for 3 atom types (Si, H, O).  Of course,
you would still need to specify coefficients for H/Si and O/Si
interactions in your input script to have a complete pairwise
interaction model.

\sphinxAtStartPar
An alternative to using the \sphinxstyleemphasis{extra} keywords with the read\_data
command, is to use the {\hyperref[\detokenize{create_box::doc}]{\sphinxcrossref{\DUrole{doc}{create\_box}}}} command to
initialize the simulation box and all the various type limits you need
via its \sphinxstyleemphasis{extra} keywords.  Then use the read\_data command one or more
times to populate the system with atoms, bonds, angles, etc, using the
\sphinxstyleemphasis{offset} keyword if desired to alter types used in the various data
files you read.


\bigskip\hrule\bigskip



\subsection{Format of a data file}
\label{\detokenize{read_data:format-of-a-data-file}}
\sphinxAtStartPar
The structure of the data file is important, though many settings and
sections are optional or can come in any order.  See the examples
directory for sample data files for different problems.

\sphinxAtStartPar
The file will be read line by line, but there is a limit of 254
characters per line and characters beyond that limit will be ignored.

\sphinxAtStartPar
A data file has a header and a body.  The header appears first.  The
first line of the header and thus of the data file is \sphinxstyleemphasis{always} skipped;
it typically contains a description of the file or a comment from the
software that created the file.

\sphinxAtStartPar
Then lines are read one line at a time.  Lines can have a trailing
comment starting with ‘\#’ that is ignored.  There \sphinxstyleemphasis{must} be at least one
blank between any valid content and the comment. If a line is blank
(i.e. contains only white\sphinxhyphen{}space after comments are deleted), it is
skipped.  If the line contains a header keyword, the corresponding
value(s) is/are read from the line.  A line that is \sphinxstyleemphasis{not} blank and does
\sphinxstyleemphasis{not} contain a header keyword begins the body of the file.

\sphinxAtStartPar
The body of the file contains zero or more sections.  The first line of
a section has only a keyword.  This line can have a trailing comment
starting with ‘\#’ that is either ignored or can be used to check for a
style match, as described below.  There must be a blank between the
keyword and any comment. The \sphinxstyleemphasis{next} line is \sphinxstyleemphasis{always} skipped.  The
remaining lines of the section contain values.  The number of lines
depends on the section keyword as described below.  Zero or more blank
lines can be used \sphinxstyleemphasis{between} sections.  Sections can appear in any order,
with a few exceptions as noted below.

\sphinxAtStartPar
The keyword \sphinxstyleemphasis{fix} can be used one or more times.  Each usage specifies
a fix that will be used to process a specific portion of the data
file.  Any header line containing \sphinxstyleemphasis{header\sphinxhyphen{}string} and any section that
is an exact match to \sphinxstyleemphasis{section\sphinxhyphen{}string} will be passed to the specified
fix.  See the {\hyperref[\detokenize{fix_property_atom::doc}]{\sphinxcrossref{\DUrole{doc}{fix property/atom}}}} command for
an example of a fix that operates in this manner.  The doc page for
the fix defines the syntax of the header line(s) and section that it
reads from the data file.  Note that the \sphinxstyleemphasis{header\sphinxhyphen{}string} can be
specified as NULL, in which case no header lines are passed to the
fix.  This means the fix can infer the length of its Section from
standard header settings, such as the number of atoms.  Also the
\sphinxstyleemphasis{section\sphinxhyphen{}string} may be specified as NULL, and in that case the fix
ID is used as section name.

\sphinxAtStartPar
The formatting of individual lines in the data file (indentation,
spacing between words and numbers) is not important except that header
and section keywords (e.g. atoms, xlo xhi, Masses, Bond Coeffs) must
be capitalized as shown and cannot have extra white\sphinxhyphen{}space between
their words \sphinxhyphen{} e.g. two spaces or a tab between the 2 words in “xlo
xhi” or the 2 words in “Bond Coeffs”, is not valid.


\bigskip\hrule\bigskip



\subsection{Format of the header of a data file}
\label{\detokenize{read_data:format-of-the-header-of-a-data-file}}
\sphinxAtStartPar
These are the recognized header keywords.  Header lines can come in
any order.  Each keyword takes a single value unless noted in this
list.  The value(s) are read from the beginning of the line.
Thus the keyword \sphinxstyleemphasis{atoms} should be in a line like “1000 atoms”; the
keyword \sphinxstyleemphasis{ylo yhi} should be in a line like “\sphinxhyphen{}10.0 10.0 ylo yhi”; the
keyword \sphinxstyleemphasis{xy xz yz} should be in a line like “0.0 5.0 6.0 xy xz yz”.

\sphinxAtStartPar
All these settings have a default value of 0, except for the
simulation box size settings; their defaults are explained below.  A
keyword line need only appear if its value is different than the
default.
\begin{itemize}
\item {} 
\sphinxAtStartPar
\sphinxstyleemphasis{atoms} = \# of atoms in system

\item {} 
\sphinxAtStartPar
\sphinxstyleemphasis{bonds} = \# of bonds in system

\item {} 
\sphinxAtStartPar
\sphinxstyleemphasis{angles} = \# of angles in system

\item {} 
\sphinxAtStartPar
\sphinxstyleemphasis{dihedrals} = \# of dihedrals in system

\item {} 
\sphinxAtStartPar
\sphinxstyleemphasis{impropers} = \# of impropers in system

\item {} 
\sphinxAtStartPar
\sphinxstyleemphasis{atom types} = \# of atom types in system

\item {} 
\sphinxAtStartPar
\sphinxstyleemphasis{bond types} = \# of bond types in system

\item {} 
\sphinxAtStartPar
\sphinxstyleemphasis{angle types} = \# of angle types in system

\item {} 
\sphinxAtStartPar
\sphinxstyleemphasis{dihedral types} = \# of dihedral types in system

\item {} 
\sphinxAtStartPar
\sphinxstyleemphasis{improper types} = \# of improper types in system

\item {} 
\sphinxAtStartPar
\sphinxstyleemphasis{extra bond per atom} = leave space for this many new bonds per atom (deprecated, use extra/bond/per/atom keyword)

\item {} 
\sphinxAtStartPar
\sphinxstyleemphasis{extra angle per atom} = leave space for this many new angles per atom (deprecated, use extra/angle/per/atom keyword)

\item {} 
\sphinxAtStartPar
\sphinxstyleemphasis{extra dihedral per atom} = leave space for this many new dihedrals per atom (deprecated, use extra/dihedral/per/atom keyword)

\item {} 
\sphinxAtStartPar
\sphinxstyleemphasis{extra improper per atom} = leave space for this many new impropers per atom (deprecated, use extra/improper/per/atom keyword)

\item {} 
\sphinxAtStartPar
\sphinxstyleemphasis{extra special per atom} = leave space for this many new special bonds per atom (deprecated, use extra/special/per/atom keyword)

\item {} 
\sphinxAtStartPar
\sphinxstyleemphasis{ellipsoids} = \# of ellipsoids in system

\item {} 
\sphinxAtStartPar
\sphinxstyleemphasis{lines} = \# of line segments in system

\item {} 
\sphinxAtStartPar
\sphinxstyleemphasis{triangles} = \# of triangles in system

\item {} 
\sphinxAtStartPar
\sphinxstyleemphasis{bodies} = \# of bodies in system

\item {} 
\sphinxAtStartPar
\sphinxstyleemphasis{xlo xhi} = simulation box boundaries in x dimension (2 values)

\item {} 
\sphinxAtStartPar
\sphinxstyleemphasis{ylo yhi} = simulation box boundaries in y dimension (2 values)

\item {} 
\sphinxAtStartPar
\sphinxstyleemphasis{zlo zhi} = simulation box boundaries in z dimension (2 values)

\item {} 
\sphinxAtStartPar
\sphinxstyleemphasis{xy xz yz} = simulation box tilt factors for triclinic system (3 values)

\item {} 
\sphinxAtStartPar
\sphinxstyleemphasis{avec} = first edge vector of a general triclinic simulation box (3 values)

\item {} 
\sphinxAtStartPar
\sphinxstyleemphasis{bvec} = second edge vector of a general triclinic simulation box (3 values)

\item {} 
\sphinxAtStartPar
\sphinxstyleemphasis{cvec} = third edge vector of a general triclinic simulation box (3 values)

\item {} 
\sphinxAtStartPar
\sphinxstyleemphasis{abc origin} = origin of a general triclinic simulation box (3 values)

\end{itemize}


\bigskip\hrule\bigskip



\subsection{Header specification of the simulation box size and shape}
\label{\detokenize{read_data:header-specification-of-the-simulation-box-size-and-shape}}
\sphinxAtStartPar
The last 8 keywords in the list of header keywords are for simulation
boxes of 3 kinds which LAMMPS supports:
\begin{itemize}
\item {} 
\sphinxAtStartPar
orthogonal box = faces are perpendicular to the xyz coordinate axes

\item {} 
\sphinxAtStartPar
restricted triclinic box = a parallelepiped defined by 3 edge vectors oriented in a constrained manner

\item {} 
\sphinxAtStartPar
general triclinic box = a parallelepiped defined by 3 arbitrary edge vectors

\end{itemize}

\sphinxAtStartPar
For restricted and general triclinic boxes, see the
{\hyperref[\detokenize{Howto_triclinic::doc}]{\sphinxcrossref{\DUrole{doc}{Howto\_triclinic}}}} doc page for a fuller
description than is given here.

\sphinxAtStartPar
The units of the values for all 8 keywords in in distance units; see
the {\hyperref[\detokenize{units::doc}]{\sphinxcrossref{\DUrole{doc}{units}}}} command for details.

\sphinxAtStartPar
For all 3 kinds of simulation boxes, the system may be periodic or
non\sphinxhyphen{}periodic in any dimension; see the {\hyperref[\detokenize{boundary::doc}]{\sphinxcrossref{\DUrole{doc}{boundary}}}}
command for details.

\sphinxAtStartPar
When the simulation box is created by the read\_data command, it is
also partitioned into a regular 3d grid of subdomains, one per
processor, based on the number of processors being used and the
settings of the {\hyperref[\detokenize{processors::doc}]{\sphinxcrossref{\DUrole{doc}{processors}}}} command.  For each kind
of simulation box the subdomains have the same shape as the simulation
box, i.e. smaller orthogonal bricks for orthogonal boxes, smaller
parallelepipeds for triclinic boxes.  The partitioning can later be
changed by the {\hyperref[\detokenize{balance::doc}]{\sphinxcrossref{\DUrole{doc}{balance}}}} or {\hyperref[\detokenize{fix_balance::doc}]{\sphinxcrossref{\DUrole{doc}{fix balance}}}} commands.

\sphinxAtStartPar
For an orthogonal box, only the \sphinxstyleemphasis{xlo xhi}, \sphinxstyleemphasis{ylo yhi}, \sphinxstyleemphasis{zlo zhi}
keywords are used.  They define the extent of the simulation box in
each dimension so that the resulting edge vectors of an orthogonal box
are:
\begin{itemize}
\item {} 
\sphinxAtStartPar
\sphinxstylestrong{A} = (xhi\sphinxhyphen{}xlo,0,0)

\item {} 
\sphinxAtStartPar
\sphinxstylestrong{B} = (0,yhi\sphinxhyphen{}ylo,0)

\item {} 
\sphinxAtStartPar
\sphinxstylestrong{C} = (0,0,zhi\sphinxhyphen{}zlo)

\end{itemize}

\sphinxAtStartPar
The origin (lower left corner) of the orthogonal box is at
(xlo,ylo,zlo).  The default values for these 3 keywords are \sphinxhyphen{}0.5 and
0.5 for each lo/hi pair.  For a 2d simulation, the zlo and zhi values
must straddle zero.  The default zlo/zhi values do this, so that
keyword is not needed in 2d.

\sphinxAtStartPar
For a restricted triclinic box, the \sphinxstyleemphasis{xy xz yz} keyword is used in
addition to the \sphinxstyleemphasis{xlo xhi}, \sphinxstyleemphasis{ylo yhi}, \sphinxstyleemphasis{zlo zhi} keywords.  The three
\sphinxstyleemphasis{xy,xz,yz} values can be 0.0 or positive or negative, and are called
“tilt factors” because they are the amount of displacement applied to
edges of faces of an orthogonal box to transform it into a restricted
triclinic parallelepiped.

\sphinxAtStartPar
The {\hyperref[\detokenize{Howto_triclinic::doc}]{\sphinxcrossref{\DUrole{doc}{Howto\_triclinic}}}} doc page discusses the
tilt factors in detail and explains that the resulting edge vectors of
a restricted triclinic box are:
\begin{itemize}
\item {} 
\sphinxAtStartPar
\sphinxstylestrong{A} = (xhi\sphinxhyphen{}xlo,0,0)

\item {} 
\sphinxAtStartPar
\sphinxstylestrong{B} = (xy,yhi\sphinxhyphen{}ylo,0)

\item {} 
\sphinxAtStartPar
\sphinxstylestrong{C} = (xz,yz,zhi\sphinxhyphen{}zlo)

\end{itemize}

\sphinxAtStartPar
This restricted form of edge vectors requires that \sphinxstylestrong{A} be in the
direction of the x\sphinxhyphen{}axis, \sphinxstylestrong{B} be in the xy plane with its y\sphinxhyphen{}component
in the +y direction, and \sphinxstylestrong{C} have its z\sphinxhyphen{}component in the +z
direction.  The origin (lower left corner) of the restricted triclinic
box is at (xlo,ylo,zlo).

\sphinxAtStartPar
For a 2d simulation, the zlo and zhi values must straddle zero.  The
default zlo/zhi values do this, so that keyword is not needed in 2d.
The xz and yz values must also be zero in 2d.  The shape of the 2d
restricted triclinic simulation box is effectively a parallelogram.

\begin{sphinxadmonition}{note}{Note:}
\sphinxAtStartPar
When a restricted triclinic box is used, the simulation domain
should normally be periodic in any dimensions that tilt is applied
to, which is given by the second dimension of the tilt factor
(e.g. y for xy tilt).  This is so that pairs of atoms interacting
across that boundary will have one of them shifted by the tilt
factor.  Periodicity is set by the {\hyperref[\detokenize{boundary::doc}]{\sphinxcrossref{\DUrole{doc}{boundary}}}}
command which also describes the shifting by the tilt factor.  For
example, if the xy tilt factor is non\sphinxhyphen{}zero, then the y dimension
should be periodic.  Similarly, the z dimension should be periodic
if xz or yz is non\sphinxhyphen{}zero.  LAMMPS does not require this periodicity,
but you may lose atoms if this is not the case.
\end{sphinxadmonition}

\begin{sphinxadmonition}{note}{Note:}
\sphinxAtStartPar
Normally, the specified tilt factors (xy,xz,yz) should not skew the
simulation box by more than half the distance of the corresponding
parallel box length for computational efficiency.  For example, if
\(x_\text{lo} = 2\) and \(x_\text{hi} = 12\), then the
\(x\) box length is 10 and the \(xy\) tilt factor should be
between \(-5\) and \(5\).  LAMMPS will issue a warning if
this is not the case.  See the last sub\sphinxhyphen{}section of the
{\hyperref[\detokenize{Howto_triclinic::doc}]{\sphinxcrossref{\DUrole{doc}{Howto\_triclinic}}}} doc page for more details.
\end{sphinxadmonition}

\begin{sphinxadmonition}{note}{Note:}
\sphinxAtStartPar
If a simulation box is initially orthogonal, but will tilt during a
simulation, e.g. via the {\hyperref[\detokenize{fix_deform::doc}]{\sphinxcrossref{\DUrole{doc}{fix deform}}}} command,
then the box should be defined as restricted triclinic with all 3
tilt factors = 0.0.  Alternatively, the {\hyperref[\detokenize{change_box::doc}]{\sphinxcrossref{\DUrole{doc}{change box}}}} command can be used to convert an orthogonal box to a
restricted triclinic box.
\end{sphinxadmonition}

\sphinxAtStartPar
For a general triclinic box, the \sphinxstyleemphasis{avec}, \sphinxstyleemphasis{bvec}, \sphinxstyleemphasis{cvec}, and \sphinxstyleemphasis{abc
origin} keywords are used.  The \sphinxstyleemphasis{xlo xhi}, \sphinxstyleemphasis{ylo yhi}, \sphinxstyleemphasis{zlo zhi}, and
\sphinxstyleemphasis{xy xz yz} keywords are NOT used.  The first 3 keywords define the 3
edge vectors \sphinxstylestrong{A}, \sphinxstylestrong{B}, \sphinxstylestrong{C} of the general triclinic box.  They
can be arbitrary vectors so long as they are distinct, non\sphinxhyphen{}zero, and
not co\sphinxhyphen{}planar.  They must also define a right\sphinxhyphen{}handed system such that
(\sphinxstylestrong{A} x \sphinxstylestrong{B}) points in the direction of \sphinxstylestrong{C}.  Note that a
left\sphinxhyphen{}handed system can be converted to a right\sphinxhyphen{}handed system by simply
swapping the order of any pair of the \sphinxstylestrong{A}, \sphinxstylestrong{B}, \sphinxstylestrong{C} vectors.
The origin of the box (origin of the 3 edge vectors) is set by the
\sphinxstyleemphasis{abc origin} keyword.

\sphinxAtStartPar
The default values for these 4 keywords are as follows:
\begin{itemize}
\item {} 
\sphinxAtStartPar
avec = (1,0,0)

\item {} 
\sphinxAtStartPar
bvec = (0,1,0)

\item {} 
\sphinxAtStartPar
cvec = (0,0,1)

\item {} 
\sphinxAtStartPar
abc origin = (0,0,0) for 3d, (0,0,\sphinxhyphen{}0.5) for 2d

\end{itemize}

\sphinxAtStartPar
For 2d simulations, \sphinxstyleemphasis{cvec} = (0,0,1) is required, and the 3rd value of
\sphinxstyleemphasis{abc origin} must be \sphinxhyphen{}0.5.  These are the default values, so the
\sphinxstyleemphasis{cvec} keyword is not needed in 2d.

\begin{sphinxadmonition}{note}{Note:}
\sphinxAtStartPar
LAMMPS allows specification of general triclinic simulation boxes
as a convenience for users who may be converting data from
solid\sphinxhyphen{}state crystallographic representations or from DFT codes for
input to LAMMPS.  However, as explained on the
{\hyperref[\detokenize{Howto_triclinic::doc}]{\sphinxcrossref{\DUrole{doc}{Howto\_triclinic}}}} doc page, internally,
LAMMPS only uses restricted triclinic simulation boxes.  This means
the box and per\sphinxhyphen{}atom information (e.g. coordinates, velocities) in
the data file are converted (rotated) from general to restricted
triclinic form when the file is read.  Other sections of the data
file must also list their per\sphinxhyphen{}atom data appropriately if vector
quantities are specified. This requirement is explained below for
the relevant sections.  The {\hyperref[\detokenize{Howto_triclinic::doc}]{\sphinxcrossref{\DUrole{doc}{Howto\_triclinic}}}} doc page also discusses other LAMMPS commands
which can input/output general triclinic representations of the
simulation box and per\sphinxhyphen{}atom data.
\end{sphinxadmonition}

\sphinxAtStartPar
The following explanations apply to all 3 kinds of simulation boxes:
orthogonal, restricted triclinic, and general triclinic.

\sphinxAtStartPar
If the system is periodic (in a dimension), then atom coordinates can
be outside the bounds (in that dimension); they will be remapped (in a
periodic sense) back inside the box.  For triclinic boxes, periodicity
in x,y,z refers to the faces of the parallelepiped defined by the
\sphinxstylestrong{A},**B**,**C** edge vectors of the simulation box.  See the
{\hyperref[\detokenize{boundary::doc}]{\sphinxcrossref{\DUrole{doc}{boundary}}}} command doc page for a fuller discussion.

\sphinxAtStartPar
Note that if the \sphinxstyleemphasis{add} option is being used to add atoms to a
simulation box that already exists, this periodic remapping will be
performed using simulation box bounds that are the union of the
existing box and the box boundaries in the new data file.

\sphinxAtStartPar
If the system is non\sphinxhyphen{}periodic (in a dimension), then an image flag for
that direction has no meaning, since there cannot be periodic images
without periodicity and the data file is therefore \sphinxhyphen{} technically speaking
\sphinxhyphen{} invalid.  This situation would happen when a data file was written
with periodic boundaries and then read back for non\sphinxhyphen{}periodic boundaries.
Accepting a non\sphinxhyphen{}zero image flag can lead to unexpected results for any
operations and computations in LAMMPS that internally use unwrapped
coordinates (for example computing the center of mass of a group of
atoms). Thus all non\sphinxhyphen{}zero image flags for non\sphinxhyphen{}periodic dimensions will
be be reset to zero on reading the data file and LAMMPS will print a
warning message, if that happens.  This is equivalent to wrapping atoms
individually back into the principal unit cell in that direction.  This
operation is equivalent to the behavior of the {\hyperref[\detokenize{change_box::doc}]{\sphinxcrossref{\DUrole{doc}{change\_box command}}}} when used to change periodicity.

\sphinxAtStartPar
If those atoms with non\sphinxhyphen{}zero image flags are involved in bonded
interactions, this reset can lead to undesired changes, when the image
flag values differ between the atoms, i.e. the bonded interaction
straddles domain boundaries.  For example a bond can become stretched
across the unit cell if one of its atoms is wrapped to one side of the
cell and the second atom to the other. In those cases the data file
needs to be pre\sphinxhyphen{}processed externally to become valid again.  This can be
done by first unwrapping coordinates and then wrapping entire molecules
instead of individual atoms back into the principal simulation cell and
finally expanding the cell dimensions in the non\sphinxhyphen{}periodic direction as
needed, so that the image flag would be zero.

\begin{sphinxadmonition}{note}{Note:}
\sphinxAtStartPar
If the system is non\sphinxhyphen{}periodic (in a dimension), then all atoms in
the data file must have coordinates (in that dimension) that are
“greater than or equal to” the lo value and “less than or equal to”
the hi value.  If the non\sphinxhyphen{}periodic dimension is of style “fixed”
(see the {\hyperref[\detokenize{boundary::doc}]{\sphinxcrossref{\DUrole{doc}{boundary}}}} command), then the atom coords
must be strictly “less than” the hi value, due to the way LAMMPS
assign atoms to processors.  Note that you should not make the
lo/hi values radically smaller/larger than the extent of the atoms.
For example, if atoms extend from 0 to 50, you should not specify
the box bounds as \sphinxhyphen{}10000 and 10000 unless you also use the
{\hyperref[\detokenize{processors::doc}]{\sphinxcrossref{\DUrole{doc}{processors command}}}}.  This is because LAMMPS
uses the specified box size to layout the 3d grid of processors.  A
huge (mostly empty) box will be sub\sphinxhyphen{}optimal for performance when
using “fixed” boundary conditions (see the {\hyperref[\detokenize{boundary::doc}]{\sphinxcrossref{\DUrole{doc}{boundary}}}} command).  When using “shrink\sphinxhyphen{}wrap” boundary conditions
(see the {\hyperref[\detokenize{boundary::doc}]{\sphinxcrossref{\DUrole{doc}{boundary}}}} command), a huge (mostly empty)
box may cause a parallel simulation to lose atoms when LAMMPS
shrink\sphinxhyphen{}wraps the box around the atoms.  The read\_data command will
generate an error in this case.
\end{sphinxadmonition}


\bigskip\hrule\bigskip



\subsection{Meaning of other header keywords}
\label{\detokenize{read_data:meaning-of-other-header-keywords}}
\sphinxAtStartPar
The “extra bond per atom” setting (angle, dihedral, improper) is only
needed if new bonds (angles, dihedrals, impropers) will be added to
the system when a simulation runs, e.g. by using the {\hyperref[\detokenize{fix_bond_create::doc}]{\sphinxcrossref{\DUrole{doc}{fix
bond/create}}}} command. Using this header flag is
deprecated; please use the \sphinxstyleemphasis{extra/bond/per/atom} keyword (and
correspondingly for angles, dihedrals and impropers) in the read\_data
command instead. Either will pre\sphinxhyphen{}allocate space in LAMMPS data
structures for storing the new bonds (angles, dihedrals, impropers).

\sphinxAtStartPar
The “extra special per atom” setting is typically only needed if new
bonds/angles/etc will be added to the system, e.g. by using the
{\hyperref[\detokenize{fix_bond_create::doc}]{\sphinxcrossref{\DUrole{doc}{fix bond/create}}}} command.  Or if entire new
molecules will be added to the system, e.g. by using the {\hyperref[\detokenize{fix_deposit::doc}]{\sphinxcrossref{\DUrole{doc}{fix
deposit}}}} or {\hyperref[\detokenize{fix_pour::doc}]{\sphinxcrossref{\DUrole{doc}{fix pour}}}} commands, which
will have more special 1\sphinxhyphen{}2,1\sphinxhyphen{}3,1\sphinxhyphen{}4 neighbors than any other molecules
defined in the data file.  Using this header flag is deprecated;
please use the \sphinxstyleemphasis{extra/special/per/atom} keyword instead.  Using this
setting will pre\sphinxhyphen{}allocate space in the LAMMPS data structures for
storing these neighbors.  See the {\hyperref[\detokenize{special_bonds::doc}]{\sphinxcrossref{\DUrole{doc}{special\_bonds}}}}
and {\hyperref[\detokenize{molecule::doc}]{\sphinxcrossref{\DUrole{doc}{molecule}}}} doc pages for more discussion of
1\sphinxhyphen{}2,1\sphinxhyphen{}3,1\sphinxhyphen{}4 neighbors.

\begin{sphinxadmonition}{note}{Note:}
\sphinxAtStartPar
All of the “extra” settings are only applied in the first data file
read and when no simulation box has yet been created; as soon as
the simulation box is created (and read\_data implies that), these
settings are \sphinxstyleemphasis{locked} and cannot be changed anymore. Please see the
description of the \sphinxstyleemphasis{add} keyword above for reading multiple data
files.  If they appear in later data files, they are ignored.
\end{sphinxadmonition}

\sphinxAtStartPar
The “ellipsoids” and “lines” and “triangles” and “bodies” settings are
only used with {\hyperref[\detokenize{atom_style::doc}]{\sphinxcrossref{\DUrole{doc}{atom\_style ellipsoid or line or tri or body}}}} and specify how many of the atoms are finite\sphinxhyphen{}size
ellipsoids or lines or triangles or bodies; the remainder are point
particles.  See the discussion of ellipsoidflag and the \sphinxstyleemphasis{Ellipsoids}
section below.  See the discussion of lineflag and the \sphinxstyleemphasis{Lines} section
below.  See the discussion of triangleflag and the \sphinxstyleemphasis{Triangles} section
below.  See the discussion of bodyflag and the \sphinxstyleemphasis{Bodies} section below.

\begin{sphinxadmonition}{note}{Note:}
\sphinxAtStartPar
For {\hyperref[\detokenize{atom_style::doc}]{\sphinxcrossref{\DUrole{doc}{atom\_style template}}}}, the molecular topology
(bonds,angles,etc) is contained in the molecule templates read\sphinxhyphen{}in
by the {\hyperref[\detokenize{molecule::doc}]{\sphinxcrossref{\DUrole{doc}{molecule}}}} command.  This means you cannot
set the \sphinxstyleemphasis{bonds}, \sphinxstyleemphasis{angles}, etc header keywords in the data file,
nor can you define \sphinxstyleemphasis{Bonds}, \sphinxstyleemphasis{Angles}, etc sections as discussed
below.  You can set the \sphinxstyleemphasis{bond types}, \sphinxstyleemphasis{angle types}, etc header
keywords, though it is not necessary.  If specified, they must
match the maximum values defined in any of the template molecules.
\end{sphinxadmonition}


\bigskip\hrule\bigskip



\subsection{Format of the body of a data file}
\label{\detokenize{read_data:format-of-the-body-of-a-data-file}}
\sphinxAtStartPar
These are the section keywords for the body of the file.
\begin{itemize}
\item {} 
\sphinxAtStartPar
\sphinxstyleemphasis{Atoms, Velocities, Masses, Ellipsoids, Lines, Triangles, Bodies} = atom\sphinxhyphen{}property sections

\item {} 
\sphinxAtStartPar
\sphinxstyleemphasis{Bonds, Angles, Dihedrals, Impropers} = molecular topology sections

\item {} 
\sphinxAtStartPar
\sphinxstyleemphasis{Atom Type Labels, Bond Type Labels, Angle Type Labels, Dihedral Type Labels, Improper Type Labels} = type label maps

\item {} 
\sphinxAtStartPar
\sphinxstyleemphasis{Pair Coeffs, PairIJ Coeffs, Bond Coeffs, Angle Coeffs, Dihedral Coeffs,    Improper Coeffs} = force field sections

\item {} 
\sphinxAtStartPar
\sphinxstyleemphasis{BondBond Coeffs, BondAngle Coeffs, MiddleBondTorsion Coeffs,    EndBondTorsion Coeffs, AngleTorsion Coeffs, AngleAngleTorsion Coeffs,    BondBond13 Coeffs, AngleAngle Coeffs} = class 2 force field sections

\end{itemize}

\sphinxAtStartPar
These keywords will check an appended comment for a match with the
currently defined style:
\begin{itemize}
\item {} 
\sphinxAtStartPar
\sphinxstyleemphasis{Atoms, Pair Coeffs, PairIJ Coeffs, Bond Coeffs, Angle Coeffs, Dihedral Coeffs, Improper Coeffs}

\end{itemize}

\sphinxAtStartPar
For example, these lines:

\begin{sphinxVerbatim}[commandchars=\\\{\}]
Atoms \PYGZsh{} sphere
Pair Coeffs \PYGZsh{} lj/cut
\end{sphinxVerbatim}

\sphinxAtStartPar
will check if the currently\sphinxhyphen{}defined {\hyperref[\detokenize{atom_style::doc}]{\sphinxcrossref{\DUrole{doc}{atom\_style}}}} is
\sphinxstyleemphasis{sphere}, and the current {\hyperref[\detokenize{pair_style::doc}]{\sphinxcrossref{\DUrole{doc}{pair\_style}}}} is \sphinxstyleemphasis{lj/cut}.
If not, LAMMPS will issue a warning to indicate that the data file
section likely does not contain the correct number or type of
parameters expected for the currently\sphinxhyphen{}defined style.

\sphinxAtStartPar
Each section is listed below in alphabetic order.  The format of each
section is described including the number of lines it must contain and
rules (if any) for where it can appear in the data file.

\sphinxAtStartPar
Any individual line in the various sections can have a trailing
comment starting with “\#” for annotation purposes. There must be at least
one blank between valid content and the comment. E.g. in the
Atoms section:

\begin{sphinxVerbatim}[commandchars=\\\{\}]
10 1 17 \PYGZhy{}1.0 10.0 5.0 6.0   \PYGZsh{} salt ion
\end{sphinxVerbatim}


\bigskip\hrule\bigskip


\sphinxAtStartPar
\sphinxstyleemphasis{Angle Coeffs} section:
\begin{itemize}
\item {} 
\sphinxAtStartPar
one line per angle type

\item {} 
\sphinxAtStartPar
line syntax: ID coeffs

\begin{sphinxVerbatim}[commandchars=\\\{\}]
ID = angle type (1\PYGZhy{}N)
coeffs = list of coeffs
\end{sphinxVerbatim}

\item {} 
\sphinxAtStartPar
example:

\begin{sphinxVerbatim}[commandchars=\\\{\}]
6 70 108.5 0 0
\end{sphinxVerbatim}

\end{itemize}

\sphinxAtStartPar
The number and meaning of the coefficients are specific to the defined
angle style.  See the {\hyperref[\detokenize{angle_style::doc}]{\sphinxcrossref{\DUrole{doc}{angle\_style}}}} and
{\hyperref[\detokenize{angle_coeff::doc}]{\sphinxcrossref{\DUrole{doc}{angle\_coeff}}}} commands for details.  Coefficients can
also be set via the {\hyperref[\detokenize{angle_coeff::doc}]{\sphinxcrossref{\DUrole{doc}{angle\_coeff}}}} command in the
input script.


\bigskip\hrule\bigskip


\sphinxAtStartPar
\sphinxstyleemphasis{Angle Type Labels} section:
\begin{itemize}
\item {} 
\sphinxAtStartPar
one line per angle type

\item {} 
\sphinxAtStartPar
line syntax: ID label

\begin{sphinxVerbatim}[commandchars=\\\{\}]
ID = angle type (1\PYGZhy{}N)
label = alphanumeric type label
\end{sphinxVerbatim}

\end{itemize}

\sphinxAtStartPar
Define alphanumeric type labels for each numeric angle type.  These
can be used in the Angles section in place of a numeric type, but only
if the this section appears before the Angles section.

\sphinxAtStartPar
See the {\hyperref[\detokenize{Howto_type_labels::doc}]{\sphinxcrossref{\DUrole{doc}{Howto type labels}}}} doc page for the
allowed syntax of type labels and a general discussion of how type
labels can be used.


\bigskip\hrule\bigskip


\sphinxAtStartPar
\sphinxstyleemphasis{AngleAngle Coeffs} section:
\begin{itemize}
\item {} 
\sphinxAtStartPar
one line per improper type

\item {} 
\sphinxAtStartPar
line syntax: ID coeffs
\begin{sphinxalltt}
ID = improper type (1\sphinxhyphen{}N)
coeffs = list of coeffs (see {\hyperref[\detokenize{improper_coeff::doc}]{\sphinxcrossref{\DUrole{doc}{improper\_coeff}}}})
\end{sphinxalltt}

\end{itemize}


\bigskip\hrule\bigskip


\sphinxAtStartPar
\sphinxstyleemphasis{AngleAngleTorsion Coeffs} section:
\begin{itemize}
\item {} 
\sphinxAtStartPar
one line per dihedral type

\item {} 
\sphinxAtStartPar
line syntax: ID coeffs
\begin{sphinxalltt}
ID = dihedral type (1\sphinxhyphen{}N)
coeffs = list of coeffs (see {\hyperref[\detokenize{dihedral_coeff::doc}]{\sphinxcrossref{\DUrole{doc}{dihedral\_coeff}}}})
\end{sphinxalltt}

\end{itemize}


\bigskip\hrule\bigskip


\sphinxAtStartPar
\sphinxstyleemphasis{Angles} section:
\begin{itemize}
\item {} 
\sphinxAtStartPar
one line per angle

\item {} 
\sphinxAtStartPar
line syntax: ID type atom1 atom2 atom3

\begin{sphinxVerbatim}[commandchars=\\\{\}]
ID = number of angle (1\PYGZhy{}Nangles)
type = angle type (1\PYGZhy{}Nangletype, or type label)
atom1,atom2,atom3 = IDs of 1st,2nd,3rd atom in angle
\end{sphinxVerbatim}

\end{itemize}

\sphinxAtStartPar
example:
\begin{quote}

\begin{sphinxVerbatim}[commandchars=\\\{\}]
2 2 17 29 430
\end{sphinxVerbatim}
\end{quote}

\sphinxAtStartPar
The three atoms are ordered linearly within the angle.  Thus the central
atom (around which the angle is computed) is the atom2 in the list.
E.g. H,O,H for a water molecule.  The \sphinxstyleemphasis{Angles} section must appear
after the \sphinxstyleemphasis{Atoms} section.

\sphinxAtStartPar
All values in this section must be integers (1, not 1.0).  However,
the type can be a numeric value or an alphanumeric label.  The latter
is only allowed if the type label has been defined by the
{\hyperref[\detokenize{labelmap::doc}]{\sphinxcrossref{\DUrole{doc}{labelmap}}}} command or an Angle Type Labels section
earlier in the data file.  See the {\hyperref[\detokenize{Howto_type_labels::doc}]{\sphinxcrossref{\DUrole{doc}{Howto type labels}}}} doc page for the allowed syntax of type labels
and a general discussion of how type labels can be used.


\bigskip\hrule\bigskip


\sphinxAtStartPar
\sphinxstyleemphasis{AngleTorsion Coeffs} section:
\begin{itemize}
\item {} 
\sphinxAtStartPar
one line per dihedral type

\item {} 
\sphinxAtStartPar
line syntax: ID coeffs
\begin{sphinxalltt}
ID = dihedral type (1\sphinxhyphen{}N)
coeffs = list of coeffs (see {\hyperref[\detokenize{dihedral_coeff::doc}]{\sphinxcrossref{\DUrole{doc}{dihedral\_coeff}}}})
\end{sphinxalltt}

\end{itemize}


\bigskip\hrule\bigskip


\sphinxAtStartPar
\sphinxstyleemphasis{Atom Type Labels} section:
\begin{itemize}
\item {} 
\sphinxAtStartPar
one line per atom type

\item {} 
\sphinxAtStartPar
line syntax: ID label

\begin{sphinxVerbatim}[commandchars=\\\{\}]
ID = numeric atom type (1\PYGZhy{}N)
label = alphanumeric type label
\end{sphinxVerbatim}

\end{itemize}

\sphinxAtStartPar
Define alphanumeric type labels for each numeric atom type.  These
can be used in the Atoms section in place of a numeric type, but only
if the Atom Type Labels section appears before the Atoms section.

\sphinxAtStartPar
See the {\hyperref[\detokenize{Howto_type_labels::doc}]{\sphinxcrossref{\DUrole{doc}{Howto type labels}}}} doc page for the
allowed syntax of type labels and a general discussion of how type
labels can be used.


\bigskip\hrule\bigskip


\sphinxAtStartPar
\sphinxstyleemphasis{Atoms} section:
\begin{itemize}
\item {} 
\sphinxAtStartPar
one line per atom

\item {} 
\sphinxAtStartPar
line syntax: depends on atom style

\end{itemize}

\sphinxAtStartPar
An \sphinxstyleemphasis{Atoms} section must appear in the data file if natoms \textgreater{} 0 in the
header section.  The atoms can be listed in any order.  These are the
line formats for each {\hyperref[\detokenize{atom_style::doc}]{\sphinxcrossref{\DUrole{doc}{atom style}}}} in LAMMPS.  As
discussed below, each line can optionally have 3 flags (nx,ny,nz)
appended to it, which indicate which image of a periodic simulation
box the atom is in.  These may be important to include for some kinds
of analysis.

\begin{sphinxadmonition}{note}{Note:}
\sphinxAtStartPar
For orthogonal and restricted and general triclinic simulation
boxes, the atom coordinates (x,y,z) listed in this section should
be inside the corresponding simulation box.  For restricted
triclinic boxes that means the parallelepiped defined by the \sphinxstyleemphasis{xlo
xhi}, \sphinxstyleemphasis{ylo yhi}, \sphinxstyleemphasis{zlo zhi}, and \sphinxstyleemphasis{xy xz yz}, keywords.  For general
triclinic boxes that means the parallelepiped defined by the 3 edge
vectors and origin specified by the \sphinxstyleemphasis{avec}, \sphinxstyleemphasis{bvec}, \sphinxstyleemphasis{cvec}, and
\sphinxstyleemphasis{abc origin} header keywords.  See the discussion in the header
section above about how atom coordinates outside the simulation box
are (or are not) remapped to be inside the box.
\end{sphinxadmonition}


\begin{savenotes}\sphinxattablestart
\sphinxthistablewithglobalstyle
\centering
\begin{tabulary}{\linewidth}[t]{TT}
\sphinxtoprule
\sphinxtableatstartofbodyhook
\sphinxAtStartPar
angle
&
\sphinxAtStartPar
atom\sphinxhyphen{}ID molecule\sphinxhyphen{}ID atom\sphinxhyphen{}type x y z
\\
\sphinxhline
\sphinxAtStartPar
atomic
&
\sphinxAtStartPar
atom\sphinxhyphen{}ID atom\sphinxhyphen{}type x y z
\\
\sphinxhline
\sphinxAtStartPar
body
&
\sphinxAtStartPar
atom\sphinxhyphen{}ID atom\sphinxhyphen{}type bodyflag mass x y z
\\
\sphinxhline
\sphinxAtStartPar
bond
&
\sphinxAtStartPar
atom\sphinxhyphen{}ID molecule\sphinxhyphen{}ID atom\sphinxhyphen{}type x y z
\\
\sphinxhline
\sphinxAtStartPar
bpm/sphere
&
\sphinxAtStartPar
atom\sphinxhyphen{}ID molecule\sphinxhyphen{}ID atom\sphinxhyphen{}type diameter density x y z
\\
\sphinxhline
\sphinxAtStartPar
charge
&
\sphinxAtStartPar
atom\sphinxhyphen{}ID atom\sphinxhyphen{}type q x y z
\\
\sphinxhline
\sphinxAtStartPar
dielectric
&
\sphinxAtStartPar
atom\sphinxhyphen{}ID atom\sphinxhyphen{}type q x y z mux muy muz area ed em epsilon curvature
\\
\sphinxhline
\sphinxAtStartPar
dipole
&
\sphinxAtStartPar
atom\sphinxhyphen{}ID atom\sphinxhyphen{}type q x y z mux muy muz
\\
\sphinxhline
\sphinxAtStartPar
dpd
&
\sphinxAtStartPar
atom\sphinxhyphen{}ID atom\sphinxhyphen{}type theta x y z
\\
\sphinxhline
\sphinxAtStartPar
edpd
&
\sphinxAtStartPar
atom\sphinxhyphen{}ID atom\sphinxhyphen{}type edpd\_temp edpd\_cv x y z
\\
\sphinxhline
\sphinxAtStartPar
electron
&
\sphinxAtStartPar
atom\sphinxhyphen{}ID atom\sphinxhyphen{}type q espin eradius x y z
\\
\sphinxhline
\sphinxAtStartPar
ellipsoid
&
\sphinxAtStartPar
atom\sphinxhyphen{}ID atom\sphinxhyphen{}type ellipsoidflag density x y z
\\
\sphinxhline
\sphinxAtStartPar
full
&
\sphinxAtStartPar
atom\sphinxhyphen{}ID molecule\sphinxhyphen{}ID atom\sphinxhyphen{}type q x y z
\\
\sphinxhline
\sphinxAtStartPar
line
&
\sphinxAtStartPar
atom\sphinxhyphen{}ID molecule\sphinxhyphen{}ID atom\sphinxhyphen{}type lineflag density x y z
\\
\sphinxhline
\sphinxAtStartPar
mdpd
&
\sphinxAtStartPar
atom\sphinxhyphen{}ID atom\sphinxhyphen{}type rho x y z
\\
\sphinxhline
\sphinxAtStartPar
molecular
&
\sphinxAtStartPar
atom\sphinxhyphen{}ID molecule\sphinxhyphen{}ID atom\sphinxhyphen{}type x y z
\\
\sphinxhline
\sphinxAtStartPar
peri
&
\sphinxAtStartPar
atom\sphinxhyphen{}ID atom\sphinxhyphen{}type volume density x y z
\\
\sphinxhline
\sphinxAtStartPar
rheo
&
\sphinxAtStartPar
atom\sphinxhyphen{}ID atom\sphinxhyphen{}type status rho x y z
\\
\sphinxhline
\sphinxAtStartPar
rheo/thermal
&
\sphinxAtStartPar
atom\sphinxhyphen{}ID atom\sphinxhyphen{}type status rho energy x y z
\\
\sphinxhline
\sphinxAtStartPar
smd
&
\sphinxAtStartPar
atom\sphinxhyphen{}ID atom\sphinxhyphen{}type molecule volume mass kradius cradius x0 y0 z0 x y z
\\
\sphinxhline
\sphinxAtStartPar
sph
&
\sphinxAtStartPar
atom\sphinxhyphen{}ID atom\sphinxhyphen{}type rho esph cv x y z
\\
\sphinxhline
\sphinxAtStartPar
sphere
&
\sphinxAtStartPar
atom\sphinxhyphen{}ID atom\sphinxhyphen{}type diameter density x y z
\\
\sphinxhline
\sphinxAtStartPar
spin
&
\sphinxAtStartPar
atom\sphinxhyphen{}ID atom\sphinxhyphen{}type x y z spx spy spz sp
\\
\sphinxhline
\sphinxAtStartPar
tdpd
&
\sphinxAtStartPar
atom\sphinxhyphen{}ID atom\sphinxhyphen{}type x y z cc1 cc2 … ccNspecies
\\
\sphinxhline
\sphinxAtStartPar
template
&
\sphinxAtStartPar
atom\sphinxhyphen{}ID atom\sphinxhyphen{}type molecule\sphinxhyphen{}ID template\sphinxhyphen{}index template\sphinxhyphen{}atom x y z
\\
\sphinxhline
\sphinxAtStartPar
tri
&
\sphinxAtStartPar
atom\sphinxhyphen{}ID molecule\sphinxhyphen{}ID atom\sphinxhyphen{}type triangleflag density x y z
\\
\sphinxhline
\sphinxAtStartPar
wavepacket
&
\sphinxAtStartPar
atom\sphinxhyphen{}ID atom\sphinxhyphen{}type charge espin eradius etag cs\_re cs\_im x y z
\\
\sphinxhline
\sphinxAtStartPar
hybrid
&
\sphinxAtStartPar
atom\sphinxhyphen{}ID atom\sphinxhyphen{}type x y z sub\sphinxhyphen{}style1 sub\sphinxhyphen{}style2 …
\\
\sphinxbottomrule
\end{tabulary}
\sphinxtableafterendhook\par
\sphinxattableend\end{savenotes}

\sphinxAtStartPar
The per\sphinxhyphen{}atom values have these meanings and units, listed alphabetically:
\begin{itemize}
\item {} 
\sphinxAtStartPar
atom\sphinxhyphen{}ID = integer ID of atom

\item {} 
\sphinxAtStartPar
atom\sphinxhyphen{}type = type of atom (1\sphinxhyphen{}Ntype, or type label)

\item {} 
\sphinxAtStartPar
bodyflag = 1 for body particles, 0 for point particles

\item {} 
\sphinxAtStartPar
ccN = chemical concentration for tDPD particles for each species (mole/volume units)

\item {} 
\sphinxAtStartPar
cradius = contact radius for SMD particles (distance units)

\item {} 
\sphinxAtStartPar
cs\_re,cs\_im = real/imaginary parts of wave packet coefficients

\item {} 
\sphinxAtStartPar
cv = heat capacity (need units) for SPH particles

\item {} 
\sphinxAtStartPar
density = density of particle (mass/distance\textasciicircum{}3 or mass/distance\textasciicircum{}2 or mass/distance units, depending on dimensionality of particle)

\item {} 
\sphinxAtStartPar
diameter = diameter of spherical atom (distance units)

\item {} 
\sphinxAtStartPar
edpd\_temp = temperature for eDPD particles (temperature units)

\item {} 
\sphinxAtStartPar
edpd\_cv = volumetric heat capacity for eDPD particles (energy/temperature/volume units)

\item {} 
\sphinxAtStartPar
ellipsoidflag = 1 for ellipsoidal particles, 0 for point particles

\item {} 
\sphinxAtStartPar
eradius = electron radius (or fixed\sphinxhyphen{}core radius)

\item {} 
\sphinxAtStartPar
esph = energy (need units) for SPH particles

\item {} 
\sphinxAtStartPar
espin = electron spin (+1/\sphinxhyphen{}1), 0 = nuclei, 2 = fixed\sphinxhyphen{}core, 3 = pseudo\sphinxhyphen{}cores (i.e. ECP)

\item {} 
\sphinxAtStartPar
etag = integer ID of electron that each wave packet belongs to

\item {} 
\sphinxAtStartPar
kradius = kernel radius for SMD particles (distance units)

\item {} 
\sphinxAtStartPar
lineflag = 1 for line segment particles, 0 for point or spherical particles

\item {} 
\sphinxAtStartPar
mass = mass of particle (mass units)

\item {} 
\sphinxAtStartPar
molecule\sphinxhyphen{}ID = integer ID of molecule the atom belongs to

\item {} 
\sphinxAtStartPar
mux,muy,muz = components of dipole moment of atom (dipole units) (see general triclinic note below)

\item {} 
\sphinxAtStartPar
q = charge on atom (charge units)

\item {} 
\sphinxAtStartPar
rho = density (need units) for SPH particles

\item {} 
\sphinxAtStartPar
sp = magnitude of magnetic spin of atom (Bohr magnetons)

\item {} 
\sphinxAtStartPar
spx,spy,spz = components of magnetic spin of atom (unit vector) (see general triclinic note below)

\item {} 
\sphinxAtStartPar
template\sphinxhyphen{}atom = which atom within a template molecule the atom is

\item {} 
\sphinxAtStartPar
template\sphinxhyphen{}index = which molecule within the molecule template the atom is part of

\item {} 
\sphinxAtStartPar
theta = internal temperature of a DPD particle

\item {} 
\sphinxAtStartPar
triangleflag = 1 for triangular particles, 0 for point or spherical particles

\item {} 
\sphinxAtStartPar
volume = volume of Peridynamic particle (distance\textasciicircum{}3 units)

\item {} 
\sphinxAtStartPar
x,y,z = coordinates of atom (distance units)

\item {} 
\sphinxAtStartPar
x0,y0,z0 = original (strain\sphinxhyphen{}free) coordinates of atom (distance
units) (see general triclinic note below)

\end{itemize}

\sphinxAtStartPar
The units for these quantities depend on the unit style; see the
{\hyperref[\detokenize{units::doc}]{\sphinxcrossref{\DUrole{doc}{units}}}} command for details.

\sphinxAtStartPar
For 2d simulations, the atom coordinate z must be specified as 0.0.
If the data file is created by another program, then z values for a 2d
simulation can be within epsilon of 0.0, and LAMMPS will force them to
zero.

\begin{sphinxadmonition}{note}{Note:}
\sphinxAtStartPar
If the data file defines a general triclinic box, then the
following per\sphinxhyphen{}atom values in the list above are per\sphinxhyphen{}atom vectors
which imply an orientation: (mux,muy,muz) and (spx,spy,spz).  This
means they should be specified consistent with the general
triclinic box and its orientation relative to the standard x,y,z
coordinate axes.  For example a dipole moment vector which will be
in the +x direction once LAMMPS converts from a general to
restricted triclinic box, should be specified in the data file in
the direction of the \sphinxstylestrong{A} edge vector.  Likewise the (x0,y0,z0)
per\sphinxhyphen{}atom strain\sphinxhyphen{}free coordinates should be inside the general
triclinic simulation box as explained in the note above.  See the
{\hyperref[\detokenize{Howto_triclinic::doc}]{\sphinxcrossref{\DUrole{doc}{Howto triclinic}}}} doc page for more details.
\end{sphinxadmonition}

\sphinxAtStartPar
The atom\sphinxhyphen{}ID is used to identify the atom throughout the simulation and
in dump files.  Normally, it is a unique value from 1 to Natoms for
each atom.  Unique values larger than Natoms can be used, but they
will cause extra memory to be allocated on each processor, if an atom
map array is used, but not if an atom map hash is used; see the
{\hyperref[\detokenize{atom_modify::doc}]{\sphinxcrossref{\DUrole{doc}{atom\_modify}}}} command for details.  If an atom map is
not used (e.g. an atomic system with no bonds), and you don’t care if
unique atom IDs appear in dump files, then the atom\sphinxhyphen{}IDs can all be set
to 0.

\sphinxAtStartPar
The atom\sphinxhyphen{}type can be a numeric value or an alphanumeric label.  The
latter is only allowed if the type label has been defined by the
{\hyperref[\detokenize{labelmap::doc}]{\sphinxcrossref{\DUrole{doc}{labelmap}}}} command or an Atom Type Labels section
earlier in the data file.  See the {\hyperref[\detokenize{Howto_type_labels::doc}]{\sphinxcrossref{\DUrole{doc}{Howto type labels}}}} doc page for the allowed syntax of type labels
and a general discussion of how type labels can be used.

\sphinxAtStartPar
The molecule ID is a second identifier attached to an atom.  Normally, it
is a number from 1 to N, identifying which molecule the atom belongs
to.  It can be 0 if it is a non\sphinxhyphen{}bonded atom or if you don’t care to
keep track of molecule assignments.

\sphinxAtStartPar
The diameter specifies the size of a finite\sphinxhyphen{}size spherical particle.
It can be set to 0.0, which means that atom is a point particle.

\sphinxAtStartPar
The ellipsoidflag, lineflag, triangleflag, and bodyflag determine
whether the particle is a finite\sphinxhyphen{}size ellipsoid or line or triangle or
body of finite size, or whether the particle is a point particle.
Additional attributes must be defined for each ellipsoid, line,
triangle, or body in the corresponding \sphinxstyleemphasis{Ellipsoids}, \sphinxstyleemphasis{Lines},
\sphinxstyleemphasis{Triangles}, or \sphinxstyleemphasis{Bodies} section.

\sphinxAtStartPar
The \sphinxstyleemphasis{template\sphinxhyphen{}index} and \sphinxstyleemphasis{template\sphinxhyphen{}atom} are only defined used by
{\hyperref[\detokenize{atom_style::doc}]{\sphinxcrossref{\DUrole{doc}{atom\_style template}}}}.  In this case the
{\hyperref[\detokenize{molecule::doc}]{\sphinxcrossref{\DUrole{doc}{molecule}}}} command is used to define a molecule template
which contains one or more molecules (as separate files).  If an atom
belongs to one of those molecules, its \sphinxstyleemphasis{template\sphinxhyphen{}index} and \sphinxstyleemphasis{template\sphinxhyphen{}atom}
are both set to positive integers; if not the values are both 0.  The
\sphinxstyleemphasis{template\sphinxhyphen{}index} is which molecule (1 to Nmols) the atom belongs to.
The \sphinxstyleemphasis{template\sphinxhyphen{}atom} is which atom (1 to Natoms) within the molecule
the atom is.

\sphinxAtStartPar
Some pair styles and fixes and computes that operate on finite\sphinxhyphen{}size
particles allow for a mixture of finite\sphinxhyphen{}size and point particles.  See
the doc pages of individual commands for details.

\sphinxAtStartPar
For finite\sphinxhyphen{}size particles, the density is used in conjunction with the
particle volume to set the mass of each particle as mass = density *
volume.  In this context, volume can be a 3d quantity (for spheres or
ellipsoids), a 2d quantity (for triangles), or a 1d quantity (for line
segments).  If the volume is 0.0, meaning a point particle, then the
density value is used as the mass.  One exception is for the body atom
style, in which case the mass of each particle (body or point
particle) is specified explicitly.  This is because the volume of the
body is unknown.

\sphinxAtStartPar
Note that for 2d simulations of spheres, this command will treat them
as spheres when converting density to mass.  However, they can also be
modeled as 2d discs (circles) if the {\hyperref[\detokenize{set::doc}]{\sphinxcrossref{\DUrole{doc}{set density/disc}}}}
command is used to reset their mass after the read\_data command is
used.  A \sphinxstyleemphasis{disc} keyword can also be used with time integration fixes,
such as {\hyperref[\detokenize{fix_nve_sphere::doc}]{\sphinxcrossref{\DUrole{doc}{fix nve/sphere}}}} and {\hyperref[\detokenize{fix_nve_sphere::doc}]{\sphinxcrossref{\DUrole{doc}{fix
nvt/sphere}}}} to time integrate their motion as 2d
discs (not 3d spheres), by changing their moment of inertia.

\sphinxAtStartPar
For atom\_style hybrid, following the 5 initial values
(ID,type,x,y,z), specific values for each sub\sphinxhyphen{}style must be listed.
The order of the sub\sphinxhyphen{}styles is the same as they were listed in the
{\hyperref[\detokenize{atom_style::doc}]{\sphinxcrossref{\DUrole{doc}{atom\_style}}}} command.  The specific values for each
sub\sphinxhyphen{}style are those that are not the 5 standard ones (ID,type,x,y,z).
For example, for the “charge” sub\sphinxhyphen{}style, a “q” value would appear.
For the “full” sub\sphinxhyphen{}style, a “molecule\sphinxhyphen{}ID” and “q” would appear.  These
are listed in the same order they appear as listed above.  Thus if

\begin{sphinxVerbatim}[commandchars=\\\{\}]
atom\PYGZus{}style hybrid charge sphere
\end{sphinxVerbatim}

\sphinxAtStartPar
were used in the input script, each atom line would have these fields:

\begin{sphinxVerbatim}[commandchars=\\\{\}]
atom\PYGZhy{}ID atom\PYGZhy{}type x y z q diameter density
\end{sphinxVerbatim}

\sphinxAtStartPar
Note that if a non\sphinxhyphen{}standard value is defined by multiple sub\sphinxhyphen{}styles,
it only appears once in the atom line.  E.g. the atom line for
atom\_style hybrid dipole full would list “q” only once, with the
dipole sub\sphinxhyphen{}style fields; “q” does not appear with the full sub\sphinxhyphen{}style
fields.

\begin{sphinxVerbatim}[commandchars=\\\{\}]
atom\PYGZhy{}ID atom\PYGZhy{}type x y z q mux muy myz molecule\PYGZhy{}ID
\end{sphinxVerbatim}

\sphinxAtStartPar
Atom lines specify the (x,y,z) coordinates of atoms.  These can be
inside or outside the simulation box.  When the data file is read,
LAMMPS wraps coordinates outside the box back into the box for
dimensions that are periodic.  As discussed above, if an atom is
outside the box in a non\sphinxhyphen{}periodic dimension, it will be lost.

\sphinxAtStartPar
LAMMPS always stores atom coordinates as values which are inside the
simulation box.  It also stores 3 flags which indicate which image of
the simulation box (in each dimension) the atom would be in if its
coordinates were unwrapped across periodic boundaries.  An image flag
of 0 means the atom is still inside the box when unwrapped.  A value
of 2 means add 2 box lengths to get the unwrapped coordinate.  A value
of \sphinxhyphen{}1 means subtract 1 box length to get the unwrapped coordinate.
LAMMPS updates these flags as atoms cross periodic boundaries during
the simulation.  The {\hyperref[\detokenize{dump::doc}]{\sphinxcrossref{\DUrole{doc}{dump}}}} command can output atom
coordinates in wrapped or unwrapped form, as well as the 3 image
flags.

\sphinxAtStartPar
In the data file, atom lines (all lines or none of them) can
optionally list 3 trailing integer values (nx,ny,nz), which are used
to initialize the atom’s image flags.  If nx,ny,nz values are not
listed in the data file, LAMMPS initializes them to 0.  Note that the
image flags are immediately updated if an atom’s coordinates need to
wrapped back into the simulation box.

\sphinxAtStartPar
It is only important to set image flags correctly in a data file if a
simulation model relies on unwrapped coordinates for some calculation;
otherwise they can be left unspecified.  Examples of LAMMPS commands
that use unwrapped coordinates internally are as follows:
\begin{itemize}
\item {} 
\sphinxAtStartPar
Atoms in a rigid body (see {\hyperref[\detokenize{fix_rigid::doc}]{\sphinxcrossref{\DUrole{doc}{fix rigid}}}}, {\hyperref[\detokenize{fix_rigid::doc}]{\sphinxcrossref{\DUrole{doc}{fix rigid/small}}}}) must have consistent image flags, so that
when the atoms are unwrapped, they are near each other, i.e. as a
single body.

\item {} 
\sphinxAtStartPar
If the {\hyperref[\detokenize{replicate::doc}]{\sphinxcrossref{\DUrole{doc}{replicate}}}} command is used to generate a larger
system, image flags must be consistent for bonded atoms when the bond
crosses a periodic boundary.  I.e. the values of the image flags
should be different by 1 (in the appropriate dimension) for the two
atoms in such a bond.

\item {} 
\sphinxAtStartPar
If you plan to {\hyperref[\detokenize{dump::doc}]{\sphinxcrossref{\DUrole{doc}{dump}}}} image flags and perform post\sphinxhyphen{}analysis
that will unwrap atom coordinates, it may be important that a
continued run (restarted from a data file) begins with image flags
that are consistent with the previous run.

\end{itemize}

\begin{sphinxadmonition}{note}{Note:}
\sphinxAtStartPar
If your system is an infinite periodic crystal with bonds then
it is impossible to have fully consistent image flags.  This is because
some bonds will cross periodic boundaries and connect two atoms with the
same image flag.
\end{sphinxadmonition}

\sphinxAtStartPar
Atom velocities and other atom quantities not defined above are set to
0.0 when the \sphinxstyleemphasis{Atoms} section is read.  Velocities can be set later by
a \sphinxstyleemphasis{Velocities} section in the data file or by a {\hyperref[\detokenize{velocity::doc}]{\sphinxcrossref{\DUrole{doc}{velocity}}}} or {\hyperref[\detokenize{set::doc}]{\sphinxcrossref{\DUrole{doc}{set}}}} command in the input script.


\bigskip\hrule\bigskip


\sphinxAtStartPar
\sphinxstyleemphasis{Bodies} section:
\begin{itemize}
\item {} 
\sphinxAtStartPar
one or more lines per body

\item {} 
\sphinxAtStartPar
first line syntax: atom\sphinxhyphen{}ID Ninteger Ndouble

\begin{sphinxVerbatim}[commandchars=\\\{\}]
Ninteger = \PYGZsh{} of integer quantities for this particle
Ndouble = \PYGZsh{} of floating\PYGZhy{}point quantities for this particle
\end{sphinxVerbatim}

\item {} 
\sphinxAtStartPar
0 or more integer lines with total of Ninteger values

\item {} 
\sphinxAtStartPar
0 or more double lines with total of Ndouble values

\item {} 
\sphinxAtStartPar
example:

\begin{sphinxVerbatim}[commandchars=\\\{\}]
12 3 6
2 3 2
1.0 2.0 3.0 1.0 2.0 4.0
\end{sphinxVerbatim}

\item {} 
\sphinxAtStartPar
example:

\begin{sphinxVerbatim}[commandchars=\\\{\}]
12 0 14
1.0 2.0 3.0 1.0 2.0 4.0 1.0
2.0 3.0 1.0 2.0 4.0 4.0 2.0
\end{sphinxVerbatim}

\end{itemize}

\sphinxAtStartPar
The \sphinxstyleemphasis{Bodies} section must appear if {\hyperref[\detokenize{atom_style::doc}]{\sphinxcrossref{\DUrole{doc}{atom\_style body}}}}
is used and any atoms listed in the \sphinxstyleemphasis{Atoms} section have a bodyflag =
1.  The number of bodies should be specified in the header section via
the “bodies” keyword.

\sphinxAtStartPar
Each body can have a variable number of integer and/or floating\sphinxhyphen{}point
values.  The number and meaning of the values is defined by the body
style, as described in the {\hyperref[\detokenize{Howto_body::doc}]{\sphinxcrossref{\DUrole{doc}{Howto body}}}} doc page.
The body style is given as an argument to the {\hyperref[\detokenize{atom_style::doc}]{\sphinxcrossref{\DUrole{doc}{atom\_style body}}}} command.

\sphinxAtStartPar
The Ninteger and Ndouble values determine how many integer and
floating\sphinxhyphen{}point values are specified for this particle.  Ninteger and
Ndouble can be as large as needed and can be different for every body.
Integer values are then listed next on subsequent lines.  Lines are
read one at a time until Ninteger values are read.  Floating\sphinxhyphen{}point
values follow on subsequent lines, Again lines are read one at a time
until Ndouble values are read.  Note that if there are no values of a
particular type, no lines appear for that type.

\sphinxAtStartPar
The \sphinxstyleemphasis{Bodies} section must appear after the \sphinxstyleemphasis{Atoms} section.


\bigskip\hrule\bigskip


\sphinxAtStartPar
\sphinxstyleemphasis{Bond Coeffs} section:
\begin{itemize}
\item {} 
\sphinxAtStartPar
one line per bond type

\item {} 
\sphinxAtStartPar
line syntax: ID coeffs

\begin{sphinxVerbatim}[commandchars=\\\{\}]
ID = bond type (1\PYGZhy{}N)
coeffs = list of coeffs
\end{sphinxVerbatim}

\item {} 
\sphinxAtStartPar
example:

\begin{sphinxVerbatim}[commandchars=\\\{\}]
4 250 1.49
\end{sphinxVerbatim}

\end{itemize}

\sphinxAtStartPar
The number and meaning of the coefficients are specific to the defined
bond style.  See the {\hyperref[\detokenize{bond_style::doc}]{\sphinxcrossref{\DUrole{doc}{bond\_style}}}} and
{\hyperref[\detokenize{bond_coeff::doc}]{\sphinxcrossref{\DUrole{doc}{bond\_coeff}}}} commands for details.  Coefficients can
also be set via the {\hyperref[\detokenize{bond_coeff::doc}]{\sphinxcrossref{\DUrole{doc}{bond\_coeff}}}} command in the input
script.


\bigskip\hrule\bigskip


\sphinxAtStartPar
\sphinxstyleemphasis{Bond Type Labels} section:
\begin{itemize}
\item {} 
\sphinxAtStartPar
one line per bond type

\item {} 
\sphinxAtStartPar
line syntax: ID label

\begin{sphinxVerbatim}[commandchars=\\\{\}]
ID = bond type (1\PYGZhy{}N)
label = alphanumeric type label
\end{sphinxVerbatim}

\end{itemize}

\sphinxAtStartPar
Define alphanumeric type labels for each numeric bond type.  These can
be used in the Bonds section in place of a numeric type, but only if
the this section appears before the Angles section.

\sphinxAtStartPar
See the {\hyperref[\detokenize{Howto_type_labels::doc}]{\sphinxcrossref{\DUrole{doc}{Howto type labels}}}} doc page for the
allowed syntax of type labels and a general discussion of how type
labels can be used.


\bigskip\hrule\bigskip


\sphinxAtStartPar
\sphinxstyleemphasis{BondAngle Coeffs} section:
\begin{itemize}
\item {} 
\sphinxAtStartPar
one line per angle type

\item {} 
\sphinxAtStartPar
line syntax: ID coeffs
\begin{sphinxalltt}
ID = angle type (1\sphinxhyphen{}N)
coeffs = list of coeffs (see class 2 section of {\hyperref[\detokenize{angle_coeff::doc}]{\sphinxcrossref{\DUrole{doc}{angle\_coeff}}}})
\end{sphinxalltt}

\end{itemize}


\bigskip\hrule\bigskip


\sphinxAtStartPar
\sphinxstyleemphasis{BondBond Coeffs} section:
\begin{itemize}
\item {} 
\sphinxAtStartPar
one line per angle type

\item {} 
\sphinxAtStartPar
line syntax: ID coeffs
\begin{sphinxalltt}
ID = angle type (1\sphinxhyphen{}N)
coeffs = list of coeffs (see class 2 section of {\hyperref[\detokenize{angle_coeff::doc}]{\sphinxcrossref{\DUrole{doc}{angle\_coeff}}}})
\end{sphinxalltt}

\end{itemize}


\bigskip\hrule\bigskip


\sphinxAtStartPar
\sphinxstyleemphasis{BondBond13 Coeffs} section:
\begin{itemize}
\item {} 
\sphinxAtStartPar
one line per dihedral type

\item {} 
\sphinxAtStartPar
line syntax: ID coeffs
\begin{sphinxalltt}
ID = dihedral type (1\sphinxhyphen{}N)
coeffs = list of coeffs (see class 2 section of {\hyperref[\detokenize{dihedral_coeff::doc}]{\sphinxcrossref{\DUrole{doc}{dihedral\_coeff}}}})
\end{sphinxalltt}

\end{itemize}


\bigskip\hrule\bigskip


\sphinxAtStartPar
\sphinxstyleemphasis{Bonds} section:
\begin{itemize}
\item {} 
\sphinxAtStartPar
one line per bond

\item {} 
\sphinxAtStartPar
line syntax: ID type atom1 atom2

\begin{sphinxVerbatim}[commandchars=\\\{\}]
ID = bond number (1\PYGZhy{}Nbonds)
type = bond type (1\PYGZhy{}Nbondtype, or type label)
atom1,atom2 = IDs of 1st,2nd atom in bond
\end{sphinxVerbatim}

\item {} 
\sphinxAtStartPar
example:

\begin{sphinxVerbatim}[commandchars=\\\{\}]
12 3 17 29
\end{sphinxVerbatim}

\end{itemize}

\sphinxAtStartPar
The \sphinxstyleemphasis{Bonds} section must appear after the \sphinxstyleemphasis{Atoms} section.

\sphinxAtStartPar
All values in this section must be integers (1, not 1.0).  However,
the type can be a numeric value or an alphanumeric label.  The latter
is only allowed if the type label has been defined by the
{\hyperref[\detokenize{labelmap::doc}]{\sphinxcrossref{\DUrole{doc}{labelmap}}}} command or a Bond Type Labels section
earlier in the data file.  See the {\hyperref[\detokenize{Howto_type_labels::doc}]{\sphinxcrossref{\DUrole{doc}{Howto type labels}}}} doc page for the allowed syntax of type labels
and a general discussion of how type labels can be used.


\bigskip\hrule\bigskip


\sphinxAtStartPar
\sphinxstyleemphasis{Dihedral Coeffs} section:
\begin{itemize}
\item {} 
\sphinxAtStartPar
one line per dihedral type

\item {} 
\sphinxAtStartPar
line syntax: ID coeffs

\begin{sphinxVerbatim}[commandchars=\\\{\}]
ID = dihedral type (1\PYGZhy{}N)
coeffs = list of coeffs
\end{sphinxVerbatim}

\item {} 
\sphinxAtStartPar
example:

\begin{sphinxVerbatim}[commandchars=\\\{\}]
3 0.6 1 0 1
\end{sphinxVerbatim}

\end{itemize}

\sphinxAtStartPar
The number and meaning of the coefficients are specific to the defined
dihedral style.  See the {\hyperref[\detokenize{dihedral_style::doc}]{\sphinxcrossref{\DUrole{doc}{dihedral\_style}}}} and
{\hyperref[\detokenize{dihedral_coeff::doc}]{\sphinxcrossref{\DUrole{doc}{dihedral\_coeff}}}} commands for details.
Coefficients can also be set via the
{\hyperref[\detokenize{dihedral_coeff::doc}]{\sphinxcrossref{\DUrole{doc}{dihedral\_coeff}}}} command in the input script.


\bigskip\hrule\bigskip


\sphinxAtStartPar
\sphinxstyleemphasis{Dihedral Type Labels} section:
\begin{itemize}
\item {} 
\sphinxAtStartPar
one line per dihedral type

\item {} 
\sphinxAtStartPar
line syntax: ID label

\begin{sphinxVerbatim}[commandchars=\\\{\}]
ID = dihedral type (1\PYGZhy{}N)
label = alphanumeric type label
\end{sphinxVerbatim}

\end{itemize}

\sphinxAtStartPar
Define alphanumeric type labels for each numeric dihedral type.  These
can be used in the Dihedrals section in place of a numeric type, but
only if the this section appears before the Dihedrals section.

\sphinxAtStartPar
See the {\hyperref[\detokenize{Howto_type_labels::doc}]{\sphinxcrossref{\DUrole{doc}{Howto type labels}}}} doc page for the
allowed syntax of type labels and a general discussion of how type
labels can be used.


\bigskip\hrule\bigskip


\sphinxAtStartPar
\sphinxstyleemphasis{Dihedrals} section:
\begin{itemize}
\item {} 
\sphinxAtStartPar
one line per dihedral

\item {} 
\sphinxAtStartPar
line syntax: ID type atom1 atom2 atom3 atom4

\begin{sphinxVerbatim}[commandchars=\\\{\}]
ID = number of dihedral (1\PYGZhy{}Ndihedrals)
type = dihedral type (1\PYGZhy{}Ndihedraltype, or type label)
atom1,atom2,atom3,atom4 = IDs of 1st,2nd,3rd,4th atom in dihedral
\end{sphinxVerbatim}

\item {} 
\sphinxAtStartPar
example:

\begin{sphinxVerbatim}[commandchars=\\\{\}]
12 4 17 29 30 21
\end{sphinxVerbatim}

\end{itemize}

\sphinxAtStartPar
The 4 atoms are ordered linearly within the dihedral.  The \sphinxstyleemphasis{Dihedrals}
section must appear after the \sphinxstyleemphasis{Atoms} section.

\sphinxAtStartPar
All values in this section must be integers (1, not 1.0).  However,
the type can be a numeric value or an alphanumeric label.  The latter
is only allowed if the type label has been defined by the
{\hyperref[\detokenize{labelmap::doc}]{\sphinxcrossref{\DUrole{doc}{labelmap}}}} command or a Dihedral Type Labels section
earlier in the data file.  See the {\hyperref[\detokenize{Howto_type_labels::doc}]{\sphinxcrossref{\DUrole{doc}{Howto type labels}}}} doc page for the allowed syntax of type labels
and a general discussion of how type labels can be used.


\bigskip\hrule\bigskip


\sphinxAtStartPar
\sphinxstyleemphasis{Ellipsoids} section:
\begin{itemize}
\item {} 
\sphinxAtStartPar
one line per ellipsoid

\item {} 
\sphinxAtStartPar
line syntax: atom\sphinxhyphen{}ID shapex shapey shapez quatw quati quatj quatk

\begin{sphinxVerbatim}[commandchars=\\\{\}]
atom\PYGZhy{}ID = ID of atom which is an ellipsoid
shapex,shapey,shapez = 3 diameters of ellipsoid (distance units)
quatw,quati,quatj,quatk = quaternion components for orientation of atom
\end{sphinxVerbatim}

\item {} 
\sphinxAtStartPar
example:

\begin{sphinxVerbatim}[commandchars=\\\{\}]
12 1 2 1 1 0 0 0
\end{sphinxVerbatim}

\end{itemize}

\sphinxAtStartPar
The \sphinxstyleemphasis{Ellipsoids} section must appear if {\hyperref[\detokenize{atom_style::doc}]{\sphinxcrossref{\DUrole{doc}{atom\_style ellipsoid}}}} is used and any atoms are listed in the \sphinxstyleemphasis{Atoms} section
with an ellipsoidflag = 1.  The number of ellipsoids should be
specified in the header section via the “ellipsoids” keyword.

\sphinxAtStartPar
The 3 shape values specify the 3 diameters or aspect ratios of a
finite\sphinxhyphen{}size ellipsoidal particle, when it is oriented along the 3
coordinate axes.  They must all be non\sphinxhyphen{}zero values.

\sphinxAtStartPar
The values \sphinxstyleemphasis{quatw}, \sphinxstyleemphasis{quati}, \sphinxstyleemphasis{quatj}, and \sphinxstyleemphasis{quatk} set the orientation
of the atom as a quaternion (4\sphinxhyphen{}vector).  Note that the shape
attributes specify the aspect ratios of an ellipsoidal particle, which
is oriented by default with its x\sphinxhyphen{}axis along the simulation box’s
x\sphinxhyphen{}axis, and similarly for y and z.  If this body is rotated (via the
right\sphinxhyphen{}hand rule) by an angle theta around a unit vector (a,b,c), then
the quaternion that represents its new orientation is given by
(cos(theta/2), a*sin(theta/2), b*sin(theta/2), c*sin(theta/2)).  These
4 components are quatw, quati, quatj, and quatk as specified above.
LAMMPS normalizes each atom’s quaternion in case (a,b,c) is not
specified as a unit vector.

\sphinxAtStartPar
If the data file defines a general triclinic box, then the quaternion
for each ellipsoid should be specified for its orientation relative to
the standard x,y,z coordinate axes.  When the system is converted to a
restricted triclinic box, the ellipsoid quaternions will be altered to
reflect the new orientation of the ellipsoid.

\sphinxAtStartPar
The \sphinxstyleemphasis{Ellipsoids} section must appear after the \sphinxstyleemphasis{Atoms} section.


\bigskip\hrule\bigskip


\sphinxAtStartPar
\sphinxstyleemphasis{EndBondTorsion Coeffs} section:
\begin{itemize}
\item {} 
\sphinxAtStartPar
one line per dihedral type

\item {} 
\sphinxAtStartPar
line syntax: ID coeffs
\begin{sphinxalltt}
ID = dihedral type (1\sphinxhyphen{}N)
coeffs = list of coeffs (see class 2 section of {\hyperref[\detokenize{dihedral_coeff::doc}]{\sphinxcrossref{\DUrole{doc}{dihedral\_coeff}}}})
\end{sphinxalltt}

\end{itemize}


\bigskip\hrule\bigskip


\sphinxAtStartPar
\sphinxstyleemphasis{Improper Coeffs} section:
\begin{itemize}
\item {} 
\sphinxAtStartPar
one line per improper type

\item {} 
\sphinxAtStartPar
line syntax: ID coeffs

\begin{sphinxVerbatim}[commandchars=\\\{\}]
ID = improper type (1\PYGZhy{}N)
coeffs = list of coeffs
\end{sphinxVerbatim}

\item {} 
\sphinxAtStartPar
example:

\begin{sphinxVerbatim}[commandchars=\\\{\}]
2 20 0.0548311
\end{sphinxVerbatim}

\end{itemize}

\sphinxAtStartPar
The number and meaning of the coefficients are specific to the defined
improper style.  See the {\hyperref[\detokenize{improper_style::doc}]{\sphinxcrossref{\DUrole{doc}{improper\_style}}}} and
{\hyperref[\detokenize{improper_coeff::doc}]{\sphinxcrossref{\DUrole{doc}{improper\_coeff}}}} commands for details.
Coefficients can also be set via the
{\hyperref[\detokenize{improper_coeff::doc}]{\sphinxcrossref{\DUrole{doc}{improper\_coeff}}}} command in the input script.


\bigskip\hrule\bigskip


\sphinxAtStartPar
\sphinxstyleemphasis{Improper Type Labels} section:
\begin{itemize}
\item {} 
\sphinxAtStartPar
one line per improper type

\item {} 
\sphinxAtStartPar
line syntax: ID label

\begin{sphinxVerbatim}[commandchars=\\\{\}]
ID = improper type (1\PYGZhy{}N)
label = alphanumeric type label
\end{sphinxVerbatim}

\end{itemize}

\sphinxAtStartPar
Define alphanumeric type labels for each numeric improper type.  These
can be used in the Impropers section in place of a numeric type, but
only if the this section appears before the Impropers section.

\sphinxAtStartPar
See the {\hyperref[\detokenize{Howto_type_labels::doc}]{\sphinxcrossref{\DUrole{doc}{Howto type labels}}}} doc page for the
allowed syntax of type labels and a general discussion of how type
labels can be used.


\bigskip\hrule\bigskip


\sphinxAtStartPar
\sphinxstyleemphasis{Impropers} section:
\begin{itemize}
\item {} 
\sphinxAtStartPar
one line per improper

\item {} 
\sphinxAtStartPar
line syntax: ID type atom1 atom2 atom3 atom4

\begin{sphinxVerbatim}[commandchars=\\\{\}]
ID = number of improper (1\PYGZhy{}Nimpropers)
type = improper type (1\PYGZhy{}Nimpropertype, or type label)
atom1,atom2,atom3,atom4 = IDs of 1st,2nd,3rd,4th atom in improper
\end{sphinxVerbatim}

\item {} 
\sphinxAtStartPar
example:

\begin{sphinxVerbatim}[commandchars=\\\{\}]
12 3 17 29 13 100
\end{sphinxVerbatim}

\end{itemize}

\sphinxAtStartPar
The ordering of the 4 atoms determines the definition of the improper
angle used in the formula for each {\hyperref[\detokenize{improper_style::doc}]{\sphinxcrossref{\DUrole{doc}{improper style}}}}.  See the doc pages for individual styles for
details.

\sphinxAtStartPar
The \sphinxstyleemphasis{Impropers} section must appear after the \sphinxstyleemphasis{Atoms} section.

\sphinxAtStartPar
All values in this section must be integers (1, not 1.0).  However,
the type can be a numeric value or an alphanumeric label.  The latter
is only allowed if the type label has been defined by the
{\hyperref[\detokenize{labelmap::doc}]{\sphinxcrossref{\DUrole{doc}{labelmap}}}} command or a Improper Type Labels section
earlier in the data file.  See the {\hyperref[\detokenize{Howto_type_labels::doc}]{\sphinxcrossref{\DUrole{doc}{Howto type labels}}}} doc page for the allowed syntax of type labels
and a general discussion of how type labels can be used.


\bigskip\hrule\bigskip


\sphinxAtStartPar
\sphinxstyleemphasis{Lines} section:
\begin{itemize}
\item {} 
\sphinxAtStartPar
one line per line segment

\item {} 
\sphinxAtStartPar
line syntax: atom\sphinxhyphen{}ID x1 y1 x2 y2

\begin{sphinxVerbatim}[commandchars=\\\{\}]
atom\PYGZhy{}ID = ID of atom which is a line segment
x1,y1 = 1st end point
x2,y2 = 2nd end point
\end{sphinxVerbatim}

\item {} 
\sphinxAtStartPar
example:

\begin{sphinxVerbatim}[commandchars=\\\{\}]
12 1.0 0.0 2.0 0.0
\end{sphinxVerbatim}

\end{itemize}

\sphinxAtStartPar
The \sphinxstyleemphasis{Lines} section must appear if {\hyperref[\detokenize{atom_style::doc}]{\sphinxcrossref{\DUrole{doc}{atom\_style line}}}}
is used and any atoms are listed in the \sphinxstyleemphasis{Atoms} section with a
lineflag = 1.  The number of lines should be specified in the header
section via the “lines” keyword.

\sphinxAtStartPar
The 2 end points are the end points of the line segment.  They should
be values close to the center point of the line segment specified in
the Atoms section of the data file, even if individual end points are
outside the simulation box.

\sphinxAtStartPar
The ordering of the 2 points should be such that using a right\sphinxhyphen{}hand
rule to cross the line segment with a unit vector in the +z direction,
gives an “outward” normal vector perpendicular to the line segment.
I.e. normal = (c2\sphinxhyphen{}c1) x (0,0,1).  This orientation may be important
for defining some interactions.

\sphinxAtStartPar
If the data file defines a general triclinic box, then the x1,y1 and
x2,y2 values for each line segment should be specified for its
orientation relative to the standard x,y,z coordinate axes.  When the
system is converted to a restricted triclinic box, the x1,y1,x2,y2
values will be altered to reflect the new orientation of the line
segment.

\sphinxAtStartPar
The \sphinxstyleemphasis{Lines} section must appear after the \sphinxstyleemphasis{Atoms} section.


\bigskip\hrule\bigskip


\sphinxAtStartPar
\sphinxstyleemphasis{Masses} section:
\begin{itemize}
\item {} 
\sphinxAtStartPar
one line per atom type

\item {} 
\sphinxAtStartPar
line syntax: ID mass

\begin{sphinxVerbatim}[commandchars=\\\{\}]
ID = atom type (1\PYGZhy{}N or atom type label)
mass = mass value
\end{sphinxVerbatim}

\item {} 
\sphinxAtStartPar
example:

\begin{sphinxVerbatim}[commandchars=\\\{\}]
3 1.01
\end{sphinxVerbatim}

\end{itemize}

\sphinxAtStartPar
This defines the mass of each atom type.  This can also be set via the
{\hyperref[\detokenize{mass::doc}]{\sphinxcrossref{\DUrole{doc}{mass}}}} command in the input script.  This section cannot be
used for atom styles that define a mass for individual atoms \sphinxhyphen{}
e.g. {\hyperref[\detokenize{atom_style::doc}]{\sphinxcrossref{\DUrole{doc}{atom\_style sphere}}}}.

\sphinxAtStartPar
Using type labels instead of atom type numbers is only allowed if the
type label has been defined by the {\hyperref[\detokenize{labelmap::doc}]{\sphinxcrossref{\DUrole{doc}{labelmap}}}} command or
a Atom Type Labels section earlier in the data file.  See the
{\hyperref[\detokenize{Howto_type_labels::doc}]{\sphinxcrossref{\DUrole{doc}{Howto type labels}}}} doc page for the allowed
syntax of type labels and a general discussion of how type labels can be
used.


\bigskip\hrule\bigskip


\sphinxAtStartPar
\sphinxstyleemphasis{MiddleBondTorsion Coeffs} section:
\begin{itemize}
\item {} 
\sphinxAtStartPar
one line per dihedral type

\item {} 
\sphinxAtStartPar
line syntax: ID coeffs
\begin{sphinxalltt}
ID = dihedral type (1\sphinxhyphen{}N)
coeffs = list of coeffs (see class 2 section of {\hyperref[\detokenize{dihedral_coeff::doc}]{\sphinxcrossref{\DUrole{doc}{dihedral\_coeff}}}})
\end{sphinxalltt}

\end{itemize}


\bigskip\hrule\bigskip


\sphinxAtStartPar
\sphinxstyleemphasis{Pair Coeffs} section:
\begin{itemize}
\item {} 
\sphinxAtStartPar
one line per atom type

\item {} 
\sphinxAtStartPar
line syntax: ID coeffs

\begin{sphinxVerbatim}[commandchars=\\\{\}]
ID = atom type (1\PYGZhy{}N)
coeffs = list of coeffs
\end{sphinxVerbatim}

\item {} 
\sphinxAtStartPar
example:

\begin{sphinxVerbatim}[commandchars=\\\{\}]
3 0.022 2.35197 0.022 2.35197
\end{sphinxVerbatim}

\end{itemize}

\sphinxAtStartPar
The number and meaning of the coefficients are specific to the defined
pair style.  See the {\hyperref[\detokenize{pair_style::doc}]{\sphinxcrossref{\DUrole{doc}{pair\_style}}}} and
{\hyperref[\detokenize{pair_coeff::doc}]{\sphinxcrossref{\DUrole{doc}{pair\_coeff}}}} commands for details.  Since pair
coefficients for types I != J are not specified, these will be
generated automatically by the pair style’s mixing rule.  See the
individual pair\_style doc pages and the {\hyperref[\detokenize{pair_modify::doc}]{\sphinxcrossref{\DUrole{doc}{pair\_modify mix}}}} command for details.  Pair coefficients can also be set
via the {\hyperref[\detokenize{pair_coeff::doc}]{\sphinxcrossref{\DUrole{doc}{pair\_coeff}}}} command in the input script.


\bigskip\hrule\bigskip


\sphinxAtStartPar
\sphinxstyleemphasis{PairIJ Coeffs} section:
\begin{itemize}
\item {} 
\sphinxAtStartPar
one line per pair of atom types for all I,J with I \textless{}= J

\item {} 
\sphinxAtStartPar
line syntax: ID1 ID2 coeffs

\begin{sphinxVerbatim}[commandchars=\\\{\}]
ID1 = atom type I = 1\PYGZhy{}N
ID2 = atom type J = I\PYGZhy{}N, with I \PYGZlt{}= J
coeffs = list of coeffs
\end{sphinxVerbatim}

\item {} 
\sphinxAtStartPar
examples:

\begin{sphinxVerbatim}[commandchars=\\\{\}]
3 3 0.022 2.35197 0.022 2.35197
3 5 0.022 2.35197 0.022 2.35197
\end{sphinxVerbatim}

\end{itemize}

\sphinxAtStartPar
This section must have N*(N+1)/2 lines where N = \# of atom types.
The number and meaning of the coefficients are specific to the defined
pair style.  See the {\hyperref[\detokenize{pair_style::doc}]{\sphinxcrossref{\DUrole{doc}{pair\_style}}}} and
{\hyperref[\detokenize{pair_coeff::doc}]{\sphinxcrossref{\DUrole{doc}{pair\_coeff}}}} commands for details.  Since pair
coefficients for types I != J are all specified, these values will
turn off the default mixing rule defined by the pair style.  See the
individual pair\_style doc pages and the {\hyperref[\detokenize{pair_modify::doc}]{\sphinxcrossref{\DUrole{doc}{pair\_modify mix}}}} command for details.  Pair coefficients can also be set
via the {\hyperref[\detokenize{pair_coeff::doc}]{\sphinxcrossref{\DUrole{doc}{pair\_coeff}}}} command in the input script.


\bigskip\hrule\bigskip


\sphinxAtStartPar
\sphinxstyleemphasis{Triangles} section:
\begin{itemize}
\item {} 
\sphinxAtStartPar
one line per triangle

\item {} 
\sphinxAtStartPar
line syntax: atom\sphinxhyphen{}ID x1 y1 z1 x2 y2 z2 x3 y3 z3

\begin{sphinxVerbatim}[commandchars=\\\{\}]
atom\PYGZhy{}ID = ID of atom which is a line segment
x1,y1,z1 = 1st corner point
x2,y2,z2 = 2nd corner point
x3,y3,z3 = 3rd corner point
\end{sphinxVerbatim}

\item {} 
\sphinxAtStartPar
example:

\begin{sphinxVerbatim}[commandchars=\\\{\}]
12 0.0 0.0 0.0 2.0 0.0 1.0 0.0 2.0 1.0
\end{sphinxVerbatim}

\end{itemize}

\sphinxAtStartPar
The \sphinxstyleemphasis{Triangles} section must appear if {\hyperref[\detokenize{atom_style::doc}]{\sphinxcrossref{\DUrole{doc}{atom\_style tri}}}} is used and any atoms are listed in the \sphinxstyleemphasis{Atoms} section
with a triangleflag = 1.  The number of lines should be specified in
the header section via the “triangles” keyword.

\sphinxAtStartPar
The 3 corner points are the corner points of the triangle.  They
should be values close to the center point of the triangle specified
in the Atoms section of the data file, even if individual corner
points are outside the simulation box.

\sphinxAtStartPar
The ordering of the 3 points should be such that using a right\sphinxhyphen{}hand
rule to go from point1 to point2 to point3 gives an “outward” normal
vector to the face of the triangle.  I.e. normal = (c2\sphinxhyphen{}c1) x (c3\sphinxhyphen{}c1).
This orientation may be important for defining some interactions.

\sphinxAtStartPar
If the data file defines a general triclinic box, then the x1,y1,z1
and x2,y2,z2 and x3,y3,z3 values for each triangle should be specified
for its orientation relative to the standard x,y,z coordinate axes.
When the system is converted to a restricted triclinic box, the
x1,y1,z1,x2,y2,z2,x3,y3,z3 values will be altered to reflect the new
orientation of the triangle.

\sphinxAtStartPar
The \sphinxstyleemphasis{Triangles} section must appear after the \sphinxstyleemphasis{Atoms} section.


\bigskip\hrule\bigskip


\sphinxAtStartPar
\sphinxstyleemphasis{Velocities} section:
\begin{itemize}
\item {} 
\sphinxAtStartPar
one line per atom

\item {} 
\sphinxAtStartPar
line syntax: depends on atom style

\end{itemize}


\begin{savenotes}\sphinxattablestart
\sphinxthistablewithglobalstyle
\centering
\begin{tabulary}{\linewidth}[t]{TT}
\sphinxtoprule
\sphinxtableatstartofbodyhook
\sphinxAtStartPar
all styles except those listed
&
\sphinxAtStartPar
atom\sphinxhyphen{}ID vx vy vz
\\
\sphinxhline
\sphinxAtStartPar
electron
&
\sphinxAtStartPar
atom\sphinxhyphen{}ID vx vy vz ervel
\\
\sphinxhline
\sphinxAtStartPar
ellipsoid
&
\sphinxAtStartPar
atom\sphinxhyphen{}ID vx vy vz lx ly lz
\\
\sphinxhline
\sphinxAtStartPar
sphere
&
\sphinxAtStartPar
atom\sphinxhyphen{}ID vx vy vz wx wy wz
\\
\sphinxhline
\sphinxAtStartPar
hybrid
&
\sphinxAtStartPar
atom\sphinxhyphen{}ID vx vy vz sub\sphinxhyphen{}style1 sub\sphinxhyphen{}style2 …
\\
\sphinxbottomrule
\end{tabulary}
\sphinxtableafterendhook\par
\sphinxattableend\end{savenotes}

\sphinxAtStartPar
where the keywords have these meanings:
\begin{quote}

\begin{sphinxVerbatim}[commandchars=\\\{\}]
vx,vy,vz = translational velocity of atom
lx,ly,lz = angular momentum of aspherical atom
wx,wy,wz = angular velocity of spherical atom
ervel = electron radial velocity (0 for fixed\PYGZhy{}core)
\end{sphinxVerbatim}
\end{quote}

\sphinxAtStartPar
The velocity lines can appear in any order.  This section can only be
used after an \sphinxstyleemphasis{Atoms} section.  This is because the \sphinxstyleemphasis{Atoms} section
must have assigned a unique atom ID to each atom so that velocities
can be assigned to them.

\sphinxAtStartPar
Vx, vy, vz, and ervel are in {\hyperref[\detokenize{units::doc}]{\sphinxcrossref{\DUrole{doc}{units}}}} of velocity.  Lx, ly,
lz are in units of angular momentum (distance\sphinxhyphen{}velocity\sphinxhyphen{}mass).  Wx, Wy,
Wz are in units of angular velocity (radians/time).

\sphinxAtStartPar
If the data file defines a general triclinic box, then each of the 3
vectors (translational velocity, angular momentum, angular velocity)
should be specified for the rotated coordinate axes of the general
triclinic box.  See the {\hyperref[\detokenize{Howto_triclinic::doc}]{\sphinxcrossref{\DUrole{doc}{Howto triclinic}}}} doc
page for more details.

\sphinxAtStartPar
For atom\_style hybrid, following the 4 initial values (ID,vx,vy,vz),
specific values for each sub\sphinxhyphen{}style must be listed.  The order of the
sub\sphinxhyphen{}styles is the same as they were listed in the
{\hyperref[\detokenize{atom_style::doc}]{\sphinxcrossref{\DUrole{doc}{atom\_style}}}} command.  The sub\sphinxhyphen{}style specific values
are those that are not the 5 standard ones (ID,vx,vy,vz).  For
example, for the “sphere” sub\sphinxhyphen{}style, “wx”, “wy”, “wz” values would
appear.  These are listed in the same order they appear as listed
above.  Thus if

\begin{sphinxVerbatim}[commandchars=\\\{\}]
\PYG{k}{atom\PYGZus{}style}\PYG{+w}{ }\PYG{n}{hybrid}\PYG{+w}{ }\PYG{n}{electron}\PYG{+w}{ }\PYG{n}{sphere}
\end{sphinxVerbatim}

\sphinxAtStartPar
were used in the input script, each velocity line would have these
fields:

\begin{sphinxVerbatim}[commandchars=\\\{\}]
atom\PYGZhy{}ID vx vy vz ervel wx wy wz
\end{sphinxVerbatim}

\sphinxAtStartPar
Translational velocities can also be (re)set by the {\hyperref[\detokenize{velocity::doc}]{\sphinxcrossref{\DUrole{doc}{velocity}}}} command in the input script.


\bigskip\hrule\bigskip



\subsection{Restrictions}
\label{\detokenize{read_data:restrictions}}
\sphinxAtStartPar
To read gzipped data files, you must compile LAMMPS with the
\sphinxhyphen{}DLAMMPS\_GZIP option.  See the {\hyperref[\detokenize{Build_settings::doc}]{\sphinxcrossref{\DUrole{doc}{Build settings}}}}
doc page for details.

\sphinxAtStartPar
Label maps are currently not supported when using the KOKKOS package.


\subsection{Related commands}
\label{\detokenize{read_data:related-commands}}
\sphinxAtStartPar
{\hyperref[\detokenize{read_dump::doc}]{\sphinxcrossref{\DUrole{doc}{read\_dump}}}}, {\hyperref[\detokenize{read_restart::doc}]{\sphinxcrossref{\DUrole{doc}{read\_restart}}}},
{\hyperref[\detokenize{create_atoms::doc}]{\sphinxcrossref{\DUrole{doc}{create\_atoms}}}}, {\hyperref[\detokenize{write_data::doc}]{\sphinxcrossref{\DUrole{doc}{write\_data}}}},
{\hyperref[\detokenize{labelmap::doc}]{\sphinxcrossref{\DUrole{doc}{labelmap}}}}


\subsection{Default}
\label{\detokenize{read_data:default}}
\sphinxAtStartPar
The default for all the \sphinxstyleemphasis{extra} keywords is 0.

\sphinxstepscope

\index{read\_dump@\spxentry{read\_dump}}\ignorespaces 

\section{read\_dump command}
\label{\detokenize{read_dump:read-dump-command}}\label{\detokenize{read_dump:index-0}}\label{\detokenize{read_dump::doc}}

\subsection{Syntax}
\label{\detokenize{read_dump:syntax}}
\begin{sphinxVerbatim}[commandchars=\\\{\}]
\PYG{k}{read\PYGZus{}dump}\PYG{+w}{ }\PYG{n}{file}\PYG{+w}{ }\PYG{n}{Nstep}\PYG{+w}{ }\PYG{n}{field1}\PYG{+w}{ }\PYG{n}{field2}\PYG{+w}{ }\PYG{n}{...}\PYG{+w}{ }\PYG{n}{keyword}\PYG{+w}{ }\PYG{n}{values}\PYG{+w}{ }\PYG{n}{...}
\end{sphinxVerbatim}
\begin{itemize}
\item {} 
\sphinxAtStartPar
file = name of dump file to read

\item {} 
\sphinxAtStartPar
Nstep = snapshot timestep to read from file

\item {} 
\sphinxAtStartPar
one or more fields may be appended
\begin{sphinxalltt}
field = \sphinxstyleemphasis{x} or \sphinxstyleemphasis{y} or \sphinxstyleemphasis{z} or \sphinxstyleemphasis{vx} or \sphinxstyleemphasis{vy} or \sphinxstyleemphasis{vz} or \sphinxstyleemphasis{q} or \sphinxstyleemphasis{ix} or \sphinxstyleemphasis{iy} or \sphinxstyleemphasis{iz} or \sphinxstyleemphasis{fx} or \sphinxstyleemphasis{fy} or \sphinxstyleemphasis{fz}
  \sphinxstyleemphasis{x},\sphinxstyleemphasis{y},\sphinxstyleemphasis{z} = atom coordinates
  \sphinxstyleemphasis{vx},\sphinxstyleemphasis{vy},\sphinxstyleemphasis{vz} = velocity components
  \sphinxstyleemphasis{q} = charge
  \sphinxstyleemphasis{ix},\sphinxstyleemphasis{iy},\sphinxstyleemphasis{iz} = image flags in each dimension
  \sphinxstyleemphasis{fx},\sphinxstyleemphasis{fy},\sphinxstyleemphasis{fz} = force components
\end{sphinxalltt}

\item {} 
\sphinxAtStartPar
zero or more keyword/value pairs may be appended

\item {} 
\sphinxAtStartPar
keyword = \sphinxstyleemphasis{nfile} or \sphinxstyleemphasis{box} or \sphinxstyleemphasis{timestep} or \sphinxstyleemphasis{replace} or \sphinxstyleemphasis{purge} or \sphinxstyleemphasis{trim} or \sphinxstyleemphasis{add} or \sphinxstyleemphasis{label} or \sphinxstyleemphasis{scaled} or \sphinxstyleemphasis{wrapped} or \sphinxstyleemphasis{format}
\begin{sphinxalltt}
\sphinxstyleemphasis{nfile} value = Nfiles = how many parallel dump files exist
\sphinxstyleemphasis{box} value = \sphinxstyleemphasis{yes} or \sphinxstyleemphasis{no} = replace simulation box with dump box
\sphinxstyleemphasis{timestep} value = \sphinxstyleemphasis{yes} or \sphinxstyleemphasis{no} = reset simulation timestep with dump timestep
\sphinxstyleemphasis{replace} value = \sphinxstyleemphasis{yes} or \sphinxstyleemphasis{no} = overwrite atoms with dump atoms
\sphinxstyleemphasis{purge} value = \sphinxstyleemphasis{yes} or \sphinxstyleemphasis{no} = delete all atoms before adding dump atoms
\sphinxstyleemphasis{trim} value = \sphinxstyleemphasis{yes} or \sphinxstyleemphasis{no} = trim atoms not in dump snapshot
\sphinxstyleemphasis{add} value = \sphinxstyleemphasis{yes} or \sphinxstyleemphasis{keep} or \sphinxstyleemphasis{no} = add new dump atoms to system
\sphinxstyleemphasis{label} value = field column
  field = one of the listed fields or \sphinxstyleemphasis{id} or \sphinxstyleemphasis{type}
  column = label on corresponding column in dump file
\sphinxstyleemphasis{scaled} value = \sphinxstyleemphasis{yes} or \sphinxstyleemphasis{no} = coords in dump file are scaled/unscaled
\sphinxstyleemphasis{wrapped} value = \sphinxstyleemphasis{yes} or \sphinxstyleemphasis{no} = coords in dump file are wrapped/unwrapped
\sphinxstyleemphasis{format} values = format of dump file, must be last keyword if used
  \sphinxstyleemphasis{native} = native LAMMPS dump file
  \sphinxstyleemphasis{xyz} = XYZ file
  \sphinxstyleemphasis{adios} {[}\sphinxstyleemphasis{timeout} value{]} = dump file written by the {\hyperref[\detokenize{dump_adios::doc}]{\sphinxcrossref{\DUrole{doc}{dump adios}}}} command
    \sphinxstyleemphasis{timeout} = specify waiting time for the arrival of the timestep when running concurrently.
              The value is a float number and is interpreted in seconds.
  \sphinxstyleemphasis{molfile} style path = VMD molfile plugin interface
    style = \sphinxstyleemphasis{dcd} or \sphinxstyleemphasis{xyz} or others supported by molfile plugins
    path = optional path for location of molfile plugins
\end{sphinxalltt}

\end{itemize}


\subsection{Examples}
\label{\detokenize{read_dump:examples}}
\begin{sphinxVerbatim}[commandchars=\\\{\}]
\PYG{k}{read\PYGZus{}dump}\PYG{+w}{ }\PYG{n}{dump.file}\PYG{+w}{ }\PYG{l+m}{5000}\PYG{+w}{ }\PYG{n}{x}\PYG{+w}{ }\PYG{n}{y}\PYG{+w}{ }\PYG{n}{z}
\PYG{k}{read\PYGZus{}dump}\PYG{+w}{ }\PYG{n}{dump.xyz}\PYG{+w}{ }\PYG{l+m}{5}\PYG{+w}{ }\PYG{n}{x}\PYG{+w}{ }\PYG{n}{y}\PYG{+w}{ }\PYG{n}{z}\PYG{+w}{ }\PYG{n}{box}\PYG{+w}{ }\PYG{n}{no}\PYG{+w}{ }\PYG{n}{format}\PYG{+w}{ }\PYG{n}{xyz}
\PYG{k}{read\PYGZus{}dump}\PYG{+w}{ }\PYG{n}{dump.xyz}\PYG{+w}{ }\PYG{l+m}{10}\PYG{+w}{ }\PYG{n}{x}\PYG{+w}{ }\PYG{n}{y}\PYG{+w}{ }\PYG{n}{z}\PYG{+w}{ }\PYG{n}{box}\PYG{+w}{ }\PYG{n}{no}\PYG{+w}{ }\PYG{n}{format}\PYG{+w}{ }\PYG{n}{molfile}\PYG{+w}{ }\PYG{n}{xyz}\PYG{+w}{ }\PYG{l+s}{\PYGZdq{}}\PYG{l+s}{../plugins}\PYG{l+s}{\PYGZdq{}}
\PYG{k}{read\PYGZus{}dump}\PYG{+w}{ }\PYG{n}{dump.dcd}\PYG{+w}{ }\PYG{l+m}{0}\PYG{+w}{ }\PYG{n}{x}\PYG{+w}{ }\PYG{n}{y}\PYG{+w}{ }\PYG{n}{z}\PYG{+w}{ }\PYG{n}{box}\PYG{+w}{ }\PYG{n}{yes}\PYG{+w}{ }\PYG{n}{format}\PYG{+w}{ }\PYG{n}{molfile}\PYG{+w}{ }\PYG{n}{dcd}
\PYG{k}{read\PYGZus{}dump}\PYG{+w}{ }\PYG{n}{dump.file}\PYG{+w}{ }\PYG{l+m}{1000}\PYG{+w}{ }\PYG{n}{x}\PYG{+w}{ }\PYG{n}{y}\PYG{+w}{ }\PYG{n}{z}\PYG{+w}{ }\PYG{n}{vx}\PYG{+w}{ }\PYG{n}{vy}\PYG{+w}{ }\PYG{n}{vz}\PYG{+w}{ }\PYG{n}{box}\PYG{+w}{ }\PYG{n}{yes}\PYG{+w}{ }\PYG{n}{format}\PYG{+w}{ }\PYG{n}{molfile}\PYG{+w}{ }\PYG{n}{lammpstrj}\PYG{+w}{ }\PYG{o}{/}\PYG{n}{usr}\PYG{o}{/}\PYG{n}{local}\PYG{o}{/}\PYG{n}{lib}\PYG{o}{/}\PYG{n}{vmd}\PYG{o}{/}\PYG{n}{plugins}\PYG{o}{/}\PYG{n}{LINUXAMD64}\PYG{o}{/}\PYG{n}{plugins}\PYG{o}{/}\PYG{n}{molfile}
\PYG{k}{read\PYGZus{}dump}\PYG{+w}{ }\PYG{n}{dump.file}\PYG{+w}{ }\PYG{l+m}{5000}\PYG{+w}{ }\PYG{n}{x}\PYG{+w}{ }\PYG{n}{y}\PYG{+w}{ }\PYG{n}{vx}\PYG{+w}{ }\PYG{n}{vy}\PYG{+w}{ }\PYG{n}{trim}\PYG{+w}{ }\PYG{n}{yes}
\PYG{k}{read\PYGZus{}dump}\PYG{+w}{ }\PYG{n}{dump.file}\PYG{+w}{ }\PYG{l+m}{5000}\PYG{+w}{ }\PYG{n}{x}\PYG{+w}{ }\PYG{n}{y}\PYG{+w}{ }\PYG{n}{vx}\PYG{+w}{ }\PYG{n}{vy}\PYG{+w}{ }\PYG{n}{add}\PYG{+w}{ }\PYG{n}{yes}\PYG{+w}{ }\PYG{n}{box}\PYG{+w}{ }\PYG{n}{no}\PYG{+w}{ }\PYG{n}{timestep}\PYG{+w}{ }\PYG{n}{no}
\PYG{k}{read\PYGZus{}dump}\PYG{+w}{ }\PYG{n}{..}\PYG{o}{/}\PYG{n}{run7}\PYG{o}{/}\PYG{n}{dump.file.gz}\PYG{+w}{ }\PYG{l+m}{10000}\PYG{+w}{ }\PYG{n}{x}\PYG{+w}{ }\PYG{n}{y}\PYG{+w}{ }\PYG{n}{z}\PYG{+w}{ }\PYG{n}{box}\PYG{+w}{ }\PYG{n}{yes}
\PYG{k}{read\PYGZus{}dump}\PYG{+w}{ }\PYG{n}{dump.xyz}\PYG{+w}{ }\PYG{l+m}{10}\PYG{+w}{ }\PYG{n}{x}\PYG{+w}{ }\PYG{n}{y}\PYG{+w}{ }\PYG{n}{z}\PYG{+w}{ }\PYG{n}{box}\PYG{+w}{ }\PYG{n}{no}\PYG{+w}{ }\PYG{n}{format}\PYG{+w}{ }\PYG{n}{molfile}\PYG{+w}{ }\PYG{n}{xyz}\PYG{+w}{ }\PYG{n}{..}\PYG{o}{/}\PYG{n}{plugins}
\PYG{k}{read\PYGZus{}dump}\PYG{+w}{ }\PYG{n}{dump.dcd}\PYG{+w}{ }\PYG{l+m}{0}\PYG{+w}{ }\PYG{n}{x}\PYG{+w}{ }\PYG{n}{y}\PYG{+w}{ }\PYG{n}{z}\PYG{+w}{ }\PYG{n}{format}\PYG{+w}{ }\PYG{n}{molfile}\PYG{+w}{ }\PYG{n}{dcd}
\PYG{k}{read\PYGZus{}dump}\PYG{+w}{ }\PYG{n}{dump.file}\PYG{+w}{ }\PYG{l+m}{1000}\PYG{+w}{ }\PYG{n}{x}\PYG{+w}{ }\PYG{n}{y}\PYG{+w}{ }\PYG{n}{z}\PYG{+w}{ }\PYG{n}{vx}\PYG{+w}{ }\PYG{n}{vy}\PYG{+w}{ }\PYG{n}{vz}\PYG{+w}{ }\PYG{n}{format}\PYG{+w}{ }\PYG{n}{molfile}\PYG{+w}{ }\PYG{n}{lammpstrj}\PYG{+w}{ }\PYG{o}{/}\PYG{n}{usr}\PYG{o}{/}\PYG{n}{local}\PYG{o}{/}\PYG{n}{lib}\PYG{o}{/}\PYG{n}{vmd}\PYG{o}{/}\PYG{n}{plugins}\PYG{o}{/}\PYG{n}{LINUXAMD64}\PYG{o}{/}\PYG{n}{plugins}\PYG{o}{/}\PYG{n}{molfile}
\PYG{k}{read\PYGZus{}dump}\PYG{+w}{ }\PYG{n}{dump.bp}\PYG{+w}{ }\PYG{l+m}{5000}\PYG{+w}{ }\PYG{n}{x}\PYG{+w}{ }\PYG{n}{y}\PYG{+w}{ }\PYG{n}{z}\PYG{+w}{ }\PYG{n}{vx}\PYG{+w}{ }\PYG{n}{vy}\PYG{+w}{ }\PYG{n}{vz}\PYG{+w}{ }\PYG{n}{format}\PYG{+w}{ }\PYG{n}{adios}
\PYG{k}{read\PYGZus{}dump}\PYG{+w}{ }\PYG{n}{dump.bp}\PYG{+w}{ }\PYG{l+m}{5000}\PYG{+w}{ }\PYG{n}{x}\PYG{+w}{ }\PYG{n}{y}\PYG{+w}{ }\PYG{n}{z}\PYG{+w}{ }\PYG{n}{vx}\PYG{+w}{ }\PYG{n}{vy}\PYG{+w}{ }\PYG{n}{vz}\PYG{+w}{ }\PYG{n}{format}\PYG{+w}{ }\PYG{n}{adios}\PYG{+w}{ }\PYG{n}{timeout}\PYG{+w}{ }\PYG{l+m}{60.0}
\end{sphinxVerbatim}


\subsection{Description}
\label{\detokenize{read_dump:description}}
\sphinxAtStartPar
Read atom information from a dump file to overwrite the current atom
coordinates, and optionally the atom velocities and image flags, the
simulation timestep, and the simulation box dimensions.  This is useful
for restarting a run from a particular snapshot in a dump file.  See the
{\hyperref[\detokenize{read_restart::doc}]{\sphinxcrossref{\DUrole{doc}{read\_restart}}}} and {\hyperref[\detokenize{read_data::doc}]{\sphinxcrossref{\DUrole{doc}{read\_data}}}}
commands for alternative methods to do this.  Also see the
{\hyperref[\detokenize{rerun::doc}]{\sphinxcrossref{\DUrole{doc}{rerun}}}} command for a means of reading multiple snapshots
from a dump file.

\sphinxAtStartPar
Note that a simulation box must already be defined before using the
read\_dump command.  This can be done by the
{\hyperref[\detokenize{create_box::doc}]{\sphinxcrossref{\DUrole{doc}{create\_box}}}}, {\hyperref[\detokenize{read_data::doc}]{\sphinxcrossref{\DUrole{doc}{read\_data}}}}, or
{\hyperref[\detokenize{read_restart::doc}]{\sphinxcrossref{\DUrole{doc}{read\_restart}}}} commands.  The read\_dump command can
reset the simulation box dimensions, as explained below.

\sphinxAtStartPar
Also note that reading per\sphinxhyphen{}atom information from a dump snapshot is
limited to the atom coordinates, velocities and image flags, as
explained below.  Other atom properties, which may be necessary to run
a valid simulation, such as atom charge, or bond topology information
for a molecular system, are not read from (or may not even be contained
in) dump files.  Thus this auxiliary information should be defined in
the usual way, e.g. in a data file read in by a {\hyperref[\detokenize{read_data::doc}]{\sphinxcrossref{\DUrole{doc}{read\_data}}}}
command, before using the read\_dump command, or by the {\hyperref[\detokenize{set::doc}]{\sphinxcrossref{\DUrole{doc}{set}}}}
command, after the dump snapshot is read.


\bigskip\hrule\bigskip


\sphinxAtStartPar
If the dump filename specified as \sphinxstyleemphasis{file} ends with “.gz”, the dump
file is read in gzipped format.

\sphinxAtStartPar
You can read dump files that were written (in parallel) to multiple
files via the “\%” wild\sphinxhyphen{}card character in the dump file name.  If any
specified dump file name contains a “\%”, they must all contain it.
See the {\hyperref[\detokenize{dump::doc}]{\sphinxcrossref{\DUrole{doc}{dump}}}} command for details.
The “\%” wild\sphinxhyphen{}card character is only supported by the \sphinxstyleemphasis{native} format
for dump files, described next.

\sphinxAtStartPar
If reading parallel dump files, you must also use the \sphinxstyleemphasis{nfile} keyword
to tell LAMMPS how many parallel files exist, via its specified
\sphinxstyleemphasis{Nfiles} value.

\sphinxAtStartPar
The format of the dump file is selected through the \sphinxstyleemphasis{format} keyword.
If specified, it must be the last keyword used, since all remaining
arguments are passed on to the dump reader.  The \sphinxstyleemphasis{native} format is
for native LAMMPS dump files, written with a {\hyperref[\detokenize{dump::doc}]{\sphinxcrossref{\DUrole{doc}{dump atom}}}}
or {\hyperref[\detokenize{dump::doc}]{\sphinxcrossref{\DUrole{doc}{dump custom}}}} command.  The \sphinxstyleemphasis{xyz} format is for generic XYZ
formatted dump files.  These formats take no additional values.

\sphinxAtStartPar
The \sphinxstyleemphasis{molfile} format supports reading data through using the \sphinxhref{https://www.ks.uiuc.edu/Research/vmd}{VMD}
molfile plugin interface. This dump reader format is only available,
if the MOLFILE package has been installed when compiling
LAMMPS.

\sphinxAtStartPar
The \sphinxstyleemphasis{molfile} format takes one or two additional values.  The \sphinxstyleemphasis{style}
value determines the file format to be used and can be any format that
the molfile plugins support, such as DCD or XYZ.  Note that DCD dump
files can be written by LAMMPS via the {\hyperref[\detokenize{dump::doc}]{\sphinxcrossref{\DUrole{doc}{dump dcd}}}} command.
The \sphinxstyleemphasis{path} value specifies a list of directories which LAMMPS will
search for the molfile plugins appropriate to the specified \sphinxstyleemphasis{style}.
The syntax of the \sphinxstyleemphasis{path} value is like other search paths: it can
contain multiple directories separated by a colon (or semicolon on
windows).  The \sphinxstyleemphasis{path} keyword is optional and defaults to “.”,
i.e. the current directory.

\sphinxAtStartPar
The \sphinxstyleemphasis{adios} format supports reading data that was written by the
{\hyperref[\detokenize{dump_adios::doc}]{\sphinxcrossref{\DUrole{doc}{dump adios}}}} command. The
entire dump is read in parallel across all the processes, dividing
the atoms evenly among the processes. The number of writers that
has written the dump file does not matter. Using the adios style for
dump and read\_dump is a convenient way to dump all atoms from \sphinxstyleemphasis{N}
writers and read it back by \sphinxstyleemphasis{M} readers. If one is running two
LAMMPS instances concurrently where one dumps data and the other is
reading it with the rerun command, the timeout option can be specified
to wait on the reader side for the arrival of the requested step.

\sphinxAtStartPar
Support for other dump format readers may be added in the future.


\bigskip\hrule\bigskip


\sphinxAtStartPar
Global information is first read from the dump file, namely timestep
and box information.

\sphinxAtStartPar
The dump file is scanned for a snapshot with a timestamp that matches
the specified \sphinxstyleemphasis{Nstep}.  This means the LAMMPS timestep the dump file
snapshot was written on for the \sphinxstyleemphasis{native} or \sphinxstyleemphasis{adios} formats.

\sphinxAtStartPar
The list of timestamps available in an adios .bp file is stored in the
variable \sphinxstyleemphasis{ntimestep}:

\begin{sphinxVerbatim}[commandchars=\\\{\}]
console

\PYGZdl{} bpls dump.bp \PYGZhy{}d ntimestep
  uint64\PYGZus{}t  ntimestep  5*scalar
    (0)    0 50 100 150 200
\end{sphinxVerbatim}

\sphinxAtStartPar
Note that the \sphinxstyleemphasis{xyz} and \sphinxstyleemphasis{molfile} formats do not store the timestep.
For these formats, timesteps are numbered logically, in a sequential
manner, starting from 0.  Thus to access the 10th snapshot in an \sphinxstyleemphasis{xyz}
or \sphinxstyleemphasis{mofile} formatted dump file, use \sphinxstyleemphasis{Nstep} = 9.

\sphinxAtStartPar
The dimensions of the simulation box for the selected snapshot are
also read; see the \sphinxstyleemphasis{box} keyword discussion below.  For the \sphinxstyleemphasis{native}
format, an error is generated if the snapshot is for a triclinic box
and the current simulation box is orthogonal or vice versa.  A warning
will be generated if the snapshot box boundary conditions (periodic,
shrink\sphinxhyphen{}wrapped, etc) do not match the current simulation boundary
conditions, but the boundary condition information in the snapshot is
otherwise ignored.  See the “boundary” command for more details. The
\sphinxstyleemphasis{adios} reader does the same as the \sphinxstyleemphasis{native} format reader.

\sphinxAtStartPar
For the \sphinxstyleemphasis{xyz} format, no information about the box is available, so
you must set the \sphinxstyleemphasis{box} flag to \sphinxstyleemphasis{no}.  See details below.

\sphinxAtStartPar
For the \sphinxstyleemphasis{molfile} format, reading simulation box information is
typically supported, but the location of the simulation box origin is
lost and no explicit information about periodicity or
orthogonal/triclinic box shape is available.  The MOLFILE package
makes a best effort to guess based on heuristics, but this may not
always work perfectly.


\bigskip\hrule\bigskip


\sphinxAtStartPar
Per\sphinxhyphen{}atom information from the dump file snapshot is then read from the
dump file snapshot.  This corresponds to the specified \sphinxstyleemphasis{fields} listed
in the read\_dump command.  It is an error to specify a z\sphinxhyphen{}dimension
field, namely \sphinxstyleemphasis{z}, \sphinxstyleemphasis{vz}, or \sphinxstyleemphasis{iz}, for a 2d simulation.

\sphinxAtStartPar
For dump files in \sphinxstyleemphasis{native} format, each column of per\sphinxhyphen{}atom data has a
text label listed in the file.  A matching label for each field must
appear, e.g. the label “vy” for the field \sphinxstyleemphasis{vy}.  For the \sphinxstyleemphasis{x}, \sphinxstyleemphasis{y}, \sphinxstyleemphasis{z}
fields any of the following labels are considered a match:
\begin{sphinxalltt}
x, xs, xu, xsu for field \sphinxstyleemphasis{x}
y, ys, yu, ysu for field \sphinxstyleemphasis{y}
z, zs, zu, zsu for field \sphinxstyleemphasis{z}
\end{sphinxalltt}

\sphinxAtStartPar
The meaning of xs (scaled), xu (unwrapped), and xsu (scaled and
unwrapped) is explained on the {\hyperref[\detokenize{dump::doc}]{\sphinxcrossref{\DUrole{doc}{dump}}}} command doc page.
These labels are searched for in the list of column labels in the dump
file, in order, until a match is found.

\sphinxAtStartPar
The dump file must also contain atom IDs, with a column label of “id”.

\sphinxAtStartPar
If the \sphinxstyleemphasis{add} keyword is specified with a value of \sphinxstyleemphasis{yes} or \sphinxstyleemphasis{keep}, as
discussed below, the dump file must contain atom types, with a column
label of “type”.

\sphinxAtStartPar
If a column label you want to read from the dump file is not a match
to a specified field, the \sphinxstyleemphasis{label} keyword can be used to specify the
specific column label from the dump file to associate with that field.
An example is if a time\sphinxhyphen{}averaged coordinate is written to the dump
file via the {\hyperref[\detokenize{fix_ave_atom::doc}]{\sphinxcrossref{\DUrole{doc}{fix ave/atom}}}} command.  The column
will then have a label corresponding to the fix\sphinxhyphen{}ID rather than “x” or
“xs”.  The \sphinxstyleemphasis{label} keyword can also be used to specify new column
labels for fields \sphinxstyleemphasis{id} and \sphinxstyleemphasis{type}.

\sphinxAtStartPar
For dump files in \sphinxstyleemphasis{xyz} format, only the \sphinxstyleemphasis{x}, \sphinxstyleemphasis{y}, and \sphinxstyleemphasis{z} fields are
supported.  The dump file does not store atom IDs, so these are
assigned consecutively to the atoms as they appear in the dump file,
starting from 1.  Thus you should ensure that order of atoms is
consistent from snapshot to snapshot in the XYZ dump file.  See
the {\hyperref[\detokenize{dump_modify::doc}]{\sphinxcrossref{\DUrole{doc}{dump\_modify sort}}}} command if the XYZ dump file
was written by LAMMPS.

\sphinxAtStartPar
For dump files in \sphinxstyleemphasis{molfile} format, the \sphinxstyleemphasis{x}, \sphinxstyleemphasis{y}, \sphinxstyleemphasis{z}, \sphinxstyleemphasis{vx}, \sphinxstyleemphasis{vy}, and
\sphinxstyleemphasis{vz} fields can be specified.  However, not all molfile formats store
velocities, or their respective plugins may not support reading of
velocities.  The molfile dump files do not store atom IDs, so these
are assigned consecutively to the atoms as they appear in the dump
file, starting from 1.  Thus you should ensure that order of atoms are
consistent from snapshot to snapshot in the molfile dump file.
See the {\hyperref[\detokenize{dump_modify::doc}]{\sphinxcrossref{\DUrole{doc}{dump\_modify sort}}}} command if the dump file
was written by LAMMPS.

\sphinxAtStartPar
The \sphinxstyleemphasis{adios} format supports all fields that the \sphinxstyleemphasis{native} format supports
except for the \sphinxstyleemphasis{q} charge field.
The list of fields stored in an adios .bp file is recorded in the attributes
\sphinxstyleemphasis{columns} (array of short strings) and \sphinxstyleemphasis{columnstr} (space\sphinxhyphen{}separated single string).

\begin{sphinxVerbatim}[commandchars=\\\{\}]
console

\PYGZdl{} bpls \PYGZhy{}la dump.bp column*
  string    columns            attr   = \PYGZob{}\PYGZdq{}id\PYGZdq{}, \PYGZdq{}type\PYGZdq{}, \PYGZdq{}x\PYGZdq{}, \PYGZdq{}y\PYGZdq{}, \PYGZdq{}z\PYGZdq{}, \PYGZdq{}vx\PYGZdq{}, \PYGZdq{}vy\PYGZdq{}, \PYGZdq{}vz\PYGZdq{}\PYGZcb{}
  string    columnstr          attr   = \PYGZdq{}id type x y z vx vy vz \PYGZdq{}
\end{sphinxVerbatim}


\bigskip\hrule\bigskip


\sphinxAtStartPar
Information from the dump file snapshot is used to overwrite or
replace properties of the current system.  There are various options
for how this is done, determined by the specified fields and optional
keywords.

\sphinxAtStartPar
\DUrole{versionmodified,changed}{Changed in version 3Aug2022.}

\sphinxAtStartPar
The timestep of the snapshot becomes the current timestep for the
simulation unless the \sphinxstyleemphasis{timestep} keyword is specified with a \sphinxstyleemphasis{no} value
(default setting is \sphinxstyleemphasis{yes}).  See the {\hyperref[\detokenize{reset_timestep::doc}]{\sphinxcrossref{\DUrole{doc}{reset\_timestep}}}}
command if you wish to change this to a different value after the dump
snapshot is read.

\sphinxAtStartPar
If the \sphinxstyleemphasis{box} keyword is specified with a \sphinxstyleemphasis{yes} value, then the current
simulation box dimensions are replaced by the dump snapshot box
dimensions.  If the \sphinxstyleemphasis{box} keyword is specified with a \sphinxstyleemphasis{no} value, the
current simulation box is unchanged.

\sphinxAtStartPar
If the \sphinxstyleemphasis{purge} keyword is specified with a \sphinxstyleemphasis{yes} value, then all
current atoms in the system are deleted before any of the operations
invoked by the \sphinxstyleemphasis{replace}, \sphinxstyleemphasis{trim}, or \sphinxstyleemphasis{add} keywords take place.

\sphinxAtStartPar
If the \sphinxstyleemphasis{replace} keyword is specified with a \sphinxstyleemphasis{yes} value, then atoms
with IDs that are in both the current system and the dump snapshot
have their properties overwritten by field values.  If the \sphinxstyleemphasis{replace}
keyword is specified with a \sphinxstyleemphasis{no} value, atoms with IDs that are in
both the current system and the dump snapshot are not modified.

\sphinxAtStartPar
If the \sphinxstyleemphasis{trim} keyword is specified with a \sphinxstyleemphasis{yes} value, then atoms with
IDs that are in the current system but not in the dump snapshot are
deleted.  These atoms are unaffected if the \sphinxstyleemphasis{trim} keyword is
specified with a \sphinxstyleemphasis{no} value.

\sphinxAtStartPar
If the \sphinxstyleemphasis{add} keyword is specified with a \sphinxstyleemphasis{no} value (default), then
dump file atoms with IDs that are not in the current system are not
added to the system.  They are simply ignored.

\sphinxAtStartPar
If a \sphinxstyleemphasis{yes} value is specified, the atoms with new IDs are added to the
system but their atom IDs are not preserved.  Instead, after all the
atoms are added, new IDs are assigned to them in the same manner as is
described for the {\hyperref[\detokenize{create_atoms::doc}]{\sphinxcrossref{\DUrole{doc}{create\_atoms}}}} command.  Basically
the largest existing atom ID in the system is identified, and all the
added atoms are assigned IDs that consecutively follow the largest ID.

\sphinxAtStartPar
If a \sphinxstyleemphasis{keep} value is specified, the atoms with new IDs are added to
the system and their atom IDs are preserved.  This may lead to
non\sphinxhyphen{}contiguous IDs for the combined system.

\sphinxAtStartPar
Note that atoms added via the \sphinxstyleemphasis{add} keyword will only have the
attributes read from the dump file due to the \sphinxstyleemphasis{field} arguments.  For
example, if \sphinxstyleemphasis{x} or \sphinxstyleemphasis{y} or \sphinxstyleemphasis{z} or \sphinxstyleemphasis{q} is not specified as a field, a
value of 0.0 is used for added atoms.  Added atoms must have an atom
type, so this value must appear in the dump file.

\sphinxAtStartPar
Any other attributes (e.g. charge or particle diameter for spherical
particles) will be set to default values, the same as if the
{\hyperref[\detokenize{create_atoms::doc}]{\sphinxcrossref{\DUrole{doc}{create\_atoms}}}} command were used.


\bigskip\hrule\bigskip


\sphinxAtStartPar
Atom coordinates read from the dump file are first converted into
unscaled coordinates, relative to the box dimensions of the snapshot.
These coordinates are then be assigned to an existing or new atom in
the current simulation.  The coordinates will then be remapped to the
simulation box, whether it is the original box or the dump snapshot
box.  If periodic boundary conditions apply, this means the atom will
be remapped back into the simulation box if necessary.  If shrink\sphinxhyphen{}wrap
boundary conditions apply, the new coordinates may change the
simulation box dimensions.  If fixed boundary conditions apply, the
atom will be lost if it is outside the simulation box.

\sphinxAtStartPar
For \sphinxstyleemphasis{native} format dump files, the 3 xyz image flags for an atom in
the dump file are set to the corresponding values appearing in the
dump file if the \sphinxstyleemphasis{ix}, \sphinxstyleemphasis{iy}, \sphinxstyleemphasis{iz} fields are specified.  If not
specified, the image flags for replaced atoms are not changed and
image flags for new atoms are set to default values.  If coordinates
read from the dump file are in unwrapped format (e.g. \sphinxstyleemphasis{xu}) then the
image flags for read\sphinxhyphen{}in atoms are also set to default values.  The
remapping procedure described in the previous paragraph will then
change images flags for all atoms (old and new) if periodic boundary
conditions are applied to remap an atom back into the simulation box.

\begin{sphinxadmonition}{note}{Note:}
\sphinxAtStartPar
If you get a warning about inconsistent image flags after
reading in a dump snapshot, it means one or more pairs of bonded atoms
now have inconsistent image flags.  As discussed on the {\hyperref[\detokenize{Errors_common::doc}]{\sphinxcrossref{\DUrole{doc}{Errors common}}}} page this may or may not cause problems
for subsequent simulations.  One way this can happen is if you read
image flag fields from the dump file but do not also use the dump file
box parameters.
\end{sphinxadmonition}

\sphinxAtStartPar
LAMMPS knows how to compute unscaled and remapped coordinates for the
snapshot column labels discussed above, e.g. \sphinxstyleemphasis{x}, \sphinxstyleemphasis{xs}, \sphinxstyleemphasis{xu}, \sphinxstyleemphasis{xsu}.
If another column label is assigned to the \sphinxstyleemphasis{x} or \sphinxstyleemphasis{y} or \sphinxstyleemphasis{z} field via
the \sphinxstyleemphasis{label} keyword, e.g. for coordinates output by the {\hyperref[\detokenize{fix_ave_atom::doc}]{\sphinxcrossref{\DUrole{doc}{fix ave/atom}}}} command, then LAMMPS needs to know whether
the coordinate information in the dump file is scaled and/or wrapped.
This can be set via the \sphinxstyleemphasis{scaled} and \sphinxstyleemphasis{wrapped} keywords.  Note that
the value of the \sphinxstyleemphasis{scaled} and \sphinxstyleemphasis{wrapped} keywords is ignored for fields
\sphinxstyleemphasis{x} or \sphinxstyleemphasis{y} or \sphinxstyleemphasis{z} if the \sphinxstyleemphasis{label} keyword is not used to assign a
column label to that field.

\sphinxAtStartPar
The scaled/unscaled and wrapped/unwrapped setting must be identical
for any of the \sphinxstyleemphasis{x}, \sphinxstyleemphasis{y}, \sphinxstyleemphasis{z} fields that are specified.  Thus you
cannot read \sphinxstyleemphasis{xs} and \sphinxstyleemphasis{yu} from the dump file.  Also, if the dump file
coordinates are scaled and the simulation box is triclinic, then all 3
of the \sphinxstyleemphasis{x}, \sphinxstyleemphasis{y}, \sphinxstyleemphasis{z} fields must be specified, since they are all
needed to generate absolute, unscaled coordinates.


\bigskip\hrule\bigskip



\subsection{Restrictions}
\label{\detokenize{read_dump:restrictions}}
\sphinxAtStartPar
To read gzipped dump files, you must compile LAMMPS with the
\sphinxhyphen{}DLAMMPS\_GZIP option.  See the {\hyperref[\detokenize{Build_settings::doc}]{\sphinxcrossref{\DUrole{doc}{Build settings}}}}
doc page for details.

\sphinxAtStartPar
The \sphinxstyleemphasis{molfile} dump file formats are part of the MOLFILE package.
They are only enabled if LAMMPS was built with that packages.  See the
{\hyperref[\detokenize{Build_package::doc}]{\sphinxcrossref{\DUrole{doc}{Build package}}}} page for more info.

\sphinxAtStartPar
To write and read adios .bp files, you must compile LAMMPS with the
{\hyperref[\detokenize{Packages_details:pkg-adios}]{\sphinxcrossref{\DUrole{std,std-ref}{ADIOS}}}} package.


\subsection{Related commands}
\label{\detokenize{read_dump:related-commands}}
\sphinxAtStartPar
{\hyperref[\detokenize{dump::doc}]{\sphinxcrossref{\DUrole{doc}{dump}}}}, {\hyperref[\detokenize{dump_molfile::doc}]{\sphinxcrossref{\DUrole{doc}{dump molfile}}}},
{\hyperref[\detokenize{dump_adios::doc}]{\sphinxcrossref{\DUrole{doc}{dump adios}}}},
{\hyperref[\detokenize{read_data::doc}]{\sphinxcrossref{\DUrole{doc}{read\_data}}}}, {\hyperref[\detokenize{read_restart::doc}]{\sphinxcrossref{\DUrole{doc}{read\_restart}}}},
{\hyperref[\detokenize{rerun::doc}]{\sphinxcrossref{\DUrole{doc}{rerun}}}}


\subsection{Default}
\label{\detokenize{read_dump:default}}
\sphinxAtStartPar
The option defaults are box = yes, timestep = yes, replace = yes, purge = no,
trim = no, add = no, scaled = no, wrapped = yes, and format = native.

\sphinxstepscope

\index{read\_restart@\spxentry{read\_restart}}\ignorespaces 

\section{read\_restart command}
\label{\detokenize{read_restart:read-restart-command}}\label{\detokenize{read_restart:index-0}}\label{\detokenize{read_restart::doc}}

\subsection{Syntax}
\label{\detokenize{read_restart:syntax}}
\begin{sphinxVerbatim}[commandchars=\\\{\}]
\PYG{k}{read\PYGZus{}restart}\PYG{+w}{ }\PYG{n}{file}
\end{sphinxVerbatim}
\begin{itemize}
\item {} 
\sphinxAtStartPar
file = name of binary restart file to read in

\end{itemize}


\subsection{Examples}
\label{\detokenize{read_restart:examples}}
\begin{sphinxVerbatim}[commandchars=\\\{\}]
\PYG{k}{read\PYGZus{}restart}\PYG{+w}{ }\PYG{n}{save.10000}
\PYG{k}{read\PYGZus{}restart}\PYG{+w}{ }\PYG{n}{restart.}\PYG{o}{*}
\end{sphinxVerbatim}


\subsection{Description}
\label{\detokenize{read_restart:description}}
\sphinxAtStartPar
Read in a previously saved system configuration from a restart file.
This allows continuation of a previous run.  Details about what
information is stored (and not stored) in a restart file is given below.
Basically this operation will re\sphinxhyphen{}create the simulation box with all its
atoms and their attributes as well as some related global settings, at
the point in time it was written to the restart file by a previous
simulation.  The simulation box will be partitioned into a regular 3d
grid of rectangular bricks, one per processor, based on the number of
processors in the current simulation and the settings of the
{\hyperref[\detokenize{processors::doc}]{\sphinxcrossref{\DUrole{doc}{processors}}}} command.  The partitioning can later be
changed by the {\hyperref[\detokenize{balance::doc}]{\sphinxcrossref{\DUrole{doc}{balance}}}} or {\hyperref[\detokenize{fix_balance::doc}]{\sphinxcrossref{\DUrole{doc}{fix balance}}}} commands.

\sphinxAtStartPar
\DUrole{versionmodified,deprecated}{Deprecated since version 23Jun2022.}

\sphinxAtStartPar
Atom coordinates that are found to be outside the simulation box when
reading the restart will be remapped back into the box and their image
flags updated accordingly.  This previously required specifying the
\sphinxstyleemphasis{remap} option, but that is no longer required.

\sphinxAtStartPar
Restart files are saved in binary format to enable exact restarts,
meaning that the trajectories of a restarted run will precisely match
those produced by the original run had it continued on.

\sphinxAtStartPar
Some information about a restart file can be gathered directly from the
command line when using LAMMPS with the {\hyperref[\detokenize{Run_options:restart2info}]{\sphinxcrossref{\DUrole{std,std-ref}{\sphinxhyphen{}restart2info}}}} command line flag.  On Unix\sphinxhyphen{}like operating systems (like
Linux or macOS), one can also {\hyperref[\detokenize{Tools:magic}]{\sphinxcrossref{\DUrole{std,std-ref}{configure the “file” command line
program}}}} to display basic information about a restart file

\sphinxAtStartPar
The binary restart file format was not designed with backward, forward,
or cross\sphinxhyphen{}platform compatibility in mind, so the files are only expected
to be read correctly by the same LAMMPS executable on the same platform.
Changes to the architecture, compilation settings, or LAMMPS version can
render a restart file unreadable or it may read the data incorrectly.
If you want a more portable format, you can use the data file format as
created by the {\hyperref[\detokenize{write_data::doc}]{\sphinxcrossref{\DUrole{doc}{write\_data}}}} command.  Binary restart
files can also be converted into a data file from the command line by
the LAMMPS executable that wrote them using the {\hyperref[\detokenize{Run_options:restart2data}]{\sphinxcrossref{\DUrole{std,std-ref}{\sphinxhyphen{}restart2data}}}} command line flag.

\sphinxAtStartPar
Several things can prevent exact restarts due to round\sphinxhyphen{}off effects, in
which case the trajectories in the 2 runs will slowly diverge.  These
include running on a different number of processors or changing
certain settings such as those set by the {\hyperref[\detokenize{newton::doc}]{\sphinxcrossref{\DUrole{doc}{newton}}}} or
{\hyperref[\detokenize{processors::doc}]{\sphinxcrossref{\DUrole{doc}{processors}}}} commands.  LAMMPS will issue a warning in
these cases.

\sphinxAtStartPar
Certain fixes will not restart exactly, though they should provide
statistically similar results.  These include {\hyperref[\detokenize{fix_shake::doc}]{\sphinxcrossref{\DUrole{doc}{fix shake}}}} and {\hyperref[\detokenize{fix_langevin::doc}]{\sphinxcrossref{\DUrole{doc}{fix langevin}}}}.

\sphinxAtStartPar
Certain pair styles will not restart exactly, though they should
provide statistically similar results.  This is because the forces
they compute depend on atom velocities, which are used at half\sphinxhyphen{}step
values every timestep when forces are computed.  When a run restarts,
forces are initially evaluated with a full\sphinxhyphen{}step velocity, which is
different than if the run had continued.  These pair styles include
{\hyperref[\detokenize{pair_gran::doc}]{\sphinxcrossref{\DUrole{doc}{granular pair styles}}}}, {\hyperref[\detokenize{pair_dpd::doc}]{\sphinxcrossref{\DUrole{doc}{pair dpd}}}}, and
{\hyperref[\detokenize{pair_lubricate::doc}]{\sphinxcrossref{\DUrole{doc}{pair lubricate}}}}.

\sphinxAtStartPar
If a restarted run is immediately different than the run which
produced the restart file, it could be a LAMMPS bug, so consider
{\hyperref[\detokenize{Errors_bugs::doc}]{\sphinxcrossref{\DUrole{doc}{reporting it}}}} if you think the behavior is a bug.

\sphinxAtStartPar
Because restart files are binary, they may not be portable to other
machines.  In this case, you can use the {\hyperref[\detokenize{Run_options::doc}]{\sphinxcrossref{\DUrole{doc}{\sphinxhyphen{}restart command\sphinxhyphen{}line
switch}}}} to convert a restart file to a data file.

\sphinxAtStartPar
Similar to how restart files are written (see the {\hyperref[\detokenize{write_restart::doc}]{\sphinxcrossref{\DUrole{doc}{write\_restart}}}} and {\hyperref[\detokenize{restart::doc}]{\sphinxcrossref{\DUrole{doc}{restart}}}} commands), the restart
filename can contain two wild\sphinxhyphen{}card characters.  If a “*” appears in the
filename, the directory is searched for all filenames that match the
pattern where “*” is replaced with a timestep value.  The file with the
largest timestep value is read in.  Thus, this effectively means, read
the latest restart file.  It’s useful if you want your script to
continue a run from where it left off.  See the {\hyperref[\detokenize{run::doc}]{\sphinxcrossref{\DUrole{doc}{run}}}} command
and its “upto” option for how to specify the run command so it does not
need to be changed either.

\sphinxAtStartPar
If a “\%” character appears in the restart filename, LAMMPS expects a
set of multiple files to exist.  The {\hyperref[\detokenize{restart::doc}]{\sphinxcrossref{\DUrole{doc}{restart}}}} and
{\hyperref[\detokenize{write_restart::doc}]{\sphinxcrossref{\DUrole{doc}{write\_restart}}}} commands explain how such sets are
created.  Read\_restart will first read a filename where “\%” is
replaced by “base”.  This file tells LAMMPS how many processors
created the set and how many files are in it.  Read\_restart then reads
the additional files.  For example, if the restart file was specified
as save.\% when it was written, then read\_restart reads the files
save.base, save.0, save.1, … save.P\sphinxhyphen{}1, where P is the number of
processors that created the restart file.

\sphinxAtStartPar
Note that P could be the total number of processors in the previous
simulation, or some subset of those processors, if the \sphinxstyleemphasis{fileper} or
\sphinxstyleemphasis{nfile} options were used when the restart file was written; see the
{\hyperref[\detokenize{restart::doc}]{\sphinxcrossref{\DUrole{doc}{restart}}}} and {\hyperref[\detokenize{write_restart::doc}]{\sphinxcrossref{\DUrole{doc}{write\_restart}}}} commands
for details.  The processors in the current LAMMPS simulation share
the work of reading these files; each reads a roughly equal subset of
the files.  The number of processors which created the set can be
different the number of processors in the current LAMMPS simulation.
This can be a fast mode of input on parallel machines that support
parallel I/O.


\bigskip\hrule\bigskip


\sphinxAtStartPar
Here is the list of information included in a restart file, which
means these quantities do not need to be re\sphinxhyphen{}specified in the input
script that reads the restart file, though you can redefine many of
these settings after the restart file is read.
\begin{itemize}
\item {} 
\sphinxAtStartPar
{\hyperref[\detokenize{units::doc}]{\sphinxcrossref{\DUrole{doc}{units}}}}

\item {} 
\sphinxAtStartPar
{\hyperref[\detokenize{newton::doc}]{\sphinxcrossref{\DUrole{doc}{newton bond}}}} (see discussion of newton command below)

\item {} 
\sphinxAtStartPar
{\hyperref[\detokenize{atom_style::doc}]{\sphinxcrossref{\DUrole{doc}{atom style}}}} and {\hyperref[\detokenize{atom_modify::doc}]{\sphinxcrossref{\DUrole{doc}{atom\_modify}}}} settings id, map, sort

\item {} 
\sphinxAtStartPar
{\hyperref[\detokenize{comm_style::doc}]{\sphinxcrossref{\DUrole{doc}{comm style}}}} and {\hyperref[\detokenize{comm_modify::doc}]{\sphinxcrossref{\DUrole{doc}{comm\_modify}}}} settings mode, cutoff, vel

\item {} 
\sphinxAtStartPar
{\hyperref[\detokenize{timestep::doc}]{\sphinxcrossref{\DUrole{doc}{timestep size}}}} and {\hyperref[\detokenize{reset_timestep::doc}]{\sphinxcrossref{\DUrole{doc}{timestep number}}}}

\item {} 
\sphinxAtStartPar
simulation box size and shape and {\hyperref[\detokenize{boundary::doc}]{\sphinxcrossref{\DUrole{doc}{boundary}}}} settings

\item {} 
\sphinxAtStartPar
atom {\hyperref[\detokenize{group::doc}]{\sphinxcrossref{\DUrole{doc}{group}}}} definitions

\item {} 
\sphinxAtStartPar
per\sphinxhyphen{}type atom settings such as {\hyperref[\detokenize{mass::doc}]{\sphinxcrossref{\DUrole{doc}{mass}}}}

\item {} 
\sphinxAtStartPar
per\sphinxhyphen{}atom attributes including their group assignments and molecular topology attributes (bonds, angles, etc)

\item {} 
\sphinxAtStartPar
force field styles ({\hyperref[\detokenize{pair_style::doc}]{\sphinxcrossref{\DUrole{doc}{pair}}}}, {\hyperref[\detokenize{bond_style::doc}]{\sphinxcrossref{\DUrole{doc}{bond}}}}, {\hyperref[\detokenize{angle_style::doc}]{\sphinxcrossref{\DUrole{doc}{angle}}}}, etc)

\item {} 
\sphinxAtStartPar
force field coefficients ({\hyperref[\detokenize{pair_coeff::doc}]{\sphinxcrossref{\DUrole{doc}{pair}}}}, {\hyperref[\detokenize{bond_coeff::doc}]{\sphinxcrossref{\DUrole{doc}{bond}}}}, {\hyperref[\detokenize{angle_coeff::doc}]{\sphinxcrossref{\DUrole{doc}{angle}}}}, etc) in some cases (see below)

\item {} 
\sphinxAtStartPar
{\hyperref[\detokenize{pair_modify::doc}]{\sphinxcrossref{\DUrole{doc}{pair\_modify}}}} settings, except the compute option

\item {} 
\sphinxAtStartPar
{\hyperref[\detokenize{special_bonds::doc}]{\sphinxcrossref{\DUrole{doc}{special\_bonds}}}} settings

\end{itemize}

\sphinxAtStartPar
Here is a list of information not stored in a restart file, which
means you must re\sphinxhyphen{}issue these commands in your input script, after
reading the restart file.
\begin{itemize}
\item {} 
\sphinxAtStartPar
{\hyperref[\detokenize{newton::doc}]{\sphinxcrossref{\DUrole{doc}{newton pair}}}} (see discussion of newton command below)

\item {} 
\sphinxAtStartPar
{\hyperref[\detokenize{fix::doc}]{\sphinxcrossref{\DUrole{doc}{fix}}}} commands (see below)

\item {} 
\sphinxAtStartPar
{\hyperref[\detokenize{compute::doc}]{\sphinxcrossref{\DUrole{doc}{compute}}}} commands (see below)

\item {} 
\sphinxAtStartPar
{\hyperref[\detokenize{variable::doc}]{\sphinxcrossref{\DUrole{doc}{variable}}}} commands

\item {} 
\sphinxAtStartPar
{\hyperref[\detokenize{region::doc}]{\sphinxcrossref{\DUrole{doc}{region}}}} commands

\item {} 
\sphinxAtStartPar
{\hyperref[\detokenize{neighbor::doc}]{\sphinxcrossref{\DUrole{doc}{neighbor list}}}} criteria including {\hyperref[\detokenize{neigh_modify::doc}]{\sphinxcrossref{\DUrole{doc}{neigh\_modify}}}} settings

\item {} 
\sphinxAtStartPar
{\hyperref[\detokenize{kspace_style::doc}]{\sphinxcrossref{\DUrole{doc}{kspace\_style}}}} and {\hyperref[\detokenize{kspace_modify::doc}]{\sphinxcrossref{\DUrole{doc}{kspace\_modify}}}} settings

\item {} 
\sphinxAtStartPar
info for {\hyperref[\detokenize{thermo_style::doc}]{\sphinxcrossref{\DUrole{doc}{thermodynamic}}}}, {\hyperref[\detokenize{dump::doc}]{\sphinxcrossref{\DUrole{doc}{dump}}}}, or {\hyperref[\detokenize{restart::doc}]{\sphinxcrossref{\DUrole{doc}{restart}}}} output

\end{itemize}

\sphinxAtStartPar
The {\hyperref[\detokenize{newton::doc}]{\sphinxcrossref{\DUrole{doc}{newton}}}} command has two settings, one for pairwise
interactions, the other for bonded.  Both settings are stored in the
restart file.  For the bond setting, the value in the file will
overwrite the current value (at the time the read\_restart command is
issued) and warn if the two values are not the same and the current
value is not the default.  For the pair setting, the value in the file
will not overwrite the current value (so that you can override the
previous run’s value), but a warning is issued if the two values are
not the same and the current value is not the default.

\sphinxAtStartPar
Note that some force field styles (pair, bond, angle, etc) do not
store their coefficient info in restart files.  Typically these are
many\sphinxhyphen{}body or tabulated potentials which read their parameters from
separate files.  In these cases you will need to re\sphinxhyphen{}specify the
{\hyperref[\detokenize{pair_coeff::doc}]{\sphinxcrossref{\DUrole{doc}{pair\_coeff}}}}, {\hyperref[\detokenize{bond_coeff::doc}]{\sphinxcrossref{\DUrole{doc}{bond\_coeff}}}}, etc
commands in your restart input script.  The doc pages for individual
force field styles mention if this is the case.  This is also true of
{\hyperref[\detokenize{pair_hybrid::doc}]{\sphinxcrossref{\DUrole{doc}{pair\_style hybrid}}}} (bond hybrid, angle hybrid, etc)
commands; they do not store coefficient info.

\sphinxAtStartPar
As indicated in the above list, the {\hyperref[\detokenize{fix::doc}]{\sphinxcrossref{\DUrole{doc}{fixes}}}} used for a
simulation are not stored in the restart file.  This means the new
input script should specify all fixes it will use.  However, note that
some fixes store an internal “state” which is written to the restart
file.  This allows the fix to continue on with its calculations in a
restarted simulation.  To re\sphinxhyphen{}enable such a fix, the fix command in the
new input script must be of the same style and use the same fix\sphinxhyphen{}ID as
was used in the input script that wrote the restart file.

\sphinxAtStartPar
If a match is found, LAMMPS prints a message indicating that the fix
is being re\sphinxhyphen{}enabled.  If no match is found before the first run or
minimization is performed by the new script, the “state” information
for the saved fix is discarded.  At the time the discard occurs,
LAMMPS will also print a list of fixes for which the information is
being discarded.  See the doc pages for individual fixes for info on
which ones can be restarted in this manner.  Note that fixes which are
created internally by other LAMMPS commands (computes, fixes, etc)
will have style names which are all\sphinxhyphen{}capitalized, and IDs which are
generated internally.

\sphinxAtStartPar
Likewise, the {\hyperref[\detokenize{fix::doc}]{\sphinxcrossref{\DUrole{doc}{computes}}}} used for a simulation are not stored
in the restart file.  This means the new input script should specify
all computes it will use.  However, some computes create a fix
internally to store “state” information that persists from timestep to
timestep.  An example is the {\hyperref[\detokenize{compute_msd::doc}]{\sphinxcrossref{\DUrole{doc}{compute msd}}}} command
which uses a fix to store a reference coordinate for each atom, so
that a displacement can be calculated at any later time.  If the
compute command in the new input script uses the same compute\sphinxhyphen{}ID and
group\sphinxhyphen{}ID as was used in the input script that wrote the restart file,
then it will create the same fix in the restarted run.  This means the
re\sphinxhyphen{}created fix will be re\sphinxhyphen{}enabled with the stored state information as
described in the previous paragraph, so that the compute can continue
its calculations in a consistent manner.

\begin{sphinxadmonition}{note}{Note:}
\sphinxAtStartPar
There are a handful of commands which can be used before or between
runs which may require a system initialization.  Examples include the
“balance”, “displace\_atoms”, “delete\_atoms”, “set” (some options),
and “velocity” (some options) commands.  This is because they can
migrate atoms to new processors.  Thus they will also discard unused
“state” information from fixes.  You will know the discard has
occurred because a list of discarded fixes will be printed to the
screen and log file, as explained above.  This means that if you wish
to retain that info in a restarted run, you must re\sphinxhyphen{}specify the
relevant fixes and computes (which create fixes) before those
commands are used.
\end{sphinxadmonition}

\sphinxAtStartPar
Some pair styles, like the {\hyperref[\detokenize{pair_gran::doc}]{\sphinxcrossref{\DUrole{doc}{granular pair styles}}}}, also
use a fix to store “state” information that persists from timestep to
timestep.  In the case of granular potentials, it is contact
information between pairs of touching particles.  This info will also
be re\sphinxhyphen{}enabled in the restart script, assuming you re\sphinxhyphen{}use the same
granular pair style.

\sphinxAtStartPar
LAMMPS allows bond interactions (angle, etc) to be turned off or
deleted in various ways, which can affect how their info is stored in
a restart file.

\sphinxAtStartPar
If bonds (angles, etc) have been turned off by the {\hyperref[\detokenize{fix_shake::doc}]{\sphinxcrossref{\DUrole{doc}{fix shake}}}} or {\hyperref[\detokenize{delete_bonds::doc}]{\sphinxcrossref{\DUrole{doc}{delete\_bonds}}}} command, their info
will be written to a restart file as if they are turned on.  This means
they will need to be turned off again in a new run after the restart
file is read.

\sphinxAtStartPar
Bonds that are broken (e.g. by a bond\sphinxhyphen{}breaking potential) are written
to the restart file as broken bonds with a type of 0.  Thus these
bonds will still be broken when the restart file is read.

\sphinxAtStartPar
Bonds that have been broken by the {\hyperref[\detokenize{fix_bond_break::doc}]{\sphinxcrossref{\DUrole{doc}{fix bond/break}}}} command have disappeared from the system.  No
information about these bonds is written to the restart file.


\bigskip\hrule\bigskip



\subsection{Restrictions}
\label{\detokenize{read_restart:restrictions}}
\sphinxAtStartPar
none


\subsection{Related commands}
\label{\detokenize{read_restart:related-commands}}
\sphinxAtStartPar
{\hyperref[\detokenize{read_data::doc}]{\sphinxcrossref{\DUrole{doc}{read\_data}}}}, {\hyperref[\detokenize{read_dump::doc}]{\sphinxcrossref{\DUrole{doc}{read\_dump}}}},
{\hyperref[\detokenize{write_restart::doc}]{\sphinxcrossref{\DUrole{doc}{write\_restart}}}}, {\hyperref[\detokenize{restart::doc}]{\sphinxcrossref{\DUrole{doc}{restart}}}}


\subsection{Default}
\label{\detokenize{read_restart:default}}
\sphinxAtStartPar
none

\sphinxstepscope

\index{region@\spxentry{region}}\ignorespaces 

\section{region command}
\label{\detokenize{region:region-command}}\label{\detokenize{region:index-0}}\label{\detokenize{region::doc}}

\subsection{Syntax}
\label{\detokenize{region:syntax}}
\begin{sphinxVerbatim}[commandchars=\\\{\}]
\PYG{k}{region }\PYG{n+nv+nvIdentifier}{ID}\PYG{+w}{ }\PYG{n}{style}\PYG{+w}{ }\PYG{n}{args}\PYG{+w}{ }\PYG{n}{keyword}\PYG{+w}{ }\PYG{n}{arg}\PYG{+w}{ }\PYG{n}{...}
\end{sphinxVerbatim}
\begin{itemize}
\item {} 
\sphinxAtStartPar
ID = user\sphinxhyphen{}assigned name for the region

\item {} 
\sphinxAtStartPar
style = \sphinxstyleemphasis{delete} or \sphinxstyleemphasis{block} or \sphinxstyleemphasis{cone} or \sphinxstyleemphasis{cylinder} or \sphinxstyleemphasis{ellipsoid} or \sphinxstyleemphasis{plane} or \sphinxstyleemphasis{prism} or \sphinxstyleemphasis{sphere} or \sphinxstyleemphasis{union} or \sphinxstyleemphasis{intersect}
\begin{sphinxalltt}
\sphinxstyleemphasis{delete} = no args
\sphinxstyleemphasis{block} args = xlo xhi ylo yhi zlo zhi
  xlo,xhi,ylo,yhi,zlo,zhi = bounds of block in all dimensions (distance units)
  xlo,xhi,ylo,yhi,zlo,zhi can be a variable
\sphinxstyleemphasis{cone} args = dim c1 c2 radlo radhi lo hi
  dim = \sphinxstyleemphasis{x} or \sphinxstyleemphasis{y} or \sphinxstyleemphasis{z} = axis of cone
  c1,c2 = coords of cone axis in other 2 dimensions (distance units)
  radlo,radhi = cone radii at lo and hi end (distance units)
  lo,hi = bounds of cone in dim (distance units)
  c1,c2,radlo,radhi,lo,hi can be a variable (see below)
\sphinxstyleemphasis{cylinder} args = dim c1 c2 radius lo hi
  dim = \sphinxstyleemphasis{x} or \sphinxstyleemphasis{y} or \sphinxstyleemphasis{z} = axis of cylinder
  c1,c2 = coords of cylinder axis in other 2 dimensions (distance units)
  radius = cylinder radius (distance units)
    c1,c2, and radius can be a variable (see below)
  lo,hi = bounds of cylinder in dim (distance units)
\sphinxstyleemphasis{ellipsoid} args = x y z a b c
  x,y,z = center of ellipsoid (distance units)
  a,b,c = half the length of the principal axes of the ellipsoid (distance units)
    x,y,z,a,b and c can be a variable (see below)
\sphinxstyleemphasis{plane} args = px py pz nx ny nz
  px,py,pz = point on the plane (distance units)
  nx,ny,nz = direction normal to plane (distance units)
\sphinxstyleemphasis{prism} args = xlo xhi ylo yhi zlo zhi xy xz yz
  xlo,xhi,ylo,yhi,zlo,zhi = bounds of untilted prism (distance units)
  xy = distance to tilt y in x direction (distance units)
  xz = distance to tilt z in x direction (distance units)
  yz = distance to tilt z in y direction (distance units)
\sphinxstyleemphasis{sphere} args = x y z radius
  x,y,z = center of sphere (distance units)
  radius = radius of sphere (distance units)
    x,y,z, and radius can be a variable (see below)
\sphinxstyleemphasis{union} args = N reg\sphinxhyphen{}ID1 reg\sphinxhyphen{}ID2 ...
  N = \# of regions to follow, must be 2 or greater
  reg\sphinxhyphen{}ID1,reg\sphinxhyphen{}ID2, ... = IDs of regions to join together
\sphinxstyleemphasis{intersect} args = N reg\sphinxhyphen{}ID1 reg\sphinxhyphen{}ID2 ...
  N = \# of regions to follow, must be 2 or greater
  reg\sphinxhyphen{}ID1,reg\sphinxhyphen{}ID2, ... = IDs of regions to intersect
\end{sphinxalltt}

\item {} 
\sphinxAtStartPar
zero or more keyword/arg pairs may be appended

\item {} 
\sphinxAtStartPar
keyword = \sphinxstyleemphasis{side} or \sphinxstyleemphasis{units} or \sphinxstyleemphasis{move} or \sphinxstyleemphasis{rotate} or \sphinxstyleemphasis{open}
\begin{sphinxalltt}
\sphinxstyleemphasis{side} value = \sphinxstyleemphasis{in} or \sphinxstyleemphasis{out}
  \sphinxstyleemphasis{in} = the region is inside the specified geometry
  \sphinxstyleemphasis{out} = the region is outside the specified geometry
\sphinxstyleemphasis{units} value = \sphinxstyleemphasis{lattice} or \sphinxstyleemphasis{box}
  \sphinxstyleemphasis{lattice} = the geometry is defined in lattice units
  \sphinxstyleemphasis{box} = the geometry is defined in simulation box units
\sphinxstyleemphasis{move} args = v\_x v\_y v\_z
  v\_x,v\_y,v\_z = equal\sphinxhyphen{}style variables for x,y,z displacement of region over time (distance units)
\sphinxstyleemphasis{rotate} args = v\_theta Px Py Pz Rx Ry Rz
  v\_theta = equal\sphinxhyphen{}style variable for rotaton of region over time (in radians)
  Px,Py,Pz = origin for axis of rotation (distance units)
  Rx,Ry,Rz = axis of rotation vector
\sphinxstyleemphasis{open} value = integer from 1\sphinxhyphen{}6 corresponding to face index (see below)
\end{sphinxalltt}

\item {} 
\sphinxAtStartPar
accelerated styles (with same args) = \sphinxstyleemphasis{block/kk}

\end{itemize}


\subsection{Examples}
\label{\detokenize{region:examples}}
\begin{sphinxVerbatim}[commandchars=\\\{\}]
\PYG{k}{region }\PYG{n+nv+nvIdentifier}{1}\PYG{+w}{ }\PYG{n}{block}\PYG{+w}{ }\PYG{o}{\PYGZhy{}}\PYG{l+m}{3.0}\PYG{+w}{ }\PYG{l+m}{5.0}\PYG{+w}{ }\PYG{n}{INF}\PYG{+w}{ }\PYG{l+m}{10.0}\PYG{+w}{ }\PYG{n}{INF}\PYG{+w}{ }\PYG{n}{INF}
\PYG{k}{region }\PYG{n+nv+nvIdentifier}{2}\PYG{+w}{ }\PYG{n}{sphere}\PYG{+w}{ }\PYG{l+m}{0.0}\PYG{+w}{ }\PYG{l+m}{0.0}\PYG{+w}{ }\PYG{l+m}{0.0}\PYG{+w}{ }\PYG{l+m}{5}\PYG{+w}{ }\PYG{n}{side}\PYG{+w}{ }\PYG{n}{out}
\PYG{k}{region }\PYG{n+nv+nvIdentifier}{void}\PYG{+w}{ }\PYG{n}{cylinder}\PYG{+w}{ }\PYG{n}{y}\PYG{+w}{ }\PYG{l+m}{2}\PYG{+w}{ }\PYG{l+m}{3}\PYG{+w}{ }\PYG{l+m}{5}\PYG{+w}{ }\PYG{o}{\PYGZhy{}}\PYG{l+m}{5.0}\PYG{+w}{ }\PYG{n}{EDGE}\PYG{+w}{ }\PYG{n}{units}\PYG{+w}{ }\PYG{n}{box}
\PYG{k}{region }\PYG{n+nv+nvIdentifier}{1}\PYG{+w}{ }\PYG{n}{prism}\PYG{+w}{ }\PYG{l+m}{0}\PYG{+w}{ }\PYG{l+m}{10}\PYG{+w}{ }\PYG{l+m}{0}\PYG{+w}{ }\PYG{l+m}{10}\PYG{+w}{ }\PYG{l+m}{0}\PYG{+w}{ }\PYG{l+m}{10}\PYG{+w}{ }\PYG{l+m}{2}\PYG{+w}{ }\PYG{l+m}{0}\PYG{+w}{ }\PYG{l+m}{0}
\PYG{k}{region }\PYG{n+nv+nvIdentifier}{outside}\PYG{+w}{ }\PYG{n}{union}\PYG{+w}{ }\PYG{l+m}{4}\PYG{+w}{ }\PYG{n}{side1}\PYG{+w}{ }\PYG{n}{side2}\PYG{+w}{ }\PYG{n}{side3}\PYG{+w}{ }\PYG{n}{side4}
\PYG{k}{region }\PYG{n+nv+nvIdentifier}{2}\PYG{+w}{ }\PYG{n}{sphere}\PYG{+w}{ }\PYG{l+m}{0.0}\PYG{+w}{ }\PYG{l+m}{0.0}\PYG{+w}{ }\PYG{l+m}{0.0}\PYG{+w}{ }\PYG{l+m}{5}\PYG{+w}{ }\PYG{n}{side}\PYG{+w}{ }\PYG{n}{out}\PYG{+w}{ }\PYG{n}{move}\PYG{+w}{ }\PYG{n}{v\PYGZus{}left}\PYG{+w}{ }\PYG{n}{v\PYGZus{}up}\PYG{+w}{ }\PYG{n}{NULL}
\PYG{k}{region }\PYG{n+nv+nvIdentifier}{openbox}\PYG{+w}{ }\PYG{n}{block}\PYG{+w}{ }\PYG{l+m}{0}\PYG{+w}{ }\PYG{l+m}{10}\PYG{+w}{ }\PYG{l+m}{0}\PYG{+w}{ }\PYG{l+m}{10}\PYG{+w}{ }\PYG{l+m}{0}\PYG{+w}{ }\PYG{l+m}{10}\PYG{+w}{ }\PYG{n}{open}\PYG{+w}{ }\PYG{l+m}{5}\PYG{+w}{ }\PYG{n}{open}\PYG{+w}{ }\PYG{l+m}{6}\PYG{+w}{ }\PYG{n}{units}\PYG{+w}{ }\PYG{n}{box}
\PYG{k}{region }\PYG{n+nv+nvIdentifier}{funnel}\PYG{+w}{ }\PYG{n}{cone}\PYG{+w}{ }\PYG{n}{z}\PYG{+w}{ }\PYG{l+m}{10}\PYG{+w}{ }\PYG{l+m}{10}\PYG{+w}{ }\PYG{l+m}{2}\PYG{+w}{ }\PYG{l+m}{5}\PYG{+w}{ }\PYG{l+m}{0}\PYG{+w}{ }\PYG{l+m}{10}\PYG{+w}{ }\PYG{n}{open}\PYG{+w}{ }\PYG{l+m}{1}\PYG{+w}{ }\PYG{n}{units}\PYG{+w}{ }\PYG{n}{box}
\end{sphinxVerbatim}


\subsection{Description}
\label{\detokenize{region:description}}
\sphinxAtStartPar
This command defines a geometric region of space.  Various other
commands use regions.  For example, the region can be filled with
atoms via the {\hyperref[\detokenize{create_atoms::doc}]{\sphinxcrossref{\DUrole{doc}{create\_atoms}}}} command.  Or a bounding
box around the region, can be used to define the simulation box via
the {\hyperref[\detokenize{create_box::doc}]{\sphinxcrossref{\DUrole{doc}{create\_box}}}} command.  Or the atoms in the region
can be identified as a group via the {\hyperref[\detokenize{group::doc}]{\sphinxcrossref{\DUrole{doc}{group}}}} command, or
deleted via the {\hyperref[\detokenize{delete_atoms::doc}]{\sphinxcrossref{\DUrole{doc}{delete\_atoms}}}} command.  Or the
surface of the region can be used as a boundary wall via the {\hyperref[\detokenize{fix_wall_region::doc}]{\sphinxcrossref{\DUrole{doc}{fix wall/region}}}} command.

\sphinxAtStartPar
Commands which use regions typically test whether an atom’s position
is contained in the region or not.  For this purpose, coordinates
exactly on the region boundary are considered to be interior to the
region.  This means, for example, for a spherical region, an atom on
the sphere surface would be part of the region if the sphere were
defined with the \sphinxstyleemphasis{side in} keyword, but would not be part of the
region if it were defined using the \sphinxstyleemphasis{side out} keyword.  See more
details on the \sphinxstyleemphasis{side} keyword below.

\sphinxAtStartPar
Normally, regions in LAMMPS are “static”, meaning their geometric
extent does not change with time.  If the \sphinxstyleemphasis{move} or \sphinxstyleemphasis{rotate} keyword
is used, as described below, the region becomes “dynamic”, meaning
it’s location or orientation changes with time.  This may be useful,
for example, when thermostatting a region, via the compute temp/region
command, or when the fix wall/region command uses a region surface as
a bounding wall on particle motion, i.e. a rotating container.

\sphinxAtStartPar
The \sphinxstyleemphasis{delete} style removes the named region.  Since there is little
overhead to defining extra regions, there is normally no need to do
this, unless you are defining and discarding large numbers of regions
in your input script.

\sphinxAtStartPar
The lo/hi values for \sphinxstyleemphasis{block} or \sphinxstyleemphasis{cone} or \sphinxstyleemphasis{cylinder} or \sphinxstyleemphasis{prism} styles
can be specified as EDGE or INF.  EDGE means they extend all the way
to the global simulation box boundary.  Note that this is the current
box boundary; if the box changes size during a simulation, the region
does not.  INF means a large negative or positive number (1.0e20), so
it should encompass the simulation box even if it changes size.  If a
region is defined before the simulation box has been created (via
{\hyperref[\detokenize{create_box::doc}]{\sphinxcrossref{\DUrole{doc}{create\_box}}}} or {\hyperref[\detokenize{read_data::doc}]{\sphinxcrossref{\DUrole{doc}{read\_data}}}} or
{\hyperref[\detokenize{read_restart::doc}]{\sphinxcrossref{\DUrole{doc}{read\_restart}}}} commands), then an EDGE or INF
parameter cannot be used.  For a \sphinxstyleemphasis{prism} region, a non\sphinxhyphen{}zero tilt
factor in any pair of dimensions cannot be used if both the lo/hi
values in either of those dimensions are INF.  E.g. if the xy tilt is
non\sphinxhyphen{}zero, then xlo and xhi cannot both be INF, nor can ylo and yhi.

\begin{sphinxadmonition}{note}{Note:}
\sphinxAtStartPar
Regions in LAMMPS do not get wrapped across periodic boundaries,
as specified by the {\hyperref[\detokenize{boundary::doc}]{\sphinxcrossref{\DUrole{doc}{boundary}}}} command.  For example, a
spherical region that is defined so that it overlaps a periodic
boundary is not treated as 2 half\sphinxhyphen{}spheres, one on either side of the
simulation box.
\end{sphinxadmonition}

\begin{sphinxadmonition}{note}{Note:}
\sphinxAtStartPar
Regions in LAMMPS are always 3d geometric objects, regardless of
whether the {\hyperref[\detokenize{dimension::doc}]{\sphinxcrossref{\DUrole{doc}{dimension}}}} of a simulation is 2d or 3d.
Thus when using regions in a 2d simulation, you should be careful to
define the region so that its intersection with the 2d x\sphinxhyphen{}y plane of
the simulation has the 2d geometric extent you want.
\end{sphinxadmonition}

\sphinxAtStartPar
For style \sphinxstyleemphasis{cone}, an axis\sphinxhyphen{}aligned cone is defined which is like a
\sphinxstyleemphasis{cylinder} except that two different radii (one at each end) can be
defined.  Either of the radii (but not both) can be 0.0.

\sphinxAtStartPar
For style \sphinxstyleemphasis{cone} and \sphinxstyleemphasis{cylinder}, the c1,c2 params are coordinates in
the 2 other dimensions besides the cylinder axis dimension.  For dim =
x, c1/c2 = y/z; for dim = y, c1/c2 = x/z; for dim = z, c1/c2 = x/y.
Thus the third example above specifies a cylinder with its axis in the
y\sphinxhyphen{}direction located at x = 2.0 and z = 3.0, with a radius of 5.0, and
extending in the y\sphinxhyphen{}direction from \sphinxhyphen{}5.0 to the upper box boundary.

\sphinxAtStartPar
\DUrole{versionmodified,added}{Added in version 4May2022.}

\sphinxAtStartPar
For style \sphinxstyleemphasis{ellipsoid}, an axis\sphinxhyphen{}aligned ellipsoid is defined. The
ellipsoid has its center at (x,y,z) and is defined by 3 axis\sphinxhyphen{}aligned
vectors given by A = (a,0,0); B = (0,b,0); C = (0,0,c).  Note that
although the ellipsoid is specified as axis\sphinxhyphen{}aligned it can be rotated
via the optional \sphinxstyleemphasis{rotate} keyword.

\sphinxAtStartPar
For style \sphinxstyleemphasis{plane}, a plane is defined which contain the point
(px,py,pz) and has a normal vector (nx,ny,nz).  The normal vector does
not have to be of unit length.  The “inside” of the plane is the
half\sphinxhyphen{}space in the direction of the normal vector; see the discussion
of the \sphinxstyleemphasis{side} option below.

\sphinxAtStartPar
For style \sphinxstyleemphasis{prism}, a parallelepiped is defined (it’s too hard to spell
parallelepiped in an input script!).  The parallelepiped has its
“origin” at (xlo,ylo,zlo) and is defined by 3 edge vectors starting
from the origin given by A = (xhi\sphinxhyphen{}xlo,0,0); B = (xy,yhi\sphinxhyphen{}ylo,0); C =
(xz,yz,zhi\sphinxhyphen{}zlo).  \sphinxstyleemphasis{Xy,xz,yz} can be 0.0 or positive or negative values
and are called “tilt factors” because they are the amount of
displacement applied to faces of an originally orthogonal box to
transform it into the parallelepiped.

\sphinxAtStartPar
A prism region that will be used with the {\hyperref[\detokenize{create_box::doc}]{\sphinxcrossref{\DUrole{doc}{create\_box}}}}
command to define a triclinic simulation box must have tilt factors
(xy,xz,yz) that do not skew the box more than half the distance of
corresponding the parallel box length.  For example, if xlo = 2 and
xhi = 12, then the x box length is 10 and the xy tilt factor must be
between \sphinxhyphen{}5 and 5.  Similarly, both xz and yz must be between
\sphinxhyphen{}(xhi\sphinxhyphen{}xlo)/2 and +(yhi\sphinxhyphen{}ylo)/2.  Note that this is not a limitation,
since if the maximum tilt factor is 5 (as in this example), then
configurations with tilt = …, \sphinxhyphen{}15, \sphinxhyphen{}5, 5, 15, 25, … are all
geometrically equivalent.

\sphinxAtStartPar
For style \sphinxstyleemphasis{sphere}, a sphere is defined with its center at (x,y,z)
and with radius as its radius.

\sphinxAtStartPar
The \sphinxstyleemphasis{radius} value for styles \sphinxstyleemphasis{sphere} and \sphinxstyleemphasis{cylinder}, and the
parameters a,b,c for style \sphinxstyleemphasis{ellipsoid}, can each be specified as an
equal\sphinxhyphen{}style {\hyperref[\detokenize{variable::doc}]{\sphinxcrossref{\DUrole{doc}{variable}}}}.  Likewise, for style \sphinxstyleemphasis{sphere}
and \sphinxstyleemphasis{ellipsoid} the x\sphinxhyphen{}, y\sphinxhyphen{}, and z\sphinxhyphen{} coordinates of the center of the
sphere/ellipsoid can be specified as an equal\sphinxhyphen{}style variable.  And for
style \sphinxstyleemphasis{cylinder} the two center positions c1 and c2 for the location
of the cylinder axes can be specified as a equal\sphinxhyphen{}style variable. For style \sphinxstyleemphasis{cone}
all properties can be defined via equal\sphinxhyphen{}style variables.

\sphinxAtStartPar
If the value is a variable, it should be specified as v\_name, where
name is the variable name.  In this case, the variable will be
evaluated each timestep, and its value used to determine the radius of
the region.

\sphinxAtStartPar
Equal\sphinxhyphen{}style variables can specify formulas with various mathematical
functions, and include {\hyperref[\detokenize{thermo_style::doc}]{\sphinxcrossref{\DUrole{doc}{thermo\_style}}}} command
keywords for the simulation box parameters and timestep and elapsed
time.  Thus it is easy to specify a time\sphinxhyphen{}dependent radius or have
a time dependent position of the sphere or cylinder region.

\sphinxAtStartPar
See the {\hyperref[\detokenize{Howto_triclinic::doc}]{\sphinxcrossref{\DUrole{doc}{Howto tricilinc}}}} page for a
geometric description of triclinic boxes, as defined by LAMMPS, and
how to transform these parameters to and from other commonly used
triclinic representations.

\sphinxAtStartPar
The \sphinxstyleemphasis{union} style creates a region consisting of the volume of all the
listed regions combined.  The \sphinxstyleemphasis{intersect} style creates a region
consisting of the volume that is common to all the listed regions.

\begin{sphinxadmonition}{note}{Note:}
\sphinxAtStartPar
The \sphinxstyleemphasis{union} and \sphinxstyleemphasis{intersect} regions operate by invoking methods
from their list of sub\sphinxhyphen{}regions.  Thus you cannot delete the
sub\sphinxhyphen{}regions after defining a \sphinxstyleemphasis{union} or \sphinxstyleemphasis{intersection} region.
\end{sphinxadmonition}


\bigskip\hrule\bigskip


\sphinxAtStartPar
The \sphinxstyleemphasis{side} keyword determines whether the region is considered to be
inside or outside of the specified geometry.  Using this keyword in
conjunction with \sphinxstyleemphasis{union} and \sphinxstyleemphasis{intersect} regions, complex geometries
can be built up.  For example, if the interior of two spheres were
each defined as regions, and a \sphinxstyleemphasis{union} style with \sphinxstyleemphasis{side} = out was
constructed listing the region\sphinxhyphen{}IDs of the 2 spheres, the resulting
region would be all the volume in the simulation box that was outside
both of the spheres.

\sphinxAtStartPar
The \sphinxstyleemphasis{units} keyword determines the meaning of the distance units used
to define the region for any argument above listed as having distance
units.  It also affects the scaling of the velocity vector specified
with the \sphinxstyleemphasis{vel} keyword, the amplitude vector specified with the
\sphinxstyleemphasis{wiggle} keyword, and the rotation point specified with the \sphinxstyleemphasis{rotate}
keyword, since they each involve a distance metric.

\sphinxAtStartPar
A \sphinxstyleemphasis{box} value selects standard distance units as defined by the
{\hyperref[\detokenize{units::doc}]{\sphinxcrossref{\DUrole{doc}{units}}}} command, e.g. Angstroms for units = real or metal.
A \sphinxstyleemphasis{lattice} value means the distance units are in lattice spacings.
The {\hyperref[\detokenize{lattice::doc}]{\sphinxcrossref{\DUrole{doc}{lattice}}}} command must have been previously used to
define the lattice spacings which are used as follows:
\begin{itemize}
\item {} 
\sphinxAtStartPar
For style \sphinxstyleemphasis{block}, the lattice spacing in dimension x is applied to
xlo and xhi, similarly the spacings in dimensions y,z are applied to
ylo/yhi and zlo/zhi.

\item {} 
\sphinxAtStartPar
For style \sphinxstyleemphasis{cone}, the lattice spacing in argument \sphinxstyleemphasis{dim} is applied to
lo and hi.  The spacings in the two radial dimensions are applied to
c1 and c2.  The two cone radii are scaled by the lattice
spacing in the dimension corresponding to c1.

\item {} 
\sphinxAtStartPar
For style \sphinxstyleemphasis{cylinder}, the lattice spacing in argument \sphinxstyleemphasis{dim} is applied
to lo and hi.  The spacings in the two radial dimensions are applied
to c1 and c2.  The cylinder radius is scaled by the lattice
spacing in the dimension corresponding to c1.

\item {} 
\sphinxAtStartPar
For style \sphinxstyleemphasis{ellipsoid}, the lattice spacing in dimensions x,y,z are
applied to the ellipsoid center x,y,z.  The spacing in dimensions
x,y,z are applied to the ellipsoid radii a,b,c respectively.

\item {} 
\sphinxAtStartPar
For style \sphinxstyleemphasis{plane}, the lattice spacing in dimension x is applied to
px and nx, similarly the spacings in dimensions y,z are applied to
py/ny and pz/nz.

\item {} 
\sphinxAtStartPar
For style \sphinxstyleemphasis{prism}, the lattice spacing in dimension x is applied to
xlo and xhi, similarly for ylo/yhi and zlo/zhi.  The lattice spacing
in dimension x is applied to xy and xz, and the spacing in dimension y
to yz.

\item {} 
\sphinxAtStartPar
For style \sphinxstyleemphasis{sphere}, the lattice spacing in dimensions x,y,z are
applied to the sphere center x,y,z.  The spacing in dimension x is
applied to the sphere radius.

\end{itemize}


\bigskip\hrule\bigskip


\sphinxAtStartPar
If the \sphinxstyleemphasis{move} or \sphinxstyleemphasis{rotate} keywords are used, the region is “dynamic”,
meaning its location or orientation changes with time.  These keywords
cannot be used with a \sphinxstyleemphasis{union} or \sphinxstyleemphasis{intersect} style region.  Instead,
the keywords should be used to make the individual sub\sphinxhyphen{}regions of the
\sphinxstyleemphasis{union} or \sphinxstyleemphasis{intersect} region dynamic.  Normally, each sub\sphinxhyphen{}region
should be “dynamic” in the same manner (e.g. rotate around the same
point), though this is not a requirement.

\sphinxAtStartPar
The \sphinxstyleemphasis{move} keyword allows one or more {\hyperref[\detokenize{variable::doc}]{\sphinxcrossref{\DUrole{doc}{equal\sphinxhyphen{}style variables}}}} to be used to specify the x,y,z displacement
of the region, typically as a function of time.  A variable is
specified as v\_name, where name is the variable name.  Any of the
three variables can be specified as NULL, in which case no
displacement is calculated in that dimension.

\sphinxAtStartPar
Note that equal\sphinxhyphen{}style variables can specify formulas with various
mathematical functions, and include {\hyperref[\detokenize{thermo_style::doc}]{\sphinxcrossref{\DUrole{doc}{thermo\_style}}}}
command keywords for the simulation box parameters and timestep and
elapsed time.  Thus it is easy to specify a region displacement that
change as a function of time or spans consecutive runs in a continuous
fashion.  For the latter, see the \sphinxstyleemphasis{start} and \sphinxstyleemphasis{stop} keywords of the
{\hyperref[\detokenize{run::doc}]{\sphinxcrossref{\DUrole{doc}{run}}}} command and the \sphinxstyleemphasis{elaplong} keyword of {\hyperref[\detokenize{thermo_style::doc}]{\sphinxcrossref{\DUrole{doc}{thermo\_style custom}}}} for details.

\sphinxAtStartPar
For example, these commands would displace a region from its initial
position, in the positive x direction, effectively at a constant
velocity:

\begin{sphinxVerbatim}[commandchars=\\\{\}]
\PYG{k}{variable }\PYG{n+nv+nvIdentifier}{dx}\PYG{+w}{ }\PYG{n}{equal}\PYG{+w}{ }\PYG{n}{ramp}\PYG{n+nv}{(}\PYG{n+nv}{0,10}\PYG{n+nv}{)}
\PYG{k}{region }\PYG{n+nv+nvIdentifier}{2}\PYG{+w}{ }\PYG{n}{sphere}\PYG{+w}{ }\PYG{l+m}{10.0}\PYG{+w}{ }\PYG{l+m}{10.0}\PYG{+w}{ }\PYG{l+m}{0.0}\PYG{+w}{ }\PYG{l+m}{5}\PYG{+w}{ }\PYG{n}{move}\PYG{+w}{ }\PYG{n}{v\PYGZus{}dx}\PYG{+w}{ }\PYG{n}{NULL}\PYG{+w}{ }\PYG{n}{NULL}
\end{sphinxVerbatim}

\sphinxAtStartPar
Note that the initial displacement is 0.0, though that is not required.

\sphinxAtStartPar
Either of these variables would “wiggle” the region back and forth in
the y direction:

\begin{sphinxVerbatim}[commandchars=\\\{\}]
\PYG{k}{variable }\PYG{n+nv+nvIdentifier}{dy}\PYG{+w}{ }\PYG{n}{equal}\PYG{+w}{ }\PYG{n}{swiggle}\PYG{n+nv}{(}\PYG{n+nv}{0,5,100}\PYG{n+nv}{)}
\PYG{k}{variable }\PYG{n+nv+nvIdentifier}{dysame}\PYG{+w}{ }\PYG{n}{equal}\PYG{+w}{ }\PYG{l+m}{5}\PYG{o}{*}\PYG{n}{sin}\PYG{n+nv}{(}\PYG{n+nv}{2*PI*elaplong*dt/100}\PYG{n+nv}{)}
\PYG{k}{region }\PYG{n+nv+nvIdentifier}{2}\PYG{+w}{ }\PYG{n}{sphere}\PYG{+w}{ }\PYG{l+m}{10.0}\PYG{+w}{ }\PYG{l+m}{10.0}\PYG{+w}{ }\PYG{l+m}{0.0}\PYG{+w}{ }\PYG{l+m}{5}\PYG{+w}{ }\PYG{n}{move}\PYG{+w}{ }\PYG{n}{NULL}\PYG{+w}{ }\PYG{n}{v\PYGZus{}dy}\PYG{+w}{ }\PYG{n}{NULL}
\end{sphinxVerbatim}

\sphinxAtStartPar
The \sphinxstyleemphasis{rotate} keyword rotates the region around a rotation axis \sphinxstyleemphasis{R} =
(Rx,Ry,Rz) that goes through a point \sphinxstyleemphasis{P} = (Px,Py,Pz).  The rotation
angle is calculated, presumably as a function of time, by a variable
specified as v\_theta, where theta is the variable name.  The variable
should generate its result in radians.  The direction of rotation for
the region around the rotation axis is consistent with the right\sphinxhyphen{}hand
rule: if your right\sphinxhyphen{}hand thumb points along \sphinxstyleemphasis{R}, then your fingers
wrap around the axis in the direction of rotation.

\sphinxAtStartPar
The \sphinxstyleemphasis{move} and \sphinxstyleemphasis{rotate} keywords can be used together.  In this case,
the displacement specified by the \sphinxstyleemphasis{move} keyword is applied to the \sphinxstyleemphasis{P}
point of the \sphinxstyleemphasis{rotate} keyword.


\bigskip\hrule\bigskip


\sphinxAtStartPar
The \sphinxstyleemphasis{open} keyword can be used (multiple times) to indicate that one
or more faces of the region are ignored for purposes of particle/wall
interactions.  This keyword is only relevant for regions used by the
\sphinxstyleemphasis{fix wall/region} and \sphinxstyleemphasis{fix wall/gran/region} commands.  It can be used
to create “open” containers where only some of the region faces are
walls.  For example, a funnel can be created with a \sphinxstyleemphasis{cone} style
region that has an open face at the smaller radius for particles to
flow out, or at the larger radius for pouring particles into the cone,
or both.

\sphinxAtStartPar
Note that using the \sphinxstyleemphasis{open} keyword partly overrides the \sphinxstyleemphasis{side}
keyword, since both exterior and interior surfaces of an open region
are tested for particle contacts.  The exception to this is a \sphinxstyleemphasis{union}
or \sphinxstyleemphasis{intersect} region which includes an open sub\sphinxhyphen{}region.  In that case
the \sphinxstyleemphasis{side} keyword is still used to define the union/intersect region
volume, and the \sphinxstyleemphasis{open} settings are only applied to the individual
sub\sphinxhyphen{}regions that use them.

\sphinxAtStartPar
The indices specified as part of the \sphinxstyleemphasis{open} keyword have the following
meanings:

\sphinxAtStartPar
For style \sphinxstyleemphasis{block}, indices 1\sphinxhyphen{}6 correspond to the xlo, xhi, ylo, yhi,
zlo, zhi surfaces of the block.  I.e. 1 is the yz plane at x = xlo, 2
is the yz\sphinxhyphen{}plane at x = xhi, 3 is the xz plane at y = ylo, 4 is the xz
plane at y = yhi, 5 is the xy plane at z = zlo, 6 is the xy plane at z
= zhi).  In the second\sphinxhyphen{}to\sphinxhyphen{}last example above, the region is a box open
at both xy planes.

\sphinxAtStartPar
For style \sphinxstyleemphasis{prism}, values 1\sphinxhyphen{}6 have the same mapping as for style
\sphinxstyleemphasis{block}.  I.e. in an untilted \sphinxstyleemphasis{prism}, \sphinxstyleemphasis{open} indices correspond to
the xlo, xhi, ylo, yhi, zlo, zhi surfaces.

\sphinxAtStartPar
For style \sphinxstyleemphasis{cylinder}, index 1 corresponds to the flat end cap at the
low coordinate along the cylinder axis, index 2 corresponds to the
high\sphinxhyphen{}coordinate flat end cap along the cylinder axis, and index 3 is
the curved cylinder surface.  For example, a \sphinxstyleemphasis{cylinder} region with
\sphinxstyleemphasis{open 1 open 2} keywords will be open at both ends (e.g. a section of
pipe), regardless of the cylinder orientation.

\sphinxAtStartPar
For style \sphinxstyleemphasis{cone}, the mapping is the same as for style \sphinxstyleemphasis{cylinder}.
Index 1 is the low\sphinxhyphen{}coordinate flat end cap, index 2 is the
high\sphinxhyphen{}coordinate flat end cap, and index 3 is the curved cone surface.
In the last example above, a \sphinxstyleemphasis{cone} region is defined along the z\sphinxhyphen{}axis
that is open at the zlo value (e.g. for use as a funnel).

\sphinxAtStartPar
For all other styles, the \sphinxstyleemphasis{open} keyword is ignored.  As indicated
above, this includes the \sphinxstyleemphasis{intersect} and \sphinxstyleemphasis{union} regions, though their
sub\sphinxhyphen{}regions can be defined with the \sphinxstyleemphasis{open} keyword.


\bigskip\hrule\bigskip


\sphinxAtStartPar
Styles with a \sphinxstyleemphasis{gpu}, \sphinxstyleemphasis{intel}, \sphinxstyleemphasis{kk}, \sphinxstyleemphasis{omp}, or \sphinxstyleemphasis{opt} suffix are
functionally the same as the corresponding style without the suffix.
They have been optimized to run faster, depending on your available
hardware, as discussed on the {\hyperref[\detokenize{Speed_packages::doc}]{\sphinxcrossref{\DUrole{doc}{Accelerator packages}}}}
page.  The accelerated styles take the same arguments and should
produce the same results, except for round\sphinxhyphen{}off and precision issues.

\sphinxAtStartPar
These accelerated styles are part of the GPU, INTEL, KOKKOS,
OPENMP, and OPT packages, respectively.  They are only enabled if
LAMMPS was built with those packages.  See the {\hyperref[\detokenize{Build_package::doc}]{\sphinxcrossref{\DUrole{doc}{Build package}}}} page for more info.

\sphinxAtStartPar
You can specify the accelerated styles explicitly in your input script
by including their suffix, or you can use the {\hyperref[\detokenize{Run_options::doc}]{\sphinxcrossref{\DUrole{doc}{\sphinxhyphen{}suffix command\sphinxhyphen{}line switch}}}} when you invoke LAMMPS, or you can use the
{\hyperref[\detokenize{suffix::doc}]{\sphinxcrossref{\DUrole{doc}{suffix}}}} command in your input script.

\sphinxAtStartPar
See the {\hyperref[\detokenize{Speed_packages::doc}]{\sphinxcrossref{\DUrole{doc}{Accelerator packages}}}} page for more
instructions on how to use the accelerated styles effectively.

\begin{sphinxadmonition}{note}{Note:}
\sphinxAtStartPar
Currently, only \sphinxstyleemphasis{block} style regions are supported by Kokkos.  The
code using the region (such as a fix or compute) must also be
supported by Kokkos or no acceleration will occur.
\end{sphinxadmonition}


\bigskip\hrule\bigskip



\subsection{Restrictions}
\label{\detokenize{region:restrictions}}
\sphinxAtStartPar
A prism cannot be of 0.0 thickness in any dimension; use a small z
thickness for 2d simulations.  For 2d simulations, the xz and yz
parameters must be 0.0.


\subsection{Related commands}
\label{\detokenize{region:related-commands}}
\sphinxAtStartPar
{\hyperref[\detokenize{lattice::doc}]{\sphinxcrossref{\DUrole{doc}{lattice}}}}, {\hyperref[\detokenize{create_atoms::doc}]{\sphinxcrossref{\DUrole{doc}{create\_atoms}}}},
{\hyperref[\detokenize{delete_atoms::doc}]{\sphinxcrossref{\DUrole{doc}{delete\_atoms}}}}, {\hyperref[\detokenize{group::doc}]{\sphinxcrossref{\DUrole{doc}{group}}}}


\subsection{Default}
\label{\detokenize{region:default}}
\sphinxAtStartPar
The option defaults are side = in, units = lattice, and no move or
rotation.

\sphinxstepscope

\index{replicate@\spxentry{replicate}}\ignorespaces 

\section{replicate command}
\label{\detokenize{replicate:replicate-command}}\label{\detokenize{replicate:index-0}}\label{\detokenize{replicate::doc}}

\subsection{Syntax}
\label{\detokenize{replicate:syntax}}
\begin{sphinxVerbatim}[commandchars=\\\{\}]
\PYG{k}{replicate}\PYG{+w}{ }\PYG{n}{nx}\PYG{+w}{ }\PYG{n}{ny}\PYG{+w}{ }\PYG{n}{nz}\PYG{+w}{ }\PYG{n}{keyword}\PYG{+w}{ }\PYG{n}{...}
\end{sphinxVerbatim}

\sphinxAtStartPar
nx,ny,nz = replication factors in each dimension
\begin{itemize}
\item {} 
\sphinxAtStartPar
zero or more keywords may be appended

\item {} 
\sphinxAtStartPar
keyword = \sphinxstyleemphasis{bbox} or \sphinxstyleemphasis{bond/periodic}
\begin{sphinxalltt}
\sphinxstyleemphasis{bbox} = use a bounding\sphinxhyphen{}box algorithm which is faster for large proc counts
\sphinxstyleemphasis{bond/periodic} = use an algorithm that correctly replicates periodic bond loops
\end{sphinxalltt}

\end{itemize}


\subsection{Examples}
\label{\detokenize{replicate:examples}}
\sphinxAtStartPar
For examples of replicating simple linear polymer chains (periodic or
non\sphinxhyphen{}periodic) or periodic carbon nanotubes, see examples/replicate.

\begin{sphinxVerbatim}[commandchars=\\\{\}]
\PYG{k}{replicate}\PYG{+w}{ }\PYG{l+m}{2}\PYG{+w}{ }\PYG{l+m}{3}\PYG{+w}{ }\PYG{l+m}{2}
\PYG{k}{replicate}\PYG{+w}{ }\PYG{l+m}{2}\PYG{+w}{ }\PYG{l+m}{3}\PYG{+w}{ }\PYG{l+m}{2}\PYG{+w}{ }\PYG{n}{bbox}
\PYG{k}{replicate}\PYG{+w}{ }\PYG{l+m}{2}\PYG{+w}{ }\PYG{l+m}{3}\PYG{+w}{ }\PYG{l+m}{2}\PYG{+w}{ }\PYG{n}{bond}\PYG{o}{/}\PYG{n}{periodic}
\end{sphinxVerbatim}


\subsection{Description}
\label{\detokenize{replicate:description}}
\sphinxAtStartPar
Replicate the current system one or more times in each dimension.  For
example, replication factors of 2,2,2 will create a simulation with 8x
as many atoms by doubling the size of the simulation box in each
dimension.  A replication factor of 1 leaves the simulation domain
unchanged in that dimension.

\sphinxAtStartPar
When the new simulation box is created it is partitioned into a
regular 3d grid of rectangular bricks, one per processor, based on the
number of processors being used and the settings of the
{\hyperref[\detokenize{processors::doc}]{\sphinxcrossref{\DUrole{doc}{processors}}}} command.  The partitioning can be
changed by subsequent {\hyperref[\detokenize{balance::doc}]{\sphinxcrossref{\DUrole{doc}{balance}}}} or {\hyperref[\detokenize{fix_balance::doc}]{\sphinxcrossref{\DUrole{doc}{fix balance}}}} commands.

\sphinxAtStartPar
All properties of each atom are replicated (except per\sphinxhyphen{}atom fix data,
see the Restrictions section below).  This includes their velocities,
which may or may not be desirable.  New atom IDs are assigned to new
atoms, as are new molecule IDs.  Bonds and other topology interactions
are created between pairs of new atoms as well as between old and new
atoms.

\begin{sphinxadmonition}{note}{Note:}
\sphinxAtStartPar
The bond discussion which follows only refers to models with
permanent covalent bonds typically defined in LAMMPS via a data
file.  It is not relevant to systems modeled with many\sphinxhyphen{}body
potentials which can define bonds on\sphinxhyphen{}the\sphinxhyphen{}fly, based on the current
positions of nearby atoms, e.g. models using the {\hyperref[\detokenize{pair_airebo::doc}]{\sphinxcrossref{\DUrole{doc}{AIREBO}}}} or {\hyperref[\detokenize{pair_reaxff::doc}]{\sphinxcrossref{\DUrole{doc}{ReaxFF}}}} potentials.
\end{sphinxadmonition}

\sphinxAtStartPar
If the \sphinxstyleemphasis{bond/periodic} keyword is not specified, bond replication is
done by using the image flag for each atom to “unwrap” it out of the
periodic box before replicating it.  After replication is performed,
atoms outside the new periodic box are wrapped back into it.  This
assigns correct images flags to all atoms in the system.  For this to
work, all original atoms in the original simulation box must have
consistent image flags.  This means that if two atoms have a bond
between them which crosses a periodic boundary, their respective image
flags will differ by 1 in that dimension.

\sphinxAtStartPar
Image flag consistency is not possible if a system has a periodic bond
loop, meaning there is a chain of bonds which crosses an entire
dimension and re\sphinxhyphen{}connects to itself across a periodic boundary.  In
this case you MUST use the \sphinxstyleemphasis{bond/periodic} keyword to correctly
replicate the system.  This option zeroes the image flags for all
atoms and uses a different algorithm to find new (nearby) bond
neighbors in the replicated system.  In the final replicated system
all image flags are zero (in each dimension).

\begin{sphinxadmonition}{note}{Note:}
\sphinxAtStartPar
LAMMPS does not check for image flag consistency before performing
the replication (it does issue a warning about this before a
simulation is run).  If the original image flags are inconsistent,
the replicated system will also have inconsistent image flags, but
will otherwise be correctly replicated.  This is NOT the case if
there is a periodic bond loop.  See the next note.
\end{sphinxadmonition}

\begin{sphinxadmonition}{note}{Note:}
\sphinxAtStartPar
LAMMPS does not check for periodic bond loops.  If you use the
\sphinxstyleemphasis{bond/periodic} keyword for a system without periodic bond loops,
the system will be correctly replicated, but image flag information
will be lost (which may or may not be important to your model).  If
you do not use the \sphinxstyleemphasis{bond/periodic} keyword for a system with
periodic bond loops, the replicated system will have invalid bonds
(typically very long), resulting in bad dynamics.
\end{sphinxadmonition}

\sphinxAtStartPar
If possible, the \sphinxstyleemphasis{bbox} keyword should be used when running on a large
number of processors, as it can result in a substantial speed\sphinxhyphen{}up for
the replication operation.  It uses a bounding box to only check atoms
in replicas that overlap with each processor’s new subdomain when
assigning atoms to processors.  It also preserves image flag
information.  The only drawback to the \sphinxstyleemphasis{bbox} option is that it
requires a temporary use of more memory.  Each processor must be able
to store all atoms (and their per\sphinxhyphen{}atom data) in the original system,
before it is replicated.

\begin{sphinxadmonition}{note}{Note:}
\sphinxAtStartPar
The algorithm used by the \sphinxstyleemphasis{bond/periodic} keyword builds on the
algorithm used by the \sphinxstyleemphasis{bbox} keyword and thus has the same memory
requirements.  If you specify only the \sphinxstyleemphasis{bond/peridoic} keyword it
will internally set the \sphinxstyleemphasis{bbox} keyword as well.
\end{sphinxadmonition}


\bigskip\hrule\bigskip



\subsection{Restrictions}
\label{\detokenize{replicate:restrictions}}
\sphinxAtStartPar
A 2d simulation cannot be replicated in the z dimension.

\sphinxAtStartPar
If a simulation is non\sphinxhyphen{}periodic in a dimension, care should be used
when replicating it in that dimension, as it may generate atoms nearly
on top of each other.

\sphinxAtStartPar
If the current simulation was read in from a restart file (before a
run is performed), there must not be any fix information stored in the
file for individual atoms.  Similarly, no fixes can be defined at the
time the replicate command is used that require vectors of atom
information to be stored.  This is because the replicate command does
not know how to replicate that information for new atoms it creates.

\sphinxAtStartPar
To work around this restriction two options are possible.  (1) Fixes
which use the stored data in the restart file can be defined before
replication and then deleted via the {\hyperref[\detokenize{unfix::doc}]{\sphinxcrossref{\DUrole{doc}{unfix}}}} command and
re\sphinxhyphen{}defined after it.  Or (2) the restart file can be converted to a
data file (which deletes the stored fix information) and fixes defined
after the replicate command.  In both these scenarios, the per\sphinxhyphen{}atom
fix information in the restart file is lost.


\subsection{Related commands}
\label{\detokenize{replicate:related-commands}}
\sphinxAtStartPar
none


\subsection{Default}
\label{\detokenize{replicate:default}}
\sphinxAtStartPar
No settings for using the \sphinxstyleemphasis{bbox} or \sphinxstyleemphasis{bond/periodic} algorithms.

\sphinxstepscope

\index{rerun@\spxentry{rerun}}\ignorespaces 

\section{rerun command}
\label{\detokenize{rerun:rerun-command}}\label{\detokenize{rerun:index-0}}\label{\detokenize{rerun::doc}}

\subsection{Syntax}
\label{\detokenize{rerun:syntax}}
\begin{sphinxVerbatim}[commandchars=\\\{\}]
\PYG{k}{rerun}\PYG{+w}{ }\PYG{n}{file1}\PYG{+w}{ }\PYG{n}{file2}\PYG{+w}{ }\PYG{n}{...}\PYG{+w}{ }\PYG{n}{keyword}\PYG{+w}{ }\PYG{n}{args}\PYG{+w}{ }\PYG{n}{...}
\end{sphinxVerbatim}
\begin{itemize}
\item {} 
\sphinxAtStartPar
file1,file2,… = dump file(s) to read

\item {} 
\sphinxAtStartPar
one or more keywords may be appended, keyword \sphinxstyleemphasis{dump} must appear and be last
\begin{sphinxalltt}
keyword = \sphinxstyleemphasis{first} or \sphinxstyleemphasis{last} or \sphinxstyleemphasis{every} or \sphinxstyleemphasis{skip} or \sphinxstyleemphasis{start} or \sphinxstyleemphasis{stop} or \sphinxstyleemphasis{post} or \sphinxstyleemphasis{dump}
 \sphinxstyleemphasis{first} args = Nfirst
   Nfirst = dump timestep to start on
 \sphinxstyleemphasis{last} args = Nlast
   Nlast = dumptimestep to stop on
 \sphinxstyleemphasis{every} args = Nevery
   Nevery = read snapshots matching every this many timesteps
 \sphinxstyleemphasis{skip} args = Nskip
   Nskip = read one out of every Nskip snapshots
 \sphinxstyleemphasis{start} args = Nstart
   Nstart = timestep on which pseudo run will start
 \sphinxstyleemphasis{stop} args = Nstop
   Nstop = timestep to which pseudo run will end
 \sphinxstyleemphasis{post} value = \sphinxstyleemphasis{yes} or \sphinxstyleemphasis{no}
 \sphinxstyleemphasis{dump} args = same as {\hyperref[\detokenize{read_dump::doc}]{\sphinxcrossref{\DUrole{doc}{read\_dump}}}} command starting with its field arguments
\end{sphinxalltt}

\end{itemize}


\subsection{Examples}
\label{\detokenize{rerun:examples}}
\begin{sphinxVerbatim}[commandchars=\\\{\}]
\PYG{k}{rerun}\PYG{+w}{ }\PYG{n}{dump.file}\PYG{+w}{ }\PYG{k}{dump }\PYG{n+nv+nvIdentifier}{x}\PYG{+w}{ }\PYG{n+nv+nvIdentifier}{y}\PYG{+w}{ }\PYG{n}{z}\PYG{+w}{ }\PYG{n}{vx}\PYG{+w}{ }\PYG{n}{vy}\PYG{+w}{ }\PYG{n}{vz}
\PYG{k}{rerun}\PYG{+w}{ }\PYG{n}{dump1.txt}\PYG{+w}{ }\PYG{n}{dump2.txt}\PYG{+w}{ }\PYG{n}{first}\PYG{+w}{ }\PYG{l+m}{10000}\PYG{+w}{ }\PYG{n}{every}\PYG{+w}{ }\PYG{l+m}{1000}\PYG{+w}{ }\PYG{k}{dump }\PYG{n+nv+nvIdentifier}{x}\PYG{+w}{ }\PYG{n+nv+nvIdentifier}{y}\PYG{+w}{ }\PYG{n}{z}
\PYG{k}{rerun}\PYG{+w}{ }\PYG{n}{dump.vels}\PYG{+w}{ }\PYG{k}{dump }\PYG{n+nv+nvIdentifier}{x}\PYG{+w}{ }\PYG{n+nv+nvIdentifier}{y}\PYG{+w}{ }\PYG{n}{z}\PYG{+w}{ }\PYG{n}{vx}\PYG{+w}{ }\PYG{n}{vy}\PYG{+w}{ }\PYG{n}{vz}\PYG{+w}{ }\PYG{n}{box}\PYG{+w}{ }\PYG{n}{yes}\PYG{+w}{ }\PYG{n}{format}\PYG{+w}{ }\PYG{n}{molfile}\PYG{+w}{ }\PYG{n}{lammpstrj}
\PYG{k}{rerun}\PYG{+w}{ }\PYG{n}{dump.dcd}\PYG{+w}{ }\PYG{k}{dump }\PYG{n+nv+nvIdentifier}{x}\PYG{+w}{ }\PYG{n+nv+nvIdentifier}{y}\PYG{+w}{ }\PYG{n}{z}\PYG{+w}{ }\PYG{n}{box}\PYG{+w}{ }\PYG{n}{no}\PYG{+w}{ }\PYG{n}{format}\PYG{+w}{ }\PYG{n}{molfile}\PYG{+w}{ }\PYG{n}{dcd}
\PYG{k}{rerun}\PYG{+w}{ }\PYG{n}{..}\PYG{o}{/}\PYG{n}{run7}\PYG{o}{/}\PYG{n}{dump.file.gz}\PYG{+w}{ }\PYG{n}{skip}\PYG{+w}{ }\PYG{l+m}{2}\PYG{+w}{ }\PYG{k}{dump }\PYG{n+nv+nvIdentifier}{x}\PYG{+w}{ }\PYG{n+nv+nvIdentifier}{y}\PYG{+w}{ }\PYG{n}{z}\PYG{+w}{ }\PYG{n}{box}\PYG{+w}{ }\PYG{n}{yes}
\PYG{k}{rerun}\PYG{+w}{ }\PYG{n}{dump.bp}\PYG{+w}{ }\PYG{k}{dump }\PYG{n+nv+nvIdentifier}{x}\PYG{+w}{ }\PYG{n+nv+nvIdentifier}{y}\PYG{+w}{ }\PYG{n}{z}\PYG{+w}{ }\PYG{n}{box}\PYG{+w}{ }\PYG{n}{no}\PYG{+w}{ }\PYG{n}{format}\PYG{+w}{ }\PYG{n}{adios}
\PYG{k}{rerun}\PYG{+w}{ }\PYG{n}{dump.bp}\PYG{+w}{ }\PYG{k}{dump }\PYG{n+nv+nvIdentifier}{x}\PYG{+w}{ }\PYG{n+nv+nvIdentifier}{y}\PYG{+w}{ }\PYG{n}{z}\PYG{+w}{ }\PYG{n}{vx}\PYG{+w}{ }\PYG{n}{vy}\PYG{+w}{ }\PYG{n}{vz}\PYG{+w}{ }\PYG{n}{format}\PYG{+w}{ }\PYG{n}{adios}\PYG{+w}{ }\PYG{n}{timeout}\PYG{+w}{ }\PYG{l+m}{10.0}
\end{sphinxVerbatim}


\subsection{Description}
\label{\detokenize{rerun:description}}
\sphinxAtStartPar
Perform a pseudo simulation run where atom information is read one
snapshot at a time from a dump file(s), and energies and forces are
computed on the shapshot to produce thermodynamic or other output.

\sphinxAtStartPar
This can be useful in the following kinds of scenarios, after an
initial simulation produced the dump file:
\begin{itemize}
\item {} 
\sphinxAtStartPar
Compute the energy and forces of snapshots using a different potential.

\item {} 
\sphinxAtStartPar
Calculate one or more diagnostic quantities on the snapshots that
were not computed in the initial run.  These can also be computed with
settings not used in the initial run, e.g. computing an RDF via the
{\hyperref[\detokenize{compute_rdf::doc}]{\sphinxcrossref{\DUrole{doc}{compute rdf}}}} command with a longer cutoff than was
used initially.

\item {} 
\sphinxAtStartPar
Calculate the portion of per\sphinxhyphen{}atom forces resulting from a subset of
the potential.  E.g. compute only Coulombic forces.  This can be done
by only defining only a Coulombic pair style in the rerun script.
Doing this in the original script would result in different (bad)
dynamics.

\end{itemize}

\sphinxAtStartPar
Conceptually, using the rerun command is like running an input script
that has a loop in it (see the {\hyperref[\detokenize{next::doc}]{\sphinxcrossref{\DUrole{doc}{next}}}} and {\hyperref[\detokenize{jump::doc}]{\sphinxcrossref{\DUrole{doc}{jump}}}}
commands).  Each iteration of the loop reads one snapshot from the
dump file via the {\hyperref[\detokenize{read_dump::doc}]{\sphinxcrossref{\DUrole{doc}{read\_dump}}}} command, sets the
timestep to the appropriate value, and then invokes a {\hyperref[\detokenize{run::doc}]{\sphinxcrossref{\DUrole{doc}{run}}}}
command for zero timesteps to simply compute energy and forces, and
any other {\hyperref[\detokenize{thermo_style::doc}]{\sphinxcrossref{\DUrole{doc}{thermodynamic output}}}} or diagnostic info
you have defined.  This computation also invokes any fixes you have
defined that apply constraints to the system, such as {\hyperref[\detokenize{fix_shake::doc}]{\sphinxcrossref{\DUrole{doc}{fix shake}}}} or {\hyperref[\detokenize{fix_indent::doc}]{\sphinxcrossref{\DUrole{doc}{fix indent}}}}.

\sphinxAtStartPar
Note that a simulation box must already be defined before using the
rerun command.  This can be done by the {\hyperref[\detokenize{create_box::doc}]{\sphinxcrossref{\DUrole{doc}{create\_box}}}},
{\hyperref[\detokenize{read_data::doc}]{\sphinxcrossref{\DUrole{doc}{read\_data}}}}, or {\hyperref[\detokenize{read_restart::doc}]{\sphinxcrossref{\DUrole{doc}{read\_restart}}}}
commands.

\sphinxAtStartPar
Also note that reading per\sphinxhyphen{}atom information from dump snapshots is
limited to the atom coordinates, velocities and image flags as
explained in the {\hyperref[\detokenize{read_dump::doc}]{\sphinxcrossref{\DUrole{doc}{read\_dump}}}} command.  Other atom
properties, which may be necessary to compute energies and forces,
such as atom charge, or bond topology information for a molecular
system, are not read from (or even contained in) dump files.  Thus
this auxiliary information should be defined in the usual way, e.g. in
a data file read in by a {\hyperref[\detokenize{read_data::doc}]{\sphinxcrossref{\DUrole{doc}{read\_data}}}} command, before
using the rerun command.

\sphinxAtStartPar
Also note that the frequency of thermodynamic or dump output from the
rerun simulation will depend on settings made in the rerun script, the
same as for output from any LAMMPS simulation.  See further info below
as to what that means if the timesteps for snapshots read from dump
files do not match the specified output frequency.


\bigskip\hrule\bigskip


\sphinxAtStartPar
If more than one dump file is specified, the dump files are read one
after the other in the order specified.  It is assumed that snapshot
timesteps will be in ascending order.  If a snapshot is encountered that
is not in ascending order, it will skip the snapshot until it reads one
that is.
This allows skipping of a duplicate snapshot (same timestep),
e.g. that appeared at the end of one file and beginning of the next.
However if you specify a series of dump files in an incorrect order
(with respect to the timesteps they contain), you may skip large
numbers of snapshots.

\sphinxAtStartPar
Note that the dump files specified as part of the \sphinxstyleemphasis{dump} keyword can be
parallel files, i.e. written as multiple files either per processor
and/or per snapshot.  If that is the case they will also be read in
parallel which can make the rerun command operate dramatically faster
for large systems.  See the page for the {\hyperref[\detokenize{read_dump::doc}]{\sphinxcrossref{\DUrole{doc}{read\_dump}}}} and {\hyperref[\detokenize{dump::doc}]{\sphinxcrossref{\DUrole{doc}{dump}}}} commands which describe how to read
and write parallel dump files.

\sphinxAtStartPar
The \sphinxstyleemphasis{first}, \sphinxstyleemphasis{last}, \sphinxstyleemphasis{every}, \sphinxstyleemphasis{skip} keywords determine which
snapshots are read from the dump file(s).  Snapshots are skipped until
they have a timestep \textgreater{}= \sphinxstyleemphasis{Nfirst}.  When a snapshot with a timestep \textgreater{}
\sphinxstyleemphasis{Nlast} is encountered, the rerun command finishes.  Note that
the defaults for \sphinxstyleemphasis{first} and \sphinxstyleemphasis{last} are to read all snapshots.  If the
\sphinxstyleemphasis{every} keyword is set to a value \textgreater{} 0, then only snapshots with
timesteps that are a multiple of \sphinxstyleemphasis{Nevery} are read (the first
snapshot is always read).  If \sphinxstyleemphasis{Nevery} = 0, then this criterion is
ignored, i.e. every snapshot is read that meets the other criteria.
If the \sphinxstyleemphasis{skip} keyword is used, then after the first snapshot is read,
every Nth snapshot is read, where N = \sphinxstyleemphasis{Nskip}.  E.g. if \sphinxstyleemphasis{Nskip} = 3,
then only 1 out of every 3 snapshots is read, assuming the snapshot
timestep is also consistent with the other criteria.

\begin{sphinxadmonition}{note}{Note:}
\sphinxAtStartPar
Not all dump formats contain the timestep and not all dump readers
support reading it.  In that case individual snapshots are assigned
consecutive timestep numbers starting at 1.
\end{sphinxadmonition}

\sphinxAtStartPar
The \sphinxstyleemphasis{start} and \sphinxstyleemphasis{stop} keywords do not affect which snapshots are read
from the dump file(s).  Rather, they have the same meaning that they
do for the {\hyperref[\detokenize{run::doc}]{\sphinxcrossref{\DUrole{doc}{run}}}} command.  They only need to be defined if
(a) you are using a {\hyperref[\detokenize{fix::doc}]{\sphinxcrossref{\DUrole{doc}{fix}}}} command that changes some value
over time, and (b) you want the reference point for elapsed time (from
start to stop) to be different than the \sphinxstyleemphasis{first} and \sphinxstyleemphasis{last} settings.
See the page for individual fixes to see which ones can be used
with the \sphinxstyleemphasis{start/stop} keywords.  Note that if you define neither of
the \sphinxstyleemphasis{start}/\sphinxstyleemphasis{stop} or \sphinxstyleemphasis{first}/\sphinxstyleemphasis{last} keywords, then LAMMPS treats the
pseudo run as going from 0 to a huge value (effectively infinity).
This means that any quantity that a fix scales as a fraction of
elapsed time in the run, will essentially remain at its initial value.
Also note that an error will occur if you read a snapshot from the
dump file with a timestep value larger than the \sphinxstyleemphasis{stop} setting you
have specified.

\sphinxAtStartPar
The \sphinxstyleemphasis{post} keyword can be used to minimize the output to the screen that
happens after a \sphinxstyleemphasis{rerun} command, similar to the post keyword of the
{\hyperref[\detokenize{run::doc}]{\sphinxcrossref{\DUrole{doc}{run command}}}}. It is set to \sphinxstyleemphasis{no} by default.

\sphinxAtStartPar
The \sphinxstyleemphasis{dump} keyword is required and must be the last keyword specified.
Its arguments are passed internally to the {\hyperref[\detokenize{read_dump::doc}]{\sphinxcrossref{\DUrole{doc}{read\_dump}}}}
command.  The first argument following the \sphinxstyleemphasis{dump} keyword should be
the \sphinxstyleemphasis{field1} argument of the {\hyperref[\detokenize{read_dump::doc}]{\sphinxcrossref{\DUrole{doc}{read\_dump}}}} command.  See
the {\hyperref[\detokenize{read_dump::doc}]{\sphinxcrossref{\DUrole{doc}{read\_dump}}}} page for details on the various
options it allows for extracting information from the dump file
snapshots, and for using that information to alter the LAMMPS
simulation.


\bigskip\hrule\bigskip


\sphinxAtStartPar
In general, a LAMMPS input script that uses a rerun command can
include and perform all the usual operations of an input script that
uses the {\hyperref[\detokenize{run::doc}]{\sphinxcrossref{\DUrole{doc}{run}}}} command.  There are a few exceptions and
points to consider, as discussed here.

\sphinxAtStartPar
Fixes that perform time integration, such as {\hyperref[\detokenize{fix_nve::doc}]{\sphinxcrossref{\DUrole{doc}{fix nve}}}}
or {\hyperref[\detokenize{fix_nh::doc}]{\sphinxcrossref{\DUrole{doc}{fix npt}}}} are not invoked, since no time integration
is performed.  Fixes that perturb or constrain the forces on atoms
will be invoked, just as they would during a normal run.  Examples are
{\hyperref[\detokenize{fix_indent::doc}]{\sphinxcrossref{\DUrole{doc}{fix indent}}}} and {\hyperref[\detokenize{fix_langevin::doc}]{\sphinxcrossref{\DUrole{doc}{fix langevin}}}}.
So you should think carefully as to whether that makes sense for the
manner in which you are reprocessing the dump snapshots.

\sphinxAtStartPar
If you only want the rerun script to perform an analysis that does not
involve pair interactions, such as use compute msd to calculated
displacements over time, you do not need to define a {\hyperref[\detokenize{pair_style::doc}]{\sphinxcrossref{\DUrole{doc}{pair style}}}}, which may also mean neighbor lists will not need to be
calculated which saves time.  The {\hyperref[\detokenize{comm_modify::doc}]{\sphinxcrossref{\DUrole{doc}{comm\_modify cutoff}}}} command can also be used to ensure ghost atoms are
acquired from far enough away for operations like bond and angle
evaluations, if no pair style is being used.

\sphinxAtStartPar
Every time a snapshot is read, the timestep for the simulation is reset,
as if the {\hyperref[\detokenize{reset_timestep::doc}]{\sphinxcrossref{\DUrole{doc}{reset\_timestep}}}} command were used.
This command has some restrictions as to what fixes can be defined.  See
its documentation page for details.  For example, the {\hyperref[\detokenize{fix_deposit::doc}]{\sphinxcrossref{\DUrole{doc}{fix deposit}}}} and {\hyperref[\detokenize{fix_dt_reset::doc}]{\sphinxcrossref{\DUrole{doc}{fix dt/reset}}}} fixes are in this
category.  They also make no sense to use with a rerun command.

\sphinxAtStartPar
If time\sphinxhyphen{}averaging fixes like {\hyperref[\detokenize{fix_ave_time::doc}]{\sphinxcrossref{\DUrole{doc}{fix ave/time}}}} are
used, they are invoked on timesteps that are a function of their
\sphinxstyleemphasis{Nevery}, \sphinxstyleemphasis{Nrepeat}, and \sphinxstyleemphasis{Nfreq} settings.  As an example, see the
{\hyperref[\detokenize{fix_ave_time::doc}]{\sphinxcrossref{\DUrole{doc}{fix ave/time}}}} page for details.  You must
ensure those settings are consistent with the snapshot timestamps that
are read from the dump file(s).  If an averaging fix is not invoked on
a timestep it expects to be, LAMMPS will flag an error.

\sphinxAtStartPar
The various forms of LAMMPS output, as defined by the
{\hyperref[\detokenize{thermo_style::doc}]{\sphinxcrossref{\DUrole{doc}{thermo\_style}}}}, {\hyperref[\detokenize{thermo::doc}]{\sphinxcrossref{\DUrole{doc}{thermo}}}},
{\hyperref[\detokenize{dump::doc}]{\sphinxcrossref{\DUrole{doc}{dump}}}}, and {\hyperref[\detokenize{restart::doc}]{\sphinxcrossref{\DUrole{doc}{restart}}}} commands occur with
specified frequency, e.g. every N steps.  If the timestep for a dump
snapshot is not a multiple of N, then it will be read and processed,
but no output will be produced.  If you want output for every dump
snapshot, you can simply use N=1 for an output frequency, e.g. for
thermodynamic output or new dump file output.


\bigskip\hrule\bigskip



\subsection{Restrictions}
\label{\detokenize{rerun:restrictions}}
\sphinxAtStartPar
The \sphinxstyleemphasis{rerun} command is subject to all restrictions of
the {\hyperref[\detokenize{read_dump::doc}]{\sphinxcrossref{\DUrole{doc}{read\_dump}}}} command.


\subsection{Related commands}
\label{\detokenize{rerun:related-commands}}
\sphinxAtStartPar
{\hyperref[\detokenize{read_dump::doc}]{\sphinxcrossref{\DUrole{doc}{read\_dump}}}}


\subsection{Default}
\label{\detokenize{rerun:default}}
\sphinxAtStartPar
The option defaults are first = 0, last = a huge value (effectively
infinity), start = same as first, stop = same as last, every = 0, skip
= 1, post = no;

\sphinxstepscope

\index{reset\_atoms@\spxentry{reset\_atoms}}\ignorespaces 

\section{reset\_atoms command}
\label{\detokenize{reset_atoms:reset-atoms-command}}\label{\detokenize{reset_atoms:index-0}}\label{\detokenize{reset_atoms::doc}}

\subsection{Syntax}
\label{\detokenize{reset_atoms:syntax}}
\begin{sphinxVerbatim}[commandchars=\\\{\}]
\PYG{k}{reset\PYGZus{}atoms}\PYG{+w}{ }\PYG{n}{property}\PYG{+w}{ }\PYG{n}{arguments}\PYG{+w}{ }\PYG{n}{...}
\end{sphinxVerbatim}
\begin{itemize}
\item {} 
\sphinxAtStartPar
property = \sphinxstyleemphasis{id} or \sphinxstyleemphasis{image} or \sphinxstyleemphasis{mol}

\item {} 
\sphinxAtStartPar
additional arguments depend on the property

\begin{sphinxVerbatim}[commandchars=\\\{\}]
\PYG{k}{reset\PYGZus{}atoms}\PYG{+w}{ }\PYG{n}{id}\PYG{+w}{ }\PYG{n}{keyword}\PYG{+w}{ }\PYG{n}{value}\PYG{+w}{ }\PYG{n}{...}
\end{sphinxVerbatim}
\begin{itemize}
\item {} 
\sphinxAtStartPar
zero or more keyword/value pairs can be appended

\item {} 
\sphinxAtStartPar
keyword = \sphinxstyleemphasis{sort}
\begin{sphinxalltt}
\sphinxstyleemphasis{sort} value = \sphinxstyleemphasis{yes} or \sphinxstyleemphasis{no}
\end{sphinxalltt}

\end{itemize}

\begin{sphinxVerbatim}[commandchars=\\\{\}]
\PYG{k}{reset\PYGZus{}atoms}\PYG{+w}{ }\PYG{n}{image}\PYG{+w}{ }\PYG{n}{group}\PYG{o}{\PYGZhy{}}\PYG{n}{ID}
\end{sphinxVerbatim}
\begin{itemize}
\item {} 
\sphinxAtStartPar
group\sphinxhyphen{}ID = ID of group of atoms whose image flags will be reset

\end{itemize}

\begin{sphinxVerbatim}[commandchars=\\\{\}]
\PYG{k}{reset\PYGZus{}atoms}\PYG{+w}{ }\PYG{n}{mol}\PYG{+w}{ }\PYG{n}{group}\PYG{o}{\PYGZhy{}}\PYG{n}{ID}\PYG{+w}{ }\PYG{n}{keyword}\PYG{+w}{ }\PYG{n}{value}\PYG{+w}{ }\PYG{n}{...}
\end{sphinxVerbatim}
\begin{itemize}
\item {} 
\sphinxAtStartPar
group\sphinxhyphen{}ID = ID of group of atoms whose molecule IDs will be reset

\item {} 
\sphinxAtStartPar
zero or more keyword/value pairs can be appended

\item {} 
\sphinxAtStartPar
keyword = \sphinxstyleemphasis{compress} or \sphinxstyleemphasis{offset} or \sphinxstyleemphasis{single}
\begin{sphinxalltt}
\sphinxstyleemphasis{compress} value = \sphinxstyleemphasis{yes} or \sphinxstyleemphasis{no}
\sphinxstyleemphasis{offset} value = \sphinxstyleemphasis{Noffset} \textgreater{}= \sphinxhyphen{}1
\sphinxstyleemphasis{single} value = \sphinxstyleemphasis{yes} or \sphinxstyleemphasis{no} to treat single atoms (no bonds) as molecules
\end{sphinxalltt}

\end{itemize}

\end{itemize}


\subsection{Examples}
\label{\detokenize{reset_atoms:examples}}
\begin{sphinxVerbatim}[commandchars=\\\{\}]
\PYG{k}{reset\PYGZus{}atoms}\PYG{+w}{ }\PYG{n}{id}
\PYG{k}{reset\PYGZus{}atoms}\PYG{+w}{ }\PYG{n}{id}\PYG{+w}{ }\PYG{n}{sort}\PYG{+w}{ }\PYG{n}{yes}
\PYG{k}{reset\PYGZus{}atoms}\PYG{+w}{ }\PYG{n}{image}\PYG{+w}{ }\PYG{n}{all}
\PYG{k}{reset\PYGZus{}atoms}\PYG{+w}{ }\PYG{n}{image}\PYG{+w}{ }\PYG{n}{mobile}
\PYG{k}{reset\PYGZus{}atoms}\PYG{+w}{ }\PYG{n}{mol}\PYG{+w}{ }\PYG{n}{all}
\PYG{k}{reset\PYGZus{}atoms}\PYG{+w}{ }\PYG{n}{mol}\PYG{+w}{ }\PYG{n}{all}\PYG{+w}{ }\PYG{n}{offset}\PYG{+w}{ }\PYG{l+m}{10}\PYG{+w}{ }\PYG{n}{single}\PYG{+w}{ }\PYG{n}{yes}
\PYG{k}{reset\PYGZus{}atoms}\PYG{+w}{ }\PYG{n}{mol}\PYG{+w}{ }\PYG{n}{solvent}\PYG{+w}{ }\PYG{n}{compress}\PYG{+w}{ }\PYG{n}{yes}\PYG{+w}{ }\PYG{n}{offset}\PYG{+w}{ }\PYG{l+m}{100}
\PYG{k}{reset\PYGZus{}atoms}\PYG{+w}{ }\PYG{n}{mol}\PYG{+w}{ }\PYG{n}{solvent}\PYG{+w}{ }\PYG{n}{compress}\PYG{+w}{ }\PYG{n}{no}
\end{sphinxVerbatim}


\subsection{Description}
\label{\detokenize{reset_atoms:description}}
\sphinxAtStartPar
\DUrole{versionmodified,added}{Added in version 22Dec2022.}

\sphinxAtStartPar
The \sphinxstyleemphasis{reset\_atoms} command resets the values of a specified atom
property.  In contrast to the \sphinxstyleemphasis{set} command, it does this in a
collective manner which resets the values for many atoms in a
self\sphinxhyphen{}consistent way.  This command is often useful when the simulated
system has undergone significant modifications like adding or removing
atoms or molecules, joining data files, changing bonds, or large\sphinxhyphen{}scale
diffusion.

\sphinxAtStartPar
The new values can be thought of as a \sphinxstyleemphasis{reset}, similar to values atoms
would have if a new data file were being read or a new simulation
performed.  Note that the \sphinxstyleemphasis{set} command also resets atom properties to
new values, but it treats each atom independently.

\sphinxAtStartPar
The \sphinxstyleemphasis{property} setting can be \sphinxstyleemphasis{id} or \sphinxstyleemphasis{image} or \sphinxstyleemphasis{mol}.  For \sphinxstyleemphasis{id}, the
IDs of all the atoms are reset to contiguous values.  For \sphinxstyleemphasis{image}, the
image flags of atoms in the specified \sphinxstyleemphasis{group\sphinxhyphen{}ID} are reset so that at
least one atom in each molecule is in the simulation box (image flag =
0).  For \sphinxstyleemphasis{mol}, the molecule IDs of all atoms are reset to contiguous
values.

\sphinxAtStartPar
More details on these operations and their arguments or optional
keyword/value settings are given below.


\bigskip\hrule\bigskip


\sphinxAtStartPar
Property: \sphinxstyleemphasis{id}

\sphinxAtStartPar
Reset atom IDs for the entire system, including all the global IDs
stored for bond, angle, dihedral, improper topology data.  This will
create a set of IDs that are numbered contiguously from 1 to N for a N
atoms system.

\sphinxAtStartPar
This can be useful to do after performing a “delete\_atoms” command for
a molecular system.  The delete\_atoms compress yes option will not
perform this operation due to the existence of bond topology.  It can
also be useful to do after any simulation which has lost atoms,
e.g. due to atoms moving outside a simulation box with fixed
boundaries (see the “boundary command”), or due to evaporation (see
the “fix evaporate” command).

\sphinxAtStartPar
If the \sphinxstyleemphasis{sort} keyword is used with a setting of \sphinxstyleemphasis{yes}, then the
assignment of new atom IDs will be the same no matter how many
processors LAMMPS is running on.  This is done by first doing a
spatial sort of all the atoms into bins and sorting them within each
bin.  Because the set of bins is independent of the number of
processors, this enables a consistent assignment of new IDs to each
atom.

\sphinxAtStartPar
This can be useful to do after using the “create\_atoms” command and/or
“replicate” command.  In general those commands do not guarantee
assignment of the same atom ID to the same physical atom when LAMMPS
is run on different numbers of processors.  Enforcing consistent IDs
can be useful for debugging or comparing output from two different
runs.

\sphinxAtStartPar
Note that the spatial sort requires communication of atom IDs and
coordinates between processors in an all\sphinxhyphen{}to\sphinxhyphen{}all manner.  This is done
efficiently in LAMMPS, but it is more expensive than how atom IDs are
reset without sorting.

\sphinxAtStartPar
Note that whether sorting or not, the resetting of IDs is not a
compression, where gaps in atom IDs are removed by decrementing atom
IDs that are larger.  Instead the IDs for all atoms are erased, and
new IDs are assigned so that the atoms owned by an individual
processor have consecutive IDs, as the {\hyperref[\detokenize{create_atoms::doc}]{\sphinxcrossref{\DUrole{doc}{create\_atoms}}}} command explains.

\begin{sphinxadmonition}{note}{Note:}
\sphinxAtStartPar
If this command is used before a {\hyperref[\detokenize{pair_style::doc}]{\sphinxcrossref{\DUrole{doc}{pair style}}}} is
defined, an error about bond topology atom IDs not being found may
result.  This is because the cutoff distance for ghost atom
communication was not sufficient to find atoms in bonds, angles, etc
that are owned by other processors.  The {\hyperref[\detokenize{comm_modify::doc}]{\sphinxcrossref{\DUrole{doc}{comm\_modify cutoff}}}} command can be used to correct this issue.  Or you can
define a pair style before using this command.  If you do the former,
you should unset the \sphinxstyleemphasis{comm\_modify cutoff} after using \sphinxstyleemphasis{reset
atoms id} so that subsequent communication is not inefficient.
\end{sphinxadmonition}


\bigskip\hrule\bigskip


\sphinxAtStartPar
Property: \sphinxstyleemphasis{image}

\sphinxAtStartPar
Reset the image flags of atoms so that at least one atom in each
molecule has an image flag of 0.  Molecular topology is respected so
that if the molecule straddles a periodic simulation box boundary, the
images flags of all atoms in the molecule will be consistent.  This
avoids inconsistent image flags that could result from resetting all
image flags to zero with the {\hyperref[\detokenize{set::doc}]{\sphinxcrossref{\DUrole{doc}{set}}}} command.

\begin{sphinxadmonition}{note}{Note:}
\sphinxAtStartPar
If the system has no bonds, there is no reason to use this command,
since image flags for different atoms do not need to be
consistent.  Use the {\hyperref[\detokenize{set::doc}]{\sphinxcrossref{\DUrole{doc}{set}}}} command with its \sphinxstyleemphasis{image}
keyword instead.
\end{sphinxadmonition}

\sphinxAtStartPar
Only image flags for atoms in the specified \sphinxstyleemphasis{group\sphinxhyphen{}ID} are reset; all
others remain unchanged.  No check is made for whether the group
covers complete molecule fragments and thus whether the command will
result in inconsistent image flags.

\sphinxAtStartPar
Molecular fragments are identified by the algorithm used by the
{\hyperref[\detokenize{compute_cluster_atom::doc}]{\sphinxcrossref{\DUrole{doc}{compute fragment/atom}}}} command.  For each
fragment the average of the largest and the smallest image flag in
each direction across all atoms in the fragment is computed and
subtracted from the current image flag in the same direction.

\sphinxAtStartPar
This can be a useful operation to perform after running longer
equilibration runs of mobile systems where molecules would pass
through the system multiple times and thus produce non\sphinxhyphen{}zero image
flags.

\begin{sphinxadmonition}{note}{Note:}
\sphinxAtStartPar
Same as explained for the {\hyperref[\detokenize{compute_cluster_atom::doc}]{\sphinxcrossref{\DUrole{doc}{compute fragment/atom}}}} command, molecules are identified using the
current bond topology.  This will \sphinxstylestrong{not} account for bonds broken by
the {\hyperref[\detokenize{bond_quartic::doc}]{\sphinxcrossref{\DUrole{doc}{bond\_style quartic}}}} command, because this
bond style does not perform a full update of the bond topology data
structures within LAMMPS.  In that case, using the {\hyperref[\detokenize{delete_bonds::doc}]{\sphinxcrossref{\DUrole{doc}{delete\_bonds
all bond 0 remove}}}} will permanently delete such
broken bonds and should thus be used first.
\end{sphinxadmonition}


\bigskip\hrule\bigskip


\sphinxAtStartPar
Property: \sphinxstyleemphasis{mol}

\sphinxAtStartPar
Reset molecule IDs for a specified group of atoms based on current
bond connectivity.  This will typically create a new set of molecule
IDs for atoms in the group.  Only molecule IDs for atoms in the
specified \sphinxstyleemphasis{group\sphinxhyphen{}ID} are reset; molecule IDs for atoms not in the
group are not changed.

\sphinxAtStartPar
For purposes of this operation, molecules are identified by the current
bond connectivity in the system, which may or may not be consistent with
the current molecule IDs.  A molecule in this context is a set of atoms
connected to each other with explicit bonds.  The specific algorithm
used is the one of {\hyperref[\detokenize{compute_cluster_atom::doc}]{\sphinxcrossref{\DUrole{doc}{compute fragment/atom}}}}.
Once the molecules are identified and a new molecule ID computed for
each, this command will update the current molecule ID for all atoms in
the group with the new molecule ID.  Note that if the group excludes
atoms within molecules, one (physical) molecule may become two or more
(logical) molecules.  For example if the group excludes atoms in the
middle of a linear chain, then each end of the chain is considered an
independent molecule and will be assigned a different molecule ID.

\sphinxAtStartPar
This can be a useful operation to perform after running reactive
molecular dynamics run with {\hyperref[\detokenize{fix_bond_react::doc}]{\sphinxcrossref{\DUrole{doc}{fix bond/react}}}},
{\hyperref[\detokenize{fix_bond_create::doc}]{\sphinxcrossref{\DUrole{doc}{fix bond/create}}}}, or {\hyperref[\detokenize{fix_bond_break::doc}]{\sphinxcrossref{\DUrole{doc}{fix bond/break}}}}, all of which can change molecule topologies. It can
also be useful after molecules have been deleted with the
{\hyperref[\detokenize{delete_atoms::doc}]{\sphinxcrossref{\DUrole{doc}{delete\_atoms}}}} command or after a simulation which
has lost molecules, e.g. via the {\hyperref[\detokenize{fix_evaporate::doc}]{\sphinxcrossref{\DUrole{doc}{fix evaporate}}}}
command.

\sphinxAtStartPar
The \sphinxstyleemphasis{compress} keyword determines how new molecule IDs are computed.  If
the setting is \sphinxstyleemphasis{yes} (the default) and there are N molecules in the
group, the new molecule IDs will be a set of N contiguous values.  See
the \sphinxstyleemphasis{offset} keyword for details on selecting the range of these values.
If the setting is \sphinxstyleemphasis{no}, the molecule ID of every atom in the molecule
will be set to the smallest atom ID of any atom in the molecule.

\sphinxAtStartPar
The \sphinxstyleemphasis{single} keyword determines whether single atoms (not bonded to
another atom) are treated as one\sphinxhyphen{}atom molecules or not, based on the
\sphinxstyleemphasis{yes} or \sphinxstyleemphasis{no} setting.  If the setting is \sphinxstyleemphasis{no} (the default), their
molecule IDs are set to 0.  This setting can be important if the new
molecule IDs will be used as input to other commands such as
{\hyperref[\detokenize{compute_chunk_atom::doc}]{\sphinxcrossref{\DUrole{doc}{compute chunk/atom molecule}}}} or {\hyperref[\detokenize{fix_rigid::doc}]{\sphinxcrossref{\DUrole{doc}{fix
rigid molecule}}}}.

\sphinxAtStartPar
The \sphinxstyleemphasis{offset} keyword is only used if the \sphinxstyleemphasis{compress} setting is \sphinxstyleemphasis{yes}.
Its default value is \sphinxstyleemphasis{Noffset} = \sphinxhyphen{}1.  In that case, if the specified
group is \sphinxstyleemphasis{all}, then the new compressed molecule IDs will range from 1
to N.  If the specified group is not \sphinxstyleemphasis{all} and the largest molecule ID
of atoms outside that group is M, then the new compressed molecule IDs will
range from M+1 to M+N, to avoid collision with existing molecule
IDs.  If an \sphinxstyleemphasis{Noffset} \textgreater{}= 0 is specified, then the new compressed
molecule IDs will range from \sphinxstyleemphasis{Noffset}+1 to \sphinxstyleemphasis{Noffset}+N.  If the group
is not \sphinxstyleemphasis{all} there may be collisions with the molecule IDs of other atoms.

\begin{sphinxadmonition}{note}{Note:}
\sphinxAtStartPar
Same as explained for the {\hyperref[\detokenize{compute_cluster_atom::doc}]{\sphinxcrossref{\DUrole{doc}{compute fragment/atom}}}} command, molecules are identified using the
current bond topology.  This will \sphinxstylestrong{not} account for bonds broken by
the {\hyperref[\detokenize{bond_quartic::doc}]{\sphinxcrossref{\DUrole{doc}{bond\_style quartic}}}} command, because this
bond style does not perform a full update of the bond topology data
structures within LAMMPS.  In that case, using the {\hyperref[\detokenize{delete_bonds::doc}]{\sphinxcrossref{\DUrole{doc}{delete\_bonds
all bond 0 remove}}}} will permanently delete such broken
bonds and should thus be used first.
\end{sphinxadmonition}


\subsection{Restrictions}
\label{\detokenize{reset_atoms:restrictions}}
\sphinxAtStartPar
The \sphinxstyleemphasis{image} property can only be used when the atom style supports bonds.


\subsection{Related commands}
\label{\detokenize{reset_atoms:related-commands}}
\sphinxAtStartPar
{\hyperref[\detokenize{compute_cluster_atom::doc}]{\sphinxcrossref{\DUrole{doc}{compute fragment/atom}}}},
{\hyperref[\detokenize{fix_bond_react::doc}]{\sphinxcrossref{\DUrole{doc}{fix bond/react}}}},
{\hyperref[\detokenize{fix_bond_create::doc}]{\sphinxcrossref{\DUrole{doc}{fix bond/create}}}},
{\hyperref[\detokenize{fix_bond_break::doc}]{\sphinxcrossref{\DUrole{doc}{fix bond/break}}}},
{\hyperref[\detokenize{fix_evaporate::doc}]{\sphinxcrossref{\DUrole{doc}{fix evaporate}}}},
{\hyperref[\detokenize{delete_atoms::doc}]{\sphinxcrossref{\DUrole{doc}{delete\_atoms}}}},
{\hyperref[\detokenize{delete_bonds::doc}]{\sphinxcrossref{\DUrole{doc}{delete\_bonds}}}}


\subsection{Defaults}
\label{\detokenize{reset_atoms:defaults}}
\sphinxAtStartPar
For property \sphinxstyleemphasis{id}, the default keyword setting is sort = no.

\sphinxAtStartPar
For property \sphinxstyleemphasis{mol}, the default keyword settings are compress = yes,
single = no, and offset = \sphinxhyphen{}1.

\sphinxstepscope

\index{reset\_timestep@\spxentry{reset\_timestep}}\ignorespaces 

\section{reset\_timestep command}
\label{\detokenize{reset_timestep:reset-timestep-command}}\label{\detokenize{reset_timestep:index-0}}\label{\detokenize{reset_timestep::doc}}

\subsection{Syntax}
\label{\detokenize{reset_timestep:syntax}}
\begin{sphinxVerbatim}[commandchars=\\\{\}]
\PYG{k}{reset\PYGZus{}timestep}\PYG{+w}{ }\PYG{n}{N}\PYG{+w}{ }\PYG{n}{keyword}\PYG{+w}{ }\PYG{n}{values}\PYG{+w}{ }\PYG{n}{...}
\end{sphinxVerbatim}
\begin{itemize}
\item {} 
\sphinxAtStartPar
N = timestep number

\item {} 
\sphinxAtStartPar
zero or more keyword/value pairs may be appended

\item {} 
\sphinxAtStartPar
keyword = \sphinxstyleemphasis{time}
\begin{sphinxalltt}
\sphinxstyleemphasis{time} value = atime
   atime = accumulated simulation time
\end{sphinxalltt}

\end{itemize}


\subsection{Examples}
\label{\detokenize{reset_timestep:examples}}
\begin{sphinxVerbatim}[commandchars=\\\{\}]
\PYG{k}{reset\PYGZus{}timestep}\PYG{+w}{ }\PYG{l+m}{0}
\PYG{k}{reset\PYGZus{}timestep}\PYG{+w}{ }\PYG{l+m}{4000000}
\PYG{k}{reset\PYGZus{}timestep}\PYG{+w}{ }\PYG{l+m}{1000}\PYG{+w}{ }\PYG{n}{time}\PYG{+w}{ }\PYG{l+m}{100.0}
\end{sphinxVerbatim}


\subsection{Description}
\label{\detokenize{reset_timestep:description}}
\sphinxAtStartPar
Set the timestep counter to the specified value.  This command
usually comes after the timestep has been set by reading a restart
file via the {\hyperref[\detokenize{read_restart::doc}]{\sphinxcrossref{\DUrole{doc}{read\_restart}}}} command, or a previous
simulation run or minimization advanced the timestep.

\sphinxAtStartPar
The optional \sphinxstyleemphasis{time} keyword allows to also set the accumulated
simulation time.  This is usually the number of timesteps times
the size of the timestep, but when using variable size timesteps
with {\hyperref[\detokenize{fix_dt_reset::doc}]{\sphinxcrossref{\DUrole{doc}{fix dt/reset}}}} it can differ.

\sphinxAtStartPar
The {\hyperref[\detokenize{read_data::doc}]{\sphinxcrossref{\DUrole{doc}{read\_data}}}} and {\hyperref[\detokenize{create_box::doc}]{\sphinxcrossref{\DUrole{doc}{create\_box}}}}
commands set the timestep to 0; the {\hyperref[\detokenize{read_restart::doc}]{\sphinxcrossref{\DUrole{doc}{read\_restart}}}}
command sets the timestep to the value it had when the restart file
was written.  The same applies to the accumulated simulation time.


\subsection{Restrictions}
\label{\detokenize{reset_timestep:restrictions}}
\sphinxAtStartPar
This command cannot be used when any fixes are defined that keep track
of elapsed time to perform certain kinds of time\sphinxhyphen{}dependent operations.
Examples are the {\hyperref[\detokenize{fix_deposit::doc}]{\sphinxcrossref{\DUrole{doc}{fix deposit}}}} and {\hyperref[\detokenize{fix_dt_reset::doc}]{\sphinxcrossref{\DUrole{doc}{fix dt/reset}}}} commands.  The former adds atoms on specific timesteps.
The latter keeps track of accumulated time.

\sphinxAtStartPar
Various fixes use the current timestep to calculate related quantities.
If the timestep is reset, this may produce unexpected behavior, but
LAMMPS allows the fixes to be defined even if the timestep is reset.
For example, commands which thermostat the system, e.g. {\hyperref[\detokenize{fix_nh::doc}]{\sphinxcrossref{\DUrole{doc}{fix nvt}}}}, allow you to specify a target temperature which ramps from
Tstart to Tstop which may persist over several runs.  If you change the
timestep, you may induce an instantaneous change in the target
temperature.

\sphinxAtStartPar
Resetting the timestep clears flags for {\hyperref[\detokenize{compute::doc}]{\sphinxcrossref{\DUrole{doc}{computes}}}} that
may have calculated some quantity from a previous run.  This means these
quantity cannot be accessed by a variable in between runs until a new
run is performed.  See the {\hyperref[\detokenize{variable::doc}]{\sphinxcrossref{\DUrole{doc}{variable}}}} command for more
details.


\subsection{Related commands}
\label{\detokenize{reset_timestep:related-commands}}
\sphinxAtStartPar
{\hyperref[\detokenize{rerun::doc}]{\sphinxcrossref{\DUrole{doc}{rerun}}}}, {\hyperref[\detokenize{timestep::doc}]{\sphinxcrossref{\DUrole{doc}{timestep}}}},
{\hyperref[\detokenize{fix_dt_reset::doc}]{\sphinxcrossref{\DUrole{doc}{fix dt/reset}}}}


\subsection{Default}
\label{\detokenize{reset_timestep:default}}
\sphinxAtStartPar
none

\sphinxstepscope

\index{restart@\spxentry{restart}}\ignorespaces 

\section{restart command}
\label{\detokenize{restart:restart-command}}\label{\detokenize{restart:index-0}}\label{\detokenize{restart::doc}}

\subsection{Syntax}
\label{\detokenize{restart:syntax}}
\begin{sphinxVerbatim}[commandchars=\\\{\}]
\PYG{k}{restart}\PYG{+w}{ }\PYG{l+m}{0}
\PYG{k}{restart}\PYG{+w}{ }\PYG{n}{N}\PYG{+w}{ }\PYG{n}{root}\PYG{+w}{ }\PYG{n}{keyword}\PYG{+w}{ }\PYG{n}{value}\PYG{+w}{ }\PYG{n}{...}
\PYG{k}{restart}\PYG{+w}{ }\PYG{n}{N}\PYG{+w}{ }\PYG{n}{file1}\PYG{+w}{ }\PYG{n}{file2}\PYG{+w}{ }\PYG{n}{keyword}\PYG{+w}{ }\PYG{n}{value}\PYG{+w}{ }\PYG{n}{...}
\end{sphinxVerbatim}
\begin{itemize}
\item {} 
\sphinxAtStartPar
N = write a restart file on timesteps which are multiples of N

\item {} 
\sphinxAtStartPar
N can be a variable (see below)

\item {} 
\sphinxAtStartPar
root = filename to which timestep \# is appended

\item {} 
\sphinxAtStartPar
file1,file2 = two full filenames, toggle between them when writing file

\item {} 
\sphinxAtStartPar
zero or more keyword/value pairs may be appended

\item {} 
\sphinxAtStartPar
keyword = \sphinxstyleemphasis{fileper} or \sphinxstyleemphasis{nfile}
\begin{sphinxalltt}
\sphinxstyleemphasis{fileper} arg = Np
  Np = write one file for every this many processors
\sphinxstyleemphasis{nfile} arg = Nf
  Nf = write this many files, one from each of Nf processors
\end{sphinxalltt}

\end{itemize}


\subsection{Examples}
\label{\detokenize{restart:examples}}
\begin{sphinxVerbatim}[commandchars=\\\{\}]
\PYG{k}{restart}\PYG{+w}{ }\PYG{l+m}{0}
\PYG{k}{restart}\PYG{+w}{ }\PYG{l+m}{1000}\PYG{+w}{ }\PYG{n}{poly.restart}
\PYG{k}{restart}\PYG{+w}{ }\PYG{l+m}{1000}\PYG{+w}{ }\PYG{n}{restart.}\PYG{o}{*}\PYG{n}{.equil}
\PYG{k}{restart}\PYG{+w}{ }\PYG{l+m}{10000}\PYG{+w}{ }\PYG{n}{poly.}\PYG{o}{\PYGZpc{}}\PYG{n}{.1}\PYG{+w}{ }\PYG{n}{poly.}\PYG{o}{\PYGZpc{}}\PYG{n}{.2}\PYG{+w}{ }\PYG{n}{nfile}\PYG{+w}{ }\PYG{l+m}{10}
\PYG{k}{restart}\PYG{+w}{ }\PYG{n}{v\PYGZus{}mystep}\PYG{+w}{ }\PYG{n}{poly.restart}
\end{sphinxVerbatim}


\subsection{Description}
\label{\detokenize{restart:description}}
\sphinxAtStartPar
Write out a binary restart file with the current state of the
simulation on timesteps which are a multiple of N.  A value of N = 0
means do not write out any restart files, which is the default.
Restart files are written in one (or both) of two modes as a run
proceeds.  If one filename is specified, a series of filenames will be
created which include the timestep in the filename.  If two filenames
are specified, only 2 restart files will be created, with those names.
LAMMPS will toggle between the 2 names as it writes successive restart
files.

\sphinxAtStartPar
Note that you can specify the restart command twice, once with a
single filename and once with two filenames.  This would allow you,
for example, to write out archival restart files every 100000 steps
using a single filename, and more frequent temporary restart files
every 1000 steps, using two filenames.  Using restart 0 will turn off
both modes of output.

\sphinxAtStartPar
Similar to {\hyperref[\detokenize{dump::doc}]{\sphinxcrossref{\DUrole{doc}{dump}}}} files, the restart filename(s) can contain
two wild\sphinxhyphen{}card characters.

\sphinxAtStartPar
If a “*” appears in the single filename, it is replaced with the
current timestep value.  This is only recognized when a single
filename is used (not when toggling back and forth).  Thus, the third
example above creates restart files as follows: restart.1000.equil,
restart.2000.equil, etc.  If a single filename is used with no “*”,
then the timestep value is appended.  E.g. the second example above
creates restart files as follows: poly.restart.1000,
poly.restart.2000, etc.

\sphinxAtStartPar
If a “\%” character appears in the restart filename(s), then one file
is written for each processor and the “\%” character is replaced with
the processor ID from 0 to P\sphinxhyphen{}1.  An additional file with the “\%”
replaced by “base” is also written, which contains global information.
For example, the files written on step 1000 for filename restart.\%
would be restart.base.1000, restart.0.1000, restart.1.1000, …,
restart.P\sphinxhyphen{}1.1000.  This creates smaller files and can be a fast mode
of output and subsequent input on parallel machines that support
parallel I/O.  The optional \sphinxstyleemphasis{fileper} and \sphinxstyleemphasis{nfile} keywords discussed
below can alter the number of files written.

\sphinxAtStartPar
Restart files are written on timesteps that are a multiple of N but
not on the first timestep of a run or minimization.  You can use the
{\hyperref[\detokenize{write_restart::doc}]{\sphinxcrossref{\DUrole{doc}{write\_restart}}}} command to write a restart file
before a run begins.  A restart file is not written on the last
timestep of a run unless it is a multiple of N.  A restart file is
written on the last timestep of a minimization if N \textgreater{} 0 and the
minimization converges.

\sphinxAtStartPar
Instead of a numeric value, N can be specified as an {\hyperref[\detokenize{variable::doc}]{\sphinxcrossref{\DUrole{doc}{equal\sphinxhyphen{}style
variable}}}}, which should be specified as v\_name, where name is
the variable name.  In this case, the variable is evaluated at the
beginning of a run to determine the next timestep at which a restart
file will be written out.  On that timestep, the variable will be
evaluated again to determine the next timestep, etc.  Thus the variable
should return timestep values.  See the stagger() and logfreq() and
stride() math functions for {\hyperref[\detokenize{variable::doc}]{\sphinxcrossref{\DUrole{doc}{equal\sphinxhyphen{}style variables}}}}, as
examples of useful functions to use in this context.  Other similar math
functions could easily be added as options for {\hyperref[\detokenize{variable::doc}]{\sphinxcrossref{\DUrole{doc}{equal\sphinxhyphen{}style
variables}}}}.

\sphinxAtStartPar
For example, the following commands will write restart files
every step from 1100 to 1200, and could be useful for debugging
a simulation where something goes wrong at step 1163:

\begin{sphinxVerbatim}[commandchars=\\\{\}]
\PYG{k}{variable       }\PYG{n+nv+nvIdentifier}{s}\PYG{+w}{ }\PYG{n}{equal}\PYG{+w}{ }\PYG{n}{stride}\PYG{n+nv}{(}\PYG{n+nv}{1100,1200,1}\PYG{n+nv}{)}
\PYG{k}{restart}\PYG{+w}{        }\PYG{n}{v\PYGZus{}s}\PYG{+w}{ }\PYG{n}{tmp.restart}
\end{sphinxVerbatim}


\bigskip\hrule\bigskip


\sphinxAtStartPar
See the {\hyperref[\detokenize{read_restart::doc}]{\sphinxcrossref{\DUrole{doc}{read\_restart}}}} command for information about
what is stored in a restart file.

\sphinxAtStartPar
Restart files can be read by a {\hyperref[\detokenize{read_restart::doc}]{\sphinxcrossref{\DUrole{doc}{read\_restart}}}}
command to restart a simulation from a particular state.  Because the
file is binary (to enable exact restarts), it may not be readable on
another machine.  In this case, you can use the {\hyperref[\detokenize{Run_options::doc}]{\sphinxcrossref{\DUrole{doc}{\sphinxhyphen{}r command\sphinxhyphen{}line switch}}}} to convert a restart file to a data file.

\begin{sphinxadmonition}{note}{Note:}
\sphinxAtStartPar
Although the purpose of restart files is to enable restarting a
simulation from where it left off, not all information about a
simulation is stored in the file.  For example, the list of fixes that
were specified during the initial run is not stored, which means the
new input script must specify any fixes you want to use.  Even when
restart information is stored in the file, as it is for some fixes,
commands may need to be re\sphinxhyphen{}specified in the new input script, in order
to re\sphinxhyphen{}use that information.  See the {\hyperref[\detokenize{read_restart::doc}]{\sphinxcrossref{\DUrole{doc}{read\_restart}}}}
command for information about what is stored in a restart file.
\end{sphinxadmonition}


\bigskip\hrule\bigskip


\sphinxAtStartPar
The optional \sphinxstyleemphasis{nfile} or \sphinxstyleemphasis{fileper} keywords can be used in conjunction
with the “\%” wildcard character in the specified restart file name(s).
As explained above, the “\%” character causes the restart file to be
written in pieces, one piece for each of P processors.  By default P =
the number of processors the simulation is running on.  The \sphinxstyleemphasis{nfile} or
\sphinxstyleemphasis{fileper} keyword can be used to set P to a smaller value, which can
be more efficient when running on a large number of processors.

\sphinxAtStartPar
The \sphinxstyleemphasis{nfile} keyword sets P to the specified Nf value.  For example, if
Nf = 4, and the simulation is running on 100 processors, 4 files will
be written, by processors 0,25,50,75.  Each will collect information
from itself and the next 24 processors and write it to a restart file.

\sphinxAtStartPar
For the \sphinxstyleemphasis{fileper} keyword, the specified value of Np means write one
file for every Np processors.  For example, if Np = 4, every fourth
processor (0,4,8,12,etc) will collect information from itself and the
next 3 processors and write it to a restart file.


\bigskip\hrule\bigskip



\subsection{Restrictions}
\label{\detokenize{restart:restrictions}}
\sphinxAtStartPar
none


\subsection{Related commands}
\label{\detokenize{restart:related-commands}}
\sphinxAtStartPar
{\hyperref[\detokenize{write_restart::doc}]{\sphinxcrossref{\DUrole{doc}{write\_restart}}}}, {\hyperref[\detokenize{read_restart::doc}]{\sphinxcrossref{\DUrole{doc}{read\_restart}}}}


\subsection{Default}
\label{\detokenize{restart:default}}
\begin{sphinxVerbatim}[commandchars=\\\{\}]
\PYG{k}{restart}\PYG{+w}{ }\PYG{l+m}{0}
\end{sphinxVerbatim}

\sphinxstepscope

\index{run@\spxentry{run}}\ignorespaces 

\section{run command}
\label{\detokenize{run:run-command}}\label{\detokenize{run:index-0}}\label{\detokenize{run::doc}}

\subsection{Syntax}
\label{\detokenize{run:syntax}}
\begin{sphinxVerbatim}[commandchars=\\\{\}]
\PYG{k}{run}\PYG{+w}{ }\PYG{n}{N}\PYG{+w}{ }\PYG{n}{keyword}\PYG{+w}{ }\PYG{n}{values}\PYG{+w}{ }\PYG{n}{...}
\end{sphinxVerbatim}
\begin{itemize}
\item {} 
\sphinxAtStartPar
N = \# of timesteps

\item {} 
\sphinxAtStartPar
zero or more keyword/value pairs may be appended

\item {} 
\sphinxAtStartPar
keyword = \sphinxstyleemphasis{upto} or \sphinxstyleemphasis{start} or \sphinxstyleemphasis{stop} or \sphinxstyleemphasis{pre} or \sphinxstyleemphasis{post} or \sphinxstyleemphasis{every}
\begin{sphinxalltt}
\sphinxstyleemphasis{upto} value = none
\sphinxstyleemphasis{start} value = N1
  N1 = timestep at which 1st run started
\sphinxstyleemphasis{stop} value = N2
  N2 = timestep at which last run will end
\sphinxstyleemphasis{pre} value = \sphinxstyleemphasis{no} or \sphinxstyleemphasis{yes}
\sphinxstyleemphasis{post} value = \sphinxstyleemphasis{no} or \sphinxstyleemphasis{yes}
\sphinxstyleemphasis{every} values = M c1 c2 ...
  M = break the run into M\sphinxhyphen{}timestep segments and invoke one or more commands between each segment
  c1,c2,...,cN = one or more LAMMPS commands, each enclosed in quotes
  c1 = NULL means no command will be invoked
\end{sphinxalltt}

\end{itemize}


\subsection{Examples}
\label{\detokenize{run:examples}}
\begin{sphinxVerbatim}[commandchars=\\\{\}]
\PYG{k}{run}\PYG{+w}{ }\PYG{l+m}{10000}
\PYG{k}{run}\PYG{+w}{ }\PYG{l+m}{1000000}\PYG{+w}{ }\PYG{n}{upto}
\PYG{k}{run}\PYG{+w}{ }\PYG{l+m}{100}\PYG{+w}{ }\PYG{n}{start}\PYG{+w}{ }\PYG{l+m}{0}\PYG{+w}{ }\PYG{n}{stop}\PYG{+w}{ }\PYG{l+m}{1000}
\PYG{k}{run}\PYG{+w}{ }\PYG{l+m}{1000}\PYG{+w}{ }\PYG{n}{pre}\PYG{+w}{ }\PYG{n}{no}\PYG{+w}{ }\PYG{n}{post}\PYG{+w}{ }\PYG{n}{yes}
\PYG{k}{run}\PYG{+w}{ }\PYG{l+m}{100000}\PYG{+w}{ }\PYG{n}{start}\PYG{+w}{ }\PYG{l+m}{0}\PYG{+w}{ }\PYG{n}{stop}\PYG{+w}{ }\PYG{l+m}{1000000}\PYG{+w}{ }\PYG{n}{every}\PYG{+w}{ }\PYG{l+m}{1000}\PYG{+w}{ }\PYG{l+s}{\PYGZdq{}}\PYG{l+s}{print \PYGZsq{}Protein Rg = \PYGZdl{}r\PYGZsq{}}\PYG{l+s}{\PYGZdq{}}
\PYG{k}{run}\PYG{+w}{ }\PYG{l+m}{100000}\PYG{+w}{ }\PYG{n}{every}\PYG{+w}{ }\PYG{l+m}{1000}\PYG{+w}{ }\PYG{n}{NULL}
\end{sphinxVerbatim}


\subsection{Description}
\label{\detokenize{run:description}}
\sphinxAtStartPar
Run or continue dynamics for a specified number of timesteps.

\sphinxAtStartPar
When the {\hyperref[\detokenize{run_style::doc}]{\sphinxcrossref{\DUrole{doc}{run style}}}} is \sphinxstyleemphasis{respa}, N refers to outer
loop (largest) timesteps.

\sphinxAtStartPar
A value of N = 0 is acceptable; only the thermodynamics of the system
are computed and printed without taking a timestep.

\sphinxAtStartPar
The \sphinxstyleemphasis{upto} keyword means to perform a run starting at the current
timestep up to the specified timestep.  E.g. if the current timestep
is 10,000 and “run 100000 upto” is used, then an additional 90,000
timesteps will be run.  This can be useful for very long runs on a
machine that allocates chunks of time and terminate your job when time
is exceeded.  If you need to restart your script multiple times
(reading in the last restart file), you can keep restarting your
script with the same run command until the simulation finally
completes.

\sphinxAtStartPar
The \sphinxstyleemphasis{start} or \sphinxstyleemphasis{stop} keywords can be used if multiple runs are being
performed and you want a {\hyperref[\detokenize{fix::doc}]{\sphinxcrossref{\DUrole{doc}{fix}}}} command that changes some
value over time (e.g. temperature) to make the change across the
entire set of runs and not just a single run.  See the page for
individual fixes to see which ones can be used with the \sphinxstyleemphasis{start/stop}
keywords.

\sphinxAtStartPar
For example, consider this fix followed by 10 run commands:

\begin{sphinxVerbatim}[commandchars=\\\{\}]
\PYG{k}{fix         }\PYG{n+nv+nvIdentifier}{1}\PYG{+w}{ }\PYG{n+nv+nvIdentifier}{all}\PYG{+w}{ }\PYG{n}{nvt}\PYG{+w}{ }\PYG{l+m}{200.0}\PYG{+w}{ }\PYG{l+m}{300.0}\PYG{+w}{ }\PYG{l+m}{1.0}
\PYG{k}{run}\PYG{+w}{         }\PYG{l+m}{1000}\PYG{+w}{ }\PYG{n}{start}\PYG{+w}{ }\PYG{l+m}{0}\PYG{+w}{ }\PYG{n}{stop}\PYG{+w}{ }\PYG{l+m}{10000}
\PYG{k}{run}\PYG{+w}{         }\PYG{l+m}{1000}\PYG{+w}{ }\PYG{n}{start}\PYG{+w}{ }\PYG{l+m}{0}\PYG{+w}{ }\PYG{n}{stop}\PYG{+w}{ }\PYG{l+m}{10000}
\PYG{n}{...}
\PYG{k}{run}\PYG{+w}{         }\PYG{l+m}{1000}\PYG{+w}{ }\PYG{n}{start}\PYG{+w}{ }\PYG{l+m}{0}\PYG{+w}{ }\PYG{n}{stop}\PYG{+w}{ }\PYG{l+m}{10000}
\end{sphinxVerbatim}

\sphinxAtStartPar
The NVT fix ramps the target temperature from 200.0 to 300.0 during a
run.  If the run commands did not have the start/stop keywords (just
“run 1000”), then the temperature would ramp from 200.0 to 300.0
during the 1000 steps of each run.  With the start/stop keywords, the
ramping takes place over the 10000 steps of all runs together.

\sphinxAtStartPar
The \sphinxstyleemphasis{pre} and \sphinxstyleemphasis{post} keywords can be used to streamline the setup,
clean\sphinxhyphen{}up, and associated output to the screen that happens before and
after a run.  This can be useful if you wish to do many short runs in
succession (e.g. LAMMPS is being called as a library which is doing
other computations between successive short LAMMPS runs).

\sphinxAtStartPar
By default (pre and post = yes), LAMMPS creates neighbor lists,
computes forces, and imposes fix constraints before every run.  And
after every run it gathers and prints timings statistics.  If a run is
just a continuation of a previous run (i.e. no settings are changed),
the initial computation is not necessary; the old neighbor list is
still valid as are the forces.  So if \sphinxstyleemphasis{pre} is specified as “no” then
the initial setup is skipped, except for printing thermodynamic info.
Note that if \sphinxstyleemphasis{pre} is set to “no” for the very first run LAMMPS
performs, then it is overridden, since the initial setup computations
must be done.

\begin{sphinxadmonition}{note}{Note:}
\sphinxAtStartPar
If your input script changes the system between 2 runs, then the
initial setup must be performed to ensure the change is recognized by
all parts of the code that are affected.  Examples are adding a
{\hyperref[\detokenize{fix::doc}]{\sphinxcrossref{\DUrole{doc}{fix}}}} or {\hyperref[\detokenize{dump::doc}]{\sphinxcrossref{\DUrole{doc}{dump}}}} or {\hyperref[\detokenize{compute::doc}]{\sphinxcrossref{\DUrole{doc}{compute}}}}, changing
a {\hyperref[\detokenize{neigh_modify::doc}]{\sphinxcrossref{\DUrole{doc}{neighbor}}}} list parameter, or writing restart file
which can migrate atoms between processors.  LAMMPS has no easy way to
check if this has happened, but it is an error to use the \sphinxstyleemphasis{pre no}
option in this case.
\end{sphinxadmonition}

\sphinxAtStartPar
If \sphinxstyleemphasis{post} is specified as “no”, the full timing summary is skipped;
only a one\sphinxhyphen{}line summary timing is printed.

\sphinxAtStartPar
The \sphinxstyleemphasis{every} keyword provides a means of breaking a LAMMPS run into a
series of shorter runs.  Optionally, one or more LAMMPS commands (c1,
c2, …, cN) will be executed in between the short runs.  If used, the
\sphinxstyleemphasis{every} keyword must be the last keyword, since it has a variable
number of arguments.  Each of the trailing arguments is a single
LAMMPS command, and each command should be enclosed in quotes, so that
the entire command will be treated as a single argument.  This will
also prevent any variables in the command from being evaluated until
it is executed multiple times during the run.  Note that if a command
itself needs one of its arguments quoted (e.g. the {\hyperref[\detokenize{print::doc}]{\sphinxcrossref{\DUrole{doc}{print}}}}
command), then you can use a combination of single and double quotes,
as in the example above or below.

\sphinxAtStartPar
The \sphinxstyleemphasis{every} keyword is a means to avoid listing a long series of runs
and interleaving commands in your input script.  For example, a
{\hyperref[\detokenize{print::doc}]{\sphinxcrossref{\DUrole{doc}{print}}}} command could be invoked or a {\hyperref[\detokenize{fix::doc}]{\sphinxcrossref{\DUrole{doc}{fix}}}} could
be redefined, e.g. to reset a thermostat temperature.  Or this could
be useful for invoking a command you have added to LAMMPS that wraps
some other code (e.g. as a library) to perform a computation
periodically during a long LAMMPS run.  See the {\hyperref[\detokenize{Modify::doc}]{\sphinxcrossref{\DUrole{doc}{Modify}}}}
doc page for info about how to add new commands to LAMMPS.  See the
{\hyperref[\detokenize{Howto_couple::doc}]{\sphinxcrossref{\DUrole{doc}{Howto couple}}}} page for ideas about how to
couple LAMMPS to other codes.

\sphinxAtStartPar
With the \sphinxstyleemphasis{every} option, N total steps are simulated, in shorter runs
of M steps each.  After each M\sphinxhyphen{}length run, the specified commands are
invoked.  If only a single command is specified as NULL, then no
command is invoked.  Thus these lines:

\begin{sphinxVerbatim}[commandchars=\\\{\}]
\PYG{k}{variable }\PYG{n+nv+nvIdentifier}{q}\PYG{+w}{ }\PYG{n}{equal}\PYG{+w}{ }\PYG{n}{x[100]}
\PYG{k}{run}\PYG{+w}{ }\PYG{l+m}{6000}\PYG{+w}{ }\PYG{n}{every}\PYG{+w}{ }\PYG{l+m}{2000}\PYG{+w}{ }\PYG{l+s}{\PYGZdq{}}\PYG{l+s}{print \PYGZsq{}Coord = \PYGZdl{}q\PYGZsq{}}\PYG{l+s}{\PYGZdq{}}
\end{sphinxVerbatim}

\sphinxAtStartPar
are the equivalent of:

\begin{sphinxVerbatim}[commandchars=\\\{\}]
\PYG{k}{variable }\PYG{n+nv+nvIdentifier}{q}\PYG{+w}{ }\PYG{n}{equal}\PYG{+w}{ }\PYG{n}{x[100]}
\PYG{k}{run}\PYG{+w}{ }\PYG{l+m}{2000}
\PYG{k}{print}\PYG{+w}{ }\PYG{l+s}{\PYGZdq{}}\PYG{l+s}{Coord = \PYGZdl{}q}\PYG{l+s}{\PYGZdq{}}
\PYG{k}{run}\PYG{+w}{ }\PYG{l+m}{2000}
\PYG{k}{print}\PYG{+w}{ }\PYG{l+s}{\PYGZdq{}}\PYG{l+s}{Coord = \PYGZdl{}q}\PYG{l+s}{\PYGZdq{}}
\PYG{k}{run}\PYG{+w}{ }\PYG{l+m}{2000}
\PYG{k}{print}\PYG{+w}{ }\PYG{l+s}{\PYGZdq{}}\PYG{l+s}{Coord = \PYGZdl{}q}\PYG{l+s}{\PYGZdq{}}
\end{sphinxVerbatim}

\sphinxAtStartPar
which does 3 runs of 2000 steps and prints the x\sphinxhyphen{}coordinate of a
particular atom between runs.  Note that the variable “\$q” will
be evaluated afresh each time the print command is executed.

\sphinxAtStartPar
Note that by using the line continuation character “\&”, the run every
command can be spread across many lines, though it is still a single
command:

\begin{sphinxVerbatim}[commandchars=\\\{\}]
\PYG{k}{run}\PYG{+w}{ }\PYG{l+m}{100000}\PYG{+w}{ }\PYG{n}{every}\PYG{+w}{ }\PYG{l+m}{1000}\PYG{+w}{ }\PYG{o}{\PYGZam{}}
\PYG{+w}{  }\PYG{l+s}{\PYGZdq{}}\PYG{l+s}{print \PYGZsq{}Minimum value = \PYGZdl{}a\PYGZsq{}}\PYG{l+s}{\PYGZdq{}}\PYG{+w}{ }\PYG{o}{\PYGZam{}}
\PYG{+w}{  }\PYG{l+s}{\PYGZdq{}}\PYG{l+s}{print \PYGZsq{}Maximum value = \PYGZdl{}b\PYGZsq{}}\PYG{l+s}{\PYGZdq{}}\PYG{+w}{ }\PYG{o}{\PYGZam{}}
\PYG{+w}{  }\PYG{l+s}{\PYGZdq{}}\PYG{l+s}{print \PYGZsq{}Temp = \PYGZdl{}c\PYGZsq{}}\PYG{l+s}{\PYGZdq{}}\PYG{+w}{ }\PYG{o}{\PYGZam{}}
\PYG{+w}{  }\PYG{l+s}{\PYGZdq{}}\PYG{l+s}{print \PYGZsq{}Press = \PYGZdl{}d\PYGZsq{}}\PYG{l+s}{\PYGZdq{}}
\end{sphinxVerbatim}

\sphinxAtStartPar
If the \sphinxstyleemphasis{pre} and \sphinxstyleemphasis{post} options are set to “no” when used with the
\sphinxstyleemphasis{every} keyword, then the first run will do the full setup and the last
run will print the full timing summary, but these operations will be
skipped for intermediate runs.

\begin{sphinxadmonition}{note}{Note:}
\sphinxAtStartPar
You might wish to specify a command that exits the run by
jumping out of the loop, e.g.
\end{sphinxadmonition}

\begin{sphinxVerbatim}[commandchars=\\\{\}]
\PYG{k}{variable }\PYG{n+nv+nvIdentifier}{t}\PYG{+w}{ }\PYG{n}{equal}\PYG{+w}{ }\PYG{n}{temp}
\PYG{k}{run}\PYG{+w}{ }\PYG{l+m}{10000}\PYG{+w}{ }\PYG{n}{every}\PYG{+w}{ }\PYG{l+m}{100}\PYG{+w}{ }\PYG{l+s}{\PYGZdq{}}\PYG{l+s}{if \PYGZsq{}\PYGZdl{}t \PYGZlt{} 300.0\PYGZsq{} then \PYGZsq{}jump SELF afterrun\PYGZsq{}}\PYG{l+s}{\PYGZdq{}}
\end{sphinxVerbatim}

\sphinxAtStartPar
However, this will not work.  The run command simply executes each
command one at a time each time it pauses, then continues the run.

\sphinxAtStartPar
Instead, you should use the {\hyperref[\detokenize{fix_halt::doc}]{\sphinxcrossref{\DUrole{doc}{fix halt}}}} command, which
has additional options for how to exit the run.


\subsection{Restrictions}
\label{\detokenize{run:restrictions}}
\sphinxAtStartPar
When not using the \sphinxstyleemphasis{upto} keyword, the number of specified timesteps N
must fit in a signed 32\sphinxhyphen{}bit integer, so you are limited to slightly
more than 2 billion steps (2\textasciicircum{}31) in a single run.  When using \sphinxstyleemphasis{upto},
N can be larger than a signed 32\sphinxhyphen{}bit integer, however the difference
between N and the current timestep must still be no larger than
2\textasciicircum{}31 steps.

\sphinxAtStartPar
However, with or without the \sphinxstyleemphasis{upto} keyword, you can perform
successive runs to run a simulation for any number of steps (ok, up to
2\textasciicircum{}63 total steps).  I.e. the timestep counter within LAMMPS is a
64\sphinxhyphen{}bit signed integer.


\subsection{Related commands}
\label{\detokenize{run:related-commands}}
\sphinxAtStartPar
{\hyperref[\detokenize{minimize::doc}]{\sphinxcrossref{\DUrole{doc}{minimize}}}}, {\hyperref[\detokenize{run_style::doc}]{\sphinxcrossref{\DUrole{doc}{run\_style}}}},
{\hyperref[\detokenize{temper::doc}]{\sphinxcrossref{\DUrole{doc}{temper}}}}, {\hyperref[\detokenize{fix_halt::doc}]{\sphinxcrossref{\DUrole{doc}{fix halt}}}}


\subsection{Default}
\label{\detokenize{run:default}}
\sphinxAtStartPar
The option defaults are start = the current timestep, stop = current
timestep + N, pre = yes, and post = yes.

\sphinxstepscope

\index{run\_style@\spxentry{run\_style}}\ignorespaces 

\section{run\_style command}
\label{\detokenize{run_style:run-style-command}}\label{\detokenize{run_style:index-0}}\label{\detokenize{run_style::doc}}

\subsection{Syntax}
\label{\detokenize{run_style:syntax}}
\begin{sphinxVerbatim}[commandchars=\\\{\}]
\PYG{k}{run\PYGZus{}style}\PYG{+w}{ }\PYG{n}{style}\PYG{+w}{ }\PYG{n}{args}
\end{sphinxVerbatim}
\begin{itemize}
\item {} 
\sphinxAtStartPar
style = \sphinxstyleemphasis{verlet} or \sphinxstyleemphasis{verlet/split} or \sphinxstyleemphasis{respa} or \sphinxstyleemphasis{respa/omp}
\begin{sphinxalltt}
\sphinxstyleemphasis{verlet} args = none
\sphinxstyleemphasis{verlet/split} args = none
\sphinxstyleemphasis{respa} args = N n1 n2 ... keyword values ...
  N = \# of levels of rRESPA
  n1, n2, ... = loop factors between rRESPA levels (N\sphinxhyphen{}1 values)
  zero or more keyword/value pairings may be appended to the loop factors
  keyword = \sphinxstyleemphasis{bond} or \sphinxstyleemphasis{angle} or \sphinxstyleemphasis{dihedral} or \sphinxstyleemphasis{improper} or
            \sphinxstyleemphasis{pair} or \sphinxstyleemphasis{inner} or \sphinxstyleemphasis{middle} or \sphinxstyleemphasis{outer} or \sphinxstyleemphasis{hybrid} or \sphinxstyleemphasis{kspace}
    \sphinxstyleemphasis{bond} value = M
      M = which level (1\sphinxhyphen{}N) to compute bond forces in
    \sphinxstyleemphasis{angle} value = M
      M = which level (1\sphinxhyphen{}N) to compute angle forces in
    \sphinxstyleemphasis{dihedral} value = M
      M = which level (1\sphinxhyphen{}N) to compute dihedral forces in
    \sphinxstyleemphasis{improper} value = M
      M = which level (1\sphinxhyphen{}N) to compute improper forces in
    \sphinxstyleemphasis{pair} value = M
      M = which level (1\sphinxhyphen{}N) to compute pair forces in
    \sphinxstyleemphasis{inner} values = M cut1 cut2
      M = which level (1\sphinxhyphen{}N) to compute pair inner forces in
      cut1 = inner cutoff between pair inner and
             pair middle or outer  (distance units)
      cut2 = outer cutoff between pair inner and
             pair middle or outer  (distance units)
    \sphinxstyleemphasis{middle} values = M cut1 cut2
      M = which level (1\sphinxhyphen{}N) to compute pair middle forces in
      cut1 = inner cutoff between pair middle and pair outer (distance units)
      cut2 = outer cutoff between pair middle and pair outer (distance units)
    \sphinxstyleemphasis{outer} value = M
      M = which level (1\sphinxhyphen{}N) to compute pair outer forces in
    \sphinxstyleemphasis{hybrid} values = M1 M2 ... (as many values as there are hybrid sub\sphinxhyphen{}styles
      M1 = which level (1\sphinxhyphen{}N) to compute the first pair\_style hybrid sub\sphinxhyphen{}style in
      M2 = which level (1\sphinxhyphen{}N) to compute the second pair\_style hybrid sub\sphinxhyphen{}style in
      M3,etc
    \sphinxstyleemphasis{kspace} value = M
      M = which level (1\sphinxhyphen{}N) to compute kspace forces in
\end{sphinxalltt}

\end{itemize}


\subsection{Examples}
\label{\detokenize{run_style:examples}}
\begin{sphinxVerbatim}[commandchars=\\\{\}]
\PYG{k}{run\PYGZus{}style}\PYG{+w}{ }\PYG{n}{verlet}
\PYG{k}{run\PYGZus{}style}\PYG{+w}{ }\PYG{n}{respa}\PYG{+w}{ }\PYG{l+m}{4}\PYG{+w}{ }\PYG{l+m}{2}\PYG{+w}{ }\PYG{l+m}{2}\PYG{+w}{ }\PYG{l+m}{2}\PYG{+w}{ }\PYG{n}{bond}\PYG{+w}{ }\PYG{l+m}{1}\PYG{+w}{ }\PYG{n}{dihedral}\PYG{+w}{ }\PYG{l+m}{2}\PYG{+w}{ }\PYG{n}{pair}\PYG{+w}{ }\PYG{l+m}{3}\PYG{+w}{ }\PYG{n}{kspace}\PYG{+w}{ }\PYG{l+m}{4}
\PYG{k}{run\PYGZus{}style}\PYG{+w}{ }\PYG{n}{respa}\PYG{+w}{ }\PYG{l+m}{4}\PYG{+w}{ }\PYG{l+m}{2}\PYG{+w}{ }\PYG{l+m}{2}\PYG{+w}{ }\PYG{l+m}{2}\PYG{+w}{ }\PYG{n}{bond}\PYG{+w}{ }\PYG{l+m}{1}\PYG{+w}{ }\PYG{n}{dihedral}\PYG{+w}{ }\PYG{l+m}{2}\PYG{+w}{ }\PYG{n}{inner}\PYG{+w}{ }\PYG{l+m}{3}\PYG{+w}{ }\PYG{l+m}{5.0}\PYG{+w}{ }\PYG{l+m}{6.0}\PYG{+w}{ }\PYG{n}{outer}\PYG{+w}{ }\PYG{l+m}{4}\PYG{+w}{ }\PYG{n}{kspace}\PYG{+w}{ }\PYG{l+m}{4}
\PYG{k}{run\PYGZus{}style}\PYG{+w}{ }\PYG{n}{respa}\PYG{+w}{ }\PYG{l+m}{3}\PYG{+w}{ }\PYG{l+m}{4}\PYG{+w}{ }\PYG{l+m}{2}\PYG{+w}{ }\PYG{n}{bond}\PYG{+w}{ }\PYG{l+m}{1}\PYG{+w}{ }\PYG{n}{hybrid}\PYG{+w}{ }\PYG{l+m}{2}\PYG{+w}{ }\PYG{l+m}{2}\PYG{+w}{ }\PYG{l+m}{1}\PYG{+w}{ }\PYG{n}{kspace}\PYG{+w}{ }\PYG{l+m}{3}
\end{sphinxVerbatim}


\subsection{Description}
\label{\detokenize{run_style:description}}
\sphinxAtStartPar
Choose the style of time integrator used for molecular dynamics
simulations performed by LAMMPS.

\sphinxAtStartPar
The \sphinxstyleemphasis{verlet} style is the velocity form of the
Stoermer\sphinxhyphen{}Verlet time integration algorithm (velocity\sphinxhyphen{}Verlet)


\bigskip\hrule\bigskip


\sphinxAtStartPar
The \sphinxstyleemphasis{verlet/split} style is also a velocity\sphinxhyphen{}Verlet integrator, but it
splits the force calculation within each timestep over 2 partitions of
processors.  See the {\hyperref[\detokenize{Run_options::doc}]{\sphinxcrossref{\DUrole{doc}{\sphinxhyphen{}partition command\sphinxhyphen{}line switch}}}}
for info on how to run LAMMPS with multiple partitions.

\sphinxAtStartPar
Specifically, this style performs all computation except the
{\hyperref[\detokenize{kspace_style::doc}]{\sphinxcrossref{\DUrole{doc}{kspace\_style}}}} portion of the force field on the
first partition.  This include the {\hyperref[\detokenize{pair_style::doc}]{\sphinxcrossref{\DUrole{doc}{pair style}}}},
{\hyperref[\detokenize{bond_style::doc}]{\sphinxcrossref{\DUrole{doc}{bond style}}}}, {\hyperref[\detokenize{neighbor::doc}]{\sphinxcrossref{\DUrole{doc}{neighbor list building}}}}, {\hyperref[\detokenize{fix::doc}]{\sphinxcrossref{\DUrole{doc}{fixes}}}} including time integration, and output.
The {\hyperref[\detokenize{kspace_style::doc}]{\sphinxcrossref{\DUrole{doc}{kspace\_style}}}} portion of the calculation is
performed on the second partition.

\sphinxAtStartPar
This can lead to a significant speedup, if the number of processors can
be easily increased and the fraction of time is spent in computing
Kspace interactions is significant, too.  The two partitions may have a
different number of processors.  This is most useful for the PPPM
kspace\_style when its performance on a large number of processors
degrades due to the cost of communication in its 3d FFTs.  In this
scenario, splitting your P total processors into 2 subsets of
processors, P1 in the first partition and P2 in the second partition,
can enable your simulation to run faster.  This is because the
long\sphinxhyphen{}range forces in PPPM can be calculated at the same time as pairwise
and bonded forces are being calculated \sphinxstyleemphasis{and} the parallel 3d FFTs can be
faster to compute when running on fewer processors.  Please note that
the scenario of using fewer MPI processes to reduce communication
overhead can also be implemented through using MPI with OpenMP threads
via the INTEL, KOKKOS, or OPENMP package.  This alternative option is
typically more effective in case of a fixed number of available
processors and less complex to execute.

\sphinxAtStartPar
To use the \sphinxstyleemphasis{verlet/split} style, you must define 2 partitions with the
{\hyperref[\detokenize{Run_options::doc}]{\sphinxcrossref{\DUrole{doc}{\sphinxhyphen{}partition command\sphinxhyphen{}line switch}}}}, where partition P1
is either the same size or an integer multiple of the size of the
partition P2.  Typically having P1 be 3x larger than P2 is a good
choice, since the (serial) performance of LAMMPS is often best if the
time spent in the \sphinxcode{\sphinxupquote{Pair}} computation versus \sphinxcode{\sphinxupquote{Kspace}} is a 3:1 split.
The 3d processor layouts in each partition must overlay in the following
sense.  If P1 is a Px1 by Py1 by Pz1 grid, and P2 = Px2 by Py2 by Pz2,
then Px1 must be an integer multiple of Px2, and similarly for Py1 a
multiple of Py2, and Pz1 a multiple of Pz2.

\sphinxAtStartPar
Typically the best way to do this is to let the first partition choose
its own optimal layout, then require the second partition’s layout to
match the integer multiple constraint.  See the
{\hyperref[\detokenize{processors::doc}]{\sphinxcrossref{\DUrole{doc}{processors}}}} command with its \sphinxstyleemphasis{part} keyword for a way
to control this, e.g.

\begin{sphinxVerbatim}[commandchars=\\\{\}]
\PYG{k}{processors}\PYG{+w}{ }\PYG{o}{*}\PYG{+w}{ }\PYG{o}{*}\PYG{+w}{ }\PYG{o}{*}\PYG{+w}{ }\PYG{n}{part}\PYG{+w}{ }\PYG{l+m}{1}\PYG{+w}{ }\PYG{l+m}{2}\PYG{+w}{ }\PYG{n}{multiple}
\end{sphinxVerbatim}

\sphinxAtStartPar
You can also use the {\hyperref[\detokenize{partition::doc}]{\sphinxcrossref{\DUrole{doc}{partition}}}} command to explicitly
specify the processor layout on each partition.  E.g. for 2 partitions
of 60 and 15 processors each:

\begin{sphinxVerbatim}[commandchars=\\\{\}]
\PYG{k}{partition}\PYG{+w}{ }\PYG{n}{yes}\PYG{+w}{ }\PYG{l+m}{1}\PYG{+w}{ }\PYG{n}{processors}\PYG{+w}{ }\PYG{l+m}{3}\PYG{+w}{ }\PYG{l+m}{4}\PYG{+w}{ }\PYG{l+m}{5}
\PYG{k}{partition}\PYG{+w}{ }\PYG{n}{yes}\PYG{+w}{ }\PYG{l+m}{2}\PYG{+w}{ }\PYG{n}{processors}\PYG{+w}{ }\PYG{l+m}{3}\PYG{+w}{ }\PYG{l+m}{1}\PYG{+w}{ }\PYG{l+m}{5}
\end{sphinxVerbatim}

\sphinxAtStartPar
When you run in 2\sphinxhyphen{}partition mode with the \sphinxstyleemphasis{verlet/split} style, the
thermodynamic data for the entire simulation will be output to the log
and screen file of the first partition, which are log.lammps.0 and
screen.0 by default; see the {\hyperref[\detokenize{Run_options::doc}]{\sphinxcrossref{\DUrole{doc}{\sphinxhyphen{}plog and \sphinxhyphen{}pscreen command\sphinxhyphen{}line
switches}}}} to change this.  The log and screen file for the
second partition will not contain thermodynamic output beyond the first
timestep of the run.

\sphinxAtStartPar
See the {\hyperref[\detokenize{Speed_packages::doc}]{\sphinxcrossref{\DUrole{doc}{Accelerator packages}}}} page for
performance details of the speed\sphinxhyphen{}up offered by the \sphinxstyleemphasis{verlet/split}
style.  One important performance consideration is the assignment of
logical processors in the 2 partitions to the physical cores of a
parallel machine.  The {\hyperref[\detokenize{processors::doc}]{\sphinxcrossref{\DUrole{doc}{processors}}}} command has
options to support this, and strategies are discussed in {\hyperref[\detokenize{Speed::doc}]{\sphinxcrossref{\DUrole{doc}{Section
5}}}} of the manual.


\bigskip\hrule\bigskip


\sphinxAtStartPar
The \sphinxstyleemphasis{respa} style implements the rRESPA multi\sphinxhyphen{}timescale integrator
{\hyperref[\detokenize{run_style:tuckerman3}]{\sphinxcrossref{\DUrole{std,std-ref}{(Tuckerman)}}}} with N hierarchical levels, where level
1 is the innermost loop (shortest timestep) and level N is the outermost
loop (largest timestep).  The loop factor arguments specify what the
looping factor is between levels.  N1 specifies the number of iterations
of level 1 for a single iteration of level 2, N2 is the iterations of
level 2 per iteration of level 3, etc.  N\sphinxhyphen{}1 looping parameters must be
specified.

\sphinxAtStartPar
Thus with a 4\sphinxhyphen{}level respa setting of “2 2 2” for the 3 loop factors, you
could choose to have bond interactions computed 8x per large timestep,
angle interactions computed 4x, pair interactions computed 2x, and
long\sphinxhyphen{}range interactions once per large timestep.

\sphinxAtStartPar
The {\hyperref[\detokenize{timestep::doc}]{\sphinxcrossref{\DUrole{doc}{timestep}}}} command sets the large timestep for the
outermost rRESPA level.  Thus if the 3 loop factors are “2 2 2” for
4\sphinxhyphen{}level rRESPA, and the outer timestep is set to 4.0 fs, then the inner
timestep would be 8x smaller or 0.5 fs.  All other LAMMPS commands that
specify number of timesteps (e.g. {\hyperref[\detokenize{thermo::doc}]{\sphinxcrossref{\DUrole{doc}{thermo}}}} for thermo
output every N steps, {\hyperref[\detokenize{neigh_modify::doc}]{\sphinxcrossref{\DUrole{doc}{neigh\_modify delay/every}}}}
parameters, {\hyperref[\detokenize{dump::doc}]{\sphinxcrossref{\DUrole{doc}{dump}}}} every N steps, etc) refer to the
outermost timesteps.

\sphinxAtStartPar
The rRESPA keywords enable you to specify at what level of the hierarchy
various forces will be computed.  If not specified, the defaults are
that bond forces are computed at level 1 (innermost loop), angle forces
are computed where bond forces are, dihedral forces are computed where
angle forces are, improper forces are computed where dihedral forces
are, pair forces are computed at the outermost level, and kspace forces
are computed where pair forces are.  The inner, middle, outer forces
have no defaults.

\sphinxAtStartPar
For fixes that support it, the rRESPA level at which a given fix is
active, can be selected through the {\hyperref[\detokenize{fix_modify::doc}]{\sphinxcrossref{\DUrole{doc}{fix\_modify}}}}
command.

\sphinxAtStartPar
The \sphinxstyleemphasis{inner} and \sphinxstyleemphasis{middle} keywords take additional arguments for cutoffs
that are used by the pairwise force computations.  If the 2 cutoffs for
\sphinxstyleemphasis{inner} are 5.0 and 6.0, this means that all pairs up to 6.0 apart are
computed by the inner force.  Those between 5.0 and 6.0 have their force
go ramped to 0.0 so the overlap with the next regime (middle or outer)
is smooth.  The next regime (middle or outer) will compute forces for
all pairs from 5.0 outward, with those from 5.0 to 6.0 having their
value ramped in an inverse manner.

\sphinxAtStartPar
Note that you can use \sphinxstyleemphasis{inner} and \sphinxstyleemphasis{outer} without using \sphinxstyleemphasis{middle} to
split the pairwise computations into two portions instead of three.
Unless you are using a very long pairwise cutoff, a 2\sphinxhyphen{}way split is often
faster than a 3\sphinxhyphen{}way split, since it avoids too much duplicate
computation of pairwise interactions near the intermediate cutoffs.

\sphinxAtStartPar
Also note that only a few pair potentials support the use of the \sphinxstyleemphasis{inner}
and \sphinxstyleemphasis{middle} and \sphinxstyleemphasis{outer} keywords.  If not, only the \sphinxstyleemphasis{pair} keyword can
be used with that pair style, meaning all pairwise forces are computed
at the same rRESPA level.  See the doc pages for individual pair styles
for details.

\sphinxAtStartPar
Another option for using pair potentials with rRESPA is with the
\sphinxstyleemphasis{hybrid} keyword, which requires the use of the {\hyperref[\detokenize{pair_hybrid::doc}]{\sphinxcrossref{\DUrole{doc}{pair\_style hybrid
or hybrid/overlay}}}} command.  In this scenario, different
sub\sphinxhyphen{}styles of the hybrid pair style are evaluated at different rRESPA
levels.  This can be useful, for example, to set different timesteps for
hybrid coarse\sphinxhyphen{}grained/all\sphinxhyphen{}atom models.  The \sphinxstyleemphasis{hybrid} keyword requires as
many level assignments as there are hybrid sub\sphinxhyphen{}styles, which assigns
each sub\sphinxhyphen{}style to a rRESPA level, following their order of definition in
the pair\_style command. Since the \sphinxstyleemphasis{hybrid} keyword operates on pair
style computations, it is mutually exclusive with either the \sphinxstyleemphasis{pair} or
the \sphinxstyleemphasis{inner}/\sphinxstyleemphasis{middle}/\sphinxstyleemphasis{outer} keywords.

\sphinxAtStartPar
When using rRESPA (or for any MD simulation) care must be taken to
choose a timestep size(s) that ensures the Hamiltonian for the chosen
ensemble is conserved.  For the constant NVE ensemble, total energy
must be conserved.  Unfortunately, it is difficult to know \sphinxstyleemphasis{a priori}
how well energy will be conserved, and a fairly long test simulation
(\textasciitilde{}10 ps) is usually necessary in order to verify that no long\sphinxhyphen{}term
drift in energy occurs with the trial set of parameters.

\sphinxAtStartPar
With that caveat, a few rules\sphinxhyphen{}of\sphinxhyphen{}thumb may be useful in selecting
\sphinxstyleemphasis{respa} settings.  The following applies mostly to biomolecular
simulations using the CHARMM or a similar all\sphinxhyphen{}atom force field, but
the concepts are adaptable to other problems.  Without SHAKE, bonds
involving hydrogen atoms exhibit high\sphinxhyphen{}frequency vibrations and require
a timestep on the order of 0.5 fs in order to conserve energy.  The
relatively inexpensive force computations for the bonds, angles,
impropers, and dihedrals can be computed on this innermost 0.5 fs
step.  The outermost timestep cannot be greater than 4.0 fs without
risking energy drift.  Smooth switching of forces between the levels
of the rRESPA hierarchy is also necessary to avoid drift, and a 1\sphinxhyphen{}2
Angstrom “healing distance” (the distance between the outer and inner
cutoffs) works reasonably well.  We thus recommend the following
settings for use of the \sphinxstyleemphasis{respa} style without SHAKE in biomolecular
simulations:

\begin{sphinxVerbatim}[commandchars=\\\{\}]
\PYG{k}{timestep}\PYG{+w}{  }\PYG{l+m}{4.0}
\PYG{k}{run\PYGZus{}style}\PYG{+w}{ }\PYG{n}{respa}\PYG{+w}{ }\PYG{l+m}{4}\PYG{+w}{ }\PYG{l+m}{2}\PYG{+w}{ }\PYG{l+m}{2}\PYG{+w}{ }\PYG{l+m}{2}\PYG{+w}{ }\PYG{n}{inner}\PYG{+w}{ }\PYG{l+m}{2}\PYG{+w}{ }\PYG{l+m}{4.5}\PYG{+w}{ }\PYG{l+m}{6.0}\PYG{+w}{ }\PYG{n}{middle}\PYG{+w}{ }\PYG{l+m}{3}\PYG{+w}{ }\PYG{l+m}{8.0}\PYG{+w}{ }\PYG{l+m}{10.0}\PYG{+w}{ }\PYG{n}{outer}\PYG{+w}{ }\PYG{l+m}{4}
\end{sphinxVerbatim}

\sphinxAtStartPar
With these settings, users can expect good energy conservation and
roughly a 2.5 fold speedup over the \sphinxstyleemphasis{verlet} style with a 0.5 fs
timestep.

\sphinxAtStartPar
If SHAKE is used with the \sphinxstyleemphasis{respa} style, time reversibility is lost,
but substantially longer time steps can be achieved.  For biomolecular
simulations using the CHARMM or similar all\sphinxhyphen{}atom force field, bonds
involving hydrogen atoms exhibit high frequency vibrations and require
a time step on the order of 0.5 fs in order to conserve energy.
These high frequency modes also limit the outer time step sizes since
the modes are coupled.  It is therefore desirable to use SHAKE with
respa in order to freeze out these high frequency motions and increase
the size of the time steps in the respa hierarchy.  The following
settings can be used for biomolecular simulations with SHAKE and
rRESPA:

\begin{sphinxVerbatim}[commandchars=\\\{\}]
\PYG{k}{fix             }\PYG{n+nv+nvIdentifier}{2}\PYG{+w}{ }\PYG{n+nv+nvIdentifier}{all}\PYG{+w}{ }\PYG{n}{shake}\PYG{+w}{ }\PYG{l+m}{0.000001}\PYG{+w}{ }\PYG{l+m}{500}\PYG{+w}{ }\PYG{l+m}{0}\PYG{+w}{ }\PYG{n}{m}\PYG{+w}{ }\PYG{l+m}{1.0}\PYG{+w}{ }\PYG{n}{a}\PYG{+w}{ }\PYG{l+m}{1}
\PYG{k}{timestep}\PYG{+w}{        }\PYG{l+m}{4.0}
\PYG{k}{run\PYGZus{}style}\PYG{+w}{       }\PYG{n}{respa}\PYG{+w}{ }\PYG{l+m}{2}\PYG{+w}{ }\PYG{l+m}{2}\PYG{+w}{ }\PYG{n}{inner}\PYG{+w}{ }\PYG{l+m}{1}\PYG{+w}{ }\PYG{l+m}{4.0}\PYG{+w}{ }\PYG{l+m}{5.0}\PYG{+w}{ }\PYG{n}{outer}\PYG{+w}{ }\PYG{l+m}{2}
\end{sphinxVerbatim}

\sphinxAtStartPar
With these settings, users can expect good energy conservation and
roughly a 1.5 fold speedup over the \sphinxstyleemphasis{verlet} style with SHAKE and a
2.0 fs timestep.

\sphinxAtStartPar
For non\sphinxhyphen{}biomolecular simulations, the \sphinxstyleemphasis{respa} style can be
advantageous if there is a clear separation of time scales \sphinxhyphen{} fast and
slow modes in the simulation.  For example, a system of slowly\sphinxhyphen{}moving
charged polymer chains could be setup as follows:

\begin{sphinxVerbatim}[commandchars=\\\{\}]
\PYG{k}{timestep}\PYG{+w}{ }\PYG{l+m}{4.0}
\PYG{k}{run\PYGZus{}style}\PYG{+w}{ }\PYG{n}{respa}\PYG{+w}{ }\PYG{l+m}{2}\PYG{+w}{ }\PYG{l+m}{8}
\end{sphinxVerbatim}

\sphinxAtStartPar
This is two\sphinxhyphen{}level rRESPA with an 8x difference between the short and
long timesteps.  The bonds, angles, dihedrals will be computed every
0.5 fs (assuming real units), while the pair and kspace interactions
will be computed once every 4 fs.  These are the default settings for
each kind of interaction, so no additional keywords are necessary.

\sphinxAtStartPar
Even a LJ system can benefit from rRESPA if the interactions are
divided by the inner, middle and outer keywords.  A 2\sphinxhyphen{}fold or more
speedup can be obtained while maintaining good energy conservation.
In real units, for a pure LJ fluid at liquid density, with a sigma of
3.0 Angstroms, and epsilon of 0.1 kcal/mol, the following settings
seem to work well:

\begin{sphinxVerbatim}[commandchars=\\\{\}]
\PYG{k}{timestep}\PYG{+w}{  }\PYG{l+m}{36.0}
\PYG{k}{run\PYGZus{}style}\PYG{+w}{ }\PYG{n}{respa}\PYG{+w}{ }\PYG{l+m}{3}\PYG{+w}{ }\PYG{l+m}{3}\PYG{+w}{ }\PYG{l+m}{4}\PYG{+w}{ }\PYG{n}{inner}\PYG{+w}{ }\PYG{l+m}{1}\PYG{+w}{ }\PYG{l+m}{3.0}\PYG{+w}{ }\PYG{l+m}{4.0}\PYG{+w}{ }\PYG{n}{middle}\PYG{+w}{ }\PYG{l+m}{2}\PYG{+w}{ }\PYG{l+m}{6.0}\PYG{+w}{ }\PYG{l+m}{7.0}\PYG{+w}{ }\PYG{n}{outer}\PYG{+w}{ }\PYG{l+m}{3}
\end{sphinxVerbatim}


\bigskip\hrule\bigskip


\sphinxAtStartPar
The \sphinxstyleemphasis{respa/omp} style is a variant of \sphinxstyleemphasis{respa} adapted for use with
pair, bond, angle, dihedral, improper, or kspace styles with an \sphinxstyleemphasis{omp}
suffix. It is functionally equivalent to \sphinxstyleemphasis{respa} but performs
additional operations required for managing \sphinxstyleemphasis{omp} styles.  For more on
\sphinxstyleemphasis{omp} styles see the {\hyperref[\detokenize{Speed_omp::doc}]{\sphinxcrossref{\DUrole{doc}{Speed omp}}}} doc page.  Accelerated
styles take the same arguments and should produce the same results,
except for round\sphinxhyphen{}off and precision issues.

\sphinxAtStartPar
You can specify \sphinxstyleemphasis{respa/omp} explicitly in your input script, or you
can use the {\hyperref[\detokenize{Run_options::doc}]{\sphinxcrossref{\DUrole{doc}{\sphinxhyphen{}suffix command\sphinxhyphen{}line switch}}}} when you
invoke LAMMPS, or you can use the {\hyperref[\detokenize{suffix::doc}]{\sphinxcrossref{\DUrole{doc}{suffix}}}} command in
your input script.

\sphinxAtStartPar
See the {\hyperref[\detokenize{Speed_packages::doc}]{\sphinxcrossref{\DUrole{doc}{Accelerator packages}}}} page for more
instructions on how to use the accelerated styles effectively.


\bigskip\hrule\bigskip



\subsection{Restrictions}
\label{\detokenize{run_style:restrictions}}
\sphinxAtStartPar
The \sphinxstyleemphasis{verlet/split} style can only be used if LAMMPS was built with the
REPLICA package.  Correspondingly the \sphinxstyleemphasis{respa/omp} style is available
only if the OPENMP package was included. See the {\hyperref[\detokenize{Build_package::doc}]{\sphinxcrossref{\DUrole{doc}{Build package}}}} page for more info.

\sphinxAtStartPar
Run style \sphinxstyleemphasis{verlet/split} is not compatible with kspace styles from
the INTEL package and it is not compatible with any tip4p, dipole,
or spin kspace styles.

\sphinxAtStartPar
Whenever using rRESPA, the user should experiment with trade\sphinxhyphen{}offs in
speed and accuracy for their system, and verify that they are
conserving energy to adequate precision.


\subsection{Related commands}
\label{\detokenize{run_style:related-commands}}
\sphinxAtStartPar
{\hyperref[\detokenize{timestep::doc}]{\sphinxcrossref{\DUrole{doc}{timestep}}}}, {\hyperref[\detokenize{run::doc}]{\sphinxcrossref{\DUrole{doc}{run}}}}


\subsection{Default}
\label{\detokenize{run_style:default}}
\begin{sphinxVerbatim}[commandchars=\\\{\}]
\PYG{k}{run\PYGZus{}style}\PYG{+w}{ }\PYG{n}{verlet}
\end{sphinxVerbatim}

\sphinxAtStartPar
For run\_style respa, the default assignment of interactions
to rRESPA levels is as follows:
\begin{itemize}
\item {} 
\sphinxAtStartPar
bond forces = level 1 (innermost loop)

\item {} 
\sphinxAtStartPar
angle forces = same level as bond forces

\item {} 
\sphinxAtStartPar
dihedral forces = same level as angle forces

\item {} 
\sphinxAtStartPar
improper forces = same level as dihedral forces

\item {} 
\sphinxAtStartPar
pair forces = level N (outermost level)

\item {} 
\sphinxAtStartPar
kspace forces = same level as pair forces

\item {} 
\sphinxAtStartPar
inner, middle, outer forces = no default

\end{itemize}


\bigskip\hrule\bigskip

\phantomsection\label{\detokenize{run_style:tuckerman3}}
\sphinxAtStartPar
\sphinxstylestrong{(Tuckerman)} Tuckerman, Berne and Martyna, J Chem Phys, 97, p 1990
(1992).

\sphinxstepscope

\index{set@\spxentry{set}}\ignorespaces 

\section{set command}
\label{\detokenize{set:set-command}}\label{\detokenize{set:index-0}}\label{\detokenize{set::doc}}

\subsection{Syntax}
\label{\detokenize{set:syntax}}
\begin{sphinxVerbatim}[commandchars=\\\{\}]
\PYG{k}{set}\PYG{+w}{ }\PYG{n}{style}\PYG{+w}{ }\PYG{n}{ID}\PYG{+w}{ }\PYG{n}{keyword}\PYG{+w}{ }\PYG{n}{values}\PYG{+w}{ }\PYG{n}{...}
\end{sphinxVerbatim}
\begin{itemize}
\item {} 
\sphinxAtStartPar
style = \sphinxstyleemphasis{atom} or \sphinxstyleemphasis{type} or \sphinxstyleemphasis{mol} or \sphinxstyleemphasis{group} or \sphinxstyleemphasis{region}

\item {} 
\sphinxAtStartPar
ID = depends on style

\end{itemize}
\begin{sphinxalltt}
for style = \sphinxstyleemphasis{atom}, ID = a range of atom IDs
for style = \sphinxstyleemphasis{type}, ID = a range of numeric types or a single type label
for style = \sphinxstyleemphasis{mol}, ID = a range of molecule IDs
for style = \sphinxstyleemphasis{group}, ID = a group ID
for style = \sphinxstyleemphasis{region}, ID = a region ID
\end{sphinxalltt}
\begin{itemize}
\item {} 
\sphinxAtStartPar
one or more keyword/value pairs may be appended

\item {} 
\sphinxAtStartPar
keyword = \sphinxstyleemphasis{type} or \sphinxstyleemphasis{type/fraction} or \sphinxstyleemphasis{type/ratio} or \sphinxstyleemphasis{type/subset}
or \sphinxstyleemphasis{mol} or \sphinxstyleemphasis{x} or \sphinxstyleemphasis{y} or \sphinxstyleemphasis{z} or \sphinxstyleemphasis{vx} or \sphinxstyleemphasis{vy} or \sphinxstyleemphasis{vz} or \sphinxstyleemphasis{charge} or
\sphinxstyleemphasis{dipole} or \sphinxstyleemphasis{dipole/random} or \sphinxstyleemphasis{quat} or \sphinxstyleemphasis{spin/atom} or \sphinxstyleemphasis{spin/atom/random} or
\sphinxstyleemphasis{spin/electron} or \sphinxstyleemphasis{radius/electron} or
\sphinxstyleemphasis{quat} or \sphinxstyleemphasis{quat/random} or \sphinxstyleemphasis{diameter} or \sphinxstyleemphasis{shape} or \sphinxstyleemphasis{length} or \sphinxstyleemphasis{tri} or
\sphinxstyleemphasis{theta} or \sphinxstyleemphasis{theta/random} or \sphinxstyleemphasis{angmom} or \sphinxstyleemphasis{omega} or
\sphinxstyleemphasis{mass} or \sphinxstyleemphasis{density} or \sphinxstyleemphasis{density/disc} or \sphinxstyleemphasis{temperature} or
\sphinxstyleemphasis{volume} or \sphinxstyleemphasis{image} or \sphinxstyleemphasis{bond} or \sphinxstyleemphasis{angle} or \sphinxstyleemphasis{dihedral} or
\sphinxstyleemphasis{improper} or \sphinxstyleemphasis{sph/e} or \sphinxstyleemphasis{sph/cv} or \sphinxstyleemphasis{sph/rho} or
\sphinxstyleemphasis{smd/contact/radius} or \sphinxstyleemphasis{smd/mass/density} or \sphinxstyleemphasis{dpd/theta} or
\sphinxstyleemphasis{edpd/temp} or \sphinxstyleemphasis{edpd/cv} or \sphinxstyleemphasis{cc} or \sphinxstyleemphasis{epsilon} or
\sphinxstyleemphasis{i\_name} or \sphinxstyleemphasis{d\_name} or \sphinxstyleemphasis{i2\_name} or \sphinxstyleemphasis{d2\_name}
\begin{sphinxalltt}
\sphinxstyleemphasis{type} value = numeric atom type or type label
  value can be an atom\sphinxhyphen{}style variable (see below)
\sphinxstyleemphasis{type/fraction} values = type fraction seed
  type = numeric atom type or type label
  fraction = approximate fraction of selected atoms to set to new atom type
  seed = random \# seed (positive integer)
\sphinxstyleemphasis{type/ratio} values = type fraction seed
  type = numeric atom type or type label
  fraction = exact fraction of selected atoms to set to new atom type
  seed = random \# seed (positive integer)
\sphinxstyleemphasis{type/subset} values = type Nsubset seed
  type = numeric atom type or type label
  Nsubset = exact number of selected atoms to set to new atom type
  seed = random \# seed (positive integer)
\sphinxstyleemphasis{mol} value = molecule ID
value can be an atom\sphinxhyphen{}style variable (see below)
\sphinxstyleemphasis{x},\sphinxstyleemphasis{y},\sphinxstyleemphasis{z} value = atom coordinate (distance units)
  value can be an atom\sphinxhyphen{}style variable (see below)
\sphinxstyleemphasis{vx},\sphinxstyleemphasis{vy},\sphinxstyleemphasis{vz} value = atom velocity (velocity units)
  value can be an atom\sphinxhyphen{}style variable (see below)
\sphinxstyleemphasis{charge} value = atomic charge (charge units)
  value can be an atom\sphinxhyphen{}style variable (see below)
\sphinxstyleemphasis{dipole} values = x y z
  x,y,z = orientation of dipole moment vector
  any of x,y,z can be an atom\sphinxhyphen{}style variable (see below)
\sphinxstyleemphasis{dipole/random} value = seed Dlen
  seed = random \# seed (positive integer) for dipole moment orientations
  Dlen = magnitude of dipole moment (dipole units)
\sphinxstyleemphasis{spin/atom} values = g x y z
  g = magnitude of magnetic spin vector (in Bohr magneton\textquotesingle{}s unit)
  x,y,z = orientation of magnetic spin vector
  any of x,y,z can be an atom\sphinxhyphen{}style variable (see below)
\sphinxstyleemphasis{spin/atom/random} value = seed Dlen
  seed = random \# seed (positive integer) for magnetic spin orientations
  Dlen = magnitude of magnetic spin vector (in Bohr magneton\textquotesingle{}s unit)
\sphinxstyleemphasis{radius/electron} values = eradius
  eradius = electron radius (or fixed\sphinxhyphen{}core radius) (distance units)
\sphinxstyleemphasis{spin/electron} value = espin
  espin = electron spin (+1/\sphinxhyphen{}1), 0 = nuclei, 2 = fixed\sphinxhyphen{}core, 3 = pseudo\sphinxhyphen{}cores (i.e. ECP)
\sphinxstyleemphasis{quat} values = a b c theta
  a,b,c = unit vector to rotate particle around via right\sphinxhyphen{}hand rule
  theta = rotation angle (degrees)
  any of a,b,c,theta can be an atom\sphinxhyphen{}style variable (see below)
\sphinxstyleemphasis{quat/random} value = seed
  seed = random \# seed (positive integer) for quaternion orientations
\sphinxstyleemphasis{diameter} value = diameter of spherical particle (distance units)
  value can be an atom\sphinxhyphen{}style variable (see below)
\sphinxstyleemphasis{shape} value = Sx Sy Sz
  Sx,Sy,Sz = 3 diameters of ellipsoid (distance units)
\sphinxstyleemphasis{length} value = len
  len = length of line segment (distance units)
  len can be an atom\sphinxhyphen{}style variable (see below)
\sphinxstyleemphasis{tri} value = side
  side = side length of equilateral triangle (distance units)
  side can be an atom\sphinxhyphen{}style variable (see below)
\sphinxstyleemphasis{theta} value = angle (degrees)
  angle = orientation of line segment with respect to x\sphinxhyphen{}axis
  angle can be an atom\sphinxhyphen{}style variable (see below)
\sphinxstyleemphasis{theta/random} value = seed
  seed = random \# seed (positive integer) for line segment orienations
\sphinxstyleemphasis{angmom} values = Lx Ly Lz
  Lx,Ly,Lz = components of angular momentum vector (distance\sphinxhyphen{}mass\sphinxhyphen{}velocity units)
  any of Lx,Ly,Lz can be an atom\sphinxhyphen{}style variable (see below)
\sphinxstyleemphasis{omega} values = Wx Wy Wz
  Wx,Wy,Wz = components of angular velocity vector (radians/time units)
  any of wx,wy,wz can be an atom\sphinxhyphen{}style variable (see below)
\sphinxstyleemphasis{mass} value = per\sphinxhyphen{}atom mass (mass units)
  value can be an atom\sphinxhyphen{}style variable (see below)
\sphinxstyleemphasis{density} value = particle density for a sphere or ellipsoid (mass/distance\textasciicircum{}3 units), or for a triangle (mass/distance\textasciicircum{}2 units) or line (mass/distance units) particle
  value can be an atom\sphinxhyphen{}style variable (see below)
\sphinxstyleemphasis{density/disc} value = particle density for a 2d disc or ellipse (mass/distance\textasciicircum{}2 units)
  value can be an atom\sphinxhyphen{}style variable (see below)
\sphinxstyleemphasis{temperature} value = temperature for finite\sphinxhyphen{}size particles (temperature units)
  value can be an atom\sphinxhyphen{}style variable (see below)
\sphinxstyleemphasis{volume} value = particle volume for Peridynamic particle (distance\textasciicircum{}3 units)
  value can be an atom\sphinxhyphen{}style variable (see below)
\sphinxstyleemphasis{image} nx ny nz
  nx,ny,nz = which periodic image of the simulation box the atom is in
  any of nx,ny,nz can be an atom\sphinxhyphen{}style variable (see below)
\sphinxstyleemphasis{bond} value = numeric bond type or bond type label, for all bonds between selected atoms
\sphinxstyleemphasis{angle} value = numeric angle type or angle type label, for all angles between selected atoms
\sphinxstyleemphasis{dihedral} value = numeric dihedral type or dihedral type label, for all dihedrals between selected atoms
\sphinxstyleemphasis{improper} value = numeric improper type or improper type label, for all impropers between selected atoms
\sphinxstyleemphasis{rheo/rho} value = density of RHEO particles (mass/distance\textasciicircum{}3)
\sphinxstyleemphasis{rheo/status} value = status or phase of RHEO particles (unitless)
\sphinxstyleemphasis{sph/e} value = energy of SPH particles (need units)
  value can be an atom\sphinxhyphen{}style variable (see below)
\sphinxstyleemphasis{sph/cv} value = heat capacity of SPH particles (need units)
  value can be an atom\sphinxhyphen{}style variable (see below)
\sphinxstyleemphasis{sph/rho} value = density of SPH particles (need units)
  value can be an atom\sphinxhyphen{}style variable (see below)
\sphinxstyleemphasis{smd/contact/radius} = radius for short range interactions, i.e. contact and friction
  value can be an atom\sphinxhyphen{}style variable (see below)
\sphinxstyleemphasis{smd/mass/density} = set particle mass based on volume by providing a mass density
  value can be an atom\sphinxhyphen{}style variable (see below)
\sphinxstyleemphasis{dpd/theta} value = internal temperature of DPD particles (temperature units)
  value can be an atom\sphinxhyphen{}style variable (see below)
  value can be NULL which sets internal temp of each particle to KE temp
\sphinxstyleemphasis{edpd/temp} value = temperature of eDPD particles (temperature units)
  value can be an atom\sphinxhyphen{}style variable (see below)
\sphinxstyleemphasis{edpd/cv} value = volumetric heat capacity of eDPD particles (energy/temperature/volume units)
  value can be an atom\sphinxhyphen{}style variable (see below)
\sphinxstyleemphasis{cc} values = index cc
  index = index of a chemical species (1 to Nspecies)
  cc = chemical concentration of tDPD particles for a species (mole/volume units)
\sphinxstyleemphasis{epsilon} value = dielectric constant of the medium where the atoms reside
\sphinxstyleemphasis{i\_name} value = custom integer vector with name
\sphinxstyleemphasis{d\_name} value = custom floating\sphinxhyphen{}point vector with name
\sphinxstyleemphasis{i2\_name} value = column of a custom integer array with name
                  column specified as i2\_name{[}N{]} where N is 1 to Ncol
\sphinxstyleemphasis{d2\_name} value = column of a custom floating\sphinxhyphen{}point array with name
                  column specified as d2\_name{[}N{]} where N is 1 to Ncol
\end{sphinxalltt}

\end{itemize}


\subsection{Examples}
\label{\detokenize{set:examples}}
\begin{sphinxVerbatim}[commandchars=\\\{\}]
\PYG{k}{set}\PYG{+w}{ }\PYG{k}{group }\PYG{n+nv+nvIdentifier}{solvent}\PYG{+w}{ }\PYG{n}{type}\PYG{+w}{ }\PYG{l+m}{2}
\PYG{k}{set}\PYG{+w}{ }\PYG{k}{group }\PYG{n+nv+nvIdentifier}{solvent}\PYG{+w}{ }\PYG{n}{type}\PYG{+w}{ }\PYG{n}{C}
\PYG{k}{set}\PYG{+w}{ }\PYG{k}{group }\PYG{n+nv+nvIdentifier}{solvent}\PYG{+w}{ }\PYG{n}{type}\PYG{o}{/}\PYG{n}{fraction}\PYG{+w}{ }\PYG{l+m}{2}\PYG{+w}{ }\PYG{l+m}{0.5}\PYG{+w}{ }\PYG{l+m}{12393}
\PYG{k}{set}\PYG{+w}{ }\PYG{k}{group }\PYG{n+nv+nvIdentifier}{solvent}\PYG{+w}{ }\PYG{n}{type}\PYG{o}{/}\PYG{n}{fraction}\PYG{+w}{ }\PYG{n}{C}\PYG{+w}{ }\PYG{l+m}{0.5}\PYG{+w}{ }\PYG{l+m}{12393}
\PYG{k}{set}\PYG{+w}{ }\PYG{k}{group }\PYG{n+nv+nvIdentifier}{edge}\PYG{+w}{ }\PYG{n}{bond}\PYG{+w}{ }\PYG{l+m}{4}
\PYG{k}{set}\PYG{+w}{ }\PYG{k}{region }\PYG{n+nv+nvIdentifier}{half}\PYG{+w}{ }\PYG{n}{charge}\PYG{+w}{ }\PYG{l+m}{0.5}
\PYG{k}{set}\PYG{+w}{ }\PYG{n}{type}\PYG{+w}{ }\PYG{l+m}{3}\PYG{+w}{ }\PYG{n}{charge}\PYG{+w}{ }\PYG{l+m}{0.5}
\PYG{k}{set}\PYG{+w}{ }\PYG{n}{type}\PYG{+w}{ }\PYG{n}{H}\PYG{+w}{ }\PYG{n}{charge}\PYG{+w}{ }\PYG{l+m}{0.5}
\PYG{k}{set}\PYG{+w}{ }\PYG{n}{type}\PYG{+w}{ }\PYG{l+m}{1}\PYG{o}{*}\PYG{l+m}{3}\PYG{+w}{ }\PYG{n}{charge}\PYG{+w}{ }\PYG{l+m}{0.5}
\PYG{k}{set}\PYG{+w}{ }\PYG{n}{atom}\PYG{+w}{ }\PYG{o}{*}\PYG{+w}{ }\PYG{n}{charge}\PYG{+w}{ }\PYG{n}{v\PYGZus{}atomfile}
\PYG{k}{set}\PYG{+w}{ }\PYG{n}{atom}\PYG{+w}{ }\PYG{l+m}{100}\PYG{o}{*}\PYG{l+m}{200}\PYG{+w}{ }\PYG{n}{x}\PYG{+w}{ }\PYG{l+m}{0.5}\PYG{+w}{ }\PYG{n}{y}\PYG{+w}{ }\PYG{l+m}{1.0}
\PYG{k}{set}\PYG{+w}{ }\PYG{n}{atom}\PYG{+w}{ }\PYG{l+m}{100}\PYG{+w}{ }\PYG{n}{vx}\PYG{+w}{ }\PYG{l+m}{0.0}\PYG{+w}{ }\PYG{n}{vy}\PYG{+w}{ }\PYG{l+m}{0.0}\PYG{+w}{ }\PYG{n}{vz}\PYG{+w}{ }\PYG{o}{\PYGZhy{}}\PYG{l+m}{1.0}
\PYG{k}{set}\PYG{+w}{ }\PYG{n}{atom}\PYG{+w}{ }\PYG{l+m}{1492}\PYG{+w}{ }\PYG{n}{type}\PYG{+w}{ }\PYG{l+m}{3}
\PYG{k}{set}\PYG{+w}{ }\PYG{n}{atom}\PYG{+w}{ }\PYG{l+m}{1492}\PYG{+w}{ }\PYG{n}{type}\PYG{+w}{ }\PYG{n}{H}
\PYG{k}{set}\PYG{+w}{ }\PYG{n}{atom}\PYG{+w}{ }\PYG{o}{*}\PYG{+w}{ }\PYG{n}{i\PYGZus{}myVal}\PYG{+w}{ }\PYG{l+m}{5}
\PYG{k}{set}\PYG{+w}{ }\PYG{n}{atom}\PYG{+w}{ }\PYG{o}{*}\PYG{+w}{ }\PYG{n}{d2\PYGZus{}Sxyz[1]}\PYG{+w}{ }\PYG{l+m}{6.4}
\end{sphinxVerbatim}


\subsection{Description}
\label{\detokenize{set:description}}
\sphinxAtStartPar
Set one or more properties of one or more atoms.  Since atom
properties are initially assigned by the {\hyperref[\detokenize{read_data::doc}]{\sphinxcrossref{\DUrole{doc}{read\_data}}}},
{\hyperref[\detokenize{read_restart::doc}]{\sphinxcrossref{\DUrole{doc}{read\_restart}}}} or {\hyperref[\detokenize{create_atoms::doc}]{\sphinxcrossref{\DUrole{doc}{create\_atoms}}}}
commands, this command changes those assignments.  This can be useful
for overriding the default values assigned by the
{\hyperref[\detokenize{create_atoms::doc}]{\sphinxcrossref{\DUrole{doc}{create\_atoms}}}} command (e.g. charge = 0.0).  It can
be useful for altering pairwise and molecular force interactions,
since force\sphinxhyphen{}field coefficients are defined in terms of types.  It can
be used to change the labeling of atoms by atom type or molecule ID
when they are output in {\hyperref[\detokenize{dump::doc}]{\sphinxcrossref{\DUrole{doc}{dump}}}} files.  It can also be useful
for debugging purposes; i.e. positioning an atom at a precise location
to compute subsequent forces or energy.

\sphinxAtStartPar
Note that the \sphinxstyleemphasis{style} and \sphinxstyleemphasis{ID} arguments determine which atoms have
their properties reset.  The remaining keywords specify which
properties to reset and what the new values are.  Some strings like
\sphinxstyleemphasis{type} or \sphinxstyleemphasis{mol} can be used as a style and/or a keyword.


\bigskip\hrule\bigskip


\sphinxAtStartPar
This section describes how to select which atoms to change
the properties of, via the \sphinxstyleemphasis{style} and \sphinxstyleemphasis{ID} arguments.

\sphinxAtStartPar
\DUrole{versionmodified,changed}{Changed in version 28Mar2023: }Support for type labels was added for selecting atoms by type

\sphinxAtStartPar
The style \sphinxstyleemphasis{atom} selects all the atoms in a range of atom IDs.

\sphinxAtStartPar
The style \sphinxstyleemphasis{type} selects all the atoms in a range of types or type
labels.  The style \sphinxstyleemphasis{type} selects atoms in one of two ways.  A range
of numeric atom types can be specified.  Or a single atom type label
can be specified, e.g. “C”.  The style \sphinxstyleemphasis{mol} selects all the atoms in
a range of molecule IDs.

\sphinxAtStartPar
In each of the range cases, the range can be specified as a single
numeric value, or a wildcard asterisk can be used to specify a range
of values.  This takes the form “*” or “*n” or “n*” or “m*n”.  For
example, for the style \sphinxstyleemphasis{type}, if N = the number of atom types, then
an asterisk with no numeric values means all types from 1 to N.  A
leading asterisk means all types from 1 to n (inclusive).  A trailing
asterisk means all types from n to N (inclusive).  A middle asterisk
means all types from m to n (inclusive).  For all the styles except
\sphinxstyleemphasis{mol}, the lowest value for the wildcard is 1; for \sphinxstyleemphasis{mol} it is 0.

\sphinxAtStartPar
The style \sphinxstyleemphasis{group} selects all the atoms in the specified group.  The
style \sphinxstyleemphasis{region} selects all the atoms in the specified geometric
region.  See the {\hyperref[\detokenize{group::doc}]{\sphinxcrossref{\DUrole{doc}{group}}}} and {\hyperref[\detokenize{region::doc}]{\sphinxcrossref{\DUrole{doc}{region}}}} commands
for details of how to specify a group or region.


\bigskip\hrule\bigskip


\sphinxAtStartPar
This section describes the keyword options for which properties to
change, for the selected atoms.

\sphinxAtStartPar
Note that except where explicitly prohibited below, all of the
keywords allow an {\hyperref[\detokenize{variable::doc}]{\sphinxcrossref{\DUrole{doc}{atom\sphinxhyphen{}style or atomfile\sphinxhyphen{}style variable}}}} to be used as the specified value(s).  If the value is a
variable, it should be specified as v\_name, where name is the
variable name.  In this case, the variable will be evaluated, and its
resulting per\sphinxhyphen{}atom value used to determine the value assigned to each
selected atom.  Note that the per\sphinxhyphen{}atom value from the variable will be
ignored for atoms that are not selected via the \sphinxstyleemphasis{style} and \sphinxstyleemphasis{ID}
settings explained above.  A simple way to use per\sphinxhyphen{}atom values from
the variable to reset a property for all atoms is to use style \sphinxstyleemphasis{atom}
with \sphinxstyleemphasis{ID} = “*”; this selects all atom IDs.

\sphinxAtStartPar
Atom\sphinxhyphen{}style variables can specify formulas with various mathematical
functions, and include {\hyperref[\detokenize{thermo_style::doc}]{\sphinxcrossref{\DUrole{doc}{thermo\_style}}}} command
keywords for the simulation box parameters and timestep and elapsed
time.  They can also include per\sphinxhyphen{}atom values, such as atom
coordinates.  Thus it is easy to specify a time\sphinxhyphen{}dependent or
spatially\sphinxhyphen{}dependent set of per\sphinxhyphen{}atom values.  As explained on the
{\hyperref[\detokenize{variable::doc}]{\sphinxcrossref{\DUrole{doc}{variable}}}} doc page, atomfile\sphinxhyphen{}style variables can be
used in place of atom\sphinxhyphen{}style variables, and thus as arguments to the
set command.  Atomfile\sphinxhyphen{}style variables read their per\sphinxhyphen{}atoms values
from a file.

\begin{sphinxadmonition}{note}{Note:}
\sphinxAtStartPar
Atom\sphinxhyphen{}style and atomfile\sphinxhyphen{}style variables return floating point
per\sphinxhyphen{}atom values.  If the values are assigned to an integer variable,
such as the molecule ID, then the floating point value is truncated to
its integer portion, e.g. a value of 2.6 would become 2.
\end{sphinxadmonition}

\sphinxAtStartPar
\DUrole{versionmodified,changed}{Changed in version 28Mar2023: }Support for type labels was added for setting atom, bond, angle,
dihedral, and improper types

\sphinxAtStartPar
Keyword \sphinxstyleemphasis{type} sets the atom type for all selected atoms.  A specified
value can be either a numeric atom type or an atom type label. When
using a numeric type, the specified value must be from 1 to ntypes,
where ntypes was set by the {\hyperref[\detokenize{create_box::doc}]{\sphinxcrossref{\DUrole{doc}{create\_box}}}} command or
the \sphinxstyleemphasis{atom types} field in the header of the data file read by the
{\hyperref[\detokenize{read_data::doc}]{\sphinxcrossref{\DUrole{doc}{read\_data}}}} command.  When using a type label it must
have been defined previously.  See the {\hyperref[\detokenize{Howto_type_labels::doc}]{\sphinxcrossref{\DUrole{doc}{Howto type labels}}}} doc page for the allowed syntax of type labels
and a general discussion of how type labels can be used.

\sphinxAtStartPar
Keyword \sphinxstyleemphasis{type/fraction} sets the atom type for a fraction of the selected
atoms.  The actual number of atoms changed is not guaranteed
to be exactly the specified fraction (0 \textless{}= \sphinxstyleemphasis{fraction} \textless{}= 1), but
should be statistically close.  Random numbers are used in such a way
that a particular atom is changed or not changed, regardless of how
many processors are being used.  This keyword does not allow use of an
atom\sphinxhyphen{}style variable.

\sphinxAtStartPar
Keywords \sphinxstyleemphasis{type/ratio} and \sphinxstyleemphasis{type/subset} also set the atom type for a
fraction of the selected atoms.  The actual number of atoms changed
will be exactly the requested number.  For \sphinxstyleemphasis{type/ratio} the specified
fraction (0 \textless{}= \sphinxstyleemphasis{fraction} \textless{}= 1) determines the number.  For
\sphinxstyleemphasis{type/subset}, the specified \sphinxstyleemphasis{Nsubset} is the number.  An iterative
algorithm is used which ensures the correct number of atoms are
selected, in a perfectly random fashion.  Which atoms are selected
will change with the number of processors used.  These keywords do not
allow use of an atom\sphinxhyphen{}style variable.

\sphinxAtStartPar
Keyword \sphinxstyleemphasis{mol} sets the molecule ID for all selected atoms.  The
{\hyperref[\detokenize{atom_style::doc}]{\sphinxcrossref{\DUrole{doc}{atom style}}}} being used must support the use of
molecule IDs.

\sphinxAtStartPar
Keywords \sphinxstyleemphasis{x}, \sphinxstyleemphasis{y}, \sphinxstyleemphasis{z}, and \sphinxstyleemphasis{charge} set the coordinates or
charge of all selected atoms.  For \sphinxstyleemphasis{charge}, the {\hyperref[\detokenize{atom_style::doc}]{\sphinxcrossref{\DUrole{doc}{atom style}}}} being used must support the use of atomic
charge. Keywords \sphinxstyleemphasis{vx}, \sphinxstyleemphasis{vy}, and \sphinxstyleemphasis{vz} set the velocities of all
selected atoms.

\sphinxAtStartPar
Keyword \sphinxstyleemphasis{dipole} uses the specified x,y,z values as components of a
vector to set as the orientation of the dipole moment vectors of the
selected atoms.  The magnitude of the dipole moment is set by the
length of this orientation vector.

\sphinxAtStartPar
Keyword \sphinxstyleemphasis{dipole/random} randomizes the orientation of the dipole
moment vectors for the selected atoms and sets the magnitude of each
to the specified \sphinxstyleemphasis{Dlen} value.  For 2d systems, the z component of the
orientation is set to 0.0.  Random numbers are used in such a way that
the orientation of a particular atom is the same, regardless of how
many processors are being used.  This keyword does not allow use of an
atom\sphinxhyphen{}style variable.

\sphinxAtStartPar
\DUrole{versionmodified,changed}{Changed in version 15Sep2022.}

\sphinxAtStartPar
Keyword \sphinxstyleemphasis{spin/atom} uses the specified g value to set the magnitude of the
magnetic spin vectors, and the x,y,z values as components of a vector
to set as the orientation of the magnetic spin vectors of the selected
atoms.  This keyword was previously called \sphinxstyleemphasis{spin}.

\sphinxAtStartPar
\DUrole{versionmodified,changed}{Changed in version 15Sep2022.}

\sphinxAtStartPar
Keyword \sphinxstyleemphasis{spin/atom/random} randomizes the orientation of the magnetic spin
vectors for the selected atoms and sets the magnitude of each to the
specified \sphinxstyleemphasis{Dlen} value.  This keyword was previously called \sphinxstyleemphasis{spin/random}.

\sphinxAtStartPar
\DUrole{versionmodified,added}{Added in version 15Sep2022.}

\sphinxAtStartPar
Keyword \sphinxstyleemphasis{radius/electron} uses the specified value to set the radius of
electrons or fixed cores.

\sphinxAtStartPar
\DUrole{versionmodified,added}{Added in version 15Sep2022.}

\sphinxAtStartPar
Keyword \sphinxstyleemphasis{spin/electron} sets the spin of an electron (+/\sphinxhyphen{} 1) or indicates
nuclei (=0), fixed\sphinxhyphen{}cores (=2), or pseudo\sphinxhyphen{}cores (= 3).

\sphinxAtStartPar
Keyword \sphinxstyleemphasis{quat} uses the specified values to create a quaternion
(4\sphinxhyphen{}vector) that represents the orientation of the selected atoms.  The
particles must define a quaternion for their orientation
(e.g. ellipsoids, triangles, body particles) as defined by the
{\hyperref[\detokenize{atom_style::doc}]{\sphinxcrossref{\DUrole{doc}{atom\_style}}}} command.  Note that particles defined by
{\hyperref[\detokenize{atom_style::doc}]{\sphinxcrossref{\DUrole{doc}{atom\_style ellipsoid}}}} have 3 shape parameters.  The 3
values must be non\sphinxhyphen{}zero for each particle set by this command.  They
are used to specify the aspect ratios of an ellipsoidal particle,
which is oriented by default with its x\sphinxhyphen{}axis along the simulation
box’s x\sphinxhyphen{}axis, and similarly for y and z.  If this body is rotated (via
the right\sphinxhyphen{}hand rule) by an angle theta around a unit rotation vector
(a,b,c), then the quaternion that represents its new orientation is
given by (cos(theta/2), a*sin(theta/2), b*sin(theta/2),
c*sin(theta/2)).  The theta and a,b,c values are the arguments to the
\sphinxstyleemphasis{quat} keyword.  LAMMPS normalizes the quaternion in case (a,b,c) was
not specified as a unit vector.  For 2d systems, the a,b,c values are
ignored, since a rotation vector of (0,0,1) is the only valid choice.

\sphinxAtStartPar
Keyword \sphinxstyleemphasis{quat/random} randomizes the orientation of the quaternion for
the selected atoms.  The particles must define a quaternion for their
orientation (e.g. ellipsoids, triangles, body particles) as defined by
the {\hyperref[\detokenize{atom_style::doc}]{\sphinxcrossref{\DUrole{doc}{atom\_style}}}} command.  Random numbers are used in
such a way that the orientation of a particular atom is the same,
regardless of how many processors are being used.  For 2d systems,
only orientations in the xy plane are generated.  As with keyword
\sphinxstyleemphasis{quat}, for ellipsoidal particles, the 3 shape values must be non\sphinxhyphen{}zero
for each particle set by this command.  This keyword does not allow
use of an atom\sphinxhyphen{}style variable.

\sphinxAtStartPar
Keyword \sphinxstyleemphasis{diameter} sets the size of the selected atoms.  The particles
must be finite\sphinxhyphen{}size spheres as defined by the {\hyperref[\detokenize{atom_style::doc}]{\sphinxcrossref{\DUrole{doc}{atom\_style sphere}}}} command.  The diameter of a particle can be set to 0.0,
which means they will be treated as point particles.  Note that this
command does not adjust the particle mass, even if it was defined with
a density, e.g. via the {\hyperref[\detokenize{read_data::doc}]{\sphinxcrossref{\DUrole{doc}{read\_data}}}} command.

\sphinxAtStartPar
Keyword \sphinxstyleemphasis{shape} sets the size and shape of the selected atoms.  The
particles must be ellipsoids as defined by the {\hyperref[\detokenize{atom_style::doc}]{\sphinxcrossref{\DUrole{doc}{atom\_style
ellipsoid}}}} command.  The \sphinxstyleemphasis{Sx}, \sphinxstyleemphasis{Sy}, \sphinxstyleemphasis{Sz} settings
are the 3 diameters of the ellipsoid in each direction.  All 3 can be
set to the same value, which means the ellipsoid is effectively a
sphere.  They can also all be set to 0.0 which means the particle will
be treated as a point particle.  Note that this command does not
adjust the particle mass, even if it was defined with a density,
e.g. via the {\hyperref[\detokenize{read_data::doc}]{\sphinxcrossref{\DUrole{doc}{read\_data}}}} command.

\sphinxAtStartPar
Keyword \sphinxstyleemphasis{length} sets the length of selected atoms.  The particles
must be line segments as defined by the {\hyperref[\detokenize{atom_style::doc}]{\sphinxcrossref{\DUrole{doc}{atom\_style line}}}} command.  If the specified value is non\sphinxhyphen{}zero the line
segment is (re)set to a length = the specified value, centered around
the particle position, with an orientation along the x\sphinxhyphen{}axis.  If the
specified value is 0.0, the particle will become a point particle.
Note that this command does not adjust the particle mass, even if it
was defined with a density, e.g. via the {\hyperref[\detokenize{read_data::doc}]{\sphinxcrossref{\DUrole{doc}{read\_data}}}}
command.

\sphinxAtStartPar
Keyword \sphinxstyleemphasis{tri} sets the size of selected atoms.  The particles must be
triangles as defined by the {\hyperref[\detokenize{atom_style::doc}]{\sphinxcrossref{\DUrole{doc}{atom\_style tri}}}} command.
If the specified value is non\sphinxhyphen{}zero the triangle is (re)set to be an
equilateral triangle in the xy plane with side length = the specified
value, with a centroid at the particle position, with its base
parallel to the x axis, and the y\sphinxhyphen{}axis running from the center of the
base to the top point of the triangle.  If the specified value is 0.0,
the particle will become a point particle.  Note that this command
does not adjust the particle mass, even if it was defined with a
density, e.g. via the {\hyperref[\detokenize{read_data::doc}]{\sphinxcrossref{\DUrole{doc}{read\_data}}}} command.

\sphinxAtStartPar
Keyword \sphinxstyleemphasis{theta} sets the orientation of selected atoms.  The particles
must be line segments as defined by the {\hyperref[\detokenize{atom_style::doc}]{\sphinxcrossref{\DUrole{doc}{atom\_style line}}}} command.  The specified value is used to set the
orientation angle of the line segments with respect to the x axis.

\sphinxAtStartPar
Keyword \sphinxstyleemphasis{theta/random} randomizes the orientation of theta for the
selected atoms.  The particles must be line segments as defined by the
{\hyperref[\detokenize{atom_style::doc}]{\sphinxcrossref{\DUrole{doc}{atom\_style line}}}} command.  Random numbers are used in
such a way that the orientation of a particular atom is the same,
regardless of how many processors are being used.  This keyword does
not allow use of an atom\sphinxhyphen{}style variable.

\sphinxAtStartPar
Keyword \sphinxstyleemphasis{angmom} sets the angular momentum of selected atoms.  The
particles must be ellipsoids as defined by the {\hyperref[\detokenize{atom_style::doc}]{\sphinxcrossref{\DUrole{doc}{atom\_style
ellipsoid}}}} command or triangles as defined by the
{\hyperref[\detokenize{atom_style::doc}]{\sphinxcrossref{\DUrole{doc}{atom\_style tri}}}} command.  The angular momentum
vector of the particles is set to the 3 specified components.

\sphinxAtStartPar
Keyword \sphinxstyleemphasis{omega} sets the angular velocity of selected atoms.  The
particles must be spheres as defined by the {\hyperref[\detokenize{atom_style::doc}]{\sphinxcrossref{\DUrole{doc}{atom\_style sphere}}}} command.  The angular velocity vector of the particles is
set to the 3 specified components.

\sphinxAtStartPar
Keyword \sphinxstyleemphasis{mass} sets the mass of all selected particles.  The particles
must have a per\sphinxhyphen{}atom mass attribute, as defined by the {\hyperref[\detokenize{atom_style::doc}]{\sphinxcrossref{\DUrole{doc}{atom\_style}}}} command.  See the “mass” command for how to set mass
values on a per\sphinxhyphen{}type basis.

\sphinxAtStartPar
Keyword \sphinxstyleemphasis{density} or \sphinxstyleemphasis{density/disc} also sets the mass of all selected
particles, but in a different way.  The particles must have a per\sphinxhyphen{}atom
mass attribute, as defined by the {\hyperref[\detokenize{atom_style::doc}]{\sphinxcrossref{\DUrole{doc}{atom\_style}}}}
command.  If the atom has a radius attribute (see {\hyperref[\detokenize{atom_style::doc}]{\sphinxcrossref{\DUrole{doc}{atom\_style
sphere}}}}) and its radius is non\sphinxhyphen{}zero, its mass is set from
the density and particle volume for 3d systems (the input density is
assumed to be in mass/distance\textasciicircum{}3 units).  For 2d, the default is for
LAMMPS to model particles with a radius attribute as spheres.  However,
if the \sphinxstyleemphasis{density/disc} keyword is used, then they can be modeled as 2d
discs (circles).  Their mass is set from the density and particle area
(the input density is assumed to be in mass/distance\textasciicircum{}2 units).

\sphinxAtStartPar
If the atom has a shape attribute (see {\hyperref[\detokenize{atom_style::doc}]{\sphinxcrossref{\DUrole{doc}{atom\_style ellipsoid}}}}) and its 3 shape parameters are non\sphinxhyphen{}zero, then its mass is
set from the density and particle volume (the input density is assumed
to be in mass/distance\textasciicircum{}3 units).  The \sphinxstyleemphasis{density/disc} keyword has no
effect; it does not (yet) treat 3d ellipsoids as 2d ellipses.

\sphinxAtStartPar
If the atom has a length attribute (see {\hyperref[\detokenize{atom_style::doc}]{\sphinxcrossref{\DUrole{doc}{atom\_style line}}}}) and its length is non\sphinxhyphen{}zero, then its mass is set from the
density and line segment length (the input density is assumed to be in
mass/distance units).  If the atom has an area attribute (see
{\hyperref[\detokenize{atom_style::doc}]{\sphinxcrossref{\DUrole{doc}{atom\_style tri}}}}) and its area is non\sphinxhyphen{}zero, then its
mass is set from the density and triangle area (the input density is
assumed to be in mass/distance\textasciicircum{}2 units).

\sphinxAtStartPar
If none of these cases are valid, then the mass is set to the density
value directly (the input density is assumed to be in mass units).

\sphinxAtStartPar
Keyword \sphinxstyleemphasis{temperature} sets the temperature of a finite\sphinxhyphen{}size particle.
Currently, only the GRANULAR package supports this attribute. The
temperature must be added using an instance of
{\hyperref[\detokenize{fix_property_atom::doc}]{\sphinxcrossref{\DUrole{doc}{fix property/atom}}}} The values for the
temperature must be positive.

\sphinxAtStartPar
Keyword \sphinxstyleemphasis{volume} sets the volume of all selected particles.  Currently,
only the {\hyperref[\detokenize{atom_style::doc}]{\sphinxcrossref{\DUrole{doc}{atom\_style peri}}}} command defines particles
with a volume attribute.  Note that this command does not adjust the
particle mass.

\sphinxAtStartPar
Keyword \sphinxstyleemphasis{image} sets which image of the simulation box the atom is
considered to be in.  An image of 0 means it is inside the box as
defined.  A value of 2 means add 2 box lengths to get the true value.  A
value of \sphinxhyphen{}1 means subtract 1 box length to get the true value.  LAMMPS
updates these flags as atoms cross periodic boundaries during the
simulation.  The flags can be output with atom snapshots via the
{\hyperref[\detokenize{dump::doc}]{\sphinxcrossref{\DUrole{doc}{dump}}}} command.  If a value of NULL is specified for any of
nx,ny,nz, then the current image value for that dimension is unchanged.
For non\sphinxhyphen{}periodic dimensions only a value of 0 can be specified.  This
command can be useful after a system has been equilibrated and atoms
have diffused one or more box lengths in various directions.  This
command can then reset the image values for atoms so that they are
effectively inside the simulation box, e.g if a diffusion coefficient is
about to be measured via the {\hyperref[\detokenize{compute_msd::doc}]{\sphinxcrossref{\DUrole{doc}{compute msd}}}} command.
Care should be taken not to reset the image flags of two atoms in a bond
to the same value if the bond straddles a periodic boundary (rather they
should be different by +/\sphinxhyphen{} 1).  This will not affect the dynamics of a
simulation, but may mess up analysis of the trajectories if a LAMMPS
diagnostic or your own analysis relies on the image flags to unwrap a
molecule which straddles the periodic box.

\sphinxAtStartPar
Keywords \sphinxstyleemphasis{bond}, \sphinxstyleemphasis{angle}, \sphinxstyleemphasis{dihedral}, and \sphinxstyleemphasis{improper}, set the bond
type (angle type, etc) of all bonds (angles, etc) of selected atoms to
the specified value.  The value can be a numeric type from 1 to
nbondtypes (nangletypes, etc).  Or it can be a type label (bond type
label, angle type label, etc).  See the {\hyperref[\detokenize{Howto_type_labels::doc}]{\sphinxcrossref{\DUrole{doc}{Howto type labels}}}} doc page for the allowed syntax of type labels
and a general discussion of how type labels can be used.  All atoms in
a particular bond (angle, etc) must be selected atoms in order for the
change to be made.  The value of nbondtypes (nangletypes, etc) was set
by the \sphinxstyleemphasis{bond types} (\sphinxstyleemphasis{angle types}, etc) field in the header of the
data file read by the {\hyperref[\detokenize{read_data::doc}]{\sphinxcrossref{\DUrole{doc}{read\_data}}}} command.  These
keywords do not allow use of an atom\sphinxhyphen{}style variable.

\sphinxAtStartPar
Keywords \sphinxstyleemphasis{rheo/rho} and \sphinxstyleemphasis{rheo/status} set the density and the status of
rheo particles. In particular, one can only set the phase in the status
as described by the {\hyperref[\detokenize{Howto_rheo::doc}]{\sphinxcrossref{\DUrole{doc}{RHEO howto page}}}}.

\sphinxAtStartPar
Keywords \sphinxstyleemphasis{sph/e}, \sphinxstyleemphasis{sph/cv}, and \sphinxstyleemphasis{sph/rho} set the energy, heat capacity,
and density of smoothed particle hydrodynamics (SPH) particles.  See
\sphinxhref{PDF/SPH\_LAMMPS\_userguide.pdf}{this PDF guide} to using SPH in LAMMPS.

\sphinxAtStartPar
Keyword \sphinxstyleemphasis{smd/mass/density} sets the mass of all selected particles, but
it is only applicable to the Smooth Mach Dynamics package MACHDYN.  It
assumes that the particle volume has already been correctly set and
calculates particle mass from the provided mass density value.

\sphinxAtStartPar
Keyword \sphinxstyleemphasis{smd/contact/radius} only applies to simulations with the Smooth
Mach Dynamics package MACHDYN.  Itsets an interaction radius for
computing short\sphinxhyphen{}range interactions, e.g. repulsive forces to prevent
different individual physical bodies from penetrating each other. Note
that the SPH smoothing kernel diameter used for computing long range,
nonlocal interactions, is set using the \sphinxstyleemphasis{diameter} keyword.

\sphinxAtStartPar
Keyword \sphinxstyleemphasis{dpd/theta} sets the internal temperature of a DPD particle as
defined by the DPD\sphinxhyphen{}REACT package.  If the specified value is a number it
must be \textgreater{}= 0.0.  If the specified value is NULL, then the kinetic
temperature Tkin of each particle is computed as 3/2 k Tkin = KE = 1/2 m
v\textasciicircum{}2 = 1/2 m (vx*vx+vy*vy+vz*vz).  Each particle’s internal
temperature is set to Tkin.  If the specified value is an atom\sphinxhyphen{}style
variable, then the variable is evaluated for each particle.  If a value
\textgreater{}= 0.0, the internal temperature is set to that value.  If it is \textless{} 0.0,
the computation of Tkin is performed and the internal temperature is set
to that value.

\sphinxAtStartPar
Keywords \sphinxstyleemphasis{edpd/temp} and \sphinxstyleemphasis{edpd/cv} set the temperature and volumetric
heat capacity of an eDPD particle as defined by the DPD\sphinxhyphen{}MESO package.
Currently, only {\hyperref[\detokenize{atom_style::doc}]{\sphinxcrossref{\DUrole{doc}{atom\_style edpd}}}} defines particles
with these attributes. The values for the temperature and heat capacity
must be positive.

\sphinxAtStartPar
Keyword \sphinxstyleemphasis{cc} sets the chemical concentration of a tDPD particle for a
specified species as defined by the DPD\sphinxhyphen{}MESO package. Currently, only
{\hyperref[\detokenize{atom_style::doc}]{\sphinxcrossref{\DUrole{doc}{atom\_style tdpd}}}} defines particles with this
attribute. An integer for “index” selects a chemical species (1 to
Nspecies) where Nspecies is set by the atom\_style command. The value for
the chemical concentration must be \textgreater{}= 0.0.

\sphinxAtStartPar
Keyword \sphinxstyleemphasis{epsilon} sets the dielectric constant of a particle, precisely
of the medium where the particle resides as defined by the DIELECTRIC
package. Currently, only {\hyperref[\detokenize{atom_style::doc}]{\sphinxcrossref{\DUrole{doc}{atom\_style dielectric}}}}
defines particles with this attribute. The value for the dielectric
constant must be \textgreater{}= 0.0.  Note that the set command with this keyword
will rescale the particle charge accordingly so that the real charge
(e.g., as read from a data file) stays intact. To change the real
charges, one needs to use the set command with the \sphinxstyleemphasis{charge}
keyword. Care must be taken to ensure that the real and scaled charges,
and dielectric constants are consistent.

\sphinxAtStartPar
Keywords \sphinxstyleemphasis{i\_name}, \sphinxstyleemphasis{d\_name}, \sphinxstyleemphasis{i2\_name}, \sphinxstyleemphasis{d2\_name} refer to custom
per\sphinxhyphen{}atom integer and floating\sphinxhyphen{}point vectors or arrays that have been
added via the {\hyperref[\detokenize{fix_property_atom::doc}]{\sphinxcrossref{\DUrole{doc}{fix property/atom}}}} command.
When that command is used specific names are given to each attribute
which are the “name” portion of these keywords.  For arrays \sphinxstyleemphasis{i2\_name}
and \sphinxstyleemphasis{d2\_name}, the column of the array must also be included following
the name in brackets: e.g. d2\_xyz{[}2{]}, i2\_mySpin{[}3{]}.


\subsection{Restrictions}
\label{\detokenize{set:restrictions}}
\sphinxAtStartPar
You cannot set an atom attribute (e.g. \sphinxstyleemphasis{mol} or \sphinxstyleemphasis{q} or \sphinxstyleemphasis{volume}) if
the {\hyperref[\detokenize{atom_style::doc}]{\sphinxcrossref{\DUrole{doc}{atom\_style}}}} does not have that attribute.

\sphinxAtStartPar
This command requires inter\sphinxhyphen{}processor communication to coordinate the
setting of bond types (angle types, etc).  This means that your system
must be ready to perform a simulation before using one of these
keywords (force fields set, atom mass set, etc).  This is not
necessary for other keywords.

\sphinxAtStartPar
Using the \sphinxstyleemphasis{region} style with the bond (angle, etc) keywords can give
unpredictable results if there are bonds (angles, etc) that straddle
periodic boundaries.  This is because the region may only extend up to
the boundary and partner atoms in the bond (angle, etc) may have
coordinates outside the simulation box if they are ghost atoms.


\subsection{Related commands}
\label{\detokenize{set:related-commands}}
\sphinxAtStartPar
{\hyperref[\detokenize{create_box::doc}]{\sphinxcrossref{\DUrole{doc}{create\_box}}}}, {\hyperref[\detokenize{create_atoms::doc}]{\sphinxcrossref{\DUrole{doc}{create\_atoms}}}},
{\hyperref[\detokenize{read_data::doc}]{\sphinxcrossref{\DUrole{doc}{read\_data}}}}


\subsection{Default}
\label{\detokenize{set:default}}
\sphinxAtStartPar
none

\sphinxstepscope

\index{shell@\spxentry{shell}}\ignorespaces 

\section{shell command}
\label{\detokenize{shell:shell-command}}\label{\detokenize{shell:index-0}}\label{\detokenize{shell::doc}}

\subsection{Syntax}
\label{\detokenize{shell:syntax}}
\begin{sphinxVerbatim}[commandchars=\\\{\}]
\PYG{k}{shell}\PYG{+w}{ }\PYG{n}{command}\PYG{+w}{ }\PYG{n}{args}
\end{sphinxVerbatim}
\begin{itemize}
\item {} 
\sphinxAtStartPar
command = \sphinxstyleemphasis{cd} or \sphinxstyleemphasis{mkdir} or \sphinxstyleemphasis{mv} or \sphinxstyleemphasis{rm} or \sphinxstyleemphasis{rmdir} or \sphinxstyleemphasis{putenv} or arbitrary command
\begin{sphinxalltt}
\sphinxstyleemphasis{cd} arg = dir
  dir = directory to change to
\sphinxstyleemphasis{mkdir} args = dir1 dir2 ...
  dir1,dir2 = one or more directories to create
\sphinxstyleemphasis{mv} args = old new
  old = old filename
  new = new filename or destination folder
\sphinxstyleemphasis{rm} args = {[}\sphinxhyphen{}f{]} file1 file2 ...
  \sphinxhyphen{}f = turn off warnings (optional)
  file1,file2 = one or more filenames to delete
\sphinxstyleemphasis{rmdir} args = dir1 dir2 ...
  dir1,dir2 = one or more directories to delete
\sphinxstyleemphasis{putenv} args = var1=value1 var2=value2
  var=value = one of more definitions of environment variables
anything else is passed as a command to the shell for direct execution
\end{sphinxalltt}

\end{itemize}


\subsection{Examples}
\label{\detokenize{shell:examples}}
\begin{sphinxVerbatim}[commandchars=\\\{\}]
\PYG{k}{shell}\PYG{+w}{ }\PYG{n}{cd}\PYG{+w}{ }\PYG{n}{sub1}
\PYG{k}{shell}\PYG{+w}{ }\PYG{n}{cd}\PYG{+w}{ }\PYG{n}{..}
\PYG{k}{shell}\PYG{+w}{ }\PYG{n}{mkdir}\PYG{+w}{ }\PYG{n}{tmp1}\PYG{+w}{ }\PYG{n}{tmp2}\PYG{o}{/}\PYG{n}{tmp3}
\PYG{k}{shell}\PYG{+w}{ }\PYG{n}{rmdir}\PYG{+w}{ }\PYG{n}{tmp1}\PYG{+w}{ }\PYG{n}{tmp2}
\PYG{k}{shell}\PYG{+w}{ }\PYG{n}{mv}\PYG{+w}{ }\PYG{n}{log.lammps}\PYG{+w}{ }\PYG{n}{hold}\PYG{o}{/}\PYG{n}{log.1}
\PYG{k}{shell}\PYG{+w}{ }\PYG{n}{rm}\PYG{+w}{ }\PYG{n}{TMP}\PYG{o}{/}\PYG{n}{file1}\PYG{+w}{ }\PYG{n}{TMP}\PYG{o}{/}\PYG{n}{file2}
\PYG{k}{shell}\PYG{+w}{ }\PYG{n}{putenv}\PYG{+w}{ }\PYG{n}{LAMMPS\PYGZus{}POTENTIALS}\PYG{o}{=}\PYG{n}{..}\PYG{o}{/}\PYG{n}{..}\PYG{o}{/}\PYG{n}{potentials}
\PYG{k}{shell}\PYG{+w}{ }\PYG{n}{my\PYGZus{}setup}\PYG{+w}{ }\PYG{n}{file1}\PYG{+w}{ }\PYG{l+m}{10}\PYG{+w}{ }\PYG{n}{file2}
\PYG{k}{shell}\PYG{+w}{ }\PYG{n}{my\PYGZus{}post\PYGZus{}process}\PYG{+w}{ }\PYG{l+m}{100}\PYG{+w}{ }\PYG{n}{dump.out}
\end{sphinxVerbatim}


\subsection{Description}
\label{\detokenize{shell:description}}
\sphinxAtStartPar
Execute a shell command.  A few simple file\sphinxhyphen{}based shell commands are
supported directly, in Unix\sphinxhyphen{}style syntax.  Any command not listed
above is passed as\sphinxhyphen{}is to the C\sphinxhyphen{}library system() call, which invokes
the command in a shell.  To use the external executable instead of
the built\sphinxhyphen{}in version one needs to use a full path, for example
\sphinxstyleemphasis{/bin/rm} instead of \sphinxstyleemphasis{rm}.  The built\sphinxhyphen{}in commands will also work
on operating systems, that do not \sphinxhyphen{} by default \sphinxhyphen{} provide the
corresponding external executables (like \sphinxstyleemphasis{mkdir} on Windows).

\sphinxAtStartPar
This command provides a ways to invoke custom commands or executables
from your input script.  For example, you can move files around in
preparation for the next section of the input script.  Or you can run a
program that pre\sphinxhyphen{}processes data for input into LAMMPS.  Or you can run a
program that post\sphinxhyphen{}processes LAMMPS output data.

\sphinxAtStartPar
With the exception of \sphinxstyleemphasis{cd}, all commands, including ones invoked via a
system() call, are executed by only a single processor, so that
files/directories are not being manipulated by multiple processors
concurrently which may result in unexpected errors or corrupted files.

\sphinxAtStartPar
The \sphinxstyleemphasis{cd} command changes the current working directory similar to
the \sphinxcode{\sphinxupquote{cd}} command.  All subsequent LAMMPS commands that read/write files
will use the new directory.  All processors execute this command.

\sphinxAtStartPar
The \sphinxstyleemphasis{mkdir} command creates directories similar to the Unix \sphinxcode{\sphinxupquote{mkdir \sphinxhyphen{}p}}
command.  That is, it will attempt to create the entire path of
subdirectories if they do not exist yet.

\sphinxAtStartPar
The \sphinxstyleemphasis{mv} command renames a file and/or moves it to a new directory.
It cannot rename files across filesystem boundaries or between drives.

\sphinxAtStartPar
The \sphinxstyleemphasis{rm} command deletes file similar to the Unix \sphinxcode{\sphinxupquote{rm}} command.

\sphinxAtStartPar
The \sphinxstyleemphasis{rmdir} command deletes directories similar to Unix \sphinxcode{\sphinxupquote{rmdir}} command.
If a directory is not empty, its contents are also removed recursively
similar to the Unix \sphinxcode{\sphinxupquote{rm \sphinxhyphen{}r}} command.

\sphinxAtStartPar
The \sphinxstyleemphasis{putenv} command defines or updates an environment variable directly.
Since this command does not pass through the shell, no shell variable
expansion or globbing is performed, only the usual substitution for
LAMMPS variables defined with the {\hyperref[\detokenize{variable::doc}]{\sphinxcrossref{\DUrole{doc}{variable}}}} command is
performed.  The resulting string is then used literally.

\sphinxAtStartPar
Any other command is passed as\sphinxhyphen{}is to the shell along with its arguments as
one string, invoked by the C\sphinxhyphen{}library system() call.  For example,
these lines in your input script:

\begin{sphinxVerbatim}[commandchars=\\\{\}]
\PYG{k}{variable }\PYG{n+nv+nvIdentifier}{n}\PYG{+w}{ }\PYG{n}{equal}\PYG{+w}{ }\PYG{l+m}{10}
\PYG{k}{variable }\PYG{n+nv+nvIdentifier}{foo}\PYG{+w}{ }\PYG{n}{string}\PYG{+w}{ }\PYG{n}{file2}
\PYG{k}{shell}\PYG{+w}{ }\PYG{n}{my\PYGZus{}setup}\PYG{+w}{ }\PYG{n}{file1}\PYG{+w}{ }\PYG{n+nv}{\PYGZdl{}n}\PYG{+w}{ }\PYG{n+nv}{\PYGZdl{}\PYGZob{}}\PYG{n+nv}{foo}\PYG{n+nv}{\PYGZcb{}}
\end{sphinxVerbatim}

\sphinxAtStartPar
would be the same as invoking

\begin{sphinxVerbatim}[commandchars=\\\{\}]
my\PYGZus{}setup\PYG{+w}{ }file1\PYG{+w}{ }\PYG{l+m}{10}\PYG{+w}{ }file2
\end{sphinxVerbatim}

\sphinxAtStartPar
from a command\sphinxhyphen{}line prompt.  The executable program “my\_setup” is run
with 3 arguments: file1 10 file2.


\subsection{Restrictions}
\label{\detokenize{shell:restrictions}}
\sphinxAtStartPar
LAMMPS will do a best effort to detect errors and print suitable
warnings, but due to the nature of delegating commands to the C\sphinxhyphen{}library
system() call, this is not always reliable.


\subsection{Related commands}
\label{\detokenize{shell:related-commands}}
\sphinxAtStartPar
none


\subsection{Default}
\label{\detokenize{shell:default}}
\sphinxAtStartPar
none

\sphinxstepscope

\index{special\_bonds@\spxentry{special\_bonds}}\ignorespaces 

\section{special\_bonds command}
\label{\detokenize{special_bonds:special-bonds-command}}\label{\detokenize{special_bonds:index-0}}\label{\detokenize{special_bonds::doc}}

\subsection{Syntax}
\label{\detokenize{special_bonds:syntax}}
\begin{sphinxVerbatim}[commandchars=\\\{\}]
\PYG{k}{special\PYGZus{}bonds}\PYG{+w}{ }\PYG{n}{keyword}\PYG{+w}{ }\PYG{n}{values}\PYG{+w}{ }\PYG{n}{...}
\end{sphinxVerbatim}
\begin{itemize}
\item {} 
\sphinxAtStartPar
one or more keyword/value pairs may be appended

\item {} 
\sphinxAtStartPar
keyword = \sphinxstyleemphasis{amber} or \sphinxstyleemphasis{charmm} or \sphinxstyleemphasis{dreiding} or \sphinxstyleemphasis{fene} or \sphinxstyleemphasis{lj/coul} or \sphinxstyleemphasis{lj} or \sphinxstyleemphasis{coul} or \sphinxstyleemphasis{angle} or \sphinxstyleemphasis{dihedral} or \sphinxstyleemphasis{one/five}
\begin{sphinxalltt}
\sphinxstyleemphasis{amber} values = none
\sphinxstyleemphasis{charmm} values = none
\sphinxstyleemphasis{dreiding} values = none
\sphinxstyleemphasis{fene} values = none
\sphinxstyleemphasis{lj/coul} values = w1,w2,w3
  w1,w2,w3 = weights (0.0 to 1.0) on pairwise Lennard\sphinxhyphen{}Jones and Coulombic interactions
\sphinxstyleemphasis{lj} values = w1,w2,w3
  w1,w2,w3 = weights (0.0 to 1.0) on pairwise Lennard\sphinxhyphen{}Jones interactions
\sphinxstyleemphasis{coul} values = w1,w2,w3
  w1,w2,w3 = weights (0.0 to 1.0) on pairwise Coulombic interactions
\sphinxstyleemphasis{angle} value = \sphinxstyleemphasis{yes} or \sphinxstyleemphasis{no}
\sphinxstyleemphasis{dihedral} value = \sphinxstyleemphasis{yes} or \sphinxstyleemphasis{no}
\sphinxstyleemphasis{one/five} value = \sphinxstyleemphasis{yes} or \sphinxstyleemphasis{no}
\end{sphinxalltt}

\end{itemize}


\subsection{Examples}
\label{\detokenize{special_bonds:examples}}
\begin{sphinxVerbatim}[commandchars=\\\{\}]
\PYG{k}{special\PYGZus{}bonds}\PYG{+w}{ }\PYG{n}{amber}
\PYG{k}{special\PYGZus{}bonds}\PYG{+w}{ }\PYG{n}{charmm}
\PYG{k}{special\PYGZus{}bonds}\PYG{+w}{ }\PYG{n}{fene}\PYG{+w}{ }\PYG{n}{dihedral}\PYG{+w}{ }\PYG{n}{no}
\PYG{k}{special\PYGZus{}bonds}\PYG{+w}{ }\PYG{n}{lj}\PYG{o}{/}\PYG{n}{coul}\PYG{+w}{ }\PYG{l+m}{0.0}\PYG{+w}{ }\PYG{l+m}{0.0}\PYG{+w}{ }\PYG{l+m}{0.5}\PYG{+w}{ }\PYG{n}{angle}\PYG{+w}{ }\PYG{n}{yes}\PYG{+w}{ }\PYG{n}{dihedral}\PYG{+w}{ }\PYG{n}{yes}
\PYG{k}{special\PYGZus{}bonds}\PYG{+w}{ }\PYG{n}{lj}\PYG{+w}{ }\PYG{l+m}{0.0}\PYG{+w}{ }\PYG{l+m}{0.0}\PYG{+w}{ }\PYG{l+m}{0.5}\PYG{+w}{ }\PYG{n}{coul}\PYG{+w}{ }\PYG{l+m}{0.0}\PYG{+w}{ }\PYG{l+m}{0.0}\PYG{+w}{ }\PYG{l+m}{0.0}\PYG{+w}{ }\PYG{n}{dihedral}\PYG{+w}{ }\PYG{n}{yes}
\end{sphinxVerbatim}


\subsection{Description}
\label{\detokenize{special_bonds:description}}
\sphinxAtStartPar
Set weighting coefficients for pairwise energy and force contributions
between pairs of atoms that are also permanently bonded to each other,
either directly or via one or two intermediate bonds.  These weighting
factors are used by nearly all {\hyperref[\detokenize{pair_style::doc}]{\sphinxcrossref{\DUrole{doc}{pair styles}}}} in
LAMMPS that compute simple pairwise interactions.  Permanent bonds
between atoms are specified by defining the bond topology in the data
file read by the {\hyperref[\detokenize{read_data::doc}]{\sphinxcrossref{\DUrole{doc}{read\_data}}}} command.  Typically a
{\hyperref[\detokenize{bond_style::doc}]{\sphinxcrossref{\DUrole{doc}{bond\_style}}}} command is also used to define a bond
potential.  The rationale for using these weighting factors is that
the interaction between a pair of bonded atoms is all (or mostly)
specified by the bond, angle, dihedral potentials, and thus the
non\sphinxhyphen{}bonded Lennard\sphinxhyphen{}Jones or Coulombic interaction between the pair of
atoms should be excluded (or reduced by a weighting factor).

\begin{sphinxadmonition}{note}{Note:}
\sphinxAtStartPar
These weighting factors are NOT used by {\hyperref[\detokenize{pair_style::doc}]{\sphinxcrossref{\DUrole{doc}{pair styles}}}} that compute many\sphinxhyphen{}body interactions, since the
“bonds” that result from such interactions are not permanent, but
are created and broken dynamically as atom conformations change.
Examples of pair styles in this category are EAM, MEAM,
Stillinger\sphinxhyphen{}Weber, Tersoff, COMB, AIREBO, and ReaxFF.  In fact, it
generally makes no sense to define permanent bonds between atoms
that interact via these potentials, though such bonds may exist
elsewhere in your system, e.g. when using the {\hyperref[\detokenize{pair_hybrid::doc}]{\sphinxcrossref{\DUrole{doc}{pair\_style
hybrid}}}} command.  Thus LAMMPS ignores special\_bonds
settings when many\sphinxhyphen{}body potentials are calculated.  Please note,
that the existence of explicit bonds for atoms that are described
by a many\sphinxhyphen{}body potential will alter the neighbor list and thus can
render the computation of those interactions invalid, since those
pairs are not only used to determine direct pairwise interactions
but also neighbors of neighbors and more.  The recommended course
of action is to remove such bonds, or \sphinxhyphen{} if that is not possible \sphinxhyphen{}
use a special bonds setting of 1.0 1.0 1.0.
\end{sphinxadmonition}

\begin{sphinxadmonition}{note}{Note:}
\sphinxAtStartPar
Unlike some commands in LAMMPS, you cannot use this command
multiple times in an incremental fashion: e.g. to first set the LJ
settings and then the Coulombic ones.  Each time you use this
command it sets all the coefficients to default values and only
overrides the one you specify, so you should set all the options
you need each time you use it.  See more details at the bottom of
this page.
\end{sphinxadmonition}

\sphinxAtStartPar
The Coulomb factors are applied to any Coulomb (charge interaction)
term that the potential calculates.  The LJ factors are applied to the
remaining terms that the potential calculates, whether they represent
LJ interactions or not.  The weighting factors are a scaling
prefactor on the energy and force between the pair of atoms.

\sphinxAtStartPar
A value of 1.0 means include the full interaction without flagging the
pair as a “special pair”; a value of 0.0 means exclude the pair
completely from the neighbor list, except for pair styles that require a
{\hyperref[\detokenize{kspace_style::doc}]{\sphinxcrossref{\DUrole{doc}{kspace style}}}} and pair styles {\hyperref[\detokenize{pair_amoeba::doc}]{\sphinxcrossref{\DUrole{doc}{amoeba}}}}, {\hyperref[\detokenize{pair_amoeba::doc}]{\sphinxcrossref{\DUrole{doc}{hippo}}}}, {\hyperref[\detokenize{pair_thole::doc}]{\sphinxcrossref{\DUrole{doc}{thole}}}},
{\hyperref[\detokenize{pair_coul::doc}]{\sphinxcrossref{\DUrole{doc}{coul/exclude}}}}, and pair styles that include
“coul/dsf” or “coul/wolf”.

\begin{sphinxadmonition}{note}{Note:}
\sphinxAtStartPar
To include pairs that would otherwise be excluded (so they are
included in the neighbor list for certain analysis compute styles),
you can use a very small but non\sphinxhyphen{}zero value like 1.0e\sphinxhyphen{}100 instead of
0.0.  Due to using floating\sphinxhyphen{}point math, the computed force, energy,
and virial contributions from the pairs will be too small to cause
differences.
\end{sphinxadmonition}

\sphinxAtStartPar
The first of the 3 coefficients (LJ or Coulombic) is the weighting
factor on 1\sphinxhyphen{}2 atom pairs, which are pairs of atoms directly bonded to
each other.  The second coefficient is the weighting factor on 1\sphinxhyphen{}3
atom pairs which are those separated by 2 bonds (e.g. the two H atoms
in a water molecule).  The third coefficient is the weighting factor
on 1\sphinxhyphen{}4 atom pairs which are those separated by 3 bonds (e.g. the first
and fourth atoms in a dihedral interaction).  Thus if the 1\sphinxhyphen{}2
coefficient is set to 0.0, then the pairwise interaction is
effectively turned off for all pairs of atoms bonded to each other.
If it is set to 1.0, then that interaction will be at full strength.

\begin{sphinxadmonition}{note}{Note:}
\sphinxAtStartPar
For purposes of computing weighted pairwise interactions, 1\sphinxhyphen{}3
and 1\sphinxhyphen{}4 interactions are not defined from the list of angles or
dihedrals used by the simulation.  Rather, they are inferred
topologically from the set of bonds specified when the simulation is
defined from a data or restart file (see {\hyperref[\detokenize{read_data::doc}]{\sphinxcrossref{\DUrole{doc}{read\_data}}}} or
{\hyperref[\detokenize{read_restart::doc}]{\sphinxcrossref{\DUrole{doc}{read\_restart}}}} commands).  Thus the set of
1\sphinxhyphen{}2,1\sphinxhyphen{}3,1\sphinxhyphen{}4 interactions that the weights apply to is the same whether
angle and dihedral potentials are computed or not, and remains the
same even if bonds are constrained, or turned off, or removed during a
simulation.
\end{sphinxadmonition}

\sphinxAtStartPar
The two exceptions to this rule are (a) if the \sphinxstyleemphasis{angle} or \sphinxstyleemphasis{dihedral}
keywords are set to \sphinxstyleemphasis{yes} (see below), or (b) if the
{\hyperref[\detokenize{delete_bonds::doc}]{\sphinxcrossref{\DUrole{doc}{delete\_bonds}}}} command is used with the \sphinxstyleemphasis{special}
option that re\sphinxhyphen{}computes the 1\sphinxhyphen{}2,1\sphinxhyphen{}3,1\sphinxhyphen{}4 topologies after bonds are
deleted; see the {\hyperref[\detokenize{delete_bonds::doc}]{\sphinxcrossref{\DUrole{doc}{delete\_bonds}}}} command for more
details.

\sphinxAtStartPar
The \sphinxstyleemphasis{amber} keyword sets the 3 coefficients to 0.0, 0.0, 0.5 for LJ
interactions and to 0.0, 0.0, 0.8333 for Coulombic interactions, which
is the default for a commonly used version of the AMBER force field,
where the last value is really 5/6.  See {\hyperref[\detokenize{special_bonds:cornell}]{\sphinxcrossref{\DUrole{std,std-ref}{(Cornell)}}}} for a
description of the AMBER force field.

\sphinxAtStartPar
The \sphinxstyleemphasis{charmm} keyword sets the 3 coefficients to 0.0, 0.0, 0.0 for both
LJ and Coulombic interactions, which is the default for a commonly
used version of the CHARMM force field.  Note that in pair styles
\sphinxstyleemphasis{lj/charmm/coul/charmm} and \sphinxstyleemphasis{lj/charmm/coul/long} the 1\sphinxhyphen{}4 coefficients
are defined explicitly, and these pairwise contributions are computed
as part of the charmm dihedral style \sphinxhyphen{} see the
{\hyperref[\detokenize{pair_coeff::doc}]{\sphinxcrossref{\DUrole{doc}{pair\_coeff}}}} and {\hyperref[\detokenize{dihedral_style::doc}]{\sphinxcrossref{\DUrole{doc}{dihedral\_style}}}}
commands for more information.  See {\hyperref[\detokenize{special_bonds:mackerell}]{\sphinxcrossref{\DUrole{std,std-ref}{(MacKerell)}}}} for a
description of the CHARMM force field.

\sphinxAtStartPar
The \sphinxstyleemphasis{dreiding} keyword sets the 3 coefficients to 0.0, 0.0, 1.0 for both
LJ and Coulombic interactions, which is the default for the Dreiding
force field, as discussed in {\hyperref[\detokenize{special_bonds:mayo}]{\sphinxcrossref{\DUrole{std,std-ref}{(Mayo)}}}}.

\sphinxAtStartPar
The \sphinxstyleemphasis{fene} keyword sets the 3 coefficients to 0.0, 1.0, 1.0 for both
LJ and Coulombic interactions, which is consistent with a
coarse\sphinxhyphen{}grained polymer model with {\hyperref[\detokenize{bond_fene::doc}]{\sphinxcrossref{\DUrole{doc}{FENE bonds}}}}.  See
{\hyperref[\detokenize{special_bonds:kremer}]{\sphinxcrossref{\DUrole{std,std-ref}{(Kremer)}}}} for a description of FENE bonds.

\sphinxAtStartPar
The \sphinxstyleemphasis{lj/coul}, \sphinxstyleemphasis{lj}, and \sphinxstyleemphasis{coul} keywords allow the 3 coefficients to
be set explicitly.  The \sphinxstyleemphasis{lj/coul} keyword sets both the LJ and
Coulombic coefficients to the same 3 values.  The \sphinxstyleemphasis{lj} and \sphinxstyleemphasis{coul}
keywords only set either the LJ or Coulombic coefficients.  Use both
of them if you wish to set the LJ coefficients to different values
than the Coulombic coefficients.

\sphinxAtStartPar
The \sphinxstyleemphasis{angle} keyword allows the 1\sphinxhyphen{}3 weighting factor to be ignored for
individual atom pairs if they are not listed as the first and last
atoms in any angle defined in the simulation or as 1,3 or 2,4 atoms in
any dihedral defined in the simulation.  For example, imagine the 1\sphinxhyphen{}3
weighting factor is set to 0.5 and you have a linear molecule with 4
atoms and bonds as follows: 1\sphinxhyphen{}2\sphinxhyphen{}3\sphinxhyphen{}4.  If your data file defines 1\sphinxhyphen{}2\sphinxhyphen{}3
as an angle, but does not define 2\sphinxhyphen{}3\sphinxhyphen{}4 as an angle or 1\sphinxhyphen{}2\sphinxhyphen{}3\sphinxhyphen{}4 as a
dihedral, then the pairwise interaction between atoms 1 and 3 will
always be weighted by 0.5, but different force fields use different
rules for weighting the pairwise interaction between atoms 2 and 4.
If the \sphinxstyleemphasis{angle} keyword is specified as \sphinxstyleemphasis{yes}, then the pairwise
interaction between atoms 2 and 4 will be unaffected (full weighting
of 1.0).  If the \sphinxstyleemphasis{angle} keyword is specified as \sphinxstyleemphasis{no} which is the
default, then the 2,4 interaction will also be weighted by 0.5.

\sphinxAtStartPar
The \sphinxstyleemphasis{dihedral} keyword allows the 1\sphinxhyphen{}4 weighting factor to be ignored
for individual atom pairs if they are not listed as the first and last
atoms in any dihedral defined in the simulation.  For example, imagine
the 1\sphinxhyphen{}4 weighting factor is set to 0.5 and you have a linear molecule
with 5 atoms and bonds as follows: 1\sphinxhyphen{}2\sphinxhyphen{}3\sphinxhyphen{}4\sphinxhyphen{}5.  If your data file
defines 1\sphinxhyphen{}2\sphinxhyphen{}3\sphinxhyphen{}4 as a dihedral, but does not define 2\sphinxhyphen{}3\sphinxhyphen{}4\sphinxhyphen{}5 as a
dihedral, then the pairwise interaction between atoms 1 and 4 will
always be weighted by 0.5, but different force fields use different
rules for weighting the pairwise interaction between atoms 2 and 5.
If the \sphinxstyleemphasis{dihedral} keyword is specified as \sphinxstyleemphasis{yes}, then the pairwise
interaction between atoms 2 and 5 will be unaffected (full weighting
of 1.0).  If the \sphinxstyleemphasis{dihedral} keyword is specified as \sphinxstyleemphasis{no} which is the
default, then the 2,5 interaction will also be weighted by 0.5.

\sphinxAtStartPar
The \sphinxstyleemphasis{one/five} keyword enable calculation and storage of a list of 1\sphinxhyphen{}5
neighbors in the molecular topology for each atom.  It is required by
some pair styles, such as {\hyperref[\detokenize{pair_style::doc}]{\sphinxcrossref{\DUrole{doc}{pair\_style amoeba}}}} and
{\hyperref[\detokenize{pair_style::doc}]{\sphinxcrossref{\DUrole{doc}{pair\_style hippo}}}}.


\bigskip\hrule\bigskip


\begin{sphinxadmonition}{note}{Note:}
\sphinxAtStartPar
LAMMPS stores and maintains a data structure with a list of the
first, second, and third neighbors of each atom (within the bond
topology of the system).  If new bonds are created (or molecules
added containing atoms with more special neighbors), the size of
this list needs to grow.  Note that adding a single bond always
adds a new first neighbor but may also induce *many* new second
and third neighbors, depending on the molecular topology of your
system.  Using the \sphinxstyleemphasis{extra/special/per/atom} keyword to either
{\hyperref[\detokenize{read_data::doc}]{\sphinxcrossref{\DUrole{doc}{read\_data}}}} or {\hyperref[\detokenize{create_box::doc}]{\sphinxcrossref{\DUrole{doc}{create\_box}}}}
reserves empty space in the list for this N additional first,
second, or third neighbors to be added.  If you do not do this, you
may get an error when bonds (or molecules) are added.
\end{sphinxadmonition}


\bigskip\hrule\bigskip


\begin{sphinxadmonition}{note}{Note:}
\sphinxAtStartPar
If you reuse this command in an input script, you should set all
the options you need each time.  This command cannot be used a second
time incrementally.  E.g. these two commands:
\end{sphinxadmonition}

\begin{sphinxVerbatim}[commandchars=\\\{\}]
\PYG{k}{special\PYGZus{}bonds}\PYG{+w}{ }\PYG{n}{lj}\PYG{+w}{ }\PYG{l+m}{0.0}\PYG{+w}{ }\PYG{l+m}{1.0}\PYG{+w}{ }\PYG{l+m}{1.0}
\PYG{k}{special\PYGZus{}bonds}\PYG{+w}{ }\PYG{n}{coul}\PYG{+w}{ }\PYG{l+m}{0.0}\PYG{+w}{ }\PYG{l+m}{0.0}\PYG{+w}{ }\PYG{l+m}{1.0}
\end{sphinxVerbatim}

\sphinxAtStartPar
are not the same as

\begin{sphinxVerbatim}[commandchars=\\\{\}]
\PYG{k}{special\PYGZus{}bonds}\PYG{+w}{ }\PYG{n}{lj}\PYG{+w}{ }\PYG{l+m}{0.0}\PYG{+w}{ }\PYG{l+m}{1.0}\PYG{+w}{ }\PYG{l+m}{1.0}\PYG{+w}{ }\PYG{n}{coul}\PYG{+w}{ }\PYG{l+m}{0.0}\PYG{+w}{ }\PYG{l+m}{0.0}\PYG{+w}{ }\PYG{l+m}{1.0}
\end{sphinxVerbatim}

\sphinxAtStartPar
In the first case you end up with (after the second command):

\begin{sphinxVerbatim}[commandchars=\\\{\}]
LJ: 0.0 0.0 0.0
Coul: 0.0 0.0 1.0
\end{sphinxVerbatim}

\sphinxAtStartPar
while only in the second case do you get the desired settings of:

\begin{sphinxVerbatim}[commandchars=\\\{\}]
LJ: 0.0 1.0 1.0
Coul: 0.0 0.0 1.0
\end{sphinxVerbatim}

\sphinxAtStartPar
This happens because the LJ (and Coul) settings are reset to their
default values before modifying them, each time the \sphinxstyleemphasis{special\_bonds}
command is issued.


\subsection{Restrictions}
\label{\detokenize{special_bonds:restrictions}}
\sphinxAtStartPar
none


\subsection{Related commands}
\label{\detokenize{special_bonds:related-commands}}
\sphinxAtStartPar
{\hyperref[\detokenize{delete_bonds::doc}]{\sphinxcrossref{\DUrole{doc}{delete\_bonds}}}}, {\hyperref[\detokenize{fix_bond_create::doc}]{\sphinxcrossref{\DUrole{doc}{fix bond/create}}}}


\subsection{Default}
\label{\detokenize{special_bonds:default}}
\sphinxAtStartPar
All 3 Lennard\sphinxhyphen{}Jones and 3 Coulombic weighting coefficients = 0.0,
angle = no, dihedral = no.


\bigskip\hrule\bigskip

\phantomsection\label{\detokenize{special_bonds:cornell}}
\sphinxAtStartPar
\sphinxstylestrong{(Cornell)} Cornell, Cieplak, Bayly, Gould, Merz, Ferguson,
Spellmeyer, Fox, Caldwell, Kollman, JACS 117, 5179\sphinxhyphen{}5197 (1995).

\phantomsection\label{\detokenize{special_bonds:kremer}}
\sphinxAtStartPar
\sphinxstylestrong{(Kremer)} Kremer, Grest, J Chem Phys, 92, 5057 (1990).

\phantomsection\label{\detokenize{special_bonds:mackerell}}
\sphinxAtStartPar
\sphinxstylestrong{(MacKerell)} MacKerell, Bashford, Bellott, Dunbrack, Evanseck, Field,
Fischer, Gao, Guo, Ha, et al, J Phys Chem, 102, 3586 (1998).

\phantomsection\label{\detokenize{special_bonds:mayo}}
\sphinxAtStartPar
\sphinxstylestrong{(Mayo)} Mayo, Olfason, Goddard III, J Phys Chem, 94, 8897\sphinxhyphen{}8909
(1990).

\sphinxstepscope

\index{suffix@\spxentry{suffix}}\ignorespaces 

\section{suffix command}
\label{\detokenize{suffix:suffix-command}}\label{\detokenize{suffix:index-0}}\label{\detokenize{suffix::doc}}

\subsection{Syntax}
\label{\detokenize{suffix:syntax}}
\begin{sphinxVerbatim}[commandchars=\\\{\}]
\PYG{k}{suffix}\PYG{+w}{ }\PYG{n}{style}\PYG{+w}{ }\PYG{n}{args}
\end{sphinxVerbatim}
\begin{itemize}
\item {} 
\sphinxAtStartPar
style = \sphinxstyleemphasis{off} or \sphinxstyleemphasis{on} or \sphinxstyleemphasis{gpu} or \sphinxstyleemphasis{intel} or \sphinxstyleemphasis{kk} or \sphinxstyleemphasis{omp} or \sphinxstyleemphasis{opt} or \sphinxstyleemphasis{hybrid}

\item {} 
\sphinxAtStartPar
args = for hybrid style, default suffix to be used and alternative suffix

\end{itemize}


\subsection{Examples}
\label{\detokenize{suffix:examples}}
\begin{sphinxVerbatim}[commandchars=\\\{\}]
\PYG{k}{suffix}\PYG{+w}{ }\PYG{n}{off}
\PYG{k}{suffix}\PYG{+w}{ }\PYG{n}{on}
\PYG{k}{suffix}\PYG{+w}{ }\PYG{n}{gpu}
\PYG{k}{suffix}\PYG{+w}{ }\PYG{n}{intel}
\PYG{k}{suffix}\PYG{+w}{ }\PYG{n}{hybrid}\PYG{+w}{ }\PYG{n}{intel}\PYG{+w}{ }\PYG{n}{omp}
\PYG{k}{suffix}\PYG{+w}{ }\PYG{n}{kk}
\end{sphinxVerbatim}


\subsection{Description}
\label{\detokenize{suffix:description}}
\sphinxAtStartPar
This command allows you to use variants of various styles if they
exist.  In that respect it operates the same as the {\hyperref[\detokenize{Run_options::doc}]{\sphinxcrossref{\DUrole{doc}{\sphinxhyphen{}suffix command\sphinxhyphen{}line switch}}}}.  It also has options to turn
off or back on any suffix setting made via the command line.

\sphinxAtStartPar
The specified style can be \sphinxstyleemphasis{gpu}, \sphinxstyleemphasis{intel}, \sphinxstyleemphasis{kk}, \sphinxstyleemphasis{omp}, \sphinxstyleemphasis{opt} or
\sphinxstyleemphasis{hybrid}. These refer to optional packages that LAMMPS can be built
with, as described on the {\hyperref[\detokenize{Build_package::doc}]{\sphinxcrossref{\DUrole{doc}{Build package}}}} doc page.
The “gpu” style corresponds to the GPU package, the “intel” style to
the INTEL package, the “kk” style to the KOKKOS package, the
“omp” style to the OPENMP package, and the “opt” style to the OPT
package.

\sphinxAtStartPar
These are the variants these packages provide:
\begin{itemize}
\item {} 
\sphinxAtStartPar
GPU = a handful of pair styles and the PPPM kspace\_style, optimized to
run on one or more GPUs or multicore CPU/GPU nodes

\item {} 
\sphinxAtStartPar
INTEL = a collection of pair styles and neighbor routines
optimized to run in single, mixed, or double precision on CPUs and
Intel(R) Xeon Phi(TM) co\sphinxhyphen{}processors.

\item {} 
\sphinxAtStartPar
KOKKOS = a collection of atom, pair, and fix styles optimized to run
using the Kokkos library on various kinds of hardware, including GPUs
via CUDA and many\sphinxhyphen{}core chips via OpenMP or threading.

\item {} 
\sphinxAtStartPar
OPENMP = a collection of pair, bond, angle, dihedral, improper,
kspace, compute, and fix styles with support for OpenMP
multi\sphinxhyphen{}threading

\item {} 
\sphinxAtStartPar
OPT = a handful of pair styles, cache\sphinxhyphen{}optimized for faster CPU
performance

\item {} 
\sphinxAtStartPar
HYBRID = a combination of two packages can be specified (see below)

\end{itemize}

\sphinxAtStartPar
As an example, all of the packages provide a {\hyperref[\detokenize{pair_lj::doc}]{\sphinxcrossref{\DUrole{doc}{pair\_style lj/cut}}}} variant, with style names lj/cut/opt, lj/cut/omp,
lj/cut/gpu, lj/cut/intel, or lj/cut/kk.  A variant styles
can be specified explicitly in your input script, e.g. pair\_style
lj/cut/gpu. If the suffix command is used with the appropriate style,
you do not need to modify your input script.  The specified suffix
(opt,omp,gpu,intel,kk) is automatically appended whenever your
input script command creates a new {\hyperref[\detokenize{atom_style::doc}]{\sphinxcrossref{\DUrole{doc}{atom}}}},
{\hyperref[\detokenize{pair_style::doc}]{\sphinxcrossref{\DUrole{doc}{pair}}}}, {\hyperref[\detokenize{bond_style::doc}]{\sphinxcrossref{\DUrole{doc}{bond}}}},
{\hyperref[\detokenize{angle_style::doc}]{\sphinxcrossref{\DUrole{doc}{angle}}}}, {\hyperref[\detokenize{dihedral_style::doc}]{\sphinxcrossref{\DUrole{doc}{dihedral}}}},
{\hyperref[\detokenize{improper_style::doc}]{\sphinxcrossref{\DUrole{doc}{improper}}}}, {\hyperref[\detokenize{kspace_style::doc}]{\sphinxcrossref{\DUrole{doc}{kspace}}}},
{\hyperref[\detokenize{fix::doc}]{\sphinxcrossref{\DUrole{doc}{fix}}}}, {\hyperref[\detokenize{compute::doc}]{\sphinxcrossref{\DUrole{doc}{compute}}}}, or {\hyperref[\detokenize{run_style::doc}]{\sphinxcrossref{\DUrole{doc}{run}}}} style.
If the variant version does not exist, the standard version is
created.

\sphinxAtStartPar
For “hybrid”, two packages are specified. The first is used whenever
available. If a style with the first suffix is not available, the style
with the suffix for the second package will be used if available. For
example, “hybrid intel omp” will use styles from the INTEL package
as a first choice and styles from the OPENMP package as a second choice
if no INTEL variant is available.

\sphinxAtStartPar
If the specified style is \sphinxstyleemphasis{off}, then any previously specified suffix
is temporarily disabled, whether it was specified by a command\sphinxhyphen{}line
switch or a previous suffix command.  If the specified style is \sphinxstyleemphasis{on},
a disabled suffix is turned back on.  The use of these 2 commands lets
your input script use a standard LAMMPS style (i.e. a non\sphinxhyphen{}accelerated
variant), which can be useful for testing or benchmarking purposes.
Of course this is also possible by not using any suffix commands, and
explicitly appending or not appending the suffix to the relevant
commands in your input script.

\begin{sphinxadmonition}{note}{Note:}
\sphinxAtStartPar
The default {\hyperref[\detokenize{run_style::doc}]{\sphinxcrossref{\DUrole{doc}{run\_style}}}} verlet is invoked prior to
reading the input script and is therefore not affected by a suffix command
in the input script. The KOKKOS package requires “run\_style verlet/kk”,
so when using the KOKKOS package it is necessary to either use the command
line “\sphinxhyphen{}sf kk” command or add an explicit “run\_style verlet” command to the
input script.
\end{sphinxadmonition}


\subsection{Restrictions}
\label{\detokenize{suffix:restrictions}}\begin{quote}

\sphinxAtStartPar
none
\end{quote}


\subsection{Related commands}
\label{\detokenize{suffix:related-commands}}
\sphinxAtStartPar
{\hyperref[\detokenize{Run_options::doc}]{\sphinxcrossref{\DUrole{doc}{\sphinxhyphen{}suffix command\sphinxhyphen{}line switch}}}}


\subsection{Default}
\label{\detokenize{suffix:default}}
\sphinxAtStartPar
none

\sphinxstepscope

\index{tad@\spxentry{tad}}\ignorespaces 

\section{tad command}
\label{\detokenize{tad:tad-command}}\label{\detokenize{tad:index-0}}\label{\detokenize{tad::doc}}

\subsection{Syntax}
\label{\detokenize{tad:syntax}}
\begin{sphinxVerbatim}[commandchars=\\\{\}]
\PYG{k}{tad}\PYG{+w}{ }\PYG{n}{N}\PYG{+w}{ }\PYG{n}{t\PYGZus{}event}\PYG{+w}{ }\PYG{n}{T\PYGZus{}lo}\PYG{+w}{ }\PYG{n}{T\PYGZus{}hi}\PYG{+w}{ }\PYG{n}{delta}\PYG{+w}{ }\PYG{n}{tmax}\PYG{+w}{ }\PYG{n}{compute}\PYG{o}{\PYGZhy{}}\PYG{n}{ID}\PYG{+w}{ }\PYG{n}{keyword}\PYG{+w}{ }\PYG{n}{value}\PYG{+w}{ }\PYG{n}{...}
\end{sphinxVerbatim}
\begin{itemize}
\item {} 
\sphinxAtStartPar
N = \# of timesteps to run (not including dephasing/quenching)

\item {} 
\sphinxAtStartPar
t\_event = timestep interval between event checks

\item {} 
\sphinxAtStartPar
T\_lo = temperature at which event times are desired

\item {} 
\sphinxAtStartPar
T\_hi = temperature at which MD simulation is performed

\item {} 
\sphinxAtStartPar
delta = desired confidence level for stopping criterion

\item {} 
\sphinxAtStartPar
tmax = reciprocal of lowest expected pre\sphinxhyphen{}exponential factor (time units)

\item {} 
\sphinxAtStartPar
compute\sphinxhyphen{}ID = ID of the compute used for event detection

\item {} 
\sphinxAtStartPar
zero or more keyword/value pairs may be appended

\item {} 
\sphinxAtStartPar
keyword = \sphinxstyleemphasis{min} or \sphinxstyleemphasis{neb} or \sphinxstyleemphasis{neb\_style} or \sphinxstyleemphasis{neb\_step} or \sphinxstyleemphasis{neb\_log}
\begin{sphinxalltt}
\sphinxstyleemphasis{min} values = etol ftol maxiter maxeval
  etol = stopping tolerance for energy (energy units)
  ftol = stopping tolerance for force (force units)
  maxiter = max iterations of minimize
  maxeval = max number of force/energy evaluations
\sphinxstyleemphasis{neb} values = ftol N1 N2 Nevery
  etol = stopping tolerance for energy (energy units)
  ftol = stopping tolerance for force (force units)
  N1 = max \# of iterations (timesteps) to run initial NEB
  N2 = max \# of iterations (timesteps) to run barrier\sphinxhyphen{}climbing NEB
  Nevery = print NEB statistics every this many timesteps
\sphinxstyleemphasis{neb\_style} value = \sphinxstyleemphasis{quickmin} or \sphinxstyleemphasis{fire}
\sphinxstyleemphasis{neb\_step} value = dtneb
  dtneb = timestep for NEB damped dynamics minimization
\sphinxstyleemphasis{neb\_log} value = file where NEB statistics are printed
\end{sphinxalltt}

\end{itemize}


\subsection{Examples}
\label{\detokenize{tad:examples}}
\begin{sphinxVerbatim}[commandchars=\\\{\}]
\PYG{k}{tad}\PYG{+w}{ }\PYG{l+m}{2000}\PYG{+w}{ }\PYG{l+m}{50}\PYG{+w}{ }\PYG{l+m}{1800}\PYG{+w}{ }\PYG{l+m}{2300}\PYG{+w}{ }\PYG{l+m}{0.01}\PYG{+w}{ }\PYG{l+m}{0.01}\PYG{+w}{ }\PYG{n}{event}
\PYG{k}{tad}\PYG{+w}{ }\PYG{l+m}{2000}\PYG{+w}{ }\PYG{l+m}{50}\PYG{+w}{ }\PYG{l+m}{1800}\PYG{+w}{ }\PYG{l+m}{2300}\PYG{+w}{ }\PYG{l+m}{0.01}\PYG{+w}{ }\PYG{l+m}{0.01}\PYG{+w}{ }\PYG{n}{event}\PYG{+w}{ }\PYG{o}{\PYGZam{}}
\PYG{+w}{    }\PYG{n}{min}\PYG{+w}{ }\PYG{l+m}{1e\PYGZhy{}05}\PYG{+w}{ }\PYG{l+m}{1e\PYGZhy{}05}\PYG{+w}{ }\PYG{l+m}{100}\PYG{+w}{ }\PYG{l+m}{100}\PYG{+w}{ }\PYG{o}{\PYGZam{}}
\PYG{+w}{    }\PYG{n}{neb}\PYG{+w}{ }\PYG{l+m}{0.0}\PYG{+w}{ }\PYG{l+m}{0.01}\PYG{+w}{ }\PYG{l+m}{200}\PYG{+w}{ }\PYG{l+m}{200}\PYG{+w}{ }\PYG{l+m}{20}\PYG{+w}{ }\PYG{o}{\PYGZam{}}
\PYG{+w}{    }\PYG{n}{min\PYGZus{}style}\PYG{+w}{ }\PYG{n}{cg}\PYG{+w}{ }\PYG{o}{\PYGZam{}}
\PYG{+w}{    }\PYG{n}{neb\PYGZus{}style}\PYG{+w}{ }\PYG{n}{fire}\PYG{+w}{ }\PYG{o}{\PYGZam{}}
\PYG{+w}{    }\PYG{n}{neb\PYGZus{}log}\PYG{+w}{ }\PYG{n}{log.neb}
\end{sphinxVerbatim}


\subsection{Description}
\label{\detokenize{tad:description}}
\sphinxAtStartPar
Run a temperature accelerated dynamics (TAD) simulation. This method
requires two or more partitions to perform NEB transition state
searches.

\sphinxAtStartPar
TAD is described in {\hyperref[\detokenize{tad:voter2000}]{\sphinxcrossref{\DUrole{std,std-ref}{this paper}}}} by Art Voter.  It is
a method that uses accelerated dynamics at an elevated temperature to
generate results at a specified lower temperature.  A good overview of
accelerated dynamics methods (AMD) for such systems is given in
{\hyperref[\detokenize{tad:voter2002}]{\sphinxcrossref{\DUrole{std,std-ref}{this review paper}}}} from the same group.  To quote
from the review paper: “The dynamical evolution is characterized by
vibrational excursions within a potential basin, punctuated by
occasional transitions between basins.  The transition probability is
characterized by p(t) = k*exp(\sphinxhyphen{}kt) where k is the rate constant.”

\sphinxAtStartPar
TAD is a suitable AMD method for infrequent\sphinxhyphen{}event systems, where in
addition, the transition kinetics are well\sphinxhyphen{}approximated by harmonic
transition state theory (hTST).  In hTST, the temperature dependence
of transition rates follows the Arrhenius relation.  As a consequence
a set of event times generated in a high\sphinxhyphen{}temperature simulation can be
mapped to a set of much longer estimated times in the low\sphinxhyphen{}temperature
system. However, because this mapping involves the energy barrier of
the transition event, which is different for each event, the first
event at the high temperature may not be the earliest event at the low
temperature. TAD handles this by first generating a set of possible
events from the current basin. After each event, the simulation is
reflected backwards into the current basin.  This is repeated until
the stopping criterion is satisfied, at which point the event with the
earliest low\sphinxhyphen{}temperature occurrence time is selected.  The stopping
criterion is that the confidence measure be greater than 1\sphinxhyphen{}\sphinxstyleemphasis{delta}. The confidence measure is the probability that no earlier
low\sphinxhyphen{}temperature event will occur at some later time in the
high\sphinxhyphen{}temperature simulation.  hTST provides an lower bound for this
probability, based on the user\sphinxhyphen{}specified minimum pre\sphinxhyphen{}exponential
factor (reciprocal of \sphinxstyleemphasis{tmax}).

\sphinxAtStartPar
In order to estimate the energy barrier for each event, the TAD method
invokes the {\hyperref[\detokenize{neb::doc}]{\sphinxcrossref{\DUrole{doc}{NEB}}}} method. Each NEB replica runs on a
partition of processors. The current NEB implementation in LAMMPS
restricts you to having exactly one processor per replica. For more
information, see the documentation for the {\hyperref[\detokenize{neb::doc}]{\sphinxcrossref{\DUrole{doc}{neb}}}} command.  In
the current LAMMPS implementation of TAD, all the non\sphinxhyphen{}NEB TAD
operations are performed on the first partition, while the other
partitions remain idle. See the {\hyperref[\detokenize{Howto_replica::doc}]{\sphinxcrossref{\DUrole{doc}{Howto replica}}}} doc
page for further discussion of multi\sphinxhyphen{}replica simulations.

\sphinxAtStartPar
A TAD run has several stages, which are repeated each time an event is
performed.  The logic for a TAD run is as follows:

\begin{sphinxVerbatim}[commandchars=\\\{\}]
while (time remains):
  while (time \PYGZlt{} tstop):
    until (event occurs):
      run dynamics for t\PYGZus{}event steps
      quench
    run neb calculation using all replicas
    compute tlo from energy barrier
    update earliest event
    update tstop
    reflect back into current basin
  execute earliest event
\end{sphinxVerbatim}

\sphinxAtStartPar
Before this outer loop begins, the initial potential energy basin is
identified by quenching (an energy minimization, see below) the
initial state and storing the resulting coordinates for reference.

\sphinxAtStartPar
Inside the inner loop, dynamics is run continuously according to
whatever integrator has been specified by the user, stopping every
\sphinxstyleemphasis{t\_event} steps to check if a transition event has occurred.  This
check is performed by quenching the system and comparing the resulting
atom coordinates to the coordinates from the previous basin.

\sphinxAtStartPar
A quench is an energy minimization and is performed by whichever
algorithm has been defined by the {\hyperref[\detokenize{min_style::doc}]{\sphinxcrossref{\DUrole{doc}{min\_style}}}} command;
its default is the CG minimizer.  The tolerances and limits for each
quench can be set by the \sphinxstyleemphasis{min} keyword.  Note that typically, you do
not need to perform a highly\sphinxhyphen{}converged minimization to detect a
transition event.

\sphinxAtStartPar
The event check is performed by a compute with the specified
\sphinxstyleemphasis{compute\sphinxhyphen{}ID}.  Currently there is only one compute that works with the
TAD command, which is the {\hyperref[\detokenize{compute_event_displace::doc}]{\sphinxcrossref{\DUrole{doc}{compute event/displace}}}} command.  Other
event\sphinxhyphen{}checking computes may be added.  {\hyperref[\detokenize{compute_event_displace::doc}]{\sphinxcrossref{\DUrole{doc}{Compute event/displace}}}} checks whether any atom in
the compute group has moved further than a specified threshold
distance.  If so, an “event” has occurred.

\sphinxAtStartPar
The NEB calculation is similar to that invoked by the {\hyperref[\detokenize{neb::doc}]{\sphinxcrossref{\DUrole{doc}{neb}}}}
command, except that the final state is generated internally, instead
of being read in from a file.  The style of minimization performed by
NEB is determined by the \sphinxstyleemphasis{neb\_style} keyword and must be a damped
dynamics minimizer.  The tolerances and limits for each NEB
calculation can be set by the \sphinxstyleemphasis{neb} keyword.  As discussed on the
{\hyperref[\detokenize{neb::doc}]{\sphinxcrossref{\DUrole{doc}{neb}}}}, it is often advantageous to use a larger timestep for
NEB than for normal dynamics.  Since the size of the timestep set by
the {\hyperref[\detokenize{timestep::doc}]{\sphinxcrossref{\DUrole{doc}{timestep}}}} command is used by TAD for performing
dynamics, there is a \sphinxstyleemphasis{neb\_step} keyword which can be used to set a
larger timestep for each NEB calculation if desired.


\bigskip\hrule\bigskip


\sphinxAtStartPar
A key aspect of the TAD method is setting the stopping criterion
appropriately.  If this criterion is too conservative, then many
events must be generated before one is finally executed.  Conversely,
if this criterion is too aggressive, high\sphinxhyphen{}entropy high\sphinxhyphen{}barrier events
will be over\sphinxhyphen{}sampled, while low\sphinxhyphen{}entropy low\sphinxhyphen{}barrier events will be
under\sphinxhyphen{}sampled. If the lowest pre\sphinxhyphen{}exponential factor is known fairly
accurately, then it can be used to estimate \sphinxstyleemphasis{tmax}, and the value of
\sphinxstyleemphasis{delta} can be set to the desired confidence level e.g. \sphinxstyleemphasis{delta} = 0.05
corresponds to 95\% confidence. However, for systems where the dynamics
are not well characterized (the most common case), it will be
necessary to experiment with the values of \sphinxstyleemphasis{delta} and \sphinxstyleemphasis{tmax} to get a
good trade\sphinxhyphen{}off between accuracy and performance.

\sphinxAtStartPar
A second key aspect is the choice of \sphinxstyleemphasis{t\_hi}. A larger value greatly
increases the rate at which new events are generated.  However, too
large a value introduces errors due to anharmonicity (not accounted
for within hTST). Once again, for any given system, experimentation is
necessary to determine the best value of \sphinxstyleemphasis{t\_hi}.


\bigskip\hrule\bigskip


\sphinxAtStartPar
Five kinds of output can be generated during a TAD run: event
statistics, NEB statistics, thermodynamic output by each replica, dump
files, and restart files.

\sphinxAtStartPar
Event statistics are printed to the screen and master log.lammps file
each time an event is executed. The quantities are the timestep, CPU
time, global event number \sphinxstyleemphasis{N}, local event number \sphinxstyleemphasis{M}, event status,
energy barrier, time margin, \sphinxstyleemphasis{t\_lo} and \sphinxstyleemphasis{delt\_lo}.  The timestep is
the usual LAMMPS timestep, which corresponds to the high\sphinxhyphen{}temperature
time at which the event was detected, in units of timestep.  The CPU
time is the total processor time since the start of the TAD run.  The
global event number \sphinxstyleemphasis{N} is a counter that increments with each
executed event. The local event number \sphinxstyleemphasis{M} is a counter that resets to
zero upon entering each new basin.  The event status is \sphinxstyleemphasis{E} when an
event is executed, and is \sphinxstyleemphasis{D} for an event that is detected, while
\sphinxstyleemphasis{DF} is for a detected event that is also the earliest (first) event
at the low temperature.

\sphinxAtStartPar
The time margin is the ratio of the high temperature time in the
current basin to the stopping time. This last number can be used to
judge whether the stopping time is too short or too long (see above).

\sphinxAtStartPar
\sphinxstyleemphasis{t\_lo} is the low\sphinxhyphen{}temperature event time when the current basin was
entered, in units of timestep.  del*t\_lo* is the time of each detected
event, measured relative to \sphinxstyleemphasis{t\_lo}.  \sphinxstyleemphasis{delt\_lo} is equal to the
high\sphinxhyphen{}temperature time since entering the current basin, scaled by an
exponential factor that depends on the hi/lo temperature ratio and the
energy barrier for that event.

\sphinxAtStartPar
On lines for executed events, with status \sphinxstyleemphasis{E}, the global event number
is incremented by one,
the local event number and time margin are reset to zero,
while the global event number, energy barrier, and
\sphinxstyleemphasis{delt\_lo} match the last event with status \sphinxstyleemphasis{DF}
in the immediately preceding block of detected events.
The low\sphinxhyphen{}temperature event time \sphinxstyleemphasis{t\_lo} is incremented by \sphinxstyleemphasis{delt\_lo}.

\sphinxAtStartPar
NEB statistics are written to the file specified by the \sphinxstyleemphasis{neb\_log}
keyword. If the keyword value is “none”, then no NEB statistics are
printed out. The statistics are written every \sphinxstyleemphasis{Nevery} timesteps.  See
the {\hyperref[\detokenize{neb::doc}]{\sphinxcrossref{\DUrole{doc}{neb}}}} command for a full description of the NEB
statistics. When invoked from TAD, NEB statistics are never printed to
the screen.

\sphinxAtStartPar
Because the NEB calculation must run on multiple partitions, LAMMPS
produces additional screen and log files for each partition,
e.g. log.lammps.0, log.lammps.1, etc. For the TAD command, these
contain the thermodynamic output of each NEB replica. In addition, the
log file for the first partition, log.lammps.0, will contain
thermodynamic output from short runs and minimizations corresponding
to the dynamics and quench operations, as well as a line for each new
detected event, as described above.

\sphinxAtStartPar
After the TAD command completes, timing statistics for the TAD run are
printed in each replica’s log file, giving a breakdown of how much CPU
time was spent in each stage (NEB, dynamics, quenching, etc).

\sphinxAtStartPar
Any {\hyperref[\detokenize{dump::doc}]{\sphinxcrossref{\DUrole{doc}{dump files}}}} defined in the input script will be written
to during a TAD run at timesteps when an event is executed.  This
means the requested dump frequency in the {\hyperref[\detokenize{dump::doc}]{\sphinxcrossref{\DUrole{doc}{dump}}}} command
is ignored.  There will be one dump file (per dump command) created
for all partitions.  The atom coordinates of the dump snapshot are
those of the minimum energy configuration resulting from quenching
following the executed event.  The timesteps written into the dump
files correspond to the timestep at which the event occurred and NOT
the clock.  A dump snapshot corresponding to the initial minimum state
used for event detection is written to the dump file at the beginning
of each TAD run.

\sphinxAtStartPar
If the {\hyperref[\detokenize{restart::doc}]{\sphinxcrossref{\DUrole{doc}{restart}}}} command is used, a single restart file
for all the partitions is generated, which allows a TAD run to be
continued by a new input script in the usual manner.  The restart file
is generated after an event is executed. The restart file contains a
snapshot of the system in the new quenched state, including the event
number and the low\sphinxhyphen{}temperature time.  The restart frequency specified
in the {\hyperref[\detokenize{restart::doc}]{\sphinxcrossref{\DUrole{doc}{restart}}}} command is interpreted differently when
performing a TAD run.  It does not mean the timestep interval between
restart files.  Instead it means an event interval for executed
events.  Thus a frequency of 1 means write a restart file every time
an event is executed.  A frequency of 10 means write a restart file
every 10th executed event.  When an input script reads a restart file
from a previous TAD run, the new script can be run on a different
number of replicas or processors.

\sphinxAtStartPar
Note that within a single state, the dynamics will typically
temporarily continue beyond the event that is ultimately chosen, until
the stopping criterion is satisfied.  When the event is eventually
executed, the timestep counter is reset to the value when the event
was detected. Similarly, after each quench and NEB minimization, the
timestep counter is reset to the value at the start of the
minimization. This means that the timesteps listed in the replica log
files do not always increase monotonically. However, the timestep
values printed to the master log file, dump files, and restart files
are always monotonically increasing.


\bigskip\hrule\bigskip



\subsection{Restrictions}
\label{\detokenize{tad:restrictions}}
\sphinxAtStartPar
This command can only be used if LAMMPS was built with the REPLICA
package.  See the {\hyperref[\detokenize{Build_package::doc}]{\sphinxcrossref{\DUrole{doc}{Build package}}}} doc
page for more info.

\sphinxAtStartPar
\sphinxstyleemphasis{N} setting must be integer multiple of \sphinxstyleemphasis{t\_event}.

\sphinxAtStartPar
Runs restarted from restart files written during a TAD run will only
produce identical results if the user\sphinxhyphen{}specified integrator supports
exact restarts. So {\hyperref[\detokenize{fix_nh::doc}]{\sphinxcrossref{\DUrole{doc}{fix nvt}}}} will produce an exact
restart, but {\hyperref[\detokenize{fix_langevin::doc}]{\sphinxcrossref{\DUrole{doc}{fix langevin}}}} will not.

\sphinxAtStartPar
This command cannot be used when any fixes are defined that keep track
of elapsed time to perform time\sphinxhyphen{}dependent operations.  Examples
include the “ave” fixes such as {\hyperref[\detokenize{fix_ave_chunk::doc}]{\sphinxcrossref{\DUrole{doc}{fix ave/chunk}}}}.
Also {\hyperref[\detokenize{fix_dt_reset::doc}]{\sphinxcrossref{\DUrole{doc}{fix dt/reset}}}} and {\hyperref[\detokenize{fix_deposit::doc}]{\sphinxcrossref{\DUrole{doc}{fix deposit}}}}.


\subsection{Related commands}
\label{\detokenize{tad:related-commands}}
\sphinxAtStartPar
{\hyperref[\detokenize{compute_event_displace::doc}]{\sphinxcrossref{\DUrole{doc}{compute event/displace}}}},
{\hyperref[\detokenize{min_modify::doc}]{\sphinxcrossref{\DUrole{doc}{min\_modify}}}}, {\hyperref[\detokenize{min_style::doc}]{\sphinxcrossref{\DUrole{doc}{min\_style}}}},
{\hyperref[\detokenize{run_style::doc}]{\sphinxcrossref{\DUrole{doc}{run\_style}}}}, {\hyperref[\detokenize{minimize::doc}]{\sphinxcrossref{\DUrole{doc}{minimize}}}},
{\hyperref[\detokenize{temper::doc}]{\sphinxcrossref{\DUrole{doc}{temper}}}}, {\hyperref[\detokenize{neb::doc}]{\sphinxcrossref{\DUrole{doc}{neb}}}},
{\hyperref[\detokenize{prd::doc}]{\sphinxcrossref{\DUrole{doc}{prd}}}}


\subsection{Default}
\label{\detokenize{tad:default}}
\sphinxAtStartPar
The option defaults are \sphinxstyleemphasis{min} = 0.1 0.1 40 50, \sphinxstyleemphasis{neb} = 0.01 100 100
10, \sphinxstyleemphasis{neb\_style} = \sphinxstyleemphasis{quickmin}, \sphinxstyleemphasis{neb\_step} = the same timestep set by
the {\hyperref[\detokenize{timestep::doc}]{\sphinxcrossref{\DUrole{doc}{timestep}}}} command, and \sphinxstyleemphasis{neb\_log} = “none”.


\bigskip\hrule\bigskip

\phantomsection\label{\detokenize{tad:voter2000}}
\sphinxAtStartPar
\sphinxstylestrong{(Voter2000)} Sorensen and Voter, J Chem Phys, 112, 9599 (2000)

\phantomsection\label{\detokenize{tad:voter2002}}
\sphinxAtStartPar
\sphinxstylestrong{(Voter2002)} Voter, Montalenti, Germann, Annual Review of Materials
Research 32, 321 (2002).

\sphinxstepscope

\index{temper@\spxentry{temper}}\ignorespaces 

\section{temper command}
\label{\detokenize{temper:temper-command}}\label{\detokenize{temper:index-0}}\label{\detokenize{temper::doc}}

\subsection{Syntax}
\label{\detokenize{temper:syntax}}
\begin{sphinxVerbatim}[commandchars=\\\{\}]
\PYG{k}{temper}\PYG{+w}{ }\PYG{n}{N}\PYG{+w}{ }\PYG{n}{M}\PYG{+w}{ }\PYG{n}{temp}\PYG{+w}{ }\PYG{n}{fix}\PYG{o}{\PYGZhy{}}\PYG{n}{ID}\PYG{+w}{ }\PYG{n}{seed1}\PYG{+w}{ }\PYG{n}{seed2}\PYG{+w}{ }\PYG{n}{index}
\end{sphinxVerbatim}
\begin{itemize}
\item {} 
\sphinxAtStartPar
N = total \# of timesteps to run

\item {} 
\sphinxAtStartPar
M = attempt a tempering swap every this many steps

\item {} 
\sphinxAtStartPar
temp = initial temperature for this ensemble

\item {} 
\sphinxAtStartPar
fix\sphinxhyphen{}ID = ID of the fix that will control temperature during the run

\item {} 
\sphinxAtStartPar
seed1 = random \# seed used to decide on adjacent temperature to partner with

\item {} 
\sphinxAtStartPar
seed2 = random \# seed for Boltzmann factor in Metropolis swap

\item {} 
\sphinxAtStartPar
index = which temperature (0 to N\sphinxhyphen{}1) I am simulating (optional)

\end{itemize}


\subsection{Examples}
\label{\detokenize{temper:examples}}
\begin{sphinxVerbatim}[commandchars=\\\{\}]
\PYG{k}{temper}\PYG{+w}{ }\PYG{l+m}{100000}\PYG{+w}{ }\PYG{l+m}{100}\PYG{+w}{ }\PYG{n+nv}{\PYGZdl{}t}\PYG{+w}{ }\PYG{n}{tempfix}\PYG{+w}{ }\PYG{l+m}{0}\PYG{+w}{ }\PYG{l+m}{58728}
\PYG{k}{temper}\PYG{+w}{ }\PYG{l+m}{40000}\PYG{+w}{ }\PYG{l+m}{100}\PYG{+w}{ }\PYG{n+nv}{\PYGZdl{}t}\PYG{+w}{ }\PYG{n}{tempfix}\PYG{+w}{ }\PYG{l+m}{0}\PYG{+w}{ }\PYG{l+m}{32285}\PYG{+w}{ }\PYG{n+nv}{\PYGZdl{}w}
\end{sphinxVerbatim}


\subsection{Description}
\label{\detokenize{temper:description}}
\sphinxAtStartPar
Run a parallel tempering or replica exchange simulation using multiple
replicas (ensembles) of a system.  Two or more replicas must be used.

\sphinxAtStartPar
Each replica runs on a partition of one or more processors.  Processor
partitions are defined at run\sphinxhyphen{}time using the {\hyperref[\detokenize{Run_options::doc}]{\sphinxcrossref{\DUrole{doc}{\sphinxhyphen{}partition command\sphinxhyphen{}line switch}}}}.  Note that if you have MPI installed, you
can run a multi\sphinxhyphen{}replica simulation with more replicas (partitions)
than you have physical processors, e.g you can run a 10\sphinxhyphen{}replica
simulation on one or two processors.  You will simply not get the
performance speed\sphinxhyphen{}up you would see with one or more physical
processors per replica.  See the {\hyperref[\detokenize{Howto_replica::doc}]{\sphinxcrossref{\DUrole{doc}{Howto replica}}}}
doc page for further discussion.

\sphinxAtStartPar
Each replica’s temperature is controlled at a different value by a fix
with \sphinxstyleemphasis{fix\sphinxhyphen{}ID} that controls temperature. Most thermostat fix styles
(with and without included time integration) are supported. The command
will print an error message and abort, if the chosen fix is unsupported.
The desired temperature is specified by \sphinxstyleemphasis{temp}, which is typically a
variable previously set in the input script, so that each partition is
assigned a different temperature.  See the {\hyperref[\detokenize{variable::doc}]{\sphinxcrossref{\DUrole{doc}{variable}}}}
command for more details.  For example:

\begin{sphinxVerbatim}[commandchars=\\\{\}]
\PYG{k}{variable }\PYG{n+nv+nvIdentifier}{t}\PYG{+w}{ }\PYG{n}{world}\PYG{+w}{ }\PYG{l+m}{300.0}\PYG{+w}{ }\PYG{l+m}{310.0}\PYG{+w}{ }\PYG{l+m}{320.0}\PYG{+w}{ }\PYG{l+m}{330.0}
\PYG{k}{fix }\PYG{n+nv+nvIdentifier}{myfix}\PYG{+w}{ }\PYG{n+nv+nvIdentifier}{all}\PYG{+w}{ }\PYG{n}{nvt}\PYG{+w}{ }\PYG{n}{temp}\PYG{+w}{ }\PYG{n+nv}{\PYGZdl{}t}\PYG{+w}{ }\PYG{n+nv}{\PYGZdl{}t}\PYG{+w}{ }\PYG{l+m}{100.0}
\PYG{k}{temper}\PYG{+w}{ }\PYG{l+m}{100000}\PYG{+w}{ }\PYG{l+m}{100}\PYG{+w}{ }\PYG{n+nv}{\PYGZdl{}t}\PYG{+w}{ }\PYG{n}{myfix}\PYG{+w}{ }\PYG{l+m}{3847}\PYG{+w}{ }\PYG{l+m}{58382}
\end{sphinxVerbatim}

\sphinxAtStartPar
would define 4 temperatures, and assign one of them to the thermostat
used by each replica, and to the temper command.

\sphinxAtStartPar
As the tempering simulation runs for \sphinxstyleemphasis{N} timesteps, a temperature swap
between adjacent ensembles will be attempted every \sphinxstyleemphasis{M} timesteps.  If
\sphinxstyleemphasis{seed1} is 0, then the swap attempts will alternate between odd and
even pairings.  If \sphinxstyleemphasis{seed1} is non\sphinxhyphen{}zero then it is used as a seed in a
random number generator to randomly choose an odd or even pairing each
time.  Each attempted swap of temperatures is either accepted or
rejected based on a Boltzmann\sphinxhyphen{}weighted Metropolis criterion which uses
\sphinxstyleemphasis{seed2} in the random number generator.

\sphinxAtStartPar
As a tempering run proceeds, multiple log files and screen output
files are created, one per replica.  By default these files are named
log.lammps.M and screen.M where M is the replica number from 0 to N\sphinxhyphen{}1,
with N = \# of replicas.  See the {\hyperref[\detokenize{Run_options::doc}]{\sphinxcrossref{\DUrole{doc}{\sphinxhyphen{}log and \sphinxhyphen{}screen command\sphinxhyphen{}line swiches}}}} for info on how to change these names.

\sphinxAtStartPar
The main screen and log file (log.lammps) will list information about
which temperature is assigned to each replica at each thermodynamic
output timestep.  E.g. for a simulation with 16 replicas:

\begin{sphinxVerbatim}[commandchars=\\\{\}]
Running on 16 partitions of processors
Step T0 T1 T2 T3 T4 T5 T6 T7 T8 T9 T10 T11 T12 T13 T14 T15
0    0 1 2 3 4 5 6 7 8 9 10 11 12 13 14 15
500  1 0 3 2 5 4 6 7 8 9 10 11 12 13 14 15
1000 2 0 4 1 5 3 6 7 8 9 10 11 12 14 13 15
1500 2 1 4 0 5 3 6 7 9 8 10 11 12 14 13 15
2000 2 1 3 0 6 4 5 7 10 8 9 11 12 14 13 15
2500 2 1 3 0 6 4 5 7 11 8 9 10 12 14 13 15
...
\end{sphinxVerbatim}

\sphinxAtStartPar
The column headings T0 to TN\sphinxhyphen{}1 mean which temperature is currently
assigned to the replica 0 to N\sphinxhyphen{}1.  Thus the columns represent replicas
and the value in each column is its temperature (also numbered 0 to
N\sphinxhyphen{}1).  For example, a 0 in the fourth column (column T3, step 2500) means
that the fourth replica is assigned temperature 0, i.e. the lowest
temperature.  You can verify this time sequence of temperature
assignments for the Nth replica by comparing the Nth column of screen
output to the thermodynamic data in the corresponding log.lammps.N or
screen.N files as time proceeds.

\sphinxAtStartPar
You can have each replica create its own dump file in the following
manner:

\begin{sphinxVerbatim}[commandchars=\\\{\}]
\PYG{k}{variable }\PYG{n+nv+nvIdentifier}{rep}\PYG{+w}{ }\PYG{n}{world}\PYG{+w}{ }\PYG{l+m}{0}\PYG{+w}{ }\PYG{l+m}{1}\PYG{+w}{ }\PYG{l+m}{2}\PYG{+w}{ }\PYG{l+m}{3}\PYG{+w}{ }\PYG{l+m}{4}\PYG{+w}{ }\PYG{l+m}{5}\PYG{+w}{ }\PYG{l+m}{6}\PYG{+w}{ }\PYG{l+m}{7}
\PYG{k}{dump }\PYG{n+nv+nvIdentifier}{1}\PYG{+w}{ }\PYG{n+nv+nvIdentifier}{all}\PYG{+w}{ }\PYG{n}{atom}\PYG{+w}{ }\PYG{l+m}{1000}\PYG{+w}{ }\PYG{n}{dump.temper.}\PYG{n+nv}{\PYGZdl{}\PYGZob{}}\PYG{n+nv}{rep}\PYG{n+nv}{\PYGZcb{}}
\end{sphinxVerbatim}

\begin{sphinxadmonition}{note}{Note:}
\sphinxAtStartPar
Each replica’s dump file will contain a continuous trajectory
for its atoms where the temperature varies over time as swaps take
place involving that replica.  If you want a series of dump files,
each with snapshots (from all replicas) that are all at a single
temperature, then you will need to post\sphinxhyphen{}process the dump files using
the information from the log.lammps file.  E.g. you could produce one
dump file with snapshots at 300K (from all replicas), another with
snapshots at 310K, etc.  Note that these new dump files will not
contain “continuous trajectories” for individual atoms, because two
successive snapshots (in time) may be from different replicas. The
reorder\_remd\_traj python script can do the reordering for you
(and additionally also calculated configurational log\sphinxhyphen{}weights of
trajectory snapshots in the canonical ensemble). The script can be found
in the tools/replica directory while instructions on how to use it is
available in doc/Tools (in brief) and as a README file in tools/replica
(in detail).
\end{sphinxadmonition}

\sphinxAtStartPar
The last argument \sphinxstyleemphasis{index} in the temper command is optional and is
used when restarting a tempering run from a set of restart files (one
for each replica) which had previously swapped to new temperatures.
The \sphinxstyleemphasis{index} value (from 0 to N\sphinxhyphen{}1, where N is the \# of replicas)
identifies which temperature the replica was simulating on the
timestep the restart files were written.  Obviously, this argument
must be a variable so that each partition has the correct value.  Set
the variable to the \sphinxstyleemphasis{N} values listed in the log file for the previous
run for the replica temperatures at that timestep.  For example if the
log file listed the following for a simulation with 5 replicas:

\begin{sphinxVerbatim}[commandchars=\\\{\}]
500000 2 4 0 1 3
\end{sphinxVerbatim}

\sphinxAtStartPar
then a setting of

\begin{sphinxVerbatim}[commandchars=\\\{\}]
\PYG{k}{variable }\PYG{n+nv+nvIdentifier}{w}\PYG{+w}{ }\PYG{n}{world}\PYG{+w}{ }\PYG{l+m}{2}\PYG{+w}{ }\PYG{l+m}{4}\PYG{+w}{ }\PYG{l+m}{0}\PYG{+w}{ }\PYG{l+m}{1}\PYG{+w}{ }\PYG{l+m}{3}
\end{sphinxVerbatim}

\sphinxAtStartPar
would be used to restart the run with a tempering command like the
example above with \$w as the last argument.


\bigskip\hrule\bigskip



\subsection{Restrictions}
\label{\detokenize{temper:restrictions}}
\sphinxAtStartPar
This command can only be used if LAMMPS was built with the REPLICA
package.  See the {\hyperref[\detokenize{Build_package::doc}]{\sphinxcrossref{\DUrole{doc}{Build package}}}} doc
page for more info.


\subsection{Related commands}
\label{\detokenize{temper:related-commands}}
\sphinxAtStartPar
{\hyperref[\detokenize{variable::doc}]{\sphinxcrossref{\DUrole{doc}{variable}}}}, {\hyperref[\detokenize{prd::doc}]{\sphinxcrossref{\DUrole{doc}{prd}}}}, {\hyperref[\detokenize{neb::doc}]{\sphinxcrossref{\DUrole{doc}{neb}}}}


\subsection{Default}
\label{\detokenize{temper:default}}
\sphinxAtStartPar
none

\sphinxstepscope

\index{temper/grem@\spxentry{temper/grem}}\ignorespaces 

\section{temper/grem command}
\label{\detokenize{temper_grem:temper-grem-command}}\label{\detokenize{temper_grem:index-0}}\label{\detokenize{temper_grem::doc}}

\subsection{Syntax}
\label{\detokenize{temper_grem:syntax}}
\begin{sphinxVerbatim}[commandchars=\\\{\}]
\PYG{k}{temper/grem}\PYG{+w}{ }\PYG{n}{N}\PYG{+w}{ }\PYG{n}{M}\PYG{+w}{ }\PYG{n}{lambda}\PYG{+w}{ }\PYG{n}{fix}\PYG{o}{\PYGZhy{}}\PYG{n}{ID}\PYG{+w}{ }\PYG{n}{thermostat}\PYG{o}{\PYGZhy{}}\PYG{n}{ID}\PYG{+w}{ }\PYG{n}{seed1}\PYG{+w}{ }\PYG{n}{seed2}\PYG{+w}{ }\PYG{n}{index}
\end{sphinxVerbatim}
\begin{itemize}
\item {} 
\sphinxAtStartPar
N = total \# of timesteps to run

\item {} 
\sphinxAtStartPar
M = attempt a tempering swap every this many steps

\item {} 
\sphinxAtStartPar
lambda = initial lambda for this ensemble

\item {} 
\sphinxAtStartPar
fix\sphinxhyphen{}ID = ID of \sphinxstyleemphasis{fix grem}

\item {} 
\sphinxAtStartPar
thermostat\sphinxhyphen{}ID = ID of the thermostat that controls kinetic temperature

\item {} 
\sphinxAtStartPar
seed1 = random \# seed used to decide on adjacent temperature to partner with

\item {} 
\sphinxAtStartPar
seed2 = random \# seed for Boltzmann factor in Metropolis swap

\item {} 
\sphinxAtStartPar
index = which temperature (0 to N\sphinxhyphen{}1) I am simulating (optional)

\end{itemize}


\subsection{Examples}
\label{\detokenize{temper_grem:examples}}
\begin{sphinxVerbatim}[commandchars=\\\{\}]
\PYG{k}{temper/grem}\PYG{+w}{ }\PYG{l+m}{100000}\PYG{+w}{ }\PYG{l+m}{1000}\PYG{+w}{ }\PYG{n+nv}{\PYGZdl{}\PYGZob{}}\PYG{n+nv}{lambda}\PYG{n+nv}{\PYGZcb{}}\PYG{+w}{ }\PYG{n}{fxgREM}\PYG{+w}{ }\PYG{n}{fxnvt}\PYG{+w}{ }\PYG{l+m}{0}\PYG{+w}{ }\PYG{l+m}{58728}
\PYG{k}{temper/grem}\PYG{+w}{ }\PYG{l+m}{40000}\PYG{+w}{ }\PYG{l+m}{100}\PYG{+w}{ }\PYG{n+nv}{\PYGZdl{}\PYGZob{}}\PYG{n+nv}{lambda}\PYG{n+nv}{\PYGZcb{}}\PYG{+w}{ }\PYG{n}{fxgREM}\PYG{+w}{ }\PYG{n}{fxnpt}\PYG{+w}{ }\PYG{l+m}{0}\PYG{+w}{ }\PYG{l+m}{32285}\PYG{+w}{ }\PYG{n+nv}{\PYGZdl{}\PYGZob{}}\PYG{n+nv}{walkers}\PYG{n+nv}{\PYGZcb{}}
\end{sphinxVerbatim}


\subsection{Description}
\label{\detokenize{temper_grem:description}}
\sphinxAtStartPar
Run a parallel tempering or replica exchange simulation in LAMMPS
partition mode using multiple generalized replicas (ensembles) of a
system defined by {\hyperref[\detokenize{fix_grem::doc}]{\sphinxcrossref{\DUrole{doc}{fix grem}}}}, which stands for the
generalized replica exchange method (gREM) originally developed by
{\hyperref[\detokenize{temper_grem:kimstraub}]{\sphinxcrossref{\DUrole{std,std-ref}{(Kim)}}}}.  It uses non\sphinxhyphen{}Boltzmann ensembles to sample over first
order phase transitions. The is done by defining replicas with an
enthalpy dependent effective temperature

\sphinxAtStartPar
Two or more replicas must be used.  See the {\hyperref[\detokenize{temper::doc}]{\sphinxcrossref{\DUrole{doc}{temper}}}}
command for an explanation of how to run replicas on multiple
partitions of one or more processors.

\sphinxAtStartPar
This command is a modification of the {\hyperref[\detokenize{temper::doc}]{\sphinxcrossref{\DUrole{doc}{temper}}}} command and
has the same dependencies, restraints, and input variables which are
discussed there in greater detail.

\sphinxAtStartPar
Instead of temperature, this command performs replica exchanges in
lambda as per the generalized ensemble enforced by {\hyperref[\detokenize{fix_grem::doc}]{\sphinxcrossref{\DUrole{doc}{fix grem}}}}.  The desired lambda is specified by \sphinxstyleemphasis{lambda}, which is
typically a variable previously set in the input script, so that each
partition is assigned a different temperature.  See the {\hyperref[\detokenize{variable::doc}]{\sphinxcrossref{\DUrole{doc}{variable}}}} command for more details.  For example:

\begin{sphinxVerbatim}[commandchars=\\\{\}]
\PYG{k}{variable }\PYG{n+nv+nvIdentifier}{lambda}\PYG{+w}{ }\PYG{n}{world}\PYG{+w}{ }\PYG{l+m}{400}\PYG{+w}{ }\PYG{l+m}{420}\PYG{+w}{ }\PYG{l+m}{440}\PYG{+w}{ }\PYG{l+m}{460}
\PYG{k}{fix }\PYG{n+nv+nvIdentifier}{fxnvt}\PYG{+w}{ }\PYG{n+nv+nvIdentifier}{all}\PYG{+w}{ }\PYG{n}{nvt}\PYG{+w}{ }\PYG{n}{temp}\PYG{+w}{ }\PYG{l+m}{300.0}\PYG{+w}{ }\PYG{l+m}{300.0}\PYG{+w}{ }\PYG{l+m}{100.0}
\PYG{k}{fix }\PYG{n+nv+nvIdentifier}{fxgREM}\PYG{+w}{ }\PYG{n+nv+nvIdentifier}{all}\PYG{+w}{ }\PYG{n}{grem}\PYG{+w}{ }\PYG{n+nv}{\PYGZdl{}\PYGZob{}}\PYG{n+nv}{lambda}\PYG{n+nv}{\PYGZcb{}}\PYG{+w}{ }\PYG{o}{\PYGZhy{}}\PYG{l+m}{0.05}\PYG{+w}{ }\PYG{o}{\PYGZhy{}}\PYG{l+m}{50000}\PYG{+w}{ }\PYG{n}{fxnvt}
\PYG{k}{temper/grem}\PYG{+w}{ }\PYG{l+m}{100000}\PYG{+w}{ }\PYG{l+m}{100}\PYG{+w}{ }\PYG{n+nv}{\PYGZdl{}\PYGZob{}}\PYG{n+nv}{lambda}\PYG{n+nv}{\PYGZcb{}}\PYG{+w}{ }\PYG{n}{fxgREM}\PYG{+w}{ }\PYG{n}{fxnvt}\PYG{+w}{ }\PYG{l+m}{3847}\PYG{+w}{ }\PYG{l+m}{58382}
\end{sphinxVerbatim}

\sphinxAtStartPar
would define 4 lambdas with constant kinetic temperature but unique
generalized temperature, and assign one of them to {\hyperref[\detokenize{fix_grem::doc}]{\sphinxcrossref{\DUrole{doc}{fix grem}}}} used by each replica, and to the grem command.

\sphinxAtStartPar
As the gREM simulation runs for \sphinxstyleemphasis{N} timesteps, a swap between adjacent
ensembles will be attempted every \sphinxstyleemphasis{M} timesteps.  If \sphinxstyleemphasis{seed1} is 0,
then the swap attempts will alternate between odd and even pairings.
If \sphinxstyleemphasis{seed1} is non\sphinxhyphen{}zero then it is used as a seed in a random number
generator to randomly choose an odd or even pairing each time.  Each
attempted swap of temperatures is either accepted or rejected based on
a Metropolis criterion, derived for gREM by {\hyperref[\detokenize{Build_extras:kim}]{\sphinxcrossref{\DUrole{std,std-ref}{(Kim)}}}}, which uses
\sphinxstyleemphasis{seed2} in the random number generator.

\sphinxAtStartPar
File management works identical to the {\hyperref[\detokenize{temper::doc}]{\sphinxcrossref{\DUrole{doc}{temper}}}} command.
Dump files created by this fix contain continuous trajectories and
require post\sphinxhyphen{}processing to obtain per\sphinxhyphen{}replica information.

\sphinxAtStartPar
The last argument \sphinxstyleemphasis{index} in the grem command is optional and is used
when restarting a run from a set of restart files (one for each
replica) which had previously swapped to new lambda.  This is done
using a variable. For example if the log file listed the following for
a simulation with 5 replicas:

\begin{sphinxVerbatim}[commandchars=\\\{\}]
500000 2 4 0 1 3
\end{sphinxVerbatim}

\sphinxAtStartPar
then a setting of

\begin{sphinxVerbatim}[commandchars=\\\{\}]
\PYG{k}{variable }\PYG{n+nv+nvIdentifier}{walkers}\PYG{+w}{ }\PYG{n}{world}\PYG{+w}{ }\PYG{l+m}{2}\PYG{+w}{ }\PYG{l+m}{4}\PYG{+w}{ }\PYG{l+m}{0}\PYG{+w}{ }\PYG{l+m}{1}\PYG{+w}{ }\PYG{l+m}{3}
\end{sphinxVerbatim}

\sphinxAtStartPar
would be used to restart the run with a grem command like the example
above with \$\{walkers\} as the last argument. This functionality is
identical to {\hyperref[\detokenize{temper::doc}]{\sphinxcrossref{\DUrole{doc}{temper}}}}.


\bigskip\hrule\bigskip



\subsection{Restrictions}
\label{\detokenize{temper_grem:restrictions}}
\sphinxAtStartPar
This command can only be used if LAMMPS was built with the REPLICA
package.  See the {\hyperref[\detokenize{Build_package::doc}]{\sphinxcrossref{\DUrole{doc}{Build package}}}} doc
page for more info.

\sphinxAtStartPar
This command must be used with {\hyperref[\detokenize{fix_grem::doc}]{\sphinxcrossref{\DUrole{doc}{fix grem}}}}.


\subsection{Related commands}
\label{\detokenize{temper_grem:related-commands}}
\sphinxAtStartPar
{\hyperref[\detokenize{fix_grem::doc}]{\sphinxcrossref{\DUrole{doc}{fix grem}}}}, {\hyperref[\detokenize{temper::doc}]{\sphinxcrossref{\DUrole{doc}{temper}}}}, {\hyperref[\detokenize{variable::doc}]{\sphinxcrossref{\DUrole{doc}{variable}}}}


\subsection{Default}
\label{\detokenize{temper_grem:default}}
\sphinxAtStartPar
none

\phantomsection\label{\detokenize{temper_grem:kimstraub}}
\sphinxAtStartPar
\sphinxstylestrong{(Kim)} Kim, Keyes, Straub, J Chem Phys, 132, 224107 (2010).

\sphinxstepscope

\index{temper/npt@\spxentry{temper/npt}}\ignorespaces 

\section{temper/npt command}
\label{\detokenize{temper_npt:temper-npt-command}}\label{\detokenize{temper_npt:index-0}}\label{\detokenize{temper_npt::doc}}

\subsection{Syntax}
\label{\detokenize{temper_npt:syntax}}
\begin{sphinxVerbatim}[commandchars=\\\{\}]
\PYG{k}{temper/npt}\PYG{+w}{  }\PYG{n}{N}\PYG{+w}{ }\PYG{n}{M}\PYG{+w}{ }\PYG{n}{temp}\PYG{+w}{ }\PYG{n}{fix}\PYG{o}{\PYGZhy{}}\PYG{n}{ID}\PYG{+w}{ }\PYG{n}{seed1}\PYG{+w}{ }\PYG{n}{seed2}\PYG{+w}{ }\PYG{n}{pressure}\PYG{+w}{ }\PYG{n}{index}
\end{sphinxVerbatim}
\begin{itemize}
\item {} 
\sphinxAtStartPar
N = total \# of timesteps to run

\item {} 
\sphinxAtStartPar
M = attempt a tempering swap every this many steps

\item {} 
\sphinxAtStartPar
temp = initial temperature for this ensemble

\item {} 
\sphinxAtStartPar
fix\sphinxhyphen{}ID = ID of the fix that will control temperature and pressure during the run

\item {} 
\sphinxAtStartPar
seed1 = random \# seed used to decide on adjacent temperature to partner with

\item {} 
\sphinxAtStartPar
seed2 = random \# seed for Boltzmann factor in Metropolis swap

\item {} 
\sphinxAtStartPar
pressure = setpoint pressure for the ensemble

\item {} 
\sphinxAtStartPar
index = which temperature (0 to N\sphinxhyphen{}1) I am simulating (optional)

\end{itemize}


\subsection{Examples}
\label{\detokenize{temper_npt:examples}}
\begin{sphinxVerbatim}[commandchars=\\\{\}]
\PYG{k}{temper/npt}\PYG{+w}{ }\PYG{l+m}{100000}\PYG{+w}{ }\PYG{l+m}{100}\PYG{+w}{ }\PYG{n+nv}{\PYGZdl{}t}\PYG{+w}{ }\PYG{n}{nptfix}\PYG{+w}{ }\PYG{l+m}{0}\PYG{+w}{ }\PYG{l+m}{58728}\PYG{+w}{ }\PYG{l+m}{1}
\PYG{k}{temper/npt}\PYG{+w}{ }\PYG{l+m}{2500000}\PYG{+w}{ }\PYG{l+m}{1000}\PYG{+w}{ }\PYG{l+m}{300}\PYG{+w}{ }\PYG{n}{nptfix}\PYG{+w}{  }\PYG{l+m}{0}\PYG{+w}{ }\PYG{l+m}{32285}\PYG{+w}{ }\PYG{n+nv}{\PYGZdl{}p}
\PYG{k}{temper/npt}\PYG{+w}{ }\PYG{l+m}{5000000}\PYG{+w}{ }\PYG{l+m}{2000}\PYG{+w}{ }\PYG{n+nv}{\PYGZdl{}t}\PYG{+w}{ }\PYG{n}{nptfix}\PYG{+w}{ }\PYG{l+m}{0}\PYG{+w}{ }\PYG{l+m}{12523}\PYG{+w}{ }\PYG{l+m}{1}\PYG{+w}{ }\PYG{n+nv}{\PYGZdl{}w}
\end{sphinxVerbatim}


\subsection{Description}
\label{\detokenize{temper_npt:description}}
\sphinxAtStartPar
Run a parallel tempering or replica exchange simulation using multiple
replicas (ensembles) of a system in the isothermal\sphinxhyphen{}isobaric (NPT)
ensemble.  The command temper/npt works like {\hyperref[\detokenize{temper::doc}]{\sphinxcrossref{\DUrole{doc}{temper}}}} but
requires running replicas in the NPT ensemble instead of the canonical
(NVT) ensemble and allows for pressure to be set in the ensembles.
These multiple ensembles can run in parallel at different temperatures
or different pressures.  The acceptance criteria for temper/npt is
specific to the NPT ensemble and can be found in references
{\hyperref[\detokenize{temper_npt:okabe2}]{\sphinxcrossref{\DUrole{std,std-ref}{(Okabe)}}}} and {\hyperref[\detokenize{temper_npt:mori2}]{\sphinxcrossref{\DUrole{std,std-ref}{(Mori)}}}}.

\sphinxAtStartPar
Apart from the difference in acceptance criteria and the specification
of pressure, this command works much like the {\hyperref[\detokenize{temper::doc}]{\sphinxcrossref{\DUrole{doc}{temper}}}}
command. See the documentation on {\hyperref[\detokenize{temper::doc}]{\sphinxcrossref{\DUrole{doc}{temper}}}} for information
on how the parallel tempering is handled in general.


\bigskip\hrule\bigskip



\subsection{Restrictions}
\label{\detokenize{temper_npt:restrictions}}
\sphinxAtStartPar
This command can only be used if LAMMPS was built with the REPLICA
package.  See the {\hyperref[\detokenize{Build_package::doc}]{\sphinxcrossref{\DUrole{doc}{Build package}}}} page for more
info.

\sphinxAtStartPar
This command should be used with a fix that maintains the
isothermal\sphinxhyphen{}isobaric (NPT) ensemble.


\subsection{Related commands}
\label{\detokenize{temper_npt:related-commands}}
\sphinxAtStartPar
{\hyperref[\detokenize{temper::doc}]{\sphinxcrossref{\DUrole{doc}{temper}}}}, {\hyperref[\detokenize{variable::doc}]{\sphinxcrossref{\DUrole{doc}{variable}}}}, {\hyperref[\detokenize{fix_nh::doc}]{\sphinxcrossref{\DUrole{doc}{fix\_npt}}}}


\subsection{Default}
\label{\detokenize{temper_npt:default}}
\sphinxAtStartPar
none

\phantomsection\label{\detokenize{temper_npt:okabe2}}
\sphinxAtStartPar
\sphinxstylestrong{(Okabe)} T. Okabe, M. Kawata, Y. Okamoto, M. Masuhiro, Chem. Phys. Lett., 335, 435\sphinxhyphen{}439 (2001).

\phantomsection\label{\detokenize{temper_npt:mori2}}
\sphinxAtStartPar
\sphinxstylestrong{(Mori)} Y. Mori, Y. Okamoto, J. Phys. Soc. Jpn., 7, 074003 (2010).

\sphinxstepscope

\index{thermo@\spxentry{thermo}}\ignorespaces 

\section{thermo command}
\label{\detokenize{thermo:thermo-command}}\label{\detokenize{thermo:index-0}}\label{\detokenize{thermo::doc}}

\subsection{Syntax}
\label{\detokenize{thermo:syntax}}
\begin{sphinxVerbatim}[commandchars=\\\{\}]
\PYG{k}{thermo}\PYG{+w}{ }\PYG{n}{N}
\end{sphinxVerbatim}
\begin{itemize}
\item {} 
\sphinxAtStartPar
N = output thermodynamics every N timesteps

\item {} 
\sphinxAtStartPar
N can be a variable (see below)

\end{itemize}


\subsection{Examples}
\label{\detokenize{thermo:examples}}
\begin{sphinxVerbatim}[commandchars=\\\{\}]
\PYG{k}{thermo}\PYG{+w}{ }\PYG{l+m}{100}
\end{sphinxVerbatim}


\subsection{Description}
\label{\detokenize{thermo:description}}
\sphinxAtStartPar
Compute and print thermodynamic info (e.g. temperature, energy,
pressure) on timesteps that are a multiple of N and at the beginning
and end of a simulation.  A value of 0 will only print thermodynamics
at the beginning and end.

\sphinxAtStartPar
The content and format of what is printed is controlled by the
{\hyperref[\detokenize{thermo_style::doc}]{\sphinxcrossref{\DUrole{doc}{thermo\_style}}}} and
{\hyperref[\detokenize{thermo_modify::doc}]{\sphinxcrossref{\DUrole{doc}{thermo\_modify}}}} commands.

\sphinxAtStartPar
Instead of a numeric value, N can be specified as an {\hyperref[\detokenize{variable::doc}]{\sphinxcrossref{\DUrole{doc}{equal\sphinxhyphen{}style
variable}}}}, which should be specified as v\_name, where name is
the variable name.  In this case, the variable is evaluated at the
beginning of a run to determine the next timestep at which thermodynamic
info will be written out.  On that timestep, the variable will be
evaluated again to determine the next timestep, etc.  Thus the variable
should return timestep values.  See the stagger() and logfreq() and
stride() math functions for {\hyperref[\detokenize{variable::doc}]{\sphinxcrossref{\DUrole{doc}{equal\sphinxhyphen{}style variables}}}}, as
examples of useful functions to use in this context.  Other similar math
functions could easily be added as options for {\hyperref[\detokenize{variable::doc}]{\sphinxcrossref{\DUrole{doc}{equal\sphinxhyphen{}style
variables}}}}.

\sphinxAtStartPar
For example, the following commands will output thermodynamic info at
timesteps 0, 10, 20, 30, 100, 200, 300, 1000, 2000, \sphinxstyleemphasis{etc}:

\begin{sphinxVerbatim}[commandchars=\\\{\}]
\PYG{k}{variable        }\PYG{n+nv+nvIdentifier}{s}\PYG{+w}{ }\PYG{n}{equal}\PYG{+w}{ }\PYG{n}{logfreq}\PYG{n+nv}{(}\PYG{n+nv}{10,3,10}\PYG{n+nv}{)}
\PYG{k}{thermo}\PYG{+w}{          }\PYG{n}{v\PYGZus{}s}
\end{sphinxVerbatim}


\subsection{Restrictions}
\label{\detokenize{thermo:restrictions}}\begin{quote}

\sphinxAtStartPar
none
\end{quote}


\subsection{Related commands}
\label{\detokenize{thermo:related-commands}}
\sphinxAtStartPar
{\hyperref[\detokenize{thermo_style::doc}]{\sphinxcrossref{\DUrole{doc}{thermo\_style}}}}, {\hyperref[\detokenize{thermo_modify::doc}]{\sphinxcrossref{\DUrole{doc}{thermo\_modify}}}}


\subsection{Default}
\label{\detokenize{thermo:default}}
\begin{sphinxVerbatim}[commandchars=\\\{\}]
\PYG{k}{thermo}\PYG{+w}{ }\PYG{l+m}{0}
\end{sphinxVerbatim}

\sphinxstepscope

\index{thermo\_modify@\spxentry{thermo\_modify}}\ignorespaces 

\section{thermo\_modify command}
\label{\detokenize{thermo_modify:thermo-modify-command}}\label{\detokenize{thermo_modify:index-0}}\label{\detokenize{thermo_modify::doc}}

\subsection{Syntax}
\label{\detokenize{thermo_modify:syntax}}
\begin{sphinxVerbatim}[commandchars=\\\{\}]
\PYG{k}{thermo\PYGZus{}modify}\PYG{+w}{ }\PYG{n}{keyword}\PYG{+w}{ }\PYG{n}{value}\PYG{+w}{ }\PYG{n}{...}
\end{sphinxVerbatim}
\begin{itemize}
\item {} 
\sphinxAtStartPar
one or more keyword/value pairs may be listed

\item {} 
\sphinxAtStartPar
keyword = \sphinxstyleemphasis{lost} or \sphinxstyleemphasis{lost/bond} or \sphinxstyleemphasis{warn} or \sphinxstyleemphasis{norm} or \sphinxstyleemphasis{flush} or \sphinxstyleemphasis{line} or \sphinxstyleemphasis{colname} or \sphinxstyleemphasis{format} or \sphinxstyleemphasis{temp} or \sphinxstyleemphasis{press} or \sphinxstyleemphasis{triclinic/general}
\begin{sphinxalltt}
\sphinxstyleemphasis{lost} value = \sphinxstyleemphasis{error} or \sphinxstyleemphasis{warn} or \sphinxstyleemphasis{ignore}
\sphinxstyleemphasis{lost/bond} value = \sphinxstyleemphasis{error} or \sphinxstyleemphasis{warn} or \sphinxstyleemphasis{ignore}
\sphinxstyleemphasis{warn} value = \sphinxstyleemphasis{ignore} or \sphinxstyleemphasis{reset} or \sphinxstyleemphasis{default} or a number
\sphinxstyleemphasis{norm} value = \sphinxstyleemphasis{yes} or \sphinxstyleemphasis{no}
\sphinxstyleemphasis{flush} value = \sphinxstyleemphasis{yes} or \sphinxstyleemphasis{no}
\sphinxstyleemphasis{line} value = \sphinxstyleemphasis{one} or \sphinxstyleemphasis{multi} or \sphinxstyleemphasis{yaml}
\sphinxstyleemphasis{colname} values =  ID string, or \sphinxstyleemphasis{default}
  string = new column header name
  ID = integer from 1 to N, or integer from \sphinxhyphen{}1 to \sphinxhyphen{}N, where N = \# of quantities being output
       \sphinxstyleemphasis{or} a thermo keyword or reference to compute, fix, property or variable.
\sphinxstyleemphasis{format} values = \sphinxstyleemphasis{line} string, \sphinxstyleemphasis{int} string, \sphinxstyleemphasis{float} string, ID string, or \sphinxstyleemphasis{none}
  string = C\sphinxhyphen{}style format string
  ID = integer from 1 to N, or integer from \sphinxhyphen{}1 to \sphinxhyphen{}N, where N = \# of quantities being output
       \sphinxstyleemphasis{or} an integer range such as 2*6 (negative values are not allowed)
       \sphinxstyleemphasis{or} a thermo keyword or reference to compute, fix, property or variable.
\sphinxstyleemphasis{temp} value = compute ID that calculates a temperature
\sphinxstyleemphasis{press} value = compute ID that calculates a pressure
\sphinxstyleemphasis{triclinic/general} arg = \sphinxstyleemphasis{yes} or \sphinxstyleemphasis{no}
\end{sphinxalltt}

\end{itemize}


\subsection{Examples}
\label{\detokenize{thermo_modify:examples}}
\begin{sphinxVerbatim}[commandchars=\\\{\}]
\PYG{k}{thermo\PYGZus{}modify}\PYG{+w}{ }\PYG{n}{lost}\PYG{+w}{ }\PYG{n}{ignore}\PYG{+w}{ }\PYG{n}{flush}\PYG{+w}{ }\PYG{n}{yes}
\PYG{k}{thermo\PYGZus{}modify}\PYG{+w}{ }\PYG{n}{temp}\PYG{+w}{ }\PYG{n}{myTemp}\PYG{+w}{ }\PYG{n}{format}\PYG{+w}{ }\PYG{l+m}{3}\PYG{+w}{ }\PYG{o}{\PYGZpc{}}\PYG{l+m}{15.8}\PYG{n}{g}
\PYG{k}{thermo\PYGZus{}modify}\PYG{+w}{ }\PYG{n}{temp}\PYG{+w}{ }\PYG{n}{myTemp}\PYG{+w}{ }\PYG{n}{format}\PYG{+w}{ }\PYG{n}{line}\PYG{+w}{ }\PYG{l+s}{\PYGZdq{}}\PYG{l+s}{\PYGZpc{}ld \PYGZpc{}g \PYGZpc{}g \PYGZpc{}15.8g}\PYG{l+s}{\PYGZdq{}}
\PYG{k}{thermo\PYGZus{}modify}\PYG{+w}{ }\PYG{n}{line}\PYG{+w}{ }\PYG{n}{multi}\PYG{+w}{ }\PYG{n}{format}\PYG{+w}{ }\PYG{n}{float}\PYG{+w}{ }\PYG{o}{\PYGZpc{}}\PYG{n}{g}
\PYG{k}{thermo\PYGZus{}modify}\PYG{+w}{ }\PYG{n}{line}\PYG{+w}{ }\PYG{n}{yaml}\PYG{+w}{ }\PYG{n}{format}\PYG{+w}{ }\PYG{n}{none}
\PYG{k}{thermo\PYGZus{}modify}\PYG{+w}{ }\PYG{n}{colname}\PYG{+w}{ }\PYG{l+m}{1}\PYG{+w}{ }\PYG{n}{Timestep}\PYG{+w}{ }\PYG{n}{colname}\PYG{+w}{ }\PYG{o}{\PYGZhy{}}\PYG{l+m}{2}\PYG{+w}{ }\PYG{n}{Pressure}\PYG{+w}{ }\PYG{n}{colname}\PYG{+w}{ }\PYG{n}{f\PYGZus{}1[1]}\PYG{+w}{ }\PYG{n}{AvgDensity}
\end{sphinxVerbatim}


\subsection{Description}
\label{\detokenize{thermo_modify:description}}
\sphinxAtStartPar
Set options for how thermodynamic information is computed and printed by
LAMMPS.

\begin{sphinxadmonition}{note}{Note:}
\sphinxAtStartPar
These options apply to the \sphinxstyleemphasis{currently defined} thermo style.  When
you specify a {\hyperref[\detokenize{thermo_style::doc}]{\sphinxcrossref{\DUrole{doc}{thermo\_style}}}} command, all
thermodynamic settings are restored to their default values,
including those previously reset by a thermo\_modify command.  Thus if
your input script specifies a thermo\_style command, you should use
the thermo\_modify command \sphinxstylestrong{after} it.
\end{sphinxadmonition}

\sphinxAtStartPar
The \sphinxstyleemphasis{lost} keyword determines whether LAMMPS checks for lost atoms each
time it computes thermodynamics and what it does if atoms are lost.  An
atom can be “lost” if it moves across a non\sphinxhyphen{}periodic simulation box
{\hyperref[\detokenize{boundary::doc}]{\sphinxcrossref{\DUrole{doc}{boundary}}}} or if it moves more than a box length outside
the simulation domain (or more than a processor subdomain length)
before reneighboring occurs.  The latter case is typically due to bad
dynamics (e.g., too large a time step and/or huge forces and velocities).  If
the value is \sphinxstyleemphasis{ignore}, LAMMPS does not check for lost atoms.  If the
value is \sphinxstyleemphasis{error} or \sphinxstyleemphasis{warn}, LAMMPS checks and either issues an error or
warning.  The simulation will exit with an error and continue with a warning.
A warning will only be issued once, the first time an atom is lost.
This can be a useful debugging option.

\sphinxAtStartPar
The \sphinxstyleemphasis{lost/bond} keyword determines whether LAMMPS throws an error or
not if an atom in a bonded interaction (bond, angle, etc) cannot be
found when it creates bonded neighbor lists.  By default this is a
fatal error.  However in some scenarios it may be desirable to only
issue a warning or ignore it and skip the computation of the missing
bond, angle, etc.  An example would be when gas molecules in a vapor
are drifting out of the box through a fixed boundary condition (see
the {\hyperref[\detokenize{boundary::doc}]{\sphinxcrossref{\DUrole{doc}{boundary}}}} command).  In this case one atom may be
deleted before the rest of the molecule is, on a later timestep.

\sphinxAtStartPar
The \sphinxstyleemphasis{warn} keyword allows you to control whether LAMMPS will print
warning messages and how many of them.  Most warning messages are only
printed by MPI rank 0.  They are usually pointing out important issues
that should be investigated, but LAMMPS cannot determine for
certain whether they are an indication of an error.

\sphinxAtStartPar
Some warning messages are printed during a run (or immediately before)
each time a specific MPI rank encounters the issue (e.g., bonds that are
stretched too far or dihedrals in extreme configurations). These number
of these can quickly blow up the size of the log file and screen output.
Thus, a limit of 100 warning messages is applied by default.  The warning
count is applied to the entire input unless reset with a \sphinxcode{\sphinxupquote{thermo\_modify
warn reset}} command.  If there are more warnings than the limit, LAMMPS
will print one final warning that it will not print any additional
warning messages.

\begin{sphinxadmonition}{note}{Note:}
\sphinxAtStartPar
The warning limit is enforced on either the per\sphinxhyphen{}processor count or
the total count across all processors. For efficiency reasons,
however, the total count is only updated at steps with thermodynamic
output. Thus when running on a large number of processors in
parallel, the total number of warnings printed can be significantly
larger than the given limit.
\end{sphinxadmonition}

\sphinxAtStartPar
Any number after the keyword \sphinxstyleemphasis{warn} will change the warning limit
accordingly.  With the value \sphinxstyleemphasis{ignore} all warnings will be suppressed,
with the value \sphinxstyleemphasis{always} no limit will be applied and warnings will
always be printed, with the value \sphinxstyleemphasis{reset} the internal warning counter
will be reset to zero, and with the value \sphinxstyleemphasis{default}, the counter is
reset and the limit set to 100.  An example usage of either \sphinxstyleemphasis{reset} or
\sphinxstyleemphasis{default} would be to re\sphinxhyphen{}enable warnings that were disabled or have
reached the limit during equilibration, where the warnings would be
acceptable while the system is still adjusting, but then change to all
warnings for the production run, where they would indicate problems that
would require a closer look at what is causing them.

\sphinxAtStartPar
The \sphinxstyleemphasis{norm} keyword determines whether various thermodynamic output
values are normalized by the number of atoms or not, depending on
whether it is set to \sphinxstyleemphasis{yes} or \sphinxstyleemphasis{no}.  Different unit styles have
different defaults for this setting (see below).  Even if \sphinxstyleemphasis{norm} is set
to \sphinxstyleemphasis{yes}, a value is only normalized if it is an “extensive” quantity,
meaning that it scales with the number of atoms in the system.  For the
thermo keywords described by the page for the {\hyperref[\detokenize{thermo_style::doc}]{\sphinxcrossref{\DUrole{doc}{thermo\_style}}}} command, all energy\sphinxhyphen{}related keywords are extensive, such
as \sphinxstyleemphasis{pe} or \sphinxstyleemphasis{ebond} or \sphinxstyleemphasis{enthalpy}.  Other keywords such as \sphinxstyleemphasis{temp} or
\sphinxstyleemphasis{press} are “intensive” meaning their value is independent (in a
statistical sense) of the number of atoms in the system and thus are
never normalized.  For thermodynamic output values extracted from fixes
and computes in a {\hyperref[\detokenize{thermo_style::doc}]{\sphinxcrossref{\DUrole{doc}{thermo\_style custom}}}} command, the
page for the individual {\hyperref[\detokenize{fix::doc}]{\sphinxcrossref{\DUrole{doc}{fix}}}} or {\hyperref[\detokenize{compute::doc}]{\sphinxcrossref{\DUrole{doc}{compute}}}}
lists whether the value is “extensive” or “intensive” and thus whether
it is normalized.  Thermodynamic output values calculated by a variable
formula are assumed to be “intensive” and thus are never normalized.
You can always include a divide by the number of atoms in the variable
formula if this is not the case.

\sphinxAtStartPar
The \sphinxstyleemphasis{flush} keyword invokes a flush operation after thermodynamic info
is written to the screen and log file.  This ensures the output is
updated and not buffered (by the application) even if LAMMPS halts
before the simulation completes.  Please note that this does not affect
buffering by the OS or devices, so you may still lose data in case the
simulation stops due to a hardware failure.

\sphinxAtStartPar
The \sphinxstyleemphasis{line} keyword determines whether thermodynamics will be output as a
series of numeric values on one line (“one”), in a multi\sphinxhyphen{}line format
with 3 quantities with text strings per line and a dashed\sphinxhyphen{}line header
containing the timestep and CPU time (“multi”), or in a YAML format
block (“yaml”).  This modify option overrides the \sphinxstyleemphasis{one}, \sphinxstyleemphasis{multi}, or
\sphinxstyleemphasis{yaml} thermo\_style settings.

\sphinxAtStartPar
\DUrole{versionmodified,added}{Added in version 4May2022.}

\sphinxAtStartPar
The \sphinxstyleemphasis{colname} keyword can be used to change the default header keyword
for a column or field of thermodynamic output.  The setting for \sphinxstyleemphasis{ID
string} replaces the default text with the provided string.  \sphinxstyleemphasis{ID} can be
a positive integer when it represents the column number counting from
the left, a negative integer when it represents the column number from
the right (i.e., \(-1\) is the last column/keyword), or a thermo keyword
(or compute, fix, property, or variable reference) and then it replaces the
string for that specific thermo keyword.

\sphinxAtStartPar
The \sphinxstyleemphasis{colname} keyword can be used multiple times. If multiple \sphinxstyleemphasis{colname}
settings refer to the same keyword, the last setting has precedence.  A
setting of \sphinxstyleemphasis{default} clears all previous settings, reverting all values
to their default values.

\sphinxAtStartPar
The \sphinxstyleemphasis{format} keyword can be used to change the default numeric format of
any of quantities the {\hyperref[\detokenize{thermo_style::doc}]{\sphinxcrossref{\DUrole{doc}{thermo\_style}}}} command
outputs.  All the specified format strings are C\sphinxhyphen{}style formats (i.e., as
used by the C/C++ printf() command).  The \sphinxstyleemphasis{line} keyword takes a single
argument which is the format string for the entire line of thermo
output, with \(N\) fields, which you must enclose in quotes if it is more
than one field.  The \sphinxstyleemphasis{int} and \sphinxstyleemphasis{float} keywords take a single format
argument and are applied to all integer or floating\sphinxhyphen{}point quantities
output.  The setting for \sphinxstyleemphasis{ID string} also takes a single format argument
that is used for the indexed value in each line.  The interpretation is
the same as for \sphinxstyleemphasis{colname} (i.e., a positive integer is the n\sphinxhyphen{}th value
corresponding to the n\sphinxhyphen{}th thermo keyword, a negative integer is counting
backwards, and a string matches the entry with the thermo keyword).
For example, the fifth column is output in high precision for
“format 5 \%20.15g”, and the pair energy for “format epair \%20.15g”.
The \sphinxstyleemphasis{ID} field can be a range, such as “3*6”, “*”, “2*”, or “*3”;
in such cases, all fields in the range (inclusive) are set to the specified
format string. Ranges containing negative numbers are not supported.

\sphinxAtStartPar
The \sphinxstyleemphasis{format} keyword can be used multiple times.  The precedence is
that for each value in a line of output, the \sphinxstyleemphasis{ID} format (if specified)
is used, else the \sphinxstyleemphasis{int} or \sphinxstyleemphasis{float} setting (if specified) is used,
else the \sphinxstyleemphasis{line} setting (if specified) for that value is used, else
the default setting is used.  A setting of \sphinxstyleemphasis{none} clears all previous
settings, reverting all values to their default format.

\begin{sphinxadmonition}{note}{Note:}
\sphinxAtStartPar
The thermo output values \sphinxstyleemphasis{step} and \sphinxstyleemphasis{atoms} are stored internally as
8\sphinxhyphen{}byte signed integers, rather than the usual 4\sphinxhyphen{}byte signed integers.
When specifying the \sphinxstyleemphasis{format int} option you can use a “\%d”\sphinxhyphen{}style
format identifier in the format string and LAMMPS will convert this
to the corresponding 8\sphinxhyphen{}byte form when it is applied to those
keywords.  However, when specifying the \sphinxstyleemphasis{line} option or \sphinxstyleemphasis{format ID
string} option for \sphinxstyleemphasis{step} and \sphinxstyleemphasis{natoms}, you should specify a format
string appropriate for an 8\sphinxhyphen{}byte signed integer (i.e., one with “\%ld”
or “\%lld”, depending on the platform).
\end{sphinxadmonition}

\sphinxAtStartPar
The \sphinxstyleemphasis{temp} keyword is used to determine how thermodynamic temperature is
calculated, which is used by all thermo quantities that require a
temperature (“temp”, “press”, “ke”, “etotal”, “enthalpy”, “pxx”, etc).
The specified compute ID must have been previously defined by the user
via the {\hyperref[\detokenize{compute::doc}]{\sphinxcrossref{\DUrole{doc}{compute}}}} command and it must be a style of
compute that calculates a temperature.  As described in the
{\hyperref[\detokenize{thermo_style::doc}]{\sphinxcrossref{\DUrole{doc}{thermo\_style}}}} command, thermo output uses a default
compute for temperature with ID = \sphinxstyleemphasis{thermo\_temp}.  This option allows the
user to override the default.

\sphinxAtStartPar
The \sphinxstyleemphasis{press} keyword is used to determine how thermodynamic pressure is
calculated, which is used by all thermo quantities that require a
pressure (“press”, “enthalpy”, “pxx”, etc).  The specified compute ID
must have been previously defined by the user via the {\hyperref[\detokenize{compute::doc}]{\sphinxcrossref{\DUrole{doc}{compute}}}} command and it must be a style of compute that calculates a
pressure.  As described in the {\hyperref[\detokenize{thermo_style::doc}]{\sphinxcrossref{\DUrole{doc}{thermo\_style}}}}
command, thermo output uses a default compute for pressure with ID =
\sphinxstyleemphasis{thermo\_press}.  This option allows the user to override the default.

\begin{sphinxadmonition}{note}{Note:}
\sphinxAtStartPar
If both the \sphinxstyleemphasis{temp} and \sphinxstyleemphasis{press} keywords are used in a single
thermo\_modify command (or in two separate commands), then the order
in which the keywords are specified is important.  Note that a
{\hyperref[\detokenize{compute_pressure::doc}]{\sphinxcrossref{\DUrole{doc}{pressure compute}}}} defines its own
temperature compute as an argument when it is specified.  The \sphinxstyleemphasis{temp}
keyword will override this (for the pressure compute being used by
thermodynamics), but only if the \sphinxstyleemphasis{temp} keyword comes after the
\sphinxstyleemphasis{press} keyword.  If the \sphinxstyleemphasis{temp} keyword comes before the \sphinxstyleemphasis{press}
keyword, then the new pressure compute specified by the \sphinxstyleemphasis{press}
keyword will be unaffected by the \sphinxstyleemphasis{temp} setting.
\end{sphinxadmonition}

\sphinxAtStartPar
The \sphinxstyleemphasis{triclinic/general} keyword can only be used with a value of \sphinxstyleemphasis{yes}
if the simulation box was created as a general triclinic box.  See the
{\hyperref[\detokenize{Howto_triclinic::doc}]{\sphinxcrossref{\DUrole{doc}{Howto\_triclinic}}}} doc page for a detailed
explanation of orthogonal, restricted triclinic, and general triclinic
simulation boxes.

\sphinxAtStartPar
If this keyword is \sphinxstyleemphasis{yes}, the output of the simulation box edge
vectors and the pressure tensor components for the system are
affected.  These are specified by the \sphinxstyleemphasis{avec,bvec,cvec} and
\sphinxstyleemphasis{pxx,pyy,pzz,pxy,pxz,pyz} keywords of the {\hyperref[\detokenize{thermo_style::doc}]{\sphinxcrossref{\DUrole{doc}{thermo\_style}}}} command.  See the {\hyperref[\detokenize{thermo_style::doc}]{\sphinxcrossref{\DUrole{doc}{thermo\_style}}}}
doc page for details.


\subsection{Restrictions}
\label{\detokenize{thermo_modify:restrictions}}
\sphinxAtStartPar
none


\subsection{Related commands}
\label{\detokenize{thermo_modify:related-commands}}
\sphinxAtStartPar
{\hyperref[\detokenize{thermo::doc}]{\sphinxcrossref{\DUrole{doc}{thermo}}}}, {\hyperref[\detokenize{thermo_style::doc}]{\sphinxcrossref{\DUrole{doc}{thermo\_style}}}}


\subsection{Default}
\label{\detokenize{thermo_modify:default}}
\sphinxAtStartPar
The option defaults are lost = error, warn = 100, norm = yes for unit
style of \sphinxstyleemphasis{lj}, norm = no for unit style of \sphinxstyleemphasis{real} and \sphinxstyleemphasis{metal}, flush =
no, temp/press = compute IDs defined by thermo\_style, and
triclinic/general = no.

\sphinxAtStartPar
The defaults for the line and format options depend on the thermo style.
For styles “one” and “custom”, the line and format defaults are “one”,
“\%10d”, and “\%14.8g”.  For style “multi”, the line and format defaults
are “multi”, “\%14d”, and “\%14.4f”. For style “yaml”, the line and format
defaults are “\%d” and “\%.15g”.

\sphinxstepscope

\index{thermo\_style@\spxentry{thermo\_style}}\ignorespaces 

\section{thermo\_style command}
\label{\detokenize{thermo_style:thermo-style-command}}\label{\detokenize{thermo_style:index-0}}\label{\detokenize{thermo_style::doc}}

\subsection{Syntax}
\label{\detokenize{thermo_style:syntax}}
\begin{sphinxVerbatim}[commandchars=\\\{\}]
\PYG{k}{thermo\PYGZus{}style}\PYG{+w}{ }\PYG{n}{style}\PYG{+w}{ }\PYG{n}{args}
\end{sphinxVerbatim}
\begin{itemize}
\item {} 
\sphinxAtStartPar
style = \sphinxstyleemphasis{one} or \sphinxstyleemphasis{multi} or \sphinxstyleemphasis{yaml} or \sphinxstyleemphasis{custom}

\item {} 
\sphinxAtStartPar
args = list of arguments for a particular style
\begin{sphinxalltt}
\sphinxstyleemphasis{one} args = none
\sphinxstyleemphasis{multi} args = none
\sphinxstyleemphasis{yaml} args = none
\sphinxstyleemphasis{custom} args = list of keywords
  possible keywords = step, elapsed, elaplong, dt, time,
                      cpu, tpcpu, spcpu, cpuremain, part, timeremain,
                      atoms, temp, press, pe, ke, etotal,
                      evdwl, ecoul, epair, ebond, eangle, edihed, eimp,
                      emol, elong, etail,
                      enthalpy, ecouple, econserve,
                      vol, density,
                      xlo, xhi, ylo, yhi, zlo, zhi,
                      xy, xz, yz,
                      avecx, avecy, avecz,
                      bvecx, bvecy, bvecz,
                      cvecx, cvecy, cvecz,
                      lx, ly, lz,
                      xlat, ylat, zlat,
                      cella, cellb, cellc, cellalpha, cellbeta, cellgamma,
                      pxx, pyy, pzz, pxy, pxz, pyz,
                      bonds, angles, dihedrals, impropers,
                      fmax, fnorm, nbuild, ndanger,
                      c\_ID, c\_ID{[}I{]}, c\_ID{[}I{]}{[}J{]},
                      f\_ID, f\_ID{[}I{]}, f\_ID{[}I{]}{[}J{]},
                      v\_name, v\_name{[}I{]}
    step = timestep
    elapsed = timesteps since start of this run
    elaplong = timesteps since start of initial run in a series of runs
    dt = timestep size
    time = simulation time
    cpu = elapsed CPU time in seconds since start of this run
    tpcpu = time per CPU second
    spcpu = timesteps per CPU second
    cpuremain = estimated CPU time remaining in run
    part = which partition (0 to Npartition\sphinxhyphen{}1) this is
    timeremain = remaining time in seconds on timer timeout.
    atoms = \# of atoms
    temp = temperature
    press = pressure
    pe = total potential energy
    ke = kinetic energy
    etotal = total energy (pe + ke)
    evdwl = van der Waals pairwise energy (includes etail)
    ecoul = Coulombic pairwise energy
    epair = pairwise energy (evdwl + ecoul + elong)
    ebond = bond energy
    eangle = angle energy
    edihed = dihedral energy
    eimp = improper energy
    emol = molecular energy (ebond + eangle + edihed + eimp)
    elong = long\sphinxhyphen{}range kspace energy
    etail = van der Waals energy long\sphinxhyphen{}range tail correction
    enthalpy = enthalpy (etotal + press*vol)
    ecouple = cumulative energy change due to thermo/baro statting fixes
    econserve = pe + ke + ecouple = etotal + ecouple
    vol = volume
    density = mass density of system
    xlo,xhi,ylo,yhi,zlo,zhi = box boundaries
    xy,xz,yz = box tilt for restricted triclinic (non\sphinxhyphen{}orthogonal) simulation boxes
    avecx,avecy,avecz = components of edge vector A of the simulation box
    bvecx,bvecy,bvecz = components of edge vector B of the simulation box
    cvecx,cvecy,cvecz = components of edge vector C of the simulation box
    lx,ly,lz = box lengths in x,y,z
    xlat,ylat,zlat = lattice spacings as calculated by {\hyperref[\detokenize{lattice::doc}]{\sphinxcrossref{\DUrole{doc}{lattice}}}} command
    cella,cellb,cellc = periodic cell lattice constants a,b,c
    cellalpha, cellbeta, cellgamma = periodic cell angles alpha,beta,gamma
    pxx,pyy,pzz,pxy,pxz,pyz = 6 components of pressure tensor
    bonds,angles,dihedrals,impropers = \# of these interactions defined
    fmax = max component of force on any atom in any dimension
    fnorm = length of force vector for all atoms
    nbuild = \# of neighbor list builds
    ndanger = \# of dangerous neighbor list builds
    c\_ID = global scalar value calculated by a compute with ID
    c\_ID{[}I{]} = Ith component of global vector calculated by a compute with ID, I can include wildcard (see below)
    c\_ID{[}I{]}{[}J{]} = I,J component of global array calculated by a compute with ID
    f\_ID = global scalar value calculated by a fix with ID
    f\_ID{[}I{]} = Ith component of global vector calculated by a fix with ID, I can include wildcard (see below)
    f\_ID{[}I{]}{[}J{]} = I,J component of global array calculated by a fix with ID
    v\_name = value calculated by an equal\sphinxhyphen{}style variable with name
    v\_name{[}I{]} = value calculated by a vector\sphinxhyphen{}style variable with name, I can include wildcard (see below)
\end{sphinxalltt}

\end{itemize}


\subsection{Examples}
\label{\detokenize{thermo_style:examples}}
\begin{sphinxVerbatim}[commandchars=\\\{\}]
\PYG{k}{thermo\PYGZus{}style}\PYG{+w}{ }\PYG{n}{multi}
\PYG{k}{thermo\PYGZus{}style}\PYG{+w}{ }\PYG{n}{yaml}
\PYG{k}{thermo\PYGZus{}style}\PYG{+w}{ }\PYG{n}{one}
\PYG{k}{thermo\PYGZus{}style}\PYG{+w}{ }\PYG{n}{custom}\PYG{+w}{ }\PYG{n}{step}\PYG{+w}{ }\PYG{n}{temp}\PYG{+w}{ }\PYG{n}{pe}\PYG{+w}{ }\PYG{n}{etotal}\PYG{+w}{ }\PYG{n}{press}\PYG{+w}{ }\PYG{n}{vol}
\PYG{k}{thermo\PYGZus{}style}\PYG{+w}{ }\PYG{n}{custom}\PYG{+w}{ }\PYG{n}{step}\PYG{+w}{ }\PYG{n}{temp}\PYG{+w}{ }\PYG{n}{etotal}\PYG{+w}{ }\PYG{n}{c\PYGZus{}myTemp}\PYG{+w}{ }\PYG{n}{v\PYGZus{}abc}
\PYG{k}{thermo\PYGZus{}style}\PYG{+w}{ }\PYG{n}{custom}\PYG{+w}{ }\PYG{n}{step}\PYG{+w}{ }\PYG{n}{temp}\PYG{+w}{ }\PYG{n}{etotal}\PYG{+w}{ }\PYG{n}{c\PYGZus{}myTemp[}\PYG{o}{*}\PYG{n}{]}\PYG{+w}{ }\PYG{n}{v\PYGZus{}abc}
\end{sphinxVerbatim}


\subsection{Description}
\label{\detokenize{thermo_style:description}}
\sphinxAtStartPar
Set the style and content for printing thermodynamic data to the
screen and log files.  The units for each column of output
corresponding to the list of keywords is determined by the {\hyperref[\detokenize{units::doc}]{\sphinxcrossref{\DUrole{doc}{units}}}} command for the simulation.  E.g. energies will be in energy
units, temperature in temperature units, pressure in pressure units.

\sphinxAtStartPar
Style \sphinxstyleemphasis{one} prints a single line of thermodynamic info that is the
equivalent of “thermo\_style custom step temp epair emol etotal press”.
The line contains only numeric values.

\sphinxAtStartPar
Style \sphinxstyleemphasis{multi} prints a multiple\sphinxhyphen{}line listing of thermodynamic info
that is the equivalent of “thermo\_style custom etotal ke temp pe ebond
eangle edihed eimp evdwl ecoul elong press”.  The listing contains
numeric values and a string ID for each quantity.

\sphinxAtStartPar
\DUrole{versionmodified,added}{Added in version 24Mar2022.}

\sphinxAtStartPar
Style \sphinxstyleemphasis{yaml} is similar to style \sphinxstyleemphasis{one} but prints the output in \sphinxhref{https://yaml.org/}{YAML} format which can be easily read by a variety of
script languages and data handling packages.  Since LAMMPS may print
other output before, after, or in between thermodynamic output, the
YAML format content needs to be separated from the rest.  All YAML
format thermodynamic output can be matched with a regular expression
and can thus be extracted with commands like \sphinxcode{\sphinxupquote{egrep}} as follows:

\begin{sphinxVerbatim}[commandchars=\\\{\}]
egrep\PYG{+w}{  }\PYG{l+s+s1}{\PYGZsq{}\PYGZca{}(keywords:|data:\PYGZdl{}|\PYGZhy{}\PYGZhy{}\PYGZhy{}\PYGZdl{}|\PYGZbs{}.\PYGZbs{}.\PYGZbs{}.\PYGZdl{}|  \PYGZhy{} \PYGZbs{}[)\PYGZsq{}}\PYG{+w}{ }log.lammps\PYG{+w}{ }\PYGZgt{}\PYG{+w}{ }log.yaml
\end{sphinxVerbatim}

\sphinxAtStartPar
Information about processing such YAML files is in the {\hyperref[\detokenize{Howto_structured_data::doc}]{\sphinxcrossref{\DUrole{doc}{structured
data output howto}}}}.

\sphinxAtStartPar
Style \sphinxstyleemphasis{custom} is the most general setting and allows you to specify
which of the keywords listed above you want printed on each
thermodynamic timestep.  Note that the keywords c\_ID, f\_ID, v\_name are
references to {\hyperref[\detokenize{compute::doc}]{\sphinxcrossref{\DUrole{doc}{computes}}}}, {\hyperref[\detokenize{fix::doc}]{\sphinxcrossref{\DUrole{doc}{fixes}}}}, and
equal\sphinxhyphen{}style {\hyperref[\detokenize{variable::doc}]{\sphinxcrossref{\DUrole{doc}{variables}}}} that have been defined elsewhere
in the input script or can even be new styles which users have added
to LAMMPS.  See the {\hyperref[\detokenize{Modify::doc}]{\sphinxcrossref{\DUrole{doc}{Modify}}}} page for details on the
latter.  Thus the \sphinxstyleemphasis{custom} style provides a flexible means of
outputting essentially any desired quantity as a simulation proceeds.

\sphinxAtStartPar
All styles except \sphinxstyleemphasis{custom} have \sphinxstyleemphasis{vol} appended to their list of
outputs if the simulation box volume changes during the simulation.

\sphinxAtStartPar
The values printed by the various keywords are instantaneous values,
calculated on the current timestep.  Time\sphinxhyphen{}averaged quantities, which
include values from previous timesteps, can be output by using the
f\_ID keyword and accessing a fix that does time\sphinxhyphen{}averaging such as the
{\hyperref[\detokenize{fix_ave_time::doc}]{\sphinxcrossref{\DUrole{doc}{fix ave/time}}}} command.

\sphinxAtStartPar
Options invoked by the {\hyperref[\detokenize{thermo_modify::doc}]{\sphinxcrossref{\DUrole{doc}{thermo\_modify}}}} command can
be used to set the one\sphinxhyphen{} or multi\sphinxhyphen{}line format of the print\sphinxhyphen{}out, the
normalization of thermodynamic output (total values versus per\sphinxhyphen{}atom
values for extensive quantities (ones which scale with the number of
atoms in the system), and the numeric precision of each printed value.

\begin{sphinxadmonition}{note}{Note:}
\sphinxAtStartPar
When you use a “thermo\_style” command, all thermodynamic
settings are restored to their default values, including those
previously set by a {\hyperref[\detokenize{thermo_modify::doc}]{\sphinxcrossref{\DUrole{doc}{thermo\_modify}}}} command.  Thus
if your input script specifies a thermo\_style command, you should use
the thermo\_modify command after it.
\end{sphinxadmonition}


\bigskip\hrule\bigskip


\sphinxAtStartPar
Several of the thermodynamic quantities require a temperature to be
computed: “temp”, “press”, “ke”, “etotal”, “enthalpy”, “pxx”, etc.  By
default this is done by using a \sphinxstyleemphasis{temperature} compute which is created
when LAMMPS starts up, as if this command had been issued:

\begin{sphinxVerbatim}[commandchars=\\\{\}]
\PYG{k}{compute }\PYG{n+nv+nvIdentifier}{thermo\PYGZus{}temp}\PYG{+w}{ }\PYG{n+nv+nvIdentifier}{all}\PYG{+w}{ }\PYG{n}{temp}
\end{sphinxVerbatim}

\sphinxAtStartPar
See the {\hyperref[\detokenize{compute_temp::doc}]{\sphinxcrossref{\DUrole{doc}{compute temp}}}} command for details.  Note
that the ID of this compute is \sphinxstyleemphasis{thermo\_temp} and the group is \sphinxstyleemphasis{all}.
You can change the attributes of this temperature (e.g. its
degrees\sphinxhyphen{}of\sphinxhyphen{}freedom) via the {\hyperref[\detokenize{compute_modify::doc}]{\sphinxcrossref{\DUrole{doc}{compute\_modify}}}}
command.  Alternatively, you can directly assign a new compute (that
calculates temperature) which you have defined, to be used for
calculating any thermodynamic quantity that requires a temperature.
This is done via the {\hyperref[\detokenize{thermo_modify::doc}]{\sphinxcrossref{\DUrole{doc}{thermo\_modify}}}} command.

\sphinxAtStartPar
Several of the thermodynamic quantities require a pressure to be
computed: “press”, “enthalpy”, “pxx”, etc.  By default this is done by
using a \sphinxstyleemphasis{pressure} compute which is created when LAMMPS starts up, as
if this command had been issued:

\begin{sphinxVerbatim}[commandchars=\\\{\}]
\PYG{k}{compute }\PYG{n+nv+nvIdentifier}{thermo\PYGZus{}press}\PYG{+w}{ }\PYG{n+nv+nvIdentifier}{all}\PYG{+w}{ }\PYG{n}{pressure}\PYG{+w}{ }\PYG{n}{thermo\PYGZus{}temp}
\end{sphinxVerbatim}

\sphinxAtStartPar
See the {\hyperref[\detokenize{compute_pressure::doc}]{\sphinxcrossref{\DUrole{doc}{compute pressure}}}} command for details.
Note that the ID of this compute is \sphinxstyleemphasis{thermo\_press} and the group is
\sphinxstyleemphasis{all}.  You can change the attributes of this pressure via the
{\hyperref[\detokenize{compute_modify::doc}]{\sphinxcrossref{\DUrole{doc}{compute\_modify}}}} command.  Alternatively, you can
directly assign a new compute (that calculates pressure) which you
have defined, to be used for calculating any thermodynamic quantity
that requires a pressure.  This is done via the
{\hyperref[\detokenize{thermo_modify::doc}]{\sphinxcrossref{\DUrole{doc}{thermo\_modify}}}} command.

\sphinxAtStartPar
Several of the thermodynamic quantities require a potential energy to
be computed: “pe”, “etotal”, “ebond”, etc.  This is done by using a
\sphinxstyleemphasis{pe} compute which is created when LAMMPS starts up, as if this
command had been issued:

\begin{sphinxVerbatim}[commandchars=\\\{\}]
\PYG{k}{compute }\PYG{n+nv+nvIdentifier}{thermo\PYGZus{}pe}\PYG{+w}{ }\PYG{n+nv+nvIdentifier}{all}\PYG{+w}{ }\PYG{n}{pe}
\end{sphinxVerbatim}

\sphinxAtStartPar
See the {\hyperref[\detokenize{compute_pe::doc}]{\sphinxcrossref{\DUrole{doc}{compute pe}}}} command for details.  Note that
the ID of this compute is \sphinxstyleemphasis{thermo\_pe} and the group is \sphinxstyleemphasis{all}.  You can
change the attributes of this potential energy via the
{\hyperref[\detokenize{compute_modify::doc}]{\sphinxcrossref{\DUrole{doc}{compute\_modify}}}} command.


\bigskip\hrule\bigskip


\sphinxAtStartPar
The kinetic energy of the system \sphinxstyleemphasis{ke} is inferred from the temperature
of the system with \(\frac{1}{2} k_B T\) of energy for each degree
of freedom.  Thus, using different {\hyperref[\detokenize{compute::doc}]{\sphinxcrossref{\DUrole{doc}{compute commands}}}}
for calculating temperature, via the {\hyperref[\detokenize{thermo_modify::doc}]{\sphinxcrossref{\DUrole{doc}{thermo\_modify temp}}}} command, may yield different kinetic energies, since
different computes that calculate temperature can subtract out
different non\sphinxhyphen{}thermal components of velocity and/or include different
degrees of freedom (translational, rotational, etc).

\sphinxAtStartPar
The potential energy of the system \sphinxstyleemphasis{pe} will include contributions
from fixes if the {\hyperref[\detokenize{fix_modify::doc}]{\sphinxcrossref{\DUrole{doc}{fix\_modify energy yes}}}} option is
set for a fix that calculates such a contribution.  For example, the
{\hyperref[\detokenize{fix_wall::doc}]{\sphinxcrossref{\DUrole{doc}{fix wall/lj93}}}} fix calculates the energy of atoms
interacting with the wall.  See the doc pages for “individual fixes”
to see which ones contribute and whether their default
{\hyperref[\detokenize{fix_modify::doc}]{\sphinxcrossref{\DUrole{doc}{fix\_modify energy}}}} setting is \sphinxstyleemphasis{yes} or \sphinxstyleemphasis{no}.

\sphinxAtStartPar
A long\sphinxhyphen{}range tail correction \sphinxstyleemphasis{etail} for the van der Waals pairwise
energy will be non\sphinxhyphen{}zero only if the {\hyperref[\detokenize{pair_modify::doc}]{\sphinxcrossref{\DUrole{doc}{pair\_modify tail}}}} option is turned on.  The \sphinxstyleemphasis{etail} contribution is
included in \sphinxstyleemphasis{evdwl}, \sphinxstyleemphasis{epair}, \sphinxstyleemphasis{pe}, and \sphinxstyleemphasis{etotal}, and the
corresponding tail correction to the pressure is included in \sphinxstyleemphasis{press}
and \sphinxstyleemphasis{pxx}, \sphinxstyleemphasis{pyy}, etc.


\bigskip\hrule\bigskip


\sphinxAtStartPar
Here is more information on other keywords whose meaning may not be
clear.

\sphinxAtStartPar
The \sphinxstyleemphasis{step}, \sphinxstyleemphasis{elapsed}, and \sphinxstyleemphasis{elaplong} keywords refer to timestep
count.  \sphinxstyleemphasis{Step} is the current timestep, or iteration count when a
{\hyperref[\detokenize{minimize::doc}]{\sphinxcrossref{\DUrole{doc}{minimization}}}} is being performed.  \sphinxstyleemphasis{Elapsed} is the
number of timesteps elapsed since the beginning of this run.
\sphinxstyleemphasis{Elaplong} is the number of timesteps elapsed since the beginning of
an initial run in a series of runs.  See the \sphinxstyleemphasis{start} and \sphinxstyleemphasis{stop}
keywords for the {\hyperref[\detokenize{run::doc}]{\sphinxcrossref{\DUrole{doc}{run}}}} for info on how to invoke a series of
runs that keep track of an initial starting time.  If these keywords
are not used, then \sphinxstyleemphasis{elapsed} and \sphinxstyleemphasis{elaplong} are the same value.

\sphinxAtStartPar
The \sphinxstyleemphasis{dt} keyword is the current timestep size in time {\hyperref[\detokenize{units::doc}]{\sphinxcrossref{\DUrole{doc}{units}}}}.  The \sphinxstyleemphasis{time} keyword is the current elapsed simulation time,
also in time {\hyperref[\detokenize{units::doc}]{\sphinxcrossref{\DUrole{doc}{units}}}}, which is simply (step*dt) if the
timestep size has not changed and the timestep has not been reset.  If
the timestep has changed (e.g. via {\hyperref[\detokenize{fix_dt_reset::doc}]{\sphinxcrossref{\DUrole{doc}{fix dt/reset}}}})
or the timestep has been reset (e.g. via the “reset\_timestep”
command), then the simulation time is effectively a cumulative value
up to the current point.

\sphinxAtStartPar
The \sphinxstyleemphasis{cpu} keyword is elapsed CPU seconds since the beginning of this
run.  The \sphinxstyleemphasis{tpcpu} and \sphinxstyleemphasis{spcpu} keywords are measures of how fast your
simulation is currently running.  The \sphinxstyleemphasis{tpcpu} keyword is simulation
time per CPU second, where simulation time is in time
{\hyperref[\detokenize{units::doc}]{\sphinxcrossref{\DUrole{doc}{units}}}}.  E.g. for metal units, the \sphinxstyleemphasis{tpcpu} value would be
picoseconds per CPU second.  The \sphinxstyleemphasis{spcpu} keyword is the number of
timesteps per CPU second.  Both quantities are on\sphinxhyphen{}the\sphinxhyphen{}fly metrics,
measured relative to the last time they were invoked.  Thus if you are
printing out thermodynamic output every 100 timesteps, the two keywords
will continually output the time and timestep rate for the last 100
steps.  The \sphinxstyleemphasis{tpcpu} keyword does not attempt to track any changes in
timestep size, e.g. due to using the {\hyperref[\detokenize{fix_dt_reset::doc}]{\sphinxcrossref{\DUrole{doc}{fix dt/reset}}}}
command.

\sphinxAtStartPar
The \sphinxstyleemphasis{cpuremain} keyword estimates the CPU time remaining in the
current run, based on the time elapsed thus far.  It will only be a
good estimate if the CPU time/timestep for the rest of the run is
similar to the preceding timesteps.  On the initial timestep the value
will be 0.0 since there is no history to estimate from.  For a
minimization run performed by the “minimize” command, the estimate is
based on the \sphinxstyleemphasis{maxiter} parameter, assuming the minimization will
proceed for the maximum number of allowed iterations.

\sphinxAtStartPar
The \sphinxstyleemphasis{part} keyword is useful for multi\sphinxhyphen{}replica or multi\sphinxhyphen{}partition
simulations to indicate which partition this output and this file
corresponds to, or for use in a {\hyperref[\detokenize{variable::doc}]{\sphinxcrossref{\DUrole{doc}{variable}}}} to append to
a filename for output specific to this partition.  See discussion of
the {\hyperref[\detokenize{Run_options::doc}]{\sphinxcrossref{\DUrole{doc}{\sphinxhyphen{}partition command\sphinxhyphen{}line switch}}}} for details on
running in multi\sphinxhyphen{}partition mode.

\sphinxAtStartPar
The \sphinxstyleemphasis{timeremain} keyword is the seconds remaining when a timeout has
been configured via the {\hyperref[\detokenize{timer::doc}]{\sphinxcrossref{\DUrole{doc}{timer timeout}}}} command.  If the
timeout timer is inactive, the value of this keyword is 0.0 and if the
timer is expired, it is negative. This allows for example to exit
loops cleanly, if the timeout is expired with:

\begin{sphinxVerbatim}[commandchars=\\\{\}]
\PYG{k}{if}\PYG{+w}{ }\PYG{l+s}{\PYGZdq{}}\PYG{l+s}{\PYGZdl{}(timeremain) \PYGZlt{} 0.0}\PYG{l+s}{\PYGZdq{}}\PYG{+w}{ }\PYG{k}{then}\PYG{+w}{ }\PYG{l+s}{\PYGZdq{}}\PYG{l+s}{quit 0}\PYG{l+s}{\PYGZdq{}}
\end{sphinxVerbatim}

\sphinxAtStartPar
The \sphinxstyleemphasis{ecouple} keyword is cumulative energy change in the system due to
any thermostatting or barostatting fixes that are being used.  A
positive value means that energy has been subtracted from the system
(added to the coupling reservoir).  See the \sphinxstyleemphasis{econserve} keyword for an
explanation of why this sign choice makes sense.

\sphinxAtStartPar
The \sphinxstyleemphasis{econserve} keyword is the sum of the potential and kinetic energy
of the system as well as the energy that has been transferred by
thermostatting or barostatting to their coupling reservoirs \textendash{} that is,
\sphinxstyleemphasis{econserve} = \sphinxstyleemphasis{pe} + \sphinxstyleemphasis{ke} + \sphinxstyleemphasis{ecouple}.  Ideally, for a simulation in
the NVT, NPH, or NPT ensembles, the \sphinxstyleemphasis{econserve} quantity should remain
constant over time even though \sphinxstyleemphasis{etotal} may change.

\sphinxAtStartPar
In LAMMPS, the simulation box can be defined as orthogonal or
triclinic (non\sphinxhyphen{}orthogonal).  See the {\hyperref[\detokenize{Howto_triclinic::doc}]{\sphinxcrossref{\DUrole{doc}{Howto\_triclinic}}}} doc page for a detailed explanation of orthogonal,
restricted triclinic, and general triclinic simulation boxes and how
LAMMPS rotates a general triclinic box to be restricted triclinic
internally.

\sphinxAtStartPar
The \sphinxstyleemphasis{lx}, \sphinxstyleemphasis{ly}, \sphinxstyleemphasis{lz} keywords are the extent of the simulation box in
each dimension.  The \sphinxstyleemphasis{xlo}, \sphinxstyleemphasis{xhi}, \sphinxstyleemphasis{ylo}, \sphinxstyleemphasis{yhi}, \sphinxstyleemphasis{zlo}, \sphinxstyleemphasis{zhi} keywords
are the lower and upper bounds of the simulation box in each dimension.
I.e. \sphinxstyleemphasis{lx} = \sphinxstyleemphasis{xhi} \sphinxhyphen{} \sphinxstyleemphasis{xlo}).  These 9 values are the same for all 3 kinds
of boxes.  I.e. for a restricted triclinic box, they are the values as
if the box were not tilted.  For a general triclinic box, they are the
values after it is internally rotated to be a restricted triclinic box.

\sphinxAtStartPar
The \sphinxstyleemphasis{xy}, \sphinxstyleemphasis{xz}, \sphinxstyleemphasis{yz} are the current tilt factors for a triclinic box.
They are the same for restricted and general triclinic boxes.

\sphinxAtStartPar
The \sphinxstyleemphasis{avecx}, \sphinxstyleemphasis{avecy}, \sphinxstyleemphasis{avecz}, \sphinxstyleemphasis{bvecx}, \sphinxstyleemphasis{bvecy}, \sphinxstyleemphasis{bvecz}, \sphinxstyleemphasis{cvecx},
\sphinxstyleemphasis{cvecy}, \sphinxstyleemphasis{cvecz} are the components of the 3 edge vectors of the
current general simulation box.  If it is an orthogonal box the
vectors are along the x, y, z coordinate axes.  If it is a restricted
triclinic box, the \sphinxstylestrong{A} vector is along the x axis, the \sphinxstylestrong{B} vector
is in the xy plane with a +y coordinate, and the \sphinxstylestrong{C} vector has a +z
coordinate, as explained on the {\hyperref[\detokenize{Howto_triclinic::doc}]{\sphinxcrossref{\DUrole{doc}{Howto\_triclinic}}}} doc page.  If the {\hyperref[\detokenize{thermo_modify::doc}]{\sphinxcrossref{\DUrole{doc}{thermo\_modify
triclinic/general}}}} option is set then they are the
\sphinxstylestrong{A}, \sphinxstylestrong{B}, \sphinxstylestrong{C} vector which define the general triclinic box.

\sphinxAtStartPar
The \sphinxstyleemphasis{cella}, \sphinxstyleemphasis{cellb}, \sphinxstyleemphasis{cellc}, \sphinxstyleemphasis{cellalpha}, \sphinxstyleemphasis{cellbeta}, \sphinxstyleemphasis{cellgamma}
keywords correspond to the usual crystallographic quantities that
define the periodic simulation box of a crystalline system.  See the
{\hyperref[\detokenize{Howto_triclinic::doc}]{\sphinxcrossref{\DUrole{doc}{Howto triclinic}}}} page for a precise definition
of these quantities in terms of the LAMMPS representation of a
restricted triclinic simulation box via \sphinxstyleemphasis{lx}, \sphinxstyleemphasis{ly}, \sphinxstyleemphasis{lz}, \sphinxstyleemphasis{yz}, \sphinxstyleemphasis{xz},
\sphinxstyleemphasis{xy}.

\sphinxAtStartPar
The \sphinxstyleemphasis{pxx,pyy,pzz,pxy,pxz,pyz} keywords are the 6 components of the
symmetric pressure tensor for the system.  See the {\hyperref[\detokenize{compute_pressure::doc}]{\sphinxcrossref{\DUrole{doc}{compute
pressure}}}} command doc page for details of how it is
calculated.

\sphinxAtStartPar
If the {\hyperref[\detokenize{thermo_modify::doc}]{\sphinxcrossref{\DUrole{doc}{thermo\_modify triclinic/general}}}} option
is set then the 6 components will be output as values consistent with
the orientation of the general triclinic box relative to the standard
xyz coordinate axes.  If this keyword is not used, the values will be
consistent with the orientation of the restricted triclinic box (which
aligns with the xyz coordinate axes).  As explained on the
{\hyperref[\detokenize{Howto_triclinic::doc}]{\sphinxcrossref{\DUrole{doc}{Howto\_triclinic}}}} doc page, even if the
simulation box is created as a general triclinic box, internally
LAMMPS uses a restricted triclinic box.

\sphinxAtStartPar
Note that because the pressure tensor components are computed using
force vectors and atom coordinates, both of which are rotated in the
general versus restricted triclinic representation, the values will
typically be different for the two cases.

\sphinxAtStartPar
The \sphinxstyleemphasis{fmax} and \sphinxstyleemphasis{fnorm} keywords are useful for monitoring the progress
of an {\hyperref[\detokenize{minimize::doc}]{\sphinxcrossref{\DUrole{doc}{energy minimization}}}}.  The \sphinxstyleemphasis{fmax} keyword
calculates the maximum force in any dimension on any atom in the
system, or the infinity\sphinxhyphen{}norm of the force vector for the system.  The
\sphinxstyleemphasis{fnorm} keyword calculates the 2\sphinxhyphen{}norm or length of the force vector.

\sphinxAtStartPar
The \sphinxstyleemphasis{nbuild} and \sphinxstyleemphasis{ndanger} keywords are useful for monitoring neighbor
list builds during a run.  Note that both these values are also
printed with the end\sphinxhyphen{}of\sphinxhyphen{}run statistics.  The \sphinxstyleemphasis{nbuild} keyword is the
number of re\sphinxhyphen{}builds during the current run.  The \sphinxstyleemphasis{ndanger} keyword is
the number of re\sphinxhyphen{}builds that LAMMPS considered potentially
“dangerous”.  If atom movement triggered neighbor list rebuilding (see
the {\hyperref[\detokenize{neigh_modify::doc}]{\sphinxcrossref{\DUrole{doc}{neigh\_modify}}}} command), then dangerous
reneighborings are those that were triggered on the first timestep
atom movement was checked for.  If this count is non\sphinxhyphen{}zero you may wish
to reduce the delay factor to ensure no force interactions are missed
by atoms moving beyond the neighbor skin distance before a rebuild
takes place.


\bigskip\hrule\bigskip


\sphinxAtStartPar
For output values from a compute or fix or variable, the bracketed
index I used to index a vector, as in \sphinxstyleemphasis{c\_ID{[}I{]}} or \sphinxstyleemphasis{f\_ID{[}I{]}} or
\sphinxstyleemphasis{v\_name{[}I{]}}, can be specified using a wildcard asterisk with the index
to effectively specify multiple values.  This takes the form “*” or
“*n” or “n*” or “m*n”.  If N = the size of the vector, then an
asterisk with no numeric values means all indices from 1 to N.  A
leading asterisk means all indices from 1 to n (inclusive).  A
trailing asterisk means all indices from n to N (inclusive).  A middle
asterisk means all indices from m to n (inclusive).

\sphinxAtStartPar
Using a wildcard is the same as if the individual elements of the
vector had been listed one by one.  E.g. these 2 thermo\_style commands
are equivalent, since the {\hyperref[\detokenize{compute_temp::doc}]{\sphinxcrossref{\DUrole{doc}{compute temp}}}} command
creates a global vector with 6 values.

\begin{sphinxVerbatim}[commandchars=\\\{\}]
\PYG{k}{compute }\PYG{n+nv+nvIdentifier}{myTemp}\PYG{+w}{ }\PYG{n+nv+nvIdentifier}{all}\PYG{+w}{ }\PYG{n}{temp}
\PYG{k}{thermo\PYGZus{}style}\PYG{+w}{ }\PYG{n}{custom}\PYG{+w}{ }\PYG{n}{step}\PYG{+w}{ }\PYG{n}{temp}\PYG{+w}{ }\PYG{n}{etotal}\PYG{+w}{ }\PYG{n}{c\PYGZus{}myTemp[}\PYG{o}{*}\PYG{n}{]}
\PYG{k}{thermo\PYGZus{}style}\PYG{+w}{ }\PYG{n}{custom}\PYG{+w}{ }\PYG{n}{step}\PYG{+w}{ }\PYG{n}{temp}\PYG{+w}{ }\PYG{n}{etotal}\PYG{+w}{ }\PYG{o}{\PYGZam{}}
\PYG{+w}{             }\PYG{n}{c\PYGZus{}myTemp[1]}\PYG{+w}{ }\PYG{n}{c\PYGZus{}myTemp[2]}\PYG{+w}{ }\PYG{n}{c\PYGZus{}myTemp[3]}\PYG{+w}{ }\PYG{o}{\PYGZam{}}
\PYG{+w}{             }\PYG{n}{c\PYGZus{}myTemp[4]}\PYG{+w}{ }\PYG{n}{c\PYGZus{}myTemp[5]}\PYG{+w}{ }\PYG{n}{c\PYGZus{}myTemp[6]}
\end{sphinxVerbatim}

\begin{sphinxadmonition}{note}{Note:}
\sphinxAtStartPar
For a vector\sphinxhyphen{}style variable, only the wildcard forms “*n” or
“m*n” are allowed.  You must specify the upper bound, because
vector\sphinxhyphen{}style variable lengths are not determined until the variable
is evaluated.  If n is specified larger than the vector length
turns out to be, zeroes are output for missing vector values.
\end{sphinxadmonition}


\bigskip\hrule\bigskip


\sphinxAtStartPar
The \sphinxstyleemphasis{c\_ID} and \sphinxstyleemphasis{c\_ID{[}I{]}} and \sphinxstyleemphasis{c\_ID{[}I{]}{[}J{]}} keywords allow global values
calculated by a compute to be output.  As discussed on the
{\hyperref[\detokenize{compute::doc}]{\sphinxcrossref{\DUrole{doc}{compute}}}} doc page, computes can calculate global,
per\sphinxhyphen{}atom, local, and per\sphinxhyphen{}grid values.  Only global values can be
referenced by this command.  However, per\sphinxhyphen{}atom compute values for an
individual atom can be referenced in a {\hyperref[\detokenize{variable::doc}]{\sphinxcrossref{\DUrole{doc}{equal\sphinxhyphen{}style variable}}}} and the variable referenced by thermo\_style custom, as
discussed below.  See the discussion above for how the I in \sphinxstyleemphasis{c\_ID{[}I{]}}
can be specified with a wildcard asterisk to effectively specify
multiple values from a global compute vector.

\sphinxAtStartPar
The ID in the keyword should be replaced by the actual ID of a compute
that has been defined elsewhere in the input script.  See the
{\hyperref[\detokenize{compute::doc}]{\sphinxcrossref{\DUrole{doc}{compute}}}} command for details.  If the compute
calculates a global scalar, vector, or array, then the keyword formats
with 0, 1, or 2 brackets will reference a scalar value from the
compute.

\sphinxAtStartPar
Note that some computes calculate “intensive” global quantities like
temperature; others calculate “extensive” global quantities like
kinetic energy that are summed over all atoms in the compute group.
Intensive quantities are printed directly without normalization by
thermo\_style custom.  Extensive quantities may be normalized by the
total number of atoms in the simulation (NOT the number of atoms in
the compute group) when output, depending on the {\hyperref[\detokenize{thermo_modify::doc}]{\sphinxcrossref{\DUrole{doc}{thermo\_modify
norm}}}} option being used.

\sphinxAtStartPar
The \sphinxstyleemphasis{f\_ID} and \sphinxstyleemphasis{f\_ID{[}I{]}} and \sphinxstyleemphasis{f\_ID{[}I{]}{[}J{]}} keywords allow global values
calculated by a fix to be output.  As discussed on the {\hyperref[\detokenize{fix::doc}]{\sphinxcrossref{\DUrole{doc}{fix}}}} doc page, fixes can calculate global, per\sphinxhyphen{}atom, local, and
per\sphinxhyphen{}grid values.  Only global values can be referenced by this
command.  However, per\sphinxhyphen{}atom fix values can be referenced for an
individual atom in a {\hyperref[\detokenize{variable::doc}]{\sphinxcrossref{\DUrole{doc}{equal\sphinxhyphen{}style variable}}}} and the
variable referenced by thermo\_style custom, as discussed below.  See
the discussion above for how the I in \sphinxstyleemphasis{f\_ID{[}I{]}} can be specified with
a wildcard asterisk to effectively specify multiple values from a
global fix vector.

\sphinxAtStartPar
The ID in the keyword should be replaced by the actual ID of a fix
that has been defined elsewhere in the input script.  See the
{\hyperref[\detokenize{fix::doc}]{\sphinxcrossref{\DUrole{doc}{fix}}}} command for details.  If the fix calculates a global
scalar, vector, or array, then the keyword formats with 0, 1, or 2
brackets will reference a scalar value from the fix.

\sphinxAtStartPar
Note that some fixes calculate “intensive” global quantities like
timestep size; others calculate “extensive” global quantities like
energy that are summed over all atoms in the fix group.  Intensive
quantities are printed directly without normalization by thermo\_style
custom.  Extensive quantities may be normalized by the total number of
atoms in the simulation (NOT the number of atoms in the fix group)
when output, depending on the {\hyperref[\detokenize{thermo_modify::doc}]{\sphinxcrossref{\DUrole{doc}{thermo\_modify norm}}}} option being used.

\sphinxAtStartPar
The \sphinxstyleemphasis{v\_name} keyword allow the current value of a variable to be
output.  The name in the keyword should be replaced by the variable
name that has been defined elsewhere in the input script.  Only
equal\sphinxhyphen{}style and vector\sphinxhyphen{}style variables can be referenced; the latter
requires a bracketed term to specify the Ith element of the vector
calculated by the variable.  However, an equal\sphinxhyphen{}style variable can use
an atom\sphinxhyphen{}style variable in its formula indexed by the ID of an
individual atom.  This is a way to output a specific atom’s per\sphinxhyphen{}atom
coordinates or other per\sphinxhyphen{}atom properties in thermo output.  See the
{\hyperref[\detokenize{variable::doc}]{\sphinxcrossref{\DUrole{doc}{variable}}}} command for details.  Note that variables
of style \sphinxstyleemphasis{equal} and \sphinxstyleemphasis{vector} and \sphinxstyleemphasis{atom} define a formula which can
reference per\sphinxhyphen{}atom properties or thermodynamic keywords, or they can
invoke other computes, fixes, or variables when evaluated, so this is
a very general means of creating thermodynamic output.

\sphinxAtStartPar
Note that equal\sphinxhyphen{}style and vector\sphinxhyphen{}style variables are assumed to
produce “intensive” global quantities, which are thus printed as\sphinxhyphen{}is,
without normalization by thermo\_style custom.  You can include a
division by “natoms” in the variable formula if this is not the case.


\bigskip\hrule\bigskip



\subsection{Restrictions}
\label{\detokenize{thermo_style:restrictions}}
\sphinxAtStartPar
This command must come after the simulation box is defined by a
{\hyperref[\detokenize{read_data::doc}]{\sphinxcrossref{\DUrole{doc}{read\_data}}}}, {\hyperref[\detokenize{read_restart::doc}]{\sphinxcrossref{\DUrole{doc}{read\_restart}}}}, or
{\hyperref[\detokenize{create_box::doc}]{\sphinxcrossref{\DUrole{doc}{create\_box}}}} command.


\subsection{Related commands}
\label{\detokenize{thermo_style:related-commands}}
\sphinxAtStartPar
{\hyperref[\detokenize{thermo::doc}]{\sphinxcrossref{\DUrole{doc}{thermo}}}}, {\hyperref[\detokenize{thermo_modify::doc}]{\sphinxcrossref{\DUrole{doc}{thermo\_modify}}}},
{\hyperref[\detokenize{fix_modify::doc}]{\sphinxcrossref{\DUrole{doc}{fix\_modify}}}}, {\hyperref[\detokenize{compute_temp::doc}]{\sphinxcrossref{\DUrole{doc}{compute temp}}}},
{\hyperref[\detokenize{compute_pressure::doc}]{\sphinxcrossref{\DUrole{doc}{compute pressure}}}}


\subsection{Default}
\label{\detokenize{thermo_style:default}}
\begin{sphinxVerbatim}[commandchars=\\\{\}]
\PYG{k}{thermo\PYGZus{}style}\PYG{+w}{ }\PYG{n}{one}
\end{sphinxVerbatim}

\sphinxstepscope

\index{third\_order@\spxentry{third\_order}}\ignorespaces 
\index{third\_order/kk@\spxentry{third\_order/kk}}\ignorespaces 

\section{third\_order command}
\label{\detokenize{third_order:third-order-command}}\label{\detokenize{third_order:index-1}}\label{\detokenize{third_order:index-0}}\label{\detokenize{third_order::doc}}
\sphinxAtStartPar
Accelerator Variant: third\_order/kk


\subsection{Syntax}
\label{\detokenize{third_order:syntax}}
\begin{sphinxVerbatim}[commandchars=\\\{\}]
\PYG{k}{third\PYGZus{}order}\PYG{+w}{ }\PYG{n}{group}\PYG{o}{\PYGZhy{}}\PYG{n}{ID}\PYG{+w}{ }\PYG{n}{style}\PYG{+w}{ }\PYG{n}{delta}\PYG{+w}{ }\PYG{n}{args}\PYG{+w}{ }\PYG{n}{keyword}\PYG{+w}{ }\PYG{n}{value}\PYG{+w}{ }\PYG{n}{...}
\end{sphinxVerbatim}
\begin{itemize}
\item {} 
\sphinxAtStartPar
group\sphinxhyphen{}ID = ID of group of atoms to displace

\item {} 
\sphinxAtStartPar
style = \sphinxstyleemphasis{regular} or \sphinxstyleemphasis{eskm}

\item {} 
\sphinxAtStartPar
delta = finite different displacement length (distance units)

\item {} 
\sphinxAtStartPar
one or more keyword/arg pairs may be appended
\begin{sphinxalltt}
keyword = \sphinxstyleemphasis{file} or \sphinxstyleemphasis{binary}
  \sphinxstyleemphasis{file} name = name of output file for the third order tensor
  \sphinxstyleemphasis{binary} arg = \sphinxstyleemphasis{yes} or \sphinxstyleemphasis{no} or \sphinxstyleemphasis{gzip}
\end{sphinxalltt}

\end{itemize}


\subsection{Examples}
\label{\detokenize{third_order:examples}}
\begin{sphinxVerbatim}[commandchars=\\\{\}]
\PYG{k}{third\PYGZus{}order}\PYG{+w}{ }\PYG{l+m}{1}\PYG{+w}{ }\PYG{n}{regular}\PYG{+w}{ }\PYG{l+m}{0.000001}
\PYG{k}{third\PYGZus{}order}\PYG{+w}{ }\PYG{l+m}{1}\PYG{+w}{ }\PYG{n}{eskm}\PYG{+w}{ }\PYG{l+m}{0.000001}
\PYG{k}{third\PYGZus{}order}\PYG{+w}{ }\PYG{l+m}{3}\PYG{+w}{ }\PYG{n}{regular}\PYG{+w}{ }\PYG{l+m}{0.00004}\PYG{+w}{ }\PYG{n}{file}\PYG{+w}{ }\PYG{n}{third\PYGZus{}order.dat}
\PYG{k}{third\PYGZus{}order}\PYG{+w}{ }\PYG{l+m}{5}\PYG{+w}{ }\PYG{n}{eskm}\PYG{+w}{ }\PYG{l+m}{0.00000001}\PYG{+w}{ }\PYG{n}{file}\PYG{+w}{ }\PYG{n}{third\PYGZus{}order.dat}\PYG{+w}{ }\PYG{n}{binary}\PYG{+w}{ }\PYG{n}{yes}
\end{sphinxVerbatim}


\subsection{Description}
\label{\detokenize{third_order:description}}
\sphinxAtStartPar
Calculate the third order force constant tensor by finite difference of the selected group,
\begin{equation*}
\begin{split}\Phi^{\alpha\beta\gamma}_{ijk} = \frac{\partial^3 U}{\partial x_{i,\alpha} \partial x_{j,\beta} \partial x_{k, \gamma}}\end{split}
\end{equation*}
\sphinxAtStartPar
where Phi is the third order force constant tensor.

\sphinxAtStartPar
The output of the command is the tensor, three elements at a time. The
three elements correspond to the three gamma elements for a specific i/alpha/j/beta/k.
The initial five numbers are i, alpha, j, beta, and k respectively.

\sphinxAtStartPar
If the style eskm is selected, the tensor will be using energy units of 10 J/mol.
These units conform to eskm style from the dynamical\_matrix command, which
will simplify operations using dynamical matrices with third order tensors.


\bigskip\hrule\bigskip


\sphinxAtStartPar
Styles with a \sphinxstyleemphasis{gpu}, \sphinxstyleemphasis{intel}, \sphinxstyleemphasis{kk}, \sphinxstyleemphasis{omp}, or \sphinxstyleemphasis{opt} suffix are
functionally the same as the corresponding style without the suffix.
They have been optimized to run faster, depending on your available
hardware, as discussed on the {\hyperref[\detokenize{Speed_packages::doc}]{\sphinxcrossref{\DUrole{doc}{Accelerator packages}}}}
page.  The accelerated styles take the same arguments and should
produce the same results, except for round\sphinxhyphen{}off and precision issues.

\sphinxAtStartPar
These accelerated styles are part of the GPU, INTEL, KOKKOS,
OPENMP, and OPT packages, respectively.  They are only enabled if
LAMMPS was built with those packages.  See the {\hyperref[\detokenize{Build_package::doc}]{\sphinxcrossref{\DUrole{doc}{Build package}}}} page for more info.

\sphinxAtStartPar
You can specify the accelerated styles explicitly in your input script
by including their suffix, or you can use the {\hyperref[\detokenize{Run_options::doc}]{\sphinxcrossref{\DUrole{doc}{\sphinxhyphen{}suffix command\sphinxhyphen{}line switch}}}} when you invoke LAMMPS, or you can use the
{\hyperref[\detokenize{suffix::doc}]{\sphinxcrossref{\DUrole{doc}{suffix}}}} command in your input script.

\sphinxAtStartPar
See the {\hyperref[\detokenize{Speed_packages::doc}]{\sphinxcrossref{\DUrole{doc}{Accelerator packages}}}} page for more
instructions on how to use the accelerated styles effectively.


\bigskip\hrule\bigskip



\subsection{Restrictions}
\label{\detokenize{third_order:restrictions}}
\sphinxAtStartPar
The command collects a 9 times the number of atoms in the group on every single MPI rank,
so the memory requirements can be very significant for large systems.

\sphinxAtStartPar
This command is part of the PHONON package.  It is only enabled if
LAMMPS was built with that package.  See the {\hyperref[\detokenize{Build_package::doc}]{\sphinxcrossref{\DUrole{doc}{Build package}}}} page for more info.


\subsection{Related commands}
\label{\detokenize{third_order:related-commands}}
\sphinxAtStartPar
{\hyperref[\detokenize{fix_phonon::doc}]{\sphinxcrossref{\DUrole{doc}{fix phonon}}}} {\hyperref[\detokenize{dynamical_matrix::doc}]{\sphinxcrossref{\DUrole{doc}{dynamical\_matrix}}}}


\subsection{Default}
\label{\detokenize{third_order:default}}
\sphinxAtStartPar
The default settings are file = “third\_order.dat”, binary = no

\sphinxstepscope

\index{timer@\spxentry{timer}}\ignorespaces 

\section{timer command}
\label{\detokenize{timer:timer-command}}\label{\detokenize{timer:index-0}}\label{\detokenize{timer::doc}}

\subsection{Syntax}
\label{\detokenize{timer:syntax}}
\begin{sphinxVerbatim}[commandchars=\\\{\}]
\PYG{k}{timer}\PYG{+w}{ }\PYG{n}{args}
\end{sphinxVerbatim}
\begin{itemize}
\item {} 
\sphinxAtStartPar
\sphinxstyleemphasis{args} = one or more of \sphinxstyleemphasis{off} or \sphinxstyleemphasis{loop} or \sphinxstyleemphasis{normal} or \sphinxstyleemphasis{full} or \sphinxstyleemphasis{sync} or \sphinxstyleemphasis{nosync} or \sphinxstyleemphasis{timeout} or \sphinxstyleemphasis{every}

\end{itemize}
\begin{sphinxalltt}
\sphinxstyleemphasis{off} = do not collect or print any timing information
\sphinxstyleemphasis{loop} = collect only the total time for the simulation loop
\sphinxstyleemphasis{normal} = collect timer information broken down by sections (default)
\sphinxstyleemphasis{full} = like \sphinxstyleemphasis{normal} but also include CPU and thread utilization
\sphinxstyleemphasis{sync} = explicitly synchronize MPI tasks between sections
\sphinxstyleemphasis{nosync} = do not synchronize MPI tasks between sections (default)
\sphinxstyleemphasis{timeout} elapse = set wall time limit to \sphinxstyleemphasis{elapse}
\sphinxstyleemphasis{every} Ncheck = perform timeout check every \sphinxstyleemphasis{Ncheck} steps
\end{sphinxalltt}


\subsection{Examples}
\label{\detokenize{timer:examples}}
\begin{sphinxVerbatim}[commandchars=\\\{\}]
\PYG{k}{timer}\PYG{+w}{ }\PYG{n}{full}\PYG{+w}{ }\PYG{n}{sync}
\PYG{k}{timer}\PYG{+w}{ }\PYG{n}{timeout}\PYG{+w}{ }\PYG{l+m}{2:00:00}\PYG{+w}{ }\PYG{n}{every}\PYG{+w}{ }\PYG{l+m}{100}
\PYG{k}{timer}\PYG{+w}{ }\PYG{n}{loop}
\end{sphinxVerbatim}


\subsection{Description}
\label{\detokenize{timer:description}}
\sphinxAtStartPar
Select the level of detail at which LAMMPS performs its CPU timings.
Multiple keywords can be specified with the \sphinxstyleemphasis{timer} command.  For
keywords that are mutually exclusive, the last one specified takes
precedence.

\sphinxAtStartPar
During a simulation run LAMMPS collects information about how much
time is spent in different sections of the code and thus can provide
information for determining performance and load imbalance problems.
This can be done at different levels of detail and accuracy.  For more
information about the timing output, see the {\hyperref[\detokenize{Run_output::doc}]{\sphinxcrossref{\DUrole{doc}{Run output}}}} doc page.

\sphinxAtStartPar
The \sphinxstyleemphasis{off} setting will turn all time measurements off. The \sphinxstyleemphasis{loop}
setting will only measure the total time for a run and not collect any
detailed per section information.  With the \sphinxstyleemphasis{normal} setting, timing
information for portions of the timestep (pairwise calculations,
neighbor list construction, output, etc) are collected as well as
information about load imbalances for those sections across
processors.  The \sphinxstyleemphasis{full} setting adds information about CPU
utilization and thread utilization, when multi\sphinxhyphen{}threading is enabled.

\sphinxAtStartPar
With the \sphinxstyleemphasis{sync} setting, all MPI tasks are synchronized at each timer
call which measures load imbalance for each section more accurately,
though it can also slow down the simulation by prohibiting overlapping
independent computations on different MPI ranks  Using the \sphinxstyleemphasis{nosync}
setting (which is the default) turns this synchronization off.

\sphinxAtStartPar
With the \sphinxstyleemphasis{timeout} keyword a wall time limit can be imposed, that
affects the {\hyperref[\detokenize{run::doc}]{\sphinxcrossref{\DUrole{doc}{run}}}} and {\hyperref[\detokenize{minimize::doc}]{\sphinxcrossref{\DUrole{doc}{minimize}}}} commands.
This can be convenient when calculations have to comply with execution
time limits, e.g. when running under a batch system when you want to
maximize the utilization of the batch time slot, especially for runs
where the time per timestep varies much and thus it becomes difficult
to predict how many steps a simulation can perform for a given wall time
limit. This also applies for difficult to converge minimizations.
The timeout \sphinxstyleemphasis{elapse} value should be somewhat smaller than the maximum
wall time requested from the batch system, as there is usually
some overhead to launch jobs, and it is advisable to write
out a restart after terminating a run due to a timeout.

\sphinxAtStartPar
The timeout timer starts when the command is issued. When the time
limit is reached, the run or energy minimization will exit on the
next step or iteration that is a multiple of the \sphinxstyleemphasis{Ncheck} value
which can be set with the \sphinxstyleemphasis{every} keyword. Default is checking
every 10 steps. After the timer timeout has expired all subsequent
run or minimize commands in the input script will be skipped.
The remaining time or timer status can be accessed with the
{\hyperref[\detokenize{thermo_style::doc}]{\sphinxcrossref{\DUrole{doc}{thermo}}}} variable \sphinxstyleemphasis{timeremain}, which will be
zero, if the timeout is inactive (default setting), it will be
negative, if the timeout time is expired and positive if there
is time remaining and in this case the value of the variable are
the number of seconds remaining.

\sphinxAtStartPar
When the \sphinxstyleemphasis{timeout} key word is used a second time, the timer is
restarted with a new time limit. The timeout \sphinxstyleemphasis{elapse} value can
be specified as \sphinxstyleemphasis{off} or \sphinxstyleemphasis{unlimited} to impose a no timeout condition
(which is the default).  The \sphinxstyleemphasis{elapse} setting can be specified as
a single number for seconds, two numbers separated by a colon (MM:SS)
for minutes and seconds, or as three numbers separated by colons for
hours, minutes, and seconds (H:MM:SS).

\sphinxAtStartPar
The \sphinxstyleemphasis{every} keyword sets how frequently during a run or energy
minimization the wall clock will be checked.  This check count applies
to the outer iterations or time steps during minimizations or {\hyperref[\detokenize{run_style::doc}]{\sphinxcrossref{\DUrole{doc}{r\sphinxhyphen{}RESPA runs}}}}, respectively.  Checking for timeout too often,
can slow a calculation down.  Checking too infrequently can make the
timeout measurement less accurate, with the run being stopped later
than desired.

\begin{sphinxadmonition}{note}{Note:}
\sphinxAtStartPar
Using the \sphinxstyleemphasis{full} and \sphinxstyleemphasis{sync} options provides the most detailed
and accurate timing information, but can also have a negative
performance impact due to the overhead of the many required system
calls. It is thus recommended to use these settings only when testing
tests to identify performance bottlenecks. For calculations with few
atoms or a very large number of processors, even the \sphinxstyleemphasis{normal} setting
can have a measurable negative performance impact. In those cases you
can just use the \sphinxstyleemphasis{loop} or \sphinxstyleemphasis{off} setting.
\end{sphinxadmonition}


\subsection{Restrictions}
\label{\detokenize{timer:restrictions}}\begin{quote}

\sphinxAtStartPar
none
\end{quote}


\subsection{Related commands}
\label{\detokenize{timer:related-commands}}
\sphinxAtStartPar
{\hyperref[\detokenize{run::doc}]{\sphinxcrossref{\DUrole{doc}{run post no}}}}, {\hyperref[\detokenize{kspace_modify::doc}]{\sphinxcrossref{\DUrole{doc}{kspace\_modify fftbench}}}}


\subsection{Default}
\label{\detokenize{timer:default}}
\begin{sphinxVerbatim}[commandchars=\\\{\}]
\PYG{k}{timer}\PYG{+w}{ }\PYG{n}{normal}\PYG{+w}{ }\PYG{n}{nosync}
\PYG{k}{timer}\PYG{+w}{ }\PYG{n}{timeout}\PYG{+w}{ }\PYG{n}{off}
\PYG{k}{timer}\PYG{+w}{ }\PYG{n}{every}\PYG{+w}{ }\PYG{l+m}{10}
\end{sphinxVerbatim}

\sphinxstepscope

\index{timestep@\spxentry{timestep}}\ignorespaces 

\section{timestep command}
\label{\detokenize{timestep:timestep-command}}\label{\detokenize{timestep:index-0}}\label{\detokenize{timestep::doc}}

\subsection{Syntax}
\label{\detokenize{timestep:syntax}}
\begin{sphinxVerbatim}[commandchars=\\\{\}]
\PYG{k}{timestep}\PYG{+w}{ }\PYG{n}{dt}
\end{sphinxVerbatim}
\begin{itemize}
\item {} 
\sphinxAtStartPar
dt = timestep size (time units)

\end{itemize}


\subsection{Examples}
\label{\detokenize{timestep:examples}}
\begin{sphinxVerbatim}[commandchars=\\\{\}]
\PYG{k}{timestep}\PYG{+w}{ }\PYG{l+m}{2.0}
\PYG{k}{timestep}\PYG{+w}{ }\PYG{l+m}{0.003}
\end{sphinxVerbatim}


\subsection{Description}
\label{\detokenize{timestep:description}}
\sphinxAtStartPar
Set the timestep size for subsequent molecular dynamics simulations.
See the {\hyperref[\detokenize{units::doc}]{\sphinxcrossref{\DUrole{doc}{units}}}} command for the time units associated with
each choice of units that LAMMPS supports.

\sphinxAtStartPar
The default value for the timestep size also depends on the choice of
units for the simulation; see the default values below.

\sphinxAtStartPar
When the {\hyperref[\detokenize{run_style::doc}]{\sphinxcrossref{\DUrole{doc}{run style}}}} is \sphinxstyleemphasis{respa}, dt is the timestep for
the outer loop (largest) timestep.


\subsection{Restrictions}
\label{\detokenize{timestep:restrictions}}\begin{quote}

\sphinxAtStartPar
none
\end{quote}


\subsection{Related commands}
\label{\detokenize{timestep:related-commands}}
\sphinxAtStartPar
{\hyperref[\detokenize{fix_dt_reset::doc}]{\sphinxcrossref{\DUrole{doc}{fix dt/reset}}}}, {\hyperref[\detokenize{run::doc}]{\sphinxcrossref{\DUrole{doc}{run}}}},
{\hyperref[\detokenize{run_style::doc}]{\sphinxcrossref{\DUrole{doc}{run\_style}}}} respa, {\hyperref[\detokenize{units::doc}]{\sphinxcrossref{\DUrole{doc}{units}}}}


\subsection{Default}
\label{\detokenize{timestep:default}}

\begin{savenotes}\sphinxattablestart
\sphinxthistablewithglobalstyle
\centering
\begin{tabulary}{\linewidth}[t]{TTT}
\sphinxtoprule
\sphinxtableatstartofbodyhook
\sphinxAtStartPar
choice of {\hyperref[\detokenize{units::doc}]{\sphinxcrossref{\DUrole{doc}{units}}}}
&
\sphinxAtStartPar
time units
&
\sphinxAtStartPar
default timestep size
\\
\sphinxhline
\sphinxAtStartPar
lj
&
\sphinxAtStartPar
\(\tau\)
&
\sphinxAtStartPar
0.005 \(\tau\)
\\
\sphinxhline
\sphinxAtStartPar
real
&
\sphinxAtStartPar
fs
&
\sphinxAtStartPar
1.0 fs
\\
\sphinxhline
\sphinxAtStartPar
metal
&
\sphinxAtStartPar
ps
&
\sphinxAtStartPar
0.001 ps
\\
\sphinxhline
\sphinxAtStartPar
si
&
\sphinxAtStartPar
s
&
\sphinxAtStartPar
1.0e\sphinxhyphen{}8 s (10 ns)
\\
\sphinxhline
\sphinxAtStartPar
cgs
&
\sphinxAtStartPar
s
&
\sphinxAtStartPar
1.0e\sphinxhyphen{}8 s (10 ns)
\\
\sphinxhline
\sphinxAtStartPar
electron
&
\sphinxAtStartPar
fs
&
\sphinxAtStartPar
0.001 fs
\\
\sphinxhline
\sphinxAtStartPar
micro
&
\sphinxAtStartPar
\(\mu\)s
&
\sphinxAtStartPar
2.0 \(\mu\)s
\\
\sphinxhline
\sphinxAtStartPar
nano
&
\sphinxAtStartPar
ns
&
\sphinxAtStartPar
0.00045 ns
\\
\sphinxbottomrule
\end{tabulary}
\sphinxtableafterendhook\par
\sphinxattableend\end{savenotes}

\sphinxstepscope

\index{uncompute@\spxentry{uncompute}}\ignorespaces 

\section{uncompute command}
\label{\detokenize{uncompute:uncompute-command}}\label{\detokenize{uncompute:index-0}}\label{\detokenize{uncompute::doc}}

\subsection{Syntax}
\label{\detokenize{uncompute:syntax}}
\begin{sphinxVerbatim}[commandchars=\\\{\}]
\PYG{k}{uncompute }\PYG{n+nv+nvIdentifier}{compute\PYGZhy{}ID}
\end{sphinxVerbatim}
\begin{itemize}
\item {} 
\sphinxAtStartPar
compute\sphinxhyphen{}ID = ID of a previously defined compute

\end{itemize}


\subsection{Examples}
\label{\detokenize{uncompute:examples}}
\begin{sphinxVerbatim}[commandchars=\\\{\}]
\PYG{k}{uncompute }\PYG{n+nv+nvIdentifier}{2}
\PYG{k}{uncompute }\PYG{n+nv+nvIdentifier}{lower\PYGZhy{}boundary}
\end{sphinxVerbatim}


\subsection{Description}
\label{\detokenize{uncompute:description}}
\sphinxAtStartPar
Delete a compute that was previously defined with a {\hyperref[\detokenize{compute::doc}]{\sphinxcrossref{\DUrole{doc}{compute}}}}
command.  This also wipes out any additional changes made to the compute
via the {\hyperref[\detokenize{compute_modify::doc}]{\sphinxcrossref{\DUrole{doc}{compute\_modify}}}} command.


\subsection{Restrictions}
\label{\detokenize{uncompute:restrictions}}\begin{quote}

\sphinxAtStartPar
none
\end{quote}


\subsection{Related commands}
\label{\detokenize{uncompute:related-commands}}
\sphinxAtStartPar
{\hyperref[\detokenize{compute::doc}]{\sphinxcrossref{\DUrole{doc}{compute}}}}


\subsection{Default}
\label{\detokenize{uncompute:default}}
\sphinxAtStartPar
none

\sphinxstepscope

\index{undump@\spxentry{undump}}\ignorespaces 

\section{undump command}
\label{\detokenize{undump:undump-command}}\label{\detokenize{undump:index-0}}\label{\detokenize{undump::doc}}

\subsection{Syntax}
\label{\detokenize{undump:syntax}}
\begin{sphinxVerbatim}[commandchars=\\\{\}]
\PYG{k}{undump }\PYG{n+nv+nvIdentifier}{dump\PYGZhy{}ID}
\end{sphinxVerbatim}
\begin{itemize}
\item {} 
\sphinxAtStartPar
dump\sphinxhyphen{}ID = ID of previously defined dump

\end{itemize}


\subsection{Examples}
\label{\detokenize{undump:examples}}
\begin{sphinxVerbatim}[commandchars=\\\{\}]
\PYG{k}{undump }\PYG{n+nv+nvIdentifier}{mine}
\PYG{k}{undump }\PYG{n+nv+nvIdentifier}{2}
\end{sphinxVerbatim}


\subsection{Description}
\label{\detokenize{undump:description}}
\sphinxAtStartPar
Turn off a previously defined dump so that it is no longer active.
This closes the file associated with the dump.


\subsection{Restrictions}
\label{\detokenize{undump:restrictions}}\begin{quote}

\sphinxAtStartPar
none
\end{quote}


\subsection{Related commands}
\label{\detokenize{undump:related-commands}}
\sphinxAtStartPar
{\hyperref[\detokenize{dump::doc}]{\sphinxcrossref{\DUrole{doc}{dump}}}}


\subsection{Default}
\label{\detokenize{undump:default}}
\sphinxAtStartPar
none

\sphinxstepscope

\index{unfix@\spxentry{unfix}}\ignorespaces 

\section{unfix command}
\label{\detokenize{unfix:unfix-command}}\label{\detokenize{unfix:index-0}}\label{\detokenize{unfix::doc}}

\subsection{Syntax}
\label{\detokenize{unfix:syntax}}
\begin{sphinxVerbatim}[commandchars=\\\{\}]
\PYG{k}{unfix }\PYG{n+nv+nvIdentifier}{fix\PYGZhy{}ID}
\end{sphinxVerbatim}
\begin{itemize}
\item {} 
\sphinxAtStartPar
fix\sphinxhyphen{}ID = ID of a previously defined fix

\end{itemize}


\subsection{Examples}
\label{\detokenize{unfix:examples}}
\begin{sphinxVerbatim}[commandchars=\\\{\}]
\PYG{k}{unfix }\PYG{n+nv+nvIdentifier}{2}
\PYG{k}{unfix }\PYG{n+nv+nvIdentifier}{lower\PYGZhy{}boundary}
\end{sphinxVerbatim}


\subsection{Description}
\label{\detokenize{unfix:description}}
\sphinxAtStartPar
Delete a fix that was previously defined with a {\hyperref[\detokenize{fix::doc}]{\sphinxcrossref{\DUrole{doc}{fix}}}}
command.  This also wipes out any additional changes made to the fix
via the {\hyperref[\detokenize{fix_modify::doc}]{\sphinxcrossref{\DUrole{doc}{fix\_modify}}}} command.


\subsection{Restrictions}
\label{\detokenize{unfix:restrictions}}\begin{quote}

\sphinxAtStartPar
none
\end{quote}


\subsection{Related commands}
\label{\detokenize{unfix:related-commands}}
\sphinxAtStartPar
{\hyperref[\detokenize{fix::doc}]{\sphinxcrossref{\DUrole{doc}{fix}}}}


\subsection{Default}
\label{\detokenize{unfix:default}}
\sphinxAtStartPar
none

\sphinxstepscope

\index{units@\spxentry{units}}\ignorespaces 

\section{units command}
\label{\detokenize{units:units-command}}\label{\detokenize{units:index-0}}\label{\detokenize{units::doc}}

\subsection{Syntax}
\label{\detokenize{units:syntax}}
\begin{sphinxVerbatim}[commandchars=\\\{\}]
\PYG{k}{units}\PYG{+w}{ }\PYG{n}{style}
\end{sphinxVerbatim}
\begin{itemize}
\item {} 
\sphinxAtStartPar
style = \sphinxstyleemphasis{lj} or \sphinxstyleemphasis{real} or \sphinxstyleemphasis{metal} or \sphinxstyleemphasis{si} or \sphinxstyleemphasis{cgs} or \sphinxstyleemphasis{electron} or \sphinxstyleemphasis{micro} or \sphinxstyleemphasis{nano}

\end{itemize}


\subsection{Examples}
\label{\detokenize{units:examples}}
\begin{sphinxVerbatim}[commandchars=\\\{\}]
\PYG{k}{units}\PYG{+w}{ }\PYG{n}{metal}
\PYG{k}{units}\PYG{+w}{ }\PYG{n}{lj}
\end{sphinxVerbatim}


\subsection{Description}
\label{\detokenize{units:description}}
\sphinxAtStartPar
This command sets the style of units used for a simulation.  It
determines the units of all quantities specified in the input script
and data file, as well as quantities output to the screen, log file,
and dump files.  Typically, this command is used at the very beginning
of an input script.

\sphinxAtStartPar
For all units except \sphinxstyleemphasis{lj}, LAMMPS uses physical constants from
www.physics.nist.gov.  For the definition of kcal in real units,
LAMMPS uses the thermochemical calorie = 4.184 J.

\sphinxAtStartPar
The choice you make for units simply sets some internal conversion
factors within LAMMPS.  This means that any simulation you perform for
one choice of units can be duplicated with any other unit setting
LAMMPS supports.  In this context “duplicate” means the particles will
have identical trajectories and all output generated by the simulation
will be identical.  This will be the case for some number of timesteps
until round\sphinxhyphen{}off effects accumulate, since the conversion factors for
two different unit systems are not identical to infinite precision.

\sphinxAtStartPar
To perform the same simulation in a different set of units you must
change all the unit\sphinxhyphen{}based input parameters in your input script and
other input files (data file, potential files, etc) correctly to the
new units.  And you must correctly convert all output from the new
units to the old units when comparing to the original results.  That
is often not simple to do.

\sphinxAtStartPar
Potential or table files may have a \sphinxcode{\sphinxupquote{UNITS:}} tag included in the
first line indicating the unit style those files were created for.
If the tag exists, its value will be compared to the chosen unit style
and LAMMPS will stop with an error message if there is a mismatch.
In some select cases and for specific combinations of unit styles,
LAMMPS is capable of automatically converting potential parameters
from a file. In those cases, a warning message signaling that an
automatic conversion has happened is printed to the screen.


\bigskip\hrule\bigskip


\sphinxAtStartPar
For style \sphinxstyleemphasis{lj}, all quantities are unitless.  Without loss of
generality, LAMMPS sets the fundamental quantities mass, \(\sigma\),
\(\epsilon\), and the Boltzmann constant \(k_B = 1\).  The
masses, distances, energies you specify are multiples of these
fundamental values.  The formulas relating the reduced or unitless
quantity (with an asterisk) to the same quantity with units is also
given.  Thus you can use the mass, \(\sigma\), and \(\epsilon\)
values for a specific material and convert the results from a unitless
LJ simulation into physical quantities.  Please note that using
these three properties as base, your unit of time has to conform
to the relation \(\epsilon = \frac{m \sigma^2}{\tau^2}\) since
energy is a derived unit (in SI units you equivalently have the relation
\(1\mathsf{J} = 1\frac{\mathsf{kg}\cdot\mathsf{m}^2}{\mathsf{s}^2}\)).
\begin{itemize}
\item {} 
\sphinxAtStartPar
mass = mass or \(m\), where \(M^* = \frac{M}{m}\)

\item {} 
\sphinxAtStartPar
distance = \(\sigma\), where \(x^* = \frac{x}{\sigma}\)

\item {} 
\sphinxAtStartPar
time = \(\tau\), where \(\tau^* = \tau \sqrt{\frac{\epsilon}{m \sigma^2}}\)

\item {} 
\sphinxAtStartPar
energy = \(\epsilon\), where \(E^* = \frac{E}{\epsilon}\)

\item {} 
\sphinxAtStartPar
velocity = \(\frac{\sigma}{\tau}\), where \(v^* = v \frac{\tau}{\sigma}\)

\item {} 
\sphinxAtStartPar
force = \(\frac{\epsilon}{\sigma}\), where \(f^* = f \frac{\sigma}{\epsilon}\)

\item {} 
\sphinxAtStartPar
torque = \(\epsilon\), where \(t^* = \frac{t}{\epsilon}\)

\item {} 
\sphinxAtStartPar
temperature = reduced LJ temperature, where \(T^* = \frac{T k_B}{\epsilon}\)

\item {} 
\sphinxAtStartPar
pressure = reduced LJ pressure, where \(p^* = p \frac{\sigma^3}{\epsilon}\)

\item {} 
\sphinxAtStartPar
dynamic viscosity = reduced LJ viscosity, where \(\eta^* = \eta \frac{\sigma^3}{\epsilon\tau}\)

\item {} 
\sphinxAtStartPar
charge = reduced LJ charge, where \(q^* = q \frac{1}{\sqrt{4 \pi \varepsilon_0 \sigma \epsilon}}\)

\item {} 
\sphinxAtStartPar
dipole = reduced LJ dipole, moment where \(\mu^* = \mu \frac{1}{\sqrt{4 \pi \varepsilon_0 \sigma^3 \epsilon}}\)

\item {} 
\sphinxAtStartPar
electric field = force/charge, where \(E^* = E \frac{\sqrt{4 \pi \varepsilon_0 \sigma \epsilon} \sigma}{\epsilon}\)

\item {} 
\sphinxAtStartPar
density = mass/volume, where \(\rho^* = \rho \frac{\sigma^{dim}}{m}\)

\end{itemize}

\sphinxAtStartPar
Note that for LJ units, the default mode of thermodynamic output via
the {\hyperref[\detokenize{thermo_style::doc}]{\sphinxcrossref{\DUrole{doc}{thermo\_style}}}} command is to normalize all
extensive quantities by the number of atoms.  E.g. potential energy is
extensive because it is summed over atoms, so it is output as
energy/atom.  Temperature is intensive since it is already normalized
by the number of atoms, so it is output as\sphinxhyphen{}is.  This behavior can be
changed via the {\hyperref[\detokenize{thermo_modify::doc}]{\sphinxcrossref{\DUrole{doc}{thermo\_modify norm}}}} command.

\sphinxAtStartPar
For style \sphinxstyleemphasis{real}, these are the units:
\begin{itemize}
\item {} 
\sphinxAtStartPar
mass = grams/mole

\item {} 
\sphinxAtStartPar
distance = Angstroms

\item {} 
\sphinxAtStartPar
time = femtoseconds

\item {} 
\sphinxAtStartPar
energy = kcal/mol

\item {} 
\sphinxAtStartPar
velocity = Angstroms/femtosecond

\item {} 
\sphinxAtStartPar
force = (kcal/mol)/Angstrom

\item {} 
\sphinxAtStartPar
torque = kcal/mol

\item {} 
\sphinxAtStartPar
temperature = Kelvin

\item {} 
\sphinxAtStartPar
pressure = atmospheres

\item {} 
\sphinxAtStartPar
dynamic viscosity = Poise

\item {} 
\sphinxAtStartPar
charge = multiple of electron charge (1.0 is a proton)

\item {} 
\sphinxAtStartPar
dipole = charge*Angstroms

\item {} 
\sphinxAtStartPar
electric field = volts/Angstrom

\item {} 
\sphinxAtStartPar
density = g/cm\textasciicircum{}dim

\end{itemize}

\sphinxAtStartPar
For style \sphinxstyleemphasis{metal}, these are the units:
\begin{itemize}
\item {} 
\sphinxAtStartPar
mass = grams/mole

\item {} 
\sphinxAtStartPar
distance = Angstroms

\item {} 
\sphinxAtStartPar
time = picoseconds

\item {} 
\sphinxAtStartPar
energy = eV

\item {} 
\sphinxAtStartPar
velocity = Angstroms/picosecond

\item {} 
\sphinxAtStartPar
force = eV/Angstrom

\item {} 
\sphinxAtStartPar
torque = eV

\item {} 
\sphinxAtStartPar
temperature = Kelvin

\item {} 
\sphinxAtStartPar
pressure = bars

\item {} 
\sphinxAtStartPar
dynamic viscosity = Poise

\item {} 
\sphinxAtStartPar
charge = multiple of electron charge (1.0 is a proton)

\item {} 
\sphinxAtStartPar
dipole = charge*Angstroms

\item {} 
\sphinxAtStartPar
electric field = volts/Angstrom

\item {} 
\sphinxAtStartPar
density = gram/cm\textasciicircum{}dim

\end{itemize}

\sphinxAtStartPar
For style \sphinxstyleemphasis{si}, these are the units:
\begin{itemize}
\item {} 
\sphinxAtStartPar
mass = kilograms

\item {} 
\sphinxAtStartPar
distance = meters

\item {} 
\sphinxAtStartPar
time = seconds

\item {} 
\sphinxAtStartPar
energy = Joules

\item {} 
\sphinxAtStartPar
velocity = meters/second

\item {} 
\sphinxAtStartPar
force = Newtons

\item {} 
\sphinxAtStartPar
torque = Newton\sphinxhyphen{}meters

\item {} 
\sphinxAtStartPar
temperature = Kelvin

\item {} 
\sphinxAtStartPar
pressure = Pascals

\item {} 
\sphinxAtStartPar
dynamic viscosity = Pascal*second

\item {} 
\sphinxAtStartPar
charge = Coulombs (1.6021765e\sphinxhyphen{}19 is a proton)

\item {} 
\sphinxAtStartPar
dipole = Coulombs*meters

\item {} 
\sphinxAtStartPar
electric field = volts/meter

\item {} 
\sphinxAtStartPar
density = kilograms/meter\textasciicircum{}dim

\end{itemize}

\sphinxAtStartPar
For style \sphinxstyleemphasis{cgs}, these are the units:
\begin{itemize}
\item {} 
\sphinxAtStartPar
mass = grams

\item {} 
\sphinxAtStartPar
distance = centimeters

\item {} 
\sphinxAtStartPar
time = seconds

\item {} 
\sphinxAtStartPar
energy = ergs

\item {} 
\sphinxAtStartPar
velocity = centimeters/second

\item {} 
\sphinxAtStartPar
force = dynes

\item {} 
\sphinxAtStartPar
torque = dyne\sphinxhyphen{}centimeters

\item {} 
\sphinxAtStartPar
temperature = Kelvin

\item {} 
\sphinxAtStartPar
pressure = dyne/cm\textasciicircum{}2 or barye = 1.0e\sphinxhyphen{}6 bars

\item {} 
\sphinxAtStartPar
dynamic viscosity = Poise

\item {} 
\sphinxAtStartPar
charge = statcoulombs or esu (4.8032044e\sphinxhyphen{}10 is a proton)

\item {} 
\sphinxAtStartPar
dipole = statcoul\sphinxhyphen{}cm = 10\textasciicircum{}18 debye

\item {} 
\sphinxAtStartPar
electric field = statvolt/cm or dyne/esu

\item {} 
\sphinxAtStartPar
density = grams/cm\textasciicircum{}dim

\end{itemize}

\sphinxAtStartPar
For style \sphinxstyleemphasis{electron}, these are the units:
\begin{itemize}
\item {} 
\sphinxAtStartPar
mass = atomic mass units

\item {} 
\sphinxAtStartPar
distance = Bohr

\item {} 
\sphinxAtStartPar
time = femtoseconds

\item {} 
\sphinxAtStartPar
energy = Hartrees

\item {} 
\sphinxAtStartPar
velocity = Bohr/atomic time units {[}1.03275e\sphinxhyphen{}15 seconds{]}

\item {} 
\sphinxAtStartPar
force = Hartrees/Bohr

\item {} 
\sphinxAtStartPar
temperature = Kelvin

\item {} 
\sphinxAtStartPar
pressure = Pascals

\item {} 
\sphinxAtStartPar
charge = multiple of electron charge (1.0 is a proton)

\item {} 
\sphinxAtStartPar
dipole moment = Debye

\item {} 
\sphinxAtStartPar
electric field = volts/cm

\end{itemize}

\sphinxAtStartPar
For style \sphinxstyleemphasis{micro}, these are the units:
\begin{itemize}
\item {} 
\sphinxAtStartPar
mass = picograms

\item {} 
\sphinxAtStartPar
distance = micrometers

\item {} 
\sphinxAtStartPar
time = microseconds

\item {} 
\sphinxAtStartPar
energy = picogram\sphinxhyphen{}micrometer\textasciicircum{}2/microsecond\textasciicircum{}2

\item {} 
\sphinxAtStartPar
velocity = micrometers/microsecond

\item {} 
\sphinxAtStartPar
force = picogram\sphinxhyphen{}micrometer/microsecond\textasciicircum{}2

\item {} 
\sphinxAtStartPar
torque = picogram\sphinxhyphen{}micrometer\textasciicircum{}2/microsecond\textasciicircum{}2

\item {} 
\sphinxAtStartPar
temperature = Kelvin

\item {} 
\sphinxAtStartPar
pressure = picogram/(micrometer\sphinxhyphen{}microsecond\textasciicircum{}2)

\item {} 
\sphinxAtStartPar
dynamic viscosity = picogram/(micrometer\sphinxhyphen{}microsecond)

\item {} 
\sphinxAtStartPar
charge = picocoulombs (1.6021765e\sphinxhyphen{}7 is a proton)

\item {} 
\sphinxAtStartPar
dipole = picocoulomb\sphinxhyphen{}micrometer

\item {} 
\sphinxAtStartPar
electric field = volt/micrometer

\item {} 
\sphinxAtStartPar
density = picograms/micrometer\textasciicircum{}dim

\end{itemize}

\sphinxAtStartPar
For style \sphinxstyleemphasis{nano}, these are the units:
\begin{itemize}
\item {} 
\sphinxAtStartPar
mass = attograms

\item {} 
\sphinxAtStartPar
distance = nanometers

\item {} 
\sphinxAtStartPar
time = nanoseconds

\item {} 
\sphinxAtStartPar
energy = attogram\sphinxhyphen{}nanometer\textasciicircum{}2/nanosecond\textasciicircum{}2

\item {} 
\sphinxAtStartPar
velocity = nanometers/nanosecond

\item {} 
\sphinxAtStartPar
force = attogram\sphinxhyphen{}nanometer/nanosecond\textasciicircum{}2

\item {} 
\sphinxAtStartPar
torque = attogram\sphinxhyphen{}nanometer\textasciicircum{}2/nanosecond\textasciicircum{}2

\item {} 
\sphinxAtStartPar
temperature = Kelvin

\item {} 
\sphinxAtStartPar
pressure = attogram/(nanometer\sphinxhyphen{}nanosecond\textasciicircum{}2)

\item {} 
\sphinxAtStartPar
dynamic viscosity = attogram/(nanometer\sphinxhyphen{}nanosecond)

\item {} 
\sphinxAtStartPar
charge = multiple of electron charge (1.0 is a proton)

\item {} 
\sphinxAtStartPar
dipole = charge\sphinxhyphen{}nanometer

\item {} 
\sphinxAtStartPar
electric field = volt/nanometer

\item {} 
\sphinxAtStartPar
density = attograms/nanometer\textasciicircum{}dim

\end{itemize}

\sphinxAtStartPar
The units command also sets the timestep size and neighbor skin
distance to default values for each style:
\begin{itemize}
\item {} 
\sphinxAtStartPar
For style \sphinxstyleemphasis{lj} these are dt = 0.005 \(\tau\) and skin = 0.3 \(\sigma\).

\item {} 
\sphinxAtStartPar
For style \sphinxstyleemphasis{real} these are dt = 1.0 femtoseconds and skin = 2.0 Angstroms.

\item {} 
\sphinxAtStartPar
For style \sphinxstyleemphasis{metal} these are dt = 0.001 picoseconds and skin = 2.0 Angstroms.

\item {} 
\sphinxAtStartPar
For style \sphinxstyleemphasis{si} these are dt = 1.0e\sphinxhyphen{}8 seconds and skin = 0.001 meters.

\item {} 
\sphinxAtStartPar
For style \sphinxstyleemphasis{cgs} these are dt = 1.0e\sphinxhyphen{}8 seconds and skin = 0.1 centimeters.

\item {} 
\sphinxAtStartPar
For style \sphinxstyleemphasis{electron} these are dt = 0.001 femtoseconds and skin = 2.0 Bohr.

\item {} 
\sphinxAtStartPar
For style \sphinxstyleemphasis{micro} these are dt = 2.0 microseconds and skin = 0.1 micrometers.

\item {} 
\sphinxAtStartPar
For style \sphinxstyleemphasis{nano} these are dt = 0.00045 nanoseconds and skin = 0.1 nanometers.

\end{itemize}


\subsection{Restrictions}
\label{\detokenize{units:restrictions}}
\sphinxAtStartPar
This command cannot be used after the simulation box is defined by a
{\hyperref[\detokenize{read_data::doc}]{\sphinxcrossref{\DUrole{doc}{read\_data}}}} or {\hyperref[\detokenize{create_box::doc}]{\sphinxcrossref{\DUrole{doc}{create\_box}}}} command.


\subsection{Related commands}
\label{\detokenize{units:related-commands}}
\sphinxAtStartPar
none


\subsection{Default}
\label{\detokenize{units:default}}
\begin{sphinxVerbatim}[commandchars=\\\{\}]
\PYG{k}{units}\PYG{+w}{ }\PYG{n}{lj}
\end{sphinxVerbatim}

\sphinxstepscope

\index{variable@\spxentry{variable}}\ignorespaces 

\section{variable command}
\label{\detokenize{variable:variable-command}}\label{\detokenize{variable:index-0}}\label{\detokenize{variable::doc}}

\subsection{Syntax}
\label{\detokenize{variable:syntax}}
\begin{sphinxVerbatim}[commandchars=\\\{\}]
\PYG{k}{variable }\PYG{n+nv+nvIdentifier}{name}\PYG{+w}{ }\PYG{n}{style}\PYG{+w}{ }\PYG{n}{args}\PYG{+w}{ }\PYG{n}{...}
\end{sphinxVerbatim}
\begin{itemize}
\item {} 
\sphinxAtStartPar
name = name of variable to define

\item {} 
\sphinxAtStartPar
style = \sphinxstyleemphasis{delete} or \sphinxstyleemphasis{atomfile} or \sphinxstyleemphasis{file} or \sphinxstyleemphasis{format} or \sphinxstyleemphasis{getenv} or \sphinxstyleemphasis{index} or \sphinxstyleemphasis{internal} or \sphinxstyleemphasis{loop} or \sphinxstyleemphasis{python} or \sphinxstyleemphasis{string} or \sphinxstyleemphasis{timer} or \sphinxstyleemphasis{uloop} or \sphinxstyleemphasis{universe} or \sphinxstyleemphasis{world} or \sphinxstyleemphasis{equal} or \sphinxstyleemphasis{vector} or \sphinxstyleemphasis{atom}
\begin{sphinxalltt}
\sphinxstyleemphasis{delete} = no args
\sphinxstyleemphasis{atomfile} arg = filename
\sphinxstyleemphasis{file} arg = filename
\sphinxstyleemphasis{format} args = vname fstr
  vname = name of equal\sphinxhyphen{}style variable to evaluate
  fstr = C\sphinxhyphen{}style format string
\sphinxstyleemphasis{getenv} arg = one string
\sphinxstyleemphasis{index} args = one or more strings
\sphinxstyleemphasis{internal} arg = numeric value
\sphinxstyleemphasis{loop} args = N
  N = integer size of loop, loop from 1 to N inclusive
\sphinxstyleemphasis{loop} args = N pad
  N = integer size of loop, loop from 1 to N inclusive
  pad = all values will be same length, e.g. 001, 002, ..., 100
\sphinxstyleemphasis{loop} args = N1 N2
  N1,N2 = loop from N1 to N2 inclusive
\sphinxstyleemphasis{loop} args = N1 N2 pad
  N1,N2 = loop from N1 to N2 inclusive
  pad = all values will be same length, e.g. 050, 051, ..., 100
\sphinxstyleemphasis{python} arg = function
\sphinxstyleemphasis{string} arg = one string
\sphinxstyleemphasis{timer} arg = no arguments
\sphinxstyleemphasis{uloop} args = N
  N = integer size of loop
\sphinxstyleemphasis{uloop} args = N pad
  N = integer size of loop
  pad = all values will be same length, e.g. 001, 002, ..., 100
\sphinxstyleemphasis{universe} args = one or more strings
\sphinxstyleemphasis{world} args = one string for each partition of processors

\sphinxstyleemphasis{equal} or \sphinxstyleemphasis{vector} or \sphinxstyleemphasis{atom} args = one formula containing numbers, thermo keywords, math operations, built\sphinxhyphen{}in functions, atom values and vectors, compute/fix/variable references
  numbers = 0.0, 100, \sphinxhyphen{}5.4, 2.8e\sphinxhyphen{}4, etc
  constants = PI, version, on, off, true, false, yes, no
  thermo keywords = vol, ke, press, etc from {\hyperref[\detokenize{thermo_style::doc}]{\sphinxcrossref{\DUrole{doc}{thermo\_style}}}}
  math operators = (), \sphinxhyphen{}x, x+y, x\sphinxhyphen{}y, x*y, x/y, x\textasciicircum{}y, x\%y,
                   x == y, x != y, x \textless{} y, x \textless{}= y, x \textgreater{} y, x \textgreater{}= y, x \&\& y, x || y, x |\textasciicircum{} y, !x
  math functions = sqrt(x), exp(x), ln(x), log(x), abs(x),
                   sin(x), cos(x), tan(x), asin(x), acos(x), atan(x), atan2(y,x),
                   random(x,y,z), normal(x,y,z), ceil(x), floor(x), round(x), ternary(x,y,z),
                   ramp(x,y), stagger(x,y), logfreq(x,y,z), logfreq2(x,y,z),
                   logfreq3(x,y,z), stride(x,y,z), stride2(x,y,z,a,b,c),
                   vdisplace(x,y), swiggle(x,y,z), cwiggle(x,y,z)
  group functions = count(group), mass(group), charge(group),
                    xcm(group,dim), vcm(group,dim), fcm(group,dim),
                    bound(group,dir), gyration(group), ke(group),
                    angmom(group,dim), torque(group,dim),
                    inertia(group,dimdim), omega(group,dim)
  region functions = count(group,region), mass(group,region), charge(group,region),
                    xcm(group,dim,region), vcm(group,dim,region), fcm(group,dim,region),
                    bound(group,dir,region), gyration(group,region), ke(group,reigon),
                    angmom(group,dim,region), torque(group,dim,region),
                    inertia(group,dimdim,region), omega(group,dim,region)
  special functions = sum(x), min(x), max(x), ave(x), trap(x), slope(x), sort(x), rsort(x), gmask(x), rmask(x), grmask(x,y), next(x), is\_file(name), is\_os(name), extract\_setting(name), label2type(kind,label), is\_typelabel(kind,label), is\_timeout()
  feature functions = is\_available(category,feature), is\_active(category,feature), is\_defined(category,id)
  atom value = id{[}i{]}, mass{[}i{]}, type{[}i{]}, mol{[}i{]}, x{[}i{]}, y{[}i{]}, z{[}i{]}, vx{[}i{]}, vy{[}i{]}, vz{[}i{]}, fx{[}i{]}, fy{[}i{]}, fz{[}i{]}, q{[}i{]}
  atom vector = id, mass, type, mol, radius, q, x, y, z, vx, vy, vz, fx, fy, fz
  custom atom property = i\_name, d\_name, i\_name{[}i{]}, d\_name{[}i{]}, i2\_name{[}i{]}, d2\_name{[}i{]}, i2\_name{[}i{]}{[}j{]}, d\_name{[}i{]}{[}j{]}
  compute references = c\_ID, c\_ID{[}i{]}, c\_ID{[}i{]}{[}j{]}, C\_ID, C\_ID{[}i{]}
  fix references = f\_ID, f\_ID{[}i{]}, f\_ID{[}i{]}{[}j{]}, F\_ID, F\_ID{[}i{]}
  variable references = v\_name, v\_name{[}i{]}
  vector initialization = {[}1,3,7,10{]} (for \sphinxstyleemphasis{vector} variables only)
\end{sphinxalltt}

\end{itemize}


\subsection{Examples}
\label{\detokenize{variable:examples}}
\begin{sphinxVerbatim}[commandchars=\\\{\}]
\PYG{k}{variable }\PYG{n+nv+nvIdentifier}{x}\PYG{+w}{ }\PYG{n}{index}\PYG{+w}{ }\PYG{n}{run1}\PYG{+w}{ }\PYG{n}{run2}\PYG{+w}{ }\PYG{n}{run3}\PYG{+w}{ }\PYG{n}{run4}\PYG{+w}{ }\PYG{n}{run5}\PYG{+w}{ }\PYG{n}{run6}\PYG{+w}{ }\PYG{n}{run7}\PYG{+w}{ }\PYG{n}{run8}
\PYG{k}{variable }\PYG{n+nv+nvIdentifier}{LoopVar}\PYG{+w}{ }\PYG{n}{loop}\PYG{+w}{ }\PYG{n+nv}{\PYGZdl{}n}
\PYG{k}{variable }\PYG{n+nv+nvIdentifier}{beta}\PYG{+w}{ }\PYG{n}{equal}\PYG{+w}{ }\PYG{n}{temp}\PYG{o}{/}\PYG{l+m}{3.0}
\PYG{k}{variable }\PYG{n+nv+nvIdentifier}{b1}\PYG{+w}{ }\PYG{n}{equal}\PYG{+w}{ }\PYG{n}{x[234]}\PYG{o}{+}\PYG{l+m}{0.5}\PYG{o}{*}\PYG{n}{vol}
\PYG{k}{variable }\PYG{n+nv+nvIdentifier}{b1}\PYG{+w}{ }\PYG{n}{equal}\PYG{+w}{ }\PYG{l+s}{\PYGZdq{}}\PYG{l+s}{x[234] + 0.5*vol}\PYG{l+s}{\PYGZdq{}}
\PYG{k}{variable }\PYG{n+nv+nvIdentifier}{b}\PYG{+w}{ }\PYG{n}{equal}\PYG{+w}{ }\PYG{n}{xcm}\PYG{n+nv}{(}\PYG{n+nv}{mol1,x}\PYG{n+nv}{)}\PYG{o}{/}\PYG{l+m}{2.0}
\PYG{k}{variable }\PYG{n+nv+nvIdentifier}{b}\PYG{+w}{ }\PYG{n}{equal}\PYG{+w}{ }\PYG{n}{c\PYGZus{}myTemp}
\PYG{k}{variable }\PYG{n+nv+nvIdentifier}{b}\PYG{+w}{ }\PYG{n}{atom}\PYG{+w}{ }\PYG{n}{x}\PYG{o}{*}\PYG{n}{y}\PYG{o}{/}\PYG{n}{vol}
\PYG{k}{variable }\PYG{n+nv+nvIdentifier}{foo}\PYG{+w}{ }\PYG{n}{string}\PYG{+w}{ }\PYG{n}{myfile}
\PYG{k}{variable }\PYG{n+nv+nvIdentifier}{foo}\PYG{+w}{ }\PYG{n}{internal}\PYG{+w}{ }\PYG{l+m}{3.5}
\PYG{k}{variable }\PYG{n+nv+nvIdentifier}{myPy}\PYG{+w}{ }\PYG{n}{python}\PYG{+w}{ }\PYG{n}{increase}
\PYG{k}{variable }\PYG{n+nv+nvIdentifier}{f}\PYG{+w}{ }\PYG{n}{file}\PYG{+w}{ }\PYG{n}{values.txt}
\PYG{k}{variable }\PYG{n+nv+nvIdentifier}{temp}\PYG{+w}{ }\PYG{n}{world}\PYG{+w}{ }\PYG{l+m}{300.0}\PYG{+w}{ }\PYG{l+m}{310.0}\PYG{+w}{ }\PYG{l+m}{320.0}\PYG{+w}{ }\PYG{n+nv}{\PYGZdl{}\PYGZob{}}\PYG{n+nv}{Tfinal}\PYG{n+nv}{\PYGZcb{}}
\PYG{k}{variable }\PYG{n+nv+nvIdentifier}{x}\PYG{+w}{ }\PYG{n}{universe}\PYG{+w}{ }\PYG{l+m}{1}\PYG{+w}{ }\PYG{l+m}{2}\PYG{+w}{ }\PYG{l+m}{3}\PYG{+w}{ }\PYG{l+m}{4}\PYG{+w}{ }\PYG{l+m}{5}\PYG{+w}{ }\PYG{l+m}{6}\PYG{+w}{ }\PYG{l+m}{7}\PYG{+w}{ }\PYG{l+m}{8}\PYG{+w}{ }\PYG{l+m}{9}\PYG{+w}{ }\PYG{l+m}{10}\PYG{+w}{ }\PYG{l+m}{11}\PYG{+w}{ }\PYG{l+m}{12}\PYG{+w}{ }\PYG{l+m}{13}\PYG{+w}{ }\PYG{l+m}{14}\PYG{+w}{ }\PYG{l+m}{15}
\PYG{k}{variable }\PYG{n+nv+nvIdentifier}{x}\PYG{+w}{ }\PYG{n}{uloop}\PYG{+w}{ }\PYG{l+m}{15}\PYG{+w}{ }\PYG{n}{pad}
\PYG{k}{variable }\PYG{n+nv+nvIdentifier}{str}\PYG{+w}{ }\PYG{n}{format}\PYG{+w}{ }\PYG{n}{x}\PYG{+w}{ }\PYG{o}{\PYGZpc{}}\PYG{n}{.6g}
\PYG{k}{variable }\PYG{n+nv+nvIdentifier}{myvec}\PYG{+w}{ }\PYG{n}{vector}\PYG{+w}{ }\PYG{n}{[1},3,7,10\PYG{n}{]}
\PYG{k}{variable }\PYG{n+nv+nvIdentifier}{x}\PYG{+w}{ }\PYG{n}{delete}
\end{sphinxVerbatim}

\begin{sphinxVerbatim}[commandchars=\\\{\}]
\PYG{k}{variable }\PYG{n+nv+nvIdentifier}{start}\PYG{+w}{ }\PYG{n}{timer}
\PYG{n}{other}\PYG{+w}{ }\PYG{n}{commands}
\PYG{k}{variable }\PYG{n+nv+nvIdentifier}{stop}\PYG{+w}{ }\PYG{n}{timer}
\PYG{k}{print}\PYG{+w}{ }\PYG{l+s}{\PYGZdq{}}\PYG{l+s}{Elapsed time: \PYGZdl{}(v\PYGZus{}stop\PYGZhy{}v\PYGZus{}start:\PYGZpc{}.6f)}\PYG{l+s}{\PYGZdq{}}
\end{sphinxVerbatim}


\subsection{Description}
\label{\detokenize{variable:description}}
\sphinxAtStartPar
This command assigns one or more strings to a variable name for
evaluation later in the input script or during a simulation.

\sphinxAtStartPar
Variables can thus be useful in several contexts.  A variable can be
defined and then referenced elsewhere in an input script to become
part of a new input command.  For variable styles that store multiple
strings, the {\hyperref[\detokenize{next::doc}]{\sphinxcrossref{\DUrole{doc}{next}}}} command can be used to increment which
string is assigned to the variable.  Variables of style \sphinxstyleemphasis{equal} store
a formula which when evaluated produces a single numeric value which
can be output either directly (see the {\hyperref[\detokenize{print::doc}]{\sphinxcrossref{\DUrole{doc}{print}}}}, {\hyperref[\detokenize{fix_print::doc}]{\sphinxcrossref{\DUrole{doc}{fix
print}}}}, and {\hyperref[\detokenize{run::doc}]{\sphinxcrossref{\DUrole{doc}{run every}}}} commands) or as part of
thermodynamic output (see the {\hyperref[\detokenize{thermo_style::doc}]{\sphinxcrossref{\DUrole{doc}{thermo\_style}}}}
command), or used as input to an averaging fix (see the {\hyperref[\detokenize{fix_ave_time::doc}]{\sphinxcrossref{\DUrole{doc}{fix
ave/time}}}} command).  Variables of style \sphinxstyleemphasis{vector} store
a formula which produces a vector of such values which can be used as
input to various averaging fixes, or elements of which can be part of
thermodynamic output.  Variables of style \sphinxstyleemphasis{atom} store a formula which
when evaluated produces one numeric value per atom which can be output
to a dump file (see the {\hyperref[\detokenize{dump::doc}]{\sphinxcrossref{\DUrole{doc}{dump custom}}}} command) or used as
input to an averaging fix (see the {\hyperref[\detokenize{fix_ave_chunk::doc}]{\sphinxcrossref{\DUrole{doc}{fix ave/chunk}}}} and {\hyperref[\detokenize{fix_ave_atom::doc}]{\sphinxcrossref{\DUrole{doc}{fix ave/atom}}}} commands).
Variables of style \sphinxstyleemphasis{atomfile} can be used anywhere in an input script
that atom\sphinxhyphen{}style variables are used; they get their per\sphinxhyphen{}atom values
from a file rather than from a formula.  Variables of style \sphinxstyleemphasis{python}
can be hooked to Python functions using code you provide, so that the
variable gets its value from the evaluation of the Python code.
Variables of style \sphinxstyleemphasis{internal} are used by a few commands which set
their value directly.

\begin{sphinxadmonition}{note}{Note:}
\sphinxAtStartPar
As discussed on the {\hyperref[\detokenize{Commands_parse::doc}]{\sphinxcrossref{\DUrole{doc}{Commands parse}}}} doc
page, an input script can use “immediate” variables, specified as
\$(formula) with parenthesis, where the numeric formula has the same
syntax as equal\sphinxhyphen{}style variables described on this page.  This is a
convenient way to evaluate a formula immediately without using the
variable command to define a named variable and then evaluate that
variable.  The formula can include a trailing colon and format
string which determines the precision with which the numeric value
is generated.  This is also explained on the {\hyperref[\detokenize{Commands_parse::doc}]{\sphinxcrossref{\DUrole{doc}{Commands parse}}}} doc page.
\end{sphinxadmonition}

\sphinxAtStartPar
In the discussion that follows, the “name” of the variable is the
arbitrary string that is the first argument in the variable command.
This name can only contain alphanumeric characters and underscores.
The “string” is one or more of the subsequent arguments.  The “string”
can be simple text as in the first example above, it can contain other
variables as in the second example, or it can be a formula as in the third
example.  The “value” is the numeric quantity resulting from
evaluation of the string.  Note that the same string can generate
different values when it is evaluated at different times during a
simulation.

\begin{sphinxadmonition}{note}{Note:}
\sphinxAtStartPar
When an input script line is encountered that defines a variable
of style \sphinxstyleemphasis{equal} or \sphinxstyleemphasis{vector} or \sphinxstyleemphasis{atom} or \sphinxstyleemphasis{python} that contains a
formula or Python code, the formula is NOT immediately evaluated.  It
will be evaluated every time when the variable is \sphinxstylestrong{used} instead.  If
you simply want to evaluate a formula in place you can use as
so\sphinxhyphen{}called. See the section below about “Immediate Evaluation of
Variables” for more details on the topic.  This is also true of a
\sphinxstyleemphasis{format} style variable since it evaluates another variable when it is
invoked.
\end{sphinxadmonition}

\sphinxAtStartPar
Variables of style \sphinxstyleemphasis{equal} and \sphinxstyleemphasis{vector} and \sphinxstyleemphasis{atom} can be used as
inputs to various other commands which evaluate their formulas as
needed, e.g. at different timesteps during a {\hyperref[\detokenize{run::doc}]{\sphinxcrossref{\DUrole{doc}{run}}}}.  In
this context, variables of style \sphinxstyleemphasis{timer} or \sphinxstyleemphasis{internal} or \sphinxstyleemphasis{python} can
be used in place of an equal\sphinxhyphen{}style variable, with the following two
caveats.

\sphinxAtStartPar
First, internal\sphinxhyphen{}style variables can be used except by commands that
set the value stored by the internal variable.  When the LAMMPS
command evaluates the internal\sphinxhyphen{}style variable, it will use the value
set (internally) by another command.  Second, python\sphinxhyphen{}style variables
can be used so long as the associated Python function, as defined by
the {\hyperref[\detokenize{python::doc}]{\sphinxcrossref{\DUrole{doc}{python}}}} command, returns a numeric value.  When the
LAMMPS command evaluates the python\sphinxhyphen{}style variable, the Python
function will be executed.

\begin{sphinxadmonition}{note}{Note:}
\sphinxAtStartPar
When a variable command is encountered in the input script and
the variable name has already been specified, the command is ignored.
This means variables can NOT be re\sphinxhyphen{}defined in an input script (with
two exceptions, read further).  This is to allow an input script to be
processed multiple times without resetting the variables; see the
{\hyperref[\detokenize{jump::doc}]{\sphinxcrossref{\DUrole{doc}{jump}}}} or {\hyperref[\detokenize{include::doc}]{\sphinxcrossref{\DUrole{doc}{include}}}} commands.  It also means
that using the {\hyperref[\detokenize{Run_options::doc}]{\sphinxcrossref{\DUrole{doc}{command\sphinxhyphen{}line switch}}}} \sphinxhyphen{}var will
override a corresponding index variable setting in the input script.
\end{sphinxadmonition}

\sphinxAtStartPar
There are two exceptions to this rule.  First, variables of style
\sphinxstyleemphasis{string}, \sphinxstyleemphasis{getenv}, \sphinxstyleemphasis{internal}, \sphinxstyleemphasis{equal}, \sphinxstyleemphasis{vector}, \sphinxstyleemphasis{atom}, and
\sphinxstyleemphasis{python} ARE redefined each time the command is encountered.  This
allows these style of variables to be redefined multiple times in an
input script.  In a loop, this means the formula associated with an
\sphinxstyleemphasis{equal} or \sphinxstyleemphasis{atom} style variable can change if it contains a
substitution for another variable, e.g. \$x or v\_x.

\sphinxAtStartPar
Second, as described below, if a variable is iterated on to the end of
its list of strings via the {\hyperref[\detokenize{next::doc}]{\sphinxcrossref{\DUrole{doc}{next}}}} command, it is removed
from the list of active variables, and is thus available to be
re\sphinxhyphen{}defined in a subsequent variable command.  The \sphinxstyleemphasis{delete} style does
the same thing.

\sphinxAtStartPar
Variables are \sphinxstylestrong{not} deleted by the {\hyperref[\detokenize{clear::doc}]{\sphinxcrossref{\DUrole{doc}{clear}}}} command with
the exception of atomfile\sphinxhyphen{}style variables.


\bigskip\hrule\bigskip


\sphinxAtStartPar
The {\hyperref[\detokenize{Commands_parse::doc}]{\sphinxcrossref{\DUrole{doc}{Commands parse}}}} page explains how
occurrences of a variable name in an input script line are replaced by
the variable’s string.  The variable name can be referenced as \$x if
the name “x” is a single character, or as \$\{LoopVar\} if the name
“LoopVar” is one or more characters.

\sphinxAtStartPar
As described below, for variable styles \sphinxstyleemphasis{index}, \sphinxstyleemphasis{loop}, \sphinxstyleemphasis{file},
\sphinxstyleemphasis{universe}, and \sphinxstyleemphasis{uloop}, which string is assigned to a variable can be
incremented via the {\hyperref[\detokenize{next::doc}]{\sphinxcrossref{\DUrole{doc}{next}}}} command.  When there are no more
strings to assign, the variable is exhausted and a flag is set that
causes the next {\hyperref[\detokenize{jump::doc}]{\sphinxcrossref{\DUrole{doc}{jump}}}} command encountered in the input
script to be skipped.  This enables the construction of simple loops
in the input script that are iterated over and then exited from.

\sphinxAtStartPar
As explained above, an exhausted variable can be re\sphinxhyphen{}used in an input
script.  The \sphinxstyleemphasis{delete} style also removes the variable, the same as if
it were exhausted, allowing it to be redefined later in the input
script or when the input script is looped over.  This can be useful
when breaking out of a loop via the {\hyperref[\detokenize{if::doc}]{\sphinxcrossref{\DUrole{doc}{if}}}} and {\hyperref[\detokenize{jump::doc}]{\sphinxcrossref{\DUrole{doc}{jump}}}}
commands before the variable would become exhausted.  For example,

\begin{sphinxVerbatim}[commandchars=\\\{\}]
\PYG{k}{label       }\PYG{l+s+sc}{loop}
\PYG{k}{variable    }\PYG{n+nv+nvIdentifier}{a}\PYG{+w}{ }\PYG{n}{loop}\PYG{+w}{ }\PYG{l+m}{5}
\PYG{k}{print}\PYG{+w}{       }\PYG{l+s}{\PYGZdq{}}\PYG{l+s}{A = \PYGZdl{}a}\PYG{l+s}{\PYGZdq{}}
\PYG{k}{if}\PYG{+w}{          }\PYG{l+s}{\PYGZdq{}}\PYG{l+s}{\PYGZdl{}a \PYGZgt{} 2}\PYG{l+s}{\PYGZdq{}}\PYG{+w}{ }\PYG{k}{then}\PYG{+w}{ }\PYG{l+s}{\PYGZdq{}}\PYG{l+s}{jump in.script break}\PYG{l+s}{\PYGZdq{}}
\PYG{k}{next        }\PYG{n+nv+nvIdentifier}{a}
\PYG{k}{jump        }\PYG{l+s+sc}{in.script}\PYG{+w}{ }\PYG{n}{loop}
\PYG{k}{label       }\PYG{l+s+sc}{break}
\PYG{k}{variable    }\PYG{n+nv+nvIdentifier}{a}\PYG{+w}{ }\PYG{n}{delete}
\end{sphinxVerbatim}


\bigskip\hrule\bigskip


\sphinxAtStartPar
The next sections describe in how all the various variable styles are
defined and what they store.  The styles are listed alphabetically,
except for the \sphinxstyleemphasis{equal} and \sphinxstyleemphasis{vector} and \sphinxstyleemphasis{atom} styles, which are
explained together after all the others.

\sphinxAtStartPar
Many of the styles store one or more strings.  Note that a single
string can contain spaces (multiple words), if it is enclosed in
quotes in the variable command.  When the variable is substituted for
in another input script command, its returned string will then be
interpreted as multiple arguments in the expanded command.


\bigskip\hrule\bigskip


\sphinxAtStartPar
For the \sphinxstyleemphasis{atomfile} style, a filename is provided which contains one or
more sets of values, to assign on a per\sphinxhyphen{}atom basis to the variable.
The format of the file is described below.

\sphinxAtStartPar
When an atomfile\sphinxhyphen{}style variable is defined, the file is opened and the
first set of per\sphinxhyphen{}atom values are read and stored with the variable.
This means the variable can then be evaluated as many times as desired
and will return those values.  There are two ways to cause the next
set of per\sphinxhyphen{}atom values from the file to be read: use the
{\hyperref[\detokenize{next::doc}]{\sphinxcrossref{\DUrole{doc}{next}}}} command or the next() function in an atom\sphinxhyphen{}style
variable, as discussed below.  Unlike most variable styles, which
remain defined, atomfile\sphinxhyphen{}style variables are \sphinxstylestrong{deleted} during a
{\hyperref[\detokenize{clear::doc}]{\sphinxcrossref{\DUrole{doc}{clear}}}} command.

\sphinxAtStartPar
The rules for formatting the file are as follows.  Each time a set of
per\sphinxhyphen{}atom values is read, a non\sphinxhyphen{}blank line is searched for in the file.
The file is read line by line but only up to 254 characters are used.
The rest are ignored.  A comment character “\#” can be used anywhere
on a line and all text following and the “\#” character are ignored;
text starting with the comment character is stripped.  Blank lines
are skipped.  The first non\sphinxhyphen{}blank line is expected to contain a single
integer number as the count \sphinxstyleemphasis{N} of per\sphinxhyphen{}atom lines to follow.  \sphinxstyleemphasis{N} can
be the total number of atoms in the system or less, indicating that data
for a subset is read.  The next N lines must consist of two numbers,
the atom\sphinxhyphen{}ID of the atom for which a value is set followed by a floating
point number with the value.  The atom\sphinxhyphen{}IDs may be listed in any order.

\begin{sphinxadmonition}{note}{Note:}
\sphinxAtStartPar
Every time a set of per\sphinxhyphen{}atom lines is read, the value of the atomfile
variable for \sphinxstylestrong{all} atoms is first initialized to 0.0.  Thus values
for atoms whose ID do not appear in the set in the file will remain
at 0.0.
\end{sphinxadmonition}

\sphinxAtStartPar
Below is a small example for the atomfile variable file format:
\begin{quote}

\begin{sphinxVerbatim}[commandchars=\\\{\}]
\PYGZsh{} first set
4
\PYGZsh{} atom\PYGZhy{}ID value
3 1
4 \PYGZhy{}4
1 0.5
2 \PYGZhy{}0.5

\PYGZsh{} second set
2

2  1.0
4 \PYGZhy{}1.0
\end{sphinxVerbatim}
\end{quote}


\bigskip\hrule\bigskip


\sphinxAtStartPar
For the \sphinxstyleemphasis{file} style, a filename is provided which contains a list of
strings to assign to the variable, one per line.  The strings can be
numeric values if desired.  See the discussion of the next() function
below for equal\sphinxhyphen{}style variables, which will convert the string of a
file\sphinxhyphen{}style variable into a numeric value in a formula.

\sphinxAtStartPar
When a file\sphinxhyphen{}style variable is defined, the file is opened and the
string on the first line is read and stored with the variable.  This
means the variable can then be evaluated as many times as desired and
will return that string.  There are two ways to cause the next string
from the file to be read: use the {\hyperref[\detokenize{next::doc}]{\sphinxcrossref{\DUrole{doc}{next}}}} command or the
next() function in an equal\sphinxhyphen{} or atom\sphinxhyphen{}style variable, as discussed
below.

\sphinxAtStartPar
The rules for formatting the file are as follows.  A comment character
“\#” can be used anywhere on a line; text starting with the comment
character is stripped.  Blank lines are skipped.  The first “word” of
a non\sphinxhyphen{}blank line, delimited by white\sphinxhyphen{}space, is the “string” assigned
to the variable.


\bigskip\hrule\bigskip


\sphinxAtStartPar
For the \sphinxstyleemphasis{format} style, an equal\sphinxhyphen{}style or compatible variable is
specified along with a C\sphinxhyphen{}style format string, e.g. “\%f” or “\%.10g”,
which must be appropriate for formatting a double\sphinxhyphen{}precision
floating\sphinxhyphen{}point value and may not have extra characters.  The default
format is “\%.15g”.  This variable style allows an equal\sphinxhyphen{}style variable
to be formatted precisely when it is evaluated.

\sphinxAtStartPar
Note that if you simply wish to print a variable value with desired
precision to the screen or logfile via the {\hyperref[\detokenize{print::doc}]{\sphinxcrossref{\DUrole{doc}{print}}}} or
{\hyperref[\detokenize{fix_print::doc}]{\sphinxcrossref{\DUrole{doc}{fix print}}}} commands, you can also do this by
specifying an “immediate” variable with a trailing colon and format
string, as part of the string argument of those commands.  This is
explained on the {\hyperref[\detokenize{Commands_parse::doc}]{\sphinxcrossref{\DUrole{doc}{Commands parse}}}} doc page.


\bigskip\hrule\bigskip


\sphinxAtStartPar
For the \sphinxstyleemphasis{getenv} style, a single string is assigned to the variable
which should be the name of an environment variable.  When the
variable is evaluated, it returns the value of the environment
variable, or an empty string if it not defined.  This style of
variable can be used to adapt the behavior of LAMMPS input scripts via
environment variable settings, or to retrieve information that has
been previously stored with the {\hyperref[\detokenize{shell::doc}]{\sphinxcrossref{\DUrole{doc}{shell putenv}}}} command.
Note that because environment variable settings are stored by the
operating systems, they persist even if the corresponding \sphinxstyleemphasis{getenv}
style variable is deleted, and also are set for sub\sphinxhyphen{}shells executed
by the {\hyperref[\detokenize{shell::doc}]{\sphinxcrossref{\DUrole{doc}{shell}}}} command.


\bigskip\hrule\bigskip


\sphinxAtStartPar
For the \sphinxstyleemphasis{index} style, one or more strings are specified.  Initially,
the first string is assigned to the variable.  Each time a
{\hyperref[\detokenize{next::doc}]{\sphinxcrossref{\DUrole{doc}{next}}}} command is used with the variable name, the next
string is assigned.  All processors assign the same string to the
variable.

\sphinxAtStartPar
Index\sphinxhyphen{}style variables with a single string value can also be set by
using the {\hyperref[\detokenize{Run_options::doc}]{\sphinxcrossref{\DUrole{doc}{command\sphinxhyphen{}line switch \sphinxhyphen{}var}}}}.


\bigskip\hrule\bigskip


\sphinxAtStartPar
For the \sphinxstyleemphasis{internal} style a numeric value is provided.  This value will
be assigned to the variable until a LAMMPS command sets it to a new
value.  There are currently only two LAMMPS commands that require
\sphinxstyleemphasis{internal} variables as inputs, because they reset them:
{\hyperref[\detokenize{create_atoms::doc}]{\sphinxcrossref{\DUrole{doc}{create\_atoms}}}} and {\hyperref[\detokenize{fix_controller::doc}]{\sphinxcrossref{\DUrole{doc}{fix controller}}}}.  As mentioned above, an internal\sphinxhyphen{}style variable can
be used in place of an equal\sphinxhyphen{}style variable anywhere else in an input
script, e.g. as an argument to another command that allows for
equal\sphinxhyphen{}style variables.


\bigskip\hrule\bigskip


\sphinxAtStartPar
The \sphinxstyleemphasis{loop} style is identical to the \sphinxstyleemphasis{index} style except that the
strings are the integers from 1 to N inclusive, if only one argument N
is specified.  This allows generation of a long list of runs
(e.g. 1000) without having to list N strings in the input script.
Initially, the string “1” is assigned to the variable.  Each time a
{\hyperref[\detokenize{next::doc}]{\sphinxcrossref{\DUrole{doc}{next}}}} command is used with the variable name, the next
string (“2”, “3”, etc) is assigned.  All processors assign the same
string to the variable.  The \sphinxstyleemphasis{loop} style can also be specified with
two arguments N1 and N2.  In this case the loop runs from N1 to N2
inclusive, and the string N1 is initially assigned to the variable.
N1 \textless{}= N2 and N2 \textgreater{}= 0 is required.


\bigskip\hrule\bigskip


\sphinxAtStartPar
For the \sphinxstyleemphasis{python} style a Python function name is provided.  This needs
to match a function name specified in a {\hyperref[\detokenize{python::doc}]{\sphinxcrossref{\DUrole{doc}{python}}}} command
which returns a value to this variable as defined by its \sphinxstyleemphasis{return}
keyword.  For example these two commands would be self\sphinxhyphen{}consistent:

\begin{sphinxVerbatim}[commandchars=\\\{\}]
\PYG{k}{variable }\PYG{n+nv+nvIdentifier}{foo}\PYG{+w}{ }\PYG{n}{python}\PYG{+w}{ }\PYG{n}{myMultiply}
\PYG{k}{python}\PYG{+w}{ }\PYG{n}{myMultiply}\PYG{+w}{ }\PYG{n}{return}\PYG{+w}{ }\PYG{n}{v\PYGZus{}foo}\PYG{+w}{ }\PYG{n}{format}\PYG{+w}{ }\PYG{n}{f}\PYG{+w}{ }\PYG{n}{file}\PYG{+w}{ }\PYG{n}{funcs.py}
\end{sphinxVerbatim}

\sphinxAtStartPar
The two commands can appear in either order so long as both are
specified before the Python function is invoked for the first time.

\sphinxAtStartPar
Each time the variable is evaluated, the associated Python function is
invoked, and the value it returns is also returned by the variable.
Since the Python function can use other LAMMPS variables as input, or
query interal LAMMPS quantities to perform its computation, this means
the variable can return a different value each time it is evaluated.

\sphinxAtStartPar
The type of value stored in the variable is determined by the \sphinxstyleemphasis{format}
keyword of the {\hyperref[\detokenize{python::doc}]{\sphinxcrossref{\DUrole{doc}{python}}}} command.  It can be an integer
(i), floating point (f), or string (s) value.  As mentioned above, if
it is a numeric value (integer or floating point), then the
python\sphinxhyphen{}style variable can be used in place of an equal\sphinxhyphen{}style variable
anywhere in an input script, e.g. as an argument to another command
that allows for equal\sphinxhyphen{}style variables.


\bigskip\hrule\bigskip


\sphinxAtStartPar
For the \sphinxstyleemphasis{string} style, a single string is assigned to the variable.
Two differences between this style and using the \sphinxstyleemphasis{index} style exist:
a variable with \sphinxstyleemphasis{string} style can be redefined, e.g. by another command later
in the input script, or if the script is read again in a loop. The other
difference is that \sphinxstyleemphasis{string} performs variable substitution even if the
string parameter is quoted.


\bigskip\hrule\bigskip


\sphinxAtStartPar
The \sphinxstyleemphasis{uloop} style is identical to the \sphinxstyleemphasis{universe} style except that the
strings are the integers from 1 to N.  This allows generation of long
list of runs (e.g. 1000) without having to list N strings in the input
script.


\bigskip\hrule\bigskip


\sphinxAtStartPar
For the \sphinxstyleemphasis{universe} style, one or more strings are specified.  There
must be at least as many strings as there are processor partitions or
“worlds”.  LAMMPS can be run with multiple partitions via the
{\hyperref[\detokenize{Run_options::doc}]{\sphinxcrossref{\DUrole{doc}{\sphinxhyphen{}partition command\sphinxhyphen{}line switch}}}}.  This variable
command initially assigns one string to each world.  When a
{\hyperref[\detokenize{next::doc}]{\sphinxcrossref{\DUrole{doc}{next}}}} command is encountered using this variable, the first
processor partition to encounter it, is assigned the next available
string.  This continues until all the variable strings are consumed.
Thus, this command can be used to run 50 simulations on 8 processor
partitions.  The simulations will be run one after the other on
whatever partition becomes available, until they are all finished.
Universe\sphinxhyphen{}style variables are incremented using the files
“tmp.lammps.variable” and “tmp.lammps.variable.lock” which you will
see in your directory during such a LAMMPS run.


\bigskip\hrule\bigskip


\sphinxAtStartPar
For the \sphinxstyleemphasis{world} style, one or more strings are specified.  There must
be one string for each processor partition or “world”.  LAMMPS can be
run with multiple partitions via the {\hyperref[\detokenize{Run_options::doc}]{\sphinxcrossref{\DUrole{doc}{\sphinxhyphen{}partition command\sphinxhyphen{}line
switch}}}}.  This variable command assigns one string to
each world.  All processors in the world are assigned the same string.
The next command cannot be used with equal\sphinxhyphen{}style variables, since
there is only one value per world.  This style of variable is useful
when you wish to run different simulations on different partitions, or
when performing a parallel tempering simulation (see the {\hyperref[\detokenize{temper::doc}]{\sphinxcrossref{\DUrole{doc}{temper}}}} command), to assign different temperatures to different
partitions.


\bigskip\hrule\bigskip


\sphinxAtStartPar
For the \sphinxstyleemphasis{equal} and \sphinxstyleemphasis{vector} and \sphinxstyleemphasis{atom} styles, a single string is
specified which represents a formula that will be evaluated afresh
each time the variable is used.  If you want spaces in the string,
enclose it in double quotes so the parser will treat it as a single
argument.  For \sphinxstyleemphasis{equal}\sphinxhyphen{}style variables the formula computes a scalar
quantity, which becomes the value of the variable whenever it is
evaluated.  For \sphinxstyleemphasis{vector}\sphinxhyphen{}style variables the formula must compute a
vector of quantities, which becomes the value of the variable whenever
it is evaluated.  The calculated vector can be of length one, but it
cannot be a simple scalar value like that produced by an equal\sphinxhyphen{}style
compute.  I.e. the formula for a vector\sphinxhyphen{}style variable must have at
least one quantity in it that refers to a global vector produced by a
compute, fix, or other vector\sphinxhyphen{}style variable.  For \sphinxstyleemphasis{atom}\sphinxhyphen{}style
variables the formula computes one quantity for each atom whenever it
is evaluated.

\sphinxAtStartPar
Note that \sphinxstyleemphasis{equal}, \sphinxstyleemphasis{vector}, and \sphinxstyleemphasis{atom} variables can produce
different values at different stages of the input script or at
different times during a run.  For example, if an \sphinxstyleemphasis{equal} variable is
used in a {\hyperref[\detokenize{fix_print::doc}]{\sphinxcrossref{\DUrole{doc}{fix print}}}} command, different values could
be printed each timestep it was invoked.  If you want a variable to be
evaluated immediately, so that the result is stored by the variable
instead of the string, see the section below on “Immediate Evaluation
of Variables”.

\sphinxAtStartPar
The next command cannot be used with \sphinxstyleemphasis{equal} or \sphinxstyleemphasis{vector} or \sphinxstyleemphasis{atom}
style variables, since there is only one string.

\sphinxAtStartPar
The formula for an \sphinxstyleemphasis{equal}, \sphinxstyleemphasis{vector}, or \sphinxstyleemphasis{atom} variable can contain a
variety of quantities.  The syntax for each kind of quantity is
simple, but multiple quantities can be nested and combined in various
ways to build up formulas of arbitrary complexity.  For example, this
is a valid (though strange) variable formula:

\begin{sphinxVerbatim}[commandchars=\\\{\}]
\PYG{k}{variable }\PYG{n+nv+nvIdentifier}{x}\PYG{+w}{ }\PYG{n}{equal}\PYG{+w}{ }\PYG{l+s}{\PYGZdq{}}\PYG{l+s}{pe + c\PYGZus{}MyTemp / vol\PYGZca{}(1/3)}\PYG{l+s}{\PYGZdq{}}
\end{sphinxVerbatim}

\sphinxAtStartPar
Specifically, a formula can contain numbers, constants, thermo
keywords, math operators, math functions, group functions, region
functions, special functions, feature functions, atom values, atom
vectors, custom atom properties, compute references, fix references, and references to other
variables.


\begin{savenotes}\sphinxattablestart
\sphinxthistablewithglobalstyle
\centering
\begin{tabulary}{\linewidth}[t]{TT}
\sphinxtoprule
\sphinxtableatstartofbodyhook
\sphinxAtStartPar
Number
&
\sphinxAtStartPar
0.2, 100, 1.0e20, \sphinxhyphen{}15.4, etc
\\
\sphinxhline
\sphinxAtStartPar
Constant
&
\sphinxAtStartPar
PI, version, on, off, true, false, yes, no
\\
\sphinxhline
\sphinxAtStartPar
Thermo keywords
&
\sphinxAtStartPar
vol, pe, ebond, etc
\\
\sphinxhline
\sphinxAtStartPar
Math operators
&
\sphinxAtStartPar
(), \sphinxhyphen{}x, x+y, x\sphinxhyphen{}y, x*y, x/y, x\textasciicircum{}y, x\%y, x == y, x != y, x \textless{} y, x \textless{}= y, x \textgreater{} y, x \textgreater{}= y, x \&\& y, x || y, x |\textasciicircum{} y, !x
\\
\sphinxhline
\sphinxAtStartPar
Math functions
&
\sphinxAtStartPar
sqrt(x), exp(x), ln(x), log(x), abs(x), sin(x), cos(x), tan(x), asin(x), acos(x), atan(x), atan2(y,x), random(x,y,z), normal(x,y,z), ceil(x), floor(x), round(x), ternary(x,y,z), ramp(x,y), stagger(x,y), logfreq(x,y,z), logfreq2(x,y,z), logfreq3(x,y,z), stride(x,y,z), stride2(x,y,z,a,b,c), vdisplace(x,y), swiggle(x,y,z), cwiggle(x,y,z)
\\
\sphinxhline
\sphinxAtStartPar
Group functions
&
\sphinxAtStartPar
count(ID), mass(ID), charge(ID), xcm(ID,dim), vcm(ID,dim), fcm(ID,dim), bound(ID,dir), gyration(ID), ke(ID), angmom(ID,dim), torque(ID,dim), inertia(ID,dimdim), omega(ID,dim)
\\
\sphinxhline
\sphinxAtStartPar
Region functions
&
\sphinxAtStartPar
count(ID,IDR), mass(ID,IDR), charge(ID,IDR), xcm(ID,dim,IDR), vcm(ID,dim,IDR), fcm(ID,dim,IDR), bound(ID,dir,IDR), gyration(ID,IDR), ke(ID,IDR), angmom(ID,dim,IDR), torque(ID,dim,IDR), inertia(ID,dimdim,IDR), omega(ID,dim,IDR)
\\
\sphinxhline
\sphinxAtStartPar
Special functions
&
\sphinxAtStartPar
sum(x), min(x), max(x), ave(x), trap(x), slope(x), sort(x), rsort(x), gmask(x), rmask(x), grmask(x,y), next(x), is\_file(name), is\_os(name), extract\_setting(name), label2type(kind,label), is\_typelabel(kind,label), is\_timeout()
\\
\sphinxhline
\sphinxAtStartPar
Feature functions
&
\sphinxAtStartPar
is\_available(category,feature), is\_active(category,feature), is\_defined(category,id)
\\
\sphinxhline
\sphinxAtStartPar
Atom values
&
\sphinxAtStartPar
id{[}i{]}, mass{[}i{]}, type{[}i{]}, mol{[}i{]}, x{[}i{]}, y{[}i{]}, z{[}i{]}, vx{[}i{]}, vy{[}i{]}, vz{[}i{]}, fx{[}i{]}, fy{[}i{]}, fz{[}i{]}, q{[}i{]}
\\
\sphinxhline
\sphinxAtStartPar
Atom vectors
&
\sphinxAtStartPar
id, mass, type, mol, x, y, z, vx, vy, vz, fx, fy, fz, q
\\
\sphinxhline
\sphinxAtStartPar
Custom atom properties
&
\sphinxAtStartPar
i\_name, d\_name, i\_name{[}i{]}, d\_name{[}i{]}, i2\_name{[}i{]}, d2\_name{[}i{]}, i2\_name{[}i{]}{[}j{]}, d\_name{[}i{]}{[}j{]}
\\
\sphinxhline
\sphinxAtStartPar
Compute references
&
\sphinxAtStartPar
c\_ID, c\_ID{[}i{]}, c\_ID{[}i{]}{[}j{]}, C\_ID, C\_ID{[}i{]}
\\
\sphinxhline
\sphinxAtStartPar
Fix references
&
\sphinxAtStartPar
f\_ID, f\_ID{[}i{]}, f\_ID{[}i{]}{[}j{]}, F\_ID, F\_ID{[}i{]}
\\
\sphinxhline
\sphinxAtStartPar
Other variables
&
\sphinxAtStartPar
v\_name, v\_name{[}i{]}
\\
\sphinxbottomrule
\end{tabulary}
\sphinxtableafterendhook\par
\sphinxattableend\end{savenotes}

\sphinxAtStartPar
Most of the formula elements produce a scalar value.  Some produce a
global or per\sphinxhyphen{}atom vector of values.  Global vectors can be produced
by computes or fixes or by other vector\sphinxhyphen{}style variables.  Per\sphinxhyphen{}atom
vectors are produced by atom vectors, computes or fixes which output a
per\sphinxhyphen{}atom vector or array, and variables that are atom\sphinxhyphen{}style variables.
Math functions that operate on scalar values produce a scalar value;
math function that operate on global or per\sphinxhyphen{}atom vectors do so
element\sphinxhyphen{}by\sphinxhyphen{}element and produce a global or per\sphinxhyphen{}atom vector.

\sphinxAtStartPar
A formula for equal\sphinxhyphen{}style variables cannot use any formula element
that produces a global or per\sphinxhyphen{}atom vector.  A formula for a
vector\sphinxhyphen{}style variable can use formula elements that produce either a
scalar value or a global vector value, but cannot use a formula
element that produces a per\sphinxhyphen{}atom vector.  A formula for an atom\sphinxhyphen{}style
variable can use formula elements that produce either a scalar value
or a per\sphinxhyphen{}atom vector, but not one that produces a global vector.

\sphinxAtStartPar
Atom\sphinxhyphen{}style variables are evaluated by other commands that define a
{\hyperref[\detokenize{group::doc}]{\sphinxcrossref{\DUrole{doc}{group}}}} on which they operate, e.g. a {\hyperref[\detokenize{dump::doc}]{\sphinxcrossref{\DUrole{doc}{dump}}}}
or {\hyperref[\detokenize{compute::doc}]{\sphinxcrossref{\DUrole{doc}{compute}}}} or {\hyperref[\detokenize{fix::doc}]{\sphinxcrossref{\DUrole{doc}{fix}}}} command.  When they
invoke the atom\sphinxhyphen{}style variable, only atoms in the group are included
in the formula evaluation.  The variable evaluates to 0.0 for atoms
not in the group.


\bigskip\hrule\bigskip



\subsubsection{Numbers, constants, and thermo keywords}
\label{\detokenize{variable:numbers-constants-and-thermo-keywords}}
\sphinxAtStartPar
Numbers can contain digits, scientific notation
(3.0e20,3.0e\sphinxhyphen{}20,3.0E20,3.0E\sphinxhyphen{}20), and leading minus signs.

\sphinxAtStartPar
Constants are set at compile time and cannot be changed. \sphinxstyleemphasis{PI} will
return the number 3.14159265358979323846; \sphinxstyleemphasis{on}, \sphinxstyleemphasis{true} or \sphinxstyleemphasis{yes} will
return 1.0; \sphinxstyleemphasis{off}, \sphinxstyleemphasis{false} or \sphinxstyleemphasis{no} will return 0.0; \sphinxstyleemphasis{version} will
return a numeric version code of the current LAMMPS version (e.g.
version 2 Sep 2015 will return the number 20150902). The corresponding
value for newer versions of LAMMPS will be larger, for older versions
of LAMMPS will be smaller. This can be used to have input scripts
adapt automatically to LAMMPS versions, when non\sphinxhyphen{}backwards compatible
syntax changes are introduced. Here is an illustrative example (which
will not work, since the \sphinxstyleemphasis{version} has been introduced more recently):

\begin{sphinxVerbatim}[commandchars=\\\{\}]
\PYG{k}{if}\PYG{+w}{ }\PYG{n+nv}{\PYGZdl{}(}\PYG{n+nv}{version\PYGZlt{}20140513}\PYG{n+nv}{)}\PYG{+w}{ }\PYG{k}{then}\PYG{+w}{ }\PYG{l+s}{\PYGZdq{}}\PYG{l+s}{communicate vel yes}\PYG{l+s}{\PYGZdq{}}\PYG{+w}{ }\PYG{k}{else}\PYG{+w}{ }\PYG{l+s}{\PYGZdq{}}\PYG{l+s}{comm\PYGZus{}modify vel yes}\PYG{l+s}{\PYGZdq{}}
\end{sphinxVerbatim}

\sphinxAtStartPar
The thermo keywords allowed in a formula are those defined by the
{\hyperref[\detokenize{thermo_style::doc}]{\sphinxcrossref{\DUrole{doc}{thermo\_style custom}}}} command.  Thermo keywords
that require a {\hyperref[\detokenize{compute::doc}]{\sphinxcrossref{\DUrole{doc}{compute}}}} to calculate their values such
as “temp” or “press”, use computes stored and invoked by the
{\hyperref[\detokenize{thermo_style::doc}]{\sphinxcrossref{\DUrole{doc}{thermo\_style}}}} command.  This means that you can
only use those keywords in a variable if the style you are using with
the thermo\_style command (and the thermo keywords associated with that
style) also define and use the needed compute.  Note that some thermo
keywords use a compute indirectly to calculate their value (e.g. the
enthalpy keyword uses temp, pe, and pressure).  If a variable is
evaluated directly in an input script (not during a run), then the
values accessed by the thermo keyword must be current.  See the
discussion below about “Variable Accuracy”.


\bigskip\hrule\bigskip



\subsubsection{Math Operators}
\label{\detokenize{variable:math-operators}}
\sphinxAtStartPar
Math operators are written in the usual way, where the “x” and “y” in
the examples can themselves be arbitrarily complex formulas, as in the
examples above.  In this syntax, “x” and “y” can be scalar values or
per\sphinxhyphen{}atom vectors.  For example, “ke/natoms” is the division of two
scalars, where “vy+vz” is the element\sphinxhyphen{}by\sphinxhyphen{}element sum of two per\sphinxhyphen{}atom
vectors of y and z velocities.

\sphinxAtStartPar
Operators are evaluated left to right and have the usual C\sphinxhyphen{}style
precedence: unary minus and unary logical NOT operator “!” have the
highest precedence, exponentiation “\textasciicircum{}” is next; multiplication and
division and the modulo operator “\%” are next; addition and
subtraction are next; the 4 relational operators “\textless{}”, “\textless{}=”, “\textgreater{}”, and
“\textgreater{}=” are next; the two remaining relational operators “==” and “!=”
are next; then the logical AND operator “\&\&”; and finally the logical
OR operator “||” and logical XOR (exclusive or) operator “|\textasciicircum{}” have the
lowest precedence.  Parenthesis can be used to group one or more
portions of a formula and/or enforce a different order of evaluation
than what would occur with the default precedence.

\begin{sphinxadmonition}{note}{Note:}
\sphinxAtStartPar
Because a unary minus is higher precedence than exponentiation,
the formula “\sphinxhyphen{}2\textasciicircum{}2” will evaluate to 4, not \sphinxhyphen{}4.  This convention is
compatible with some programming languages, but not others.  As
mentioned, this behavior can be easily overridden with parenthesis;
the formula “\sphinxhyphen{}(2\textasciicircum{}2)” will evaluate to \sphinxhyphen{}4.
\end{sphinxadmonition}

\sphinxAtStartPar
The 6 relational operators return either a 1.0 or 0.0 depending on
whether the relationship between x and y is TRUE or FALSE.  For
example the expression x\textless{}10.0 in an atom\sphinxhyphen{}style variable formula will
return 1.0 for all atoms whose x\sphinxhyphen{}coordinate is less than 10.0, and 0.0
for the others.  The logical AND operator will return 1.0 if both its
arguments are non\sphinxhyphen{}zero, else it returns 0.0.  The logical OR operator
will return 1.0 if either of its arguments is non\sphinxhyphen{}zero, else it
returns 0.0.  The logical XOR operator will return 1.0 if one of its
arguments is zero and the other non\sphinxhyphen{}zero, else it returns 0.0.  The
logical NOT operator returns 1.0 if its argument is 0.0, else it
returns 0.0.

\sphinxAtStartPar
These relational and logical operators can be used as a masking or
selection operation in a formula.  For example, the number of atoms
whose properties satisfy one or more criteria could be calculated by
taking the returned per\sphinxhyphen{}atom vector of ones and zeroes and passing it
to the {\hyperref[\detokenize{compute_reduce::doc}]{\sphinxcrossref{\DUrole{doc}{compute reduce}}}} command.


\bigskip\hrule\bigskip



\subsubsection{Math Functions}
\label{\detokenize{variable:math-functions}}
\sphinxAtStartPar
Math functions are specified as keywords followed by one or more
parenthesized arguments “x”, “y”, “z”, each of which can themselves be
arbitrarily complex formulas.  In this syntax, the arguments can
represent scalar values or global vectors or per\sphinxhyphen{}atom vectors.  In the
latter case, the math operation is performed on each element of the
vector.  For example, “sqrt(natoms)” is the sqrt() of a scalar, where
“sqrt(y*z)” yields a per\sphinxhyphen{}atom vector with each element being the
sqrt() of the product of one atom’s y and z coordinates.

\sphinxAtStartPar
Most of the math functions perform obvious operations.  The ln() is
the natural log; log() is the base 10 log.

\sphinxAtStartPar
The random(x,y,z) function takes 3 arguments: x = lo, y = hi, and z =
seed.  It generates a uniform random number between lo and hi.  The
normal(x,y,z) function also takes 3 arguments: x = mu, y = sigma, and
z = seed.  It generates a Gaussian variate centered on mu with
variance sigma\textasciicircum{}2.  In both cases the seed is used the first time the
internal random number generator is invoked, to initialize it.  For
equal\sphinxhyphen{}style and vector\sphinxhyphen{}style variables, every processor uses the same
seed so that they each generate the same sequence of random numbers.
For atom\sphinxhyphen{}style variables, a unique seed is created for each processor,
based on the specified seed.  This effectively generates a different
random number for each atom being looped over in the atom\sphinxhyphen{}style
variable.

\begin{sphinxadmonition}{note}{Note:}
\sphinxAtStartPar
Internally, there is just one random number generator for all
equal\sphinxhyphen{}style and vector\sphinxhyphen{}style variables and another one for all
atom\sphinxhyphen{}style variables.  If you define multiple variables (of each
style) which use the random() or normal() math functions, then the
internal random number generators will only be initialized once, which
means only one of the specified seeds will determine the sequence of
generated random numbers.
\end{sphinxadmonition}

\sphinxAtStartPar
The ceil(), floor(), and round() functions are those in the C math
library.  Ceil() is the smallest integer not less than its argument.
Floor() if the largest integer not greater than its argument.  Round()
is the nearest integer to its argument.

\sphinxAtStartPar
\DUrole{versionmodified,added}{Added in version 7Feb2024.}

\sphinxAtStartPar
The ternary(x,y,z) function is the equivalent of the ternary operator
(? and :) in C or C++.  It takes 3 arguments.  The first argument is a
conditional.  The result of the function is y if x evaluates to true
(non\sphinxhyphen{}zero).  The result is z if x evaluates to false (zero).

\sphinxAtStartPar
The ramp(x,y) function uses the current timestep to generate a value
linearly interpolated between the specified x,y values over the course
of a run, according to this formula:
\begin{sphinxalltt}
value = x + (y\sphinxhyphen{}x) * (timestep\sphinxhyphen{}startstep) / (stopstep\sphinxhyphen{}startstep)
\end{sphinxalltt}

\sphinxAtStartPar
The run begins on startstep and ends on stopstep.  Startstep and
stopstep can span multiple runs, using the \sphinxstyleemphasis{start} and \sphinxstyleemphasis{stop} keywords
of the {\hyperref[\detokenize{run::doc}]{\sphinxcrossref{\DUrole{doc}{run}}}} command.  See the {\hyperref[\detokenize{run::doc}]{\sphinxcrossref{\DUrole{doc}{run}}}} command for
details of how to do this.  If called in between runs or during a
{\hyperref[\detokenize{run::doc}]{\sphinxcrossref{\DUrole{doc}{run 0}}}} command, the ramp(x,y) function will return the
value of x.

\sphinxAtStartPar
The stagger(x,y) function uses the current timestep to generate a new
timestep.  X,y \textgreater{} 0 and x \textgreater{} y are required.  The generated timesteps
increase in a staggered fashion, as the sequence
x,x+y,2x,2x+y,3x,3x+y,etc.  For any current timestep, the next
timestep in the sequence is returned.  Thus if stagger(1000,100) is
used in a variable by the {\hyperref[\detokenize{dump_modify::doc}]{\sphinxcrossref{\DUrole{doc}{dump\_modify every}}}}
command, it will generate the sequence of output timesteps:

\begin{sphinxVerbatim}[commandchars=\\\{\}]
100,1000,1100,2000,2100,3000,etc
\end{sphinxVerbatim}

\sphinxAtStartPar
The logfreq(x,y,z) function uses the current timestep to generate a
new timestep.  X,y,z \textgreater{} 0 and y \textless{} z are required.  The generated
timesteps are on a base\sphinxhyphen{}z logarithmic scale, starting with x, and the
y value is how many of the z\sphinxhyphen{}1 possible timesteps within one
logarithmic interval are generated.  I.e. the timesteps follow the
sequence
x,2x,3x,…y*x,x*z,2x*z,3x*z,…y*x*z,x*z\textasciicircum{}2,2x*z\textasciicircum{}2,etc.  For
any current timestep, the next timestep in the sequence is returned.
Thus if logfreq(100,4,10) is used in a variable by the
{\hyperref[\detokenize{dump_modify::doc}]{\sphinxcrossref{\DUrole{doc}{dump\_modify every}}}} command, it will generate this
sequence of output timesteps:

\begin{sphinxVerbatim}[commandchars=\\\{\}]
100,200,300,400,1000,2000,3000,4000,10000,20000,etc
\end{sphinxVerbatim}

\sphinxAtStartPar
The logfreq2(x,y,z) function is similar to logfreq, except a single
logarithmic interval is divided into y equally\sphinxhyphen{}spaced timesteps and
all of them are output.  Y \textless{} z is not required.  Thus, if
logfreq2(100,18,10) is used in a variable by the {\hyperref[\detokenize{dump_modify::doc}]{\sphinxcrossref{\DUrole{doc}{dump\_modify
every}}}} command, then the interval between 100 and 1000
is divided as 900/18 = 50 steps, and it will generate the sequence of
output timesteps:

\begin{sphinxVerbatim}[commandchars=\\\{\}]
100,150,200,...950,1000,1500,2000,...9500,10000,15000,etc
\end{sphinxVerbatim}

\sphinxAtStartPar
The logfreq3(x,y,z) function generates y points between x and z (inclusive),
that are separated by a multiplicative ratio: (z/x)\textasciicircum{}(1/(y\sphinxhyphen{}1)). Constraints
are: x,z \textgreater{} 0, y \textgreater{} 1, z\sphinxhyphen{}x \textgreater{}= y\sphinxhyphen{}1. For eg., if logfreq3(10,25,1000) is used in
a variable by the {\hyperref[\detokenize{fix_print::doc}]{\sphinxcrossref{\DUrole{doc}{fix print}}}} command, then the interval
between 10 and 1000 is divided into 24 parts with a multiplicative
separation of \textasciitilde{}1.21, and it will generate the following sequence of output
timesteps:

\begin{sphinxVerbatim}[commandchars=\\\{\}]
10, 13, 15, 18, 22, 27, 32,...384, 465, 563, 682, 826, 1000
\end{sphinxVerbatim}

\sphinxAtStartPar
The stride(x,y,z) function uses the current timestep to generate a new
timestep.  X,y \textgreater{}= 0 and z \textgreater{} 0 and x \textless{}= y are required.  The generated
timesteps increase in increments of z, from x to y, i.e. it generates
the sequence x,x+z,x+2z,…,y.  If y\sphinxhyphen{}x is not a multiple of z, then
similar to the way a for loop operates, the last value will be one
that does not exceed y.  For any current timestep, the next timestep
in the sequence is returned.  Thus if stride(1000,2000,100) is used
in a variable by the {\hyperref[\detokenize{dump_modify::doc}]{\sphinxcrossref{\DUrole{doc}{dump\_modify every}}}} command, it
will generate the sequence of output timesteps:

\begin{sphinxVerbatim}[commandchars=\\\{\}]
1000,1100,1200, ... ,1900,2000
\end{sphinxVerbatim}

\sphinxAtStartPar
The stride2(x,y,z,a,b,c) function is similar to the stride() function
except it generates two sets of strided timesteps, one at a coarser
level and one at a finer level.  Thus it is useful for debugging,
e.g. to produce output every timestep at the point in simulation when
a problem occurs.  X,y \textgreater{}= 0 and z \textgreater{} 0 and x \textless{}= y are required, as are
a,b \textgreater{}= 0 and c \textgreater{} 0 and a \textless{} b.  Also, a \textgreater{}= x and b \textless{}= y are required so
that the second stride is inside the first.  The generated timesteps
increase in increments of z, starting at x, until a is reached.  At
that point the timestep increases in increments of c, from a to b,
then after b, increments by z are resumed until y is reached.  For any
current timestep, the next timestep in the sequence is returned.  Thus
if stride2(1000,2000,100,1350,1360,1) is used in a variable by the
{\hyperref[\detokenize{dump_modify::doc}]{\sphinxcrossref{\DUrole{doc}{dump\_modify every}}}} command, it will generate the
sequence of output timesteps:

\begin{sphinxVerbatim}[commandchars=\\\{\}]
1000,1100,1200,1300,1350,1351,1352, ... 1359,1360,1400,1500, ... ,2000
\end{sphinxVerbatim}

\sphinxAtStartPar
The vdisplace(x,y) function takes 2 arguments: x = value0 and y =
velocity, and uses the elapsed time to change the value by a linear
displacement due to the applied velocity over the course of a run,
according to this formula:
\begin{sphinxalltt}
value = value0 + velocity*(timestep\sphinxhyphen{}startstep)*dt
\end{sphinxalltt}

\sphinxAtStartPar
where dt = the timestep size.

\sphinxAtStartPar
The run begins on startstep.  Startstep can span multiple runs, using
the \sphinxstyleemphasis{start} keyword of the {\hyperref[\detokenize{run::doc}]{\sphinxcrossref{\DUrole{doc}{run}}}} command.  See the {\hyperref[\detokenize{run::doc}]{\sphinxcrossref{\DUrole{doc}{run}}}} command for details of how to do this.  Note that the
{\hyperref[\detokenize{thermo_style::doc}]{\sphinxcrossref{\DUrole{doc}{thermo\_style}}}} keyword elaplong = timestep\sphinxhyphen{}startstep.
If used between runs this function will return
the value according to the end of the last run or the value of x if
used before \sphinxstyleemphasis{any} runs.  This function assumes the length of the time
step does not change and thus may not be used in combination with
{\hyperref[\detokenize{fix_dt_reset::doc}]{\sphinxcrossref{\DUrole{doc}{fix dt/reset}}}}.

\sphinxAtStartPar
The swiggle(x,y,z) and cwiggle(x,y,z) functions each take 3 arguments:
x = value0, y = amplitude, z = period.  They use the elapsed time to
oscillate the value by a sin() or cos() function over the course of a
run, according to one of these formulas, where omega = 2 PI / period:
\begin{sphinxalltt}
value = value0 + Amplitude * sin(omega*(timestep\sphinxhyphen{}startstep)*dt)
value = value0 + Amplitude * (1 \sphinxhyphen{} cos(omega*(timestep\sphinxhyphen{}startstep)*dt))
\end{sphinxalltt}

\sphinxAtStartPar
where dt = the timestep size.

\sphinxAtStartPar
The run begins on startstep.  Startstep can span multiple runs, using
the \sphinxstyleemphasis{start} keyword of the {\hyperref[\detokenize{run::doc}]{\sphinxcrossref{\DUrole{doc}{run}}}} command.  See the {\hyperref[\detokenize{run::doc}]{\sphinxcrossref{\DUrole{doc}{run}}}} command for details of how to do this.  Note that the
{\hyperref[\detokenize{thermo_style::doc}]{\sphinxcrossref{\DUrole{doc}{thermo\_style}}}} keyword elaplong = timestep\sphinxhyphen{}startstep.
If used between runs these functions will return
the value according to the end of the last run or the value of x if
used before \sphinxstyleemphasis{any} runs.  These functions assume the length of the time
step does not change and thus may not be used in combination with
{\hyperref[\detokenize{fix_dt_reset::doc}]{\sphinxcrossref{\DUrole{doc}{fix dt/reset}}}}.


\bigskip\hrule\bigskip



\subsubsection{Group and Region Functions}
\label{\detokenize{variable:group-and-region-functions}}
\sphinxAtStartPar
Group functions are specified as keywords followed by one or two
parenthesized arguments.  The first argument \sphinxstyleemphasis{ID} is the group\sphinxhyphen{}ID.
The \sphinxstyleemphasis{dim} argument, if it exists, is \sphinxstyleemphasis{x} or \sphinxstyleemphasis{y} or \sphinxstyleemphasis{z}.  The \sphinxstyleemphasis{dir}
argument, if it exists, is \sphinxstyleemphasis{xmin}, \sphinxstyleemphasis{xmax}, \sphinxstyleemphasis{ymin}, \sphinxstyleemphasis{ymax}, \sphinxstyleemphasis{zmin}, or
\sphinxstyleemphasis{zmax}.  The \sphinxstyleemphasis{dimdim} argument, if it exists, is \sphinxstyleemphasis{xx} or \sphinxstyleemphasis{yy} or \sphinxstyleemphasis{zz}
or \sphinxstyleemphasis{xy} or \sphinxstyleemphasis{yz} or \sphinxstyleemphasis{xz}.

\sphinxAtStartPar
The group function count() is the number of atoms in the group.  The
group functions mass() and charge() are the total mass and charge of
the group.  Xcm() and vcm() return components of the position and
velocity of the center of mass of the group.  Fcm() returns a
component of the total force on the group of atoms.  Bound() returns
the min/max of a particular coordinate for all atoms in the group.
Gyration() computes the radius\sphinxhyphen{}of\sphinxhyphen{}gyration of the group of atoms.  See
the {\hyperref[\detokenize{compute_gyration::doc}]{\sphinxcrossref{\DUrole{doc}{compute gyration}}}} command for a definition
of the formula.  Angmom() returns components of the angular momentum
of the group of atoms around its center of mass.  Torque() returns
components of the torque on the group of atoms around its center of
mass, based on current forces on the atoms.  Inertia() returns one of
6 components of the symmetric inertia tensor of the group of atoms
around its center of mass, ordered as Ixx,Iyy,Izz,Ixy,Iyz,Ixz.
Omega() returns components of the angular velocity of the group of
atoms around its center of mass.

\sphinxAtStartPar
Region functions are specified exactly the same way as group functions
except they take an extra final argument \sphinxstyleemphasis{IDR} which is the region ID.
The function is computed for all atoms that are in both the group and
the region.  If the group is “all”, then the only criteria for atom
inclusion is that it be in the region.


\bigskip\hrule\bigskip



\subsubsection{Special Functions}
\label{\detokenize{variable:special-functions}}
\sphinxAtStartPar
Special functions take specific kinds of arguments, meaning their
arguments cannot be formulas themselves.

\sphinxAtStartPar
The sum(x), min(x), max(x), ave(x), trap(x), slope(x), sort(x), and
rsort(x) functions each take 1 argument which is of the form “c\_ID” or
“c\_ID{[}N{]}” or “f\_ID” or “f\_ID{[}N{]}” or “v\_name”.  The first two are
computes and the second two are fixes; the ID in the reference should be
replaced by the ID of a compute or fix defined elsewhere in the input
script.  The compute or fix must produce either a global vector or
array.  If it produces a global vector, then the notation without “{[}N{]}”
should be used.  If it produces a global array, then the notation with
“{[}N{]}” should be used, where N is an integer, to specify which column of
the global array is being referenced.  The last form of argument
“v\_name” is for a vector\sphinxhyphen{}style variable where “name” is replaced by the
name of the variable.

\sphinxAtStartPar
The sum(x), min(x), max(x), ave(x), trap(x), and slope(x) functions
operate on a global vector of inputs and reduce it to a single scalar
value.  This is analogous to the operation of the {\hyperref[\detokenize{compute_reduce::doc}]{\sphinxcrossref{\DUrole{doc}{compute reduce}}}} command, which performs similar operations on per\sphinxhyphen{}atom
and local vectors.

\sphinxAtStartPar
The sort(x) and rsort(x) functions operate on a global vector of inputs
and return a global vector of the same length.

\sphinxAtStartPar
The sum() function calculates the sum of all the vector elements.  The
min() and max() functions find the minimum and maximum element
respectively.  The ave() function is the same as sum() except that it
divides the result by the length of the vector.

\sphinxAtStartPar
The trap() function is the same as sum() except the first and last
elements are multiplied by a weighting factor of 1/2 when performing
the sum.  This effectively implements an integration via the
trapezoidal rule on the global vector of data.  I.e. consider a set of
points, equally spaced by 1 in their x coordinate: (1,V1), (2,V2),
…, (N,VN), where the Vi are the values in the global vector of
length N.  The integral from 1 to N of these points is trap().  When
appropriately normalized by the timestep size, this function is useful
for calculating integrals of time\sphinxhyphen{}series data, like that generated by
the {\hyperref[\detokenize{fix_ave_correlate::doc}]{\sphinxcrossref{\DUrole{doc}{fix ave/correlate}}}} command.

\sphinxAtStartPar
The slope() function uses linear regression to fit a line to the set
of points, equally spaced by 1 in their x coordinate: (1,V1), (2,V2),
…, (N,VN), where the Vi are the values in the global vector of
length N.  The returned value is the slope of the line.  If the line
has a single point or is vertical, it returns 1.0e20.

\sphinxAtStartPar
\DUrole{versionmodified,added}{Added in version 27June2024.}

\sphinxAtStartPar
The sort(x) and rsort(x) functions sort the data of the input vector by
their numeric value: sort(x) sorts in ascending order, rsort(x) sorts
in descending order.

\sphinxAtStartPar
The gmask(x) function takes 1 argument which is a group ID.  It
can only be used in atom\sphinxhyphen{}style variables.  It returns a 1 for
atoms that are in the group, and a 0 for atoms that are not.

\sphinxAtStartPar
The rmask(x) function takes 1 argument which is a region ID.  It can
only be used in atom\sphinxhyphen{}style variables.  It returns a 1 for atoms that
are in the geometric region, and a 0 for atoms that are not.

\sphinxAtStartPar
The grmask(x,y) function takes 2 arguments.  The first is a group ID,
and the second is a region ID.  It can only be used in atom\sphinxhyphen{}style
variables.  It returns a 1 for atoms that are in both the group and
region, and a 0 for atoms that are not in both.

\sphinxAtStartPar
The next(x) function takes 1 argument which is a variable ID (not
“v\_foo”, just “foo”).  It must be for a file\sphinxhyphen{}style or atomfile\sphinxhyphen{}style
variable.  Each time the next() function is invoked (i.e. each time
the equal\sphinxhyphen{}style or atom\sphinxhyphen{}style variable is evaluated), the following
steps occur.

\sphinxAtStartPar
For file\sphinxhyphen{}style variables, the current string value stored by the
file\sphinxhyphen{}style variable is converted to a numeric value and returned by
the function.  And the next string value in the file is read and
stored.  Note that if the line previously read from the file was not a
numeric string, then it will typically evaluate to 0.0, which is
likely not what you want.

\sphinxAtStartPar
For atomfile\sphinxhyphen{}style variables, the current per\sphinxhyphen{}atom values stored by
the atomfile\sphinxhyphen{}style variable are returned by the function.  And the
next set of per\sphinxhyphen{}atom values in the file is read and stored.

\sphinxAtStartPar
Since file\sphinxhyphen{}style and atomfile\sphinxhyphen{}style variables read and store the first
line of the file or first set of per\sphinxhyphen{}atoms values when they are
defined in the input script, these are the value(s) that will be
returned the first time the next() function is invoked.  If next() is
invoked more times than there are lines or sets of lines in the file,
the variable is deleted, similar to how the {\hyperref[\detokenize{next::doc}]{\sphinxcrossref{\DUrole{doc}{next}}}} command
operates.

\sphinxAtStartPar
The is\_file(name) function is a test whether \sphinxstyleemphasis{name} is a (readable) file
and returns 1 in this case, otherwise it returns 0.  For that \sphinxstyleemphasis{name}
is taken as a literal string and must not have any blanks in it.

\sphinxAtStartPar
The is\_os(name) function is a test whether \sphinxstyleemphasis{name} is part of the OS
information that LAMMPS collects and provides in the
{\hyperref[\detokenize{Developer_platform:_CPPv4N9LAMMPS_NS8platform7os_infoEv}]{\sphinxcrossref{\sphinxcode{\sphinxupquote{platform::os\_info()}}}}} function.
The argument \sphinxstyleemphasis{name} is interpreted as a regular expression as documented
for the {\hyperref[\detokenize{Developer_utils:_CPPv4N9LAMMPS_NS5utils8strmatchERKNSt6stringERKNSt6stringE}]{\sphinxcrossref{\sphinxcode{\sphinxupquote{utils::strmatch()}}}}}
function. This allows to adapt LAMMPS inputs to the OS it runs on:

\begin{sphinxVerbatim}[commandchars=\\\{\}]
\PYG{k}{if}\PYG{+w}{ }\PYG{n+nv}{\PYGZdl{}(}\PYG{n+nv}{is\PYGZus{}os}\PYG{n+nv}{(}\PYG{n+nv}{\PYGZca{}Windows}\PYG{n+nv}{)}\PYG{n+nv}{)}\PYG{+w}{ }\PYG{k}{then}\PYG{+w}{ }\PYG{o}{\PYGZam{}}
\PYG{+w}{  }\PYG{l+s}{\PYGZdq{}}\PYG{l+s}{shell copy \PYGZdl{}\PYGZob{}input\PYGZus{}dir\PYGZcb{}\PYGZbs{}some\PYGZus{}file.txt .}\PYG{l+s}{\PYGZdq{}}\PYG{+w}{ }\PYG{o}{\PYGZam{}}
\PYG{k}{else}\PYG{+w}{ }\PYG{o}{\PYGZam{}}
\PYG{+w}{  }\PYG{l+s}{\PYGZdq{}}\PYG{l+s}{shell cp \PYGZdl{}\PYGZob{}input\PYGZus{}dir\PYGZcb{}/some\PYGZus{}file.txt .}\PYG{l+s}{\PYGZdq{}}
\end{sphinxVerbatim}

\sphinxAtStartPar
The extract\_setting(name) function enables access to basic settings for
the LAMMPS executable and the running simulation via calling the
{\hyperref[\detokenize{Library_properties:_CPPv422lammps_extract_settingPvPKc}]{\sphinxcrossref{\sphinxcode{\sphinxupquote{lammps\_extract\_setting()}}}}} library function.  For example, the
number of processors (MPI ranks) being used by the simulation or the MPI
process ID (for this processor) can be queried, or the number of atom
types, bond types and so on. For the full list of available keywords
\sphinxstyleemphasis{name} and their meaning, see the documentation for extract\_setting()
via the link in this paragraph.

\sphinxAtStartPar
The label2type(kind,label) function converts type labels into numeric
types, using label maps created by the {\hyperref[\detokenize{labelmap::doc}]{\sphinxcrossref{\DUrole{doc}{labelmap}}}} or
{\hyperref[\detokenize{read_data::doc}]{\sphinxcrossref{\DUrole{doc}{read\_data}}}} commands.  The first argument is the label
map kind (atom, bond, angle, dihedral, or improper) and the second
argument is the label.  The function returns the corresponding numeric
type or triggers an error if the queried label does not exist.

\sphinxAtStartPar
\DUrole{versionmodified,added}{Added in version 15Jun2023.}

\sphinxAtStartPar
The is\_typelabel(kind,label) function has the same arguments as
label2type(), but returns 1 if the type label has been assigned,
otherwise it returns 0.  This function can be used to check if a
particular type label already exists in the simulation.

\sphinxAtStartPar
\DUrole{versionmodified,added}{Added in version 29Aug2024.}

\sphinxAtStartPar
The is\_timeout() function returns 1 when the {\hyperref[\detokenize{timer::doc}]{\sphinxcrossref{\DUrole{doc}{timer timeout}}}} has expired otherwise it returns 0.  This function can be used
to check inputs in combination with the {\hyperref[\detokenize{if::doc}]{\sphinxcrossref{\DUrole{doc}{if command}}}} to
execute commands after the timer has expired. Example:

\begin{sphinxVerbatim}[commandchars=\\\{\}]
\PYG{k}{variable }\PYG{n+nv+nvIdentifier}{timeout}\PYG{+w}{ }\PYG{n}{equal}\PYG{+w}{ }\PYG{n}{is\PYGZus{}timeout}\PYG{n+nv}{(}\PYG{n+nv}{)}
\PYG{k}{timer}\PYG{+w}{ }\PYG{n}{timeout}\PYG{+w}{ }\PYG{l+m}{0:10:00}\PYG{+w}{ }\PYG{n}{every}\PYG{+w}{ }\PYG{l+m}{10}
\PYG{k}{run}\PYG{+w}{ }\PYG{l+m}{10000}
\PYG{k}{if}\PYG{+w}{ }\PYG{n+nv}{\PYGZdl{}\PYGZob{}}\PYG{n+nv}{timeout}\PYG{n+nv}{\PYGZcb{}}\PYG{+w}{ }\PYG{k}{then}\PYG{+w}{ }\PYG{l+s}{\PYGZdq{}}\PYG{l+s}{print \PYGZsq{}Timer has expired\PYGZsq{}}\PYG{l+s}{\PYGZdq{}}
\end{sphinxVerbatim}


\bigskip\hrule\bigskip



\subsubsection{Feature Functions}
\label{\detokenize{variable:feature-functions}}
\sphinxAtStartPar
Feature functions allow probing of the running LAMMPS executable for
whether specific features are available, active, or defined.  All 3 of
the functions take two arguments, a \sphinxstyleemphasis{category} and a category\sphinxhyphen{}specific
second argument.  Both are strings and thus cannot be formulas
themselves; only \$\sphinxhyphen{}style immediate variable expansion is possible.
The return value of the functions is either 1.0 or 0.0 depending on
whether the function evaluates to true or false, respectively.

\sphinxAtStartPar
The \sphinxstyleemphasis{is\_available(category,name)} function queries whether a specific
feature is available in the LAMMPS executable that is being run, i.e
whether it was included or enabled at compile time.

\sphinxAtStartPar
This supports the following categories: \sphinxstyleemphasis{command}, \sphinxstyleemphasis{compute}, \sphinxstyleemphasis{fix},
\sphinxstyleemphasis{pair\_style} and \sphinxstyleemphasis{feature}.  For all the categories except \sphinxstyleemphasis{feature}
the \sphinxstyleemphasis{name} is a style name, e.g. \sphinxstyleemphasis{nve} for the \sphinxstyleemphasis{fix} category.  Note
that many LAMMPS input script commands such as \sphinxstyleemphasis{create\_atoms} are
actually instances of a command style which LAMMPS defines, as opposed
to built\sphinxhyphen{}in commands.  For all of these styles except \sphinxstyleemphasis{command},
appending of active suffixes is also tried before reporting failure.

\sphinxAtStartPar
The \sphinxstyleemphasis{feature} category checks the availability of the following
compile\sphinxhyphen{}time enabled features: GZIP support, PNG support, JPEG support,
FFMPEG support, and C++ exceptions for error handling. Corresponding
names are \sphinxstyleemphasis{gzip}, \sphinxstyleemphasis{png}, \sphinxstyleemphasis{jpeg}, \sphinxstyleemphasis{ffmpeg} and \sphinxstyleemphasis{exceptions}.

\sphinxAtStartPar
Example: Only dump in a given format if the compiled binary supports it.

\begin{sphinxVerbatim}[commandchars=\\\{\}]
\PYG{k}{if}\PYG{+w}{ }\PYG{l+s}{\PYGZdq{}}\PYG{l+s}{\PYGZdl{}(is\PYGZus{}available(feature,png))}\PYG{l+s}{\PYGZdq{}}\PYG{+w}{ }\PYG{k}{then}\PYG{+w}{ }\PYG{l+s}{\PYGZdq{}}\PYG{l+s}{print \PYGZsq{}PNG supported\PYGZsq{}}\PYG{l+s}{\PYGZdq{}}\PYG{+w}{ }\PYG{k}{else}\PYG{+w}{ }\PYG{l+s}{\PYGZdq{}}\PYG{l+s}{print \PYGZsq{}PNG not supported\PYGZsq{}}\PYG{l+s}{\PYGZdq{}}
\PYG{k}{if}\PYG{+w}{ }\PYG{l+s}{\PYGZdq{}}\PYG{l+s}{\PYGZdl{}(is\PYGZus{}available(feature,ffmpeg)}\PYG{l+s}{\PYGZdq{}}\PYG{+w}{ }\PYG{k}{then}\PYG{+w}{ }\PYG{l+s}{\PYGZdq{}}\PYG{l+s}{dump 3 all movie 25 movie.mp4 type type zoom 1.6 adiam 1.0}\PYG{l+s}{\PYGZdq{}}
\end{sphinxVerbatim}

\sphinxAtStartPar
The \sphinxstyleemphasis{is\_active(category,feature)} function queries whether a specific
feature is currently active within LAMMPS.  The features are grouped
by categories.  Supported categories and features are:
\begin{itemize}
\item {} 
\sphinxAtStartPar
\sphinxstyleemphasis{package}: features = \sphinxstyleemphasis{gpu} or \sphinxstyleemphasis{intel} or \sphinxstyleemphasis{kokkos} or \sphinxstyleemphasis{omp}

\item {} 
\sphinxAtStartPar
\sphinxstyleemphasis{newton}: features = \sphinxstyleemphasis{pair} or \sphinxstyleemphasis{bond} or \sphinxstyleemphasis{any}

\item {} 
\sphinxAtStartPar
\sphinxstyleemphasis{pair}: features = \sphinxstyleemphasis{single} or \sphinxstyleemphasis{respa} or \sphinxstyleemphasis{manybody} or \sphinxstyleemphasis{tail} or \sphinxstyleemphasis{shift}

\item {} 
\sphinxAtStartPar
\sphinxstyleemphasis{comm\_style}: features = \sphinxstyleemphasis{brick} or \sphinxstyleemphasis{tiled}

\item {} 
\sphinxAtStartPar
\sphinxstyleemphasis{min\_style}: features = a minimizer style name

\item {} 
\sphinxAtStartPar
\sphinxstyleemphasis{run\_style}: features = a run style name

\item {} 
\sphinxAtStartPar
\sphinxstyleemphasis{atom\_style}: features = an atom style name

\item {} 
\sphinxAtStartPar
\sphinxstyleemphasis{pair\_style}: features = a pair style name

\item {} 
\sphinxAtStartPar
\sphinxstyleemphasis{bond\_style}: features = a bond style name

\item {} 
\sphinxAtStartPar
\sphinxstyleemphasis{angle\_style}: features = an angle style name

\item {} 
\sphinxAtStartPar
\sphinxstyleemphasis{dihedral\_style}: features = a dihedral style name

\item {} 
\sphinxAtStartPar
\sphinxstyleemphasis{improper\_style}: features = an improper style name

\item {} 
\sphinxAtStartPar
\sphinxstyleemphasis{kspace\_style}: features = a kspace style name

\end{itemize}

\sphinxAtStartPar
Most of the settings are self\sphinxhyphen{}explanatory.  For the \sphinxstyleemphasis{package}
category, a package may have been included in the LAMMPS build, but
not have enabled by any input script command, and hence be inactive.
The \sphinxstyleemphasis{single} feature in the \sphinxstyleemphasis{pair} category checks whether the
currently defined pair style supports a Pair::single() function as
needed by compute group/group and others features or LAMMPS.
Similarly, the \sphinxstyleemphasis{respa} feature checks whether the inner/middle/outer
mode of r\sphinxhyphen{}RESPA is supported by the current pair style.

\sphinxAtStartPar
For the categories with \sphinxstyleemphasis{style} in their name, only a single instance
of the style is ever active at any time in a LAMMPS simulation.  Thus
the check is whether the currently active style matches the specified
name.  This check is also done using suffix flags, if available and
enabled.

\sphinxAtStartPar
Example 1: Disable use of suffix for PPPM when using GPU package
(i.e. run it on the CPU concurrently while running the pair style on
the GPU), but do use the suffix otherwise (e.g. with OPENMP).

\begin{sphinxVerbatim}[commandchars=\\\{\}]
\PYG{k}{pair\PYGZus{}style}\PYG{+w}{ }\PYG{n}{lj}\PYG{o}{/}\PYG{n}{cut}\PYG{o}{/}\PYG{n}{coul}\PYG{o}{/}\PYG{n}{long}\PYG{+w}{ }\PYG{l+m}{14.0}
\PYG{k}{if}\PYG{+w}{ }\PYG{n+nv}{\PYGZdl{}(}\PYG{n+nv}{is\PYGZus{}active}\PYG{n+nv}{(}\PYG{n+nv}{package,gpu}\PYG{n+nv}{)}\PYG{n+nv}{)}\PYG{+w}{ }\PYG{k}{then}\PYG{+w}{ }\PYG{l+s}{\PYGZdq{}}\PYG{l+s}{suffix off}\PYG{l+s}{\PYGZdq{}}
\PYG{k}{kspace\PYGZus{}style}\PYG{+w}{ }\PYG{n}{pppm}
\end{sphinxVerbatim}

\sphinxAtStartPar
Example 2: Use r\sphinxhyphen{}RESPA with inner/outer cutoff, if supported by the
current pair style, otherwise fall back to using r\sphinxhyphen{}RESPA with simply
the pair keyword and reducing the outer time step.

\begin{sphinxVerbatim}[commandchars=\\\{\}]
\PYG{k}{timestep}\PYG{+w}{ }\PYG{n+nv}{\PYGZdl{}(}\PYG{n+nv}{2.0*}\PYG{n+nv}{(}\PYG{n+nv}{1.0+2.0*is\PYGZus{}active}\PYG{n+nv}{(}\PYG{n+nv}{pair,respa}\PYG{n+nv}{)}\PYG{n+nv}{)}\PYG{n+nv}{)}
\PYG{k}{if}\PYG{+w}{ }\PYG{n+nv}{\PYGZdl{}(}\PYG{n+nv}{is\PYGZus{}active}\PYG{n+nv}{(}\PYG{n+nv}{pair,respa}\PYG{n+nv}{)}\PYG{n+nv}{)}\PYG{+w}{ }\PYG{k}{then}\PYG{+w}{ }\PYG{l+s}{\PYGZdq{}}\PYG{l+s}{run\PYGZus{}style respa 4 3 2 2 improper 1 inner 2 5.5 7.0 outer 3 kspace 4}\PYG{l+s}{\PYGZdq{}}\PYG{+w}{ }\PYG{k}{else}\PYG{+w}{ }\PYG{l+s}{\PYGZdq{}}\PYG{l+s}{run\PYGZus{}style respa 3 3 2 improper 1 pair 2 kspace 3}\PYG{l+s}{\PYGZdq{}}
\end{sphinxVerbatim}

\sphinxAtStartPar
The \sphinxstyleemphasis{is\_defined(category,id)} function checks whether an instance of a
style or variable with a specific ID or name is currently defined
within LAMMPS.  The supported categories are \sphinxstyleemphasis{compute}, \sphinxstyleemphasis{dump},
\sphinxstyleemphasis{fix}, \sphinxstyleemphasis{group}, \sphinxstyleemphasis{region}, and \sphinxstyleemphasis{variable}.  Each of these styles (as
well as the variable command) can be specified multiple times within
LAMMPS, each with a unique \sphinxstyleemphasis{id}.  This function checks whether the
specified \sphinxstyleemphasis{id} exists.  For category \sphinxstyleemphasis{variable”, the *id} is the
variable name.


\bigskip\hrule\bigskip



\subsubsection{Atom Values and Vectors}
\label{\detokenize{variable:atom-values-and-vectors}}
\sphinxAtStartPar
Atom values take an integer argument I from 1 to N, where I is the
atom\sphinxhyphen{}ID, e.g. x{[}243{]}, which means use the x coordinate of the atom
with ID = 243.  Or they can take a variable name, specified as v\_name,
where name is the name of the variable, like x{[}v\_myIndex{]}.  The
variable can be of any style except \sphinxstyleemphasis{vector} or \sphinxstyleemphasis{atom} or \sphinxstyleemphasis{atomfile}
variables.  The variable is evaluated and the result is expected to be
numeric and is cast to an integer (i.e. 3.4 becomes 3), to use an
index, which must be a value from 1 to N.  Note that a “formula”
cannot be used as the argument between the brackets, e.g. x{[}243+10{]}
or x{[}v\_myIndex+1{]} are not allowed.  To do this a single variable can
be defined that contains the needed formula.

\sphinxAtStartPar
Note that the 0 \textless{} atom\sphinxhyphen{}ID \textless{}= N, where N is the largest atom ID
in the system.  If an ID is specified for an atom that does not
currently exist, then the generated value is 0.0.

\sphinxAtStartPar
Atom vectors generate one value per atom, so that a reference like
“vx” means the x\sphinxhyphen{}component of each atom’s velocity will be used when
evaluating the variable.

\sphinxAtStartPar
The meaning of the different atom values and vectors is mostly
self\sphinxhyphen{}explanatory.  \sphinxstyleemphasis{Mol} refers to the molecule ID of an atom, and is
only defined if an {\hyperref[\detokenize{atom_style::doc}]{\sphinxcrossref{\DUrole{doc}{atom\_style}}}} is being used that
defines molecule IDs.

\sphinxAtStartPar
Note that many other atom attributes can be used as inputs to a
variable by using the {\hyperref[\detokenize{compute_property_atom::doc}]{\sphinxcrossref{\DUrole{doc}{compute property/atom}}}} command and then referencing that compute.


\bigskip\hrule\bigskip



\subsubsection{Custom atom properties}
\label{\detokenize{variable:custom-atom-properties}}
\sphinxAtStartPar
\DUrole{versionmodified,added}{Added in version 7Feb2024.}

\sphinxAtStartPar
Custom atom properties refer to per\sphinxhyphen{}atom integer and floating point
vectors or arrays that have been added via the {\hyperref[\detokenize{fix_property_atom::doc}]{\sphinxcrossref{\DUrole{doc}{fix property/atom}}}} command.  When that command is used specific
names are given to each attribute which are the “name” portion of
these references.  References beginning with \sphinxstyleemphasis{i} and \sphinxstyleemphasis{d} refer to
integer and floating point properties respectively.  Per\sphinxhyphen{}atom vectors
are referenced by \sphinxstyleemphasis{i\_name} and \sphinxstyleemphasis{d\_name}; per\sphinxhyphen{}atom arrays are
referenced by \sphinxstyleemphasis{i2\_name} and \sphinxstyleemphasis{d2\_name}.

\sphinxAtStartPar
The various allowed references to integer custom atom properties in
the variable formulas for equal\sphinxhyphen{}, vector\sphinxhyphen{}, and atom\sphinxhyphen{}style variables
are listed in the following table.  References to floating point
custom atom properties are the same; just replace the leading “i” with
“d”.


\begin{savenotes}\sphinxattablestart
\sphinxthistablewithglobalstyle
\centering
\begin{tabulary}{\linewidth}[t]{TTT}
\sphinxtoprule
\sphinxtableatstartofbodyhook
\sphinxAtStartPar
equal
&
\sphinxAtStartPar
i\_name{[}I{]}
&
\sphinxAtStartPar
element of per\sphinxhyphen{}atom vector (I = atom ID)
\\
\sphinxhline
\sphinxAtStartPar
equal
&
\sphinxAtStartPar
i2\_name{[}I{]}{[}J{]}
&
\sphinxAtStartPar
element of per\sphinxhyphen{}atom array (I = atom ID)
\\
\sphinxhline&&\\
\sphinxhline
\sphinxAtStartPar
vector
&
\sphinxAtStartPar
i\_name{[}I{]}
&
\sphinxAtStartPar
element of per\sphinxhyphen{}atom vector (I = atom ID)
\\
\sphinxhline
\sphinxAtStartPar
vector
&
\sphinxAtStartPar
i2\_name{[}I{]}{[}J{]}
&
\sphinxAtStartPar
element of per\sphinxhyphen{}atom array (I = atom ID)
\\
\sphinxhline&&\\
\sphinxhline
\sphinxAtStartPar
atom
&
\sphinxAtStartPar
i\_name
&
\sphinxAtStartPar
per\sphinxhyphen{}atom vector
\\
\sphinxhline
\sphinxAtStartPar
atom
&
\sphinxAtStartPar
i2\_name{[}I{]}
&
\sphinxAtStartPar
column of per\sphinxhyphen{}atom array
\\
\sphinxbottomrule
\end{tabulary}
\sphinxtableafterendhook\par
\sphinxattableend\end{savenotes}

\sphinxAtStartPar
The I and J indices in these custom atom property references can be
integers or can be a variable name, specified as v\_name, where name is
the name of the variable.  The rules for this syntax are the same as
for indices in the “Atom Values and Vectors” discussion above.


\bigskip\hrule\bigskip



\subsubsection{Compute References}
\label{\detokenize{variable:compute-references}}
\sphinxAtStartPar
Compute references access quantities calculated by a {\hyperref[\detokenize{compute::doc}]{\sphinxcrossref{\DUrole{doc}{compute}}}}.  The ID in the reference should be replaced by the ID of a
compute defined elsewhere in the input script.

\sphinxAtStartPar
As discussed on the page for the {\hyperref[\detokenize{compute::doc}]{\sphinxcrossref{\DUrole{doc}{compute}}}} command,
computes can produce global, per\sphinxhyphen{}atom, local, and per\sphinxhyphen{}grid values.
Only global and per\sphinxhyphen{}atom values can be used in a variable.  Computes
can also produce scalars (global only), vectors, and arrays.  See the
doc pages for individual computes to see what different kinds of data
they produce.

\sphinxAtStartPar
An equal\sphinxhyphen{}style variable can only use scalar values, either from global
or per\sphinxhyphen{}atom data.  In the case of per\sphinxhyphen{}atom data, this would be a value
for a specific atom.

\sphinxAtStartPar
A vector\sphinxhyphen{}style variable can use scalar values (same as for equal\sphinxhyphen{}style
variables), or global vectors of values.  The latter can also be a
column of a global array.

\sphinxAtStartPar
Atom\sphinxhyphen{}style variables can use scalar values (same as for equal\sphinxhyphen{}style
variables), or per\sphinxhyphen{}atom vectors of values.  The latter can also be a
column of a per\sphinxhyphen{}atom array.

\sphinxAtStartPar
The various allowed compute references in the variable formulas for
equal\sphinxhyphen{}, vector\sphinxhyphen{}, and atom\sphinxhyphen{}style variables are listed in the following
table:


\begin{savenotes}\sphinxattablestart
\sphinxthistablewithglobalstyle
\centering
\begin{tabulary}{\linewidth}[t]{TTT}
\sphinxtoprule
\sphinxtableatstartofbodyhook
\sphinxAtStartPar
equal
&
\sphinxAtStartPar
c\_ID
&
\sphinxAtStartPar
global scalar
\\
\sphinxhline
\sphinxAtStartPar
equal
&
\sphinxAtStartPar
c\_ID{[}I{]}
&
\sphinxAtStartPar
element of global vector
\\
\sphinxhline
\sphinxAtStartPar
equal
&
\sphinxAtStartPar
c\_ID{[}I{]}{[}J{]}
&
\sphinxAtStartPar
element of global array
\\
\sphinxhline
\sphinxAtStartPar
equal
&
\sphinxAtStartPar
C\_ID{[}I{]}
&
\sphinxAtStartPar
element of per\sphinxhyphen{}atom vector (I = atom ID)
\\
\sphinxhline
\sphinxAtStartPar
equal
&
\sphinxAtStartPar
C\_ID{[}I{]}{[}J{]}
&
\sphinxAtStartPar
element of per\sphinxhyphen{}atom array (I = atom ID)
\\
\sphinxhline&&\\
\sphinxhline
\sphinxAtStartPar
vector
&
\sphinxAtStartPar
c\_ID
&
\sphinxAtStartPar
global vector
\\
\sphinxhline
\sphinxAtStartPar
vector
&
\sphinxAtStartPar
c\_ID{[}I{]}
&
\sphinxAtStartPar
column of global array
\\
\sphinxhline&&\\
\sphinxhline
\sphinxAtStartPar
atom
&
\sphinxAtStartPar
c\_ID
&
\sphinxAtStartPar
per\sphinxhyphen{}atom vector
\\
\sphinxhline
\sphinxAtStartPar
atom
&
\sphinxAtStartPar
c\_ID{[}I{]}
&
\sphinxAtStartPar
column of per\sphinxhyphen{}atom array
\\
\sphinxbottomrule
\end{tabulary}
\sphinxtableafterendhook\par
\sphinxattableend\end{savenotes}

\sphinxAtStartPar
Note that if an equal\sphinxhyphen{}style variable formula wishes to access per\sphinxhyphen{}atom
data from a compute, it must use capital “C” as the ID prefix and not
lower\sphinxhyphen{}case “c”.

\sphinxAtStartPar
Also note that if a vector\sphinxhyphen{} or atom\sphinxhyphen{}style variable formula needs to
access a scalar value from a compute (i.e. the 5 kinds of values in
the first 5 lines of the table), it can not do so directly.  Instead,
it can use a reference to an equal\sphinxhyphen{}style variable which stores the
scalar value from the compute.

\sphinxAtStartPar
The I and J indices in these compute references can be integers or can
be a variable name, specified as v\_name, where name is the name of the
variable.  The rules for this syntax are the same as for indices in
the “Atom Values and Vectors” discussion above.

\sphinxAtStartPar
If a variable containing a compute is evaluated directly in an input
script (not during a run), then the values accessed by the compute
should be current.  See the discussion below about “Variable
Accuracy”.


\bigskip\hrule\bigskip



\subsubsection{Fix References}
\label{\detokenize{variable:fix-references}}
\sphinxAtStartPar
Fix references access quantities calculated by a {\hyperref[\detokenize{fix::doc}]{\sphinxcrossref{\DUrole{doc}{fix}}}}.
The ID in the reference should be replaced by the ID of a fix defined
elsewhere in the input script.

\sphinxAtStartPar
As discussed on the page for the {\hyperref[\detokenize{fix::doc}]{\sphinxcrossref{\DUrole{doc}{fix}}}} command, fixes can
produce global, per\sphinxhyphen{}atom, local, and per\sphinxhyphen{}grid values.  Only global and
per\sphinxhyphen{}atom values can be used in a variable.  Fixes can also produce
scalars (global only), vectors, and arrays.  See the doc pages for
individual fixes to see what different kinds of data they produce.

\sphinxAtStartPar
An equal\sphinxhyphen{}style variable can only use scalar values, either from global
or per\sphinxhyphen{}atom data.  In the case of per\sphinxhyphen{}atom data, this would be a value
for a specific atom.

\sphinxAtStartPar
A vector\sphinxhyphen{}style variable can use scalar values (same as for equal\sphinxhyphen{}style
variables), or global vectors of values.  The latter can also be a
column of a global array.

\sphinxAtStartPar
Atom\sphinxhyphen{}style variables can use scalar values (same as for equal\sphinxhyphen{}style
variables), or per\sphinxhyphen{}atom vectors of values.  The latter can also be a
column of a per\sphinxhyphen{}atom array.

\sphinxAtStartPar
The allowed fix references in variable formulas for equal\sphinxhyphen{}, vector\sphinxhyphen{},
and atom\sphinxhyphen{}style variables are listed in the following table:


\begin{savenotes}\sphinxattablestart
\sphinxthistablewithglobalstyle
\centering
\begin{tabulary}{\linewidth}[t]{TTT}
\sphinxtoprule
\sphinxtableatstartofbodyhook
\sphinxAtStartPar
equal
&
\sphinxAtStartPar
f\_ID
&
\sphinxAtStartPar
global scalar
\\
\sphinxhline
\sphinxAtStartPar
equal
&
\sphinxAtStartPar
f\_ID{[}I{]}
&
\sphinxAtStartPar
element of global vector
\\
\sphinxhline
\sphinxAtStartPar
equal
&
\sphinxAtStartPar
f\_ID{[}I{]}{[}J{]}
&
\sphinxAtStartPar
element of global array
\\
\sphinxhline
\sphinxAtStartPar
equal
&
\sphinxAtStartPar
F\_ID{[}I{]}
&
\sphinxAtStartPar
element of per\sphinxhyphen{}atom vector (I = atom ID)
\\
\sphinxhline
\sphinxAtStartPar
equal
&
\sphinxAtStartPar
F\_ID{[}I{]}{[}J{]}
&
\sphinxAtStartPar
element of per\sphinxhyphen{}atom array (I = atom ID)
\\
\sphinxhline&&\\
\sphinxhline
\sphinxAtStartPar
vector
&
\sphinxAtStartPar
f\_ID
&
\sphinxAtStartPar
global vector
\\
\sphinxhline
\sphinxAtStartPar
vector
&
\sphinxAtStartPar
f\_ID{[}I{]}
&
\sphinxAtStartPar
column of global array
\\
\sphinxhline&&\\
\sphinxhline
\sphinxAtStartPar
atom
&
\sphinxAtStartPar
f\_ID
&
\sphinxAtStartPar
per\sphinxhyphen{}atom vector
\\
\sphinxhline
\sphinxAtStartPar
atom
&
\sphinxAtStartPar
f\_ID{[}I{]}
&
\sphinxAtStartPar
column of per\sphinxhyphen{}atom array
\\
\sphinxbottomrule
\end{tabulary}
\sphinxtableafterendhook\par
\sphinxattableend\end{savenotes}

\sphinxAtStartPar
Note that if an equal\sphinxhyphen{}style variable formula wishes to access per\sphinxhyphen{}atom
data from a fix, it must use capital “F” as the ID prefix and not
lower\sphinxhyphen{}case “f”.

\sphinxAtStartPar
Also note that if a vector\sphinxhyphen{} or atom\sphinxhyphen{}style variable formula needs to
access a scalar value from a fix (i.e. the 5 kinds of values in the
first 5 lines of the table), it can not do so directly.  Instead, it
can use a reference to an equal\sphinxhyphen{}style variable which stores the scalar
value from the fix.

\sphinxAtStartPar
The I and J indices in these fix references can be integers or can be
a variable name, specified as v\_name, where name is the name of the
variable.  The rules for this syntax are the same as for indices in
the “Atom Values and Vectors” discussion above.

\sphinxAtStartPar
Note that some fixes only generate quantities on certain timesteps.
If a variable attempts to access the fix on non\sphinxhyphen{}allowed timesteps, an
error is generated.  For example, the {\hyperref[\detokenize{fix_ave_time::doc}]{\sphinxcrossref{\DUrole{doc}{fix ave/time}}}}
command may only generate averaged quantities every 100 steps.  See
the doc pages for individual fix commands for details.

\sphinxAtStartPar
If a variable containing a fix is evaluated directly in an input
script (not during a run), then the values accessed by the fix should
be current.  See the discussion below about “Variable Accuracy”.


\bigskip\hrule\bigskip



\subsubsection{Variable References}
\label{\detokenize{variable:variable-references}}
\sphinxAtStartPar
Variable references access quantities stored or calculated by other
variables, which will cause those variables to be evaluated.  The name
in the reference should be replaced by the name of a variable defined
elsewhere in the input script.

\sphinxAtStartPar
As discussed on this doc page, equal\sphinxhyphen{}style variables generate a single
global numeric value, vector\sphinxhyphen{}style variables generate a vector of
global numeric values, and atom\sphinxhyphen{}style and atomfile\sphinxhyphen{}style variables
generate a per\sphinxhyphen{}atom vector of numeric values.  All other variables
store one or more strings.

\sphinxAtStartPar
The formula for an equal\sphinxhyphen{}style variable can use any style of variable
including a vector\_style or atom\sphinxhyphen{}style or atomfile\sphinxhyphen{}style.  For these
3 styles, a subscript must be used to access a single value from
the vector\sphinxhyphen{}, atom\sphinxhyphen{}, or atomfile\sphinxhyphen{}style variable.  If a string\sphinxhyphen{}storing
variable is used, the string is converted to a numeric value.  Note
that this will typically produce a 0.0 if the string is not a numeric
string, which is likely not what you want.

\sphinxAtStartPar
The formula for a vector\sphinxhyphen{}style variable can use any style of variable,
including atom\sphinxhyphen{}style or atomfile\sphinxhyphen{}style variables.  For these 2 styles,
a subscript must be used to access a single value from the atom\sphinxhyphen{}, or
atomfile\sphinxhyphen{}style variable.

\sphinxAtStartPar
The formula for an atom\sphinxhyphen{}style variable can use any style of variable,
including other atom\sphinxhyphen{}style or atomfile\sphinxhyphen{}style variables.  If it uses a
vector\sphinxhyphen{}style variable, a subscript must be used to access a single
value from the vector\sphinxhyphen{}style variable.

\sphinxAtStartPar
The allowed variable references in variable formulas for equal\sphinxhyphen{},
vector\sphinxhyphen{}, and atom\sphinxhyphen{}style variables are listed in the following table.
Note that there is no ambiguity as to what a reference means, since
referenced variables produce only a global scalar or global vector or
per\sphinxhyphen{}atom vector.


\begin{savenotes}\sphinxattablestart
\sphinxthistablewithglobalstyle
\centering
\begin{tabulary}{\linewidth}[t]{TTT}
\sphinxtoprule
\sphinxtableatstartofbodyhook
\sphinxAtStartPar
equal
&
\sphinxAtStartPar
v\_name
&
\sphinxAtStartPar
global scalar from an equal\sphinxhyphen{}style variable
\\
\sphinxhline
\sphinxAtStartPar
equal
&
\sphinxAtStartPar
v\_name{[}I{]}
&
\sphinxAtStartPar
element of global vector from a vector\sphinxhyphen{}style variable
\\
\sphinxhline
\sphinxAtStartPar
equal
&
\sphinxAtStartPar
v\_name{[}I{]}
&
\sphinxAtStartPar
element of per\sphinxhyphen{}atom vector (I = atom ID) from an atom\sphinxhyphen{} or atomfile\sphinxhyphen{}style variable
\\
\sphinxhline&&\\
\sphinxhline
\sphinxAtStartPar
vector
&
\sphinxAtStartPar
v\_name
&
\sphinxAtStartPar
global scalar from an equal\sphinxhyphen{}style variable
\\
\sphinxhline
\sphinxAtStartPar
vector
&
\sphinxAtStartPar
v\_name
&
\sphinxAtStartPar
global vector from a vector\sphinxhyphen{}style variable
\\
\sphinxhline
\sphinxAtStartPar
vector
&
\sphinxAtStartPar
v\_name{[}I{]}
&
\sphinxAtStartPar
element of global vector from a vector\sphinxhyphen{}style variable
\\
\sphinxhline
\sphinxAtStartPar
vector
&
\sphinxAtStartPar
v\_name{[}I{]}
&
\sphinxAtStartPar
element of per\sphinxhyphen{}atom vector (I = atom ID) from an atom\sphinxhyphen{} or atomfile\sphinxhyphen{}style variable
\\
\sphinxhline&&\\
\sphinxhline
\sphinxAtStartPar
atom
&
\sphinxAtStartPar
v\_name
&
\sphinxAtStartPar
global scalar from an equal\sphinxhyphen{}style variable
\\
\sphinxhline
\sphinxAtStartPar
atom
&
\sphinxAtStartPar
v\_name
&
\sphinxAtStartPar
per\sphinxhyphen{}atom vector from an atom\sphinxhyphen{}style or atomfile\sphinxhyphen{}style variable
\\
\sphinxhline
\sphinxAtStartPar
atom
&
\sphinxAtStartPar
v\_name{[}I{]}
&
\sphinxAtStartPar
element of global vector from a vector\sphinxhyphen{}style variable
\\
\sphinxhline
\sphinxAtStartPar
atom
&
\sphinxAtStartPar
v\_name{[}I{]}
&
\sphinxAtStartPar
element of per\sphinxhyphen{}atom vector (I = atom ID) from an atom\sphinxhyphen{} or atomfile\sphinxhyphen{}style variable
\\
\sphinxbottomrule
\end{tabulary}
\sphinxtableafterendhook\par
\sphinxattableend\end{savenotes}

\sphinxAtStartPar
For the I index, an integer can be specified or a variable name,
specified as v\_name, where name is the name of the variable.  The
rules for this syntax are the same as for indices in the “Atom Values
and Vectors” discussion above.


\bigskip\hrule\bigskip



\subsubsection{Vector Initialization}
\label{\detokenize{variable:vector-initialization}}
\sphinxAtStartPar
\DUrole{versionmodified,added}{Added in version 15Jun2023.}

\sphinxAtStartPar
\sphinxstyleemphasis{Vector}\sphinxhyphen{}style variables only can be initialized with a special
syntax, instead of using a formula.  The syntax is a bracketed,
comma\sphinxhyphen{}separated syntax like the following:

\begin{sphinxVerbatim}[commandchars=\\\{\}]
\PYG{k}{variable }\PYG{n+nv+nvIdentifier}{myvec}\PYG{+w}{ }\PYG{n}{vector}\PYG{+w}{ }\PYG{n}{[1},3.5,7,10.2\PYG{n}{]}
\end{sphinxVerbatim}

\sphinxAtStartPar
The 3rd argument formula is replaced by the vector values in brackets,
separated by commas.  This example creates a 4\sphinxhyphen{}length vector with
specific numeric values, each of which can be specified as an integer
or floating point value.  Note that while whitespace can be added
before or after individual values, no other mathematical operations
can be specified.  E.g. “3*10” or “3*v\_abc” are not valid vector
elements, nor is “10*{[}1,2,3,4{]}” valid for the entire vector.

\sphinxAtStartPar
Unlike vector variables specified with formulas, this vector variable
is static; its length and values never changes.  Its values can be
used in other commands (including vector\sphinxhyphen{}style variables specified
with formulas) via the usual syntax for accessing individual vector
elements or the entire vector.


\bigskip\hrule\bigskip



\subsection{Immediate Evaluation of Variables}
\label{\detokenize{variable:immediate-evaluation-of-variables}}
\sphinxAtStartPar
If you want an equal\sphinxhyphen{}style variable to be evaluated immediately, it
may be the case that you do not need to define a variable at all.  See
the {\hyperref[\detokenize{Commands_parse::doc}]{\sphinxcrossref{\DUrole{doc}{Commands parse}}}} page for info on how to
use “immediate” variables in an input script, specified as \$(formula)
with parenthesis, where the formula has the same syntax as equal\sphinxhyphen{}style
variables described on this page.  This effectively evaluates a
formula immediately without using the variable command to define a
named variable.

\sphinxAtStartPar
More generally, there is a difference between referencing a variable
with a leading \$ sign (e.g. \$x or \$\{abc\}) versus with a leading “v\_”
(e.g. v\_x or v\_abc).  The former can be used in any input script
command, including a variable command.  The input script parser
evaluates the reference variable immediately and substitutes its value
into the command.  As explained on the {\hyperref[\detokenize{Commands_parse::doc}]{\sphinxcrossref{\DUrole{doc}{Commands parse}}}} doc page, you can also use un\sphinxhyphen{}named “immediate”
variables for this purpose.  For example, a string like this
\$((xlo+xhi)/2+sqrt(v\_area)) in an input script command evaluates the
string between the parenthesis as an equal\sphinxhyphen{}style variable formula.

\sphinxAtStartPar
Referencing a variable with a leading “v\_” is an optional or required
kind of argument for some commands (e.g. the {\hyperref[\detokenize{fix_ave_chunk::doc}]{\sphinxcrossref{\DUrole{doc}{fix ave/chunk}}}} or {\hyperref[\detokenize{dump::doc}]{\sphinxcrossref{\DUrole{doc}{dump custom}}}} or {\hyperref[\detokenize{thermo_style::doc}]{\sphinxcrossref{\DUrole{doc}{thermo\_style}}}} commands) if you wish it to evaluate a variable
periodically during a run.  It can also be used in a variable formula
if you wish to reference a second variable.  The second variable will
be evaluated whenever the first variable is evaluated.

\sphinxAtStartPar
As an example, suppose you use this command in your input script to
define the variable “v” as

\begin{sphinxVerbatim}[commandchars=\\\{\}]
\PYG{k}{variable }\PYG{n+nv+nvIdentifier}{v}\PYG{+w}{ }\PYG{n}{equal}\PYG{+w}{ }\PYG{n}{vol}
\end{sphinxVerbatim}

\sphinxAtStartPar
before a run where the simulation box size changes.  You might think
this will assign the initial volume to the variable “v”.  That is not
the case.  Rather it assigns a formula which evaluates the volume
(using the thermo\_style keyword “vol”) to the variable “v”.  If you
use the variable “v” in some other command like {\hyperref[\detokenize{fix_ave_time::doc}]{\sphinxcrossref{\DUrole{doc}{fix ave/time}}}} then the current volume of the box will be evaluated
continuously during the run.

\sphinxAtStartPar
If you want to store the initial volume of the system, you can do it
this way:

\begin{sphinxVerbatim}[commandchars=\\\{\}]
\PYG{k}{variable }\PYG{n+nv+nvIdentifier}{v}\PYG{+w}{ }\PYG{n}{equal}\PYG{+w}{ }\PYG{n}{vol}
\PYG{k}{variable }\PYG{n+nv+nvIdentifier}{v0}\PYG{+w}{ }\PYG{n}{equal}\PYG{+w}{ }\PYG{n+nv}{\PYGZdl{}v}
\end{sphinxVerbatim}

\sphinxAtStartPar
The second command will force “v” to be evaluated (yielding the
initial volume) and assign that value to the variable “v0”.  Thus the
command

\begin{sphinxVerbatim}[commandchars=\\\{\}]
\PYG{k}{thermo\PYGZus{}style}\PYG{+w}{ }\PYG{n}{custom}\PYG{+w}{ }\PYG{n}{step}\PYG{+w}{ }\PYG{n}{v\PYGZus{}v}\PYG{+w}{ }\PYG{n}{v\PYGZus{}v0}
\end{sphinxVerbatim}

\sphinxAtStartPar
would print out both the current and initial volume periodically
during the run.

\sphinxAtStartPar
Note that it is a mistake to enclose a variable formula in double
quotes if it contains variables preceded by \$ signs.  For example,

\begin{sphinxVerbatim}[commandchars=\\\{\}]
\PYG{k}{variable }\PYG{n+nv+nvIdentifier}{vratio}\PYG{+w}{ }\PYG{n}{equal}\PYG{+w}{ }\PYG{l+s}{\PYGZdq{}}\PYG{l+s}{\PYGZdl{}\PYGZob{}vfinal\PYGZcb{}/\PYGZdl{}\PYGZob{}v0\PYGZcb{}}\PYG{l+s}{\PYGZdq{}}
\end{sphinxVerbatim}

\sphinxAtStartPar
This is because the quotes prevent variable substitution (explained on
the {\hyperref[\detokenize{Commands_parse::doc}]{\sphinxcrossref{\DUrole{doc}{Commands parse}}}} doc page), and thus an error
will occur when the formula for “vratio” is evaluated later.


\bigskip\hrule\bigskip



\subsection{Variable Accuracy}
\label{\detokenize{variable:variable-accuracy}}
\sphinxAtStartPar
Obviously, LAMMPS attempts to evaluate variables which contain
formulas (\sphinxstyleemphasis{equal} and \sphinxstyleemphasis{vector} and \sphinxstyleemphasis{atom} style variables)
accurately whenever the evaluation is performed.  Depending on what is
included in the formula, this may require invoking a {\hyperref[\detokenize{compute::doc}]{\sphinxcrossref{\DUrole{doc}{compute}}}}, either directly or indirectly via a thermo keyword, or
accessing a value previously calculated by a compute, or accessing a
value calculated and stored by a {\hyperref[\detokenize{fix::doc}]{\sphinxcrossref{\DUrole{doc}{fix}}}}.  If the compute is
one that calculates the energy or pressure of the system, then the
corresponding energy or virial quantities need to be tallied during
the evaluation of the interatomic potentials (pair, bond, etc) on any
timestep that the variable needs the tallies.  An input script can
also request variables be evaluated before or after or in between
runs, e.g. by including them in a {\hyperref[\detokenize{print::doc}]{\sphinxcrossref{\DUrole{doc}{print}}}} command.

\sphinxAtStartPar
LAMMPS keeps track of all of this as it performs a {\hyperref[\detokenize{run::doc}]{\sphinxcrossref{\DUrole{doc}{run}}}} or
{\hyperref[\detokenize{minimize::doc}]{\sphinxcrossref{\DUrole{doc}{minimize}}}} simulation, as well as in between
simulations.  An error will be generated if you attempt to evaluate a
variable when LAMMPS knows it cannot produce accurate values.  For
example, if a {\hyperref[\detokenize{thermo_style::doc}]{\sphinxcrossref{\DUrole{doc}{thermo\_style custom}}}} command prints
a variable which accesses values stored by a {\hyperref[\detokenize{fix_ave_time::doc}]{\sphinxcrossref{\DUrole{doc}{fix ave/time}}}} command and the timesteps on which thermo output is
generated are not multiples of the averaging frequency used in the fix
command, then an error will occur.

\sphinxAtStartPar
However, there are two special cases to be aware when a variable
requires invocation of a compute (directly or indirectly).  The first
is if the variable is evaluated before the first {\hyperref[\detokenize{run::doc}]{\sphinxcrossref{\DUrole{doc}{run}}}} or
{\hyperref[\detokenize{minimize::doc}]{\sphinxcrossref{\DUrole{doc}{minimize}}}} command in the input script.  In this case,
LAMMPS will generate an error.  This is because many computes require
initializations which have not yet taken place.  One example is the
calculation of degrees of freedom for temperature computes.  Another
example are the computes mentioned above which require tallying of
energy or virial quantities; these values are not tallied until the
first simulation begins.

\sphinxAtStartPar
The second special case is when a variable that depends on a compute
is evaluated in between {\hyperref[\detokenize{run::doc}]{\sphinxcrossref{\DUrole{doc}{run}}}} or {\hyperref[\detokenize{minimize::doc}]{\sphinxcrossref{\DUrole{doc}{minimize}}}}
commands.  It is possible for other input script commands issued
following the previous run, but before the variable is evaluated, to
change the system.  For example, the {\hyperref[\detokenize{delete_atoms::doc}]{\sphinxcrossref{\DUrole{doc}{delete\_atoms}}}}
command could be used to remove atoms.  Since the compute will not
re\sphinxhyphen{}initialize itself until the next simulation or it may depend on
energy/virial computations performed before the system was changed, it
will potentially generate an incorrect answer when evaluated.  Note
that LAMMPS will not generate an error in this case; the evaluated
variable may simply be incorrect.

\sphinxAtStartPar
The way to get around both of these special cases is to perform a
0\sphinxhyphen{}timestep run before evaluating the variable.  For example, these
commands

\begin{sphinxVerbatim}[commandchars=\\\{\}]
\PYG{c}{\PYGZsh{} delete\PYGZus{}atoms random fraction 0.5 yes all NULL 49839}
\PYG{c}{\PYGZsh{} run 0 post no}
\PYG{k}{variable }\PYG{n+nv+nvIdentifier}{t}\PYG{+w}{ }\PYG{n}{equal}\PYG{+w}{ }\PYG{n}{temp}\PYG{+w}{    }\PYG{c}{\PYGZsh{} this thermo keyword invokes a temperature compute}
\PYG{k}{print}\PYG{+w}{ }\PYG{l+s}{\PYGZdq{}}\PYG{l+s}{Temperature of system = \PYGZdl{}t}\PYG{l+s}{\PYGZdq{}}
\PYG{k}{run}\PYG{+w}{ }\PYG{l+m}{1000}
\end{sphinxVerbatim}

\sphinxAtStartPar
will generate an error if the “run 1000” command is the first
simulation in the input script.  If there were a previous run, these
commands will print the correct temperature of the system.  But if the
{\hyperref[\detokenize{delete_atoms::doc}]{\sphinxcrossref{\DUrole{doc}{delete\_atoms}}}} command is uncommented, the printed
temperature will be incorrect, because information stored by
temperature compute is no longer valid.

\sphinxAtStartPar
Both these issues are resolved, if the “run 0” command is uncommented.
This is because the “run 0” simulation will initialize (or
re\sphinxhyphen{}initialize) the temperature compute correctly.


\bigskip\hrule\bigskip



\subsection{Restrictions}
\label{\detokenize{variable:restrictions}}
\sphinxAtStartPar
Indexing any formula element by global atom ID, such as an atom value,
requires the {\hyperref[\detokenize{atom_style::doc}]{\sphinxcrossref{\DUrole{doc}{atom style}}}} to use a global mapping in
order to look up the vector indices.  By default, only atom styles
with molecular information create global maps.  The {\hyperref[\detokenize{atom_modify::doc}]{\sphinxcrossref{\DUrole{doc}{atom\_modify
map}}}} command can override the default, e.g. for
atomic\sphinxhyphen{}style atom styles.

\sphinxAtStartPar
All \sphinxstyleemphasis{universe}\sphinxhyphen{} and \sphinxstyleemphasis{uloop}\sphinxhyphen{}style variables defined in an input
script must have the same number of values.


\subsection{Related commands}
\label{\detokenize{variable:related-commands}}
\sphinxAtStartPar
{\hyperref[\detokenize{next::doc}]{\sphinxcrossref{\DUrole{doc}{next}}}}, {\hyperref[\detokenize{jump::doc}]{\sphinxcrossref{\DUrole{doc}{jump}}}}, {\hyperref[\detokenize{include::doc}]{\sphinxcrossref{\DUrole{doc}{include}}}},
{\hyperref[\detokenize{temper::doc}]{\sphinxcrossref{\DUrole{doc}{temper}}}}, {\hyperref[\detokenize{fix_print::doc}]{\sphinxcrossref{\DUrole{doc}{fix print}}}}, {\hyperref[\detokenize{print::doc}]{\sphinxcrossref{\DUrole{doc}{print}}}}


\subsection{Default}
\label{\detokenize{variable:default}}
\sphinxAtStartPar
none

\sphinxstepscope

\index{velocity@\spxentry{velocity}}\ignorespaces 

\section{velocity command}
\label{\detokenize{velocity:velocity-command}}\label{\detokenize{velocity:index-0}}\label{\detokenize{velocity::doc}}

\subsection{Syntax}
\label{\detokenize{velocity:syntax}}
\begin{sphinxVerbatim}[commandchars=\\\{\}]
\PYG{k}{velocity }\PYG{n+nv+nvIdentifier}{group\PYGZhy{}ID}\PYG{+w}{ }\PYG{n}{style}\PYG{+w}{ }\PYG{n}{args}\PYG{+w}{ }\PYG{n}{keyword}\PYG{+w}{ }\PYG{n}{value}\PYG{+w}{ }\PYG{n}{...}
\end{sphinxVerbatim}
\begin{itemize}
\item {} 
\sphinxAtStartPar
group\sphinxhyphen{}ID = ID of group of atoms whose velocity will be changed

\item {} 
\sphinxAtStartPar
style = \sphinxstyleemphasis{create} or \sphinxstyleemphasis{set} or \sphinxstyleemphasis{scale} or \sphinxstyleemphasis{ramp} or \sphinxstyleemphasis{zero}
\begin{sphinxalltt}
\sphinxstyleemphasis{create} args = temp seed
  temp = temperature value (temperature units)
  seed = random \# seed (positive integer)
\sphinxstyleemphasis{set} args = vx vy vz
  vx,vy,vz = velocity value or NULL (velocity units)
  any of vx,vy,vz can be a variable (see below)
\sphinxstyleemphasis{scale} arg = temp
  temp = temperature value (temperature units)
\sphinxstyleemphasis{ramp} args = vdim vlo vhi dim clo chi
  vdim = \sphinxstyleemphasis{vx} or \sphinxstyleemphasis{vy} or \sphinxstyleemphasis{vz}
  vlo,vhi = lower and upper velocity value (velocity units)
  dim = \sphinxstyleemphasis{x} or \sphinxstyleemphasis{y} or \sphinxstyleemphasis{z}
  clo,chi = lower and upper coordinate bound (distance units)
\sphinxstyleemphasis{zero} arg = \sphinxstyleemphasis{linear} or \sphinxstyleemphasis{angular}
  \sphinxstyleemphasis{linear} = zero the linear momentum
  \sphinxstyleemphasis{angular} = zero the angular momentum
\end{sphinxalltt}

\item {} 
\sphinxAtStartPar
zero or more keyword/value pairs may be appended

\item {} 
\sphinxAtStartPar
keyword = \sphinxstyleemphasis{dist} or \sphinxstyleemphasis{sum} or \sphinxstyleemphasis{mom} or \sphinxstyleemphasis{rot} or \sphinxstyleemphasis{temp} or \sphinxstyleemphasis{bias} or \sphinxstyleemphasis{loop} or \sphinxstyleemphasis{rigid} or \sphinxstyleemphasis{units}
\begin{sphinxalltt}
\sphinxstyleemphasis{dist} value = \sphinxstyleemphasis{uniform} or \sphinxstyleemphasis{gaussian}
\sphinxstyleemphasis{sum} value = \sphinxstyleemphasis{no} or \sphinxstyleemphasis{yes}
\sphinxstyleemphasis{mom} value = \sphinxstyleemphasis{no} or \sphinxstyleemphasis{yes}
\sphinxstyleemphasis{rot} value = \sphinxstyleemphasis{no} or \sphinxstyleemphasis{yes}
\sphinxstyleemphasis{temp} value = temperature compute ID
\sphinxstyleemphasis{bias} value = \sphinxstyleemphasis{no} or \sphinxstyleemphasis{yes}
\sphinxstyleemphasis{loop} value = \sphinxstyleemphasis{all} or \sphinxstyleemphasis{local} or \sphinxstyleemphasis{geom}
\sphinxstyleemphasis{rigid} value = fix\sphinxhyphen{}ID
  fix\sphinxhyphen{}ID = ID of rigid body fix
\sphinxstyleemphasis{units} value = \sphinxstyleemphasis{box} or \sphinxstyleemphasis{lattice}
\end{sphinxalltt}

\end{itemize}


\subsection{Examples}
\label{\detokenize{velocity:examples}}
\begin{sphinxVerbatim}[commandchars=\\\{\}]
\PYG{k}{velocity }\PYG{n+nv+nvIdentifier}{all}\PYG{+w}{ }\PYG{n}{create}\PYG{+w}{ }\PYG{l+m}{300.0}\PYG{+w}{ }\PYG{l+m}{4928459}\PYG{+w}{ }\PYG{n}{rot}\PYG{+w}{ }\PYG{n}{yes}\PYG{+w}{ }\PYG{n}{dist}\PYG{+w}{ }\PYG{n}{gaussian}
\PYG{k}{velocity }\PYG{n+nv+nvIdentifier}{border}\PYG{+w}{ }\PYG{n}{set}\PYG{+w}{ }\PYG{n}{NULL}\PYG{+w}{ }\PYG{l+m}{4.0}\PYG{+w}{ }\PYG{n}{v\PYGZus{}vz}\PYG{+w}{ }\PYG{n}{sum}\PYG{+w}{ }\PYG{n}{yes}\PYG{+w}{ }\PYG{n}{units}\PYG{+w}{ }\PYG{n}{box}
\PYG{k}{velocity }\PYG{n+nv+nvIdentifier}{flow}\PYG{+w}{ }\PYG{n}{scale}\PYG{+w}{ }\PYG{l+m}{300.0}
\PYG{k}{velocity }\PYG{n+nv+nvIdentifier}{flow}\PYG{+w}{ }\PYG{n}{ramp}\PYG{+w}{ }\PYG{n}{vx}\PYG{+w}{ }\PYG{l+m}{0.0}\PYG{+w}{ }\PYG{l+m}{5.0}\PYG{+w}{ }\PYG{n}{y}\PYG{+w}{ }\PYG{l+m}{5}\PYG{+w}{ }\PYG{l+m}{25}\PYG{+w}{ }\PYG{n}{temp}\PYG{+w}{ }\PYG{n}{mytemp}
\PYG{k}{velocity }\PYG{n+nv+nvIdentifier}{all}\PYG{+w}{ }\PYG{n}{zero}\PYG{+w}{ }\PYG{n}{linear}
\end{sphinxVerbatim}


\subsection{Description}
\label{\detokenize{velocity:description}}
\sphinxAtStartPar
Set or change the velocities of a group of atoms in one of several
styles.  For each style, there are required arguments and optional
keyword/value parameters.  Not all options are used by each style.
Each option has a default as listed below.

\sphinxAtStartPar
The \sphinxstyleemphasis{create} style generates an ensemble of velocities using a random
number generator with the specified seed at the specified temperature.

\sphinxAtStartPar
The \sphinxstyleemphasis{set} style sets the velocities of all atoms in the group to the
specified values.  If any component is specified as NULL, then it is
not set.  Any of the vx,vy,vz velocity components can be specified as
an equal\sphinxhyphen{}style or atom\sphinxhyphen{}style {\hyperref[\detokenize{variable::doc}]{\sphinxcrossref{\DUrole{doc}{variable}}}}.  If the value
is a variable, it should be specified as v\_name, where name is the
variable name.  In this case, the variable will be evaluated, and its
value used to determine the velocity component.  Note that if a
variable is used, the velocity it calculates must be in box units, not
lattice units; see the discussion of the \sphinxstyleemphasis{units} keyword below.

\sphinxAtStartPar
Equal\sphinxhyphen{}style variables can specify formulas with various mathematical
functions, and include {\hyperref[\detokenize{thermo_style::doc}]{\sphinxcrossref{\DUrole{doc}{thermo\_style}}}} command
keywords for the simulation box parameters or other parameters.

\sphinxAtStartPar
Atom\sphinxhyphen{}style variables can specify the same formulas as equal\sphinxhyphen{}style
variables but can also include per\sphinxhyphen{}atom values, such as atom
coordinates.  Thus it is easy to specify a spatially\sphinxhyphen{}dependent
velocity field.

\sphinxAtStartPar
The \sphinxstyleemphasis{scale} style computes the current temperature of the group of
atoms and then rescales the velocities to the specified temperature.

\sphinxAtStartPar
The \sphinxstyleemphasis{ramp} style is similar to that used by the {\hyperref[\detokenize{compute_temp_ramp::doc}]{\sphinxcrossref{\DUrole{doc}{compute
temp/ramp}}}} command.  Velocities ramped uniformly
from vlo to vhi are applied to dimension vx, or vy, or vz.  The value
assigned to a particular atom depends on its relative coordinate value
(in dim) from clo to chi.  For the example above, an atom with
y\sphinxhyphen{}coordinate of 10 (1/4 of the way from 5 to 25), would be assigned a
x\sphinxhyphen{}velocity of 1.25 (1/4 of the way from 0.0 to 5.0).  Atoms outside
the coordinate bounds (less than 5 or greater than 25 in this case),
are assigned velocities equal to vlo or vhi (0.0 or 5.0 in this case).

\sphinxAtStartPar
The \sphinxstyleemphasis{zero} style adjusts the velocities of the group of atoms so that
the aggregate linear or angular momentum is zero.  No other changes
are made to the velocities of the atoms.  If the \sphinxstyleemphasis{rigid} option is
specified (see below), then the zeroing is performed on individual
rigid bodies, as defined by the {\hyperref[\detokenize{fix_rigid::doc}]{\sphinxcrossref{\DUrole{doc}{fix rigid or fix rigid/small}}}} commands.  In other words, zero linear will set the
linear momentum of each rigid body to zero, and zero angular will set
the angular momentum of each rigid body to zero.  This is done by
adjusting the velocities of the atoms in each rigid body.

\sphinxAtStartPar
All temperatures specified in the velocity command are in temperature
units; see the {\hyperref[\detokenize{units::doc}]{\sphinxcrossref{\DUrole{doc}{units}}}} command.  The units of velocities and
coordinates depend on whether the \sphinxstyleemphasis{units} keyword is set to \sphinxstyleemphasis{box} or
\sphinxstyleemphasis{lattice}, as discussed below.

\sphinxAtStartPar
For all styles, no atoms are assigned z\sphinxhyphen{}component velocities if the
simulation is 2d; see the {\hyperref[\detokenize{dimension::doc}]{\sphinxcrossref{\DUrole{doc}{dimension}}}} command.


\bigskip\hrule\bigskip


\sphinxAtStartPar
The keyword/value options are used in the following ways by the
various styles.

\sphinxAtStartPar
The \sphinxstyleemphasis{dist} keyword is used by \sphinxstyleemphasis{create}.  The ensemble of generated
velocities can be a \sphinxstyleemphasis{uniform} distribution from some minimum to
maximum value, scaled to produce the requested temperature.  Or it can
be a \sphinxstyleemphasis{gaussian} distribution with a mean of 0.0 and a sigma scaled to
produce the requested temperature.

\sphinxAtStartPar
The \sphinxstyleemphasis{sum} keyword is used by all styles, except \sphinxstyleemphasis{zero}.  The new
velocities will be added to the existing ones if sum = yes, or will
replace them if sum = no.

\sphinxAtStartPar
The \sphinxstyleemphasis{mom} and \sphinxstyleemphasis{rot} keywords are used by \sphinxstyleemphasis{create}.  If mom = yes, the
linear momentum of the newly created ensemble of velocities is zeroed;
if rot = yes, the angular momentum is zeroed.


\bigskip\hrule\bigskip


\sphinxAtStartPar
If specified, the \sphinxstyleemphasis{temp} keyword is used by \sphinxstyleemphasis{create} and \sphinxstyleemphasis{scale} to
specify a {\hyperref[\detokenize{compute::doc}]{\sphinxcrossref{\DUrole{doc}{compute}}}} that calculates temperature in a
desired way, e.g. by first subtracting out a velocity bias, as
discussed on the {\hyperref[\detokenize{Howto_thermostat::doc}]{\sphinxcrossref{\DUrole{doc}{Howto thermostat}}}} doc page.
If this keyword is not specified, \sphinxstyleemphasis{create} and \sphinxstyleemphasis{scale} calculate
temperature using a compute that is defined internally as follows:

\begin{sphinxVerbatim}[commandchars=\\\{\}]
\PYG{k}{compute }\PYG{n+nv+nvIdentifier}{velocity\PYGZus{}temp}\PYG{+w}{ }\PYG{n+nv+nvIdentifier}{group\PYGZhy{}ID}\PYG{+w}{ }\PYG{n}{temp}
\end{sphinxVerbatim}

\sphinxAtStartPar
where group\sphinxhyphen{}ID is the same ID used in the velocity command. i.e. the
group of atoms whose velocity is being altered.  This compute is
deleted when the velocity command is finished.  See the {\hyperref[\detokenize{compute_temp::doc}]{\sphinxcrossref{\DUrole{doc}{compute
temp}}}} command for details.  If the calculated
temperature should have degrees\sphinxhyphen{}of\sphinxhyphen{}freedom removed due to fix
constraints (e.g. SHAKE or rigid\sphinxhyphen{}body constraints), then the
appropriate fix command must be specified before the velocity command
is issued.

\sphinxAtStartPar
The \sphinxstyleemphasis{bias} keyword with a \sphinxstyleemphasis{yes} setting is used by \sphinxstyleemphasis{create} and
\sphinxstyleemphasis{scale}, but only if the \sphinxstyleemphasis{temp} keyword is also used to specify a
{\hyperref[\detokenize{compute::doc}]{\sphinxcrossref{\DUrole{doc}{compute}}}} that calculates temperature in a desired way.
If the temperature compute also calculates a velocity bias, the bias
is subtracted from atom velocities before the \sphinxstyleemphasis{create} and \sphinxstyleemphasis{scale}
operations are performed.  After the operations, the bias is added
back to the atom velocities.  See the {\hyperref[\detokenize{Howto_thermostat::doc}]{\sphinxcrossref{\DUrole{doc}{Howto thermostat}}}} page for more discussion of temperature computes
with biases.  Note that the velocity bias is only applied to atoms in
the temperature compute specified with the \sphinxstyleemphasis{temp} keyword.

\sphinxAtStartPar
As an example, assume atoms are currently streaming in a flow
direction (which could be separately initialized with the \sphinxstyleemphasis{ramp}
style), and you wish to initialize their thermal velocity to a desired
temperature.  In this context thermal velocity means the per\sphinxhyphen{}particle
velocity that remains when the streaming velocity is subtracted.  This
can be done using the \sphinxstyleemphasis{create} style with the \sphinxstyleemphasis{temp} keyword
specifying the ID of a {\hyperref[\detokenize{compute_temp_ramp::doc}]{\sphinxcrossref{\DUrole{doc}{compute temp/ramp}}}} or
{\hyperref[\detokenize{compute_temp_profile::doc}]{\sphinxcrossref{\DUrole{doc}{compute temp/profile}}}} command, and the
\sphinxstyleemphasis{bias} keyword set to a \sphinxstyleemphasis{yes} value.


\bigskip\hrule\bigskip


\sphinxAtStartPar
The \sphinxstyleemphasis{loop} keyword is used by \sphinxstyleemphasis{create} in the following ways.

\sphinxAtStartPar
If loop = all, then each processor loops over all atoms in the
simulation to create velocities, but only stores velocities for atoms
it owns.  This can be a slow loop for a large simulation.  If atoms
were read from a data file, the velocity assigned to a particular atom
will be the same, independent of how many processors are being used.
This will not be the case if atoms were created using the
{\hyperref[\detokenize{create_atoms::doc}]{\sphinxcrossref{\DUrole{doc}{create\_atoms}}}} command, since atom IDs will likely
be assigned to atoms differently.

\sphinxAtStartPar
If loop = local, then each processor loops over only its atoms to
produce velocities.  The random number seed is adjusted to give a
different set of velocities on each processor.  This is a fast loop,
but the velocity assigned to a particular atom will depend on which
processor owns it.  Thus the results will always be different when a
simulation is run on a different number of processors.

\sphinxAtStartPar
If loop = geom, then each processor loops over only its atoms.  For
each atom a unique random number seed is created, based on the atom’s
xyz coordinates.  A velocity is generated using that seed.  This is a
fast loop and the velocity assigned to a particular atom will be the
same, independent of how many processors are used.  However, the set
of generated velocities may be more correlated than if the \sphinxstyleemphasis{all} or
\sphinxstyleemphasis{local} keywords are used.

\sphinxAtStartPar
Note that the \sphinxstyleemphasis{loop geom} keyword will not necessarily assign
identical velocities for two simulations run on different machines.
This is because the computations based on xyz coordinates are
sensitive to tiny differences in the double\sphinxhyphen{}precision value for a
coordinate as stored on a particular machine.


\bigskip\hrule\bigskip


\sphinxAtStartPar
The \sphinxstyleemphasis{rigid} keyword only has meaning when used with the \sphinxstyleemphasis{zero} style.
It allows specification of a fix\sphinxhyphen{}ID for one of the {\hyperref[\detokenize{fix_rigid::doc}]{\sphinxcrossref{\DUrole{doc}{rigid\sphinxhyphen{}body
fix}}}} variants which defines a set of rigid bodies.  The
zeroing of linear or angular momentum is then performed for each rigid
body defined by the fix, as described above.

\sphinxAtStartPar
The \sphinxstyleemphasis{units} keyword is used by \sphinxstyleemphasis{set} and \sphinxstyleemphasis{ramp}.  If units = box,
the velocities and coordinates specified in the velocity command are
in the standard units described by the {\hyperref[\detokenize{units::doc}]{\sphinxcrossref{\DUrole{doc}{units}}}} command
(e.g. Angstroms/fs for real units).  If units = lattice, velocities
are in units of lattice spacings per time (e.g. spacings/fs) and
coordinates are in lattice spacings.  The {\hyperref[\detokenize{lattice::doc}]{\sphinxcrossref{\DUrole{doc}{lattice}}}}
command must have been previously used to define the lattice spacing.


\bigskip\hrule\bigskip



\subsection{Restrictions}
\label{\detokenize{velocity:restrictions}}
\sphinxAtStartPar
Assigning a temperature via the \sphinxstyleemphasis{create} style to a system with
{\hyperref[\detokenize{fix_rigid::doc}]{\sphinxcrossref{\DUrole{doc}{rigid bodies}}}} or {\hyperref[\detokenize{fix_shake::doc}]{\sphinxcrossref{\DUrole{doc}{SHAKE constraints}}}} may not have the desired outcome for two reasons.  First,
the velocity command can be invoked before all of the relevant fixes
are created and initialized and the number of adjusted degrees of
freedom (DOFs) is known.  Thus it is not possible to compute the
target temperature correctly.  Second, the assigned velocities may be
partially canceled when constraints are first enforced, leading to a
different temperature than desired.  A workaround for this is to
perform a {\hyperref[\detokenize{run::doc}]{\sphinxcrossref{\DUrole{doc}{run 0}}}} command, which ensures all DOFs are
accounted for properly, and then rescale the temperature to the
desired value before performing a simulation.  For example:

\begin{sphinxVerbatim}[commandchars=\\\{\}]
\PYG{k}{velocity }\PYG{n+nv+nvIdentifier}{all}\PYG{+w}{ }\PYG{n}{create}\PYG{+w}{ }\PYG{l+m}{300.0}\PYG{+w}{ }\PYG{l+m}{12345}
\PYG{k}{run}\PYG{+w}{ }\PYG{l+m}{0}\PYG{+w}{                             }\PYG{c}{\PYGZsh{} temperature may not be 300K}
\PYG{k}{velocity }\PYG{n+nv+nvIdentifier}{all}\PYG{+w}{ }\PYG{n}{scale}\PYG{+w}{ }\PYG{l+m}{300.0}\PYG{+w}{          }\PYG{c}{\PYGZsh{} now it should be}
\end{sphinxVerbatim}


\subsection{Related commands}
\label{\detokenize{velocity:related-commands}}
\sphinxAtStartPar
{\hyperref[\detokenize{fix_rigid::doc}]{\sphinxcrossref{\DUrole{doc}{fix rigid}}}}, {\hyperref[\detokenize{fix_shake::doc}]{\sphinxcrossref{\DUrole{doc}{fix shake}}}},
{\hyperref[\detokenize{lattice::doc}]{\sphinxcrossref{\DUrole{doc}{lattice}}}}


\subsection{Default}
\label{\detokenize{velocity:default}}
\sphinxAtStartPar
The keyword defaults are dist = uniform, sum = no, mom = yes, rot =
no, bias = no, loop = all, and units = lattice.  The temp and rigid
keywords are not defined by default.

\sphinxstepscope

\index{write\_coeff@\spxentry{write\_coeff}}\ignorespaces 

\section{write\_coeff command}
\label{\detokenize{write_coeff:write-coeff-command}}\label{\detokenize{write_coeff:index-0}}\label{\detokenize{write_coeff::doc}}

\subsection{Syntax}
\label{\detokenize{write_coeff:syntax}}
\begin{sphinxVerbatim}[commandchars=\\\{\}]
\PYG{k}{write\PYGZus{}coeff }\PYG{n+nv+nvIdentifier}{file}
\end{sphinxVerbatim}
\begin{itemize}
\item {} 
\sphinxAtStartPar
file = name of data file to write out

\end{itemize}


\subsection{Examples}
\label{\detokenize{write_coeff:examples}}
\begin{sphinxVerbatim}[commandchars=\\\{\}]
\PYG{k}{write\PYGZus{}coeff }\PYG{n+nv+nvIdentifier}{polymer.coeff}
\end{sphinxVerbatim}


\subsection{Description}
\label{\detokenize{write_coeff:description}}
\sphinxAtStartPar
Write a text format file with the currently defined force field
coefficients in a way, that it can be read by LAMMPS with the
{\hyperref[\detokenize{include::doc}]{\sphinxcrossref{\DUrole{doc}{include}}}} command. In combination with the nocoeff
option of {\hyperref[\detokenize{write_data::doc}]{\sphinxcrossref{\DUrole{doc}{write\_data}}}} this can be used to move
the Coeffs sections from a data file into a separate file.

\begin{sphinxadmonition}{note}{Note:}
\sphinxAtStartPar
The write\_coeff command is not yet fully implemented as
some pair styles do not output their coefficient information.
This means you will need to add/copy this information manually.
\end{sphinxadmonition}


\bigskip\hrule\bigskip



\subsection{Restrictions}
\label{\detokenize{write_coeff:restrictions}}
\sphinxAtStartPar
none


\subsection{Related commands}
\label{\detokenize{write_coeff:related-commands}}
\sphinxAtStartPar
{\hyperref[\detokenize{read_data::doc}]{\sphinxcrossref{\DUrole{doc}{read\_data}}}}, {\hyperref[\detokenize{write_restart::doc}]{\sphinxcrossref{\DUrole{doc}{write\_restart}}}},
{\hyperref[\detokenize{write_data::doc}]{\sphinxcrossref{\DUrole{doc}{write\_data}}}}

\sphinxstepscope

\index{write\_data@\spxentry{write\_data}}\ignorespaces 

\section{write\_data command}
\label{\detokenize{write_data:write-data-command}}\label{\detokenize{write_data:index-0}}\label{\detokenize{write_data::doc}}

\subsection{Syntax}
\label{\detokenize{write_data:syntax}}
\begin{sphinxVerbatim}[commandchars=\\\{\}]
\PYG{k}{write\PYGZus{}data }\PYG{n+nv+nvIdentifier}{file}\PYG{+w}{ }\PYG{n}{keyword}\PYG{+w}{ }\PYG{n}{value}\PYG{+w}{ }\PYG{n}{...}
\end{sphinxVerbatim}
\begin{itemize}
\item {} 
\sphinxAtStartPar
file = name of data file to write out

\item {} 
\sphinxAtStartPar
zero or more keyword/value pairs may be appended

\item {} 
\sphinxAtStartPar
keyword = \sphinxstyleemphasis{nocoeff} or \sphinxstyleemphasis{nofix} or \sphinxstyleemphasis{nolabelmap} or \sphinxstyleemphasis{triclinic/general} or \sphinxstyleemphasis{types} or \sphinxstyleemphasis{pair}
\begin{sphinxalltt}
\sphinxstyleemphasis{nocoeff} = do not write out force field info
\sphinxstyleemphasis{nofix} = do not write out extra sections read by fixes
\sphinxstyleemphasis{nolabelmap} = do not write out type labels
\sphinxstyleemphasis{triclinic/general} = write data file in general triclinic format
\sphinxstyleemphasis{types} value = \sphinxstyleemphasis{numeric} or \sphinxstyleemphasis{labels}
\sphinxstyleemphasis{pair} value = \sphinxstyleemphasis{ii} or \sphinxstyleemphasis{ij}
  \sphinxstyleemphasis{ii} = write one line of pair coefficient info per atom type
  \sphinxstyleemphasis{ij} = write one line of pair coefficient info per IJ atom type pair
\end{sphinxalltt}

\end{itemize}


\subsection{Examples}
\label{\detokenize{write_data:examples}}
\begin{sphinxVerbatim}[commandchars=\\\{\}]
\PYG{k}{write\PYGZus{}data }\PYG{n+nv+nvIdentifier}{data.polymer}
\PYG{k}{write\PYGZus{}data }\PYG{n+nv+nvIdentifier}{data.}\PYG{o}{*}
\PYG{k}{write\PYGZus{}data }\PYG{n+nv+nvIdentifier}{data.solid}\PYG{+w}{ }\PYG{n}{triclinic}\PYG{o}{/}\PYG{n}{general}
\end{sphinxVerbatim}


\subsection{Description}
\label{\detokenize{write_data:description}}
\sphinxAtStartPar
Write a data file in text format of the current state of the
simulation.  Data files can be read by the {\hyperref[\detokenize{read_data::doc}]{\sphinxcrossref{\DUrole{doc}{read data}}}}
command to begin a simulation.  The {\hyperref[\detokenize{read_data::doc}]{\sphinxcrossref{\DUrole{doc}{read\_data}}}} command
also describes their format.

\sphinxAtStartPar
Similar to {\hyperref[\detokenize{dump::doc}]{\sphinxcrossref{\DUrole{doc}{dump}}}} files, the data filename can contain a “*”
wild\sphinxhyphen{}card character.  The “*” is replaced with the current timestep
value.

\begin{sphinxadmonition}{note}{Data in Coeff sections}

\sphinxAtStartPar
The write\_data command may not always write all coefficient settings
to the corresponding Coeff sections of the data file.  This can have
one of multiple reasons.  1) The style may be a hybrid style. In that
case \sphinxstyleemphasis{no} coeff information is written.  2) A few styles may be
missing the code that would write those sections (This is rare these
days, but if you come across one, please notify the LAMMPS
developers).  3) Some pair styles require a single pair\_coeff
statement and those are not compatible with data files.  4) The
default for write\_data is to write a PairCoeff section, which has
only entries for atom types i == j. The remaining coefficients would
be inferred through the currently selected mixing rule.  If there has
been a pair\_coeff command with i != j, this setting would be lost.
LAMMPS will detect this and print a warning message unless \sphinxstyleemphasis{pair ij}
is appended to the write\_data command.  This will request writing a
PairIJCoeff section which has information for all pairs of atom types.
In cases where the coefficient data in the data file is incomplete,
you will need to re\sphinxhyphen{}specify that information in your input script
that reads the data file.
\end{sphinxadmonition}

\sphinxAtStartPar
Because a data file is in text format, if you use a data file written
out by this command to restart a simulation, the initial state of the
new run will be slightly different than the final state of the old run
(when the file was written) which was represented internally by LAMMPS
in binary format.  A new simulation which reads the data file will
thus typically diverge from a simulation that continued in the
original input script.

\sphinxAtStartPar
If you want to do more exact restarts, using binary files, see the
{\hyperref[\detokenize{restart::doc}]{\sphinxcrossref{\DUrole{doc}{restart}}}}, {\hyperref[\detokenize{write_restart::doc}]{\sphinxcrossref{\DUrole{doc}{write\_restart}}}}, and
{\hyperref[\detokenize{read_restart::doc}]{\sphinxcrossref{\DUrole{doc}{read\_restart}}}} commands.  You can also convert
binary restart files to text data files, after a simulation has run,
using the {\hyperref[\detokenize{Run_options::doc}]{\sphinxcrossref{\DUrole{doc}{\sphinxhyphen{}r command\sphinxhyphen{}line switch}}}}.

\begin{sphinxadmonition}{note}{Note:}
\sphinxAtStartPar
Only limited information about a simulation is stored in a data
file.  For example, no information about atom {\hyperref[\detokenize{group::doc}]{\sphinxcrossref{\DUrole{doc}{groups}}}} and
{\hyperref[\detokenize{fix::doc}]{\sphinxcrossref{\DUrole{doc}{fixes}}}} are stored.  {\hyperref[\detokenize{read_restart::doc}]{\sphinxcrossref{\DUrole{doc}{Binary restart files}}}}
store more information.
\end{sphinxadmonition}

\sphinxAtStartPar
Bond interactions (angle, etc) that have been turned off by the
{\hyperref[\detokenize{fix_shake::doc}]{\sphinxcrossref{\DUrole{doc}{fix shake}}}} or {\hyperref[\detokenize{delete_bonds::doc}]{\sphinxcrossref{\DUrole{doc}{delete\_bonds}}}}
command will be written to a data file as if they are turned on.  This
means they will need to be turned off again in a new run after the
data file is read.

\sphinxAtStartPar
Bonds that are broken (e.g. by a bond\sphinxhyphen{}breaking potential) are not
written to the data file.  Thus these bonds will not exist when the
data file is read.


\bigskip\hrule\bigskip


\sphinxAtStartPar
Use of the \sphinxstyleemphasis{nocoeff} keyword means no force field parameters are
written to the data file. This can be helpful, for example, if you
want to make significant changes to the force field or if the force
field parameters are read in separately, e.g. from an include file.

\sphinxAtStartPar
Use of the \sphinxstyleemphasis{nofix} keyword means no extra sections read by fixes are
written to the data file (see the \sphinxstyleemphasis{fix} option of the {\hyperref[\detokenize{read_data::doc}]{\sphinxcrossref{\DUrole{doc}{read\_data}}}} command for details). For example, this option excludes
sections for user\sphinxhyphen{}created per\sphinxhyphen{}atom properties from {\hyperref[\detokenize{fix_property_atom::doc}]{\sphinxcrossref{\DUrole{doc}{fix
property/atom}}}}.

\sphinxAtStartPar
The \sphinxstyleemphasis{nolabelmap} and \sphinxstyleemphasis{types} keywords refer to type labels that may be
defined for numeric atom types, bond types, angle types, etc.  The
label map can be defined in two ways, either by the {\hyperref[\detokenize{labelmap::doc}]{\sphinxcrossref{\DUrole{doc}{labelmap}}}} command or in data files read by the {\hyperref[\detokenize{read_data::doc}]{\sphinxcrossref{\DUrole{doc}{read\_data}}}} command which have sections for Atom Type Labels, Bond
Type Labels, Angle Type Labels, etc.  See the {\hyperref[\detokenize{Howto_type_labels::doc}]{\sphinxcrossref{\DUrole{doc}{Howto type labels}}}} doc page for the allowed syntax of type labels
and a general discussion of how type labels can be used.

\sphinxAtStartPar
Use of the \sphinxstyleemphasis{nolabelmap} keyword means that even if type labels exist
for a given type\sphinxhyphen{}kind (Atoms, Bonds, Angles, etc.), type labels are
not written to the data file.  By default, they are written if they
exist.  A type label must be defined for every numeric type (within a
given type\sphinxhyphen{}kind) to be written to the data file.

\sphinxAtStartPar
Use of the \sphinxstyleemphasis{triclinic/general} keyword will output a data file which
specifies a general triclinic simulation box as well as per\sphinxhyphen{}atom
quantities consistent with the general triclinic box.  The latter means
that per\sphinxhyphen{}atom vectors, such as velocities and dipole moments will be
oriented consistent with the 3d rotation implied by the general
triclinic box (relative to the associated restricted triclinic box).

\sphinxAtStartPar
This option can only be requested if the simulation box was initially
defined to be general triclinic.  If if was and the
\sphinxstyleemphasis{triclinic/general} keyword is not used, then the data file will
specify a restricted triclinic box, since that is the internal format
LAMMPS uses for both general and restricted triclinic simulations.
See the {\hyperref[\detokenize{Howto_triclinic::doc}]{\sphinxcrossref{\DUrole{doc}{Howto triclinic}}}} doc page for more
explanation of how general triclinic simulation boxes are supported by
LAMMPS.  And see the {\hyperref[\detokenize{read_data::doc}]{\sphinxcrossref{\DUrole{doc}{read\_data}}}} doc page for details
of how the format is altered for general triclinic data files.

\sphinxAtStartPar
The \sphinxstyleemphasis{types} keyword determines how atom types, bond types, angle
types, etc are written into these data file sections: Atoms, Bonds,
Angles, etc.  The default is the \sphinxstyleemphasis{numeric} setting, even if type label
maps exist.  If the \sphinxstyleemphasis{labels} setting is used, type labels will be
written to the data file, if the corresponding label map exists.  Note
that when using \sphinxstyleemphasis{types labels}, the \sphinxstyleemphasis{nolabelmap} keyword cannot be
used.

\sphinxAtStartPar
The \sphinxstyleemphasis{pair} keyword lets you specify in what format the pair
coefficient information is written into the data file.  If the value
is specified as \sphinxstyleemphasis{ii}, then one line per atom type is written, to
specify the coefficients for each of the I=J interactions.  This means
that no cross\sphinxhyphen{}interactions for I != J will be specified in the data
file and the pair style will apply its mixing rule, as documented on
individual {\hyperref[\detokenize{pair_style::doc}]{\sphinxcrossref{\DUrole{doc}{pair\_style}}}} doc pages.  Of course this
behavior can be overridden in the input script after reading the data
file, by specifying additional {\hyperref[\detokenize{pair_coeff::doc}]{\sphinxcrossref{\DUrole{doc}{pair\_coeff}}}} commands
for any desired I,J pairs.

\sphinxAtStartPar
If the value is specified as \sphinxstyleemphasis{ij}, then one line of coefficients is
written for all I,J pairs where I \textless{}= J.  These coefficients will
include any specific settings made in the input script up to that
point.  The presence of these I != J coefficients in the data file
will effectively turn off the default mixing rule for the pair style.
Again, the coefficient values in the data file can be overridden
in the input script after reading the data file, by specifying
additional {\hyperref[\detokenize{pair_coeff::doc}]{\sphinxcrossref{\DUrole{doc}{pair\_coeff}}}} commands for any desired I,J
pairs.


\bigskip\hrule\bigskip



\subsection{Restrictions}
\label{\detokenize{write_data:restrictions}}
\sphinxAtStartPar
This command requires inter\sphinxhyphen{}processor communication to migrate atoms
before the data file is written.  This means that your system must be
ready to perform a simulation before using this command (force fields
setup, atom masses initialized, etc).


\subsection{Related commands}
\label{\detokenize{write_data:related-commands}}
\sphinxAtStartPar
{\hyperref[\detokenize{read_data::doc}]{\sphinxcrossref{\DUrole{doc}{read\_data}}}}, {\hyperref[\detokenize{write_restart::doc}]{\sphinxcrossref{\DUrole{doc}{write\_restart}}}}


\subsection{Default}
\label{\detokenize{write_data:default}}
\sphinxAtStartPar
The option defaults are pair = ii and types = numeric.

\sphinxstepscope

\index{write\_dump@\spxentry{write\_dump}}\ignorespaces 

\section{write\_dump command}
\label{\detokenize{write_dump:write-dump-command}}\label{\detokenize{write_dump:index-0}}\label{\detokenize{write_dump::doc}}

\subsection{Syntax}
\label{\detokenize{write_dump:syntax}}
\begin{sphinxVerbatim}[commandchars=\\\{\}]
\PYG{k}{write\PYGZus{}dump }\PYG{n+nv+nvIdentifier}{group\PYGZhy{}ID}\PYG{+w}{ }\PYG{n}{style}\PYG{+w}{ }\PYG{n}{file}\PYG{+w}{ }\PYG{n}{dump}\PYG{o}{\PYGZhy{}}\PYG{n}{args}\PYG{+w}{ }\PYG{n}{modify}\PYG{+w}{ }\PYG{n}{dump\PYGZus{}modify}\PYG{o}{\PYGZhy{}}\PYG{n}{args}
\end{sphinxVerbatim}
\begin{itemize}
\item {} 
\sphinxAtStartPar
group\sphinxhyphen{}ID = ID of the group of atoms to be dumped

\item {} 
\sphinxAtStartPar
style = any of the supported {\hyperref[\detokenize{dump::doc}]{\sphinxcrossref{\DUrole{doc}{dump styles}}}}

\item {} 
\sphinxAtStartPar
file = name of file to write dump info to

\item {} 
\sphinxAtStartPar
dump\sphinxhyphen{}args = any additional args needed for a particular {\hyperref[\detokenize{dump::doc}]{\sphinxcrossref{\DUrole{doc}{dump style}}}}

\item {} 
\sphinxAtStartPar
modify = all args after this keyword are passed to {\hyperref[\detokenize{dump_modify::doc}]{\sphinxcrossref{\DUrole{doc}{dump\_modify}}}} (optional)

\item {} 
\sphinxAtStartPar
dump\sphinxhyphen{}modify\sphinxhyphen{}args = args for {\hyperref[\detokenize{dump_modify::doc}]{\sphinxcrossref{\DUrole{doc}{dump\_modify}}}} (optional)

\end{itemize}


\subsection{Examples}
\label{\detokenize{write_dump:examples}}
\begin{sphinxVerbatim}[commandchars=\\\{\}]
\PYG{k}{write\PYGZus{}dump }\PYG{n+nv+nvIdentifier}{all}\PYG{+w}{ }\PYG{n}{atom}\PYG{+w}{ }\PYG{n}{dump.atom}
\PYG{k}{write\PYGZus{}dump }\PYG{n+nv+nvIdentifier}{subgroup}\PYG{+w}{ }\PYG{n}{atom}\PYG{+w}{ }\PYG{n}{dump.run.bin}
\PYG{k}{write\PYGZus{}dump }\PYG{n+nv+nvIdentifier}{all}\PYG{+w}{ }\PYG{n}{custom}\PYG{+w}{ }\PYG{n}{dump.myforce.}\PYG{o}{*}\PYG{+w}{ }\PYG{n}{id}\PYG{+w}{ }\PYG{n}{type}\PYG{+w}{ }\PYG{n}{x}\PYG{+w}{ }\PYG{n}{y}\PYG{+w}{ }\PYG{n}{vx}\PYG{+w}{ }\PYG{n}{fx}
\PYG{k}{write\PYGZus{}dump }\PYG{n+nv+nvIdentifier}{flow}\PYG{+w}{ }\PYG{n}{custom}\PYG{+w}{ }\PYG{n}{dump.}\PYG{o}{\PYGZpc{}}\PYG{n}{.myforce}\PYG{+w}{ }\PYG{n}{id}\PYG{+w}{ }\PYG{n}{type}\PYG{+w}{ }\PYG{n}{c\PYGZus{}myF[3]}\PYG{+w}{ }\PYG{n}{v\PYGZus{}ke}\PYG{+w}{ }\PYG{n}{modify}\PYG{+w}{ }\PYG{n}{sort}\PYG{+w}{ }\PYG{n}{id}
\PYG{k}{write\PYGZus{}dump }\PYG{n+nv+nvIdentifier}{all}\PYG{+w}{ }\PYG{n}{xyz}\PYG{+w}{ }\PYG{n}{system.xyz}\PYG{+w}{ }\PYG{n}{modify}\PYG{+w}{ }\PYG{n}{sort}\PYG{+w}{ }\PYG{n}{id}\PYG{+w}{ }\PYG{n}{element}\PYG{+w}{ }\PYG{n}{O}\PYG{+w}{ }\PYG{n}{H}
\PYG{k}{write\PYGZus{}dump }\PYG{n+nv+nvIdentifier}{all}\PYG{+w}{ }\PYG{n}{image}\PYG{+w}{ }\PYG{n}{snap}\PYG{o}{*}\PYG{n}{.jpg}\PYG{+w}{ }\PYG{n}{type}\PYG{+w}{ }\PYG{n}{type}\PYG{+w}{ }\PYG{n}{size}\PYG{+w}{ }\PYG{l+m}{960}\PYG{+w}{ }\PYG{l+m}{960}\PYG{+w}{ }\PYG{n}{modify}\PYG{+w}{ }\PYG{n}{backcolor}\PYG{+w}{ }\PYG{n}{white}
\PYG{k}{write\PYGZus{}dump }\PYG{n+nv+nvIdentifier}{all}\PYG{+w}{ }\PYG{n}{image}\PYG{+w}{ }\PYG{n}{snap}\PYG{o}{*}\PYG{n}{.jpg}\PYG{+w}{ }\PYG{n}{element}\PYG{+w}{ }\PYG{n}{element}\PYG{+w}{ }\PYG{o}{\PYGZam{}}
\PYG{+w}{   }\PYG{n}{bond}\PYG{+w}{ }\PYG{n}{atom}\PYG{+w}{ }\PYG{l+m}{0.3}\PYG{+w}{ }\PYG{n}{shiny}\PYG{+w}{ }\PYG{l+m}{0.1}\PYG{+w}{ }\PYG{n}{ssao}\PYG{+w}{ }\PYG{n}{yes}\PYG{+w}{ }\PYG{l+m}{6345}\PYG{+w}{ }\PYG{l+m}{0.2}\PYG{+w}{ }\PYG{n}{size}\PYG{+w}{ }\PYG{l+m}{1600}\PYG{+w}{ }\PYG{l+m}{1600}\PYG{+w}{  }\PYG{o}{\PYGZam{}}
\PYG{+w}{   }\PYG{n}{modify}\PYG{+w}{ }\PYG{n}{backcolor}\PYG{+w}{ }\PYG{n}{white}\PYG{+w}{ }\PYG{n}{element}\PYG{+w}{ }\PYG{n}{C}\PYG{+w}{ }\PYG{n}{C}\PYG{+w}{ }\PYG{n}{O}\PYG{+w}{ }\PYG{n}{H}\PYG{+w}{ }\PYG{n}{N}\PYG{+w}{ }\PYG{n}{C}\PYG{+w}{ }\PYG{n}{C}\PYG{+w}{ }\PYG{n}{C}\PYG{+w}{ }\PYG{n}{O}\PYG{+w}{ }\PYG{n}{H}\PYG{+w}{ }\PYG{n}{H}\PYG{+w}{ }\PYG{n}{S}\PYG{+w}{ }\PYG{n}{O}\PYG{+w}{ }\PYG{n}{H}
\PYG{k}{write\PYGZus{}dump }\PYG{n+nv+nvIdentifier}{all}\PYG{+w}{ }\PYG{n}{atom}\PYG{o}{/}\PYG{n}{gz}\PYG{+w}{ }\PYG{n}{dump.atom.gz}\PYG{+w}{ }\PYG{n}{modify}\PYG{+w}{ }\PYG{n}{compression\PYGZus{}level}\PYG{+w}{ }\PYG{l+m}{9}
\PYG{k}{write\PYGZus{}dump }\PYG{n+nv+nvIdentifier}{flow}\PYG{+w}{ }\PYG{n}{custom}\PYG{o}{/}\PYG{n}{zstd}\PYG{+w}{ }\PYG{n}{dump.}\PYG{o}{\PYGZpc{}}\PYG{n}{.myforce.zst}\PYG{+w}{ }\PYG{o}{\PYGZam{}}
\PYG{+w}{                }\PYG{n}{id}\PYG{+w}{ }\PYG{n}{type}\PYG{+w}{ }\PYG{n}{c\PYGZus{}myF[3]}\PYG{+w}{ }\PYG{n}{v\PYGZus{}ke}\PYG{+w}{ }\PYG{o}{\PYGZam{}}
\PYG{+w}{                }\PYG{n}{modify}\PYG{+w}{ }\PYG{n}{sort}\PYG{+w}{ }\PYG{n}{id}\PYG{+w}{ }\PYG{o}{\PYGZam{}}
\PYG{+w}{                }\PYG{n}{compression\PYGZus{}level}\PYG{+w}{ }\PYG{l+m}{15}
\end{sphinxVerbatim}


\subsection{Description}
\label{\detokenize{write_dump:description}}
\sphinxAtStartPar
Dump a single snapshot of atom quantities to one or more files for the
current state of the system.  This is a one\sphinxhyphen{}time immediate operation,
in contrast to the {\hyperref[\detokenize{dump::doc}]{\sphinxcrossref{\DUrole{doc}{dump}}}} command which will will set up a
dump style to write out snapshots periodically during a running
simulation.

\sphinxAtStartPar
The syntax for this command is mostly identical to that of the
{\hyperref[\detokenize{dump::doc}]{\sphinxcrossref{\DUrole{doc}{dump}}}} and {\hyperref[\detokenize{dump_modify::doc}]{\sphinxcrossref{\DUrole{doc}{dump\_modify}}}} commands as if
they were concatenated together, with the following exceptions: There
is no need for a dump ID or dump frequency and the keyword \sphinxstyleemphasis{modify} is
added.  The latter is so that the full range of
{\hyperref[\detokenize{dump_modify::doc}]{\sphinxcrossref{\DUrole{doc}{dump\_modify}}}} options can be specified for the single
snapshot, just as they can be for multiple snapshots.  The \sphinxstyleemphasis{modify}
keyword separates the arguments that would normally be passed to the
\sphinxstyleemphasis{dump} command from those that would be given the \sphinxstyleemphasis{dump\_modify}.  Both
support optional arguments and thus LAMMPS needs to be able to cleanly
separate the two sets of args.

\sphinxAtStartPar
Note that if the specified filename uses wildcard characters “*” or
“\%”, as supported by the {\hyperref[\detokenize{dump::doc}]{\sphinxcrossref{\DUrole{doc}{dump}}}} command, they will operate
in the same fashion to create the new filename(s).  Normally, {\hyperref[\detokenize{dump_image::doc}]{\sphinxcrossref{\DUrole{doc}{dump image}}}} files require a filename with a “*” character
for the timestep.  That is not the case for the write\_dump command; no
wildcard “*” character is necessary.


\bigskip\hrule\bigskip



\subsection{Restrictions}
\label{\detokenize{write_dump:restrictions}}
\sphinxAtStartPar
All restrictions for the {\hyperref[\detokenize{dump::doc}]{\sphinxcrossref{\DUrole{doc}{dump}}}} and
{\hyperref[\detokenize{dump_modify::doc}]{\sphinxcrossref{\DUrole{doc}{dump\_modify}}}} commands apply to this command as well,
with the exception of the {\hyperref[\detokenize{dump_image::doc}]{\sphinxcrossref{\DUrole{doc}{dump image}}}} filename not
requiring a wildcard “*” character, as noted above.

\sphinxAtStartPar
Since dumps are normally written during a {\hyperref[\detokenize{run::doc}]{\sphinxcrossref{\DUrole{doc}{run}}}} or {\hyperref[\detokenize{minimize::doc}]{\sphinxcrossref{\DUrole{doc}{energy minimization}}}}, the simulation has to be ready to run
before this command can be used.  Similarly, if the dump requires
information from a compute, fix, or variable, the information needs to
have been calculated for the current timestep (e.g. by a prior run),
else LAMMPS will generate an error message.

\sphinxAtStartPar
For example, it is not possible to dump per\sphinxhyphen{}atom energy with this
command before a run has been performed, since no energies and forces
have yet been calculated.  See the {\hyperref[\detokenize{variable::doc}]{\sphinxcrossref{\DUrole{doc}{variable}}}} doc page
section on Variable Accuracy for more information on this topic.


\subsection{Related commands}
\label{\detokenize{write_dump:related-commands}}
\sphinxAtStartPar
{\hyperref[\detokenize{dump::doc}]{\sphinxcrossref{\DUrole{doc}{dump}}}}, {\hyperref[\detokenize{dump_image::doc}]{\sphinxcrossref{\DUrole{doc}{dump image}}}},
{\hyperref[\detokenize{dump_modify::doc}]{\sphinxcrossref{\DUrole{doc}{dump\_modify}}}}


\subsection{Default}
\label{\detokenize{write_dump:default}}
\sphinxAtStartPar
The defaults are listed on the doc pages for the {\hyperref[\detokenize{dump::doc}]{\sphinxcrossref{\DUrole{doc}{dump}}}} and
{\hyperref[\detokenize{dump_image::doc}]{\sphinxcrossref{\DUrole{doc}{dump image}}}} and {\hyperref[\detokenize{dump_modify::doc}]{\sphinxcrossref{\DUrole{doc}{dump\_modify}}}}
commands.

\sphinxstepscope

\index{write\_restart@\spxentry{write\_restart}}\ignorespaces 

\section{write\_restart command}
\label{\detokenize{write_restart:write-restart-command}}\label{\detokenize{write_restart:index-0}}\label{\detokenize{write_restart::doc}}

\subsection{Syntax}
\label{\detokenize{write_restart:syntax}}
\begin{sphinxVerbatim}[commandchars=\\\{\}]
\PYG{k}{write\PYGZus{}restart }\PYG{n+nv+nvIdentifier}{file}\PYG{+w}{ }\PYG{n}{keyword}\PYG{+w}{ }\PYG{n}{value}\PYG{+w}{ }\PYG{n}{...}
\end{sphinxVerbatim}
\begin{itemize}
\item {} 
\sphinxAtStartPar
file = name of file to write restart information to

\item {} 
\sphinxAtStartPar
zero or more keyword/value pairs may be appended

\item {} 
\sphinxAtStartPar
keyword = \sphinxstyleemphasis{fileper} or \sphinxstyleemphasis{nfile}
\begin{sphinxalltt}
\sphinxstyleemphasis{fileper} arg = Np
  Np = write one file for every this many processors
\sphinxstyleemphasis{nfile} arg = Nf
  Nf = write this many files, one from each of Nf processors
\end{sphinxalltt}

\end{itemize}


\subsection{Examples}
\label{\detokenize{write_restart:examples}}
\begin{sphinxVerbatim}[commandchars=\\\{\}]
\PYG{k}{write\PYGZus{}restart }\PYG{n+nv+nvIdentifier}{restart.equil}
\PYG{k}{write\PYGZus{}restart }\PYG{n+nv+nvIdentifier}{poly.}\PYG{o}{\PYGZpc{}}\PYG{n}{.}\PYG{o}{*}\PYG{+w}{ }\PYG{n}{nfile}\PYG{+w}{ }\PYG{l+m}{10}
\end{sphinxVerbatim}


\subsection{Description}
\label{\detokenize{write_restart:description}}
\sphinxAtStartPar
Write a binary restart file of the current state of the simulation.

\sphinxAtStartPar
During a long simulation, the {\hyperref[\detokenize{restart::doc}]{\sphinxcrossref{\DUrole{doc}{restart}}}} command is
typically used to output restart files periodically.  The
write\_restart command is useful after a minimization or whenever you
wish to write out a single current restart file.

\sphinxAtStartPar
Similar to {\hyperref[\detokenize{dump::doc}]{\sphinxcrossref{\DUrole{doc}{dump}}}} files, the restart filename can contain
two wild\sphinxhyphen{}card characters.  If a “*” appears in the filename, it is
replaced with the current timestep value.  If a “\%” character appears
in the filename, then one file is written by each processor and the
“\%” character is replaced with the processor ID from 0 to P\sphinxhyphen{}1.  An
additional file with the “\%” replaced by “base” is also written, which
contains global information.  For example, the files written for
filename restart.\% would be restart.base, restart.0, restart.1, …
restart.P\sphinxhyphen{}1.  This creates smaller files and can be a fast mode of
output and subsequent input on parallel machines that support parallel
I/O.  The optional \sphinxstyleemphasis{fileper} and \sphinxstyleemphasis{nfile} keywords discussed below can
alter the number of files written.

\sphinxAtStartPar
Restart files can be read by a {\hyperref[\detokenize{read_restart::doc}]{\sphinxcrossref{\DUrole{doc}{read\_restart}}}}
command to restart a simulation from a particular state.  Because the
file is binary (to enable exact restarts), it may not be readable on
another machine.  In this case, you can use the {\hyperref[\detokenize{Run_options::doc}]{\sphinxcrossref{\DUrole{doc}{\sphinxhyphen{}r command\sphinxhyphen{}line
switch}}}} to convert a restart file to a data file.

\begin{sphinxadmonition}{note}{Note:}
\sphinxAtStartPar
Although the purpose of restart files is to enable restarting a
simulation from where it left off, not all information about a
simulation is stored in the file.  For example, the list of fixes
that were specified during the initial run is not stored, which
means the new input script must specify any fixes you want to use.
Even when restart information is stored in the file, as it is for
some fixes, commands may need to be re\sphinxhyphen{}specified in the new input
script, in order to re\sphinxhyphen{}use that information. Details are usually
given in the documentation of the respective command. Also, see the
{\hyperref[\detokenize{read_restart::doc}]{\sphinxcrossref{\DUrole{doc}{read\_restart}}}} command for general information
about what is stored in a restart file.
\end{sphinxadmonition}


\bigskip\hrule\bigskip


\sphinxAtStartPar
The optional \sphinxstyleemphasis{nfile} or \sphinxstyleemphasis{fileper} keywords can be used in conjunction
with the “\%” wildcard character in the specified restart file name.
As explained above, the “\%” character causes the restart file to be
written in pieces, one piece for each of P processors.  By default P =
the number of processors the simulation is running on.  The \sphinxstyleemphasis{nfile} or
\sphinxstyleemphasis{fileper} keyword can be used to set P to a smaller value, which can
be more efficient when running on a large number of processors.

\sphinxAtStartPar
The \sphinxstyleemphasis{nfile} keyword sets P to the specified Nf value.  For example, if
Nf = 4, and the simulation is running on 100 processors, 4 files will
be written, by processors 0,25,50,75.  Each will collect information
from itself and the next 24 processors and write it to a restart file.

\sphinxAtStartPar
For the \sphinxstyleemphasis{fileper} keyword, the specified value of Np means write one
file for every Np processors.  For example, if Np = 4, every fourth
processor (0,4,8,12,etc) will collect information from itself and the
next 3 processors and write it to a restart file.


\bigskip\hrule\bigskip



\subsection{Restrictions}
\label{\detokenize{write_restart:restrictions}}
\sphinxAtStartPar
This command requires inter\sphinxhyphen{}processor communication to migrate atoms
before the restart file is written.  This means that your system must
be ready to perform a simulation before using this command (force
fields setup, atom masses initialized, etc).


\subsection{Related commands}
\label{\detokenize{write_restart:related-commands}}
\sphinxAtStartPar
{\hyperref[\detokenize{restart::doc}]{\sphinxcrossref{\DUrole{doc}{restart}}}}, {\hyperref[\detokenize{read_restart::doc}]{\sphinxcrossref{\DUrole{doc}{read\_restart}}}},
{\hyperref[\detokenize{write_data::doc}]{\sphinxcrossref{\DUrole{doc}{write\_data}}}}


\subsection{Default}
\label{\detokenize{write_restart:default}}
\sphinxAtStartPar
none

\sphinxstepscope


\chapter{Fixes}
\label{\detokenize{fixes:fixes}}\label{\detokenize{fixes::doc}}
\sphinxstepscope

\index{fix accelerate/cos@\spxentry{fix accelerate/cos}}\ignorespaces 

\section{fix accelerate/cos command}
\label{\detokenize{fix_accelerate_cos:fix-accelerate-cos-command}}\label{\detokenize{fix_accelerate_cos:index-0}}\label{\detokenize{fix_accelerate_cos::doc}}

\subsection{Syntax}
\label{\detokenize{fix_accelerate_cos:syntax}}
\begin{sphinxVerbatim}[commandchars=\\\{\}]
\PYG{k}{fix }\PYG{n+nv+nvIdentifier}{ID}\PYG{+w}{ }\PYG{n+nv+nvIdentifier}{group\PYGZhy{}ID}\PYG{+w}{ }\PYG{n}{accelerate}\PYG{+w}{ }\PYG{n}{value}
\end{sphinxVerbatim}
\begin{itemize}
\item {} 
\sphinxAtStartPar
ID, group\sphinxhyphen{}ID are documented in {\hyperref[\detokenize{fix::doc}]{\sphinxcrossref{\DUrole{doc}{fix}}}} command

\item {} 
\sphinxAtStartPar
accelerate/cos = style name of this fix command

\item {} 
\sphinxAtStartPar
value = amplitude of acceleration (in unit of velocity/time)

\end{itemize}


\subsection{Examples}
\label{\detokenize{fix_accelerate_cos:examples}}
\begin{sphinxVerbatim}[commandchars=\\\{\}]
\PYG{k}{fix }\PYG{n+nv+nvIdentifier}{1}\PYG{+w}{ }\PYG{n+nv+nvIdentifier}{all}\PYG{+w}{ }\PYG{n}{accelerate}\PYG{o}{/}\PYG{n}{cos}\PYG{+w}{ }\PYG{l+m}{2.0e\PYGZhy{}7}
\end{sphinxVerbatim}


\subsection{Description}
\label{\detokenize{fix_accelerate_cos:description}}
\sphinxAtStartPar
Give each atom a acceleration in x\sphinxhyphen{}direction based on its z coordinate.
The acceleration is a periodic function along the z\sphinxhyphen{}direction:
\begin{equation*}
\begin{split}a_{x}(z) = A \cos \left(\frac{2 \pi z}{l_{z}}\right)\end{split}
\end{equation*}
\sphinxAtStartPar
where \(A\) is the acceleration amplitude, \(l_z\) is the
\(z\)\sphinxhyphen{}length of the simulation box.
At steady state, the acceleration generates a velocity profile:
\begin{equation*}
\begin{split}v_{x}(z) = V \cos \left(\frac{2 \pi z}{l_{z}}\right)\end{split}
\end{equation*}
\sphinxAtStartPar
The generated velocity amplitude \(V\) is related to the
shear viscosity \(\eta\) by:
\begin{equation*}
\begin{split}V = \frac{A \rho}{\eta}\left(\frac{l_{z}}{2 \pi}\right)^{2}\end{split}
\end{equation*}
\sphinxAtStartPar
and it can be obtained from ensemble average of the velocity profile:
\begin{equation*}
\begin{split}V = \frac{\sum\limits_i 2 m_{i} v_{i, x} \cos \left(\frac{2 \pi z_i}{l_{z}}\right)}{\sum\limits_i m_{i}},\end{split}
\end{equation*}
\sphinxAtStartPar
where \(m_i\), \(v_{i,x}\), and \(z_i\) are the mass,
\(x\)\sphinxhyphen{}component velocity, and \(z\)\sphinxhyphen{}coordinate of a particle,
respectively.

\sphinxAtStartPar
The velocity amplitude \(V\) can be calculated with {\hyperref[\detokenize{compute_viscosity_cos::doc}]{\sphinxcrossref{\DUrole{doc}{compute
viscosity/cos}}}}, which enables viscosity
calculation with periodic perturbation method, as described by
{\hyperref[\detokenize{fix_accelerate_cos:hess2}]{\sphinxcrossref{\DUrole{std,std-ref}{Hess}}}}.  Because the applied acceleration drives the system
away from equilibration, the calculated shear viscosity is lower than
the intrinsic viscosity due to the shear\sphinxhyphen{}thinning effect.  Extrapolation
to zero acceleration should generally be performed to predict the
zero\sphinxhyphen{}shear viscosity.  As the shear stress decreases, the
signal\sphinxhyphen{}to\sphinxhyphen{}noise ratio decreases rapidly, and the simulation time must be
extended accordingly to get converged results.

\sphinxAtStartPar
In order to get meaningful results, the group ID of this fix should be all.


\bigskip\hrule\bigskip



\subsection{Restart, fix\_modify, output, run start/stop, minimize info}
\label{\detokenize{fix_accelerate_cos:restart-fix-modify-output-run-start-stop-minimize-info}}
\sphinxAtStartPar
No information about this fix is written to binary restart files.
None of the fix\_modify options are relevant to this fix.
No global or per\sphinxhyphen{}atom quantities are stored by this fix for access by various
output commands.  No parameter of this fix can be used with the start/stop
keywords of the run command.

\sphinxAtStartPar
This fix is not invoked during energy minimization.


\subsection{Restrictions}
\label{\detokenize{fix_accelerate_cos:restrictions}}
\sphinxAtStartPar
This fix is part of the MISC package.  It is only enabled if LAMMPS was
built with that package.  See the {\hyperref[\detokenize{Build_package::doc}]{\sphinxcrossref{\DUrole{doc}{Build package}}}}
page for more info.

\sphinxAtStartPar
Since this fix depends on the \(z\)\sphinxhyphen{}coordinate of atoms, it cannot be used
in 2d simulations.


\subsection{Related commands}
\label{\detokenize{fix_accelerate_cos:related-commands}}
\sphinxAtStartPar
{\hyperref[\detokenize{compute_viscosity_cos::doc}]{\sphinxcrossref{\DUrole{doc}{compute viscosity/cos}}}}


\subsection{Default}
\label{\detokenize{fix_accelerate_cos:default}}
\sphinxAtStartPar
none


\bigskip\hrule\bigskip

\phantomsection\label{\detokenize{fix_accelerate_cos:hess2}}
\sphinxAtStartPar
\sphinxstylestrong{(Hess)} Hess, B. Journal of Chemical Physics 2002, 116 (1), 209\textendash{}217.

\sphinxstepscope

\index{fix acks2/reaxff@\spxentry{fix acks2/reaxff}}\ignorespaces 
\index{fix acks2/reaxff/kk@\spxentry{fix acks2/reaxff/kk}}\ignorespaces 

\section{fix acks2/reaxff command}
\label{\detokenize{fix_acks2_reaxff:fix-acks2-reaxff-command}}\label{\detokenize{fix_acks2_reaxff:index-1}}\label{\detokenize{fix_acks2_reaxff:index-0}}\label{\detokenize{fix_acks2_reaxff::doc}}
\sphinxAtStartPar
Accelerator Variants: \sphinxstyleemphasis{acks2/reaxff/kk}


\subsection{Syntax}
\label{\detokenize{fix_acks2_reaxff:syntax}}
\begin{sphinxVerbatim}[commandchars=\\\{\}]
\PYG{k}{fix }\PYG{n+nv+nvIdentifier}{ID}\PYG{+w}{ }\PYG{n+nv+nvIdentifier}{group\PYGZhy{}ID}\PYG{+w}{ }\PYG{n}{acks2}\PYG{o}{/}\PYG{n}{reaxff}\PYG{+w}{ }\PYG{n}{Nevery}\PYG{+w}{ }\PYG{n}{cutlo}\PYG{+w}{ }\PYG{n}{cuthi}\PYG{+w}{ }\PYG{n}{tolerance}\PYG{+w}{ }\PYG{n}{params}\PYG{+w}{ }\PYG{n}{args}
\end{sphinxVerbatim}
\begin{itemize}
\item {} 
\sphinxAtStartPar
ID, group\sphinxhyphen{}ID are documented in {\hyperref[\detokenize{fix::doc}]{\sphinxcrossref{\DUrole{doc}{fix}}}} command

\item {} 
\sphinxAtStartPar
acks2/reaxff = style name of this fix command

\item {} 
\sphinxAtStartPar
Nevery = perform ACKS2 every this many steps

\item {} 
\sphinxAtStartPar
cutlo,cuthi = lo and hi cutoff for Taper radius

\item {} 
\sphinxAtStartPar
tolerance = precision to which charges will be equilibrated

\item {} 
\sphinxAtStartPar
params = reaxff or a filename

\item {} 
\sphinxAtStartPar
one or more keywords or keyword/value pairs may be appended
\begin{sphinxalltt}
keyword = \sphinxstyleemphasis{maxiter}
  \sphinxstyleemphasis{maxiter} N = limit the number of iterations to \sphinxstyleemphasis{N}
\end{sphinxalltt}

\end{itemize}


\subsection{Examples}
\label{\detokenize{fix_acks2_reaxff:examples}}
\begin{sphinxVerbatim}[commandchars=\\\{\}]
\PYG{k}{fix }\PYG{n+nv+nvIdentifier}{1}\PYG{+w}{ }\PYG{n+nv+nvIdentifier}{all}\PYG{+w}{ }\PYG{n}{acks2}\PYG{o}{/}\PYG{n}{reaxff}\PYG{+w}{ }\PYG{l+m}{1}\PYG{+w}{ }\PYG{l+m}{0.0}\PYG{+w}{ }\PYG{l+m}{10.0}\PYG{+w}{ }\PYG{l+m}{1.0e\PYGZhy{}6}\PYG{+w}{ }\PYG{n}{reaxff}
\PYG{k}{fix }\PYG{n+nv+nvIdentifier}{1}\PYG{+w}{ }\PYG{n+nv+nvIdentifier}{all}\PYG{+w}{ }\PYG{n}{acks2}\PYG{o}{/}\PYG{n}{reaxff}\PYG{+w}{ }\PYG{l+m}{1}\PYG{+w}{ }\PYG{l+m}{0.0}\PYG{+w}{ }\PYG{l+m}{10.0}\PYG{+w}{ }\PYG{l+m}{1.0e\PYGZhy{}6}\PYG{+w}{ }\PYG{n}{param.acks2}\PYG{+w}{ }\PYG{n}{maxiter}\PYG{+w}{ }\PYG{l+m}{500}
\end{sphinxVerbatim}


\subsection{Description}
\label{\detokenize{fix_acks2_reaxff:description}}
\sphinxAtStartPar
Perform the atom\sphinxhyphen{}condensed Kohn\textendash{}Sham DFT to second order (ACKS2) charge
equilibration method as described in {\hyperref[\detokenize{fix_acks2_reaxff:verstraelen}]{\sphinxcrossref{\DUrole{std,std-ref}{(Verstraelen)}}}}.
ACKS2 impedes unphysical long\sphinxhyphen{}range charge transfer sometimes seen with
QEq (e.g., for dissociation of molecules), at increased computational
cost.  It is typically used in conjunction with the ReaxFF force field
model as implemented in the {\hyperref[\detokenize{pair_reaxff::doc}]{\sphinxcrossref{\DUrole{doc}{pair\_style reaxff}}}}
command, but it can be used with any potential in LAMMPS, so long as it
defines and uses charges on each atom.  For more technical details about
the charge equilibration performed by fix acks2/reaxff, see the
{\hyperref[\detokenize{fix_acks2_reaxff:o-hearn}]{\sphinxcrossref{\DUrole{std,std-ref}{(O’Hearn)}}}} paper.

\sphinxAtStartPar
The ACKS2 method minimizes the electrostatic energy of the system by
adjusting the partial charge on individual atoms based on interactions
with their neighbors.  It requires some parameters for each atom type.
If the \sphinxstyleemphasis{params} setting above is the word “reaxff”, then these are
extracted from the {\hyperref[\detokenize{pair_reaxff::doc}]{\sphinxcrossref{\DUrole{doc}{pair\_style reaxff}}}} command and
the ReaxFF force field file it reads in.  If a file name is specified
for \sphinxstyleemphasis{params}, then the parameters are taken from the specified file
and the file must contain one line for each atom type.  The latter
form must be used when performing QeQ with a non\sphinxhyphen{}ReaxFF potential.
The lines should be formatted as follows:

\begin{sphinxVerbatim}[commandchars=\\\{\}]
bond\PYGZus{}softness
itype chi eta gamma bcut
\end{sphinxVerbatim}

\sphinxAtStartPar
where the first line is the global parameter \sphinxstyleemphasis{bond\_softness}. The
remaining 1 to Ntypes lines include \sphinxstyleemphasis{itype}, the atom type from 1 to
Ntypes, \sphinxstyleemphasis{chi}, the electronegativity in eV, \sphinxstyleemphasis{eta}, the self\sphinxhyphen{}Coulomb
potential in eV, \sphinxstyleemphasis{gamma}, the valence orbital exponent, and \sphinxstyleemphasis{bcut}, the
bond cutoff distance.  Note that these 4 quantities are also in the
ReaxFF potential file, except that eta is defined here as twice the eta
value in the ReaxFF file. Note that unlike the rest of LAMMPS, the units
of this fix are hard\sphinxhyphen{}coded to be \(\AA\), eV, and
electronic charge.

\sphinxAtStartPar
The optional \sphinxstyleemphasis{maxiter} keyword allows changing the max number
of iterations in the linear solver. The default value is 200.

\begin{sphinxadmonition}{note}{Note:}
\sphinxAtStartPar
In order to solve the self\sphinxhyphen{}consistent equations for electronegativity
equalization, LAMMPS imposes the additional constraint that all the
charges in the fix group must add up to zero.  The initial charge
assignments should also satisfy this constraint.  LAMMPS will print a
warning if that is not the case.
\end{sphinxadmonition}


\subsection{Restart, fix\_modify, output, run start/stop, minimize info}
\label{\detokenize{fix_acks2_reaxff:restart-fix-modify-output-run-start-stop-minimize-info}}
\sphinxAtStartPar
No information about this fix is written to {\hyperref[\detokenize{restart::doc}]{\sphinxcrossref{\DUrole{doc}{binary restart files}}}}.  This fix computes a global scalar (the number of
iterations) for access by various {\hyperref[\detokenize{Howto_output::doc}]{\sphinxcrossref{\DUrole{doc}{output commands}}}}.
No parameter of this fix can be used with the \sphinxstyleemphasis{start/stop} keywords of
the {\hyperref[\detokenize{run::doc}]{\sphinxcrossref{\DUrole{doc}{run}}}} command.

\sphinxAtStartPar
This fix is invoked during {\hyperref[\detokenize{minimize::doc}]{\sphinxcrossref{\DUrole{doc}{energy minimization}}}}.


\bigskip\hrule\bigskip


\sphinxAtStartPar
Styles with a \sphinxstyleemphasis{gpu}, \sphinxstyleemphasis{intel}, \sphinxstyleemphasis{kk}, \sphinxstyleemphasis{omp}, or \sphinxstyleemphasis{opt} suffix are
functionally the same as the corresponding style without the suffix.
They have been optimized to run faster, depending on your available
hardware, as discussed on the {\hyperref[\detokenize{Speed_packages::doc}]{\sphinxcrossref{\DUrole{doc}{Accelerator packages}}}}
page.  The accelerated styles take the same arguments and should
produce the same results, except for round\sphinxhyphen{}off and precision issues.

\sphinxAtStartPar
These accelerated styles are part of the GPU, INTEL, KOKKOS,
OPENMP, and OPT packages, respectively.  They are only enabled if
LAMMPS was built with those packages.  See the {\hyperref[\detokenize{Build_package::doc}]{\sphinxcrossref{\DUrole{doc}{Build package}}}} page for more info.

\sphinxAtStartPar
You can specify the accelerated styles explicitly in your input script
by including their suffix, or you can use the {\hyperref[\detokenize{Run_options::doc}]{\sphinxcrossref{\DUrole{doc}{\sphinxhyphen{}suffix command\sphinxhyphen{}line switch}}}} when you invoke LAMMPS, or you can use the
{\hyperref[\detokenize{suffix::doc}]{\sphinxcrossref{\DUrole{doc}{suffix}}}} command in your input script.

\sphinxAtStartPar
See the {\hyperref[\detokenize{Speed_packages::doc}]{\sphinxcrossref{\DUrole{doc}{Accelerator packages}}}} page for more
instructions on how to use the accelerated styles effectively.


\bigskip\hrule\bigskip



\subsection{Restrictions}
\label{\detokenize{fix_acks2_reaxff:restrictions}}
\sphinxAtStartPar
This fix is part of the REAXFF package.  It is only enabled if
LAMMPS was built with that package. See the {\hyperref[\detokenize{Build_package::doc}]{\sphinxcrossref{\DUrole{doc}{Build package}}}} page for more info.

\sphinxAtStartPar
This fix does not correctly handle interactions involving multiple
periodic images of the same atom.  Hence, it should not be used for
periodic cell dimensions less than \(10~\AA\).

\sphinxAtStartPar
This fix may be used in combination with {\hyperref[\detokenize{fix_efield::doc}]{\sphinxcrossref{\DUrole{doc}{fix efield}}}}
and will apply the external electric field during charge equilibration,
but there may be only one fix efield instance used, it may only use a
constant electric field, and the electric field vector may only have
components in non\sphinxhyphen{}periodic directions.


\subsection{Related commands}
\label{\detokenize{fix_acks2_reaxff:related-commands}}
\sphinxAtStartPar
{\hyperref[\detokenize{pair_reaxff::doc}]{\sphinxcrossref{\DUrole{doc}{pair\_style reaxff}}}}, {\hyperref[\detokenize{fix_qeq_reaxff::doc}]{\sphinxcrossref{\DUrole{doc}{fix qeq/reaxff}}}}


\subsection{Default}
\label{\detokenize{fix_acks2_reaxff:default}}
\sphinxAtStartPar
maxiter 200


\bigskip\hrule\bigskip

\phantomsection\label{\detokenize{fix_acks2_reaxff:o-hearn}}
\sphinxAtStartPar
\sphinxstylestrong{(O’Hearn)} O’Hearn, Alperen, Aktulga, SIAM J. Sci. Comput., 42(1), C1\textendash{}C22 (2020).

\phantomsection\label{\detokenize{fix_acks2_reaxff:verstraelen}}
\sphinxAtStartPar
\sphinxstylestrong{(Verstraelen)} Verstraelen, Ayers, Speybroeck, Waroquier, J. Chem. Phys. 138, 074108 (2013).

\sphinxstepscope

\index{fix adapt@\spxentry{fix adapt}}\ignorespaces 

\section{fix adapt command}
\label{\detokenize{fix_adapt:fix-adapt-command}}\label{\detokenize{fix_adapt:index-0}}\label{\detokenize{fix_adapt::doc}}

\subsection{Syntax}
\label{\detokenize{fix_adapt:syntax}}
\begin{sphinxVerbatim}[commandchars=\\\{\}]
\PYG{k}{fix }\PYG{n+nv+nvIdentifier}{ID}\PYG{+w}{ }\PYG{n+nv+nvIdentifier}{group\PYGZhy{}ID}\PYG{+w}{ }\PYG{n}{adapt}\PYG{+w}{ }\PYG{n}{N}\PYG{+w}{ }\PYG{n}{attribute}\PYG{+w}{ }\PYG{n}{args}\PYG{+w}{ }\PYG{n}{...}\PYG{+w}{ }\PYG{n}{keyword}\PYG{+w}{ }\PYG{n}{value}\PYG{+w}{ }\PYG{n}{...}
\end{sphinxVerbatim}
\begin{itemize}
\item {} 
\sphinxAtStartPar
ID, group\sphinxhyphen{}ID are documented in {\hyperref[\detokenize{fix::doc}]{\sphinxcrossref{\DUrole{doc}{fix}}}} command

\item {} 
\sphinxAtStartPar
adapt = style name of this fix command

\item {} 
\sphinxAtStartPar
N = adapt simulation settings every this many timesteps

\item {} 
\sphinxAtStartPar
one or more attribute/arg pairs may be appended

\item {} 
\sphinxAtStartPar
attribute = \sphinxstyleemphasis{pair} or \sphinxstyleemphasis{bond} or \sphinxstyleemphasis{angle} or \sphinxstyleemphasis{kspace} or \sphinxstyleemphasis{atom}
\begin{sphinxalltt}
\sphinxstyleemphasis{pair} args = pstyle pparam I J v\_name
  pstyle = pair style name (e.g., lj/cut)
  pparam = parameter to adapt over time
  I,J = type pair(s) to set parameter for (integer or type label)
  v\_name = variable with name that calculates value of pparam
\sphinxstyleemphasis{bond} args = bstyle bparam I v\_name
  bstyle = bond style name (e.g., harmonic)
  bparam = parameter to adapt over time
  I = type bond to set parameter for (integer or type label)
  v\_name = variable with name that calculates value of bparam
\sphinxstyleemphasis{angle} args = astyle aparam I v\_name
  astyle = angle style name (e.g., harmonic)
  aparam = parameter to adapt over time
  I = type angle to set parameter for (integer or type label)
  v\_name = variable with name that calculates value of aparam
\sphinxstyleemphasis{kspace} arg = v\_name
  v\_name = variable with name that calculates scale factor on \(k\)\sphinxhyphen{}space terms
\sphinxstyleemphasis{atom} args = atomparam v\_name
  atomparam = \sphinxstyleemphasis{charge} or \sphinxstyleemphasis{diameter} or \sphinxstyleemphasis{diameter/disc} = parameter to adapt over time
  v\_name = variable with name that calculates value of atomparam
\end{sphinxalltt}

\item {} 
\sphinxAtStartPar
zero or more keyword/value pairs may be appended

\item {} 
\sphinxAtStartPar
keyword = \sphinxstyleemphasis{scale} or \sphinxstyleemphasis{reset} or \sphinxstyleemphasis{mass}
\begin{sphinxalltt}
\sphinxstyleemphasis{scale} value = \sphinxstyleemphasis{no} or \sphinxstyleemphasis{yes}
  \sphinxstyleemphasis{no} = the variable value is the new setting
  \sphinxstyleemphasis{yes} = the variable value multiplies the original setting
\sphinxstyleemphasis{reset} value = \sphinxstyleemphasis{no} or \sphinxstyleemphasis{yes}
  \sphinxstyleemphasis{no} = values will remain altered at the end of a run
  \sphinxstyleemphasis{yes} = reset altered values to their original values at the end of a run
\sphinxstyleemphasis{mass} value = \sphinxstyleemphasis{no} or \sphinxstyleemphasis{yes}
  \sphinxstyleemphasis{no} = mass is not altered by changes in diameter
  \sphinxstyleemphasis{yes} = mass is altered by changes in diameter
\end{sphinxalltt}

\end{itemize}


\subsection{Examples}
\label{\detokenize{fix_adapt:examples}}
\begin{sphinxVerbatim}[commandchars=\\\{\}]
\PYG{k}{fix }\PYG{n+nv+nvIdentifier}{1}\PYG{+w}{ }\PYG{n+nv+nvIdentifier}{all}\PYG{+w}{ }\PYG{n}{adapt}\PYG{+w}{ }\PYG{l+m}{1}\PYG{+w}{ }\PYG{n}{pair}\PYG{+w}{ }\PYG{n}{soft}\PYG{+w}{ }\PYG{n}{a}\PYG{+w}{ }\PYG{l+m}{1}\PYG{+w}{ }\PYG{l+m}{1}\PYG{+w}{ }\PYG{n}{v\PYGZus{}prefactor}
\PYG{k}{fix }\PYG{n+nv+nvIdentifier}{1}\PYG{+w}{ }\PYG{n+nv+nvIdentifier}{all}\PYG{+w}{ }\PYG{n}{adapt}\PYG{+w}{ }\PYG{l+m}{1}\PYG{+w}{ }\PYG{n}{pair}\PYG{+w}{ }\PYG{n}{soft}\PYG{+w}{ }\PYG{n}{a}\PYG{+w}{ }\PYG{l+m}{2}\PYG{o}{*}\PYG{+w}{ }\PYG{l+m}{3}\PYG{+w}{ }\PYG{n}{v\PYGZus{}prefactor}
\PYG{k}{fix }\PYG{n+nv+nvIdentifier}{1}\PYG{+w}{ }\PYG{n+nv+nvIdentifier}{all}\PYG{+w}{ }\PYG{n}{adapt}\PYG{+w}{ }\PYG{l+m}{1}\PYG{+w}{ }\PYG{n}{pair}\PYG{+w}{ }\PYG{n}{lj}\PYG{o}{/}\PYG{n}{cut}\PYG{+w}{ }\PYG{n}{epsilon}\PYG{+w}{ }\PYG{o}{*}\PYG{+w}{ }\PYG{o}{*}\PYG{+w}{ }\PYG{n}{v\PYGZus{}scale1}\PYG{+w}{ }\PYG{n}{pair}\PYG{+w}{ }\PYG{n}{coul}\PYG{o}{/}\PYG{n}{cut}\PYG{+w}{ }\PYG{n}{scale}\PYG{+w}{ }\PYG{l+m}{3}\PYG{+w}{ }\PYG{l+m}{3}\PYG{+w}{ }\PYG{n}{v\PYGZus{}scale2}\PYG{+w}{ }\PYG{n}{scale}\PYG{+w}{ }\PYG{n}{yes}\PYG{+w}{ }\PYG{n}{reset}\PYG{+w}{ }\PYG{n}{yes}
\PYG{k}{fix }\PYG{n+nv+nvIdentifier}{1}\PYG{+w}{ }\PYG{n+nv+nvIdentifier}{all}\PYG{+w}{ }\PYG{n}{adapt}\PYG{+w}{ }\PYG{l+m}{10}\PYG{+w}{ }\PYG{n}{atom}\PYG{+w}{ }\PYG{n}{diameter}\PYG{+w}{ }\PYG{n}{v\PYGZus{}size}

\PYG{k}{variable }\PYG{n+nv+nvIdentifier}{ramp\PYGZus{}up}\PYG{+w}{ }\PYG{n}{equal}\PYG{+w}{ }\PYG{l+s}{\PYGZdq{}}\PYG{l+s}{ramp(0.01,0.5)}\PYG{l+s}{\PYGZdq{}}
\PYG{k}{fix }\PYG{n+nv+nvIdentifier}{stretch}\PYG{+w}{ }\PYG{n+nv+nvIdentifier}{all}\PYG{+w}{ }\PYG{n}{adapt}\PYG{+w}{ }\PYG{l+m}{1}\PYG{+w}{ }\PYG{n}{bond}\PYG{+w}{ }\PYG{n}{harmonic}\PYG{+w}{ }\PYG{n}{r0}\PYG{+w}{ }\PYG{l+m}{1}\PYG{+w}{ }\PYG{n}{v\PYGZus{}ramp\PYGZus{}up}

\PYG{k}{labelmap}\PYG{+w}{ }\PYG{n}{atom}\PYG{+w}{ }\PYG{l+m}{1}\PYG{+w}{ }\PYG{n}{c1}
\PYG{k}{fix }\PYG{n+nv+nvIdentifier}{1}\PYG{+w}{ }\PYG{n+nv+nvIdentifier}{all}\PYG{+w}{ }\PYG{n}{adapt}\PYG{+w}{ }\PYG{l+m}{1}\PYG{+w}{ }\PYG{n}{pair}\PYG{+w}{ }\PYG{n}{soft}\PYG{+w}{ }\PYG{n}{a}\PYG{+w}{ }\PYG{n}{c1}\PYG{+w}{ }\PYG{n}{c1}\PYG{+w}{ }\PYG{n}{v\PYGZus{}prefactor}
\end{sphinxVerbatim}


\subsection{Description}
\label{\detokenize{fix_adapt:description}}
\sphinxAtStartPar
Change or adapt one or more specific simulation attributes or settings over
time as a simulation runs.  Pair potential and \(k\)\sphinxhyphen{}space and atom
attributes which can be varied by this fix are discussed below.  Many other
fixes can also be used to time\sphinxhyphen{}vary simulation parameters (e.g., the
{\hyperref[\detokenize{fix_deform::doc}]{\sphinxcrossref{\DUrole{doc}{fix deform}}}} command will change the simulation box
size/shape and the {\hyperref[\detokenize{fix_move::doc}]{\sphinxcrossref{\DUrole{doc}{fix move}}}} command will change atom
positions and velocities in a prescribed manner).  Also note that many commands
allow variables as arguments for specific parameters, if described in that
manner on their doc pages.  An equal\sphinxhyphen{}style variable can calculate a
time\sphinxhyphen{}dependent quantity, so this is another way to vary a simulation parameter
over time.

\sphinxAtStartPar
If \(N\) is specified as 0, the specified attributes are only changed
once, before the simulation begins.  This is all that is needed if the
associated variables are not time\sphinxhyphen{}dependent.  If \(N > 0\), then changes
are made every \(N\) steps during the simulation, presumably with a
variable that is time\sphinxhyphen{}dependent.

\sphinxAtStartPar
Depending on the value of the \sphinxstyleemphasis{reset} keyword, attributes changed by
this fix will or will not be reset back to their original values at
the end of a simulation.  Even if \sphinxstyleemphasis{reset} is specified as \sphinxstyleemphasis{yes}, a
restart file written during a simulation will contain the modified
settings.

\sphinxAtStartPar
If the \sphinxstyleemphasis{scale} keyword is set to \sphinxstyleemphasis{no}, which is the default, then
the value of the altered parameter will be whatever the variable
generates.  If the \sphinxstyleemphasis{scale} keyword is set to \sphinxstyleemphasis{yes}, then the value
of the altered parameter will be the initial value of that parameter
multiplied by whatever the variable generates (i.e., the variable is
now a “scale factor” applied in (presumably) a time\sphinxhyphen{}varying fashion to
the parameter).

\sphinxAtStartPar
Note that whether scale is \sphinxstyleemphasis{no} or \sphinxstyleemphasis{yes}, internally, the parameters
themselves are actually altered by this fix.  Make sure you use the
\sphinxstyleemphasis{reset yes} option if you want the parameters to be restored to their
initial values after the run.


\bigskip\hrule\bigskip


\sphinxAtStartPar
The \sphinxstyleemphasis{pair} keyword enables various parameters of potentials defined by
the {\hyperref[\detokenize{pair_style::doc}]{\sphinxcrossref{\DUrole{doc}{pair\_style}}}} command to be changed, if the pair
style supports it.  Note that the {\hyperref[\detokenize{pair_style::doc}]{\sphinxcrossref{\DUrole{doc}{pair\_style}}}} and
{\hyperref[\detokenize{pair_coeff::doc}]{\sphinxcrossref{\DUrole{doc}{pair\_coeff}}}} commands must be used in the usual manner
to specify these parameters initially; the fix adapt command simply
overrides the parameters.

\sphinxAtStartPar
The \sphinxstyleemphasis{pstyle} argument is the name of the pair style.  If
{\hyperref[\detokenize{pair_hybrid::doc}]{\sphinxcrossref{\DUrole{doc}{pair\_style hybrid or hybrid/overlay}}}} is used,
\sphinxstyleemphasis{pstyle} should be a sub\sphinxhyphen{}style name.  If there are multiple
sub\sphinxhyphen{}styles using the same pair style, then \sphinxstyleemphasis{pstyle} should be specified
as “style:N”, where \sphinxstyleemphasis{N} is which instance of the pair style you wish to
adapt (e.g., the first or second).  For example, \sphinxstyleemphasis{pstyle} could be
specified as “soft” or “lubricate” or “lj/cut:1” or “lj/cut:2”.  The
\sphinxstyleemphasis{pparam} argument is the name of the parameter to change.  This is the
current list of pair styles and parameters that can be varied by this
fix.  See the doc pages for individual pair styles and their energy
formulas for the meaning of these parameters:


\begin{savenotes}
\sphinxatlongtablestart
\sphinxthistablewithglobalstyle
\makeatletter
  \LTleft \@totalleftmargin plus1fill
  \LTright\dimexpr\columnwidth-\@totalleftmargin-\linewidth\relax plus1fill
\makeatother
\begin{longtable}{lll}
\sphinxtoprule
\endfirsthead

\multicolumn{3}{c}{\sphinxnorowcolor
    \makebox[0pt]{\sphinxtablecontinued{\tablename\ \thetable{} \textendash{} continued from previous page}}%
}\\
\sphinxtoprule
\endhead

\sphinxbottomrule
\multicolumn{3}{r}{\sphinxnorowcolor
    \makebox[0pt][r]{\sphinxtablecontinued{continues on next page}}%
}\\
\endfoot

\endlastfoot
\sphinxtableatstartofbodyhook

\sphinxAtStartPar
{\hyperref[\detokenize{pair_born::doc}]{\sphinxcrossref{\DUrole{doc}{born}}}}
&
\sphinxAtStartPar
a,b,c
&
\sphinxAtStartPar
type pairs
\\
\sphinxhline
\sphinxAtStartPar
{\hyperref[\detokenize{pair_born::doc}]{\sphinxcrossref{\DUrole{doc}{born/coul/long, born/coul/msm}}}}
&
\sphinxAtStartPar
coulombic\_cutoff
&
\sphinxAtStartPar
type global
\\
\sphinxhline
\sphinxAtStartPar
{\hyperref[\detokenize{pair_born_gauss::doc}]{\sphinxcrossref{\DUrole{doc}{born/gauss}}}}
&
\sphinxAtStartPar
biga0,biga1,r0
&
\sphinxAtStartPar
type pairs
\\
\sphinxhline
\sphinxAtStartPar
{\hyperref[\detokenize{pair_buck::doc}]{\sphinxcrossref{\DUrole{doc}{buck, buck/coul/cut}}}}
&
\sphinxAtStartPar
a,c
&
\sphinxAtStartPar
type pairs
\\
\sphinxhline
\sphinxAtStartPar
{\hyperref[\detokenize{pair_buck::doc}]{\sphinxcrossref{\DUrole{doc}{buck/coul/long, buck/coul/msm}}}}
&
\sphinxAtStartPar
a,c,coulombic\_cutoff
&
\sphinxAtStartPar
type pairs
\\
\sphinxhline
\sphinxAtStartPar
{\hyperref[\detokenize{pair_mdf::doc}]{\sphinxcrossref{\DUrole{doc}{buck/mdf}}}}
&
\sphinxAtStartPar
a,c
&
\sphinxAtStartPar
type pairs
\\
\sphinxhline
\sphinxAtStartPar
{\hyperref[\detokenize{pair_coul::doc}]{\sphinxcrossref{\DUrole{doc}{coul/cut, coul/cut/global}}}}
&
\sphinxAtStartPar
scale
&
\sphinxAtStartPar
type pairs
\\
\sphinxhline
\sphinxAtStartPar
{\hyperref[\detokenize{pair_fep_soft::doc}]{\sphinxcrossref{\DUrole{doc}{coul/cut/soft}}}}
&
\sphinxAtStartPar
lambda
&
\sphinxAtStartPar
type pairs
\\
\sphinxhline
\sphinxAtStartPar
{\hyperref[\detokenize{pair_coul::doc}]{\sphinxcrossref{\DUrole{doc}{coul/debye}}}}
&
\sphinxAtStartPar
scale
&
\sphinxAtStartPar
type pairs
\\
\sphinxhline
\sphinxAtStartPar
{\hyperref[\detokenize{pair_coul::doc}]{\sphinxcrossref{\DUrole{doc}{coul/dsf}}}}
&
\sphinxAtStartPar
coulombic\_cutoff
&
\sphinxAtStartPar
type global
\\
\sphinxhline
\sphinxAtStartPar
{\hyperref[\detokenize{pair_coul::doc}]{\sphinxcrossref{\DUrole{doc}{coul/long, coul/msm}}}}
&
\sphinxAtStartPar
coulombic\_cutoff, scale
&
\sphinxAtStartPar
type pairs
\\
\sphinxhline
\sphinxAtStartPar
{\hyperref[\detokenize{pair_fep_soft::doc}]{\sphinxcrossref{\DUrole{doc}{coul/long/soft}}}}
&
\sphinxAtStartPar
scale, lambda, coulombic\_cutoff
&
\sphinxAtStartPar
type pairs
\\
\sphinxhline
\sphinxAtStartPar
{\hyperref[\detokenize{pair_coul_slater::doc}]{\sphinxcrossref{\DUrole{doc}{coul/slater/long}}}}
&
\sphinxAtStartPar
scale
&
\sphinxAtStartPar
type pairs
\\
\sphinxhline
\sphinxAtStartPar
{\hyperref[\detokenize{pair_coul::doc}]{\sphinxcrossref{\DUrole{doc}{coul/streitz}}}}
&
\sphinxAtStartPar
scale
&
\sphinxAtStartPar
type pairs
\\
\sphinxhline
\sphinxAtStartPar
{\hyperref[\detokenize{pair_eam::doc}]{\sphinxcrossref{\DUrole{doc}{eam, eam/alloy, eam/fs}}}}
&
\sphinxAtStartPar
scale
&
\sphinxAtStartPar
type pairs
\\
\sphinxhline
\sphinxAtStartPar
{\hyperref[\detokenize{pair_gauss::doc}]{\sphinxcrossref{\DUrole{doc}{gauss}}}}
&
\sphinxAtStartPar
a
&
\sphinxAtStartPar
type pairs
\\
\sphinxhline
\sphinxAtStartPar
{\hyperref[\detokenize{pair_harmonic_cut::doc}]{\sphinxcrossref{\DUrole{doc}{harmonic/cut}}}}
&
\sphinxAtStartPar
k, cutoff
&
\sphinxAtStartPar
type pairs
\\
\sphinxhline
\sphinxAtStartPar
{\hyperref[\detokenize{pair_kim::doc}]{\sphinxcrossref{\DUrole{doc}{kim}}}}
&
\sphinxAtStartPar
scale
&
\sphinxAtStartPar
type global
\\
\sphinxhline
\sphinxAtStartPar
{\hyperref[\detokenize{pair_mdf::doc}]{\sphinxcrossref{\DUrole{doc}{lennard/mdf}}}}
&
\sphinxAtStartPar
A,B
&
\sphinxAtStartPar
type pairs
\\
\sphinxhline
\sphinxAtStartPar
{\hyperref[\detokenize{pair_class2::doc}]{\sphinxcrossref{\DUrole{doc}{lj/class2}}}}
&
\sphinxAtStartPar
epsilon,sigma
&
\sphinxAtStartPar
type pairs
\\
\sphinxhline
\sphinxAtStartPar
{\hyperref[\detokenize{pair_class2::doc}]{\sphinxcrossref{\DUrole{doc}{lj/class2/coul/cut, lj/class2/coul/long}}}}
&
\sphinxAtStartPar
epsilon,sigma,coulombic\_cutoff
&
\sphinxAtStartPar
type pairs
\\
\sphinxhline
\sphinxAtStartPar
{\hyperref[\detokenize{pair_lj::doc}]{\sphinxcrossref{\DUrole{doc}{lj/cut}}}}
&
\sphinxAtStartPar
epsilon,sigma
&
\sphinxAtStartPar
type pairs
\\
\sphinxhline
\sphinxAtStartPar
{\hyperref[\detokenize{pair_lj_cut_coul::doc}]{\sphinxcrossref{\DUrole{doc}{lj/cut/coul/cut, lj/cut/coul/long, lj/cut/coul/msm}}}}
&
\sphinxAtStartPar
epsilon,sigma,coulombic\_cutoff
&
\sphinxAtStartPar
type pairs
\\
\sphinxhline
\sphinxAtStartPar
{\hyperref[\detokenize{pair_fep_soft::doc}]{\sphinxcrossref{\DUrole{doc}{lj/cut/coul/cut/soft, lj/cut/coul/long/soft}}}}
&
\sphinxAtStartPar
epsilon,sigma,lambda,coulombic\_cutoff
&
\sphinxAtStartPar
type pairs
\\
\sphinxhline
\sphinxAtStartPar
{\hyperref[\detokenize{pair_lj_cut_coul::doc}]{\sphinxcrossref{\DUrole{doc}{lj/cut/coul/dsf}}}}
&
\sphinxAtStartPar
cutoff
&
\sphinxAtStartPar
type global
\\
\sphinxhline
\sphinxAtStartPar
{\hyperref[\detokenize{pair_lj_cut_tip4p::doc}]{\sphinxcrossref{\DUrole{doc}{lj/cut/tip4p/cut}}}}
&
\sphinxAtStartPar
epsilon,sigma,coulombic\_cutoff
&
\sphinxAtStartPar
type pairs
\\
\sphinxhline
\sphinxAtStartPar
{\hyperref[\detokenize{pair_fep_soft::doc}]{\sphinxcrossref{\DUrole{doc}{lj/cut/soft}}}}
&
\sphinxAtStartPar
epsilon,sigma,lambda
&
\sphinxAtStartPar
type pairs
\\
\sphinxhline
\sphinxAtStartPar
{\hyperref[\detokenize{pair_lj_expand::doc}]{\sphinxcrossref{\DUrole{doc}{lj/expand}}}}
&
\sphinxAtStartPar
epsilon,sigma,delta
&
\sphinxAtStartPar
type pairs
\\
\sphinxhline
\sphinxAtStartPar
{\hyperref[\detokenize{pair_mdf::doc}]{\sphinxcrossref{\DUrole{doc}{lj/mdf}}}}
&
\sphinxAtStartPar
epsilon,sigma
&
\sphinxAtStartPar
type pairs
\\
\sphinxhline
\sphinxAtStartPar
{\hyperref[\detokenize{pair_dipole::doc}]{\sphinxcrossref{\DUrole{doc}{lj/sf/dipole/sf}}}}
&
\sphinxAtStartPar
epsilon,sigma,scale
&
\sphinxAtStartPar
type pairs
\\
\sphinxhline
\sphinxAtStartPar
{\hyperref[\detokenize{pair_lubricate::doc}]{\sphinxcrossref{\DUrole{doc}{lubricate}}}}
&
\sphinxAtStartPar
mu
&
\sphinxAtStartPar
global
\\
\sphinxhline
\sphinxAtStartPar
{\hyperref[\detokenize{pair_meam::doc}]{\sphinxcrossref{\DUrole{doc}{meam}}}}
&
\sphinxAtStartPar
scale
&
\sphinxAtStartPar
type pairs
\\
\sphinxhline
\sphinxAtStartPar
{\hyperref[\detokenize{pair_mie::doc}]{\sphinxcrossref{\DUrole{doc}{mie/cut}}}}
&
\sphinxAtStartPar
epsilon,sigma,gamma\_repulsive,gamma\_attractive
&
\sphinxAtStartPar
type pairs
\\
\sphinxhline
\sphinxAtStartPar
{\hyperref[\detokenize{pair_morse::doc}]{\sphinxcrossref{\DUrole{doc}{morse, morse/smooth/linear}}}}
&
\sphinxAtStartPar
D0,R0,alpha
&
\sphinxAtStartPar
type pairs
\\
\sphinxhline
\sphinxAtStartPar
{\hyperref[\detokenize{pair_morse::doc}]{\sphinxcrossref{\DUrole{doc}{morse/soft}}}}
&
\sphinxAtStartPar
D0,R0,alpha,lambda
&
\sphinxAtStartPar
type pairs
\\
\sphinxhline
\sphinxAtStartPar
{\hyperref[\detokenize{pair_nm::doc}]{\sphinxcrossref{\DUrole{doc}{nm/cut}}}}
&
\sphinxAtStartPar
E0,R0,m,n
&
\sphinxAtStartPar
type pairs
\\
\sphinxhline
\sphinxAtStartPar
{\hyperref[\detokenize{pair_nm::doc}]{\sphinxcrossref{\DUrole{doc}{nm/cut/coul/cut, nm/cut/coul/long}}}}
&
\sphinxAtStartPar
E0,R0,m,n,coulombic\_cutoff
&
\sphinxAtStartPar
type pairs
\\
\sphinxhline
\sphinxAtStartPar
{\hyperref[\detokenize{pair_pace::doc}]{\sphinxcrossref{\DUrole{doc}{pace, pace/extrapolation}}}}
&
\sphinxAtStartPar
scale
&
\sphinxAtStartPar
type pairs
\\
\sphinxhline
\sphinxAtStartPar
{\hyperref[\detokenize{pair_quip::doc}]{\sphinxcrossref{\DUrole{doc}{quip}}}}
&
\sphinxAtStartPar
scale
&
\sphinxAtStartPar
type global
\\
\sphinxhline
\sphinxAtStartPar
{\hyperref[\detokenize{pair_snap::doc}]{\sphinxcrossref{\DUrole{doc}{snap}}}}
&
\sphinxAtStartPar
scale
&
\sphinxAtStartPar
type pairs
\\
\sphinxhline
\sphinxAtStartPar
{\hyperref[\detokenize{pair_spin_dmi::doc}]{\sphinxcrossref{\DUrole{doc}{spin/dmi}}}}
&
\sphinxAtStartPar
coulombic\_cutoff
&
\sphinxAtStartPar
type global
\\
\sphinxhline
\sphinxAtStartPar
{\hyperref[\detokenize{pair_spin_exchange::doc}]{\sphinxcrossref{\DUrole{doc}{spin/exchange}}}}
&
\sphinxAtStartPar
coulombic\_cutoff
&
\sphinxAtStartPar
type global
\\
\sphinxhline
\sphinxAtStartPar
{\hyperref[\detokenize{pair_spin_magelec::doc}]{\sphinxcrossref{\DUrole{doc}{spin/magelec}}}}
&
\sphinxAtStartPar
coulombic\_cutoff
&
\sphinxAtStartPar
type global
\\
\sphinxhline
\sphinxAtStartPar
{\hyperref[\detokenize{pair_spin_neel::doc}]{\sphinxcrossref{\DUrole{doc}{spin/neel}}}}
&
\sphinxAtStartPar
coulombic\_cutoff
&
\sphinxAtStartPar
type global
\\
\sphinxhline
\sphinxAtStartPar
{\hyperref[\detokenize{pair_soft::doc}]{\sphinxcrossref{\DUrole{doc}{soft}}}}
&
\sphinxAtStartPar
a
&
\sphinxAtStartPar
type pairs
\\
\sphinxhline
\sphinxAtStartPar
{\hyperref[\detokenize{pair_table::doc}]{\sphinxcrossref{\DUrole{doc}{table}}}}
&
\sphinxAtStartPar
table\_cutoff
&
\sphinxAtStartPar
type pairs
\\
\sphinxhline
\sphinxAtStartPar
{\hyperref[\detokenize{pair_ufm::doc}]{\sphinxcrossref{\DUrole{doc}{ufm}}}}
&
\sphinxAtStartPar
epsilon,sigma,scale
&
\sphinxAtStartPar
type pairs
\\
\sphinxhline
\sphinxAtStartPar
{\hyperref[\detokenize{pair_wf_cut::doc}]{\sphinxcrossref{\DUrole{doc}{wf/cut}}}}
&
\sphinxAtStartPar
epsilon,sigma,nu,mu
&
\sphinxAtStartPar
type pairs
\\
\sphinxbottomrule
\end{longtable}
\sphinxtableafterendhook
\sphinxatlongtableend
\end{savenotes}

\begin{sphinxadmonition}{note}{Note:}
\sphinxAtStartPar
It is easy to add new pairwise potentials and their parameters
to this list.  All it typically takes is adding an extract() method to
the pair\_*.cpp file associated with the potential.
\end{sphinxadmonition}

\sphinxAtStartPar
Some parameters are global settings for the pair style (e.g., the
viscosity setting “mu” for {\hyperref[\detokenize{pair_lubricate::doc}]{\sphinxcrossref{\DUrole{doc}{pair\_style lubricate}}}}).
Other parameters apply to atom type pairs within the pair style (e.g., the
prefactor \(a\) for {\hyperref[\detokenize{pair_soft::doc}]{\sphinxcrossref{\DUrole{doc}{pair\_style soft}}}}).

\sphinxAtStartPar
Note that for many of the potentials, the parameter that can be varied
is effectively a prefactor on the entire energy expression for the
potential (e.g., the lj/cut epsilon).  The parameters listed as “scale”
are exactly that, since the energy expression for the
{\hyperref[\detokenize{pair_coul::doc}]{\sphinxcrossref{\DUrole{doc}{coul/cut}}}} potential (for example) has no labeled
prefactor in its formula.  To apply an effective prefactor to some
potentials, multiple parameters need to be altered.  For example, the
{\hyperref[\detokenize{pair_buck::doc}]{\sphinxcrossref{\DUrole{doc}{Buckingham potential}}}} needs both the \(A\) and
\(C\) terms altered together.  To scale the Buckingham potential, you
should thus list the pair style twice, once for \(A\) and once for
\(C\).

\sphinxAtStartPar
If a type pair parameter is specified, the \(I\) and \(J\) settings
should be specified to indicate which type pairs to apply it to.  If a global
parameter is specified, the \(I\) and \(J\) settings still need to be
specified, but are ignored.

\sphinxAtStartPar
Similar to the {\hyperref[\detokenize{pair_coeff::doc}]{\sphinxcrossref{\DUrole{doc}{pair\_coeff command}}}}, \(I\) and
\(J\) can be specified in one of several ways.  Explicit numeric values
can be used for each, as in the first example above.  Or, one or both of
the types in the I,J pair can be a {\hyperref[\detokenize{Howto_type_labels::doc}]{\sphinxcrossref{\DUrole{doc}{type label}}}}.
LAMMPS sets the coefficients for the symmetric \(J,I\) interaction to
the same values.

\sphinxAtStartPar
A wild\sphinxhyphen{}card asterisk can be used in place of or in conjunction with
the \(I,J\) arguments to set the coefficients for multiple pairs of atom
types.  This takes the form “*” or “*n” or “m*” or “m*n”.  If \(N\)
is the number of atom types, then an asterisk with no numeric values
means all types from 1 to \(N\).  A leading asterisk means all types from
1 to n (inclusive).  A trailing asterisk means all types from m to \(N\)
(inclusive).  A middle asterisk means all types from m to n
(inclusive).  For the asterisk syntax, note that only type pairs with
\(I \le J\) are considered; if asterisks imply type pairs where
\(J < I\), they are ignored.

\sphinxAtStartPar
IMPORTANT NOTE: If {\hyperref[\detokenize{pair_hybrid::doc}]{\sphinxcrossref{\DUrole{doc}{pair\_style hybrid or hybrid/overlay}}}} is being used, then the \sphinxstyleemphasis{pstyle} will be a sub\sphinxhyphen{}style
name.  You must specify \(I,J\) arguments that correspond to type pair
values defined (via the {\hyperref[\detokenize{pair_coeff::doc}]{\sphinxcrossref{\DUrole{doc}{pair\_coeff}}}} command) for
that sub\sphinxhyphen{}style.

\sphinxAtStartPar
The \sphinxstyleemphasis{v\_name} argument for keyword \sphinxstyleemphasis{pair} is the name of an
{\hyperref[\detokenize{variable::doc}]{\sphinxcrossref{\DUrole{doc}{equal\sphinxhyphen{}style variable}}}} which will be evaluated each time
this fix is invoked to set the parameter to a new value.  It should be
specified as v\_name, where name is the variable name.  Equal\sphinxhyphen{}style
variables can specify formulas with various mathematical functions, and
include {\hyperref[\detokenize{thermo_style::doc}]{\sphinxcrossref{\DUrole{doc}{thermo\_style}}}} command keywords for the
simulation box parameters and timestep and elapsed time.  Thus it is
easy to specify parameters that change as a function of time or span
consecutive runs in a continuous fashion.  For the latter, see the
\sphinxstyleemphasis{start} and \sphinxstyleemphasis{stop} keywords of the {\hyperref[\detokenize{run::doc}]{\sphinxcrossref{\DUrole{doc}{run}}}} command and the
\sphinxstyleemphasis{elaplong} keyword of {\hyperref[\detokenize{thermo_style::doc}]{\sphinxcrossref{\DUrole{doc}{thermo\_style custom}}}} for
details.

\sphinxAtStartPar
For example, these commands would change the prefactor coefficient of
the {\hyperref[\detokenize{pair_soft::doc}]{\sphinxcrossref{\DUrole{doc}{pair\_style soft}}}} potential from 10.0 to 30.0 in a
linear fashion over the course of a simulation:

\begin{sphinxVerbatim}[commandchars=\\\{\}]
\PYG{k}{variable }\PYG{n+nv+nvIdentifier}{prefactor}\PYG{+w}{ }\PYG{n}{equal}\PYG{+w}{ }\PYG{n}{ramp}\PYG{n+nv}{(}\PYG{n+nv}{10,30}\PYG{n+nv}{)}
\PYG{k}{fix }\PYG{n+nv+nvIdentifier}{1}\PYG{+w}{ }\PYG{n+nv+nvIdentifier}{all}\PYG{+w}{ }\PYG{n}{adapt}\PYG{+w}{ }\PYG{l+m}{1}\PYG{+w}{ }\PYG{n}{pair}\PYG{+w}{ }\PYG{n}{soft}\PYG{+w}{ }\PYG{n}{a}\PYG{+w}{ }\PYG{o}{*}\PYG{+w}{ }\PYG{o}{*}\PYG{+w}{ }\PYG{n}{v\PYGZus{}prefactor}
\end{sphinxVerbatim}


\bigskip\hrule\bigskip


\sphinxAtStartPar
The \sphinxstyleemphasis{bond} keyword uses the specified variable to change the value of
a bond coefficient over time, very similar to how the \sphinxstyleemphasis{pair} keyword
operates. The only difference is that now a bond coefficient for a
given bond type is adapted.

\sphinxAtStartPar
A wild\sphinxhyphen{}card asterisk can be used in place of or in conjunction with the
bond type argument to set the coefficients for multiple bond types.
This takes the form “*” or “*n” or “m*” or “m*n”.  If \(N\) is
the number of bond types, then an asterisk with no numeric values means
all types from 1 to \(N\).  A leading asterisk means all types from
1 to n (inclusive).  A trailing asterisk means all types from m to
\(N\) (inclusive).  A middle asterisk means all types from m to n
(inclusive).

\sphinxAtStartPar
Currently \sphinxstyleemphasis{bond} does not support bond\_style hybrid nor bond\_style
hybrid/overlay as bond styles. The bond styles that currently work
with fix\_adapt are


\begin{savenotes}\sphinxattablestart
\sphinxthistablewithglobalstyle
\centering
\begin{tabulary}{\linewidth}[t]{TTT}
\sphinxtoprule
\sphinxtableatstartofbodyhook
\sphinxAtStartPar
{\hyperref[\detokenize{bond_class2::doc}]{\sphinxcrossref{\DUrole{doc}{class2}}}}
&
\sphinxAtStartPar
r0
&
\sphinxAtStartPar
type bonds
\\
\sphinxhline
\sphinxAtStartPar
{\hyperref[\detokenize{bond_fene::doc}]{\sphinxcrossref{\DUrole{doc}{fene}}}}
&
\sphinxAtStartPar
k,r0
&
\sphinxAtStartPar
type bonds
\\
\sphinxhline
\sphinxAtStartPar
{\hyperref[\detokenize{bond_fene::doc}]{\sphinxcrossref{\DUrole{doc}{fene/nm}}}}
&
\sphinxAtStartPar
k,r0
&
\sphinxAtStartPar
type bonds
\\
\sphinxhline
\sphinxAtStartPar
{\hyperref[\detokenize{bond_gromos::doc}]{\sphinxcrossref{\DUrole{doc}{gromos}}}}
&
\sphinxAtStartPar
k,r0
&
\sphinxAtStartPar
type bonds
\\
\sphinxhline
\sphinxAtStartPar
{\hyperref[\detokenize{bond_harmonic::doc}]{\sphinxcrossref{\DUrole{doc}{harmonic}}}}
&
\sphinxAtStartPar
k,r0
&
\sphinxAtStartPar
type bonds
\\
\sphinxhline
\sphinxAtStartPar
{\hyperref[\detokenize{bond_morse::doc}]{\sphinxcrossref{\DUrole{doc}{morse}}}}
&
\sphinxAtStartPar
r0
&
\sphinxAtStartPar
type bonds
\\
\sphinxhline
\sphinxAtStartPar
{\hyperref[\detokenize{bond_nonlinear::doc}]{\sphinxcrossref{\DUrole{doc}{nonlinear}}}}
&
\sphinxAtStartPar
epsilon,r0
&
\sphinxAtStartPar
type bonds
\\
\sphinxbottomrule
\end{tabulary}
\sphinxtableafterendhook\par
\sphinxattableend\end{savenotes}


\bigskip\hrule\bigskip


\sphinxAtStartPar
\DUrole{versionmodified,added}{Added in version 4May2022.}

\sphinxAtStartPar
The \sphinxstyleemphasis{angle} keyword uses the specified variable to change the value of
an angle coefficient over time, very similar to how the \sphinxstyleemphasis{pair} keyword
operates. The only difference is that now an angle coefficient for a
given angle type is adapted.

\sphinxAtStartPar
A wild\sphinxhyphen{}card asterisk can be used in place of or in conjunction with the
angle type argument to set the coefficients for multiple angle types.
This takes the form “*” or “*n” or “m*” or “m*n”.  If \(N\) is
the number of angle types, then an asterisk with no numeric values means
all types from 1 to \(N\).  A leading asterisk means all types from
1 to n (inclusive).  A trailing asterisk means all types from m to
\(N\) (inclusive).  A middle asterisk means all types from m to n
(inclusive).

\sphinxAtStartPar
Currently \sphinxstyleemphasis{angle} does not support angle\_style hybrid nor angle\_style
hybrid/overlay as angle styles. The angle styles that currently work
with fix\_adapt are


\begin{savenotes}\sphinxattablestart
\sphinxthistablewithglobalstyle
\centering
\begin{tabulary}{\linewidth}[t]{TTT}
\sphinxtoprule
\sphinxtableatstartofbodyhook
\sphinxAtStartPar
{\hyperref[\detokenize{angle_harmonic::doc}]{\sphinxcrossref{\DUrole{doc}{harmonic}}}}
&
\sphinxAtStartPar
k,theta0
&
\sphinxAtStartPar
type angles
\\
\sphinxhline
\sphinxAtStartPar
{\hyperref[\detokenize{angle_cosine::doc}]{\sphinxcrossref{\DUrole{doc}{cosine}}}}
&
\sphinxAtStartPar
k
&
\sphinxAtStartPar
type angles
\\
\sphinxbottomrule
\end{tabulary}
\sphinxtableafterendhook\par
\sphinxattableend\end{savenotes}

\sphinxAtStartPar
Note that internally, theta0 is stored in radians, so the variable
this fix uses to reset theta0 needs to generate values in radians.


\bigskip\hrule\bigskip


\sphinxAtStartPar
The \sphinxstyleemphasis{kspace} keyword used the specified variable as a scale factor on
the energy, forces, virial calculated by whatever \(k\)\sphinxhyphen{}space solver is
defined by the {\hyperref[\detokenize{kspace_style::doc}]{\sphinxcrossref{\DUrole{doc}{kspace\_style}}}} command.  If the
variable has a value of 1.0, then the solver is unaltered.

\sphinxAtStartPar
The \sphinxstyleemphasis{kspace} keyword works this way whether the \sphinxstyleemphasis{scale} keyword
is set to \sphinxstyleemphasis{no} or \sphinxstyleemphasis{yes}.


\bigskip\hrule\bigskip


\sphinxAtStartPar
The \sphinxstyleemphasis{atom} keyword enables various atom properties to be changed.  The
\sphinxstyleemphasis{aparam} argument is the name of the parameter to change.  This is the
current list of atom parameters that can be varied by this fix:
\begin{itemize}
\item {} 
\sphinxAtStartPar
charge = charge on particle

\item {} 
\sphinxAtStartPar
diameter or diameter/disc = diameter of particle

\end{itemize}

\sphinxAtStartPar
The \sphinxstyleemphasis{v\_name} argument of the \sphinxstyleemphasis{atom} keyword is the name of an
{\hyperref[\detokenize{variable::doc}]{\sphinxcrossref{\DUrole{doc}{equal\sphinxhyphen{}style variable}}}} which will be evaluated each
time this fix is invoked to set, or scale the parameter to a new
value.  It should be specified as v\_name, where name is the variable
name.  See the discussion above describing the formulas associated
with equal\sphinxhyphen{}style variables.  The new value is assigned to the
corresponding attribute for all atoms in the fix group.

\sphinxAtStartPar
If the atom parameter is \sphinxstyleemphasis{diameter} and per\sphinxhyphen{}atom density and per\sphinxhyphen{}atom
mass are defined for particles (e.g., {\hyperref[\detokenize{atom_style::doc}]{\sphinxcrossref{\DUrole{doc}{atom\_style granular}}}}), then the mass of each particle is, by default, also
changed when the diameter changes. The mass is set from the particle
volume for 3d systems (density is assumed to stay constant). For 2d,
the default is for LAMMPS to model particles with a radius attribute
as spheres. However, if the atom parameter is \sphinxstyleemphasis{diameter/disc}, then the
mass is set from the particle area (the density is assumed to be in
mass/distance\(^2\) units). The mass of the particle may also be kept
constant if the \sphinxstyleemphasis{mass} keyword is set to \sphinxstyleemphasis{no}. This can be useful to account
for diameter changes that do not involve mass changes (e.g., thermal
expansion).

\sphinxAtStartPar
For example, these commands would shrink the diameter of all granular
particles in the “center” group from 1.0 to 0.1 in a linear fashion
over the course of a 1000\sphinxhyphen{}step simulation:

\begin{sphinxVerbatim}[commandchars=\\\{\}]
\PYG{k}{variable }\PYG{n+nv+nvIdentifier}{size}\PYG{+w}{ }\PYG{n}{equal}\PYG{+w}{ }\PYG{n}{ramp}\PYG{n+nv}{(}\PYG{n+nv}{1.0,0.1}\PYG{n+nv}{)}
\PYG{k}{fix }\PYG{n+nv+nvIdentifier}{1}\PYG{+w}{ }\PYG{n+nv+nvIdentifier}{center}\PYG{+w}{ }\PYG{n}{adapt}\PYG{+w}{ }\PYG{l+m}{10}\PYG{+w}{ }\PYG{n}{atom}\PYG{+w}{ }\PYG{n}{diameter}\PYG{+w}{ }\PYG{n}{v\PYGZus{}size}
\end{sphinxVerbatim}


\bigskip\hrule\bigskip


\sphinxAtStartPar
This fix can be used in long simulations which are restarted one or
more times to continuously adapt simulation parameters, but it must be
done carefully.  There are two issues to consider.  The first is how
to adapt the parameters in a continuous manner from one simulation to
the next.  The second is how, if desired, to reset the parameters to
their original values at the end of the last restarted run.

\sphinxAtStartPar
Note that all the parameters changed by this fix are written into a
restart file in their current changed state.  A new restarted
simulation does not know the original time=0 values, unless the
input script explicitly resets the parameters (after the restart file
is read) to the original values.

\sphinxAtStartPar
Also note that the time\sphinxhyphen{}dependent variable(s) used in the restart
script should typically be written as a function of time elapsed since
the original simulation began.

\sphinxAtStartPar
With this in mind, if the \sphinxstyleemphasis{scale} keyword is set to \sphinxstyleemphasis{no} (the default)
in a restarted simulation, original parameters are not needed.  The
adapted parameters should seamlessly continue their variation relative
to the preceding simulation.

\sphinxAtStartPar
If the \sphinxstyleemphasis{scale} keyword is set to \sphinxstyleemphasis{yes}, then the input script should
typically reset the parameters being adapted to their original values,
so that the scaling formula specified by the variable will operate
correctly.  An exception is if the \sphinxstyleemphasis{atom} keyword is being used with
\sphinxstyleemphasis{scale yes}.  In this case, information is added to the restart file
so that per\sphinxhyphen{}atom properties in the new run will automatically be
scaled relative to their original values.  This will only work if the
fix adapt command specified in the restart script has the same ID as
the one used in the original script.

\sphinxAtStartPar
In a restarted run, if the \sphinxstyleemphasis{reset} keyword is set to \sphinxstyleemphasis{yes}, and the
run ends in this script (as opposed to just writing more restart
files), parameters will be restored to the values they were at the
beginning of the run command in the restart script, which as
explained above, may or may not be the original values of the
parameters.  Again, an exception is if the \sphinxstyleemphasis{atom} keyword is being
used with \sphinxstyleemphasis{reset yes} (in all the runs). In that case, the original
per\sphinxhyphen{}atom parameters are stored in the restart file, and will be
restored when the restarted run finally completes.


\bigskip\hrule\bigskip



\subsection{Restart, fix\_modify, output, run start/stop, minimize info}
\label{\detokenize{fix_adapt:restart-fix-modify-output-run-start-stop-minimize-info}}
\sphinxAtStartPar
If the \sphinxstyleemphasis{atom} keyword is used and the \sphinxstyleemphasis{scale} or \sphinxstyleemphasis{reset} keyword is
set to \sphinxstyleemphasis{yes}, then this fix writes information to a restart file so
that in a restarted run scaling can continue in a seamless manner
and/or the per\sphinxhyphen{}atom values can be restored, as explained above.

\sphinxAtStartPar
None of the {\hyperref[\detokenize{fix_modify::doc}]{\sphinxcrossref{\DUrole{doc}{fix\_modify}}}} options are relevant to
this fix.  No global or per\sphinxhyphen{}atom quantities are stored by this fix for
access by various {\hyperref[\detokenize{Howto_output::doc}]{\sphinxcrossref{\DUrole{doc}{output commands}}}}.  No parameter
of this fix can be used with the \sphinxstyleemphasis{start/stop} keywords of the
{\hyperref[\detokenize{run::doc}]{\sphinxcrossref{\DUrole{doc}{run}}}} command.  This fix is not invoked during {\hyperref[\detokenize{minimize::doc}]{\sphinxcrossref{\DUrole{doc}{energy
minimization}}}}.

\sphinxAtStartPar
For {\hyperref[\detokenize{run_style::doc}]{\sphinxcrossref{\DUrole{doc}{rRESPA time integration}}}}, this fix changes
parameters on the outermost rRESPA level.


\subsection{Restrictions}
\label{\detokenize{fix_adapt:restrictions}}\begin{quote}

\sphinxAtStartPar
none
\end{quote}


\subsection{Related commands}
\label{\detokenize{fix_adapt:related-commands}}
\sphinxAtStartPar
{\hyperref[\detokenize{compute_ti::doc}]{\sphinxcrossref{\DUrole{doc}{compute ti}}}}, {\hyperref[\detokenize{fix_adapt_fep::doc}]{\sphinxcrossref{\DUrole{doc}{fix adapt/fep}}}}


\subsection{Default}
\label{\detokenize{fix_adapt:default}}
\sphinxAtStartPar
The option defaults are scale = no, reset = no, mass = yes.

\sphinxstepscope

\index{fix adapt/fep@\spxentry{fix adapt/fep}}\ignorespaces 

\section{fix adapt/fep command}
\label{\detokenize{fix_adapt_fep:fix-adapt-fep-command}}\label{\detokenize{fix_adapt_fep:index-0}}\label{\detokenize{fix_adapt_fep::doc}}

\subsection{Syntax}
\label{\detokenize{fix_adapt_fep:syntax}}
\begin{sphinxVerbatim}[commandchars=\\\{\}]
\PYG{k}{fix }\PYG{n+nv+nvIdentifier}{ID}\PYG{+w}{ }\PYG{n+nv+nvIdentifier}{group\PYGZhy{}ID}\PYG{+w}{ }\PYG{n}{adapt}\PYG{o}{/}\PYG{n}{fep}\PYG{+w}{ }\PYG{n}{N}\PYG{+w}{ }\PYG{n}{attribute}\PYG{+w}{ }\PYG{n}{args}\PYG{+w}{ }\PYG{n}{...}\PYG{+w}{ }\PYG{n}{keyword}\PYG{+w}{ }\PYG{n}{value}\PYG{+w}{ }\PYG{n}{...}
\end{sphinxVerbatim}
\begin{itemize}
\item {} 
\sphinxAtStartPar
ID, group\sphinxhyphen{}ID are documented in {\hyperref[\detokenize{fix::doc}]{\sphinxcrossref{\DUrole{doc}{fix}}}} command

\item {} 
\sphinxAtStartPar
adapt/fep = style name of this fix command

\item {} 
\sphinxAtStartPar
N = adapt simulation settings every this many timesteps

\item {} 
\sphinxAtStartPar
one or more attribute/arg pairs may be appended

\item {} 
\sphinxAtStartPar
attribute = \sphinxstyleemphasis{pair} or \sphinxstyleemphasis{kspace} or \sphinxstyleemphasis{atom}
\begin{sphinxalltt}
\sphinxstyleemphasis{pair} args = pstyle pparam I J v\_name
  pstyle = pair style name (e.g., lj/cut)
  pparam = parameter to adapt over time
  I,J = type pair(s) to set parameter for (integer or type label)
  v\_name = variable with name that calculates value of pparam
\sphinxstyleemphasis{kspace} arg = v\_name
  v\_name = variable with name that calculates scale factor on K\sphinxhyphen{}space terms
\sphinxstyleemphasis{atom} args = aparam v\_name
  aparam = parameter to adapt over time
  I = type(s) to set parameter for (integer or type label)
  v\_name = variable with name that calculates value of aparam
\end{sphinxalltt}

\item {} 
\sphinxAtStartPar
zero or more keyword/value pairs may be appended

\item {} 
\sphinxAtStartPar
keyword = \sphinxstyleemphasis{scale} or \sphinxstyleemphasis{reset} or \sphinxstyleemphasis{after}
\begin{sphinxalltt}
\sphinxstyleemphasis{scale} value = \sphinxstyleemphasis{no} or \sphinxstyleemphasis{yes}
  \sphinxstyleemphasis{no} = the variable value is the new setting
  \sphinxstyleemphasis{yes} = the variable value multiplies the original setting
\sphinxstyleemphasis{reset} value = \sphinxstyleemphasis{no} or \sphinxstyleemphasis{yes}
  \sphinxstyleemphasis{no} = values will remain altered at the end of a run
  \sphinxstyleemphasis{yes} = reset altered values to their original values at the end
  of a run
\sphinxstyleemphasis{after} value = \sphinxstyleemphasis{no} or \sphinxstyleemphasis{yes}
  \sphinxstyleemphasis{no} = parameters are adapted at timestep N
  \sphinxstyleemphasis{yes} = parameters are adapted one timestep after N
\end{sphinxalltt}

\end{itemize}


\subsection{Examples}
\label{\detokenize{fix_adapt_fep:examples}}
\begin{sphinxVerbatim}[commandchars=\\\{\}]
\PYG{k}{fix }\PYG{n+nv+nvIdentifier}{1}\PYG{+w}{ }\PYG{n+nv+nvIdentifier}{all}\PYG{+w}{ }\PYG{n}{adapt}\PYG{o}{/}\PYG{n}{fep}\PYG{+w}{ }\PYG{l+m}{1}\PYG{+w}{ }\PYG{n}{pair}\PYG{+w}{ }\PYG{n}{soft}\PYG{+w}{ }\PYG{n}{a}\PYG{+w}{ }\PYG{l+m}{1}\PYG{+w}{ }\PYG{l+m}{1}\PYG{+w}{ }\PYG{n}{v\PYGZus{}prefactor}
\PYG{k}{fix }\PYG{n+nv+nvIdentifier}{1}\PYG{+w}{ }\PYG{n+nv+nvIdentifier}{all}\PYG{+w}{ }\PYG{n}{adapt}\PYG{o}{/}\PYG{n}{fep}\PYG{+w}{ }\PYG{l+m}{1}\PYG{+w}{ }\PYG{n}{pair}\PYG{+w}{ }\PYG{n}{soft}\PYG{+w}{ }\PYG{n}{a}\PYG{+w}{ }\PYG{l+m}{2}\PYG{o}{*}\PYG{+w}{ }\PYG{l+m}{3}\PYG{+w}{ }\PYG{n}{v\PYGZus{}prefactor}
\PYG{k}{fix }\PYG{n+nv+nvIdentifier}{1}\PYG{+w}{ }\PYG{n+nv+nvIdentifier}{all}\PYG{+w}{ }\PYG{n}{adapt}\PYG{o}{/}\PYG{n}{fep}\PYG{+w}{ }\PYG{l+m}{1}\PYG{+w}{ }\PYG{n}{pair}\PYG{+w}{ }\PYG{n}{lj}\PYG{o}{/}\PYG{n}{cut}\PYG{+w}{ }\PYG{n}{epsilon}\PYG{+w}{ }\PYG{o}{*}\PYG{+w}{ }\PYG{o}{*}\PYG{+w}{ }\PYG{n}{v\PYGZus{}scale1}\PYG{+w}{ }\PYG{n}{coul}\PYG{o}{/}\PYG{n}{cut}\PYG{+w}{ }\PYG{n}{scale}\PYG{+w}{ }\PYG{l+m}{3}\PYG{+w}{ }\PYG{l+m}{3}\PYG{+w}{ }\PYG{n}{v\PYGZus{}scale2}\PYG{+w}{ }\PYG{n}{scale}\PYG{+w}{ }\PYG{n}{yes}\PYG{+w}{ }\PYG{n}{reset}\PYG{+w}{ }\PYG{n}{yes}
\PYG{k}{fix }\PYG{n+nv+nvIdentifier}{1}\PYG{+w}{ }\PYG{n+nv+nvIdentifier}{all}\PYG{+w}{ }\PYG{n}{adapt}\PYG{o}{/}\PYG{n}{fep}\PYG{+w}{ }\PYG{l+m}{10}\PYG{+w}{ }\PYG{n}{atom}\PYG{+w}{ }\PYG{n}{diameter}\PYG{+w}{ }\PYG{l+m}{1}\PYG{+w}{ }\PYG{n}{v\PYGZus{}size}

\PYG{k}{labelmap}\PYG{+w}{ }\PYG{n}{atom}\PYG{+w}{ }\PYG{l+m}{1}\PYG{+w}{ }\PYG{n}{c1}
\PYG{k}{fix }\PYG{n+nv+nvIdentifier}{1}\PYG{+w}{ }\PYG{n+nv+nvIdentifier}{all}\PYG{+w}{ }\PYG{n}{adapt}\PYG{o}{/}\PYG{n}{fep}\PYG{+w}{ }\PYG{l+m}{1}\PYG{+w}{ }\PYG{n}{pair}\PYG{+w}{ }\PYG{n}{soft}\PYG{+w}{ }\PYG{n}{a}\PYG{+w}{ }\PYG{n}{c1}\PYG{+w}{ }\PYG{n}{c1}\PYG{+w}{ }\PYG{n}{v\PYGZus{}prefactor}
\end{sphinxVerbatim}

\sphinxAtStartPar
Example input scripts available: examples/PACKAGES/fep


\subsection{Description}
\label{\detokenize{fix_adapt_fep:description}}
\sphinxAtStartPar
Change or adapt one or more specific simulation attributes or settings
over time as a simulation runs.

\sphinxAtStartPar
This is an enhanced version of the {\hyperref[\detokenize{fix_adapt::doc}]{\sphinxcrossref{\DUrole{doc}{fix adapt}}}} command
with two differences:
\begin{itemize}
\item {} 
\sphinxAtStartPar
It is possible to modify the charges of chosen atom types only,
instead of scaling all the charges in the system.

\item {} 
\sphinxAtStartPar
There is a new option \sphinxstyleemphasis{after} for better compatibility with
{\hyperref[\detokenize{fix_ave_time::doc}]{\sphinxcrossref{\DUrole{doc}{fix ave/time}}}}.

\end{itemize}

\sphinxAtStartPar
This version is suited for free energy calculations using
{\hyperref[\detokenize{compute_ti::doc}]{\sphinxcrossref{\DUrole{doc}{compute ti}}}} or {\hyperref[\detokenize{compute_fep::doc}]{\sphinxcrossref{\DUrole{doc}{compute fep}}}}.

\sphinxAtStartPar
If \sphinxstyleemphasis{N} is specified as 0, the specified attributes are only changed
once, before the simulation begins.  This is all that is needed if the
associated variables are not time\sphinxhyphen{}dependent.  If \sphinxstyleemphasis{N} \textgreater{} 0, then changes
are made every \sphinxstyleemphasis{N} steps during the simulation, presumably with a
variable that is time\sphinxhyphen{}dependent.

\sphinxAtStartPar
Depending on the value of the \sphinxstyleemphasis{reset} keyword, attributes changed by
this fix will or will not be reset back to their original values at
the end of a simulation.  Even if \sphinxstyleemphasis{reset} is specified as \sphinxstyleemphasis{yes}, a
restart file written during a simulation will contain the modified
settings.

\sphinxAtStartPar
If the \sphinxstyleemphasis{scale} keyword is set to \sphinxstyleemphasis{no}, then the value the parameter is
set to will be whatever the variable generates.  If the \sphinxstyleemphasis{scale}
keyword is set to \sphinxstyleemphasis{yes}, then the value of the altered parameter will
be the initial value of that parameter multiplied by whatever the
variable generates (i.e., the variable is now a “scale factor” applied
in (presumably) a time\sphinxhyphen{}varying fashion to the parameter).  Internally,
the parameters themselves are actually altered; make sure you use the
\sphinxstyleemphasis{reset yes} option if you want the parameters to be restored to their
initial values after the run.

\sphinxAtStartPar
If the \sphinxstyleemphasis{after} keyword is set to \sphinxstyleemphasis{yes}, then the parameters are
changed one timestep after the multiple of N. In this manner, if a fix
such as “fix ave/time” is used to calculate averages at every N
timesteps, all the contributions to the average will be obtained with
the same values of the parameters.


\bigskip\hrule\bigskip


\sphinxAtStartPar
The \sphinxstyleemphasis{pair} keyword enables various parameters of potentials defined by
the {\hyperref[\detokenize{pair_style::doc}]{\sphinxcrossref{\DUrole{doc}{pair\_style}}}} command to be changed, if the pair
style supports it.  Note that the {\hyperref[\detokenize{pair_style::doc}]{\sphinxcrossref{\DUrole{doc}{pair\_style}}}} and
{\hyperref[\detokenize{pair_coeff::doc}]{\sphinxcrossref{\DUrole{doc}{pair\_coeff}}}} commands must be used in the usual manner
to specify these parameters initially; the fix adapt command simply
overrides the parameters.

\sphinxAtStartPar
The \sphinxstyleemphasis{pstyle} argument is the name of the pair style.  If {\hyperref[\detokenize{pair_hybrid::doc}]{\sphinxcrossref{\DUrole{doc}{pair\_style hybrid or hybrid/overlay}}}} is used, \sphinxstyleemphasis{pstyle} should be
a sub\sphinxhyphen{}style name.  For example, \sphinxstyleemphasis{pstyle} could be specified as “soft”
or “lubricate”.  The \sphinxstyleemphasis{pparam} argument is the name of the parameter to
change.  This is the current list of pair styles and parameters that
can be varied by this fix.  See the doc pages for individual pair
styles and their energy formulas for the meaning of these parameters:


\begin{savenotes}
\sphinxatlongtablestart
\sphinxthistablewithglobalstyle
\makeatletter
  \LTleft \@totalleftmargin plus1fill
  \LTright\dimexpr\columnwidth-\@totalleftmargin-\linewidth\relax plus1fill
\makeatother
\begin{longtable}{lll}
\sphinxtoprule
\endfirsthead

\multicolumn{3}{c}{\sphinxnorowcolor
    \makebox[0pt]{\sphinxtablecontinued{\tablename\ \thetable{} \textendash{} continued from previous page}}%
}\\
\sphinxtoprule
\endhead

\sphinxbottomrule
\multicolumn{3}{r}{\sphinxnorowcolor
    \makebox[0pt][r]{\sphinxtablecontinued{continues on next page}}%
}\\
\endfoot

\endlastfoot
\sphinxtableatstartofbodyhook

\sphinxAtStartPar
{\hyperref[\detokenize{pair_born::doc}]{\sphinxcrossref{\DUrole{doc}{born}}}}
&
\sphinxAtStartPar
a,b,c
&
\sphinxAtStartPar
type pairs
\\
\sphinxhline
\sphinxAtStartPar
{\hyperref[\detokenize{pair_born_gauss::doc}]{\sphinxcrossref{\DUrole{doc}{born/gauss}}}}
&
\sphinxAtStartPar
biga0,biga1,r0
&
\sphinxAtStartPar
type pairs
\\
\sphinxhline
\sphinxAtStartPar
{\hyperref[\detokenize{pair_buck::doc}]{\sphinxcrossref{\DUrole{doc}{buck, buck/coul/cut, buck/coul/long, buck/coul/msm}}}}
&
\sphinxAtStartPar
a,c
&
\sphinxAtStartPar
type pairs
\\
\sphinxhline
\sphinxAtStartPar
{\hyperref[\detokenize{pair_mdf::doc}]{\sphinxcrossref{\DUrole{doc}{buck/mdf}}}}
&
\sphinxAtStartPar
a,c
&
\sphinxAtStartPar
type pairs
\\
\sphinxhline
\sphinxAtStartPar
{\hyperref[\detokenize{pair_coul::doc}]{\sphinxcrossref{\DUrole{doc}{coul/cut, coul/cut/global}}}}
&
\sphinxAtStartPar
scale
&
\sphinxAtStartPar
type pairs
\\
\sphinxhline
\sphinxAtStartPar
{\hyperref[\detokenize{pair_fep_soft::doc}]{\sphinxcrossref{\DUrole{doc}{coul/cut/soft}}}}
&
\sphinxAtStartPar
lambda
&
\sphinxAtStartPar
type pairs
\\
\sphinxhline
\sphinxAtStartPar
{\hyperref[\detokenize{pair_coul::doc}]{\sphinxcrossref{\DUrole{doc}{coul/debye}}}}
&
\sphinxAtStartPar
scale
&
\sphinxAtStartPar
type pairs
\\
\sphinxhline
\sphinxAtStartPar
{\hyperref[\detokenize{pair_coul::doc}]{\sphinxcrossref{\DUrole{doc}{coul/long, coul/msm}}}}
&
\sphinxAtStartPar
scale
&
\sphinxAtStartPar
type pairs
\\
\sphinxhline
\sphinxAtStartPar
{\hyperref[\detokenize{pair_fep_soft::doc}]{\sphinxcrossref{\DUrole{doc}{coul/long/soft}}}}
&
\sphinxAtStartPar
scale, lambda
&
\sphinxAtStartPar
type pairs
\\
\sphinxhline
\sphinxAtStartPar
{\hyperref[\detokenize{pair_coul_slater::doc}]{\sphinxcrossref{\DUrole{doc}{coul/slater/long}}}}
&
\sphinxAtStartPar
scale
&
\sphinxAtStartPar
type pairs
\\
\sphinxhline
\sphinxAtStartPar
{\hyperref[\detokenize{pair_coul::doc}]{\sphinxcrossref{\DUrole{doc}{coul/streitz}}}}
&
\sphinxAtStartPar
scale
&
\sphinxAtStartPar
type pairs
\\
\sphinxhline
\sphinxAtStartPar
{\hyperref[\detokenize{pair_eam::doc}]{\sphinxcrossref{\DUrole{doc}{eam, eam/alloy, eam/fs}}}}
&
\sphinxAtStartPar
scale
&
\sphinxAtStartPar
type pairs
\\
\sphinxhline
\sphinxAtStartPar
{\hyperref[\detokenize{pair_harmonic_cut::doc}]{\sphinxcrossref{\DUrole{doc}{harmonic/cut}}}}
&
\sphinxAtStartPar
k
&
\sphinxAtStartPar
type pairs
\\
\sphinxhline
\sphinxAtStartPar
{\hyperref[\detokenize{pair_gauss::doc}]{\sphinxcrossref{\DUrole{doc}{gauss}}}}
&
\sphinxAtStartPar
a
&
\sphinxAtStartPar
type pairs
\\
\sphinxhline
\sphinxAtStartPar
{\hyperref[\detokenize{pair_mdf::doc}]{\sphinxcrossref{\DUrole{doc}{lennard/mdf}}}}
&
\sphinxAtStartPar
a,b
&
\sphinxAtStartPar
type pairs
\\
\sphinxhline
\sphinxAtStartPar
{\hyperref[\detokenize{pair_class2::doc}]{\sphinxcrossref{\DUrole{doc}{lj/class2}}}}
&
\sphinxAtStartPar
epsilon,sigma
&
\sphinxAtStartPar
type pairs
\\
\sphinxhline
\sphinxAtStartPar
{\hyperref[\detokenize{pair_class2::doc}]{\sphinxcrossref{\DUrole{doc}{lj/class2/coul/cut, lj/class2/coul/long}}}}
&
\sphinxAtStartPar
epsilon,sigma
&
\sphinxAtStartPar
type pairs
\\
\sphinxhline
\sphinxAtStartPar
{\hyperref[\detokenize{pair_lj::doc}]{\sphinxcrossref{\DUrole{doc}{lj/cut}}}}
&
\sphinxAtStartPar
epsilon,sigma
&
\sphinxAtStartPar
type pairs
\\
\sphinxhline
\sphinxAtStartPar
{\hyperref[\detokenize{pair_fep_soft::doc}]{\sphinxcrossref{\DUrole{doc}{lj/cut/soft}}}}
&
\sphinxAtStartPar
epsilon,sigma,lambda
&
\sphinxAtStartPar
type pairs
\\
\sphinxhline
\sphinxAtStartPar
{\hyperref[\detokenize{pair_lj_cut_coul::doc}]{\sphinxcrossref{\DUrole{doc}{lj/cut/coul/cut, lj/cut/coul/long, lj/cut/coul/msm}}}}
&
\sphinxAtStartPar
epsilon,sigma
&
\sphinxAtStartPar
type pairs
\\
\sphinxhline
\sphinxAtStartPar
{\hyperref[\detokenize{pair_fep_soft::doc}]{\sphinxcrossref{\DUrole{doc}{lj/cut/coul/cut/soft, lj/cut/coul/long/soft}}}}
&
\sphinxAtStartPar
epsilon,sigma,lambda
&
\sphinxAtStartPar
type pairs
\\
\sphinxhline
\sphinxAtStartPar
{\hyperref[\detokenize{pair_lj_cut_tip4p::doc}]{\sphinxcrossref{\DUrole{doc}{lj/cut/tip4p/cut, lj/cut/tip4p/long}}}}
&
\sphinxAtStartPar
epsilon,sigma
&
\sphinxAtStartPar
type pairs
\\
\sphinxhline
\sphinxAtStartPar
{\hyperref[\detokenize{pair_fep_soft::doc}]{\sphinxcrossref{\DUrole{doc}{lj/cut/tip4p/long/soft}}}}
&
\sphinxAtStartPar
epsilon,sigma,lambda
&
\sphinxAtStartPar
type pairs
\\
\sphinxhline
\sphinxAtStartPar
{\hyperref[\detokenize{pair_lj_expand::doc}]{\sphinxcrossref{\DUrole{doc}{lj/expand}}}}
&
\sphinxAtStartPar
epsilon,sigma,delta
&
\sphinxAtStartPar
type pairs
\\
\sphinxhline
\sphinxAtStartPar
{\hyperref[\detokenize{pair_mdf::doc}]{\sphinxcrossref{\DUrole{doc}{lj/mdf}}}}
&
\sphinxAtStartPar
epsilon,sigma
&
\sphinxAtStartPar
type pairs
\\
\sphinxhline
\sphinxAtStartPar
{\hyperref[\detokenize{pair_dipole::doc}]{\sphinxcrossref{\DUrole{doc}{lj/sf/dipole/sf}}}}
&
\sphinxAtStartPar
epsilon,sigma,scale
&
\sphinxAtStartPar
type pairs
\\
\sphinxhline
\sphinxAtStartPar
{\hyperref[\detokenize{pair_meam::doc}]{\sphinxcrossref{\DUrole{doc}{meam}}}}
&
\sphinxAtStartPar
scale
&
\sphinxAtStartPar
type pairs
\\
\sphinxhline
\sphinxAtStartPar
{\hyperref[\detokenize{pair_mie::doc}]{\sphinxcrossref{\DUrole{doc}{mie/cut}}}}
&
\sphinxAtStartPar
epsilon,sigma,gamR,gamA
&
\sphinxAtStartPar
type pairs
\\
\sphinxhline
\sphinxAtStartPar
{\hyperref[\detokenize{pair_morse::doc}]{\sphinxcrossref{\DUrole{doc}{morse, morse/smooth/linear}}}}
&
\sphinxAtStartPar
d0,r0,alpha
&
\sphinxAtStartPar
type pairs
\\
\sphinxhline
\sphinxAtStartPar
{\hyperref[\detokenize{pair_morse::doc}]{\sphinxcrossref{\DUrole{doc}{morse/soft}}}}
&
\sphinxAtStartPar
d0,r0,alpha,lambda
&
\sphinxAtStartPar
type pairs
\\
\sphinxhline
\sphinxAtStartPar
{\hyperref[\detokenize{pair_nm::doc}]{\sphinxcrossref{\DUrole{doc}{nm/cut}}}}
&
\sphinxAtStartPar
e0,r0,nn,mm
&
\sphinxAtStartPar
type pairs
\\
\sphinxhline
\sphinxAtStartPar
{\hyperref[\detokenize{pair_nm::doc}]{\sphinxcrossref{\DUrole{doc}{nm/cut/coul/cut, nm/cut/coul/long}}}}
&
\sphinxAtStartPar
e0,r0,nn,mm
&
\sphinxAtStartPar
type pairs
\\
\sphinxhline
\sphinxAtStartPar
{\hyperref[\detokenize{pair_pace::doc}]{\sphinxcrossref{\DUrole{doc}{pace, pace/extrapolation}}}}
&
\sphinxAtStartPar
scale
&
\sphinxAtStartPar
type pairs
\\
\sphinxhline
\sphinxAtStartPar
{\hyperref[\detokenize{pair_snap::doc}]{\sphinxcrossref{\DUrole{doc}{snap}}}}
&
\sphinxAtStartPar
scale
&
\sphinxAtStartPar
type pairs
\\
\sphinxhline
\sphinxAtStartPar
{\hyperref[\detokenize{pair_soft::doc}]{\sphinxcrossref{\DUrole{doc}{soft}}}}
&
\sphinxAtStartPar
a
&
\sphinxAtStartPar
type pairs
\\
\sphinxhline
\sphinxAtStartPar
{\hyperref[\detokenize{pair_ufm::doc}]{\sphinxcrossref{\DUrole{doc}{ufm}}}}
&
\sphinxAtStartPar
epsilon,sigma,scale
&
\sphinxAtStartPar
type pairs
\\
\sphinxhline
\sphinxAtStartPar
{\hyperref[\detokenize{pair_wf_cut::doc}]{\sphinxcrossref{\DUrole{doc}{wf/cut}}}}
&
\sphinxAtStartPar
epsilon,sigma,nu,mu
&
\sphinxAtStartPar
type pairs
\\
\sphinxbottomrule
\end{longtable}
\sphinxtableafterendhook
\sphinxatlongtableend
\end{savenotes}

\begin{sphinxadmonition}{note}{Note:}
\sphinxAtStartPar
It is easy to add new potentials and their parameters to this
list.  All it typically takes is adding an extract() method to the
pair\_*.cpp file associated with the potential.
\end{sphinxadmonition}

\sphinxAtStartPar
Note that for many of the potentials, the parameter that can be varied
is effectively a prefactor on the entire energy expression for the
potential (e.g., the lj/cut epsilon).  The parameters listed as “scale”
are exactly that, since the energy expression for the
{\hyperref[\detokenize{pair_coul::doc}]{\sphinxcrossref{\DUrole{doc}{coul/cut}}}} potential (for example) has no labeled
prefactor in its formula.  To apply an effective prefactor to some
potentials, multiple parameters need to be altered.  For example, the
{\hyperref[\detokenize{pair_buck::doc}]{\sphinxcrossref{\DUrole{doc}{Buckingham potential}}}} needs both the A and C terms
altered together.  To scale the Buckingham potential, you should thus
list the pair style twice, once for A and once for C.

\sphinxAtStartPar
If a type pair parameter is specified, the \sphinxstyleemphasis{I} and \sphinxstyleemphasis{J} settings should
be specified to indicate which type pairs to apply it to.  If a global
parameter is specified, the \sphinxstyleemphasis{I} and \sphinxstyleemphasis{J} settings still need to be
specified, but are ignored.

\sphinxAtStartPar
Similar to the {\hyperref[\detokenize{pair_coeff::doc}]{\sphinxcrossref{\DUrole{doc}{pair\_coeff command}}}}, \(I\) and
\(J\) can be specified in one of several ways.  Explicit numeric values
can be used for each, as in the first example above.  Or, one or both of
the types in the I,J pair can be a {\hyperref[\detokenize{Howto_type_labels::doc}]{\sphinxcrossref{\DUrole{doc}{type label}}}}.
LAMMPS sets the coefficients for the symmetric \(J,I\) interaction to
the same values.

\sphinxAtStartPar
A wild\sphinxhyphen{}card asterisk can be used in place of or in conjunction with
the \(I,J\) arguments to set the coefficients for multiple pairs of atom
types.  This takes the form “*” or “*n” or “m*” or “m*n”.  If \(N\) is
the number of atom types, then an asterisk with no numeric values means
all types from 1 to \(N\).  A leading asterisk means all types from 1 to n
(inclusive).  A trailing asterisk means all types from m to \(N\)
(inclusive).  A middle asterisk means all types from m to n
(inclusive).  For the asterisk syntax, note that only type pairs with
\(I \le J\) are considered; if asterisks imply type pairs where
\(J < I\), they are ignored.

\sphinxAtStartPar
IMPROTANT NOTE: If {\hyperref[\detokenize{pair_hybrid::doc}]{\sphinxcrossref{\DUrole{doc}{pair\_style hybrid or hybrid/overlay}}}} is
being used, then the \sphinxstyleemphasis{pstyle} will be a sub\sphinxhyphen{}style name.  You must specify
\(I,J\) arguments that correspond to type pair values defined (via the
{\hyperref[\detokenize{pair_coeff::doc}]{\sphinxcrossref{\DUrole{doc}{pair\_coeff}}}} command) for that sub\sphinxhyphen{}style.

\sphinxAtStartPar
The \sphinxstyleemphasis{v\_name} argument for keyword \sphinxstyleemphasis{pair} is the name of an
{\hyperref[\detokenize{variable::doc}]{\sphinxcrossref{\DUrole{doc}{equal\sphinxhyphen{}style variable}}}} which will be evaluated each time
this fix is invoked to set the parameter to a new value.  It should be
specified as v\_name, where name is the variable name.  Equal\sphinxhyphen{}style
variables can specify formulas with various mathematical functions,
and include {\hyperref[\detokenize{thermo_style::doc}]{\sphinxcrossref{\DUrole{doc}{thermo\_style}}}} command keywords for the
simulation box parameters and timestep and elapsed time.  Thus it is
easy to specify parameters that change as a function of time or span
consecutive runs in a continuous fashion.  For the latter, see the
\sphinxstyleemphasis{start} and \sphinxstyleemphasis{stop} keywords of the {\hyperref[\detokenize{run::doc}]{\sphinxcrossref{\DUrole{doc}{run}}}} command and the
\sphinxstyleemphasis{elaplong} keyword of {\hyperref[\detokenize{thermo_style::doc}]{\sphinxcrossref{\DUrole{doc}{thermo\_style custom}}}} for details.

\sphinxAtStartPar
For example, these commands would change the prefactor coefficient of
the {\hyperref[\detokenize{pair_soft::doc}]{\sphinxcrossref{\DUrole{doc}{pair\_style soft}}}} potential from 10.0 to 30.0 in a
linear fashion over the course of a simulation:

\begin{sphinxVerbatim}[commandchars=\\\{\}]
\PYG{k}{variable }\PYG{n+nv+nvIdentifier}{prefactor}\PYG{+w}{ }\PYG{n}{equal}\PYG{+w}{ }\PYG{n}{ramp}\PYG{n+nv}{(}\PYG{n+nv}{10,30}\PYG{n+nv}{)}
\PYG{k}{fix }\PYG{n+nv+nvIdentifier}{1}\PYG{+w}{ }\PYG{n+nv+nvIdentifier}{all}\PYG{+w}{ }\PYG{n}{adapt}\PYG{+w}{ }\PYG{l+m}{1}\PYG{+w}{ }\PYG{n}{pair}\PYG{+w}{ }\PYG{n}{soft}\PYG{+w}{ }\PYG{n}{a}\PYG{+w}{ }\PYG{o}{*}\PYG{+w}{ }\PYG{o}{*}\PYG{+w}{ }\PYG{n}{v\PYGZus{}prefactor}
\end{sphinxVerbatim}


\bigskip\hrule\bigskip


\sphinxAtStartPar
The \sphinxstyleemphasis{kspace} keyword used the specified variable as a scale factor on
the energy, forces, virial calculated by whatever \(k\)\sphinxhyphen{}space solver is
defined by the {\hyperref[\detokenize{kspace_style::doc}]{\sphinxcrossref{\DUrole{doc}{kspace\_style}}}} command.  If the
variable has a value of 1.0, then the solver is unaltered.

\sphinxAtStartPar
The \sphinxstyleemphasis{kspace} keyword works this way whether the \sphinxstyleemphasis{scale} keyword
is set to \sphinxstyleemphasis{no} or \sphinxstyleemphasis{yes}.


\bigskip\hrule\bigskip


\sphinxAtStartPar
The \sphinxstyleemphasis{atom} keyword enables various atom properties to be changed.  The
\sphinxstyleemphasis{aparam} argument is the name of the parameter to change.  This is the
current list of atom parameters that can be varied by this fix:
\begin{itemize}
\item {} 
\sphinxAtStartPar
charge = charge on particle

\item {} 
\sphinxAtStartPar
diameter = diameter of particle

\end{itemize}

\sphinxAtStartPar
The \(I\) argument indicates which atom types are affected. A wild\sphinxhyphen{}card
asterisk can be used in place of or in conjunction with the \(I\) argument
to set the coefficients for multiple atom types.

\sphinxAtStartPar
The \sphinxstyleemphasis{v\_name} argument of the \sphinxstyleemphasis{atom} keyword is the name of an
{\hyperref[\detokenize{variable::doc}]{\sphinxcrossref{\DUrole{doc}{equal\sphinxhyphen{}style variable}}}} which will be evaluated each time
this fix is invoked to set the parameter to a new value.  It should be
specified as v\_name, where name is the variable name.  See the
discussion above describing the formulas associated with equal\sphinxhyphen{}style
variables.  The new value is assigned to the corresponding attribute
for all atoms in the fix group.

\sphinxAtStartPar
If the atom parameter is \sphinxstyleemphasis{diameter} and per\sphinxhyphen{}atom density and per\sphinxhyphen{}atom
mass are defined for particles (e.g., {\hyperref[\detokenize{atom_style::doc}]{\sphinxcrossref{\DUrole{doc}{atom\_style granular}}}}),
then the mass of each particle is also changed when the diameter changes
(density is assumed to stay constant).

\sphinxAtStartPar
For example, these commands would shrink the diameter of all granular
particles in the “center” group from 1.0 to 0.1 in a linear fashion
over the course of a 1000\sphinxhyphen{}step simulation:

\begin{sphinxVerbatim}[commandchars=\\\{\}]
\PYG{k}{variable }\PYG{n+nv+nvIdentifier}{size}\PYG{+w}{ }\PYG{n}{equal}\PYG{+w}{ }\PYG{n}{ramp}\PYG{n+nv}{(}\PYG{n+nv}{1.0,0.1}\PYG{n+nv}{)}
\PYG{k}{fix }\PYG{n+nv+nvIdentifier}{1}\PYG{+w}{ }\PYG{n+nv+nvIdentifier}{center}\PYG{+w}{ }\PYG{n}{adapt}\PYG{+w}{ }\PYG{l+m}{10}\PYG{+w}{ }\PYG{n}{atom}\PYG{+w}{ }\PYG{n}{diameter}\PYG{+w}{ }\PYG{o}{*}\PYG{+w}{ }\PYG{n}{v\PYGZus{}size}
\end{sphinxVerbatim}

\sphinxAtStartPar
For {\hyperref[\detokenize{run_style::doc}]{\sphinxcrossref{\DUrole{doc}{rRESPA time integration}}}}, this fix changes
parameters on the outermost rRESPA level.


\bigskip\hrule\bigskip



\subsection{Restart, fix\_modify, output, run start/stop, minimize info}
\label{\detokenize{fix_adapt_fep:restart-fix-modify-output-run-start-stop-minimize-info}}
\sphinxAtStartPar
No information about this fix is written to
{\hyperref[\detokenize{restart::doc}]{\sphinxcrossref{\DUrole{doc}{binary restart files}}}}.
None of the {\hyperref[\detokenize{fix_modify::doc}]{\sphinxcrossref{\DUrole{doc}{fix\_modify}}}} options
are relevant to this fix.  No global or per\sphinxhyphen{}atom quantities are stored
by this fix for access by various {\hyperref[\detokenize{Howto_output::doc}]{\sphinxcrossref{\DUrole{doc}{output commands}}}}.
No parameter of this fix can be used with the \sphinxstyleemphasis{start/stop} keywords of
the {\hyperref[\detokenize{run::doc}]{\sphinxcrossref{\DUrole{doc}{run}}}} command.  This fix is not invoked during
{\hyperref[\detokenize{minimize::doc}]{\sphinxcrossref{\DUrole{doc}{energy minimization}}}}.


\subsection{Restrictions}
\label{\detokenize{fix_adapt_fep:restrictions}}
\sphinxAtStartPar
The keyword “scale yes” is not supported for scaling per\sphinxhyphen{}atom parameters
diameter and change. You can use {\hyperref[\detokenize{fix_adapt::doc}]{\sphinxcrossref{\DUrole{doc}{fix adapt}}}} for those.


\subsection{Related commands}
\label{\detokenize{fix_adapt_fep:related-commands}}
\sphinxAtStartPar
{\hyperref[\detokenize{compute_fep::doc}]{\sphinxcrossref{\DUrole{doc}{compute fep}}}}, {\hyperref[\detokenize{fix_adapt::doc}]{\sphinxcrossref{\DUrole{doc}{fix adapt}}}},
{\hyperref[\detokenize{compute_ti::doc}]{\sphinxcrossref{\DUrole{doc}{compute ti}}}}, {\hyperref[\detokenize{pair_fep_soft::doc}]{\sphinxcrossref{\DUrole{doc}{pair\_style */soft}}}}


\subsection{Default}
\label{\detokenize{fix_adapt_fep:default}}
\sphinxAtStartPar
The option defaults are scale = no, reset = no, after = no.

\sphinxstepscope

\index{fix add/heat@\spxentry{fix add/heat}}\ignorespaces 

\section{fix add/heat command}
\label{\detokenize{fix_add_heat:fix-add-heat-command}}\label{\detokenize{fix_add_heat:index-0}}\label{\detokenize{fix_add_heat::doc}}

\subsection{Syntax}
\label{\detokenize{fix_add_heat:syntax}}
\begin{sphinxVerbatim}[commandchars=\\\{\}]
\PYG{k}{fix }\PYG{n+nv+nvIdentifier}{ID}\PYG{+w}{ }\PYG{n+nv+nvIdentifier}{group\PYGZhy{}ID}\PYG{+w}{ }\PYG{n}{add}\PYG{o}{/}\PYG{n}{heat}\PYG{+w}{ }\PYG{n}{style}\PYG{+w}{ }\PYG{n}{args}\PYG{+w}{ }\PYG{n}{keyword}\PYG{+w}{ }\PYG{n}{values}\PYG{+w}{ }\PYG{n}{...}
\end{sphinxVerbatim}
\begin{itemize}
\item {} 
\sphinxAtStartPar
ID, group\sphinxhyphen{}ID are documented in {\hyperref[\detokenize{fix::doc}]{\sphinxcrossref{\DUrole{doc}{fix}}}} command

\item {} 
\sphinxAtStartPar
add/heat = style name of this fix command

\item {} 
\sphinxAtStartPar
style = \sphinxstyleemphasis{constant} or \sphinxstyleemphasis{linear} or \sphinxstyleemphasis{quartic}
\begin{sphinxalltt}
\sphinxstyleemphasis{constant} args = \sphinxstyleemphasis{rate}
  \sphinxstyleemphasis{rate} = rate of heat flow (energy/time units)
\sphinxstyleemphasis{linear} args = \sphinxstyleemphasis{Ttarget} \sphinxstyleemphasis{k}
  \sphinxstyleemphasis{Ttarget} = target temperature (temperature units)
  \sphinxstyleemphasis{k} = prefactor (energy/(time*temperature) units)
\sphinxstyleemphasis{quartic} args = \sphinxstyleemphasis{Ttarget} \sphinxstyleemphasis{k}
  \sphinxstyleemphasis{Ttarget} = target temperature (temperature units)
  \sphinxstyleemphasis{k} = prefactor (energy/(time*temperature\textasciicircum{}4) units)
\end{sphinxalltt}

\item {} 
\sphinxAtStartPar
zero or more keyword/value pairs may be appended to args

\item {} 
\sphinxAtStartPar
keyword = \sphinxstyleemphasis{overwrite}
\begin{sphinxalltt}
\sphinxstyleemphasis{overwrite} value = \sphinxstyleemphasis{yes} or \sphinxstyleemphasis{no}
  \sphinxstyleemphasis{yes} = sets current heat flow of particle
  \sphinxstyleemphasis{no} = adds to current heat flow of particle
\end{sphinxalltt}

\end{itemize}


\subsection{Examples}
\label{\detokenize{fix_add_heat:examples}}
\begin{sphinxVerbatim}[commandchars=\\\{\}]
\PYG{k}{fix }\PYG{n+nv+nvIdentifier}{1}\PYG{+w}{ }\PYG{n+nv+nvIdentifier}{all}\PYG{+w}{ }\PYG{n}{add}\PYG{o}{/}\PYG{n}{heat}\PYG{+w}{ }\PYG{n}{constant}\PYG{+w}{ }\PYG{n}{v\PYGZus{}heat}
\PYG{k}{fix }\PYG{n+nv+nvIdentifier}{1}\PYG{+w}{ }\PYG{n+nv+nvIdentifier}{all}\PYG{+w}{ }\PYG{n}{add}\PYG{o}{/}\PYG{n}{heat}\PYG{+w}{ }\PYG{n}{linear}\PYG{+w}{ }\PYG{l+m}{10.0}\PYG{+w}{ }\PYG{l+m}{1.0}\PYG{+w}{ }\PYG{n}{overwrite}\PYG{+w}{ }\PYG{n}{yes}
\end{sphinxVerbatim}


\subsection{Description}
\label{\detokenize{fix_add_heat:description}}
\sphinxAtStartPar
This fix adds heat to particles with the temperature attribute every timestep
at a given rate. Note that this is an internal temperature of a particle intended
for use with non\sphinxhyphen{}atomistic models like the discrete element method.

\sphinxAtStartPar
For the \sphinxstyleemphasis{constant} style, heat is added at the specified rate. For the \sphinxstyleemphasis{linear} style,
heat is added at a rate of \(k (T_{target} - T)\) where \(k\) is the
specified prefactor, \(T_{target}\) is the specified target temperature, and
\(T\) is the temperature of the atom. This may be more representative of a
conductive process. For the \sphinxstyleemphasis{quartic} style, heat is added at a rate of
\(k (T_{target}^4 - T^4)\), akin to radiative heat transfer.

\sphinxAtStartPar
The rate or temperature can be can be specified as an equal\sphinxhyphen{}style or atom\sphinxhyphen{}style
{\hyperref[\detokenize{variable::doc}]{\sphinxcrossref{\DUrole{doc}{variable}}}}.  If the value is a variable, it should be
specified as v\_name, where name is the variable name.  In this case, the
variable will be evaluated each time step, and its value will be used to
determine the rate of heat added.

\sphinxAtStartPar
Equal\sphinxhyphen{}style variables can specify formulas with various mathematical
functions and include {\hyperref[\detokenize{thermo_style::doc}]{\sphinxcrossref{\DUrole{doc}{thermo\_style}}}} command
keywords for the simulation box parameters, time step, and elapsed time
to specify time\sphinxhyphen{}dependent heating.

\sphinxAtStartPar
Atom\sphinxhyphen{}style variables can specify the same formulas as equal\sphinxhyphen{}style
variables but can also include per\sphinxhyphen{}atom values, such as atom
coordinates to specify spatially\sphinxhyphen{}dependent heating.

\sphinxAtStartPar
If the \sphinxstyleemphasis{overwrite} keyword is set to \sphinxstyleemphasis{yes}, this fix will set the total
heat flow on a particle every timestep, overwriting contributions from pair
styles or other fixes. If \sphinxstyleemphasis{overwrite} is \sphinxstyleemphasis{no}, this fix will add heat on
top of other contributions.


\bigskip\hrule\bigskip



\subsection{Restart, fix\_modify, output, run start/stop, minimize info}
\label{\detokenize{fix_add_heat:restart-fix-modify-output-run-start-stop-minimize-info}}
\sphinxAtStartPar
No information about this fix is written to {\hyperref[\detokenize{restart::doc}]{\sphinxcrossref{\DUrole{doc}{binary restart files}}}}.
None of the {\hyperref[\detokenize{fix_modify::doc}]{\sphinxcrossref{\DUrole{doc}{fix\_modify}}}} options are relevant to this fix.
No global or per\sphinxhyphen{}atom quantities are stored by this fix for access by various
{\hyperref[\detokenize{Howto_output::doc}]{\sphinxcrossref{\DUrole{doc}{output commands}}}}. No parameter of this fix can be used
with the \sphinxstyleemphasis{start/stop} keywords of the {\hyperref[\detokenize{run::doc}]{\sphinxcrossref{\DUrole{doc}{run}}}} command.  This fix is
not invoked during {\hyperref[\detokenize{minimize::doc}]{\sphinxcrossref{\DUrole{doc}{energy minimization}}}}.


\subsection{Restrictions}
\label{\detokenize{fix_add_heat:restrictions}}
\sphinxAtStartPar
This pair style is part of the GRANULAR package.  It is
only enabled if LAMMPS was built with that package.
See the {\hyperref[\detokenize{Build_package::doc}]{\sphinxcrossref{\DUrole{doc}{Build package}}}} page for more info.

\sphinxAtStartPar
This fix requires that atoms store temperature and heat flow
as defined by the {\hyperref[\detokenize{fix_property_atom::doc}]{\sphinxcrossref{\DUrole{doc}{fix property/atom}}}} command or
included in certain atom styles, such as atom\_style rheo/thermal.


\subsection{Related commands}
\label{\detokenize{fix_add_heat:related-commands}}
\sphinxAtStartPar
{\hyperref[\detokenize{fix_heat_flow::doc}]{\sphinxcrossref{\DUrole{doc}{fix heat/flow}}}},
{\hyperref[\detokenize{fix_property_atom::doc}]{\sphinxcrossref{\DUrole{doc}{fix property/atom}}}},
{\hyperref[\detokenize{fix_rheo_thermal::doc}]{\sphinxcrossref{\DUrole{doc}{fix rheo/thermal}}}}


\subsection{Default}
\label{\detokenize{fix_add_heat:default}}
\sphinxAtStartPar
The default for the \sphinxstyleemphasis{overwrite} keyword is \sphinxstyleemphasis{no}

\sphinxstepscope

\index{fix addforce@\spxentry{fix addforce}}\ignorespaces 

\section{fix addforce command}
\label{\detokenize{fix_addforce:fix-addforce-command}}\label{\detokenize{fix_addforce:index-0}}\label{\detokenize{fix_addforce::doc}}

\subsection{Syntax}
\label{\detokenize{fix_addforce:syntax}}
\begin{sphinxVerbatim}[commandchars=\\\{\}]
\PYG{k}{fix }\PYG{n+nv+nvIdentifier}{ID}\PYG{+w}{ }\PYG{n+nv+nvIdentifier}{group\PYGZhy{}ID}\PYG{+w}{ }\PYG{n}{addforce}\PYG{+w}{ }\PYG{n}{fx}\PYG{+w}{ }\PYG{n}{fy}\PYG{+w}{ }\PYG{n}{fz}\PYG{+w}{ }\PYG{n}{keyword}\PYG{+w}{ }\PYG{n}{value}\PYG{+w}{ }\PYG{n}{...}
\end{sphinxVerbatim}
\begin{itemize}
\item {} 
\sphinxAtStartPar
ID, group\sphinxhyphen{}ID are documented in {\hyperref[\detokenize{fix::doc}]{\sphinxcrossref{\DUrole{doc}{fix}}}} command

\item {} 
\sphinxAtStartPar
addforce = style name of this fix command

\item {} 
\sphinxAtStartPar
fx,fy,fz = force component values (force units)

\begin{sphinxVerbatim}[commandchars=\\\{\}]
any of fx,fy,fz can be a variable (see below)
\end{sphinxVerbatim}

\item {} 
\sphinxAtStartPar
zero or more keyword/value pairs may be appended to args

\item {} 
\sphinxAtStartPar
keyword = \sphinxstyleemphasis{every} or \sphinxstyleemphasis{region} or \sphinxstyleemphasis{energy}
\begin{sphinxalltt}
\sphinxstyleemphasis{every} value = Nevery
  Nevery = add force every this many time steps
\sphinxstyleemphasis{region} value = region\sphinxhyphen{}ID
  region\sphinxhyphen{}ID = ID of region atoms must be in to have added force
\sphinxstyleemphasis{energy} value = v\_name
  v\_name = variable with name that calculates the potential energy of each atom in the added force field
\end{sphinxalltt}

\end{itemize}


\subsection{Examples}
\label{\detokenize{fix_addforce:examples}}
\begin{sphinxVerbatim}[commandchars=\\\{\}]
\PYG{k}{fix }\PYG{n+nv+nvIdentifier}{kick}\PYG{+w}{ }\PYG{n+nv+nvIdentifier}{flow}\PYG{+w}{ }\PYG{n}{addforce}\PYG{+w}{ }\PYG{l+m}{1.0}\PYG{+w}{ }\PYG{l+m}{0.0}\PYG{+w}{ }\PYG{l+m}{0.0}
\PYG{k}{fix }\PYG{n+nv+nvIdentifier}{kick}\PYG{+w}{ }\PYG{n+nv+nvIdentifier}{flow}\PYG{+w}{ }\PYG{n}{addforce}\PYG{+w}{ }\PYG{l+m}{1.0}\PYG{+w}{ }\PYG{l+m}{0.0}\PYG{+w}{ }\PYG{n}{v\PYGZus{}oscillate}
\PYG{k}{fix }\PYG{n+nv+nvIdentifier}{ff}\PYG{+w}{ }\PYG{n+nv+nvIdentifier}{boundary}\PYG{+w}{ }\PYG{n}{addforce}\PYG{+w}{ }\PYG{l+m}{0.0}\PYG{+w}{ }\PYG{l+m}{0.0}\PYG{+w}{ }\PYG{n}{v\PYGZus{}push}\PYG{+w}{ }\PYG{n}{energy}\PYG{+w}{ }\PYG{n}{v\PYGZus{}espace}
\end{sphinxVerbatim}


\subsection{Description}
\label{\detokenize{fix_addforce:description}}
\sphinxAtStartPar
Add \((f_x,f_y,f_z)\) to the corresponding component of the force for each
atom in the group.  This command can be used to give an additional push to
atoms in a simulation, such as for a simulation of Poiseuille flow in
a channel.

\sphinxAtStartPar
Any of the three quantities defining the force components, namely \(f_x\),
\(f_y\), and \(f_z\), can be specified as an equal\sphinxhyphen{}style or atom\sphinxhyphen{}style
{\hyperref[\detokenize{variable::doc}]{\sphinxcrossref{\DUrole{doc}{variable}}}}.  If the value is a variable, it should be specified
as v\_name, where name is the variable name.  In this case, the variable
will be evaluated each time step, and its value(s) will be used to determine
the force component(s).

\sphinxAtStartPar
Equal\sphinxhyphen{}style variables can specify formulas with various mathematical
functions and include {\hyperref[\detokenize{thermo_style::doc}]{\sphinxcrossref{\DUrole{doc}{thermo\_style}}}} command
keywords for the simulation box parameters, time step, and elapsed time.
Thus, it is easy to specify a time\sphinxhyphen{}dependent force field.

\sphinxAtStartPar
Atom\sphinxhyphen{}style variables can specify the same formulas as equal\sphinxhyphen{}style
variables but can also include per\sphinxhyphen{}atom values, such as atom
coordinates.  Thus, it is easy to specify a spatially\sphinxhyphen{}dependent force
field with optional time\sphinxhyphen{}dependence as well.

\sphinxAtStartPar
If the \sphinxstyleemphasis{every} keyword is used, the \sphinxstyleemphasis{Nevery} setting determines how
often the forces are applied.  The default value is 1, for every
time step.

\sphinxAtStartPar
If the \sphinxstyleemphasis{region} keyword is used, the atom must also be in the
specified geometric {\hyperref[\detokenize{region::doc}]{\sphinxcrossref{\DUrole{doc}{region}}}} in order to have force added
to it.


\bigskip\hrule\bigskip


\sphinxAtStartPar
Adding a force to atoms implies a change in their potential energy as
they move due to the applied force field.  For dynamics via the “run”
command, this energy can be optionally added to the system’s potential
energy for thermodynamic output (see below).  For energy minimization
via the “minimize” command, this energy must be added to the system’s
potential energy to formulate a self\sphinxhyphen{}consistent minimization problem
(see below).

\sphinxAtStartPar
The \sphinxstyleemphasis{energy} keyword is not allowed if the added force is a constant
vector \(\vec F = (f_x,f_y,f_z)\), with all components defined as numeric
constants and not as variables.  This is because LAMMPS can compute
the energy for each atom directly as
\begin{equation*}
\begin{split}E = -\vec x \cdot \vec F = -(x f_x + y f_y + z f_z),\end{split}
\end{equation*}
\sphinxAtStartPar
so that \(-\vec\nabla E = \vec F\).

\sphinxAtStartPar
The \sphinxstyleemphasis{energy} keyword is optional if the added force is defined with
one or more variables, and if you are performing dynamics via the
{\hyperref[\detokenize{run::doc}]{\sphinxcrossref{\DUrole{doc}{run}}}} command.  If the keyword is not used, LAMMPS will set
the energy to 0.0, which is typically fine for dynamics.

\sphinxAtStartPar
The \sphinxstyleemphasis{energy} keyword is required if the added force is defined with
one or more variables, and you are performing energy minimization via
the “minimize” command.  The keyword specifies the name of an
atom\sphinxhyphen{}style {\hyperref[\detokenize{variable::doc}]{\sphinxcrossref{\DUrole{doc}{variable}}}} which is used to compute the
energy of each atom as function of its position.  Like variables used
for \(f_x\), \(f_y\), \(f_z\), the energy variable is specified as
v\_name, where name is the variable name.

\sphinxAtStartPar
Note that when the \sphinxstyleemphasis{energy} keyword is used during an energy
minimization, you must ensure that the formula defined for the
atom\sphinxhyphen{}style {\hyperref[\detokenize{variable::doc}]{\sphinxcrossref{\DUrole{doc}{variable}}}} is consistent with the force
variable formulas (i.e., that \(-\vec\nabla E = \vec F\)).
For example, if the force were a spring\sphinxhyphen{}like, \(\vec F = -k\vec x\), then
the energy formula should be \(E = \frac12 kx^2\).  If you do not do this
correctly, the minimization will not converge properly.


\bigskip\hrule\bigskip


\sphinxAtStartPar
Styles with a \sphinxstyleemphasis{gpu}, \sphinxstyleemphasis{intel}, \sphinxstyleemphasis{kk}, \sphinxstyleemphasis{omp}, or \sphinxstyleemphasis{opt} suffix are
functionally the same as the corresponding style without the suffix.
They have been optimized to run faster, depending on your available
hardware, as discussed on the {\hyperref[\detokenize{Speed_packages::doc}]{\sphinxcrossref{\DUrole{doc}{Accelerator packages}}}}
page.  The accelerated styles take the same arguments and should
produce the same results, except for round\sphinxhyphen{}off and precision issues.

\sphinxAtStartPar
These accelerated styles are part of the GPU, INTEL, KOKKOS,
OPENMP, and OPT packages, respectively.  They are only enabled if
LAMMPS was built with those packages.  See the {\hyperref[\detokenize{Build_package::doc}]{\sphinxcrossref{\DUrole{doc}{Build package}}}} page for more info.

\sphinxAtStartPar
You can specify the accelerated styles explicitly in your input script
by including their suffix, or you can use the {\hyperref[\detokenize{Run_options::doc}]{\sphinxcrossref{\DUrole{doc}{\sphinxhyphen{}suffix command\sphinxhyphen{}line switch}}}} when you invoke LAMMPS, or you can use the
{\hyperref[\detokenize{suffix::doc}]{\sphinxcrossref{\DUrole{doc}{suffix}}}} command in your input script.

\sphinxAtStartPar
See the {\hyperref[\detokenize{Speed_packages::doc}]{\sphinxcrossref{\DUrole{doc}{Accelerator packages}}}} page for more
instructions on how to use the accelerated styles effectively.


\bigskip\hrule\bigskip



\subsection{Restart, fix\_modify, output, run start/stop, minimize info}
\label{\detokenize{fix_addforce:restart-fix-modify-output-run-start-stop-minimize-info}}
\sphinxAtStartPar
No information about this fix is written to {\hyperref[\detokenize{restart::doc}]{\sphinxcrossref{\DUrole{doc}{binary restart files}}}}.

\sphinxAtStartPar
The {\hyperref[\detokenize{fix_modify::doc}]{\sphinxcrossref{\DUrole{doc}{fix\_modify}}}} \sphinxstyleemphasis{energy} option is supported by
this fix to add the potential energy inferred by the added force to
the global potential energy of the system as part of
{\hyperref[\detokenize{thermo_style::doc}]{\sphinxcrossref{\DUrole{doc}{thermodynamic output}}}}.  The default setting for
this fix is {\hyperref[\detokenize{fix_modify::doc}]{\sphinxcrossref{\DUrole{doc}{fix\_modify energy no}}}}.  Note that this
energy is a fictitious quantity but is needed so that the
{\hyperref[\detokenize{minimize::doc}]{\sphinxcrossref{\DUrole{doc}{minimize}}}} command can include the forces added by
this fix in a consistent manner (i.e., there is a decrease in
potential energy when atoms move in the direction of the added force).

\sphinxAtStartPar
The {\hyperref[\detokenize{fix_modify::doc}]{\sphinxcrossref{\DUrole{doc}{fix\_modify}}}} \sphinxstyleemphasis{virial} option is supported by
this fix to add the contribution due to the added forces on atoms to
both the global pressure and per\sphinxhyphen{}atom stress of the system via the
{\hyperref[\detokenize{compute_pressure::doc}]{\sphinxcrossref{\DUrole{doc}{compute pressure}}}} and {\hyperref[\detokenize{compute_stress_atom::doc}]{\sphinxcrossref{\DUrole{doc}{compute
stress/atom}}}} commands.  The former can be
accessed by {\hyperref[\detokenize{thermo_style::doc}]{\sphinxcrossref{\DUrole{doc}{thermodynamic output}}}}.  The default
setting for this fix is {\hyperref[\detokenize{fix_modify::doc}]{\sphinxcrossref{\DUrole{doc}{fix\_modify virial no}}}}.

\sphinxAtStartPar
The {\hyperref[\detokenize{fix_modify::doc}]{\sphinxcrossref{\DUrole{doc}{fix\_modify}}}} \sphinxstyleemphasis{respa} option is supported by this
fix. This allows to set at which level of the {\hyperref[\detokenize{run_style::doc}]{\sphinxcrossref{\DUrole{doc}{r\sphinxhyphen{}RESPA}}}} integrator the fix is adding its forces. Default is the
outermost level.

\sphinxAtStartPar
This fix computes a global scalar and a global three\sphinxhyphen{}vector of forces,
which can be accessed by various {\hyperref[\detokenize{Howto_output::doc}]{\sphinxcrossref{\DUrole{doc}{output commands}}}}.  The scalar is the potential energy discussed above.
The vector is the total force on the group of atoms before the forces
on individual atoms are changed by the fix.  The scalar and vector
values calculated by this fix are “extensive”.

\sphinxAtStartPar
No parameter of this fix can be used with the \sphinxstyleemphasis{start/stop} keywords of
the {\hyperref[\detokenize{run::doc}]{\sphinxcrossref{\DUrole{doc}{run}}}} command.

\sphinxAtStartPar
The forces due to this fix are imposed during an energy minimization,
invoked by the {\hyperref[\detokenize{minimize::doc}]{\sphinxcrossref{\DUrole{doc}{minimize}}}} command.  You should not
specify force components with a variable that has time\sphinxhyphen{}dependence for
use with a minimizer, since the minimizer increments the time step as
the iteration count during the minimization.

\begin{sphinxadmonition}{note}{Note:}
\sphinxAtStartPar
If you want the fictitious potential energy associated with the
added forces to be included in the total potential energy of the
system (the quantity being minimized), you MUST enable the
{\hyperref[\detokenize{fix_modify::doc}]{\sphinxcrossref{\DUrole{doc}{fix\_modify}}}} \sphinxstyleemphasis{energy} option for this fix.
\end{sphinxadmonition}


\subsection{Restrictions}
\label{\detokenize{fix_addforce:restrictions}}\begin{quote}

\sphinxAtStartPar
none
\end{quote}


\subsection{Related commands}
\label{\detokenize{fix_addforce:related-commands}}
\sphinxAtStartPar
{\hyperref[\detokenize{fix_setforce::doc}]{\sphinxcrossref{\DUrole{doc}{fix setforce}}}}, {\hyperref[\detokenize{fix_aveforce::doc}]{\sphinxcrossref{\DUrole{doc}{fix aveforce}}}}


\subsection{Default}
\label{\detokenize{fix_addforce:default}}
\sphinxAtStartPar
The option default for the every keyword is every = 1.

\sphinxstepscope

\index{fix addtorque@\spxentry{fix addtorque}}\ignorespaces 

\section{fix addtorque command}
\label{\detokenize{fix_addtorque:fix-addtorque-command}}\label{\detokenize{fix_addtorque:index-0}}\label{\detokenize{fix_addtorque::doc}}

\subsection{Syntax}
\label{\detokenize{fix_addtorque:syntax}}
\begin{sphinxVerbatim}[commandchars=\\\{\}]
\PYG{k}{fix }\PYG{n+nv+nvIdentifier}{ID}\PYG{+w}{ }\PYG{n+nv+nvIdentifier}{group\PYGZhy{}ID}\PYG{+w}{ }\PYG{n}{addtorque}\PYG{+w}{ }\PYG{n}{Tx}\PYG{+w}{ }\PYG{n}{Ty}\PYG{+w}{ }\PYG{n}{Tz}
\end{sphinxVerbatim}
\begin{itemize}
\item {} 
\sphinxAtStartPar
ID, group\sphinxhyphen{}ID are documented in {\hyperref[\detokenize{fix::doc}]{\sphinxcrossref{\DUrole{doc}{fix}}}} command

\item {} 
\sphinxAtStartPar
addtorque = style name of this fix command

\item {} 
\sphinxAtStartPar
Tx,Ty,Tz = torque component values (torque units)

\item {} 
\sphinxAtStartPar
any of Tx,Ty,Tz can be a variable (see below)

\end{itemize}


\subsection{Examples}
\label{\detokenize{fix_addtorque:examples}}
\begin{sphinxVerbatim}[commandchars=\\\{\}]
\PYG{k}{fix }\PYG{n+nv+nvIdentifier}{kick}\PYG{+w}{ }\PYG{n+nv+nvIdentifier}{bead}\PYG{+w}{ }\PYG{n}{addtorque}\PYG{+w}{ }\PYG{l+m}{2.0}\PYG{+w}{ }\PYG{l+m}{3.0}\PYG{+w}{ }\PYG{l+m}{5.0}
\PYG{k}{fix }\PYG{n+nv+nvIdentifier}{kick}\PYG{+w}{ }\PYG{n+nv+nvIdentifier}{bead}\PYG{+w}{ }\PYG{n}{addtorque}\PYG{+w}{ }\PYG{l+m}{0.0}\PYG{+w}{ }\PYG{l+m}{0.0}\PYG{+w}{ }\PYG{n}{v\PYGZus{}oscillate}
\end{sphinxVerbatim}


\subsection{Description}
\label{\detokenize{fix_addtorque:description}}
\sphinxAtStartPar
Add a set of forces to each atom in
the group such that:
\begin{itemize}
\item {} 
\sphinxAtStartPar
the components of the total torque applied on the group (around its
center of mass) are \(T_x\), \(T_y\), and \(T_z\)

\item {} 
\sphinxAtStartPar
the group would move as a rigid body in the absence of other
forces.

\end{itemize}

\sphinxAtStartPar
This command can be used to drive a group of atoms into rotation.

\sphinxAtStartPar
Any of the three quantities defining the torque components can be specified
as an equal\sphinxhyphen{}style {\hyperref[\detokenize{variable::doc}]{\sphinxcrossref{\DUrole{doc}{variable}}}}, namely \sphinxstyleemphasis{Tx},
\sphinxstyleemphasis{Ty}, \sphinxstyleemphasis{Tz}.  If the value is a variable, it should be specified as
v\_name, where name is the variable name.  In this case, the variable
will be evaluated each timestep, and its value used to determine the
torque component.

\sphinxAtStartPar
Equal\sphinxhyphen{}style variables can specify formulas with various mathematical
functions, and include {\hyperref[\detokenize{thermo_style::doc}]{\sphinxcrossref{\DUrole{doc}{thermo\_style}}}} command
keywords for the simulation box parameters and timestep and elapsed
time.  Thus it is easy to specify a time\sphinxhyphen{}dependent torque.


\bigskip\hrule\bigskip



\subsection{Restart, fix\_modify, output, run start/stop, minimize info}
\label{\detokenize{fix_addtorque:restart-fix-modify-output-run-start-stop-minimize-info}}
\sphinxAtStartPar
No information about this fix is written to
{\hyperref[\detokenize{restart::doc}]{\sphinxcrossref{\DUrole{doc}{binary restart files}}}}.

\sphinxAtStartPar
The {\hyperref[\detokenize{fix_modify::doc}]{\sphinxcrossref{\DUrole{doc}{fix\_modify}}}} \sphinxstyleemphasis{energy} option is supported by
this fix to add the potential “energy” inferred by the added torques
to the global potential energy of the system as part of
{\hyperref[\detokenize{thermo_style::doc}]{\sphinxcrossref{\DUrole{doc}{thermodynamic output}}}}.  The default setting for
this fix is {\hyperref[\detokenize{fix_modify::doc}]{\sphinxcrossref{\DUrole{doc}{fix\_modify energy no}}}}.  Note that this
is a fictitious quantity but is needed so that the {\hyperref[\detokenize{minimize::doc}]{\sphinxcrossref{\DUrole{doc}{minimize}}}} command can include the forces added by this fix in a
consistent manner (i.e., there is a decrease in potential energy when
atoms move in the direction of the added forces).

\sphinxAtStartPar
The {\hyperref[\detokenize{fix_modify::doc}]{\sphinxcrossref{\DUrole{doc}{fix\_modify}}}} \sphinxstyleemphasis{respa} option is supported by
this fix. This allows to set at which level of the {\hyperref[\detokenize{run_style::doc}]{\sphinxcrossref{\DUrole{doc}{r\sphinxhyphen{}RESPA}}}}
integrator the fix is adding its torque. Default is the outermost level.

\sphinxAtStartPar
This fix computes a global scalar and a global 3\sphinxhyphen{}vector, which can be
accessed by various {\hyperref[\detokenize{Howto_output::doc}]{\sphinxcrossref{\DUrole{doc}{output commands}}}}.  The scalar
is the potential energy discussed above.  The vector is the total
torque on the group of atoms before the forces on individual atoms are
changed by the fix.  The scalar and vector values calculated by this
fix are “extensive”.

\sphinxAtStartPar
No parameter of this fix can be used with the \sphinxstyleemphasis{start/stop} keywords of
the {\hyperref[\detokenize{run::doc}]{\sphinxcrossref{\DUrole{doc}{run}}}} command.

\sphinxAtStartPar
The forces due to this fix are imposed during an energy minimization,
invoked by the {\hyperref[\detokenize{minimize::doc}]{\sphinxcrossref{\DUrole{doc}{minimize}}}} command.

\begin{sphinxadmonition}{note}{Note:}
\sphinxAtStartPar
If you want the fictitious potential energy associated with the
added forces to be included in the total potential energy of the
system (the quantity being minimized), you MUST enable the
{\hyperref[\detokenize{fix_modify::doc}]{\sphinxcrossref{\DUrole{doc}{fix\_modify}}}} \sphinxstyleemphasis{energy} option for this fix.
\end{sphinxadmonition}

\begin{sphinxadmonition}{note}{Note:}
\sphinxAtStartPar
You should not specify force components with a variable that has
time\sphinxhyphen{}dependence for use with a minimizer, since the minimizer
increments the timestep as the iteration count during the
minimization.
\end{sphinxadmonition}


\subsection{Restrictions}
\label{\detokenize{fix_addtorque:restrictions}}
\sphinxAtStartPar
This fix is part of the EXTRA\sphinxhyphen{}FIX package.  It is only enabled if LAMMPS was
built with that package.  See the {\hyperref[\detokenize{Build_package::doc}]{\sphinxcrossref{\DUrole{doc}{Build package}}}} page for
more info.


\subsection{Related commands}
\label{\detokenize{fix_addtorque:related-commands}}
\sphinxAtStartPar
{\hyperref[\detokenize{fix_addforce::doc}]{\sphinxcrossref{\DUrole{doc}{fix addforce}}}}


\subsection{Default}
\label{\detokenize{fix_addtorque:default}}
\sphinxAtStartPar
none

\sphinxstepscope

\index{fix alchemy@\spxentry{fix alchemy}}\ignorespaces 

\section{fix alchemy command}
\label{\detokenize{fix_alchemy:fix-alchemy-command}}\label{\detokenize{fix_alchemy:index-0}}\label{\detokenize{fix_alchemy::doc}}

\subsection{Syntax}
\label{\detokenize{fix_alchemy:syntax}}
\begin{sphinxVerbatim}[commandchars=\\\{\}]
\PYG{k}{fix }\PYG{n+nv+nvIdentifier}{ID}\PYG{+w}{ }\PYG{n+nv+nvIdentifier}{group\PYGZhy{}ID}\PYG{+w}{ }\PYG{n}{alchemy}\PYG{+w}{ }\PYG{n}{v\PYGZus{}name}
\end{sphinxVerbatim}
\begin{itemize}
\item {} 
\sphinxAtStartPar
ID, group\sphinxhyphen{}ID are documented in {\hyperref[\detokenize{fix::doc}]{\sphinxcrossref{\DUrole{doc}{fix}}}} command

\item {} 
\sphinxAtStartPar
alchemy = style name of this fix command

\item {} 
\sphinxAtStartPar
v\_name = variable with name that determines the \(\lambda_R\) value

\end{itemize}


\subsection{Examples}
\label{\detokenize{fix_alchemy:examples}}
\begin{sphinxVerbatim}[commandchars=\\\{\}]
\PYG{k}{fix }\PYG{n+nv+nvIdentifier}{trans}\PYG{+w}{ }\PYG{n+nv+nvIdentifier}{all}\PYG{+w}{ }\PYG{n}{alchemy}\PYG{+w}{ }\PYG{n}{v\PYGZus{}ramp}
\end{sphinxVerbatim}


\subsection{Description}
\label{\detokenize{fix_alchemy:description}}
\sphinxAtStartPar
\DUrole{versionmodified,added}{Added in version 28Mar2023.}

\sphinxAtStartPar
This fix command enables an “alchemical transformation” to be performed
between two systems, whereby one system slowly transforms into the other
over the course of a molecular dynamics run.  This is useful for
measuring thermodynamic differences between two different systems.  It
also allows transformations that are not easily possible with the
{\hyperref[\detokenize{pair_hybrid::doc}]{\sphinxcrossref{\DUrole{doc}{pair style hybrid/scaled}}}}, {\hyperref[\detokenize{fix_adapt::doc}]{\sphinxcrossref{\DUrole{doc}{fix adapt}}}} or {\hyperref[\detokenize{fix_adapt_fep::doc}]{\sphinxcrossref{\DUrole{doc}{fix adapt/fep}}}} commands.

\sphinxAtStartPar
Example inputs are included in the \sphinxcode{\sphinxupquote{examples/PACKAGES/alchemy}}
directory for (a) transforming a pure copper system into a
copper/aluminum bronze alloy and (b) transforming two water molecules
in a box of water into a hydronium and a hydroxyl ion.

\sphinxAtStartPar
The two systems must be defined as {\hyperref[\detokenize{Howto_replica::doc}]{\sphinxcrossref{\DUrole{doc}{separate replica}}}} and run in separate partitions of processors using the
{\hyperref[\detokenize{Run_options::doc}]{\sphinxcrossref{\DUrole{doc}{\sphinxhyphen{}partition}}}} command\sphinxhyphen{}line switch.  Exactly two
partitions must be specified, and each partition must use the same number
of processors and the same domain decomposition.

\sphinxAtStartPar
Because the forces applied to the atoms are the same mix of the forces
from each partition and the simulation starts with the same atom
positions across both partitions, they will generate the same trajectory
of coordinates for each atom, and the same simulation box size and
shape.  The latter two conditions are \sphinxstyleemphasis{enforced} by this fix; it
exchanges coordinates and box information between the replicas.  This is
not strictly required, but since MD simulations are an example of a
chaotic system, even the tiniest random difference will eventually grow
exponentially into an unwanted divergence.

\sphinxAtStartPar
Otherwise, the properties of each atom (type, charge, bond and angle
partners, etc.), as well as energy and forces between interacting atoms
(pair, bond, angle styles, etc.) can be different in the two systems.

\sphinxAtStartPar
This can be initialized in the same input script by using commands which
only apply to one or the other replica.  The example scripts use a
world\sphinxhyphen{}style {\hyperref[\detokenize{variable::doc}]{\sphinxcrossref{\DUrole{doc}{variable}}}} command along with
{\hyperref[\detokenize{if::doc}]{\sphinxcrossref{\DUrole{doc}{if/then/else}}}} commands for this purpose.  The
{\hyperref[\detokenize{partition::doc}]{\sphinxcrossref{\DUrole{doc}{partition}}}} command can also be used.

\begin{sphinxVerbatim}[commandchars=\\\{\}]
\PYG{k}{create\PYGZus{}box }\PYG{n+nv+nvIdentifier}{2}\PYG{+w}{ }\PYG{n+nv+nvIdentifier}{box}
\PYG{k}{create\PYGZus{}atoms}\PYG{+w}{ }\PYG{l+m}{1}\PYG{+w}{ }\PYG{n}{box}
\PYG{k}{pair\PYGZus{}style}\PYG{+w}{ }\PYG{n}{eam}\PYG{o}{/}\PYG{n}{alloy}
\PYG{k}{pair\PYGZus{}coeff}\PYG{+w}{ }\PYG{o}{*}\PYG{+w}{ }\PYG{o}{*}\PYG{+w}{ }\PYG{n}{AlCu.eam.alloy}\PYG{+w}{ }\PYG{n}{Cu}\PYG{+w}{ }\PYG{n}{Al}

\PYG{c}{\PYGZsh{} replace 5\PYGZpc{} of copper with aluminum on the second partition only}

\PYG{k}{variable }\PYG{n+nv+nvIdentifier}{name}\PYG{+w}{ }\PYG{n}{world}\PYG{+w}{ }\PYG{n}{pure}\PYG{+w}{ }\PYG{n}{alloy}
\PYG{k}{if}\PYG{+w}{ }\PYG{l+s}{\PYGZdq{}}\PYG{l+s}{\PYGZdl{}\PYGZob{}name\PYGZcb{} == alloy}\PYG{l+s}{\PYGZdq{}}\PYG{+w}{ }\PYG{k}{then}\PYG{+w}{ }\PYG{o}{\PYGZam{}}
\PYG{+w}{  }\PYG{l+s}{\PYGZdq{}}\PYG{l+s}{set type 1 type/fraction 2 0.05 6745234}\PYG{l+s}{\PYGZdq{}}
\end{sphinxVerbatim}

\sphinxAtStartPar
Both replicas must define an instance of this fix, but with a different
\sphinxstyleemphasis{v\_name} variable.  The named variable must be an equal\sphinxhyphen{}style or
equivalent {\hyperref[\detokenize{variable::doc}]{\sphinxcrossref{\DUrole{doc}{variable}}}}.  The two variables should be
defined so that one ramps \sphinxstyleemphasis{down} from 1.0 to 0.0 for the \sphinxstyleemphasis{first} replica
(\sphinxstyleemphasis{R=0}) and the other ramps \sphinxstyleemphasis{up} from 0.0 to 1.0 for the \sphinxstyleemphasis{second}
replica (\sphinxstyleemphasis{R=1}).  A simple way is to do this is linearly, which can be
done using the ramp() function of the {\hyperref[\detokenize{variable::doc}]{\sphinxcrossref{\DUrole{doc}{variable}}}}
command.  You could also define a variable which returns a value between
0.0 and 1.0 as a non\sphinxhyphen{}linear function of the timestep.  Here is a linear
example:

\begin{sphinxVerbatim}[commandchars=\\\{\}]
\PYG{k}{partition}\PYG{+w}{ }\PYG{n}{yes}\PYG{+w}{ }\PYG{l+m}{1}\PYG{+w}{ }\PYG{k}{variable }\PYG{n+nv+nvIdentifier}{ramp}\PYG{+w}{ }\PYG{n}{equal}\PYG{+w}{ }\PYG{n}{ramp}\PYG{n+nv}{(}\PYG{n+nv}{1.0,0.0}\PYG{n+nv}{)}
\PYG{k}{partition}\PYG{+w}{ }\PYG{n}{yes}\PYG{+w}{ }\PYG{l+m}{2}\PYG{+w}{ }\PYG{k}{variable }\PYG{n+nv+nvIdentifier}{ramp}\PYG{+w}{ }\PYG{n}{equal}\PYG{+w}{ }\PYG{n}{ramp}\PYG{n+nv}{(}\PYG{n+nv}{0.0,1.0}\PYG{n+nv}{)}
\PYG{k}{fix }\PYG{n+nv+nvIdentifier}{2}\PYG{+w}{ }\PYG{n+nv+nvIdentifier}{all}\PYG{+w}{ }\PYG{n}{alchemy}\PYG{+w}{ }\PYG{n}{v\PYGZus{}ramp}
\end{sphinxVerbatim}

\begin{sphinxadmonition}{note}{Note:}
\sphinxAtStartPar
For an alchemical transformation, the two variables should sum to
exactly 1.0 at any timestep.  LAMMPS does \sphinxstyleemphasis{NOT} check that this is
the case.
\end{sphinxadmonition}

\sphinxAtStartPar
If you use the \sphinxcode{\sphinxupquote{ramp()}} function to define the two variables, this fix
can easily be used across successive runs in the same input script by
ensuring each instance of the {\hyperref[\detokenize{run::doc}]{\sphinxcrossref{\DUrole{doc}{run}}}} command specifies the
appropriate \sphinxstyleemphasis{start} or \sphinxstyleemphasis{stop} options.

\sphinxAtStartPar
At each timestep of an MD run, the two instances of this fix evaluate
their respective variables as a \(\lambda_R\) factor, where \sphinxstyleemphasis{R} = 0
or 1 for each replica.  The forces used by each system for the
propagation of their atoms is set to the sum of the forces for the two
systems, each scaled by their respective \(\lambda_R\) factor.  Thus,
during the MD run, the system will transform incrementally from the
first system to the second system.

\begin{sphinxadmonition}{note}{Note:}
\sphinxAtStartPar
As mentioned above, the coordinates of the atoms and box size/shape
must be exactly the same in the two replicas.  Therefore, it is
generally not a good idea to initialize the two replicas by reading
different data files or creating them individually from scratch.
Rather, a single system should be initialized and then desired
modifications applied to the system to either replica.  If your
input script somehow induces the two systems to become different
(e.g. by performing {\hyperref[\detokenize{atom_modify::doc}]{\sphinxcrossref{\DUrole{doc}{atom\_modify sort}}}}
differently, or by adding or depositing a different number of atoms),
then LAMMPS will detect the mismatch and generate an error.  This is
done by ensuring that each step the number and ordering of atoms is
identical within each pair of processors in the two replicas.
\end{sphinxadmonition}


\bigskip\hrule\bigskip



\subsection{Restart, fix\_modify, output, run start/stop, minimize info}
\label{\detokenize{fix_alchemy:restart-fix-modify-output-run-start-stop-minimize-info}}
\sphinxAtStartPar
No information about this fix is written to {\hyperref[\detokenize{restart::doc}]{\sphinxcrossref{\DUrole{doc}{binary restart files}}}}.  None of the {\hyperref[\detokenize{fix_modify::doc}]{\sphinxcrossref{\DUrole{doc}{fix\_modify}}}} options are
relevant to this fix.

\sphinxAtStartPar
This fix stores a global scalar (the current value of \(\lambda_R\))
and a global vector of length 3 which contains the potential energy of
the first partition, the second partition and the combined value,
respectively. The global scalar is unitless and “intensive”, the vector
is in {\hyperref[\detokenize{units::doc}]{\sphinxcrossref{\DUrole{doc}{energy units}}}} and “extensive”.  These values can be
used by any command that uses a global value from a fix as input.  See
the {\hyperref[\detokenize{Howto_output::doc}]{\sphinxcrossref{\DUrole{doc}{output howto}}}} page for an overview of LAMMPS
output options.

\sphinxAtStartPar
This fix is not invoked during {\hyperref[\detokenize{minimize::doc}]{\sphinxcrossref{\DUrole{doc}{energy minimization}}}}.


\subsection{Restrictions}
\label{\detokenize{fix_alchemy:restrictions}}
\sphinxAtStartPar
This fix is part of the REPLICA package.  It is only enabled if LAMMPS
was built with that package.  See the {\hyperref[\detokenize{Build_package::doc}]{\sphinxcrossref{\DUrole{doc}{Build package}}}} page for more info.

\sphinxAtStartPar
There may be only one instance of this fix in use at a time within
each replica.


\subsection{Related commands}
\label{\detokenize{fix_alchemy:related-commands}}
\sphinxAtStartPar
{\hyperref[\detokenize{compute_pressure_alchemy::doc}]{\sphinxcrossref{\DUrole{doc}{compute pressure/alchemy}}}} command,
{\hyperref[\detokenize{fix_adapt::doc}]{\sphinxcrossref{\DUrole{doc}{fix adapt}}}} command, {\hyperref[\detokenize{fix_adapt_fep::doc}]{\sphinxcrossref{\DUrole{doc}{fix adapt/fep}}}}
command, {\hyperref[\detokenize{pair_hybrid::doc}]{\sphinxcrossref{\DUrole{doc}{pair\_style hybrid/scaled}}}} command.


\subsection{Default}
\label{\detokenize{fix_alchemy:default}}
\sphinxAtStartPar
none

\sphinxstepscope

\index{fix amoeba/bitorsion@\spxentry{fix amoeba/bitorsion}}\ignorespaces 

\section{fix amoeba/bitorsion command}
\label{\detokenize{fix_amoeba_bitorsion:fix-amoeba-bitorsion-command}}\label{\detokenize{fix_amoeba_bitorsion:index-0}}\label{\detokenize{fix_amoeba_bitorsion::doc}}

\subsection{Syntax}
\label{\detokenize{fix_amoeba_bitorsion:syntax}}
\begin{sphinxVerbatim}[commandchars=\\\{\}]
\PYG{k}{fix }\PYG{n+nv+nvIdentifier}{ID}\PYG{+w}{ }\PYG{n+nv+nvIdentifier}{group\PYGZhy{}ID}\PYG{+w}{ }\PYG{n}{ameoba}\PYG{o}{/}\PYG{n}{bitorsion}\PYG{+w}{ }\PYG{n}{filename}
\end{sphinxVerbatim}
\begin{itemize}
\item {} 
\sphinxAtStartPar
ID, group\sphinxhyphen{}ID are documented in {\hyperref[\detokenize{fix::doc}]{\sphinxcrossref{\DUrole{doc}{fix}}}} command

\item {} 
\sphinxAtStartPar
amoeba/bitorsion = style name of this fix command

\item {} 
\sphinxAtStartPar
filename = force\sphinxhyphen{}field file with AMOEBA bitorsion coefficients

\end{itemize}


\subsection{Examples}
\label{\detokenize{fix_amoeba_bitorsion:examples}}
\begin{sphinxVerbatim}[commandchars=\\\{\}]
\PYG{k}{fix            }\PYG{n+nv+nvIdentifier}{bit}\PYG{+w}{ }\PYG{n+nv+nvIdentifier}{all}\PYG{+w}{ }\PYG{n}{amoeba}\PYG{o}{/}\PYG{n}{bitorsion}\PYG{+w}{ }\PYG{n}{bitorsion.ubiquitin.data}
\PYG{k}{read\PYGZus{}data}\PYG{+w}{      }\PYG{n}{proteinX.data}\PYG{+w}{ }\PYG{k}{fix }\PYG{n+nv+nvIdentifier}{bit}\PYG{+w}{ }\PYG{n+nv+nvIdentifier}{bitorsions}\PYG{+w}{ }\PYG{n}{BiTorsions}
\PYG{k}{fix\PYGZus{}modify     }\PYG{n+nv+nvIdentifier}{bit}\PYG{+w}{ }\PYG{n}{energy}\PYG{+w}{ }\PYG{n}{yes}
\end{sphinxVerbatim}


\subsection{Description}
\label{\detokenize{fix_amoeba_bitorsion:description}}
\sphinxAtStartPar
This command enables 5\sphinxhyphen{}body torsion/torsion interactions to be added
to simulations which use the AMOEBA and HIPPO force fields.  It
matches how the Tinker MD code computes its torsion/torsion
interactions for the AMOEBA and HIPPO force fields.  See the
{\hyperref[\detokenize{Howto_amoeba::doc}]{\sphinxcrossref{\DUrole{doc}{Howto amoeba}}}} doc page for more information about
the implementation of AMOEBA and HIPPO in LAMMPS.

\sphinxAtStartPar
Bitorsion interactions add additional potential energy contributions
to pairs of overlapping phi\sphinxhyphen{}psi dihedrals of amino\sphinxhyphen{}acids, which are
important to properly represent their conformational behavior.  Each
bitorsion interaction is thus defined for a 5\sphinxhyphen{}tuple of atoms
\(IJKLM\) with bonds between successive atoms in the list,
i.e. two overlapping dihedral interactions for atoms \(IJKL\) and
\(JKLM\).

\sphinxAtStartPar
The examples/amoeba directory has a sample input script and data file
for ubiquitin, which illustrates use of the fix amoeba/bitorsion
command.

\sphinxAtStartPar
As in the example above, this fix should be used before reading a data
file that contains a listing of bitorsion interactions.  The
\sphinxstyleemphasis{filename} specified should contain the bitorsion parameters for the
AMOEBA or HIPPO force field.

\sphinxAtStartPar
The data file read by the {\hyperref[\detokenize{read_data::doc}]{\sphinxcrossref{\DUrole{doc}{read\_data}}}} command must
contain the topology of all the bitorsion interactions, similar to the
topology data for bonds, angles, dihedrals, etc.  Specifically it
should have a line like this in its header section:

\begin{sphinxVerbatim}[commandchars=\\\{\}]
N bitorsions
\end{sphinxVerbatim}

\sphinxAtStartPar
where \(N\) is the number of bitorsion 5\sphinxhyphen{}body interactions.  It should
also have a section in the body of the data file like this with \(N\)
lines:

\begin{sphinxVerbatim}[commandchars=\\\{\}]
BiTorsions

       1       1       8      10      12      18      20
       2       5      18      20      22      25      27
       [...]
       N       3     314     315     317      318    330
\end{sphinxVerbatim}

\sphinxAtStartPar
The first column is an index from 1 to \(N\) to enumerate the
bitorsion 5\sphinxhyphen{}atom tuples; it is ignored by LAMMPS.  The second column
is the \sphinxstyleemphasis{type} of the interaction; it is an index into the bitorsion
force field file.  The remaining 5 columns are the atom IDs of the
atoms (in order) for the 5\sphinxhyphen{}tuple \(IJKLM\), as described above.

\sphinxAtStartPar
Note that the \sphinxstyleemphasis{bitorsions} and \sphinxstyleemphasis{BiTorsions} keywords for the header
and body sections match those specified in the {\hyperref[\detokenize{read_data::doc}]{\sphinxcrossref{\DUrole{doc}{read\_data}}}} command following the data file name.

\sphinxAtStartPar
The data file should be generated by using the
tools/tinker/tinker2lmp.py conversion script which creates a LAMMPS
data file from Tinker input files, including its PRM file which
contains the parameters necessary for computing bitorsion
interactions.  The script must be invoked with the optional
“\sphinxhyphen{}bitorsion” flag to do this; see the example for the ubiquitin system
in the tools/tinker/README file.  The same conversion script also
creates the file of bitorsion coefficient data which is read by this
command.

\sphinxAtStartPar
The potential energy associated with bitorsion interactions can be
output as described below.  It can also be included in the total
potential energy of the system, as output by the {\hyperref[\detokenize{thermo_style::doc}]{\sphinxcrossref{\DUrole{doc}{thermo\_style}}}} command, if the {\hyperref[\detokenize{fix_modify::doc}]{\sphinxcrossref{\DUrole{doc}{fix\_modify energy}}}}
command is used, as in the example above.  See the note below about
how to include the bitorsion energy when performing an {\hyperref[\detokenize{minimize::doc}]{\sphinxcrossref{\DUrole{doc}{energy
minimization}}}}.


\bigskip\hrule\bigskip



\subsection{Restart, fix\_modify, output, run start/stop, minimize info}
\label{\detokenize{fix_amoeba_bitorsion:restart-fix-modify-output-run-start-stop-minimize-info}}
\sphinxAtStartPar
This fix writes the list of bitorsion interactions to {\hyperref[\detokenize{restart::doc}]{\sphinxcrossref{\DUrole{doc}{binary
restart files}}}}.  See the {\hyperref[\detokenize{read_restart::doc}]{\sphinxcrossref{\DUrole{doc}{read\_restart}}}}
command for info on how to re\sphinxhyphen{}specify a fix in an input script that
reads a restart file, so that the operation of the fix continues in an
uninterrupted fashion.

\sphinxAtStartPar
The {\hyperref[\detokenize{fix_modify::doc}]{\sphinxcrossref{\DUrole{doc}{fix\_modify}}}} \sphinxstyleemphasis{energy} option is supported by
this fix to add the potential energy of the bitorsion interactions to
both the global potential energy and peratom potential energies of the
system as part of {\hyperref[\detokenize{thermo_style::doc}]{\sphinxcrossref{\DUrole{doc}{thermodynamic output}}}} or output
by the {\hyperref[\detokenize{compute_pe_atom::doc}]{\sphinxcrossref{\DUrole{doc}{compute pe/atom}}}} command.  The default
setting for this fix is {\hyperref[\detokenize{fix_modify::doc}]{\sphinxcrossref{\DUrole{doc}{fix\_modify energy yes}}}}.

\sphinxAtStartPar
The {\hyperref[\detokenize{fix_modify::doc}]{\sphinxcrossref{\DUrole{doc}{fix\_modify}}}} \sphinxstyleemphasis{virial} option is supported by
this fix to add the contribution due to the bitorsion interactions to
both the global pressure and per\sphinxhyphen{}atom stress of the system via the
{\hyperref[\detokenize{compute_pressure::doc}]{\sphinxcrossref{\DUrole{doc}{compute pressure}}}} and {\hyperref[\detokenize{compute_stress_atom::doc}]{\sphinxcrossref{\DUrole{doc}{compute
stress/atom}}}} commands.  The former can be
accessed by {\hyperref[\detokenize{thermo_style::doc}]{\sphinxcrossref{\DUrole{doc}{thermodynamic output}}}}.  The default
setting for this fix is {\hyperref[\detokenize{fix_modify::doc}]{\sphinxcrossref{\DUrole{doc}{fix\_modify virial yes}}}}.

\sphinxAtStartPar
This fix computes a global scalar which can be accessed by various
{\hyperref[\detokenize{Howto_output::doc}]{\sphinxcrossref{\DUrole{doc}{output commands}}}}.  The scalar is the potential
energy discussed above.  The scalar value calculated by this fix is
“extensive”.

\sphinxAtStartPar
No parameter of this fix can be used with the \sphinxstyleemphasis{start/stop} keywords of
the {\hyperref[\detokenize{run::doc}]{\sphinxcrossref{\DUrole{doc}{run}}}} command.

\sphinxAtStartPar
The forces due to this fix are imposed during an energy minimization,
invoked by the {\hyperref[\detokenize{minimize::doc}]{\sphinxcrossref{\DUrole{doc}{minimize}}}} command.

\sphinxAtStartPar
The {\hyperref[\detokenize{fix_modify::doc}]{\sphinxcrossref{\DUrole{doc}{fix\_modify}}}} \sphinxstyleemphasis{respa} option is supported by this
fix. This allows to set at which level of the {\hyperref[\detokenize{run_style::doc}]{\sphinxcrossref{\DUrole{doc}{r\sphinxhyphen{}RESPA}}}} integrator the fix is adding its forces. Default is the
outermost level.

\begin{sphinxadmonition}{note}{Note:}
\sphinxAtStartPar
For energy minimization, if you want the potential energy
associated with the bitorsion terms forces to be included in the
total potential energy of the system (the quantity being
minimized), you MUST not disable the {\hyperref[\detokenize{fix_modify::doc}]{\sphinxcrossref{\DUrole{doc}{fix\_modify}}}}
\sphinxstyleemphasis{energy} option for this fix.
\end{sphinxadmonition}


\subsection{Restrictions}
\label{\detokenize{fix_amoeba_bitorsion:restrictions}}
\sphinxAtStartPar
To function as expected this fix command must be issued \sphinxstyleemphasis{before} a
{\hyperref[\detokenize{read_data::doc}]{\sphinxcrossref{\DUrole{doc}{read\_data}}}} command but \sphinxstyleemphasis{after} a {\hyperref[\detokenize{read_restart::doc}]{\sphinxcrossref{\DUrole{doc}{read\_restart}}}} command.

\sphinxAtStartPar
This fix can only be used if LAMMPS was built with the AMOEBA package.
See the {\hyperref[\detokenize{Build_package::doc}]{\sphinxcrossref{\DUrole{doc}{Build package}}}} page for more info.


\subsection{Related commands}
\label{\detokenize{fix_amoeba_bitorsion:related-commands}}
\sphinxAtStartPar
{\hyperref[\detokenize{fix_modify::doc}]{\sphinxcrossref{\DUrole{doc}{fix\_modify}}}}, {\hyperref[\detokenize{read_data::doc}]{\sphinxcrossref{\DUrole{doc}{read\_data}}}}


\subsection{Default}
\label{\detokenize{fix_amoeba_bitorsion:default}}
\sphinxAtStartPar
none

\sphinxstepscope

\index{fix amoeba/pitorsion@\spxentry{fix amoeba/pitorsion}}\ignorespaces 

\section{fix amoeba/pitorsion command}
\label{\detokenize{fix_amoeba_pitorsion:fix-amoeba-pitorsion-command}}\label{\detokenize{fix_amoeba_pitorsion:index-0}}\label{\detokenize{fix_amoeba_pitorsion::doc}}

\subsection{Syntax}
\label{\detokenize{fix_amoeba_pitorsion:syntax}}
\begin{sphinxVerbatim}[commandchars=\\\{\}]
\PYG{k}{fix }\PYG{n+nv+nvIdentifier}{ID}\PYG{+w}{ }\PYG{n+nv+nvIdentifier}{group\PYGZhy{}ID}\PYG{+w}{ }\PYG{n}{ameoba}\PYG{o}{/}\PYG{n}{pitorsion}
\end{sphinxVerbatim}
\begin{itemize}
\item {} 
\sphinxAtStartPar
ID, group\sphinxhyphen{}ID are documented in {\hyperref[\detokenize{fix::doc}]{\sphinxcrossref{\DUrole{doc}{fix}}}} command

\item {} 
\sphinxAtStartPar
amoeba/pitorsion = style name of this fix command

\end{itemize}


\subsection{Examples}
\label{\detokenize{fix_amoeba_pitorsion:examples}}
\begin{sphinxVerbatim}[commandchars=\\\{\}]
\PYG{k}{fix            }\PYG{n+nv+nvIdentifier}{pit}\PYG{+w}{ }\PYG{n+nv+nvIdentifier}{all}\PYG{+w}{ }\PYG{n}{amoeba}\PYG{o}{/}\PYG{n}{pitorsion}
\PYG{k}{read\PYGZus{}data}\PYG{+w}{      }\PYG{n}{proteinX.data}\PYG{+w}{ }\PYG{k}{fix }\PYG{n+nv+nvIdentifier}{pit}\PYG{+w}{ }\PYG{l+s}{\PYGZdq{}}\PYG{l+s}{pitorsion types}\PYG{l+s}{\PYGZdq{}}\PYG{+w}{ }\PYG{l+s}{\PYGZdq{}}\PYG{l+s}{PiTorsion Coeffs}\PYG{l+s}{\PYGZdq{}}\PYG{+w}{ }\PYG{o}{\PYGZam{}}
\PYG{+w}{               }\PYG{k}{fix }\PYG{n+nv+nvIdentifier}{pit}\PYG{+w}{ }\PYG{n+nv+nvIdentifier}{pitorsions}\PYG{+w}{ }\PYG{n}{PiTorsions}
\PYG{k}{fix\PYGZus{}modify     }\PYG{n+nv+nvIdentifier}{pit}\PYG{+w}{ }\PYG{n}{energy}\PYG{+w}{ }\PYG{n}{yes}
\end{sphinxVerbatim}


\subsection{Description}
\label{\detokenize{fix_amoeba_pitorsion:description}}
\sphinxAtStartPar
This command enables 6\sphinxhyphen{}body pitorsion interactions to be added to
simulations which use the AMOEBA and HIPPO force fields.  It matches
how the Tinker MD code computes its pitorsion interactions for the
AMOEBA and HIPPO force fields.  See the {\hyperref[\detokenize{Howto_amoeba::doc}]{\sphinxcrossref{\DUrole{doc}{Howto amoeba}}}}
doc page for more information about the implementation of AMOEBA and HIPPO in
LAMMPS.

\sphinxAtStartPar
Pitorsion interactions add additional potential energy contributions
to 6\sphinxhyphen{}tuples of atoms \(IJKLMN\) that have a bond between atoms
\(K\) and \(L\), where both \(K\) and \(L\) are additionally
bonded to exactly two other atoms. Namely, \(K\) is also bonded to
\(I\) and \(J\), and \(L\) is also bonded to \(M\) and
\(N\).

\sphinxAtStartPar
The examples/amoeba directory has a sample input script and data file
for ubiquitin, which illustrates use of the fix amoeba/pitorsion
command.

\sphinxAtStartPar
As in the example above, this fix should be used before reading a data
file that contains a listing of pitorsion interactions.

\sphinxAtStartPar
The data file read by the {\hyperref[\detokenize{read_data::doc}]{\sphinxcrossref{\DUrole{doc}{read\_data}}}} command must
contain the topology of all the pitorsion interactions, similar to the
topology data for bonds, angles, dihedrals, etc.  Specifically, it
should have two lines like these in its header section:

\begin{sphinxVerbatim}[commandchars=\\\{\}]
M pitorsion types
N pitorsions
\end{sphinxVerbatim}

\sphinxAtStartPar
where \(N\) is the number of pitorsion 6\sphinxhyphen{}body interactions and \(M\) is
the number of pitorsion types.  It should also have two sections in the body
of the data file like these with \(M\) and \(N\) lines each:

\begin{sphinxVerbatim}[commandchars=\\\{\}]
PiTorsion Coeffs

       1 6.85
       2 10.2
       [...]
       M 6.85
\end{sphinxVerbatim}

\begin{sphinxVerbatim}[commandchars=\\\{\}]
PiTorsions

       1       1       2      4      3      20      21      24
       2       5      21     23     22      37      38      41
       [...]
       N       7      27     29     28      30      35      36
\end{sphinxVerbatim}

\sphinxAtStartPar
For PiTorsion Coeffs, the first column is an index from 1 to \(M\)
to enumerate the pitorsion types.  The second column is the single
prefactor coefficient needed for each type.

\sphinxAtStartPar
For PiTorsions, the first column is an index from 1 to \(N\) to
enumerate the pitorsion 6\sphinxhyphen{}atom tuples; it is ignored by LAMMPS.  The
second column is the “type” of the interaction; it is an index into
the PiTorsion Coeffs.  The remaining 6 columns are the atom IDs of the
atoms (in order) for the 6\sphinxhyphen{}tuple \(IJKLMN\), as described above.

\sphinxAtStartPar
Note that the \sphinxstyleemphasis{pitorsion types} and \sphinxstyleemphasis{pitorsions} and \sphinxstyleemphasis{PiTorsion
Coeffs} and \sphinxstyleemphasis{PiTorsions} keywords for the header and body sections of
the data file match those specified in the {\hyperref[\detokenize{read_data::doc}]{\sphinxcrossref{\DUrole{doc}{read\_data}}}}
command following the data file name.

\sphinxAtStartPar
The data file should be generated by using the
tools/tinker/tinker2lmp.py conversion script which creates a LAMMPS
data file from Tinker input files, including its PRM file which
contains the parameters necessary for computing pitorsion
interactions.

\sphinxAtStartPar
The potential energy associated with pitorsion interactions can be
output as described below.  It can also be included in the total
potential energy of the system, as output by the {\hyperref[\detokenize{thermo_style::doc}]{\sphinxcrossref{\DUrole{doc}{thermo\_style}}}} command, if the {\hyperref[\detokenize{fix_modify::doc}]{\sphinxcrossref{\DUrole{doc}{fix\_modify energy}}}}
command is used, as in the example above.  See the note below about
how to include the pitorsion energy when performing an {\hyperref[\detokenize{minimize::doc}]{\sphinxcrossref{\DUrole{doc}{energy
minimization}}}}.


\bigskip\hrule\bigskip



\subsection{Restart, fix\_modify, output, run start/stop, minimize info}
\label{\detokenize{fix_amoeba_pitorsion:restart-fix-modify-output-run-start-stop-minimize-info}}
\sphinxAtStartPar
This fix writes the list of pitorsion interactions to {\hyperref[\detokenize{restart::doc}]{\sphinxcrossref{\DUrole{doc}{binary
restart files}}}}.  See the {\hyperref[\detokenize{read_restart::doc}]{\sphinxcrossref{\DUrole{doc}{read\_restart}}}}
command for info on how to re\sphinxhyphen{}specify a fix in an input script that
reads a restart file, so that the operation of the fix continues in an
uninterrupted fashion.

\sphinxAtStartPar
The {\hyperref[\detokenize{fix_modify::doc}]{\sphinxcrossref{\DUrole{doc}{fix\_modify}}}} \sphinxstyleemphasis{energy} option is supported by
this fix to add the potential energy of the pitorsion interactions to
both the global potential energy and peratom potential energies of the
system as part of {\hyperref[\detokenize{thermo_style::doc}]{\sphinxcrossref{\DUrole{doc}{thermodynamic output}}}} or output
by the {\hyperref[\detokenize{compute_pe_atom::doc}]{\sphinxcrossref{\DUrole{doc}{compute pe/atom}}}} command.  The default
setting for this fix is {\hyperref[\detokenize{fix_modify::doc}]{\sphinxcrossref{\DUrole{doc}{fix\_modify energy yes}}}}.

\sphinxAtStartPar
The {\hyperref[\detokenize{fix_modify::doc}]{\sphinxcrossref{\DUrole{doc}{fix\_modify}}}} \sphinxstyleemphasis{virial} option is supported by
this fix to add the contribution due to the pitorsion interactions to
both the global pressure and per\sphinxhyphen{}atom stress of the system via the
{\hyperref[\detokenize{compute_pressure::doc}]{\sphinxcrossref{\DUrole{doc}{compute pressure}}}} and {\hyperref[\detokenize{compute_stress_atom::doc}]{\sphinxcrossref{\DUrole{doc}{compute
stress/atom}}}} commands.  The former can be
accessed by {\hyperref[\detokenize{thermo_style::doc}]{\sphinxcrossref{\DUrole{doc}{thermodynamic output}}}}.  The default
setting for this fix is {\hyperref[\detokenize{fix_modify::doc}]{\sphinxcrossref{\DUrole{doc}{fix\_modify virial yes}}}}.

\sphinxAtStartPar
This fix computes a global scalar which can be accessed by various
{\hyperref[\detokenize{Howto_output::doc}]{\sphinxcrossref{\DUrole{doc}{output commands}}}}.  The scalar is the potential
energy discussed above.  The scalar value calculated by this fix is
“extensive”.

\sphinxAtStartPar
No parameter of this fix can be used with the \sphinxstyleemphasis{start/stop} keywords of
the {\hyperref[\detokenize{run::doc}]{\sphinxcrossref{\DUrole{doc}{run}}}} command.

\sphinxAtStartPar
The forces due to this fix are imposed during an energy minimization,
invoked by the {\hyperref[\detokenize{minimize::doc}]{\sphinxcrossref{\DUrole{doc}{minimize}}}} command.

\sphinxAtStartPar
The {\hyperref[\detokenize{fix_modify::doc}]{\sphinxcrossref{\DUrole{doc}{fix\_modify}}}} \sphinxstyleemphasis{respa} option is supported by this
fix. This allows to set at which level of the {\hyperref[\detokenize{run_style::doc}]{\sphinxcrossref{\DUrole{doc}{r\sphinxhyphen{}RESPA}}}} integrator the fix is adding its forces. Default is the
outermost level.

\begin{sphinxadmonition}{note}{Note:}
\sphinxAtStartPar
For energy minimization, if you want the potential energy
associated with the pitorsion terms forces to be included in the
total potential energy of the system (the quantity being
minimized), you MUST not disable the {\hyperref[\detokenize{fix_modify::doc}]{\sphinxcrossref{\DUrole{doc}{fix\_modify}}}}
\sphinxstyleemphasis{energy} option for this fix.
\end{sphinxadmonition}


\subsection{Restrictions}
\label{\detokenize{fix_amoeba_pitorsion:restrictions}}
\sphinxAtStartPar
To function as expected this fix command must be issued \sphinxstyleemphasis{before} a
{\hyperref[\detokenize{read_data::doc}]{\sphinxcrossref{\DUrole{doc}{read\_data}}}} command but \sphinxstyleemphasis{after} a
{\hyperref[\detokenize{read_restart::doc}]{\sphinxcrossref{\DUrole{doc}{read\_restart}}}} command.

\sphinxAtStartPar
This fix can only be used if LAMMPS was built with the AMOEBA package.
See the {\hyperref[\detokenize{Build_package::doc}]{\sphinxcrossref{\DUrole{doc}{Build package}}}} page for more info.


\subsection{Related commands}
\label{\detokenize{fix_amoeba_pitorsion:related-commands}}
\sphinxAtStartPar
{\hyperref[\detokenize{fix_modify::doc}]{\sphinxcrossref{\DUrole{doc}{fix\_modify}}}}, {\hyperref[\detokenize{read_data::doc}]{\sphinxcrossref{\DUrole{doc}{read\_data}}}}


\subsection{Default}
\label{\detokenize{fix_amoeba_pitorsion:default}}
\sphinxAtStartPar
none

\sphinxstepscope

\index{fix append/atoms@\spxentry{fix append/atoms}}\ignorespaces 

\section{fix append/atoms command}
\label{\detokenize{fix_append_atoms:fix-append-atoms-command}}\label{\detokenize{fix_append_atoms:index-0}}\label{\detokenize{fix_append_atoms::doc}}

\subsection{Syntax}
\label{\detokenize{fix_append_atoms:syntax}}
\begin{sphinxVerbatim}[commandchars=\\\{\}]
\PYG{k}{fix }\PYG{n+nv+nvIdentifier}{ID}\PYG{+w}{ }\PYG{n+nv+nvIdentifier}{group\PYGZhy{}ID}\PYG{+w}{ }\PYG{n}{append}\PYG{o}{/}\PYG{n}{atoms}\PYG{+w}{ }\PYG{n}{face}\PYG{+w}{ }\PYG{n}{...}\PYG{+w}{ }\PYG{n}{keyword}\PYG{+w}{ }\PYG{n}{value}\PYG{+w}{ }\PYG{n}{...}
\end{sphinxVerbatim}
\begin{itemize}
\item {} 
\sphinxAtStartPar
ID, group\sphinxhyphen{}ID are documented in {\hyperref[\detokenize{fix::doc}]{\sphinxcrossref{\DUrole{doc}{fix}}}} command

\item {} 
\sphinxAtStartPar
append/atoms = style name of this fix command

\item {} 
\sphinxAtStartPar
face = \sphinxstyleemphasis{zhi}

\item {} 
\sphinxAtStartPar
zero or more keyword/value pairs may be appended

\item {} 
\sphinxAtStartPar
keyword = \sphinxstyleemphasis{basis} or \sphinxstyleemphasis{size} or \sphinxstyleemphasis{freq} or \sphinxstyleemphasis{temp} or \sphinxstyleemphasis{random} or \sphinxstyleemphasis{units}
\begin{sphinxalltt}
\sphinxstyleemphasis{basis} values = M itype
  M = which basis atom
  itype = atom type (1\sphinxhyphen{}N) to assign to this basis atom
\sphinxstyleemphasis{size} args = Lz
  Lz = z size of lattice region appended in a single event(distance units)
\sphinxstyleemphasis{freq} args = freq
  freq = the number of timesteps between append events
\sphinxstyleemphasis{temp} args = target damp seed extent
  target = target temperature for the region between zhi\sphinxhyphen{}extent and zhi (temperature units)
  damp = damping parameter (time units)
  seed = random number seed for langevin kicks
  extent = extent of thermostatted region (distance units)
\sphinxstyleemphasis{random} args = xmax ymax zmax seed
  \sphinxstyleemphasis{xmax}, \sphinxstyleemphasis{ymax}, \sphinxstyleemphasis{zmax} = maximum displacement in particular direction (distance units)
  \sphinxstyleemphasis{seed} = random number seed for random displacement
\sphinxstyleemphasis{units} value = \sphinxstyleemphasis{lattice} or \sphinxstyleemphasis{box}
  \sphinxstyleemphasis{lattice} = the wall position is defined in lattice units
  \sphinxstyleemphasis{box} = the wall position is defined in simulation box units
\end{sphinxalltt}

\end{itemize}


\subsection{Examples}
\label{\detokenize{fix_append_atoms:examples}}
\begin{sphinxVerbatim}[commandchars=\\\{\}]
\PYG{k}{fix }\PYG{n+nv+nvIdentifier}{1}\PYG{+w}{ }\PYG{n+nv+nvIdentifier}{all}\PYG{+w}{ }\PYG{n}{append}\PYG{o}{/}\PYG{n}{atoms}\PYG{+w}{ }\PYG{n}{zhi}\PYG{+w}{ }\PYG{n}{size}\PYG{+w}{ }\PYG{l+m}{5.0}\PYG{+w}{ }\PYG{n}{freq}\PYG{+w}{ }\PYG{l+m}{295}\PYG{+w}{ }\PYG{n}{units}\PYG{+w}{ }\PYG{n}{lattice}
\PYG{k}{fix }\PYG{n+nv+nvIdentifier}{4}\PYG{+w}{ }\PYG{n+nv+nvIdentifier}{all}\PYG{+w}{ }\PYG{n}{append}\PYG{o}{/}\PYG{n}{atoms}\PYG{+w}{ }\PYG{n}{zhi}\PYG{+w}{ }\PYG{n}{size}\PYG{+w}{ }\PYG{l+m}{15.0}\PYG{+w}{ }\PYG{n}{freq}\PYG{+w}{ }\PYG{l+m}{5}\PYG{+w}{ }\PYG{n}{units}\PYG{+w}{ }\PYG{n}{box}
\PYG{k}{fix }\PYG{n+nv+nvIdentifier}{A}\PYG{+w}{ }\PYG{n+nv+nvIdentifier}{all}\PYG{+w}{ }\PYG{n}{append}\PYG{o}{/}\PYG{n}{atoms}\PYG{+w}{ }\PYG{n}{zhi}\PYG{+w}{ }\PYG{n}{size}\PYG{+w}{ }\PYG{l+m}{1.0}\PYG{+w}{ }\PYG{n}{freq}\PYG{+w}{ }\PYG{l+m}{1000}\PYG{+w}{ }\PYG{n}{units}\PYG{+w}{ }\PYG{n}{lattice}
\end{sphinxVerbatim}


\subsection{Description}
\label{\detokenize{fix_append_atoms:description}}
\sphinxAtStartPar
This fix creates atoms on a lattice, appended on the zhi edge of the
system box.  This can be useful when a shock or wave is propagating
from zlo.  This allows the system to grow with time to accommodate an
expanding wave.  A simulation box must already exist, which is
typically created via the {\hyperref[\detokenize{create_box::doc}]{\sphinxcrossref{\DUrole{doc}{create\_box}}}} command.
Before using this command, a lattice must also be defined using the
{\hyperref[\detokenize{lattice::doc}]{\sphinxcrossref{\DUrole{doc}{lattice}}}} command.

\sphinxAtStartPar
This fix will automatically freeze atoms on the zhi edge of the
system, so that overlaps are avoided when new atoms are appended.

\sphinxAtStartPar
The \sphinxstyleemphasis{basis} keyword specifies an atom type that will be assigned to
specific basis atoms as they are created.  See the
{\hyperref[\detokenize{lattice::doc}]{\sphinxcrossref{\DUrole{doc}{lattice}}}} command for specifics on how basis atoms are
defined for the unit cell of the lattice.  By default, all created
atoms are assigned type = 1 unless this keyword specifies differently.

\sphinxAtStartPar
The \sphinxstyleemphasis{size} keyword defines the size in \(z\) of the chunk of material to
be added.

\sphinxAtStartPar
The \sphinxstyleemphasis{random} keyword will give the atoms random displacements around
their lattice points to simulate some initial temperature.

\sphinxAtStartPar
The \sphinxstyleemphasis{temp} keyword will cause a region to be thermostatted with a
Langevin thermostat on the zhi boundary.  The size of the region is
measured from zhi and is set with the \sphinxstyleemphasis{extent} argument.

\sphinxAtStartPar
The \sphinxstyleemphasis{units} keyword determines the meaning of the distance units used
to define a wall position, but only when a numeric constant is used.
A \sphinxstyleemphasis{box} value selects standard distance units as defined by the
{\hyperref[\detokenize{units::doc}]{\sphinxcrossref{\DUrole{doc}{units}}}} command (e.g., \(\AA\)
for units = real or metal.
A \sphinxstyleemphasis{lattice} value means the distance units are in lattice spacings.
The {\hyperref[\detokenize{lattice::doc}]{\sphinxcrossref{\DUrole{doc}{lattice}}}} command must have been previously used to
define the lattice spacings.


\bigskip\hrule\bigskip



\subsection{Restart, fix\_modify, output, run start/stop, minimize info}
\label{\detokenize{fix_append_atoms:restart-fix-modify-output-run-start-stop-minimize-info}}
\sphinxAtStartPar
No information about this fix is written to
{\hyperref[\detokenize{restart::doc}]{\sphinxcrossref{\DUrole{doc}{binary restart files}}}}.  None of the
{\hyperref[\detokenize{fix_modify::doc}]{\sphinxcrossref{\DUrole{doc}{fix\_modify}}}} options
are relevant to this fix.  No global or per\sphinxhyphen{}atom quantities are stored
by this fix for access by various {\hyperref[\detokenize{Howto_output::doc}]{\sphinxcrossref{\DUrole{doc}{output commands}}}}.
No parameter of this fix can be used with the \sphinxstyleemphasis{start/stop} keywords of
the {\hyperref[\detokenize{run::doc}]{\sphinxcrossref{\DUrole{doc}{run}}}} command.  This fix is not invoked during
{\hyperref[\detokenize{minimize::doc}]{\sphinxcrossref{\DUrole{doc}{energy minimization}}}}.


\subsection{Restrictions}
\label{\detokenize{fix_append_atoms:restrictions}}
\sphinxAtStartPar
This fix style is part of the SHOCK package.  It is only enabled if
LAMMPS was built with that package. See the
{\hyperref[\detokenize{Build_package::doc}]{\sphinxcrossref{\DUrole{doc}{Build package}}}} page for more info.

\sphinxAtStartPar
The boundary on which atoms are added with append/atoms must be
shrink/minimum.  The opposite boundary may be any boundary type other
than periodic.


\subsection{Related commands}
\label{\detokenize{fix_append_atoms:related-commands}}
\sphinxAtStartPar
{\hyperref[\detokenize{fix_wall_piston::doc}]{\sphinxcrossref{\DUrole{doc}{fix wall/piston}}}} command


\subsection{Default}
\label{\detokenize{fix_append_atoms:default}}
\sphinxAtStartPar
The keyword defaults are size = 0.0, freq = 0, units = lattice.  All
added atoms are of type 1 unless the basis keyword is used.

\sphinxstepscope

\index{fix atc@\spxentry{fix atc}}\ignorespaces 

\section{fix atc command}
\label{\detokenize{fix_atc:fix-atc-command}}\label{\detokenize{fix_atc:index-0}}\label{\detokenize{fix_atc::doc}}

\subsection{Syntax}
\label{\detokenize{fix_atc:syntax}}
\begin{sphinxVerbatim}[commandchars=\\\{\}]
\PYG{k}{fix }\PYG{o}{\PYGZlt{}}\PYG{n}{fixID}\PYG{o}{\PYGZgt{}}\PYG{+w}{ }\PYG{o}{\PYGZlt{}}\PYG{n}{group}\PYG{o}{\PYGZgt{}}\PYG{+w}{ }\PYG{n}{atc}\PYG{+w}{ }\PYG{o}{\PYGZlt{}}\PYG{n}{type}\PYG{o}{\PYGZgt{}}\PYG{+w}{ }\PYG{o}{\PYGZlt{}}\PYG{n}{parameter\PYGZus{}file}\PYG{o}{\PYGZgt{}}
\end{sphinxVerbatim}
\begin{itemize}
\item {} 
\sphinxAtStartPar
fixID = name of fix

\item {} 
\sphinxAtStartPar
group = name of group fix is to be applied

\item {} 
\sphinxAtStartPar
type = \sphinxstyleemphasis{thermal} or \sphinxstyleemphasis{two\_temperature} or \sphinxstyleemphasis{hardy} or \sphinxstyleemphasis{field}

\end{itemize}
\begin{sphinxalltt}
\sphinxstyleemphasis{thermal} = thermal coupling with fields: temperature
\sphinxstyleemphasis{two\_temperature} = electron\sphinxhyphen{}phonon coupling with field: temperature and electron\_temperature
\sphinxstyleemphasis{hardy} = on\sphinxhyphen{}the\sphinxhyphen{}fly post\sphinxhyphen{}processing using kernel localization functions
\sphinxstyleemphasis{field} = on\sphinxhyphen{}the\sphinxhyphen{}fly post\sphinxhyphen{}processing using mesh\sphinxhyphen{}based localization functions
\end{sphinxalltt}
\begin{itemize}
\item {} 
\sphinxAtStartPar
parameter\_file = name of the file with material parameters. Note: Neither hardy nor field requires a parameter file

\end{itemize}


\subsection{Examples}
\label{\detokenize{fix_atc:examples}}
\begin{sphinxVerbatim}[commandchars=\\\{\}]
\PYG{k}{fix }\PYG{n+nv+nvIdentifier}{AtC}\PYG{+w}{ }\PYG{n+nv+nvIdentifier}{internal}\PYG{+w}{ }\PYG{n}{atc}\PYG{+w}{ }\PYG{n}{thermal}\PYG{+w}{ }\PYG{n}{Ar\PYGZus{}thermal.dat}
\PYG{k}{fix }\PYG{n+nv+nvIdentifier}{AtC}\PYG{+w}{ }\PYG{n+nv+nvIdentifier}{internal}\PYG{+w}{ }\PYG{n}{atc}\PYG{+w}{ }\PYG{n}{two\PYGZus{}temperature}\PYG{+w}{ }\PYG{n}{Ar\PYGZus{}ttm.mat}
\PYG{k}{fix }\PYG{n+nv+nvIdentifier}{AtC}\PYG{+w}{ }\PYG{n+nv+nvIdentifier}{internal}\PYG{+w}{ }\PYG{n}{atc}\PYG{+w}{ }\PYG{n}{hardy}
\PYG{k}{fix }\PYG{n+nv+nvIdentifier}{AtC}\PYG{+w}{ }\PYG{n+nv+nvIdentifier}{internal}\PYG{+w}{ }\PYG{n}{atc}\PYG{+w}{ }\PYG{n}{field}
\end{sphinxVerbatim}


\subsection{Description}
\label{\detokenize{fix_atc:description}}
\sphinxAtStartPar
This fix is the beginning to creating a coupled FE/MD simulation and/or
an on\sphinxhyphen{}the\sphinxhyphen{}fly estimation of continuum fields. The coupled versions of
this fix do Verlet integration and the post\sphinxhyphen{}processing does not. After
instantiating this fix, several other fix\_modify commands will be
needed to set up the problem (i.e., define the finite element mesh and
prescribe initial and boundary conditions).

\noindent{\hspace*{\fill}\sphinxincludegraphics{{atc_nanotube}.jpg}\hspace*{\fill}}

\sphinxAtStartPar
The following coupling example is typical, but non\sphinxhyphen{}exhaustive:

\begin{sphinxVerbatim}[commandchars=\\\{\}]
\PYG{c}{\PYGZsh{} ... commands to create and initialize the MD system}

\PYG{c}{\PYGZsh{} initial fix to designate coupling type and group to apply it to}
\PYG{c}{\PYGZsh{} tag group physics material\PYGZus{}file}
\PYG{k}{fix }\PYG{n+nv+nvIdentifier}{AtC}\PYG{+w}{ }\PYG{n+nv+nvIdentifier}{internal}\PYG{+w}{ }\PYG{n}{atc}\PYG{+w}{ }\PYG{n}{thermal}\PYG{+w}{ }\PYG{n}{Ar\PYGZus{}thermal.mat}

\PYG{c}{\PYGZsh{} create a uniform 12 x 2 x 2 mesh that covers region contain the group}
\PYG{c}{\PYGZsh{} nx ny nz region periodicity}
\PYG{k}{fix\PYGZus{}modify }\PYG{n+nv+nvIdentifier}{AtC}\PYG{+w}{ }\PYG{n}{mesh}\PYG{+w}{ }\PYG{n}{create}\PYG{+w}{ }\PYG{l+m}{12}\PYG{+w}{ }\PYG{l+m}{2}\PYG{+w}{ }\PYG{l+m}{2}\PYG{+w}{ }\PYG{n}{mdRegion}\PYG{+w}{ }\PYG{n}{f}\PYG{+w}{ }\PYG{n}{p}\PYG{+w}{ }\PYG{n}{p}

\PYG{c}{\PYGZsh{} specify the control method for the type of coupling}
\PYG{c}{\PYGZsh{} physics control\PYGZus{}type}
\PYG{k}{fix\PYGZus{}modify }\PYG{n+nv+nvIdentifier}{AtC}\PYG{+w}{ }\PYG{n}{thermal}\PYG{+w}{ }\PYG{n}{control}\PYG{+w}{ }\PYG{n}{flux}

\PYG{c}{\PYGZsh{} specify the initial values for the empirical field \PYGZdq{}temperature\PYGZdq{}}
\PYG{c}{\PYGZsh{} field node\PYGZus{}group value}
\PYG{k}{fix\PYGZus{}modify }\PYG{n+nv+nvIdentifier}{AtC}\PYG{+w}{ }\PYG{n}{initial}\PYG{+w}{ }\PYG{n}{temperature}\PYG{+w}{ }\PYG{n}{all}\PYG{+w}{ }\PYG{l+m}{30}

\PYG{c}{\PYGZsh{} create an output stream for nodal fields}
\PYG{c}{\PYGZsh{} filename output\PYGZus{}frequency}
\PYG{k}{fix\PYGZus{}modify }\PYG{n+nv+nvIdentifier}{AtC}\PYG{+w}{ }\PYG{n}{output}\PYG{+w}{ }\PYG{n}{atc\PYGZus{}fe\PYGZus{}output}\PYG{+w}{ }\PYG{l+m}{100}

\PYG{k}{run}\PYG{+w}{ }\PYG{l+m}{1000}
\end{sphinxVerbatim}

\sphinxAtStartPar
likewise for this post\sphinxhyphen{}processing example:

\begin{sphinxVerbatim}[commandchars=\\\{\}]
\PYG{c}{\PYGZsh{} ... commands to create and initialize the MD system}

\PYG{c}{\PYGZsh{} initial fix to designate post\PYGZhy{}processing and the group to apply it to}
\PYG{c}{\PYGZsh{} no material file is allowed nor required}
\PYG{k}{fix }\PYG{n+nv+nvIdentifier}{AtC}\PYG{+w}{ }\PYG{n+nv+nvIdentifier}{internal}\PYG{+w}{ }\PYG{n}{atc}\PYG{+w}{ }\PYG{n}{hardy}

\PYG{c}{\PYGZsh{} for hardy fix, specific kernel function (function type and range) to \PYGZsh{} be used as a localization function}
\PYG{k}{fix }\PYG{n+nv+nvIdentifier}{AtC}\PYG{+w}{ }\PYG{n+nv+nvIdentifier}{kernel}\PYG{+w}{ }\PYG{n}{quartic\PYGZus{}sphere}\PYG{+w}{ }\PYG{l+m}{10.0}

\PYG{c}{\PYGZsh{} create a uniform 1 x 1 x 1 mesh that covers region contain the group}
\PYG{c}{\PYGZsh{} with periodicity this effectively creates a system average}
\PYG{k}{fix\PYGZus{}modify }\PYG{n+nv+nvIdentifier}{AtC}\PYG{+w}{ }\PYG{n}{mesh}\PYG{+w}{ }\PYG{n}{create}\PYG{+w}{ }\PYG{l+m}{1}\PYG{+w}{ }\PYG{l+m}{1}\PYG{+w}{ }\PYG{l+m}{1}\PYG{+w}{ }\PYG{n}{box}\PYG{+w}{ }\PYG{n}{p}\PYG{+w}{ }\PYG{n}{p}\PYG{+w}{ }\PYG{n}{p}

\PYG{c}{\PYGZsh{} change from default lagrangian map to eulerian}
\PYG{c}{\PYGZsh{} refreshed every 100 steps}
\PYG{k}{fix\PYGZus{}modify }\PYG{n+nv+nvIdentifier}{AtC}\PYG{+w}{ }\PYG{n}{atom\PYGZus{}element\PYGZus{}map}\PYG{+w}{ }\PYG{n}{eulerian}\PYG{+w}{ }\PYG{l+m}{100}

\PYG{c}{\PYGZsh{} start with no field defined}
\PYG{c}{\PYGZsh{} add mass density, potential energy density, stress and temperature}
\PYG{k}{fix\PYGZus{}modify }\PYG{n+nv+nvIdentifier}{AtC}\PYG{+w}{ }\PYG{n}{fields}\PYG{+w}{ }\PYG{n}{add}\PYG{+w}{ }\PYG{n}{density}\PYG{+w}{ }\PYG{n}{energy}\PYG{+w}{ }\PYG{n}{stress}\PYG{+w}{ }\PYG{n}{temperature}

\PYG{c}{\PYGZsh{} create an output stream for nodal fields}
\PYG{c}{\PYGZsh{} filename output\PYGZus{}frequency}
\PYG{k}{fix\PYGZus{}modify }\PYG{n+nv+nvIdentifier}{AtC}\PYG{+w}{ }\PYG{n}{output}\PYG{+w}{ }\PYG{n}{nvtFE}\PYG{+w}{ }\PYG{l+m}{100}\PYG{+w}{ }\PYG{n}{text}

\PYG{k}{run}\PYG{+w}{ }\PYG{l+m}{1000}
\end{sphinxVerbatim}

\sphinxAtStartPar
the mesh’s linear interpolation functions can be used as the localization function
by using the field option:

\begin{sphinxVerbatim}[commandchars=\\\{\}]
\PYG{k}{fix }\PYG{n+nv+nvIdentifier}{AtC}\PYG{+w}{ }\PYG{n+nv+nvIdentifier}{internal}\PYG{+w}{ }\PYG{n}{atc}\PYG{+w}{ }\PYG{n}{field}
\PYG{k}{fix\PYGZus{}modify }\PYG{n+nv+nvIdentifier}{AtC}\PYG{+w}{ }\PYG{n}{mesh}\PYG{+w}{ }\PYG{n}{create}\PYG{+w}{ }\PYG{l+m}{1}\PYG{+w}{ }\PYG{l+m}{1}\PYG{+w}{ }\PYG{l+m}{1}\PYG{+w}{ }\PYG{n}{box}\PYG{+w}{ }\PYG{n}{p}\PYG{+w}{ }\PYG{n}{p}\PYG{+w}{ }\PYG{n}{p}
\PYG{n}{...}
\end{sphinxVerbatim}

\sphinxAtStartPar
Note coupling and post\sphinxhyphen{}processing can be combined in the same simulations using separate fixes.


\bigskip\hrule\bigskip



\subsection{Restart, fix\_modify, output, run start/stop, minimize info}
\label{\detokenize{fix_atc:restart-fix-modify-output-run-start-stop-minimize-info}}
\sphinxAtStartPar
No information about this fix is written to {\hyperref[\detokenize{restart::doc}]{\sphinxcrossref{\DUrole{doc}{binary restart files}}}}.

\sphinxAtStartPar
The {\hyperref[\detokenize{fix_modify::doc}]{\sphinxcrossref{\DUrole{doc}{fix\_modify}}}} \sphinxstyleemphasis{energy} option is not supported by
this fix, but this fix does add the kinetic energy imparted to atoms
by the momentum coupling mode of the AtC package to the global
potential energy of the system as part of {\hyperref[\detokenize{thermo_style::doc}]{\sphinxcrossref{\DUrole{doc}{thermodynamic output}}}}.

\sphinxAtStartPar
Additional {\hyperref[\detokenize{fix_modify::doc}]{\sphinxcrossref{\DUrole{doc}{fix\_modify}}}} options relevant to this
fix are listed below.

\sphinxAtStartPar
This fix computes a global scalar which can be accessed by various
{\hyperref[\detokenize{Howto_output::doc}]{\sphinxcrossref{\DUrole{doc}{output commands}}}}.  The scalar is the energy
discussed in the previous paragraph.  The scalar value is “extensive”.

\sphinxAtStartPar
No parameter of this fix can be used with the
\sphinxstyleemphasis{start/stop} keywords of the {\hyperref[\detokenize{run::doc}]{\sphinxcrossref{\DUrole{doc}{run}}}} command.  This fix is not
invoked during {\hyperref[\detokenize{minimize::doc}]{\sphinxcrossref{\DUrole{doc}{energy minimization}}}}.


\subsection{Restrictions}
\label{\detokenize{fix_atc:restrictions}}
\sphinxAtStartPar
Thermal and two\_temperature (coupling) types use a Verlet
time\sphinxhyphen{}integration algorithm. The hardy type does not contain its own
time\sphinxhyphen{}integrator and must be used with a separate fix that does contain
one (e.g., nve, nvt). In addition, currently:
\begin{itemize}
\item {} 
\sphinxAtStartPar
the coupling is restricted to thermal physics

\item {} 
\sphinxAtStartPar
the FE computations are done in serial on each processor.

\end{itemize}


\subsection{Related commands}
\label{\detokenize{fix_atc:related-commands}}\label{\detokenize{fix_atc:atc-fix-modify}}
\sphinxAtStartPar
After specifying this fix in your input script, several
{\hyperref[\detokenize{fix_modify::doc}]{\sphinxcrossref{\DUrole{doc}{fix\_modify AtC}}}} commands are used to setup the
problem (e.g., define the finite element mesh and prescribe initial and
boundary conditions).  Each of these options has its own doc page.

\sphinxAtStartPar
\sphinxstyleemphasis{fix\_modify} commands for setup:
\begin{itemize}
\item {} 
\sphinxAtStartPar
{\hyperref[\detokenize{atc_mesh_create::doc}]{\sphinxcrossref{\DUrole{doc}{fix\_modify AtC mesh create}}}}

\item {} 
\sphinxAtStartPar
{\hyperref[\detokenize{atc_mesh_quadrature::doc}]{\sphinxcrossref{\DUrole{doc}{fix\_modify AtC mesh quadrature}}}}

\item {} 
\sphinxAtStartPar
{\hyperref[\detokenize{atc_mesh_read::doc}]{\sphinxcrossref{\DUrole{doc}{fix\_modify AtC mesh read}}}}

\item {} 
\sphinxAtStartPar
{\hyperref[\detokenize{atc_mesh_write::doc}]{\sphinxcrossref{\DUrole{doc}{fix\_modify AtC mesh write}}}}

\item {} 
\sphinxAtStartPar
{\hyperref[\detokenize{atc_mesh_create_nodeset::doc}]{\sphinxcrossref{\DUrole{doc}{fix\_modify AtC mesh create\_nodeset}}}}

\item {} 
\sphinxAtStartPar
{\hyperref[\detokenize{atc_mesh_add_to_nodeset::doc}]{\sphinxcrossref{\DUrole{doc}{fix\_modify AtC mesh add\_to\_nodeset}}}}

\item {} 
\sphinxAtStartPar
{\hyperref[\detokenize{atc_mesh_create_faceset_box::doc}]{\sphinxcrossref{\DUrole{doc}{fix\_modify AtC mesh create\_faceset box}}}}

\item {} 
\sphinxAtStartPar
{\hyperref[\detokenize{atc_mesh_create_faceset_plane::doc}]{\sphinxcrossref{\DUrole{doc}{fix\_modify AtC mesh create\_faceset plane}}}}

\item {} 
\sphinxAtStartPar
{\hyperref[\detokenize{atc_mesh_create_elementset::doc}]{\sphinxcrossref{\DUrole{doc}{fix\_modify AtC mesh create\_elementset}}}}

\item {} 
\sphinxAtStartPar
{\hyperref[\detokenize{atc_mesh_delete_elements::doc}]{\sphinxcrossref{\DUrole{doc}{fix\_modify AtC mesh delete\_elements}}}}

\item {} 
\sphinxAtStartPar
{\hyperref[\detokenize{atc_mesh_nodeset_to_elementset::doc}]{\sphinxcrossref{\DUrole{doc}{fix\_modify AtC mesh nodeset\_to\_elementset}}}}

\item {} 
\sphinxAtStartPar
{\hyperref[\detokenize{atc_boundary_type::doc}]{\sphinxcrossref{\DUrole{doc}{fix\_modify AtC boundary type}}}}

\item {} 
\sphinxAtStartPar
{\hyperref[\detokenize{atc_internal_quadrature::doc}]{\sphinxcrossref{\DUrole{doc}{fix\_modify AtC internal\_quadrature}}}}

\item {} 
\sphinxAtStartPar
{\hyperref[\detokenize{atc_time_integration::doc}]{\sphinxcrossref{\DUrole{doc}{fix\_modify AtC time\_integration}}}}

\item {} 
\sphinxAtStartPar
{\hyperref[\detokenize{atc_electron_integration::doc}]{\sphinxcrossref{\DUrole{doc}{fix\_modify AtC extrinsic electron\_integration}}}}

\item {} 
\sphinxAtStartPar
{\hyperref[\detokenize{atc_internal_element_set::doc}]{\sphinxcrossref{\DUrole{doc}{fix\_modify AtC internal\_element\_set}}}}

\item {} 
\sphinxAtStartPar
{\hyperref[\detokenize{atc_decomposition::doc}]{\sphinxcrossref{\DUrole{doc}{fix\_modify AtC decomposition}}}}

\end{itemize}

\sphinxAtStartPar
\sphinxstyleemphasis{fix\_modify} commands for boundary and initial conditions:
\begin{itemize}
\item {} 
\sphinxAtStartPar
{\hyperref[\detokenize{atc_initial::doc}]{\sphinxcrossref{\DUrole{doc}{fix\_modify AtC initial}}}}

\item {} 
\sphinxAtStartPar
{\hyperref[\detokenize{atc_fix::doc}]{\sphinxcrossref{\DUrole{doc}{fix\_modify AtC fix}}}}

\item {} 
\sphinxAtStartPar
{\hyperref[\detokenize{atc_unfix::doc}]{\sphinxcrossref{\DUrole{doc}{fix\_modify AtC unfix}}}}

\item {} 
\sphinxAtStartPar
{\hyperref[\detokenize{atc_fix_flux::doc}]{\sphinxcrossref{\DUrole{doc}{fix\_modify AtC fix\_flux}}}}

\item {} 
\sphinxAtStartPar
{\hyperref[\detokenize{atc_unfix_flux::doc}]{\sphinxcrossref{\DUrole{doc}{fix\_modify AtC unfix\_flux}}}}

\item {} 
\sphinxAtStartPar
{\hyperref[\detokenize{atc_source::doc}]{\sphinxcrossref{\DUrole{doc}{fix\_modify AtC source}}}}

\item {} 
\sphinxAtStartPar
{\hyperref[\detokenize{atc_remove_source::doc}]{\sphinxcrossref{\DUrole{doc}{fix\_modify AtC remove\_source}}}}

\end{itemize}

\sphinxAtStartPar
\sphinxstyleemphasis{fix\_modify} commands for control and filtering:
\begin{itemize}
\item {} 
\sphinxAtStartPar
{\hyperref[\detokenize{atc_control_thermal::doc}]{\sphinxcrossref{\DUrole{doc}{fix\_modify AtC control thermal}}}}

\item {} 
\sphinxAtStartPar
{\hyperref[\detokenize{atc_control_momentum::doc}]{\sphinxcrossref{\DUrole{doc}{fix\_modify AtC control momentum}}}}

\item {} 
\sphinxAtStartPar
{\hyperref[\detokenize{atc_control_localized_lambda::doc}]{\sphinxcrossref{\DUrole{doc}{fix\_modify AtC control localized\_lambda}}}}

\item {} 
\sphinxAtStartPar
{\hyperref[\detokenize{atc_lumped_lambda_solve::doc}]{\sphinxcrossref{\DUrole{doc}{fix\_modify AtC control lumped\_lambda\_solve}}}}

\item {} 
\sphinxAtStartPar
{\hyperref[\detokenize{atc_mask_direction::doc}]{\sphinxcrossref{\DUrole{doc}{fix\_modify AtC control mask\_direction}}}}

\item {} 
\sphinxAtStartPar
{\hyperref[\detokenize{atc_time_filter::doc}]{\sphinxcrossref{\DUrole{doc}{fix\_modify AtC filter}}}}

\item {} 
\sphinxAtStartPar
{\hyperref[\detokenize{atc_filter_scale::doc}]{\sphinxcrossref{\DUrole{doc}{fix\_modify AtC filter scale}}}}

\item {} 
\sphinxAtStartPar
{\hyperref[\detokenize{atc_filter_type::doc}]{\sphinxcrossref{\DUrole{doc}{fix\_modify AtC filter type}}}}

\item {} 
\sphinxAtStartPar
{\hyperref[\detokenize{atc_equilibrium_start::doc}]{\sphinxcrossref{\DUrole{doc}{fix\_modify AtC equilibrium\_start}}}}

\item {} 
\sphinxAtStartPar
{\hyperref[\detokenize{atc_extrinsic_exchange::doc}]{\sphinxcrossref{\DUrole{doc}{fix\_modify AtC extrinsic exchange}}}}

\item {} 
\sphinxAtStartPar
{\hyperref[\detokenize{atc_poisson_solver::doc}]{\sphinxcrossref{\DUrole{doc}{fix\_modify AtC poisson\_solver}}}}

\end{itemize}

\sphinxAtStartPar
\sphinxstyleemphasis{fix\_modify} commands for output:
\begin{itemize}
\item {} 
\sphinxAtStartPar
{\hyperref[\detokenize{atc_output::doc}]{\sphinxcrossref{\DUrole{doc}{fix\_modify AtC output}}}}

\item {} 
\sphinxAtStartPar
{\hyperref[\detokenize{atc_output_nodeset::doc}]{\sphinxcrossref{\DUrole{doc}{fix\_modify AtC output nodeset}}}}

\item {} 
\sphinxAtStartPar
{\hyperref[\detokenize{atc_output_volume_integral::doc}]{\sphinxcrossref{\DUrole{doc}{fix\_modify AtC output volume\_integral}}}}

\item {} 
\sphinxAtStartPar
{\hyperref[\detokenize{atc_output_boundary_integral::doc}]{\sphinxcrossref{\DUrole{doc}{fix\_modify AtC output boundary\_integral}}}}

\item {} 
\sphinxAtStartPar
{\hyperref[\detokenize{atc_output_contour_integral::doc}]{\sphinxcrossref{\DUrole{doc}{fix\_modify AtC output contour\_integral}}}}

\item {} 
\sphinxAtStartPar
{\hyperref[\detokenize{atc_mesh_output::doc}]{\sphinxcrossref{\DUrole{doc}{fix\_modify AtC mesh output}}}}

\item {} 
\sphinxAtStartPar
{\hyperref[\detokenize{atc_write_restart::doc}]{\sphinxcrossref{\DUrole{doc}{fix\_modify AtC write\_restart}}}}

\item {} 
\sphinxAtStartPar
{\hyperref[\detokenize{atc_read_restart::doc}]{\sphinxcrossref{\DUrole{doc}{fix\_modify AtC read\_restart}}}}

\end{itemize}

\sphinxAtStartPar
\sphinxstyleemphasis{fix\_modify} commands for post\sphinxhyphen{}processing:
\begin{itemize}
\item {} 
\sphinxAtStartPar
{\hyperref[\detokenize{atc_hardy_kernel::doc}]{\sphinxcrossref{\DUrole{doc}{fix\_modify AtC kernel}}}}

\item {} 
\sphinxAtStartPar
{\hyperref[\detokenize{atc_hardy_fields::doc}]{\sphinxcrossref{\DUrole{doc}{fix\_modify AtC fields}}}}

\item {} 
\sphinxAtStartPar
{\hyperref[\detokenize{atc_hardy_gradients::doc}]{\sphinxcrossref{\DUrole{doc}{fix\_modify AtC gradients}}}}

\item {} 
\sphinxAtStartPar
{\hyperref[\detokenize{atc_hardy_rates::doc}]{\sphinxcrossref{\DUrole{doc}{fix\_modify AtC rates}}}}

\item {} 
\sphinxAtStartPar
{\hyperref[\detokenize{atc_hardy_computes::doc}]{\sphinxcrossref{\DUrole{doc}{fix\_modify AtC computes}}}}

\item {} 
\sphinxAtStartPar
{\hyperref[\detokenize{atc_hardy_on_the_fly::doc}]{\sphinxcrossref{\DUrole{doc}{fix\_modify AtC on\_the\_fly}}}}

\item {} 
\sphinxAtStartPar
{\hyperref[\detokenize{atc_pair_interactions::doc}]{\sphinxcrossref{\DUrole{doc}{fix\_modify AtC pair/bond\_interactions}}}}

\item {} 
\sphinxAtStartPar
{\hyperref[\detokenize{atc_sample_frequency::doc}]{\sphinxcrossref{\DUrole{doc}{fix\_modify AtC sample\_frequency}}}}

\item {} 
\sphinxAtStartPar
{\hyperref[\detokenize{atc_set_reference_pe::doc}]{\sphinxcrossref{\DUrole{doc}{fix\_modify AtC set}}}}

\end{itemize}

\sphinxAtStartPar
miscellaneous \sphinxstyleemphasis{fix\_modify} commands:
\begin{itemize}
\item {} 
\sphinxAtStartPar
{\hyperref[\detokenize{atc_atom_element_map::doc}]{\sphinxcrossref{\DUrole{doc}{fix\_modify AtC atom\_element\_map}}}}

\item {} 
\sphinxAtStartPar
{\hyperref[\detokenize{atc_atom_weight::doc}]{\sphinxcrossref{\DUrole{doc}{fix\_modify AtC atom\_weight}}}}

\item {} 
\sphinxAtStartPar
{\hyperref[\detokenize{atc_write_atom_weights::doc}]{\sphinxcrossref{\DUrole{doc}{fix\_modify AtC write\_atom\_weights}}}}

\item {} 
\sphinxAtStartPar
{\hyperref[\detokenize{atc_kernel_bandwidth::doc}]{\sphinxcrossref{\DUrole{doc}{fix\_modify AtC kernel\_bandwidth}}}}

\item {} 
\sphinxAtStartPar
{\hyperref[\detokenize{atc_reset_time::doc}]{\sphinxcrossref{\DUrole{doc}{fix\_modify AtC reset\_time}}}}

\item {} 
\sphinxAtStartPar
{\hyperref[\detokenize{atc_reset_atomic_reference::doc}]{\sphinxcrossref{\DUrole{doc}{fix\_modify AtC reset\_atomic\_reference\_positions}}}}

\item {} 
\sphinxAtStartPar
{\hyperref[\detokenize{atc_fe_md_boundary::doc}]{\sphinxcrossref{\DUrole{doc}{fix\_modify AtC fe\_md\_boundary}}}}

\item {} 
\sphinxAtStartPar
{\hyperref[\detokenize{atc_boundary_faceset::doc}]{\sphinxcrossref{\DUrole{doc}{fix\_modify AtC boundary\_faceset}}}}

\item {} 
\sphinxAtStartPar
{\hyperref[\detokenize{atc_consistent_fe_initialization::doc}]{\sphinxcrossref{\DUrole{doc}{fix\_modify AtC consistent\_fe\_initialization}}}}

\item {} 
\sphinxAtStartPar
{\hyperref[\detokenize{atc_mass_matrix::doc}]{\sphinxcrossref{\DUrole{doc}{fix\_modify AtC mass\_matrix}}}}

\item {} 
\sphinxAtStartPar
{\hyperref[\detokenize{atc_material::doc}]{\sphinxcrossref{\DUrole{doc}{fix\_modify AtC material}}}}

\item {} 
\sphinxAtStartPar
{\hyperref[\detokenize{atc_atomic_charge::doc}]{\sphinxcrossref{\DUrole{doc}{fix\_modify AtC atomic\_charge}}}}

\item {} 
\sphinxAtStartPar
{\hyperref[\detokenize{atc_source_integration::doc}]{\sphinxcrossref{\DUrole{doc}{fix\_modify AtC source\_integration}}}}

\item {} 
\sphinxAtStartPar
{\hyperref[\detokenize{atc_temperature_definition::doc}]{\sphinxcrossref{\DUrole{doc}{fix\_modify AtC temperature\_definition}}}}

\item {} 
\sphinxAtStartPar
{\hyperref[\detokenize{atc_track_displacement::doc}]{\sphinxcrossref{\DUrole{doc}{fix\_modify AtC track\_displacement}}}}

\item {} 
\sphinxAtStartPar
{\hyperref[\detokenize{atc_boundary_dynamics::doc}]{\sphinxcrossref{\DUrole{doc}{fix\_modify AtC boundary\_dynamics}}}}

\item {} 
\sphinxAtStartPar
{\hyperref[\detokenize{atc_add_species::doc}]{\sphinxcrossref{\DUrole{doc}{fix\_modify AtC add\_species}}}}

\item {} 
\sphinxAtStartPar
{\hyperref[\detokenize{atc_add_molecule::doc}]{\sphinxcrossref{\DUrole{doc}{fix\_modify AtC add\_molecule}}}}

\item {} 
\sphinxAtStartPar
{\hyperref[\detokenize{atc_remove_species::doc}]{\sphinxcrossref{\DUrole{doc}{fix\_modify AtC remove\_species}}}}

\item {} 
\sphinxAtStartPar
{\hyperref[\detokenize{atc_remove_molecule::doc}]{\sphinxcrossref{\DUrole{doc}{fix\_modify AtC remove\_molecule}}}}

\end{itemize}

\sphinxAtStartPar
Note: a set of example input files with the attendant material files
are included in the \sphinxcode{\sphinxupquote{examples/PACKAGES/atc}} folders.


\subsection{Default}
\label{\detokenize{fix_atc:default}}
\sphinxAtStartPar
None


\bigskip\hrule\bigskip


\sphinxAtStartPar
For detailed exposition of the theory and algorithms please see:

\phantomsection\label{\detokenize{fix_atc:wagner}}
\sphinxAtStartPar
\sphinxstylestrong{(Wagner)} Wagner, GJ; Jones, RE; Templeton, JA; Parks, MA, “An
atomistic\sphinxhyphen{}to\sphinxhyphen{}continuum coupling method for heat transfer in solids.”
Special Issue of Computer Methods and Applied Mechanics (2008)
197:3351.

\phantomsection\label{\detokenize{fix_atc:zimmeman2004}}
\sphinxAtStartPar
\sphinxstylestrong{(Zimmerman2004)} Zimmerman, JA; Webb, EB; Hoyt, JJ;. Jones, RE;
Klein, PA; Bammann, DJ, “Calculation of stress in atomistic
simulation.” Special Issue of Modelling and Simulation in Materials
Science and Engineering (2004), 12:S319.

\phantomsection\label{\detokenize{fix_atc:zimmerman2010}}
\sphinxAtStartPar
\sphinxstylestrong{(Zimmerman2010)} Zimmerman, JA; Jones, RE; Templeton, JA, “A
material frame approach for evaluating continuum variables in
atomistic simulations.” Journal of Computational Physics (2010),
229:2364.

\phantomsection\label{\detokenize{fix_atc:templeton2010}}
\sphinxAtStartPar
\sphinxstylestrong{(Templeton2010)} Templeton, JA; Jones, RE; Wagner, GJ, “Application
of a field\sphinxhyphen{}based method to spatially varying thermal transport
problems in molecular dynamics.” Modelling and Simulation in
Materials Science and Engineering (2010), 18:085007.

\phantomsection\label{\detokenize{fix_atc:jones}}
\sphinxAtStartPar
\sphinxstylestrong{(Jones)} Jones, RE; Templeton, JA; Wagner, GJ; Olmsted, D; Modine,
JA, “Electron transport enhanced molecular dynamics for metals and
semi\sphinxhyphen{}metals.” International Journal for Numerical Methods in
Engineering (2010), 83:940.

\phantomsection\label{\detokenize{fix_atc:templeton2011}}
\sphinxAtStartPar
\sphinxstylestrong{(Templeton2011)} Templeton, JA; Jones, RE; Lee, JW; Zimmerman, JA;
Wong, BM, “A long\sphinxhyphen{}range electric field solver for molecular dynamics
based on atomistic\sphinxhyphen{}to\sphinxhyphen{}continuum modeling.” Journal of Chemical Theory
and Computation (2011), 7:1736.

\phantomsection\label{\detokenize{fix_atc:mandadapu}}
\sphinxAtStartPar
\sphinxstylestrong{(Mandadapu)} Mandadapu, KK; Templeton, JA; Lee, JW, “Polarization
as a field variable from molecular dynamics simulations.” Journal of
Chemical Physics (2013), 139:054115.

\sphinxAtStartPar
Please refer to the standard finite element (FE) texts (e.g., T.J.R.
Hughes, \sphinxstyleemphasis{The Finite Element Method,} Dover 2003) for the basics of FE
simulations.

\sphinxstepscope

\index{fix atom/swap@\spxentry{fix atom/swap}}\ignorespaces 

\section{fix atom/swap command}
\label{\detokenize{fix_atom_swap:fix-atom-swap-command}}\label{\detokenize{fix_atom_swap:index-0}}\label{\detokenize{fix_atom_swap::doc}}

\subsection{Syntax}
\label{\detokenize{fix_atom_swap:syntax}}
\begin{sphinxVerbatim}[commandchars=\\\{\}]
\PYG{k}{fix }\PYG{n+nv+nvIdentifier}{ID}\PYG{+w}{ }\PYG{n+nv+nvIdentifier}{group\PYGZhy{}ID}\PYG{+w}{ }\PYG{n}{atom}\PYG{o}{/}\PYG{n}{swap}\PYG{+w}{ }\PYG{n}{N}\PYG{+w}{ }\PYG{n}{X}\PYG{+w}{ }\PYG{n}{seed}\PYG{+w}{ }\PYG{n}{T}\PYG{+w}{ }\PYG{n}{keyword}\PYG{+w}{ }\PYG{n}{values}\PYG{+w}{ }\PYG{n}{...}
\end{sphinxVerbatim}
\begin{itemize}
\item {} 
\sphinxAtStartPar
ID, group\sphinxhyphen{}ID are documented in {\hyperref[\detokenize{fix::doc}]{\sphinxcrossref{\DUrole{doc}{fix}}}} command

\item {} 
\sphinxAtStartPar
atom/swap = style name of this fix command

\item {} 
\sphinxAtStartPar
N = invoke this fix every N steps

\item {} 
\sphinxAtStartPar
X = number of swaps to attempt every N steps

\item {} 
\sphinxAtStartPar
seed = random \# seed (positive integer)

\item {} 
\sphinxAtStartPar
T = scaling temperature of the MC swaps (temperature units)

\item {} 
\sphinxAtStartPar
one or more keyword/value pairs may be appended to args

\item {} 
\sphinxAtStartPar
keyword = \sphinxstyleemphasis{types} or \sphinxstyleemphasis{mu} or \sphinxstyleemphasis{ke} or \sphinxstyleemphasis{semi\sphinxhyphen{}grand} or \sphinxstyleemphasis{region}
\begin{sphinxalltt}
\sphinxstyleemphasis{types} values = two or more atom types (1\sphinxhyphen{}Ntypes or type label)
\sphinxstyleemphasis{mu} values = chemical potential of swap types (energy units)
\sphinxstyleemphasis{ke} value = \sphinxstyleemphasis{no} or \sphinxstyleemphasis{yes}
  \sphinxstyleemphasis{no} = no conservation of kinetic energy after atom swaps
  \sphinxstyleemphasis{yes} = kinetic energy is conserved after atom swaps
\sphinxstyleemphasis{semi\sphinxhyphen{}grand} value = \sphinxstyleemphasis{no} or \sphinxstyleemphasis{yes}
  \sphinxstyleemphasis{no} = particle type counts and fractions conserved
  \sphinxstyleemphasis{yes} = semi\sphinxhyphen{}grand canonical ensemble, particle fractions not conserved
\sphinxstyleemphasis{region} value = region\sphinxhyphen{}ID
  region\sphinxhyphen{}ID = ID of region to use as an exchange/move volume
\end{sphinxalltt}

\end{itemize}


\subsection{Examples}
\label{\detokenize{fix_atom_swap:examples}}
\begin{sphinxVerbatim}[commandchars=\\\{\}]
\PYG{k}{fix }\PYG{n+nv+nvIdentifier}{2}\PYG{+w}{ }\PYG{n+nv+nvIdentifier}{all}\PYG{+w}{ }\PYG{n}{atom}\PYG{o}{/}\PYG{n}{swap}\PYG{+w}{ }\PYG{l+m}{1}\PYG{+w}{ }\PYG{l+m}{1}\PYG{+w}{ }\PYG{l+m}{29494}\PYG{+w}{ }\PYG{l+m}{300.0}\PYG{+w}{ }\PYG{n}{ke}\PYG{+w}{ }\PYG{n}{no}\PYG{+w}{ }\PYG{n}{types}\PYG{+w}{ }\PYG{l+m}{1}\PYG{+w}{ }\PYG{l+m}{2}
\PYG{k}{fix }\PYG{n+nv+nvIdentifier}{myFix}\PYG{+w}{ }\PYG{n+nv+nvIdentifier}{all}\PYG{+w}{ }\PYG{n}{atom}\PYG{o}{/}\PYG{n}{swap}\PYG{+w}{ }\PYG{l+m}{100}\PYG{+w}{ }\PYG{l+m}{1}\PYG{+w}{ }\PYG{l+m}{12345}\PYG{+w}{ }\PYG{l+m}{298.0}\PYG{+w}{ }\PYG{k}{region }\PYG{n+nv+nvIdentifier}{my\PYGZus{}swap\PYGZus{}region}\PYG{+w}{ }\PYG{n}{types}\PYG{+w}{ }\PYG{l+m}{5}\PYG{+w}{ }\PYG{l+m}{6}
\PYG{k}{fix }\PYG{n+nv+nvIdentifier}{SGMC}\PYG{+w}{ }\PYG{n+nv+nvIdentifier}{all}\PYG{+w}{ }\PYG{n}{atom}\PYG{o}{/}\PYG{n}{swap}\PYG{+w}{ }\PYG{l+m}{1}\PYG{+w}{ }\PYG{l+m}{100}\PYG{+w}{ }\PYG{l+m}{345}\PYG{+w}{ }\PYG{l+m}{1.0}\PYG{+w}{ }\PYG{n}{semi}\PYG{o}{\PYGZhy{}}\PYG{n}{grand}\PYG{+w}{ }\PYG{n}{yes}\PYG{+w}{ }\PYG{n}{types}\PYG{+w}{ }\PYG{l+m}{1}\PYG{+w}{ }\PYG{l+m}{2}\PYG{+w}{ }\PYG{l+m}{3}\PYG{+w}{ }\PYG{n}{mu}\PYG{+w}{ }\PYG{l+m}{0.0}\PYG{+w}{ }\PYG{l+m}{4.3}\PYG{+w}{ }\PYG{o}{\PYGZhy{}}\PYG{l+m}{5.0}
\end{sphinxVerbatim}


\subsection{Description}
\label{\detokenize{fix_atom_swap:description}}
\sphinxAtStartPar
This fix performs Monte Carlo swaps of atoms of one given atom type with
atoms of the other given atom types.  The specified scaling temperature
\sphinxstyleemphasis{T} is used in the Metropolis criterion dictating swap probabilities.

\sphinxAtStartPar
Perform \sphinxstyleemphasis{X} swaps of atoms of one type with atoms of another type
according to a Monte Carlo probability. Swap candidates must be in the
fix group, must be in the region (if specified), and must be of one of
the listed types. Swaps are attempted between candidates that are chosen
randomly with equal probability among the candidate atoms. Swaps are not
attempted between atoms of the same type since nothing would happen.

\sphinxAtStartPar
All atoms in the simulation domain can be moved using regular time
integration displacements (e.g., via {\hyperref[\detokenize{fix_nh::doc}]{\sphinxcrossref{\DUrole{doc}{fix nvt}}}}), resulting
in a hybrid MC+MD simulation. A smaller\sphinxhyphen{}than\sphinxhyphen{}usual timestep size may be
needed when running such a hybrid simulation, especially if the swapped
atoms are not well equilibrated.

\sphinxAtStartPar
The \sphinxstyleemphasis{types} keyword is required. At least two atom types must be
specified. If not using \sphinxstyleemphasis{semi\sphinxhyphen{}grand}, exactly two atom types are
required.

\sphinxAtStartPar
The \sphinxstyleemphasis{ke} keyword can be set to \sphinxstyleemphasis{no} to turn off kinetic energy
conservation for swaps. The default is \sphinxstyleemphasis{yes}, which means that swapped
atoms have their velocities scaled by the ratio of the masses of the
swapped atom types. This ensures that the kinetic energy of each atom is
the same after the swap as it was before the swap, even though the atom
masses have changed.

\sphinxAtStartPar
The \sphinxstyleemphasis{semi\sphinxhyphen{}grand} keyword can be set to \sphinxstyleemphasis{yes} to switch to the semi\sphinxhyphen{}grand
canonical ensemble as discussed in {\hyperref[\detokenize{fix_atom_swap:sadigh}]{\sphinxcrossref{\DUrole{std,std-ref}{(Sadigh)}}}}. This means
that the total number of each particle type does not need to be
conserved. The default is \sphinxstyleemphasis{no}, which means that the only kind of swap
allowed exchanges an atom of one type with an atom of a different given
type. In other words, the relative mole fractions of the swapped atoms
remains constant. Whereas in the semi\sphinxhyphen{}grand canonical ensemble, the
composition of the system can change. Note that when using \sphinxstyleemphasis{semi\sphinxhyphen{}grand},
atoms in the fix group whose type is not listed in the \sphinxstyleemphasis{types} keyword
are ineligible for attempted conversion. An attempt is made to switch
the selected atom (if eligible) to one of the other listed types with
equal probability.  Acceptance of each attempt depends upon the
Metropolis criterion.

\sphinxAtStartPar
The \sphinxstyleemphasis{mu} keyword allows users to specify chemical potentials. This is
required and allowed only when using \sphinxstyleemphasis{semi\sphinxhyphen{}grand}.  All chemical
potentials are absolute, so there is one for each swap type listed
following the \sphinxstyleemphasis{types} keyword.  In semi\sphinxhyphen{}grand canonical ensemble
simulations the chemical composition of the system is controlled by the
difference in these values. So shifting all values by a constant amount
will have no effect on the simulation.

\sphinxAtStartPar
This command may optionally use the \sphinxstyleemphasis{region} keyword to define swap
volume.  The specified region must have been previously defined with a
{\hyperref[\detokenize{region::doc}]{\sphinxcrossref{\DUrole{doc}{region}}}} command.  It must be defined with side = \sphinxstyleemphasis{in}.
Swap attempts occur only between atoms that are both within the
specified region. Swaps are not otherwise attempted.

\sphinxAtStartPar
You should ensure you do not swap atoms belonging to a molecule, or
LAMMPS will eventually generate an error when it tries to find those
atoms.  LAMMPS will warn you if any of the atoms eligible for swapping
have a non\sphinxhyphen{}zero molecule ID, but does not check for this at the time of
swapping.

\sphinxAtStartPar
If not using \sphinxstyleemphasis{semi\sphinxhyphen{}grand} this fix checks to ensure all atoms of the
given types have the same atomic charge. LAMMPS does not enforce this in
general, but it is needed for this fix to simplify the swapping
procedure. Successful swaps will swap the atom type and charge of the
swapped atoms. Conversely, when using \sphinxstyleemphasis{semi\sphinxhyphen{}grand}, it is assumed that
all the atom types involved in switches have the same charge. Otherwise,
charge would not be conserved. As a consequence, no checks on atomic
charges are performed, and successful switches update the atom type but
not the atom charge. While it is possible to use \sphinxstyleemphasis{semi\sphinxhyphen{}grand} with
groups of atoms that have different charges, these charges will not be
changed when the atom types change.

\sphinxAtStartPar
Since this fix computes total potential energies before and after
proposed swaps, so even complicated potential energy calculations are
OK, including the following:
\begin{itemize}
\item {} 
\sphinxAtStartPar
long\sphinxhyphen{}range electrostatics (\(k\)\sphinxhyphen{}space)

\item {} 
\sphinxAtStartPar
many body pair styles

\item {} 
\sphinxAtStartPar
hybrid pair styles

\item {} 
\sphinxAtStartPar
eam pair styles

\item {} 
\sphinxAtStartPar
triclinic systems

\item {} 
\sphinxAtStartPar
need to include potential energy contributions from other fixes

\end{itemize}

\sphinxAtStartPar
Some fixes have an associated potential energy. Examples of such fixes
include: {\hyperref[\detokenize{fix_efield::doc}]{\sphinxcrossref{\DUrole{doc}{efield}}}}, {\hyperref[\detokenize{fix_gravity::doc}]{\sphinxcrossref{\DUrole{doc}{gravity}}}},
{\hyperref[\detokenize{fix_addforce::doc}]{\sphinxcrossref{\DUrole{doc}{addforce}}}}, {\hyperref[\detokenize{fix_langevin::doc}]{\sphinxcrossref{\DUrole{doc}{langevin}}}},
{\hyperref[\detokenize{fix_restrain::doc}]{\sphinxcrossref{\DUrole{doc}{restrain}}}}, {\hyperref[\detokenize{fix_temp_berendsen::doc}]{\sphinxcrossref{\DUrole{doc}{temp/berendsen}}}}, {\hyperref[\detokenize{fix_temp_rescale::doc}]{\sphinxcrossref{\DUrole{doc}{temp/rescale}}}}, and
{\hyperref[\detokenize{fix_wall::doc}]{\sphinxcrossref{\DUrole{doc}{wall fixes}}}}.  For that energy to be included in the
total potential energy of the system (the quantity used when
performing GCMC moves), you \sphinxstylestrong{must} enable the {\hyperref[\detokenize{fix_modify::doc}]{\sphinxcrossref{\DUrole{doc}{fix\_modify}}}} \sphinxstyleemphasis{energy} option for that fix.  The doc pages for
individual {\hyperref[\detokenize{fix::doc}]{\sphinxcrossref{\DUrole{doc}{fix}}}} commands specify if this should be done.


\subsection{Restart, fix\_modify, output, run start/stop, minimize info}
\label{\detokenize{fix_atom_swap:restart-fix-modify-output-run-start-stop-minimize-info}}
\sphinxAtStartPar
This fix writes the state of the fix to {\hyperref[\detokenize{restart::doc}]{\sphinxcrossref{\DUrole{doc}{binary restart files}}}}.  This includes information about the random number generator
seed, the next timestep for MC exchanges, the number of exchange
attempts and successes, etc.  See the {\hyperref[\detokenize{read_restart::doc}]{\sphinxcrossref{\DUrole{doc}{read\_restart}}}}
command for info on how to re\sphinxhyphen{}specify a fix in an input script that
reads a restart file, so that the operation of the fix continues in an
uninterrupted fashion.

\begin{sphinxadmonition}{note}{Note:}
\sphinxAtStartPar
For this to work correctly, the timestep must \sphinxstylestrong{not} be changed
after reading the restart with {\hyperref[\detokenize{reset_timestep::doc}]{\sphinxcrossref{\DUrole{doc}{reset\_timestep}}}}.  The fix will try to detect it and stop with an
error.
\end{sphinxadmonition}

\sphinxAtStartPar
None of the {\hyperref[\detokenize{fix_modify::doc}]{\sphinxcrossref{\DUrole{doc}{fix\_modify}}}} options are relevant to this
fix.

\sphinxAtStartPar
This fix computes a global vector of length 2, which can be accessed
by various {\hyperref[\detokenize{Howto_output::doc}]{\sphinxcrossref{\DUrole{doc}{output commands}}}}.  The vector values are
the following global cumulative quantities:
\begin{enumerate}
\sphinxsetlistlabels{\arabic}{enumi}{enumii}{}{.}%
\item {} 
\sphinxAtStartPar
swap attempts

\item {} 
\sphinxAtStartPar
swap accepts

\end{enumerate}

\sphinxAtStartPar
The vector values calculated by this fix are “intensive”.

\sphinxAtStartPar
No parameter of this fix can be used with the \sphinxstyleemphasis{start/stop} keywords of
the {\hyperref[\detokenize{run::doc}]{\sphinxcrossref{\DUrole{doc}{run}}}} command.  This fix is not invoked during
{\hyperref[\detokenize{minimize::doc}]{\sphinxcrossref{\DUrole{doc}{energy minimization}}}}.


\subsection{Restrictions}
\label{\detokenize{fix_atom_swap:restrictions}}
\sphinxAtStartPar
This fix is part of the MC package.  It is only enabled if LAMMPS was
built with that package.  See the {\hyperref[\detokenize{Build_package::doc}]{\sphinxcrossref{\DUrole{doc}{Build package}}}}
doc page for more info.

\sphinxAtStartPar
This fix cannot be used with systems that do not have per\sphinxhyphen{}type masses
(e.g. atom style sphere) since the implemented algorithm pre\sphinxhyphen{}computes
velocity rescaling factors from per\sphinxhyphen{}type masses and ignores any per\sphinxhyphen{}atom
masses, if present.  In case both, per\sphinxhyphen{}type and per\sphinxhyphen{}atom masses are
present, a warning is printed.


\subsection{Related commands}
\label{\detokenize{fix_atom_swap:related-commands}}
\sphinxAtStartPar
{\hyperref[\detokenize{fix_nh::doc}]{\sphinxcrossref{\DUrole{doc}{fix nvt}}}}, {\hyperref[\detokenize{neighbor::doc}]{\sphinxcrossref{\DUrole{doc}{neighbor}}}},
{\hyperref[\detokenize{fix_deposit::doc}]{\sphinxcrossref{\DUrole{doc}{fix deposit}}}}, {\hyperref[\detokenize{fix_evaporate::doc}]{\sphinxcrossref{\DUrole{doc}{fix evaporate}}}},
{\hyperref[\detokenize{delete_atoms::doc}]{\sphinxcrossref{\DUrole{doc}{delete\_atoms}}}}, {\hyperref[\detokenize{fix_gcmc::doc}]{\sphinxcrossref{\DUrole{doc}{fix gcmc}}}},
{\hyperref[\detokenize{fix_mol_swap::doc}]{\sphinxcrossref{\DUrole{doc}{fix mol/swap}}}}, {\hyperref[\detokenize{fix_sgcmc::doc}]{\sphinxcrossref{\DUrole{doc}{fix sgcmc}}}}


\subsection{Default}
\label{\detokenize{fix_atom_swap:default}}
\sphinxAtStartPar
The option defaults are \sphinxstyleemphasis{ke} = yes, \sphinxstyleemphasis{semi\sphinxhyphen{}grand} = no, \sphinxstyleemphasis{mu} = 0.0 for
all atom types.


\bigskip\hrule\bigskip

\phantomsection\label{\detokenize{fix_atom_swap:sadigh}}
\sphinxAtStartPar
\sphinxstylestrong{(Sadigh)} B Sadigh, P Erhart, A Stukowski, A Caro, E Martinez, and
L Zepeda\sphinxhyphen{}Ruiz, Phys. Rev. B, 85, 184203 (2012).

\sphinxstepscope

\index{fix ave/atom@\spxentry{fix ave/atom}}\ignorespaces 

\section{fix ave/atom command}
\label{\detokenize{fix_ave_atom:fix-ave-atom-command}}\label{\detokenize{fix_ave_atom:index-0}}\label{\detokenize{fix_ave_atom::doc}}

\subsection{Syntax}
\label{\detokenize{fix_ave_atom:syntax}}
\begin{sphinxVerbatim}[commandchars=\\\{\}]
\PYG{k}{fix }\PYG{n+nv+nvIdentifier}{ID}\PYG{+w}{ }\PYG{n+nv+nvIdentifier}{group\PYGZhy{}ID}\PYG{+w}{ }\PYG{n}{ave}\PYG{o}{/}\PYG{n}{atom}\PYG{+w}{ }\PYG{n}{Nevery}\PYG{+w}{ }\PYG{n}{Nrepeat}\PYG{+w}{ }\PYG{n}{Nfreq}\PYG{+w}{ }\PYG{n}{value1}\PYG{+w}{ }\PYG{n}{value2}\PYG{+w}{ }\PYG{n}{...}
\end{sphinxVerbatim}
\begin{itemize}
\item {} 
\sphinxAtStartPar
ID, group\sphinxhyphen{}ID are documented in {\hyperref[\detokenize{fix::doc}]{\sphinxcrossref{\DUrole{doc}{fix}}}} command

\item {} 
\sphinxAtStartPar
ave/atom = style name of this fix command

\item {} 
\sphinxAtStartPar
Nevery = use input values every this many timesteps

\item {} 
\sphinxAtStartPar
Nrepeat = \# of times to use input values for calculating averages

\item {} 
\sphinxAtStartPar
Nfreq = calculate averages every this many timesteps

\item {} 
\sphinxAtStartPar
one or more input values can be listed

\item {} 
\sphinxAtStartPar
value = \sphinxstyleemphasis{x}, \sphinxstyleemphasis{y}, \sphinxstyleemphasis{z}, \sphinxstyleemphasis{vx}, \sphinxstyleemphasis{vy}, \sphinxstyleemphasis{vz}, \sphinxstyleemphasis{fx}, \sphinxstyleemphasis{fy}, \sphinxstyleemphasis{fz}, c\_ID, c\_ID{[}i{]}, f\_ID, f\_ID{[}i{]}, v\_name

\begin{sphinxVerbatim}[commandchars=\\\{\}]
x,y,z,vx,vy,vz,fx,fy,fz = atom attribute (position, velocity, force component)
c\PYGZus{}ID = per\PYGZhy{}atom vector calculated by a compute with ID
c\PYGZus{}ID[I] = Ith column of per\PYGZhy{}atom array calculated by a compute with ID, I can include wildcard (see below)
f\PYGZus{}ID = per\PYGZhy{}atom vector calculated by a fix with ID
f\PYGZus{}ID[I] = Ith column of per\PYGZhy{}atom array calculated by a fix with ID, I can include wildcard (see below)
v\PYGZus{}name = per\PYGZhy{}atom vector calculated by an atom\PYGZhy{}style variable with name
\end{sphinxVerbatim}

\end{itemize}


\subsection{Examples}
\label{\detokenize{fix_ave_atom:examples}}
\begin{sphinxVerbatim}[commandchars=\\\{\}]
\PYG{k}{fix }\PYG{n+nv+nvIdentifier}{1}\PYG{+w}{ }\PYG{n+nv+nvIdentifier}{all}\PYG{+w}{ }\PYG{n}{ave}\PYG{o}{/}\PYG{n}{atom}\PYG{+w}{ }\PYG{l+m}{1}\PYG{+w}{ }\PYG{l+m}{100}\PYG{+w}{ }\PYG{l+m}{100}\PYG{+w}{ }\PYG{n}{vx}\PYG{+w}{ }\PYG{n}{vy}\PYG{+w}{ }\PYG{n}{vz}
\PYG{k}{fix }\PYG{n+nv+nvIdentifier}{1}\PYG{+w}{ }\PYG{n+nv+nvIdentifier}{all}\PYG{+w}{ }\PYG{n}{ave}\PYG{o}{/}\PYG{n}{atom}\PYG{+w}{ }\PYG{l+m}{10}\PYG{+w}{ }\PYG{l+m}{20}\PYG{+w}{ }\PYG{l+m}{1000}\PYG{+w}{ }\PYG{n}{c\PYGZus{}my\PYGZus{}stress[1]}
\PYG{k}{fix }\PYG{n+nv+nvIdentifier}{1}\PYG{+w}{ }\PYG{n+nv+nvIdentifier}{all}\PYG{+w}{ }\PYG{n}{ave}\PYG{o}{/}\PYG{n}{atom}\PYG{+w}{ }\PYG{l+m}{10}\PYG{+w}{ }\PYG{l+m}{20}\PYG{+w}{ }\PYG{l+m}{1000}\PYG{+w}{ }\PYG{n}{c\PYGZus{}my\PYGZus{}stress[}\PYG{o}{*}\PYG{n}{]}
\end{sphinxVerbatim}


\subsection{Description}
\label{\detokenize{fix_ave_atom:description}}
\sphinxAtStartPar
Use one or more per\sphinxhyphen{}atom vectors as inputs every few timesteps, and
average them atom by atom over longer timescales.  The resulting
per\sphinxhyphen{}atom averages can be used by other {\hyperref[\detokenize{Howto_output::doc}]{\sphinxcrossref{\DUrole{doc}{output commands}}}}
such as the {\hyperref[\detokenize{fix_ave_chunk::doc}]{\sphinxcrossref{\DUrole{doc}{fix ave/chunk}}}} or {\hyperref[\detokenize{dump::doc}]{\sphinxcrossref{\DUrole{doc}{dump custom}}}}
commands.

\sphinxAtStartPar
The group specified with the command means only atoms within the group
have their averages computed.  Results are set to 0.0 for atoms not in
the group.

\sphinxAtStartPar
Each input value can be an atom attribute (position, velocity, force
component) or can be the result of a {\hyperref[\detokenize{compute::doc}]{\sphinxcrossref{\DUrole{doc}{compute}}}} or
{\hyperref[\detokenize{fix::doc}]{\sphinxcrossref{\DUrole{doc}{fix}}}} or the evaluation of an atom\sphinxhyphen{}style
{\hyperref[\detokenize{variable::doc}]{\sphinxcrossref{\DUrole{doc}{variable}}}}.  In the latter cases, the compute, fix, or
variable must produce a per\sphinxhyphen{}atom vector, not a global quantity or
local quantity.  If you wish to time\sphinxhyphen{}average global quantities from a
compute, fix, or variable, then see the {\hyperref[\detokenize{fix_ave_time::doc}]{\sphinxcrossref{\DUrole{doc}{fix ave/time}}}}
command.

\sphinxAtStartPar
Each per\sphinxhyphen{}atom value of each input vector is averaged independently.

\sphinxAtStartPar
{\hyperref[\detokenize{compute::doc}]{\sphinxcrossref{\DUrole{doc}{Computes}}}} that produce per\sphinxhyphen{}atom vectors or arrays are
those which have the word \sphinxstyleemphasis{atom} in their style name.  See the doc
pages for individual {\hyperref[\detokenize{fix::doc}]{\sphinxcrossref{\DUrole{doc}{fixes}}}} to determine which ones produce
per\sphinxhyphen{}atom vectors or arrays.  {\hyperref[\detokenize{variable::doc}]{\sphinxcrossref{\DUrole{doc}{Variables}}}} of style \sphinxstyleemphasis{atom}
are the only ones that can be used with this fix since they produce
per\sphinxhyphen{}atom vectors.

\sphinxAtStartPar
Note that for values from a compute or fix, the bracketed index I can
be specified using a wildcard asterisk with the index to effectively
specify multiple values.  This takes the form “*” or “*n” or “m*” or
“m*n”.  If \(N\) is the size of the vector (for \sphinxstyleemphasis{mode} = scalar) or the
number of columns in the array (for \sphinxstyleemphasis{mode} = vector), then an asterisk
with no numeric values means all indices from 1 to \(N\).  A leading
asterisk means all indices from 1 to n (inclusive).  A trailing
asterisk means all indices from m to \(N\) (inclusive).  A middle asterisk
means all indices from m to n (inclusive).

\sphinxAtStartPar
Using a wildcard is the same as if the individual columns of the array
had been listed one by one.  For example, these two fix ave/atom commands are
equivalent, since the {\hyperref[\detokenize{compute_stress_atom::doc}]{\sphinxcrossref{\DUrole{doc}{compute stress/atom}}}}
command creates a per\sphinxhyphen{}atom array with six columns:

\begin{sphinxVerbatim}[commandchars=\\\{\}]
\PYG{k}{compute }\PYG{n+nv+nvIdentifier}{my\PYGZus{}stress}\PYG{+w}{ }\PYG{n+nv+nvIdentifier}{all}\PYG{+w}{ }\PYG{n}{stress}\PYG{o}{/}\PYG{n}{atom}\PYG{+w}{ }\PYG{n}{NULL}
\PYG{k}{fix }\PYG{n+nv+nvIdentifier}{1}\PYG{+w}{ }\PYG{n+nv+nvIdentifier}{all}\PYG{+w}{ }\PYG{n}{ave}\PYG{o}{/}\PYG{n}{atom}\PYG{+w}{ }\PYG{l+m}{10}\PYG{+w}{ }\PYG{l+m}{20}\PYG{+w}{ }\PYG{l+m}{1000}\PYG{+w}{ }\PYG{n}{c\PYGZus{}my\PYGZus{}stress[}\PYG{o}{*}\PYG{n}{]}
\PYG{k}{fix }\PYG{n+nv+nvIdentifier}{1}\PYG{+w}{ }\PYG{n+nv+nvIdentifier}{all}\PYG{+w}{ }\PYG{n}{ave}\PYG{o}{/}\PYG{n}{atom}\PYG{+w}{ }\PYG{l+m}{10}\PYG{+w}{ }\PYG{l+m}{20}\PYG{+w}{ }\PYG{l+m}{1000}\PYG{+w}{ }\PYG{n}{c\PYGZus{}my\PYGZus{}stress[1]}\PYG{+w}{ }\PYG{n}{c\PYGZus{}my\PYGZus{}stress[2]}\PYG{+w}{ }\PYG{o}{\PYGZam{}}
\PYG{+w}{                              }\PYG{n}{c\PYGZus{}my\PYGZus{}stress[3]}\PYG{+w}{ }\PYG{n}{c\PYGZus{}my\PYGZus{}stress[4]}\PYG{+w}{ }\PYG{o}{\PYGZam{}}
\PYG{+w}{                              }\PYG{n}{c\PYGZus{}my\PYGZus{}stress[5]}\PYG{+w}{ }\PYG{n}{c\PYGZus{}my\PYGZus{}stress[6]}
\end{sphinxVerbatim}


\bigskip\hrule\bigskip


\sphinxAtStartPar
The \(N_\text{every}\), \(N_\text{repeat}\), and \(N_\text{freq}\)
arguments specify on what timesteps the input values will be used in order to
contribute to the average.  The final averaged quantities are generated on
timesteps that are a multiple of \(N_\text{freq}\).  The average is over
\(N_\text{repeat}\) quantities, computed in the preceding portion of the
simulation every \(N_\text{every}\) timesteps.  \(N_\text{freq}\) must
be a multiple of \(N_\text{every}\) and \(N_\text{every}\) must be
non\sphinxhyphen{}zero even if \(N_\text{repeat}\) is 1.  Also, the timesteps
contributing to the average value cannot overlap; that is,
\(N_\text{repeat} \times N_\text{every}\) cannot exceed \(N_\text{freq}\).

\sphinxAtStartPar
For example, if \(N_\text{every}=2\), \(N_\text{repeat}=6\), and
\(N_\text{freq}=100\), then values on timesteps 90, 92, 94, 96, 98, and 100
will be used to compute the final average on time step 100.  Similarly for
timesteps 190, 192, 194, 196, 198, and 200 on time step 200, etc.


\bigskip\hrule\bigskip


\sphinxAtStartPar
The atom attribute values (\sphinxstyleemphasis{x}, \sphinxstyleemphasis{y}, \sphinxstyleemphasis{z}, \sphinxstyleemphasis{vx}, \sphinxstyleemphasis{vy}, \sphinxstyleemphasis{vz}, \sphinxstyleemphasis{fx}, \sphinxstyleemphasis{fy}, and
\sphinxstyleemphasis{fz}) are self\sphinxhyphen{}explanatory.  Note that other atom attributes can be used as
inputs to this fix by using the
{\hyperref[\detokenize{compute_property_atom::doc}]{\sphinxcrossref{\DUrole{doc}{compute property/atom}}}} command and then
specifying an input value from that compute.

\begin{sphinxadmonition}{note}{Note:}
\sphinxAtStartPar
The \sphinxstyleemphasis{x}, \sphinxstyleemphasis{y}, and \sphinxstyleemphasis{z} attributes are values that are re\sphinxhyphen{}wrapped inside
the periodic box whenever an atom crosses a periodic boundary.  Thus, if
you time\sphinxhyphen{}average an atom that spends half of its time on either side of
the periodic box, you will get a value in the middle of the box.  If
this is not what you want, consider averaging unwrapped coordinates,
which can be provided by the
{\hyperref[\detokenize{compute_property_atom::doc}]{\sphinxcrossref{\DUrole{doc}{compute property/atom}}}}
command via its \sphinxstyleemphasis{xu}, \sphinxstyleemphasis{yu}, and \sphinxstyleemphasis{zu} attributes.
\end{sphinxadmonition}

\sphinxAtStartPar
If a value begins with “c\_”, a compute ID must follow which has been
previously defined in the input script.  If no bracketed term is
appended, the per\sphinxhyphen{}atom vector calculated by the compute is used.  If a
bracketed term containing an index \(I\) is appended, the
\(I^\text{th}\) column of the per\sphinxhyphen{}atom array calculated by the compute is
used.  Users can also write code for their own compute styles and
{\hyperref[\detokenize{Modify::doc}]{\sphinxcrossref{\DUrole{doc}{add them to LAMMPS}}}}.  See the discussion above for how
\(I\) can be specified with a wildcard asterisk to effectively specify
multiple values.

\sphinxAtStartPar
If a value begins with “f\_”, a fix ID must follow which has been previously
defined in the input script.  If no bracketed term is appended, the per\sphinxhyphen{}atom
vector calculated by the fix is used.  If a bracketed term containing an index
\(I\) is appended, the \(I^\text{th}\) column of the per\sphinxhyphen{}atom array
calculated by the fix is used.  Note that some fixes only produce their values
on certain timesteps, which must be compatible with \(N_\text{every}\),
else an error will result.  Users can also write code for their own fix styles
and {\hyperref[\detokenize{Modify::doc}]{\sphinxcrossref{\DUrole{doc}{add them to LAMMPS}}}}.  See the discussion above for how
\(I\) can be specified with a wildcard asterisk to effectively specify
multiple values.

\sphinxAtStartPar
If a value begins with “v\_”, a variable name must follow which has
been previously defined in the input script as an
{\hyperref[\detokenize{variable::doc}]{\sphinxcrossref{\DUrole{doc}{atom\sphinxhyphen{}style variable}}}}. Variables of style \sphinxstyleemphasis{atom} can reference
thermodynamic keywords or invoke other computes, fixes, or variables
when they are evaluated, so this is a very general means of generating
per\sphinxhyphen{}atom quantities to time average.


\bigskip\hrule\bigskip



\subsection{Restart, fix\_modify, output, run start/stop, minimize info}
\label{\detokenize{fix_ave_atom:restart-fix-modify-output-run-start-stop-minimize-info}}
\sphinxAtStartPar
No information about this fix is written to
{\hyperref[\detokenize{restart::doc}]{\sphinxcrossref{\DUrole{doc}{binary restart files}}}}.  None of the
{\hyperref[\detokenize{fix_modify::doc}]{\sphinxcrossref{\DUrole{doc}{fix\_modify}}}} options are relevant to this fix.
No global scalar or vector quantities are stored by this fix for access by
various {\hyperref[\detokenize{Howto_output::doc}]{\sphinxcrossref{\DUrole{doc}{output commands}}}}.

\sphinxAtStartPar
This fix produces a per\sphinxhyphen{}atom vector or array which can be accessed by
various {\hyperref[\detokenize{Howto_output::doc}]{\sphinxcrossref{\DUrole{doc}{output commands}}}}.  A vector is produced if
only a single quantity is averaged by this fix.  If two or more
quantities are averaged, then an array of values is produced.  The
per\sphinxhyphen{}atom values can only be accessed on timesteps that are multiples
of \(N_\text{freq}\) since that is when averaging is performed.

\sphinxAtStartPar
No parameter of this fix can be used with the \sphinxstyleemphasis{start/stop} keywords of
the {\hyperref[\detokenize{run::doc}]{\sphinxcrossref{\DUrole{doc}{run}}}} command.  This fix is not invoked during
{\hyperref[\detokenize{minimize::doc}]{\sphinxcrossref{\DUrole{doc}{energy minimization}}}}.


\subsection{Restrictions}
\label{\detokenize{fix_ave_atom:restrictions}}\begin{quote}

\sphinxAtStartPar
none
\end{quote}


\subsection{Related commands}
\label{\detokenize{fix_ave_atom:related-commands}}
\sphinxAtStartPar
{\hyperref[\detokenize{compute::doc}]{\sphinxcrossref{\DUrole{doc}{compute}}}}, {\hyperref[\detokenize{fix_ave_histo::doc}]{\sphinxcrossref{\DUrole{doc}{fix ave/histo}}}},
{\hyperref[\detokenize{fix_ave_chunk::doc}]{\sphinxcrossref{\DUrole{doc}{fix ave/chunk}}}}, {\hyperref[\detokenize{fix_ave_time::doc}]{\sphinxcrossref{\DUrole{doc}{fix ave/time}}}},
{\hyperref[\detokenize{variable::doc}]{\sphinxcrossref{\DUrole{doc}{variable}}}},


\subsection{Default}
\label{\detokenize{fix_ave_atom:default}}
\sphinxAtStartPar
none

\sphinxstepscope

\index{fix ave/chunk@\spxentry{fix ave/chunk}}\ignorespaces 

\section{fix ave/chunk command}
\label{\detokenize{fix_ave_chunk:fix-ave-chunk-command}}\label{\detokenize{fix_ave_chunk:index-0}}\label{\detokenize{fix_ave_chunk::doc}}

\subsection{Syntax}
\label{\detokenize{fix_ave_chunk:syntax}}
\begin{sphinxVerbatim}[commandchars=\\\{\}]
\PYG{k}{fix }\PYG{n+nv+nvIdentifier}{ID}\PYG{+w}{ }\PYG{n+nv+nvIdentifier}{group\PYGZhy{}ID}\PYG{+w}{ }\PYG{n}{ave}\PYG{o}{/}\PYG{n}{chunk}\PYG{+w}{ }\PYG{n}{Nevery}\PYG{+w}{ }\PYG{n}{Nrepeat}\PYG{+w}{ }\PYG{n}{Nfreq}\PYG{+w}{ }\PYG{n}{chunkID}\PYG{+w}{ }\PYG{n}{value1}\PYG{+w}{ }\PYG{n}{value2}\PYG{+w}{ }\PYG{n}{...}\PYG{+w}{ }\PYG{n}{keyword}\PYG{+w}{ }\PYG{n}{args}\PYG{+w}{ }\PYG{n}{...}
\end{sphinxVerbatim}
\begin{itemize}
\item {} 
\sphinxAtStartPar
ID, group\sphinxhyphen{}ID are documented in {\hyperref[\detokenize{fix::doc}]{\sphinxcrossref{\DUrole{doc}{fix}}}} command

\item {} 
\sphinxAtStartPar
ave/chunk = style name of this fix command

\item {} 
\sphinxAtStartPar
Nevery = use input values every this many timesteps

\item {} 
\sphinxAtStartPar
Nrepeat = \# of times to use input values for calculating averages

\item {} 
\sphinxAtStartPar
Nfreq = calculate averages every this many timesteps

\item {} 
\sphinxAtStartPar
chunkID = ID of {\hyperref[\detokenize{compute_chunk_atom::doc}]{\sphinxcrossref{\DUrole{doc}{compute chunk/atom}}}} command

\item {} 
\sphinxAtStartPar
one or more input values can be listed

\item {} 
\sphinxAtStartPar
value = \sphinxstyleemphasis{vx}, \sphinxstyleemphasis{vy}, \sphinxstyleemphasis{vz}, \sphinxstyleemphasis{fx}, \sphinxstyleemphasis{fy}, \sphinxstyleemphasis{fz}, \sphinxstyleemphasis{density/mass}, \sphinxstyleemphasis{density/number}, \sphinxstyleemphasis{mass}, \sphinxstyleemphasis{temp}, c\_ID, c\_ID{[}I{]}, f\_ID, f\_ID{[}I{]}, v\_name

\begin{sphinxVerbatim}[commandchars=\\\{\}]
vx,vy,vz,fx,fy,fz,mass = atom attribute (velocity, force component, mass)
density/number, density/mass = number or mass density (per volume)
temp = temperature
c\PYGZus{}ID = per\PYGZhy{}atom vector calculated by a compute with ID
c\PYGZus{}ID[I] = Ith column of per\PYGZhy{}atom array calculated by a compute with ID, I can include wildcard (see below)
f\PYGZus{}ID = per\PYGZhy{}atom vector calculated by a fix with ID
f\PYGZus{}ID[I] = Ith column of per\PYGZhy{}atom array calculated by a fix with ID, I can include wildcard (see below)
v\PYGZus{}name = per\PYGZhy{}atom vector calculated by an atom\PYGZhy{}style variable with name
\end{sphinxVerbatim}

\item {} 
\sphinxAtStartPar
zero or more keyword/arg pairs may be appended

\item {} 
\sphinxAtStartPar
keyword = \sphinxstyleemphasis{norm} or \sphinxstyleemphasis{ave} or \sphinxstyleemphasis{bias} or \sphinxstyleemphasis{adof} or \sphinxstyleemphasis{cdof} or \sphinxstyleemphasis{file} or \sphinxstyleemphasis{append} or \sphinxstyleemphasis{overwrite} or \sphinxstyleemphasis{format} or \sphinxstyleemphasis{title1} or \sphinxstyleemphasis{title2} or \sphinxstyleemphasis{title3}
\begin{sphinxalltt}
\sphinxstyleemphasis{norm} arg = \sphinxstyleemphasis{all} or \sphinxstyleemphasis{sample} or \sphinxstyleemphasis{none} = how output on \sphinxstyleemphasis{Nfreq} steps is normalized
  all = output is sum of atoms across all \sphinxstyleemphasis{Nrepeat} samples, divided by atom count
  sample = output is sum of \sphinxstyleemphasis{Nrepeat} sample averages, divided by \sphinxstyleemphasis{Nrepeat}
  none = output is sum of \sphinxstyleemphasis{Nrepeat} sample sums, divided by \sphinxstyleemphasis{Nrepeat}
\sphinxstyleemphasis{ave} args = \sphinxstyleemphasis{one} or \sphinxstyleemphasis{running} or \sphinxstyleemphasis{window M}
  one = output new average value every Nfreq steps
  running = output cumulative average of all previous Nfreq steps
  window M = output average of M most recent Nfreq steps
\sphinxstyleemphasis{bias} arg = bias\sphinxhyphen{}ID
  bias\sphinxhyphen{}ID = ID of a temperature compute that removes a velocity bias for temperature calculation
\sphinxstyleemphasis{adof} value = dof\_per\_atom
  dof\_per\_atom = define this many degrees\sphinxhyphen{}of\sphinxhyphen{}freedom per atom for temperature calculation
\sphinxstyleemphasis{cdof} value = dof\_per\_chunk
  dof\_per\_chunk = define this many degrees\sphinxhyphen{}of\sphinxhyphen{}freedom per chunk for temperature calculation
\sphinxstyleemphasis{file} arg = filename
  filename = file to write results to
\sphinxstyleemphasis{append} arg = filename
  filename = file to append results to
\sphinxstyleemphasis{overwrite} arg = none = overwrite output file with only latest output
\sphinxstyleemphasis{format} arg = string
  string = C\sphinxhyphen{}style format string
\sphinxstyleemphasis{title1} arg = string
  string = text to print as 1st line of output file
\sphinxstyleemphasis{title2} arg = string
  string = text to print as 2nd line of output file
\sphinxstyleemphasis{title3} arg = string
  string = text to print as 3rd line of output file
\end{sphinxalltt}

\end{itemize}


\subsection{Examples}
\label{\detokenize{fix_ave_chunk:examples}}
\begin{sphinxVerbatim}[commandchars=\\\{\}]
\PYG{k}{fix }\PYG{n+nv+nvIdentifier}{1}\PYG{+w}{ }\PYG{n+nv+nvIdentifier}{all}\PYG{+w}{ }\PYG{n}{ave}\PYG{o}{/}\PYG{n}{chunk}\PYG{+w}{ }\PYG{l+m}{10000}\PYG{+w}{ }\PYG{l+m}{1}\PYG{+w}{ }\PYG{l+m}{10000}\PYG{+w}{ }\PYG{n}{binchunk}\PYG{+w}{ }\PYG{n}{c\PYGZus{}myCentro}\PYG{+w}{ }\PYG{n}{title1}\PYG{+w}{ }\PYG{l+s}{\PYGZdq{}}\PYG{l+s}{My output values}\PYG{l+s}{\PYGZdq{}}
\PYG{k}{fix }\PYG{n+nv+nvIdentifier}{1}\PYG{+w}{ }\PYG{n+nv+nvIdentifier}{flow}\PYG{+w}{ }\PYG{n}{ave}\PYG{o}{/}\PYG{n}{chunk}\PYG{+w}{ }\PYG{l+m}{100}\PYG{+w}{ }\PYG{l+m}{10}\PYG{+w}{ }\PYG{l+m}{1000}\PYG{+w}{ }\PYG{n}{molchunk}\PYG{+w}{ }\PYG{n}{vx}\PYG{+w}{ }\PYG{n}{vz}\PYG{+w}{ }\PYG{n}{norm}\PYG{+w}{ }\PYG{n}{sample}\PYG{+w}{ }\PYG{n}{file}\PYG{+w}{ }\PYG{n}{vel.profile}
\PYG{k}{fix }\PYG{n+nv+nvIdentifier}{1}\PYG{+w}{ }\PYG{n+nv+nvIdentifier}{flow}\PYG{+w}{ }\PYG{n}{ave}\PYG{o}{/}\PYG{n}{chunk}\PYG{+w}{ }\PYG{l+m}{100}\PYG{+w}{ }\PYG{l+m}{5}\PYG{+w}{ }\PYG{l+m}{1000}\PYG{+w}{ }\PYG{n}{binchunk}\PYG{+w}{ }\PYG{n}{density}\PYG{o}{/}\PYG{n}{mass}\PYG{+w}{ }\PYG{n}{ave}\PYG{+w}{ }\PYG{n}{running}
\PYG{k}{fix }\PYG{n+nv+nvIdentifier}{1}\PYG{+w}{ }\PYG{n+nv+nvIdentifier}{flow}\PYG{+w}{ }\PYG{n}{ave}\PYG{o}{/}\PYG{n}{chunk}\PYG{+w}{ }\PYG{l+m}{100}\PYG{+w}{ }\PYG{l+m}{5}\PYG{+w}{ }\PYG{l+m}{1000}\PYG{+w}{ }\PYG{n}{binchunk}\PYG{+w}{ }\PYG{n}{density}\PYG{o}{/}\PYG{n}{mass}\PYG{+w}{ }\PYG{n}{ave}\PYG{+w}{ }\PYG{n}{running}
\end{sphinxVerbatim}

\begin{sphinxadmonition}{note}{Note:}
\sphinxAtStartPar
\DUrole{versionmodified,changed}{Changed in version 31May2016.}

\sphinxAtStartPar
If you are trying to replace a deprecated fix ave/spatial command
with the newer, more flexible fix ave/chunk and {\hyperref[\detokenize{compute_chunk_atom::doc}]{\sphinxcrossref{\DUrole{doc}{compute
chunk/atom}}}} commands, you simply need to split
the fix ave/spatial arguments across the two new commands.  For
example, this command:

\begin{sphinxVerbatim}[commandchars=\\\{\}]
\PYG{k}{fix }\PYG{n+nv+nvIdentifier}{1}\PYG{+w}{ }\PYG{n+nv+nvIdentifier}{flow}\PYG{+w}{ }\PYG{n}{ave}\PYG{o}{/}\PYG{n}{spatial}\PYG{+w}{ }\PYG{l+m}{100}\PYG{+w}{ }\PYG{l+m}{10}\PYG{+w}{ }\PYG{l+m}{1000}\PYG{+w}{ }\PYG{n}{y}\PYG{+w}{ }\PYG{l+m}{0.0}\PYG{+w}{ }\PYG{l+m}{1.0}\PYG{+w}{ }\PYG{n}{vx}\PYG{+w}{ }\PYG{n}{vz}\PYG{+w}{ }\PYG{n}{norm}\PYG{+w}{ }\PYG{n}{sample}\PYG{+w}{ }\PYG{n}{file}\PYG{+w}{ }\PYG{n}{vel.profile}
\end{sphinxVerbatim}

\sphinxAtStartPar
could be replaced by:

\begin{sphinxVerbatim}[commandchars=\\\{\}]
\PYG{k}{compute }\PYG{n+nv+nvIdentifier}{cc1}\PYG{+w}{ }\PYG{n+nv+nvIdentifier}{flow}\PYG{+w}{ }\PYG{n}{chunk}\PYG{o}{/}\PYG{n}{atom}\PYG{+w}{ }\PYG{n}{bin}\PYG{o}{/}\PYG{l+m}{1}\PYG{n}{d}\PYG{+w}{ }\PYG{n}{y}\PYG{+w}{ }\PYG{l+m}{0.0}\PYG{+w}{ }\PYG{l+m}{1.0}
\PYG{k}{fix }\PYG{n+nv+nvIdentifier}{1}\PYG{+w}{ }\PYG{n+nv+nvIdentifier}{flow}\PYG{+w}{ }\PYG{n}{ave}\PYG{o}{/}\PYG{n}{chunk}\PYG{+w}{ }\PYG{l+m}{100}\PYG{+w}{ }\PYG{l+m}{10}\PYG{+w}{ }\PYG{l+m}{1000}\PYG{+w}{ }\PYG{n}{cc1}\PYG{+w}{ }\PYG{n}{vx}\PYG{+w}{ }\PYG{n}{vz}\PYG{+w}{ }\PYG{n}{norm}\PYG{+w}{ }\PYG{n}{sample}\PYG{+w}{ }\PYG{n}{file}\PYG{+w}{ }\PYG{n}{vel.profile}
\end{sphinxVerbatim}
\end{sphinxadmonition}


\subsection{Description}
\label{\detokenize{fix_ave_chunk:description}}
\sphinxAtStartPar
Use one or more per\sphinxhyphen{}atom vectors as inputs every few timesteps, sum
the values over the atoms in each chunk at each timestep, then average
the per\sphinxhyphen{}chunk values over longer timescales.  The resulting chunk
averages can be used by other {\hyperref[\detokenize{Howto_output::doc}]{\sphinxcrossref{\DUrole{doc}{output commands}}}} such
as {\hyperref[\detokenize{thermo_style::doc}]{\sphinxcrossref{\DUrole{doc}{thermo\_style custom}}}}, and can also be written to
a file.

\sphinxAtStartPar
In LAMMPS, chunks are collections of atoms defined by a {\hyperref[\detokenize{compute_chunk_atom::doc}]{\sphinxcrossref{\DUrole{doc}{compute
chunk/atom}}}} command, which assigns each atom to a
single chunk (or no chunk).  The ID for this command is specified as
chunkID.  For example, a single chunk could be the atoms in a molecule
or atoms in a spatial bin.  See the {\hyperref[\detokenize{compute_chunk_atom::doc}]{\sphinxcrossref{\DUrole{doc}{compute chunk/atom}}}} page and the {\hyperref[\detokenize{Howto_chunk::doc}]{\sphinxcrossref{\DUrole{doc}{Howto chunk}}}}
page for details of how chunks can be defined and examples of how they
can be used to measure properties of a system.

\sphinxAtStartPar
Note that if the {\hyperref[\detokenize{compute_chunk_atom::doc}]{\sphinxcrossref{\DUrole{doc}{compute chunk/atom}}}}
command defines spatial bins, the fix ave/chunk command performs a
similar computation as the {\hyperref[\detokenize{fix_ave_grid::doc}]{\sphinxcrossref{\DUrole{doc}{fix ave/grid}}}} command.
However, the per\sphinxhyphen{}bin outputs from the fix ave/chunk command are
global; each processor stores a copy of the entire set of bin data.
By contrast, the {\hyperref[\detokenize{fix_ave_grid::doc}]{\sphinxcrossref{\DUrole{doc}{fix ave/grid}}}} command uses a
distributed grid where each processor owns a subset of the bins.  Thus
it is more efficient to use the {\hyperref[\detokenize{fix_ave_grid::doc}]{\sphinxcrossref{\DUrole{doc}{fix ave/grid}}}}
command when the grid is large and a simulation is run on many
processors.

\sphinxAtStartPar
Note that only atoms in the specified group contribute to the summing
and averaging calculations.  The {\hyperref[\detokenize{compute_chunk_atom::doc}]{\sphinxcrossref{\DUrole{doc}{compute chunk/atom}}}} command defines its own group as well as an
optional region.  Atoms will have a chunk ID = 0, meaning they belong
to no chunk, if they are not in that group or region.  Thus you can
specify the “all” group for this command if you simply want to use the
chunk definitions provided by chunkID.

\sphinxAtStartPar
Each specified per\sphinxhyphen{}atom value can be an atom attribute (position,
velocity, force component), a number or mass density, a mass or
temperature, or the result of a {\hyperref[\detokenize{compute::doc}]{\sphinxcrossref{\DUrole{doc}{compute}}}} or {\hyperref[\detokenize{fix::doc}]{\sphinxcrossref{\DUrole{doc}{fix}}}} or the evaluation of an atom\sphinxhyphen{}style {\hyperref[\detokenize{variable::doc}]{\sphinxcrossref{\DUrole{doc}{variable}}}}.
In the latter cases, the compute, fix, or variable must produce a
per\sphinxhyphen{}atom quantity, not a global quantity.  Note that the {\hyperref[\detokenize{compute_property_atom::doc}]{\sphinxcrossref{\DUrole{doc}{compute
property/atom}}}} command provides access to any
attribute defined and stored by atoms.  If you wish to time\sphinxhyphen{}average
global quantities from a compute, fix, or variable, then see the
{\hyperref[\detokenize{fix_ave_time::doc}]{\sphinxcrossref{\DUrole{doc}{fix ave/time}}}} command.

\sphinxAtStartPar
The per\sphinxhyphen{}atom values of each input vector are summed and averaged
independently of the per\sphinxhyphen{}atom values in other input vectors.

\sphinxAtStartPar
{\hyperref[\detokenize{compute::doc}]{\sphinxcrossref{\DUrole{doc}{Computes}}}} that produce per\sphinxhyphen{}atom quantities are those
which have the word \sphinxstyleemphasis{atom} in their style name.  See the doc pages for
individual {\hyperref[\detokenize{fix::doc}]{\sphinxcrossref{\DUrole{doc}{fixes}}}} to determine which ones produce per\sphinxhyphen{}atom
quantities.  {\hyperref[\detokenize{variable::doc}]{\sphinxcrossref{\DUrole{doc}{Variables}}}} of style \sphinxstyleemphasis{atom} are the only
ones that can be used with this fix since all other styles of variable
produce global quantities.

\sphinxAtStartPar
Note that for values from a compute or fix that produces a per\sphinxhyphen{}atom
array (multiple values per atom), the bracketed index I can be
specified using a wildcard asterisk with the index to effectively
specify multiple values.  This takes the form “*” or “*n” or “n*”
or “m*n”.  If \(N\) = the size of the vector (for \sphinxstyleemphasis{mode} = scalar) or the
number of columns in the array (for \sphinxstyleemphasis{mode} = vector), then an asterisk
with no numeric values means all indices from 1 to \(N\).  A leading
asterisk means all indices from 1 to n (inclusive).  A trailing
asterisk means all indices from m to \(N\) (inclusive).  A middle asterisk
means all indices from m to n (inclusive).

\sphinxAtStartPar
Using a wildcard is the same as if the individual columns of the array
had been listed one by one.  For example, these two fix ave/chunk commands are
equivalent, since the {\hyperref[\detokenize{compute_property_atom::doc}]{\sphinxcrossref{\DUrole{doc}{compute property/atom}}}} command creates, in this case, a per\sphinxhyphen{}atom
array with three columns:

\begin{sphinxVerbatim}[commandchars=\\\{\}]
\PYG{k}{compute }\PYG{n+nv+nvIdentifier}{myAng}\PYG{+w}{ }\PYG{n+nv+nvIdentifier}{all}\PYG{+w}{ }\PYG{n}{property}\PYG{o}{/}\PYG{n}{atom}\PYG{+w}{ }\PYG{n}{angmomx}\PYG{+w}{ }\PYG{n}{angmomy}\PYG{+w}{ }\PYG{n}{angmomz}
\PYG{k}{fix }\PYG{n+nv+nvIdentifier}{1}\PYG{+w}{ }\PYG{n+nv+nvIdentifier}{all}\PYG{+w}{ }\PYG{n}{ave}\PYG{o}{/}\PYG{n}{chunk}\PYG{+w}{ }\PYG{l+m}{100}\PYG{+w}{ }\PYG{l+m}{1}\PYG{+w}{ }\PYG{l+m}{100}\PYG{+w}{ }\PYG{n}{cc1}\PYG{+w}{ }\PYG{n}{c\PYGZus{}myAng[}\PYG{o}{*}\PYG{n}{]}\PYG{+w}{ }\PYG{n}{file}\PYG{+w}{ }\PYG{n}{tmp.angmom}
\PYG{k}{fix }\PYG{n+nv+nvIdentifier}{2}\PYG{+w}{ }\PYG{n+nv+nvIdentifier}{all}\PYG{+w}{ }\PYG{n}{ave}\PYG{o}{/}\PYG{n}{chunk}\PYG{+w}{ }\PYG{l+m}{100}\PYG{+w}{ }\PYG{l+m}{1}\PYG{+w}{ }\PYG{l+m}{100}\PYG{+w}{ }\PYG{n}{cc1}\PYG{+w}{ }\PYG{n}{c\PYGZus{}myAng[1]}\PYG{+w}{ }\PYG{n}{c\PYGZus{}myAng[2]}\PYG{+w}{ }\PYG{n}{c\PYGZus{}myAng[3]}\PYG{+w}{ }\PYG{n}{file}\PYG{+w}{ }\PYG{n}{tmp.angmom}
\end{sphinxVerbatim}

\begin{sphinxadmonition}{note}{Note:}
\sphinxAtStartPar
This fix works by creating an array of size
\(N_\text{chunk} \times N_\text{values}\) on each processor.
\(N_\text{chunk}\) is the number of chunks, which is defined by the
{\hyperref[\detokenize{compute_chunk_atom::doc}]{\sphinxcrossref{\DUrole{doc}{compute chunk/atom}}}} command.
\(N_\text{values}\) is the number of input values specified.
Each processor loops over its atoms, tallying its values to the appropriate
chunk.  Then the entire array is summed across all processors.  This means
that using a large number of chunks will incur an overhead in memory and
computational cost (summing across processors), so be careful to
define a reasonable number of chunks.
\end{sphinxadmonition}


\bigskip\hrule\bigskip


\sphinxAtStartPar
The \(N_\text{every}\), \(N_\text{repeat}\), and \(N_\text{freq}\)
arguments specify on what time steps the input values will be accessed and
contribute to the average.  The final averaged quantities are generated on time
steps that are a multiples of \(N_\text{freq}\).  The average is over
\(N_\text{repeat}\) quantities, computed in the preceding portion of the
simulation every \(N_\text{every}\) time steps.  \(N_\text{freq}\)
must be a multiple of \(N_\text{every}\) and \(N_\text{every}\) must be
non\sphinxhyphen{}zero even if \(N_\text{repeat} = 1\).  Also, the time steps
contributing to the average value cannot overlap (i.e.,
\(N_\text{repeat} \times N_\text{every}\) cannot exceed \(N_\text{freq}\)).

\sphinxAtStartPar
For example, if \(N_\text{every}=2\), \(N_\text{repeat}=6\), and
\(N_\text{freq}=100\), then values on
time steps 90, 92, 94, 96, 98, 100 will be used to compute the final average
on time step 100.  Similarly for time steps 190, 192, 194, 196, 198, 200 on
time step 200, etc.  If \(N_\text{repeat}=1\) and
\(N_\text{freq} = 100\), then no time averaging is done; values are simply
generated on time steps 100, 200, etc.

\sphinxAtStartPar
Each input value can also be averaged over the atoms in each chunk.
The way the averaging is done across the \(N_\text{repeat}\) time steps to
produce output on the \(N_\text{freq}\) time steps, and across multiple
\(N_\text{freq}\) outputs, is determined by the \sphinxstyleemphasis{norm} and \sphinxstyleemphasis{ave} keyword
settings, as discussed below.

\begin{sphinxadmonition}{note}{Note:}
\sphinxAtStartPar
To perform per\sphinxhyphen{}chunk averaging within a \(N_\text{freq}\) time window,
the number of chunks \(N_\text{chunk}\) defined by the
{\hyperref[\detokenize{compute_chunk_atom::doc}]{\sphinxcrossref{\DUrole{doc}{compute chunk/atom}}}} command must remain
constant.  If the \sphinxstyleemphasis{ave} keyword is set to \sphinxstyleemphasis{running} or \sphinxstyleemphasis{window} then
\(N_\text{chunk}\) must remain constant for the duration of the
simulation.  This fix forces the chunk/atom compute specified by chunkID to
hold \(N_\text{chunk}\) constant for the appropriate time windows,
by not allowing it to re\sphinxhyphen{}calculate \(N_\text{chunk}\), which can also
affect how it assigns chunk IDs to atoms.  This is particularly important to
understand if the chunks defined by the {\hyperref[\detokenize{compute_chunk_atom::doc}]{\sphinxcrossref{\DUrole{doc}{compute chunk/atom}}}} command are spatial bins.  If its \sphinxstyleemphasis{units}
keyword is set to \sphinxstyleemphasis{box} or \sphinxstyleemphasis{lattice}, then the number of bins
\(N_\text{chunk}\) and size of each bin will be fixed over the
\(N_\text{freq}\) time window, which can affect which atoms are
discarded if the simulation box size changes.  If its \sphinxstyleemphasis{units} keyword is set
to \sphinxstyleemphasis{reduced}, then the number of bins \(N_\text{chunk}\) will still be
fixed, but the size of each bin can vary at each time step if the
simulation box size changes (e.g., for an NPT simulation).
\end{sphinxadmonition}


\bigskip\hrule\bigskip


\sphinxAtStartPar
The atom attribute values (\sphinxstyleemphasis{vx}, \sphinxstyleemphasis{vy}, \sphinxstyleemphasis{vz}, \sphinxstyleemphasis{fx}, \sphinxstyleemphasis{fy}, \sphinxstyleemphasis{fz}, \sphinxstyleemphasis{mass}) are
self\sphinxhyphen{}explanatory.  As noted above, any other atom attributes can be
used as input values to this fix by using the {\hyperref[\detokenize{compute_property_atom::doc}]{\sphinxcrossref{\DUrole{doc}{compute
property/atom}}}} command and then specifying an
input value from that compute.

\sphinxAtStartPar
The \sphinxstyleemphasis{density/number} value means the number density is computed for
each chunk (i.e., number/volume).  The \sphinxstyleemphasis{density/mass} value means the
mass density is computed for each chunk (i.e., total\sphinxhyphen{}mass/volume).  The
output values are in units of 1/volume or mass density (mass/volume).  See
the {\hyperref[\detokenize{units::doc}]{\sphinxcrossref{\DUrole{doc}{units}}}} command page for the definition of density
for each choice of units (e.g., g/cm\(^3\)).  If the chunks defined by
the {\hyperref[\detokenize{compute_chunk_atom::doc}]{\sphinxcrossref{\DUrole{doc}{compute chunk/atom}}}} command are spatial
bins, the volume is the bin volume.  Otherwise, it is the volume of the
entire simulation box.

\sphinxAtStartPar
The \sphinxstyleemphasis{temp} value means the temperature is computed for each chunk,
by the formula
\begin{equation*}
\begin{split}\text{KE} = \frac{\text{DOF}}{2} k_B T,\end{split}
\end{equation*}
\sphinxAtStartPar
where KE is the total kinetic energy of the chunk of atoms (sum of
\(\frac{1}{2} m v^2\)), DOF is the the total number of degrees of freedom
for all atoms in the chunk, \(k_B\) is the Boltzmann constant, and
\(T\) is the absolute temperature.

\sphinxAtStartPar
The DOF is calculated as \(N\)*adof + cdof, where \(N\) is the
number of atoms in the chunk, adof is the number of degrees of freedom per
atom, and cdof is the number of degrees of freedom per chunk.  By default,
adof = 2 or 3 = dimensionality of system,
as set via the {\hyperref[\detokenize{dimension::doc}]{\sphinxcrossref{\DUrole{doc}{dimension}}}} command, and cdof = 0.0.
This gives the usual formula for temperature.

\sphinxAtStartPar
Note that currently this temperature only includes translational
degrees of freedom for each atom.  No rotational degrees of freedom
are included for finite\sphinxhyphen{}size particles.  Also, no degrees of freedom
are subtracted for any velocity bias or constraints that are applied,
such as {\hyperref[\detokenize{compute_temp_partial::doc}]{\sphinxcrossref{\DUrole{doc}{compute temp/partial}}}},
{\hyperref[\detokenize{fix_shake::doc}]{\sphinxcrossref{\DUrole{doc}{fix shake}}}}, or {\hyperref[\detokenize{fix_rigid::doc}]{\sphinxcrossref{\DUrole{doc}{fix rigid}}}}.  This is
because those degrees of freedom (e.g., a constrained bond) could apply
to sets of atoms that are both included and excluded from a specific
chunk, and hence the concept is somewhat ill\sphinxhyphen{}defined.  In some cases,
you can use the \sphinxstyleemphasis{adof} and \sphinxstyleemphasis{cdof} keywords to adjust the calculated
degrees of freedom appropriately, as explained below.

\sphinxAtStartPar
Also note that a bias can be subtracted from atom velocities before
they are used in the above formula for KE, by using the \sphinxstyleemphasis{bias}
keyword.  This allows, for example, a thermal temperature to be
computed after removal of a flow velocity profile.

\sphinxAtStartPar
Note that the per\sphinxhyphen{}chunk temperature calculated by this fix and the
{\hyperref[\detokenize{compute_temp_chunk::doc}]{\sphinxcrossref{\DUrole{doc}{compute temp/chunk}}}} command can be
different.  The compute calculates the temperature for each chunk for
a single snapshot.  This fix can do that but can also time average
those values over many snapshots, or it can compute a temperature as
if the atoms in the chunk on different time steps were collected
together as one set of atoms to calculate their temperature.  The
compute allows the center\sphinxhyphen{}of\sphinxhyphen{}mass velocity of each chunk to be
subtracted before calculating the temperature; this fix does not.

\sphinxAtStartPar
If a value begins with “c\_”, a compute ID must follow which has been
previously defined in the input script.  If no bracketed integer is
appended, the per\sphinxhyphen{}atom vector calculated by the compute is used.  If a
bracketed integer is appended, the Ith column of the per\sphinxhyphen{}atom array
calculated by the compute is used.  Users can also write code for
their own compute styles and {\hyperref[\detokenize{Modify::doc}]{\sphinxcrossref{\DUrole{doc}{add them to LAMMPS}}}}.
See the discussion above for how I can be specified with a wildcard
asterisk to effectively specify multiple values.

\sphinxAtStartPar
If a value begins with “f\_”, a fix ID must follow which has been
previously defined in the input script.  If no bracketed integer is
appended, the per\sphinxhyphen{}atom vector calculated by the fix is used.  If a
bracketed integer is appended, the Ith column of the per\sphinxhyphen{}atom array
calculated by the fix is used.  Note that some fixes only produce
their values on certain time steps, which must be compatible with
\(N_\text{every}\), else an error results.  Users can also write code for
their own fix styles and {\hyperref[\detokenize{Modify::doc}]{\sphinxcrossref{\DUrole{doc}{add them to LAMMPS}}}}.  See the
discussion above for how I can be specified with a wildcard asterisk
to effectively specify multiple values.

\sphinxAtStartPar
If a value begins with “v\_”, a variable name must follow which has
been previously defined in the input script.  Variables of style
\sphinxstyleemphasis{atom} can reference thermodynamic keywords and various per\sphinxhyphen{}atom
attributes, or invoke other computes, fixes, or variables when they
are evaluated, so this is a very general means of generating per\sphinxhyphen{}atom
quantities to average within chunks.


\bigskip\hrule\bigskip


\sphinxAtStartPar
Additional optional keywords also affect the operation of this fix
and its outputs.

\sphinxAtStartPar
The \sphinxstyleemphasis{norm} keyword affects how averaging is done for the per\sphinxhyphen{}chunk
values that are output every \(N_\text{freq}\) time steps.

\sphinxAtStartPar
It the \sphinxstyleemphasis{norm} setting is \sphinxstyleemphasis{all}, which is the default, a chunk value is summed
over all atoms in all \(N_\text{repeat}\) samples, as is the count of
atoms in the chunk.  The averaged output value for the chunk on the
\(N_\text{freq}\) time steps is Total\sphinxhyphen{}sum / Total\sphinxhyphen{}count.  In other words it
is an average over atoms across the entire \(N_\text{freq}\) timescale.
For the \sphinxstyleemphasis{density/number} and \sphinxstyleemphasis{density/mass} values, the volume (bin volume or
system volume) used in the final normalization will be the volume at
the final \(N_\text{freq}\) time step. For the \sphinxstyleemphasis{temp} values, degrees of
freedom and kinetic energy are summed separately across the entire
\(N_\text{freq}\) timescale, and the output value is calculated by dividing
those two sums.

\sphinxAtStartPar
If the \sphinxstyleemphasis{norm} setting is \sphinxstyleemphasis{sample}, the chunk value is summed over
atoms for each sample, as is the count, and an “average sample value”
is computed for each sample (i.e., Sample\sphinxhyphen{}sum / Sample\sphinxhyphen{}count).  The
output value for the chunk on the \(N_\text{freq}\) time steps is the
average of the \(N_\text{repeat}\) “average sample values” (i.e., the sum
of \(N_\text{repeat}\) “average sample values” divided by
\(N_\text{repeat}\)).  In other words, it is an average of an average.
For the \sphinxstyleemphasis{density/number} and \sphinxstyleemphasis{density/mass} values, the volume (bin volume or
system volume) used in the per\sphinxhyphen{}sample normalization will be the current volume
at each sampling step.

\sphinxAtStartPar
If the \sphinxstyleemphasis{norm} setting is \sphinxstyleemphasis{none}, a similar computation as for the
\sphinxstyleemphasis{sample} setting is done, except the individual “average sample
values” are “summed sample values”.  A summed sample value is simply
the chunk value summed over atoms in the sample, without dividing by
the number of atoms in the sample.  The output value for the chunk on
the \(N_\text{freq}\) timesteps is the average of the
\(N_\text{repeat}\) “summed sample values” (i.e., the sum of
\(N_\text{repeat}\) “summed sample values” divided by
\(N_\text{repeat}\)).
For the \sphinxstyleemphasis{density/number} and \sphinxstyleemphasis{density/mass} values, the
volume (bin volume or system volume) used in the per\sphinxhyphen{}sample sum
normalization will be the current volume at each sampling step.


\bigskip\hrule\bigskip


\sphinxAtStartPar
The \sphinxstyleemphasis{ave} keyword determines how the per\sphinxhyphen{}chunk values produced every
\(N_\text{freq}\) steps are averaged with values produced on previous steps
that were multiples of \(N_\text{freq}\), before they are accessed by
another output command or written to a file.

\sphinxAtStartPar
If the \sphinxstyleemphasis{ave} setting is \sphinxstyleemphasis{one}, which is the default, then the chunk
values produced on timesteps that are multiples of \(N_\text{freq}\) are
independent of each other; they are output as\sphinxhyphen{}is without further averaging.

\sphinxAtStartPar
If the \sphinxstyleemphasis{ave} setting is \sphinxstyleemphasis{running}, then the chunk values produced on
timesteps that are multiples of \(N_\text{freq}\) are summed and averaged
in a cumulative sense before being output.  Each output chunk value is thus
the average of the chunk value produced on that timestep with all
preceding values for the same chunk.  This running average begins when
the fix is defined; it can only be restarted by deleting the fix via
the {\hyperref[\detokenize{unfix::doc}]{\sphinxcrossref{\DUrole{doc}{unfix}}}} command, or re\sphinxhyphen{}defining the fix by re\sphinxhyphen{}specifying it.

\sphinxAtStartPar
If the \sphinxstyleemphasis{ave} setting is \sphinxstyleemphasis{window}, then the chunk values produced on
timesteps that are multiples of \(N_\text{freq}\) are summed and averaged
within a moving “window” of time, so that the last \(M\) values for the
same chunk are used to produce the output.  For example, if \(M = 3\) and
\(N_\text{freq} = 1000\), then the output on step 10000 will be the average
of the individual chunk values on time steps 8000, 9000, and 10000.  Outputs on
early steps will average over less than \(M\) values if they are not
available.


\bigskip\hrule\bigskip


\sphinxAtStartPar
The \sphinxstyleemphasis{bias} keyword specifies the ID of a temperature compute that
removes a “bias” velocity from each atom, specified as \sphinxstyleemphasis{bias\sphinxhyphen{}ID}.
It is only used when the \sphinxstyleemphasis{temp} value is calculated, to compute the
thermal temperature of each chunk after the translational kinetic
energy components have been altered in a prescribed way (e.g., to
remove a flow velocity profile).  See the doc pages for individual
computes that calculate a temperature to see which ones implement a bias.

\sphinxAtStartPar
The \sphinxstyleemphasis{adof} and \sphinxstyleemphasis{cdof} keywords define the values used in the degree of
freedom (DOF) formula described above for temperature calculation
for each chunk.  They are only used when the \sphinxstyleemphasis{temp} value is
calculated.  They can be used to calculate a more appropriate
temperature for some kinds of chunks.  Here are three examples:

\sphinxAtStartPar
If spatially binned chunks contain some number of water molecules and
{\hyperref[\detokenize{fix_shake::doc}]{\sphinxcrossref{\DUrole{doc}{fix shake}}}} is used to make each molecule rigid, then
you could calculate a temperature with six degrees of freedom (DOF) (three
translational, three rotational) per molecule by setting \sphinxstyleemphasis{adof} to 2.0.

\sphinxAtStartPar
If {\hyperref[\detokenize{compute_temp_partial::doc}]{\sphinxcrossref{\DUrole{doc}{compute temp/partial}}}} is used with the
\sphinxstyleemphasis{bias} keyword to only allow the \(x\) component of velocity to contribute
to the temperature, then \sphinxstyleemphasis{adof} = 1.0 would be appropriate.

\sphinxAtStartPar
If each chunk consists of a large molecule, with some number of its
bonds constrained by {\hyperref[\detokenize{fix_shake::doc}]{\sphinxcrossref{\DUrole{doc}{fix shake}}}} or the entire molecule
by {\hyperref[\detokenize{fix_rigid::doc}]{\sphinxcrossref{\DUrole{doc}{fix rigid/small}}}}, \sphinxstyleemphasis{adof} = 0.0 and \sphinxstyleemphasis{cdof} could be
set to the remaining degrees of freedom for the entire molecule
(entire chunk in this case), that is, 6 for 3d or 3 for 2d for a rigid
molecule.


\bigskip\hrule\bigskip


\sphinxAtStartPar
\DUrole{versionmodified,added}{Added in version 17Apr2024: }new keyword \sphinxstyleemphasis{append}

\sphinxAtStartPar
The \sphinxstyleemphasis{file} or \sphinxstyleemphasis{append} keywords allow a filename to be specified.  If
\sphinxstyleemphasis{file} is used, then the filename is overwritten if it already exists.
If \sphinxstyleemphasis{append} is used, then the filename is appended to if it already
exists, or created if it does not exist.  Every \(N_\text{freq}\)
timesteps, a section of chunk info will be written to a text file in the
following format.  A line with the timestep and number of chunks is
written.  Then one line per chunk is written, containing the chunk ID
\((1-N_\text{chunk}),\) an optional original ID value, optional
coordinate values for chunks that represent spatial bins, the number of
atoms in the chunk, and one or more calculated values.  More explanation
of the optional values is given below.  The number of values in each
line corresponds to the number of values specified in the fix ave/chunk
command.  The number of atoms and the value(s) are summed or average
quantities, as explained above.

\sphinxAtStartPar
The \sphinxstyleemphasis{overwrite} keyword will continuously overwrite the output file
with the latest output, so that it only contains one timestep worth of
output.  This option can only be used with the \sphinxstyleemphasis{ave running} setting.

\sphinxAtStartPar
The \sphinxstyleemphasis{format} keyword sets the numeric format of each value when it is
printed to a file via the \sphinxstyleemphasis{file} keyword.  Note that all values are
floating point quantities.  The default format is \%g.  You can specify
a higher precision if desired (e.g., \%20.16g).

\sphinxAtStartPar
The \sphinxstyleemphasis{title1} and \sphinxstyleemphasis{title2} and \sphinxstyleemphasis{title3} keywords allow specification of
the strings that will be printed as the first three lines of the output
file, assuming the \sphinxstyleemphasis{file} keyword was used.  LAMMPS uses default
values for each of these, so they do not need to be specified.

\sphinxAtStartPar
By default, these header lines are as follows:

\begin{sphinxVerbatim}[commandchars=\\\{\}]
\PYGZsh{} Chunk\PYGZhy{}averaged data for fix ID and group name
\PYGZsh{} Timestep Number\PYGZhy{}of\PYGZhy{}chunks
\PYGZsh{} Chunk (OrigID) (Coord1) (Coord2) (Coord3) Ncount value1 value2 ...
\end{sphinxVerbatim}

\sphinxAtStartPar
In the first line, ID and name are replaced with the fix\sphinxhyphen{}ID and group
name.  The second line describes the two values that are printed at
the first of each section of output.  In the third line the values are
replaced with the appropriate value names (e.g., \sphinxstyleemphasis{fx} or c\_myCompute{[}2{]}).

\sphinxAtStartPar
The words in parenthesis only appear with corresponding columns if the
chunk style specified for the {\hyperref[\detokenize{compute_chunk_atom::doc}]{\sphinxcrossref{\DUrole{doc}{compute chunk/atom}}}} command supports them.  The OrigID column is
only used if the \sphinxstyleemphasis{compress} keyword was set to \sphinxstyleemphasis{yes} for the
{\hyperref[\detokenize{compute_chunk_atom::doc}]{\sphinxcrossref{\DUrole{doc}{compute chunk/atom}}}} command.  This means
that the original chunk IDs (e.g., molecule IDs) will have been
compressed to remove chunk IDs with no atoms assigned to them.  Thus a
compressed chunk ID of 3 may correspond to an original chunk ID or
molecule ID of 415.  The OrigID column will list 415 for the third chunk.

\sphinxAtStartPar
The CoordN columns only appear if a \sphinxstyleemphasis{binning} style was used in the
{\hyperref[\detokenize{compute_chunk_atom::doc}]{\sphinxcrossref{\DUrole{doc}{compute chunk/atom}}}} command.  For \sphinxstyleemphasis{bin/1d},
\sphinxstyleemphasis{bin/2d}, and \sphinxstyleemphasis{bin/3d} styles the column values are the center point
of the bin in the corresponding dimension.  Just Coord1 is used for
\sphinxstyleemphasis{bin/1d}, Coord2 is added for \sphinxstyleemphasis{bin/2d}, Coord3 is added for \sphinxstyleemphasis{bin/3d}.
For \sphinxstyleemphasis{bin/sphere}, just Coord1 is used, and it is the radial
coordinate.  For \sphinxstyleemphasis{bin/cylinder}, Coord1 and Coord2 are used.  Coord1
is the radial coordinate (away from the cylinder axis), and coord2 is
the coordinate along the cylinder axis.

\sphinxAtStartPar
Note that if the value of the \sphinxstyleemphasis{units} keyword used in the
{\hyperref[\detokenize{compute_chunk_atom::doc}]{\sphinxcrossref{\DUrole{doc}{compute chunk/atom command}}}} is \sphinxstyleemphasis{box} or
\sphinxstyleemphasis{lattice}, the coordinate values will be in distance {\hyperref[\detokenize{units::doc}]{\sphinxcrossref{\DUrole{doc}{units}}}}.
If the value of the \sphinxstyleemphasis{units} keyword is \sphinxstyleemphasis{reduced}, the
coordinate values will be in unitless reduced units (0\textendash{}1).  This is
not true for the Coord1 value of style \sphinxstyleemphasis{bin/sphere} or \sphinxstyleemphasis{bin/cylinder}
which both represent radial dimensions.  Those values are always in
distance {\hyperref[\detokenize{units::doc}]{\sphinxcrossref{\DUrole{doc}{units}}}}.


\bigskip\hrule\bigskip



\subsection{Restart, fix\_modify, output, run start/stop, minimize info}
\label{\detokenize{fix_ave_chunk:restart-fix-modify-output-run-start-stop-minimize-info}}
\sphinxAtStartPar
No information about this fix is written to {\hyperref[\detokenize{restart::doc}]{\sphinxcrossref{\DUrole{doc}{binary restart files}}}}.  None of the {\hyperref[\detokenize{fix_modify::doc}]{\sphinxcrossref{\DUrole{doc}{fix\_modify}}}} options are
relevant to this fix.

\sphinxAtStartPar
This fix computes a global array of values which can be accessed by
various {\hyperref[\detokenize{Howto_output::doc}]{\sphinxcrossref{\DUrole{doc}{output commands}}}}.  The values can only be
accessed on timesteps that are multiples of \(N_\text{freq}\), since
that is when averaging is performed.  The global array has \# of rows =
the number of chunks \(N_\text{chunk}\), as calculated by the
specified {\hyperref[\detokenize{compute_chunk_atom::doc}]{\sphinxcrossref{\DUrole{doc}{compute chunk/atom}}}} command.  The \#
of columns is \(M+1+N_\text{values}\), where \(M \in
\{1,\dotsc,4\}\), depending on whether the optional columns for OrigID
and CoordN are used, as explained above.  Following the optional
columns, the next column contains the count of atoms in the chunk, and
the remaining columns are the Nvalue quantities.  When the array is
accessed with a row \(I\) that exceeds the current number of chunks,
than a 0.0 is returned by the fix instead of an error, since the number
of chunks can vary as a simulation runs depending on how that value is
computed by the compute chunk/atom command.

\sphinxAtStartPar
The array values calculated by this fix are treated as “intensive”,
since they are typically already normalized by the count of atoms in
each chunk.

\sphinxAtStartPar
No parameter of this fix can be used with the \sphinxstyleemphasis{start/stop} keywords of
the {\hyperref[\detokenize{run::doc}]{\sphinxcrossref{\DUrole{doc}{run}}}} command.  This fix is not invoked during
{\hyperref[\detokenize{minimize::doc}]{\sphinxcrossref{\DUrole{doc}{energy minimization}}}}.


\subsection{Restrictions}
\label{\detokenize{fix_ave_chunk:restrictions}}\begin{quote}

\sphinxAtStartPar
none
\end{quote}


\subsection{Related commands}
\label{\detokenize{fix_ave_chunk:related-commands}}
\sphinxAtStartPar
{\hyperref[\detokenize{compute::doc}]{\sphinxcrossref{\DUrole{doc}{compute}}}}, {\hyperref[\detokenize{fix_ave_atom::doc}]{\sphinxcrossref{\DUrole{doc}{fix ave/atom}}}},
{\hyperref[\detokenize{fix_ave_histo::doc}]{\sphinxcrossref{\DUrole{doc}{fix ave/histo}}}}, {\hyperref[\detokenize{fix_ave_time::doc}]{\sphinxcrossref{\DUrole{doc}{fix ave/time}}}},
{\hyperref[\detokenize{variable::doc}]{\sphinxcrossref{\DUrole{doc}{variable}}}}, {\hyperref[\detokenize{fix_ave_correlate::doc}]{\sphinxcrossref{\DUrole{doc}{fix ave/correlate}}}},
{\hyperref[\detokenize{fix_ave_grid::doc}]{\sphinxcrossref{\DUrole{doc}{fix ave/grid}}}}


\subsection{Default}
\label{\detokenize{fix_ave_chunk:default}}
\sphinxAtStartPar
The option defaults are norm = all, ave = one, bias = none, no file
output, and title 1,2,3 = strings as described above.

\sphinxstepscope

\index{fix ave/correlate@\spxentry{fix ave/correlate}}\ignorespaces 

\section{fix ave/correlate command}
\label{\detokenize{fix_ave_correlate:fix-ave-correlate-command}}\label{\detokenize{fix_ave_correlate:index-0}}\label{\detokenize{fix_ave_correlate::doc}}

\subsection{Syntax}
\label{\detokenize{fix_ave_correlate:syntax}}
\begin{sphinxVerbatim}[commandchars=\\\{\}]
\PYG{k}{fix }\PYG{n+nv+nvIdentifier}{ID}\PYG{+w}{ }\PYG{n+nv+nvIdentifier}{group\PYGZhy{}ID}\PYG{+w}{ }\PYG{n}{ave}\PYG{o}{/}\PYG{n}{correlate}\PYG{+w}{ }\PYG{n}{Nevery}\PYG{+w}{ }\PYG{n}{Nrepeat}\PYG{+w}{ }\PYG{n}{Nfreq}\PYG{+w}{ }\PYG{n}{value1}\PYG{+w}{ }\PYG{n}{value2}\PYG{+w}{ }\PYG{n}{...}\PYG{+w}{ }\PYG{n}{keyword}\PYG{+w}{ }\PYG{n}{args}\PYG{+w}{ }\PYG{n}{...}
\end{sphinxVerbatim}
\begin{itemize}
\item {} 
\sphinxAtStartPar
ID, group\sphinxhyphen{}ID are documented in {\hyperref[\detokenize{fix::doc}]{\sphinxcrossref{\DUrole{doc}{fix}}}} command

\item {} 
\sphinxAtStartPar
ave/correlate = style name of this fix command

\item {} 
\sphinxAtStartPar
Nevery = use input values every this many timesteps

\item {} 
\sphinxAtStartPar
Nrepeat = \# of correlation time windows to accumulate

\item {} 
\sphinxAtStartPar
Nfreq = calculate time window averages every this many timesteps

\item {} 
\sphinxAtStartPar
one or more input values can be listed

\item {} 
\sphinxAtStartPar
value = c\_ID, c\_ID{[}N{]}, f\_ID, f\_ID{[}N{]}, v\_name

\begin{sphinxVerbatim}[commandchars=\\\{\}]
c\PYGZus{}ID = global scalar calculated by a compute with ID
c\PYGZus{}ID[I] = Ith component of global vector calculated by a compute with ID, I can include wildcard (see below)
f\PYGZus{}ID = global scalar calculated by a fix with ID
f\PYGZus{}ID[I] = Ith component of global vector calculated by a fix with ID, I can include wildcard (see below)
v\PYGZus{}name = global value calculated by an equal\PYGZhy{}style variable with name
v\PYGZus{}name[I] = Ith component of a vector\PYGZhy{}style variable with name, I can include wildcard (see below)
\end{sphinxVerbatim}

\item {} 
\sphinxAtStartPar
zero or more keyword/arg pairs may be appended

\item {} 
\sphinxAtStartPar
keyword = \sphinxstyleemphasis{type} or \sphinxstyleemphasis{ave} or \sphinxstyleemphasis{start} or \sphinxstyleemphasis{prefactor} or \sphinxstyleemphasis{file} or \sphinxstyleemphasis{overwrite} or \sphinxstyleemphasis{title1} or \sphinxstyleemphasis{title2} or \sphinxstyleemphasis{title3}
\begin{sphinxalltt}
\sphinxstyleemphasis{type} arg = \sphinxstyleemphasis{auto} or \sphinxstyleemphasis{upper} or \sphinxstyleemphasis{lower} or \sphinxstyleemphasis{auto/upper} or \sphinxstyleemphasis{auto/lower} or \sphinxstyleemphasis{full}
  auto = correlate each value with itself
  upper = correlate each value with each succeeding value
  lower = correlate each value with each preceding value
  auto/upper = auto + upper
  auto/lower = auto + lower
  full = correlate each value with every other value, including itself = auto + upper + lower
\sphinxstyleemphasis{ave} args = \sphinxstyleemphasis{one} or \sphinxstyleemphasis{running}
  one = zero the correlation accumulation every Nfreq steps
  running = accumulate correlations continuously
\sphinxstyleemphasis{start} args = Nstart
  Nstart = start accumulating correlations on this timestep
\sphinxstyleemphasis{prefactor} args = value
  value = prefactor to scale all the correlation data by
\sphinxstyleemphasis{file} arg = filename
  filename = name of file to output correlation data to
\sphinxstyleemphasis{overwrite} arg = none = overwrite output file with only latest output
\sphinxstyleemphasis{title1} arg = string
  string = text to print as 1st line of output file
\sphinxstyleemphasis{title2} arg = string
  string = text to print as 2nd line of output file
\sphinxstyleemphasis{title3} arg = string
  string = text to print as 3rd line of output file
\end{sphinxalltt}

\end{itemize}


\subsection{Examples}
\label{\detokenize{fix_ave_correlate:examples}}
\begin{sphinxVerbatim}[commandchars=\\\{\}]
\PYG{k}{fix }\PYG{n+nv+nvIdentifier}{1}\PYG{+w}{ }\PYG{n+nv+nvIdentifier}{all}\PYG{+w}{ }\PYG{n}{ave}\PYG{o}{/}\PYG{n}{correlate}\PYG{+w}{ }\PYG{l+m}{5}\PYG{+w}{ }\PYG{l+m}{100}\PYG{+w}{ }\PYG{l+m}{1000}\PYG{+w}{ }\PYG{n}{c\PYGZus{}myTemp}\PYG{+w}{ }\PYG{n}{file}\PYG{+w}{ }\PYG{n}{temp.correlate}
\PYG{k}{fix }\PYG{n+nv+nvIdentifier}{1}\PYG{+w}{ }\PYG{n+nv+nvIdentifier}{all}\PYG{+w}{ }\PYG{n}{ave}\PYG{o}{/}\PYG{n}{correlate}\PYG{+w}{ }\PYG{l+m}{1}\PYG{+w}{ }\PYG{l+m}{50}\PYG{+w}{ }\PYG{l+m}{10000}\PYG{+w}{ }\PYG{o}{\PYGZam{}}
\PYG{+w}{          }\PYG{n}{c\PYGZus{}thermo\PYGZus{}press[1]}\PYG{+w}{ }\PYG{n}{c\PYGZus{}thermo\PYGZus{}press[2]}\PYG{+w}{ }\PYG{n}{c\PYGZus{}thermo\PYGZus{}press[3]}\PYG{+w}{ }\PYG{o}{\PYGZam{}}
\PYG{+w}{          }\PYG{n}{type}\PYG{+w}{ }\PYG{n}{upper}\PYG{+w}{ }\PYG{n}{ave}\PYG{+w}{ }\PYG{n}{running}\PYG{+w}{ }\PYG{n}{title1}\PYG{+w}{ }\PYG{l+s}{\PYGZdq{}}\PYG{l+s}{My correlation data}\PYG{l+s}{\PYGZdq{}}
\PYG{k}{fix }\PYG{n+nv+nvIdentifier}{1}\PYG{+w}{ }\PYG{n+nv+nvIdentifier}{all}\PYG{+w}{ }\PYG{n}{ave}\PYG{o}{/}\PYG{n}{correlate}\PYG{+w}{ }\PYG{l+m}{1}\PYG{+w}{ }\PYG{l+m}{50}\PYG{+w}{ }\PYG{l+m}{10000}\PYG{+w}{ }\PYG{n}{c\PYGZus{}thermo\PYGZus{}press[}\PYG{o}{*}\PYG{n}{]}
\end{sphinxVerbatim}


\subsection{Description}
\label{\detokenize{fix_ave_correlate:description}}
\sphinxAtStartPar
Use one or more global scalar values as inputs every few timesteps,
calculate time correlations between them at varying time intervals,
and average the correlation data over longer timescales.  The resulting
correlation values can be time integrated by
{\hyperref[\detokenize{variable::doc}]{\sphinxcrossref{\DUrole{doc}{variables}}}} or used by other
{\hyperref[\detokenize{Howto_output::doc}]{\sphinxcrossref{\DUrole{doc}{output commands}}}} such as
{\hyperref[\detokenize{thermo_style::doc}]{\sphinxcrossref{\DUrole{doc}{thermo\_style custom}}}}, and can also be written to a file.
See the {\hyperref[\detokenize{fix_ave_correlate_long::doc}]{\sphinxcrossref{\DUrole{doc}{fix ave/correlate/long}}}} command for an
alternate method for computing correlation functions efficiently over
very long time windows.

\sphinxAtStartPar
The group specified with this command is ignored.  However, note that
specified values may represent calculations performed by computes and
fixes which store their own “group” definitions.

\sphinxAtStartPar
Each listed value can be the result of a {\hyperref[\detokenize{compute::doc}]{\sphinxcrossref{\DUrole{doc}{compute}}}} or
{\hyperref[\detokenize{fix::doc}]{\sphinxcrossref{\DUrole{doc}{fix}}}} or the evaluation of an equal\sphinxhyphen{}style or vector\sphinxhyphen{}style
{\hyperref[\detokenize{variable::doc}]{\sphinxcrossref{\DUrole{doc}{variable}}}}.  In each case, the compute, fix, or variable
must produce a global quantity, not a per\sphinxhyphen{}atom or local quantity.  If
you wish to spatial\sphinxhyphen{} or time\sphinxhyphen{}average or histogram per\sphinxhyphen{}atom quantities
from a compute, fix, or variable, then see the
{\hyperref[\detokenize{fix_ave_chunk::doc}]{\sphinxcrossref{\DUrole{doc}{fix ave/chunk}}}}, {\hyperref[\detokenize{fix_ave_atom::doc}]{\sphinxcrossref{\DUrole{doc}{fix ave/atom}}}}, or
{\hyperref[\detokenize{fix_ave_histo::doc}]{\sphinxcrossref{\DUrole{doc}{fix ave/histo}}}} commands.  If you wish to convert a
per\sphinxhyphen{}atom quantity into a single global value, see the
{\hyperref[\detokenize{compute_reduce::doc}]{\sphinxcrossref{\DUrole{doc}{compute reduce}}}} command.

\sphinxAtStartPar
The input values must be all scalars.  What kinds of
correlations between input values are calculated is determined by the
\sphinxstyleemphasis{type} keyword as discussed below.

\sphinxAtStartPar
{\hyperref[\detokenize{compute::doc}]{\sphinxcrossref{\DUrole{doc}{Computes}}}} that produce global quantities are those which
do not have the word \sphinxstyleemphasis{atom} in their style name.  Only a few
{\hyperref[\detokenize{fix::doc}]{\sphinxcrossref{\DUrole{doc}{fixes}}}} produce global quantities.  See the doc pages for
individual fixes for info on which ones produce such values.
{\hyperref[\detokenize{variable::doc}]{\sphinxcrossref{\DUrole{doc}{Variables}}}} of style \sphinxstyleemphasis{equal} and \sphinxstyleemphasis{vector} are the only
ones that can be used with this fix.  Variables of style \sphinxstyleemphasis{atom} cannot
be used, since they produce per\sphinxhyphen{}atom values.


\bigskip\hrule\bigskip


\sphinxAtStartPar
For input values from a compute or fix or variable , the bracketed
index I can be specified using a wildcard asterisk with the index to
effectively specify multiple values.  This takes the form “*” or
“*n” or “m*” or “m*n”.  If \(N\) is the size of the vector, then an
asterisk with no numeric values means all indices from 1 to \(N\).  A
leading asterisk means all indices from 1 to n (inclusive).  A
trailing asterisk means all indices from m to \(N\) (inclusive).
A middle asterisk means all indices from m to n (inclusive).

\sphinxAtStartPar
Using a wildcard is the same as if the individual elements of the
vector had been listed one by one.  For example, the following two fix
ave/correlate commands are equivalent, since the {\hyperref[\detokenize{compute_pressure::doc}]{\sphinxcrossref{\DUrole{doc}{compute pressure}}}} command creates a global vector with six values:

\begin{sphinxVerbatim}[commandchars=\\\{\}]
\PYG{k}{compute }\PYG{n+nv+nvIdentifier}{myPress}\PYG{+w}{ }\PYG{n+nv+nvIdentifier}{all}\PYG{+w}{ }\PYG{n}{pressure}\PYG{+w}{ }\PYG{n}{NULL}
\PYG{k}{fix }\PYG{n+nv+nvIdentifier}{1}\PYG{+w}{ }\PYG{n+nv+nvIdentifier}{all}\PYG{+w}{ }\PYG{n}{ave}\PYG{o}{/}\PYG{n}{correlate}\PYG{+w}{ }\PYG{l+m}{1}\PYG{+w}{ }\PYG{l+m}{50}\PYG{+w}{ }\PYG{l+m}{10000}\PYG{+w}{ }\PYG{n}{c\PYGZus{}myPress[}\PYG{o}{*}\PYG{n}{]}
\PYG{k}{fix }\PYG{n+nv+nvIdentifier}{1}\PYG{+w}{ }\PYG{n+nv+nvIdentifier}{all}\PYG{+w}{ }\PYG{n}{ave}\PYG{o}{/}\PYG{n}{correlate}\PYG{+w}{ }\PYG{l+m}{1}\PYG{+w}{ }\PYG{l+m}{50}\PYG{+w}{ }\PYG{l+m}{10000}\PYG{+w}{ }\PYG{o}{\PYGZam{}}
\PYG{+w}{          }\PYG{n}{c\PYGZus{}myPress[1]}\PYG{+w}{ }\PYG{n}{c\PYGZus{}myPress[2]}\PYG{+w}{ }\PYG{n}{c\PYGZus{}myPress[3]}\PYG{+w}{ }\PYG{o}{\PYGZam{}}
\PYG{+w}{          }\PYG{n}{c\PYGZus{}myPress[4]}\PYG{+w}{ }\PYG{n}{c\PYGZus{}myPress[5]}\PYG{+w}{ }\PYG{n}{c\PYGZus{}myPress[6]}
\end{sphinxVerbatim}

\begin{sphinxadmonition}{note}{Note:}
\sphinxAtStartPar
For a vector\sphinxhyphen{}style variable, only the wildcard forms “*n” or
“m*n” are allowed.  You must specify the upper bound, because
vector\sphinxhyphen{}style variable lengths are not determined until the variable
is evaluated.  If n is specified larger than the vector length
turns out to be, zeroes are output for missing vector values.
\end{sphinxadmonition}


\bigskip\hrule\bigskip


\sphinxAtStartPar
The \(N_\text{every}\), \(N_\text{repeat}\), and \(N_\text{freq}\)
arguments specify on what timesteps the input values will be used to calculate
correlation data.  The input values are sampled every \(N_\text{every}\)
time steps.  The correlation data for the preceding samples is computed on
time steps that are a multiple of \(N_\text{freq}\).  Consider a set of
samples from some initial time up to an output timestep.  The initial time
could be the beginning of the simulation or the last output time; see the \sphinxstyleemphasis{ave}
keyword for options.  For the set of samples, the correlation value
\(C_{ij}\) is calculated as:
\begin{equation*}
\begin{split}C_{ij}(\Delta t) = \left\langle V_i(t) V_j(t+\Delta t)\right\rangle,\end{split}
\end{equation*}
\sphinxAtStartPar
which is the correlation value between input values \(V_i\) and
\(V_j\), separated by time \(\Delta t\).  Note that the second value
\(V_j\) in the pair is always the one sampled at the later time.  The
average is an average over every pair of samples in the set that are separated
by time \(\Delta t\).  The maximum \(\Delta t\) used is of size
\((N_\text{repeat} - 1) N_\text{every}\).
Thus the correlation between a pair of input values yields
\(N_\text{repeat}\) correlation data:
\begin{equation*}
\begin{split}C_{ij}(0), C_{ij}(N_\text{every}), C_{ij}(2N_\text{every}), \dotsc,
  C_{ij}\bigl((N_\text{repeat}-1) N_\text{every}\bigr)\end{split}
\end{equation*}
\sphinxAtStartPar
For example, if \(N_\text{every}=5\), \(N_\text{repeat}=6\), and
\(N_\text{freq}=100\), then values on time steps
\(0, 5, 10, 15,\dotsc,100\) will be used to compute the final averages
on time step 100.  Six averages will be computed: \(C_{ij}(0)\),
\(C_{ij}(5)\), \(C_{ij}(10)\), \(C_{ij}(15)\), \(C_{ij}(20)\),
and \(C_{ij}(25)\).  \(C_{ij}(10)\) on time step 100 will
be the average of 19 samples, namely \(V_i(0) V_j(10)\),
\(V_i(5) V_j(15)\), \(V_i(10) V_j(20)\),
\(V_i(15) V_j(25), \dotsc,\)
\(V_i(85) V_j(95)\), and \(V_i(90) V_j(100)\).

\sphinxAtStartPar
\(N_\text{freq}\) must be a multiple of \(N_\text{every}\);
\(N_\text{every}\) and \(N_\text{repeat}\) must be non\sphinxhyphen{}zero.
Also, if the \sphinxstyleemphasis{ave} keyword is set to \sphinxstyleemphasis{one} which is the default, then
\(N_\text{freq} \ge (N_\text{repeat} -1) N_\text{every}\) is required.


\bigskip\hrule\bigskip


\sphinxAtStartPar
If a value begins with “c\_”, a compute ID must follow which has been
previously defined in the input script.  If no bracketed term is
appended, the global scalar calculated by the compute is used.  If a
bracketed term is appended, the \(I^\text{th}\) element of the global
vector calculated by the compute is used.  See the discussion above for how
\(I\) can be specified with a wildcard asterisk to effectively specify
multiple values.

\sphinxAtStartPar
Note that there is a {\hyperref[\detokenize{compute_reduce::doc}]{\sphinxcrossref{\DUrole{doc}{compute reduce}}}} command
that can sum per\sphinxhyphen{}atom quantities into a global scalar or vector which
can then be accessed by fix ave/correlate.  It can also be a compute defined
not in your input script, but by {\hyperref[\detokenize{thermo_style::doc}]{\sphinxcrossref{\DUrole{doc}{thermodynamic output}}}}
or other fixes such as {\hyperref[\detokenize{fix_nh::doc}]{\sphinxcrossref{\DUrole{doc}{fix nvt}}}}
or {\hyperref[\detokenize{fix_temp_rescale::doc}]{\sphinxcrossref{\DUrole{doc}{fix temp/rescale}}}}.  See the doc pages for
these commands which give the IDs of these computes.  Users can also
write code for their own compute styles and {\hyperref[\detokenize{Modify::doc}]{\sphinxcrossref{\DUrole{doc}{add them to LAMMPS}}}}.

\sphinxAtStartPar
If a value begins with “f\_”, a fix ID must follow which has been
previously defined in the input script.  If no bracketed term is
appended, the global scalar calculated by the fix is used.  If a
bracketed term is appended, the \(I^\text{th}\) element of the global
vector calculated by the fix is used.  See the discussion above for how
\(I\) can be specified with a wildcard asterisk to effectively specify
multiple values.

\sphinxAtStartPar
Note that some fixes only produce their values on certain timesteps,
which must be compatible with \(N_\text{every}\), else an error will
result.  Users can also write code for their own fix styles and
{\hyperref[\detokenize{Modify::doc}]{\sphinxcrossref{\DUrole{doc}{add them to LAMMPS}}}}.

\sphinxAtStartPar
If a value begins with “v\_”, a variable name must follow which has been
previously defined in the input script.  Only equal\sphinxhyphen{}style or vector\sphinxhyphen{}style
variables can be referenced; the latter requires a bracketed term to specify
the \(I^\text{th}\) element of the vector calculated by the variable.
See the {\hyperref[\detokenize{variable::doc}]{\sphinxcrossref{\DUrole{doc}{variable}}}} command for details. Note that variables of
style \sphinxstyleemphasis{equal} or \sphinxstyleemphasis{vector} define a formula which can reference individual atom
properties or thermodynamic keywords, or they can invoke other computes, fixes,
or variables when they are evaluated, so this is a very general means of
specifying quantities to time correlate.


\bigskip\hrule\bigskip


\sphinxAtStartPar
Additional optional keywords also affect the operation of this fix.

\sphinxAtStartPar
The \sphinxstyleemphasis{type} keyword determines which pairs of input values are
correlated with each other.  For \(N\) input values \(V_i\),
with \(i \in \{1,\dotsc,N\}\), let the number of pairs be
\(N_\text{pair}\).  Note that the second value in the
pair, \(V_i(t) V_j(t+\Delta t)\), is always the one sampled at the later
time.
\begin{itemize}
\item {} 
\sphinxAtStartPar
If \sphinxstyleemphasis{type} is set to \sphinxstyleemphasis{auto} then each input value is correlated with
itself (i.e., \(C_{ii} = V_i^2\) for \(i \in \{1,\dotsc,N\}\),
so \(N_\text{pair} = N\)).

\item {} 
\sphinxAtStartPar
If \sphinxstyleemphasis{type} is set to \sphinxstyleemphasis{upper} then each input value is correlated with every
succeeding value (i.e., \(C_{ij} = V_i V_j\) for \(i < j\), so
\(N_\text{pair} = N (N-1)/2\)).

\item {} 
\sphinxAtStartPar
If \sphinxstyleemphasis{type} is set to \sphinxstyleemphasis{lower} then each input value is correlated with every
preceding value (i.e., \(C_{ij} = V_i V_j\) for \(i > j\), so
\(N_\text{pair} = N(N-1)/2\)).

\item {} 
\sphinxAtStartPar
If \sphinxstyleemphasis{type} is set to \sphinxstyleemphasis{auto/upper} then each input value is correlated
with itself and every succeeding value (i.e., \(C_{ij} = V_i V_j\)
for \(i \ge j\), so \(N_\text{pair} = N(N+1)/2\)).

\item {} 
\sphinxAtStartPar
If \sphinxstyleemphasis{type} is set to \sphinxstyleemphasis{auto/lower} then each input value is correlated
with itself and every preceding value (i.e., \(C_{ij} = V_i V_j\)
for \(i \le j\), so \(N_\text{pair} = N(N+1)/2\)).

\item {} 
\sphinxAtStartPar
If \sphinxstyleemphasis{type} is set to \sphinxstyleemphasis{full} then each input value is correlated with
itself and every other value (i.e., \(C_{ij} = V_i V_j\) for
\(\{i,j\} = \{1,N\}\), so \(N_\text{pair} = N^2\)).

\end{itemize}

\sphinxAtStartPar
The \sphinxstyleemphasis{ave} keyword determines what happens to the accumulation of correlation
samples every \(N_\text{freq}\) timesteps.  If the \sphinxstyleemphasis{ave} setting is \sphinxstyleemphasis{one},
then the accumulation is restarted or zeroed every \(N_\text{freq}\)
timesteps.  Thus the outputs on successive \(N_\text{freq}\) timesteps are
essentially independent of each other.  The exception is that the
\(C_{ij}(0) = V_i(t) V_j(t)\) value at a time step \(t,\) where
\(t\) is a multiple of \(N_\text{freq}\), contributes to the
correlation output both at time \(t\) and at time \(t+N_\text{freq}\).

\sphinxAtStartPar
If the \sphinxstyleemphasis{ave} setting is \sphinxstyleemphasis{running}, then the accumulation is never zeroed.
Thus the output of correlation data at any timestep is the average over samples
accumulated every \(N_\text{every}\) steps since the fix was defined.
It can only be restarted by deleting the fix via the {\hyperref[\detokenize{unfix::doc}]{\sphinxcrossref{\DUrole{doc}{unfix}}}}
command, or by re\sphinxhyphen{}defining the fix by re\sphinxhyphen{}specifying it.

\sphinxAtStartPar
The \sphinxstyleemphasis{start} keyword specifies what time step the accumulation of
correlation samples will begin on.  The default is step 0.  Setting it
to a larger value can avoid adding non\sphinxhyphen{}equilibrated data to the
correlation averages.

\sphinxAtStartPar
The \sphinxstyleemphasis{prefactor} keyword specifies a constant which will be used as a multiplier
on the correlation data after it is averaged.  It is effectively a scale factor
on \(V_i V_j\), which can be used to account for the size of the time
window or other unit conversions.

\sphinxAtStartPar
The \sphinxstyleemphasis{file} keyword allows a filename to be specified.  Every
\(N_\text{freq}\) steps, an array of correlation data is written to the
file.  The number of rows is \(N_\text{repeat}\), as described above.
The number of columns is \(N_\text{pair}+2\), also as described above.
Thus the file ends up to be a series of these array sections.

\sphinxAtStartPar
The \sphinxstyleemphasis{overwrite} keyword will continuously overwrite the output file
with the latest output, so that it only contains one timestep worth of
output.  This option can only be used with the \sphinxstyleemphasis{ave running} setting.

\sphinxAtStartPar
The \sphinxstyleemphasis{title1}, \sphinxstyleemphasis{title2}, and \sphinxstyleemphasis{title3} keywords allow specification of
the strings that will be printed as the first three lines of the output file,
assuming the \sphinxstyleemphasis{file} keyword was used.  LAMMPS uses default values for each of
these, so they do not need to be specified.

\sphinxAtStartPar
By default, these header lines are as follows:
\begin{sphinxalltt}
\# Time\sphinxhyphen{}correlated data for fix ID
\# TimeStep Number\sphinxhyphen{}of\sphinxhyphen{}time\sphinxhyphen{}windows
\# Index TimeDelta Ncount valueI*valueJ valueI*valueJ ...
\end{sphinxalltt}

\sphinxAtStartPar
In the first line, ID is replaced with the fix\sphinxhyphen{}ID.  The second line
describes the two values that are printed at the first of each section
of output.  In the third line the value pairs are replaced with the
appropriate fields from the fix ave/correlate command.


\bigskip\hrule\bigskip


\sphinxAtStartPar
Let \(S_{ij}\) be a set of time correlation data for input values
\(I\) and \(J\), namely the \(N_\text{repeat}\) values:
\begin{equation*}
\begin{split}S_{ij} = C_{ij}(0), C_{ij}(N_\text{every}), C_{ij}(2N_\text{every}),
 \dotsc, C_{ijI}\bigl((N_\text{repeat}-1) N_\text{every}\bigr)\end{split}
\end{equation*}
\sphinxAtStartPar
As explained below, these data are output as one column of a global
array, which is effectively the correlation matrix.

\sphinxAtStartPar
The \sphinxstyleemphasis{trap} function defined for {\hyperref[\detokenize{variable::doc}]{\sphinxcrossref{\DUrole{doc}{equal\sphinxhyphen{}style variables}}}}
can be used to perform a time integration of this vector of data,
using a trapezoidal rule.  This is useful for calculating various
quantities which can be derived from time correlation data.  If a
normalization factor is needed for the time integration, it can be
included in the variable formula or via the \sphinxstyleemphasis{prefactor} keyword.


\bigskip\hrule\bigskip



\subsection{Restart, fix\_modify, output, run start/stop, minimize info}
\label{\detokenize{fix_ave_correlate:restart-fix-modify-output-run-start-stop-minimize-info}}
\sphinxAtStartPar
No information about this fix is written to
{\hyperref[\detokenize{restart::doc}]{\sphinxcrossref{\DUrole{doc}{binary restart files}}}}.  None of the
{\hyperref[\detokenize{fix_modify::doc}]{\sphinxcrossref{\DUrole{doc}{fix\_modify}}}} options are relevant to this fix.

\sphinxAtStartPar
This fix computes a global array of values which can be accessed by
various {\hyperref[\detokenize{Howto_output::doc}]{\sphinxcrossref{\DUrole{doc}{output commands}}}}.  The values can only be
accessed on timesteps that are multiples of \(N_\text{freq}\) since that is
when averaging is performed.  The global array has \# of rows
\(N_\text{repeat}\) and \# of columns \(N_\text{pair}+2\).  The first
column has the time \(\Delta t\) (in time steps) between the pairs of input
values used to calculate the correlation, as described above.  The second
column has the number of samples contributing to the correlation average, as
described above.  The remaining Npair columns are for \(I,J\) pairs of the
\(N\) input values, as determined by the \sphinxstyleemphasis{type} keyword, as described
above.
\begin{itemize}
\item {} 
\sphinxAtStartPar
For \sphinxstyleemphasis{type} = \sphinxstyleemphasis{auto}, the \(N_\text{pair} = N\) columns are ordered:
\(C_{11}, C_{22}, \dotsc, C_{NN}\)

\item {} 
\sphinxAtStartPar
For \sphinxstyleemphasis{type} = \sphinxstyleemphasis{upper}, the \(N_\text{pair} = N(N-1)/2\) columns are
ordered: \(C_{12}, C_{13}, \dotsc, C_{1N}, C_{23}, \dotsc, C_{2N},
C_{34}, \dotsc, C_{N-1,N}\)

\item {} 
\sphinxAtStartPar
For \sphinxstyleemphasis{type} = \sphinxstyleemphasis{lower}, the \(N_\text{pair} = N(N-1)/2\) columns are
ordered: \(C_{21}, C_{31}, C_{32}, C_{41}, C_{42}, C_{43I}, \dotsc,
C_{N1}, C_{N2}, \dotsc, C_{N,N-1}\)

\item {} 
\sphinxAtStartPar
For \sphinxstyleemphasis{type} = \sphinxstyleemphasis{auto/upper}, the \(N_\text{pair} = N(N+1)/2\) columns are
ordered: \(C_{11}, C_{12}, C_{13}, \dotsc, C_{1N}, C_{22}, C_{23},
\dotsc, C_{2N}, C_{33}, C_{34}, \dotsc, C_{N-1,N}, C_{NN}\)

\item {} 
\sphinxAtStartPar
For \sphinxstyleemphasis{type} = \sphinxstyleemphasis{auto/lower}, the \(N_\text{pair} = N(N+1)/2\) columns are
ordered: \(C_{11}, C_{21}, C_{22}, C_{31}, C_{32}, C_{33}, C_{41},
\dotsc, C_{44}, C_{N1}, C_{N2}, \dotsc, C_{N,N-1}, C_{NN}\)

\item {} 
\sphinxAtStartPar
For \sphinxstyleemphasis{type} = \sphinxstyleemphasis{full}, the \(N_\text{pair} = N^2\) columns are ordered:
\(C_{11}, C_{12}, \dotsc, C_{1N}, C_{21}, C_{22}, \dotsc, C_{2N},
C_{31}, \dotsc, C_{3N}, \dotsc, C_{N1}, \dotsc, C_{N,N-1}, C_{NN}\)

\end{itemize}

\sphinxAtStartPar
The array values calculated by this fix are treated as extensive.  If
you need to divide them by the number of atoms, you must do this in a
later processing step (e.g., when using them in a {\hyperref[\detokenize{variable::doc}]{\sphinxcrossref{\DUrole{doc}{variable}}}}).

\sphinxAtStartPar
No parameter of this fix can be used with the \sphinxstyleemphasis{start/stop} keywords of
the {\hyperref[\detokenize{run::doc}]{\sphinxcrossref{\DUrole{doc}{run}}}} command.  This fix is not invoked during {\hyperref[\detokenize{minimize::doc}]{\sphinxcrossref{\DUrole{doc}{energy minimization}}}}.


\subsection{Restrictions}
\label{\detokenize{fix_ave_correlate:restrictions}}\begin{quote}

\sphinxAtStartPar
none
\end{quote}


\subsection{Related commands}
\label{\detokenize{fix_ave_correlate:related-commands}}
\sphinxAtStartPar
{\hyperref[\detokenize{fix_ave_correlate_long::doc}]{\sphinxcrossref{\DUrole{doc}{fix ave/correlate/long}}}},
{\hyperref[\detokenize{compute::doc}]{\sphinxcrossref{\DUrole{doc}{compute}}}}, {\hyperref[\detokenize{fix_ave_time::doc}]{\sphinxcrossref{\DUrole{doc}{fix ave/time}}}},
{\hyperref[\detokenize{fix_ave_atom::doc}]{\sphinxcrossref{\DUrole{doc}{fix ave/atom}}}}, {\hyperref[\detokenize{fix_ave_chunk::doc}]{\sphinxcrossref{\DUrole{doc}{fix ave/chunk}}}},
{\hyperref[\detokenize{fix_ave_histo::doc}]{\sphinxcrossref{\DUrole{doc}{fix ave/histo}}}}, {\hyperref[\detokenize{variable::doc}]{\sphinxcrossref{\DUrole{doc}{variable}}}}


\subsection{Default}
\label{\detokenize{fix_ave_correlate:default}}
\sphinxAtStartPar
none

\sphinxAtStartPar
The option defaults are ave = one, type = auto, start = 0, no file
output, title 1,2,3 = strings as described above, and prefactor = 1.0.

\sphinxstepscope

\index{fix ave/correlate/long@\spxentry{fix ave/correlate/long}}\ignorespaces 

\section{fix ave/correlate/long command}
\label{\detokenize{fix_ave_correlate_long:fix-ave-correlate-long-command}}\label{\detokenize{fix_ave_correlate_long:index-0}}\label{\detokenize{fix_ave_correlate_long::doc}}

\subsection{Syntax}
\label{\detokenize{fix_ave_correlate_long:syntax}}
\begin{sphinxVerbatim}[commandchars=\\\{\}]
\PYG{k}{fix }\PYG{n+nv+nvIdentifier}{ID}\PYG{+w}{ }\PYG{n+nv+nvIdentifier}{group\PYGZhy{}ID}\PYG{+w}{ }\PYG{n}{ave}\PYG{o}{/}\PYG{n}{correlate}\PYG{o}{/}\PYG{n}{long}\PYG{+w}{ }\PYG{n}{Nevery}\PYG{+w}{ }\PYG{n}{Nfreq}\PYG{+w}{ }\PYG{n}{value1}\PYG{+w}{ }\PYG{n}{value2}\PYG{+w}{ }\PYG{n}{...}\PYG{+w}{ }\PYG{n}{keyword}\PYG{+w}{ }\PYG{n}{args}\PYG{+w}{ }\PYG{n}{...}
\end{sphinxVerbatim}
\begin{itemize}
\item {} 
\sphinxAtStartPar
ID, group\sphinxhyphen{}ID are documented in {\hyperref[\detokenize{fix::doc}]{\sphinxcrossref{\DUrole{doc}{fix}}}} command

\item {} 
\sphinxAtStartPar
ave/correlate/long = style name of this fix command

\item {} 
\sphinxAtStartPar
Nevery = use input values every this many time steps

\item {} 
\sphinxAtStartPar
Nfreq = save state of the time correlation functions every this many time steps

\item {} 
\sphinxAtStartPar
one or more input values can be listed

\item {} 
\sphinxAtStartPar
value = c\_ID, c\_ID{[}N{]}, f\_ID, f\_ID{[}N{]}, v\_name, v\_name{[}I{]}

\begin{sphinxVerbatim}[commandchars=\\\{\}]
c\PYGZus{}ID = global scalar calculated by a compute with ID
c\PYGZus{}ID[I] = Ith component of global vector calculated by a compute with ID, I can include wildcard (see below)
f\PYGZus{}ID = global scalar calculated by a fix with ID
f\PYGZus{}ID[I] = Ith component of global vector calculated by a fix with ID, I can include wildcard (see below)
v\PYGZus{}name = global value calculated by an equal\PYGZhy{}style variable with name
v\PYGZus{}name[I] = Ith component of a vector\PYGZhy{}style variable with name, I can include wildcard (see below)
\end{sphinxVerbatim}

\item {} 
\sphinxAtStartPar
zero or more keyword/arg pairs may be appended

\item {} 
\sphinxAtStartPar
keyword = \sphinxstyleemphasis{type} or \sphinxstyleemphasis{start} or \sphinxstyleemphasis{file} or \sphinxstyleemphasis{overwrite} or \sphinxstyleemphasis{title1} or \sphinxstyleemphasis{title2} or \sphinxstyleemphasis{ncorr} or \sphinxstyleemphasis{nlen} or \sphinxstyleemphasis{ncount}
\begin{sphinxalltt}
\sphinxstyleemphasis{type} arg = \sphinxstyleemphasis{auto} or \sphinxstyleemphasis{upper} or \sphinxstyleemphasis{lower} or \sphinxstyleemphasis{auto/upper} or \sphinxstyleemphasis{auto/lower} or \sphinxstyleemphasis{full}
  auto = correlate each value with itself
  upper = correlate each value with each succeeding value
  lower = correlate each value with each preceding value
  auto/upper = auto + upper
  auto/lower = auto + lower
  full = correlate each value with every other value, including itself = auto + upper + lower
\sphinxstyleemphasis{start} args = Nstart
  Nstart = start accumulating correlations on this time step
\sphinxstyleemphasis{file} arg = filename
  filename = name of file to output correlation data to
\sphinxstyleemphasis{overwrite} arg = none = overwrite output file with only latest output
\sphinxstyleemphasis{title1} arg = string
  string = text to print as 1st line of output file
\sphinxstyleemphasis{title2} arg = string
  string = text to print as 2nd line of output file
\sphinxstyleemphasis{ncorr} arg = Ncorrelators
  Ncorrelators = number of correlators to store
\sphinxstyleemphasis{nlen} args = Nlen
  Nlen = length of each correlator
\sphinxstyleemphasis{ncount} args = Ncount
  Ncount = number of values over which successive correlators are averaged
\end{sphinxalltt}

\end{itemize}


\subsection{Examples}
\label{\detokenize{fix_ave_correlate_long:examples}}
\begin{sphinxVerbatim}[commandchars=\\\{\}]
\PYG{k}{fix }\PYG{n+nv+nvIdentifier}{1}\PYG{+w}{ }\PYG{n+nv+nvIdentifier}{all}\PYG{+w}{ }\PYG{n}{ave}\PYG{o}{/}\PYG{n}{correlate}\PYG{o}{/}\PYG{n}{long}\PYG{+w}{ }\PYG{l+m}{5}\PYG{+w}{ }\PYG{l+m}{1000}\PYG{+w}{ }\PYG{n}{c\PYGZus{}myTemp}\PYG{+w}{ }\PYG{n}{file}\PYG{+w}{ }\PYG{n}{temp.correlate}
\PYG{k}{fix }\PYG{n+nv+nvIdentifier}{1}\PYG{+w}{ }\PYG{n+nv+nvIdentifier}{all}\PYG{+w}{ }\PYG{n}{ave}\PYG{o}{/}\PYG{n}{correlate}\PYG{o}{/}\PYG{n}{long}\PYG{+w}{ }\PYG{l+m}{1}\PYG{+w}{ }\PYG{l+m}{10000}\PYG{+w}{ }\PYG{o}{\PYGZam{}}
\PYG{+w}{          }\PYG{n}{c\PYGZus{}thermo\PYGZus{}press[1]}\PYG{+w}{ }\PYG{n}{c\PYGZus{}thermo\PYGZus{}press[2]}\PYG{+w}{ }\PYG{n}{c\PYGZus{}thermo\PYGZus{}press[3]}\PYG{+w}{ }\PYG{o}{\PYGZam{}}
\PYG{+w}{          }\PYG{n}{type}\PYG{+w}{ }\PYG{n}{upper}\PYG{+w}{ }\PYG{n}{title1}\PYG{+w}{ }\PYG{l+s}{\PYGZdq{}}\PYG{l+s}{My correlation data}\PYG{l+s}{\PYGZdq{}}\PYG{+w}{ }\PYG{n}{nlen}\PYG{+w}{ }\PYG{l+m}{15}\PYG{+w}{ }\PYG{n}{ncount}\PYG{+w}{ }\PYG{l+m}{3}
\PYG{k}{fix }\PYG{n+nv+nvIdentifier}{1}\PYG{+w}{ }\PYG{n+nv+nvIdentifier}{all}\PYG{+w}{ }\PYG{n}{ave}\PYG{o}{/}\PYG{n}{correlate}\PYG{o}{/}\PYG{n}{long}\PYG{+w}{ }\PYG{l+m}{1}\PYG{+w}{ }\PYG{l+m}{10000}\PYG{+w}{ }\PYG{n}{c\PYGZus{}thermo\PYGZus{}press[}\PYG{o}{*}\PYG{n}{]}
\end{sphinxVerbatim}


\subsection{Description}
\label{\detokenize{fix_ave_correlate_long:description}}
\sphinxAtStartPar
This fix is similar in spirit and syntax to the
{\hyperref[\detokenize{fix_ave_correlate::doc}]{\sphinxcrossref{\DUrole{doc}{fix ave/correlate}}}}.
However, this fix allows the efficient calculation of time correlation
functions on\sphinxhyphen{}the\sphinxhyphen{}fly over extremely long time windows with little
additional CPU overhead, using a multiple\sphinxhyphen{}\(\tau\) method
{\hyperref[\detokenize{fix_ave_correlate_long:ramirez}]{\sphinxcrossref{\DUrole{std,std-ref}{(Ramirez)}}}} that decreases the resolution of the stored
correlation function with time.  It is not a full drop\sphinxhyphen{}in replacement.

\sphinxAtStartPar
The group specified with this command is ignored.  However, note that
specified values may represent calculations performed by computes and
fixes which store their own “group” definitions.

\sphinxAtStartPar
Each listed value can be the result of a compute or fix or the
evaluation of an equal\sphinxhyphen{}style or vector\sphinxhyphen{}style variable.  For
vector\sphinxhyphen{}style variables, the specified indices can include a wildcard
character.  See the {\hyperref[\detokenize{fix_ave_correlate::doc}]{\sphinxcrossref{\DUrole{doc}{fix ave/correlate}}}} page
for details.

\sphinxAtStartPar
The \sphinxstyleemphasis{Nevery} and \sphinxstyleemphasis{Nfreq} arguments specify on what time steps the input
values will be used to calculate correlation data and the frequency
with which the time correlation functions will be output to a file,
respectively.
Note that there is no \sphinxstyleemphasis{Nrepeat} argument, unlike the
{\hyperref[\detokenize{fix_ave_correlate::doc}]{\sphinxcrossref{\DUrole{doc}{fix ave/correlate}}}} command.

\sphinxAtStartPar
The optional keywords \sphinxstyleemphasis{ncorr}, \sphinxstyleemphasis{nlen}, and \sphinxstyleemphasis{ncount} are unique to this
command and determine the number of correlation points calculated and
the memory and CPU overhead used by this calculation. \sphinxstyleemphasis{Nlen} and
\sphinxstyleemphasis{ncount} determine the amount of averaging done at longer correlation
times.  The default values \sphinxstyleemphasis{nlen} = 16 and \sphinxstyleemphasis{ncount} = 2 ensure that the
systematic error of the multiple\sphinxhyphen{}\(\tau\) correlator is always below the
level of the statistical error of a typical simulation (which depends
on the ensemble size and the simulation length).

\sphinxAtStartPar
The maximum correlation time (in time steps) that can be reached is
given by the formula \((nlen-1) ncount^{(ncorr-1)}\).  Longer correlation
times are discarded and not calculated.  With the default values of
the parameters (\(ncorr=20\), \(nlen=16\) and \(ncount=2\)),
this corresponds to 7864320 time steps.  If longer correlation times are
needed, the value of ncorr should be increased. Using \(nlen=16\) and
\(ncount=2\), with \(ncorr=30\), the maximum number of steps that can
be correlated is 80530636808.  If \(ncorr=40\), correlation times in excess
of \(8\times 10^{12}\) time steps can be calculated.

\sphinxAtStartPar
The total memory needed for each correlation pair is roughly
\(4 \times ncorr\times nlen \times 8\) bytes.
With the default values of the parameters, this corresponds to about 10 KB.

\sphinxAtStartPar
For the meaning of the additional optional keywords, see the
{\hyperref[\detokenize{fix_ave_correlate::doc}]{\sphinxcrossref{\DUrole{doc}{fix ave/correlate}}}} doc page.


\subsection{Restart, fix\_modify, output, run start/stop, minimize info}
\label{\detokenize{fix_ave_correlate_long:restart-fix-modify-output-run-start-stop-minimize-info}}
\sphinxAtStartPar
Contrary to {\hyperref[\detokenize{fix_ave_correlate::doc}]{\sphinxcrossref{\DUrole{doc}{fix ave/correlate}}}} this fix
does \sphinxstylestrong{not} provide access to its internal data to various output
options. Since this fix in intended for the calculation of time
correlation functions over very long MD simulations, the information
about this fix is written automatically to binary restart files, so
that the time correlation calculation can continue in subsequent
simulations. None of the fix\_modify options are relevant to this fix.

\sphinxAtStartPar
No parameter of this fix can be used with the start/stop keywords of
the run command. This fix is not invoked during energy minimization.


\subsection{Restrictions}
\label{\detokenize{fix_ave_correlate_long:restrictions}}
\sphinxAtStartPar
This compute is part of the EXTRA\sphinxhyphen{}FIX package.  It is only enabled if
LAMMPS was built with that package.  See the
{\hyperref[\detokenize{Build_package::doc}]{\sphinxcrossref{\DUrole{doc}{Build package}}}} page for more info.


\subsection{Related commands}
\label{\detokenize{fix_ave_correlate_long:related-commands}}
\sphinxAtStartPar
{\hyperref[\detokenize{fix_ave_correlate::doc}]{\sphinxcrossref{\DUrole{doc}{fix ave/correlate}}}}


\subsection{Default}
\label{\detokenize{fix_ave_correlate_long:default}}
\sphinxAtStartPar
none

\sphinxAtStartPar
The option defaults for keywords that are also keywords for the
{\hyperref[\detokenize{fix_ave_correlate::doc}]{\sphinxcrossref{\DUrole{doc}{fix ave/correlate}}}} command are as follows:
type = auto, start = 0, no file output, title 1,2 = strings as described on
the {\hyperref[\detokenize{fix_ave_correlate::doc}]{\sphinxcrossref{\DUrole{doc}{fix ave/correlate}}}} doc page.

\sphinxAtStartPar
The option defaults for keywords unique to this command are as
follows: ncorr=20, nlen=16, ncount=2.


\bigskip\hrule\bigskip

\phantomsection\label{\detokenize{fix_ave_correlate_long:ramirez}}
\sphinxAtStartPar
\sphinxstylestrong{(Ramirez)} J. Ramirez, S.K. Sukumaran, B. Vorselaars and
A.E. Likhtman, J. Chem. Phys. 133, 154103 (2010).

\sphinxstepscope

\index{fix ave/grid@\spxentry{fix ave/grid}}\ignorespaces 

\section{fix ave/grid command}
\label{\detokenize{fix_ave_grid:fix-ave-grid-command}}\label{\detokenize{fix_ave_grid:index-0}}\label{\detokenize{fix_ave_grid::doc}}

\subsection{Syntax}
\label{\detokenize{fix_ave_grid:syntax}}
\begin{sphinxVerbatim}[commandchars=\\\{\}]
\PYG{k}{fix }\PYG{n+nv+nvIdentifier}{ID}\PYG{+w}{ }\PYG{n+nv+nvIdentifier}{group\PYGZhy{}ID}\PYG{+w}{ }\PYG{n}{ave}\PYG{o}{/}\PYG{n}{grid}\PYG{+w}{ }\PYG{n}{Nevery}\PYG{+w}{ }\PYG{n}{Nrepeat}\PYG{+w}{ }\PYG{n}{Nfreq}\PYG{+w}{ }\PYG{n}{Nx}\PYG{+w}{ }\PYG{n}{Ny}\PYG{+w}{ }\PYG{n}{Nz}\PYG{+w}{ }\PYG{n}{value1}\PYG{+w}{ }\PYG{n}{value2}\PYG{+w}{ }\PYG{n}{...}\PYG{+w}{ }\PYG{n}{keyword}\PYG{+w}{ }\PYG{n}{args}\PYG{+w}{ }\PYG{n}{...}
\end{sphinxVerbatim}
\begin{itemize}
\item {} 
\sphinxAtStartPar
ID, group\sphinxhyphen{}ID are documented in {\hyperref[\detokenize{fix::doc}]{\sphinxcrossref{\DUrole{doc}{fix}}}} command

\item {} 
\sphinxAtStartPar
ave/grid = style name of this fix command

\item {} 
\sphinxAtStartPar
Nevery = use input values every this many timesteps

\item {} 
\sphinxAtStartPar
Nrepeat = \# of times to use input values for calculating averages

\item {} 
\sphinxAtStartPar
Nfreq = calculate averages every this many timesteps

\item {} 
\sphinxAtStartPar
Nx, Ny, Nz = grid size in each dimension

\item {} 
\sphinxAtStartPar
one or more per\sphinxhyphen{}atom or per\sphinxhyphen{}grid input values can be listed

\item {} 
\sphinxAtStartPar
per\sphinxhyphen{}atom value = vx, vy, vz, fx, fy, fz, density/mass, density/number, mass, temp, c\_ID, c\_ID{[}I{]}, f\_ID, f\_ID{[}I{]}, v\_name

\begin{sphinxVerbatim}[commandchars=\\\{\}]
vx,vy,vz,fx,fy,fz,mass = atom attribute (velocity, force component, mass)
density/number, density/mass = number or mass density (per volume)
temp = temperature
c\PYGZus{}ID = per\PYGZhy{}atom vector calculated by a compute with ID
c\PYGZus{}ID[I] = Ith column of per\PYGZhy{}atom array calculated by a compute with ID, I can include wildcard (see below)
f\PYGZus{}ID = per\PYGZhy{}atom vector calculated by a fix with ID
f\PYGZus{}ID[I] = Ith column of per\PYGZhy{}atom array calculated by a fix with ID, I can include wildcard (see below)
v\PYGZus{}name = per\PYGZhy{}atom vector calculated by an atom\PYGZhy{}style variable with name
\end{sphinxVerbatim}

\item {} 
\sphinxAtStartPar
per\sphinxhyphen{}grid value = c\_ID:gname:dname, c\_ID:gname:dname{[}I{]}, f\_ID:gname:dname, f\_ID:gname:dname{[}I{]}

\begin{sphinxVerbatim}[commandchars=\\\{\}]
gname = name of grid defined by compute or fix
dname = name of data field defined by compute or fix
c\PYGZus{}ID = per\PYGZhy{}grid vector calculated by a compute with ID
c\PYGZus{}ID[I] = Ith column of per\PYGZhy{}grid array calculated by a compute with ID, I can include wildcard (see below)
f\PYGZus{}ID = per\PYGZhy{}grid vector calculated by a fix with ID
f\PYGZus{}ID[I] = Ith column of per\PYGZhy{}grid array calculated by a fix with ID, I can include wildcard (see below)
\end{sphinxVerbatim}

\item {} 
\sphinxAtStartPar
zero or more keyword/arg pairs may be appended

\item {} 
\sphinxAtStartPar
keyword = \sphinxstyleemphasis{discard} or \sphinxstyleemphasis{norm} or \sphinxstyleemphasis{ave} or \sphinxstyleemphasis{bias} or \sphinxstyleemphasis{adof} or \sphinxstyleemphasis{cdof}
\begin{sphinxalltt}
\sphinxstyleemphasis{discard} arg = \sphinxstyleemphasis{yes} or \sphinxstyleemphasis{no}
  yes = discard an atom outside grid in a non\sphinxhyphen{}periodic dimension
  no = remap an atom outside grid in a non\sphinxhyphen{}periodic dimension to first or last grid cell
\sphinxstyleemphasis{norm} arg = \sphinxstyleemphasis{all} or \sphinxstyleemphasis{sample} or \sphinxstyleemphasis{none} = how output on \sphinxstyleemphasis{Nfreq} steps is normalized
  all = output is sum of atoms across all \sphinxstyleemphasis{Nrepeat} samples, divided by atom count
  sample = output is sum of \sphinxstyleemphasis{Nrepeat} sample averages, divided by \sphinxstyleemphasis{Nrepeat}
  none = output is sum of \sphinxstyleemphasis{Nrepeat} sample sums, divided by \sphinxstyleemphasis{Nrepeat}
\sphinxstyleemphasis{ave} args = \sphinxstyleemphasis{one} or \sphinxstyleemphasis{running} or \sphinxstyleemphasis{window M}
  one = output new average value every Nfreq steps
  running = output cumulative average of all previous Nfreq steps
  window M = output average of M most recent Nfreq steps
\sphinxstyleemphasis{bias} arg = bias\sphinxhyphen{}ID
  bias\sphinxhyphen{}ID = ID of a temperature compute that removes a velocity bias for temperature calculation
\sphinxstyleemphasis{adof} value = dof\_per\_atom
  dof\_per\_atom = define this many degrees\sphinxhyphen{}of\sphinxhyphen{}freedom per atom for temperature calculation
\sphinxstyleemphasis{cdof} value = dof\_per\_grid\_cell
  dof\_per\_grid\_cell = add this many degrees\sphinxhyphen{}of\sphinxhyphen{}freedom per grid\_cell for temperature calculation
\end{sphinxalltt}

\end{itemize}


\subsection{Examples}
\label{\detokenize{fix_ave_grid:examples}}
\begin{sphinxVerbatim}[commandchars=\\\{\}]
\PYG{k}{fix }\PYG{n+nv+nvIdentifier}{1}\PYG{+w}{ }\PYG{n+nv+nvIdentifier}{all}\PYG{+w}{ }\PYG{n}{ave}\PYG{o}{/}\PYG{n}{grid}\PYG{+w}{ }\PYG{l+m}{10000}\PYG{+w}{ }\PYG{l+m}{1}\PYG{+w}{ }\PYG{l+m}{10000}\PYG{+w}{ }\PYG{l+m}{10}\PYG{+w}{ }\PYG{l+m}{10}\PYG{+w}{ }\PYG{l+m}{10}\PYG{+w}{ }\PYG{n}{fx}\PYG{+w}{ }\PYG{n}{fy}\PYG{+w}{ }\PYG{n}{fz}\PYG{+w}{ }\PYG{n}{c\PYGZus{}myMSD[}\PYG{o}{*}\PYG{n}{]}
\PYG{k}{fix }\PYG{n+nv+nvIdentifier}{1}\PYG{+w}{ }\PYG{n+nv+nvIdentifier}{flow}\PYG{+w}{ }\PYG{n}{ave}\PYG{o}{/}\PYG{n}{grid}\PYG{+w}{ }\PYG{l+m}{100}\PYG{+w}{ }\PYG{l+m}{10}\PYG{+w}{ }\PYG{l+m}{1000}\PYG{+w}{ }\PYG{l+m}{20}\PYG{+w}{ }\PYG{l+m}{20}\PYG{+w}{ }\PYG{l+m}{30}\PYG{+w}{ }\PYG{n}{f\PYGZus{}TTM}:grid:data
\end{sphinxVerbatim}


\subsection{Description}
\label{\detokenize{fix_ave_grid:description}}
\sphinxAtStartPar
Overlay the 2d or 3d simulation box with a uniformly spaced 2d or 3d
grid and use it to either (a) time\sphinxhyphen{}average per\sphinxhyphen{}atom quantities for the
atoms in each grid cell, or to (b) time\sphinxhyphen{}average per\sphinxhyphen{}grid quantities
produced by other computes or fixes.  This fix operates in either
“per\sphinxhyphen{}atom mode” (all input values are per\sphinxhyphen{}atom) or in “per\sphinxhyphen{}grid mode”
(all input values are per\sphinxhyphen{}grid).  You cannot use both per\sphinxhyphen{}atom and
per\sphinxhyphen{}grid inputs in the same command.

\sphinxAtStartPar
The grid created by this command is distributed; each processor owns
the grid points that are within its subdomain.  This is similar to
the {\hyperref[\detokenize{fix_ave_chunk::doc}]{\sphinxcrossref{\DUrole{doc}{fix ave/chunk}}}} command when it uses chunks
from the {\hyperref[\detokenize{compute_chunk_atom::doc}]{\sphinxcrossref{\DUrole{doc}{compute chunk/atom}}}} command which
are 2d or 3d regular bins.  However, the per\sphinxhyphen{}bin outputs in that case
are global; each processor stores a copy of the entire set of bin
data.  Thus it more efficient to use the fix ave/grid command when the
grid is large and a simulation is run on many processors.

\sphinxAtStartPar
For per\sphinxhyphen{}atom mode, only atoms in the specified group contribute to the
summing and averaging calculations.  For per\sphinxhyphen{}grid mode, the specified
group is ignored.


\bigskip\hrule\bigskip


\sphinxAtStartPar
The \(N_\text{every}\), \(N_\text{repeat}\), and \(N_\text{freq}\)
arguments specify on what time steps the input values will be accessed and
contribute to the average.  The final averaged quantities are generated on time
steps that are a multiples of \(N_\text{freq}\).  The average is over
\(N_\text{repeat}\) quantities, computed in the preceding portion of the
simulation every \(N_\text{every}\) time steps.  \(N_\text{freq}\)
must be a multiple of \(N_\text{every}\) and \(N_\text{every}\) must be
non\sphinxhyphen{}zero even if \(N_\text{repeat} = 1\).  Also, the time steps
contributing to the average value cannot overlap (i.e.,
\(N_\text{repeat} \times N_\text{every}\) cannot exceed \(N_\text{freq}\)).

\sphinxAtStartPar
For example, if \(N_\text{every}=2\), \(N_\text{repeat}=6\), and
\(N_\text{freq}=100\), then values on
time steps 90,92,94,96,98,100 will be used to compute the final average
on timestep 100.  Similarly for timesteps 190,192,194,196,198,200 on
timestep 200, etc.  If \(N_\text{repeat}=1\) and \(N_\text{freq} = 100\), then no time
averaging is done; values are simply generated on timesteps
100,200,etc.

\sphinxAtStartPar
In per\sphinxhyphen{}atom mode, each input value can also be averaged over the atoms
in each grid cell.  The way the averaging is done across the \(N_\text{repeat}\)
timesteps to produce output on the \(N_\text{freq}\) timesteps, and across
multiple \(N_\text{freq}\) outputs, is determined by the \sphinxstyleemphasis{norm} and \sphinxstyleemphasis{ave}
keyword settings, as discussed below.


\bigskip\hrule\bigskip


\sphinxAtStartPar
The \sphinxstyleemphasis{Nx}, \sphinxstyleemphasis{Ny}, and \sphinxstyleemphasis{Nz} arguments specify the size of the grid that
overlays the simulation box.  For 2d simulations, \sphinxstyleemphasis{Nz} must be 1.  The
\sphinxstyleemphasis{Nx}, \sphinxstyleemphasis{Ny}, \sphinxstyleemphasis{Nz} values can be any positive integer.  The grid can be
very coarse compared to the particle count, or very fine.  If one or
more of the values = 1, then bins are 2d planes or 1d slices of the
simulation domain.  Note that if the total number of grid cells is
small, it may be more efficient to use the {\hyperref[\detokenize{fix_ave_chunk::doc}]{\sphinxcrossref{\DUrole{doc}{fix ave/chunk}}}} command which can treat a grid defined by the
{\hyperref[\detokenize{compute_chunk_atom::doc}]{\sphinxcrossref{\DUrole{doc}{compute chunk/atom}}}} command as a global
grid where each processor owns a copy of all the grid cells.  If \sphinxstyleemphasis{Nx}
= \sphinxstyleemphasis{Ny} = \sphinxstyleemphasis{Nz} = 1 is used, the same calculation would be more
efficiently performed by the {\hyperref[\detokenize{fix_ave_atom::doc}]{\sphinxcrossref{\DUrole{doc}{fix ave/atom}}}}
command.

\sphinxAtStartPar
If the simulation box size or shape changes during a simulation, the
grid always conforms to the size/shape of the current simulation box.
If one more dimensions have non\sphinxhyphen{}periodic shrink\sphinxhyphen{}wrapped boundary
conditions, as defined by the {\hyperref[\detokenize{boundary::doc}]{\sphinxcrossref{\DUrole{doc}{boundary}}}} command, then
the grid will extend over the (dynamic) shrink\sphinxhyphen{}wrapped extent in each
dimension.  If the box shape is triclinic, as explained in {\hyperref[\detokenize{Howto_triclinic::doc}]{\sphinxcrossref{\DUrole{doc}{Howto
triclinic}}}}, then the grid is also triclinic; each
grid cell is a small triclinic cell with the same shape as the
simulation box.


\bigskip\hrule\bigskip


\sphinxAtStartPar
In both per\sphinxhyphen{}atom and per\sphinxhyphen{}grid mode, input values from a compute or fix
that produces an array of values (multiple values per atom or per grid
point), the bracketed index I can be specified using a wildcard
asterisk with the index to effectively specify multiple values.  This
takes the form “*” or “*n” or “n*” or “m*n”.  If N = the number of
columns in the array (for \sphinxstyleemphasis{mode} = vector), then an asterisk with no
numeric values means all indices from 1 to N.  A leading asterisk
means all indices from 1 to n (inclusive).  A trailing asterisk means
all indices from n to N (inclusive).  A middle asterisk means all
indices from m to n (inclusive).

\sphinxAtStartPar
Using a wildcard is the same as if the individual columns of the array
had been listed one by one.  E.g. if there were a compute fft/grid
command which produced 3 values for each grid point, these two fix
ave/grid commands would be equivalent:

\begin{sphinxVerbatim}[commandchars=\\\{\}]
\PYG{k}{compute }\PYG{n+nv+nvIdentifier}{myFFT}\PYG{+w}{ }\PYG{n+nv+nvIdentifier}{all}\PYG{+w}{ }\PYG{n}{fft}\PYG{o}{/}\PYG{n}{grid}\PYG{+w}{ }\PYG{l+m}{10}\PYG{+w}{ }\PYG{l+m}{10}\PYG{+w}{ }\PYG{l+m}{10}\PYG{+w}{ }\PYG{n}{...}
\PYG{k}{fix }\PYG{n+nv+nvIdentifier}{1}\PYG{+w}{ }\PYG{n+nv+nvIdentifier}{all}\PYG{+w}{ }\PYG{n}{ave}\PYG{o}{/}\PYG{n}{grid}\PYG{+w}{ }\PYG{l+m}{100}\PYG{+w}{ }\PYG{l+m}{1}\PYG{+w}{ }\PYG{l+m}{100}\PYG{+w}{ }\PYG{l+m}{10}\PYG{+w}{ }\PYG{l+m}{10}\PYG{+w}{ }\PYG{l+m}{10}\PYG{+w}{ }\PYG{n}{c\PYGZus{}myFFT}:grid:data\PYG{n}{[}\PYG{o}{*}\PYG{n}{]}
\PYG{k}{fix }\PYG{n+nv+nvIdentifier}{2}\PYG{+w}{ }\PYG{n+nv+nvIdentifier}{all}\PYG{+w}{ }\PYG{n}{ave}\PYG{o}{/}\PYG{n}{grid}\PYG{+w}{ }\PYG{l+m}{100}\PYG{+w}{ }\PYG{l+m}{1}\PYG{+w}{ }\PYG{l+m}{100}\PYG{+w}{ }\PYG{l+m}{10}\PYG{+w}{ }\PYG{l+m}{10}\PYG{+w}{ }\PYG{l+m}{10}\PYG{+w}{ }\PYG{n}{c\PYGZus{}myFFT}:grid:data\PYG{n}{[}\PYG{o}{*}\PYG{n}{][1]}\PYG{+w}{ }\PYG{n}{c\PYGZus{}myFFT}:grid:data\PYG{n}{[}\PYG{o}{*}\PYG{n}{][2]}\PYG{+w}{ }\PYG{n}{c\PYGZus{}myFFT}:grid:data\PYG{n}{[3]}
\end{sphinxVerbatim}


\bigskip\hrule\bigskip


\sphinxAtStartPar
\sphinxstyleemphasis{Per\sphinxhyphen{}atom mode}:

\sphinxAtStartPar
Each specified per\sphinxhyphen{}atom value can be an atom attribute (velocity,
force component), a number or mass density, a mass or temperature, or
the result of a {\hyperref[\detokenize{compute::doc}]{\sphinxcrossref{\DUrole{doc}{compute}}}} or {\hyperref[\detokenize{fix::doc}]{\sphinxcrossref{\DUrole{doc}{fix}}}} or the
evaluation of an atom\sphinxhyphen{}style {\hyperref[\detokenize{variable::doc}]{\sphinxcrossref{\DUrole{doc}{variable}}}}.  In the latter
cases, the compute, fix, or variable must produce a per\sphinxhyphen{}atom quantity,
not a global quantity.  Note that the {\hyperref[\detokenize{compute_property_atom::doc}]{\sphinxcrossref{\DUrole{doc}{compute property/atom}}}} command provides access to any attribute
defined and stored by atoms.

\sphinxAtStartPar
The per\sphinxhyphen{}atom values of each input vector are summed and averaged
independently of the per\sphinxhyphen{}atom values in other input vectors.

\sphinxAtStartPar
{\hyperref[\detokenize{compute::doc}]{\sphinxcrossref{\DUrole{doc}{Computes}}}} that produce per\sphinxhyphen{}atom quantities are those
which have the word \sphinxstyleemphasis{atom} in their style name.  See the doc pages for
individual {\hyperref[\detokenize{fix::doc}]{\sphinxcrossref{\DUrole{doc}{fixes}}}} to determine which ones produce per\sphinxhyphen{}atom
quantities.  {\hyperref[\detokenize{variable::doc}]{\sphinxcrossref{\DUrole{doc}{Variables}}}} of style \sphinxstyleemphasis{atom} are the only
ones that can be used with this fix since all other styles of variable
produce global quantities.


\bigskip\hrule\bigskip


\sphinxAtStartPar
The atom attribute values (vx,vy,vz,fx,fy,fz,mass) are
self\sphinxhyphen{}explanatory.  As noted above, any other atom attributes can be
used as input values to this fix by using the {\hyperref[\detokenize{compute_property_atom::doc}]{\sphinxcrossref{\DUrole{doc}{compute
property/atom}}}} command and then specifying an
input value from that compute.

\sphinxAtStartPar
The \sphinxstyleemphasis{density/number} value means the number density is computed for
each grid cell, i.e. number/volume.  The \sphinxstyleemphasis{density/mass} value means
the mass density is computed for each grid/cell,
i.e. total\sphinxhyphen{}mass/volume.  The output values are in units of 1/volume or
density (mass/volume).  See the {\hyperref[\detokenize{units::doc}]{\sphinxcrossref{\DUrole{doc}{units}}}} command page for
the definition of density for each choice of units, e.g. gram/cm\textasciicircum{}3.

\sphinxAtStartPar
The \sphinxstyleemphasis{temp} value computes the temperature for each grid cell, by the
formula
\begin{equation*}
\begin{split}\text{KE} = \frac{\text{DOF}}{2} k_B T,\end{split}
\end{equation*}
\sphinxAtStartPar
where KE = total kinetic energy of the atoms in the grid cell (
\(\frac{1}{2} m v^2\)), DOF = the total number of degrees of
freedom for all atoms in the grid cell, \(k_B\) = Boltzmann
constant, and \(T\) = temperature.

\sphinxAtStartPar
The DOF is calculated as N*adof + cdof, where N = number of atoms in
the grid cell, adof = degrees of freedom per atom, and cdof = degrees
of freedom per grid cell.  By default adof = 2 or 3 = dimensionality
of system, as set via the {\hyperref[\detokenize{dimension::doc}]{\sphinxcrossref{\DUrole{doc}{dimension}}}} command, and
cdof = 0.0.  This gives the usual formula for temperature.

\sphinxAtStartPar
Note that currently this temperature only includes translational
degrees of freedom for each atom.  No rotational degrees of freedom
are included for finite\sphinxhyphen{}size particles.  Also no degrees of freedom
are subtracted for any velocity bias or constraints that are applied,
such as {\hyperref[\detokenize{compute_temp_partial::doc}]{\sphinxcrossref{\DUrole{doc}{compute temp/partial}}}}, or
{\hyperref[\detokenize{fix_shake::doc}]{\sphinxcrossref{\DUrole{doc}{fix shake}}}} or {\hyperref[\detokenize{fix_rigid::doc}]{\sphinxcrossref{\DUrole{doc}{fix rigid}}}}.  This is
because those degrees of freedom (e.g. a constrained bond) could apply
to sets of atoms that are both inside and outside a specific grid
cell, and hence the concept is somewhat ill\sphinxhyphen{}defined.  In some cases,
you can use the \sphinxstyleemphasis{adof} and \sphinxstyleemphasis{cdof} keywords to adjust the calculated
degrees of freedom appropriately, as explained below.

\sphinxAtStartPar
Also note that a bias can be subtracted from atom velocities before
they are used in the above formula for KE, by using the \sphinxstyleemphasis{bias}
keyword.  This allows, for example, a thermal temperature to be
computed after removal of a flow velocity profile.

\sphinxAtStartPar
Note that the per\sphinxhyphen{}grid\sphinxhyphen{}cell temperature calculated by this fix and the
{\hyperref[\detokenize{compute_temp_chunk::doc}]{\sphinxcrossref{\DUrole{doc}{compute temp/chunk}}}} command (using bins)
can be different.  The compute calculates the temperature for each
chunk for a single snapshot.  This fix can do that but can also time
average those values over many snapshots, or it can compute a
temperature as if the atoms in the grid cell on different timesteps
were collected together as one set of atoms to calculate their
temperature.  The compute allows the center\sphinxhyphen{}of\sphinxhyphen{}mass velocity of each
chunk to be subtracted before calculating the temperature; this fix
does not.

\sphinxAtStartPar
If a value begins with “c\_”, a compute ID must follow which has been
previously defined in the input script.  If no bracketed integer is
appended, the per\sphinxhyphen{}atom vector calculated by the compute is used.  If a
bracketed integer is appended, the Ith column of the per\sphinxhyphen{}atom array
calculated by the compute is used.  Users can also write code for
their own compute styles and {\hyperref[\detokenize{Modify::doc}]{\sphinxcrossref{\DUrole{doc}{add them to LAMMPS}}}}.  See
the discussion above for how I can be specified with a wildcard
asterisk to effectively specify multiple values.

\sphinxAtStartPar
If a value begins with “f\_”, a fix ID must follow which has been
previously defined in the input script.  If no bracketed integer is
appended, the per\sphinxhyphen{}atom vector calculated by the fix is used.  If a
bracketed integer is appended, the Ith column of the per\sphinxhyphen{}atom array
calculated by the fix is used.  Note that some fixes only produce
their values on certain timesteps, which must be compatible with
\(N_\text{every}\), else an error results.  Users can also write code for their
own fix styles and {\hyperref[\detokenize{Modify::doc}]{\sphinxcrossref{\DUrole{doc}{add them to LAMMPS}}}}.  See the
discussion above for how I can be specified with a wildcard asterisk
to effectively specify multiple values.

\sphinxAtStartPar
If a value begins with “v\_”, a variable name must follow which has
been previously defined in the input script.  Variables of style
\sphinxstyleemphasis{atom} can reference thermodynamic keywords and various per\sphinxhyphen{}atom
attributes, or invoke other computes, fixes, or variables when they
are evaluated, so this is a very general means of generating per\sphinxhyphen{}atom
quantities to average within grid cells.


\bigskip\hrule\bigskip


\sphinxAtStartPar
\sphinxstyleemphasis{Per\sphinxhyphen{}grid mode}:

\sphinxAtStartPar
The attributes that begin with \sphinxstyleemphasis{c\_ID} and \sphinxstyleemphasis{f\_ID} both take
colon\sphinxhyphen{}separated fields \sphinxstyleemphasis{gname} and \sphinxstyleemphasis{dname}.  These refer to a grid
name and data field name which is defined by the compute or fix.  Note
that a compute or fix can define one or more grids (of different
sizes) and one or more data fields for each of those grids.  The sizes
of all grids used as values for one instance of this fix must be the
same.

\sphinxAtStartPar
The \sphinxstyleemphasis{c\_ID:gname:dname} and \sphinxstyleemphasis{c\_ID:gname:dname{[}I{]}} attributes allow
per\sphinxhyphen{}grid vectors or arrays calculated by a {\hyperref[\detokenize{compute::doc}]{\sphinxcrossref{\DUrole{doc}{compute}}}} to
be accessed.  The ID in the attribute should be replaced by the actual
ID of the compute that has been defined previously in the input
script.

\sphinxAtStartPar
If \sphinxstyleemphasis{c\_ID:gname:dname} is used as a attribute, then the per\sphinxhyphen{}grid vector
calculated by the compute is accessed.  If \sphinxstyleemphasis{c\_ID:gname:dname{[}I{]}} is
used, then I must be in the range from 1\sphinxhyphen{}M, which will access the Ith
column of the per\sphinxhyphen{}grid array with M columns calculated by the compute.
See the discussion above for how I can be specified with a wildcard
asterisk to effectively specify multiple values.

\sphinxAtStartPar
The \sphinxstyleemphasis{f\_ID:gname:dname} and \sphinxstyleemphasis{f\_ID:gname:dname{[}I{]}} attributes allow
per\sphinxhyphen{}grid vectors or arrays calculated by a {\hyperref[\detokenize{fix::doc}]{\sphinxcrossref{\DUrole{doc}{fix}}}} to be
output.  The ID in the attribute should be replaced by the actual ID
of the fix that has been defined previously in the input script.

\sphinxAtStartPar
If \sphinxstyleemphasis{f\_ID:gname:dname} is used as a attribute, then the per\sphinxhyphen{}grid vector
calculated by the fix is printed.  If \sphinxstyleemphasis{f\_ID:gname:dname{[}I{]}} is used,
then I must be in the range from 1\sphinxhyphen{}M, which will print the Ith column
of the per\sphinxhyphen{}grid with M columns calculated by the fix.  See the
discussion above for how I can be specified with a wildcard asterisk
to effectively specify multiple values.


\bigskip\hrule\bigskip


\sphinxAtStartPar
Additional optional keywords also affect the operation of this fix and
its outputs.  Some are only applicable to per\sphinxhyphen{}atom mode.  Some are
applicable to both per\sphinxhyphen{}atom and per\sphinxhyphen{}grid mode.

\sphinxAtStartPar
The \sphinxstyleemphasis{discard} keyword is only applicable to per\sphinxhyphen{}atom mode.  If a
dimension of the system is non\sphinxhyphen{}periodic, then grid cells will only
span the box dimension (fixed or shrink\sphinxhyphen{}wrap boundaries as set by the
{\hyperref[\detokenize{boundary::doc}]{\sphinxcrossref{\DUrole{doc}{boundary command}}}} command).  An atom may thus be slightly outside the
range of grid cells on a particular timestep.  If \sphinxstyleemphasis{discard} is set to
\sphinxstyleemphasis{yes} (the default), then the atom will be assigned to the closest
grid cell (lowest or highest) in that dimension.  If \sphinxstyleemphasis{discard} is set
to \sphinxstyleemphasis{no} the atom will be ignored.


\bigskip\hrule\bigskip


\sphinxAtStartPar
The \sphinxstyleemphasis{norm} keyword is only applicable to per\sphinxhyphen{}atom mode.  In per\sphinxhyphen{}grid
mode, the \sphinxstyleemphasis{norm} keyword setting is ignored.  The output grid value on
an \(N_\text{freq}\) timestep is the sum of the grid values in each of the
\(N_\text{repeat}\) samples, divided by \(N_\text{repeat}\).

\sphinxAtStartPar
In per\sphinxhyphen{}atom mode, the \sphinxstyleemphasis{norm} keyword affects how averaging is done for
the per\sphinxhyphen{}grid values that are output on an \(N_\text{freq}\) timestep. \(N_\text{repeat}\)
samples contribute to the output.  The \sphinxstyleemphasis{norm} keyword has 3 possible
settings: \sphinxstyleemphasis{all} or \sphinxstyleemphasis{sample} or \sphinxstyleemphasis{none}.  \sphinxstyleemphasis{All} is the default.

\sphinxAtStartPar
In the formulas that follow, SumI is the sum of a per\sphinxhyphen{}atom property
over the CountI atoms in a grid cell for a single sample I, where I
varies from 1 to N, and N = Nrepeat.  These formulas are used for any
per\sphinxhyphen{}atom input value listed above, except \sphinxstyleemphasis{density/number},
\sphinxstyleemphasis{density/mass}, and \sphinxstyleemphasis{temp}.  Those input values are discussed below.

\sphinxAtStartPar
In per\sphinxhyphen{}atom mode, for \sphinxstyleemphasis{norm all} the output grid value on the \(N_\text{freq}\)
timestep is an average over atoms across the entire \(N_\text{freq}\) timescale:

\sphinxAtStartPar
Output = (Sum1 + Sum2 + … + SumN) / (Count1 + Count2 + … + CountN)

\sphinxAtStartPar
In per\sphinxhyphen{}atom mode, for \sphinxstyleemphasis{norm sample} the output grid value on the
\(N_\text{freq}\) timestep is an average of an average:

\sphinxAtStartPar
Output = (Sum1/Count1 + Sum2/Count2 + … + SumN/CountN) / Nrepeat

\sphinxAtStartPar
In per\sphinxhyphen{}atom mode, for \sphinxstyleemphasis{norm none} the output grid value on the
\(N_\text{freq}\) timestep is not normalized by the atom counts:

\sphinxAtStartPar
Output = (Sum1 + Sum2 + … SumN) / Nrepeat

\sphinxAtStartPar
For \sphinxstyleemphasis{density/number} and \sphinxstyleemphasis{density/mass}, the output value is the same
as in the formulas above for \sphinxstyleemphasis{norm all} and \sphinxstyleemphasis{norm sample}, except that
the result is also divided by the grid cell volume.  For \sphinxstyleemphasis{norm all},
this will be the volume at the final \(N_\text{freq}\) timestep.  For \sphinxstyleemphasis{norm
sample}, the divide\sphinxhyphen{}by\sphinxhyphen{}volume is done for each sample, using the grid
cell volume at the sample timestep.  For \sphinxstyleemphasis{norm none}, the output is
the same as for \sphinxstyleemphasis{norm all}.

\sphinxAtStartPar
For \sphinxstyleemphasis{temp}, the output temperature uses the formula for kinetic energy
KE listed above, and is normalized similarly to the formulas above for
\sphinxstyleemphasis{norm all} and \sphinxstyleemphasis{norm sample}, except for the way the degrees of
freedom (DOF) are calculated.  For \sphinxstyleemphasis{norm none}, the output is the same
as for \sphinxstyleemphasis{norm all}.

\sphinxAtStartPar
For \sphinxstyleemphasis{norm all}, the DOF = \(N_\text{repeat} \times\) \sphinxstyleemphasis{cdof} plus \sphinxstyleemphasis{Count} times
\sphinxstyleemphasis{adof}, where \sphinxstyleemphasis{Count} = (Count1 + Count2 + … + CountN).  The \sphinxstyleemphasis{cdof}
and \sphinxstyleemphasis{adof} keywords are discussed below.  The output temperature is
computed with all atoms across all samples contributing.

\sphinxAtStartPar
For \sphinxstyleemphasis{norm sample}, the DOF for a single sample = \sphinxstyleemphasis{cdof} plus \sphinxstyleemphasis{Count}
times \sphinxstyleemphasis{adof}, where \sphinxstyleemphasis{Count} = CountI for a single sample.  The output
temperature is the average of \sphinxstyleemphasis{Nsample} temperatures calculated for
each sample.

\sphinxAtStartPar
Finally, for all 3 \sphinxstyleemphasis{norm} settings the output count of atoms per grid
cell is:

\sphinxAtStartPar
Output count = (Count1 + Count2 + … CountN) / Nrepeat

\sphinxAtStartPar
This count is the same for all per\sphinxhyphen{}atom input values, including
\sphinxstyleemphasis{density/number}, \sphinxstyleemphasis{density/mass}, and \sphinxstyleemphasis{temp}.


\bigskip\hrule\bigskip


\sphinxAtStartPar
The \sphinxstyleemphasis{ave} keyword is applied to both per\sphinxhyphen{}atom and per\sphinxhyphen{}grid mode.  It
determines how the per\sphinxhyphen{}grid values produced once every \(N_\text{freq}\) steps
are averaged with values produced on previous steps that were
multiples of \(N_\text{freq}\), before they are accessed by another output
command.

\sphinxAtStartPar
If the \sphinxstyleemphasis{ave} setting is \sphinxstyleemphasis{one}, which is the default, then the grid
values produced on \(N_\text{freq}\) timesteps are independent of each other;
they are output as\sphinxhyphen{}is without further averaging.

\sphinxAtStartPar
If the \sphinxstyleemphasis{ave} setting is \sphinxstyleemphasis{running}, then the grid values produced on
\(N_\text{freq}\) timesteps are summed and averaged in a cumulative sense before
being output.  Each output grid value is thus the average of the grid
value produced on that timestep with all preceding values for the same
grid value.  This running average begins when the fix is defined; it
can only be restarted by deleting the fix via the {\hyperref[\detokenize{unfix::doc}]{\sphinxcrossref{\DUrole{doc}{unfix}}}}
command, or re\sphinxhyphen{}defining the fix by re\sphinxhyphen{}specifying it.

\sphinxAtStartPar
If the \sphinxstyleemphasis{ave} setting is \sphinxstyleemphasis{window}, then the grid values produced on
\(N_\text{freq}\) timesteps are summed and averaged within a moving “window” of
time, so that the last M values for the same grid are used to produce
the output.  E.g. if M = 3 and Nfreq = 1000, then the grid value
output on step 10000 will be the average of the grid values on steps
8000,9000,10000.  Outputs on early steps will average over less than M
values if they are not available.


\bigskip\hrule\bigskip


\sphinxAtStartPar
The \sphinxstyleemphasis{bias}, \sphinxstyleemphasis{adof}, and \sphinxstyleemphasis{cdof} keywords are only applicable to
per\sphinxhyphen{}atom mode.

\sphinxAtStartPar
The \sphinxstyleemphasis{bias} keyword specifies the ID of a temperature compute that
removes a “bias” velocity from each atom, specified as \sphinxstyleemphasis{bias\sphinxhyphen{}ID}.
It is only used when the \sphinxstyleemphasis{temp} value is calculated, to compute the
thermal temperature of each grid cell after the translational kinetic
energy components have been altered in a prescribed way, e.g.  to
remove a flow velocity profile.  See the doc pages for individual
computes that calculate a temperature to see which ones implement a
bias.

\sphinxAtStartPar
The \sphinxstyleemphasis{adof} and \sphinxstyleemphasis{cdof} keywords define the values used in the degree of
freedom (DOF) formula described above for temperature calculation for
each grid cell.  They are only used when the \sphinxstyleemphasis{temp} value is
calculated.  They can be used to calculate a more appropriate
temperature in some cases.  Here are 3 examples:

\sphinxAtStartPar
If grid cells contain some number of water molecules and {\hyperref[\detokenize{fix_shake::doc}]{\sphinxcrossref{\DUrole{doc}{fix
shake}}}} is used to make each molecule rigid, then you could
calculate a temperature with 6 degrees of freedom (DOF) (3
translational, 3 rotational) per molecule by setting \sphinxstyleemphasis{adof} to 2.0.

\sphinxAtStartPar
If {\hyperref[\detokenize{compute_temp_partial::doc}]{\sphinxcrossref{\DUrole{doc}{compute temp/partial}}}} is used with the
\sphinxstyleemphasis{bias} keyword to only allow the x component of velocity to contribute
to the temperature, then \sphinxstyleemphasis{adof} = 1.0 would be appropriate.

\sphinxAtStartPar
Using \sphinxstyleemphasis{cdof} = \sphinxhyphen{}2 or \sphinxhyphen{}3 (for 2d or 3d simulations) will subtract out 2
or 3 degrees of freedom for each grid cell, similar to how the
{\hyperref[\detokenize{compute_temp::doc}]{\sphinxcrossref{\DUrole{doc}{compute temp}}}} command subtracts out 3 DOF for the
entire system.


\bigskip\hrule\bigskip



\subsection{Restart, fix\_modify, output, run start/stop, minimize info}
\label{\detokenize{fix_ave_grid:restart-fix-modify-output-run-start-stop-minimize-info}}
\sphinxAtStartPar
No information about this fix is written to {\hyperref[\detokenize{restart::doc}]{\sphinxcrossref{\DUrole{doc}{binary restart files}}}}.  None of the {\hyperref[\detokenize{fix_modify::doc}]{\sphinxcrossref{\DUrole{doc}{fix\_modify}}}} options are
relevant to this fix.

\sphinxAtStartPar
This fix calculates a per\sphinxhyphen{}grid array which has one column for each of
the specified input values.  The units for each column with be in the
{\hyperref[\detokenize{units::doc}]{\sphinxcrossref{\DUrole{doc}{units}}}} for the per\sphinxhyphen{}atom or per\sphinxhyphen{}grid quantity for the
corresponding input value.  If the fix is used in per\sphinxhyphen{}atom mode, it
also calculates a per\sphinxhyphen{}grid vector with the count of atoms in each grid
cell.  The number of rows in the per\sphinxhyphen{}grid array and number of values
in the per\sphinxhyphen{}grid vector (distributed across all processors) is Nx *
Ny * Nz.

\sphinxAtStartPar
For access by other commands, the name of the single grid produced by
this fix is “grid”.  The names of its two per\sphinxhyphen{}grid datums are “data”
for the per\sphinxhyphen{}grid array and “count” for the per\sphinxhyphen{}grid vector (if using
per\sphinxhyphen{}atom values).  Both datums can be accessed by various {\hyperref[\detokenize{Howto_output::doc}]{\sphinxcrossref{\DUrole{doc}{output
commands}}}}.

\sphinxAtStartPar
In per\sphinxhyphen{}atom mode, the per\sphinxhyphen{}grid array values calculated by this fix are
treated as “intensive”, since they are typically already normalized by
the count of atoms in each grid cell.

\sphinxAtStartPar
No parameter of this fix can be used with the \sphinxstyleemphasis{start/stop} keywords of
the {\hyperref[\detokenize{run::doc}]{\sphinxcrossref{\DUrole{doc}{run}}}} command.  This fix is not invoked during
{\hyperref[\detokenize{minimize::doc}]{\sphinxcrossref{\DUrole{doc}{energy minimization}}}}.


\subsection{Restrictions}
\label{\detokenize{fix_ave_grid:restrictions}}
\sphinxAtStartPar
none


\subsection{Related commands}
\label{\detokenize{fix_ave_grid:related-commands}}
\sphinxAtStartPar
{\hyperref[\detokenize{fix_ave_atom::doc}]{\sphinxcrossref{\DUrole{doc}{fix ave/atom}}}}, {\hyperref[\detokenize{fix_ave_chunk::doc}]{\sphinxcrossref{\DUrole{doc}{fix ave/chunk}}}}


\subsection{Default}
\label{\detokenize{fix_ave_grid:default}}
\sphinxAtStartPar
The option defaults are discard = yes, norm = all, ave = one, and bias
= none.

\sphinxstepscope

\index{fix ave/histo@\spxentry{fix ave/histo}}\ignorespaces 
\index{fix ave/histo/weight@\spxentry{fix ave/histo/weight}}\ignorespaces 

\section{fix ave/histo command}
\label{\detokenize{fix_ave_histo:fix-ave-histo-command}}\label{\detokenize{fix_ave_histo:index-1}}\label{\detokenize{fix_ave_histo:index-0}}\label{\detokenize{fix_ave_histo::doc}}

\section{fix ave/histo/weight command}
\label{\detokenize{fix_ave_histo:fix-ave-histo-weight-command}}

\subsection{Syntax}
\label{\detokenize{fix_ave_histo:syntax}}
\begin{sphinxVerbatim}[commandchars=\\\{\}]
\PYG{k}{fix }\PYG{n+nv+nvIdentifier}{ID}\PYG{+w}{ }\PYG{n+nv+nvIdentifier}{group\PYGZhy{}ID}\PYG{+w}{ }\PYG{n}{style}\PYG{+w}{ }\PYG{n}{Nevery}\PYG{+w}{ }\PYG{n}{Nrepeat}\PYG{+w}{ }\PYG{n}{Nfreq}\PYG{+w}{ }\PYG{n}{lo}\PYG{+w}{ }\PYG{n}{hi}\PYG{+w}{ }\PYG{n}{Nbin}\PYG{+w}{ }\PYG{n}{value1}\PYG{+w}{ }\PYG{n}{value2}\PYG{+w}{ }\PYG{n}{...}\PYG{+w}{ }\PYG{n}{keyword}\PYG{+w}{ }\PYG{n}{args}\PYG{+w}{ }\PYG{n}{...}
\end{sphinxVerbatim}
\begin{itemize}
\item {} 
\sphinxAtStartPar
ID, group\sphinxhyphen{}ID are documented in {\hyperref[\detokenize{fix::doc}]{\sphinxcrossref{\DUrole{doc}{fix}}}} command

\item {} 
\sphinxAtStartPar
style = \sphinxstyleemphasis{ave/histo} or \sphinxstyleemphasis{ave/histo/weight} = style name of this fix command

\item {} 
\sphinxAtStartPar
Nevery = use input values every this many timesteps

\item {} 
\sphinxAtStartPar
Nrepeat = \# of times to use input values for calculating histogram

\item {} 
\sphinxAtStartPar
Nfreq = calculate histogram every this many timesteps

\item {} 
\sphinxAtStartPar
lo,hi = lo/hi bounds within which to histogram

\item {} 
\sphinxAtStartPar
Nbin = \# of histogram bins

\item {} 
\sphinxAtStartPar
one or more input values can be listed

\item {} 
\sphinxAtStartPar
value = \sphinxstyleemphasis{x}, \sphinxstyleemphasis{y}, \sphinxstyleemphasis{z}, \sphinxstyleemphasis{vx}, \sphinxstyleemphasis{vy}, \sphinxstyleemphasis{vz}, \sphinxstyleemphasis{fx}, \sphinxstyleemphasis{fy}, \sphinxstyleemphasis{fz}, c\_ID, c\_ID{[}N{]}, f\_ID, f\_ID{[}N{]}, v\_name

\begin{sphinxVerbatim}[commandchars=\\\{\}]
x,y,z,vx,vy,vz,fx,fy,fz = atom attribute (position, velocity, force component)
c\PYGZus{}ID = scalar or vector calculated by a compute with ID
c\PYGZus{}ID[I] = Ith component of vector or Ith column of array calculated by a compute with ID, I can include wildcard (see below)
f\PYGZus{}ID = scalar or vector calculated by a fix with ID
f\PYGZus{}ID[I] = Ith component of vector or Ith column of array calculated by a fix with ID, I can include wildcard (see below)
v\PYGZus{}name = value(s) calculated by an equal\PYGZhy{}style or vector\PYGZhy{}style or atom\PYGZhy{}style variable with name
v\PYGZus{}name[I] = value calculated by a vector\PYGZhy{}style variable with name, I can include wildcard (see below)
\end{sphinxVerbatim}

\item {} 
\sphinxAtStartPar
zero or more keyword/arg pairs may be appended

\item {} 
\sphinxAtStartPar
keyword = \sphinxstyleemphasis{mode} or \sphinxstyleemphasis{kind} or \sphinxstyleemphasis{file} or \sphinxstyleemphasis{append} or \sphinxstyleemphasis{ave} or \sphinxstyleemphasis{start} or \sphinxstyleemphasis{beyond} or \sphinxstyleemphasis{overwrite} or \sphinxstyleemphasis{title1} or \sphinxstyleemphasis{title2} or \sphinxstyleemphasis{title3}
\begin{sphinxalltt}
\sphinxstyleemphasis{mode} arg = \sphinxstyleemphasis{scalar} or \sphinxstyleemphasis{vector}
  scalar = all input values are scalars
  vector = all input values are vectors
\sphinxstyleemphasis{kind} arg = \sphinxstyleemphasis{global} or \sphinxstyleemphasis{peratom} or \sphinxstyleemphasis{local}
\sphinxstyleemphasis{file} arg = filename
  filename = name of file to output histogram(s) to
\sphinxstyleemphasis{append} arg = filename
  filename = name of file to append histogram(s) to
\sphinxstyleemphasis{ave} args = \sphinxstyleemphasis{one} or \sphinxstyleemphasis{running} or \sphinxstyleemphasis{window}
  one = output a new average value every Nfreq steps
  running = output cumulative average of all previous Nfreq steps
  window M = output average of M most recent Nfreq steps
\sphinxstyleemphasis{start} args = Nstart
  Nstart = start averaging on this timestep
\sphinxstyleemphasis{beyond} arg = \sphinxstyleemphasis{ignore} or \sphinxstyleemphasis{end} or \sphinxstyleemphasis{extra}
  ignore = ignore values outside histogram lo/hi bounds
  end = count values outside histogram lo/hi bounds in end bins
  extra = create 2 extra bins for value outside histogram lo/hi bounds
\sphinxstyleemphasis{overwrite} arg = none = overwrite output file with only latest output
\sphinxstyleemphasis{title1} arg = string
  string = text to print as 1st line of output file
\sphinxstyleemphasis{title2} arg = string
  string = text to print as 2nd line of output file
\sphinxstyleemphasis{title3} arg = string
  string = text to print as 3rd line of output file, only for vector mode
\end{sphinxalltt}

\end{itemize}


\subsection{Examples}
\label{\detokenize{fix_ave_histo:examples}}
\begin{sphinxVerbatim}[commandchars=\\\{\}]
\PYG{k}{fix }\PYG{n+nv+nvIdentifier}{1}\PYG{+w}{ }\PYG{n+nv+nvIdentifier}{all}\PYG{+w}{ }\PYG{n}{ave}\PYG{o}{/}\PYG{n}{histo}\PYG{+w}{ }\PYG{l+m}{100}\PYG{+w}{ }\PYG{l+m}{5}\PYG{+w}{ }\PYG{l+m}{1000}\PYG{+w}{ }\PYG{l+m}{0.5}\PYG{+w}{ }\PYG{l+m}{1.5}\PYG{+w}{ }\PYG{l+m}{50}\PYG{+w}{ }\PYG{n}{c\PYGZus{}myTemp}\PYG{+w}{ }\PYG{n}{file}\PYG{+w}{ }\PYG{n}{temp.histo}\PYG{+w}{ }\PYG{n}{ave}\PYG{+w}{ }\PYG{n}{running}
\PYG{k}{fix }\PYG{n+nv+nvIdentifier}{1}\PYG{+w}{ }\PYG{n+nv+nvIdentifier}{all}\PYG{+w}{ }\PYG{n}{ave}\PYG{o}{/}\PYG{n}{histo}\PYG{+w}{ }\PYG{l+m}{100}\PYG{+w}{ }\PYG{l+m}{5}\PYG{+w}{ }\PYG{l+m}{1000}\PYG{+w}{ }\PYG{o}{\PYGZhy{}}\PYG{l+m}{5}\PYG{+w}{ }\PYG{l+m}{5}\PYG{+w}{ }\PYG{l+m}{100}\PYG{+w}{ }\PYG{n}{c\PYGZus{}thermo\PYGZus{}press[2]}\PYG{+w}{ }\PYG{n}{c\PYGZus{}thermo\PYGZus{}press[3]}\PYG{+w}{ }\PYG{n}{title1}\PYG{+w}{ }\PYG{l+s}{\PYGZdq{}}\PYG{l+s}{My output values}\PYG{l+s}{\PYGZdq{}}
\PYG{k}{fix }\PYG{n+nv+nvIdentifier}{1}\PYG{+w}{ }\PYG{n+nv+nvIdentifier}{all}\PYG{+w}{ }\PYG{n}{ave}\PYG{o}{/}\PYG{n}{histo}\PYG{+w}{ }\PYG{l+m}{100}\PYG{+w}{ }\PYG{l+m}{5}\PYG{+w}{ }\PYG{l+m}{1000}\PYG{+w}{ }\PYG{o}{\PYGZhy{}}\PYG{l+m}{5}\PYG{+w}{ }\PYG{l+m}{5}\PYG{+w}{ }\PYG{l+m}{100}\PYG{+w}{ }\PYG{n}{c\PYGZus{}thermo\PYGZus{}press[}\PYG{o}{*}\PYG{n}{]}
\PYG{k}{fix }\PYG{n+nv+nvIdentifier}{1}\PYG{+w}{ }\PYG{n+nv+nvIdentifier}{all}\PYG{+w}{ }\PYG{n}{ave}\PYG{o}{/}\PYG{n}{histo}\PYG{+w}{ }\PYG{l+m}{1}\PYG{+w}{ }\PYG{l+m}{100}\PYG{+w}{ }\PYG{l+m}{1000}\PYG{+w}{ }\PYG{o}{\PYGZhy{}}\PYG{l+m}{2.0}\PYG{+w}{ }\PYG{l+m}{2.0}\PYG{+w}{ }\PYG{l+m}{18}\PYG{+w}{ }\PYG{n}{vx}\PYG{+w}{ }\PYG{n}{vy}\PYG{+w}{ }\PYG{n}{vz}\PYG{+w}{ }\PYG{n}{mode}\PYG{+w}{ }\PYG{n}{vector}\PYG{+w}{ }\PYG{n}{ave}\PYG{+w}{ }\PYG{n}{running}\PYG{+w}{ }\PYG{n}{beyond}\PYG{+w}{ }\PYG{n}{extra}
\PYG{k}{fix }\PYG{n+nv+nvIdentifier}{1}\PYG{+w}{ }\PYG{n+nv+nvIdentifier}{all}\PYG{+w}{ }\PYG{n}{ave}\PYG{o}{/}\PYG{n}{histo}\PYG{o}{/}\PYG{n}{weight}\PYG{+w}{ }\PYG{l+m}{1}\PYG{+w}{ }\PYG{l+m}{1}\PYG{+w}{ }\PYG{l+m}{1}\PYG{+w}{ }\PYG{l+m}{10}\PYG{+w}{ }\PYG{l+m}{100}\PYG{+w}{ }\PYG{l+m}{2000}\PYG{+w}{ }\PYG{n}{c\PYGZus{}XRD[1]}\PYG{+w}{ }\PYG{n}{c\PYGZus{}XRD[2]}
\end{sphinxVerbatim}


\subsection{Description}
\label{\detokenize{fix_ave_histo:description}}
\sphinxAtStartPar
Use one or more values as inputs every few timesteps to create a
single histogram.  The histogram can then be averaged over longer
timescales.  The resulting histogram can be used by other {\hyperref[\detokenize{Howto_output::doc}]{\sphinxcrossref{\DUrole{doc}{output
commands}}}}, and can also be written to a file.  The fix
ave/histo/weight command has identical syntax to fix ave/histo, except
that exactly two values must be specified.  See details below.

\sphinxAtStartPar
The group specified with this command is ignored for global and local
input values.  For per\sphinxhyphen{}atom input values, only atoms in the group
contribute to the histogram.  Note that regardless of the specified
group, specified values may represent calculations performed by
computes and fixes which store their own “group” definition.

\sphinxAtStartPar
A histogram is simply a count of the number of values that fall within
a histogram bin.  \sphinxstyleemphasis{Nbins} are defined, with even spacing between \sphinxstyleemphasis{lo}
and \sphinxstyleemphasis{hi}.  Values that fall outside the lo/hi bounds can be treated in
different ways; see the discussion of the \sphinxstyleemphasis{beyond} keyword below.

\sphinxAtStartPar
Each input value can be an atom attribute (position, velocity, force
component) or can be the result of a {\hyperref[\detokenize{compute::doc}]{\sphinxcrossref{\DUrole{doc}{compute}}}} or
{\hyperref[\detokenize{fix::doc}]{\sphinxcrossref{\DUrole{doc}{fix}}}} or the evaluation of an equal\sphinxhyphen{}style or vector\sphinxhyphen{}style
or atom\sphinxhyphen{}style {\hyperref[\detokenize{variable::doc}]{\sphinxcrossref{\DUrole{doc}{variable}}}}.  The set of input values can
be either all global, all per\sphinxhyphen{}atom, or all local quantities.  Inputs
of different kinds (e.g. global and per\sphinxhyphen{}atom) cannot be mixed.  Atom
attributes are per\sphinxhyphen{}atom vector values.  See the page for individual
“compute” and “fix” commands to see what kinds of quantities they
generate.

\sphinxAtStartPar
Note that a compute or fix can produce multiple kinds of data (global,
per\sphinxhyphen{}atom, local).  If LAMMPS cannot unambiguously determine which kind
of data to use, the optional \sphinxstyleemphasis{kind} keyword discussed below can force
the desired disambiguation.

\sphinxAtStartPar
Note that the output of this command is a single histogram for all
input values combined together, not one histogram per input value.
See below for details on the format of the output of this fix.

\sphinxAtStartPar
The input values must either be all scalars or all vectors (or
arrays), depending on the setting of the \sphinxstyleemphasis{mode} keyword.

\sphinxAtStartPar
If \sphinxstyleemphasis{mode} = scalar, then the input values must be scalars, or vectors
with a bracketed term appended, indicating the Ith value of the vector
is used.

\sphinxAtStartPar
If \sphinxstyleemphasis{mode} = vector, then the input values must be vectors, or arrays
with a bracketed term appended, indicating the Ith column of the array
is used.

\sphinxAtStartPar
If the fix ave/histo/weight command is used, exactly two values must
be specified.  If the values are vectors, they must be the same
length.  The first value (a scalar or vector) is what is histogrammed
into bins, in the same manner the fix ave/histo command operates.  The
second value (a scalar or vector) is used as a “weight”.  This means
that instead of each value tallying a “1” to its bin, the
corresponding weight is tallied.  For example, the \(N^\text{th}\) entry
(weight) in the second vector is tallied to the bin corresponding to the
\(N^\text{th}\) entry in the first vector.


\bigskip\hrule\bigskip


\sphinxAtStartPar
For input values from a compute or fix or variable, the bracketed
index I can be specified using a wildcard asterisk with the index to
effectively specify multiple values.  This takes the form “*” or
“*n” or “m*” or “m*n”.  If \(N\) is the size of the vector
(for \sphinxstyleemphasis{mode} = scalar) or the number of columns in the array
(for \sphinxstyleemphasis{mode} = vector), then an asterisk with no numeric values means all
indices from 1 to \(N\).
A leading asterisk means all indices from 1 to n (inclusive).  A
trailing asterisk means all indices from m to \(N\) (inclusive).  A middle
asterisk means all indices from m to n (inclusive).

\sphinxAtStartPar
Using a wildcard is the same as if the individual elements of the
vector or columns of the array had been listed one by one.  For example, the
following two fix ave/histo commands are equivalent, since the {\hyperref[\detokenize{compute_com_chunk::doc}]{\sphinxcrossref{\DUrole{doc}{compute
com/chunk}}}} command creates a global array with three
columns:

\begin{sphinxVerbatim}[commandchars=\\\{\}]
\PYG{k}{compute }\PYG{n+nv+nvIdentifier}{myCOM}\PYG{+w}{ }\PYG{n+nv+nvIdentifier}{all}\PYG{+w}{ }\PYG{n}{com}\PYG{o}{/}\PYG{n}{chunk}
\PYG{k}{fix }\PYG{n+nv+nvIdentifier}{1}\PYG{+w}{ }\PYG{n+nv+nvIdentifier}{all}\PYG{+w}{ }\PYG{n}{ave}\PYG{o}{/}\PYG{n}{histo}\PYG{+w}{ }\PYG{l+m}{100}\PYG{+w}{ }\PYG{l+m}{1}\PYG{+w}{ }\PYG{l+m}{100}\PYG{+w}{ }\PYG{o}{\PYGZhy{}}\PYG{l+m}{10.0}\PYG{+w}{ }\PYG{l+m}{10.0}\PYG{+w}{ }\PYG{l+m}{100}\PYG{+w}{ }\PYG{n}{c\PYGZus{}myCOM[}\PYG{o}{*}\PYG{n}{]}\PYG{+w}{ }\PYG{n}{file}\PYG{+w}{ }\PYG{n}{tmp1.com}\PYG{+w}{ }\PYG{n}{mode}\PYG{+w}{ }\PYG{n}{vector}
\PYG{k}{fix }\PYG{n+nv+nvIdentifier}{2}\PYG{+w}{ }\PYG{n+nv+nvIdentifier}{all}\PYG{+w}{ }\PYG{n}{ave}\PYG{o}{/}\PYG{n}{histo}\PYG{+w}{ }\PYG{l+m}{100}\PYG{+w}{ }\PYG{l+m}{1}\PYG{+w}{ }\PYG{l+m}{100}\PYG{+w}{ }\PYG{o}{\PYGZhy{}}\PYG{l+m}{10.0}\PYG{+w}{ }\PYG{l+m}{10.0}\PYG{+w}{ }\PYG{l+m}{100}\PYG{+w}{ }\PYG{n}{c\PYGZus{}myCOM[1]}\PYG{+w}{ }\PYG{n}{c\PYGZus{}myCOM[2]}\PYG{+w}{ }\PYG{n}{c\PYGZus{}myCOM[3]}\PYG{+w}{ }\PYG{n}{file}\PYG{+w}{ }\PYG{n}{tmp2.com}\PYG{+w}{ }\PYG{n}{mode}\PYG{+w}{ }\PYG{n}{vector}
\end{sphinxVerbatim}

\begin{sphinxadmonition}{note}{Note:}
\sphinxAtStartPar
For a vector\sphinxhyphen{}style variable, only the wildcard forms “*n” or
“m*n” are allowed.  You must specify the upper bound, because
vector\sphinxhyphen{}style variable lengths are not determined until the variable
is evaluated.  If n is specified larger than the vector length
turns out to be, zeroes are output for missing vector values.
\end{sphinxadmonition}


\bigskip\hrule\bigskip


\sphinxAtStartPar
The \(N_\text{every}\), \(N_\text{repeat}\), and \(N_\text{freq}\)
arguments specify on what time steps the input values will be used in order to
contribute to the histogram.  The final histogram is generated on time steps
that are multiple of \(N_\text{freq}\).  It is averaged over
\(N_\text{repeat}\) histograms, computed in the preceding portion of the
simulation every \(N_\text{every}\) time steps.
\(N_\text{freq}\) must be a multiple of \(N_\text{every}\) and
\(N_\text{every}\) must be non\sphinxhyphen{}zero even if \(N_\text{repeat}\) is 1.
Also, the time steps contributing to the histogram value cannot overlap
(i.e., \(N_\text{repeat}\times N_\text{every}\) cannot exceed
\(N_\text{freq}\)).

\sphinxAtStartPar
For example, if \(N_\text{every}=2\), \(N_\text{repeat}=6\), and
\(N_\text{freq}=100\), then input values on time steps 90, 92, 94, 96, 98,
and 100 will be used to compute the final histogram on timestep 100.
Similarly for timesteps 190, 192, 194, 196, 198, and 200 on timestep 200, etc.
If \(N_\text{repeat}=1\) and \(N_\text{freq} = 100\), then no time
averaging of the histogram is done; a histogram is simply generated on
timesteps 100, 200, etc.


\bigskip\hrule\bigskip


\sphinxAtStartPar
The atom attribute values (\sphinxstyleemphasis{x}, \sphinxstyleemphasis{y}, \sphinxstyleemphasis{z}, \sphinxstyleemphasis{vx}, \sphinxstyleemphasis{vy}, \sphinxstyleemphasis{vz}, \sphinxstyleemphasis{fx}, \sphinxstyleemphasis{fy}, and
\sphinxstyleemphasis{fz}) are self\sphinxhyphen{}explanatory.  Note that other atom attributes can be used as
inputs to this fix by using the
{\hyperref[\detokenize{compute_property_atom::doc}]{\sphinxcrossref{\DUrole{doc}{compute property/atom}}}} command and then
specifying an input value from that compute.

\sphinxAtStartPar
If a value begins with “c\_”, a compute ID must follow which has been
previously defined in the input script.  If \sphinxstyleemphasis{mode} = scalar, then if
no bracketed term is appended, the global scalar calculated by the
compute is used.  If a bracketed term is appended, the Ith element of
the global vector calculated by the compute is used.  If \sphinxstyleemphasis{mode} =
vector, then if no bracketed term is appended, the global or per\sphinxhyphen{}atom
or local vector calculated by the compute is used.  If a bracketed
term is appended, the Ith column of the global or per\sphinxhyphen{}atom or local
array calculated by the compute is used.  See the discussion above for
how I can be specified with a wildcard asterisk to effectively specify
multiple values.

\sphinxAtStartPar
Note that there is a {\hyperref[\detokenize{compute_reduce::doc}]{\sphinxcrossref{\DUrole{doc}{compute reduce}}}} command
that can sum per\sphinxhyphen{}atom quantities into a global scalar or vector, which
can then be accessed by fix ave/histo.  It can also be a compute defined
not in your input script, but by {\hyperref[\detokenize{thermo_style::doc}]{\sphinxcrossref{\DUrole{doc}{thermodynamic output}}}}
or other fixes such as {\hyperref[\detokenize{fix_nh::doc}]{\sphinxcrossref{\DUrole{doc}{fix nvt}}}}
or {\hyperref[\detokenize{fix_temp_rescale::doc}]{\sphinxcrossref{\DUrole{doc}{fix temp/rescale}}}}.  See the doc pages for
these commands which give the IDs of these computes.  Users can also
write code for their own compute styles and
{\hyperref[\detokenize{Modify::doc}]{\sphinxcrossref{\DUrole{doc}{add them to LAMMPS}}}}.

\sphinxAtStartPar
If a value begins with “f\_”, a fix ID must follow which has been
previously defined in the input script.  If \sphinxstyleemphasis{mode} = scalar, then if
no bracketed term is appended, the global scalar calculated by the fix
is used.  If a bracketed term is appended, the Ith element of the
global vector calculated by the fix is used.  If \sphinxstyleemphasis{mode} = vector, then
if no bracketed term is appended, the global or per\sphinxhyphen{}atom or local
vector calculated by the fix is used.  If a bracketed term is
appended, the \(I^\text{th}\) column of the global or per\sphinxhyphen{}atom or local
array calculated by the fix is used.  See the discussion above for how
\(I\) can be specified with a wildcard asterisk to effectively specify
multiple values.

\sphinxAtStartPar
Note that some fixes only produce their values on certain timesteps,
which must be compatible with \(N_\text{every}\), else an error will
result.  Users can also write code for their own fix styles and
{\hyperref[\detokenize{Modify::doc}]{\sphinxcrossref{\DUrole{doc}{add them to LAMMPS}}}}.

\sphinxAtStartPar
If a value begins with “v\_”, a variable name must follow which has
been previously defined in the input script.  If \sphinxstyleemphasis{mode} = scalar, then
only equal\sphinxhyphen{}style or vector\sphinxhyphen{}style variables can be used, which both
produce global values.  In this mode, a vector\sphinxhyphen{}style variable requires
a bracketed term to specify the \(I^\text{th}\) element of the vector
calculated by the variable.  If \sphinxstyleemphasis{mode} = vector, then only vector\sphinxhyphen{}style or
atom\sphinxhyphen{}style variables can be used, which produce a global or per\sphinxhyphen{}atom
vector respectively.  The vector\sphinxhyphen{}style variable must be used without a
bracketed term.  See the {\hyperref[\detokenize{variable::doc}]{\sphinxcrossref{\DUrole{doc}{variable}}}} command for details.

\sphinxAtStartPar
Note that variables of style \sphinxstyleemphasis{equal}, \sphinxstyleemphasis{vector}, and \sphinxstyleemphasis{atom} define a
formula which can reference individual atom properties or
thermodynamic keywords, or they can invoke other computes, fixes, or
variables when they are evaluated, so this is a very general means of
specifying quantities to histogram.


\bigskip\hrule\bigskip


\sphinxAtStartPar
Additional optional keywords also affect the operation of this fix.

\sphinxAtStartPar
If the \sphinxstyleemphasis{mode} keyword is set to \sphinxstyleemphasis{scalar}, then all input values must
be global scalars, or elements of global vectors.  If the \sphinxstyleemphasis{mode}
keyword is set to \sphinxstyleemphasis{vector}, then all input values must be global or
per\sphinxhyphen{}atom or local vectors, or columns of global or per\sphinxhyphen{}atom or local
arrays.

\sphinxAtStartPar
The \sphinxstyleemphasis{kind} keyword only needs to be used if any of the specified input
computes or fixes produce more than one kind of output (global,
per\sphinxhyphen{}atom, local).  If not, LAMMPS will determine the kind of data all
the inputs produce and verify it is all the same kind.  If not, an
error will be triggered.  If a compute or fix produces more than one
kind of output, the \sphinxstyleemphasis{kind} keyword should be used to specify which
output will be used.  The other input arguments must still be
consistent.

\sphinxAtStartPar
The \sphinxstyleemphasis{beyond} keyword determines how input values that fall outside the
\sphinxstyleemphasis{lo} to \sphinxstyleemphasis{hi} bounds are treated.  Values such that \sphinxstyleemphasis{lo} \(\le\) value
\(\le\) \sphinxstyleemphasis{hi} are assigned to one bin.  Values on a bin boundary are
assigned to the lower of the two bins.  If \sphinxstyleemphasis{beyond} is set to \sphinxstyleemphasis{ignore} then
values \(<\) \sphinxstyleemphasis{lo} and values \(>\) \sphinxstyleemphasis{hi} are ignored (i.e., they are not
binned). If \sphinxstyleemphasis{beyond} is set to \sphinxstyleemphasis{end}, then values \(<\) \sphinxstyleemphasis{lo} are counted in
the first bin and values \(>\) \sphinxstyleemphasis{hi} are counted in the last bin.
If \sphinxstyleemphasis{beyond} is set to \sphinxstyleemphasis{extend}, then two extra bins are created so that there
are \(N_\text{bins}+2\) total bins.  Values \(<\) \sphinxstyleemphasis{lo} are counted in
the first bin and values \(>\) \sphinxstyleemphasis{hi} are counted in the last bin
\((N_\text{bins}+2)\).  Values between
\sphinxstyleemphasis{lo} and \sphinxstyleemphasis{hi} (inclusive) are counted in bins 2 through
\(N_\text{bins}+1\).  The “coordinate” stored and printed for these two
extra bins is \sphinxstyleemphasis{lo} and \sphinxstyleemphasis{hi}.

\sphinxAtStartPar
The \sphinxstyleemphasis{ave} keyword determines how the histogram produced every
\(N_\text{freq}\) steps are averaged with histograms produced on previous
steps that were multiples of \(N_\text{freq}\), before they are accessed by
another output command or written to a file.

\sphinxAtStartPar
If the \sphinxstyleemphasis{ave} setting is \sphinxstyleemphasis{one}, then the histograms produced on
timesteps that are multiples of \(N_\text{freq}\) are independent of each
other; they are output as\sphinxhyphen{}is without further averaging.

\sphinxAtStartPar
If the \sphinxstyleemphasis{ave} setting is \sphinxstyleemphasis{running}, then the histograms produced on
timesteps that are multiples of \(N_\text{freq}\) are summed and averaged
in a cumulative sense before being output.  Each bin value in the histogram
is thus the average of the bin value produced on that timestep with all
preceding values for the same bin.  This running average begins when the fix is
defined; it can only be restarted by deleting the fix via the
{\hyperref[\detokenize{unfix::doc}]{\sphinxcrossref{\DUrole{doc}{unfix}}}} command, or by re\sphinxhyphen{}defining the fix by re\sphinxhyphen{}specifying it.

\sphinxAtStartPar
If the \sphinxstyleemphasis{ave} setting is \sphinxstyleemphasis{window}, then the histograms produced on
timesteps that are multiples of \(N_\text{freq}\) are summed within a
moving “window” of time, so that the last \(M\) histograms are used to
produce the output (e.g., if \(M = 3\) and \(N_\text{freq} = 1000\),
then the output on step 10000 will be the combined histogram of the individual
histograms on steps 8000, 9000, and 10000.  Outputs on early steps will be sums
over less than \(M\) histograms if they are not available.

\sphinxAtStartPar
The \sphinxstyleemphasis{start} keyword specifies what timestep histogramming will begin
on.  The default is step 0.  Often input values can be 0.0 at time 0,
so setting \sphinxstyleemphasis{start} to a larger value can avoid including a 0.0 in
a running or windowed histogram.

\sphinxAtStartPar
\DUrole{versionmodified,added}{Added in version 17Apr2024: }new keyword \sphinxstyleemphasis{append}

\sphinxAtStartPar
The \sphinxstyleemphasis{file} or \sphinxstyleemphasis{append} keywords allow a filename to be specified.  If
\sphinxstyleemphasis{file} is used, then the filename is overwritten if it already exists.
If \sphinxstyleemphasis{append} is used, then the filename is appended to if it already
exists, or created if it does not exist.  Every \sphinxstyleemphasis{Nfreq} steps, one
histogram is written to the file.  This includes a leading line that
contains the timestep, number of bins, the total count of values
contributing to the histogram, the count of values that were not
histogrammed (see the \sphinxstyleemphasis{beyond} keyword), the minimum value encountered,
and the maximum value encountered.  The min/max values include values
that were not histogrammed.  Following the leading line, one line per
bin is written into the file.  Each line contains the bin \#, the
coordinate for the center of the bin (between \sphinxstyleemphasis{lo} and \sphinxstyleemphasis{hi}), the
count of values in the bin, and the normalized count.  The normalized
count is the bin count divided by the total count (not including values
not histogrammed), so that the normalized values sum to 1.0 across all
bins.

\sphinxAtStartPar
The \sphinxstyleemphasis{overwrite} keyword will continuously overwrite the output file
with the latest output, so that it only contains one timestep worth of
output.  This option can only be used with the \sphinxstyleemphasis{ave running} setting.

\sphinxAtStartPar
The \sphinxstyleemphasis{title1}, \sphinxstyleemphasis{title2}, and \sphinxstyleemphasis{title3} keywords allow specification of
the strings that will be printed as the first three lines of the output
file, assuming the \sphinxstyleemphasis{file} keyword was used.  LAMMPS uses default
values for each of these, so they do not need to be specified.

\sphinxAtStartPar
By default, these header lines are as follows:

\begin{sphinxVerbatim}[commandchars=\\\{\}]
\PYGZsh{} Histogram for fix ID
\PYGZsh{} TimeStep Number\PYGZhy{}of\PYGZhy{}bins Total\PYGZhy{}counts Missing\PYGZhy{}counts Min\PYGZhy{}value Max\PYGZhy{}value
\PYGZsh{} Bin Coord Count Count/Total
\end{sphinxVerbatim}

\sphinxAtStartPar
In the first line, ID is replaced with the fix\sphinxhyphen{}ID.  The second line
describes the six values that are printed at the first of each section
of output.  The third describes the four values printed for each bin in
the histogram.


\bigskip\hrule\bigskip



\subsection{Restart, fix\_modify, output, run start/stop, minimize info}
\label{\detokenize{fix_ave_histo:restart-fix-modify-output-run-start-stop-minimize-info}}
\sphinxAtStartPar
No information about this fix is written to
{\hyperref[\detokenize{restart::doc}]{\sphinxcrossref{\DUrole{doc}{binary restart files}}}}.
None of the {\hyperref[\detokenize{fix_modify::doc}]{\sphinxcrossref{\DUrole{doc}{fix\_modify}}}} options are relevant to this fix.

\sphinxAtStartPar
This fix produces a global vector and global array which can be
accessed by various {\hyperref[\detokenize{Howto_output::doc}]{\sphinxcrossref{\DUrole{doc}{output commands}}}}.  The values
can only be accessed on timesteps that are multiples of \(N_\text{freq}\)
since that is when a histogram is generated.  The global vector has four
values:
\begin{enumerate}
\sphinxsetlistlabels{\arabic}{enumi}{enumii}{}{.}%
\item {} 
\sphinxAtStartPar
total counts in the histogram

\item {} 
\sphinxAtStartPar
values that were not histogrammed (see \sphinxstyleemphasis{beyond} keyword)

\item {} 
\sphinxAtStartPar
min value of all input values, including ones not histogrammed

\item {} 
\sphinxAtStartPar
max value of all input values, including ones not histogrammed

\end{enumerate}

\sphinxAtStartPar
The global array has \(N_\text{bins}\) rows and three columns.  The
first column has the bin coordinate, the second column has the count of
values in that histogram bin, and the third column has the bin count
divided by the total count (not including missing counts), so that the
values in the third column sum to 1.0.

\sphinxAtStartPar
The vector and array values calculated by this fix are all treated as
intensive.  If this is not the case (e.g., due to histogramming
per\sphinxhyphen{}atom input values), then you will need to account for that when
interpreting the values produced by this fix.

\sphinxAtStartPar
No parameter of this fix can be used with the \sphinxstyleemphasis{start/stop} keywords of
the {\hyperref[\detokenize{run::doc}]{\sphinxcrossref{\DUrole{doc}{run}}}} command.
This fix is not invoked during {\hyperref[\detokenize{minimize::doc}]{\sphinxcrossref{\DUrole{doc}{energy minimization}}}}.


\subsection{Restrictions}
\label{\detokenize{fix_ave_histo:restrictions}}\begin{quote}

\sphinxAtStartPar
none
\end{quote}


\subsection{Related commands}
\label{\detokenize{fix_ave_histo:related-commands}}
\sphinxAtStartPar
{\hyperref[\detokenize{compute::doc}]{\sphinxcrossref{\DUrole{doc}{compute}}}}, {\hyperref[\detokenize{fix_ave_atom::doc}]{\sphinxcrossref{\DUrole{doc}{fix ave/atom}}}},
{\hyperref[\detokenize{fix_ave_chunk::doc}]{\sphinxcrossref{\DUrole{doc}{fix ave/chunk}}}}, {\hyperref[\detokenize{fix_ave_time::doc}]{\sphinxcrossref{\DUrole{doc}{fix ave/time}}}},
{\hyperref[\detokenize{variable::doc}]{\sphinxcrossref{\DUrole{doc}{variable}}}}, {\hyperref[\detokenize{fix_ave_correlate::doc}]{\sphinxcrossref{\DUrole{doc}{fix ave/correlate}}}},


\subsection{Default}
\label{\detokenize{fix_ave_histo:default}}
\sphinxAtStartPar
none

\sphinxAtStartPar
The option defaults are mode = scalar, kind = figured out from input
arguments, ave = one, start = 0, no file output, beyond = ignore, and
title 1,2,3 = strings as described above.

\sphinxstepscope


\section{fix ave/spatial command}
\label{\detokenize{fix_ave_spatial:fix-ave-spatial-command}}\label{\detokenize{fix_ave_spatial::doc}}
\sphinxAtStartPar
\DUrole{versionmodified,deprecated}{Deprecated since version 11Dec2015.}

\sphinxAtStartPar
The \sphinxtitleref{fix ave/spatial} command has been superseded by {\hyperref[\detokenize{fix_ave_chunk::doc}]{\sphinxcrossref{\DUrole{doc}{fix ave/chunk}}}}.

\sphinxstepscope


\section{fix ave/spatial/sphere command}
\label{\detokenize{fix_ave_spatial_sphere:fix-ave-spatial-sphere-command}}\label{\detokenize{fix_ave_spatial_sphere::doc}}
\sphinxAtStartPar
\DUrole{versionmodified,deprecated}{Deprecated since version 11Dec2015.}

\sphinxAtStartPar
The \sphinxtitleref{fix ave/spatial/sphere} command has been superseded by {\hyperref[\detokenize{fix_ave_chunk::doc}]{\sphinxcrossref{\DUrole{doc}{fix ave/chunk}}}}.

\sphinxstepscope

\index{fix ave/time@\spxentry{fix ave/time}}\ignorespaces 

\section{fix ave/time command}
\label{\detokenize{fix_ave_time:fix-ave-time-command}}\label{\detokenize{fix_ave_time:index-0}}\label{\detokenize{fix_ave_time::doc}}

\subsection{Syntax}
\label{\detokenize{fix_ave_time:syntax}}
\begin{sphinxVerbatim}[commandchars=\\\{\}]
\PYG{k}{fix }\PYG{n+nv+nvIdentifier}{ID}\PYG{+w}{ }\PYG{n+nv+nvIdentifier}{group\PYGZhy{}ID}\PYG{+w}{ }\PYG{n}{ave}\PYG{o}{/}\PYG{n}{time}\PYG{+w}{ }\PYG{n}{Nevery}\PYG{+w}{ }\PYG{n}{Nrepeat}\PYG{+w}{ }\PYG{n}{Nfreq}\PYG{+w}{ }\PYG{n}{value1}\PYG{+w}{ }\PYG{n}{value2}\PYG{+w}{ }\PYG{n}{...}\PYG{+w}{ }\PYG{n}{keyword}\PYG{+w}{ }\PYG{n}{args}\PYG{+w}{ }\PYG{n}{...}
\end{sphinxVerbatim}
\begin{itemize}
\item {} 
\sphinxAtStartPar
ID, group\sphinxhyphen{}ID are documented in {\hyperref[\detokenize{fix::doc}]{\sphinxcrossref{\DUrole{doc}{fix}}}} command

\item {} 
\sphinxAtStartPar
ave/time = style name of this fix command

\item {} 
\sphinxAtStartPar
Nevery = use input values every this many time steps

\item {} 
\sphinxAtStartPar
Nrepeat = \# of times to use input values for calculating averages

\item {} 
\sphinxAtStartPar
Nfreq = calculate averages every this many time steps

\item {} 
\sphinxAtStartPar
one or more input values can be listed

\item {} 
\sphinxAtStartPar
value = c\_ID, c\_ID{[}N{]}, f\_ID, f\_ID{[}N{]}, v\_name

\begin{sphinxVerbatim}[commandchars=\\\{\}]
c\PYGZus{}ID = global scalar or vector calculated by a compute with ID
c\PYGZus{}ID[I] = Ith component of global vector or Ith column of global array calculated by a compute with ID, I can include wildcard (see below)
f\PYGZus{}ID = global scalar or vector calculated by a fix with ID
f\PYGZus{}ID[I] = Ith component of global vector or Ith column of global array calculated by a fix with ID, I can include wildcard (see below)
v\PYGZus{}name = value(s) calculated by an equal\PYGZhy{}style or vector\PYGZhy{}style variable with name
v\PYGZus{}name[I] = value calculated by a vector\PYGZhy{}style variable with name, I can include wildcard (see below)
\end{sphinxVerbatim}

\item {} 
\sphinxAtStartPar
zero or more keyword/arg pairs may be appended

\item {} 
\sphinxAtStartPar
keyword = \sphinxstyleemphasis{mode} or \sphinxstyleemphasis{file} or \sphinxstyleemphasis{append} or \sphinxstyleemphasis{ave} or \sphinxstyleemphasis{start} or \sphinxstyleemphasis{off} or \sphinxstyleemphasis{overwrite} or \sphinxstyleemphasis{format} or \sphinxstyleemphasis{title1} or \sphinxstyleemphasis{title2} or \sphinxstyleemphasis{title3}
\begin{sphinxalltt}
\sphinxstyleemphasis{mode} arg = \sphinxstyleemphasis{scalar} or \sphinxstyleemphasis{vector}
  scalar = all input values are global scalars
  vector = all input values are global vectors or global arrays
\sphinxstyleemphasis{ave} args = \sphinxstyleemphasis{one} or \sphinxstyleemphasis{running} or \sphinxstyleemphasis{window M}
  one = output a new average value every Nfreq steps
  running = output cumulative average of all previous Nfreq steps
  window M = output average of M most recent Nfreq steps
\sphinxstyleemphasis{start} args = Nstart
  Nstart = start averaging on this time step
\sphinxstyleemphasis{off} arg = M = do not average this value
  M = value \# from 1 to Nvalues
\sphinxstyleemphasis{file} arg = filename
  filename = name of file to output time averages to
\sphinxstyleemphasis{append} arg = filename
  filename = name of file to append time averages to
\sphinxstyleemphasis{overwrite} arg = none = overwrite output file with only latest output
\sphinxstyleemphasis{format} arg = string
  string = C\sphinxhyphen{}style format string
\sphinxstyleemphasis{title1} arg = string
  string = text to print as 1st line of output file
\sphinxstyleemphasis{title2} arg = string
  string = text to print as 2nd line of output file
\sphinxstyleemphasis{title3} arg = string
  string = text to print as 3rd line of output file, only for vector mode
\end{sphinxalltt}

\end{itemize}


\subsection{Examples}
\label{\detokenize{fix_ave_time:examples}}
\begin{sphinxVerbatim}[commandchars=\\\{\}]
\PYG{k}{fix }\PYG{n+nv+nvIdentifier}{1}\PYG{+w}{ }\PYG{n+nv+nvIdentifier}{all}\PYG{+w}{ }\PYG{n}{ave}\PYG{o}{/}\PYG{n}{time}\PYG{+w}{ }\PYG{l+m}{100}\PYG{+w}{ }\PYG{l+m}{5}\PYG{+w}{ }\PYG{l+m}{1000}\PYG{+w}{ }\PYG{n}{c\PYGZus{}myTemp}\PYG{+w}{ }\PYG{n}{c\PYGZus{}thermo\PYGZus{}temp}\PYG{+w}{ }\PYG{n}{file}\PYG{+w}{ }\PYG{n}{temp.profile}
\PYG{k}{fix }\PYG{n+nv+nvIdentifier}{1}\PYG{+w}{ }\PYG{n+nv+nvIdentifier}{all}\PYG{+w}{ }\PYG{n}{ave}\PYG{o}{/}\PYG{n}{time}\PYG{+w}{ }\PYG{l+m}{100}\PYG{+w}{ }\PYG{l+m}{5}\PYG{+w}{ }\PYG{l+m}{1000}\PYG{+w}{ }\PYG{n}{c\PYGZus{}thermo\PYGZus{}press[2]}\PYG{+w}{ }\PYG{n}{ave}\PYG{+w}{ }\PYG{n}{window}\PYG{+w}{ }\PYG{l+m}{20}\PYG{+w}{ }\PYG{o}{\PYGZam{}}
\PYG{+w}{                              }\PYG{n}{title1}\PYG{+w}{ }\PYG{l+s}{\PYGZdq{}}\PYG{l+s}{My output values}\PYG{l+s}{\PYGZdq{}}
\PYG{k}{fix }\PYG{n+nv+nvIdentifier}{1}\PYG{+w}{ }\PYG{n+nv+nvIdentifier}{all}\PYG{+w}{ }\PYG{n}{ave}\PYG{o}{/}\PYG{n}{time}\PYG{+w}{ }\PYG{l+m}{100}\PYG{+w}{ }\PYG{l+m}{5}\PYG{+w}{ }\PYG{l+m}{1000}\PYG{+w}{ }\PYG{n}{c\PYGZus{}thermo\PYGZus{}press[}\PYG{o}{*}\PYG{n}{]}
\PYG{k}{fix }\PYG{n+nv+nvIdentifier}{1}\PYG{+w}{ }\PYG{n+nv+nvIdentifier}{all}\PYG{+w}{ }\PYG{n}{ave}\PYG{o}{/}\PYG{n}{time}\PYG{+w}{ }\PYG{l+m}{1}\PYG{+w}{ }\PYG{l+m}{100}\PYG{+w}{ }\PYG{l+m}{1000}\PYG{+w}{ }\PYG{n}{f\PYGZus{}indent}\PYG{+w}{ }\PYG{n}{f\PYGZus{}indent[1]}\PYG{+w}{ }\PYG{n}{file}\PYG{+w}{ }\PYG{n}{temp.indent}\PYG{+w}{ }\PYG{n}{off}\PYG{+w}{ }\PYG{l+m}{1}
\end{sphinxVerbatim}


\subsection{Description}
\label{\detokenize{fix_ave_time:description}}
\sphinxAtStartPar
Use one or more global values as inputs every few time steps, and
average them over longer timescales.  The resulting averages can be
used by other {\hyperref[\detokenize{Howto_output::doc}]{\sphinxcrossref{\DUrole{doc}{output commands}}}} such as
{\hyperref[\detokenize{thermo_style::doc}]{\sphinxcrossref{\DUrole{doc}{thermo\_style custom}}}}, and can also be written to a
file.  Note that if no time averaging is done, this command can be
used as a convenient way to simply output one or more global values to
a file.

\sphinxAtStartPar
The group specified with this command is ignored.  However, note that
specified values may represent calculations performed by computes and
fixes which store their own “group” definitions.

\sphinxAtStartPar
Each listed value can be the result of a {\hyperref[\detokenize{compute::doc}]{\sphinxcrossref{\DUrole{doc}{compute}}}} or
{\hyperref[\detokenize{fix::doc}]{\sphinxcrossref{\DUrole{doc}{fix}}}} or the evaluation of an equal\sphinxhyphen{}style or vector\sphinxhyphen{}style
{\hyperref[\detokenize{variable::doc}]{\sphinxcrossref{\DUrole{doc}{variable}}}}.  In each case, the compute, fix, or variable
must produce a global quantity, not a per\sphinxhyphen{}atom or local quantity.  If
you wish to spatial\sphinxhyphen{} or time\sphinxhyphen{}average or histogram per\sphinxhyphen{}atom quantities
from a compute, fix, or variable, then see the
{\hyperref[\detokenize{fix_ave_chunk::doc}]{\sphinxcrossref{\DUrole{doc}{fix ave/chunk}}}}, {\hyperref[\detokenize{fix_ave_atom::doc}]{\sphinxcrossref{\DUrole{doc}{fix ave/atom}}}},
or {\hyperref[\detokenize{fix_ave_histo::doc}]{\sphinxcrossref{\DUrole{doc}{fix ave/histo}}}} commands.  If you wish to sum a
per\sphinxhyphen{}atom quantity into a single global quantity, see the
{\hyperref[\detokenize{compute_reduce::doc}]{\sphinxcrossref{\DUrole{doc}{compute reduce}}}} command.

\sphinxAtStartPar
{\hyperref[\detokenize{compute::doc}]{\sphinxcrossref{\DUrole{doc}{Computes}}}} that produce global quantities are those which
do not have the word \sphinxstyleemphasis{atom} in their style name.  Only a few
{\hyperref[\detokenize{fix::doc}]{\sphinxcrossref{\DUrole{doc}{fixes}}}} produce global quantities.  See the doc pages for
individual fixes for info on which ones produce such values.
{\hyperref[\detokenize{variable::doc}]{\sphinxcrossref{\DUrole{doc}{Variables}}}} of style \sphinxstyleemphasis{equal} and \sphinxstyleemphasis{vector} are the only
ones that can be used with this fix.  Variables of style \sphinxstyleemphasis{atom} cannot
be used, since they produce per\sphinxhyphen{}atom values.

\sphinxAtStartPar
The input values must either be all scalars or all vectors depending
on the setting of the \sphinxstyleemphasis{mode} keyword.  In both cases, the averaging is
performed independently on each input value (i.e., each input scalar
is averaged independently or each element of each input vector is
averaged independently).

\sphinxAtStartPar
If \sphinxstyleemphasis{mode} = scalar, then the input values must be scalars, or vectors
with a bracketed term appended, indicating the \(I^\text{th}\) value of the
vector is used.

\sphinxAtStartPar
If \sphinxstyleemphasis{mode} = vector, then the input values must be vectors, or arrays
with a bracketed term appended, indicating the Ith column of the array
is used.  All vectors must be the same length, which is the length of
the vector or number of rows in the array.


\bigskip\hrule\bigskip


\sphinxAtStartPar
For input values from a compute or fix or variable, the bracketed
index I can be specified using a wildcard asterisk with the index to
effectively specify multiple values.  This takes the form “*” or
“*n” or “m*” or “m*n”.  If \(N\) is the size of the vector (for \sphinxstyleemphasis{mode} =
scalar) or the number of columns in the array (for \sphinxstyleemphasis{mode} = vector),
then an asterisk with no numeric values means all indices from 1 to \(N\).
A leading asterisk means all indices from 1 to n (inclusive).  A trailing
asterisk means all indices from n to \(N\) (inclusive).  A middle asterisk
means all indices from m to n (inclusive).

\sphinxAtStartPar
Using a wildcard is the same as if the individual elements of the
vector or columns of the array had been listed one by one.  For example, the
following two fix ave/time commands are equivalent, since the {\hyperref[\detokenize{compute_rdf::doc}]{\sphinxcrossref{\DUrole{doc}{compute rdf}}}} command creates, in this case, a global array with three
columns, each of length 50:

\begin{sphinxVerbatim}[commandchars=\\\{\}]
\PYG{k}{compute }\PYG{n+nv+nvIdentifier}{myRDF}\PYG{+w}{ }\PYG{n+nv+nvIdentifier}{all}\PYG{+w}{ }\PYG{n}{rdf}\PYG{+w}{ }\PYG{l+m}{50}\PYG{+w}{ }\PYG{l+m}{1}\PYG{+w}{ }\PYG{l+m}{2}
\PYG{k}{fix }\PYG{n+nv+nvIdentifier}{1}\PYG{+w}{ }\PYG{n+nv+nvIdentifier}{all}\PYG{+w}{ }\PYG{n}{ave}\PYG{o}{/}\PYG{n}{time}\PYG{+w}{ }\PYG{l+m}{100}\PYG{+w}{ }\PYG{l+m}{1}\PYG{+w}{ }\PYG{l+m}{100}\PYG{+w}{ }\PYG{n}{c\PYGZus{}myRDF[}\PYG{o}{*}\PYG{n}{]}\PYG{+w}{ }\PYG{n}{file}\PYG{+w}{ }\PYG{n}{tmp1.rdf}\PYG{+w}{ }\PYG{n}{mode}\PYG{+w}{ }\PYG{n}{vector}
\PYG{k}{fix }\PYG{n+nv+nvIdentifier}{2}\PYG{+w}{ }\PYG{n+nv+nvIdentifier}{all}\PYG{+w}{ }\PYG{n}{ave}\PYG{o}{/}\PYG{n}{time}\PYG{+w}{ }\PYG{l+m}{100}\PYG{+w}{ }\PYG{l+m}{1}\PYG{+w}{ }\PYG{l+m}{100}\PYG{+w}{ }\PYG{n}{c\PYGZus{}myRDF[1]}\PYG{+w}{ }\PYG{n}{c\PYGZus{}myRDF[2]}\PYG{+w}{ }\PYG{n}{c\PYGZus{}myRDF[3]}\PYG{+w}{ }\PYG{n}{file}\PYG{+w}{ }\PYG{n}{tmp2.rdf}\PYG{+w}{ }\PYG{n}{mode}\PYG{+w}{ }\PYG{n}{vector}
\end{sphinxVerbatim}

\begin{sphinxadmonition}{note}{Note:}
\sphinxAtStartPar
For a vector\sphinxhyphen{}style variable, only the wildcard forms “*n” or
“m*n” are allowed.  You must specify the upper bound, because
vector\sphinxhyphen{}style variable lengths are not determined until the variable
is evaluated.  If n is specified larger than the vector length
turns out to be, zeroes are output for missing vector values.
\end{sphinxadmonition}


\bigskip\hrule\bigskip


\sphinxAtStartPar
The \(N_\text{every}\), \(N_\text{repeat}\), and \(N_\text{freq}\)
arguments specify on what time steps the input values will be used in order to
contribute to the average.  The final averaged quantities are generated on
time steps that are a multiple of \(N_\text{freq}\).  The average is over
\(N_\text{repeat}\) quantities, computed in the preceding portion of the
simulation every \(N_\text{every}\) time steps.  \(N_\text{freq}\) must
be a multiple of \(N_\text{every}\) and \(N_\text{every}\) must be
non\sphinxhyphen{}zero even if \(N_\text{repeat} = 1\). Also, the time steps
contributing to the average value cannot overlap (i.e.,
\(N_\text{repeat} \times N_\text{every}\) cannot exceed \(N_\text{freq}\)).

\sphinxAtStartPar
For example, if \(N_\text{every}=2\), \(N_\text{repeat}=6\), and
\(N_\text{freq}=100\), then values on time steps 90, 92, 94, 96, 98, and
100 will be used to compute the final average on time step 100.  Similarly for
time steps 190, 192, 194, 196, 198, and 200 on time step 200, etc.
If \(N_\text{repeat}=1\) and \(N_\text{freq} = 100\), then no time
averaging is done; values are simply generated on time steps
100, 200, etc.


\bigskip\hrule\bigskip


\sphinxAtStartPar
If a value begins with “c\_”, a compute ID must follow which has been
previously defined in the input script.  If \sphinxstyleemphasis{mode} = scalar, then if
no bracketed term is appended, the global scalar calculated by the
compute is used.  If a bracketed term is appended, the Ith element of
the global vector calculated by the compute is used.  If \sphinxstyleemphasis{mode} =
vector, then if no bracketed term is appended, the global vector
calculated by the compute is used.  If a bracketed term is appended,
the Ith column of the global array calculated by the compute is used.
See the discussion above for how I can be specified with a wildcard
asterisk to effectively specify multiple values.

\sphinxAtStartPar
Note that there is a {\hyperref[\detokenize{compute_reduce::doc}]{\sphinxcrossref{\DUrole{doc}{compute reduce}}}} command
that can sum per\sphinxhyphen{}atom quantities into a global scalar or vector, which
can then be accessed by fix ave/time.  It can also be a compute defined
not in your input script, but by {\hyperref[\detokenize{thermo_style::doc}]{\sphinxcrossref{\DUrole{doc}{thermodynamic output}}}} or other fixes such as {\hyperref[\detokenize{fix_nh::doc}]{\sphinxcrossref{\DUrole{doc}{fix nvt}}}} or
{\hyperref[\detokenize{fix_temp_rescale::doc}]{\sphinxcrossref{\DUrole{doc}{fix temp/rescale}}}}.  See the doc pages for
these commands which give the IDs of these computes.  Users can also
write code for their own compute styles and {\hyperref[\detokenize{Modify::doc}]{\sphinxcrossref{\DUrole{doc}{add them to LAMMPS}}}}.

\sphinxAtStartPar
If a value begins with “f\_”, a fix ID must follow which has been
previously defined in the input script.  If \sphinxstyleemphasis{mode} = scalar, then if
no bracketed term is appended, the global scalar calculated by the fix
is used.  If a bracketed term is appended, the Ith element of the
global vector calculated by the fix is used.  If \sphinxstyleemphasis{mode} = vector, then
if no bracketed term is appended, the global vector calculated by the
fix is used.  If a bracketed term is appended, the Ith column of the
global array calculated by the fix is used.  See the discussion above
for how I can be specified with a wildcard asterisk to effectively
specify multiple values.

\sphinxAtStartPar
Note that some fixes only produce their values on certain time steps,
which must be compatible with \sphinxstyleemphasis{Nevery}, else an error will result.
Users can also write code for their own fix styles and {\hyperref[\detokenize{Modify::doc}]{\sphinxcrossref{\DUrole{doc}{add them to LAMMPS}}}}.

\sphinxAtStartPar
If a value begins with “v\_”, a variable name must follow which has
been previously defined in the input script.  If \sphinxstyleemphasis{mode} = scalar, then
only equal\sphinxhyphen{}style or vector\sphinxhyphen{}style variables can be used, which both
produce global values.  In this mode, a vector\sphinxhyphen{}style variable requires
a bracketed term to specify the Ith element of the vector calculated
by the variable.  If \sphinxstyleemphasis{mode} = vector, then only a vector\sphinxhyphen{}style
variable can be used, without a bracketed term.  See the
{\hyperref[\detokenize{variable::doc}]{\sphinxcrossref{\DUrole{doc}{variable}}}} command for details.

\sphinxAtStartPar
Note that variables of style \sphinxstyleemphasis{equal} and \sphinxstyleemphasis{vector} define a formula
which can reference individual atom properties or thermodynamic
keywords, or they can invoke other computes, fixes, or variables when
they are evaluated, so this is a very general means of specifying
quantities to time average.


\bigskip\hrule\bigskip


\sphinxAtStartPar
Additional optional keywords also affect the operation of this fix.

\sphinxAtStartPar
If the \sphinxstyleemphasis{mode} keyword is set to \sphinxstyleemphasis{scalar}, then all input values must
be global scalars, or elements of global vectors.  If the \sphinxstyleemphasis{mode}
keyword is set to \sphinxstyleemphasis{vector}, then all input values must be global
vectors, or columns of global arrays.  They can also be global arrays,
which are converted into a series of global vectors (one per column),
as explained above.

\sphinxAtStartPar
The \sphinxstyleemphasis{ave} keyword determines how the values produced every
\(N_\text{freq}\) steps are averaged with values produced on previous steps
that were multiples of \(N_\text{freq}\), before they are accessed by
another output command or written to a file.

\sphinxAtStartPar
If the \sphinxstyleemphasis{ave} setting is \sphinxstyleemphasis{one}, then the values produced on time steps
that are multiples of \(N_\text{freq}\) are independent of each other; they
are output as\sphinxhyphen{}is without further averaging.

\sphinxAtStartPar
If the \sphinxstyleemphasis{ave} setting is \sphinxstyleemphasis{running}, then the values produced on
time steps that are multiples of \(N_\text{freq}\) are summed and averaged
in a cumulative sense before being output.  Each output value is thus the
average of the value produced on that time step with all preceding
values.  This running average begins when the fix is defined; it can
only be restarted by deleting the fix via the {\hyperref[\detokenize{unfix::doc}]{\sphinxcrossref{\DUrole{doc}{unfix}}}}
command, or by re\sphinxhyphen{}defining the fix by re\sphinxhyphen{}specifying it.

\sphinxAtStartPar
If the \sphinxstyleemphasis{ave} setting is \sphinxstyleemphasis{window}, then the values produced on
time steps that are multiples of \sphinxstyleemphasis{Nfreq} are summed and averaged within
a moving “window” of time, so that the last M values are used to
produce the output.  For example, if \(M = 3\) and
\(N_\text{freq} = 1000\), then the output on step 10000 will be the average
of the individual values on steps 8000, 9000, and 10000.  Outputs on early
steps will average over less than \(M\) values if they are not available.

\sphinxAtStartPar
The \sphinxstyleemphasis{start} keyword specifies what time step averaging will begin on.
The default is step 0.  Often input values can be 0.0 at time 0, so
setting \sphinxstyleemphasis{start} to a larger value can avoid including a 0.0 in a
running or windowed average.

\sphinxAtStartPar
The \sphinxstyleemphasis{off} keyword can be used to flag any of the input values.  If a
value is flagged, it will not be time averaged.  Instead the most
recent input value will always be stored and output.  This is useful
if one of more of the inputs produced by a compute or fix or variable
are effectively constant or are simply current values (e.g., they are
being written to a file with other time\sphinxhyphen{}averaged values for purposes
of creating well\sphinxhyphen{}formatted output).

\sphinxAtStartPar
\DUrole{versionmodified,added}{Added in version 17Apr2024: }new keyword \sphinxstyleemphasis{append}

\sphinxAtStartPar
The \sphinxstyleemphasis{file} or \sphinxstyleemphasis{append} keywords allow a filename to be specified.  If
\sphinxstyleemphasis{file} is used, then the filename is overwritten if it already exists.
If \sphinxstyleemphasis{append} is used, then the filename is appended to if it already
exists, or created if it does not exist.  Every \sphinxstyleemphasis{Nfreq} steps, one
quantity or vector of quantities is written to the file for each input
value specified in the fix ave/time command.  For \sphinxstyleemphasis{mode} = scalar, this
means a single line is written each time output is performed.  Thus the
file ends up to be a series of lines, i.e. one column of numbers for
each input value.  For \sphinxstyleemphasis{mode} = vector, an array of numbers is written
each time output is performed.  The number of rows is the length of the
input vectors, and the number of columns is the number of values.  Thus
the file ends up to be a series of these array sections.

\sphinxAtStartPar
\DUrole{versionmodified,added}{Added in version 4May2022.}

\sphinxAtStartPar
If the filename ends in ‘.yaml’ or ‘.yml’ then the output format
conforms to the \sphinxhref{https://yaml.org/}{YAML standard} which allows
easy import that data into tools and scripts that support reading YAML
files. The {\hyperref[\detokenize{Howto_structured_data::doc}]{\sphinxcrossref{\DUrole{doc}{structured data Howto}}}} contains
examples for parsing and plotting such data with very little programming
effort in Python using the \sphinxstyleemphasis{pyyaml}, \sphinxstyleemphasis{pandas}, and \sphinxstyleemphasis{matplotlib}
packages.

\sphinxAtStartPar
The \sphinxstyleemphasis{overwrite} keyword will continuously overwrite the output file
with the latest output, so that it only contains one time step worth of
output.  This option can only be used with the \sphinxstyleemphasis{ave running} setting.

\sphinxAtStartPar
The \sphinxstyleemphasis{format} keyword sets the numeric format of each value when it is
printed to a file via the \sphinxstyleemphasis{file} keyword.  Note that all values are
floating point quantities.  The default format is \%g.  You can specify
a higher precision if desired (e.g., \%20.16g).

\sphinxAtStartPar
The \sphinxstyleemphasis{title1} and \sphinxstyleemphasis{title2} and \sphinxstyleemphasis{title3} keywords allow specification of
the strings that will be printed as the first 2 or 3 lines of the
output file, assuming the \sphinxstyleemphasis{file} keyword was used.  LAMMPS uses
default values for each of these, so they do not need to be specified.

\sphinxAtStartPar
By default, these header lines are as follows for \sphinxstyleemphasis{mode} = scalar:

\begin{sphinxVerbatim}[commandchars=\\\{\}]
\PYGZsh{} Time\PYGZhy{}averaged data for fix ID
\PYGZsh{} TimeStep value1 value2 ...
\end{sphinxVerbatim}

\sphinxAtStartPar
In the first line, ID is replaced with the fix\sphinxhyphen{}ID.  In the second line
the values are replaced with the appropriate fields from the fix
ave/time command.  There is no third line in the header of the file,
so the \sphinxstyleemphasis{title3} setting is ignored when \sphinxstyleemphasis{mode} = scalar.

\sphinxAtStartPar
By default, these header lines are as follows for \sphinxstyleemphasis{mode} = vector:

\begin{sphinxVerbatim}[commandchars=\\\{\}]
\PYGZsh{} Time\PYGZhy{}averaged data for fix ID
\PYGZsh{} TimeStep Number\PYGZhy{}of\PYGZhy{}rows
\PYGZsh{} Row value1 value2 ...
\end{sphinxVerbatim}

\sphinxAtStartPar
In the first line, ID is replaced with the fix\sphinxhyphen{}ID.  The second line
describes the two values that are printed at the first of each section
of output.  In the third line the values are replaced with the
appropriate fields from the fix ave/time command.


\bigskip\hrule\bigskip



\subsection{Restart, fix\_modify, output, run start/stop, minimize info}
\label{\detokenize{fix_ave_time:restart-fix-modify-output-run-start-stop-minimize-info}}
\sphinxAtStartPar
\DUrole{versionmodified,added}{Added in version 4May2022.}

\sphinxAtStartPar
No information about this fix is written to {\hyperref[\detokenize{restart::doc}]{\sphinxcrossref{\DUrole{doc}{binary restart files}}}}.  The {\hyperref[\detokenize{fix_modify::doc}]{\sphinxcrossref{\DUrole{doc}{fix\_modify colname}}}} option can be
used to change the name of the column in the output file.  When writing
a YAML format file this name will be in the list of keywords.

\sphinxAtStartPar
This fix produces a global scalar or global vector or global array
which can be accessed by various {\hyperref[\detokenize{Howto_output::doc}]{\sphinxcrossref{\DUrole{doc}{output commands}}}}.
The values can only be accessed on time steps that are multiples of
\(N_\text{freq}\) since that is when averaging is performed.

\sphinxAtStartPar
A scalar is produced if only a single input value is averaged and
\sphinxstyleemphasis{mode} = scalar.  A vector is produced if multiple input values are
averaged for \sphinxstyleemphasis{mode} = scalar, or a single input value for \sphinxstyleemphasis{mode} =
vector.  In the first case, the length of the vector is the number of
inputs.  In the second case, the length of the vector is the same as
the length of the input vector.  An array is produced if multiple
input values are averaged and \sphinxstyleemphasis{mode} = vector.  The global array has \#
of rows = length of the input vectors and \# of columns = number of
inputs.

\sphinxAtStartPar
If the fix produces a scalar or vector, then the scalar and each
element of the vector can be either “intensive” or “extensive”,
depending on whether the values contributing to the scalar or vector
element are “intensive” or “extensive”.  If the fix produces an array,
then all elements in the array must be the same, either “intensive” or
“extensive”.  If a compute or fix provides the value being time
averaged, then the compute or fix determines whether the value is
intensive or extensive; see the page for that compute or fix for
further info.  Values produced by a variable are treated as intensive.

\sphinxAtStartPar
No parameter of this fix can be used with the \sphinxstyleemphasis{start/stop} keywords of
the {\hyperref[\detokenize{run::doc}]{\sphinxcrossref{\DUrole{doc}{run}}}} command.  This fix is not invoked during
{\hyperref[\detokenize{minimize::doc}]{\sphinxcrossref{\DUrole{doc}{energy minimization}}}}.


\subsection{Restrictions}
\label{\detokenize{fix_ave_time:restrictions}}\begin{quote}

\sphinxAtStartPar
none
\end{quote}


\subsection{Related commands}
\label{\detokenize{fix_ave_time:related-commands}}
\sphinxAtStartPar
{\hyperref[\detokenize{compute::doc}]{\sphinxcrossref{\DUrole{doc}{compute}}}}, {\hyperref[\detokenize{fix_ave_atom::doc}]{\sphinxcrossref{\DUrole{doc}{fix ave/atom}}}},
{\hyperref[\detokenize{fix_ave_chunk::doc}]{\sphinxcrossref{\DUrole{doc}{fix ave/chunk}}}}, {\hyperref[\detokenize{fix_ave_histo::doc}]{\sphinxcrossref{\DUrole{doc}{fix ave/histo}}}},
{\hyperref[\detokenize{variable::doc}]{\sphinxcrossref{\DUrole{doc}{variable}}}}, {\hyperref[\detokenize{fix_ave_correlate::doc}]{\sphinxcrossref{\DUrole{doc}{fix ave/correlate}}}},


\subsection{Default}
\label{\detokenize{fix_ave_time:default}}
\sphinxAtStartPar
The option defaults are mode = scalar, ave = one, start = 0, no file
output, format = \%g, title 1,2,3 = strings as described above, and no
off settings for any input values.

\sphinxstepscope

\index{fix aveforce@\spxentry{fix aveforce}}\ignorespaces 

\section{fix aveforce command}
\label{\detokenize{fix_aveforce:fix-aveforce-command}}\label{\detokenize{fix_aveforce:index-0}}\label{\detokenize{fix_aveforce::doc}}

\subsection{Syntax}
\label{\detokenize{fix_aveforce:syntax}}
\begin{sphinxVerbatim}[commandchars=\\\{\}]
\PYG{k}{fix }\PYG{n+nv+nvIdentifier}{ID}\PYG{+w}{ }\PYG{n+nv+nvIdentifier}{group\PYGZhy{}ID}\PYG{+w}{ }\PYG{n}{aveforce}\PYG{+w}{ }\PYG{n}{fx}\PYG{+w}{ }\PYG{n}{fy}\PYG{+w}{ }\PYG{n}{fz}\PYG{+w}{ }\PYG{n}{keyword}\PYG{+w}{ }\PYG{n}{value}\PYG{+w}{ }\PYG{n}{...}
\end{sphinxVerbatim}
\begin{itemize}
\item {} 
\sphinxAtStartPar
ID, group\sphinxhyphen{}ID are documented in {\hyperref[\detokenize{fix::doc}]{\sphinxcrossref{\DUrole{doc}{fix}}}} command

\item {} 
\sphinxAtStartPar
aveforce = style name of this fix command

\item {} 
\sphinxAtStartPar
fx,fy,fz = force component values (force units)

\begin{sphinxVerbatim}[commandchars=\\\{\}]
any of fx,fy,fz can be a variable (see below)
\end{sphinxVerbatim}

\item {} 
\sphinxAtStartPar
zero or more keyword/value pairs may be appended to args

\item {} 
\sphinxAtStartPar
keyword = \sphinxstyleemphasis{region}
\begin{sphinxalltt}
\sphinxstyleemphasis{region} value = region\sphinxhyphen{}ID
  region\sphinxhyphen{}ID = ID of region atoms must be in to have added force
\end{sphinxalltt}

\end{itemize}


\subsection{Examples}
\label{\detokenize{fix_aveforce:examples}}
\begin{sphinxVerbatim}[commandchars=\\\{\}]
\PYG{k}{fix }\PYG{n+nv+nvIdentifier}{pressdown}\PYG{+w}{ }\PYG{n+nv+nvIdentifier}{topwall}\PYG{+w}{ }\PYG{n}{aveforce}\PYG{+w}{ }\PYG{l+m}{0.0}\PYG{+w}{ }\PYG{o}{\PYGZhy{}}\PYG{l+m}{1.0}\PYG{+w}{ }\PYG{l+m}{0.0}
\PYG{k}{fix }\PYG{n+nv+nvIdentifier}{2}\PYG{+w}{ }\PYG{n+nv+nvIdentifier}{bottomwall}\PYG{+w}{ }\PYG{n}{aveforce}\PYG{+w}{ }\PYG{n}{NULL}\PYG{+w}{ }\PYG{o}{\PYGZhy{}}\PYG{l+m}{1.0}\PYG{+w}{ }\PYG{l+m}{0.0}\PYG{+w}{ }\PYG{k}{region }\PYG{n+nv+nvIdentifier}{top}
\PYG{k}{fix }\PYG{n+nv+nvIdentifier}{2}\PYG{+w}{ }\PYG{n+nv+nvIdentifier}{bottomwall}\PYG{+w}{ }\PYG{n}{aveforce}\PYG{+w}{ }\PYG{n}{NULL}\PYG{+w}{ }\PYG{o}{\PYGZhy{}}\PYG{l+m}{1.0}\PYG{+w}{ }\PYG{n}{v\PYGZus{}oscillate}\PYG{+w}{ }\PYG{k}{region }\PYG{n+nv+nvIdentifier}{top}
\end{sphinxVerbatim}


\subsection{Description}
\label{\detokenize{fix_aveforce:description}}
\sphinxAtStartPar
Apply an additional external force to a group of atoms in such a way
that every atom experiences the same force.  This is useful for
pushing on wall or boundary atoms so that the structure of the wall
does not change over time.

\sphinxAtStartPar
The existing force is averaged for the group of atoms, component by
component.  The actual force on each atom is then set to the average
value plus the component specified in this command.  This means each
atom in the group receives the same force.

\sphinxAtStartPar
Any of the \sphinxstyleemphasis{fx}, \sphinxstyleemphasis{fy}, or \sphinxstyleemphasis{fz} values can be specified as \sphinxcode{\sphinxupquote{NULL}}, which
means the force in that dimension is not changed.  Note that this is not the
same as specifying a 0.0 value, since that sets all forces to the same
average value without adding in any additional force.

\sphinxAtStartPar
Any of the three quantities defining the force components, namely \sphinxstyleemphasis{fx}, \sphinxstyleemphasis{fy},
and \sphinxstyleemphasis{fz}, can be specified as an equal\sphinxhyphen{}style {\hyperref[\detokenize{variable::doc}]{\sphinxcrossref{\DUrole{doc}{variable}}}}.
If the value is a variable, it should be specified as v\_name, where
name is the variable name.  In this case, the variable will be
evaluated each timestep, and its value used to determine the average
force.

\sphinxAtStartPar
Equal\sphinxhyphen{}style variables can specify formulas with various mathematical
functions, and include {\hyperref[\detokenize{thermo_style::doc}]{\sphinxcrossref{\DUrole{doc}{thermo\_style}}}} command
keywords for the simulation box parameters and timestep and elapsed
time.  Thus it is easy to specify a time\sphinxhyphen{}dependent average force.

\sphinxAtStartPar
If the \sphinxstyleemphasis{region} keyword is used, the atom must also be in the
specified geometric {\hyperref[\detokenize{region::doc}]{\sphinxcrossref{\DUrole{doc}{region}}}} in order to have force added
to it.


\bigskip\hrule\bigskip


\sphinxAtStartPar
Styles with a \sphinxstyleemphasis{gpu}, \sphinxstyleemphasis{intel}, \sphinxstyleemphasis{kk}, \sphinxstyleemphasis{omp}, or \sphinxstyleemphasis{opt} suffix are
functionally the same as the corresponding style without the suffix.
They have been optimized to run faster, depending on your available
hardware, as discussed on the {\hyperref[\detokenize{Speed_packages::doc}]{\sphinxcrossref{\DUrole{doc}{Accelerator packages}}}}
page.  The accelerated styles take the same arguments and should
produce the same results, except for round\sphinxhyphen{}off and precision issues.

\sphinxAtStartPar
These accelerated styles are part of the GPU, INTEL, KOKKOS,
OPENMP, and OPT packages, respectively.  They are only enabled if
LAMMPS was built with those packages.  See the {\hyperref[\detokenize{Build_package::doc}]{\sphinxcrossref{\DUrole{doc}{Build package}}}} page for more info.

\sphinxAtStartPar
You can specify the accelerated styles explicitly in your input script
by including their suffix, or you can use the {\hyperref[\detokenize{Run_options::doc}]{\sphinxcrossref{\DUrole{doc}{\sphinxhyphen{}suffix command\sphinxhyphen{}line switch}}}} when you invoke LAMMPS, or you can use the
{\hyperref[\detokenize{suffix::doc}]{\sphinxcrossref{\DUrole{doc}{suffix}}}} command in your input script.

\sphinxAtStartPar
See the {\hyperref[\detokenize{Speed_packages::doc}]{\sphinxcrossref{\DUrole{doc}{Accelerator packages}}}} page for more
instructions on how to use the accelerated styles effectively.


\bigskip\hrule\bigskip



\subsection{Restart, fix\_modify, output, run start/stop, minimize info}
\label{\detokenize{fix_aveforce:restart-fix-modify-output-run-start-stop-minimize-info}}
\sphinxAtStartPar
No information about this fix is written to
{\hyperref[\detokenize{restart::doc}]{\sphinxcrossref{\DUrole{doc}{binary restart files}}}}.

\sphinxAtStartPar
The {\hyperref[\detokenize{fix_modify::doc}]{\sphinxcrossref{\DUrole{doc}{fix\_modify}}}} \sphinxstyleemphasis{respa} option is supported by this
fix. This allows to set at which level of the {\hyperref[\detokenize{run_style::doc}]{\sphinxcrossref{\DUrole{doc}{r\sphinxhyphen{}RESPA}}}}
integrator the fix is adding its forces. Default is the outermost level.

\sphinxAtStartPar
This fix computes a global three\sphinxhyphen{}vector of forces, which can be accessed
by various {\hyperref[\detokenize{Howto_output::doc}]{\sphinxcrossref{\DUrole{doc}{output commands}}}}.  This is the total
force on the group of atoms before the forces on individual atoms are
changed by the fix.  The vector values calculated by this fix are “extensive”.

\sphinxAtStartPar
No parameter of this fix can be used with the \sphinxstyleemphasis{start/stop} keywords of
the {\hyperref[\detokenize{run::doc}]{\sphinxcrossref{\DUrole{doc}{run}}}} command.

\sphinxAtStartPar
The forces due to this fix are imposed during an energy minimization,
invoked by the {\hyperref[\detokenize{minimize::doc}]{\sphinxcrossref{\DUrole{doc}{minimize}}}} command.  You should not
specify force components with a variable that has time\sphinxhyphen{}dependence for
use with a minimizer, since the minimizer increments the timestep as
the iteration count during the minimization.


\subsection{Restrictions}
\label{\detokenize{fix_aveforce:restrictions}}\begin{quote}

\sphinxAtStartPar
none
\end{quote}


\subsection{Related commands}
\label{\detokenize{fix_aveforce:related-commands}}
\sphinxAtStartPar
{\hyperref[\detokenize{fix_setforce::doc}]{\sphinxcrossref{\DUrole{doc}{fix setforce}}}}, {\hyperref[\detokenize{fix_addforce::doc}]{\sphinxcrossref{\DUrole{doc}{fix addforce}}}}


\subsection{Default}
\label{\detokenize{fix_aveforce:default}}
\sphinxAtStartPar
none

\sphinxstepscope

\index{fix balance@\spxentry{fix balance}}\ignorespaces 

\section{fix balance command}
\label{\detokenize{fix_balance:fix-balance-command}}\label{\detokenize{fix_balance:index-0}}\label{\detokenize{fix_balance::doc}}

\subsection{Syntax}
\label{\detokenize{fix_balance:syntax}}
\begin{sphinxVerbatim}[commandchars=\\\{\}]
\PYG{k}{fix }\PYG{n+nv+nvIdentifier}{ID}\PYG{+w}{ }\PYG{n+nv+nvIdentifier}{group\PYGZhy{}ID}\PYG{+w}{ }\PYG{n}{balance}\PYG{+w}{ }\PYG{n}{Nfreq}\PYG{+w}{ }\PYG{n}{thresh}\PYG{+w}{ }\PYG{n}{style}\PYG{+w}{ }\PYG{n}{args}\PYG{+w}{ }\PYG{n}{keyword}\PYG{+w}{ }\PYG{n}{args}\PYG{+w}{ }\PYG{n}{...}
\end{sphinxVerbatim}
\begin{itemize}
\item {} 
\sphinxAtStartPar
ID, group\sphinxhyphen{}ID are documented in {\hyperref[\detokenize{fix::doc}]{\sphinxcrossref{\DUrole{doc}{fix}}}} command

\item {} 
\sphinxAtStartPar
balance = style name of this fix command

\item {} 
\sphinxAtStartPar
Nfreq = perform dynamic load balancing every this many steps

\item {} 
\sphinxAtStartPar
thresh = imbalance threshold that must be exceeded to perform a re\sphinxhyphen{}balance

\item {} 
\sphinxAtStartPar
style = \sphinxstyleemphasis{shift} or \sphinxstyleemphasis{rcb} or \sphinxstyleemphasis{report}
\begin{sphinxalltt}
\sphinxstyleemphasis{shift} args = dimstr Niter stopthresh
  dimstr = sequence of letters containing \sphinxstyleemphasis{x} or \sphinxstyleemphasis{y} or \sphinxstyleemphasis{z}, each not more than once
  Niter = \# of times to iterate within each dimension of dimstr sequence
  stopthresh = stop balancing when this imbalance threshold is reached
\sphinxstyleemphasis{rcb} args = none
\sphinxstyleemphasis{report} args = none
\end{sphinxalltt}

\item {} 
\sphinxAtStartPar
zero or more keyword/arg pairs may be appended

\item {} 
\sphinxAtStartPar
keyword = \sphinxstyleemphasis{weight} or \sphinxstyleemphasis{out}
\begin{sphinxalltt}
\sphinxstyleemphasis{weight} style args = use weighted particle counts for the balancing
  \sphinxstyleemphasis{style} = \sphinxstyleemphasis{group} or \sphinxstyleemphasis{neigh} or \sphinxstyleemphasis{time} or \sphinxstyleemphasis{var} or \sphinxstyleemphasis{store}
    \sphinxstyleemphasis{group} args = Ngroup group1 weight1 group2 weight2 ...
      Ngroup = number of groups with assigned weights
      group1, group2, ... = group IDs
      weight1, weight2, ...   = corresponding weight factors
    \sphinxstyleemphasis{neigh} factor = compute weight based on number of neighbors
      factor = scaling factor (\textgreater{} 0)
    \sphinxstyleemphasis{time} factor = compute weight based on time spend computing
      factor = scaling factor (\textgreater{} 0)
    \sphinxstyleemphasis{var} name = take weight from atom\sphinxhyphen{}style variable
      name = name of the atom\sphinxhyphen{}style variable
    \sphinxstyleemphasis{store} name = store weight in custom atom property defined by {\hyperref[\detokenize{fix_property_atom::doc}]{\sphinxcrossref{\DUrole{doc}{fix property/atom}}}} command
      name = atom property name (without d\_ prefix)
\sphinxstyleemphasis{sort} arg = \sphinxstyleemphasis{no} or \sphinxstyleemphasis{yes}
\sphinxstyleemphasis{out} arg = filename
  filename = write each processor\textquotesingle{}s subdomain to a file, at each re\sphinxhyphen{}balancing
\end{sphinxalltt}

\end{itemize}


\subsection{Examples}
\label{\detokenize{fix_balance:examples}}
\begin{sphinxVerbatim}[commandchars=\\\{\}]
\PYG{k}{fix }\PYG{n+nv+nvIdentifier}{2}\PYG{+w}{ }\PYG{n+nv+nvIdentifier}{all}\PYG{+w}{ }\PYG{n}{balance}\PYG{+w}{ }\PYG{l+m}{1000}\PYG{+w}{ }\PYG{l+m}{1.05}\PYG{+w}{ }\PYG{n}{shift}\PYG{+w}{ }\PYG{n}{x}\PYG{+w}{ }\PYG{l+m}{10}\PYG{+w}{ }\PYG{l+m}{1.05}
\PYG{k}{fix }\PYG{n+nv+nvIdentifier}{2}\PYG{+w}{ }\PYG{n+nv+nvIdentifier}{all}\PYG{+w}{ }\PYG{n}{balance}\PYG{+w}{ }\PYG{l+m}{100}\PYG{+w}{ }\PYG{l+m}{0.9}\PYG{+w}{ }\PYG{n}{shift}\PYG{+w}{ }\PYG{n}{xy}\PYG{+w}{ }\PYG{l+m}{20}\PYG{+w}{ }\PYG{l+m}{1.1}\PYG{+w}{ }\PYG{n}{out}\PYG{+w}{ }\PYG{n}{tmp.balance}
\PYG{k}{fix }\PYG{n+nv+nvIdentifier}{2}\PYG{+w}{ }\PYG{n+nv+nvIdentifier}{all}\PYG{+w}{ }\PYG{n}{balance}\PYG{+w}{ }\PYG{l+m}{100}\PYG{+w}{ }\PYG{l+m}{0.9}\PYG{+w}{ }\PYG{n}{shift}\PYG{+w}{ }\PYG{n}{xy}\PYG{+w}{ }\PYG{l+m}{20}\PYG{+w}{ }\PYG{l+m}{1.1}\PYG{+w}{ }\PYG{n}{weight}\PYG{+w}{ }\PYG{k}{group }\PYG{n+nv+nvIdentifier}{3}\PYG{+w}{ }\PYG{n}{substrate}\PYG{+w}{ }\PYG{l+m}{3.0}\PYG{+w}{ }\PYG{n}{solvent}\PYG{+w}{ }\PYG{l+m}{1.0}\PYG{+w}{ }\PYG{n}{solute}\PYG{+w}{ }\PYG{l+m}{0.8}\PYG{+w}{ }\PYG{n}{out}\PYG{+w}{ }\PYG{n}{tmp.balance}
\PYG{k}{fix }\PYG{n+nv+nvIdentifier}{2}\PYG{+w}{ }\PYG{n+nv+nvIdentifier}{all}\PYG{+w}{ }\PYG{n}{balance}\PYG{+w}{ }\PYG{l+m}{100}\PYG{+w}{ }\PYG{l+m}{1.0}\PYG{+w}{ }\PYG{n}{shift}\PYG{+w}{ }\PYG{n}{x}\PYG{+w}{ }\PYG{l+m}{10}\PYG{+w}{ }\PYG{l+m}{1.1}\PYG{+w}{ }\PYG{n}{weight}\PYG{+w}{ }\PYG{n}{time}\PYG{+w}{ }\PYG{l+m}{0.8}
\PYG{k}{fix }\PYG{n+nv+nvIdentifier}{2}\PYG{+w}{ }\PYG{n+nv+nvIdentifier}{all}\PYG{+w}{ }\PYG{n}{balance}\PYG{+w}{ }\PYG{l+m}{100}\PYG{+w}{ }\PYG{l+m}{1.0}\PYG{+w}{ }\PYG{n}{shift}\PYG{+w}{ }\PYG{n}{xy}\PYG{+w}{ }\PYG{l+m}{5}\PYG{+w}{ }\PYG{l+m}{1.1}\PYG{+w}{ }\PYG{n}{weight}\PYG{+w}{ }\PYG{n}{var}\PYG{+w}{ }\PYG{n}{myweight}\PYG{+w}{ }\PYG{n}{weight}\PYG{+w}{ }\PYG{n}{neigh}\PYG{+w}{ }\PYG{l+m}{0.6}\PYG{+w}{ }\PYG{n}{weight}\PYG{+w}{ }\PYG{n}{store}\PYG{+w}{ }\PYG{n}{allweight}
\PYG{k}{fix }\PYG{n+nv+nvIdentifier}{2}\PYG{+w}{ }\PYG{n+nv+nvIdentifier}{all}\PYG{+w}{ }\PYG{n}{balance}\PYG{+w}{ }\PYG{l+m}{1000}\PYG{+w}{ }\PYG{l+m}{1.1}\PYG{+w}{ }\PYG{n}{rcb}
\end{sphinxVerbatim}


\subsection{Description}
\label{\detokenize{fix_balance:description}}
\sphinxAtStartPar
This command adjusts the size and shape of processor subdomains
within the simulation box, to attempt to balance the number of
particles and thus the computational cost (load) evenly across
processors.  The load balancing is “dynamic” in the sense that
re\sphinxhyphen{}balancing is performed periodically during the simulation.  To
perform “static” balancing, before or between runs, see the
{\hyperref[\detokenize{balance::doc}]{\sphinxcrossref{\DUrole{doc}{balance}}}} command.

\sphinxAtStartPar
\DUrole{versionmodified,added}{Added in version 17Apr2024.}

\sphinxAtStartPar
The \sphinxstyleemphasis{report} balance style only computes the load imbalance but
does not attempt any re\sphinxhyphen{}balancing.  This way the load imbalance
information can be used otherwise, for instance for stopping a
run with {\hyperref[\detokenize{fix_halt::doc}]{\sphinxcrossref{\DUrole{doc}{fix halt}}}}.

\sphinxAtStartPar
Load\sphinxhyphen{}balancing is typically most useful if the particles in the
simulation box have a spatially\sphinxhyphen{}varying density distribution or
where the computational cost varies significantly between different
atoms (e.g., a model of a vapor/liquid interface, or a solid with
an irregular\sphinxhyphen{}shaped geometry containing void regions, or
{\hyperref[\detokenize{pair_hybrid::doc}]{\sphinxcrossref{\DUrole{doc}{hybrid pair style simulations}}}} that combine
pair styles with different computational cost).  In these cases, the
LAMMPS default of dividing the simulation box volume into a
regular\sphinxhyphen{}spaced grid of 3d bricks, with one equal\sphinxhyphen{}volume subdomain
per processor, may assign numbers of particles per processor in a
way that the computational effort varies significantly.  This can
lead to poor performance when the simulation is run in parallel.

\sphinxAtStartPar
The balancing can be performed with or without per\sphinxhyphen{}particle weighting.
With no weighting, the balancing attempts to assign an equal number of
particles to each processor.  With weighting, the balancing attempts
to assign an equal aggregate computational weight to each processor,
which typically induces a different number of atoms assigned to each
processor.

\begin{sphinxadmonition}{note}{Note:}
\sphinxAtStartPar
The weighting options listed above are documented with the
{\hyperref[\detokenize{balance::doc}]{\sphinxcrossref{\DUrole{doc}{balance}}}} command in {\hyperref[\detokenize{balance:weighted-balance}]{\sphinxcrossref{\DUrole{std,std-ref}{this section of the balance
command}}}} doc page.  That section
describes the various weighting options and gives a few examples of
how they can be used.  The weighting options are the same for both the
fix balance and {\hyperref[\detokenize{balance::doc}]{\sphinxcrossref{\DUrole{doc}{balance}}}} commands.
\end{sphinxadmonition}

\sphinxAtStartPar
Note that the {\hyperref[\detokenize{processors::doc}]{\sphinxcrossref{\DUrole{doc}{processors}}}} command allows some control
over how the box volume is split across processors.  Specifically, for
a \(P_x \times P_y \times P_z\) grid of processors, it allows choices of
\(P_x\), \(P_y\), and \(P_z\) subject to the constraint that
\(P_x P_y P_z = P\), the total number of processors.
This is sufficient to achieve good load\sphinxhyphen{}balance for
some problems on some processor counts.  However, all the processor
subdomains will still have the same shape and the same volume.

\sphinxAtStartPar
On a particular time step, a load\sphinxhyphen{}balancing operation is only performed
if the current “imbalance factor” in particles owned by each processor
exceeds the specified \sphinxstyleemphasis{thresh} parameter.  The imbalance factor is
defined as the maximum number of particles (or weight) owned by any
processor, divided by the average number of particles (or weight) per
processor.  Thus, an imbalance factor of 1.0 is perfect balance.

\sphinxAtStartPar
As an example, for 10000 particles running on 10 processors, if the
most heavily loaded processor has 1200 particles, then the imbalance factor is
1.2, meaning there is a 20\% imbalance.  Note that re\sphinxhyphen{}balances can be
forced even if the current balance is perfect (1.0) be specifying a
\sphinxstyleemphasis{thresh} \textless{} 1.0.

\begin{sphinxadmonition}{note}{Note:}
\sphinxAtStartPar
This command attempts to minimize the imbalance factor, as
defined above.  But depending on the method a perfect balance (1.0)
may not be achieved.  For example, “grid” methods (defined below) that
create a logical 3d grid cannot achieve perfect balance for many
irregular distributions of particles.  Likewise, if a portion of the
system is a perfect, non\sphinxhyphen{}rotated lattice (e.g., the initial system is
generated by the {\hyperref[\detokenize{create_atoms::doc}]{\sphinxcrossref{\DUrole{doc}{create\_atoms}}}} command with no
rotations), then “grid” methods may be unable to achieve exact balance.
This is because entire lattice planes will be owned or not owned by a single
processor.
\end{sphinxadmonition}

\begin{sphinxadmonition}{note}{Note:}
\sphinxAtStartPar
The imbalance factor is also an estimate of the maximum speed\sphinxhyphen{}up
you can hope to achieve by running a perfectly balanced simulation
versus an imbalanced one.  In the example above, the 10000\sphinxhyphen{}particle
simulation could run up to 20\% faster if it were perfectly balanced,
versus when imbalanced.  However, computational cost is not strictly
proportional to particle count, and changing the relative size and
shape of processor subdomains may lead to additional computational
and communication overheads (e.g., in the PPPM solver used via the
{\hyperref[\detokenize{kspace_style::doc}]{\sphinxcrossref{\DUrole{doc}{kspace\_style}}}} command).  Thus, you should benchmark
the run times of a simulation before and after balancing.
\end{sphinxadmonition}


\bigskip\hrule\bigskip


\sphinxAtStartPar
The method used to perform a load balance is specified by one of the
listed styles, which are described in detail below.  There are two kinds
of styles.

\sphinxAtStartPar
The \sphinxstyleemphasis{shift} style is a “grid” method which produces a logical 3d grid
of processors.  It operates by changing the cutting planes (or lines)
between processors in 3d (or 2d), to adjust the volume (area in 2d)
assigned to each processor, as in the following 2d diagram where
processor subdomains are shown and atoms are colored by the processor
that owns them.

\sphinxAtStartPar
\sphinxincludegraphics[width=0.320\linewidth]{{balance_uniform}.jpg}  \sphinxincludegraphics[width=0.320\linewidth]{{balance_nonuniform}.jpg}  \sphinxincludegraphics[width=0.320\linewidth]{{balance_rcb}.jpg}

\sphinxAtStartPar
The leftmost diagram is the default partitioning of the simulation box
across processors (one sub\sphinxhyphen{}box for each of 16 processors); the middle
diagram is after a “grid” method has been applied. The \sphinxstyleemphasis{rcb} style is a
“tiling” method which does not produce a logical 3d grid of processors.
Rather it tiles the simulation domain with rectangular sub\sphinxhyphen{}boxes of
varying size and shape in an irregular fashion so as to have equal
numbers of particles (or weight) in each sub\sphinxhyphen{}box, as in the rightmost
diagram above.

\sphinxAtStartPar
The “grid” methods can be used with either of the
{\hyperref[\detokenize{comm_style::doc}]{\sphinxcrossref{\DUrole{doc}{comm\_style}}}} command options, \sphinxstyleemphasis{brick} or \sphinxstyleemphasis{tiled}.  The
“tiling” methods can only be used with {\hyperref[\detokenize{comm_style::doc}]{\sphinxcrossref{\DUrole{doc}{comm\_style tiled}}}}.

\sphinxAtStartPar
When a “grid” method is specified, the current domain partitioning can
be either a logical 3d grid or a tiled partitioning.  In the former
case, the current logical 3d grid is used as a starting point and
changes are made to improve the imbalance factor.  In the latter case,
the tiled partitioning is discarded and a logical 3d grid is created
with uniform spacing in all dimensions.  This is the starting point
for the balancing operation.

\sphinxAtStartPar
When a “tiling” method is specified, the current domain partitioning
(“grid” or “tiled”) is ignored, and a new partitioning is computed
from scratch.


\bigskip\hrule\bigskip


\sphinxAtStartPar
The \sphinxstyleemphasis{group\sphinxhyphen{}ID} is ignored.  However the impact of balancing on
different groups of atoms can be affected by using the \sphinxstyleemphasis{group} weight
style as described below.

\sphinxAtStartPar
The \(N_\text{freq}\) setting determines how often a re\sphinxhyphen{}balance is
performed.  If \(N_\text{freq} > 0\), then re\sphinxhyphen{}balancing will occur every
\(N_\text{freq}\) steps.  Each time a re\sphinxhyphen{}balance occurs, a reneighboring is
triggered, so \(N_\text{freq}\) should not be too small.  If
\(N_\text{freq} = 0\), then re\sphinxhyphen{}balancing will be done every time
reneighboring normally occurs, as determined by the
the {\hyperref[\detokenize{neighbor::doc}]{\sphinxcrossref{\DUrole{doc}{neighbor}}}} and {\hyperref[\detokenize{neigh_modify::doc}]{\sphinxcrossref{\DUrole{doc}{neigh\_modify}}}}
command settings.

\sphinxAtStartPar
On re\sphinxhyphen{}balance steps, re\sphinxhyphen{}balancing will only be attempted if the current
imbalance factor, as defined above, exceeds the \sphinxstyleemphasis{thresh} setting.


\bigskip\hrule\bigskip


\sphinxAtStartPar
The \sphinxstyleemphasis{shift} style invokes a “grid” method for balancing, as described
above.  It changes the positions of cutting planes between processors
in an iterative fashion, seeking to reduce the imbalance factor.

\sphinxAtStartPar
The \sphinxstyleemphasis{dimstr} argument is a string of characters, each of which must be
\sphinxstyleemphasis{x} or \sphinxstyleemphasis{y} or \sphinxstyleemphasis{z}.  Each character can appear zero or one time,
since there is no advantage to balancing on a dimension more than
once.  You should normally only list dimensions where you expect there
to be a density variation in the particles.

\sphinxAtStartPar
Balancing proceeds by adjusting the cutting planes in each of the
dimensions listed in \sphinxstyleemphasis{dimstr}, one dimension at a time.  For a single
dimension, the balancing operation (described below) is iterated on up
to \(N_\text{iter}\) times.  After each dimension finishes, the imbalance
factor is re\sphinxhyphen{}computed, and the balancing operation halts if the \sphinxstyleemphasis{stopthresh}
criterion is met.

\sphinxAtStartPar
A re\sphinxhyphen{}balance operation in a single dimension is performed using a
density\sphinxhyphen{}dependent recursive multisectioning algorithm, where the
position of each cutting plane (line in 2d) in the dimension is
adjusted independently.  This is similar to a recursive bisectioning
for a single value, except that the bounds used for each bisectioning
take advantage of information from neighboring cuts if possible, as
well as counts of particles at the bounds on either side of each cuts,
which themselves were cuts in previous iterations.  The latter is used
to infer a density of particles near each of the current cuts.  At
each iteration, the count of particles on either side of each plane is
tallied.  If the counts do not match the target value for the plane,
the position of the cut is adjusted based on the local density.  The
low and high bounds are adjusted on each iteration, using new count
information, so that they become closer together over time.  Thus as
the recursion progresses, the count of particles on either side of the
plane gets closer to the target value.

\sphinxAtStartPar
The density\sphinxhyphen{}dependent part of this algorithm is often an advantage
when you re\sphinxhyphen{}balance a system that is already nearly balanced.  It
typically converges more quickly than the geometric bisectioning
algorithm used by the {\hyperref[\detokenize{balance::doc}]{\sphinxcrossref{\DUrole{doc}{balance}}}} command.  However, if can
be a disadvantage if you attempt to re\sphinxhyphen{}balance a system that is far
from balanced, and converge more slowly.  In this case you probably
want to use the {\hyperref[\detokenize{balance::doc}]{\sphinxcrossref{\DUrole{doc}{balance}}}} command before starting a run,
so that you begin the run with a balanced system.

\sphinxAtStartPar
Once the re\sphinxhyphen{}balancing is complete and final processor subdomains
assigned, particles migrate to their new owning processor as part of
the normal reneighboring procedure.

\begin{sphinxadmonition}{note}{Note:}
\sphinxAtStartPar
At each re\sphinxhyphen{}balance operation, the bisectioning for each cutting
plane (line in 2d) typically starts with low and high bounds separated
by the extent of a processor’s subdomain in one dimension.  The size
of this bracketing region shrinks based on the local density, as
described above, which should typically be 1/2 or more every
iteration.  Thus if \(N_\text{iter}\) is specified as 10, the cutting
plane will typically be positioned to better than 1 part in 1000 accuracy
(relative to the perfect target position).  For \(N_\text{iter} = 20\),
it will be accurate to better than 1 part in a million.  Thus there is no
need to set \(N_\text{iter}\) to a large value.  This is especially true
if you are re\sphinxhyphen{}balancing often enough that each time you expect only an
incremental adjustment in the cutting planes is necessary.  LAMMPS will
check if the threshold accuracy is reached (in a dimension) is less
iterations than \(N_\text{iter}\) and exit early.
\end{sphinxadmonition}


\bigskip\hrule\bigskip


\sphinxAtStartPar
The \sphinxstyleemphasis{rcb} style invokes a “tiled” method for balancing, as described
above.  It performs a recursive coordinate bisectioning (RCB) of the
simulation domain. The basic idea is as follows.

\sphinxAtStartPar
The simulation domain is cut into two boxes by an axis\sphinxhyphen{}aligned cut in
the longest dimension, leaving one new box on either side of the cut.
All the processors are also partitioned into two groups, half assigned
to the box on the lower side of the cut, and half to the box on the
upper side.  If the processor count is odd, one side gets an extra
processor.  The cut is positioned so that the number of atoms in the
lower box is exactly the number that the processors assigned to that
box should own for load balance to be perfect.  This also makes load
balance for the upper box perfect.  The positioning is done
iteratively, by a bisectioning method.  Note that counting atoms on
either side of the cut requires communication between all processors
at each iteration.

\sphinxAtStartPar
That is the procedure for the first cut.  Subsequent cuts are made
recursively, in exactly the same manner.  The subset of processors
assigned to each box make a new cut in the longest dimension of that
box, splitting the box, the subset of processors, and the atoms in
the box in two.  The recursion continues until every processor is
assigned a sub\sphinxhyphen{}box of the entire simulation domain, and owns the atoms
in that sub\sphinxhyphen{}box.


\bigskip\hrule\bigskip


\sphinxAtStartPar
The \sphinxstyleemphasis{sort} keyword determines whether the communication of per\sphinxhyphen{}atom
data to other processors during load\sphinxhyphen{}balancing will be random or
deterministic.  Random is generally faster; deterministic will ensure
the new ordering of atoms on each processor is the same each time the
same simulation is run.  This can be useful for debugging purposes.
Since the fix balance command is performed during timestepping, the
default is \sphinxstyleemphasis{no} so that sorting is not performed.

\sphinxAtStartPar
The \sphinxstyleemphasis{out} keyword writes text to the specified \sphinxstyleemphasis{filename} with the
results of each re\sphinxhyphen{}balancing operation.  The file contains the bounds
of the subdomain for each processor after the balancing operation
completes.  The format of the file is compatible with the
\sphinxhref{https://lammps.github.io/pizza}{Pizza.py} \sphinxstyleemphasis{mdump} tool which has support for manipulating and
visualizing mesh files.  An example is shown here for a balancing by four
processors for a 2d problem:

\begin{sphinxVerbatim}[commandchars=\\\{\}]
ITEM: TIMESTEP
0
ITEM: NUMBER OF NODES
16
ITEM: BOX BOUNDS
0 10
0 10
0 10
ITEM: NODES
1 1 0 0 0
2 1 5 0 0
3 1 5 5 0
4 1 0 5 0
5 1 5 0 0
6 1 10 0 0
7 1 10 5 0
8 1 5 5 0
9 1 0 5 0
10 1 5 5 0
11 1 5 10 0
12 1 10 5 0
13 1 5 5 0
14 1 10 5 0
15 1 10 10 0
16 1 5 10 0
ITEM: TIMESTEP
0
ITEM: NUMBER OF SQUARES
4
ITEM: SQUARES
1 1 1 2 3 4
2 1 5 6 7 8
3 1 9 10 11 12
4 1 13 14 15 16
\end{sphinxVerbatim}

\sphinxAtStartPar
The coordinates of all the vertices are listed in the NODES section, five
per processor.  Note that the four subdomains share vertices, so there
will be duplicate nodes in the list.

\sphinxAtStartPar
The “SQUARES” section lists the node IDs of the four vertices in a
rectangle for each processor (1 to 4).

\sphinxAtStartPar
For a 3d problem, the syntax is similar but with eight vertices listed for
each processor instead of four, and “SQUARES” replaced by “CUBES”.


\bigskip\hrule\bigskip



\subsection{Restart, fix\_modify, output, run start/stop, minimize info}
\label{\detokenize{fix_balance:restart-fix-modify-output-run-start-stop-minimize-info}}
\sphinxAtStartPar
No information about this fix is written to
{\hyperref[\detokenize{restart::doc}]{\sphinxcrossref{\DUrole{doc}{binary restart files}}}}.  None of the
{\hyperref[\detokenize{fix_modify::doc}]{\sphinxcrossref{\DUrole{doc}{fix\_modify}}}} options are relevant to this fix.

\sphinxAtStartPar
This fix computes a global scalar which is the imbalance factor
after the most recent re\sphinxhyphen{}balance and a global vector of length 3 with
additional information about the most recent re\sphinxhyphen{}balancing.  The three
values in the vector are as follows:
\begin{enumerate}
\sphinxsetlistlabels{\arabic}{enumi}{enumii}{}{.}%
\item {} 
\sphinxAtStartPar
max \# of particles per processor

\item {} 
\sphinxAtStartPar
total \# iterations performed in last re\sphinxhyphen{}balance

\item {} 
\sphinxAtStartPar
imbalance factor right before the last re\sphinxhyphen{}balance was performed

\end{enumerate}

\sphinxAtStartPar
As explained above, the imbalance factor is the ratio of the maximum
number of particles (or total weight) on any processor to the average
number of particles (or total weight) per processor.

\sphinxAtStartPar
These quantities can be accessed by various
{\hyperref[\detokenize{Howto_output::doc}]{\sphinxcrossref{\DUrole{doc}{output commands}}}}.  The scalar and vector values calculated
by this fix are “intensive”.

\sphinxAtStartPar
No parameter of this fix can be used with the \sphinxstyleemphasis{start/stop} keywords of
the {\hyperref[\detokenize{run::doc}]{\sphinxcrossref{\DUrole{doc}{run}}}} command.  This fix is not invoked during
{\hyperref[\detokenize{minimize::doc}]{\sphinxcrossref{\DUrole{doc}{energy minimization}}}}.


\bigskip\hrule\bigskip



\subsection{Restrictions}
\label{\detokenize{fix_balance:restrictions}}
\sphinxAtStartPar
For 2d simulations, the \sphinxstyleemphasis{z} style cannot be used, nor can \sphinxstyleemphasis{z}
appear in \sphinxstyleemphasis{dimstr} for the \sphinxstyleemphasis{shift} style.

\sphinxAtStartPar
Balancing through recursive bisectioning (\sphinxstyleemphasis{rcb} style) requires
{\hyperref[\detokenize{comm_style::doc}]{\sphinxcrossref{\DUrole{doc}{comm\_style tiled}}}}.


\subsection{Related commands}
\label{\detokenize{fix_balance:related-commands}}
\sphinxAtStartPar
{\hyperref[\detokenize{group::doc}]{\sphinxcrossref{\DUrole{doc}{group}}}}, {\hyperref[\detokenize{processors::doc}]{\sphinxcrossref{\DUrole{doc}{processors}}}}, {\hyperref[\detokenize{balance::doc}]{\sphinxcrossref{\DUrole{doc}{balance}}}},
{\hyperref[\detokenize{comm_style::doc}]{\sphinxcrossref{\DUrole{doc}{comm\_style}}}}


\subsection{Default}
\label{\detokenize{fix_balance:default}}
\sphinxAtStartPar
The default setting is sort = no.

\sphinxstepscope

\index{fix bocs@\spxentry{fix bocs}}\ignorespaces 

\section{fix bocs command}
\label{\detokenize{fix_bocs:fix-bocs-command}}\label{\detokenize{fix_bocs:index-0}}\label{\detokenize{fix_bocs::doc}}

\subsection{Syntax}
\label{\detokenize{fix_bocs:syntax}}
\begin{sphinxVerbatim}[commandchars=\\\{\}]
\PYG{k}{fix }\PYG{n+nv+nvIdentifier}{ID}\PYG{+w}{ }\PYG{n+nv+nvIdentifier}{group\PYGZhy{}ID}\PYG{+w}{ }\PYG{n}{bocs}\PYG{+w}{ }\PYG{n}{keyword}\PYG{+w}{ }\PYG{n}{values}\PYG{+w}{ }\PYG{n}{...}
\end{sphinxVerbatim}
\begin{itemize}
\item {} 
\sphinxAtStartPar
ID, group\sphinxhyphen{}ID are documented in {\hyperref[\detokenize{fix::doc}]{\sphinxcrossref{\DUrole{doc}{fix}}}} command

\item {} 
\sphinxAtStartPar
bocs = style name of this fix command

\item {} 
\sphinxAtStartPar
two or more keyword/value pairs may be appended

\item {} 
\sphinxAtStartPar
keyword = \sphinxstyleemphasis{temp} or \sphinxstyleemphasis{cgiso} or \sphinxstyleemphasis{tchain} or \sphinxstyleemphasis{pchain} or \sphinxstyleemphasis{mtk} or \sphinxstyleemphasis{tloop} or \sphinxstyleemphasis{ploop}
\begin{sphinxalltt}
\sphinxstyleemphasis{temp} values = Tstart Tstop Tdamp
\sphinxstyleemphasis{cgiso} values = Pstart Pstop Pdamp basis\_set args
  basis\_set = \sphinxstyleemphasis{analytic} or \sphinxstyleemphasis{linear\_spline} or \sphinxstyleemphasis{cubic\_spline}
    \sphinxstyleemphasis{analytic} args = V\_avg N\_particles N\_coeff Coeff\_1 Coeff\_2 ... Coeff\_N
    \sphinxstyleemphasis{linear\_spline} args = input\_filename
    \sphinxstyleemphasis{cubic\_spline} args = input\_filename
\sphinxstyleemphasis{tchain} value = N = length of thermostat chain (1 = single thermostat)
\sphinxstyleemphasis{pchain} value = N = length of thermostat on barostat (0 = no thermostat)
\sphinxstyleemphasis{mtk} value = \sphinxstyleemphasis{yes} or \sphinxstyleemphasis{no} = add MTK adjustment term or not
\sphinxstyleemphasis{tloop} value = M = number of sub\sphinxhyphen{}cycles to perform on thermostat
\sphinxstyleemphasis{ploop} value = M = number of sub\sphinxhyphen{}cycles to perform on barostat
\end{sphinxalltt}

\end{itemize}


\subsection{Examples}
\label{\detokenize{fix_bocs:examples}}
\begin{sphinxVerbatim}[commandchars=\\\{\}]
\PYG{k}{fix }\PYG{n+nv+nvIdentifier}{1}\PYG{+w}{ }\PYG{n+nv+nvIdentifier}{all}\PYG{+w}{ }\PYG{n}{bocs}\PYG{+w}{ }\PYG{n}{temp}\PYG{+w}{ }\PYG{l+m}{300.0}\PYG{+w}{ }\PYG{l+m}{300.0}\PYG{+w}{ }\PYG{l+m}{100.0}\PYG{+w}{ }\PYG{n}{cgiso}\PYG{+w}{ }\PYG{l+m}{0.986}\PYG{+w}{ }\PYG{l+m}{0.986}\PYG{+w}{ }\PYG{l+m}{1000.0}\PYG{+w}{ }\PYG{n}{analytic}\PYG{+w}{ }\PYG{l+m}{66476.015}\PYG{+w}{ }\PYG{l+m}{968}\PYG{+w}{ }\PYG{l+m}{2}\PYG{+w}{ }\PYG{l+m}{245030.10}\PYG{+w}{ }\PYG{l+m}{8962.20}
\PYG{k}{fix }\PYG{n+nv+nvIdentifier}{1}\PYG{+w}{ }\PYG{n+nv+nvIdentifier}{all}\PYG{+w}{ }\PYG{n}{bocs}\PYG{+w}{ }\PYG{n}{temp}\PYG{+w}{ }\PYG{l+m}{300.0}\PYG{+w}{ }\PYG{l+m}{300.0}\PYG{+w}{ }\PYG{l+m}{100.0}\PYG{+w}{ }\PYG{n}{cgiso}\PYG{+w}{ }\PYG{l+m}{0.986}\PYG{+w}{ }\PYG{l+m}{0.986}\PYG{+w}{ }\PYG{l+m}{1000.0}\PYG{+w}{ }\PYG{n}{cubic\PYGZus{}spline}\PYG{+w}{ }\PYG{n}{input\PYGZus{}Fv.dat}
\PYG{k}{thermo\PYGZus{}modify}\PYG{+w}{ }\PYG{n}{press}\PYG{+w}{ }\PYG{l+m}{1}\PYG{n}{\PYGZus{}press}
\end{sphinxVerbatim}


\subsection{Description}
\label{\detokenize{fix_bocs:description}}
\sphinxAtStartPar
These commands incorporate a pressure correction as described by
Dunn and Noid {\hyperref[\detokenize{fix_bocs:bocs-dunn1}]{\sphinxcrossref{\DUrole{std,std-ref}{(Dunn1)}}}} to the standard MTK
barostat by Martyna et al. {\hyperref[\detokenize{fix_bocs:bocs-martyna}]{\sphinxcrossref{\DUrole{std,std-ref}{(Martyna)}}}}.
The first half of the command mimics a standard {\hyperref[\detokenize{fix_nh::doc}]{\sphinxcrossref{\DUrole{doc}{fix npt}}}}
command:

\begin{sphinxVerbatim}[commandchars=\\\{\}]
\PYG{k}{fix }\PYG{n+nv+nvIdentifier}{1}\PYG{+w}{ }\PYG{n+nv+nvIdentifier}{all}\PYG{+w}{ }\PYG{n}{bocs}\PYG{+w}{ }\PYG{n}{temp}\PYG{+w}{ }\PYG{n}{Tstart}\PYG{+w}{ }\PYG{n}{Tstop}\PYG{+w}{ }\PYG{n}{Tcoupl}\PYG{+w}{ }\PYG{n}{cgiso}\PYG{+w}{ }\PYG{n}{Pstart}\PYG{+w}{ }\PYG{n}{Pstop}\PYG{+w}{ }\PYG{n}{Pdamp}
\end{sphinxVerbatim}

\sphinxAtStartPar
The two differences are replacing \sphinxstyleemphasis{npt} with \sphinxstyleemphasis{bocs}, and replacing
\sphinxstyleemphasis{iso}/\sphinxstyleemphasis{aniso}/etc. with \sphinxstyleemphasis{cgiso}.
The rest of the command details what form you would like to use for
the pressure correction equation. The choices are: \sphinxstyleemphasis{analytic}, \sphinxstyleemphasis{linear\_spline},
or \sphinxstyleemphasis{cubic\_spline}.

\sphinxAtStartPar
With either spline method, the only argument that needs to follow it
is the name of a file that contains the desired pressure correction
as a function of volume. The file must be formatted so each line has:

\begin{sphinxVerbatim}[commandchars=\\\{\}]
Volume\PYGZus{}i, PressureCorrection\PYGZus{}i
\end{sphinxVerbatim}

\sphinxAtStartPar
Note both the COMMA and the SPACE separating the volume’s
value and its corresponding pressure correction. The volumes in the file
must be uniformly spaced. Both the volumes and the pressure corrections
should be provided in the proper units (e.g., if you are using \sphinxstyleemphasis{units real},
the volumes should all be in \(\mathrm{\mathring{A}}^3\) and the pressure
corrections should all be in atm). Furthermore, the table should start/end at a
volume considerably smaller/larger than you expect your system to sample
during the simulation. If the system ever reaches a volume outside of the
range provided, the simulation will stop.

\sphinxAtStartPar
With the \sphinxstyleemphasis{analytic} option, the arguments are as follows:

\begin{sphinxVerbatim}[commandchars=\\\{\}]
... analytic V\PYGZus{}avg N\PYGZus{}particles N\PYGZus{}coeff Coeff\PYGZus{}1 Coeff\PYGZus{}2 ... Coeff\PYGZus{}N
\end{sphinxVerbatim}

\sphinxAtStartPar
Note that \sphinxstyleemphasis{V\_avg} and \sphinxstyleemphasis{Coeff\_i} should all be in the proper units (e.g., if you
are using \sphinxstyleemphasis{units real}, \sphinxstyleemphasis{V\_avg} should be in \(\mathrm{\mathring{A}^3}\)
and the coefficients should all be in
\(\mathrm{atm}\cdot\mathrm{\mathring{A}^3}\)).


\bigskip\hrule\bigskip



\subsection{Restart, fix\_modify, output, run start/stop, minimize info}
\label{\detokenize{fix_bocs:restart-fix-modify-output-run-start-stop-minimize-info}}
\sphinxAtStartPar
This fix writes the cumulative global energy change to {\hyperref[\detokenize{restart::doc}]{\sphinxcrossref{\DUrole{doc}{binary
restart files}}}}.  See the {\hyperref[\detokenize{read_restart::doc}]{\sphinxcrossref{\DUrole{doc}{read\_restart}}}}
command for info on how to re\sphinxhyphen{}specify a fix in an input script that
reads a restart file, so that the fix continues in an uninterrupted
fashion.

\sphinxAtStartPar
The {\hyperref[\detokenize{fix_modify::doc}]{\sphinxcrossref{\DUrole{doc}{fix\_modify}}}} \sphinxstyleemphasis{temp} option is supported by this
fix.  You can use it to assign a temperature {\hyperref[\detokenize{compute::doc}]{\sphinxcrossref{\DUrole{doc}{compute}}}}
you have defined to this fix which will be used in its thermostatting
procedure, as described above.  For consistency, the group used by
this fix and by the compute should be the same.

\sphinxAtStartPar
The cumulative energy change in the system imposed by this fix is
included in the {\hyperref[\detokenize{thermo_style::doc}]{\sphinxcrossref{\DUrole{doc}{thermodynamic output}}}} keywords
\sphinxstyleemphasis{ecouple} and \sphinxstyleemphasis{econserve}.  See the {\hyperref[\detokenize{thermo_style::doc}]{\sphinxcrossref{\DUrole{doc}{thermo\_style}}}}
doc page for details.

\sphinxAtStartPar
This fix computes a global scalar which can be accessed by various
{\hyperref[\detokenize{Howto_output::doc}]{\sphinxcrossref{\DUrole{doc}{output commands}}}}.  The scalar is the same
cumulative energy change due to this fix described in the previous
paragraph.  The scalar value calculated by this fix is “extensive”.

\sphinxAtStartPar
This fix can ramp its target temperature over multiple runs, using the
\sphinxstyleemphasis{start} and \sphinxstyleemphasis{stop} keywords of the {\hyperref[\detokenize{run::doc}]{\sphinxcrossref{\DUrole{doc}{run}}}} command.  See the
{\hyperref[\detokenize{run::doc}]{\sphinxcrossref{\DUrole{doc}{run}}}} command for details of how to do this.

\sphinxAtStartPar
This fix is not invoked during {\hyperref[\detokenize{minimize::doc}]{\sphinxcrossref{\DUrole{doc}{energy minimization}}}}.


\subsection{Restrictions}
\label{\detokenize{fix_bocs:restrictions}}
\sphinxAtStartPar
This fix is part of the BOCS package.  It is only enabled if
LAMMPS was built with that package.  See the
{\hyperref[\detokenize{Build_package::doc}]{\sphinxcrossref{\DUrole{doc}{Build package}}}} page for more info.

\sphinxAtStartPar
As this is computing a (modified) pressure, group\sphinxhyphen{}ID should be \sphinxstyleemphasis{all}.

\sphinxAtStartPar
The pressure correction has only been tested for use with an isotropic
pressure coupling in 3 dimensions.

\sphinxAtStartPar
By default, LAMMPS will still report the normal value for the pressure
if the pressure is printed via a \sphinxstyleemphasis{thermo} command, or if the pressures
are written to a file every so often. In order to have LAMMPS report the
modified pressure, you must include the \sphinxstyleemphasis{thermo\_modify} command given in
the examples. For the last argument in the command, you should put
XXXX\_press, where XXXX is the ID given to the fix bocs command (in the
example, the ID of the fix bocs command is 1).


\subsection{Further information}
\label{\detokenize{fix_bocs:further-information}}
\sphinxAtStartPar
For more details about the pressure correction and the entire BOCS software
package, visit the \sphinxhref{https://github.com/noid-group/BOCS}{BOCS package on GitHub} and read the release
paper by Dunn et al. {\hyperref[\detokenize{fix_bocs:bocs-dunn2}]{\sphinxcrossref{\DUrole{std,std-ref}{(Dunn2)}}}} .


\bigskip\hrule\bigskip

\phantomsection\label{\detokenize{fix_bocs:bocs-dunn1}}
\sphinxAtStartPar
\sphinxstylestrong{(Dunn1)} Dunn and Noid, J Chem Phys, 143, 243148 (2015).

\phantomsection\label{\detokenize{fix_bocs:bocs-martyna}}
\sphinxAtStartPar
\sphinxstylestrong{(Martyna)} Martyna, Tobias, and Klein, J Chem Phys, 101, 4177 (1994).

\phantomsection\label{\detokenize{fix_bocs:bocs-dunn2}}
\sphinxAtStartPar
\sphinxstylestrong{(Dunn2)} Dunn, Lebold, DeLyser, Rudzinski, and Noid, J. Phys. Chem. B, 122, 3363 (2018).

\sphinxstepscope

\index{fix bond/break@\spxentry{fix bond/break}}\ignorespaces 

\section{fix bond/break command}
\label{\detokenize{fix_bond_break:fix-bond-break-command}}\label{\detokenize{fix_bond_break:index-0}}\label{\detokenize{fix_bond_break::doc}}

\subsection{Syntax}
\label{\detokenize{fix_bond_break:syntax}}
\begin{sphinxVerbatim}[commandchars=\\\{\}]
\PYG{k}{fix }\PYG{n+nv+nvIdentifier}{ID}\PYG{+w}{ }\PYG{n+nv+nvIdentifier}{group\PYGZhy{}ID}\PYG{+w}{ }\PYG{n}{bond}\PYG{o}{/}\PYG{n}{break}\PYG{+w}{ }\PYG{n}{Nevery}\PYG{+w}{ }\PYG{n}{bondtype}\PYG{+w}{ }\PYG{n}{Rmax}\PYG{+w}{ }\PYG{n}{keyword}\PYG{+w}{ }\PYG{n}{values}\PYG{+w}{ }\PYG{n}{...}
\end{sphinxVerbatim}
\begin{itemize}
\item {} 
\sphinxAtStartPar
ID, group\sphinxhyphen{}ID are documented in {\hyperref[\detokenize{fix::doc}]{\sphinxcrossref{\DUrole{doc}{fix}}}} command

\item {} 
\sphinxAtStartPar
bond/break = style name of this fix command

\item {} 
\sphinxAtStartPar
Nevery = attempt bond breaking every this many steps

\item {} 
\sphinxAtStartPar
bondtype = type of bonds to break (integer or type label)

\item {} 
\sphinxAtStartPar
Rmax = bond longer than Rmax can break (distance units)

\item {} 
\sphinxAtStartPar
zero or more keyword/value pairs may be appended

\item {} 
\sphinxAtStartPar
keyword = \sphinxstyleemphasis{prob}
\begin{sphinxalltt}
\sphinxstyleemphasis{prob} values = fraction seed
  fraction = break a bond with this probability if otherwise eligible
  seed = random number seed (positive integer)
\end{sphinxalltt}

\end{itemize}


\subsection{Examples}
\label{\detokenize{fix_bond_break:examples}}
\begin{sphinxVerbatim}[commandchars=\\\{\}]
\PYG{k}{fix }\PYG{n+nv+nvIdentifier}{5}\PYG{+w}{ }\PYG{n+nv+nvIdentifier}{all}\PYG{+w}{ }\PYG{n}{bond}\PYG{o}{/}\PYG{n}{break}\PYG{+w}{ }\PYG{l+m}{10}\PYG{+w}{ }\PYG{l+m}{2}\PYG{+w}{ }\PYG{l+m}{1.2}
\PYG{k}{fix }\PYG{n+nv+nvIdentifier}{5}\PYG{+w}{ }\PYG{n+nv+nvIdentifier}{polymer}\PYG{+w}{ }\PYG{n}{bond}\PYG{o}{/}\PYG{n}{break}\PYG{+w}{ }\PYG{l+m}{1}\PYG{+w}{ }\PYG{l+m}{1}\PYG{+w}{ }\PYG{l+m}{2.0}\PYG{+w}{ }\PYG{n}{prob}\PYG{+w}{ }\PYG{l+m}{0.5}\PYG{+w}{ }\PYG{l+m}{49829}
\end{sphinxVerbatim}


\subsection{Description}
\label{\detokenize{fix_bond_break:description}}
\sphinxAtStartPar
Break bonds between pairs of atoms as a simulation runs according to
specified criteria.  This can be used to model the dissolution of a
polymer network due to stretching of the simulation box or other
deformations.  In this context, a bond means an interaction between a
pair of atoms computed by the {\hyperref[\detokenize{bond_style::doc}]{\sphinxcrossref{\DUrole{doc}{bond\_style}}}} command.
Once the bond is broken it will be permanently deleted, as will all
angle, dihedral, and improper interactions that bond is part of.

\sphinxAtStartPar
This is different than a {\hyperref[\detokenize{pair_style::doc}]{\sphinxcrossref{\DUrole{doc}{pair\sphinxhyphen{}wise}}}} bond\sphinxhyphen{}order
potential such as Tersoff or AIREBO which infers bonds and many\sphinxhyphen{}body
interactions based on the current geometry of a small cluster of atoms
and effectively creates and destroys bonds and higher\sphinxhyphen{}order many\sphinxhyphen{}body
interactions from timestep to timestep as atoms move.

\sphinxAtStartPar
A check for possible bond breakage is performed every \sphinxstyleemphasis{Nevery}
timesteps.  If two bonded atoms \(i\) and \(j\) are farther than the
distance \sphinxstyleemphasis{Rmax} from each other, the bond is of type \sphinxstyleemphasis{bondtype}, and both
\(i\) and \(j\) are in the specified fix group, then the bond between
\(i\) and \(j\) is labeled as a “possible” bond to break.

\sphinxAtStartPar
If several bonds involving an atom are stretched, it may have multiple
possible bonds to break.  Every atom checks its list of possible bonds
to break and labels the longest such bond as its “sole” bond to break.
After this is done, if atom \(i\) is bonded to atom \(j\) in its sole
bond, and atom \(j\) is bonded to atom \(j\) in its sole bond, then the
bond between \(i\) and \(j\) is “eligible” to be broken.

\sphinxAtStartPar
Note that these rules mean an atom will only be part of at most one
broken bond on a given time step.  It also means that if atom \(i\) chooses
atom \(j\) as its sole partner, but atom \(j\) chooses atom \(k\)
as its sole partner (because \(R_{jk} > R_{ij}\)), then this means atom
\(i\) will not be part of a broken bond on this time step, even if it has
other possible bond partners.

\sphinxAtStartPar
The \sphinxstyleemphasis{prob} keyword can effect whether an eligible bond is actually
broken.  The \sphinxstyleemphasis{fraction} setting must be a value between 0.0 and 1.0.
A uniform random number between 0.0 and 1.0 is generated and the
eligible bond is only broken if the random number is less than \sphinxstyleemphasis{fraction}.

\sphinxAtStartPar
When a bond is broken, data structures within LAMMPS that store bond
topologies are updated to reflect the breakage.  Likewise, if the bond
is part of a 3\sphinxhyphen{}body (angle) or 4\sphinxhyphen{}body (dihedral, improper)
interaction, that interaction is removed as well.  These changes
typically affect pair\sphinxhyphen{}wise interactions between atoms that used to be
part of bonds, angles, etc.

\begin{sphinxadmonition}{note}{Note:}
\sphinxAtStartPar
One data structure that is not updated when a bond breaks are
the molecule IDs stored by each atom.  Even though one molecule
becomes two molecules due to the broken bond, all atoms in both new
molecules retain their original molecule IDs.
\end{sphinxadmonition}

\sphinxAtStartPar
Computationally, each time step this fix is invoked, it loops over all
the bonds in the system and computes distances between pairs of bonded
atoms.  It also communicates between neighboring processors to
coordinate which bonds are broken.  Moreover, if any bonds are broken,
neighbor lists must be immediately updated on the same time step.  This
is to ensure that any pair\sphinxhyphen{}wise interactions that should be turned “on”
due to a bond breaking, because they are no longer excluded by the
presence of the bond and the settings of the
{\hyperref[\detokenize{special_bonds::doc}]{\sphinxcrossref{\DUrole{doc}{special\_bonds}}}} command, will be immediately
recognized.  All of these operations increase the cost of a time step.
Thus, you should be cautious about invoking this fix too frequently.

\sphinxAtStartPar
You can dump out snapshots of the current bond topology via the {\hyperref[\detokenize{dump::doc}]{\sphinxcrossref{\DUrole{doc}{dump local}}}} command.

\begin{sphinxadmonition}{note}{Note:}
\sphinxAtStartPar
Breaking a bond typically alters the energy of a system.  You
should be careful not to choose bond breaking criteria that induce a
dramatic change in energy.  For example, if you define a very stiff
harmonic bond and break it when two atoms are separated by a distance
far from the equilibrium bond length, then the two atoms will be
dramatically released when the bond is broken.  More generally, you
may need to thermostat your system to compensate for energy changes
resulting from broken bonds (as well as angles, dihedrals, and impropers).
\end{sphinxadmonition}

\sphinxAtStartPar
See the {\hyperref[\detokenize{Howto_broken_bonds::doc}]{\sphinxcrossref{\DUrole{doc}{Howto}}}} page on broken bonds for more
information on related features in LAMMPS.


\bigskip\hrule\bigskip



\subsection{Restart, fix\_modify, output, run start/stop, minimize info}
\label{\detokenize{fix_bond_break:restart-fix-modify-output-run-start-stop-minimize-info}}
\sphinxAtStartPar
No information about this fix is written to {\hyperref[\detokenize{restart::doc}]{\sphinxcrossref{\DUrole{doc}{binary restart files}}}}.  None of the {\hyperref[\detokenize{fix_modify::doc}]{\sphinxcrossref{\DUrole{doc}{fix\_modify}}}} options
are relevant to this fix.

\sphinxAtStartPar
This fix computes two statistics, which it stores in a global vector of
length 2. This vector can be accessed by various {\hyperref[\detokenize{Howto_output::doc}]{\sphinxcrossref{\DUrole{doc}{output commands}}}}.  The vector values calculated by this fix are “intensive”.

\sphinxAtStartPar
The two quantities in the global vector are
\begin{enumerate}
\sphinxsetlistlabels{\arabic}{enumi}{enumii}{(}{)}%
\item {} 
\sphinxAtStartPar
number of bonds broken on the most recent breakage time step

\item {} 
\sphinxAtStartPar
cumulative number of bonds broken

\end{enumerate}

\sphinxAtStartPar
No parameter of this fix can be used with the \sphinxstyleemphasis{start/stop} keywords of
the {\hyperref[\detokenize{run::doc}]{\sphinxcrossref{\DUrole{doc}{run}}}} command.  This fix is not invoked during {\hyperref[\detokenize{minimize::doc}]{\sphinxcrossref{\DUrole{doc}{energy minimization}}}}.


\subsection{Restrictions}
\label{\detokenize{fix_bond_break:restrictions}}
\sphinxAtStartPar
This fix is part of the MC package.  It is only enabled if LAMMPS was
built with that package.  See the {\hyperref[\detokenize{Build_package::doc}]{\sphinxcrossref{\DUrole{doc}{Build package}}}}
doc page for more info.


\subsection{Related commands}
\label{\detokenize{fix_bond_break:related-commands}}
\sphinxAtStartPar
{\hyperref[\detokenize{fix_bond_create::doc}]{\sphinxcrossref{\DUrole{doc}{fix bond/create}}}}, {\hyperref[\detokenize{fix_bond_react::doc}]{\sphinxcrossref{\DUrole{doc}{fix bond/react}}}}, {\hyperref[\detokenize{fix_bond_swap::doc}]{\sphinxcrossref{\DUrole{doc}{fix bond/swap}}}},
{\hyperref[\detokenize{dump::doc}]{\sphinxcrossref{\DUrole{doc}{dump local}}}}, {\hyperref[\detokenize{special_bonds::doc}]{\sphinxcrossref{\DUrole{doc}{special\_bonds}}}}


\subsection{Default}
\label{\detokenize{fix_bond_break:default}}
\sphinxAtStartPar
The option defaults are prob = 1.0.

\sphinxstepscope

\index{fix bond/create@\spxentry{fix bond/create}}\ignorespaces 
\index{fix bond/create/angle@\spxentry{fix bond/create/angle}}\ignorespaces 

\section{fix bond/create command}
\label{\detokenize{fix_bond_create:fix-bond-create-command}}\label{\detokenize{fix_bond_create:index-1}}\label{\detokenize{fix_bond_create:index-0}}\label{\detokenize{fix_bond_create::doc}}

\section{fix bond/create/angle command}
\label{\detokenize{fix_bond_create:fix-bond-create-angle-command}}

\subsection{Syntax}
\label{\detokenize{fix_bond_create:syntax}}
\begin{sphinxVerbatim}[commandchars=\\\{\}]
\PYG{k}{fix }\PYG{n+nv+nvIdentifier}{ID}\PYG{+w}{ }\PYG{n+nv+nvIdentifier}{group\PYGZhy{}ID}\PYG{+w}{ }\PYG{n}{style}\PYG{+w}{ }\PYG{n}{Nevery}\PYG{+w}{ }\PYG{n}{itype}\PYG{+w}{ }\PYG{n}{jtype}\PYG{+w}{ }\PYG{n}{Rmin}\PYG{+w}{ }\PYG{n}{bondtype}\PYG{+w}{ }\PYG{n}{keyword}\PYG{+w}{ }\PYG{n}{values}\PYG{+w}{ }\PYG{n}{...}
\end{sphinxVerbatim}
\begin{itemize}
\item {} 
\sphinxAtStartPar
ID, group\sphinxhyphen{}ID are documented in {\hyperref[\detokenize{fix::doc}]{\sphinxcrossref{\DUrole{doc}{fix}}}} command

\item {} 
\sphinxAtStartPar
style = \sphinxstyleemphasis{bond/create} or \sphinxstyleemphasis{bond/create/angle}

\item {} 
\sphinxAtStartPar
Nevery = attempt bond creation every this many steps

\item {} 
\sphinxAtStartPar
itype,jtype = atoms of itype can bond to atoms of jtype (1\sphinxhyphen{}Ntypes or type label)

\item {} 
\sphinxAtStartPar
Rmin = two atoms separated by less than Rmin can bond (distance units)

\item {} 
\sphinxAtStartPar
bondtype = type of created bonds (integer or type label)

\item {} 
\sphinxAtStartPar
zero or more keyword/value pairs may be appended to args

\item {} 
\sphinxAtStartPar
keyword = \sphinxstyleemphasis{iparam} or \sphinxstyleemphasis{jparam} or \sphinxstyleemphasis{prob} or \sphinxstyleemphasis{atype} or \sphinxstyleemphasis{dtype} or \sphinxstyleemphasis{itype} or \sphinxstyleemphasis{aconstrain}
\begin{sphinxalltt}
\sphinxstyleemphasis{iparam} values = maxbond, newtype
  maxbond = max \# of bonds of bondtype the itype atom can have
  newtype = change the itype atom to this type when maxbonds exist (1\sphinxhyphen{}Ntypes or type label)
\sphinxstyleemphasis{jparam} values = maxbond, newtype
  maxbond = max \# of bonds of bondtype the jtype atom can have
  newtype = change the jtype atom to this type when maxbonds exist (1\sphinxhyphen{}Ntypes or type label)
\sphinxstyleemphasis{prob} values = fraction seed
  fraction = create a bond with this probability if otherwise eligible
  seed = random number seed (positive integer)
\sphinxstyleemphasis{atype} value = angletype
  angletype = type of created angles (integer or type label)
\sphinxstyleemphasis{dtype} value = dihedraltype
  dihedraltype = type of created dihedrals (integer or type label)
\sphinxstyleemphasis{itype} value = impropertype
  impropertype = type of created impropers (integer or type label)
\sphinxstyleemphasis{aconstrain} value = amin amax
  amin = minimal angle at which new bonds can be created
  amax = maximal angle at which new bonds can be created
\end{sphinxalltt}

\end{itemize}


\subsection{Examples}
\label{\detokenize{fix_bond_create:examples}}
\begin{sphinxVerbatim}[commandchars=\\\{\}]
\PYG{k}{fix }\PYG{n+nv+nvIdentifier}{5}\PYG{+w}{ }\PYG{n+nv+nvIdentifier}{all}\PYG{+w}{ }\PYG{n}{bond}\PYG{o}{/}\PYG{n}{create}\PYG{+w}{ }\PYG{l+m}{10}\PYG{+w}{ }\PYG{l+m}{1}\PYG{+w}{ }\PYG{l+m}{2}\PYG{+w}{ }\PYG{l+m}{0.8}\PYG{+w}{ }\PYG{l+m}{1}
\PYG{k}{fix }\PYG{n+nv+nvIdentifier}{5}\PYG{+w}{ }\PYG{n+nv+nvIdentifier}{all}\PYG{+w}{ }\PYG{n}{bond}\PYG{o}{/}\PYG{n}{create}\PYG{+w}{ }\PYG{l+m}{1}\PYG{+w}{ }\PYG{l+m}{3}\PYG{+w}{ }\PYG{l+m}{3}\PYG{+w}{ }\PYG{l+m}{0.8}\PYG{+w}{ }\PYG{l+m}{1}\PYG{+w}{ }\PYG{n}{prob}\PYG{+w}{ }\PYG{l+m}{0.5}\PYG{+w}{ }\PYG{l+m}{85784}\PYG{+w}{ }\PYG{n}{iparam}\PYG{+w}{ }\PYG{l+m}{2}\PYG{+w}{ }\PYG{l+m}{3}
\PYG{k}{fix }\PYG{n+nv+nvIdentifier}{5}\PYG{+w}{ }\PYG{n+nv+nvIdentifier}{all}\PYG{+w}{ }\PYG{n}{bond}\PYG{o}{/}\PYG{n}{create}\PYG{+w}{ }\PYG{l+m}{1}\PYG{+w}{ }\PYG{l+m}{3}\PYG{+w}{ }\PYG{l+m}{3}\PYG{+w}{ }\PYG{l+m}{0.8}\PYG{+w}{ }\PYG{l+m}{1}\PYG{+w}{ }\PYG{n}{prob}\PYG{+w}{ }\PYG{l+m}{0.5}\PYG{+w}{ }\PYG{l+m}{85784}\PYG{+w}{ }\PYG{n}{iparam}\PYG{+w}{ }\PYG{l+m}{2}\PYG{+w}{ }\PYG{l+m}{3}\PYG{+w}{ }\PYG{n}{atype}\PYG{+w}{ }\PYG{l+m}{1}\PYG{+w}{ }\PYG{n}{dtype}\PYG{+w}{ }\PYG{l+m}{2}
\PYG{k}{fix }\PYG{n+nv+nvIdentifier}{5}\PYG{+w}{ }\PYG{n+nv+nvIdentifier}{all}\PYG{+w}{ }\PYG{n}{bond}\PYG{o}{/}\PYG{n}{create}\PYG{o}{/}\PYG{n}{angle}\PYG{+w}{ }\PYG{l+m}{10}\PYG{+w}{ }\PYG{l+m}{1}\PYG{+w}{ }\PYG{l+m}{2}\PYG{+w}{ }\PYG{l+m}{1.122}\PYG{+w}{ }\PYG{l+m}{1}\PYG{+w}{ }\PYG{n}{aconstrain}\PYG{+w}{ }\PYG{l+m}{120}\PYG{+w}{ }\PYG{l+m}{180}\PYG{+w}{ }\PYG{n}{prob}\PYG{+w}{ }\PYG{l+m}{1}\PYG{+w}{ }\PYG{l+m}{4928459}\PYG{+w}{ }\PYG{n}{iparam}\PYG{+w}{ }\PYG{l+m}{2}\PYG{+w}{ }\PYG{l+m}{1}\PYG{+w}{ }\PYG{n}{jparam}\PYG{+w}{ }\PYG{l+m}{2}\PYG{+w}{ }\PYG{l+m}{2}

\PYG{k}{labelmap}\PYG{+w}{ }\PYG{n}{atom}\PYG{+w}{ }\PYG{l+m}{1}\PYG{+w}{ }\PYG{n}{c1}\PYG{+w}{ }\PYG{l+m}{2}\PYG{+w}{ }\PYG{n}{n2}
\PYG{k}{labelmap}\PYG{+w}{ }\PYG{n}{bond}\PYG{+w}{ }\PYG{l+m}{1}\PYG{+w}{ }\PYG{n}{c1}\PYG{o}{\PYGZhy{}}\PYG{n}{n2}
\PYG{k}{fix }\PYG{n+nv+nvIdentifier}{5}\PYG{+w}{ }\PYG{n+nv+nvIdentifier}{all}\PYG{+w}{ }\PYG{n}{bond}\PYG{o}{/}\PYG{n}{create}\PYG{+w}{ }\PYG{l+m}{10}\PYG{+w}{ }\PYG{n}{c1}\PYG{+w}{ }\PYG{n}{n2}\PYG{+w}{ }\PYG{l+m}{0.8}\PYG{+w}{ }\PYG{n}{c1}\PYG{o}{\PYGZhy{}}\PYG{n}{n2}
\end{sphinxVerbatim}


\subsection{Description}
\label{\detokenize{fix_bond_create:description}}
\sphinxAtStartPar
Create bonds between pairs of atoms as a simulation runs according to
specified criteria.  This can be used to model the cross\sphinxhyphen{}linking of
polymers, the formation of a percolation network, etc.  In this
context, a bond means an interaction between a pair of atoms computed
by the {\hyperref[\detokenize{bond_style::doc}]{\sphinxcrossref{\DUrole{doc}{bond\_style}}}} command.  Once the bond is created
it will be permanently in place.  Optionally, the creation of a bond
can also create angle, dihedral, and improper interactions that the bond
is part of.  See the discussion of the \sphinxstyleemphasis{atype}, \sphinxstyleemphasis{dtype}, and \sphinxstyleemphasis{itype}
keywords below.

\sphinxAtStartPar
This process is different than a {\hyperref[\detokenize{pair_style::doc}]{\sphinxcrossref{\DUrole{doc}{pair\sphinxhyphen{}wise}}}} bond\sphinxhyphen{}order
potential such as Tersoff or AIREBO, which infer bonds and many\sphinxhyphen{}body
interactions based on the current geometry of a small cluster of atoms
and effectively create and destroy bonds and higher\sphinxhyphen{}order many\sphinxhyphen{}body
interactions from time step to time step as the atoms move.

\sphinxAtStartPar
A check for possible new bonds is performed every \sphinxstyleemphasis{Nevery} time steps.
If two atoms \(i\) and \(j\) are within a distance \sphinxstyleemphasis{Rmin} of each
other, atom \(i\) is of type \sphinxstyleemphasis{itype}, atom \(j\) is of type \sphinxstyleemphasis{jtype},
and both \(i\) and \(j\) are in the specified fix group, then if a bond
does not already exist between atoms \(i\) and \(j\), and if both
\(i\) and \(j\) meet their respective \sphinxstyleemphasis{maxbond} requirements (explained
below), then \(i\) and \(j\) are labeled as a “possible” bond pair.

\sphinxAtStartPar
If several atoms are close to an atom, it may have multiple possible
bond partners.  Every atom checks its list of possible bond partners
and labels the closest such partner as its “sole” bond partner.  After
this is done, if atom \(i\) has atom \(j\) as its sole partner and
atom \(j\) has atom \(i\) as its sole partner, then the
\(i,j\) bond is “eligible” to be formed.

\sphinxAtStartPar
Note that these rules mean that an atom will only be part of at most one
created bond on a given time step.  It also means that if atom \(i\)
chooses atom \(j\) as its sole partner, but atom \(j\) chooses atom
\(k\) as its sole partner (because \(R_{jk} < R_{ij}\)), then atom
\(i\) will not form a bond on this time step, even if it has other possible
bond partners.

\sphinxAtStartPar
It is permissible to have \sphinxstyleemphasis{itype} = \sphinxstyleemphasis{jtype}.  \sphinxstyleemphasis{Rmin} must be \(\leq\) the
pair\sphinxhyphen{}wise cutoff distance between \sphinxstyleemphasis{itype} and \sphinxstyleemphasis{jtype} atoms, as defined
by the {\hyperref[\detokenize{pair_style::doc}]{\sphinxcrossref{\DUrole{doc}{pair\_style}}}} command.

\sphinxAtStartPar
The \sphinxstyleemphasis{iparam} and \sphinxstyleemphasis{jparam} keywords can be used to limit the bonding
functionality of the participating atoms.  Each atom keeps track of
how many bonds of \sphinxstyleemphasis{bondtype} it already has.  If atom \(i\) of type
\sphinxstyleemphasis{itype} already has \sphinxstyleemphasis{maxbond} bonds (as set by the \sphinxstyleemphasis{iparam}
keyword), then it will not form any more, and likewise for atom \(j\).
If \sphinxstyleemphasis{maxbond} is set to 0, then there is no limit on the number of bonds
that can be formed with that atom.

\sphinxAtStartPar
The \sphinxstyleemphasis{newtype} value for \sphinxstyleemphasis{iparam} and \sphinxstyleemphasis{jparam} can be used to change
the atom type of atom \(i\) or \(j\) when it reaches \sphinxstyleemphasis{maxbond} number
of bonds of type \sphinxstyleemphasis{bondtype}.  This means it can now interact in a pair\sphinxhyphen{}wise
fashion with other atoms in a different way by specifying different
{\hyperref[\detokenize{pair_coeff::doc}]{\sphinxcrossref{\DUrole{doc}{pair\_coeff}}}} coefficients.  If you do not wish the
atom type to change, simply specify \sphinxstyleemphasis{newtype} as \sphinxstyleemphasis{itype} or \sphinxstyleemphasis{jtype}.

\sphinxAtStartPar
The \sphinxstyleemphasis{prob} keyword can also affect whether an eligible bond is
actually created.  The \sphinxstyleemphasis{fraction} setting must be a value between 0.0
and 1.0.  A uniform random number between 0.0 and 1.0 is generated and
the eligible bond is only created if the random number is less than \sphinxstyleemphasis{fraction}.

\sphinxAtStartPar
The \sphinxstyleemphasis{aconstrain} keyword is only available with the fix
bond/create/angle command.  It allows one to specify minimum and maximum
angles \sphinxstyleemphasis{amin} and \sphinxstyleemphasis{amax}, respectively, between the two prospective bonding
partners and a third particle that is already bonded to one of the two
partners. Such a criterion can be important when new angles are defined
together with the formation of a new bond.  Without a restriction on the
permissible angle, and for stiffer angle potentials, very large energies
can arise and lead to unphysical behavior.

\sphinxAtStartPar
Any bond that is created is assigned a bond type of \sphinxstyleemphasis{bondtype}.

\sphinxAtStartPar
When a bond is created, data structures within LAMMPS that store bond
topologies are updated to reflect the creation.  If the bond is part of
new 3\sphinxhyphen{}body (angle) or 4\sphinxhyphen{}body (dihedral, improper) interactions, you
can choose to create new angles, dihedrals, and impropers as well using
the \sphinxstyleemphasis{atype}, \sphinxstyleemphasis{dtype}, and \sphinxstyleemphasis{itype} keywords.  All of these changes
typically affect pair\sphinxhyphen{}wise interactions between atoms that are now part
of new bonds, angles, etc.

\begin{sphinxadmonition}{note}{Note:}
\sphinxAtStartPar
One data structure that is not updated when a bond breaks are
the molecule IDs stored by each atom.  Even though two molecules
become one molecule due to the created bond, all atoms in the new
molecule retain their original molecule IDs.
\end{sphinxadmonition}

\sphinxAtStartPar
If the \sphinxstyleemphasis{atype} keyword is used and if an angle potential is defined
via the {\hyperref[\detokenize{angle_style::doc}]{\sphinxcrossref{\DUrole{doc}{angle\_style}}}} command, then any new 3\sphinxhyphen{}body
interactions inferred by the creation of a bond will create new angles
of type \sphinxstyleemphasis{angletype}, with parameters assigned by the corresponding
{\hyperref[\detokenize{angle_coeff::doc}]{\sphinxcrossref{\DUrole{doc}{angle\_coeff}}}} command.  Likewise, the \sphinxstyleemphasis{dtype} and
\sphinxstyleemphasis{itype} keywords will create new dihedrals and impropers of type
\sphinxstyleemphasis{dihedraltype} and \sphinxstyleemphasis{impropertype}.

\begin{sphinxadmonition}{note}{Note:}
\sphinxAtStartPar
To create a new bond, the internal LAMMPS data structures that
store this information must have space for it.  When LAMMPS is
initialized from a data file, the list of bonds is scanned and the
maximum number of bonds per atom is tallied.  If some atom will
acquire more bonds than this limit as this fix operates, then the
“extra bond per atom” parameter must be set to allow for it.  Ditto
for “extra angle per atom”, “extra dihedral per atom”, and “extra
improper per atom” if angles, dihedrals, or impropers are being added
when bonds are created.  See the {\hyperref[\detokenize{read_data::doc}]{\sphinxcrossref{\DUrole{doc}{read\_data}}}} or
{\hyperref[\detokenize{create_box::doc}]{\sphinxcrossref{\DUrole{doc}{create\_box}}}} command for more details.  Note that a
data file with no atoms can be used if you wish to add non\sphinxhyphen{}bonded
atoms via the {\hyperref[\detokenize{create_atoms::doc}]{\sphinxcrossref{\DUrole{doc}{create atoms}}}} command (e.g., for a
percolation simulation).
\end{sphinxadmonition}

\begin{sphinxadmonition}{note}{Note:}
\sphinxAtStartPar
LAMMPS stores and maintains a data structure with a list of the
first, second, and third neighbors of each atom (within the bond topology of
the system) for use in weighting pair\sphinxhyphen{}wise interactions for bonded
atoms.  Note that adding a single bond always adds a new first neighbor
but may also induce \sphinxstylestrong{many} new second and third neighbors, depending on the
molecular topology of your system.  The “extra special per atom”
parameter must typically be set to allow for the new maximum total
size (first + second + third neighbors) of this per\sphinxhyphen{}atom list.  There are two
ways to do this.  See the {\hyperref[\detokenize{read_data::doc}]{\sphinxcrossref{\DUrole{doc}{read\_data}}}} or
{\hyperref[\detokenize{create_box::doc}]{\sphinxcrossref{\DUrole{doc}{create\_box}}}} commands for details.
\end{sphinxadmonition}

\begin{sphinxadmonition}{note}{Note:}
\sphinxAtStartPar
Even if you do not use the \sphinxstyleemphasis{atype}, \sphinxstyleemphasis{dtype}, or \sphinxstyleemphasis{itype}
keywords, the list of topological neighbors is updated for atoms
affected by the new bond.  This in turn affects which neighbors are
considered for pair\sphinxhyphen{}wise interactions, using the weighting rules set by
the {\hyperref[\detokenize{special_bonds::doc}]{\sphinxcrossref{\DUrole{doc}{special\_bonds}}}} command.  Consider a new bond
created between atoms \(i\) and \(j\).  If \(j\) has a bonded
neighbor \(k\), then \(k\) becomes a second neighbor of \(i\).
Even if the \sphinxstyleemphasis{atype} keyword is not used to create angle \(\angle ijk\),
the pair\sphinxhyphen{}wise interaction between \(i\) and \(k\) could potentially
be turned off or weighted by the 1\textendash{}3 weighting specified
by the {\hyperref[\detokenize{special_bonds::doc}]{\sphinxcrossref{\DUrole{doc}{special\_bonds}}}} command.  This is the case
even if the “angle yes” option was used with that command.  The same
is true for third neighbors (1\textendash{}4 interactions), the \sphinxstyleemphasis{dtype} keyword, and
the “dihedral yes” option used with the
{\hyperref[\detokenize{special_bonds::doc}]{\sphinxcrossref{\DUrole{doc}{special\_bonds}}}} command.
\end{sphinxadmonition}

\sphinxAtStartPar
Note that even if your simulation starts with no bonds, you must
define a {\hyperref[\detokenize{bond_style::doc}]{\sphinxcrossref{\DUrole{doc}{bond\_style}}}} and use the
{\hyperref[\detokenize{bond_coeff::doc}]{\sphinxcrossref{\DUrole{doc}{bond\_coeff}}}} command to specify coefficients for the
\sphinxstyleemphasis{bondtype}.  Similarly, if new atom types are specified by the
\sphinxstyleemphasis{iparam} or \sphinxstyleemphasis{jparam} keywords, they must be within the range of atom
types allowed by the simulation and pair\sphinxhyphen{}wise coefficients must be
specified for the new types.

\sphinxAtStartPar
Computationally, each time step this fix is invoked, it loops over
neighbor lists and computes distances between pairs of atoms in the
list.  It also communicates between neighboring processors to
coordinate which bonds are created.  Moreover, if any bonds are
created, neighbor lists must be immediately updated on the same
time step.  This is to ensure that any pair\sphinxhyphen{}wise interactions that
should be turned “off” due to a bond creation, because they are now
excluded by the presence of the bond and the settings of the
{\hyperref[\detokenize{special_bonds::doc}]{\sphinxcrossref{\DUrole{doc}{special\_bonds}}}} command, will be immediately
recognized.  All of these operations increase the cost of a time step.
Thus, you should be cautious about invoking this fix too frequently.

\sphinxAtStartPar
You can dump out snapshots of the current bond topology via the {\hyperref[\detokenize{dump::doc}]{\sphinxcrossref{\DUrole{doc}{dump local}}}} command.

\begin{sphinxadmonition}{note}{Note:}
\sphinxAtStartPar
Creating a bond typically alters the energy of a system.  You
should be careful not to choose bond creation criteria that induce a
dramatic change in energy.  For example, if you define a very stiff
harmonic bond and create it when two atoms are separated by a distance
far from the equilibrium bond length, then the two atoms will oscillate
dramatically when the bond is formed.  More generally, you may need to
thermostat your system to compensate for energy changes resulting from
created bonds (and angles, dihedrals, impropers).
\end{sphinxadmonition}


\bigskip\hrule\bigskip



\subsection{Restart, fix\_modify, output, run start/stop, minimize info}
\label{\detokenize{fix_bond_create:restart-fix-modify-output-run-start-stop-minimize-info}}
\sphinxAtStartPar
No information about this fix is written to {\hyperref[\detokenize{restart::doc}]{\sphinxcrossref{\DUrole{doc}{binary restart files}}}}.  None of the {\hyperref[\detokenize{fix_modify::doc}]{\sphinxcrossref{\DUrole{doc}{fix\_modify}}}} options are
relevant to this fix.

\sphinxAtStartPar
This fix computes two statistics which it stores in a global vector of
length 2, which can be accessed by various {\hyperref[\detokenize{Howto_output::doc}]{\sphinxcrossref{\DUrole{doc}{output commands}}}}.  The vector values calculated by this fix are
“intensive”.

\sphinxAtStartPar
The two quantities in the global vector are
\begin{enumerate}
\sphinxsetlistlabels{\arabic}{enumi}{enumii}{(}{)}%
\item {} 
\sphinxAtStartPar
number of bonds created on the most recent creation time step

\item {} 
\sphinxAtStartPar
cumulative number of bonds created

\end{enumerate}

\sphinxAtStartPar
No parameter of this fix can be used with the \sphinxstyleemphasis{start/stop} keywords of
the {\hyperref[\detokenize{run::doc}]{\sphinxcrossref{\DUrole{doc}{run}}}} command.  This fix is not invoked during {\hyperref[\detokenize{minimize::doc}]{\sphinxcrossref{\DUrole{doc}{energy minimization}}}}.


\subsection{Restrictions}
\label{\detokenize{fix_bond_create:restrictions}}
\sphinxAtStartPar
This fix is part of the MC package.  It is only enabled if LAMMPS was
built with that package.  See the {\hyperref[\detokenize{Build_package::doc}]{\sphinxcrossref{\DUrole{doc}{Build package}}}}
doc page for more info.


\subsection{Related commands}
\label{\detokenize{fix_bond_create:related-commands}}
\sphinxAtStartPar
{\hyperref[\detokenize{fix_bond_break::doc}]{\sphinxcrossref{\DUrole{doc}{fix bond/break}}}}, {\hyperref[\detokenize{fix_bond_react::doc}]{\sphinxcrossref{\DUrole{doc}{fix bond/react}}}}, {\hyperref[\detokenize{fix_bond_swap::doc}]{\sphinxcrossref{\DUrole{doc}{fix bond/swap}}}},
{\hyperref[\detokenize{dump::doc}]{\sphinxcrossref{\DUrole{doc}{dump local}}}}, {\hyperref[\detokenize{special_bonds::doc}]{\sphinxcrossref{\DUrole{doc}{special\_bonds}}}}


\subsection{Default}
\label{\detokenize{fix_bond_create:default}}
\sphinxAtStartPar
The option defaults are iparam = (0,itype), jparam = (0,jtype), and
prob = 1.0.

\sphinxstepscope

\index{fix bond/react@\spxentry{fix bond/react}}\ignorespaces 

\section{fix bond/react command}
\label{\detokenize{fix_bond_react:fix-bond-react-command}}\label{\detokenize{fix_bond_react:index-0}}\label{\detokenize{fix_bond_react::doc}}

\subsection{Syntax}
\label{\detokenize{fix_bond_react:syntax}}
\begin{sphinxVerbatim}[commandchars=\\\{\}]
\PYG{k}{fix }\PYG{n+nv+nvIdentifier}{ID}\PYG{+w}{ }\PYG{n+nv+nvIdentifier}{group\PYGZhy{}ID}\PYG{+w}{ }\PYG{n}{bond}\PYG{o}{/}\PYG{n}{react}\PYG{+w}{ }\PYG{n}{common\PYGZus{}keyword}\PYG{+w}{ }\PYG{n}{values}\PYG{+w}{ }\PYG{o}{\PYGZam{}}
\PYG{+w}{  }\PYG{n}{react}\PYG{+w}{ }\PYG{n}{react}\PYG{o}{\PYGZhy{}}\PYG{n}{ID}\PYG{+w}{ }\PYG{n}{react}\PYG{o}{\PYGZhy{}}\PYG{n}{group}\PYG{o}{\PYGZhy{}}\PYG{n}{ID}\PYG{+w}{ }\PYG{n}{Nevery}\PYG{+w}{ }\PYG{n}{Rmin}\PYG{+w}{ }\PYG{n}{Rmax}\PYG{+w}{ }\PYG{n}{template}\PYG{o}{\PYGZhy{}}\PYG{n}{ID}\PYG{n+nv}{(}\PYG{n+nv}{pre\PYGZhy{}reacted}\PYG{n+nv}{)}\PYG{+w}{ }\PYG{n}{template}\PYG{o}{\PYGZhy{}}\PYG{n}{ID}\PYG{n+nv}{(}\PYG{n+nv}{post\PYGZhy{}reacted}\PYG{n+nv}{)}\PYG{+w}{ }\PYG{n}{map\PYGZus{}file}\PYG{+w}{ }\PYG{n}{individual\PYGZus{}keyword}\PYG{+w}{ }\PYG{n}{values}\PYG{+w}{ }\PYG{o}{\PYGZam{}}
\PYG{+w}{  }\PYG{n}{react}\PYG{+w}{ }\PYG{n}{react}\PYG{o}{\PYGZhy{}}\PYG{n}{ID}\PYG{+w}{ }\PYG{n}{react}\PYG{o}{\PYGZhy{}}\PYG{n}{group}\PYG{o}{\PYGZhy{}}\PYG{n}{ID}\PYG{+w}{ }\PYG{n}{Nevery}\PYG{+w}{ }\PYG{n}{Rmin}\PYG{+w}{ }\PYG{n}{Rmax}\PYG{+w}{ }\PYG{n}{template}\PYG{o}{\PYGZhy{}}\PYG{n}{ID}\PYG{n+nv}{(}\PYG{n+nv}{pre\PYGZhy{}reacted}\PYG{n+nv}{)}\PYG{+w}{ }\PYG{n}{template}\PYG{o}{\PYGZhy{}}\PYG{n}{ID}\PYG{n+nv}{(}\PYG{n+nv}{post\PYGZhy{}reacted}\PYG{n+nv}{)}\PYG{+w}{ }\PYG{n}{map\PYGZus{}file}\PYG{+w}{ }\PYG{n}{individual\PYGZus{}keyword}\PYG{+w}{ }\PYG{n}{values}\PYG{+w}{ }\PYG{o}{\PYGZam{}}
\PYG{+w}{  }\PYG{n}{react}\PYG{+w}{ }\PYG{n}{react}\PYG{o}{\PYGZhy{}}\PYG{n}{ID}\PYG{+w}{ }\PYG{n}{react}\PYG{o}{\PYGZhy{}}\PYG{n}{group}\PYG{o}{\PYGZhy{}}\PYG{n}{ID}\PYG{+w}{ }\PYG{n}{Nevery}\PYG{+w}{ }\PYG{n}{Rmin}\PYG{+w}{ }\PYG{n}{Rmax}\PYG{+w}{ }\PYG{n}{template}\PYG{o}{\PYGZhy{}}\PYG{n}{ID}\PYG{n+nv}{(}\PYG{n+nv}{pre\PYGZhy{}reacted}\PYG{n+nv}{)}\PYG{+w}{ }\PYG{n}{template}\PYG{o}{\PYGZhy{}}\PYG{n}{ID}\PYG{n+nv}{(}\PYG{n+nv}{post\PYGZhy{}reacted}\PYG{n+nv}{)}\PYG{+w}{ }\PYG{n}{map\PYGZus{}file}\PYG{+w}{ }\PYG{n}{individual\PYGZus{}keyword}\PYG{+w}{ }\PYG{n}{values}\PYG{+w}{ }\PYG{o}{\PYGZam{}}
\PYG{+w}{  }\PYG{n}{...}
\end{sphinxVerbatim}
\begin{itemize}
\item {} 
\sphinxAtStartPar
ID, group\sphinxhyphen{}ID are documented in {\hyperref[\detokenize{fix::doc}]{\sphinxcrossref{\DUrole{doc}{fix}}}} command.

\item {} 
\sphinxAtStartPar
bond/react = style name of this fix command

\item {} 
\sphinxAtStartPar
the common keyword/values may be appended directly after ‘bond/react’

\item {} 
\sphinxAtStartPar
common keywords apply to all reaction specifications

\item {} 
\sphinxAtStartPar
common\_keyword = \sphinxstyleemphasis{stabilization} or \sphinxstyleemphasis{reset\_mol\_ids}
\begin{sphinxalltt}
\sphinxstyleemphasis{stabilization} values = stabilize group\_prefix xmax
  stabilize = \sphinxstyleemphasis{yes} or \sphinxstyleemphasis{no}
    \sphinxstyleemphasis{yes} = perform reaction site stabilization
    \sphinxstyleemphasis{no} = no reaction site stabilization (default)
  group\_prefix = user\sphinxhyphen{}assigned prefix for the dynamic group of atoms not currently involved in a reaction
  xmax = value that is used by an internally\sphinxhyphen{}created {\hyperref[\detokenize{fix_nve_limit::doc}]{\sphinxcrossref{\DUrole{doc}{nve/limit}}}} integrator
\sphinxstyleemphasis{reset\_mol\_ids} values = \sphinxstyleemphasis{yes} or \sphinxstyleemphasis{no}
  \sphinxstyleemphasis{yes} = update molecule IDs based on new global topology (default)
  \sphinxstyleemphasis{no} = do not update molecule IDs
\end{sphinxalltt}

\item {} 
\sphinxAtStartPar
react = mandatory argument indicating new reaction specification

\item {} 
\sphinxAtStartPar
react\sphinxhyphen{}ID = user\sphinxhyphen{}assigned name for the reaction

\item {} 
\sphinxAtStartPar
react\sphinxhyphen{}group\sphinxhyphen{}ID = only atoms in this group are considered for the reaction

\item {} 
\sphinxAtStartPar
Nevery = attempt reaction every this many steps

\item {} 
\sphinxAtStartPar
Rmin = initiator atoms must be separated by more than Rmin to initiate reaction (distance units)

\item {} 
\sphinxAtStartPar
Rmax = initiator atoms must be separated by less than Rmax to initiate reaction (distance units)

\item {} 
\sphinxAtStartPar
template\sphinxhyphen{}ID(pre\sphinxhyphen{}reacted) = ID of a molecule template containing pre\sphinxhyphen{}reaction topology

\item {} 
\sphinxAtStartPar
template\sphinxhyphen{}ID(post\sphinxhyphen{}reacted) = ID of a molecule template containing post\sphinxhyphen{}reaction topology

\item {} 
\sphinxAtStartPar
map\_file = name of file specifying corresponding atom\sphinxhyphen{}IDs in the pre\sphinxhyphen{} and post\sphinxhyphen{}reacted templates

\item {} 
\sphinxAtStartPar
zero or more individual keyword/value pairs may be appended to each react argument

\item {} 
\sphinxAtStartPar
individual\_keyword = \sphinxstyleemphasis{prob} or \sphinxstyleemphasis{rate\_limit} or \sphinxstyleemphasis{max\_rxn} or \sphinxstyleemphasis{stabilize\_steps} or \sphinxstyleemphasis{custom\_charges} or \sphinxstyleemphasis{rescale\_charges} or \sphinxstyleemphasis{molecule} or \sphinxstyleemphasis{modify\_create}
\begin{sphinxalltt}
\sphinxstyleemphasis{prob} values = fraction seed
  fraction = initiate reaction with this probability if otherwise eligible
  seed = random number seed (positive integer)
\sphinxstyleemphasis{rate\_limit} = Nlimit Nsteps
  Nlimit = maximum number of reactions allowed to occur within interval
  Nsteps = the interval (number of timesteps) over which to count reactions
\sphinxstyleemphasis{max\_rxn} value = N
  N = maximum number of reactions allowed to occur
\sphinxstyleemphasis{stabilize\_steps} value = timesteps
  timesteps = number of time steps to apply the internally\sphinxhyphen{}created {\hyperref[\detokenize{fix_nve_limit::doc}]{\sphinxcrossref{\DUrole{doc}{nve/limit}}}} fix to reacting atoms
\sphinxstyleemphasis{custom\_charges} value = \sphinxstyleemphasis{no} or fragment\sphinxhyphen{}ID
  \sphinxstyleemphasis{no} = update all atomic charges (default)
  fragment\sphinxhyphen{}ID = ID of molecule fragment whose charges are updated
\sphinxstyleemphasis{rescale\_charges} value = \sphinxstyleemphasis{no} or \sphinxstyleemphasis{yes}
  \sphinxstyleemphasis{no} = do not rescale atomic charges (default)
  \sphinxstyleemphasis{yes} = rescale charges such that total charge does not change during reaction
\sphinxstyleemphasis{molecule} value = \sphinxstyleemphasis{off} or \sphinxstyleemphasis{inter} or \sphinxstyleemphasis{intra}
  \sphinxstyleemphasis{off} = allow both inter\sphinxhyphen{} and intramolecular reactions (default)
  \sphinxstyleemphasis{inter} = search for reactions between molecules with different IDs
  \sphinxstyleemphasis{intra} = search for reactions within the same molecule
\sphinxstyleemphasis{modify\_create} values = keyword arg
  \sphinxstyleemphasis{fit} arg = \sphinxstyleemphasis{all} or fragment\sphinxhyphen{}ID
    \sphinxstyleemphasis{all} = use all eligible atoms for create\sphinxhyphen{}atoms fit (default)
    fragment\sphinxhyphen{}ID = ID of molecule fragment used for create\sphinxhyphen{}atoms fit
  \sphinxstyleemphasis{overlap} value = R
    R = only insert atom/molecule if further than R from existing particles (distance units)
\end{sphinxalltt}

\end{itemize}


\subsection{Examples}
\label{\detokenize{fix_bond_react:examples}}
\sphinxAtStartPar
For unabridged example scripts and files, see examples/PACKAGES/reaction.

\begin{sphinxVerbatim}[commandchars=\\\{\}]
\PYG{k}{molecule}\PYG{+w}{ }\PYG{n}{mol1}\PYG{+w}{ }\PYG{n}{pre\PYGZus{}reacted\PYGZus{}topology.txt}
\PYG{k}{molecule}\PYG{+w}{ }\PYG{n}{mol2}\PYG{+w}{ }\PYG{n}{post\PYGZus{}reacted\PYGZus{}topology.txt}
\PYG{k}{fix }\PYG{n+nv+nvIdentifier}{5}\PYG{+w}{ }\PYG{n+nv+nvIdentifier}{all}\PYG{+w}{ }\PYG{n}{bond}\PYG{o}{/}\PYG{n}{react}\PYG{+w}{ }\PYG{n}{react}\PYG{+w}{ }\PYG{n}{myrxn1}\PYG{+w}{ }\PYG{n}{all}\PYG{+w}{ }\PYG{l+m}{1}\PYG{+w}{ }\PYG{l+m}{0}\PYG{+w}{ }\PYG{l+m}{3.25}\PYG{+w}{ }\PYG{n}{mol1}\PYG{+w}{ }\PYG{n}{mol2}\PYG{+w}{ }\PYG{n}{map\PYGZus{}file.txt}

\PYG{k}{molecule}\PYG{+w}{ }\PYG{n}{mol1}\PYG{+w}{ }\PYG{n}{pre\PYGZus{}reacted\PYGZus{}rxn1.txt}
\PYG{k}{molecule}\PYG{+w}{ }\PYG{n}{mol2}\PYG{+w}{ }\PYG{n}{post\PYGZus{}reacted\PYGZus{}rxn1.txt}
\PYG{k}{molecule}\PYG{+w}{ }\PYG{n}{mol3}\PYG{+w}{ }\PYG{n}{pre\PYGZus{}reacted\PYGZus{}rxn2.txt}
\PYG{k}{molecule}\PYG{+w}{ }\PYG{n}{mol4}\PYG{+w}{ }\PYG{n}{post\PYGZus{}reacted\PYGZus{}rxn2.txt}
\PYG{k}{fix }\PYG{n+nv+nvIdentifier}{5}\PYG{+w}{ }\PYG{n+nv+nvIdentifier}{all}\PYG{+w}{ }\PYG{n}{bond}\PYG{o}{/}\PYG{n}{react}\PYG{+w}{ }\PYG{n}{stabilization}\PYG{+w}{ }\PYG{n}{yes}\PYG{+w}{ }\PYG{n}{nvt\PYGZus{}grp}\PYG{+w}{ }\PYG{n}{.03}\PYG{+w}{ }\PYG{o}{\PYGZam{}}
\PYG{+w}{  }\PYG{n}{react}\PYG{+w}{ }\PYG{n}{myrxn1}\PYG{+w}{ }\PYG{n}{all}\PYG{+w}{ }\PYG{l+m}{1}\PYG{+w}{ }\PYG{l+m}{0}\PYG{+w}{ }\PYG{l+m}{3.25}\PYG{+w}{ }\PYG{n}{mol1}\PYG{+w}{ }\PYG{n}{mol2}\PYG{+w}{ }\PYG{n}{map\PYGZus{}file\PYGZus{}rxn1.txt}\PYG{+w}{ }\PYG{n}{prob}\PYG{+w}{ }\PYG{l+m}{0.50}\PYG{+w}{ }\PYG{l+m}{12345}\PYG{+w}{ }\PYG{o}{\PYGZam{}}
\PYG{+w}{  }\PYG{n}{react}\PYG{+w}{ }\PYG{n}{myrxn2}\PYG{+w}{ }\PYG{n}{all}\PYG{+w}{ }\PYG{l+m}{1}\PYG{+w}{ }\PYG{l+m}{0}\PYG{+w}{ }\PYG{l+m}{2.75}\PYG{+w}{ }\PYG{n}{mol3}\PYG{+w}{ }\PYG{n}{mol4}\PYG{+w}{ }\PYG{n}{map\PYGZus{}file\PYGZus{}rxn2.txt}\PYG{+w}{ }\PYG{n}{prob}\PYG{+w}{ }\PYG{l+m}{0.25}\PYG{+w}{ }\PYG{l+m}{12345}
\PYG{k}{fix }\PYG{n+nv+nvIdentifier}{6}\PYG{+w}{ }\PYG{n+nv+nvIdentifier}{nvt\PYGZus{}grp\PYGZus{}REACT}\PYG{+w}{ }\PYG{n}{nvt}\PYG{+w}{ }\PYG{n}{temp}\PYG{+w}{ }\PYG{l+m}{300}\PYG{+w}{ }\PYG{l+m}{300}\PYG{+w}{ }\PYG{l+m}{100}\PYG{+w}{ }\PYG{c}{\PYGZsh{} set thermostat after bond/react}
\end{sphinxVerbatim}


\subsection{Description}
\label{\detokenize{fix_bond_react:description}}
\sphinxAtStartPar
Initiate complex covalent bonding (topology) changes. These topology
changes will be referred to as ‘reactions’ throughout this
documentation. Topology changes are defined in pre\sphinxhyphen{} and post\sphinxhyphen{}reaction
molecule templates and can include creation and deletion of bonds,
angles, dihedrals, impropers, bond types, angle types, dihedral types,
atom types, or atomic charges. In addition, reaction by\sphinxhyphen{}products or
other molecules can be identified and deleted. Finally, atoms can be
created and inserted at specific positions relative to the reaction
site.

\sphinxAtStartPar
Fix bond/react does not use quantum mechanical (e.g., {\hyperref[\detokenize{fix_qmmm::doc}]{\sphinxcrossref{\DUrole{doc}{fix qmmm}}}}) or
pairwise bond\sphinxhyphen{}order potential (e.g., {\hyperref[\detokenize{pair_tersoff::doc}]{\sphinxcrossref{\DUrole{doc}{Tersoff}}}} or
{\hyperref[\detokenize{pair_airebo::doc}]{\sphinxcrossref{\DUrole{doc}{AIREBO}}}}) methods to
determine bonding changes a priori. Rather, it uses a distance\sphinxhyphen{}based
probabilistic criteria to effect predetermined topology changes in
simulations using standard force fields.

\sphinxAtStartPar
This fix was created to facilitate the dynamic creation of polymeric,
amorphous or highly cross\sphinxhyphen{}linked systems. A suggested workflow for
using this fix is
\begin{enumerate}
\sphinxsetlistlabels{\arabic}{enumi}{enumii}{(}{)}%
\item {} 
\sphinxAtStartPar
identify a reaction to be simulated

\item {} 
\sphinxAtStartPar
build a molecule template of the reaction site before the reaction has occurred

\item {} 
\sphinxAtStartPar
build a molecule template of the reaction site after the reaction has occurred

\item {} 
\sphinxAtStartPar
create a map that relates the template\sphinxhyphen{}atom\sphinxhyphen{}IDs of each atom between pre\sphinxhyphen{} and post\sphinxhyphen{}reaction molecule templates

\item {} 
\sphinxAtStartPar
fill a simulation box with molecules and run a simulation with fix bond/react.

\end{enumerate}

\begin{sphinxadmonition}{note}{Note:}
\sphinxAtStartPar
\DUrole{versionmodified,added}{Added in version 15Sep2022.}

\sphinxAtStartPar
{\hyperref[\detokenize{Howto_type_labels::doc}]{\sphinxcrossref{\DUrole{doc}{Type labels}}}} allow for molecule templates
and data files to use alphanumeric atom types that match those of
a force field. Input files that use type labels are inherently
compatible with each other and portable between different
simulations. Therefore, it is highly recommended to use type labels
to specify atom, bond, etc. types when using fix bond/react.
\end{sphinxadmonition}

\sphinxAtStartPar
Only one ‘fix bond/react’ command can be used at a time. Multiple
reactions can be simultaneously applied by specifying multiple \sphinxstyleemphasis{react}
arguments to a single ‘fix bond/react’ command. This syntax is
necessary because the “common” keywords are applied to all reactions.

\sphinxAtStartPar
The \sphinxstyleemphasis{stabilization} keyword enables reaction site stabilization.
Reaction site stabilization is performed by including reacting atoms
in an internally\sphinxhyphen{}created fix {\hyperref[\detokenize{fix_nve_limit::doc}]{\sphinxcrossref{\DUrole{doc}{nve/limit}}}} time
integrator for a set number of time steps given by the
\sphinxstyleemphasis{stabilize\_steps} keyword. While reacting atoms are being time
integrated by the internal nve/limit, they are prevented from being
involved in any new reactions. The \sphinxstyleemphasis{xmax} value keyword should
typically be set to the maximum distance that non\sphinxhyphen{}reacting atoms move
during the simulation.

\sphinxAtStartPar
Fix bond/react creates and maintains two important dynamic groups of
atoms when using the \sphinxstyleemphasis{stabilization} keyword. The first group contains
all atoms currently involved in a reaction; this group is
automatically time\sphinxhyphen{}integrated by an internally\sphinxhyphen{}created
{\hyperref[\detokenize{fix_nve_limit::doc}]{\sphinxcrossref{\DUrole{doc}{nve/limit}}}} integrator. The second group contains
all atoms currently not involved in a reaction. This group should be
controlled by a thermostat in order to time integrate the system. The name
of this group of non\sphinxhyphen{}reacting atoms is created by appending ‘\_REACT’
to the group\sphinxhyphen{}ID argument of the \sphinxstyleemphasis{stabilization} keyword, as shown in
the second example above.

\begin{sphinxadmonition}{note}{Note:}
\sphinxAtStartPar
When using reaction stabilization, you should generally \sphinxstylestrong{not} have
a separate thermostat that acts on the “all” group.
\end{sphinxadmonition}

\sphinxAtStartPar
The group\sphinxhyphen{}ID set using the \sphinxstyleemphasis{stabilization} keyword can be an existing
static group or a previously\sphinxhyphen{}unused group\sphinxhyphen{}ID. It cannot be specified
as “all”. If the group\sphinxhyphen{}ID is previously unused, the fix bond/react
command creates a {\hyperref[\detokenize{group::doc}]{\sphinxcrossref{\DUrole{doc}{dynamic group}}}} that is initialized to
include all atoms. If the group\sphinxhyphen{}ID is that of an existing static
group, the group is used as the parent group of new,
internally\sphinxhyphen{}created dynamic group. In both cases, this new dynamic
group is named by appending ‘\_REACT’ to the group\sphinxhyphen{}ID (e.g.,
nvt\_grp\_REACT). By specifying an existing group, you may thermostat
constant\sphinxhyphen{}topology parts of your system separately. The dynamic group
contains only atoms not involved in a reaction at a given time step,
and therefore should be used by a subsequent system\sphinxhyphen{}wide time
integrator such as {\hyperref[\detokenize{fix_nh::doc}]{\sphinxcrossref{\DUrole{doc}{fix nvt}}}}, {\hyperref[\detokenize{fix_nh::doc}]{\sphinxcrossref{\DUrole{doc}{fix npt}}}}, or
{\hyperref[\detokenize{fix_nve::doc}]{\sphinxcrossref{\DUrole{doc}{fix nve}}}}, as shown in the second example
above (full examples can be found in examples/PACKAGES/reaction). The time
integration command should be placed after the fix bond/react command
due to the internal dynamic grouping performed by fix bond/react.

\begin{sphinxadmonition}{note}{Note:}
\sphinxAtStartPar
If the group\sphinxhyphen{}ID is an existing static group, react\sphinxhyphen{}group\sphinxhyphen{}IDs
should also be specified as this static group or a subset.
\end{sphinxadmonition}

\sphinxAtStartPar
The \sphinxstyleemphasis{reset\_mol\_ids} keyword invokes the {\hyperref[\detokenize{reset_atoms::doc}]{\sphinxcrossref{\DUrole{doc}{reset\_atoms mol}}}} command after a reaction occurs, to ensure that
molecule IDs are consistent with the new bond topology. The group\sphinxhyphen{}ID
used for {\hyperref[\detokenize{reset_atoms::doc}]{\sphinxcrossref{\DUrole{doc}{reset\_atoms mol}}}} is the group\sphinxhyphen{}ID for this
fix.  Resetting molecule IDs is necessarily a global operation, so it
can be slow for very large systems.

\sphinxAtStartPar
The following comments pertain to each \sphinxstyleemphasis{react} argument (in other
words, they can be customized for each reaction, or reaction step):

\sphinxAtStartPar
A check for possible new reaction sites is performed every \sphinxstyleemphasis{Nevery}
time steps. \sphinxstyleemphasis{Nevery} can be specified with an equal\sphinxhyphen{}style
{\hyperref[\detokenize{variable::doc}]{\sphinxcrossref{\DUrole{doc}{variable}}}}, whose value is rounded up to the nearest
integer.

\sphinxAtStartPar
Three physical conditions must be met for a reaction to occur. First,
an initiator atom pair must be identified within the reaction distance
cutoffs. Second, the topology surrounding the initiator atom pair must
match the topology of the pre\sphinxhyphen{}reaction template. Only atom types and
bond connectivity are used to identify a valid reaction site (not bond
types, etc.). Finally, any reaction constraints listed in the map file
(see below) must be satisfied. If all of these conditions are met, the
reaction site is eligible to be modified to match the post\sphinxhyphen{}reaction
template.

\sphinxAtStartPar
An initiator atom pair will be identified if several conditions are
met. First, a pair of atoms \(i\) and \(j\) within the specified
react\sphinxhyphen{}group\sphinxhyphen{}ID of type \sphinxstyleemphasis{itype} and \sphinxstyleemphasis{jtype} must be separated by a distance
between \sphinxstyleemphasis{Rmin} and \sphinxstyleemphasis{Rmax}. \sphinxstyleemphasis{Rmin} and \sphinxstyleemphasis{Rmax} can be specified with
equal\sphinxhyphen{}style {\hyperref[\detokenize{variable::doc}]{\sphinxcrossref{\DUrole{doc}{variables}}}}. For example, these reaction cutoffs
can be functions of the reaction conversion using the following commands:

\begin{sphinxVerbatim}[commandchars=\\\{\}]
\PYG{k}{variable }\PYG{n+nv+nvIdentifier}{rmax}\PYG{+w}{ }\PYG{n}{equal}\PYG{+w}{ }\PYG{l+m}{0}\PYG{+w}{ }\PYG{c}{\PYGZsh{} initialize variable before bond/react}
\PYG{k}{fix }\PYG{n+nv+nvIdentifier}{myrxn}\PYG{+w}{ }\PYG{n+nv+nvIdentifier}{all}\PYG{+w}{ }\PYG{n}{bond}\PYG{o}{/}\PYG{n}{react}\PYG{+w}{ }\PYG{n}{react}\PYG{+w}{ }\PYG{n}{myrxn1}\PYG{+w}{ }\PYG{n}{all}\PYG{+w}{ }\PYG{l+m}{1}\PYG{+w}{ }\PYG{l+m}{0}\PYG{+w}{ }\PYG{n}{v\PYGZus{}rmax}\PYG{+w}{ }\PYG{n}{mol1}\PYG{+w}{ }\PYG{n}{mol2}\PYG{+w}{ }\PYG{n}{map\PYGZus{}file.txt}
\PYG{k}{variable }\PYG{n+nv+nvIdentifier}{rmax}\PYG{+w}{ }\PYG{n}{equal}\PYG{+w}{ }\PYG{l+m}{3}\PYG{o}{+}\PYG{n}{f\PYGZus{}myrxn[1]}\PYG{o}{/}\PYG{l+m}{100}\PYG{+w}{ }\PYG{c}{\PYGZsh{} arbitrary function of reaction count}
\end{sphinxVerbatim}

\sphinxAtStartPar
The following criteria are used if multiple candidate initiator atom
pairs are identified within the cutoff distance:
\begin{enumerate}
\sphinxsetlistlabels{\arabic}{enumi}{enumii}{(}{)}%
\item {} 
\sphinxAtStartPar
If the initiator atoms in the pre\sphinxhyphen{}reaction template are not 1\textendash{}2
neighbors (i.e., not directly bonded) the closest potential partner is
chosen.

\item {} 
\sphinxAtStartPar
Otherwise, if the initiator atoms in the pre\sphinxhyphen{}reaction template are 1\textendash{}2
neighbors (i.e. directly bonded) the farthest potential partner is
chosen.

\item {} 
\sphinxAtStartPar
Then, if both an atom \(i\) and atom \(j\) have each other as
initiator partners, these two atoms are identified as the initiator atom
pair of the reaction site.

\end{enumerate}

\sphinxAtStartPar
Note that it can be helpful to select unique atom types for the
initiator atoms: if an initiator atom pair is identified, as described
in the previous steps, but it does not correspond to the same pair
specified in the pre\sphinxhyphen{}reaction template, an otherwise eligible reaction
could be prevented from occurring. Once this unique initiator atom
pair is identified for each reaction, there could be two or more
reactions that involve the same atom on the same time step. If this is
the case, only one such reaction is permitted to occur. This reaction
is chosen randomly from all potential reactions involving the
overlapping atom. This capability allows, for example, different
reaction pathways to proceed from identical reaction sites with
user\sphinxhyphen{}specified probabilities.

\sphinxAtStartPar
The pre\sphinxhyphen{}reacted molecule template is specified by a molecule command.
This molecule template file contains a sample reaction site and its
surrounding topology. As described below, the initiator atom pairs of
the pre\sphinxhyphen{}reacted template are specified by atom ID in the map file. The
pre\sphinxhyphen{}reacted molecule template should contain as few atoms as possible
while still completely describing the topology of all atoms affected
by the reaction (which includes all atoms that change atom type or
connectivity, and all bonds that change bond type). For example, if
the force field contains dihedrals, the pre\sphinxhyphen{}reacted template should
contain any atom within three bonds of reacting atoms.

\sphinxAtStartPar
Some atoms in the pre\sphinxhyphen{}reacted template that are not reacting may have
missing topology with respect to the simulation. For example, the
pre\sphinxhyphen{}reacted template may contain an atom that, in the simulation, is
currently connected to the rest of a long polymer chain. These are
referred to as edge atoms, and are also specified in the map file. All
pre\sphinxhyphen{}reaction template atoms should be linked to an initiator atom, via
at least one path that does not involve edge atoms. When the
pre\sphinxhyphen{}reaction template contains edge atoms, not all atoms, bonds,
charges, etc. specified in the reaction templates will be updated.
Specifically, topology that involves only atoms that are “too near” to
template edges will not be updated. The definition of “too near the
edge” depends on which interactions are defined in the simulation. If
the simulation has defined dihedrals, atoms within two bonds of edge
atoms are considered “too near the edge.” If the simulation defines
angles, but not dihedrals, atoms within one bond of edge atoms are
considered “too near the edge.” If just bonds are defined, only edge
atoms are considered “too near the edge.”

\begin{sphinxadmonition}{note}{Note:}
\sphinxAtStartPar
Small molecules (i.e., ones that have all their atoms contained
within the reaction templates) never have edge atoms.
\end{sphinxadmonition}

\sphinxAtStartPar
Note that some care must be taken when a building a molecule template
for a given simulation. All atom types in the pre\sphinxhyphen{}reacted template
must be the same as those of a potential reaction site in the
simulation. A detailed discussion of matching molecule template atom
types with the simulation is provided on the {\hyperref[\detokenize{molecule::doc}]{\sphinxcrossref{\DUrole{doc}{molecule}}}}
command page. It is highly recommended to use {\hyperref[\detokenize{Howto_type_labels::doc}]{\sphinxcrossref{\DUrole{doc}{Type labels}}}}
(added in version 15Sep2022) in both molecule templates and data
files, which automates the process of syncing atom types between
different input files.

\sphinxAtStartPar
The post\sphinxhyphen{}reacted molecule template contains a sample of the reaction
site and its surrounding topology after the reaction has occurred. It
must contain the same number of atoms as the pre\sphinxhyphen{}reacted template
(unless there are created atoms). A one\sphinxhyphen{}to\sphinxhyphen{}one correspondence between
the atom IDs in the pre\sphinxhyphen{} and post\sphinxhyphen{}reacted templates is specified in
the map file as described below. Note that during a reaction, an atom,
bond, etc. type may change to one that was previously not present in
the simulation. These new types must also be defined during the setup
of a given simulation. A discussion of correctly handling this is also
provided on the {\hyperref[\detokenize{molecule::doc}]{\sphinxcrossref{\DUrole{doc}{molecule}}}} command page.

\begin{sphinxadmonition}{note}{Note:}
\sphinxAtStartPar
When a reaction occurs, it is possible that the resulting
topology/atom (e.g., special bonds, dihedrals) exceeds that of
the existing system and reaction templates. As when inserting
molecules, enough space for this increased topology/atom must be
reserved by using the relevant “extra” keywords to the
{\hyperref[\detokenize{read_data::doc}]{\sphinxcrossref{\DUrole{doc}{read\_data}}}} or {\hyperref[\detokenize{create_box::doc}]{\sphinxcrossref{\DUrole{doc}{create\_box}}}} commands.
\end{sphinxadmonition}

\sphinxAtStartPar
The map file is a text document with the following format:

\sphinxAtStartPar
A map file has a header and a body. The header of map file the
contains one mandatory keyword and five optional keywords. The
mandatory keyword is \sphinxstyleemphasis{equivalences}:
\begin{sphinxalltt}
N \sphinxstyleemphasis{equivalences} = \# of atoms N in the reaction molecule templates
\end{sphinxalltt}

\sphinxAtStartPar
The optional keywords are \sphinxstyleemphasis{edgeIDs}, \sphinxstyleemphasis{deleteIDs}, \sphinxstyleemphasis{chiralIDs}, and
\sphinxstyleemphasis{constraints}:
\begin{sphinxalltt}
N \sphinxstyleemphasis{edgeIDs} = \# of edge atoms N in the pre\sphinxhyphen{}reacted molecule template
N \sphinxstyleemphasis{deleteIDs} = \# of atoms N that are deleted
N \sphinxstyleemphasis{createIDs} = \# of atoms N that are created
N \sphinxstyleemphasis{chiralIDs} = \# of chiral centers N
N \sphinxstyleemphasis{constraints} = \# of reaction constraints N
\end{sphinxalltt}

\sphinxAtStartPar
The body of the map file contains two mandatory sections and five
optional sections. The first mandatory section begins with the keyword
“InitiatorIDs” and lists the two atom IDs of the initiator atom pair
in the pre\sphinxhyphen{}reacted molecule template. The second mandatory section
begins with the keyword “Equivalences” and lists a one\sphinxhyphen{}to\sphinxhyphen{}one
correspondence between atom IDs of the pre\sphinxhyphen{} and post\sphinxhyphen{}reacted
templates. The first column is an atom ID of the pre\sphinxhyphen{}reacted molecule
template, and the second column is the corresponding atom ID of the
post\sphinxhyphen{}reacted molecule template. The first optional section begins with
the keyword “EdgeIDs” and lists the atom IDs of edge atoms in the
pre\sphinxhyphen{}reacted molecule template. The second optional section begins with
the keyword “DeleteIDs” and lists the atom IDs of pre\sphinxhyphen{}reaction
template atoms to delete. The third optional section begins with the
keyword “CreateIDs” and lists the atom IDs of the post\sphinxhyphen{}reaction
template atoms to create. The fourth optional section begins with the
keyword “ChiralIDs” lists the atom IDs of chiral atoms whose
handedness should be enforced. The fifth optional section begins with
the keyword “Constraints” and lists additional criteria that must be
satisfied in order for the reaction to occur. Currently, there are
six types of constraints available, as discussed below: “distance”,
“angle”, “dihedral”, “arrhenius”, “rmsd”, and “custom”.

\sphinxAtStartPar
A sample map file is given below:


\bigskip\hrule\bigskip


\begin{sphinxVerbatim}[commandchars=\\\{\}]
\PYGZsh{} this is a map file

7 equivalences
2 edgeIDs

InitiatorIDs

3
5

EdgeIDs

1
7

Equivalences

1   1
2   2
3   3
4   4
5   5
6   6
7   7
\end{sphinxVerbatim}


\bigskip\hrule\bigskip


\sphinxAtStartPar
A user\sphinxhyphen{}specified set of atoms can be deleted by listing their
pre\sphinxhyphen{}reaction template IDs in the DeleteIDs section. A deleted atom
must still be included in the post\sphinxhyphen{}reaction molecule template, in
which it cannot be bonded to an atom that is not deleted. In addition
to deleting unwanted reaction by\sphinxhyphen{}products, this feature can be used to
remove specific topologies, such as small rings, that may be otherwise
indistinguishable.

\sphinxAtStartPar
Atoms can be created by listing their post\sphinxhyphen{}reaction template IDs in
the CreateIDs section. A created atom should not be included in the
pre\sphinxhyphen{}reaction template. The inserted positions of created atoms are
determined by the coordinates of the post\sphinxhyphen{}reaction template, after
optimal translation and rotation of the post\sphinxhyphen{}reaction template to the
reaction site (using a fit with atoms that are neither created nor
deleted). The \sphinxstyleemphasis{modify\_create} keyword can be used to modify the
default behavior when creating atoms. The \sphinxstyleemphasis{modify\_create} keyword has
two sub\sphinxhyphen{}keywords, \sphinxstyleemphasis{fit} and \sphinxstyleemphasis{overlap}. One or more of the sub\sphinxhyphen{}keywords
may be used after the \sphinxstyleemphasis{modify\_create} keyword. The \sphinxstyleemphasis{fit} sub\sphinxhyphen{}keyword
can be used to specify which post\sphinxhyphen{}reaction atoms are used for the
optimal translation and rotation of the post\sphinxhyphen{}reaction template. The
fragment\sphinxhyphen{}ID value of the \sphinxstyleemphasis{fit} sub\sphinxhyphen{}keyword must be the name of a
molecule fragment defined in the post\sphinxhyphen{}reaction {\hyperref[\detokenize{molecule::doc}]{\sphinxcrossref{\DUrole{doc}{molecule}}}} template, and only atoms in this fragment are used for the
fit. Atoms are created only if no current atom in the simulation is
within a distance \(R\) of any created atom, including the effect of
periodic boundary conditions if applicable. \(R\) is defined by the
\sphinxstyleemphasis{overlap} sub\sphinxhyphen{}keyword. Note that the default value for \(R\) is 0.0, which
will allow atoms to strongly overlap if you are inserting where other
atoms are present. The velocity of each created atom is initialized in
a random direction with a magnitude calculated from the instantaneous
temperature of the reaction site.

\begin{sphinxadmonition}{note}{Note:}
\sphinxAtStartPar
The ‘Coords’ section must be included in the post\sphinxhyphen{}reaction template
when creating atoms because these coordinates are used to determine
where new atoms are inserted.
\end{sphinxadmonition}

\sphinxAtStartPar
The handedness of atoms that are chiral centers can be enforced by
listing their IDs in the ChiralIDs section. A chiral atom must be
bonded to four atoms with mutually different atom types. This feature
uses the coordinates and types of the involved atoms in the
pre\sphinxhyphen{}reaction template to determine handedness. Three atoms bonded to
the chiral center are arbitrarily chosen, to define an oriented plane,
and the relative position of the fourth bonded atom determines the
chiral center’s handedness.

\sphinxAtStartPar
Any number of additional constraints may be specified in the
Constraints section of the map file. The constraint of type “distance”
has syntax as follows:
\begin{sphinxalltt}
distance \sphinxstyleemphasis{ID1} \sphinxstyleemphasis{ID2} \sphinxstyleemphasis{rmin} \sphinxstyleemphasis{rmax}
\end{sphinxalltt}

\sphinxAtStartPar
where “distance” is the required keyword, \sphinxstyleemphasis{ID1} and \sphinxstyleemphasis{ID2} are
pre\sphinxhyphen{}reaction atom IDs (or molecule\sphinxhyphen{}fragment IDs, see below), and these
two atoms must be separated by a distance between \sphinxstyleemphasis{rmin} and \sphinxstyleemphasis{rmax}
for the reaction to occur.

\sphinxAtStartPar
The constraint of type “angle” has the following syntax:
\begin{sphinxalltt}
angle \sphinxstyleemphasis{ID1} \sphinxstyleemphasis{ID2} \sphinxstyleemphasis{ID3} \sphinxstyleemphasis{amin} \sphinxstyleemphasis{amax}
\end{sphinxalltt}

\sphinxAtStartPar
where “angle” is the required keyword, \sphinxstyleemphasis{ID1}, \sphinxstyleemphasis{ID2} and \sphinxstyleemphasis{ID3} are
pre\sphinxhyphen{}reaction atom IDs (or molecule\sphinxhyphen{}fragment IDs, see below), and these
three atoms must form an angle between \sphinxstyleemphasis{amin} and \sphinxstyleemphasis{amax} for the
reaction to occur (where \sphinxstyleemphasis{ID2} is the central atom). Angles must be
specified in degrees. This constraint can be used to enforce a certain
orientation between reacting molecules.

\sphinxAtStartPar
The constraint of type “dihedral” has the following syntax:
\begin{sphinxalltt}
dihedral \sphinxstyleemphasis{ID1} \sphinxstyleemphasis{ID2} \sphinxstyleemphasis{ID3} \sphinxstyleemphasis{ID4} \sphinxstyleemphasis{amin} \sphinxstyleemphasis{amax} \sphinxstyleemphasis{amin2} \sphinxstyleemphasis{amax2}
\end{sphinxalltt}

\sphinxAtStartPar
where “dihedral” is the required keyword, and \sphinxstyleemphasis{ID1}, \sphinxstyleemphasis{ID2}, \sphinxstyleemphasis{ID3}
and \sphinxstyleemphasis{ID4} are pre\sphinxhyphen{}reaction atom IDs (or molecule\sphinxhyphen{}fragment IDs, see
below). Dihedral angles are calculated in the interval \((-180^\circ,180^\circ]\).
Refer to the {\hyperref[\detokenize{dihedral_style::doc}]{\sphinxcrossref{\DUrole{doc}{dihedral style}}}} documentation for
further details on convention. If \sphinxstyleemphasis{amin} is less than \sphinxstyleemphasis{amax}, these
four atoms must form a dihedral angle greater than \sphinxstyleemphasis{amin} \sphinxstylestrong{and} less
than \sphinxstyleemphasis{amax} for the reaction to occur. If \sphinxstyleemphasis{amin} is greater than
\sphinxstyleemphasis{amax}, these four atoms must form a dihedral angle greater than
\sphinxstyleemphasis{amin} \sphinxstylestrong{or} less than \sphinxstyleemphasis{amax} for the reaction to occur. Angles must
be specified in degrees. Optionally, a second range of permissible
angles \sphinxstyleemphasis{amin2} to \sphinxstyleemphasis{amax2} can be specified.

\sphinxAtStartPar
For the ‘distance’, ‘angle’, and ‘dihedral’ constraints (explained
above), atom IDs can be replaced by pre\sphinxhyphen{}reaction molecule\sphinxhyphen{}fragment
IDs. The molecule\sphinxhyphen{}fragment ID must begin with a letter. The location
of the ID is the geometric center of all atom positions in the
fragment. The molecule fragment must have been defined in the
{\hyperref[\detokenize{molecule::doc}]{\sphinxcrossref{\DUrole{doc}{molecule}}}} command for the pre\sphinxhyphen{}reaction template.

\sphinxAtStartPar
The constraint of type ‘arrhenius’ imposes an additional reaction
probability according to the modified Arrhenius equation,
\begin{equation*}
\begin{split}k = AT^{n}e^{-E_{a}/k_{B}T}.\end{split}
\end{equation*}
\sphinxAtStartPar
The Arrhenius constraint has the following syntax:
\begin{sphinxalltt}
arrhenius \sphinxstyleemphasis{A} \sphinxstyleemphasis{n} \sphinxstyleemphasis{E\_a} \sphinxstyleemphasis{seed}
\end{sphinxalltt}

\sphinxAtStartPar
where “arrhenius” is the required keyword, \sphinxstyleemphasis{A} is the pre\sphinxhyphen{}exponential
factor, \sphinxstyleemphasis{n} is the exponent of the temperature dependence, \(E_a\)
is the activation energy ({\hyperref[\detokenize{units::doc}]{\sphinxcrossref{\DUrole{doc}{units}}}} of energy), and \sphinxstyleemphasis{seed} is a
random number seed. The temperature is defined as the instantaneous
temperature averaged over all atoms in the reaction site and is
calculated in the same manner as for example
{\hyperref[\detokenize{compute_temp_chunk::doc}]{\sphinxcrossref{\DUrole{doc}{compute temp/chunk}}}}. Currently, there are no
options for additional temperature averaging or velocity\sphinxhyphen{}biased
temperature calculations. A uniform random number between 0 and 1 is
generated using \sphinxstyleemphasis{seed}; if this number is less than the result of the
Arrhenius equation above, the reaction is permitted to occur.

\sphinxAtStartPar
The constraint of type ‘rmsd’ has the following syntax:
\begin{sphinxalltt}
rmsd \sphinxstyleemphasis{RMSDmax} \sphinxstyleemphasis{molfragment}
\end{sphinxalltt}

\sphinxAtStartPar
where “rmsd” is the required keyword, and \sphinxstyleemphasis{RMSDmax} is the maximum
root\sphinxhyphen{}mean\sphinxhyphen{}square deviation between atom positions of the pre\sphinxhyphen{}reaction
template and the local reaction site (distance units), after optimal
translation and rotation of the pre\sphinxhyphen{}reaction template. Optionally, the
name of a molecule fragment (of the pre\sphinxhyphen{}reaction template) can be
specified by \sphinxstyleemphasis{molfragment}. If a molecule fragment is specified,
only atoms that are part of this molecule fragment are used to
determine the RMSD. A molecule fragment must have been defined in the
{\hyperref[\detokenize{molecule::doc}]{\sphinxcrossref{\DUrole{doc}{molecule}}}} command for the pre\sphinxhyphen{}reaction template. For
example, the molecule fragment could consist of only the backbone
atoms of a polymer chain. This constraint can be used to enforce a
specific relative position and orientation between reacting molecules.

\sphinxAtStartPar
\DUrole{versionmodified,changed}{Changed in version 22Dec2022.}

\sphinxAtStartPar
The constraint of type “custom” has the following syntax:
\begin{sphinxalltt}
custom \sphinxstyleemphasis{varstring}
\end{sphinxalltt}

\sphinxAtStartPar
where ‘custom’ is the required keyword, and \sphinxstyleemphasis{varstring} is a variable
expression. The expression must be a valid equal\sphinxhyphen{}style variable
formula that can be read by the {\hyperref[\detokenize{variable::doc}]{\sphinxcrossref{\DUrole{doc}{variable}}}} command,
after any special reaction functions are evaluated. If the resulting
expression is zero, the reaction is prevented from occurring;
otherwise, it is permitted to occur. There are three special reaction
functions available, ‘rxnbond’, ‘rxnsum’, and ‘rxnave’. The ‘rxnbond’
function allows per\sphinxhyphen{}bond values to be included in the variable strings
of the custom constraint. The ‘rxnbond’ function has two mandatory
arguments. The first argument is the ID of a previously defined
‘compute bond/local’ command. This ‘compute bond/local’ must compute
only one value, e.g. bond force. This value is returned by the
‘rxnbond’ function. The second argument is the name of a molecule
fragment in the pre\sphinxhyphen{}reaction template. The fragment must contain
exactly two atoms, corresponding to the atoms involved in the bond
whose value should be calculated. An example of a constraint that uses
the force experienced by a bond is provided below. The ‘rxnsum’ and
‘rxnave’ functions operate over the atoms in a given reaction site,
and have one mandatory argument and one optional argument. The
mandatory argument is the identifier for an atom\sphinxhyphen{}style variable. The
second, optional argument is the name of a molecule fragment in the
pre\sphinxhyphen{}reaction template, and can be used to operate over a subset of
atoms in the reaction site. The ‘rxnsum’ function sums the atom\sphinxhyphen{}style
variable over the reaction site, while the ‘rxnave’ returns the
average value. For example, a constraint on the total potential energy
of atoms involved in the reaction can be imposed as follows:

\begin{sphinxVerbatim}[commandchars=\\\{\}]
\PYG{k}{compute }\PYG{n+nv+nvIdentifier}{1}\PYG{+w}{ }\PYG{n+nv+nvIdentifier}{all}\PYG{+w}{ }\PYG{n}{pe}\PYG{o}{/}\PYG{n}{atom}\PYG{+w}{ }\PYG{c}{\PYGZsh{} in LAMMPS input script}
\PYG{k}{variable }\PYG{n+nv+nvIdentifier}{my\PYGZus{}pe}\PYG{+w}{ }\PYG{n}{atom}\PYG{+w}{ }\PYG{n}{c\PYGZus{}1}\PYG{+w}{ }\PYG{c}{\PYGZsh{} in LAMMPS input script}
\end{sphinxVerbatim}

\begin{sphinxVerbatim}[commandchars=\\\{\}]
\PYG{n}{custom}\PYG{+w}{ }\PYG{l+s}{\PYGZdq{}}\PYG{l+s}{rxnsum(v\PYGZus{}my\PYGZus{}pe) \PYGZgt{} 100}\PYG{l+s}{\PYGZdq{}}\PYG{+w}{ }\PYG{c}{\PYGZsh{} in Constraints section of map file}
\end{sphinxVerbatim}

\sphinxAtStartPar
The above example prevents the reaction from occurring unless the
total potential energy of the reaction site is above 100. As a second
example, this time using the ‘rxnbond’ function, consider a modified
Arrhenius constraint that depends on the bond force of a specific bond:

\begin{sphinxVerbatim}[commandchars=\\\{\}]
\PYG{c}{\PYGZsh{} in LAMMPS input script}

\PYG{k}{compute }\PYG{n+nv+nvIdentifier}{bondforce}\PYG{+w}{ }\PYG{n+nv+nvIdentifier}{all}\PYG{+w}{ }\PYG{n}{bond}\PYG{o}{/}\PYG{n}{local}\PYG{+w}{ }\PYG{n}{force}

\PYG{k}{compute }\PYG{n+nv+nvIdentifier}{ke\PYGZus{}atom}\PYG{+w}{ }\PYG{n+nv+nvIdentifier}{all}\PYG{+w}{ }\PYG{n}{ke}\PYG{o}{/}\PYG{n}{atom}
\PYG{k}{variable }\PYG{n+nv+nvIdentifier}{ke}\PYG{+w}{ }\PYG{n}{atom}\PYG{+w}{ }\PYG{n}{c\PYGZus{}ke\PYGZus{}atom}

\PYG{k}{variable }\PYG{n+nv+nvIdentifier}{E\PYGZus{}a}\PYG{+w}{ }\PYG{n}{equal}\PYG{+w}{ }\PYG{l+m}{100.0}\PYG{+w}{ }\PYG{c}{\PYGZsh{} activation energy}
\PYG{k}{variable }\PYG{n+nv+nvIdentifier}{l0}\PYG{+w}{ }\PYG{n}{equal}\PYG{+w}{ }\PYG{l+m}{1.0}\PYG{+w}{ }\PYG{c}{\PYGZsh{} characteristic length}
\end{sphinxVerbatim}

\begin{sphinxVerbatim}[commandchars=\\\{\}]
\PYG{c}{\PYGZsh{} in Constraints section of map file}

\PYG{n}{custom}\PYG{+w}{ }\PYG{l+s}{\PYGZdq{}}\PYG{l+s}{exp(\PYGZhy{}(v\PYGZus{}E\PYGZus{}a\PYGZhy{}rxnbond(c\PYGZus{}bondforce,bond1frag)*v\PYGZus{}l0)/(2/3*rxnave(v\PYGZus{}ke))) \PYGZlt{} random(0,1,12345)}\PYG{l+s}{\PYGZdq{}}
\end{sphinxVerbatim}

\sphinxAtStartPar
By using an inequality and the ‘random(x,y,z)’ function, the left\sphinxhyphen{}hand
side can be interpreted as the probability of the reaction occurring,
similar to the ‘arrhenius’ constraint above.

\sphinxAtStartPar
By default, all constraints must be satisfied for the reaction to
occur. In other words, constraints are evaluated as a series of
logical values using the logical AND operator “\&\&”. More complex logic
can be achieved by explicitly adding the logical AND operator “\&\&” or
the logical OR operator “||” after a given constraint command. If a
logical operator is specified after a constraint, it must be placed
after all constraint parameters, on the same line as the constraint
(one per line). Similarly, parentheses can be used to group
constraints. The expression that results from concatenating all
constraints should be a valid logical expression that can be read by
the {\hyperref[\detokenize{variable::doc}]{\sphinxcrossref{\DUrole{doc}{variable}}}} command after converting each
constraint to a logical value. Because exactly one constraint is
allowed per line, having a valid logical expression implies that left
parentheses “(” should only appear before a constraint, and right
parentheses “)” should only appear after a constraint and before any
logical operator.

\sphinxAtStartPar
Once a reaction site has been successfully identified, data structures
within LAMMPS that store bond topology are updated to reflect the
post\sphinxhyphen{}reacted molecule template. All force fields with fixed bonds,
angles, dihedrals or impropers are supported.

\sphinxAtStartPar
A few capabilities to note:
\begin{enumerate}
\sphinxsetlistlabels{\arabic}{enumi}{enumii}{(}{)}%
\item {} 
\sphinxAtStartPar
You may specify as many \sphinxstyleemphasis{react} arguments as desired. For example, you
could break down a complicated reaction mechanism into several reaction
steps, each defined by its own \sphinxstyleemphasis{react} argument.

\item {} 
\sphinxAtStartPar
While typically a bond is formed or removed between the initiator atoms
specified in the pre\sphinxhyphen{}reacted molecule template, this is not required.

\item {} 
\sphinxAtStartPar
By reversing the order of the pre\sphinxhyphen{} and post\sphinxhyphen{}reacted molecule templates in
another \sphinxstyleemphasis{react} argument, you can allow for the possibility of one or
more reverse reactions.

\end{enumerate}

\sphinxAtStartPar
The optional keywords deal with the probability of a given reaction
occurring as well as the stable equilibration of each reaction site as
it occurs.

\sphinxAtStartPar
The \sphinxstyleemphasis{prob} keyword can affect whether or not an eligible reaction
actually occurs. The fraction setting must be a value between 0.0 and
1.0, and can be specified with an equal\sphinxhyphen{}style {\hyperref[\detokenize{variable::doc}]{\sphinxcrossref{\DUrole{doc}{variable}}}}.
A uniform random number between 0.0 and 1.0 is generated and the
eligible reaction only occurs if the random number is less than the
fraction. Up to \(N\) reactions are permitted to occur, as optionally
specified by the \sphinxstyleemphasis{max\_rxn} keyword.

\sphinxAtStartPar
\DUrole{versionmodified,added}{Added in version 22Dec2022.}

\sphinxAtStartPar
The \sphinxstyleemphasis{rate\_limit} keyword can enforce an upper limit on the overall
rate of the reaction. The number of reaction occurrences is limited to
Nlimit within an interval of Nsteps timesteps. No reactions are
permitted to occur within the first Nsteps timesteps of the first run
after reading a data file. Nlimit can be specified with an equal\sphinxhyphen{}style
{\hyperref[\detokenize{variable::doc}]{\sphinxcrossref{\DUrole{doc}{variable}}}}.

\sphinxAtStartPar
The \sphinxstyleemphasis{stabilize\_steps} keyword allows for the specification of how many
time steps a reaction site is stabilized before being returned to the
overall system thermostat. In order to produce the most physical
behavior, this “reaction site equilibration time” should be tuned to
be as small as possible while retaining stability for a given system
or reaction step. After a limited number of case studies, this number
has been set to a default of 60 time steps. Ideally, it should be
individually tuned for each fix reaction step. Note that in some
situations, decreasing rather than increasing this parameter will
result in an increase in stability.

\sphinxAtStartPar
The \sphinxstyleemphasis{custom\_charges} keyword can be used to specify which atoms’
atomic charges are updated. When the value is set to \sphinxstyleemphasis{no}, all atomic
charges are updated to those specified by the post\sphinxhyphen{}reaction template
(default). Otherwise, the value should be the name of a molecule
fragment defined in the pre\sphinxhyphen{}reaction molecule template. In this case,
only the atomic charges of atoms in the molecule fragment are updated.

\sphinxAtStartPar
\DUrole{versionmodified,added}{Added in version 22Dec2022.}

\sphinxAtStartPar
The \sphinxstyleemphasis{rescale\_charges} keyword can be used to ensure the total charge
of the system does not change as reactions occur. When the argument is
set to \sphinxstyleemphasis{yes}, a fixed value is added to the charges of post\sphinxhyphen{}reaction
atoms such that their total charge equals that of the pre\sphinxhyphen{}reaction
site. If only a subset of atomic charges are updated via the
\sphinxstyleemphasis{custom\_charges} keyword, this rescaling is applied to the subset.
This keyword could be useful for systems that contain different
molecules with the same reactive site, if the partial charges on the
reaction site vary from molecule to molecule, or when removing
reaction by\sphinxhyphen{}products.

\sphinxAtStartPar
The \sphinxstyleemphasis{molecule} keyword can be used to force the reaction to be
intermolecular, intramolecular or either. When the value is set to
\sphinxstyleemphasis{off}, molecule IDs are not considered when searching for reactions
(default). When the value is set to \sphinxstyleemphasis{inter}, the initiator atoms must
have different molecule IDs in order to be considered for the
reaction. When the value is set to \sphinxstyleemphasis{intra}, only initiator atoms with
the same molecule ID are considered for the reaction.

\sphinxAtStartPar
A few other considerations:

\sphinxAtStartPar
Optionally, you can enforce additional behaviors on reacting atoms.
For example, it may be beneficial to force reacting atoms to remain at
a certain temperature. For this, you can use the internally\sphinxhyphen{}created
dynamic group named “bond\_react\_MASTER\_group”, which consists of all
atoms currently involved in a reaction. For example, adding the
following command would add an additional thermostat to the group of
all currently\sphinxhyphen{}reacting atoms:

\begin{sphinxVerbatim}[commandchars=\\\{\}]
\PYG{k}{fix }\PYG{n+nv+nvIdentifier}{1}\PYG{+w}{ }\PYG{n+nv+nvIdentifier}{bond\PYGZus{}react\PYGZus{}MASTER\PYGZus{}group}\PYG{+w}{ }\PYG{n}{temp}\PYG{o}{/}\PYG{n}{rescale}\PYG{+w}{ }\PYG{l+m}{1}\PYG{+w}{ }\PYG{l+m}{300}\PYG{+w}{ }\PYG{l+m}{300}\PYG{+w}{ }\PYG{l+m}{10}\PYG{+w}{ }\PYG{l+m}{1}
\end{sphinxVerbatim}

\begin{sphinxadmonition}{note}{Note:}
\sphinxAtStartPar
This command must be added after the fix bond/react command, and
will apply to all reactions.
\end{sphinxadmonition}

\sphinxAtStartPar
Computationally, each time step this fix is invoked, it loops over
neighbor lists (for bond\sphinxhyphen{}forming reactions) and computes distances
between pairs of atoms in the list. It also communicates between
neighboring processors to coordinate which bonds are created and/or
removed. All of these operations increase the cost of a time step. Thus,
you should be cautious about invoking this fix too frequently.

\sphinxAtStartPar
You can dump out snapshots of the current bond topology via the
{\hyperref[\detokenize{dump::doc}]{\sphinxcrossref{\DUrole{doc}{dump local}}}} command.


\bigskip\hrule\bigskip



\subsection{Restart, fix\_modify, output, run start/stop, minimize info}
\label{\detokenize{fix_bond_react:restart-fix-modify-output-run-start-stop-minimize-info}}
\sphinxAtStartPar
Cumulative reaction counts for each reaction are written to {\hyperref[\detokenize{restart::doc}]{\sphinxcrossref{\DUrole{doc}{binary restart files}}}}.
These values are associated with the reaction name (react\sphinxhyphen{}ID).
Additionally, internally\sphinxhyphen{}created per\sphinxhyphen{}atom properties are stored to
allow for smooth restarts. None of the {\hyperref[\detokenize{fix_modify::doc}]{\sphinxcrossref{\DUrole{doc}{fix\_modify}}}}
options are relevant to this fix.

\sphinxAtStartPar
This fix computes one statistic for each \sphinxstyleemphasis{react} argument that it
stores in a global vector, of length (number of react arguments), that
can be accessed by various {\hyperref[\detokenize{Howto_output::doc}]{\sphinxcrossref{\DUrole{doc}{output commands}}}}. The
vector values calculated by this fix are “intensive”.

\sphinxAtStartPar
There is one quantity in the global vector for each \sphinxstyleemphasis{react} argument:
\begin{enumerate}
\sphinxsetlistlabels{\arabic}{enumi}{enumii}{(}{)}%
\item {} 
\sphinxAtStartPar
cumulative number of reactions that occurred

\end{enumerate}

\sphinxAtStartPar
No parameter of this fix can be used with the \sphinxstyleemphasis{start/stop} keywords
of the {\hyperref[\detokenize{run::doc}]{\sphinxcrossref{\DUrole{doc}{run}}}} command.  This fix is not invoked during {\hyperref[\detokenize{minimize::doc}]{\sphinxcrossref{\DUrole{doc}{energy minimization}}}}.

\sphinxAtStartPar
When fix bond/react is “{\hyperref[\detokenize{unfix::doc}]{\sphinxcrossref{\DUrole{doc}{unfixed}}}}”, all internally\sphinxhyphen{}created
groups are deleted. Therefore, fix bond/react can only be unfixed after
unfixing all other fixes that use any group created by fix bond/react.


\subsection{Restrictions}
\label{\detokenize{fix_bond_react:restrictions}}
\sphinxAtStartPar
This fix is part of the REACTION package.  It is only enabled if
LAMMPS was built with that package.  See the
{\hyperref[\detokenize{Build_package::doc}]{\sphinxcrossref{\DUrole{doc}{Build package}}}} page for more info.


\subsection{Related commands}
\label{\detokenize{fix_bond_react:related-commands}}
\sphinxAtStartPar
{\hyperref[\detokenize{fix_bond_create::doc}]{\sphinxcrossref{\DUrole{doc}{fix bond/create}}}},
{\hyperref[\detokenize{fix_bond_break::doc}]{\sphinxcrossref{\DUrole{doc}{fix bond/break}}}},
{\hyperref[\detokenize{fix_bond_swap::doc}]{\sphinxcrossref{\DUrole{doc}{fix bond/swap}}}},
{\hyperref[\detokenize{dump::doc}]{\sphinxcrossref{\DUrole{doc}{dump local}}}}, {\hyperref[\detokenize{special_bonds::doc}]{\sphinxcrossref{\DUrole{doc}{special\_bonds}}}}


\subsection{Default}
\label{\detokenize{fix_bond_react:default}}
\sphinxAtStartPar
The option defaults are stabilization = no, prob = 1.0, stabilize\_steps = 60,
reset\_mol\_ids = yes, custom\_charges = no, molecule = off, modify\_create = \sphinxstyleemphasis{fit all}


\bigskip\hrule\bigskip

\phantomsection\label{\detokenize{fix_bond_react:gissinger}}
\sphinxAtStartPar
\sphinxstylestrong{(Gissinger2017)} Gissinger, Jensen and Wise, Polymer, 128, 211\sphinxhyphen{}217 (2017).

\phantomsection\label{\detokenize{fix_bond_react:gissinger2020}}
\sphinxAtStartPar
\sphinxstylestrong{(Gissinger2020)} Gissinger, Jensen and Wise, Macromolecules, 53, 22, 9953\sphinxhyphen{}9961 (2020).

\phantomsection\label{\detokenize{fix_bond_react:gissinger2024}}
\sphinxAtStartPar
\sphinxstylestrong{(Gissinger2024)} Gissinger, Jensen and Wise, Computer Physics Communications, 304, 109287 (2024).

\sphinxstepscope

\index{fix bond/swap@\spxentry{fix bond/swap}}\ignorespaces 

\section{fix bond/swap command}
\label{\detokenize{fix_bond_swap:fix-bond-swap-command}}\label{\detokenize{fix_bond_swap:index-0}}\label{\detokenize{fix_bond_swap::doc}}

\subsection{Syntax}
\label{\detokenize{fix_bond_swap:syntax}}
\begin{sphinxVerbatim}[commandchars=\\\{\}]
\PYG{k}{fix }\PYG{n+nv+nvIdentifier}{ID}\PYG{+w}{ }\PYG{n+nv+nvIdentifier}{group\PYGZhy{}ID}\PYG{+w}{ }\PYG{n}{bond}\PYG{o}{/}\PYG{n}{swap}\PYG{+w}{ }\PYG{n}{Nevery}\PYG{+w}{ }\PYG{n}{fraction}\PYG{+w}{ }\PYG{n}{cutoff}\PYG{+w}{ }\PYG{n}{seed}
\end{sphinxVerbatim}
\begin{itemize}
\item {} 
\sphinxAtStartPar
ID, group\sphinxhyphen{}ID are documented in {\hyperref[\detokenize{fix::doc}]{\sphinxcrossref{\DUrole{doc}{fix}}}} command

\item {} 
\sphinxAtStartPar
bond/swap = style name of this fix command

\item {} 
\sphinxAtStartPar
Nevery = attempt bond swapping every this many steps

\item {} 
\sphinxAtStartPar
fraction = fraction of group atoms to consider for swapping

\item {} 
\sphinxAtStartPar
cutoff = distance at which swapping will be considered (distance units)

\item {} 
\sphinxAtStartPar
seed = random \# seed (positive integer)

\end{itemize}


\subsection{Examples}
\label{\detokenize{fix_bond_swap:examples}}
\begin{sphinxVerbatim}[commandchars=\\\{\}]
\PYG{k}{fix }\PYG{n+nv+nvIdentifier}{1}\PYG{+w}{ }\PYG{n+nv+nvIdentifier}{all}\PYG{+w}{ }\PYG{n}{bond}\PYG{o}{/}\PYG{n}{swap}\PYG{+w}{ }\PYG{l+m}{50}\PYG{+w}{ }\PYG{l+m}{0.5}\PYG{+w}{ }\PYG{l+m}{1.3}\PYG{+w}{ }\PYG{l+m}{598934}
\end{sphinxVerbatim}


\subsection{Description}
\label{\detokenize{fix_bond_swap:description}}
\sphinxAtStartPar
In a simulation of polymer chains this command attempts to swap a pair
of bonds, as illustrated below.  This is done via Monte Carlo rules
using the Boltzmann acceptance criterion, typically with the goal of
equilibrating the polymer system more quickly.  This fix is designed
for use with idealized bead\sphinxhyphen{}spring polymer chains where each polymer
is a linear chain of monomers, but LAMMPS does not check that is the
case for your system.

\sphinxAtStartPar
Here are two use cases for this fix.

\sphinxAtStartPar
The first use case is for swapping bonds on two different chains,
effectively grafting the end of one chain onto the other chain and
vice versa.  The purpose is to equilibrate the polymer chain
conformations more rapidly than dynamics alone would do it, by
enabling instantaneous large conformational changes in a dense polymer
melt.  The polymer chains should thus more rapidly converge to the
proper end\sphinxhyphen{}to\sphinxhyphen{}end distances and radii of gyration.

\sphinxAtStartPar
A schematic of the kinds of bond swaps that can occur in this use case
is shown here:

\noindent{\hspace*{\fill}\sphinxincludegraphics{{bondswap}.jpg}\hspace*{\fill}}

\sphinxAtStartPar
On the left, the red and blue chains have two monomers A1 and B1 close
to each other, which are currently bonded to monomers A2 and B2
respectively within their own chains.  The bond swap operation will
attempt to delete the A1\sphinxhyphen{}A2 and B1\sphinxhyphen{}B2 bonds and replace them with
A1\sphinxhyphen{}B2 and B1\sphinxhyphen{}A2 bonds.  If the swap is energetically favorable, the
two chains on the right are the result and each polymer chain has
undergone a dramatic conformational change.  This reference,
{\hyperref[\detokenize{fix_bond_swap:sides}]{\sphinxcrossref{\DUrole{std,std-ref}{(Sides)}}}} provides more details on the algorithm’s
effectiveness for this use case.

\sphinxAtStartPar
The second use case is a collection of polymer chains with some
fraction of their sites identified as “sticker” sites.  Initially each
polymer chain is isolated from the others in a topological sense, and
there is an intra\sphinxhyphen{}chain bond between every pair of sticker sites on
the same chain.  Over time, bonds swap so that inter\sphinxhyphen{}molecular sticker
bonds are created.  This models a vitrification\sphinxhyphen{}style process whereby
the polymer chains all become interconnected.  For this use case, if
angles are defined they should not include bonds between sticker
sites.

\begin{sphinxadmonition}{note}{Note:}
\sphinxAtStartPar
For the sticker site model, you should set the newton flag for
bonds to “off”, via the {\hyperref[\detokenize{newton::doc}]{\sphinxcrossref{\DUrole{doc}{newton on off}}}} command (“on”
is the default for the 2nd argument).  This is to ensure
appropriate randomness in bond selection because the I,J bond will
be stored by both atom I and atom J.  LAMMPS cannot check for this,
because it is not aware that a sticker site model is being used.
\end{sphinxadmonition}


\bigskip\hrule\bigskip


\sphinxAtStartPar
The bond swapping operation is invoked once every \sphinxstyleemphasis{Nevery} timesteps.
If any bond in the entire system is swapped, a re\sphinxhyphen{}build of the
neighbor lists is triggered, since a swap alters the list of which
neighbors are considered for pairwise interaction.  At each
invocation, each processor considers a random specified \sphinxstyleemphasis{fraction} of
its atoms as potential swapping monomers for this timestep.  Choosing
a small \sphinxstyleemphasis{fraction} value can reduce the likelihood of a reverse swap
occurring soon after an initial swap.

\sphinxAtStartPar
For each monomer A1, its neighbors are looped over as B1 monomers.
For each A1,B1 an additional double loop of bond partners A2 of A1,
and bond partners B2 of B1 a is performed.  For each pair of A1\sphinxhyphen{}A2 and
B1\sphinxhyphen{}B2 bonds to be eligible for swapping, the following 4 criteria must
be met:
\begin{enumerate}
\sphinxsetlistlabels{\arabic}{enumi}{enumii}{}{.}%
\item {} 
\sphinxAtStartPar
All 4 monomers must be in the fix group.

\item {} 
\sphinxAtStartPar
All 4 monomers must be owned by the processor (not ghost atoms).
This ensures that another processor does not attempt to swap bonds
involving the same atoms on the same timestep.  Note that this also
means that bond pairs which straddle processor boundaries are not
eligible for swapping on this step.

\item {} 
\sphinxAtStartPar
The distances between 4 pairs of atoms \textendash{} (A1,A2), (B1,B2), (A1,B2),
(B1,A2) \textendash{} must all be less than the specified \sphinxstyleemphasis{cutoff}.

\item {} 
\sphinxAtStartPar
The molecule IDs of A1 and B1 must be the same (see below).

\end{enumerate}

\sphinxAtStartPar
If an eligible B1 partner is found, the energy change due to swapping
the two bonds is computed.  This includes changes in pairwise, bond, and
angle energies due to the altered connectivity of the 2 chains.
Dihedral and improper interactions are not allowed to be defined when
this fix is used.

\sphinxAtStartPar
If the energy decreases due to the swap operation, the bond swap is
accepted.  If the energy increases it is accepted with probability
exp(\sphinxhyphen{}delta/kT) where delta is the increase in energy, k is the
Boltzmann constant, and T is the current temperature of the system.

\begin{sphinxadmonition}{note}{Note:}
\sphinxAtStartPar
Whether the swap is accepted or rejected, no other swaps are
attempted by this processor on this timestep.  No other eligible
4\sphinxhyphen{}tuples of atoms are considered.  This means that each processor
will perform either a single swap or none on timesteps this fix is
invoked.
\end{sphinxadmonition}


\bigskip\hrule\bigskip


\sphinxAtStartPar
The criterion for matching molecule IDs is how the first use case
described above can be simulated while conserving chain lengths.  This
is done by setting up the molecule IDs for the polymer chains in a
specific way, typically in the data file, read by the {\hyperref[\detokenize{read_data::doc}]{\sphinxcrossref{\DUrole{doc}{read\_data}}}} command.

\sphinxAtStartPar
Consider a system of 6\sphinxhyphen{}mer chains.  You have 2 choices.  If the
molecule IDs for monomers on each chain are set to 1,2,3,4,5,6 then
swaps will conserve chain length.  For a particular monomer there will
be only one other monomer on another chain which is a potential swap
partner.  If the molecule IDs for monomers on each chain are set to
1,2,3,3,2,1 then swaps will conserve chain length but swaps will be
able to occur at either end of a chain.  Thus for a particular monomer
there will be 2 possible swap partners on another chain.  In this
scenario, swaps can also occur within a single chain, i.e. the two
ends of a chain swap with each other.

\begin{sphinxadmonition}{note}{Note:}
\sphinxAtStartPar
If your simulation uses molecule IDs in the usual way, where all
monomers on a single chain are assigned the same ID (different for
each chain), then swaps will only occur within the same chain.  If you
assign the same molecule ID to all monomers in all chains then
inter\sphinxhyphen{}chain swaps will occur, but they will not conserve chain length.
Neither of these scenarios is probably what you want for this fix.
\end{sphinxadmonition}

\begin{sphinxadmonition}{note}{Note:}
\sphinxAtStartPar
When a bond swap occurs the image flags of monomers in the new
polymer chains can become inconsistent.  See the {\hyperref[\detokenize{dump::doc}]{\sphinxcrossref{\DUrole{doc}{dump}}}}
command for a discussion of image flags.  This is not an issue for
running dynamics, but can affect calculation of some diagnostic
quantities or the printing of unwrapped coordinates to a dump file.
\end{sphinxadmonition}

\sphinxAtStartPar
For the second use case described above, the molecule IDs for all
sticker sites should be the same.


\bigskip\hrule\bigskip


\sphinxAtStartPar
This fix computes a temperature each time it is invoked for use by the
Boltzmann criterion.  To do this, the fix creates its own compute of
style \sphinxstyleemphasis{temp}, as if this command had been issued:

\begin{sphinxVerbatim}[commandchars=\\\{\}]
\PYG{k}{compute }\PYG{n+nv+nvIdentifier}{fix\PYGZhy{}ID\PYGZus{}temp}\PYG{+w}{ }\PYG{n+nv+nvIdentifier}{all}\PYG{+w}{ }\PYG{n}{temp}
\end{sphinxVerbatim}

\sphinxAtStartPar
See the {\hyperref[\detokenize{compute_temp::doc}]{\sphinxcrossref{\DUrole{doc}{compute temp}}}} command for details.  Note
that the ID of the new compute is the fix\sphinxhyphen{}ID with underscore + “temp”
appended and the group for the new compute is “all”, so that the
temperature of the entire system is used.

\sphinxAtStartPar
Note that this is NOT the compute used by thermodynamic output (see
the {\hyperref[\detokenize{thermo_style::doc}]{\sphinxcrossref{\DUrole{doc}{thermo\_style}}}} command) with ID =
\sphinxstyleemphasis{thermo\_temp}.  This means you can change the attributes of this fix’s
temperature (e.g. its degrees\sphinxhyphen{}of\sphinxhyphen{}freedom) via the {\hyperref[\detokenize{compute_modify::doc}]{\sphinxcrossref{\DUrole{doc}{compute\_modify}}}} command or print this temperature during
thermodynamic output via the {\hyperref[\detokenize{thermo_style::doc}]{\sphinxcrossref{\DUrole{doc}{thermo\_style custom}}}}
command using the appropriate compute\sphinxhyphen{}ID.  It also means that changing
attributes of \sphinxstyleemphasis{thermo\_temp} will have no effect on this fix.


\bigskip\hrule\bigskip



\subsection{Restart, fix\_modify, output, run start/stop, minimize info}
\label{\detokenize{fix_bond_swap:restart-fix-modify-output-run-start-stop-minimize-info}}
\sphinxAtStartPar
No information about this fix is written to {\hyperref[\detokenize{restart::doc}]{\sphinxcrossref{\DUrole{doc}{binary restart files}}}}.  Because the state of the random number generator is not
saved in restart files, this means you cannot do “exact” restarts with
this fix, where the simulation continues on the same as if no restart
had taken place.  However, in a statistical sense, a restarted
simulation should produce the same behavior.  Also note that each
processor generates possible swaps independently of other processors.
Thus if you repeat the same simulation on a different number of
processors, the specific swaps performed will be different.

\sphinxAtStartPar
The {\hyperref[\detokenize{fix_modify::doc}]{\sphinxcrossref{\DUrole{doc}{fix\_modify}}}} \sphinxstyleemphasis{temp} option is supported by this
fix.  You can use it to assign a {\hyperref[\detokenize{compute::doc}]{\sphinxcrossref{\DUrole{doc}{compute}}}} you have
defined to this fix which will be used to compute the temperature for
the Boltzmann criterion.

\sphinxAtStartPar
This fix computes two statistical quantities as a global 2\sphinxhyphen{}vector of
output, which can be accessed by various {\hyperref[\detokenize{Howto_output::doc}]{\sphinxcrossref{\DUrole{doc}{output commands}}}}.  The first component of the vector is the cumulative
number of swaps performed by all processors.  The second component of
the vector is the cumulative number of swaps attempted (whether
accepted or rejected).  Note that a swap “attempt” only occurs when
swap partners meeting the criteria described above are found on a
particular timestep.  The vector values calculated by this fix are
“intensive”.

\sphinxAtStartPar
No parameter of this fix can be used with the \sphinxstyleemphasis{start/stop} keywords of
the {\hyperref[\detokenize{run::doc}]{\sphinxcrossref{\DUrole{doc}{run}}}} command.  This fix is not invoked during
{\hyperref[\detokenize{minimize::doc}]{\sphinxcrossref{\DUrole{doc}{energy minimization}}}}.


\subsection{Restrictions}
\label{\detokenize{fix_bond_swap:restrictions}}
\sphinxAtStartPar
This fix is part of the MC package.  It is only enabled if LAMMPS was
built with that package.  See the {\hyperref[\detokenize{Build_package::doc}]{\sphinxcrossref{\DUrole{doc}{Build package}}}}
doc page for more info.

\sphinxAtStartPar
The settings of the “special\_bond” command must be 0,1,1 in order to
use this fix, which is typical of bead\sphinxhyphen{}spring chains with FENE or
harmonic bonds.  This means that pairwise interactions between bonded
atoms are turned off, but are turned on between atoms two or three
hops away along the chain backbone.

\sphinxAtStartPar
Currently, energy changes in dihedral and improper interactions due to
a bond swap are not considered.  Thus a simulation that uses this fix
cannot use a dihedral or improper potential.


\subsection{Related commands}
\label{\detokenize{fix_bond_swap:related-commands}}
\sphinxAtStartPar
{\hyperref[\detokenize{fix_atom_swap::doc}]{\sphinxcrossref{\DUrole{doc}{fix atom/swap}}}}


\subsection{Default}
\label{\detokenize{fix_bond_swap:default}}
\sphinxAtStartPar
none


\bigskip\hrule\bigskip

\phantomsection\label{\detokenize{fix_bond_swap:sides}}
\sphinxAtStartPar
\sphinxstylestrong{(Sides)} Sides, Grest, Stevens, Plimpton, J Polymer Science B, 42,
199\sphinxhyphen{}208 (2004).

\sphinxstepscope

\index{fix box/relax@\spxentry{fix box/relax}}\ignorespaces 

\section{fix box/relax command}
\label{\detokenize{fix_box_relax:fix-box-relax-command}}\label{\detokenize{fix_box_relax:index-0}}\label{\detokenize{fix_box_relax::doc}}

\subsection{Syntax}
\label{\detokenize{fix_box_relax:syntax}}
\begin{sphinxVerbatim}[commandchars=\\\{\}]
\PYG{k}{fix }\PYG{n+nv+nvIdentifier}{ID}\PYG{+w}{ }\PYG{n+nv+nvIdentifier}{group\PYGZhy{}ID}\PYG{+w}{ }\PYG{n}{box}\PYG{o}{/}\PYG{n}{relax}\PYG{+w}{ }\PYG{n}{keyword}\PYG{+w}{ }\PYG{n}{value}\PYG{+w}{ }\PYG{n}{...}
\end{sphinxVerbatim}
\begin{itemize}
\item {} 
\sphinxAtStartPar
ID, group\sphinxhyphen{}ID are documented in {\hyperref[\detokenize{fix::doc}]{\sphinxcrossref{\DUrole{doc}{fix}}}} command

\item {} 
\sphinxAtStartPar
box/relax = style name of this fix command
\begin{sphinxalltt}
one or more keyword value pairs may be appended
keyword = \sphinxstyleemphasis{iso} or \sphinxstyleemphasis{aniso} or \sphinxstyleemphasis{tri} or \sphinxstyleemphasis{x} or \sphinxstyleemphasis{y} or \sphinxstyleemphasis{z} or \sphinxstyleemphasis{xy} or \sphinxstyleemphasis{yz} or \sphinxstyleemphasis{xz} or \sphinxstyleemphasis{couple} or \sphinxstyleemphasis{nreset} or \sphinxstyleemphasis{vmax} or \sphinxstyleemphasis{dilate} or \sphinxstyleemphasis{scaleyz} or \sphinxstyleemphasis{scalexz} or \sphinxstyleemphasis{scalexy} or \sphinxstyleemphasis{fixedpoint}
  \sphinxstyleemphasis{iso} or \sphinxstyleemphasis{aniso} or \sphinxstyleemphasis{tri} value = Ptarget = desired pressure (pressure units)
  \sphinxstyleemphasis{x} or \sphinxstyleemphasis{y} or \sphinxstyleemphasis{z} or \sphinxstyleemphasis{xy} or \sphinxstyleemphasis{yz} or \sphinxstyleemphasis{xz} value = Ptarget = desired pressure (pressure units)
  \sphinxstyleemphasis{couple} = \sphinxstyleemphasis{none} or \sphinxstyleemphasis{xyz} or \sphinxstyleemphasis{xy} or \sphinxstyleemphasis{yz} or \sphinxstyleemphasis{xz}
  \sphinxstyleemphasis{nreset} value = reset reference cell every this many minimizer iterations
  \sphinxstyleemphasis{vmax} value = fraction = max allowed volume change in one iteration
  \sphinxstyleemphasis{dilate} value = \sphinxstyleemphasis{all} or \sphinxstyleemphasis{partial}
  \sphinxstyleemphasis{scaleyz} value = \sphinxstyleemphasis{yes} or \sphinxstyleemphasis{no} = scale yz with lz
  \sphinxstyleemphasis{scalexz} value = \sphinxstyleemphasis{yes} or \sphinxstyleemphasis{no} = scale xz with lz
  \sphinxstyleemphasis{scalexy} value = \sphinxstyleemphasis{yes} or \sphinxstyleemphasis{no} = scale xy with ly
  \sphinxstyleemphasis{fixedpoint} values = x y z
    x,y,z = perform relaxation dilation/contraction around this point (distance units)
\end{sphinxalltt}

\end{itemize}


\subsection{Examples}
\label{\detokenize{fix_box_relax:examples}}
\begin{sphinxVerbatim}[commandchars=\\\{\}]
\PYG{k}{fix }\PYG{n+nv+nvIdentifier}{1}\PYG{+w}{ }\PYG{n+nv+nvIdentifier}{all}\PYG{+w}{ }\PYG{n}{box}\PYG{o}{/}\PYG{n}{relax}\PYG{+w}{ }\PYG{n}{iso}\PYG{+w}{ }\PYG{l+m}{0.0}\PYG{+w}{ }\PYG{n}{vmax}\PYG{+w}{ }\PYG{l+m}{0.001}
\PYG{k}{fix }\PYG{n+nv+nvIdentifier}{2}\PYG{+w}{ }\PYG{n+nv+nvIdentifier}{water}\PYG{+w}{ }\PYG{n}{box}\PYG{o}{/}\PYG{n}{relax}\PYG{+w}{ }\PYG{n}{aniso}\PYG{+w}{ }\PYG{l+m}{0.0}\PYG{+w}{ }\PYG{n}{dilate}\PYG{+w}{ }\PYG{n}{partial}
\PYG{k}{fix }\PYG{n+nv+nvIdentifier}{2}\PYG{+w}{ }\PYG{n+nv+nvIdentifier}{ice}\PYG{+w}{ }\PYG{n}{box}\PYG{o}{/}\PYG{n}{relax}\PYG{+w}{ }\PYG{n}{tri}\PYG{+w}{ }\PYG{l+m}{0.0}\PYG{+w}{ }\PYG{n}{couple}\PYG{+w}{ }\PYG{n}{xy}\PYG{+w}{ }\PYG{n}{nreset}\PYG{+w}{ }\PYG{l+m}{100}
\end{sphinxVerbatim}


\subsection{Description}
\label{\detokenize{fix_box_relax:description}}
\sphinxAtStartPar
Apply an external pressure or stress tensor to the simulation box
during an {\hyperref[\detokenize{minimize::doc}]{\sphinxcrossref{\DUrole{doc}{energy minimization}}}}.  This allows the box
size and shape to vary during the iterations of the minimizer so that
the final configuration will be both an energy minimum for the
potential energy of the atoms, and the system pressure tensor will be
close to the specified external tensor.  Conceptually, specifying a
positive pressure is like squeezing on the simulation box; a negative
pressure typically allows the box to expand.


\bigskip\hrule\bigskip


\sphinxAtStartPar
The external pressure tensor is specified using one or more of the
\sphinxstyleemphasis{iso}, \sphinxstyleemphasis{aniso}, \sphinxstyleemphasis{tri}, \sphinxstyleemphasis{x}, \sphinxstyleemphasis{y}, \sphinxstyleemphasis{z}, \sphinxstyleemphasis{xy}, \sphinxstyleemphasis{xz}, \sphinxstyleemphasis{yz}, and \sphinxstyleemphasis{couple}
keywords.  These keywords give you the ability to specify all 6
components of an external stress tensor, and to couple various of
these components together so that the dimensions they represent are
varied together during the minimization.

\sphinxAtStartPar
Orthogonal simulation boxes have 3 adjustable dimensions (x,y,z).
Triclinic (non\sphinxhyphen{}orthogonal) simulation boxes have 6 adjustable
dimensions (x,y,z,xy,xz,yz).  The {\hyperref[\detokenize{create_box::doc}]{\sphinxcrossref{\DUrole{doc}{create\_box}}}}, {\hyperref[\detokenize{read_data::doc}]{\sphinxcrossref{\DUrole{doc}{read data}}}}, and {\hyperref[\detokenize{read_restart::doc}]{\sphinxcrossref{\DUrole{doc}{read\_restart}}}} commands
specify whether the simulation box is orthogonal or non\sphinxhyphen{}orthogonal
(triclinic) and explain the meaning of the xy,xz,yz tilt factors.

\sphinxAtStartPar
The target pressures \sphinxstyleemphasis{Ptarget} for each of the 6 components of the
stress tensor can be specified independently via the \sphinxstyleemphasis{x}, \sphinxstyleemphasis{y}, \sphinxstyleemphasis{z},
\sphinxstyleemphasis{xy}, \sphinxstyleemphasis{xz}, \sphinxstyleemphasis{yz} keywords, which correspond to the 6 simulation box
dimensions.  For example, if the \sphinxstyleemphasis{y} keyword is used, the y\sphinxhyphen{}box length
will change during the minimization.  If the \sphinxstyleemphasis{xy} keyword is used, the
xy tilt factor will change.  A box dimension will not change if that
component is not specified.

\sphinxAtStartPar
Note that in order to use the \sphinxstyleemphasis{xy}, \sphinxstyleemphasis{xz}, or \sphinxstyleemphasis{yz} keywords, the
simulation box must be triclinic, even if its initial tilt factors are
0.0.

\sphinxAtStartPar
When the size of the simulation box changes, all atoms are re\sphinxhyphen{}scaled
to new positions, unless the keyword \sphinxstyleemphasis{dilate} is specified with a
value of \sphinxstyleemphasis{partial}, in which case only the atoms in the fix group are
re\sphinxhyphen{}scaled.  This can be useful for leaving the coordinates of atoms in
a solid substrate unchanged and controlling the pressure of a
surrounding fluid.

\sphinxAtStartPar
The \sphinxstyleemphasis{scaleyz}, \sphinxstyleemphasis{scalexz}, and \sphinxstyleemphasis{scalexy} keywords control whether or
not the corresponding tilt factors are scaled with the associated box
dimensions when relaxing triclinic periodic cells.  The default
values \sphinxstyleemphasis{yes} will turn on scaling, which corresponds to adjusting the
linear dimensions of the cell while preserving its shape.  Choosing
\sphinxstyleemphasis{no} ensures that the tilt factors are not scaled with the box
dimensions. See below for restrictions and default values in different
situations. In older versions of LAMMPS, scaling of tilt factors was
not performed. The old behavior can be recovered by setting all three
scale keywords to \sphinxstyleemphasis{no}.

\sphinxAtStartPar
The \sphinxstyleemphasis{fixedpoint} keyword specifies the fixed point for cell relaxation.
By default, it is the center of the box.  Whatever point is
chosen will not move during the simulation.  For example, if the lower
periodic boundaries pass through (0,0,0), and this point is provided
to \sphinxstyleemphasis{fixedpoint}, then the lower periodic boundaries will remain at
(0,0,0), while the upper periodic boundaries will move twice as
far. In all cases, the particle positions at each iteration are
unaffected by the chosen value, except that all particles are
displaced by the same amount, different on each iteration.

\begin{sphinxadmonition}{note}{Note:}
\sphinxAtStartPar
Applying an external pressure to tilt dimensions \sphinxstyleemphasis{xy}, \sphinxstyleemphasis{xz}, \sphinxstyleemphasis{yz}
can sometimes result in arbitrarily large values of the tilt factors,
i.e. a dramatically deformed simulation box.  This typically indicates
that there is something badly wrong with how the simulation was
constructed.  The two most common sources of this error are applying a
shear stress to a liquid system or specifying an external shear stress
tensor that exceeds the yield stress of the solid.  In either case the
minimization may converge to a bogus conformation or not converge at
all.  Also note that if the box shape tilts to an extreme shape,
LAMMPS will run less efficiently, due to the large volume of
communication needed to acquire ghost atoms around a processor’s
irregular\sphinxhyphen{}shaped subdomain.  For extreme values of tilt, LAMMPS may
also lose atoms and generate an error.
\end{sphinxadmonition}

\begin{sphinxadmonition}{note}{Note:}
\sphinxAtStartPar
Performing a minimization with this fix is not a mathematically
well\sphinxhyphen{}defined minimization problem.  This is because the objective
function being minimized changes if the box size/shape changes.  In
practice this means the minimizer can get “stuck” before you have
reached the desired tolerance.  The solution to this is to restart the
minimizer from the new adjusted box size/shape, since that creates a
new objective function valid for the new box size/shape.  Repeat as
necessary until the box size/shape has reached its new equilibrium.
\end{sphinxadmonition}


\bigskip\hrule\bigskip


\sphinxAtStartPar
The \sphinxstyleemphasis{couple} keyword allows two or three of the diagonal components of
the pressure tensor to be “coupled” together.  The value specified
with the keyword determines which are coupled.  For example, \sphinxstyleemphasis{xz}
means the \sphinxstyleemphasis{Pxx} and \sphinxstyleemphasis{Pzz} components of the stress tensor are coupled.
\sphinxstyleemphasis{Xyz} means all 3 diagonal components are coupled.  Coupling means two
things: the instantaneous stress will be computed as an average of the
corresponding diagonal components, and the coupled box dimensions will
be changed together in lockstep, meaning coupled dimensions will be
dilated or contracted by the same percentage every timestep.  The
\sphinxstyleemphasis{Ptarget} values for any coupled dimensions must be identical.
\sphinxstyleemphasis{Couple xyz} can be used for a 2d simulation; the \sphinxstyleemphasis{z} dimension is
simply ignored.


\bigskip\hrule\bigskip


\sphinxAtStartPar
The \sphinxstyleemphasis{iso}, \sphinxstyleemphasis{aniso}, and \sphinxstyleemphasis{tri} keywords are simply shortcuts that are
equivalent to specifying several other keywords together.

\sphinxAtStartPar
The keyword \sphinxstyleemphasis{iso} means couple all 3 diagonal components together when
pressure is computed (hydrostatic pressure), and dilate/contract the
dimensions together.  Using “iso Ptarget” is the same as specifying
these 4 keywords:

\begin{sphinxVerbatim}[commandchars=\\\{\}]
x Ptarget y Ptarget z Ptarget couple xyz
\end{sphinxVerbatim}

\sphinxAtStartPar
The keyword \sphinxstyleemphasis{aniso} means \sphinxstyleemphasis{x}, \sphinxstyleemphasis{y}, and \sphinxstyleemphasis{z} dimensions are controlled
independently using the \sphinxstyleemphasis{Pxx}, \sphinxstyleemphasis{Pyy}, and \sphinxstyleemphasis{Pzz} components of the
stress tensor as the driving forces, and the specified scalar external
pressure.  Using “aniso Ptarget” is the same as specifying these 4
keywords:

\begin{sphinxVerbatim}[commandchars=\\\{\}]
x Ptarget y Ptarget z Ptarget couple none
\end{sphinxVerbatim}

\sphinxAtStartPar
The keyword \sphinxstyleemphasis{tri} means \sphinxstyleemphasis{x}, \sphinxstyleemphasis{y}, \sphinxstyleemphasis{z}, \sphinxstyleemphasis{xy}, \sphinxstyleemphasis{xz}, and \sphinxstyleemphasis{yz} dimensions
are controlled independently using their individual stress components
as the driving forces, and the specified scalar pressure as the
external normal stress.  Using “tri Ptarget” is the same as specifying
these 7 keywords:

\begin{sphinxVerbatim}[commandchars=\\\{\}]
x Ptarget y Ptarget z Ptarget xy 0.0 yz 0.0 xz 0.0 couple none
\end{sphinxVerbatim}


\bigskip\hrule\bigskip


\sphinxAtStartPar
The \sphinxstyleemphasis{vmax} keyword can be used to limit the fractional change in the
volume of the simulation box that can occur in one iteration of the
minimizer.  If the pressure is not settling down during the
minimization this can be because the volume is fluctuating too much.
The specified fraction must be greater than 0.0 and should be \textless{}\textless{} 1.0.
A value of 0.001 means the volume cannot change by more than 1/10 of a
percent in one iteration when \sphinxstyleemphasis{couple xyz} has been specified.  For
any other case it means no linear dimension of the simulation box can
change by more than 1/10 of a percent.


\bigskip\hrule\bigskip


\sphinxAtStartPar
With this fix, the potential energy used by the minimizer is augmented
by an additional energy provided by the fix. The overall objective
function then is:
\begin{equation*}
\begin{split}E = U + P_t \left(V-V_0 \right) + E_{strain}\end{split}
\end{equation*}
\sphinxAtStartPar
where \sphinxstyleemphasis{U} is the system potential energy, \(P_t\) is the desired
hydrostatic pressure, \(V\) and \(V_0\) are the system and reference
volumes, respectively.  \(E_{strain}\) is the strain energy expression
proposed by Parrinello and Rahman {\hyperref[\detokenize{fix_box_relax:parrinello1981}]{\sphinxcrossref{\DUrole{std,std-ref}{(Parrinello1981)}}}}.
Taking derivatives of \sphinxstyleemphasis{E} w.r.t. the box dimensions, and setting these
to zero, we find that at the minimum of the objective function, the
global system stress tensor \sphinxstylestrong{P} will satisfy the relation:
\begin{equation*}
\begin{split}\mathbf P = P_t \mathbf I + {\mathbf S_t} \left( \mathbf h_0^{-1} \right)^t \mathbf h_{0d}\end{split}
\end{equation*}
\sphinxAtStartPar
where \sphinxstylestrong{I} is the identity matrix, \(\mathbf{h_0}\) is the box
dimension tensor of the reference cell, and :\(\mathbf{h_{0d}}\)
is the diagonal part of \(\mathbf{h_0}\). \(\mathbf{S_t}\)
is a symmetric stress tensor that is chosen by LAMMPS
so that the upper\sphinxhyphen{}triangular components of \sphinxstylestrong{P} equal the stress tensor
specified by the user.

\sphinxAtStartPar
This equation only applies when the box dimensions are equal to those
of the reference dimensions. If this is not the case, then the
converged stress tensor will not equal that specified by the user.  We
can resolve this problem by periodically resetting the reference
dimensions. The keyword \sphinxstyleemphasis{nreset} controls how often this is done.  If
this keyword is not used, or is given a value of zero, then the
reference dimensions are set to those of the initial simulation domain
and are never changed. A value of \sphinxstyleemphasis{nstep} means that every \sphinxstyleemphasis{nstep}
minimization steps, the reference dimensions are set to those of the
current simulation domain. Note that resetting the reference
dimensions changes the objective function and gradients, which
sometimes causes the minimization to fail. This can be resolved by
changing the value of \sphinxstyleemphasis{nreset}, or simply continuing the minimization
from a restart file.

\begin{sphinxadmonition}{note}{Note:}
\sphinxAtStartPar
As normally computed, pressure includes a kinetic\sphinxhyphen{} energy or
temperature\sphinxhyphen{}dependent component; see the {\hyperref[\detokenize{compute_pressure::doc}]{\sphinxcrossref{\DUrole{doc}{compute pressure}}}} command.  However, atom velocities are
ignored during a minimization, and the applied pressure(s) specified
with this command are assumed to only be the virial component of the
pressure (the non\sphinxhyphen{}kinetic portion).  Thus if atoms have a non\sphinxhyphen{}zero
temperature and you print the usual thermodynamic pressure, it may not
appear the system is converging to your specified pressure.  The
solution for this is to either (a) zero the velocities of all atoms
before performing the minimization, or (b) make sure you are
monitoring the pressure without its kinetic component.  The latter can
be done by outputting the pressure from the pressure compute this
command creates (see below) or a pressure compute you define yourself.
\end{sphinxadmonition}

\begin{sphinxadmonition}{note}{Note:}
\sphinxAtStartPar
Because pressure is often a very sensitive function of volume,
it can be difficult for the minimizer to equilibrate the system the
desired pressure with high precision, particularly for solids.  Some
techniques that seem to help are (a) use the “min\_modify line
quadratic” option when minimizing with box relaxations, (b) minimize
several times in succession if need be, to drive the pressure closer
to the target pressure, (c) relax the atom positions before relaxing
the box, and (d) relax the box to the target hydrostatic pressure
before relaxing to a target shear stress state. Also note that some
systems (e.g. liquids) will not sustain a non\sphinxhyphen{}hydrostatic applied
pressure, which means the minimizer will not converge.
\end{sphinxadmonition}


\bigskip\hrule\bigskip


\sphinxAtStartPar
This fix computes a temperature and pressure each timestep.  The
temperature is used to compute the kinetic contribution to the
pressure, even though this is subsequently ignored by default.  To do
this, the fix creates its own computes of style “temp” and “pressure”,
as if these commands had been issued:

\begin{sphinxVerbatim}[commandchars=\\\{\}]
\PYG{k}{compute }\PYG{n+nv+nvIdentifier}{fix\PYGZhy{}ID\PYGZus{}temp}\PYG{+w}{ }\PYG{n+nv+nvIdentifier}{group\PYGZhy{}ID}\PYG{+w}{ }\PYG{n}{temp}
\PYG{k}{compute }\PYG{n+nv+nvIdentifier}{fix\PYGZhy{}ID\PYGZus{}press}\PYG{+w}{ }\PYG{n+nv+nvIdentifier}{group\PYGZhy{}ID}\PYG{+w}{ }\PYG{n}{pressure}\PYG{+w}{ }\PYG{n}{fix}\PYG{o}{\PYGZhy{}}\PYG{n}{ID\PYGZus{}temp}\PYG{+w}{ }\PYG{n}{virial}
\end{sphinxVerbatim}

\sphinxAtStartPar
See the {\hyperref[\detokenize{compute_temp::doc}]{\sphinxcrossref{\DUrole{doc}{compute temp}}}} and {\hyperref[\detokenize{compute_pressure::doc}]{\sphinxcrossref{\DUrole{doc}{compute pressure}}}} commands for details.  Note that the
IDs of the new computes are the fix\sphinxhyphen{}ID + underscore + “temp” or fix\_ID
+ underscore + “press”, and the group for the new computes is the same
as the fix group.  Also note that the pressure compute does not
include a kinetic component.

\sphinxAtStartPar
Note that these are NOT the computes used by thermodynamic output (see
the {\hyperref[\detokenize{thermo_style::doc}]{\sphinxcrossref{\DUrole{doc}{thermo\_style}}}} command) with ID = \sphinxstyleemphasis{thermo\_temp}
and \sphinxstyleemphasis{thermo\_press}.  This means you can change the attributes of this
fix’s temperature or pressure via the
{\hyperref[\detokenize{compute_modify::doc}]{\sphinxcrossref{\DUrole{doc}{compute\_modify}}}} command or print this temperature
or pressure during thermodynamic output via the {\hyperref[\detokenize{thermo_style::doc}]{\sphinxcrossref{\DUrole{doc}{thermo\_style custom}}}} command using the appropriate compute\sphinxhyphen{}ID.
It also means that changing attributes of \sphinxstyleemphasis{thermo\_temp} or
\sphinxstyleemphasis{thermo\_press} will have no effect on this fix.


\bigskip\hrule\bigskip



\subsection{Restart, fix\_modify, output, run start/stop, minimize info}
\label{\detokenize{fix_box_relax:restart-fix-modify-output-run-start-stop-minimize-info}}
\sphinxAtStartPar
No information about this fix is written to {\hyperref[\detokenize{restart::doc}]{\sphinxcrossref{\DUrole{doc}{binary restart files}}}}.

\sphinxAtStartPar
The {\hyperref[\detokenize{fix_modify::doc}]{\sphinxcrossref{\DUrole{doc}{fix\_modify}}}} \sphinxstyleemphasis{temp} and \sphinxstyleemphasis{press} options are
supported by this fix.  You can use them to assign a
{\hyperref[\detokenize{compute::doc}]{\sphinxcrossref{\DUrole{doc}{compute}}}} you have defined to this fix which will be used
in its temperature and pressure calculation, as described above.  Note
that as described above, if you assign a pressure compute to this fix
that includes a kinetic energy component it will affect the
minimization, most likely in an undesirable way.

\begin{sphinxadmonition}{note}{Note:}
\sphinxAtStartPar
If both the \sphinxstyleemphasis{temp} and \sphinxstyleemphasis{press} keywords are used in a single
thermo\_modify command (or in two separate commands), then the order in
which the keywords are specified is important.  Note that a {\hyperref[\detokenize{compute_pressure::doc}]{\sphinxcrossref{\DUrole{doc}{pressure compute}}}} defines its own temperature compute as
an argument when it is specified.  The \sphinxstyleemphasis{temp} keyword will override
this (for the pressure compute being used by fix box/relax), but only if the
\sphinxstyleemphasis{temp} keyword comes after the \sphinxstyleemphasis{press} keyword.  If the \sphinxstyleemphasis{temp} keyword
comes before the \sphinxstyleemphasis{press} keyword, then the new pressure compute
specified by the \sphinxstyleemphasis{press} keyword will be unaffected by the \sphinxstyleemphasis{temp}
setting.
\end{sphinxadmonition}

\sphinxAtStartPar
This fix computes a global scalar which can be accessed by various
{\hyperref[\detokenize{Howto_output::doc}]{\sphinxcrossref{\DUrole{doc}{output commands}}}}. The scalar is the pressure\sphinxhyphen{}volume
energy, plus the strain energy, if it exists, as described above.  The
energy values reported at the end of a minimization run under
“Minimization stats” include this energy, and so differ from what
LAMMPS normally reports as potential energy. This fix does not support
the {\hyperref[\detokenize{fix_modify::doc}]{\sphinxcrossref{\DUrole{doc}{fix\_modify}}}} \sphinxstyleemphasis{energy} option, because that would
result in double\sphinxhyphen{}counting of the fix energy in the minimization
energy. Instead, the fix energy can be explicitly added to the
potential energy using one of these two variants:

\begin{sphinxVerbatim}[commandchars=\\\{\}]
\PYG{k}{variable }\PYG{n+nv+nvIdentifier}{emin}\PYG{+w}{ }\PYG{n}{equal}\PYG{+w}{ }\PYG{n}{pe}\PYG{o}{+}\PYG{n}{f\PYGZus{}1}

\PYG{k}{variable }\PYG{n+nv+nvIdentifier}{emin}\PYG{+w}{ }\PYG{n}{equal}\PYG{+w}{ }\PYG{n}{pe}\PYG{o}{+}\PYG{n}{f\PYGZus{}1}\PYG{o}{/}\PYG{n}{atoms}
\end{sphinxVerbatim}

\sphinxAtStartPar
No parameter of this fix can be used with the \sphinxstyleemphasis{start/stop} keywords of
the {\hyperref[\detokenize{run::doc}]{\sphinxcrossref{\DUrole{doc}{run}}}} command.

\sphinxAtStartPar
This fix is invoked during {\hyperref[\detokenize{minimize::doc}]{\sphinxcrossref{\DUrole{doc}{energy minimization}}}}, but
not for the purpose of adding a contribution to the energy or forces
being minimized.  Instead it alters the simulation box geometry as
described above.


\subsection{Restrictions}
\label{\detokenize{fix_box_relax:restrictions}}
\sphinxAtStartPar
Only dimensions that are available can be adjusted by this fix.
Non\sphinxhyphen{}periodic dimensions are not available.  \sphinxstyleemphasis{z}, \sphinxstyleemphasis{xz}, and \sphinxstyleemphasis{yz}, are
not available for 2D simulations. \sphinxstyleemphasis{xy}, \sphinxstyleemphasis{xz}, and \sphinxstyleemphasis{yz} are only
available if the simulation domain is non\sphinxhyphen{}orthogonal.  The
{\hyperref[\detokenize{create_box::doc}]{\sphinxcrossref{\DUrole{doc}{create\_box}}}}, {\hyperref[\detokenize{read_data::doc}]{\sphinxcrossref{\DUrole{doc}{read data}}}}, and
{\hyperref[\detokenize{read_restart::doc}]{\sphinxcrossref{\DUrole{doc}{read\_restart}}}} commands specify whether the
simulation box is orthogonal or non\sphinxhyphen{}orthogonal (triclinic) and explain
the meaning of the xy,xz,yz tilt factors.

\sphinxAtStartPar
The \sphinxstyleemphasis{scaleyz yes} and \sphinxstyleemphasis{scalexz yes} keyword/value pairs can not be used
for 2D simulations. \sphinxstyleemphasis{scaleyz yes}, \sphinxstyleemphasis{scalexz yes}, and \sphinxstyleemphasis{scalexy yes} options
can only be used if the second dimension in the keyword is periodic,
and if the tilt factor is not coupled to the barostat via keywords
\sphinxstyleemphasis{tri}, \sphinxstyleemphasis{yz}, \sphinxstyleemphasis{xz}, and \sphinxstyleemphasis{xy}.


\subsection{Related commands}
\label{\detokenize{fix_box_relax:related-commands}}
\sphinxAtStartPar
{\hyperref[\detokenize{fix_nh::doc}]{\sphinxcrossref{\DUrole{doc}{fix npt}}}}, {\hyperref[\detokenize{minimize::doc}]{\sphinxcrossref{\DUrole{doc}{minimize}}}}


\subsection{Default}
\label{\detokenize{fix_box_relax:default}}
\sphinxAtStartPar
The keyword defaults are dilate = all, vmax = 0.0001, nreset = 0.


\bigskip\hrule\bigskip

\phantomsection\label{\detokenize{fix_box_relax:parrinello1981}}
\sphinxAtStartPar
\sphinxstylestrong{(Parrinello1981)} Parrinello and Rahman, J Appl Phys, 52, 7182 (1981).

\sphinxstepscope

\index{fix brownian@\spxentry{fix brownian}}\ignorespaces 
\index{fix brownian/sphere@\spxentry{fix brownian/sphere}}\ignorespaces 
\index{fix brownian/asphere@\spxentry{fix brownian/asphere}}\ignorespaces 

\section{fix brownian command}
\label{\detokenize{fix_brownian:fix-brownian-command}}\label{\detokenize{fix_brownian:index-2}}\label{\detokenize{fix_brownian:index-1}}\label{\detokenize{fix_brownian:index-0}}\label{\detokenize{fix_brownian::doc}}

\section{fix brownian/sphere command}
\label{\detokenize{fix_brownian:fix-brownian-sphere-command}}

\section{fix brownian/asphere command}
\label{\detokenize{fix_brownian:fix-brownian-asphere-command}}

\subsection{Syntax}
\label{\detokenize{fix_brownian:syntax}}
\begin{sphinxVerbatim}[commandchars=\\\{\}]
\PYG{k}{fix }\PYG{n+nv+nvIdentifier}{ID}\PYG{+w}{ }\PYG{n+nv+nvIdentifier}{group\PYGZhy{}ID}\PYG{+w}{ }\PYG{n}{style\PYGZus{}name}\PYG{+w}{ }\PYG{n}{temp}\PYG{+w}{ }\PYG{n}{seed}\PYG{+w}{ }\PYG{n}{keyword}\PYG{+w}{ }\PYG{n}{args}
\end{sphinxVerbatim}
\begin{itemize}
\item {} 
\sphinxAtStartPar
ID, group\sphinxhyphen{}ID are documented in {\hyperref[\detokenize{fix::doc}]{\sphinxcrossref{\DUrole{doc}{fix}}}} command

\item {} 
\sphinxAtStartPar
style\_name = \sphinxstyleemphasis{brownian} or \sphinxstyleemphasis{brownian/sphere} or \sphinxstyleemphasis{brownian/asphere}

\item {} 
\sphinxAtStartPar
temp = temperature

\item {} 
\sphinxAtStartPar
seed = random number generator seed

\item {} 
\sphinxAtStartPar
one or more keyword/value pairs may be appended

\item {} 
\sphinxAtStartPar
keyword = \sphinxstyleemphasis{rng} or \sphinxstyleemphasis{dipole} or \sphinxstyleemphasis{gamma\_r\_eigen} or \sphinxstyleemphasis{gamma\_t\_eigen} or \sphinxstyleemphasis{gamma\_r} or \sphinxstyleemphasis{gamma\_t} or \sphinxstyleemphasis{rotation\_temp} or \sphinxstyleemphasis{planar\_rotation}
\begin{sphinxalltt}
\sphinxstyleemphasis{rng} value = \sphinxstyleemphasis{uniform} or \sphinxstyleemphasis{gaussian} or \sphinxstyleemphasis{none}
  \sphinxstyleemphasis{uniform} = use uniform random number generator
  \sphinxstyleemphasis{gaussian} = use gaussian random number generator
  \sphinxstyleemphasis{none} = turn off noise
\sphinxstyleemphasis{dipole} value = \sphinxstyleemphasis{mux} and \sphinxstyleemphasis{muy} and \sphinxstyleemphasis{muz} for \sphinxstyleemphasis{brownian/asphere}
  \sphinxstyleemphasis{mux}, \sphinxstyleemphasis{muy}, and \sphinxstyleemphasis{muz} = update orientation of dipole having direction (\sphinxstyleemphasis{mux},*muy*,*muz*) in body frame of rigid body
\sphinxstyleemphasis{gamma\_r\_eigen} values = \sphinxstyleemphasis{gr1} and \sphinxstyleemphasis{gr2} and \sphinxstyleemphasis{gr3} for \sphinxstyleemphasis{brownian/asphere}
  \sphinxstyleemphasis{gr1}, \sphinxstyleemphasis{gr2}, and \sphinxstyleemphasis{gr3} = diagonal entries of body frame rotational friction tensor
\sphinxstyleemphasis{gamma\_r} values = \sphinxstyleemphasis{gr} for \sphinxstyleemphasis{brownian/sphere}
  \sphinxstyleemphasis{gr} = magnitude of the (isotropic) rotational friction tensor
\sphinxstyleemphasis{gamma\_t\_eigen} values = \sphinxstyleemphasis{gt1} and \sphinxstyleemphasis{gt2} and \sphinxstyleemphasis{gt3} for \sphinxstyleemphasis{brownian/asphere}
  \sphinxstyleemphasis{gt1}, \sphinxstyleemphasis{gt2}, and \sphinxstyleemphasis{gt3} = diagonal entries of body frame translational friction tensor
\sphinxstyleemphasis{gamma\_t} values = \sphinxstyleemphasis{gt} for \sphinxstyleemphasis{brownian} and \sphinxstyleemphasis{brownian/sphere}
   \sphinxstyleemphasis{gt} = magnitude of the (isotropic) translational friction tensor
\sphinxstyleemphasis{rotation\_temp} values = \sphinxstyleemphasis{T} for \sphinxstyleemphasis{brownian/sphere} and \sphinxstyleemphasis{brownian/asphere}
   \sphinxstyleemphasis{T} = rotation temperature, which can be different then \sphinxstyleemphasis{temp} when out of equilibrium
\sphinxstyleemphasis{planar\_rotation} values = none (constrains rotational diffusion to be in xy plane if in 3D)
\end{sphinxalltt}

\end{itemize}


\subsection{Examples}
\label{\detokenize{fix_brownian:examples}}
\begin{sphinxVerbatim}[commandchars=\\\{\}]
\PYG{k}{fix }\PYG{n+nv+nvIdentifier}{1}\PYG{+w}{ }\PYG{n+nv+nvIdentifier}{all}\PYG{+w}{ }\PYG{n}{brownian}\PYG{+w}{ }\PYG{l+m}{1.0}\PYG{+w}{ }\PYG{l+m}{12908410}\PYG{+w}{ }\PYG{n}{gamma\PYGZus{}t}\PYG{+w}{ }\PYG{l+m}{1.0}
\PYG{k}{fix }\PYG{n+nv+nvIdentifier}{1}\PYG{+w}{ }\PYG{n+nv+nvIdentifier}{all}\PYG{+w}{ }\PYG{n}{brownian}\PYG{+w}{ }\PYG{l+m}{1.0}\PYG{+w}{ }\PYG{l+m}{12908410}\PYG{+w}{ }\PYG{n}{gamma\PYGZus{}t}\PYG{+w}{ }\PYG{l+m}{3.0}\PYG{+w}{ }\PYG{n}{rng}\PYG{+w}{ }\PYG{n}{gaussian}
\PYG{k}{fix }\PYG{n+nv+nvIdentifier}{1}\PYG{+w}{ }\PYG{n+nv+nvIdentifier}{all}\PYG{+w}{ }\PYG{n}{brownian}\PYG{o}{/}\PYG{n}{sphere}\PYG{+w}{ }\PYG{l+m}{1.0}\PYG{+w}{ }\PYG{l+m}{1294019}\PYG{+w}{ }\PYG{n}{gamma\PYGZus{}t}\PYG{+w}{ }\PYG{l+m}{3.0}\PYG{+w}{ }\PYG{n}{gamma\PYGZus{}r}\PYG{+w}{ }\PYG{l+m}{1.0}
\PYG{k}{fix }\PYG{n+nv+nvIdentifier}{1}\PYG{+w}{ }\PYG{n+nv+nvIdentifier}{all}\PYG{+w}{ }\PYG{n}{brownian}\PYG{o}{/}\PYG{n}{sphere}\PYG{+w}{ }\PYG{l+m}{1.0}\PYG{+w}{ }\PYG{l+m}{19581092}\PYG{+w}{ }\PYG{n}{gamma\PYGZus{}t}\PYG{+w}{ }\PYG{l+m}{1.0}\PYG{+w}{ }\PYG{n}{gamma\PYGZus{}r}\PYG{+w}{ }\PYG{l+m}{0.3}\PYG{+w}{  }\PYG{n}{rng}\PYG{+w}{ }\PYG{n}{none}
\PYG{k}{fix }\PYG{n+nv+nvIdentifier}{1}\PYG{+w}{ }\PYG{n+nv+nvIdentifier}{all}\PYG{+w}{ }\PYG{n}{brownian}\PYG{o}{/}\PYG{n}{asphere}\PYG{+w}{ }\PYG{l+m}{1.0}\PYG{+w}{ }\PYG{l+m}{1294019}\PYG{+w}{ }\PYG{n}{gamma\PYGZus{}t\PYGZus{}eigen}\PYG{+w}{ }\PYG{l+m}{1.0}\PYG{+w}{ }\PYG{l+m}{2.0}\PYG{+w}{ }\PYG{l+m}{3.0}\PYG{+w}{ }\PYG{n}{gamma\PYGZus{}r\PYGZus{}eigen}\PYG{+w}{ }\PYG{l+m}{4.0}\PYG{+w}{ }\PYG{l+m}{7.0}\PYG{+w}{ }\PYG{l+m}{8.0}\PYG{+w}{ }\PYG{n}{rng}\PYG{+w}{ }\PYG{n}{gaussian}
\PYG{k}{fix }\PYG{n+nv+nvIdentifier}{1}\PYG{+w}{ }\PYG{n+nv+nvIdentifier}{all}\PYG{+w}{ }\PYG{n}{brownian}\PYG{o}{/}\PYG{n}{asphere}\PYG{+w}{ }\PYG{l+m}{1.0}\PYG{+w}{ }\PYG{l+m}{1294019}\PYG{+w}{ }\PYG{n}{gamma\PYGZus{}t\PYGZus{}eigen}\PYG{+w}{ }\PYG{l+m}{1.0}\PYG{+w}{ }\PYG{l+m}{2.0}\PYG{+w}{ }\PYG{l+m}{3.0}\PYG{+w}{ }\PYG{n}{gamma\PYGZus{}r\PYGZus{}eigen}\PYG{+w}{ }\PYG{l+m}{4.0}\PYG{+w}{ }\PYG{l+m}{7.0}\PYG{+w}{ }\PYG{l+m}{8.0}\PYG{+w}{ }\PYG{n}{dipole}\PYG{+w}{ }\PYG{l+m}{1.0}\PYG{+w}{ }\PYG{l+m}{0.0}\PYG{+w}{ }\PYG{l+m}{0.0}
\end{sphinxVerbatim}


\subsection{Description}
\label{\detokenize{fix_brownian:description}}
\sphinxAtStartPar
Perform Brownian Dynamics time integration to update position, velocity,
dipole orientation (for spheres) and quaternion orientation (for
ellipsoids, with optional dipole update as well) of all particles in the
fix group in each timestep.  Brownian Dynamics uses Newton’s laws of
motion in the limit that inertial forces are negligible compared to
viscous forces. The stochastic equation of motion for the center of mass
positions is
\begin{equation*}
\begin{split}d\mathbf{r} = \boldsymbol{\gamma}_t^{-1}\mathbf{F}dt
+ \sqrt{2k_B T}\boldsymbol{\gamma}_t^{-1/2}d\mathbf{W}_t,\end{split}
\end{equation*}
\sphinxAtStartPar
in the lab\sphinxhyphen{}frame (i.e., \(\boldsymbol{\gamma}_t\) is not diagonal, but
only depends on orientation and so the noise is still additive).

\sphinxAtStartPar
The rotational motion for the spherical and ellipsoidal particles is not
as simple an expression, but is chosen to replicate the Boltzmann
distribution for the case of conservative torques (see {\hyperref[\detokenize{fix_brownian:ilie1}]{\sphinxcrossref{\DUrole{std,std-ref}{(Ilie)}}}} or {\hyperref[\detokenize{fix_brownian:delong1}]{\sphinxcrossref{\DUrole{std,std-ref}{(Delong)}}}}).

\sphinxAtStartPar
For the style \sphinxstyleemphasis{brownian}, only the positions of the particles are
updated. This is therefore suitable for point particle simulations.

\sphinxAtStartPar
For the style \sphinxstyleemphasis{brownian/sphere}, the positions of the particles are
updated, and a dipole slaved to the spherical orientation is also
updated. This style therefore requires the hybrid atom style
{\hyperref[\detokenize{atom_style::doc}]{\sphinxcrossref{\DUrole{doc}{atom\_style dipole}}}} and {\hyperref[\detokenize{atom_style::doc}]{\sphinxcrossref{\DUrole{doc}{atom\_style sphere}}}}. The equation of motion for the dipole is
\begin{equation*}
\begin{split}\boldsymbol{\mu}(t+dt) = \frac{\boldsymbol{\mu}(t) + \boldsymbol{\omega} \times \boldsymbol{\mu}dt
}{|\boldsymbol{\mu}(t) + \boldsymbol{\omega} \times \boldsymbol{\mu}|}\end{split}
\end{equation*}
\sphinxAtStartPar
which correctly reproduces a Boltzmann distribution of orientations and
rotational diffusion moments (see {\hyperref[\detokenize{fix_brownian:ilie1}]{\sphinxcrossref{\DUrole{std,std-ref}{(Ilie)}}}}) when
\begin{equation*}
\begin{split}\boldsymbol{\omega} = \frac{\mathbf{T}}{\gamma_r} + \sqrt{\frac{2 k_B T_{rot}}{\gamma_r}\frac{d\mathbf{W}}{dt}},\end{split}
\end{equation*}
\sphinxAtStartPar
with \(d\mathbf{W}\) being a random number with zero mean and variance \(dt\)
and \(T_{rot}\) is \sphinxstyleemphasis{rotation\_temp}.

\sphinxAtStartPar
For the style \sphinxstyleemphasis{brownian/asphere}, the center of mass positions and the
quaternions of ellipsoidal particles are updated. This fix style is
suitable for equations of motion where the rotational and translational
friction tensors can be diagonalized in a certain (body) reference
frame. In this case, the rotational equation of motion is updated via
the quaternion
\begin{equation*}
\begin{split}\mathbf{q}(t+dt) = \frac{\mathbf{q}(t) + d\mathbf{q}}{\lVert\mathbf{q}(t) + d\mathbf{q}\rVert}\end{split}
\end{equation*}
\sphinxAtStartPar
which correctly reproduces a Boltzmann distribution of orientations and rotational
diffusion moments {[}see {\hyperref[\detokenize{fix_brownian:ilie1}]{\sphinxcrossref{\DUrole{std,std-ref}{(Ilie)}}}}{]} when the quaternion step is given by
\begin{equation*}
\begin{split}d\mathbf{q} = \boldsymbol{\Psi}\boldsymbol{\omega}dt\end{split}
\end{equation*}
\sphinxAtStartPar
where \(\boldsymbol{\Psi}\) has rows \((-q_1,-q_2,-q_3)\),
\((q_0,-q_3,q_2)\), \((q_3,q_0,-q_1)\), and \((-q_2,q_1,q_0)\).
\(\boldsymbol{\omega}\) is evaluated in the body frame of reference where the
friction tensor is diagonal.  See {\hyperref[\detokenize{fix_brownian:delong1}]{\sphinxcrossref{\DUrole{std,std-ref}{(Delong)}}}} for more details of
a similar algorithm.


\bigskip\hrule\bigskip


\begin{sphinxadmonition}{note}{Note:}
\sphinxAtStartPar
This integrator does not by default assume a relationship between the
rotational and translational friction tensors, though such a
relationship should exist in the case of no\sphinxhyphen{}slip boundary conditions
between the particles and the surrounding (implicit) solvent. For example,
in the case of spherical particles, the condition
\(\gamma_t=3\gamma_r/\sigma^2\) must be explicitly accounted for
by setting \sphinxstyleemphasis{gamma\_t} to 3x and \sphinxstyleemphasis{gamma\_r} to x (where \(\sigma\)
is the sphere’s diameter). A similar (though more complex)
relationship holds for ellipsoids and rod\sphinxhyphen{}like particles. The
translational diffusion and rotational diffusion are given by
\sphinxstyleemphasis{temp/gamma\_t} and \sphinxstyleemphasis{rotation\_temp/gamma\_r}.
\end{sphinxadmonition}


\bigskip\hrule\bigskip


\begin{sphinxadmonition}{note}{Note:}
\sphinxAtStartPar
Temperature computation using the {\hyperref[\detokenize{compute_temp::doc}]{\sphinxcrossref{\DUrole{doc}{compute temp}}}}
will not correctly compute the temperature of these overdamped dynamics
since we are explicitly neglecting inertial effects.  Furthermore,
this time integrator does not add the stochastic terms or viscous
terms to the force and/or torques.  Rather, they are just added in to
the equations of motion to update the degrees of freedom.
\end{sphinxadmonition}


\bigskip\hrule\bigskip


\sphinxAtStartPar
If the \sphinxstyleemphasis{rng} keyword is used with the \sphinxstyleemphasis{uniform} value, then the noise
is generated from a uniform distribution (see
{\hyperref[\detokenize{fix_brownian:dunweg7}]{\sphinxcrossref{\DUrole{std,std-ref}{(Dunweg)}}}} for why this works). This is the same method
of noise generation as used in {\hyperref[\detokenize{fix_langevin::doc}]{\sphinxcrossref{\DUrole{doc}{fix\_langevin}}}}.

\sphinxAtStartPar
If the \sphinxstyleemphasis{rng} keyword is used with the \sphinxstyleemphasis{gaussian} value, then the noise
is generated from a Gaussian distribution. Typically this added
complexity is unnecessary, and one should be fine using the \sphinxstyleemphasis{uniform}
value for reasons argued in {\hyperref[\detokenize{fix_brownian:dunweg7}]{\sphinxcrossref{\DUrole{std,std-ref}{(Dunweg)}}}}.

\sphinxAtStartPar
If the \sphinxstyleemphasis{rng} keyword is used with the \sphinxstyleemphasis{none} value, then the noise
terms are set to zero.

\sphinxAtStartPar
The \sphinxstyleemphasis{gamma\_t} keyword sets the (isotropic) translational viscous damping.
Required for (and only compatible with) \sphinxstyleemphasis{brownian} and \sphinxstyleemphasis{brownian/sphere}.
The units of \sphinxstyleemphasis{gamma\_t} are mass/time.

\sphinxAtStartPar
The \sphinxstyleemphasis{gamma\_r} keyword sets the (isotropic) rotational viscous damping.
Required for (and only compatible with) \sphinxstyleemphasis{brownian/sphere}.
The units of \sphinxstyleemphasis{gamma\_r} are mass*length**2/time.

\sphinxAtStartPar
The \sphinxstyleemphasis{gamma\_r\_eigen}, and \sphinxstyleemphasis{gamma\_t\_eigen} keywords are the eigenvalues of
the rotational and viscous damping tensors (having the same units as
their isotropic counterparts). Required for (and only compatible with)
\sphinxstyleemphasis{brownian/asphere}. For a 2D system, the first two values of
\sphinxstyleemphasis{gamma\_r\_eigen} must be \sphinxstyleemphasis{inf} (only rotation in \sphinxstyleemphasis{x}\textendash{}\sphinxstyleemphasis{y} plane), and the third
value of \sphinxstyleemphasis{gamma\_t\_eigen} must be \sphinxstyleemphasis{inf} (only diffusion in the \sphinxstyleemphasis{x}\textendash{}\sphinxstyleemphasis{y} plane).

\sphinxAtStartPar
If the \sphinxstyleemphasis{dipole} keyword is used, then the dipole moments of the particles
are updated as described above. Only compatible with \sphinxstyleemphasis{brownian/asphere}
(as \sphinxstyleemphasis{brownian/sphere} updates dipoles automatically).

\sphinxAtStartPar
If the \sphinxstyleemphasis{rotation\_temp} keyword is used, then the rotational diffusion
will be occur at this prescribed temperature instead of \sphinxstyleemphasis{temp}. Only
compatible with \sphinxstyleemphasis{brownian/sphere} and \sphinxstyleemphasis{brownian/asphere}.

\sphinxAtStartPar
If the \sphinxstyleemphasis{planar\_rotation} keyword is used, then rotation is constrained
to the \sphinxstyleemphasis{x}\textendash{} \sphinxstyleemphasis{y} plane in a 3D simulation. Only compatible with
\sphinxstyleemphasis{brownian/sphere} and \sphinxstyleemphasis{brownian/asphere} in 3D.


\bigskip\hrule\bigskip


\begin{sphinxadmonition}{note}{Note:}
\sphinxAtStartPar
For style \sphinxstyleemphasis{brownian/asphere}, the components \sphinxstyleemphasis{gamma\_t\_eigen} = (x,x,x) and
\sphinxstyleemphasis{gamma\_r\_eigen} = (y,y,y), the dynamics will replicate those of the
\sphinxstyleemphasis{brownian/sphere} style with \sphinxstyleemphasis{gamma\_t} = x and \sphinxstyleemphasis{gamma\_r} = y.
\end{sphinxadmonition}


\bigskip\hrule\bigskip



\subsection{Restart, fix\_modify, output, run start/stop, minimize info}
\label{\detokenize{fix_brownian:restart-fix-modify-output-run-start-stop-minimize-info}}
\sphinxAtStartPar
No information about this fix is written to {\hyperref[\detokenize{restart::doc}]{\sphinxcrossref{\DUrole{doc}{binary restart files}}}}.  No global or per\sphinxhyphen{}atom quantities are stored by this fix for
access by various {\hyperref[\detokenize{Howto_output::doc}]{\sphinxcrossref{\DUrole{doc}{output commands}}}}.

\sphinxAtStartPar
No parameter of this fix can be used with the \sphinxstyleemphasis{start/stop} keywords of
the {\hyperref[\detokenize{run::doc}]{\sphinxcrossref{\DUrole{doc}{run}}}} command.  This fix is not invoked during
{\hyperref[\detokenize{minimize::doc}]{\sphinxcrossref{\DUrole{doc}{energy minimization}}}}.


\subsection{Restrictions}
\label{\detokenize{fix_brownian:restrictions}}
\sphinxAtStartPar
The style \sphinxstyleemphasis{brownian/sphere} fix requires that atoms store torque and
angular velocity (omega) as defined by the {\hyperref[\detokenize{atom_style::doc}]{\sphinxcrossref{\DUrole{doc}{atom\_style sphere}}}} command.  The style \sphinxstyleemphasis{brownian/asphere} fix requires that
atoms store torque and quaternions as defined by the {\hyperref[\detokenize{atom_style::doc}]{\sphinxcrossref{\DUrole{doc}{atom\_style
ellipsoid}}}} command.  If the \sphinxstyleemphasis{dipole} keyword is used, they
must also store a dipole moment as defined by the {\hyperref[\detokenize{atom_style::doc}]{\sphinxcrossref{\DUrole{doc}{atom\_style
dipole}}}} command.

\sphinxAtStartPar
This fix is part of the BROWNIAN package.  It is only enabled if LAMMPS
was built with that package.  See the {\hyperref[\detokenize{Build_package::doc}]{\sphinxcrossref{\DUrole{doc}{Build package}}}} doc page for more info.


\subsection{Related commands}
\label{\detokenize{fix_brownian:related-commands}}
\sphinxAtStartPar
{\hyperref[\detokenize{fix_propel_self::doc}]{\sphinxcrossref{\DUrole{doc}{fix propel/self}}}},
{\hyperref[\detokenize{fix_langevin::doc}]{\sphinxcrossref{\DUrole{doc}{fix langevin}}}}, {\hyperref[\detokenize{fix_nve_sphere::doc}]{\sphinxcrossref{\DUrole{doc}{fix nve/sphere}}}},


\subsection{Default}
\label{\detokenize{fix_brownian:default}}
\sphinxAtStartPar
The default for \sphinxstyleemphasis{rng} is \sphinxstyleemphasis{uniform}. The default for the rotational and
translational friction tensors are the identity tensor.


\bigskip\hrule\bigskip

\phantomsection\label{\detokenize{fix_brownian:ilie1}}
\sphinxAtStartPar
\sphinxstylestrong{(Ilie)} Ilie, Briels, den Otter, Journal of Chemical Physics, 142, 114103 (2015).

\phantomsection\label{\detokenize{fix_brownian:delong1}}
\sphinxAtStartPar
\sphinxstylestrong{(Delong)} Delong, Usabiaga, Donev, Journal of Chemical Physics. 143, 144107 (2015)

\phantomsection\label{\detokenize{fix_brownian:dunweg7}}
\sphinxAtStartPar
\sphinxstylestrong{(Dunweg)} Dunweg and Paul, Int J of Modern Physics C, 2, 817\sphinxhyphen{}27 (1991).

\sphinxstepscope

\index{fix charge/regulation@\spxentry{fix charge/regulation}}\ignorespaces 

\section{fix charge/regulation command}
\label{\detokenize{fix_charge_regulation:fix-charge-regulation-command}}\label{\detokenize{fix_charge_regulation:index-0}}\label{\detokenize{fix_charge_regulation::doc}}

\subsection{Syntax}
\label{\detokenize{fix_charge_regulation:syntax}}
\begin{sphinxVerbatim}[commandchars=\\\{\}]
\PYG{k}{fix }\PYG{n+nv+nvIdentifier}{ID}\PYG{+w}{ }\PYG{n+nv+nvIdentifier}{group\PYGZhy{}ID}\PYG{+w}{ }\PYG{n}{charge}\PYG{o}{/}\PYG{n}{regulation}\PYG{+w}{ }\PYG{n}{cation\PYGZus{}type}\PYG{+w}{ }\PYG{n}{anion\PYGZus{}type}\PYG{+w}{ }\PYG{n}{keyword}\PYG{+w}{ }\PYG{n}{value}\PYG{n+nv}{(}\PYG{n+nv}{s}\PYG{n+nv}{)}
\end{sphinxVerbatim}
\begin{itemize}
\item {} 
\sphinxAtStartPar
ID, group\sphinxhyphen{}ID are documented in fix command

\item {} 
\sphinxAtStartPar
charge/regulation = style name of this fix command

\item {} 
\sphinxAtStartPar
cation\_type = atom type of free cations (integer or type label)

\item {} 
\sphinxAtStartPar
anion\_type = atom type of free anions (integer or type label)

\item {} 
\sphinxAtStartPar
zero or more keyword/value pairs may be appended
\begin{sphinxalltt}
keyword = \sphinxstyleemphasis{pH}, \sphinxstyleemphasis{pKa}, \sphinxstyleemphasis{pKb}, \sphinxstyleemphasis{pIp}, \sphinxstyleemphasis{pIm}, \sphinxstyleemphasis{pKs}, \sphinxstyleemphasis{acid\_type}, \sphinxstyleemphasis{base\_type}, \sphinxstyleemphasis{lunit\_nm}, \sphinxstyleemphasis{temp}, \sphinxstyleemphasis{tempfixid}, \sphinxstyleemphasis{nevery}, \sphinxstyleemphasis{nmc}, \sphinxstyleemphasis{rxd}, \sphinxstyleemphasis{seed}, \sphinxstyleemphasis{tag}, \sphinxstyleemphasis{group}, \sphinxstyleemphasis{onlysalt}, \sphinxstyleemphasis{pmcmoves}
\sphinxstyleemphasis{pH} value = pH of the solution (can be specified as an equal\sphinxhyphen{}style variable)
\sphinxstyleemphasis{pKa} value = acid dissociation constant (in the \sphinxhyphen{}log10 representation)
\sphinxstyleemphasis{pKb} value = base dissociation constant (in the \sphinxhyphen{}log10 representation)
\sphinxstyleemphasis{pIp} value = activity (effective concentration) of free cations (in the \sphinxhyphen{}log10 representation)
\sphinxstyleemphasis{pIm} value = activity (effective concentration) of free anions (in the \sphinxhyphen{}log10 representation)
\sphinxstyleemphasis{pKs} value = solvent self\sphinxhyphen{}dissociation constant (in the \sphinxhyphen{}log10 representation)
\sphinxstyleemphasis{acid\_type} = atom type of acid groups (integer or type label)
\sphinxstyleemphasis{base\_type}  = atom type of base groups (integer or type label)
\sphinxstyleemphasis{lunit\_nm} value = unit length used by LAMMPS (\# in the units of nanometers)
\sphinxstyleemphasis{temp} value = temperature
\sphinxstyleemphasis{tempfixid} value = fix ID of temperature thermostat
\sphinxstyleemphasis{nevery} value = invoke this fix every nevery steps
\sphinxstyleemphasis{nmc} value = number of charge regulation MC moves to attempt every nevery steps
\sphinxstyleemphasis{rxd} value = cutoff distance for acid/base reaction
\sphinxstyleemphasis{seed} value = random \# seed (positive integer)
\sphinxstyleemphasis{tag} value = yes or no (yes: The code assign unique tags to inserted ions; no: The tag of all inserted ions is "0")
\sphinxstyleemphasis{group} value = group\sphinxhyphen{}ID, inserted ions are assigned to group group\sphinxhyphen{}ID. Can be used multiple times to assign inserted ions to multiple groups.
\sphinxstyleemphasis{onlysalt} values = flag charge\_cation charge\_anion.
   flag = yes or no (yes: the fix performs only ion insertion/deletion, no: perform acid/base dissociation and ion insertion/deletion)
   charge\_cation, charge\_anion = value of cation/anion charge, must be an integer (only specify if flag = yes)
\sphinxstyleemphasis{pmcmoves} values = pmcA pmcB pmcI \sphinxhyphen{}  MC move fractions for acid ionization (pmcA), base ionization (pmcB) and free ion exchange (pmcI)
\end{sphinxalltt}

\end{itemize}


\subsection{Examples}
\label{\detokenize{fix_charge_regulation:examples}}
\begin{sphinxVerbatim}[commandchars=\\\{\}]
\PYG{k}{fix }\PYG{n+nv+nvIdentifier}{chareg}\PYG{+w}{ }\PYG{n+nv+nvIdentifier}{all}\PYG{+w}{ }\PYG{n}{charge}\PYG{o}{/}\PYG{n}{regulation}\PYG{+w}{ }\PYG{l+m}{1}\PYG{+w}{ }\PYG{l+m}{2}\PYG{+w}{ }\PYG{n}{acid\PYGZus{}type}\PYG{+w}{ }\PYG{l+m}{3}\PYG{+w}{ }\PYG{n}{base\PYGZus{}type}\PYG{+w}{ }\PYG{l+m}{4}\PYG{+w}{ }\PYG{n}{pKa}\PYG{+w}{ }\PYG{l+m}{5.0}\PYG{+w}{ }\PYG{n}{pKb}\PYG{+w}{ }\PYG{l+m}{6.0}\PYG{+w}{ }\PYG{n}{pH}\PYG{+w}{ }\PYG{l+m}{7.0}\PYG{+w}{ }\PYG{n}{pIp}\PYG{+w}{ }\PYG{l+m}{3.0}\PYG{+w}{ }\PYG{n}{pIm}\PYG{+w}{ }\PYG{l+m}{3.0}\PYG{+w}{ }\PYG{n}{nevery}\PYG{+w}{ }\PYG{l+m}{200}\PYG{+w}{ }\PYG{n}{nmc}\PYG{+w}{ }\PYG{l+m}{200}\PYG{+w}{ }\PYG{n}{seed}\PYG{+w}{ }\PYG{l+m}{123}\PYG{+w}{ }\PYG{n}{tempfixid}\PYG{+w}{ }\PYG{n}{fT}
\PYG{k}{fix }\PYG{n+nv+nvIdentifier}{chareg}\PYG{+w}{ }\PYG{n+nv+nvIdentifier}{all}\PYG{+w}{ }\PYG{n}{charge}\PYG{o}{/}\PYG{n}{regulation}\PYG{+w}{ }\PYG{l+m}{1}\PYG{+w}{ }\PYG{l+m}{2}\PYG{+w}{ }\PYG{n}{pIp}\PYG{+w}{ }\PYG{l+m}{3}\PYG{+w}{ }\PYG{n}{pIm}\PYG{+w}{ }\PYG{l+m}{3}\PYG{+w}{ }\PYG{n}{onlysalt}\PYG{+w}{ }\PYG{n}{yes}\PYG{+w}{ }\PYG{l+m}{2}\PYG{+w}{ }\PYG{o}{\PYGZhy{}}\PYG{l+m}{1}\PYG{+w}{ }\PYG{n}{seed}\PYG{+w}{ }\PYG{l+m}{123}\PYG{+w}{ }\PYG{n}{tag}\PYG{+w}{ }\PYG{n}{yes}\PYG{+w}{ }\PYG{n}{temp}\PYG{+w}{ }\PYG{l+m}{1.0}

\PYG{k}{labelmap}\PYG{+w}{ }\PYG{n}{atom}\PYG{+w}{ }\PYG{l+m}{1}\PYG{+w}{ }\PYG{n}{H}\PYG{o}{+}\PYG{+w}{ }\PYG{l+m}{2}\PYG{+w}{ }\PYG{n}{OH}\PYG{o}{\PYGZhy{}}
\PYG{k}{fix }\PYG{n+nv+nvIdentifier}{chareg}\PYG{+w}{ }\PYG{n+nv+nvIdentifier}{all}\PYG{+w}{ }\PYG{n}{charge}\PYG{o}{/}\PYG{n}{regulation}\PYG{+w}{ }\PYG{n}{H}\PYG{o}{+}\PYG{+w}{ }\PYG{n}{OH}\PYG{o}{\PYGZhy{}}\PYG{+w}{ }\PYG{n}{pIp}\PYG{+w}{ }\PYG{l+m}{3}\PYG{+w}{ }\PYG{n}{pIm}\PYG{+w}{ }\PYG{l+m}{3}\PYG{+w}{ }\PYG{n}{onlysalt}\PYG{+w}{ }\PYG{n}{yes}\PYG{+w}{ }\PYG{l+m}{2}\PYG{+w}{ }\PYG{o}{\PYGZhy{}}\PYG{l+m}{1}\PYG{+w}{ }\PYG{n}{seed}\PYG{+w}{ }\PYG{l+m}{123}\PYG{+w}{ }\PYG{n}{tag}\PYG{+w}{ }\PYG{n}{yes}\PYG{+w}{ }\PYG{n}{temp}\PYG{+w}{ }\PYG{l+m}{1.0}
\end{sphinxVerbatim}


\subsection{Description}
\label{\detokenize{fix_charge_regulation:description}}
\sphinxAtStartPar
This fix performs Monte Carlo (MC) sampling of charge regulation and
exchange of ions with a reservoir as discussed in {\hyperref[\detokenize{fix_charge_regulation:curk1}]{\sphinxcrossref{\DUrole{std,std-ref}{(Curk1)}}}}
and {\hyperref[\detokenize{fix_charge_regulation:curk2}]{\sphinxcrossref{\DUrole{std,std-ref}{(Curk2)}}}}.  The implemented method is largely analogous
to the grand\sphinxhyphen{}reaction ensemble method in {\hyperref[\detokenize{fix_charge_regulation:landsgesell}]{\sphinxcrossref{\DUrole{std,std-ref}{(Landsgesell)}}}}.  The implementation is parallelized, compatible with
existing LAMMPS functionalities, and applicable to any system utilizing
discrete, ionizable groups or surface sites.

\sphinxAtStartPar
Specifically, the fix implements the following three types of MC moves,
which discretely change the charge state of individual particles and
insert ions into the systems: \(\mathrm{A} \rightleftharpoons
\mathrm{A}^-+\mathrm{X}^+\), \(\mathrm{B} \rightleftharpoons
\mathrm{B}^++\mathrm{X}^-\), and \(\emptyset \rightleftharpoons
Z^-\mathrm{X}^{Z^+}+Z^+\mathrm{X}^{-Z^-}\).  In the former two types of
reactions, Monte Carlo moves alter the charge value of specific atoms
(\(\mathrm{A}\), \(\mathrm{B}\)) and simultaneously insert a
counterion to preserve the charge neutrality of the system, modeling the
dissociation/association process.  The last type of reaction performs
grand canonical MC exchange of ion pairs with a (fictitious) reservoir.

\sphinxAtStartPar
In our implementation “acid” refers to particles that can attain charge
\(q=\{0,-1\}\) and “base” to particles with \(q=\{0,1\}\),
whereas the MC exchange of free ions allows any integer charge values of
\({Z^+}\) and \({Z^-}\).

\sphinxAtStartPar
Here we provide several practical examples for modeling charge
regulation effects in solvated systems.  An acid ionization reaction
(\(\mathrm{A} \rightleftharpoons \mathrm{A}^-+\mathrm{H}^+\)) can be
defined via a single line in the input file

\begin{sphinxVerbatim}[commandchars=\\\{\}]
\PYG{k}{fix }\PYG{n+nv+nvIdentifier}{acid\PYGZus{}reaction}\PYG{+w}{ }\PYG{n+nv+nvIdentifier}{all}\PYG{+w}{ }\PYG{n}{charge}\PYG{o}{/}\PYG{n}{regulation}\PYG{+w}{ }\PYG{l+m}{2}\PYG{+w}{ }\PYG{l+m}{3}\PYG{+w}{ }\PYG{n}{acid\PYGZus{}type}\PYG{+w}{ }\PYG{l+m}{1}\PYG{+w}{ }\PYG{n}{pH}\PYG{+w}{ }\PYG{l+m}{7.0}\PYG{+w}{ }\PYG{n}{pKa}\PYG{+w}{ }\PYG{l+m}{5.0}\PYG{+w}{ }\PYG{n}{pIp}\PYG{+w}{ }\PYG{l+m}{7.0}\PYG{+w}{ }\PYG{n}{pIm}\PYG{+w}{ }\PYG{l+m}{7.0}
\end{sphinxVerbatim}

\sphinxAtStartPar
where the fix attempts to charge \(\mathrm{A}\) (discharge
\(\mathrm{A}^-\)) to \(\mathrm{A}^-\) (\(\mathrm{A}\)) and
insert (delete) a proton (atom type 2). Besides, the fix implements
self\sphinxhyphen{}ionization reaction of water \(\emptyset \rightleftharpoons
\mathrm{H}^++\mathrm{OH}^-\).

\sphinxAtStartPar
However, this approach is highly inefficient at \(\mathrm{pH}
\approx 7\) when the concentration of both protons and hydroxyl ions is
low, resulting in a relatively low acceptance rate of MC moves.

\sphinxAtStartPar
A more efficient way is to allow salt ions to participate in ionization
reactions, which can be easily achieved via

\begin{sphinxVerbatim}[commandchars=\\\{\}]
\PYG{k}{fix }\PYG{n+nv+nvIdentifier}{acid\PYGZus{}reaction2}\PYG{+w}{ }\PYG{n+nv+nvIdentifier}{all}\PYG{+w}{ }\PYG{n}{charge}\PYG{o}{/}\PYG{n}{regulation}\PYG{+w}{ }\PYG{l+m}{4}\PYG{+w}{ }\PYG{l+m}{5}\PYG{+w}{ }\PYG{n}{acid\PYGZus{}type}\PYG{+w}{ }\PYG{l+m}{1}\PYG{+w}{ }\PYG{n}{pH}\PYG{+w}{ }\PYG{l+m}{7.0}\PYG{+w}{ }\PYG{n}{pKa}\PYG{+w}{ }\PYG{l+m}{5.0}\PYG{+w}{ }\PYG{n}{pIp}\PYG{+w}{ }\PYG{l+m}{2.0}\PYG{+w}{ }\PYG{n}{pIm}\PYG{+w}{ }\PYG{l+m}{2.0}
\end{sphinxVerbatim}

\sphinxAtStartPar
where particles of atom type 4 and 5 are the salt cations and anions,
both at activity (effective concentration) of \(10^{-2}\) mol/l, see
{\hyperref[\detokenize{fix_charge_regulation:curk1}]{\sphinxcrossref{\DUrole{std,std-ref}{(Curk1)}}}} and {\hyperref[\detokenize{fix_charge_regulation:landsgesell}]{\sphinxcrossref{\DUrole{std,std-ref}{(Landsgesell)}}}} for more
details.

\sphinxAtStartPar
We could have simultaneously added a base ionization reaction
(\(\mathrm{B} \rightleftharpoons \mathrm{B}^++\mathrm{OH}^-\))

\begin{sphinxVerbatim}[commandchars=\\\{\}]
\PYG{k}{fix }\PYG{n+nv+nvIdentifier}{acid\PYGZus{}base\PYGZus{}reaction}\PYG{+w}{ }\PYG{n+nv+nvIdentifier}{all}\PYG{+w}{ }\PYG{n}{charge}\PYG{o}{/}\PYG{n}{regulation}\PYG{+w}{ }\PYG{l+m}{2}\PYG{+w}{ }\PYG{l+m}{3}\PYG{+w}{ }\PYG{n}{acid\PYGZus{}type}\PYG{+w}{ }\PYG{l+m}{1}\PYG{+w}{ }\PYG{n}{base\PYGZus{}type}\PYG{+w}{ }\PYG{l+m}{6}\PYG{+w}{ }\PYG{n}{pH}\PYG{+w}{ }\PYG{l+m}{7.0}\PYG{+w}{ }\PYG{n}{pKa}\PYG{+w}{ }\PYG{l+m}{5.0}\PYG{+w}{ }\PYG{n}{pKb}\PYG{+w}{ }\PYG{l+m}{6.0}\PYG{+w}{ }\PYG{n}{pIp}\PYG{+w}{ }\PYG{l+m}{7.0}\PYG{+w}{ }\PYG{n}{pIm}\PYG{+w}{ }\PYG{l+m}{7.0}
\end{sphinxVerbatim}

\sphinxAtStartPar
where the fix will attempt to charge \(\mathrm{B}\) (discharge
\(\mathrm{B}^+\)) to \(\mathrm{B}^+\) (\(\mathrm{B}\)) and
insert (delete) a hydroxyl ion \(\mathrm{OH}^-\) of atom type 3.

\sphinxAtStartPar
Dissociated ions and salt ions can be combined into a single particle
type, which reduces the number of necessary MC moves and increases
sampling performance, see {\hyperref[\detokenize{fix_charge_regulation:curk1}]{\sphinxcrossref{\DUrole{std,std-ref}{(Curk1)}}}}. The
\(\mathrm{H}^+\) and monovalent salt cation (\(\mathrm{S}^+\))
are combined into a single particle type, \(\mathrm{X}^+ =
\{\mathrm{H}^+, \mathrm{S}^+\}\). In this case “pIp” refers to the
effective concentration of the combined cation type \(\mathrm{X}^+\)
and its value is determined by \(10^{-\mathrm{pIp}} =
10^{-\mathrm{pH}} + 10^{-\mathrm{pSp}}\), where
\(10^{-\mathrm{pSp}}\) is the effective concentration of salt
cations. For example, at pH=7 and pSp=6 we would find pIp\textasciitilde{}5.958 and the
command that performs reactions with combined ions could read,

\begin{sphinxVerbatim}[commandchars=\\\{\}]
\PYG{k}{fix }\PYG{n+nv+nvIdentifier}{acid\PYGZus{}reaction\PYGZus{}combined}\PYG{+w}{ }\PYG{n+nv+nvIdentifier}{all}\PYG{+w}{ }\PYG{n}{charge}\PYG{o}{/}\PYG{n}{regulation}\PYG{+w}{ }\PYG{l+m}{2}\PYG{+w}{ }\PYG{l+m}{3}\PYG{+w}{ }\PYG{n}{acid\PYGZus{}type}\PYG{+w}{ }\PYG{l+m}{1}\PYG{+w}{ }\PYG{n}{pH}\PYG{+w}{ }\PYG{l+m}{7.0}\PYG{+w}{ }\PYG{n}{pKa}\PYG{+w}{ }\PYG{l+m}{5.0}\PYG{+w}{ }\PYG{n}{pIp}\PYG{+w}{ }\PYG{l+m}{5.958}\PYG{+w}{ }\PYG{n}{pIm}\PYG{+w}{ }\PYG{l+m}{5.958}
\end{sphinxVerbatim}

\sphinxAtStartPar
If neither the acid or the base type is specified, for example,

\begin{sphinxVerbatim}[commandchars=\\\{\}]
\PYG{k}{fix }\PYG{n+nv+nvIdentifier}{salt\PYGZus{}reaction}\PYG{+w}{ }\PYG{n+nv+nvIdentifier}{all}\PYG{+w}{ }\PYG{n}{charge}\PYG{o}{/}\PYG{n}{regulation}\PYG{+w}{ }\PYG{l+m}{4}\PYG{+w}{ }\PYG{l+m}{5}\PYG{+w}{ }\PYG{n}{pIp}\PYG{+w}{ }\PYG{l+m}{2.0}\PYG{+w}{ }\PYG{n}{pIm}\PYG{+w}{ }\PYG{l+m}{2.0}
\end{sphinxVerbatim}

\sphinxAtStartPar
the fix simply inserts or deletes an ion pair of a free cation (atom
type 4) and a free anion (atom type 5) as done in a conventional
grand\sphinxhyphen{}canonical MC simulation. Multivalent ions can be inserted
(deleted) by using the \sphinxstyleemphasis{onlysalt} keyword.

\sphinxAtStartPar
This fix is compatible with LAMMPS packages such as MOLECULE or
RIGID. The acid and base particles can be part of larger molecules or
rigid bodies. Free ions that are inserted to or deleted from the system
must be defined as single particles (no bonded interactions allowed) and
cannot be part of larger molecules or rigid bodies. If an atom style
with molecule IDs is used, all inserted ions have a molecule ID equal to
zero.

\sphinxAtStartPar
Note that LAMMPS implicitly assumes a constant number of particles
(degrees of freedom). Since using this fix alters the total number of
particles during the simulation, any thermostat used by LAMMPS, such as
NVT or Langevin, must use a dynamic calculation of system
temperature. This can be achieved by specifying a dynamic temperature
compute (e.g. dtemp) and using it with the desired thermostat, e.g. a
Langevin thermostat:

\begin{sphinxVerbatim}[commandchars=\\\{\}]
\PYG{k}{compute }\PYG{n+nv+nvIdentifier}{dtemp}\PYG{+w}{ }\PYG{n+nv+nvIdentifier}{all}\PYG{+w}{ }\PYG{n}{temp}
\PYG{k}{compute\PYGZus{}modify }\PYG{n+nv+nvIdentifier}{dtemp}\PYG{+w}{ }\PYG{n}{dynamic}\PYG{o}{/}\PYG{n}{dof}\PYG{+w}{ }\PYG{n}{yes}
\PYG{k}{fix }\PYG{n+nv+nvIdentifier}{fT}\PYG{+w}{ }\PYG{n+nv+nvIdentifier}{all}\PYG{+w}{ }\PYG{n}{langevin}\PYG{+w}{ }\PYG{l+m}{1.0}\PYG{+w}{ }\PYG{l+m}{1.0}\PYG{+w}{ }\PYG{l+m}{1.0}\PYG{+w}{ }\PYG{l+m}{123}
\PYG{k}{fix\PYGZus{}modify }\PYG{n+nv+nvIdentifier}{fT}\PYG{+w}{ }\PYG{n}{temp}\PYG{+w}{ }\PYG{n}{dtemp}
\end{sphinxVerbatim}

\sphinxAtStartPar
The units of pH, pKa, pKb, pIp, pIm are considered to be in the
standard \sphinxhyphen{}log10 representation assuming reference concentration
\(\rho_0 = \mathrm{mol}/\mathrm{l}\).  For example, in the dilute
ideal solution limit, the concentration of free cations will be
\(c_\mathrm{I} = 10^{-\mathrm{pIp}}\mathrm{mol}/\mathrm{l}\). To
perform the internal unit conversion, the the value of the LAMMPS unit
length must be specified in nanometers via \sphinxstyleemphasis{lunit\_nm}. The default value
is set to the Bjerrum length in water at room temperature (0.71 nm),
\sphinxstyleemphasis{lunit\_nm} = 0.71.

\sphinxAtStartPar
The temperature used in MC acceptance probability is set by \sphinxstyleemphasis{temp}. This
temperature should be the same as the temperature set by the molecular
dynamics thermostat. For most purposes, it is probably best to use
\sphinxstyleemphasis{tempfixid} keyword which dynamically sets the temperature equal to the
chosen MD thermostat temperature, in the example above we assumed the
thermostat fix\sphinxhyphen{}ID is \sphinxstyleemphasis{fT}. The inserted particles attain a random
velocity corresponding to the specified temperature. Using \sphinxstyleemphasis{tempfixid}
overrides any fixed temperature set by \sphinxstyleemphasis{temp}.

\sphinxAtStartPar
The \sphinxstyleemphasis{rxd} keyword can be used to restrict the inserted/deleted
counterions to a specific radial distance from an acid or base particle
that is currently participating in a reaction. This can be used to
simulate more realist reaction dynamics. If \sphinxstyleemphasis{rxd} = 0 or \sphinxstyleemphasis{rxd} \textgreater{} \sphinxstyleemphasis{L} /
2, where \sphinxstyleemphasis{L} is the smallest box dimension, the radial restriction is
automatically turned off and free ion can be inserted or deleted
anywhere in the simulation box.

\sphinxAtStartPar
If the \sphinxstyleemphasis{tag yes} is used, every inserted atom gets a unique tag ID,
otherwise, the tag of every inserted atom is set to 0. \sphinxstyleemphasis{tag yes} might
cause an integer overflow in very long simulations since the tags are
unique to every particle and thus increase with every successful
particle insertion.

\sphinxAtStartPar
The \sphinxstyleemphasis{pmcmoves} keyword sets the relative probability of attempting the
three types of MC moves (reactions): acid charging, base charging, and
ion pair exchange.  The fix only attempts to perform particle charging
MC moves if \sphinxstyleemphasis{acid\_type} or \sphinxstyleemphasis{base\_type} is defined. Otherwise fix only
performs free ion insertion/deletion. For example, if \sphinxstyleemphasis{acid\_type} is not
defined, \sphinxstyleemphasis{pmcA} is automatically set to 0. The vector \sphinxstyleemphasis{pmcmoves} is
automatically normalized, for example, if set to \sphinxstyleemphasis{pmcmoves} 0 0.33 0.33,
the vector would be normalized to {[}0,0.5,0.5{]}.

\sphinxAtStartPar
The \sphinxstyleemphasis{only\_salt} option can be used to perform multivalent
grand\sphinxhyphen{}canonical ion\sphinxhyphen{}exchange moves. If \sphinxstyleemphasis{only\_salt yes} is used, no
charge exchange is performed, only ion insertion/deletion (\sphinxstyleemphasis{pmcmoves} is
set to {[}0,0,1{]}), but ions can be multivalent. In the example above, an
MC move would consist of three ion insertion/deletion to preserve the
charge neutrality of the system.

\sphinxAtStartPar
The \sphinxstyleemphasis{group} keyword can be used to add inserted particles to a specific
group\sphinxhyphen{}ID. All inserted particles are automatically added to group \sphinxstyleemphasis{all}.


\subsection{Output}
\label{\detokenize{fix_charge_regulation:output}}
\sphinxAtStartPar
This fix computes a global vector of length 8, which can be accessed by
various output commands. The vector values are the following global
quantities:
\begin{enumerate}
\sphinxsetlistlabels{\arabic}{enumi}{enumii}{}{.}%
\item {} 
\sphinxAtStartPar
cumulative MC attempts

\item {} 
\sphinxAtStartPar
cumulative MC successes

\item {} 
\sphinxAtStartPar
current \# of neutral acid atoms

\item {} 
\sphinxAtStartPar
current \# of \sphinxhyphen{}1 charged acid atoms

\item {} 
\sphinxAtStartPar
current \# of neutral base atoms

\item {} 
\sphinxAtStartPar
current \# of +1 charged base atoms

\item {} 
\sphinxAtStartPar
current \# of free cations

\item {} 
\sphinxAtStartPar
current \# of free anions

\end{enumerate}


\subsection{Restrictions}
\label{\detokenize{fix_charge_regulation:restrictions}}
\sphinxAtStartPar
This fix is part of the MC package. It is only enabled if LAMMPS was
built with that package.  See the {\hyperref[\detokenize{Build_package::doc}]{\sphinxcrossref{\DUrole{doc}{Build package}}}}
page for more info.

\sphinxAtStartPar
The {\hyperref[\detokenize{atom_style::doc}]{\sphinxcrossref{\DUrole{doc}{atom\_style}}}}, used must contain the charge
property and have per atom type masses, for example, the style could be
\sphinxstyleemphasis{charge} or \sphinxstyleemphasis{full}. Only usable for 3D simulations.  Atoms specified as
free ions cannot be part of rigid bodies or molecules and cannot have
bonding interactions.  The scheme is limited to integer charges, any
atoms with non\sphinxhyphen{}integer charges will not be considered by the fix.

\sphinxAtStartPar
All interaction potentials used must be continuous, otherwise the MD
integration and the particle exchange MC moves do not correspond to the
same equilibrium ensemble. For example, if an lj/cut pair style is used,
the LJ potential must be shifted so that it vanishes at the cutoff. This
can be easily achieved using the {\hyperref[\detokenize{pair_modify::doc}]{\sphinxcrossref{\DUrole{doc}{pair\_modify}}}}
command, i.e., by using: \sphinxstyleemphasis{pair\_modify shift yes}.

\begin{sphinxadmonition}{note}{Note:}
\sphinxAtStartPar
Region restrictions are not yet implemented.
\end{sphinxadmonition}


\subsection{Related commands}
\label{\detokenize{fix_charge_regulation:related-commands}}
\sphinxAtStartPar
{\hyperref[\detokenize{fix_gcmc::doc}]{\sphinxcrossref{\DUrole{doc}{fix gcmc}}}},
{\hyperref[\detokenize{fix_atom_swap::doc}]{\sphinxcrossref{\DUrole{doc}{fix atom/swap}}}}


\subsection{Default}
\label{\detokenize{fix_charge_regulation:default}}
\sphinxAtStartPar
pH = 7.0; pKa = 100.0; pKb = 100.0; pIp = 5.0; pIm = 5.0; pKs = 14.0;
acid\_type = \sphinxhyphen{}1; base\_type = \sphinxhyphen{}1; lunit\_nm = 0.71; temp = 1.0; nevery =
100; nmc = 100; rxd = 0; seed = 0; tag = no; onlysalt = no, pmcmoves =
{[}1/3, 1/3, 1/3{]}, group\sphinxhyphen{}ID = all


\bigskip\hrule\bigskip

\phantomsection\label{\detokenize{fix_charge_regulation:curk1}}
\sphinxAtStartPar
\sphinxstylestrong{(Curk1)} T. Curk, J. Yuan, and E. Luijten, “Accelerated simulation method for charge regulation effects”, JCP 156 (2022).

\phantomsection\label{\detokenize{fix_charge_regulation:curk2}}
\sphinxAtStartPar
\sphinxstylestrong{(Curk2)} T. Curk and E. Luijten, “Charge\sphinxhyphen{}regulation effects in nanoparticle self\sphinxhyphen{}assembly”, PRL 126 (2021)

\phantomsection\label{\detokenize{fix_charge_regulation:landsgesell}}
\sphinxAtStartPar
\sphinxstylestrong{(Landsgesell)} J. Landsgesell, P. Hebbeker, O. Rud, R. Lunkad, P. Kosovan, and C. Holm, “Grand\sphinxhyphen{}reaction method for simulations of ionization equilibria coupled to ion partitioning”, Macromolecules 53, 3007\sphinxhyphen{}3020 (2020).

\sphinxstepscope

\index{fix cmap@\spxentry{fix cmap}}\ignorespaces 

\section{fix cmap command}
\label{\detokenize{fix_cmap:fix-cmap-command}}\label{\detokenize{fix_cmap:index-0}}\label{\detokenize{fix_cmap::doc}}

\subsection{Syntax}
\label{\detokenize{fix_cmap:syntax}}
\begin{sphinxVerbatim}[commandchars=\\\{\}]
\PYG{k}{fix }\PYG{n+nv+nvIdentifier}{ID}\PYG{+w}{ }\PYG{n+nv+nvIdentifier}{group\PYGZhy{}ID}\PYG{+w}{ }\PYG{n}{cmap}\PYG{+w}{ }\PYG{n}{filename}
\end{sphinxVerbatim}
\begin{itemize}
\item {} 
\sphinxAtStartPar
ID, group\sphinxhyphen{}ID are documented in {\hyperref[\detokenize{fix::doc}]{\sphinxcrossref{\DUrole{doc}{fix}}}} command

\item {} 
\sphinxAtStartPar
cmap = style name of this fix command

\item {} 
\sphinxAtStartPar
filename = force\sphinxhyphen{}field file with CMAP coefficients

\end{itemize}


\subsection{Examples}
\label{\detokenize{fix_cmap:examples}}
\begin{sphinxVerbatim}[commandchars=\\\{\}]
\PYG{k}{fix            }\PYG{n+nv+nvIdentifier}{myCMAP}\PYG{+w}{ }\PYG{n+nv+nvIdentifier}{all}\PYG{+w}{ }\PYG{n}{cmap}\PYG{+w}{ }\PYG{n}{..}\PYG{o}{/}\PYG{n}{potentials}\PYG{o}{/}\PYG{n}{cmap36.data}
\PYG{k}{read\PYGZus{}data}\PYG{+w}{      }\PYG{n}{proteinX.data}\PYG{+w}{ }\PYG{k}{fix }\PYG{n+nv+nvIdentifier}{myCMAP}\PYG{+w}{ }\PYG{n+nv+nvIdentifier}{crossterm}\PYG{+w}{ }\PYG{n}{CMAP}
\PYG{k}{fix\PYGZus{}modify     }\PYG{n+nv+nvIdentifier}{myCMAP}\PYG{+w}{ }\PYG{n}{energy}\PYG{+w}{ }\PYG{n}{yes}
\end{sphinxVerbatim}


\subsection{Description}
\label{\detokenize{fix_cmap:description}}
\sphinxAtStartPar
This command enables CMAP 5\sphinxhyphen{}body interactions to be added to
simulations which use the CHARMM force field.  These are relevant for
any CHARMM model of a peptide or protein sequences that is 3 or more
amino\sphinxhyphen{}acid residues long; see {\hyperref[\detokenize{fix_cmap:buck}]{\sphinxcrossref{\DUrole{std,std-ref}{(Buck)}}}} and {\hyperref[\detokenize{fix_cmap:brooks2}]{\sphinxcrossref{\DUrole{std,std-ref}{(Brooks)}}}} for details, including the analytic energy expressions for
CMAP interactions.  The CMAP 5\sphinxhyphen{}body terms add additional potential
energy contributions to pairs of overlapping phi\sphinxhyphen{}psi dihedrals of
amino\sphinxhyphen{}acids, which are important to properly represent their
conformational behavior.

\sphinxAtStartPar
The examples/cmap directory has a sample input script and data file
for a small peptide, that illustrates use of the fix cmap command.

\sphinxAtStartPar
As in the example above, this fix should be used before reading a data
file that contains a listing of CMAP interactions.  The \sphinxstyleemphasis{filename}
specified should contain the CMAP parameters for a particular version
of the CHARMM force field.  Two such files are including in the
lammps/potentials directory: charmm22.cmap and charmm36.cmap.

\sphinxAtStartPar
The data file read by the “read\_data” must contain the topology of all
the CMAP interactions, similar to the topology data for bonds, angles,
dihedrals, etc.  Specially it should have a line like this in its
header section:

\begin{sphinxVerbatim}[commandchars=\\\{\}]
N crossterms
\end{sphinxVerbatim}

\sphinxAtStartPar
where N is the number of CMAP 5\sphinxhyphen{}body interactions.  It should also
have a section in the body of the data file like this with N lines:

\begin{sphinxVerbatim}[commandchars=\\\{\}]
CMAP

       1       1       8      10      12      18      20
       2       5      18      20      22      25      27
       [...]
       N       3     314     315     317      318    330
\end{sphinxVerbatim}

\sphinxAtStartPar
The first column is an index from 1 to N to enumerate the CMAP 5\sphinxhyphen{}atom
tuples; it is ignored by LAMMPS.  The second column is the “type” of
the interaction; it is an index into the CMAP force field file.  The
remaining 5 columns are the atom IDs of the atoms in the two 4\sphinxhyphen{}atom
dihedrals that overlap to create the CMAP interaction.  Note that the
“crossterm” and “CMAP” keywords for the header and body sections match
those specified in the read\_data command following the data file name;
see the {\hyperref[\detokenize{read_data::doc}]{\sphinxcrossref{\DUrole{doc}{read\_data}}}} page for more details.

\sphinxAtStartPar
A data file containing CMAP 5\sphinxhyphen{}body interactions can be generated from
a PDB file using the charmm2lammps.pl script in the tools/ch2lmp
directory of the LAMMPS distribution.  The script must be invoked with
the optional “\sphinxhyphen{}cmap” flag to do this; see the tools/ch2lmp/README file
for more information.  The same conversion script also creates the
file of CMAP coefficient data which is read by this command.

\sphinxAtStartPar
The potential energy associated with CMAP interactions can be output
as described below.  It can also be included in the total potential
energy of the system, as output by the {\hyperref[\detokenize{thermo_style::doc}]{\sphinxcrossref{\DUrole{doc}{thermo\_style}}}} command, if the {\hyperref[\detokenize{fix_modify::doc}]{\sphinxcrossref{\DUrole{doc}{fix\_modify energy}}}}
command is used, as in the example above.  See the note below about
how to include the CMAP energy when performing an {\hyperref[\detokenize{minimize::doc}]{\sphinxcrossref{\DUrole{doc}{energy
minimization}}}}.


\bigskip\hrule\bigskip



\subsection{Restart, fix\_modify, output, run start/stop, minimize info}
\label{\detokenize{fix_cmap:restart-fix-modify-output-run-start-stop-minimize-info}}
\sphinxAtStartPar
This fix writes the list of CMAP cross\sphinxhyphen{}terms to {\hyperref[\detokenize{restart::doc}]{\sphinxcrossref{\DUrole{doc}{binary restart
files}}}}.  See the {\hyperref[\detokenize{read_restart::doc}]{\sphinxcrossref{\DUrole{doc}{read\_restart}}}} command
for info on how to re\sphinxhyphen{}specify a fix in an input script that reads a
restart file, so that the operation of the fix continues in an
uninterrupted fashion.

\sphinxAtStartPar
The {\hyperref[\detokenize{fix_modify::doc}]{\sphinxcrossref{\DUrole{doc}{fix\_modify}}}} \sphinxstyleemphasis{energy} option is supported by
this fix to add the potential energy of the CMAP interactions to both
the global potential energy and peratom potential energies of the
system as part of {\hyperref[\detokenize{thermo_style::doc}]{\sphinxcrossref{\DUrole{doc}{thermodynamic output}}}} or
output by the {\hyperref[\detokenize{compute_pe_atom::doc}]{\sphinxcrossref{\DUrole{doc}{compute pe/atom}}}} command.  The
default setting for this fix is {\hyperref[\detokenize{fix_modify::doc}]{\sphinxcrossref{\DUrole{doc}{fix\_modify energy yes}}}}.

\sphinxAtStartPar
The {\hyperref[\detokenize{fix_modify::doc}]{\sphinxcrossref{\DUrole{doc}{fix\_modify}}}} \sphinxstyleemphasis{virial} option is supported by
this fix to add the contribution due to the CMAP interactions to both
the global pressure and per\sphinxhyphen{}atom stress of the system via the
{\hyperref[\detokenize{compute_pressure::doc}]{\sphinxcrossref{\DUrole{doc}{compute pressure}}}} and {\hyperref[\detokenize{compute_stress_atom::doc}]{\sphinxcrossref{\DUrole{doc}{compute
stress/atom}}}} commands.  The former can be
accessed by {\hyperref[\detokenize{thermo_style::doc}]{\sphinxcrossref{\DUrole{doc}{thermodynamic output}}}}.  The default
setting for this fix is {\hyperref[\detokenize{fix_modify::doc}]{\sphinxcrossref{\DUrole{doc}{fix\_modify virial yes}}}}.

\sphinxAtStartPar
This fix computes a global scalar which can be accessed by various
{\hyperref[\detokenize{Howto_output::doc}]{\sphinxcrossref{\DUrole{doc}{output commands}}}}.  The scalar is the potential
energy discussed above.  The scalar value calculated by this fix is
“extensive”.

\sphinxAtStartPar
No parameter of this fix can be used with the \sphinxstyleemphasis{start/stop} keywords of
the {\hyperref[\detokenize{run::doc}]{\sphinxcrossref{\DUrole{doc}{run}}}} command.

\sphinxAtStartPar
The forces due to this fix are imposed during an energy minimization,
invoked by the {\hyperref[\detokenize{minimize::doc}]{\sphinxcrossref{\DUrole{doc}{minimize}}}} command.

\sphinxAtStartPar
The {\hyperref[\detokenize{fix_modify::doc}]{\sphinxcrossref{\DUrole{doc}{fix\_modify}}}} \sphinxstyleemphasis{respa} option is supported by this
fix. This allows to set at which level of the {\hyperref[\detokenize{run_style::doc}]{\sphinxcrossref{\DUrole{doc}{r\sphinxhyphen{}RESPA}}}} integrator the fix is adding its forces. Default is the
outermost level.

\begin{sphinxadmonition}{note}{Note:}
\sphinxAtStartPar
For energy minimization, if you want the potential energy
associated with the CMAP terms forces to be included in the total
potential energy of the system (the quantity being minimized), you
MUST not disable the {\hyperref[\detokenize{fix_modify::doc}]{\sphinxcrossref{\DUrole{doc}{fix\_modify}}}} \sphinxstyleemphasis{energy} option
for this fix.
\end{sphinxadmonition}


\subsection{Restrictions}
\label{\detokenize{fix_cmap:restrictions}}
\sphinxAtStartPar
To function as expected this fix command must be issued \sphinxstyleemphasis{before} a
{\hyperref[\detokenize{read_data::doc}]{\sphinxcrossref{\DUrole{doc}{read\_data}}}} command but \sphinxstyleemphasis{after} a
{\hyperref[\detokenize{read_restart::doc}]{\sphinxcrossref{\DUrole{doc}{read\_restart}}}} command.

\sphinxAtStartPar
This fix can only be used if LAMMPS was built with the MOLECULE
package.  See the {\hyperref[\detokenize{Build_package::doc}]{\sphinxcrossref{\DUrole{doc}{Build package}}}} page for more
info.


\subsection{Related commands}
\label{\detokenize{fix_cmap:related-commands}}
\sphinxAtStartPar
{\hyperref[\detokenize{fix_modify::doc}]{\sphinxcrossref{\DUrole{doc}{fix\_modify}}}}, {\hyperref[\detokenize{read_data::doc}]{\sphinxcrossref{\DUrole{doc}{read\_data}}}}


\subsection{Default}
\label{\detokenize{fix_cmap:default}}
\sphinxAtStartPar
none


\bigskip\hrule\bigskip

\phantomsection\label{\detokenize{fix_cmap:buck}}
\sphinxAtStartPar
\sphinxstylestrong{(Buck)} Buck, Bouguet\sphinxhyphen{}Bonnet, Pastor, MacKerell Jr., Biophys J, 90, L36
(2006).

\phantomsection\label{\detokenize{fix_cmap:brooks2}}
\sphinxAtStartPar
\sphinxstylestrong{(Brooks)} Brooks, Brooks, MacKerell Jr., J Comput Chem, 30, 1545 (2009).

\sphinxstepscope

\index{fix colvars@\spxentry{fix colvars}}\ignorespaces 

\section{fix colvars command}
\label{\detokenize{fix_colvars:fix-colvars-command}}\label{\detokenize{fix_colvars:index-0}}\label{\detokenize{fix_colvars::doc}}

\subsection{Syntax}
\label{\detokenize{fix_colvars:syntax}}
\begin{sphinxVerbatim}[commandchars=\\\{\}]
\PYG{k}{fix }\PYG{n+nv+nvIdentifier}{ID}\PYG{+w}{ }\PYG{n+nv+nvIdentifier}{group\PYGZhy{}ID}\PYG{+w}{ }\PYG{n}{colvars}\PYG{+w}{ }\PYG{o}{*}\PYG{n}{configfile}\PYG{o}{*}\PYG{+w}{ }\PYG{n}{keyword}\PYG{+w}{ }\PYG{n}{value}\PYG{+w}{ }\PYG{n}{...}
\end{sphinxVerbatim}
\begin{itemize}
\item {} 
\sphinxAtStartPar
\sphinxstyleemphasis{ID}, \sphinxstyleemphasis{group\sphinxhyphen{}ID} are documented in {\hyperref[\detokenize{fix::doc}]{\sphinxcrossref{\DUrole{doc}{fix}}}} command

\item {} 
\sphinxAtStartPar
“colvars” = style name of this fix command

\item {} 
\sphinxAtStartPar
\sphinxstyleemphasis{configfile} = configuration file for Colvars (use “\sphinxstyleemphasis{none}” to provide it inline)

\item {} 
\sphinxAtStartPar
keyword = \sphinxstyleemphasis{output} or \sphinxstyleemphasis{input} or \sphinxstyleemphasis{unwrap} or \sphinxstyleemphasis{tstat} or \sphinxstyleemphasis{seed}
\begin{sphinxalltt}
\sphinxstyleemphasis{output} value = state filename/prefix for Colvars (default: "out")
\sphinxstyleemphasis{input} value = input state filename/prefix for Colvars (optional, default: "NULL")
\sphinxstyleemphasis{unwrap} value = "yes" or "no" (default: "yes")
\sphinxstyleemphasis{tstat} value = fix ID of thermostat applied to relevant atoms (default: "NULL")
\sphinxstyleemphasis{seed} value = seed for random number generator (default: 1966)
\end{sphinxalltt}

\end{itemize}


\subsection{Examples}
\label{\detokenize{fix_colvars:examples}}
\begin{sphinxVerbatim}[commandchars=\\\{\}]
\PYG{c}{\PYGZsh{} Create the fix using a config file, set prefix for its output files}
\PYG{k}{fix }\PYG{n+nv+nvIdentifier}{Colvars}\PYG{+w}{ }\PYG{n+nv+nvIdentifier}{all}\PYG{+w}{ }\PYG{n}{colvars}\PYG{+w}{ }\PYG{n}{colvars.inp}\PYG{+w}{ }\PYG{n}{output}\PYG{+w}{ }\PYG{n+nv}{\PYGZdl{}\PYGZob{}}\PYG{n+nv}{JOB}\PYG{n+nv}{\PYGZcb{}}

\PYG{c}{\PYGZsh{} Communicate the LAMMPS target temperature to the Colvars module}
\PYG{k}{fix\PYGZus{}modify }\PYG{n+nv+nvIdentifier}{Colvars}\PYG{+w}{ }\PYG{n}{tstat}\PYG{+w}{ }\PYG{n}{NPT}

\PYG{c}{\PYGZsh{} Add a new restraint specific to this LAMMPS run}
\PYG{k}{fix\PYGZus{}modify }\PYG{n+nv+nvIdentifier}{Colvars}\PYG{+w}{ }\PYG{n}{config}\PYG{+w}{ }\PYG{l+s}{\PYGZdq{}}\PYG{l+s}{\PYGZdq{}}\PYG{l+s}{\PYGZdq{}}
\PYG{l+s}{harmonic \PYGZob{}}
\PYG{l+s}{    name restraint}
\PYG{l+s}{    colvars distance1 distance2}
\PYG{l+s}{    centers \PYGZdl{}\PYGZob{}ref1\PYGZcb{} \PYGZdl{}\PYGZob{}ref2\PYGZcb{}}
\PYG{l+s}{    forceConstant \PYGZdl{}\PYGZob{}kappa\PYGZcb{}}
\PYG{l+s}{\PYGZcb{}}\PYG{l+s}{\PYGZdq{}}\PYG{l+s}{\PYGZdq{}}\PYG{l+s}{\PYGZdq{}}
\end{sphinxVerbatim}


\subsection{Description}
\label{\detokenize{fix_colvars:description}}
\sphinxAtStartPar
This fix interfaces LAMMPS to the collective variables \sphinxhref{https://colvars.github.io}{Colvars} library, which allows to accelerate
sampling of rare events and the computation of free energy surfaces
and potentials of mean force (PMFs) for any set of collective
variables using a variety of sampling methods (e.g. umbrella\sphinxhyphen{}sampling,
metadynamics, ABF…).

\sphinxAtStartPar
This documentation describes only the “fix colvars” command itself in a
LAMMPS script.  The Colvars library is fully documented in the included
\sphinxhref{PDF/colvars-refman-lammps.pdf}{PDF manual} or in the webpage
\sphinxurl{https://colvars.github.io/colvars-refman-lammps/colvars-refman-lammps.html}.

\sphinxAtStartPar
The Colvars library is developed at \sphinxhref{https://github.com/colvars/colvars}{https://github.com/Colvars/colvars}.  A detailed discussion of its
implementation is in {\hyperref[\detokenize{fix_colvars:fiorin}]{\sphinxcrossref{\DUrole{std,std-ref}{(Fiorin)}}}}; additional citations for
specific features are printed at runtime if these features are used.

\sphinxAtStartPar
There are example scripts on the \sphinxhref{https://colvars.github.io}{Colvars website} as well as in the
\sphinxcode{\sphinxupquote{examples/PACKAGES/colvars}} directory in the LAMMPS source tree.


\bigskip\hrule\bigskip


\sphinxAtStartPar
The only required argument to the fix is the name of the Colvars
configuration file.  The contents of this file are independent from the
MD engine in which the Colvars library has been integrated, save for the
units that are specific to each engine.  In LAMMPS, the units used by
Colvars are consistent with those specified by the {\hyperref[\detokenize{units::doc}]{\sphinxcrossref{\DUrole{doc}{units}}}}
command.

\sphinxAtStartPar
\DUrole{versionmodified,added}{Added in version Colvars\_2023\sphinxhyphen{}06\sphinxhyphen{}04: }The special value “\sphinxstyleemphasis{none}” (lowercase) initializes an empty Colvars
module, which allows loading configuration dynamically using
{\hyperref[\detokenize{fix_modify::doc}]{\sphinxcrossref{\DUrole{doc}{fix\_modify}}}} (see below).

\sphinxAtStartPar
The \sphinxstyleemphasis{group\sphinxhyphen{}ID} entry is ignored.  “fix colvars” will always apply to the
entire system, but specific atoms will be selected based on selection
keywords in the Colvars configuration file or files.  There is no need
to define multiple “fix colvars” instances and it is not allowed.

\sphinxAtStartPar
The “output” keyword allows to specify the prefix of output files
generated by Colvars, for example “\sphinxstyleemphasis{output}.colvars.traj” or
“output.pmf”.  Supplying an empty string suppresses any file output from
Colvars to file, except for data saved into the LAMMPS {\hyperref[\detokenize{restart::doc}]{\sphinxcrossref{\DUrole{doc}{binary
restart}}}} files.

\sphinxAtStartPar
The “input” keyword allows to specify an optional state file that
contains the restart information needed to continue a previous
simulation state.  However, because “fix colvars” records its state in
LAMMPS {\hyperref[\detokenize{restart::doc}]{\sphinxcrossref{\DUrole{doc}{binary restart}}}} files, this is usually not needed
when using the {\hyperref[\detokenize{read_restart::doc}]{\sphinxcrossref{\DUrole{doc}{read\_restart}}}} command.

\sphinxAtStartPar
The \sphinxstyleemphasis{unwrap} keyword controls whether wrapped or unwrapped coordinates
are passed to the Colvars library for calculation of the collective
variables and the resulting forces.  The default is \sphinxstyleemphasis{yes}, i.e. the
image flags are used to reconstruct the absolute atom positions.
Setting this to \sphinxstyleemphasis{no} will use the current local coordinates that are
wrapped back into the simulation cell at each re\sphinxhyphen{}neighboring step
instead.  For information about when and how this affects results,
please see
\sphinxurl{https://colvars.github.io/colvars-refman-lammps/colvars-refman-lammps.html\#sec:colvar\_atom\_groups\_wrapping}.

\sphinxAtStartPar
The \sphinxstyleemphasis{tstat} keyword can be either “NULL” or the label of a
thermostatting fix that thermostats all atoms in the fix colvars
group. This will be used to provide the colvars module with the current
thermostat target temperature.

\sphinxAtStartPar
The \sphinxstyleemphasis{seed} keyword contains the seed for the random number generator
that will be used in the colvars module.


\subsection{Restarting}
\label{\detokenize{fix_colvars:restarting}}
\sphinxAtStartPar
This fix writes the current state of the Colvars module into
{\hyperref[\detokenize{restart::doc}]{\sphinxcrossref{\DUrole{doc}{binary restart files}}}}.  This is in addition to the
text\sphinxhyphen{}mode “.colvars.state” state file that is written by the Colvars
module itself.  The information contained in both files is identical,
and the binary LAMMPS restart file is also used by fix colvars when
{\hyperref[\detokenize{read_restart::doc}]{\sphinxcrossref{\DUrole{doc}{read\_restart}}}} is called in a LAMMPS script.  In
that case, there is typically no need to specify the \sphinxstyleemphasis{input} keyword.

\sphinxAtStartPar
As long as LAMMPS binary restarts will be used to continue a simulation,
it is safe to delete the “.colvars.state” files to save space.  However,
when a LAMMPS simulation is restarted using {\hyperref[\detokenize{read_data::doc}]{\sphinxcrossref{\DUrole{doc}{read\_data}}}}, the Colvars state file must be available and loaded via
the “input” keyword or via a “fix\_modify Colvars load” command (see
below).

\sphinxAtStartPar
When restarting, the fix and the Colvars module should be created and
configured using the original configuration file(s).


\subsection{Output}
\label{\detokenize{fix_colvars:output}}
\sphinxAtStartPar
This fix computes a global scalar which can be accessed by various
{\hyperref[\detokenize{Howto_output::doc}]{\sphinxcrossref{\DUrole{doc}{output commands}}}}.  The scalar is the energy due to
all external potentials defined in the Colvars configuration.  The
scalar value calculated by this fix is “extensive”.

\sphinxAtStartPar
Aside from the state information in a “.colvars.state” file, other
\sphinxhref{https://colvars.github.io/colvars-refman-lammps/colvars-refman-lammps.html\#sec:colvars\_output}{output files}
are produced by Colvars depending on the type of simulation.  For this
reason, the “output” keyword is required for fix colvars.


\subsection{Controlling Colvars via \sphinxtitleref{fix\_modify}}
\label{\detokenize{fix_colvars:controlling-colvars-via-fix-modify}}
\sphinxAtStartPar
\DUrole{versionmodified,added}{Added in version Colvars\_2023\sphinxhyphen{}06\sphinxhyphen{}04.}

\sphinxAtStartPar
The {\hyperref[\detokenize{fix_modify::doc}]{\sphinxcrossref{\DUrole{doc}{fix\_modify}}}} command may be used on “fix colvars”
in either one of two ways:
\begin{enumerate}
\sphinxsetlistlabels{\arabic}{enumi}{enumii}{(}{)}%
\item {} 
\sphinxAtStartPar
Provide updated values for the fix parameters, such as \sphinxstyleemphasis{output},
\sphinxstyleemphasis{input}, \sphinxstyleemphasis{unwrap}, \sphinxstyleemphasis{tstat} and \sphinxstyleemphasis{seed}.  Additionally, the
{\hyperref[\detokenize{fix_modify::doc}]{\sphinxcrossref{\DUrole{doc}{fix\_modify}}}} \sphinxstyleemphasis{energy} keyword is supported by this fix
to add the energy change from the biasing force added by Colvars to the
global potential energy of the system as part of {\hyperref[\detokenize{thermo_style::doc}]{\sphinxcrossref{\DUrole{doc}{thermodynamic
output}}}} (the default is {\hyperref[\detokenize{fix_modify::doc}]{\sphinxcrossref{\DUrole{doc}{fix\_modify energy no}}}}).  For example, in a multi\sphinxhyphen{}step LAMMPS script involving
multiple thermostats (e.g. fix nvt followed by fix npt), Colvars can
read a new thermostat’s target temperature like this:

\end{enumerate}
\begin{quote}

\begin{sphinxVerbatim}[commandchars=\\\{\}]
\PYG{k}{fix }\PYG{n+nv+nvIdentifier}{NVT}\PYG{+w}{ }\PYG{n+nv+nvIdentifier}{all}\PYG{+w}{ }\PYG{n}{nvt}\PYG{+w}{ }\PYG{n}{...}
\PYG{k}{fix }\PYG{n+nv+nvIdentifier}{Colvars}\PYG{+w}{ }\PYG{n+nv+nvIdentifier}{all}\PYG{+w}{ }\PYG{n}{colvars}\PYG{+w}{ }\PYG{o}{\PYGZlt{}}\PYG{n}{configfile}\PYG{o}{\PYGZgt{}}\PYG{+w}{ }\PYG{n}{output}\PYG{+w}{ }\PYG{n}{equil1}\PYG{+w}{ }\PYG{n}{tstat}\PYG{+w}{ }\PYG{n}{NVT}
\PYG{k}{run}\PYG{+w}{ }\PYG{o}{\PYGZlt{}}\PYG{n}{NUMSTEPS}\PYG{o}{\PYGZgt{}}
\PYG{k}{unfix }\PYG{n+nv+nvIdentifier}{nvt}
\PYG{k}{fix }\PYG{n+nv+nvIdentifier}{NPT}\PYG{+w}{ }\PYG{n+nv+nvIdentifier}{all}\PYG{+w}{ }\PYG{n}{n}\PYG{+w}{ }\PYG{n}{...}
\PYG{k}{fix\PYGZus{}modify }\PYG{n+nv+nvIdentifier}{Colvars}\PYG{+w}{ }\PYG{n}{tstat}\PYG{+w}{ }\PYG{n}{NPT}
\PYG{k}{fix\PYGZus{}modify }\PYG{n+nv+nvIdentifier}{Colvars}\PYG{+w}{ }\PYG{n}{output}\PYG{+w}{ }\PYG{n}{equil2}
\end{sphinxVerbatim}
\end{quote}
\begin{enumerate}
\sphinxsetlistlabels{\arabic}{enumi}{enumii}{(}{)}%
\setcounter{enumi}{1}
\item {} 
\sphinxAtStartPar
Call one of the scripting functions provided by the Colvars module
itself (a full list is available in the Colvars doc).  The arguments
to these functions are provided as strings and passed to Colvars.

\sphinxAtStartPar
LAMMPS variables referenced by their string representation
“\$\{variable\}” will be expanded immediately.  Note also that this
variable expansion \sphinxstyleemphasis{will also happen within quotes}, similar to what
the {\hyperref[\detokenize{mdi::doc}]{\sphinxcrossref{\DUrole{doc}{mdi}}}} command provides.  This feature makes it
possible to use the values of certain LAMMPS variables in Colvars
configuration strings.  For example, to synchronize the LAMMPS and
Colvars dump frequencies:

\end{enumerate}
\begin{quote}

\begin{sphinxVerbatim}[commandchars=\\\{\}]
\PYG{k}{variable }\PYG{n+nv+nvIdentifier}{freq}\PYG{+w}{ }\PYG{n}{index}\PYG{+w}{ }\PYG{l+m}{10000}
\PYG{k}{dump }\PYG{n+nv+nvIdentifier}{myDump}\PYG{+w}{ }\PYG{n+nv+nvIdentifier}{all}\PYG{+w}{ }\PYG{n}{atom}\PYG{o}{/}\PYG{n}{zstd}\PYG{+w}{ }\PYG{n+nv}{\PYGZdl{}\PYGZob{}}\PYG{n+nv}{freq}\PYG{n+nv}{\PYGZcb{}}\PYG{+w}{ }\PYG{n}{dump.atom.zstd}
\PYG{k}{fix\PYGZus{}modify }\PYG{n+nv+nvIdentifier}{Colvars}\PYG{+w}{ }\PYG{n}{config}\PYG{+w}{ }\PYG{l+s}{\PYGZdq{}}\PYG{l+s}{colvarsTrajFrequency \PYGZdl{}\PYGZob{}freq\PYGZcb{}}\PYG{l+s}{\PYGZdq{}}
\end{sphinxVerbatim}
\end{quote}

\begin{sphinxadmonition}{note}{Note:}
\sphinxAtStartPar
Although it is possible to use {\hyperref[\detokenize{fix_modify::doc}]{\sphinxcrossref{\DUrole{doc}{fix\_modify}}}} at any
time, its results will only reflect the state of the Colvars module
at the end of the most recent “run” or “minimize” command.  Any new
configuration added via “fix\_modify Colvars configfile” or
“fix\_modify Colvars config” will only be loaded when the simulation
resumes.  Configuration files or strings will be parsed in the same
sequence as they were provided in the LAMMPS script.
\end{sphinxadmonition}


\subsection{Restrictions}
\label{\detokenize{fix_colvars:restrictions}}
\sphinxAtStartPar
This fix is provided by the COLVARS package and is only available if
LAMMPS was built with that package (default in most builds).  Some of
the features also require code available from the LEPTON package.  See
the {\hyperref[\detokenize{Build_package::doc}]{\sphinxcrossref{\DUrole{doc}{Build package}}}} page for more info.

\sphinxAtStartPar
There can only be one Colvars instance defined at a time.  Since the
interface communicates only the minimum required amount of information,
and the Colvars module itself can handle an arbitrary number of
collective variables, this is not a limitation of functionality.


\subsection{Related commands}
\label{\detokenize{fix_colvars:related-commands}}
\sphinxAtStartPar
{\hyperref[\detokenize{fix_smd::doc}]{\sphinxcrossref{\DUrole{doc}{fix smd}}}}, {\hyperref[\detokenize{fix_spring::doc}]{\sphinxcrossref{\DUrole{doc}{fix spring}}}},
{\hyperref[\detokenize{fix_plumed::doc}]{\sphinxcrossref{\DUrole{doc}{fix plumed}}}}


\bigskip\hrule\bigskip

\phantomsection\label{\detokenize{fix_colvars:fiorin}}
\sphinxAtStartPar
\sphinxstylestrong{(Fiorin)} Fiorin, Klein, Henin, Mol. Phys. 111, 3345 (2013) \sphinxurl{https://doi.org/10.1080/00268976.2013.813594}

\phantomsection\label{\detokenize{fix_colvars:colvars-lammps-doc}}
\sphinxAtStartPar
\sphinxurl{https://colvars.github.io/colvars-refman-lammps/colvars-refman-lammps.html}

\sphinxstepscope

\index{fix controller@\spxentry{fix controller}}\ignorespaces 

\section{fix controller command}
\label{\detokenize{fix_controller:fix-controller-command}}\label{\detokenize{fix_controller:index-0}}\label{\detokenize{fix_controller::doc}}

\subsection{Syntax}
\label{\detokenize{fix_controller:syntax}}
\begin{sphinxVerbatim}[commandchars=\\\{\}]
\PYG{k}{fix }\PYG{n+nv+nvIdentifier}{ID}\PYG{+w}{ }\PYG{n+nv+nvIdentifier}{group\PYGZhy{}ID}\PYG{+w}{ }\PYG{n}{controller}\PYG{+w}{ }\PYG{n}{Nevery}\PYG{+w}{ }\PYG{n}{alpha}\PYG{+w}{ }\PYG{n}{Kp}\PYG{+w}{ }\PYG{n}{Ki}\PYG{+w}{ }\PYG{n}{Kd}\PYG{+w}{ }\PYG{n}{pvar}\PYG{+w}{ }\PYG{n}{setpoint}\PYG{+w}{ }\PYG{n}{cvar}
\end{sphinxVerbatim}
\begin{itemize}
\item {} 
\sphinxAtStartPar
ID, group\sphinxhyphen{}ID are documented in {\hyperref[\detokenize{fix::doc}]{\sphinxcrossref{\DUrole{doc}{fix}}}} command

\item {} 
\sphinxAtStartPar
controller = style name of this fix command

\item {} 
\sphinxAtStartPar
Nevery = invoke controller every this many timesteps

\item {} 
\sphinxAtStartPar
alpha = coupling constant for PID equation (see units discussion below)

\item {} 
\sphinxAtStartPar
Kp = proportional gain in PID equation (unitless)

\item {} 
\sphinxAtStartPar
Ki = integral gain in PID equation (unitless)

\item {} 
\sphinxAtStartPar
Kd = derivative gain in PID equation (unitless)

\item {} 
\sphinxAtStartPar
pvar = process variable of form c\_ID, c\_ID{[}I{]}, f\_ID, f\_ID{[}I{]}, or v\_name

\begin{sphinxVerbatim}[commandchars=\\\{\}]
c\PYGZus{}ID = global scalar calculated by a compute with ID
c\PYGZus{}ID[I] = Ith component of global vector calculated by a compute with ID
f\PYGZus{}ID = global scalar calculated by a fix with ID
f\PYGZus{}ID[I] = Ith component of global vector calculated by a fix with ID
v\PYGZus{}name = value calculated by an equal\PYGZhy{}style variable with name
\end{sphinxVerbatim}

\item {} 
\sphinxAtStartPar
setpoint = desired value of process variable (same units as process variable)

\item {} 
\sphinxAtStartPar
cvar = name of control variable

\end{itemize}


\subsection{Examples}
\label{\detokenize{fix_controller:examples}}
\begin{sphinxVerbatim}[commandchars=\\\{\}]
\PYG{k}{fix }\PYG{n+nv+nvIdentifier}{1}\PYG{+w}{ }\PYG{n+nv+nvIdentifier}{all}\PYG{+w}{ }\PYG{n}{controller}\PYG{+w}{ }\PYG{l+m}{100}\PYG{+w}{ }\PYG{l+m}{1.0}\PYG{+w}{ }\PYG{l+m}{0.5}\PYG{+w}{ }\PYG{l+m}{0.0}\PYG{+w}{ }\PYG{l+m}{0.0}\PYG{+w}{ }\PYG{n}{c\PYGZus{}thermo\PYGZus{}temp}\PYG{+w}{ }\PYG{l+m}{1.5}\PYG{+w}{ }\PYG{n}{tcontrol}
\PYG{k}{fix }\PYG{n+nv+nvIdentifier}{1}\PYG{+w}{ }\PYG{n+nv+nvIdentifier}{all}\PYG{+w}{ }\PYG{n}{controller}\PYG{+w}{ }\PYG{l+m}{100}\PYG{+w}{ }\PYG{l+m}{0.2}\PYG{+w}{ }\PYG{l+m}{0.5}\PYG{+w}{ }\PYG{l+m}{0}\PYG{+w}{ }\PYG{l+m}{100.0}\PYG{+w}{ }\PYG{n}{v\PYGZus{}pxxwall}\PYG{+w}{ }\PYG{l+m}{1.01325}\PYG{+w}{ }\PYG{n}{xwall}
\PYG{k}{fix }\PYG{n+nv+nvIdentifier}{1}\PYG{+w}{ }\PYG{n+nv+nvIdentifier}{all}\PYG{+w}{ }\PYG{n}{controller}\PYG{+w}{ }\PYG{l+m}{10000}\PYG{+w}{ }\PYG{l+m}{0.2}\PYG{+w}{ }\PYG{l+m}{0.5}\PYG{+w}{ }\PYG{l+m}{0}\PYG{+w}{ }\PYG{l+m}{2000}\PYG{+w}{ }\PYG{n}{v\PYGZus{}avpe}\PYG{+w}{ }\PYG{o}{\PYGZhy{}}\PYG{l+m}{3.785}\PYG{+w}{ }\PYG{n}{tcontrol}
\end{sphinxVerbatim}


\subsection{Description}
\label{\detokenize{fix_controller:description}}
\sphinxAtStartPar
This fix enables control of a LAMMPS simulation using a control loop
feedback mechanism known as a proportional\sphinxhyphen{}integral\sphinxhyphen{}derivative (PID)
controller.  The basic idea is to define a “process variable” which is
a quantity that can be monitored during a running simulation.  A
desired target value is chosen for the process variable.  A “control
variable” is also defined which is an adjustable attribute of the
running simulation, which the process variable will respond to.  The
PID controller continuously adjusts the control variable based on the
difference between the process variable and the target.

\sphinxAtStartPar
Here are examples of ways in which this fix can be used.  The
examples/pid directory contains a script that implements the simple
thermostat.


\begin{savenotes}\sphinxattablestart
\sphinxthistablewithglobalstyle
\centering
\begin{tabulary}{\linewidth}[t]{TTT}
\sphinxtoprule
\sphinxtableatstartofbodyhook
\sphinxAtStartPar
Goal
&
\sphinxAtStartPar
process variable
&
\sphinxAtStartPar
control variable
\\
\sphinxhline
\sphinxAtStartPar
Simple thermostat
&
\sphinxAtStartPar
instantaneous T
&
\sphinxAtStartPar
thermostat target T
\\
\sphinxhline
\sphinxAtStartPar
Find melting temperature
&
\sphinxAtStartPar
average PE per atom
&
\sphinxAtStartPar
thermostat target T
\\
\sphinxhline
\sphinxAtStartPar
Control pressure in non\sphinxhyphen{}periodic system
&
\sphinxAtStartPar
force on wall
&
\sphinxAtStartPar
position of wall
\\
\sphinxhline&&\\
\sphinxbottomrule
\end{tabulary}
\sphinxtableafterendhook\par
\sphinxattableend\end{savenotes}

\begin{sphinxadmonition}{note}{Note:}
\sphinxAtStartPar
For this fix to work, the control variable must actually induce
a change in a running LAMMPS simulation.  Typically this will only
occur if there is some other command (e.g. a thermostat fix) which
uses the control variable as an input parameter.  This could be done
directly or indirectly, e.g. the other command uses a variable as
input whose formula uses the control variable.  The other command
should alter its behavior dynamically as the variable changes.
\end{sphinxadmonition}

\begin{sphinxadmonition}{note}{Note:}
\sphinxAtStartPar
If there is a command you think could be used in this fashion,
but does not currently allow a variable as an input parameter, please
notify the LAMMPS developers.  It is often not difficult to enable a
command to use a variable as an input parameter.
\end{sphinxadmonition}

\sphinxAtStartPar
The group specified with this command is ignored.  However, note that
the process variable may be defined by calculations performed by
computes and fixes which store their own “group” definitions.

\sphinxAtStartPar
The PID controller is invoked once each \sphinxstyleemphasis{Nevery} timesteps.

\sphinxAtStartPar
The PID controller is implemented as a discretized version of
the following dynamic equation:
\begin{equation*}
\begin{split}\frac{dc}{dt}  = -\alpha (K_p e + K_i \int_0^t e \, dt + K_d \frac{de}{dt} )\end{split}
\end{equation*}
\sphinxAtStartPar
where \sphinxstyleemphasis{c} is the continuous time analog of the control variable,
\sphinxstyleemphasis{e} =\sphinxstyleemphasis{pvar}\sphinxhyphen{}\sphinxstyleemphasis{setpoint} is the error in the process variable, and
\(\alpha\), \(K_p\), \(K_i\) , and \(K_d\) are constants
set by the corresponding
keywords described above. The discretized version of this equation is:
\begin{equation*}
\begin{split}c_n  = c_{n-1} -\alpha \left( K_p \tau e_n + K_i \tau^2 \sum_{i=1}^n e_i + K_d (e_n - e_{n-1}) \right)\end{split}
\end{equation*}
\sphinxAtStartPar
where \(\tau = \mathtt{Nevery} \cdot \mathtt{timestep}\) is the time
interval between updates,
and the subscripted variables indicate the values of \sphinxstyleemphasis{c} and \sphinxstyleemphasis{e} at
successive updates.

\sphinxAtStartPar
From the first equation, it is clear that if the three gain values
\(K_p\), \(K_i\), \(K_d\) are dimensionless constants,
then \(\alpha\) must have
units of {[}unit \sphinxstyleemphasis{cvar}{]}/{[}unit \sphinxstyleemphasis{pvar}{]}/{[}unit time{]} e.g. {[} eV/K/ps
{]}. The advantage of this unit scheme is that the value of the
constants should be invariant under a change of either the MD timestep
size or the value of \sphinxstyleemphasis{Nevery}. Similarly, if the LAMMPS {\hyperref[\detokenize{units::doc}]{\sphinxcrossref{\DUrole{doc}{unit style}}}} is changed, it should only be necessary to change
the value of \(\alpha\) to reflect this, while leaving \(K_p\),
\(K_i\), and \(K_d\) unaltered.

\sphinxAtStartPar
When choosing the values of the four constants, it is best to first
pick a value and sign for \(\alpha\) that is consistent with the
magnitudes and signs of \sphinxstyleemphasis{pvar} and \sphinxstyleemphasis{cvar}.  The magnitude of \(K_p\)
should then be tested over a large positive range keeping \(K_i = K_d =0\).
A good value for \(K_p\) will produce a fast response in \sphinxstyleemphasis{pvar},
without overshooting the \sphinxstyleemphasis{setpoint}.  For many applications, proportional
feedback is sufficient, and so \(K_i = K_d =0\) can be used. In cases
where there is a substantial lag time in the response of \sphinxstyleemphasis{pvar} to a change
in \sphinxstyleemphasis{cvar}, this can be counteracted by increasing \(K_d\). In situations
where \sphinxstyleemphasis{pvar} plateaus without reaching \sphinxstyleemphasis{setpoint}, this can be
counteracted by increasing \(K_i\).  In the language of Charles Dickens,
\(K_p\) represents the error of the present, \(K_i\) the error of
the past, and \(K_d\) the error yet to come.

\sphinxAtStartPar
Because this fix updates \sphinxstyleemphasis{cvar}, but does not initialize its value,
the initial value \(c_0\) is that assigned by the user in the input script via
the {\hyperref[\detokenize{variable::doc}]{\sphinxcrossref{\DUrole{doc}{internal\sphinxhyphen{}style variable}}}} command.  This value is
used (by every other LAMMPS command that uses the variable) until this
fix performs its first update of \sphinxstyleemphasis{cvar} after \sphinxstyleemphasis{Nevery} timesteps.  On
the first update, the value of the derivative term is set to zero,
because the value of \(e_{n-1}\) is not yet defined.


\bigskip\hrule\bigskip


\sphinxAtStartPar
The process variable \sphinxstyleemphasis{pvar} can be specified as the output of a
{\hyperref[\detokenize{compute::doc}]{\sphinxcrossref{\DUrole{doc}{compute}}}} or {\hyperref[\detokenize{fix::doc}]{\sphinxcrossref{\DUrole{doc}{fix}}}} or the evaluation of a
{\hyperref[\detokenize{variable::doc}]{\sphinxcrossref{\DUrole{doc}{variable}}}}.  In each case, the compute, fix, or variable
must produce a global quantity, not a per\sphinxhyphen{}atom or local quantity.

\sphinxAtStartPar
If \sphinxstyleemphasis{pvar} begins with “c\_”, a compute ID must follow which has been
previously defined in the input script and which generates a global
scalar or vector.  See the individual {\hyperref[\detokenize{compute::doc}]{\sphinxcrossref{\DUrole{doc}{compute}}}} doc page
for details.  If no bracketed integer is appended, the scalar
calculated by the compute is used.  If a bracketed integer is
appended, the Ith value of the vector calculated by the compute is
used.  Users can also write code for their own compute styles and {\hyperref[\detokenize{Modify::doc}]{\sphinxcrossref{\DUrole{doc}{add them to LAMMPS}}}}.

\sphinxAtStartPar
If \sphinxstyleemphasis{pvar} begins with “f\_”, a fix ID must follow which has been
previously defined in the input script and which generates a global
scalar or vector.  See the individual {\hyperref[\detokenize{fix::doc}]{\sphinxcrossref{\DUrole{doc}{fix}}}} page for
details.  Note that some fixes only produce their values on certain
timesteps, which must be compatible with when fix controller
references the values, or else an error results.  If no bracketed integer
is appended, the scalar calculated by the fix is used.  If a bracketed
integer is appended, the Ith value of the vector calculated by the fix
is used.  Users can also write code for their own fix style and {\hyperref[\detokenize{Modify::doc}]{\sphinxcrossref{\DUrole{doc}{add them to LAMMPS}}}}.

\sphinxAtStartPar
If \sphinxstyleemphasis{pvar} begins with “v\_”, a variable name must follow which has been
previously defined in the input script.  Only equal\sphinxhyphen{}style variables
can be referenced.  See the {\hyperref[\detokenize{variable::doc}]{\sphinxcrossref{\DUrole{doc}{variable}}}} command for
details.  Note that variables of style \sphinxstyleemphasis{equal} define a formula which
can reference individual atom properties or thermodynamic keywords, or
they can invoke other computes, fixes, or variables when they are
evaluated, so this is a very general means of specifying the process
variable.

\sphinxAtStartPar
The target value \sphinxstyleemphasis{setpoint} for the process variable must be a numeric
value, in whatever units \sphinxstyleemphasis{pvar} is defined for.

\sphinxAtStartPar
The control variable \sphinxstyleemphasis{cvar} must be the name of an {\hyperref[\detokenize{variable::doc}]{\sphinxcrossref{\DUrole{doc}{internal\sphinxhyphen{}style variable}}}} previously defined in the input script.  Note
that it is not specified with a “v\_” prefix, just the name of the
variable.  It must be an internal\sphinxhyphen{}style variable, because this fix
updates its value directly.  Note that other commands can use an
equal\sphinxhyphen{}style versus internal\sphinxhyphen{}style variable interchangeably.


\bigskip\hrule\bigskip



\subsection{Restart, fix\_modify, output, run start/stop, minimize info}
\label{\detokenize{fix_controller:restart-fix-modify-output-run-start-stop-minimize-info}}
\sphinxAtStartPar
Currently, no information about this fix is written to {\hyperref[\detokenize{restart::doc}]{\sphinxcrossref{\DUrole{doc}{binary restart files}}}}.  None of the {\hyperref[\detokenize{fix_modify::doc}]{\sphinxcrossref{\DUrole{doc}{fix\_modify}}}} options
are relevant to this fix.

\sphinxAtStartPar
This fix produces a global vector with 3 values which can be accessed
by various {\hyperref[\detokenize{Howto_output::doc}]{\sphinxcrossref{\DUrole{doc}{output commands}}}}.  The values can be
accessed on any timestep, though they are only updated on timesteps
that are a multiple of \sphinxstyleemphasis{Nevery}.

\sphinxAtStartPar
The three values are the most recent updates made to the control
variable by each of the 3 terms in the PID equation above.  The first
value is the proportional term, the second is the integral term, the
third is the derivative term.

\sphinxAtStartPar
The units of the vector values will be whatever units the control
variable is in.  The vector values calculated by this fix are
“extensive”.

\sphinxAtStartPar
No parameter of this fix can be used with the \sphinxstyleemphasis{start/stop} keywords of
the {\hyperref[\detokenize{run::doc}]{\sphinxcrossref{\DUrole{doc}{run}}}} command.  This fix is not invoked during {\hyperref[\detokenize{minimize::doc}]{\sphinxcrossref{\DUrole{doc}{energy minimization}}}}.


\subsection{Restrictions}
\label{\detokenize{fix_controller:restrictions}}
\sphinxAtStartPar
This fix is part of the EXTRA\sphinxhyphen{}FIX package.  It is only enabled if LAMMPS
was built with that package.  See the {\hyperref[\detokenize{Build_package::doc}]{\sphinxcrossref{\DUrole{doc}{Build package}}}} page for more info.


\subsection{Related commands}
\label{\detokenize{fix_controller:related-commands}}
\sphinxAtStartPar
{\hyperref[\detokenize{fix_adapt::doc}]{\sphinxcrossref{\DUrole{doc}{fix adapt}}}}


\subsection{Default}
\label{\detokenize{fix_controller:default}}
\sphinxAtStartPar
none

\sphinxstepscope

\index{fix damping/cundall@\spxentry{fix damping/cundall}}\ignorespaces 

\section{fix damping/cundall command}
\label{\detokenize{fix_damping_cundall:fix-damping-cundall-command}}\label{\detokenize{fix_damping_cundall:index-0}}\label{\detokenize{fix_damping_cundall::doc}}

\subsection{Syntax}
\label{\detokenize{fix_damping_cundall:syntax}}
\begin{sphinxVerbatim}[commandchars=\\\{\}]
\PYG{k}{fix }\PYG{n+nv+nvIdentifier}{ID}\PYG{+w}{ }\PYG{n+nv+nvIdentifier}{group\PYGZhy{}ID}\PYG{+w}{ }\PYG{n}{damping}\PYG{o}{/}\PYG{n}{cundall}\PYG{+w}{ }\PYG{n}{gamma\PYGZus{}l}\PYG{+w}{ }\PYG{n}{gamma\PYGZus{}a}\PYG{+w}{ }\PYG{n}{keyword}\PYG{+w}{ }\PYG{n}{values}\PYG{+w}{ }\PYG{n}{...}
\end{sphinxVerbatim}
\begin{itemize}
\item {} 
\sphinxAtStartPar
ID, group\sphinxhyphen{}ID are documented in {\hyperref[\detokenize{fix::doc}]{\sphinxcrossref{\DUrole{doc}{fix}}}} command

\item {} 
\sphinxAtStartPar
damping/cundall = style name of this fix command

\item {} 
\sphinxAtStartPar
gamma\_l = linear damping coefficient (dimensionless)

\item {} 
\sphinxAtStartPar
gamma\_a = angular damping coefficient (dimensionless)

\item {} 
\sphinxAtStartPar
zero or more keyword/value pairs may be appended
\begin{sphinxalltt}
keyword = \sphinxstyleemphasis{scale}
  \sphinxstyleemphasis{scale} values = \sphinxstyleemphasis{type ratio} or \sphinxstyleemphasis{v\_name}
    type = atom type (1\sphinxhyphen{}N)
    ratio = factor to scale the damping coefficients by
    v\_name = reference to atom style variable \sphinxstyleemphasis{name}
\end{sphinxalltt}

\end{itemize}


\subsection{Examples}
\label{\detokenize{fix_damping_cundall:examples}}
\begin{sphinxVerbatim}[commandchars=\\\{\}]
\PYG{k}{fix }\PYG{n+nv+nvIdentifier}{1}\PYG{+w}{ }\PYG{n+nv+nvIdentifier}{all}\PYG{+w}{ }\PYG{n}{damping}\PYG{o}{/}\PYG{n}{cundall}\PYG{+w}{ }\PYG{l+m}{0.8}\PYG{+w}{ }\PYG{l+m}{0.8}
\PYG{k}{fix }\PYG{n+nv+nvIdentifier}{1}\PYG{+w}{ }\PYG{n+nv+nvIdentifier}{all}\PYG{+w}{ }\PYG{n}{damping}\PYG{o}{/}\PYG{n}{cundall}\PYG{+w}{ }\PYG{l+m}{0.8}\PYG{+w}{ }\PYG{l+m}{0.5}\PYG{+w}{ }\PYG{n}{scale}\PYG{+w}{ }\PYG{l+m}{3}\PYG{+w}{ }\PYG{l+m}{2.5}
\PYG{k}{fix }\PYG{n+nv+nvIdentifier}{a}\PYG{+w}{ }\PYG{n+nv+nvIdentifier}{all}\PYG{+w}{ }\PYG{n}{damping}\PYG{o}{/}\PYG{n}{cundall}\PYG{+w}{ }\PYG{l+m}{0.8}\PYG{+w}{ }\PYG{l+m}{0.5}\PYG{+w}{ }\PYG{n}{scale}\PYG{+w}{ }\PYG{n}{v\PYGZus{}radscale}
\end{sphinxVerbatim}


\subsection{Description}
\label{\detokenize{fix_damping_cundall:description}}
\sphinxAtStartPar
Add damping force and torque to finite\sphinxhyphen{}size spherical particles in the group
following the model of {\hyperref[\detokenize{fix_damping_cundall:cundall1987}]{\sphinxcrossref{\DUrole{std,std-ref}{Cundall, 1987}}}}, as implemented in other
granular physics code (e.g., {\hyperref[\detokenize{fix_damping_cundall:yadedem}]{\sphinxcrossref{\DUrole{std,std-ref}{Yade\sphinxhyphen{}DEM}}}}, {\hyperref[\detokenize{fix_damping_cundall:pfc}]{\sphinxcrossref{\DUrole{std,std-ref}{PFC}}}}).

\sphinxAtStartPar
The damping is constructed to always have negative mechanical power with respect
to the current velocity/angular velocity to ensure dissipation of kinetic energy.
If used without additional thermostatting (to add kinetic energy to the system),
it has the effect of slowly (or rapidly) freezing the system; hence it can also
be used as a simple energy minimization technique.

\sphinxAtStartPar
The magnitude of the damping force/torque \(F_d\)/\(T_d\) is a fraction
\(\gamma \in [0;1]\) of the current force/torque \(F\)/\(T\) on the
particle. Damping is applied component\sphinxhyphen{}by\sphinxhyphen{}component in each direction
\(k\in\{x, y, z\}\):
\begin{equation*}
\begin{split}{F_d}_k = - \gamma_l \, F_k \, \mathrm{sign}(F_k v_k)\end{split}
\end{equation*}\begin{equation*}
\begin{split}{T_d}_k = - \gamma_a \, T_k \, \mathrm{sign}(T_k \omega_k)\end{split}
\end{equation*}
\sphinxAtStartPar
The larger the coefficients, the faster the kinetic energy is reduced.

\sphinxAtStartPar
If the optional keyword \sphinxstyleemphasis{scale} is used, \(\gamma_l\) and \(\gamma_a\)
can be scaled up or down by the specified factor for atoms.  This factor can be
set for different atom types and thus the \sphinxstyleemphasis{scale} keyword used multiple times
followed by the atom type and the associated scale factor.  Alternately the
scaling factor can be computed for each atom (e.g. based on its radius) by
using an {\hyperref[\detokenize{variable::doc}]{\sphinxcrossref{\DUrole{doc}{atom\sphinxhyphen{}style variable}}}}.

\begin{sphinxadmonition}{note}{Note:}
\sphinxAtStartPar
The damping force/torque is computed based on the force/torque at the moment
this fix is invoked. Any force/torque added after this fix, e.g., by
{\hyperref[\detokenize{fix_addforce::doc}]{\sphinxcrossref{\DUrole{doc}{fix addforce}}}} or {\hyperref[\detokenize{fix_addtorque::doc}]{\sphinxcrossref{\DUrole{doc}{fix addtorque}}}}
will not be damped. When performing simulations with gravity, invoking
{\hyperref[\detokenize{fix_gravity::doc}]{\sphinxcrossref{\DUrole{doc}{fix gravity}}}} after this fix will maintain the specified
gravitational acceleration.
\end{sphinxadmonition}

\begin{sphinxadmonition}{note}{Note:}
\sphinxAtStartPar
This scheme is dependent on the coordinates system and does not correspond to
realistic physical processes. It is constructed for numerical convenience and
efficacy.
\end{sphinxadmonition}

\sphinxAtStartPar
This non\sphinxhyphen{}viscous damping presents the following advantages:
\begin{enumerate}
\sphinxsetlistlabels{\arabic}{enumi}{enumii}{}{.}%
\item {} 
\sphinxAtStartPar
damping is independent of velocity, equally damping regions with
distinct natural frequencies,

\item {} 
\sphinxAtStartPar
damping affects acceleration and vanishes for steady uniform motion
of the particles,

\item {} 
\sphinxAtStartPar
damping parameter \(\gamma\) is dimensionless and does not
require scaling.

\end{enumerate}


\bigskip\hrule\bigskip



\subsection{Restart, fix\_modify, output, run start/stop, minimize info}
\label{\detokenize{fix_damping_cundall:restart-fix-modify-output-run-start-stop-minimize-info}}
\sphinxAtStartPar
No information about this fix is written to {\hyperref[\detokenize{restart::doc}]{\sphinxcrossref{\DUrole{doc}{binary restart files}}}}.  No global or per\sphinxhyphen{}atom quantities are stored by this fix for
access by various {\hyperref[\detokenize{Howto_output::doc}]{\sphinxcrossref{\DUrole{doc}{output commands}}}}.  No parameter
of this fix can be used with the \sphinxstyleemphasis{start/stop} keywords of the {\hyperref[\detokenize{run::doc}]{\sphinxcrossref{\DUrole{doc}{run}}}} command.

\sphinxAtStartPar
The {\hyperref[\detokenize{fix_modify::doc}]{\sphinxcrossref{\DUrole{doc}{fix\_modify}}}} \sphinxstyleemphasis{respa} option is supported by this
fix. This allows to set at which level of the {\hyperref[\detokenize{run_style::doc}]{\sphinxcrossref{\DUrole{doc}{r\sphinxhyphen{}RESPA}}}}
integrator the fix is modifying forces/torques.  Default is the outermost
level.

\sphinxAtStartPar
The forces/torques due to this fix are imposed during an energy
minimization, invoked by the {\hyperref[\detokenize{minimize::doc}]{\sphinxcrossref{\DUrole{doc}{minimize}}}} command.  This
fix should only be used with damped dynamics minimizers that allow for
non\sphinxhyphen{}conservative forces.  See the {\hyperref[\detokenize{min_style::doc}]{\sphinxcrossref{\DUrole{doc}{min\_style}}}} command
for details.


\subsection{Restrictions}
\label{\detokenize{fix_damping_cundall:restrictions}}
\sphinxAtStartPar
This fix is part of the GRANULAR package.  It is only enabled if LAMMPS
was built with that package.  See the {\hyperref[\detokenize{Build_package::doc}]{\sphinxcrossref{\DUrole{doc}{Build package}}}} page for more info.

\sphinxAtStartPar
This fix requires that atoms store torque and a radius as defined by the
{\hyperref[\detokenize{atom_style::doc}]{\sphinxcrossref{\DUrole{doc}{atom\_style sphere}}}} command.


\subsection{Related commands}
\label{\detokenize{fix_damping_cundall:related-commands}}
\sphinxAtStartPar
{\hyperref[\detokenize{fix_viscous::doc}]{\sphinxcrossref{\DUrole{doc}{fix viscous}}}}, {\hyperref[\detokenize{fix_viscous_sphere::doc}]{\sphinxcrossref{\DUrole{doc}{fix viscous/sphere}}}}


\subsection{Default}
\label{\detokenize{fix_damping_cundall:default}}
\sphinxAtStartPar
none


\subsection{References}
\label{\detokenize{fix_damping_cundall:references}}\phantomsection\label{\detokenize{fix_damping_cundall:cundall1987}}
\sphinxAtStartPar
\sphinxstylestrong{(Cundall, 1987)} Cundall, P. A. Distinct Element Models of Rock and Soil
Structure, in Analytical and Computational Methods in Engineering Rock
Mechanics, Ch. 4, pp. 129\sphinxhyphen{}163. E. T. Brown, ed. London: Allen \& Unwin., 1987.

\phantomsection\label{\detokenize{fix_damping_cundall:pfc}}
\sphinxAtStartPar
\sphinxstylestrong{(PFC)} PFC Particle Flow Code 6.0 Documentation. Itasca Consulting Group.

\phantomsection\label{\detokenize{fix_damping_cundall:yadedem}}
\sphinxAtStartPar
\sphinxstylestrong{(Yade\sphinxhyphen{}DEM)} V. Smilauer et al. (2021), Yade Documentation 3rd ed.
The Yade Project. DOI:10.5281/zenodo.5705394 (https://yade\sphinxhyphen{}dem.org/doc/)

\sphinxstepscope

\index{fix deform@\spxentry{fix deform}}\ignorespaces 
\index{fix deform/kk@\spxentry{fix deform/kk}}\ignorespaces 

\section{fix deform command}
\label{\detokenize{fix_deform:fix-deform-command}}\label{\detokenize{fix_deform:index-1}}\label{\detokenize{fix_deform:index-0}}\label{\detokenize{fix_deform::doc}}
\sphinxAtStartPar
Accelerator Variants: \sphinxstyleemphasis{deform/kk}


\subsection{Syntax}
\label{\detokenize{fix_deform:syntax}}
\begin{sphinxVerbatim}[commandchars=\\\{\}]
\PYG{k}{fix }\PYG{n+nv+nvIdentifier}{ID}\PYG{+w}{ }\PYG{n+nv+nvIdentifier}{group\PYGZhy{}ID}\PYG{+w}{ }\PYG{n}{deform}\PYG{+w}{ }\PYG{n}{N}\PYG{+w}{ }\PYG{n}{parameter}\PYG{+w}{ }\PYG{n}{style}\PYG{+w}{ }\PYG{n}{args}\PYG{+w}{ }\PYG{n}{...}\PYG{+w}{ }\PYG{n}{keyword}\PYG{+w}{ }\PYG{n}{value}\PYG{+w}{ }\PYG{n}{...}
\end{sphinxVerbatim}
\begin{itemize}
\item {} 
\sphinxAtStartPar
ID, group\sphinxhyphen{}ID are documented in {\hyperref[\detokenize{fix::doc}]{\sphinxcrossref{\DUrole{doc}{fix}}}} command

\item {} 
\sphinxAtStartPar
N = perform box deformation every this many timesteps

\item {} 
\sphinxAtStartPar
one or more parameter/args sequences may be appended
\begin{sphinxalltt}
parameter = \sphinxstyleemphasis{x} or \sphinxstyleemphasis{y} or \sphinxstyleemphasis{z} or \sphinxstyleemphasis{xy} or \sphinxstyleemphasis{xz} or \sphinxstyleemphasis{yz}
  \sphinxstyleemphasis{x}, \sphinxstyleemphasis{y}, \sphinxstyleemphasis{z} args = style value(s)
    style = \sphinxstyleemphasis{final} or \sphinxstyleemphasis{delta} or \sphinxstyleemphasis{scale} or \sphinxstyleemphasis{vel} or \sphinxstyleemphasis{erate} or \sphinxstyleemphasis{trate} or \sphinxstyleemphasis{volume} or \sphinxstyleemphasis{wiggle} or \sphinxstyleemphasis{variable} or \sphinxstyleemphasis{pressure} or \sphinxstyleemphasis{pressure/mean}
      \sphinxstyleemphasis{final} values = lo hi
        lo hi = box boundaries at end of run (distance units)
      \sphinxstyleemphasis{delta} values = dlo dhi
        dlo dhi = change in box boundaries at end of run (distance units)
      \sphinxstyleemphasis{scale} values = factor
        factor = multiplicative factor for change in box length at end of run
      \sphinxstyleemphasis{vel} value = V
        V = change box length at this velocity (distance/time units),
            effectively an engineering strain rate
      \sphinxstyleemphasis{erate} value = R
        R = engineering strain rate (1/time units)
      \sphinxstyleemphasis{trate} value = R
        R = true strain rate (1/time units)
      \sphinxstyleemphasis{volume} value = none = adjust this dim to preserve volume of system
      \sphinxstyleemphasis{wiggle} values = A Tp
        A = amplitude of oscillation (distance units)
        Tp = period of oscillation (time units)
      \sphinxstyleemphasis{variable} values = v\_name1 v\_name2
        v\_name1 = variable with name1 for box length change as function of time
        v\_name2 = variable with name2 for change rate as function of time

  \sphinxstyleemphasis{xy}, \sphinxstyleemphasis{xz}, \sphinxstyleemphasis{yz} args = style value
    style = \sphinxstyleemphasis{final} or \sphinxstyleemphasis{delta} or \sphinxstyleemphasis{vel} or \sphinxstyleemphasis{erate} or \sphinxstyleemphasis{trate} or \sphinxstyleemphasis{wiggle} or \sphinxstyleemphasis{variable}
      \sphinxstyleemphasis{final} value = tilt
        tilt = tilt factor at end of run (distance units)
      \sphinxstyleemphasis{delta} value = dtilt
        dtilt = change in tilt factor at end of run (distance units)
      \sphinxstyleemphasis{vel} value = V
        V = change tilt factor at this velocity (distance/time units),
            effectively an engineering shear strain rate
      \sphinxstyleemphasis{erate} value = R
        R = engineering shear strain rate (1/time units)
      \sphinxstyleemphasis{trate} value = R
        R = true shear strain rate (1/time units)
      \sphinxstyleemphasis{wiggle} values = A Tp
        A = amplitude of oscillation (distance units)
        Tp = period of oscillation (time units)
      \sphinxstyleemphasis{variable} values = v\_name1 v\_name2
        v\_name1 = variable with name1 for tilt change as function of time
        v\_name2 = variable with name2 for change rate as function of time
\end{sphinxalltt}

\item {} 
\sphinxAtStartPar
zero or more keyword/value pairs may be appended

\item {} 
\sphinxAtStartPar
keyword = \sphinxstyleemphasis{remap} or \sphinxstyleemphasis{flip} or \sphinxstyleemphasis{units} or \sphinxstyleemphasis{couple} or \sphinxstyleemphasis{vol/balance/p} or \sphinxstyleemphasis{max/rate} or \sphinxstyleemphasis{normalize/pressure}
\begin{sphinxalltt}
\sphinxstyleemphasis{remap} value = \sphinxstyleemphasis{x} or \sphinxstyleemphasis{v} or \sphinxstyleemphasis{none}
  x = remap coords of atoms in group into deforming box
  v = remap velocities of atoms in group when they cross periodic boundaries
  none = no remapping of x or v
\sphinxstyleemphasis{flip} value = \sphinxstyleemphasis{yes} or \sphinxstyleemphasis{no}
  allow or disallow box flips when it becomes highly skewed
\sphinxstyleemphasis{units} value = \sphinxstyleemphasis{lattice} or \sphinxstyleemphasis{box}
  lattice = distances are defined in lattice units
  box = distances are defined in simulation box units
\end{sphinxalltt}

\end{itemize}


\subsection{Examples}
\label{\detokenize{fix_deform:examples}}
\begin{sphinxVerbatim}[commandchars=\\\{\}]
\PYG{k}{fix }\PYG{n+nv+nvIdentifier}{1}\PYG{+w}{ }\PYG{n+nv+nvIdentifier}{all}\PYG{+w}{ }\PYG{n}{deform}\PYG{+w}{ }\PYG{l+m}{1}\PYG{+w}{ }\PYG{n}{x}\PYG{+w}{ }\PYG{n}{final}\PYG{+w}{ }\PYG{l+m}{0.0}\PYG{+w}{ }\PYG{l+m}{9.0}\PYG{+w}{ }\PYG{n}{z}\PYG{+w}{ }\PYG{n}{final}\PYG{+w}{ }\PYG{l+m}{0.0}\PYG{+w}{ }\PYG{l+m}{5.0}\PYG{+w}{ }\PYG{n}{units}\PYG{+w}{ }\PYG{n}{box}
\PYG{k}{fix }\PYG{n+nv+nvIdentifier}{1}\PYG{+w}{ }\PYG{n+nv+nvIdentifier}{all}\PYG{+w}{ }\PYG{n}{deform}\PYG{+w}{ }\PYG{l+m}{1}\PYG{+w}{ }\PYG{n}{x}\PYG{+w}{ }\PYG{n}{trate}\PYG{+w}{ }\PYG{l+m}{0.1}\PYG{+w}{ }\PYG{n}{y}\PYG{+w}{ }\PYG{n}{volume}\PYG{+w}{ }\PYG{n}{z}\PYG{+w}{ }\PYG{n}{volume}
\PYG{k}{fix }\PYG{n+nv+nvIdentifier}{1}\PYG{+w}{ }\PYG{n+nv+nvIdentifier}{all}\PYG{+w}{ }\PYG{n}{deform}\PYG{+w}{ }\PYG{l+m}{1}\PYG{+w}{ }\PYG{n}{xy}\PYG{+w}{ }\PYG{n}{erate}\PYG{+w}{ }\PYG{l+m}{0.001}\PYG{+w}{ }\PYG{n}{remap}\PYG{+w}{ }\PYG{n}{v}
\PYG{k}{fix }\PYG{n+nv+nvIdentifier}{1}\PYG{+w}{ }\PYG{n+nv+nvIdentifier}{all}\PYG{+w}{ }\PYG{n}{deform}\PYG{+w}{ }\PYG{l+m}{10}\PYG{+w}{ }\PYG{n}{y}\PYG{+w}{ }\PYG{n}{delta}\PYG{+w}{ }\PYG{o}{\PYGZhy{}}\PYG{l+m}{0.5}\PYG{+w}{ }\PYG{l+m}{0.5}\PYG{+w}{ }\PYG{n}{xz}\PYG{+w}{ }\PYG{n}{vel}\PYG{+w}{ }\PYG{l+m}{1.0}
\end{sphinxVerbatim}


\subsection{Description}
\label{\detokenize{fix_deform:description}}
\sphinxAtStartPar
Change the volume and/or shape of the simulation box during a dynamics
run.  Orthogonal simulation boxes have 3 adjustable parameters
(x,y,z).  Triclinic (non\sphinxhyphen{}orthogonal) simulation boxes have 6
adjustable parameters (x,y,z,xy,xz,yz).  Any or all of them can be
adjusted independently and simultaneously.

\sphinxAtStartPar
The {\hyperref[\detokenize{fix_deform_pressure::doc}]{\sphinxcrossref{\DUrole{doc}{fix deform/pressure}}}} command extends
this command with additional keywords and arguments.  The rest of this
page explains the options common to both commands.  The {\hyperref[\detokenize{fix_deform_pressure::doc}]{\sphinxcrossref{\DUrole{doc}{fix
deform/pressure}}}} page explains the options
available ONLY with the fix deform/pressure command.  Note that a
simulation can define only a single deformation command: fix deform or
fix deform/pressure.

\sphinxAtStartPar
Both these fixes can be used to perform non\sphinxhyphen{}equilibrium MD (NEMD)
simulations of a continuously strained system.  See the {\hyperref[\detokenize{fix_nvt_sllod::doc}]{\sphinxcrossref{\DUrole{doc}{fix
nvt/sllod}}}} and {\hyperref[\detokenize{compute_temp_deform::doc}]{\sphinxcrossref{\DUrole{doc}{compute temp/deform}}}} commands for more details.  Note that
simulation of a continuously extended system (extensional flow) can be
modeled using the {\hyperref[\detokenize{Packages_details:pkg-uef}]{\sphinxcrossref{\DUrole{std,std-ref}{UEF package}}}} and its {\hyperref[\detokenize{fix_nh_uef::doc}]{\sphinxcrossref{\DUrole{doc}{fix
commands}}}}.

\begin{sphinxadmonition}{note}{Inconsistent trajectories due to image flags}

\sphinxAtStartPar
When running long simulations while shearing the box or using a high
shearing rate, it is possible that the image flags used for storing
unwrapped atom positions will “wrap around”.  When LAMMPS is compiled
with the default settings, case image flags are limited to a range of
\(-512 \le i \le 511\), which will overflow when atoms starting
at zero image flag value have passed through a periodic box dimension
more than 512 times.

\sphinxAtStartPar
Changing the {\hyperref[\detokenize{Build_settings:size}]{\sphinxcrossref{\DUrole{std,std-ref}{size of LAMMPS integer types}}}} to the
“bigbig” setting can make this overflow much less likely, since it
increases the image flag value range to \(- 1,048,576 \le i \le
1\,048\,575\)
\end{sphinxadmonition}


\bigskip\hrule\bigskip


\sphinxAtStartPar
For the \sphinxstyleemphasis{x}, \sphinxstyleemphasis{y}, \sphinxstyleemphasis{z} parameters, the associated dimension cannot be
shrink\sphinxhyphen{}wrapped.  For the \sphinxstyleemphasis{xy}, \sphinxstyleemphasis{yz}, \sphinxstyleemphasis{xz} parameters, the associated
second dimension cannot be shrink\sphinxhyphen{}wrapped.  Dimensions not varied by
this command can be periodic or non\sphinxhyphen{}periodic.  Dimensions
corresponding to unspecified parameters can also be controlled by a
{\hyperref[\detokenize{fix_nh::doc}]{\sphinxcrossref{\DUrole{doc}{fix npt}}}} or {\hyperref[\detokenize{fix_nh::doc}]{\sphinxcrossref{\DUrole{doc}{fix nph}}}} command.

\sphinxAtStartPar
The size and shape of the simulation box at the beginning of the
simulation run were either specified by the {\hyperref[\detokenize{create_box::doc}]{\sphinxcrossref{\DUrole{doc}{create\_box}}}} or {\hyperref[\detokenize{read_data::doc}]{\sphinxcrossref{\DUrole{doc}{read\_data}}}} or {\hyperref[\detokenize{read_restart::doc}]{\sphinxcrossref{\DUrole{doc}{read\_restart}}}} command used to setup the simulation initially if it
is the first run, or they are the values from the end of the previous
run.  The {\hyperref[\detokenize{create_box::doc}]{\sphinxcrossref{\DUrole{doc}{create\_box}}}}, {\hyperref[\detokenize{read_data::doc}]{\sphinxcrossref{\DUrole{doc}{read data}}}}, and {\hyperref[\detokenize{read_restart::doc}]{\sphinxcrossref{\DUrole{doc}{read\_restart}}}} commands specify
whether the simulation box is orthogonal or non\sphinxhyphen{}orthogonal (triclinic)
and explain the meaning of the xy,xz,yz tilt factors.  If fix deform
changes the xy,xz,yz tilt factors, then the simulation box must be
triclinic, even if its initial tilt factors are 0.0.

\sphinxAtStartPar
As described below, the desired simulation box size and shape at the
end of the run are determined by the parameters of the fix deform
command.  Every Nth timestep during the run, the simulation box is
expanded, contracted, or tilted to ramped values between the initial
and final values.


\bigskip\hrule\bigskip


\sphinxAtStartPar
For the \sphinxstyleemphasis{x}, \sphinxstyleemphasis{y}, and \sphinxstyleemphasis{z} parameters, this is the meaning of their
styles and values.

\sphinxAtStartPar
The \sphinxstyleemphasis{final}, \sphinxstyleemphasis{delta}, \sphinxstyleemphasis{scale}, \sphinxstyleemphasis{vel}, and \sphinxstyleemphasis{erate} styles all change
the specified dimension of the box via “constant displacement” which
is effectively a “constant engineering strain rate”.  This means the
box dimension changes linearly with time from its initial to final
value.

\sphinxAtStartPar
For style \sphinxstyleemphasis{final}, the final lo and hi box boundaries of a dimension
are specified.  The values can be in lattice or box distance units.
See the discussion of the units keyword below.

\sphinxAtStartPar
For style \sphinxstyleemphasis{delta}, plus or minus changes in the lo/hi box boundaries
of a dimension are specified.  The values can be in lattice or box
distance units.  See the discussion of the units keyword below.

\sphinxAtStartPar
For style \sphinxstyleemphasis{scale}, a multiplicative factor to apply to the box length
of a dimension is specified.  For example, if the initial box length
is 10, and the factor is 1.1, then the final box length will be 11.  A
factor less than 1.0 means compression.

\sphinxAtStartPar
For style \sphinxstyleemphasis{vel}, a velocity at which the box length changes is
specified in units of distance/time.  This is effectively a “constant
engineering strain rate”, where rate = V/L0 and L0 is the initial box
length.  The distance can be in lattice or box distance units.  See
the discussion of the units keyword below.  For example, if the
initial box length is 100 Angstroms, and V is 10 Angstroms/ps, then
after 10 ps, the box length will have doubled.  After 20 ps, it
will have tripled.

\sphinxAtStartPar
The \sphinxstyleemphasis{erate} style changes a dimension of the box at a “constant
engineering strain rate”.  The units of the specified strain rate are
1/time.  See the {\hyperref[\detokenize{units::doc}]{\sphinxcrossref{\DUrole{doc}{units}}}} command for the time units
associated with different choices of simulation units,
e.g. picoseconds for “metal” units).  Tensile strain is unitless and
is defined as delta/L0, where L0 is the original box length and delta
is the change relative to the original length.  The box length L as a
function of time will change as
\begin{sphinxalltt}
L(t) = L0 (1 + erate*dt)
\end{sphinxalltt}

\sphinxAtStartPar
where dt is the elapsed time (in time units).  Thus if \sphinxstyleemphasis{erate} R is
specified as 0.1 and time units are picoseconds, this means the box
length will increase by 10\% of its original length every picosecond.
I.e. strain after 1 ps = 0.1, strain after 2 ps = 0.2, etc.  R =
\sphinxhyphen{}0.01 means the box length will shrink by 1\% of its original length
every picosecond.  Note that for an “engineering” rate the change is
based on the original box length, so running with R = 1 for 10
picoseconds expands the box length by a factor of 11 (strain of 10),
which is different that what the \sphinxstyleemphasis{trate} style would induce.

\sphinxAtStartPar
The \sphinxstyleemphasis{trate} style changes a dimension of the box at a “constant true
strain rate”.  Note that this is not an “engineering strain rate”, as
the other styles are.  Rather, for a “true” rate, the rate of change
is constant, which means the box dimension changes non\sphinxhyphen{}linearly with
time from its initial to final value.  The units of the specified
strain rate are 1/time.  See the {\hyperref[\detokenize{units::doc}]{\sphinxcrossref{\DUrole{doc}{units}}}} command for the
time units associated with different choices of simulation units,
e.g. picoseconds for “metal” units).  Tensile strain is unitless and
is defined as delta/L0, where L0 is the original box length and delta
is the change relative to the original length.

\sphinxAtStartPar
The box length L as a function of time will change as
\begin{sphinxalltt}
L(t) = L0 exp(trate*dt)
\end{sphinxalltt}

\sphinxAtStartPar
where dt is the elapsed time (in time units).  Thus if \sphinxstyleemphasis{trate} R is
specified as ln(1.1) and time units are picoseconds, this means the
box length will increase by 10\% of its current (not original) length
every picosecond.  I.e. strain after 1 ps = 0.1, strain after 2 ps
= 0.21, etc.  R = ln(2) or ln(3) means the box length will double or
triple every picosecond.  R = ln(0.99) means the box length will
shrink by 1\% of its current length every picosecond.  Note that for a
“true” rate the change is continuous and based on the current length,
so running with R = ln(2) for 10 picoseconds does not expand the box
length by a factor of 11 as it would with \sphinxstyleemphasis{erate}, but by a factor of
1024 since the box length will double every picosecond.

\sphinxAtStartPar
Note that to change the volume (or cross\sphinxhyphen{}sectional area) of the
simulation box at a constant rate, you can change multiple dimensions
via \sphinxstyleemphasis{erate} or \sphinxstyleemphasis{trate}.  E.g. to double the box volume in a picosecond
picosecond, you could set “x erate M”, “y erate M”, “z erate M”, with
M = pow(2,1/3) \sphinxhyphen{} 1 = 0.26, since if each box dimension grows by 26\%,
the box volume doubles.  Or you could set “x trate M”, “y trate M”, “z
trate M”, with M = ln(1.26) = 0.231, and the box volume would double
every picosecond.

\sphinxAtStartPar
The \sphinxstyleemphasis{volume} style changes the specified dimension in such a way that
the box volume remains constant while other box dimensions are changed
explicitly via the styles discussed above.  For example, “x scale 1.1
y scale 1.1 z volume” will shrink the z box length as the x,y box
lengths increase, to keep the volume constant (product of x,y,z
lengths).  If “x scale 1.1 z volume” is specified and parameter \sphinxstyleemphasis{y} is
unspecified, then the z box length will shrink as x increases to keep
the product of x,z lengths constant.  If “x scale 1.1 y volume z
volume” is specified, then both the y,z box lengths will shrink as x
increases to keep the volume constant (product of x,y,z lengths).  In
this case, the y,z box lengths shrink so as to keep their relative
aspect ratio constant.

\sphinxAtStartPar
For solids or liquids, note that when one dimension of the box is
expanded via fix deform (i.e. tensile strain), it may be physically
undesirable to hold the other 2 box lengths constant (unspecified by
fix deform) since that implies a density change.  Using the \sphinxstyleemphasis{volume}
style for those 2 dimensions to keep the box volume constant may make
more physical sense, but may also not be correct for materials and
potentials whose Poisson ratio is not 0.5.  An alternative is to use
{\hyperref[\detokenize{fix_nh::doc}]{\sphinxcrossref{\DUrole{doc}{fix npt aniso}}}} with zero applied pressure on those 2
dimensions, so that they respond to the tensile strain dynamically.

\sphinxAtStartPar
The \sphinxstyleemphasis{wiggle} style oscillates the specified box length dimension
sinusoidally with the specified amplitude and period.  I.e. the box
length L as a function of time is given by
\begin{sphinxalltt}
L(t) = L0 + A sin(2*pi t/Tp)
\end{sphinxalltt}

\sphinxAtStartPar
where L0 is its initial length.  If the amplitude A is a positive
number the box initially expands, then contracts, etc.  If A is
negative then the box initially contracts, then expands, etc.  The
amplitude can be in lattice or box distance units.  See the discussion
of the units keyword below.

\sphinxAtStartPar
The \sphinxstyleemphasis{variable} style changes the specified box length dimension by
evaluating a variable, which presumably is a function of time.  The
variable with \sphinxstyleemphasis{name1} must be an {\hyperref[\detokenize{variable::doc}]{\sphinxcrossref{\DUrole{doc}{equal\sphinxhyphen{}style variable}}}} and should calculate a change in box length in units of
distance.  Note that this distance is in box units, not lattice units;
see the discussion of the \sphinxstyleemphasis{units} keyword below.  The formula
associated with variable \sphinxstyleemphasis{name1} can reference the current timestep.
Note that it should return the “change” in box length, not the
absolute box length.  This means it should evaluate to 0.0 when
invoked on the initial timestep of the run following the definition of
fix deform.  It should evaluate to a value \textgreater{} 0.0 to dilate the box at
future times, or a value \textless{} 0.0 to compress the box.

\sphinxAtStartPar
The variable \sphinxstyleemphasis{name2} must also be an {\hyperref[\detokenize{variable::doc}]{\sphinxcrossref{\DUrole{doc}{equal\sphinxhyphen{}style variable}}}} and should calculate the rate of box length change, in
units of distance/time, i.e. the time\sphinxhyphen{}derivative of the \sphinxstyleemphasis{name1}
variable.  This quantity is used internally by LAMMPS to reset atom
velocities when they cross periodic boundaries.  It is computed
internally for the other styles, but you must provide it when using an
arbitrary variable.

\sphinxAtStartPar
Here is an example of using the \sphinxstyleemphasis{variable} style to perform the same
box deformation as the \sphinxstyleemphasis{wiggle} style formula listed above, where we
assume that the current timestep = 0.

\begin{sphinxVerbatim}[commandchars=\\\{\}]
\PYG{k}{variable }\PYG{n+nv+nvIdentifier}{A}\PYG{+w}{ }\PYG{n}{equal}\PYG{+w}{ }\PYG{l+m}{5.0}
\PYG{k}{variable }\PYG{n+nv+nvIdentifier}{Tp}\PYG{+w}{ }\PYG{n}{equal}\PYG{+w}{ }\PYG{l+m}{10.0}
\PYG{k}{variable }\PYG{n+nv+nvIdentifier}{displace}\PYG{+w}{ }\PYG{n}{equal}\PYG{+w}{ }\PYG{l+s}{\PYGZdq{}}\PYG{l+s}{v\PYGZus{}A * sin(2*PI * step*dt/v\PYGZus{}Tp)}\PYG{l+s}{\PYGZdq{}}
\PYG{k}{variable }\PYG{n+nv+nvIdentifier}{rate}\PYG{+w}{ }\PYG{n}{equal}\PYG{+w}{ }\PYG{l+s}{\PYGZdq{}}\PYG{l+s}{2*PI*v\PYGZus{}A/v\PYGZus{}Tp * cos(2*PI * step*dt/v\PYGZus{}Tp)}\PYG{l+s}{\PYGZdq{}}
\PYG{k}{fix }\PYG{n+nv+nvIdentifier}{2}\PYG{+w}{ }\PYG{n+nv+nvIdentifier}{all}\PYG{+w}{ }\PYG{n}{deform}\PYG{+w}{ }\PYG{l+m}{1}\PYG{+w}{ }\PYG{n}{x}\PYG{+w}{ }\PYG{k}{variable }\PYG{n+nv+nvIdentifier}{v\PYGZus{}displace}\PYG{+w}{ }\PYG{n}{v\PYGZus{}rate}\PYG{+w}{ }\PYG{n}{remap}\PYG{+w}{ }\PYG{n}{v}
\end{sphinxVerbatim}

\sphinxAtStartPar
For the \sphinxstyleemphasis{scale}, \sphinxstyleemphasis{vel}, \sphinxstyleemphasis{erate}, \sphinxstyleemphasis{trate}, \sphinxstyleemphasis{volume}, \sphinxstyleemphasis{wiggle}, and
\sphinxstyleemphasis{variable} styles, the box length is expanded or compressed around its
mid point.


\bigskip\hrule\bigskip


\sphinxAtStartPar
For the \sphinxstyleemphasis{xy}, \sphinxstyleemphasis{xz}, and \sphinxstyleemphasis{yz} parameters, this is the meaning of their
styles and values.  Note that changing the tilt factors of a triclinic
box does not change its volume.

\sphinxAtStartPar
The \sphinxstyleemphasis{final}, \sphinxstyleemphasis{delta}, \sphinxstyleemphasis{vel}, and \sphinxstyleemphasis{erate} styles all change the shear
strain at a “constant engineering shear strain rate”.  This means the
tilt factor changes linearly with time from its initial to final
value.

\sphinxAtStartPar
For style \sphinxstyleemphasis{final}, the final tilt factor is specified.  The value
can be in lattice or box distance units.  See the discussion of the
units keyword below.

\sphinxAtStartPar
For style \sphinxstyleemphasis{delta}, a plus or minus change in the tilt factor is
specified.  The value can be in lattice or box distance units.  See
the discussion of the units keyword below.

\sphinxAtStartPar
For style \sphinxstyleemphasis{vel}, a velocity at which the tilt factor changes is
specified in units of distance/time.  This is effectively an
“engineering shear strain rate”, where rate = V/L0 and L0 is the
initial box length perpendicular to the direction of shear.  The
distance can be in lattice or box distance units.  See the discussion
of the units keyword below.  For example, if the initial tilt factor
is 5 Angstroms, and the V is 10 Angstroms/ps, then after 1 ps, the
tilt factor will be 15 Angstroms.  After 2 ps, it will be 25
Angstroms.

\sphinxAtStartPar
The \sphinxstyleemphasis{erate} style changes a tilt factor at a “constant engineering
shear strain rate”.  The units of the specified shear strain rate are
1/time.  See the {\hyperref[\detokenize{units::doc}]{\sphinxcrossref{\DUrole{doc}{units}}}} command for the time units
associated with different choices of simulation units,
e.g. picoseconds for “metal” units).  Shear strain is unitless and is
defined as offset/length, where length is the box length perpendicular
to the shear direction (e.g. y box length for xy deformation) and
offset is the displacement distance in the shear direction (e.g. x
direction for xy deformation) from the unstrained orientation.

\sphinxAtStartPar
The tilt factor T as a function of time will change as
\begin{sphinxalltt}
T(t) = T0 + L0*erate*dt
\end{sphinxalltt}

\sphinxAtStartPar
where T0 is the initial tilt factor, L0 is the original length of the
box perpendicular to the shear direction (e.g. y box length for xy
deformation), and dt is the elapsed time (in time units).  Thus if
\sphinxstyleemphasis{erate} R is specified as 0.1 and time units are picoseconds, this
means the shear strain will increase by 0.1 every picosecond.  I.e. if
the xy shear strain was initially 0.0, then strain after 1 ps = 0.1,
strain after 2 ps = 0.2, etc.  Thus the tilt factor would be 0.0 at
time 0, 0.1*ybox at 1 ps, 0.2*ybox at 2 ps, etc, where ybox is the
original y box length.  R = 1 or 2 means the tilt factor will increase
by 1 or 2 every picosecond.  R = \sphinxhyphen{}0.01 means a decrease in shear
strain by 0.01 every picosecond.

\sphinxAtStartPar
The \sphinxstyleemphasis{trate} style changes a tilt factor at a “constant true shear
strain rate”.  Note that this is not an “engineering shear strain
rate”, as the other styles are.  Rather, for a “true” rate, the rate
of change is constant, which means the tilt factor changes
non\sphinxhyphen{}linearly with time from its initial to final value.  The units of
the specified shear strain rate are 1/time.  See the
{\hyperref[\detokenize{units::doc}]{\sphinxcrossref{\DUrole{doc}{units}}}} command for the time units associated with
different choices of simulation units, e.g. picoseconds for “metal”
units).  Shear strain is unitless and is defined as offset/length,
where length is the box length perpendicular to the shear direction
(e.g. y box length for xy deformation) and offset is the displacement
distance in the shear direction (e.g. x direction for xy deformation)
from the unstrained orientation.

\sphinxAtStartPar
The tilt factor T as a function of time will change as
\begin{sphinxalltt}
T(t) = T0 exp(trate*dt)
\end{sphinxalltt}

\sphinxAtStartPar
where T0 is the initial tilt factor and dt is the elapsed time (in
time units).  Thus if \sphinxstyleemphasis{trate} R is specified as ln(1.1) and time units
are picoseconds, this means the shear strain or tilt factor will
increase by 10\% every picosecond.  I.e. if the xy shear strain was
initially 0.1, then strain after 1 ps = 0.11, strain after 2 ps =
0.121, etc.  R = ln(2) or ln(3) means the tilt factor will double or
triple every picosecond.  R = ln(0.99) means the tilt factor will
shrink by 1\% every picosecond.  Note that the change is continuous, so
running with R = ln(2) for 10 picoseconds does not change the tilt
factor by a factor of 10, but by a factor of 1024 since it doubles
every picosecond.  Note that the initial tilt factor must be non\sphinxhyphen{}zero
to use the \sphinxstyleemphasis{trate} option.

\sphinxAtStartPar
Note that shear strain is defined as the tilt factor divided by the
perpendicular box length.  The \sphinxstyleemphasis{erate} and \sphinxstyleemphasis{trate} styles control the
tilt factor, but assume the perpendicular box length remains constant.
If this is not the case (e.g. it changes due to another fix deform
parameter), then this effect on the shear strain is ignored.

\sphinxAtStartPar
The \sphinxstyleemphasis{wiggle} style oscillates the specified tilt factor sinusoidally
with the specified amplitude and period.  I.e. the tilt factor T as a
function of time is given by
\begin{sphinxalltt}
T(t) = T0 + A sin(2*pi t/Tp)
\end{sphinxalltt}

\sphinxAtStartPar
where T0 is its initial value.  If the amplitude A is a positive
number the tilt factor initially becomes more positive, then more
negative, etc.  If A is negative then the tilt factor initially
becomes more negative, then more positive, etc.  The amplitude can be
in lattice or box distance units.  See the discussion of the units
keyword below.

\sphinxAtStartPar
The \sphinxstyleemphasis{variable} style changes the specified tilt factor by evaluating a
variable, which presumably is a function of time.  The variable with
\sphinxstyleemphasis{name1} must be an {\hyperref[\detokenize{variable::doc}]{\sphinxcrossref{\DUrole{doc}{equal\sphinxhyphen{}style variable}}}} and should
calculate a change in tilt in units of distance.  Note that this
distance is in box units, not lattice units; see the discussion of the
\sphinxstyleemphasis{units} keyword below.  The formula associated with variable \sphinxstyleemphasis{name1}
can reference the current timestep.  Note that it should return the
“change” in tilt factor, not the absolute tilt factor.  This means it
should evaluate to 0.0 when invoked on the initial timestep of the run
following the definition of fix deform.

\sphinxAtStartPar
The variable \sphinxstyleemphasis{name2} must also be an {\hyperref[\detokenize{variable::doc}]{\sphinxcrossref{\DUrole{doc}{equal\sphinxhyphen{}style variable}}}} and should calculate the rate of tilt change, in units of
distance/time, i.e. the time\sphinxhyphen{}derivative of the \sphinxstyleemphasis{name1} variable.  This
quantity is used internally by LAMMPS to reset atom velocities when
they cross periodic boundaries.  It is computed internally for the
other styles, but you must provide it when using an arbitrary
variable.

\sphinxAtStartPar
Here is an example of using the \sphinxstyleemphasis{variable} style to perform the same
box deformation as the \sphinxstyleemphasis{wiggle} style formula listed above, where we
assume that the current timestep = 0.

\begin{sphinxVerbatim}[commandchars=\\\{\}]
\PYG{k}{variable }\PYG{n+nv+nvIdentifier}{A}\PYG{+w}{ }\PYG{n}{equal}\PYG{+w}{ }\PYG{l+m}{5.0}
\PYG{k}{variable }\PYG{n+nv+nvIdentifier}{Tp}\PYG{+w}{ }\PYG{n}{equal}\PYG{+w}{ }\PYG{l+m}{10.0}
\PYG{k}{variable }\PYG{n+nv+nvIdentifier}{displace}\PYG{+w}{ }\PYG{n}{equal}\PYG{+w}{ }\PYG{l+s}{\PYGZdq{}}\PYG{l+s}{v\PYGZus{}A * sin(2*PI * step*dt/v\PYGZus{}Tp)}\PYG{l+s}{\PYGZdq{}}
\PYG{k}{variable }\PYG{n+nv+nvIdentifier}{rate}\PYG{+w}{ }\PYG{n}{equal}\PYG{+w}{ }\PYG{l+s}{\PYGZdq{}}\PYG{l+s}{2*PI*v\PYGZus{}A/v\PYGZus{}Tp * cos(2*PI * step*dt/v\PYGZus{}Tp)}\PYG{l+s}{\PYGZdq{}}
\PYG{k}{fix }\PYG{n+nv+nvIdentifier}{2}\PYG{+w}{ }\PYG{n+nv+nvIdentifier}{all}\PYG{+w}{ }\PYG{n}{deform}\PYG{+w}{ }\PYG{l+m}{1}\PYG{+w}{ }\PYG{n}{xy}\PYG{+w}{ }\PYG{k}{variable }\PYG{n+nv+nvIdentifier}{v\PYGZus{}displace}\PYG{+w}{ }\PYG{n}{v\PYGZus{}rate}\PYG{+w}{ }\PYG{n}{remap}\PYG{+w}{ }\PYG{n}{v}
\end{sphinxVerbatim}


\bigskip\hrule\bigskip


\sphinxAtStartPar
All of the tilt styles change the xy, xz, yz tilt factors during a
simulation.  In LAMMPS, tilt factors (xy,xz,yz) for triclinic boxes
are normally bounded by half the distance of the parallel box length.
See the discussion of the \sphinxstyleemphasis{flip} keyword below, to allow this bound to
be exceeded, if desired.

\sphinxAtStartPar
For example, if xlo = 2 and xhi = 12, then the x box length is 10 and
the xy tilt factor must be between \sphinxhyphen{}5 and 5.  Similarly, both xz and
yz must be between \sphinxhyphen{}(xhi\sphinxhyphen{}xlo)/2 and +(yhi\sphinxhyphen{}ylo)/2.  Note that this is
not a limitation, since if the maximum tilt factor is 5 (as in this
example), then configurations with tilt = …, \sphinxhyphen{}15, \sphinxhyphen{}5, 5, 15, 25,
… are all equivalent.

\sphinxAtStartPar
To obey this constraint and allow for large shear deformations to be
applied via the \sphinxstyleemphasis{xy}, \sphinxstyleemphasis{xz}, or \sphinxstyleemphasis{yz} parameters, the following
algorithm is used.  If \sphinxstyleemphasis{prd} is the associated parallel box length (10
in the example above), then if the tilt factor exceeds the accepted
range of \sphinxhyphen{}5 to 5 during the simulation, then the box is flipped to the
other limit (an equivalent box) and the simulation continues.  Thus
for this example, if the initial xy tilt factor was 0.0 and “xy final
100.0” was specified, then during the simulation the xy tilt factor
would increase from 0.0 to 5.0, the box would be flipped so that the
tilt factor becomes \sphinxhyphen{}5.0, the tilt factor would increase from \sphinxhyphen{}5.0 to
5.0, the box would be flipped again, etc.  The flip occurs 10 times
and the final tilt factor at the end of the simulation would be 0.0.
During each flip event, atoms are remapped into the new box in the
appropriate manner.

\sphinxAtStartPar
The one exception to this rule is if the first dimension in the tilt
factor (x for xy) is non\sphinxhyphen{}periodic.  In that case, the limits on the
tilt factor are not enforced, since flipping the box in that dimension
does not change the atom positions due to non\sphinxhyphen{}periodicity.  In this
mode, if you tilt the system to extreme angles, the simulation will
simply become inefficient due to the highly skewed simulation box.


\bigskip\hrule\bigskip


\sphinxAtStartPar
Each time the box size or shape is changed, the \sphinxstyleemphasis{remap} keyword
determines whether atom positions are remapped to the new box.  If
\sphinxstyleemphasis{remap} is set to \sphinxstyleemphasis{x} (the default), atoms in the fix group are
remapped; otherwise they are not.  Note that their velocities are not
changed, just their positions are altered.  If \sphinxstyleemphasis{remap} is set to \sphinxstyleemphasis{v},
then any atom in the fix group that crosses a periodic boundary will
have a delta added to its velocity equal to the difference in
velocities between the lo and hi boundaries.  Note that this velocity
difference can include tilt components, e.g. a delta in the x velocity
when an atom crosses the y periodic boundary.  If \sphinxstyleemphasis{remap} is set to
\sphinxstyleemphasis{none}, then neither of these remappings take place.

\sphinxAtStartPar
Conceptually, setting \sphinxstyleemphasis{remap} to \sphinxstyleemphasis{x} forces the atoms to deform via an
affine transformation that exactly matches the box deformation.  This
setting is typically appropriate for solids.  Note that though the
atoms are effectively “moving” with the box over time, it is not due
to their having a velocity that tracks the box change, but only due to
the remapping.  By contrast, setting \sphinxstyleemphasis{remap} to \sphinxstyleemphasis{v} is typically
appropriate for fluids, where you want the atoms to respond to the
change in box size/shape on their own and acquire a velocity that
matches the box change, so that their motion will naturally track the
box without explicit remapping of their coordinates.

\begin{sphinxadmonition}{note}{Note:}
\sphinxAtStartPar
When non\sphinxhyphen{}equilibrium MD (NEMD) simulations are performed using
this fix, the option “remap v” should normally be used.  This is
because {\hyperref[\detokenize{fix_nvt_sllod::doc}]{\sphinxcrossref{\DUrole{doc}{fix nvt/sllod}}}} adjusts the atom positions
and velocities to induce a velocity profile that matches the changing
box size/shape.  Thus atom coordinates should NOT be remapped by fix
deform, but velocities SHOULD be when atoms cross periodic boundaries,
since that is consistent with maintaining the velocity profile already
created by fix nvt/sllod.  LAMMPS will warn you if the \sphinxstyleemphasis{remap} setting
is not consistent with fix nvt/sllod.
\end{sphinxadmonition}

\begin{sphinxadmonition}{note}{Note:}
\sphinxAtStartPar
For non\sphinxhyphen{}equilibrium MD (NEMD) simulations using “remap v” it is
usually desirable that the fluid (or flowing material,
e.g. granular particles) stream with a velocity profile consistent
with the deforming box.  As mentioned above, using a thermostat
such as {\hyperref[\detokenize{fix_nvt_sllod::doc}]{\sphinxcrossref{\DUrole{doc}{fix nvt/sllod}}}} or {\hyperref[\detokenize{fix_langevin::doc}]{\sphinxcrossref{\DUrole{doc}{fix lavgevin}}}} (with a bias provided by {\hyperref[\detokenize{compute_temp_deform::doc}]{\sphinxcrossref{\DUrole{doc}{compute temp/deform}}}}), will typically accomplish that.  If you do
not use a thermostat, then there is no driving force pushing the
atoms to flow in a manner consistent with the deforming box.
E.g. for a shearing system the box deformation velocity may vary
from 0 at the bottom to 10 at the top of the box.  But the stream
velocity profile of the atoms may vary from \sphinxhyphen{}5 at the bottom to +5
at the top.  You can monitor these effects using the {\hyperref[\detokenize{fix_ave_chunk::doc}]{\sphinxcrossref{\DUrole{doc}{fix
ave/chunk}}}}, {\hyperref[\detokenize{compute_temp_deform::doc}]{\sphinxcrossref{\DUrole{doc}{compute temp/deform}}}}, and {\hyperref[\detokenize{compute_temp_profile::doc}]{\sphinxcrossref{\DUrole{doc}{compute temp/profile}}}} commands.  One way to induce atoms to
stream consistent with the box deformation is to give them an
initial velocity profile, via the {\hyperref[\detokenize{velocity::doc}]{\sphinxcrossref{\DUrole{doc}{velocity ramp}}}}
command, that matches the box deformation rate.  This also
typically helps the system come to equilibrium more quickly, even
if a thermostat is used.
\end{sphinxadmonition}

\begin{sphinxadmonition}{note}{Note:}
\sphinxAtStartPar
If a {\hyperref[\detokenize{fix_rigid::doc}]{\sphinxcrossref{\DUrole{doc}{fix rigid}}}} is defined for rigid bodies, and
\sphinxstyleemphasis{remap} is set to \sphinxstyleemphasis{x}, then the center\sphinxhyphen{}of\sphinxhyphen{}mass coordinates of rigid
bodies will be remapped to the changing simulation box.  This will
be done regardless of whether atoms in the rigid bodies are in the
fix deform group or not.  The velocity of the centers of mass are
not remapped even if \sphinxstyleemphasis{remap} is set to \sphinxstyleemphasis{v}, since {\hyperref[\detokenize{fix_nvt_sllod::doc}]{\sphinxcrossref{\DUrole{doc}{fix
nvt/sllod}}}} does not currently do anything special
for rigid particles.  If you wish to perform a NEMD simulation of
rigid particles, you can either thermostat them independently or
include a background fluid and thermostat the fluid via {\hyperref[\detokenize{fix_nvt_sllod::doc}]{\sphinxcrossref{\DUrole{doc}{fix
nvt/sllod}}}}.
\end{sphinxadmonition}

\sphinxAtStartPar
The \sphinxstyleemphasis{flip} keyword allows the tilt factors for a triclinic box to
exceed half the distance of the parallel box length, as discussed
above.  If the \sphinxstyleemphasis{flip} value is set to \sphinxstyleemphasis{yes}, the bound is enforced by
flipping the box when it is exceeded.  If the \sphinxstyleemphasis{flip} value is set to
\sphinxstyleemphasis{no}, the tilt will continue to change without flipping.  Note that if
you apply large deformations, this means the box shape can tilt
dramatically LAMMPS will run less efficiently, due to the large volume
of communication needed to acquire ghost atoms around a processor’s
irregular\sphinxhyphen{}shaped subdomain.  For extreme values of tilt, LAMMPS may
also lose atoms and generate an error.

\sphinxAtStartPar
The \sphinxstyleemphasis{units} keyword determines the meaning of the distance units used
to define various arguments.  A \sphinxstyleemphasis{box} value selects standard distance
units as defined by the {\hyperref[\detokenize{units::doc}]{\sphinxcrossref{\DUrole{doc}{units}}}} command, e.g. Angstroms for
units = real or metal.  A \sphinxstyleemphasis{lattice} value means the distance units are
in lattice spacings.  The {\hyperref[\detokenize{lattice::doc}]{\sphinxcrossref{\DUrole{doc}{lattice}}}} command must have
been previously used to define the lattice spacing.  Note that the
units choice also affects the \sphinxstyleemphasis{vel} style parameters since it is
defined in terms of distance/time.  Also note that the units keyword
does not affect the \sphinxstyleemphasis{variable} style.  You should use the \sphinxstyleemphasis{xlat},
\sphinxstyleemphasis{ylat}, \sphinxstyleemphasis{zlat} keywords of the {\hyperref[\detokenize{thermo_style::doc}]{\sphinxcrossref{\DUrole{doc}{thermo\_style}}}}
command if you want to include lattice spacings in a variable formula.


\bigskip\hrule\bigskip


\sphinxAtStartPar
Styles with a \sphinxstyleemphasis{gpu}, \sphinxstyleemphasis{intel}, \sphinxstyleemphasis{kk}, \sphinxstyleemphasis{omp}, or \sphinxstyleemphasis{opt} suffix are
functionally the same as the corresponding style without the suffix.
They have been optimized to run faster, depending on your available
hardware, as discussed on the {\hyperref[\detokenize{Speed_packages::doc}]{\sphinxcrossref{\DUrole{doc}{Accelerator packages}}}}
page.  The accelerated styles take the same arguments and should
produce the same results, except for round\sphinxhyphen{}off and precision issues.

\sphinxAtStartPar
These accelerated styles are part of the GPU, INTEL, KOKKOS,
OPENMP, and OPT packages, respectively.  They are only enabled if
LAMMPS was built with those packages.  See the {\hyperref[\detokenize{Build_package::doc}]{\sphinxcrossref{\DUrole{doc}{Build package}}}} page for more info.

\sphinxAtStartPar
You can specify the accelerated styles explicitly in your input script
by including their suffix, or you can use the {\hyperref[\detokenize{Run_options::doc}]{\sphinxcrossref{\DUrole{doc}{\sphinxhyphen{}suffix command\sphinxhyphen{}line switch}}}} when you invoke LAMMPS, or you can use the
{\hyperref[\detokenize{suffix::doc}]{\sphinxcrossref{\DUrole{doc}{suffix}}}} command in your input script.

\sphinxAtStartPar
See the {\hyperref[\detokenize{Speed_packages::doc}]{\sphinxcrossref{\DUrole{doc}{Accelerator packages}}}} page for more
instructions on how to use the accelerated styles effectively.


\subsection{Restart, fix\_modify, output, run start/stop, minimize info}
\label{\detokenize{fix_deform:restart-fix-modify-output-run-start-stop-minimize-info}}
\sphinxAtStartPar
This fix will restore the initial box settings from {\hyperref[\detokenize{restart::doc}]{\sphinxcrossref{\DUrole{doc}{binary
restart files}}}}, which allows the fix to be properly continue
deformation, when using the start/stop options of the {\hyperref[\detokenize{run::doc}]{\sphinxcrossref{\DUrole{doc}{run}}}}
command.  None of the {\hyperref[\detokenize{fix_modify::doc}]{\sphinxcrossref{\DUrole{doc}{fix\_modify}}}} options are
relevant to this fix.  No global or per\sphinxhyphen{}atom quantities are stored by
this fix for access by various {\hyperref[\detokenize{Howto_output::doc}]{\sphinxcrossref{\DUrole{doc}{output commands}}}}.

\sphinxAtStartPar
This fix can perform deformation over multiple runs, using the \sphinxstyleemphasis{start}
and \sphinxstyleemphasis{stop} keywords of the {\hyperref[\detokenize{run::doc}]{\sphinxcrossref{\DUrole{doc}{run}}}} command.  See the
{\hyperref[\detokenize{run::doc}]{\sphinxcrossref{\DUrole{doc}{run}}}} command for details of how to do this.

\sphinxAtStartPar
This fix is not invoked during {\hyperref[\detokenize{minimize::doc}]{\sphinxcrossref{\DUrole{doc}{energy minimization}}}}.


\subsection{Restrictions}
\label{\detokenize{fix_deform:restrictions}}
\sphinxAtStartPar
You cannot apply x, y, or z deformations to a dimension that is
shrink\sphinxhyphen{}wrapped via the {\hyperref[\detokenize{boundary::doc}]{\sphinxcrossref{\DUrole{doc}{boundary}}}} command.

\sphinxAtStartPar
You cannot apply xy, yz, or xz deformations to a second dimension (y
in xy) that is shrink\sphinxhyphen{}wrapped via the {\hyperref[\detokenize{boundary::doc}]{\sphinxcrossref{\DUrole{doc}{boundary}}}}
command.


\subsection{Related commands}
\label{\detokenize{fix_deform:related-commands}}
\sphinxAtStartPar
{\hyperref[\detokenize{fix_deform_pressure::doc}]{\sphinxcrossref{\DUrole{doc}{fix deform/pressure}}}},
{\hyperref[\detokenize{change_box::doc}]{\sphinxcrossref{\DUrole{doc}{change\_box}}}}


\subsection{Default}
\label{\detokenize{fix_deform:default}}
\sphinxAtStartPar
The option defaults are remap = x, flip = yes, and units = lattice.

\sphinxstepscope

\index{fix deform/pressure@\spxentry{fix deform/pressure}}\ignorespaces 

\section{fix deform/pressure command}
\label{\detokenize{fix_deform_pressure:fix-deform-pressure-command}}\label{\detokenize{fix_deform_pressure:index-0}}\label{\detokenize{fix_deform_pressure::doc}}

\subsection{Syntax}
\label{\detokenize{fix_deform_pressure:syntax}}
\begin{sphinxVerbatim}[commandchars=\\\{\}]
fix ID group\PYGZhy{}ID deform/pressure N parameter style args ... keyword value ...
\end{sphinxVerbatim}
\begin{itemize}
\item {} 
\sphinxAtStartPar
ID, group\sphinxhyphen{}ID are documented in {\hyperref[\detokenize{fix::doc}]{\sphinxcrossref{\DUrole{doc}{fix}}}} command

\item {} 
\sphinxAtStartPar
deform/pressure = style name of this fix command

\item {} 
\sphinxAtStartPar
N = perform box deformation every this many timesteps

\item {} 
\sphinxAtStartPar
one or more parameter/args sequences may be appended
\begin{sphinxalltt}
parameter = \sphinxstyleemphasis{x} or \sphinxstyleemphasis{y} or \sphinxstyleemphasis{z} or \sphinxstyleemphasis{xy} or \sphinxstyleemphasis{xz} or \sphinxstyleemphasis{yz} or \sphinxstyleemphasis{box}
  \sphinxstyleemphasis{x}, \sphinxstyleemphasis{y}, \sphinxstyleemphasis{z} args = style value(s)
    style = \sphinxstyleemphasis{final} or \sphinxstyleemphasis{delta} or \sphinxstyleemphasis{scale} or \sphinxstyleemphasis{vel} or \sphinxstyleemphasis{erate} or \sphinxstyleemphasis{trate} or \sphinxstyleemphasis{volume} or \sphinxstyleemphasis{wiggle} or \sphinxstyleemphasis{variable} or \sphinxstyleemphasis{pressure} or \sphinxstyleemphasis{pressure/mean}
      \sphinxstyleemphasis{final} values = lo hi
        lo hi = box boundaries at end of run (distance units)
      \sphinxstyleemphasis{delta} values = dlo dhi
        dlo dhi = change in box boundaries at end of run (distance units)
      \sphinxstyleemphasis{scale} values = factor
        factor = multiplicative factor for change in box length at end of run
      \sphinxstyleemphasis{vel} value = V
        V = change box length at this velocity (distance/time units),
            effectively an engineering strain rate
      \sphinxstyleemphasis{erate} value = R
        R = engineering strain rate (1/time units)
      \sphinxstyleemphasis{trate} value = R
        R = true strain rate (1/time units)
      \sphinxstyleemphasis{volume} value = none = adjust this dim to preserve volume of system
      \sphinxstyleemphasis{wiggle} values = A Tp
        A = amplitude of oscillation (distance units)
        Tp = period of oscillation (time units)
      \sphinxstyleemphasis{variable} values = v\_name1 v\_name2
        v\_name1 = variable with name1 for box length change as function of time
        v\_name2 = variable with name2 for change rate as function of time
      \sphinxstyleemphasis{pressure} values = target gain
        target = target pressure (pressure units)
        gain = proportional gain constant (1/(time * pressure) or 1/time units)
      \sphinxstyleemphasis{pressure/mean} values = target gain
        target = target pressure (pressure units)
        gain = proportional gain constant (1/(time * pressure) or 1/time units)

  \sphinxstyleemphasis{xy}, \sphinxstyleemphasis{xz}, \sphinxstyleemphasis{yz} args = style value
    style = \sphinxstyleemphasis{final} or \sphinxstyleemphasis{delta} or \sphinxstyleemphasis{vel} or \sphinxstyleemphasis{erate} or \sphinxstyleemphasis{trate} or \sphinxstyleemphasis{wiggle} or \sphinxstyleemphasis{variable} or \sphinxstyleemphasis{pressure} or \sphinxstyleemphasis{erate/rescale}
      \sphinxstyleemphasis{final} value = tilt
        tilt = tilt factor at end of run (distance units)
      \sphinxstyleemphasis{delta} value = dtilt
        dtilt = change in tilt factor at end of run (distance units)
      \sphinxstyleemphasis{vel} value = V
        V = change tilt factor at this velocity (distance/time units),
            effectively an engineering shear strain rate
      \sphinxstyleemphasis{erate} value = R
        R = engineering shear strain rate (1/time units)
      \sphinxstyleemphasis{erate/rescale} value = R
        R = engineering shear strain rate (1/time units)
      \sphinxstyleemphasis{trate} value = R
        R = true shear strain rate (1/time units)
      \sphinxstyleemphasis{wiggle} values = A Tp
        A = amplitude of oscillation (distance units)
        Tp = period of oscillation (time units)
      \sphinxstyleemphasis{variable} values = v\_name1 v\_name2
        v\_name1 = variable with name1 for tilt change as function of time
        v\_name2 = variable with name2 for change rate as function of time
      \sphinxstyleemphasis{pressure} values = target gain
        target = target pressure (pressure units)
        gain = proportional gain constant (1/(time * pressure) or 1/time units)
      \sphinxstyleemphasis{erate/rescale} value = R
        R = engineering shear strain rate (1/time units)

  \sphinxstyleemphasis{box} = style value
    style = \sphinxstyleemphasis{volume} or \sphinxstyleemphasis{pressure}
      \sphinxstyleemphasis{volume} value = none = isotropically adjust system to preserve volume of system
      \sphinxstyleemphasis{pressure} values = target gain
        target = target mean pressure (pressure units)
        gain = proportional gain constant (1/(time * pressure) or 1/time units)
\end{sphinxalltt}

\item {} 
\sphinxAtStartPar
zero or more keyword/value pairs may be appended

\item {} 
\sphinxAtStartPar
keyword = \sphinxstyleemphasis{remap} or \sphinxstyleemphasis{flip} or \sphinxstyleemphasis{units} or \sphinxstyleemphasis{couple} or \sphinxstyleemphasis{vol/balance/p} or \sphinxstyleemphasis{max/rate} or \sphinxstyleemphasis{normalize/pressure}
\begin{sphinxalltt}
\sphinxstyleemphasis{remap} value = \sphinxstyleemphasis{x} or \sphinxstyleemphasis{v} or \sphinxstyleemphasis{none}
  x = remap coords of atoms in group into deforming box
  v = remap velocities of atoms in group when they cross periodic boundaries
  none = no remapping of x or v
\sphinxstyleemphasis{flip} value = \sphinxstyleemphasis{yes} or \sphinxstyleemphasis{no}
  allow or disallow box flips when it becomes highly skewed
\sphinxstyleemphasis{units} value = \sphinxstyleemphasis{lattice} or \sphinxstyleemphasis{box}
  lattice = distances are defined in lattice units
  box = distances are defined in simulation box units
\sphinxstyleemphasis{couple} value = \sphinxstyleemphasis{none} or \sphinxstyleemphasis{xyz} or \sphinxstyleemphasis{xy} or \sphinxstyleemphasis{yz} or \sphinxstyleemphasis{xz}
  couple pressure values of various dimensions
\sphinxstyleemphasis{vol/balance/p} value = \sphinxstyleemphasis{yes} or \sphinxstyleemphasis{no}
  Modifies the behavior of the \sphinxstyleemphasis{volume} option to try and balance pressures
\sphinxstyleemphasis{max/rate} value = \sphinxstyleemphasis{rate}
  rate = maximum strain rate for pressure control
\sphinxstyleemphasis{normalize/pressure} value = \sphinxstyleemphasis{yes} or \sphinxstyleemphasis{no}
  Modifies pressure controls such that the deviation in pressure is normalized by the target pressure
\end{sphinxalltt}

\end{itemize}


\subsection{Examples}
\label{\detokenize{fix_deform_pressure:examples}}
\begin{sphinxVerbatim}[commandchars=\\\{\}]
\PYG{k}{fix }\PYG{n+nv+nvIdentifier}{1}\PYG{+w}{ }\PYG{n+nv+nvIdentifier}{all}\PYG{+w}{ }\PYG{n}{deform}\PYG{o}{/}\PYG{n}{pressure}\PYG{+w}{ }\PYG{l+m}{1}\PYG{+w}{ }\PYG{n}{x}\PYG{+w}{ }\PYG{n}{pressure}\PYG{+w}{ }\PYG{l+m}{2.0}\PYG{+w}{ }\PYG{l+m}{0.1}\PYG{+w}{ }\PYG{n}{normalize}\PYG{o}{/}\PYG{n}{pressure}\PYG{+w}{ }\PYG{n}{yes}\PYG{+w}{ }\PYG{n}{max}\PYG{o}{/}\PYG{n}{rate}\PYG{+w}{ }\PYG{l+m}{0.001}
\PYG{k}{fix }\PYG{n+nv+nvIdentifier}{1}\PYG{+w}{ }\PYG{n+nv+nvIdentifier}{all}\PYG{+w}{ }\PYG{n}{deform}\PYG{o}{/}\PYG{n}{pressure}\PYG{+w}{ }\PYG{l+m}{1}\PYG{+w}{ }\PYG{n}{x}\PYG{+w}{ }\PYG{n}{trate}\PYG{+w}{ }\PYG{l+m}{0.1}\PYG{+w}{ }\PYG{n}{y}\PYG{+w}{ }\PYG{n}{volume}\PYG{+w}{ }\PYG{n}{z}\PYG{+w}{ }\PYG{n}{volume}\PYG{+w}{ }\PYG{n}{vol}\PYG{o}{/}\PYG{n}{balance}\PYG{o}{/}\PYG{n}{p}\PYG{+w}{ }\PYG{n}{yes}
\PYG{k}{fix }\PYG{n+nv+nvIdentifier}{1}\PYG{+w}{ }\PYG{n+nv+nvIdentifier}{all}\PYG{+w}{ }\PYG{n}{deform}\PYG{o}{/}\PYG{n}{pressure}\PYG{+w}{ }\PYG{l+m}{1}\PYG{+w}{ }\PYG{n}{x}\PYG{+w}{ }\PYG{n}{trate}\PYG{+w}{ }\PYG{l+m}{0.1}\PYG{+w}{ }\PYG{n}{y}\PYG{+w}{ }\PYG{n}{pressure}\PYG{o}{/}\PYG{n}{mean}\PYG{+w}{ }\PYG{l+m}{0.0}\PYG{+w}{ }\PYG{l+m}{1.0}\PYG{+w}{ }\PYG{n}{z}\PYG{+w}{ }\PYG{n}{pressure}\PYG{o}{/}\PYG{n}{mean}\PYG{+w}{ }\PYG{l+m}{0.0}\PYG{+w}{ }\PYG{l+m}{1.0}
\end{sphinxVerbatim}


\subsection{Description}
\label{\detokenize{fix_deform_pressure:description}}
\sphinxAtStartPar
\DUrole{versionmodified,added}{Added in version 17Apr2024.}

\sphinxAtStartPar
This fix is an extension of the {\hyperref[\detokenize{fix_deform::doc}]{\sphinxcrossref{\DUrole{doc}{fix deform}}}}
command, which allows all of its options to be used as well as new
pressure\sphinxhyphen{}based controls implemented by this command.

\sphinxAtStartPar
All arguments described on the {\hyperref[\detokenize{fix_deform::doc}]{\sphinxcrossref{\DUrole{doc}{fix deform}}}} doc page
also apply to this fix unless otherwise noted below.  The rest of this
page explains the arguments specific to this fix only.  Note that a
simulation can define only a single deformation command: fix deform or
fix deform/pressure.

\begin{sphinxadmonition}{note}{Inconsistent trajectories due to image flags}

\sphinxAtStartPar
When running long simulations while shearing the box or using a high
shearing rate, it is possible that the image flags used for storing
unwrapped atom positions will “wrap around”.  When LAMMPS is compiled
with the default settings, case image flags are limited to a range of
\(-512 \le i \le 511\), which will overflow when atoms starting
at zero image flag value have passed through a periodic box dimension
more than 512 times.

\sphinxAtStartPar
Changing the {\hyperref[\detokenize{Build_settings:size}]{\sphinxcrossref{\DUrole{std,std-ref}{size of LAMMPS integer types}}}} to the
“bigbig” setting can make this overflow much less likely, since it
increases the image flag value range to \(- 1,048,576 \le i \le
1\,048\,575\)
\end{sphinxadmonition}


\bigskip\hrule\bigskip


\sphinxAtStartPar
For the \sphinxstyleemphasis{x}, \sphinxstyleemphasis{y}, and \sphinxstyleemphasis{z} parameters, this is the meaning of the
styles and values provided by this fix.

\sphinxAtStartPar
The \sphinxstyleemphasis{pressure} style adjusts a dimension’s box length to control the
corresponding component of the pressure tensor. This option attempts to
maintain a specified target pressure using a linear controller where the
box length \(L\) evolves according to the equation
\begin{equation*}
\begin{split}\frac{d L(t)}{dt} = L(t) k (P_t - P)\end{split}
\end{equation*}
\sphinxAtStartPar
where \(k\) is a proportional gain constant, \(P_t\) is the target
pressure, and \(P\) is the current pressure along that dimension. This
approach is similar to the method used to control the pressure by
{\hyperref[\detokenize{fix_press_berendsen::doc}]{\sphinxcrossref{\DUrole{doc}{fix press/berendsen}}}}. The target pressure
accepts either a constant numeric value or a LAMMPS {\hyperref[\detokenize{Howto_output:variable}]{\sphinxcrossref{\DUrole{std,std-ref}{variable}}}}.
Notably, this variable can be a function of time or other components of
the pressure tensor. By default, \(k\) has units of 1/(time * pressure)
although this will change if the \sphinxstyleemphasis{normalize/pressure} option is set as
{\hyperref[\detokenize{fix_deform_pressure:deform-normalize}]{\sphinxcrossref{\DUrole{std,std-ref}{discussed below}}}}. There is no proven method
to choosing an appropriate value of \(k\) as it will depend on the
specific details of a simulation. Testing different values is recommended.

\sphinxAtStartPar
By default, there is no limit on the resulting strain rate in any dimension.
A maximum limit can be applied using the {\hyperref[\detokenize{fix_deform_pressure:deform-max-rate}]{\sphinxcrossref{\DUrole{std,std-ref}{max/rate}}}}
option. Akin to {\hyperref[\detokenize{fix_nh::doc}]{\sphinxcrossref{\DUrole{doc}{fix npt and nph}}}}, pressures in different
dimensions can be coupled using the {\hyperref[\detokenize{fix_deform_pressure:deform-couple}]{\sphinxcrossref{\DUrole{std,std-ref}{couple}}}} option.
This means the instantaneous pressure along coupled dimensions are averaged
and the box strains identically along the coupled dimensions.

\sphinxAtStartPar
The \sphinxstyleemphasis{pressure/mean} style changes a dimension’s box length to maintain
a constant mean pressure defined as the trace of the pressure tensor.
This option has identical arguments to the \sphinxstyleemphasis{pressure} style and a similar
functional equation, except the current and target pressures refer to the
mean trace of the pressure tensor. All options for the \sphinxstyleemphasis{pressure} style
also apply to the \sphinxstyleemphasis{pressure/mean} style except for the
{\hyperref[\detokenize{fix_deform_pressure:deform-couple}]{\sphinxcrossref{\DUrole{std,std-ref}{couple}}}} option.

\sphinxAtStartPar
Note that while this style can be identical to coupled \sphinxstyleemphasis{pressure} styles,
it is generally not the same. For instance in 2D, a coupled \sphinxstyleemphasis{pressure}
style in the \sphinxstyleemphasis{x} and \sphinxstyleemphasis{y} dimensions would be equivalent to using the
\sphinxstyleemphasis{pressure/mean} style with identical settings in each dimension. However,
it would not be the same if settings (e.g. gain constants) were used in
the \sphinxstyleemphasis{x} and \sphinxstyleemphasis{y} dimensions or if the \sphinxstyleemphasis{pressure/mean} command was only applied
along one dimension.


\bigskip\hrule\bigskip


\sphinxAtStartPar
For the \sphinxstyleemphasis{xy}, \sphinxstyleemphasis{xz}, and \sphinxstyleemphasis{yz} parameters, this is the meaning of the
styles and values provided by this fix.  Note that changing the
tilt factors of a triclinic box does not change its volume.

\sphinxAtStartPar
The \sphinxstyleemphasis{pressure} style adjusts a tilt factor to control the corresponding
off\sphinxhyphen{}diagonal component of the pressure tensor. This option attempts to
maintain a specified target value using a linear controller where the
tilt factor T evolves according to the equation
\begin{equation*}
\begin{split}\frac{d T(t)}{dt} = L(t) k (P - P_t)\end{split}
\end{equation*}
\sphinxAtStartPar
where \(k\) is a proportional gain constant, \(P_t\) is the
target pressure, \(P\) is the current pressure, and \(L\) is
the perpendicular box length. The target pressure accepts either a
constant numeric value or a LAMMPS {\hyperref[\detokenize{Howto_output:variable}]{\sphinxcrossref{\DUrole{std,std-ref}{variable}}}}. Notably, this variable can be a function of time or other
components of the pressure tensor. By default, \(k\) has units of
1/(time * pressure) although this will change if the
\sphinxstyleemphasis{normalize/pessure} option is set as {\hyperref[\detokenize{fix_deform_pressure:deform-normalize}]{\sphinxcrossref{\DUrole{std,std-ref}{discussed below}}}}.  There is no proven method to choosing an
appropriate value of \(k\) as it will depend on the specific
details of a simulation and testing different values is
recommended. One can also apply a maximum limit to the magnitude of
the applied strain using the {\hyperref[\detokenize{fix_deform_pressure:deform-max-rate}]{\sphinxcrossref{\DUrole{std,std-ref}{max/rate}}}} option.

\sphinxAtStartPar
The \sphinxstyleemphasis{erate/rescale} style operates similarly to the \sphinxstyleemphasis{erate} style with
a specified strain rate in units of 1/time. The difference is that
the change in the tilt factor will depend on the current length of
the box perpendicular to the shear direction, \sphinxstyleemphasis{L}, instead of the
original length, \sphinxstyleemphasis{L0}. The tilt factor T as a function of time will
change as
\begin{equation*}
\begin{split}T(t) = T(t-1) + L \cdot \mathrm{erate} \cdot \Delta t\end{split}
\end{equation*}
\sphinxAtStartPar
where T(t\sphinxhyphen{}1) is the tilt factor on the previous timestep and
\(\Delta t\) is the timestep size. This option may be useful
in scenarios where \sphinxstyleemphasis{L} changes in time.


\bigskip\hrule\bigskip


\sphinxAtStartPar
The \sphinxstyleemphasis{box} parameter provides an additional control over the \sphinxstyleemphasis{x}, \sphinxstyleemphasis{y},
and \sphinxstyleemphasis{z} box lengths by isotropically dilating or contracting the box
to either maintain a fixed mean pressure or volume. This isotropic
scaling is applied after the box is deformed by the above \sphinxstyleemphasis{x}, \sphinxstyleemphasis{y},
\sphinxstyleemphasis{z}, \sphinxstyleemphasis{xy}, \sphinxstyleemphasis{xz}, and \sphinxstyleemphasis{yz} styles, acting as a second deformation
step. This parameter will change the overall strain rate in the \sphinxstyleemphasis{x},
\sphinxstyleemphasis{y}, or \sphinxstyleemphasis{z} dimensions.  This parameter can only be used in
combination with the \sphinxstyleemphasis{x}, \sphinxstyleemphasis{y}, or \sphinxstyleemphasis{z} commands: \sphinxstyleemphasis{vel}, \sphinxstyleemphasis{erate},
\sphinxstyleemphasis{trate}, \sphinxstyleemphasis{pressure}, or \sphinxstyleemphasis{wiggle}. This is the meaning of its styles
and values.

\sphinxAtStartPar
The \sphinxstyleemphasis{volume} style isotropically scales box lengths to maintain a constant
box volume in response to deformation from other parameters. This style
may be useful in scenarios where one wants to apply a constant deviatoric
pressure using \sphinxstyleemphasis{pressure} styles in the \sphinxstyleemphasis{x}, \sphinxstyleemphasis{y}, and \sphinxstyleemphasis{z} dimensions (
deforming the shape of the box), while maintaining a constant volume.

\sphinxAtStartPar
The \sphinxstyleemphasis{pressure} style isotropically scales box lengths in an attempt to
maintain a target mean pressure (the trace of the pressure tensor) of the
system. This is accomplished by isotropically scaling all box lengths
\(L\) by an additional factor of \(k (P_t - P_m)\) where \(k\)
is the proportional gain constant, \(P_t\) is the target pressure, and
\(P_m\) is the current mean pressure. This style may be useful in
scenarios where one wants to apply a constant deviatoric strain rate
using various strain\sphinxhyphen{}based styles (e.g. \sphinxstyleemphasis{trate}) along the \sphinxstyleemphasis{x}, \sphinxstyleemphasis{y}, and \sphinxstyleemphasis{z}
dimensions (deforming the shape of the box), while maintaining a mean pressure.


\bigskip\hrule\bigskip


\sphinxAtStartPar
The optional keywords provided by this fix are described below.

\phantomsection\label{\detokenize{fix_deform_pressure:deform-normalize}}
\sphinxAtStartPar
The \sphinxstyleemphasis{normalize/pressure} keyword changes how box dimensions evolve when
using the \sphinxstyleemphasis{pressure} or \sphinxstyleemphasis{pressure/mean} deformation styles. If the
\sphinxstyleemphasis{deform/normalize} value is set to \sphinxstyleemphasis{yes}, then the deviation from the
target pressure is normalized by the absolute value of the target
pressure such that the proportional gain constant scales a percentage
error and has units of 1/time. If the target pressure is ever zero, this
will produce an error unless the \sphinxstyleemphasis{max/rate} keyword is defined,
described below, which will cap the divergence.

\phantomsection\label{\detokenize{fix_deform_pressure:deform-max-rate}}
\sphinxAtStartPar
The \sphinxstyleemphasis{max/rate} keyword sets an upper threshold, \sphinxstyleemphasis{rate}, that limits the
maximum magnitude of the instantaneous strain rate applied in any dimension.
This keyword only applies to the \sphinxstyleemphasis{pressure} and \sphinxstyleemphasis{pressure/mean} options. If
a pressure\sphinxhyphen{}controlled rate is used for both \sphinxstyleemphasis{box} and either \sphinxstyleemphasis{x}, \sphinxstyleemphasis{y}, or
\sphinxstyleemphasis{z}, then this threshold will apply separately to each individual controller
such that the cumulative strain rate on a box dimension may be up to twice
the value of \sphinxstyleemphasis{rate}.

\phantomsection\label{\detokenize{fix_deform_pressure:deform-couple}}
\sphinxAtStartPar
The \sphinxstyleemphasis{couple} keyword allows two or three of the diagonal components of
the pressure tensor to be “coupled” together for the \sphinxstyleemphasis{pressure} option.
The value specified with the keyword determines which are coupled. For
example, \sphinxstyleemphasis{xz} means the \sphinxstyleemphasis{Pxx} and \sphinxstyleemphasis{Pzz} components of the stress tensor
are coupled. \sphinxstyleemphasis{Xyz} means all 3 diagonal components are coupled. Coupling
means two things: the instantaneous stress will be computed as an average
of the corresponding diagonal components, and the coupled box dimensions
will be changed together in lockstep, meaning coupled dimensions will be
dilated or contracted by the same percentage every timestep. If a \sphinxstyleemphasis{pressure}
style is defined for more than one coupled dimension, the target pressures
and gain constants must be identical. Alternatively, if a \sphinxstyleemphasis{pressure}
style is only defined for one of the coupled dimensions, its settings are
copied to other dimensions with undefined styles. \sphinxstyleemphasis{Couple xyz} can be used
for a 2d simulation; the \sphinxstyleemphasis{z} dimension is simply ignored.

\phantomsection\label{\detokenize{fix_deform_pressure:deform-balance}}
\sphinxAtStartPar
The \sphinxstyleemphasis{vol/balance/p} keyword modifies the behavior of the \sphinxstyleemphasis{volume} style when
applied to two of the \sphinxstyleemphasis{x}, \sphinxstyleemphasis{y}, and \sphinxstyleemphasis{z} dimensions. Instead of straining
the two dimensions in lockstep, the two dimensions are allowed to
separately dilate or contract in a manner to maintain a constant
volume while simultaneously trying to keep the pressure along each
dimension equal using a method described in {\hyperref[\detokenize{fix_deform_pressure:huang2014}]{\sphinxcrossref{\DUrole{std,std-ref}{(Huang2014)}}}}.


\bigskip\hrule\bigskip


\sphinxAtStartPar
If any pressure controls are used, this fix computes a temperature and
pressure each timestep. To do this, the fix creates its own computes
of style “temp” and “pressure”, as if these commands had been issued:

\begin{sphinxVerbatim}[commandchars=\\\{\}]
\PYG{k}{compute }\PYG{n+nv+nvIdentifier}{fix\PYGZhy{}ID\PYGZus{}temp}\PYG{+w}{ }\PYG{n+nv+nvIdentifier}{group\PYGZhy{}ID}\PYG{+w}{ }\PYG{n}{temp}
\PYG{k}{compute }\PYG{n+nv+nvIdentifier}{fix\PYGZhy{}ID\PYGZus{}press}\PYG{+w}{ }\PYG{n+nv+nvIdentifier}{group\PYGZhy{}ID}\PYG{+w}{ }\PYG{n}{pressure}\PYG{+w}{ }\PYG{n}{fix}\PYG{o}{\PYGZhy{}}\PYG{n}{ID\PYGZus{}temp}
\end{sphinxVerbatim}

\sphinxAtStartPar
See the {\hyperref[\detokenize{compute_temp::doc}]{\sphinxcrossref{\DUrole{doc}{compute temp}}}} and {\hyperref[\detokenize{compute_pressure::doc}]{\sphinxcrossref{\DUrole{doc}{compute pressure}}}} commands for details.  Note that the IDs of the
new computes are the fix\sphinxhyphen{}ID + underscore + “temp” or fix\_ID
+ underscore + “press”, and the group for the new computes is the same
as the fix group.

\sphinxAtStartPar
Note that these are NOT the computes used by thermodynamic output (see
the {\hyperref[\detokenize{thermo_style::doc}]{\sphinxcrossref{\DUrole{doc}{thermo\_style}}}} command) with ID =
\sphinxstyleemphasis{thermo\_temp} and \sphinxstyleemphasis{thermo\_press}.  This means you can change the
attributes of this fix’s temperature or pressure via the
{\hyperref[\detokenize{compute_modify::doc}]{\sphinxcrossref{\DUrole{doc}{compute\_modify}}}} command or print this
temperature or pressure during thermodynamic output via the
{\hyperref[\detokenize{thermo_style::doc}]{\sphinxcrossref{\DUrole{doc}{thermo\_style custom}}}} command using the
appropriate compute\sphinxhyphen{}ID. It also means that changing attributes of
\sphinxstyleemphasis{thermo\_temp} or \sphinxstyleemphasis{thermo\_press} will have no effect on this fix.


\subsection{Restart, fix\_modify, output, run start/stop, minimize info}
\label{\detokenize{fix_deform_pressure:restart-fix-modify-output-run-start-stop-minimize-info}}
\sphinxAtStartPar
This fix will restore the initial box settings from {\hyperref[\detokenize{restart::doc}]{\sphinxcrossref{\DUrole{doc}{binary
restart files}}}}, which allows the fix to be properly continue
deformation, when using the start/stop options of the {\hyperref[\detokenize{run::doc}]{\sphinxcrossref{\DUrole{doc}{run}}}}
command.  No global or per\sphinxhyphen{}atom quantities are stored by this fix for
access by various {\hyperref[\detokenize{Howto_output::doc}]{\sphinxcrossref{\DUrole{doc}{output commands}}}}.

\sphinxAtStartPar
If any pressure controls are used, the {\hyperref[\detokenize{fix_modify::doc}]{\sphinxcrossref{\DUrole{doc}{fix\_modify}}}}
\sphinxstyleemphasis{temp} and \sphinxstyleemphasis{press} options are supported by this fix, unlike in
{\hyperref[\detokenize{fix_deform::doc}]{\sphinxcrossref{\DUrole{doc}{fix deform}}}}.  You can use them to assign a
{\hyperref[\detokenize{compute::doc}]{\sphinxcrossref{\DUrole{doc}{compute}}}} you have defined to this fix which will be
used in its temperature and pressure calculations.  If you do this,
note that the kinetic energy derived from the compute temperature
should be consistent with the virial term computed using all atoms for
the pressure.  LAMMPS will warn you if you choose to compute
temperature on a subset of atoms.

\sphinxAtStartPar
This fix can perform deformation over multiple runs, using the \sphinxstyleemphasis{start}
and \sphinxstyleemphasis{stop} keywords of the {\hyperref[\detokenize{run::doc}]{\sphinxcrossref{\DUrole{doc}{run}}}} command.  See the
{\hyperref[\detokenize{run::doc}]{\sphinxcrossref{\DUrole{doc}{run}}}} command for details of how to do this.

\sphinxAtStartPar
This fix is not invoked during {\hyperref[\detokenize{minimize::doc}]{\sphinxcrossref{\DUrole{doc}{energy minimization}}}}.


\subsection{Restrictions}
\label{\detokenize{fix_deform_pressure:restrictions}}
\sphinxAtStartPar
This fix is part of the EXTRA\sphinxhyphen{}FIX package.  It is only enabled if LAMMPS
was built with that package.  See the {\hyperref[\detokenize{Build_package::doc}]{\sphinxcrossref{\DUrole{doc}{Build package}}}}
page for more info.

\sphinxAtStartPar
You cannot apply x, y, or z deformations to a dimension that is
shrink\sphinxhyphen{}wrapped via the {\hyperref[\detokenize{boundary::doc}]{\sphinxcrossref{\DUrole{doc}{boundary}}}} command.

\sphinxAtStartPar
You cannot apply xy, yz, or xz deformations to a second dimension (y
in xy) that is shrink\sphinxhyphen{}wrapped via the {\hyperref[\detokenize{boundary::doc}]{\sphinxcrossref{\DUrole{doc}{boundary}}}}
command.


\subsection{Related commands}
\label{\detokenize{fix_deform_pressure:related-commands}}
\sphinxAtStartPar
{\hyperref[\detokenize{fix_deform::doc}]{\sphinxcrossref{\DUrole{doc}{fix deform}}}}, {\hyperref[\detokenize{change_box::doc}]{\sphinxcrossref{\DUrole{doc}{change\_box}}}}


\subsection{Default}
\label{\detokenize{fix_deform_pressure:default}}
\sphinxAtStartPar
The option defaults are normalize/pressure = no.


\bigskip\hrule\bigskip

\phantomsection\label{\detokenize{fix_deform_pressure:huang2014}}
\sphinxAtStartPar
\sphinxstylestrong{(Huang2014)} X. Huang, “Exploring critical\sphinxhyphen{}state behavior using DEM”, Doctoral dissertation, Imperial College. (2014). \sphinxurl{https://doi.org/10.25560/25316}

\sphinxstepscope

\index{fix deposit@\spxentry{fix deposit}}\ignorespaces 

\section{fix deposit command}
\label{\detokenize{fix_deposit:fix-deposit-command}}\label{\detokenize{fix_deposit:index-0}}\label{\detokenize{fix_deposit::doc}}

\subsection{Syntax}
\label{\detokenize{fix_deposit:syntax}}
\begin{sphinxVerbatim}[commandchars=\\\{\}]
\PYG{k}{fix }\PYG{n+nv+nvIdentifier}{ID}\PYG{+w}{ }\PYG{n+nv+nvIdentifier}{group\PYGZhy{}ID}\PYG{+w}{ }\PYG{n}{deposit}\PYG{+w}{ }\PYG{n}{N}\PYG{+w}{ }\PYG{n}{type}\PYG{+w}{ }\PYG{n}{M}\PYG{+w}{ }\PYG{n}{seed}\PYG{+w}{ }\PYG{n}{keyword}\PYG{+w}{ }\PYG{n}{values}\PYG{+w}{ }\PYG{n}{...}
\end{sphinxVerbatim}
\begin{itemize}
\item {} 
\sphinxAtStartPar
ID, group\sphinxhyphen{}ID are documented in {\hyperref[\detokenize{fix::doc}]{\sphinxcrossref{\DUrole{doc}{fix}}}} command

\item {} 
\sphinxAtStartPar
deposit = style name of this fix command

\item {} 
\sphinxAtStartPar
N = \# of atoms or molecules to insert

\item {} 
\sphinxAtStartPar
type = atom type (1\sphinxhyphen{}Ntypes or type label) to assign to inserted atoms (offset for molecule insertion)

\item {} 
\sphinxAtStartPar
M = insert a single atom or molecule every M steps

\item {} 
\sphinxAtStartPar
seed = random \# seed (positive integer)

\item {} 
\sphinxAtStartPar
one or more keyword/value pairs may be appended to args

\item {} 
\sphinxAtStartPar
keyword = \sphinxstyleemphasis{region} or \sphinxstyleemphasis{var} or \sphinxstyleemphasis{set} or \sphinxstyleemphasis{id} or \sphinxstyleemphasis{global} or \sphinxstyleemphasis{local} or \sphinxstyleemphasis{near} or \sphinxstyleemphasis{gaussian} or \sphinxstyleemphasis{attempt} or \sphinxstyleemphasis{rate} or \sphinxstyleemphasis{vx} or \sphinxstyleemphasis{vy} or \sphinxstyleemphasis{vz} or \sphinxstyleemphasis{target} or \sphinxstyleemphasis{mol} or \sphinxstyleemphasis{molfrac} or \sphinxstyleemphasis{rigid} or \sphinxstyleemphasis{shake} or \sphinxstyleemphasis{orient} or \sphinxstyleemphasis{units}
\begin{sphinxalltt}
\sphinxstyleemphasis{region} value = region\sphinxhyphen{}ID
  region\sphinxhyphen{}ID = ID of region to use as insertion volume
\sphinxstyleemphasis{var} value = name = variable name to evaluate for test of atom creation
\sphinxstyleemphasis{set} values = dim name
  dim = \sphinxstyleemphasis{x} or \sphinxstyleemphasis{y} or \sphinxstyleemphasis{z}
  name = name of variable to set with x, y, or z atom position
\sphinxstyleemphasis{id} value = \sphinxstyleemphasis{max} or \sphinxstyleemphasis{next}
  max = atom ID for new atom(s) is max ID of all current atoms plus one
  next = atom ID for new atom(s) increments by one for every deposition
\sphinxstyleemphasis{global} values = lo hi
  lo,hi = put new atom/molecule a distance lo\sphinxhyphen{}hi above all other atoms (distance units)
\sphinxstyleemphasis{local} values = lo hi delta
  lo,hi = put new atom/molecule a distance lo\sphinxhyphen{}hi above any nearby atom beneath it (distance units)
  delta = lateral distance within which a neighbor is considered "nearby" (distance units)
\sphinxstyleemphasis{near} value = R
  R = only insert atom/molecule if further than R from existing particles (distance units)
\sphinxstyleemphasis{gaussian} values = xmid ymid zmid sigma
  xmid,ymid,zmid = center of the gaussian distribution (distance units)
  sigma = width of gaussian distribution (distance units)
\sphinxstyleemphasis{attempt} value = Q
  Q = attempt a single insertion up to Q times
\sphinxstyleemphasis{rate} value = V
  V = z velocity (y in 2d) at which insertion volume moves (velocity units)
\sphinxstyleemphasis{vx} values = vxlo vxhi
  vxlo,vxhi = range of x velocities for inserted atom/molecule (velocity units)
\sphinxstyleemphasis{vy} values = vylo vyhi
  vylo,vyhi = range of y velocities for inserted atom/molecule (velocity units)
\sphinxstyleemphasis{vz} values = vzlo vzhi
  vzlo,vzhi = range of z velocities for inserted atom/molecule (velocity units)
\sphinxstyleemphasis{target} values = tx ty tz
  tx,ty,tz = location of target point (distance units)
\sphinxstyleemphasis{mol} value = template\sphinxhyphen{}ID
  template\sphinxhyphen{}ID = ID of molecule template specified in a separate {\hyperref[\detokenize{molecule::doc}]{\sphinxcrossref{\DUrole{doc}{molecule}}}} command
\sphinxstyleemphasis{molfrac} values = f1 f2 ... fN
  f1 to fN = relative probability of creating each of N molecules in template\sphinxhyphen{}ID
\sphinxstyleemphasis{rigid} value = fix\sphinxhyphen{}ID
  fix\sphinxhyphen{}ID = ID of {\hyperref[\detokenize{fix_rigid::doc}]{\sphinxcrossref{\DUrole{doc}{fix rigid/small}}}} command
\sphinxstyleemphasis{shake} value = fix\sphinxhyphen{}ID
  fix\sphinxhyphen{}ID = ID of {\hyperref[\detokenize{fix_shake::doc}]{\sphinxcrossref{\DUrole{doc}{fix shake}}}} command
\sphinxstyleemphasis{orient} values = rx ry rz
  rx,ry,rz = vector to randomly rotate an inserted molecule around
\sphinxstyleemphasis{units} value = \sphinxstyleemphasis{lattice} or \sphinxstyleemphasis{box}
  lattice = the geometry is defined in lattice units
  box = the geometry is defined in simulation box units
\end{sphinxalltt}

\end{itemize}


\subsection{Examples}
\label{\detokenize{fix_deposit:examples}}
\begin{sphinxVerbatim}[commandchars=\\\{\}]
\PYG{k}{fix }\PYG{n+nv+nvIdentifier}{3}\PYG{+w}{ }\PYG{n+nv+nvIdentifier}{all}\PYG{+w}{ }\PYG{n}{deposit}\PYG{+w}{ }\PYG{l+m}{1000}\PYG{+w}{ }\PYG{l+m}{2}\PYG{+w}{ }\PYG{l+m}{100}\PYG{+w}{ }\PYG{l+m}{29494}\PYG{+w}{ }\PYG{k}{region }\PYG{n+nv+nvIdentifier}{myblock}\PYG{+w}{ }\PYG{n}{local}\PYG{+w}{ }\PYG{l+m}{1.0}\PYG{+w}{ }\PYG{l+m}{1.0}\PYG{+w}{ }\PYG{l+m}{1.0}\PYG{+w}{ }\PYG{n}{units}\PYG{+w}{ }\PYG{n}{box}
\PYG{k}{fix }\PYG{n+nv+nvIdentifier}{2}\PYG{+w}{ }\PYG{n+nv+nvIdentifier}{newatoms}\PYG{+w}{ }\PYG{n}{deposit}\PYG{+w}{ }\PYG{l+m}{10000}\PYG{+w}{ }\PYG{l+m}{1}\PYG{+w}{ }\PYG{l+m}{500}\PYG{+w}{ }\PYG{l+m}{12345}\PYG{+w}{ }\PYG{k}{region }\PYG{n+nv+nvIdentifier}{disk}\PYG{+w}{ }\PYG{n}{near}\PYG{+w}{ }\PYG{l+m}{2.0}\PYG{+w}{ }\PYG{n}{vz}\PYG{+w}{ }\PYG{o}{\PYGZhy{}}\PYG{l+m}{1.0}\PYG{+w}{ }\PYG{o}{\PYGZhy{}}\PYG{l+m}{0.8}
\PYG{k}{fix }\PYG{n+nv+nvIdentifier}{4}\PYG{+w}{ }\PYG{n+nv+nvIdentifier}{sputter}\PYG{+w}{ }\PYG{n}{deposit}\PYG{+w}{ }\PYG{l+m}{1000}\PYG{+w}{ }\PYG{l+m}{2}\PYG{+w}{ }\PYG{l+m}{500}\PYG{+w}{ }\PYG{l+m}{12235}\PYG{+w}{ }\PYG{k}{region }\PYG{n+nv+nvIdentifier}{sphere}\PYG{+w}{ }\PYG{n}{vz}\PYG{+w}{ }\PYG{o}{\PYGZhy{}}\PYG{l+m}{1.0}\PYG{+w}{ }\PYG{o}{\PYGZhy{}}\PYG{l+m}{1.0}\PYG{+w}{ }\PYG{n}{target}\PYG{+w}{ }\PYG{l+m}{5.0}\PYG{+w}{ }\PYG{l+m}{5.0}\PYG{+w}{ }\PYG{l+m}{0.0}\PYG{+w}{ }\PYG{n}{units}\PYG{+w}{ }\PYG{n}{lattice}
\PYG{k}{fix }\PYG{n+nv+nvIdentifier}{5}\PYG{+w}{ }\PYG{n+nv+nvIdentifier}{insert}\PYG{+w}{ }\PYG{n}{deposit}\PYG{+w}{ }\PYG{l+m}{200}\PYG{+w}{ }\PYG{l+m}{2}\PYG{+w}{ }\PYG{l+m}{100}\PYG{+w}{ }\PYG{l+m}{777}\PYG{+w}{ }\PYG{k}{region }\PYG{n+nv+nvIdentifier}{disk}\PYG{+w}{ }\PYG{n}{gaussian}\PYG{+w}{ }\PYG{l+m}{5.0}\PYG{+w}{ }\PYG{l+m}{5.0}\PYG{+w}{ }\PYG{l+m}{9.0}\PYG{+w}{ }\PYG{l+m}{1.0}\PYG{+w}{ }\PYG{n}{units}\PYG{+w}{ }\PYG{n}{box}

\PYG{k}{labelmap}\PYG{+w}{ }\PYG{n}{atom}\PYG{+w}{ }\PYG{l+m}{1}\PYG{+w}{ }\PYG{n}{Au}
\PYG{k}{fix }\PYG{n+nv+nvIdentifier}{4}\PYG{+w}{ }\PYG{n+nv+nvIdentifier}{sputter}\PYG{+w}{ }\PYG{n}{deposit}\PYG{+w}{ }\PYG{l+m}{1000}\PYG{+w}{ }\PYG{n}{Au}\PYG{+w}{ }\PYG{l+m}{500}\PYG{+w}{ }\PYG{l+m}{12235}\PYG{+w}{ }\PYG{k}{region }\PYG{n+nv+nvIdentifier}{sphere}\PYG{+w}{ }\PYG{n}{vz}\PYG{+w}{ }\PYG{o}{\PYGZhy{}}\PYG{l+m}{1.0}\PYG{+w}{ }\PYG{o}{\PYGZhy{}}\PYG{l+m}{1.0}\PYG{+w}{ }\PYG{n}{target}\PYG{+w}{ }\PYG{l+m}{5.0}\PYG{+w}{ }\PYG{l+m}{5.0}\PYG{+w}{ }\PYG{l+m}{0.0}\PYG{+w}{ }\PYG{n}{units}\PYG{+w}{ }\PYG{n}{lattice}
\end{sphinxVerbatim}


\subsection{Description}
\label{\detokenize{fix_deposit:description}}
\sphinxAtStartPar
Insert a single atom or molecule into the simulation domain every M
timesteps until N atoms or molecules have been inserted.  This is
useful for simulating deposition onto a surface.  For the remainder of
this doc page, a single inserted atom or molecule is referred to as a
“particle”.

\sphinxAtStartPar
If inserted particles are individual atoms, they are assigned the
specified atom type.  If they are molecules, the type of each atom in
the inserted molecule is specified in the file read by the
{\hyperref[\detokenize{molecule::doc}]{\sphinxcrossref{\DUrole{doc}{molecule}}}} command, and those values are added to the
specified atom type.  E.g. if the file specifies atom types 1,2,3, and
those are the atom types you want for inserted molecules, then specify
\sphinxstyleemphasis{type} = 0.  If you specify \sphinxstyleemphasis{type} = 2, the in the inserted molecule
will have atom types 3,4,5.

\sphinxAtStartPar
All atoms in the inserted particle are assigned to two groups: the
default group “all” and the group specified in the fix deposit command
(which can also be “all”).

\sphinxAtStartPar
If you are computing temperature values which include inserted
particles, you will want to use the {\hyperref[\detokenize{compute_modify::doc}]{\sphinxcrossref{\DUrole{doc}{compute\_modify dynamic/dof yes}}}} option, which ensures the current number of atoms is
used as a normalizing factor each time the temperature is computed.

\sphinxAtStartPar
Care must be taken that inserted particles are not too near existing
atoms, using the options described below.  When inserting particles
above a surface in a non\sphinxhyphen{}periodic box (see the
{\hyperref[\detokenize{boundary::doc}]{\sphinxcrossref{\DUrole{doc}{boundary}}}} command), the possibility of a particle
escaping the surface and flying upward should be considered, since the
particle may be lost or the box size may grow infinitely large.  A
{\hyperref[\detokenize{fix_wall_reflect::doc}]{\sphinxcrossref{\DUrole{doc}{fix wall/reflect}}}} command can be used to
prevent this behavior.  Note that if a shrink\sphinxhyphen{}wrap boundary is used,
it is OK to insert the new particle outside the box, however the box
will immediately be expanded to include the new particle. When
simulating a sputtering experiment it is probably more realistic to
ignore those atoms using the {\hyperref[\detokenize{thermo_modify::doc}]{\sphinxcrossref{\DUrole{doc}{thermo\_modify}}}}
command with the \sphinxstyleemphasis{lost ignore} option and a fixed
{\hyperref[\detokenize{boundary::doc}]{\sphinxcrossref{\DUrole{doc}{boundary}}}}.

\sphinxAtStartPar
The fix deposit command must use the \sphinxstyleemphasis{region} keyword to define an
insertion volume.  The specified region must have been previously
defined with a {\hyperref[\detokenize{region::doc}]{\sphinxcrossref{\DUrole{doc}{region}}}} command.  It must be defined with
side = \sphinxstyleemphasis{in}.

\begin{sphinxadmonition}{note}{Note:}
\sphinxAtStartPar
LAMMPS checks that the specified region is wholly inside the
simulation box.  It can do this correctly for orthonormal simulation
boxes.  However for {\hyperref[\detokenize{Howto_triclinic::doc}]{\sphinxcrossref{\DUrole{doc}{triclinic boxes}}}}, it only
tests against the larger orthonormal box that bounds the tilted
simulation box.  If the specified region includes volume outside the
tilted box, then an insertion will likely fail, leading to a “lost
atoms” error.  Thus for triclinic boxes you should ensure the
specified region is wholly inside the simulation box.
\end{sphinxadmonition}

\sphinxAtStartPar
The locations of inserted particles are taken from uniform distributed
random numbers, unless the \sphinxstyleemphasis{gaussian} keyword is used. Then the
individual coordinates are taken from a gaussian distribution of
width \sphinxstyleemphasis{sigma} centered on \sphinxstyleemphasis{xmid,ymid,zmid}.

\sphinxAtStartPar
Individual atoms are inserted, unless the \sphinxstyleemphasis{mol} keyword is used.  It
specifies a \sphinxstyleemphasis{template\sphinxhyphen{}ID} previously defined using the
{\hyperref[\detokenize{molecule::doc}]{\sphinxcrossref{\DUrole{doc}{molecule}}}} command, which reads files that define one or
more molecules.  The coordinates, atom types, charges, etc, as well as
any bond/angle/etc and special neighbor information for the molecule
can be specified in the molecule file.  See the
{\hyperref[\detokenize{molecule::doc}]{\sphinxcrossref{\DUrole{doc}{molecule}}}} command for details.  The only settings
required to be in each file are the coordinates and types of atoms in
the molecule.

\sphinxAtStartPar
If the molecule template contains more than one molecule, the relative
probability of depositing each molecule can be specified by the
\sphinxstyleemphasis{molfrac} keyword.  N relative probabilities, each from 0.0 to 1.0, are
specified, where N is the number of molecules in the template.  Each
time a molecule is deposited, a random number is used to sample from
the list of relative probabilities.  The N values must sum to 1.0.

\sphinxAtStartPar
If you wish to insert molecules via the \sphinxstyleemphasis{mol} keyword, that will be
treated as rigid bodies, use the \sphinxstyleemphasis{rigid} keyword, specifying as its
value the ID of a separate {\hyperref[\detokenize{fix_rigid::doc}]{\sphinxcrossref{\DUrole{doc}{fix rigid/small}}}}
command which also appears in your input script.

\begin{sphinxadmonition}{note}{Note:}
\sphinxAtStartPar
If you wish the new rigid molecules (and other rigid molecules) to be
thermostatted correctly via {\hyperref[\detokenize{fix_rigid::doc}]{\sphinxcrossref{\DUrole{doc}{fix rigid/small/nvt}}}} or
{\hyperref[\detokenize{fix_rigid::doc}]{\sphinxcrossref{\DUrole{doc}{fix rigid/small/npt}}}}, then you need to use the
{\hyperref[\detokenize{fix_modify::doc}]{\sphinxcrossref{\DUrole{doc}{fix\_modify dynamic/dof yes}}}} command for the rigid
fix.  This is to inform that fix that the molecule count will vary
dynamically.
\end{sphinxadmonition}

\sphinxAtStartPar
If you wish to insert molecules via the \sphinxstyleemphasis{mol} keyword, that will have
their bonds or angles constrained via SHAKE, use the \sphinxstyleemphasis{shake} keyword,
specifying as its value the ID of a separate {\hyperref[\detokenize{fix_shake::doc}]{\sphinxcrossref{\DUrole{doc}{fix shake}}}} command which also appears in your input script.

\sphinxAtStartPar
Each timestep a particle is inserted, the coordinates for its atoms
are chosen as follows.  For insertion of individual atoms, the
“position” referred to in the following description is the coordinate
of the atom.  For insertion of molecule, the “position” is the
geometric center of the molecule; see the {\hyperref[\detokenize{molecule::doc}]{\sphinxcrossref{\DUrole{doc}{molecule}}}} doc
page for details.  A random rotation of the molecule around its center
point is performed, which determines the coordinates all the
individual atoms.

\sphinxAtStartPar
A random position within the region insertion volume is generated.  If
neither the \sphinxstyleemphasis{global} or \sphinxstyleemphasis{local} keyword is used, the random position
is the trial position.  If the \sphinxstyleemphasis{global} keyword is used, the random
x,y values are used, but the z position of the new particle is set
above the highest current atom in the simulation by a distance
randomly chosen between lo/hi.  (For a 2d simulation, this is done for
the y position.)  If the \sphinxstyleemphasis{local} keyword is used, the z position is
set a distance between lo/hi above the highest current atom in the
simulation that is “nearby” the chosen x,y position.  In this context,
“nearby” means the lateral distance (in x,y) between the new and old
particles is less than the \sphinxstyleemphasis{delta} setting.

\sphinxAtStartPar
Once a trial x,y,z position has been selected, the insertion is only
performed if both the \sphinxstyleemphasis{near} and \sphinxstyleemphasis{var} keywords are satisfied (see below).
If either the \sphinxstyleemphasis{near} or the \sphinxstyleemphasis{var} keyword is not satisfied, a new random
position within the insertion volume is chosen and another trial is made.
Up to Q attempts are made.  If one or more particle insertions are not
successful, LAMMPS prints a warning message.

\sphinxAtStartPar
The \sphinxstyleemphasis{near} keyword ensures that no current atom in the simulation is within
a distance R of any atom in the new particle, including the effect of
periodic boundary conditions if applicable.  Note that the default value
for R is 0.0, which will allow atoms to strongly overlap if you are
inserting where other atoms are present.  This distance test is performed
independently for each atom in an inserted molecule, based on the randomly
rotated configuration of the molecule.

\begin{sphinxadmonition}{note}{Note:}
\sphinxAtStartPar
If you are inserting finite size particles or a molecule or
rigid body consisting of finite\sphinxhyphen{}size particles, then you should
typically set R larger than the distance at which any inserted
particle may overlap with either a previously inserted particle or an
existing particle.  LAMMPS will issue a warning if R is smaller than
this value, based on the radii of existing and inserted particles.
\end{sphinxadmonition}

\sphinxAtStartPar
\DUrole{versionmodified,added}{Added in version 21Nov2023.}

\sphinxAtStartPar
The \sphinxstyleemphasis{var} and \sphinxstyleemphasis{set} keywords can be used together to provide a criterion
for accepting or rejecting the addition of an individual atom, based on its
coordinates.  The \sphinxstyleemphasis{name} specified for the \sphinxstyleemphasis{var} keyword is the name of an
{\hyperref[\detokenize{variable::doc}]{\sphinxcrossref{\DUrole{doc}{equal\sphinxhyphen{}style variable}}}} that should evaluate to a zero or
non\sphinxhyphen{}zero value based on one or two or three variables that will store the
\sphinxstyleemphasis{x}, \sphinxstyleemphasis{y}, or \sphinxstyleemphasis{z} coordinates of an atom (one variable per coordinate).  If
used, these other variables must be {\hyperref[\detokenize{variable::doc}]{\sphinxcrossref{\DUrole{doc}{internal\sphinxhyphen{}style variables}}}} defined in the input script; their initial numeric value can be
anything. They must be internal\sphinxhyphen{}style variables, because this command
resets their values directly.  The \sphinxstyleemphasis{set} keyword is used to identify the
names of these other variables, one variable for the \sphinxstyleemphasis{x}\sphinxhyphen{}coordinate of a
created atom, one for \sphinxstyleemphasis{y}, and one for \sphinxstyleemphasis{z}.  When an atom is created, its
\((x,y,z)\) coordinates become the values for any \sphinxstyleemphasis{set} variable that
is defined.  The \sphinxstyleemphasis{var} variable is then evaluated.  If the returned value
is 0.0, the atom is not created.  If it is non\sphinxhyphen{}zero, the atom is created.
For an example of how to use these keywords, see the
{\hyperref[\detokenize{create_atoms::doc}]{\sphinxcrossref{\DUrole{doc}{create\_atoms}}}} command.

\sphinxAtStartPar
The \sphinxstyleemphasis{rate} option moves the insertion volume in the z direction (3d)
or y direction (2d).  This enables particles to be inserted from a
successively higher height over time.  Note that this parameter is
ignored if the \sphinxstyleemphasis{global} or \sphinxstyleemphasis{local} keywords are used, since those
options choose a z\sphinxhyphen{}coordinate for insertion independently.

\sphinxAtStartPar
The vx, vy, and vz components of velocity for the inserted particle
are set by sampling a uniform distribution between the bounds set by
the values specified for the \sphinxstyleemphasis{vx}, \sphinxstyleemphasis{vy}, and \sphinxstyleemphasis{vz} keywords. Note that
normally, new particles should be a assigned a negative vertical
velocity so that they move towards the surface.  For molecules, the
same velocity is given to every particle (no rotation or bond vibration).

\sphinxAtStartPar
If the \sphinxstyleemphasis{target} option is used, the velocity vector of the inserted
particle is changed so that it points from the insertion position
towards the specified target point.  The magnitude of the velocity is
unchanged.  This can be useful, for example, for simulating a
sputtering process.  E.g. the target point can be far away, so that
all incident particles strike the surface as if they are in an
incident beam of particles at a prescribed angle.

\sphinxAtStartPar
The \sphinxstyleemphasis{orient} keyword is only used when molecules are deposited.  By
default, each molecule is inserted at a random orientation.  If this
keyword is specified, then (rx,ry,rz) is used as an orientation
vector, and each inserted molecule is rotated around that vector with
a random value from zero to 2*PI.  For a 2d simulation, rx = ry = 0.0
is required, since rotations can only be performed around the z axis.

\sphinxAtStartPar
The \sphinxstyleemphasis{id} keyword determines how atom IDs and molecule IDs are assigned
to newly deposited particles.  Molecule IDs are only assigned if
molecules are being inserted.  For the \sphinxstyleemphasis{max} setting, the atom and
molecule IDs of all current atoms are checked.  Atoms in the new
particle are assigned IDs starting with the current maximum plus one.
If a molecule is inserted it is assigned an ID = current maximum plus
one.  This means that if particles leave the system, the new IDs may
replace the lost ones.  For the \sphinxstyleemphasis{next} setting, the maximum ID of any
atom and molecule is stored at the time the fix is defined.  Each time
a new particle is added, this value is incremented to assign IDs to
the new atom(s) or molecule.  Thus atom and molecule IDs for deposited
particles will be consecutive even if particles leave the system over
time.

\sphinxAtStartPar
The \sphinxstyleemphasis{units} keyword determines the meaning of the distance units used
for the other deposition parameters.  A \sphinxstyleemphasis{box} value selects standard
distance units as defined by the {\hyperref[\detokenize{units::doc}]{\sphinxcrossref{\DUrole{doc}{units}}}} command,
e.g. Angstroms for units = real or metal.  A \sphinxstyleemphasis{lattice} value means the
distance units are in lattice spacings.  The {\hyperref[\detokenize{lattice::doc}]{\sphinxcrossref{\DUrole{doc}{lattice}}}}
command must have been previously used to define the lattice spacing.
Note that the units choice affects all the keyword values that have
units of distance or velocity.

\begin{sphinxadmonition}{note}{Note:}
\sphinxAtStartPar
If you are monitoring the temperature of a system where the atom
count is changing due to adding particles, you typically should use
the {\hyperref[\detokenize{compute_modify::doc}]{\sphinxcrossref{\DUrole{doc}{compute\_modify dynamic/dof yes}}}} command
for the temperature compute you are using.
\end{sphinxadmonition}


\subsection{Restart, fix\_modify, output, run start/stop, minimize info}
\label{\detokenize{fix_deposit:restart-fix-modify-output-run-start-stop-minimize-info}}
\sphinxAtStartPar
This fix writes the state of the deposition to {\hyperref[\detokenize{restart::doc}]{\sphinxcrossref{\DUrole{doc}{binary restart files}}}}.  This includes information about how many
particles have been deposited, the random number generator seed, the
next timestep for deposition, etc.  See the
{\hyperref[\detokenize{read_restart::doc}]{\sphinxcrossref{\DUrole{doc}{read\_restart}}}} command for info on how to re\sphinxhyphen{}specify
a fix in an input script that reads a restart file, so that the
operation of the fix continues in an uninterrupted fashion.

\begin{sphinxadmonition}{note}{Note:}
\sphinxAtStartPar
For this to work correctly, the timestep must \sphinxstylestrong{not} be changed
after reading the restart with {\hyperref[\detokenize{reset_timestep::doc}]{\sphinxcrossref{\DUrole{doc}{reset\_timestep}}}}.
The fix will try to detect it and stop with an error.
\end{sphinxadmonition}

\sphinxAtStartPar
None of the {\hyperref[\detokenize{fix_modify::doc}]{\sphinxcrossref{\DUrole{doc}{fix\_modify}}}} options are relevant to this
fix.  This fix computes a global scalar, which can be accessed by various
output commands.  The scalar is the cumulative number of insertions.  The
scalar value calculated by this fix is “intensive”.  No parameter of this
fix can be used with the \sphinxstyleemphasis{start/stop} keywords of the {\hyperref[\detokenize{run::doc}]{\sphinxcrossref{\DUrole{doc}{run}}}}
command.  This fix is not invoked during {\hyperref[\detokenize{minimize::doc}]{\sphinxcrossref{\DUrole{doc}{energy minimization}}}}.


\subsection{Restrictions}
\label{\detokenize{fix_deposit:restrictions}}
\sphinxAtStartPar
The specified insertion region cannot be a “dynamic” region, as
defined by the {\hyperref[\detokenize{region::doc}]{\sphinxcrossref{\DUrole{doc}{region}}}} command.


\subsection{Related commands}
\label{\detokenize{fix_deposit:related-commands}}
\sphinxAtStartPar
{\hyperref[\detokenize{fix_pour::doc}]{\sphinxcrossref{\DUrole{doc}{fix pour}}}}, {\hyperref[\detokenize{region::doc}]{\sphinxcrossref{\DUrole{doc}{region}}}}


\subsection{Default}
\label{\detokenize{fix_deposit:default}}
\sphinxAtStartPar
Insertions are performed for individual atoms, i.e. no \sphinxstyleemphasis{mol} setting
is defined.  If the \sphinxstyleemphasis{mol} keyword is used, the default for \sphinxstyleemphasis{molfrac}
is an equal probabilities for all molecules in the template.
Additional option defaults are id = max, delta = 0.0, near = 0.0,
attempt = 10, rate = 0.0, vx = 0.0 0.0, vy = 0.0 0.0, vz = 0.0 0.0,
and units = lattice.

\sphinxstepscope

\index{fix dpd/energy@\spxentry{fix dpd/energy}}\ignorespaces 
\index{fix dpd/energy/kk@\spxentry{fix dpd/energy/kk}}\ignorespaces 

\section{fix dpd/energy command}
\label{\detokenize{fix_dpd_energy:fix-dpd-energy-command}}\label{\detokenize{fix_dpd_energy:index-1}}\label{\detokenize{fix_dpd_energy:index-0}}\label{\detokenize{fix_dpd_energy::doc}}
\sphinxAtStartPar
Accelerator Variants: \sphinxstyleemphasis{dpd/energy/kk}


\subsection{Syntax}
\label{\detokenize{fix_dpd_energy:syntax}}
\begin{sphinxVerbatim}[commandchars=\\\{\}]
\PYG{k}{fix }\PYG{n+nv+nvIdentifier}{ID}\PYG{+w}{ }\PYG{n+nv+nvIdentifier}{group\PYGZhy{}ID}\PYG{+w}{ }\PYG{n}{dpd}\PYG{o}{/}\PYG{n}{energy}
\end{sphinxVerbatim}
\begin{itemize}
\item {} 
\sphinxAtStartPar
ID, group\sphinxhyphen{}ID are documented in {\hyperref[\detokenize{fix::doc}]{\sphinxcrossref{\DUrole{doc}{fix}}}} command

\item {} 
\sphinxAtStartPar
dpd/energy = style name of this fix command

\end{itemize}


\subsection{Examples}
\label{\detokenize{fix_dpd_energy:examples}}
\begin{sphinxVerbatim}[commandchars=\\\{\}]
\PYG{k}{fix }\PYG{n+nv+nvIdentifier}{1}\PYG{+w}{ }\PYG{n+nv+nvIdentifier}{all}\PYG{+w}{ }\PYG{n}{dpd}\PYG{o}{/}\PYG{n}{energy}
\end{sphinxVerbatim}


\subsection{Description}
\label{\detokenize{fix_dpd_energy:description}}
\sphinxAtStartPar
Perform constant energy dissipative particle dynamics (DPD\sphinxhyphen{}E)
integration.  This fix updates the internal energies for particles in
the group at each timestep.  It must be used in conjunction with a
deterministic integrator (e.g. {\hyperref[\detokenize{fix_nve::doc}]{\sphinxcrossref{\DUrole{doc}{fix nve}}}}) that updates
the particle positions and velocities.

\sphinxAtStartPar
For fix \sphinxstyleemphasis{dpd/energy}, the particle internal temperature is related to
the particle internal energy through a mesoparticle equation of state.
An additional fix must be specified that defines the equation of state
for each particle, e.g. {\hyperref[\detokenize{fix_eos_cv::doc}]{\sphinxcrossref{\DUrole{doc}{fix eos/cv}}}}.

\sphinxAtStartPar
This fix must be used with the {\hyperref[\detokenize{pair_style::doc}]{\sphinxcrossref{\DUrole{doc}{pair\_style dpd/fdt/energy}}}} command.

\sphinxAtStartPar
Note that numerous variants of DPD can be specified by choosing an
appropriate combination of the integrator and {\hyperref[\detokenize{pair_style::doc}]{\sphinxcrossref{\DUrole{doc}{pair\_style dpd/fdt/energy}}}} command.  DPD under isoenergetic conditions
can be specified by using fix \sphinxstyleemphasis{dpd/energy}, fix \sphinxstyleemphasis{nve} and pair\_style
\sphinxstyleemphasis{dpd/fdt/energy}.  DPD under isoenthalpic conditions can
be specified by using fix \sphinxstyleemphasis{dpd/energy}, fix \sphinxstyleemphasis{nph} and pair\_style
\sphinxstyleemphasis{dpd/fdt/energy}.  Examples of each DPD variant are provided in the
examples/PACKAGES/dpd\sphinxhyphen{}react directory.


\bigskip\hrule\bigskip


\sphinxAtStartPar
Styles with a \sphinxstyleemphasis{gpu}, \sphinxstyleemphasis{intel}, \sphinxstyleemphasis{kk}, \sphinxstyleemphasis{omp}, or \sphinxstyleemphasis{opt} suffix are
functionally the same as the corresponding style without the suffix.
They have been optimized to run faster, depending on your available
hardware, as discussed on the {\hyperref[\detokenize{Speed_packages::doc}]{\sphinxcrossref{\DUrole{doc}{Accelerator packages}}}}
page.  The accelerated styles take the same arguments and should
produce the same results, except for round\sphinxhyphen{}off and precision issues.

\sphinxAtStartPar
These accelerated styles are part of the GPU, INTEL, KOKKOS,
OPENMP, and OPT packages, respectively.  They are only enabled if
LAMMPS was built with those packages.  See the {\hyperref[\detokenize{Build_package::doc}]{\sphinxcrossref{\DUrole{doc}{Build package}}}} page for more info.

\sphinxAtStartPar
You can specify the accelerated styles explicitly in your input script
by including their suffix, or you can use the {\hyperref[\detokenize{Run_options::doc}]{\sphinxcrossref{\DUrole{doc}{\sphinxhyphen{}suffix command\sphinxhyphen{}line switch}}}} when you invoke LAMMPS, or you can use the
{\hyperref[\detokenize{suffix::doc}]{\sphinxcrossref{\DUrole{doc}{suffix}}}} command in your input script.

\sphinxAtStartPar
See the {\hyperref[\detokenize{Speed_packages::doc}]{\sphinxcrossref{\DUrole{doc}{Accelerator packages}}}} page for more
instructions on how to use the accelerated styles effectively.


\bigskip\hrule\bigskip



\subsection{Restrictions}
\label{\detokenize{fix_dpd_energy:restrictions}}
\sphinxAtStartPar
This command is part of the DPD\sphinxhyphen{}REACT package.  It is only enabled if
LAMMPS was built with that package.  See the {\hyperref[\detokenize{Build_package::doc}]{\sphinxcrossref{\DUrole{doc}{Build package}}}} page for more info.

\sphinxAtStartPar
This fix must be used with an additional fix that specifies time
integration, e.g. {\hyperref[\detokenize{fix_nve::doc}]{\sphinxcrossref{\DUrole{doc}{fix nve}}}}.

\sphinxAtStartPar
The fix \sphinxstyleemphasis{dpd/energy} requires the \sphinxstyleemphasis{dpd} {\hyperref[\detokenize{atom_style::doc}]{\sphinxcrossref{\DUrole{doc}{atom\_style}}}}
to be used in order to properly account for the particle internal
energies and temperature.

\sphinxAtStartPar
The fix \sphinxstyleemphasis{dpd/energy} must be used with an additional fix that specifies the
mesoparticle equation of state for each particle.


\subsection{Related commands}
\label{\detokenize{fix_dpd_energy:related-commands}}
\sphinxAtStartPar
{\hyperref[\detokenize{fix_nve::doc}]{\sphinxcrossref{\DUrole{doc}{fix nve}}}} {\hyperref[\detokenize{fix_eos_cv::doc}]{\sphinxcrossref{\DUrole{doc}{fix eos/cv}}}}


\subsection{Default}
\label{\detokenize{fix_dpd_energy:default}}
\sphinxAtStartPar
none


\bigskip\hrule\bigskip

\phantomsection\label{\detokenize{fix_dpd_energy:lisal1}}
\sphinxAtStartPar
\sphinxstylestrong{(Lisal)} M. Lisal, J.K. Brennan, J. Bonet Avalos, J. Chem. Phys., 135, 204105 (2011).

\phantomsection\label{\detokenize{fix_dpd_energy:larentzos3}}
\sphinxAtStartPar
\sphinxstylestrong{(Larentzos)} J.P. Larentzos, J.K. Brennan, J.D. Moore, and W.D. Mattson, ARL\sphinxhyphen{}TR\sphinxhyphen{}6863, U.S. Army Research Laboratory, Aberdeen Proving Ground, MD (2014).

\sphinxstepscope

\index{fix edpd/source@\spxentry{fix edpd/source}}\ignorespaces 
\index{fix tdpd/source@\spxentry{fix tdpd/source}}\ignorespaces 

\section{fix edpd/source command}
\label{\detokenize{fix_dpd_source:fix-edpd-source-command}}\label{\detokenize{fix_dpd_source:index-1}}\label{\detokenize{fix_dpd_source:index-0}}\label{\detokenize{fix_dpd_source::doc}}

\section{fix tdpd/source command}
\label{\detokenize{fix_dpd_source:fix-tdpd-source-command}}

\subsection{Syntax}
\label{\detokenize{fix_dpd_source:syntax}}
\begin{sphinxVerbatim}[commandchars=\\\{\}]
\PYG{k}{fix }\PYG{n+nv+nvIdentifier}{ID}\PYG{+w}{ }\PYG{n+nv+nvIdentifier}{group\PYGZhy{}ID}\PYG{+w}{ }\PYG{n}{edpd}\PYG{o}{/}\PYG{n}{source}\PYG{+w}{ }\PYG{n}{keyword}\PYG{+w}{ }\PYG{n}{values}\PYG{+w}{ }\PYG{n}{...}
\PYG{k}{fix }\PYG{n+nv+nvIdentifier}{ID}\PYG{+w}{ }\PYG{n+nv+nvIdentifier}{group\PYGZhy{}ID}\PYG{+w}{ }\PYG{n}{tdpd}\PYG{o}{/}\PYG{n}{source}\PYG{+w}{ }\PYG{n}{cc\PYGZus{}index}\PYG{+w}{ }\PYG{n}{keyword}\PYG{+w}{ }\PYG{n}{values}\PYG{+w}{ }\PYG{n}{...}
\end{sphinxVerbatim}
\begin{itemize}
\item {} 
\sphinxAtStartPar
ID, group\sphinxhyphen{}ID are documented in {\hyperref[\detokenize{fix::doc}]{\sphinxcrossref{\DUrole{doc}{fix}}}} command

\item {} 
\sphinxAtStartPar
edpd/source or tdpd/source = style name of this fix command

\item {} 
\sphinxAtStartPar
index (only specified for tdpd/source) = index of chemical species (1 to Nspecies)

\item {} 
\sphinxAtStartPar
keyword = \sphinxstyleemphasis{sphere} or \sphinxstyleemphasis{cuboid} or \sphinxstyleemphasis{region}
\begin{sphinxalltt}
\sphinxstyleemphasis{sphere} args = cx cy cz radius source
  cx,cy,cz = x,y,z center of spherical domain (distance units)
  radius = radius of a spherical domain (distance units)
  source = heat source or concentration source (flux units, see below)
\sphinxstyleemphasis{cuboid} values = cx cy cz dLx dLy dLz source
  cx,cy,cz = x,y,z center of a cuboid domain (distance units)
  dLx,dLy,dLz = x,y,z side length of a cuboid domain (distance units)
  source = heat source or concentration source (flux units, see below)
\sphinxstyleemphasis{region} values = region\sphinxhyphen{}ID source
  region = ID of region for heat or concentration source
  source = heat source or concentration source (flux units, see below)
\end{sphinxalltt}

\end{itemize}


\subsection{Examples}
\label{\detokenize{fix_dpd_source:examples}}
\begin{sphinxVerbatim}[commandchars=\\\{\}]
\PYG{k}{fix }\PYG{n+nv+nvIdentifier}{1}\PYG{+w}{ }\PYG{n+nv+nvIdentifier}{all}\PYG{+w}{ }\PYG{n}{edpd}\PYG{o}{/}\PYG{n}{source}\PYG{+w}{ }\PYG{n}{sphere}\PYG{+w}{ }\PYG{l+m}{0.0}\PYG{+w}{ }\PYG{l+m}{0.0}\PYG{+w}{ }\PYG{l+m}{0.0}\PYG{+w}{ }\PYG{l+m}{5.0}\PYG{+w}{ }\PYG{l+m}{0.01}
\PYG{k}{fix }\PYG{n+nv+nvIdentifier}{1}\PYG{+w}{ }\PYG{n+nv+nvIdentifier}{all}\PYG{+w}{ }\PYG{n}{edpd}\PYG{o}{/}\PYG{n}{source}\PYG{+w}{ }\PYG{n}{cuboid}\PYG{+w}{ }\PYG{l+m}{0.0}\PYG{+w}{ }\PYG{l+m}{0.0}\PYG{+w}{ }\PYG{l+m}{0.0}\PYG{+w}{ }\PYG{l+m}{20.0}\PYG{+w}{ }\PYG{l+m}{10.0}\PYG{+w}{ }\PYG{l+m}{10.0}\PYG{+w}{ }\PYG{o}{\PYGZhy{}}\PYG{l+m}{0.01}
\PYG{k}{fix }\PYG{n+nv+nvIdentifier}{1}\PYG{+w}{ }\PYG{n+nv+nvIdentifier}{all}\PYG{+w}{ }\PYG{n}{tdpd}\PYG{o}{/}\PYG{n}{source}\PYG{+w}{ }\PYG{l+m}{1}\PYG{+w}{ }\PYG{n}{sphere}\PYG{+w}{ }\PYG{l+m}{5.0}\PYG{+w}{ }\PYG{l+m}{0.0}\PYG{+w}{ }\PYG{l+m}{0.0}\PYG{+w}{ }\PYG{l+m}{5.0}\PYG{+w}{ }\PYG{l+m}{0.01}
\PYG{k}{fix }\PYG{n+nv+nvIdentifier}{1}\PYG{+w}{ }\PYG{n+nv+nvIdentifier}{all}\PYG{+w}{ }\PYG{n}{tdpd}\PYG{o}{/}\PYG{n}{source}\PYG{+w}{ }\PYG{l+m}{2}\PYG{+w}{ }\PYG{n}{cuboid}\PYG{+w}{ }\PYG{l+m}{0.0}\PYG{+w}{ }\PYG{l+m}{0.0}\PYG{+w}{ }\PYG{l+m}{0.0}\PYG{+w}{ }\PYG{l+m}{20.0}\PYG{+w}{ }\PYG{l+m}{10.0}\PYG{+w}{ }\PYG{l+m}{10.0}\PYG{+w}{ }\PYG{l+m}{0.01}
\PYG{k}{fix }\PYG{n+nv+nvIdentifier}{1}\PYG{+w}{ }\PYG{n+nv+nvIdentifier}{all}\PYG{+w}{ }\PYG{n}{tdpd}\PYG{o}{/}\PYG{n}{source}\PYG{+w}{ }\PYG{l+m}{1}\PYG{+w}{ }\PYG{k}{region }\PYG{n+nv+nvIdentifier}{lower}\PYG{+w}{ }\PYG{o}{\PYGZhy{}}\PYG{l+m}{0.01}
\end{sphinxVerbatim}


\subsection{Description}
\label{\detokenize{fix_dpd_source:description}}
\sphinxAtStartPar
Fix \sphinxstyleemphasis{edpd/source} adds a heat source as an external heat flux to each
atom in a spherical or cuboid domain, where the \sphinxstyleemphasis{source} is in units
of energy/time.  Fix \sphinxstyleemphasis{tdpd/source} adds an external concentration
source of the chemical species specified by \sphinxstyleemphasis{index} as an external
concentration flux for each atom in a spherical or cuboid domain,
where the \sphinxstyleemphasis{source} is in units of mole/volume/time.

\sphinxAtStartPar
This command can be used to give an additional heat/concentration
source term to atoms in a simulation, such as for a simulation of a
heat conduction with a source term (see Fig.12 in {\hyperref[\detokenize{fix_dpd_source:li2014b}]{\sphinxcrossref{\DUrole{std,std-ref}{(Li2014)}}}})
or diffusion with a source term (see Fig.1 in {\hyperref[\detokenize{fix_dpd_source:li2015b}]{\sphinxcrossref{\DUrole{std,std-ref}{(Li2015)}}}}), as
an analog of a periodic Poiseuille flow problem.

\sphinxAtStartPar
\DUrole{versionmodified,deprecated}{Deprecated since version 15Jun2023: }The \sphinxstyleemphasis{sphere} and \sphinxstyleemphasis{cuboid} keywords will be removed in a future version
of LAMMPS.  The same functionality and more can be achieved with a region.

\sphinxAtStartPar
If the \sphinxstyleemphasis{sphere} keyword is used, the \sphinxstyleemphasis{cx, cy, cz, radius} values define
a spherical domain to apply the source flux to.

\sphinxAtStartPar
If the \sphinxstyleemphasis{cuboid} keyword is used, the \sphinxstyleemphasis{cx, cy, cz, dLx, dLy, dLz} define
a cuboid domain to apply the source flux to.

\sphinxAtStartPar
If the \sphinxstyleemphasis{region} keyword is used, the \sphinxstyleemphasis{region\sphinxhyphen{}ID} selects which
{\hyperref[\detokenize{region::doc}]{\sphinxcrossref{\DUrole{doc}{region}}}} to apply the source flux to.


\bigskip\hrule\bigskip



\subsection{Restart, fix\_modify, output, run start/stop, minimize info}
\label{\detokenize{fix_dpd_source:restart-fix-modify-output-run-start-stop-minimize-info}}
\sphinxAtStartPar
No information of these fixes is written to {\hyperref[\detokenize{restart::doc}]{\sphinxcrossref{\DUrole{doc}{binary restart files}}}}.  None of the {\hyperref[\detokenize{fix_modify::doc}]{\sphinxcrossref{\DUrole{doc}{fix\_modify}}}} options are
relevant to these fixes.  No global or per\sphinxhyphen{}atom quantities are stored by
these fixes for access by various {\hyperref[\detokenize{Howto_output::doc}]{\sphinxcrossref{\DUrole{doc}{output commands}}}}.
No parameter of these fixes can be used with the \sphinxstyleemphasis{start/stop} keywords
of the {\hyperref[\detokenize{run::doc}]{\sphinxcrossref{\DUrole{doc}{run}}}} command.  These fixes are not invoked during
{\hyperref[\detokenize{minimize::doc}]{\sphinxcrossref{\DUrole{doc}{energy minimization}}}}.


\subsection{Restrictions}
\label{\detokenize{fix_dpd_source:restrictions}}
\sphinxAtStartPar
These fixes are part of the DPD\sphinxhyphen{}MESO package.  They are only enabled if
LAMMPS was built with that package.  See the {\hyperref[\detokenize{Build_package::doc}]{\sphinxcrossref{\DUrole{doc}{Build package}}}} page for more info.

\sphinxAtStartPar
Fix \sphinxstyleemphasis{edpd/source} must be used with the {\hyperref[\detokenize{pair_mesodpd::doc}]{\sphinxcrossref{\DUrole{doc}{pair\_style edpd}}}} command.  Fix \sphinxstyleemphasis{tdpd/source} must be used with the
{\hyperref[\detokenize{pair_mesodpd::doc}]{\sphinxcrossref{\DUrole{doc}{pair\_style tdpd}}}} command.


\subsection{Related commands}
\label{\detokenize{fix_dpd_source:related-commands}}
\sphinxAtStartPar
{\hyperref[\detokenize{pair_mesodpd::doc}]{\sphinxcrossref{\DUrole{doc}{pair\_style edpd}}}}, {\hyperref[\detokenize{pair_mesodpd::doc}]{\sphinxcrossref{\DUrole{doc}{pair\_style tdpd}}}},
{\hyperref[\detokenize{compute_edpd_temp_atom::doc}]{\sphinxcrossref{\DUrole{doc}{compute edpd/temp/atom}}}},
{\hyperref[\detokenize{compute_tdpd_cc_atom::doc}]{\sphinxcrossref{\DUrole{doc}{compute tdpd/cc/atom}}}}


\subsection{Default}
\label{\detokenize{fix_dpd_source:default}}
\sphinxAtStartPar
none


\bigskip\hrule\bigskip

\phantomsection\label{\detokenize{fix_dpd_source:li2014b}}
\sphinxAtStartPar
\sphinxstylestrong{(Li2014)} Z. Li, Y.\sphinxhyphen{}H. Tang, H. Lei, B. Caswell and G.E. Karniadakis,
“Energy\sphinxhyphen{}conserving dissipative particle dynamics with
temperature\sphinxhyphen{}dependent properties”, J. Comput. Phys., 265: 113\sphinxhyphen{}127
(2014). DOI: 10.1016/j.jcp.2014.02.003

\phantomsection\label{\detokenize{fix_dpd_source:li2015b}}
\sphinxAtStartPar
\sphinxstylestrong{(Li2015)} Z. Li, A. Yazdani, A. Tartakovsky and G.E. Karniadakis,
“Transport dissipative particle dynamics model for mesoscopic
advection\sphinxhyphen{}diffusion\sphinxhyphen{}reaction problems”, J. Chem. Phys., 143: 014101
(2015).  DOI: 10.1063/1.4923254

\sphinxstepscope

\index{fix drag@\spxentry{fix drag}}\ignorespaces 

\section{fix drag command}
\label{\detokenize{fix_drag:fix-drag-command}}\label{\detokenize{fix_drag:index-0}}\label{\detokenize{fix_drag::doc}}

\subsection{Syntax}
\label{\detokenize{fix_drag:syntax}}
\begin{sphinxVerbatim}[commandchars=\\\{\}]
\PYG{k}{fix }\PYG{n+nv+nvIdentifier}{ID}\PYG{+w}{ }\PYG{n+nv+nvIdentifier}{group\PYGZhy{}ID}\PYG{+w}{ }\PYG{n}{drag}\PYG{+w}{ }\PYG{n}{x}\PYG{+w}{ }\PYG{n}{y}\PYG{+w}{ }\PYG{n}{z}\PYG{+w}{ }\PYG{n}{fmag}\PYG{+w}{ }\PYG{n}{delta}
\end{sphinxVerbatim}
\begin{itemize}
\item {} 
\sphinxAtStartPar
ID, group\sphinxhyphen{}ID are documented in {\hyperref[\detokenize{fix::doc}]{\sphinxcrossref{\DUrole{doc}{fix}}}} command

\item {} 
\sphinxAtStartPar
drag = style name of this fix command

\item {} 
\sphinxAtStartPar
x,y,z = coord to drag atoms towards

\item {} 
\sphinxAtStartPar
fmag = magnitude of force to apply to each atom (force units)

\item {} 
\sphinxAtStartPar
delta = cutoff distance inside of which force         is not applied (distance units)

\end{itemize}


\subsection{Examples}
\label{\detokenize{fix_drag:examples}}
\begin{sphinxVerbatim}[commandchars=\\\{\}]
\PYG{k}{fix }\PYG{n+nv+nvIdentifier}{center}\PYG{+w}{ }\PYG{n+nv+nvIdentifier}{small\PYGZhy{}molecule}\PYG{+w}{ }\PYG{n}{drag}\PYG{+w}{ }\PYG{l+m}{0.0}\PYG{+w}{ }\PYG{l+m}{10.0}\PYG{+w}{ }\PYG{l+m}{0.0}\PYG{+w}{ }\PYG{l+m}{5.0}\PYG{+w}{ }\PYG{l+m}{2.0}
\end{sphinxVerbatim}


\subsection{Description}
\label{\detokenize{fix_drag:description}}
\sphinxAtStartPar
Apply a force to each atom in a group to drag it towards the point
(x,y,z).  The magnitude of the force is specified by fmag.  If an atom
is closer than a distance delta to the point, then the force is not
applied.

\sphinxAtStartPar
Any of the x,y,z values can be specified as NULL which means do not
include that dimension in the distance calculation or force
application.

\sphinxAtStartPar
This command can be used to steer one or more atoms to a new location
in the simulation.


\subsection{Restart, fix\_modify, output, run start/stop, minimize info}
\label{\detokenize{fix_drag:restart-fix-modify-output-run-start-stop-minimize-info}}
\sphinxAtStartPar
No information about this fix is written to {\hyperref[\detokenize{restart::doc}]{\sphinxcrossref{\DUrole{doc}{binary restart files}}}}.

\sphinxAtStartPar
The {\hyperref[\detokenize{fix_modify::doc}]{\sphinxcrossref{\DUrole{doc}{fix\_modify}}}} \sphinxstyleemphasis{respa} option is supported by this
fix. This allows to set at which level of the {\hyperref[\detokenize{run_style::doc}]{\sphinxcrossref{\DUrole{doc}{r\sphinxhyphen{}RESPA}}}}
integrator the fix is adding its forces. Default is the outermost level.

\sphinxAtStartPar
This fix computes a global 3\sphinxhyphen{}vector of forces, which can be accessed
by various {\hyperref[\detokenize{Howto_output::doc}]{\sphinxcrossref{\DUrole{doc}{output commands}}}}.  This is the total
force on the group of atoms by the drag force.  The vector values
calculated by this fix are “extensive”.

\sphinxAtStartPar
No parameter of this fix can be used with the \sphinxstyleemphasis{start/stop} keywords of
the {\hyperref[\detokenize{run::doc}]{\sphinxcrossref{\DUrole{doc}{run}}}} command.  This fix is not invoked during {\hyperref[\detokenize{minimize::doc}]{\sphinxcrossref{\DUrole{doc}{energy minimization}}}}.


\subsection{Restrictions}
\label{\detokenize{fix_drag:restrictions}}
\sphinxAtStartPar
This fix is part of the EXTRA\sphinxhyphen{}FIX package.  It is only enabled if LAMMPS
was built with that package.  See the {\hyperref[\detokenize{Build_package::doc}]{\sphinxcrossref{\DUrole{doc}{Build package}}}} page for more info.


\subsection{Related commands}
\label{\detokenize{fix_drag:related-commands}}
\sphinxAtStartPar
{\hyperref[\detokenize{fix_spring::doc}]{\sphinxcrossref{\DUrole{doc}{fix spring}}}}, {\hyperref[\detokenize{fix_spring_self::doc}]{\sphinxcrossref{\DUrole{doc}{fix spring/self}}}},
{\hyperref[\detokenize{fix_spring_rg::doc}]{\sphinxcrossref{\DUrole{doc}{fix spring/rg}}}}, {\hyperref[\detokenize{fix_smd::doc}]{\sphinxcrossref{\DUrole{doc}{fix smd}}}}


\subsection{Default}
\label{\detokenize{fix_drag:default}}
\sphinxAtStartPar
none

\sphinxstepscope

\index{fix drude@\spxentry{fix drude}}\ignorespaces 

\section{fix drude command}
\label{\detokenize{fix_drude:fix-drude-command}}\label{\detokenize{fix_drude:index-0}}\label{\detokenize{fix_drude::doc}}

\subsection{Syntax}
\label{\detokenize{fix_drude:syntax}}
\begin{sphinxVerbatim}[commandchars=\\\{\}]
\PYG{k}{fix }\PYG{n+nv+nvIdentifier}{ID}\PYG{+w}{ }\PYG{n+nv+nvIdentifier}{group\PYGZhy{}ID}\PYG{+w}{ }\PYG{n}{drude}\PYG{+w}{ }\PYG{n}{flag1}\PYG{+w}{ }\PYG{n}{flag2}\PYG{+w}{ }\PYG{n}{...}\PYG{+w}{ }\PYG{n}{flagN}
\end{sphinxVerbatim}
\begin{itemize}
\item {} 
\sphinxAtStartPar
ID, group\sphinxhyphen{}ID are documented in {\hyperref[\detokenize{fix::doc}]{\sphinxcrossref{\DUrole{doc}{fix}}}} command

\item {} 
\sphinxAtStartPar
drude = style name of this fix command

\item {} 
\sphinxAtStartPar
flag1 flag2 … flagN = Drude flag for each atom type (1 to N) in the system

\end{itemize}


\subsection{Examples}
\label{\detokenize{fix_drude:examples}}
\begin{sphinxVerbatim}[commandchars=\\\{\}]
\PYG{k}{fix }\PYG{n+nv+nvIdentifier}{1}\PYG{+w}{ }\PYG{n+nv+nvIdentifier}{all}\PYG{+w}{ }\PYG{n}{drude}\PYG{+w}{ }\PYG{l+m}{1}\PYG{+w}{ }\PYG{l+m}{1}\PYG{+w}{ }\PYG{l+m}{0}\PYG{+w}{ }\PYG{l+m}{1}\PYG{+w}{ }\PYG{l+m}{0}\PYG{+w}{ }\PYG{l+m}{2}\PYG{+w}{ }\PYG{l+m}{2}\PYG{+w}{ }\PYG{l+m}{2}
\PYG{k}{fix }\PYG{n+nv+nvIdentifier}{1}\PYG{+w}{ }\PYG{n+nv+nvIdentifier}{all}\PYG{+w}{ }\PYG{n}{drude}\PYG{+w}{ }\PYG{n}{C}\PYG{+w}{ }\PYG{n}{C}\PYG{+w}{ }\PYG{n}{N}\PYG{+w}{ }\PYG{n}{C}\PYG{+w}{ }\PYG{n}{N}\PYG{+w}{ }\PYG{n}{D}\PYG{+w}{ }\PYG{n}{D}\PYG{+w}{ }\PYG{n}{D}
\end{sphinxVerbatim}

\sphinxAtStartPar
Example input scripts available: examples/PACKAGES/drude


\subsection{Description}
\label{\detokenize{fix_drude:description}}
\sphinxAtStartPar
Assign each atom type in the system to be one of 3 kinds of atoms
within the Drude polarization model. This fix is designed to be used
with the {\hyperref[\detokenize{Howto_drude::doc}]{\sphinxcrossref{\DUrole{doc}{thermalized Drude oscillator model}}}}.
Polarizable models in LAMMPS are described on the {\hyperref[\detokenize{Howto_polarizable::doc}]{\sphinxcrossref{\DUrole{doc}{Howto polarizable}}}} doc page.

\sphinxAtStartPar
The three possible types can be designated with an integer (0,1,2)
or capital letter (N,C,D):
\begin{itemize}
\item {} 
\sphinxAtStartPar
0 or N = non\sphinxhyphen{}polarizable atom (not part of Drude model)

\item {} 
\sphinxAtStartPar
1 or C = Drude core

\item {} 
\sphinxAtStartPar
2 or D = Drude electron

\end{itemize}


\subsection{Restrictions}
\label{\detokenize{fix_drude:restrictions}}
\sphinxAtStartPar
This fix should be invoked before any other commands that implement
the Drude oscillator model, such as {\hyperref[\detokenize{fix_langevin_drude::doc}]{\sphinxcrossref{\DUrole{doc}{fix langevin/drude}}}}, {\hyperref[\detokenize{fix_tgnh_drude::doc}]{\sphinxcrossref{\DUrole{doc}{fix tgnvt/drude}}}}, {\hyperref[\detokenize{fix_drude_transform::doc}]{\sphinxcrossref{\DUrole{doc}{fix drude/transform}}}}, {\hyperref[\detokenize{compute_temp_drude::doc}]{\sphinxcrossref{\DUrole{doc}{compute temp/drude}}}}, {\hyperref[\detokenize{pair_thole::doc}]{\sphinxcrossref{\DUrole{doc}{pair\_style thole}}}}.


\subsection{Related commands}
\label{\detokenize{fix_drude:related-commands}}
\sphinxAtStartPar
{\hyperref[\detokenize{fix_langevin_drude::doc}]{\sphinxcrossref{\DUrole{doc}{fix langevin/drude}}}}, {\hyperref[\detokenize{fix_tgnh_drude::doc}]{\sphinxcrossref{\DUrole{doc}{fix tgnvt/drude}}}}, {\hyperref[\detokenize{fix_drude_transform::doc}]{\sphinxcrossref{\DUrole{doc}{fix drude/transform}}}}, {\hyperref[\detokenize{compute_temp_drude::doc}]{\sphinxcrossref{\DUrole{doc}{compute temp/drude}}}}, {\hyperref[\detokenize{pair_thole::doc}]{\sphinxcrossref{\DUrole{doc}{pair\_style thole}}}}


\subsection{Default}
\label{\detokenize{fix_drude:default}}
\sphinxAtStartPar
none

\sphinxstepscope

\index{fix drude/transform/direct@\spxentry{fix drude/transform/direct}}\ignorespaces 
\index{fix drude/transform/inverse@\spxentry{fix drude/transform/inverse}}\ignorespaces 

\section{fix drude/transform/direct command}
\label{\detokenize{fix_drude_transform:fix-drude-transform-direct-command}}\label{\detokenize{fix_drude_transform:index-1}}\label{\detokenize{fix_drude_transform:index-0}}\label{\detokenize{fix_drude_transform::doc}}

\section{fix drude/transform/inverse command}
\label{\detokenize{fix_drude_transform:fix-drude-transform-inverse-command}}

\subsection{Syntax}
\label{\detokenize{fix_drude_transform:syntax}}
\begin{sphinxVerbatim}[commandchars=\\\{\}]
\PYG{k}{fix }\PYG{n+nv+nvIdentifier}{ID}\PYG{+w}{ }\PYG{n+nv+nvIdentifier}{group\PYGZhy{}ID}\PYG{+w}{ }\PYG{n}{style}\PYG{+w}{ }\PYG{n}{keyword}\PYG{+w}{ }\PYG{n}{value}\PYG{+w}{ }\PYG{n}{...}
\end{sphinxVerbatim}
\begin{itemize}
\item {} 
\sphinxAtStartPar
ID, group\sphinxhyphen{}ID are documented in {\hyperref[\detokenize{fix::doc}]{\sphinxcrossref{\DUrole{doc}{fix}}}} command

\item {} 
\sphinxAtStartPar
style = \sphinxstyleemphasis{drude/transform/direct} or \sphinxstyleemphasis{drude/transform/inverse}

\end{itemize}


\subsection{Examples}
\label{\detokenize{fix_drude_transform:examples}}
\begin{sphinxVerbatim}[commandchars=\\\{\}]
\PYG{k}{fix }\PYG{n+nv+nvIdentifier}{3}\PYG{+w}{ }\PYG{n+nv+nvIdentifier}{all}\PYG{+w}{ }\PYG{n}{drude}\PYG{o}{/}\PYG{n}{transform}\PYG{o}{/}\PYG{n}{direct}
\PYG{k}{fix }\PYG{n+nv+nvIdentifier}{1}\PYG{+w}{ }\PYG{n+nv+nvIdentifier}{all}\PYG{+w}{ }\PYG{n}{drude}\PYG{o}{/}\PYG{n}{transform}\PYG{o}{/}\PYG{n}{inverse}
\end{sphinxVerbatim}

\sphinxAtStartPar
Example input scripts available: examples/PACKAGES/drude


\subsection{Description}
\label{\detokenize{fix_drude_transform:description}}
\sphinxAtStartPar
Transform the coordinates of Drude oscillators from real to reduced
and back for thermalizing the Drude oscillators as described in
{\hyperref[\detokenize{fix_drude_transform:lamoureux1}]{\sphinxcrossref{\DUrole{std,std-ref}{(Lamoureux)}}}} using a Nose\sphinxhyphen{}Hoover thermostat.  This fix is
designed to be used with the {\hyperref[\detokenize{Howto_drude::doc}]{\sphinxcrossref{\DUrole{doc}{thermalized Drude oscillator model}}}}.  Polarizable models in LAMMPS are described
on the {\hyperref[\detokenize{Howto_polarizable::doc}]{\sphinxcrossref{\DUrole{doc}{Howto polarizable}}}} doc page.

\sphinxAtStartPar
Drude oscillators are a pair of atoms representing a single
polarizable atom.  Ideally, the mass of Drude particles would vanish
and their positions would be determined self\sphinxhyphen{}consistently by iterative
minimization of the energy, the cores’ positions being fixed.  It is
however more efficient and it yields comparable results, if the Drude
oscillators (the motion of the Drude particle relative to the core)
are thermalized at a low temperature.  In that case, the Drude
particles need a small mass.

\sphinxAtStartPar
The thermostats act on the reduced degrees of freedom, which are
defined by the following equations.  Note that in these equations
upper case denotes atomic or center of mass values and lower case
denotes Drude particle or dipole values. Primes denote the transformed
(reduced) values, while bare letters denote the original values.

\sphinxAtStartPar
Masses:
\begin{equation*}
\begin{split}M' = M + m\end{split}
\end{equation*}\begin{equation*}
\begin{split}m' = \frac {M\, m } {M'}\end{split}
\end{equation*}
\sphinxAtStartPar
Positions:
\begin{equation*}
\begin{split}X' = \frac {M\, X + m\, x} {M'}\end{split}
\end{equation*}\begin{equation*}
\begin{split}x' = x - X\end{split}
\end{equation*}
\sphinxAtStartPar
Velocities:
\begin{equation*}
\begin{split}V' = \frac {M\, V + m\, v} {M'}\end{split}
\end{equation*}\begin{equation*}
\begin{split}v' = v - V\end{split}
\end{equation*}
\sphinxAtStartPar
Forces:
\begin{equation*}
\begin{split}F' = F + f\end{split}
\end{equation*}\begin{equation*}
\begin{split}f' = \frac { M\, f - m\, F} {M'}\end{split}
\end{equation*}
\sphinxAtStartPar
This transform conserves the total kinetic energy
\begin{equation*}
\begin{split} \frac 1 2 \, (M\, V^2\ + m\, v^2)
= \frac 1 2 \, (M'\, V'^2\ + m'\, v'^2)\end{split}
\end{equation*}
\sphinxAtStartPar
and the virial defined with absolute positions
\begin{equation*}
\begin{split}X\, F + x\, f = X'\, F' + x'\, f'\end{split}
\end{equation*}

\bigskip\hrule\bigskip


\sphinxAtStartPar
This fix requires each atom know whether it is a Drude particle or
not.  You must therefore use the {\hyperref[\detokenize{fix_drude::doc}]{\sphinxcrossref{\DUrole{doc}{fix drude}}}} command to
specify the Drude status of each atom type.

\begin{sphinxadmonition}{note}{Note:}
\sphinxAtStartPar
only the Drude core atoms need to be in the group specified for
this fix. A Drude electron will be transformed together with its core
even if it is not itself in the group.  It is safe to include Drude
electrons or non\sphinxhyphen{}polarizable atoms in the group. The non\sphinxhyphen{}polarizable
atoms will simply not be transformed.
\end{sphinxadmonition}


\bigskip\hrule\bigskip


\sphinxAtStartPar
This fix does NOT perform time integration. It only transform masses,
coordinates, velocities and forces. Thus you must use separate time
integration fixes, like {\hyperref[\detokenize{fix_nve::doc}]{\sphinxcrossref{\DUrole{doc}{fix nve}}}} or {\hyperref[\detokenize{fix_nh::doc}]{\sphinxcrossref{\DUrole{doc}{fix npt}}}} to actually update the velocities and positions of
atoms.  In order to thermalize the reduced degrees of freedom at
different temperatures, two Nose\sphinxhyphen{}Hoover thermostats must be defined,
acting on two distinct groups.

\begin{sphinxadmonition}{note}{Note:}
\sphinxAtStartPar
The \sphinxstyleemphasis{fix drude/transform/direct} command must appear before any
Nose\sphinxhyphen{}Hoover thermostatting fixes.  The \sphinxstyleemphasis{fix drude/transform/inverse}
command must appear after any Nose\sphinxhyphen{}Hoover thermostatting fixes.
\end{sphinxadmonition}

\sphinxAtStartPar
Example:

\begin{sphinxVerbatim}[commandchars=\\\{\}]
\PYG{k}{fix }\PYG{n+nv+nvIdentifier}{fDIRECT}\PYG{+w}{ }\PYG{n+nv+nvIdentifier}{all}\PYG{+w}{ }\PYG{n}{drude}\PYG{o}{/}\PYG{n}{transform}\PYG{o}{/}\PYG{n}{direct}
\PYG{k}{fix }\PYG{n+nv+nvIdentifier}{fNVT}\PYG{+w}{ }\PYG{n+nv+nvIdentifier}{gCORES}\PYG{+w}{ }\PYG{n}{nvt}\PYG{+w}{ }\PYG{n}{temp}\PYG{+w}{ }\PYG{l+m}{300.0}\PYG{+w}{ }\PYG{l+m}{300.0}\PYG{+w}{ }\PYG{l+m}{100}
\PYG{k}{fix }\PYG{n+nv+nvIdentifier}{fNVT}\PYG{+w}{ }\PYG{n+nv+nvIdentifier}{gDRUDES}\PYG{+w}{ }\PYG{n}{nvt}\PYG{+w}{ }\PYG{n}{temp}\PYG{+w}{ }\PYG{l+m}{1.0}\PYG{+w}{ }\PYG{l+m}{1.0}\PYG{+w}{ }\PYG{l+m}{100}
\PYG{k}{fix }\PYG{n+nv+nvIdentifier}{fINVERSE}\PYG{+w}{ }\PYG{n+nv+nvIdentifier}{all}\PYG{+w}{ }\PYG{n}{drude}\PYG{o}{/}\PYG{n}{transform}\PYG{o}{/}\PYG{n}{inverse}
\PYG{k}{compute }\PYG{n+nv+nvIdentifier}{TDRUDE}\PYG{+w}{ }\PYG{n+nv+nvIdentifier}{all}\PYG{+w}{ }\PYG{n}{temp}\PYG{o}{/}\PYG{n}{drude}
\PYG{k}{thermo\PYGZus{}style}\PYG{+w}{ }\PYG{n}{custom}\PYG{+w}{ }\PYG{n}{step}\PYG{+w}{ }\PYG{n}{cpu}\PYG{+w}{ }\PYG{n}{etotal}\PYG{+w}{ }\PYG{n}{ke}\PYG{+w}{ }\PYG{n}{pe}\PYG{+w}{ }\PYG{n}{ebond}\PYG{+w}{ }\PYG{n}{ecoul}\PYG{+w}{ }\PYG{n}{elong}\PYG{+w}{ }\PYG{n}{press}\PYG{+w}{ }\PYG{n}{vol}\PYG{+w}{ }\PYG{n}{temp}\PYG{+w}{ }\PYG{n}{c\PYGZus{}TDRUDE[1]}\PYG{+w}{ }\PYG{n}{c\PYGZus{}TDRUDE[2]}
\end{sphinxVerbatim}

\sphinxAtStartPar
In this example, \sphinxstyleemphasis{gCORES} is the group of the atom cores and \sphinxstyleemphasis{gDRUDES}
is the group of the Drude particles (electrons). The centers of mass
of the Drude oscillators will be thermostatted at 300.0 and the
internal degrees of freedom will be thermostatted at 1.0.  The
temperatures of cores and Drude particles, in center\sphinxhyphen{}of\sphinxhyphen{}mass and
relative coordinates, are calculated using {\hyperref[\detokenize{compute_temp_drude::doc}]{\sphinxcrossref{\DUrole{doc}{compute temp/drude}}}}

\sphinxAtStartPar
In addition, if you want to use a barostat to simulate a system at
constant pressure, only one of the Nose\sphinxhyphen{}Hoover fixes must be \sphinxstyleemphasis{npt},
the other one should be \sphinxstyleemphasis{nvt}. You must add a \sphinxstyleemphasis{compute temp/com} and a
\sphinxstyleemphasis{fix\_modify} command so that the temperature of the \sphinxstyleemphasis{npt} fix be just
that of its group (the Drude cores) but the pressure be the overall
pressure \sphinxstyleemphasis{thermo\_press}.

\sphinxAtStartPar
Example:

\begin{sphinxVerbatim}[commandchars=\\\{\}]
\PYG{k}{compute }\PYG{n+nv+nvIdentifier}{cTEMP\PYGZus{}CORE}\PYG{+w}{ }\PYG{n+nv+nvIdentifier}{gCORES}\PYG{+w}{ }\PYG{n}{temp}\PYG{o}{/}\PYG{n}{com}
\PYG{k}{fix }\PYG{n+nv+nvIdentifier}{fDIRECT}\PYG{+w}{ }\PYG{n+nv+nvIdentifier}{all}\PYG{+w}{ }\PYG{n}{drude}\PYG{o}{/}\PYG{n}{transform}\PYG{o}{/}\PYG{n}{direct}
\PYG{k}{fix }\PYG{n+nv+nvIdentifier}{fNPT}\PYG{+w}{ }\PYG{n+nv+nvIdentifier}{gCORES}\PYG{+w}{ }\PYG{n}{npt}\PYG{+w}{ }\PYG{n}{temp}\PYG{+w}{ }\PYG{l+m}{298.0}\PYG{+w}{ }\PYG{l+m}{298.0}\PYG{+w}{ }\PYG{l+m}{100}\PYG{+w}{ }\PYG{n}{iso}\PYG{+w}{ }\PYG{l+m}{1.0}\PYG{+w}{ }\PYG{l+m}{1.0}\PYG{+w}{ }\PYG{l+m}{500}
\PYG{k}{fix\PYGZus{}modify }\PYG{n+nv+nvIdentifier}{fNPT}\PYG{+w}{ }\PYG{n}{temp}\PYG{+w}{ }\PYG{n}{cTEMP\PYGZus{}CORE}\PYG{+w}{ }\PYG{n}{press}\PYG{+w}{ }\PYG{n}{thermo\PYGZus{}press}
\PYG{k}{fix }\PYG{n+nv+nvIdentifier}{fNVT}\PYG{+w}{ }\PYG{n+nv+nvIdentifier}{gDRUDES}\PYG{+w}{ }\PYG{n}{nvt}\PYG{+w}{ }\PYG{n}{temp}\PYG{+w}{ }\PYG{l+m}{5.0}\PYG{+w}{ }\PYG{l+m}{5.0}\PYG{+w}{ }\PYG{l+m}{100}
\PYG{k}{fix }\PYG{n+nv+nvIdentifier}{fINVERSE}\PYG{+w}{ }\PYG{n+nv+nvIdentifier}{all}\PYG{+w}{ }\PYG{n}{drude}\PYG{o}{/}\PYG{n}{transform}\PYG{o}{/}\PYG{n}{inverse}
\end{sphinxVerbatim}

\sphinxAtStartPar
In this example, \sphinxstyleemphasis{gCORES} is the group of the atom cores and \sphinxstyleemphasis{gDRUDES}
is the group of the Drude particles. The centers of mass of the Drude
oscillators will be thermostatted at 298.0 and the internal degrees of
freedom will be thermostatted at 5.0. The whole system will be
barostatted at 1.0.

\sphinxAtStartPar
In order to avoid the flying ice cube problem (irreversible transfer
of linear momentum to the center of mass of the system), you may need
to add a \sphinxstyleemphasis{fix momentum} command:

\begin{sphinxVerbatim}[commandchars=\\\{\}]
\PYG{k}{fix }\PYG{n+nv+nvIdentifier}{fMOMENTUM}\PYG{+w}{ }\PYG{n+nv+nvIdentifier}{all}\PYG{+w}{ }\PYG{n}{momentum}\PYG{+w}{ }\PYG{l+m}{100}\PYG{+w}{ }\PYG{n}{linear}\PYG{+w}{ }\PYG{l+m}{1}\PYG{+w}{ }\PYG{l+m}{1}\PYG{+w}{ }\PYG{l+m}{1}
\end{sphinxVerbatim}


\bigskip\hrule\bigskip



\subsection{Restart, fix\_modify, output, run start/stop, minimize info}
\label{\detokenize{fix_drude_transform:restart-fix-modify-output-run-start-stop-minimize-info}}
\sphinxAtStartPar
No information about this fix is written to {\hyperref[\detokenize{restart::doc}]{\sphinxcrossref{\DUrole{doc}{binary restart files}}}}.


\subsection{Restrictions}
\label{\detokenize{fix_drude_transform:restrictions}}\begin{quote}

\sphinxAtStartPar
none
\end{quote}


\subsection{Related commands}
\label{\detokenize{fix_drude_transform:related-commands}}
\sphinxAtStartPar
{\hyperref[\detokenize{fix_drude::doc}]{\sphinxcrossref{\DUrole{doc}{fix drude}}}},
{\hyperref[\detokenize{fix_langevin_drude::doc}]{\sphinxcrossref{\DUrole{doc}{fix langevin/drude}}}},
{\hyperref[\detokenize{compute_temp_drude::doc}]{\sphinxcrossref{\DUrole{doc}{compute temp/drude}}}},
{\hyperref[\detokenize{pair_thole::doc}]{\sphinxcrossref{\DUrole{doc}{pair\_style thole}}}}


\subsection{Default}
\label{\detokenize{fix_drude_transform:default}}
\sphinxAtStartPar
none


\bigskip\hrule\bigskip

\phantomsection\label{\detokenize{fix_drude_transform:lamoureux1}}
\sphinxAtStartPar
\sphinxstylestrong{(Lamoureux)} Lamoureux and Roux, J Chem Phys, 119, 3025\sphinxhyphen{}3039 (2003).

\sphinxstepscope

\index{fix dt/reset@\spxentry{fix dt/reset}}\ignorespaces 
\index{fix dt/reset/kk@\spxentry{fix dt/reset/kk}}\ignorespaces 

\section{fix dt/reset command}
\label{\detokenize{fix_dt_reset:fix-dt-reset-command}}\label{\detokenize{fix_dt_reset:index-1}}\label{\detokenize{fix_dt_reset:index-0}}\label{\detokenize{fix_dt_reset::doc}}
\sphinxAtStartPar
Accelerator Variants: \sphinxstyleemphasis{dt/reset/kk}


\subsection{Syntax}
\label{\detokenize{fix_dt_reset:syntax}}
\begin{sphinxVerbatim}[commandchars=\\\{\}]
\PYG{k}{fix }\PYG{n+nv+nvIdentifier}{ID}\PYG{+w}{ }\PYG{n+nv+nvIdentifier}{group\PYGZhy{}ID}\PYG{+w}{ }\PYG{n}{dt}\PYG{o}{/}\PYG{n}{reset}\PYG{+w}{ }\PYG{n}{N}\PYG{+w}{ }\PYG{n}{Tmin}\PYG{+w}{ }\PYG{n}{Tmax}\PYG{+w}{ }\PYG{n}{Xmax}\PYG{+w}{ }\PYG{n}{keyword}\PYG{+w}{ }\PYG{n}{values}\PYG{+w}{ }\PYG{n}{...}
\end{sphinxVerbatim}
\begin{itemize}
\item {} 
\sphinxAtStartPar
ID, group\sphinxhyphen{}ID are documented in {\hyperref[\detokenize{fix::doc}]{\sphinxcrossref{\DUrole{doc}{fix}}}} command

\item {} 
\sphinxAtStartPar
dt/reset = style name of this fix command

\item {} 
\sphinxAtStartPar
N = re\sphinxhyphen{}compute dt every N timesteps

\item {} 
\sphinxAtStartPar
Tmin = minimum dt allowed which can be NULL (time units)

\item {} 
\sphinxAtStartPar
Tmax = maximum dt allowed which can be NULL (time units)

\item {} 
\sphinxAtStartPar
Xmax = maximum distance for an atom to move in one timestep (distance units)

\item {} 
\sphinxAtStartPar
zero or more keyword/value pairs may be appended

\item {} 
\sphinxAtStartPar
keyword = \sphinxstyleemphasis{emax} or \sphinxstyleemphasis{units}

\end{itemize}
\begin{sphinxalltt}
\sphinxstyleemphasis{emax} value = Emax
  Emax = maximum kinetic energy change for an atom in one timestep (energy units)
\sphinxstyleemphasis{units} value = \sphinxstyleemphasis{lattice} or \sphinxstyleemphasis{box}
  lattice = Xmax is defined in lattice units
  box = Xmax is defined in simulation box units
\end{sphinxalltt}


\subsection{Examples}
\label{\detokenize{fix_dt_reset:examples}}
\begin{sphinxVerbatim}[commandchars=\\\{\}]
\PYG{k}{fix }\PYG{n+nv+nvIdentifier}{5}\PYG{+w}{ }\PYG{n+nv+nvIdentifier}{all}\PYG{+w}{ }\PYG{n}{dt}\PYG{o}{/}\PYG{n}{reset}\PYG{+w}{ }\PYG{l+m}{10}\PYG{+w}{ }\PYG{l+m}{1.0e\PYGZhy{}5}\PYG{+w}{ }\PYG{l+m}{0.01}\PYG{+w}{ }\PYG{l+m}{0.1}
\PYG{k}{fix }\PYG{n+nv+nvIdentifier}{5}\PYG{+w}{ }\PYG{n+nv+nvIdentifier}{all}\PYG{+w}{ }\PYG{n}{dt}\PYG{o}{/}\PYG{n}{reset}\PYG{+w}{ }\PYG{l+m}{10}\PYG{+w}{ }\PYG{l+m}{0.01}\PYG{+w}{ }\PYG{l+m}{2.0}\PYG{+w}{ }\PYG{l+m}{0.2}\PYG{+w}{ }\PYG{n}{units}\PYG{+w}{ }\PYG{n}{box}
\PYG{k}{fix }\PYG{n+nv+nvIdentifier}{5}\PYG{+w}{ }\PYG{n+nv+nvIdentifier}{all}\PYG{+w}{ }\PYG{n}{dt}\PYG{o}{/}\PYG{n}{reset}\PYG{+w}{ }\PYG{l+m}{5}\PYG{+w}{ }\PYG{n}{NULL}\PYG{+w}{ }\PYG{l+m}{0.001}\PYG{+w}{ }\PYG{l+m}{0.5}\PYG{+w}{ }\PYG{n}{emax}\PYG{+w}{ }\PYG{l+m}{30}\PYG{+w}{ }\PYG{n}{units}\PYG{+w}{ }\PYG{n}{box}
\end{sphinxVerbatim}


\subsection{Description}
\label{\detokenize{fix_dt_reset:description}}
\sphinxAtStartPar
Reset the timestep size every N steps during a run, so that no atom
moves further than the specified \sphinxstyleemphasis{Xmax} distance, based on current
atom velocities and forces.  Optionally an additional criterion is
imposed by the \sphinxstyleemphasis{emax} keyword, so that no atom’s kinetic energy
changes by more than the specified \sphinxstyleemphasis{Emax}.

\sphinxAtStartPar
This can be useful when starting from a configuration with overlapping
atoms, where forces will be large.  Or it can be useful when running
an impact simulation where one or more high\sphinxhyphen{}energy atoms collide with
a solid, causing a damage cascade.

\sphinxAtStartPar
This fix overrides the timestep size setting made by the
{\hyperref[\detokenize{timestep::doc}]{\sphinxcrossref{\DUrole{doc}{timestep}}}} command.  The new timestep size \sphinxstyleemphasis{dt} is
computed in the following manner.

\sphinxAtStartPar
For each atom, the timestep is computed that would cause it to
displace \sphinxstyleemphasis{Xmax} on the next integration step, as a function of its
current velocity and force.  Since performing this calculation exactly
would require the solution to a quartic equation, a cheaper estimate
is generated.  The estimate is conservative in that the atom’s
displacement is guaranteed not to exceed \sphinxstyleemphasis{Xmax}, though it may be
smaller.

\sphinxAtStartPar
In addition if the \sphinxstyleemphasis{emax} keyword is used, the specified \sphinxstyleemphasis{Emax} value
is enforced as a limit on how much an atom’s kinetic energy can
change.  If the timestep required is even smaller than for the \sphinxstyleemphasis{Xmax}
displacement, then the smaller timestep is used.

\sphinxAtStartPar
Given this putative timestep for each atom, the minimum timestep value
across all atoms is computed.  Then the \sphinxstyleemphasis{Tmin} and \sphinxstyleemphasis{Tmax} bounds are
applied, if specified.  If one (or both) is specified as NULL, it is
not applied.

\sphinxAtStartPar
When the {\hyperref[\detokenize{run_style::doc}]{\sphinxcrossref{\DUrole{doc}{run style}}}} is \sphinxstyleemphasis{respa}, this fix resets the
outer loop (largest) timestep, which is the same timestep that the
{\hyperref[\detokenize{timestep::doc}]{\sphinxcrossref{\DUrole{doc}{timestep}}}} command sets.

\sphinxAtStartPar
Note that the cumulative simulation time (in time units), which
accounts for changes in the timestep size as a simulation proceeds,
can be accessed by the {\hyperref[\detokenize{thermo_style::doc}]{\sphinxcrossref{\DUrole{doc}{thermo\_style time}}}}
keyword.

\sphinxAtStartPar
Also note that the {\hyperref[\detokenize{dump_modify::doc}]{\sphinxcrossref{\DUrole{doc}{dump\_modify every/time}}}} option
allows dump files to be written at intervals specified by simulation
time, rather than by timesteps.  Simulation time is in time units;
see the {\hyperref[\detokenize{units::doc}]{\sphinxcrossref{\DUrole{doc}{units}}}} doc page for details.


\bigskip\hrule\bigskip


\sphinxAtStartPar
Styles with a \sphinxstyleemphasis{gpu}, \sphinxstyleemphasis{intel}, \sphinxstyleemphasis{kk}, \sphinxstyleemphasis{omp}, or \sphinxstyleemphasis{opt} suffix are
functionally the same as the corresponding style without the suffix.
They have been optimized to run faster, depending on your available
hardware, as discussed on the {\hyperref[\detokenize{Speed_packages::doc}]{\sphinxcrossref{\DUrole{doc}{Accelerator packages}}}}
page.  The accelerated styles take the same arguments and should
produce the same results, except for round\sphinxhyphen{}off and precision issues.

\sphinxAtStartPar
These accelerated styles are part of the GPU, INTEL, KOKKOS,
OPENMP, and OPT packages, respectively.  They are only enabled if
LAMMPS was built with those packages.  See the {\hyperref[\detokenize{Build_package::doc}]{\sphinxcrossref{\DUrole{doc}{Build package}}}} page for more info.

\sphinxAtStartPar
You can specify the accelerated styles explicitly in your input script
by including their suffix, or you can use the {\hyperref[\detokenize{Run_options::doc}]{\sphinxcrossref{\DUrole{doc}{\sphinxhyphen{}suffix command\sphinxhyphen{}line switch}}}} when you invoke LAMMPS, or you can use the
{\hyperref[\detokenize{suffix::doc}]{\sphinxcrossref{\DUrole{doc}{suffix}}}} command in your input script.

\sphinxAtStartPar
See the {\hyperref[\detokenize{Speed_packages::doc}]{\sphinxcrossref{\DUrole{doc}{Accelerator packages}}}} page for more
instructions on how to use the accelerated styles effectively.


\subsection{Restart, fix\_modify, output, run start/stop, minimize info}
\label{\detokenize{fix_dt_reset:restart-fix-modify-output-run-start-stop-minimize-info}}
\sphinxAtStartPar
No information about this fix is written to {\hyperref[\detokenize{restart::doc}]{\sphinxcrossref{\DUrole{doc}{binary restart files}}}}.  None of the {\hyperref[\detokenize{fix_modify::doc}]{\sphinxcrossref{\DUrole{doc}{fix\_modify}}}} options are
relevant to this fix.

\sphinxAtStartPar
This fix computes a global scalar which can be accessed by various
{\hyperref[\detokenize{Howto_output::doc}]{\sphinxcrossref{\DUrole{doc}{output commands}}}}.  The scalar stores the last
timestep on which the timestep was reset to a new value.

\sphinxAtStartPar
The scalar value calculated by this fix is “intensive”.

\sphinxAtStartPar
No parameter of this fix can be used with the \sphinxstyleemphasis{start/stop} keywords of
the {\hyperref[\detokenize{run::doc}]{\sphinxcrossref{\DUrole{doc}{run}}}} command.  This fix is not invoked during
{\hyperref[\detokenize{minimize::doc}]{\sphinxcrossref{\DUrole{doc}{energy minimization}}}}.


\subsection{Restrictions}
\label{\detokenize{fix_dt_reset:restrictions}}\begin{quote}

\sphinxAtStartPar
none
\end{quote}


\subsection{Related commands}
\label{\detokenize{fix_dt_reset:related-commands}}
\sphinxAtStartPar
{\hyperref[\detokenize{timestep::doc}]{\sphinxcrossref{\DUrole{doc}{timestep}}}}, {\hyperref[\detokenize{dump_modify::doc}]{\sphinxcrossref{\DUrole{doc}{dump\_modify every/time}}}}


\subsection{Default}
\label{\detokenize{fix_dt_reset:default}}
\sphinxAtStartPar
The option defaults are units = lattice, and no emax kinetic energy
limit.

\sphinxstepscope

\index{fix efield@\spxentry{fix efield}}\ignorespaces 
\index{fix efield/kk@\spxentry{fix efield/kk}}\ignorespaces 
\index{fix efield/tip4p@\spxentry{fix efield/tip4p}}\ignorespaces 

\section{fix efield command}
\label{\detokenize{fix_efield:fix-efield-command}}\label{\detokenize{fix_efield:index-2}}\label{\detokenize{fix_efield:index-1}}\label{\detokenize{fix_efield:index-0}}\label{\detokenize{fix_efield::doc}}

\section{fix efield/tip4p command}
\label{\detokenize{fix_efield:fix-efield-tip4p-command}}

\subsection{Syntax}
\label{\detokenize{fix_efield:syntax}}
\begin{sphinxVerbatim}[commandchars=\\\{\}]
\PYG{k}{fix }\PYG{n+nv+nvIdentifier}{ID}\PYG{+w}{ }\PYG{n+nv+nvIdentifier}{group\PYGZhy{}ID}\PYG{+w}{ }\PYG{n}{style}\PYG{+w}{ }\PYG{n}{ex}\PYG{+w}{ }\PYG{n}{ey}\PYG{+w}{ }\PYG{n}{ez}\PYG{+w}{ }\PYG{n}{keyword}\PYG{+w}{ }\PYG{n}{value}\PYG{+w}{ }\PYG{n}{...}
\end{sphinxVerbatim}
\begin{itemize}
\item {} 
\sphinxAtStartPar
ID, group\sphinxhyphen{}ID are documented in {\hyperref[\detokenize{fix::doc}]{\sphinxcrossref{\DUrole{doc}{fix}}}} command

\item {} 
\sphinxAtStartPar
style = \sphinxstyleemphasis{efield} or \sphinxstyleemphasis{efield/tip4p}

\item {} 
\sphinxAtStartPar
ex,ey,ez = E\sphinxhyphen{}field component values (electric field units)

\item {} 
\sphinxAtStartPar
any of ex,ey,ez can be a variable (see below)

\item {} 
\sphinxAtStartPar
zero or more keyword/value pairs may be appended to args

\item {} 
\sphinxAtStartPar
keyword = \sphinxstyleemphasis{region} or \sphinxstyleemphasis{energy} or \sphinxstyleemphasis{potential}
\begin{sphinxalltt}
\sphinxstyleemphasis{region} value = region\sphinxhyphen{}ID
  region\sphinxhyphen{}ID = ID of region atoms must be in to have added force
\sphinxstyleemphasis{energy} value = v\_name
  v\_name = variable with name that calculates the potential energy of each atom in the added E\sphinxhyphen{}field
\sphinxstyleemphasis{potential} value = v\_name
  v\_name = variable with name that calculates the electric potential of each atom in the added E\sphinxhyphen{}field
\end{sphinxalltt}

\end{itemize}


\subsection{Examples}
\label{\detokenize{fix_efield:examples}}
\begin{sphinxVerbatim}[commandchars=\\\{\}]
\PYG{k}{fix }\PYG{n+nv+nvIdentifier}{kick}\PYG{+w}{ }\PYG{n+nv+nvIdentifier}{external\PYGZhy{}field}\PYG{+w}{ }\PYG{n}{efield}\PYG{+w}{ }\PYG{l+m}{1.0}\PYG{+w}{ }\PYG{l+m}{0.0}\PYG{+w}{ }\PYG{l+m}{0.0}
\PYG{k}{fix }\PYG{n+nv+nvIdentifier}{kick}\PYG{+w}{ }\PYG{n+nv+nvIdentifier}{external\PYGZhy{}field}\PYG{+w}{ }\PYG{n}{efield}\PYG{+w}{ }\PYG{l+m}{0.0}\PYG{+w}{ }\PYG{l+m}{0.0}\PYG{+w}{ }\PYG{n}{v\PYGZus{}oscillate}
\PYG{k}{fix }\PYG{n+nv+nvIdentifier}{kick}\PYG{+w}{ }\PYG{n+nv+nvIdentifier}{external\PYGZhy{}field}\PYG{+w}{ }\PYG{n}{efield}\PYG{o}{/}\PYG{n}{tip4p}\PYG{+w}{ }\PYG{l+m}{1.0}\PYG{+w}{ }\PYG{l+m}{0.0}\PYG{+w}{ }\PYG{l+m}{0.0}
\end{sphinxVerbatim}


\subsection{Description}
\label{\detokenize{fix_efield:description}}
\sphinxAtStartPar
Add a force \(\vec{F} = q\vec{E}\) to each charged atom in the group due to an
external electric field being applied to the system.  If the system
contains point\sphinxhyphen{}dipoles, also add a torque on the dipoles due to the
external electric field.

\sphinxAtStartPar
\DUrole{versionmodified,added}{Added in version 28Mar2023.}

\sphinxAtStartPar
When the \sphinxstyleemphasis{efield/tip4p} style is used, the E\sphinxhyphen{}field will be applied to
the position of the virtual charge site M of a TIP4P molecule instead of
the oxygen position as it is defined by a corresponding {\hyperref[\detokenize{pair_lj_cut_tip4p::doc}]{\sphinxcrossref{\DUrole{doc}{TIP4P pair
style}}}}.  The forces on the M site due to the
external field are projected on the oxygen and hydrogen atoms of the
TIP4P molecules.

\sphinxAtStartPar
For charges, any of the 3 quantities defining the E\sphinxhyphen{}field components can
be specified as an equal\sphinxhyphen{}style or atom\sphinxhyphen{}style {\hyperref[\detokenize{variable::doc}]{\sphinxcrossref{\DUrole{doc}{variable}}}},
namely \sphinxstyleemphasis{ex}, \sphinxstyleemphasis{ey}, \sphinxstyleemphasis{ez}.  If the value is a variable, it should be
specified as v\_name, where name is the variable name.  In this case, the
variable will be evaluated each timestep, and its value used to
determine the E\sphinxhyphen{}field component.

\sphinxAtStartPar
For point\sphinxhyphen{}dipoles, equal\sphinxhyphen{}style variables can be used, but atom\sphinxhyphen{}style
variables are not currently supported, since they imply a spatial
gradient in the electric field which means additional terms with
gradients of the field are required for the force and torque on dipoles.

\sphinxAtStartPar
Equal\sphinxhyphen{}style variables can specify formulas with various mathematical
functions, and include {\hyperref[\detokenize{thermo_style::doc}]{\sphinxcrossref{\DUrole{doc}{thermo\_style}}}} command
keywords for the simulation box parameters and timestep and elapsed
time.  Thus it is easy to specify a time\sphinxhyphen{}dependent E\sphinxhyphen{}field.

\sphinxAtStartPar
Atom\sphinxhyphen{}style variables can specify the same formulas as equal\sphinxhyphen{}style
variables but can also include per\sphinxhyphen{}atom values, such as atom
coordinates.  Thus it is easy to specify a spatially\sphinxhyphen{}dependent E\sphinxhyphen{}field
with optional time\sphinxhyphen{}dependence as well.

\sphinxAtStartPar
If the \sphinxstyleemphasis{region} keyword is used, the atom must also be in the
specified geometric {\hyperref[\detokenize{region::doc}]{\sphinxcrossref{\DUrole{doc}{region}}}} in order to have force added
to it.


\bigskip\hrule\bigskip


\sphinxAtStartPar
Adding a force or torque to atoms implies a change in their potential
energy as they move or rotate due to the applied E\sphinxhyphen{}field.

\sphinxAtStartPar
For dynamics via the “run” command, this energy can be optionally
added to the system’s potential energy for thermodynamic output (see
below).  For energy minimization via the “minimize” command, this
energy must be added to the system’s potential energy to formulate a
self\sphinxhyphen{}consistent minimization problem (see below).

\sphinxAtStartPar
The \sphinxstyleemphasis{energy} keyword is not allowed if the added field is a constant
vector (ex,ey,ez), with all components defined as numeric constants
and not as variables.  This is because LAMMPS can compute the energy
for each charged particle directly as
\begin{equation*}
\begin{split}U_{efield} = -\vec{x} \cdot q\vec{E} = -q (x\cdot E_x + y\cdot E_y + z\cdot Ez),\end{split}
\end{equation*}
\sphinxAtStartPar
so that \(-\nabla U_{efield} = \vec{F}\).  Similarly for point\sphinxhyphen{}dipole particles
the energy can be computed as
\begin{equation*}
\begin{split}U_{efield} = -\vec{\mu} \cdot \vec{E} = -\mu_x\cdot E_x + \mu_y\cdot E_y + \mu_z\cdot E_z\end{split}
\end{equation*}
\sphinxAtStartPar
The \sphinxstyleemphasis{energy} keyword is optional if the added force is defined with
one or more variables, and if you are performing dynamics via the
{\hyperref[\detokenize{run::doc}]{\sphinxcrossref{\DUrole{doc}{run}}}} command.  If the keyword is not used, LAMMPS will set
the energy to 0.0, which is typically fine for dynamics.

\sphinxAtStartPar
The \sphinxstyleemphasis{energy} keyword (or \sphinxstyleemphasis{potential} keyword, described below)
is required if the added force is defined with
one or more variables, and you are performing energy minimization via
the “minimize” command for charged particles.  It is not required for
point\sphinxhyphen{}dipoles, but a warning is issued since the minimizer in LAMMPS
does not rotate dipoles, so you should not expect to be able to
minimize the orientation of dipoles in an applied electric field.

\sphinxAtStartPar
The \sphinxstyleemphasis{energy} keyword specifies the name of an atom\sphinxhyphen{}style
{\hyperref[\detokenize{variable::doc}]{\sphinxcrossref{\DUrole{doc}{variable}}}} which is used to compute the energy of each
atom as function of its position.  Like variables used for \sphinxstyleemphasis{ex},
\sphinxstyleemphasis{ey}, \sphinxstyleemphasis{ez}, the energy variable is specified as “v\_name”, where “name”
is the variable name.

\sphinxAtStartPar
Note that when the \sphinxstyleemphasis{energy} keyword is used during an energy
minimization, you must ensure that the formula defined for the
atom\sphinxhyphen{}style {\hyperref[\detokenize{variable::doc}]{\sphinxcrossref{\DUrole{doc}{variable}}}} is consistent with the force
variable formulas, i.e. that \sphinxhyphen{}Grad(E) = F.  For example, if the force
due to the electric field were a spring\sphinxhyphen{}like F = kx, then the energy
formula should be E = \sphinxhyphen{}0.5kx\textasciicircum{}2.  If you don’t do this correctly, the
minimization will not converge properly.

\sphinxAtStartPar
\DUrole{versionmodified,added}{Added in version 15Jun2023.}

\sphinxAtStartPar
The \sphinxstyleemphasis{potential} keyword can be used as an alternative to the \sphinxstyleemphasis{energy} keyword
to specify the name of an atom\sphinxhyphen{}style variable, which is used to compute the
added electric potential to each atom as a function of its position.  The
variable should have units of electric field multiplied by distance (that is,
in \sphinxtitleref{units real}, the potential should be in volts). As with the \sphinxstyleemphasis{energy}
keyword, the variable name is specified as “v\_name”. The energy added by this
fix is then calculated as the electric potential multiplied by charge.

\sphinxAtStartPar
The \sphinxstyleemphasis{potential} keyword is mainly intended for correct charge
equilibration in simulations with {\hyperref[\detokenize{fix_qeq_reaxff::doc}]{\sphinxcrossref{\DUrole{doc}{fix qeq/reaxff}}}},
since with variable charges the electric potential can be known
beforehand but the energy cannot.  A small additional benefit is that
the \sphinxstyleemphasis{energy} keyword requires an additional conversion to energy units
which the \sphinxstyleemphasis{potential} keyword avoids.  Thus, when the \sphinxstyleemphasis{potential}
keyword is specified, the \sphinxstyleemphasis{energy} keyword must not be used.  As with
\sphinxstyleemphasis{energy}, the \sphinxstyleemphasis{potential} keyword is not allowed if the added field is a
constant vector.  The \sphinxstyleemphasis{potential} keyword is not supported by \sphinxstyleemphasis{fix
efield/tip4p}.


\bigskip\hrule\bigskip



\subsection{Restart, fix\_modify, output, run start/stop, minimize info}
\label{\detokenize{fix_efield:restart-fix-modify-output-run-start-stop-minimize-info}}
\sphinxAtStartPar
No information about this fix is written to {\hyperref[\detokenize{restart::doc}]{\sphinxcrossref{\DUrole{doc}{binary restart files}}}}.

\sphinxAtStartPar
The {\hyperref[\detokenize{fix_modify::doc}]{\sphinxcrossref{\DUrole{doc}{fix\_modify}}}} \sphinxstyleemphasis{energy} option is supported by this
fix to add the potential energy inferred by the added force due to the
electric field to the global potential energy of the system as part of
{\hyperref[\detokenize{thermo_style::doc}]{\sphinxcrossref{\DUrole{doc}{thermodynamic output}}}}.  The default setting for
this fix is {\hyperref[\detokenize{fix_modify::doc}]{\sphinxcrossref{\DUrole{doc}{fix\_modify energy no}}}}.  Note that this
energy is a fictitious quantity but is needed so that the {\hyperref[\detokenize{minimize::doc}]{\sphinxcrossref{\DUrole{doc}{minimize}}}} command can include the forces added by this fix in a
consistent manner.  I.e. there is a decrease in potential energy when
atoms move in the direction of the added force due to the electric
field.

\sphinxAtStartPar
The {\hyperref[\detokenize{fix_modify::doc}]{\sphinxcrossref{\DUrole{doc}{fix\_modify}}}} \sphinxstyleemphasis{virial} option is supported by this
fix to add the contribution due to the added forces on atoms to both the
global pressure and per\sphinxhyphen{}atom stress of the system via the {\hyperref[\detokenize{compute_pressure::doc}]{\sphinxcrossref{\DUrole{doc}{compute
pressure}}}} and {\hyperref[\detokenize{compute_stress_atom::doc}]{\sphinxcrossref{\DUrole{doc}{compute stress/atom}}}} commands.  The former can be accessed by
{\hyperref[\detokenize{thermo_style::doc}]{\sphinxcrossref{\DUrole{doc}{thermodynamic output}}}}.  The default setting for
this fix is {\hyperref[\detokenize{fix_modify::doc}]{\sphinxcrossref{\DUrole{doc}{fix\_modify virial no}}}}.

\sphinxAtStartPar
The {\hyperref[\detokenize{fix_modify::doc}]{\sphinxcrossref{\DUrole{doc}{fix\_modify}}}} \sphinxstyleemphasis{respa} option is supported by this
fix. This allows to set at which level of the {\hyperref[\detokenize{run_style::doc}]{\sphinxcrossref{\DUrole{doc}{r\sphinxhyphen{}RESPA}}}}
integrator the fix adding its forces. Default is the outermost level.

\sphinxAtStartPar
This fix computes a global scalar and a global 3\sphinxhyphen{}vector of forces,
which can be accessed by various {\hyperref[\detokenize{Howto_output::doc}]{\sphinxcrossref{\DUrole{doc}{output commands}}}}.  The scalar is the potential energy discussed above.
The vector is the total force added to the group of atoms.  The scalar
and vector values calculated by this fix are “extensive”.

\sphinxAtStartPar
No parameter of this fix can be used with the \sphinxstyleemphasis{start/stop} keywords of
the {\hyperref[\detokenize{run::doc}]{\sphinxcrossref{\DUrole{doc}{run}}}} command.

\sphinxAtStartPar
The forces due to this fix are imposed during an energy minimization,
invoked by the {\hyperref[\detokenize{minimize::doc}]{\sphinxcrossref{\DUrole{doc}{minimize}}}} command.  You should not
specify force components with a variable that has time\sphinxhyphen{}dependence for
use with a minimizer, since the minimizer increments the timestep as
the iteration count during the minimization.

\begin{sphinxadmonition}{note}{Note:}
\sphinxAtStartPar
If you want the fictitious potential energy associated with the
added forces to be included in the total potential energy of the
system (the quantity being minimized), you MUST enable the
{\hyperref[\detokenize{fix_modify::doc}]{\sphinxcrossref{\DUrole{doc}{fix\_modify}}}} \sphinxstyleemphasis{energy} option for this fix.
\end{sphinxadmonition}


\bigskip\hrule\bigskip


\sphinxAtStartPar
Styles with a \sphinxstyleemphasis{gpu}, \sphinxstyleemphasis{intel}, \sphinxstyleemphasis{kk}, \sphinxstyleemphasis{omp}, or \sphinxstyleemphasis{opt} suffix are
functionally the same as the corresponding style without the suffix.
They have been optimized to run faster, depending on your available
hardware, as discussed on the {\hyperref[\detokenize{Speed_packages::doc}]{\sphinxcrossref{\DUrole{doc}{Accelerator packages}}}}
page.  The accelerated styles take the same arguments and should
produce the same results, except for round\sphinxhyphen{}off and precision issues.

\sphinxAtStartPar
These accelerated styles are part of the GPU, INTEL, KOKKOS,
OPENMP, and OPT packages, respectively.  They are only enabled if
LAMMPS was built with those packages.  See the {\hyperref[\detokenize{Build_package::doc}]{\sphinxcrossref{\DUrole{doc}{Build package}}}} page for more info.

\sphinxAtStartPar
You can specify the accelerated styles explicitly in your input script
by including their suffix, or you can use the {\hyperref[\detokenize{Run_options::doc}]{\sphinxcrossref{\DUrole{doc}{\sphinxhyphen{}suffix command\sphinxhyphen{}line switch}}}} when you invoke LAMMPS, or you can use the
{\hyperref[\detokenize{suffix::doc}]{\sphinxcrossref{\DUrole{doc}{suffix}}}} command in your input script.

\sphinxAtStartPar
See the {\hyperref[\detokenize{Speed_packages::doc}]{\sphinxcrossref{\DUrole{doc}{Accelerator packages}}}} page for more
instructions on how to use the accelerated styles effectively.


\bigskip\hrule\bigskip



\subsection{Restrictions}
\label{\detokenize{fix_efield:restrictions}}
\sphinxAtStartPar
Fix style \sphinxstyleemphasis{efield/tip4p} is part of the EXTRA\sphinxhyphen{}FIX package. It is only
enabled if LAMMPS was built with that package.  See the {\hyperref[\detokenize{Build_package::doc}]{\sphinxcrossref{\DUrole{doc}{Build
package}}}} page for more info.

\sphinxAtStartPar
Fix style \sphinxstyleemphasis{efield/tip4p} can only be used with tip4p pair styles.


\subsection{Related commands}
\label{\detokenize{fix_efield:related-commands}}
\sphinxAtStartPar
{\hyperref[\detokenize{fix_addforce::doc}]{\sphinxcrossref{\DUrole{doc}{fix addforce}}}}


\subsection{Default}
\label{\detokenize{fix_efield:default}}
\sphinxAtStartPar
none

\sphinxstepscope

\index{fix ehex@\spxentry{fix ehex}}\ignorespaces 

\section{fix ehex command}
\label{\detokenize{fix_ehex:fix-ehex-command}}\label{\detokenize{fix_ehex:index-0}}\label{\detokenize{fix_ehex::doc}}

\subsection{Syntax}
\label{\detokenize{fix_ehex:syntax}}
\begin{sphinxVerbatim}[commandchars=\\\{\}]
\PYG{k}{fix }\PYG{n+nv+nvIdentifier}{ID}\PYG{+w}{ }\PYG{n+nv+nvIdentifier}{group\PYGZhy{}ID}\PYG{+w}{ }\PYG{n}{ehex}\PYG{+w}{ }\PYG{n}{nevery}\PYG{+w}{ }\PYG{n}{F}\PYG{+w}{ }\PYG{n}{keyword}\PYG{+w}{ }\PYG{n}{value}
\end{sphinxVerbatim}
\begin{itemize}
\item {} 
\sphinxAtStartPar
ID, group\sphinxhyphen{}ID are documented in {\hyperref[\detokenize{fix::doc}]{\sphinxcrossref{\DUrole{doc}{fix}}}} command

\item {} 
\sphinxAtStartPar
ehex  = style name of this fix command

\item {} 
\sphinxAtStartPar
nevery = add/subtract heat every this many timesteps

\item {} 
\sphinxAtStartPar
F = energy flux into the reservoir (energy/time units)

\item {} 
\sphinxAtStartPar
zero or more keyword/value pairs may be appended to args

\item {} 
\sphinxAtStartPar
keyword = \sphinxstyleemphasis{region} or \sphinxstyleemphasis{constrain} or \sphinxstyleemphasis{com} or \sphinxstyleemphasis{hex}
\begin{sphinxalltt}
\sphinxstyleemphasis{region} value = region\sphinxhyphen{}ID
  region\sphinxhyphen{}ID = ID of region (reservoir) atoms must be in for added thermostatting force
\sphinxstyleemphasis{constrain} value = none
  apply the constraint algorithm (SHAKE or RATTLE) again at the end of the timestep
\sphinxstyleemphasis{com} value = none
  rescale all sites of a constrained cluster of atom if its COM is in the reservoir
\sphinxstyleemphasis{hex} value = none
  omit the coordinate correction to recover the HEX algorithm
\end{sphinxalltt}

\end{itemize}


\subsection{Examples}
\label{\detokenize{fix_ehex:examples}}
\begin{sphinxVerbatim}[commandchars=\\\{\}]
\PYG{c}{\PYGZsh{} Lennard\PYGZhy{}Jones, from examples/in.ehex.lj}

\PYG{k}{fix }\PYG{n+nv+nvIdentifier}{fnve}\PYG{+w}{ }\PYG{n+nv+nvIdentifier}{all}\PYG{+w}{ }\PYG{n}{nve}
\PYG{c}{\PYGZsh{} specify regions rhot and rcold}
\PYG{n}{...}
\PYG{k}{fix }\PYG{n+nv+nvIdentifier}{fhot}\PYG{+w}{ }\PYG{n+nv+nvIdentifier}{all}\PYG{+w}{ }\PYG{n}{ehex}\PYG{+w}{ }\PYG{l+m}{1}\PYG{+w}{ }\PYG{l+m}{0.15}\PYG{+w}{ }\PYG{k}{region }\PYG{n+nv+nvIdentifier}{rhot}
\PYG{k}{fix }\PYG{n+nv+nvIdentifier}{fcold}\PYG{+w}{ }\PYG{n+nv+nvIdentifier}{all}\PYG{+w}{ }\PYG{n}{ehex}\PYG{+w}{ }\PYG{l+m}{1}\PYG{+w}{ }\PYG{o}{\PYGZhy{}}\PYG{l+m}{0.15}\PYG{+w}{ }\PYG{k}{region }\PYG{n+nv+nvIdentifier}{rcold}

\PYG{c}{\PYGZsh{} SPC/E water, from examples/in.ehex.spce}
\PYG{k}{fix }\PYG{n+nv+nvIdentifier}{fnve}\PYG{+w}{ }\PYG{n+nv+nvIdentifier}{all}\PYG{+w}{ }\PYG{n}{nve}
\PYG{c}{\PYGZsh{} specify regions rhot and rcold}
\PYG{n}{...}
\PYG{k}{fix }\PYG{n+nv+nvIdentifier}{fhot}\PYG{+w}{ }\PYG{n+nv+nvIdentifier}{all}\PYG{+w}{ }\PYG{n}{ehex}\PYG{+w}{ }\PYG{l+m}{1}\PYG{+w}{ }\PYG{l+m}{0.075}\PYG{+w}{ }\PYG{k}{region }\PYG{n+nv+nvIdentifier}{rhot}\PYG{+w}{ }\PYG{n}{constrain}\PYG{+w}{ }\PYG{n}{com}
\PYG{k}{fix }\PYG{n+nv+nvIdentifier}{fcold}\PYG{+w}{ }\PYG{n+nv+nvIdentifier}{all}\PYG{+w}{ }\PYG{n}{ehex}\PYG{+w}{ }\PYG{l+m}{1}\PYG{+w}{ }\PYG{o}{\PYGZhy{}}\PYG{l+m}{0.075}\PYG{+w}{ }\PYG{k}{region }\PYG{n+nv+nvIdentifier}{rcold}\PYG{+w}{ }\PYG{n}{constrain}\PYG{+w}{ }\PYG{n}{com}
\PYG{k}{fix }\PYG{n+nv+nvIdentifier}{frattle}\PYG{+w}{ }\PYG{n+nv+nvIdentifier}{all}\PYG{+w}{ }\PYG{n}{rattle}\PYG{+w}{ }\PYG{l+m}{1e\PYGZhy{}10}\PYG{+w}{ }\PYG{l+m}{400}\PYG{+w}{ }\PYG{l+m}{0}\PYG{+w}{ }\PYG{n}{b}\PYG{+w}{ }\PYG{l+m}{1}\PYG{+w}{ }\PYG{n}{a}\PYG{+w}{ }\PYG{l+m}{1}
\end{sphinxVerbatim}


\subsection{Description}
\label{\detokenize{fix_ehex:description}}
\sphinxAtStartPar
This fix implements the asymmetric version of the enhanced heat
exchange algorithm {\hyperref[\detokenize{fix_ehex:wirnsberger}]{\sphinxcrossref{\DUrole{std,std-ref}{(Wirnsberger)}}}}. The eHEX algorithm is
an extension of the heat exchange algorithm {\hyperref[\detokenize{fix_ehex:ikeshoji}]{\sphinxcrossref{\DUrole{std,std-ref}{(Ikeshoji)}}}} and
adds an additional coordinate integration to account for higher\sphinxhyphen{}order
truncation terms in the operator splitting.  The original HEX
algorithm (implemented as {\hyperref[\detokenize{fix_heat::doc}]{\sphinxcrossref{\DUrole{doc}{fix heat}}}}) is known to
exhibit a slight energy drift limiting the accessible simulation times
to a few nanoseconds.  This issue is greatly improved by the new
algorithm decreasing the energy drift by at least a factor of a
hundred (LJ and SPC/E water) with little computational overhead.

\sphinxAtStartPar
In both algorithms (non\sphinxhyphen{}translational) kinetic energy is constantly
swapped between regions (reservoirs) to impose a heat flux onto the
system.  The equations of motion are therefore modified if a particle
\(i\) is located inside a reservoir \(\Gamma_k\) where \(k>0\).  We
use \(\Gamma_0\) to label those parts of the simulation box which
are not thermostatted.)  The input parameter \sphinxstyleemphasis{region\sphinxhyphen{}ID} of this fix
corresponds to \(k\).  The energy swap is modelled by introducing an
additional thermostatting force to the equations of motion, such that
the time evolution of coordinates and momenta of particle \(i\)
becomes {\hyperref[\detokenize{fix_ehex:wirnsberger}]{\sphinxcrossref{\DUrole{std,std-ref}{(Wirnsberger)}}}}
\begin{equation*}
\begin{split}\dot{\mathbf r}_i &= \mathbf v_i,  \\
\dot{\mathbf v}_i &= \frac{\mathbf f_i}{m_i} + \frac{\mathbf g_i}{m_i}.\end{split}
\end{equation*}
\sphinxAtStartPar
The thermostatting force is given by
\begin{equation*}
\begin{split}\mathbf g_i = \begin{cases}
\frac{m_i}{2}   \frac{ F_{\Gamma_{k(\mathbf r_i)}}}{ K_{\Gamma_{k(\mathbf r_i)}}}
\left(\mathbf v_i -  \mathbf v_{\Gamma_{k(\mathbf r_i)}} \right) &  \mbox{$k(\mathbf r_i)> 0$ (inside a reservoir),} \\
 0                                     &  \mbox{otherwise, }
\end{cases}\end{split}
\end{equation*}
\sphinxAtStartPar
where \(m_i\) is the mass and \(k(\mathbf r_i)\) maps the particle
position to the respective reservoir. The quantity
\(F_{\Gamma_{k(\mathbf r_i)}}\) corresponds to the input parameter
\sphinxstyleemphasis{F}, which is the energy flux into the reservoir. Furthermore,
\(K_{\Gamma_{k(\mathbf r_i)}}\) and \(v_{\Gamma_{k(\mathbf r_i)}}\)
denote the non\sphinxhyphen{}translational kinetic energy and the center of mass
velocity of that reservoir. The thermostatting force does not affect
the center of mass velocities of the individual reservoirs and the
entire simulation box. A derivation of the equations and details on
the numerical implementation with velocity Verlet in LAMMPS can be
found in reference “(Wirnsberger)”\#\_Wirnsberger.

\begin{sphinxadmonition}{note}{Note:}
\sphinxAtStartPar
This fix only integrates the thermostatting force and must be
combined with another integrator, such as {\hyperref[\detokenize{fix_nve::doc}]{\sphinxcrossref{\DUrole{doc}{fix nve}}}}, to
solve the full equations of motion.
\end{sphinxadmonition}

\sphinxAtStartPar
This fix is different from a thermostat such as {\hyperref[\detokenize{fix_nh::doc}]{\sphinxcrossref{\DUrole{doc}{fix nvt}}}}
or {\hyperref[\detokenize{fix_temp_rescale::doc}]{\sphinxcrossref{\DUrole{doc}{fix temp/rescale}}}} in that energy is
added/subtracted continually.  Thus if there is not another mechanism
in place to counterbalance this effect, the entire system will heat or
cool continuously.

\begin{sphinxadmonition}{note}{Note:}
\sphinxAtStartPar
If heat is subtracted from the system too aggressively so that
the group’s kinetic energy would go to zero, then LAMMPS will halt
with an error message. Increasing the value of \sphinxstyleemphasis{nevery} means that
heat is added/subtracted less frequently but in larger portions.  The
resulting temperature profile will therefore be the same.
\end{sphinxadmonition}

\sphinxAtStartPar
This fix will default to {\hyperref[\detokenize{fix_heat::doc}]{\sphinxcrossref{\DUrole{doc}{fix\_heat}}}} (HEX algorithm) if
the keyword \sphinxstyleemphasis{hex} is specified.


\bigskip\hrule\bigskip


\sphinxAtStartPar
\sphinxstylestrong{Compatibility with SHAKE and RATTLE (rigid molecules)}:

\sphinxAtStartPar
This fix is compatible with {\hyperref[\detokenize{fix_shake::doc}]{\sphinxcrossref{\DUrole{doc}{fix shake}}}} and {\hyperref[\detokenize{fix_shake::doc}]{\sphinxcrossref{\DUrole{doc}{fix rattle}}}}. If either of these constraining algorithms is
specified in the input script and the keyword \sphinxstyleemphasis{constrain} is set, the
bond distances will be corrected a second time at the end of the
integration step.  It is recommended to specify the keyword \sphinxstyleemphasis{com} in
addition to the keyword \sphinxstyleemphasis{constrain}. With this option all sites of a
constrained cluster are rescaled, if its center of mass is located
inside the region. Rescaling all sites of a cluster by the same factor
does not introduce any velocity components along fixed bonds. No
rescaling takes place if the center of mass lies outside the region.

\begin{sphinxadmonition}{note}{Note:}
\sphinxAtStartPar
You can only use the keyword \sphinxstyleemphasis{com} along with \sphinxstyleemphasis{constrain}.
\end{sphinxadmonition}

\sphinxAtStartPar
To achieve the highest accuracy it is recommended to use {\hyperref[\detokenize{fix_shake::doc}]{\sphinxcrossref{\DUrole{doc}{fix rattle}}}} with the keywords \sphinxstyleemphasis{constrain} and \sphinxstyleemphasis{com} as
shown in the second example. Only if RATTLE is employed, the velocity
constraints will be satisfied.

\begin{sphinxadmonition}{note}{Note:}
\sphinxAtStartPar
Even if RATTLE is used and the keywords \sphinxstyleemphasis{com} and \sphinxstyleemphasis{constrain}
are both set, the coordinate constraints will not necessarily be
satisfied up to the target precision. The velocity constraints are
satisfied as long as all sites of a cluster are rescaled (keyword
\sphinxstyleemphasis{com}) and the cluster does not span adjacent reservoirs. The current
implementation of the eHEX algorithm introduces a small error in the
bond distances, which goes to zero with order three in the
timestep. For example, in a simulation of SPC/E water with a timestep
of 2 fs the maximum relative error in the bond distances was found to
be on the order of \(10^{-7}\) for relatively large
temperature gradients.  A higher precision can be achieved by
decreasing the timestep.
\end{sphinxadmonition}


\bigskip\hrule\bigskip



\subsection{Restart, fix\_modify, output, run start/stop, minimize info}
\label{\detokenize{fix_ehex:restart-fix-modify-output-run-start-stop-minimize-info}}
\sphinxAtStartPar
No information about this fix is written to {\hyperref[\detokenize{restart::doc}]{\sphinxcrossref{\DUrole{doc}{binary restart files}}}}.  None of the {\hyperref[\detokenize{fix_modify::doc}]{\sphinxcrossref{\DUrole{doc}{fix\_modify}}}} options
are relevant to this fix.

\sphinxAtStartPar
No parameter of this fix can be used with the \sphinxstyleemphasis{start/stop} keywords of
the {\hyperref[\detokenize{run::doc}]{\sphinxcrossref{\DUrole{doc}{run}}}} command.  This fix is not invoked during {\hyperref[\detokenize{minimize::doc}]{\sphinxcrossref{\DUrole{doc}{energy minimization}}}}.


\subsection{Restrictions}
\label{\detokenize{fix_ehex:restrictions}}
\sphinxAtStartPar
This fix is part of the RIGID package.  It is only enabled if LAMMPS
was built with that package.  See the {\hyperref[\detokenize{Build_package::doc}]{\sphinxcrossref{\DUrole{doc}{Build package}}}} page for more info.


\subsection{Related commands}
\label{\detokenize{fix_ehex:related-commands}}
\sphinxAtStartPar
{\hyperref[\detokenize{fix_heat::doc}]{\sphinxcrossref{\DUrole{doc}{fix heat}}}}, {\hyperref[\detokenize{fix_thermal_conductivity::doc}]{\sphinxcrossref{\DUrole{doc}{fix thermal/conductivity}}}}, {\hyperref[\detokenize{compute_temp::doc}]{\sphinxcrossref{\DUrole{doc}{compute temp}}}}, {\hyperref[\detokenize{compute_temp_region::doc}]{\sphinxcrossref{\DUrole{doc}{compute temp/region}}}}


\subsection{Default}
\label{\detokenize{fix_ehex:default}}
\sphinxAtStartPar
none


\bigskip\hrule\bigskip

\phantomsection\label{\detokenize{fix_ehex:ikeshoji}}
\sphinxAtStartPar
\sphinxstylestrong{(Ikeshoji)} Ikeshoji and Hafskjold, Molecular Physics, 81, 251\sphinxhyphen{}261 (1994).

\phantomsection\label{\detokenize{fix_ehex:wirnsberger}}
\sphinxAtStartPar
\sphinxstylestrong{(Wirnsberger)} Wirnsberger, Frenkel, and Dellago, J Chem Phys, 143,
124104 (2015).

\sphinxstepscope

\index{fix electrode/conp@\spxentry{fix electrode/conp}}\ignorespaces 
\index{fix electrode/conq@\spxentry{fix electrode/conq}}\ignorespaces 
\index{fix electrode/thermo@\spxentry{fix electrode/thermo}}\ignorespaces 
\index{fix electrode/conp/intel@\spxentry{fix electrode/conp/intel}}\ignorespaces 
\index{fix electrode/conq/intel@\spxentry{fix electrode/conq/intel}}\ignorespaces 
\index{fix electrode/thermo/intel@\spxentry{fix electrode/thermo/intel}}\ignorespaces 

\section{fix electrode/conp command}
\label{\detokenize{fix_electrode:fix-electrode-conp-command}}\label{\detokenize{fix_electrode:index-5}}\label{\detokenize{fix_electrode:index-4}}\label{\detokenize{fix_electrode:index-3}}\label{\detokenize{fix_electrode:index-2}}\label{\detokenize{fix_electrode:index-1}}\label{\detokenize{fix_electrode:index-0}}\label{\detokenize{fix_electrode::doc}}
\sphinxAtStartPar
Accelerator Variant: \sphinxstyleemphasis{electrode/conp/intel}


\section{fix electrode/conq command}
\label{\detokenize{fix_electrode:fix-electrode-conq-command}}
\sphinxAtStartPar
Accelerator Variant: \sphinxstyleemphasis{electrode/conq/intel}


\section{fix electrode/thermo command}
\label{\detokenize{fix_electrode:fix-electrode-thermo-command}}
\sphinxAtStartPar
Accelerator Variant: \sphinxstyleemphasis{electrode/thermo/intel}


\subsection{Syntax}
\label{\detokenize{fix_electrode:syntax}}
\begin{sphinxVerbatim}[commandchars=\\\{\}]
\PYG{k}{fix }\PYG{n+nv+nvIdentifier}{ID}\PYG{+w}{ }\PYG{n+nv+nvIdentifier}{group\PYGZhy{}ID}\PYG{+w}{ }\PYG{n}{style}\PYG{+w}{ }\PYG{n}{args}\PYG{+w}{ }\PYG{n}{keyword}\PYG{+w}{ }\PYG{n}{value}\PYG{+w}{ }\PYG{n}{...}
\end{sphinxVerbatim}
\begin{itemize}
\item {} 
\sphinxAtStartPar
ID, group\sphinxhyphen{}ID are documented in {\hyperref[\detokenize{fix::doc}]{\sphinxcrossref{\DUrole{doc}{fix}}}} command

\item {} 
\sphinxAtStartPar
style = \sphinxstyleemphasis{electrode/conp} or \sphinxstyleemphasis{electrode/conq} or \sphinxstyleemphasis{electrode/thermo}

\item {} 
\sphinxAtStartPar
args = arguments used by a particular style
\begin{sphinxalltt}
\sphinxstyleemphasis{electrode/conp} args = potential eta
\sphinxstyleemphasis{electrode/conq} args = charge eta
\sphinxstyleemphasis{electrode/thermo} args = potential eta \sphinxstyleemphasis{temp} values
     potential = electrode potential
     charge = electrode charge
     eta = reciprocal width of electrode charge smearing (can be NULL if eta keyword is used)
     \sphinxstyleemphasis{temp} values = T\_v tau\_v rng\_v
         T\_v = temperature of thermo\sphinxhyphen{}potentiostat
         tau\_v = time constant of thermo\sphinxhyphen{}potentiostat
         rng\_v = integer used to initialize random number generator
\end{sphinxalltt}

\item {} 
\sphinxAtStartPar
zero or more keyword/value pairs may be appended

\item {} 
\sphinxAtStartPar
keyword = \sphinxstyleemphasis{algo} or \sphinxstyleemphasis{symm} or \sphinxstyleemphasis{couple} or \sphinxstyleemphasis{etypes} or \sphinxstyleemphasis{ffield} or \sphinxstyleemphasis{write\_mat} or \sphinxstyleemphasis{write\_inv} or \sphinxstyleemphasis{read\_mat} or \sphinxstyleemphasis{read\_inv} or \sphinxstyleemphasis{qtotal} or \sphinxstyleemphasis{eta}

\end{itemize}
\begin{sphinxalltt}
\sphinxstyleemphasis{algo} values = \sphinxstyleemphasis{mat\_inv} or \sphinxstyleemphasis{mat\_cg} tol or \sphinxstyleemphasis{cg} tol
    specify the algorithm used to compute the electrode charges
\sphinxstyleemphasis{symm} value = \sphinxstyleemphasis{on} or \sphinxstyleemphasis{off}
    turn on/off charge neutrality constraint for the electrodes
\sphinxstyleemphasis{couple} values = group\sphinxhyphen{}ID val
    group\sphinxhyphen{}ID = group of atoms treated as additional electrode
    val = electric potential or charge on this electrode
\sphinxstyleemphasis{etypes} value = \sphinxstyleemphasis{on} or \sphinxstyleemphasis{off}
    turn on/off type\sphinxhyphen{}based optimized neighbor lists (electrode and electrolyte types may not overlap)
\sphinxstyleemphasis{ffield} value = \sphinxstyleemphasis{on} or \sphinxstyleemphasis{off}
    turn on/off finite\sphinxhyphen{}field implementation
\sphinxstyleemphasis{write\_mat} value = filename
    filename = file to which to write elastance matrix
\sphinxstyleemphasis{write\_inv} value = filename
    filename = file to which to write inverted matrix
\sphinxstyleemphasis{read\_mat} value = filename
    filename = file from which to read elastance matrix
\sphinxstyleemphasis{read\_inv} value = filename
    filename = file from which to read inverted matrix
\sphinxstyleemphasis{qtotal} value = number or \sphinxstyleemphasis{v\_} equal\sphinxhyphen{}style variable
    add overall potential so that all electrode charges add up to \sphinxstyleemphasis{qtotal}
\sphinxstyleemphasis{eta} value = d\_propname
    d\_propname = a custom double vector defined via fix property/atom
\end{sphinxalltt}


\subsection{Examples}
\label{\detokenize{fix_electrode:examples}}
\begin{sphinxVerbatim}[commandchars=\\\{\}]
\PYG{k}{fix }\PYG{n+nv+nvIdentifier}{fxconp}\PYG{+w}{ }\PYG{n+nv+nvIdentifier}{bot}\PYG{+w}{ }\PYG{n}{electrode}\PYG{o}{/}\PYG{n}{conp}\PYG{+w}{ }\PYG{o}{\PYGZhy{}}\PYG{l+m}{1.0}\PYG{+w}{ }\PYG{l+m}{1.805}\PYG{+w}{ }\PYG{n}{couple}\PYG{+w}{ }\PYG{n}{top}\PYG{+w}{ }\PYG{l+m}{1.0}\PYG{+w}{ }\PYG{n}{couple}\PYG{+w}{ }\PYG{n}{ref}\PYG{+w}{ }\PYG{l+m}{0.0}\PYG{+w}{ }\PYG{n}{write\PYGZus{}inv}\PYG{+w}{ }\PYG{n}{inv.csv}\PYG{+w}{ }\PYG{n}{symm}\PYG{+w}{ }\PYG{n}{on}
\PYG{k}{fix }\PYG{n+nv+nvIdentifier}{fxconp}\PYG{+w}{ }\PYG{n+nv+nvIdentifier}{electrodes}\PYG{+w}{ }\PYG{n}{electrode}\PYG{o}{/}\PYG{n}{conq}\PYG{+w}{ }\PYG{l+m}{0.0}\PYG{+w}{ }\PYG{l+m}{1.805}\PYG{+w}{ }\PYG{n}{algo}\PYG{+w}{ }\PYG{n}{cg}\PYG{+w}{ }\PYG{l+m}{1e\PYGZhy{}5}
\PYG{k}{fix }\PYG{n+nv+nvIdentifier}{fxconp}\PYG{+w}{ }\PYG{n+nv+nvIdentifier}{bot}\PYG{+w}{ }\PYG{n}{electrode}\PYG{o}{/}\PYG{n}{thermo}\PYG{+w}{ }\PYG{o}{\PYGZhy{}}\PYG{l+m}{1.0}\PYG{+w}{ }\PYG{l+m}{1.805}\PYG{+w}{ }\PYG{n}{temp}\PYG{+w}{ }\PYG{l+m}{298}\PYG{+w}{ }\PYG{l+m}{100}\PYG{+w}{ }\PYG{n}{couple}\PYG{+w}{ }\PYG{n}{top}\PYG{+w}{ }\PYG{l+m}{1.0}
\end{sphinxVerbatim}


\subsection{Description}
\label{\detokenize{fix_electrode:description}}
\sphinxAtStartPar
The \sphinxstyleemphasis{electrode} fixes implement the constant potential method (CPM)
({\hyperref[\detokenize{fix_electrode:siepmann}]{\sphinxcrossref{\DUrole{std,std-ref}{Siepmann}}}}, {\hyperref[\detokenize{fix_electrode:reed3}]{\sphinxcrossref{\DUrole{std,std-ref}{Reed}}}}), and modern variants,
to accurately model electrified, conductive electrodes. This is
primarily useful for studying electrode\sphinxhyphen{}electrolyte interfaces,
especially at high potential differences or ionicities, with non\sphinxhyphen{}planar
electrodes such as nanostructures or nanopores, and to study dynamic
phenomena such as charging or discharging time scales or conductivity or
ionic diffusivities.

\sphinxAtStartPar
Each \sphinxstyleemphasis{electrode} fix allows users to set additional electrostatic
relationships between the specified groups which model useful
electrostatic configurations:
\begin{itemize}
\item {} 
\sphinxAtStartPar
\sphinxstyleemphasis{electrode/conp} sets potentials or potential differences between electrodes
\begin{itemize}
\item {} 
\sphinxAtStartPar
(resulting in changing electrode total charges)

\end{itemize}

\item {} 
\sphinxAtStartPar
\sphinxstyleemphasis{electrode/conq} sets the total charge on each electrode
\begin{itemize}
\item {} 
\sphinxAtStartPar
(resulting in changing electrode potentials)

\end{itemize}

\item {} 
\sphinxAtStartPar
\sphinxstyleemphasis{electrode/thermo} sets a thermopotentiostat
{\hyperref[\detokenize{fix_electrode:deissenbeck}]{\sphinxcrossref{\DUrole{std,std-ref}{(Deissenbeck)}}}} between two electrodes
\begin{itemize}
\item {} 
\sphinxAtStartPar
(resulting in changing charges and potentials with appropriate
average potential difference and thermal variance)

\end{itemize}

\end{itemize}

\sphinxAtStartPar
The first group\sphinxhyphen{}ID provided to each fix specifies the first electrode
group, and more group(s) are added using the \sphinxstyleemphasis{couple} keyword for each
additional group.  While \sphinxstyleemphasis{electrode/thermo} only accepts two groups,
\sphinxstyleemphasis{electrode/conp} and \sphinxstyleemphasis{electrode/conq} accept any number of groups, up to
LAMMPS’s internal restrictions (see Restrictions below). Electrode
groups must not overlap, i.e.  the fix will issue an error if any
particle is detected to belong to at least two electrode groups.

\sphinxAtStartPar
CPM involves updating charges on groups of electrode particles, per time
step, so that the system’s total energy is minimized with respect to
those charges.  From basic electrostatics, this is equivalent to making
each group conductive, or imposing an equal electrostatic potential on
every particle in the same group (hence the name CPM).  The charges are
usually modelled as a Gaussian distribution to make the charge\sphinxhyphen{}charge
interaction matrix invertible ({\hyperref[\detokenize{fix_electrode:gingrich}]{\sphinxcrossref{\DUrole{std,std-ref}{Gingrich}}}}).  The keyword
\sphinxstyleemphasis{eta} specifies the distribution’s width in units of inverse length.

\sphinxAtStartPar
\DUrole{versionmodified,added}{Added in version 22Dec2022.}

\sphinxAtStartPar
Three algorithms are available to minimize the energy, varying in how
matrices are pre\sphinxhyphen{}calculated before a run to provide computational
speedup. These algorithms can be selected using the keyword \sphinxstyleemphasis{algo}:
\begin{itemize}
\item {} 
\sphinxAtStartPar
\sphinxstyleemphasis{algo mat\_inv} pre\sphinxhyphen{}calculates the capacitance matrix and obtains the
charge configuration in one matrix\sphinxhyphen{}vector calculation per time step

\item {} 
\sphinxAtStartPar
\sphinxstyleemphasis{algo mat\_cg} pre\sphinxhyphen{}calculates the elastance matrix (inverse of
capacitance matrix) and obtains the charge configuration using a
conjugate gradient solver in multiple matrix\sphinxhyphen{}vector calculations per
time step

\item {} 
\sphinxAtStartPar
\sphinxstyleemphasis{algo cg} does not perform any pre\sphinxhyphen{}calculation and obtains the charge
configuration using a conjugate gradient solver and multiple
calculations of the electric potential per time step.

\end{itemize}

\sphinxAtStartPar
For both \sphinxstyleemphasis{cg} methods, the command must specify the conjugate gradient
tolerance. \sphinxstyleemphasis{fix electrode/thermo} currently only supports the \sphinxstyleemphasis{mat\_inv}
algorithm.

\sphinxAtStartPar
The keyword \sphinxstyleemphasis{symm} can be set \sphinxstyleemphasis{on} (or \sphinxstyleemphasis{off}) to turn on (or turn off)
the capacitance matrix constraint that sets total electrode charge to be
zero.  This has slightly different effects for each \sphinxstyleemphasis{fix electrode}
variant.  For \sphinxstyleemphasis{fix electrode/conp}, with \sphinxstyleemphasis{symm off}, the potentials
specified are absolute potentials, but the charge configurations
satisfying them may add up to an overall non\sphinxhyphen{}zero, varying charge for
the electrodes (and thus the simulation box). With \sphinxstyleemphasis{symm on}, the total
charge over all electrode groups is constrained to zero, and potential
differences rather than absolute potentials are the physically relevant
quantities.

\sphinxAtStartPar
For \sphinxstyleemphasis{fix electrode/conq}, with \sphinxstyleemphasis{symm off}, overall neutrality is
explicitly obeyed or violated by the user input (which is not
checked!). With \sphinxstyleemphasis{symm on}, overall neutrality is ensured by ignoring the
user\sphinxhyphen{}input charge for the last listed electrode (instead, its charge
will always be minus the total sum of all other electrode charges). For
\sphinxstyleemphasis{fix electrode/thermo}, overall neutrality is always automatically
imposed for any setting of \sphinxstyleemphasis{symm}, but \sphinxstyleemphasis{symm on} allows finite\sphinxhyphen{}field
mode (\sphinxstyleemphasis{ffield on}, described below) for faster simulations.

\sphinxAtStartPar
For all three fixes, any potential (or charge for \sphinxstyleemphasis{conq}) can be
specified as an equal\sphinxhyphen{}style variable prefixed with “v\_”. For example,
the following code will ramp the potential difference between electrodes
from 0.0V to 2.0V over the course of the simulation:

\begin{sphinxVerbatim}[commandchars=\\\{\}]
\PYG{k}{fix }\PYG{n+nv+nvIdentifier}{fxconp}\PYG{+w}{ }\PYG{n+nv+nvIdentifier}{bot}\PYG{+w}{ }\PYG{n}{electrode}\PYG{o}{/}\PYG{n}{conp}\PYG{+w}{ }\PYG{l+m}{0.0}\PYG{+w}{ }\PYG{l+m}{1.805}\PYG{+w}{ }\PYG{n}{couple}\PYG{+w}{ }\PYG{n}{top}\PYG{+w}{ }\PYG{n}{v\PYGZus{}v}\PYG{+w}{ }\PYG{n}{symm}\PYG{+w}{ }\PYG{n}{on}
\PYG{k}{variable }\PYG{n+nv+nvIdentifier}{v}\PYG{+w}{ }\PYG{n}{equal}\PYG{+w}{ }\PYG{n}{ramp}\PYG{n+nv}{(}\PYG{n+nv}{0.0, 2.0}\PYG{n+nv}{)}
\end{sphinxVerbatim}

\sphinxAtStartPar
Note that these fixes only parse their supplied variable name when
starting a run, and so these fixes will accept equal\sphinxhyphen{}style variables
defined \sphinxstyleemphasis{after} the fix definition, including variables dependent on the
fix’s own output. This is useful, for example, in the fix’s internal
finite\sphinxhyphen{}field commands (see below).  For an advanced example of this see
the in.conq2 input file in the directory
\sphinxcode{\sphinxupquote{examples/PACKAGES/electrode/graph\sphinxhyphen{}il}}.

\sphinxAtStartPar
This fix necessitates the use of a long range solver that calculates and
provides the matrix of electrode\sphinxhyphen{}electrode interactions and a vector of
electrode\sphinxhyphen{}electrolyte interactions.  The Kspace styles
\sphinxstyleemphasis{ewald/electrode}, \sphinxstyleemphasis{pppm/electrode} and \sphinxstyleemphasis{pppm/electrode/intel} are
created specifically for this task {\hyperref[\detokenize{fix_electrode:ahrens-iwers}]{\sphinxcrossref{\DUrole{std,std-ref}{(Ahrens\sphinxhyphen{}Iwers)}}}}.

\sphinxAtStartPar
For systems with non\sphinxhyphen{}periodic boundaries in one or two directions dipole
corrections are available with the {\hyperref[\detokenize{kspace_modify::doc}]{\sphinxcrossref{\DUrole{doc}{kspace\_modify}}}}.
For ewald/electrode a two\sphinxhyphen{}dimensional Ewald summation {\hyperref[\detokenize{fix_electrode:hu}]{\sphinxcrossref{\DUrole{std,std-ref}{(Hu)}}}}
can be used by setting “slab ew2d”:

\begin{sphinxVerbatim}[commandchars=\\\{\}]
\PYG{k}{kspace\PYGZus{}modify}\PYG{+w}{ }\PYG{n}{slab}\PYG{+w}{ }\PYG{o}{\PYGZlt{}}\PYG{n}{slab\PYGZus{}factor}\PYG{o}{\PYGZgt{}}
\PYG{k}{kspace\PYGZus{}modify}\PYG{+w}{ }\PYG{n}{wire}\PYG{+w}{ }\PYG{o}{\PYGZlt{}}\PYG{n}{wire\PYGZus{}factor}\PYG{o}{\PYGZgt{}}
\PYG{k}{kspace\PYGZus{}modify}\PYG{+w}{ }\PYG{n}{slab}\PYG{+w}{ }\PYG{n}{ew2d}
\end{sphinxVerbatim}

\sphinxAtStartPar
Two implementations for the calculation of the elastance matrix are
available with pppm and can be selected using the \sphinxstyleemphasis{amat onestep/twostep}
keyword.  \sphinxstyleemphasis{onestep} is the default; \sphinxstyleemphasis{twostep} can be faster for large
electrodes and a moderate mesh size but requires more memory.

\begin{sphinxVerbatim}[commandchars=\\\{\}]
\PYG{k}{kspace\PYGZus{}modify}\PYG{+w}{ }\PYG{n}{amat}\PYG{+w}{ }\PYG{n}{onestep}\PYG{o}{/}\PYG{n}{twostep}
\end{sphinxVerbatim}

\sphinxAtStartPar
For all versions of the fix, the keyword\sphinxhyphen{}value \sphinxstyleemphasis{ffield on} enables the
finite\sphinxhyphen{}field mode ({\hyperref[\detokenize{fix_electrode:dufils}]{\sphinxcrossref{\DUrole{std,std-ref}{Dufils}}}}, {\hyperref[\detokenize{fix_electrode:tee}]{\sphinxcrossref{\DUrole{std,std-ref}{Tee}}}}), which uses
an electric field across a periodic cell instead of non\sphinxhyphen{}periodic
boundary conditions to impose a potential difference between the two
electrodes bounding the cell. The fix (with name \sphinxstyleemphasis{fix\sphinxhyphen{}ID}) detects which
of the two electrodes is “on top” (has the larger maximum \sphinxstyleemphasis{z}\sphinxhyphen{}coordinate
among all particles).  Assuming the first electrode group is on top, it
then issues the following commands internally:

\begin{sphinxVerbatim}[commandchars=\\\{\}]
\PYG{k}{variable }\PYG{n+nv+nvIdentifier}{fix\PYGZhy{}ID\PYGZus{}ffield\PYGZus{}zfield}\PYG{+w}{ }\PYG{n}{equal}\PYG{+w}{ }\PYG{n+nv}{(}\PYG{n+nv}{f\PYGZus{}fix\PYGZhy{}ID[2]\PYGZhy{}f\PYGZus{}fix\PYGZhy{}ID[1]}\PYG{n+nv}{)}\PYG{o}{/}\PYG{n}{lz}
\PYG{n}{efield}\PYG{+w}{ }\PYG{n}{fix}\PYG{o}{\PYGZhy{}}\PYG{n}{ID\PYGZus{}efield}\PYG{+w}{ }\PYG{n}{all}\PYG{+w}{ }\PYG{n}{efield}\PYG{+w}{ }\PYG{l+m}{0.0}\PYG{+w}{ }\PYG{l+m}{0.0}\PYG{+w}{ }\PYG{n}{v\PYGZus{}fix}\PYG{o}{\PYGZhy{}}\PYG{n}{ID\PYGZus{}ffield\PYGZus{}zfield}
\end{sphinxVerbatim}

\sphinxAtStartPar
which implements the required electric field as the potential difference
divided by cell length.  The internal commands use variable so that the
electric field will correctly vary with changing potentials in the
correct way (for example with equal\sphinxhyphen{}style potential difference or with
\sphinxstyleemphasis{fix electrode/conq}).  This keyword requires two electrodes and will
issue an error with any other number of electrodes. This keyword
requires electroneutrality to be imposed (\sphinxstyleemphasis{symm on}) and will issue an
error otherwise.

\sphinxAtStartPar
\DUrole{versionmodified,changed}{Changed in version 22Dec2022.}

\sphinxAtStartPar
For all versions of the fix, the keyword\sphinxhyphen{}value \sphinxstyleemphasis{etypes on} enables
type\sphinxhyphen{}based optimized neighbor lists. With this feature enabled, LAMMPS
provides the fix with an occasional neighbor list restricted to
electrode\sphinxhyphen{}electrode interactions for calculating the electrode matrix,
and a perpetual neighbor list restricted to electrode\sphinxhyphen{}electrolyte
interactions for calculating the electrode potentials, using particle
types to list only desired interactions, and typically resulting in
5\textendash{}10\% less computational time.  Without this feature the fix will
simply use the active pair style’s neighbor list.  This feature cannot
be enabled if any electrode particle has the same type as any
electrolyte particle (which would be unusual in a typical simulation)
and the fix will issue an error in that case.

\sphinxAtStartPar
\DUrole{versionmodified,added}{Added in version 17Apr2024.}

\sphinxAtStartPar
The keyword \sphinxstyleemphasis{qtotal} causes \sphinxstyleemphasis{fix electrode/conp} and \sphinxstyleemphasis{fix
electrode/thermo} to add an overall potential to all electrodes so that
the total charge on the electrodes is a specified amount (which may be
an equal\sphinxhyphen{}style variable).  For example, if a user wanted to simulate a
solution of excess cations such that the total electrolyte charge is +2,
setting \sphinxstyleemphasis{qtotal \sphinxhyphen{}2} would cause the total electrode charge to be \sphinxhyphen{}2, so
that the simulation box remains overall electroneutral. Since \sphinxstyleemphasis{fix
electrode/conq} constrains the total charges of individual electrodes,
and since \sphinxstyleemphasis{symm on} constrains the total charge of all electrodes to be
zero, either option is incompatible with the \sphinxstyleemphasis{qtotal} keyword (even if
\sphinxstyleemphasis{qtotal} is set to zero).

\sphinxAtStartPar
\DUrole{versionmodified,added}{Added in version 17Apr2024.}

\sphinxAtStartPar
The keyword \sphinxstyleemphasis{eta} takes the name of a custom double vector defined via
fix property/atom.  The values will be used instead of the standard eta
value.  The property/atom fix must be for vector of double values and
use the \sphinxstyleemphasis{ghost on} option.


\subsection{Restart, fix\_modify, output, run start/stop, minimize info}
\label{\detokenize{fix_electrode:restart-fix-modify-output-run-start-stop-minimize-info}}
\sphinxAtStartPar
This fix currently does not write any information to restart files.

\sphinxAtStartPar
The \sphinxstyleemphasis{fix\_modify tf} option enables the Thomas\sphinxhyphen{}Fermi metallicity model
({\hyperref[\detokenize{fix_electrode:scalfi}]{\sphinxcrossref{\DUrole{std,std-ref}{Scalfi}}}}) and allows parameters to be set for each atom type.

\begin{sphinxVerbatim}[commandchars=\\\{\}]
\PYG{k}{fix\PYGZus{}modify }\PYG{n+nv+nvIdentifier}{ID}\PYG{+w}{ }\PYG{n}{tf}\PYG{+w}{ }\PYG{n}{type}\PYG{+w}{ }\PYG{n}{length}\PYG{+w}{ }\PYG{n}{voronoi}
\end{sphinxVerbatim}

\sphinxAtStartPar
If this option is used, these two parameters must be set for
all atom types of the electrode:
\begin{itemize}
\item {} 
\sphinxAtStartPar
\sphinxtitleref{tf} is the Thomas\sphinxhyphen{}Fermi length \(l_{TF}\)

\item {} 
\sphinxAtStartPar
\sphinxtitleref{voronoi} is the Voronoi volume per atom in units of length cubed

\end{itemize}

\sphinxAtStartPar
Different types may have different \sphinxtitleref{tf} and \sphinxtitleref{voronoi} values.
The following self\sphinxhyphen{}energy term is then added for all electrode atoms:
\begin{equation*}
\begin{split}A_{ii} += \frac{1}{4 \pi \epsilon_0} \times \frac{4 \pi l_{TF}^2}{\mathrm{Voronoi volume}}\end{split}
\end{equation*}
\sphinxAtStartPar
The \sphinxstyleemphasis{fix\_modify timer} option turns on (off) additional timer outputs in the log
file, for code developers to track optimization.

\begin{sphinxVerbatim}[commandchars=\\\{\}]
\PYG{k}{fix\PYGZus{}modify }\PYG{n+nv+nvIdentifier}{ID}\PYG{+w}{ }\PYG{n}{timer}\PYG{+w}{ }\PYG{n}{on}\PYG{o}{/}\PYG{n}{off}
\end{sphinxVerbatim}


\bigskip\hrule\bigskip


\sphinxAtStartPar
These fixes compute a global (extensive) scalar, a global (intensive)
vector, and a global array, which can be accessed by various
{\hyperref[\detokenize{Howto_output::doc}]{\sphinxcrossref{\DUrole{doc}{output commands}}}}.

\sphinxAtStartPar
The global scalar outputs the energy added to the system by this fix,
which is the negative of the total charge on each electrode multiplied
by that electrode’s potential.

\sphinxAtStartPar
The global vector outputs the potential on each electrode (and thus has
\sphinxstyleemphasis{N} entries if the fix manages \sphinxstyleemphasis{N} electrode groups), in {\hyperref[\detokenize{units::doc}]{\sphinxcrossref{\DUrole{doc}{units}}}} of electric field multiplied by distance (thus volts for \sphinxstyleemphasis{real}
and \sphinxstyleemphasis{metal} units).  The electrode groups’ ordering follows the order in
which they were input in the fix command using \sphinxstyleemphasis{couple}. The global
vector output is useful for \sphinxstyleemphasis{fix electrode/conq} and \sphinxstyleemphasis{fix
electrode/thermo}, where potential is dynamically updated based on
electrolyte configuration instead of being directly set.

\sphinxAtStartPar
The global array has \sphinxstyleemphasis{N} rows and \sphinxstyleemphasis{2N+1} columns, where the fix manages
\sphinxstyleemphasis{N} electrode groups managed by the fix. For the \sphinxstyleemphasis{I}\sphinxhyphen{}th row of the
array, the elements are:
\begin{itemize}
\item {} 
\sphinxAtStartPar
array{[}I{]}{[}1{]} = total charge that group \sphinxstyleemphasis{I} would have had \sphinxstyleemphasis{if it were
at 0 V applied potential}

\item {} 
\sphinxAtStartPar
array{[}I{]}{[}2 to \sphinxstyleemphasis{N} + 1{]} = the \sphinxstyleemphasis{N} entries
of the \sphinxstyleemphasis{I}\sphinxhyphen{}th row of the electrode capacitance matrix (definition
follows)

\item {} 
\sphinxAtStartPar
array{[}I{]}{[}\sphinxstyleemphasis{N} + 2 to \sphinxstyleemphasis{2N} + 1{]} = the \sphinxstyleemphasis{N} entries of the
\sphinxstyleemphasis{I}\sphinxhyphen{}th row of the electrode elastance matrix (the inverse of the
electrode capacitance matrix)

\end{itemize}

\sphinxAtStartPar
The \(N \times N\) electrode capacitance matrix, denoted \(\mathbf{C}\)
in the following equation, summarizes how the total charge induced on each
electrode (\(\mathbf{Q}\) as an \sphinxstyleemphasis{N}\sphinxhyphen{}vector) is related to the potential on
each electrode, \(\mathbf{V}\), and the charge\sphinxhyphen{}at\sphinxhyphen{}0V \(\mathbf{Q}_{0V}\)
(which is influenced by the local electrolyte structure):
\begin{equation*}
\begin{split}\mathbf{Q} = \mathbf{Q}_{0V} + \mathbf{C} \cdot \mathbf{V}\end{split}
\end{equation*}
\sphinxAtStartPar
The charge\sphinxhyphen{}at\sphinxhyphen{}0V, electrode capacitance and elastance matrices are internally
used to calculate the potentials required to induce the specified total
electrode charges in \sphinxstyleemphasis{fix electrode/conq} and \sphinxstyleemphasis{fix electrode/thermo}. With the
\sphinxstyleemphasis{symm on} option, the electrode capacitance matrix would be singular, and thus
its last row is replaced with \sphinxstyleemphasis{N} copies of its top\sphinxhyphen{}left entry
(\(\mathbf{C}_{11}\)) for invertibility.

\sphinxAtStartPar
The global array output is mainly useful for quickly determining the ‘vacuum
capacitance’ of the system (capacitance with only electrodes, no electrolyte),
and can also be used for advanced simulations setting the potential as some
function of the charge\sphinxhyphen{}at\sphinxhyphen{}0V (such as the \sphinxcode{\sphinxupquote{in.conq2}} example mentioned above).

\sphinxAtStartPar
Please cite {\hyperref[\detokenize{fix_electrode:ahrens-iwers2}]{\sphinxcrossref{\DUrole{std,std-ref}{(Ahrens\sphinxhyphen{}Iwers2022)}}}} in any publication that
uses this implementation.  Please cite also the publication on the combination
of the CPM with PPPM if you use \sphinxstyleemphasis{pppm/electrode} {\hyperref[\detokenize{fix_electrode:ahrens-iwers}]{\sphinxcrossref{\DUrole{std,std-ref}{(Ahrens\sphinxhyphen{}Iwers)}}}}.


\bigskip\hrule\bigskip



\subsection{Restrictions}
\label{\detokenize{fix_electrode:restrictions}}
\sphinxAtStartPar
For algorithms that use a matrix for the electrode\sphinxhyphen{}electrode
interactions, positions of electrode particles have to be immobilized at
all times.

\sphinxAtStartPar
With \sphinxstyleemphasis{ffield off} (i.e. the default), the box geometry is expected to be
\sphinxstyleemphasis{z}\sphinxhyphen{}non\sphinxhyphen{}periodic (i.e. \sphinxstyleemphasis{boundary p p f}), and this fix will issue an
error if the box is \sphinxstyleemphasis{z}\sphinxhyphen{}periodic. With \sphinxstyleemphasis{ffield on}, the box geometry is
expected to be \sphinxstyleemphasis{z}\sphinxhyphen{}periodic, and this fix will issue an error if the box
is \sphinxstyleemphasis{z}\sphinxhyphen{}non\sphinxhyphen{}periodic.

\sphinxAtStartPar
The parallelization for the fix works best if electrode atoms are evenly
distributed across processors. For a system with two electrodes at the bottom
and top of the cell this can be achieved with \sphinxstyleemphasis{processors * * 2}, or with the
line

\begin{sphinxVerbatim}[commandchars=\\\{\}]
\PYG{k}{if}\PYG{+w}{ }\PYG{l+s}{\PYGZdq{}}\PYG{l+s}{\PYGZdl{}(extract\PYGZus{}setting(world\PYGZus{}size) \PYGZpc{} 2) == 0}\PYG{l+s}{\PYGZdq{}}\PYG{+w}{ }\PYG{k}{then}\PYG{+w}{ }\PYG{l+s}{\PYGZdq{}}\PYG{l+s}{processors * * 2}\PYG{l+s}{\PYGZdq{}}
\end{sphinxVerbatim}

\sphinxAtStartPar
which avoids an error if the script is run on an odd number of
processors (such as on just one processor for testing).

\sphinxAtStartPar
The fix creates an additional group named \sphinxstyleemphasis{{[}fix\sphinxhyphen{}ID{]}\_group} which is the
union of all electrode groups supplied to LAMMPS. This additional group
counts towards LAMMPS’s limitation on the total number of groups
(currently 32), which may not allow scripts that use that many groups to
run with this fix.

\sphinxAtStartPar
The matrix\sphinxhyphen{}based algorithms (\sphinxstyleemphasis{algo mat\_inv} and \sphinxstyleemphasis{algo mat\_cg}) currently
store an interaction matrix (either elastance or capacitance) of \sphinxstyleemphasis{N} by
\sphinxstyleemphasis{N} doubles for each MPI process. This memory requirement may be
prohibitive for large electrode groups.  The fix will issue a warning if
it expects to use more than 0.5 GiB of memory.


\subsection{Default}
\label{\detokenize{fix_electrode:default}}
\sphinxAtStartPar
The default keyword\sphinxhyphen{}option settings are \sphinxstyleemphasis{algo mat\_inv}, \sphinxstyleemphasis{symm off},
\sphinxstyleemphasis{etypes off} and \sphinxstyleemphasis{ffield off}.


\bigskip\hrule\bigskip


\sphinxAtStartPar
Styles with a \sphinxstyleemphasis{gpu}, \sphinxstyleemphasis{intel}, \sphinxstyleemphasis{kk}, \sphinxstyleemphasis{omp}, or \sphinxstyleemphasis{opt} suffix are
functionally the same as the corresponding style without the suffix.
They have been optimized to run faster, depending on your available
hardware, as discussed on the {\hyperref[\detokenize{Speed_packages::doc}]{\sphinxcrossref{\DUrole{doc}{Accelerator packages}}}}
page.  The accelerated styles take the same arguments and should
produce the same results, except for round\sphinxhyphen{}off and precision issues.

\sphinxAtStartPar
These accelerated styles are part of the GPU, INTEL, KOKKOS,
OPENMP, and OPT packages, respectively.  They are only enabled if
LAMMPS was built with those packages.  See the {\hyperref[\detokenize{Build_package::doc}]{\sphinxcrossref{\DUrole{doc}{Build package}}}} page for more info.

\sphinxAtStartPar
You can specify the accelerated styles explicitly in your input script
by including their suffix, or you can use the {\hyperref[\detokenize{Run_options::doc}]{\sphinxcrossref{\DUrole{doc}{\sphinxhyphen{}suffix command\sphinxhyphen{}line switch}}}} when you invoke LAMMPS, or you can use the
{\hyperref[\detokenize{suffix::doc}]{\sphinxcrossref{\DUrole{doc}{suffix}}}} command in your input script.

\sphinxAtStartPar
See the {\hyperref[\detokenize{Speed_packages::doc}]{\sphinxcrossref{\DUrole{doc}{Accelerator packages}}}} page for more
instructions on how to use the accelerated styles effectively.


\bigskip\hrule\bigskip

\phantomsection\label{\detokenize{fix_electrode:siepmann}}
\sphinxAtStartPar
\sphinxstylestrong{(Siepmann)} Siepmann and Sprik, J. Chem. Phys. 102, 511 (1995).

\phantomsection\label{\detokenize{fix_electrode:reed3}}
\sphinxAtStartPar
\sphinxstylestrong{(Reed)} Reed \sphinxstyleemphasis{et al.}, J. Chem. Phys. 126, 084704 (2007).

\phantomsection\label{\detokenize{fix_electrode:deissenbeck}}
\sphinxAtStartPar
\sphinxstylestrong{(Deissenbeck)} Deissenbeck \sphinxstyleemphasis{et al.}, Phys. Rev. Letters 126, 136803 (2021).

\phantomsection\label{\detokenize{fix_electrode:gingrich}}
\sphinxAtStartPar
\sphinxstylestrong{(Gingrich)} Gingrich, \sphinxtitleref{MSc thesis} \textless{}\sphinxurl{https://gingrich.chem.northwestern.edu/papers/ThesiswCorrections.pdf}\textgreater{}\textasciigrave{} (2010).

\phantomsection\label{\detokenize{fix_electrode:ahrens-iwers}}
\sphinxAtStartPar
\sphinxstylestrong{(Ahrens\sphinxhyphen{}Iwers)} Ahrens\sphinxhyphen{}Iwers and Meissner, J. Chem. Phys. 155, 104104 (2021).

\phantomsection\label{\detokenize{fix_electrode:hu}}
\sphinxAtStartPar
\sphinxstylestrong{(Hu)} Hu, J. Chem. Theory Comput. 10, 5254 (2014).

\phantomsection\label{\detokenize{fix_electrode:dufils}}
\sphinxAtStartPar
\sphinxstylestrong{(Dufils)} Dufils \sphinxstyleemphasis{et al.}, Phys. Rev. Letters 123, 195501 (2019).

\phantomsection\label{\detokenize{fix_electrode:tee}}
\sphinxAtStartPar
\sphinxstylestrong{(Tee)} Tee and Searles, J. Chem. Phys. 156, 184101 (2022).

\phantomsection\label{\detokenize{fix_electrode:scalfi}}
\sphinxAtStartPar
\sphinxstylestrong{(Scalfi)} Scalfi \sphinxstyleemphasis{et al.}, J. Chem. Phys., 153, 174704 (2020).

\phantomsection\label{\detokenize{fix_electrode:ahrens-iwers2}}
\sphinxAtStartPar
\sphinxstylestrong{(Ahrens\sphinxhyphen{}Iwers2022)} Ahrens\sphinxhyphen{}Iwers \sphinxstyleemphasis{et al.}, J. Chem. Phys. 157, 084801 (2022).

\sphinxstepscope

\index{fix electron/stopping@\spxentry{fix electron/stopping}}\ignorespaces 
\index{fix electron/stopping/fit@\spxentry{fix electron/stopping/fit}}\ignorespaces 

\section{fix electron/stopping command}
\label{\detokenize{fix_electron_stopping:fix-electron-stopping-command}}\label{\detokenize{fix_electron_stopping:index-1}}\label{\detokenize{fix_electron_stopping:index-0}}\label{\detokenize{fix_electron_stopping::doc}}

\section{fix electron/stopping/fit command}
\label{\detokenize{fix_electron_stopping:fix-electron-stopping-fit-command}}

\subsection{Syntax}
\label{\detokenize{fix_electron_stopping:syntax}}
\begin{sphinxVerbatim}[commandchars=\\\{\}]
\PYG{k}{fix }\PYG{n+nv+nvIdentifier}{ID}\PYG{+w}{ }\PYG{n+nv+nvIdentifier}{group\PYGZhy{}ID}\PYG{+w}{ }\PYG{n}{style}\PYG{+w}{ }\PYG{n}{args}
\end{sphinxVerbatim}
\begin{itemize}
\item {} 
\sphinxAtStartPar
ID, group\sphinxhyphen{}ID are documented in {\hyperref[\detokenize{fix::doc}]{\sphinxcrossref{\DUrole{doc}{fix}}}} command

\item {} 
\sphinxAtStartPar
style = \sphinxstyleemphasis{electron/stopping} or \sphinxstyleemphasis{electron/stopping/fit}
\begin{sphinxalltt}
\sphinxstyleemphasis{electron/stopping} args = Ecut file keyword value ...
  Ecut  = minimum kinetic energy for electronic stopping (energy units)
  file  = name of the file containing the electronic stopping power table

\sphinxstyleemphasis{electron/stopping/fit} args = Ecut c1 c2 ...
  Ecut  = minimum kinetic energy for electronic stopping (energy units)
  c1 c2 = linear and quadratic coefficients for the fitted quadratic polynomial
\end{sphinxalltt}

\item {} 
\sphinxAtStartPar
zero or more keyword/value pairs may be appended to args for style = \sphinxstyleemphasis{electron/stopping}
\begin{sphinxalltt}
keyword = \sphinxstyleemphasis{region} or \sphinxstyleemphasis{minneigh}
   \sphinxstyleemphasis{region} value = region\sphinxhyphen{}ID
      region\sphinxhyphen{}ID = region whose atoms will be affected by this fix
   \sphinxstyleemphasis{minneigh} value = minneigh
      minneigh = minimum number of neighbors an atom to have stopping applied
\end{sphinxalltt}

\end{itemize}


\subsection{Examples}
\label{\detokenize{fix_electron_stopping:examples}}
\begin{sphinxVerbatim}[commandchars=\\\{\}]
\PYG{k}{fix }\PYG{n+nv+nvIdentifier}{el}\PYG{+w}{ }\PYG{n+nv+nvIdentifier}{all}\PYG{+w}{ }\PYG{n}{electron}\PYG{o}{/}\PYG{n}{stopping}\PYG{+w}{ }\PYG{l+m}{10.0}\PYG{+w}{ }\PYG{n}{elstop}\PYG{o}{\PYGZhy{}}\PYG{n}{table.txt}
\PYG{k}{fix }\PYG{n+nv+nvIdentifier}{el}\PYG{+w}{ }\PYG{n+nv+nvIdentifier}{all}\PYG{+w}{ }\PYG{n}{electron}\PYG{o}{/}\PYG{n}{stopping}\PYG{+w}{ }\PYG{l+m}{10.0}\PYG{+w}{ }\PYG{n}{elstop}\PYG{o}{\PYGZhy{}}\PYG{n}{table.txt}\PYG{+w}{ }\PYG{n}{minneigh}\PYG{+w}{ }\PYG{l+m}{3}
\PYG{k}{fix }\PYG{n+nv+nvIdentifier}{el}\PYG{+w}{ }\PYG{n+nv+nvIdentifier}{mygroup}\PYG{+w}{ }\PYG{n}{electron}\PYG{o}{/}\PYG{n}{stopping}\PYG{+w}{ }\PYG{l+m}{1.0}\PYG{+w}{ }\PYG{n}{elstop}\PYG{o}{\PYGZhy{}}\PYG{n}{table.txt}\PYG{+w}{ }\PYG{k}{region }\PYG{n+nv+nvIdentifier}{bulk}
\PYG{k}{fix }\PYG{n+nv+nvIdentifier}{1}\PYG{+w}{ }\PYG{n+nv+nvIdentifier}{all}\PYG{+w}{ }\PYG{n}{electron}\PYG{o}{/}\PYG{n}{stopping}\PYG{o}{/}\PYG{n}{fit}\PYG{+w}{ }\PYG{l+m}{4.63}\PYG{+w}{ }\PYG{l+m}{3.3e\PYGZhy{}3}\PYG{+w}{ }\PYG{l+m}{4.0e\PYGZhy{}8}
\PYG{k}{fix }\PYG{n+nv+nvIdentifier}{1}\PYG{+w}{ }\PYG{n+nv+nvIdentifier}{all}\PYG{+w}{ }\PYG{n}{electron}\PYG{o}{/}\PYG{n}{stopping}\PYG{o}{/}\PYG{n}{fit}\PYG{+w}{ }\PYG{l+m}{3.49}\PYG{+w}{ }\PYG{l+m}{1.8e\PYGZhy{}3}\PYG{+w}{ }\PYG{l+m}{9.0e\PYGZhy{}8}\PYG{+w}{ }\PYG{l+m}{7.57}\PYG{+w}{ }\PYG{l+m}{4.2e\PYGZhy{}3}\PYG{+w}{ }\PYG{l+m}{5.0e\PYGZhy{}8}
\end{sphinxVerbatim}


\subsection{Description}
\label{\detokenize{fix_electron_stopping:description}}
\sphinxAtStartPar
This fix implements inelastic energy loss for fast projectiles in solids. It
applies a friction force to fast moving atoms to slow them down due to
{\hyperref[\detokenize{fix_electron_stopping:elstopping}]{\sphinxcrossref{\DUrole{std,std-ref}{electronic stopping}}}} (energy lost via electronic collisions per
unit of distance). This fix should be used for simulation of irradiation
damage or ion implantation, where the ions can lose noticeable amounts of
energy from electron excitations. If the electronic stopping power is not
considered, the simulated range of the ions can be severely overestimated
({\hyperref[\detokenize{fix_electron_stopping:nordlund98}]{\sphinxcrossref{\DUrole{std,std-ref}{Nordlund98}}}}, {\hyperref[\detokenize{fix_electron_stopping:nordlund95}]{\sphinxcrossref{\DUrole{std,std-ref}{Nordlund95}}}}).

\sphinxAtStartPar
The electronic stopping is implemented by applying a friction force
to each atom as:
\begin{equation*}
\begin{split}\vec{F}_i = \vec{F}^0_i - \frac{\vec{v}_i}{\|\vec{v}_i\|} \cdot S_e\end{split}
\end{equation*}
\sphinxAtStartPar
where \(\vec{F}_i\) is the resulting total force on the atom.
\(\vec{F}^0_i\) is the original force applied to the atom, \(\vec{v}_i\) is
its velocity and \(S_e\) is the stopping power of the ion.

\begin{sphinxadmonition}{note}{Note:}
\sphinxAtStartPar
In addition to electronic stopping, atomic cascades and irradiation
simulations require the use of an adaptive timestep (see
{\hyperref[\detokenize{fix_dt_reset::doc}]{\sphinxcrossref{\DUrole{doc}{fix dt/reset}}}}) and the repulsive ZBL potential (see
{\hyperref[\detokenize{pair_zbl::doc}]{\sphinxcrossref{\DUrole{doc}{ZBL}}}} potential) or similar. Without these settings the
interaction between the ion and the target atoms will be faulty. It is also
common to use in such simulations a thermostat ({\hyperref[\detokenize{fix_nh::doc}]{\sphinxcrossref{\DUrole{doc}{fix\_nvt}}}}) in
the borders of the simulation cell.
\end{sphinxadmonition}

\begin{sphinxadmonition}{note}{Note:}
\sphinxAtStartPar
This fix removes energy from fast projectiles without depositing it as a
heat to the simulation cell. Such implementation might lead to the unphysical
results when the amount of energy deposited to the electronic system is large,
e.g. simulations of Swift Heavy Ions (energy per nucleon of 100 keV/amu or
higher) or multiple projectiles. You could compensate energy loss by coupling
bulk atoms with some thermostat or control heat transfer between electronic and
atomic subsystems with the two\sphinxhyphen{}temperature model ({\hyperref[\detokenize{fix_ttm::doc}]{\sphinxcrossref{\DUrole{doc}{fix\_ttm}}}}).
\end{sphinxadmonition}

\sphinxAtStartPar
At low velocities the electronic stopping is negligible. The electronic
friction is not applied to atoms whose kinetic energy is smaller than \sphinxstyleemphasis{Ecut},
or smaller than the lowest energy value given in the table in \sphinxstyleemphasis{file}.
Electronic stopping should be applied only when a projectile reaches bulk
material. This fix scans neighbor list and excludes atoms with fewer than
\sphinxstyleemphasis{minneigh} neighbors (by default one). If the pair potential cutoff is large,
minneigh should be increased, though not above the number of nearest neighbors
in bulk material. An alternative is to disable the check for neighbors by
setting \sphinxstyleemphasis{minneigh} to zero and using the \sphinxstyleemphasis{region} keyword. This is necessary
when running simulations of cluster bombardment.

\sphinxAtStartPar
If the \sphinxstyleemphasis{region} keyword is used, the atom must also be in the specified
geometric {\hyperref[\detokenize{region::doc}]{\sphinxcrossref{\DUrole{doc}{region}}}} in order to have electronic stopping applied to
it. This is useful if the position of the bulk material is fixed. By default
the electronic stopping is applied everywhere in the simulation cell.


\bigskip\hrule\bigskip


\sphinxAtStartPar
The energy ranges and stopping powers are read from the file \sphinxstyleemphasis{file}.
Lines starting with \sphinxstyleemphasis{\#} and empty lines are ignored. Otherwise each
line must contain exactly \sphinxstylestrong{N+1} numbers, where \sphinxstylestrong{N} is the number of atom
types in the simulation.

\sphinxAtStartPar
The first column is the energy for which the stopping powers on that
line apply. The energies must be sorted from the smallest to the largest.
The other columns are the stopping powers \(S_e\) for each atom type,
in ascending order, in force {\hyperref[\detokenize{units::doc}]{\sphinxcrossref{\DUrole{doc}{units}}}}. The stopping powers for
intermediate energy values are calculated with linear interpolation between
2 nearest points.

\sphinxAtStartPar
For example:

\begin{sphinxVerbatim}[commandchars=\\\{\}]
\PYGZsh{} This is a comment
\PYGZsh{}       atom\PYGZhy{}1    atom\PYGZhy{}2
\PYGZsh{} eV    eV/Ang    eV/Ang  \PYGZsh{} units metal
 10        0        0
250       60       80
750      100      150
\end{sphinxVerbatim}

\sphinxAtStartPar
If an atom which would have electronic stopping applied to it has a
kinetic energy higher than the largest energy given in \sphinxstyleemphasis{file}, LAMMPS
will exit with an error message.

\sphinxAtStartPar
The stopping power depends on the energy of the ion and the target
material. The electronic stopping table can be obtained from
scientific publications, experimental databases or by using
{\hyperref[\detokenize{fix_electron_stopping:srim}]{\sphinxcrossref{\DUrole{std,std-ref}{SRIM}}}} software. Other programs such as {\hyperref[\detokenize{fix_electron_stopping:casp}]{\sphinxcrossref{\DUrole{std,std-ref}{CasP}}}} or
{\hyperref[\detokenize{fix_electron_stopping:pass}]{\sphinxcrossref{\DUrole{std,std-ref}{PASS}}}} can calculate the energy deposited as a function
of the impact parameter of the ion; these results can be used
to derive the stopping power.


\bigskip\hrule\bigskip


\sphinxAtStartPar
Style \sphinxstyleemphasis{electron/stopping/fit} calculates the electronic stopping power
and cumulative energy lost to the electron gas via a quadratic functional
and applies a drag force to the classical equations\sphinxhyphen{}of\sphinxhyphen{}motion for all
atoms moving above some minimum cutoff velocity (i.e., kinetic energy).
These coefficients can be determined by fitting a quadratic polynomial to
electronic stopping data predicted by, for example, SRIM or TD\sphinxhyphen{}DFT. Multiple
‘Ecut c1 c2’ values can be provided for multi\sphinxhyphen{}species simulations in the order
of the atom types. There is an examples/PACKAGES/electron\_stopping/ directory,
which illustrates uses of this command. Details of this implementation are
further described in {\hyperref[\detokenize{fix_electron_stopping:stewart2018}]{\sphinxcrossref{\DUrole{std,std-ref}{Stewart2018}}}} and {\hyperref[\detokenize{fix_electron_stopping:lee2020}]{\sphinxcrossref{\DUrole{std,std-ref}{Lee2020}}}}.


\subsection{Restart, fix\_modify, output, run start/stop, minimize info}
\label{\detokenize{fix_electron_stopping:restart-fix-modify-output-run-start-stop-minimize-info}}
\sphinxAtStartPar
No information about this fix is written to {\hyperref[\detokenize{restart::doc}]{\sphinxcrossref{\DUrole{doc}{binary restart files}}}}.

\sphinxAtStartPar
The {\hyperref[\detokenize{fix_modify::doc}]{\sphinxcrossref{\DUrole{doc}{fix\_modify}}}} options are not supported.

\sphinxAtStartPar
This fix computes a global scalar, which can be accessed by various
{\hyperref[\detokenize{Howto_output::doc}]{\sphinxcrossref{\DUrole{doc}{output commands}}}}. The scalar is the total energy
loss from electronic stopping applied by this fix since the start of
the latest run. It is considered “intensive”.

\sphinxAtStartPar
The \sphinxstyleemphasis{start/stop} keywords of the {\hyperref[\detokenize{run::doc}]{\sphinxcrossref{\DUrole{doc}{run}}}} command have no effect
on this fix.


\subsection{Restrictions}
\label{\detokenize{fix_electron_stopping:restrictions}}
\sphinxAtStartPar
This pair style is part of the EXTRA\sphinxhyphen{}FIX package. It is only enabled if
LAMMPS was built with that package. See the {\hyperref[\detokenize{Build_package::doc}]{\sphinxcrossref{\DUrole{doc}{Build package}}}}
doc page for more info.


\subsection{Default}
\label{\detokenize{fix_electron_stopping:default}}
\sphinxAtStartPar
The default is no limitation by region, and minneigh = 1.


\bigskip\hrule\bigskip

\phantomsection\label{\detokenize{fix_electron_stopping:elstopping}}
\sphinxAtStartPar
\sphinxstylestrong{(electronic stopping)} Wikipedia \sphinxhyphen{} Electronic Stopping Power: \sphinxurl{https://en.wikipedia.org/wiki/Stopping\_power\_\%28particle\_radiation\%29}

\phantomsection\label{\detokenize{fix_electron_stopping:nordlund98}}
\sphinxAtStartPar
\sphinxstylestrong{(Nordlund98)} Nordlund, Kai, et al.  Physical Review B 57.13 (1998): 7556.

\phantomsection\label{\detokenize{fix_electron_stopping:nordlund95}}
\sphinxAtStartPar
\sphinxstylestrong{(Nordlund95)} Nordlund, Kai. Computational materials science 3.4 (1995): 448\sphinxhyphen{}456.

\phantomsection\label{\detokenize{fix_electron_stopping:srim}}
\sphinxAtStartPar
\sphinxstylestrong{(SRIM)} SRIM webpage: \sphinxurl{http://www.srim.org/}

\phantomsection\label{\detokenize{fix_electron_stopping:casp}}
\sphinxAtStartPar
\sphinxstylestrong{(CasP)} CasP webpage: \sphinxurl{http://www.casp-program.org/}

\phantomsection\label{\detokenize{fix_electron_stopping:pass}}
\sphinxAtStartPar
\sphinxstylestrong{(PASS)} PASS webpage: \sphinxurl{https://www.sdu.dk/en/DPASS}

\phantomsection\label{\detokenize{fix_electron_stopping:stewart2018}}
\sphinxAtStartPar
\sphinxstylestrong{(Stewart2018)} J.A. Stewart, et al. (2018) Journal of Applied Physics, 123(16), 165902.

\phantomsection\label{\detokenize{fix_electron_stopping:lee2020}}
\sphinxAtStartPar
\sphinxstylestrong{(Lee2020)} C.W. Lee, et al. (2020) Physical Review B, 102(2), 024107.

\sphinxstepscope

\index{fix enforce2d@\spxentry{fix enforce2d}}\ignorespaces 
\index{fix enforce2d/kk@\spxentry{fix enforce2d/kk}}\ignorespaces 

\section{fix enforce2d command}
\label{\detokenize{fix_enforce2d:fix-enforce2d-command}}\label{\detokenize{fix_enforce2d:index-1}}\label{\detokenize{fix_enforce2d:index-0}}\label{\detokenize{fix_enforce2d::doc}}
\sphinxAtStartPar
Accelerator Variants: \sphinxstyleemphasis{enforce2d/kk}


\subsection{Syntax}
\label{\detokenize{fix_enforce2d:syntax}}
\begin{sphinxVerbatim}[commandchars=\\\{\}]
\PYG{k}{fix }\PYG{n+nv+nvIdentifier}{ID}\PYG{+w}{ }\PYG{n+nv+nvIdentifier}{group\PYGZhy{}ID}\PYG{+w}{ }\PYG{n}{enforce2d}
\end{sphinxVerbatim}
\begin{itemize}
\item {} 
\sphinxAtStartPar
ID, group\sphinxhyphen{}ID are documented in {\hyperref[\detokenize{fix::doc}]{\sphinxcrossref{\DUrole{doc}{fix}}}} command

\item {} 
\sphinxAtStartPar
enforce2d = style name of this fix command

\end{itemize}


\subsection{Examples}
\label{\detokenize{fix_enforce2d:examples}}
\begin{sphinxVerbatim}[commandchars=\\\{\}]
\PYG{k}{fix }\PYG{n+nv+nvIdentifier}{5}\PYG{+w}{ }\PYG{n+nv+nvIdentifier}{all}\PYG{+w}{ }\PYG{n}{enforce2d}
\end{sphinxVerbatim}


\subsection{Description}
\label{\detokenize{fix_enforce2d:description}}
\sphinxAtStartPar
Zero out the z\sphinxhyphen{}dimension velocity and force on each atom in the group.
This is useful when running a 2d simulation to ensure that atoms do
not move from their initial z coordinate.


\bigskip\hrule\bigskip


\sphinxAtStartPar
Styles with a \sphinxstyleemphasis{gpu}, \sphinxstyleemphasis{intel}, \sphinxstyleemphasis{kk}, \sphinxstyleemphasis{omp}, or \sphinxstyleemphasis{opt} suffix are
functionally the same as the corresponding style without the suffix.
They have been optimized to run faster, depending on your available
hardware, as discussed on the {\hyperref[\detokenize{Speed_packages::doc}]{\sphinxcrossref{\DUrole{doc}{Accelerator packages}}}}
page.  The accelerated styles take the same arguments and should
produce the same results, except for round\sphinxhyphen{}off and precision issues.

\sphinxAtStartPar
These accelerated styles are part of the GPU, INTEL, KOKKOS,
OPENMP, and OPT packages, respectively.  They are only enabled if
LAMMPS was built with those packages.  See the {\hyperref[\detokenize{Build_package::doc}]{\sphinxcrossref{\DUrole{doc}{Build package}}}} page for more info.

\sphinxAtStartPar
You can specify the accelerated styles explicitly in your input script
by including their suffix, or you can use the {\hyperref[\detokenize{Run_options::doc}]{\sphinxcrossref{\DUrole{doc}{\sphinxhyphen{}suffix command\sphinxhyphen{}line switch}}}} when you invoke LAMMPS, or you can use the
{\hyperref[\detokenize{suffix::doc}]{\sphinxcrossref{\DUrole{doc}{suffix}}}} command in your input script.

\sphinxAtStartPar
See the {\hyperref[\detokenize{Speed_packages::doc}]{\sphinxcrossref{\DUrole{doc}{Accelerator packages}}}} page for more
instructions on how to use the accelerated styles effectively.


\bigskip\hrule\bigskip



\subsection{Restart, fix\_modify, output, run start/stop, minimize info}
\label{\detokenize{fix_enforce2d:restart-fix-modify-output-run-start-stop-minimize-info}}
\sphinxAtStartPar
No information about this fix is written to {\hyperref[\detokenize{restart::doc}]{\sphinxcrossref{\DUrole{doc}{binary restart files}}}}.  None of the {\hyperref[\detokenize{fix_modify::doc}]{\sphinxcrossref{\DUrole{doc}{fix\_modify}}}} options
are relevant to this fix.  No global or per\sphinxhyphen{}atom quantities are stored
by this fix for access by various {\hyperref[\detokenize{Howto_output::doc}]{\sphinxcrossref{\DUrole{doc}{output commands}}}}.
No parameter of this fix can be used with the \sphinxstyleemphasis{start/stop} keywords of
the {\hyperref[\detokenize{run::doc}]{\sphinxcrossref{\DUrole{doc}{run}}}} command.

\sphinxAtStartPar
The forces due to this fix are imposed during an energy minimization,
invoked by the {\hyperref[\detokenize{minimize::doc}]{\sphinxcrossref{\DUrole{doc}{minimize}}}} command.


\subsection{Restrictions}
\label{\detokenize{fix_enforce2d:restrictions}}\begin{quote}

\sphinxAtStartPar
none
\end{quote}


\subsection{Related commands}
\label{\detokenize{fix_enforce2d:related-commands}}
\sphinxAtStartPar
none


\subsection{Default}
\label{\detokenize{fix_enforce2d:default}}
\sphinxAtStartPar
none

\sphinxstepscope

\index{fix eos/cv@\spxentry{fix eos/cv}}\ignorespaces 

\section{fix eos/cv command}
\label{\detokenize{fix_eos_cv:fix-eos-cv-command}}\label{\detokenize{fix_eos_cv:index-0}}\label{\detokenize{fix_eos_cv::doc}}

\subsection{Syntax}
\label{\detokenize{fix_eos_cv:syntax}}
\begin{sphinxVerbatim}[commandchars=\\\{\}]
\PYG{k}{fix }\PYG{n+nv+nvIdentifier}{ID}\PYG{+w}{ }\PYG{n+nv+nvIdentifier}{group\PYGZhy{}ID}\PYG{+w}{ }\PYG{n}{eos}\PYG{o}{/}\PYG{n}{cv}\PYG{+w}{ }\PYG{n}{cv}
\end{sphinxVerbatim}
\begin{itemize}
\item {} 
\sphinxAtStartPar
ID, group\sphinxhyphen{}ID are documented in {\hyperref[\detokenize{fix::doc}]{\sphinxcrossref{\DUrole{doc}{fix}}}} command

\item {} 
\sphinxAtStartPar
eos/cv = style name of this fix command

\item {} 
\sphinxAtStartPar
cv = constant\sphinxhyphen{}volume heat capacity (energy/temperature units)

\end{itemize}


\subsection{Examples}
\label{\detokenize{fix_eos_cv:examples}}
\begin{sphinxVerbatim}[commandchars=\\\{\}]
\PYG{k}{fix }\PYG{n+nv+nvIdentifier}{1}\PYG{+w}{ }\PYG{n+nv+nvIdentifier}{all}\PYG{+w}{ }\PYG{n}{eos}\PYG{o}{/}\PYG{n}{cv}\PYG{+w}{ }\PYG{l+m}{0.01}
\end{sphinxVerbatim}


\subsection{Description}
\label{\detokenize{fix_eos_cv:description}}
\sphinxAtStartPar
Fix \sphinxstyleemphasis{eos/cv} applies a mesoparticle equation of state to relate the
particle internal energy (\(u_i\)) to the particle internal temperature
(\(\theta_i\)).  The \sphinxstyleemphasis{eos/cv} mesoparticle equation of state requires
the constant\sphinxhyphen{}volume heat capacity, and is defined as follows:
\begin{equation*}
\begin{split}u_{i} = u^{mech}_{i} + u^{cond}_{i} = C_{V} \theta_{i}\end{split}
\end{equation*}
\sphinxAtStartPar
where \(C_V\) is the constant\sphinxhyphen{}volume heat capacity, \(u^{cond}\)
is the internal conductive energy, and \(u^{mech}\) is the internal
mechanical energy.  Note that alternative definitions of the mesoparticle
equation of state are possible.


\bigskip\hrule\bigskip



\subsection{Restrictions}
\label{\detokenize{fix_eos_cv:restrictions}}
\sphinxAtStartPar
This command is part of the DPD\sphinxhyphen{}REACT package.  It is only enabled if
LAMMPS was built with that package.  See the {\hyperref[\detokenize{Build_package::doc}]{\sphinxcrossref{\DUrole{doc}{Build package}}}} page for more info.

\sphinxAtStartPar
This command also requires use of the {\hyperref[\detokenize{atom_style::doc}]{\sphinxcrossref{\DUrole{doc}{atom\_style dpd}}}}
command.


\subsection{Related commands}
\label{\detokenize{fix_eos_cv:related-commands}}
\sphinxAtStartPar
{\hyperref[\detokenize{fix_shardlow::doc}]{\sphinxcrossref{\DUrole{doc}{fix shardlow}}}}, {\hyperref[\detokenize{pair_dpd_fdt::doc}]{\sphinxcrossref{\DUrole{doc}{pair dpd/fdt}}}}


\subsection{Default}
\label{\detokenize{fix_eos_cv:default}}
\sphinxAtStartPar
none


\bigskip\hrule\bigskip

\phantomsection\label{\detokenize{fix_eos_cv:larentzos4}}
\sphinxAtStartPar
\sphinxstylestrong{(Larentzos)} J.P. Larentzos, J.K. Brennan, J.D. Moore, and
W.D. Mattson, “LAMMPS Implementation of Constant Energy Dissipative
Particle Dynamics (DPD\sphinxhyphen{}E)”, ARL\sphinxhyphen{}TR\sphinxhyphen{}6863, U.S. Army Research
Laboratory, Aberdeen Proving Ground, MD (2014).

\sphinxstepscope

\index{fix eos/table@\spxentry{fix eos/table}}\ignorespaces 

\section{fix eos/table command}
\label{\detokenize{fix_eos_table:fix-eos-table-command}}\label{\detokenize{fix_eos_table:index-0}}\label{\detokenize{fix_eos_table::doc}}

\subsection{Syntax}
\label{\detokenize{fix_eos_table:syntax}}
\begin{sphinxVerbatim}[commandchars=\\\{\}]
\PYG{k}{fix }\PYG{n+nv+nvIdentifier}{ID}\PYG{+w}{ }\PYG{n+nv+nvIdentifier}{group\PYGZhy{}ID}\PYG{+w}{ }\PYG{n}{eos}\PYG{o}{/}\PYG{n}{table}\PYG{+w}{ }\PYG{n}{style}\PYG{+w}{ }\PYG{n}{file}\PYG{+w}{ }\PYG{n}{N}\PYG{+w}{ }\PYG{n}{keyword}
\end{sphinxVerbatim}
\begin{itemize}
\item {} 
\sphinxAtStartPar
ID, group\sphinxhyphen{}ID are documented in {\hyperref[\detokenize{fix::doc}]{\sphinxcrossref{\DUrole{doc}{fix}}}} command

\item {} 
\sphinxAtStartPar
eos/table = style name of this fix command

\item {} 
\sphinxAtStartPar
style = \sphinxstyleemphasis{linear} = method of interpolation

\item {} 
\sphinxAtStartPar
file = filename containing the tabulated equation of state

\item {} 
\sphinxAtStartPar
N = use N values in \sphinxstyleemphasis{linear} tables

\item {} 
\sphinxAtStartPar
keyword = name of table keyword corresponding to table file

\end{itemize}


\subsection{Examples}
\label{\detokenize{fix_eos_table:examples}}
\begin{sphinxVerbatim}[commandchars=\\\{\}]
\PYG{k}{fix }\PYG{n+nv+nvIdentifier}{1}\PYG{+w}{ }\PYG{n+nv+nvIdentifier}{all}\PYG{+w}{ }\PYG{n}{eos}\PYG{o}{/}\PYG{n}{table}\PYG{+w}{ }\PYG{n}{linear}\PYG{+w}{ }\PYG{n}{eos.table}\PYG{+w}{ }\PYG{l+m}{100000}\PYG{+w}{ }\PYG{n}{KEYWORD}
\end{sphinxVerbatim}


\subsection{Description}
\label{\detokenize{fix_eos_table:description}}
\sphinxAtStartPar
Fix \sphinxstyleemphasis{eos/table} applies a tabulated mesoparticle equation of state to
relate the particle internal energy (u\_i) to the particle internal
temperature (dpdTheta\_i).

\sphinxAtStartPar
Fix \sphinxstyleemphasis{eos/table} creates interpolation tables of length \sphinxstyleemphasis{N} from
internal energy values listed in a file as a function of internal
temperature.

\sphinxAtStartPar
The interpolation tables are created by fitting cubic splines to the
file values and interpolating energy values at each of \sphinxstyleemphasis{N} internal
temperatures, and vice versa.  During a simulation, these tables are
used to interpolate internal energy or temperature values as needed.
The interpolation is done with the \sphinxstyleemphasis{linear} style.

\sphinxAtStartPar
For the \sphinxstyleemphasis{linear} style, the internal temperature is used to find 2
surrounding table values from which an internal energy is computed by
linear interpolation, and vice versa.

\sphinxAtStartPar
The filename specifies a file containing tabulated internal
temperature and internal energy values.  The keyword specifies a
section of the file.  The format of this file is described below.


\bigskip\hrule\bigskip


\sphinxAtStartPar
The format of a tabulated file is as follows (without the
parenthesized comments):

\begin{sphinxVerbatim}[commandchars=\\\{\}]
\PYGZsh{} EOS TABLE                (one or more comment or blank lines)

KEYWORD                    (keyword is first text on line)
N 500                      (N  parameter)
                           (blank)
1   1.00 0.000             (index, internal temperature, internal energy)
2   1.02 0.001
...
500 10.0 0.500
\end{sphinxVerbatim}

\sphinxAtStartPar
A section begins with a non\sphinxhyphen{}blank line whose first character is not a
“\#”; blank lines or lines starting with “\#” can be used as comments
between sections.  The first line begins with a keyword which
identifies the section.  The line can contain additional text, but the
initial text must match the argument specified in the fix command.

\sphinxAtStartPar
The next line lists the number of table entries.  The parameter “N” is
required and its value is the number of table entries that follow.
Note that this may be different than the \sphinxstyleemphasis{N} specified in the {\hyperref[\detokenize{fix_eos_table::doc}]{\sphinxcrossref{\DUrole{doc}{fix eos/table}}}} command.  Let Ntable = \sphinxstyleemphasis{N} in the fix
command, and Nfile = “N” in the tabulated file.  What LAMMPS does is a
preliminary interpolation by creating splines using the Nfile
tabulated values as nodal points.  It uses these to interpolate as
needed to generate energy and temperature values at Ntable different
points.  The resulting tables of length Ntable are then used as
described above, when computing energy and temperature relationships.
This means that if you want the interpolation tables of length Ntable
to match exactly what is in the tabulated file (with effectively no
preliminary interpolation), you should set Ntable = Nfile.

\sphinxAtStartPar
Following a blank line, the next N lines list the tabulated values.
On each line, the first value is the index from 1 to N, the second value is
the internal temperature (in temperature units), the third value is the
internal energy (in energy units).

\sphinxAtStartPar
Note that the internal temperature and internal energy values must
increase from one line to the next.

\sphinxAtStartPar
Note that one file can contain many sections, each with a tabulated
potential.  LAMMPS reads the file section by section until it finds
one that matches the specified keyword.


\bigskip\hrule\bigskip



\subsection{Restrictions}
\label{\detokenize{fix_eos_table:restrictions}}
\sphinxAtStartPar
This command is part of the DPD\sphinxhyphen{}REACT package.  It is only enabled if
LAMMPS was built with that package.  See the {\hyperref[\detokenize{Build_package::doc}]{\sphinxcrossref{\DUrole{doc}{Build package}}}} page for more info.

\sphinxAtStartPar
This command also requires use of the {\hyperref[\detokenize{atom_style::doc}]{\sphinxcrossref{\DUrole{doc}{atom\_style dpd}}}}
command.

\sphinxAtStartPar
The equation of state must be a monotonically increasing function.

\sphinxAtStartPar
An error will occur if the internal temperature or internal energies
are not within the table cutoffs.


\subsection{Related commands}
\label{\detokenize{fix_eos_table:related-commands}}
\sphinxAtStartPar
{\hyperref[\detokenize{fix_shardlow::doc}]{\sphinxcrossref{\DUrole{doc}{fix shardlow}}}}, {\hyperref[\detokenize{pair_dpd_fdt::doc}]{\sphinxcrossref{\DUrole{doc}{pair dpd/fdt}}}}


\subsection{Default}
\label{\detokenize{fix_eos_table:default}}
\sphinxAtStartPar
none

\sphinxstepscope

\index{fix eos/table/rx@\spxentry{fix eos/table/rx}}\ignorespaces 
\index{fix eos/table/rx/kk@\spxentry{fix eos/table/rx/kk}}\ignorespaces 

\section{fix eos/table/rx command}
\label{\detokenize{fix_eos_table_rx:fix-eos-table-rx-command}}\label{\detokenize{fix_eos_table_rx:index-1}}\label{\detokenize{fix_eos_table_rx:index-0}}\label{\detokenize{fix_eos_table_rx::doc}}
\sphinxAtStartPar
Accelerator Variants: \sphinxstyleemphasis{eos/table/rx/kk}


\subsection{Syntax}
\label{\detokenize{fix_eos_table_rx:syntax}}
\begin{sphinxVerbatim}[commandchars=\\\{\}]
\PYG{k}{fix }\PYG{n+nv+nvIdentifier}{ID}\PYG{+w}{ }\PYG{n+nv+nvIdentifier}{group\PYGZhy{}ID}\PYG{+w}{ }\PYG{n}{eos}\PYG{o}{/}\PYG{n}{table}\PYG{o}{/}\PYG{n}{rx}\PYG{+w}{ }\PYG{n}{style}\PYG{+w}{ }\PYG{n}{file1}\PYG{+w}{ }\PYG{n}{N}\PYG{+w}{ }\PYG{n}{keyword}\PYG{+w}{ }\PYG{n}{...}
\end{sphinxVerbatim}
\begin{itemize}
\item {} 
\sphinxAtStartPar
ID, group\sphinxhyphen{}ID are documented in {\hyperref[\detokenize{fix::doc}]{\sphinxcrossref{\DUrole{doc}{fix}}}} command

\item {} 
\sphinxAtStartPar
eos/table/rx = style name of this fix command

\item {} 
\sphinxAtStartPar
style = \sphinxstyleemphasis{linear} = method of interpolation

\item {} 
\sphinxAtStartPar
file1 = filename containing the tabulated equation of state

\item {} 
\sphinxAtStartPar
N = use N values in \sphinxstyleemphasis{linear} tables

\item {} 
\sphinxAtStartPar
keyword = name of table keyword corresponding to table file

\item {} 
\sphinxAtStartPar
file2 = filename containing the heats of formation of each species (optional)

\item {} 
\sphinxAtStartPar
deltaHf = heat of formation for a single species in energy units (optional)

\item {} 
\sphinxAtStartPar
energyCorr = energy correction in energy units (optional)

\item {} 
\sphinxAtStartPar
tempCorrCoeff = temperature correction coefficient (optional)

\end{itemize}


\subsection{Examples}
\label{\detokenize{fix_eos_table_rx:examples}}
\begin{sphinxVerbatim}[commandchars=\\\{\}]
\PYG{k}{fix }\PYG{n+nv+nvIdentifier}{1}\PYG{+w}{ }\PYG{n+nv+nvIdentifier}{all}\PYG{+w}{ }\PYG{n}{eos}\PYG{o}{/}\PYG{n}{table}\PYG{o}{/}\PYG{n}{rx}\PYG{+w}{ }\PYG{n}{linear}\PYG{+w}{ }\PYG{n}{eos.table}\PYG{+w}{ }\PYG{l+m}{10000}\PYG{+w}{ }\PYG{n}{KEYWORD}\PYG{+w}{ }\PYG{n}{thermo.table}
\PYG{k}{fix }\PYG{n+nv+nvIdentifier}{1}\PYG{+w}{ }\PYG{n+nv+nvIdentifier}{all}\PYG{+w}{ }\PYG{n}{eos}\PYG{o}{/}\PYG{n}{table}\PYG{o}{/}\PYG{n}{rx}\PYG{+w}{ }\PYG{n}{linear}\PYG{+w}{ }\PYG{n}{eos.table}\PYG{+w}{ }\PYG{l+m}{10000}\PYG{+w}{ }\PYG{n}{KEYWORD}\PYG{+w}{ }\PYG{l+m}{1.5}
\PYG{k}{fix }\PYG{n+nv+nvIdentifier}{1}\PYG{+w}{ }\PYG{n+nv+nvIdentifier}{all}\PYG{+w}{ }\PYG{n}{eos}\PYG{o}{/}\PYG{n}{table}\PYG{o}{/}\PYG{n}{rx}\PYG{+w}{ }\PYG{n}{linear}\PYG{+w}{ }\PYG{n}{eos.table}\PYG{+w}{ }\PYG{l+m}{10000}\PYG{+w}{ }\PYG{n}{KEYWORD}\PYG{+w}{ }\PYG{l+m}{1.5}\PYG{+w}{ }\PYG{l+m}{0.025}\PYG{+w}{ }\PYG{l+m}{0.0}
\end{sphinxVerbatim}


\subsection{Description}
\label{\detokenize{fix_eos_table_rx:description}}
\sphinxAtStartPar
Fix \sphinxstyleemphasis{eos/table/rx} applies a tabulated mesoparticle equation
of state to relate the concentration\sphinxhyphen{}dependent particle internal
energy (\(u_i\)) to the particle internal temperature (\(\theta_i\)).

\sphinxAtStartPar
The concentration\sphinxhyphen{}dependent particle internal energy (\(u_i\)) is
computed according to the following relation:
\begin{equation*}
\begin{split}U_{i} = \displaystyle\sum_{j=1}^{m} c_{i,j}(u_{j} + \Delta H_{f,j}) + \frac{3k_{B}T}{2} + Nk_{B}T \\\end{split}
\end{equation*}
\sphinxAtStartPar
where \sphinxstyleemphasis{m} is the number of species, \(c_{i,j}\) is the
concentration of species \sphinxstyleemphasis{j} in particle \sphinxstyleemphasis{i}, \(u_j\) is the
internal energy of species j, \(\Delta H_{f,j} is the heat of
formation of species *j*, N is the number of molecules represented
by the coarse-grained particle, :math:`k_B\) is the Boltzmann constant,
and \(T\) is the temperature of the system.  Additionally, it is
possible to modify the concentration\sphinxhyphen{}dependent particle internal
energy relation by adding an energy correction, temperature\sphinxhyphen{}dependent
correction, and/or a molecule\sphinxhyphen{}dependent correction.  An energy
correction can be specified as a constant (in energy units).  A
temperature correction can be specified by multiplying a temperature
correction coefficient by the internal temperature.  A molecular
correction can be specified by by multiplying a molecule correction
coefficient by the average number of product gas particles in the
coarse\sphinxhyphen{}grain particle.

\sphinxAtStartPar
Fix \sphinxstyleemphasis{eos/table/rx} creates interpolation tables of length \sphinxstyleemphasis{N} from \sphinxstyleemphasis{m}
internal energy values of each species \(u_j\) listed in a file as a
function of internal temperature.  During a simulation, these tables
are used to interpolate internal energy or temperature values as needed.
The interpolation is done with the \sphinxstyleemphasis{linear} style.  For the \sphinxstyleemphasis{linear} style,
the internal temperature is used to find 2 surrounding table values from
which an internal energy is computed by linear interpolation.  A secant
solver is used to determine the internal temperature from the internal energy.

\sphinxAtStartPar
The first filename specifies a file containing tabulated internal
temperature and \sphinxstyleemphasis{m} internal energy values for each species \(u_j\).
The keyword specifies a section of the file.  The format of this
file is described below.

\sphinxAtStartPar
The second filename specifies a file containing heat of formation
\(\Delta H_{f,j}\) for each species.

\sphinxAtStartPar
In cases where the coarse\sphinxhyphen{}grain particle represents a single molecular
species (i.e., no reactions occur and fix \sphinxstyleemphasis{rx} is not present in the
input file), fix \sphinxstyleemphasis{eos/table/rx} can be applied in a similar manner to
fix \sphinxstyleemphasis{eos/table} within a non\sphinxhyphen{}reactive DPD simulation.  In this case,
the heat of formation filename is replaced with the heat of formation
value for the single species.  Additionally, the energy correction and
temperature correction coefficients may also be specified as fix
arguments.


\bigskip\hrule\bigskip


\sphinxAtStartPar
The format of a tabulated file is as follows (without the
parenthesized comments):

\begin{sphinxVerbatim}[commandchars=\\\{\}]
\PYGZsh{} EOS TABLE                (one or more comment or blank lines)

KEYWORD                    (keyword is first text on line)
N 500 h2 no2 n2 ... no     (N  parameter species1 species2 ... speciesN)
                           (blank)
1   1.00 0.000 ... 0.0000  (index, internal temperature, internal energy of species 1, ..., internal energy of species m)
2   1.02 0.001 ... 0.0002
...
500 10.0 0.500 ... 1.0000
\end{sphinxVerbatim}

\sphinxAtStartPar
A section begins with a non\sphinxhyphen{}blank line whose first character is not a
“\#”; blank lines or lines starting with “\#” can be used as comments
between sections.  The first line begins with a keyword which
identifies the section.  The line can contain additional text, but the
initial text must match the argument specified in the fix command.

\sphinxAtStartPar
The next line lists the number of table entries and the species names
that correspond with all the species listed in the reaction equations
through the \sphinxstyleemphasis{fix rx} command.
The parameter “N” is required and its value is the number of table
entries that follow.  Let Nfile = “N” in the tabulated file.
What LAMMPS does is a preliminary interpolation by creating splines
using the Nfile tabulated values as nodal points.

\sphinxAtStartPar
Following a blank line, the next N lines list the tabulated values.
On each line, the first value is the index from 1 to N, the second value is
the internal temperature (in temperature units), the third value until
the \sphinxstyleemphasis{m+3} value are the internal energies of the m species (in energy units).

\sphinxAtStartPar
Note that all internal temperature and internal energy values must
increase from one line to the next.

\sphinxAtStartPar
Note that one file can contain many sections, each with a tabulated
potential.  LAMMPS reads the file section by section until it finds
one that matches the specified keyword.


\bigskip\hrule\bigskip


\sphinxAtStartPar
The format of a heat of formation file is as follows (without the
parenthesized comments):

\begin{sphinxVerbatim}[commandchars=\\\{\}]
\PYGZsh{} HEAT OF FORMATION TABLE  (one or more comment or blank lines)

                           (blank)
h2      0.00               (species name, heat of formation)
no2     0.34
n2      0.00
...
no      0.93
\end{sphinxVerbatim}

\sphinxAtStartPar
Note that the species can be listed in any order.  The tag that is
used as the species name must correspond with the tags used to define
the reactions with the {\hyperref[\detokenize{fix_rx::doc}]{\sphinxcrossref{\DUrole{doc}{fix rx}}}} command.

\sphinxAtStartPar
Alternatively, corrections to the EOS can be included by specifying
three additional columns that correspond to the energy correction,
the temperature correction coefficient and molecule correction
coefficient.  In this case, the format of the file is as follows:

\begin{sphinxVerbatim}[commandchars=\\\{\}]
\PYGZsh{} HEAT OF FORMATION TABLE     (one or more comment or blank lines)

                              (blank)
h2      0.00 1.23 0.025  0.0  (species name, heat of formation, energy correction, temperature correction coefficient, molecule correction coefficient)
no2     0.34 0.00 0.000 \PYGZhy{}1.76
n2      0.00 0.00 0.000 \PYGZhy{}1.76
...
no      0.93 0.00 0.000 \PYGZhy{}1.76
\end{sphinxVerbatim}


\bigskip\hrule\bigskip


\sphinxAtStartPar
Styles with a \sphinxstyleemphasis{gpu}, \sphinxstyleemphasis{intel}, \sphinxstyleemphasis{kk}, \sphinxstyleemphasis{omp}, or \sphinxstyleemphasis{opt} suffix are
functionally the same as the corresponding style without the suffix.
They have been optimized to run faster, depending on your available
hardware, as discussed on the {\hyperref[\detokenize{Speed_packages::doc}]{\sphinxcrossref{\DUrole{doc}{Accelerator packages}}}}
page.  The accelerated styles take the same arguments and should
produce the same results, except for round\sphinxhyphen{}off and precision issues.

\sphinxAtStartPar
These accelerated styles are part of the GPU, INTEL, KOKKOS,
OPENMP, and OPT packages, respectively.  They are only enabled if
LAMMPS was built with those packages.  See the {\hyperref[\detokenize{Build_package::doc}]{\sphinxcrossref{\DUrole{doc}{Build package}}}} page for more info.

\sphinxAtStartPar
You can specify the accelerated styles explicitly in your input script
by including their suffix, or you can use the {\hyperref[\detokenize{Run_options::doc}]{\sphinxcrossref{\DUrole{doc}{\sphinxhyphen{}suffix command\sphinxhyphen{}line switch}}}} when you invoke LAMMPS, or you can use the
{\hyperref[\detokenize{suffix::doc}]{\sphinxcrossref{\DUrole{doc}{suffix}}}} command in your input script.

\sphinxAtStartPar
See the {\hyperref[\detokenize{Speed_packages::doc}]{\sphinxcrossref{\DUrole{doc}{Accelerator packages}}}} page for more
instructions on how to use the accelerated styles effectively.


\bigskip\hrule\bigskip



\subsection{Restrictions}
\label{\detokenize{fix_eos_table_rx:restrictions}}
\sphinxAtStartPar
This command is part of the DPD\sphinxhyphen{}REACT package.  It is only enabled if
LAMMPS was built with that package.  See the {\hyperref[\detokenize{Build_package::doc}]{\sphinxcrossref{\DUrole{doc}{Build package}}}} page for more info.

\sphinxAtStartPar
This command also requires use of the {\hyperref[\detokenize{atom_style::doc}]{\sphinxcrossref{\DUrole{doc}{atom\_style dpd}}}}
command.

\sphinxAtStartPar
The equation of state must be a monotonically increasing function.

\sphinxAtStartPar
An error will occur if the internal temperature or internal energies
are not within the table cutoffs.


\subsection{Related commands}
\label{\detokenize{fix_eos_table_rx:related-commands}}
\sphinxAtStartPar
{\hyperref[\detokenize{fix_rx::doc}]{\sphinxcrossref{\DUrole{doc}{fix rx}}}},
{\hyperref[\detokenize{pair_dpd_fdt::doc}]{\sphinxcrossref{\DUrole{doc}{pair dpd/fdt}}}}


\subsection{Default}
\label{\detokenize{fix_eos_table_rx:default}}
\sphinxAtStartPar
none

\sphinxstepscope

\index{fix evaporate@\spxentry{fix evaporate}}\ignorespaces 

\section{fix evaporate command}
\label{\detokenize{fix_evaporate:fix-evaporate-command}}\label{\detokenize{fix_evaporate:index-0}}\label{\detokenize{fix_evaporate::doc}}

\subsection{Syntax}
\label{\detokenize{fix_evaporate:syntax}}
\begin{sphinxVerbatim}[commandchars=\\\{\}]
\PYG{k}{fix }\PYG{n+nv+nvIdentifier}{ID}\PYG{+w}{ }\PYG{n+nv+nvIdentifier}{group\PYGZhy{}ID}\PYG{+w}{ }\PYG{n}{evaporate}\PYG{+w}{ }\PYG{n}{N}\PYG{+w}{ }\PYG{n}{M}\PYG{+w}{ }\PYG{n}{region}\PYG{o}{\PYGZhy{}}\PYG{n}{ID}\PYG{+w}{ }\PYG{n}{seed}
\end{sphinxVerbatim}
\begin{itemize}
\item {} 
\sphinxAtStartPar
ID, group\sphinxhyphen{}ID are documented in {\hyperref[\detokenize{fix::doc}]{\sphinxcrossref{\DUrole{doc}{fix}}}} command

\item {} 
\sphinxAtStartPar
evaporate = style name of this fix command

\item {} 
\sphinxAtStartPar
N = delete atoms every this many timesteps

\item {} 
\sphinxAtStartPar
M = number of atoms to delete each time

\item {} 
\sphinxAtStartPar
region\sphinxhyphen{}ID = ID of region within which to perform deletions

\item {} 
\sphinxAtStartPar
seed = random number seed to use for choosing atoms to delete

\item {} 
\sphinxAtStartPar
zero or more keyword/value pairs may be appended
\begin{sphinxalltt}
keyword = \sphinxstyleemphasis{molecule}
  \sphinxstyleemphasis{molecule} value = \sphinxstyleemphasis{no} or \sphinxstyleemphasis{yes}
\end{sphinxalltt}

\end{itemize}


\subsection{Examples}
\label{\detokenize{fix_evaporate:examples}}
\begin{sphinxVerbatim}[commandchars=\\\{\}]
\PYG{k}{fix }\PYG{n+nv+nvIdentifier}{1}\PYG{+w}{ }\PYG{n+nv+nvIdentifier}{solvent}\PYG{+w}{ }\PYG{n}{evaporate}\PYG{+w}{ }\PYG{l+m}{1000}\PYG{+w}{ }\PYG{l+m}{10}\PYG{+w}{ }\PYG{n}{surface}\PYG{+w}{ }\PYG{l+m}{49892}
\PYG{k}{fix }\PYG{n+nv+nvIdentifier}{1}\PYG{+w}{ }\PYG{n+nv+nvIdentifier}{solvent}\PYG{+w}{ }\PYG{n}{evaporate}\PYG{+w}{ }\PYG{l+m}{1000}\PYG{+w}{ }\PYG{l+m}{10}\PYG{+w}{ }\PYG{n}{surface}\PYG{+w}{ }\PYG{l+m}{38277}\PYG{+w}{ }\PYG{n}{molecule}\PYG{+w}{ }\PYG{n}{yes}
\end{sphinxVerbatim}


\subsection{Description}
\label{\detokenize{fix_evaporate:description}}
\sphinxAtStartPar
Remove M atoms from the simulation every N steps.  This can be used,
for example, to model evaporation of solvent particles or molecules
(i.e. drying) of a system.  Every N steps, the number of atoms in the
fix group and within the specified region are counted.  M of these are
chosen at random and deleted.  If there are less than M eligible
particles, then all of them are deleted.

\sphinxAtStartPar
If the setting for the \sphinxstyleemphasis{molecule} keyword is \sphinxstyleemphasis{no}, then only single
atoms are deleted.  In this case, you should ensure you do not delete
only a portion of a molecule (only some of its atoms), or LAMMPS will
soon generate an error when it tries to find those atoms.  LAMMPS will
warn you if any of the atoms eligible for deletion have a non\sphinxhyphen{}zero
molecule ID, but does not check for this at the time of deletion.

\sphinxAtStartPar
If the setting for the \sphinxstyleemphasis{molecule} keyword is \sphinxstyleemphasis{yes}, then when an atom
is chosen for deletion, the entire molecule it is part of is deleted.
The count of deleted atoms is incremented by the number of atoms in
the molecule, which may make it exceed \sphinxstyleemphasis{M}.  If the molecule ID of the
chosen atom is 0, then it is assumed to not be part of a molecule, and
just the single atom is deleted.

\sphinxAtStartPar
As an example, if you wish to delete 10 water molecules every \sphinxstyleemphasis{N}
steps, you should set \sphinxstyleemphasis{M} to 30.  If only the water’s oxygen atoms
were in the fix group, then two hydrogen atoms would be deleted when
an oxygen atom is selected for deletion, whether the hydrogen atoms
are inside the evaporation region or not.

\sphinxAtStartPar
Note that neighbor lists are re\sphinxhyphen{}built on timesteps that atoms are
removed.  Thus you should not remove atoms too frequently or you will
incur overhead due to the cost of building neighbor lists.

\begin{sphinxadmonition}{note}{Note:}
\sphinxAtStartPar
If you are monitoring the temperature of a system where the atom
count is changing due to evaporation, you typically should use the
{\hyperref[\detokenize{compute_modify::doc}]{\sphinxcrossref{\DUrole{doc}{compute\_modify dynamic/dof yes}}}} command for the
temperature compute you are using.
\end{sphinxadmonition}


\subsection{Restart, fix\_modify, output, run start/stop, minimize info}
\label{\detokenize{fix_evaporate:restart-fix-modify-output-run-start-stop-minimize-info}}
\sphinxAtStartPar
No information about this fix is written to {\hyperref[\detokenize{restart::doc}]{\sphinxcrossref{\DUrole{doc}{binary restart files}}}}.  None of the {\hyperref[\detokenize{fix_modify::doc}]{\sphinxcrossref{\DUrole{doc}{fix\_modify}}}} options are
relevant to this fix.

\sphinxAtStartPar
This fix computes a global scalar, which can be accessed by various
{\hyperref[\detokenize{Howto_output::doc}]{\sphinxcrossref{\DUrole{doc}{output commands}}}}.  The scalar is the cumulative
number of deleted atoms.  The scalar value calculated by this fix is
“intensive”.

\sphinxAtStartPar
No parameter of this fix can be used with the \sphinxstyleemphasis{start/stop} keywords of
the {\hyperref[\detokenize{run::doc}]{\sphinxcrossref{\DUrole{doc}{run}}}} command.  This fix is not invoked during
{\hyperref[\detokenize{minimize::doc}]{\sphinxcrossref{\DUrole{doc}{energy minimization}}}}.


\subsection{Restrictions}
\label{\detokenize{fix_evaporate:restrictions}}
\sphinxAtStartPar
None


\subsection{Related commands}
\label{\detokenize{fix_evaporate:related-commands}}
\sphinxAtStartPar
{\hyperref[\detokenize{fix_deposit::doc}]{\sphinxcrossref{\DUrole{doc}{fix deposit}}}}


\subsection{Default}
\label{\detokenize{fix_evaporate:default}}
\sphinxAtStartPar
The option defaults are molecule = no.

\sphinxstepscope

\index{fix external@\spxentry{fix external}}\ignorespaces 

\section{fix external command}
\label{\detokenize{fix_external:fix-external-command}}\label{\detokenize{fix_external:index-0}}\label{\detokenize{fix_external::doc}}

\subsection{Syntax}
\label{\detokenize{fix_external:syntax}}
\begin{sphinxVerbatim}[commandchars=\\\{\}]
\PYG{k}{fix }\PYG{n+nv+nvIdentifier}{ID}\PYG{+w}{ }\PYG{n+nv+nvIdentifier}{group\PYGZhy{}ID}\PYG{+w}{ }\PYG{n}{external}\PYG{+w}{ }\PYG{n}{mode}\PYG{+w}{ }\PYG{n}{args}
\end{sphinxVerbatim}
\begin{itemize}
\item {} 
\sphinxAtStartPar
ID, group\sphinxhyphen{}ID are documented in {\hyperref[\detokenize{fix::doc}]{\sphinxcrossref{\DUrole{doc}{fix}}}} command

\item {} 
\sphinxAtStartPar
external = style name of this fix command

\item {} 
\sphinxAtStartPar
mode = \sphinxstyleemphasis{pf/callback} or \sphinxstyleemphasis{pf/array}
\begin{sphinxalltt}
\sphinxstyleemphasis{pf/callback} args = Ncall Napply
  Ncall = make callback every Ncall steps
  Napply = apply callback forces every Napply steps
\sphinxstyleemphasis{pf/array} args = Napply
  Napply = apply array forces every Napply steps
\end{sphinxalltt}

\end{itemize}


\subsection{Examples}
\label{\detokenize{fix_external:examples}}
\begin{sphinxVerbatim}[commandchars=\\\{\}]
\PYG{k}{fix }\PYG{n+nv+nvIdentifier}{1}\PYG{+w}{ }\PYG{n+nv+nvIdentifier}{all}\PYG{+w}{ }\PYG{n}{external}\PYG{+w}{ }\PYG{n}{pf}\PYG{o}{/}\PYG{n}{callback}\PYG{+w}{ }\PYG{l+m}{1}\PYG{+w}{ }\PYG{l+m}{1}
\PYG{k}{fix }\PYG{n+nv+nvIdentifier}{1}\PYG{+w}{ }\PYG{n+nv+nvIdentifier}{all}\PYG{+w}{ }\PYG{n}{external}\PYG{+w}{ }\PYG{n}{pf}\PYG{o}{/}\PYG{n}{callback}\PYG{+w}{ }\PYG{l+m}{100}\PYG{+w}{ }\PYG{l+m}{1}
\PYG{k}{fix }\PYG{n+nv+nvIdentifier}{1}\PYG{+w}{ }\PYG{n+nv+nvIdentifier}{all}\PYG{+w}{ }\PYG{n}{external}\PYG{+w}{ }\PYG{n}{pf}\PYG{o}{/}\PYG{n}{array}\PYG{+w}{ }\PYG{l+m}{10}
\end{sphinxVerbatim}


\subsection{Description}
\label{\detokenize{fix_external:description}}
\sphinxAtStartPar
This fix allows external programs that are running LAMMPS through its
{\hyperref[\detokenize{Howto_library::doc}]{\sphinxcrossref{\DUrole{doc}{library interface}}}} to modify certain LAMMPS
properties on specific timesteps, similar to the way other fixes do.
The external driver can be a {\hyperref[\detokenize{Howto_library::doc}]{\sphinxcrossref{\DUrole{doc}{C/C++ or Fortran program}}}} or a {\hyperref[\detokenize{Python_head::doc}]{\sphinxcrossref{\DUrole{doc}{Python script}}}}.


\bigskip\hrule\bigskip


\sphinxAtStartPar
If mode is \sphinxstyleemphasis{pf/callback} then the fix will make a callback every
\sphinxstyleemphasis{Ncall} timesteps or minimization iterations to the external program.
The external program computes forces on atoms by setting values in an
array owned by the fix.  The fix then adds these forces to each atom
in the group, once every \sphinxstyleemphasis{Napply} steps, similar to the way the {\hyperref[\detokenize{fix_addforce::doc}]{\sphinxcrossref{\DUrole{doc}{fix addforce}}}} command works.  Note that if \sphinxstyleemphasis{Ncall} \textgreater{}
\sphinxstyleemphasis{Napply}, the force values produced by one callback will persist, and
be used multiple times to update atom forces.

\sphinxAtStartPar
The callback function “foo” is invoked by the fix as:

\begin{sphinxVerbatim}[commandchars=\\\{\}]
\PYG{n}{foo}\PYG{p}{(}\PYG{k+kt}{void}\PYG{+w}{ }\PYG{o}{*}\PYG{n}{ptr}\PYG{p}{,}\PYG{+w}{ }\PYG{n}{bigint}\PYG{+w}{ }\PYG{n}{timestep}\PYG{p}{,}\PYG{+w}{ }\PYG{k+kt}{int}\PYG{+w}{ }\PYG{n}{nlocal}\PYG{p}{,}\PYG{+w}{ }\PYG{n}{tagint}\PYG{+w}{ }\PYG{o}{*}\PYG{n}{ids}\PYG{p}{,}\PYG{+w}{ }\PYG{k+kt}{double}\PYG{+w}{ }\PYG{o}{*}\PYG{o}{*}\PYG{n}{x}\PYG{p}{,}\PYG{+w}{ }\PYG{k+kt}{double}\PYG{+w}{ }\PYG{o}{*}\PYG{o}{*}\PYG{n}{fexternal}\PYG{p}{)}\PYG{p}{;}
\end{sphinxVerbatim}

\sphinxAtStartPar
The arguments are as follows:
\begin{itemize}
\item {} 
\sphinxAtStartPar
\sphinxstyleemphasis{ptr} = pointer provided by and simply passed back to external driver

\item {} 
\sphinxAtStartPar
\sphinxstyleemphasis{timestep} = current LAMMPS timestep

\item {} 
\sphinxAtStartPar
\sphinxstyleemphasis{nlocal} = \# of atoms on this processor

\item {} 
\sphinxAtStartPar
\sphinxstyleemphasis{ids} = list of atom IDs on this processor

\item {} 
\sphinxAtStartPar
\sphinxstyleemphasis{x} = coordinates of atoms on this processor

\item {} 
\sphinxAtStartPar
\sphinxstyleemphasis{fexternal} = forces to add to atoms on this processor

\end{itemize}

\sphinxAtStartPar
Note that \sphinxstyleemphasis{timestep} is a “bigint” which is defined in src/lmptype.h,
typically as a 64\sphinxhyphen{}bit integer. And \sphinxstyleemphasis{ids} is a pointer to type “tagint”
which is typically a 32\sphinxhyphen{}bit integer unless LAMMPS is compiled with
\sphinxhyphen{}DLAMMPS\_BIGBIG. For more info please see the {\hyperref[\detokenize{Build_settings:size}]{\sphinxcrossref{\DUrole{std,std-ref}{build settings}}}} section of the manual.  Finally, \sphinxstyleemphasis{fexternal} are the forces
returned by the driver program.

\sphinxAtStartPar
The fix has a set\_callback() method which the external driver can call
to pass a pointer to its foo() function.  See the
couple/lammps\_quest/lmpqst.cpp file in the LAMMPS distribution for an
example of how this is done.  This sample application performs
classical MD using quantum forces computed by a density functional
code \sphinxhref{https://dft.sandia.gov/Quest}{Quest}.


\bigskip\hrule\bigskip


\sphinxAtStartPar
If mode is \sphinxstyleemphasis{pf/array} then the fix simply stores force values in an
array.  The fix adds these forces to each atom in the group, once
every \sphinxstyleemphasis{Napply} steps, similar to the way the {\hyperref[\detokenize{fix_addforce::doc}]{\sphinxcrossref{\DUrole{doc}{fix addforce}}}} command works.

\sphinxAtStartPar
The name of the public force array provided by the FixExternal
class is

\begin{sphinxVerbatim}[commandchars=\\\{\}]
\PYG{k+kt}{double}\PYG{+w}{ }\PYG{o}{*}\PYG{o}{*}\PYG{n}{fexternal}\PYG{p}{;}
\end{sphinxVerbatim}

\sphinxAtStartPar
It is allocated by the FixExternal class as an (N,3) array where N is
the number of atoms owned by a processor.  The 3 corresponds to the
fx, fy, fz components of force.

\sphinxAtStartPar
It is up to the external program to set the values in this array to
the desired quantities, as often as desired.  For example, the driver
program might perform an MD run in stages of 1000 timesteps each.  In
between calls to the LAMMPS {\hyperref[\detokenize{run::doc}]{\sphinxcrossref{\DUrole{doc}{run}}}} command, it could retrieve
atom coordinates from LAMMPS, compute forces, set values in fexternal,
etc.


\bigskip\hrule\bigskip


\sphinxAtStartPar
To use this fix during energy minimization, the energy corresponding
to the added forces must also be set so as to be consistent with the
added forces.  Otherwise the minimization will not converge correctly.
Correspondingly, the global virial needs to be updated to be use this
fix with variable cell calculations (e.g. {\hyperref[\detokenize{fix_box_relax::doc}]{\sphinxcrossref{\DUrole{doc}{fix box/relax}}}}
or {\hyperref[\detokenize{fix_nh::doc}]{\sphinxcrossref{\DUrole{doc}{fix npt}}}}).

\sphinxAtStartPar
This can be done from the external driver by calling these public
methods of the FixExternal class:

\begin{sphinxVerbatim}[commandchars=\\\{\}]
\PYG{k+kt}{void}\PYG{+w}{ }\PYG{n+nf}{set\PYGZus{}energy\PYGZus{}global}\PYG{p}{(}\PYG{k+kt}{double}\PYG{+w}{ }\PYG{n}{eng}\PYG{p}{)}\PYG{p}{;}
\PYG{k+kt}{void}\PYG{+w}{ }\PYG{n+nf}{set\PYGZus{}virial\PYGZus{}global}\PYG{p}{(}\PYG{k+kt}{double}\PYG{+w}{ }\PYG{o}{*}\PYG{n}{virial}\PYG{p}{)}\PYG{p}{;}
\end{sphinxVerbatim}

\sphinxAtStartPar
where \sphinxstyleemphasis{eng} is the potential energy, and \sphinxstyleemphasis{virial} an array of the 6
stress tensor components.  Eng is an extensive quantity,
meaning it should be the sum over per\sphinxhyphen{}atom energies of all affected
atoms.  It should also be provided in {\hyperref[\detokenize{units::doc}]{\sphinxcrossref{\DUrole{doc}{energy units}}}}
consistent with the simulation.  See the details below for how to
ensure this energy setting is used appropriately in a minimization.

\sphinxAtStartPar
Additional public methods that the caller can use to update system
properties are:

\begin{sphinxVerbatim}[commandchars=\\\{\}]
\PYG{k+kt}{void}\PYG{+w}{ }\PYG{n+nf}{set\PYGZus{}energy\PYGZus{}peratom}\PYG{p}{(}\PYG{k+kt}{double}\PYG{+w}{ }\PYG{o}{*}\PYG{n}{eng}\PYG{p}{)}\PYG{p}{;}
\PYG{k+kt}{void}\PYG{+w}{ }\PYG{n+nf}{set\PYGZus{}virial\PYGZus{}peratom}\PYG{p}{(}\PYG{k+kt}{double}\PYG{+w}{ }\PYG{o}{*}\PYG{o}{*}\PYG{n}{virial}\PYG{p}{)}\PYG{p}{;}
\PYG{k+kt}{void}\PYG{+w}{ }\PYG{n+nf}{set\PYGZus{}vector\PYGZus{}length}\PYG{p}{(}\PYG{k+kt}{int}\PYG{+w}{ }\PYG{n}{n}\PYG{p}{)}\PYG{p}{;}
\PYG{k+kt}{void}\PYG{+w}{ }\PYG{n+nf}{set\PYGZus{}vector}\PYG{p}{(}\PYG{k+kt}{int}\PYG{+w}{ }\PYG{n}{idx}\PYG{p}{,}\PYG{+w}{ }\PYG{k+kt}{double}\PYG{+w}{ }\PYG{n}{val}\PYG{p}{)}\PYG{p}{;}
\end{sphinxVerbatim}

\sphinxAtStartPar
These allow to set per\sphinxhyphen{}atom energy contributions, per\sphinxhyphen{}atom stress
contributions, the length and individual values of a global vector
of properties that the caller code may want to communicate  to LAMMPS
(e.g. for use in {\hyperref[\detokenize{fix_ave_time::doc}]{\sphinxcrossref{\DUrole{doc}{fix ave/time}}}} or in
{\hyperref[\detokenize{variable::doc}]{\sphinxcrossref{\DUrole{doc}{equal\sphinxhyphen{}style variables}}}} or for
{\hyperref[\detokenize{thermo_style::doc}]{\sphinxcrossref{\DUrole{doc}{custom thermo output}}}}.


\bigskip\hrule\bigskip



\subsection{Restart, fix\_modify, output, run start/stop, minimize info}
\label{\detokenize{fix_external:restart-fix-modify-output-run-start-stop-minimize-info}}
\sphinxAtStartPar
No information about this fix is written to {\hyperref[\detokenize{restart::doc}]{\sphinxcrossref{\DUrole{doc}{binary restart files}}}}.

\sphinxAtStartPar
The {\hyperref[\detokenize{fix_modify::doc}]{\sphinxcrossref{\DUrole{doc}{fix\_modify}}}} \sphinxstyleemphasis{energy} option is supported by
this fix to add the potential energy set by the external driver to
both the global potential energy and peratom potential energies of the
system as part of {\hyperref[\detokenize{thermo_style::doc}]{\sphinxcrossref{\DUrole{doc}{thermodynamic output}}}} or output
by the {\hyperref[\detokenize{compute_pe_atom::doc}]{\sphinxcrossref{\DUrole{doc}{compute pe/atom}}}} command.  The default
setting for this fix is {\hyperref[\detokenize{fix_modify::doc}]{\sphinxcrossref{\DUrole{doc}{fix\_modify energy yes}}}}.
Note that this energy may be a fictitious quantity but it is needed so
that the {\hyperref[\detokenize{minimize::doc}]{\sphinxcrossref{\DUrole{doc}{minimize}}}} command can include the forces
added by this fix in a consistent manner.  I.e. there is a decrease in
potential energy when atoms move in the direction of the added force.

\sphinxAtStartPar
The {\hyperref[\detokenize{fix_modify::doc}]{\sphinxcrossref{\DUrole{doc}{fix\_modify}}}} \sphinxstyleemphasis{virial} option is supported by
this fix to add the contribution computed by the external program to
both the global pressure and per\sphinxhyphen{}atom stress of the system via the
{\hyperref[\detokenize{compute_pressure::doc}]{\sphinxcrossref{\DUrole{doc}{compute pressure}}}} and {\hyperref[\detokenize{compute_stress_atom::doc}]{\sphinxcrossref{\DUrole{doc}{compute
stress/atom}}}} commands.  The former can be
accessed by {\hyperref[\detokenize{thermo_style::doc}]{\sphinxcrossref{\DUrole{doc}{thermodynamic output}}}}.  The default
setting for this fix is {\hyperref[\detokenize{fix_modify::doc}]{\sphinxcrossref{\DUrole{doc}{fix\_modify virial yes}}}}.

\sphinxAtStartPar
This fix computes a global scalar which can be accessed by various
{\hyperref[\detokenize{Howto_output::doc}]{\sphinxcrossref{\DUrole{doc}{output commands}}}}.  The scalar is the potential
energy discussed above.  The scalar stored by this fix is “extensive”.

\sphinxAtStartPar
No parameter of this fix can be used with the \sphinxstyleemphasis{start/stop} keywords of
the {\hyperref[\detokenize{run::doc}]{\sphinxcrossref{\DUrole{doc}{run}}}} command.

\sphinxAtStartPar
The forces due to this fix are imposed during an energy minimization,
invoked by the {\hyperref[\detokenize{minimize::doc}]{\sphinxcrossref{\DUrole{doc}{minimize}}}} command.

\begin{sphinxadmonition}{note}{Note:}
\sphinxAtStartPar
If you want the fictitious potential energy associated with the
added forces to be included in the total potential energy of the
system (the quantity being minimized), you MUST not disable the
{\hyperref[\detokenize{fix_modify::doc}]{\sphinxcrossref{\DUrole{doc}{fix\_modify}}}} \sphinxstyleemphasis{energy} option for this fix.
\end{sphinxadmonition}


\subsection{Restrictions}
\label{\detokenize{fix_external:restrictions}}\begin{quote}

\sphinxAtStartPar
none
\end{quote}


\subsection{Related commands}
\label{\detokenize{fix_external:related-commands}}
\sphinxAtStartPar
none


\subsection{Default}
\label{\detokenize{fix_external:default}}
\sphinxAtStartPar
none

\sphinxstepscope

\index{fix ffl@\spxentry{fix ffl}}\ignorespaces 

\section{fix ffl command}
\label{\detokenize{fix_ffl:fix-ffl-command}}\label{\detokenize{fix_ffl:index-0}}\label{\detokenize{fix_ffl::doc}}

\subsection{Syntax}
\label{\detokenize{fix_ffl:syntax}}
\begin{sphinxVerbatim}[commandchars=\\\{\}]
\PYG{k}{fix }\PYG{n+nv+nvIdentifier}{ID}\PYG{+w}{ }\PYG{n+nv+nvIdentifier}{id\PYGZhy{}group}\PYG{+w}{ }\PYG{n}{ffl}\PYG{+w}{ }\PYG{n}{tau}\PYG{+w}{ }\PYG{n}{Tstart}\PYG{+w}{ }\PYG{n}{Tstop}\PYG{+w}{ }\PYG{n}{seed}\PYG{+w}{ }\PYG{n}{[flip}\PYG{o}{\PYGZhy{}}\PYG{n}{type]}
\end{sphinxVerbatim}
\begin{itemize}
\item {} 
\sphinxAtStartPar
ID, group\sphinxhyphen{}ID are documented in {\hyperref[\detokenize{fix::doc}]{\sphinxcrossref{\DUrole{doc}{fix}}}} command

\item {} 
\sphinxAtStartPar
ffl = style name of this fix command

\item {} 
\sphinxAtStartPar
tau = thermostat parameter (positive real)

\item {} 
\sphinxAtStartPar
Tstart, Tstop = temperature ramp during the run

\item {} 
\sphinxAtStartPar
seed = random number seed to use for generating noise (positive integer)

\item {} 
\sphinxAtStartPar
one more value may be appended

\begin{sphinxVerbatim}[commandchars=\\\{\}]
flip\PYGZhy{}type  = determines the flipping type, can be chosen between rescale \PYGZhy{} no\PYGZus{}flip \PYGZhy{} hard \PYGZhy{} soft, if no flip type is given, rescale will be chosen by default
\end{sphinxVerbatim}

\end{itemize}


\subsection{Examples}
\label{\detokenize{fix_ffl:examples}}
\begin{sphinxVerbatim}[commandchars=\\\{\}]
\PYG{k}{fix }\PYG{n+nv+nvIdentifier}{3}\PYG{+w}{ }\PYG{n+nv+nvIdentifier}{boundary}\PYG{+w}{ }\PYG{n}{ffl}\PYG{+w}{ }\PYG{l+m}{10}\PYG{+w}{ }\PYG{l+m}{300}\PYG{+w}{ }\PYG{l+m}{300}\PYG{+w}{ }\PYG{l+m}{31415}
\PYG{k}{fix }\PYG{n+nv+nvIdentifier}{1}\PYG{+w}{ }\PYG{n+nv+nvIdentifier}{all}\PYG{+w}{ }\PYG{n}{ffl}\PYG{+w}{ }\PYG{l+m}{100}\PYG{+w}{ }\PYG{l+m}{500}\PYG{+w}{ }\PYG{l+m}{500}\PYG{+w}{ }\PYG{l+m}{9265}\PYG{+w}{ }\PYG{n}{soft}
\end{sphinxVerbatim}


\subsection{Description}
\label{\detokenize{fix_ffl:description}}
\sphinxAtStartPar
Apply a Fast\sphinxhyphen{}Forward Langevin Equation (FFL) thermostat as described
in {\hyperref[\detokenize{fix_ffl:hijazi}]{\sphinxcrossref{\DUrole{std,std-ref}{(Hijazi)}}}}. Contrary to
{\hyperref[\detokenize{fix_langevin::doc}]{\sphinxcrossref{\DUrole{doc}{fix langevin}}}}, this fix performs both
thermostatting and evolution of the Hamiltonian equations of motion, so it
should not be used together with {\hyperref[\detokenize{fix_nve::doc}]{\sphinxcrossref{\DUrole{doc}{fix nve}}}} \textendash{} at least not
on the same atom groups.

\sphinxAtStartPar
The time\sphinxhyphen{}evolution of a single particle undergoing Langevin dynamics is described
by the equations
\begin{equation*}
\begin{split}\frac {dq}{dt} = \frac{p}{m},\end{split}
\end{equation*}\begin{equation*}
\begin{split}\frac {dp}{dt} = -\gamma p + W + F,\end{split}
\end{equation*}
\sphinxAtStartPar
where \(F\) is the physical force, \(\gamma\) is the friction coefficient, and \(W\) is a
Gaussian random force.

\sphinxAtStartPar
The friction coefficient is the inverse of the thermostat parameter : \(\gamma = 1/\tau\), with \(\tau\) the thermostat parameter \sphinxstyleemphasis{tau}.
The thermostat parameter is given in the time units, \(\gamma\) is in inverse time units.

\sphinxAtStartPar
Equilibrium sampling a temperature T is obtained by specifying the
target value as the \sphinxstyleemphasis{Tstart} and \sphinxstyleemphasis{Tstop} arguments, so that the internal
constants depending on the temperature are computed automatically.

\sphinxAtStartPar
The random number \sphinxstyleemphasis{seed} must be a positive integer.  A Marsaglia random
number generator is used.  Each processor uses the input seed to
generate its own unique seed and its own stream of random numbers.
Thus the dynamics of the system will not be identical on two runs on
different numbers of processors.

\sphinxAtStartPar
The flipping type \sphinxstyleemphasis{flip\sphinxhyphen{}type} can be chosen between 4 types described in
{\hyperref[\detokenize{fix_ffl:hijazi}]{\sphinxcrossref{\DUrole{std,std-ref}{(Hijazi)}}}}. The flipping operation occurs during the thermostatting
step and it flips the momenta of the atoms. If no\_flip is chosen, no flip
will be executed and the integration will be the same as a standard
Langevin thermostat {\hyperref[\detokenize{fix_ffl:bussi3}]{\sphinxcrossref{\DUrole{std,std-ref}{(Bussi)}}}}. The other flipping types are : rescale \sphinxhyphen{} hard \sphinxhyphen{} soft.


\subsection{Restart, fix\_modify, output, run start/stop, minimize info}
\label{\detokenize{fix_ffl:restart-fix-modify-output-run-start-stop-minimize-info}}
\sphinxAtStartPar
The instantaneous values of the extended variables are written to
{\hyperref[\detokenize{restart::doc}]{\sphinxcrossref{\DUrole{doc}{binary restart files}}}}.  Because the state of the
random number generator is not saved in restart files, this means you
cannot do “exact” restarts with this fix, where the simulation
continues on the same as if no restart had taken place. However, in a
statistical sense, a restarted simulation should produce the same
behavior.  Note however that you should use a different seed each time
you restart, otherwise the same sequence of random numbers will be
used each time, which might lead to stochastic synchronization and
subtle artifacts in the sampling.

\sphinxAtStartPar
The cumulative energy change in the system imposed by this fix is
included in the {\hyperref[\detokenize{thermo_style::doc}]{\sphinxcrossref{\DUrole{doc}{thermodynamic output}}}} keywords
\sphinxstyleemphasis{ecouple} and \sphinxstyleemphasis{econserve}.  See the {\hyperref[\detokenize{thermo_style::doc}]{\sphinxcrossref{\DUrole{doc}{thermo\_style}}}}
doc page for details.

\sphinxAtStartPar
This fix computes a global scalar which can be accessed by various
{\hyperref[\detokenize{Howto_output::doc}]{\sphinxcrossref{\DUrole{doc}{output commands}}}}.  The scalar is the same
cumulative energy change due to this fix described in the previous
paragraph.  The scalar value calculated by this fix is “extensive”.

\sphinxAtStartPar
This fix can ramp its target temperature over multiple runs, using the
\sphinxstyleemphasis{start} and \sphinxstyleemphasis{stop} keywords of the {\hyperref[\detokenize{run::doc}]{\sphinxcrossref{\DUrole{doc}{run}}}} command.  See the
{\hyperref[\detokenize{run::doc}]{\sphinxcrossref{\DUrole{doc}{run}}}} command for details of how to do this.

\sphinxAtStartPar
This fix is not invoked during {\hyperref[\detokenize{minimize::doc}]{\sphinxcrossref{\DUrole{doc}{energy minimization}}}}.


\subsection{Restrictions}
\label{\detokenize{fix_ffl:restrictions}}
\sphinxAtStartPar
In order to perform constant\sphinxhyphen{}pressure simulations please use
{\hyperref[\detokenize{fix_press_berendsen::doc}]{\sphinxcrossref{\DUrole{doc}{fix press/berendsen}}}}, rather than
{\hyperref[\detokenize{fix_nh::doc}]{\sphinxcrossref{\DUrole{doc}{fix npt}}}}, to avoid duplicate integration of the
equations of motion.

\sphinxAtStartPar
This fix is part of the EXTRA\sphinxhyphen{}FIX package.  It is only enabled if
LAMMPS was built with that package.  See the {\hyperref[\detokenize{Build_package::doc}]{\sphinxcrossref{\DUrole{doc}{Build package}}}} page for more info.


\subsection{Related commands}
\label{\detokenize{fix_ffl:related-commands}}
\sphinxAtStartPar
{\hyperref[\detokenize{fix_nh::doc}]{\sphinxcrossref{\DUrole{doc}{fix nvt}}}}, {\hyperref[\detokenize{fix_temp_rescale::doc}]{\sphinxcrossref{\DUrole{doc}{fix temp/rescale}}}}, {\hyperref[\detokenize{fix_viscous::doc}]{\sphinxcrossref{\DUrole{doc}{fix viscous}}}}, {\hyperref[\detokenize{fix_nh::doc}]{\sphinxcrossref{\DUrole{doc}{fix nvt}}}}, {\hyperref[\detokenize{pair_dpd::doc}]{\sphinxcrossref{\DUrole{doc}{pair\_style dpd/tstat}}}}, {\hyperref[\detokenize{fix_gld::doc}]{\sphinxcrossref{\DUrole{doc}{fix gld}}}}, {\hyperref[\detokenize{fix_gle::doc}]{\sphinxcrossref{\DUrole{doc}{fix gle}}}}


\bigskip\hrule\bigskip

\phantomsection\label{\detokenize{fix_ffl:hijazi}}\phantomsection\label{\detokenize{fix_ffl:bussi3}}
\sphinxAtStartPar
\sphinxstylestrong{(Hijazi)} M. Hijazi, D. M. Wilkins, M. Ceriotti, J. Chem. Phys. 148, 184109 (2018)

\sphinxAtStartPar
\sphinxstylestrong{(Bussi)} G. Bussi, M. Parrinello, Phs. Rev. E 75, 056707 (2007)

\sphinxstepscope

\index{fix filter/corotate@\spxentry{fix filter/corotate}}\ignorespaces 

\section{fix filter/corotate command}
\label{\detokenize{fix_filter_corotate:fix-filter-corotate-command}}\label{\detokenize{fix_filter_corotate:index-0}}\label{\detokenize{fix_filter_corotate::doc}}

\subsection{Syntax}
\label{\detokenize{fix_filter_corotate:syntax}}
\begin{sphinxVerbatim}[commandchars=\\\{\}]
\PYG{k}{fix }\PYG{n+nv+nvIdentifier}{ID}\PYG{+w}{ }\PYG{n+nv+nvIdentifier}{group\PYGZhy{}ID}\PYG{+w}{ }\PYG{n}{filter}\PYG{o}{/}\PYG{n}{corotate}\PYG{+w}{ }\PYG{n}{keyword}\PYG{+w}{ }\PYG{n}{value}\PYG{+w}{ }\PYG{n}{...}
\end{sphinxVerbatim}
\begin{itemize}
\item {} 
\sphinxAtStartPar
ID, group\sphinxhyphen{}ID are documented in {\hyperref[\detokenize{fix::doc}]{\sphinxcrossref{\DUrole{doc}{fix}}}} command

\item {} 
\sphinxAtStartPar
one or more constraint/value pairs are appended

\item {} 
\sphinxAtStartPar
constraint = \sphinxstyleemphasis{b} or \sphinxstyleemphasis{a} or \sphinxstyleemphasis{t} or \sphinxstyleemphasis{m}
\begin{sphinxalltt}
\sphinxstyleemphasis{b} values = one or more bond types
\sphinxstyleemphasis{a} values = one or more angle types
\sphinxstyleemphasis{t} values = one or more atom types
\sphinxstyleemphasis{m} value = one or more mass values
\end{sphinxalltt}

\end{itemize}


\subsection{Examples}
\label{\detokenize{fix_filter_corotate:examples}}
\begin{sphinxVerbatim}[commandchars=\\\{\}]
\PYG{k}{timestep}\PYG{+w}{ }\PYG{l+m}{8}
\PYG{k}{run\PYGZus{}style}\PYG{+w}{ }\PYG{n}{respa}\PYG{+w}{ }\PYG{l+m}{3}\PYG{+w}{ }\PYG{l+m}{2}\PYG{+w}{ }\PYG{l+m}{8}\PYG{+w}{ }\PYG{n}{bond}\PYG{+w}{ }\PYG{l+m}{1}\PYG{+w}{ }\PYG{n}{pair}\PYG{+w}{ }\PYG{l+m}{2}\PYG{+w}{ }\PYG{n}{kspace}\PYG{+w}{ }\PYG{l+m}{3}
\PYG{k}{fix }\PYG{n+nv+nvIdentifier}{cor}\PYG{+w}{ }\PYG{n+nv+nvIdentifier}{all}\PYG{+w}{ }\PYG{n}{filter}\PYG{o}{/}\PYG{n}{corotate}\PYG{+w}{ }\PYG{n}{m}\PYG{+w}{ }\PYG{l+m}{1.0}

\PYG{k}{fix }\PYG{n+nv+nvIdentifier}{cor}\PYG{+w}{ }\PYG{n+nv+nvIdentifier}{all}\PYG{+w}{ }\PYG{n}{filter}\PYG{o}{/}\PYG{n}{corotate}\PYG{+w}{ }\PYG{n}{b}\PYG{+w}{ }\PYG{l+m}{4}\PYG{+w}{ }\PYG{l+m}{19}\PYG{+w}{ }\PYG{n}{a}\PYG{+w}{ }\PYG{l+m}{3}\PYG{+w}{ }\PYG{l+m}{5}\PYG{+w}{ }\PYG{l+m}{2}
\end{sphinxVerbatim}


\subsection{Description}
\label{\detokenize{fix_filter_corotate:description}}
\sphinxAtStartPar
This fix implements a corotational filter for a mollified impulse
method. In biomolecular simulations, it allows the usage of larger
timesteps for long\sphinxhyphen{}range electrostatic interactions.  For details, see
{\hyperref[\detokenize{fix_filter_corotate:fath2017}]{\sphinxcrossref{\DUrole{std,std-ref}{(Fath)}}}}.

\sphinxAtStartPar
When using {\hyperref[\detokenize{run_style::doc}]{\sphinxcrossref{\DUrole{doc}{run\_style respa}}}} for a biomolecular
simulation with high\sphinxhyphen{}frequency covalent bonds, the outer time\sphinxhyphen{}step is
restricted to below \textasciitilde{} 4fs due to resonance problems. This fix filters
the outer stage of the respa and thus a larger (outer) time\sphinxhyphen{}step can
be used. Since in large biomolecular simulations the computation of
the long\sphinxhyphen{}range electrostatic contributions poses a major bottleneck,
this can significantly accelerate the simulation.

\sphinxAtStartPar
The filter computes a cluster decomposition of the molecular structure
following the criteria indicated by the options a, b, t and m. This
process is similar to the approach in {\hyperref[\detokenize{fix_shake::doc}]{\sphinxcrossref{\DUrole{doc}{fix shake}}}},
however, the clusters are not kept constrained. Instead, the position
is slightly modified only for the computation of long\sphinxhyphen{}range forces. A
good cluster decomposition constitutes in building clusters which
contain the fastest covalent bonds inside clusters.

\sphinxAtStartPar
If the clusters are chosen suitably, the {\hyperref[\detokenize{run_style::doc}]{\sphinxcrossref{\DUrole{doc}{run\_style respa}}}} is stable for outer timesteps of at least 8fs.


\bigskip\hrule\bigskip



\subsection{Restart, fix\_modify, output, run start/stop, minimize info}
\label{\detokenize{fix_filter_corotate:restart-fix-modify-output-run-start-stop-minimize-info}}
\sphinxAtStartPar
No information about these fixes is written to {\hyperref[\detokenize{restart::doc}]{\sphinxcrossref{\DUrole{doc}{binary restart files}}}}.  None of the {\hyperref[\detokenize{fix_modify::doc}]{\sphinxcrossref{\DUrole{doc}{fix\_modify}}}} options
are relevant to these fixes.  No global or per\sphinxhyphen{}atom quantities are
stored by these fixes for access by various {\hyperref[\detokenize{Howto_output::doc}]{\sphinxcrossref{\DUrole{doc}{output commands}}}}.  No parameter of these fixes can be used
with the \sphinxstyleemphasis{start/stop} keywords of the {\hyperref[\detokenize{run::doc}]{\sphinxcrossref{\DUrole{doc}{run}}}} command.  These
fixes are not invoked during {\hyperref[\detokenize{minimize::doc}]{\sphinxcrossref{\DUrole{doc}{energy minimization}}}}.


\subsection{Restrictions}
\label{\detokenize{fix_filter_corotate:restrictions}}
\sphinxAtStartPar
This fix is part of the EXTRA\sphinxhyphen{}FIX package. It is only enabled if
LAMMPS was built with that package. See the {\hyperref[\detokenize{Build_package::doc}]{\sphinxcrossref{\DUrole{doc}{Build package}}}} page for more info.

\sphinxAtStartPar
Currently, it does not support {\hyperref[\detokenize{molecule::doc}]{\sphinxcrossref{\DUrole{doc}{molecule templates}}}}.


\subsection{Related commands}
\label{\detokenize{fix_filter_corotate:related-commands}}

\subsection{Default}
\label{\detokenize{fix_filter_corotate:default}}
\sphinxAtStartPar
none


\bigskip\hrule\bigskip

\phantomsection\label{\detokenize{fix_filter_corotate:fath2017}}
\sphinxAtStartPar
\sphinxstylestrong{(Fath)} Fath, Hochbruck, Singh, J Comp Phys, 333, 180\sphinxhyphen{}198 (2017).

\sphinxstepscope

\index{fix flow/gauss@\spxentry{fix flow/gauss}}\ignorespaces 

\section{fix flow/gauss command}
\label{\detokenize{fix_flow_gauss:fix-flow-gauss-command}}\label{\detokenize{fix_flow_gauss:index-0}}\label{\detokenize{fix_flow_gauss::doc}}

\subsection{Syntax}
\label{\detokenize{fix_flow_gauss:syntax}}
\begin{sphinxVerbatim}[commandchars=\\\{\}]
\PYG{k}{fix }\PYG{n+nv+nvIdentifier}{ID}\PYG{+w}{ }\PYG{n+nv+nvIdentifier}{group\PYGZhy{}ID}\PYG{+w}{ }\PYG{n}{flow}\PYG{o}{/}\PYG{n}{gauss}\PYG{+w}{ }\PYG{n}{xflag}\PYG{+w}{ }\PYG{n}{yflag}\PYG{+w}{ }\PYG{n}{zflag}\PYG{+w}{ }\PYG{n}{keyword}
\end{sphinxVerbatim}
\begin{itemize}
\item {} 
\sphinxAtStartPar
ID, group\sphinxhyphen{}ID are documented in {\hyperref[\detokenize{fix::doc}]{\sphinxcrossref{\DUrole{doc}{fix}}}} command

\item {} 
\sphinxAtStartPar
flow/gauss = style name of this fix command

\item {} 
\sphinxAtStartPar
xflag,yflag,zflag = 0 or 1

\begin{sphinxVerbatim}[commandchars=\\\{\}]
0 = do not conserve current in this dimension
1 = conserve current in this dimension
\end{sphinxVerbatim}

\item {} 
\sphinxAtStartPar
zero or more keyword/value pairs may be appended

\item {} 
\sphinxAtStartPar
keyword = \sphinxstyleemphasis{energy}
\begin{sphinxalltt}
\sphinxstyleemphasis{energy} value = \sphinxstyleemphasis{no} or \sphinxstyleemphasis{yes}
  \sphinxstyleemphasis{no} = do not compute work done by this fix
  \sphinxstyleemphasis{yes} = compute work done by this fix
\end{sphinxalltt}

\end{itemize}


\subsection{Examples}
\label{\detokenize{fix_flow_gauss:examples}}
\begin{sphinxVerbatim}[commandchars=\\\{\}]
\PYG{k}{fix }\PYG{n+nv+nvIdentifier}{GD}\PYG{+w}{ }\PYG{n+nv+nvIdentifier}{fluid}\PYG{+w}{ }\PYG{n}{flow}\PYG{o}{/}\PYG{n}{gauss}\PYG{+w}{ }\PYG{l+m}{1}\PYG{+w}{ }\PYG{l+m}{0}\PYG{+w}{ }\PYG{l+m}{0}
\PYG{k}{fix }\PYG{n+nv+nvIdentifier}{GD}\PYG{+w}{ }\PYG{n+nv+nvIdentifier}{fluid}\PYG{+w}{ }\PYG{n}{flow}\PYG{o}{/}\PYG{n}{gauss}\PYG{+w}{ }\PYG{l+m}{1}\PYG{+w}{ }\PYG{l+m}{1}\PYG{+w}{ }\PYG{l+m}{1}\PYG{+w}{ }\PYG{n}{energy}\PYG{+w}{ }\PYG{n}{yes}
\end{sphinxVerbatim}


\subsection{Description}
\label{\detokenize{fix_flow_gauss:description}}
\sphinxAtStartPar
This fix implements the Gaussian dynamics (GD) method to simulate a
system at constant mass flux {\hyperref[\detokenize{fix_flow_gauss:strong}]{\sphinxcrossref{\DUrole{std,std-ref}{(Strong)}}}}. GD is a
nonequilibrium molecular dynamics simulation method that can be used
to study fluid flows through pores, pipes, and channels. In its
original implementation GD was used to compute the pressure required
to achieve a fixed mass flux through an opening.  The flux can be
conserved in any combination of the directions, x, y, or z, using
xflag,yflag,zflag. This fix does not initialize a net flux through a
system, it only conserves the center\sphinxhyphen{}of\sphinxhyphen{}mass momentum that is present
when the fix is declared in the input script. Use the
{\hyperref[\detokenize{velocity::doc}]{\sphinxcrossref{\DUrole{doc}{velocity}}}} command to generate an initial center\sphinxhyphen{}of\sphinxhyphen{}mass
momentum.

\sphinxAtStartPar
GD applies an external fluctuating gravitational field that acts as a
driving force to keep the system away from equilibrium. To maintain
steady state, a profile\sphinxhyphen{}unbiased thermostat must be implemented to
dissipate the heat that is added by the driving force. {\hyperref[\detokenize{compute_temp_profile::doc}]{\sphinxcrossref{\DUrole{doc}{Compute
temp/profile}}}} can be used to implement a
profile\sphinxhyphen{}unbiased thermostat.

\sphinxAtStartPar
A common use of this fix is to compute a pressure drop across a pipe,
pore, or membrane. The pressure profile can be computed in LAMMPS with
{\hyperref[\detokenize{compute_stress_atom::doc}]{\sphinxcrossref{\DUrole{doc}{compute stress/atom}}}} and {\hyperref[\detokenize{fix_ave_chunk::doc}]{\sphinxcrossref{\DUrole{doc}{fix
ave/chunk}}}}, or with the hardy method in {\hyperref[\detokenize{fix_atc::doc}]{\sphinxcrossref{\DUrole{doc}{fix atc}}}}. Note that the simple {\hyperref[\detokenize{compute_stress_atom::doc}]{\sphinxcrossref{\DUrole{doc}{compute stress/atom}}}} method is only accurate away from
inhomogeneities in the fluid, such as fixed wall atoms. Further, the
computed pressure profile must be corrected for the acceleration
applied by GD before computing a pressure drop or comparing it to
other methods, such as the pump method {\hyperref[\detokenize{fix_flow_gauss:zhu}]{\sphinxcrossref{\DUrole{std,std-ref}{(Zhu)}}}}. The
pressure correction is discussed and described in {\hyperref[\detokenize{fix_flow_gauss:strong}]{\sphinxcrossref{\DUrole{std,std-ref}{(Strong)}}}}.

\sphinxAtStartPar
For a complete example including the considerations discussed
above, see the examples/PACKAGES/flow\_gauss directory.

\begin{sphinxadmonition}{note}{Note:}
\sphinxAtStartPar
Only the flux of the atoms in group\sphinxhyphen{}ID will be conserved. If the
velocities of the group\sphinxhyphen{}ID atoms are coupled to the velocities of
other atoms in the simulation, the flux will not be conserved. For
example, in a simulation with fluid atoms and harmonically constrained
wall atoms, if a single thermostat is applied to group \sphinxstyleemphasis{all}, the
fluid atom velocities will be coupled to the wall atom velocities, and
the flux will not be conserved. This issue can be avoided by
thermostatting the fluid and wall groups separately.
\end{sphinxadmonition}

\sphinxAtStartPar
Adding an acceleration to atoms does work on the system. This added
energy can be optionally subtracted from the potential energy for the
thermodynamic output (see below) to check that the timestep is small
enough to conserve energy. Since the applied acceleration is
fluctuating in time, the work cannot be computed from a potential. As
a result, computing the work is slightly more computationally
expensive than usual, so it is not performed by default. To invoke the
work calculation, use the \sphinxstyleemphasis{energy} keyword. The
{\hyperref[\detokenize{fix_modify::doc}]{\sphinxcrossref{\DUrole{doc}{fix\_modify}}}} \sphinxstyleemphasis{energy} option also invokes the work
calculation, and overrides an \sphinxstyleemphasis{energy no} setting here. If neither
\sphinxstyleemphasis{energy yes} or \sphinxstyleemphasis{fix\_modify energy yes} are set, the global scalar
computed by the fix will return zero.

\begin{sphinxadmonition}{note}{Note:}
\sphinxAtStartPar
In order to check energy conservation, any other fixes that do
work on the system must have \sphinxstyleemphasis{fix\_modify energy yes} set as well. This
includes thermostat fixes and any constraints that hold the positions
of wall atoms fixed, such as {\hyperref[\detokenize{fix_spring_self::doc}]{\sphinxcrossref{\DUrole{doc}{fix spring/self}}}}.
\end{sphinxadmonition}

\sphinxAtStartPar
If this fix is used in a simulation with the {\hyperref[\detokenize{run_style::doc}]{\sphinxcrossref{\DUrole{doc}{rRESPA}}}}
integrator, the applied acceleration must be computed and applied at
the same rRESPA level as the interactions between the flowing fluid
and the obstacle.  The rRESPA level at which the acceleration is
applied can be changed using the {\hyperref[\detokenize{fix_modify::doc}]{\sphinxcrossref{\DUrole{doc}{fix\_modify}}}}
\sphinxstyleemphasis{respa} option discussed below. If the flowing fluid and the obstacle
interact through multiple interactions that are computed at different
rRESPA levels, then there must be a separate flow/gauss fix for each
level. For example, if the flowing fluid and obstacle interact through
pairwise and long\sphinxhyphen{}range Coulomb interactions, which are computed at
rRESPA levels 3 and 4, respectively, then there must be two separate
flow/gauss fixes, one that specifies \sphinxstyleemphasis{fix\_modify respa 3} and one with
\sphinxstyleemphasis{fix\_modify respa 4}.


\bigskip\hrule\bigskip



\subsection{Restart, fix\_modify, output, run start/stop, minimize info}
\label{\detokenize{fix_flow_gauss:restart-fix-modify-output-run-start-stop-minimize-info}}
\sphinxAtStartPar
No information about this fix is written to {\hyperref[\detokenize{restart::doc}]{\sphinxcrossref{\DUrole{doc}{binary restart files}}}}.

\sphinxAtStartPar
The {\hyperref[\detokenize{fix_modify::doc}]{\sphinxcrossref{\DUrole{doc}{fix\_modify}}}} \sphinxstyleemphasis{energy} option is supported by
this fix to add the potential energy added by the fix to the global
potential energy of the system as part of {\hyperref[\detokenize{thermo_style::doc}]{\sphinxcrossref{\DUrole{doc}{thermodynamic output}}}}.  The default setting for this fix is {\hyperref[\detokenize{fix_modify::doc}]{\sphinxcrossref{\DUrole{doc}{fix\_modify
energy no}}}}.

\sphinxAtStartPar
The {\hyperref[\detokenize{fix_modify::doc}]{\sphinxcrossref{\DUrole{doc}{fix\_modify}}}} \sphinxstyleemphasis{respa} option is supported by this
fix. This allows the user to set at which level of the {\hyperref[\detokenize{run_style::doc}]{\sphinxcrossref{\DUrole{doc}{rRESPA}}}} integrator the fix computes and adds the external
acceleration. Default is the outermost level.

\sphinxAtStartPar
This fix computes a global scalar and a global 3\sphinxhyphen{}vector of forces,
which can be accessed by various {\hyperref[\detokenize{Howto_output::doc}]{\sphinxcrossref{\DUrole{doc}{output commands}}}}.  The scalar is the negative of the work done on the
system, see the discussion above.  It is only calculated if the
\sphinxstyleemphasis{energy} keyword is enabled or {\hyperref[\detokenize{fix_modify::doc}]{\sphinxcrossref{\DUrole{doc}{fix\_modify energy yes}}}} is set.

\sphinxAtStartPar
The vector is the total force that this fix applied to the group of
atoms on the current timestep.  The scalar and vector values
calculated by this fix are “extensive”.

\sphinxAtStartPar
No parameter of this fix can be used with the \sphinxstyleemphasis{start/stop} keywords of
the {\hyperref[\detokenize{run::doc}]{\sphinxcrossref{\DUrole{doc}{run}}}} command.

\sphinxAtStartPar
This fix is not invoked during {\hyperref[\detokenize{minimize::doc}]{\sphinxcrossref{\DUrole{doc}{energy minimization}}}}.


\subsection{Restrictions}
\label{\detokenize{fix_flow_gauss:restrictions}}
\sphinxAtStartPar
This fix is part of the EXTRA\sphinxhyphen{}FIX package.  It is only enabled if
LAMMPS was built with that package.  See the {\hyperref[\detokenize{Build_package::doc}]{\sphinxcrossref{\DUrole{doc}{Build package}}}} page for more info.


\subsection{Related commands}
\label{\detokenize{fix_flow_gauss:related-commands}}
\sphinxAtStartPar
{\hyperref[\detokenize{fix_addforce::doc}]{\sphinxcrossref{\DUrole{doc}{fix addforce}}}},
{\hyperref[\detokenize{compute_temp_profile::doc}]{\sphinxcrossref{\DUrole{doc}{compute temp/profile}}}},
{\hyperref[\detokenize{velocity::doc}]{\sphinxcrossref{\DUrole{doc}{velocity}}}}


\subsection{Default}
\label{\detokenize{fix_flow_gauss:default}}
\sphinxAtStartPar
The option default for the \sphinxstyleemphasis{energy} keyword is energy = no.


\bigskip\hrule\bigskip

\phantomsection\label{\detokenize{fix_flow_gauss:strong}}
\sphinxAtStartPar
\sphinxstylestrong{(Strong)} Strong and Eaves, J. Phys. Chem. B 121, 189 (2017).

\phantomsection\label{\detokenize{fix_flow_gauss:evans2}}
\sphinxAtStartPar
\sphinxstylestrong{(Evans)} Evans and Morriss, Phys. Rev. Lett. 56, 2172 (1986).

\phantomsection\label{\detokenize{fix_flow_gauss:zhu}}
\sphinxAtStartPar
\sphinxstylestrong{(Zhu)} Zhu, Tajkhorshid, and Schulten, Biophys. J. 83, 154 (2002).

\sphinxstepscope

\index{fix freeze@\spxentry{fix freeze}}\ignorespaces 
\index{fix freeze/kk@\spxentry{fix freeze/kk}}\ignorespaces 

\section{fix freeze command}
\label{\detokenize{fix_freeze:fix-freeze-command}}\label{\detokenize{fix_freeze:index-1}}\label{\detokenize{fix_freeze:index-0}}\label{\detokenize{fix_freeze::doc}}
\sphinxAtStartPar
Accelerator Variants: \sphinxstyleemphasis{freeze/kk}


\subsection{Syntax}
\label{\detokenize{fix_freeze:syntax}}
\begin{sphinxVerbatim}[commandchars=\\\{\}]
\PYG{k}{fix }\PYG{n+nv+nvIdentifier}{ID}\PYG{+w}{ }\PYG{n+nv+nvIdentifier}{group\PYGZhy{}ID}\PYG{+w}{ }\PYG{n}{freeze}
\end{sphinxVerbatim}
\begin{itemize}
\item {} 
\sphinxAtStartPar
ID, group\sphinxhyphen{}ID are documented in {\hyperref[\detokenize{fix::doc}]{\sphinxcrossref{\DUrole{doc}{fix}}}} command

\item {} 
\sphinxAtStartPar
freeze = style name of this fix command

\end{itemize}


\subsection{Examples}
\label{\detokenize{fix_freeze:examples}}
\begin{sphinxVerbatim}[commandchars=\\\{\}]
\PYG{k}{fix }\PYG{n+nv+nvIdentifier}{2}\PYG{+w}{ }\PYG{n+nv+nvIdentifier}{bottom}\PYG{+w}{ }\PYG{n}{freeze}
\end{sphinxVerbatim}


\subsection{Description}
\label{\detokenize{fix_freeze:description}}
\sphinxAtStartPar
Zero out the force and torque on a granular particle.  This is useful
for preventing certain particles from moving in a simulation.  The
{\hyperref[\detokenize{pair_gran::doc}]{\sphinxcrossref{\DUrole{doc}{granular pair styles}}}} also detect if this fix has been
defined and compute interactions between frozen and non\sphinxhyphen{}frozen
particles appropriately, as if the frozen particle has infinite mass.
A similar functionality for normal (point) particles can be obtained
using {\hyperref[\detokenize{fix_setforce::doc}]{\sphinxcrossref{\DUrole{doc}{fix setforce}}}}.


\bigskip\hrule\bigskip


\sphinxAtStartPar
Styles with a \sphinxstyleemphasis{gpu}, \sphinxstyleemphasis{intel}, \sphinxstyleemphasis{kk}, \sphinxstyleemphasis{omp}, or \sphinxstyleemphasis{opt} suffix are
functionally the same as the corresponding style without the suffix.
They have been optimized to run faster, depending on your available
hardware, as discussed on the {\hyperref[\detokenize{Speed_packages::doc}]{\sphinxcrossref{\DUrole{doc}{Accelerator packages}}}}
page.  The accelerated styles take the same arguments and should
produce the same results, except for round\sphinxhyphen{}off and precision issues.

\sphinxAtStartPar
These accelerated styles are part of the GPU, INTEL, KOKKOS,
OPENMP, and OPT packages, respectively.  They are only enabled if
LAMMPS was built with those packages.  See the {\hyperref[\detokenize{Build_package::doc}]{\sphinxcrossref{\DUrole{doc}{Build package}}}} page for more info.

\sphinxAtStartPar
You can specify the accelerated styles explicitly in your input script
by including their suffix, or you can use the {\hyperref[\detokenize{Run_options::doc}]{\sphinxcrossref{\DUrole{doc}{\sphinxhyphen{}suffix command\sphinxhyphen{}line switch}}}} when you invoke LAMMPS, or you can use the
{\hyperref[\detokenize{suffix::doc}]{\sphinxcrossref{\DUrole{doc}{suffix}}}} command in your input script.

\sphinxAtStartPar
See the {\hyperref[\detokenize{Speed_packages::doc}]{\sphinxcrossref{\DUrole{doc}{Accelerator packages}}}} page for more
instructions on how to use the accelerated styles effectively.


\bigskip\hrule\bigskip



\subsection{Restart, fix\_modify, output, run start/stop, minimize info}
\label{\detokenize{fix_freeze:restart-fix-modify-output-run-start-stop-minimize-info}}
\sphinxAtStartPar
No information about this fix is written to {\hyperref[\detokenize{restart::doc}]{\sphinxcrossref{\DUrole{doc}{binary restart files}}}}.  None of the {\hyperref[\detokenize{fix_modify::doc}]{\sphinxcrossref{\DUrole{doc}{fix\_modify}}}} options
are relevant to this fix.

\sphinxAtStartPar
This fix computes a global 3\sphinxhyphen{}vector of forces, which can be accessed
by various {\hyperref[\detokenize{Howto_output::doc}]{\sphinxcrossref{\DUrole{doc}{output commands}}}}.  This is the total
force on the group of atoms before the forces on individual atoms are
changed by the fix.  The vector values calculated by this fix are
“extensive”.

\sphinxAtStartPar
No parameter of this fix can be used with the \sphinxstyleemphasis{start/stop} keywords of
the {\hyperref[\detokenize{run::doc}]{\sphinxcrossref{\DUrole{doc}{run}}}} command.  This fix is not invoked during {\hyperref[\detokenize{minimize::doc}]{\sphinxcrossref{\DUrole{doc}{energy minimization}}}}.


\subsection{Restrictions}
\label{\detokenize{fix_freeze:restrictions}}
\sphinxAtStartPar
This fix is part of the GRANULAR package.  It is only enabled if
LAMMPS was built with that package.  See the {\hyperref[\detokenize{Build_package::doc}]{\sphinxcrossref{\DUrole{doc}{Build package}}}} page for more info.

\sphinxAtStartPar
There can only be a single freeze fix defined.  This is because other
the {\hyperref[\detokenize{pair_gran::doc}]{\sphinxcrossref{\DUrole{doc}{granular pair styles}}}} treat frozen particles
differently and need to be able to reference a single group to which
this fix is applied.


\subsection{Related commands}
\label{\detokenize{fix_freeze:related-commands}}
\sphinxAtStartPar
{\hyperref[\detokenize{atom_style::doc}]{\sphinxcrossref{\DUrole{doc}{atom\_style sphere}}}}, {\hyperref[\detokenize{fix_setforce::doc}]{\sphinxcrossref{\DUrole{doc}{fix setforce}}}}


\subsection{Default}
\label{\detokenize{fix_freeze:default}}
\sphinxAtStartPar
none

\sphinxstepscope

\index{fix gcmc@\spxentry{fix gcmc}}\ignorespaces 

\section{fix gcmc command}
\label{\detokenize{fix_gcmc:fix-gcmc-command}}\label{\detokenize{fix_gcmc:index-0}}\label{\detokenize{fix_gcmc::doc}}

\subsection{Syntax}
\label{\detokenize{fix_gcmc:syntax}}
\begin{sphinxVerbatim}[commandchars=\\\{\}]
\PYG{k}{fix }\PYG{n+nv+nvIdentifier}{ID}\PYG{+w}{ }\PYG{n+nv+nvIdentifier}{group\PYGZhy{}ID}\PYG{+w}{ }\PYG{n}{gcmc}\PYG{+w}{ }\PYG{n}{N}\PYG{+w}{ }\PYG{n}{X}\PYG{+w}{ }\PYG{n}{M}\PYG{+w}{ }\PYG{n}{type}\PYG{+w}{ }\PYG{n}{seed}\PYG{+w}{ }\PYG{n}{T}\PYG{+w}{ }\PYG{n}{mu}\PYG{+w}{ }\PYG{n}{displace}\PYG{+w}{ }\PYG{n}{keyword}\PYG{+w}{ }\PYG{n}{values}\PYG{+w}{ }\PYG{n}{...}
\end{sphinxVerbatim}
\begin{itemize}
\item {} 
\sphinxAtStartPar
ID, group\sphinxhyphen{}ID are documented in {\hyperref[\detokenize{fix::doc}]{\sphinxcrossref{\DUrole{doc}{fix}}}} command

\item {} 
\sphinxAtStartPar
gcmc = style name of this fix command

\item {} 
\sphinxAtStartPar
N = invoke this fix every N steps

\item {} 
\sphinxAtStartPar
X = average number of GCMC exchanges to attempt every N steps

\item {} 
\sphinxAtStartPar
M = average number of MC moves to attempt every N steps

\item {} 
\sphinxAtStartPar
type = atom type (1\sphinxhyphen{}Ntypes or type label) for inserted atoms (must be 0 if mol keyword used)

\item {} 
\sphinxAtStartPar
seed = random \# seed (positive integer)

\item {} 
\sphinxAtStartPar
T = temperature of the ideal gas reservoir (temperature units)

\item {} 
\sphinxAtStartPar
mu = chemical potential of the ideal gas reservoir (energy units)

\item {} 
\sphinxAtStartPar
displace = maximum Monte Carlo translation distance (length units)

\item {} 
\sphinxAtStartPar
zero or more keyword/value pairs may be appended to args
\begin{sphinxalltt}
keyword = \sphinxstyleemphasis{mol}, \sphinxstyleemphasis{region}, \sphinxstyleemphasis{maxangle}, \sphinxstyleemphasis{pressure}, \sphinxstyleemphasis{fugacity\_coeff}, \sphinxstyleemphasis{full\_energy}, \sphinxstyleemphasis{charge}, \sphinxstyleemphasis{group}, \sphinxstyleemphasis{grouptype}, \sphinxstyleemphasis{intra\_energy}, \sphinxstyleemphasis{tfac\_insert}, or \sphinxstyleemphasis{overlap\_cutoff}
  \sphinxstyleemphasis{mol} value = template\sphinxhyphen{}ID
    template\sphinxhyphen{}ID = ID of molecule template specified in a separate {\hyperref[\detokenize{molecule::doc}]{\sphinxcrossref{\DUrole{doc}{molecule}}}} command
  \sphinxstyleemphasis{mcmoves} values = Patomtrans Pmoltrans Pmolrotate
    Patomtrans = proportion of atom translation MC moves
    Pmoltrans = proportion of molecule translation MC moves
    Pmolrotate = proportion of molecule rotation MC moves
  \sphinxstyleemphasis{rigid} value = fix\sphinxhyphen{}ID
    fix\sphinxhyphen{}ID = ID of {\hyperref[\detokenize{fix_rigid::doc}]{\sphinxcrossref{\DUrole{doc}{fix rigid/small}}}} command
  \sphinxstyleemphasis{shake} value = fix\sphinxhyphen{}ID
    fix\sphinxhyphen{}ID = ID of {\hyperref[\detokenize{fix_shake::doc}]{\sphinxcrossref{\DUrole{doc}{fix shake}}}} command
  \sphinxstyleemphasis{region} value = region\sphinxhyphen{}ID
    region\sphinxhyphen{}ID = ID of region where GCMC exchanges and MC moves are allowed
  \sphinxstyleemphasis{maxangle} value = maximum molecular rotation angle (degrees)
  \sphinxstyleemphasis{pressure} value = pressure of the gas reservoir (pressure units)
  \sphinxstyleemphasis{fugacity\_coeff} value = fugacity coefficient of the gas reservoir (unitless)
  \sphinxstyleemphasis{full\_energy} = compute the entire system energy when performing GCMC exchanges and MC moves
  \sphinxstyleemphasis{charge} value = charge of inserted atoms (charge units)
  \sphinxstyleemphasis{group} value = group\sphinxhyphen{}ID
    group\sphinxhyphen{}ID = group\sphinxhyphen{}ID for inserted atoms (string)
  \sphinxstyleemphasis{grouptype} values = type group\sphinxhyphen{}ID
    type = atom type (1\sphinxhyphen{}Ntypes or type label)
    group\sphinxhyphen{}ID = group\sphinxhyphen{}ID for inserted atoms (string)
  \sphinxstyleemphasis{intra\_energy} value = intramolecular energy (energy units)
  \sphinxstyleemphasis{tfac\_insert} value = scale up/down temperature of inserted atoms (unitless)
  \sphinxstyleemphasis{overlap\_cutoff} value = maximum pair distance for overlap rejection (distance units)
  \sphinxstyleemphasis{max} value = Maximum number of molecules allowed in the system
  \sphinxstyleemphasis{min} value = Minimum number of molecules allowed in the system
\end{sphinxalltt}

\end{itemize}


\subsection{Examples}
\label{\detokenize{fix_gcmc:examples}}
\begin{sphinxVerbatim}[commandchars=\\\{\}]
\PYG{k}{fix }\PYG{n+nv+nvIdentifier}{2}\PYG{+w}{ }\PYG{n+nv+nvIdentifier}{gas}\PYG{+w}{ }\PYG{n}{gcmc}\PYG{+w}{ }\PYG{l+m}{10}\PYG{+w}{ }\PYG{l+m}{1000}\PYG{+w}{ }\PYG{l+m}{1000}\PYG{+w}{ }\PYG{l+m}{2}\PYG{+w}{ }\PYG{l+m}{29494}\PYG{+w}{ }\PYG{l+m}{298.0}\PYG{+w}{ }\PYG{o}{\PYGZhy{}}\PYG{l+m}{0.5}\PYG{+w}{ }\PYG{l+m}{0.01}
\PYG{k}{fix }\PYG{n+nv+nvIdentifier}{3}\PYG{+w}{ }\PYG{n+nv+nvIdentifier}{water}\PYG{+w}{ }\PYG{n}{gcmc}\PYG{+w}{ }\PYG{l+m}{10}\PYG{+w}{ }\PYG{l+m}{100}\PYG{+w}{ }\PYG{l+m}{100}\PYG{+w}{ }\PYG{l+m}{0}\PYG{+w}{ }\PYG{l+m}{3456543}\PYG{+w}{ }\PYG{l+m}{3.0}\PYG{+w}{ }\PYG{o}{\PYGZhy{}}\PYG{l+m}{2.5}\PYG{+w}{ }\PYG{l+m}{0.1}\PYG{+w}{ }\PYG{n}{mol}\PYG{+w}{ }\PYG{n}{my\PYGZus{}one\PYGZus{}water}\PYG{+w}{ }\PYG{n}{maxangle}\PYG{+w}{ }\PYG{l+m}{180}\PYG{+w}{ }\PYG{n}{full\PYGZus{}energy}
\PYG{k}{fix }\PYG{n+nv+nvIdentifier}{4}\PYG{+w}{ }\PYG{n+nv+nvIdentifier}{my\PYGZus{}gas}\PYG{+w}{ }\PYG{n}{gcmc}\PYG{+w}{ }\PYG{l+m}{1}\PYG{+w}{ }\PYG{l+m}{10}\PYG{+w}{ }\PYG{l+m}{10}\PYG{+w}{ }\PYG{l+m}{1}\PYG{+w}{ }\PYG{l+m}{123456543}\PYG{+w}{ }\PYG{l+m}{300.0}\PYG{+w}{ }\PYG{o}{\PYGZhy{}}\PYG{l+m}{12.5}\PYG{+w}{ }\PYG{l+m}{1.0}\PYG{+w}{ }\PYG{k}{region }\PYG{n+nv+nvIdentifier}{disk}

\PYG{k}{labelmap}\PYG{+w}{ }\PYG{n}{atom}\PYG{+w}{ }\PYG{l+m}{1}\PYG{+w}{ }\PYG{n}{Li}
\PYG{k}{fix }\PYG{n+nv+nvIdentifier}{2}\PYG{+w}{ }\PYG{n+nv+nvIdentifier}{ion}\PYG{+w}{ }\PYG{n}{gcmc}\PYG{+w}{ }\PYG{l+m}{10}\PYG{+w}{ }\PYG{l+m}{1000}\PYG{+w}{ }\PYG{l+m}{1000}\PYG{+w}{ }\PYG{n}{Li}\PYG{+w}{ }\PYG{l+m}{29494}\PYG{+w}{ }\PYG{l+m}{298.0}\PYG{+w}{ }\PYG{o}{\PYGZhy{}}\PYG{l+m}{0.5}\PYG{+w}{ }\PYG{l+m}{0.01}
\end{sphinxVerbatim}


\subsection{Description}
\label{\detokenize{fix_gcmc:description}}
\sphinxAtStartPar
This fix performs grand canonical Monte Carlo (GCMC) exchanges of atoms or
molecules with an imaginary ideal gas reservoir at the specified T and
chemical potential (mu) as discussed in {\hyperref[\detokenize{fix_gcmc:frenkel2}]{\sphinxcrossref{\DUrole{std,std-ref}{(Frenkel)}}}}.  It
also attempts Monte Carlo (MC) moves (translations and molecule rotations)
within the simulation cell or region.  If used with the {\hyperref[\detokenize{fix_nh::doc}]{\sphinxcrossref{\DUrole{doc}{fix nvt}}}}
command, simulations in the grand canonical ensemble (muVT, constant
chemical potential, constant volume, and constant temperature) can be
performed.  Specific uses include computing isotherms in microporous
materials, or computing vapor\sphinxhyphen{}liquid coexistence curves.

\sphinxAtStartPar
Every N timesteps the fix attempts both GCMC exchanges (insertions or
deletions) and MC moves of gas atoms or molecules.  On those timesteps, the
average number of attempted GCMC exchanges is X, while the average number
of attempted MC moves is M.  For GCMC exchanges of either molecular or
atomic gasses, these exchanges can be either deletions or insertions, with
equal probability.

\sphinxAtStartPar
The possible choices for MC moves are translation of an atom, translation
of a molecule, and rotation of a molecule.  The relative amounts of each are
determined by the optional \sphinxstyleemphasis{mcmoves} keyword (see below).  The default
behavior is as follows. If the \sphinxstyleemphasis{mol} keyword is used, only molecule
translations and molecule rotations are performed with equal probability.
Conversely, if the \sphinxstyleemphasis{mol} keyword is not used, only atom translations are
performed.  M should typically be chosen to be approximately equal to the
expected number of gas atoms or molecules of the given type within the
simulation cell or region, which will result in roughly one MC move per
atom or molecule per MC cycle.

\sphinxAtStartPar
All inserted particles are always added to two groups: the default group
“all” and the fix group specified in the fix command.  In addition,
particles are also added to any groups specified by the \sphinxstyleemphasis{group} and
\sphinxstyleemphasis{grouptype} keywords.  If inserted particles are individual atoms, they are
assigned the atom type given by the type argument.  If they are molecules,
the type argument has no effect and must be set to zero. Instead, the type
of each atom in the inserted molecule is specified in the file read by the
{\hyperref[\detokenize{molecule::doc}]{\sphinxcrossref{\DUrole{doc}{molecule}}}} command.

\begin{sphinxadmonition}{note}{Note:}
\sphinxAtStartPar
Care should be taken to apply fix gcmc only to
a group that contains only those atoms and molecules
that you wish to manipulate using Monte Carlo.
Hence it is generally not a good idea to specify
the default group “all” in the fix command, although it is allowed.
\end{sphinxadmonition}

\sphinxAtStartPar
This fix cannot be used to perform GCMC insertions of gas atoms or
molecules other than the exchanged type, but GCMC deletions,
and MC translations, and rotations can be performed on any atom/molecule in
the fix group.  All atoms in the simulation cell can be moved using
regular time integration translations, e.g. via {\hyperref[\detokenize{fix_nh::doc}]{\sphinxcrossref{\DUrole{doc}{fix nvt}}}},
resulting in a hybrid GCMC+MD simulation. A smaller\sphinxhyphen{}than\sphinxhyphen{}usual
timestep size may be needed when running such a hybrid simulation,
especially if the inserted molecules are not well equilibrated.

\sphinxAtStartPar
This command may optionally use the \sphinxstyleemphasis{region} keyword to define an
exchange and move volume.  The specified region must have been
previously defined with a {\hyperref[\detokenize{region::doc}]{\sphinxcrossref{\DUrole{doc}{region}}}} command.  It must be
defined with side = \sphinxstyleemphasis{in}.  Insertion attempts occur only within the
specified region. For non\sphinxhyphen{}rectangular regions, random trial points are
generated within the rectangular bounding box until a point is found
that lies inside the region. If no valid point is generated after 1000
trials, no insertion is performed, but it is counted as an attempted
insertion.  Move and deletion attempt candidates are selected from gas
atoms or molecules within the region. If there are no candidates, no
move or deletion is performed, but it is counted as an attempt move or
deletion. If an attempted move places the atom or molecule
center\sphinxhyphen{}of\sphinxhyphen{}mass outside the specified region, a new attempted move is
generated. This process is repeated until the atom or molecule
center\sphinxhyphen{}of\sphinxhyphen{}mass is inside the specified region.

\sphinxAtStartPar
If used with {\hyperref[\detokenize{fix_nh::doc}]{\sphinxcrossref{\DUrole{doc}{fix nvt}}}}, the temperature of the imaginary
reservoir, T, should be set to be equivalent to the target temperature
used in fix nvt. Otherwise, the imaginary reservoir will not be in
thermal equilibrium with the simulation cell. Also, it is important that
the temperature used by \sphinxstyleemphasis{fix nvt} is dynamically updated, which can be
achieved as follows:

\begin{sphinxVerbatim}[commandchars=\\\{\}]
\PYG{k}{compute }\PYG{n+nv+nvIdentifier}{mdtemp}\PYG{+w}{ }\PYG{n+nv+nvIdentifier}{mdatoms}\PYG{+w}{ }\PYG{n}{temp}
\PYG{k}{compute\PYGZus{}modify }\PYG{n+nv+nvIdentifier}{mdtemp}\PYG{+w}{ }\PYG{n}{dynamic}\PYG{o}{/}\PYG{n}{dof}\PYG{+w}{ }\PYG{n}{yes}
\PYG{k}{fix }\PYG{n+nv+nvIdentifier}{mdnvt}\PYG{+w}{ }\PYG{n+nv+nvIdentifier}{mdatoms}\PYG{+w}{ }\PYG{n}{nvt}\PYG{+w}{ }\PYG{n}{temp}\PYG{+w}{ }\PYG{l+m}{300.0}\PYG{+w}{ }\PYG{l+m}{300.0}\PYG{+w}{ }\PYG{l+m}{10.0}
\PYG{k}{fix\PYGZus{}modify }\PYG{n+nv+nvIdentifier}{mdnvt}\PYG{+w}{ }\PYG{n}{temp}\PYG{+w}{ }\PYG{n}{mdtemp}
\end{sphinxVerbatim}

\sphinxAtStartPar
Note that neighbor lists are re\sphinxhyphen{}built every timestep that this fix is
invoked, so you should not set N to be too small.  However, periodic
rebuilds are necessary in order to avoid dangerous rebuilds and missed
interactions. Specifically, avoid performing so many MC translations
per timestep that atoms can move beyond the neighbor list skin
distance. See the {\hyperref[\detokenize{neighbor::doc}]{\sphinxcrossref{\DUrole{doc}{neighbor}}}} command for details.

\sphinxAtStartPar
When an atom or molecule is to be inserted, its coordinates are chosen
at a random position within the current simulation cell or region, and
new atom velocities are randomly chosen from the specified temperature
distribution given by T. The effective temperature for new atom
velocities can be increased or decreased using the optional keyword
\sphinxstyleemphasis{tfac\_insert} (see below). Relative coordinates for atoms in a
molecule are taken from the template molecule provided by the
user. The center of mass of the molecule is placed at the insertion
point. The orientation of the molecule is chosen at random by rotating
about this point.

\sphinxAtStartPar
Individual atoms are inserted, unless the \sphinxstyleemphasis{mol} keyword is used.  It
specifies a \sphinxstyleemphasis{template\sphinxhyphen{}ID} previously defined using the
{\hyperref[\detokenize{molecule::doc}]{\sphinxcrossref{\DUrole{doc}{molecule}}}} command, which reads a file that defines the
molecule.  The coordinates, atom types, charges, etc., as well as any
bonding and special neighbor information for the molecule can
be specified in the molecule file.  See the {\hyperref[\detokenize{molecule::doc}]{\sphinxcrossref{\DUrole{doc}{molecule}}}}
command for details.  The only settings required to be in this file
are the coordinates and types of atoms in the molecule.

\sphinxAtStartPar
When not using the \sphinxstyleemphasis{mol} keyword, you should ensure you do not delete
atoms that are bonded to other atoms, or LAMMPS will soon generate an
error when it tries to find bonded neighbors.  LAMMPS will warn you if
any of the atoms eligible for deletion have a non\sphinxhyphen{}zero molecule ID,
but does not check for this at the time of deletion.

\sphinxAtStartPar
If you wish to insert molecules using the \sphinxstyleemphasis{mol} keyword that will be
treated as rigid bodies, use the \sphinxstyleemphasis{rigid} keyword, specifying as its
value the ID of a separate {\hyperref[\detokenize{fix_rigid::doc}]{\sphinxcrossref{\DUrole{doc}{fix rigid/small}}}} command
which also appears in your input script.

\begin{sphinxadmonition}{note}{Note:}
\sphinxAtStartPar
If you wish the new rigid molecules (and other rigid molecules) to be
thermostatted correctly via {\hyperref[\detokenize{fix_rigid::doc}]{\sphinxcrossref{\DUrole{doc}{fix rigid/small/nvt}}}} or
{\hyperref[\detokenize{fix_rigid::doc}]{\sphinxcrossref{\DUrole{doc}{fix rigid/small/npt}}}}, then you need to use the
{\hyperref[\detokenize{fix_modify::doc}]{\sphinxcrossref{\DUrole{doc}{fix\_modify dynamic/dof yes}}}} command for the rigid
fix.  This is to inform that fix that the molecule count will vary
dynamically.
\end{sphinxadmonition}

\sphinxAtStartPar
If you wish to insert molecules via the \sphinxstyleemphasis{mol} keyword, that will have
their bonds or angles constrained via SHAKE, use the \sphinxstyleemphasis{shake} keyword,
specifying as its value the ID of a separate {\hyperref[\detokenize{fix_shake::doc}]{\sphinxcrossref{\DUrole{doc}{fix shake}}}} command which also appears in your input script.

\sphinxAtStartPar
Optionally, users may specify the relative amounts of different MC
moves using the \sphinxstyleemphasis{mcmoves} keyword. The values \sphinxstyleemphasis{Patomtrans},
\sphinxstyleemphasis{Pmoltrans}, \sphinxstyleemphasis{Pmolrotate} specify the average proportion of
atom translations, molecule translations, and molecule rotations,
respectively. The values must be non\sphinxhyphen{}negative integers or real
numbers, with at least one non\sphinxhyphen{}zero value. For example, (10,30,0)
would result in 25\% of the MC moves being atomic translations, 75\%
molecular translations, and no molecular rotations.

\sphinxAtStartPar
Optionally, users may specify the maximum rotation angle for molecular
rotations using the \sphinxstyleemphasis{maxangle} keyword and specifying the angle in
degrees. Rotations are performed by generating a random point on the
unit sphere and a random rotation angle on the range
{[}0,maxangle). The molecule is then rotated by that angle about an
axis passing through the molecule center of mass. The axis is parallel
to the unit vector defined by the point on the unit sphere.  The same
procedure is used for randomly rotating molecules when they are
inserted, except that the maximum angle is 360 degrees.

\sphinxAtStartPar
Note that fix gcmc does not use configurational bias MC or any other
kind of sampling of intramolecular degrees of freedom.  Inserted
molecules can have different orientations, but they will all have the
same intramolecular configuration, which was specified in the molecule
command input.

\sphinxAtStartPar
For atomic gasses, inserted atoms have the specified atom type, but
deleted atoms are any atoms that have been inserted or that already
belong to the fix group. For molecular gasses, exchanged
molecules use the same atom types as in the template molecule supplied
by the user.  In both cases, exchanged atoms/molecules are assigned to
two groups: the default group “all” and the fix group
(which can also be “all”).

\sphinxAtStartPar
The chemical potential is a user\sphinxhyphen{}specified input parameter defined
as:
\begin{equation*}
\begin{split}\mu = \mu^{id} + \mu^{ex}\end{split}
\end{equation*}
\sphinxAtStartPar
The second term mu\_ex is the excess chemical potential due to
energetic interactions and is formally zero for the fictitious gas
reservoir but is non\sphinxhyphen{}zero for interacting systems. So, while the
chemical potential of the reservoir and the simulation cell are equal,
mu\_ex is not, and as a result, the densities of the two are generally
quite different.  The first term mu\_id is the ideal gas contribution
to the chemical potential.  mu\_id can be related to the density or
pressure of the fictitious gas reservoir by:
\begin{equation*}
\begin{split}\mu^{id}  = & k T \ln{\rho \Lambda^3} \\
          = & k T \ln{\frac{\phi P \Lambda^3}{k_B T}}\end{split}
\end{equation*}
\sphinxAtStartPar
where \(k_B\) is the Boltzmann constant, \(T\) is the user\sphinxhyphen{}specified
temperature, \(\rho\) is the number density, \sphinxstyleemphasis{P} is the pressure,
and \(\phi\) is the fugacity coefficient.  The constant
\(\Lambda\) is required for dimensional consistency.  For all unit
styles except \sphinxstyleemphasis{lj} it is defined as the thermal de Broglie wavelength
\begin{equation*}
\begin{split}\Lambda = \sqrt{ \frac{h^2}{2 \pi m k_B T}}\end{split}
\end{equation*}
\sphinxAtStartPar
where \sphinxstyleemphasis{h} is Planck’s constant, and \sphinxstyleemphasis{m} is the mass of the exchanged atom
or molecule.  For unit style \sphinxstyleemphasis{lj}, \(\Lambda\) is simply set to
unity. Note that prior to March 2017, \(\Lambda\) for unit style \sphinxstyleemphasis{lj}
was calculated using the above formula with \sphinxstyleemphasis{h} set to the rather specific
value of 0.18292026.  Chemical potential under the old definition can
be converted to an equivalent value under the new definition by
subtracting \(3 k T \ln(\Lambda_{old})\).

\sphinxAtStartPar
As an alternative to specifying mu directly, the ideal gas reservoir
can be defined by its pressure \sphinxstyleemphasis{P} using the \sphinxstyleemphasis{pressure} keyword, in
which case the user\sphinxhyphen{}specified chemical potential is ignored. The user
may also specify the fugacity coefficient \(\phi\) using the
\sphinxstyleemphasis{fugacity\_coeff} keyword, which defaults to unity.

\sphinxAtStartPar
The \sphinxstyleemphasis{full\_energy} option means that the fix calculates the total
potential energy of the entire simulated system, instead of just
the energy of the part that is changed. The total system
energy before and after the proposed GCMC exchange or MC move
is then used in the
Metropolis criterion to determine whether or not to accept the
proposed change. By default, this option is off,
in which case only
partial energies are computed to determine the energy difference
due to the proposed change.

\sphinxAtStartPar
The \sphinxstyleemphasis{full\_energy} option is needed for systems with complicated
potential energy calculations, including the following:
\begin{itemize}
\item {} 
\sphinxAtStartPar
long\sphinxhyphen{}range electrostatics (kspace)

\item {} 
\sphinxAtStartPar
many\sphinxhyphen{}body pair styles

\item {} 
\sphinxAtStartPar
hybrid pair styles

\item {} 
\sphinxAtStartPar
eam pair styles

\item {} 
\sphinxAtStartPar
tail corrections

\item {} 
\sphinxAtStartPar
need to include potential energy contributions from other fixes

\end{itemize}

\sphinxAtStartPar
In these cases, LAMMPS will automatically apply the \sphinxstyleemphasis{full\_energy}
keyword and issue a warning message.

\sphinxAtStartPar
When the \sphinxstyleemphasis{mol} keyword is used, the \sphinxstyleemphasis{full\_energy} option also includes
the intramolecular energy of inserted and deleted molecules, whereas
this energy is not included when \sphinxstyleemphasis{full\_energy} is not used. If this
is not desired, the \sphinxstyleemphasis{intra\_energy} keyword can be used to define an
amount of energy that is subtracted from the final energy when a
molecule is inserted, and subtracted from the initial energy when a molecule
is deleted. For molecules that have a non\sphinxhyphen{}zero intramolecular energy,
this will ensure roughly the same behavior whether or not the
\sphinxstyleemphasis{full\_energy} option is used.

\sphinxAtStartPar
Inserted atoms and molecules are assigned random velocities based on
the specified temperature \(T\). Because the relative velocity of all
atoms in the molecule is zero, this may result in inserted molecules
that are systematically too cold. In addition, the intramolecular
potential energy of the inserted molecule may cause the kinetic energy
of the molecule to quickly increase or decrease after insertion.  The
\sphinxstyleemphasis{tfac\_insert} keyword allows the user to counteract these effects by
changing the temperature used to assign velocities to inserted atoms
and molecules by a constant factor. For a particular application, some
experimentation may be required to find a value of \sphinxstyleemphasis{tfac\_insert} that
results in inserted molecules that equilibrate quickly to the correct
temperature.

\sphinxAtStartPar
Some fixes have an associated potential energy. Examples of such fixes
include: {\hyperref[\detokenize{fix_efield::doc}]{\sphinxcrossref{\DUrole{doc}{efield}}}}, {\hyperref[\detokenize{fix_gravity::doc}]{\sphinxcrossref{\DUrole{doc}{gravity}}}},
{\hyperref[\detokenize{fix_addforce::doc}]{\sphinxcrossref{\DUrole{doc}{addforce}}}}, {\hyperref[\detokenize{fix_langevin::doc}]{\sphinxcrossref{\DUrole{doc}{langevin}}}},
{\hyperref[\detokenize{fix_restrain::doc}]{\sphinxcrossref{\DUrole{doc}{restrain}}}}, {\hyperref[\detokenize{fix_temp_berendsen::doc}]{\sphinxcrossref{\DUrole{doc}{temp/berendsen}}}}, {\hyperref[\detokenize{fix_temp_rescale::doc}]{\sphinxcrossref{\DUrole{doc}{temp/rescale}}}}, and
{\hyperref[\detokenize{fix_wall::doc}]{\sphinxcrossref{\DUrole{doc}{wall fixes}}}}.  For that energy to be included in the
total potential energy of the system (the quantity used when performing
GCMC exchange and MC moves), you MUST enable the {\hyperref[\detokenize{fix_modify::doc}]{\sphinxcrossref{\DUrole{doc}{fix\_modify}}}} \sphinxstyleemphasis{energy} option for that fix.  The doc pages for
individual {\hyperref[\detokenize{fix::doc}]{\sphinxcrossref{\DUrole{doc}{fix}}}} commands specify if this should be done.

\sphinxAtStartPar
Use the \sphinxstyleemphasis{charge} option to insert atoms with a user\sphinxhyphen{}specified point
charge. Note that doing so will cause the system to become non\sphinxhyphen{}neutral.
LAMMPS issues a warning when using long\sphinxhyphen{}range electrostatics (kspace)
with non\sphinxhyphen{}neutral systems. See the {\hyperref[\detokenize{compute_group_group::doc}]{\sphinxcrossref{\DUrole{doc}{compute group/group}}}} documentation for more details about simulating
non\sphinxhyphen{}neutral systems with kspace on.

\sphinxAtStartPar
Use of this fix typically will cause the number of atoms to fluctuate,
therefore, you will want to use the {\hyperref[\detokenize{compute_modify::doc}]{\sphinxcrossref{\DUrole{doc}{compute\_modify dynamic/dof}}}} command to ensure that the current number of atoms is
used as a normalizing factor each time temperature is computed. A simple
example of this is:

\begin{sphinxVerbatim}[commandchars=\\\{\}]
\PYG{k}{compute\PYGZus{}modify }\PYG{n+nv+nvIdentifier}{thermo\PYGZus{}temp}\PYG{+w}{ }\PYG{n}{dynamic}\PYG{o}{/}\PYG{n}{dof}\PYG{+w}{ }\PYG{n}{yes}
\end{sphinxVerbatim}

\sphinxAtStartPar
A more complicated example is listed earlier on this page
in the context of NVT dynamics.

\begin{sphinxadmonition}{note}{Note:}
\sphinxAtStartPar
If the density of the cell is initially very small or zero, and
increases to a much larger density after a period of equilibration,
then certain quantities that are only calculated once at the start
(kspace parameters) may no longer be accurate.  The solution is to
start a new simulation after the equilibrium density has been
reached.
\end{sphinxadmonition}

\sphinxAtStartPar
With some pair\_styles, such as {\hyperref[\detokenize{pair_buck::doc}]{\sphinxcrossref{\DUrole{doc}{Buckingham}}}},
{\hyperref[\detokenize{pair_born::doc}]{\sphinxcrossref{\DUrole{doc}{Born\sphinxhyphen{}Mayer\sphinxhyphen{}Huggins}}}} and {\hyperref[\detokenize{pair_reaxff::doc}]{\sphinxcrossref{\DUrole{doc}{ReaxFF}}}},
two atoms placed close to each other may have an arbitrary large,
negative potential energy due to the functional form of the potential.
While these unphysical configurations are inaccessible to typical
dynamical trajectories, they can be generated by Monte Carlo moves. The
\sphinxstyleemphasis{overlap\_cutoff} keyword suppresses these moves by effectively assigning
an infinite positive energy to all new configurations that place any
pair of atoms closer than the specified overlap cutoff distance.

\sphinxAtStartPar
The \sphinxstyleemphasis{max} and \sphinxstyleemphasis{min} keywords allow for the restriction of the number of
atoms in the simulation. They automatically reject all insertion or
deletion moves that would take the system beyond the set boundaries.
Should the system already be beyond the boundary, only moves that bring
the system closer to the bounds may be accepted.

\sphinxAtStartPar
The \sphinxstyleemphasis{group} keyword adds all inserted atoms to the {\hyperref[\detokenize{group::doc}]{\sphinxcrossref{\DUrole{doc}{group}}}}
of the group\sphinxhyphen{}ID value. The \sphinxstyleemphasis{grouptype} keyword adds all inserted atoms
of the specified type to the {\hyperref[\detokenize{group::doc}]{\sphinxcrossref{\DUrole{doc}{group}}}} of the group\sphinxhyphen{}ID value.


\subsection{Restart, fix\_modify, output, run start/stop, minimize info}
\label{\detokenize{fix_gcmc:restart-fix-modify-output-run-start-stop-minimize-info}}
\sphinxAtStartPar
This fix writes the state of the fix to {\hyperref[\detokenize{restart::doc}]{\sphinxcrossref{\DUrole{doc}{binary restart files}}}}.  This includes information about the random number
generator seed, the next timestep for MC exchanges, the number of MC
step attempts and successes etc.  See the {\hyperref[\detokenize{read_restart::doc}]{\sphinxcrossref{\DUrole{doc}{read\_restart}}}} command for info on how to re\sphinxhyphen{}specify a fix in an
input script that reads a restart file, so that the operation of the
fix continues in an uninterrupted fashion.

\begin{sphinxadmonition}{note}{Note:}
\sphinxAtStartPar
For this to work correctly, the timestep must \sphinxstylestrong{not} be changed
after reading the restart with {\hyperref[\detokenize{reset_timestep::doc}]{\sphinxcrossref{\DUrole{doc}{reset\_timestep}}}}.
The fix will try to detect it and stop with an error.
\end{sphinxadmonition}

\sphinxAtStartPar
None of the {\hyperref[\detokenize{fix_modify::doc}]{\sphinxcrossref{\DUrole{doc}{fix\_modify}}}} options are relevant to
this fix.

\sphinxAtStartPar
This fix computes a global vector of length 8, which can be accessed
by various {\hyperref[\detokenize{Howto_output::doc}]{\sphinxcrossref{\DUrole{doc}{output commands}}}}.  The vector values are
the following global cumulative quantities:
\begin{enumerate}
\sphinxsetlistlabels{\arabic}{enumi}{enumii}{}{.}%
\item {} 
\sphinxAtStartPar
translation attempts

\item {} 
\sphinxAtStartPar
translation successes

\item {} 
\sphinxAtStartPar
insertion attempts

\item {} 
\sphinxAtStartPar
insertion successes

\item {} 
\sphinxAtStartPar
deletion attempts

\item {} 
\sphinxAtStartPar
deletion successes

\item {} 
\sphinxAtStartPar
rotation attempts

\item {} 
\sphinxAtStartPar
rotation successes

\end{enumerate}

\sphinxAtStartPar
The vector values calculated by this fix are “intensive”.

\sphinxAtStartPar
No parameter of this fix can be used with the \sphinxstyleemphasis{start/stop} keywords of
the {\hyperref[\detokenize{run::doc}]{\sphinxcrossref{\DUrole{doc}{run}}}} command.  This fix is not invoked during
{\hyperref[\detokenize{minimize::doc}]{\sphinxcrossref{\DUrole{doc}{energy minimization}}}}.


\subsection{Restrictions}
\label{\detokenize{fix_gcmc:restrictions}}
\sphinxAtStartPar
This fix is part of the MC package.  It is only enabled if LAMMPS was
built with that package.  See the {\hyperref[\detokenize{Build_package::doc}]{\sphinxcrossref{\DUrole{doc}{Build package}}}}
doc page for more info.

\sphinxAtStartPar
This fix style requires an {\hyperref[\detokenize{atom_style::doc}]{\sphinxcrossref{\DUrole{doc}{atom style}}}} with per atom
type masses.

\sphinxAtStartPar
Do not set “neigh\_modify once yes” or else this fix will never be
called.  Reneighboring is \sphinxstylestrong{required}.

\sphinxAtStartPar
Only usable for 3D simulations.

\sphinxAtStartPar
This fix can be run in parallel, but aspects of the GCMC part will not
scale well in parallel.  Currently, molecule translations and rotations
are not supported with more than one MPI process.  It is still possible
to do parallel molecule exchange without translation and rotation moves
by setting MC moves to zero and/or by using the \sphinxstyleemphasis{mcmoves} keyword with
\sphinxstyleemphasis{Pmoltrans} = \sphinxstyleemphasis{Pmolrotate} = 0 .

\sphinxAtStartPar
When using fix gcmc in combination with fix shake or fix rigid, only
GCMC exchange moves are supported, so the argument \sphinxstyleemphasis{M} must be zero.

\sphinxAtStartPar
When using fix gcmc in combination with fix rigid, deletion of the last
remaining molecule is not allowed for technical reasons, and so the
molecule count will never drop below 1, regardless of the specified
chemical potential.

\sphinxAtStartPar
Note that very lengthy simulations involving insertions/deletions of
billions of gas molecules may run out of atom or molecule IDs and
trigger an error, so it is better to run multiple shorter\sphinxhyphen{}duration
simulations. Likewise, very large molecules have not been tested and
may turn out to be problematic.

\sphinxAtStartPar
Use of multiple fix gcmc commands in the same input script can be
problematic if using a template molecule. The issue is that the
user\sphinxhyphen{}referenced template molecule in the second fix gcmc command may
no longer exist since it might have been deleted by the first fix gcmc
command. An existing template molecule will need to be referenced by
the user for each subsequent fix gcmc command.


\subsection{Related commands}
\label{\detokenize{fix_gcmc:related-commands}}
\sphinxAtStartPar
{\hyperref[\detokenize{fix_atom_swap::doc}]{\sphinxcrossref{\DUrole{doc}{fix atom/swap}}}},
{\hyperref[\detokenize{fix_nh::doc}]{\sphinxcrossref{\DUrole{doc}{fix nvt}}}}, {\hyperref[\detokenize{neighbor::doc}]{\sphinxcrossref{\DUrole{doc}{neighbor}}}},
{\hyperref[\detokenize{fix_deposit::doc}]{\sphinxcrossref{\DUrole{doc}{fix deposit}}}}, {\hyperref[\detokenize{fix_evaporate::doc}]{\sphinxcrossref{\DUrole{doc}{fix evaporate}}}},
{\hyperref[\detokenize{delete_atoms::doc}]{\sphinxcrossref{\DUrole{doc}{delete\_atoms}}}}


\subsection{Default}
\label{\detokenize{fix_gcmc:default}}
\sphinxAtStartPar
The option defaults are mol = no, maxangle = 10, overlap\_cutoff = 0.0,
fugacity\_coeff = 1.0, intra\_energy = 0.0, tfac\_insert = 1.0.
(Patomtrans, Pmoltrans, Pmolrotate) = (1, 0, 0) for mol = no and
(0, 1, 1) for mol = yes. full\_energy = no,
except for the situations where full\_energy is required, as
listed above.


\bigskip\hrule\bigskip

\phantomsection\label{\detokenize{fix_gcmc:frenkel2}}
\sphinxAtStartPar
\sphinxstylestrong{(Frenkel)} Frenkel and Smit, Understanding Molecular Simulation,
Academic Press, London, 2002.

\sphinxstepscope

\index{fix gld@\spxentry{fix gld}}\ignorespaces 

\section{fix gld command}
\label{\detokenize{fix_gld:fix-gld-command}}\label{\detokenize{fix_gld:index-0}}\label{\detokenize{fix_gld::doc}}

\subsection{Syntax}
\label{\detokenize{fix_gld:syntax}}
\begin{sphinxVerbatim}[commandchars=\\\{\}]
\PYG{k}{fix }\PYG{n+nv+nvIdentifier}{ID}\PYG{+w}{ }\PYG{n+nv+nvIdentifier}{group\PYGZhy{}ID}\PYG{+w}{ }\PYG{n}{gld}\PYG{+w}{ }\PYG{n}{Tstart}\PYG{+w}{ }\PYG{n}{Tstop}\PYG{+w}{ }\PYG{n}{N\PYGZus{}k}\PYG{+w}{ }\PYG{n}{seed}\PYG{+w}{ }\PYG{n}{series}\PYG{+w}{ }\PYG{n}{c\PYGZus{}1}\PYG{+w}{ }\PYG{n}{tau\PYGZus{}1}\PYG{+w}{ }\PYG{n}{...}\PYG{+w}{ }\PYG{n}{c\PYGZus{}N\PYGZus{}k}\PYG{+w}{ }\PYG{n}{tau\PYGZus{}N\PYGZus{}k}\PYG{+w}{ }\PYG{n}{keyword}\PYG{+w}{ }\PYG{n}{values}\PYG{+w}{ }\PYG{n}{...}
\end{sphinxVerbatim}
\begin{itemize}
\item {} 
\sphinxAtStartPar
ID, group\sphinxhyphen{}ID are documented in {\hyperref[\detokenize{fix::doc}]{\sphinxcrossref{\DUrole{doc}{fix}}}} command

\item {} 
\sphinxAtStartPar
gld = style name of this fix command

\item {} 
\sphinxAtStartPar
Tstart,Tstop = desired temperature at start/end of run (temperature units)

\item {} 
\sphinxAtStartPar
N\_k = number of terms in the Prony series representation of the memory kernel

\item {} 
\sphinxAtStartPar
seed = random number seed to use for white noise (positive integer)

\item {} 
\sphinxAtStartPar
series = \sphinxstyleemphasis{pprony} is presently the only available option

\item {} 
\sphinxAtStartPar
c\_k = the weight of the kth term in the Prony series (mass per time units)

\item {} 
\sphinxAtStartPar
tau\_k = the time constant of the kth term in the Prony series (time units)

\item {} 
\sphinxAtStartPar
zero or more keyword/value pairs may be appended
\begin{sphinxalltt}
keyword = \sphinxstyleemphasis{frozen} or \sphinxstyleemphasis{zero}
  \sphinxstyleemphasis{frozen} value = \sphinxstyleemphasis{no} or \sphinxstyleemphasis{yes}
    \sphinxstyleemphasis{no} = initialize extended variables using values drawn from equilibrium distribution at Tstart
    \sphinxstyleemphasis{yes} = initialize extended variables to zero (i.e., from equilibrium distribution at zero temperature)
  \sphinxstyleemphasis{zero} value = \sphinxstyleemphasis{no} or \sphinxstyleemphasis{yes}
    \sphinxstyleemphasis{no} = do not set total random force to zero
    \sphinxstyleemphasis{yes} = set total random force to zero
\end{sphinxalltt}

\end{itemize}


\subsection{Examples}
\label{\detokenize{fix_gld:examples}}
\begin{sphinxVerbatim}[commandchars=\\\{\}]
\PYG{k}{fix }\PYG{n+nv+nvIdentifier}{1}\PYG{+w}{ }\PYG{n+nv+nvIdentifier}{all}\PYG{+w}{ }\PYG{n}{gld}\PYG{+w}{ }\PYG{l+m}{1.0}\PYG{+w}{ }\PYG{l+m}{1.0}\PYG{+w}{ }\PYG{l+m}{2}\PYG{+w}{ }\PYG{l+m}{82885}\PYG{+w}{ }\PYG{n}{pprony}\PYG{+w}{ }\PYG{l+m}{0.5}\PYG{+w}{ }\PYG{l+m}{1.0}\PYG{+w}{ }\PYG{l+m}{1.0}\PYG{+w}{ }\PYG{l+m}{2.0}\PYG{+w}{ }\PYG{n}{frozen}\PYG{+w}{ }\PYG{n}{yes}\PYG{+w}{ }\PYG{n}{zero}\PYG{+w}{ }\PYG{n}{yes}
\PYG{k}{fix }\PYG{n+nv+nvIdentifier}{3}\PYG{+w}{ }\PYG{n+nv+nvIdentifier}{rouse}\PYG{+w}{ }\PYG{n}{gld}\PYG{+w}{ }\PYG{l+m}{7.355}\PYG{+w}{ }\PYG{l+m}{7.355}\PYG{+w}{ }\PYG{l+m}{4}\PYG{+w}{ }\PYG{l+m}{48823}\PYG{+w}{ }\PYG{n}{pprony}\PYG{+w}{ }\PYG{l+m}{107.1}\PYG{+w}{ }\PYG{l+m}{0.02415}\PYG{+w}{ }\PYG{l+m}{186.0}\PYG{+w}{ }\PYG{l+m}{0.04294}\PYG{+w}{ }\PYG{l+m}{428.6}\PYG{+w}{ }\PYG{l+m}{0.09661}\PYG{+w}{ }\PYG{l+m}{1714}\PYG{+w}{ }\PYG{l+m}{0.38643}
\end{sphinxVerbatim}


\subsection{Description}
\label{\detokenize{fix_gld:description}}
\sphinxAtStartPar
Applies Generalized Langevin Dynamics to a group of atoms, as
described in {\hyperref[\detokenize{fix_gld:baczewski}]{\sphinxcrossref{\DUrole{std,std-ref}{(Baczewski)}}}}.  This is intended to model the
effect of an implicit solvent with a temporally non\sphinxhyphen{}local dissipative
force and a colored Gaussian random force, consistent with the
Fluctuation\sphinxhyphen{}Dissipation Theorem.  The functional form of the memory
kernel associated with the temporally non\sphinxhyphen{}local force is constrained
to be a Prony series.

\begin{sphinxadmonition}{note}{Note:}
\sphinxAtStartPar
While this fix bears many similarities to {\hyperref[\detokenize{fix_langevin::doc}]{\sphinxcrossref{\DUrole{doc}{fix langevin}}}}, it has one significant
difference. Namely, {\hyperref[\detokenize{fix_gld::doc}]{\sphinxcrossref{\DUrole{doc}{fix gld}}}} performs time integration,
whereas {\hyperref[\detokenize{fix_langevin::doc}]{\sphinxcrossref{\DUrole{doc}{fix langevin}}}} does NOT. To this end, the
specification of another fix to perform time integration, such as {\hyperref[\detokenize{fix_nve::doc}]{\sphinxcrossref{\DUrole{doc}{fix nve}}}}, is NOT necessary.
\end{sphinxadmonition}

\sphinxAtStartPar
With this fix active, the force on the \sphinxstyleemphasis{j}th atom is given as
\begin{equation*}
\begin{split}{\bf F}_{j}(t) = & {\bf F}^C_j(t)-\int \limits_{0}^{t} \Gamma_j(t-s) {\bf v}_j(s)~\text{d}s + {\bf F}^R_j(t) \\
\Gamma_j(t-s) = & \sum \limits_{k=1}^{N_k} \frac{c_k}{\tau_k} e^{-(t-s)/\tau_k} \\
\langle{\bf F}^R_j(t),{\bf F}^R_j(s)\rangle = & \text{k$_\text{B}$T} ~\Gamma_j(t-s)\end{split}
\end{equation*}
\sphinxAtStartPar
Here, the first term is representative of all conservative (pairwise,
bonded, etc) forces external to this fix, the second is the temporally
non\sphinxhyphen{}local dissipative force given as a Prony series, and the third is
the colored Gaussian random force.

\sphinxAtStartPar
The Prony series form of the memory kernel is chosen to enable an
extended variable formalism, with a number of exemplary mathematical
features discussed in {\hyperref[\detokenize{fix_gld:baczewski}]{\sphinxcrossref{\DUrole{std,std-ref}{(Baczewski)}}}}. In particular, \(3N_k\)
extended variables are added to each atom, which effect the action of
the memory kernel without having to explicitly evaluate the integral
over time in the second term of the force. This also has the benefit
of requiring the generation of uncorrelated random forces, rather than
correlated random forces as specified in the third term of the force.

\sphinxAtStartPar
Presently, the Prony series coefficients are limited to being greater
than or equal to zero, and the time constants are limited to being
greater than zero. To this end, the value of series MUST be set to
\sphinxstyleemphasis{pprony}, for now. Future updates will allow for negative coefficients
and other representations of the memory kernel. It is with these
updates in mind that the series option was included.

\sphinxAtStartPar
The units of the Prony series coefficients are chosen to be mass per
time to ensure that the numerical integration scheme stably approaches
the Newtonian and Langevin limits. Details of these limits, and the
associated numerical concerns are discussed in
{\hyperref[\detokenize{fix_gld:baczewski}]{\sphinxcrossref{\DUrole{std,std-ref}{(Baczewski)}}}}.

\sphinxAtStartPar
The desired temperature at each timestep is ramped from \sphinxstyleemphasis{Tstart} to
\sphinxstyleemphasis{Tstop} over the course of the next run.

\sphinxAtStartPar
The random \# \sphinxstyleemphasis{seed} must be a positive integer. A Marsaglia random
number generator is used. Each processor uses the input seed to
generate its own unique seed and its own stream of random
numbers. Thus the dynamics of the system will not be identical on two
runs on different numbers of processors.


\bigskip\hrule\bigskip


\sphinxAtStartPar
The keyword/value option pairs are used in the following ways.

\sphinxAtStartPar
The keyword \sphinxstyleemphasis{frozen} can be used to specify how the extended variables
associated with the GLD memory kernel are initialized. Specifying no
(the default), the initial values are drawn at random from an
equilibrium distribution at \sphinxstyleemphasis{Tstart}, consistent with the
Fluctuation\sphinxhyphen{}Dissipation Theorem. Specifying yes, initializes the
extended variables to zero.

\sphinxAtStartPar
The keyword \sphinxstyleemphasis{zero} can be used to eliminate drift due to the
thermostat. Because the random forces on different atoms are
independent, they do not sum exactly to zero. As a result, this fix
applies a small random force to the entire system, and the
center\sphinxhyphen{}of\sphinxhyphen{}mass of the system undergoes a slow random walk. If the
keyword \sphinxstyleemphasis{zero} is set to \sphinxstyleemphasis{yes}, the total random force is set exactly
to zero by subtracting off an equal part of it from each atom in the
group. As a result, the center\sphinxhyphen{}of\sphinxhyphen{}mass of a system with zero initial
momentum will not drift over time.


\bigskip\hrule\bigskip



\subsection{Restart, fix\_modify, output, run start/stop, minimize info}
\label{\detokenize{fix_gld:restart-fix-modify-output-run-start-stop-minimize-info}}
\sphinxAtStartPar
The instantaneous values of the extended variables are written to
{\hyperref[\detokenize{restart::doc}]{\sphinxcrossref{\DUrole{doc}{binary restart files}}}}.  Because the state of the random
number generator is not saved in restart files, this means you cannot
do “exact” restarts with this fix, where the simulation continues on
the same as if no restart had taken place. However, in a statistical
sense, a restarted simulation should produce the same behavior.

\sphinxAtStartPar
None of the {\hyperref[\detokenize{fix_modify::doc}]{\sphinxcrossref{\DUrole{doc}{fix\_modify}}}} options are relevant to this
fix.  No global or per\sphinxhyphen{}atom quantities are stored by this fix for
access by various {\hyperref[\detokenize{Howto_output::doc}]{\sphinxcrossref{\DUrole{doc}{output commands}}}}.

\sphinxAtStartPar
This fix can ramp its target temperature over multiple runs, using the
\sphinxstyleemphasis{start} and \sphinxstyleemphasis{stop} keywords of the {\hyperref[\detokenize{run::doc}]{\sphinxcrossref{\DUrole{doc}{run}}}} command.  See the
{\hyperref[\detokenize{run::doc}]{\sphinxcrossref{\DUrole{doc}{run}}}} command for details of how to do this.

\sphinxAtStartPar
This fix is not invoked during {\hyperref[\detokenize{minimize::doc}]{\sphinxcrossref{\DUrole{doc}{energy minimization}}}}.


\subsection{Restrictions}
\label{\detokenize{fix_gld:restrictions}}
\sphinxAtStartPar
This fix is part of the EXTRA\sphinxhyphen{}FIX package.  It is only enabled if LAMMPS
was built with that package.  See the {\hyperref[\detokenize{Build_package::doc}]{\sphinxcrossref{\DUrole{doc}{Build package}}}} page for more info.


\subsection{Related commands}
\label{\detokenize{fix_gld:related-commands}}
\sphinxAtStartPar
{\hyperref[\detokenize{fix_langevin::doc}]{\sphinxcrossref{\DUrole{doc}{fix langevin}}}}, {\hyperref[\detokenize{fix_viscous::doc}]{\sphinxcrossref{\DUrole{doc}{fix viscous}}}},
{\hyperref[\detokenize{pair_dpd::doc}]{\sphinxcrossref{\DUrole{doc}{pair\_style dpd/tstat}}}}


\subsection{Default}
\label{\detokenize{fix_gld:default}}
\sphinxAtStartPar
The option defaults are frozen = no, zero = no.


\bigskip\hrule\bigskip

\phantomsection\label{\detokenize{fix_gld:baczewski}}
\sphinxAtStartPar
\sphinxstylestrong{(Baczewski)} A.D. Baczewski and S.D. Bond, J. Chem. Phys. 139, 044107 (2013).

\sphinxstepscope

\index{fix gle@\spxentry{fix gle}}\ignorespaces 

\section{fix gle command}
\label{\detokenize{fix_gle:fix-gle-command}}\label{\detokenize{fix_gle:index-0}}\label{\detokenize{fix_gle::doc}}

\subsection{Syntax}
\label{\detokenize{fix_gle:syntax}}
\begin{sphinxVerbatim}[commandchars=\\\{\}]
\PYG{k}{fix }\PYG{n+nv+nvIdentifier}{ID}\PYG{+w}{ }\PYG{n+nv+nvIdentifier}{id\PYGZhy{}group}\PYG{+w}{ }\PYG{n}{gle}\PYG{+w}{ }\PYG{n}{Ns}\PYG{+w}{ }\PYG{n}{Tstart}\PYG{+w}{ }\PYG{n}{Tstop}\PYG{+w}{ }\PYG{n}{seed}\PYG{+w}{ }\PYG{n}{Amatrix}\PYG{+w}{ }\PYG{n}{[noneq}\PYG{+w}{ }\PYG{n}{Cmatrix]}\PYG{+w}{ }\PYG{n}{[every}\PYG{+w}{ }\PYG{n}{stride]}
\end{sphinxVerbatim}
\begin{itemize}
\item {} 
\sphinxAtStartPar
ID, group\sphinxhyphen{}ID are documented in {\hyperref[\detokenize{fix::doc}]{\sphinxcrossref{\DUrole{doc}{fix}}}} command

\item {} 
\sphinxAtStartPar
gle = style name of this fix command

\item {} 
\sphinxAtStartPar
Ns = number of additional fictitious momenta

\item {} 
\sphinxAtStartPar
Tstart, Tstop = temperature ramp during the run

\item {} 
\sphinxAtStartPar
Amatrix = file to read the drift matrix A from

\item {} 
\sphinxAtStartPar
seed = random number seed to use for generating noise (positive integer)

\item {} 
\sphinxAtStartPar
zero or more keyword/value pairs may be appended
\begin{sphinxalltt}
keyword = \sphinxstyleemphasis{noneq} or \sphinxstyleemphasis{every}
  \sphinxstyleemphasis{noneq} Cmatrix  = file to read the non\sphinxhyphen{}equilibrium covariance matrix from
  \sphinxstyleemphasis{every} stride   = apply the GLE once every time steps. Reduces the accuracy
      of the integration of the GLE, but has *no effect* on the accuracy of equilibrium
      sampling. It might change sampling properties when used together with \sphinxstyleemphasis{noneq}.
\end{sphinxalltt}

\end{itemize}


\subsection{Examples}
\label{\detokenize{fix_gle:examples}}
\begin{sphinxVerbatim}[commandchars=\\\{\}]
\PYG{k}{fix }\PYG{n+nv+nvIdentifier}{3}\PYG{+w}{ }\PYG{n+nv+nvIdentifier}{boundary}\PYG{+w}{ }\PYG{n}{gle}\PYG{+w}{ }\PYG{l+m}{6}\PYG{+w}{ }\PYG{l+m}{300}\PYG{+w}{ }\PYG{l+m}{300}\PYG{+w}{ }\PYG{l+m}{31415}\PYG{+w}{ }\PYG{n}{smart.A}
\PYG{k}{fix }\PYG{n+nv+nvIdentifier}{1}\PYG{+w}{ }\PYG{n+nv+nvIdentifier}{all}\PYG{+w}{ }\PYG{n}{gle}\PYG{+w}{ }\PYG{l+m}{6}\PYG{+w}{ }\PYG{l+m}{300}\PYG{+w}{ }\PYG{l+m}{300}\PYG{+w}{ }\PYG{l+m}{31415}\PYG{+w}{ }\PYG{n}{qt}\PYG{o}{\PYGZhy{}}\PYG{l+m}{300}\PYG{n}{k.A}\PYG{+w}{ }\PYG{n}{noneq}\PYG{+w}{ }\PYG{n}{qt}\PYG{o}{\PYGZhy{}}\PYG{l+m}{300}\PYG{n}{k.C}
\end{sphinxVerbatim}


\subsection{Description}
\label{\detokenize{fix_gle:description}}
\sphinxAtStartPar
Apply a Generalized Langevin Equation (GLE) thermostat as described
in {\hyperref[\detokenize{fix_gle:ceriotti}]{\sphinxcrossref{\DUrole{std,std-ref}{(Ceriotti)}}}}. The formalism allows one to obtain a number
of different effects ranging from efficient sampling of all
vibrational modes in the system to inexpensive (approximate)
modelling of nuclear quantum effects. Contrary to
{\hyperref[\detokenize{fix_langevin::doc}]{\sphinxcrossref{\DUrole{doc}{fix langevin}}}}, this fix performs both
thermostatting and evolution of the Hamiltonian equations of motion, so it
should not be used together with {\hyperref[\detokenize{fix_nve::doc}]{\sphinxcrossref{\DUrole{doc}{fix nve}}}} \textendash{} at least not
on the same atom groups.

\sphinxAtStartPar
Each degree of freedom in the thermostatted group is supplemented
with Ns additional degrees of freedom s, and the equations of motion
become

\begin{sphinxVerbatim}[commandchars=\\\{\}]
dq/dt=p/m
d(p,s)/dt=(F,0) \PYGZhy{} A(p,s) + B dW/dt
\end{sphinxVerbatim}

\sphinxAtStartPar
where F is the physical force, A is the drift matrix (that generalizes
the friction in Langevin dynamics), B is the diffusion term and dW/dt
un\sphinxhyphen{}correlated Gaussian random forces. The A matrix couples the physical
(q,p) dynamics with that of the additional degrees of freedom,
and makes it possible to obtain effectively a history\sphinxhyphen{}dependent
noise and friction kernel.

\sphinxAtStartPar
The drift matrix should be given as an external file \sphinxstyleemphasis{Afile},
as a (Ns+1 x Ns+1) matrix in inverse time units. Matrices that are
optimal for a given application and the system of choice can be
obtained from {\hyperref[\detokenize{fix_gle:gle4md}]{\sphinxcrossref{\DUrole{std,std-ref}{(GLE4MD)}}}}.

\sphinxAtStartPar
Equilibrium sampling a temperature T is obtained by specifying the
target value as the \sphinxstyleemphasis{Tstart} and \sphinxstyleemphasis{Tstop} arguments, so that the diffusion
matrix that gives canonical sampling for a given A is computed automatically.
However, the GLE framework also allow for non\sphinxhyphen{}equilibrium sampling, that
can be used for instance to model inexpensively zero\sphinxhyphen{}point energy
effects {\hyperref[\detokenize{fix_gle:ceriotti2}]{\sphinxcrossref{\DUrole{std,std-ref}{(Ceriotti2)}}}}. This is achieved specifying the \sphinxstyleemphasis{noneq}
keyword followed by the name of the file that contains the static covariance
matrix for the non\sphinxhyphen{}equilibrium dynamics.  Please note, that the covariance
matrix is expected to be given in \sphinxstylestrong{temperature units}.

\sphinxAtStartPar
Since integrating GLE dynamics can be costly when used together with
simple potentials, one can use the \sphinxstyleemphasis{every} optional keyword to
apply the Langevin terms only once every several MD steps, in a
multiple time\sphinxhyphen{}step fashion. This should be used with care when doing
non\sphinxhyphen{}equilibrium sampling, but should have no effect on equilibrium
averages when using canonical sampling.

\sphinxAtStartPar
The random number \sphinxstyleemphasis{seed} must be a positive integer.  A Marsaglia random
number generator is used.  Each processor uses the input seed to
generate its own unique seed and its own stream of random numbers.
Thus the dynamics of the system will not be identical on two runs on
different numbers of processors.

\sphinxAtStartPar
Note also that the Generalized Langevin Dynamics scheme that is
implemented by the {\hyperref[\detokenize{fix_gld::doc}]{\sphinxcrossref{\DUrole{doc}{fix gld}}}} scheme is closely related
to the present one. In fact, it should be always possible to cast the
Prony series form of the memory kernel used by GLD into an appropriate
input matrix for {\hyperref[\detokenize{fix_gle::doc}]{\sphinxcrossref{\DUrole{doc}{fix gle}}}}. While the GLE scheme is more
general, the form used by {\hyperref[\detokenize{fix_gld::doc}]{\sphinxcrossref{\DUrole{doc}{fix gld}}}} can be more directly
related to the representation of an implicit solvent environment.


\subsection{Restart, fix\_modify, output, run start/stop, minimize info}
\label{\detokenize{fix_gle:restart-fix-modify-output-run-start-stop-minimize-info}}
\sphinxAtStartPar
The instantaneous values of the extended variables are written to
{\hyperref[\detokenize{restart::doc}]{\sphinxcrossref{\DUrole{doc}{binary restart files}}}}.  Because the state of the
random number generator is not saved in restart files, this means you
cannot do “exact” restarts with this fix, where the simulation
continues on the same as if no restart had taken place. However, in a
statistical sense, a restarted simulation should produce the same
behavior.  Note however that you should use a different seed each time
you restart, otherwise the same sequence of random numbers will be
used each time, which might lead to stochastic synchronization and
subtle artifacts in the sampling.

\sphinxAtStartPar
The cumulative energy change in the system imposed by this fix is
included in the {\hyperref[\detokenize{thermo_style::doc}]{\sphinxcrossref{\DUrole{doc}{thermodynamic output}}}} keywords
\sphinxstyleemphasis{ecouple} and \sphinxstyleemphasis{econserve}.  See the {\hyperref[\detokenize{thermo_style::doc}]{\sphinxcrossref{\DUrole{doc}{thermo\_style}}}}
doc page for details.

\sphinxAtStartPar
This fix computes a global scalar which can be accessed by various
{\hyperref[\detokenize{Howto_output::doc}]{\sphinxcrossref{\DUrole{doc}{output commands}}}}.  The scalar is the same
cumulative energy change due to this fix described in the previous
paragraph.  The scalar value calculated by this fix is “extensive”.

\sphinxAtStartPar
This fix can ramp its target temperature over multiple runs, using the
\sphinxstyleemphasis{start} and \sphinxstyleemphasis{stop} keywords of the {\hyperref[\detokenize{run::doc}]{\sphinxcrossref{\DUrole{doc}{run}}}} command.  See the
{\hyperref[\detokenize{run::doc}]{\sphinxcrossref{\DUrole{doc}{run}}}} command for details of how to do this.

\sphinxAtStartPar
This fix is not invoked during {\hyperref[\detokenize{minimize::doc}]{\sphinxcrossref{\DUrole{doc}{energy minimization}}}}.


\subsection{Restrictions}
\label{\detokenize{fix_gle:restrictions}}
\sphinxAtStartPar
The GLE thermostat in its current implementation should not be used
with rigid bodies, SHAKE or RATTLE. It is expected that all the
thermostatted degrees of freedom are fully flexible, and the sampled
ensemble will not be correct otherwise.

\sphinxAtStartPar
In order to perform constant\sphinxhyphen{}pressure simulations please use
{\hyperref[\detokenize{fix_press_berendsen::doc}]{\sphinxcrossref{\DUrole{doc}{fix press/berendsen}}}}, rather than
{\hyperref[\detokenize{fix_nh::doc}]{\sphinxcrossref{\DUrole{doc}{fix npt}}}}, to avoid duplicate integration of the
equations of motion.

\sphinxAtStartPar
This fix is part of the EXTRA\sphinxhyphen{}FIX package.  It is only enabled if
LAMMPS was built with that package.  See the {\hyperref[\detokenize{Build_package::doc}]{\sphinxcrossref{\DUrole{doc}{Build package}}}} page for more info.


\subsection{Related commands}
\label{\detokenize{fix_gle:related-commands}}
\sphinxAtStartPar
{\hyperref[\detokenize{fix_nh::doc}]{\sphinxcrossref{\DUrole{doc}{fix nvt}}}}, {\hyperref[\detokenize{fix_temp_rescale::doc}]{\sphinxcrossref{\DUrole{doc}{fix temp/rescale}}}}, {\hyperref[\detokenize{fix_viscous::doc}]{\sphinxcrossref{\DUrole{doc}{fix viscous}}}}, {\hyperref[\detokenize{fix_nh::doc}]{\sphinxcrossref{\DUrole{doc}{fix nvt}}}}, {\hyperref[\detokenize{pair_dpd::doc}]{\sphinxcrossref{\DUrole{doc}{pair\_style dpd/tstat}}}}, {\hyperref[\detokenize{fix_gld::doc}]{\sphinxcrossref{\DUrole{doc}{fix gld}}}}


\bigskip\hrule\bigskip

\phantomsection\label{\detokenize{fix_gle:ceriotti}}
\sphinxAtStartPar
\sphinxstylestrong{(Ceriotti)} Ceriotti, Bussi and Parrinello, J Chem Theory Comput 6,
1170\sphinxhyphen{}80 (2010)

\phantomsection\label{\detokenize{fix_gle:gle4md}}
\sphinxAtStartPar
\sphinxstylestrong{(GLE4MD)} \sphinxurl{https://gle4md.org/}

\phantomsection\label{\detokenize{fix_gle:ceriotti2}}
\sphinxAtStartPar
\sphinxstylestrong{(Ceriotti2)} Ceriotti, Bussi and Parrinello, Phys Rev Lett 103,
030603 (2009)

\sphinxstepscope

\index{fix gravity@\spxentry{fix gravity}}\ignorespaces 
\index{fix gravity/omp@\spxentry{fix gravity/omp}}\ignorespaces 
\index{fix gravity/kk@\spxentry{fix gravity/kk}}\ignorespaces 

\section{fix gravity command}
\label{\detokenize{fix_gravity:fix-gravity-command}}\label{\detokenize{fix_gravity:index-2}}\label{\detokenize{fix_gravity:index-1}}\label{\detokenize{fix_gravity:index-0}}\label{\detokenize{fix_gravity::doc}}
\sphinxAtStartPar
Accelerator Variants: \sphinxstyleemphasis{gravity/omp}, \sphinxstyleemphasis{gravity/kk}


\subsection{Syntax}
\label{\detokenize{fix_gravity:syntax}}
\begin{sphinxVerbatim}[commandchars=\\\{\}]
\PYG{k}{fix }\PYG{n+nv+nvIdentifier}{ID}\PYG{+w}{ }\PYG{n+nv+nvIdentifier}{group}\PYG{+w}{ }\PYG{n}{gravity}\PYG{+w}{ }\PYG{n}{magnitude}\PYG{+w}{ }\PYG{n}{style}\PYG{+w}{ }\PYG{n}{args}
\end{sphinxVerbatim}
\begin{itemize}
\item {} 
\sphinxAtStartPar
ID, group are documented in {\hyperref[\detokenize{fix::doc}]{\sphinxcrossref{\DUrole{doc}{fix}}}} command

\item {} 
\sphinxAtStartPar
gravity = style name of this fix command

\item {} 
\sphinxAtStartPar
magnitude = size of acceleration (force/mass units)

\item {} 
\sphinxAtStartPar
magnitude can be a variable (see below)

\item {} 
\sphinxAtStartPar
style = \sphinxstyleemphasis{chute} or \sphinxstyleemphasis{spherical} or \sphinxstyleemphasis{gradient} or \sphinxstyleemphasis{vector}
\begin{sphinxalltt}
\sphinxstyleemphasis{chute} args = angle
  angle = angle in +x away from \sphinxhyphen{}z or \sphinxhyphen{}y axis in 3d/2d (in degrees)
  angle can be a variable (see below)
\sphinxstyleemphasis{spherical} args = phi theta
  phi = azimuthal angle from +x axis (in degrees)
  theta = angle from +z or +y axis in 3d/2d (in degrees)
  phi or theta can be a variable (see below)
\sphinxstyleemphasis{vector} args = x y z
  x y z = vector direction to apply the acceleration
  x or y or z can be a variable (see below)
\end{sphinxalltt}

\end{itemize}


\subsection{Examples}
\label{\detokenize{fix_gravity:examples}}
\begin{sphinxVerbatim}[commandchars=\\\{\}]
\PYG{k}{fix }\PYG{n+nv+nvIdentifier}{1}\PYG{+w}{ }\PYG{n+nv+nvIdentifier}{all}\PYG{+w}{ }\PYG{n}{gravity}\PYG{+w}{ }\PYG{l+m}{1.0}\PYG{+w}{ }\PYG{n}{chute}\PYG{+w}{ }\PYG{l+m}{24.0}
\PYG{k}{fix }\PYG{n+nv+nvIdentifier}{1}\PYG{+w}{ }\PYG{n+nv+nvIdentifier}{all}\PYG{+w}{ }\PYG{n}{gravity}\PYG{+w}{ }\PYG{n}{v\PYGZus{}increase}\PYG{+w}{ }\PYG{n}{chute}\PYG{+w}{ }\PYG{l+m}{24.0}
\PYG{k}{fix }\PYG{n+nv+nvIdentifier}{1}\PYG{+w}{ }\PYG{n+nv+nvIdentifier}{all}\PYG{+w}{ }\PYG{n}{gravity}\PYG{+w}{ }\PYG{l+m}{1.0}\PYG{+w}{ }\PYG{n}{spherical}\PYG{+w}{ }\PYG{l+m}{0.0}\PYG{+w}{ }\PYG{o}{\PYGZhy{}}\PYG{l+m}{180.0}
\PYG{k}{fix }\PYG{n+nv+nvIdentifier}{1}\PYG{+w}{ }\PYG{n+nv+nvIdentifier}{all}\PYG{+w}{ }\PYG{n}{gravity}\PYG{+w}{ }\PYG{l+m}{10.0}\PYG{+w}{ }\PYG{n}{spherical}\PYG{+w}{ }\PYG{n}{v\PYGZus{}phi}\PYG{+w}{ }\PYG{n}{v\PYGZus{}theta}
\PYG{k}{fix }\PYG{n+nv+nvIdentifier}{1}\PYG{+w}{ }\PYG{n+nv+nvIdentifier}{all}\PYG{+w}{ }\PYG{n}{gravity}\PYG{+w}{ }\PYG{l+m}{100.0}\PYG{+w}{ }\PYG{n}{vector}\PYG{+w}{ }\PYG{l+m}{1}\PYG{+w}{ }\PYG{l+m}{1}\PYG{+w}{ }\PYG{l+m}{0}
\end{sphinxVerbatim}


\subsection{Description}
\label{\detokenize{fix_gravity:description}}
\sphinxAtStartPar
Impose an additional acceleration on each particle in the group.  This
fix is typically used with granular systems to include a “gravity”
term acting on the macroscopic particles.  More generally, it can
represent any kind of driving field, e.g. a pressure gradient inducing
a Poiseuille flow in a fluid.  Note that this fix operates differently
than the {\hyperref[\detokenize{fix_addforce::doc}]{\sphinxcrossref{\DUrole{doc}{fix addforce}}}} command.  The addforce fix
adds the same force to each atom, independent of its mass.  This
command imparts the same acceleration to each atom (force/mass).

\sphinxAtStartPar
The \sphinxstyleemphasis{magnitude} of the acceleration is specified in force/mass units.
For granular systems (LJ units) this is typically 1.0.  See the
{\hyperref[\detokenize{units::doc}]{\sphinxcrossref{\DUrole{doc}{units}}}} command for details.

\sphinxAtStartPar
Style \sphinxstyleemphasis{chute} is typically used for simulations of chute flow where
the specified \sphinxstyleemphasis{angle} is the chute angle, with flow occurring in the +x
direction.  For 3d systems, the tilt is away from the z axis; for 2d
systems, the tilt is away from the y axis.

\sphinxAtStartPar
Style \sphinxstyleemphasis{spherical} allows an arbitrary 3d direction to be specified for
the acceleration vector.  \sphinxstyleemphasis{Phi} and \sphinxstyleemphasis{theta} are defined in the usual
spherical coordinates.  Thus for acceleration acting in the \sphinxhyphen{}z
direction, \sphinxstyleemphasis{theta} would be 180.0 (or \sphinxhyphen{}180.0).  \sphinxstyleemphasis{Theta} = 90.0 and
\sphinxstyleemphasis{phi} = \sphinxhyphen{}90.0 would mean acceleration acts in the \sphinxhyphen{}y direction.  For
2d systems, \sphinxstyleemphasis{phi} is ignored and \sphinxstyleemphasis{theta} is an angle in the xy plane
where \sphinxstyleemphasis{theta} = 0.0 is the y\sphinxhyphen{}axis.

\sphinxAtStartPar
Style \sphinxstyleemphasis{vector} imposes an acceleration in the vector direction given
by (x,y,z).  Only the direction of the vector is important; it’s
length is ignored.  For 2d systems, the \sphinxstyleemphasis{z} component is ignored.

\sphinxAtStartPar
Any of the quantities \sphinxstyleemphasis{magnitude}, \sphinxstyleemphasis{angle}, \sphinxstyleemphasis{phi}, \sphinxstyleemphasis{theta}, \sphinxstyleemphasis{x}, \sphinxstyleemphasis{y},
\sphinxstyleemphasis{z} which define the gravitational magnitude and direction, can be
specified as an equal\sphinxhyphen{}style {\hyperref[\detokenize{variable::doc}]{\sphinxcrossref{\DUrole{doc}{variable}}}}.  If the value is
a variable, it should be specified as v\_name, where name is the
variable name.  In this case, the variable will be evaluated each
timestep, and its value used to determine the quantity.  You should
ensure that the variable calculates a result in the appropriate units,
e.g. force/mass or degrees.

\sphinxAtStartPar
Equal\sphinxhyphen{}style variables can specify formulas with various mathematical
functions, and include {\hyperref[\detokenize{thermo_style::doc}]{\sphinxcrossref{\DUrole{doc}{thermo\_style}}}} command
keywords for the simulation box parameters and timestep and elapsed
time.  Thus it is easy to specify a time\sphinxhyphen{}dependent gravitational
field.


\bigskip\hrule\bigskip


\sphinxAtStartPar
Styles with a \sphinxstyleemphasis{gpu}, \sphinxstyleemphasis{intel}, \sphinxstyleemphasis{kk}, \sphinxstyleemphasis{omp}, or \sphinxstyleemphasis{opt} suffix are
functionally the same as the corresponding style without the suffix.
They have been optimized to run faster, depending on your available
hardware, as discussed on the {\hyperref[\detokenize{Speed_packages::doc}]{\sphinxcrossref{\DUrole{doc}{Accelerator packages}}}}
page.  The accelerated styles take the same arguments and should
produce the same results, except for round\sphinxhyphen{}off and precision issues.

\sphinxAtStartPar
These accelerated styles are part of the GPU, INTEL, KOKKOS,
OPENMP, and OPT packages, respectively.  They are only enabled if
LAMMPS was built with those packages.  See the {\hyperref[\detokenize{Build_package::doc}]{\sphinxcrossref{\DUrole{doc}{Build package}}}} page for more info.

\sphinxAtStartPar
You can specify the accelerated styles explicitly in your input script
by including their suffix, or you can use the {\hyperref[\detokenize{Run_options::doc}]{\sphinxcrossref{\DUrole{doc}{\sphinxhyphen{}suffix command\sphinxhyphen{}line switch}}}} when you invoke LAMMPS, or you can use the
{\hyperref[\detokenize{suffix::doc}]{\sphinxcrossref{\DUrole{doc}{suffix}}}} command in your input script.

\sphinxAtStartPar
See the {\hyperref[\detokenize{Speed_packages::doc}]{\sphinxcrossref{\DUrole{doc}{Accelerator packages}}}} page for more
instructions on how to use the accelerated styles effectively.


\bigskip\hrule\bigskip



\subsection{Restart, fix\_modify, output, run start/stop, minimize info}
\label{\detokenize{fix_gravity:restart-fix-modify-output-run-start-stop-minimize-info}}
\sphinxAtStartPar
No information about this fix is written to {\hyperref[\detokenize{restart::doc}]{\sphinxcrossref{\DUrole{doc}{binary restart files}}}}.

\sphinxAtStartPar
The {\hyperref[\detokenize{fix_modify::doc}]{\sphinxcrossref{\DUrole{doc}{fix\_modify}}}} \sphinxstyleemphasis{energy} option is supported by
this fix to add the gravitational potential energy of the system to
the global potential energy of the system as part of
{\hyperref[\detokenize{thermo_style::doc}]{\sphinxcrossref{\DUrole{doc}{thermodynamic output}}}}.  The default setting for
this fix is {\hyperref[\detokenize{fix_modify::doc}]{\sphinxcrossref{\DUrole{doc}{fix\_modify energy no}}}}.

\sphinxAtStartPar
The {\hyperref[\detokenize{fix_modify::doc}]{\sphinxcrossref{\DUrole{doc}{fix\_modify}}}} \sphinxstyleemphasis{respa} option is supported by this
fix. This allows to set at which level of the {\hyperref[\detokenize{run_style::doc}]{\sphinxcrossref{\DUrole{doc}{r\sphinxhyphen{}RESPA}}}} integrator the fix is adding its forces. Default is the
outermost level.

\sphinxAtStartPar
This fix computes a global scalar which can be accessed by various
{\hyperref[\detokenize{Howto_output::doc}]{\sphinxcrossref{\DUrole{doc}{output commands}}}}.  This scalar is the
gravitational potential energy of the particles in the defined field,
namely mass * (g dot x) for each particles, where x and mass are the
particles position and mass, and g is the gravitational field.  The
scalar value calculated by this fix is “extensive”.

\sphinxAtStartPar
No parameter of this fix can be used with the \sphinxstyleemphasis{start/stop} keywords of
the {\hyperref[\detokenize{run::doc}]{\sphinxcrossref{\DUrole{doc}{run}}}} command.  This fix is not invoked during
{\hyperref[\detokenize{minimize::doc}]{\sphinxcrossref{\DUrole{doc}{energy minimization}}}}.


\subsection{Restrictions}
\label{\detokenize{fix_gravity:restrictions}}\begin{quote}

\sphinxAtStartPar
none
\end{quote}


\subsection{Related commands}
\label{\detokenize{fix_gravity:related-commands}}
\sphinxAtStartPar
{\hyperref[\detokenize{atom_style::doc}]{\sphinxcrossref{\DUrole{doc}{atom\_style sphere}}}}, {\hyperref[\detokenize{fix_addforce::doc}]{\sphinxcrossref{\DUrole{doc}{fix addforce}}}}


\subsection{Default}
\label{\detokenize{fix_gravity:default}}
\sphinxAtStartPar
none

\sphinxstepscope

\index{fix grem@\spxentry{fix grem}}\ignorespaces 

\section{fix grem command}
\label{\detokenize{fix_grem:fix-grem-command}}\label{\detokenize{fix_grem:index-0}}\label{\detokenize{fix_grem::doc}}

\subsection{Syntax}
\label{\detokenize{fix_grem:syntax}}
\begin{sphinxVerbatim}[commandchars=\\\{\}]
\PYG{k}{fix }\PYG{n+nv+nvIdentifier}{ID}\PYG{+w}{ }\PYG{n+nv+nvIdentifier}{group\PYGZhy{}ID}\PYG{+w}{ }\PYG{n}{grem}\PYG{+w}{ }\PYG{n}{lambda}\PYG{+w}{ }\PYG{n}{eta}\PYG{+w}{ }\PYG{n}{H0}\PYG{+w}{ }\PYG{n}{thermostat}\PYG{o}{\PYGZhy{}}\PYG{n}{ID}
\end{sphinxVerbatim}
\begin{itemize}
\item {} 
\sphinxAtStartPar
ID, group\sphinxhyphen{}ID are documented in {\hyperref[\detokenize{fix::doc}]{\sphinxcrossref{\DUrole{doc}{fix}}}} command

\item {} 
\sphinxAtStartPar
grem = style name of this fix command

\item {} 
\sphinxAtStartPar
lambda = intercept parameter of linear effective temperature function

\item {} 
\sphinxAtStartPar
eta = slope parameter of linear effective temperature function

\item {} 
\sphinxAtStartPar
H0 = shift parameter of linear effective temperature function

\item {} 
\sphinxAtStartPar
thermostat\sphinxhyphen{}ID = ID of Nose\sphinxhyphen{}Hoover thermostat or barostat used in simulation

\end{itemize}


\subsection{Examples}
\label{\detokenize{fix_grem:examples}}
\begin{sphinxVerbatim}[commandchars=\\\{\}]
\PYG{k}{fix             }\PYG{n+nv+nvIdentifier}{fxgREM}\PYG{+w}{ }\PYG{n+nv+nvIdentifier}{all}\PYG{+w}{ }\PYG{n}{grem}\PYG{+w}{ }\PYG{l+m}{400}\PYG{+w}{ }\PYG{o}{\PYGZhy{}}\PYG{l+m}{0.01}\PYG{+w}{ }\PYG{o}{\PYGZhy{}}\PYG{l+m}{30000}\PYG{+w}{ }\PYG{n}{fxnpt}
\PYG{k}{thermo\PYGZus{}modify}\PYG{+w}{   }\PYG{n}{press}\PYG{+w}{ }\PYG{n}{fxgREM\PYGZus{}press}

\PYG{k}{fix             }\PYG{n+nv+nvIdentifier}{fxgREM}\PYG{+w}{ }\PYG{n+nv+nvIdentifier}{all}\PYG{+w}{ }\PYG{n}{grem}\PYG{+w}{ }\PYG{l+m}{502}\PYG{+w}{ }\PYG{o}{\PYGZhy{}}\PYG{l+m}{0.15}\PYG{+w}{ }\PYG{o}{\PYGZhy{}}\PYG{l+m}{80000}\PYG{+w}{ }\PYG{n}{fxnvt}
\end{sphinxVerbatim}


\subsection{Description}
\label{\detokenize{fix_grem:description}}
\sphinxAtStartPar
This fix implements the molecular dynamics version of the generalized
replica exchange method (gREM) originally developed by {\hyperref[\detokenize{fix_grem:kim2010}]{\sphinxcrossref{\DUrole{std,std-ref}{(Kim)}}}},
which uses non\sphinxhyphen{}Boltzmann ensembles to sample over first order phase
transitions. The is done by defining replicas with an enthalpy
dependent effective temperature
\begin{equation*}
\begin{split}T_{eff} = \lambda + \eta (H - H_0)\end{split}
\end{equation*}
\sphinxAtStartPar
with \(\eta\) negative and steep enough to only intersect the
characteristic microcanonical temperature (Ts) of the system once,
ensuring a unimodal enthalpy distribution in that replica.
\(\lambda\) is the intercept and effects the generalized ensemble
similar to how temperature effects a Boltzmann ensemble. \(H_0\)
is a reference enthalpy, and is typically set as the lowest desired
sampled enthalpy.  Further explanation can be found in our recent
papers {\hyperref[\detokenize{fix_grem:malolepsza}]{\sphinxcrossref{\DUrole{std,std-ref}{(Malolepsza)}}}}.

\sphinxAtStartPar
This fix requires a Nose\sphinxhyphen{}Hoover thermostat fix reference passed to the
grem as \sphinxstyleemphasis{thermostat\sphinxhyphen{}ID}. Two distinct temperatures exist in this
generalized ensemble, the effective temperature defined above, and a
kinetic temperature that controls the velocity distribution of
particles as usual. Either constant volume or constant pressure
algorithms can be used.

\sphinxAtStartPar
The fix enforces a generalized ensemble in a single replica
only. Typically, this ideology is combined with replica exchange with
replicas differing by \(\lambda\) only for simplicity, but this is not
required. A multi\sphinxhyphen{}replica simulation can be run within the LAMMPS
environment using the {\hyperref[\detokenize{temper_grem::doc}]{\sphinxcrossref{\DUrole{doc}{temper/grem}}}} command. This
utilizes LAMMPS partition mode and requires the number of available
processors be on the order of the number of desired replicas. A
100\sphinxhyphen{}replica simulation would require at least 100 processors (1 per
world at minimum). If many replicas are needed on a small number of
processors, multi\sphinxhyphen{}replica runs can be run outside of LAMMPS.  An
example of this can be found in examples/PACKAGES/grem and has no
limit on the number of replicas per processor. However, this is very
inefficient and error prone and should be avoided if possible.

\sphinxAtStartPar
In general, defining the generalized ensembles is unique for every
system. When starting a many\sphinxhyphen{}replica simulation without any knowledge
of the underlying microcanonical temperature, there are several tricks
we have utilized to optimize the process.  Choosing a less\sphinxhyphen{}steep
\(\eta\) yields broader distributions, requiring fewer replicas to
map the microcanonical temperature.  While this likely struggles from
the same sampling problems gREM was built to avoid, it provides quick
insight to Ts.  Initially using an evenly\sphinxhyphen{}spaced \(\lambda\)
distribution identifies regions where small changes in enthalpy lead
to large temperature changes. Replicas are easily added where needed.


\bigskip\hrule\bigskip



\subsection{Restart, fix\_modify, output, run start/stop, minimize info}
\label{\detokenize{fix_grem:restart-fix-modify-output-run-start-stop-minimize-info}}
\sphinxAtStartPar
No information about this fix is written to {\hyperref[\detokenize{restart::doc}]{\sphinxcrossref{\DUrole{doc}{binary restart files}}}}.

\sphinxAtStartPar
The {\hyperref[\detokenize{thermo_modify::doc}]{\sphinxcrossref{\DUrole{doc}{thermo\_modify}}}} \sphinxstyleemphasis{press} option is supported
by this fix to add the rescaled kinetic pressure as part of
{\hyperref[\detokenize{thermo_style::doc}]{\sphinxcrossref{\DUrole{doc}{thermodynamic output}}}}.

\sphinxAtStartPar
This fix computes a global scalar which can be accessed by various
{\hyperref[\detokenize{Howto_output::doc}]{\sphinxcrossref{\DUrole{doc}{output commands}}}}.  The scalar is the effective
temperature \(T_{eff}\).  The scalar value calculated by this
fix is “intensive”.


\subsection{Restrictions}
\label{\detokenize{fix_grem:restrictions}}
\sphinxAtStartPar
This fix is part of the REPLICA package. It is only enabled if
LAMMPS was built with that package. See the {\hyperref[\detokenize{Build_package::doc}]{\sphinxcrossref{\DUrole{doc}{Build package}}}} page for more info.


\subsection{Related commands}
\label{\detokenize{fix_grem:related-commands}}
\sphinxAtStartPar
{\hyperref[\detokenize{temper_grem::doc}]{\sphinxcrossref{\DUrole{doc}{temper/grem}}}}, {\hyperref[\detokenize{fix_nh::doc}]{\sphinxcrossref{\DUrole{doc}{fix nvt}}}}, {\hyperref[\detokenize{fix_nh::doc}]{\sphinxcrossref{\DUrole{doc}{fix npt}}}}, {\hyperref[\detokenize{thermo_modify::doc}]{\sphinxcrossref{\DUrole{doc}{thermo\_modify}}}}


\subsection{Default}
\label{\detokenize{fix_grem:default}}
\sphinxAtStartPar
none


\bigskip\hrule\bigskip

\phantomsection\label{\detokenize{fix_grem:kim2010}}
\sphinxAtStartPar
\sphinxstylestrong{(Kim)} Kim, Keyes, Straub, J Chem. Phys, 132, 224107 (2010).

\phantomsection\label{\detokenize{fix_grem:malolepsza}}
\sphinxAtStartPar
\sphinxstylestrong{(Malolepsza)} Malolepsza, Secor, Keyes, J Phys Chem B 119 (42),
13379\sphinxhyphen{}13384 (2015).

\sphinxstepscope

\index{fix halt@\spxentry{fix halt}}\ignorespaces 

\section{fix halt command}
\label{\detokenize{fix_halt:fix-halt-command}}\label{\detokenize{fix_halt:index-0}}\label{\detokenize{fix_halt::doc}}

\subsection{Syntax}
\label{\detokenize{fix_halt:syntax}}
\begin{sphinxVerbatim}[commandchars=\\\{\}]
\PYG{k}{fix }\PYG{n+nv+nvIdentifier}{ID}\PYG{+w}{ }\PYG{n+nv+nvIdentifier}{group\PYGZhy{}ID}\PYG{+w}{ }\PYG{n}{halt}\PYG{+w}{ }\PYG{n}{N}\PYG{+w}{ }\PYG{n}{attribute}\PYG{+w}{ }\PYG{n}{operator}\PYG{+w}{ }\PYG{n}{avalue}\PYG{+w}{ }\PYG{n}{keyword}\PYG{+w}{ }\PYG{n}{value}\PYG{+w}{ }\PYG{n}{...}
\end{sphinxVerbatim}
\begin{itemize}
\item {} 
\sphinxAtStartPar
ID, group\sphinxhyphen{}ID are documented in {\hyperref[\detokenize{fix::doc}]{\sphinxcrossref{\DUrole{doc}{fix}}}} command

\item {} 
\sphinxAtStartPar
halt = style name of this fix command

\item {} 
\sphinxAtStartPar
N = check halt condition every N steps

\item {} 
\sphinxAtStartPar
attribute = \sphinxstyleemphasis{bondmax} or \sphinxstyleemphasis{tlimit} or v\_name
\begin{sphinxalltt}
bondmax = length of longest bond in the system (in length units)
tlimit = elapsed CPU time (in seconds)
diskfree = free disk space (in MBytes)
v\_name = name of {\hyperref[\detokenize{variable::doc}]{\sphinxcrossref{\DUrole{doc}{equal\sphinxhyphen{}style variable}}}}
\end{sphinxalltt}

\item {} 
\sphinxAtStartPar
operator = “\textless{}” or “\textless{}=” or “\textgreater{}” or “\textgreater{}=” or “==” or “!=” or “|\textasciicircum{}”

\item {} 
\sphinxAtStartPar
avalue = numeric value to compare attribute to

\item {} 
\sphinxAtStartPar
zero or more keyword/value pairs may be appended

\item {} 
\sphinxAtStartPar
keyword = \sphinxstyleemphasis{error} or \sphinxstyleemphasis{message} or \sphinxstyleemphasis{path}
\begin{sphinxalltt}
\sphinxstyleemphasis{error} value = \sphinxstyleemphasis{hard} or \sphinxstyleemphasis{soft} or \sphinxstyleemphasis{continue}
\sphinxstyleemphasis{message} value = \sphinxstyleemphasis{yes} or \sphinxstyleemphasis{no}
\sphinxstyleemphasis{path} value = path to check for free space (may be in quotes)
\end{sphinxalltt}

\end{itemize}


\subsection{Examples}
\label{\detokenize{fix_halt:examples}}
\begin{sphinxVerbatim}[commandchars=\\\{\}]
\PYG{k}{fix }\PYG{n+nv+nvIdentifier}{10}\PYG{+w}{ }\PYG{n+nv+nvIdentifier}{all}\PYG{+w}{ }\PYG{n}{halt}\PYG{+w}{ }\PYG{l+m}{1}\PYG{+w}{ }\PYG{n}{bondmax}\PYG{+w}{ }\PYG{o}{\PYGZgt{}}\PYG{+w}{ }\PYG{l+m}{1.5}
\PYG{k}{fix }\PYG{n+nv+nvIdentifier}{10}\PYG{+w}{ }\PYG{n+nv+nvIdentifier}{all}\PYG{+w}{ }\PYG{n}{halt}\PYG{+w}{ }\PYG{l+m}{10}\PYG{+w}{ }\PYG{n}{v\PYGZus{}myCheck}\PYG{+w}{ }\PYG{o}{!}\PYG{o}{=}\PYG{+w}{ }\PYG{l+m}{0}\PYG{+w}{ }\PYG{n}{error}\PYG{+w}{ }\PYG{n}{soft}
\PYG{k}{fix }\PYG{n+nv+nvIdentifier}{10}\PYG{+w}{ }\PYG{n+nv+nvIdentifier}{all}\PYG{+w}{ }\PYG{n}{halt}\PYG{+w}{ }\PYG{l+m}{100}\PYG{+w}{ }\PYG{n}{diskfree}\PYG{+w}{ }\PYG{o}{\PYGZlt{}}\PYG{+w}{ }\PYG{l+m}{100000.0}\PYG{+w}{ }\PYG{n}{path}\PYG{+w}{ }\PYG{l+s}{\PYGZdq{}}\PYG{l+s}{dump storage/.}\PYG{l+s}{\PYGZdq{}}
\end{sphinxVerbatim}


\subsection{Description}
\label{\detokenize{fix_halt:description}}
\sphinxAtStartPar
Check a condition every N steps during a simulation run.  N must be \textgreater{}=1.
If the condition is met, exit the run.  In this context a “run” can be
dynamics or minimization iterations, as specified by the {\hyperref[\detokenize{run::doc}]{\sphinxcrossref{\DUrole{doc}{run}}}} or {\hyperref[\detokenize{minimize::doc}]{\sphinxcrossref{\DUrole{doc}{minimize}}}} command.

\sphinxAtStartPar
The specified group\sphinxhyphen{}ID is ignored by this fix.

\sphinxAtStartPar
The specified \sphinxstyleemphasis{attribute} can be one of the options listed above, namely
\sphinxstyleemphasis{bondmax}, \sphinxstyleemphasis{tlimit}, \sphinxstyleemphasis{diskfree}, or an {\hyperref[\detokenize{variable::doc}]{\sphinxcrossref{\DUrole{doc}{equal\sphinxhyphen{}style variable}}}} referenced as \sphinxstyleemphasis{v\_name}, where “name” is the name of a
variable that has been defined previously in the input script.

\sphinxAtStartPar
The \sphinxstyleemphasis{bondmax} attribute will loop over all bonds in the system,
compute their current lengths, and set \sphinxstyleemphasis{attribute} to the longest bond
distance.

\sphinxAtStartPar
The \sphinxstyleemphasis{tlimit} attribute queries the elapsed CPU time (in seconds) since
the current run began, and sets \sphinxstyleemphasis{attribute} to that value.  This is an
alternative way to limit the length of a simulation run, similar to
the {\hyperref[\detokenize{timer::doc}]{\sphinxcrossref{\DUrole{doc}{timer}}}} timeout command.  There are two differences in
using this method versus the timer command option.  The first is that
the clock starts at the beginning of the current run (not when the
timer or fix command is specified), so that any setup time for the run
is not included in the elapsed time.  The second is that the timer
invocation and syncing across all processors (via MPI\_Allreduce) is
not performed once every \sphinxstyleemphasis{N} steps by this command.  Instead it is
performed (typically) only a small number of times and the elapsed
times are used to predict when the end\sphinxhyphen{}of\sphinxhyphen{}the\sphinxhyphen{}run will be.  Both of
these attributes can be useful when performing benchmark calculations
for a desired length of time with minimal overhead.  For example, if
a run is performing 1000s of timesteps/sec, the overhead for syncing
the timer frequently across a large number of processors may be
non\sphinxhyphen{}negligible.

\sphinxAtStartPar
The \sphinxstyleemphasis{diskfree} attribute will check for available disk space (in
MBytes) on supported operating systems. By default it will
check the file system of the current working directory.  This
can be changed with the optional \sphinxstyleemphasis{path} keyword, which will take
the path to a file or folder on the file system to be checked
as argument.  This path must be given with single or double quotes,
if it contains blanks or other special characters (like \$).

\sphinxAtStartPar
Equal\sphinxhyphen{}style variables evaluate to a numeric value.  See the
{\hyperref[\detokenize{variable::doc}]{\sphinxcrossref{\DUrole{doc}{variable}}}} command for a description.  They calculate
formulas which can involve mathematical operations, atom properties,
group properties, thermodynamic properties, global values calculated
by a {\hyperref[\detokenize{compute::doc}]{\sphinxcrossref{\DUrole{doc}{compute}}}} or {\hyperref[\detokenize{fix::doc}]{\sphinxcrossref{\DUrole{doc}{fix}}}}, or references to other
{\hyperref[\detokenize{variable::doc}]{\sphinxcrossref{\DUrole{doc}{variables}}}}.  Thus they are a very general means of
computing some attribute of the current system.  For example, the
following “bondmax” variable will calculate the same quantity as the
hstyle = bondmax option.

\begin{sphinxVerbatim}[commandchars=\\\{\}]
\PYG{k}{compute         }\PYG{n+nv+nvIdentifier}{bdist}\PYG{+w}{ }\PYG{n+nv+nvIdentifier}{all}\PYG{+w}{ }\PYG{n}{bond}\PYG{o}{/}\PYG{n}{local}\PYG{+w}{ }\PYG{n}{dist}
\PYG{k}{compute         }\PYG{n+nv+nvIdentifier}{bmax}\PYG{+w}{ }\PYG{n+nv+nvIdentifier}{all}\PYG{+w}{ }\PYG{n}{reduce}\PYG{+w}{ }\PYG{n}{max}\PYG{+w}{ }\PYG{n}{c\PYGZus{}bdist}
\PYG{k}{variable        }\PYG{n+nv+nvIdentifier}{bondmax}\PYG{+w}{ }\PYG{n}{equal}\PYG{+w}{ }\PYG{n}{c\PYGZus{}bmax}
\end{sphinxVerbatim}

\sphinxAtStartPar
Thus these two versions of a fix halt command will do the same thing:

\begin{sphinxVerbatim}[commandchars=\\\{\}]
\PYG{k}{fix }\PYG{n+nv+nvIdentifier}{10}\PYG{+w}{ }\PYG{n+nv+nvIdentifier}{all}\PYG{+w}{ }\PYG{n}{halt}\PYG{+w}{ }\PYG{l+m}{1}\PYG{+w}{ }\PYG{n}{bondmax}\PYG{+w}{ }\PYG{o}{\PYGZgt{}}\PYG{+w}{ }\PYG{l+m}{1.5}
\PYG{k}{fix }\PYG{n+nv+nvIdentifier}{10}\PYG{+w}{ }\PYG{n+nv+nvIdentifier}{all}\PYG{+w}{ }\PYG{n}{halt}\PYG{+w}{ }\PYG{l+m}{1}\PYG{+w}{ }\PYG{n}{v\PYGZus{}bondmax}\PYG{+w}{ }\PYG{o}{\PYGZgt{}}\PYG{+w}{ }\PYG{l+m}{1.5}
\end{sphinxVerbatim}

\sphinxAtStartPar
The version with “bondmax” will just run somewhat faster, due to less
overhead in computing bond lengths and not storing them in a separate
compute.

\sphinxAtStartPar
A variable can be used to implement a large variety of conditions,
including to stop when a specific file exists.  Example:

\begin{sphinxVerbatim}[commandchars=\\\{\}]
\PYG{k}{variable }\PYG{n+nv+nvIdentifier}{exit}\PYG{+w}{ }\PYG{n}{equal}\PYG{+w}{ }\PYG{n}{is\PYGZus{}file}\PYG{n+nv}{(}\PYG{n+nv}{EXIT}\PYG{n+nv}{)}
\PYG{k}{fix }\PYG{n+nv+nvIdentifier}{10}\PYG{+w}{ }\PYG{n+nv+nvIdentifier}{all}\PYG{+w}{ }\PYG{n}{halt}\PYG{+w}{ }\PYG{l+m}{100}\PYG{+w}{ }\PYG{n}{v\PYGZus{}exit}\PYG{+w}{ }\PYG{o}{!}\PYG{o}{=}\PYG{+w}{ }\PYG{l+m}{0}\PYG{+w}{ }\PYG{n}{error}\PYG{+w}{ }\PYG{n}{soft}
\end{sphinxVerbatim}

\sphinxAtStartPar
Will stop the current run command when a file \sphinxcode{\sphinxupquote{EXIT}} is created
in the current working directory.  The condition can be cleared
by removing the file through the {\hyperref[\detokenize{shell::doc}]{\sphinxcrossref{\DUrole{doc}{shell}}}} command.

\sphinxAtStartPar
The choice of operators listed above are the usual comparison
operators.  The XOR operation (exclusive or) is also included as “|\textasciicircum{}”.
In this context, XOR means that if either the attribute or avalue is
0.0 and the other is non\sphinxhyphen{}zero, then the result is “true”.  Otherwise
it is “false”.

\sphinxAtStartPar
The specified \sphinxstyleemphasis{avalue} must be a numeric value.


\bigskip\hrule\bigskip


\sphinxAtStartPar
The optional \sphinxstyleemphasis{error} keyword determines how the current run is halted.
If its value is \sphinxstyleemphasis{hard}, then LAMMPS will stop with an error message.

\sphinxAtStartPar
If its value is \sphinxstyleemphasis{soft}, LAMMPS will exit the current run, but continue
to execute subsequent commands in the input script.  However,
additional {\hyperref[\detokenize{run::doc}]{\sphinxcrossref{\DUrole{doc}{run}}}} or {\hyperref[\detokenize{minimize::doc}]{\sphinxcrossref{\DUrole{doc}{minimize}}}} commands will be
skipped.  For example, this allows a script to output the current
state of the system, e.g. via a {\hyperref[\detokenize{write_dump::doc}]{\sphinxcrossref{\DUrole{doc}{write\_dump}}}} or
{\hyperref[\detokenize{write_restart::doc}]{\sphinxcrossref{\DUrole{doc}{write\_restart}}}} command.

\sphinxAtStartPar
If its value is \sphinxstyleemphasis{continue}, the behavior is the same as for \sphinxstyleemphasis{soft},
except subsequent {\hyperref[\detokenize{run::doc}]{\sphinxcrossref{\DUrole{doc}{run}}}} or {\hyperref[\detokenize{minimize::doc}]{\sphinxcrossref{\DUrole{doc}{minimize}}}} commands
are executed.  This allows your script to remedy the condition that
triggered the halt, if necessary.  Note that you may wish use the
{\hyperref[\detokenize{unfix::doc}]{\sphinxcrossref{\DUrole{doc}{unfix}}}} command on the fix halt ID, so that the same
condition is not immediately triggered in a subsequent run.

\sphinxAtStartPar
The optional \sphinxstyleemphasis{message} keyword determines whether a message is printed
to the screen and logfile when the halt condition is triggered.  If
\sphinxstyleemphasis{message} is set to yes, a one line message with the values that
triggered the halt is printed.  If \sphinxstyleemphasis{message} is set to no, no message
is printed; the run simply exits.  The latter may be desirable for
post\sphinxhyphen{}processing tools that extract thermodynamic information from log
files.


\subsection{Restart, fix\_modify, output, run start/stop, minimize info}
\label{\detokenize{fix_halt:restart-fix-modify-output-run-start-stop-minimize-info}}
\sphinxAtStartPar
No information about this fix is written to {\hyperref[\detokenize{restart::doc}]{\sphinxcrossref{\DUrole{doc}{binary restart files}}}}.  None of the {\hyperref[\detokenize{fix_modify::doc}]{\sphinxcrossref{\DUrole{doc}{fix\_modify}}}} options
are relevant to this fix.  No global or per\sphinxhyphen{}atom quantities are stored
by this fix for access by various {\hyperref[\detokenize{Howto_output::doc}]{\sphinxcrossref{\DUrole{doc}{output commands}}}}.
No parameter of this fix can be used with the \sphinxstyleemphasis{start/stop} keywords of
the {\hyperref[\detokenize{run::doc}]{\sphinxcrossref{\DUrole{doc}{run}}}} command.


\subsection{Restrictions}
\label{\detokenize{fix_halt:restrictions}}
\sphinxAtStartPar
The \sphinxstyleemphasis{diskfree} attribute is currently only supported on Linux, macOS, and *BSD.


\subsection{Related commands}
\label{\detokenize{fix_halt:related-commands}}
\sphinxAtStartPar
{\hyperref[\detokenize{variable::doc}]{\sphinxcrossref{\DUrole{doc}{variable}}}}


\subsection{Default}
\label{\detokenize{fix_halt:default}}
\sphinxAtStartPar
The option defaults are error = soft, message = yes, and path = “.”.

\sphinxstepscope

\index{fix heat@\spxentry{fix heat}}\ignorespaces 

\section{fix heat command}
\label{\detokenize{fix_heat:fix-heat-command}}\label{\detokenize{fix_heat:index-0}}\label{\detokenize{fix_heat::doc}}

\subsection{Syntax}
\label{\detokenize{fix_heat:syntax}}
\begin{sphinxVerbatim}[commandchars=\\\{\}]
\PYG{k}{fix }\PYG{n+nv+nvIdentifier}{ID}\PYG{+w}{ }\PYG{n+nv+nvIdentifier}{group\PYGZhy{}ID}\PYG{+w}{ }\PYG{n}{heat}\PYG{+w}{ }\PYG{n}{N}\PYG{+w}{ }\PYG{n}{eflux}
\end{sphinxVerbatim}
\begin{itemize}
\item {} 
\sphinxAtStartPar
ID, group\sphinxhyphen{}ID are documented in {\hyperref[\detokenize{fix::doc}]{\sphinxcrossref{\DUrole{doc}{fix}}}} command

\item {} 
\sphinxAtStartPar
heat = style name of this fix command

\item {} 
\sphinxAtStartPar
N = add/subtract heat every this many timesteps

\item {} 
\sphinxAtStartPar
eflux = rate of heat addition or subtraction (energy/time units)

\item {} 
\sphinxAtStartPar
eflux can be a variable (see below)

\item {} 
\sphinxAtStartPar
zero or more keyword/value pairs may be appended to args

\item {} 
\sphinxAtStartPar
keyword = \sphinxstyleemphasis{region}
\begin{sphinxalltt}
\sphinxstyleemphasis{region} value = region\sphinxhyphen{}ID
  region\sphinxhyphen{}ID = ID of region atoms must be in to have added force
\end{sphinxalltt}

\end{itemize}


\subsection{Examples}
\label{\detokenize{fix_heat:examples}}
\begin{sphinxVerbatim}[commandchars=\\\{\}]
\PYG{k}{fix }\PYG{n+nv+nvIdentifier}{3}\PYG{+w}{ }\PYG{n+nv+nvIdentifier}{qin}\PYG{+w}{ }\PYG{n}{heat}\PYG{+w}{ }\PYG{l+m}{1}\PYG{+w}{ }\PYG{l+m}{1.0}
\PYG{k}{fix }\PYG{n+nv+nvIdentifier}{3}\PYG{+w}{ }\PYG{n+nv+nvIdentifier}{qin}\PYG{+w}{ }\PYG{n}{heat}\PYG{+w}{ }\PYG{l+m}{10}\PYG{+w}{ }\PYG{n}{v\PYGZus{}flux}
\PYG{k}{fix }\PYG{n+nv+nvIdentifier}{4}\PYG{+w}{ }\PYG{n+nv+nvIdentifier}{qout}\PYG{+w}{ }\PYG{n}{heat}\PYG{+w}{ }\PYG{l+m}{1}\PYG{+w}{ }\PYG{o}{\PYGZhy{}}\PYG{l+m}{1.0}\PYG{+w}{ }\PYG{k}{region }\PYG{n+nv+nvIdentifier}{top}
\end{sphinxVerbatim}


\subsection{Description}
\label{\detokenize{fix_heat:description}}
\sphinxAtStartPar
Add non\sphinxhyphen{}translational kinetic energy (heat) to a group of atoms in a
manner that conserves their aggregate momentum.  Two of these fixes
can be used to establish a temperature gradient across a simulation
domain by adding heat (energy) to one group of atoms (hot reservoir)
and subtracting heat from another (cold reservoir).  E.g. a simulation
sampling from the McDLT ensemble.

\sphinxAtStartPar
If the \sphinxstyleemphasis{region} keyword is used, the atom must be in both the group
and the specified geometric {\hyperref[\detokenize{region::doc}]{\sphinxcrossref{\DUrole{doc}{region}}}} in order to have
energy added or subtracted to it.  If not specified, then the atoms in
the group are affected wherever they may move to.

\sphinxAtStartPar
Heat addition/subtraction is performed every N timesteps.

\sphinxAtStartPar
The \sphinxstyleemphasis{eflux} parameter can be specified as a numeric constant or as an
equal\sphinxhyphen{} or atom\sphinxhyphen{}style {\hyperref[\detokenize{variable::doc}]{\sphinxcrossref{\DUrole{doc}{variable}}}}.  If the value is a
variable, it should be specified as v\_name, where \sphinxstyleemphasis{name} is the variable
name.  In this case, the variable will be evaluated each timestep, and
its current value(s) used to determine the flux.

\sphinxAtStartPar
If \sphinxstyleemphasis{eflux} is a numeric constant or equal\sphinxhyphen{}style variable which evaluates
to a scalar value, then \sphinxstyleemphasis{eflux} determines the change in aggregate energy
of the entire group of atoms per unit time, e.g. in eV/ps for
{\hyperref[\detokenize{units::doc}]{\sphinxcrossref{\DUrole{doc}{metal units}}}}.  In this case it is an “extensive” quantity,
meaning its magnitude should be scaled with the number of atoms in the
group.  Note that since \sphinxstyleemphasis{eflux} also has per\sphinxhyphen{}time units (i.e. it is a
flux), this means that a larger value of N will add/subtract a larger
amount of energy each time the fix is invoked.

\begin{sphinxadmonition}{note}{Note:}
\sphinxAtStartPar
The heat\sphinxhyphen{}exchange (HEX) algorithm implemented by this fix is
known to exhibit a pronounced energy drift. An improved algorithm
(eHEX) is available as a {\hyperref[\detokenize{fix_ehex::doc}]{\sphinxcrossref{\DUrole{doc}{fix ehex}}}} command and might be
preferable if energy conservation is important.
\end{sphinxadmonition}

\sphinxAtStartPar
If \sphinxstyleemphasis{eflux} is specified as an atom\sphinxhyphen{}style variable (see below), then
the variable computes one value per atom.  In this case, each value is
the energy flux for a single atom, again in units of energy per unit
time.  In this case, each value is an “intensive” quantity, which need
not be scaled with the number of atoms in the group.

\sphinxAtStartPar
Equal\sphinxhyphen{}style variables can specify formulas with various mathematical
functions, and include {\hyperref[\detokenize{thermo_style::doc}]{\sphinxcrossref{\DUrole{doc}{thermo\_style}}}} command
keywords for the simulation box parameters and timestep and elapsed
time.  Thus it is easy to specify a time\sphinxhyphen{}dependent flux.

\sphinxAtStartPar
Atom\sphinxhyphen{}style variables can specify the same formulas as equal\sphinxhyphen{}style
variables but can also include per\sphinxhyphen{}atom values, such as atom
coordinates.  Thus it is easy to specify a spatially\sphinxhyphen{}dependent flux
with optional time\sphinxhyphen{}dependence as well.

\begin{sphinxadmonition}{note}{Note:}
\sphinxAtStartPar
If heat is subtracted from the system too aggressively so that
the group’s kinetic energy would go to zero, or any individual atom’s
kinetic energy would go to zero for the case where \sphinxstyleemphasis{eflux} is an
atom\sphinxhyphen{}style variable, then LAMMPS will halt with an error message.
\end{sphinxadmonition}

\sphinxAtStartPar
Fix heat is different from a thermostat such as {\hyperref[\detokenize{fix_nh::doc}]{\sphinxcrossref{\DUrole{doc}{fix nvt}}}}
or {\hyperref[\detokenize{fix_temp_rescale::doc}]{\sphinxcrossref{\DUrole{doc}{fix temp/rescale}}}} in that energy is
added/subtracted continually.  Thus if there is not another mechanism
in place to counterbalance this effect, the entire system will heat or
cool continuously.  You can use multiple heat fixes so that the net
energy change is 0.0 or use {\hyperref[\detokenize{fix_viscous::doc}]{\sphinxcrossref{\DUrole{doc}{fix viscous}}}} to drain
energy from the system.

\sphinxAtStartPar
This fix does not change the coordinates of its atoms; it only scales
their velocities.  Thus you must still use an integration fix
(e.g. {\hyperref[\detokenize{fix_nve::doc}]{\sphinxcrossref{\DUrole{doc}{fix nve}}}}) on the affected atoms.  This fix should
not normally be used on atoms that have their temperature controlled
by another fix \sphinxhyphen{} e.g. {\hyperref[\detokenize{fix_nh::doc}]{\sphinxcrossref{\DUrole{doc}{fix nvt}}}} or {\hyperref[\detokenize{fix_langevin::doc}]{\sphinxcrossref{\DUrole{doc}{fix langevin}}}} fix.


\subsection{Restart, fix\_modify, output, run start/stop, minimize info}
\label{\detokenize{fix_heat:restart-fix-modify-output-run-start-stop-minimize-info}}
\sphinxAtStartPar
No information about this fix is written to {\hyperref[\detokenize{restart::doc}]{\sphinxcrossref{\DUrole{doc}{binary restart files}}}}.  None of the {\hyperref[\detokenize{fix_modify::doc}]{\sphinxcrossref{\DUrole{doc}{fix\_modify}}}} options
are relevant to this fix.

\sphinxAtStartPar
This fix computes a global scalar which can be accessed by various
{\hyperref[\detokenize{Howto_output::doc}]{\sphinxcrossref{\DUrole{doc}{output commands}}}}.  This scalar is the most recent
value by which velocities were scaled.  The scalar value calculated by
this fix is “intensive”.  If \sphinxstyleemphasis{eflux} is specified as an atom\sphinxhyphen{}style
variable, this fix computes the average value by which the velocities
were scaled for all of the atoms that had their velocities scaled.

\sphinxAtStartPar
No parameter of this fix can be used with the \sphinxstyleemphasis{start/stop} keywords of
the {\hyperref[\detokenize{run::doc}]{\sphinxcrossref{\DUrole{doc}{run}}}} command.  This fix is not invoked during {\hyperref[\detokenize{minimize::doc}]{\sphinxcrossref{\DUrole{doc}{energy minimization}}}}.


\subsection{Restrictions}
\label{\detokenize{fix_heat:restrictions}}\begin{quote}

\sphinxAtStartPar
none
\end{quote}


\subsection{Related commands}
\label{\detokenize{fix_heat:related-commands}}
\sphinxAtStartPar
{\hyperref[\detokenize{fix_ehex::doc}]{\sphinxcrossref{\DUrole{doc}{fix ehex}}}}, {\hyperref[\detokenize{compute_temp::doc}]{\sphinxcrossref{\DUrole{doc}{compute temp}}}}, {\hyperref[\detokenize{compute_temp_region::doc}]{\sphinxcrossref{\DUrole{doc}{compute temp/region}}}}


\subsection{Default}
\label{\detokenize{fix_heat:default}}
\sphinxAtStartPar
none

\sphinxstepscope

\index{fix heat/flow@\spxentry{fix heat/flow}}\ignorespaces 

\section{fix heat/flow command}
\label{\detokenize{fix_heat_flow:fix-heat-flow-command}}\label{\detokenize{fix_heat_flow:index-0}}\label{\detokenize{fix_heat_flow::doc}}

\subsection{Syntax}
\label{\detokenize{fix_heat_flow:syntax}}
\begin{sphinxVerbatim}[commandchars=\\\{\}]
\PYG{k}{fix }\PYG{n+nv+nvIdentifier}{ID}\PYG{+w}{ }\PYG{n+nv+nvIdentifier}{group\PYGZhy{}ID}\PYG{+w}{ }\PYG{n}{heat}\PYG{o}{/}\PYG{n}{flow}\PYG{+w}{ }\PYG{n}{style}\PYG{+w}{ }\PYG{n}{values}\PYG{+w}{ }\PYG{n}{...}
\end{sphinxVerbatim}
\begin{itemize}
\item {} 
\sphinxAtStartPar
ID, group\sphinxhyphen{}ID are documented in {\hyperref[\detokenize{fix::doc}]{\sphinxcrossref{\DUrole{doc}{fix}}}} command

\item {} 
\sphinxAtStartPar
heat/flow = style name of this fix command

\item {} 
\sphinxAtStartPar
one style with corresponding value(s) needs to be listed
\begin{sphinxalltt}
style = \sphinxstyleemphasis{constant} or \sphinxstyleemphasis{type}
  \sphinxstyleemphasis{constant} = cp
    cp = value of specifc heat (energy/(mass * temperature) units)
  \sphinxstyleemphasis{type} = cp1 ... cpN
    cpN = value of specifc heat for type N (energy/(mass * temperature) units)
\end{sphinxalltt}

\item {} 
\end{itemize}


\subsection{Examples}
\label{\detokenize{fix_heat_flow:examples}}
\begin{sphinxVerbatim}[commandchars=\\\{\}]
\PYG{k}{fix }\PYG{n+nv+nvIdentifier}{1}\PYG{+w}{ }\PYG{n+nv+nvIdentifier}{all}\PYG{+w}{ }\PYG{n}{heat}\PYG{o}{/}\PYG{n}{flow}\PYG{+w}{ }\PYG{n}{constant}\PYG{+w}{ }\PYG{l+m}{1.0}
\PYG{k}{fix }\PYG{n+nv+nvIdentifier}{1}\PYG{+w}{ }\PYG{n+nv+nvIdentifier}{all}\PYG{+w}{ }\PYG{n}{heat}\PYG{o}{/}\PYG{n}{flow}\PYG{+w}{ }\PYG{n}{type}\PYG{+w}{ }\PYG{l+m}{1.0}\PYG{+w}{ }\PYG{l+m}{0.5}
\end{sphinxVerbatim}


\subsection{Description}
\label{\detokenize{fix_heat_flow:description}}
\sphinxAtStartPar
Perform plain time integration to update temperature for atoms in the
group each timestep. The specific heat of atoms can be defined using either
the \sphinxstyleemphasis{constant} or \sphinxstyleemphasis{type} keywords. For style \sphinxstyleemphasis{constant}, the specific heat
is a constant value \sphinxstyleemphasis{cp} for all atoms. For style \sphinxstyleemphasis{type}, \sphinxstyleemphasis{N} different values
of the specific heat are defined, one for each of the \sphinxstyleemphasis{N} types of atoms.


\bigskip\hrule\bigskip



\subsection{Restart, fix\_modify, output, run start/stop, minimize info}
\label{\detokenize{fix_heat_flow:restart-fix-modify-output-run-start-stop-minimize-info}}
\sphinxAtStartPar
No information about this fix is written to {\hyperref[\detokenize{restart::doc}]{\sphinxcrossref{\DUrole{doc}{binary restart files}}}}.
None of the {\hyperref[\detokenize{fix_modify::doc}]{\sphinxcrossref{\DUrole{doc}{fix\_modify}}}} options are relevant to this fix.
No global or per\sphinxhyphen{}atom quantities are stored by this fix for access by various
{\hyperref[\detokenize{Howto_output::doc}]{\sphinxcrossref{\DUrole{doc}{output commands}}}}. No parameter of this fix can be used
with the \sphinxstyleemphasis{start/stop} keywords of the {\hyperref[\detokenize{run::doc}]{\sphinxcrossref{\DUrole{doc}{run}}}} command.  This fix is
not invoked during {\hyperref[\detokenize{minimize::doc}]{\sphinxcrossref{\DUrole{doc}{energy minimization}}}}.


\subsection{Restrictions}
\label{\detokenize{fix_heat_flow:restrictions}}
\sphinxAtStartPar
This pair style is part of the GRANULAR package.  It is
only enabled if LAMMPS was built with that package.
See the {\hyperref[\detokenize{Build_package::doc}]{\sphinxcrossref{\DUrole{doc}{Build package}}}} page for more info.

\sphinxAtStartPar
This fix requires that atoms store temperature and heat flow
as defined by the {\hyperref[\detokenize{fix_property_atom::doc}]{\sphinxcrossref{\DUrole{doc}{fix property/atom}}}} command.


\subsection{Related commands}
\label{\detokenize{fix_heat_flow:related-commands}}
\sphinxAtStartPar
{\hyperref[\detokenize{pair_granular::doc}]{\sphinxcrossref{\DUrole{doc}{pair granular}}}},
{\hyperref[\detokenize{fix_add_heat::doc}]{\sphinxcrossref{\DUrole{doc}{fix add/heat}}}},
{\hyperref[\detokenize{fix_property_atom::doc}]{\sphinxcrossref{\DUrole{doc}{fix property/atom}}}}


\subsection{Default}
\label{\detokenize{fix_heat_flow:default}}
\sphinxAtStartPar
none

\sphinxstepscope

\index{fix hyper/global@\spxentry{fix hyper/global}}\ignorespaces 

\section{fix hyper/global command}
\label{\detokenize{fix_hyper_global:fix-hyper-global-command}}\label{\detokenize{fix_hyper_global:index-0}}\label{\detokenize{fix_hyper_global::doc}}

\subsection{Syntax}
\label{\detokenize{fix_hyper_global:syntax}}
\begin{sphinxVerbatim}[commandchars=\\\{\}]
\PYG{k}{fix }\PYG{n+nv+nvIdentifier}{ID}\PYG{+w}{ }\PYG{n+nv+nvIdentifier}{group\PYGZhy{}ID}\PYG{+w}{ }\PYG{n}{hyper}\PYG{o}{/}\PYG{n}{global}\PYG{+w}{ }\PYG{n}{cutbond}\PYG{+w}{ }\PYG{n}{qfactor}\PYG{+w}{ }\PYG{n}{Vmax}\PYG{+w}{ }\PYG{n}{Tequil}
\end{sphinxVerbatim}
\begin{itemize}
\item {} 
\sphinxAtStartPar
ID, group\sphinxhyphen{}ID are documented in {\hyperref[\detokenize{fix::doc}]{\sphinxcrossref{\DUrole{doc}{fix}}}} command

\item {} 
\sphinxAtStartPar
hyper/global = style name of this fix command

\item {} 
\sphinxAtStartPar
cutbond = max distance at which a pair of atoms is considered bonded (distance units)

\item {} 
\sphinxAtStartPar
qfactor = max strain at which bias potential goes to 0.0 (unitless)

\item {} 
\sphinxAtStartPar
Vmax = height of bias potential (energy units)

\item {} 
\sphinxAtStartPar
Tequil = equilibration temperature (temperature units)

\end{itemize}


\subsection{Examples}
\label{\detokenize{fix_hyper_global:examples}}
\begin{sphinxVerbatim}[commandchars=\\\{\}]
\PYG{k}{fix }\PYG{n+nv+nvIdentifier}{1}\PYG{+w}{ }\PYG{n+nv+nvIdentifier}{all}\PYG{+w}{ }\PYG{n}{hyper}\PYG{o}{/}\PYG{n}{global}\PYG{+w}{ }\PYG{l+m}{1.0}\PYG{+w}{ }\PYG{l+m}{0.3}\PYG{+w}{ }\PYG{l+m}{0.8}\PYG{+w}{ }\PYG{l+m}{300.0}
\end{sphinxVerbatim}


\subsection{Description}
\label{\detokenize{fix_hyper_global:description}}
\sphinxAtStartPar
This fix is meant to be used with the {\hyperref[\detokenize{hyper::doc}]{\sphinxcrossref{\DUrole{doc}{hyper}}}} command to
perform a bond\sphinxhyphen{}boost global hyperdynamics (GHD) simulation.  The role
of this fix is to a select a single pair of atoms in the system at
each timestep to add a global bias potential to, which will alter the
dynamics of the system in a manner that effectively accelerates time.
This is in contrast to the {\hyperref[\detokenize{fix_hyper_local::doc}]{\sphinxcrossref{\DUrole{doc}{fix hyper/local}}}}
command, which can be user to perform a local hyperdynamics (LHD)
simulation, by adding a local bias potential to multiple pairs of
atoms at each timestep.  GHD can time accelerate a small simulation
with up to a few 100 atoms.  For larger systems, LHD is needed to
achieve good time acceleration.

\sphinxAtStartPar
For a system that undergoes rare transition events, where one or more
atoms move over an energy barrier to a new potential energy basin, the
effect of the bias potential is to induce more rapid transitions.
This can lead to a dramatic speed\sphinxhyphen{}up in the rate at which events
occurs, without altering their relative frequencies, thus leading to
an overall increase in the elapsed real time of the simulation as
compared to running for the same number of timesteps with normal MD.
See the {\hyperref[\detokenize{hyper::doc}]{\sphinxcrossref{\DUrole{doc}{hyper}}}} page for a more general discussion of
hyperdynamics and citations that explain both GHD and LHD.

\sphinxAtStartPar
The equations and logic used by this fix and described here to perform
GHD follow the description given in {\hyperref[\detokenize{fix_hyper_global:voter2013ghd}]{\sphinxcrossref{\DUrole{std,std-ref}{(Voter2013)}}}}.  The
bond\sphinxhyphen{}boost form of a bias potential for HD is due to Miron and
Fichthorn as described in {\hyperref[\detokenize{fix_hyper_global:mironghd}]{\sphinxcrossref{\DUrole{std,std-ref}{(Miron)}}}}.  In LAMMPS we use a
simplified version of bond\sphinxhyphen{}boost GHD where a single bond in the system
is biased at any one timestep.

\sphinxAtStartPar
Bonds are defined between each pair of atoms \sphinxstyleemphasis{ij}, whose \(R^0_{ij}\)
distance is less than \sphinxstyleemphasis{cutbond}, when the system is in a quenched state
(minimum) energy.  Note that these are not “bonds” in a covalent
sense.  A bond is simply any pair of atoms that meet the distance
criterion.  \sphinxstyleemphasis{Cutbond} is an argument to this fix; it is discussed
below.  A bond is only formed if one or both of the \sphinxstyleemphasis{ij} atoms are in
the specified group.

\sphinxAtStartPar
The current strain of bond \sphinxstyleemphasis{ij} (when running dynamics) is defined as
\begin{equation*}
\begin{split}E_{ij} = \frac{R_{ij} - R^0_{ij}}{R^0_{ij}}\end{split}
\end{equation*}
\sphinxAtStartPar
where \(R_{ij}\) is the current distance between atoms \sphinxstyleemphasis{i} and \sphinxstyleemphasis{j},
and \(R^0_{ij}\) is the equilibrium distance in the quenched state.

\sphinxAtStartPar
The bias energy \(V_{ij}\) of any bond between atoms \sphinxstyleemphasis{i} and \sphinxstyleemphasis{j}
is defined as
\begin{equation*}
\begin{split}V_{ij} = V^{max} \cdot \left( 1 - \left(\frac{E_{ij}}{q}\right)^2 \right) \textrm{ for } \left|E_{ij}\right| < qfactor \textrm{ or } 0 \textrm{ otherwise}\end{split}
\end{equation*}
\sphinxAtStartPar
where the prefactor \(V^{max}\) and the cutoff \sphinxstyleemphasis{qfactor} are arguments to
this fix; they are discussed below.  This functional form is an
inverse parabola centered at 0.0 with height \(V^{max}\) and
which goes to 0.0 at +/\sphinxhyphen{} qfactor.

\sphinxAtStartPar
Let \(E^{max}\) be the maximum of \(\left| E_{ij} \right|\)
for all \sphinxstyleemphasis{ij} bonds in the system on a
given timestep.  On that step, \(V_{ij}\) is added as a bias potential
to only the single bond with strain \(E^{max}\), call it
\(V^{max}_{ij}\).  Note that \(V^{max}_{ij}\) will be 0.0
if \(E^{max} >= \textrm{qfactor}\) on that timestep.  Also note
that \(V^{max}_{ij}\) is added to the normal interatomic potential
that is computed between all atoms in the system at every step.

\sphinxAtStartPar
The derivative of \(V^{max}_{ij}\) with respect to the position of
each atom in the \(E^{max}\) bond gives a bias force
\(F^{max}_{ij}\) acting on the bond as
\begin{equation*}
\begin{split}F^{max}_{ij} = - \frac{dV^{max}_{ij}}{dE_{ij}} = \frac{2 V^{max} E-{ij}}{\textrm{qfactor}^2}   \textrm{ for } \left|E_{ij}\right| < \textrm{qfactor} \textrm{ or } 0 \textrm{ otherwise}\end{split}
\end{equation*}
\sphinxAtStartPar
which can be decomposed into an equal and opposite force acting on
only the two \sphinxstyleemphasis{ij} atoms in the \(E^{max}\) bond.

\sphinxAtStartPar
The time boost factor for the system is given each timestep I by
\begin{equation*}
\begin{split}B_i = e^{\beta V^{max}_{ij}}\end{split}
\end{equation*}
\sphinxAtStartPar
where \(\beta = \frac{1}{kT_{equil}}\), and \(T_{equil}\) is the
temperature of the system and an argument to this fix.  Note that
\(B_i >= 1\) at every step.

\begin{sphinxadmonition}{note}{Note:}
\sphinxAtStartPar
To run a GHD simulation, the input script must also use the {\hyperref[\detokenize{fix_langevin::doc}]{\sphinxcrossref{\DUrole{doc}{fix
langevin}}}} command to thermostat the atoms at the same
\(T_{equil}\) as specified by this fix, so that the system is
running constant\sphinxhyphen{}temperature (NVT) dynamics.  LAMMPS does not check
that this is done.
\end{sphinxadmonition}

\sphinxAtStartPar
The elapsed time \(t_{hyper}\) for a GHD simulation running for \sphinxstyleemphasis{N}
timesteps is simply
\begin{equation*}
\begin{split}t_{hyper} = \sum_{i=1,N} B-i \cdot dt\end{split}
\end{equation*}
\sphinxAtStartPar
where \sphinxstyleemphasis{dt} is the timestep size defined by the {\hyperref[\detokenize{timestep::doc}]{\sphinxcrossref{\DUrole{doc}{timestep}}}}
command.  The effective time acceleration due to GHD is thus t\_hyper /
N*dt, where N*dt is elapsed time for a normal MD run of N timesteps.

\sphinxAtStartPar
Note that in GHD, the boost factor varies from timestep to timestep.
Likewise, which bond has \(E^{max}\) strain and thus which pair of
atoms the bias potential is added to, will also vary from timestep to timestep.
This is in contrast to local hyperdynamics (LHD) where the boost
factor is an input parameter; see the {\hyperref[\detokenize{fix_hyper_local::doc}]{\sphinxcrossref{\DUrole{doc}{fix hyper/local}}}} page for details.


\bigskip\hrule\bigskip


\sphinxAtStartPar
Here is additional information on the input parameters for GHD.

\sphinxAtStartPar
The \sphinxstyleemphasis{cutbond} argument is the cutoff distance for defining bonds between
pairs of nearby atoms.  A pair of \sphinxstyleemphasis{ij} atoms in their equilibrium,
minimum\sphinxhyphen{}energy configuration, which are separated by a distance
\(R_{ij} < cutbond\), are flagged as a bonded pair.  Setting
\sphinxstyleemphasis{cubond} to be \textasciitilde{}25\% larger than the nearest\sphinxhyphen{}neighbor distance in a
crystalline lattice is a typical choice for solids, so that bonds exist
only between nearest neighbor pairs.

\sphinxAtStartPar
The \sphinxstyleemphasis{qfactor} argument is the limiting strain at which the bias
potential goes to 0.0.  It is dimensionless, so a value of 0.3 means a
bond distance can be up to 30\% larger or 30\% smaller than the
equilibrium (quenched) \(R^0_{ij}\) distance and the two atoms in
the bond could still experience a non\sphinxhyphen{}zero bias force.

\sphinxAtStartPar
If \sphinxstyleemphasis{qfactor} is set too large, then transitions from one energy basin
to another are affected because the bias potential is non\sphinxhyphen{}zero at the
transition state (e.g. saddle point).  If \sphinxstyleemphasis{qfactor} is set too small
than little boost is achieved because the \(E_{ij}\) strain of some bond in
the system will (nearly) always exceed \sphinxstyleemphasis{qfactor}.  A value of 0.3 for
\sphinxstyleemphasis{qfactor} is typically reasonable.

\sphinxAtStartPar
The \sphinxstyleemphasis{Vmax} argument is the prefactor on the bias potential.  Ideally,
tt should be set to a value slightly less than the smallest barrier
height for an event to occur.  Otherwise the applied bias potential
may be large enough (when added to the interatomic potential) to
produce a local energy basin with a maxima in the center.  This can
produce artificial energy minima in the same basin that trap an atom.
Or if \sphinxstyleemphasis{Vmax} is even larger, it may induce an atom(s) to rapidly
transition to another energy basin.  Both cases are “bad dynamics”
which violate the assumptions of GHD that guarantee an accelerated
time\sphinxhyphen{}accurate trajectory of the system.

\sphinxAtStartPar
Note that if \sphinxstyleemphasis{Vmax} is set too small, the GHD simulation will run
correctly.  There will just be fewer events because the hyper time
(t\_hyper equation above) will be shorter.

\begin{sphinxadmonition}{note}{Note:}
\sphinxAtStartPar
If you have no physical intuition as to the smallest barrier
height in your system, a reasonable strategy to determine the largest
\sphinxstyleemphasis{Vmax} you can use for a GHD model, is to run a sequence of
simulations with smaller and smaller \sphinxstyleemphasis{Vmax} values, until the event
rate does not change (as a function of hyper time).
\end{sphinxadmonition}

\sphinxAtStartPar
The \sphinxstyleemphasis{Tequil} argument is the temperature at which the system is
simulated; see the comment above about the {\hyperref[\detokenize{fix_langevin::doc}]{\sphinxcrossref{\DUrole{doc}{fix langevin}}}} thermostatting.  It is also part of the beta term in the
exponential factor that determines how much boost is achieved as a
function of the bias potential.

\sphinxAtStartPar
In general, the lower the value of \sphinxstyleemphasis{Tequil} and the higher the value
of \sphinxstyleemphasis{Vmax}, the more time boost will be achievable by the GHD
algorithm.


\bigskip\hrule\bigskip



\subsection{Restart, fix\_modify, output, run start/stop, minimize info}
\label{\detokenize{fix_hyper_global:restart-fix-modify-output-run-start-stop-minimize-info}}
\sphinxAtStartPar
No information about this fix is written to {\hyperref[\detokenize{restart::doc}]{\sphinxcrossref{\DUrole{doc}{binary restart files}}}}.

\sphinxAtStartPar
The {\hyperref[\detokenize{fix_modify::doc}]{\sphinxcrossref{\DUrole{doc}{fix\_modify}}}} \sphinxstyleemphasis{energy} option is supported by
this fix to add the energy of the bias potential to the global
potential energy of the system as part of {\hyperref[\detokenize{thermo_style::doc}]{\sphinxcrossref{\DUrole{doc}{thermodynamic output}}}}.  The default setting for this fix is {\hyperref[\detokenize{fix_modify::doc}]{\sphinxcrossref{\DUrole{doc}{fix\_modify
energy no}}}}.

\sphinxAtStartPar
This fix computes a global scalar and global vector of length 12,
which can be accessed by various {\hyperref[\detokenize{Howto_output::doc}]{\sphinxcrossref{\DUrole{doc}{output commands}}}}.  The scalar is the magnitude of the bias potential
(energy units) applied on the current timestep.  The vector stores the
following quantities:
\begin{enumerate}
\sphinxsetlistlabels{\arabic}{enumi}{enumii}{}{.}%
\item {} 
\sphinxAtStartPar
boost factor on this step (unitless)

\item {} 
\sphinxAtStartPar
max strain \(E_{ij}\) of any bond on this step (absolute value, unitless)

\item {} 
\sphinxAtStartPar
ID of first atom in the max\sphinxhyphen{}strain bond

\item {} 
\sphinxAtStartPar
ID of second atom in the max\sphinxhyphen{}strain bond

\item {} 
\sphinxAtStartPar
average \# of bonds/atom on this step

\item {} 
\sphinxAtStartPar
fraction of timesteps where the biased bond has bias = 0.0 during this run

\item {} 
\sphinxAtStartPar
fraction of timesteps where the biased bond has negative strain during this run

\item {} 
\sphinxAtStartPar
max drift distance of any atom during this run (distance units)

\item {} 
\sphinxAtStartPar
max bond length during this run (distance units)

\item {} 
\sphinxAtStartPar
cumulative hyper time since fix was defined (time units)

\item {} 
\sphinxAtStartPar
cumulative count of event timesteps since fix was defined

\item {} 
\sphinxAtStartPar
cumulative count of atoms in events since fix was defined

\end{enumerate}

\sphinxAtStartPar
The first 5 quantities are for the current timestep.  Quantities 6\sphinxhyphen{}9
are for the current hyper run.  They are reset each time a new hyper
run is performed.  Quantities 19\sphinxhyphen{}12 are cumulative across multiple
runs (since the point in the input script the fix was defined).

\sphinxAtStartPar
For value 8, drift is the distance an atom moves between two quenched
states when the second quench determines an event has occurred.  Atoms
involved in an event will typically move the greatest distance since
others typically remain near their original quenched position.

\sphinxAtStartPar
For value 11, events are checked for by the {\hyperref[\detokenize{hyper::doc}]{\sphinxcrossref{\DUrole{doc}{hyper}}}} command
once every \sphinxstyleemphasis{Nevent} timesteps.  This value is the count of those
timesteps on which one (or more) events was detected.  It is NOT the
number of distinct events, since more than one event may occur in the
same \sphinxstyleemphasis{Nevent} time window.

\sphinxAtStartPar
For value 12, each time the {\hyperref[\detokenize{hyper::doc}]{\sphinxcrossref{\DUrole{doc}{hyper}}}} command checks for an
event, it invokes a compute to flag zero or more atoms as
participating in one or more events.  E.g. atoms that have displaced
more than some distance from the previous quench state.  Value 11 is
the cumulative count of the number of atoms participating in any of
the events that were found.

\sphinxAtStartPar
The scalar and vector values calculated by this fix are all
“intensive”.

\sphinxAtStartPar
No parameter of this fix can be used with the \sphinxstyleemphasis{start/stop} keywords of
the {\hyperref[\detokenize{run::doc}]{\sphinxcrossref{\DUrole{doc}{run}}}} command.  This fix is not invoked during
{\hyperref[\detokenize{minimize::doc}]{\sphinxcrossref{\DUrole{doc}{energy minimization}}}}.


\subsection{Restrictions}
\label{\detokenize{fix_hyper_global:restrictions}}
\sphinxAtStartPar
This command can only be used if LAMMPS was built with the REPLICA
package.  See the {\hyperref[\detokenize{Build_package::doc}]{\sphinxcrossref{\DUrole{doc}{Build package}}}} page for more
info.


\subsection{Related commands}
\label{\detokenize{fix_hyper_global:related-commands}}
\sphinxAtStartPar
{\hyperref[\detokenize{hyper::doc}]{\sphinxcrossref{\DUrole{doc}{hyper}}}}, {\hyperref[\detokenize{fix_hyper_local::doc}]{\sphinxcrossref{\DUrole{doc}{fix hyper/local}}}}


\subsection{Default}
\label{\detokenize{fix_hyper_global:default}}
\sphinxAtStartPar
none


\bigskip\hrule\bigskip

\phantomsection\label{\detokenize{fix_hyper_global:voter2013ghd}}
\sphinxAtStartPar
\sphinxstylestrong{(Voter2013)} S. Y. Kim, D. Perez, A. F. Voter, J Chem Phys, 139,
144110 (2013).

\phantomsection\label{\detokenize{fix_hyper_global:mironghd}}
\sphinxAtStartPar
\sphinxstylestrong{(Miron)} R. A. Miron and K. A. Fichthorn, J Chem Phys, 119, 6210 (2003).

\sphinxstepscope

\index{fix hyper/local@\spxentry{fix hyper/local}}\ignorespaces 

\section{fix hyper/local command}
\label{\detokenize{fix_hyper_local:fix-hyper-local-command}}\label{\detokenize{fix_hyper_local:index-0}}\label{\detokenize{fix_hyper_local::doc}}

\subsection{Syntax}
\label{\detokenize{fix_hyper_local:syntax}}
\begin{sphinxVerbatim}[commandchars=\\\{\}]
\PYG{k}{fix }\PYG{n+nv+nvIdentifier}{ID}\PYG{+w}{ }\PYG{n+nv+nvIdentifier}{group\PYGZhy{}ID}\PYG{+w}{ }\PYG{n}{hyper}\PYG{o}{/}\PYG{n}{local}\PYG{+w}{ }\PYG{n}{cutbond}\PYG{+w}{ }\PYG{n}{qfactor}\PYG{+w}{ }\PYG{n}{Vmax}\PYG{+w}{ }\PYG{n}{Tequil}\PYG{+w}{ }\PYG{n}{Dcut}\PYG{+w}{ }\PYG{n}{alpha}\PYG{+w}{ }\PYG{n}{Btarget}
\end{sphinxVerbatim}
\begin{itemize}
\item {} 
\sphinxAtStartPar
ID, group\sphinxhyphen{}ID are documented in {\hyperref[\detokenize{fix::doc}]{\sphinxcrossref{\DUrole{doc}{fix}}}} command

\item {} 
\sphinxAtStartPar
hyper/local = style name of this fix command

\item {} 
\sphinxAtStartPar
cutbond = max distance at which a pair of atoms is considered bonded (distance units)

\item {} 
\sphinxAtStartPar
qfactor = max strain at which bias potential goes to 0.0 (unitless)

\item {} 
\sphinxAtStartPar
Vmax = estimated height of bias potential (energy units)

\item {} 
\sphinxAtStartPar
Tequil = equilibration temperature (temperature units)

\item {} 
\sphinxAtStartPar
Dcut = minimum distance between boosted bonds (distance units)

\item {} 
\sphinxAtStartPar
alpha = boostostat relaxation time (time units)

\item {} 
\sphinxAtStartPar
Btarget = desired time boost factor (unitless)

\item {} 
\sphinxAtStartPar
zero or more keyword/value pairs may be appended

\item {} 
\sphinxAtStartPar
keyword = \sphinxstyleemphasis{bound} or \sphinxstyleemphasis{reset} or \sphinxstyleemphasis{check/ghost} or \sphinxstyleemphasis{check/bias}
\begin{sphinxalltt}
\sphinxstyleemphasis{bound} value = Bfrac
  Bfrac =  \sphinxhyphen{}1 or a value \textgreater{}= 0.0
\sphinxstyleemphasis{reset} value = Rfreq
  Rfreq = \sphinxhyphen{}1 or 0 or timestep value \textgreater{} 0
\sphinxstyleemphasis{check/ghost} values = none
\sphinxstyleemphasis{check/bias} values = Nevery error/warn/ignore
\end{sphinxalltt}

\end{itemize}


\subsection{Examples}
\label{\detokenize{fix_hyper_local:examples}}
\begin{sphinxVerbatim}[commandchars=\\\{\}]
\PYG{k}{fix }\PYG{n+nv+nvIdentifier}{1}\PYG{+w}{ }\PYG{n+nv+nvIdentifier}{all}\PYG{+w}{ }\PYG{n}{hyper}\PYG{o}{/}\PYG{n}{local}\PYG{+w}{ }\PYG{l+m}{1.0}\PYG{+w}{ }\PYG{l+m}{0.3}\PYG{+w}{ }\PYG{l+m}{0.8}\PYG{+w}{ }\PYG{l+m}{300.0}
\PYG{k}{fix }\PYG{n+nv+nvIdentifier}{1}\PYG{+w}{ }\PYG{n+nv+nvIdentifier}{all}\PYG{+w}{ }\PYG{n}{hyper}\PYG{o}{/}\PYG{n}{local}\PYG{+w}{ }\PYG{l+m}{1.0}\PYG{+w}{ }\PYG{l+m}{0.3}\PYG{+w}{ }\PYG{l+m}{0.8}\PYG{+w}{ }\PYG{l+m}{300.0}\PYG{+w}{ }\PYG{n}{bound}\PYG{+w}{ }\PYG{l+m}{0.1}\PYG{+w}{ }\PYG{n}{reset}\PYG{+w}{ }\PYG{l+m}{0}
\end{sphinxVerbatim}


\subsection{Description}
\label{\detokenize{fix_hyper_local:description}}
\sphinxAtStartPar
This fix is meant to be used with the {\hyperref[\detokenize{hyper::doc}]{\sphinxcrossref{\DUrole{doc}{hyper}}}} command to
perform a bond\sphinxhyphen{}boost local hyperdynamics (LHD) simulation.  The role
of this fix is to a select multiple pairs of atoms in the system at
each timestep to add a local bias potential to, which will alter the
dynamics of the system in a manner that effectively accelerates time.
This is in contrast to the {\hyperref[\detokenize{fix_hyper_global::doc}]{\sphinxcrossref{\DUrole{doc}{fix hyper/global}}}}
command, which can be user to perform a global hyperdynamics (GHD)
simulation, by adding a global bias potential to a single pair of
atoms at each timestep.  GHD can time accelerate a small simulation
with up to a few 100 atoms.  For larger systems, LHD is needed to
achieve good time acceleration.

\sphinxAtStartPar
For a system that undergoes rare transition events, where one or more
atoms move over an energy barrier to a new potential energy basin, the
effect of the bias potential is to induce more rapid transitions.
This can lead to a dramatic speed\sphinxhyphen{}up in the rate at which events
occurs, without altering their relative frequencies, thus leading to
an overall increase in the elapsed real time of the simulation as
compared to running for the same number of timesteps with normal MD.
See the {\hyperref[\detokenize{hyper::doc}]{\sphinxcrossref{\DUrole{doc}{hyper}}}} page for a more general discussion of
hyperdynamics and citations that explain both GHD and LHD.

\sphinxAtStartPar
The equations and logic used by this fix and described here to perform
LHD follow the description given in {\hyperref[\detokenize{fix_hyper_local:voter2013lhd}]{\sphinxcrossref{\DUrole{std,std-ref}{(Voter2013)}}}}.  The
bond\sphinxhyphen{}boost form of a bias potential for HD is due to Miron and
Fichthorn as described in {\hyperref[\detokenize{fix_hyper_local:mironlhd}]{\sphinxcrossref{\DUrole{std,std-ref}{(Miron)}}}}.

\sphinxAtStartPar
To understand this description, you should first read the description
of the GHD algorithm on the {\hyperref[\detokenize{fix_hyper_global::doc}]{\sphinxcrossref{\DUrole{doc}{fix hyper/global}}}}
doc page.  This description of LHD builds on the GHD description.

\sphinxAtStartPar
The definition of bonds and \(E_{ij}\) are the same for GHD and LHD.
The formulas for \(V^{max}_{ij}\) and \(F^{max}_{ij}\) are also
the same except for a prefactor \(C_{ij}\), explained below.

\sphinxAtStartPar
The bias energy \(V_{ij}\) applied to a bond \sphinxstyleemphasis{ij} with maximum strain is
\begin{equation*}
\begin{split}V^{max}_{ij} = C_{ij} \cdot V^{max} \cdot \left(1 - \left(\frac{E_{ij}}{q}\right)^2\right) \textrm{ for } \left|E_{ij}\right| < qfactor \textrm{ or } 0 \textrm{ otherwise}\end{split}
\end{equation*}
\sphinxAtStartPar
The derivative of \(V^{max}_{ij}\) with respect to the position of
each atom in the \sphinxstyleemphasis{ij} bond gives a bias force \(F^{max}_{ij}\) acting
on the bond as
\begin{equation*}
\begin{split}F^{max}_{ij} = - \frac{dV^{max}_{ij}}{dE_{ij}} = 2 C_{ij} V^{max} \frac{E_{ij}}{qfactor^2} \textrm{ for } \left|E_{ij}\right| < qfactor \textrm{ or } 0 \textrm{ otherwise}\end{split}
\end{equation*}
\sphinxAtStartPar
which can be decomposed into an equal and opposite force acting on
only the two atoms \sphinxstyleemphasis{i} and \sphinxstyleemphasis{j} in the \sphinxstyleemphasis{ij} bond.

\sphinxAtStartPar
The key difference is that in GHD a bias energy and force is added (on
a particular timestep) to only one bond (pair of atoms) in the system,
which is the bond with maximum strain \(E^{max}\).

\sphinxAtStartPar
In LHD, a bias energy and force can be added to multiple bonds
separated by the specified \sphinxstyleemphasis{Dcut} distance or more.  A bond \sphinxstyleemphasis{ij} is
biased if it is the maximum strain bond within its local
“neighborhood”, which is defined as the bond \sphinxstyleemphasis{ij} plus any neighbor
bonds within a distance \sphinxstyleemphasis{Dcut} from \sphinxstyleemphasis{ij}.  The “distance” between bond
\sphinxstyleemphasis{ij} and bond \sphinxstyleemphasis{kl} is the minimum distance between any of the \sphinxstyleemphasis{ik}, \sphinxstyleemphasis{il},
\sphinxstyleemphasis{jk}, and \sphinxstyleemphasis{jl} pairs of atoms.

\sphinxAtStartPar
For a large system, multiple bonds will typically meet this
requirement, and thus a bias potential \(V^{max}_{ij}\) will be
applied to many bonds on the same timestep.

\sphinxAtStartPar
In LHD, all bonds store a \(C_{ij}\) prefactor which appears in
the \(V^{max}_{ij}\) and \(F^{max}_{ij}equations above.  Note
that the :math:`C_{ij}\) factor scales the strength of the bias energy
and forces whenever bond \sphinxstyleemphasis{ij} is the maximum strain bond in its neighborhood.

\sphinxAtStartPar
\(C_{ij}\) is initialized to 1.0 when a bond between the \sphinxstyleemphasis{ij} atoms
is first defined.  The specified \sphinxstyleemphasis{Btarget} factor is then used to adjust the
\(C_{ij}\) prefactors for each bond every timestep in the following manner.

\sphinxAtStartPar
An instantaneous boost factor \(B_{ij}\) is computed each timestep
for each bond, as
\begin{equation*}
\begin{split}B_{ij} = e^{\beta V^{max}_{kl}}\end{split}
\end{equation*}
\sphinxAtStartPar
where \(V^{max}_{kl}\) is the bias energy of the maxstrain bond \sphinxstyleemphasis{kl}
within bond \sphinxstyleemphasis{ij}‘s neighborhood, \(\beta = \frac{1}{kT_{equil}}\),
and \(T_{equil}\) is the temperature of the system and an argument
to this fix.

\begin{sphinxadmonition}{note}{Note:}
\sphinxAtStartPar
To run an LHD simulation, the input script must also use the
{\hyperref[\detokenize{fix_langevin::doc}]{\sphinxcrossref{\DUrole{doc}{fix langevin}}}} command to thermostat the atoms at
the same \(T_{equil}\) as specified by this fix, so that the
system is running constant\sphinxhyphen{}temperature (NVT) dynamics.  LAMMPS does
not check that this is done.
\end{sphinxadmonition}

\sphinxAtStartPar
Note that if \sphinxstyleemphasis{ij}== \sphinxstyleemphasis{kl}, then bond \sphinxstyleemphasis{ij} is a biased bond on that
timestep, otherwise it is not.  But regardless, the boost factor
\(B_{ij}\) can be thought of an estimate of time boost currently
being applied within a local region centered on bond \sphinxstyleemphasis{ij}.  For LHD,
we want this to be the specified \sphinxstyleemphasis{Btarget} value everywhere in the
simulation domain.

\sphinxAtStartPar
To accomplish this, if \(B_{ij} < B_{target}\), the \(C_{ij}\)
prefactor for bond \sphinxstyleemphasis{ij} is incremented on the current timestep by an
amount proportional to the inverse of the specified \(\alpha\) and
the difference (\(B_{ij} - B_{target}\)).  Conversely if
\(B_{ij} > B_{target}\), \(C_{ij}\) is decremented by the same
amount.  This procedure is termed “boostostatting” in {\hyperref[\detokenize{fix_hyper_local:voter2013lhd}]{\sphinxcrossref{\DUrole{std,std-ref}{(Voter2013)}}}}.  It drives all of the individual \(C_{ij}\) to
values such that when \(V^{max}_{ij}\) is applied as a bias to bond
\sphinxstyleemphasis{ij}, the resulting boost factor \(B_{ij}\) will be close to
\(B_{target}\) on average.  Thus the LHD time acceleration factor
for the overall system is effectively \sphinxstyleemphasis{Btarget}.

\sphinxAtStartPar
Note that in LHD, the boost factor \(B_{target}\) is specified by the user.
This is in contrast to global hyperdynamics (GHD) where the boost
factor varies each timestep and is computed as a function of \(V_{max}\),
\(E_{max}\), and \(T_{equil}\); see the
{\hyperref[\detokenize{fix_hyper_global::doc}]{\sphinxcrossref{\DUrole{doc}{fix hyper/global}}}} page for details.


\bigskip\hrule\bigskip


\sphinxAtStartPar
Here is additional information on the input parameters for LHD.

\sphinxAtStartPar
Note that the \sphinxstyleemphasis{cutbond}, \sphinxstyleemphasis{qfactor}, and \sphinxstyleemphasis{Tequil} arguments have the
same meaning as for GHD.  The \sphinxstyleemphasis{Vmax} argument is slightly different.
The \sphinxstyleemphasis{Dcut}, \sphinxstyleemphasis{alpha}, and \sphinxstyleemphasis{Btarget} parameters are unique to LHD.

\sphinxAtStartPar
The \sphinxstyleemphasis{cutbond} argument is the cutoff distance for defining bonds
between pairs of nearby atoms.  A pair of I,J atoms in their
equilibrium, minimum\sphinxhyphen{}energy configuration, which are separated by a
distance \(R_{ij} < cutbond\), are flagged as a bonded pair.  Setting
\sphinxstyleemphasis{cubond} to be \textasciitilde{}25\% larger than the nearest\sphinxhyphen{}neighbor distance in a
crystalline lattice is a typical choice for solids, so that bonds
exist only between nearest neighbor pairs.

\sphinxAtStartPar
The \sphinxstyleemphasis{qfactor} argument is the limiting strain at which the bias
potential goes to 0.0.  It is dimensionless, so a value of 0.3 means a
bond distance can be up to 30\% larger or 30\% smaller than the
equilibrium (quenched) \(R^0_{ij}\) distance and the two atoms in the bond
could still experience a non\sphinxhyphen{}zero bias force.

\sphinxAtStartPar
If \sphinxstyleemphasis{qfactor} is set too large, then transitions from one energy basin
to another are affected because the bias potential is non\sphinxhyphen{}zero at the
transition state (e.g. saddle point).  If \sphinxstyleemphasis{qfactor} is set too small
than little boost can be achieved because the \(E_{ij}\) strain of
some bond in
the system will (nearly) always exceed \sphinxstyleemphasis{qfactor}.  A value of 0.3 for
\sphinxstyleemphasis{qfactor} is typically a reasonable value.

\sphinxAtStartPar
The \sphinxstyleemphasis{Vmax} argument is a fixed prefactor on the bias potential.  There
is a also a dynamic prefactor \(C_{ij}\), driven by the choice of
\sphinxstyleemphasis{Btarget} as discussed above.  The product of these should be a value less than
the smallest barrier height for an event to occur.  Otherwise the
applied bias potential may be large enough (when added to the
interatomic potential) to produce a local energy basin with a maxima
in the center.  This can produce artificial energy minima in the same
basin that trap an atom.  Or if \(C_{ij} \cdot V^{max}\) is even
larger, it may
induce an atom(s) to rapidly transition to another energy basin.  Both
cases are “bad dynamics” which violate the assumptions of LHD that
guarantee an accelerated time\sphinxhyphen{}accurate trajectory of the system.

\begin{sphinxadmonition}{note}{Note:}
\sphinxAtStartPar
It may seem that \(V^{max}\) can be set to any value, and
\(C_{ij}\) will compensate to reduce the overall prefactor
if necessary.  However the \(C_{ij}\) are initialized to 1.0
and the boostostatting procedure typically operates slowly enough
that there can be a time period of bad dynamics if \(V^{max}\)
is set too large.  A better strategy is to set \(V^{max}\) to the
slightly smaller than the lowest barrier height for an event (the same
as for GHD), so that the \(C_{ij}\) remain near unity.
\end{sphinxadmonition}

\sphinxAtStartPar
The \sphinxstyleemphasis{Tequil} argument is the temperature at which the system is
simulated; see the comment above about the {\hyperref[\detokenize{fix_langevin::doc}]{\sphinxcrossref{\DUrole{doc}{fix langevin}}}} thermostatting.  It is also part of the beta term in the
exponential factor that determines how much boost is achieved as a
function of the bias potential.  See the discussion of the \sphinxstyleemphasis{Btarget}
argument below.

\sphinxAtStartPar
As discussed above, the \sphinxstyleemphasis{Dcut} argument is the distance required
between two locally maxstrain bonds for them to both be selected as
biased bonds on the same timestep.  Computationally, the larger \sphinxstyleemphasis{Dcut}
is, the more work (computation and communication) must be done each
timestep within the LHD algorithm.  And the fewer bonds can be
simultaneously biased, which may mean the specified \sphinxstyleemphasis{Btarget} time
acceleration cannot be achieved.

\sphinxAtStartPar
Physically \sphinxstyleemphasis{Dcut} should be a long enough distance that biasing two
pairs of atoms that close together will not influence the dynamics of
each pair.  E.g. something like 2x the cutoff of the interatomic
potential.  In practice a \sphinxstyleemphasis{Dcut} value of \textasciitilde{}10 Angstroms seems to work
well for many solid\sphinxhyphen{}state systems.

\begin{sphinxadmonition}{note}{Note:}
\sphinxAtStartPar
You should ensure that ghost atom communication is performed for
a distance of at least \sphinxstyleemphasis{Dcut} + \sphinxstyleemphasis{cutevent} = the distance one or more
atoms move (between quenched states) to be considered an “event”.  It
is an argument to the “compute event/displace” command used to detect
events.  By default the ghost communication distance is set by the
pair\_style cutoff, which will typically be \textless{} \sphinxstyleemphasis{Dcut}.  The {\hyperref[\detokenize{comm_modify::doc}]{\sphinxcrossref{\DUrole{doc}{comm\_modify cutoff}}}} command should be used to override the ghost
cutoff explicitly, e.g.
\end{sphinxadmonition}

\begin{sphinxVerbatim}[commandchars=\\\{\}]
\PYG{k}{comm\PYGZus{}modify}\PYG{+w}{ }\PYG{n}{cutoff}\PYG{+w}{ }\PYG{l+m}{12.0}
\end{sphinxVerbatim}

\sphinxAtStartPar
Note that this fix does not know the \sphinxstyleemphasis{cutevent} parameter, but uses
half the \sphinxstyleemphasis{cutbond} parameter as an estimate to warn if the ghost
cutoff is not long enough.

\sphinxAtStartPar
As described above the \sphinxstyleemphasis{alpha} argument is a prefactor in the
boostostat update equation for each bond’s \(C_{ij}\) prefactor.
\sphinxstyleemphasis{Alpha} is specified in time units, similar to other thermostat or barostat
damping parameters.  It is roughly the physical time it will take the
boostostat to adjust a \(C_{ij}\) value from a too high (or too low)
value to a correct one.  An \sphinxstyleemphasis{alpha} setting of a few ps is typically good for
solid\sphinxhyphen{}state systems.  Note that the \sphinxstyleemphasis{alpha} argument here is the
inverse of the alpha parameter discussed in
{\hyperref[\detokenize{fix_hyper_local:voter2013lhd}]{\sphinxcrossref{\DUrole{std,std-ref}{(Voter2013)}}}}.

\sphinxAtStartPar
The \sphinxstyleemphasis{Btarget} argument is the desired time boost factor (a value \textgreater{} 1)
that all the atoms in the system will experience.  The elapsed time
t\_hyper for an LHD simulation running for \sphinxstyleemphasis{N} timesteps is simply
\begin{equation*}
\begin{split}t_{hyper} = B_{target} \cdot N \cdot dt\end{split}
\end{equation*}
\sphinxAtStartPar
where \sphinxstyleemphasis{dt} is the timestep size defined by the {\hyperref[\detokenize{timestep::doc}]{\sphinxcrossref{\DUrole{doc}{timestep}}}}
command.  The effective time acceleration due to LHD is thus
\(\frac{t_{hyper}}{N\cdot dt} = B_{target}\), where \(N\cdot dt\)
is the elapsed time for a normal MD run of N timesteps.

\sphinxAtStartPar
You cannot choose an arbitrarily large setting for \sphinxstyleemphasis{Btarget}.  The
maximum value you should choose is
\begin{equation*}
\begin{split}B_{target} = e^{\beta V_{small}}\end{split}
\end{equation*}
\sphinxAtStartPar
where \(V_{small}\) is the smallest event barrier height in your
system, \(\beta = \frac{1}{kT_{equil}}\), and \(T_{equil}\)
is the specified temperature of the system
(both by this fix and the Langevin thermostat).

\sphinxAtStartPar
Note that if \sphinxstyleemphasis{Btarget} is set smaller than this, the LHD simulation
will run correctly.  There will just be fewer events because the hyper
time (t\_hyper equation above) will be shorter.

\begin{sphinxadmonition}{note}{Note:}
\sphinxAtStartPar
If you have no physical intuition as to the smallest barrier
height in your system, a reasonable strategy to determine the largest
\sphinxstyleemphasis{Btarget} you can use for an LHD model, is to run a sequence of
simulations with smaller and smaller \sphinxstyleemphasis{Btarget} values, until the event
rate does not change (as a function of hyper time).
\end{sphinxadmonition}


\bigskip\hrule\bigskip


\sphinxAtStartPar
Here is additional information on the optional keywords for this fix.

\sphinxAtStartPar
The \sphinxstyleemphasis{bound} keyword turns on min/max bounds for bias coefficients
\(C_{ij}\) for all bonds.  \(C_{ij}\) is a prefactor for each bond on
the bias potential of maximum strength \(V^{max}\).  Depending on the
choice of \sphinxstyleemphasis{alpha} and \sphinxstyleemphasis{Btarget} and \sphinxstyleemphasis{Vmax}, the boostostatting can cause
individual \(C_{ij}\) values to fluctuate.  If the fluctuations are too
large \(C_{ij} \cdot V^{max}\) can exceed low barrier heights and induce
bad event dynamics.  Bounding the \(C_{ij}\) values is a way to prevent
this.  If \sphinxstyleemphasis{Bfrac} is set to \sphinxhyphen{}1 or any negative value (the default) then no
bounds are enforced on \(C_{ij}\) values (except they must always
be \textgreater{}= 0.0).  A \sphinxstyleemphasis{Bfrac} setting \textgreater{}= 0.0
sets a lower bound of 1.0 \sphinxhyphen{} Bfrac and upper bound of 1.0 + Bfrac on each
\(C_{ij}\) value.  Note that all \(C_{ij}\) values are initialized
to 1.0 when a bond is created for the first time.  Thus \sphinxstyleemphasis{Bfrac} limits the
bias potential height to \sphinxstyleemphasis{Vmax} +/\sphinxhyphen{} \sphinxstyleemphasis{Bfrac}*\sphinxstyleemphasis{Vmax}.

\sphinxAtStartPar
The \sphinxstyleemphasis{reset} keyword allow \sphinxstyleemphasis{Vmax} to be adjusted dynamically depending on the
average value of all \(C_{ij}\) prefactors.  This can be useful if you
are unsure what value of \sphinxstyleemphasis{Vmax} will match the \sphinxstyleemphasis{Btarget} boost for the
system.  The \(C_{ij}\) values will then adjust in aggregate (up or down)
so that \(C_{ij} \cdot V^{max}\) produces a boost of \sphinxstyleemphasis{Btarget}, but this
may conflict with the \sphinxstyleemphasis{bound} keyword settings.  By using \sphinxstyleemphasis{bound} and \sphinxstyleemphasis{reset}
together, \(V^{max}\) itself can be reset, and desired bounds still applied
to the \(C_{ij}\) values.

\sphinxAtStartPar
A setting for \sphinxstyleemphasis{Rfreq} of \sphinxhyphen{}1 (the default) means \sphinxstyleemphasis{Vmax} never changes.
A setting of 0 means \(V^{max}\) is adjusted every time an event occurs and
bond pairs are recalculated.  A setting of N \textgreater{} 0 timesteps means
\(V^{max}\) is adjusted on the first time an event occurs on a timestep \textgreater{}=
N steps after the previous adjustment.  The adjustment to \(V^{max}\) is
computed as follows.  The current average of all \(C_{ij} \cdot V^{max}\)
values is computed and the \(V^{max}\) is reset to that value.  All
\(C_{ij}\) values are changed to new prefactors such the new
\(C_{ij} \cdot V^{max}\) is the same as it was previously.  If the
\sphinxstyleemphasis{bound} keyword was used, those bounds are enforced on the new \(C_{ij}\)
values.  Henceforth, new bonds are assigned a \(C_{ij} = 1.0\), which
means their bias potential magnitude is the new \(V^{max}\).

\sphinxAtStartPar
The \sphinxstyleemphasis{check/ghost} keyword turns on extra computation each timestep to
compute statistics about ghost atoms used to determine which bonds to
bias.  The output of these stats are the vector values 14 and 15,
described below.  If this keyword is not enabled, the output
of the stats will be zero.

\sphinxAtStartPar
The \sphinxstyleemphasis{check/bias} keyword turns on extra computation and communication
to check if any biased bonds are closer than \sphinxstyleemphasis{Dcut} to each other,
which should not be the case if LHD is operating correctly.  Thus it
is a debugging check.  The \sphinxstyleemphasis{Nevery} setting determines how often the
check is made.  The \sphinxstyleemphasis{error}, \sphinxstyleemphasis{warn}, or \sphinxstyleemphasis{ignore} setting determines
what is done if the count of too\sphinxhyphen{}close bonds is not zero.  Either the
code will exit, or issue a warning, or silently tally the count.  The
count can be output as vector value 17, as described below.  If this
keyword is not enabled, the output of that statistic will be 0.

\sphinxAtStartPar
Note that both of these computations are costly, hence they are only
enabled by these keywords.


\bigskip\hrule\bigskip



\subsection{Restart, fix\_modify, output, run start/stop, minimize info}
\label{\detokenize{fix_hyper_local:restart-fix-modify-output-run-start-stop-minimize-info}}
\sphinxAtStartPar
No information about this fix is written to {\hyperref[\detokenize{restart::doc}]{\sphinxcrossref{\DUrole{doc}{binary restart files}}}}.

\sphinxAtStartPar
The {\hyperref[\detokenize{fix_modify::doc}]{\sphinxcrossref{\DUrole{doc}{fix\_modify}}}} \sphinxstyleemphasis{energy} option is supported by
this fix to add the energy of the bias potential to the global
potential energy of the system as part of {\hyperref[\detokenize{thermo_style::doc}]{\sphinxcrossref{\DUrole{doc}{thermodynamic output}}}}.  The default setting for this fix is {\hyperref[\detokenize{fix_modify::doc}]{\sphinxcrossref{\DUrole{doc}{fix\_modify
energy no}}}}.

\sphinxAtStartPar
This fix computes a global scalar and global vector of length 28,
which can be accessed by various {\hyperref[\detokenize{Howto_output::doc}]{\sphinxcrossref{\DUrole{doc}{output commands}}}}.  The scalar is the magnitude of the bias potential
(energy units) applied on the current timestep, summed over all biased
bonds.  The vector stores the following quantities:
\begin{enumerate}
\sphinxsetlistlabels{\arabic}{enumi}{enumii}{}{.}%
\item {} 
\sphinxAtStartPar
average boost for all bonds on this step (unitless)

\item {} 
\sphinxAtStartPar
\# of biased bonds on this step

\item {} 
\sphinxAtStartPar
max strain \(E_{ij}\) of any bond on this step (absolute value, unitless)

\item {} 
\sphinxAtStartPar
value of \(V^{max}\) on this step (energy units)

\item {} 
\sphinxAtStartPar
average bias coeff for all bonds on this step (unitless)

\item {} 
\sphinxAtStartPar
min bias coeff for all bonds on this step (unitless)

\item {} 
\sphinxAtStartPar
max bias coeff for all bonds on this step (unitless)

\item {} 
\sphinxAtStartPar
average \# of bonds/atom on this step

\item {} 
\sphinxAtStartPar
average neighbor bonds/bond on this step within \sphinxstyleemphasis{Dcut}

\item {} 
\sphinxAtStartPar
average boost for all bonds during this run (unitless)

\item {} 
\sphinxAtStartPar
average \# of biased bonds/step during this run

\item {} 
\sphinxAtStartPar
fraction of biased bonds with no bias during this run

\item {} 
\sphinxAtStartPar
fraction of biased bonds with negative strain during this run

\item {} 
\sphinxAtStartPar
max bond length during this run (distance units)

\item {} 
\sphinxAtStartPar
average bias coeff for all bonds during this run (unitless)

\item {} 
\sphinxAtStartPar
min bias coeff for any bond during this run (unitless)

\item {} 
\sphinxAtStartPar
max bias coeff for any bond during this run (unitless)

\item {} 
\sphinxAtStartPar
max drift distance of any bond atom during this run (distance units)

\item {} 
\sphinxAtStartPar
max distance from proc subbox of any ghost atom with maxstrain \textless{} qfactor during this run (distance units)

\item {} 
\sphinxAtStartPar
max distance outside my box of any ghost atom with any maxstrain during this run (distance units)

\item {} 
\sphinxAtStartPar
count of ghost atoms that could not be found on reneighbor steps during this run

\item {} 
\sphinxAtStartPar
count of bias overlaps (\textless{} \sphinxstyleemphasis{Dcut}) found during this run

\item {} 
\sphinxAtStartPar
cumulative hyper time since fix created (time units)

\item {} 
\sphinxAtStartPar
cumulative count of event timesteps since fix created

\item {} 
\sphinxAtStartPar
cumulative count of atoms in events since fix created

\item {} 
\sphinxAtStartPar
cumulative \# of new bonds formed since fix created

\item {} 
\sphinxAtStartPar
average boost for biased bonds on this step (unitless)

\item {} 
\sphinxAtStartPar
\# of bonds with absolute strain \textgreater{}= q on this step

\end{enumerate}

\sphinxAtStartPar
Quantities 1\sphinxhyphen{}9 are for the current timestep.  Quantities 10\sphinxhyphen{}22
are for the current hyper run.  They are reset each time a new
hyper run is performed.  Quantities 23\sphinxhyphen{}26 are cumulative across
multiple runs (since the point in the input script the fix was
defined).

\sphinxAtStartPar
For value 10, each bond instantaneous boost factor is given by the
equation for \(B_{ij}\) above.  The total system boost (average across all
bonds) fluctuates, but should average to a value close to the
specified \(B_{target}\).

\sphinxAtStartPar
For value 12, the numerator is a count of all biased bonds on each
timestep whose bias energy = 0.0 due to \(E_{ij} >= qfactor\).  The
denominator is the count of all biased bonds on all timesteps.

\sphinxAtStartPar
For value 13, the numerator is a count of all biased bonds on each
timestep with negative strain.  The denominator is the count of all
biased bonds on all timesteps.

\sphinxAtStartPar
Values 18\sphinxhyphen{}22 are mostly useful for debugging and diagnostic purposes.

\sphinxAtStartPar
For value 18, drift is the distance an atom moves between two quenched
states when the second quench determines an event has occurred.  Atoms
involved in an event will typically move the greatest distance since
others typically remain near their original quenched position.

\sphinxAtStartPar
For values 19\sphinxhyphen{}21, neighbor atoms in the full neighbor list with cutoff
\sphinxstyleemphasis{Dcut} may be ghost atoms outside a processor’s sub\sphinxhyphen{}box.  Before the
next event occurs they may move further than \sphinxstyleemphasis{Dcut} away from the
sub\sphinxhyphen{}box boundary.  Value 19 is the furthest (from the sub\sphinxhyphen{}box) any
ghost atom in the neighbor list with maxstrain \textless{} \sphinxstyleemphasis{qfactor} was
accessed during the run.  Value 20 is the same except that the ghost
atom’s maxstrain may be \textgreater{}= \sphinxstyleemphasis{qfactor}, which may mean it is about to
participate in an event.  Value 21 is a count of how many ghost atoms
could not be found on reneighbor steps, presumably because they moved
too far away due to their participation in an event (which will likely
be detected at the next quench).

\sphinxAtStartPar
Typical values for 19 and 20 should be slightly larger than \sphinxstyleemphasis{Dcut},
which accounts for ghost atoms initially at a \sphinxstyleemphasis{Dcut} distance moving
thermally before the next event takes place.

\sphinxAtStartPar
Note that for values 19 and 20 to be computed, the optional keyword
\sphinxstyleemphasis{check/ghost} must be specified.  Otherwise these values will be zero.
This is because computing them incurs overhead, so the values are only
computed if requested.

\sphinxAtStartPar
Value 21 should be zero or small.  As explained above a small count
likely means some ghost atoms were participating in their own events
and moved a longer distance.  If the value is large, it likely means
the communication cutoff for ghosts is too close to \sphinxstyleemphasis{Dcut} leading to
many not\sphinxhyphen{}found ghost atoms before the next event.  This may lead to a
reduced number of bonds being selected for biasing, since the code
assumes those atoms are part of highly strained bonds.  As explained
above, the {\hyperref[\detokenize{comm_modify::doc}]{\sphinxcrossref{\DUrole{doc}{comm\_modify cutoff}}}} command can be used
to set a longer cutoff.

\sphinxAtStartPar
For value 22, no two bonds should be biased if they are within a
\sphinxstyleemphasis{Dcut} distance of each other.  This value should be zero, indicating
that no pair of biased bonds are closer than \sphinxstyleemphasis{Dcut} from each other.

\sphinxAtStartPar
Note that for value 22 to be computed, the optional keyword
\sphinxstyleemphasis{check/bias} must be specified and it determines how often this check
is performed.  This is because performing the check incurs overhead,
so if only computed as often as requested.

\sphinxAtStartPar
The result at the end of the run is the cumulative total from every
timestep the check was made.  Note that the value is a count of atoms
in bonds which found other atoms in bonds too close, so it is almost
always an over\sphinxhyphen{}count of the number of too\sphinxhyphen{}close bonds.

\sphinxAtStartPar
Value 23 is simply the specified \sphinxstyleemphasis{boost} factor times the number of
timesteps times the timestep size.

\sphinxAtStartPar
For value 24, events are checked for by the {\hyperref[\detokenize{hyper::doc}]{\sphinxcrossref{\DUrole{doc}{hyper}}}} command
once every \sphinxstyleemphasis{Nevent} timesteps.  This value is the count of those
timesteps on which one (or more) events was detected.  It is NOT the
number of distinct events, since more than one event may occur in the
same \sphinxstyleemphasis{Nevent} time window.

\sphinxAtStartPar
For value 25, each time the {\hyperref[\detokenize{hyper::doc}]{\sphinxcrossref{\DUrole{doc}{hyper}}}} command checks for an
event, it invokes a compute to flag zero or more atoms as
participating in one or more events.  E.g. atoms that have displaced
more than some distance from the previous quench state.  Value 25 is
the cumulative count of the number of atoms participating in any of
the events that were found.

\sphinxAtStartPar
Value 26 tallies the number of new bonds created by the bond reset
operation.  Bonds between a specific I,J pair of atoms may persist for
the entire hyperdynamics simulation if neither I or J are involved in
an event.

\sphinxAtStartPar
Value 27 computes the average boost for biased bonds only on this step.

\sphinxAtStartPar
Value 28 is the count of bonds with an absolute value of strain \textgreater{}= q
on this step.

\sphinxAtStartPar
The scalar value and vector values are all “intensive”.

\sphinxAtStartPar
This fix also computes a local vector of length the number of bonds
currently in the system.  The value for each bond is its \(C_{ij}\)
prefactor (bias coefficient).  These values can be can be accessed by various
{\hyperref[\detokenize{Howto_output::doc}]{\sphinxcrossref{\DUrole{doc}{output commands}}}}.  A particularly useful one is the
{\hyperref[\detokenize{fix_ave_histo::doc}]{\sphinxcrossref{\DUrole{doc}{fix ave/histo}}}} command which can be used to
histogram the Cij values to see if they are distributed reasonably
close to 1.0, which indicates a good choice of \(V^{max}\).

\sphinxAtStartPar
The local values calculated by this fix are unitless.

\sphinxAtStartPar
No parameter of this fix can be used with the \sphinxstyleemphasis{start/stop} keywords of
the {\hyperref[\detokenize{run::doc}]{\sphinxcrossref{\DUrole{doc}{run}}}} command.  This fix is not invoked during
{\hyperref[\detokenize{minimize::doc}]{\sphinxcrossref{\DUrole{doc}{energy minimization}}}}.


\subsection{Restrictions}
\label{\detokenize{fix_hyper_local:restrictions}}
\sphinxAtStartPar
This fix is part of the REPLICA package.  It is only enabled if LAMMPS
was built with that package.  See the {\hyperref[\detokenize{Build_package::doc}]{\sphinxcrossref{\DUrole{doc}{Build package}}}}
doc page for more info.


\subsection{Related commands}
\label{\detokenize{fix_hyper_local:related-commands}}
\sphinxAtStartPar
{\hyperref[\detokenize{hyper::doc}]{\sphinxcrossref{\DUrole{doc}{hyper}}}}, {\hyperref[\detokenize{fix_hyper_global::doc}]{\sphinxcrossref{\DUrole{doc}{fix hyper/global}}}}


\subsection{Default}
\label{\detokenize{fix_hyper_local:default}}
\sphinxAtStartPar
The default settings for optimal keywords are bounds = \sphinxhyphen{}1 and reset =
\sphinxhyphen{}1.  The check/ghost and check/bias keywords are not enabled by
default.


\bigskip\hrule\bigskip

\phantomsection\label{\detokenize{fix_hyper_local:voter2013lhd}}
\sphinxAtStartPar
\sphinxstylestrong{(Voter2013)} S. Y. Kim, D. Perez, A. F. Voter, J Chem Phys, 139,
144110 (2013).

\phantomsection\label{\detokenize{fix_hyper_local:mironlhd}}
\sphinxAtStartPar
\sphinxstylestrong{(Miron)} R. A. Miron and K. A. Fichthorn, J Chem Phys, 119, 6210 (2003).

\sphinxstepscope

\index{fix imd@\spxentry{fix imd}}\ignorespaces 

\section{fix imd command}
\label{\detokenize{fix_imd:fix-imd-command}}\label{\detokenize{fix_imd:index-0}}\label{\detokenize{fix_imd::doc}}

\subsection{Syntax}
\label{\detokenize{fix_imd:syntax}}
\begin{sphinxVerbatim}[commandchars=\\\{\}]
\PYG{k}{fix }\PYG{n+nv+nvIdentifier}{ID}\PYG{+w}{ }\PYG{n+nv+nvIdentifier}{group\PYGZhy{}ID}\PYG{+w}{ }\PYG{n}{imd}\PYG{+w}{ }\PYG{n}{trate}\PYG{+w}{ }\PYG{n}{port}\PYG{+w}{ }\PYG{n}{keyword}\PYG{+w}{ }\PYG{n}{values}\PYG{+w}{ }\PYG{n}{...}
\end{sphinxVerbatim}
\begin{itemize}
\item {} 
\sphinxAtStartPar
ID, group\sphinxhyphen{}ID are documented in {\hyperref[\detokenize{fix::doc}]{\sphinxcrossref{\DUrole{doc}{fix}}}} command

\item {} 
\sphinxAtStartPar
imd = style name of this fix command

\item {} 
\sphinxAtStartPar
port = port number on which the fix listens for an IMD client

\item {} 
\sphinxAtStartPar
keyword = \sphinxstyleemphasis{unwrap} or \sphinxstyleemphasis{fscale} or \sphinxstyleemphasis{trate} or \sphinxstyleemphasis{nowait}
\begin{sphinxalltt}
\sphinxstyleemphasis{unwrap} arg = \sphinxstyleemphasis{on} or \sphinxstyleemphasis{off}
  off = coordinates are wrapped back into the principal unit cell (default)
  on = "unwrapped" coordinates using the image flags used
\sphinxstyleemphasis{fscale} arg = factor
  factor = floating point number to scale IMD forces (default: 1.0)
\sphinxstyleemphasis{trate} arg = transmission rate of coordinate data sets (default: 1)
\sphinxstyleemphasis{nowait} arg = \sphinxstyleemphasis{on} or \sphinxstyleemphasis{off}
  off = LAMMPS waits to be connected to an IMD client before continuing (default)
  on = LAMMPS listens for an IMD client, but continues with the run
\end{sphinxalltt}

\end{itemize}


\subsection{Examples}
\label{\detokenize{fix_imd:examples}}
\begin{sphinxVerbatim}[commandchars=\\\{\}]
\PYG{k}{fix }\PYG{n+nv+nvIdentifier}{vmd}\PYG{+w}{ }\PYG{n+nv+nvIdentifier}{all}\PYG{+w}{ }\PYG{n}{imd}\PYG{+w}{ }\PYG{l+m}{5678}
\PYG{k}{fix }\PYG{n+nv+nvIdentifier}{comm}\PYG{+w}{ }\PYG{n+nv+nvIdentifier}{all}\PYG{+w}{ }\PYG{n}{imd}\PYG{+w}{ }\PYG{l+m}{8888}\PYG{+w}{ }\PYG{n}{trate}\PYG{+w}{ }\PYG{l+m}{5}\PYG{+w}{ }\PYG{n}{unwrap}\PYG{+w}{ }\PYG{n}{on}\PYG{+w}{ }\PYG{n}{fscale}\PYG{+w}{ }\PYG{l+m}{10.0}
\end{sphinxVerbatim}


\subsection{Description}
\label{\detokenize{fix_imd:description}}
\sphinxAtStartPar
This fix implements the “Interactive MD” (IMD) protocol which allows
realtime visualization and manipulation of MD simulations through the
IMD protocol, as initially implemented in VMD and NAMD.  Specifically
it allows LAMMPS to connect an IMD client, for example the \sphinxhref{https://www.ks.uiuc.edu/Research/vmd}{VMD visualization program}, so that it can monitor the progress of the
simulation and interactively apply forces to selected atoms.

\sphinxAtStartPar
If LAMMPS is compiled with the pre\sphinxhyphen{}processor flag \sphinxhyphen{}DLAMMPS\_ASYNC\_IMD
then fix imd will use POSIX threads to spawn a IMD communication
thread on MPI rank 0 in order to offload data reading and writing
from the main execution thread and potentially lower the inferred
latencies for slow communication links. This feature has only been
tested under linux.

\sphinxAtStartPar
The source code for this fix includes code developed by the Theoretical
and Computational Biophysics Group in the Beckman Institute for Advanced
Science and Technology at the University of Illinois at
Urbana\sphinxhyphen{}Champaign.  We thank them for providing a software interface that
allows codes like LAMMPS to hook to \sphinxhref{https://www.ks.uiuc.edu/Research/vmd}{VMD}.

\sphinxAtStartPar
Upon initialization of the fix, it will open a communication port on
the node with MPI task 0 and wait for an incoming connection.  As soon
as an IMD client is connected, the simulation will continue and the
fix will send the current coordinates of the fix’s group to the IMD
client at every trate MD step. When using r\sphinxhyphen{}RESPA, trate applies to
the steps of the outmost RESPA level.  During a run with an active IMD
connection also the IMD client can request to apply forces to selected
atoms of the fix group.

\sphinxAtStartPar
The port number selected must be an available network port number.  On
many machines, port numbers \textless{} 1024 are reserved for accounts with
system manager privilege and specific applications. If multiple imd
fixes would be active at the same time, each needs to use a different
port number.

\sphinxAtStartPar
The \sphinxstyleemphasis{nowait} keyword controls the behavior of the fix when no IMD
client is connected. With the default setting of \sphinxstyleemphasis{off}, LAMMPS will
wait until a connection is made before continuing with the
execution. Setting \sphinxstyleemphasis{nowait} to \sphinxstyleemphasis{on} will have the LAMMPS code be ready
to connect to a client, but continue with the simulation. This can for
example be used to monitor the progress of an ongoing calculation
without the need to be permanently connected or having to download a
trajectory file.

\sphinxAtStartPar
The \sphinxstyleemphasis{trate} keyword allows to select how often the coordinate data is
sent to the IMD client. It can also be changed on request of the IMD
client through an IMD protocol message.  The \sphinxstyleemphasis{unwrap} keyword allows
to send “unwrapped” coordinates to the IMD client that undo the
wrapping back of coordinates into the principle unit cell, as done by
default in LAMMPS.  The \sphinxstyleemphasis{fscale} keyword allows to apply a scaling
factor to forces transmitted by the IMD client. The IMD protocols
stipulates that forces are transferred in kcal/mol/Angstrom under the
assumption that coordinates are given in Angstrom. For LAMMPS runs
with different units or as a measure to tweak the forces generated by
the manipulation of the IMD client, this option allows to make
adjustments.

\sphinxAtStartPar
To connect VMD to a listening LAMMPS simulation on the same machine
with fix imd enabled, one needs to start VMD and load a coordinate or
topology file that matches the fix group.  When the VMD command
prompts appears, one types the command line:

\begin{sphinxVerbatim}[commandchars=\\\{\}]
imd connect localhost 5678
\end{sphinxVerbatim}

\sphinxAtStartPar
This assumes that \sphinxstyleemphasis{fix imd} was started with 5678 as a port
number for the IMD protocol.

\sphinxAtStartPar
The steps to do interactive manipulation of a running simulation in
VMD are the following:

\sphinxAtStartPar
In the Mouse menu of the VMD Main window, select “Mouse \sphinxhyphen{}\textgreater{} Force \sphinxhyphen{}\textgreater{}
Atom”.  You may alternately select “Residue”, or “Fragment” to apply
forces to whole residues or fragments. Your mouse can now be used to
apply forces to your simulation. Click on an atom, residue, or fragment
and drag to apply a force. Click quickly without moving the mouse to
turn the force off. You can also use a variety of 3D position trackers
to apply forces to your simulation. Game controllers or haptic devices
with force\sphinxhyphen{}feedback such as the Novint Falcon or Sensable PHANTOM allow
you to feel the resistance due to inertia or interactions with neighbors
that the atoms experience you are trying to move, as if they were real
objects. See the \sphinxhref{https://www.ks.uiuc.edu/Research/vmd/imd/}{VMD IMD Homepage} for more details.

\sphinxAtStartPar
If IMD control messages are received, a line of text describing the
message and its effect will be printed to the LAMMPS output screen, if
screen output is active.


\subsection{Restart, fix\_modify, output, run start/stop, minimize info}
\label{\detokenize{fix_imd:restart-fix-modify-output-run-start-stop-minimize-info}}
\sphinxAtStartPar
No information about this fix is written to {\hyperref[\detokenize{restart::doc}]{\sphinxcrossref{\DUrole{doc}{binary restart files}}}}.  None of the {\hyperref[\detokenize{fix_modify::doc}]{\sphinxcrossref{\DUrole{doc}{fix\_modify}}}} options are
relevant to this fix.  No global scalar or vector or per\sphinxhyphen{}atom quantities
are stored by this fix for access by various {\hyperref[\detokenize{Howto_output::doc}]{\sphinxcrossref{\DUrole{doc}{output commands}}}}.  No parameter of this fix can be used with the
\sphinxstyleemphasis{start/stop} keywords of the {\hyperref[\detokenize{run::doc}]{\sphinxcrossref{\DUrole{doc}{run}}}} command.  This fix is not
invoked during {\hyperref[\detokenize{minimize::doc}]{\sphinxcrossref{\DUrole{doc}{energy minimization}}}}.


\subsection{Restrictions}
\label{\detokenize{fix_imd:restrictions}}
\sphinxAtStartPar
This fix is part of the MISC package.  It is only enabled if LAMMPS was
built with that package.  See the {\hyperref[\detokenize{Build_package::doc}]{\sphinxcrossref{\DUrole{doc}{Build package}}}}
page for more info.

\sphinxAtStartPar
When used in combination with VMD, a topology or coordinate file has
to be loaded, which matches (in number and ordering of atoms) the
group the fix is applied to. The fix internally sorts atom IDs by
ascending integer value; in VMD (and thus the IMD protocol) those will
be assigned 0\sphinxhyphen{}based consecutive index numbers.

\sphinxAtStartPar
When using multiple active IMD connections at the same time, each
needs to use a different port number.


\subsection{Related commands}
\label{\detokenize{fix_imd:related-commands}}
\sphinxAtStartPar
none


\subsection{Default}
\label{\detokenize{fix_imd:default}}
\sphinxAtStartPar
none

\sphinxstepscope

\index{fix indent@\spxentry{fix indent}}\ignorespaces 

\section{fix indent command}
\label{\detokenize{fix_indent:fix-indent-command}}\label{\detokenize{fix_indent:index-0}}\label{\detokenize{fix_indent::doc}}

\subsection{Syntax}
\label{\detokenize{fix_indent:syntax}}
\begin{sphinxVerbatim}[commandchars=\\\{\}]
\PYG{k}{fix }\PYG{n+nv+nvIdentifier}{ID}\PYG{+w}{ }\PYG{n+nv+nvIdentifier}{group\PYGZhy{}ID}\PYG{+w}{ }\PYG{n}{indent}\PYG{+w}{ }\PYG{n}{K}\PYG{+w}{ }\PYG{n}{gstyle}\PYG{+w}{ }\PYG{n}{args}\PYG{+w}{ }\PYG{n}{keyword}\PYG{+w}{ }\PYG{n}{value}\PYG{+w}{ }\PYG{n}{...}
\end{sphinxVerbatim}
\begin{itemize}
\item {} 
\sphinxAtStartPar
ID, group\sphinxhyphen{}ID are documented in {\hyperref[\detokenize{fix::doc}]{\sphinxcrossref{\DUrole{doc}{fix}}}} command

\item {} 
\sphinxAtStartPar
indent = style name of this fix command

\item {} 
\sphinxAtStartPar
K = force constant for indenter surface (force/distance\textasciicircum{}2 units)

\item {} 
\sphinxAtStartPar
gstyle = \sphinxstyleemphasis{sphere} or \sphinxstyleemphasis{cylinder} or \sphinxstyleemphasis{cone} or \sphinxstyleemphasis{plane}
\begin{sphinxalltt}
\sphinxstyleemphasis{sphere} args = x y z R
  x, y, z = position of center of indenter (distance units)
  R = sphere radius of indenter (distance units)
  any of x, y, z, R can be a variable (see below)
\sphinxstyleemphasis{cylinder} args = dim c1 c2 R
  dim = \sphinxstyleemphasis{x} or \sphinxstyleemphasis{y} or \sphinxstyleemphasis{z} = axis of cylinder
  c1, c2 = coords of cylinder axis in other 2 dimensions (distance units)
  R = cylinder radius of indenter (distance units)
  any of c1,c2,R can be a variable (see below)
\sphinxstyleemphasis{cone} args = dim c1 c2 radlo radhi lo hi
  dim = \sphinxstyleemphasis{x} or \sphinxstyleemphasis{y} or \sphinxstyleemphasis{z} = axis of cone
  c1, c2 = coords of cone axis in other 2 dimensions (distance units)
  radlo,radhi = cone radii at lo and hi end (distance units)
  lo,hi = bounds of cone in dim (distance units)
  any of c1, c2, radlo, radhi, lo, hi can be a variable (see below)
\sphinxstyleemphasis{plane} args = dim pos side
  dim = \sphinxstyleemphasis{x} or \sphinxstyleemphasis{y} or \sphinxstyleemphasis{z} = plane perpendicular to this dimension
  pos = position of plane in dimension x, y, or z (distance units)
  pos can be a variable (see below)
  side = \sphinxstyleemphasis{lo} or \sphinxstyleemphasis{hi}
\end{sphinxalltt}

\item {} 
\sphinxAtStartPar
zero or more keyword/value pairs may be appended

\item {} 
\sphinxAtStartPar
keyword = \sphinxstyleemphasis{side} or \sphinxstyleemphasis{units}
\begin{sphinxalltt}
\sphinxstyleemphasis{side} value = \sphinxstyleemphasis{in} or \sphinxstyleemphasis{out}
  \sphinxstyleemphasis{in} = the indenter acts on particles inside the sphere or cylinder or cone
  \sphinxstyleemphasis{out} = the indenter acts on particles outside the sphere or cylinder or cone
\sphinxstyleemphasis{units} value = \sphinxstyleemphasis{lattice} or \sphinxstyleemphasis{box}
  lattice = the geometry is defined in lattice units
  box = the geometry is defined in simulation box units
\end{sphinxalltt}

\end{itemize}


\subsection{Examples}
\label{\detokenize{fix_indent:examples}}
\begin{sphinxVerbatim}[commandchars=\\\{\}]
\PYG{k}{fix }\PYG{n+nv+nvIdentifier}{1}\PYG{+w}{ }\PYG{n+nv+nvIdentifier}{all}\PYG{+w}{ }\PYG{n}{indent}\PYG{+w}{ }\PYG{l+m}{10.0}\PYG{+w}{ }\PYG{n}{sphere}\PYG{+w}{ }\PYG{l+m}{0.0}\PYG{+w}{ }\PYG{l+m}{0.0}\PYG{+w}{ }\PYG{l+m}{15.0}\PYG{+w}{ }\PYG{l+m}{3.0}
\PYG{k}{fix }\PYG{n+nv+nvIdentifier}{1}\PYG{+w}{ }\PYG{n+nv+nvIdentifier}{all}\PYG{+w}{ }\PYG{n}{indent}\PYG{+w}{ }\PYG{l+m}{10.0}\PYG{+w}{ }\PYG{n}{sphere}\PYG{+w}{ }\PYG{n}{v\PYGZus{}x}\PYG{+w}{ }\PYG{n}{v\PYGZus{}y}\PYG{+w}{ }\PYG{l+m}{0.0}\PYG{+w}{ }\PYG{n}{v\PYGZus{}radius}\PYG{+w}{ }\PYG{n}{side}\PYG{+w}{ }\PYG{n}{in}
\PYG{k}{fix }\PYG{n+nv+nvIdentifier}{2}\PYG{+w}{ }\PYG{n+nv+nvIdentifier}{flow}\PYG{+w}{ }\PYG{n}{indent}\PYG{+w}{ }\PYG{l+m}{10.0}\PYG{+w}{ }\PYG{n}{cylinder}\PYG{+w}{ }\PYG{n}{z}\PYG{+w}{ }\PYG{l+m}{0.0}\PYG{+w}{ }\PYG{l+m}{0.0}\PYG{+w}{ }\PYG{l+m}{10.0}\PYG{+w}{ }\PYG{n}{units}\PYG{+w}{ }\PYG{n}{box}
\end{sphinxVerbatim}


\subsection{Description}
\label{\detokenize{fix_indent:description}}
\sphinxAtStartPar
Insert an indenter within a simulation box.  The indenter repels all
atoms in the group that touch it, so it can be used to push into a
material or as an obstacle in a flow.  Alternatively, it can be used as a
constraining wall around a simulation; see the discussion of the
\sphinxstyleemphasis{side} keyword below.

\sphinxAtStartPar
The \sphinxstyleemphasis{gstyle} keyword selects the geometry of the indenter and it can
either have the value of \sphinxstyleemphasis{sphere}, \sphinxstyleemphasis{cylinder}, \sphinxstyleemphasis{cone}, or \sphinxstyleemphasis{plane}.

\sphinxAtStartPar
A spherical indenter (\sphinxstyleemphasis{gstyle} = \sphinxstyleemphasis{sphere}) exerts a force of magnitude
\begin{equation*}
\begin{split}F(r) = - K \left( r - R \right)^2\end{split}
\end{equation*}
\sphinxAtStartPar
on each atom where \sphinxstyleemphasis{K} is the specified force constant, \sphinxstyleemphasis{r} is the
distance from the atom to the center of the indenter, and \sphinxstyleemphasis{R} is the
radius of the indenter.  The force is repulsive and F(r) = 0 for \sphinxstyleemphasis{r} \textgreater{}
\sphinxstyleemphasis{R}.

\sphinxAtStartPar
A cylindrical indenter (\sphinxstyleemphasis{gstyle} = \sphinxstyleemphasis{cylinder}) follows the same formula
for the force as a sphere, except that \sphinxstyleemphasis{r} is defined the distance
from the atom to the center axis of the cylinder.  The cylinder extends
infinitely along its axis.

\sphinxAtStartPar
\DUrole{versionmodified,added}{Added in version 17April2024.}

\sphinxAtStartPar
A conical indenter (\sphinxstyleemphasis{gstyle} = \sphinxstyleemphasis{cone}) is similar to a cylindrical indenter
except that it has a finite length (between \sphinxstyleemphasis{lo} and \sphinxstyleemphasis{hi}), and that two
different radii (one at each end, \sphinxstyleemphasis{radlo} and \sphinxstyleemphasis{radhi}) can be defined.

\sphinxAtStartPar
Spherical, cylindrical, and conical indenters account for periodic
boundaries in two ways.  First, the center point of a spherical
indenter (x,y,z) or axis of a cylindrical/conical indenter (c1,c2) is
remapped back into the simulation box, if the box is periodic in a
particular dimension.  This occurs every timestep if the indenter
geometry is specified with a variable (see below), e.g. it is moving
over time.  Second, the calculation of distance to the indenter center
or axis accounts for periodic boundaries.  Both of these mean that an
indenter can effectively move through and straddle one or more
periodic boundaries.

\sphinxAtStartPar
A planar indenter (\sphinxstyleemphasis{gstyle} = \sphinxstyleemphasis{plane}) behaves like an axis\sphinxhyphen{}aligned
infinite\sphinxhyphen{}extent wall with the same force expression on atoms in the
system as before, but where \sphinxstyleemphasis{R} is the position of the plane and \sphinxstyleemphasis{r\sphinxhyphen{}R}
is the distance of an from the plane.  If the \sphinxstyleemphasis{side} parameter of the
plane is specified as \sphinxstyleemphasis{lo} then it will indent from the lo end of the
simulation box, meaning that atoms with a coordinate less than the
plane’s current position will be pushed towards the hi end of the box
and atoms with a coordinate higher than the plane’s current position
will feel no force.  Vice versa if \sphinxstyleemphasis{side} is specified as \sphinxstyleemphasis{hi}.

\sphinxAtStartPar
Any of the 4 quantities defining a spherical indenter’s geometry can
be specified as an equal\sphinxhyphen{}style {\hyperref[\detokenize{variable::doc}]{\sphinxcrossref{\DUrole{doc}{variable}}}}, namely \sphinxstyleemphasis{x},
\sphinxstyleemphasis{y}, \sphinxstyleemphasis{z}, or \sphinxstyleemphasis{R}.  For a cylindrical indenter, any of the 3
quantities \sphinxstyleemphasis{c1}, \sphinxstyleemphasis{c2}, or \sphinxstyleemphasis{R}, can be a variable.  For a conical
indenter, any of the 6 quantities \sphinxstyleemphasis{c1}, \sphinxstyleemphasis{c2}, \sphinxstyleemphasis{radlo}, \sphinxstyleemphasis{radhi}, \sphinxstyleemphasis{lo},
or \sphinxstyleemphasis{hi} can be a variable.  For a planar indenter, the single value
\sphinxstyleemphasis{pos} can be a variable.

\sphinxAtStartPar
If any of these values is a variable, it should be specified as
v\_name, where name is the variable name.  In this case, the variable
will be evaluated each timestep, and its value used to define the
indenter geometry.

\sphinxAtStartPar
Note that equal\sphinxhyphen{}style variables can specify formulas with various
mathematical functions, and include {\hyperref[\detokenize{thermo_style::doc}]{\sphinxcrossref{\DUrole{doc}{thermo\_style}}}}
command keywords for the simulation box parameters and timestep and
elapsed time.  Thus it is easy to specify indenter properties that
change as a function of time or span consecutive runs in a continuous
fashion.  For the latter, see the \sphinxstyleemphasis{start} and \sphinxstyleemphasis{stop} keywords of the
{\hyperref[\detokenize{run::doc}]{\sphinxcrossref{\DUrole{doc}{run}}}} command and the \sphinxstyleemphasis{elaplong} keyword of
{\hyperref[\detokenize{thermo_style::doc}]{\sphinxcrossref{\DUrole{doc}{thermo\_style custom}}}} for details.

\sphinxAtStartPar
For example, if a spherical indenter’s x\sphinxhyphen{}position is specified as v\_x,
then this variable definition will keep it’s center at a relative
position in the simulation box, 1/4 of the way from the left edge to
the right edge, even if the box size changes:

\begin{sphinxVerbatim}[commandchars=\\\{\}]
\PYG{k}{variable }\PYG{n+nv+nvIdentifier}{x}\PYG{+w}{ }\PYG{n}{equal}\PYG{+w}{ }\PYG{l+s}{\PYGZdq{}}\PYG{l+s}{xlo + 0.25*lx}\PYG{l+s}{\PYGZdq{}}
\end{sphinxVerbatim}

\sphinxAtStartPar
Similarly, either of these variable definitions will move the indenter
from an initial position at 2.5 at a constant velocity of 5:

\begin{sphinxVerbatim}[commandchars=\\\{\}]
\PYG{k}{variable }\PYG{n+nv+nvIdentifier}{x}\PYG{+w}{ }\PYG{n}{equal}\PYG{+w}{ }\PYG{l+s}{\PYGZdq{}}\PYG{l+s}{2.5 + 5*elaplong*dt}\PYG{l+s}{\PYGZdq{}}
\PYG{k}{variable }\PYG{n+nv+nvIdentifier}{x}\PYG{+w}{ }\PYG{n}{equal}\PYG{+w}{ }\PYG{n}{vdisplace}\PYG{n+nv}{(}\PYG{n+nv}{2.5,5}\PYG{n+nv}{)}
\end{sphinxVerbatim}

\sphinxAtStartPar
If a spherical indenter’s radius is specified as v\_r, then these
variable definitions will grow the size of the indenter at a specified
rate.

\begin{sphinxVerbatim}[commandchars=\\\{\}]
\PYG{k}{variable }\PYG{n+nv+nvIdentifier}{r0}\PYG{+w}{ }\PYG{n}{equal}\PYG{+w}{ }\PYG{l+m}{0.0}
\PYG{k}{variable }\PYG{n+nv+nvIdentifier}{rate}\PYG{+w}{ }\PYG{n}{equal}\PYG{+w}{ }\PYG{l+m}{1.0}
\PYG{k}{variable }\PYG{n+nv+nvIdentifier}{r}\PYG{+w}{ }\PYG{n}{equal}\PYG{+w}{ }\PYG{l+s}{\PYGZdq{}}\PYG{l+s}{v\PYGZus{}r0 + step*dt*v\PYGZus{}rate}\PYG{l+s}{\PYGZdq{}}
\end{sphinxVerbatim}

\sphinxAtStartPar
If the \sphinxstyleemphasis{side} keyword is specified as \sphinxstyleemphasis{out}, which is the default,
then particles outside the indenter are pushed away from its outer
surface, as described above.  This only applies to spherical,
cylindrical, and conical indenters.  If the \sphinxstyleemphasis{side} keyword is
specified as \sphinxstyleemphasis{in}, the action of the indenter is reversed.  Particles
inside the indenter are pushed away from its inner surface.  In other
words, the indenter is now a containing wall that traps the particles
inside it.  If the radius shrinks over time, it will squeeze the
particles.

\sphinxAtStartPar
The \sphinxstyleemphasis{units} keyword determines the meaning of the distance units used
to define the indenter geometry.  A \sphinxstyleemphasis{box} value selects standard
distance units as defined by the {\hyperref[\detokenize{units::doc}]{\sphinxcrossref{\DUrole{doc}{units}}}} command,
e.g. Angstroms for units = real or metal.  A \sphinxstyleemphasis{lattice} value means the
distance units are in lattice spacings.  The {\hyperref[\detokenize{lattice::doc}]{\sphinxcrossref{\DUrole{doc}{lattice}}}}
command must have been previously used to define the lattice spacing.
The (x,y,z) coords of the indenter position are scaled by the x,y,z
lattice spacings respectively.  The radius of a spherical or
cylindrical indenter is scaled by the x lattice spacing.

\sphinxAtStartPar
Note that the units keyword only affects indenter geometry parameters
specified directly with numbers, not those specified as variables.  In
the latter case, you should use the \sphinxstyleemphasis{xlat}, \sphinxstyleemphasis{ylat}, \sphinxstyleemphasis{zlat} keywords of
the {\hyperref[\detokenize{thermo_style::doc}]{\sphinxcrossref{\DUrole{doc}{thermo\_style}}}} command if you want to include
lattice spacings in a variable formula.

\sphinxAtStartPar
The force constant \sphinxstyleemphasis{K} is not affected by the \sphinxstyleemphasis{units} keyword.  It is
always in force/distance\textasciicircum{}2 units where force and distance are defined
by the {\hyperref[\detokenize{units::doc}]{\sphinxcrossref{\DUrole{doc}{units}}}} command.  If you wish K to be scaled by
the lattice spacing, you can define K with a variable whose formula
contains \sphinxstyleemphasis{xlat}, \sphinxstyleemphasis{ylat}, \sphinxstyleemphasis{zlat} keywords of the {\hyperref[\detokenize{thermo_style::doc}]{\sphinxcrossref{\DUrole{doc}{thermo\_style}}}} command, e.g.

\begin{sphinxVerbatim}[commandchars=\\\{\}]
\PYG{k}{variable }\PYG{n+nv+nvIdentifier}{k}\PYG{+w}{ }\PYG{n}{equal}\PYG{+w}{ }\PYG{l+m}{100.0}\PYG{o}{/}\PYG{n}{xlat}\PYG{o}{/}\PYG{n}{xlat}
\PYG{k}{fix }\PYG{n+nv+nvIdentifier}{1}\PYG{+w}{ }\PYG{n+nv+nvIdentifier}{all}\PYG{+w}{ }\PYG{n}{indent}\PYG{+w}{ }\PYG{n+nv}{\PYGZdl{}k}\PYG{+w}{ }\PYG{n}{sphere}\PYG{+w}{ }\PYG{n}{...}
\end{sphinxVerbatim}


\subsection{Restart, fix\_modify, output, run start/stop, minimize info}
\label{\detokenize{fix_indent:restart-fix-modify-output-run-start-stop-minimize-info}}
\sphinxAtStartPar
No information about this fix is written to {\hyperref[\detokenize{restart::doc}]{\sphinxcrossref{\DUrole{doc}{binary restart files}}}}.

\sphinxAtStartPar
The {\hyperref[\detokenize{fix_modify::doc}]{\sphinxcrossref{\DUrole{doc}{fix\_modify}}}} \sphinxstyleemphasis{energy} option is supported by
this fix to add the energy of interaction between atoms and the
indenter to the global potential energy of the system as part of
{\hyperref[\detokenize{thermo_style::doc}]{\sphinxcrossref{\DUrole{doc}{thermodynamic output}}}}.  The default setting for
this fix is {\hyperref[\detokenize{fix_modify::doc}]{\sphinxcrossref{\DUrole{doc}{fix\_modify energy no}}}}. The energy of
each particle interacting with the indenter is K/3 (r \sphinxhyphen{} R)\textasciicircum{}3.

\sphinxAtStartPar
The {\hyperref[\detokenize{fix_modify::doc}]{\sphinxcrossref{\DUrole{doc}{fix\_modify}}}} \sphinxstyleemphasis{respa} option is supported by this
fix. This allows to set at which level of the {\hyperref[\detokenize{run_style::doc}]{\sphinxcrossref{\DUrole{doc}{r\sphinxhyphen{}RESPA}}}} integrator the fix is adding its forces. Default is the
outermost level.

\sphinxAtStartPar
This fix computes a global scalar energy and a global 3\sphinxhyphen{}vector of
forces (on the indenter), which can be accessed by various
{\hyperref[\detokenize{Howto_output::doc}]{\sphinxcrossref{\DUrole{doc}{output commands}}}}.  The scalar and vector values
calculated by this fix are “extensive”.

\sphinxAtStartPar
The forces due to this fix are imposed during an energy minimization,
invoked by the {\hyperref[\detokenize{minimize::doc}]{\sphinxcrossref{\DUrole{doc}{minimize}}}} command.  Note that if you
define the indenter geometry with a variable using a time\sphinxhyphen{}dependent
formula, LAMMPS uses the iteration count in the minimizer as the
timestep.  But it is almost certainly a bad idea to have the indenter
change its position or size during a minimization.  LAMMPS does not
check if you have done this.

\begin{sphinxadmonition}{note}{Note:}
\sphinxAtStartPar
If you want the atom/indenter interaction energy to be included in
the total potential energy of the system (the quantity being
minimized), you must enable the {\hyperref[\detokenize{fix_modify::doc}]{\sphinxcrossref{\DUrole{doc}{fix\_modify}}}}
\sphinxstyleemphasis{energy} option for this fix.
\end{sphinxadmonition}


\subsection{Restrictions}
\label{\detokenize{fix_indent:restrictions}}\begin{quote}

\sphinxAtStartPar
none
\end{quote}


\subsection{Related commands}
\label{\detokenize{fix_indent:related-commands}}
\sphinxAtStartPar
none


\subsection{Default}
\label{\detokenize{fix_indent:default}}
\sphinxAtStartPar
The option defaults are side = out and units = lattice.

\sphinxstepscope

\index{fix ipi@\spxentry{fix ipi}}\ignorespaces 

\section{fix ipi command}
\label{\detokenize{fix_ipi:fix-ipi-command}}\label{\detokenize{fix_ipi:index-0}}\label{\detokenize{fix_ipi::doc}}

\subsection{Syntax}
\label{\detokenize{fix_ipi:syntax}}
\begin{sphinxVerbatim}[commandchars=\\\{\}]
\PYG{k}{fix }\PYG{n+nv+nvIdentifier}{ID}\PYG{+w}{ }\PYG{n+nv+nvIdentifier}{group\PYGZhy{}ID}\PYG{+w}{ }\PYG{n}{ipi}\PYG{+w}{ }\PYG{n}{address}\PYG{+w}{ }\PYG{n}{port}\PYG{+w}{ }\PYG{n}{[unix]}\PYG{+w}{ }\PYG{n}{[reset]}
\end{sphinxVerbatim}
\begin{itemize}
\item {} 
\sphinxAtStartPar
ID, group\sphinxhyphen{}ID are documented in {\hyperref[\detokenize{fix::doc}]{\sphinxcrossref{\DUrole{doc}{fix}}}} command

\item {} 
\sphinxAtStartPar
ipi = style name of this fix command

\item {} 
\sphinxAtStartPar
address = internet address (FQDN or IP), or UNIX socket name

\item {} 
\sphinxAtStartPar
port = port number (ignored for UNIX sockets)

\item {} 
\sphinxAtStartPar
zero or more keywords may be appended

\item {} 
\sphinxAtStartPar
keyword = \sphinxstyleemphasis{unix} or \sphinxstyleemphasis{reset}
\begin{sphinxalltt}
\sphinxstyleemphasis{unix} args = none = use a unix socket
\sphinxstyleemphasis{reset} args = none = reset electrostatics at each call
\end{sphinxalltt}

\end{itemize}


\subsection{Examples}
\label{\detokenize{fix_ipi:examples}}
\begin{sphinxVerbatim}[commandchars=\\\{\}]
\PYG{k}{fix }\PYG{n+nv+nvIdentifier}{1}\PYG{+w}{ }\PYG{n+nv+nvIdentifier}{all}\PYG{+w}{ }\PYG{n}{ipi}\PYG{+w}{ }\PYG{n}{my.server.com}\PYG{+w}{ }\PYG{l+m}{12345}
\PYG{k}{fix }\PYG{n+nv+nvIdentifier}{1}\PYG{+w}{ }\PYG{n+nv+nvIdentifier}{all}\PYG{+w}{ }\PYG{n}{ipi}\PYG{+w}{ }\PYG{n}{mysocket}\PYG{+w}{ }\PYG{l+m}{666}\PYG{+w}{ }\PYG{n}{unix}\PYG{+w}{ }\PYG{n}{reset}
\end{sphinxVerbatim}


\subsection{Description}
\label{\detokenize{fix_ipi:description}}
\sphinxAtStartPar
This fix enables LAMMPS to be run as a client for the i\sphinxhyphen{}PI Python
wrapper {\hyperref[\detokenize{fix_ipi:ipihome}]{\sphinxcrossref{\DUrole{std,std-ref}{(IPI)}}}}. i\sphinxhyphen{}PI is a universal force engine,
designed to perform advanced molecular simulations, with a special
focus on path integral molecular dynamics (PIMD) simulation.
The philosophy behind i\sphinxhyphen{}PI is to separate the evaluation of the
energy and forces, which is delegated to the client, and the evolution
of the dynamics, that is the responsibility of i\sphinxhyphen{}PI. This approach also
simplifies combining energies computed from different codes, which
can for instance be useful to mix first\sphinxhyphen{}principles calculations,
empirical force fields or machine\sphinxhyphen{}learning potentials.
The following publication {\hyperref[\detokenize{fix_ipi:ipicpc}]{\sphinxcrossref{\DUrole{std,std-ref}{(IPI\sphinxhyphen{}CPC\sphinxhyphen{}2014)}}}} discusses the
overall implementation of i\sphinxhyphen{}PI, and focuses on path\sphinxhyphen{}integral techniques,
while a later release {\hyperref[\detokenize{fix_ipi:ipicpc2}]{\sphinxcrossref{\DUrole{std,std-ref}{(IPI\sphinxhyphen{}CPC\sphinxhyphen{}2019)}}}} introduces several
additional features and simulation schemes.

\sphinxAtStartPar
The communication between i\sphinxhyphen{}PI and LAMMPS takes place using sockets,
and is reduced to the bare minimum. All the parameters of the dynamics
are specified in the input of i\sphinxhyphen{}PI, and all the parameters of the force
field must be specified as LAMMPS inputs, preceding the \sphinxstyleemphasis{fix ipi} command.

\sphinxAtStartPar
The server address must be specified by the \sphinxstyleemphasis{address} argument, and
can be either the IP address, the fully\sphinxhyphen{}qualified name of the server,
or the name of a UNIX socket for local, faster communication. In the
case of internet sockets, the \sphinxstyleemphasis{port} argument specifies the port
number on which i\sphinxhyphen{}PI is listening, while the \sphinxstyleemphasis{unix} optional switch
specifies that the socket is a UNIX socket.

\sphinxAtStartPar
Note that there is no check of data integrity, or that the atomic
configurations make sense. It is assumed that the species in the i\sphinxhyphen{}PI
input are listed in the same order as in the data file of LAMMPS. The
initial configuration is ignored, as it will be substituted with the
coordinates received from i\sphinxhyphen{}PI before forces are ever evaluated.

\sphinxAtStartPar
A note of caution when using potentials that contain long\sphinxhyphen{}range
electrostatics, or that contain parameters that depend on box size:
all of these options will be initialized based on the cell size in the
LAMMPS\sphinxhyphen{}side initial configuration and kept constant during the run.
This is required to e.g. obtain reproducible and conserved forces.
If the cell varies too wildly, it may be advisable to re\sphinxhyphen{}initialize
these interactions at each call. This behavior can be requested by
setting the \sphinxstyleemphasis{reset} switch.


\subsection{Obtaining i\sphinxhyphen{}PI}
\label{\detokenize{fix_ipi:obtaining-i-pi}}
\sphinxAtStartPar
Here are the commands to set up a virtual environment and install
i\sphinxhyphen{}PI into it with all its dependencies via the PyPI repository and
the pip package manager.

\begin{sphinxVerbatim}[commandchars=\\\{\}]
python\PYG{+w}{ }\PYGZhy{}m\PYG{+w}{ }venv\PYG{+w}{ }ipienv
\PYG{n+nb}{source}\PYG{+w}{ }ipienv/bin/activate
pip\PYG{+w}{ }install\PYG{+w}{ }\PYGZhy{}\PYGZhy{}upgrade\PYG{+w}{ }pip
pip\PYG{+w}{ }install\PYG{+w}{ }ipi
\end{sphinxVerbatim}


\subsection{Restart, fix\_modify, output, run start/stop, minimize info}
\label{\detokenize{fix_ipi:restart-fix-modify-output-run-start-stop-minimize-info}}
\sphinxAtStartPar
There is no restart information associated with this fix, since all
the dynamical parameters are dealt with by i\sphinxhyphen{}PI.


\subsection{Restrictions}
\label{\detokenize{fix_ipi:restrictions}}
\sphinxAtStartPar
Using this fix on anything other than all atoms requires particular
care, since i\sphinxhyphen{}PI will know nothing on atoms that are not those whose
coordinates are transferred. However, one could use this strategy to
define an external potential acting on the atoms that are moved by
i\sphinxhyphen{}PI.

\sphinxAtStartPar
Since the i\sphinxhyphen{}PI code uses atomic units internally, this fix needs to
convert LAMMPS data to and from its {\hyperref[\detokenize{units::doc}]{\sphinxcrossref{\DUrole{doc}{specified units}}}}
accordingly when communicating with i\sphinxhyphen{}PI.  This is not possible for
reduced units (“units lj”) and thus \sphinxstyleemphasis{fix ipi} will stop with an error in
this case.

\sphinxAtStartPar
This fix is part of the MISC package.  It is only enabled if
LAMMPS was built with that package.  See the
{\hyperref[\detokenize{Build_package::doc}]{\sphinxcrossref{\DUrole{doc}{Build package}}}} page for more info.
Because of the use of UNIX domain sockets, this fix will only
work in a UNIX environment.


\subsection{Related commands}
\label{\detokenize{fix_ipi:related-commands}}
\sphinxAtStartPar
{\hyperref[\detokenize{fix_nve::doc}]{\sphinxcrossref{\DUrole{doc}{fix nve}}}}


\bigskip\hrule\bigskip

\phantomsection\label{\detokenize{fix_ipi:ipicpc}}
\sphinxAtStartPar
\sphinxstylestrong{(IPI\sphinxhyphen{}CPC\sphinxhyphen{}2014)} Ceriotti, More and Manolopoulos, Comp Phys Comm 185,
1019\sphinxhyphen{}1026 (2014).

\phantomsection\label{\detokenize{fix_ipi:ipicpc2}}
\sphinxAtStartPar
\sphinxstylestrong{(IPI\sphinxhyphen{}CPC\sphinxhyphen{}2019)} Kapil et al., Comp Phys Comm 236, 214\sphinxhyphen{}223 (2019).

\phantomsection\label{\detokenize{fix_ipi:ipihome}}
\sphinxAtStartPar
\sphinxstylestrong{(IPI)}
\sphinxurl{https://ipi-code.org}

\sphinxstepscope

\index{fix langevin@\spxentry{fix langevin}}\ignorespaces 
\index{fix langevin/kk@\spxentry{fix langevin/kk}}\ignorespaces 

\section{fix langevin command}
\label{\detokenize{fix_langevin:fix-langevin-command}}\label{\detokenize{fix_langevin:index-1}}\label{\detokenize{fix_langevin:index-0}}\label{\detokenize{fix_langevin::doc}}
\sphinxAtStartPar
Accelerator Variants: \sphinxstyleemphasis{langevin/kk}


\subsection{Syntax}
\label{\detokenize{fix_langevin:syntax}}
\begin{sphinxVerbatim}[commandchars=\\\{\}]
\PYG{k}{fix }\PYG{n+nv+nvIdentifier}{ID}\PYG{+w}{ }\PYG{n+nv+nvIdentifier}{group\PYGZhy{}ID}\PYG{+w}{ }\PYG{n}{langevin}\PYG{+w}{ }\PYG{n}{Tstart}\PYG{+w}{ }\PYG{n}{Tstop}\PYG{+w}{ }\PYG{n}{damp}\PYG{+w}{ }\PYG{n}{seed}\PYG{+w}{ }\PYG{n}{keyword}\PYG{+w}{ }\PYG{n}{values}\PYG{+w}{ }\PYG{n}{...}
\end{sphinxVerbatim}
\begin{itemize}
\item {} 
\sphinxAtStartPar
ID, group\sphinxhyphen{}ID are documented in {\hyperref[\detokenize{fix::doc}]{\sphinxcrossref{\DUrole{doc}{fix}}}} command

\item {} 
\sphinxAtStartPar
langevin = style name of this fix command

\item {} 
\sphinxAtStartPar
Tstart,Tstop = desired temperature at start/end of run (temperature units)

\item {} 
\sphinxAtStartPar
Tstart can be a variable (see below)

\item {} 
\sphinxAtStartPar
damp = damping parameter (time units)

\item {} 
\sphinxAtStartPar
seed = random number seed to use for white noise (positive integer)

\item {} 
\sphinxAtStartPar
zero or more keyword/value pairs may be appended

\item {} 
\sphinxAtStartPar
keyword = \sphinxstyleemphasis{angmom} or \sphinxstyleemphasis{gjf} or \sphinxstyleemphasis{omega} or \sphinxstyleemphasis{scale} or \sphinxstyleemphasis{tally} or \sphinxstyleemphasis{zero}
\begin{sphinxalltt}
\sphinxstyleemphasis{angmom} value = \sphinxstyleemphasis{no} or factor
  \sphinxstyleemphasis{no} = do not thermostat rotational degrees of freedom via the angular momentum
  factor = do thermostat rotational degrees of freedom via the angular momentum and apply numeric scale factor as discussed below
\sphinxstyleemphasis{gjf} value = \sphinxstyleemphasis{no} or \sphinxstyleemphasis{vfull} or \sphinxstyleemphasis{vhalf}
  \sphinxstyleemphasis{no} = use standard formulation
  \sphinxstyleemphasis{vfull} = use Gronbech\sphinxhyphen{}Jensen/Farago formulation
  \sphinxstyleemphasis{vhalf} = use 2GJ formulation
\sphinxstyleemphasis{omega} value = \sphinxstyleemphasis{no} or \sphinxstyleemphasis{yes}
  \sphinxstyleemphasis{no} = do not thermostat rotational degrees of freedom via the angular velocity
  \sphinxstyleemphasis{yes} = do thermostat rotational degrees of freedom via the angular velocity
\sphinxstyleemphasis{scale} values = type ratio
  type = atom type (1\sphinxhyphen{}N)
  ratio = factor by which to scale the damping coefficient
\sphinxstyleemphasis{tally} value = \sphinxstyleemphasis{no} or \sphinxstyleemphasis{yes}
  \sphinxstyleemphasis{no} = do not tally the energy added/subtracted to atoms
  \sphinxstyleemphasis{yes} = do tally the energy added/subtracted to atoms
\sphinxstyleemphasis{zero} value = \sphinxstyleemphasis{no} or \sphinxstyleemphasis{yes}
  \sphinxstyleemphasis{no} = do not set total random force to zero
  \sphinxstyleemphasis{yes} = set total random force to zero
\end{sphinxalltt}

\end{itemize}


\subsection{Examples}
\label{\detokenize{fix_langevin:examples}}
\begin{sphinxVerbatim}[commandchars=\\\{\}]
\PYG{k}{fix }\PYG{n+nv+nvIdentifier}{3}\PYG{+w}{ }\PYG{n+nv+nvIdentifier}{boundary}\PYG{+w}{ }\PYG{n}{langevin}\PYG{+w}{ }\PYG{l+m}{1.0}\PYG{+w}{ }\PYG{l+m}{1.0}\PYG{+w}{ }\PYG{l+m}{1000.0}\PYG{+w}{ }\PYG{l+m}{699483}
\PYG{k}{fix }\PYG{n+nv+nvIdentifier}{1}\PYG{+w}{ }\PYG{n+nv+nvIdentifier}{all}\PYG{+w}{ }\PYG{n}{langevin}\PYG{+w}{ }\PYG{l+m}{1.0}\PYG{+w}{ }\PYG{l+m}{1.1}\PYG{+w}{ }\PYG{l+m}{100.0}\PYG{+w}{ }\PYG{l+m}{48279}\PYG{+w}{ }\PYG{n}{scale}\PYG{+w}{ }\PYG{l+m}{3}\PYG{+w}{ }\PYG{l+m}{1.5}
\PYG{k}{fix }\PYG{n+nv+nvIdentifier}{1}\PYG{+w}{ }\PYG{n+nv+nvIdentifier}{all}\PYG{+w}{ }\PYG{n}{langevin}\PYG{+w}{ }\PYG{l+m}{1.0}\PYG{+w}{ }\PYG{l+m}{1.1}\PYG{+w}{ }\PYG{l+m}{100.0}\PYG{+w}{ }\PYG{l+m}{48279}\PYG{+w}{ }\PYG{n}{angmom}\PYG{+w}{ }\PYG{l+m}{3.333}
\end{sphinxVerbatim}


\subsection{Description}
\label{\detokenize{fix_langevin:description}}
\sphinxAtStartPar
Apply a Langevin thermostat as described in {\hyperref[\detokenize{fix_langevin:schneider1}]{\sphinxcrossref{\DUrole{std,std-ref}{(Schneider)}}}}
to a group of atoms which models an interaction with a background
implicit solvent.  Used with {\hyperref[\detokenize{fix_nve::doc}]{\sphinxcrossref{\DUrole{doc}{fix nve}}}}, this command
performs Brownian dynamics (BD), since the total force on each atom
will have the form:
\begin{equation*}
\begin{split}F = & F_c + F_f + F_r \\
F_f = & - \frac{m}{\mathrm{damp}} v \\
F_r \propto & \sqrt{\frac{k_B T m}{dt~\mathrm{damp}}}\end{split}
\end{equation*}
\sphinxAtStartPar
\(F_c\) is the conservative force computed via the usual
inter\sphinxhyphen{}particle interactions ({\hyperref[\detokenize{pair_style::doc}]{\sphinxcrossref{\DUrole{doc}{pair\_style}}}},
{\hyperref[\detokenize{bond_style::doc}]{\sphinxcrossref{\DUrole{doc}{bond\_style}}}}, etc).  The \(F_f\) and \(F_r\)
terms are added by this fix on a per\sphinxhyphen{}particle basis.  See the
{\hyperref[\detokenize{pair_dpd::doc}]{\sphinxcrossref{\DUrole{doc}{pair\_style dpd/tstat}}}} command for a thermostatting
option that adds similar terms on a pairwise basis to pairs of
interacting particles.

\sphinxAtStartPar
\(F_f\) is a frictional drag or viscous damping term proportional to
the particle’s velocity.  The proportionality constant for each atom is
computed as \(\frac{m}{\mathrm{damp}}\), where \sphinxstyleemphasis{m} is the mass of the
particle and damp is the damping factor specified by the user.

\sphinxAtStartPar
\(F_r\) is a force due to solvent atoms at a temperature \(T\)
randomly bumping into the particle.  As derived from the
fluctuation/dissipation theorem, its magnitude as shown above is
proportional to \(\sqrt{\frac{k_B T m}{dt~\mathrm{damp}}}\), where
\(k_B\) is the Boltzmann constant, \(T\) is the desired temperature,
\sphinxstyleemphasis{m} is the mass of the particle, \sphinxstyleemphasis{dt} is the timestep size, and damp is
the damping factor.  Random numbers are used to randomize the direction
and magnitude of this force as described in {\hyperref[\detokenize{fix_langevin:dunweg1}]{\sphinxcrossref{\DUrole{std,std-ref}{(Dunweg)}}}},
where a uniform random number is used (instead of a Gaussian random
number) for speed.

\sphinxAtStartPar
Note that unless you use the \sphinxstyleemphasis{omega} or \sphinxstyleemphasis{angmom} keywords, the
thermostat effect of this fix is applied to only the translational
degrees of freedom for the particles, which is an important
consideration for finite\sphinxhyphen{}size particles, which have rotational degrees
of freedom, are being thermostatted.  The translational degrees of
freedom can also have a bias velocity removed from them before
thermostatting takes place; see the description below.

\begin{sphinxadmonition}{note}{Note:}
\sphinxAtStartPar
Unlike the {\hyperref[\detokenize{fix_nh::doc}]{\sphinxcrossref{\DUrole{doc}{fix nvt}}}} command which performs Nose/Hoover
thermostatting AND time integration, this fix does NOT perform time
integration.  It only modifies forces to effect thermostatting.  Thus
you must use a separate time integration fix, like {\hyperref[\detokenize{fix_nve::doc}]{\sphinxcrossref{\DUrole{doc}{fix nve}}}} to actually update the velocities and positions of atoms
using the modified forces.  Likewise, this fix should not normally be
used on atoms that also have their temperature controlled by another
fix \sphinxhyphen{} e.g. by {\hyperref[\detokenize{fix_nh::doc}]{\sphinxcrossref{\DUrole{doc}{fix nvt}}}} or {\hyperref[\detokenize{fix_temp_rescale::doc}]{\sphinxcrossref{\DUrole{doc}{fix temp/rescale}}}} commands.
\end{sphinxadmonition}

\sphinxAtStartPar
See the {\hyperref[\detokenize{Howto_thermostat::doc}]{\sphinxcrossref{\DUrole{doc}{Howto thermostat}}}} page for
a discussion of different ways to compute temperature and perform
thermostatting.

\sphinxAtStartPar
The desired temperature at each timestep is a ramped value during the
run from \sphinxstyleemphasis{Tstart} to \sphinxstyleemphasis{Tstop}.

\sphinxAtStartPar
\sphinxstyleemphasis{Tstart} can be specified as an equal\sphinxhyphen{}style or atom\sphinxhyphen{}style
{\hyperref[\detokenize{variable::doc}]{\sphinxcrossref{\DUrole{doc}{variable}}}}.  In this case, the \sphinxstyleemphasis{Tstop} setting is
ignored.  If the value is a variable, it should be specified as
v\_name, where name is the variable name.  In this case, the variable
will be evaluated each timestep, and its value used to determine the
target temperature.

\sphinxAtStartPar
Equal\sphinxhyphen{}style variables can specify formulas with various mathematical
functions, and include {\hyperref[\detokenize{thermo_style::doc}]{\sphinxcrossref{\DUrole{doc}{thermo\_style}}}} command
keywords for the simulation box parameters and timestep and elapsed
time.  Thus it is easy to specify a time\sphinxhyphen{}dependent temperature.

\sphinxAtStartPar
Atom\sphinxhyphen{}style variables can specify the same formulas as equal\sphinxhyphen{}style
variables but can also include per\sphinxhyphen{}atom values, such as atom
coordinates.  Thus it is easy to specify a spatially\sphinxhyphen{}dependent
temperature with optional time\sphinxhyphen{}dependence as well.

\sphinxAtStartPar
Like other fixes that perform thermostatting, this fix can be used
with {\hyperref[\detokenize{compute::doc}]{\sphinxcrossref{\DUrole{doc}{compute commands}}}} that remove a “bias” from the
atom velocities.  E.g. to apply the thermostat only to atoms within a
spatial {\hyperref[\detokenize{region::doc}]{\sphinxcrossref{\DUrole{doc}{region}}}}, or to remove the center\sphinxhyphen{}of\sphinxhyphen{}mass
velocity from a group of atoms, or to remove the x\sphinxhyphen{}component of
velocity from the calculation.

\sphinxAtStartPar
This is not done by default, but only if the {\hyperref[\detokenize{fix_modify::doc}]{\sphinxcrossref{\DUrole{doc}{fix\_modify}}}} command is used to assign a temperature compute to this
fix that includes such a bias term.  See the doc pages for individual
{\hyperref[\detokenize{compute::doc}]{\sphinxcrossref{\DUrole{doc}{compute temp commands}}}} to determine which ones include
a bias.  In this case, the thermostat works in the following manner:
bias is removed from each atom, thermostatting is performed on the
remaining thermal degrees of freedom, and the bias is added back in.

\sphinxAtStartPar
The \sphinxstyleemphasis{damp} parameter is specified in time units and determines how
rapidly the temperature is relaxed.  For example, a value of 100.0 means
to relax the temperature in a timespan of (roughly) 100 time units
(\(\tau\) or fs or ps \sphinxhyphen{} see the {\hyperref[\detokenize{units::doc}]{\sphinxcrossref{\DUrole{doc}{units}}}} command).  The
damp factor can be thought of as inversely related to the viscosity of
the solvent.  I.e. a small relaxation time implies a high\sphinxhyphen{}viscosity
solvent and vice versa.  See the discussion about \(\gamma\) and
viscosity in the documentation for the {\hyperref[\detokenize{fix_viscous::doc}]{\sphinxcrossref{\DUrole{doc}{fix viscous}}}}
command for more details.

\sphinxAtStartPar
The random \# \sphinxstyleemphasis{seed} must be a positive integer.  A Marsaglia random
number generator is used.  Each processor uses the input seed to
generate its own unique seed and its own stream of random numbers.
Thus the dynamics of the system will not be identical on two runs on
different numbers of processors.


\bigskip\hrule\bigskip


\sphinxAtStartPar
The keyword/value option pairs are used in the following ways.

\sphinxAtStartPar
The keyword \sphinxstyleemphasis{angmom} and \sphinxstyleemphasis{omega} keywords enable thermostatting of
rotational degrees of freedom in addition to the usual translational
degrees of freedom.  This can only be done for finite\sphinxhyphen{}size particles.

\sphinxAtStartPar
A simulation using atom\_style sphere defines an omega for finite\sphinxhyphen{}size
spheres.  A simulation using atom\_style ellipsoid defines a finite
size and shape for aspherical particles and an angular momentum.
The Langevin formulas for thermostatting the rotational degrees of
freedom are the same as those above, where force is replaced by
torque, m is replaced by the moment of inertia I, and v is replaced by
omega (which is derived from the angular momentum in the case of
aspherical particles).

\sphinxAtStartPar
The rotational temperature of the particles can be monitored by the
{\hyperref[\detokenize{compute_temp_sphere::doc}]{\sphinxcrossref{\DUrole{doc}{compute temp/sphere}}}} and {\hyperref[\detokenize{compute_temp_asphere::doc}]{\sphinxcrossref{\DUrole{doc}{compute
temp/asphere}}}} commands with their rotate
options.

\sphinxAtStartPar
For the \sphinxstyleemphasis{omega} keyword there is also a scale factor of
\(\frac{10.0}{3.0}\) that is applied as a multiplier on the
\(F_f\) (damping) term in the equation above and of
\(\sqrt{\frac{10.0}{3.0}}\) as a multiplier on the \(F_r\) term.
This does not affect the thermostatting behavior of the Langevin
formalism but ensures that the randomized rotational diffusivity of
spherical particles is correct.

\sphinxAtStartPar
For the \sphinxstyleemphasis{angmom} keyword a similar scale factor is needed which is
\(\frac{10.0}{3.0}\) for spherical particles, but is anisotropic for
aspherical particles (e.g. ellipsoids).  Currently LAMMPS only applies
an isotropic scale factor, and you can choose its magnitude as the
specified value of the \sphinxstyleemphasis{angmom} keyword.  If your aspherical particles
are (nearly) spherical than a value of \(\frac{10.0}{3.0} =
3.\overline{3}\) is a good choice.  If they are highly aspherical, a
value of 1.0 is as good a choice as any, since the effects on rotational
diffusivity of the particles will be incorrect regardless.  Note that
for any reasonable scale factor, the thermostatting effect of the
\sphinxstyleemphasis{angmom} keyword on the rotational temperature of the aspherical
particles should still be valid.

\sphinxAtStartPar
The keyword \sphinxstyleemphasis{scale} allows the damp factor to be scaled up or down by
the specified factor for atoms of that type.  This can be useful when
different atom types have different sizes or masses.  It can be used
multiple times to adjust damp for several atom types.  Note that
specifying a ratio of 2 increases the relaxation time which is
equivalent to the solvent’s viscosity acting on particles with
\(\frac{1}{2}\) the diameter.  This is the opposite effect of scale
factors used by the {\hyperref[\detokenize{fix_viscous::doc}]{\sphinxcrossref{\DUrole{doc}{fix viscous}}}} command, since the
damp factor in fix \sphinxstyleemphasis{langevin} is inversely related to the \(\gamma\)
factor in fix \sphinxstyleemphasis{viscous}.  Also note that the damping factor in fix
\sphinxstyleemphasis{langevin} includes the particle mass in Ff, unlike fix \sphinxstyleemphasis{viscous}.
Thus the mass and size of different atom types should be accounted for
in the choice of ratio values.

\sphinxAtStartPar
The keyword \sphinxstyleemphasis{tally} enables the calculation of the cumulative energy
added/subtracted to the atoms as they are thermostatted.  Effectively
it is the energy exchanged between the infinite thermal reservoir and
the particles.  As described below, this energy can then be printed
out or added to the potential energy of the system to monitor energy
conservation.

\begin{sphinxadmonition}{note}{Note:}
\sphinxAtStartPar
This accumulated energy does NOT include kinetic energy removed
by the \sphinxstyleemphasis{zero} flag. LAMMPS will print a warning when both options are
active.
\end{sphinxadmonition}

\sphinxAtStartPar
The keyword \sphinxstyleemphasis{zero} can be used to eliminate drift due to the
thermostat. Because the random forces on different atoms are
independent, they do not sum exactly to zero.  As a result, this fix
applies a small random force to the entire system, and the
center\sphinxhyphen{}of\sphinxhyphen{}mass of the system undergoes a slow random walk.  If the
keyword \sphinxstyleemphasis{zero} is set to \sphinxstyleemphasis{yes}, the total random force is set exactly
to zero by subtracting off an equal part of it from each atom in the
group.  As a result, the center\sphinxhyphen{}of\sphinxhyphen{}mass of a system with zero initial
momentum will not drift over time.

\sphinxAtStartPar
The keyword \sphinxstyleemphasis{gjf} can be used to run the {\hyperref[\detokenize{fix_langevin:gronbech-jensen}]{\sphinxcrossref{\DUrole{std,std-ref}{Gronbech\sphinxhyphen{}Jensen/Farago}}}} time\sphinxhyphen{}discretization of the Langevin model.  As
described in the papers cited below, the purpose of this method is to
enable longer timesteps to be used (up to the numerical stability
limit of the integrator), while still producing the correct Boltzmann
distribution of atom positions.

\sphinxAtStartPar
The current implementation provides the user with the option to output
the velocity in one of two forms: \sphinxstyleemphasis{vfull} or \sphinxstyleemphasis{vhalf}, which replaces
the outdated option \sphinxstyleemphasis{yes}. The \sphinxstyleemphasis{gjf} option \sphinxstyleemphasis{vfull} outputs the
on\sphinxhyphen{}site velocity given in {\hyperref[\detokenize{fix_langevin:gronbech-jensen}]{\sphinxcrossref{\DUrole{std,std-ref}{Gronbech\sphinxhyphen{}Jensen/Farago}}}}; this velocity is shown to be systematically lower
than the target temperature by a small amount, which grows
quadratically with the timestep.  The \sphinxstyleemphasis{gjf} option \sphinxstyleemphasis{vhalf} outputs the
2GJ half\sphinxhyphen{}step velocity given in {\hyperref[\detokenize{fix_langevin:id1}]{\sphinxcrossref{\DUrole{std,std-ref}{Gronbech Jensen/Gronbech\sphinxhyphen{}Jensen}}}}; for linear systems, this velocity is shown to not
have any statistical errors for any stable time step.  An overview of
statistically correct Boltzmann and Maxwell\sphinxhyphen{}Boltzmann sampling of true
on\sphinxhyphen{}site and true half\sphinxhyphen{}step velocities is given in
{\hyperref[\detokenize{fix_langevin:id2}]{\sphinxcrossref{\DUrole{std,std-ref}{Gronbech\sphinxhyphen{}Jensen}}}}.  Regardless of the choice
of output velocity, the sampling of the configurational distribution
of atom positions is the same, and linearly consistent with the target
temperature.


\bigskip\hrule\bigskip


\sphinxAtStartPar
Styles with a \sphinxstyleemphasis{gpu}, \sphinxstyleemphasis{intel}, \sphinxstyleemphasis{kk}, \sphinxstyleemphasis{omp}, or \sphinxstyleemphasis{opt} suffix are
functionally the same as the corresponding style without the suffix.
They have been optimized to run faster, depending on your available
hardware, as discussed on the {\hyperref[\detokenize{Speed_packages::doc}]{\sphinxcrossref{\DUrole{doc}{Accelerator packages}}}}
page.  The accelerated styles take the same arguments and should
produce the same results, except for round\sphinxhyphen{}off and precision issues.

\sphinxAtStartPar
These accelerated styles are part of the GPU, INTEL, KOKKOS,
OPENMP, and OPT packages, respectively.  They are only enabled if
LAMMPS was built with those packages.  See the {\hyperref[\detokenize{Build_package::doc}]{\sphinxcrossref{\DUrole{doc}{Build package}}}} page for more info.

\sphinxAtStartPar
You can specify the accelerated styles explicitly in your input script
by including their suffix, or you can use the {\hyperref[\detokenize{Run_options::doc}]{\sphinxcrossref{\DUrole{doc}{\sphinxhyphen{}suffix command\sphinxhyphen{}line switch}}}} when you invoke LAMMPS, or you can use the
{\hyperref[\detokenize{suffix::doc}]{\sphinxcrossref{\DUrole{doc}{suffix}}}} command in your input script.

\sphinxAtStartPar
See the {\hyperref[\detokenize{Speed_packages::doc}]{\sphinxcrossref{\DUrole{doc}{Accelerator packages}}}} page for more
instructions on how to use the accelerated styles effectively.


\bigskip\hrule\bigskip



\subsection{Restart, fix\_modify, output, run start/stop, minimize info}
\label{\detokenize{fix_langevin:restart-fix-modify-output-run-start-stop-minimize-info}}
\sphinxAtStartPar
No information about this fix is written to {\hyperref[\detokenize{restart::doc}]{\sphinxcrossref{\DUrole{doc}{binary restart files}}}}.  Because the state of the random number generator
is not saved in restart files, this means you cannot do “exact”
restarts with this fix, where the simulation continues on the same as
if no restart had taken place.  However, in a statistical sense, a
restarted simulation should produce the same behavior.

\sphinxAtStartPar
The {\hyperref[\detokenize{fix_modify::doc}]{\sphinxcrossref{\DUrole{doc}{fix\_modify}}}} \sphinxstyleemphasis{temp} option is supported by this
fix.  You can use it to assign a temperature {\hyperref[\detokenize{compute::doc}]{\sphinxcrossref{\DUrole{doc}{compute}}}}
you have defined to this fix which will be used in its thermostatting
procedure, as described above.  For consistency, the group used by
this fix and by the compute should be the same.

\sphinxAtStartPar
The cumulative energy change in the system imposed by this fix is
included in the {\hyperref[\detokenize{thermo_style::doc}]{\sphinxcrossref{\DUrole{doc}{thermodynamic output}}}} keywords
\sphinxstyleemphasis{ecouple} and \sphinxstyleemphasis{econserve}, but only if the \sphinxstyleemphasis{tally} keyword to set to
\sphinxstyleemphasis{yes}.  See the {\hyperref[\detokenize{thermo_style::doc}]{\sphinxcrossref{\DUrole{doc}{thermo\_style}}}} page for
details.

\sphinxAtStartPar
This fix computes a global scalar which can be accessed by various
{\hyperref[\detokenize{Howto_output::doc}]{\sphinxcrossref{\DUrole{doc}{output commands}}}}.  The scalar is the same
cumulative energy change due to this fix described in the previous
paragraph.  The scalar value calculated by this fix is “extensive”.
Note that calculation of this quantity also requires setting the
\sphinxstyleemphasis{tally} keyword to \sphinxstyleemphasis{yes}.

\sphinxAtStartPar
This fix can ramp its target temperature over multiple runs, using the
\sphinxstyleemphasis{start} and \sphinxstyleemphasis{stop} keywords of the {\hyperref[\detokenize{run::doc}]{\sphinxcrossref{\DUrole{doc}{run}}}} command.  See the
{\hyperref[\detokenize{run::doc}]{\sphinxcrossref{\DUrole{doc}{run}}}} command for details of how to do this.

\sphinxAtStartPar
This fix is not invoked during {\hyperref[\detokenize{minimize::doc}]{\sphinxcrossref{\DUrole{doc}{energy minimization}}}}.


\subsection{Restrictions}
\label{\detokenize{fix_langevin:restrictions}}
\sphinxAtStartPar
For \sphinxstyleemphasis{gjf} do not choose damp=dt/2. \sphinxstyleemphasis{gjf} is not compatible
with run\_style respa.


\subsection{Related commands}
\label{\detokenize{fix_langevin:related-commands}}
\sphinxAtStartPar
{\hyperref[\detokenize{fix_nh::doc}]{\sphinxcrossref{\DUrole{doc}{fix nvt}}}}, {\hyperref[\detokenize{fix_temp_rescale::doc}]{\sphinxcrossref{\DUrole{doc}{fix temp/rescale}}}}, {\hyperref[\detokenize{fix_viscous::doc}]{\sphinxcrossref{\DUrole{doc}{fix viscous}}}}, {\hyperref[\detokenize{fix_nh::doc}]{\sphinxcrossref{\DUrole{doc}{fix nvt}}}}, {\hyperref[\detokenize{pair_dpd::doc}]{\sphinxcrossref{\DUrole{doc}{pair\_style dpd/tstat}}}}


\subsection{Default}
\label{\detokenize{fix_langevin:default}}
\sphinxAtStartPar
The option defaults are angmom = no, omega = no, scale = 1.0 for all
types, tally = no, zero = no, gjf = no.


\bigskip\hrule\bigskip

\phantomsection\label{\detokenize{fix_langevin:dunweg1}}
\sphinxAtStartPar
\sphinxstylestrong{(Dunweg)} Dunweg and Paul, Int J of Modern Physics C, 2, 817\sphinxhyphen{}27 (1991).

\phantomsection\label{\detokenize{fix_langevin:schneider1}}
\sphinxAtStartPar
\sphinxstylestrong{(Schneider)} Schneider and Stoll, Phys Rev B, 17, 1302 (1978).

\phantomsection\label{\detokenize{fix_langevin:gronbech-jensen}}
\sphinxAtStartPar
\sphinxstylestrong{(Gronbech\sphinxhyphen{}Jensen)} Gronbech\sphinxhyphen{}Jensen and Farago, Mol Phys, 111, 983
(2013); Gronbech\sphinxhyphen{}Jensen, Hayre, and Farago, Comp Phys Comm, 185, 524 (2014)

\phantomsection\label{\detokenize{fix_langevin:id1}}
\sphinxAtStartPar
\sphinxstylestrong{(Gronbech\sphinxhyphen{}Jensen)} Gronbech Jensen and Gronbech\sphinxhyphen{}Jensen, Mol Phys, 117, 2511 (2019)

\phantomsection\label{\detokenize{fix_langevin:id2}}
\sphinxAtStartPar
\sphinxstylestrong{(Gronbech\sphinxhyphen{}Jensen)} Gronbech\sphinxhyphen{}Jensen, Mol Phys (2019); \sphinxurl{https://doi.org/10.1080/00268976.2019.1662506}

\sphinxstepscope

\index{fix langevin/drude@\spxentry{fix langevin/drude}}\ignorespaces 

\section{fix langevin/drude command}
\label{\detokenize{fix_langevin_drude:fix-langevin-drude-command}}\label{\detokenize{fix_langevin_drude:index-0}}\label{\detokenize{fix_langevin_drude::doc}}

\subsection{Syntax}
\label{\detokenize{fix_langevin_drude:syntax}}
\begin{sphinxVerbatim}[commandchars=\\\{\}]
\PYG{k}{fix }\PYG{n+nv+nvIdentifier}{ID}\PYG{+w}{ }\PYG{n+nv+nvIdentifier}{group\PYGZhy{}ID}\PYG{+w}{ }\PYG{n}{langevin}\PYG{o}{/}\PYG{n}{drude}\PYG{+w}{ }\PYG{n}{Tcom}\PYG{+w}{ }\PYG{n}{damp\PYGZus{}com}\PYG{+w}{ }\PYG{n}{seed\PYGZus{}com}\PYG{+w}{ }\PYG{n}{Tdrude}\PYG{+w}{ }\PYG{n}{damp\PYGZus{}drude}\PYG{+w}{ }\PYG{n}{seed\PYGZus{}drude}\PYG{+w}{ }\PYG{n}{keyword}\PYG{+w}{ }\PYG{n}{values}\PYG{+w}{ }\PYG{n}{...}
\end{sphinxVerbatim}
\begin{itemize}
\item {} 
\sphinxAtStartPar
ID, group\sphinxhyphen{}ID are documented in {\hyperref[\detokenize{fix::doc}]{\sphinxcrossref{\DUrole{doc}{fix}}}} command

\item {} 
\sphinxAtStartPar
langevin/drude = style name of this fix command

\item {} 
\sphinxAtStartPar
Tcom = desired temperature of the centers of mass (temperature units)

\item {} 
\sphinxAtStartPar
damp\_com = damping parameter for the thermostat on centers of mass (time units)

\item {} 
\sphinxAtStartPar
seed\_com = random number seed to use for white noise of the thermostat on centers of mass (positive integer)

\item {} 
\sphinxAtStartPar
Tdrude = desired temperature of the Drude oscillators (temperature units)

\item {} 
\sphinxAtStartPar
damp\_drude = damping parameter for the thermostat on Drude oscillators (time units)

\item {} 
\sphinxAtStartPar
seed\_drude = random number seed to use for white noise of the thermostat on Drude oscillators (positive integer)

\item {} 
\sphinxAtStartPar
zero or more keyword/value pairs may be appended

\item {} 
\sphinxAtStartPar
keyword = \sphinxstyleemphasis{zero}
\begin{sphinxalltt}
\sphinxstyleemphasis{zero} value = \sphinxstyleemphasis{no} or \sphinxstyleemphasis{yes}
  \sphinxstyleemphasis{no} = do not set total random force on centers of mass to zero
  \sphinxstyleemphasis{yes} = set total random force on centers of mass to zero
\end{sphinxalltt}

\end{itemize}


\subsection{Examples}
\label{\detokenize{fix_langevin_drude:examples}}
\begin{sphinxVerbatim}[commandchars=\\\{\}]
\PYG{k}{fix }\PYG{n+nv+nvIdentifier}{3}\PYG{+w}{ }\PYG{n+nv+nvIdentifier}{all}\PYG{+w}{ }\PYG{n}{langevin}\PYG{o}{/}\PYG{n}{drude}\PYG{+w}{ }\PYG{l+m}{300.0}\PYG{+w}{ }\PYG{l+m}{100.0}\PYG{+w}{ }\PYG{l+m}{19377}\PYG{+w}{ }\PYG{l+m}{1.0}\PYG{+w}{ }\PYG{l+m}{20.0}\PYG{+w}{ }\PYG{l+m}{83451}
\PYG{k}{fix }\PYG{n+nv+nvIdentifier}{1}\PYG{+w}{ }\PYG{n+nv+nvIdentifier}{all}\PYG{+w}{ }\PYG{n}{langevin}\PYG{o}{/}\PYG{n}{drude}\PYG{+w}{ }\PYG{l+m}{298.15}\PYG{+w}{ }\PYG{l+m}{100.0}\PYG{+w}{ }\PYG{l+m}{19377}\PYG{+w}{ }\PYG{l+m}{5.0}\PYG{+w}{ }\PYG{l+m}{10.0}\PYG{+w}{ }\PYG{l+m}{83451}\PYG{+w}{ }\PYG{n}{zero}\PYG{+w}{ }\PYG{n}{yes}
\end{sphinxVerbatim}

\sphinxAtStartPar
Example input scripts available: examples/PACKAGES/drude


\subsection{Description}
\label{\detokenize{fix_langevin_drude:description}}
\sphinxAtStartPar
Apply two Langevin thermostats as described in {\hyperref[\detokenize{fix_langevin_drude:jiang1}]{\sphinxcrossref{\DUrole{std,std-ref}{(Jiang1)}}}} for
thermalizing the reduced degrees of freedom of Drude oscillators.
This link describes how to use the {\hyperref[\detokenize{Howto_drude::doc}]{\sphinxcrossref{\DUrole{doc}{thermalized Drude oscillator model}}}} in LAMMPS and polarizable models in LAMMPS
are discussed on the {\hyperref[\detokenize{Howto_polarizable::doc}]{\sphinxcrossref{\DUrole{doc}{Howto polarizable}}}} doc
page.

\sphinxAtStartPar
Drude oscillators are a way to simulate polarizables atoms, by
splitting them into a core and a Drude particle bound by a harmonic
bond.  The thermalization works by transforming the particles degrees
of freedom by these equations.  In these equations upper case denotes
atomic or center of mass values and lower case denotes Drude particle
or dipole values. Primes denote the transformed (reduced) values,
while bare letters denote the original values.

\sphinxAtStartPar
Velocities:
\begin{equation*}
\begin{split}V' = \frac {M\, V + m\, v} {M'}\end{split}
\end{equation*}\begin{equation*}
\begin{split}v' = v - V\end{split}
\end{equation*}
\sphinxAtStartPar
Masses:
\begin{equation*}
\begin{split}M' = M + m\end{split}
\end{equation*}\begin{equation*}
\begin{split}m' = \frac {M\, m } {M'}\end{split}
\end{equation*}
\sphinxAtStartPar
The Langevin forces are computed as
\begin{equation*}
\begin{split}F' = - \frac {M'} {\mathtt{damp_com}}\, V' + F_r'\end{split}
\end{equation*}\begin{equation*}
\begin{split}f' = - \frac {m'} {\mathtt{damp_drude}}\, v' + f_r'\end{split}
\end{equation*}
\sphinxAtStartPar
\(F_r'\) is a random force proportional to
\(\sqrt { \frac {2\, k_B \mathtt{Tcom}\, m'}                  {\mathrm dt\, \mathtt{damp_com} }         }\).
\(f_r'\) is a random force proportional to
\(\sqrt { \frac {2\, k_B \mathtt{Tdrude}\, m'}                  {\mathrm dt\, \mathtt{damp_drude} }         }\).
Then the real forces acting on the particles are computed from the inverse
transform:
\begin{equation*}
\begin{split}F = \frac M {M'}\, F' - f'\end{split}
\end{equation*}\begin{equation*}
\begin{split}f = \frac m {M'}\, F' + f'\end{split}
\end{equation*}
\sphinxAtStartPar
This fix also thermostats non\sphinxhyphen{}polarizable atoms in the group at
temperature \sphinxstyleemphasis{Tcom}, as if they had a massless Drude partner.  The
Drude particles themselves need not be in the group. The center of
mass and the dipole are thermostatted iff the core atom is in the
group.

\sphinxAtStartPar
Note that the thermostat effect of this fix is applied to only the
translational degrees of freedom of the particles, which is an
important consideration if finite\sphinxhyphen{}size particles, which have
rotational degrees of freedom, are being thermostatted. The
translational degrees of freedom can also have a bias velocity removed
from them before thermostatting takes place; see the description below.

\begin{sphinxadmonition}{note}{Note:}
\sphinxAtStartPar
Like the {\hyperref[\detokenize{fix_langevin::doc}]{\sphinxcrossref{\DUrole{doc}{fix langevin}}}} command, this fix does
NOT perform time integration. It only modifies forces to effect
thermostatting. Thus you must use a separate time integration fix, like
{\hyperref[\detokenize{fix_nve::doc}]{\sphinxcrossref{\DUrole{doc}{fix nve}}}} or {\hyperref[\detokenize{fix_nh::doc}]{\sphinxcrossref{\DUrole{doc}{fix nph}}}} to actually update the
velocities and positions of atoms using the modified forces.
Likewise, this fix should not normally be used on atoms that also have
their temperature controlled by another fix \sphinxhyphen{} e.g. by {\hyperref[\detokenize{fix_nh::doc}]{\sphinxcrossref{\DUrole{doc}{fix nvt}}}} or {\hyperref[\detokenize{fix_temp_rescale::doc}]{\sphinxcrossref{\DUrole{doc}{fix temp/rescale}}}} commands.
\end{sphinxadmonition}

\sphinxAtStartPar
See the {\hyperref[\detokenize{Howto_thermostat::doc}]{\sphinxcrossref{\DUrole{doc}{Howto thermostat}}}} page for a
discussion of different ways to compute temperature and perform
thermostatting.


\bigskip\hrule\bigskip


\sphinxAtStartPar
This fix requires each atom know whether it is a Drude particle or
not.  You must therefore use the {\hyperref[\detokenize{fix_drude::doc}]{\sphinxcrossref{\DUrole{doc}{fix drude}}}} command to
specify the Drude status of each atom type.

\begin{sphinxadmonition}{note}{Note:}
\sphinxAtStartPar
only the Drude core atoms need to be in the group specified for
this fix. A Drude electron will be transformed together with its cores
even if it is not itself in the group.  It is safe to include Drude
electrons or non\sphinxhyphen{}polarizable atoms in the group. The non\sphinxhyphen{}polarizable
atoms will simply be thermostatted as if they had a massless Drude
partner (electron).
\end{sphinxadmonition}

\begin{sphinxadmonition}{note}{Note:}
\sphinxAtStartPar
Ghost atoms need to know their velocity for this fix to act
correctly.  You must use the {\hyperref[\detokenize{comm_modify::doc}]{\sphinxcrossref{\DUrole{doc}{comm\_modify}}}} command to
enable this, e.g.
\end{sphinxadmonition}

\begin{sphinxVerbatim}[commandchars=\\\{\}]
\PYG{k}{comm\PYGZus{}modify}\PYG{+w}{ }\PYG{n}{vel}\PYG{+w}{ }\PYG{n}{yes}
\end{sphinxVerbatim}


\bigskip\hrule\bigskip


\sphinxAtStartPar
\sphinxstyleemphasis{Tcom} is the target temperature of the centers of mass, which would
be used to thermostat the non\sphinxhyphen{}polarizable atoms.  \sphinxstyleemphasis{Tdrude} is the
(normally low) target temperature of the core\sphinxhyphen{}Drude particle pairs
(dipoles).  \sphinxstyleemphasis{Tcom} and \sphinxstyleemphasis{Tdrude} can be specified as an equal\sphinxhyphen{}style
{\hyperref[\detokenize{variable::doc}]{\sphinxcrossref{\DUrole{doc}{variable}}}}.  If the value is a variable, it should be
specified as v\_name, where name is the variable name. In this case,
the variable will be evaluated each timestep, and its value used to
determine the target temperature.

\sphinxAtStartPar
Equal\sphinxhyphen{}style variables can specify formulas with various mathematical
functions, and include {\hyperref[\detokenize{thermo_style::doc}]{\sphinxcrossref{\DUrole{doc}{thermo\_style}}}} command
keywords for the simulation box parameters and timestep and elapsed
time.  Thus it is easy to specify a time\sphinxhyphen{}dependent temperature.

\sphinxAtStartPar
Like other fixes that perform thermostatting, this fix can be used
with {\hyperref[\detokenize{compute::doc}]{\sphinxcrossref{\DUrole{doc}{compute commands}}}} that remove a “bias” from the
atom velocities.  E.g. to apply the thermostat only to atoms within a
spatial {\hyperref[\detokenize{region::doc}]{\sphinxcrossref{\DUrole{doc}{region}}}}, or to remove the center\sphinxhyphen{}of\sphinxhyphen{}mass
velocity from a group of atoms, or to remove the x\sphinxhyphen{}component of
velocity from the calculation.

\sphinxAtStartPar
This is not done by default, but only if the {\hyperref[\detokenize{fix_modify::doc}]{\sphinxcrossref{\DUrole{doc}{fix\_modify}}}} command is used to assign a temperature compute to this
fix that includes such a bias term.  See the doc pages for individual
{\hyperref[\detokenize{compute::doc}]{\sphinxcrossref{\DUrole{doc}{compute temp commands}}}} to determine which ones include
a bias.  In this case, the thermostat works in the following manner:
bias is removed from each atom, thermostatting is performed on the
remaining thermal degrees of freedom, and the bias is added back in.

\sphinxAtStartPar
Note: The temperature thermostatting the core\sphinxhyphen{}Drude particle pairs
should be chosen low enough, so as to mimic as closely as possible the
self\sphinxhyphen{}consistent minimization. It must however be high enough, so that
the dipoles can follow the local electric field exerted by the
neighboring atoms. The optimal value probably depends on the
temperature of the centers of mass and on the mass of the Drude
particles.

\sphinxAtStartPar
\sphinxstyleemphasis{damp\_com} is the characteristic time for reaching thermal equilibrium
of the centers of mass.  For example, a value of 100.0 means to relax
the temperature of the centers of mass in a timespan of (roughly) 100
time units (tau or fs or ps \sphinxhyphen{} see the {\hyperref[\detokenize{units::doc}]{\sphinxcrossref{\DUrole{doc}{units}}}}
command).  \sphinxstyleemphasis{damp\_drude} is the characteristic time for reaching
thermal equilibrium of the dipoles. It is typically a few timesteps.

\sphinxAtStartPar
The number \sphinxstyleemphasis{seed\_com} and \sphinxstyleemphasis{seed\_drude} are positive integers. They set
the seeds of the Marsaglia random number generators used for
generating the random forces on centers of mass and on the
dipoles. Each processor uses the input seed to generate its own unique
seed and its own stream of random numbers.  Thus the dynamics of the
system will not be identical on two runs on different numbers of
processors.

\sphinxAtStartPar
The keyword \sphinxstyleemphasis{zero} can be used to eliminate drift due to the
thermostat on centers of mass. Because the random forces on different
centers of mass are independent, they do not sum exactly to zero.  As
a result, this fix applies a small random force to the entire system,
and the momentum of the total center of mass of the system undergoes a
slow random walk.  If the keyword \sphinxstyleemphasis{zero} is set to \sphinxstyleemphasis{yes}, the total
random force on the centers of mass is set exactly to zero by
subtracting off an equal part of it from each center of mass in the
group. As a result, the total center of mass of a system with zero
initial momentum will not drift over time.

\sphinxAtStartPar
The actual temperatures of cores and Drude particles, in
center\sphinxhyphen{}of\sphinxhyphen{}mass and relative coordinates, respectively, can be
calculated using the {\hyperref[\detokenize{compute_temp_drude::doc}]{\sphinxcrossref{\DUrole{doc}{compute temp/drude}}}}
command.


\bigskip\hrule\bigskip


\sphinxAtStartPar
Usage example for rigid bodies in the NPT ensemble:

\begin{sphinxVerbatim}[commandchars=\\\{\}]
\PYG{k}{comm\PYGZus{}modify}\PYG{+w}{ }\PYG{n}{vel}\PYG{+w}{ }\PYG{n}{yes}
\PYG{k}{fix }\PYG{n+nv+nvIdentifier}{TEMP}\PYG{+w}{ }\PYG{n+nv+nvIdentifier}{all}\PYG{+w}{ }\PYG{n}{langevin}\PYG{o}{/}\PYG{n}{drude}\PYG{+w}{ }\PYG{l+m}{300}\PYG{n}{.}\PYG{+w}{ }\PYG{l+m}{100}\PYG{n}{.}\PYG{+w}{ }\PYG{l+m}{1256}\PYG{+w}{ }\PYG{l+m}{1}\PYG{n}{.}\PYG{+w}{ }\PYG{l+m}{20}\PYG{n}{.}\PYG{+w}{ }\PYG{l+m}{13977}\PYG{+w}{ }\PYG{n}{zero}\PYG{+w}{ }\PYG{n}{yes}
\PYG{k}{fix }\PYG{n+nv+nvIdentifier}{NPH}\PYG{+w}{ }\PYG{n+nv+nvIdentifier}{ATOMS}\PYG{+w}{ }\PYG{n}{rigid}\PYG{o}{/}\PYG{n}{nph}\PYG{o}{/}\PYG{n}{small}\PYG{+w}{ }\PYG{n}{molecule}\PYG{+w}{ }\PYG{n}{iso}\PYG{+w}{ }\PYG{l+m}{1}\PYG{n}{.}\PYG{+w}{ }\PYG{l+m}{1}\PYG{n}{.}\PYG{+w}{ }\PYG{l+m}{500}\PYG{n}{.}
\PYG{k}{fix }\PYG{n+nv+nvIdentifier}{NVE}\PYG{+w}{ }\PYG{n+nv+nvIdentifier}{DRUDES}\PYG{+w}{ }\PYG{n}{nve}
\PYG{k}{compute }\PYG{n+nv+nvIdentifier}{TDRUDE}\PYG{+w}{ }\PYG{n+nv+nvIdentifier}{all}\PYG{+w}{ }\PYG{n}{temp}\PYG{o}{/}\PYG{n}{drude}
\PYG{k}{thermo\PYGZus{}style}\PYG{+w}{ }\PYG{n}{custom}\PYG{+w}{ }\PYG{n}{step}\PYG{+w}{ }\PYG{n}{cpu}\PYG{+w}{ }\PYG{n}{etotal}\PYG{+w}{ }\PYG{n}{ke}\PYG{+w}{ }\PYG{n}{pe}\PYG{+w}{ }\PYG{n}{ebond}\PYG{+w}{ }\PYG{n}{ecoul}\PYG{+w}{ }\PYG{n}{elong}\PYG{+w}{ }\PYG{n}{press}\PYG{+w}{ }\PYG{n}{vol}\PYG{+w}{ }\PYG{n}{temp}\PYG{+w}{ }\PYG{n}{c\PYGZus{}TDRUDE[1]}\PYG{+w}{ }\PYG{n}{c\PYGZus{}TDRUDE[2]}
\end{sphinxVerbatim}

\sphinxAtStartPar
Comments:
\begin{itemize}
\item {} 
\sphinxAtStartPar
Drude particles should not be in the rigid group, otherwise the Drude
oscillators will be frozen and the system will lose its
polarizability.

\item {} 
\sphinxAtStartPar
\sphinxstyleemphasis{zero yes} avoids a drift of the center of mass of
the system, but is a bit slower.

\item {} 
\sphinxAtStartPar
Use two different random seeds to avoid unphysical correlations.

\item {} 
\sphinxAtStartPar
Temperature is controlled by the fix \sphinxstyleemphasis{langevin/drude}, so the
time\sphinxhyphen{}integration fixes do not thermostat.  Don’t forget to
time\sphinxhyphen{}integrate both cores and Drude particles.

\item {} 
\sphinxAtStartPar
Pressure is time\sphinxhyphen{}integrated only once by using \sphinxstyleemphasis{nve} for Drude
particles and \sphinxstyleemphasis{nph} for atoms/cores (or vice versa). Do not use \sphinxstyleemphasis{nph}
for both.

\item {} 
\sphinxAtStartPar
The temperatures of cores and Drude particles are calculated by
{\hyperref[\detokenize{compute_temp_drude::doc}]{\sphinxcrossref{\DUrole{doc}{compute temp/drude}}}}

\item {} 
\sphinxAtStartPar
Contrary to the alternative thermostatting using Nose\sphinxhyphen{}Hoover thermostat
fix \sphinxstyleemphasis{npt} and {\hyperref[\detokenize{fix_drude_transform::doc}]{\sphinxcrossref{\DUrole{doc}{fix drude/transform}}}}, the
\sphinxstyleemphasis{fix\_modify} command is not required here, because the fix \sphinxstyleemphasis{nph}
computes the global pressure even if its group is \sphinxstyleemphasis{ATOMS}. This is
what we want. If we thermostatted \sphinxstyleemphasis{ATOMS} using \sphinxstyleemphasis{npt}, the pressure
should be the global one, but the temperature should be only that of
the cores. That’s why the command \sphinxstyleemphasis{fix\_modify} should be called in
that case.

\end{itemize}


\bigskip\hrule\bigskip



\subsection{Restart, fix\_modify, output, run start/stop, minimize info}
\label{\detokenize{fix_langevin_drude:restart-fix-modify-output-run-start-stop-minimize-info}}
\sphinxAtStartPar
No information about this fix is written to {\hyperref[\detokenize{restart::doc}]{\sphinxcrossref{\DUrole{doc}{binary restart files}}}}.  Because the state of the random number generator
is not saved in restart files, this means you cannot do “exact”
restarts with this fix, where the simulation continues on the same as
if no restart had taken place.  However, in a statistical sense, a
restarted simulation should produce the same behavior.

\sphinxAtStartPar
The {\hyperref[\detokenize{fix_modify::doc}]{\sphinxcrossref{\DUrole{doc}{fix\_modify}}}} \sphinxstyleemphasis{temp} option is supported by this
fix.  You can use it to assign a temperature {\hyperref[\detokenize{compute::doc}]{\sphinxcrossref{\DUrole{doc}{compute}}}}
you have defined to this fix which will be used in its thermostatting
procedure, as described above. For consistency, the group used by the
compute should include the group of this fix and the Drude particles.

\sphinxAtStartPar
This fix is not invoked during {\hyperref[\detokenize{minimize::doc}]{\sphinxcrossref{\DUrole{doc}{energy minimization}}}}.


\subsection{Restrictions}
\label{\detokenize{fix_langevin_drude:restrictions}}\begin{quote}

\sphinxAtStartPar
none
\end{quote}


\subsection{Related commands}
\label{\detokenize{fix_langevin_drude:related-commands}}
\sphinxAtStartPar
{\hyperref[\detokenize{fix_langevin::doc}]{\sphinxcrossref{\DUrole{doc}{fix langevin}}}},
{\hyperref[\detokenize{fix_drude::doc}]{\sphinxcrossref{\DUrole{doc}{fix drude}}}},
{\hyperref[\detokenize{fix_drude_transform::doc}]{\sphinxcrossref{\DUrole{doc}{fix drude/transform}}}},
{\hyperref[\detokenize{compute_temp_drude::doc}]{\sphinxcrossref{\DUrole{doc}{compute temp/drude}}}},
{\hyperref[\detokenize{pair_thole::doc}]{\sphinxcrossref{\DUrole{doc}{pair\_style thole}}}}


\subsection{Default}
\label{\detokenize{fix_langevin_drude:default}}
\sphinxAtStartPar
The option defaults are zero = no.


\bigskip\hrule\bigskip

\phantomsection\label{\detokenize{fix_langevin_drude:jiang1}}
\sphinxAtStartPar
\sphinxstylestrong{(Jiang1)} Jiang, Hardy, Phillips, MacKerell, Schulten, and Roux, J Phys Chem Lett, 2, 87\sphinxhyphen{}92 (2011).

\sphinxstepscope

\index{fix langevin/eff@\spxentry{fix langevin/eff}}\ignorespaces 

\section{fix langevin/eff command}
\label{\detokenize{fix_langevin_eff:fix-langevin-eff-command}}\label{\detokenize{fix_langevin_eff:index-0}}\label{\detokenize{fix_langevin_eff::doc}}

\subsection{Syntax}
\label{\detokenize{fix_langevin_eff:syntax}}
\begin{sphinxVerbatim}[commandchars=\\\{\}]
\PYG{k}{fix }\PYG{n+nv+nvIdentifier}{ID}\PYG{+w}{ }\PYG{n+nv+nvIdentifier}{group\PYGZhy{}ID}\PYG{+w}{ }\PYG{n}{langevin}\PYG{o}{/}\PYG{n}{eff}\PYG{+w}{ }\PYG{n}{Tstart}\PYG{+w}{ }\PYG{n}{Tstop}\PYG{+w}{ }\PYG{n}{damp}\PYG{+w}{ }\PYG{n}{seed}\PYG{+w}{ }\PYG{n}{keyword}\PYG{+w}{ }\PYG{n}{values}\PYG{+w}{ }\PYG{n}{...}
\end{sphinxVerbatim}
\begin{itemize}
\item {} 
\sphinxAtStartPar
ID, group\sphinxhyphen{}ID are documented in {\hyperref[\detokenize{fix::doc}]{\sphinxcrossref{\DUrole{doc}{fix}}}} command

\item {} 
\sphinxAtStartPar
langevin/eff = style name of this fix command

\item {} 
\sphinxAtStartPar
Tstart,Tstop = desired temperature at start/end of run (temperature units)

\item {} 
\sphinxAtStartPar
damp = damping parameter (time units)

\item {} 
\sphinxAtStartPar
seed = random number seed to use for white noise (positive integer)

\item {} 
\sphinxAtStartPar
zero or more keyword/value pairs may be appended
\begin{sphinxalltt}
keyword = \sphinxstyleemphasis{scale} or \sphinxstyleemphasis{tally} or \sphinxstyleemphasis{zero}
  \sphinxstyleemphasis{scale} values = type ratio
    type = atom type (1\sphinxhyphen{}N)
    ratio = factor by which to scale the damping coefficient
  \sphinxstyleemphasis{tally} values = \sphinxstyleemphasis{no} or \sphinxstyleemphasis{yes}
    \sphinxstyleemphasis{no} = do not tally the energy added/subtracted to atoms
    \sphinxstyleemphasis{yes} = do tally the energy added/subtracted to atoms
\end{sphinxalltt}
\begin{sphinxalltt}
\sphinxstyleemphasis{zero} value = \sphinxstyleemphasis{no} or \sphinxstyleemphasis{yes}
  \sphinxstyleemphasis{no} = do not set total random force to zero
  \sphinxstyleemphasis{yes} = set total random force to zero
\end{sphinxalltt}

\end{itemize}


\subsection{Examples}
\label{\detokenize{fix_langevin_eff:examples}}
\begin{sphinxVerbatim}[commandchars=\\\{\}]
\PYG{k}{fix }\PYG{n+nv+nvIdentifier}{3}\PYG{+w}{ }\PYG{n+nv+nvIdentifier}{boundary}\PYG{+w}{ }\PYG{n}{langevin}\PYG{o}{/}\PYG{n}{eff}\PYG{+w}{ }\PYG{l+m}{1.0}\PYG{+w}{ }\PYG{l+m}{1.0}\PYG{+w}{ }\PYG{l+m}{10.0}\PYG{+w}{ }\PYG{l+m}{699483}
\PYG{k}{fix }\PYG{n+nv+nvIdentifier}{1}\PYG{+w}{ }\PYG{n+nv+nvIdentifier}{all}\PYG{+w}{ }\PYG{n}{langevin}\PYG{o}{/}\PYG{n}{eff}\PYG{+w}{ }\PYG{l+m}{1.0}\PYG{+w}{ }\PYG{l+m}{1.1}\PYG{+w}{ }\PYG{l+m}{10.0}\PYG{+w}{ }\PYG{l+m}{48279}\PYG{+w}{ }\PYG{n}{scale}\PYG{+w}{ }\PYG{l+m}{3}\PYG{+w}{ }\PYG{l+m}{1.5}
\end{sphinxVerbatim}


\subsection{Description}
\label{\detokenize{fix_langevin_eff:description}}
\sphinxAtStartPar
Apply a Langevin thermostat as described in {\hyperref[\detokenize{fix_langevin_eff:schneider2}]{\sphinxcrossref{\DUrole{std,std-ref}{(Schneider)}}}}
to a group of nuclei and electrons in the {\hyperref[\detokenize{pair_eff::doc}]{\sphinxcrossref{\DUrole{doc}{electron force field}}}} model.  Used with {\hyperref[\detokenize{fix_nve_eff::doc}]{\sphinxcrossref{\DUrole{doc}{fix nve/eff}}}},
this command performs Brownian dynamics (BD), since the total force on
each atom will have the form:
\begin{equation*}
\begin{split}F   = & F_c + F_f + F_r \\
F_f = & - \frac{m}{\mathrm{damp}} v \\
F_r \propto &  \sqrt{\frac{k_B T m}{dt~\mathrm{damp}}}\end{split}
\end{equation*}
\sphinxAtStartPar
\(F_c\) is the conservative force computed via the usual
inter\sphinxhyphen{}particle interactions ({\hyperref[\detokenize{pair_style::doc}]{\sphinxcrossref{\DUrole{doc}{pair\_style}}}}).
The \(F_f\) and \(F_r\) terms are added by this fix on a
per\sphinxhyphen{}particle basis.

\sphinxAtStartPar
The operation of this fix is exactly like that described by the
{\hyperref[\detokenize{fix_langevin::doc}]{\sphinxcrossref{\DUrole{doc}{fix langevin}}}} command, except that the
thermostatting is also applied to the radial electron velocity for
electron particles.


\subsection{Restart, fix\_modify, output, run start/stop, minimize info}
\label{\detokenize{fix_langevin_eff:restart-fix-modify-output-run-start-stop-minimize-info}}
\sphinxAtStartPar
No information about this fix is written to {\hyperref[\detokenize{restart::doc}]{\sphinxcrossref{\DUrole{doc}{binary restart files}}}}.  Because the state of the random number generator is not
saved in restart files, this means you cannot do “exact” restarts with
this fix, where the simulation continues on the same as if no restart
had taken place.  However, in a statistical sense, a restarted
simulation should produce the same behavior.

\sphinxAtStartPar
The {\hyperref[\detokenize{fix_modify::doc}]{\sphinxcrossref{\DUrole{doc}{fix\_modify}}}} \sphinxstyleemphasis{temp} option is supported by this
fix.  You can use it to assign a temperature {\hyperref[\detokenize{compute::doc}]{\sphinxcrossref{\DUrole{doc}{compute}}}}
you have defined to this fix which will be used in its thermostatting
procedure, as described above.  For consistency, the group used by
this fix and by the compute should be the same.

\sphinxAtStartPar
The cumulative energy change in the system imposed by this fix is
included in the {\hyperref[\detokenize{thermo_style::doc}]{\sphinxcrossref{\DUrole{doc}{thermodynamic output}}}} keywords
\sphinxstyleemphasis{ecouple} and \sphinxstyleemphasis{econserve}, but only if the \sphinxstyleemphasis{tally} keyword to set to
\sphinxstyleemphasis{yes}.  See the {\hyperref[\detokenize{thermo_style::doc}]{\sphinxcrossref{\DUrole{doc}{thermo\_style}}}} page for
details.

\sphinxAtStartPar
This fix computes a global scalar which can be accessed by various
{\hyperref[\detokenize{Howto_output::doc}]{\sphinxcrossref{\DUrole{doc}{output commands}}}}.  The scalar is the same
cumulative energy change due to this fix described in the previous
paragraph.  The scalar value calculated by this fix is “extensive”.
Note that calculation of this quantity also requires setting the
\sphinxstyleemphasis{tally} keyword to \sphinxstyleemphasis{yes}.

\sphinxAtStartPar
This fix can ramp its target temperature over multiple runs, using the
\sphinxstyleemphasis{start} and \sphinxstyleemphasis{stop} keywords of the {\hyperref[\detokenize{run::doc}]{\sphinxcrossref{\DUrole{doc}{run}}}} command.  See the
{\hyperref[\detokenize{run::doc}]{\sphinxcrossref{\DUrole{doc}{run}}}} command for details of how to do this.

\sphinxAtStartPar
This fix is not invoked during {\hyperref[\detokenize{minimize::doc}]{\sphinxcrossref{\DUrole{doc}{energy minimization}}}}.


\subsection{Restrictions}
\label{\detokenize{fix_langevin_eff:restrictions}}\begin{quote}

\sphinxAtStartPar
none
\end{quote}

\sphinxAtStartPar
This fix is part of the EFF package.  It is only enabled if
LAMMPS was built with that package.  See the {\hyperref[\detokenize{Build_package::doc}]{\sphinxcrossref{\DUrole{doc}{Build package}}}} page for more info.


\subsection{Related commands}
\label{\detokenize{fix_langevin_eff:related-commands}}
\sphinxAtStartPar
{\hyperref[\detokenize{fix_langevin::doc}]{\sphinxcrossref{\DUrole{doc}{fix langevin}}}}


\subsection{Default}
\label{\detokenize{fix_langevin_eff:default}}
\sphinxAtStartPar
The option defaults are scale = 1.0 for all types and tally = no.


\bigskip\hrule\bigskip

\phantomsection\label{\detokenize{fix_langevin_eff:dunweg2}}
\sphinxAtStartPar
\sphinxstylestrong{(Dunweg)} Dunweg and Paul, Int J of Modern Physics C, 2, 817\sphinxhyphen{}27 (1991).

\phantomsection\label{\detokenize{fix_langevin_eff:schneider2}}
\sphinxAtStartPar
\sphinxstylestrong{(Schneider)} Schneider and Stoll, Phys Rev B, 17, 1302 (1978).

\sphinxstepscope

\index{fix langevin/spin@\spxentry{fix langevin/spin}}\ignorespaces 

\section{fix langevin/spin command}
\label{\detokenize{fix_langevin_spin:fix-langevin-spin-command}}\label{\detokenize{fix_langevin_spin:index-0}}\label{\detokenize{fix_langevin_spin::doc}}

\subsection{Syntax}
\label{\detokenize{fix_langevin_spin:syntax}}
\begin{sphinxVerbatim}[commandchars=\\\{\}]
\PYG{k}{fix }\PYG{n+nv+nvIdentifier}{ID}\PYG{+w}{ }\PYG{n+nv+nvIdentifier}{group\PYGZhy{}ID}\PYG{+w}{ }\PYG{n}{langevin}\PYG{o}{/}\PYG{n}{spin}\PYG{+w}{ }\PYG{n}{T}\PYG{+w}{ }\PYG{n}{Tdamp}\PYG{+w}{ }\PYG{n}{seed}
\end{sphinxVerbatim}
\begin{itemize}
\item {} 
\sphinxAtStartPar
ID, group\sphinxhyphen{}ID are documented in {\hyperref[\detokenize{fix::doc}]{\sphinxcrossref{\DUrole{doc}{fix}}}} command

\item {} 
\sphinxAtStartPar
langevin/spin = style name of this fix command

\item {} 
\sphinxAtStartPar
T = desired temperature of the bath (temperature units, K in metal units)

\item {} 
\sphinxAtStartPar
Tdamp = transverse magnetic damping parameter (adim)

\item {} 
\sphinxAtStartPar
seed = random number seed to use for white noise (positive integer)

\end{itemize}


\subsection{Examples}
\label{\detokenize{fix_langevin_spin:examples}}
\begin{sphinxVerbatim}[commandchars=\\\{\}]
\PYG{k}{fix }\PYG{n+nv+nvIdentifier}{2}\PYG{+w}{ }\PYG{n+nv+nvIdentifier}{all}\PYG{+w}{ }\PYG{n}{langevin}\PYG{o}{/}\PYG{n}{spin}\PYG{+w}{ }\PYG{l+m}{300.0}\PYG{+w}{ }\PYG{l+m}{0.01}\PYG{+w}{ }\PYG{l+m}{21}
\end{sphinxVerbatim}


\subsection{Description}
\label{\detokenize{fix_langevin_spin:description}}
\sphinxAtStartPar
Apply a Langevin thermostat as described in {\hyperref[\detokenize{fix_langevin_spin:mayergoyz1}]{\sphinxcrossref{\DUrole{std,std-ref}{(Mayergoyz)}}}} to the
magnetic spins associated to the atoms.
Used with {\hyperref[\detokenize{fix_nve_spin::doc}]{\sphinxcrossref{\DUrole{doc}{fix nve/spin}}}}, this command performs
Brownian dynamics (BD).
A random torque and a transverse dissipation are applied to each spin i according to
the following stochastic differential equation:
\begin{equation*}
\begin{split} \frac{d \vec{s}_{i}}{dt} = \frac{1}{\left(1+\lambda^2 \right)} \left( \left(
 \vec{\omega}_{i} +\vec{\eta} \right) \times \vec{s}_{i} + \lambda\, \vec{s}_{i}
\times\left( \vec{\omega}_{i} \times\vec{s}_{i} \right) \right)\end{split}
\end{equation*}
\sphinxAtStartPar
with \(\lambda\) the transverse damping, and \(\eta\) a random vector.
This equation is referred to as the stochastic Landau\sphinxhyphen{}Lifshitz (sLL)
equation.

\sphinxAtStartPar
The components of \(\eta\) are drawn from a Gaussian probability
law. Their amplitude is defined as a proportion of the temperature of
the external thermostat T (in K in metal units).

\sphinxAtStartPar
More details about this implementation are reported in {\hyperref[\detokenize{fix_langevin_spin:tranchida2}]{\sphinxcrossref{\DUrole{std,std-ref}{(Tranchida)}}}}.

\sphinxAtStartPar
Note: due to the form of the sLL equation, this fix has to be defined just
before the nve/spin fix (and after all other magnetic fixes).
As an example:

\begin{sphinxVerbatim}[commandchars=\\\{\}]
\PYG{k}{fix }\PYG{n+nv+nvIdentifier}{1}\PYG{+w}{ }\PYG{n+nv+nvIdentifier}{all}\PYG{+w}{ }\PYG{n}{precession}\PYG{o}{/}\PYG{n}{spin}\PYG{+w}{ }\PYG{n}{zeeman}\PYG{+w}{ }\PYG{l+m}{0.01}\PYG{+w}{ }\PYG{l+m}{0.0}\PYG{+w}{ }\PYG{l+m}{0.0}\PYG{+w}{ }\PYG{l+m}{1.0}
\PYG{k}{fix }\PYG{n+nv+nvIdentifier}{2}\PYG{+w}{ }\PYG{n+nv+nvIdentifier}{all}\PYG{+w}{ }\PYG{n}{langevin}\PYG{o}{/}\PYG{n}{spin}\PYG{+w}{ }\PYG{l+m}{300.0}\PYG{+w}{ }\PYG{l+m}{0.01}\PYG{+w}{ }\PYG{l+m}{21}
\PYG{k}{fix }\PYG{n+nv+nvIdentifier}{3}\PYG{+w}{ }\PYG{n+nv+nvIdentifier}{all}\PYG{+w}{ }\PYG{n}{nve}\PYG{o}{/}\PYG{n}{spin}\PYG{+w}{ }\PYG{n}{lattice}\PYG{+w}{ }\PYG{n}{moving}
\end{sphinxVerbatim}

\sphinxAtStartPar
is correct, but defining a force/spin command after the langevin/spin command
would give an error message.

\sphinxAtStartPar
Note: The random \# \sphinxstyleemphasis{seed} must be a positive integer.  A Marsaglia random
number generator is used.  Each processor uses the input seed to
generate its own unique seed and its own stream of random numbers.
Thus the dynamics of the system will not be identical on two runs on
different numbers of processors.


\bigskip\hrule\bigskip



\subsection{Restart, fix\_modify, output, run start/stop, minimize info}
\label{\detokenize{fix_langevin_spin:restart-fix-modify-output-run-start-stop-minimize-info}}
\sphinxAtStartPar
No information about this fix is written to {\hyperref[\detokenize{restart::doc}]{\sphinxcrossref{\DUrole{doc}{binary restart files}}}}.  Because the state of the random number generator
is not saved in restart files, this means you cannot do “exact”
restarts with this fix, where the simulation continues on the same as
if no restart had taken place.  However, in a statistical sense, a
restarted simulation should produce the same behavior.

\sphinxAtStartPar
This fix is not invoked during {\hyperref[\detokenize{minimize::doc}]{\sphinxcrossref{\DUrole{doc}{energy minimization}}}}.


\subsection{Restrictions}
\label{\detokenize{fix_langevin_spin:restrictions}}
\sphinxAtStartPar
The \sphinxstyleemphasis{langevin/spin} fix is part of the SPIN package.  This style is
only enabled if LAMMPS was built with this package.  See the {\hyperref[\detokenize{Build_package::doc}]{\sphinxcrossref{\DUrole{doc}{Build package}}}} page for more info.

\sphinxAtStartPar
The numerical integration has to be performed with \sphinxstyleemphasis{fix nve/spin}
when \sphinxstyleemphasis{fix langevin/spin} is enabled.

\sphinxAtStartPar
This fix has to be the last defined magnetic fix before the time
integration fix (e.g. \sphinxstyleemphasis{fix nve/spin}).


\subsection{Related commands}
\label{\detokenize{fix_langevin_spin:related-commands}}
\sphinxAtStartPar
{\hyperref[\detokenize{fix_nve_spin::doc}]{\sphinxcrossref{\DUrole{doc}{fix nve/spin}}}}, {\hyperref[\detokenize{fix_precession_spin::doc}]{\sphinxcrossref{\DUrole{doc}{fix precession/spin}}}}


\subsection{Default}
\label{\detokenize{fix_langevin_spin:default}}
\sphinxAtStartPar
none


\bigskip\hrule\bigskip

\phantomsection\label{\detokenize{fix_langevin_spin:mayergoyz1}}
\sphinxAtStartPar
\sphinxstylestrong{(Mayergoyz)} I.D. Mayergoyz, G. Bertotti, C. Serpico (2009). Elsevier (2009)

\phantomsection\label{\detokenize{fix_langevin_spin:tranchida2}}
\sphinxAtStartPar
\sphinxstylestrong{(Tranchida)} Tranchida, Plimpton, Thibaudeau and Thompson,
Journal of Computational Physics, 372, 406\sphinxhyphen{}425, (2018).

\sphinxstepscope

\index{fix lb/fluid@\spxentry{fix lb/fluid}}\ignorespaces 

\section{fix lb/fluid command}
\label{\detokenize{fix_lb_fluid:fix-lb-fluid-command}}\label{\detokenize{fix_lb_fluid:index-0}}\label{\detokenize{fix_lb_fluid::doc}}

\subsection{Syntax}
\label{\detokenize{fix_lb_fluid:syntax}}
\begin{sphinxVerbatim}[commandchars=\\\{\}]
\PYG{k}{fix }\PYG{n+nv+nvIdentifier}{ID}\PYG{+w}{ }\PYG{n+nv+nvIdentifier}{group\PYGZhy{}ID}\PYG{+w}{ }\PYG{n}{lb}\PYG{o}{/}\PYG{n}{fluid}\PYG{+w}{ }\PYG{n}{nevery}\PYG{+w}{ }\PYG{n}{viscosity}\PYG{+w}{ }\PYG{n}{density}\PYG{+w}{ }\PYG{n}{keyword}\PYG{+w}{ }\PYG{n}{values}\PYG{+w}{ }\PYG{n}{...}
\end{sphinxVerbatim}
\begin{itemize}
\item {} 
\sphinxAtStartPar
ID, group\sphinxhyphen{}ID are documented in {\hyperref[\detokenize{fix::doc}]{\sphinxcrossref{\DUrole{doc}{fix}}}} command

\item {} 
\sphinxAtStartPar
lb/fluid = style name of this fix command

\item {} 
\sphinxAtStartPar
nevery = update the lattice\sphinxhyphen{}Boltzmann fluid every this many timesteps (should normally be 1)

\item {} 
\sphinxAtStartPar
viscosity = the fluid viscosity (units of mass/(time*length)).

\item {} 
\sphinxAtStartPar
density = the fluid density.

\item {} 
\sphinxAtStartPar
zero or more keyword/value pairs may be appended

\item {} 
\sphinxAtStartPar
keyword = \sphinxstyleemphasis{dx} or \sphinxstyleemphasis{dm} or \sphinxstyleemphasis{noise} or \sphinxstyleemphasis{stencil} or \sphinxstyleemphasis{read\_restart} or \sphinxstyleemphasis{write\_restart} or \sphinxstyleemphasis{zwall\_velocity} or \sphinxstyleemphasis{pressurebcx} or \sphinxstyleemphasis{bodyforce} or \sphinxstyleemphasis{D3Q19} or \sphinxstyleemphasis{dumpxdmf} or \sphinxstyleemphasis{linearInit} or \sphinxstyleemphasis{dof} or \sphinxstyleemphasis{scaleGamma} or \sphinxstyleemphasis{a0} or \sphinxstyleemphasis{npits} or \sphinxstyleemphasis{wp} or \sphinxstyleemphasis{sw}
\begin{sphinxalltt}
\sphinxstyleemphasis{dx} values = dx\_LB = the lattice spacing.
\sphinxstyleemphasis{dm} values = dm\_LB = the lattice\sphinxhyphen{}Boltzmann mass unit.
\sphinxstyleemphasis{noise} values = Temperature seed
    Temperature = fluid temperature.
    seed = random number generator seed (positive integer)
\sphinxstyleemphasis{stencil} values = 2 (trilinear stencil, the default), 3 (3\sphinxhyphen{}point immersed boundary stencil), or 4 (4\sphinxhyphen{}point Keys\textquotesingle{} interpolation stencil)
\sphinxstyleemphasis{read\_restart} values = restart file = name of the restart file to use to restart a fluid run.
\sphinxstyleemphasis{write\_restart} values = N = write a restart file every N MD timesteps.
\sphinxstyleemphasis{zwall\_velocity} values = velocity\_bottom velocity\_top = velocities along the y\sphinxhyphen{}direction of the bottom and top walls (located at z=zmin and z=zmax).
\sphinxstyleemphasis{pressurebcx} values = pgradav = imposes a pressure jump at the (periodic) x\sphinxhyphen{}boundary of pgradav*Lx*1000.
\sphinxstyleemphasis{bodyforce} values = bodyforcex bodyforcey bodyforcez = the x,y and z components of a constant body force added to the fluid.
\sphinxstyleemphasis{D3Q19} values = none (used to switch from the default D3Q15, 15 velocity lattice, to the D3Q19, 19 velocity lattice).
\sphinxstyleemphasis{dumpxdmf} values = N file timeI
    N = output the force and torque every N timesteps
    file = output file name
    timeI = 1 (use simulation time to index xdmf file), 0 (use output frame number to index xdmf file)
\sphinxstyleemphasis{linearInit} values = none = initialize density and velocity using linear interpolation (default is uniform density, no velocities)
\sphinxstyleemphasis{dof} values = dof = specify the number of degrees of freedom for temperature calculation
\sphinxstyleemphasis{scaleGamma} values = type gammaFactor
    type = atom type (1\sphinxhyphen{}N)
    gammaFactor = factor to scale the \sphinxstyleemphasis{setGamma} gamma value by, for the specified atom type.
\sphinxstyleemphasis{a0} values = a\_0\_real = the square of the speed of sound in the fluid.
\sphinxstyleemphasis{npits} values = npits h\_p l\_p l\_pp l\_e
    npits = number of pit regions
    h\_p = z\sphinxhyphen{}height of pit regions (floor to bottom of slit)
    l\_p = x\sphinxhyphen{}length of pit regions
    l\_pp = x\sphinxhyphen{}length of slit regions between consecutive pits
    l\_e = x\sphinxhyphen{}length of slit regions at ends
\sphinxstyleemphasis{wp} values = w\_p = y\sphinxhyphen{}width of slit regions (defaults to full width if not present or if sw active)
\sphinxstyleemphasis{sw} values = none  (turns on y\sphinxhyphen{}sidewalls (in xz plane) if npits option active)
\end{sphinxalltt}

\end{itemize}


\subsection{Examples}
\label{\detokenize{fix_lb_fluid:examples}}
\begin{sphinxVerbatim}[commandchars=\\\{\}]
\PYG{k}{fix }\PYG{n+nv+nvIdentifier}{1}\PYG{+w}{ }\PYG{n+nv+nvIdentifier}{all}\PYG{+w}{ }\PYG{n}{lb}\PYG{o}{/}\PYG{n}{fluid}\PYG{+w}{ }\PYG{l+m}{1}\PYG{+w}{ }\PYG{l+m}{1.0}\PYG{+w}{ }\PYG{l+m}{0.0009982071}\PYG{+w}{ }\PYG{n}{dx}\PYG{+w}{ }\PYG{l+m}{1.2}\PYG{+w}{ }\PYG{n}{dm}\PYG{+w}{ }\PYG{l+m}{0.001}
\PYG{k}{fix }\PYG{n+nv+nvIdentifier}{1}\PYG{+w}{ }\PYG{n+nv+nvIdentifier}{all}\PYG{+w}{ }\PYG{n}{lb}\PYG{o}{/}\PYG{n}{fluid}\PYG{+w}{ }\PYG{l+m}{1}\PYG{+w}{ }\PYG{l+m}{1.0}\PYG{+w}{ }\PYG{l+m}{0.0009982071}\PYG{+w}{ }\PYG{n}{dx}\PYG{+w}{ }\PYG{l+m}{1.2}\PYG{+w}{ }\PYG{n}{dm}\PYG{+w}{ }\PYG{l+m}{0.001}\PYG{+w}{ }\PYG{n}{noise}\PYG{+w}{ }\PYG{l+m}{300.0}\PYG{+w}{ }\PYG{l+m}{2761}
\PYG{k}{fix }\PYG{n+nv+nvIdentifier}{1}\PYG{+w}{ }\PYG{n+nv+nvIdentifier}{all}\PYG{+w}{ }\PYG{n}{lb}\PYG{o}{/}\PYG{n}{fluid}\PYG{+w}{ }\PYG{l+m}{1}\PYG{+w}{ }\PYG{l+m}{1.0}\PYG{+w}{ }\PYG{l+m}{1.0}\PYG{+w}{ }\PYG{n}{dx}\PYG{+w}{ }\PYG{l+m}{4.0}\PYG{+w}{ }\PYG{n}{dm}\PYG{+w}{ }\PYG{l+m}{10.0}\PYG{+w}{ }\PYG{n}{dumpxdmf}\PYG{+w}{ }\PYG{l+m}{500}\PYG{+w}{ }\PYG{n}{fflow}\PYG{+w}{ }\PYG{l+m}{0}\PYG{+w}{ }\PYG{n}{pressurebcx}\PYG{+w}{ }\PYG{l+m}{0.01}\PYG{+w}{ }\PYG{n}{npits}\PYG{+w}{ }\PYG{l+m}{2}\PYG{+w}{ }\PYG{l+m}{20}\PYG{+w}{ }\PYG{l+m}{40}\PYG{+w}{ }\PYG{l+m}{5}\PYG{+w}{ }\PYG{l+m}{0}\PYG{+w}{ }\PYG{n}{wp}\PYG{+w}{ }\PYG{l+m}{30}
\end{sphinxVerbatim}


\subsection{Description}
\label{\detokenize{fix_lb_fluid:description}}
\sphinxAtStartPar
\DUrole{versionmodified,changed}{Changed in version 24Mar2022.}

\sphinxAtStartPar
Implement a lattice\sphinxhyphen{}Boltzmann fluid on a uniform mesh covering the
LAMMPS simulation domain.  Note that this fix was updated in 2022 and is
not backward compatible with the previous version.  If you need the
previous version, please download an older version of LAMMPS.  The MD
particles described by \sphinxstyleemphasis{group\sphinxhyphen{}ID} apply a velocity dependent force to
the fluid.

\sphinxAtStartPar
The lattice\sphinxhyphen{}Boltzmann algorithm solves for the fluid motion governed by
the Navier Stokes equations,
\begin{equation*}
\begin{split}\partial_t \rho + \partial_{\beta}\left(\rho u_{\beta}\right)= & 0 \\
\partial_t\left(\rho u_{\alpha}\right) + \partial_{\beta}\left(\rho u_{\alpha} u_{\beta}\right) = & \partial_{\beta}\sigma_{\alpha \beta} + F_{\alpha} + \partial_{\beta}\left(\eta_{\alpha \beta \gamma \nu}\partial_{\gamma} u_{\nu}\right)\end{split}
\end{equation*}
\sphinxAtStartPar
with,
\begin{equation*}
\begin{split}\eta_{\alpha \beta \gamma \nu} = \eta\left[\delta_{\alpha \gamma}\delta_{\beta \nu} + \delta_{\alpha \nu}\delta_{\beta \gamma} - \frac{2}{3}\delta_{\alpha \beta}\delta_{\gamma \nu}\right] + \Lambda \delta_{\alpha \beta}\delta_{\gamma \nu}\end{split}
\end{equation*}
\sphinxAtStartPar
where \(\rho\) is the fluid density, \sphinxstyleemphasis{u} is the local
fluid velocity, \(\sigma\) is the stress tensor, \sphinxstyleemphasis{F} is a local external
force, and \(\eta\) and \(\Lambda\) are the shear and bulk viscosities
respectively.  Here, we have implemented
\begin{equation*}
\begin{split}\sigma_{\alpha \beta} = -P_{\alpha \beta} = -\rho a_0 \delta_{\alpha \beta}\end{split}
\end{equation*}
\sphinxAtStartPar
with \(a_0\) set to \(\frac{1}{3} \frac{dx}{dt}^2\) by default.
You should not normally need to change this default.

\sphinxAtStartPar
The algorithm involves tracking the time evolution of a set of partial
distribution functions which evolve according to a velocity discretized
version of the Boltzmann equation,
\begin{equation*}
\begin{split}\left(\partial_t + e_{i\alpha}\partial_{\alpha}\right)f_i = -\frac{1}{\tau}\left(f_i - f_i^{eq}\right) + W_i\end{split}
\end{equation*}
\sphinxAtStartPar
where the first term on the right hand side represents a single time
relaxation towards the equilibrium distribution function, and
\(\tau\) is a parameter physically related to the viscosity.  On a
technical note, we have implemented a 15 velocity model (D3Q15) as
default; however, the user can switch to a 19 velocity model (D3Q19)
through the use of the \sphinxstyleemphasis{D3Q19} keyword.  Physical variables are then
defined in terms of moments of the distribution functions,
\begin{equation*}
\begin{split}\rho = & \displaystyle\sum\limits_{i} f_i \\
\rho u_{\alpha} = & \displaystyle\sum\limits_{i} f_i e_{i\alpha}\end{split}
\end{equation*}
\sphinxAtStartPar
Full details of the lattice\sphinxhyphen{}Boltzmann algorithm used can be found in
{\hyperref[\detokenize{fix_lb_fluid:fluid-denniston}]{\sphinxcrossref{\DUrole{std,std-ref}{Denniston et al.}}}}.

\sphinxAtStartPar
The fluid is coupled to the MD particles described by \sphinxstyleemphasis{group\sphinxhyphen{}ID} through
a velocity dependent force.  The contribution to the fluid force on a
given lattice mesh site j due to MD particle \(\alpha\) is
calculated as:
\begin{equation*}
\begin{split}{\bf F}_{j \alpha} = \gamma \left({\bf v}_n - {\bf u}_f \right) \zeta_{j\alpha}\end{split}
\end{equation*}
\sphinxAtStartPar
where \(\mathbf{v}_n\) is the velocity of the MD particle,
\(\mathbf{u}_f\) is the fluid velocity interpolated to the particle
location, and \(\gamma\) is the force coupling constant.  This
force, as with most forces in LAMMPS, and hence the velocities, are
calculated at the half\sphinxhyphen{}time step. \(\zeta\) is a weight assigned to
the grid point, obtained by distributing the particle to the nearest
lattice sites.

\sphinxAtStartPar
The force coupling constant, \(\gamma\), is calculated
according to
\begin{equation*}
\begin{split}\gamma = \frac{2m_um_v}{m_u+m_v}\left(\frac{1}{\Delta t}\right)\end{split}
\end{equation*}
\sphinxAtStartPar
Here, \(m_v\) is the mass of the MD particle, \(m_u\) is a
representative fluid mass at the particle location, and \(\Delta t\)
is the time step.  The fluid mass \(m_u\) that the MD particle
interacts with is calculated internally.  This coupling is chosen to
constrain the particle and associated fluid velocity to match at the end
of the time step.  As with other constraints, such as {\hyperref[\detokenize{fix_shake::doc}]{\sphinxcrossref{\DUrole{doc}{shake}}}}, this constraint can remove degrees of freedom from the
simulation which are accounted for internally in the algorithm.

\begin{sphinxadmonition}{note}{Note:}
\sphinxAtStartPar
While this fix applies the force of the particles on the fluid, it
does not apply the force of the fluid to the particles.  There is
only one option to include this hydrodynamic force on the particles,
and that is through the use of the {\hyperref[\detokenize{fix_lb_viscous::doc}]{\sphinxcrossref{\DUrole{doc}{lb/viscous}}}}
fix.  This fix adds the hydrodynamic force to the total force acting
on the particles, after which any of the built\sphinxhyphen{}in LAMMPS integrators
can be used to integrate the particle motion.  If the
{\hyperref[\detokenize{fix_lb_viscous::doc}]{\sphinxcrossref{\DUrole{doc}{lb/viscous}}}} fix is NOT used to add the
hydrodynamic force to the total force acting on the particles, this
physically corresponds to a situation in which an infinitely massive
particle is moving through the fluid (since collisions between the
particle and the fluid do not act to change the particle’s velocity).
In this case, setting \sphinxstyleemphasis{scaleGamma} to \sphinxhyphen{}1 for the corresponding
particle type will explicitly take this limit (of infinite particle
mass) in computing the force coupling for the fluid force.
\end{sphinxadmonition}


\bigskip\hrule\bigskip


\sphinxAtStartPar
Physical parameters describing the fluid are specified through
\sphinxstyleemphasis{viscosity} and \sphinxstyleemphasis{density}.  These parameters should all be given in
terms of the mass, distance, and time units chosen for the main LAMMPS
run, as they are scaled by the LB timestep, lattice spacing, and mass
unit, inside the fix.

\sphinxAtStartPar
The \sphinxstyleemphasis{dx} keyword allows the user to specify a value for the LB grid
spacing and the \sphinxstyleemphasis{dm} keyword allows the user to specify the LB mass
unit.  Inside the fix, parameters are scaled by the lattice\sphinxhyphen{}Boltzmann
timestep, \(dt_{LB}\), grid spacing, \(dx_{LB}\), and mass unit,
\(dm_{LB}\).  \(dt_{LB}\) is set equal to
\(\mathrm{nevery}\cdot dt_{MD}\), where \(dt_{MD}\) is the MD
timestep.  By default, \(dm_{LB}\) is set equal to 1.0, and
\(dx_{LB}\) is chosen so that \(\frac{\tau}{dt} = \frac{3\eta
dt}{\rho dx^2}\) is approximately equal to 1.
\begin{quote}

\begin{sphinxadmonition}{note}{Note:}
\sphinxAtStartPar
Care must be taken when choosing both a value for \(dx_{LB}\),
and a simulation domain size.  This fix uses the same subdivision of
the simulation domain among processors as the main LAMMPS program.  In
order to uniformly cover the simulation domain with lattice sites, the
lengths of the individual LAMMPS subdomains must all be evenly
divisible by \(dx_{LB}\).  If the simulation domain size is cubic,
with equal lengths in all dimensions, and the default value for
\(dx_{LB}\) is used, this will automatically be satisfied.
\end{sphinxadmonition}
\end{quote}

\sphinxAtStartPar
If the \sphinxstyleemphasis{noise} keyword is used, followed by a positive temperature
value, and a positive integer random number seed, the thermal LB algorithm
of {\hyperref[\detokenize{fix_lb_fluid:adhikari}]{\sphinxcrossref{\DUrole{std,std-ref}{Adhikari et al.}}}} is used.

\sphinxAtStartPar
If the keyword \sphinxstyleemphasis{stencil} is used, the value sets the number of
interpolation points used in each direction.  For this, the user has the
choice between a trilinear stencil (\sphinxstyleemphasis{stencil} 2), which provides a
support of 8 lattice sites, or the 3\sphinxhyphen{}point immersed boundary method
stencil (\sphinxstyleemphasis{stencil} 3), which provides a support of 27 lattice sites, or
the 4\sphinxhyphen{}point Keys’ interpolation stencil (stencil 4), which provides a
support of 64 lattice sites.  The trilinear stencil is the default as it
is better suited for simulation of objects close to walls or other
objects, due to its smaller support.  The 3\sphinxhyphen{}point stencil provides
smoother motion of the lattice and is suitable for particles not likely
to be to close to walls or other objects.

\sphinxAtStartPar
If the keyword \sphinxstyleemphasis{write\_restart} is used, followed by a positive integer,
N, a binary restart file is printed every N LB timesteps.  This restart
file only contains information about the fluid.  Therefore, a LAMMPS
restart file should also be written in order to print out full details
of the simulation.

\begin{sphinxadmonition}{note}{Note:}
\sphinxAtStartPar
When a large number of lattice grid points are used, the restart
files may become quite large.
\end{sphinxadmonition}

\sphinxAtStartPar
In order to restart the fluid portion of the simulation, the keyword
\sphinxstyleemphasis{read\_restart} is specified, followed by the name of the binary
lb\_fluid restart file to be used.

\sphinxAtStartPar
If the \sphinxstyleemphasis{zwall\_velocity} keyword is used y\sphinxhyphen{}velocities are assigned to
the lower and upper walls.  This keyword requires the presence of
walls in the z\sphinxhyphen{}direction.  This is set by assigning fixed boundary
conditions in the z\sphinxhyphen{}direction.  If fixed boundary conditions are
present in the z\sphinxhyphen{}direction, and this keyword is not used, the walls
are assumed to be stationary.

\sphinxAtStartPar
If the \sphinxstyleemphasis{pressurebcx} keyword is used, a pressure jump (implemented by a
step jump in density) is imposed at the (periodic) x\sphinxhyphen{}boundary.  The
value set specifies what would be the resulting equilibrium average
pressure gradient in the x\sphinxhyphen{}direction if the system had a constant
cross\sphinxhyphen{}section (i.e. resistance to flow).  It is converted to a pressure
jump by multiplication by the system size in the x\sphinxhyphen{}direction.  As this
value should normally be quite small, it is also assumed to be scaled
by 1000.

\sphinxAtStartPar
If the \sphinxstyleemphasis{bodyforce} keyword is used, a constant body force is added to
the fluid, defined by it’s x, y and z components.

\sphinxAtStartPar
If the keyword \sphinxstyleemphasis{D3Q19} is used, the 19 velocity (D3Q19) lattice is
used by the lattice\sphinxhyphen{}Boltzmann algorithm.  By default, the 15 velocity
(D3Q15) lattice is used.

\sphinxAtStartPar
If the \sphinxstyleemphasis{dumpxdmf} keyword is used, followed by a positive integer, N,
and a file name, the fluid densities and velocities at each lattice site
are output to an xdmf file every N timesteps.  This is a binary file
format that can be read by visualization packages such as \sphinxhref{https://www.paraview.org/}{Paraview} .  The xdmf file format contains a time
index for each frame dump and the value timeI = 1 uses simulation time
while 0 uses the output frame number to index xdmf file.  The later can
be useful if the {\hyperref[\detokenize{dump_vtk::doc}]{\sphinxcrossref{\DUrole{doc}{dump vtk}}}} command is used to output
the particle positions at the same timesteps and you want to visualize
both the fluid and particle data together in \sphinxhref{https://www.paraview.org/}{Paraview} .

\sphinxAtStartPar
The \sphinxstyleemphasis{scaleGamma} keyword allows the user to scale the \(\gamma\)
value by a factor, gammaFactor, for a given atom type.  Setting
\sphinxstyleemphasis{scaleGamma} to \sphinxhyphen{}1 for the corresponding particle type will explicitly
take the limit of infinite particle mass in computing the force coupling
for the fluid force (see note above).

\sphinxAtStartPar
If the \sphinxstyleemphasis{a0} keyword is used, the value specified is used for the square
of the speed of sound in the fluid.  If this keyword is not present, the
speed of sound squared is set equal to
\(\frac{1}{3}\left(\frac{dx_{LB}}{dt_{LB}}\right)^2\).  Setting
\(a0 > (\frac{dx_{LB}}{dt_{LB}})^2\) is not allowed, as this may
lead to instabilities.  As the speed of sound should usually be much
larger than any fluid velocity of interest, its value does not normally
have a significant impact on the results.  As such, it is usually best
to use the default for this option.

\sphinxAtStartPar
The \sphinxstyleemphasis{npits} keyword (followed by integer arguments: npits, h\_p, l\_p,
l\_pp, l\_e) sets the fluid domain to the pits geometry.  These arguments
should only be used if you actually want something more complex than a
rectangular/cubic geometry.  The npits value sets the number of pits
regions (arranged along x).  The remaining arguments are sizes measured
in multiples of dx\_lb: h\_p is the z\sphinxhyphen{}height of the pit regions, l\_p is
the x\sphinxhyphen{}length of the pit regions, l\_pp is the length of the region
between consecutive pits (referred to as a “slit” region), and l\_e is
the x\sphinxhyphen{}length of the slit regions at each end of the channel.  The pit
geometry must fill the system in the x\sphinxhyphen{}direction but can be longer, in
which case it is truncated (which enables asymmetric entrance/exit end
sections).  The additional \sphinxstyleemphasis{wp} keyword allows the width (in
y\sphinxhyphen{}direction) of the pit to be specified (the default is full width) and
the \sphinxstyleemphasis{sw} keyword indicates that there should be sidewalls in the
y\sphinxhyphen{}direction (default is periodic in y\sphinxhyphen{}direction).  These parameters are
illustrated below:

\begin{sphinxVerbatim}[commandchars=\\\{\}]
Sideview (in xz plane) of pit geometry:
\PYGZus{}\PYGZus{}\PYGZus{}\PYGZus{}\PYGZus{}\PYGZus{}\PYGZus{}\PYGZus{}\PYGZus{}\PYGZus{}\PYGZus{}\PYGZus{}\PYGZus{}\PYGZus{}\PYGZus{}\PYGZus{}\PYGZus{}\PYGZus{}\PYGZus{}\PYGZus{}\PYGZus{}\PYGZus{}\PYGZus{}\PYGZus{}\PYGZus{}\PYGZus{}\PYGZus{}\PYGZus{}\PYGZus{}\PYGZus{}\PYGZus{}\PYGZus{}\PYGZus{}\PYGZus{}\PYGZus{}\PYGZus{}\PYGZus{}\PYGZus{}\PYGZus{}\PYGZus{}\PYGZus{}\PYGZus{}\PYGZus{}\PYGZus{}\PYGZus{}\PYGZus{}\PYGZus{}\PYGZus{}\PYGZus{}\PYGZus{}\PYGZus{}\PYGZus{}\PYGZus{}\PYGZus{}\PYGZus{}\PYGZus{}\PYGZus{}\PYGZus{}\PYGZus{}\PYGZus{}\PYGZus{}\PYGZus{}\PYGZus{}\PYGZus{}\PYGZus{}\PYGZus{}\PYGZus{}\PYGZus{}\PYGZus{}\PYGZus{}
  slit                          slit                          slit     \PYGZca{}
                                                                       |
\PYGZlt{}\PYGZhy{}\PYGZhy{}\PYGZhy{}le\PYGZhy{}\PYGZhy{}\PYGZhy{}\PYGZgt{}\PYGZlt{}\PYGZhy{}\PYGZhy{}\PYGZhy{}\PYGZhy{}\PYGZhy{}\PYGZhy{}\PYGZhy{}\PYGZhy{}\PYGZhy{}lp\PYGZhy{}\PYGZhy{}\PYGZhy{}\PYGZhy{}\PYGZhy{}\PYGZhy{}\PYGZhy{}\PYGZgt{}\PYGZlt{}\PYGZhy{}\PYGZhy{}\PYGZhy{}lpp\PYGZhy{}\PYGZhy{}\PYGZhy{}\PYGZgt{}\PYGZlt{}\PYGZhy{}\PYGZhy{}\PYGZhy{}\PYGZhy{}\PYGZhy{}\PYGZhy{}\PYGZhy{}lp\PYGZhy{}\PYGZhy{}\PYGZhy{}\PYGZhy{}\PYGZhy{}\PYGZhy{}\PYGZhy{}\PYGZhy{}\PYGZgt{}\PYGZlt{}\PYGZhy{}\PYGZhy{}\PYGZhy{}le\PYGZhy{}\PYGZhy{}\PYGZhy{}\PYGZgt{} hs = (Nbz\PYGZhy{}1) \PYGZhy{} hp
                                                                       |
\PYGZus{}\PYGZus{}\PYGZus{}\PYGZus{}\PYGZus{}\PYGZus{}\PYGZus{}\PYGZus{}\PYGZus{}\PYGZus{}                    \PYGZus{}\PYGZus{}\PYGZus{}\PYGZus{}\PYGZus{}\PYGZus{}\PYGZus{}\PYGZus{}\PYGZus{}\PYGZus{}                    \PYGZus{}\PYGZus{}\PYGZus{}\PYGZus{}\PYGZus{}\PYGZus{}\PYGZus{}\PYGZus{}\PYGZus{}\PYGZus{} v
          |                  |          |                  |           \PYGZca{}       z
          |                  |          |                  |           |       |
          |       pit        |          |       pit        |           hp      +\PYGZhy{}x
          |                  |          |                  |           |
          |\PYGZus{}\PYGZus{}\PYGZus{}\PYGZus{}\PYGZus{}\PYGZus{}\PYGZus{}\PYGZus{}\PYGZus{}\PYGZus{}\PYGZus{}\PYGZus{}\PYGZus{}\PYGZus{}\PYGZus{}\PYGZus{}\PYGZus{}\PYGZus{}|          |\PYGZus{}\PYGZus{}\PYGZus{}\PYGZus{}\PYGZus{}\PYGZus{}\PYGZus{}\PYGZus{}\PYGZus{}\PYGZus{}\PYGZus{}\PYGZus{}\PYGZus{}\PYGZus{}\PYGZus{}\PYGZus{}\PYGZus{}\PYGZus{}|           v

Endview (in yz plane) of pit geometry (no sw so wp is active):
\PYGZus{}\PYGZus{}\PYGZus{}\PYGZus{}\PYGZus{}\PYGZus{}\PYGZus{}\PYGZus{}\PYGZus{}\PYGZus{}\PYGZus{}\PYGZus{}\PYGZus{}\PYGZus{}\PYGZus{}\PYGZus{}\PYGZus{}\PYGZus{}\PYGZus{}\PYGZus{}\PYGZus{}
                      \PYGZca{}
                      |
                      hs
                      |
\PYGZus{}\PYGZus{}\PYGZus{}\PYGZus{}\PYGZus{}\PYGZus{}\PYGZus{}\PYGZus{}\PYGZus{}\PYGZus{}\PYGZus{}\PYGZus{}\PYGZus{}\PYGZus{}\PYGZus{}\PYGZus{}\PYGZus{}\PYGZus{}\PYGZus{}\PYGZus{}\PYGZus{} v
    |          |      \PYGZca{}
    |          |      |          z
    |\PYGZlt{}\PYGZhy{}\PYGZhy{}\PYGZhy{}wp\PYGZhy{}\PYGZhy{}\PYGZhy{}\PYGZgt{}|      hp         |
    |          |      |          +\PYGZhy{}y
    |\PYGZus{}\PYGZus{}\PYGZus{}\PYGZus{}\PYGZus{}\PYGZus{}\PYGZus{}\PYGZus{}\PYGZus{}\PYGZus{}|      v
\end{sphinxVerbatim}


\bigskip\hrule\bigskip


\sphinxAtStartPar
For further details, as well as descriptions and results of several test
runs, see {\hyperref[\detokenize{fix_lb_fluid:fluid-denniston}]{\sphinxcrossref{\DUrole{std,std-ref}{Denniston et al.}}}}.  Please include a
citation to this paper if the lb\_fluid fix is used in work contributing
to published research.


\bigskip\hrule\bigskip



\subsection{Restart, fix\_modify, output, run start/stop, minimize info}
\label{\detokenize{fix_lb_fluid:restart-fix-modify-output-run-start-stop-minimize-info}}
\sphinxAtStartPar
Due to the large size of the fluid data, this fix writes it’s own
binary restart files, if requested, independent of the main LAMMPS
{\hyperref[\detokenize{restart::doc}]{\sphinxcrossref{\DUrole{doc}{binary restart files}}}}; no information about \sphinxstyleemphasis{lb\_fluid}
is written to the main LAMMPS {\hyperref[\detokenize{restart::doc}]{\sphinxcrossref{\DUrole{doc}{binary restart files}}}}.

\sphinxAtStartPar
None of the {\hyperref[\detokenize{fix_modify::doc}]{\sphinxcrossref{\DUrole{doc}{fix\_modify}}}} options are relevant to this
fix.

\sphinxAtStartPar
The fix computes a global scalar which can be accessed by various
{\hyperref[\detokenize{Howto_output::doc}]{\sphinxcrossref{\DUrole{doc}{output commands}}}}.  The scalar is the current
temperature of the group of particles described by \sphinxstyleemphasis{group\sphinxhyphen{}ID} along with
the fluid constrained to move with them. The temperature is computed via
the kinetic energy of the group and fluid constrained to move with them
and the total number of degrees of freedom (calculated internally).  If
the particles are not integrated independently (such as via {\hyperref[\detokenize{fix_nve::doc}]{\sphinxcrossref{\DUrole{doc}{fix
NVE}}}}) but have additional constraints imposed on them (such as
via integration using {\hyperref[\detokenize{fix_rigid::doc}]{\sphinxcrossref{\DUrole{doc}{fix rigid}}}}) the degrees of
freedom removed from these additional constraints will not be properly
accounted for.  In this case, the user can specify the total degrees of
freedom independently using the \sphinxstyleemphasis{dof} keyword.

\sphinxAtStartPar
The fix also computes a global array of values which can be accessed by
various {\hyperref[\detokenize{Howto_output::doc}]{\sphinxcrossref{\DUrole{doc}{output commands}}}}.  There are 5 entries in
the array.  The first entry is the temperature of the fluid, the second
entry is the total mass of the fluid plus particles, the third through
fifth entries give the x, y, and z total momentum of the fluid plus
particles.

\sphinxAtStartPar
No parameter of this fix can be used with the \sphinxstyleemphasis{start/stop} keywords of
the {\hyperref[\detokenize{run::doc}]{\sphinxcrossref{\DUrole{doc}{run}}}} command.  This fix is not invoked during
{\hyperref[\detokenize{minimize::doc}]{\sphinxcrossref{\DUrole{doc}{energy minimization}}}}.


\subsection{Restrictions}
\label{\detokenize{fix_lb_fluid:restrictions}}
\sphinxAtStartPar
This fix is part of the LATBOLTZ package.  It is only enabled if LAMMPS
was built with that package.  See the {\hyperref[\detokenize{Build_package::doc}]{\sphinxcrossref{\DUrole{doc}{Build package}}}} page for more info.

\sphinxAtStartPar
This fix can only be used with an orthogonal simulation domain.

\sphinxAtStartPar
The boundary conditions for the fluid are specified independently to the
particles.  However, these should normally be specified consistently via
the main LAMMPS {\hyperref[\detokenize{boundary::doc}]{\sphinxcrossref{\DUrole{doc}{boundary}}}} command (p p p, p p f, and p
f f are the only consistent possibilities).  Shrink\sphinxhyphen{}wrapped boundary
conditions are not permitted with this fix.

\sphinxAtStartPar
This fix must be used before any of {\hyperref[\detokenize{fix_lb_viscous::doc}]{\sphinxcrossref{\DUrole{doc}{fix lb/viscous}}}} and {\hyperref[\detokenize{fix_lb_momentum::doc}]{\sphinxcrossref{\DUrole{doc}{fix lb/momentum}}}} as the
fluid needs to be initialized before any of these routines try to access
its properties.  In addition, in order for the hydrodynamic forces to be
added to the particles, this fix must be used in conjunction with the
{\hyperref[\detokenize{fix_lb_viscous::doc}]{\sphinxcrossref{\DUrole{doc}{lb/viscous}}}} fix.

\sphinxAtStartPar
This fix needs to be used in conjunction with a standard LAMMPS
integrator such as {\hyperref[\detokenize{fix_nve::doc}]{\sphinxcrossref{\DUrole{doc}{fix NVE}}}} or {\hyperref[\detokenize{fix_rigid::doc}]{\sphinxcrossref{\DUrole{doc}{fix rigid}}}}.


\subsection{Related commands}
\label{\detokenize{fix_lb_fluid:related-commands}}
\sphinxAtStartPar
{\hyperref[\detokenize{fix_lb_viscous::doc}]{\sphinxcrossref{\DUrole{doc}{fix lb/viscous}}}}, {\hyperref[\detokenize{fix_lb_momentum::doc}]{\sphinxcrossref{\DUrole{doc}{fix lb/momentum}}}}


\subsection{Default}
\label{\detokenize{fix_lb_fluid:default}}
\sphinxAtStartPar
\sphinxstyleemphasis{dx} is chosen such that \(\frac{\tau}{dt_{LB}} = \frac{3\eta dt_{LB}}{\rho dx_{LB}^2}\) is approximately equal to 1.
\sphinxstyleemphasis{dm} is set equal to 1.0.
\sphinxstyleemphasis{a0} is set equal to \(\frac{1}{3}\left(\frac{dx_{LB}}{dt_{LB}}\right)^2\).
The trilinear stencil is used as the default interpolation method.
The D3Q15 lattice is used for the lattice\sphinxhyphen{}Boltzmann algorithm.


\bigskip\hrule\bigskip

\phantomsection\label{\detokenize{fix_lb_fluid:fluid-denniston}}
\sphinxAtStartPar
\sphinxstylestrong{(Denniston et al.)} Denniston, C., Afrasiabian, N., Cole\sphinxhyphen{}Andre, M.G., Mackay, F. E., Ollila, S.T.T., and Whitehead, T., LAMMPS lb/fluid fix version 2: Improved Hydrodynamic Forces Implemented into LAMMPS through a lattice\sphinxhyphen{}Boltzmann fluid, Computer Physics Communications 275 (2022) \sphinxhref{https://doi.org/10.1016/j.cpc.2022.108318}{108318} .

\phantomsection\label{\detokenize{fix_lb_fluid:mackay2}}
\sphinxAtStartPar
\sphinxstylestrong{(Mackay and Denniston)} Mackay, F. E., and Denniston, C., Coupling MD particles to a lattice\sphinxhyphen{}Boltzmann fluid through the use of conservative forces, J. Comput. Phys. 237 (2013) 289\sphinxhyphen{}298.

\phantomsection\label{\detokenize{fix_lb_fluid:adhikari}}
\sphinxAtStartPar
\sphinxstylestrong{(Adhikari et al.)} Adhikari, R., Stratford, K.,  Cates, M. E., and Wagner, A. J., Fluctuating lattice Boltzmann, Europhys. Lett. 71 (2005) 473\sphinxhyphen{}479.

\sphinxstepscope

\index{fix lb/momentum@\spxentry{fix lb/momentum}}\ignorespaces 

\section{fix lb/momentum command}
\label{\detokenize{fix_lb_momentum:fix-lb-momentum-command}}\label{\detokenize{fix_lb_momentum:index-0}}\label{\detokenize{fix_lb_momentum::doc}}

\subsection{Syntax}
\label{\detokenize{fix_lb_momentum:syntax}}
\begin{sphinxVerbatim}[commandchars=\\\{\}]
\PYG{k}{fix }\PYG{n+nv+nvIdentifier}{ID}\PYG{+w}{ }\PYG{n+nv+nvIdentifier}{group\PYGZhy{}ID}\PYG{+w}{ }\PYG{n}{lb}\PYG{o}{/}\PYG{n}{momentum}\PYG{+w}{ }\PYG{n}{nevery}\PYG{+w}{ }\PYG{n}{keyword}\PYG{+w}{ }\PYG{n}{values}\PYG{+w}{ }\PYG{n}{...}
\end{sphinxVerbatim}
\begin{itemize}
\item {} 
\sphinxAtStartPar
ID, group\sphinxhyphen{}ID are documented in the {\hyperref[\detokenize{fix::doc}]{\sphinxcrossref{\DUrole{doc}{fix}}}} command

\item {} 
\sphinxAtStartPar
lb/momentum = style name of this fix command

\item {} 
\sphinxAtStartPar
nevery = adjust the momentum every this many timesteps

\item {} 
\sphinxAtStartPar
zero or more keyword/value pairs may be appended

\item {} 
\sphinxAtStartPar
keyword = \sphinxstyleemphasis{linear}
\begin{sphinxalltt}
\sphinxstyleemphasis{linear} values = xflag yflag zflag
  xflag,yflag,zflag = 0/1 to exclude/include each dimension.
\end{sphinxalltt}

\end{itemize}


\subsection{Examples}
\label{\detokenize{fix_lb_momentum:examples}}
\begin{sphinxVerbatim}[commandchars=\\\{\}]
\PYG{k}{fix }\PYG{n+nv+nvIdentifier}{1}\PYG{+w}{ }\PYG{n+nv+nvIdentifier}{sphere}\PYG{+w}{ }\PYG{n}{lb}\PYG{o}{/}\PYG{n}{momentum}
\PYG{k}{fix }\PYG{n+nv+nvIdentifier}{1}\PYG{+w}{ }\PYG{n+nv+nvIdentifier}{all}\PYG{+w}{ }\PYG{n}{lb}\PYG{o}{/}\PYG{n}{momentum}\PYG{+w}{ }\PYG{n}{linear}\PYG{+w}{ }\PYG{l+m}{1}\PYG{+w}{ }\PYG{l+m}{1}\PYG{+w}{ }\PYG{l+m}{0}
\end{sphinxVerbatim}


\subsection{Description}
\label{\detokenize{fix_lb_momentum:description}}
\sphinxAtStartPar
This fix is based on the {\hyperref[\detokenize{fix_momentum::doc}]{\sphinxcrossref{\DUrole{doc}{fix momentum}}}} command, and
was created to be used in place of that command, when a
lattice\sphinxhyphen{}Boltzmann fluid is present.

\sphinxAtStartPar
Zero the total linear momentum of the system, including both the atoms
specified by group\sphinxhyphen{}ID and the lattice\sphinxhyphen{}Boltzmann fluid every nevery
timesteps.  If there are no atoms specified by group\sphinxhyphen{}ID only the fluid momentum is affected.  This is accomplished by adjusting the particle velocities
and the fluid velocities at each lattice site.

\begin{sphinxadmonition}{note}{Note:}
\sphinxAtStartPar
This fix only considers the linear momentum of the system.
\end{sphinxadmonition}

\sphinxAtStartPar
By default, the subtraction is performed for each dimension.  This can
be changed by specifying the keyword \sphinxstyleemphasis{linear}, along with a set of
three flags set to 0/1 in order to exclude/ include the corresponding
dimension.


\subsection{Restart, fix\_modify, output, run start/stop, minimize info}
\label{\detokenize{fix_lb_momentum:restart-fix-modify-output-run-start-stop-minimize-info}}
\sphinxAtStartPar
No information about this fix is written to {\hyperref[\detokenize{restart::doc}]{\sphinxcrossref{\DUrole{doc}{binary restart files}}}}.  None of the {\hyperref[\detokenize{fix_modify::doc}]{\sphinxcrossref{\DUrole{doc}{fix\_modify}}}} options
are relevant to this fix.  No global or per\sphinxhyphen{}atom quantities are stored
by this fix for access by various {\hyperref[\detokenize{Howto_output::doc}]{\sphinxcrossref{\DUrole{doc}{output commands}}}}.
No parameter of this fix can be used with the \sphinxstyleemphasis{start/stop} keywords of
the {\hyperref[\detokenize{run::doc}]{\sphinxcrossref{\DUrole{doc}{run}}}} command.  This fix is not invoked during {\hyperref[\detokenize{minimize::doc}]{\sphinxcrossref{\DUrole{doc}{energy minimization}}}}.


\subsection{Restrictions}
\label{\detokenize{fix_lb_momentum:restrictions}}
\sphinxAtStartPar
Can only be used if a lattice\sphinxhyphen{}Boltzmann fluid has been created via the
{\hyperref[\detokenize{fix_lb_fluid::doc}]{\sphinxcrossref{\DUrole{doc}{fix lb/fluid}}}} command, and must come after this
command.

\sphinxAtStartPar
This fix is part of the LATBOLTZ package.  It is only enabled if LAMMPS
was built with that package.  See the {\hyperref[\detokenize{Build_package::doc}]{\sphinxcrossref{\DUrole{doc}{Build package}}}} page for more info.


\subsection{Related commands}
\label{\detokenize{fix_lb_momentum:related-commands}}
\sphinxAtStartPar
{\hyperref[\detokenize{fix_momentum::doc}]{\sphinxcrossref{\DUrole{doc}{fix momentum}}}}, {\hyperref[\detokenize{fix_lb_fluid::doc}]{\sphinxcrossref{\DUrole{doc}{fix lb/fluid}}}}


\subsection{Default}
\label{\detokenize{fix_lb_momentum:default}}
\sphinxAtStartPar
Zeros the total system linear momentum in each dimension.

\sphinxstepscope

\index{fix lb/viscous@\spxentry{fix lb/viscous}}\ignorespaces 

\section{fix lb/viscous command}
\label{\detokenize{fix_lb_viscous:fix-lb-viscous-command}}\label{\detokenize{fix_lb_viscous:index-0}}\label{\detokenize{fix_lb_viscous::doc}}

\subsection{Syntax}
\label{\detokenize{fix_lb_viscous:syntax}}
\begin{sphinxVerbatim}[commandchars=\\\{\}]
\PYG{k}{fix }\PYG{n+nv+nvIdentifier}{ID}\PYG{+w}{ }\PYG{n+nv+nvIdentifier}{group\PYGZhy{}ID}\PYG{+w}{ }\PYG{n}{lb}\PYG{o}{/}\PYG{n}{viscous}
\end{sphinxVerbatim}
\begin{itemize}
\item {} 
\sphinxAtStartPar
ID, group\sphinxhyphen{}ID are documented in {\hyperref[\detokenize{fix::doc}]{\sphinxcrossref{\DUrole{doc}{fix}}}} command

\item {} 
\sphinxAtStartPar
lb/viscous = style name of this fix command

\end{itemize}


\subsection{Examples}
\label{\detokenize{fix_lb_viscous:examples}}
\begin{sphinxVerbatim}[commandchars=\\\{\}]
\PYG{k}{fix }\PYG{n+nv+nvIdentifier}{1}\PYG{+w}{ }\PYG{n+nv+nvIdentifier}{flow}\PYG{+w}{ }\PYG{n}{lb}\PYG{o}{/}\PYG{n}{viscous}
\end{sphinxVerbatim}


\subsection{Description}
\label{\detokenize{fix_lb_viscous:description}}
\sphinxAtStartPar
This fix is similar to the {\hyperref[\detokenize{fix_viscous::doc}]{\sphinxcrossref{\DUrole{doc}{fix viscous}}}} command, and
is to be used in place of that command when a lattice\sphinxhyphen{}Boltzmann fluid
is present using the {\hyperref[\detokenize{fix_lb_fluid::doc}]{\sphinxcrossref{\DUrole{doc}{fix lb/fluid}}}}.  This should be used in conjunction with one of the built\sphinxhyphen{}in LAMMPS integrators, such as {\hyperref[\detokenize{fix_nve::doc}]{\sphinxcrossref{\DUrole{doc}{fix NVE}}}} or {\hyperref[\detokenize{fix_rigid::doc}]{\sphinxcrossref{\DUrole{doc}{fix rigid}}}}.

\sphinxAtStartPar
This fix adds a viscous force to each atom to cause it move with the same velocity as the fluid (an equal and opposite force is applied to the fluid via {\hyperref[\detokenize{fix_lb_fluid::doc}]{\sphinxcrossref{\DUrole{doc}{fix lb/fluid}}}}).  When {\hyperref[\detokenize{fix_lb_fluid::doc}]{\sphinxcrossref{\DUrole{doc}{fix lb/fluid}}}} is called with the noise option, the atoms will also experience random forces which will thermalize them to the same temperature as the fluid.  In this way, the combination of this fix with {\hyperref[\detokenize{fix_lb_fluid::doc}]{\sphinxcrossref{\DUrole{doc}{fix lb/fluid}}}} and a LAMMPS integrator like {\hyperref[\detokenize{fix_nve::doc}]{\sphinxcrossref{\DUrole{doc}{fix NVE}}}} is analogous to {\hyperref[\detokenize{fix_langevin::doc}]{\sphinxcrossref{\DUrole{doc}{fix langevin}}}} except here the fluid is explicit.  The temperature of the particles can be monitored via the scalar output of {\hyperref[\detokenize{fix_lb_fluid::doc}]{\sphinxcrossref{\DUrole{doc}{fix lb/fluid}}}}.


\bigskip\hrule\bigskip


\sphinxAtStartPar
For details of this fix, as well as descriptions and results of several
test runs, see {\hyperref[\detokenize{fix_lb_viscous:fluid-denniston2}]{\sphinxcrossref{\DUrole{std,std-ref}{Denniston et al.}}}}.  Please include a citation to
this paper if this fix is used in work contributing to published
research.


\bigskip\hrule\bigskip



\subsection{Restart, fix\_modify, output, run start/stop, minimize info}
\label{\detokenize{fix_lb_viscous:restart-fix-modify-output-run-start-stop-minimize-info}}
\sphinxAtStartPar
As described in the {\hyperref[\detokenize{fix_viscous::doc}]{\sphinxcrossref{\DUrole{doc}{fix viscous}}}} documentation:

\sphinxAtStartPar
“No information about this fix is written to {\hyperref[\detokenize{restart::doc}]{\sphinxcrossref{\DUrole{doc}{binary restart files}}}}.  None of the {\hyperref[\detokenize{fix_modify::doc}]{\sphinxcrossref{\DUrole{doc}{fix\_modify}}}} options
are relevant to this fix.  No global or per\sphinxhyphen{}atom quantities are stored
by this fix for access by various {\hyperref[\detokenize{Howto_output::doc}]{\sphinxcrossref{\DUrole{doc}{output commands}}}}.
No parameter of this fix can be used with the \sphinxstyleemphasis{start/stop} keywords of
the {\hyperref[\detokenize{run::doc}]{\sphinxcrossref{\DUrole{doc}{run}}}} command.

\sphinxAtStartPar
The forces due to this fix are imposed during an energy minimization,
invoked by the {\hyperref[\detokenize{minimize::doc}]{\sphinxcrossref{\DUrole{doc}{minimize}}}} command.  This fix should only
be used with damped dynamics minimizers that allow for
non\sphinxhyphen{}conservative forces.  See the {\hyperref[\detokenize{min_style::doc}]{\sphinxcrossref{\DUrole{doc}{min\_style}}}} command
for details.”


\subsection{Restrictions}
\label{\detokenize{fix_lb_viscous:restrictions}}
\sphinxAtStartPar
This fix is part of the LATBOLTZ package.  It is only enabled if LAMMPS
was built with that package.  See the {\hyperref[\detokenize{Build_package::doc}]{\sphinxcrossref{\DUrole{doc}{Build package}}}} page for more info.

\sphinxAtStartPar
Can only be used if a lattice\sphinxhyphen{}Boltzmann fluid has been created via the
{\hyperref[\detokenize{fix_lb_fluid::doc}]{\sphinxcrossref{\DUrole{doc}{fix lb/fluid}}}} command, and must come after this
command.


\subsection{Related commands}
\label{\detokenize{fix_lb_viscous:related-commands}}
\sphinxAtStartPar
{\hyperref[\detokenize{fix_lb_fluid::doc}]{\sphinxcrossref{\DUrole{doc}{fix lb/fluid}}}}


\subsection{Default}
\label{\detokenize{fix_lb_viscous:default}}
\sphinxAtStartPar
none


\bigskip\hrule\bigskip

\phantomsection\label{\detokenize{fix_lb_viscous:fluid-denniston2}}
\sphinxAtStartPar
\sphinxstylestrong{(Denniston et al.)} Denniston, C., Afrasiabian, N., Cole\sphinxhyphen{}Andre, M.G., Mackay, F. E., Ollila, S.T.T., and Whitehead, T., LAMMPS lb/fluid fix version 2: Improved Hydrodynamic Forces Implemented into LAMMPS through a lattice\sphinxhyphen{}Boltzmann fluid, Computer Physics Communications 275 (2022) \sphinxhref{https://doi.org/10.1016/j.cpc.2022.108318}{108318} .

\sphinxstepscope

\index{fix lineforce@\spxentry{fix lineforce}}\ignorespaces 

\section{fix lineforce command}
\label{\detokenize{fix_lineforce:fix-lineforce-command}}\label{\detokenize{fix_lineforce:index-0}}\label{\detokenize{fix_lineforce::doc}}

\subsection{Syntax}
\label{\detokenize{fix_lineforce:syntax}}
\begin{sphinxVerbatim}[commandchars=\\\{\}]
\PYG{k}{fix }\PYG{n+nv+nvIdentifier}{ID}\PYG{+w}{ }\PYG{n+nv+nvIdentifier}{group\PYGZhy{}ID}\PYG{+w}{ }\PYG{n}{lineforce}\PYG{+w}{ }\PYG{n}{x}\PYG{+w}{ }\PYG{n}{y}\PYG{+w}{ }\PYG{n}{z}
\end{sphinxVerbatim}
\begin{itemize}
\item {} 
\sphinxAtStartPar
ID, group\sphinxhyphen{}ID are documented in {\hyperref[\detokenize{fix::doc}]{\sphinxcrossref{\DUrole{doc}{fix}}}} command

\item {} 
\sphinxAtStartPar
lineforce = style name of this fix command

\item {} 
\sphinxAtStartPar
x y z = direction of line as a 3\sphinxhyphen{}vector

\end{itemize}


\subsection{Examples}
\label{\detokenize{fix_lineforce:examples}}
\begin{sphinxVerbatim}[commandchars=\\\{\}]
\PYG{k}{fix }\PYG{n+nv+nvIdentifier}{hold}\PYG{+w}{ }\PYG{n+nv+nvIdentifier}{boundary}\PYG{+w}{ }\PYG{n}{lineforce}\PYG{+w}{ }\PYG{l+m}{0.0}\PYG{+w}{ }\PYG{l+m}{1.0}\PYG{+w}{ }\PYG{l+m}{1.0}
\end{sphinxVerbatim}


\subsection{Description}
\label{\detokenize{fix_lineforce:description}}
\sphinxAtStartPar
Adjust the forces on each atom in the group so that only the component
of force along the linear direction specified by the vector (x,y,z)
remains.  This is done by subtracting out components of force in the
plane perpendicular to the line.

\sphinxAtStartPar
If the initial velocity of the atom is 0.0 (or along the line), then
it should continue to move along the line thereafter.


\subsection{Restart, fix\_modify, output, run start/stop, minimize info}
\label{\detokenize{fix_lineforce:restart-fix-modify-output-run-start-stop-minimize-info}}
\sphinxAtStartPar
No information about this fix is written to {\hyperref[\detokenize{restart::doc}]{\sphinxcrossref{\DUrole{doc}{binary restart files}}}}.  None of the {\hyperref[\detokenize{fix_modify::doc}]{\sphinxcrossref{\DUrole{doc}{fix\_modify}}}} options
are relevant to this fix.  No global or per\sphinxhyphen{}atom quantities are stored
by this fix for access by various {\hyperref[\detokenize{Howto_output::doc}]{\sphinxcrossref{\DUrole{doc}{output commands}}}}.
No parameter of this fix can be used with the \sphinxstyleemphasis{start/stop} keywords of
the {\hyperref[\detokenize{run::doc}]{\sphinxcrossref{\DUrole{doc}{run}}}} command.

\sphinxAtStartPar
The forces due to this fix are imposed during an energy minimization,
invoked by the {\hyperref[\detokenize{minimize::doc}]{\sphinxcrossref{\DUrole{doc}{minimize}}}} command.


\subsection{Restrictions}
\label{\detokenize{fix_lineforce:restrictions}}\begin{quote}

\sphinxAtStartPar
none
\end{quote}


\subsection{Related commands}
\label{\detokenize{fix_lineforce:related-commands}}
\sphinxAtStartPar
{\hyperref[\detokenize{fix_planeforce::doc}]{\sphinxcrossref{\DUrole{doc}{fix planeforce}}}}


\subsection{Default}
\label{\detokenize{fix_lineforce:default}}
\sphinxAtStartPar
none

\sphinxstepscope

\index{fix manifoldforce@\spxentry{fix manifoldforce}}\ignorespaces 

\section{fix manifoldforce command}
\label{\detokenize{fix_manifoldforce:fix-manifoldforce-command}}\label{\detokenize{fix_manifoldforce:index-0}}\label{\detokenize{fix_manifoldforce::doc}}

\subsection{Syntax}
\label{\detokenize{fix_manifoldforce:syntax}}
\begin{sphinxVerbatim}[commandchars=\\\{\}]
\PYG{k}{fix }\PYG{n+nv+nvIdentifier}{ID}\PYG{+w}{ }\PYG{n+nv+nvIdentifier}{group\PYGZhy{}ID}\PYG{+w}{ }\PYG{n}{manifoldforce}\PYG{+w}{ }\PYG{n}{manifold}\PYG{+w}{ }\PYG{n}{manifold}\PYG{o}{\PYGZhy{}}\PYG{n}{args}\PYG{+w}{ }\PYG{n}{...}
\end{sphinxVerbatim}
\begin{itemize}
\item {} 
\sphinxAtStartPar
ID, group\sphinxhyphen{}ID are documented in {\hyperref[\detokenize{fix::doc}]{\sphinxcrossref{\DUrole{doc}{fix}}}} command

\item {} 
\sphinxAtStartPar
manifold = name of the manifold

\item {} 
\sphinxAtStartPar
manifold\sphinxhyphen{}args = parameters for the manifold

\end{itemize}


\subsection{Examples}
\label{\detokenize{fix_manifoldforce:examples}}
\begin{sphinxVerbatim}[commandchars=\\\{\}]
\PYG{k}{fix }\PYG{n+nv+nvIdentifier}{constrain}\PYG{+w}{ }\PYG{n+nv+nvIdentifier}{all}\PYG{+w}{ }\PYG{n}{manifoldforce}\PYG{+w}{ }\PYG{n}{sphere}\PYG{+w}{ }\PYG{l+m}{5.0}
\end{sphinxVerbatim}


\subsection{Description}
\label{\detokenize{fix_manifoldforce:description}}
\sphinxAtStartPar
This fix subtracts each time step from the force the component along
the normal of the specified {\hyperref[\detokenize{Howto_manifold::doc}]{\sphinxcrossref{\DUrole{doc}{manifold}}}}.  This can be
used in combination with {\hyperref[\detokenize{minimize::doc}]{\sphinxcrossref{\DUrole{doc}{minimize}}}} to remove overlap
between particles while keeping them (roughly) constrained to the
given manifold, e.g. to set up a run with {\hyperref[\detokenize{fix_nve_manifold_rattle::doc}]{\sphinxcrossref{\DUrole{doc}{fix nve/manifold/rattle}}}}.  I have found that
only \sphinxstyleemphasis{hftn} and \sphinxstyleemphasis{quickmin} with a very small time step perform
adequately though.


\bigskip\hrule\bigskip



\subsection{Restart, fix\_modify, output, run start/stop, minimize info}
\label{\detokenize{fix_manifoldforce:restart-fix-modify-output-run-start-stop-minimize-info}}
\sphinxAtStartPar
No information about this fix is written to {\hyperref[\detokenize{restart::doc}]{\sphinxcrossref{\DUrole{doc}{binary restart files}}}}.  None of the {\hyperref[\detokenize{fix_modify::doc}]{\sphinxcrossref{\DUrole{doc}{fix\_modify}}}} options
are relevant to this fix.  No global or per\sphinxhyphen{}atom quantities are stored
by this fix for access by various {\hyperref[\detokenize{Howto_output::doc}]{\sphinxcrossref{\DUrole{doc}{output commands}}}}.
No parameter of this fix can be used with the \sphinxstyleemphasis{start/stop} keywords of
the {\hyperref[\detokenize{run::doc}]{\sphinxcrossref{\DUrole{doc}{run}}}} command.  This fix is invoked during {\hyperref[\detokenize{minimize::doc}]{\sphinxcrossref{\DUrole{doc}{energy minimization}}}}.


\bigskip\hrule\bigskip



\subsection{Restrictions}
\label{\detokenize{fix_manifoldforce:restrictions}}
\sphinxAtStartPar
This fix is part of the MANIFOLD package. It is only enabled if
LAMMPS was built with that package.  See the {\hyperref[\detokenize{Build_package::doc}]{\sphinxcrossref{\DUrole{doc}{Build package}}}} page for more info.

\sphinxAtStartPar
Only use this with \sphinxstyleemphasis{min\_style hftn} or \sphinxstyleemphasis{min\_style quickmin}. If not,
the constraints will not be satisfied very well at all. A warning is
generated if the \sphinxstyleemphasis{min\_style} is incompatible but no error.


\bigskip\hrule\bigskip



\subsection{Related commands}
\label{\detokenize{fix_manifoldforce:related-commands}}
\sphinxAtStartPar
{\hyperref[\detokenize{fix_nve_manifold_rattle::doc}]{\sphinxcrossref{\DUrole{doc}{fix nve/manifold/rattle}}}}, {\hyperref[\detokenize{fix_nvt_manifold_rattle::doc}]{\sphinxcrossref{\DUrole{doc}{fix nvt/manifold/rattle}}}}

\sphinxstepscope

\index{fix mdi/qm@\spxentry{fix mdi/qm}}\ignorespaces 

\section{fix mdi/qm command}
\label{\detokenize{fix_mdi_qm:fix-mdi-qm-command}}\label{\detokenize{fix_mdi_qm:index-0}}\label{\detokenize{fix_mdi_qm::doc}}

\subsection{Syntax}
\label{\detokenize{fix_mdi_qm:syntax}}
\begin{sphinxVerbatim}[commandchars=\\\{\}]
\PYG{k}{fix }\PYG{n+nv+nvIdentifier}{ID}\PYG{+w}{ }\PYG{n+nv+nvIdentifier}{group\PYGZhy{}ID}\PYG{+w}{ }\PYG{n}{mdi}\PYG{o}{/}\PYG{n}{qm}\PYG{+w}{ }\PYG{n}{keyword}\PYG{+w}{ }\PYG{n}{value}\PYG{n+nv}{(}\PYG{n+nv}{s}\PYG{n+nv}{)}\PYG{+w}{ }\PYG{n}{keyword}\PYG{+w}{ }\PYG{n}{value}\PYG{n+nv}{(}\PYG{n+nv}{s}\PYG{n+nv}{)}\PYG{+w}{ }\PYG{n}{...}
\end{sphinxVerbatim}
\begin{itemize}
\item {} 
\sphinxAtStartPar
ID, group\sphinxhyphen{}ID are documented in {\hyperref[\detokenize{fix::doc}]{\sphinxcrossref{\DUrole{doc}{fix}}}} command

\item {} 
\sphinxAtStartPar
mdi/qm = style name of this fix command

\item {} 
\sphinxAtStartPar
zero or more keyword/value pairs may be appended

\item {} 
\sphinxAtStartPar
keyword = \sphinxstyleemphasis{virial} or \sphinxstyleemphasis{add} or \sphinxstyleemphasis{every} or \sphinxstyleemphasis{connect} or \sphinxstyleemphasis{elements} or \sphinxstyleemphasis{mc}
\begin{sphinxalltt}
\sphinxstyleemphasis{virial} args = \sphinxstyleemphasis{yes} or \sphinxstyleemphasis{no}
  yes = request virial tensor from server code
  no = do not request virial tensor from server code
\sphinxstyleemphasis{add} args = \sphinxstyleemphasis{yes} or \sphinxstyleemphasis{no}
  yes = add returned value from server code to LAMMPS quantities
  no = do not add returned values to LAMMPS quantities
\sphinxstyleemphasis{every} args = Nevery
  Nevery = request values from server code once every Nevery steps
\sphinxstyleemphasis{connect} args = \sphinxstyleemphasis{yes} or \sphinxstyleemphasis{no}
  yes = perform a one\sphinxhyphen{}time connection to the MDI engine code
  no = do not perform the connection operation
\sphinxstyleemphasis{elements} args = N\_1 N\_2 ... N\_ntypes
  N\_1,N\_2,...N\_ntypes = chemical symbol for each of ntypes LAMMPS atom types
\sphinxstyleemphasis{mc} args = mcfixID
  mcfixID = ID of a Monte Carlo fix designed to work with this fix
\end{sphinxalltt}

\end{itemize}


\subsection{Examples}
\label{\detokenize{fix_mdi_qm:examples}}
\begin{sphinxVerbatim}[commandchars=\\\{\}]
\PYG{k}{fix }\PYG{n+nv+nvIdentifier}{1}\PYG{+w}{ }\PYG{n+nv+nvIdentifier}{all}\PYG{+w}{ }\PYG{n}{mdi}\PYG{o}{/}\PYG{n}{qm}
\PYG{k}{fix }\PYG{n+nv+nvIdentifier}{1}\PYG{+w}{ }\PYG{n+nv+nvIdentifier}{all}\PYG{+w}{ }\PYG{n}{mdi}\PYG{o}{/}\PYG{n}{qm}\PYG{+w}{ }\PYG{n}{virial}\PYG{+w}{ }\PYG{n}{yes}
\PYG{k}{fix }\PYG{n+nv+nvIdentifier}{1}\PYG{+w}{ }\PYG{n+nv+nvIdentifier}{all}\PYG{+w}{ }\PYG{n}{mdi}\PYG{o}{/}\PYG{n}{qm}\PYG{+w}{ }\PYG{n}{add}\PYG{+w}{ }\PYG{n}{no}\PYG{+w}{ }\PYG{n}{every}\PYG{+w}{ }\PYG{l+m}{100}\PYG{+w}{ }\PYG{n}{elements}\PYG{+w}{ }\PYG{n}{C}\PYG{+w}{ }\PYG{n}{C}\PYG{+w}{ }\PYG{n}{H}\PYG{+w}{ }\PYG{n}{O}
\end{sphinxVerbatim}


\subsection{Description}
\label{\detokenize{fix_mdi_qm:description}}
\sphinxAtStartPar
\DUrole{versionmodified,added}{Added in version 3Aug2022.}

\sphinxAtStartPar
This command enables LAMMPS to act as a client with another server
code that will compute the total energy, per\sphinxhyphen{}atom forces, and total
virial for atom conformations and simulation box size/shapes that
LAMMPS sends it.

\sphinxAtStartPar
Typically the server code will be a quantum mechanics (QM) code, hence
the name of the fix.  However this is not required, the server code
could be another classical molecular dynamics code or LAMMPS itself.
The server code must support use of the \sphinxhref{https://molssi-mdi.github.io/MDI\_Library/html/index.html}{MDI Library} as
explained below.

\sphinxAtStartPar
Typically, to use this fix, the input script should not define any
other classical force field components, e.g. a pair style, bond style,
etc.

\sphinxAtStartPar
These are example use cases for this fix, discussed further below:
\begin{itemize}
\item {} 
\sphinxAtStartPar
perform an ab initio MD (AIMD) simulation with quantum forces

\item {} 
\sphinxAtStartPar
perform an energy minimization with quantum forces

\item {} 
\sphinxAtStartPar
perform a nudged elastic band (NEB) calculation with quantum forces

\item {} 
\sphinxAtStartPar
perform a QM calculation for a series of independent systems which
LAMMPS reads or generates once

\item {} 
\sphinxAtStartPar
run a classical MD simulation and calculate QM energy/forces once
every N steps on the current configuration

\end{itemize}

\sphinxAtStartPar
More generally any command which calculates per\sphinxhyphen{}atom forces can instead
use quantum forces by defining this fix.  Examples are the Monte Carlo
commands {\hyperref[\detokenize{fix_gcmc::doc}]{\sphinxcrossref{\DUrole{doc}{fix gcmc}}}} and {\hyperref[\detokenize{fix_atom_swap::doc}]{\sphinxcrossref{\DUrole{doc}{fix atom/swap}}}}, as well as the {\hyperref[\detokenize{compute_born_matrix::doc}]{\sphinxcrossref{\DUrole{doc}{compute born/matrix}}}} command.  The only requirement is that internally
the command invokes the post\_force() method of fixes such as this one,
which will trigger the quantum calculation.

\sphinxAtStartPar
The code coupling performed by this command is done via the \sphinxhref{https://molssi-mdi.github.io/MDI\_Library/html/index.html}{MDI
Library}.
LAMMPS runs as an MDI driver (client), and sends MDI commands to an
external MDI engine code (server), e.g. a QM code which has support
for MDI.  See the {\hyperref[\detokenize{Howto_mdi::doc}]{\sphinxcrossref{\DUrole{doc}{Howto mdi}}}} page for more
information about how LAMMPS can operate as either an MDI driver or
engine.

\sphinxAtStartPar
The \sphinxcode{\sphinxupquote{examples/mdi}} directory contains input scripts using this fix in
the various use cases discussed below.  In each case, two instances of
LAMMPS are used, once as an MDI driver, once as an MDI engine (surrogate
for a QM code).  The \sphinxcode{\sphinxupquote{examples/mdi/README}} file explains how to launch
two codes so that they communicate via the MDI library using either MPI
or sockets.  Any QM code that supports MDI could be used in place of
LAMMPS acting as a QM surrogate.  See the {\hyperref[\detokenize{Howto_mdi::doc}]{\sphinxcrossref{\DUrole{doc}{Howto mdi}}}}
page for a current list (March 2022) of such QM codes.  The
\sphinxcode{\sphinxupquote{examples/QUANTUM}} directory has examples for coupling LAMMPS to 3 QM
codes either via this fix or the {\hyperref[\detokenize{fix_mdi_qmmm::doc}]{\sphinxcrossref{\DUrole{doc}{fix mdi/qmmm}}}}
command.

\sphinxAtStartPar
Note that an engine code can support MDI in either or both of two modes.
It can be used as a stand\sphinxhyphen{}alone code, launched at the same time as
LAMMPS.  Or it can be used as a plugin library, which LAMMPS loads.  See
the {\hyperref[\detokenize{mdi::doc}]{\sphinxcrossref{\DUrole{doc}{mdi plugin}}}} command for how to trigger LAMMPS to load a
plugin library.  The \sphinxcode{\sphinxupquote{examples/mdi/README}} file and
\sphinxcode{\sphinxupquote{examples/QUANTUM/QM\sphinxhyphen{}code/README}} files explain how to launch the two
codes in either mode.


\bigskip\hrule\bigskip


\sphinxAtStartPar
The \sphinxstyleemphasis{virial} keyword setting of yes or no determines whether LAMMPS
will request the QM code to also compute and return the QM
contribution to a stress tensor for the system which LAMMPS will
convert to a 6\sphinxhyphen{}element symmetric virial tensor.

\sphinxAtStartPar
The \sphinxstyleemphasis{add} keyword setting of \sphinxstyleemphasis{yes} or \sphinxstyleemphasis{no} determines whether the
energy and forces and virial returned by the QM code will be added to
the LAMMPS internal energy and forces and virial or not.  If the
setting is \sphinxstyleemphasis{no} then the default {\hyperref[\detokenize{fix_modify::doc}]{\sphinxcrossref{\DUrole{doc}{fix\_modify energy}}}}
and {\hyperref[\detokenize{fix_modify::doc}]{\sphinxcrossref{\DUrole{doc}{fix\_modify virial}}}} settings are also set to
\sphinxstyleemphasis{no} and your input scripts should not set them to yes.  See more
details on these fix\_modify settings below.

\sphinxAtStartPar
Whatever the setting for the \sphinxstyleemphasis{add} keyword, the QM energy, forces, and
virial will be stored by the fix, so they can be accessed by other
commands.  See details below.

\sphinxAtStartPar
The \sphinxstyleemphasis{every} keyword determines how often the QM code will be invoked
during a dynamics run with the current LAMMPS simulation box and
configuration of atoms.  The QM code will be called once every
\sphinxstyleemphasis{Nevery} timesteps.  By default \sphinxstyleemphasis{Nevery} = 1.

\sphinxAtStartPar
The \sphinxstyleemphasis{connect} keyword determines whether this fix performs a one\sphinxhyphen{}time
connection to the QM code.  The default is \sphinxstyleemphasis{yes}.  The only time a \sphinxstyleemphasis{no}
is needed is if this command is used multiple times in an input script
and the MDI coupling is between two stand\sphinxhyphen{}alone codes (not plugin mode).
E.g. if it used inside a loop which also uses the {\hyperref[\detokenize{clear::doc}]{\sphinxcrossref{\DUrole{doc}{clear}}}}
command to destroy the system (including this fix).  See the
\sphinxcode{\sphinxupquote{examples/mdi/in.series.driver}} script as an example of this, where
LAMMPS is using the QM code to compute energy and forces for a series of
system configurations.  In this use case \sphinxstyleemphasis{connect no} is used along with
the {\hyperref[\detokenize{mdi::doc}]{\sphinxcrossref{\DUrole{doc}{mdi connect and exit}}}} command to one\sphinxhyphen{}time
initiate/terminate the connection outside the loop.

\sphinxAtStartPar
The \sphinxstyleemphasis{elements} keyword allows specification of what element each
LAMMPS atom type corresponds to.  This is specified by the chemical
symbol of the element, e.g. C or Al or Si.  A symbol must be specified
for each of the ntypes LAMMPS atom types.  Multiple LAMMPS types can
represent the same element.  Ntypes is typically specified via the
{\hyperref[\detokenize{create_box::doc}]{\sphinxcrossref{\DUrole{doc}{create\_box}}}} command or in the data file read by the
{\hyperref[\detokenize{read_data::doc}]{\sphinxcrossref{\DUrole{doc}{read\_data}}}} command.

\sphinxAtStartPar
If this keyword is specified, then this fix will send the MDI
“\textgreater{}ELEMENTS” command to the engine, to ensure the two codes are
consistent in their definition of atomic species.  If this keyword is
not specified, then this fix will send the MDI \textgreater{}TYPES command to the
engine.  This is fine if both the LAMMPS driver and the MDI engine are
initialized so that the atom type values are consistent in both codes.

\sphinxAtStartPar
The \sphinxstyleemphasis{mc} keyword enables this fix to be used with a Monte Carlo (MC)
fix to calculate before/after quantum energies as part of the MC
accept/reject criterion.  The {\hyperref[\detokenize{fix_gcmc::doc}]{\sphinxcrossref{\DUrole{doc}{fix gcmc}}}} and {\hyperref[\detokenize{fix_atom_swap::doc}]{\sphinxcrossref{\DUrole{doc}{fix
atom/swap}}}} commands can be used in this manner.
Specify the ID of the MC fix following the \sphinxstyleemphasis{mc} keyword.  This allows
the two fixes to coordinate when MC events are being calculated versus
MD timesteps between the MC events.


\bigskip\hrule\bigskip


\sphinxAtStartPar
The following 3 example use cases are illustrated in the
\sphinxcode{\sphinxupquote{examples/mdi}} directory.  See its README file for more details.

\sphinxAtStartPar
(1) To run an ab initio MD (AIMD) dynamics simulation, or an energy
minimization with QM forces, or a multi\sphinxhyphen{}replica NEB calculation, use
\sphinxstyleemphasis{add yes} and \sphinxstyleemphasis{every 1} (the defaults).  This is so that every time
LAMMPS needs energy and forces, the QM code will be invoked.

\sphinxAtStartPar
Both LAMMPS and the QM code should define the same system (simulation
box, atoms and their types) in their respective input scripts.  Note
that on this scenario, it may not be necessary for LAMMPS to define a
pair style or use a neighbor list.

\sphinxAtStartPar
LAMMPS will then perform the timestepping or minimization iterations
for the simulation.  At the point in each timestep or iteration when
LAMMPS needs the force on each atom, it communicates with the engine
code.  It sends the current simulation box size and shape (if they
change dynamically, e.g. during an NPT simulation), and the current
atom coordinates.  The engine code computes quantum forces on each
atom and the total energy of the system and returns them to LAMMPS.

\sphinxAtStartPar
Note that if the AIMD simulation is an NPT or NPH model, or the energy
minimization includes {\hyperref[\detokenize{fix_box_relax::doc}]{\sphinxcrossref{\DUrole{doc}{fix box relax}}}} to
equilibrate the box size/shape, then LAMMPS computes a pressure.  This
means the \sphinxstyleemphasis{virial} keyword should be set to \sphinxstyleemphasis{yes} so that the QM
contribution to the pressure can be included.

\sphinxAtStartPar
(2) To run dynamics with a LAMMPS interatomic potential, and evaluate
the QM energy and forces once every 1000 steps, use \sphinxstyleemphasis{add no} and
\sphinxstyleemphasis{every 1000}.  This could be useful for using an MD run to generate
randomized configurations which are then passed to the QM code to
produce training data for a machine learning potential.  A {\hyperref[\detokenize{dump::doc}]{\sphinxcrossref{\DUrole{doc}{dump
custom}}}} command could be invoked every 1000 steps to dump the
atom coordinates and QM forces to a file.  Likewise the QM energy and
virial could be output with the {\hyperref[\detokenize{thermo_style::doc}]{\sphinxcrossref{\DUrole{doc}{thermo\_style custom}}}} command.

\sphinxAtStartPar
(3) To do a QM evaluation of energy and forces for a series of \sphinxstyleemphasis{N}
independent systems (simulation box and atoms), use \sphinxstyleemphasis{add no} and
\sphinxstyleemphasis{every 1}.  Write a LAMMPS input script which loops over the \sphinxstyleemphasis{N}
systems.  See the {\hyperref[\detokenize{Howto_multiple::doc}]{\sphinxcrossref{\DUrole{doc}{Howto multiple}}}} doc page for
details on looping and removing old systems.  The series of systems
could be initialized by reading them from data files with
{\hyperref[\detokenize{read_data::doc}]{\sphinxcrossref{\DUrole{doc}{read\_data}}}} commands.  Or, for example, by using the
{\hyperref[\detokenize{lattice::doc}]{\sphinxcrossref{\DUrole{doc}{lattice}}}} , {\hyperref[\detokenize{create_atoms::doc}]{\sphinxcrossref{\DUrole{doc}{create\_atoms}}}},
{\hyperref[\detokenize{delete_atoms::doc}]{\sphinxcrossref{\DUrole{doc}{delete\_atoms}}}}, and/or {\hyperref[\detokenize{displace_atoms::doc}]{\sphinxcrossref{\DUrole{doc}{displace\_atoms
random}}}} commands to generate a series of different
systems.  At the end of the loop perform {\hyperref[\detokenize{run::doc}]{\sphinxcrossref{\DUrole{doc}{run 0}}}} and
{\hyperref[\detokenize{write_dump::doc}]{\sphinxcrossref{\DUrole{doc}{write\_dump}}}} commands to invoke the QM code and
output the QM energy and forces.  As in (2) this be useful to produce
QM data for training a machine learning potential.


\bigskip\hrule\bigskip



\subsection{Restart, fix\_modify, output, run start/stop, minimize info}
\label{\detokenize{fix_mdi_qm:restart-fix-modify-output-run-start-stop-minimize-info}}
\sphinxAtStartPar
No information about this fix is written to {\hyperref[\detokenize{restart::doc}]{\sphinxcrossref{\DUrole{doc}{binary restart files}}}}.

\sphinxAtStartPar
The {\hyperref[\detokenize{fix_modify::doc}]{\sphinxcrossref{\DUrole{doc}{fix\_modify}}}} \sphinxstyleemphasis{energy} option is supported by
this fix to add the potential energy computed by the QM code to the
global potential energy of the system as part of {\hyperref[\detokenize{thermo_style::doc}]{\sphinxcrossref{\DUrole{doc}{thermodynamic
output}}}}.  The default setting for this fix is
{\hyperref[\detokenize{fix_modify::doc}]{\sphinxcrossref{\DUrole{doc}{fix\_modify energy yes}}}}, unless the \sphinxstyleemphasis{add} keyword is
set to \sphinxstyleemphasis{no}, in which case the default setting is \sphinxstyleemphasis{no}.

\sphinxAtStartPar
The {\hyperref[\detokenize{fix_modify::doc}]{\sphinxcrossref{\DUrole{doc}{fix\_modify}}}} \sphinxstyleemphasis{virial} option is supported by
this fix to add the contribution computed by the QM code to the global
pressure of the system as part of {\hyperref[\detokenize{thermo_style::doc}]{\sphinxcrossref{\DUrole{doc}{thermodynamic output}}}}.  The default setting for this fix is {\hyperref[\detokenize{fix_modify::doc}]{\sphinxcrossref{\DUrole{doc}{fix\_modify
virial yes}}}}, unless the \sphinxstyleemphasis{add} keyword is set to \sphinxstyleemphasis{no}, in
which case the default setting is \sphinxstyleemphasis{no}.

\sphinxAtStartPar
This fix computes a global scalar which can be accessed by various
{\hyperref[\detokenize{Howto_output::doc}]{\sphinxcrossref{\DUrole{doc}{output commands}}}}.  The scalar is the energy
returned by the QM code.  The scalar value calculated by this fix is
“extensive”.

\sphinxAtStartPar
This fix also computes a global vector with of length 6 which contains
the symmetric virial tensor values returned by the QM code.  It can
likewise be accessed by various {\hyperref[\detokenize{Howto_output::doc}]{\sphinxcrossref{\DUrole{doc}{output commands}}}}.

\sphinxAtStartPar
The ordering of values in the symmetric virial tensor is as follows:
vxx, vyy, vzz, vxy, vxz, vyz.  The values will be in pressure
{\hyperref[\detokenize{units::doc}]{\sphinxcrossref{\DUrole{doc}{units}}}}.

\sphinxAtStartPar
This fix also computes a peratom array with 3 columns which contains
the peratom forces returned by the QM code.  It can likewise be
accessed by various {\hyperref[\detokenize{Howto_output::doc}]{\sphinxcrossref{\DUrole{doc}{output commands}}}}.

\sphinxAtStartPar
No parameter of this fix can be used with the \sphinxstyleemphasis{start/stop} keywords of
the {\hyperref[\detokenize{run::doc}]{\sphinxcrossref{\DUrole{doc}{run}}}} command.

\sphinxAtStartPar
Assuming the \sphinxstyleemphasis{add} keyword is set to \sphinxstyleemphasis{yes} (the default), the forces
computed by the QM code are used during an energy minimization,
invoked by the {\hyperref[\detokenize{minimize::doc}]{\sphinxcrossref{\DUrole{doc}{minimize}}}} command.

\begin{sphinxadmonition}{note}{Note:}
\sphinxAtStartPar
If you want the potential energy associated with the QM forces to
be included in the total potential energy of the system (the
quantity being minimized), you MUST not disable the
{\hyperref[\detokenize{fix_modify::doc}]{\sphinxcrossref{\DUrole{doc}{fix\_modify}}}} \sphinxstyleemphasis{energy} option for this fix, which
means the \sphinxstyleemphasis{add} keyword should also be set to \sphinxstyleemphasis{yes} (the default).
\end{sphinxadmonition}


\subsection{Restrictions}
\label{\detokenize{fix_mdi_qm:restrictions}}
\sphinxAtStartPar
This fix is part of the MDI package.  It is only enabled if LAMMPS
was built with that package.  See the {\hyperref[\detokenize{Build_package::doc}]{\sphinxcrossref{\DUrole{doc}{Build package}}}} page for more info.

\sphinxAtStartPar
To use LAMMPS as an MDI driver in conjunction with other MDI\sphinxhyphen{}enabled
codes (MD or QM codes), the {\hyperref[\detokenize{units::doc}]{\sphinxcrossref{\DUrole{doc}{units}}}} command should be used
to specify \sphinxstyleemphasis{real} or \sphinxstyleemphasis{metal} units.  This will ensure the correct unit
conversions between LAMMPS and MDI units.  The other code will also
perform similar unit conversions into its preferred units.

\sphinxAtStartPar
LAMMPS can also be used as an MDI driver in other unit choices it
supports, e.g. \sphinxstyleemphasis{lj}, but then no unit conversion to MDI units is
performed.

\sphinxAtStartPar
If this fix is used in conjunction with a QM code that does not support
periodic boundary conditions (more specifically, a QM code that does not
support the \sphinxcode{\sphinxupquote{\textgreater{}CELL}} MDI command), the LAMMPS system must be fully
non\sphinxhyphen{}periodic.  I.e. no dimension of the system can be periodic.


\subsection{Related commands}
\label{\detokenize{fix_mdi_qm:related-commands}}
\sphinxAtStartPar
{\hyperref[\detokenize{mdi::doc}]{\sphinxcrossref{\DUrole{doc}{mdi plugin}}}},
{\hyperref[\detokenize{mdi::doc}]{\sphinxcrossref{\DUrole{doc}{mdi engine}}}},
{\hyperref[\detokenize{fix_mdi_qmmm::doc}]{\sphinxcrossref{\DUrole{doc}{fix mdi/qmmm}}}}


\subsection{Default}
\label{\detokenize{fix_mdi_qm:default}}
\sphinxAtStartPar
The default for the optional keywords are virial = no, add = yes,
every = 1, connect = yes.

\sphinxstepscope

\index{fix mdi/qmmm@\spxentry{fix mdi/qmmm}}\ignorespaces 

\section{fix mdi/qmmm command}
\label{\detokenize{fix_mdi_qmmm:fix-mdi-qmmm-command}}\label{\detokenize{fix_mdi_qmmm:index-0}}\label{\detokenize{fix_mdi_qmmm::doc}}

\subsection{Syntax}
\label{\detokenize{fix_mdi_qmmm:syntax}}
\begin{sphinxVerbatim}[commandchars=\\\{\}]
\PYG{k}{fix }\PYG{n+nv+nvIdentifier}{ID}\PYG{+w}{ }\PYG{n+nv+nvIdentifier}{group\PYGZhy{}ID}\PYG{+w}{ }\PYG{n}{mdi}\PYG{o}{/}\PYG{n}{qmmm}\PYG{+w}{ }\PYG{n}{mode}\PYG{+w}{ }\PYG{n}{keyword}\PYG{+w}{ }\PYG{n}{value}\PYG{n+nv}{(}\PYG{n+nv}{s}\PYG{n+nv}{)}\PYG{+w}{ }\PYG{n}{keyword}\PYG{+w}{ }\PYG{n}{value}\PYG{n+nv}{(}\PYG{n+nv}{s}\PYG{n+nv}{)}\PYG{+w}{ }\PYG{n}{...}
\end{sphinxVerbatim}
\begin{itemize}
\item {} 
\sphinxAtStartPar
ID, group\sphinxhyphen{}ID are documented in {\hyperref[\detokenize{fix::doc}]{\sphinxcrossref{\DUrole{doc}{fix}}}} command

\item {} 
\sphinxAtStartPar
mdi/qmmm = style name of this fix command

\item {} 
\sphinxAtStartPar
mode = \sphinxstyleemphasis{direct} or \sphinxstyleemphasis{potential}

\item {} 
\sphinxAtStartPar
zero or more keyword/value pairs may be appended

\item {} 
\sphinxAtStartPar
keyword = \sphinxstyleemphasis{virial} or \sphinxstyleemphasis{add} or \sphinxstyleemphasis{every} or \sphinxstyleemphasis{connect} or \sphinxstyleemphasis{elements}
\begin{sphinxalltt}
\sphinxstyleemphasis{virial} args = \sphinxstyleemphasis{yes} or \sphinxstyleemphasis{no}
  yes = request virial tensor from server code
  no = do not request virial tensor from server code
\sphinxstyleemphasis{connect} args = \sphinxstyleemphasis{yes} or \sphinxstyleemphasis{no}
  yes = perform a one\sphinxhyphen{}time connection to the MDI engine code
  no = do not perform the connection operation
\sphinxstyleemphasis{elements} args = N\_1 N\_2 ... N\_ntypes
  N\_1,N\_2,...N\_ntypes = chemical symbol for each of ntypes LAMMPS atom types
\end{sphinxalltt}

\end{itemize}


\subsection{Examples}
\label{\detokenize{fix_mdi_qmmm:examples}}
\begin{sphinxVerbatim}[commandchars=\\\{\}]
\PYG{k}{fix }\PYG{n+nv+nvIdentifier}{1}\PYG{+w}{ }\PYG{n+nv+nvIdentifier}{all}\PYG{+w}{ }\PYG{n}{mdi}\PYG{o}{/}\PYG{n}{qmmm}\PYG{+w}{ }\PYG{n}{direct}
\PYG{k}{fix }\PYG{n+nv+nvIdentifier}{1}\PYG{+w}{ }\PYG{n+nv+nvIdentifier}{all}\PYG{+w}{ }\PYG{n}{mdi}\PYG{o}{/}\PYG{n}{qmmm}\PYG{+w}{ }\PYG{n}{potential}\PYG{+w}{ }\PYG{n}{virial}\PYG{+w}{ }\PYG{n}{yes}
\PYG{k}{fix }\PYG{n+nv+nvIdentifier}{1}\PYG{+w}{ }\PYG{n+nv+nvIdentifier}{all}\PYG{+w}{ }\PYG{n}{mdi}\PYG{o}{/}\PYG{n}{qmmm}\PYG{+w}{ }\PYG{n}{potential}\PYG{+w}{ }\PYG{n}{virial}\PYG{+w}{ }\PYG{n}{yes}\PYG{+w}{ }\PYG{n}{elements}\PYG{+w}{ }\PYG{l+m}{13}\PYG{+w}{ }\PYG{l+m}{29}
\end{sphinxVerbatim}


\subsection{Description}
\label{\detokenize{fix_mdi_qmmm:description}}
\sphinxAtStartPar
\DUrole{versionmodified,added}{Added in version 28Mar2023.}

\sphinxAtStartPar
This command enables LAMMPS to act as a client with another server code
to perform a coupled QM/MM (quantum\sphinxhyphen{}mechanics/molecular\sphinxhyphen{}mechanics)
simulation.  LAMMPS will perform classical MD (molecular mechanics
or MM) for the (typically larger) MM portion of the system.  A quantum
mechanics code will calculate quantum energy and forces for the QM
portion of the system.  The two codes work together to calculate the
energy and forces due to the cross interactions between QM and MM atoms.
The QM server code must support use of the \sphinxhref{https://molssi-mdi.github.io/MDI\_Library/html/index.html}{MDI Library} as
explained below.

\sphinxAtStartPar
The partitioning of the system between QM and MM atoms is as follows.
Atoms in the specified group are QM atoms; the remaining atoms are MM
atoms.  The input script should thus define this partitioning.
See additional information below about other requirements for an input
script to use this fix and perform a QM/MM simulation.

\sphinxAtStartPar
The code coupling performed by this command is done via the \sphinxhref{https://molssi-mdi.github.io/MDI\_Library/html/index.html}{MDI
Library}.
LAMMPS runs as an MDI driver (client), and sends MDI commands to an
external MDI engine code (server), in this case a QM code which has
support for MDI.  See the {\hyperref[\detokenize{Howto_mdi::doc}]{\sphinxcrossref{\DUrole{doc}{Howto mdi}}}} page for more
information about how LAMMPS can operate as either an MDI driver or
engine.

\sphinxAtStartPar
The \sphinxcode{\sphinxupquote{examples/QUANTUM}} directory has sub\sphinxhyphen{}directories with example
input scripts using this fix in tandem with different QM codes.  The
README files in the sub\sphinxhyphen{}directories explain how to download and build
the various QM codes.  They also explain how to launch LAMMPS and the QM
code so that they communicate via the MDI library using either MPI or
sockets.  Any QM code that supports MDI could be used in addition to
those discussed in the sub\sphinxhyphen{}directories.  See the {\hyperref[\detokenize{Howto_mdi::doc}]{\sphinxcrossref{\DUrole{doc}{Howto mdi}}}} page for a current list (March 2022) of such QM codes.

\sphinxAtStartPar
Note that an engine code can support MDI in either or both of two modes.
It can be used as a stand\sphinxhyphen{}alone code, launched at the same time as
LAMMPS.  Or it can be used as a plugin library, which LAMMPS loads.  See
the {\hyperref[\detokenize{mdi::doc}]{\sphinxcrossref{\DUrole{doc}{mdi plugin}}}} command for how to trigger LAMMPS to load a
plugin library.  The \sphinxcode{\sphinxupquote{examples/QUANTUM}} sub\sphinxhyphen{}directory README files
explains how to launch the two codes in either mode.


\bigskip\hrule\bigskip


\sphinxAtStartPar
The \sphinxstyleemphasis{mode} setting determines which QM/MM coupling algorithm is used.
LAMMPS currently supports \sphinxstyleemphasis{direct} and \sphinxstyleemphasis{potential} algorithms, based
on the \sphinxstyleemphasis{mode} setting.  Both algorithms should give reasonably
accurate results, but some QM codes support only one of the two modes.
E.g. in the \sphinxcode{\sphinxupquote{examples/QUANTUM}} directory, PySCF supports only \sphinxstyleemphasis{direct},
NWChem supports only \sphinxstyleemphasis{potential}, and LATTE currently supports
neither, so it cannot be used for QM/MM simulations using this fix.

\sphinxAtStartPar
The \sphinxstyleemphasis{direct} option passes the coordinates and charges of each MM atom
to the quantum code, in addition to the coordinates of each QM atom.
The quantum code returns forces on each QM atom as well as forces on
each MM atom.  The latter is effectively the force on MM atoms due to
the QM atoms.

\sphinxAtStartPar
The input script for performing a \sphinxstyleemphasis{direct} mode QM/MM simulation should
do the following:
\begin{itemize}
\item {} 
\sphinxAtStartPar
delete all bonds (angles, dihedrals, etc) between QM atoms

\item {} 
\sphinxAtStartPar
set the charge on each QM atom to zero

\item {} 
\sphinxAtStartPar
define no bonds (angles, dihedrals, etc) which involve both QM and MM atoms

\item {} 
\sphinxAtStartPar
define a force field (pair, bonds, angles, optional kspace) for the entire system

\end{itemize}

\sphinxAtStartPar
The first two bullet can be performed using the {\hyperref[\detokenize{delete_bonds::doc}]{\sphinxcrossref{\DUrole{doc}{delete\_bonds}}}} and {\hyperref[\detokenize{set::doc}]{\sphinxcrossref{\DUrole{doc}{set}}}} commands.

\sphinxAtStartPar
The third bullet is required to have a consistent model, but is not
checked by LAMMPS.

\sphinxAtStartPar
The fourth bullet implies that non\sphinxhyphen{}bonded non\sphinxhyphen{}Coulombic interactions
(e.g. van der Waals) between QM/QM and QM/MM pairs of atoms are
computed by LAMMPS.

\sphinxAtStartPar
See the \sphinxcode{\sphinxupquote{examples/QUANTUM/PySCF/in.*}} files for examples of input
scripts for QM/MM simulations using the \sphinxstyleemphasis{direct} mode.

\sphinxAtStartPar
The \sphinxstyleemphasis{potential} option passes the coordinates of each QM atom and a
Coulomb potential for each QM atom to the quantum code.  The latter is
calculated by performing a Coulombics\sphinxhyphen{}only calculation for the entire
system, subtracting all QM/QM pairwise Coulombic terms, and dividing
the Coulomb energy on each QM atom by the charge of the QM atom.  The
potential value represents the Coulombic influence of all the MM atoms
on each QM atom.

\sphinxAtStartPar
The quantum code returns forces and charge on each QM atom.  The new
charges on the QM atom are used to re\sphinxhyphen{}calculate the MM force field,
resulting in altered forces on the MM atoms.

\sphinxAtStartPar
The input script for performing a \sphinxstyleemphasis{potential} mode QM/MM simulation
should do the following:
\begin{itemize}
\item {} 
\sphinxAtStartPar
delete all bonds (angles, dihedrals, etc) between QM atoms

\item {} 
\sphinxAtStartPar
define a hybrid pair style which includes a Coulomb\sphinxhyphen{}only pair sub\sphinxhyphen{}style

\item {} 
\sphinxAtStartPar
define no bonds (angles, dihedrals, etc) which involve both QM and MM atoms

\item {} 
\sphinxAtStartPar
define a force field (pair, bonds, angles, optional kspace) for the entire system

\end{itemize}

\sphinxAtStartPar
The first operation can be performed using the {\hyperref[\detokenize{delete_bonds::doc}]{\sphinxcrossref{\DUrole{doc}{delete\_bonds}}}} command.  See the \sphinxcode{\sphinxupquote{examples/QUANTUM/NWChem/in.*}} files
for examples of how to do this.

\sphinxAtStartPar
The second operation is necessary so that this fix can calculate the
Coulomb potential for the QM atoms.

\sphinxAtStartPar
The third bullet is required to have a consistent model, but is not
checked by LAMMPS.

\sphinxAtStartPar
The fourth bullet implies that non\sphinxhyphen{}bonded non\sphinxhyphen{}Coulombic interactions
(e.g. van der Waals) between QM/QM and QM/MM pairs of atoms are computed
by LAMMPS.  However, some QM codes do not want the MM code (LAMMPS) to
compute QM/QM van der Waals interactions.  NWChem is an example.  In
this case, the coefficients for those interactions need to be turned
off, which typically requires the atom types for the QM atoms be
different than those for the MM atoms.

\sphinxAtStartPar
See the \sphinxcode{\sphinxupquote{examples/QUANTUM/NWChem/in.*}} files for examples of input
scripts for QM/MM simulations using the \sphinxstyleemphasis{potential} mode.  Those scripts
also illustrate how to turn off QM/QM van der Waals interactions.


\bigskip\hrule\bigskip


\sphinxAtStartPar
The \sphinxstyleemphasis{virial} keyword setting of yes or no determines whether LAMMPS
will request the QM code to also compute and return the QM
contribution to a stress tensor for the system which LAMMPS will
convert to a 6\sphinxhyphen{}element symmetric virial tensor.

\sphinxAtStartPar
The \sphinxstyleemphasis{connect} keyword determines whether this fix performs a one\sphinxhyphen{}time
connection to the QM code.  The default is \sphinxstyleemphasis{yes}.  The only time a
\sphinxstyleemphasis{no} is needed is if this command is used multiple times in an input
script.  E.g. if it used inside a loop which also uses the {\hyperref[\detokenize{clear::doc}]{\sphinxcrossref{\DUrole{doc}{clear}}}} command to destroy the system (including this fix).  As
example would be a script which loop over a series of independent QM/MM
simulations, e.g. each with their own data file.  In this use case
\sphinxstyleemphasis{connect no} could be used along with the {\hyperref[\detokenize{mdi::doc}]{\sphinxcrossref{\DUrole{doc}{mdi connect and exit}}}} command to one\sphinxhyphen{}time initiate/terminate the connection outside
the loop.

\sphinxAtStartPar
The \sphinxstyleemphasis{elements} keyword allows specification of what element each
LAMMPS atom type corresponds to.  This is specified by the chemical
symbol of the element, e.g. C or Al or Si.  A symbol must be specified
for each of the ntypes LAMMPS atom types.  Multiple LAMMPS types can
represent the same element.  Ntypes is typically specified via the
{\hyperref[\detokenize{create_box::doc}]{\sphinxcrossref{\DUrole{doc}{create\_box}}}} command or in the data file read by the
{\hyperref[\detokenize{read_data::doc}]{\sphinxcrossref{\DUrole{doc}{read\_data}}}} command.

\sphinxAtStartPar
If this keyword is specified, then this fix will send the MDI
“\textgreater{}ELEMENTS” command to the engine, to insure the two codes are
consistent in their definition of atomic species.  If this keyword is
not specified, then this fix will send the MDI \textgreater{}TYPES command to the
engine.  This is fine if both the LAMMPS driver and the MDI engine are
initialized so that the atom type values are consistent in both codes.


\bigskip\hrule\bigskip



\subsection{Restart, fix\_modify, output, run start/stop, minimize info}
\label{\detokenize{fix_mdi_qmmm:restart-fix-modify-output-run-start-stop-minimize-info}}
\sphinxAtStartPar
No information about this fix is written to {\hyperref[\detokenize{restart::doc}]{\sphinxcrossref{\DUrole{doc}{binary restart files}}}}.

\sphinxAtStartPar
The {\hyperref[\detokenize{fix_modify::doc}]{\sphinxcrossref{\DUrole{doc}{fix\_modify}}}} \sphinxstyleemphasis{energy} option is supported by
this fix to add the potential energy computed by the QM code to the
global potential energy of the system as part of {\hyperref[\detokenize{thermo_style::doc}]{\sphinxcrossref{\DUrole{doc}{thermodynamic
output}}}}.  The default setting for this fix is
{\hyperref[\detokenize{fix_modify::doc}]{\sphinxcrossref{\DUrole{doc}{fix\_modify energy yes}}}}.

\sphinxAtStartPar
The {\hyperref[\detokenize{fix_modify::doc}]{\sphinxcrossref{\DUrole{doc}{fix\_modify}}}} \sphinxstyleemphasis{virial} option is supported by
this fix to add the contribution computed by the QM code to the global
pressure of the system as part of {\hyperref[\detokenize{thermo_style::doc}]{\sphinxcrossref{\DUrole{doc}{thermodynamic output}}}}.  The default setting for this fix is {\hyperref[\detokenize{fix_modify::doc}]{\sphinxcrossref{\DUrole{doc}{fix\_modify
virial yes}}}}.

\sphinxAtStartPar
This fix computes a global scalar which can be accessed by various
{\hyperref[\detokenize{Howto_output::doc}]{\sphinxcrossref{\DUrole{doc}{output commands}}}}.  The scalar is the energy
returned by the QM code.  The scalar value calculated by this fix is
“extensive”.

\sphinxAtStartPar
This fix also computes a global vector with of length 6 which contains
the symmetric virial tensor values returned by the QM code.  It can
likewise be accessed by various {\hyperref[\detokenize{Howto_output::doc}]{\sphinxcrossref{\DUrole{doc}{output commands}}}}.

\sphinxAtStartPar
The ordering of values in the symmetric virial tensor is as follows:
vxx, vyy, vzz, vxy, vxz, vyz.  The values will be in pressure
{\hyperref[\detokenize{units::doc}]{\sphinxcrossref{\DUrole{doc}{units}}}}.

\sphinxAtStartPar
This fix also computes a peratom array with 3 columns which contains
the peratom forces returned by the QM code.  It can likewise be
accessed by various {\hyperref[\detokenize{Howto_output::doc}]{\sphinxcrossref{\DUrole{doc}{output commands}}}}.  Note that
for \sphinxstyleemphasis{direct} mode this will be quantum forces on both QM and MM atoms.
For \sphinxstyleemphasis{potential} mode it will only be quantum forces on QM atoms; the
forces for MM atoms will be zero.

\sphinxAtStartPar
No parameter of this fix can be used with the \sphinxstyleemphasis{start/stop} keywords of
the {\hyperref[\detokenize{run::doc}]{\sphinxcrossref{\DUrole{doc}{run}}}} command.

\sphinxAtStartPar
The forces computed by the QM code are used during an energy
minimization, invoked by the {\hyperref[\detokenize{minimize::doc}]{\sphinxcrossref{\DUrole{doc}{minimize}}}} command.

\begin{sphinxadmonition}{note}{Note:}
\sphinxAtStartPar
If you want the potential energy associated with the QM forces to
be included in the total potential energy of the system (the
quantity being minimized), you MUST not disable the
{\hyperref[\detokenize{fix_modify::doc}]{\sphinxcrossref{\DUrole{doc}{fix\_modify}}}} \sphinxstyleemphasis{energy} option for this fix.
\end{sphinxadmonition}


\subsection{Restrictions}
\label{\detokenize{fix_mdi_qmmm:restrictions}}
\sphinxAtStartPar
This command is part of the MDI package.  It is only enabled if LAMMPS
was built with that package.  See the {\hyperref[\detokenize{Build_package::doc}]{\sphinxcrossref{\DUrole{doc}{Build package}}}} page for more info.

\sphinxAtStartPar
To use LAMMPS as an MDI driver in conjunction with other MDI\sphinxhyphen{}enabled
codes (MD or QM codes), the {\hyperref[\detokenize{units::doc}]{\sphinxcrossref{\DUrole{doc}{units}}}} command should be used
to specify \sphinxstyleemphasis{real} or \sphinxstyleemphasis{metal} units.  This will ensure the correct unit
conversions between LAMMPS and MDI units.  The other code will also
perform similar unit conversions into its preferred units.

\sphinxAtStartPar
If this fix is used in conjunction with a QM code that does not support
periodic boundary conditions (more specifically, a QM code that does not
support the \sphinxcode{\sphinxupquote{\textgreater{}CELL}} MDI command), the LAMMPS system must be fully
non\sphinxhyphen{}periodic.  I.e. no dimension of the system can be periodic.


\subsection{Related commands}
\label{\detokenize{fix_mdi_qmmm:related-commands}}
\sphinxAtStartPar
{\hyperref[\detokenize{mdi::doc}]{\sphinxcrossref{\DUrole{doc}{mdi plugin}}}},
{\hyperref[\detokenize{mdi::doc}]{\sphinxcrossref{\DUrole{doc}{mdi engine}}}},
{\hyperref[\detokenize{fix_mdi_qm::doc}]{\sphinxcrossref{\DUrole{doc}{fix mdi/qm}}}}


\subsection{Default}
\label{\detokenize{fix_mdi_qmmm:default}}
\sphinxAtStartPar
The default for the optional keywords are virial = no and connect = yes.

\sphinxstepscope

\index{fix meso/move@\spxentry{fix meso/move}}\ignorespaces 

\section{fix meso/move command}
\label{\detokenize{fix_meso_move:fix-meso-move-command}}\label{\detokenize{fix_meso_move:index-0}}\label{\detokenize{fix_meso_move::doc}}

\subsection{Syntax}
\label{\detokenize{fix_meso_move:syntax}}
\begin{sphinxVerbatim}[commandchars=\\\{\}]
\PYG{k}{fix }\PYG{n+nv+nvIdentifier}{ID}\PYG{+w}{ }\PYG{n+nv+nvIdentifier}{group\PYGZhy{}ID}\PYG{+w}{ }\PYG{n}{meso}\PYG{o}{/}\PYG{n}{move}\PYG{+w}{ }\PYG{n}{style}\PYG{+w}{ }\PYG{n}{args}\PYG{+w}{ }\PYG{n}{keyword}\PYG{+w}{ }\PYG{n}{values}\PYG{+w}{ }\PYG{n}{...}
\end{sphinxVerbatim}
\begin{itemize}
\item {} 
\sphinxAtStartPar
ID, group\sphinxhyphen{}ID are documented in {\hyperref[\detokenize{fix::doc}]{\sphinxcrossref{\DUrole{doc}{fix}}}} command

\item {} 
\sphinxAtStartPar
meso/move = style name of this fix command

\item {} 
\sphinxAtStartPar
style = \sphinxstyleemphasis{linear} or \sphinxstyleemphasis{wiggle} or \sphinxstyleemphasis{rotate} or \sphinxstyleemphasis{variable}
\begin{sphinxalltt}
\sphinxstyleemphasis{linear} args = Vx Vy Vz
  Vx,Vy,Vz = components of velocity vector (velocity units), any component can be specified as NULL
\sphinxstyleemphasis{wiggle} args = Ax Ay Az period
  Ax,Ay,Az = components of amplitude vector (distance units), any component can be specified as NULL
  period = period of oscillation (time units)
\sphinxstyleemphasis{rotate} args = Px Py Pz Rx Ry Rz period
  Px,Py,Pz = origin point of axis of rotation (distance units)
  Rx,Ry,Rz = axis of rotation vector
  period = period of rotation (time units)
\sphinxstyleemphasis{variable} args = v\_dx v\_dy v\_dz v\_vx v\_vy v\_vz
  v\_dx,v\_dy,v\_dz = 3 variable names that calculate x,y,z displacement as function of time, any component can be specified as NULL
  v\_vx,v\_vy,v\_vz = 3 variable names that calculate x,y,z velocity as function of time, any component can be specified as NULL
\end{sphinxalltt}

\item {} 
\sphinxAtStartPar
zero or more keyword/value pairs may be appended

\item {} 
\sphinxAtStartPar
keyword = \sphinxstyleemphasis{units}
\begin{sphinxalltt}
\sphinxstyleemphasis{units} value = \sphinxstyleemphasis{box} or \sphinxstyleemphasis{lattice}
\end{sphinxalltt}

\end{itemize}


\subsection{Examples}
\label{\detokenize{fix_meso_move:examples}}
\begin{sphinxVerbatim}[commandchars=\\\{\}]
\PYG{k}{fix }\PYG{n+nv+nvIdentifier}{1}\PYG{+w}{ }\PYG{n+nv+nvIdentifier}{boundary}\PYG{+w}{ }\PYG{n}{meso}\PYG{o}{/}\PYG{n}{move}\PYG{+w}{ }\PYG{n}{wiggle}\PYG{+w}{ }\PYG{l+m}{3.0}\PYG{+w}{ }\PYG{l+m}{0.0}\PYG{+w}{ }\PYG{l+m}{0.0}\PYG{+w}{ }\PYG{l+m}{1.0}\PYG{+w}{ }\PYG{n}{units}\PYG{+w}{ }\PYG{n}{box}
\PYG{k}{fix }\PYG{n+nv+nvIdentifier}{2}\PYG{+w}{ }\PYG{n+nv+nvIdentifier}{boundary}\PYG{+w}{ }\PYG{n}{meso}\PYG{o}{/}\PYG{n}{move}\PYG{+w}{ }\PYG{n}{rotate}\PYG{+w}{ }\PYG{l+m}{0.0}\PYG{+w}{ }\PYG{l+m}{0.0}\PYG{+w}{ }\PYG{l+m}{0.0}\PYG{+w}{ }\PYG{l+m}{0.0}\PYG{+w}{ }\PYG{l+m}{0.0}\PYG{+w}{ }\PYG{l+m}{1.0}\PYG{+w}{ }\PYG{l+m}{5.0}
\PYG{k}{fix }\PYG{n+nv+nvIdentifier}{2}\PYG{+w}{ }\PYG{n+nv+nvIdentifier}{boundary}\PYG{+w}{ }\PYG{n}{meso}\PYG{o}{/}\PYG{n}{move}\PYG{+w}{ }\PYG{k}{variable }\PYG{n+nv+nvIdentifier}{v\PYGZus{}myx}\PYG{+w}{ }\PYG{n}{v\PYGZus{}myy}\PYG{+w}{ }\PYG{n}{NULL}\PYG{+w}{ }\PYG{n}{v\PYGZus{}VX}\PYG{+w}{ }\PYG{n}{v\PYGZus{}VY}\PYG{+w}{ }\PYG{n}{NULL}
\end{sphinxVerbatim}


\subsection{Description}
\label{\detokenize{fix_meso_move:description}}
\sphinxAtStartPar
Perform updates of position, velocity, internal energy and local
density for mesoscopic particles in the group each timestep using the
specified settings or formulas, without regard to forces on the
particles. This can be useful for boundary, solid bodies or other
particles, whose movement can influence nearby particles.

\sphinxAtStartPar
The operation of this fix is exactly like that described by the
{\hyperref[\detokenize{fix_move::doc}]{\sphinxcrossref{\DUrole{doc}{fix move}}}} command, except that particles’ density,
internal energy and extrapolated velocity are also updated.

\begin{sphinxadmonition}{note}{Note:}
\sphinxAtStartPar
The particles affected by this fix should not be time integrated by
other fixes (e.g. {\hyperref[\detokenize{fix_sph::doc}]{\sphinxcrossref{\DUrole{doc}{fix sph}}}}, {\hyperref[\detokenize{fix_sph_stationary::doc}]{\sphinxcrossref{\DUrole{doc}{fix
sph/stationary}}}}), since that will change their
positions and velocities twice.
\end{sphinxadmonition}

\begin{sphinxadmonition}{note}{Note:}
\sphinxAtStartPar
As particles move due to this fix, they will pass through periodic
boundaries and be remapped to the other side of the simulation box,
just as they would during normal time integration (e.g. via the
{\hyperref[\detokenize{fix_sph::doc}]{\sphinxcrossref{\DUrole{doc}{fix sph}}}} command).  It is up to you to decide
whether periodic boundaries are appropriate with the kind of
particle motion you are prescribing with this fix.
\end{sphinxadmonition}

\begin{sphinxadmonition}{note}{Note:}
\sphinxAtStartPar
As discussed below, particles are moved relative to their initial
position at the time the fix is specified.  These initial coordinates
are stored by the fix in “unwrapped” form, by using the image flags
associated with each particle.  See the {\hyperref[\detokenize{dump::doc}]{\sphinxcrossref{\DUrole{doc}{dump custom}}}} command
for a discussion of “unwrapped” coordinates.  See the Atoms section of
the {\hyperref[\detokenize{read_data::doc}]{\sphinxcrossref{\DUrole{doc}{read\_data}}}} command for a discussion of image flags
and how they are set for each particle.  You can reset the image flags
(e.g. to 0) before invoking this fix by using the {\hyperref[\detokenize{set::doc}]{\sphinxcrossref{\DUrole{doc}{set image}}}}
command.
\end{sphinxadmonition}


\bigskip\hrule\bigskip


\sphinxAtStartPar
The \sphinxstyleemphasis{linear} style moves particles at a constant velocity, so that their
position \sphinxstyleemphasis{X} = (x,y,z) as a function of time is given in vector
notation as
\begin{sphinxalltt}
X(t) = X0 + V * delta
\end{sphinxalltt}

\sphinxAtStartPar
where \sphinxstyleemphasis{X0} = (x0,y0,z0) is their position at the time the fix is
specified, \sphinxstyleemphasis{V} is the specified velocity vector with components
(Vx,Vy,Vz), and \sphinxstyleemphasis{delta} is the time elapsed since the fix was
specified.  This style also sets the velocity of each particle to V =
(Vx,Vy,Vz).  If any of the velocity components is specified as NULL,
then the position and velocity of that component is time integrated
the same as the {\hyperref[\detokenize{fix_sph::doc}]{\sphinxcrossref{\DUrole{doc}{fix sph}}}} command would perform, using
the corresponding force component on the particle.

\sphinxAtStartPar
Note that the \sphinxstyleemphasis{linear} style is identical to using the \sphinxstyleemphasis{variable}
style with an {\hyperref[\detokenize{variable::doc}]{\sphinxcrossref{\DUrole{doc}{equal\sphinxhyphen{}style variable}}}} that uses the
vdisplace() function.  E.g.

\begin{sphinxVerbatim}[commandchars=\\\{\}]
\PYG{k}{variable }\PYG{n+nv+nvIdentifier}{V}\PYG{+w}{ }\PYG{n}{equal}\PYG{+w}{ }\PYG{l+m}{10.0}
\PYG{k}{variable }\PYG{n+nv+nvIdentifier}{x}\PYG{+w}{ }\PYG{n}{equal}\PYG{+w}{ }\PYG{n}{vdisplace}\PYG{n+nv}{(}\PYG{n+nv}{0.0,\PYGZdl{}V}\PYG{n+nv}{)}
\PYG{k}{fix }\PYG{n+nv+nvIdentifier}{1}\PYG{+w}{ }\PYG{n+nv+nvIdentifier}{boundary}\PYG{+w}{ }\PYG{n}{move}\PYG{+w}{ }\PYG{k}{variable }\PYG{n+nv+nvIdentifier}{v\PYGZus{}x}\PYG{+w}{ }\PYG{n}{NULL}\PYG{+w}{ }\PYG{n}{NULL}\PYG{+w}{ }\PYG{n}{v\PYGZus{}V}\PYG{+w}{ }\PYG{n}{NULL}\PYG{+w}{ }\PYG{n}{NULL}
\end{sphinxVerbatim}

\sphinxAtStartPar
The \sphinxstyleemphasis{wiggle} style moves particles in an oscillatory fashion, so that
their position \sphinxstyleemphasis{X} = (x,y,z) as a function of time is given in vector
notation as
\begin{sphinxalltt}
X(t) = X0 + A sin(omega*delta)
\end{sphinxalltt}

\sphinxAtStartPar
where \sphinxstyleemphasis{X0} = (x0,y0,z0) is their position at the time the fix is
specified, \sphinxstyleemphasis{A} is the specified amplitude vector with components
(Ax,Ay,Az), \sphinxstyleemphasis{omega} is 2 PI / \sphinxstyleemphasis{period}, and \sphinxstyleemphasis{delta} is the time
elapsed since the fix was specified.  This style also sets the
velocity of each particle to the time derivative of this expression.
If any of the amplitude components is specified as NULL, then the
position and velocity of that component is time integrated the same as
the {\hyperref[\detokenize{fix_sph::doc}]{\sphinxcrossref{\DUrole{doc}{fix sph}}}} command would perform, using the
corresponding force component on the particle.

\sphinxAtStartPar
Note that the \sphinxstyleemphasis{wiggle} style is identical to using the \sphinxstyleemphasis{variable}
style with {\hyperref[\detokenize{variable::doc}]{\sphinxcrossref{\DUrole{doc}{equal\sphinxhyphen{}style variables}}}} that use the
swiggle() and cwiggle() functions.  E.g.

\begin{sphinxVerbatim}[commandchars=\\\{\}]
\PYG{k}{variable }\PYG{n+nv+nvIdentifier}{A}\PYG{+w}{ }\PYG{n}{equal}\PYG{+w}{ }\PYG{l+m}{10.0}
\PYG{k}{variable }\PYG{n+nv+nvIdentifier}{T}\PYG{+w}{ }\PYG{n}{equal}\PYG{+w}{ }\PYG{l+m}{5.0}
\PYG{k}{variable }\PYG{n+nv+nvIdentifier}{omega}\PYG{+w}{ }\PYG{n}{equal}\PYG{+w}{ }\PYG{l+m}{2.0}\PYG{o}{*}\PYG{n}{PI}\PYG{o}{/}\PYG{n+nv}{\PYGZdl{}T}
\PYG{k}{variable }\PYG{n+nv+nvIdentifier}{x}\PYG{+w}{ }\PYG{n}{equal}\PYG{+w}{ }\PYG{n}{swiggle}\PYG{n+nv}{(}\PYG{n+nv}{0.0,\PYGZdl{}A,\PYGZdl{}T}\PYG{n+nv}{)}
\PYG{k}{variable }\PYG{n+nv+nvIdentifier}{v}\PYG{+w}{ }\PYG{n}{equal}\PYG{+w}{ }\PYG{n}{v\PYGZus{}omega}\PYG{o}{*}\PYG{n+nv}{(}\PYG{n+nv}{\PYGZdl{}A\PYGZhy{}cwiggle}\PYG{n+nv}{(}\PYG{n+nv}{0.0,\PYGZdl{}A,\PYGZdl{}T}\PYG{n+nv}{)}\PYG{n+nv}{)}
\PYG{k}{fix }\PYG{n+nv+nvIdentifier}{1}\PYG{+w}{ }\PYG{n+nv+nvIdentifier}{boundary}\PYG{+w}{ }\PYG{n}{move}\PYG{+w}{ }\PYG{k}{variable }\PYG{n+nv+nvIdentifier}{v\PYGZus{}x}\PYG{+w}{ }\PYG{n}{NULL}\PYG{+w}{ }\PYG{n}{NULL}\PYG{+w}{ }\PYG{n}{v\PYGZus{}v}\PYG{+w}{ }\PYG{n}{NULL}\PYG{+w}{ }\PYG{n}{NULL}
\end{sphinxVerbatim}

\sphinxAtStartPar
The \sphinxstyleemphasis{rotate} style rotates particles around a rotation axis \sphinxstyleemphasis{R} =
(Rx,Ry,Rz) that goes through a point \sphinxstyleemphasis{P} = (Px,Py,Pz).  The \sphinxstyleemphasis{period} of
the rotation is also specified.  The direction of rotation for the
particles around the rotation axis is consistent with the right\sphinxhyphen{}hand
rule: if your right\sphinxhyphen{}hand thumb points along \sphinxstyleemphasis{R}, then your fingers wrap
around the axis in the direction of rotation.

\sphinxAtStartPar
This style also sets the velocity of each particle to (omega cross
Rperp) where omega is its angular velocity around the rotation axis and
Rperp is a perpendicular vector from the rotation axis to the particle.

\sphinxAtStartPar
The \sphinxstyleemphasis{variable} style allows the position and velocity components of
each particle to be set by formulas specified via the
{\hyperref[\detokenize{variable::doc}]{\sphinxcrossref{\DUrole{doc}{variable}}}} command.  Each of the 6 variables is
specified as an argument to the fix as v\_name, where name is the
variable name that is defined elsewhere in the input script.

\sphinxAtStartPar
Each variable must be of either the \sphinxstyleemphasis{equal} or \sphinxstyleemphasis{atom} style.
\sphinxstyleemphasis{Equal}\sphinxhyphen{}style variables compute a single numeric quantity, that can be
a function of the timestep as well as of other simulation values.
\sphinxstyleemphasis{Atom}\sphinxhyphen{}style variables compute a numeric quantity for each particle, that
can be a function per\sphinxhyphen{}atom quantities, such as the particle’s position, as
well as of the timestep and other simulation values.  Note that this
fix stores the original coordinates of each particle (see note below) so
that per\sphinxhyphen{}atom quantity can be used in an atom\sphinxhyphen{}style variable formula.
See the {\hyperref[\detokenize{variable::doc}]{\sphinxcrossref{\DUrole{doc}{variable}}}} command for details.

\sphinxAtStartPar
The first 3 variables (v\_dx,v\_dy,v\_dz) specified for the \sphinxstyleemphasis{variable}
style are used to calculate a displacement from the particle’s original
position at the time the fix was specified.  The second 3 variables
(v\_vx,v\_vy,v\_vz) specified are used to compute a velocity for each
particle.

\sphinxAtStartPar
Any of the 6 variables can be specified as NULL.  If both the
displacement and velocity variables for a particular x,y,z component
are specified as NULL, then the position and velocity of that
component is time integrated the same as the {\hyperref[\detokenize{fix_sph::doc}]{\sphinxcrossref{\DUrole{doc}{fix sph}}}}
command would perform, using the corresponding force component on the
particle.  If only the velocity variable for a component is specified
as NULL, then the displacement variable will be used to set the
position of the particle, and its velocity component will not be
changed. If only the displacement variable for a component is
specified as NULL, then the velocity variable will be used to set the
velocity of the particle, and the position of the particle will be
time integrated using that velocity.

\sphinxAtStartPar
The \sphinxstyleemphasis{units} keyword determines the meaning of the distance units used
to define the \sphinxstyleemphasis{linear} velocity and \sphinxstyleemphasis{wiggle} amplitude and \sphinxstyleemphasis{rotate}
origin.  This setting is ignored for the \sphinxstyleemphasis{variable} style.  A \sphinxstyleemphasis{box}
value selects standard units as defined by the {\hyperref[\detokenize{units::doc}]{\sphinxcrossref{\DUrole{doc}{units}}}}
command, e.g. velocity in Angstroms/fs and amplitude and position
in Angstroms for units = real.  A \sphinxstyleemphasis{lattice} value means the velocity
units are in lattice spacings per time and the amplitude and position
are in lattice spacings.  The {\hyperref[\detokenize{lattice::doc}]{\sphinxcrossref{\DUrole{doc}{lattice}}}} command must have
been previously used to define the lattice spacing.  Each of these 3
quantities may be dependent on the x,y,z dimension, since the lattice
spacings can be different in x,y,z.


\bigskip\hrule\bigskip



\subsection{Restart, fix\_modify, output, run start/stop, minimize info}
\label{\detokenize{fix_meso_move:restart-fix-modify-output-run-start-stop-minimize-info}}
\sphinxAtStartPar
This fix writes the original coordinates of moving particles to {\hyperref[\detokenize{restart::doc}]{\sphinxcrossref{\DUrole{doc}{binary restart files}}}}, as well as the initial timestep, so that
the motion can be continuous in a restarted simulation.  See the
{\hyperref[\detokenize{read_restart::doc}]{\sphinxcrossref{\DUrole{doc}{read\_restart}}}} command for info on how to re\sphinxhyphen{}specify
a fix in an input script that reads a restart file, so that the
operation of the fix continues in an uninterrupted fashion.

\begin{sphinxadmonition}{note}{Note:}
\sphinxAtStartPar
Because the move positions are a function of the current
timestep and the initial timestep, you cannot reset the timestep to a
different value after reading a restart file, if you expect a fix move
command to work in an uninterrupted fashion.
\end{sphinxadmonition}

\sphinxAtStartPar
None of the {\hyperref[\detokenize{fix_modify::doc}]{\sphinxcrossref{\DUrole{doc}{fix\_modify}}}} options are relevant to this
fix.

\sphinxAtStartPar
This fix produces a per\sphinxhyphen{}atom array which can be accessed by various
{\hyperref[\detokenize{Howto_output::doc}]{\sphinxcrossref{\DUrole{doc}{output commands}}}}.  The number of columns for each
atom is 3, and the columns store the original unwrapped x,y,z coords
of each particle.  The per\sphinxhyphen{}atom values can be accessed on any timestep.

\sphinxAtStartPar
No parameter of this fix can be used with the \sphinxstyleemphasis{start/stop} keywords of
the {\hyperref[\detokenize{run::doc}]{\sphinxcrossref{\DUrole{doc}{run}}}} command.

\sphinxAtStartPar
This fix is not invoked during {\hyperref[\detokenize{minimize::doc}]{\sphinxcrossref{\DUrole{doc}{energy minimization}}}}.


\subsection{Restrictions}
\label{\detokenize{fix_meso_move:restrictions}}
\sphinxAtStartPar
This fix is part of the DPD\sphinxhyphen{}SMOOTH package.  It is only enabled if
LAMMPS was built with that package. See the {\hyperref[\detokenize{Build_package::doc}]{\sphinxcrossref{\DUrole{doc}{Build package}}}} page for more info.

\sphinxAtStartPar
This fix requires that atoms store density and internal energy as
defined by the {\hyperref[\detokenize{atom_style::doc}]{\sphinxcrossref{\DUrole{doc}{atom\_style sph}}}} command.

\sphinxAtStartPar
All particles in the group must be mesoscopic SPH/SDPD particles.

\sphinxAtStartPar
\DUrole{versionmodified,changed}{Changed in version 29Aug2024.}

\sphinxAtStartPar
This fix is incompatible with deformation controls that remap velocity,
for instance the \sphinxstyleemphasis{remap v} option of {\hyperref[\detokenize{fix_deform::doc}]{\sphinxcrossref{\DUrole{doc}{fix deform}}}}.


\subsection{Related commands}
\label{\detokenize{fix_meso_move:related-commands}}
\sphinxAtStartPar
{\hyperref[\detokenize{fix_move::doc}]{\sphinxcrossref{\DUrole{doc}{fix move}}}}, {\hyperref[\detokenize{fix_sph::doc}]{\sphinxcrossref{\DUrole{doc}{fix sph}}}},
{\hyperref[\detokenize{displace_atoms::doc}]{\sphinxcrossref{\DUrole{doc}{displace\_atoms}}}}


\subsection{Default}
\label{\detokenize{fix_meso_move:default}}
\sphinxAtStartPar
The option default is units = lattice.

\sphinxstepscope


\section{fix\_modify AtC commands}
\label{\detokenize{fix_modify_atc_commands:fix-modify-atc-commands}}\label{\detokenize{fix_modify_atc_commands::doc}}
\sphinxstepscope

\index{fix\_modify AtC add\_molecule@\spxentry{fix\_modify AtC add\_molecule}}\ignorespaces 

\subsection{fix\_modify AtC add\_molecule command}
\label{\detokenize{atc_add_molecule:fix-modify-atc-add-molecule-command}}\label{\detokenize{atc_add_molecule:index-0}}\label{\detokenize{atc_add_molecule::doc}}

\subsubsection{Syntax}
\label{\detokenize{atc_add_molecule:syntax}}
\begin{sphinxVerbatim}[commandchars=\\\{\}]
\PYG{k}{fix\PYGZus{}modify }\PYG{o}{\PYGZlt{}}\PYG{n}{AtC}\PYG{+w}{ }\PYG{n}{fixID}\PYG{o}{\PYGZgt{}}\PYG{+w}{ }\PYG{n}{add\PYGZus{}molecule}\PYG{+w}{ }\PYG{o}{\PYGZlt{}}\PYG{n}{small}\PYG{o}{|}\PYG{n}{large}\PYG{o}{\PYGZgt{}}\PYG{+w}{ }\PYG{o}{\PYGZlt{}}\PYG{n}{tag}\PYG{o}{\PYGZgt{}}\PYG{+w}{ }\PYG{o}{\PYGZlt{}}\PYG{n}{group}\PYG{o}{\PYGZhy{}}\PYG{n}{ID}\PYG{o}{\PYGZgt{}}
\end{sphinxVerbatim}
\begin{itemize}
\item {} 
\sphinxAtStartPar
AtC fixID = ID of {\hyperref[\detokenize{fix_atc::doc}]{\sphinxcrossref{\DUrole{doc}{fix atc}}}} instance

\item {} 
\sphinxAtStartPar
add\_molecule = name of the AtC sub\sphinxhyphen{}command

\item {} 
\sphinxAtStartPar
\sphinxstyleemphasis{small} or \sphinxstyleemphasis{large} = can be \sphinxstyleemphasis{small} if molecule size \textless{} cutoff radius, must be \sphinxstyleemphasis{large} otherwise

\item {} 
\sphinxAtStartPar
tag = tag for tracking a molecule

\item {} 
\sphinxAtStartPar
\sphinxstyleemphasis{group\sphinxhyphen{}ID} = LAMMPS defined group\sphinxhyphen{}ID

\end{itemize}


\subsubsection{Examples}
\label{\detokenize{atc_add_molecule:examples}}
\begin{sphinxVerbatim}[commandchars=\\\{\}]
\PYG{k}{group }\PYG{n+nv+nvIdentifier}{WATERGROUP}\PYG{+w}{ }\PYG{n}{type}\PYG{+w}{ }\PYG{l+m}{1}\PYG{+w}{ }\PYG{l+m}{2}
\PYG{k}{fix\PYGZus{}modify }\PYG{n+nv+nvIdentifier}{AtC}\PYG{+w}{ }\PYG{n}{add\PYGZus{}molecule}\PYG{+w}{ }\PYG{n}{small}\PYG{+w}{ }\PYG{n}{water}\PYG{+w}{ }\PYG{n}{WATERGROUP}
\end{sphinxVerbatim}


\subsubsection{Description}
\label{\detokenize{atc_add_molecule:description}}
\sphinxAtStartPar
Associates a tag with all molecules corresponding to a specified group.


\subsubsection{Restrictions}
\label{\detokenize{atc_add_molecule:restrictions}}
\sphinxAtStartPar
None.


\subsubsection{Related AtC commands}
\label{\detokenize{atc_add_molecule:related-atc-commands}}\begin{itemize}
\item {} 
\sphinxAtStartPar
{\hyperref[\detokenize{fix_atc:atc-fix-modify}]{\sphinxcrossref{\DUrole{std,std-ref}{fix\_modify AtC command overview}}}}

\item {} 
\sphinxAtStartPar
{\hyperref[\detokenize{atc_add_species::doc}]{\sphinxcrossref{\DUrole{doc}{fix\_modify AtC add\_species}}}}

\item {} 
\sphinxAtStartPar
{\hyperref[\detokenize{atc_remove_species::doc}]{\sphinxcrossref{\DUrole{doc}{fix\_modify AtC remove\_species}}}}

\item {} 
\sphinxAtStartPar
{\hyperref[\detokenize{atc_remove_molecule::doc}]{\sphinxcrossref{\DUrole{doc}{fix\_modify AtC remove\_molecule}}}}

\end{itemize}


\subsubsection{Default}
\label{\detokenize{atc_add_molecule:default}}
\sphinxAtStartPar
None.

\sphinxstepscope

\index{fix\_modify AtC add\_species@\spxentry{fix\_modify AtC add\_species}}\ignorespaces 

\subsection{fix\_modify AtC add\_species command}
\label{\detokenize{atc_add_species:fix-modify-atc-add-species-command}}\label{\detokenize{atc_add_species:index-0}}\label{\detokenize{atc_add_species::doc}}

\subsubsection{Syntax}
\label{\detokenize{atc_add_species:syntax}}
\begin{sphinxVerbatim}[commandchars=\\\{\}]
\PYG{k}{fix\PYGZus{}modify }\PYG{o}{\PYGZlt{}}\PYG{n}{AtC}\PYG{+w}{ }\PYG{n}{fixID}\PYG{o}{\PYGZgt{}}\PYG{+w}{ }\PYG{n}{add\PYGZus{}species}\PYG{+w}{ }\PYG{o}{\PYGZlt{}}\PYG{n}{tag}\PYG{o}{\PYGZgt{}}\PYG{+w}{ }\PYG{o}{\PYGZlt{}}\PYG{n}{group}\PYG{o}{|}\PYG{n}{type}\PYG{o}{\PYGZgt{}}\PYG{+w}{ }\PYG{o}{\PYGZlt{}}\PYG{n}{ID}\PYG{o}{\PYGZgt{}}
\end{sphinxVerbatim}
\begin{itemize}
\item {} 
\sphinxAtStartPar
AtC fixID = ID of {\hyperref[\detokenize{fix_atc::doc}]{\sphinxcrossref{\DUrole{doc}{fix atc}}}} instance

\item {} 
\sphinxAtStartPar
add\_species = name of the AtC sub\sphinxhyphen{}command

\item {} 
\sphinxAtStartPar
tag = tag for tracking a species

\item {} 
\sphinxAtStartPar
\sphinxstyleemphasis{group} or \sphinxstyleemphasis{type} = LAMMPS defined group or type of atoms

\item {} 
\sphinxAtStartPar
ID = name of group or type number

\end{itemize}


\subsubsection{Examples}
\label{\detokenize{atc_add_species:examples}}
\begin{sphinxVerbatim}[commandchars=\\\{\}]
\PYG{k}{fix\PYGZus{}modify }\PYG{n+nv+nvIdentifier}{AtC}\PYG{+w}{ }\PYG{n}{add\PYGZus{}species}\PYG{+w}{ }\PYG{n}{gold}\PYG{+w}{ }\PYG{n}{type}\PYG{+w}{ }\PYG{l+m}{1}
\PYG{k}{group }\PYG{n+nv+nvIdentifier}{GOLDGROUP}\PYG{+w}{ }\PYG{n}{type}\PYG{+w}{ }\PYG{l+m}{1}
\PYG{k}{fix\PYGZus{}modify }\PYG{n+nv+nvIdentifier}{AtC}\PYG{+w}{ }\PYG{n}{add\PYGZus{}species}\PYG{+w}{ }\PYG{n}{gold}\PYG{+w}{ }\PYG{k}{group }\PYG{n+nv+nvIdentifier}{GOLDGROUP}
\end{sphinxVerbatim}


\subsubsection{Description}
\label{\detokenize{atc_add_species:description}}
\sphinxAtStartPar
Associates a tag with all atoms of a specified type or within a specified group.


\subsubsection{Restrictions}
\label{\detokenize{atc_add_species:restrictions}}
\sphinxAtStartPar
None.


\subsubsection{Related AtC commands}
\label{\detokenize{atc_add_species:related-atc-commands}}\begin{itemize}
\item {} 
\sphinxAtStartPar
{\hyperref[\detokenize{fix_atc:atc-fix-modify}]{\sphinxcrossref{\DUrole{std,std-ref}{fix\_modify AtC command overview}}}}

\item {} 
\sphinxAtStartPar
{\hyperref[\detokenize{atc_add_molecule::doc}]{\sphinxcrossref{\DUrole{doc}{fix\_modify AtC add\_molecule}}}}

\item {} 
\sphinxAtStartPar
{\hyperref[\detokenize{atc_remove_species::doc}]{\sphinxcrossref{\DUrole{doc}{fix\_modify AtC remove\_species}}}}

\item {} 
\sphinxAtStartPar
{\hyperref[\detokenize{atc_remove_molecule::doc}]{\sphinxcrossref{\DUrole{doc}{fix\_modify AtC remove\_molecule}}}}

\end{itemize}


\subsubsection{Default}
\label{\detokenize{atc_add_species:default}}
\sphinxAtStartPar
None.

\sphinxstepscope

\index{fix\_modify AtC atom\_element\_map@\spxentry{fix\_modify AtC atom\_element\_map}}\ignorespaces 

\subsection{fix\_modify AtC atom\_element\_map command}
\label{\detokenize{atc_atom_element_map:fix-modify-atc-atom-element-map-command}}\label{\detokenize{atc_atom_element_map:index-0}}\label{\detokenize{atc_atom_element_map::doc}}

\subsubsection{Syntax}
\label{\detokenize{atc_atom_element_map:syntax}}
\begin{sphinxVerbatim}[commandchars=\\\{\}]
\PYG{k}{fix\PYGZus{}modify }\PYG{o}{\PYGZlt{}}\PYG{n}{AtC}\PYG{+w}{ }\PYG{n}{fixID}\PYG{o}{\PYGZgt{}}\PYG{+w}{ }\PYG{n}{atom\PYGZus{}element\PYGZus{}map}\PYG{+w}{ }\PYG{o}{\PYGZlt{}}\PYG{n}{eulerian}\PYG{o}{|}\PYG{n}{lagrangian}\PYG{o}{\PYGZgt{}}\PYG{+w}{ }\PYG{n}{[}\PYG{o}{\PYGZlt{}}\PYG{n}{frequency}\PYG{o}{\PYGZgt{}}\PYG{n}{]}
\end{sphinxVerbatim}
\begin{itemize}
\item {} 
\sphinxAtStartPar
AtC fixID = ID of {\hyperref[\detokenize{fix_atc::doc}]{\sphinxcrossref{\DUrole{doc}{fix atc}}}} instance

\item {} 
\sphinxAtStartPar
atom\_element\_map = name of the AtC sub\sphinxhyphen{}command

\item {} 
\sphinxAtStartPar
\sphinxstyleemphasis{eulerian} or \sphinxstyleemphasis{lagrangian} = frame of reference

\item {} 
\sphinxAtStartPar
frequency = frequency of updating atom\sphinxhyphen{}to\sphinxhyphen{}continuum maps based on the current configuration \sphinxhyphen{} (only for eulerian)

\end{itemize}


\subsubsection{Examples}
\label{\detokenize{atc_atom_element_map:examples}}
\begin{sphinxVerbatim}[commandchars=\\\{\}]
\PYG{k}{fix\PYGZus{}modify }\PYG{n+nv+nvIdentifier}{AtC}\PYG{+w}{ }\PYG{n}{atom\PYGZus{}element\PYGZus{}map}\PYG{+w}{ }\PYG{n}{eulerian}\PYG{+w}{ }\PYG{l+m}{100}
\end{sphinxVerbatim}


\subsubsection{Description}
\label{\detokenize{atc_atom_element_map:description}}
\sphinxAtStartPar
Changes frame of reference from \sphinxstyleemphasis{eulerian} to \sphinxstyleemphasis{lagrangian} or vice versa
and sets the frequency for which the map from atoms to elements is
reformed and all the attendant data is recalculated.


\subsubsection{Restrictions}
\label{\detokenize{atc_atom_element_map:restrictions}}
\sphinxAtStartPar
Cannot change map type after initialization.


\subsubsection{Related AtC commands}
\label{\detokenize{atc_atom_element_map:related-atc-commands}}\begin{itemize}
\item {} 
\sphinxAtStartPar
{\hyperref[\detokenize{fix_atc:atc-fix-modify}]{\sphinxcrossref{\DUrole{std,std-ref}{fix\_modify AtC command overview}}}}

\end{itemize}


\subsubsection{Default}
\label{\detokenize{atc_atom_element_map:default}}
\sphinxAtStartPar
\sphinxstyleemphasis{lagrangian}

\sphinxstepscope

\index{fix\_modify AtC atom\_weight@\spxentry{fix\_modify AtC atom\_weight}}\ignorespaces 

\subsection{fix\_modify AtC atom\_weight command}
\label{\detokenize{atc_atom_weight:fix-modify-atc-atom-weight-command}}\label{\detokenize{atc_atom_weight:index-0}}\label{\detokenize{atc_atom_weight::doc}}

\subsubsection{Syntax}
\label{\detokenize{atc_atom_weight:syntax}}
\begin{sphinxVerbatim}[commandchars=\\\{\}]
\PYG{k}{fix\PYGZus{}modify }\PYG{o}{\PYGZlt{}}\PYG{n}{AtC}\PYG{+w}{ }\PYG{n}{fixID}\PYG{o}{\PYGZgt{}}\PYG{+w}{ }\PYG{n}{atom\PYGZus{}weight}\PYG{+w}{ }\PYG{o}{\PYGZlt{}}\PYG{n}{method}\PYG{o}{\PYGZgt{}}\PYG{+w}{ }\PYG{o}{\PYGZlt{}}\PYG{n}{args}\PYG{o}{\PYGZgt{}}
\end{sphinxVerbatim}
\begin{itemize}
\item {} 
\sphinxAtStartPar
AtC fixID = ID of {\hyperref[\detokenize{fix_atc::doc}]{\sphinxcrossref{\DUrole{doc}{fix atc}}}} instance

\item {} 
\sphinxAtStartPar
atom\_weight = name of the AtC sub\sphinxhyphen{}command

\item {} 
\sphinxAtStartPar
method = \sphinxstyleemphasis{constant} or \sphinxstyleemphasis{lattice} or \sphinxstyleemphasis{element} or \sphinxstyleemphasis{region} or \sphinxstyleemphasis{group} or \sphinxstyleemphasis{read\_in}
\begin{itemize}
\item {} 
\sphinxAtStartPar
\sphinxstyleemphasis{constant} \textless{}group\sphinxhyphen{}ID\textgreater{} \textless{}value\textgreater{}: atoms in specified group are assigned the constant value given

\item {} 
\sphinxAtStartPar
\sphinxstyleemphasis{lattice}: volume per atom for specified lattice type (e.g. fcc) and parameter

\item {} 
\sphinxAtStartPar
\sphinxstyleemphasis{element}: element volume divided among atoms within element

\item {} 
\sphinxAtStartPar
\sphinxstyleemphasis{region}: volume per atom determined based on the atom count in the MD regions and their volumes. Note: meaningful only if atoms completely fill all the regions.

\item {} 
\sphinxAtStartPar
\sphinxstyleemphasis{group}: volume per atom determined based on the atom count in a group and its volume

\item {} 
\sphinxAtStartPar
\sphinxstyleemphasis{node}: (undocumented)

\item {} 
\sphinxAtStartPar
\sphinxstyleemphasis{node\_element}: (undocumented)

\item {} 
\sphinxAtStartPar
\sphinxstyleemphasis{read\_in}\textless{}filename\textgreater{}: list of values for atoms are read\sphinxhyphen{}in from specified file

\end{itemize}

\end{itemize}


\subsubsection{Examples}
\label{\detokenize{atc_atom_weight:examples}}
\begin{sphinxVerbatim}[commandchars=\\\{\}]
\PYG{k}{fix\PYGZus{}modify }\PYG{n+nv+nvIdentifier}{AtC}\PYG{+w}{ }\PYG{n}{atom\PYGZus{}weight}\PYG{+w}{ }\PYG{n}{constant}\PYG{+w}{ }\PYG{n}{myatoms}\PYG{+w}{ }\PYG{l+m}{11.8}
\PYG{k}{fix\PYGZus{}modify }\PYG{n+nv+nvIdentifier}{AtC}\PYG{+w}{ }\PYG{n}{atom\PYGZus{}weight}\PYG{+w}{ }\PYG{n}{lattice}
\PYG{k}{fix\PYGZus{}modify }\PYG{n+nv+nvIdentifier}{AtC}\PYG{+w}{ }\PYG{n}{atom\PYGZus{}weight}\PYG{+w}{ }\PYG{n}{read}\PYG{o}{\PYGZhy{}}\PYG{n}{in}\PYG{+w}{ }\PYG{n}{atm\PYGZus{}wt\PYGZus{}file.txt}
\end{sphinxVerbatim}


\subsubsection{Description}
\label{\detokenize{atc_atom_weight:description}}
\sphinxAtStartPar
Command for assigning the value of atomic weights used for atomic
integration in atom\sphinxhyphen{}continuum coupled simulations.


\subsubsection{Restrictions}
\label{\detokenize{atc_atom_weight:restrictions}}
\sphinxAtStartPar
The use of the lattice option requires a lattice type and parameter is already specified.


\subsubsection{Related AtC commands}
\label{\detokenize{atc_atom_weight:related-atc-commands}}\begin{itemize}
\item {} 
\sphinxAtStartPar
{\hyperref[\detokenize{fix_atc:atc-fix-modify}]{\sphinxcrossref{\DUrole{std,std-ref}{fix\_modify AtC command overview}}}}

\end{itemize}


\subsubsection{Default}
\label{\detokenize{atc_atom_weight:default}}
\sphinxAtStartPar
\sphinxstyleemphasis{lattice}

\sphinxstepscope

\index{fix\_modify AtC atomic\_charge@\spxentry{fix\_modify AtC atomic\_charge}}\ignorespaces 

\subsection{fix\_modify AtC atomic\_charge command}
\label{\detokenize{atc_atomic_charge:fix-modify-atc-atomic-charge-command}}\label{\detokenize{atc_atomic_charge:index-0}}\label{\detokenize{atc_atomic_charge::doc}}

\subsubsection{Syntax}
\label{\detokenize{atc_atomic_charge:syntax}}
\begin{sphinxVerbatim}[commandchars=\\\{\}]
\PYG{k}{fix\PYGZus{}modify }\PYG{o}{\PYGZlt{}}\PYG{n}{AtC}\PYG{+w}{ }\PYG{n}{fixID}\PYG{o}{\PYGZgt{}}\PYG{+w}{ }\PYG{o}{\PYGZlt{}}\PYG{n}{include}\PYG{o}{|}\PYG{n}{omit}\PYG{o}{\PYGZgt{}}\PYG{+w}{ }\PYG{n}{atomic\PYGZus{}charge}
\end{sphinxVerbatim}
\begin{itemize}
\item {} 
\sphinxAtStartPar
AtC fixID = ID of {\hyperref[\detokenize{fix_atc::doc}]{\sphinxcrossref{\DUrole{doc}{fix atc}}}} instance

\item {} 
\sphinxAtStartPar
\sphinxstyleemphasis{include} or \sphinxstyleemphasis{omit} = switch to activate/deactivate inclusion of intrinsic atomic charge in ATC

\item {} 
\sphinxAtStartPar
atomic\_charge = name of the AtC sub\sphinxhyphen{}command

\end{itemize}


\subsubsection{Examples}
\label{\detokenize{atc_atomic_charge:examples}}
\begin{sphinxVerbatim}[commandchars=\\\{\}]
\PYG{k}{fix\PYGZus{}modify }\PYG{n+nv+nvIdentifier}{AtC}\PYG{+w}{ }\PYG{n}{include}\PYG{+w}{ }\PYG{n}{atomic\PYGZus{}charge}
\end{sphinxVerbatim}


\subsubsection{Description}
\label{\detokenize{atc_atomic_charge:description}}
\sphinxAtStartPar
Determines whether AtC tracks the total charge as a finite element
field.


\subsubsection{Restrictions}
\label{\detokenize{atc_atomic_charge:restrictions}}
\sphinxAtStartPar
Required for: \sphinxstyleemphasis{electrostatics}


\subsubsection{Related AtC commands}
\label{\detokenize{atc_atomic_charge:related-atc-commands}}\begin{itemize}
\item {} 
\sphinxAtStartPar
{\hyperref[\detokenize{fix_atc:atc-fix-modify}]{\sphinxcrossref{\DUrole{std,std-ref}{fix\_modify AtC command overview}}}}

\end{itemize}


\subsubsection{Default}
\label{\detokenize{atc_atomic_charge:default}}
\sphinxAtStartPar
If the atom charge is defined, default is on, otherwise default is off.

\sphinxstepscope

\index{fix\_modify AtC boundary\_dynamics@\spxentry{fix\_modify AtC boundary\_dynamics}}\ignorespaces 

\subsection{fix\_modify AtC boundary\_dynamics command}
\label{\detokenize{atc_boundary_dynamics:fix-modify-atc-boundary-dynamics-command}}\label{\detokenize{atc_boundary_dynamics:index-0}}\label{\detokenize{atc_boundary_dynamics::doc}}

\subsubsection{Syntax}
\label{\detokenize{atc_boundary_dynamics:syntax}}
\begin{sphinxVerbatim}[commandchars=\\\{\}]
\PYG{k}{fix\PYGZus{}modify }\PYG{o}{\PYGZlt{}}\PYG{n}{AtC}\PYG{+w}{ }\PYG{n}{fixID}\PYG{o}{\PYGZgt{}}\PYG{+w}{ }\PYG{n}{boundary\PYGZus{}dynamics}\PYG{+w}{ }\PYG{o}{\PYGZlt{}}\PYG{n}{on}\PYG{o}{|}\PYG{n}{damped\PYGZus{}harmonic}\PYG{o}{|}\PYG{n}{prescribed}\PYG{o}{|}\PYG{n}{coupled}\PYG{o}{|}\PYG{n}{none}\PYG{o}{\PYGZgt{}}
\end{sphinxVerbatim}
\begin{itemize}
\item {} 
\sphinxAtStartPar
AtC fixID = ID of {\hyperref[\detokenize{fix_atc::doc}]{\sphinxcrossref{\DUrole{doc}{fix atc}}}} instance

\item {} 
\sphinxAtStartPar
boundary\_dynamics = name of the AtC sub\sphinxhyphen{}command

\item {} 
\sphinxAtStartPar
\sphinxstyleemphasis{on} or \sphinxstyleemphasis{damped\_harmonic} \sphinxstyleemphasis{prescribed} \sphinxstyleemphasis{coupled} \sphinxstyleemphasis{none}

\end{itemize}


\subsubsection{Description}
\label{\detokenize{atc_boundary_dynamics:description}}
\sphinxAtStartPar
Sets different schemes for controlling boundary atoms.  \sphinxstyleemphasis{on} will integrate
the boundary atoms using the velocity\sphinxhyphen{}verlet algorithm.  \sphinxstyleemphasis{damped\_harmonic}
uses a mass/spring/dashpot for the boundary atoms with added arguments
of the damping and spring constants followed by the ratio of the
boundary type mass to the desired mass.  \sphinxstyleemphasis{prescribed} forces the boundary
atoms to follow the finite element displacement.  \sphinxstyleemphasis{coupled} does the same.


\subsubsection{Restrictions}
\label{\detokenize{atc_boundary_dynamics:restrictions}}
\sphinxAtStartPar
Boundary atoms must be specified.  When using swaps between internal and
boundary atoms, the initial configuration must have already correctly
partitioned the two.


\subsubsection{Related AtC commands}
\label{\detokenize{atc_boundary_dynamics:related-atc-commands}}\begin{itemize}
\item {} 
\sphinxAtStartPar
{\hyperref[\detokenize{fix_atc:atc-fix-modify}]{\sphinxcrossref{\DUrole{std,std-ref}{fix\_modify AtC command overview}}}}

\end{itemize}


\subsubsection{Default}
\label{\detokenize{atc_boundary_dynamics:default}}
\sphinxAtStartPar
\sphinxstyleemphasis{prescribed on}

\sphinxstepscope

\index{fix\_modify AtC boundary\_faceset@\spxentry{fix\_modify AtC boundary\_faceset}}\ignorespaces 

\subsection{fix\_modify AtC boundary\_faceset command}
\label{\detokenize{atc_boundary_faceset:fix-modify-atc-boundary-faceset-command}}\label{\detokenize{atc_boundary_faceset:index-0}}\label{\detokenize{atc_boundary_faceset::doc}}

\subsubsection{Syntax}
\label{\detokenize{atc_boundary_faceset:syntax}}
\begin{sphinxVerbatim}[commandchars=\\\{\}]
\PYG{k}{fix\PYGZus{}modify }\PYG{o}{\PYGZlt{}}\PYG{n}{AtC}\PYG{+w}{ }\PYG{n}{fixID}\PYG{o}{\PYGZgt{}}\PYG{+w}{ }\PYG{n}{boundary\PYGZus{}faceset}\PYG{+w}{ }\PYG{o}{\PYGZlt{}}\PYG{n}{is}\PYG{o}{|}\PYG{n}{add}\PYG{o}{\PYGZgt{}}\PYG{+w}{ }\PYG{o}{\PYGZlt{}}\PYG{n}{faceset\PYGZus{}name}\PYG{o}{\PYGZgt{}}
\end{sphinxVerbatim}
\begin{itemize}
\item {} 
\sphinxAtStartPar
AtC fixID = ID of {\hyperref[\detokenize{fix_atc::doc}]{\sphinxcrossref{\DUrole{doc}{fix atc}}}} instance

\item {} 
\sphinxAtStartPar
boundary\_faceset = name of the AtC sub\sphinxhyphen{}command

\item {} 
\sphinxAtStartPar
\sphinxstyleemphasis{is} or \sphinxstyleemphasis{add} = select whether to select or add a faceset

\item {} 
\sphinxAtStartPar
faceset\_name = name of the faceset

\end{itemize}


\subsubsection{Examples}
\label{\detokenize{atc_boundary_faceset:examples}}
\begin{sphinxVerbatim}[commandchars=\\\{\}]
\PYG{k}{fix\PYGZus{}modify }\PYG{n+nv+nvIdentifier}{AtC}\PYG{+w}{ }\PYG{n}{boundary\PYGZus{}faceset}\PYG{+w}{ }\PYG{n}{is}\PYG{+w}{ }\PYG{n}{obndy}
\end{sphinxVerbatim}


\subsubsection{Description}
\label{\detokenize{atc_boundary_faceset:description}}
\sphinxAtStartPar
This command species the faceset name when using a faceset to compute
the MD/FE boundary fluxes. The faceset must already exist.


\subsubsection{Restrictions}
\label{\detokenize{atc_boundary_faceset:restrictions}}
\sphinxAtStartPar
This is only valid when \sphinxstyleemphasis{fe\_md\_boundary} is set to \sphinxstyleemphasis{faceset}.


\subsubsection{Related AtC commands}
\label{\detokenize{atc_boundary_faceset:related-atc-commands}}\begin{itemize}
\item {} 
\sphinxAtStartPar
{\hyperref[\detokenize{fix_atc:atc-fix-modify}]{\sphinxcrossref{\DUrole{std,std-ref}{fix\_modify AtC command overview}}}}

\item {} 
\sphinxAtStartPar
{\hyperref[\detokenize{atc_fe_md_boundary::doc}]{\sphinxcrossref{\DUrole{doc}{fix\_modify AtC fe\_md\_boundary}}}}

\item {} 
\sphinxAtStartPar
{\hyperref[\detokenize{atc_mesh_create_faceset_box::doc}]{\sphinxcrossref{\DUrole{doc}{fix\_modify AtC mesh create\_faceset box}}}}

\item {} 
\sphinxAtStartPar
{\hyperref[\detokenize{atc_mesh_create_faceset_plane::doc}]{\sphinxcrossref{\DUrole{doc}{fix\_modify AtC mesh create\_faceset plane}}}}

\end{itemize}


\subsubsection{Default}
\label{\detokenize{atc_boundary_faceset:default}}
\sphinxAtStartPar
None.

\sphinxstepscope

\index{fix\_modify AtC boundary type@\spxentry{fix\_modify AtC boundary type}}\ignorespaces 

\subsection{fix\_modify AtC boundary type command}
\label{\detokenize{atc_boundary_type:fix-modify-atc-boundary-type-command}}\label{\detokenize{atc_boundary_type:index-0}}\label{\detokenize{atc_boundary_type::doc}}

\subsubsection{Syntax}
\label{\detokenize{atc_boundary_type:syntax}}
\begin{sphinxVerbatim}[commandchars=\\\{\}]
\PYG{k}{fix\PYGZus{}modify }\PYG{o}{\PYGZlt{}}\PYG{n}{AtC}\PYG{+w}{ }\PYG{n}{fixID}\PYG{o}{\PYGZgt{}}\PYG{+w}{ }\PYG{n}{boundary}\PYG{+w}{ }\PYG{n}{type}\PYG{+w}{ }\PYG{o}{\PYGZlt{}}\PYG{n}{atom}\PYG{o}{\PYGZhy{}}\PYG{n}{type}\PYG{o}{\PYGZhy{}}\PYG{n}{id}\PYG{o}{\PYGZgt{}}
\end{sphinxVerbatim}
\begin{itemize}
\item {} 
\sphinxAtStartPar
AtC fixID = ID of {\hyperref[\detokenize{fix_atc::doc}]{\sphinxcrossref{\DUrole{doc}{fix atc}}}} instance

\item {} 
\sphinxAtStartPar
boundary type = name of the AtC sub\sphinxhyphen{}command

\item {} 
\sphinxAtStartPar
atom\sphinxhyphen{}type\sphinxhyphen{}id =  type id for atoms that represent a fictitious boundary internal to the FE mesh

\end{itemize}


\subsubsection{Examples}
\label{\detokenize{atc_boundary_type:examples}}
\begin{sphinxVerbatim}[commandchars=\\\{\}]
\PYG{k}{fix\PYGZus{}modify }\PYG{n+nv+nvIdentifier}{AtC}\PYG{+w}{ }\PYG{n}{boundary}\PYG{+w}{ }\PYG{n}{type}\PYG{+w}{ }\PYG{n}{ghost\PYGZus{}atoms}
\end{sphinxVerbatim}


\subsubsection{Description}
\label{\detokenize{atc_boundary_type:description}}
\sphinxAtStartPar
Command to define the atoms that represent the fictitious boundary
internal to the FE mesh. For fully overlapped MD/FE domains with
periodic boundary conditions no boundary atoms should be defined.


\subsubsection{Restrictions}
\label{\detokenize{atc_boundary_type:restrictions}}
\sphinxAtStartPar
None.


\subsubsection{Related AtC commands}
\label{\detokenize{atc_boundary_type:related-atc-commands}}\begin{itemize}
\item {} 
\sphinxAtStartPar
{\hyperref[\detokenize{fix_atc:atc-fix-modify}]{\sphinxcrossref{\DUrole{std,std-ref}{fix\_modify AtC command overview}}}}

\end{itemize}


\subsubsection{Default}
\label{\detokenize{atc_boundary_type:default}}
\sphinxAtStartPar
None.

\sphinxstepscope

\index{fix\_modify AtC consistent\_fe\_initialization@\spxentry{fix\_modify AtC consistent\_fe\_initialization}}\ignorespaces 

\subsection{fix\_modify AtC consistent\_fe\_initialization command}
\label{\detokenize{atc_consistent_fe_initialization:fix-modify-atc-consistent-fe-initialization-command}}\label{\detokenize{atc_consistent_fe_initialization:index-0}}\label{\detokenize{atc_consistent_fe_initialization::doc}}

\subsubsection{Syntax}
\label{\detokenize{atc_consistent_fe_initialization:syntax}}
\begin{sphinxVerbatim}[commandchars=\\\{\}]
\PYG{k}{fix\PYGZus{}modify }\PYG{o}{\PYGZlt{}}\PYG{n}{AtC}\PYG{+w}{ }\PYG{n}{fixID}\PYG{o}{\PYGZgt{}}\PYG{+w}{ }\PYG{n}{consistent\PYGZus{}fe\PYGZus{}initialization}\PYG{+w}{ }\PYG{o}{\PYGZlt{}}\PYG{n}{on}\PYG{o}{|}\PYG{n}{off}\PYG{o}{\PYGZgt{}}
\end{sphinxVerbatim}
\begin{itemize}
\item {} 
\sphinxAtStartPar
AtC fixID = ID of {\hyperref[\detokenize{fix_atc::doc}]{\sphinxcrossref{\DUrole{doc}{fix atc}}}} instance

\item {} 
\sphinxAtStartPar
consistent\_fe\_initialization = name of the AtC sub\sphinxhyphen{}command

\item {} 
\sphinxAtStartPar
\sphinxstyleemphasis{on} or \sphinxstyleemphasis{off} = switch to activate/deactivate the initial setting of the FE intrinsic field to match the projected MD field

\end{itemize}


\subsubsection{Examples}
\label{\detokenize{atc_consistent_fe_initialization:examples}}
\begin{sphinxVerbatim}[commandchars=\\\{\}]
\PYG{k}{fix\PYGZus{}modify }\PYG{n+nv+nvIdentifier}{AtC}\PYG{+w}{ }\PYG{n}{consistent\PYGZus{}fe\PYGZus{}initialization}\PYG{+w}{ }\PYG{n}{on}
\end{sphinxVerbatim}


\subsubsection{Description}
\label{\detokenize{atc_consistent_fe_initialization:description}}
\sphinxAtStartPar
Determines whether AtC initializes FE intrinsic fields (e.g.,
temperature) to match the projected MD values. This is particularly
useful for fully overlapping simulations.


\subsubsection{Restrictions}
\label{\detokenize{atc_consistent_fe_initialization:restrictions}}
\sphinxAtStartPar
Can be used with: \sphinxstyleemphasis{thermal}, \sphinxstyleemphasis{two\_temperature}.
Cannot be used with time filtering on.
Does not include boundary nodes.


\subsubsection{Related AtC commands}
\label{\detokenize{atc_consistent_fe_initialization:related-atc-commands}}\begin{itemize}
\item {} 
\sphinxAtStartPar
{\hyperref[\detokenize{fix_atc:atc-fix-modify}]{\sphinxcrossref{\DUrole{std,std-ref}{fix\_modify AtC command overview}}}}

\end{itemize}


\subsubsection{Default}
\label{\detokenize{atc_consistent_fe_initialization:default}}
\sphinxAtStartPar
Default is \sphinxstyleemphasis{off}

\sphinxstepscope

\index{fix\_modify AtC control localized\_lambda@\spxentry{fix\_modify AtC control localized\_lambda}}\ignorespaces 

\subsection{fix\_modify AtC control localized\_lambda command}
\label{\detokenize{atc_control_localized_lambda:fix-modify-atc-control-localized-lambda-command}}\label{\detokenize{atc_control_localized_lambda:index-0}}\label{\detokenize{atc_control_localized_lambda::doc}}

\subsubsection{Syntax}
\label{\detokenize{atc_control_localized_lambda:syntax}}
\begin{sphinxVerbatim}[commandchars=\\\{\}]
\PYG{k}{fix\PYGZus{}modify }\PYG{o}{\PYGZlt{}}\PYG{n}{AtC}\PYG{+w}{ }\PYG{n}{fixID}\PYG{o}{\PYGZgt{}}\PYG{+w}{ }\PYG{n}{control}\PYG{+w}{ }\PYG{n}{localized\PYGZus{}lambda}\PYG{+w}{ }\PYG{o}{\PYGZlt{}}\PYG{n}{on}\PYG{o}{|}\PYG{n}{off}\PYG{o}{\PYGZgt{}}
\end{sphinxVerbatim}
\begin{itemize}
\item {} 
\sphinxAtStartPar
AtC fixID = ID of {\hyperref[\detokenize{fix_atc::doc}]{\sphinxcrossref{\DUrole{doc}{fix atc}}}} instance

\item {} 
\sphinxAtStartPar
control localized\_lambda = name of the AtC sub\sphinxhyphen{}command

\item {} 
\sphinxAtStartPar
\sphinxstyleemphasis{on} or \sphinxstyleemphasis{off} = Toggles state of localization algorithm

\end{itemize}


\subsubsection{Examples}
\label{\detokenize{atc_control_localized_lambda:examples}}
\begin{sphinxVerbatim}[commandchars=\\\{\}]
\PYG{k}{fix\PYGZus{}modify }\PYG{n+nv+nvIdentifier}{AtC}\PYG{+w}{ }\PYG{n}{control}\PYG{+w}{ }\PYG{n}{localized\PYGZus{}lambda}\PYG{+w}{ }\PYG{n}{on}
\end{sphinxVerbatim}


\subsubsection{Description}
\label{\detokenize{atc_control_localized_lambda:description}}
\sphinxAtStartPar
Turns the localization algorithms \sphinxstyleemphasis{on} or \sphinxstyleemphasis{off} for control algorithms
to restrict the influence of FE coupling or boundary conditions to a
region near the boundary of the MD region.  Control algorithms will not
affect atoms in elements not possessing faces on the boundary of the
region.  Flux\sphinxhyphen{}based control is localized via row\sphinxhyphen{}sum lumping while
quantity control is done by solving a truncated matrix equation.


\subsubsection{Restrictions}
\label{\detokenize{atc_control_localized_lambda:restrictions}}
\sphinxAtStartPar
None.


\subsubsection{Related AtC commands}
\label{\detokenize{atc_control_localized_lambda:related-atc-commands}}\begin{itemize}
\item {} 
\sphinxAtStartPar
{\hyperref[\detokenize{fix_atc:atc-fix-modify}]{\sphinxcrossref{\DUrole{std,std-ref}{fix\_modify AtC command overview}}}}

\end{itemize}


\subsubsection{Default}
\label{\detokenize{atc_control_localized_lambda:default}}
\sphinxAtStartPar
off

\sphinxstepscope

\index{fix\_modify AtC control momentum@\spxentry{fix\_modify AtC control momentum}}\ignorespaces 

\subsection{fix\_modify AtC control momentum command}
\label{\detokenize{atc_control_momentum:fix-modify-atc-control-momentum-command}}\label{\detokenize{atc_control_momentum:index-0}}\label{\detokenize{atc_control_momentum::doc}}

\subsubsection{Syntax}
\label{\detokenize{atc_control_momentum:syntax}}
\begin{sphinxVerbatim}[commandchars=\\\{\}]
\PYG{k}{fix\PYGZus{}modify }\PYG{o}{\PYGZlt{}}\PYG{n}{AtC}\PYG{+w}{ }\PYG{n}{fixID}\PYG{o}{\PYGZgt{}}\PYG{+w}{ }\PYG{n}{control}\PYG{+w}{ }\PYG{o}{\PYGZlt{}}\PYG{n}{physics\PYGZus{}type}\PYG{o}{\PYGZgt{}}\PYG{+w}{ }\PYG{o}{\PYGZlt{}}\PYG{n}{solution\PYGZus{}parameter}\PYG{o}{\PYGZgt{}}\PYG{+w}{ }\PYG{o}{\PYGZlt{}}\PYG{n}{value}\PYG{o}{\PYGZgt{}}
\PYG{k}{fix\PYGZus{}modify }\PYG{n+nv+nvIdentifier}{AtC}\PYG{+w}{ }\PYG{n}{control}\PYG{+w}{ }\PYG{n}{momentum}\PYG{+w}{ }\PYG{n}{none}
\PYG{k}{fix\PYGZus{}modify }\PYG{n+nv+nvIdentifier}{AtC}\PYG{+w}{ }\PYG{n}{control}\PYG{+w}{ }\PYG{n}{momentum}\PYG{+w}{ }\PYG{n}{rescale}\PYG{+w}{ }\PYG{o}{\PYGZlt{}}\PYG{n}{frequency}\PYG{o}{\PYGZgt{}}
\PYG{k}{fix\PYGZus{}modify }\PYG{n+nv+nvIdentifier}{AtC}\PYG{+w}{ }\PYG{n}{control}\PYG{+w}{ }\PYG{n}{momentum}\PYG{+w}{ }\PYG{n}{glc\PYGZus{}displacement}
\PYG{k}{fix\PYGZus{}modify }\PYG{n+nv+nvIdentifier}{AtC}\PYG{+w}{ }\PYG{n}{control}\PYG{+w}{ }\PYG{n}{momentum}\PYG{+w}{ }\PYG{n}{glc\PYGZus{}velocity}
\PYG{k}{fix\PYGZus{}modify }\PYG{n+nv+nvIdentifier}{AtC}\PYG{+w}{ }\PYG{n}{control}\PYG{+w}{ }\PYG{n}{momentum}\PYG{+w}{ }\PYG{n}{hoover}
\PYG{k}{fix\PYGZus{}modify }\PYG{n+nv+nvIdentifier}{AtC}\PYG{+w}{ }\PYG{n}{control}\PYG{+w}{ }\PYG{n}{momentum}\PYG{+w}{ }\PYG{n}{flux}\PYG{+w}{ }\PYG{n}{[faceset}\PYG{+w}{ }\PYG{n}{face\PYGZus{}set\PYGZus{}id},\PYG{+w}{ }\PYG{n}{interpolate]}
\end{sphinxVerbatim}
\begin{itemize}
\item {} 
\sphinxAtStartPar
AtC fixID = ID of {\hyperref[\detokenize{fix_atc::doc}]{\sphinxcrossref{\DUrole{doc}{fix atc}}}} instance

\item {} 
\sphinxAtStartPar
control = name of the AtC sub\sphinxhyphen{}command

\item {} 
\sphinxAtStartPar
physics\_type = \sphinxstyleemphasis{thermal} or \sphinxstyleemphasis{momentum}

\item {} 
\sphinxAtStartPar
solution\_parameter = \sphinxstyleemphasis{max\_iterations} or \sphinxstyleemphasis{tolerance}

\item {} 
\sphinxAtStartPar
value = solution\_parameter value

\item {} 
\sphinxAtStartPar
momentum option = \sphinxstyleemphasis{none} or \sphinxstyleemphasis{rescale} or \sphinxstyleemphasis{glc\_displacement} or \sphinxstyleemphasis{glc\_velocity} \sphinxstyleemphasis{hoover} or \sphinxstyleemphasis{flux}

\item {} 
\sphinxAtStartPar
frequency = time step frequency for applying displacement and velocity rescaling

\item {} 
\sphinxAtStartPar
faceset\_id = id of boundary face set (optional, only for \sphinxstyleemphasis{faceset})

\end{itemize}


\subsubsection{Examples}
\label{\detokenize{atc_control_momentum:examples}}
\begin{sphinxVerbatim}[commandchars=\\\{\}]
\PYG{k}{fix\PYGZus{}modify }\PYG{n+nv+nvIdentifier}{AtC}\PYG{+w}{ }\PYG{n}{control}\PYG{+w}{ }\PYG{n}{momentum}\PYG{+w}{ }\PYG{n}{none}
\PYG{k}{fix\PYGZus{}modify }\PYG{n+nv+nvIdentifier}{AtC}\PYG{+w}{ }\PYG{n}{control}\PYG{+w}{ }\PYG{n}{momentum}\PYG{+w}{ }\PYG{n}{flux}\PYG{+w}{ }\PYG{n}{faceset}\PYG{+w}{ }\PYG{n}{bndy\PYGZus{}faces}
\PYG{k}{fix\PYGZus{}modify }\PYG{n+nv+nvIdentifier}{AtC}\PYG{+w}{ }\PYG{n}{control}\PYG{+w}{ }\PYG{n}{momentum}\PYG{+w}{ }\PYG{n}{glc\PYGZus{}velocity}
\end{sphinxVerbatim}


\subsubsection{Description}
\label{\detokenize{atc_control_momentum:description}}
\sphinxAtStartPar
The general version of \sphinxstyleemphasis{control} sets the numerical parameters for the
matrix solvers used in the specified control algorithm.  Many solution
approaches require iterative solvers, and these methods enable users to
provide the maximum number of iterations and the relative tolerance.

\sphinxAtStartPar
The \sphinxstyleemphasis{control momentum} version sets the momentum exchange mechanism from
the finite elements to the atoms, managed through a control algorithm.
\sphinxstyleemphasis{rescale} computes a scale factor for each atom to match the finite
element temperature.  \sphinxstyleemphasis{hoover} is a Gaussian least\sphinxhyphen{}constraint isokinetic
thermostat enforces that the nodal restricted atomic temperature matches
the finite element temperature.  \sphinxstyleemphasis{flux} is a similar mode, but rather
adds energy to the atoms based on conservation of energy.

\sphinxAtStartPar
\sphinxstyleemphasis{correction\_max\_iterations} sets the maximum number of iterations to
compute the second order in time correction term for lambda with the
fractional step method. The method uses the same tolerance as the
controller’s matrix solver.


\subsubsection{Restrictions}
\label{\detokenize{atc_control_momentum:restrictions}}
\sphinxAtStartPar
Only for be used with the specific controllers \sphinxstyleemphasis{thermal} or \sphinxstyleemphasis{momentum}.
They are ignored if a lumped solution is requested.

\sphinxAtStartPar
\sphinxstyleemphasis{control momentum} is only for be used with specific transfers: elastic
\sphinxstyleemphasis{rescale} not valid with time filtering activated


\subsubsection{Related AtC commands}
\label{\detokenize{atc_control_momentum:related-atc-commands}}\begin{itemize}
\item {} 
\sphinxAtStartPar
{\hyperref[\detokenize{fix_atc:atc-fix-modify}]{\sphinxcrossref{\DUrole{std,std-ref}{fix\_modify AtC command overview}}}}

\item {} 
\sphinxAtStartPar
{\hyperref[\detokenize{atc_control_thermal::doc}]{\sphinxcrossref{\DUrole{doc}{fix\_modify AtC control thermal}}}}

\end{itemize}


\subsubsection{Default}
\label{\detokenize{atc_control_momentum:default}}\begin{itemize}
\item {} 
\sphinxAtStartPar
\sphinxstyleemphasis{max\_iterations} is the number of rows in the matrix.

\item {} 
\sphinxAtStartPar
\sphinxstyleemphasis{tolerance} is 1.0e\sphinxhyphen{}10.

\end{itemize}

\sphinxstepscope

\index{fix\_modify AtC control thermal@\spxentry{fix\_modify AtC control thermal}}\ignorespaces 

\subsection{fix\_modify AtC control thermal command}
\label{\detokenize{atc_control_thermal:fix-modify-atc-control-thermal-command}}\label{\detokenize{atc_control_thermal:index-0}}\label{\detokenize{atc_control_thermal::doc}}

\subsubsection{Syntax}
\label{\detokenize{atc_control_thermal:syntax}}
\begin{sphinxVerbatim}[commandchars=\\\{\}]
\PYG{k}{fix\PYGZus{}modify }\PYG{o}{\PYGZlt{}}\PYG{n}{AtC}\PYG{+w}{ }\PYG{n}{fixID}\PYG{o}{\PYGZgt{}}\PYG{+w}{ }\PYG{n}{control}\PYG{+w}{ }\PYG{o}{\PYGZlt{}}\PYG{n}{physics\PYGZus{}type}\PYG{o}{\PYGZgt{}}\PYG{+w}{ }\PYG{o}{\PYGZlt{}}\PYG{n}{solution\PYGZus{}parameter}\PYG{o}{\PYGZgt{}}\PYG{+w}{ }\PYG{o}{\PYGZlt{}}\PYG{n}{value}\PYG{o}{\PYGZgt{}}
\PYG{k}{fix\PYGZus{}modify }\PYG{o}{\PYGZlt{}}\PYG{n}{AtC}\PYG{+w}{ }\PYG{n}{fixID}\PYG{o}{\PYGZgt{}}\PYG{+w}{ }\PYG{n}{control}\PYG{+w}{ }\PYG{n}{thermal}\PYG{+w}{ }\PYG{o}{\PYGZlt{}}\PYG{n}{control\PYGZus{}type}\PYG{o}{\PYGZgt{}}\PYG{+w}{ }\PYG{o}{\PYGZlt{}}\PYG{n}{optional\PYGZus{}args}\PYG{o}{\PYGZgt{}}
\PYG{k}{fix\PYGZus{}modify }\PYG{o}{\PYGZlt{}}\PYG{n}{AtC}\PYG{+w}{ }\PYG{n}{fixID}\PYG{o}{\PYGZgt{}}\PYG{+w}{ }\PYG{n}{control}\PYG{+w}{ }\PYG{n}{thermal}\PYG{+w}{ }\PYG{n}{rescale}\PYG{+w}{ }\PYG{o}{\PYGZlt{}}\PYG{n}{frequency}\PYG{o}{\PYGZgt{}}
\PYG{k}{fix\PYGZus{}modify }\PYG{o}{\PYGZlt{}}\PYG{n}{AtC}\PYG{+w}{ }\PYG{n}{fixID}\PYG{o}{\PYGZgt{}}\PYG{+w}{ }\PYG{n}{control}\PYG{+w}{ }\PYG{n}{thermal}\PYG{+w}{ }\PYG{n}{flux}\PYG{+w}{ }\PYG{o}{\PYGZlt{}}\PYG{n}{boundary\PYGZus{}integration\PYGZus{}type}\PYG{o}{\PYGZgt{}}\PYG{+w}{ }\PYG{o}{\PYGZlt{}}\PYG{n}{faceset\PYGZus{}id}\PYG{o}{\PYGZgt{}}
\PYG{k}{fix\PYGZus{}modify }\PYG{o}{\PYGZlt{}}\PYG{n}{AtC}\PYG{+w}{ }\PYG{n}{fixID}\PYG{o}{\PYGZgt{}}\PYG{+w}{ }\PYG{n}{control}\PYG{+w}{ }\PYG{n}{thermal}\PYG{+w}{ }\PYG{n}{correction\PYGZus{}max\PYGZus{}iterations}\PYG{+w}{ }\PYG{o}{\PYGZlt{}}\PYG{n}{max\PYGZus{}iterations}\PYG{o}{\PYGZgt{}}
\end{sphinxVerbatim}
\begin{itemize}
\item {} 
\sphinxAtStartPar
AtC fixID = ID of {\hyperref[\detokenize{fix_atc::doc}]{\sphinxcrossref{\DUrole{doc}{fix atc}}}} instance

\item {} 
\sphinxAtStartPar
control = name of the AtC sub\sphinxhyphen{}command

\item {} 
\sphinxAtStartPar
physics\_type = \sphinxstyleemphasis{thermal} or \sphinxstyleemphasis{momentum}

\item {} 
\sphinxAtStartPar
solution\_parameter = \sphinxstyleemphasis{max\_iterations} or \sphinxstyleemphasis{tolerance}

\item {} 
\sphinxAtStartPar
value = solution\_parameter value

\item {} 
\sphinxAtStartPar
thermal control\_type = \sphinxstyleemphasis{none} or \sphinxstyleemphasis{rescale} or \sphinxstyleemphasis{hoover} or \sphinxstyleemphasis{flux}

\item {} 
\sphinxAtStartPar
frequency = time step frequency for applying velocity rescaling

\item {} 
\sphinxAtStartPar
boundary\_integration\_type = \sphinxstyleemphasis{faceset} or \sphinxstyleemphasis{interpolate} (optional)

\item {} 
\sphinxAtStartPar
faceset\_id = id of boundary face set (optional, only for \sphinxstyleemphasis{faceset})

\item {} 
\sphinxAtStartPar
correction\_max\_iterations = maximum number of iterations that will be used by iterative matrix solvers for \sphinxstyleemphasis{thermal} physics type

\end{itemize}


\subsubsection{Examples}
\label{\detokenize{atc_control_thermal:examples}}
\begin{sphinxVerbatim}[commandchars=\\\{\}]
\PYG{k}{fix\PYGZus{}modify }\PYG{n+nv+nvIdentifier}{AtC}\PYG{+w}{ }\PYG{n}{control}\PYG{+w}{ }\PYG{n}{thermal}\PYG{+w}{ }\PYG{n}{none}
\PYG{k}{fix\PYGZus{}modify }\PYG{n+nv+nvIdentifier}{AtC}\PYG{+w}{ }\PYG{n}{control}\PYG{+w}{ }\PYG{n}{thermal}\PYG{+w}{ }\PYG{n}{rescale}\PYG{+w}{ }\PYG{l+m}{10}
\PYG{k}{fix\PYGZus{}modify }\PYG{n+nv+nvIdentifier}{AtC}\PYG{+w}{ }\PYG{n}{control}\PYG{+w}{ }\PYG{n}{thermal}\PYG{+w}{ }\PYG{n}{hoover}
\PYG{k}{fix\PYGZus{}modify }\PYG{n+nv+nvIdentifier}{AtC}\PYG{+w}{ }\PYG{n}{control}\PYG{+w}{ }\PYG{n}{thermal}\PYG{+w}{ }\PYG{n}{flux}
\PYG{k}{fix\PYGZus{}modify }\PYG{n+nv+nvIdentifier}{AtC}\PYG{+w}{ }\PYG{n}{control}\PYG{+w}{ }\PYG{n}{thermal}\PYG{+w}{ }\PYG{n}{flux}\PYG{+w}{ }\PYG{n}{faceset}\PYG{+w}{ }\PYG{n}{bndy\PYGZus{}faces}
\PYG{k}{fix\PYGZus{}modify }\PYG{n+nv+nvIdentifier}{AtC}\PYG{+w}{ }\PYG{n}{control}\PYG{+w}{ }\PYG{n}{thermal}\PYG{+w}{ }\PYG{n}{correction\PYGZus{}max\PYGZus{}iterations}\PYG{+w}{ }\PYG{l+m}{10}
\end{sphinxVerbatim}


\subsubsection{Description}
\label{\detokenize{atc_control_thermal:description}}
\sphinxAtStartPar
The general version of \sphinxstyleemphasis{control} sets the numerical parameters for the
matrix solvers used in the specified control algorithm.  Many solution
approaches require iterative solvers, and these methods enable users to
provide the maximum number of iterations and the relative tolerance.

\sphinxAtStartPar
The \sphinxstyleemphasis{control thermal} version sets the energy exchange mechanism from
the finite elements to the atoms, managed through a control algorithm.
\sphinxstyleemphasis{rescale} computes a scale factor for each atom to match the finite
element temperature.  \sphinxstyleemphasis{hoover} is a Gaussian least\sphinxhyphen{}constraint isokinetic
thermostat enforces that the nodal restricted atomic temperature matches
the finite element temperature.  \sphinxstyleemphasis{flux} is a similar mode, but rather
adds energy to the atoms based on conservation of energy. \sphinxstyleemphasis{hoover} and
\sphinxstyleemphasis{flux} allow the prescription of sources or fixed temperatures on the
atoms.

\sphinxAtStartPar
\sphinxstyleemphasis{correction\_max\_iterations} sets the maximum number of iterations to
compute the second order in time correction term for lambda with the
fractional step method. The method uses the same tolerance as the
controller’s matrix solver.


\subsubsection{Restrictions}
\label{\detokenize{atc_control_thermal:restrictions}}
\sphinxAtStartPar
Only for be used with the specific controllers \sphinxstyleemphasis{thermal} or \sphinxstyleemphasis{momentum}.
They are ignored if a lumped solution is requested.

\sphinxAtStartPar
\sphinxstyleemphasis{control thermal} is only for be used with specific transfers: thermal (\sphinxstyleemphasis{rescale}, \sphinxstyleemphasis{hoover}, \sphinxstyleemphasis{flux}), \sphinxstyleemphasis{two\_temperature} (\sphinxstyleemphasis{flux}).
\sphinxstyleemphasis{rescale} not valid with time filtering activated

\sphinxAtStartPar
\sphinxstyleemphasis{correction\_max\_iterations} is only for use with \sphinxstyleemphasis{thermal} physics using
the fractional step method.


\subsubsection{Related AtC commands}
\label{\detokenize{atc_control_thermal:related-atc-commands}}\begin{itemize}
\item {} 
\sphinxAtStartPar
{\hyperref[\detokenize{fix_atc:atc-fix-modify}]{\sphinxcrossref{\DUrole{std,std-ref}{fix\_modify AtC command overview}}}}

\item {} 
\sphinxAtStartPar
{\hyperref[\detokenize{atc_control_momentum::doc}]{\sphinxcrossref{\DUrole{doc}{fix\_modify AtC control momentum}}}}

\end{itemize}


\subsubsection{Default}
\label{\detokenize{atc_control_thermal:default}}\begin{itemize}
\item {} 
\sphinxAtStartPar
\sphinxstyleemphasis{max\_iterations} is the number of rows in the matrix.

\item {} 
\sphinxAtStartPar
\sphinxstyleemphasis{tolerance} is 1.0e\sphinxhyphen{}10.

\item {} 
\sphinxAtStartPar
\sphinxstyleemphasis{rescale} frequency is 1

\item {} 
\sphinxAtStartPar
\sphinxstyleemphasis{flux} boundary\_integration\_type is \sphinxstyleemphasis{interpolate}

\item {} 
\sphinxAtStartPar
\sphinxstyleemphasis{correction\_max\_iterations} is 20

\end{itemize}

\sphinxstepscope

\index{fix\_modify AtC decomposition@\spxentry{fix\_modify AtC decomposition}}\ignorespaces 

\subsection{fix\_modify AtC decomposition command}
\label{\detokenize{atc_decomposition:fix-modify-atc-decomposition-command}}\label{\detokenize{atc_decomposition:index-0}}\label{\detokenize{atc_decomposition::doc}}

\subsubsection{Syntax}
\label{\detokenize{atc_decomposition:syntax}}
\begin{sphinxVerbatim}[commandchars=\\\{\}]
\PYG{k}{fix\PYGZus{}modify }\PYG{o}{\PYGZlt{}}\PYG{n}{AtC}\PYG{+w}{ }\PYG{n}{fixID}\PYG{o}{\PYGZgt{}}\PYG{+w}{ }\PYG{n}{decomposition}\PYG{+w}{ }\PYG{o}{\PYGZlt{}}\PYG{n}{type}\PYG{o}{\PYGZgt{}}
\end{sphinxVerbatim}
\begin{itemize}
\item {} 
\sphinxAtStartPar
AtC fixID = ID of {\hyperref[\detokenize{fix_atc::doc}]{\sphinxcrossref{\DUrole{doc}{fix atc}}}} instance

\item {} 
\sphinxAtStartPar
decomposition = name of the AtC sub\sphinxhyphen{}command

\item {} 
\sphinxAtStartPar
type =  \sphinxstyleemphasis{replicated\_memory} or \sphinxstyleemphasis{distributed\_memory}

\end{itemize}


\subsubsection{Examples}
\label{\detokenize{atc_decomposition:examples}}
\begin{sphinxVerbatim}[commandchars=\\\{\}]
\PYG{k}{fix\PYGZus{}modify }\PYG{n+nv+nvIdentifier}{AtC}\PYG{+w}{ }\PYG{n}{decomposition}\PYG{+w}{ }\PYG{n}{distributed\PYGZus{}memory}
\end{sphinxVerbatim}


\subsubsection{Description}
\label{\detokenize{atc_decomposition:description}}
\sphinxAtStartPar
Command for assigning the distribution of work and memory for parallel
runs.  With \sphinxstyleemphasis{replicated\_memory} the nodal information is replicated on
each processor, and with \sphinxstyleemphasis{distributed\_memory} only the owned nodal
information kept on each processor.  The \sphinxstyleemphasis{replicated\_memory} option
is most appropriate for simulations were the number of nodes is much
smaller than the number of atoms.


\subsubsection{Restrictions}
\label{\detokenize{atc_decomposition:restrictions}}
\sphinxAtStartPar
None.


\subsubsection{Related AtC commands}
\label{\detokenize{atc_decomposition:related-atc-commands}}\begin{itemize}
\item {} 
\sphinxAtStartPar
{\hyperref[\detokenize{fix_atc:atc-fix-modify}]{\sphinxcrossref{\DUrole{std,std-ref}{fix\_modify AtC command overview}}}}

\end{itemize}


\subsubsection{Default}
\label{\detokenize{atc_decomposition:default}}
\sphinxAtStartPar
replicated\_memory

\sphinxstepscope

\index{fix\_modify AtC extrinsic electron\_integration@\spxentry{fix\_modify AtC extrinsic electron\_integration}}\ignorespaces 

\subsection{fix\_modify AtC extrinsic electron\_integration command}
\label{\detokenize{atc_electron_integration:fix-modify-atc-extrinsic-electron-integration-command}}\label{\detokenize{atc_electron_integration:index-0}}\label{\detokenize{atc_electron_integration::doc}}

\subsubsection{Syntax}
\label{\detokenize{atc_electron_integration:syntax}}
\begin{sphinxVerbatim}[commandchars=\\\{\}]
\PYG{k}{fix\PYGZus{}modify }\PYG{o}{\PYGZlt{}}\PYG{n}{AtC}\PYG{+w}{ }\PYG{n}{fixID}\PYG{o}{\PYGZgt{}}\PYG{+w}{ }\PYG{n}{extrinsic}\PYG{+w}{ }\PYG{n}{electron\PYGZus{}integration}\PYG{+w}{ }\PYG{o}{\PYGZlt{}}\PYG{n}{integration\PYGZus{}type}\PYG{o}{\PYGZgt{}}\PYG{+w}{ }\PYG{n}{[}\PYG{o}{\PYGZlt{}}\PYG{n}{num\PYGZus{}subcycle\PYGZus{}steps}\PYG{o}{\PYGZgt{}}\PYG{n}{]}
\end{sphinxVerbatim}
\begin{itemize}
\item {} 
\sphinxAtStartPar
AtC fixID = ID of {\hyperref[\detokenize{fix_atc::doc}]{\sphinxcrossref{\DUrole{doc}{fix atc}}}} instance

\item {} 
\sphinxAtStartPar
extrinsic electron\_integration = name of the AtC sub\sphinxhyphen{}command

\item {} 
\sphinxAtStartPar
integration\_type = \sphinxstyleemphasis{explicit} or \sphinxstyleemphasis{implicit} or \sphinxstyleemphasis{steady}

\item {} 
\sphinxAtStartPar
num\_subcycle\_steps = number of subcycle steps for the electron time integration (optional)

\end{itemize}


\subsubsection{Examples}
\label{\detokenize{atc_electron_integration:examples}}
\begin{sphinxVerbatim}[commandchars=\\\{\}]
\PYG{k}{fix\PYGZus{}modify }\PYG{n+nv+nvIdentifier}{AtC}\PYG{+w}{ }\PYG{n}{extrinsic}\PYG{+w}{ }\PYG{n}{electron\PYGZus{}integration}\PYG{+w}{ }\PYG{n}{implicit}
\PYG{k}{fix\PYGZus{}modify }\PYG{n+nv+nvIdentifier}{AtC}\PYG{+w}{ }\PYG{n}{extrinsic}\PYG{+w}{ }\PYG{n}{electron\PYGZus{}integration}\PYG{+w}{ }\PYG{n}{explicit}\PYG{+w}{ }\PYG{l+m}{100}
\end{sphinxVerbatim}


\subsubsection{Description}
\label{\detokenize{atc_electron_integration:description}}
\sphinxAtStartPar
Switches between integration schemes for the electron temperature. The
number of subcycling steps used to integrate the electron temperature for
one LAMMPS timestep can be manually adjusted to capture fast electron
dynamics.


\subsubsection{Restrictions}
\label{\detokenize{atc_electron_integration:restrictions}}
\sphinxAtStartPar
For use only with the two\_temperature type of the AtC fix (see
{\hyperref[\detokenize{fix_atc::doc}]{\sphinxcrossref{\DUrole{doc}{fix atc}}}} command)


\subsubsection{Related AtC commands}
\label{\detokenize{atc_electron_integration:related-atc-commands}}\begin{itemize}
\item {} 
\sphinxAtStartPar
{\hyperref[\detokenize{fix_atc:atc-fix-modify}]{\sphinxcrossref{\DUrole{std,std-ref}{fix\_modify AtC command overview}}}}

\end{itemize}


\subsubsection{Default}
\label{\detokenize{atc_electron_integration:default}}
\sphinxAtStartPar
\sphinxstyleemphasis{implicit} and \sphinxstyleemphasis{subcycle\_steps} = 1

\sphinxstepscope

\index{fix\_modify AtC equilibrium\_start@\spxentry{fix\_modify AtC equilibrium\_start}}\ignorespaces 

\subsection{fix\_modify AtC equilibrium\_start command}
\label{\detokenize{atc_equilibrium_start:fix-modify-atc-equilibrium-start-command}}\label{\detokenize{atc_equilibrium_start:index-0}}\label{\detokenize{atc_equilibrium_start::doc}}

\subsubsection{Syntax}
\label{\detokenize{atc_equilibrium_start:syntax}}
\begin{sphinxVerbatim}[commandchars=\\\{\}]
\PYG{k}{fix\PYGZus{}modify }\PYG{o}{\PYGZlt{}}\PYG{n}{AtC}\PYG{+w}{ }\PYG{n}{fixID}\PYG{o}{\PYGZgt{}}\PYG{+w}{ }\PYG{n}{equilibrium\PYGZus{}start}\PYG{+w}{ }\PYG{o}{\PYGZlt{}}\PYG{n}{on}\PYG{o}{|}\PYG{n}{off}\PYG{o}{\PYGZgt{}}
\end{sphinxVerbatim}
\begin{itemize}
\item {} 
\sphinxAtStartPar
AtC fixID = ID of {\hyperref[\detokenize{fix_atc::doc}]{\sphinxcrossref{\DUrole{doc}{fix atc}}}} instance

\item {} 
\sphinxAtStartPar
equilibrium\_start = name of the AtC sub\sphinxhyphen{}command

\item {} 
\sphinxAtStartPar
\sphinxstyleemphasis{exponential} or \sphinxstyleemphasis{step} or \sphinxstyleemphasis{no\_filter} = select type of filter

\end{itemize}


\subsubsection{Examples}
\label{\detokenize{atc_equilibrium_start:examples}}
\begin{sphinxVerbatim}[commandchars=\\\{\}]
\PYG{k}{fix\PYGZus{}modify }\PYG{n+nv+nvIdentifier}{AtC}\PYG{+w}{ }\PYG{n}{equilibrium\PYGZus{}start}\PYG{+w}{ }\PYG{n}{on}
\end{sphinxVerbatim}


\subsubsection{Description}
\label{\detokenize{atc_equilibrium_start:description}}
\sphinxAtStartPar
Starts filtered calculations assuming they start in equilibrium,
i.e. perfect finite element force balance.


\subsubsection{Restrictions}
\label{\detokenize{atc_equilibrium_start:restrictions}}
\sphinxAtStartPar
Only for use with these specific transfers: thermal, two\_temperature


\subsubsection{Related AtC commands}
\label{\detokenize{atc_equilibrium_start:related-atc-commands}}\begin{itemize}
\item {} 
\sphinxAtStartPar
{\hyperref[\detokenize{fix_atc:atc-fix-modify}]{\sphinxcrossref{\DUrole{std,std-ref}{fix\_modify AtC command overview}}}}

\item {} 
\sphinxAtStartPar
{\hyperref[\detokenize{atc_time_filter::doc}]{\sphinxcrossref{\DUrole{doc}{fix\_modify AtC filter}}}}

\item {} 
\sphinxAtStartPar
{\hyperref[\detokenize{atc_filter_scale::doc}]{\sphinxcrossref{\DUrole{doc}{fix\_modify AtC filter scale}}}}

\end{itemize}


\subsubsection{Default}
\label{\detokenize{atc_equilibrium_start:default}}
\sphinxAtStartPar
None.

\sphinxstepscope

\index{fix\_modify AtC extrinsic exchange@\spxentry{fix\_modify AtC extrinsic exchange}}\ignorespaces 

\subsection{fix\_modify AtC extrinsic exchange command}
\label{\detokenize{atc_extrinsic_exchange:fix-modify-atc-extrinsic-exchange-command}}\label{\detokenize{atc_extrinsic_exchange:index-0}}\label{\detokenize{atc_extrinsic_exchange::doc}}

\subsubsection{Syntax}
\label{\detokenize{atc_extrinsic_exchange:syntax}}
\begin{sphinxVerbatim}[commandchars=\\\{\}]
\PYG{k}{fix\PYGZus{}modify }\PYG{o}{\PYGZlt{}}\PYG{n}{AtC}\PYG{+w}{ }\PYG{n}{fixID}\PYG{o}{\PYGZgt{}}\PYG{+w}{ }\PYG{n}{extrinsic}\PYG{+w}{ }\PYG{n}{exchange}\PYG{+w}{ }\PYG{o}{\PYGZlt{}}\PYG{n}{on}\PYG{o}{|}\PYG{n}{off}\PYG{o}{\PYGZgt{}}
\end{sphinxVerbatim}
\begin{itemize}
\item {} 
\sphinxAtStartPar
AtC fixID = ID of {\hyperref[\detokenize{fix_atc::doc}]{\sphinxcrossref{\DUrole{doc}{fix atc}}}} instance

\item {} 
\sphinxAtStartPar
extrinsic exchange = name of the AtC sub\sphinxhyphen{}command

\item {} 
\sphinxAtStartPar
\sphinxstyleemphasis{on} or \sphinxstyleemphasis{off} = set state of energy exchange

\end{itemize}


\subsubsection{Examples}
\label{\detokenize{atc_extrinsic_exchange:examples}}
\begin{sphinxVerbatim}[commandchars=\\\{\}]
\PYG{k}{fix\PYGZus{}modify }\PYG{n+nv+nvIdentifier}{AtC}\PYG{+w}{ }\PYG{n}{extrinsic}\PYG{+w}{ }\PYG{n}{exchange}\PYG{+w}{ }\PYG{n}{on}
\end{sphinxVerbatim}


\subsubsection{Description}
\label{\detokenize{atc_extrinsic_exchange:description}}
\sphinxAtStartPar
Switches energy exchange between the MD system and the electron system
on or off


\subsubsection{Restrictions}
\label{\detokenize{atc_extrinsic_exchange:restrictions}}
\sphinxAtStartPar
For use only with the two\_temperature type of the AtC fix (see
{\hyperref[\detokenize{fix_atc::doc}]{\sphinxcrossref{\DUrole{doc}{fix atc}}}} command)


\subsubsection{Related AtC commands}
\label{\detokenize{atc_extrinsic_exchange:related-atc-commands}}\begin{itemize}
\item {} 
\sphinxAtStartPar
{\hyperref[\detokenize{fix_atc:atc-fix-modify}]{\sphinxcrossref{\DUrole{std,std-ref}{fix\_modify AtC command overview}}}}

\end{itemize}


\subsubsection{Default}
\label{\detokenize{atc_extrinsic_exchange:default}}
\sphinxAtStartPar
\sphinxstyleemphasis{on}

\sphinxstepscope

\index{fix\_modify AtC fe\_md\_boundary@\spxentry{fix\_modify AtC fe\_md\_boundary}}\ignorespaces 

\subsection{fix\_modify AtC fe\_md\_boundary command}
\label{\detokenize{atc_fe_md_boundary:fix-modify-atc-fe-md-boundary-command}}\label{\detokenize{atc_fe_md_boundary:index-0}}\label{\detokenize{atc_fe_md_boundary::doc}}

\subsubsection{Syntax}
\label{\detokenize{atc_fe_md_boundary:syntax}}
\begin{sphinxVerbatim}[commandchars=\\\{\}]
\PYG{k}{fix\PYGZus{}modify }\PYG{o}{\PYGZlt{}}\PYG{n}{AtC}\PYG{+w}{ }\PYG{n}{fixID}\PYG{o}{\PYGZgt{}}\PYG{+w}{ }\PYG{n}{fe\PYGZus{}md\PYGZus{}boundary}\PYG{+w}{ }\PYG{o}{\PYGZlt{}}\PYG{n}{faceset}\PYG{o}{|}\PYG{n}{interpolate}\PYG{o}{|}\PYG{n}{no\PYGZus{}boundary}\PYG{o}{\PYGZgt{}}
\end{sphinxVerbatim}
\begin{itemize}
\item {} 
\sphinxAtStartPar
AtC fixID = ID of {\hyperref[\detokenize{fix_atc::doc}]{\sphinxcrossref{\DUrole{doc}{fix atc}}}} instance

\item {} 
\sphinxAtStartPar
fe\_md\_boundary = name of the AtC sub\sphinxhyphen{}command

\item {} 
\sphinxAtStartPar
\sphinxstyleemphasis{faceset} or \sphinxstyleemphasis{interpolate} or \sphinxstyleemphasis{no\_boundary}

\end{itemize}


\subsubsection{Examples}
\label{\detokenize{atc_fe_md_boundary:examples}}
\begin{sphinxVerbatim}[commandchars=\\\{\}]
\PYG{k}{fix\PYGZus{}modify }\PYG{n+nv+nvIdentifier}{AtC}\PYG{+w}{ }\PYG{n}{fe\PYGZus{}md\PYGZus{}boundary}\PYG{+w}{ }\PYG{n}{interpolate}
\end{sphinxVerbatim}


\subsubsection{Description}
\label{\detokenize{atc_fe_md_boundary:description}}
\sphinxAtStartPar
Specifies different methods for computing fluxes between between the MD
and FE integration regions.  Faceset defines a faceset separating the MD
and FE regions and uses finite element face quadrature to compute the
flux.  Interpolate uses a reconstruction scheme to approximate the flux,
which is more robust but less accurate if the MD/FE boundary does
correspond to a faceset.  No boundary results in no fluxes between the
systems being computed.


\subsubsection{Restrictions}
\label{\detokenize{atc_fe_md_boundary:restrictions}}
\sphinxAtStartPar
If \sphinxstyleemphasis{faceset} is used, all the AtC non\sphinxhyphen{}boundary atoms must lie within and
completely fill the domain enclosed by the faceset.


\subsubsection{Related AtC commands}
\label{\detokenize{atc_fe_md_boundary:related-atc-commands}}\begin{itemize}
\item {} 
\sphinxAtStartPar
{\hyperref[\detokenize{fix_atc:atc-fix-modify}]{\sphinxcrossref{\DUrole{std,std-ref}{fix\_modify AtC command overview}}}}

\item {} 
\sphinxAtStartPar
{\hyperref[\detokenize{atc_mesh_create_faceset_box::doc}]{\sphinxcrossref{\DUrole{doc}{fix\_modify AtC mesh create\_faceset box}}}}

\item {} 
\sphinxAtStartPar
{\hyperref[\detokenize{atc_mesh_create_faceset_plane::doc}]{\sphinxcrossref{\DUrole{doc}{fix\_modify AtC mesh create\_faceset plane}}}}

\end{itemize}


\subsubsection{Default}
\label{\detokenize{atc_fe_md_boundary:default}}
\sphinxAtStartPar
\sphinxstyleemphasis{interpolate}

\sphinxstepscope

\index{fix\_modify AtC filter scale@\spxentry{fix\_modify AtC filter scale}}\ignorespaces 

\subsection{fix\_modify AtC filter scale command}
\label{\detokenize{atc_filter_scale:fix-modify-atc-filter-scale-command}}\label{\detokenize{atc_filter_scale:index-0}}\label{\detokenize{atc_filter_scale::doc}}

\subsubsection{Syntax}
\label{\detokenize{atc_filter_scale:syntax}}
\begin{sphinxVerbatim}[commandchars=\\\{\}]
\PYG{k}{fix\PYGZus{}modify }\PYG{o}{\PYGZlt{}}\PYG{n}{AtC}\PYG{+w}{ }\PYG{n}{fixID}\PYG{o}{\PYGZgt{}}\PYG{+w}{ }\PYG{n}{filter}\PYG{+w}{ }\PYG{n}{scale}\PYG{+w}{ }\PYG{o}{\PYGZlt{}}\PYG{n}{scale}\PYG{o}{\PYGZgt{}}
\end{sphinxVerbatim}
\begin{itemize}
\item {} 
\sphinxAtStartPar
AtC fixID = ID of {\hyperref[\detokenize{fix_atc::doc}]{\sphinxcrossref{\DUrole{doc}{fix atc}}}} instance

\item {} 
\sphinxAtStartPar
filter scale = name of the AtC sub\sphinxhyphen{}command

\item {} 
\sphinxAtStartPar
scale = characteristic times scale of the filter

\end{itemize}


\subsubsection{Examples}
\label{\detokenize{atc_filter_scale:examples}}
\begin{sphinxVerbatim}[commandchars=\\\{\}]
\PYG{k}{fix\PYGZus{}modify }\PYG{n+nv+nvIdentifier}{AtC}\PYG{+w}{ }\PYG{n}{filter}\PYG{+w}{ }\PYG{n}{scale}\PYG{+w}{ }\PYG{l+m}{10.0}
\end{sphinxVerbatim}


\subsubsection{Description}
\label{\detokenize{atc_filter_scale:description}}
\sphinxAtStartPar
Sets the time scale for MD dynamics filter to construct a more
appropriate continuous field.


\subsubsection{Restrictions}
\label{\detokenize{atc_filter_scale:restrictions}}
\sphinxAtStartPar
Only for use with these specific transfers: thermal, two\_temperature


\subsubsection{Related AtC commands}
\label{\detokenize{atc_filter_scale:related-atc-commands}}\begin{itemize}
\item {} 
\sphinxAtStartPar
{\hyperref[\detokenize{fix_atc:atc-fix-modify}]{\sphinxcrossref{\DUrole{std,std-ref}{fix\_modify AtC command overview}}}}

\item {} 
\sphinxAtStartPar
{\hyperref[\detokenize{atc_time_filter::doc}]{\sphinxcrossref{\DUrole{doc}{fix\_modify AtC filter}}}}

\item {} 
\sphinxAtStartPar
{\hyperref[\detokenize{atc_filter_type::doc}]{\sphinxcrossref{\DUrole{doc}{fix\_modify AtC filter type}}}}

\end{itemize}


\subsubsection{Default}
\label{\detokenize{atc_filter_scale:default}}
\sphinxAtStartPar
0.0

\sphinxstepscope

\index{fix\_modify AtC filter type@\spxentry{fix\_modify AtC filter type}}\ignorespaces 

\subsection{fix\_modify AtC filter type command}
\label{\detokenize{atc_filter_type:fix-modify-atc-filter-type-command}}\label{\detokenize{atc_filter_type:index-0}}\label{\detokenize{atc_filter_type::doc}}

\subsubsection{Syntax}
\label{\detokenize{atc_filter_type:syntax}}
\begin{sphinxVerbatim}[commandchars=\\\{\}]
\PYG{k}{fix\PYGZus{}modify }\PYG{o}{\PYGZlt{}}\PYG{n}{AtC}\PYG{+w}{ }\PYG{n}{fixID}\PYG{o}{\PYGZgt{}}\PYG{+w}{ }\PYG{n}{filter}\PYG{+w}{ }\PYG{n}{type}\PYG{+w}{ }\PYG{o}{\PYGZlt{}}\PYG{n}{exponential}\PYG{o}{|}\PYG{n}{step}\PYG{o}{|}\PYG{n}{no\PYGZus{}filter}\PYG{o}{\PYGZgt{}}
\end{sphinxVerbatim}
\begin{itemize}
\item {} 
\sphinxAtStartPar
AtC fixID = ID of {\hyperref[\detokenize{fix_atc::doc}]{\sphinxcrossref{\DUrole{doc}{fix atc}}}} instance

\item {} 
\sphinxAtStartPar
filter type = name of the AtC sub\sphinxhyphen{}command

\item {} 
\sphinxAtStartPar
\sphinxstyleemphasis{exponential} or \sphinxstyleemphasis{step} or \sphinxstyleemphasis{no\_filter} = select type of filter

\end{itemize}


\subsubsection{Examples}
\label{\detokenize{atc_filter_type:examples}}
\begin{sphinxVerbatim}[commandchars=\\\{\}]
\PYG{k}{fix\PYGZus{}modify }\PYG{n+nv+nvIdentifier}{AtC}\PYG{+w}{ }\PYG{n}{filter}\PYG{+w}{ }\PYG{n}{type}\PYG{+w}{ }\PYG{n}{exponential}
\end{sphinxVerbatim}


\subsubsection{Description}
\label{\detokenize{atc_filter_type:description}}
\sphinxAtStartPar
Specifies the type of time filter used.


\subsubsection{Restrictions}
\label{\detokenize{atc_filter_type:restrictions}}
\sphinxAtStartPar
Only for use with these specific transfers: thermal, two\_temperature


\subsubsection{Related AtC commands}
\label{\detokenize{atc_filter_type:related-atc-commands}}\begin{itemize}
\item {} 
\sphinxAtStartPar
{\hyperref[\detokenize{fix_atc:atc-fix-modify}]{\sphinxcrossref{\DUrole{std,std-ref}{fix\_modify AtC command overview}}}}

\item {} 
\sphinxAtStartPar
{\hyperref[\detokenize{atc_time_filter::doc}]{\sphinxcrossref{\DUrole{doc}{fix\_modify AtC filter}}}}

\item {} 
\sphinxAtStartPar
{\hyperref[\detokenize{atc_filter_scale::doc}]{\sphinxcrossref{\DUrole{doc}{fix\_modify AtC filter scale}}}}

\end{itemize}


\subsubsection{Default}
\label{\detokenize{atc_filter_type:default}}
\sphinxAtStartPar
None.

\sphinxstepscope

\index{fix\_modify AtC fix@\spxentry{fix\_modify AtC fix}}\ignorespaces 

\subsection{fix\_modify AtC fix command}
\label{\detokenize{atc_fix:fix-modify-atc-fix-command}}\label{\detokenize{atc_fix:index-0}}\label{\detokenize{atc_fix::doc}}

\subsubsection{Syntax}
\label{\detokenize{atc_fix:syntax}}
\begin{sphinxVerbatim}[commandchars=\\\{\}]
\PYG{k}{fix\PYGZus{}modify }\PYG{o}{\PYGZlt{}}\PYG{n}{AtC}\PYG{+w}{ }\PYG{n}{fixID}\PYG{o}{\PYGZgt{}}\PYG{+w}{ }\PYG{k}{fix }\PYG{o}{\PYGZlt{}}\PYG{n}{field}\PYG{o}{\PYGZgt{}}\PYG{+w}{ }\PYG{o}{\PYGZlt{}}\PYG{n}{nodeset}\PYG{o}{\PYGZgt{}}\PYG{+w}{ }\PYG{o}{\PYGZlt{}}\PYG{n}{constant}\PYG{o}{|}\PYG{n}{function}\PYG{o}{\PYGZgt{}}
\end{sphinxVerbatim}
\begin{itemize}
\item {} 
\sphinxAtStartPar
AtC fixID = ID of {\hyperref[\detokenize{fix_atc::doc}]{\sphinxcrossref{\DUrole{doc}{fix atc}}}} instance

\item {} 
\sphinxAtStartPar
fix = name of the AtC sub\sphinxhyphen{}command

\item {} 
\sphinxAtStartPar
field = field kind name valid for type of physics: \sphinxstyleemphasis{temperature} or \sphinxstyleemphasis{electron\_temperature}

\item {} 
\sphinxAtStartPar
nodeset = name of set of nodes to apply boundary condition

\item {} 
\sphinxAtStartPar
\sphinxstyleemphasis{constant} or \sphinxstyleemphasis{function} = value or name of function followed by its parameters

\end{itemize}


\subsubsection{Examples}
\label{\detokenize{atc_fix:examples}}
\begin{sphinxVerbatim}[commandchars=\\\{\}]
\PYG{k}{fix\PYGZus{}modify }\PYG{n+nv+nvIdentifier}{AtC}\PYG{+w}{ }\PYG{k}{fix }\PYG{n+nv+nvIdentifier}{temperature}\PYG{+w}{ }\PYG{n+nv+nvIdentifier}{groupNAME}\PYG{+w}{ }\PYG{l+m}{10}\PYG{n}{.}
\PYG{k}{fix\PYGZus{}modify }\PYG{n+nv+nvIdentifier}{AtC}\PYG{+w}{ }\PYG{k}{fix }\PYG{n+nv+nvIdentifier}{temperature}\PYG{+w}{ }\PYG{n+nv+nvIdentifier}{groupNAME}\PYG{+w}{ }\PYG{l+m}{0}\PYG{+w}{ }\PYG{l+m}{0}\PYG{+w}{ }\PYG{l+m}{0}\PYG{+w}{ }\PYG{l+m}{10.0}\PYG{+w}{ }\PYG{l+m}{0}\PYG{+w}{ }\PYG{l+m}{0}\PYG{+w}{ }\PYG{l+m}{1.0}
\end{sphinxVerbatim}


\subsubsection{Description}
\label{\detokenize{atc_fix:description}}
\sphinxAtStartPar
Creates a constraint on the values of the specified field at specified nodes.


\subsubsection{Restrictions}
\label{\detokenize{atc_fix:restrictions}}
\sphinxAtStartPar
The keyword \sphinxstyleemphasis{all} is reserved and thus not available as nodeset name.


\subsubsection{Related AtC commands}
\label{\detokenize{atc_fix:related-atc-commands}}\begin{itemize}
\item {} 
\sphinxAtStartPar
{\hyperref[\detokenize{fix_atc:atc-fix-modify}]{\sphinxcrossref{\DUrole{std,std-ref}{fix\_modify AtC command overview}}}}

\item {} 
\sphinxAtStartPar
{\hyperref[\detokenize{atc_unfix::doc}]{\sphinxcrossref{\DUrole{doc}{fix\_modify AtC unfix}}}}

\end{itemize}


\subsubsection{Default}
\label{\detokenize{atc_fix:default}}
\sphinxAtStartPar
None.

\sphinxstepscope

\index{fix\_modify AtC fix\_flux@\spxentry{fix\_modify AtC fix\_flux}}\ignorespaces 

\subsection{fix\_modify AtC fix\_flux command}
\label{\detokenize{atc_fix_flux:fix-modify-atc-fix-flux-command}}\label{\detokenize{atc_fix_flux:index-0}}\label{\detokenize{atc_fix_flux::doc}}

\subsubsection{Syntax}
\label{\detokenize{atc_fix_flux:syntax}}
\begin{sphinxVerbatim}[commandchars=\\\{\}]
\PYG{k}{fix\PYGZus{}modify }\PYG{o}{\PYGZlt{}}\PYG{n}{AtC}\PYG{+w}{ }\PYG{n}{fixID}\PYG{o}{\PYGZgt{}}\PYG{+w}{ }\PYG{n}{fix\PYGZus{}flux}\PYG{+w}{ }\PYG{o}{\PYGZlt{}}\PYG{n}{field}\PYG{o}{\PYGZgt{}}\PYG{+w}{ }\PYG{o}{\PYGZlt{}}\PYG{n}{face\PYGZus{}set}\PYG{o}{\PYGZgt{}}\PYG{+w}{ }\PYG{o}{\PYGZlt{}}\PYG{n}{value}\PYG{o}{|}\PYG{n}{function}\PYG{o}{\PYGZgt{}}
\end{sphinxVerbatim}
\begin{itemize}
\item {} 
\sphinxAtStartPar
AtC fixID = ID of {\hyperref[\detokenize{fix_atc::doc}]{\sphinxcrossref{\DUrole{doc}{fix atc}}}} instance

\item {} 
\sphinxAtStartPar
fix\_flux = name of the AtC sub\sphinxhyphen{}command

\item {} 
\sphinxAtStartPar
field = field kind name valid for type of physics: \sphinxstyleemphasis{temperature} or \sphinxstyleemphasis{electron\_temperature}

\item {} 
\sphinxAtStartPar
face\_set = name of set of element faces

\item {} 
\sphinxAtStartPar
\sphinxstyleemphasis{value} or \sphinxstyleemphasis{function} = value or name of function followed by its parameters

\end{itemize}


\subsubsection{Examples}
\label{\detokenize{atc_fix_flux:examples}}
\begin{sphinxVerbatim}[commandchars=\\\{\}]
\PYG{k}{fix\PYGZus{}modify }\PYG{n+nv+nvIdentifier}{AtC}\PYG{+w}{ }\PYG{n}{fix\PYGZus{}flux}\PYG{+w}{ }\PYG{n}{temperature}\PYG{+w}{ }\PYG{n}{faceSet}\PYG{+w}{ }\PYG{l+m}{10.0}
\end{sphinxVerbatim}


\subsubsection{Description}
\label{\detokenize{atc_fix_flux:description}}
\sphinxAtStartPar
Command for fixing normal fluxes e.g. heat\_flux. This command only
prescribes the normal component of the physical flux, e.g. heat (energy)
flux. The units are in AtC units, i.e. derived from the LAMMPS length,
time, and mass scales.


\subsubsection{Restrictions}
\label{\detokenize{atc_fix_flux:restrictions}}
\sphinxAtStartPar
Only normal fluxes (Neumann data) can be prescribed.


\subsubsection{Related AtC commands}
\label{\detokenize{atc_fix_flux:related-atc-commands}}\begin{itemize}
\item {} 
\sphinxAtStartPar
{\hyperref[\detokenize{fix_atc:atc-fix-modify}]{\sphinxcrossref{\DUrole{std,std-ref}{fix\_modify AtC command overview}}}}

\item {} 
\sphinxAtStartPar
{\hyperref[\detokenize{atc_unfix_flux::doc}]{\sphinxcrossref{\DUrole{doc}{fix\_modify AtC unfix\_flux}}}}

\end{itemize}


\subsubsection{Default}
\label{\detokenize{atc_fix_flux:default}}
\sphinxAtStartPar
None.

\sphinxstepscope

\index{fix\_modify AtC computes@\spxentry{fix\_modify AtC computes}}\ignorespaces 

\subsection{fix\_modify AtC computes command}
\label{\detokenize{atc_hardy_computes:fix-modify-atc-computes-command}}\label{\detokenize{atc_hardy_computes:index-0}}\label{\detokenize{atc_hardy_computes::doc}}

\subsubsection{Syntax}
\label{\detokenize{atc_hardy_computes:syntax}}
\begin{sphinxVerbatim}[commandchars=\\\{\}]
\PYG{k}{fix\PYGZus{}modify }\PYG{o}{\PYGZlt{}}\PYG{n}{AtC}\PYG{+w}{ }\PYG{n}{fixID}\PYG{o}{\PYGZgt{}}\PYG{+w}{ }\PYG{n}{computes}\PYG{+w}{ }\PYG{o}{\PYGZlt{}}\PYG{n}{add}\PYG{o}{|}\PYG{n}{delete}\PYG{o}{\PYGZgt{}}\PYG{+w}{ }\PYG{o}{\PYGZlt{}}\PYG{n}{per}\PYG{o}{\PYGZhy{}}\PYG{n}{atom}\PYG{+w}{ }\PYG{n}{compute}\PYG{o}{\PYGZhy{}}\PYG{n}{ID}\PYG{o}{\PYGZgt{}}\PYG{+w}{ }\PYG{o}{\PYGZlt{}}\PYG{n}{volume}\PYG{o}{|}\PYG{n}{number}\PYG{o}{\PYGZgt{}}
\end{sphinxVerbatim}
\begin{itemize}
\item {} 
\sphinxAtStartPar
AtC fixID = ID of {\hyperref[\detokenize{fix_atc::doc}]{\sphinxcrossref{\DUrole{doc}{fix atc}}}} instance

\item {} 
\sphinxAtStartPar
computes = name of the AtC sub\sphinxhyphen{}command

\item {} 
\sphinxAtStartPar
\sphinxstyleemphasis{add} or \sphinxstyleemphasis{delete} = add or delete the calculation of an equivalent continuum field for the specified per\sphinxhyphen{}atom compute as volume or number density quantity

\item {} 
\sphinxAtStartPar
per\sphinxhyphen{}atom compute\sphinxhyphen{}ID = ID of a per\sphinxhyphen{}atom compute; fields can be calculated for all per\sphinxhyphen{}atom computes available in LAMMPS

\item {} 
\sphinxAtStartPar
\sphinxstyleemphasis{volume} or \sphinxstyleemphasis{number} = select whether the created field is a per\sphinxhyphen{}unit\sphinxhyphen{}volume quantity or a per\sphinxhyphen{}atom quantity as weighted by kernel functions

\end{itemize}


\subsubsection{Examples}
\label{\detokenize{atc_hardy_computes:examples}}
\begin{sphinxVerbatim}[commandchars=\\\{\}]
\PYG{k}{compute }\PYG{n+nv+nvIdentifier}{virial}\PYG{+w}{ }\PYG{n+nv+nvIdentifier}{all}\PYG{+w}{ }\PYG{n}{stress}\PYG{o}{/}\PYG{n}{atom}
\PYG{k}{fix\PYGZus{}modify }\PYG{n+nv+nvIdentifier}{AtC}\PYG{+w}{ }\PYG{n}{computes}\PYG{+w}{ }\PYG{n}{add}\PYG{+w}{ }\PYG{n}{virial}\PYG{+w}{ }\PYG{n}{volume}
\PYG{k}{fix\PYGZus{}modify }\PYG{n+nv+nvIdentifier}{AtC}\PYG{+w}{ }\PYG{n}{computes}\PYG{+w}{ }\PYG{n}{delete}\PYG{+w}{ }\PYG{n}{virial}

\PYG{k}{compute }\PYG{n+nv+nvIdentifier}{centrosymmetry}\PYG{+w}{ }\PYG{n+nv+nvIdentifier}{all}\PYG{+w}{ }\PYG{n}{centro}\PYG{o}{/}\PYG{n}{atom}
\PYG{k}{fix\PYGZus{}modify }\PYG{n+nv+nvIdentifier}{AtC}\PYG{+w}{ }\PYG{n}{computes}\PYG{+w}{ }\PYG{n}{add}\PYG{+w}{ }\PYG{n}{centrosymmetry}\PYG{+w}{ }\PYG{n}{number}
\end{sphinxVerbatim}


\subsubsection{Description}
\label{\detokenize{atc_hardy_computes:description}}
\sphinxAtStartPar
Calculates continuum fields corresponding to specified per\sphinxhyphen{}atom
{\hyperref[\detokenize{compute::doc}]{\sphinxcrossref{\DUrole{doc}{computes}}}} created by LAMMPS.


\subsubsection{Restrictions}
\label{\detokenize{atc_hardy_computes:restrictions}}
\sphinxAtStartPar
Must be used with {\hyperref[\detokenize{fix_atc::doc}]{\sphinxcrossref{\DUrole{doc}{fix atc hardy}}}}.  The per\sphinxhyphen{}atom compute
must be specified before the corresponding continuum field can be
requested.


\subsubsection{Related AtC commands}
\label{\detokenize{atc_hardy_computes:related-atc-commands}}\begin{itemize}
\item {} 
\sphinxAtStartPar
{\hyperref[\detokenize{fix_atc:atc-fix-modify}]{\sphinxcrossref{\DUrole{std,std-ref}{fix\_modify AtC command overview}}}}

\item {} 
\sphinxAtStartPar
{\hyperref[\detokenize{atc_hardy_fields::doc}]{\sphinxcrossref{\DUrole{doc}{fix\_modify AtC fields}}}}

\item {} 
\sphinxAtStartPar
{\hyperref[\detokenize{compute::doc}]{\sphinxcrossref{\DUrole{doc}{compute}}}}

\end{itemize}


\subsubsection{Default}
\label{\detokenize{atc_hardy_computes:default}}
\sphinxAtStartPar
None.

\sphinxstepscope

\index{fix\_modify AtC fields@\spxentry{fix\_modify AtC fields}}\ignorespaces 

\subsection{fix\_modify AtC fields command}
\label{\detokenize{atc_hardy_fields:fix-modify-atc-fields-command}}\label{\detokenize{atc_hardy_fields:index-0}}\label{\detokenize{atc_hardy_fields::doc}}

\subsubsection{Syntax}
\label{\detokenize{atc_hardy_fields:syntax}}
\begin{sphinxVerbatim}[commandchars=\\\{\}]
\PYG{k}{fix\PYGZus{}modify }\PYG{o}{\PYGZlt{}}\PYG{n}{AtC}\PYG{+w}{ }\PYG{n}{fixID}\PYG{o}{\PYGZgt{}}\PYG{+w}{ }\PYG{n}{fields}\PYG{+w}{ }\PYG{o}{\PYGZlt{}}\PYG{n}{all}\PYG{o}{|}\PYG{n}{none}\PYG{o}{\PYGZgt{}}
\PYG{k}{fix\PYGZus{}modify }\PYG{o}{\PYGZlt{}}\PYG{n}{AtC}\PYG{+w}{ }\PYG{n}{fixID}\PYG{o}{\PYGZgt{}}\PYG{+w}{ }\PYG{n}{fields}\PYG{+w}{ }\PYG{o}{\PYGZlt{}}\PYG{n}{add}\PYG{o}{|}\PYG{n}{delete}\PYG{o}{\PYGZgt{}}\PYG{+w}{ }\PYG{o}{\PYGZlt{}}\PYG{n}{list\PYGZus{}of\PYGZus{}fields}\PYG{o}{\PYGZgt{}}
\end{sphinxVerbatim}
\begin{itemize}
\item {} 
\sphinxAtStartPar
AtC fixID = ID of {\hyperref[\detokenize{fix_atc::doc}]{\sphinxcrossref{\DUrole{doc}{fix atc}}}} instance

\item {} 
\sphinxAtStartPar
fields = name of the AtC sub\sphinxhyphen{}command

\item {} 
\sphinxAtStartPar
\sphinxstyleemphasis{all} or \sphinxstyleemphasis{none} = output all or no fields

\item {} 
\sphinxAtStartPar
\sphinxstyleemphasis{add} or \sphinxstyleemphasis{delete} = add or delete the listed output fields

\item {} 
\sphinxAtStartPar
list\_of\_fields = one or more of the fields listed below:
\begin{itemize}
\item {} 
\sphinxAtStartPar
density : mass per unit volume

\item {} 
\sphinxAtStartPar
displacement : displacement vector

\item {} 
\sphinxAtStartPar
momentum : momentum per unit volume

\item {} 
\sphinxAtStartPar
velocity : defined by momentum divided by density

\item {} 
\sphinxAtStartPar
projected\_velocity : simple kernel estimation of atomic velocities

\item {} 
\sphinxAtStartPar
temperature : temperature derived from the relative atomic kinetic energy

\item {} 
\sphinxAtStartPar
kinetic\_temperature : temperature derived from the full kinetic energy

\item {} 
\sphinxAtStartPar
number\_density : simple kernel estimation of number of atoms per unit volume

\item {} 
\sphinxAtStartPar
stress : Cauchy stress tensor for eulerian analysis (atom\_element\_map), or first Piola\sphinxhyphen{}Kirchhoff stress tensor for lagrangian analysis

\item {} 
\sphinxAtStartPar
transformed\_stress : first Piola\sphinxhyphen{}Kirchhoff stress tensor for eulerian analysis (atom\_element\_map), or Cauchy stress tensor for lagrangian analysis

\item {} 
\sphinxAtStartPar
heat\_flux : spatial heat flux vector for eulerian, or referential heat flux vector for lagrangian

\item {} 
\sphinxAtStartPar
potential\_energy : potential energy per unit volume

\item {} 
\sphinxAtStartPar
kinetic\_energy : kinetic energy per unit volume

\item {} 
\sphinxAtStartPar
thermal\_energy : thermal energy (kinetic energy \sphinxhyphen{} continuum kinetic energy) per unit volume

\item {} 
\sphinxAtStartPar
internal\_energy : total internal energy (potential + thermal) per unit volume

\item {} 
\sphinxAtStartPar
energy : total energy (potential + kinetic) per unit volume

\item {} 
\sphinxAtStartPar
number\_density : number of atoms per unit volume

\item {} 
\sphinxAtStartPar
eshelby\_stress : configurational stress (energy\sphinxhyphen{}momentum) tensor defined by {\hyperref[\detokenize{atc_hardy_fields:eshelby1}]{\sphinxcrossref{\DUrole{std,std-ref}{(Eshelby)}}}}

\item {} 
\sphinxAtStartPar
vacancy\_concentration : volume fraction of vacancy content

\item {} 
\sphinxAtStartPar
type\_concentration : volume fraction of a specific atom type

\end{itemize}

\end{itemize}


\subsubsection{Examples}
\label{\detokenize{atc_hardy_fields:examples}}
\begin{sphinxVerbatim}[commandchars=\\\{\}]
\PYG{k}{fix\PYGZus{}modify }\PYG{n+nv+nvIdentifier}{AtC}\PYG{+w}{ }\PYG{n}{fields}\PYG{+w}{ }\PYG{n}{add}\PYG{+w}{ }\PYG{k}{velocity }\PYG{n+nv+nvIdentifier}{temperature}
\end{sphinxVerbatim}


\subsubsection{Description}
\label{\detokenize{atc_hardy_fields:description}}
\sphinxAtStartPar
Allows modification of the fields calculated and output by the AtC
transfer class.  The commands are cumulative, e.g.:

\begin{sphinxVerbatim}[commandchars=\\\{\}]
\PYG{k}{fix\PYGZus{}modify }\PYG{n+nv+nvIdentifier}{AtC}\PYG{+w}{ }\PYG{n}{fields}\PYG{+w}{ }\PYG{n}{none}
\PYG{k}{fix\PYGZus{}modify }\PYG{n+nv+nvIdentifier}{AtC}\PYG{+w}{ }\PYG{n}{fields}\PYG{+w}{ }\PYG{n}{add}\PYG{+w}{ }\PYG{k}{velocity }\PYG{n+nv+nvIdentifier}{temperature}
\end{sphinxVerbatim}

\sphinxAtStartPar
will only output the velocity and temperature fields.


\subsubsection{Restrictions}
\label{\detokenize{atc_hardy_fields:restrictions}}
\sphinxAtStartPar
Must be used with {\hyperref[\detokenize{fix_atc::doc}]{\sphinxcrossref{\DUrole{doc}{fix atc hardy}}}}.  Currently, the stress
and heat flux formulas are only correct for central force potentials,
e.g. Lennard\sphinxhyphen{}Jones and EAM but not Stillinger\sphinxhyphen{}Weber.


\subsubsection{Related AtC commands}
\label{\detokenize{atc_hardy_fields:related-atc-commands}}\begin{itemize}
\item {} 
\sphinxAtStartPar
{\hyperref[\detokenize{fix_atc:atc-fix-modify}]{\sphinxcrossref{\DUrole{std,std-ref}{fix\_modify AtC command overview}}}}

\item {} 
\sphinxAtStartPar
{\hyperref[\detokenize{atc_hardy_gradients::doc}]{\sphinxcrossref{\DUrole{doc}{fix\_modify AtC gradients}}}}

\item {} 
\sphinxAtStartPar
{\hyperref[\detokenize{atc_hardy_rates::doc}]{\sphinxcrossref{\DUrole{doc}{fix\_modify AtC rates}}}}

\item {} 
\sphinxAtStartPar
{\hyperref[\detokenize{atc_hardy_computes::doc}]{\sphinxcrossref{\DUrole{doc}{fix\_modify AtC computes}}}}

\end{itemize}


\subsubsection{Default}
\label{\detokenize{atc_hardy_fields:default}}
\sphinxAtStartPar
By default, no fields are output.


\subsubsection{References}
\label{\detokenize{atc_hardy_fields:references}}\phantomsection\label{\detokenize{atc_hardy_fields:eshelby1}}
\sphinxAtStartPar
\sphinxstylestrong{(Eshelby)} J.D. Eshelby, Philos. Trans. Royal Soc. London A, Math. Phys. Sci., Vol. 244, No. 877 (1951) pp. 87\sphinxhyphen{}112; J. Elasticity, Vol. 5, Nos. 3\sphinxhyphen{}4 (1975) pp. 321\sphinxhyphen{}335{]}

\sphinxstepscope

\index{fix\_modify AtC gradients@\spxentry{fix\_modify AtC gradients}}\ignorespaces 

\subsection{fix\_modify AtC gradients command}
\label{\detokenize{atc_hardy_gradients:fix-modify-atc-gradients-command}}\label{\detokenize{atc_hardy_gradients:index-0}}\label{\detokenize{atc_hardy_gradients::doc}}

\subsubsection{Syntax}
\label{\detokenize{atc_hardy_gradients:syntax}}
\begin{sphinxVerbatim}[commandchars=\\\{\}]
\PYG{k}{fix\PYGZus{}modify }\PYG{o}{\PYGZlt{}}\PYG{n}{AtC}\PYG{+w}{ }\PYG{n}{fixID}\PYG{o}{\PYGZgt{}}\PYG{+w}{ }\PYG{n}{gradients}\PYG{+w}{ }\PYG{o}{\PYGZlt{}}\PYG{n}{add}\PYG{o}{|}\PYG{n}{delete}\PYG{o}{\PYGZgt{}}\PYG{+w}{ }\PYG{o}{\PYGZlt{}}\PYG{n}{list\PYGZus{}of\PYGZus{}fields}\PYG{o}{\PYGZgt{}}
\end{sphinxVerbatim}
\begin{itemize}
\item {} 
\sphinxAtStartPar
AtC fixID = ID of {\hyperref[\detokenize{fix_atc::doc}]{\sphinxcrossref{\DUrole{doc}{fix atc}}}} instance

\item {} 
\sphinxAtStartPar
gradients = name of the AtC sub\sphinxhyphen{}command

\item {} 
\sphinxAtStartPar
\sphinxstyleemphasis{add} or \sphinxstyleemphasis{delete} = select whether to add or delete calculation of gradients for the listed output fields

\item {} 
\sphinxAtStartPar
list\_of\_fields = one or more of the fields listed below:
\begin{itemize}
\item {} 
\sphinxAtStartPar
density : mass per unit volume

\item {} 
\sphinxAtStartPar
displacement : displacement vector

\item {} 
\sphinxAtStartPar
momentum : momentum per unit volume

\item {} 
\sphinxAtStartPar
velocity : defined by momentum divided by density

\item {} 
\sphinxAtStartPar
projected\_velocity : simple kernel estimation of atomic velocities

\item {} 
\sphinxAtStartPar
temperature : temperature derived from the relative atomic kinetic energy

\item {} 
\sphinxAtStartPar
kinetic\_temperature : temperature derived from the full kinetic energy

\item {} 
\sphinxAtStartPar
number\_density : simple kernel estimation of number of atoms per unit volume

\item {} 
\sphinxAtStartPar
stress : Cauchy stress tensor for eulerian analysis (atom\_element\_map), or first Piola\sphinxhyphen{}Kirchhoff stress tensor for lagrangian analysis

\item {} 
\sphinxAtStartPar
transformed\_stress : first Piola\sphinxhyphen{}Kirchhoff stress tensor for eulerian analysis (atom\_element\_map), or Cauchy stress tensor for lagrangian analysis

\item {} 
\sphinxAtStartPar
heat\_flux : spatial heat flux vector for eulerian, or referential heat flux vector for lagrangian

\item {} 
\sphinxAtStartPar
potential\_energy : potential energy per unit volume

\item {} 
\sphinxAtStartPar
kinetic\_energy : kinetic energy per unit volume

\item {} 
\sphinxAtStartPar
thermal\_energy : thermal energy (kinetic energy \sphinxhyphen{} continuum kinetic energy) per unit volume

\item {} 
\sphinxAtStartPar
internal\_energy : total internal energy (potential + thermal) per unit volume

\item {} 
\sphinxAtStartPar
energy : total energy (potential + kinetic) per unit volume

\item {} 
\sphinxAtStartPar
number\_density : number of atoms per unit volume

\item {} 
\sphinxAtStartPar
eshelby\_stress : configurational stress (energy\sphinxhyphen{}momentum) tensor defined by {\hyperref[\detokenize{atc_hardy_gradients:eshelby3}]{\sphinxcrossref{\DUrole{std,std-ref}{(Eshelby)}}}}

\item {} 
\sphinxAtStartPar
vacancy\_concentration : volume fraction of vacancy content

\item {} 
\sphinxAtStartPar
type\_concentration : volume fraction of a specific atom type

\end{itemize}

\end{itemize}


\subsubsection{Examples}
\label{\detokenize{atc_hardy_gradients:examples}}
\begin{sphinxVerbatim}[commandchars=\\\{\}]
\PYG{k}{fix\PYGZus{}modify }\PYG{n+nv+nvIdentifier}{AtC}\PYG{+w}{ }\PYG{n}{gradients}\PYG{+w}{ }\PYG{n}{add}\PYG{+w}{ }\PYG{n}{temperature}\PYG{+w}{ }\PYG{k}{velocity }\PYG{n+nv+nvIdentifier}{stress}
\PYG{k}{fix\PYGZus{}modify }\PYG{n+nv+nvIdentifier}{AtC}\PYG{+w}{ }\PYG{n}{gradients}\PYG{+w}{ }\PYG{n}{delete}\PYG{+w}{ }\PYG{k}{velocity}
\end{sphinxVerbatim}


\subsubsection{Description}
\label{\detokenize{atc_hardy_gradients:description}}
\sphinxAtStartPar
Requests calculation and output of gradients of the fields from the AtC
transfer class.  These gradients will be with regard to spatial or
material coordinate for Eulerian or Lagrangian analysis, respectively,
as specified by {\hyperref[\detokenize{atc_atom_element_map::doc}]{\sphinxcrossref{\DUrole{doc}{fix\_modify AtC atom\_element\_map}}}}


\subsubsection{Restrictions}
\label{\detokenize{atc_hardy_gradients:restrictions}}
\sphinxAtStartPar
Must be used with {\hyperref[\detokenize{fix_atc::doc}]{\sphinxcrossref{\DUrole{doc}{fix atc hardy}}}}.


\subsubsection{Related AtC commands}
\label{\detokenize{atc_hardy_gradients:related-atc-commands}}\begin{itemize}
\item {} 
\sphinxAtStartPar
{\hyperref[\detokenize{fix_atc:atc-fix-modify}]{\sphinxcrossref{\DUrole{std,std-ref}{fix\_modify AtC command overview}}}}

\item {} 
\sphinxAtStartPar
{\hyperref[\detokenize{atc_atom_element_map::doc}]{\sphinxcrossref{\DUrole{doc}{fix\_modify AtC atom\_element\_map}}}}

\item {} 
\sphinxAtStartPar
{\hyperref[\detokenize{atc_hardy_fields::doc}]{\sphinxcrossref{\DUrole{doc}{fix\_modify AtC fields}}}}

\item {} 
\sphinxAtStartPar
{\hyperref[\detokenize{atc_hardy_rates::doc}]{\sphinxcrossref{\DUrole{doc}{fix\_modify AtC rates}}}}

\end{itemize}


\subsubsection{Default}
\label{\detokenize{atc_hardy_gradients:default}}
\sphinxAtStartPar
None.


\subsubsection{References}
\label{\detokenize{atc_hardy_gradients:references}}\phantomsection\label{\detokenize{atc_hardy_gradients:eshelby3}}
\sphinxAtStartPar
\sphinxstylestrong{(Eshelby)} J.D. Eshelby, Philos. Trans. Royal Soc. London A, Math. Phys. Sci., Vol. 244, No. 877 (1951) pp. 87\sphinxhyphen{}112; J. Elasticity, Vol. 5, Nos. 3\sphinxhyphen{}4 (1975) pp. 321\sphinxhyphen{}335{]}

\sphinxstepscope

\index{fix\_modify AtC kernel@\spxentry{fix\_modify AtC kernel}}\ignorespaces 

\subsection{fix\_modify AtC kernel command}
\label{\detokenize{atc_hardy_kernel:fix-modify-atc-kernel-command}}\label{\detokenize{atc_hardy_kernel:index-0}}\label{\detokenize{atc_hardy_kernel::doc}}

\subsubsection{Syntax}
\label{\detokenize{atc_hardy_kernel:syntax}}
\begin{sphinxVerbatim}[commandchars=\\\{\}]
\PYG{k}{fix\PYGZus{}modify }\PYG{o}{\PYGZlt{}}\PYG{n}{AtC}\PYG{+w}{ }\PYG{n}{fixID}\PYG{o}{\PYGZgt{}}\PYG{+w}{ }\PYG{n}{kernel}\PYG{+w}{ }\PYG{o}{\PYGZlt{}}\PYG{n}{type}\PYG{o}{\PYGZgt{}}\PYG{+w}{ }\PYG{o}{\PYGZlt{}}\PYG{n}{parameters}\PYG{o}{\PYGZgt{}}
\end{sphinxVerbatim}
\begin{itemize}
\item {} 
\sphinxAtStartPar
AtC fixID = ID of {\hyperref[\detokenize{fix_atc::doc}]{\sphinxcrossref{\DUrole{doc}{fix atc}}}} instance

\item {} 
\sphinxAtStartPar
kernel = name of the AtC sub\sphinxhyphen{}command

\item {} 
\sphinxAtStartPar
type = \sphinxstyleemphasis{step} or \sphinxstyleemphasis{cell} or \sphinxstyleemphasis{cubic\_bar} or \sphinxstyleemphasis{cubic\_cylinder} or
\sphinxstyleemphasis{cubic\_sphere} or \sphinxstyleemphasis{quartic\_bar} or \sphinxstyleemphasis{quartic\_cylinder} or
\sphinxstyleemphasis{quartic\_sphere}

\item {} 
\sphinxAtStartPar
the following parameter(s) are required for each kernel:
\begin{itemize}
\item {} 
\sphinxAtStartPar
\sphinxstyleemphasis{step} : \textless{}radius\textgreater{}

\item {} 
\sphinxAtStartPar
\sphinxstyleemphasis{cell} : \textless{}hx\textgreater{} \textless{}hy\textgreater{} \textless{}hz\textgreater{} or \textless{}h\textgreater{}

\item {} 
\sphinxAtStartPar
\sphinxstyleemphasis{cubic\_bar} : \textless{}half\_width\textgreater{}

\item {} 
\sphinxAtStartPar
\sphinxstyleemphasis{cubic\_cylinder} : \textless{}radius\textgreater{}

\item {} 
\sphinxAtStartPar
\sphinxstyleemphasis{cubic\_sphere} : \textless{}radius\textgreater{}

\item {} 
\sphinxAtStartPar
\sphinxstyleemphasis{quartic\_bar} : \textless{}half\_width\textgreater{}

\item {} 
\sphinxAtStartPar
\sphinxstyleemphasis{quartic\_cylinder} : \textless{}radius\textgreater{}

\item {} 
\sphinxAtStartPar
\sphinxstyleemphasis{quartic\_sphere} : \textless{}radius\textgreater{}

\end{itemize}

\end{itemize}


\subsubsection{Examples}
\label{\detokenize{atc_hardy_kernel:examples}}
\begin{sphinxVerbatim}[commandchars=\\\{\}]
\PYG{k}{fix\PYGZus{}modify }\PYG{n+nv+nvIdentifier}{AtC}\PYG{+w}{ }\PYG{n}{kernel}\PYG{+w}{ }\PYG{n}{cell}\PYG{+w}{ }\PYG{l+m}{1.0}\PYG{+w}{ }\PYG{l+m}{1.0}\PYG{+w}{ }\PYG{l+m}{1.0}
\PYG{k}{fix\PYGZus{}modify }\PYG{n+nv+nvIdentifier}{AtC}\PYG{+w}{ }\PYG{n}{kernel}\PYG{+w}{ }\PYG{n}{quartic\PYGZus{}sphere}\PYG{+w}{ }\PYG{l+m}{10.0}
\end{sphinxVerbatim}


\subsubsection{Description}
\label{\detokenize{atc_hardy_kernel:description}}
\sphinxAtStartPar
Sets the localization kernel type and parameters for {\hyperref[\detokenize{fix_atc::doc}]{\sphinxcrossref{\DUrole{doc}{fix atc hardy}}}}.


\subsubsection{Restrictions}
\label{\detokenize{atc_hardy_kernel:restrictions}}
\sphinxAtStartPar
Must be used with {\hyperref[\detokenize{fix_atc::doc}]{\sphinxcrossref{\DUrole{doc}{fix atc hardy}}}}.  For bar kernel types,
half\sphinxhyphen{}width oriented along x\sphinxhyphen{}direction.  For cylinder kernel types,
cylindrical axis is assumed to be in z\sphinxhyphen{}direction.


\subsubsection{Related AtC commands}
\label{\detokenize{atc_hardy_kernel:related-atc-commands}}\begin{itemize}
\item {} 
\sphinxAtStartPar
{\hyperref[\detokenize{fix_atc:atc-fix-modify}]{\sphinxcrossref{\DUrole{std,std-ref}{fix\_modify AtC command overview}}}}

\item {} 
\sphinxAtStartPar
{\hyperref[\detokenize{atc_hardy_fields::doc}]{\sphinxcrossref{\DUrole{doc}{fix\_modify AtC fields}}}}

\item {} 
\sphinxAtStartPar
{\hyperref[\detokenize{atc_hardy_gradients::doc}]{\sphinxcrossref{\DUrole{doc}{fix\_modify AtC gradients}}}}

\item {} 
\sphinxAtStartPar
{\hyperref[\detokenize{atc_hardy_rates::doc}]{\sphinxcrossref{\DUrole{doc}{fix\_modify AtC rates}}}}

\item {} 
\sphinxAtStartPar
{\hyperref[\detokenize{atc_hardy_computes::doc}]{\sphinxcrossref{\DUrole{doc}{fix\_modify AtC computes}}}}

\end{itemize}


\subsubsection{Default}
\label{\detokenize{atc_hardy_kernel:default}}
\sphinxAtStartPar
None.

\sphinxstepscope

\index{fix\_modify AtC on\_the\_fly@\spxentry{fix\_modify AtC on\_the\_fly}}\ignorespaces 

\subsection{fix\_modify AtC on\_the\_fly command}
\label{\detokenize{atc_hardy_on_the_fly:fix-modify-atc-on-the-fly-command}}\label{\detokenize{atc_hardy_on_the_fly:index-0}}\label{\detokenize{atc_hardy_on_the_fly::doc}}

\subsubsection{Syntax}
\label{\detokenize{atc_hardy_on_the_fly:syntax}}
\begin{sphinxVerbatim}[commandchars=\\\{\}]
\PYG{k}{fix\PYGZus{}modify }\PYG{o}{\PYGZlt{}}\PYG{n}{AtC}\PYG{+w}{ }\PYG{n}{fixID}\PYG{o}{\PYGZgt{}}\PYG{+w}{ }\PYG{n}{on\PYGZus{}the\PYGZus{}fly}\PYG{+w}{ }\PYG{o}{\PYGZlt{}}\PYG{n}{bond}\PYG{o}{|}\PYG{n}{kernel}\PYG{o}{\PYGZgt{}}\PYG{+w}{ }\PYG{o}{\PYGZlt{}}\PYG{n}{on}\PYG{o}{|}\PYG{n}{off}\PYG{o}{\PYGZgt{}}
\end{sphinxVerbatim}
\begin{itemize}
\item {} 
\sphinxAtStartPar
AtC fixID = ID of {\hyperref[\detokenize{fix_atc::doc}]{\sphinxcrossref{\DUrole{doc}{fix atc}}}} instance

\item {} 
\sphinxAtStartPar
on\_the\_fly = name of the AtC sub\sphinxhyphen{}command

\item {} 
\sphinxAtStartPar
\sphinxstyleemphasis{bond} or \sphinxstyleemphasis{kernel} = specifies on\sphinxhyphen{}the\sphinxhyphen{}fly calculation of \sphinxstyleemphasis{bond} or \sphinxstyleemphasis{kernel} matrix elements

\item {} 
\sphinxAtStartPar
\sphinxstyleemphasis{on} or \sphinxstyleemphasis{off} = activate or discontinue on\sphinxhyphen{}the\sphinxhyphen{}fly mode

\end{itemize}


\subsubsection{Examples}
\label{\detokenize{atc_hardy_on_the_fly:examples}}
\begin{sphinxVerbatim}[commandchars=\\\{\}]
\PYG{k}{fix\PYGZus{}modify }\PYG{n+nv+nvIdentifier}{AtC}\PYG{+w}{ }\PYG{n}{on\PYGZus{}the\PYGZus{}fly}\PYG{+w}{ }\PYG{n}{bond}\PYG{+w}{ }\PYG{n}{on}
\PYG{k}{fix\PYGZus{}modify }\PYG{n+nv+nvIdentifier}{AtC}\PYG{+w}{ }\PYG{n}{on\PYGZus{}the\PYGZus{}fly}\PYG{+w}{ }\PYG{n}{kernel}
\PYG{k}{fix\PYGZus{}modify }\PYG{n+nv+nvIdentifier}{AtC}\PYG{+w}{ }\PYG{n}{on\PYGZus{}the\PYGZus{}fly}\PYG{+w}{ }\PYG{n}{kernel}\PYG{+w}{ }\PYG{n}{off}
\end{sphinxVerbatim}


\subsubsection{Description}
\label{\detokenize{atc_hardy_on_the_fly:description}}
\sphinxAtStartPar
Overrides normal mode of pre\sphinxhyphen{}calculating and storing bond pair\sphinxhyphen{}to\sphinxhyphen{}node a
nd kernel atom\sphinxhyphen{}to\sphinxhyphen{}node matrices. If activated, it will calculate elements
of these matrices during repeated calls of field computations
(i.e. “on\sphinxhyphen{}the\sphinxhyphen{}fly”) and not store them for future use.  The \sphinxstyleemphasis{on} flag is
optional \sphinxhyphen{} if omitted, on\_the\_fly will be activated for the specified
matrix.  Can be deactivated using the \sphinxstyleemphasis{off} flag.


\subsubsection{Restrictions}
\label{\detokenize{atc_hardy_on_the_fly:restrictions}}
\sphinxAtStartPar
Must be used with {\hyperref[\detokenize{fix_atc::doc}]{\sphinxcrossref{\DUrole{doc}{fix atc hardy}}}}.


\subsubsection{Related AtC commands}
\label{\detokenize{atc_hardy_on_the_fly:related-atc-commands}}\begin{itemize}
\item {} 
\sphinxAtStartPar
{\hyperref[\detokenize{fix_atc:atc-fix-modify}]{\sphinxcrossref{\DUrole{std,std-ref}{fix\_modify AtC command overview}}}}

\end{itemize}


\subsubsection{Default}
\label{\detokenize{atc_hardy_on_the_fly:default}}
\sphinxAtStartPar
By default, on\sphinxhyphen{}the\sphinxhyphen{}fly calculation is not active (i.e. off). However,
THE code does a memory allocation check to determine if it can store all
needed bond and kernel matrix elements. If this allocation fails,
on\sphinxhyphen{}the\sphinxhyphen{}fly will be activated.

\sphinxstepscope

\index{fix\_modify AtC rates@\spxentry{fix\_modify AtC rates}}\ignorespaces 

\subsection{fix\_modify AtC rates command}
\label{\detokenize{atc_hardy_rates:fix-modify-atc-rates-command}}\label{\detokenize{atc_hardy_rates:index-0}}\label{\detokenize{atc_hardy_rates::doc}}

\subsubsection{Syntax}
\label{\detokenize{atc_hardy_rates:syntax}}
\begin{sphinxVerbatim}[commandchars=\\\{\}]
\PYG{k}{fix\PYGZus{}modify }\PYG{o}{\PYGZlt{}}\PYG{n}{AtC}\PYG{+w}{ }\PYG{n}{fixID}\PYG{o}{\PYGZgt{}}\PYG{+w}{ }\PYG{n}{rates}\PYG{+w}{ }\PYG{o}{\PYGZlt{}}\PYG{n}{add}\PYG{o}{|}\PYG{n}{delete}\PYG{o}{\PYGZgt{}}\PYG{+w}{ }\PYG{o}{\PYGZlt{}}\PYG{n}{list\PYGZus{}of\PYGZus{}fields}\PYG{o}{\PYGZgt{}}
\end{sphinxVerbatim}
\begin{itemize}
\item {} 
\sphinxAtStartPar
AtC fixID = ID of {\hyperref[\detokenize{fix_atc::doc}]{\sphinxcrossref{\DUrole{doc}{fix atc}}}} instance

\item {} 
\sphinxAtStartPar
rates = name of the AtC sub\sphinxhyphen{}command

\item {} 
\sphinxAtStartPar
\sphinxstyleemphasis{add} or \sphinxstyleemphasis{delete} = select whether to add or delete calculation of rates for the listed output fields

\item {} 
\sphinxAtStartPar
list\_of\_fields = one or more of the fields listed below:
\begin{itemize}
\item {} 
\sphinxAtStartPar
density : mass per unit volume

\item {} 
\sphinxAtStartPar
displacement : displacement vector

\item {} 
\sphinxAtStartPar
momentum : momentum per unit volume

\item {} 
\sphinxAtStartPar
velocity : defined by momentum divided by density

\item {} 
\sphinxAtStartPar
projected\_velocity : simple kernel estimation of atomic velocities

\item {} 
\sphinxAtStartPar
temperature : temperature derived from the relative atomic kinetic energy

\item {} 
\sphinxAtStartPar
kinetic\_temperature : temperature derived from the full kinetic energy

\item {} 
\sphinxAtStartPar
number\_density : simple kernel estimation of number of atoms per unit volume

\item {} 
\sphinxAtStartPar
stress : Cauchy stress tensor for eulerian analysis (atom\_element\_map), or first Piola\sphinxhyphen{}Kirchhoff stress tensor for lagrangian analysis

\item {} 
\sphinxAtStartPar
transformed\_stress : first Piola\sphinxhyphen{}Kirchhoff stress tensor for eulerian analysis (atom\_element\_map), or Cauchy stress tensor for lagrangian analysis

\item {} 
\sphinxAtStartPar
heat\_flux : spatial heat flux vector for eulerian, or referential heat flux vector for lagrangian

\item {} 
\sphinxAtStartPar
potential\_energy : potential energy per unit volume

\item {} 
\sphinxAtStartPar
kinetic\_energy : kinetic energy per unit volume

\item {} 
\sphinxAtStartPar
thermal\_energy : thermal energy (kinetic energy \sphinxhyphen{} continuum kinetic energy) per unit volume

\item {} 
\sphinxAtStartPar
internal\_energy : total internal energy (potential + thermal) per unit volume

\item {} 
\sphinxAtStartPar
energy : total energy (potential + kinetic) per unit volume

\item {} 
\sphinxAtStartPar
number\_density : number of atoms per unit volume

\item {} 
\sphinxAtStartPar
eshelby\_stress : configurational stress (energy\sphinxhyphen{}momentum) tensor defined by {\hyperref[\detokenize{atc_hardy_rates:eshelby2}]{\sphinxcrossref{\DUrole{std,std-ref}{(Eshelby)}}}}

\item {} 
\sphinxAtStartPar
vacancy\_concentration : volume fraction of vacancy content

\item {} 
\sphinxAtStartPar
type\_concentration : volume fraction of a specific atom type

\end{itemize}

\end{itemize}


\subsubsection{Examples}
\label{\detokenize{atc_hardy_rates:examples}}
\begin{sphinxVerbatim}[commandchars=\\\{\}]
\PYG{k}{fix\PYGZus{}modify }\PYG{n+nv+nvIdentifier}{AtC}\PYG{+w}{ }\PYG{n}{rates}\PYG{+w}{ }\PYG{n}{add}\PYG{+w}{ }\PYG{n}{temperature}\PYG{+w}{ }\PYG{k}{velocity }\PYG{n+nv+nvIdentifier}{stress}
\PYG{k}{fix\PYGZus{}modify }\PYG{n+nv+nvIdentifier}{AtC}\PYG{+w}{ }\PYG{n}{rates}\PYG{+w}{ }\PYG{n}{delete}\PYG{+w}{ }\PYG{n}{stress}
\end{sphinxVerbatim}


\subsubsection{Description}
\label{\detokenize{atc_hardy_rates:description}}
\sphinxAtStartPar
Requests calculation and output of rates (time derivatives) of the
fields from the AtC transfer class.  For Eulerian analysis (see
{\hyperref[\detokenize{atc_atom_element_map::doc}]{\sphinxcrossref{\DUrole{doc}{fix\_modify AtC atom\_element\_map}}}}) these
rates are the partial time derivatives of the nodal fields, not the full
(material) time derivatives.


\subsubsection{Restrictions}
\label{\detokenize{atc_hardy_rates:restrictions}}
\sphinxAtStartPar
Must be used with {\hyperref[\detokenize{fix_atc::doc}]{\sphinxcrossref{\DUrole{doc}{fix atc hardy}}}}.


\subsubsection{Related AtC commands}
\label{\detokenize{atc_hardy_rates:related-atc-commands}}\begin{itemize}
\item {} 
\sphinxAtStartPar
{\hyperref[\detokenize{fix_atc:atc-fix-modify}]{\sphinxcrossref{\DUrole{std,std-ref}{fix\_modify AtC command overview}}}}

\item {} 
\sphinxAtStartPar
{\hyperref[\detokenize{atc_atom_element_map::doc}]{\sphinxcrossref{\DUrole{doc}{fix\_modify AtC atom\_element\_map}}}}

\item {} 
\sphinxAtStartPar
{\hyperref[\detokenize{atc_hardy_fields::doc}]{\sphinxcrossref{\DUrole{doc}{fix\_modify AtC fields}}}}

\item {} 
\sphinxAtStartPar
{\hyperref[\detokenize{atc_hardy_gradients::doc}]{\sphinxcrossref{\DUrole{doc}{fix\_modify AtC fields}}}}

\end{itemize}


\subsubsection{Default}
\label{\detokenize{atc_hardy_rates:default}}
\sphinxAtStartPar
None.


\subsubsection{References}
\label{\detokenize{atc_hardy_rates:references}}\phantomsection\label{\detokenize{atc_hardy_rates:eshelby2}}
\sphinxAtStartPar
\sphinxstylestrong{(Eshelby)} J.D. Eshelby, Philos. Trans. Royal Soc. London A, Math. Phys. Sci., Vol. 244, No. 877 (1951) pp. 87\sphinxhyphen{}112; J. Elasticity, Vol. 5, Nos. 3\sphinxhyphen{}4 (1975) pp. 321\sphinxhyphen{}335{]}

\sphinxstepscope

\index{fix\_modify AtC initial@\spxentry{fix\_modify AtC initial}}\ignorespaces 

\subsection{fix\_modify AtC initial command}
\label{\detokenize{atc_initial:fix-modify-atc-initial-command}}\label{\detokenize{atc_initial:index-0}}\label{\detokenize{atc_initial::doc}}

\subsubsection{Syntax}
\label{\detokenize{atc_initial:syntax}}
\begin{sphinxVerbatim}[commandchars=\\\{\}]
\PYG{k}{fix\PYGZus{}modify }\PYG{o}{\PYGZlt{}}\PYG{n}{AtC}\PYG{+w}{ }\PYG{n}{fixID}\PYG{o}{\PYGZgt{}}\PYG{+w}{ }\PYG{n}{initial}\PYG{+w}{ }\PYG{o}{\PYGZlt{}}\PYG{n}{field}\PYG{o}{\PYGZgt{}}\PYG{+w}{ }\PYG{o}{\PYGZlt{}}\PYG{n}{nodeset}\PYG{o}{\PYGZgt{}}\PYG{+w}{ }\PYG{o}{\PYGZlt{}}\PYG{n}{constant}\PYG{o}{|}\PYG{n}{function}\PYG{o}{\PYGZgt{}}
\end{sphinxVerbatim}
\begin{itemize}
\item {} 
\sphinxAtStartPar
AtC fixID = ID of {\hyperref[\detokenize{fix_atc::doc}]{\sphinxcrossref{\DUrole{doc}{fix atc}}}} instance

\item {} 
\sphinxAtStartPar
initial = name of the AtC sub\sphinxhyphen{}command

\item {} 
\sphinxAtStartPar
field = field kind name valid for type of physics: \sphinxstyleemphasis{temperature} or \sphinxstyleemphasis{electron\_temperature}

\item {} 
\sphinxAtStartPar
nodeset = name of set of nodes to apply initial condition

\item {} 
\sphinxAtStartPar
\sphinxstyleemphasis{constant} or \sphinxstyleemphasis{function} = value or name of function followed by its parameters

\end{itemize}


\subsubsection{Examples}
\label{\detokenize{atc_initial:examples}}
\begin{sphinxVerbatim}[commandchars=\\\{\}]
\PYG{k}{fix\PYGZus{}modify }\PYG{n+nv+nvIdentifier}{AtC}\PYG{+w}{ }\PYG{n}{initial}\PYG{+w}{ }\PYG{n}{temperature}\PYG{+w}{ }\PYG{n}{groupNAME}\PYG{+w}{ }\PYG{l+m}{10}\PYG{n}{.}
\end{sphinxVerbatim}


\subsubsection{Description}
\label{\detokenize{atc_initial:description}}
\sphinxAtStartPar
Sets the initial values for the specified field at the specified nodes.


\subsubsection{Restrictions}
\label{\detokenize{atc_initial:restrictions}}
\sphinxAtStartPar
The keyword \sphinxstyleemphasis{all} is reserved and thus not available as nodeset name.


\subsubsection{Related AtC commands}
\label{\detokenize{atc_initial:related-atc-commands}}\begin{itemize}
\item {} 
\sphinxAtStartPar
{\hyperref[\detokenize{fix_atc:atc-fix-modify}]{\sphinxcrossref{\DUrole{std,std-ref}{fix\_modify AtC command overview}}}}

\end{itemize}


\subsubsection{Default}
\label{\detokenize{atc_initial:default}}
\sphinxAtStartPar
None.

\sphinxstepscope

\index{fix\_modify AtC internal\_element\_set@\spxentry{fix\_modify AtC internal\_element\_set}}\ignorespaces 

\subsection{fix\_modify AtC internal\_element\_set command}
\label{\detokenize{atc_internal_element_set:fix-modify-atc-internal-element-set-command}}\label{\detokenize{atc_internal_element_set:index-0}}\label{\detokenize{atc_internal_element_set::doc}}

\subsubsection{Syntax}
\label{\detokenize{atc_internal_element_set:syntax}}
\begin{sphinxVerbatim}[commandchars=\\\{\}]
\PYG{k}{fix\PYGZus{}modify }\PYG{o}{\PYGZlt{}}\PYG{n}{AtC}\PYG{+w}{ }\PYG{n}{fixID}\PYG{o}{\PYGZgt{}}\PYG{+w}{ }\PYG{n}{internal\PYGZus{}element\PYGZus{}set}\PYG{+w}{ }\PYG{o}{\PYGZlt{}}\PYG{n}{element\PYGZus{}set\PYGZus{}name}\PYG{o}{\PYGZgt{}}
\end{sphinxVerbatim}
\begin{itemize}
\item {} 
\sphinxAtStartPar
AtC fixID = ID of {\hyperref[\detokenize{fix_atc::doc}]{\sphinxcrossref{\DUrole{doc}{fix atc}}}} instance

\item {} 
\sphinxAtStartPar
internal\_element\_set = name of the AtC sub\sphinxhyphen{}command

\item {} 
\sphinxAtStartPar
element\_set\_name = name of element set defining internal region, or \sphinxstyleemphasis{off}

\end{itemize}


\subsubsection{Examples}
\label{\detokenize{atc_internal_element_set:examples}}
\begin{sphinxVerbatim}[commandchars=\\\{\}]
\PYG{k}{fix\PYGZus{}modify }\PYG{n+nv+nvIdentifier}{AtC}\PYG{+w}{ }\PYG{n}{internal\PYGZus{}element\PYGZus{}set}\PYG{+w}{ }\PYG{n}{myElementSet}
\PYG{k}{fix\PYGZus{}modify }\PYG{n+nv+nvIdentifier}{AtC}\PYG{+w}{ }\PYG{n}{internal\PYGZus{}element\PYGZus{}set}\PYG{+w}{ }\PYG{n}{off}
\end{sphinxVerbatim}


\subsubsection{Description}
\label{\detokenize{atc_internal_element_set:description}}
\sphinxAtStartPar
Enables AtC to base the region for internal atoms to be an element
set. If no ghost atoms are used, all the AtC atoms must be constrained
to remain in this element set by the user, e.g., with walls. If boundary
atoms are used in conjunction with Eulerian atom maps AtC will partition
all atoms of a boundary or internal type to be of type internal if they
are in the internal region or to be of type boundary otherwise.


\subsubsection{Restrictions}
\label{\detokenize{atc_internal_element_set:restrictions}}
\sphinxAtStartPar
If boundary atoms are used in conjunction with Eulerian atom maps, the
Eulerian reset frequency must be an integer multiple of the Lammps
reneighbor frequency.


\subsubsection{Related AtC commands}
\label{\detokenize{atc_internal_element_set:related-atc-commands}}\begin{itemize}
\item {} 
\sphinxAtStartPar
{\hyperref[\detokenize{fix_atc:atc-fix-modify}]{\sphinxcrossref{\DUrole{std,std-ref}{fix\_modify AtC command overview}}}}

\item {} 
\sphinxAtStartPar
{\hyperref[\detokenize{atc_atom_element_map::doc}]{\sphinxcrossref{\DUrole{doc}{fix\_modify AtC atom\_element\_map}}}}

\item {} 
\sphinxAtStartPar
{\hyperref[\detokenize{atc_boundary_type::doc}]{\sphinxcrossref{\DUrole{doc}{fix\_modify AtC boundary type}}}}

\end{itemize}


\subsubsection{Default}
\label{\detokenize{atc_internal_element_set:default}}
\sphinxAtStartPar
\sphinxstyleemphasis{off}

\sphinxstepscope

\index{fix\_modify AtC internal\_quadrature@\spxentry{fix\_modify AtC internal\_quadrature}}\ignorespaces 

\subsection{fix\_modify AtC internal\_quadrature command}
\label{\detokenize{atc_internal_quadrature:fix-modify-atc-internal-quadrature-command}}\label{\detokenize{atc_internal_quadrature:index-0}}\label{\detokenize{atc_internal_quadrature::doc}}

\subsubsection{Syntax}
\label{\detokenize{atc_internal_quadrature:syntax}}
\begin{sphinxVerbatim}[commandchars=\\\{\}]
\PYG{k}{fix\PYGZus{}modify }\PYG{o}{\PYGZlt{}}\PYG{n}{AtC}\PYG{+w}{ }\PYG{n}{fixID}\PYG{o}{\PYGZgt{}}\PYG{+w}{ }\PYG{n}{internal\PYGZus{}quadrature}\PYG{+w}{ }\PYG{o}{\PYGZlt{}}\PYG{n}{on}\PYG{o}{|}\PYG{n}{off}\PYG{o}{\PYGZgt{}}\PYG{+w}{ }\PYG{n}{[region]}
\end{sphinxVerbatim}
\begin{itemize}
\item {} 
\sphinxAtStartPar
AtC fixID = ID of {\hyperref[\detokenize{fix_atc::doc}]{\sphinxcrossref{\DUrole{doc}{fix atc}}}} instance

\item {} 
\sphinxAtStartPar
internal\_quadrature = name of the AtC sub\sphinxhyphen{}command

\item {} 
\sphinxAtStartPar
on or off = select whether internal quadrature is enabled or not

\item {} 
\sphinxAtStartPar
region = treat finite elements as within MD region (optional)

\end{itemize}


\subsubsection{Examples}
\label{\detokenize{atc_internal_quadrature:examples}}
\begin{sphinxVerbatim}[commandchars=\\\{\}]
\PYG{k}{fix\PYGZus{}modify }\PYG{n+nv+nvIdentifier}{AtC}\PYG{+w}{ }\PYG{n}{internal\PYGZus{}quadrature}\PYG{+w}{ }\PYG{n}{off}
\end{sphinxVerbatim}


\subsubsection{Description}
\label{\detokenize{atc_internal_quadrature:description}}
\sphinxAtStartPar
Command to use or not use atomic quadrature on internal elements fully
filled with atoms. By turning the internal quadrature off these elements
do not contribute to the governing PDE and the fields at the internal
nodes follow the weighted averages of the atomic data.

\sphinxAtStartPar
Optional region tag specifies which finite element nodes will be treated
as being within the MD region. This option is only valid with
internal\_quadrature off.


\subsubsection{Related AtC commands}
\label{\detokenize{atc_internal_quadrature:related-atc-commands}}\begin{itemize}
\item {} 
\sphinxAtStartPar
{\hyperref[\detokenize{fix_atc:atc-fix-modify}]{\sphinxcrossref{\DUrole{std,std-ref}{fix\_modify AtC command overview}}}}

\end{itemize}


\subsubsection{Default}
\label{\detokenize{atc_internal_quadrature:default}}
\sphinxAtStartPar
on.

\sphinxstepscope

\index{fix\_modify AtC kernel\_bandwidth@\spxentry{fix\_modify AtC kernel\_bandwidth}}\ignorespaces 

\subsection{fix\_modify AtC kernel\_bandwidth command}
\label{\detokenize{atc_kernel_bandwidth:fix-modify-atc-kernel-bandwidth-command}}\label{\detokenize{atc_kernel_bandwidth:index-0}}\label{\detokenize{atc_kernel_bandwidth::doc}}

\subsubsection{Syntax}
\label{\detokenize{atc_kernel_bandwidth:syntax}}
\begin{sphinxVerbatim}[commandchars=\\\{\}]
\PYG{k}{fix\PYGZus{}modify }\PYG{o}{\PYGZlt{}}\PYG{n}{AtC}\PYG{+w}{ }\PYG{n}{fixID}\PYG{o}{\PYGZgt{}}\PYG{+w}{ }\PYG{n}{kernel\PYGZus{}bandwidth}\PYG{+w}{ }\PYG{o}{\PYGZlt{}}\PYG{n}{value}\PYG{o}{\PYGZgt{}}
\end{sphinxVerbatim}
\begin{itemize}
\item {} 
\sphinxAtStartPar
AtC fixID = ID of {\hyperref[\detokenize{fix_atc::doc}]{\sphinxcrossref{\DUrole{doc}{fix atc}}}} instance

\item {} 
\sphinxAtStartPar
kernel\_bandwidth = name of the AtC sub\sphinxhyphen{}command

\item {} 
\sphinxAtStartPar
value = new bandwidth value

\end{itemize}


\subsubsection{Examples}
\label{\detokenize{atc_kernel_bandwidth:examples}}
\begin{sphinxVerbatim}[commandchars=\\\{\}]
\PYG{k}{fix\PYGZus{}modify }\PYG{n+nv+nvIdentifier}{AtC}\PYG{+w}{ }\PYG{n}{kernel\PYGZus{}bandwidth}\PYG{+w}{ }\PYG{l+m}{8}
\end{sphinxVerbatim}


\subsubsection{Description}
\label{\detokenize{atc_kernel_bandwidth:description}}
\sphinxAtStartPar
Sets a maximum parallel bandwidth for the kernel functions during
parallel communication. If the command is not issued, the default will
be to assume the bandwidth of the kernel matrix corresponds to the
number of sampling locations.


\subsubsection{Restrictions}
\label{\detokenize{atc_kernel_bandwidth:restrictions}}
\sphinxAtStartPar
Only is used if kernel functions are being used.


\subsubsection{Related AtC commands}
\label{\detokenize{atc_kernel_bandwidth:related-atc-commands}}\begin{itemize}
\item {} 
\sphinxAtStartPar
{\hyperref[\detokenize{fix_atc:atc-fix-modify}]{\sphinxcrossref{\DUrole{std,std-ref}{fix\_modify AtC command overview}}}}

\end{itemize}


\subsubsection{Default}
\label{\detokenize{atc_kernel_bandwidth:default}}
\sphinxAtStartPar
Number of sample locations.

\sphinxstepscope

\index{fix\_modify AtC control lumped\_lambda\_solve@\spxentry{fix\_modify AtC control lumped\_lambda\_solve}}\ignorespaces 

\subsection{fix\_modify AtC control lumped\_lambda\_solve command}
\label{\detokenize{atc_lumped_lambda_solve:fix-modify-atc-control-lumped-lambda-solve-command}}\label{\detokenize{atc_lumped_lambda_solve:index-0}}\label{\detokenize{atc_lumped_lambda_solve::doc}}

\subsubsection{Syntax}
\label{\detokenize{atc_lumped_lambda_solve:syntax}}
\begin{sphinxVerbatim}[commandchars=\\\{\}]
\PYG{k}{fix\PYGZus{}modify }\PYG{o}{\PYGZlt{}}\PYG{n}{AtC}\PYG{+w}{ }\PYG{n}{fixID}\PYG{o}{\PYGZgt{}}\PYG{+w}{ }\PYG{n}{control}\PYG{+w}{ }\PYG{n}{lumped\PYGZus{}lambda\PYGZus{}solve}\PYG{+w}{ }\PYG{o}{\PYGZlt{}}\PYG{n}{on}\PYG{o}{|}\PYG{n}{off}\PYG{o}{\PYGZgt{}}
\end{sphinxVerbatim}
\begin{itemize}
\item {} 
\sphinxAtStartPar
AtC fixID = ID of {\hyperref[\detokenize{fix_atc::doc}]{\sphinxcrossref{\DUrole{doc}{fix atc}}}} instance

\item {} 
\sphinxAtStartPar
control lumped\_lambda\_solve = name of the AtC sub\sphinxhyphen{}command

\item {} 
\sphinxAtStartPar
\sphinxstyleemphasis{on} or \sphinxstyleemphasis{off} = Toggles state of lumped matrix

\end{itemize}


\subsubsection{Examples}
\label{\detokenize{atc_lumped_lambda_solve:examples}}
\begin{sphinxVerbatim}[commandchars=\\\{\}]
\PYG{k}{fix\PYGZus{}modify }\PYG{n+nv+nvIdentifier}{AtC}\PYG{+w}{ }\PYG{n}{control}\PYG{+w}{ }\PYG{n}{lumped\PYGZus{}lambda\PYGZus{}solve}\PYG{+w}{ }\PYG{n}{on}
\end{sphinxVerbatim}


\subsubsection{Description}
\label{\detokenize{atc_lumped_lambda_solve:description}}
\sphinxAtStartPar
Command select whether to use or not use lumped matrix for lambda solve.


\subsubsection{Restrictions}
\label{\detokenize{atc_lumped_lambda_solve:restrictions}}
\sphinxAtStartPar
None.


\subsubsection{Related AtC commands}
\label{\detokenize{atc_lumped_lambda_solve:related-atc-commands}}\begin{itemize}
\item {} 
\sphinxAtStartPar
{\hyperref[\detokenize{fix_atc:atc-fix-modify}]{\sphinxcrossref{\DUrole{std,std-ref}{fix\_modify AtC command overview}}}}

\end{itemize}


\subsubsection{Default}
\label{\detokenize{atc_lumped_lambda_solve:default}}
\sphinxAtStartPar
off

\sphinxstepscope

\index{fix\_modify AtC control mask\_direction@\spxentry{fix\_modify AtC control mask\_direction}}\ignorespaces 

\subsection{fix\_modify AtC control mask\_direction command}
\label{\detokenize{atc_mask_direction:fix-modify-atc-control-mask-direction-command}}\label{\detokenize{atc_mask_direction:index-0}}\label{\detokenize{atc_mask_direction::doc}}

\subsubsection{Syntax}
\label{\detokenize{atc_mask_direction:syntax}}
\begin{sphinxVerbatim}[commandchars=\\\{\}]
\PYG{k}{fix\PYGZus{}modify }\PYG{o}{\PYGZlt{}}\PYG{n}{AtC}\PYG{+w}{ }\PYG{n}{fixID}\PYG{o}{\PYGZgt{}}\PYG{+w}{ }\PYG{n}{control}\PYG{+w}{ }\PYG{n}{mask\PYGZus{}direction}\PYG{+w}{ }\PYG{o}{\PYGZlt{}}\PYG{n}{direction}\PYG{o}{\PYGZgt{}}\PYG{+w}{ }\PYG{o}{\PYGZlt{}}\PYG{n}{on}\PYG{o}{|}\PYG{n}{off}\PYG{o}{\PYGZgt{}}
\end{sphinxVerbatim}
\begin{itemize}
\item {} 
\sphinxAtStartPar
AtC fixID = ID of {\hyperref[\detokenize{fix_atc::doc}]{\sphinxcrossref{\DUrole{doc}{fix atc}}}} instance

\item {} 
\sphinxAtStartPar
control mask\_direction = name of the AtC sub\sphinxhyphen{}command

\item {} 
\sphinxAtStartPar
direction = select direction

\item {} 
\sphinxAtStartPar
\sphinxstyleemphasis{on} or \sphinxstyleemphasis{off} = Toggles state

\end{itemize}


\subsubsection{Examples}
\label{\detokenize{atc_mask_direction:examples}}
\begin{sphinxVerbatim}[commandchars=\\\{\}]
\PYG{k}{fix\PYGZus{}modify }\PYG{n+nv+nvIdentifier}{AtC}\PYG{+w}{ }\PYG{n}{control}\PYG{+w}{ }\PYG{n}{mask\PYGZus{}direction}\PYG{+w}{ }\PYG{l+m}{0}\PYG{+w}{ }\PYG{n}{on}
\end{sphinxVerbatim}


\subsubsection{Description}
\label{\detokenize{atc_mask_direction:description}}
\sphinxAtStartPar
Command to mask out certain dimensions from the atomic regulator


\subsubsection{Restrictions}
\label{\detokenize{atc_mask_direction:restrictions}}
\sphinxAtStartPar
None.


\subsubsection{Related AtC commands}
\label{\detokenize{atc_mask_direction:related-atc-commands}}\begin{itemize}
\item {} 
\sphinxAtStartPar
{\hyperref[\detokenize{fix_atc:atc-fix-modify}]{\sphinxcrossref{\DUrole{std,std-ref}{fix\_modify AtC command overview}}}}

\end{itemize}

\sphinxstepscope

\index{fix\_modify AtC mass\_matrix@\spxentry{fix\_modify AtC mass\_matrix}}\ignorespaces 

\subsection{fix\_modify AtC mass\_matrix command}
\label{\detokenize{atc_mass_matrix:fix-modify-atc-mass-matrix-command}}\label{\detokenize{atc_mass_matrix:index-0}}\label{\detokenize{atc_mass_matrix::doc}}

\subsubsection{Syntax}
\label{\detokenize{atc_mass_matrix:syntax}}
\begin{sphinxVerbatim}[commandchars=\\\{\}]
\PYG{k}{fix\PYGZus{}modify }\PYG{o}{\PYGZlt{}}\PYG{n}{AtC}\PYG{+w}{ }\PYG{n}{fixID}\PYG{o}{\PYGZgt{}}\PYG{+w}{ }\PYG{n}{mass\PYGZus{}matrix}\PYG{+w}{ }\PYG{o}{\PYGZlt{}}\PYG{n}{fe}\PYG{o}{|}\PYG{n}{md\PYGZus{}fe}\PYG{o}{\PYGZgt{}}
\end{sphinxVerbatim}
\begin{itemize}
\item {} 
\sphinxAtStartPar
AtC fixID = ID of {\hyperref[\detokenize{fix_atc::doc}]{\sphinxcrossref{\DUrole{doc}{fix atc}}}} instance

\item {} 
\sphinxAtStartPar
mass\_matrix = name of the AtC sub\sphinxhyphen{}command

\item {} 
\sphinxAtStartPar
\sphinxstyleemphasis{fe} or \sphinxstyleemphasis{md\_fe} = activate/deactivate using the FE mass matrix in the MD region

\end{itemize}


\subsubsection{Examples}
\label{\detokenize{atc_mass_matrix:examples}}
\begin{sphinxVerbatim}[commandchars=\\\{\}]
\PYG{k}{fix\PYGZus{}modify }\PYG{n+nv+nvIdentifier}{AtC}\PYG{+w}{ }\PYG{n}{mass\PYGZus{}matrix}\PYG{+w}{ }\PYG{n}{fe}
\end{sphinxVerbatim}


\subsubsection{Description}
\label{\detokenize{atc_mass_matrix:description}}
\sphinxAtStartPar
Determines whether AtC uses the FE mass matrix based on Gaussian
quadrature or based on atomic quadrature in the MD region. This is
useful for fully overlapping simulations to improve efficiency.


\subsubsection{Restrictions}
\label{\detokenize{atc_mass_matrix:restrictions}}
\sphinxAtStartPar
Should not be used unless the FE region is contained within the MD
region, otherwise the method will be unstable and inaccurate.


\subsubsection{Related AtC commands}
\label{\detokenize{atc_mass_matrix:related-atc-commands}}\begin{itemize}
\item {} 
\sphinxAtStartPar
{\hyperref[\detokenize{fix_atc:atc-fix-modify}]{\sphinxcrossref{\DUrole{std,std-ref}{fix\_modify AtC command overview}}}}

\end{itemize}


\subsubsection{Default}
\label{\detokenize{atc_mass_matrix:default}}
\sphinxAtStartPar
\sphinxstyleemphasis{md\_fe}

\sphinxstepscope

\index{fix\_modify AtC material@\spxentry{fix\_modify AtC material}}\ignorespaces 

\subsection{fix\_modify AtC material command}
\label{\detokenize{atc_material:fix-modify-atc-material-command}}\label{\detokenize{atc_material:index-0}}\label{\detokenize{atc_material::doc}}

\subsubsection{Syntax}
\label{\detokenize{atc_material:syntax}}
\begin{sphinxVerbatim}[commandchars=\\\{\}]
\PYG{k}{fix\PYGZus{}modify }\PYG{o}{\PYGZlt{}}\PYG{n}{AtC}\PYG{+w}{ }\PYG{n}{fixID}\PYG{o}{\PYGZgt{}}\PYG{+w}{ }\PYG{n}{material}\PYG{+w}{ }\PYG{o}{\PYGZlt{}}\PYG{n}{elementset\PYGZus{}name}\PYG{o}{\PYGZgt{}}\PYG{+w}{ }\PYG{o}{\PYGZlt{}}\PYG{n}{material\PYGZus{}id}\PYG{o}{\PYGZgt{}}
\end{sphinxVerbatim}
\begin{itemize}
\item {} 
\sphinxAtStartPar
AtC fixID = ID of {\hyperref[\detokenize{fix_atc::doc}]{\sphinxcrossref{\DUrole{doc}{fix atc}}}} instance

\item {} 
\sphinxAtStartPar
material = name of the AtC sub\sphinxhyphen{}command

\item {} 
\sphinxAtStartPar
elementset\_name = name of the elementset

\item {} 
\sphinxAtStartPar
material\_id = ID of the material

\end{itemize}


\subsubsection{Examples}
\label{\detokenize{atc_material:examples}}
\begin{sphinxVerbatim}[commandchars=\\\{\}]
\PYG{k}{fix\PYGZus{}modify }\PYG{n+nv+nvIdentifier}{AtC}\PYG{+w}{ }\PYG{n}{material}\PYG{+w}{ }\PYG{n}{gap\PYGZus{}region}\PYG{+w}{ }\PYG{l+m}{1}
\end{sphinxVerbatim}


\subsubsection{Description}
\label{\detokenize{atc_material:description}}
\sphinxAtStartPar
Sets the material model in \sphinxstyleemphasis{elementset\_name} to be of type \sphinxstyleemphasis{material\_id}.


\subsubsection{Restrictions}
\label{\detokenize{atc_material:restrictions}}
\sphinxAtStartPar
The element set must already be created and the material must be
specified in the material file given the the atc fix on construction


\subsubsection{Related AtC commands}
\label{\detokenize{atc_material:related-atc-commands}}\begin{itemize}
\item {} 
\sphinxAtStartPar
{\hyperref[\detokenize{fix_atc:atc-fix-modify}]{\sphinxcrossref{\DUrole{std,std-ref}{fix\_modify AtC command overview}}}}

\end{itemize}


\subsubsection{Default}
\label{\detokenize{atc_material:default}}
\sphinxAtStartPar
All elements default to the first material in the material file.

\sphinxstepscope

\index{fix\_modify AtC mesh add\_to\_nodeset@\spxentry{fix\_modify AtC mesh add\_to\_nodeset}}\ignorespaces 

\subsection{fix\_modify AtC mesh add\_to\_nodeset command}
\label{\detokenize{atc_mesh_add_to_nodeset:fix-modify-atc-mesh-add-to-nodeset-command}}\label{\detokenize{atc_mesh_add_to_nodeset:index-0}}\label{\detokenize{atc_mesh_add_to_nodeset::doc}}

\subsubsection{Syntax}
\label{\detokenize{atc_mesh_add_to_nodeset:syntax}}
\begin{sphinxVerbatim}[commandchars=\\\{\}]
\PYG{k}{fix\PYGZus{}modify }\PYG{o}{\PYGZlt{}}\PYG{n}{AtC}\PYG{+w}{ }\PYG{n}{fixID}\PYG{o}{\PYGZgt{}}\PYG{+w}{ }\PYG{n}{mesh}\PYG{+w}{ }\PYG{n}{add\PYGZus{}to\PYGZus{}nodeset}\PYG{+w}{ }\PYG{o}{\PYGZlt{}}\PYG{n}{id}\PYG{o}{\PYGZgt{}}\PYG{+w}{ }\PYG{o}{\PYGZlt{}}\PYG{n}{xmin}\PYG{o}{\PYGZgt{}}\PYG{+w}{ }\PYG{o}{\PYGZlt{}}\PYG{n}{xmax}\PYG{o}{\PYGZgt{}}\PYG{+w}{ }\PYG{o}{\PYGZlt{}}\PYG{n}{ymin}\PYG{o}{\PYGZgt{}}\PYG{+w}{ }\PYG{o}{\PYGZlt{}}\PYG{n}{ymax}\PYG{o}{\PYGZgt{}}\PYG{+w}{ }\PYG{o}{\PYGZlt{}}\PYG{n}{zmin}\PYG{o}{\PYGZgt{}}\PYG{+w}{ }\PYG{o}{\PYGZlt{}}\PYG{n}{zmax}\PYG{o}{\PYGZgt{}}
\end{sphinxVerbatim}
\begin{itemize}
\item {} 
\sphinxAtStartPar
AtC fixID = ID of {\hyperref[\detokenize{fix_atc::doc}]{\sphinxcrossref{\DUrole{doc}{fix atc}}}} instance

\item {} 
\sphinxAtStartPar
mesh create\_nodeset = name of the AtC sub\sphinxhyphen{}command

\item {} 
\sphinxAtStartPar
id = id to assign to the collection of FE nodes

\item {} 
\sphinxAtStartPar
\textless{}xmin\textgreater{} \textless{}xmax\textgreater{} \textless{}ymin\textgreater{} \textless{}ymax\textgreater{} \textless{}zmin\textgreater{} \textless{}zmax\textgreater{} = coordinates of the bounding box that contains the desired nodes to be added

\end{itemize}


\subsubsection{Examples}
\label{\detokenize{atc_mesh_add_to_nodeset:examples}}
\begin{sphinxVerbatim}[commandchars=\\\{\}]
\PYG{k}{fix\PYGZus{}modify }\PYG{n+nv+nvIdentifier}{AtC}\PYG{+w}{ }\PYG{n}{mesh}\PYG{+w}{ }\PYG{n}{add\PYGZus{}to\PYGZus{}nodeset}\PYG{+w}{ }\PYG{n}{lbc}\PYG{+w}{ }\PYG{o}{\PYGZhy{}}\PYG{l+m}{12.1}\PYG{+w}{ }\PYG{o}{\PYGZhy{}}\PYG{l+m}{11.9}\PYG{+w}{ }\PYG{o}{\PYGZhy{}}\PYG{l+m}{12}\PYG{+w}{ }\PYG{l+m}{12}\PYG{+w}{ }\PYG{o}{\PYGZhy{}}\PYG{l+m}{12}\PYG{+w}{ }\PYG{l+m}{12}
\end{sphinxVerbatim}


\subsubsection{Description}
\label{\detokenize{atc_mesh_add_to_nodeset:description}}
\sphinxAtStartPar
Command to add nodes to an already existing FE nodeset.


\subsubsection{Restrictions}
\label{\detokenize{atc_mesh_add_to_nodeset:restrictions}}
\sphinxAtStartPar
None


\subsubsection{Related AtC commands}
\label{\detokenize{atc_mesh_add_to_nodeset:related-atc-commands}}\begin{itemize}
\item {} 
\sphinxAtStartPar
{\hyperref[\detokenize{fix_atc:atc-fix-modify}]{\sphinxcrossref{\DUrole{std,std-ref}{fix\_modify AtC command overview}}}}

\item {} 
\sphinxAtStartPar
{\hyperref[\detokenize{atc_mesh_create_nodeset::doc}]{\sphinxcrossref{\DUrole{doc}{fix\_modify AtC mesh create\_nodeset}}}}

\end{itemize}


\subsubsection{Default}
\label{\detokenize{atc_mesh_add_to_nodeset:default}}
\sphinxAtStartPar
Coordinates are assumed to be in lattice units.

\sphinxstepscope

\index{fix\_modify AtC mesh create@\spxentry{fix\_modify AtC mesh create}}\ignorespaces 

\subsection{fix\_modify AtC mesh create command}
\label{\detokenize{atc_mesh_create:fix-modify-atc-mesh-create-command}}\label{\detokenize{atc_mesh_create:index-0}}\label{\detokenize{atc_mesh_create::doc}}

\subsubsection{Syntax}
\label{\detokenize{atc_mesh_create:syntax}}
\begin{sphinxVerbatim}[commandchars=\\\{\}]
\PYG{k}{fix\PYGZus{}modify }\PYG{o}{\PYGZlt{}}\PYG{n}{AtC}\PYG{+w}{ }\PYG{n}{fixID}\PYG{o}{\PYGZgt{}}\PYG{+w}{ }\PYG{n}{mesh}\PYG{+w}{ }\PYG{n}{create}\PYG{+w}{ }\PYG{o}{\PYGZlt{}}\PYG{n}{nx}\PYG{o}{\PYGZgt{}}\PYG{+w}{ }\PYG{o}{\PYGZlt{}}\PYG{n}{ny}\PYG{o}{\PYGZgt{}}\PYG{+w}{ }\PYG{o}{\PYGZlt{}}\PYG{n}{nz}\PYG{o}{\PYGZgt{}}\PYG{+w}{ }\PYG{o}{\PYGZlt{}}\PYG{n}{region}\PYG{o}{\PYGZhy{}}\PYG{n}{ID}\PYG{o}{\PYGZgt{}}\PYG{+w}{ }\PYG{o}{\PYGZlt{}}\PYG{n}{f}\PYG{o}{|}\PYG{n}{p}\PYG{o}{\PYGZgt{}}\PYG{+w}{ }\PYG{o}{\PYGZlt{}}\PYG{n}{f}\PYG{o}{|}\PYG{n}{p}\PYG{o}{\PYGZgt{}}\PYG{+w}{ }\PYG{o}{\PYGZlt{}}\PYG{n}{f}\PYG{o}{|}\PYG{n}{p}\PYG{o}{\PYGZgt{}}
\end{sphinxVerbatim}
\begin{itemize}
\item {} 
\sphinxAtStartPar
AtC fixID = ID of {\hyperref[\detokenize{fix_atc::doc}]{\sphinxcrossref{\DUrole{doc}{fix atc}}}} instance

\item {} 
\sphinxAtStartPar
mesh create = name of the AtC sub\sphinxhyphen{}command

\item {} 
\sphinxAtStartPar
nx ny nz = number of elements in x\sphinxhyphen{}, y\sphinxhyphen{}, and z\sphinxhyphen{}direction

\item {} 
\sphinxAtStartPar
region\sphinxhyphen{}ID = ID of region that is to be meshed

\item {} 
\sphinxAtStartPar
f or p = periodicity flags for x\sphinxhyphen{}, y\sphinxhyphen{}, and z\sphinxhyphen{}direction

\end{itemize}


\subsubsection{Examples}
\label{\detokenize{atc_mesh_create:examples}}
\begin{sphinxVerbatim}[commandchars=\\\{\}]
\PYG{k}{fix\PYGZus{}modify }\PYG{n+nv+nvIdentifier}{AtC}\PYG{+w}{ }\PYG{n}{mesh}\PYG{+w}{ }\PYG{n}{create}\PYG{+w}{ }\PYG{l+m}{10}\PYG{+w}{ }\PYG{l+m}{1}\PYG{+w}{ }\PYG{l+m}{1}\PYG{+w}{ }\PYG{n}{feRegion}\PYG{+w}{ }\PYG{n}{p}\PYG{+w}{ }\PYG{n}{p}\PYG{+w}{ }\PYG{n}{p}
\end{sphinxVerbatim}


\subsubsection{Description}
\label{\detokenize{atc_mesh_create:description}}
\sphinxAtStartPar
Creates a uniform mesh in a rectangular region.


\subsubsection{Restrictions}
\label{\detokenize{atc_mesh_create:restrictions}}
\sphinxAtStartPar
Creates only uniform rectangular grids in a rectangular region


\subsubsection{Related AtC commands}
\label{\detokenize{atc_mesh_create:related-atc-commands}}\begin{itemize}
\item {} 
\sphinxAtStartPar
{\hyperref[\detokenize{fix_atc:atc-fix-modify}]{\sphinxcrossref{\DUrole{std,std-ref}{fix\_modify AtC command overview}}}}

\item {} 
\sphinxAtStartPar
{\hyperref[\detokenize{atc_mesh_quadrature::doc}]{\sphinxcrossref{\DUrole{doc}{fix\_modify AtC mesh quadrature}}}}

\end{itemize}


\subsubsection{Default}
\label{\detokenize{atc_mesh_create:default}}
\sphinxAtStartPar
When created, the mesh defaults to gauss2 (2\sphinxhyphen{}point Gaussian)
quadrature. Use the {\hyperref[\detokenize{atc_mesh_quadrature::doc}]{\sphinxcrossref{\DUrole{doc}{fix\_modify AtC mesh quadrature}}}} command to change the quadrature style.

\sphinxstepscope

\index{fix\_modify AtC mesh create\_elementset@\spxentry{fix\_modify AtC mesh create\_elementset}}\ignorespaces 

\subsection{fix\_modify AtC mesh create\_elementset command}
\label{\detokenize{atc_mesh_create_elementset:fix-modify-atc-mesh-create-elementset-command}}\label{\detokenize{atc_mesh_create_elementset:index-0}}\label{\detokenize{atc_mesh_create_elementset::doc}}

\subsubsection{Syntax}
\label{\detokenize{atc_mesh_create_elementset:syntax}}
\begin{sphinxVerbatim}[commandchars=\\\{\}]
\PYG{k}{fix\PYGZus{}modify }\PYG{o}{\PYGZlt{}}\PYG{n}{AtC}\PYG{+w}{ }\PYG{n}{fixID}\PYG{o}{\PYGZgt{}}\PYG{+w}{ }\PYG{n}{mesh}\PYG{+w}{ }\PYG{n}{create\PYGZus{}elementset}\PYG{+w}{ }\PYG{o}{\PYGZlt{}}\PYG{n}{id}\PYG{o}{\PYGZgt{}}\PYG{+w}{ }\PYG{o}{\PYGZlt{}}\PYG{n}{xmin}\PYG{o}{\PYGZgt{}}\PYG{+w}{ }\PYG{o}{\PYGZlt{}}\PYG{n}{xmax}\PYG{o}{\PYGZgt{}}\PYG{+w}{ }\PYG{o}{\PYGZlt{}}\PYG{n}{ymin}\PYG{o}{\PYGZgt{}}\PYG{+w}{ }\PYG{o}{\PYGZlt{}}\PYG{n}{ymax}\PYG{o}{\PYGZgt{}}\PYG{+w}{ }\PYG{o}{\PYGZlt{}}\PYG{n}{zmin}\PYG{o}{\PYGZgt{}}\PYG{+w}{ }\PYG{o}{\PYGZlt{}}\PYG{n}{zmax}\PYG{o}{\PYGZgt{}}
\end{sphinxVerbatim}
\begin{itemize}
\item {} 
\sphinxAtStartPar
AtC fixID = ID of {\hyperref[\detokenize{fix_atc::doc}]{\sphinxcrossref{\DUrole{doc}{fix atc}}}} instance

\item {} 
\sphinxAtStartPar
mesh create\_elementset = name of the AtC sub\sphinxhyphen{}command

\item {} 
\sphinxAtStartPar
id = id to assign to the collection of FE nodes

\item {} 
\sphinxAtStartPar
\textless{}xmin\textgreater{} \textless{}xmax\textgreater{} \textless{}ymin\textgreater{} \textless{}ymax\textgreater{} \textless{}zmin\textgreater{} \textless{}zmax\textgreater{} = coordinates of the bounding box that contains only the desired elements

\end{itemize}


\subsubsection{Examples}
\label{\detokenize{atc_mesh_create_elementset:examples}}
\begin{sphinxVerbatim}[commandchars=\\\{\}]
\PYG{k}{fix\PYGZus{}modify }\PYG{n+nv+nvIdentifier}{AtC}\PYG{+w}{ }\PYG{n}{mesh}\PYG{+w}{ }\PYG{n}{create\PYGZus{}elementset}\PYG{+w}{ }\PYG{n}{middle}\PYG{+w}{ }\PYG{o}{\PYGZhy{}}\PYG{l+m}{4.1}\PYG{+w}{ }\PYG{l+m}{4.1}\PYG{+w}{ }\PYG{o}{\PYGZhy{}}\PYG{l+m}{100}\PYG{+w}{ }\PYG{l+m}{100}\PYG{+w}{ }\PYG{o}{\PYGZhy{}}\PYG{l+m}{100}\PYG{+w}{ }\PYG{l+m}{1100}
\end{sphinxVerbatim}


\subsubsection{Description}
\label{\detokenize{atc_mesh_create_elementset:description}}
\sphinxAtStartPar
Command to assign an id to a set of FE elements to be used subsequently
in defining material and mesh\sphinxhyphen{}based operations.


\subsubsection{Restrictions}
\label{\detokenize{atc_mesh_create_elementset:restrictions}}
\sphinxAtStartPar
Only viable for rectangular grids.


\subsubsection{Related AtC commands}
\label{\detokenize{atc_mesh_create_elementset:related-atc-commands}}\begin{itemize}
\item {} 
\sphinxAtStartPar
{\hyperref[\detokenize{fix_atc:atc-fix-modify}]{\sphinxcrossref{\DUrole{std,std-ref}{fix\_modify AtC command overview}}}}

\item {} 
\sphinxAtStartPar
{\hyperref[\detokenize{atc_mesh_delete_elements::doc}]{\sphinxcrossref{\DUrole{doc}{fix\_modify AtC mesh delete\_elements}}}}

\item {} 
\sphinxAtStartPar
{\hyperref[\detokenize{atc_mesh_nodeset_to_elementset::doc}]{\sphinxcrossref{\DUrole{doc}{fix\_modify AtC mesh nodeset\_to\_elementset}}}}

\end{itemize}


\subsubsection{Default}
\label{\detokenize{atc_mesh_create_elementset:default}}
\sphinxAtStartPar
Coordinates are assumed to be in lattice units.

\sphinxstepscope

\index{fix\_modify AtC mesh create\_faceset box@\spxentry{fix\_modify AtC mesh create\_faceset box}}\ignorespaces 

\subsection{fix\_modify AtC mesh create\_faceset box command}
\label{\detokenize{atc_mesh_create_faceset_box:fix-modify-atc-mesh-create-faceset-box-command}}\label{\detokenize{atc_mesh_create_faceset_box:index-0}}\label{\detokenize{atc_mesh_create_faceset_box::doc}}

\subsubsection{Syntax}
\label{\detokenize{atc_mesh_create_faceset_box:syntax}}
\begin{sphinxVerbatim}[commandchars=\\\{\}]
\PYG{k}{fix\PYGZus{}modify }\PYG{o}{\PYGZlt{}}\PYG{n}{AtC}\PYG{+w}{ }\PYG{n}{fixID}\PYG{o}{\PYGZgt{}}\PYG{+w}{ }\PYG{n}{mesh}\PYG{+w}{ }\PYG{n}{create\PYGZus{}faceset}\PYG{+w}{ }\PYG{o}{\PYGZlt{}}\PYG{n}{id}\PYG{o}{\PYGZgt{}}\PYG{+w}{ }\PYG{n}{box}\PYG{+w}{ }\PYG{o}{\PYGZlt{}}\PYG{n}{xmin}\PYG{o}{\PYGZgt{}}\PYG{+w}{ }\PYG{o}{\PYGZlt{}}\PYG{n}{xmax}\PYG{o}{\PYGZgt{}}\PYG{+w}{ }\PYG{o}{\PYGZlt{}}\PYG{n}{ymin}\PYG{o}{\PYGZgt{}}\PYG{+w}{ }\PYG{o}{\PYGZlt{}}\PYG{n}{ymax}\PYG{o}{\PYGZgt{}}\PYG{+w}{ }\PYG{o}{\PYGZlt{}}\PYG{n}{zmin}\PYG{o}{\PYGZgt{}}\PYG{+w}{ }\PYG{o}{\PYGZlt{}}\PYG{n}{zmax}\PYG{o}{\PYGZgt{}}\PYG{+w}{ }\PYG{o}{\PYGZlt{}}\PYG{n}{in}\PYG{o}{|}\PYG{n}{out}\PYG{o}{\PYGZgt{}}\PYG{+w}{ }\PYG{n}{[units]}
\end{sphinxVerbatim}
\begin{itemize}
\item {} 
\sphinxAtStartPar
AtC fixID = ID of {\hyperref[\detokenize{fix_atc::doc}]{\sphinxcrossref{\DUrole{doc}{fix atc}}}} instance

\item {} 
\sphinxAtStartPar
mesh create\_faceset = name of the AtC sub\sphinxhyphen{}command

\item {} 
\sphinxAtStartPar
id = id to assign to the collection of FE faces

\item {} 
\sphinxAtStartPar
box = use bounding box to define FE faces

\item {} 
\sphinxAtStartPar
\textless{}xmin\textgreater{} \textless{}xmax\textgreater{} \textless{}ymin\textgreater{} \textless{}ymax\textgreater{} \textless{}zmin\textgreater{} \textless{}zmax\textgreater{} = coordinates of the bounding box that is coincident with the desired FE faces

\item {} 
\sphinxAtStartPar
\textless{}in|out\textgreater{} = “in” gives inner faces to the box, “out” gives the outer faces to the box

\item {} 
\sphinxAtStartPar
units = option to specify real as opposed to lattice units

\end{itemize}


\subsubsection{Examples}
\label{\detokenize{atc_mesh_create_faceset_box:examples}}
\begin{sphinxVerbatim}[commandchars=\\\{\}]
\PYG{k}{fix\PYGZus{}modify }\PYG{n+nv+nvIdentifier}{AtC}\PYG{+w}{ }\PYG{n}{mesh}\PYG{+w}{ }\PYG{n}{create\PYGZus{}faceset}\PYG{+w}{ }\PYG{n}{obndy}\PYG{+w}{ }\PYG{n}{box}\PYG{+w}{ }\PYG{o}{\PYGZhy{}}\PYG{l+m}{4.0}\PYG{+w}{ }\PYG{l+m}{4.0}\PYG{+w}{ }\PYG{o}{\PYGZhy{}}\PYG{l+m}{12}\PYG{+w}{ }\PYG{l+m}{12}\PYG{+w}{ }\PYG{o}{\PYGZhy{}}\PYG{l+m}{12}\PYG{+w}{ }\PYG{l+m}{12}\PYG{+w}{ }\PYG{n}{out}
\end{sphinxVerbatim}


\subsubsection{Description}
\label{\detokenize{atc_mesh_create_faceset_box:description}}
\sphinxAtStartPar
Command to assign an id to a set of FE faces.


\subsubsection{Restrictions}
\label{\detokenize{atc_mesh_create_faceset_box:restrictions}}
\sphinxAtStartPar
Only viable for rectangular grids.


\subsubsection{Related AtC commands}
\label{\detokenize{atc_mesh_create_faceset_box:related-atc-commands}}\begin{itemize}
\item {} 
\sphinxAtStartPar
{\hyperref[\detokenize{fix_atc:atc-fix-modify}]{\sphinxcrossref{\DUrole{std,std-ref}{fix\_modify AtC command overview}}}}

\item {} 
\sphinxAtStartPar
{\hyperref[\detokenize{atc_mesh_create_faceset_plane::doc}]{\sphinxcrossref{\DUrole{doc}{fix\_modify AtC mesh create\_faceset plane}}}}

\end{itemize}


\subsubsection{Default}
\label{\detokenize{atc_mesh_create_faceset_box:default}}
\sphinxAtStartPar
The default options are units = lattice and the use of outer faces.

\sphinxstepscope

\index{fix\_modify AtC mesh create\_faceset plane@\spxentry{fix\_modify AtC mesh create\_faceset plane}}\ignorespaces 

\subsection{fix\_modify AtC mesh create\_faceset plane command}
\label{\detokenize{atc_mesh_create_faceset_plane:fix-modify-atc-mesh-create-faceset-plane-command}}\label{\detokenize{atc_mesh_create_faceset_plane:index-0}}\label{\detokenize{atc_mesh_create_faceset_plane::doc}}

\subsubsection{Syntax}
\label{\detokenize{atc_mesh_create_faceset_plane:syntax}}
\begin{sphinxVerbatim}[commandchars=\\\{\}]
\PYG{k}{fix\PYGZus{}modify }\PYG{o}{\PYGZlt{}}\PYG{n}{AtC}\PYG{+w}{ }\PYG{n}{fixID}\PYG{o}{\PYGZgt{}}\PYG{+w}{ }\PYG{n}{mesh}\PYG{+w}{ }\PYG{n}{create\PYGZus{}faceset}\PYG{+w}{ }\PYG{o}{\PYGZlt{}}\PYG{n}{id}\PYG{o}{\PYGZgt{}}\PYG{+w}{ }\PYG{n}{plane}\PYG{+w}{ }\PYG{o}{\PYGZlt{}}\PYG{n}{x}\PYG{o}{|}\PYG{n}{y}\PYG{o}{|}\PYG{n}{z}\PYG{o}{\PYGZgt{}}\PYG{+w}{ }\PYG{o}{\PYGZlt{}}\PYG{n}{val1}\PYG{o}{\PYGZgt{}}\PYG{+w}{ }\PYG{o}{\PYGZlt{}}\PYG{n}{x}\PYG{o}{|}\PYG{n}{y}\PYG{o}{|}\PYG{n}{z}\PYG{o}{\PYGZgt{}}\PYG{+w}{ }\PYG{o}{\PYGZlt{}}\PYG{n}{lval2}\PYG{o}{\PYGZgt{}}\PYG{+w}{ }\PYG{o}{\PYGZlt{}}\PYG{n}{uval2}\PYG{o}{\PYGZgt{}}\PYG{+w}{ }\PYG{n}{[units]}
\end{sphinxVerbatim}
\begin{itemize}
\item {} 
\sphinxAtStartPar
AtC fixID = ID of {\hyperref[\detokenize{fix_atc::doc}]{\sphinxcrossref{\DUrole{doc}{fix atc}}}} instance

\item {} 
\sphinxAtStartPar
mesh create\_faceset = name of the AtC sub\sphinxhyphen{}command

\item {} 
\sphinxAtStartPar
id = id to assign to the collection of FE faces

\item {} 
\sphinxAtStartPar
plane = use plane to define faceset

\item {} 
\sphinxAtStartPar
\textless{}val1\textgreater{},\textless{}lval2\textgreater{},\textless{}uval2\textgreater{} = plane is specified as the x|y|z=val1 plane bounded by the segments x|y|z = {[}lval2,uval2{]}

\item {} 
\sphinxAtStartPar
units = option to specify real as opposed to lattice units

\end{itemize}


\subsubsection{Examples}
\label{\detokenize{atc_mesh_create_faceset_plane:examples}}
\begin{sphinxVerbatim}[commandchars=\\\{\}]
\PYG{k}{fix\PYGZus{}modify }\PYG{n+nv+nvIdentifier}{AtC}\PYG{+w}{ }\PYG{n}{mesh}\PYG{+w}{ }\PYG{n}{create\PYGZus{}faceset}\PYG{+w}{ }\PYG{n}{xyplane}\PYG{+w}{ }\PYG{n}{plane}\PYG{+w}{ }\PYG{n}{y}\PYG{+w}{ }\PYG{l+m}{0}\PYG{+w}{ }\PYG{n}{x}\PYG{+w}{ }\PYG{o}{\PYGZhy{}}\PYG{l+m}{4}\PYG{+w}{ }\PYG{l+m}{0}
\end{sphinxVerbatim}


\subsubsection{Description}
\label{\detokenize{atc_mesh_create_faceset_plane:description}}
\sphinxAtStartPar
Command to assign an id to a set of FE faces.


\subsubsection{Restrictions}
\label{\detokenize{atc_mesh_create_faceset_plane:restrictions}}
\sphinxAtStartPar
Only viable for rectangular grids.


\subsubsection{Related AtC commands}
\label{\detokenize{atc_mesh_create_faceset_plane:related-atc-commands}}\begin{itemize}
\item {} 
\sphinxAtStartPar
{\hyperref[\detokenize{fix_atc:atc-fix-modify}]{\sphinxcrossref{\DUrole{std,std-ref}{fix\_modify AtC command overview}}}}

\item {} 
\sphinxAtStartPar
{\hyperref[\detokenize{atc_mesh_create_faceset_box::doc}]{\sphinxcrossref{\DUrole{doc}{fix\_modify AtC mesh create\_faceset box}}}}

\end{itemize}


\subsubsection{Default}
\label{\detokenize{atc_mesh_create_faceset_plane:default}}
\sphinxAtStartPar
The default options are units = lattice.

\sphinxstepscope

\index{fix\_modify AtC mesh create\_nodeset@\spxentry{fix\_modify AtC mesh create\_nodeset}}\ignorespaces 

\subsection{fix\_modify AtC mesh create\_nodeset command}
\label{\detokenize{atc_mesh_create_nodeset:fix-modify-atc-mesh-create-nodeset-command}}\label{\detokenize{atc_mesh_create_nodeset:index-0}}\label{\detokenize{atc_mesh_create_nodeset::doc}}

\subsubsection{Syntax}
\label{\detokenize{atc_mesh_create_nodeset:syntax}}
\begin{sphinxVerbatim}[commandchars=\\\{\}]
\PYG{k}{fix\PYGZus{}modify }\PYG{o}{\PYGZlt{}}\PYG{n}{AtC}\PYG{+w}{ }\PYG{n}{fixID}\PYG{o}{\PYGZgt{}}\PYG{+w}{ }\PYG{n}{mesh}\PYG{+w}{ }\PYG{n}{create\PYGZus{}nodeset}\PYG{+w}{ }\PYG{o}{\PYGZlt{}}\PYG{n}{id}\PYG{o}{\PYGZgt{}}\PYG{+w}{ }\PYG{o}{\PYGZlt{}}\PYG{n}{xmin}\PYG{o}{\PYGZgt{}}\PYG{+w}{ }\PYG{o}{\PYGZlt{}}\PYG{n}{xmax}\PYG{o}{\PYGZgt{}}\PYG{+w}{ }\PYG{o}{\PYGZlt{}}\PYG{n}{ymin}\PYG{o}{\PYGZgt{}}\PYG{+w}{ }\PYG{o}{\PYGZlt{}}\PYG{n}{ymax}\PYG{o}{\PYGZgt{}}\PYG{+w}{ }\PYG{o}{\PYGZlt{}}\PYG{n}{zmin}\PYG{o}{\PYGZgt{}}\PYG{+w}{ }\PYG{o}{\PYGZlt{}}\PYG{n}{zmax}\PYG{o}{\PYGZgt{}}
\end{sphinxVerbatim}
\begin{itemize}
\item {} 
\sphinxAtStartPar
AtC fixID = ID of {\hyperref[\detokenize{fix_atc::doc}]{\sphinxcrossref{\DUrole{doc}{fix atc}}}} instance

\item {} 
\sphinxAtStartPar
mesh create\_nodeset = name of the AtC sub\sphinxhyphen{}command

\item {} 
\sphinxAtStartPar
id = id to assign to the collection of FE nodes

\item {} 
\sphinxAtStartPar
\textless{}xmin\textgreater{} \textless{}xmax\textgreater{} \textless{}ymin\textgreater{} \textless{}ymax\textgreater{} \textless{}zmin\textgreater{} \textless{}zmax\textgreater{} = coordinates of the bounding box that contains only the desired nodes

\end{itemize}


\subsubsection{Examples}
\label{\detokenize{atc_mesh_create_nodeset:examples}}
\begin{sphinxVerbatim}[commandchars=\\\{\}]
\PYG{k}{fix\PYGZus{}modify }\PYG{n+nv+nvIdentifier}{AtC}\PYG{+w}{ }\PYG{n}{mesh}\PYG{+w}{ }\PYG{n}{create\PYGZus{}nodeset}\PYG{+w}{ }\PYG{n}{lbc}\PYG{+w}{ }\PYG{o}{\PYGZhy{}}\PYG{l+m}{12.1}\PYG{+w}{ }\PYG{o}{\PYGZhy{}}\PYG{l+m}{11.9}\PYG{+w}{ }\PYG{o}{\PYGZhy{}}\PYG{l+m}{12}\PYG{+w}{ }\PYG{l+m}{12}\PYG{+w}{ }\PYG{o}{\PYGZhy{}}\PYG{l+m}{12}\PYG{+w}{ }\PYG{l+m}{12}
\end{sphinxVerbatim}


\subsubsection{Description}
\label{\detokenize{atc_mesh_create_nodeset:description}}
\sphinxAtStartPar
Command to assign an id to a set of FE nodes to be used subsequently in
defining boundary conditions.


\subsubsection{Restrictions}
\label{\detokenize{atc_mesh_create_nodeset:restrictions}}
\sphinxAtStartPar
None


\subsubsection{Related AtC commands}
\label{\detokenize{atc_mesh_create_nodeset:related-atc-commands}}\begin{itemize}
\item {} 
\sphinxAtStartPar
{\hyperref[\detokenize{fix_atc:atc-fix-modify}]{\sphinxcrossref{\DUrole{std,std-ref}{fix\_modify AtC command overview}}}}

\item {} 
\sphinxAtStartPar
{\hyperref[\detokenize{atc_mesh_add_to_nodeset::doc}]{\sphinxcrossref{\DUrole{doc}{fix\_modify AtC mesh add\_to\_nodeset}}}}

\end{itemize}


\subsubsection{Default}
\label{\detokenize{atc_mesh_create_nodeset:default}}
\sphinxAtStartPar
Coordinates are assumed to be in lattice units.

\sphinxstepscope

\index{fix\_modify AtC mesh delete\_elements@\spxentry{fix\_modify AtC mesh delete\_elements}}\ignorespaces 

\subsection{fix\_modify AtC mesh delete\_elements command}
\label{\detokenize{atc_mesh_delete_elements:fix-modify-atc-mesh-delete-elements-command}}\label{\detokenize{atc_mesh_delete_elements:index-0}}\label{\detokenize{atc_mesh_delete_elements::doc}}

\subsubsection{Syntax}
\label{\detokenize{atc_mesh_delete_elements:syntax}}
\begin{sphinxVerbatim}[commandchars=\\\{\}]
\PYG{k}{fix\PYGZus{}modify }\PYG{o}{\PYGZlt{}}\PYG{n}{AtC}\PYG{+w}{ }\PYG{n}{fixID}\PYG{o}{\PYGZgt{}}\PYG{+w}{ }\PYG{n}{mesh}\PYG{+w}{ }\PYG{n}{delete\PYGZus{}elements}\PYG{+w}{ }\PYG{o}{\PYGZlt{}}\PYG{n}{id}\PYG{o}{\PYGZgt{}}
\end{sphinxVerbatim}
\begin{itemize}
\item {} 
\sphinxAtStartPar
AtC fixID = ID of {\hyperref[\detokenize{fix_atc::doc}]{\sphinxcrossref{\DUrole{doc}{fix atc}}}} instance

\item {} 
\sphinxAtStartPar
mesh create\_elementset = name of the AtC sub\sphinxhyphen{}command

\item {} 
\sphinxAtStartPar
id = id of the element set

\end{itemize}


\subsubsection{Examples}
\label{\detokenize{atc_mesh_delete_elements:examples}}
\begin{sphinxVerbatim}[commandchars=\\\{\}]
\PYG{k}{fix\PYGZus{}modify }\PYG{n+nv+nvIdentifier}{AtC}\PYG{+w}{ }\PYG{n}{mesh}\PYG{+w}{ }\PYG{n}{delete\PYGZus{}elements}\PYG{+w}{ }\PYG{n}{gap}
\end{sphinxVerbatim}


\subsubsection{Description}
\label{\detokenize{atc_mesh_delete_elements:description}}
\sphinxAtStartPar
Deletes a group of elements from the mesh.


\subsubsection{Restrictions}
\label{\detokenize{atc_mesh_delete_elements:restrictions}}
\sphinxAtStartPar
None.


\subsubsection{Related AtC commands}
\label{\detokenize{atc_mesh_delete_elements:related-atc-commands}}\begin{itemize}
\item {} 
\sphinxAtStartPar
{\hyperref[\detokenize{fix_atc:atc-fix-modify}]{\sphinxcrossref{\DUrole{std,std-ref}{fix\_modify AtC command overview}}}}

\item {} 
\sphinxAtStartPar
{\hyperref[\detokenize{atc_mesh_create_elementset::doc}]{\sphinxcrossref{\DUrole{doc}{fix\_modify AtC mesh create\_elementset}}}}

\item {} 
\sphinxAtStartPar
{\hyperref[\detokenize{atc_mesh_nodeset_to_elementset::doc}]{\sphinxcrossref{\DUrole{doc}{fix\_modify AtC mesh nodeset\_to\_elementset}}}}

\end{itemize}


\subsubsection{Default}
\label{\detokenize{atc_mesh_delete_elements:default}}
\sphinxAtStartPar
None.

\sphinxstepscope

\index{fix\_modify AtC mesh nodeset\_to\_elementset@\spxentry{fix\_modify AtC mesh nodeset\_to\_elementset}}\ignorespaces 

\subsection{fix\_modify AtC mesh nodeset\_to\_elementset command}
\label{\detokenize{atc_mesh_nodeset_to_elementset:fix-modify-atc-mesh-nodeset-to-elementset-command}}\label{\detokenize{atc_mesh_nodeset_to_elementset:index-0}}\label{\detokenize{atc_mesh_nodeset_to_elementset::doc}}

\subsubsection{Syntax}
\label{\detokenize{atc_mesh_nodeset_to_elementset:syntax}}
\begin{sphinxVerbatim}[commandchars=\\\{\}]
\PYG{k}{fix\PYGZus{}modify }\PYG{o}{\PYGZlt{}}\PYG{n}{AtC}\PYG{+w}{ }\PYG{n}{fixID}\PYG{o}{\PYGZgt{}}\PYG{+w}{ }\PYG{n}{mesh}\PYG{+w}{ }\PYG{n}{nodeset\PYGZus{}to\PYGZus{}elementset}\PYG{+w}{ }\PYG{o}{\PYGZlt{}}\PYG{n}{nodeset\PYGZus{}id}\PYG{o}{\PYGZgt{}}\PYG{+w}{ }\PYG{o}{\PYGZlt{}}\PYG{n}{elementset\PYGZus{}id}\PYG{o}{\PYGZgt{}}\PYG{+w}{ }\PYG{o}{\PYGZlt{}}\PYG{n}{max}\PYG{o}{/}\PYG{n}{min}\PYG{o}{\PYGZgt{}}
\end{sphinxVerbatim}
\begin{itemize}
\item {} 
\sphinxAtStartPar
AtC fixID = ID of {\hyperref[\detokenize{fix_atc::doc}]{\sphinxcrossref{\DUrole{doc}{fix atc}}}} instance

\item {} 
\sphinxAtStartPar
mesh nodeset\_to\_elementset = name of the AtC sub\sphinxhyphen{}command

\item {} 
\sphinxAtStartPar
nodeset\_id = id of desired nodeset from which to create the elementset

\item {} 
\sphinxAtStartPar
elementset\_id = id to assign to the collection of FE elements

\item {} 
\sphinxAtStartPar
\textless{}max/min\textgreater{} = flag to choose either the maximal or minimal elementset

\end{itemize}


\subsubsection{Examples}
\label{\detokenize{atc_mesh_nodeset_to_elementset:examples}}
\begin{sphinxVerbatim}[commandchars=\\\{\}]
\PYG{k}{fix\PYGZus{}modify }\PYG{n+nv+nvIdentifier}{AtC}\PYG{+w}{ }\PYG{n}{mesh}\PYG{+w}{ }\PYG{n}{nodeset\PYGZus{}to\PYGZus{}elementset}\PYG{+w}{ }\PYG{n}{myNodeset}\PYG{+w}{ }\PYG{n}{myElementset}\PYG{+w}{ }\PYG{n}{min}
\end{sphinxVerbatim}


\subsubsection{Description}
\label{\detokenize{atc_mesh_nodeset_to_elementset:description}}
\sphinxAtStartPar
Command to create an elementset from an existing nodeset. Either the
minimal element set of elements with all nodes in the set, or maximal
element set with all elements with at least one node in the set, can be
created.


\subsubsection{Restrictions}
\label{\detokenize{atc_mesh_nodeset_to_elementset:restrictions}}
\sphinxAtStartPar
None.


\subsubsection{Related AtC commands}
\label{\detokenize{atc_mesh_nodeset_to_elementset:related-atc-commands}}\begin{itemize}
\item {} 
\sphinxAtStartPar
{\hyperref[\detokenize{fix_atc:atc-fix-modify}]{\sphinxcrossref{\DUrole{std,std-ref}{fix\_modify AtC command overview}}}}

\item {} 
\sphinxAtStartPar
{\hyperref[\detokenize{atc_mesh_create_elementset::doc}]{\sphinxcrossref{\DUrole{doc}{fix\_modify AtC mesh create\_elementset}}}}

\item {} 
\sphinxAtStartPar
{\hyperref[\detokenize{atc_mesh_delete_elements::doc}]{\sphinxcrossref{\DUrole{doc}{fix\_modify AtC mesh delete\_elements}}}}

\end{itemize}


\subsubsection{Default}
\label{\detokenize{atc_mesh_nodeset_to_elementset:default}}
\sphinxAtStartPar
Unless specified, the maximal element set is created.

\sphinxstepscope

\index{fix\_modify AtC mesh output@\spxentry{fix\_modify AtC mesh output}}\ignorespaces 

\subsection{fix\_modify AtC mesh output command}
\label{\detokenize{atc_mesh_output:fix-modify-atc-mesh-output-command}}\label{\detokenize{atc_mesh_output:index-0}}\label{\detokenize{atc_mesh_output::doc}}

\subsubsection{Syntax}
\label{\detokenize{atc_mesh_output:syntax}}
\begin{sphinxVerbatim}[commandchars=\\\{\}]
\PYG{k}{fix\PYGZus{}modify }\PYG{o}{\PYGZlt{}}\PYG{n}{AtC}\PYG{+w}{ }\PYG{n}{fixID}\PYG{o}{\PYGZgt{}}\PYG{+w}{ }\PYG{n}{mesh}\PYG{+w}{ }\PYG{n}{output}\PYG{+w}{ }\PYG{o}{\PYGZlt{}}\PYG{n}{file\PYGZus{}prefix}\PYG{o}{\PYGZgt{}}
\end{sphinxVerbatim}
\begin{itemize}
\item {} 
\sphinxAtStartPar
AtC fixID = ID of {\hyperref[\detokenize{fix_atc::doc}]{\sphinxcrossref{\DUrole{doc}{fix atc}}}} instance

\item {} 
\sphinxAtStartPar
mesh output = name of the AtC sub\sphinxhyphen{}command

\item {} 
\sphinxAtStartPar
file\_prefix = prefix of various generated output files

\end{itemize}


\subsubsection{Examples}
\label{\detokenize{atc_mesh_output:examples}}
\begin{sphinxVerbatim}[commandchars=\\\{\}]
\PYG{k}{fix\PYGZus{}modify }\PYG{n+nv+nvIdentifier}{AtC}\PYG{+w}{ }\PYG{n}{mesh}\PYG{+w}{ }\PYG{n}{output}\PYG{+w}{ }\PYG{n}{meshData}
\end{sphinxVerbatim}


\subsubsection{Description}
\label{\detokenize{atc_mesh_output:description}}
\sphinxAtStartPar
Command to output mesh and associated data: nodesets, facesets, and
elementsets. This data is only output once upon initialization since
currently the mesh is static. Creates binary (EnSight, “gold” format)
output of mesh data.


\subsubsection{Restrictions}
\label{\detokenize{atc_mesh_output:restrictions}}
\sphinxAtStartPar
None.


\subsubsection{Related AtC commands}
\label{\detokenize{atc_mesh_output:related-atc-commands}}\begin{itemize}
\item {} 
\sphinxAtStartPar
{\hyperref[\detokenize{fix_atc:atc-fix-modify}]{\sphinxcrossref{\DUrole{std,std-ref}{fix\_modify AtC command overview}}}}

\end{itemize}


\subsubsection{Default}
\label{\detokenize{atc_mesh_output:default}}
\sphinxAtStartPar
None.

\sphinxstepscope

\index{fix\_modify AtC mesh quadrature@\spxentry{fix\_modify AtC mesh quadrature}}\ignorespaces 

\subsection{fix\_modify AtC mesh quadrature command}
\label{\detokenize{atc_mesh_quadrature:fix-modify-atc-mesh-quadrature-command}}\label{\detokenize{atc_mesh_quadrature:index-0}}\label{\detokenize{atc_mesh_quadrature::doc}}

\subsubsection{Syntax}
\label{\detokenize{atc_mesh_quadrature:syntax}}
\begin{sphinxVerbatim}[commandchars=\\\{\}]
\PYG{k}{fix\PYGZus{}modify }\PYG{o}{\PYGZlt{}}\PYG{n}{AtC}\PYG{+w}{ }\PYG{n}{fixID}\PYG{o}{\PYGZgt{}}\PYG{+w}{ }\PYG{n}{mesh}\PYG{+w}{ }\PYG{n}{quatrature}\PYG{+w}{ }\PYG{o}{\PYGZlt{}}\PYG{n}{quad}\PYG{o}{\PYGZgt{}}
\end{sphinxVerbatim}
\begin{itemize}
\item {} 
\sphinxAtStartPar
AtC fixID = ID of {\hyperref[\detokenize{fix_atc::doc}]{\sphinxcrossref{\DUrole{doc}{fix atc}}}} instance

\item {} 
\sphinxAtStartPar
mesh quadrature = name of the AtC sub\sphinxhyphen{}command

\item {} 
\sphinxAtStartPar
quad = \sphinxstyleemphasis{nodal} or \sphinxstyleemphasis{gauss1} or \sphinxstyleemphasis{gauss2} or \sphinxstyleemphasis{gauss3} or \sphinxstyleemphasis{face}

\end{itemize}


\subsubsection{Examples}
\label{\detokenize{atc_mesh_quadrature:examples}}
\begin{sphinxVerbatim}[commandchars=\\\{\}]
\PYG{k}{fix\PYGZus{}modify }\PYG{n+nv+nvIdentifier}{AtC}\PYG{+w}{ }\PYG{n}{mesh}\PYG{+w}{ }\PYG{n}{quadrature}\PYG{+w}{ }\PYG{n}{face}
\end{sphinxVerbatim}


\subsubsection{Description}
\label{\detokenize{atc_mesh_quadrature:description}}
\sphinxAtStartPar
(Re\sphinxhyphen{})assigns the quadrature style for an existing mesh.  When a mesh is
created its quadrature method defaults to gauss2.  Use this call to
change it after the fact.


\subsubsection{Restrictions}
\label{\detokenize{atc_mesh_quadrature:restrictions}}
\sphinxAtStartPar
None.


\subsubsection{Related AtC commands}
\label{\detokenize{atc_mesh_quadrature:related-atc-commands}}\begin{itemize}
\item {} 
\sphinxAtStartPar
{\hyperref[\detokenize{fix_atc:atc-fix-modify}]{\sphinxcrossref{\DUrole{std,std-ref}{fix\_modify AtC command overview}}}}

\item {} 
\sphinxAtStartPar
{\hyperref[\detokenize{atc_mesh_create::doc}]{\sphinxcrossref{\DUrole{doc}{fix\_modify AtC mesh create}}}}

\end{itemize}


\subsubsection{Default}
\label{\detokenize{atc_mesh_quadrature:default}}
\sphinxAtStartPar
None.

\sphinxstepscope

\index{fix\_modify AtC mesh read@\spxentry{fix\_modify AtC mesh read}}\ignorespaces 

\subsection{fix\_modify AtC mesh read command}
\label{\detokenize{atc_mesh_read:fix-modify-atc-mesh-read-command}}\label{\detokenize{atc_mesh_read:index-0}}\label{\detokenize{atc_mesh_read::doc}}

\subsubsection{Syntax}
\label{\detokenize{atc_mesh_read:syntax}}
\begin{sphinxVerbatim}[commandchars=\\\{\}]
\PYG{k}{fix\PYGZus{}modify }\PYG{o}{\PYGZlt{}}\PYG{n}{AtC}\PYG{+w}{ }\PYG{n}{fixID}\PYG{o}{\PYGZgt{}}\PYG{+w}{ }\PYG{n}{mesh}\PYG{+w}{ }\PYG{n}{read}\PYG{+w}{ }\PYG{o}{\PYGZlt{}}\PYG{n}{f}\PYG{o}{|}\PYG{n}{p}\PYG{o}{\PYGZgt{}}\PYG{+w}{ }\PYG{o}{\PYGZlt{}}\PYG{n}{f}\PYG{o}{|}\PYG{n}{p}\PYG{o}{\PYGZgt{}}\PYG{+w}{ }\PYG{o}{\PYGZlt{}}\PYG{n}{f}\PYG{o}{|}\PYG{n}{p}\PYG{o}{\PYGZgt{}}
\end{sphinxVerbatim}
\begin{itemize}
\item {} 
\sphinxAtStartPar
AtC fixID = ID of {\hyperref[\detokenize{fix_atc::doc}]{\sphinxcrossref{\DUrole{doc}{fix atc}}}} instance

\item {} 
\sphinxAtStartPar
mesh read = name of the AtC sub\sphinxhyphen{}command

\item {} 
\sphinxAtStartPar
filename = name of the file containing the mesh to be read

\item {} 
\sphinxAtStartPar
f or p = periodicity flags for x\sphinxhyphen{}, y\sphinxhyphen{}, and z\sphinxhyphen{}direction (optional)

\end{itemize}


\subsubsection{Examples}
\label{\detokenize{atc_mesh_read:examples}}
\begin{sphinxVerbatim}[commandchars=\\\{\}]
\PYG{k}{fix\PYGZus{}modify }\PYG{n+nv+nvIdentifier}{AtC}\PYG{+w}{ }\PYG{n}{mesh}\PYG{+w}{ }\PYG{n}{read}\PYG{+w}{ }\PYG{n}{myComponent.mesh}\PYG{+w}{ }\PYG{n}{p}\PYG{+w}{ }\PYG{n}{p}\PYG{+w}{ }\PYG{n}{p}
\PYG{k}{fix\PYGZus{}modify }\PYG{n+nv+nvIdentifier}{AtC}\PYG{+w}{ }\PYG{n}{mesh}\PYG{+w}{ }\PYG{n}{read}\PYG{+w}{ }\PYG{n}{myOtherComponent.exo}
\end{sphinxVerbatim}


\subsubsection{Description}
\label{\detokenize{atc_mesh_read:description}}
\sphinxAtStartPar
Reads a mesh from a text or exodus file, and assigns periodic boundary
conditions if needed.


\subsubsection{Restrictions}
\label{\detokenize{atc_mesh_read:restrictions}}
\sphinxAtStartPar
None


\subsubsection{Related AtC commands}
\label{\detokenize{atc_mesh_read:related-atc-commands}}\begin{itemize}
\item {} 
\sphinxAtStartPar
{\hyperref[\detokenize{fix_atc:atc-fix-modify}]{\sphinxcrossref{\DUrole{std,std-ref}{fix\_modify AtC command overview}}}}

\item {} 
\sphinxAtStartPar
{\hyperref[\detokenize{atc_mesh_create::doc}]{\sphinxcrossref{\DUrole{doc}{fix\_modify AtC mesh create}}}}

\item {} 
\sphinxAtStartPar
{\hyperref[\detokenize{atc_mesh_write::doc}]{\sphinxcrossref{\DUrole{doc}{fix\_modify AtC mesh write}}}}

\end{itemize}


\subsubsection{Default}
\label{\detokenize{atc_mesh_read:default}}
\sphinxAtStartPar
Periodicity flags are set to false (f) by default.

\sphinxstepscope

\index{fix\_modify AtC mesh write@\spxentry{fix\_modify AtC mesh write}}\ignorespaces 

\subsection{fix\_modify AtC mesh write command}
\label{\detokenize{atc_mesh_write:fix-modify-atc-mesh-write-command}}\label{\detokenize{atc_mesh_write:index-0}}\label{\detokenize{atc_mesh_write::doc}}

\subsubsection{Syntax}
\label{\detokenize{atc_mesh_write:syntax}}
\begin{sphinxVerbatim}[commandchars=\\\{\}]
\PYG{k}{fix\PYGZus{}modify }\PYG{o}{\PYGZlt{}}\PYG{n}{AtC}\PYG{+w}{ }\PYG{n}{fixID}\PYG{o}{\PYGZgt{}}\PYG{+w}{ }\PYG{n}{mesh}\PYG{+w}{ }\PYG{n}{write}\PYG{+w}{ }\PYG{o}{\PYGZlt{}}\PYG{n}{f}\PYG{o}{|}\PYG{n}{p}\PYG{o}{\PYGZgt{}}\PYG{+w}{ }\PYG{o}{\PYGZlt{}}\PYG{n}{f}\PYG{o}{|}\PYG{n}{p}\PYG{o}{\PYGZgt{}}\PYG{+w}{ }\PYG{o}{\PYGZlt{}}\PYG{n}{f}\PYG{o}{|}\PYG{n}{p}\PYG{o}{\PYGZgt{}}
\end{sphinxVerbatim}
\begin{itemize}
\item {} 
\sphinxAtStartPar
AtC fixID = ID of {\hyperref[\detokenize{fix_atc::doc}]{\sphinxcrossref{\DUrole{doc}{fix atc}}}} instance

\item {} 
\sphinxAtStartPar
mesh write = name of the AtC sub\sphinxhyphen{}command

\item {} 
\sphinxAtStartPar
filename = name of the file containing the mesh to be write

\end{itemize}


\subsubsection{Examples}
\label{\detokenize{atc_mesh_write:examples}}
\begin{sphinxVerbatim}[commandchars=\\\{\}]
\PYG{k}{fix\PYGZus{}modify }\PYG{n+nv+nvIdentifier}{AtC}\PYG{+w}{ }\PYG{n}{mesh}\PYG{+w}{ }\PYG{n}{write}\PYG{+w}{ }\PYG{n}{myMesh.mesh}
\end{sphinxVerbatim}


\subsubsection{Description}
\label{\detokenize{atc_mesh_write:description}}
\sphinxAtStartPar
Writes a mesh to a text file.


\subsubsection{Restrictions}
\label{\detokenize{atc_mesh_write:restrictions}}
\sphinxAtStartPar
None


\subsubsection{Related AtC commands}
\label{\detokenize{atc_mesh_write:related-atc-commands}}\begin{itemize}
\item {} 
\sphinxAtStartPar
{\hyperref[\detokenize{fix_atc:atc-fix-modify}]{\sphinxcrossref{\DUrole{std,std-ref}{fix\_modify AtC command overview}}}}

\item {} 
\sphinxAtStartPar
{\hyperref[\detokenize{atc_mesh_create::doc}]{\sphinxcrossref{\DUrole{doc}{fix\_modify AtC mesh create}}}}

\item {} 
\sphinxAtStartPar
{\hyperref[\detokenize{atc_mesh_read::doc}]{\sphinxcrossref{\DUrole{doc}{fix\_modify AtC mesh read}}}}

\end{itemize}


\subsubsection{Default}
\label{\detokenize{atc_mesh_write:default}}
\sphinxAtStartPar
None.

\sphinxstepscope

\index{fix\_modify AtC output@\spxentry{fix\_modify AtC output}}\ignorespaces 

\subsection{fix\_modify AtC output command}
\label{\detokenize{atc_output:fix-modify-atc-output-command}}\label{\detokenize{atc_output:index-0}}\label{\detokenize{atc_output::doc}}

\subsubsection{Syntax}
\label{\detokenize{atc_output:syntax}}
\begin{sphinxVerbatim}[commandchars=\\\{\}]
\PYG{k}{fix\PYGZus{}modify }\PYG{o}{\PYGZlt{}}\PYG{n}{AtC}\PYG{+w}{ }\PYG{n}{fixID}\PYG{o}{\PYGZgt{}}\PYG{+w}{ }\PYG{n}{output}\PYG{+w}{ }\PYG{o}{\PYGZlt{}}\PYG{n}{filename\PYGZus{}prefix}\PYG{o}{\PYGZgt{}}\PYG{+w}{ }\PYG{o}{\PYGZlt{}}\PYG{n}{frequency}\PYG{o}{\PYGZgt{}}\PYG{+w}{ }\PYG{n}{[text}\PYG{o}{|}\PYG{n}{full\PYGZus{}text}\PYG{o}{|}\PYG{n}{binary}\PYG{o}{|}\PYG{n}{vector\PYGZus{}components}\PYG{o}{|}\PYG{n}{tensor\PYGZus{}components]}
\PYG{k}{fix\PYGZus{}modify }\PYG{o}{\PYGZlt{}}\PYG{n}{AtC}\PYG{+w}{ }\PYG{n}{fixID}\PYG{o}{\PYGZgt{}}\PYG{+w}{ }\PYG{n}{output}\PYG{+w}{ }\PYG{n}{index}\PYG{+w}{ }\PYG{n}{[step}\PYG{o}{|}\PYG{n}{time]}
\end{sphinxVerbatim}
\begin{itemize}
\item {} 
\sphinxAtStartPar
AtC fixID = ID of {\hyperref[\detokenize{fix_atc::doc}]{\sphinxcrossref{\DUrole{doc}{fix atc}}}} instance

\item {} 
\sphinxAtStartPar
\sphinxstyleemphasis{output} or \sphinxstyleemphasis{output index} = name of the AtC sub\sphinxhyphen{}command

\item {} 
\sphinxAtStartPar
filename\_prefix = prefix for data files (for \sphinxstyleemphasis{output})

\item {} 
\sphinxAtStartPar
frequency = frequency of output in timesteps (for \sphinxstyleemphasis{output})

\item {} 
\sphinxAtStartPar
optional keywords for \sphinxstyleemphasis{output}:
\begin{itemize}
\item {} 
\sphinxAtStartPar
text = creates text output of index, step and nodal variable values for unique nodes

\item {} 
\sphinxAtStartPar
full\_text = creates text output index, nodal id, step, nodal coordinates and nodal variable values for unique and image nodes

\item {} 
\sphinxAtStartPar
binary = creates binary EnSight output

\item {} 
\sphinxAtStartPar
vector\_components = outputs vectors as scalar components

\item {} 
\sphinxAtStartPar
tensor\_components = outputs tensor as scalar components (for use with ParaView)

\end{itemize}

\item {} 
\sphinxAtStartPar
\sphinxstyleemphasis{step} or \sphinxstyleemphasis{time} = index output by step or by time (for \sphinxstyleemphasis{output index})

\end{itemize}


\subsubsection{Examples}
\label{\detokenize{atc_output:examples}}
\begin{sphinxVerbatim}[commandchars=\\\{\}]
\PYG{k}{fix\PYGZus{}modify }\PYG{n+nv+nvIdentifier}{AtC}\PYG{+w}{ }\PYG{n}{output}\PYG{+w}{ }\PYG{n}{heatFE}\PYG{+w}{ }\PYG{l+m}{100}
\PYG{k}{fix\PYGZus{}modify }\PYG{n+nv+nvIdentifier}{AtC}\PYG{+w}{ }\PYG{n}{output}\PYG{+w}{ }\PYG{n}{hardyFE}\PYG{+w}{ }\PYG{l+m}{1}\PYG{+w}{ }\PYG{n}{text}\PYG{+w}{ }\PYG{n}{tensor\PYGZus{}components}
\PYG{k}{fix\PYGZus{}modify }\PYG{n+nv+nvIdentifier}{AtC}\PYG{+w}{ }\PYG{n}{output}\PYG{+w}{ }\PYG{n}{hardyFE}\PYG{+w}{ }\PYG{l+m}{10}\PYG{+w}{ }\PYG{n}{text}\PYG{+w}{ }\PYG{n}{binary}\PYG{+w}{ }\PYG{n}{tensor\PYGZus{}components}
\PYG{k}{fix\PYGZus{}modify }\PYG{n+nv+nvIdentifier}{AtC}\PYG{+w}{ }\PYG{n}{output}\PYG{+w}{ }\PYG{n}{index}\PYG{+w}{ }\PYG{n}{step}
\end{sphinxVerbatim}


\subsubsection{Description}
\label{\detokenize{atc_output:description}}
\sphinxAtStartPar
Creates text and/or binary (EnSight, “gold” format) output of nodal/mesh
data which is transfer/physics specific. Output indexing by step or time
is possible.


\subsubsection{Restrictions}
\label{\detokenize{atc_output:restrictions}}
\sphinxAtStartPar
None.


\subsubsection{Related AtC commands}
\label{\detokenize{atc_output:related-atc-commands}}\begin{itemize}
\item {} 
\sphinxAtStartPar
{\hyperref[\detokenize{fix_atc:atc-fix-modify}]{\sphinxcrossref{\DUrole{std,std-ref}{fix\_modify AtC command overview}}}}

\item {} 
\sphinxAtStartPar
{\hyperref[\detokenize{fix_atc::doc}]{\sphinxcrossref{\DUrole{doc}{fix atc command}}}}

\end{itemize}


\subsubsection{Default}
\label{\detokenize{atc_output:default}}
\sphinxAtStartPar
No default format. Output indexed by time.

\sphinxstepscope

\index{fix\_modify AtC output boundary\_integral@\spxentry{fix\_modify AtC output boundary\_integral}}\ignorespaces 

\subsection{fix\_modify AtC output boundary\_integral command}
\label{\detokenize{atc_output_boundary_integral:fix-modify-atc-output-boundary-integral-command}}\label{\detokenize{atc_output_boundary_integral:index-0}}\label{\detokenize{atc_output_boundary_integral::doc}}

\subsubsection{Syntax}
\label{\detokenize{atc_output_boundary_integral:syntax}}
\begin{sphinxVerbatim}[commandchars=\\\{\}]
\PYG{k}{fix\PYGZus{}modify }\PYG{o}{\PYGZlt{}}\PYG{n}{AtC}\PYG{+w}{ }\PYG{n}{fixID}\PYG{o}{\PYGZgt{}}\PYG{+w}{ }\PYG{n}{output}\PYG{+w}{ }\PYG{n}{boundary\PYGZus{}integral}\PYG{+w}{ }\PYG{o}{\PYGZlt{}}\PYG{n}{fieldname}\PYG{o}{\PYGZgt{}}\PYG{+w}{ }\PYG{n}{faceset}\PYG{+w}{ }\PYG{n}{[name]}
\end{sphinxVerbatim}
\begin{itemize}
\item {} 
\sphinxAtStartPar
AtC fixID = ID of {\hyperref[\detokenize{fix_atc::doc}]{\sphinxcrossref{\DUrole{doc}{fix atc}}}} instance

\item {} 
\sphinxAtStartPar
output boundary\_integral = name of the AtC sub\sphinxhyphen{}command

\item {} 
\sphinxAtStartPar
fieldname = name of hardy field

\item {} 
\sphinxAtStartPar
faceset = required keyword

\item {} 
\sphinxAtStartPar
name= name of faceset

\end{itemize}


\subsubsection{Examples}
\label{\detokenize{atc_output_boundary_integral:examples}}
\begin{sphinxVerbatim}[commandchars=\\\{\}]
\PYG{k}{fix\PYGZus{}modify }\PYG{n+nv+nvIdentifier}{AtC}\PYG{+w}{ }\PYG{n}{output}\PYG{+w}{ }\PYG{n}{boundary\PYGZus{}integral}\PYG{+w}{ }\PYG{n}{stress}\PYG{+w}{ }\PYG{n}{faceset}\PYG{+w}{ }\PYG{n}{loop1}
\end{sphinxVerbatim}


\subsubsection{Description}
\label{\detokenize{atc_output_boundary_integral:description}}
\sphinxAtStartPar
Calculates a surface integral of the given field dotted with the outward
normal of the faces and puts output in the “GLOBALS” file.


\subsubsection{Restrictions}
\label{\detokenize{atc_output_boundary_integral:restrictions}}
\sphinxAtStartPar
Must be used with the hardy/field type of {\hyperref[\detokenize{fix_atc::doc}]{\sphinxcrossref{\DUrole{doc}{fix atc}}}}


\subsubsection{Related AtC commands}
\label{\detokenize{atc_output_boundary_integral:related-atc-commands}}\begin{itemize}
\item {} 
\sphinxAtStartPar
{\hyperref[\detokenize{fix_atc:atc-fix-modify}]{\sphinxcrossref{\DUrole{std,std-ref}{fix\_modify AtC command overview}}}}

\end{itemize}


\subsubsection{Default}
\label{\detokenize{atc_output_boundary_integral:default}}
\sphinxAtStartPar
None.

\sphinxstepscope

\index{fix\_modify AtC output contour\_integral@\spxentry{fix\_modify AtC output contour\_integral}}\ignorespaces 

\subsection{fix\_modify AtC output contour\_integral command}
\label{\detokenize{atc_output_contour_integral:fix-modify-atc-output-contour-integral-command}}\label{\detokenize{atc_output_contour_integral:index-0}}\label{\detokenize{atc_output_contour_integral::doc}}

\subsubsection{Syntax}
\label{\detokenize{atc_output_contour_integral:syntax}}
\begin{sphinxVerbatim}[commandchars=\\\{\}]
\PYG{k}{fix\PYGZus{}modify }\PYG{o}{\PYGZlt{}}\PYG{n}{AtC}\PYG{+w}{ }\PYG{n}{fixID}\PYG{o}{\PYGZgt{}}\PYG{+w}{ }\PYG{n}{output}\PYG{+w}{ }\PYG{n}{contour\PYGZus{}integral}\PYG{+w}{ }\PYG{o}{\PYGZlt{}}\PYG{n}{fieldname}\PYG{o}{\PYGZgt{}}\PYG{+w}{ }\PYG{n}{faceset}\PYG{+w}{ }\PYG{o}{\PYGZlt{}}\PYG{n}{name}\PYG{o}{\PYGZgt{}}\PYG{+w}{ }\PYG{n}{[axis}\PYG{+w}{ }\PYG{n}{[x}\PYG{o}{|}\PYG{n}{y}\PYG{o}{|}\PYG{n}{z]]}
\end{sphinxVerbatim}
\begin{itemize}
\item {} 
\sphinxAtStartPar
AtC fixID = ID of {\hyperref[\detokenize{fix_atc::doc}]{\sphinxcrossref{\DUrole{doc}{fix atc}}}} instance

\item {} 
\sphinxAtStartPar
output contour\_integral = name of the AtC sub\sphinxhyphen{}command

\item {} 
\sphinxAtStartPar
fieldname = name of hardy field

\item {} 
\sphinxAtStartPar
faceset = required keyword

\item {} 
\sphinxAtStartPar
name = name of faceset

\item {} 
\sphinxAtStartPar
\sphinxstyleemphasis{axis x} or \sphinxstyleemphasis{axis y} or \sphinxstyleemphasis{axis z} = (optional)

\end{itemize}


\subsubsection{Examples}
\label{\detokenize{atc_output_contour_integral:examples}}
\begin{sphinxVerbatim}[commandchars=\\\{\}]
\PYG{k}{fix\PYGZus{}modify }\PYG{n+nv+nvIdentifier}{AtC}\PYG{+w}{ }\PYG{n}{output}\PYG{+w}{ }\PYG{n}{contour\PYGZus{}integral}\PYG{+w}{ }\PYG{n}{stress}\PYG{+w}{ }\PYG{n}{faceset}\PYG{+w}{ }\PYG{n}{loop1}
\end{sphinxVerbatim}


\subsubsection{Description}
\label{\detokenize{atc_output_contour_integral:description}}
\sphinxAtStartPar
Calculates a surface integral of the given field dotted with the outward
normal of the faces and puts output in the “GLOBALS” file.


\subsubsection{Restrictions}
\label{\detokenize{atc_output_contour_integral:restrictions}}
\sphinxAtStartPar
Must be used with the hardy/field type of {\hyperref[\detokenize{fix_atc::doc}]{\sphinxcrossref{\DUrole{doc}{fix atc}}}}


\subsubsection{Related AtC commands}
\label{\detokenize{atc_output_contour_integral:related-atc-commands}}\begin{itemize}
\item {} 
\sphinxAtStartPar
{\hyperref[\detokenize{fix_atc:atc-fix-modify}]{\sphinxcrossref{\DUrole{std,std-ref}{fix\_modify AtC command overview}}}}

\end{itemize}


\subsubsection{Default}
\label{\detokenize{atc_output_contour_integral:default}}
\sphinxAtStartPar
None.

\sphinxstepscope

\index{fix\_modify AtC output nodeset@\spxentry{fix\_modify AtC output nodeset}}\ignorespaces 

\subsection{fix\_modify AtC output nodeset command}
\label{\detokenize{atc_output_nodeset:fix-modify-atc-output-nodeset-command}}\label{\detokenize{atc_output_nodeset:index-0}}\label{\detokenize{atc_output_nodeset::doc}}

\subsubsection{Syntax}
\label{\detokenize{atc_output_nodeset:syntax}}
\begin{sphinxVerbatim}[commandchars=\\\{\}]
\PYG{k}{fix\PYGZus{}modify }\PYG{o}{\PYGZlt{}}\PYG{n}{AtC}\PYG{+w}{ }\PYG{n}{fixID}\PYG{o}{\PYGZgt{}}\PYG{+w}{ }\PYG{n}{output}\PYG{+w}{ }\PYG{n}{nodeset}\PYG{+w}{ }\PYG{o}{\PYGZlt{}}\PYG{n}{nodeset\PYGZus{}name}\PYG{o}{\PYGZgt{}}\PYG{+w}{ }\PYG{o}{\PYGZlt{}}\PYG{n}{operation}\PYG{o}{\PYGZgt{}}
\end{sphinxVerbatim}
\begin{itemize}
\item {} 
\sphinxAtStartPar
AtC fixID = ID of {\hyperref[\detokenize{fix_atc::doc}]{\sphinxcrossref{\DUrole{doc}{fix atc}}}} instance

\item {} 
\sphinxAtStartPar
output nodeset = name of the AtC sub\sphinxhyphen{}command

\item {} 
\sphinxAtStartPar
nodeset\_name= name of nodeset to be operated on

\item {} 
\sphinxAtStartPar
operation = \sphinxstyleemphasis{sum}
\begin{itemize}
\item {} 
\sphinxAtStartPar
\sphinxstyleemphasis{sum} = creates nodal sum over nodes in specified nodeset

\end{itemize}

\end{itemize}


\subsubsection{Examples}
\label{\detokenize{atc_output_nodeset:examples}}
\begin{sphinxVerbatim}[commandchars=\\\{\}]
\PYG{k}{fix\PYGZus{}modify }\PYG{n+nv+nvIdentifier}{AtC}\PYG{+w}{ }\PYG{n}{output}\PYG{+w}{ }\PYG{n}{nodeset}\PYG{+w}{ }\PYG{n}{nset1}\PYG{+w}{ }\PYG{n}{sum}
\end{sphinxVerbatim}


\subsubsection{Description}
\label{\detokenize{atc_output_nodeset:description}}
\sphinxAtStartPar
Performs operation over the nodes belonging to specified nodeset and
outputs resulting variable values to GLOBALS file.


\subsubsection{Restrictions}
\label{\detokenize{atc_output_nodeset:restrictions}}
\sphinxAtStartPar
None.


\subsubsection{Related AtC commands}
\label{\detokenize{atc_output_nodeset:related-atc-commands}}\begin{itemize}
\item {} 
\sphinxAtStartPar
{\hyperref[\detokenize{fix_atc:atc-fix-modify}]{\sphinxcrossref{\DUrole{std,std-ref}{fix\_modify AtC command overview}}}}

\item {} 
\sphinxAtStartPar
{\hyperref[\detokenize{fix_atc::doc}]{\sphinxcrossref{\DUrole{doc}{fix atc command}}}}

\end{itemize}


\subsubsection{Default}
\label{\detokenize{atc_output_nodeset:default}}
\sphinxAtStartPar
None.

\sphinxstepscope

\index{fix\_modify AtC output volume\_integral@\spxentry{fix\_modify AtC output volume\_integral}}\ignorespaces 

\subsection{fix\_modify AtC output volume\_integral command}
\label{\detokenize{atc_output_volume_integral:fix-modify-atc-output-volume-integral-command}}\label{\detokenize{atc_output_volume_integral:index-0}}\label{\detokenize{atc_output_volume_integral::doc}}

\subsubsection{Syntax}
\label{\detokenize{atc_output_volume_integral:syntax}}
\begin{sphinxVerbatim}[commandchars=\\\{\}]
\PYG{k}{fix\PYGZus{}modify }\PYG{o}{\PYGZlt{}}\PYG{n}{AtC}\PYG{+w}{ }\PYG{n}{fixID}\PYG{o}{\PYGZgt{}}\PYG{+w}{ }\PYG{n}{output}\PYG{+w}{ }\PYG{n}{volume\PYGZus{}integral}\PYG{+w}{ }\PYG{o}{\PYGZlt{}}\PYG{n}{elementset\PYGZus{}name}\PYG{o}{\PYGZgt{}}\PYG{+w}{ }\PYG{o}{\PYGZlt{}}\PYG{n}{field}\PYG{o}{\PYGZgt{}}
\end{sphinxVerbatim}
\begin{itemize}
\item {} 
\sphinxAtStartPar
AtC fixID = ID of {\hyperref[\detokenize{fix_atc::doc}]{\sphinxcrossref{\DUrole{doc}{fix atc}}}} instance

\item {} 
\sphinxAtStartPar
output volume\_integral = name of the AtC sub\sphinxhyphen{}command

\item {} 
\sphinxAtStartPar
elementset\_name= name of elementset to be integrated over

\item {} 
\sphinxAtStartPar
fieldname = name of field to integrate

\end{itemize}


\subsubsection{Examples}
\label{\detokenize{atc_output_volume_integral:examples}}
\begin{sphinxVerbatim}[commandchars=\\\{\}]
\PYG{k}{fix\PYGZus{}modify }\PYG{n+nv+nvIdentifier}{AtC}\PYG{+w}{ }\PYG{n}{output}\PYG{+w}{ }\PYG{n}{volume\PYGZus{}integral}\PYG{+w}{ }\PYG{n}{eset1}\PYG{+w}{ }\PYG{n}{mass\PYGZus{}density}
\end{sphinxVerbatim}


\subsubsection{Description}
\label{\detokenize{atc_output_volume_integral:description}}
\sphinxAtStartPar
Performs volume integration of specified field over elementset and
outputs resulting variable values to GLOBALS file.


\subsubsection{Restrictions}
\label{\detokenize{atc_output_volume_integral:restrictions}}
\sphinxAtStartPar
None.


\subsubsection{Related AtC commands}
\label{\detokenize{atc_output_volume_integral:related-atc-commands}}\begin{itemize}
\item {} 
\sphinxAtStartPar
{\hyperref[\detokenize{fix_atc:atc-fix-modify}]{\sphinxcrossref{\DUrole{std,std-ref}{fix\_modify AtC command overview}}}}

\item {} 
\sphinxAtStartPar
{\hyperref[\detokenize{fix_atc::doc}]{\sphinxcrossref{\DUrole{doc}{fix atc command}}}}

\end{itemize}


\subsubsection{Default}
\label{\detokenize{atc_output_volume_integral:default}}
\sphinxAtStartPar
None.

\sphinxstepscope

\index{fix\_modify AtC pair\_interactions@\spxentry{fix\_modify AtC pair\_interactions}}\ignorespaces 

\subsection{fix\_modify AtC pair\_interactions command}
\label{\detokenize{atc_pair_interactions:fix-modify-atc-pair-interactions-command}}\label{\detokenize{atc_pair_interactions:index-0}}\label{\detokenize{atc_pair_interactions::doc}}

\subsection{fix\_modify AtC bond\_interactions command}
\label{\detokenize{atc_pair_interactions:fix-modify-atc-bond-interactions-command}}

\subsubsection{Syntax}
\label{\detokenize{atc_pair_interactions:syntax}}
\begin{sphinxVerbatim}[commandchars=\\\{\}]
\PYG{k}{fix\PYGZus{}modify }\PYG{o}{\PYGZlt{}}\PYG{n}{AtC}\PYG{+w}{ }\PYG{n}{fixID}\PYG{o}{\PYGZgt{}}\PYG{+w}{ }\PYG{n}{pair\PYGZus{}interactions}\PYG{+w}{ }\PYG{o}{\PYGZlt{}}\PYG{n}{on}\PYG{o}{|}\PYG{n}{off}\PYG{o}{\PYGZgt{}}
\PYG{k}{fix\PYGZus{}modify }\PYG{o}{\PYGZlt{}}\PYG{n}{AtC}\PYG{+w}{ }\PYG{n}{fixID}\PYG{o}{\PYGZgt{}}\PYG{+w}{ }\PYG{n}{bond\PYGZus{}interactions}\PYG{+w}{ }\PYG{o}{\PYGZlt{}}\PYG{n}{on}\PYG{o}{|}\PYG{n}{off}\PYG{o}{\PYGZgt{}}
\end{sphinxVerbatim}
\begin{itemize}
\item {} 
\sphinxAtStartPar
AtC fixID = ID of {\hyperref[\detokenize{fix_atc::doc}]{\sphinxcrossref{\DUrole{doc}{fix atc}}}} instance

\item {} 
\sphinxAtStartPar
\sphinxstyleemphasis{pair\_interactions} or \sphinxstyleemphasis{bond\_interactions} = name of the AtC sub\sphinxhyphen{}command

\item {} 
\sphinxAtStartPar
\sphinxstyleemphasis{on} or \sphinxstyleemphasis{off} = activate or deactivate

\end{itemize}


\subsubsection{Examples}
\label{\detokenize{atc_pair_interactions:examples}}
\begin{sphinxVerbatim}[commandchars=\\\{\}]
\PYG{k}{fix\PYGZus{}modify }\PYG{n+nv+nvIdentifier}{AtC}\PYG{+w}{ }\PYG{n}{pair\PYGZus{}interactions}\PYG{+w}{ }\PYG{n}{off}
\PYG{k}{fix\PYGZus{}modify }\PYG{n+nv+nvIdentifier}{AtC}\PYG{+w}{ }\PYG{n}{bond\PYGZus{}interactions}\PYG{+w}{ }\PYG{n}{on}
\end{sphinxVerbatim}


\subsubsection{Description}
\label{\detokenize{atc_pair_interactions:description}}
\sphinxAtStartPar
Include bonds and/or pairs in stress and heat flux computations.


\subsubsection{Restrictions}
\label{\detokenize{atc_pair_interactions:restrictions}}
\sphinxAtStartPar
None.


\subsubsection{Related AtC commands}
\label{\detokenize{atc_pair_interactions:related-atc-commands}}\begin{itemize}
\item {} 
\sphinxAtStartPar
{\hyperref[\detokenize{fix_atc:atc-fix-modify}]{\sphinxcrossref{\DUrole{std,std-ref}{fix\_modify AtC command overview}}}}

\end{itemize}


\subsubsection{Default}
\label{\detokenize{atc_pair_interactions:default}}
\sphinxAtStartPar
\sphinxstyleemphasis{pair\_interactions}: on, \sphinxstyleemphasis{bond\_interactions}: off

\sphinxstepscope

\index{fix\_modify AtC poisson\_solver@\spxentry{fix\_modify AtC poisson\_solver}}\ignorespaces 

\subsection{fix\_modify AtC poisson\_solver command}
\label{\detokenize{atc_poisson_solver:fix-modify-atc-poisson-solver-command}}\label{\detokenize{atc_poisson_solver:index-0}}\label{\detokenize{atc_poisson_solver::doc}}

\subsubsection{Syntax}
\label{\detokenize{atc_poisson_solver:syntax}}
\begin{sphinxVerbatim}[commandchars=\\\{\}]
\PYG{k}{fix\PYGZus{}modify }\PYG{o}{\PYGZlt{}}\PYG{n}{AtC}\PYG{+w}{ }\PYG{n}{fixID}\PYG{o}{\PYGZgt{}}\PYG{+w}{ }\PYG{n}{poisson\PYGZus{}solver}\PYG{+w}{ }\PYG{n}{mesh}\PYG{+w}{ }\PYG{n}{create}\PYG{+w}{ }\PYG{o}{\PYGZlt{}}\PYG{n}{nx}\PYG{o}{\PYGZgt{}}\PYG{+w}{ }\PYG{o}{\PYGZlt{}}\PYG{n}{ny}\PYG{o}{\PYGZgt{}}\PYG{+w}{ }\PYG{o}{\PYGZlt{}}\PYG{n}{nz}\PYG{o}{\PYGZgt{}}\PYG{+w}{ }\PYG{o}{\PYGZlt{}}\PYG{n}{region}\PYG{o}{\PYGZhy{}}\PYG{n}{ID}\PYG{o}{\PYGZgt{}}\PYG{+w}{ }\PYG{o}{\PYGZlt{}}\PYG{n}{f}\PYG{o}{|}\PYG{n}{p}\PYG{o}{\PYGZgt{}}\PYG{+w}{ }\PYG{o}{\PYGZlt{}}\PYG{n}{f}\PYG{o}{|}\PYG{n}{p}\PYG{o}{\PYGZgt{}}\PYG{+w}{ }\PYG{o}{\PYGZlt{}}\PYG{n}{f}\PYG{o}{|}\PYG{n}{p}\PYG{o}{\PYGZgt{}}
\end{sphinxVerbatim}
\begin{itemize}
\item {} 
\sphinxAtStartPar
AtC fixID = ID of {\hyperref[\detokenize{fix_atc::doc}]{\sphinxcrossref{\DUrole{doc}{fix atc}}}} instance

\item {} 
\sphinxAtStartPar
poisson\_solver = name of the AtC sub\sphinxhyphen{}command

\item {} 
\sphinxAtStartPar
\sphinxstyleemphasis{nx} \sphinxstyleemphasis{ny} \sphinxstyleemphasis{nz} = number of elements in x, y, and z

\item {} 
\sphinxAtStartPar
region\sphinxhyphen{}id = id of region to be meshed

\item {} 
\sphinxAtStartPar
\sphinxstyleemphasis{f} or \sphinxstyleemphasis{p} = periodicity flags for x, y, and z

\end{itemize}


\subsubsection{Examples}
\label{\detokenize{atc_poisson_solver:examples}}
\begin{sphinxVerbatim}[commandchars=\\\{\}]
\PYG{k}{fix\PYGZus{}modify }\PYG{n+nv+nvIdentifier}{AtC}\PYG{+w}{ }\PYG{n}{poisson\PYGZus{}solver}\PYG{+w}{ }\PYG{n}{mesh}\PYG{+w}{ }\PYG{n}{create}\PYG{+w}{ }\PYG{l+m}{10}\PYG{+w}{ }\PYG{l+m}{1}\PYG{+w}{ }\PYG{l+m}{1}\PYG{+w}{ }\PYG{n}{feRegion}\PYG{+w}{ }\PYG{n}{p}\PYG{+w}{ }\PYG{n}{p}\PYG{+w}{ }\PYG{n}{p}
\end{sphinxVerbatim}


\subsubsection{Description}
\label{\detokenize{atc_poisson_solver:description}}
\sphinxAtStartPar
Creates a uniform mesh in a rectangular region.


\subsubsection{Restrictions}
\label{\detokenize{atc_poisson_solver:restrictions}}
\sphinxAtStartPar
Creates only uniform rectangular grids in rectangular regions.


\subsubsection{Related AtC commands}
\label{\detokenize{atc_poisson_solver:related-atc-commands}}\begin{itemize}
\item {} 
\sphinxAtStartPar
{\hyperref[\detokenize{fix_atc:atc-fix-modify}]{\sphinxcrossref{\DUrole{std,std-ref}{fix\_modify AtC command overview}}}}

\end{itemize}


\subsubsection{Default}
\label{\detokenize{atc_poisson_solver:default}}
\sphinxAtStartPar
None.

\sphinxstepscope

\index{fix\_modify AtC read\_restart@\spxentry{fix\_modify AtC read\_restart}}\ignorespaces 

\subsection{fix\_modify AtC read\_restart command}
\label{\detokenize{atc_read_restart:fix-modify-atc-read-restart-command}}\label{\detokenize{atc_read_restart:index-0}}\label{\detokenize{atc_read_restart::doc}}

\subsubsection{Syntax}
\label{\detokenize{atc_read_restart:syntax}}
\begin{sphinxVerbatim}[commandchars=\\\{\}]
\PYG{k}{fix\PYGZus{}modify }\PYG{o}{\PYGZlt{}}\PYG{n}{AtC}\PYG{+w}{ }\PYG{n}{fixID}\PYG{o}{\PYGZgt{}}\PYG{+w}{ }\PYG{n}{read\PYGZus{}restart}\PYG{+w}{ }\PYG{o}{\PYGZlt{}}\PYG{n}{file\PYGZus{}name}\PYG{o}{\PYGZgt{}}
\end{sphinxVerbatim}
\begin{itemize}
\item {} 
\sphinxAtStartPar
AtC fixID = ID of {\hyperref[\detokenize{fix_atc::doc}]{\sphinxcrossref{\DUrole{doc}{fix atc}}}} instance

\item {} 
\sphinxAtStartPar
read\_restart = name of the AtC sub\sphinxhyphen{}command

\item {} 
\sphinxAtStartPar
file\_name = name of AtC restart file

\end{itemize}


\subsubsection{Examples}
\label{\detokenize{atc_read_restart:examples}}
\begin{sphinxVerbatim}[commandchars=\\\{\}]
\PYG{k}{fix\PYGZus{}modify }\PYG{n+nv+nvIdentifier}{AtC}\PYG{+w}{ }\PYG{n}{read\PYGZus{}restart}\PYG{+w}{ }\PYG{n}{restart.mydata.AtC}
\end{sphinxVerbatim}


\subsubsection{Description}
\label{\detokenize{atc_read_restart:description}}
\sphinxAtStartPar
Reads the current state of the AtC fields from a named text\sphinxhyphen{}based restart file.


\subsubsection{Restrictions}
\label{\detokenize{atc_read_restart:restrictions}}
\sphinxAtStartPar
The restart file only contains fields and their time derivatives.  The
reference positions of the atoms and the commands that initialize the
fix are not saved e.g. an identical mesh containing the same atoms will
have to be recreated.


\subsubsection{Related AtC commands}
\label{\detokenize{atc_read_restart:related-atc-commands}}\begin{itemize}
\item {} 
\sphinxAtStartPar
{\hyperref[\detokenize{fix_atc:atc-fix-modify}]{\sphinxcrossref{\DUrole{std,std-ref}{fix\_modify AtC command overview}}}}

\item {} 
\sphinxAtStartPar
{\hyperref[\detokenize{atc_write_restart::doc}]{\sphinxcrossref{\DUrole{doc}{fix\_modify AtC write\_restart}}}}

\end{itemize}


\subsubsection{Default}
\label{\detokenize{atc_read_restart:default}}
\sphinxAtStartPar
None.

\sphinxstepscope

\index{fix\_modify AtC remove\_molecule@\spxentry{fix\_modify AtC remove\_molecule}}\ignorespaces 

\subsection{fix\_modify AtC remove\_molecule command}
\label{\detokenize{atc_remove_molecule:fix-modify-atc-remove-molecule-command}}\label{\detokenize{atc_remove_molecule:index-0}}\label{\detokenize{atc_remove_molecule::doc}}

\subsubsection{Syntax}
\label{\detokenize{atc_remove_molecule:syntax}}
\begin{sphinxVerbatim}[commandchars=\\\{\}]
\PYG{k}{fix\PYGZus{}modify }\PYG{o}{\PYGZlt{}}\PYG{n}{AtC}\PYG{+w}{ }\PYG{n}{fixID}\PYG{o}{\PYGZgt{}}\PYG{+w}{ }\PYG{n}{remove\PYGZus{}molecule}\PYG{+w}{ }\PYG{o}{\PYGZlt{}}\PYG{n}{tag}\PYG{o}{\PYGZgt{}}
\end{sphinxVerbatim}
\begin{itemize}
\item {} 
\sphinxAtStartPar
AtC fixID = ID of {\hyperref[\detokenize{fix_atc::doc}]{\sphinxcrossref{\DUrole{doc}{fix atc}}}} instance

\item {} 
\sphinxAtStartPar
remove\_molecule = name of the AtC sub\sphinxhyphen{}command

\item {} 
\sphinxAtStartPar
tag = tag for tracking a molecule

\end{itemize}


\subsubsection{Examples}
\label{\detokenize{atc_remove_molecule:examples}}
\begin{sphinxVerbatim}[commandchars=\\\{\}]
\PYG{k}{fix\PYGZus{}modify }\PYG{n+nv+nvIdentifier}{AtC}\PYG{+w}{ }\PYG{n}{remove\PYGZus{}molecule}\PYG{+w}{ }\PYG{n}{water}
\end{sphinxVerbatim}


\subsubsection{Description}
\label{\detokenize{atc_remove_molecule:description}}
\sphinxAtStartPar
Removes tag designated for tracking a specified set of molecules.


\subsubsection{Restrictions}
\label{\detokenize{atc_remove_molecule:restrictions}}
\sphinxAtStartPar
None.


\subsubsection{Related AtC commands}
\label{\detokenize{atc_remove_molecule:related-atc-commands}}\begin{itemize}
\item {} 
\sphinxAtStartPar
{\hyperref[\detokenize{fix_atc:atc-fix-modify}]{\sphinxcrossref{\DUrole{std,std-ref}{fix\_modify AtC command overview}}}}

\item {} 
\sphinxAtStartPar
{\hyperref[\detokenize{atc_add_species::doc}]{\sphinxcrossref{\DUrole{doc}{fix\_modify AtC add\_species}}}}

\item {} 
\sphinxAtStartPar
{\hyperref[\detokenize{atc_add_molecule::doc}]{\sphinxcrossref{\DUrole{doc}{fix\_modify AtC add\_molecule}}}}

\item {} 
\sphinxAtStartPar
{\hyperref[\detokenize{atc_remove_species::doc}]{\sphinxcrossref{\DUrole{doc}{fix\_modify AtC remove\_species}}}}

\end{itemize}


\subsubsection{Default}
\label{\detokenize{atc_remove_molecule:default}}
\sphinxAtStartPar
None.

\sphinxstepscope

\index{fix\_modify AtC remove\_source@\spxentry{fix\_modify AtC remove\_source}}\ignorespaces 

\subsection{fix\_modify AtC remove\_source command}
\label{\detokenize{atc_remove_source:fix-modify-atc-remove-source-command}}\label{\detokenize{atc_remove_source:index-0}}\label{\detokenize{atc_remove_source::doc}}

\subsubsection{Syntax}
\label{\detokenize{atc_remove_source:syntax}}
\begin{sphinxVerbatim}[commandchars=\\\{\}]
\PYG{k}{fix\PYGZus{}modify }\PYG{o}{\PYGZlt{}}\PYG{n}{AtC}\PYG{+w}{ }\PYG{n}{fixID}\PYG{o}{\PYGZgt{}}\PYG{+w}{ }\PYG{n}{remove\PYGZus{}source}\PYG{+w}{ }\PYG{o}{\PYGZlt{}}\PYG{n}{field}\PYG{o}{\PYGZgt{}}\PYG{+w}{ }\PYG{o}{\PYGZlt{}}\PYG{n}{element\PYGZus{}set}\PYG{o}{\PYGZgt{}}
\end{sphinxVerbatim}
\begin{itemize}
\item {} 
\sphinxAtStartPar
AtC fixID = ID of {\hyperref[\detokenize{fix_atc::doc}]{\sphinxcrossref{\DUrole{doc}{fix atc}}}} instance

\item {} 
\sphinxAtStartPar
remove\_source = name of the AtC sub\sphinxhyphen{}command

\item {} 
\sphinxAtStartPar
field = field kind name valid for type of physics: \sphinxstyleemphasis{temperature} or \sphinxstyleemphasis{electron\_temperature}

\item {} 
\sphinxAtStartPar
element\_set = name of set of elements

\end{itemize}


\subsubsection{Examples}
\label{\detokenize{atc_remove_source:examples}}
\begin{sphinxVerbatim}[commandchars=\\\{\}]
\PYG{k}{fix\PYGZus{}modify }\PYG{n+nv+nvIdentifier}{AtC}\PYG{+w}{ }\PYG{n}{remove\PYGZus{}source}\PYG{+w}{ }\PYG{n}{temperature}\PYG{+w}{ }\PYG{n}{groupNAME}
\end{sphinxVerbatim}


\subsubsection{Description}
\label{\detokenize{atc_remove_source:description}}
\sphinxAtStartPar
Remove a domain source.


\subsubsection{Restrictions}
\label{\detokenize{atc_remove_source:restrictions}}
\sphinxAtStartPar
The keyword \sphinxstyleemphasis{all} is reserved and thus not available as element\_set name.


\subsubsection{Related AtC commands}
\label{\detokenize{atc_remove_source:related-atc-commands}}\begin{itemize}
\item {} 
\sphinxAtStartPar
{\hyperref[\detokenize{fix_atc:atc-fix-modify}]{\sphinxcrossref{\DUrole{std,std-ref}{fix\_modify AtC command overview}}}}

\item {} 
\sphinxAtStartPar
{\hyperref[\detokenize{atc_source::doc}]{\sphinxcrossref{\DUrole{doc}{fix\_modify AtC source}}}}

\end{itemize}


\subsubsection{Default}
\label{\detokenize{atc_remove_source:default}}
\sphinxAtStartPar
None.

\sphinxstepscope

\index{fix\_modify AtC remove\_species@\spxentry{fix\_modify AtC remove\_species}}\ignorespaces 

\subsection{fix\_modify AtC remove\_species command}
\label{\detokenize{atc_remove_species:fix-modify-atc-remove-species-command}}\label{\detokenize{atc_remove_species:index-0}}\label{\detokenize{atc_remove_species::doc}}

\subsubsection{Syntax}
\label{\detokenize{atc_remove_species:syntax}}
\begin{sphinxVerbatim}[commandchars=\\\{\}]
\PYG{k}{fix\PYGZus{}modify }\PYG{o}{\PYGZlt{}}\PYG{n}{AtC}\PYG{+w}{ }\PYG{n}{fixID}\PYG{o}{\PYGZgt{}}\PYG{+w}{ }\PYG{n}{remove\PYGZus{}species}\PYG{+w}{ }\PYG{o}{\PYGZlt{}}\PYG{n}{tag}\PYG{o}{\PYGZgt{}}
\end{sphinxVerbatim}
\begin{itemize}
\item {} 
\sphinxAtStartPar
AtC fixID = ID of {\hyperref[\detokenize{fix_atc::doc}]{\sphinxcrossref{\DUrole{doc}{fix atc}}}} instance

\item {} 
\sphinxAtStartPar
remove\_species = name of the AtC sub\sphinxhyphen{}command

\item {} 
\sphinxAtStartPar
tag = tag for tracking a species

\end{itemize}


\subsubsection{Examples}
\label{\detokenize{atc_remove_species:examples}}
\begin{sphinxVerbatim}[commandchars=\\\{\}]
\PYG{k}{fix\PYGZus{}modify }\PYG{n+nv+nvIdentifier}{AtC}\PYG{+w}{ }\PYG{n}{remove\PYGZus{}species}\PYG{+w}{ }\PYG{n}{gold}
\end{sphinxVerbatim}


\subsubsection{Description}
\label{\detokenize{atc_remove_species:description}}
\sphinxAtStartPar
Removes tag designated for tracking a specified species.


\subsubsection{Restrictions}
\label{\detokenize{atc_remove_species:restrictions}}
\sphinxAtStartPar
None.


\subsubsection{Related AtC commands}
\label{\detokenize{atc_remove_species:related-atc-commands}}\begin{itemize}
\item {} 
\sphinxAtStartPar
{\hyperref[\detokenize{fix_atc:atc-fix-modify}]{\sphinxcrossref{\DUrole{std,std-ref}{fix\_modify AtC command overview}}}}

\item {} 
\sphinxAtStartPar
{\hyperref[\detokenize{atc_add_species::doc}]{\sphinxcrossref{\DUrole{doc}{fix\_modify AtC add\_species}}}}

\item {} 
\sphinxAtStartPar
{\hyperref[\detokenize{atc_add_molecule::doc}]{\sphinxcrossref{\DUrole{doc}{fix\_modify AtC add\_molecule}}}}

\item {} 
\sphinxAtStartPar
{\hyperref[\detokenize{atc_remove_molecule::doc}]{\sphinxcrossref{\DUrole{doc}{fix\_modify AtC remove\_molecule}}}}

\end{itemize}


\subsubsection{Default}
\label{\detokenize{atc_remove_species:default}}
\sphinxAtStartPar
None.

\sphinxstepscope

\index{fix\_modify AtC reset\_atomic\_reference\_positions@\spxentry{fix\_modify AtC reset\_atomic\_reference\_positions}}\ignorespaces 

\subsection{fix\_modify AtC reset\_atomic\_reference\_positions command}
\label{\detokenize{atc_reset_atomic_reference:fix-modify-atc-reset-atomic-reference-positions-command}}\label{\detokenize{atc_reset_atomic_reference:index-0}}\label{\detokenize{atc_reset_atomic_reference::doc}}

\subsubsection{Syntax}
\label{\detokenize{atc_reset_atomic_reference:syntax}}
\begin{sphinxVerbatim}[commandchars=\\\{\}]
\PYG{k}{fix\PYGZus{}modify }\PYG{o}{\PYGZlt{}}\PYG{n}{AtC}\PYG{+w}{ }\PYG{n}{fixID}\PYG{o}{\PYGZgt{}}\PYG{+w}{ }\PYG{n}{reset\PYGZus{}atomic\PYGZus{}reference\PYGZus{}positions}
\end{sphinxVerbatim}
\begin{itemize}
\item {} 
\sphinxAtStartPar
AtC fixID = ID of {\hyperref[\detokenize{fix_atc::doc}]{\sphinxcrossref{\DUrole{doc}{fix atc}}}} instance

\item {} 
\sphinxAtStartPar
reset\_atomic\_reference\_positions = name of the AtC sub\sphinxhyphen{}command

\end{itemize}


\subsubsection{Examples}
\label{\detokenize{atc_reset_atomic_reference:examples}}
\begin{sphinxVerbatim}[commandchars=\\\{\}]
\PYG{k}{fix\PYGZus{}modify }\PYG{n+nv+nvIdentifier}{AtC}\PYG{+w}{ }\PYG{n}{reset\PYGZus{}atomic\PYGZus{}reference\PYGZus{}positions}
\end{sphinxVerbatim}


\subsubsection{Description}
\label{\detokenize{atc_reset_atomic_reference:description}}
\sphinxAtStartPar
Resets the atomic positions ATC uses to perform point to field
operations. In can be used to use perfect lattice sites in ATC but a
thermalized or deformed lattice in LAMMPS.


\subsubsection{Restrictions}
\label{\detokenize{atc_reset_atomic_reference:restrictions}}
\sphinxAtStartPar
None.


\subsubsection{Related AtC commands}
\label{\detokenize{atc_reset_atomic_reference:related-atc-commands}}\begin{itemize}
\item {} 
\sphinxAtStartPar
{\hyperref[\detokenize{fix_atc:atc-fix-modify}]{\sphinxcrossref{\DUrole{std,std-ref}{fix\_modify AtC command overview}}}}

\end{itemize}


\subsubsection{Default}
\label{\detokenize{atc_reset_atomic_reference:default}}
\sphinxAtStartPar
None

\sphinxstepscope

\index{fix\_modify AtC reset\_time@\spxentry{fix\_modify AtC reset\_time}}\ignorespaces 

\subsection{fix\_modify AtC reset\_time command}
\label{\detokenize{atc_reset_time:fix-modify-atc-reset-time-command}}\label{\detokenize{atc_reset_time:index-0}}\label{\detokenize{atc_reset_time::doc}}

\subsubsection{Syntax}
\label{\detokenize{atc_reset_time:syntax}}
\begin{sphinxVerbatim}[commandchars=\\\{\}]
\PYG{k}{fix\PYGZus{}modify }\PYG{o}{\PYGZlt{}}\PYG{n}{AtC}\PYG{+w}{ }\PYG{n}{fixID}\PYG{o}{\PYGZgt{}}\PYG{+w}{ }\PYG{n}{reset\PYGZus{}time}\PYG{+w}{ }\PYG{o}{\PYGZlt{}}\PYG{n}{value}\PYG{o}{\PYGZgt{}}
\end{sphinxVerbatim}
\begin{itemize}
\item {} 
\sphinxAtStartPar
AtC fixID = ID of {\hyperref[\detokenize{fix_atc::doc}]{\sphinxcrossref{\DUrole{doc}{fix atc}}}} instance

\item {} 
\sphinxAtStartPar
reset\_time = name of the AtC sub\sphinxhyphen{}command

\item {} 
\sphinxAtStartPar
value = new time value

\end{itemize}


\subsubsection{Examples}
\label{\detokenize{atc_reset_time:examples}}
\begin{sphinxVerbatim}[commandchars=\\\{\}]
\PYG{k}{fix\PYGZus{}modify }\PYG{n+nv+nvIdentifier}{AtC}\PYG{+w}{ }\PYG{n}{reset\PYGZus{}time}\PYG{+w}{ }\PYG{l+m}{0.0}
\end{sphinxVerbatim}


\subsubsection{Description}
\label{\detokenize{atc_reset_time:description}}
\sphinxAtStartPar
Resets the simulation time counter.


\subsubsection{Restrictions}
\label{\detokenize{atc_reset_time:restrictions}}
\sphinxAtStartPar
None.


\subsubsection{Related AtC commands}
\label{\detokenize{atc_reset_time:related-atc-commands}}\begin{itemize}
\item {} 
\sphinxAtStartPar
{\hyperref[\detokenize{fix_atc:atc-fix-modify}]{\sphinxcrossref{\DUrole{std,std-ref}{fix\_modify AtC command overview}}}}

\end{itemize}


\subsubsection{Default}
\label{\detokenize{atc_reset_time:default}}
\sphinxAtStartPar
None

\sphinxstepscope

\index{fix\_modify AtC sample\_frequency@\spxentry{fix\_modify AtC sample\_frequency}}\ignorespaces 

\subsection{fix\_modify AtC sample\_frequency command}
\label{\detokenize{atc_sample_frequency:fix-modify-atc-sample-frequency-command}}\label{\detokenize{atc_sample_frequency:index-0}}\label{\detokenize{atc_sample_frequency::doc}}

\subsubsection{Syntax}
\label{\detokenize{atc_sample_frequency:syntax}}
\begin{sphinxVerbatim}[commandchars=\\\{\}]
\PYG{k}{fix\PYGZus{}modify }\PYG{o}{\PYGZlt{}}\PYG{n}{AtC}\PYG{+w}{ }\PYG{n}{fixID}\PYG{o}{\PYGZgt{}}\PYG{+w}{ }\PYG{n}{sample\PYGZus{}frequency}\PYG{+w}{ }\PYG{o}{\PYGZlt{}}\PYG{n}{freq}\PYG{o}{\PYGZgt{}}
\end{sphinxVerbatim}
\begin{itemize}
\item {} 
\sphinxAtStartPar
AtC fixID = ID of {\hyperref[\detokenize{fix_atc::doc}]{\sphinxcrossref{\DUrole{doc}{fix atc}}}} instance

\item {} 
\sphinxAtStartPar
sample\_frequency = name of the AtC sub\sphinxhyphen{}command

\item {} 
\sphinxAtStartPar
freq = frequency to sample fields in number of steps

\end{itemize}


\subsubsection{Examples}
\label{\detokenize{atc_sample_frequency:examples}}
\begin{sphinxVerbatim}[commandchars=\\\{\}]
\PYG{k}{fix\PYGZus{}modify }\PYG{n+nv+nvIdentifier}{AtC}\PYG{+w}{ }\PYG{n}{sample\PYGZus{}frequency}\PYG{+w}{ }\PYG{l+m}{10}
\end{sphinxVerbatim}


\subsubsection{Description}
\label{\detokenize{atc_sample_frequency:description}}
\sphinxAtStartPar
Specifies a frequency at which fields are computed for the case where
time filters are being applied.


\subsubsection{Restrictions}
\label{\detokenize{atc_sample_frequency:restrictions}}
\sphinxAtStartPar
Must be used with {\hyperref[\detokenize{fix_atc::doc}]{\sphinxcrossref{\DUrole{doc}{fix atc hardy}}}} and is only relevant when time filters are being used.


\subsubsection{Related AtC commands}
\label{\detokenize{atc_sample_frequency:related-atc-commands}}\begin{itemize}
\item {} 
\sphinxAtStartPar
{\hyperref[\detokenize{fix_atc:atc-fix-modify}]{\sphinxcrossref{\DUrole{std,std-ref}{fix\_modify AtC command overview}}}}

\end{itemize}


\subsubsection{Default}
\label{\detokenize{atc_sample_frequency:default}}
\sphinxAtStartPar
None.

\sphinxstepscope

\index{fix\_modify AtC set reference\_potential\_energy@\spxentry{fix\_modify AtC set reference\_potential\_energy}}\ignorespaces 

\subsection{fix\_modify AtC set reference\_potential\_energy command}
\label{\detokenize{atc_set_reference_pe:fix-modify-atc-set-reference-potential-energy-command}}\label{\detokenize{atc_set_reference_pe:index-0}}\label{\detokenize{atc_set_reference_pe::doc}}

\subsubsection{Syntax}
\label{\detokenize{atc_set_reference_pe:syntax}}
\begin{sphinxVerbatim}[commandchars=\\\{\}]
\PYG{k}{fix\PYGZus{}modify }\PYG{o}{\PYGZlt{}}\PYG{n}{AtC}\PYG{+w}{ }\PYG{n}{fixID}\PYG{o}{\PYGZgt{}}\PYG{+w}{ }\PYG{n}{set}\PYG{+w}{ }\PYG{n}{reference\PYGZus{}potential\PYGZus{}energy}\PYG{+w}{ }\PYG{n}{[}\PYG{o}{\PYGZlt{}}\PYG{n}{value}\PYG{o}{|}\PYG{n}{filename}\PYG{o}{\PYGZgt{}}\PYG{n}{]}
\end{sphinxVerbatim}
\begin{itemize}
\item {} 
\sphinxAtStartPar
AtC fixID = ID of {\hyperref[\detokenize{fix_atc::doc}]{\sphinxcrossref{\DUrole{doc}{fix atc}}}} instance

\item {} 
\sphinxAtStartPar
set reference\_potential\_energy = name of the AtC sub\sphinxhyphen{}command

\item {} 
\sphinxAtStartPar
value = optional user specified zero point for PE in native LAMMPS energy units

\item {} 
\sphinxAtStartPar
filename = optional user specified string for file of nodal PE values to be read\sphinxhyphen{}in

\end{itemize}


\subsubsection{Examples}
\label{\detokenize{atc_set_reference_pe:examples}}
\begin{sphinxVerbatim}[commandchars=\\\{\}]
\PYG{k}{fix\PYGZus{}modify }\PYG{n+nv+nvIdentifier}{AtC}\PYG{+w}{ }\PYG{n}{set}\PYG{+w}{ }\PYG{n}{reference\PYGZus{}potential\PYGZus{}energy}
\PYG{k}{fix\PYGZus{}modify }\PYG{n+nv+nvIdentifier}{AtC}\PYG{+w}{ }\PYG{n}{set}\PYG{+w}{ }\PYG{n}{reference\PYGZus{}potential\PYGZus{}energy}\PYG{+w}{ }\PYG{o}{\PYGZhy{}}\PYG{l+m}{0.05}
\PYG{k}{fix\PYGZus{}modify }\PYG{n+nv+nvIdentifier}{AtC}\PYG{+w}{ }\PYG{n}{set}\PYG{+w}{ }\PYG{n}{reference\PYGZus{}potential\PYGZus{}energy}\PYG{+w}{ }\PYG{n}{myPEvalues}
\end{sphinxVerbatim}


\subsubsection{Description}
\label{\detokenize{atc_set_reference_pe:description}}
\sphinxAtStartPar
Used to set various quantities for the post\sphinxhyphen{}processing algorithms. It
sets the zero point for the potential energy density using the value
provided for all nodes, or from the current configuration of the lattice
if no value is provided, or values provided within the specified
filename.


\subsubsection{Restrictions}
\label{\detokenize{atc_set_reference_pe:restrictions}}
\sphinxAtStartPar
Must be used with {\hyperref[\detokenize{fix_atc::doc}]{\sphinxcrossref{\DUrole{doc}{fix atc hardy}}}} or {\hyperref[\detokenize{fix_atc::doc}]{\sphinxcrossref{\DUrole{doc}{fix atc field}}}}.


\subsubsection{Related AtC commands}
\label{\detokenize{atc_set_reference_pe:related-atc-commands}}\begin{itemize}
\item {} 
\sphinxAtStartPar
{\hyperref[\detokenize{fix_atc:atc-fix-modify}]{\sphinxcrossref{\DUrole{std,std-ref}{fix\_modify AtC command overview}}}}

\end{itemize}


\subsubsection{Default}
\label{\detokenize{atc_set_reference_pe:default}}
\sphinxAtStartPar
Defaults to the LAMMPS zero point i.e. isolated atoms.

\sphinxstepscope

\index{fix\_modify AtC source@\spxentry{fix\_modify AtC source}}\ignorespaces 

\subsection{fix\_modify AtC source command}
\label{\detokenize{atc_source:fix-modify-atc-source-command}}\label{\detokenize{atc_source:index-0}}\label{\detokenize{atc_source::doc}}

\subsubsection{Syntax}
\label{\detokenize{atc_source:syntax}}
\begin{sphinxVerbatim}[commandchars=\\\{\}]
\PYG{k}{fix\PYGZus{}modify }\PYG{o}{\PYGZlt{}}\PYG{n}{AtC}\PYG{+w}{ }\PYG{n}{fixID}\PYG{o}{\PYGZgt{}}\PYG{+w}{ }\PYG{n}{source}\PYG{+w}{ }\PYG{o}{\PYGZlt{}}\PYG{n}{field}\PYG{o}{\PYGZgt{}}\PYG{+w}{ }\PYG{o}{\PYGZlt{}}\PYG{n}{element\PYGZus{}set}\PYG{o}{\PYGZgt{}}\PYG{+w}{ }\PYG{o}{\PYGZlt{}}\PYG{n}{value}\PYG{o}{|}\PYG{n}{function}\PYG{o}{\PYGZgt{}}
\end{sphinxVerbatim}
\begin{itemize}
\item {} 
\sphinxAtStartPar
AtC fixID = ID of {\hyperref[\detokenize{fix_atc::doc}]{\sphinxcrossref{\DUrole{doc}{fix atc}}}} instance

\item {} 
\sphinxAtStartPar
source = name of the AtC sub\sphinxhyphen{}command

\item {} 
\sphinxAtStartPar
field = field kind name valid for type of physics: \sphinxstyleemphasis{temperature} or \sphinxstyleemphasis{electron\_temperature}

\item {} 
\sphinxAtStartPar
element\_set = name of set of elements

\item {} 
\sphinxAtStartPar
\sphinxstyleemphasis{value} or \sphinxstyleemphasis{function} = value or name of function followed by its parameters

\end{itemize}


\subsubsection{Examples}
\label{\detokenize{atc_source:examples}}
\begin{sphinxVerbatim}[commandchars=\\\{\}]
\PYG{k}{fix\PYGZus{}modify }\PYG{n+nv+nvIdentifier}{AtC}\PYG{+w}{ }\PYG{n}{source}\PYG{+w}{ }\PYG{n}{temperature}\PYG{+w}{ }\PYG{n}{middle}\PYG{+w}{ }\PYG{n}{temporal\PYGZus{}ramp}\PYG{+w}{ }\PYG{l+m}{10.0}\PYG{+w}{ }\PYG{l+m}{0.0}
\end{sphinxVerbatim}


\subsubsection{Description}
\label{\detokenize{atc_source:description}}
\sphinxAtStartPar
Add domain sources to the mesh. The units are consistent with LAMMPS’s
units for mass, length and time and are defined by the PDE being solved,
e.g. for thermal transfer the balance equation is for energy and source
is energy per time.


\subsubsection{Restrictions}
\label{\detokenize{atc_source:restrictions}}
\sphinxAtStartPar
The keyword \sphinxstyleemphasis{all} is reserved and thus not available as element\_set name.


\subsubsection{Related AtC commands}
\label{\detokenize{atc_source:related-atc-commands}}\begin{itemize}
\item {} 
\sphinxAtStartPar
{\hyperref[\detokenize{fix_atc:atc-fix-modify}]{\sphinxcrossref{\DUrole{std,std-ref}{fix\_modify AtC command overview}}}}

\item {} 
\sphinxAtStartPar
{\hyperref[\detokenize{atc_remove_source::doc}]{\sphinxcrossref{\DUrole{doc}{fix\_modify AtC remove\_source}}}}

\end{itemize}


\subsubsection{Default}
\label{\detokenize{atc_source:default}}
\sphinxAtStartPar
None.

\sphinxstepscope

\index{fix\_modify AtC source\_integration@\spxentry{fix\_modify AtC source\_integration}}\ignorespaces 

\subsection{fix\_modify AtC source\_integration command}
\label{\detokenize{atc_source_integration:fix-modify-atc-source-integration-command}}\label{\detokenize{atc_source_integration:index-0}}\label{\detokenize{atc_source_integration::doc}}

\subsubsection{Syntax}
\label{\detokenize{atc_source_integration:syntax}}
\begin{sphinxVerbatim}[commandchars=\\\{\}]
\PYG{k}{fix\PYGZus{}modify }\PYG{o}{\PYGZlt{}}\PYG{n}{AtC}\PYG{+w}{ }\PYG{n}{fixID}\PYG{o}{\PYGZgt{}}\PYG{+w}{ }\PYG{n}{source\PYGZus{}integration}\PYG{+w}{ }\PYG{o}{\PYGZlt{}}\PYG{n}{fe}\PYG{o}{|}\PYG{n}{atom}\PYG{o}{\PYGZgt{}}
\end{sphinxVerbatim}
\begin{itemize}
\item {} 
\sphinxAtStartPar
AtC fixID = ID of {\hyperref[\detokenize{fix_atc::doc}]{\sphinxcrossref{\DUrole{doc}{fix atc}}}} instance

\item {} 
\sphinxAtStartPar
source\_integration = name of the AtC sub\sphinxhyphen{}command

\item {} 
\sphinxAtStartPar
\sphinxstyleemphasis{fe} or \sphinxstyleemphasis{atom} = (undocumented)

\end{itemize}


\subsubsection{Examples}
\label{\detokenize{atc_source_integration:examples}}
\begin{sphinxVerbatim}[commandchars=\\\{\}]
\PYG{k}{fix\PYGZus{}modify }\PYG{n+nv+nvIdentifier}{AtC}\PYG{+w}{ }\PYG{n}{source\PYGZus{}integration}\PYG{+w}{ }\PYG{n}{atom}
\end{sphinxVerbatim}


\subsubsection{Description}
\label{\detokenize{atc_source_integration:description}}
\sphinxAtStartPar
(undocumented)


\subsubsection{Restrictions}
\label{\detokenize{atc_source_integration:restrictions}}
\sphinxAtStartPar
None.


\subsubsection{Related AtC commands}
\label{\detokenize{atc_source_integration:related-atc-commands}}\begin{itemize}
\item {} 
\sphinxAtStartPar
{\hyperref[\detokenize{fix_atc:atc-fix-modify}]{\sphinxcrossref{\DUrole{std,std-ref}{fix\_modify AtC command overview}}}}

\end{itemize}


\subsubsection{Default}
\label{\detokenize{atc_source_integration:default}}
\sphinxAtStartPar
Default is \sphinxstyleemphasis{fe}

\sphinxstepscope

\index{fix\_modify AtC temperature\_definition@\spxentry{fix\_modify AtC temperature\_definition}}\ignorespaces 

\subsection{fix\_modify AtC temperature\_definition command}
\label{\detokenize{atc_temperature_definition:fix-modify-atc-temperature-definition-command}}\label{\detokenize{atc_temperature_definition:index-0}}\label{\detokenize{atc_temperature_definition::doc}}

\subsubsection{Syntax}
\label{\detokenize{atc_temperature_definition:syntax}}
\begin{sphinxVerbatim}[commandchars=\\\{\}]
\PYG{k}{fix\PYGZus{}modify }\PYG{o}{\PYGZlt{}}\PYG{n}{AtC}\PYG{+w}{ }\PYG{n}{fixID}\PYG{o}{\PYGZgt{}}\PYG{+w}{ }\PYG{n}{temperature\PYGZus{}definition}\PYG{+w}{ }\PYG{o}{\PYGZlt{}}\PYG{n}{kinetic}\PYG{o}{|}\PYG{n}{total}\PYG{o}{\PYGZgt{}}
\end{sphinxVerbatim}
\begin{itemize}
\item {} 
\sphinxAtStartPar
AtC fixID = ID of {\hyperref[\detokenize{fix_atc::doc}]{\sphinxcrossref{\DUrole{doc}{fix atc}}}} instance

\item {} 
\sphinxAtStartPar
temperature\_definition = name of the AtC sub\sphinxhyphen{}command

\item {} 
\sphinxAtStartPar
\sphinxstyleemphasis{kinetic} or \sphinxstyleemphasis{total} = (undocumented)

\end{itemize}


\subsubsection{Examples}
\label{\detokenize{atc_temperature_definition:examples}}
\begin{sphinxVerbatim}[commandchars=\\\{\}]
\PYG{k}{fix\PYGZus{}modify }\PYG{n+nv+nvIdentifier}{AtC}\PYG{+w}{ }\PYG{n}{temperature\PYGZus{}definition}\PYG{+w}{ }\PYG{n}{kinetic}
\end{sphinxVerbatim}


\subsubsection{Description}
\label{\detokenize{atc_temperature_definition:description}}
\sphinxAtStartPar
Change the definition for the atomic temperature used to create the
finite element temperature.  The kinetic option is based only on the
kinetic energy of the atoms while the total option uses the total energy
(kinetic + potential) of an atom.


\subsubsection{Restrictions}
\label{\detokenize{atc_temperature_definition:restrictions}}
\sphinxAtStartPar
This command is only valid when using thermal coupling.  Also, while not
a formal restriction, the user should ensure that associating a
potential energy with each atom makes physical sense for the total
option to be meaningful.


\subsubsection{Related AtC commands}
\label{\detokenize{atc_temperature_definition:related-atc-commands}}\begin{itemize}
\item {} 
\sphinxAtStartPar
{\hyperref[\detokenize{fix_atc:atc-fix-modify}]{\sphinxcrossref{\DUrole{std,std-ref}{fix\_modify AtC command overview}}}}

\end{itemize}


\subsubsection{Default}
\label{\detokenize{atc_temperature_definition:default}}
\sphinxAtStartPar
\sphinxstyleemphasis{kinetic}

\sphinxstepscope

\index{fix\_modify AtC filter@\spxentry{fix\_modify AtC filter}}\ignorespaces 

\subsection{fix\_modify AtC filter command}
\label{\detokenize{atc_time_filter:fix-modify-atc-filter-command}}\label{\detokenize{atc_time_filter:index-0}}\label{\detokenize{atc_time_filter::doc}}

\subsubsection{Syntax}
\label{\detokenize{atc_time_filter:syntax}}
\begin{sphinxVerbatim}[commandchars=\\\{\}]
\PYG{k}{fix\PYGZus{}modify }\PYG{o}{\PYGZlt{}}\PYG{n}{AtC}\PYG{+w}{ }\PYG{n}{fixID}\PYG{o}{\PYGZgt{}}\PYG{+w}{ }\PYG{n}{filter}\PYG{+w}{ }\PYG{o}{\PYGZlt{}}\PYG{n}{on}\PYG{o}{|}\PYG{n}{off}\PYG{o}{|}\PYG{n}{equilibrate}\PYG{o}{\PYGZgt{}}
\end{sphinxVerbatim}
\begin{itemize}
\item {} 
\sphinxAtStartPar
AtC fixID = ID of {\hyperref[\detokenize{fix_atc::doc}]{\sphinxcrossref{\DUrole{doc}{fix atc}}}} instance

\item {} 
\sphinxAtStartPar
filter = name of the AtC sub\sphinxhyphen{}command

\item {} 
\sphinxAtStartPar
\sphinxstyleemphasis{on} or \sphinxstyleemphasis{off} or \sphinxstyleemphasis{equilibrate} = Select state of filter

\end{itemize}


\subsubsection{Examples}
\label{\detokenize{atc_time_filter:examples}}
\begin{sphinxVerbatim}[commandchars=\\\{\}]
\PYG{k}{fix\PYGZus{}modify }\PYG{n+nv+nvIdentifier}{AtC}\PYG{+w}{ }\PYG{n}{filter}\PYG{+w}{ }\PYG{n}{on}
\end{sphinxVerbatim}


\subsubsection{Description}
\label{\detokenize{atc_time_filter:description}}
\sphinxAtStartPar
Filters the MD dynamics to construct a more appropriate continuous
field. Equilibrating first filters the time derivatives without changing
the dynamics to provide a better initial condition to the filtered
dynamics.


\subsubsection{Restrictions}
\label{\detokenize{atc_time_filter:restrictions}}
\sphinxAtStartPar
Only for use with these specific transfers: thermal, two\_temperature


\subsubsection{Related AtC commands}
\label{\detokenize{atc_time_filter:related-atc-commands}}\begin{itemize}
\item {} 
\sphinxAtStartPar
{\hyperref[\detokenize{fix_atc:atc-fix-modify}]{\sphinxcrossref{\DUrole{std,std-ref}{fix\_modify AtC command overview}}}}

\item {} 
\sphinxAtStartPar
{\hyperref[\detokenize{atc_filter_scale::doc}]{\sphinxcrossref{\DUrole{doc}{fix\_modify AtC filter scale}}}}

\item {} 
\sphinxAtStartPar
{\hyperref[\detokenize{atc_equilibrium_start::doc}]{\sphinxcrossref{\DUrole{doc}{fix\_modify AtC equilibrium\_start}}}}

\end{itemize}


\subsubsection{Default}
\label{\detokenize{atc_time_filter:default}}
\sphinxAtStartPar
off

\sphinxstepscope

\index{fix\_modify AtC time\_integration@\spxentry{fix\_modify AtC time\_integration}}\ignorespaces 

\subsection{fix\_modify AtC time\_integration command}
\label{\detokenize{atc_time_integration:fix-modify-atc-time-integration-command}}\label{\detokenize{atc_time_integration:index-0}}\label{\detokenize{atc_time_integration::doc}}

\subsubsection{Syntax}
\label{\detokenize{atc_time_integration:syntax}}
\begin{sphinxVerbatim}[commandchars=\\\{\}]
\PYG{k}{fix\PYGZus{}modify }\PYG{o}{\PYGZlt{}}\PYG{n}{AtC}\PYG{+w}{ }\PYG{n}{fixID}\PYG{o}{\PYGZgt{}}\PYG{+w}{ }\PYG{n}{time\PYGZus{}integration}\PYG{+w}{ }\PYG{o}{\PYGZlt{}}\PYG{n}{descriptor}\PYG{o}{\PYGZgt{}}
\end{sphinxVerbatim}
\begin{itemize}
\item {} 
\sphinxAtStartPar
AtC fixID = ID of {\hyperref[\detokenize{fix_atc::doc}]{\sphinxcrossref{\DUrole{doc}{fix atc}}}} instance

\item {} 
\sphinxAtStartPar
time\_integration = name of the AtC sub\sphinxhyphen{}command

\item {} 
\sphinxAtStartPar
descriptor =  \sphinxstyleemphasis{gear} or \sphinxstyleemphasis{fractional\_step} or \sphinxstyleemphasis{verlet}

\end{itemize}


\subsubsection{Examples}
\label{\detokenize{atc_time_integration:examples}}
\begin{sphinxVerbatim}[commandchars=\\\{\}]
\PYG{k}{fix\PYGZus{}modify }\PYG{n+nv+nvIdentifier}{AtC}\PYG{+w}{ }\PYG{n}{time\PYGZus{}integration}\PYG{+w}{ }\PYG{n}{fractional\PYGZus{}step}
\end{sphinxVerbatim}


\subsubsection{Description}
\label{\detokenize{atc_time_integration:description}}
\sphinxAtStartPar
Command to select the thermal or momentum time integration.


\bigskip\hrule\bigskip


\sphinxAtStartPar
Options for thermal time integration:
\begin{description}
\sphinxlineitem{\sphinxstyleemphasis{gear}}
\sphinxAtStartPar
atomic velocity update with second order Verlet, nodal temperature update
with third or fourth order Gear, thermostats based on controlling power

\sphinxlineitem{\sphinxstyleemphasis{fractional\_step}}
\sphinxAtStartPar
atomic velocity update with second order Verlet, mixed nodal temperature
update, 3/4 Gear for continuum and 2 Verlet for atomic contributions,
thermostats based on controlling discrete energy changes

\end{description}


\bigskip\hrule\bigskip


\sphinxAtStartPar
Options for momentum time integration:
\begin{description}
\sphinxlineitem{\sphinxstyleemphasis{verlet}}
\sphinxAtStartPar
atomic velocity update with second order Verlet, nodal temperature update
with second order Verlet, kinetostats based on controlling force

\sphinxlineitem{\sphinxstyleemphasis{fractional\_step}}
\sphinxAtStartPar
atomic velocity update with second order Verlet, mixed nodal momentum
update, second order Verlet for continuum and exact second order Verlet for
atomic contributions, kinetostats based on controlling discrete
momentum changes

\sphinxlineitem{\sphinxstyleemphasis{gear}}
\sphinxAtStartPar
atomic velocity update with second order Verlet, nodal temperature update
with third or fourth order Gear, kinetostats based on controlling power.

\end{description}


\bigskip\hrule\bigskip



\subsubsection{Restrictions}
\label{\detokenize{atc_time_integration:restrictions}}
\sphinxAtStartPar
None.


\subsubsection{Related AtC commands}
\label{\detokenize{atc_time_integration:related-atc-commands}}\begin{itemize}
\item {} 
\sphinxAtStartPar
{\hyperref[\detokenize{fix_atc:atc-fix-modify}]{\sphinxcrossref{\DUrole{std,std-ref}{fix\_modify AtC command overview}}}}

\end{itemize}


\subsubsection{Default}
\label{\detokenize{atc_time_integration:default}}
\sphinxAtStartPar
None.

\sphinxstepscope

\index{fix\_modify AtC track\_displacement@\spxentry{fix\_modify AtC track\_displacement}}\ignorespaces 

\subsection{fix\_modify AtC track\_displacement command}
\label{\detokenize{atc_track_displacement:fix-modify-atc-track-displacement-command}}\label{\detokenize{atc_track_displacement:index-0}}\label{\detokenize{atc_track_displacement::doc}}

\subsubsection{Syntax}
\label{\detokenize{atc_track_displacement:syntax}}
\begin{sphinxVerbatim}[commandchars=\\\{\}]
\PYG{k}{fix\PYGZus{}modify }\PYG{o}{\PYGZlt{}}\PYG{n}{AtC}\PYG{+w}{ }\PYG{n}{fixID}\PYG{o}{\PYGZgt{}}\PYG{+w}{ }\PYG{n}{track\PYGZus{}displacement}\PYG{+w}{ }\PYG{o}{\PYGZlt{}}\PYG{n}{on}\PYG{o}{|}\PYG{n}{off}\PYG{o}{\PYGZgt{}}
\end{sphinxVerbatim}
\begin{itemize}
\item {} 
\sphinxAtStartPar
AtC fixID = ID of {\hyperref[\detokenize{fix_atc::doc}]{\sphinxcrossref{\DUrole{doc}{fix atc}}}} instance

\item {} 
\sphinxAtStartPar
track\_displacement = name of the AtC sub\sphinxhyphen{}command

\item {} 
\sphinxAtStartPar
\sphinxstyleemphasis{on} or \sphinxstyleemphasis{off} = (undocumented)

\end{itemize}


\subsubsection{Examples}
\label{\detokenize{atc_track_displacement:examples}}
\begin{sphinxVerbatim}[commandchars=\\\{\}]
\PYG{k}{fix\PYGZus{}modify }\PYG{n+nv+nvIdentifier}{AtC}\PYG{+w}{ }\PYG{n}{track\PYGZus{}displacement}\PYG{+w}{ }\PYG{n}{on}
\end{sphinxVerbatim}


\subsubsection{Description}
\label{\detokenize{atc_track_displacement:description}}
\sphinxAtStartPar
Determines whether displacement is tracked or not. For solids problems
this is a useful quantity, but for fluids it is not relevant.


\subsubsection{Restrictions}
\label{\detokenize{atc_track_displacement:restrictions}}
\sphinxAtStartPar
Some constitutive models require the displacement field.


\subsubsection{Related AtC commands}
\label{\detokenize{atc_track_displacement:related-atc-commands}}\begin{itemize}
\item {} 
\sphinxAtStartPar
{\hyperref[\detokenize{fix_atc:atc-fix-modify}]{\sphinxcrossref{\DUrole{std,std-ref}{fix\_modify AtC command overview}}}}

\end{itemize}


\subsubsection{Default}
\label{\detokenize{atc_track_displacement:default}}
\sphinxAtStartPar
\sphinxstyleemphasis{on}

\sphinxstepscope

\index{fix\_modify AtC unfix@\spxentry{fix\_modify AtC unfix}}\ignorespaces 

\subsection{fix\_modify AtC unfix command}
\label{\detokenize{atc_unfix:fix-modify-atc-unfix-command}}\label{\detokenize{atc_unfix:index-0}}\label{\detokenize{atc_unfix::doc}}

\subsubsection{Syntax}
\label{\detokenize{atc_unfix:syntax}}
\begin{sphinxVerbatim}[commandchars=\\\{\}]
\PYG{k}{fix\PYGZus{}modify }\PYG{o}{\PYGZlt{}}\PYG{n}{AtC}\PYG{+w}{ }\PYG{n}{fixID}\PYG{o}{\PYGZgt{}}\PYG{+w}{ }\PYG{k}{unfix }\PYG{o}{\PYGZlt{}}\PYG{n}{field}\PYG{o}{\PYGZgt{}}\PYG{+w}{ }\PYG{o}{\PYGZlt{}}\PYG{n}{nodeset}\PYG{o}{\PYGZgt{}}
\end{sphinxVerbatim}
\begin{itemize}
\item {} 
\sphinxAtStartPar
AtC fixID = ID of {\hyperref[\detokenize{fix_atc::doc}]{\sphinxcrossref{\DUrole{doc}{fix atc}}}} instance

\item {} 
\sphinxAtStartPar
unfix = name of the AtC sub\sphinxhyphen{}command

\item {} 
\sphinxAtStartPar
field = field kind name valid for type of physics: \sphinxstyleemphasis{temperature} or \sphinxstyleemphasis{electron\_temperature}

\item {} 
\sphinxAtStartPar
nodeset = name of set of nodes to apply boundary condition

\end{itemize}


\subsubsection{Examples}
\label{\detokenize{atc_unfix:examples}}
\begin{sphinxVerbatim}[commandchars=\\\{\}]
\PYG{k}{fix\PYGZus{}modify }\PYG{n+nv+nvIdentifier}{AtC}\PYG{+w}{ }\PYG{k}{unfix }\PYG{n+nv+nvIdentifier}{temperature}\PYG{+w}{ }\PYG{n}{groupNAME}
\end{sphinxVerbatim}


\subsubsection{Description}
\label{\detokenize{atc_unfix:description}}
\sphinxAtStartPar
Removes constraint on field values for specified nodes.


\subsubsection{Restrictions}
\label{\detokenize{atc_unfix:restrictions}}
\sphinxAtStartPar
The keyword \sphinxstyleemphasis{all} is reserved and thus not available as nodeset name.


\subsubsection{Related AtC commands}
\label{\detokenize{atc_unfix:related-atc-commands}}\begin{itemize}
\item {} 
\sphinxAtStartPar
{\hyperref[\detokenize{fix_atc:atc-fix-modify}]{\sphinxcrossref{\DUrole{std,std-ref}{fix\_modify AtC command overview}}}}

\item {} 
\sphinxAtStartPar
{\hyperref[\detokenize{atc_fix::doc}]{\sphinxcrossref{\DUrole{doc}{fix\_modify AtC fix}}}}

\end{itemize}


\subsubsection{Default}
\label{\detokenize{atc_unfix:default}}
\sphinxAtStartPar
None.

\sphinxstepscope

\index{fix\_modify AtC unfix\_flux@\spxentry{fix\_modify AtC unfix\_flux}}\ignorespaces 

\subsection{fix\_modify AtC unfix\_flux command}
\label{\detokenize{atc_unfix_flux:fix-modify-atc-unfix-flux-command}}\label{\detokenize{atc_unfix_flux:index-0}}\label{\detokenize{atc_unfix_flux::doc}}

\subsubsection{Syntax}
\label{\detokenize{atc_unfix_flux:syntax}}
\begin{sphinxVerbatim}[commandchars=\\\{\}]
\PYG{k}{fix\PYGZus{}modify }\PYG{o}{\PYGZlt{}}\PYG{n}{AtC}\PYG{+w}{ }\PYG{n}{fixID}\PYG{o}{\PYGZgt{}}\PYG{+w}{ }\PYG{n}{unfix\PYGZus{}flux}\PYG{+w}{ }\PYG{o}{\PYGZlt{}}\PYG{n}{field}\PYG{o}{\PYGZgt{}}\PYG{+w}{ }\PYG{o}{\PYGZlt{}}\PYG{n}{face\PYGZus{}set}\PYG{o}{\PYGZgt{}}\PYG{+w}{ }\PYG{o}{\PYGZlt{}}\PYG{n}{value}\PYG{o}{|}\PYG{n}{function}\PYG{o}{\PYGZgt{}}
\end{sphinxVerbatim}
\begin{itemize}
\item {} 
\sphinxAtStartPar
AtC fixID = ID of {\hyperref[\detokenize{fix_atc::doc}]{\sphinxcrossref{\DUrole{doc}{fix atc}}}} instance

\item {} 
\sphinxAtStartPar
unfix\_flux = name of the AtC sub\sphinxhyphen{}command

\item {} 
\sphinxAtStartPar
field = field kind name valid for type of physics: \sphinxstyleemphasis{temperature} or \sphinxstyleemphasis{electron\_temperature}

\item {} 
\sphinxAtStartPar
face\_set = name of set of element faces

\end{itemize}


\subsubsection{Examples}
\label{\detokenize{atc_unfix_flux:examples}}
\begin{sphinxVerbatim}[commandchars=\\\{\}]
\PYG{k}{fix\PYGZus{}modify }\PYG{n+nv+nvIdentifier}{AtC}\PYG{+w}{ }\PYG{n}{unfix\PYGZus{}flux}\PYG{+w}{ }\PYG{n}{temperature}\PYG{+w}{ }\PYG{n}{faceSet}
\end{sphinxVerbatim}


\subsubsection{Description}
\label{\detokenize{atc_unfix_flux:description}}
\sphinxAtStartPar
Command for removing prescribed normal fluxes e.g. heat\_flux, stress.


\subsubsection{Restrictions}
\label{\detokenize{atc_unfix_flux:restrictions}}
\sphinxAtStartPar
None.


\subsubsection{Related AtC commands}
\label{\detokenize{atc_unfix_flux:related-atc-commands}}\begin{itemize}
\item {} 
\sphinxAtStartPar
{\hyperref[\detokenize{fix_atc:atc-fix-modify}]{\sphinxcrossref{\DUrole{std,std-ref}{fix\_modify AtC command overview}}}}

\item {} 
\sphinxAtStartPar
{\hyperref[\detokenize{atc_fix_flux::doc}]{\sphinxcrossref{\DUrole{doc}{fix\_modify AtC fix\_flux}}}}

\end{itemize}


\subsubsection{Default}
\label{\detokenize{atc_unfix_flux:default}}
\sphinxAtStartPar
None.

\sphinxstepscope

\index{fix\_modify AtC write\_atom\_weights@\spxentry{fix\_modify AtC write\_atom\_weights}}\ignorespaces 

\subsection{fix\_modify AtC write\_atom\_weights command}
\label{\detokenize{atc_write_atom_weights:fix-modify-atc-write-atom-weights-command}}\label{\detokenize{atc_write_atom_weights:index-0}}\label{\detokenize{atc_write_atom_weights::doc}}

\subsubsection{Syntax}
\label{\detokenize{atc_write_atom_weights:syntax}}
\begin{sphinxVerbatim}[commandchars=\\\{\}]
\PYG{k}{fix\PYGZus{}modify }\PYG{o}{\PYGZlt{}}\PYG{n}{AtC}\PYG{+w}{ }\PYG{n}{fixID}\PYG{o}{\PYGZgt{}}\PYG{+w}{ }\PYG{n}{write\PYGZus{}atom\PYGZus{}weights}\PYG{+w}{ }\PYG{o}{\PYGZlt{}}\PYG{n}{filename}\PYG{o}{\PYGZgt{}}\PYG{+w}{ }\PYG{o}{\PYGZlt{}}\PYG{n}{frequency}\PYG{o}{\PYGZgt{}}
\end{sphinxVerbatim}
\begin{itemize}
\item {} 
\sphinxAtStartPar
AtC fixID = ID of {\hyperref[\detokenize{fix_atc::doc}]{\sphinxcrossref{\DUrole{doc}{fix atc}}}} instance

\item {} 
\sphinxAtStartPar
write\_atom\_weights = name of the AtC sub\sphinxhyphen{}command

\item {} 
\sphinxAtStartPar
filename = name of file that atomic weights are written to

\item {} 
\sphinxAtStartPar
frequency = how often writes will occur

\end{itemize}


\subsubsection{Examples}
\label{\detokenize{atc_write_atom_weights:examples}}
\begin{sphinxVerbatim}[commandchars=\\\{\}]
\PYG{k}{fix\PYGZus{}modify }\PYG{n+nv+nvIdentifier}{AtC}\PYG{+w}{ }\PYG{n}{write\PYGZus{}atom\PYGZus{}weights}\PYG{+w}{ }\PYG{n}{atm\PYGZus{}wt\PYGZus{}file.txt}\PYG{+w}{ }\PYG{l+m}{10}
\end{sphinxVerbatim}


\subsubsection{Description}
\label{\detokenize{atc_write_atom_weights:description}}
\sphinxAtStartPar
Command for writing the values of atomic weights to a specified file.


\subsubsection{Restrictions}
\label{\detokenize{atc_write_atom_weights:restrictions}}
\sphinxAtStartPar
None.


\subsubsection{Related AtC commands}
\label{\detokenize{atc_write_atom_weights:related-atc-commands}}\begin{itemize}
\item {} 
\sphinxAtStartPar
{\hyperref[\detokenize{fix_atc:atc-fix-modify}]{\sphinxcrossref{\DUrole{std,std-ref}{fix\_modify AtC command overview}}}}

\end{itemize}


\subsubsection{Default}
\label{\detokenize{atc_write_atom_weights:default}}
\sphinxAtStartPar
None

\sphinxstepscope

\index{fix\_modify AtC write\_restart@\spxentry{fix\_modify AtC write\_restart}}\ignorespaces 

\subsection{fix\_modify AtC write\_restart command}
\label{\detokenize{atc_write_restart:fix-modify-atc-write-restart-command}}\label{\detokenize{atc_write_restart:index-0}}\label{\detokenize{atc_write_restart::doc}}

\subsubsection{Syntax}
\label{\detokenize{atc_write_restart:syntax}}
\begin{sphinxVerbatim}[commandchars=\\\{\}]
\PYG{k}{fix\PYGZus{}modify }\PYG{o}{\PYGZlt{}}\PYG{n}{AtC}\PYG{+w}{ }\PYG{n}{fixID}\PYG{o}{\PYGZgt{}}\PYG{+w}{ }\PYG{k}{write\PYGZus{}restart }\PYG{o}{\PYGZlt{}}\PYG{n}{file\PYGZus{}name}\PYG{o}{\PYGZgt{}}
\end{sphinxVerbatim}
\begin{itemize}
\item {} 
\sphinxAtStartPar
AtC fixID = ID of {\hyperref[\detokenize{fix_atc::doc}]{\sphinxcrossref{\DUrole{doc}{fix atc}}}} instance

\item {} 
\sphinxAtStartPar
write\_restart = name of the AtC sub\sphinxhyphen{}command

\item {} 
\sphinxAtStartPar
file\_name = name of AtC restart file

\end{itemize}


\subsubsection{Examples}
\label{\detokenize{atc_write_restart:examples}}
\begin{sphinxVerbatim}[commandchars=\\\{\}]
\PYG{k}{fix\PYGZus{}modify }\PYG{n+nv+nvIdentifier}{AtC}\PYG{+w}{ }\PYG{k}{write\PYGZus{}restart }\PYG{n+nv+nvIdentifier}{restart.mydata.AtC}
\end{sphinxVerbatim}


\subsubsection{Description}
\label{\detokenize{atc_write_restart:description}}
\sphinxAtStartPar
Dumps the current state of the fields to a named text\sphinxhyphen{}based restart
file.  This done when the command is invoked and not repeated, unlike
the otherwise similar LAMMPS command.


\subsubsection{Restrictions}
\label{\detokenize{atc_write_restart:restrictions}}
\sphinxAtStartPar
None.


\subsubsection{Related AtC commands}
\label{\detokenize{atc_write_restart:related-atc-commands}}\begin{itemize}
\item {} 
\sphinxAtStartPar
{\hyperref[\detokenize{fix_atc:atc-fix-modify}]{\sphinxcrossref{\DUrole{std,std-ref}{fix\_modify AtC command overview}}}}

\item {} 
\sphinxAtStartPar
{\hyperref[\detokenize{atc_read_restart::doc}]{\sphinxcrossref{\DUrole{doc}{fix\_modify AtC read\_restart}}}}

\end{itemize}


\subsubsection{Default}
\label{\detokenize{atc_write_restart:default}}
\sphinxAtStartPar
None.

\sphinxstepscope

\index{fix mol/swap@\spxentry{fix mol/swap}}\ignorespaces 

\section{fix mol/swap command}
\label{\detokenize{fix_mol_swap:fix-mol-swap-command}}\label{\detokenize{fix_mol_swap:index-0}}\label{\detokenize{fix_mol_swap::doc}}

\subsection{Syntax}
\label{\detokenize{fix_mol_swap:syntax}}
\begin{sphinxVerbatim}[commandchars=\\\{\}]
\PYG{k}{fix }\PYG{n+nv+nvIdentifier}{ID}\PYG{+w}{ }\PYG{n+nv+nvIdentifier}{group\PYGZhy{}ID}\PYG{+w}{ }\PYG{n}{mol}\PYG{o}{/}\PYG{n}{swap}\PYG{+w}{ }\PYG{n}{N}\PYG{+w}{ }\PYG{n}{X}\PYG{+w}{ }\PYG{n}{itype}\PYG{+w}{ }\PYG{n}{jtype}\PYG{+w}{ }\PYG{n}{seed}\PYG{+w}{ }\PYG{n}{T}\PYG{+w}{ }\PYG{n}{keyword}\PYG{+w}{ }\PYG{n}{value}\PYG{+w}{ }\PYG{n}{...}
\end{sphinxVerbatim}
\begin{itemize}
\item {} 
\sphinxAtStartPar
ID, group\sphinxhyphen{}ID are documented in {\hyperref[\detokenize{fix::doc}]{\sphinxcrossref{\DUrole{doc}{fix}}}} command

\item {} 
\sphinxAtStartPar
atom/swap = style name of this fix command

\item {} 
\sphinxAtStartPar
N = invoke this fix every N steps

\item {} 
\sphinxAtStartPar
X = number of swaps to attempt every N steps

\item {} 
\sphinxAtStartPar
itype,jtype = two atom types (1\sphinxhyphen{}Ntypes or type label) to swap with each other

\item {} 
\sphinxAtStartPar
seed = random \# seed (positive integer)

\item {} 
\sphinxAtStartPar
T = scaling temperature of the MC swaps (temperature units)

\item {} 
\sphinxAtStartPar
zero or more keyword/value pairs may be appended to args

\item {} 
\sphinxAtStartPar
keyword = \sphinxstyleemphasis{ke}
\begin{sphinxalltt}
\sphinxstyleemphasis{ke} value = \sphinxstyleemphasis{no} or \sphinxstyleemphasis{yes}
  \sphinxstyleemphasis{no} = no conservation of kinetic energy after atom swaps
  \sphinxstyleemphasis{yes} = kinetic energy is conserved after atom swaps
\end{sphinxalltt}

\end{itemize}


\subsection{Examples}
\label{\detokenize{fix_mol_swap:examples}}
\begin{sphinxVerbatim}[commandchars=\\\{\}]
\PYG{k}{fix }\PYG{n+nv+nvIdentifier}{2}\PYG{+w}{ }\PYG{n+nv+nvIdentifier}{all}\PYG{+w}{ }\PYG{n}{mol}\PYG{o}{/}\PYG{n}{swap}\PYG{+w}{ }\PYG{l+m}{100}\PYG{+w}{ }\PYG{l+m}{1}\PYG{+w}{ }\PYG{l+m}{2}\PYG{+w}{ }\PYG{l+m}{3}\PYG{+w}{ }\PYG{l+m}{29494}\PYG{+w}{ }\PYG{l+m}{300.0}\PYG{+w}{ }\PYG{n}{ke}\PYG{+w}{ }\PYG{n}{no}
\PYG{k}{fix }\PYG{n+nv+nvIdentifier}{mySwap}\PYG{+w}{ }\PYG{n+nv+nvIdentifier}{fluid}\PYG{+w}{ }\PYG{n}{mol}\PYG{o}{/}\PYG{n}{swap}\PYG{+w}{ }\PYG{l+m}{500}\PYG{+w}{ }\PYG{l+m}{10}\PYG{+w}{ }\PYG{l+m}{1}\PYG{+w}{ }\PYG{l+m}{2}\PYG{+w}{ }\PYG{l+m}{482798}\PYG{+w}{ }\PYG{l+m}{1.0}

\PYG{k}{labelmap}\PYG{+w}{ }\PYG{n}{atom}\PYG{+w}{ }\PYG{l+m}{1}\PYG{+w}{ }\PYG{n}{A}\PYG{+w}{ }\PYG{l+m}{2}\PYG{+w}{ }\PYG{n}{B}
\PYG{k}{fix }\PYG{n+nv+nvIdentifier}{mySwap}\PYG{+w}{ }\PYG{n+nv+nvIdentifier}{fluid}\PYG{+w}{ }\PYG{n}{mol}\PYG{o}{/}\PYG{n}{swap}\PYG{+w}{ }\PYG{l+m}{500}\PYG{+w}{ }\PYG{l+m}{10}\PYG{+w}{ }\PYG{n}{A}\PYG{+w}{ }\PYG{n}{B}\PYG{+w}{ }\PYG{l+m}{482798}\PYG{+w}{ }\PYG{l+m}{1.0}
\end{sphinxVerbatim}


\subsection{Description}
\label{\detokenize{fix_mol_swap:description}}
\sphinxAtStartPar
This fix performs Monte Carlo swaps of two specified atom types within
a randomly selected molecule.  Two possible use cases are as follows.

\sphinxAtStartPar
First, consider a mixture of some molecules with atoms of itype and
other molecules with atoms of jtype.  The fix will select a random
molecule and attempt to swap all the itype atoms to jtype for the
first kind of molecule, or all the jtype atoms to itype for the second
kind.  Because the swap will only take place if it is energetically
favorable, the fix can be used to determine the miscibility of 2
different kinds of molecules much more quickly than just dynamics
would do it.

\sphinxAtStartPar
Second, consider diblock co\sphinxhyphen{}polymers with two types of monomers itype
and jtype.  The fix will select a random molecule and attempt to do a
itype \textless{}\textendash{}\textgreater{} jtype swap of all those monomers within the molecule.  Thus
the fix can be used to find the energetically favorable fractions of
two flavors of diblock co\sphinxhyphen{}polymers.

\sphinxAtStartPar
Intra\sphinxhyphen{}molecular swaps of atom types are attempted every N timesteps.  On
that timestep, X swaps are attempted.  For each attempt a single
molecule ID is randomly selected.  The range of possible molecule IDs
from loID to hiID is pre\sphinxhyphen{}computed before each run begins.  The
loID/hiID is set for the molecule with the smallest/largest ID which
has any itype or jtype atoms in it.  Note that if you define a system
with many molecule IDs between loID and hiID which have no itype or
jtype atoms, then the fix will be inefficient at performing swaps.
Also note that if atoms with molecule ID = 0 exist, they are not
considered molecules by this fix; they are assumed to be solvent atoms
or molecules.

\sphinxAtStartPar
Candidate atoms for swapping must also be in the fix group.  Atoms
within the selected molecule which are not itype or jtype are ignored.

\sphinxAtStartPar
When an atom is swapped from itype to jtype (or vice versa), if
charges are defined, the charge values for itype versus jtype atoms
are also swapped.  This requires that all itype atoms in the system
have the same charge value.  Likewise all jtype atoms in the system
must have the same charge value.  If this is not the case, LAMMPS
issues a warning that it cannot swap charge values.

\sphinxAtStartPar
If the \sphinxstyleemphasis{ke} keyword is set to yes, which is the default, and the
masses of itype and jtype atoms are different, then when a swap
occurs, the velocity of the swapped atom is rescaled by the sqrt of
the mass ratio, so as to conserve the kinetic energy of the atom.


\bigskip\hrule\bigskip


\sphinxAtStartPar
The potential energy of the entire system is computed before and after
each swap is performed within a single molecule.  The specified
temperature T is used in the Metropolis criterion to accept or reject
the attempted swap.  If the swap is rejected all swapped values are
reversed.

\sphinxAtStartPar
The potential energy calculations can include systems and models with
the following features:
\begin{itemize}
\item {} 
\sphinxAtStartPar
manybody pair styles, including EAM

\item {} 
\sphinxAtStartPar
hybrid pair styles

\item {} 
\sphinxAtStartPar
long\sphinxhyphen{}range electrostatics (kspace)

\item {} 
\sphinxAtStartPar
triclinic systems

\item {} 
\sphinxAtStartPar
potential energy contributions from other fixes

\end{itemize}

\sphinxAtStartPar
For the last bullet point, fixes can have an associated potential
energy. Examples of such fixes include: {\hyperref[\detokenize{fix_efield::doc}]{\sphinxcrossref{\DUrole{doc}{efield}}}},
{\hyperref[\detokenize{fix_gravity::doc}]{\sphinxcrossref{\DUrole{doc}{gravity}}}}, {\hyperref[\detokenize{fix_addforce::doc}]{\sphinxcrossref{\DUrole{doc}{addforce}}}},
{\hyperref[\detokenize{fix_langevin::doc}]{\sphinxcrossref{\DUrole{doc}{langevin}}}}, {\hyperref[\detokenize{fix_restrain::doc}]{\sphinxcrossref{\DUrole{doc}{restrain}}}},
{\hyperref[\detokenize{fix_temp_berendsen::doc}]{\sphinxcrossref{\DUrole{doc}{temp/berendsen}}}}, {\hyperref[\detokenize{fix_temp_rescale::doc}]{\sphinxcrossref{\DUrole{doc}{temp/rescale}}}}, and {\hyperref[\detokenize{fix_wall::doc}]{\sphinxcrossref{\DUrole{doc}{wall fixes}}}}.  For that
energy to be included in the total potential energy of the system (the
quantity used for the swap accept/reject decision), you MUST enable
the {\hyperref[\detokenize{fix_modify::doc}]{\sphinxcrossref{\DUrole{doc}{fix\_modify}}}} \sphinxstyleemphasis{energy} option for that fix.  The
doc pages for individual {\hyperref[\detokenize{fix::doc}]{\sphinxcrossref{\DUrole{doc}{fix}}}} commands specify if this
should be done.

\begin{sphinxadmonition}{note}{Note:}
\sphinxAtStartPar
One comment on computational efficiency.  If the cutoff lengths
defined for the pair style are different for itype versus jtype
atoms (for any of their interactions with any other atom type), then
a new neighbor list needs to be generated for every attempted swap.
This is potentially expensive if N is small or X is large.
\end{sphinxadmonition}


\subsection{Restart, fix\_modify, output, run start/stop, minimize info}
\label{\detokenize{fix_mol_swap:restart-fix-modify-output-run-start-stop-minimize-info}}
\sphinxAtStartPar
This fix writes the state of the fix to {\hyperref[\detokenize{restart::doc}]{\sphinxcrossref{\DUrole{doc}{binary restart files}}}}.  This includes information about the random number
generator seed, the next timestep for MC exchanges, the number of
exchange attempts and successes etc.  See the {\hyperref[\detokenize{read_restart::doc}]{\sphinxcrossref{\DUrole{doc}{read\_restart}}}} command for info on how to re\sphinxhyphen{}specify a fix in an
input script that reads a restart file, so that the operation of the
fix continues in an uninterrupted fashion.

\begin{sphinxadmonition}{note}{Note:}
\sphinxAtStartPar
For this to work correctly, the timestep must \sphinxstylestrong{not} be changed
after reading the restart with {\hyperref[\detokenize{reset_timestep::doc}]{\sphinxcrossref{\DUrole{doc}{reset\_timestep}}}}.
The fix will try to detect it and stop with an error.
\end{sphinxadmonition}

\sphinxAtStartPar
None of the {\hyperref[\detokenize{fix_modify::doc}]{\sphinxcrossref{\DUrole{doc}{fix\_modify}}}} options are relevant to this
fix.

\sphinxAtStartPar
This fix computes a global vector of length 2, which can be accessed
by various {\hyperref[\detokenize{Howto_output::doc}]{\sphinxcrossref{\DUrole{doc}{output commands}}}}.  The vector values are
the following global cumulative quantities:
\begin{enumerate}
\sphinxsetlistlabels{\arabic}{enumi}{enumii}{}{.}%
\item {} 
\sphinxAtStartPar
swap attempts

\item {} 
\sphinxAtStartPar
swap accepts

\end{enumerate}

\sphinxAtStartPar
The vector values calculated by this fix are “intensive”.

\sphinxAtStartPar
No parameter of this fix can be used with the \sphinxstyleemphasis{start/stop} keywords of
the {\hyperref[\detokenize{run::doc}]{\sphinxcrossref{\DUrole{doc}{run}}}} command.  This fix is not invoked during
{\hyperref[\detokenize{minimize::doc}]{\sphinxcrossref{\DUrole{doc}{energy minimization}}}}.


\subsection{Restrictions}
\label{\detokenize{fix_mol_swap:restrictions}}
\sphinxAtStartPar
This fix is part of the MC package.  It is only enabled if LAMMPS was
built with that package.  See the {\hyperref[\detokenize{Build_package::doc}]{\sphinxcrossref{\DUrole{doc}{Build package}}}}
doc page for more info.


\subsection{Related commands}
\label{\detokenize{fix_mol_swap:related-commands}}
\sphinxAtStartPar
{\hyperref[\detokenize{fix_atom_swap::doc}]{\sphinxcrossref{\DUrole{doc}{fix atom/swap}}}}, {\hyperref[\detokenize{fix_gcmc::doc}]{\sphinxcrossref{\DUrole{doc}{fix gcmc}}}}


\subsection{Default}
\label{\detokenize{fix_mol_swap:default}}
\sphinxAtStartPar
The option default is ke = yes.

\sphinxstepscope

\index{fix momentum@\spxentry{fix momentum}}\ignorespaces 
\index{fix momentum/kk@\spxentry{fix momentum/kk}}\ignorespaces 
\index{fix momentum/chunk@\spxentry{fix momentum/chunk}}\ignorespaces 

\section{fix momentum command}
\label{\detokenize{fix_momentum:fix-momentum-command}}\label{\detokenize{fix_momentum:index-2}}\label{\detokenize{fix_momentum:index-1}}\label{\detokenize{fix_momentum:index-0}}\label{\detokenize{fix_momentum::doc}}
\sphinxAtStartPar
Accelerator Variants: \sphinxstyleemphasis{momentum/kk}


\section{fix momentum/chunk command}
\label{\detokenize{fix_momentum:fix-momentum-chunk-command}}

\subsection{Syntax}
\label{\detokenize{fix_momentum:syntax}}
\begin{sphinxVerbatim}[commandchars=\\\{\}]
\PYG{k}{fix }\PYG{n+nv+nvIdentifier}{ID}\PYG{+w}{ }\PYG{n+nv+nvIdentifier}{group\PYGZhy{}ID}\PYG{+w}{ }\PYG{n}{momentum}\PYG{+w}{ }\PYG{n}{N}\PYG{+w}{ }\PYG{n}{keyword}\PYG{+w}{ }\PYG{n}{values}\PYG{+w}{ }\PYG{n}{...}
\end{sphinxVerbatim}
\begin{itemize}
\item {} 
\sphinxAtStartPar
ID, group\sphinxhyphen{}ID are documented in {\hyperref[\detokenize{fix::doc}]{\sphinxcrossref{\DUrole{doc}{fix}}}} command

\item {} 
\sphinxAtStartPar
momentum = style name of this fix command

\item {} 
\sphinxAtStartPar
N = adjust the momentum every this many timesteps
one or more keyword/value pairs may be appended

\end{itemize}

\begin{sphinxVerbatim}[commandchars=\\\{\}]
fix ID group\PYGZhy{}ID momentum/chunk N chunkID keyword values ...
\end{sphinxVerbatim}
\begin{itemize}
\item {} 
\sphinxAtStartPar
ID, group\sphinxhyphen{}ID are documented in {\hyperref[\detokenize{fix::doc}]{\sphinxcrossref{\DUrole{doc}{fix}}}} command

\item {} 
\sphinxAtStartPar
momentum/chunk = style name of this fix command

\item {} 
\sphinxAtStartPar
N = adjust the momentum per chunk every this many timesteps

\item {} 
\sphinxAtStartPar
chunkID = ID of {\hyperref[\detokenize{compute_chunk_atom::doc}]{\sphinxcrossref{\DUrole{doc}{compute chunk/atom}}}} command

\sphinxAtStartPar
one or more keyword/value settings may be appended to each of the fix commands:

\item {} 
\sphinxAtStartPar
keyword = \sphinxstyleemphasis{linear} or \sphinxstyleemphasis{angular} or \sphinxstyleemphasis{rescale}
\begin{sphinxalltt}
\sphinxstyleemphasis{linear} values = xflag yflag zflag
  xflag,yflag,zflag = 0/1 to exclude/include each dimension
\sphinxstyleemphasis{angular} values = none
\sphinxstyleemphasis{rescale} values = none
\end{sphinxalltt}

\end{itemize}


\subsection{Examples}
\label{\detokenize{fix_momentum:examples}}
\begin{sphinxVerbatim}[commandchars=\\\{\}]
\PYG{k}{fix }\PYG{n+nv+nvIdentifier}{1}\PYG{+w}{ }\PYG{n+nv+nvIdentifier}{all}\PYG{+w}{ }\PYG{n}{momentum}\PYG{+w}{ }\PYG{l+m}{1}\PYG{+w}{ }\PYG{n}{linear}\PYG{+w}{ }\PYG{l+m}{1}\PYG{+w}{ }\PYG{l+m}{1}\PYG{+w}{ }\PYG{l+m}{0}
\PYG{k}{fix }\PYG{n+nv+nvIdentifier}{1}\PYG{+w}{ }\PYG{n+nv+nvIdentifier}{all}\PYG{+w}{ }\PYG{n}{momentum}\PYG{+w}{ }\PYG{l+m}{1}\PYG{+w}{ }\PYG{n}{linear}\PYG{+w}{ }\PYG{l+m}{1}\PYG{+w}{ }\PYG{l+m}{1}\PYG{+w}{ }\PYG{l+m}{1}\PYG{+w}{ }\PYG{n}{rescale}
\PYG{k}{fix }\PYG{n+nv+nvIdentifier}{1}\PYG{+w}{ }\PYG{n+nv+nvIdentifier}{all}\PYG{+w}{ }\PYG{n}{momentum}\PYG{+w}{ }\PYG{l+m}{100}\PYG{+w}{ }\PYG{n}{linear}\PYG{+w}{ }\PYG{l+m}{1}\PYG{+w}{ }\PYG{l+m}{1}\PYG{+w}{ }\PYG{l+m}{1}\PYG{+w}{ }\PYG{n}{angular}
\PYG{k}{fix }\PYG{n+nv+nvIdentifier}{1}\PYG{+w}{ }\PYG{n+nv+nvIdentifier}{all}\PYG{+w}{ }\PYG{n}{momentum}\PYG{o}{/}\PYG{n}{chunk}\PYG{+w}{ }\PYG{l+m}{100}\PYG{+w}{ }\PYG{n}{molchunk}\PYG{+w}{ }\PYG{n}{linear}\PYG{+w}{ }\PYG{l+m}{1}\PYG{+w}{ }\PYG{l+m}{1}\PYG{+w}{ }\PYG{l+m}{1}\PYG{+w}{ }\PYG{n}{angular}
\end{sphinxVerbatim}


\subsection{Description}
\label{\detokenize{fix_momentum:description}}
\sphinxAtStartPar
Fix momentum zeroes the linear and/or angular momentum of the group of
atoms every N timesteps by adjusting the velocities of the atoms.
Fix momentum/chunk works equivalently, but operates on a per\sphinxhyphen{}chunk basis.

\sphinxAtStartPar
One (or both) of the \sphinxstyleemphasis{linear} or \sphinxstyleemphasis{angular} keywords \sphinxstylestrong{must} be specified.

\sphinxAtStartPar
If the \sphinxstyleemphasis{linear} keyword is used, the linear momentum is zeroed by
subtracting the center\sphinxhyphen{}of\sphinxhyphen{}mass velocity of the group or chunk from each
atom.  This does not change the relative velocity of any pair of atoms.
One or more dimensions can be excluded from this operation by setting
the corresponding flag to 0.

\sphinxAtStartPar
If the \sphinxstyleemphasis{angular} keyword is used, the angular momentum is zeroed by
subtracting a rotational component from each atom.

\sphinxAtStartPar
This command can be used to ensure the entire collection of atoms (or
a subset of them) does not drift or rotate during the simulation due
to random perturbations (e.g. {\hyperref[\detokenize{fix_langevin::doc}]{\sphinxcrossref{\DUrole{doc}{fix langevin}}}}
thermostatting).

\sphinxAtStartPar
The \sphinxstyleemphasis{rescale} keyword enables conserving the kinetic energy of the group
or chunk of atoms by rescaling the velocities after the momentum was removed.

\sphinxAtStartPar
Note that the {\hyperref[\detokenize{velocity::doc}]{\sphinxcrossref{\DUrole{doc}{velocity}}}} command can be used to create
initial velocities with zero aggregate linear and/or angular momentum.


\bigskip\hrule\bigskip


\sphinxAtStartPar
Styles with a \sphinxstyleemphasis{gpu}, \sphinxstyleemphasis{intel}, \sphinxstyleemphasis{kk}, \sphinxstyleemphasis{omp}, or \sphinxstyleemphasis{opt} suffix are
functionally the same as the corresponding style without the suffix.
They have been optimized to run faster, depending on your available
hardware, as discussed on the {\hyperref[\detokenize{Speed_packages::doc}]{\sphinxcrossref{\DUrole{doc}{Accelerator packages}}}}
page.  The accelerated styles take the same arguments and should
produce the same results, except for round\sphinxhyphen{}off and precision issues.

\sphinxAtStartPar
These accelerated styles are part of the GPU, INTEL, KOKKOS,
OPENMP, and OPT packages, respectively.  They are only enabled if
LAMMPS was built with those packages.  See the {\hyperref[\detokenize{Build_package::doc}]{\sphinxcrossref{\DUrole{doc}{Build package}}}} page for more info.

\sphinxAtStartPar
You can specify the accelerated styles explicitly in your input script
by including their suffix, or you can use the {\hyperref[\detokenize{Run_options::doc}]{\sphinxcrossref{\DUrole{doc}{\sphinxhyphen{}suffix command\sphinxhyphen{}line switch}}}} when you invoke LAMMPS, or you can use the
{\hyperref[\detokenize{suffix::doc}]{\sphinxcrossref{\DUrole{doc}{suffix}}}} command in your input script.

\sphinxAtStartPar
See the {\hyperref[\detokenize{Speed_packages::doc}]{\sphinxcrossref{\DUrole{doc}{Accelerator packages}}}} page for more
instructions on how to use the accelerated styles effectively.


\subsection{Restart, fix\_modify, output, run start/stop, minimize info}
\label{\detokenize{fix_momentum:restart-fix-modify-output-run-start-stop-minimize-info}}
\sphinxAtStartPar
No information about this fix is written to {\hyperref[\detokenize{restart::doc}]{\sphinxcrossref{\DUrole{doc}{binary restart files}}}}.  None of the {\hyperref[\detokenize{fix_modify::doc}]{\sphinxcrossref{\DUrole{doc}{fix\_modify}}}} options are
relevant to this fix.  No global or per\sphinxhyphen{}atom quantities are stored by
this fix for access by various {\hyperref[\detokenize{Howto_output::doc}]{\sphinxcrossref{\DUrole{doc}{output commands}}}}.
No parameter of this fix can be used with the \sphinxstyleemphasis{start/stop} keywords of
the {\hyperref[\detokenize{run::doc}]{\sphinxcrossref{\DUrole{doc}{run}}}} command.  This fix is not invoked during
{\hyperref[\detokenize{minimize::doc}]{\sphinxcrossref{\DUrole{doc}{energy minimization}}}}.


\subsection{Restrictions}
\label{\detokenize{fix_momentum:restrictions}}
\sphinxAtStartPar
Fix momentum/chunk is part of the EXTRA\sphinxhyphen{}FIX package.  It is only enabled
if LAMMPS was built with that package.  See the {\hyperref[\detokenize{Build_package::doc}]{\sphinxcrossref{\DUrole{doc}{Build package}}}} page for more info.


\subsection{Related commands}
\label{\detokenize{fix_momentum:related-commands}}
\sphinxAtStartPar
{\hyperref[\detokenize{fix_recenter::doc}]{\sphinxcrossref{\DUrole{doc}{fix recenter}}}}, {\hyperref[\detokenize{velocity::doc}]{\sphinxcrossref{\DUrole{doc}{velocity}}}}


\subsection{Default}
\label{\detokenize{fix_momentum:default}}
\sphinxAtStartPar
none

\sphinxstepscope

\index{fix move@\spxentry{fix move}}\ignorespaces 

\section{fix move command}
\label{\detokenize{fix_move:fix-move-command}}\label{\detokenize{fix_move:index-0}}\label{\detokenize{fix_move::doc}}

\subsection{Syntax}
\label{\detokenize{fix_move:syntax}}
\begin{sphinxVerbatim}[commandchars=\\\{\}]
\PYG{k}{fix }\PYG{n+nv+nvIdentifier}{ID}\PYG{+w}{ }\PYG{n+nv+nvIdentifier}{group\PYGZhy{}ID}\PYG{+w}{ }\PYG{n}{move}\PYG{+w}{ }\PYG{n}{style}\PYG{+w}{ }\PYG{n}{args}\PYG{+w}{ }\PYG{n}{keyword}\PYG{+w}{ }\PYG{n}{values}\PYG{+w}{ }\PYG{n}{...}
\end{sphinxVerbatim}
\begin{itemize}
\item {} 
\sphinxAtStartPar
ID, group\sphinxhyphen{}ID are documented in {\hyperref[\detokenize{fix::doc}]{\sphinxcrossref{\DUrole{doc}{fix}}}} command

\item {} 
\sphinxAtStartPar
move = style name of this fix command

\item {} 
\sphinxAtStartPar
style = \sphinxstyleemphasis{linear} or \sphinxstyleemphasis{wiggle} or \sphinxstyleemphasis{rotate} or \sphinxstyleemphasis{transrot} or \sphinxstyleemphasis{variable}
\begin{sphinxalltt}
\sphinxstyleemphasis{linear} args = Vx Vy Vz
  Vx,Vy,Vz = components of velocity vector (velocity units), any component can be specified as NULL
\sphinxstyleemphasis{wiggle} args = Ax Ay Az period
  Ax,Ay,Az = components of amplitude vector (distance units), any component can be specified as NULL
  period = period of oscillation (time units)
\sphinxstyleemphasis{rotate} args = Px Py Pz Rx Ry Rz period
  Px,Py,Pz = origin point of axis of rotation (distance units)
  Rx,Ry,Rz = axis of rotation vector
  period = period of rotation (time units)
\sphinxstyleemphasis{transrot} args = Vx Vy Vz Px Py Pz Rx Ry Rz period
  Vx,Vy,Vz = components of velocity vector (velocity units)
  Px,Py,Pz = origin point of axis of rotation (distance units)
  Rx,Ry,Rz = axis of rotation vector
  period = period of rotation (time units)
\sphinxstyleemphasis{variable} args = v\_dx v\_dy v\_dz v\_vx v\_vy v\_vz
  v\_dx,v\_dy,v\_dz = 3 variable names that calculate x,y,z displacement as function of time, any component can be specified as NULL
  v\_vx,v\_vy,v\_vz = 3 variable names that calculate x,y,z velocity as function of time, any component can be specified as NULL
\end{sphinxalltt}

\item {} 
\sphinxAtStartPar
zero or more keyword/value pairs may be appended

\item {} 
\sphinxAtStartPar
keyword = \sphinxstyleemphasis{units}
\begin{sphinxalltt}
\sphinxstyleemphasis{units} value = \sphinxstyleemphasis{box} or \sphinxstyleemphasis{lattice}
\end{sphinxalltt}

\end{itemize}


\subsection{Examples}
\label{\detokenize{fix_move:examples}}
\begin{sphinxVerbatim}[commandchars=\\\{\}]
\PYG{k}{fix }\PYG{n+nv+nvIdentifier}{1}\PYG{+w}{ }\PYG{n+nv+nvIdentifier}{boundary}\PYG{+w}{ }\PYG{n}{move}\PYG{+w}{ }\PYG{n}{wiggle}\PYG{+w}{ }\PYG{l+m}{3.0}\PYG{+w}{ }\PYG{l+m}{0.0}\PYG{+w}{ }\PYG{l+m}{0.0}\PYG{+w}{ }\PYG{l+m}{1.0}\PYG{+w}{ }\PYG{n}{units}\PYG{+w}{ }\PYG{n}{box}
\PYG{k}{fix }\PYG{n+nv+nvIdentifier}{2}\PYG{+w}{ }\PYG{n+nv+nvIdentifier}{boundary}\PYG{+w}{ }\PYG{n}{move}\PYG{+w}{ }\PYG{n}{rotate}\PYG{+w}{ }\PYG{l+m}{0.0}\PYG{+w}{ }\PYG{l+m}{0.0}\PYG{+w}{ }\PYG{l+m}{0.0}\PYG{+w}{ }\PYG{l+m}{0.0}\PYG{+w}{ }\PYG{l+m}{0.0}\PYG{+w}{ }\PYG{l+m}{1.0}\PYG{+w}{ }\PYG{l+m}{5.0}
\PYG{k}{fix }\PYG{n+nv+nvIdentifier}{2}\PYG{+w}{ }\PYG{n+nv+nvIdentifier}{boundary}\PYG{+w}{ }\PYG{n}{move}\PYG{+w}{ }\PYG{k}{variable }\PYG{n+nv+nvIdentifier}{v\PYGZus{}myx}\PYG{+w}{ }\PYG{n}{v\PYGZus{}myy}\PYG{+w}{ }\PYG{n}{NULL}\PYG{+w}{ }\PYG{n}{v\PYGZus{}VX}\PYG{+w}{ }\PYG{n}{v\PYGZus{}VY}\PYG{+w}{ }\PYG{n}{NULL}
\PYG{k}{fix }\PYG{n+nv+nvIdentifier}{3}\PYG{+w}{ }\PYG{n+nv+nvIdentifier}{boundary}\PYG{+w}{ }\PYG{n}{move}\PYG{+w}{ }\PYG{n}{transrot}\PYG{+w}{ }\PYG{l+m}{0.1}\PYG{+w}{ }\PYG{l+m}{0.1}\PYG{+w}{ }\PYG{l+m}{0.0}\PYG{+w}{ }\PYG{l+m}{0.0}\PYG{+w}{ }\PYG{l+m}{0.0}\PYG{+w}{ }\PYG{l+m}{0.0}\PYG{+w}{ }\PYG{l+m}{0.0}\PYG{+w}{ }\PYG{l+m}{0.0}\PYG{+w}{ }\PYG{l+m}{1.0}\PYG{+w}{ }\PYG{l+m}{5.0}\PYG{+w}{ }\PYG{n}{units}\PYG{+w}{ }\PYG{n}{box}
\end{sphinxVerbatim}


\subsection{Description}
\label{\detokenize{fix_move:description}}
\sphinxAtStartPar
Perform updates of position and velocity for atoms in the group each
timestep using the specified settings or formulas, without regard to
forces on the atoms.  This can be useful for boundary or other atoms,
whose movement can influence nearby atoms.

\begin{sphinxadmonition}{note}{Note:}
\sphinxAtStartPar
The atoms affected by this fix should not normally be time integrated
by other fixes (e.g. {\hyperref[\detokenize{fix_nve::doc}]{\sphinxcrossref{\DUrole{doc}{fix nve}}}}, {\hyperref[\detokenize{fix_nh::doc}]{\sphinxcrossref{\DUrole{doc}{fix nvt}}}}), since that will change their positions and velocities
twice.
\end{sphinxadmonition}

\begin{sphinxadmonition}{note}{Note:}
\sphinxAtStartPar
As atoms move due to this fix, they will pass through periodic
boundaries and be remapped to the other side of the simulation box,
just as they would during normal time integration (e.g. via the
{\hyperref[\detokenize{fix_nve::doc}]{\sphinxcrossref{\DUrole{doc}{fix nve}}}} command).  It is up to you to decide whether
periodic boundaries are appropriate with the kind of atom motion you
are prescribing with this fix.
\end{sphinxadmonition}

\begin{sphinxadmonition}{note}{Note:}
\sphinxAtStartPar
As discussed below, atoms are moved relative to their initial
position at the time the fix is specified.  These initial coordinates
are stored by the fix in “unwrapped” form, by using the image flags
associated with each atom.  See the {\hyperref[\detokenize{dump::doc}]{\sphinxcrossref{\DUrole{doc}{dump custom}}}} command
for a discussion of “unwrapped” coordinates.  See the Atoms section
of the {\hyperref[\detokenize{read_data::doc}]{\sphinxcrossref{\DUrole{doc}{read\_data}}}} command for a discussion of image
flags and how they are set for each atom.  You can reset the image
flags (e.g. to 0) before invoking this fix by using the {\hyperref[\detokenize{set::doc}]{\sphinxcrossref{\DUrole{doc}{set
image}}}} command.
\end{sphinxadmonition}


\bigskip\hrule\bigskip


\sphinxAtStartPar
The \sphinxstyleemphasis{linear} style moves atoms at a constant velocity, so that their
position \sphinxstyleemphasis{X} = (x,y,z) as a function of time is given in vector
notation as
\begin{sphinxalltt}
X(t) = X0 + V * delta
\end{sphinxalltt}

\sphinxAtStartPar
where \sphinxstyleemphasis{X0} = (x0,y0,z0) is their position at the time the fix is
specified, \sphinxstyleemphasis{V} is the specified velocity vector with components
(Vx,Vy,Vz), and \sphinxstyleemphasis{delta} is the time elapsed since the fix was
specified.  This style also sets the velocity of each atom to V =
(Vx,Vy,Vz).  If any of the velocity components is specified as NULL,
then the position and velocity of that component is time integrated
the same as the {\hyperref[\detokenize{fix_nve::doc}]{\sphinxcrossref{\DUrole{doc}{fix nve}}}} command would perform, using
the corresponding force component on the atom.

\sphinxAtStartPar
Note that the \sphinxstyleemphasis{linear} style is identical to using the \sphinxstyleemphasis{variable}
style with an {\hyperref[\detokenize{variable::doc}]{\sphinxcrossref{\DUrole{doc}{equal\sphinxhyphen{}style variable}}}} that uses the
vdisplace() function.  E.g.

\begin{sphinxVerbatim}[commandchars=\\\{\}]
\PYG{k}{variable }\PYG{n+nv+nvIdentifier}{V}\PYG{+w}{ }\PYG{n}{equal}\PYG{+w}{ }\PYG{l+m}{10.0}
\PYG{k}{variable }\PYG{n+nv+nvIdentifier}{x}\PYG{+w}{ }\PYG{n}{equal}\PYG{+w}{ }\PYG{n}{vdisplace}\PYG{n+nv}{(}\PYG{n+nv}{0.0,\PYGZdl{}V}\PYG{n+nv}{)}
\PYG{k}{fix }\PYG{n+nv+nvIdentifier}{1}\PYG{+w}{ }\PYG{n+nv+nvIdentifier}{boundary}\PYG{+w}{ }\PYG{n}{move}\PYG{+w}{ }\PYG{k}{variable }\PYG{n+nv+nvIdentifier}{v\PYGZus{}x}\PYG{+w}{ }\PYG{n}{NULL}\PYG{+w}{ }\PYG{n}{NULL}\PYG{+w}{ }\PYG{n}{v\PYGZus{}V}\PYG{+w}{ }\PYG{n}{NULL}\PYG{+w}{ }\PYG{n}{NULL}
\end{sphinxVerbatim}

\sphinxAtStartPar
The \sphinxstyleemphasis{wiggle} style moves atoms in an oscillatory fashion, so that
their position \sphinxstyleemphasis{X} = (x,y,z) as a function of time is given in vector
notation as
\begin{sphinxalltt}
X(t) = X0 + A sin(omega*delta)
\end{sphinxalltt}

\sphinxAtStartPar
where \sphinxstyleemphasis{X0} = (x0,y0,z0) is their position at the time the fix is
specified, \sphinxstyleemphasis{A} is the specified amplitude vector with components
(Ax,Ay,Az), \sphinxstyleemphasis{omega} is 2 PI / \sphinxstyleemphasis{period}, and \sphinxstyleemphasis{delta} is the time elapsed
since the fix was specified.  This style also sets the velocity of each
atom to the time derivative of this expression.  If any of the amplitude
components is specified as NULL, then the position and velocity of that
component is time integrated the same as the {\hyperref[\detokenize{fix_nve::doc}]{\sphinxcrossref{\DUrole{doc}{fix nve}}}}
command would perform, using the corresponding force component on the
atom.

\sphinxAtStartPar
Note that the \sphinxstyleemphasis{wiggle} style is identical to using the \sphinxstyleemphasis{variable}
style with {\hyperref[\detokenize{variable::doc}]{\sphinxcrossref{\DUrole{doc}{equal\sphinxhyphen{}style variables}}}} that use the
swiggle() and cwiggle() functions.  E.g.

\begin{sphinxVerbatim}[commandchars=\\\{\}]
\PYG{k}{variable }\PYG{n+nv+nvIdentifier}{A}\PYG{+w}{ }\PYG{n}{equal}\PYG{+w}{ }\PYG{l+m}{10.0}
\PYG{k}{variable }\PYG{n+nv+nvIdentifier}{T}\PYG{+w}{ }\PYG{n}{equal}\PYG{+w}{ }\PYG{l+m}{5.0}
\PYG{k}{variable }\PYG{n+nv+nvIdentifier}{omega}\PYG{+w}{ }\PYG{n}{equal}\PYG{+w}{ }\PYG{l+m}{2.0}\PYG{o}{*}\PYG{n}{PI}\PYG{o}{/}\PYG{n+nv}{\PYGZdl{}T}
\PYG{k}{variable }\PYG{n+nv+nvIdentifier}{x}\PYG{+w}{ }\PYG{n}{equal}\PYG{+w}{ }\PYG{n}{swiggle}\PYG{n+nv}{(}\PYG{n+nv}{0.0,\PYGZdl{}A,\PYGZdl{}T}\PYG{n+nv}{)}
\PYG{k}{variable }\PYG{n+nv+nvIdentifier}{v}\PYG{+w}{ }\PYG{n}{equal}\PYG{+w}{ }\PYG{n}{v\PYGZus{}omega}\PYG{o}{*}\PYG{n+nv}{(}\PYG{n+nv}{\PYGZdl{}A\PYGZhy{}cwiggle}\PYG{n+nv}{(}\PYG{n+nv}{0.0,\PYGZdl{}A,\PYGZdl{}T}\PYG{n+nv}{)}\PYG{n+nv}{)}
\PYG{k}{fix }\PYG{n+nv+nvIdentifier}{1}\PYG{+w}{ }\PYG{n+nv+nvIdentifier}{boundary}\PYG{+w}{ }\PYG{n}{move}\PYG{+w}{ }\PYG{k}{variable }\PYG{n+nv+nvIdentifier}{v\PYGZus{}x}\PYG{+w}{ }\PYG{n}{NULL}\PYG{+w}{ }\PYG{n}{NULL}\PYG{+w}{ }\PYG{n}{v\PYGZus{}v}\PYG{+w}{ }\PYG{n}{NULL}\PYG{+w}{ }\PYG{n}{NULL}
\end{sphinxVerbatim}

\sphinxAtStartPar
The \sphinxstyleemphasis{rotate} style rotates atoms around a rotation axis \sphinxstyleemphasis{R} = (Rx,Ry,Rz)
that goes through a point \sphinxstyleemphasis{P} = (Px,Py,Pz).  The \sphinxstyleemphasis{period} of the
rotation is also specified.  The direction of rotation for the atoms
around the rotation axis is consistent with the right\sphinxhyphen{}hand rule: if your
right\sphinxhyphen{}hand thumb points along \sphinxstyleemphasis{R}, then your fingers wrap around the
axis in the direction of rotation.

\sphinxAtStartPar
This style also sets the velocity of each atom to (omega cross Rperp)
where omega is its angular velocity around the rotation axis and Rperp
is a perpendicular vector from the rotation axis to the atom.  If the
defined {\hyperref[\detokenize{atom_style::doc}]{\sphinxcrossref{\DUrole{doc}{atom\_style}}}} assigns an angular velocity or
angular momentum or orientation to each atom ({\hyperref[\detokenize{atom_style::doc}]{\sphinxcrossref{\DUrole{doc}{atom styles}}}} sphere, ellipsoid, line, tri, body), then
those properties are also updated appropriately to correspond to the
atom’s motion and rotation over time.

\sphinxAtStartPar
The \sphinxstyleemphasis{transrot} style combines the effects of \sphinxstyleemphasis{rotate} and \sphinxstyleemphasis{linear} so
that it is possible to prescribe a rotating group of atoms that also
moves at a constant velocity.  The arguments are for the translation
first and then for the rotation.  Since the rotation affects all
coordinate components, it is not possible to set any of the
translation vector components to NULL.

\sphinxAtStartPar
The \sphinxstyleemphasis{variable} style allows the position and velocity components of
each atom to be set by formulas specified via the
{\hyperref[\detokenize{variable::doc}]{\sphinxcrossref{\DUrole{doc}{variable}}}} command.  Each of the 6 variables is
specified as an argument to the fix as v\_name, where name is the
variable name that is defined elsewhere in the input script.

\sphinxAtStartPar
Each variable must be of either the \sphinxstyleemphasis{equal} or \sphinxstyleemphasis{atom} style.
\sphinxstyleemphasis{Equal}\sphinxhyphen{}style variables compute a single numeric quantity, that can be
a function of the timestep as well as of other simulation values.
\sphinxstyleemphasis{Atom}\sphinxhyphen{}style variables compute a numeric quantity for each atom, that
can be a function per\sphinxhyphen{}atom quantities, such as the atom’s position, as
well as of the timestep and other simulation values.  Note that this fix
stores the original coordinates of each atom (see note below) so that
per\sphinxhyphen{}atom quantity can be used in an atom\sphinxhyphen{}style variable formula.  See
the {\hyperref[\detokenize{variable::doc}]{\sphinxcrossref{\DUrole{doc}{variable}}}} command for details.

\sphinxAtStartPar
The first 3 variables (v\_dx,v\_dy,v\_dz) specified for the \sphinxstyleemphasis{variable}
style are used to calculate a displacement from the atom’s original
position at the time the fix was specified.  The second 3 variables
(v\_vx,v\_vy,v\_vz) specified are used to compute a velocity for each
atom.

\sphinxAtStartPar
Any of the 6 variables can be specified as NULL.  If both the
displacement and velocity variables for a particular x,y,z component
are specified as NULL, then the position and velocity of that
component is time integrated the same as the {\hyperref[\detokenize{fix_nve::doc}]{\sphinxcrossref{\DUrole{doc}{fix nve}}}}
command would perform, using the corresponding force component on the
atom.  If only the velocity variable for a component is specified as
NULL, then the displacement variable will be used to set the position
of the atom, and its velocity component will not be changed.  If only
the displacement variable for a component is specified as NULL, then
the velocity variable will be used to set the velocity of the atom,
and the position of the atom will be time integrated using that
velocity.

\sphinxAtStartPar
The \sphinxstyleemphasis{units} keyword determines the meaning of the distance units used
to define the \sphinxstyleemphasis{linear} velocity and \sphinxstyleemphasis{wiggle} amplitude and \sphinxstyleemphasis{rotate}
origin.  This setting is ignored for the \sphinxstyleemphasis{variable} style.  A \sphinxstyleemphasis{box}
value selects standard units as defined by the {\hyperref[\detokenize{units::doc}]{\sphinxcrossref{\DUrole{doc}{units}}}}
command, e.g. velocity in Angstroms/fs and amplitude and position
in Angstroms for units = real.  A \sphinxstyleemphasis{lattice} value means the velocity
units are in lattice spacings per time and the amplitude and position
are in lattice spacings.  The {\hyperref[\detokenize{lattice::doc}]{\sphinxcrossref{\DUrole{doc}{lattice}}}} command must have
been previously used to define the lattice spacing.  Each of these 3
quantities may be dependent on the x,y,z dimension, since the lattice
spacings can be different in x,y,z.


\bigskip\hrule\bigskip



\subsection{Restart, fix\_modify, output, run start/stop, minimize info}
\label{\detokenize{fix_move:restart-fix-modify-output-run-start-stop-minimize-info}}
\sphinxAtStartPar
This fix writes the original coordinates of moving atoms to {\hyperref[\detokenize{restart::doc}]{\sphinxcrossref{\DUrole{doc}{binary
restart files}}}}, as well as the initial timestep, so that the
motion can be continuous in a restarted simulation.  See the
{\hyperref[\detokenize{read_restart::doc}]{\sphinxcrossref{\DUrole{doc}{read\_restart}}}} command for info on how to re\sphinxhyphen{}specify
a fix in an input script that reads a restart file, so that the
operation of the fix continues in an uninterrupted fashion.

\begin{sphinxadmonition}{note}{Note:}
\sphinxAtStartPar
Because the move positions are a function of the current
timestep and the initial timestep, you cannot reset the timestep to a
different value after reading a restart file, if you expect a fix move
command to work in an uninterrupted fashion.
\end{sphinxadmonition}

\sphinxAtStartPar
None of the {\hyperref[\detokenize{fix_modify::doc}]{\sphinxcrossref{\DUrole{doc}{fix\_modify}}}} options are relevant to this
fix.

\sphinxAtStartPar
This fix produces a per\sphinxhyphen{}atom array which can be accessed by various
{\hyperref[\detokenize{Howto_output::doc}]{\sphinxcrossref{\DUrole{doc}{output commands}}}}.  The number of columns for each
atom is 3, and the columns store the original unwrapped x,y,z coords of
each atom.  The per\sphinxhyphen{}atom values can be accessed on any timestep.

\sphinxAtStartPar
No parameter of this fix can be used with the \sphinxstyleemphasis{start/stop} keywords of
the {\hyperref[\detokenize{run::doc}]{\sphinxcrossref{\DUrole{doc}{run}}}} command.  This fix is not invoked during
{\hyperref[\detokenize{minimize::doc}]{\sphinxcrossref{\DUrole{doc}{energy minimization}}}}.

\sphinxAtStartPar
For {\hyperref[\detokenize{run_style::doc}]{\sphinxcrossref{\DUrole{doc}{rRESPA time integration}}}}, this fix adjusts the
position and velocity of atoms on the outermost rRESPA level.


\subsection{Restrictions}
\label{\detokenize{fix_move:restrictions}}\begin{quote}

\sphinxAtStartPar
none
\end{quote}


\subsection{Related commands}
\label{\detokenize{fix_move:related-commands}}
\sphinxAtStartPar
{\hyperref[\detokenize{fix_nve::doc}]{\sphinxcrossref{\DUrole{doc}{fix nve}}}}, {\hyperref[\detokenize{displace_atoms::doc}]{\sphinxcrossref{\DUrole{doc}{displace\_atoms}}}}


\subsection{Default}
\label{\detokenize{fix_move:default}}
\sphinxAtStartPar
none

\sphinxAtStartPar
The option default is units = lattice.

\sphinxstepscope

\index{fix msst@\spxentry{fix msst}}\ignorespaces 

\section{fix msst command}
\label{\detokenize{fix_msst:fix-msst-command}}\label{\detokenize{fix_msst:index-0}}\label{\detokenize{fix_msst::doc}}

\subsection{Syntax}
\label{\detokenize{fix_msst:syntax}}
\begin{sphinxVerbatim}[commandchars=\\\{\}]
\PYG{k}{fix }\PYG{n+nv+nvIdentifier}{ID}\PYG{+w}{ }\PYG{n+nv+nvIdentifier}{group\PYGZhy{}ID}\PYG{+w}{ }\PYG{n}{msst}\PYG{+w}{ }\PYG{n}{dir}\PYG{+w}{ }\PYG{n}{shockvel}\PYG{+w}{ }\PYG{n}{keyword}\PYG{+w}{ }\PYG{n}{value}\PYG{+w}{ }\PYG{n}{...}
\end{sphinxVerbatim}
\begin{itemize}
\item {} 
\sphinxAtStartPar
ID, group\sphinxhyphen{}ID are documented in {\hyperref[\detokenize{fix::doc}]{\sphinxcrossref{\DUrole{doc}{fix}}}} command

\item {} 
\sphinxAtStartPar
msst = style name of this fix

\item {} 
\sphinxAtStartPar
dir = \sphinxstyleemphasis{x} or \sphinxstyleemphasis{y} or \sphinxstyleemphasis{z}

\item {} 
\sphinxAtStartPar
shockvel = shock velocity (strictly positive, distance/time units)

\item {} 
\sphinxAtStartPar
zero or more keyword value pairs may be appended

\item {} 
\sphinxAtStartPar
keyword = \sphinxstyleemphasis{q} or \sphinxstyleemphasis{mu} or \sphinxstyleemphasis{p0} or \sphinxstyleemphasis{v0} or \sphinxstyleemphasis{e0} or \sphinxstyleemphasis{tscale} or \sphinxstyleemphasis{beta} or \sphinxstyleemphasis{dftb}
\begin{sphinxalltt}
\sphinxstyleemphasis{q} value = cell mass\sphinxhyphen{}like parameter (mass\textasciicircum{}2/distance\textasciicircum{}4 units)
\sphinxstyleemphasis{mu} value = artificial viscosity (mass/length/time units)
\sphinxstyleemphasis{p0} value = initial pressure in the shock equations (pressure units)
\sphinxstyleemphasis{v0} value = initial simulation cell volume in the shock equations (distance\textasciicircum{}3 units)
\sphinxstyleemphasis{e0} value = initial total energy (energy units)
\sphinxstyleemphasis{tscale} value = reduction in initial temperature (unitless fraction between 0.0 and 1.0)
\sphinxstyleemphasis{dftb} value = \sphinxstyleemphasis{yes} or \sphinxstyleemphasis{no} for whether using MSST in conjunction with DFTB+
\sphinxstyleemphasis{beta} value = scale factor for improved energy conservation
\end{sphinxalltt}

\end{itemize}


\subsection{Examples}
\label{\detokenize{fix_msst:examples}}
\begin{sphinxVerbatim}[commandchars=\\\{\}]
\PYG{k}{fix }\PYG{n+nv+nvIdentifier}{1}\PYG{+w}{ }\PYG{n+nv+nvIdentifier}{all}\PYG{+w}{ }\PYG{n}{msst}\PYG{+w}{ }\PYG{n}{y}\PYG{+w}{ }\PYG{l+m}{100.0}\PYG{+w}{ }\PYG{n}{q}\PYG{+w}{ }\PYG{l+m}{1.0e5}\PYG{+w}{ }\PYG{n}{mu}\PYG{+w}{ }\PYG{l+m}{1.0e5}
\PYG{k}{fix }\PYG{n+nv+nvIdentifier}{2}\PYG{+w}{ }\PYG{n+nv+nvIdentifier}{all}\PYG{+w}{ }\PYG{n}{msst}\PYG{+w}{ }\PYG{n}{z}\PYG{+w}{ }\PYG{l+m}{50.0}\PYG{+w}{ }\PYG{n}{q}\PYG{+w}{ }\PYG{l+m}{1.0e4}\PYG{+w}{ }\PYG{n}{mu}\PYG{+w}{ }\PYG{l+m}{1.0e4}\PYG{+w}{  }\PYG{n}{v0}\PYG{+w}{ }\PYG{l+m}{4.3419}\PYG{n}{e}\PYG{o}{+}\PYG{l+m}{03}\PYG{+w}{ }\PYG{n}{p0}\PYG{+w}{ }\PYG{l+m}{3.7797}\PYG{n}{e}\PYG{o}{+}\PYG{l+m}{03}\PYG{+w}{ }\PYG{n}{e0}\PYG{+w}{ }\PYG{o}{\PYGZhy{}}\PYG{l+m}{9.72360}\PYG{n}{e}\PYG{o}{+}\PYG{l+m}{02}\PYG{+w}{ }\PYG{n}{tscale}\PYG{+w}{ }\PYG{l+m}{0.01}
\PYG{k}{fix }\PYG{n+nv+nvIdentifier}{1}\PYG{+w}{ }\PYG{n+nv+nvIdentifier}{all}\PYG{+w}{ }\PYG{n}{msst}\PYG{+w}{ }\PYG{n}{y}\PYG{+w}{ }\PYG{l+m}{100.0}\PYG{+w}{ }\PYG{n}{q}\PYG{+w}{ }\PYG{l+m}{1.0e5}\PYG{+w}{ }\PYG{n}{mu}\PYG{+w}{ }\PYG{l+m}{1.0e5}\PYG{+w}{ }\PYG{n}{dftb}\PYG{+w}{ }\PYG{n}{yes}\PYG{+w}{ }\PYG{n}{beta}\PYG{+w}{ }\PYG{l+m}{0.5}
\end{sphinxVerbatim}


\subsection{Description}
\label{\detokenize{fix_msst:description}}
\sphinxAtStartPar
This command performs the Multi\sphinxhyphen{}Scale Shock Technique (MSST)
integration to update positions and velocities each timestep to mimic
a compressive shock wave passing over the system. See {\hyperref[\detokenize{fix_msst:reed}]{\sphinxcrossref{\DUrole{std,std-ref}{(Reed)}}}}
for a detailed description of this method.  The MSST varies the cell
volume and temperature in such a way as to restrain the system to the
shock Hugoniot and the Rayleigh line. These restraints correspond to
the macroscopic conservation laws dictated by a shock
front. \sphinxstyleemphasis{shockvel} determines the steady shock velocity that will be
simulated.

\sphinxAtStartPar
To perform a simulation, choose a value of \sphinxstyleemphasis{q} that provides volume
compression on the timescale of 100 fs to 1 ps.  If the volume is not
compressing, either the shock speed is chosen to be below the material
sound speed or \sphinxstyleemphasis{p0} has been chosen inaccurately.  Volume compression
at the start can be sped up by using a non\sphinxhyphen{}zero value of \sphinxstyleemphasis{tscale}. Use
the smallest value of \sphinxstyleemphasis{tscale} that results in compression.

\sphinxAtStartPar
Under some special high\sphinxhyphen{}symmetry conditions, the pressure (volume)
and/or temperature of the system may oscillate for many cycles even
with an appropriate choice of mass\sphinxhyphen{}like parameter \sphinxstyleemphasis{q}. Such
oscillations have physical significance in some cases.  The optional
\sphinxstyleemphasis{mu} keyword adds an artificial viscosity that helps break the system
symmetry to equilibrate to the shock Hugoniot and Rayleigh line more
rapidly in such cases.

\sphinxAtStartPar
The keyword \sphinxstyleemphasis{tscale} is a factor between 0 and 1 that determines what
fraction of thermal kinetic energy is converted to compressive strain
kinetic energy at the start of the simulation.  Setting this parameter
to a non\sphinxhyphen{}zero value may assist in compression at the start of
simulations where it is slow to occur.

\sphinxAtStartPar
If keywords \sphinxstyleemphasis{e0}, \sphinxstyleemphasis{p0},or \sphinxstyleemphasis{v0} are not supplied, these quantities will
be calculated on the first step, after the energy specified by
\sphinxstyleemphasis{tscale} is removed.  The value of \sphinxstyleemphasis{e0} is not used in the dynamical
equations, but is used in calculating the deviation from the Hugoniot.

\sphinxAtStartPar
The keyword \sphinxstyleemphasis{beta} is a scaling term that can be added to the MSST
ionic equations of motion to account for drift in the conserved
quantity during long timescale simulations, similar to a Berendsen
thermostat. See {\hyperref[\detokenize{fix_msst:reed}]{\sphinxcrossref{\DUrole{std,std-ref}{(Reed)}}}} and {\hyperref[\detokenize{fix_msst:goldman2}]{\sphinxcrossref{\DUrole{std,std-ref}{(Goldman)}}}} for more
details.  The value of \sphinxstyleemphasis{beta} must be between 0.0 and 1.0 inclusive.
A value of 0.0 means no contribution, a value of 1.0 means a full
contribution.

\sphinxAtStartPar
Values of shockvel less than a critical value determined by the
material response will not have compressive solutions. This will be
reflected in lack of significant change of the volume in the MSST.

\sphinxAtStartPar
For all pressure styles, the simulation box stays orthogonal in shape.
Parrinello\sphinxhyphen{}Rahman boundary conditions (tilted box) are supported by
LAMMPS, but are not implemented for MSST.

\sphinxAtStartPar
This fix computes a temperature and pressure and potential energy each
timestep. To do this, the fix creates its own computes of style “temp”
“pressure”, and “pe”, as if these commands had been issued:

\begin{sphinxVerbatim}[commandchars=\\\{\}]
\PYG{k}{compute }\PYG{n+nv+nvIdentifier}{fix\PYGZhy{}ID\PYGZus{}MSST\PYGZus{}temp}\PYG{+w}{ }\PYG{n+nv+nvIdentifier}{all}\PYG{+w}{ }\PYG{n}{temp}
\PYG{k}{compute }\PYG{n+nv+nvIdentifier}{fix\PYGZhy{}ID\PYGZus{}MSST\PYGZus{}press}\PYG{+w}{ }\PYG{n+nv+nvIdentifier}{all}\PYG{+w}{ }\PYG{n}{pressure}\PYG{+w}{ }\PYG{n}{fix}\PYG{o}{\PYGZhy{}}\PYG{n}{ID\PYGZus{}MSST\PYGZus{}temp}

\PYG{k}{compute }\PYG{n+nv+nvIdentifier}{fix\PYGZhy{}ID\PYGZus{}MSST\PYGZus{}pe}\PYG{+w}{ }\PYG{n+nv+nvIdentifier}{all}\PYG{+w}{ }\PYG{n}{pe}
\end{sphinxVerbatim}

\sphinxAtStartPar
See the {\hyperref[\detokenize{compute_temp::doc}]{\sphinxcrossref{\DUrole{doc}{compute temp}}}} and {\hyperref[\detokenize{compute_pressure::doc}]{\sphinxcrossref{\DUrole{doc}{compute pressure}}}} commands for details.  Note that the IDs of the
new computes are the fix\sphinxhyphen{}ID + “\_MSST\_temp” or “MSST\_press” or
“\_MSST\_pe”.  The group for the new computes is “all”.


\bigskip\hrule\bigskip


\sphinxAtStartPar
The \sphinxstyleemphasis{dftb} keyword is to allow this fix to be used when LAMMPS is
being driven by DFTB+, a density\sphinxhyphen{}functional tight\sphinxhyphen{}binding code. If the
keyword \sphinxstyleemphasis{dftb} is used with a value of \sphinxstyleemphasis{yes}, then the MSST equations
are altered to account for the electron entropy contribution to the
Hugonio relations and total energy.  See {\hyperref[\detokenize{fix_msst:reed2}]{\sphinxcrossref{\DUrole{std,std-ref}{(Reed2)}}}} and
{\hyperref[\detokenize{fix_msst:goldman2}]{\sphinxcrossref{\DUrole{std,std-ref}{(Goldman)}}}} for details on this contribution.  In this case,
you must define a {\hyperref[\detokenize{fix_external::doc}]{\sphinxcrossref{\DUrole{doc}{fix external}}}} command in your
input script, which is used to callback to DFTB+ during the LAMMPS
timestepping.  DFTB+ will communicate its info to LAMMPS via that fix.


\bigskip\hrule\bigskip



\subsection{Restart, fix\_modify, output, run start/stop, minimize info}
\label{\detokenize{fix_msst:restart-fix-modify-output-run-start-stop-minimize-info}}
\sphinxAtStartPar
This fix writes the state of all internal variables to {\hyperref[\detokenize{restart::doc}]{\sphinxcrossref{\DUrole{doc}{binary
restart files}}}}.  See the {\hyperref[\detokenize{read_restart::doc}]{\sphinxcrossref{\DUrole{doc}{read\_restart}}}}
command for info on how to re\sphinxhyphen{}specify a fix in an input script that
reads a restart file, so that the operation of the fix continues in an
uninterrupted fashion.

\sphinxAtStartPar
The cumulative energy change in the system imposed by this fix is
included in the {\hyperref[\detokenize{thermo_style::doc}]{\sphinxcrossref{\DUrole{doc}{thermodynamic output}}}} keywords
\sphinxstyleemphasis{ecouple} and \sphinxstyleemphasis{econserve}.  See the {\hyperref[\detokenize{thermo_style::doc}]{\sphinxcrossref{\DUrole{doc}{thermo\_style}}}}
doc page for details.

\sphinxAtStartPar
This fix computes a global scalar which can be accessed by various
{\hyperref[\detokenize{Howto_output::doc}]{\sphinxcrossref{\DUrole{doc}{output commands}}}}.  The scalar is the same
cumulative energy change due to this fix described in the previous
paragraph.  The scalar value calculated by this fix is “extensive”.

\sphinxAtStartPar
The progress of the MSST can be monitored by printing the global
scalar and global vector quantities computed by the fix.

\sphinxAtStartPar
As mentioned above, the scalar is the cumulative energy change due to
the fix.  By monitoring the thermodynamic \sphinxstyleemphasis{econserve} output, this can
be used to test if the MD timestep is sufficiently small for accurate
integration of the dynamic equations.

\sphinxAtStartPar
The global vector contains four values in the following order.  The
vector values output by this fix are “intensive”.

\sphinxAtStartPar
{[}\sphinxstyleemphasis{dhugoniot}, \sphinxstyleemphasis{drayleigh}, \sphinxstyleemphasis{lagrangian\_speed},
\sphinxstyleemphasis{lagrangian\_position}{]}
\begin{enumerate}
\sphinxsetlistlabels{\arabic}{enumi}{enumii}{}{.}%
\item {} 
\sphinxAtStartPar
\sphinxstyleemphasis{dhugoniot} is the departure from the Hugoniot (temperature units).

\item {} 
\sphinxAtStartPar
\sphinxstyleemphasis{drayleigh} is the departure from the Rayleigh line (pressure units).

\item {} 
\sphinxAtStartPar
\sphinxstyleemphasis{lagrangian\_speed} is the laboratory\sphinxhyphen{}frame Lagrangian speed (particle velocity) of the computational cell (velocity units).

\item {} 
\sphinxAtStartPar
\sphinxstyleemphasis{lagrangian\_position} is the computational cell position in the reference frame moving at the shock speed. This is usually a good estimate of distance of the computational cell behind the shock front.

\end{enumerate}

\sphinxAtStartPar
To print these quantities to the log file with descriptive column
headers, the following LAMMPS commands are suggested:

\begin{sphinxVerbatim}[commandchars=\\\{\}]
\PYG{k}{fix              }\PYG{n+nv+nvIdentifier}{msst}\PYG{+w}{ }\PYG{n+nv+nvIdentifier}{all}\PYG{+w}{ }\PYG{n}{msst}\PYG{+w}{ }\PYG{n}{z}
\PYG{k}{variable }\PYG{n+nv+nvIdentifier}{dhug}\PYG{+w}{    }\PYG{n}{equal}\PYG{+w}{ }\PYG{n}{f\PYGZus{}msst[1]}
\PYG{k}{variable }\PYG{n+nv+nvIdentifier}{dray}\PYG{+w}{    }\PYG{n}{equal}\PYG{+w}{ }\PYG{n}{f\PYGZus{}msst[2]}
\PYG{k}{variable }\PYG{n+nv+nvIdentifier}{lgr\PYGZus{}vel}\PYG{+w}{ }\PYG{n}{equal}\PYG{+w}{ }\PYG{n}{f\PYGZus{}msst[3]}
\PYG{k}{variable }\PYG{n+nv+nvIdentifier}{lgr\PYGZus{}pos}\PYG{+w}{ }\PYG{n}{equal}\PYG{+w}{ }\PYG{n}{f\PYGZus{}msst[4]}
\PYG{k}{thermo\PYGZus{}style}\PYG{+w}{     }\PYG{n}{custom}\PYG{+w}{ }\PYG{n}{step}\PYG{+w}{ }\PYG{n}{temp}\PYG{+w}{ }\PYG{n}{ke}\PYG{+w}{ }\PYG{n}{pe}\PYG{+w}{ }\PYG{n}{lz}\PYG{+w}{ }\PYG{n}{pzz}\PYG{+w}{ }\PYG{n}{econserve}\PYG{+w}{ }\PYG{n}{v\PYGZus{}dhug}\PYG{+w}{ }\PYG{n}{v\PYGZus{}dray}\PYG{+w}{ }\PYG{n}{v\PYGZus{}lgr\PYGZus{}vel}\PYG{+w}{ }\PYG{n}{v\PYGZus{}lgr\PYGZus{}pos}\PYG{+w}{ }\PYG{n}{f\PYGZus{}msst}
\end{sphinxVerbatim}


\subsection{Restrictions}
\label{\detokenize{fix_msst:restrictions}}
\sphinxAtStartPar
This fix style is part of the SHOCK package.  It is only enabled if
LAMMPS was built with that package. See the {\hyperref[\detokenize{Build_package::doc}]{\sphinxcrossref{\DUrole{doc}{Build package}}}} page for more info.

\sphinxAtStartPar
All cell dimensions must be periodic. This fix can not be used with a
triclinic cell.  The MSST fix has been tested only for the group\sphinxhyphen{}ID
all.


\subsection{Related commands}
\label{\detokenize{fix_msst:related-commands}}
\sphinxAtStartPar
{\hyperref[\detokenize{fix_nphug::doc}]{\sphinxcrossref{\DUrole{doc}{fix nphug}}}}, {\hyperref[\detokenize{fix_deform::doc}]{\sphinxcrossref{\DUrole{doc}{fix deform}}}}


\subsection{Default}
\label{\detokenize{fix_msst:default}}
\sphinxAtStartPar
The keyword defaults are q = 10, mu = 0, tscale = 0.01, dftb = no,
beta = 0.0.  Note that p0, v0, and e0 are calculated on the first
timestep.


\bigskip\hrule\bigskip

\phantomsection\label{\detokenize{fix_msst:reed}}
\sphinxAtStartPar
\sphinxstylestrong{(Reed)} Reed, Fried, and Joannopoulos, Phys. Rev. Lett., 90, 235503
(2003).

\phantomsection\label{\detokenize{fix_msst:reed2}}
\sphinxAtStartPar
\sphinxstylestrong{(Reed2)} Reed, J. Phys. Chem. C, 116, 2205 (2012).

\phantomsection\label{\detokenize{fix_msst:goldman2}}
\sphinxAtStartPar
\sphinxstylestrong{(Goldman)} Goldman, Srinivasan, Hamel, Fried, Gaus, and Elstner,
J. Phys. Chem. C, 117, 7885 (2013).

\sphinxstepscope

\index{fix mvv/dpd@\spxentry{fix mvv/dpd}}\ignorespaces 
\index{fix mvv/edpd@\spxentry{fix mvv/edpd}}\ignorespaces 
\index{fix mvv/tdpd@\spxentry{fix mvv/tdpd}}\ignorespaces 

\section{fix mvv/dpd command}
\label{\detokenize{fix_mvv_dpd:fix-mvv-dpd-command}}\label{\detokenize{fix_mvv_dpd:index-2}}\label{\detokenize{fix_mvv_dpd:index-1}}\label{\detokenize{fix_mvv_dpd:index-0}}\label{\detokenize{fix_mvv_dpd::doc}}

\section{fix mvv/edpd command}
\label{\detokenize{fix_mvv_dpd:fix-mvv-edpd-command}}

\section{fix mvv/tdpd command}
\label{\detokenize{fix_mvv_dpd:fix-mvv-tdpd-command}}

\subsection{Syntax}
\label{\detokenize{fix_mvv_dpd:syntax}}
\begin{sphinxVerbatim}[commandchars=\\\{\}]
\PYG{k}{fix }\PYG{n+nv+nvIdentifier}{ID}\PYG{+w}{ }\PYG{n+nv+nvIdentifier}{group\PYGZhy{}ID}\PYG{+w}{ }\PYG{n}{mvv}\PYG{o}{/}\PYG{n}{dpd}\PYG{+w}{ }\PYG{n}{lambda}

\PYG{k}{fix }\PYG{n+nv+nvIdentifier}{ID}\PYG{+w}{ }\PYG{n+nv+nvIdentifier}{group\PYGZhy{}ID}\PYG{+w}{ }\PYG{n}{mvv}\PYG{o}{/}\PYG{n}{edpd}\PYG{+w}{ }\PYG{n}{lambda}

\PYG{k}{fix }\PYG{n+nv+nvIdentifier}{ID}\PYG{+w}{ }\PYG{n+nv+nvIdentifier}{group\PYGZhy{}ID}\PYG{+w}{ }\PYG{n}{mvv}\PYG{o}{/}\PYG{n}{tdpd}\PYG{+w}{ }\PYG{n}{lambda}
\end{sphinxVerbatim}
\begin{itemize}
\item {} 
\sphinxAtStartPar
ID, group\sphinxhyphen{}ID are documented in {\hyperref[\detokenize{fix::doc}]{\sphinxcrossref{\DUrole{doc}{fix}}}} command

\item {} 
\sphinxAtStartPar
mvv/dpd, mvv/edpd, mvv/tdpd = style name of this fix command

\item {} 
\sphinxAtStartPar
lambda = (optional) relaxation parameter (unitless)

\end{itemize}


\subsection{Examples}
\label{\detokenize{fix_mvv_dpd:examples}}
\begin{sphinxVerbatim}[commandchars=\\\{\}]
\PYG{k}{fix }\PYG{n+nv+nvIdentifier}{1}\PYG{+w}{ }\PYG{n+nv+nvIdentifier}{all}\PYG{+w}{ }\PYG{n}{mvv}\PYG{o}{/}\PYG{n}{dpd}
\PYG{k}{fix }\PYG{n+nv+nvIdentifier}{1}\PYG{+w}{ }\PYG{n+nv+nvIdentifier}{all}\PYG{+w}{ }\PYG{n}{mvv}\PYG{o}{/}\PYG{n}{dpd}\PYG{+w}{ }\PYG{l+m}{0.5}
\PYG{k}{fix }\PYG{n+nv+nvIdentifier}{1}\PYG{+w}{ }\PYG{n+nv+nvIdentifier}{all}\PYG{+w}{ }\PYG{n}{mvv}\PYG{o}{/}\PYG{n}{edpd}
\PYG{k}{fix }\PYG{n+nv+nvIdentifier}{1}\PYG{+w}{ }\PYG{n+nv+nvIdentifier}{all}\PYG{+w}{ }\PYG{n}{mvv}\PYG{o}{/}\PYG{n}{edpd}\PYG{+w}{ }\PYG{l+m}{0.5}
\PYG{k}{fix }\PYG{n+nv+nvIdentifier}{1}\PYG{+w}{ }\PYG{n+nv+nvIdentifier}{all}\PYG{+w}{ }\PYG{n}{mvv}\PYG{o}{/}\PYG{n}{tdpd}
\PYG{k}{fix }\PYG{n+nv+nvIdentifier}{1}\PYG{+w}{ }\PYG{n+nv+nvIdentifier}{all}\PYG{+w}{ }\PYG{n}{mvv}\PYG{o}{/}\PYG{n}{tdpd}\PYG{+w}{ }\PYG{l+m}{0.5}
\end{sphinxVerbatim}


\subsection{Description}
\label{\detokenize{fix_mvv_dpd:description}}
\sphinxAtStartPar
Perform time integration using the modified velocity\sphinxhyphen{}Verlet (MVV)
algorithm to update position and velocity (fix mvv/dpd), or position,
velocity and temperature (fix mvv/edpd), or position, velocity and
concentration (fix mvv/tdpd) for particles in the group each timestep.

\sphinxAtStartPar
The modified velocity\sphinxhyphen{}Verlet (MVV) algorithm aims to improve the
stability of the time integrator by using an extrapolated version of
the velocity for the force evaluation:
\begin{equation*}
\begin{split}v(t+\frac{\Delta t}{2}) = & v(t) + \frac{\Delta t}{2}\cdot a(t) \\
r(t+\Delta t) = & r(t) + \Delta t\cdot v(t+\frac{\Delta t}{2}) \\
a(t+\Delta t) = & \frac{1}{m}\cdot F\left[ r(t+\Delta t), v(t) +\lambda \cdot \Delta t\cdot a(t)\right] \\
v(t+\Delta t) = & v(t+\frac{\Delta t}{2}) + \frac{\Delta t}{2}\cdot a(t+\Delta t)\end{split}
\end{equation*}
\sphinxAtStartPar
where the parameter \(\lambda\) depends on the
specific choice of DPD parameters, and needs to be tuned on a
case\sphinxhyphen{}by\sphinxhyphen{}case basis.  Specification of a \sphinxstyleemphasis{lambda} value is optional.
If specified, the setting must be from 0.0 to 1.0.  If not specified,
a default value of 0.5 is used, which effectively reproduces the
standard velocity\sphinxhyphen{}Verlet (VV) scheme.  For more details, see
{\hyperref[\detokenize{fix_mvv_dpd:groot2}]{\sphinxcrossref{\DUrole{std,std-ref}{Groot}}}}.

\sphinxAtStartPar
Fix \sphinxstyleemphasis{mvv/dpd} updates the position and velocity of each atom.  It can be
used with the {\hyperref[\detokenize{pair_mesodpd::doc}]{\sphinxcrossref{\DUrole{doc}{pair\_style mdpd}}}} command or other
pair styles such as {\hyperref[\detokenize{pair_dpd::doc}]{\sphinxcrossref{\DUrole{doc}{pair dpd}}}}.

\sphinxAtStartPar
Fix \sphinxstyleemphasis{mvv/edpd} updates the per\sphinxhyphen{}atom temperature, in addition to position
and velocity, and must be used with the {\hyperref[\detokenize{pair_mesodpd::doc}]{\sphinxcrossref{\DUrole{doc}{pair\_style edpd}}}} command.

\sphinxAtStartPar
Fix \sphinxstyleemphasis{mvv/tdpd} updates the per\sphinxhyphen{}atom chemical concentration, in addition
to position and velocity, and must be used with the {\hyperref[\detokenize{pair_mesodpd::doc}]{\sphinxcrossref{\DUrole{doc}{pair\_style
tdpd}}}} command.


\bigskip\hrule\bigskip



\subsection{Restart, fix\_modify, output, run start/stop, minimize info}
\label{\detokenize{fix_mvv_dpd:restart-fix-modify-output-run-start-stop-minimize-info}}
\sphinxAtStartPar
No information about this fix is written to {\hyperref[\detokenize{restart::doc}]{\sphinxcrossref{\DUrole{doc}{binary restart files}}}}.  None of the {\hyperref[\detokenize{fix_modify::doc}]{\sphinxcrossref{\DUrole{doc}{fix\_modify}}}} options are
relevant to this fix.  No global or per\sphinxhyphen{}atom quantities are stored by
this fix for access by various {\hyperref[\detokenize{Howto_output::doc}]{\sphinxcrossref{\DUrole{doc}{output commands}}}}.
No parameter of this fix can be used with the \sphinxstyleemphasis{start/stop} keywords of
the {\hyperref[\detokenize{run::doc}]{\sphinxcrossref{\DUrole{doc}{run}}}} command.  This fix is not invoked during
{\hyperref[\detokenize{minimize::doc}]{\sphinxcrossref{\DUrole{doc}{energy minimization}}}}.


\subsection{Restrictions}
\label{\detokenize{fix_mvv_dpd:restrictions}}
\sphinxAtStartPar
These fixes are part of the DPD\sphinxhyphen{}MESO package. They are only enabled if
LAMMPS was built with that package. See the {\hyperref[\detokenize{Build_package::doc}]{\sphinxcrossref{\DUrole{doc}{Build package}}}} page for more info.

\sphinxAtStartPar
\DUrole{versionmodified,changed}{Changed in version 29Aug2024.}

\sphinxAtStartPar
This fix is incompatible with deformation controls that remap velocity,
for instance the \sphinxstyleemphasis{remap v} option of {\hyperref[\detokenize{fix_deform::doc}]{\sphinxcrossref{\DUrole{doc}{fix deform}}}}.


\subsection{Related commands}
\label{\detokenize{fix_mvv_dpd:related-commands}}
\sphinxAtStartPar
{\hyperref[\detokenize{pair_mesodpd::doc}]{\sphinxcrossref{\DUrole{doc}{pair\_style mdpd}}}}, {\hyperref[\detokenize{pair_mesodpd::doc}]{\sphinxcrossref{\DUrole{doc}{pair\_style edpd}}}},
{\hyperref[\detokenize{pair_mesodpd::doc}]{\sphinxcrossref{\DUrole{doc}{pair\_style tdpd}}}}


\subsection{Default}
\label{\detokenize{fix_mvv_dpd:default}}
\sphinxAtStartPar
The default value for the optional \sphinxstyleemphasis{lambda} parameter is 0.5.


\bigskip\hrule\bigskip

\phantomsection\label{\detokenize{fix_mvv_dpd:groot2}}
\sphinxAtStartPar
\sphinxstylestrong{(Groot)} Groot and Warren, J Chem Phys, 107: 4423\sphinxhyphen{}4435 (1997).  DOI:
10.1063/1.474784

\sphinxstepscope

\index{fix neb@\spxentry{fix neb}}\ignorespaces 

\section{fix neb command}
\label{\detokenize{fix_neb:fix-neb-command}}\label{\detokenize{fix_neb:index-0}}\label{\detokenize{fix_neb::doc}}

\subsection{Syntax}
\label{\detokenize{fix_neb:syntax}}
\begin{sphinxVerbatim}[commandchars=\\\{\}]
\PYG{k}{fix }\PYG{n+nv+nvIdentifier}{ID}\PYG{+w}{ }\PYG{n+nv+nvIdentifier}{group\PYGZhy{}ID}\PYG{+w}{ }\PYG{n}{neb}\PYG{+w}{ }\PYG{n}{Kspring}\PYG{+w}{ }\PYG{n}{keyword}\PYG{+w}{ }\PYG{n}{value}
\end{sphinxVerbatim}
\begin{itemize}
\item {} 
\sphinxAtStartPar
ID, group\sphinxhyphen{}ID are documented in {\hyperref[\detokenize{fix::doc}]{\sphinxcrossref{\DUrole{doc}{fix}}}} command

\item {} 
\sphinxAtStartPar
neb = style name of this fix command

\item {} 
\sphinxAtStartPar
Kspring = spring constant for parallel nudging force (force/distance units or force units, see parallel keyword)

\item {} 
\sphinxAtStartPar
zero or more keyword/value pairs may be appended

\item {} 
\sphinxAtStartPar
keyword = \sphinxstyleemphasis{parallel} or \sphinxstyleemphasis{perp} or \sphinxstyleemphasis{end}

\end{itemize}
\begin{sphinxalltt}
\sphinxstyleemphasis{parallel} value = \sphinxstyleemphasis{neigh} or \sphinxstyleemphasis{ideal} or \sphinxstyleemphasis{equal}
  \sphinxstyleemphasis{neigh} = parallel nudging force based on distance to neighbor replicas (Kspring = force/distance units)
  \sphinxstyleemphasis{ideal} = parallel nudging force based on interpolated ideal position (Kspring = force units)
  \sphinxstyleemphasis{equal} = parallel nudging force based on interpolated ideal position before climbing, then interpolated ideal energy whilst climbing (Kspring = force units)
\sphinxstyleemphasis{perp} value = \sphinxstyleemphasis{Kspring2}
  \sphinxstyleemphasis{Kspring2} = spring constant for perpendicular nudging force (force/distance units)
\sphinxstyleemphasis{end} values = estyle Kspring3
  \sphinxstyleemphasis{estyle} = \sphinxstyleemphasis{first} or \sphinxstyleemphasis{last} or \sphinxstyleemphasis{last/efirst} or \sphinxstyleemphasis{last/efirst/middle}
    \sphinxstyleemphasis{first} = apply force to first replica
    \sphinxstyleemphasis{last} = apply force to last replica
    \sphinxstyleemphasis{last/efirst} = apply force to last replica and set its target energy to that of first replica
    \sphinxstyleemphasis{last/efirst/middle} = same as \sphinxstyleemphasis{last/efirst} plus prevent middle replicas having lower energy than first replica
  \sphinxstyleemphasis{Kspring3} = spring constant for target energy term (1/distance units)
\end{sphinxalltt}


\subsection{Examples}
\label{\detokenize{fix_neb:examples}}
\begin{sphinxVerbatim}[commandchars=\\\{\}]
\PYG{k}{fix }\PYG{n+nv+nvIdentifier}{1}\PYG{+w}{ }\PYG{n+nv+nvIdentifier}{active}\PYG{+w}{ }\PYG{n}{neb}\PYG{+w}{ }\PYG{l+m}{10.0}
\PYG{k}{fix }\PYG{n+nv+nvIdentifier}{2}\PYG{+w}{ }\PYG{n+nv+nvIdentifier}{all}\PYG{+w}{ }\PYG{n}{neb}\PYG{+w}{ }\PYG{l+m}{1.0}\PYG{+w}{ }\PYG{n}{perp}\PYG{+w}{ }\PYG{l+m}{1.0}\PYG{+w}{ }\PYG{n}{end}\PYG{+w}{ }\PYG{n}{last}
\PYG{k}{fix }\PYG{n+nv+nvIdentifier}{2}\PYG{+w}{ }\PYG{n+nv+nvIdentifier}{all}\PYG{+w}{ }\PYG{n}{neb}\PYG{+w}{ }\PYG{l+m}{1.0}\PYG{+w}{ }\PYG{n}{perp}\PYG{+w}{ }\PYG{l+m}{1.0}\PYG{+w}{ }\PYG{n}{end}\PYG{+w}{ }\PYG{n}{first}\PYG{+w}{ }\PYG{l+m}{1.0}\PYG{+w}{ }\PYG{n}{end}\PYG{+w}{ }\PYG{n}{last}\PYG{+w}{ }\PYG{l+m}{1.0}
\PYG{k}{fix }\PYG{n+nv+nvIdentifier}{1}\PYG{+w}{ }\PYG{n+nv+nvIdentifier}{all}\PYG{+w}{ }\PYG{n}{neb}\PYG{+w}{ }\PYG{l+m}{1.0}\PYG{+w}{ }\PYG{n}{parallel}\PYG{+w}{ }\PYG{n}{ideal}\PYG{+w}{ }\PYG{n}{end}\PYG{+w}{ }\PYG{n}{last}\PYG{o}{/}\PYG{n}{efirst}\PYG{+w}{ }\PYG{l+m}{1}
\end{sphinxVerbatim}


\subsection{Description}
\label{\detokenize{fix_neb:description}}
\sphinxAtStartPar
Add nudging forces to atoms in the group for a multi\sphinxhyphen{}replica
simulation run via the {\hyperref[\detokenize{neb::doc}]{\sphinxcrossref{\DUrole{doc}{neb}}}} command to perform a nudged
elastic band (NEB) calculation for finding the transition state.
Hi\sphinxhyphen{}level explanations of NEB are given with the {\hyperref[\detokenize{neb::doc}]{\sphinxcrossref{\DUrole{doc}{neb}}}} command
and on the {\hyperref[\detokenize{Howto_replica::doc}]{\sphinxcrossref{\DUrole{doc}{Howto replica}}}} doc page.  The fix neb
command must be used with the “neb” command and defines how
inter\sphinxhyphen{}replica nudging forces are computed.  A NEB calculation is
divided in two stages. In the first stage n replicas are relaxed
toward a MEP until convergence.  In the second stage, the climbing
image scheme (see {\hyperref[\detokenize{fix_neb:henkelman2}]{\sphinxcrossref{\DUrole{std,std-ref}{(Henkelman2)}}}}) is enabled, so that the
replica having the highest energy relaxes toward the saddle point
(i.e. the point of highest energy along the MEP), and a second
relaxation is performed.

\sphinxAtStartPar
A key purpose of the nudging forces is to keep the replicas equally
spaced.  During the NEB calculation, the \(3N\)\sphinxhyphen{}length vector of
interatomic force \(F_i = -\nabla V\) for each replica \sphinxstyleemphasis{i} is
altered.  For all intermediate replicas (i.e. for \(1 < i < N\),
except the climbing replica) the force vector becomes:
\begin{equation*}
\begin{split}F_i = -\nabla V + (\nabla V \cdot T') T' + F_\parallel + F_\perp\end{split}
\end{equation*}
\sphinxAtStartPar
T’ is the unit “tangent” vector for replica \sphinxstyleemphasis{i} and is a function of
\(R_i, R_{i-1}, R_{i+1}\), and the potential energy of the 3
replicas; it points roughly in the direction of \(R_{i+i} -
R_{i-1}\); see the {\hyperref[\detokenize{fix_neb:henkelman1}]{\sphinxcrossref{\DUrole{std,std-ref}{(Henkelman1)}}}} paper for details.
\(R_i\) are the atomic coordinates of replica \sphinxstyleemphasis{i}; \(R_{i-1}\)
and \(R_{i+1}\) are the coordinates of its neighbor replicas.  The
term \(\nabla V \cdot T'\) is used to remove the component of the
gradient parallel to the path which would tend to distribute the replica
unevenly along the path.  \(F_\parallel\) is an artificial nudging
force which is applied only in the tangent direction and which maintains
the equal spacing between replicas (see below for more information).
\(F_\perp\) is an optional artificial spring which is applied in a
direction perpendicular to the tangent direction and which prevent the
paths from forming acute kinks (see below for more information).

\sphinxAtStartPar
In the second stage of the NEB calculation, the interatomic force \(F_i\)
for the climbing replica (the replica of highest energy after the
first stage) is changed to:
\begin{equation*}
\begin{split}F_i = -\nabla V + 2 (\nabla V \cdot T') T' + F_\perp\end{split}
\end{equation*}
\sphinxAtStartPar
and the relaxation procedure is continued to a new converged MEP.


\bigskip\hrule\bigskip


\sphinxAtStartPar
The keyword \sphinxstyleemphasis{parallel} specifies how the parallel nudging force is
computed.  With a value of \sphinxstyleemphasis{neigh}, the parallel nudging force is
computed as in {\hyperref[\detokenize{fix_neb:henkelman1}]{\sphinxcrossref{\DUrole{std,std-ref}{(Henkelman1)}}}} by connecting each
intermediate replica with the previous and the next image:
\begin{equation*}
\begin{split}F_\parallel = Kspring \cdot \left(\left|R_{i+1} - R_i\right| - \left|R_i - R_{i-1}\right|\right)\end{split}
\end{equation*}
\sphinxAtStartPar
Note that in this case the specified \sphinxstyleemphasis{Kspring} is in
force/distance units.

\sphinxAtStartPar
With a value of \sphinxstyleemphasis{ideal}, the spring force is computed as suggested in
{\hyperref[\detokenize{fix_neb:weinane}]{\sphinxcrossref{\DUrole{std,std-ref}{(WeinanE)}}}}
\begin{equation*}
\begin{split}F_\parallel = -Kspring \cdot (RD - RD_{ideal}) / (2 \cdot meanDist)\end{split}
\end{equation*}
\sphinxAtStartPar
where \sphinxstyleemphasis{RD} is the “reaction coordinate” see {\hyperref[\detokenize{neb::doc}]{\sphinxcrossref{\DUrole{doc}{neb}}}} section,
and \(RD_{ideal}\) is the ideal \sphinxstyleemphasis{RD} for which all the images are
equally spaced.  I.e. \(RD_{ideal} = (i-1) \cdot meanDist\) when the
climbing replica is off, where \sphinxstyleemphasis{i} is the replica number).  The
\sphinxstyleemphasis{meanDist} is the average distance between replicas.  Note that in this
case the specified \sphinxstyleemphasis{Kspring} is in force units. When the climbing
replica is on, \(RD_{ideal}\) and \(meanDist\) are calculated
separately each side of the climbing image. Note that the \sphinxstyleemphasis{ideal} form
of nudging can often be more effective at keeping the replicas equally
spaced before climbing, then equally spaced either side of the climbing
image whilst climbing.

\sphinxAtStartPar
With a value of \sphinxstyleemphasis{equal} the spring force is computed as for \sphinxstyleemphasis{ideal} when
the climbing replica is off, promoting equidistance. When the climbing
replica is on, the spring force is computed to promote equidistant
absolute differences in energy, rather than distance, each side of the
climbing image:
\begin{equation*}
\begin{split}F_\parallel = -Kspring \cdot (ED - ED_{ideal}) / (2 \cdot meanEDist)\end{split}
\end{equation*}
\sphinxAtStartPar
where \sphinxstyleemphasis{ED} is the cumulative sum of absolute energy differences:
\begin{equation*}
\begin{split}ED = \sum_{i<N} \left|E(R_{i+1}) - E(R_i)\right|,\end{split}
\end{equation*}
\sphinxAtStartPar
\sphinxstyleemphasis{meanEdist} is the average absolute energy difference between replicas
up to the climbing image or from the climbing image to the final image,
for images before or after the climbing image
respectively. \(ED_{ideal}\) is the corresponding cumulative sum of
average absolute energy differences in each case, in close analogy to
\sphinxstyleemphasis{ideal}. This form of nudging is to aid schemes which integrate forces
along, or near to, NEB pathways such as {\hyperref[\detokenize{fix_pafi::doc}]{\sphinxcrossref{\DUrole{doc}{fix\_pafi}}}}.


\bigskip\hrule\bigskip


\sphinxAtStartPar
The keyword \sphinxstyleemphasis{perp} specifies if and how a perpendicular nudging force is
computed.  It adds a spring force perpendicular to the path in order to
prevent the path from becoming too strongly kinked.  It can
significantly improve the convergence of the NEB calculation when the
resolution is poor.  I.e. when few replicas are used; see {\hyperref[\detokenize{fix_neb:maras1}]{\sphinxcrossref{\DUrole{std,std-ref}{(Maras)}}}} for details.

\sphinxAtStartPar
The perpendicular spring force is given by
\begin{equation*}
\begin{split}F_\perp = K_{spring2} \cdot F(R_{i-1},R_i,R_{i+1}) (R_{i+1} + R_{i-1} - 2 R_i)\end{split}
\end{equation*}
\sphinxAtStartPar
where \sphinxstyleemphasis{Kspring2} is the specified value.  \(F(R_{i-1}, R_i,
R_{i+1})\) is a smooth scalar function of the angle \(R_{i-1} R_i
R_{i+1}\).  It is equal to 0.0 when the path is straight and is equal to
1 when the angle \(R_{i-1} R_i R_{i+1}\) is acute.
\(F(R_{i-1}, R_i, R_{i+1})\) is defined in {\hyperref[\detokenize{fix_neb:jonsson}]{\sphinxcrossref{\DUrole{std,std-ref}{(Jonsson)}}}}.

\sphinxAtStartPar
If \sphinxstyleemphasis{Kspring2} is set to 0.0 (the default) then no perpendicular spring
force is added.


\bigskip\hrule\bigskip


\sphinxAtStartPar
By default, no additional forces act on the first and last replicas
during the NEB relaxation, so these replicas simply relax toward their
respective local minima.  By using the key word \sphinxstyleemphasis{end}, additional forces
can be applied to the first and/or last replicas, to enable them to
relax toward a MEP while constraining their energy E to the target
energy ETarget.

\sphinxAtStartPar
If \(E_{Target} > E\), the interatomic force \(F_i\) for the
specified replica becomes:
\begin{equation*}
\begin{split}F_i & = -\nabla V + (\nabla V \cdot T' + (E - E_{Target}) \cdot K_{spring3}) T', \qquad  \textrm{when} \quad \nabla V \cdot T' < 0 \\
F_i & = -\nabla V + (\nabla V \cdot T' + (E_{Target} - E) \cdot K_{spring3}) T', \qquad \textrm{when} \quad \nabla V  \cdot T' > 0\end{split}
\end{equation*}
\sphinxAtStartPar
The “spring” constant on the difference in energies is the specified
\sphinxstyleemphasis{Kspring3} value.

\sphinxAtStartPar
When \sphinxstyleemphasis{estyle} is specified as \sphinxstyleemphasis{first}, the force is applied to the first
replica.  When \sphinxstyleemphasis{estyle} is specified as \sphinxstyleemphasis{last}, the force is applied to
the last replica.  Note that the \sphinxstyleemphasis{end} keyword can be used twice to add
forces to both the first and last replicas.

\sphinxAtStartPar
For both these \sphinxstyleemphasis{estyle} settings, the target energy \sphinxstyleemphasis{ETarget} is set
to the initial energy of the replica (at the start of the NEB
calculation).

\sphinxAtStartPar
If the \sphinxstyleemphasis{estyle} is specified as \sphinxstyleemphasis{last/efirst} or \sphinxstyleemphasis{last/efirst/middle},
force is applied to the last replica, but the target energy \sphinxstyleemphasis{ETarget} is
continuously set to the energy of the first replica, as it evolves
during the NEB relaxation.

\sphinxAtStartPar
The difference between these two \sphinxstyleemphasis{estyle} options is as follows.  When
\sphinxstyleemphasis{estyle} is specified as \sphinxstyleemphasis{last/efirst}, no change is made to the
inter\sphinxhyphen{}replica force applied to the intermediate replicas (neither first
or last).  If the initial path is too far from the MEP, an intermediate
replica may relax “faster” and reach a lower energy than the last
replica.  In this case the intermediate replica will be relaxing toward
its own local minima.  This behavior can be prevented by specifying
\sphinxstyleemphasis{estyle} as \sphinxstyleemphasis{last/efirst/middle} which will alter the inter\sphinxhyphen{}replica
force applied to intermediate replicas by removing the contribution of
the gradient to the inter\sphinxhyphen{}replica force.  This will only be done if a
particular intermediate replica has a lower energy than the first
replica.  This should effectively prevent the intermediate replicas from
over\sphinxhyphen{}relaxing.

\sphinxAtStartPar
After converging a NEB calculation using an \sphinxstyleemphasis{estyle} of
\sphinxstyleemphasis{last/efirst/middle}, you should check that all intermediate replicas
have a larger energy than the first replica. If this is not the case,
the path is probably not a MEP.

\sphinxAtStartPar
Finally, note that the last replica may never reach the target energy
if it is stuck in a local minima which has a larger energy than the
target energy.


\subsection{Restart, fix\_modify, output, run start/stop, minimize info}
\label{\detokenize{fix_neb:restart-fix-modify-output-run-start-stop-minimize-info}}
\sphinxAtStartPar
No information about this fix is written to {\hyperref[\detokenize{restart::doc}]{\sphinxcrossref{\DUrole{doc}{binary restart files}}}}.  None of the {\hyperref[\detokenize{fix_modify::doc}]{\sphinxcrossref{\DUrole{doc}{fix\_modify}}}} options are
relevant to this fix.  No global or per\sphinxhyphen{}atom quantities are stored by
this fix for access by various {\hyperref[\detokenize{Howto_output::doc}]{\sphinxcrossref{\DUrole{doc}{output commands}}}}.
No parameter of this fix can be used with the \sphinxstyleemphasis{start/stop} keywords of
the {\hyperref[\detokenize{run::doc}]{\sphinxcrossref{\DUrole{doc}{run}}}} command.

\sphinxAtStartPar
The forces due to this fix are imposed during an energy minimization,
as invoked by the {\hyperref[\detokenize{minimize::doc}]{\sphinxcrossref{\DUrole{doc}{minimize}}}} command via the
{\hyperref[\detokenize{neb::doc}]{\sphinxcrossref{\DUrole{doc}{neb}}}} command.


\subsection{Restrictions}
\label{\detokenize{fix_neb:restrictions}}
\sphinxAtStartPar
This command can only be used if LAMMPS was built with the REPLICA
package.  See the {\hyperref[\detokenize{Build_package::doc}]{\sphinxcrossref{\DUrole{doc}{Build package}}}} doc
page for more info.


\subsection{Related commands}
\label{\detokenize{fix_neb:related-commands}}
\sphinxAtStartPar
{\hyperref[\detokenize{neb::doc}]{\sphinxcrossref{\DUrole{doc}{neb}}}}


\subsection{Default}
\label{\detokenize{fix_neb:default}}
\sphinxAtStartPar
The option defaults are parallel = neigh, perp = 0.0, ends is not
specified (no inter\sphinxhyphen{}replica force on the end replicas).


\bigskip\hrule\bigskip

\phantomsection\label{\detokenize{fix_neb:henkelman1}}
\sphinxAtStartPar
\sphinxstylestrong{(Henkelman1)} Henkelman and Jonsson, J Chem Phys, 113, 9978\sphinxhyphen{}9985 (2000).

\phantomsection\label{\detokenize{fix_neb:henkelman2}}
\sphinxAtStartPar
\sphinxstylestrong{(Henkelman2)} Henkelman, Uberuaga, Jonsson, J Chem Phys, 113,
9901\sphinxhyphen{}9904 (2000).

\phantomsection\label{\detokenize{fix_neb:weinane}}
\sphinxAtStartPar
\sphinxstylestrong{(WeinanE)} E, Ren, Vanden\sphinxhyphen{}Eijnden, Phys Rev B, 66, 052301 (2002).

\phantomsection\label{\detokenize{fix_neb:jonsson}}
\sphinxAtStartPar
\sphinxstylestrong{(Jonsson)} Jonsson, Mills and Jacobsen, in Classical and Quantum
Dynamics in Condensed Phase Simulations, edited by Berne, Ciccotti,
and Coker World Scientific, Singapore, 1998, p 385.

\phantomsection\label{\detokenize{fix_neb:maras1}}
\sphinxAtStartPar
\sphinxstylestrong{(Maras)} Maras, Trushin, Stukowski, Ala\sphinxhyphen{}Nissila, Jonsson,
Comp Phys Comm, 205, 13\sphinxhyphen{}21 (2016).

\sphinxstepscope

\index{fix neb/spin@\spxentry{fix neb/spin}}\ignorespaces 

\section{fix neb/spin command}
\label{\detokenize{fix_neb_spin:fix-neb-spin-command}}\label{\detokenize{fix_neb_spin:index-0}}\label{\detokenize{fix_neb_spin::doc}}

\subsection{Syntax}
\label{\detokenize{fix_neb_spin:syntax}}
\begin{sphinxVerbatim}[commandchars=\\\{\}]
\PYG{k}{fix }\PYG{n+nv+nvIdentifier}{ID}\PYG{+w}{ }\PYG{n+nv+nvIdentifier}{group\PYGZhy{}ID}\PYG{+w}{ }\PYG{n}{neb}\PYG{o}{/}\PYG{n}{spin}\PYG{+w}{ }\PYG{n}{Kspring}
\end{sphinxVerbatim}
\begin{itemize}
\item {} 
\sphinxAtStartPar
ID, group\sphinxhyphen{}ID are documented in {\hyperref[\detokenize{fix::doc}]{\sphinxcrossref{\DUrole{doc}{fix}}}} command

\item {} 
\sphinxAtStartPar
neb/spin = style name of this fix command

\end{itemize}

\begin{sphinxVerbatim}[commandchars=\\\{\}]
Kspring = spring constant for parallel nudging force
(force/distance units or force units, see parallel keyword)
\end{sphinxVerbatim}


\subsection{Examples}
\label{\detokenize{fix_neb_spin:examples}}
\begin{sphinxVerbatim}[commandchars=\\\{\}]
\PYG{k}{fix }\PYG{n+nv+nvIdentifier}{1}\PYG{+w}{ }\PYG{n+nv+nvIdentifier}{active}\PYG{+w}{ }\PYG{n}{neb}\PYG{o}{/}\PYG{n}{spin}\PYG{+w}{ }\PYG{l+m}{1.0}
\end{sphinxVerbatim}


\subsection{Description}
\label{\detokenize{fix_neb_spin:description}}
\sphinxAtStartPar
Add nudging forces to spins in the group for a multi\sphinxhyphen{}replica
simulation run via the {\hyperref[\detokenize{neb_spin::doc}]{\sphinxcrossref{\DUrole{doc}{neb/spin}}}} command to perform a
geodesic nudged elastic band (GNEB) calculation for finding the
transition state.
Hi\sphinxhyphen{}level explanations of GNEB are given with the
{\hyperref[\detokenize{neb_spin::doc}]{\sphinxcrossref{\DUrole{doc}{neb/spin}}}} command and on the
{\hyperref[\detokenize{Howto_replica::doc}]{\sphinxcrossref{\DUrole{doc}{Howto replica}}}} doc page.
The fix neb/spin command must be used with the “neb/spin” command and
defines how inter\sphinxhyphen{}replica nudging forces are computed.  A GNEB
calculation is divided in two stages. In the first stage n replicas
are relaxed toward a MEP until convergence.  In the second stage, the
climbing image scheme is enabled, so that the replica having the highest
energy relaxes toward the saddle point (i.e. the point of highest energy
along the MEP), and a second relaxation is performed.

\sphinxAtStartPar
The nudging forces are calculated as explained in
{\hyperref[\detokenize{fix_neb_spin:bessarabb}]{\sphinxcrossref{\DUrole{std,std-ref}{(Bessarab)}}}}).
See this reference for more explanation about their expression.


\subsection{Restart, fix\_modify, output, run start/stop, minimize info}
\label{\detokenize{fix_neb_spin:restart-fix-modify-output-run-start-stop-minimize-info}}
\sphinxAtStartPar
No information about this fix is written to {\hyperref[\detokenize{restart::doc}]{\sphinxcrossref{\DUrole{doc}{binary restart files}}}}.  None of the {\hyperref[\detokenize{fix_modify::doc}]{\sphinxcrossref{\DUrole{doc}{fix\_modify}}}} options
are relevant to this fix.  No global or per\sphinxhyphen{}atom quantities are stored
by this fix for access by various {\hyperref[\detokenize{Howto_output::doc}]{\sphinxcrossref{\DUrole{doc}{output commands}}}}.
No parameter of this fix can be used with the \sphinxstyleemphasis{start/stop} keywords of
the {\hyperref[\detokenize{run::doc}]{\sphinxcrossref{\DUrole{doc}{run}}}} command.

\sphinxAtStartPar
The forces due to this fix are imposed during an energy minimization,
as invoked by the {\hyperref[\detokenize{minimize::doc}]{\sphinxcrossref{\DUrole{doc}{minimize}}}} command via the
{\hyperref[\detokenize{neb_spin::doc}]{\sphinxcrossref{\DUrole{doc}{neb/spin}}}} command.


\subsection{Restrictions}
\label{\detokenize{fix_neb_spin:restrictions}}
\sphinxAtStartPar
This command can only be used if LAMMPS was built with the SPIN
package.  See the {\hyperref[\detokenize{Build_package::doc}]{\sphinxcrossref{\DUrole{doc}{Build package}}}} doc
page for more info.


\subsection{Related commands}
\label{\detokenize{fix_neb_spin:related-commands}}
\sphinxAtStartPar
{\hyperref[\detokenize{neb_spin::doc}]{\sphinxcrossref{\DUrole{doc}{neb\_spin}}}}


\subsection{Default}
\label{\detokenize{fix_neb_spin:default}}
\sphinxAtStartPar
none


\bigskip\hrule\bigskip

\phantomsection\label{\detokenize{fix_neb_spin:bessarabb}}
\sphinxAtStartPar
\sphinxstylestrong{(Bessarab)} Bessarab, Uzdin, Jonsson, Comp Phys Comm, 196,
335\sphinxhyphen{}347 (2015).

\sphinxstepscope

\index{fix nvt@\spxentry{fix nvt}}\ignorespaces 
\index{fix nvt/gpu@\spxentry{fix nvt/gpu}}\ignorespaces 
\index{fix nvt/intel@\spxentry{fix nvt/intel}}\ignorespaces 
\index{fix nvt/kk@\spxentry{fix nvt/kk}}\ignorespaces 
\index{fix nvt/omp@\spxentry{fix nvt/omp}}\ignorespaces 
\index{fix npt@\spxentry{fix npt}}\ignorespaces 
\index{fix npt/gpu@\spxentry{fix npt/gpu}}\ignorespaces 
\index{fix npt/intel@\spxentry{fix npt/intel}}\ignorespaces 
\index{fix npt/kk@\spxentry{fix npt/kk}}\ignorespaces 
\index{fix npt/omp@\spxentry{fix npt/omp}}\ignorespaces 
\index{fix nph@\spxentry{fix nph}}\ignorespaces 
\index{fix nph/kk@\spxentry{fix nph/kk}}\ignorespaces 
\index{fix nph/omp@\spxentry{fix nph/omp}}\ignorespaces 

\section{fix nvt command}
\label{\detokenize{fix_nh:fix-nvt-command}}\label{\detokenize{fix_nh:index-12}}\label{\detokenize{fix_nh:index-11}}\label{\detokenize{fix_nh:index-10}}\label{\detokenize{fix_nh:index-9}}\label{\detokenize{fix_nh:index-8}}\label{\detokenize{fix_nh:index-7}}\label{\detokenize{fix_nh:index-6}}\label{\detokenize{fix_nh:index-5}}\label{\detokenize{fix_nh:index-4}}\label{\detokenize{fix_nh:index-3}}\label{\detokenize{fix_nh:index-2}}\label{\detokenize{fix_nh:index-1}}\label{\detokenize{fix_nh:index-0}}\label{\detokenize{fix_nh::doc}}
\sphinxAtStartPar
Accelerator Variants: \sphinxstyleemphasis{nvt/gpu}, \sphinxstyleemphasis{nvt/intel}, \sphinxstyleemphasis{nvt/kk}, \sphinxstyleemphasis{nvt/omp}


\section{fix npt command}
\label{\detokenize{fix_nh:fix-npt-command}}
\sphinxAtStartPar
Accelerator Variants: \sphinxstyleemphasis{npt/gpu}, \sphinxstyleemphasis{npt/intel}, \sphinxstyleemphasis{npt/kk}, \sphinxstyleemphasis{npt/omp}


\section{fix nph command}
\label{\detokenize{fix_nh:fix-nph-command}}
\sphinxAtStartPar
Accelerator Variants: \sphinxstyleemphasis{nph/kk}, \sphinxstyleemphasis{nph/omp}


\subsection{Syntax}
\label{\detokenize{fix_nh:syntax}}
\begin{sphinxVerbatim}[commandchars=\\\{\}]
\PYG{k}{fix }\PYG{n+nv+nvIdentifier}{ID}\PYG{+w}{ }\PYG{n+nv+nvIdentifier}{group\PYGZhy{}ID}\PYG{+w}{ }\PYG{n}{style\PYGZus{}name}\PYG{+w}{ }\PYG{n}{keyword}\PYG{+w}{ }\PYG{n}{value}\PYG{+w}{ }\PYG{n}{...}
\end{sphinxVerbatim}
\begin{itemize}
\item {} 
\sphinxAtStartPar
ID, group\sphinxhyphen{}ID are documented in {\hyperref[\detokenize{fix::doc}]{\sphinxcrossref{\DUrole{doc}{fix}}}} command

\item {} 
\sphinxAtStartPar
style\_name = \sphinxstyleemphasis{nvt} or \sphinxstyleemphasis{npt} or \sphinxstyleemphasis{nph}

\item {} 
\sphinxAtStartPar
one or more keyword/value pairs may be appended
\begin{sphinxalltt}
keyword = \sphinxstyleemphasis{temp} or \sphinxstyleemphasis{iso} or \sphinxstyleemphasis{aniso} or \sphinxstyleemphasis{tri} or \sphinxstyleemphasis{x} or \sphinxstyleemphasis{y} or \sphinxstyleemphasis{z} or \sphinxstyleemphasis{xy} or \sphinxstyleemphasis{yz} or \sphinxstyleemphasis{xz} or \sphinxstyleemphasis{couple} or \sphinxstyleemphasis{tchain} or \sphinxstyleemphasis{pchain} or \sphinxstyleemphasis{mtk} or \sphinxstyleemphasis{tloop} or \sphinxstyleemphasis{ploop} or \sphinxstyleemphasis{nreset} or \sphinxstyleemphasis{drag} or \sphinxstyleemphasis{ptemp} or \sphinxstyleemphasis{dilate} or \sphinxstyleemphasis{scalexy} or \sphinxstyleemphasis{scaleyz} or \sphinxstyleemphasis{scalexz} or \sphinxstyleemphasis{flip} or \sphinxstyleemphasis{fixedpoint} or \sphinxstyleemphasis{update}
  \sphinxstyleemphasis{temp} values = Tstart Tstop Tdamp
    Tstart,Tstop = external temperature at start/end of run
    Tdamp = temperature damping parameter (time units)
  \sphinxstyleemphasis{iso} or \sphinxstyleemphasis{aniso} or \sphinxstyleemphasis{tri} values = Pstart Pstop Pdamp
    Pstart,Pstop = scalar external pressure at start/end of run (pressure units)
    Pdamp = pressure damping parameter (time units)
  \sphinxstyleemphasis{x} or \sphinxstyleemphasis{y} or \sphinxstyleemphasis{z} or \sphinxstyleemphasis{xy} or \sphinxstyleemphasis{yz} or \sphinxstyleemphasis{xz} values = Pstart Pstop Pdamp
    Pstart,Pstop = external stress tensor component at start/end of run (pressure units)
    Pdamp = stress damping parameter (time units)
  \sphinxstyleemphasis{couple} = \sphinxstyleemphasis{none} or \sphinxstyleemphasis{xyz} or \sphinxstyleemphasis{xy} or \sphinxstyleemphasis{yz} or \sphinxstyleemphasis{xz}
  \sphinxstyleemphasis{tchain} value = N
    N = length of thermostat chain (1 = single thermostat)
  \sphinxstyleemphasis{pchain} value = N
    N length of thermostat chain on barostat (0 = no thermostat)
  \sphinxstyleemphasis{mtk} value = \sphinxstyleemphasis{yes} or \sphinxstyleemphasis{no} = add in MTK adjustment term or not
  \sphinxstyleemphasis{tloop} value = M
    M = number of sub\sphinxhyphen{}cycles to perform on thermostat
  \sphinxstyleemphasis{ploop} value = M
    M = number of sub\sphinxhyphen{}cycles to perform on barostat thermostat
  \sphinxstyleemphasis{nreset} value = reset reference cell every this many timesteps
  \sphinxstyleemphasis{drag} value = Df
    Df = drag factor added to barostat/thermostat (0.0 = no drag)
  \sphinxstyleemphasis{ptemp} value = Ttarget
    Ttarget = target temperature for barostat
  \sphinxstyleemphasis{dilate} value = dilate\sphinxhyphen{}group\sphinxhyphen{}ID
    dilate\sphinxhyphen{}group\sphinxhyphen{}ID = only dilate atoms in this group due to barostat volume changes
  \sphinxstyleemphasis{scalexy} value = \sphinxstyleemphasis{yes} or \sphinxstyleemphasis{no} = scale xy with ly
  \sphinxstyleemphasis{scaleyz} value = \sphinxstyleemphasis{yes} or \sphinxstyleemphasis{no} = scale yz with lz
  \sphinxstyleemphasis{scalexz} value = \sphinxstyleemphasis{yes} or \sphinxstyleemphasis{no} = scale xz with lz
  \sphinxstyleemphasis{flip} value = \sphinxstyleemphasis{yes} or \sphinxstyleemphasis{no} = allow or disallow box flips when it becomes highly skewed
  \sphinxstyleemphasis{fixedpoint} values = x y z
    x,y,z = perform barostat dilation/contraction around this point (distance units)
  \sphinxstyleemphasis{update} value = \sphinxstyleemphasis{dipole} or \sphinxstyleemphasis{dipole/dlm}
    dipole = update dipole orientation (only for sphere variants)
    dipole/dlm = use DLM integrator to update dipole orientation (only for sphere variants)
\end{sphinxalltt}

\end{itemize}


\subsection{Examples}
\label{\detokenize{fix_nh:examples}}
\begin{sphinxVerbatim}[commandchars=\\\{\}]
\PYG{k}{fix }\PYG{n+nv+nvIdentifier}{1}\PYG{+w}{ }\PYG{n+nv+nvIdentifier}{all}\PYG{+w}{ }\PYG{n}{nvt}\PYG{+w}{ }\PYG{n}{temp}\PYG{+w}{ }\PYG{l+m}{300.0}\PYG{+w}{ }\PYG{l+m}{300.0}\PYG{+w}{ }\PYG{l+m}{100.0}
\PYG{k}{fix }\PYG{n+nv+nvIdentifier}{1}\PYG{+w}{ }\PYG{n+nv+nvIdentifier}{water}\PYG{+w}{ }\PYG{n}{npt}\PYG{+w}{ }\PYG{n}{temp}\PYG{+w}{ }\PYG{l+m}{300.0}\PYG{+w}{ }\PYG{l+m}{300.0}\PYG{+w}{ }\PYG{l+m}{100.0}\PYG{+w}{ }\PYG{n}{iso}\PYG{+w}{ }\PYG{l+m}{0.0}\PYG{+w}{ }\PYG{l+m}{0.0}\PYG{+w}{ }\PYG{l+m}{1000.0}
\PYG{k}{fix }\PYG{n+nv+nvIdentifier}{2}\PYG{+w}{ }\PYG{n+nv+nvIdentifier}{jello}\PYG{+w}{ }\PYG{n}{npt}\PYG{+w}{ }\PYG{n}{temp}\PYG{+w}{ }\PYG{l+m}{300.0}\PYG{+w}{ }\PYG{l+m}{300.0}\PYG{+w}{ }\PYG{l+m}{100.0}\PYG{+w}{ }\PYG{n}{tri}\PYG{+w}{ }\PYG{l+m}{5.0}\PYG{+w}{ }\PYG{l+m}{5.0}\PYG{+w}{ }\PYG{l+m}{1000.0}
\PYG{k}{fix }\PYG{n+nv+nvIdentifier}{2}\PYG{+w}{ }\PYG{n+nv+nvIdentifier}{ice}\PYG{+w}{ }\PYG{n}{nph}\PYG{+w}{ }\PYG{n}{x}\PYG{+w}{ }\PYG{l+m}{1.0}\PYG{+w}{ }\PYG{l+m}{1.0}\PYG{+w}{ }\PYG{l+m}{0.5}\PYG{+w}{ }\PYG{n}{y}\PYG{+w}{ }\PYG{l+m}{2.0}\PYG{+w}{ }\PYG{l+m}{2.0}\PYG{+w}{ }\PYG{l+m}{0.5}\PYG{+w}{ }\PYG{n}{z}\PYG{+w}{ }\PYG{l+m}{3.0}\PYG{+w}{ }\PYG{l+m}{3.0}\PYG{+w}{ }\PYG{l+m}{0.5}\PYG{+w}{ }\PYG{n}{yz}\PYG{+w}{ }\PYG{l+m}{0.1}\PYG{+w}{ }\PYG{l+m}{0.1}\PYG{+w}{ }\PYG{l+m}{0.5}\PYG{+w}{ }\PYG{n}{xz}\PYG{+w}{ }\PYG{l+m}{0.2}\PYG{+w}{ }\PYG{l+m}{0.2}\PYG{+w}{ }\PYG{l+m}{0.5}\PYG{+w}{ }\PYG{n}{xy}\PYG{+w}{ }\PYG{l+m}{0.3}\PYG{+w}{ }\PYG{l+m}{0.3}\PYG{+w}{ }\PYG{l+m}{0.5}\PYG{+w}{ }\PYG{n}{nreset}\PYG{+w}{ }\PYG{l+m}{1000}
\end{sphinxVerbatim}


\subsection{Description}
\label{\detokenize{fix_nh:description}}
\sphinxAtStartPar
These commands perform time integration on Nose\sphinxhyphen{}Hoover style
non\sphinxhyphen{}Hamiltonian equations of motion which are designed to generate
positions and velocities sampled from the canonical (nvt),
isothermal\sphinxhyphen{}isobaric (npt), and isenthalpic (nph) ensembles.  This
updates the position and velocity for atoms in the group each
timestep.

\sphinxAtStartPar
The thermostatting and barostatting is achieved by adding some dynamic
variables which are coupled to the particle velocities
(thermostatting) and simulation domain dimensions (barostatting).  In
addition to basic thermostatting and barostatting, these fixes can
also create a chain of thermostats coupled to the particle thermostat,
and another chain of thermostats coupled to the barostat
variables. The barostat can be coupled to the overall box volume, or
to individual dimensions, including the \sphinxstyleemphasis{xy}, \sphinxstyleemphasis{xz} and \sphinxstyleemphasis{yz} tilt
dimensions. The external pressure of the barostat can be specified as
either a scalar pressure (isobaric ensemble) or as components of a
symmetric stress tensor (constant stress ensemble).  When used
correctly, the time\sphinxhyphen{}averaged temperature and stress tensor of the
particles will match the target values specified by Tstart/Tstop and
Pstart/Pstop.

\sphinxAtStartPar
The equations of motion used are those of Shinoda et al in
{\hyperref[\detokenize{fix_nh:nh-shinoda}]{\sphinxcrossref{\DUrole{std,std-ref}{(Shinoda)}}}}, which combine the hydrostatic equations of
Martyna, Tobias and Klein in {\hyperref[\detokenize{fix_nh:nh-martyna}]{\sphinxcrossref{\DUrole{std,std-ref}{(Martyna)}}}} with the strain
energy proposed by Parrinello and Rahman in
{\hyperref[\detokenize{fix_nh:nh-parrinello}]{\sphinxcrossref{\DUrole{std,std-ref}{(Parrinello)}}}}.  The time integration schemes closely
follow the time\sphinxhyphen{}reversible measure\sphinxhyphen{}preserving Verlet and rRESPA
integrators derived by Tuckerman et al in {\hyperref[\detokenize{fix_nh:nh-tuckerman}]{\sphinxcrossref{\DUrole{std,std-ref}{(Tuckerman)}}}}.


\bigskip\hrule\bigskip


\sphinxAtStartPar
The thermostat parameters for fix styles \sphinxstyleemphasis{nvt} and \sphinxstyleemphasis{npt} are specified
using the \sphinxstyleemphasis{temp} keyword.  Other thermostat\sphinxhyphen{}related keywords are
\sphinxstyleemphasis{tchain}, \sphinxstyleemphasis{tloop} and \sphinxstyleemphasis{drag}, which are discussed below.

\sphinxAtStartPar
The thermostat is applied to only the translational degrees of freedom
for the particles.  The translational degrees of freedom can also have
a bias velocity removed before thermostatting takes place; see the
description below.  The desired temperature at each timestep is a
ramped value during the run from \sphinxstyleemphasis{Tstart} to \sphinxstyleemphasis{Tstop}.  The \sphinxstyleemphasis{Tdamp}
parameter is specified in time units and determines how rapidly the
temperature is relaxed.  For example, a value of 10.0 means to relax
the temperature in a timespan of (roughly) 10 time units (e.g. \(\tau\)
or fs or ps \sphinxhyphen{} see the {\hyperref[\detokenize{units::doc}]{\sphinxcrossref{\DUrole{doc}{units}}}} command).  The atoms in the
fix group are the only ones whose velocities and positions are updated
by the velocity/position update portion of the integration.

\begin{sphinxadmonition}{note}{Note:}
\sphinxAtStartPar
A Nose\sphinxhyphen{}Hoover thermostat will not work well for arbitrary values
of \sphinxstyleemphasis{Tdamp}.  If \sphinxstyleemphasis{Tdamp} is too small, the temperature can fluctuate
wildly; if it is too large, the temperature will take a very long time
to equilibrate.  A good choice for many models is a \sphinxstyleemphasis{Tdamp} of around
100 timesteps.  Note that this is NOT the same as 100 time units for
most {\hyperref[\detokenize{units::doc}]{\sphinxcrossref{\DUrole{doc}{units}}}} settings. A simple way to ensure this, is
via using an {\hyperref[\detokenize{variable::doc}]{\sphinxcrossref{\DUrole{doc}{immediate variable}}}} expression accessing
the thermo property ‘dt’, which is the length of the time step. Example:
\end{sphinxadmonition}

\begin{sphinxVerbatim}[commandchars=\\\{\}]
\PYG{k}{fix }\PYG{n+nv+nvIdentifier}{1}\PYG{+w}{ }\PYG{n+nv+nvIdentifier}{all}\PYG{+w}{ }\PYG{n}{nvt}\PYG{+w}{ }\PYG{n}{temp}\PYG{+w}{ }\PYG{l+m}{300.0}\PYG{+w}{ }\PYG{l+m}{300.0}\PYG{+w}{ }\PYG{n+nv}{\PYGZdl{}(}\PYG{n+nv}{100.0*dt}\PYG{n+nv}{)}
\end{sphinxVerbatim}


\bigskip\hrule\bigskip


\sphinxAtStartPar
The barostat parameters for fix styles \sphinxstyleemphasis{npt} and \sphinxstyleemphasis{nph} is specified
using one or more of the \sphinxstyleemphasis{iso}, \sphinxstyleemphasis{aniso}, \sphinxstyleemphasis{tri}, \sphinxstyleemphasis{x}, \sphinxstyleemphasis{y}, \sphinxstyleemphasis{z}, \sphinxstyleemphasis{xy},
\sphinxstyleemphasis{xz}, \sphinxstyleemphasis{yz}, and \sphinxstyleemphasis{couple} keywords.  These keywords give you the
ability to specify all 6 components of an external stress tensor, and
to couple various of these components together so that the dimensions
they represent are varied together during a constant\sphinxhyphen{}pressure
simulation.

\sphinxAtStartPar
Other barostat\sphinxhyphen{}related keywords are \sphinxstyleemphasis{pchain}, \sphinxstyleemphasis{mtk}, \sphinxstyleemphasis{ploop},
\sphinxstyleemphasis{nreset}, \sphinxstyleemphasis{drag}, and \sphinxstyleemphasis{dilate}, which are discussed below.

\sphinxAtStartPar
Orthogonal simulation boxes have 3 adjustable dimensions (x,y,z).
Triclinic (non\sphinxhyphen{}orthogonal) simulation boxes have 6 adjustable
dimensions (x,y,z,xy,xz,yz).  The {\hyperref[\detokenize{create_box::doc}]{\sphinxcrossref{\DUrole{doc}{create\_box}}}}, {\hyperref[\detokenize{read_data::doc}]{\sphinxcrossref{\DUrole{doc}{read data}}}}, and {\hyperref[\detokenize{read_restart::doc}]{\sphinxcrossref{\DUrole{doc}{read\_restart}}}} commands
specify whether the simulation box is orthogonal or non\sphinxhyphen{}orthogonal
(triclinic) and explain the meaning of the xy,xz,yz tilt factors.

\sphinxAtStartPar
The target pressures for each of the 6 components of the stress tensor
can be specified independently via the \sphinxstyleemphasis{x}, \sphinxstyleemphasis{y}, \sphinxstyleemphasis{z}, \sphinxstyleemphasis{xy}, \sphinxstyleemphasis{xz}, \sphinxstyleemphasis{yz}
keywords, which correspond to the 6 simulation box dimensions.  For
each component, the external pressure or tensor component at each
timestep is a ramped value during the run from \sphinxstyleemphasis{Pstart} to \sphinxstyleemphasis{Pstop}.
If a target pressure is specified for a component, then the
corresponding box dimension will change during a simulation.  For
example, if the \sphinxstyleemphasis{y} keyword is used, the y\sphinxhyphen{}box length will change.  If
the \sphinxstyleemphasis{xy} keyword is used, the xy tilt factor will change.  A box
dimension will not change if that component is not specified, although
you have the option to change that dimension via the {\hyperref[\detokenize{fix_deform::doc}]{\sphinxcrossref{\DUrole{doc}{fix deform}}}} command.

\sphinxAtStartPar
Note that in order to use the \sphinxstyleemphasis{xy}, \sphinxstyleemphasis{xz}, or \sphinxstyleemphasis{yz} keywords, the
simulation box must be triclinic, even if its initial tilt factors are
0.0.

\sphinxAtStartPar
For all barostat keywords, the \sphinxstyleemphasis{Pdamp} parameter operates like the
\sphinxstyleemphasis{Tdamp} parameter, determining the time scale on which pressure is
relaxed.  For example, a value of 10.0 means to relax the pressure in
a timespan of (roughly) 10 time units (e.g. \(\tau\) or fs or ps
\sphinxhyphen{} see the {\hyperref[\detokenize{units::doc}]{\sphinxcrossref{\DUrole{doc}{units}}}} command).

\begin{sphinxadmonition}{note}{Note:}
\sphinxAtStartPar
A Nose\sphinxhyphen{}Hoover barostat will not work well for arbitrary values
of \sphinxstyleemphasis{Pdamp}.  If \sphinxstyleemphasis{Pdamp} is too small, the pressure and volume can
fluctuate wildly; if it is too large, the pressure will take a very
long time to equilibrate.  A good choice for many models is a \sphinxstyleemphasis{Pdamp}
of around 1000 timesteps.  However, note that \sphinxstyleemphasis{Pdamp} is specified in
time units, and that timesteps are NOT the same as time units for most
{\hyperref[\detokenize{units::doc}]{\sphinxcrossref{\DUrole{doc}{units}}}} settings.
\end{sphinxadmonition}

\sphinxAtStartPar
The relaxation rate of the barostat is set by its inertia \(W\):
\begin{equation*}
\begin{split}W = (N + 1) k_B T_{\rm target} P_{\rm damp}^2\end{split}
\end{equation*}
\sphinxAtStartPar
where \(N\) is the number of atoms, \(k_B\) is the Boltzmann constant,
and \(T_{\rm target}\) is the target temperature of the barostat {\hyperref[\detokenize{fix_nh:nh-martyna}]{\sphinxcrossref{\DUrole{std,std-ref}{(Martyna)}}}}.
If a thermostat is defined, \(T_{\rm target}\) is the target temperature
of the thermostat. If a thermostat is not defined, \(T_{\rm target}\)
is set to the current temperature of the system when the barostat is initialized.
If this temperature is too low the simulation will quit with an error.
Note: in previous versions of LAMMPS, \(T_{\rm target}\) would default to
a value of 1.0 for \sphinxstyleemphasis{lj} units and 300.0 otherwise if the system had a temperature
of exactly zero.

\sphinxAtStartPar
If a thermostat is not specified by this fix, \(T_{\rm target}\) can be
manually specified using the \sphinxstyleemphasis{Ptemp} parameter. This may be useful if the
barostat is initialized when the current temperature does not reflect the
steady state temperature of the system. This keyword may also be useful in
athermal simulations where the temperature is not well defined.

\sphinxAtStartPar
Regardless of what atoms are in the fix group (the only atoms which
are time integrated), a global pressure or stress tensor is computed
for all atoms.  Similarly, when the size of the simulation box is
changed, all atoms are re\sphinxhyphen{}scaled to new positions, unless the keyword
\sphinxstyleemphasis{dilate} is specified with a \sphinxstyleemphasis{dilate\sphinxhyphen{}group\sphinxhyphen{}ID} for a group that
represents a subset of the atoms.  This can be useful, for example, to
leave the coordinates of atoms in a solid substrate unchanged and
controlling the pressure of a surrounding fluid.  This option should
be used with care, since it can be unphysical to dilate some atoms and
not others, because it can introduce large, instantaneous
displacements between a pair of atoms (one dilated, one not) that are
far from the dilation origin.  Also note that for atoms not in the fix
group, a separate time integration fix like {\hyperref[\detokenize{fix_nve::doc}]{\sphinxcrossref{\DUrole{doc}{fix nve}}}} or
{\hyperref[\detokenize{fix_nh::doc}]{\sphinxcrossref{\DUrole{doc}{fix nvt}}}} can be used on them, independent of whether they
are dilated or not.


\bigskip\hrule\bigskip


\sphinxAtStartPar
The \sphinxstyleemphasis{couple} keyword allows two or three of the diagonal components of
the pressure tensor to be “coupled” together.  The value specified
with the keyword determines which are coupled.  For example, \sphinxstyleemphasis{xz}
means the \sphinxstyleemphasis{Pxx} and \sphinxstyleemphasis{Pzz} components of the stress tensor are coupled.
\sphinxstyleemphasis{Xyz} means all 3 diagonal components are coupled.  Coupling means two
things: the instantaneous stress will be computed as an average of the
corresponding diagonal components, and the coupled box dimensions will
be changed together in lockstep, meaning coupled dimensions will be
dilated or contracted by the same percentage every timestep.  The
\sphinxstyleemphasis{Pstart}, \sphinxstyleemphasis{Pstop}, \sphinxstyleemphasis{Pdamp} parameters for any coupled dimensions must
be identical.  \sphinxstyleemphasis{Couple xyz} can be used for a 2d simulation; the \sphinxstyleemphasis{z}
dimension is simply ignored.


\bigskip\hrule\bigskip


\sphinxAtStartPar
The \sphinxstyleemphasis{iso}, \sphinxstyleemphasis{aniso}, and \sphinxstyleemphasis{tri} keywords are simply shortcuts that are
equivalent to specifying several other keywords together.

\sphinxAtStartPar
The keyword \sphinxstyleemphasis{iso} means couple all 3 diagonal components together when
pressure is computed (hydrostatic pressure), and dilate/contract the
dimensions together.  Using “iso Pstart Pstop Pdamp” is the same as
specifying these 4 keywords:

\begin{sphinxVerbatim}[commandchars=\\\{\}]
x Pstart Pstop Pdamp
y Pstart Pstop Pdamp
z Pstart Pstop Pdamp
couple xyz
\end{sphinxVerbatim}

\sphinxAtStartPar
The keyword \sphinxstyleemphasis{aniso} means \sphinxstyleemphasis{x}, \sphinxstyleemphasis{y}, and \sphinxstyleemphasis{z} dimensions are controlled
independently using the \sphinxstyleemphasis{Pxx}, \sphinxstyleemphasis{Pyy}, and \sphinxstyleemphasis{Pzz} components of the
stress tensor as the driving forces, and the specified scalar external
pressure.  Using “aniso Pstart Pstop Pdamp” is the same as specifying
these 4 keywords:

\begin{sphinxVerbatim}[commandchars=\\\{\}]
x Pstart Pstop Pdamp
y Pstart Pstop Pdamp
z Pstart Pstop Pdamp
couple none
\end{sphinxVerbatim}

\sphinxAtStartPar
The keyword \sphinxstyleemphasis{tri} means \sphinxstyleemphasis{x}, \sphinxstyleemphasis{y}, \sphinxstyleemphasis{z}, \sphinxstyleemphasis{xy}, \sphinxstyleemphasis{xz}, and \sphinxstyleemphasis{yz} dimensions
are controlled independently using their individual stress components
as the driving forces, and the specified scalar pressure as the
external normal stress.  Using “tri Pstart Pstop Pdamp” is the same as
specifying these 7 keywords:

\begin{sphinxVerbatim}[commandchars=\\\{\}]
x Pstart Pstop Pdamp
y Pstart Pstop Pdamp
z Pstart Pstop Pdamp
xy 0.0 0.0 Pdamp
yz 0.0 0.0 Pdamp
xz 0.0 0.0 Pdamp
couple none
\end{sphinxVerbatim}


\bigskip\hrule\bigskip


\sphinxAtStartPar
In some cases (e.g. for solids) the pressure (volume) and/or
temperature of the system can oscillate undesirably when a Nose/Hoover
barostat and thermostat is applied.  The optional \sphinxstyleemphasis{drag} keyword will
damp these oscillations, although it alters the Nose/Hoover equations.
A value of 0.0 (no drag) leaves the Nose/Hoover formalism unchanged.
A non\sphinxhyphen{}zero value adds a drag term; the larger the value specified, the
greater the damping effect.  Performing a short run and monitoring the
pressure and temperature is the best way to determine if the drag term
is working.  Typically a value between 0.2 to 2.0 is sufficient to
damp oscillations after a few periods. Note that use of the drag
keyword will interfere with energy conservation and will also change
the distribution of positions and velocities so that they do not
correspond to the nominal NVT, NPT, or NPH ensembles.

\sphinxAtStartPar
An alternative way to control initial oscillations is to use chain
thermostats. The keyword \sphinxstyleemphasis{tchain} determines the number of thermostats
in the particle thermostat. A value of 1 corresponds to the original
Nose\sphinxhyphen{}Hoover thermostat. The keyword \sphinxstyleemphasis{pchain} specifies the number of
thermostats in the chain thermostatting the barostat degrees of
freedom. A value of 0 corresponds to no thermostatting of the
barostat variables.

\sphinxAtStartPar
The \sphinxstyleemphasis{mtk} keyword controls whether or not the correction terms due to
Martyna, Tuckerman, and Klein are included in the equations of motion
{\hyperref[\detokenize{fix_nh:nh-martyna}]{\sphinxcrossref{\DUrole{std,std-ref}{(Martyna)}}}}.  Specifying \sphinxstyleemphasis{no} reproduces the original
Hoover barostat, whose volume probability distribution function
differs from the true NPT and NPH ensembles by a factor of 1/V.  Hence
using \sphinxstyleemphasis{yes} is more correct, but in many cases the difference is
negligible.

\sphinxAtStartPar
The keyword \sphinxstyleemphasis{tloop} can be used to improve the accuracy of integration
scheme at little extra cost.  The initial and final updates of the
thermostat variables are broken up into \sphinxstyleemphasis{tloop} sub\sphinxhyphen{}steps, each of
length \sphinxstyleemphasis{dt}/\sphinxstyleemphasis{tloop}. This corresponds to using a first\sphinxhyphen{}order
Suzuki\sphinxhyphen{}Yoshida scheme {\hyperref[\detokenize{fix_nh:nh-tuckerman}]{\sphinxcrossref{\DUrole{std,std-ref}{(Tuckerman)}}}}.  The keyword \sphinxstyleemphasis{ploop}
does the same thing for the barostat thermostat.

\sphinxAtStartPar
The keyword \sphinxstyleemphasis{nreset} controls how often the reference dimensions used
to define the strain energy are reset.  If this keyword is not used,
or is given a value of zero, then the reference dimensions are set to
those of the initial simulation domain and are never changed. If the
simulation domain changes significantly during the simulation, then
the final average pressure tensor will differ significantly from the
specified values of the external stress tensor.  A value of \sphinxstyleemphasis{nstep}
means that every \sphinxstyleemphasis{nstep} timesteps, the reference dimensions are set
to those of the current simulation domain.

\sphinxAtStartPar
The \sphinxstyleemphasis{scaleyz}, \sphinxstyleemphasis{scalexz}, and \sphinxstyleemphasis{scalexy} keywords control whether or
not the corresponding tilt factors are scaled with the associated box
dimensions when barostatting triclinic periodic cells.  The default
values \sphinxstyleemphasis{yes} will turn on scaling, which corresponds to adjusting the
linear dimensions of the cell while preserving its shape.  Choosing
\sphinxstyleemphasis{no} ensures that the tilt factors are not scaled with the box
dimensions. See below for restrictions and default values in different
situations. In older versions of LAMMPS, scaling of tilt factors was
not performed. The old behavior can be recovered by setting all three
scale keywords to \sphinxstyleemphasis{no}.

\sphinxAtStartPar
The \sphinxstyleemphasis{flip} keyword allows the tilt factors for a triclinic box to
exceed half the distance of the parallel box length, as discussed
below.  If the \sphinxstyleemphasis{flip} value is set to \sphinxstyleemphasis{yes}, the bound is enforced by
flipping the box when it is exceeded.  If the \sphinxstyleemphasis{flip} value is set to
\sphinxstyleemphasis{no}, the tilt will continue to change without flipping.  Note that if
applied stress induces large deformations (e.g. in a liquid), this
means the box shape can tilt dramatically and LAMMPS will run less
efficiently, due to the large volume of communication needed to
acquire ghost atoms around a processor’s irregular\sphinxhyphen{}shaped subdomain.
For extreme values of tilt, LAMMPS may also lose atoms and generate an
error.

\sphinxAtStartPar
The \sphinxstyleemphasis{fixedpoint} keyword specifies the fixed point for barostat volume
changes. By default, it is the center of the box.  Whatever point is
chosen will not move during the simulation.  For example, if the lower
periodic boundaries pass through (0,0,0), and this point is provided
to \sphinxstyleemphasis{fixedpoint}, then the lower periodic boundaries will remain at
(0,0,0), while the upper periodic boundaries will move twice as
far. In all cases, the particle trajectories are unaffected by the
chosen value, except for a time\sphinxhyphen{}dependent constant translation of
positions.

\sphinxAtStartPar
If the \sphinxstyleemphasis{update} keyword is used with the \sphinxstyleemphasis{dipole} value, then the
orientation of the dipole moment of each particle is also updated
during the time integration.  This option should be used for models
where a dipole moment is assigned to finite\sphinxhyphen{}size particles,
e.g. spheroids via use of the {\hyperref[\detokenize{atom_style::doc}]{\sphinxcrossref{\DUrole{doc}{atom\_style hybrid sphere dipole}}}} command.

\sphinxAtStartPar
The default dipole orientation integrator can be changed to the
Dullweber\sphinxhyphen{}Leimkuhler\sphinxhyphen{}McLachlan integration scheme
{\hyperref[\detokenize{fix_nh:nh-dullweber}]{\sphinxcrossref{\DUrole{std,std-ref}{(Dullweber)}}}} when using \sphinxstyleemphasis{update} with the value
\sphinxstyleemphasis{dipole/dlm}. This integrator is symplectic and time\sphinxhyphen{}reversible,
giving better energy conservation and allows slightly longer timesteps
at only a small additional computational cost.


\bigskip\hrule\bigskip


\begin{sphinxadmonition}{note}{Note:}
\sphinxAtStartPar
Using a barostat coupled to tilt dimensions \sphinxstyleemphasis{xy}, \sphinxstyleemphasis{xz}, \sphinxstyleemphasis{yz} can
sometimes result in arbitrarily large values of the tilt dimensions,
i.e. a dramatically deformed simulation box.  LAMMPS allows the tilt
factors to grow a small amount beyond the normal limit of half the box
length (0.6 times the box length), and then performs a box “flip” to
an equivalent periodic cell.  See the discussion of the \sphinxstyleemphasis{flip} keyword
above, to allow this bound to be exceeded, if desired.
\end{sphinxadmonition}

\sphinxAtStartPar
The flip operation is described in more detail in the page for
{\hyperref[\detokenize{fix_deform::doc}]{\sphinxcrossref{\DUrole{doc}{fix deform}}}}.  Both the barostat dynamics and the atom
trajectories are unaffected by this operation.  However, if a tilt
factor is incremented by a large amount (1.5 times the box length) on
a single timestep, LAMMPS can not accommodate this event and will
terminate the simulation with an error. This error typically indicates
that there is something badly wrong with how the simulation was
constructed, such as specifying values of \sphinxstyleemphasis{Pstart} that are too far
from the current stress value, or specifying a timestep that is too
large. Triclinic barostatting should be used with care. This also is
true for other barostat styles, although they tend to be more
forgiving of insults. In particular, it is important to recognize that
equilibrium liquids can not support a shear stress and that
equilibrium solids can not support shear stresses that exceed the
yield stress.

\sphinxAtStartPar
One exception to this rule is if the first dimension in the tilt factor
(x for xy) is non\sphinxhyphen{}periodic.  In that case, the limits on the tilt
factor are not enforced, since flipping the box in that dimension does
not change the atom positions due to non\sphinxhyphen{}periodicity.  In this mode,
if you tilt the system to extreme angles, the simulation will simply
become inefficient due to the highly skewed simulation box.

\begin{sphinxadmonition}{note}{Note:}
\sphinxAtStartPar
Unlike the {\hyperref[\detokenize{fix_temp_berendsen::doc}]{\sphinxcrossref{\DUrole{doc}{fix temp/berendsen}}}} command
which performs thermostatting but NO time integration, these fixes
perform thermostatting/barostatting AND time integration.  Thus you
should not use any other time integration fix, such as {\hyperref[\detokenize{fix_nve::doc}]{\sphinxcrossref{\DUrole{doc}{fix nve}}}} on atoms to which this fix is applied.  Likewise,
fix nvt and fix npt should not normally be used on atoms that also
have their temperature controlled by another fix \sphinxhyphen{} e.g. by {\hyperref[\detokenize{fix_nh::doc}]{\sphinxcrossref{\DUrole{doc}{fix langevin}}}} or {\hyperref[\detokenize{fix_temp_rescale::doc}]{\sphinxcrossref{\DUrole{doc}{fix temp/rescale}}}}
commands.
\end{sphinxadmonition}

\sphinxAtStartPar
See the {\hyperref[\detokenize{Howto_thermostat::doc}]{\sphinxcrossref{\DUrole{doc}{Howto thermostat}}}} and {\hyperref[\detokenize{Howto_barostat::doc}]{\sphinxcrossref{\DUrole{doc}{Howto barostat}}}} doc pages for a discussion of different
ways to compute temperature and perform thermostatting and
barostatting.


\bigskip\hrule\bigskip


\sphinxAtStartPar
These fixes compute a temperature and pressure each timestep.  To do
this, the thermostat and barostat fixes create their own computes of
style “temp” and “pressure”, as if one of these sets of commands had
been issued:

\sphinxAtStartPar
For fix nvt:

\begin{sphinxVerbatim}[commandchars=\\\{\}]
\PYG{k}{compute }\PYG{n+nv+nvIdentifier}{fix\PYGZhy{}ID\PYGZus{}temp}\PYG{+w}{ }\PYG{n+nv+nvIdentifier}{group\PYGZhy{}ID}\PYG{+w}{ }\PYG{n}{temp}
\end{sphinxVerbatim}

\sphinxAtStartPar
For fix npt and fix nph:

\begin{sphinxVerbatim}[commandchars=\\\{\}]
\PYG{k}{compute }\PYG{n+nv+nvIdentifier}{fix\PYGZhy{}ID\PYGZus{}temp}\PYG{+w}{ }\PYG{n+nv+nvIdentifier}{all}\PYG{+w}{ }\PYG{n}{temp}
\PYG{k}{compute }\PYG{n+nv+nvIdentifier}{fix\PYGZhy{}ID\PYGZus{}press}\PYG{+w}{ }\PYG{n+nv+nvIdentifier}{all}\PYG{+w}{ }\PYG{n}{pressure}\PYG{+w}{ }\PYG{n}{fix}\PYG{o}{\PYGZhy{}}\PYG{n}{ID\PYGZus{}temp}
\end{sphinxVerbatim}

\sphinxAtStartPar
For fix nvt, the group for the new temperature compute is the same as
the fix group.  For fix npt and fix nph, the group for both the new
temperature and pressure compute is “all” since pressure is computed
for the entire system.  In the case of fix nph, the temperature
compute is not used for thermostatting, but just for a kinetic\sphinxhyphen{}energy
contribution to the pressure.  See the {\hyperref[\detokenize{compute_temp::doc}]{\sphinxcrossref{\DUrole{doc}{compute temp}}}} and {\hyperref[\detokenize{compute_pressure::doc}]{\sphinxcrossref{\DUrole{doc}{compute pressure}}}}
commands for details.  Note that the IDs of the new computes are the
fix\sphinxhyphen{}ID + underscore + “temp” or fix\_ID + underscore + “press”.

\sphinxAtStartPar
Note that these are NOT the computes used by thermodynamic output (see
the {\hyperref[\detokenize{thermo_style::doc}]{\sphinxcrossref{\DUrole{doc}{thermo\_style}}}} command) with ID = \sphinxstyleemphasis{thermo\_temp}
and \sphinxstyleemphasis{thermo\_press}.  This means you can change the attributes of these
fix’s temperature or pressure via the
{\hyperref[\detokenize{compute_modify::doc}]{\sphinxcrossref{\DUrole{doc}{compute\_modify}}}} command.  Or you can print this
temperature or pressure during thermodynamic output via the
{\hyperref[\detokenize{thermo_style::doc}]{\sphinxcrossref{\DUrole{doc}{thermo\_style custom}}}} command using the appropriate
compute\sphinxhyphen{}ID.  It also means that changing attributes of \sphinxstyleemphasis{thermo\_temp}
or \sphinxstyleemphasis{thermo\_press} will have no effect on this fix.

\sphinxAtStartPar
Like other fixes that perform thermostatting, this fix can be used
with {\hyperref[\detokenize{compute::doc}]{\sphinxcrossref{\DUrole{doc}{compute commands}}}} that remove a “bias” from the
atom velocities.  E.g. to apply the thermostat only to atoms within a
spatial {\hyperref[\detokenize{region::doc}]{\sphinxcrossref{\DUrole{doc}{region}}}}, or to remove the center\sphinxhyphen{}of\sphinxhyphen{}mass
velocity from a group of atoms, or to remove the x\sphinxhyphen{}component of
velocity from the calculation.

\sphinxAtStartPar
This is not done by default, but only if the {\hyperref[\detokenize{fix_modify::doc}]{\sphinxcrossref{\DUrole{doc}{fix\_modify}}}} command is used to assign a temperature compute to this
fix that includes such a bias term.  See the doc pages for individual
{\hyperref[\detokenize{compute::doc}]{\sphinxcrossref{\DUrole{doc}{compute temp commands}}}} to determine which ones include
a bias.  In this case, the thermostat works in the following manner:
bias is removed from each atom, thermostatting is performed on the
remaining thermal degrees of freedom, and the bias is added back in.


\bigskip\hrule\bigskip


\sphinxAtStartPar
These fixes can be used with either the \sphinxstyleemphasis{verlet} or \sphinxstyleemphasis{respa}
{\hyperref[\detokenize{run_style::doc}]{\sphinxcrossref{\DUrole{doc}{integrators}}}}. When using one of the barostat fixes
with \sphinxstyleemphasis{respa}, LAMMPS uses an integrator constructed
according to the following factorization of the Liouville propagator
(for two rRESPA levels):
\begin{equation*}
\begin{split}\exp \left(\mathrm{i} L \Delta t \right) = & \hat{E}
\exp \left(\mathrm{i} L_{\rm T\textrm{-}baro} \frac{\Delta t}{2} \right)
\exp \left(\mathrm{i} L_{\rm T\textrm{-}part} \frac{\Delta t}{2} \right)
\exp \left(\mathrm{i} L_{\epsilon , 2} \frac{\Delta t}{2} \right)
\exp \left(\mathrm{i} L_{2}^{(2)} \frac{\Delta t}{2} \right) \\
&\times \left[
\exp \left(\mathrm{i} L_{2}^{(1)} \frac{\Delta t}{2n} \right)
\exp \left(\mathrm{i} L_{\epsilon , 1} \frac{\Delta t}{2n} \right)
\exp \left(\mathrm{i} L_1 \frac{\Delta t}{n} \right)
\exp \left(\mathrm{i} L_{\epsilon , 1} \frac{\Delta t}{2n} \right)
\exp \left(\mathrm{i} L_{2}^{(1)} \frac{\Delta t}{2n} \right)
\right]^n \\
&\times
\exp \left(\mathrm{i} L_{2}^{(2)} \frac{\Delta t}{2} \right)
\exp \left(\mathrm{i} L_{\epsilon , 2} \frac{\Delta t}{2} \right)
\exp \left(\mathrm{i} L_{\rm T\textrm{-}part} \frac{\Delta t}{2} \right)
\exp \left(\mathrm{i} L_{\rm T\textrm{-}baro} \frac{\Delta t}{2} \right) \\
&+ \mathcal{O} \left(\Delta t^3 \right)\end{split}
\end{equation*}
\sphinxAtStartPar
This factorization differs somewhat from that of Tuckerman et al, in
that the barostat is only updated at the outermost rRESPA level,
whereas Tuckerman’s factorization requires splitting the pressure into
pieces corresponding to the forces computed at each rRESPA level. In
theory, the latter method will exhibit better numerical stability. In
practice, because Pdamp is normally chosen to be a large multiple of
the outermost rRESPA timestep, the barostat dynamics are not the
limiting factor for numerical stability. Both factorizations are
time\sphinxhyphen{}reversible and can be shown to preserve the phase space measure
of the underlying non\sphinxhyphen{}Hamiltonian equations of motion.

\begin{sphinxadmonition}{note}{Note:}
\sphinxAtStartPar
This implementation has been shown to conserve linear momentum
up to machine precision under NVT dynamics. Under NPT dynamics,
for a system with zero initial total linear momentum, the total
momentum fluctuates close to zero. It may occasionally undergo brief
excursions to non\sphinxhyphen{}negligible values, before returning close to zero.
Over long simulations, this has the effect of causing the center\sphinxhyphen{}of\sphinxhyphen{}mass
to undergo a slow random walk. This can be mitigated by resetting
the momentum at infrequent intervals using the
{\hyperref[\detokenize{fix_momentum::doc}]{\sphinxcrossref{\DUrole{doc}{fix momentum}}}} command.
\end{sphinxadmonition}


\bigskip\hrule\bigskip


\sphinxAtStartPar
The fix npt and fix nph commands can be used with rigid bodies or
mixtures of rigid bodies and non\sphinxhyphen{}rigid particles (e.g. solvent).  But
there are also {\hyperref[\detokenize{fix_rigid::doc}]{\sphinxcrossref{\DUrole{doc}{fix rigid/npt}}}} and {\hyperref[\detokenize{fix_rigid::doc}]{\sphinxcrossref{\DUrole{doc}{fix rigid/nph}}}} commands, which are typically a more natural
choice.  See the page for those commands for more discussion of
the various ways to do this.


\bigskip\hrule\bigskip


\sphinxAtStartPar
Styles with a \sphinxstyleemphasis{gpu}, \sphinxstyleemphasis{intel}, \sphinxstyleemphasis{kk}, \sphinxstyleemphasis{omp}, or \sphinxstyleemphasis{opt} suffix are
functionally the same as the corresponding style without the suffix.
They have been optimized to run faster, depending on your available
hardware, as discussed on the {\hyperref[\detokenize{Speed_packages::doc}]{\sphinxcrossref{\DUrole{doc}{Accelerator packages}}}}
page.  The accelerated styles take the same arguments and should
produce the same results, except for round\sphinxhyphen{}off and precision issues.

\sphinxAtStartPar
These accelerated styles are part of the GPU, INTEL, KOKKOS,
OPENMP, and OPT packages, respectively.  They are only enabled if
LAMMPS was built with those packages.  See the {\hyperref[\detokenize{Build_package::doc}]{\sphinxcrossref{\DUrole{doc}{Build package}}}} page for more info.

\sphinxAtStartPar
You can specify the accelerated styles explicitly in your input script
by including their suffix, or you can use the {\hyperref[\detokenize{Run_options::doc}]{\sphinxcrossref{\DUrole{doc}{\sphinxhyphen{}suffix command\sphinxhyphen{}line switch}}}} when you invoke LAMMPS, or you can use the
{\hyperref[\detokenize{suffix::doc}]{\sphinxcrossref{\DUrole{doc}{suffix}}}} command in your input script.

\sphinxAtStartPar
See the {\hyperref[\detokenize{Speed_packages::doc}]{\sphinxcrossref{\DUrole{doc}{Accelerator packages}}}} page for more
instructions on how to use the accelerated styles effectively.


\bigskip\hrule\bigskip



\subsection{Restart, fix\_modify, output, run start/stop, minimize info}
\label{\detokenize{fix_nh:restart-fix-modify-output-run-start-stop-minimize-info}}
\sphinxAtStartPar
These fixes writes the state of all the thermostat and barostat
variables to {\hyperref[\detokenize{restart::doc}]{\sphinxcrossref{\DUrole{doc}{binary restart files}}}}.  See the
{\hyperref[\detokenize{read_restart::doc}]{\sphinxcrossref{\DUrole{doc}{read\_restart}}}} command for info on how to re\sphinxhyphen{}specify
a fix in an input script that reads a restart file, so that the
operation of the fix continues in an uninterrupted fashion.

\sphinxAtStartPar
The {\hyperref[\detokenize{fix_modify::doc}]{\sphinxcrossref{\DUrole{doc}{fix\_modify}}}} \sphinxstyleemphasis{temp} and \sphinxstyleemphasis{press} options are
supported by these fixes.  You can use them to assign a
{\hyperref[\detokenize{compute::doc}]{\sphinxcrossref{\DUrole{doc}{compute}}}} you have defined to this fix which will be used
in its thermostatting or barostatting procedure, as described above.
If you do this, note that the kinetic energy derived from the compute
temperature should be consistent with the virial term computed using
all atoms for the pressure.  LAMMPS will warn you if you choose to
compute temperature on a subset of atoms.

\begin{sphinxadmonition}{note}{Note:}
\sphinxAtStartPar
If both the \sphinxstyleemphasis{temp} and \sphinxstyleemphasis{press} keywords are used in a single
thermo\_modify command (or in two separate commands), then the order in
which the keywords are specified is important.  Note that a {\hyperref[\detokenize{compute_pressure::doc}]{\sphinxcrossref{\DUrole{doc}{pressure compute}}}} defines its own temperature compute as
an argument when it is specified.  The \sphinxstyleemphasis{temp} keyword will override
this (for the pressure compute being used by fix npt), but only if the
\sphinxstyleemphasis{temp} keyword comes after the \sphinxstyleemphasis{press} keyword.  If the \sphinxstyleemphasis{temp} keyword
comes before the \sphinxstyleemphasis{press} keyword, then the new pressure compute
specified by the \sphinxstyleemphasis{press} keyword will be unaffected by the \sphinxstyleemphasis{temp}
setting.
\end{sphinxadmonition}

\sphinxAtStartPar
The cumulative energy change in the system imposed by these fixes, via
either thermostatting and/or barostatting, is included in the
{\hyperref[\detokenize{thermo_style::doc}]{\sphinxcrossref{\DUrole{doc}{thermodynamic output}}}} keywords \sphinxstyleemphasis{ecouple} and
\sphinxstyleemphasis{econserve}.  See the {\hyperref[\detokenize{thermo_style::doc}]{\sphinxcrossref{\DUrole{doc}{thermo\_style}}}} page for
details.

\sphinxAtStartPar
These fixes compute a global scalar which can be accessed by various
{\hyperref[\detokenize{Howto_output::doc}]{\sphinxcrossref{\DUrole{doc}{output commands}}}}.  The scalar is the same
cumulative energy change due to this fix described in the previous
paragraph.  The scalar value calculated by this fix is “extensive”.

\sphinxAtStartPar
These fixes compute also compute a global vector of quantities, which
can be accessed by various {\hyperref[\detokenize{Howto_output::doc}]{\sphinxcrossref{\DUrole{doc}{output commands}}}}.  The
vector values are “intensive”.

\sphinxAtStartPar
The vector stores internal Nose/Hoover thermostat and barostat
variables.  The number and meaning of the vector values depends on
which fix is used and the settings for keywords \sphinxstyleemphasis{tchain} and \sphinxstyleemphasis{pchain},
which specify the number of Nose/Hoover chains for the thermostat and
barostat.  If no thermostatting is done, then \sphinxstyleemphasis{tchain} is 0.  If no
barostatting is done, then \sphinxstyleemphasis{pchain} is 0.  In the following list,
“ndof” is 0, 1, 3, or 6, and is the number of degrees of freedom in
the barostat.  Its value is 0 if no barostat is used, else its value
is 6 if any off\sphinxhyphen{}diagonal stress tensor component is barostatted, else
its value is 1 if \sphinxstyleemphasis{couple xyz} is used or \sphinxstyleemphasis{couple xy} for a 2d
simulation, otherwise its value is 3.

\sphinxAtStartPar
The order of values in the global vector and their meaning is as
follows.  The notation means there are tchain values for eta, followed
by tchain for eta\_dot, followed by ndof for omega, etc:
\begin{itemize}
\item {} 
\sphinxAtStartPar
eta{[}tchain{]} = particle thermostat displacements (unitless)

\item {} 
\sphinxAtStartPar
eta\_dot{[}tchain{]} = particle thermostat velocities (1/time units)

\item {} 
\sphinxAtStartPar
omega{[}ndof{]} = barostat displacements (unitless)

\item {} 
\sphinxAtStartPar
omega\_dot{[}ndof{]} = barostat velocities (1/time units)

\item {} 
\sphinxAtStartPar
etap{[}pchain{]} = barostat thermostat displacements (unitless)

\item {} 
\sphinxAtStartPar
etap\_dot{[}pchain{]} = barostat thermostat velocities (1/time units)

\item {} 
\sphinxAtStartPar
PE\_eta{[}tchain{]} = potential energy of each particle thermostat displacement (energy units)

\item {} 
\sphinxAtStartPar
KE\_eta\_dot{[}tchain{]} = kinetic energy of each particle thermostat velocity (energy units)

\item {} 
\sphinxAtStartPar
PE\_omega{[}ndof{]} = potential energy of each barostat displacement (energy units)

\item {} 
\sphinxAtStartPar
KE\_omega\_dot{[}ndof{]} = kinetic energy of each barostat velocity (energy units)

\item {} 
\sphinxAtStartPar
PE\_etap{[}pchain{]} = potential energy of each barostat thermostat displacement (energy units)

\item {} 
\sphinxAtStartPar
KE\_etap\_dot{[}pchain{]} = kinetic energy of each barostat thermostat velocity (energy units)

\item {} 
\sphinxAtStartPar
PE\_strain{[}1{]} = scalar strain energy (energy units)

\end{itemize}

\sphinxAtStartPar
These fixes can ramp their external temperature and pressure over
multiple runs, using the \sphinxstyleemphasis{start} and \sphinxstyleemphasis{stop} keywords of the
{\hyperref[\detokenize{run::doc}]{\sphinxcrossref{\DUrole{doc}{run}}}} command.  See the {\hyperref[\detokenize{run::doc}]{\sphinxcrossref{\DUrole{doc}{run}}}} command for details of
how to do this.

\sphinxAtStartPar
These fixes are not invoked during {\hyperref[\detokenize{minimize::doc}]{\sphinxcrossref{\DUrole{doc}{energy minimization}}}}.


\bigskip\hrule\bigskip



\subsection{Restrictions}
\label{\detokenize{fix_nh:restrictions}}
\sphinxAtStartPar
\sphinxstyleemphasis{X}, \sphinxstyleemphasis{y}, \sphinxstyleemphasis{z} cannot be barostatted if the associated dimension is not
periodic.  \sphinxstyleemphasis{Xy}, \sphinxstyleemphasis{xz}, and \sphinxstyleemphasis{yz} can only be barostatted if the
simulation domain is triclinic and the second dimension in the keyword
(\sphinxstyleemphasis{y} dimension in \sphinxstyleemphasis{xy}) is periodic.  \sphinxstyleemphasis{Z}, \sphinxstyleemphasis{xz}, and \sphinxstyleemphasis{yz}, cannot be
barostatted for 2D simulations.  The {\hyperref[\detokenize{create_box::doc}]{\sphinxcrossref{\DUrole{doc}{create\_box}}}},
{\hyperref[\detokenize{read_data::doc}]{\sphinxcrossref{\DUrole{doc}{read data}}}}, and {\hyperref[\detokenize{read_restart::doc}]{\sphinxcrossref{\DUrole{doc}{read\_restart}}}}
commands specify whether the simulation box is orthogonal or
non\sphinxhyphen{}orthogonal (triclinic) and explain the meaning of the xy,xz,yz
tilt factors.

\sphinxAtStartPar
For the \sphinxstyleemphasis{temp} keyword, the final Tstop cannot be 0.0 since it would
make the external T = 0.0 at some timestep during the simulation which
is not allowed in the Nose/Hoover formulation.

\sphinxAtStartPar
The \sphinxstyleemphasis{scaleyz yes} and \sphinxstyleemphasis{scalexz yes} keyword/value pairs can not be used
for 2D simulations. \sphinxstyleemphasis{scaleyz yes}, \sphinxstyleemphasis{scalexz yes}, and \sphinxstyleemphasis{scalexy yes} options
can only be used if the second dimension in the keyword is periodic,
and if the tilt factor is not coupled to the barostat via keywords
\sphinxstyleemphasis{tri}, \sphinxstyleemphasis{yz}, \sphinxstyleemphasis{xz}, and \sphinxstyleemphasis{xy}.

\sphinxAtStartPar
These fixes can be used with dynamic groups as defined by the
{\hyperref[\detokenize{group::doc}]{\sphinxcrossref{\DUrole{doc}{group}}}} command.  Likewise they can be used with groups to
which atoms are added or deleted over time, e.g. a deposition
simulation.  However, the conservation properties of the thermostat
and barostat are defined for systems with a static set of atoms.  You
may observe odd behavior if the atoms in a group vary dramatically
over time or the atom count becomes very small.


\subsection{Related commands}
\label{\detokenize{fix_nh:related-commands}}
\sphinxAtStartPar
{\hyperref[\detokenize{fix_nve::doc}]{\sphinxcrossref{\DUrole{doc}{fix nve}}}}, {\hyperref[\detokenize{fix_modify::doc}]{\sphinxcrossref{\DUrole{doc}{fix\_modify}}}},
{\hyperref[\detokenize{run_style::doc}]{\sphinxcrossref{\DUrole{doc}{run\_style}}}}


\subsection{Default}
\label{\detokenize{fix_nh:default}}
\sphinxAtStartPar
The keyword defaults are tchain = 3, pchain = 3, mtk = yes, tloop = 1,
ploop = 1, nreset = 0, drag = 0.0, dilate = all, couple = none,
flip = yes, scaleyz = scalexz = scalexy = yes if periodic in second
dimension and not coupled to barostat, otherwise no.


\bigskip\hrule\bigskip

\phantomsection\label{\detokenize{fix_nh:nh-martyna}}
\sphinxAtStartPar
\sphinxstylestrong{(Martyna)} Martyna, Tobias and Klein, J Chem Phys, 101, 4177 (1994).

\phantomsection\label{\detokenize{fix_nh:nh-parrinello}}
\sphinxAtStartPar
\sphinxstylestrong{(Parrinello)} Parrinello and Rahman, J Appl Phys, 52, 7182 (1981).

\phantomsection\label{\detokenize{fix_nh:nh-tuckerman}}
\sphinxAtStartPar
\sphinxstylestrong{(Tuckerman)} Tuckerman, Alejandre, Lopez\sphinxhyphen{}Rendon, Jochim, and
Martyna, J Phys A: Math Gen, 39, 5629 (2006).

\phantomsection\label{\detokenize{fix_nh:nh-shinoda}}
\sphinxAtStartPar
\sphinxstylestrong{(Shinoda)} Shinoda, Shiga, and Mikami, Phys Rev B, 69, 134103 (2004).

\phantomsection\label{\detokenize{fix_nh:nh-dullweber}}
\sphinxAtStartPar
\sphinxstylestrong{(Dullweber)} Dullweber, Leimkuhler and McLachlan, J Chem Phys, 107,
5840 (1997).

\sphinxstepscope

\index{fix nvt/eff@\spxentry{fix nvt/eff}}\ignorespaces 
\index{fix npt/eff@\spxentry{fix npt/eff}}\ignorespaces 
\index{fix nph/eff@\spxentry{fix nph/eff}}\ignorespaces 

\section{fix nvt/eff command}
\label{\detokenize{fix_nh_eff:fix-nvt-eff-command}}\label{\detokenize{fix_nh_eff:index-2}}\label{\detokenize{fix_nh_eff:index-1}}\label{\detokenize{fix_nh_eff:index-0}}\label{\detokenize{fix_nh_eff::doc}}

\section{fix npt/eff command}
\label{\detokenize{fix_nh_eff:fix-npt-eff-command}}

\section{fix nph/eff command}
\label{\detokenize{fix_nh_eff:fix-nph-eff-command}}

\subsection{Syntax}
\label{\detokenize{fix_nh_eff:syntax}}
\begin{sphinxVerbatim}[commandchars=\\\{\}]
\PYG{k}{fix }\PYG{n+nv+nvIdentifier}{ID}\PYG{+w}{ }\PYG{n+nv+nvIdentifier}{group\PYGZhy{}ID}\PYG{+w}{ }\PYG{n}{style\PYGZus{}name}\PYG{+w}{ }\PYG{n}{keyword}\PYG{+w}{ }\PYG{n}{value}\PYG{+w}{ }\PYG{n}{...}
\end{sphinxVerbatim}
\begin{itemize}
\item {} 
\sphinxAtStartPar
ID, group\sphinxhyphen{}ID are documented in {\hyperref[\detokenize{fix::doc}]{\sphinxcrossref{\DUrole{doc}{fix}}}} command

\item {} 
\sphinxAtStartPar
style\_name = \sphinxstyleemphasis{nvt/eff} or \sphinxstyleemphasis{npt/eff} or \sphinxstyleemphasis{nph/eff}
\begin{sphinxalltt}
one or more keyword value pairs may be appended
keyword = \sphinxstyleemphasis{temp} or \sphinxstyleemphasis{iso} or \sphinxstyleemphasis{aniso} or \sphinxstyleemphasis{tri} or \sphinxstyleemphasis{x} or \sphinxstyleemphasis{y} or \sphinxstyleemphasis{z} or \sphinxstyleemphasis{xy} or \sphinxstyleemphasis{yz} or \sphinxstyleemphasis{xz} or \sphinxstyleemphasis{couple} or \sphinxstyleemphasis{tchain} or \sphinxstyleemphasis{pchain} or \sphinxstyleemphasis{mtk} or \sphinxstyleemphasis{tloop} or \sphinxstyleemphasis{ploop} or \sphinxstyleemphasis{nreset} or \sphinxstyleemphasis{drag} or \sphinxstyleemphasis{dilate}
  \sphinxstyleemphasis{temp} values = Tstart Tstop Tdamp
    Tstart,Tstop = external temperature at start/end of run
    Tdamp = temperature damping parameter (time units)
  \sphinxstyleemphasis{iso} or \sphinxstyleemphasis{aniso} or \sphinxstyleemphasis{tri} values = Pstart Pstop Pdamp
    Pstart,Pstop = scalar external pressure at start/end of run (pressure units)
    Pdamp = pressure damping parameter (time units)
  \sphinxstyleemphasis{x} or \sphinxstyleemphasis{y} or \sphinxstyleemphasis{z} or \sphinxstyleemphasis{xy} or \sphinxstyleemphasis{yz} or \sphinxstyleemphasis{xz} values = Pstart Pstop Pdamp
    Pstart,Pstop = external stress tensor component at start/end of run (pressure units)
    Pdamp = stress damping parameter (time units)
  \sphinxstyleemphasis{couple} = \sphinxstyleemphasis{none} or \sphinxstyleemphasis{xyz} or \sphinxstyleemphasis{xy} or \sphinxstyleemphasis{yz} or \sphinxstyleemphasis{xz}
  \sphinxstyleemphasis{tchain} value = length of thermostat chain (1 = single thermostat)
  \sphinxstyleemphasis{pchain} values = length of thermostat chain on barostat (0 = no thermostat)
  \sphinxstyleemphasis{mtk} value = \sphinxstyleemphasis{yes} or \sphinxstyleemphasis{no} = add in MTK adjustment term or not
  \sphinxstyleemphasis{tloop} value = number of sub\sphinxhyphen{}cycles to perform on thermostat
  \sphinxstyleemphasis{ploop} value = number of sub\sphinxhyphen{}cycles to perform on barostat thermostat
  \sphinxstyleemphasis{nreset} value = reset reference cell every this many timesteps
  \sphinxstyleemphasis{drag} value = drag factor added to barostat/thermostat (0.0 = no drag)
  \sphinxstyleemphasis{dilate} value = \sphinxstyleemphasis{all} or \sphinxstyleemphasis{partial}
\end{sphinxalltt}

\end{itemize}


\subsection{Examples}
\label{\detokenize{fix_nh_eff:examples}}
\begin{sphinxVerbatim}[commandchars=\\\{\}]
\PYG{k}{fix }\PYG{n+nv+nvIdentifier}{1}\PYG{+w}{ }\PYG{n+nv+nvIdentifier}{all}\PYG{+w}{ }\PYG{n}{nvt}\PYG{o}{/}\PYG{n}{eff}\PYG{+w}{ }\PYG{n}{temp}\PYG{+w}{ }\PYG{l+m}{300.0}\PYG{+w}{ }\PYG{l+m}{300.0}\PYG{+w}{ }\PYG{l+m}{0.1}
\PYG{k}{fix }\PYG{n+nv+nvIdentifier}{1}\PYG{+w}{ }\PYG{n+nv+nvIdentifier}{part}\PYG{+w}{ }\PYG{n}{npt}\PYG{o}{/}\PYG{n}{eff}\PYG{+w}{ }\PYG{n}{temp}\PYG{+w}{ }\PYG{l+m}{300.0}\PYG{+w}{ }\PYG{l+m}{300.0}\PYG{+w}{ }\PYG{l+m}{0.1}\PYG{+w}{ }\PYG{n}{iso}\PYG{+w}{ }\PYG{l+m}{0.0}\PYG{+w}{ }\PYG{l+m}{0.0}\PYG{+w}{ }\PYG{l+m}{1.0}
\PYG{k}{fix }\PYG{n+nv+nvIdentifier}{2}\PYG{+w}{ }\PYG{n+nv+nvIdentifier}{part}\PYG{+w}{ }\PYG{n}{npt}\PYG{o}{/}\PYG{n}{eff}\PYG{+w}{ }\PYG{n}{temp}\PYG{+w}{ }\PYG{l+m}{300.0}\PYG{+w}{ }\PYG{l+m}{300.0}\PYG{+w}{ }\PYG{l+m}{0.1}\PYG{+w}{ }\PYG{n}{tri}\PYG{+w}{ }\PYG{l+m}{5.0}\PYG{+w}{ }\PYG{l+m}{5.0}\PYG{+w}{ }\PYG{l+m}{1.0}
\PYG{k}{fix }\PYG{n+nv+nvIdentifier}{2}\PYG{+w}{ }\PYG{n+nv+nvIdentifier}{ice}\PYG{+w}{ }\PYG{n}{nph}\PYG{o}{/}\PYG{n}{eff}\PYG{+w}{ }\PYG{n}{x}\PYG{+w}{ }\PYG{l+m}{1.0}\PYG{+w}{ }\PYG{l+m}{1.0}\PYG{+w}{ }\PYG{l+m}{0.5}\PYG{+w}{ }\PYG{n}{y}\PYG{+w}{ }\PYG{l+m}{2.0}\PYG{+w}{ }\PYG{l+m}{2.0}\PYG{+w}{ }\PYG{l+m}{0.5}\PYG{+w}{ }\PYG{n}{z}\PYG{+w}{ }\PYG{l+m}{3.0}\PYG{+w}{ }\PYG{l+m}{3.0}\PYG{+w}{ }\PYG{l+m}{0.5}\PYG{+w}{ }\PYG{n}{yz}\PYG{+w}{ }\PYG{l+m}{0.1}\PYG{+w}{ }\PYG{l+m}{0.1}\PYG{+w}{ }\PYG{l+m}{0.5}\PYG{+w}{ }\PYG{n}{xz}\PYG{+w}{ }\PYG{l+m}{0.2}\PYG{+w}{ }\PYG{l+m}{0.2}\PYG{+w}{ }\PYG{l+m}{0.5}\PYG{+w}{ }\PYG{n}{xy}\PYG{+w}{ }\PYG{l+m}{0.3}\PYG{+w}{ }\PYG{l+m}{0.3}\PYG{+w}{ }\PYG{l+m}{0.5}\PYG{+w}{ }\PYG{n}{nreset}\PYG{+w}{ }\PYG{l+m}{1000}
\end{sphinxVerbatim}


\subsection{Description}
\label{\detokenize{fix_nh_eff:description}}
\sphinxAtStartPar
These commands perform time integration on Nose\sphinxhyphen{}Hoover style
non\sphinxhyphen{}Hamiltonian equations of motion for nuclei and electrons in the
group for the {\hyperref[\detokenize{pair_eff::doc}]{\sphinxcrossref{\DUrole{doc}{electron force field}}}} model.  The fixes
are designed to generate positions and velocities sampled from the
canonical (nvt), isothermal\sphinxhyphen{}isobaric (npt), and isenthalpic (nph)
ensembles.  This is achieved by adding some dynamic variables which
are coupled to the particle velocities (thermostatting) and simulation
domain dimensions (barostatting).  In addition to basic thermostatting
and barostatting, these fixes can also create a chain of thermostats
coupled to the particle thermostat, and another chain of thermostats
coupled to the barostat variables. The barostat can be coupled to the
overall box volume, or to individual dimensions, including the \sphinxstyleemphasis{xy},
\sphinxstyleemphasis{xz} and \sphinxstyleemphasis{yz} tilt dimensions. The external pressure of the barostat
can be specified as either a scalar pressure (isobaric ensemble) or as
components of a symmetric stress tensor (constant stress ensemble).
When used correctly, the time\sphinxhyphen{}averaged temperature and stress tensor
of the particles will match the target values specified by
Tstart/Tstop and Pstart/Pstop.

\sphinxAtStartPar
The operation of these fixes is exactly like that described by the
{\hyperref[\detokenize{fix_nh::doc}]{\sphinxcrossref{\DUrole{doc}{fix nvt, npt, and nph}}}} commands, except that the radius
and radial velocity of electrons are also updated.  Likewise the
temperature and pressure calculated by the fix, using the computes it
creates (as discussed in the {\hyperref[\detokenize{fix_nh::doc}]{\sphinxcrossref{\DUrole{doc}{fix nvt, npt, and nph}}}}
doc page), are performed with computes that include the eFF contribution
to the temperature or kinetic energy from the electron radial velocity.

\begin{sphinxadmonition}{note}{Note:}
\sphinxAtStartPar
there are two different pressures that can be reported for eFF
when defining the pair\_style (see {\hyperref[\detokenize{pair_eff::doc}]{\sphinxcrossref{\DUrole{doc}{pair eff/cut}}}} to
understand these settings), one (default) that considers electrons do
not contribute radial virial components (i.e. electrons treated as
incompressible ‘rigid’ spheres) and one that does.  The radial
electronic contributions to the virials are only tallied if the
flexible pressure option is set, and this will affect both global and
per\sphinxhyphen{}atom quantities.  In principle, the true pressure of a system is
somewhere in between the rigid and the flexible eFF pressures, but,
for most cases, the difference between these two pressures will not be
significant over long\sphinxhyphen{}term averaged runs (i.e. even though the energy
partitioning changes, the total energy remains similar).
\end{sphinxadmonition}

\begin{sphinxadmonition}{note}{Note:}
\sphinxAtStartPar
currently, there is no available option for the user to set or
create temperature distributions that include the radial electronic
degrees of freedom with the {\hyperref[\detokenize{velocity::doc}]{\sphinxcrossref{\DUrole{doc}{velocity}}}} command, so the
the user must allow for these degrees of freedom to equilibrate
(i.e. equi\sphinxhyphen{}partitioning of energy) through time integration.
\end{sphinxadmonition}


\subsection{Restart, fix\_modify, output, run start/stop, minimize info}
\label{\detokenize{fix_nh_eff:restart-fix-modify-output-run-start-stop-minimize-info}}
\sphinxAtStartPar
See the page for the {\hyperref[\detokenize{fix_nh::doc}]{\sphinxcrossref{\DUrole{doc}{fix nvt, npt, and nph}}}} commands
for details.


\subsection{Restrictions}
\label{\detokenize{fix_nh_eff:restrictions}}
\sphinxAtStartPar
This fix is part of the EFF package.  It is only enabled if
LAMMPS was built with that package.  See the {\hyperref[\detokenize{Build_package::doc}]{\sphinxcrossref{\DUrole{doc}{Build package}}}} page for more info.

\sphinxAtStartPar
Other restriction discussed on the page for the {\hyperref[\detokenize{fix_nh::doc}]{\sphinxcrossref{\DUrole{doc}{fix nvt, npt, and nph}}}} commands also apply.

\begin{sphinxadmonition}{note}{Note:}
\sphinxAtStartPar
The temperature for systems (regions or groups) with only
electrons and no nuclei is 0.0 (i.e. not defined) in the current
temperature calculations, a practical example would be a uniform
electron gas or a very hot plasma, where electrons remain delocalized
from the nuclei.  This is because, even though electron virials are
included in the temperature calculation, these are averaged over the
nuclear degrees of freedom only.  In such cases a corrective term must
be added to the pressure to get the correct kinetic contribution.
\end{sphinxadmonition}


\subsection{Related commands}
\label{\detokenize{fix_nh_eff:related-commands}}
\sphinxAtStartPar
{\hyperref[\detokenize{fix_nh::doc}]{\sphinxcrossref{\DUrole{doc}{fix nvt}}}}, {\hyperref[\detokenize{fix_nh::doc}]{\sphinxcrossref{\DUrole{doc}{fix nph}}}}, {\hyperref[\detokenize{fix_nh::doc}]{\sphinxcrossref{\DUrole{doc}{fix npt}}}},
{\hyperref[\detokenize{fix_modify::doc}]{\sphinxcrossref{\DUrole{doc}{fix\_modify}}}}, {\hyperref[\detokenize{run_style::doc}]{\sphinxcrossref{\DUrole{doc}{run\_style}}}}


\subsection{Default}
\label{\detokenize{fix_nh_eff:default}}
\sphinxAtStartPar
The keyword defaults are tchain = 3, pchain = 3, mtk = yes, tloop =
ploop = 1, nreset = 0, drag = 0.0, dilate = all, and couple = none.


\bigskip\hrule\bigskip

\phantomsection\label{\detokenize{fix_nh_eff:martyna1}}
\sphinxAtStartPar
\sphinxstylestrong{(Martyna)} Martyna, Tobias and Klein, J Chem Phys, 101, 4177 (1994).

\phantomsection\label{\detokenize{fix_nh_eff:parrinello}}
\sphinxAtStartPar
\sphinxstylestrong{(Parrinello)} Parrinello and Rahman, J Appl Phys, 52, 7182 (1981).

\phantomsection\label{\detokenize{fix_nh_eff:tuckerman1}}
\sphinxAtStartPar
\sphinxstylestrong{(Tuckerman)} Tuckerman, Alejandre, Lopez\sphinxhyphen{}Rendon, Jochim, and
Martyna, J Phys A: Math Gen, 39, 5629 (2006).

\phantomsection\label{\detokenize{fix_nh_eff:shinoda2}}
\sphinxAtStartPar
\sphinxstylestrong{(Shinoda)} Shinoda, Shiga, and Mikami, Phys Rev B, 69, 134103 (2004).

\sphinxstepscope

\index{fix nvt/uef@\spxentry{fix nvt/uef}}\ignorespaces 
\index{fix npt/uef@\spxentry{fix npt/uef}}\ignorespaces 

\section{fix nvt/uef command}
\label{\detokenize{fix_nh_uef:fix-nvt-uef-command}}\label{\detokenize{fix_nh_uef:index-1}}\label{\detokenize{fix_nh_uef:index-0}}\label{\detokenize{fix_nh_uef::doc}}

\section{fix npt/uef command}
\label{\detokenize{fix_nh_uef:fix-npt-uef-command}}

\subsection{Syntax}
\label{\detokenize{fix_nh_uef:syntax}}
\begin{sphinxVerbatim}[commandchars=\\\{\}]
\PYG{k}{fix }\PYG{n+nv+nvIdentifier}{ID}\PYG{+w}{ }\PYG{n+nv+nvIdentifier}{group\PYGZhy{}ID}\PYG{+w}{ }\PYG{n}{style\PYGZus{}name}\PYG{+w}{ }\PYG{n}{erate}\PYG{+w}{ }\PYG{n}{edot\PYGZus{}x}\PYG{+w}{ }\PYG{n}{edot\PYGZus{}y}\PYG{+w}{ }\PYG{n}{temp}\PYG{+w}{ }\PYG{n}{Tstart}\PYG{+w}{ }\PYG{n}{Tstop}\PYG{+w}{ }\PYG{n}{Tdamp}\PYG{+w}{ }\PYG{n}{keyword}\PYG{+w}{ }\PYG{n}{value}\PYG{+w}{ }\PYG{n}{...}
\end{sphinxVerbatim}
\begin{itemize}
\item {} 
\sphinxAtStartPar
ID, group\sphinxhyphen{}ID are documented in {\hyperref[\detokenize{fix::doc}]{\sphinxcrossref{\DUrole{doc}{fix}}}} command

\item {} 
\sphinxAtStartPar
style\_name = \sphinxstyleemphasis{nvt/uef} or \sphinxstyleemphasis{npt/uef}

\item {} 
\sphinxAtStartPar
\sphinxstyleemphasis{Tstart}, \sphinxstyleemphasis{Tstop}, and \sphinxstyleemphasis{Tdamp} are documented in the {\hyperref[\detokenize{fix_nh::doc}]{\sphinxcrossref{\DUrole{doc}{fix npt}}}} command

\item {} 
\sphinxAtStartPar
\sphinxstyleemphasis{edot\_x} and \sphinxstyleemphasis{edot\_y} are the strain rates in the x and y directions (1/(time units))

\item {} 
\sphinxAtStartPar
one or more keyword/value pairs may be appended
\begin{sphinxalltt}
keyword = \sphinxstyleemphasis{erate} or \sphinxstyleemphasis{ext} or \sphinxstyleemphasis{strain} or \sphinxstyleemphasis{temp} or \sphinxstyleemphasis{iso} or \sphinxstyleemphasis{x} or \sphinxstyleemphasis{y} or \sphinxstyleemphasis{z} or \sphinxstyleemphasis{tchain} or \sphinxstyleemphasis{pchain} or \sphinxstyleemphasis{tloop} or \sphinxstyleemphasis{ploop} or \sphinxstyleemphasis{mtk}
  \sphinxstyleemphasis{erate} values = e\_x e\_y = true strain rates (required)
  \sphinxstyleemphasis{ext} value = \sphinxstyleemphasis{x} or \sphinxstyleemphasis{y} or \sphinxstyleemphasis{z} or \sphinxstyleemphasis{xy} or \sphinxstyleemphasis{yz} or \sphinxstyleemphasis{xz} = external dimensions
    sets the external dimensions used to calculate the scalar pressure
  \sphinxstyleemphasis{strain} values = e\_x e\_y = initial strain
    usually not needed, but may be needed to resume a run with a data file.
  \sphinxstyleemphasis{temp}, \sphinxstyleemphasis{iso}, \sphinxstyleemphasis{x}, \sphinxstyleemphasis{y}, \sphinxstyleemphasis{z}, \sphinxstyleemphasis{tchain}, \sphinxstyleemphasis{pchain}, \sphinxstyleemphasis{tloop}, \sphinxstyleemphasis{ploop}, \sphinxstyleemphasis{mtk}
    keywords documented by the {\hyperref[\detokenize{fix_nh::doc}]{\sphinxcrossref{\DUrole{doc}{fix npt}}}} command
\end{sphinxalltt}

\end{itemize}


\subsection{Examples}
\label{\detokenize{fix_nh_uef:examples}}
\begin{sphinxVerbatim}[commandchars=\\\{\}]
\PYG{k}{fix }\PYG{n+nv+nvIdentifier}{uniax\PYGZus{}nvt}\PYG{+w}{ }\PYG{n+nv+nvIdentifier}{all}\PYG{+w}{ }\PYG{n}{nvt}\PYG{o}{/}\PYG{n}{uef}\PYG{+w}{ }\PYG{n}{temp}\PYG{+w}{ }\PYG{l+m}{400}\PYG{+w}{ }\PYG{l+m}{400}\PYG{+w}{ }\PYG{l+m}{100}\PYG{+w}{ }\PYG{n}{erate}\PYG{+w}{ }\PYG{l+m}{0.00001}\PYG{+w}{ }\PYG{o}{\PYGZhy{}}\PYG{l+m}{0.000005}
\PYG{k}{fix }\PYG{n+nv+nvIdentifier}{biax\PYGZus{}nvt}\PYG{+w}{ }\PYG{n+nv+nvIdentifier}{all}\PYG{+w}{ }\PYG{n}{nvt}\PYG{o}{/}\PYG{n}{uef}\PYG{+w}{ }\PYG{n}{temp}\PYG{+w}{ }\PYG{l+m}{400}\PYG{+w}{ }\PYG{l+m}{400}\PYG{+w}{ }\PYG{l+m}{100}\PYG{+w}{ }\PYG{n}{erate}\PYG{+w}{ }\PYG{l+m}{0.000005}\PYG{+w}{ }\PYG{l+m}{0.000005}
\PYG{k}{fix }\PYG{n+nv+nvIdentifier}{uniax\PYGZus{}npt}\PYG{+w}{ }\PYG{n+nv+nvIdentifier}{all}\PYG{+w}{ }\PYG{n}{npt}\PYG{o}{/}\PYG{n}{uef}\PYG{+w}{ }\PYG{n}{temp}\PYG{+w}{ }\PYG{l+m}{400}\PYG{+w}{ }\PYG{l+m}{400}\PYG{+w}{ }\PYG{l+m}{300}\PYG{+w}{ }\PYG{n}{iso}\PYG{+w}{ }\PYG{l+m}{1}\PYG{+w}{ }\PYG{l+m}{1}\PYG{+w}{ }\PYG{l+m}{3000}\PYG{+w}{ }\PYG{n}{erate}\PYG{+w}{ }\PYG{l+m}{0.00001}\PYG{+w}{ }\PYG{o}{\PYGZhy{}}\PYG{l+m}{0.000005}\PYG{+w}{ }\PYG{n}{ext}\PYG{+w}{ }\PYG{n}{yz}
\PYG{k}{fix }\PYG{n+nv+nvIdentifier}{biax\PYGZus{}npt}\PYG{+w}{ }\PYG{n+nv+nvIdentifier}{all}\PYG{+w}{ }\PYG{n}{npt}\PYG{o}{/}\PYG{n}{uef}\PYG{+w}{ }\PYG{n}{temp}\PYG{+w}{ }\PYG{l+m}{400}\PYG{+w}{ }\PYG{l+m}{400}\PYG{+w}{ }\PYG{l+m}{100}\PYG{+w}{ }\PYG{n}{erate}\PYG{+w}{ }\PYG{o}{\PYGZhy{}}\PYG{l+m}{0.00001}\PYG{+w}{ }\PYG{l+m}{0.000005}\PYG{+w}{ }\PYG{n}{x}\PYG{+w}{ }\PYG{l+m}{1}\PYG{+w}{ }\PYG{l+m}{1}\PYG{+w}{ }\PYG{l+m}{3000}
\end{sphinxVerbatim}


\subsection{Description}
\label{\detokenize{fix_nh_uef:description}}
\sphinxAtStartPar
These fixes can be used to simulate non\sphinxhyphen{}equilibrium molecular dynamics
(NEMD) under diagonal flow fields, including uniaxial and bi\sphinxhyphen{}axial flow.
Simulations under continuous extensional flow may be carried out for an
indefinite amount of time.  It is an implementation of the boundary
conditions from {\hyperref[\detokenize{fix_nh_uef:dobson}]{\sphinxcrossref{\DUrole{std,std-ref}{(Dobson)}}}}, and also uses numerical
lattice reduction as was proposed by {\hyperref[\detokenize{fix_nh_uef:hunt}]{\sphinxcrossref{\DUrole{std,std-ref}{(Hunt)}}}}. The lattice
reduction algorithm is from {\hyperref[\detokenize{fix_nh_uef:semaev}]{\sphinxcrossref{\DUrole{std,std-ref}{(Semaev)}}}}. The fix is
intended for simulations of homogeneous flows, and integrates the SLLOD
equations of motion, originally proposed by Hoover and Ladd (see
{\hyperref[\detokenize{fix_nh_uef:sllod}]{\sphinxcrossref{\DUrole{std,std-ref}{(Evans and Morriss)}}}}).  Additional detail about this
implementation can be found in {\hyperref[\detokenize{fix_nh_uef:nicholson}]{\sphinxcrossref{\DUrole{std,std-ref}{(Nicholson and Rutledge)}}}}.

\sphinxAtStartPar
Note that NEMD simulations of a continuously strained system can be
performed using the {\hyperref[\detokenize{fix_deform::doc}]{\sphinxcrossref{\DUrole{doc}{fix deform}}}}, {\hyperref[\detokenize{fix_nvt_sllod::doc}]{\sphinxcrossref{\DUrole{doc}{fix nvt/sllod}}}}, and {\hyperref[\detokenize{compute_temp_deform::doc}]{\sphinxcrossref{\DUrole{doc}{compute temp/deform}}}}
commands.

\sphinxAtStartPar
The applied flow field is set by the \sphinxstyleemphasis{erate} keyword. The values
\sphinxstyleemphasis{edot\_x} and \sphinxstyleemphasis{edot\_y} correspond to the strain rates in the xx and yy
directions.  It is implicitly assumed that the flow field is
traceless, and therefore the strain rate in the zz direction is eqal
to \sphinxhyphen{}(\sphinxstyleemphasis{edot\_x} + \sphinxstyleemphasis{edot\_y}).

\begin{sphinxadmonition}{note}{Note:}
\sphinxAtStartPar
Due to an instability in the SLLOD equations under extension,
{\hyperref[\detokenize{fix_momentum::doc}]{\sphinxcrossref{\DUrole{doc}{fix momentum}}}} should be used to regularly reset the
linear momentum.
\end{sphinxadmonition}

\sphinxAtStartPar
The boundary conditions require a simulation box that does not have a
consistent alignment relative to the applied flow field. Since LAMMPS
utilizes an upper\sphinxhyphen{}triangular simulation box, it is not possible to
express the evolving simulation box in the same coordinate system as the
flow field.  These fixes keep track of two coordinate systems: the flow
frame, and the upper triangular LAMMPS frame. The coordinate systems are
related to each other through the QR decomposition, as is illustrated in
the image below.

\noindent{\hspace*{\fill}\sphinxincludegraphics{{uef_frames}.jpg}\hspace*{\fill}}

\sphinxAtStartPar
During most molecular dynamics operations, the system is represented
in the LAMMPS frame. Only when the positions and velocities are
updated is the system rotated to the flow frame, and it is rotated
back to the LAMMPS frame immediately afterwards. For this reason, all
vector\sphinxhyphen{}valued quantities (except for the tensors from
{\hyperref[\detokenize{compute_pressure_uef::doc}]{\sphinxcrossref{\DUrole{doc}{compute pressure/uef}}}} and
{\hyperref[\detokenize{compute_temp_uef::doc}]{\sphinxcrossref{\DUrole{doc}{compute temp/uef}}}}) will be computed in the
LAMMPS frame. Rotationally invariant scalar quantities like the
temperature and hydrostatic pressure are frame\sphinxhyphen{}invariant and will be
computed correctly. Additionally, the system is in the LAMMPS frame
during all of the output steps, and therefore trajectory files made
using the dump command will be in the LAMMPS frame unless the
{\hyperref[\detokenize{dump_cfg_uef::doc}]{\sphinxcrossref{\DUrole{doc}{dump cfg/uef}}}} command is used.


\bigskip\hrule\bigskip


\sphinxAtStartPar
Temperature control is achieved with the default Nose\sphinxhyphen{}Hoover style
thermostat documented in {\hyperref[\detokenize{fix_nh::doc}]{\sphinxcrossref{\DUrole{doc}{fix nvt}}}}.  When this fix is
active, only the peculiar velocity of each atom is stored, defined as
the velocity relative to the streaming velocity. This is in contrast to
{\hyperref[\detokenize{fix_nvt_sllod::doc}]{\sphinxcrossref{\DUrole{doc}{fix nvt/sllod}}}}, which uses a lab\sphinxhyphen{}frame velocity,
and removes the contribution from the streaming velocity in order to
compute the temperature.

\sphinxAtStartPar
Pressure control is achieved using the default Nose\sphinxhyphen{}Hoover barostat
documented in {\hyperref[\detokenize{fix_nh::doc}]{\sphinxcrossref{\DUrole{doc}{fix npt}}}}. There are two ways to control the
pressure using this fix. The first method involves using the \sphinxstyleemphasis{ext}
keyword along with the \sphinxstyleemphasis{iso} pressure style. With this method, the
pressure is controlled by scaling the simulation box isotropically to
achieve the average pressure only in the directions specified by
\sphinxstyleemphasis{ext}.  For example, if the \sphinxstyleemphasis{ext} value is set to \sphinxstyleemphasis{xy}, the average
pressure (Pxx+Pyy)/2 will be controlled.

\sphinxAtStartPar
This example command will control the total hydrostatic pressure under
uniaxial tension:

\begin{sphinxVerbatim}[commandchars=\\\{\}]
\PYG{k}{fix }\PYG{n+nv+nvIdentifier}{f1}\PYG{+w}{ }\PYG{n+nv+nvIdentifier}{all}\PYG{+w}{ }\PYG{n}{npt}\PYG{o}{/}\PYG{n}{uef}\PYG{+w}{ }\PYG{n}{temp}\PYG{+w}{ }\PYG{l+m}{0.7}\PYG{+w}{ }\PYG{l+m}{0.7}\PYG{+w}{ }\PYG{l+m}{0.5}\PYG{+w}{ }\PYG{n}{iso}\PYG{+w}{ }\PYG{l+m}{1}\PYG{+w}{ }\PYG{l+m}{1}\PYG{+w}{ }\PYG{l+m}{5}\PYG{+w}{ }\PYG{n}{erate}\PYG{+w}{ }\PYG{o}{\PYGZhy{}}\PYG{l+m}{0.5}\PYG{+w}{ }\PYG{o}{\PYGZhy{}}\PYG{l+m}{0.5}\PYG{+w}{ }\PYG{n}{ext}\PYG{+w}{ }\PYG{n}{xyz}
\end{sphinxVerbatim}

\sphinxAtStartPar
This example command will control the average stress in compression
directions, which would typically correspond to free surfaces under
drawing with uniaxial tension:

\begin{sphinxVerbatim}[commandchars=\\\{\}]
\PYG{k}{fix }\PYG{n+nv+nvIdentifier}{f2}\PYG{+w}{ }\PYG{n+nv+nvIdentifier}{all}\PYG{+w}{ }\PYG{n}{npt}\PYG{o}{/}\PYG{n}{uef}\PYG{+w}{ }\PYG{n}{temp}\PYG{+w}{ }\PYG{l+m}{0.7}\PYG{+w}{ }\PYG{l+m}{0.7}\PYG{+w}{ }\PYG{l+m}{0.5}\PYG{+w}{ }\PYG{n}{iso}\PYG{+w}{ }\PYG{l+m}{1}\PYG{+w}{ }\PYG{l+m}{1}\PYG{+w}{ }\PYG{l+m}{5}\PYG{+w}{ }\PYG{n}{erate}\PYG{+w}{ }\PYG{o}{\PYGZhy{}}\PYG{l+m}{0.5}\PYG{+w}{ }\PYG{o}{\PYGZhy{}}\PYG{l+m}{0.5}\PYG{+w}{ }\PYG{n}{ext}\PYG{+w}{ }\PYG{n}{xy}
\end{sphinxVerbatim}

\sphinxAtStartPar
The second method for pressure control involves setting the normal
stresses using the \sphinxstyleemphasis{x}, \sphinxstyleemphasis{y}, and/or \sphinxstyleemphasis{z} keywords. When using this
method, the same pressure must be specified via \sphinxstyleemphasis{Pstart} and \sphinxstyleemphasis{Pstop}
for all dimensions controlled. Any choice of pressure conditions that
would cause LAMMPS to compute a deviatoric stress are not permissible
and will result in an error. Additionally, all dimensions with
controlled stress must have the same applied strain rate. The \sphinxstyleemphasis{ext}
keyword must be set to the default value (\sphinxstyleemphasis{xyz}) when using this
method.

\sphinxAtStartPar
For example, the following commands will work:

\begin{sphinxVerbatim}[commandchars=\\\{\}]
\PYG{k}{fix }\PYG{n+nv+nvIdentifier}{f3}\PYG{+w}{ }\PYG{n+nv+nvIdentifier}{all}\PYG{+w}{ }\PYG{n}{npt}\PYG{o}{/}\PYG{n}{uef}\PYG{+w}{ }\PYG{n}{temp}\PYG{+w}{ }\PYG{l+m}{0.7}\PYG{+w}{ }\PYG{l+m}{0.7}\PYG{+w}{ }\PYG{l+m}{0.5}\PYG{+w}{ }\PYG{n}{x}\PYG{+w}{ }\PYG{l+m}{1}\PYG{+w}{ }\PYG{l+m}{1}\PYG{+w}{ }\PYG{l+m}{5}\PYG{+w}{ }\PYG{n}{y}\PYG{+w}{ }\PYG{l+m}{1}\PYG{+w}{ }\PYG{l+m}{1}\PYG{+w}{ }\PYG{l+m}{5}\PYG{+w}{ }\PYG{n}{erate}\PYG{+w}{ }\PYG{o}{\PYGZhy{}}\PYG{l+m}{0.5}\PYG{+w}{ }\PYG{o}{\PYGZhy{}}\PYG{l+m}{0.5}
\PYG{k}{fix }\PYG{n+nv+nvIdentifier}{f4}\PYG{+w}{ }\PYG{n+nv+nvIdentifier}{all}\PYG{+w}{ }\PYG{n}{npt}\PYG{o}{/}\PYG{n}{uef}\PYG{+w}{ }\PYG{n}{temp}\PYG{+w}{ }\PYG{l+m}{0.7}\PYG{+w}{ }\PYG{l+m}{0.7}\PYG{+w}{ }\PYG{l+m}{0.5}\PYG{+w}{ }\PYG{n}{z}\PYG{+w}{ }\PYG{l+m}{1}\PYG{+w}{ }\PYG{l+m}{1}\PYG{+w}{ }\PYG{l+m}{5}\PYG{+w}{ }\PYG{n}{erate}\PYG{+w}{ }\PYG{l+m}{0.5}\PYG{+w}{ }\PYG{l+m}{0.5}
\end{sphinxVerbatim}

\sphinxAtStartPar
The following commands will not work:

\begin{sphinxVerbatim}[commandchars=\\\{\}]
\PYG{k}{fix }\PYG{n+nv+nvIdentifier}{f5}\PYG{+w}{ }\PYG{n+nv+nvIdentifier}{all}\PYG{+w}{ }\PYG{n}{npt}\PYG{o}{/}\PYG{n}{uef}\PYG{+w}{ }\PYG{n}{temp}\PYG{+w}{ }\PYG{l+m}{0.7}\PYG{+w}{ }\PYG{l+m}{0.7}\PYG{+w}{ }\PYG{l+m}{0.5}\PYG{+w}{ }\PYG{n}{x}\PYG{+w}{ }\PYG{l+m}{1}\PYG{+w}{ }\PYG{l+m}{1}\PYG{+w}{ }\PYG{l+m}{5}\PYG{+w}{ }\PYG{n}{z}\PYG{+w}{ }\PYG{l+m}{1}\PYG{+w}{ }\PYG{l+m}{1}\PYG{+w}{ }\PYG{l+m}{5}\PYG{+w}{ }\PYG{n}{erate}\PYG{+w}{ }\PYG{o}{\PYGZhy{}}\PYG{l+m}{0.5}\PYG{+w}{ }\PYG{o}{\PYGZhy{}}\PYG{l+m}{0.5}
\PYG{k}{fix }\PYG{n+nv+nvIdentifier}{f6}\PYG{+w}{ }\PYG{n+nv+nvIdentifier}{all}\PYG{+w}{ }\PYG{n}{npt}\PYG{o}{/}\PYG{n}{uef}\PYG{+w}{ }\PYG{n}{temp}\PYG{+w}{ }\PYG{l+m}{0.7}\PYG{+w}{ }\PYG{l+m}{0.7}\PYG{+w}{ }\PYG{l+m}{0.5}\PYG{+w}{ }\PYG{n}{x}\PYG{+w}{ }\PYG{l+m}{1}\PYG{+w}{ }\PYG{l+m}{1}\PYG{+w}{ }\PYG{l+m}{5}\PYG{+w}{ }\PYG{n}{z}\PYG{+w}{ }\PYG{l+m}{2}\PYG{+w}{ }\PYG{l+m}{2}\PYG{+w}{ }\PYG{l+m}{5}\PYG{+w}{ }\PYG{n}{erate}\PYG{+w}{ }\PYG{l+m}{0.5}\PYG{+w}{ }\PYG{l+m}{0.5}
\end{sphinxVerbatim}


\bigskip\hrule\bigskip


\sphinxAtStartPar
These fixes compute a temperature and pressure each timestep.  To do
this, they create their own computes of style “temp/uef” and
“pressure/uef”, as if one of these two sets of commands had been
issued:

\begin{sphinxVerbatim}[commandchars=\\\{\}]
\PYG{k}{compute }\PYG{n+nv+nvIdentifier}{fix\PYGZhy{}ID\PYGZus{}temp}\PYG{+w}{ }\PYG{n+nv+nvIdentifier}{group\PYGZhy{}ID}\PYG{+w}{ }\PYG{n}{temp}\PYG{o}{/}\PYG{n}{uef}
\PYG{k}{compute }\PYG{n+nv+nvIdentifier}{fix\PYGZhy{}ID\PYGZus{}press}\PYG{+w}{ }\PYG{n+nv+nvIdentifier}{group\PYGZhy{}ID}\PYG{+w}{ }\PYG{n}{pressure}\PYG{o}{/}\PYG{n}{uef}\PYG{+w}{ }\PYG{n}{fix}\PYG{o}{\PYGZhy{}}\PYG{n}{ID\PYGZus{}temp}

\PYG{k}{compute }\PYG{n+nv+nvIdentifier}{fix\PYGZhy{}ID\PYGZus{}temp}\PYG{+w}{ }\PYG{n+nv+nvIdentifier}{all}\PYG{+w}{ }\PYG{n}{temp}\PYG{o}{/}\PYG{n}{uef}
\PYG{k}{compute }\PYG{n+nv+nvIdentifier}{fix\PYGZhy{}ID\PYGZus{}press}\PYG{+w}{ }\PYG{n+nv+nvIdentifier}{all}\PYG{+w}{ }\PYG{n}{pressure}\PYG{o}{/}\PYG{n}{uef}\PYG{+w}{ }\PYG{n}{fix}\PYG{o}{\PYGZhy{}}\PYG{n}{ID\PYGZus{}temp}
\end{sphinxVerbatim}

\sphinxAtStartPar
See the {\hyperref[\detokenize{compute_temp_uef::doc}]{\sphinxcrossref{\DUrole{doc}{compute temp/uef}}}} and {\hyperref[\detokenize{compute_pressure_uef::doc}]{\sphinxcrossref{\DUrole{doc}{compute
pressure/uef}}}} commands for details.  Note that
the IDs of the new computes are the fix\sphinxhyphen{}ID + underscore + “temp” or
fix\_ID + underscore + “press”.


\subsection{Restart, fix\_modify, output, run start/stop, minimize info}
\label{\detokenize{fix_nh_uef:restart-fix-modify-output-run-start-stop-minimize-info}}
\sphinxAtStartPar
The fix writes the state of all the thermostat and barostat variables,
as well as the cumulative strain applied, to {\hyperref[\detokenize{restart::doc}]{\sphinxcrossref{\DUrole{doc}{binary restart files}}}}.  See the {\hyperref[\detokenize{read_restart::doc}]{\sphinxcrossref{\DUrole{doc}{read\_restart}}}} command for info
on how to re\sphinxhyphen{}specify a fix in an input script that reads a restart file,
so that the operation of the fix continues in an uninterrupted fashion.

\begin{sphinxadmonition}{note}{Note:}
\sphinxAtStartPar
It is not necessary to set the \sphinxstyleemphasis{strain} keyword when resuming a
run from a restart file. Only for resuming from data files, which do
not contain the cumulative applied strain, will this keyword be
necessary.
\end{sphinxadmonition}

\sphinxAtStartPar
These fixes can be used with the {\hyperref[\detokenize{fix_modify::doc}]{\sphinxcrossref{\DUrole{doc}{fix\_modify}}}} \sphinxstyleemphasis{temp}
and \sphinxstyleemphasis{press} options. The temperature and pressure computes used must be
of type \sphinxstyleemphasis{temp/uef} and \sphinxstyleemphasis{pressure/uef}.

\sphinxAtStartPar
These fixes compute the same global scalar and vector quantities as
{\hyperref[\detokenize{fix_nh::doc}]{\sphinxcrossref{\DUrole{doc}{fix nvt andnpt}}}}.

\sphinxAtStartPar
These fixes are not invoked during {\hyperref[\detokenize{minimize::doc}]{\sphinxcrossref{\DUrole{doc}{energy minimization}}}}.


\subsection{Restrictions}
\label{\detokenize{fix_nh_uef:restrictions}}
\sphinxAtStartPar
These fixes are part of the UEF package. They are only enabled if LAMMPS
was built with that package. See the {\hyperref[\detokenize{Build_package::doc}]{\sphinxcrossref{\DUrole{doc}{Build package}}}} page for more info.

\sphinxAtStartPar
Due to requirements of the boundary conditions, when the \sphinxstyleemphasis{strain}
keyword is set to zero (or unset), the initial simulation box must be
cubic and have style triclinic. If the box is initially of type ortho,
use {\hyperref[\detokenize{change_box::doc}]{\sphinxcrossref{\DUrole{doc}{change\_box}}}} before invoking the fix.


\subsection{Related commands}
\label{\detokenize{fix_nh_uef:related-commands}}
\sphinxAtStartPar
{\hyperref[\detokenize{fix_nh::doc}]{\sphinxcrossref{\DUrole{doc}{fix nvt}}}}, {\hyperref[\detokenize{fix_nh::doc}]{\sphinxcrossref{\DUrole{doc}{fix npt}}}},
{\hyperref[\detokenize{fix_nvt_sllod::doc}]{\sphinxcrossref{\DUrole{doc}{fix nvt/sllod}}}},
{\hyperref[\detokenize{compute_temp_uef::doc}]{\sphinxcrossref{\DUrole{doc}{compute temp/uef}}}},
{\hyperref[\detokenize{compute_pressure_uef::doc}]{\sphinxcrossref{\DUrole{doc}{compute pressure/uef}}}},
{\hyperref[\detokenize{dump_cfg_uef::doc}]{\sphinxcrossref{\DUrole{doc}{dump cfg/uef}}}}


\subsection{Default}
\label{\detokenize{fix_nh_uef:default}}
\sphinxAtStartPar
The default keyword values specific to these fixes are exy = xyz, strain
= 0 0.  The remaining defaults are the same as for {\hyperref[\detokenize{fix_nh::doc}]{\sphinxcrossref{\DUrole{doc}{fix nvt or npt}}}} except tchain = 1.  The reason for this change is given in
{\hyperref[\detokenize{fix_nvt_sllod::doc}]{\sphinxcrossref{\DUrole{doc}{fix nvt/sllod}}}}.


\bigskip\hrule\bigskip

\phantomsection\label{\detokenize{fix_nh_uef:dobson}}
\sphinxAtStartPar
\sphinxstylestrong{(Dobson)} Dobson, J Chem Phys, 141, 184103 (2014).

\phantomsection\label{\detokenize{fix_nh_uef:hunt}}
\sphinxAtStartPar
\sphinxstylestrong{(Hunt)} Hunt, Mol Simul, 42, 347 (2016).

\phantomsection\label{\detokenize{fix_nh_uef:semaev}}
\sphinxAtStartPar
\sphinxstylestrong{(Semaev)} Semaev, Cryptography and Lattices, 181 (2001).

\phantomsection\label{\detokenize{fix_nh_uef:sllod}}
\sphinxAtStartPar
\sphinxstylestrong{(Evans and Morriss)} Evans and Morriss, Phys Rev A, 30, 1528 (1984).

\phantomsection\label{\detokenize{fix_nh_uef:nicholson}}
\sphinxAtStartPar
\sphinxstylestrong{(Nicholson and Rutledge)} Nicholson and Rutledge, J Chem Phys, 145,
244903 (2016).

\sphinxstepscope

\index{fix nonaffine/displacement@\spxentry{fix nonaffine/displacement}}\ignorespaces 

\section{fix nonaffine/displacement command}
\label{\detokenize{fix_nonaffine_displacement:fix-nonaffine-displacement-command}}\label{\detokenize{fix_nonaffine_displacement:index-0}}\label{\detokenize{fix_nonaffine_displacement::doc}}

\subsection{Syntax}
\label{\detokenize{fix_nonaffine_displacement:syntax}}
\begin{sphinxVerbatim}[commandchars=\\\{\}]
fix ID group nonaffine/displacement style args reference/style nstep keyword values
\end{sphinxVerbatim}
\begin{itemize}
\item {} 
\sphinxAtStartPar
ID, group are documented in {\hyperref[\detokenize{fix::doc}]{\sphinxcrossref{\DUrole{doc}{fix}}}} command

\item {} 
\sphinxAtStartPar
nonaffine/displacement = style name of this fix command

\item {} 
\sphinxAtStartPar
nevery = calculate nonaffine displacement every this many timesteps

\item {} 
\sphinxAtStartPar
style = \sphinxstyleemphasis{d2min} or \sphinxstyleemphasis{integrated}
\begin{sphinxalltt}
\sphinxstyleemphasis{d2min} args = cutoff args
  cutoff = \sphinxstyleemphasis{type} or \sphinxstyleemphasis{radius} or \sphinxstyleemphasis{custom}
    \sphinxstyleemphasis{type} args = none, cutoffs determined by atom types
    \sphinxstyleemphasis{radius} args = none, cutoffs determined based on atom diameters (atom style sphere)
    \sphinxstyleemphasis{custom} args = \sphinxstyleemphasis{rmax}, cutoff set by a constant numeric value \sphinxstyleemphasis{rmax} (distance units)
\sphinxstyleemphasis{integrated} args = none
\end{sphinxalltt}

\item {} 
\sphinxAtStartPar
reference/style = \sphinxstyleemphasis{fixed} or \sphinxstyleemphasis{update} or \sphinxstyleemphasis{offset}
\begin{sphinxalltt}
\sphinxstyleemphasis{fixed} = use a fixed reference frame at \sphinxstyleemphasis{nstep}
\sphinxstyleemphasis{update} = update the reference frame every \sphinxstyleemphasis{nstep} timesteps
\sphinxstyleemphasis{offset} = update the reference frame \sphinxstyleemphasis{nstep} timesteps before calculating the nonaffine displacement
\end{sphinxalltt}

\item {} 
\sphinxAtStartPar
zero or more keyword/value pairs may be appended
\begin{sphinxalltt}
\sphinxstyleemphasis{z/min} values = zmin
  zmin = minimum coordination number to calculate D2min
\end{sphinxalltt}

\end{itemize}


\subsection{Examples}
\label{\detokenize{fix_nonaffine_displacement:examples}}
\begin{sphinxVerbatim}[commandchars=\\\{\}]
\PYG{k}{fix }\PYG{n+nv+nvIdentifier}{1}\PYG{+w}{ }\PYG{n+nv+nvIdentifier}{all}\PYG{+w}{ }\PYG{n}{nonaffine}\PYG{o}{/}\PYG{n}{displacement}\PYG{+w}{ }\PYG{l+m}{100}\PYG{+w}{ }\PYG{n}{integrated}\PYG{+w}{ }\PYG{n}{update}\PYG{+w}{ }\PYG{l+m}{100}
\PYG{k}{fix }\PYG{n+nv+nvIdentifier}{1}\PYG{+w}{ }\PYG{n+nv+nvIdentifier}{all}\PYG{+w}{ }\PYG{n}{nonaffine}\PYG{o}{/}\PYG{n}{displacement}\PYG{+w}{ }\PYG{l+m}{1000}\PYG{+w}{ }\PYG{n}{d2min}\PYG{+w}{ }\PYG{n}{type}\PYG{+w}{ }\PYG{n}{fixed}\PYG{+w}{ }\PYG{l+m}{0}
\PYG{k}{fix }\PYG{n+nv+nvIdentifier}{1}\PYG{+w}{ }\PYG{n+nv+nvIdentifier}{all}\PYG{+w}{ }\PYG{n}{nonaffine}\PYG{o}{/}\PYG{n}{displacement}\PYG{+w}{ }\PYG{l+m}{1000}\PYG{+w}{ }\PYG{n}{d2min}\PYG{+w}{ }\PYG{n}{custom}\PYG{+w}{ }\PYG{l+m}{2.0}\PYG{+w}{ }\PYG{n}{offset}\PYG{+w}{ }\PYG{l+m}{100}
\end{sphinxVerbatim}


\subsection{Description}
\label{\detokenize{fix_nonaffine_displacement:description}}
\sphinxAtStartPar
\DUrole{versionmodified,added}{Added in version 7Feb2024.}

\sphinxAtStartPar
This fix computes different metrics of the nonaffine displacement of
particles. The first metric, \sphinxstyleemphasis{d2min} calculates the \(D^2_\mathrm{min}\)
nonaffine displacement by Falk and Langer in {\hyperref[\detokenize{fix_nonaffine_displacement:d2min-falk}]{\sphinxcrossref{\DUrole{std,std-ref}{(Falk)}}}}.
For each atom, the fix computes the two tensors
\begin{equation*}
\begin{split}X = \sum_{\mathrm{neighbors}} \vec{r} \left(\vec{r}_{0} \right)^T\end{split}
\end{equation*}
\sphinxAtStartPar
and
\begin{equation*}
\begin{split}Y = \sum_{\mathrm{neighbors}} \vec{r}_0 \left(\vec{r}_{0} \right)^T\end{split}
\end{equation*}
\sphinxAtStartPar
where the neighbors include all other atoms within the distance criterion
set by the cutoff option, discussed below, \(\vec{r}\) is the current
displacement between particles, and \(\vec{r}_0\) is the reference
displacement. A deformation gradient tensor is then calculated as
\(F = X Y^{-1}\) from which
\begin{equation*}
\begin{split}D^2_\mathrm{min} = \sum_{\mathrm{neighbors}} \left| \vec{r} - F \vec{r}_0 \right|^2\end{split}
\end{equation*}
\sphinxAtStartPar
and a strain tensor is calculated \(E = F F^{T} - I\) where \(I\)
is the identity tensor. This calculation is only performed on timesteps that
are a multiple of \sphinxstyleemphasis{nevery} (including timestep zero). Data accessed before
this occurs will simply be zeroed.

\sphinxAtStartPar
For particles with low coordination numbers, calculations of \(D^2_\mathrm{min}\)
may not be accurate. An optional minimum coordination number can be defined using
the \sphinxstyleemphasis{z/min} keyword. If any particle has fewer than the specified number of particles
in the cutoff distance or in contact, the above calculations will be skipped and the
corresponding peratom array entries will be zero.

\sphinxAtStartPar
The \sphinxstyleemphasis{integrated} style simply integrates the velocity of particles
every timestep to calculate a displacement. This style only works if
used in conjunction with another fix that deforms the box and displaces
atom positions such as {\hyperref[\detokenize{fix_deform::doc}]{\sphinxcrossref{\DUrole{doc}{fix deform}}}} with remap x,
{\hyperref[\detokenize{fix_press_berendsen::doc}]{\sphinxcrossref{\DUrole{doc}{fix press/berendsen}}}}, or {\hyperref[\detokenize{fix_nh::doc}]{\sphinxcrossref{\DUrole{doc}{fix nh}}}}.

\sphinxAtStartPar
Both of these methods require defining a reference state. With the \sphinxstyleemphasis{fixed} reference
style, the user picks a specific timestep \sphinxstyleemphasis{nstep} at which particle positions are saved.
If peratom data is accessed from this compute prior to this timestep, it will simply be
zeroed. The \sphinxstyleemphasis{update} reference style implies the reference state will be updated every
\sphinxstyleemphasis{nstep} timesteps. The \sphinxstyleemphasis{offset} reference will update the reference state \sphinxstyleemphasis{nstep}
timesteps before a multiple of \sphinxstyleemphasis{nevery} timesteps.


\bigskip\hrule\bigskip



\subsection{Restart, fix\_modify, output, run start/stop, minimize info}
\label{\detokenize{fix_nonaffine_displacement:restart-fix-modify-output-run-start-stop-minimize-info}}
\sphinxAtStartPar
The reference state is saved to {\hyperref[\detokenize{restart::doc}]{\sphinxcrossref{\DUrole{doc}{binary restart files}}}}.

\sphinxAtStartPar
None of the {\hyperref[\detokenize{fix_modify::doc}]{\sphinxcrossref{\DUrole{doc}{fix\_modify}}}} options are relevant to this
fix.

\sphinxAtStartPar
This fix computes a peratom array with 3 columns, which can be accessed
by indices 1\sphinxhyphen{}3 using any command that uses per\sphinxhyphen{}atom values from a fix
as input.

\sphinxAtStartPar
For the \sphinxstyleemphasis{integrated} style, the three columns are the nonaffine
displacements in the x, y, and z directions. For the \sphinxstyleemphasis{d2min} style,
the three columns are the calculated \(\sqrt{D^2_\mathrm{min}}\), the
volumetric strain, and the deviatoric strain.


\subsection{Restrictions}
\label{\detokenize{fix_nonaffine_displacement:restrictions}}
\sphinxAtStartPar
This compute is part of the EXTRA\sphinxhyphen{}FIX package.  It is only enabled if
LAMMPS was built with that package.  See the
{\hyperref[\detokenize{Build_package::doc}]{\sphinxcrossref{\DUrole{doc}{Build package}}}} page for more info.


\subsection{Related commands}
\label{\detokenize{fix_nonaffine_displacement:related-commands}}
\sphinxAtStartPar
none


\subsection{Default}
\label{\detokenize{fix_nonaffine_displacement:default}}
\sphinxAtStartPar
none


\bigskip\hrule\bigskip

\phantomsection\label{\detokenize{fix_nonaffine_displacement:d2min-falk}}
\sphinxAtStartPar
\sphinxstylestrong{(Falk)} Falk and Langer PRE, 57, 7192 (1998).

\sphinxstepscope

\index{fix nph/asphere@\spxentry{fix nph/asphere}}\ignorespaces 
\index{fix nph/asphere/omp@\spxentry{fix nph/asphere/omp}}\ignorespaces 

\section{fix nph/asphere command}
\label{\detokenize{fix_nph_asphere:fix-nph-asphere-command}}\label{\detokenize{fix_nph_asphere:index-1}}\label{\detokenize{fix_nph_asphere:index-0}}\label{\detokenize{fix_nph_asphere::doc}}
\sphinxAtStartPar
Accelerator Variants: \sphinxstyleemphasis{nph/asphere/omp}


\subsection{Syntax}
\label{\detokenize{fix_nph_asphere:syntax}}
\begin{sphinxVerbatim}[commandchars=\\\{\}]
\PYG{k}{fix }\PYG{n+nv+nvIdentifier}{ID}\PYG{+w}{ }\PYG{n+nv+nvIdentifier}{group\PYGZhy{}ID}\PYG{+w}{ }\PYG{n}{nph}\PYG{o}{/}\PYG{n}{asphere}\PYG{+w}{ }\PYG{n}{args}\PYG{+w}{ }\PYG{n}{keyword}\PYG{+w}{ }\PYG{n}{value}\PYG{+w}{ }\PYG{n}{...}
\end{sphinxVerbatim}
\begin{itemize}
\item {} 
\sphinxAtStartPar
ID, group\sphinxhyphen{}ID are documented in {\hyperref[\detokenize{fix::doc}]{\sphinxcrossref{\DUrole{doc}{fix}}}} command

\item {} 
\sphinxAtStartPar
nph/asphere = style name of this fix command

\item {} 
\sphinxAtStartPar
additional barostat related keyword/value pairs from the {\hyperref[\detokenize{fix_nh::doc}]{\sphinxcrossref{\DUrole{doc}{fix nph}}}} command can be appended

\end{itemize}


\subsection{Examples}
\label{\detokenize{fix_nph_asphere:examples}}
\begin{sphinxVerbatim}[commandchars=\\\{\}]
\PYG{k}{fix }\PYG{n+nv+nvIdentifier}{1}\PYG{+w}{ }\PYG{n+nv+nvIdentifier}{all}\PYG{+w}{ }\PYG{n}{nph}\PYG{o}{/}\PYG{n}{asphere}\PYG{+w}{ }\PYG{n}{iso}\PYG{+w}{ }\PYG{l+m}{0.0}\PYG{+w}{ }\PYG{l+m}{0.0}\PYG{+w}{ }\PYG{l+m}{1000.0}
\PYG{k}{fix }\PYG{n+nv+nvIdentifier}{2}\PYG{+w}{ }\PYG{n+nv+nvIdentifier}{all}\PYG{+w}{ }\PYG{n}{nph}\PYG{o}{/}\PYG{n}{asphere}\PYG{+w}{ }\PYG{n}{x}\PYG{+w}{ }\PYG{l+m}{5.0}\PYG{+w}{ }\PYG{l+m}{5.0}\PYG{+w}{ }\PYG{l+m}{1000.0}
\PYG{k}{fix }\PYG{n+nv+nvIdentifier}{2}\PYG{+w}{ }\PYG{n+nv+nvIdentifier}{all}\PYG{+w}{ }\PYG{n}{nph}\PYG{o}{/}\PYG{n}{asphere}\PYG{+w}{ }\PYG{n}{x}\PYG{+w}{ }\PYG{l+m}{5.0}\PYG{+w}{ }\PYG{l+m}{5.0}\PYG{+w}{ }\PYG{l+m}{1000.0}\PYG{+w}{ }\PYG{n}{drag}\PYG{+w}{ }\PYG{l+m}{0.2}
\PYG{k}{fix }\PYG{n+nv+nvIdentifier}{2}\PYG{+w}{ }\PYG{n+nv+nvIdentifier}{water}\PYG{+w}{ }\PYG{n}{nph}\PYG{o}{/}\PYG{n}{asphere}\PYG{+w}{ }\PYG{n}{aniso}\PYG{+w}{ }\PYG{l+m}{0.0}\PYG{+w}{ }\PYG{l+m}{0.0}\PYG{+w}{ }\PYG{l+m}{1000.0}\PYG{+w}{ }\PYG{n}{dilate}\PYG{+w}{ }\PYG{n}{partial}
\end{sphinxVerbatim}


\subsection{Description}
\label{\detokenize{fix_nph_asphere:description}}
\sphinxAtStartPar
Perform constant NPH integration to update position, velocity,
orientation, and angular velocity each timestep for aspherical or
ellipsoidal particles in the group using a Nose/Hoover pressure
barostat.  P is pressure; H is enthalpy.  This creates a system
trajectory consistent with the isenthalpic ensemble.

\sphinxAtStartPar
This fix differs from the {\hyperref[\detokenize{fix_nh::doc}]{\sphinxcrossref{\DUrole{doc}{fix nph}}}} command, which assumes
point particles and only updates their position and velocity.

\sphinxAtStartPar
Additional parameters affecting the barostat are specified by keywords
and values documented with the {\hyperref[\detokenize{fix_nh::doc}]{\sphinxcrossref{\DUrole{doc}{fix nph}}}} command.  See,
for example, discussion of the \sphinxstyleemphasis{aniso}, and \sphinxstyleemphasis{dilate} keywords.

\sphinxAtStartPar
The particles in the fix group are the only ones whose velocities and
positions are updated by the velocity/position update portion of the
NPH integration.

\sphinxAtStartPar
Regardless of what particles are in the fix group, a global pressure is
computed for all particles.  Similarly, when the size of the simulation
box is changed, all particles are re\sphinxhyphen{}scaled to new positions, unless the
keyword \sphinxstyleemphasis{dilate} is specified with a value of \sphinxstyleemphasis{partial}, in which case
only the particles in the fix group are re\sphinxhyphen{}scaled.  The latter can be
useful for leaving the coordinates of particles in a solid substrate
unchanged and controlling the pressure of a surrounding fluid.


\bigskip\hrule\bigskip


\sphinxAtStartPar
This fix computes a temperature and pressure each timestep.  To do
this, the fix creates its own computes of style “temp/asphere” and
“pressure”, as if these commands had been issued:

\begin{sphinxVerbatim}[commandchars=\\\{\}]
\PYG{k}{compute }\PYG{n+nv+nvIdentifier}{fix\PYGZhy{}ID\PYGZus{}temp}\PYG{+w}{ }\PYG{n+nv+nvIdentifier}{all}\PYG{+w}{ }\PYG{n}{temp}\PYG{o}{/}\PYG{n}{asphere}
\PYG{k}{compute }\PYG{n+nv+nvIdentifier}{fix\PYGZhy{}ID\PYGZus{}press}\PYG{+w}{ }\PYG{n+nv+nvIdentifier}{all}\PYG{+w}{ }\PYG{n}{pressure}\PYG{+w}{ }\PYG{n}{fix}\PYG{o}{\PYGZhy{}}\PYG{n}{ID\PYGZus{}temp}
\end{sphinxVerbatim}

\sphinxAtStartPar
See the {\hyperref[\detokenize{compute_temp_asphere::doc}]{\sphinxcrossref{\DUrole{doc}{compute temp/asphere}}}} and {\hyperref[\detokenize{compute_pressure::doc}]{\sphinxcrossref{\DUrole{doc}{compute pressure}}}} commands for details.  Note that the
IDs of the new computes are the fix\sphinxhyphen{}ID + underscore + “temp” or fix\_ID
+ underscore + “press”, and the group for the new computes is “all”
since pressure is computed for the entire system.

\sphinxAtStartPar
Note that these are NOT the computes used by thermodynamic output (see
the {\hyperref[\detokenize{thermo_style::doc}]{\sphinxcrossref{\DUrole{doc}{thermo\_style}}}} command) with ID = \sphinxstyleemphasis{thermo\_temp}
and \sphinxstyleemphasis{thermo\_press}.  This means you can change the attributes of this
fix’s temperature or pressure via the
{\hyperref[\detokenize{compute_modify::doc}]{\sphinxcrossref{\DUrole{doc}{compute\_modify}}}} command or print this temperature
or pressure during thermodynamic output via the {\hyperref[\detokenize{thermo_style::doc}]{\sphinxcrossref{\DUrole{doc}{thermo\_style custom}}}} command using the appropriate compute\sphinxhyphen{}ID.
It also means that changing attributes of \sphinxstyleemphasis{thermo\_temp} or
\sphinxstyleemphasis{thermo\_press} will have no effect on this fix.


\bigskip\hrule\bigskip


\sphinxAtStartPar
Styles with a \sphinxstyleemphasis{gpu}, \sphinxstyleemphasis{intel}, \sphinxstyleemphasis{kk}, \sphinxstyleemphasis{omp}, or \sphinxstyleemphasis{opt} suffix are
functionally the same as the corresponding style without the suffix.
They have been optimized to run faster, depending on your available
hardware, as discussed on the {\hyperref[\detokenize{Speed_packages::doc}]{\sphinxcrossref{\DUrole{doc}{Accelerator packages}}}}
page.  The accelerated styles take the same arguments and should
produce the same results, except for round\sphinxhyphen{}off and precision issues.

\sphinxAtStartPar
These accelerated styles are part of the GPU, INTEL, KOKKOS,
OPENMP, and OPT packages, respectively.  They are only enabled if
LAMMPS was built with those packages.  See the {\hyperref[\detokenize{Build_package::doc}]{\sphinxcrossref{\DUrole{doc}{Build package}}}} page for more info.

\sphinxAtStartPar
You can specify the accelerated styles explicitly in your input script
by including their suffix, or you can use the {\hyperref[\detokenize{Run_options::doc}]{\sphinxcrossref{\DUrole{doc}{\sphinxhyphen{}suffix command\sphinxhyphen{}line switch}}}} when you invoke LAMMPS, or you can use the
{\hyperref[\detokenize{suffix::doc}]{\sphinxcrossref{\DUrole{doc}{suffix}}}} command in your input script.

\sphinxAtStartPar
See the {\hyperref[\detokenize{Speed_packages::doc}]{\sphinxcrossref{\DUrole{doc}{Accelerator packages}}}} page for more
instructions on how to use the accelerated styles effectively.


\subsection{Restart, fix\_modify, output, run start/stop, minimize info}
\label{\detokenize{fix_nph_asphere:restart-fix-modify-output-run-start-stop-minimize-info}}
\sphinxAtStartPar
This fix writes the state of the Nose/Hoover barostat to {\hyperref[\detokenize{restart::doc}]{\sphinxcrossref{\DUrole{doc}{binary
restart files}}}}.  See the {\hyperref[\detokenize{read_restart::doc}]{\sphinxcrossref{\DUrole{doc}{read\_restart}}}}
command for info on how to re\sphinxhyphen{}specify a fix in an input script that
reads a restart file, so that the operation of the fix continues in an
uninterrupted fashion.

\sphinxAtStartPar
The {\hyperref[\detokenize{fix_modify::doc}]{\sphinxcrossref{\DUrole{doc}{fix\_modify}}}} \sphinxstyleemphasis{temp} and \sphinxstyleemphasis{press} options are
supported by this fix.  You can use them to assign a
{\hyperref[\detokenize{compute::doc}]{\sphinxcrossref{\DUrole{doc}{compute}}}} you have defined to this fix which will be used
in its thermostatting or barostatting procedure.  If you do this, note
that the kinetic energy derived from the compute temperature should be
consistent with the virial term computed using all atoms for the
pressure.  LAMMPS will warn you if you choose to compute temperature
on a subset of atoms.

\sphinxAtStartPar
The cumulative energy change in the system imposed by this fix is
included in the {\hyperref[\detokenize{thermo_style::doc}]{\sphinxcrossref{\DUrole{doc}{thermodynamic output}}}} keywords
\sphinxstyleemphasis{ecouple} and \sphinxstyleemphasis{econserve}.  See the {\hyperref[\detokenize{thermo_style::doc}]{\sphinxcrossref{\DUrole{doc}{thermo\_style}}}}
doc page for details.

\sphinxAtStartPar
This fix computes the same global scalar and global vector of
quantities as does the {\hyperref[\detokenize{fix_nh::doc}]{\sphinxcrossref{\DUrole{doc}{fix nph}}}} command.

\sphinxAtStartPar
This fix can ramp its target pressure over multiple runs, using the
\sphinxstyleemphasis{start} and \sphinxstyleemphasis{stop} keywords of the {\hyperref[\detokenize{run::doc}]{\sphinxcrossref{\DUrole{doc}{run}}}} command.  See the
{\hyperref[\detokenize{run::doc}]{\sphinxcrossref{\DUrole{doc}{run}}}} command for details of how to do this.

\sphinxAtStartPar
This fix is not invoked during {\hyperref[\detokenize{minimize::doc}]{\sphinxcrossref{\DUrole{doc}{energy minimization}}}}.


\subsection{Restrictions}
\label{\detokenize{fix_nph_asphere:restrictions}}
\sphinxAtStartPar
This fix is part of the ASPHERE package.  It is only enabled if LAMMPS
was built with that package.  See the {\hyperref[\detokenize{Build_package::doc}]{\sphinxcrossref{\DUrole{doc}{Build package}}}} page for more info.

\sphinxAtStartPar
This fix requires that atoms store torque and angular momentum and a
quaternion as defined by the {\hyperref[\detokenize{atom_style::doc}]{\sphinxcrossref{\DUrole{doc}{atom\_style ellipsoid}}}}
command.

\sphinxAtStartPar
All particles in the group must be finite\sphinxhyphen{}size.  They cannot be point
particles, but they can be aspherical or spherical as defined by their
shape attribute.


\subsection{Related commands}
\label{\detokenize{fix_nph_asphere:related-commands}}
\sphinxAtStartPar
{\hyperref[\detokenize{fix_nh::doc}]{\sphinxcrossref{\DUrole{doc}{fix nph}}}}, {\hyperref[\detokenize{fix_nve_asphere::doc}]{\sphinxcrossref{\DUrole{doc}{fix nve\_asphere}}}}, {\hyperref[\detokenize{fix_nvt_asphere::doc}]{\sphinxcrossref{\DUrole{doc}{fix nvt\_asphere}}}}, {\hyperref[\detokenize{fix_npt_asphere::doc}]{\sphinxcrossref{\DUrole{doc}{fix npt\_asphere}}}}, {\hyperref[\detokenize{fix_modify::doc}]{\sphinxcrossref{\DUrole{doc}{fix\_modify}}}}


\subsection{Default}
\label{\detokenize{fix_nph_asphere:default}}
\sphinxAtStartPar
none

\sphinxstepscope

\index{fix nph/body@\spxentry{fix nph/body}}\ignorespaces 

\section{fix nph/body command}
\label{\detokenize{fix_nph_body:fix-nph-body-command}}\label{\detokenize{fix_nph_body:index-0}}\label{\detokenize{fix_nph_body::doc}}

\subsection{Syntax}
\label{\detokenize{fix_nph_body:syntax}}
\begin{sphinxVerbatim}[commandchars=\\\{\}]
\PYG{k}{fix }\PYG{n+nv+nvIdentifier}{ID}\PYG{+w}{ }\PYG{n+nv+nvIdentifier}{group\PYGZhy{}ID}\PYG{+w}{ }\PYG{n}{nph}\PYG{o}{/}\PYG{n}{body}\PYG{+w}{ }\PYG{n}{args}\PYG{+w}{ }\PYG{n}{keyword}\PYG{+w}{ }\PYG{n}{value}\PYG{+w}{ }\PYG{n}{...}
\end{sphinxVerbatim}
\begin{itemize}
\item {} 
\sphinxAtStartPar
ID, group\sphinxhyphen{}ID are documented in {\hyperref[\detokenize{fix::doc}]{\sphinxcrossref{\DUrole{doc}{fix}}}} command

\item {} 
\sphinxAtStartPar
nph/body = style name of this fix command

\item {} 
\sphinxAtStartPar
additional barostat related keyword/value pairs from the {\hyperref[\detokenize{fix_nh::doc}]{\sphinxcrossref{\DUrole{doc}{fix nph}}}} command can be appended

\end{itemize}


\subsection{Examples}
\label{\detokenize{fix_nph_body:examples}}
\begin{sphinxVerbatim}[commandchars=\\\{\}]
\PYG{k}{fix }\PYG{n+nv+nvIdentifier}{1}\PYG{+w}{ }\PYG{n+nv+nvIdentifier}{all}\PYG{+w}{ }\PYG{n}{nph}\PYG{o}{/}\PYG{n}{body}\PYG{+w}{ }\PYG{n}{iso}\PYG{+w}{ }\PYG{l+m}{0.0}\PYG{+w}{ }\PYG{l+m}{0.0}\PYG{+w}{ }\PYG{l+m}{1000.0}
\PYG{k}{fix }\PYG{n+nv+nvIdentifier}{2}\PYG{+w}{ }\PYG{n+nv+nvIdentifier}{all}\PYG{+w}{ }\PYG{n}{nph}\PYG{o}{/}\PYG{n}{body}\PYG{+w}{ }\PYG{n}{x}\PYG{+w}{ }\PYG{l+m}{5.0}\PYG{+w}{ }\PYG{l+m}{5.0}\PYG{+w}{ }\PYG{l+m}{1000.0}
\PYG{k}{fix }\PYG{n+nv+nvIdentifier}{2}\PYG{+w}{ }\PYG{n+nv+nvIdentifier}{all}\PYG{+w}{ }\PYG{n}{nph}\PYG{o}{/}\PYG{n}{body}\PYG{+w}{ }\PYG{n}{x}\PYG{+w}{ }\PYG{l+m}{5.0}\PYG{+w}{ }\PYG{l+m}{5.0}\PYG{+w}{ }\PYG{l+m}{1000.0}\PYG{+w}{ }\PYG{n}{drag}\PYG{+w}{ }\PYG{l+m}{0.2}
\PYG{k}{fix }\PYG{n+nv+nvIdentifier}{2}\PYG{+w}{ }\PYG{n+nv+nvIdentifier}{water}\PYG{+w}{ }\PYG{n}{nph}\PYG{o}{/}\PYG{n}{body}\PYG{+w}{ }\PYG{n}{aniso}\PYG{+w}{ }\PYG{l+m}{0.0}\PYG{+w}{ }\PYG{l+m}{0.0}\PYG{+w}{ }\PYG{l+m}{1000.0}\PYG{+w}{ }\PYG{n}{dilate}\PYG{+w}{ }\PYG{n}{partial}
\end{sphinxVerbatim}


\subsection{Description}
\label{\detokenize{fix_nph_body:description}}
\sphinxAtStartPar
Perform constant NPH integration to update position, velocity,
orientation, and angular velocity each timestep for body
particles in the group using a Nose/Hoover pressure
barostat.  P is pressure; H is enthalpy.  This creates a system
trajectory consistent with the isenthalpic ensemble.

\sphinxAtStartPar
This fix differs from the {\hyperref[\detokenize{fix_nh::doc}]{\sphinxcrossref{\DUrole{doc}{fix nph}}}} command, which assumes
point particles and only updates their position and velocity.

\sphinxAtStartPar
Additional parameters affecting the barostat are specified by keywords
and values documented with the {\hyperref[\detokenize{fix_nh::doc}]{\sphinxcrossref{\DUrole{doc}{fix nph}}}} command.  See,
for example, discussion of the \sphinxstyleemphasis{aniso}, and \sphinxstyleemphasis{dilate} keywords.

\sphinxAtStartPar
The particles in the fix group are the only ones whose velocities and
positions are updated by the velocity/position update portion of the
NPH integration.

\sphinxAtStartPar
Regardless of what particles are in the fix group, a global pressure is
computed for all particles.  Similarly, when the size of the simulation
box is changed, all particles are re\sphinxhyphen{}scaled to new positions, unless the
keyword \sphinxstyleemphasis{dilate} is specified with a value of \sphinxstyleemphasis{partial}, in which case
only the particles in the fix group are re\sphinxhyphen{}scaled.  The latter can be
useful for leaving the coordinates of particles in a solid substrate
unchanged and controlling the pressure of a surrounding fluid.


\bigskip\hrule\bigskip


\sphinxAtStartPar
This fix computes a temperature and pressure each timestep.  To do
this, the fix creates its own computes of style “temp/body” and
“pressure”, as if these commands had been issued:

\begin{sphinxVerbatim}[commandchars=\\\{\}]
\PYG{k}{compute }\PYG{n+nv+nvIdentifier}{fix\PYGZhy{}ID\PYGZus{}temp}\PYG{+w}{ }\PYG{n+nv+nvIdentifier}{all}\PYG{+w}{ }\PYG{n}{temp}\PYG{o}{/}\PYG{n}{body}
\PYG{k}{compute }\PYG{n+nv+nvIdentifier}{fix\PYGZhy{}ID\PYGZus{}press}\PYG{+w}{ }\PYG{n+nv+nvIdentifier}{all}\PYG{+w}{ }\PYG{n}{pressure}\PYG{+w}{ }\PYG{n}{fix}\PYG{o}{\PYGZhy{}}\PYG{n}{ID\PYGZus{}temp}
\end{sphinxVerbatim}

\sphinxAtStartPar
See the {\hyperref[\detokenize{compute_temp_body::doc}]{\sphinxcrossref{\DUrole{doc}{compute temp/body}}}} and {\hyperref[\detokenize{compute_pressure::doc}]{\sphinxcrossref{\DUrole{doc}{compute pressure}}}} commands for details.  Note that the
IDs of the new computes are the fix\sphinxhyphen{}ID + underscore + “temp” or fix\_ID
+ underscore + “press”, and the group for the new computes is “all”
since pressure is computed for the entire system.

\sphinxAtStartPar
Note that these are NOT the computes used by thermodynamic output (see
the {\hyperref[\detokenize{thermo_style::doc}]{\sphinxcrossref{\DUrole{doc}{thermo\_style}}}} command) with ID = \sphinxstyleemphasis{thermo\_temp}
and \sphinxstyleemphasis{thermo\_press}.  This means you can change the attributes of this
fix’s temperature or pressure via the
{\hyperref[\detokenize{compute_modify::doc}]{\sphinxcrossref{\DUrole{doc}{compute\_modify}}}} command or print this temperature
or pressure during thermodynamic output via the {\hyperref[\detokenize{thermo_style::doc}]{\sphinxcrossref{\DUrole{doc}{thermo\_style custom}}}} command using the appropriate compute\sphinxhyphen{}ID.
It also means that changing attributes of \sphinxstyleemphasis{thermo\_temp} or
\sphinxstyleemphasis{thermo\_press} will have no effect on this fix.


\bigskip\hrule\bigskip


\sphinxAtStartPar
Styles with a \sphinxstyleemphasis{gpu}, \sphinxstyleemphasis{intel}, \sphinxstyleemphasis{kk}, \sphinxstyleemphasis{omp}, or \sphinxstyleemphasis{opt} suffix are
functionally the same as the corresponding style without the suffix.
They have been optimized to run faster, depending on your available
hardware, as discussed on the {\hyperref[\detokenize{Speed_packages::doc}]{\sphinxcrossref{\DUrole{doc}{Accelerator packages}}}}
page.  The accelerated styles take the same arguments and should
produce the same results, except for round\sphinxhyphen{}off and precision issues.

\sphinxAtStartPar
These accelerated styles are part of the GPU, INTEL, KOKKOS,
OPENMP, and OPT packages, respectively.  They are only enabled if
LAMMPS was built with those packages.  See the {\hyperref[\detokenize{Build_package::doc}]{\sphinxcrossref{\DUrole{doc}{Build package}}}} page for more info.

\sphinxAtStartPar
You can specify the accelerated styles explicitly in your input script
by including their suffix, or you can use the {\hyperref[\detokenize{Run_options::doc}]{\sphinxcrossref{\DUrole{doc}{\sphinxhyphen{}suffix command\sphinxhyphen{}line switch}}}} when you invoke LAMMPS, or you can use the
{\hyperref[\detokenize{suffix::doc}]{\sphinxcrossref{\DUrole{doc}{suffix}}}} command in your input script.

\sphinxAtStartPar
See the {\hyperref[\detokenize{Speed_packages::doc}]{\sphinxcrossref{\DUrole{doc}{Accelerator packages}}}} page for more
instructions on how to use the accelerated styles effectively.


\subsection{Restart, fix\_modify, output, run start/stop, minimize info}
\label{\detokenize{fix_nph_body:restart-fix-modify-output-run-start-stop-minimize-info}}
\sphinxAtStartPar
This fix writes the state of the Nose/Hoover barostat to {\hyperref[\detokenize{restart::doc}]{\sphinxcrossref{\DUrole{doc}{binary
restart files}}}}.  See the {\hyperref[\detokenize{read_restart::doc}]{\sphinxcrossref{\DUrole{doc}{read\_restart}}}}
command for info on how to re\sphinxhyphen{}specify a fix in an input script that
reads a restart file, so that the operation of the fix continues in an
uninterrupted fashion.

\sphinxAtStartPar
The {\hyperref[\detokenize{fix_modify::doc}]{\sphinxcrossref{\DUrole{doc}{fix\_modify}}}} \sphinxstyleemphasis{temp} and \sphinxstyleemphasis{press} options are
supported by this fix.  You can use them to assign a
{\hyperref[\detokenize{compute::doc}]{\sphinxcrossref{\DUrole{doc}{compute}}}} you have defined to this fix which will be used
in its thermostatting or barostatting procedure.  If you do this, note
that the kinetic energy derived from the compute temperature should be
consistent with the virial term computed using all atoms for the
pressure.  LAMMPS will warn you if you choose to compute temperature
on a subset of atoms.

\sphinxAtStartPar
The cumulative energy change in the system imposed by this fix is
included in the {\hyperref[\detokenize{thermo_style::doc}]{\sphinxcrossref{\DUrole{doc}{thermodynamic output}}}} keywords
\sphinxstyleemphasis{ecouple} and \sphinxstyleemphasis{econserve}.  See the {\hyperref[\detokenize{thermo_style::doc}]{\sphinxcrossref{\DUrole{doc}{thermo\_style}}}}
doc page for details.

\sphinxAtStartPar
This fix computes the same global scalar and global vector of
quantities as does the {\hyperref[\detokenize{fix_nh::doc}]{\sphinxcrossref{\DUrole{doc}{fix nph}}}} command.

\sphinxAtStartPar
This fix can ramp its target pressure over multiple runs, using the
\sphinxstyleemphasis{start} and \sphinxstyleemphasis{stop} keywords of the {\hyperref[\detokenize{run::doc}]{\sphinxcrossref{\DUrole{doc}{run}}}} command.  See the
{\hyperref[\detokenize{run::doc}]{\sphinxcrossref{\DUrole{doc}{run}}}} command for details of how to do this.

\sphinxAtStartPar
This fix is not invoked during {\hyperref[\detokenize{minimize::doc}]{\sphinxcrossref{\DUrole{doc}{energy minimization}}}}.


\subsection{Restrictions}
\label{\detokenize{fix_nph_body:restrictions}}
\sphinxAtStartPar
This fix is part of the BODY package.  It is only enabled if LAMMPS
was built with that package.  See the {\hyperref[\detokenize{Build_package::doc}]{\sphinxcrossref{\DUrole{doc}{Build package}}}} page for more info.

\sphinxAtStartPar
This fix requires that atoms store torque and angular momentum and a
quaternion as defined by the {\hyperref[\detokenize{atom_style::doc}]{\sphinxcrossref{\DUrole{doc}{atom\_style body}}}}
command.


\subsection{Related commands}
\label{\detokenize{fix_nph_body:related-commands}}
\sphinxAtStartPar
{\hyperref[\detokenize{fix_nh::doc}]{\sphinxcrossref{\DUrole{doc}{fix nph}}}}, {\hyperref[\detokenize{fix_nve_body::doc}]{\sphinxcrossref{\DUrole{doc}{fix nve\_body}}}}, {\hyperref[\detokenize{fix_nvt_body::doc}]{\sphinxcrossref{\DUrole{doc}{fix nvt\_body}}}}, {\hyperref[\detokenize{fix_npt_body::doc}]{\sphinxcrossref{\DUrole{doc}{fix npt\_body}}}}, {\hyperref[\detokenize{fix_modify::doc}]{\sphinxcrossref{\DUrole{doc}{fix\_modify}}}}


\subsection{Default}
\label{\detokenize{fix_nph_body:default}}
\sphinxAtStartPar
none

\sphinxstepscope

\index{fix nph/sphere@\spxentry{fix nph/sphere}}\ignorespaces 
\index{fix nph/sphere/omp@\spxentry{fix nph/sphere/omp}}\ignorespaces 

\section{fix nph/sphere command}
\label{\detokenize{fix_nph_sphere:fix-nph-sphere-command}}\label{\detokenize{fix_nph_sphere:index-1}}\label{\detokenize{fix_nph_sphere:index-0}}\label{\detokenize{fix_nph_sphere::doc}}
\sphinxAtStartPar
Accelerator Variants: \sphinxstyleemphasis{nph/sphere/omp}


\subsection{Syntax}
\label{\detokenize{fix_nph_sphere:syntax}}
\begin{sphinxVerbatim}[commandchars=\\\{\}]
\PYG{k}{fix }\PYG{n+nv+nvIdentifier}{ID}\PYG{+w}{ }\PYG{n+nv+nvIdentifier}{group\PYGZhy{}ID}\PYG{+w}{ }\PYG{n}{nph}\PYG{o}{/}\PYG{n}{sphere}\PYG{+w}{ }\PYG{n}{args}\PYG{+w}{ }\PYG{n}{keyword}\PYG{+w}{ }\PYG{n}{value}\PYG{+w}{ }\PYG{n}{...}
\end{sphinxVerbatim}
\begin{itemize}
\item {} 
\sphinxAtStartPar
ID, group\sphinxhyphen{}ID are documented in {\hyperref[\detokenize{fix::doc}]{\sphinxcrossref{\DUrole{doc}{fix}}}} command

\item {} 
\sphinxAtStartPar
nph/sphere = style name of this fix command

\item {} 
\sphinxAtStartPar
keyword = \sphinxstyleemphasis{disc}
\begin{sphinxalltt}
\sphinxstyleemphasis{disc} value = none = treat particles as 2d discs, not spheres
\end{sphinxalltt}

\item {} 
\sphinxAtStartPar
additional barostat related keyword/value pairs from the {\hyperref[\detokenize{fix_nh::doc}]{\sphinxcrossref{\DUrole{doc}{fix nph}}}} command can be appended

\end{itemize}


\subsection{Examples}
\label{\detokenize{fix_nph_sphere:examples}}
\begin{sphinxVerbatim}[commandchars=\\\{\}]
\PYG{k}{fix }\PYG{n+nv+nvIdentifier}{1}\PYG{+w}{ }\PYG{n+nv+nvIdentifier}{all}\PYG{+w}{ }\PYG{n}{nph}\PYG{o}{/}\PYG{n}{sphere}\PYG{+w}{ }\PYG{n}{iso}\PYG{+w}{ }\PYG{l+m}{0.0}\PYG{+w}{ }\PYG{l+m}{0.0}\PYG{+w}{ }\PYG{l+m}{1000.0}
\PYG{k}{fix }\PYG{n+nv+nvIdentifier}{2}\PYG{+w}{ }\PYG{n+nv+nvIdentifier}{all}\PYG{+w}{ }\PYG{n}{nph}\PYG{o}{/}\PYG{n}{sphere}\PYG{+w}{ }\PYG{n}{x}\PYG{+w}{ }\PYG{l+m}{5.0}\PYG{+w}{ }\PYG{l+m}{5.0}\PYG{+w}{ }\PYG{l+m}{1000.0}
\PYG{k}{fix }\PYG{n+nv+nvIdentifier}{2}\PYG{+w}{ }\PYG{n+nv+nvIdentifier}{all}\PYG{+w}{ }\PYG{n}{nph}\PYG{o}{/}\PYG{n}{sphere}\PYG{+w}{ }\PYG{n}{x}\PYG{+w}{ }\PYG{l+m}{5.0}\PYG{+w}{ }\PYG{l+m}{5.0}\PYG{+w}{ }\PYG{l+m}{1000.0}\PYG{+w}{ }\PYG{n}{disc}
\PYG{k}{fix }\PYG{n+nv+nvIdentifier}{2}\PYG{+w}{ }\PYG{n+nv+nvIdentifier}{all}\PYG{+w}{ }\PYG{n}{nph}\PYG{o}{/}\PYG{n}{sphere}\PYG{+w}{ }\PYG{n}{x}\PYG{+w}{ }\PYG{l+m}{5.0}\PYG{+w}{ }\PYG{l+m}{5.0}\PYG{+w}{ }\PYG{l+m}{1000.0}\PYG{+w}{ }\PYG{n}{drag}\PYG{+w}{ }\PYG{l+m}{0.2}
\PYG{k}{fix }\PYG{n+nv+nvIdentifier}{2}\PYG{+w}{ }\PYG{n+nv+nvIdentifier}{water}\PYG{+w}{ }\PYG{n}{nph}\PYG{o}{/}\PYG{n}{sphere}\PYG{+w}{ }\PYG{n}{aniso}\PYG{+w}{ }\PYG{l+m}{0.0}\PYG{+w}{ }\PYG{l+m}{0.0}\PYG{+w}{ }\PYG{l+m}{1000.0}\PYG{+w}{ }\PYG{n}{dilate}\PYG{+w}{ }\PYG{n}{partial}
\end{sphinxVerbatim}


\subsection{Description}
\label{\detokenize{fix_nph_sphere:description}}
\sphinxAtStartPar
Perform constant NPH integration to update position, velocity, and
angular velocity each timestep for finite\sphinxhyphen{}size spherical particles in
the group using a Nose/Hoover pressure barostat.  P is pressure; H is
enthalpy.  This creates a system trajectory consistent with the
isenthalpic ensemble.

\sphinxAtStartPar
This fix differs from the {\hyperref[\detokenize{fix_nh::doc}]{\sphinxcrossref{\DUrole{doc}{fix nph}}}} command, which assumes
point particles and only updates their position and velocity.

\sphinxAtStartPar
If the \sphinxstyleemphasis{disc} keyword is used, then each particle is treated as a 2d
disc (circle) instead of as a sphere.  This is only possible for 2d
simulations, as defined by the {\hyperref[\detokenize{dimension::doc}]{\sphinxcrossref{\DUrole{doc}{dimension}}}} keyword.
The only difference between discs and spheres in this context is their
moment of inertia, as used in the time integration.

\sphinxAtStartPar
Additional parameters affecting the barostat are specified by keywords
and values documented with the {\hyperref[\detokenize{fix_nh::doc}]{\sphinxcrossref{\DUrole{doc}{fix nph}}}} command.  See,
for example, discussion of the \sphinxstyleemphasis{aniso}, and \sphinxstyleemphasis{dilate} keywords.

\sphinxAtStartPar
The particles in the fix group are the only ones whose velocities and
positions are updated by the velocity/position update portion of the
NPH integration.

\sphinxAtStartPar
Regardless of what particles are in the fix group, a global pressure is
computed for all particles.  Similarly, when the size of the simulation
box is changed, all particles are re\sphinxhyphen{}scaled to new positions, unless the
keyword \sphinxstyleemphasis{dilate} is specified with a value of \sphinxstyleemphasis{partial}, in which case
only the particles in the fix group are re\sphinxhyphen{}scaled.  The latter can be
useful for leaving the coordinates of particles in a solid substrate
unchanged and controlling the pressure of a surrounding fluid.


\bigskip\hrule\bigskip


\sphinxAtStartPar
This fix computes a temperature and pressure each timestep.  To do
this, the fix creates its own computes of style “temp/sphere” and
“pressure”, as if these commands had been issued:

\begin{sphinxVerbatim}[commandchars=\\\{\}]
\PYG{k}{compute }\PYG{n+nv+nvIdentifier}{fix\PYGZhy{}ID\PYGZus{}temp}\PYG{+w}{ }\PYG{n+nv+nvIdentifier}{all}\PYG{+w}{ }\PYG{n}{temp}\PYG{o}{/}\PYG{n}{sphere}
\PYG{k}{compute }\PYG{n+nv+nvIdentifier}{fix\PYGZhy{}ID\PYGZus{}press}\PYG{+w}{ }\PYG{n+nv+nvIdentifier}{all}\PYG{+w}{ }\PYG{n}{pressure}\PYG{+w}{ }\PYG{n}{fix}\PYG{o}{\PYGZhy{}}\PYG{n}{ID\PYGZus{}temp}
\end{sphinxVerbatim}

\sphinxAtStartPar
See the {\hyperref[\detokenize{compute_temp_sphere::doc}]{\sphinxcrossref{\DUrole{doc}{compute temp/sphere}}}} and {\hyperref[\detokenize{compute_pressure::doc}]{\sphinxcrossref{\DUrole{doc}{compute pressure}}}} commands for details.  Note that the
IDs of the new computes are the fix\sphinxhyphen{}ID + underscore + “temp” or fix\_ID
+ underscore + “press”, and the group for the new computes is “all”
since pressure is computed for the entire system.

\sphinxAtStartPar
Note that these are NOT the computes used by thermodynamic output (see
the {\hyperref[\detokenize{thermo_style::doc}]{\sphinxcrossref{\DUrole{doc}{thermo\_style}}}} command) with ID = \sphinxstyleemphasis{thermo\_temp}
and \sphinxstyleemphasis{thermo\_press}.  This means you can change the attributes of this
fix’s temperature or pressure via the
{\hyperref[\detokenize{compute_modify::doc}]{\sphinxcrossref{\DUrole{doc}{compute\_modify}}}} command or print this temperature
or pressure during thermodynamic output via the {\hyperref[\detokenize{thermo_style::doc}]{\sphinxcrossref{\DUrole{doc}{thermo\_style custom}}}} command using the appropriate compute\sphinxhyphen{}ID.
It also means that changing attributes of \sphinxstyleemphasis{thermo\_temp} or
\sphinxstyleemphasis{thermo\_press} will have no effect on this fix.


\bigskip\hrule\bigskip


\sphinxAtStartPar
Styles with a \sphinxstyleemphasis{gpu}, \sphinxstyleemphasis{intel}, \sphinxstyleemphasis{kk}, \sphinxstyleemphasis{omp}, or \sphinxstyleemphasis{opt} suffix are
functionally the same as the corresponding style without the suffix.
They have been optimized to run faster, depending on your available
hardware, as discussed on the {\hyperref[\detokenize{Speed_packages::doc}]{\sphinxcrossref{\DUrole{doc}{Accelerator packages}}}}
page.  The accelerated styles take the same arguments and should
produce the same results, except for round\sphinxhyphen{}off and precision issues.

\sphinxAtStartPar
These accelerated styles are part of the GPU, INTEL, KOKKOS,
OPENMP, and OPT packages, respectively.  They are only enabled if
LAMMPS was built with those packages.  See the {\hyperref[\detokenize{Build_package::doc}]{\sphinxcrossref{\DUrole{doc}{Build package}}}} page for more info.

\sphinxAtStartPar
You can specify the accelerated styles explicitly in your input script
by including their suffix, or you can use the {\hyperref[\detokenize{Run_options::doc}]{\sphinxcrossref{\DUrole{doc}{\sphinxhyphen{}suffix command\sphinxhyphen{}line switch}}}} when you invoke LAMMPS, or you can use the
{\hyperref[\detokenize{suffix::doc}]{\sphinxcrossref{\DUrole{doc}{suffix}}}} command in your input script.

\sphinxAtStartPar
See the {\hyperref[\detokenize{Speed_packages::doc}]{\sphinxcrossref{\DUrole{doc}{Accelerator packages}}}} page for more
instructions on how to use the accelerated styles effectively.


\subsection{Restart, fix\_modify, output, run start/stop, minimize info}
\label{\detokenize{fix_nph_sphere:restart-fix-modify-output-run-start-stop-minimize-info}}
\sphinxAtStartPar
This fix writes the state of the Nose/Hoover barostat to {\hyperref[\detokenize{restart::doc}]{\sphinxcrossref{\DUrole{doc}{binary
restart files}}}}.  See the {\hyperref[\detokenize{read_restart::doc}]{\sphinxcrossref{\DUrole{doc}{read\_restart}}}}
command for info on how to re\sphinxhyphen{}specify a fix in an input script that
reads a restart file, so that the operation of the fix continues in an
uninterrupted fashion.

\sphinxAtStartPar
The {\hyperref[\detokenize{fix_modify::doc}]{\sphinxcrossref{\DUrole{doc}{fix\_modify}}}} \sphinxstyleemphasis{temp} and \sphinxstyleemphasis{press} options are
supported by this fix.  You can use them to assign a
{\hyperref[\detokenize{compute::doc}]{\sphinxcrossref{\DUrole{doc}{compute}}}} you have defined to this fix which will be used
in its thermostatting or barostatting procedure.  If you do this, note
that the kinetic energy derived from the compute temperature should be
consistent with the virial term computed using all atoms for the
pressure.  LAMMPS will warn you if you choose to compute temperature
on a subset of atoms.

\sphinxAtStartPar
The cumulative energy change in the system imposed by this fix is
included in the {\hyperref[\detokenize{thermo_style::doc}]{\sphinxcrossref{\DUrole{doc}{thermodynamic output}}}} keywords
\sphinxstyleemphasis{ecouple} and \sphinxstyleemphasis{econserve}.  See the {\hyperref[\detokenize{thermo_style::doc}]{\sphinxcrossref{\DUrole{doc}{thermo\_style}}}}
doc page for details.

\sphinxAtStartPar
This fix computes the same global scalar and global vector of
quantities as does the {\hyperref[\detokenize{fix_nh::doc}]{\sphinxcrossref{\DUrole{doc}{fix nph}}}} command.

\sphinxAtStartPar
This fix can ramp its target pressure over multiple runs, using the
\sphinxstyleemphasis{start} and \sphinxstyleemphasis{stop} keywords of the {\hyperref[\detokenize{run::doc}]{\sphinxcrossref{\DUrole{doc}{run}}}} command.  See the
{\hyperref[\detokenize{run::doc}]{\sphinxcrossref{\DUrole{doc}{run}}}} command for details of how to do this.

\sphinxAtStartPar
This fix is not invoked during {\hyperref[\detokenize{minimize::doc}]{\sphinxcrossref{\DUrole{doc}{energy minimization}}}}.


\subsection{Restrictions}
\label{\detokenize{fix_nph_sphere:restrictions}}
\sphinxAtStartPar
This fix requires that atoms store torque and angular velocity (omega)
and a radius as defined by the {\hyperref[\detokenize{atom_style::doc}]{\sphinxcrossref{\DUrole{doc}{atom\_style sphere}}}}
command.

\sphinxAtStartPar
All particles in the group must be finite\sphinxhyphen{}size spheres.  They cannot
be point particles.

\sphinxAtStartPar
Use of the \sphinxstyleemphasis{disc} keyword is only allowed for 2d simulations, as
defined by the {\hyperref[\detokenize{dimension::doc}]{\sphinxcrossref{\DUrole{doc}{dimension}}}} keyword.


\subsection{Related commands}
\label{\detokenize{fix_nph_sphere:related-commands}}
\sphinxAtStartPar
{\hyperref[\detokenize{fix_nh::doc}]{\sphinxcrossref{\DUrole{doc}{fix nph}}}}, {\hyperref[\detokenize{fix_nve_sphere::doc}]{\sphinxcrossref{\DUrole{doc}{fix nve\_sphere}}}}, {\hyperref[\detokenize{fix_nvt_sphere::doc}]{\sphinxcrossref{\DUrole{doc}{fix nvt\_sphere}}}}, {\hyperref[\detokenize{fix_npt_sphere::doc}]{\sphinxcrossref{\DUrole{doc}{fix npt\_sphere}}}},
{\hyperref[\detokenize{fix_modify::doc}]{\sphinxcrossref{\DUrole{doc}{fix\_modify}}}}


\subsection{Default}
\label{\detokenize{fix_nph_sphere:default}}
\sphinxAtStartPar
none

\sphinxstepscope

\index{fix nphug@\spxentry{fix nphug}}\ignorespaces 
\index{fix nphug/omp@\spxentry{fix nphug/omp}}\ignorespaces 

\section{fix nphug command}
\label{\detokenize{fix_nphug:fix-nphug-command}}\label{\detokenize{fix_nphug:index-1}}\label{\detokenize{fix_nphug:index-0}}\label{\detokenize{fix_nphug::doc}}
\sphinxAtStartPar
Accelerator Variants: \sphinxstyleemphasis{nphug/omp}


\subsection{Syntax}
\label{\detokenize{fix_nphug:syntax}}
\begin{sphinxVerbatim}[commandchars=\\\{\}]
\PYG{k}{fix }\PYG{n+nv+nvIdentifier}{ID}\PYG{+w}{ }\PYG{n+nv+nvIdentifier}{group\PYGZhy{}ID}\PYG{+w}{ }\PYG{n}{nphug}\PYG{+w}{ }\PYG{n}{keyword}\PYG{+w}{ }\PYG{n}{value}\PYG{+w}{ }\PYG{n}{...}
\end{sphinxVerbatim}
\begin{itemize}
\item {} 
\sphinxAtStartPar
ID, group\sphinxhyphen{}ID are documented in {\hyperref[\detokenize{fix::doc}]{\sphinxcrossref{\DUrole{doc}{fix}}}} command
\begin{sphinxalltt}
one or more keyword value pairs may be appended
keyword = \sphinxstyleemphasis{temp} or \sphinxstyleemphasis{iso} or \sphinxstyleemphasis{aniso} or \sphinxstyleemphasis{tri} or \sphinxstyleemphasis{x} or \sphinxstyleemphasis{y} or \sphinxstyleemphasis{z} or \sphinxstyleemphasis{couple} or \sphinxstyleemphasis{tchain} or \sphinxstyleemphasis{pchain} or \sphinxstyleemphasis{mtk} or \sphinxstyleemphasis{tloop} or \sphinxstyleemphasis{ploop} or \sphinxstyleemphasis{nreset} or \sphinxstyleemphasis{drag} or \sphinxstyleemphasis{dilate} or \sphinxstyleemphasis{scaleyz} or \sphinxstyleemphasis{scalexz} or \sphinxstyleemphasis{scalexy}
  \sphinxstyleemphasis{temp} values = Value1 Value2 Tdamp
    Value1, Value2 = Nose\sphinxhyphen{}Hoover target temperatures, ignored by Hugoniostat
    Tdamp = temperature damping parameter (time units)
  \sphinxstyleemphasis{iso} or \sphinxstyleemphasis{aniso} or \sphinxstyleemphasis{tri} values = Pstart Pstop Pdamp
    Pstart,Pstop = scalar external pressures, must be equal (pressure units)
    Pdamp = pressure damping parameter (time units)
  \sphinxstyleemphasis{x} or \sphinxstyleemphasis{y} or \sphinxstyleemphasis{z} or \sphinxstyleemphasis{xy} or \sphinxstyleemphasis{yz} or \sphinxstyleemphasis{xz} values = Pstart Pstop Pdamp
    Pstart,Pstop = external stress tensor components, must be equal (pressure units)
    Pdamp = stress damping parameter (time units)
  \sphinxstyleemphasis{couple} = \sphinxstyleemphasis{none} or \sphinxstyleemphasis{xyz} or \sphinxstyleemphasis{xy} or \sphinxstyleemphasis{yz} or \sphinxstyleemphasis{xz}
  \sphinxstyleemphasis{tchain} value = length of thermostat chain (1 = single thermostat)
  \sphinxstyleemphasis{pchain} values = length of thermostat chain on barostat (0 = no thermostat)
  \sphinxstyleemphasis{mtk} value = \sphinxstyleemphasis{yes} or \sphinxstyleemphasis{no} = add in MTK adjustment term or not
  \sphinxstyleemphasis{tloop} value = number of sub\sphinxhyphen{}cycles to perform on thermostat
  \sphinxstyleemphasis{ploop} value = number of sub\sphinxhyphen{}cycles to perform on barostat thermostat
  \sphinxstyleemphasis{nreset} value = reset reference cell every this many timesteps
  \sphinxstyleemphasis{drag} value = drag factor added to barostat/thermostat (0.0 = no drag)
  \sphinxstyleemphasis{dilate} value = \sphinxstyleemphasis{all} or \sphinxstyleemphasis{partial}
  \sphinxstyleemphasis{scaleyz} value = \sphinxstyleemphasis{yes} or \sphinxstyleemphasis{no} = scale yz with lz
  \sphinxstyleemphasis{scalexz} value = \sphinxstyleemphasis{yes} or \sphinxstyleemphasis{no} = scale xz with lz
  \sphinxstyleemphasis{scalexy} value = \sphinxstyleemphasis{yes} or \sphinxstyleemphasis{no} = scale xy with ly
\end{sphinxalltt}

\end{itemize}


\subsection{Examples}
\label{\detokenize{fix_nphug:examples}}
\begin{sphinxVerbatim}[commandchars=\\\{\}]
\PYG{k}{fix }\PYG{n+nv+nvIdentifier}{myhug}\PYG{+w}{ }\PYG{n+nv+nvIdentifier}{all}\PYG{+w}{ }\PYG{n}{nphug}\PYG{+w}{ }\PYG{n}{temp}\PYG{+w}{ }\PYG{l+m}{1.0}\PYG{+w}{ }\PYG{l+m}{1.0}\PYG{+w}{ }\PYG{l+m}{10.0}\PYG{+w}{ }\PYG{n}{z}\PYG{+w}{ }\PYG{l+m}{40.0}\PYG{+w}{ }\PYG{l+m}{40.0}\PYG{+w}{ }\PYG{l+m}{70.0}
\PYG{k}{fix }\PYG{n+nv+nvIdentifier}{myhug}\PYG{+w}{ }\PYG{n+nv+nvIdentifier}{all}\PYG{+w}{ }\PYG{n}{nphug}\PYG{+w}{ }\PYG{n}{temp}\PYG{+w}{ }\PYG{l+m}{1.0}\PYG{+w}{ }\PYG{l+m}{1.0}\PYG{+w}{ }\PYG{l+m}{10.0}\PYG{+w}{ }\PYG{n}{iso}\PYG{+w}{ }\PYG{l+m}{40.0}\PYG{+w}{ }\PYG{l+m}{40.0}\PYG{+w}{ }\PYG{l+m}{70.0}\PYG{+w}{ }\PYG{n}{drag}\PYG{+w}{ }\PYG{l+m}{200.0}\PYG{+w}{ }\PYG{n}{tchain}\PYG{+w}{ }\PYG{l+m}{1}\PYG{+w}{ }\PYG{n}{pchain}\PYG{+w}{ }\PYG{l+m}{0}
\end{sphinxVerbatim}


\subsection{Description}
\label{\detokenize{fix_nphug:description}}
\sphinxAtStartPar
This command is a variant of the Nose\sphinxhyphen{}Hoover
{\hyperref[\detokenize{fix_nh::doc}]{\sphinxcrossref{\DUrole{doc}{fix npt}}}} fix style.
It performs time integration of the Hugoniostat equations
of motion developed by Ravelo et al. {\hyperref[\detokenize{fix_nphug:ravelo1}]{\sphinxcrossref{\DUrole{std,std-ref}{(Ravelo)}}}}.
These equations compress the system to a state with average
axial stress or pressure equal to the specified target value
and that satisfies the Rankine\sphinxhyphen{}Hugoniot (RH)
jump conditions for steady shocks.

\sphinxAtStartPar
The compression can be performed
either
hydrostatically (using keyword \sphinxstyleemphasis{iso}, \sphinxstyleemphasis{aniso}, or \sphinxstyleemphasis{tri}) or uniaxially
(using keywords \sphinxstyleemphasis{x}, \sphinxstyleemphasis{y}, or \sphinxstyleemphasis{z}).  In the hydrostatic case,
the cell dimensions change dynamically so that the average axial stress
in all three directions converges towards the specified target value.
In the uniaxial case, the chosen cell dimension changes dynamically
so that the average
axial stress in that direction converges towards the target value. The
other two cell dimensions are kept fixed (zero lateral strain).

\sphinxAtStartPar
This leads to the following additional restrictions on the keywords:
\begin{itemize}
\item {} 
\sphinxAtStartPar
One and only one of the following keywords should be used: \sphinxstyleemphasis{iso}, \sphinxstyleemphasis{aniso}, \sphinxstyleemphasis{tri}, \sphinxstyleemphasis{x}, \sphinxstyleemphasis{y}, \sphinxstyleemphasis{z}

\item {} 
\sphinxAtStartPar
The specified initial and final target pressures must be the same.

\item {} 
\sphinxAtStartPar
The keywords \sphinxstyleemphasis{xy}, \sphinxstyleemphasis{xz}, \sphinxstyleemphasis{yz} may not be used.

\item {} 
\sphinxAtStartPar
The only admissible value for the couple keyword is \sphinxstyleemphasis{xyz}, which has the same effect as keyword \sphinxstyleemphasis{iso}

\item {} 
\sphinxAtStartPar
The \sphinxstyleemphasis{temp} keyword must be used to specify the time constant for kinetic energy relaxation, but initial and final target temperature values are ignored.

\end{itemize}

\sphinxAtStartPar
Essentially, a Hugoniostat simulation is an NPT simulation in which the
user\sphinxhyphen{}specified target temperature is replaced with a time\sphinxhyphen{}dependent
target temperature Tt obtained from the following equation:
\begin{equation*}
\begin{split}T_t - T = \frac{\left(\frac{1}{2}\left(P + P_0\right)\left(V_0 - V\right) + E_0 - E\right)}{N_{dof} k_B } = \Delta\end{split}
\end{equation*}
\sphinxAtStartPar
where \(T\) and \(T_t\) are the instantaneous and target
temperatures, \sphinxstyleemphasis{P} and \(P_0\) are the instantaneous and reference
pressures or axial stresses, depending on whether hydrostatic or
uniaxial compression is being performed, \sphinxstyleemphasis{V} and \(V_0\) are the
instantaneous and reference volumes, \sphinxstyleemphasis{E} and \(E_0\) are the
instantaneous and reference internal energy (potential plus kinetic),
\(N_{dof}\) is the number of degrees of freedom used in the
definition of temperature, and \(k_B\) is the Boltzmann
constant. \(\Delta\) is the negative deviation of the instantaneous
temperature from the target temperature.  When the system reaches a
stable equilibrium, the value of \(\Delta\) should fluctuate about
zero.

\sphinxAtStartPar
The values of \(E_0\), \(V_0\), and \(P_0\) are the
instantaneous values at the start of the simulation. These can be
overridden using the fix\_modify keywords \sphinxstyleemphasis{e0}, \sphinxstyleemphasis{v0}, and \sphinxstyleemphasis{p0} described
below.


\bigskip\hrule\bigskip


\begin{sphinxadmonition}{note}{Note:}
\sphinxAtStartPar
Unlike the {\hyperref[\detokenize{fix_temp_berendsen::doc}]{\sphinxcrossref{\DUrole{doc}{fix temp/berendsen}}}} command
which performs thermostatting but NO time integration, this fix
performs thermostatting/barostatting AND time integration.  Thus you
should not use any other time integration fix, such as {\hyperref[\detokenize{fix_nve::doc}]{\sphinxcrossref{\DUrole{doc}{fix nve}}}} on atoms to which this fix is applied.  Likewise,
this fix should not be used on atoms that have their temperature
controlled by another fix \sphinxhyphen{} e.g. by {\hyperref[\detokenize{fix_nh::doc}]{\sphinxcrossref{\DUrole{doc}{fix langevin}}}} or {\hyperref[\detokenize{fix_temp_rescale::doc}]{\sphinxcrossref{\DUrole{doc}{fix temp/rescale}}}} commands.
\end{sphinxadmonition}


\bigskip\hrule\bigskip


\sphinxAtStartPar
This fix computes a temperature and pressure at each timestep.  To do
this, the fix creates its own computes of style “temp” and “pressure”,
as if one of these two sets of commands had been issued:

\begin{sphinxVerbatim}[commandchars=\\\{\}]
\PYG{k}{compute }\PYG{n+nv+nvIdentifier}{fix\PYGZhy{}ID\PYGZus{}temp}\PYG{+w}{ }\PYG{n+nv+nvIdentifier}{group\PYGZhy{}ID}\PYG{+w}{ }\PYG{n}{temp}
\PYG{k}{compute }\PYG{n+nv+nvIdentifier}{fix\PYGZhy{}ID\PYGZus{}press}\PYG{+w}{ }\PYG{n+nv+nvIdentifier}{group\PYGZhy{}ID}\PYG{+w}{ }\PYG{n}{pressure}\PYG{+w}{ }\PYG{n}{fix}\PYG{o}{\PYGZhy{}}\PYG{n}{ID\PYGZus{}temp}

\PYG{k}{compute }\PYG{n+nv+nvIdentifier}{fix\PYGZhy{}ID\PYGZus{}temp}\PYG{+w}{ }\PYG{n+nv+nvIdentifier}{all}\PYG{+w}{ }\PYG{n}{temp}
\PYG{k}{compute }\PYG{n+nv+nvIdentifier}{fix\PYGZhy{}ID\PYGZus{}press}\PYG{+w}{ }\PYG{n+nv+nvIdentifier}{all}\PYG{+w}{ }\PYG{n}{pressure}\PYG{+w}{ }\PYG{n}{fix}\PYG{o}{\PYGZhy{}}\PYG{n}{ID\PYGZus{}temp}
\end{sphinxVerbatim}

\sphinxAtStartPar
See the {\hyperref[\detokenize{compute_temp::doc}]{\sphinxcrossref{\DUrole{doc}{compute temp}}}} and {\hyperref[\detokenize{compute_pressure::doc}]{\sphinxcrossref{\DUrole{doc}{compute pressure}}}} commands for details.  Note that the
IDs of the new computes are the fix\sphinxhyphen{}ID + underscore + “temp” or fix\_ID
+ underscore + “press”.  The group for
the new computes is “all” since pressure is computed for the entire
system.

\sphinxAtStartPar
Note that these are NOT the computes used by thermodynamic output (see
the {\hyperref[\detokenize{thermo_style::doc}]{\sphinxcrossref{\DUrole{doc}{thermo\_style}}}} command) with ID = \sphinxstyleemphasis{thermo\_temp}
and \sphinxstyleemphasis{thermo\_press}.  This means you can change the attributes of this
fix’s temperature or pressure via the
{\hyperref[\detokenize{compute_modify::doc}]{\sphinxcrossref{\DUrole{doc}{compute\_modify}}}} command or print this temperature
or pressure during thermodynamic output via the {\hyperref[\detokenize{thermo_style::doc}]{\sphinxcrossref{\DUrole{doc}{thermo\_style custom}}}} command using the appropriate compute\sphinxhyphen{}ID.
It also means that changing attributes of \sphinxstyleemphasis{thermo\_temp} or
\sphinxstyleemphasis{thermo\_press} will have no effect on this fix.


\bigskip\hrule\bigskip


\sphinxAtStartPar
Styles with a \sphinxstyleemphasis{gpu}, \sphinxstyleemphasis{intel}, \sphinxstyleemphasis{kk}, \sphinxstyleemphasis{omp}, or \sphinxstyleemphasis{opt} suffix are
functionally the same as the corresponding style without the suffix.
They have been optimized to run faster, depending on your available
hardware, as discussed on the {\hyperref[\detokenize{Speed_packages::doc}]{\sphinxcrossref{\DUrole{doc}{Accelerator packages}}}}
page.  The accelerated styles take the same arguments and should
produce the same results, except for round\sphinxhyphen{}off and precision issues.

\sphinxAtStartPar
These accelerated styles are part of the GPU, INTEL, KOKKOS,
OPENMP, and OPT packages, respectively.  They are only enabled if
LAMMPS was built with those packages.  See the {\hyperref[\detokenize{Build_package::doc}]{\sphinxcrossref{\DUrole{doc}{Build package}}}} page for more info.

\sphinxAtStartPar
You can specify the accelerated styles explicitly in your input script
by including their suffix, or you can use the {\hyperref[\detokenize{Run_options::doc}]{\sphinxcrossref{\DUrole{doc}{\sphinxhyphen{}suffix command\sphinxhyphen{}line switch}}}} when you invoke LAMMPS, or you can use the
{\hyperref[\detokenize{suffix::doc}]{\sphinxcrossref{\DUrole{doc}{suffix}}}} command in your input script.

\sphinxAtStartPar
See the {\hyperref[\detokenize{Speed_packages::doc}]{\sphinxcrossref{\DUrole{doc}{Accelerator packages}}}} page for more
instructions on how to use the accelerated styles effectively.


\bigskip\hrule\bigskip



\subsection{Restart, fix\_modify, output, run start/stop, minimize info}
\label{\detokenize{fix_nphug:restart-fix-modify-output-run-start-stop-minimize-info}}
\sphinxAtStartPar
This fix writes the values of \(E_0\), \(V_0\), and \(P_0\),
as well as the state of all the thermostat and barostat variables to
{\hyperref[\detokenize{restart::doc}]{\sphinxcrossref{\DUrole{doc}{binary restart files}}}}.  See the {\hyperref[\detokenize{read_restart::doc}]{\sphinxcrossref{\DUrole{doc}{read\_restart}}}} command for info on how to re\sphinxhyphen{}specify a fix in an input
script that reads a restart file, so that the operation of the fix
continues in an uninterrupted fashion.

\sphinxAtStartPar
The {\hyperref[\detokenize{fix_modify::doc}]{\sphinxcrossref{\DUrole{doc}{fix\_modify}}}} \sphinxstyleemphasis{e0}, \sphinxstyleemphasis{v0} and \sphinxstyleemphasis{p0} keywords can be
used to define the values of \(E_0\), \(V_0\), and
\(P_0\). Note the the values for \sphinxstyleemphasis{e0} and \sphinxstyleemphasis{v0} are extensive, and so
must correspond to the total energy and volume of the entire system, not
energy and volume per atom. If any of these quantities are not
specified, then the instantaneous value in the system at the start of
the simulation is used.

\sphinxAtStartPar
The {\hyperref[\detokenize{fix_modify::doc}]{\sphinxcrossref{\DUrole{doc}{fix\_modify}}}} \sphinxstyleemphasis{temp} and \sphinxstyleemphasis{press} options are
supported by this fix.  You can use them to assign a {\hyperref[\detokenize{compute::doc}]{\sphinxcrossref{\DUrole{doc}{compute}}}} you have defined to this fix which will be used in its
thermostatting or barostatting procedure, as described above.  If you
do this, note that the kinetic energy derived from the compute
temperature should be consistent with the virial term computed using
all atoms for the pressure.  LAMMPS will warn you if you choose to
compute temperature on a subset of atoms.

\sphinxAtStartPar
The cumulative energy change in the system imposed by this fix is
included in the {\hyperref[\detokenize{thermo_style::doc}]{\sphinxcrossref{\DUrole{doc}{thermodynamic output}}}} keywords
\sphinxstyleemphasis{ecouple} and \sphinxstyleemphasis{econserve}.  See the {\hyperref[\detokenize{thermo_style::doc}]{\sphinxcrossref{\DUrole{doc}{thermo\_style}}}}
doc page for details.  Note that this energy is *not* included in
the definition of internal energy E when calculating the value of
Delta in the above equation.

\sphinxAtStartPar
This fix computes a global scalar which can be accessed by various
{\hyperref[\detokenize{Howto_output::doc}]{\sphinxcrossref{\DUrole{doc}{output commands}}}}.  The scalar is the same
cumulative energy change due to this fix described in the previous
paragraph.  The scalar value calculated by this fix is “extensive”.

\sphinxAtStartPar
This fix also computes a global vector of quantities, which can be
accessed by various {\hyperref[\detokenize{Howto_output::doc}]{\sphinxcrossref{\DUrole{doc}{output commands}}}}.  The scalar
The vector values are “intensive”.

\sphinxAtStartPar
The vector stores three quantities unique to this fix (\(\Delta\),
Us, and up), followed by all the internal Nose/Hoover thermostat and
barostat variables defined for {\hyperref[\detokenize{fix_nh::doc}]{\sphinxcrossref{\DUrole{doc}{fix npt}}}}. Delta is the
deviation of the temperature from the target temperature, given by the
above equation.  Us and up are the shock and particle velocity
corresponding to a steady shock calculated from the RH
conditions. They have units of distance/time.


\subsection{Restrictions}
\label{\detokenize{fix_nphug:restrictions}}
\sphinxAtStartPar
This fix style is part of the SHOCK package.  It is only enabled if
LAMMPS was built with that package. See the {\hyperref[\detokenize{Build_package::doc}]{\sphinxcrossref{\DUrole{doc}{Build package}}}} page for more info.

\sphinxAtStartPar
All the usual restrictions for {\hyperref[\detokenize{fix_nh::doc}]{\sphinxcrossref{\DUrole{doc}{fix npt}}}} apply, plus the
additional ones mentioned above.


\subsection{Related commands}
\label{\detokenize{fix_nphug:related-commands}}
\sphinxAtStartPar
{\hyperref[\detokenize{fix_msst::doc}]{\sphinxcrossref{\DUrole{doc}{fix msst}}}}, {\hyperref[\detokenize{fix_nh::doc}]{\sphinxcrossref{\DUrole{doc}{fix npt}}}}, {\hyperref[\detokenize{fix_modify::doc}]{\sphinxcrossref{\DUrole{doc}{fix\_modify}}}}


\subsection{Default}
\label{\detokenize{fix_nphug:default}}
\sphinxAtStartPar
The keyword defaults are the same as those for {\hyperref[\detokenize{fix_nh::doc}]{\sphinxcrossref{\DUrole{doc}{fix npt}}}}


\bigskip\hrule\bigskip

\phantomsection\label{\detokenize{fix_nphug:ravelo1}}
\sphinxAtStartPar
\sphinxstylestrong{(Ravelo)} Ravelo, Holian, Germann and Lomdahl, Phys Rev B, 70, 014103 (2004).

\sphinxstepscope

\index{fix npt/asphere@\spxentry{fix npt/asphere}}\ignorespaces 
\index{fix npt/asphere/omp@\spxentry{fix npt/asphere/omp}}\ignorespaces 

\section{fix npt/asphere command}
\label{\detokenize{fix_npt_asphere:fix-npt-asphere-command}}\label{\detokenize{fix_npt_asphere:index-1}}\label{\detokenize{fix_npt_asphere:index-0}}\label{\detokenize{fix_npt_asphere::doc}}
\sphinxAtStartPar
Accelerator Variants: \sphinxstyleemphasis{npt/asphere/omp}


\subsection{Syntax}
\label{\detokenize{fix_npt_asphere:syntax}}
\begin{sphinxVerbatim}[commandchars=\\\{\}]
\PYG{k}{fix }\PYG{n+nv+nvIdentifier}{ID}\PYG{+w}{ }\PYG{n+nv+nvIdentifier}{group\PYGZhy{}ID}\PYG{+w}{ }\PYG{n}{npt}\PYG{o}{/}\PYG{n}{asphere}\PYG{+w}{ }\PYG{n}{keyword}\PYG{+w}{ }\PYG{n}{value}\PYG{+w}{ }\PYG{n}{...}
\end{sphinxVerbatim}
\begin{itemize}
\item {} 
\sphinxAtStartPar
ID, group\sphinxhyphen{}ID are documented in {\hyperref[\detokenize{fix::doc}]{\sphinxcrossref{\DUrole{doc}{fix}}}} command

\item {} 
\sphinxAtStartPar
npt/asphere = style name of this fix command

\item {} 
\sphinxAtStartPar
additional thermostat and barostat related keyword/value pairs from the {\hyperref[\detokenize{fix_nh::doc}]{\sphinxcrossref{\DUrole{doc}{fix npt}}}} command can be appended

\end{itemize}


\subsection{Examples}
\label{\detokenize{fix_npt_asphere:examples}}
\begin{sphinxVerbatim}[commandchars=\\\{\}]
\PYG{k}{fix }\PYG{n+nv+nvIdentifier}{1}\PYG{+w}{ }\PYG{n+nv+nvIdentifier}{all}\PYG{+w}{ }\PYG{n}{npt}\PYG{o}{/}\PYG{n}{asphere}\PYG{+w}{ }\PYG{n}{temp}\PYG{+w}{ }\PYG{l+m}{300.0}\PYG{+w}{ }\PYG{l+m}{300.0}\PYG{+w}{ }\PYG{l+m}{100.0}\PYG{+w}{ }\PYG{n}{iso}\PYG{+w}{ }\PYG{l+m}{0.0}\PYG{+w}{ }\PYG{l+m}{0.0}\PYG{+w}{ }\PYG{l+m}{1000.0}
\PYG{k}{fix }\PYG{n+nv+nvIdentifier}{2}\PYG{+w}{ }\PYG{n+nv+nvIdentifier}{all}\PYG{+w}{ }\PYG{n}{npt}\PYG{o}{/}\PYG{n}{asphere}\PYG{+w}{ }\PYG{n}{temp}\PYG{+w}{ }\PYG{l+m}{300.0}\PYG{+w}{ }\PYG{l+m}{300.0}\PYG{+w}{ }\PYG{l+m}{100.0}\PYG{+w}{ }\PYG{n}{x}\PYG{+w}{ }\PYG{l+m}{5.0}\PYG{+w}{ }\PYG{l+m}{5.0}\PYG{+w}{ }\PYG{l+m}{1000.0}
\PYG{k}{fix }\PYG{n+nv+nvIdentifier}{2}\PYG{+w}{ }\PYG{n+nv+nvIdentifier}{all}\PYG{+w}{ }\PYG{n}{npt}\PYG{o}{/}\PYG{n}{asphere}\PYG{+w}{ }\PYG{n}{temp}\PYG{+w}{ }\PYG{l+m}{300.0}\PYG{+w}{ }\PYG{l+m}{300.0}\PYG{+w}{ }\PYG{l+m}{100.0}\PYG{+w}{ }\PYG{n}{x}\PYG{+w}{ }\PYG{l+m}{5.0}\PYG{+w}{ }\PYG{l+m}{5.0}\PYG{+w}{ }\PYG{l+m}{1000.0}\PYG{+w}{ }\PYG{n}{drag}\PYG{+w}{ }\PYG{l+m}{0.2}
\PYG{k}{fix }\PYG{n+nv+nvIdentifier}{2}\PYG{+w}{ }\PYG{n+nv+nvIdentifier}{water}\PYG{+w}{ }\PYG{n}{npt}\PYG{o}{/}\PYG{n}{asphere}\PYG{+w}{ }\PYG{n}{temp}\PYG{+w}{ }\PYG{l+m}{300.0}\PYG{+w}{ }\PYG{l+m}{300.0}\PYG{+w}{ }\PYG{l+m}{100.0}\PYG{+w}{ }\PYG{n}{aniso}\PYG{+w}{ }\PYG{l+m}{0.0}\PYG{+w}{ }\PYG{l+m}{0.0}\PYG{+w}{ }\PYG{l+m}{1000.0}\PYG{+w}{ }\PYG{n}{dilate}\PYG{+w}{ }\PYG{n}{partial}
\end{sphinxVerbatim}


\subsection{Description}
\label{\detokenize{fix_npt_asphere:description}}
\sphinxAtStartPar
Perform constant NPT integration to update position, velocity,
orientation, and angular velocity each timestep for aspherical or
ellipsoidal particles in the group using a Nose/Hoover temperature
thermostat and Nose/Hoover pressure barostat.  P is pressure; T is
temperature.  This creates a system trajectory consistent with the
isothermal\sphinxhyphen{}isobaric ensemble.

\sphinxAtStartPar
This fix differs from the {\hyperref[\detokenize{fix_nh::doc}]{\sphinxcrossref{\DUrole{doc}{fix npt}}}} command, which
assumes point particles and only updates their position and velocity.

\sphinxAtStartPar
The thermostat is applied to both the translational and rotational
degrees of freedom for the aspherical particles, assuming a compute is
used which calculates a temperature that includes the rotational
degrees of freedom (see below).  The translational degrees of freedom
can also have a bias velocity removed from them before thermostatting
takes place; see the description below.

\sphinxAtStartPar
Additional parameters affecting the thermostat and barostat are
specified by keywords and values documented with the {\hyperref[\detokenize{fix_nh::doc}]{\sphinxcrossref{\DUrole{doc}{fix npt}}}} command.  See, for example, discussion of the \sphinxstyleemphasis{temp}, \sphinxstyleemphasis{iso},
\sphinxstyleemphasis{aniso}, and \sphinxstyleemphasis{dilate} keywords.

\sphinxAtStartPar
The particles in the fix group are the only ones whose velocities and
positions are updated by the velocity/position update portion of the
NPT integration.

\sphinxAtStartPar
Regardless of what particles are in the fix group, a global pressure is
computed for all particles.  Similarly, when the size of the simulation
box is changed, all particles are re\sphinxhyphen{}scaled to new positions, unless the
keyword \sphinxstyleemphasis{dilate} is specified with a value of \sphinxstyleemphasis{partial}, in which case
only the particles in the fix group are re\sphinxhyphen{}scaled.  The latter can be
useful for leaving the coordinates of particles in a solid substrate
unchanged and controlling the pressure of a surrounding fluid.


\bigskip\hrule\bigskip


\sphinxAtStartPar
This fix computes a temperature and pressure each timestep.  To do
this, the fix creates its own computes of style “temp/asphere” and
“pressure”, as if these commands had been issued:

\begin{sphinxVerbatim}[commandchars=\\\{\}]
\PYG{k}{compute }\PYG{n+nv+nvIdentifier}{fix\PYGZhy{}ID\PYGZus{}temp}\PYG{+w}{ }\PYG{n+nv+nvIdentifier}{all}\PYG{+w}{ }\PYG{n}{temp}\PYG{o}{/}\PYG{n}{asphere}
\PYG{k}{compute }\PYG{n+nv+nvIdentifier}{fix\PYGZhy{}ID\PYGZus{}press}\PYG{+w}{ }\PYG{n+nv+nvIdentifier}{all}\PYG{+w}{ }\PYG{n}{pressure}\PYG{+w}{ }\PYG{n}{fix}\PYG{o}{\PYGZhy{}}\PYG{n}{ID\PYGZus{}temp}
\end{sphinxVerbatim}

\sphinxAtStartPar
See the {\hyperref[\detokenize{compute_temp_asphere::doc}]{\sphinxcrossref{\DUrole{doc}{compute temp/asphere}}}} and {\hyperref[\detokenize{compute_pressure::doc}]{\sphinxcrossref{\DUrole{doc}{compute pressure}}}} commands for details.  Note that the
IDs of the new computes are the fix\sphinxhyphen{}ID + underscore + “temp” or fix\_ID
+ underscore + “press”, and the group for the new computes is “all”
since pressure is computed for the entire system.

\sphinxAtStartPar
Note that these are NOT the computes used by thermodynamic output (see
the {\hyperref[\detokenize{thermo_style::doc}]{\sphinxcrossref{\DUrole{doc}{thermo\_style}}}} command) with ID = \sphinxstyleemphasis{thermo\_temp}
and \sphinxstyleemphasis{thermo\_press}.  This means you can change the attributes of this
fix’s temperature or pressure via the
{\hyperref[\detokenize{compute_modify::doc}]{\sphinxcrossref{\DUrole{doc}{compute\_modify}}}} command or print this temperature
or pressure during thermodynamic output via the {\hyperref[\detokenize{thermo_style::doc}]{\sphinxcrossref{\DUrole{doc}{thermo\_style custom}}}} command using the appropriate compute\sphinxhyphen{}ID.
It also means that changing attributes of \sphinxstyleemphasis{thermo\_temp} or
\sphinxstyleemphasis{thermo\_press} will have no effect on this fix.

\sphinxAtStartPar
Like other fixes that perform thermostatting, this fix can be used
with {\hyperref[\detokenize{compute::doc}]{\sphinxcrossref{\DUrole{doc}{compute commands}}}} that remove a “bias” from the
atom velocities.  E.g. to apply the thermostat only to atoms within a
spatial {\hyperref[\detokenize{region::doc}]{\sphinxcrossref{\DUrole{doc}{region}}}}, or to remove the center\sphinxhyphen{}of\sphinxhyphen{}mass
velocity from a group of atoms, or to remove the x\sphinxhyphen{}component of
velocity from the calculation.

\sphinxAtStartPar
This is not done by default, but only if the {\hyperref[\detokenize{fix_modify::doc}]{\sphinxcrossref{\DUrole{doc}{fix\_modify}}}} command is used to assign a temperature compute to this
fix that includes such a bias term.  See the doc pages for individual
{\hyperref[\detokenize{compute::doc}]{\sphinxcrossref{\DUrole{doc}{compute temp commands}}}} to determine which ones include
a bias.  In this case, the thermostat works in the following manner:
bias is removed from each atom, thermostatting is performed on the
remaining thermal degrees of freedom, and the bias is added back in.


\bigskip\hrule\bigskip


\sphinxAtStartPar
Styles with a \sphinxstyleemphasis{gpu}, \sphinxstyleemphasis{intel}, \sphinxstyleemphasis{kk}, \sphinxstyleemphasis{omp}, or \sphinxstyleemphasis{opt} suffix are
functionally the same as the corresponding style without the suffix.
They have been optimized to run faster, depending on your available
hardware, as discussed on the {\hyperref[\detokenize{Speed_packages::doc}]{\sphinxcrossref{\DUrole{doc}{Accelerator packages}}}}
page.  The accelerated styles take the same arguments and should
produce the same results, except for round\sphinxhyphen{}off and precision issues.

\sphinxAtStartPar
These accelerated styles are part of the GPU, INTEL, KOKKOS,
OPENMP, and OPT packages, respectively.  They are only enabled if
LAMMPS was built with those packages.  See the {\hyperref[\detokenize{Build_package::doc}]{\sphinxcrossref{\DUrole{doc}{Build package}}}} page for more info.

\sphinxAtStartPar
You can specify the accelerated styles explicitly in your input script
by including their suffix, or you can use the {\hyperref[\detokenize{Run_options::doc}]{\sphinxcrossref{\DUrole{doc}{\sphinxhyphen{}suffix command\sphinxhyphen{}line switch}}}} when you invoke LAMMPS, or you can use the
{\hyperref[\detokenize{suffix::doc}]{\sphinxcrossref{\DUrole{doc}{suffix}}}} command in your input script.

\sphinxAtStartPar
See the {\hyperref[\detokenize{Speed_packages::doc}]{\sphinxcrossref{\DUrole{doc}{Accelerator packages}}}} page for more
instructions on how to use the accelerated styles effectively.


\subsection{Restart, fix\_modify, output, run start/stop, minimize info}
\label{\detokenize{fix_npt_asphere:restart-fix-modify-output-run-start-stop-minimize-info}}
\sphinxAtStartPar
This fix writes the state of the Nose/Hoover thermostat and barostat
to {\hyperref[\detokenize{restart::doc}]{\sphinxcrossref{\DUrole{doc}{binary restart files}}}}.  See the
{\hyperref[\detokenize{read_restart::doc}]{\sphinxcrossref{\DUrole{doc}{read\_restart}}}} command for info on how to re\sphinxhyphen{}specify
a fix in an input script that reads a restart file, so that the
operation of the fix continues in an uninterrupted fashion.

\sphinxAtStartPar
The {\hyperref[\detokenize{fix_modify::doc}]{\sphinxcrossref{\DUrole{doc}{fix\_modify}}}} \sphinxstyleemphasis{temp} and \sphinxstyleemphasis{press} options are
supported by this fix.  You can use them to assign a
{\hyperref[\detokenize{compute::doc}]{\sphinxcrossref{\DUrole{doc}{compute}}}} you have defined to this fix which will be used
in its thermostatting or barostatting procedure.  If you do this, note
that the kinetic energy derived from the compute temperature should be
consistent with the virial term computed using all atoms for the
pressure.  LAMMPS will warn you if you choose to compute temperature
on a subset of atoms.

\sphinxAtStartPar
The cumulative energy change in the system imposed by this fix is
included in the {\hyperref[\detokenize{thermo_style::doc}]{\sphinxcrossref{\DUrole{doc}{thermodynamic output}}}} keywords
\sphinxstyleemphasis{ecouple} and \sphinxstyleemphasis{econserve}.  See the {\hyperref[\detokenize{thermo_style::doc}]{\sphinxcrossref{\DUrole{doc}{thermo\_style}}}}
doc page for details.

\sphinxAtStartPar
This fix computes the same global scalar and global vector of
quantities as does the {\hyperref[\detokenize{fix_nh::doc}]{\sphinxcrossref{\DUrole{doc}{fix npt}}}} command.

\sphinxAtStartPar
This fix can ramp its target temperature and pressure over multiple
runs, using the \sphinxstyleemphasis{start} and \sphinxstyleemphasis{stop} keywords of the {\hyperref[\detokenize{run::doc}]{\sphinxcrossref{\DUrole{doc}{run}}}}
command.  See the {\hyperref[\detokenize{run::doc}]{\sphinxcrossref{\DUrole{doc}{run}}}} command for details of how to do
this.

\sphinxAtStartPar
This fix is not invoked during {\hyperref[\detokenize{minimize::doc}]{\sphinxcrossref{\DUrole{doc}{energy minimization}}}}.


\subsection{Restrictions}
\label{\detokenize{fix_npt_asphere:restrictions}}
\sphinxAtStartPar
This fix is part of the ASPHERE package.  It is only enabled if LAMMPS
was built with that package.  See the {\hyperref[\detokenize{Build_package::doc}]{\sphinxcrossref{\DUrole{doc}{Build package}}}} page for more info.

\sphinxAtStartPar
This fix requires that atoms store torque and angular momentum and a
quaternion as defined by the {\hyperref[\detokenize{atom_style::doc}]{\sphinxcrossref{\DUrole{doc}{atom\_style ellipsoid}}}}
command.

\sphinxAtStartPar
All particles in the group must be finite\sphinxhyphen{}size.  They cannot be point
particles, but they can be aspherical or spherical as defined by their
shape attribute.


\subsection{Related commands}
\label{\detokenize{fix_npt_asphere:related-commands}}
\sphinxAtStartPar
{\hyperref[\detokenize{fix_nh::doc}]{\sphinxcrossref{\DUrole{doc}{fix npt}}}}, {\hyperref[\detokenize{fix_nve_asphere::doc}]{\sphinxcrossref{\DUrole{doc}{fix nve\_asphere}}}}, {\hyperref[\detokenize{fix_nvt_asphere::doc}]{\sphinxcrossref{\DUrole{doc}{fix nvt\_asphere}}}}, {\hyperref[\detokenize{fix_modify::doc}]{\sphinxcrossref{\DUrole{doc}{fix\_modify}}}}


\subsection{Default}
\label{\detokenize{fix_npt_asphere:default}}
\sphinxAtStartPar
none

\sphinxstepscope

\index{fix npt/body@\spxentry{fix npt/body}}\ignorespaces 

\section{fix npt/body command}
\label{\detokenize{fix_npt_body:fix-npt-body-command}}\label{\detokenize{fix_npt_body:index-0}}\label{\detokenize{fix_npt_body::doc}}

\subsection{Syntax}
\label{\detokenize{fix_npt_body:syntax}}
\begin{sphinxVerbatim}[commandchars=\\\{\}]
\PYG{k}{fix }\PYG{n+nv+nvIdentifier}{ID}\PYG{+w}{ }\PYG{n+nv+nvIdentifier}{group\PYGZhy{}ID}\PYG{+w}{ }\PYG{n}{npt}\PYG{o}{/}\PYG{n}{body}\PYG{+w}{ }\PYG{n}{keyword}\PYG{+w}{ }\PYG{n}{value}\PYG{+w}{ }\PYG{n}{...}
\end{sphinxVerbatim}
\begin{itemize}
\item {} 
\sphinxAtStartPar
ID, group\sphinxhyphen{}ID are documented in {\hyperref[\detokenize{fix::doc}]{\sphinxcrossref{\DUrole{doc}{fix}}}} command

\item {} 
\sphinxAtStartPar
npt/body = style name of this fix command

\item {} 
\sphinxAtStartPar
additional thermostat and barostat related keyword/value pairs from the {\hyperref[\detokenize{fix_nh::doc}]{\sphinxcrossref{\DUrole{doc}{fix npt}}}} command can be appended

\end{itemize}


\subsection{Examples}
\label{\detokenize{fix_npt_body:examples}}
\begin{sphinxVerbatim}[commandchars=\\\{\}]
\PYG{k}{fix }\PYG{n+nv+nvIdentifier}{1}\PYG{+w}{ }\PYG{n+nv+nvIdentifier}{all}\PYG{+w}{ }\PYG{n}{npt}\PYG{o}{/}\PYG{n}{body}\PYG{+w}{ }\PYG{n}{temp}\PYG{+w}{ }\PYG{l+m}{300.0}\PYG{+w}{ }\PYG{l+m}{300.0}\PYG{+w}{ }\PYG{l+m}{100.0}\PYG{+w}{ }\PYG{n}{iso}\PYG{+w}{ }\PYG{l+m}{0.0}\PYG{+w}{ }\PYG{l+m}{0.0}\PYG{+w}{ }\PYG{l+m}{1000.0}
\PYG{k}{fix }\PYG{n+nv+nvIdentifier}{2}\PYG{+w}{ }\PYG{n+nv+nvIdentifier}{all}\PYG{+w}{ }\PYG{n}{npt}\PYG{o}{/}\PYG{n}{body}\PYG{+w}{ }\PYG{n}{temp}\PYG{+w}{ }\PYG{l+m}{300.0}\PYG{+w}{ }\PYG{l+m}{300.0}\PYG{+w}{ }\PYG{l+m}{100.0}\PYG{+w}{ }\PYG{n}{x}\PYG{+w}{ }\PYG{l+m}{5.0}\PYG{+w}{ }\PYG{l+m}{5.0}\PYG{+w}{ }\PYG{l+m}{1000.0}
\PYG{k}{fix }\PYG{n+nv+nvIdentifier}{2}\PYG{+w}{ }\PYG{n+nv+nvIdentifier}{all}\PYG{+w}{ }\PYG{n}{npt}\PYG{o}{/}\PYG{n}{body}\PYG{+w}{ }\PYG{n}{temp}\PYG{+w}{ }\PYG{l+m}{300.0}\PYG{+w}{ }\PYG{l+m}{300.0}\PYG{+w}{ }\PYG{l+m}{100.0}\PYG{+w}{ }\PYG{n}{x}\PYG{+w}{ }\PYG{l+m}{5.0}\PYG{+w}{ }\PYG{l+m}{5.0}\PYG{+w}{ }\PYG{l+m}{1000.0}\PYG{+w}{ }\PYG{n}{drag}\PYG{+w}{ }\PYG{l+m}{0.2}
\PYG{k}{fix }\PYG{n+nv+nvIdentifier}{2}\PYG{+w}{ }\PYG{n+nv+nvIdentifier}{water}\PYG{+w}{ }\PYG{n}{npt}\PYG{o}{/}\PYG{n}{body}\PYG{+w}{ }\PYG{n}{temp}\PYG{+w}{ }\PYG{l+m}{300.0}\PYG{+w}{ }\PYG{l+m}{300.0}\PYG{+w}{ }\PYG{l+m}{100.0}\PYG{+w}{ }\PYG{n}{aniso}\PYG{+w}{ }\PYG{l+m}{0.0}\PYG{+w}{ }\PYG{l+m}{0.0}\PYG{+w}{ }\PYG{l+m}{1000.0}\PYG{+w}{ }\PYG{n}{dilate}\PYG{+w}{ }\PYG{n}{partial}
\end{sphinxVerbatim}


\subsection{Description}
\label{\detokenize{fix_npt_body:description}}
\sphinxAtStartPar
Perform constant NPT integration to update position, velocity,
orientation, and angular velocity each timestep for body
particles in the group using a Nose/Hoover temperature
thermostat and Nose/Hoover pressure barostat.  P is pressure; T is
temperature.  This creates a system trajectory consistent with the
isothermal\sphinxhyphen{}isobaric ensemble.

\sphinxAtStartPar
This fix differs from the {\hyperref[\detokenize{fix_nh::doc}]{\sphinxcrossref{\DUrole{doc}{fix npt}}}} command, which
assumes point particles and only updates their position and velocity.

\sphinxAtStartPar
The thermostat is applied to both the translational and rotational
degrees of freedom for the body particles, assuming a compute is
used which calculates a temperature that includes the rotational
degrees of freedom (see below).  The translational degrees of freedom
can also have a bias velocity removed from them before thermostatting
takes place; see the description below.

\sphinxAtStartPar
Additional parameters affecting the thermostat and barostat are
specified by keywords and values documented with the {\hyperref[\detokenize{fix_nh::doc}]{\sphinxcrossref{\DUrole{doc}{fix npt}}}} command.  See, for example, discussion of the \sphinxstyleemphasis{temp},
\sphinxstyleemphasis{iso}, \sphinxstyleemphasis{aniso}, and \sphinxstyleemphasis{dilate} keywords.

\sphinxAtStartPar
The particles in the fix group are the only ones whose velocities and
positions are updated by the velocity/position update portion of the
NPT integration.

\sphinxAtStartPar
Regardless of what particles are in the fix group, a global pressure is
computed for all particles.  Similarly, when the size of the simulation
box is changed, all particles are re\sphinxhyphen{}scaled to new positions, unless the
keyword \sphinxstyleemphasis{dilate} is specified with a value of \sphinxstyleemphasis{partial}, in which case
only the particles in the fix group are re\sphinxhyphen{}scaled.  The latter can be
useful for leaving the coordinates of particles in a solid substrate
unchanged and controlling the pressure of a surrounding fluid.


\bigskip\hrule\bigskip


\sphinxAtStartPar
This fix computes a temperature and pressure each timestep.  To do
this, the fix creates its own computes of style “temp/body” and
“pressure”, as if these commands had been issued:

\begin{sphinxVerbatim}[commandchars=\\\{\}]
\PYG{k}{compute }\PYG{n+nv+nvIdentifier}{fix\PYGZhy{}ID\PYGZus{}temp}\PYG{+w}{ }\PYG{n+nv+nvIdentifier}{all}\PYG{+w}{ }\PYG{n}{temp}\PYG{o}{/}\PYG{n}{body}
\PYG{k}{compute }\PYG{n+nv+nvIdentifier}{fix\PYGZhy{}ID\PYGZus{}press}\PYG{+w}{ }\PYG{n+nv+nvIdentifier}{all}\PYG{+w}{ }\PYG{n}{pressure}\PYG{+w}{ }\PYG{n}{fix}\PYG{o}{\PYGZhy{}}\PYG{n}{ID\PYGZus{}temp}
\end{sphinxVerbatim}

\sphinxAtStartPar
See the {\hyperref[\detokenize{compute_temp_body::doc}]{\sphinxcrossref{\DUrole{doc}{compute temp/body}}}} and {\hyperref[\detokenize{compute_pressure::doc}]{\sphinxcrossref{\DUrole{doc}{compute pressure}}}} commands for details.  Note that the
IDs of the new computes are the fix\sphinxhyphen{}ID + underscore + “temp” or fix\_ID
+ underscore + “press”, and the group for the new computes is “all”
since pressure is computed for the entire system.

\sphinxAtStartPar
Note that these are NOT the computes used by thermodynamic output (see
the {\hyperref[\detokenize{thermo_style::doc}]{\sphinxcrossref{\DUrole{doc}{thermo\_style}}}} command) with ID = \sphinxstyleemphasis{thermo\_temp}
and \sphinxstyleemphasis{thermo\_press}.  This means you can change the attributes of this
fix’s temperature or pressure via the
{\hyperref[\detokenize{compute_modify::doc}]{\sphinxcrossref{\DUrole{doc}{compute\_modify}}}} command or print this temperature
or pressure during thermodynamic output via the {\hyperref[\detokenize{thermo_style::doc}]{\sphinxcrossref{\DUrole{doc}{thermo\_style custom}}}} command using the appropriate compute\sphinxhyphen{}ID.
It also means that changing attributes of \sphinxstyleemphasis{thermo\_temp} or
\sphinxstyleemphasis{thermo\_press} will have no effect on this fix.

\sphinxAtStartPar
Like other fixes that perform thermostatting, this fix can be used
with {\hyperref[\detokenize{compute::doc}]{\sphinxcrossref{\DUrole{doc}{compute commands}}}} that remove a “bias” from the
atom velocities.  E.g. to apply the thermostat only to atoms within a
spatial {\hyperref[\detokenize{region::doc}]{\sphinxcrossref{\DUrole{doc}{region}}}}, or to remove the center\sphinxhyphen{}of\sphinxhyphen{}mass
velocity from a group of atoms, or to remove the x\sphinxhyphen{}component of
velocity from the calculation.

\sphinxAtStartPar
This is not done by default, but only if the {\hyperref[\detokenize{fix_modify::doc}]{\sphinxcrossref{\DUrole{doc}{fix\_modify}}}} command is used to assign a temperature compute to this
fix that includes such a bias term.  See the doc pages for individual
{\hyperref[\detokenize{compute::doc}]{\sphinxcrossref{\DUrole{doc}{compute temp commands}}}} to determine which ones include
a bias.  In this case, the thermostat works in the following manner:
bias is removed from each atom, thermostatting is performed on the
remaining thermal degrees of freedom, and the bias is added back in.


\bigskip\hrule\bigskip


\sphinxAtStartPar
Styles with a \sphinxstyleemphasis{gpu}, \sphinxstyleemphasis{intel}, \sphinxstyleemphasis{kk}, \sphinxstyleemphasis{omp}, or \sphinxstyleemphasis{opt} suffix are
functionally the same as the corresponding style without the suffix.
They have been optimized to run faster, depending on your available
hardware, as discussed on the {\hyperref[\detokenize{Speed_packages::doc}]{\sphinxcrossref{\DUrole{doc}{Accelerator packages}}}}
page.  The accelerated styles take the same arguments and should
produce the same results, except for round\sphinxhyphen{}off and precision issues.

\sphinxAtStartPar
These accelerated styles are part of the GPU, INTEL, KOKKOS,
OPENMP, and OPT packages, respectively.  They are only enabled if
LAMMPS was built with those packages.  See the {\hyperref[\detokenize{Build_package::doc}]{\sphinxcrossref{\DUrole{doc}{Build package}}}} page for more info.

\sphinxAtStartPar
You can specify the accelerated styles explicitly in your input script
by including their suffix, or you can use the {\hyperref[\detokenize{Run_options::doc}]{\sphinxcrossref{\DUrole{doc}{\sphinxhyphen{}suffix command\sphinxhyphen{}line switch}}}} when you invoke LAMMPS, or you can use the
{\hyperref[\detokenize{suffix::doc}]{\sphinxcrossref{\DUrole{doc}{suffix}}}} command in your input script.

\sphinxAtStartPar
See the {\hyperref[\detokenize{Speed_packages::doc}]{\sphinxcrossref{\DUrole{doc}{Accelerator packages}}}} page for more
instructions on how to use the accelerated styles effectively.


\subsection{Restart, fix\_modify, output, run start/stop, minimize info}
\label{\detokenize{fix_npt_body:restart-fix-modify-output-run-start-stop-minimize-info}}
\sphinxAtStartPar
This fix writes the state of the Nose/Hoover thermostat and barostat
to {\hyperref[\detokenize{restart::doc}]{\sphinxcrossref{\DUrole{doc}{binary restart files}}}}.  See the
{\hyperref[\detokenize{read_restart::doc}]{\sphinxcrossref{\DUrole{doc}{read\_restart}}}} command for info on how to re\sphinxhyphen{}specify
a fix in an input script that reads a restart file, so that the
operation of the fix continues in an uninterrupted fashion.

\sphinxAtStartPar
The {\hyperref[\detokenize{fix_modify::doc}]{\sphinxcrossref{\DUrole{doc}{fix\_modify}}}} \sphinxstyleemphasis{temp} and \sphinxstyleemphasis{press} options are
supported by this fix.  You can use them to assign a
{\hyperref[\detokenize{compute::doc}]{\sphinxcrossref{\DUrole{doc}{compute}}}} you have defined to this fix which will be used
in its thermostatting or barostatting procedure.  If you do this, note
that the kinetic energy derived from the compute temperature should be
consistent with the virial term computed using all atoms for the
pressure.  LAMMPS will warn you if you choose to compute temperature
on a subset of atoms.

\sphinxAtStartPar
The cumulative energy change in the system imposed by this fix is
included in the {\hyperref[\detokenize{thermo_style::doc}]{\sphinxcrossref{\DUrole{doc}{thermodynamic output}}}} keywords
\sphinxstyleemphasis{ecouple} and \sphinxstyleemphasis{econserve}.  See the {\hyperref[\detokenize{thermo_style::doc}]{\sphinxcrossref{\DUrole{doc}{thermo\_style}}}}
doc page for details.

\sphinxAtStartPar
This fix computes the same global scalar and global vector of
quantities as does the {\hyperref[\detokenize{fix_nh::doc}]{\sphinxcrossref{\DUrole{doc}{fix npt}}}} command.

\sphinxAtStartPar
This fix can ramp its target temperature and pressure over multiple
runs, using the \sphinxstyleemphasis{start} and \sphinxstyleemphasis{stop} keywords of the {\hyperref[\detokenize{run::doc}]{\sphinxcrossref{\DUrole{doc}{run}}}}
command.  See the {\hyperref[\detokenize{run::doc}]{\sphinxcrossref{\DUrole{doc}{run}}}} command for details of how to do
this.

\sphinxAtStartPar
This fix is not invoked during {\hyperref[\detokenize{minimize::doc}]{\sphinxcrossref{\DUrole{doc}{energy minimization}}}}.


\subsection{Restrictions}
\label{\detokenize{fix_npt_body:restrictions}}
\sphinxAtStartPar
This fix is part of the BODY package.  It is only enabled if LAMMPS
was built with that package.  See the {\hyperref[\detokenize{Build_package::doc}]{\sphinxcrossref{\DUrole{doc}{Build package}}}} page for more info.

\sphinxAtStartPar
This fix requires that atoms store torque and angular momentum and a
quaternion as defined by the {\hyperref[\detokenize{atom_style::doc}]{\sphinxcrossref{\DUrole{doc}{atom\_style body}}}}
command.


\subsection{Related commands}
\label{\detokenize{fix_npt_body:related-commands}}
\sphinxAtStartPar
{\hyperref[\detokenize{fix_nh::doc}]{\sphinxcrossref{\DUrole{doc}{fix npt}}}}, {\hyperref[\detokenize{fix_nve_body::doc}]{\sphinxcrossref{\DUrole{doc}{fix nve\_body}}}}, {\hyperref[\detokenize{fix_nvt_body::doc}]{\sphinxcrossref{\DUrole{doc}{fix nvt\_body}}}}, {\hyperref[\detokenize{fix_modify::doc}]{\sphinxcrossref{\DUrole{doc}{fix\_modify}}}}


\subsection{Default}
\label{\detokenize{fix_npt_body:default}}
\sphinxAtStartPar
none

\sphinxstepscope

\index{fix npt/cauchy@\spxentry{fix npt/cauchy}}\ignorespaces 

\section{fix npt/cauchy command}
\label{\detokenize{fix_npt_cauchy:fix-npt-cauchy-command}}\label{\detokenize{fix_npt_cauchy:index-0}}\label{\detokenize{fix_npt_cauchy::doc}}

\subsection{Syntax}
\label{\detokenize{fix_npt_cauchy:syntax}}
\begin{sphinxVerbatim}[commandchars=\\\{\}]
\PYG{k}{fix }\PYG{n+nv+nvIdentifier}{ID}\PYG{+w}{ }\PYG{n+nv+nvIdentifier}{group\PYGZhy{}ID}\PYG{+w}{ }\PYG{n}{style\PYGZus{}name}\PYG{+w}{ }\PYG{n}{keyword}\PYG{+w}{ }\PYG{n}{value}\PYG{+w}{ }\PYG{n}{...}
\end{sphinxVerbatim}
\begin{itemize}
\item {} 
\sphinxAtStartPar
ID, group\sphinxhyphen{}ID are documented in {\hyperref[\detokenize{fix::doc}]{\sphinxcrossref{\DUrole{doc}{fix}}}} command

\item {} 
\sphinxAtStartPar
style\_name = \sphinxstyleemphasis{npt/cauchy}

\item {} 
\sphinxAtStartPar
one or more keyword/value pairs may be appended

\item {} 
\sphinxAtStartPar
keyword = \sphinxstyleemphasis{temp} or \sphinxstyleemphasis{iso} or \sphinxstyleemphasis{aniso} or \sphinxstyleemphasis{tri} or \sphinxstyleemphasis{x} or \sphinxstyleemphasis{y} or \sphinxstyleemphasis{z} or \sphinxstyleemphasis{xy} or \sphinxstyleemphasis{yz} or \sphinxstyleemphasis{xz} or \sphinxstyleemphasis{couple} or \sphinxstyleemphasis{tchain} or \sphinxstyleemphasis{pchain} or \sphinxstyleemphasis{mtk} or \sphinxstyleemphasis{tloop} or \sphinxstyleemphasis{ploop} or \sphinxstyleemphasis{nreset} or \sphinxstyleemphasis{drag} or \sphinxstyleemphasis{dilate} or \sphinxstyleemphasis{scalexy} or \sphinxstyleemphasis{scaleyz} or \sphinxstyleemphasis{scalexz} or \sphinxstyleemphasis{flip} or \sphinxstyleemphasis{alpha} or \sphinxstyleemphasis{continue} or \sphinxstyleemphasis{fixedpoint}
\begin{sphinxalltt}
\sphinxstyleemphasis{temp} values = Tstart Tstop Tdamp
  Tstart,Tstop = external temperature at start/end of run
  Tdamp = temperature damping parameter (time units)
\sphinxstyleemphasis{iso} or \sphinxstyleemphasis{aniso} or \sphinxstyleemphasis{tri} values = Pstart Pstop Pdamp
  Pstart,Pstop = scalar external pressure at start/end of run (pressure units)
  Pdamp = pressure damping parameter (time units)
\sphinxstyleemphasis{x} or \sphinxstyleemphasis{y} or \sphinxstyleemphasis{z} or \sphinxstyleemphasis{xy} or \sphinxstyleemphasis{yz} or \sphinxstyleemphasis{xz} values = Pstart Pstop Pdamp
  Pstart,Pstop = external stress tensor component at start/end of run (pressure units)
  Pdamp = stress damping parameter (time units)
\sphinxstyleemphasis{couple} = \sphinxstyleemphasis{none} or \sphinxstyleemphasis{xyz} or \sphinxstyleemphasis{xy} or \sphinxstyleemphasis{yz} or \sphinxstyleemphasis{xz}
\sphinxstyleemphasis{tchain} value = N
  N = length of thermostat chain (1 = single thermostat)
\sphinxstyleemphasis{pchain} values = N
  N length of thermostat chain on barostat (0 = no thermostat)
\sphinxstyleemphasis{mtk} value = \sphinxstyleemphasis{yes} or \sphinxstyleemphasis{no} = add in MTK adjustment term or not
\sphinxstyleemphasis{tloop} value = M
  M = number of sub\sphinxhyphen{}cycles to perform on thermostat
\sphinxstyleemphasis{ploop} value = M
  M = number of sub\sphinxhyphen{}cycles to perform on barostat thermostat
\sphinxstyleemphasis{nreset} value = reset reference cell every this many timesteps
\sphinxstyleemphasis{drag} value = Df
  Df = drag factor added to barostat/thermostat (0.0 = no drag)
\sphinxstyleemphasis{dilate} value = dilate\sphinxhyphen{}group\sphinxhyphen{}ID
  dilate\sphinxhyphen{}group\sphinxhyphen{}ID = only dilate atoms in this group due to barostat volume changes
\sphinxstyleemphasis{scalexy} value = \sphinxstyleemphasis{yes} or \sphinxstyleemphasis{no} = scale xy with ly
\sphinxstyleemphasis{scaleyz} value = \sphinxstyleemphasis{yes} or \sphinxstyleemphasis{no} = scale yz with lz
\sphinxstyleemphasis{scalexz} value = \sphinxstyleemphasis{yes} or \sphinxstyleemphasis{no} = scale xz with lz
\sphinxstyleemphasis{flip} value = \sphinxstyleemphasis{yes} or \sphinxstyleemphasis{no} = allow or disallow box flips when it becomes highly skewed
\sphinxstyleemphasis{alpha} value = strength of Cauchy stress control parameter
\sphinxstyleemphasis{continue} value = \sphinxstyleemphasis{yes} or \sphinxstyleemphasis{no} = whether of not to continue from a previous run
\sphinxstyleemphasis{fixedpoint} values = x y z
  x,y,z = perform barostat dilation/contraction around this point (distance units)
\end{sphinxalltt}

\end{itemize}


\subsection{Examples}
\label{\detokenize{fix_npt_cauchy:examples}}
\begin{sphinxVerbatim}[commandchars=\\\{\}]
\PYG{k}{fix }\PYG{n+nv+nvIdentifier}{1}\PYG{+w}{ }\PYG{n+nv+nvIdentifier}{water}\PYG{+w}{ }\PYG{n}{npt}\PYG{o}{/}\PYG{n}{cauchy}\PYG{+w}{ }\PYG{n}{temp}\PYG{+w}{ }\PYG{l+m}{300.0}\PYG{+w}{ }\PYG{l+m}{300.0}\PYG{+w}{ }\PYG{l+m}{100.0}\PYG{+w}{ }\PYG{n}{iso}\PYG{+w}{ }\PYG{l+m}{0.0}\PYG{+w}{ }\PYG{l+m}{0.0}\PYG{+w}{ }\PYG{l+m}{1000.0}\PYG{+w}{ }\PYG{n}{alpha}\PYG{+w}{ }\PYG{l+m}{0.001}
\end{sphinxVerbatim}


\subsection{Description}
\label{\detokenize{fix_npt_cauchy:description}}
\sphinxAtStartPar
This command performs time integration on Nose\sphinxhyphen{}Hoover style
non\sphinxhyphen{}Hamiltonian equations of motion which are designed to generate
positions and velocities sampled from the isothermal\sphinxhyphen{}isobaric (npt)
ensembles.  This updates the position and velocity for atoms in the
group each timestep and the box dimensions.

\sphinxAtStartPar
The thermostatting and barostatting is achieved by adding some dynamic
variables which are coupled to the particle velocities
(thermostatting) and simulation domain dimensions (barostatting).  In
addition to basic thermostatting and barostatting, this fix can
also create a chain of thermostats coupled to the particle thermostat,
and another chain of thermostats coupled to the barostat
variables. The barostat can be coupled to the overall box volume, or
to individual dimensions, including the \sphinxstyleemphasis{xy}, \sphinxstyleemphasis{xz} and \sphinxstyleemphasis{yz} tilt
dimensions. The external pressure of the barostat can be specified as
either a scalar pressure (isobaric ensemble) or as components of a
symmetric stress tensor (constant stress ensemble).  When used
correctly, the time\sphinxhyphen{}averaged temperature and stress tensor of the
particles will match the target values specified by Tstart/Tstop and
Pstart/Pstop.

\sphinxAtStartPar
The equations of motion used are those of Shinoda et al in
{\hyperref[\detokenize{fix_npt_cauchy:nc-shinoda}]{\sphinxcrossref{\DUrole{std,std-ref}{(Shinoda)}}}}, which combine the hydrostatic equations of
Martyna, Tobias and Klein in {\hyperref[\detokenize{fix_npt_cauchy:nc-martyna}]{\sphinxcrossref{\DUrole{std,std-ref}{(Martyna)}}}} with the strain
energy proposed by Parrinello and Rahman in
{\hyperref[\detokenize{fix_npt_cauchy:nc-parrinello}]{\sphinxcrossref{\DUrole{std,std-ref}{(Parrinello)}}}}.  The time integration schemes closely
follow the time\sphinxhyphen{}reversible measure\sphinxhyphen{}preserving Verlet and rRESPA
integrators derived by Tuckerman et al in {\hyperref[\detokenize{fix_npt_cauchy:nc-tuckerman}]{\sphinxcrossref{\DUrole{std,std-ref}{(Tuckerman)}}}}.


\bigskip\hrule\bigskip


\sphinxAtStartPar
The thermostat parameters are specified using the \sphinxstyleemphasis{temp} keyword.
Other thermostat\sphinxhyphen{}related keywords are \sphinxstyleemphasis{tchain}, \sphinxstyleemphasis{tloop} and \sphinxstyleemphasis{drag},
which are discussed below.

\sphinxAtStartPar
The thermostat is applied to only the translational degrees of freedom
for the particles.  The translational degrees of freedom can also have
a bias velocity removed before thermostatting takes place; see the
description below.  The desired temperature at each timestep is a
ramped value during the run from \sphinxstyleemphasis{Tstart} to \sphinxstyleemphasis{Tstop}.  The \sphinxstyleemphasis{Tdamp}
parameter is specified in time units and determines how rapidly the
temperature is relaxed.  For example, a value of 10.0 means to relax
the temperature in a timespan of (roughly) 10 time units (e.g. \(\tau\)
or fs or ps \sphinxhyphen{} see the {\hyperref[\detokenize{units::doc}]{\sphinxcrossref{\DUrole{doc}{units}}}} command).  The atoms in the
fix group are the only ones whose velocities and positions are updated
by the velocity/position update portion of the integration.

\begin{sphinxadmonition}{note}{Note:}
\sphinxAtStartPar
A Nose\sphinxhyphen{}Hoover thermostat will not work well for arbitrary values
of \sphinxstyleemphasis{Tdamp}.  If \sphinxstyleemphasis{Tdamp} is too small, the temperature can fluctuate
wildly; if it is too large, the temperature will take a very long time
to equilibrate.  A good choice for many models is a \sphinxstyleemphasis{Tdamp} of around
100 timesteps.  Note that this is NOT the same as 100 time units for
most {\hyperref[\detokenize{units::doc}]{\sphinxcrossref{\DUrole{doc}{units}}}} settings.
\end{sphinxadmonition}


\bigskip\hrule\bigskip


\sphinxAtStartPar
The barostat parameters are specified using one or more of the \sphinxstyleemphasis{iso},
\sphinxstyleemphasis{aniso}, \sphinxstyleemphasis{tri}, \sphinxstyleemphasis{x}, \sphinxstyleemphasis{y}, \sphinxstyleemphasis{z}, \sphinxstyleemphasis{xy}, \sphinxstyleemphasis{xz}, \sphinxstyleemphasis{yz}, and \sphinxstyleemphasis{couple} keywords.
These keywords give you the ability to specify all 6 components of an
external stress tensor, and to couple various of these components
together so that the dimensions they represent are varied together
during a constant\sphinxhyphen{}pressure simulation.

\sphinxAtStartPar
Other barostat\sphinxhyphen{}related keywords are \sphinxstyleemphasis{pchain}, \sphinxstyleemphasis{mtk}, \sphinxstyleemphasis{ploop},
\sphinxstyleemphasis{nreset}, \sphinxstyleemphasis{drag}, and \sphinxstyleemphasis{dilate}, which are discussed below.

\sphinxAtStartPar
Orthogonal simulation boxes have 3 adjustable dimensions (x,y,z).
Triclinic (non\sphinxhyphen{}orthogonal) simulation boxes have 6 adjustable
dimensions (x,y,z,xy,xz,yz).  The {\hyperref[\detokenize{create_box::doc}]{\sphinxcrossref{\DUrole{doc}{create\_box}}}}, {\hyperref[\detokenize{read_data::doc}]{\sphinxcrossref{\DUrole{doc}{read data}}}}, and {\hyperref[\detokenize{read_restart::doc}]{\sphinxcrossref{\DUrole{doc}{read\_restart}}}} commands
specify whether the simulation box is orthogonal or non\sphinxhyphen{}orthogonal
(triclinic) and explain the meaning of the xy,xz,yz tilt factors.

\sphinxAtStartPar
The target pressures for each of the 6 components of the stress tensor
can be specified independently via the \sphinxstyleemphasis{x}, \sphinxstyleemphasis{y}, \sphinxstyleemphasis{z}, \sphinxstyleemphasis{xy}, \sphinxstyleemphasis{xz}, \sphinxstyleemphasis{yz}
keywords, which correspond to the 6 simulation box dimensions.  For
each component, the external pressure or tensor component at each
timestep is a ramped value during the run from \sphinxstyleemphasis{Pstart} to \sphinxstyleemphasis{Pstop}.
If a target pressure is specified for a component, then the
corresponding box dimension will change during a simulation.  For
example, if the \sphinxstyleemphasis{y} keyword is used, the y\sphinxhyphen{}box length will change.  If
the \sphinxstyleemphasis{xy} keyword is used, the xy tilt factor will change.  A box
dimension will not change if that component is not specified, although
you have the option to change that dimension via the {\hyperref[\detokenize{fix_deform::doc}]{\sphinxcrossref{\DUrole{doc}{fix deform}}}} command.

\sphinxAtStartPar
Note that in order to use the \sphinxstyleemphasis{xy}, \sphinxstyleemphasis{xz}, or \sphinxstyleemphasis{yz} keywords, the
simulation box must be triclinic, even if its initial tilt factors are
0.0.

\sphinxAtStartPar
For all barostat keywords, the \sphinxstyleemphasis{Pdamp} parameter operates like the
\sphinxstyleemphasis{Tdamp} parameter, determining the time scale on which pressure is
relaxed.  For example, a value of 10.0 means to relax the pressure in
a timespan of (roughly) 10 time units (e.g. \(\tau\) or fs or ps
\sphinxhyphen{} see the {\hyperref[\detokenize{units::doc}]{\sphinxcrossref{\DUrole{doc}{units}}}} command).

\begin{sphinxadmonition}{note}{Note:}
\sphinxAtStartPar
A Nose\sphinxhyphen{}Hoover barostat will not work well for arbitrary values
of \sphinxstyleemphasis{Pdamp}.  If \sphinxstyleemphasis{Pdamp} is too small, the pressure and volume can
fluctuate wildly; if it is too large, the pressure will take a very
long time to equilibrate.  A good choice for many models is a \sphinxstyleemphasis{Pdamp}
of around 1000 timesteps.  However, note that \sphinxstyleemphasis{Pdamp} is specified in
time units, and that timesteps are NOT the same as time units for most
{\hyperref[\detokenize{units::doc}]{\sphinxcrossref{\DUrole{doc}{units}}}} settings.
\end{sphinxadmonition}

\sphinxAtStartPar
Regardless of what atoms are in the fix group (the only atoms which
are time integrated), a global pressure or stress tensor is computed
for all atoms.  Similarly, when the size of the simulation box is
changed, all atoms are re\sphinxhyphen{}scaled to new positions, unless the keyword
\sphinxstyleemphasis{dilate} is specified with a \sphinxstyleemphasis{dilate\sphinxhyphen{}group\sphinxhyphen{}ID} for a group that
represents a subset of the atoms.  This can be useful, for example, to
leave the coordinates of atoms in a solid substrate unchanged and
controlling the pressure of a surrounding fluid.  This option should
be used with care, since it can be unphysical to dilate some atoms and
not others, because it can introduce large, instantaneous
displacements between a pair of atoms (one dilated, one not) that are
far from the dilation origin.  Also note that for atoms not in the fix
group, a separate time integration fix like {\hyperref[\detokenize{fix_nve::doc}]{\sphinxcrossref{\DUrole{doc}{fix nve}}}} or
{\hyperref[\detokenize{fix_nh::doc}]{\sphinxcrossref{\DUrole{doc}{fix nvt}}}} can be used on them, independent of whether they
are dilated or not.


\bigskip\hrule\bigskip


\sphinxAtStartPar
The \sphinxstyleemphasis{couple} keyword allows two or three of the diagonal components of
the pressure tensor to be “coupled” together.  The value specified
with the keyword determines which are coupled.  For example, \sphinxstyleemphasis{xz}
means the \sphinxstyleemphasis{Pxx} and \sphinxstyleemphasis{Pzz} components of the stress tensor are coupled.
\sphinxstyleemphasis{Xyz} means all 3 diagonal components are coupled.  Coupling means two
things: the instantaneous stress will be computed as an average of the
corresponding diagonal components, and the coupled box dimensions will
be changed together in lockstep, meaning coupled dimensions will be
dilated or contracted by the same percentage every timestep.  The
\sphinxstyleemphasis{Pstart}, \sphinxstyleemphasis{Pstop}, \sphinxstyleemphasis{Pdamp} parameters for any coupled dimensions must
be identical.  \sphinxstyleemphasis{Couple xyz} can be used for a 2d simulation; the \sphinxstyleemphasis{z}
dimension is simply ignored.


\bigskip\hrule\bigskip


\sphinxAtStartPar
The \sphinxstyleemphasis{iso}, \sphinxstyleemphasis{aniso}, and \sphinxstyleemphasis{tri} keywords are simply shortcuts that are
equivalent to specifying several other keywords together.

\sphinxAtStartPar
The keyword \sphinxstyleemphasis{iso} means couple all 3 diagonal components together when
pressure is computed (hydrostatic pressure), and dilate/contract the
dimensions together.  Using “iso Pstart Pstop Pdamp” is the same as
specifying these 4 keywords:

\begin{sphinxVerbatim}[commandchars=\\\{\}]
x Pstart Pstop Pdamp
y Pstart Pstop Pdamp
z Pstart Pstop Pdamp
couple xyz
\end{sphinxVerbatim}

\sphinxAtStartPar
The keyword \sphinxstyleemphasis{aniso} means \sphinxstyleemphasis{x}, \sphinxstyleemphasis{y}, and \sphinxstyleemphasis{z} dimensions are controlled
independently using the \sphinxstyleemphasis{Pxx}, \sphinxstyleemphasis{Pyy}, and \sphinxstyleemphasis{Pzz} components of the
stress tensor as the driving forces, and the specified scalar external
pressure.  Using “aniso Pstart Pstop Pdamp” is the same as specifying
these 4 keywords:

\begin{sphinxVerbatim}[commandchars=\\\{\}]
x Pstart Pstop Pdamp
y Pstart Pstop Pdamp
z Pstart Pstop Pdamp
couple none
\end{sphinxVerbatim}

\sphinxAtStartPar
The keyword \sphinxstyleemphasis{tri} means \sphinxstyleemphasis{x}, \sphinxstyleemphasis{y}, \sphinxstyleemphasis{z}, \sphinxstyleemphasis{xy}, \sphinxstyleemphasis{xz}, and \sphinxstyleemphasis{yz} dimensions
are controlled independently using their individual stress components
as the driving forces, and the specified scalar pressure as the
external normal stress.  Using “tri Pstart Pstop Pdamp” is the same as
specifying these 7 keywords:

\begin{sphinxVerbatim}[commandchars=\\\{\}]
x Pstart Pstop Pdamp
y Pstart Pstop Pdamp
z Pstart Pstop Pdamp
xy 0.0 0.0 Pdamp
yz 0.0 0.0 Pdamp
xz 0.0 0.0 Pdamp
couple none
\end{sphinxVerbatim}


\bigskip\hrule\bigskip


\sphinxAtStartPar
In some cases (e.g. for solids) the pressure (volume) and/or
temperature of the system can oscillate undesirably when a Nose/Hoover
barostat and thermostat is applied.  The optional \sphinxstyleemphasis{drag} keyword will
damp these oscillations, although it alters the Nose/Hoover equations.
A value of 0.0 (no drag) leaves the Nose/Hoover formalism unchanged.
A non\sphinxhyphen{}zero value adds a drag term; the larger the value specified, the
greater the damping effect.  Performing a short run and monitoring the
pressure and temperature is the best way to determine if the drag term
is working.  Typically a value between 0.2 to 2.0 is sufficient to
damp oscillations after a few periods. Note that use of the drag
keyword will interfere with energy conservation and will also change
the distribution of positions and velocities so that they do not
correspond to the nominal NVT, NPT, or NPH ensembles.

\sphinxAtStartPar
An alternative way to control initial oscillations is to use chain
thermostats. The keyword \sphinxstyleemphasis{tchain} determines the number of thermostats
in the particle thermostat. A value of 1 corresponds to the original
Nose\sphinxhyphen{}Hoover thermostat. The keyword \sphinxstyleemphasis{pchain} specifies the number of
thermostats in the chain thermostatting the barostat degrees of
freedom. A value of 0 corresponds to no thermostatting of the
barostat variables.

\sphinxAtStartPar
The \sphinxstyleemphasis{mtk} keyword controls whether or not the correction terms due to
Martyna, Tuckerman, and Klein are included in the equations of motion
{\hyperref[\detokenize{fix_npt_cauchy:nc-martyna}]{\sphinxcrossref{\DUrole{std,std-ref}{(Martyna)}}}}.  Specifying \sphinxstyleemphasis{no} reproduces the original
Hoover barostat, whose volume probability distribution function
differs from the true NPT and NPH ensembles by a factor of 1/V.  Hence
using \sphinxstyleemphasis{yes} is more correct, but in many cases the difference is
negligible.

\sphinxAtStartPar
The keyword \sphinxstyleemphasis{tloop} can be used to improve the accuracy of integration
scheme at little extra cost.  The initial and final updates of the
thermostat variables are broken up into \sphinxstyleemphasis{tloop} sub\sphinxhyphen{}steps, each of
length \sphinxstyleemphasis{dt}/\sphinxstyleemphasis{tloop}. This corresponds to using a first\sphinxhyphen{}order
Suzuki\sphinxhyphen{}Yoshida scheme {\hyperref[\detokenize{fix_npt_cauchy:nc-tuckerman}]{\sphinxcrossref{\DUrole{std,std-ref}{(Tuckerman)}}}}.  The keyword \sphinxstyleemphasis{ploop}
does the same thing for the barostat thermostat.

\sphinxAtStartPar
The keyword \sphinxstyleemphasis{nreset} controls how often the reference dimensions used
to define the strain energy are reset.  If this keyword is not used,
or is given a value of zero, then the reference dimensions are set to
those of the initial simulation domain and are never changed. If the
simulation domain changes significantly during the simulation, then
the final average pressure tensor will differ significantly from the
specified values of the external stress tensor.  A value of \sphinxstyleemphasis{nstep}
means that every \sphinxstyleemphasis{nstep} timesteps, the reference dimensions are set
to those of the current simulation domain.

\sphinxAtStartPar
The \sphinxstyleemphasis{scaleyz}, \sphinxstyleemphasis{scalexz}, and \sphinxstyleemphasis{scalexy} keywords control whether or
not the corresponding tilt factors are scaled with the associated box
dimensions when barostatting triclinic periodic cells.  The default
values \sphinxstyleemphasis{yes} will turn on scaling, which corresponds to adjusting the
linear dimensions of the cell while preserving its shape.  Choosing
\sphinxstyleemphasis{no} ensures that the tilt factors are not scaled with the box
dimensions. See below for restrictions and default values in different
situations. In older versions of LAMMPS, scaling of tilt factors was
not performed. The old behavior can be recovered by setting all three
scale keywords to \sphinxstyleemphasis{no}.

\sphinxAtStartPar
The \sphinxstyleemphasis{flip} keyword allows the tilt factors for a triclinic box to
exceed half the distance of the parallel box length, as discussed
below.  If the \sphinxstyleemphasis{flip} value is set to \sphinxstyleemphasis{yes}, the bound is enforced by
flipping the box when it is exceeded.  If the \sphinxstyleemphasis{flip} value is set to
\sphinxstyleemphasis{no}, the tilt will continue to change without flipping.  Note that if
applied stress induces large deformations (e.g. in a liquid), this
means the box shape can tilt dramatically and LAMMPS will run less
efficiently, due to the large volume of communication needed to
acquire ghost atoms around a processor’s irregular\sphinxhyphen{}shaped subdomain.
For extreme values of tilt, LAMMPS may also lose atoms and generate an
error.

\sphinxAtStartPar
The \sphinxstyleemphasis{fixedpoint} keyword specifies the fixed point for barostat volume
changes. By default, it is the center of the box.  Whatever point is
chosen will not move during the simulation.  For example, if the lower
periodic boundaries pass through (0,0,0), and this point is provided
to \sphinxstyleemphasis{fixedpoint}, then the lower periodic boundaries will remain at
(0,0,0), while the upper periodic boundaries will move twice as
far. In all cases, the particle trajectories are unaffected by the
chosen value, except for a time\sphinxhyphen{}dependent constant translation of
positions.


\bigskip\hrule\bigskip


\begin{sphinxadmonition}{note}{Note:}
\sphinxAtStartPar
Using a barostat coupled to tilt dimensions \sphinxstyleemphasis{xy}, \sphinxstyleemphasis{xz}, \sphinxstyleemphasis{yz} can
sometimes result in arbitrarily large values of the tilt dimensions,
i.e. a dramatically deformed simulation box.  LAMMPS allows the tilt
factors to grow a small amount beyond the normal limit of half the box
length (0.6 times the box length), and then performs a box “flip” to
an equivalent periodic cell.  See the discussion of the \sphinxstyleemphasis{flip} keyword
above, to allow this bound to be exceeded, if desired.
\end{sphinxadmonition}

\sphinxAtStartPar
The flip operation is described in more detail in the page for
{\hyperref[\detokenize{fix_deform::doc}]{\sphinxcrossref{\DUrole{doc}{fix deform}}}}.  Both the barostat dynamics and the atom
trajectories are unaffected by this operation.  However, if a tilt
factor is incremented by a large amount (1.5 times the box length) on
a single timestep, LAMMPS can not accommodate this event and will
terminate the simulation with an error. This error typically indicates
that there is something badly wrong with how the simulation was
constructed, such as specifying values of \sphinxstyleemphasis{Pstart} that are too far
from the current stress value, or specifying a timestep that is too
large. Triclinic barostatting should be used with care. This also is
true for other barostat styles, although they tend to be more
forgiving of insults. In particular, it is important to recognize that
equilibrium liquids can not support a shear stress and that
equilibrium solids can not support shear stresses that exceed the
yield stress.

\sphinxAtStartPar
One exception to this rule is if the first dimension in the tilt factor
(x for xy) is non\sphinxhyphen{}periodic.  In that case, the limits on the tilt
factor are not enforced, since flipping the box in that dimension does
not change the atom positions due to non\sphinxhyphen{}periodicity.  In this mode,
if you tilt the system to extreme angles, the simulation will simply
become inefficient due to the highly skewed simulation box.

\begin{sphinxadmonition}{note}{Note:}
\sphinxAtStartPar
Unlike the {\hyperref[\detokenize{fix_temp_berendsen::doc}]{\sphinxcrossref{\DUrole{doc}{fix temp/berendsen}}}} command
which performs thermostatting but NO time integration, this fix
performs thermostatting/barostatting AND time integration.  Thus you
should not use any other time integration fix, such as {\hyperref[\detokenize{fix_nve::doc}]{\sphinxcrossref{\DUrole{doc}{fix nve}}}} on atoms to which this fix is applied.  Likewise,
fix npt/cauchy should not normally be used on atoms that also
have their temperature controlled by another fix \sphinxhyphen{} e.g. by {\hyperref[\detokenize{fix_nh::doc}]{\sphinxcrossref{\DUrole{doc}{fix langevin}}}} or {\hyperref[\detokenize{fix_temp_rescale::doc}]{\sphinxcrossref{\DUrole{doc}{fix temp/rescale}}}}
commands.
\end{sphinxadmonition}

\sphinxAtStartPar
See the {\hyperref[\detokenize{Howto_thermostat::doc}]{\sphinxcrossref{\DUrole{doc}{Howto thermostat}}}} and {\hyperref[\detokenize{Howto_barostat::doc}]{\sphinxcrossref{\DUrole{doc}{Howto barostat}}}} doc pages for a discussion of different
ways to compute temperature and perform thermostatting and
barostatting.


\bigskip\hrule\bigskip


\sphinxAtStartPar
This fix compute a temperature and pressure each timestep.  To do
this, the fix creates its own computes of style “temp” and “pressure”,
as if one of these sets of commands had been issued:

\begin{sphinxVerbatim}[commandchars=\\\{\}]
\PYG{k}{compute }\PYG{n+nv+nvIdentifier}{fix\PYGZhy{}ID\PYGZus{}temp}\PYG{+w}{ }\PYG{n+nv+nvIdentifier}{all}\PYG{+w}{ }\PYG{n}{temp}
\PYG{k}{compute }\PYG{n+nv+nvIdentifier}{fix\PYGZhy{}ID\PYGZus{}press}\PYG{+w}{ }\PYG{n+nv+nvIdentifier}{all}\PYG{+w}{ }\PYG{n}{pressure}\PYG{+w}{ }\PYG{n}{fix}\PYG{o}{\PYGZhy{}}\PYG{n}{ID\PYGZus{}temp}
\end{sphinxVerbatim}

\sphinxAtStartPar
The group for both the new temperature and pressure compute is “all”
since pressure is computed for the entire system.  See the {\hyperref[\detokenize{compute_temp::doc}]{\sphinxcrossref{\DUrole{doc}{compute temp}}}} and {\hyperref[\detokenize{compute_pressure::doc}]{\sphinxcrossref{\DUrole{doc}{compute pressure}}}}
commands for details.  Note that the IDs of the new computes are the
fix\sphinxhyphen{}ID + underscore + “temp” or fix\_ID + underscore + “press”.

\sphinxAtStartPar
Note that these are NOT the computes used by thermodynamic output (see
the {\hyperref[\detokenize{thermo_style::doc}]{\sphinxcrossref{\DUrole{doc}{thermo\_style}}}} command) with ID = \sphinxstyleemphasis{thermo\_temp}
and \sphinxstyleemphasis{thermo\_press}.  This means you can change the attributes of these
fix’s temperature or pressure via the
{\hyperref[\detokenize{compute_modify::doc}]{\sphinxcrossref{\DUrole{doc}{compute\_modify}}}} command.  Or you can print this
temperature or pressure during thermodynamic output via the
{\hyperref[\detokenize{thermo_style::doc}]{\sphinxcrossref{\DUrole{doc}{thermo\_style custom}}}} command using the appropriate
compute\sphinxhyphen{}ID.  It also means that changing attributes of \sphinxstyleemphasis{thermo\_temp}
or \sphinxstyleemphasis{thermo\_press} will have no effect on this fix.

\sphinxAtStartPar
Like other fixes that perform thermostatting, this fix can be used
with {\hyperref[\detokenize{compute::doc}]{\sphinxcrossref{\DUrole{doc}{compute commands}}}} that remove a “bias” from the
atom velocities.  E.g. to apply the thermostat only to atoms within a
spatial {\hyperref[\detokenize{region::doc}]{\sphinxcrossref{\DUrole{doc}{region}}}}, or to remove the center\sphinxhyphen{}of\sphinxhyphen{}mass
velocity from a group of atoms, or to remove the x\sphinxhyphen{}component of
velocity from the calculation.

\sphinxAtStartPar
This is not done by default, but only if the {\hyperref[\detokenize{fix_modify::doc}]{\sphinxcrossref{\DUrole{doc}{fix\_modify}}}} command is used to assign a temperature compute to this
fix that includes such a bias term.  See the doc pages for individual
{\hyperref[\detokenize{compute::doc}]{\sphinxcrossref{\DUrole{doc}{compute temp commands}}}} to determine which ones include
a bias.  In this case, the thermostat works in the following manner:
bias is removed from each atom, thermostatting is performed on the
remaining thermal degrees of freedom, and the bias is added back in.


\bigskip\hrule\bigskip


\sphinxAtStartPar
This fix can be used with either the \sphinxstyleemphasis{verlet} or \sphinxstyleemphasis{respa}
{\hyperref[\detokenize{run_style::doc}]{\sphinxcrossref{\DUrole{doc}{integrators}}}}. When using this fix
with \sphinxstyleemphasis{respa}, LAMMPS uses an integrator constructed
according to the following factorization of the Liouville propagator
(for two rRESPA levels):
\begin{equation*}
\begin{split}\exp \left(\mathrm{i} L \Delta t \right) = & \hat{E}
\exp \left(\mathrm{i} L_{\rm T\textrm{-}baro} \frac{\Delta t}{2} \right)
\exp \left(\mathrm{i} L_{\rm T\textrm{-}part} \frac{\Delta t}{2} \right)
\exp \left(\mathrm{i} L_{\epsilon , 2} \frac{\Delta t}{2} \right)
\exp \left(\mathrm{i} L_{2}^{(2)} \frac{\Delta t}{2} \right) \\
&\times \left[
\exp \left(\mathrm{i} L_{2}^{(1)} \frac{\Delta t}{2n} \right)
\exp \left(\mathrm{i} L_{\epsilon , 1} \frac{\Delta t}{2n} \right)
\exp \left(\mathrm{i} L_1 \frac{\Delta t}{n} \right)
\exp \left(\mathrm{i} L_{\epsilon , 1} \frac{\Delta t}{2n} \right)
\exp \left(\mathrm{i} L_{2}^{(1)} \frac{\Delta t}{2n} \right)
\right]^n \\
&\times
\exp \left(\mathrm{i} L_{2}^{(2)} \frac{\Delta t}{2} \right)
\exp \left(\mathrm{i} L_{\epsilon , 2} \frac{\Delta t}{2} \right)
\exp \left(\mathrm{i} L_{\rm T\textrm{-}part} \frac{\Delta t}{2} \right)
\exp \left(\mathrm{i} L_{\rm T\textrm{-}baro} \frac{\Delta t}{2} \right) \\
&+ \mathcal{O} \left(\Delta t^3 \right)\end{split}
\end{equation*}
\sphinxAtStartPar
This factorization differs somewhat from that of Tuckerman et al, in
that the barostat is only updated at the outermost rRESPA level,
whereas Tuckerman’s factorization requires splitting the pressure into
pieces corresponding to the forces computed at each rRESPA level. In
theory, the latter method will exhibit better numerical stability. In
practice, because Pdamp is normally chosen to be a large multiple of
the outermost rRESPA timestep, the barostat dynamics are not the
limiting factor for numerical stability. Both factorizations are
time\sphinxhyphen{}reversible and can be shown to preserve the phase space measure
of the underlying non\sphinxhyphen{}Hamiltonian equations of motion.

\begin{sphinxadmonition}{note}{Note:}
\sphinxAtStartPar
Under NPT dynamics, for a system with zero initial total linear
momentum, the total momentum fluctuates close to zero.  It may occasionally
undergo brief excursions to non\sphinxhyphen{}negligible values, before returning close
to zero.  Over long simulations, this has the effect of causing the
center\sphinxhyphen{}of\sphinxhyphen{}mass to undergo a slow random walk. This can be mitigated by
resetting the momentum at infrequent intervals using the
{\hyperref[\detokenize{fix_momentum::doc}]{\sphinxcrossref{\DUrole{doc}{fix momentum}}}} command.
\end{sphinxadmonition}


\bigskip\hrule\bigskip



\subsection{Restart, fix\_modify, output, run start/stop, minimize info}
\label{\detokenize{fix_npt_cauchy:restart-fix-modify-output-run-start-stop-minimize-info}}
\sphinxAtStartPar
This fix writes the state of all the thermostat and barostat
variables to {\hyperref[\detokenize{restart::doc}]{\sphinxcrossref{\DUrole{doc}{binary restart files}}}}.  See the
{\hyperref[\detokenize{read_restart::doc}]{\sphinxcrossref{\DUrole{doc}{read\_restart}}}} command for info on how to re\sphinxhyphen{}specify
a fix in an input script that reads a restart file, so that the
operation of the fix continues in an uninterrupted fashion.

\sphinxAtStartPar
The {\hyperref[\detokenize{fix_modify::doc}]{\sphinxcrossref{\DUrole{doc}{fix\_modify}}}} \sphinxstyleemphasis{temp} and \sphinxstyleemphasis{press} options are
supported by this fix.  You can use them to assign a
{\hyperref[\detokenize{compute::doc}]{\sphinxcrossref{\DUrole{doc}{compute}}}} you have defined to this fix which will be used
in its thermostatting or barostatting procedure, as described above.
If you do this, note that the kinetic energy derived from the compute
temperature should be consistent with the virial term computed using
all atoms for the pressure.  LAMMPS will warn you if you choose to
compute temperature on a subset of atoms.

\begin{sphinxadmonition}{note}{Note:}
\sphinxAtStartPar
If both the \sphinxstyleemphasis{temp} and \sphinxstyleemphasis{press} keywords are used in a single
thermo\_modify command (or in two separate commands), then the order
in which the keywords are specified is important.  Note that a
{\hyperref[\detokenize{compute_pressure::doc}]{\sphinxcrossref{\DUrole{doc}{pressure compute}}}} defines its own
temperature compute as an argument when it is specified.  The
\sphinxstyleemphasis{temp} keyword will override this (for the pressure compute being
used by fix npt), but only if the \sphinxstyleemphasis{temp} keyword comes after the
\sphinxstyleemphasis{press} keyword.  If the \sphinxstyleemphasis{temp} keyword comes before the \sphinxstyleemphasis{press}
keyword, then the new pressure compute specified by the \sphinxstyleemphasis{press}
keyword will be unaffected by the \sphinxstyleemphasis{temp} setting.
\end{sphinxadmonition}

\sphinxAtStartPar
The cumulative energy change in the system imposed by this fix, due to
thermostatting and/or barostatting, is included in the
{\hyperref[\detokenize{thermo_style::doc}]{\sphinxcrossref{\DUrole{doc}{thermodynamic output}}}} keywords \sphinxstyleemphasis{ecouple} and
\sphinxstyleemphasis{econserve}.  See the {\hyperref[\detokenize{thermo_style::doc}]{\sphinxcrossref{\DUrole{doc}{thermo\_style}}}} page for
details.

\sphinxAtStartPar
This fix computes a global scalar which can be accessed by various
{\hyperref[\detokenize{Howto_output::doc}]{\sphinxcrossref{\DUrole{doc}{output commands}}}}.  The scalar is the same
cumulative energy change due to this fix described in the previous
paragraph.  The scalar value calculated by this fix is “extensive”.

\sphinxAtStartPar
This fix also computes a global vector of quantities, which can be
accessed by various {\hyperref[\detokenize{Howto_output::doc}]{\sphinxcrossref{\DUrole{doc}{output commands}}}}.  Rhe vector
values are “intensive”.

\sphinxAtStartPar
The vector stores internal Nose/Hoover thermostat and barostat
variables.  The number and meaning of the vector values depends on
which fix is used and the settings for keywords \sphinxstyleemphasis{tchain} and \sphinxstyleemphasis{pchain},
which specify the number of Nose/Hoover chains for the thermostat and
barostat.  If no thermostatting is done, then \sphinxstyleemphasis{tchain} is 0.  If no
barostatting is done, then \sphinxstyleemphasis{pchain} is 0.  In the following list,
“ndof” is 0, 1, 3, or 6, and is the number of degrees of freedom in
the barostat.  Its value is 0 if no barostat is used, else its value
is 6 if any off\sphinxhyphen{}diagonal stress tensor component is barostatted, else
its value is 1 if \sphinxstyleemphasis{couple xyz} is used or \sphinxstyleemphasis{couple xy} for a 2d
simulation, otherwise its value is 3.

\sphinxAtStartPar
The order of values in the global vector and their meaning is as
follows.  The notation means there are tchain values for eta, followed
by tchain for eta\_dot, followed by ndof for omega, etc:
\begin{itemize}
\item {} 
\sphinxAtStartPar
eta{[}tchain{]} = particle thermostat displacements (unitless)

\item {} 
\sphinxAtStartPar
eta\_dot{[}tchain{]} = particle thermostat velocities (1/time units)

\item {} 
\sphinxAtStartPar
omega{[}ndof{]} = barostat displacements (unitless)

\item {} 
\sphinxAtStartPar
omega\_dot{[}ndof{]} = barostat velocities (1/time units)

\item {} 
\sphinxAtStartPar
etap{[}pchain{]} = barostat thermostat displacements (unitless)

\item {} 
\sphinxAtStartPar
etap\_dot{[}pchain{]} = barostat thermostat velocities (1/time units)

\item {} 
\sphinxAtStartPar
PE\_eta{[}tchain{]} = potential energy of each particle thermostat displacement (energy units)

\item {} 
\sphinxAtStartPar
KE\_eta\_dot{[}tchain{]} = kinetic energy of each particle thermostat velocity (energy units)

\item {} 
\sphinxAtStartPar
PE\_omega{[}ndof{]} = potential energy of each barostat displacement (energy units)

\item {} 
\sphinxAtStartPar
KE\_omega\_dot{[}ndof{]} = kinetic energy of each barostat velocity (energy units)

\item {} 
\sphinxAtStartPar
PE\_etap{[}pchain{]} = potential energy of each barostat thermostat displacement (energy units)

\item {} 
\sphinxAtStartPar
KE\_etap\_dot{[}pchain{]} = kinetic energy of each barostat thermostat velocity (energy units)

\item {} 
\sphinxAtStartPar
PE\_strain{[}1{]} = scalar strain energy (energy units)

\end{itemize}

\sphinxAtStartPar
This fix can ramp its external temperature and pressure over
multiple runs, using the \sphinxstyleemphasis{start} and \sphinxstyleemphasis{stop} keywords of the
{\hyperref[\detokenize{run::doc}]{\sphinxcrossref{\DUrole{doc}{run}}}} command.  See the {\hyperref[\detokenize{run::doc}]{\sphinxcrossref{\DUrole{doc}{run}}}} command for details of
how to do this.

\sphinxAtStartPar
This fix is not invoked during {\hyperref[\detokenize{minimize::doc}]{\sphinxcrossref{\DUrole{doc}{energy minimization}}}}.


\bigskip\hrule\bigskip



\subsection{Restrictions}
\label{\detokenize{fix_npt_cauchy:restrictions}}
\sphinxAtStartPar
This fix is part of the EXTRA\sphinxhyphen{}FIX package.  It is only enabled if
LAMMPS was built with that package.  See the {\hyperref[\detokenize{Build_package::doc}]{\sphinxcrossref{\DUrole{doc}{Build package}}}} page for more info.

\sphinxAtStartPar
\sphinxstyleemphasis{X}, \sphinxstyleemphasis{y}, \sphinxstyleemphasis{z} cannot be barostatted if the associated dimension is not
periodic.  \sphinxstyleemphasis{Xy}, \sphinxstyleemphasis{xz}, and \sphinxstyleemphasis{yz} can only be barostatted if the
simulation domain is triclinic and the second dimension in the keyword
(\sphinxstyleemphasis{y} dimension in \sphinxstyleemphasis{xy}) is periodic.  \sphinxstyleemphasis{Z}, \sphinxstyleemphasis{xz}, and \sphinxstyleemphasis{yz}, cannot be
barostatted for 2D simulations.  The {\hyperref[\detokenize{create_box::doc}]{\sphinxcrossref{\DUrole{doc}{create\_box}}}},
{\hyperref[\detokenize{read_data::doc}]{\sphinxcrossref{\DUrole{doc}{read data}}}}, and {\hyperref[\detokenize{read_restart::doc}]{\sphinxcrossref{\DUrole{doc}{read\_restart}}}}
commands specify whether the simulation box is orthogonal or
non\sphinxhyphen{}orthogonal (triclinic) and explain the meaning of the xy,xz,yz
tilt factors.

\sphinxAtStartPar
For the \sphinxstyleemphasis{temp} keyword, the final Tstop cannot be 0.0 since it would
make the external T = 0.0 at some timestep during the simulation which
is not allowed in the Nose/Hoover formulation.

\sphinxAtStartPar
The \sphinxstyleemphasis{scaleyz yes} and \sphinxstyleemphasis{scalexz yes} keyword/value pairs can not be used
for 2D simulations. \sphinxstyleemphasis{scaleyz yes}, \sphinxstyleemphasis{scalexz yes}, and \sphinxstyleemphasis{scalexy yes} options
can only be used if the second dimension in the keyword is periodic,
and if the tilt factor is not coupled to the barostat via keywords
\sphinxstyleemphasis{tri}, \sphinxstyleemphasis{yz}, \sphinxstyleemphasis{xz}, and \sphinxstyleemphasis{xy}.

\sphinxAtStartPar
The \sphinxstyleemphasis{alpha} keyword modifies the barostat as per Miller et
al. (Miller)\_”\#nc\sphinxhyphen{}Miller” so that the Cauchy stress is controlled.
\sphinxstyleemphasis{alpha} is the non\sphinxhyphen{}dimensional parameter, typically set to 0.001 or
0.01 that determines how aggressively the algorithm drives the system
towards the set Cauchy stresses.  Larger values of \sphinxstyleemphasis{alpha} will modify
the system more quickly, but can lead to instabilities.  Smaller
values will lead to longer convergence time.  Since \sphinxstyleemphasis{alpha} also
influences how much the stress fluctuations deviate from the
equilibrium fluctuations, it should be set as small as possible.

\sphinxAtStartPar
A \sphinxstyleemphasis{continue} value of \sphinxstyleemphasis{yes} indicates that the fix is subsequent to a
previous run with the npt/cauchy fix, and the intention is to continue
from the converged stress state at the end of the previous run.  This
may be required, for example, when implementing a multi\sphinxhyphen{}step loading/unloading
sequence over several fixes.

\sphinxAtStartPar
Setting \sphinxstyleemphasis{alpha} to zero is not permitted.  To “turn off” the
cauchystat control and thus restore the equilibrium stress
fluctuations, two subsequent fixes should be used.  In the first, fix
npt/cauchy is used and the simulation box equilibrates to the
correct shape for the desired stresses.  In the second,
{\hyperref[\detokenize{fix_nh::doc}]{\sphinxcrossref{\DUrole{doc}{fix npt}}}} is used instead which uses the
original Parrinello\sphinxhyphen{}Rahman algorithm, but now with the corrected
simulation box shape from using fix npt/cauchy.

\sphinxAtStartPar
This fix can be used with dynamic groups as defined by the
{\hyperref[\detokenize{group::doc}]{\sphinxcrossref{\DUrole{doc}{group}}}} command.  Likewise it can be used with groups to
which atoms are added or deleted over time, e.g. a deposition
simulation.  However, the conservation properties of the thermostat
and barostat are defined for systems with a static set of atoms.  You
may observe odd behavior if the atoms in a group vary dramatically
over time or the atom count becomes very small.


\subsection{Related commands}
\label{\detokenize{fix_npt_cauchy:related-commands}}
\sphinxAtStartPar
{\hyperref[\detokenize{fix_nve::doc}]{\sphinxcrossref{\DUrole{doc}{fix nve}}}}, {\hyperref[\detokenize{fix_modify::doc}]{\sphinxcrossref{\DUrole{doc}{fix\_modify}}}},
{\hyperref[\detokenize{run_style::doc}]{\sphinxcrossref{\DUrole{doc}{run\_style}}}}


\subsection{Default}
\label{\detokenize{fix_npt_cauchy:default}}
\sphinxAtStartPar
The keyword defaults are tchain = 3, pchain = 3, mtk = yes, tloop =
ploop = 1, nreset = 0, drag = 0.0, dilate = all, couple = none,
cauchystat = no,
scaleyz = scalexz = scalexy = yes if periodic in second dimension and
not coupled to barostat, otherwise no.


\bigskip\hrule\bigskip

\phantomsection\label{\detokenize{fix_npt_cauchy:nc-martyna}}
\sphinxAtStartPar
\sphinxstylestrong{(Martyna)} Martyna, Tobias and Klein, J Chem Phys, 101, 4177 (1994).

\phantomsection\label{\detokenize{fix_npt_cauchy:nc-parrinello}}
\sphinxAtStartPar
\sphinxstylestrong{(Parrinello)} Parrinello and Rahman, J Appl Phys, 52, 7182 (1981).

\phantomsection\label{\detokenize{fix_npt_cauchy:nc-tuckerman}}
\sphinxAtStartPar
\sphinxstylestrong{(Tuckerman)} Tuckerman, Alejandre, Lopez\sphinxhyphen{}Rendon, Jochim, and
Martyna, J Phys A: Math Gen, 39, 5629 (2006).

\phantomsection\label{\detokenize{fix_npt_cauchy:nc-shinoda}}
\sphinxAtStartPar
\sphinxstylestrong{(Shinoda)} Shinoda, Shiga, and Mikami, Phys Rev B, 69, 134103 (2004).

\phantomsection\label{\detokenize{fix_npt_cauchy:nc-miller}}
\sphinxAtStartPar
\sphinxstylestrong{(Miller)} Miller, Tadmor, Gibson, Bernstein and Pavia, J Chem Phys,
144, 184107 (2016).

\sphinxstepscope

\index{fix npt/sphere@\spxentry{fix npt/sphere}}\ignorespaces 
\index{fix npt/sphere/omp@\spxentry{fix npt/sphere/omp}}\ignorespaces 

\section{fix npt/sphere command}
\label{\detokenize{fix_npt_sphere:fix-npt-sphere-command}}\label{\detokenize{fix_npt_sphere:index-1}}\label{\detokenize{fix_npt_sphere:index-0}}\label{\detokenize{fix_npt_sphere::doc}}
\sphinxAtStartPar
Accelerator Variants: \sphinxstyleemphasis{npt/sphere/omp}


\subsection{Syntax}
\label{\detokenize{fix_npt_sphere:syntax}}
\begin{sphinxVerbatim}[commandchars=\\\{\}]
\PYG{k}{fix }\PYG{n+nv+nvIdentifier}{ID}\PYG{+w}{ }\PYG{n+nv+nvIdentifier}{group\PYGZhy{}ID}\PYG{+w}{ }\PYG{n}{npt}\PYG{o}{/}\PYG{n}{sphere}\PYG{+w}{ }\PYG{n}{keyword}\PYG{+w}{ }\PYG{n}{value}\PYG{+w}{ }\PYG{n}{...}
\end{sphinxVerbatim}
\begin{itemize}
\item {} 
\sphinxAtStartPar
ID, group\sphinxhyphen{}ID are documented in {\hyperref[\detokenize{fix::doc}]{\sphinxcrossref{\DUrole{doc}{fix}}}} command
npt/sphere = style name of this fix command
zero or more keyword/value pairs may be appended

\item {} 
\sphinxAtStartPar
keyword = \sphinxstyleemphasis{disc}

\end{itemize}
\begin{sphinxalltt}
\sphinxstyleemphasis{disc} value = none = treat particles as 2d discs, not spheres
\end{sphinxalltt}
\begin{itemize}
\item {} 
\sphinxAtStartPar
NOTE: additional thermostat and barostat and dipole related keyword/value pairs from the {\hyperref[\detokenize{fix_nh::doc}]{\sphinxcrossref{\DUrole{doc}{fix npt}}}} command can be appended

\end{itemize}


\subsection{Examples}
\label{\detokenize{fix_npt_sphere:examples}}
\begin{sphinxVerbatim}[commandchars=\\\{\}]
\PYG{k}{fix }\PYG{n+nv+nvIdentifier}{1}\PYG{+w}{ }\PYG{n+nv+nvIdentifier}{all}\PYG{+w}{ }\PYG{n}{npt}\PYG{o}{/}\PYG{n}{sphere}\PYG{+w}{ }\PYG{n}{temp}\PYG{+w}{ }\PYG{l+m}{300.0}\PYG{+w}{ }\PYG{l+m}{300.0}\PYG{+w}{ }\PYG{l+m}{100.0}\PYG{+w}{ }\PYG{n}{iso}\PYG{+w}{ }\PYG{l+m}{0.0}\PYG{+w}{ }\PYG{l+m}{0.0}\PYG{+w}{ }\PYG{l+m}{1000.0}
\PYG{k}{fix }\PYG{n+nv+nvIdentifier}{2}\PYG{+w}{ }\PYG{n+nv+nvIdentifier}{all}\PYG{+w}{ }\PYG{n}{npt}\PYG{o}{/}\PYG{n}{sphere}\PYG{+w}{ }\PYG{n}{temp}\PYG{+w}{ }\PYG{l+m}{300.0}\PYG{+w}{ }\PYG{l+m}{300.0}\PYG{+w}{ }\PYG{l+m}{100.0}\PYG{+w}{ }\PYG{n}{x}\PYG{+w}{ }\PYG{l+m}{5.0}\PYG{+w}{ }\PYG{l+m}{5.0}\PYG{+w}{ }\PYG{l+m}{1000.0}
\PYG{k}{fix }\PYG{n+nv+nvIdentifier}{2}\PYG{+w}{ }\PYG{n+nv+nvIdentifier}{all}\PYG{+w}{ }\PYG{n}{npt}\PYG{o}{/}\PYG{n}{sphere}\PYG{+w}{ }\PYG{n}{temp}\PYG{+w}{ }\PYG{l+m}{300.0}\PYG{+w}{ }\PYG{l+m}{300.0}\PYG{+w}{ }\PYG{l+m}{100.0}\PYG{+w}{ }\PYG{n}{x}\PYG{+w}{ }\PYG{l+m}{5.0}\PYG{+w}{ }\PYG{l+m}{5.0}\PYG{+w}{ }\PYG{l+m}{1000.0}\PYG{+w}{ }\PYG{n}{disc}
\PYG{k}{fix }\PYG{n+nv+nvIdentifier}{2}\PYG{+w}{ }\PYG{n+nv+nvIdentifier}{all}\PYG{+w}{ }\PYG{n}{npt}\PYG{o}{/}\PYG{n}{sphere}\PYG{+w}{ }\PYG{n}{temp}\PYG{+w}{ }\PYG{l+m}{300.0}\PYG{+w}{ }\PYG{l+m}{300.0}\PYG{+w}{ }\PYG{l+m}{100.0}\PYG{+w}{ }\PYG{n}{x}\PYG{+w}{ }\PYG{l+m}{5.0}\PYG{+w}{ }\PYG{l+m}{5.0}\PYG{+w}{ }\PYG{l+m}{1000.0}\PYG{+w}{ }\PYG{n}{drag}\PYG{+w}{ }\PYG{l+m}{0.2}
\PYG{k}{fix }\PYG{n+nv+nvIdentifier}{2}\PYG{+w}{ }\PYG{n+nv+nvIdentifier}{all}\PYG{+w}{ }\PYG{n}{npt}\PYG{o}{/}\PYG{n}{sphere}\PYG{+w}{ }\PYG{n}{temp}\PYG{+w}{ }\PYG{l+m}{300.0}\PYG{+w}{ }\PYG{l+m}{300.0}\PYG{+w}{ }\PYG{l+m}{100.0}\PYG{+w}{ }\PYG{n}{x}\PYG{+w}{ }\PYG{l+m}{5.0}\PYG{+w}{ }\PYG{l+m}{5.0}\PYG{+w}{ }\PYG{l+m}{1000.0}\PYG{+w}{ }\PYG{n}{update}\PYG{+w}{ }\PYG{n}{dipole}
\PYG{k}{fix }\PYG{n+nv+nvIdentifier}{2}\PYG{+w}{ }\PYG{n+nv+nvIdentifier}{water}\PYG{+w}{ }\PYG{n}{npt}\PYG{o}{/}\PYG{n}{sphere}\PYG{+w}{ }\PYG{n}{temp}\PYG{+w}{ }\PYG{l+m}{300.0}\PYG{+w}{ }\PYG{l+m}{300.0}\PYG{+w}{ }\PYG{l+m}{100.0}\PYG{+w}{ }\PYG{n}{aniso}\PYG{+w}{ }\PYG{l+m}{0.0}\PYG{+w}{ }\PYG{l+m}{0.0}\PYG{+w}{ }\PYG{l+m}{1000.0}\PYG{+w}{ }\PYG{n}{dilate}\PYG{+w}{ }\PYG{n}{partial}
\end{sphinxVerbatim}


\subsection{Description}
\label{\detokenize{fix_npt_sphere:description}}
\sphinxAtStartPar
Perform constant NPT integration to update position, velocity, and
angular velocity each timestep for finite\sphinxhyphen{}sizex spherical particles in
the group using a Nose/Hoover temperature thermostat and Nose/Hoover
pressure barostat.  P is pressure; T is temperature.  This creates a
system trajectory consistent with the isothermal\sphinxhyphen{}isobaric ensemble.

\sphinxAtStartPar
This fix differs from the {\hyperref[\detokenize{fix_nh::doc}]{\sphinxcrossref{\DUrole{doc}{fix npt}}}} command, which
assumes point particles and only updates their position and velocity.

\sphinxAtStartPar
The thermostat is applied to both the translational and rotational
degrees of freedom for the spherical particles, assuming a compute is
used which calculates a temperature that includes the rotational
degrees of freedom (see below).  The translational degrees of freedom
can also have a bias velocity removed from them before thermostatting
takes place; see the description below.

\sphinxAtStartPar
If the \sphinxstyleemphasis{disc} keyword is used, then each particle is treated as a 2d
disc (circle) instead of as a sphere.  This is only possible for 2d
simulations, as defined by the {\hyperref[\detokenize{dimension::doc}]{\sphinxcrossref{\DUrole{doc}{dimension}}}} keyword.
The only difference between discs and spheres in this context is their
moment of inertia, as used in the time integration.

\sphinxAtStartPar
Additional parameters affecting the thermostat and barostat are
specified by keywords and values documented with the {\hyperref[\detokenize{fix_nh::doc}]{\sphinxcrossref{\DUrole{doc}{fix npt}}}} command.  See, for example, discussion of the \sphinxstyleemphasis{temp}, \sphinxstyleemphasis{iso},
\sphinxstyleemphasis{aniso}, and \sphinxstyleemphasis{dilate} keywords.

\sphinxAtStartPar
The particles in the fix group are the only ones whose velocities and
positions are updated by the velocity/position update portion of the
NPT integration.

\sphinxAtStartPar
Regardless of what particles are in the fix group, a global pressure is
computed for all particles.  Similarly, when the size of the simulation
box is changed, all particles are re\sphinxhyphen{}scaled to new positions, unless the
keyword \sphinxstyleemphasis{dilate} is specified with a value of \sphinxstyleemphasis{partial}, in which case
only the particles in the fix group are re\sphinxhyphen{}scaled.  The latter can be
useful for leaving the coordinates of particles in a solid substrate
unchanged and controlling the pressure of a surrounding fluid.


\bigskip\hrule\bigskip


\sphinxAtStartPar
This fix computes a temperature and pressure each timestep.  To do
this, the fix creates its own computes of style “temp/sphere” and
“pressure”, as if these commands had been issued:

\begin{sphinxVerbatim}[commandchars=\\\{\}]
\PYG{k}{compute }\PYG{n+nv+nvIdentifier}{fix\PYGZhy{}ID\PYGZus{}temp}\PYG{+w}{ }\PYG{n+nv+nvIdentifier}{all}\PYG{+w}{ }\PYG{n}{temp}\PYG{o}{/}\PYG{n}{sphere}
\PYG{k}{compute }\PYG{n+nv+nvIdentifier}{fix\PYGZhy{}ID\PYGZus{}press}\PYG{+w}{ }\PYG{n+nv+nvIdentifier}{all}\PYG{+w}{ }\PYG{n}{pressure}\PYG{+w}{ }\PYG{n}{fix}\PYG{o}{\PYGZhy{}}\PYG{n}{ID\PYGZus{}temp}
\end{sphinxVerbatim}

\sphinxAtStartPar
See the {\hyperref[\detokenize{compute_temp_sphere::doc}]{\sphinxcrossref{\DUrole{doc}{compute temp/sphere}}}} and
{\hyperref[\detokenize{compute_pressure::doc}]{\sphinxcrossref{\DUrole{doc}{compute pressure}}}} commands for details.  Note
that the IDs of the new computes are the fix\sphinxhyphen{}ID + underscore + “temp”
or fix\_ID
+ underscore + “press”, and the group for the new computes is “all”
since pressure is computed for the entire system.

\sphinxAtStartPar
Note that these are NOT the computes used by thermodynamic output (see
the {\hyperref[\detokenize{thermo_style::doc}]{\sphinxcrossref{\DUrole{doc}{thermo\_style}}}} command) with ID =
\sphinxstyleemphasis{thermo\_temp} and \sphinxstyleemphasis{thermo\_press}.  This means you can change the
attributes of this fix’s temperature or pressure via the
{\hyperref[\detokenize{compute_modify::doc}]{\sphinxcrossref{\DUrole{doc}{compute\_modify}}}} command or print this
temperature or pressure during thermodynamic output via the
{\hyperref[\detokenize{thermo_style::doc}]{\sphinxcrossref{\DUrole{doc}{thermo\_style custom}}}} command using the
appropriate compute\sphinxhyphen{}ID.  It also means that changing attributes of
\sphinxstyleemphasis{thermo\_temp} or \sphinxstyleemphasis{thermo\_press} will have no effect on this fix.

\sphinxAtStartPar
Like other fixes that perform thermostatting, this fix can be used
with {\hyperref[\detokenize{compute::doc}]{\sphinxcrossref{\DUrole{doc}{compute commands}}}} that remove a “bias” from the
atom velocities.  E.g. to apply the thermostat only to atoms within a
spatial {\hyperref[\detokenize{region::doc}]{\sphinxcrossref{\DUrole{doc}{region}}}}, or to remove the center\sphinxhyphen{}of\sphinxhyphen{}mass
velocity from a group of atoms, or to remove the x\sphinxhyphen{}component of
velocity from the calculation.

\sphinxAtStartPar
This is not done by default, but only if the {\hyperref[\detokenize{fix_modify::doc}]{\sphinxcrossref{\DUrole{doc}{fix\_modify}}}} command is used to assign a temperature compute to this
fix that includes such a bias term.  See the doc pages for individual
{\hyperref[\detokenize{compute::doc}]{\sphinxcrossref{\DUrole{doc}{compute temp commands}}}} to determine which ones include
a bias.  In this case, the thermostat works in the following manner:
bias is removed from each atom, thermostatting is performed on the
remaining thermal degrees of freedom, and the bias is added back in.


\bigskip\hrule\bigskip


\sphinxAtStartPar
Styles with a \sphinxstyleemphasis{gpu}, \sphinxstyleemphasis{intel}, \sphinxstyleemphasis{kk}, \sphinxstyleemphasis{omp}, or \sphinxstyleemphasis{opt} suffix are
functionally the same as the corresponding style without the suffix.
They have been optimized to run faster, depending on your available
hardware, as discussed on the {\hyperref[\detokenize{Speed_packages::doc}]{\sphinxcrossref{\DUrole{doc}{Accelerator packages}}}}
page.  The accelerated styles take the same arguments and should
produce the same results, except for round\sphinxhyphen{}off and precision issues.

\sphinxAtStartPar
These accelerated styles are part of the GPU, INTEL, KOKKOS,
OPENMP, and OPT packages, respectively.  They are only enabled if
LAMMPS was built with those packages.  See the {\hyperref[\detokenize{Build_package::doc}]{\sphinxcrossref{\DUrole{doc}{Build package}}}} page for more info.

\sphinxAtStartPar
You can specify the accelerated styles explicitly in your input script
by including their suffix, or you can use the {\hyperref[\detokenize{Run_options::doc}]{\sphinxcrossref{\DUrole{doc}{\sphinxhyphen{}suffix command\sphinxhyphen{}line switch}}}} when you invoke LAMMPS, or you can use the
{\hyperref[\detokenize{suffix::doc}]{\sphinxcrossref{\DUrole{doc}{suffix}}}} command in your input script.

\sphinxAtStartPar
See the {\hyperref[\detokenize{Speed_packages::doc}]{\sphinxcrossref{\DUrole{doc}{Accelerator packages}}}} page for more
instructions on how to use the accelerated styles effectively.


\subsection{Restart, fix\_modify, output, run start/stop, minimize info}
\label{\detokenize{fix_npt_sphere:restart-fix-modify-output-run-start-stop-minimize-info}}
\sphinxAtStartPar
This fix writes the state of the Nose/Hoover thermostat and barostat
to {\hyperref[\detokenize{restart::doc}]{\sphinxcrossref{\DUrole{doc}{binary restart files}}}}.  See the
{\hyperref[\detokenize{read_restart::doc}]{\sphinxcrossref{\DUrole{doc}{read\_restart}}}} command for info on how to re\sphinxhyphen{}specify
a fix in an input script that reads a restart file, so that the
operation of the fix continues in an uninterrupted fashion.

\sphinxAtStartPar
The {\hyperref[\detokenize{fix_modify::doc}]{\sphinxcrossref{\DUrole{doc}{fix\_modify}}}} \sphinxstyleemphasis{temp} and \sphinxstyleemphasis{press} options are
supported by this fix.  You can use them to assign a {\hyperref[\detokenize{compute::doc}]{\sphinxcrossref{\DUrole{doc}{compute}}}} you have defined to this fix which will be used in its
thermostatting or barostatting procedure.  If you do this, note that
the kinetic energy derived from the compute temperature should be
consistent with the virial term computed using all atoms for the
pressure.  LAMMPS will warn you if you choose to compute temperature
on a subset of atoms.

\sphinxAtStartPar
The cumulative energy change in the system imposed by this fix is
included in the {\hyperref[\detokenize{thermo_style::doc}]{\sphinxcrossref{\DUrole{doc}{thermodynamic output}}}} keywords
\sphinxstyleemphasis{ecouple} and \sphinxstyleemphasis{econserve}.  See the {\hyperref[\detokenize{thermo_style::doc}]{\sphinxcrossref{\DUrole{doc}{thermo\_style}}}}
doc page for details.

\sphinxAtStartPar
This fix computes the same global scalar and global vector of
quantities as does the {\hyperref[\detokenize{fix_nh::doc}]{\sphinxcrossref{\DUrole{doc}{fix npt}}}} command.

\sphinxAtStartPar
This fix can ramp its target temperature and pressure over multiple
runs, using the \sphinxstyleemphasis{start} and \sphinxstyleemphasis{stop} keywords of the {\hyperref[\detokenize{run::doc}]{\sphinxcrossref{\DUrole{doc}{run}}}}
command.  See the {\hyperref[\detokenize{run::doc}]{\sphinxcrossref{\DUrole{doc}{run}}}} command for details of how to do
this.

\sphinxAtStartPar
This fix is not invoked during {\hyperref[\detokenize{minimize::doc}]{\sphinxcrossref{\DUrole{doc}{energy minimization}}}}.


\subsection{Restrictions}
\label{\detokenize{fix_npt_sphere:restrictions}}
\sphinxAtStartPar
This fix requires that atoms store torque and angular velocity (omega)
and a radius as defined by the {\hyperref[\detokenize{atom_style::doc}]{\sphinxcrossref{\DUrole{doc}{atom\_style sphere}}}}
command.

\sphinxAtStartPar
All particles in the group must be finite\sphinxhyphen{}size spheres.  They cannot
be point particles.

\sphinxAtStartPar
Use of the \sphinxstyleemphasis{disc} keyword is only allowed for 2d simulations, as
defined by the {\hyperref[\detokenize{dimension::doc}]{\sphinxcrossref{\DUrole{doc}{dimension}}}} keyword.


\subsection{Related commands}
\label{\detokenize{fix_npt_sphere:related-commands}}\begin{description}
\sphinxlineitem{{\hyperref[\detokenize{fix_nh::doc}]{\sphinxcrossref{\DUrole{doc}{fix npt}}}}, {\hyperref[\detokenize{fix_nve_sphere::doc}]{\sphinxcrossref{\DUrole{doc}{fix nve\_sphere}}}},}
\sphinxAtStartPar
{\hyperref[\detokenize{fix_nvt_sphere::doc}]{\sphinxcrossref{\DUrole{doc}{fix nvt\_sphere}}}}, {\hyperref[\detokenize{fix_npt_asphere::doc}]{\sphinxcrossref{\DUrole{doc}{fix npt\_asphere}}}}, {\hyperref[\detokenize{fix_modify::doc}]{\sphinxcrossref{\DUrole{doc}{fix\_modify}}}}

\end{description}


\subsection{Default}
\label{\detokenize{fix_npt_sphere:default}}
\sphinxAtStartPar
none

\sphinxstepscope

\index{fix numdiff@\spxentry{fix numdiff}}\ignorespaces 

\section{fix numdiff command}
\label{\detokenize{fix_numdiff:fix-numdiff-command}}\label{\detokenize{fix_numdiff:index-0}}\label{\detokenize{fix_numdiff::doc}}

\subsection{Syntax}
\label{\detokenize{fix_numdiff:syntax}}
\begin{sphinxVerbatim}[commandchars=\\\{\}]
\PYG{k}{fix }\PYG{n+nv+nvIdentifier}{ID}\PYG{+w}{ }\PYG{n+nv+nvIdentifier}{group\PYGZhy{}ID}\PYG{+w}{ }\PYG{n}{numdiff}\PYG{+w}{ }\PYG{n}{Nevery}\PYG{+w}{ }\PYG{n}{delta}
\end{sphinxVerbatim}
\begin{itemize}
\item {} 
\sphinxAtStartPar
ID, group\sphinxhyphen{}ID are documented in {\hyperref[\detokenize{fix::doc}]{\sphinxcrossref{\DUrole{doc}{fix}}}} command

\item {} 
\sphinxAtStartPar
numdiff = style name of this fix command

\item {} 
\sphinxAtStartPar
Nevery = calculate force by finite difference every this many timesteps

\item {} 
\sphinxAtStartPar
delta = size of atom displacements (distance units)

\end{itemize}


\subsection{Examples}
\label{\detokenize{fix_numdiff:examples}}
\begin{sphinxVerbatim}[commandchars=\\\{\}]
\PYG{k}{fix }\PYG{n+nv+nvIdentifier}{1}\PYG{+w}{ }\PYG{n+nv+nvIdentifier}{all}\PYG{+w}{ }\PYG{n}{numdiff}\PYG{+w}{ }\PYG{l+m}{10}\PYG{+w}{ }\PYG{l+m}{1e\PYGZhy{}6}
\PYG{k}{fix }\PYG{n+nv+nvIdentifier}{1}\PYG{+w}{ }\PYG{n+nv+nvIdentifier}{movegroup}\PYG{+w}{ }\PYG{n}{numdiff}\PYG{+w}{ }\PYG{l+m}{100}\PYG{+w}{ }\PYG{l+m}{0.01}
\end{sphinxVerbatim}


\subsection{Description}
\label{\detokenize{fix_numdiff:description}}
\sphinxAtStartPar
Calculate forces through finite difference calculations of energy
versus position.  These forces can be compared to analytic forces
computed by pair styles, bond styles, etc.  This can be useful for
debugging or other purposes.

\sphinxAtStartPar
The group specified with the command means only atoms within the group
have their averages computed.  Results are set to 0.0 for atoms not in
the group.

\sphinxAtStartPar
This fix performs a loop over all atoms in the group.  For each atom
and each component of force it adds \sphinxstyleemphasis{delta} to the position, and
computes the new energy of the entire system.  It then subtracts
\sphinxstyleemphasis{delta} from the original position and again computes the new energy
of the system.  It then restores the original position.  That
component of force is calculated as the difference in energy divided
by two times \sphinxstyleemphasis{delta}.

\begin{sphinxadmonition}{note}{Note:}
\sphinxAtStartPar
It is important to choose a suitable value for delta, the magnitude of
atom displacements that are used to generate finite difference
approximations to the exact forces.  For typical systems, a value in
the range of 1 part in 1e4 to 1e5 of the typical separation distance
between atoms in the liquid or solid state will be sufficient.
However, the best value will depend on a multitude of factors
including the stiffness of the interatomic potential, the thermodynamic
state of the material being probed, and so on. The only way to be sure
that you have made a good choice is to do a sensitivity study on a
representative atomic configuration, sweeping over a wide range of
values of delta.  If delta is too small, the output forces will vary
erratically due to truncation effects. If delta is increased beyond a
certain point, the output forces will start to vary smoothly with
delta, due to growing contributions from higher order derivatives. In
between these two limits, the numerical force values should be largely
independent of delta.
\end{sphinxadmonition}

\begin{sphinxadmonition}{note}{Note:}
\sphinxAtStartPar
The cost of each energy evaluation is essentially the cost of an MD
timestep.  Thus invoking this fix once for a 3d system has a cost
of 6N timesteps, where N is the total number of atoms in the system.
So this fix can be very expensive to use for large systems.
One expedient alternative is to define the fix for a group containing
only a few atoms.
\end{sphinxadmonition}


\bigskip\hrule\bigskip


\sphinxAtStartPar
The \sphinxstyleemphasis{Nevery} argument specifies on what timesteps the force will
be used calculated by finite difference.

\sphinxAtStartPar
The \sphinxstyleemphasis{delta} argument specifies the size of the displacement each
atom will undergo.


\bigskip\hrule\bigskip



\subsection{Restart, fix\_modify, output, run start/stop, minimize info}
\label{\detokenize{fix_numdiff:restart-fix-modify-output-run-start-stop-minimize-info}}
\sphinxAtStartPar
No information about this fix is written to {\hyperref[\detokenize{restart::doc}]{\sphinxcrossref{\DUrole{doc}{binary restart files}}}}.  None of the {\hyperref[\detokenize{fix_modify::doc}]{\sphinxcrossref{\DUrole{doc}{fix\_modify}}}} options are
relevant to this fix.

\sphinxAtStartPar
This fix produces a per\sphinxhyphen{}atom array which can be accessed by various
{\hyperref[\detokenize{Howto_output::doc}]{\sphinxcrossref{\DUrole{doc}{output commands}}}}, which stores the components of
the force on each atom as calculated by finite difference.  The
per\sphinxhyphen{}atom values can only be accessed on timesteps that are multiples
of \sphinxstyleemphasis{Nevery} since that is when the finite difference forces are
calculated. See the examples in \sphinxstyleemphasis{examples/numdiff} directory
to see how this fix can be used to directly compare with
the analytic forces computed by LAMMPS.

\sphinxAtStartPar
The array values calculated by this compute
will be in force {\hyperref[\detokenize{units::doc}]{\sphinxcrossref{\DUrole{doc}{units}}}}.

\sphinxAtStartPar
No parameter of this fix can be used with the \sphinxstyleemphasis{start/stop} keywords of
the {\hyperref[\detokenize{run::doc}]{\sphinxcrossref{\DUrole{doc}{run}}}} command.  This fix is invoked during {\hyperref[\detokenize{minimize::doc}]{\sphinxcrossref{\DUrole{doc}{energy
minimization}}}}.


\subsection{Restrictions}
\label{\detokenize{fix_numdiff:restrictions}}
\sphinxAtStartPar
This fix is part of the EXTRA\sphinxhyphen{}FIX package.  It is only enabled if LAMMPS
was built with that package.  See the {\hyperref[\detokenize{Build_package::doc}]{\sphinxcrossref{\DUrole{doc}{Build package}}}} page for more info.


\subsection{Related commands}
\label{\detokenize{fix_numdiff:related-commands}}
\sphinxAtStartPar
{\hyperref[\detokenize{dynamical_matrix::doc}]{\sphinxcrossref{\DUrole{doc}{dynamical\_matrix}}}}, {\hyperref[\detokenize{fix_numdiff_virial::doc}]{\sphinxcrossref{\DUrole{doc}{fix numdiff/virial}}}},


\subsection{Default}
\label{\detokenize{fix_numdiff:default}}
\sphinxAtStartPar
none

\sphinxstepscope

\index{fix numdiff/virial@\spxentry{fix numdiff/virial}}\ignorespaces 

\section{fix numdiff/virial command}
\label{\detokenize{fix_numdiff_virial:fix-numdiff-virial-command}}\label{\detokenize{fix_numdiff_virial:index-0}}\label{\detokenize{fix_numdiff_virial::doc}}

\subsection{Syntax}
\label{\detokenize{fix_numdiff_virial:syntax}}
\begin{sphinxVerbatim}[commandchars=\\\{\}]
\PYG{k}{fix }\PYG{n+nv+nvIdentifier}{ID}\PYG{+w}{ }\PYG{n+nv+nvIdentifier}{group\PYGZhy{}ID}\PYG{+w}{ }\PYG{n}{numdiff}\PYG{o}{/}\PYG{n}{virial}\PYG{+w}{ }\PYG{n}{Nevery}\PYG{+w}{ }\PYG{n}{delta}
\end{sphinxVerbatim}
\begin{itemize}
\item {} 
\sphinxAtStartPar
ID, group\sphinxhyphen{}ID are documented in {\hyperref[\detokenize{fix::doc}]{\sphinxcrossref{\DUrole{doc}{fix}}}} command

\item {} 
\sphinxAtStartPar
numdiff/virial = style name of this fix command

\item {} 
\sphinxAtStartPar
Nevery = calculate virial by finite difference every this many timesteps

\item {} 
\sphinxAtStartPar
delta = magnitude of strain fields (dimensionless)

\end{itemize}


\subsection{Examples}
\label{\detokenize{fix_numdiff_virial:examples}}
\begin{sphinxVerbatim}[commandchars=\\\{\}]
\PYG{k}{fix }\PYG{n+nv+nvIdentifier}{1}\PYG{+w}{ }\PYG{n+nv+nvIdentifier}{all}\PYG{+w}{ }\PYG{n}{numdiff}\PYG{o}{/}\PYG{n}{stress}\PYG{+w}{ }\PYG{l+m}{10}\PYG{+w}{ }\PYG{l+m}{1e\PYGZhy{}6}
\end{sphinxVerbatim}


\subsection{Description}
\label{\detokenize{fix_numdiff_virial:description}}
\sphinxAtStartPar
\DUrole{versionmodified,added}{Added in version 17Feb2022.}

\sphinxAtStartPar
Calculate the virial stress tensor through a finite difference calculation of
energy versus strain.  These values can be compared to the analytic virial
tensor computed by pair styles, bond styles, etc.  This can be useful for
debugging or other purposes. The specified group must be “all”.

\sphinxAtStartPar
This fix applies linear strain fields of magnitude \sphinxstyleemphasis{delta} to all the
atoms relative to a point at the center of the box.  The
strain fields are in six different directions, corresponding to the
six Cartesian components of the stress tensor defined by LAMMPS.
For each direction it applies the strain field in both the positive
and negative senses, and the new energy of the entire system
is calculated after each. The difference in these two energies
divided by two times \sphinxstyleemphasis{delta}, approximates the corresponding
component of the virial stress tensor, after applying
a suitable unit conversion.

\begin{sphinxadmonition}{note}{Note:}
\sphinxAtStartPar
It is important to choose a suitable value for delta, the magnitude of
strains that are used to generate finite difference
approximations to the exact virial stress.  For typical systems, a value in
the range of 1 part in 1e5 to 1e6 will be sufficient.
However, the best value will depend on a multitude of factors
including the stiffness of the interatomic potential, the thermodynamic
state of the material being probed, and so on. The only way to be sure
that you have made a good choice is to do a sensitivity study on a
representative atomic configuration, sweeping over a wide range of
values of delta.  If delta is too small, the output values will vary
erratically due to truncation effects. If delta is increased beyond a
certain point, the output values will start to vary smoothly with
delta, due to growing contributions from higher order derivatives. In
between these two limits, the numerical virial values should be largely
independent of delta.
\end{sphinxadmonition}


\bigskip\hrule\bigskip


\sphinxAtStartPar
The \sphinxstyleemphasis{Nevery} argument specifies on what timesteps the force will
be used calculated by finite difference.

\sphinxAtStartPar
The \sphinxstyleemphasis{delta} argument specifies the size of the displacement each
atom will undergo.


\bigskip\hrule\bigskip



\subsection{Restart, fix\_modify, output, run start/stop, minimize info}
\label{\detokenize{fix_numdiff_virial:restart-fix-modify-output-run-start-stop-minimize-info}}
\sphinxAtStartPar
No information about this fix is written to {\hyperref[\detokenize{restart::doc}]{\sphinxcrossref{\DUrole{doc}{binary restart files}}}}.  None of the {\hyperref[\detokenize{fix_modify::doc}]{\sphinxcrossref{\DUrole{doc}{fix\_modify}}}} options are
relevant to this fix.

\sphinxAtStartPar
This fix produces a global vector which can be accessed by various
{\hyperref[\detokenize{Howto_output::doc}]{\sphinxcrossref{\DUrole{doc}{output commands}}}}, which stores the components of
the virial stress tensor as calculated by finite difference.  The
global vector can only be accessed on timesteps that are multiples
of \sphinxstyleemphasis{Nevery} since that is when the finite difference virial is
calculated. See the examples in \sphinxstyleemphasis{examples/numdiff} directory
to see how this fix can be used to directly compare with
the analytic virial stress tensor computed by LAMMPS.

\sphinxAtStartPar
The order of the virial stress tensor components is \sphinxstyleemphasis{xx}, \sphinxstyleemphasis{yy}, \sphinxstyleemphasis{zz},
\sphinxstyleemphasis{yz}, \sphinxstyleemphasis{xz}, and \sphinxstyleemphasis{xy}, consistent with Voigt notation. Note that
the vector produced by {\hyperref[\detokenize{compute_pressure::doc}]{\sphinxcrossref{\DUrole{doc}{compute pressure}}}}
uses a different ordering, with \sphinxstyleemphasis{yz} and \sphinxstyleemphasis{xy} swapped.

\sphinxAtStartPar
The vector values calculated by this compute are
“intensive”.  The vector values will be in pressure
{\hyperref[\detokenize{units::doc}]{\sphinxcrossref{\DUrole{doc}{units}}}}.

\sphinxAtStartPar
No parameter of this fix can be used with the \sphinxstyleemphasis{start/stop} keywords of
the {\hyperref[\detokenize{run::doc}]{\sphinxcrossref{\DUrole{doc}{run}}}} command.  This fix is invoked during {\hyperref[\detokenize{minimize::doc}]{\sphinxcrossref{\DUrole{doc}{energy
minimization}}}}.


\subsection{Restrictions}
\label{\detokenize{fix_numdiff_virial:restrictions}}
\sphinxAtStartPar
This fix is part of the EXTRA\sphinxhyphen{}FIX package.  It is only enabled if LAMMPS
was built with that package.  See the {\hyperref[\detokenize{Build_package::doc}]{\sphinxcrossref{\DUrole{doc}{Build package}}}} page for more info.


\subsection{Related commands}
\label{\detokenize{fix_numdiff_virial:related-commands}}
\sphinxAtStartPar
{\hyperref[\detokenize{fix_numdiff::doc}]{\sphinxcrossref{\DUrole{doc}{fix numdiff}}}}, {\hyperref[\detokenize{compute_pressure::doc}]{\sphinxcrossref{\DUrole{doc}{compute pressure}}}}


\subsection{Default}
\label{\detokenize{fix_numdiff_virial:default}}
\sphinxAtStartPar
none

\sphinxstepscope

\index{fix nve@\spxentry{fix nve}}\ignorespaces 
\index{fix nve/gpu@\spxentry{fix nve/gpu}}\ignorespaces 
\index{fix nve/intel@\spxentry{fix nve/intel}}\ignorespaces 
\index{fix nve/kk@\spxentry{fix nve/kk}}\ignorespaces 
\index{fix nve/omp@\spxentry{fix nve/omp}}\ignorespaces 

\section{fix nve command}
\label{\detokenize{fix_nve:fix-nve-command}}\label{\detokenize{fix_nve:index-4}}\label{\detokenize{fix_nve:index-3}}\label{\detokenize{fix_nve:index-2}}\label{\detokenize{fix_nve:index-1}}\label{\detokenize{fix_nve:index-0}}\label{\detokenize{fix_nve::doc}}
\sphinxAtStartPar
Accelerator Variants: \sphinxstyleemphasis{nve/gpu}, \sphinxstyleemphasis{nve/intel}, \sphinxstyleemphasis{nve/kk}, \sphinxstyleemphasis{nve/omp}


\subsection{Syntax}
\label{\detokenize{fix_nve:syntax}}
\begin{sphinxVerbatim}[commandchars=\\\{\}]
\PYG{k}{fix }\PYG{n+nv+nvIdentifier}{ID}\PYG{+w}{ }\PYG{n+nv+nvIdentifier}{group\PYGZhy{}ID}\PYG{+w}{ }\PYG{n}{nve}
\end{sphinxVerbatim}
\begin{itemize}
\item {} 
\sphinxAtStartPar
ID, group\sphinxhyphen{}ID are documented in {\hyperref[\detokenize{fix::doc}]{\sphinxcrossref{\DUrole{doc}{fix}}}} command

\item {} 
\sphinxAtStartPar
nve = style name of this fix command

\end{itemize}


\subsection{Examples}
\label{\detokenize{fix_nve:examples}}
\begin{sphinxVerbatim}[commandchars=\\\{\}]
\PYG{k}{fix }\PYG{n+nv+nvIdentifier}{1}\PYG{+w}{ }\PYG{n+nv+nvIdentifier}{all}\PYG{+w}{ }\PYG{n}{nve}
\end{sphinxVerbatim}


\subsection{Description}
\label{\detokenize{fix_nve:description}}
\sphinxAtStartPar
Perform plain time integration to update position and velocity for
atoms in the group each timestep.  This creates a system trajectory
consistent with the microcanonical ensemble (NVE) provided there
are (full) periodic boundary conditions and no other “manipulations”
of the system (e.g. fixes that modify forces or velocities).

\sphinxAtStartPar
This fix invokes the velocity form of the
Stoermer\sphinxhyphen{}Verlet time integration algorithm (velocity\sphinxhyphen{}Verlet). Other
time integration options can be invoked using the {\hyperref[\detokenize{run_style::doc}]{\sphinxcrossref{\DUrole{doc}{run\_style}}}} command.


\bigskip\hrule\bigskip


\sphinxAtStartPar
Styles with a \sphinxstyleemphasis{gpu}, \sphinxstyleemphasis{intel}, \sphinxstyleemphasis{kk}, \sphinxstyleemphasis{omp}, or \sphinxstyleemphasis{opt} suffix are
functionally the same as the corresponding style without the suffix.
They have been optimized to run faster, depending on your available
hardware, as discussed on the {\hyperref[\detokenize{Speed_packages::doc}]{\sphinxcrossref{\DUrole{doc}{Accelerator packages}}}}
page.  The accelerated styles take the same arguments and should
produce the same results, except for round\sphinxhyphen{}off and precision issues.

\sphinxAtStartPar
These accelerated styles are part of the GPU, INTEL, KOKKOS,
OPENMP, and OPT packages, respectively.  They are only enabled if
LAMMPS was built with those packages.  See the {\hyperref[\detokenize{Build_package::doc}]{\sphinxcrossref{\DUrole{doc}{Build package}}}} page for more info.

\sphinxAtStartPar
You can specify the accelerated styles explicitly in your input script
by including their suffix, or you can use the {\hyperref[\detokenize{Run_options::doc}]{\sphinxcrossref{\DUrole{doc}{\sphinxhyphen{}suffix command\sphinxhyphen{}line switch}}}} when you invoke LAMMPS, or you can use the
{\hyperref[\detokenize{suffix::doc}]{\sphinxcrossref{\DUrole{doc}{suffix}}}} command in your input script.

\sphinxAtStartPar
See the {\hyperref[\detokenize{Speed_packages::doc}]{\sphinxcrossref{\DUrole{doc}{Accelerator packages}}}} page for more
instructions on how to use the accelerated styles effectively.


\bigskip\hrule\bigskip



\subsection{Restart, fix\_modify, output, run start/stop, minimize info}
\label{\detokenize{fix_nve:restart-fix-modify-output-run-start-stop-minimize-info}}
\sphinxAtStartPar
No information about this fix is written to {\hyperref[\detokenize{restart::doc}]{\sphinxcrossref{\DUrole{doc}{binary restart files}}}}.  None of the {\hyperref[\detokenize{fix_modify::doc}]{\sphinxcrossref{\DUrole{doc}{fix\_modify}}}} options
are relevant to this fix.  No global or per\sphinxhyphen{}atom quantities are stored
by this fix for access by various {\hyperref[\detokenize{Howto_output::doc}]{\sphinxcrossref{\DUrole{doc}{output commands}}}}.
No parameter of this fix can be used with the \sphinxstyleemphasis{start/stop} keywords of
the {\hyperref[\detokenize{run::doc}]{\sphinxcrossref{\DUrole{doc}{run}}}} command.  This fix is not invoked during {\hyperref[\detokenize{minimize::doc}]{\sphinxcrossref{\DUrole{doc}{energy minimization}}}}.


\subsection{Restrictions}
\label{\detokenize{fix_nve:restrictions}}\begin{quote}

\sphinxAtStartPar
none
\end{quote}


\subsection{Related commands}
\label{\detokenize{fix_nve:related-commands}}
\sphinxAtStartPar
{\hyperref[\detokenize{fix_nh::doc}]{\sphinxcrossref{\DUrole{doc}{fix nvt}}}}, {\hyperref[\detokenize{fix_nh::doc}]{\sphinxcrossref{\DUrole{doc}{fix npt}}}}, {\hyperref[\detokenize{run_style::doc}]{\sphinxcrossref{\DUrole{doc}{run\_style}}}}


\subsection{Default}
\label{\detokenize{fix_nve:default}}
\sphinxAtStartPar
none

\sphinxstepscope

\index{fix nve/asphere@\spxentry{fix nve/asphere}}\ignorespaces 
\index{fix nve/asphere/gpu@\spxentry{fix nve/asphere/gpu}}\ignorespaces 
\index{fix nve/asphere/intel@\spxentry{fix nve/asphere/intel}}\ignorespaces 

\section{fix nve/asphere command}
\label{\detokenize{fix_nve_asphere:fix-nve-asphere-command}}\label{\detokenize{fix_nve_asphere:index-2}}\label{\detokenize{fix_nve_asphere:index-1}}\label{\detokenize{fix_nve_asphere:index-0}}\label{\detokenize{fix_nve_asphere::doc}}
\sphinxAtStartPar
Accelerator Variants: \sphinxstyleemphasis{nve/asphere/gpu}, \sphinxstyleemphasis{nve/asphere/intel}


\subsection{Syntax}
\label{\detokenize{fix_nve_asphere:syntax}}
\begin{sphinxVerbatim}[commandchars=\\\{\}]
\PYG{k}{fix }\PYG{n+nv+nvIdentifier}{ID}\PYG{+w}{ }\PYG{n+nv+nvIdentifier}{group\PYGZhy{}ID}\PYG{+w}{ }\PYG{n}{nve}\PYG{o}{/}\PYG{n}{asphere}
\end{sphinxVerbatim}
\begin{itemize}
\item {} 
\sphinxAtStartPar
ID, group\sphinxhyphen{}ID are documented in {\hyperref[\detokenize{fix::doc}]{\sphinxcrossref{\DUrole{doc}{fix}}}} command

\item {} 
\sphinxAtStartPar
nve/asphere = style name of this fix command

\end{itemize}


\subsection{Examples}
\label{\detokenize{fix_nve_asphere:examples}}
\begin{sphinxVerbatim}[commandchars=\\\{\}]
\PYG{k}{fix }\PYG{n+nv+nvIdentifier}{1}\PYG{+w}{ }\PYG{n+nv+nvIdentifier}{all}\PYG{+w}{ }\PYG{n}{nve}\PYG{o}{/}\PYG{n}{asphere}
\end{sphinxVerbatim}


\subsection{Description}
\label{\detokenize{fix_nve_asphere:description}}
\sphinxAtStartPar
Perform constant NVE integration to update position, velocity,
orientation, and angular velocity for aspherical particles in the
group each timestep.  V is volume; E is energy.  This creates a system
trajectory consistent with the microcanonical ensemble.

\sphinxAtStartPar
This fix differs from the {\hyperref[\detokenize{fix_nve::doc}]{\sphinxcrossref{\DUrole{doc}{fix nve}}}} command, which
assumes point particles and only updates their position and velocity.


\subsection{Restart, fix\_modify, output, run start/stop, minimize info}
\label{\detokenize{fix_nve_asphere:restart-fix-modify-output-run-start-stop-minimize-info}}
\sphinxAtStartPar
No information about this fix is written to {\hyperref[\detokenize{restart::doc}]{\sphinxcrossref{\DUrole{doc}{binary restart files}}}}.  None of the {\hyperref[\detokenize{fix_modify::doc}]{\sphinxcrossref{\DUrole{doc}{fix\_modify}}}} options
are relevant to this fix.  No global or per\sphinxhyphen{}atom quantities are stored
by this fix for access by various {\hyperref[\detokenize{Howto_output::doc}]{\sphinxcrossref{\DUrole{doc}{output commands}}}}.
No parameter of this fix can be used with the \sphinxstyleemphasis{start/stop} keywords of
the {\hyperref[\detokenize{run::doc}]{\sphinxcrossref{\DUrole{doc}{run}}}} command.  This fix is not invoked during {\hyperref[\detokenize{minimize::doc}]{\sphinxcrossref{\DUrole{doc}{energy minimization}}}}.


\bigskip\hrule\bigskip


\sphinxAtStartPar
Styles with a \sphinxstyleemphasis{gpu}, \sphinxstyleemphasis{intel}, \sphinxstyleemphasis{kk}, \sphinxstyleemphasis{omp}, or \sphinxstyleemphasis{opt} suffix are
functionally the same as the corresponding style without the suffix.
They have been optimized to run faster, depending on your available
hardware, as discussed on the {\hyperref[\detokenize{Speed_packages::doc}]{\sphinxcrossref{\DUrole{doc}{Accelerator packages}}}}
page.  The accelerated styles take the same arguments and should
produce the same results, except for round\sphinxhyphen{}off and precision issues.

\sphinxAtStartPar
These accelerated styles are part of the GPU, INTEL, KOKKOS,
OPENMP, and OPT packages, respectively.  They are only enabled if
LAMMPS was built with those packages.  See the {\hyperref[\detokenize{Build_package::doc}]{\sphinxcrossref{\DUrole{doc}{Build package}}}} page for more info.

\sphinxAtStartPar
You can specify the accelerated styles explicitly in your input script
by including their suffix, or you can use the {\hyperref[\detokenize{Run_options::doc}]{\sphinxcrossref{\DUrole{doc}{\sphinxhyphen{}suffix command\sphinxhyphen{}line switch}}}} when you invoke LAMMPS, or you can use the
{\hyperref[\detokenize{suffix::doc}]{\sphinxcrossref{\DUrole{doc}{suffix}}}} command in your input script.

\sphinxAtStartPar
See the {\hyperref[\detokenize{Speed_packages::doc}]{\sphinxcrossref{\DUrole{doc}{Accelerator packages}}}} page for more
instructions on how to use the accelerated styles effectively.


\bigskip\hrule\bigskip



\subsection{Restrictions}
\label{\detokenize{fix_nve_asphere:restrictions}}
\sphinxAtStartPar
This fix is part of the ASPHERE package.  It is only enabled if LAMMPS
was built with that package.  See the {\hyperref[\detokenize{Build_package::doc}]{\sphinxcrossref{\DUrole{doc}{Build package}}}} page for more info.

\sphinxAtStartPar
This fix requires that atoms store torque and angular momentum and a
quaternion as defined by the {\hyperref[\detokenize{atom_style::doc}]{\sphinxcrossref{\DUrole{doc}{atom\_style ellipsoid}}}}
command.

\sphinxAtStartPar
All particles in the group must be finite\sphinxhyphen{}size.  They cannot be point
particles, but they can be aspherical or spherical as defined by their
shape attribute.


\subsection{Related commands}
\label{\detokenize{fix_nve_asphere:related-commands}}
\sphinxAtStartPar
{\hyperref[\detokenize{fix_nve::doc}]{\sphinxcrossref{\DUrole{doc}{fix nve}}}}, {\hyperref[\detokenize{fix_nve_sphere::doc}]{\sphinxcrossref{\DUrole{doc}{fix nve/sphere}}}}


\subsection{Default}
\label{\detokenize{fix_nve_asphere:default}}
\sphinxAtStartPar
none

\sphinxstepscope

\index{fix nve/asphere/noforce@\spxentry{fix nve/asphere/noforce}}\ignorespaces 

\section{fix nve/asphere/noforce command}
\label{\detokenize{fix_nve_asphere_noforce:fix-nve-asphere-noforce-command}}\label{\detokenize{fix_nve_asphere_noforce:index-0}}\label{\detokenize{fix_nve_asphere_noforce::doc}}

\subsection{Syntax}
\label{\detokenize{fix_nve_asphere_noforce:syntax}}
\begin{sphinxVerbatim}[commandchars=\\\{\}]
\PYG{k}{fix }\PYG{n+nv+nvIdentifier}{ID}\PYG{+w}{ }\PYG{n+nv+nvIdentifier}{group\PYGZhy{}ID}\PYG{+w}{ }\PYG{n}{nve}\PYG{o}{/}\PYG{n}{asphere}\PYG{o}{/}\PYG{n}{noforce}
\end{sphinxVerbatim}
\begin{itemize}
\item {} 
\sphinxAtStartPar
ID, group\sphinxhyphen{}ID are documented in {\hyperref[\detokenize{fix::doc}]{\sphinxcrossref{\DUrole{doc}{fix}}}} command

\item {} 
\sphinxAtStartPar
nve/asphere/noforce = style name of this fix command

\end{itemize}


\subsection{Examples}
\label{\detokenize{fix_nve_asphere_noforce:examples}}
\begin{sphinxVerbatim}[commandchars=\\\{\}]
\PYG{k}{fix }\PYG{n+nv+nvIdentifier}{1}\PYG{+w}{ }\PYG{n+nv+nvIdentifier}{all}\PYG{+w}{ }\PYG{n}{nve}\PYG{o}{/}\PYG{n}{asphere}\PYG{o}{/}\PYG{n}{noforce}
\end{sphinxVerbatim}


\subsection{Description}
\label{\detokenize{fix_nve_asphere_noforce:description}}
\sphinxAtStartPar
Perform updates of position and orientation, but not velocity or
angular momentum for atoms in the group each timestep.  In other
words, the force and torque on the atoms is ignored and their velocity
and angular momentum are not updated.  The atom velocities and
angular momenta are used to update their positions and orientation.

\sphinxAtStartPar
This is useful as an implicit time integrator for Fast Lubrication
Dynamics, since the velocity and angular momentum are updated by the
{\hyperref[\detokenize{pair_lubricateU::doc}]{\sphinxcrossref{\DUrole{doc}{pair\_style lubricuteU}}}} command.


\bigskip\hrule\bigskip



\subsection{Restart, fix\_modify, output, run start/stop, minimize info}
\label{\detokenize{fix_nve_asphere_noforce:restart-fix-modify-output-run-start-stop-minimize-info}}
\sphinxAtStartPar
No information about this fix is written to {\hyperref[\detokenize{restart::doc}]{\sphinxcrossref{\DUrole{doc}{binary restart files}}}}.  None of the {\hyperref[\detokenize{fix_modify::doc}]{\sphinxcrossref{\DUrole{doc}{fix\_modify}}}} options
are relevant to this fix.  No global or per\sphinxhyphen{}atom quantities are stored
by this fix for access by various {\hyperref[\detokenize{Howto_output::doc}]{\sphinxcrossref{\DUrole{doc}{output commands}}}}.
No parameter of this fix can be used with the \sphinxstyleemphasis{start/stop} keywords of
the {\hyperref[\detokenize{run::doc}]{\sphinxcrossref{\DUrole{doc}{run}}}} command.  This fix is not invoked during {\hyperref[\detokenize{minimize::doc}]{\sphinxcrossref{\DUrole{doc}{energy minimization}}}}.


\subsection{Restrictions}
\label{\detokenize{fix_nve_asphere_noforce:restrictions}}
\sphinxAtStartPar
This fix is part of the ASPHERE package.  It is only enabled if LAMMPS
was built with that package.  See the {\hyperref[\detokenize{Build_package::doc}]{\sphinxcrossref{\DUrole{doc}{Build package}}}} page for more info.

\sphinxAtStartPar
This fix requires that atoms store torque and angular momentum and a
quaternion as defined by the {\hyperref[\detokenize{atom_style::doc}]{\sphinxcrossref{\DUrole{doc}{atom\_style ellipsoid}}}}
command.

\sphinxAtStartPar
All particles in the group must be finite\sphinxhyphen{}size.  They cannot be point
particles, but they can be aspherical or spherical as defined by their
shape attribute.


\subsection{Related commands}
\label{\detokenize{fix_nve_asphere_noforce:related-commands}}
\sphinxAtStartPar
{\hyperref[\detokenize{fix_nve_noforce::doc}]{\sphinxcrossref{\DUrole{doc}{fix nve/noforce}}}}, {\hyperref[\detokenize{fix_nve_asphere::doc}]{\sphinxcrossref{\DUrole{doc}{fix nve/asphere}}}}


\subsection{Default}
\label{\detokenize{fix_nve_asphere_noforce:default}}
\sphinxAtStartPar
none

\sphinxstepscope

\index{fix nve/awpmd@\spxentry{fix nve/awpmd}}\ignorespaces 

\section{fix nve/awpmd command}
\label{\detokenize{fix_nve_awpmd:fix-nve-awpmd-command}}\label{\detokenize{fix_nve_awpmd:index-0}}\label{\detokenize{fix_nve_awpmd::doc}}

\subsection{Syntax}
\label{\detokenize{fix_nve_awpmd:syntax}}
\begin{sphinxVerbatim}[commandchars=\\\{\}]
\PYG{k}{fix }\PYG{n+nv+nvIdentifier}{ID}\PYG{+w}{ }\PYG{n+nv+nvIdentifier}{group\PYGZhy{}ID}\PYG{+w}{ }\PYG{n}{nve}\PYG{o}{/}\PYG{n}{awpmd}
\end{sphinxVerbatim}
\begin{itemize}
\item {} 
\sphinxAtStartPar
ID, group\sphinxhyphen{}ID are documented in {\hyperref[\detokenize{fix::doc}]{\sphinxcrossref{\DUrole{doc}{fix}}}} command

\item {} 
\sphinxAtStartPar
nve/awpmd = style name of this fix command

\end{itemize}


\subsection{Examples}
\label{\detokenize{fix_nve_awpmd:examples}}
\begin{sphinxVerbatim}[commandchars=\\\{\}]
\PYG{k}{fix }\PYG{n+nv+nvIdentifier}{1}\PYG{+w}{ }\PYG{n+nv+nvIdentifier}{all}\PYG{+w}{ }\PYG{n}{nve}\PYG{o}{/}\PYG{n}{awpmd}
\end{sphinxVerbatim}


\subsection{Description}
\label{\detokenize{fix_nve_awpmd:description}}
\sphinxAtStartPar
Perform constant NVE integration to update position and velocity for
nuclei and electrons in the group for the {\hyperref[\detokenize{pair_awpmd::doc}]{\sphinxcrossref{\DUrole{doc}{Antisymmetrized Wave Packet Molecular Dynamics}}}} model.  V is volume; E is energy.
This creates a system trajectory consistent with the microcanonical
ensemble.

\sphinxAtStartPar
The operation of this fix is exactly like that described by the {\hyperref[\detokenize{fix_nve::doc}]{\sphinxcrossref{\DUrole{doc}{fix nve}}}} command, except that the width and width\sphinxhyphen{}velocity of
the electron wave functions are also updated.


\bigskip\hrule\bigskip



\subsection{Restart, fix\_modify, output, run start/stop, minimize info}
\label{\detokenize{fix_nve_awpmd:restart-fix-modify-output-run-start-stop-minimize-info}}
\sphinxAtStartPar
No information about this fix is written to {\hyperref[\detokenize{restart::doc}]{\sphinxcrossref{\DUrole{doc}{binary restart files}}}}.  None of the {\hyperref[\detokenize{fix_modify::doc}]{\sphinxcrossref{\DUrole{doc}{fix\_modify}}}} options
are relevant to this fix.  No global or per\sphinxhyphen{}atom quantities are stored
by this fix for access by various {\hyperref[\detokenize{Howto_output::doc}]{\sphinxcrossref{\DUrole{doc}{output commands}}}}.
No parameter of this fix can be used with the \sphinxstyleemphasis{start/stop} keywords of
the {\hyperref[\detokenize{run::doc}]{\sphinxcrossref{\DUrole{doc}{run}}}} command.  This fix is not invoked during {\hyperref[\detokenize{minimize::doc}]{\sphinxcrossref{\DUrole{doc}{energy minimization}}}}.


\subsection{Restrictions}
\label{\detokenize{fix_nve_awpmd:restrictions}}
\sphinxAtStartPar
This fix is part of the AWPMD package.  It is only enabled if
LAMMPS was built with that package.  See the {\hyperref[\detokenize{Build_package::doc}]{\sphinxcrossref{\DUrole{doc}{Build package}}}} page for more info.


\subsection{Related commands}
\label{\detokenize{fix_nve_awpmd:related-commands}}
\sphinxAtStartPar
{\hyperref[\detokenize{fix_nve::doc}]{\sphinxcrossref{\DUrole{doc}{fix nve}}}}


\subsection{Default}
\label{\detokenize{fix_nve_awpmd:default}}
\sphinxAtStartPar
none

\sphinxstepscope

\index{fix nve/body@\spxentry{fix nve/body}}\ignorespaces 

\section{fix nve/body command}
\label{\detokenize{fix_nve_body:fix-nve-body-command}}\label{\detokenize{fix_nve_body:index-0}}\label{\detokenize{fix_nve_body::doc}}

\subsection{Syntax}
\label{\detokenize{fix_nve_body:syntax}}
\begin{sphinxVerbatim}[commandchars=\\\{\}]
\PYG{k}{fix }\PYG{n+nv+nvIdentifier}{ID}\PYG{+w}{ }\PYG{n+nv+nvIdentifier}{group\PYGZhy{}ID}\PYG{+w}{ }\PYG{n}{nve}\PYG{o}{/}\PYG{n}{body}
\end{sphinxVerbatim}
\begin{itemize}
\item {} 
\sphinxAtStartPar
ID, group\sphinxhyphen{}ID are documented in {\hyperref[\detokenize{fix::doc}]{\sphinxcrossref{\DUrole{doc}{fix}}}} command

\item {} 
\sphinxAtStartPar
nve/body = style name of this fix command

\end{itemize}


\subsection{Examples}
\label{\detokenize{fix_nve_body:examples}}
\begin{sphinxVerbatim}[commandchars=\\\{\}]
\PYG{k}{fix }\PYG{n+nv+nvIdentifier}{1}\PYG{+w}{ }\PYG{n+nv+nvIdentifier}{all}\PYG{+w}{ }\PYG{n}{nve}\PYG{o}{/}\PYG{n}{body}
\end{sphinxVerbatim}


\subsection{Description}
\label{\detokenize{fix_nve_body:description}}
\sphinxAtStartPar
Perform constant NVE integration to update position, velocity,
orientation, and angular velocity for body particles in the group each
timestep.  V is volume; E is energy.  This creates a system trajectory
consistent with the microcanonical ensemble.  See the {\hyperref[\detokenize{Howto_body::doc}]{\sphinxcrossref{\DUrole{doc}{Howto body}}}} page for more details on using body
particles.

\sphinxAtStartPar
This fix differs from the {\hyperref[\detokenize{fix_nve::doc}]{\sphinxcrossref{\DUrole{doc}{fix nve}}}} command, which
assumes point particles and only updates their position and velocity.


\subsection{Restart, fix\_modify, output, run start/stop, minimize info}
\label{\detokenize{fix_nve_body:restart-fix-modify-output-run-start-stop-minimize-info}}
\sphinxAtStartPar
No information about this fix is written to {\hyperref[\detokenize{restart::doc}]{\sphinxcrossref{\DUrole{doc}{binary restart files}}}}.  None of the {\hyperref[\detokenize{fix_modify::doc}]{\sphinxcrossref{\DUrole{doc}{fix\_modify}}}} options
are relevant to this fix.  No global or per\sphinxhyphen{}atom quantities are stored
by this fix for access by various {\hyperref[\detokenize{Howto_output::doc}]{\sphinxcrossref{\DUrole{doc}{output commands}}}}.
No parameter of this fix can be used with the \sphinxstyleemphasis{start/stop} keywords of
the {\hyperref[\detokenize{run::doc}]{\sphinxcrossref{\DUrole{doc}{run}}}} command.  This fix is not invoked during {\hyperref[\detokenize{minimize::doc}]{\sphinxcrossref{\DUrole{doc}{energy minimization}}}}.


\subsection{Restrictions}
\label{\detokenize{fix_nve_body:restrictions}}
\sphinxAtStartPar
This fix is part of the BODY package.  It is only enabled if LAMMPS
was built with that package.  See the {\hyperref[\detokenize{Build_package::doc}]{\sphinxcrossref{\DUrole{doc}{Build package}}}} page for more info.

\sphinxAtStartPar
This fix requires that atoms store torque and angular momentum and a
quaternion as defined by the {\hyperref[\detokenize{atom_style::doc}]{\sphinxcrossref{\DUrole{doc}{atom\_style body}}}}
command.

\sphinxAtStartPar
All particles in the group must be body particles.  They cannot be
point particles.


\subsection{Related commands}
\label{\detokenize{fix_nve_body:related-commands}}
\sphinxAtStartPar
{\hyperref[\detokenize{fix_nve::doc}]{\sphinxcrossref{\DUrole{doc}{fix nve}}}}, {\hyperref[\detokenize{fix_nve_sphere::doc}]{\sphinxcrossref{\DUrole{doc}{fix nve/sphere}}}}, {\hyperref[\detokenize{fix_nve_asphere::doc}]{\sphinxcrossref{\DUrole{doc}{fix nve/asphere}}}}


\subsection{Default}
\label{\detokenize{fix_nve_body:default}}
\sphinxAtStartPar
none

\sphinxstepscope

\index{fix nve/bpm/sphere@\spxentry{fix nve/bpm/sphere}}\ignorespaces 

\section{fix nve/bpm/sphere command}
\label{\detokenize{fix_nve_bpm_sphere:fix-nve-bpm-sphere-command}}\label{\detokenize{fix_nve_bpm_sphere:index-0}}\label{\detokenize{fix_nve_bpm_sphere::doc}}

\subsection{Syntax}
\label{\detokenize{fix_nve_bpm_sphere:syntax}}
\begin{sphinxVerbatim}[commandchars=\\\{\}]
\PYG{k}{fix }\PYG{n+nv+nvIdentifier}{ID}\PYG{+w}{ }\PYG{n+nv+nvIdentifier}{group\PYGZhy{}ID}\PYG{+w}{ }\PYG{n}{nve}\PYG{o}{/}\PYG{n}{bpm}\PYG{o}{/}\PYG{n}{sphere}
\end{sphinxVerbatim}
\begin{itemize}
\item {} 
\sphinxAtStartPar
ID, group\sphinxhyphen{}ID are documented in {\hyperref[\detokenize{fix::doc}]{\sphinxcrossref{\DUrole{doc}{fix}}}} command

\item {} 
\sphinxAtStartPar
nve/bpm/sphere = style name of this fix command

\item {} 
\sphinxAtStartPar
zero or more keyword/value pairs may be appended

\item {} 
\sphinxAtStartPar
keyword = \sphinxstyleemphasis{disc}
\begin{sphinxalltt}
\sphinxstyleemphasis{disc} value = none = treat particles as 2d discs, not spheres
\end{sphinxalltt}

\end{itemize}


\subsection{Examples}
\label{\detokenize{fix_nve_bpm_sphere:examples}}
\begin{sphinxVerbatim}[commandchars=\\\{\}]
\PYG{k}{fix }\PYG{n+nv+nvIdentifier}{1}\PYG{+w}{ }\PYG{n+nv+nvIdentifier}{all}\PYG{+w}{ }\PYG{n}{nve}\PYG{o}{/}\PYG{n}{bpm}\PYG{o}{/}\PYG{n}{sphere}
\PYG{k}{fix }\PYG{n+nv+nvIdentifier}{1}\PYG{+w}{ }\PYG{n+nv+nvIdentifier}{all}\PYG{+w}{ }\PYG{n}{nve}\PYG{o}{/}\PYG{n}{bpm}\PYG{o}{/}\PYG{n}{sphere}\PYG{+w}{ }\PYG{n}{disc}
\end{sphinxVerbatim}


\subsection{Description}
\label{\detokenize{fix_nve_bpm_sphere:description}}
\sphinxAtStartPar
\DUrole{versionmodified,added}{Added in version 4May2022.}

\sphinxAtStartPar
Perform constant NVE integration to update position, velocity, angular
velocity, and quaternion orientation for finite\sphinxhyphen{}size spherical
particles in the group each timestep.  V is volume; E is energy.  This
creates a system trajectory consistent with the microcanonical
ensemble.

\sphinxAtStartPar
This fix differs from the {\hyperref[\detokenize{fix_nve::doc}]{\sphinxcrossref{\DUrole{doc}{fix nve}}}} command, which
assumes point particles and only updates their position and velocity.
It also differs from the {\hyperref[\detokenize{fix_nve_sphere::doc}]{\sphinxcrossref{\DUrole{doc}{fix nve/sphere}}}}
command which assumes finite\sphinxhyphen{}size spheroid particles which do not
store a quaternion.  It thus does not update a particle’s orientation
or quaternion.

\sphinxAtStartPar
If the \sphinxstyleemphasis{disc} keyword is used, then each particle is treated as a 2d
disc (circle) instead of as a sphere.  This is only possible for 2d
simulations, as defined by the {\hyperref[\detokenize{dimension::doc}]{\sphinxcrossref{\DUrole{doc}{dimension}}}} keyword.
The only difference between discs and spheres in this context is their
moment of inertia, as used in the time integration.


\bigskip\hrule\bigskip



\subsection{Restart, fix\_modify, output, run start/stop, minimize info}
\label{\detokenize{fix_nve_bpm_sphere:restart-fix-modify-output-run-start-stop-minimize-info}}
\sphinxAtStartPar
No information about this fix is written to {\hyperref[\detokenize{restart::doc}]{\sphinxcrossref{\DUrole{doc}{binary restart files}}}}.  None of the {\hyperref[\detokenize{fix_modify::doc}]{\sphinxcrossref{\DUrole{doc}{fix\_modify}}}} options are
relevant to this fix.  No global or per\sphinxhyphen{}atom quantities are stored by
this fix for access by various {\hyperref[\detokenize{Howto_output::doc}]{\sphinxcrossref{\DUrole{doc}{output commands}}}}.
No parameter of this fix can be used with the \sphinxstyleemphasis{start/stop} keywords of
the {\hyperref[\detokenize{run::doc}]{\sphinxcrossref{\DUrole{doc}{run}}}} command.  This fix is not invoked during
{\hyperref[\detokenize{minimize::doc}]{\sphinxcrossref{\DUrole{doc}{energy minimization}}}}.


\subsection{Restrictions}
\label{\detokenize{fix_nve_bpm_sphere:restrictions}}
\sphinxAtStartPar
This fix is part of the BPM package.  It is only enabled if LAMMPS was
built with that package.  See the {\hyperref[\detokenize{Build_package::doc}]{\sphinxcrossref{\DUrole{doc}{Build package}}}}
page for more info.

\sphinxAtStartPar
This fix requires that atoms store torque, angular velocity (omega), a
radius, and a quaternion as defined by the {\hyperref[\detokenize{atom_style::doc}]{\sphinxcrossref{\DUrole{doc}{atom\_style bpm/sphere}}}} command.

\sphinxAtStartPar
All particles in the group must be finite\sphinxhyphen{}size spheres with
quaternions.  They cannot be point particles.

\sphinxAtStartPar
Use of the \sphinxstyleemphasis{disc} keyword is only allowed for 2d simulations, as
defined by the {\hyperref[\detokenize{dimension::doc}]{\sphinxcrossref{\DUrole{doc}{dimension}}}} keyword.


\subsection{Related commands}
\label{\detokenize{fix_nve_bpm_sphere:related-commands}}
\sphinxAtStartPar
{\hyperref[\detokenize{fix_nve::doc}]{\sphinxcrossref{\DUrole{doc}{fix nve}}}}, {\hyperref[\detokenize{fix_nve_sphere::doc}]{\sphinxcrossref{\DUrole{doc}{fix nve/sphere}}}}


\subsection{Default}
\label{\detokenize{fix_nve_bpm_sphere:default}}
\sphinxAtStartPar
none

\sphinxstepscope

\index{fix nve/dot@\spxentry{fix nve/dot}}\ignorespaces 

\section{fix nve/dot command}
\label{\detokenize{fix_nve_dot:fix-nve-dot-command}}\label{\detokenize{fix_nve_dot:index-0}}\label{\detokenize{fix_nve_dot::doc}}

\subsection{Syntax}
\label{\detokenize{fix_nve_dot:syntax}}
\begin{sphinxVerbatim}[commandchars=\\\{\}]
\PYG{k}{fix }\PYG{n+nv+nvIdentifier}{ID}\PYG{+w}{ }\PYG{n+nv+nvIdentifier}{group\PYGZhy{}ID}\PYG{+w}{ }\PYG{n}{nve}\PYG{o}{/}\PYG{n}{dot}
\end{sphinxVerbatim}
\begin{itemize}
\item {} 
\sphinxAtStartPar
ID, group\sphinxhyphen{}ID are documented in {\hyperref[\detokenize{fix::doc}]{\sphinxcrossref{\DUrole{doc}{fix}}}} command

\item {} 
\sphinxAtStartPar
nve/dot = style name of this fix command

\end{itemize}


\subsection{Examples}
\label{\detokenize{fix_nve_dot:examples}}
\begin{sphinxVerbatim}[commandchars=\\\{\}]
\PYG{k}{fix }\PYG{n+nv+nvIdentifier}{1}\PYG{+w}{ }\PYG{n+nv+nvIdentifier}{all}\PYG{+w}{ }\PYG{n}{nve}\PYG{o}{/}\PYG{n}{dot}
\end{sphinxVerbatim}


\subsection{Description}
\label{\detokenize{fix_nve_dot:description}}
\sphinxAtStartPar
Apply a rigid\sphinxhyphen{}body integrator as described in {\hyperref[\detokenize{fix_nve_dot:davidchack4}]{\sphinxcrossref{\DUrole{std,std-ref}{(Davidchack)}}}}
to a group of atoms, but without Langevin dynamics.
This command performs Molecular dynamics (MD)
via a velocity\sphinxhyphen{}Verlet algorithm and an evolution operator that rotates
the quaternion degrees of freedom, similar to the scheme outlined in {\hyperref[\detokenize{fix_nve_dot:miller4}]{\sphinxcrossref{\DUrole{std,std-ref}{(Miller)}}}}.

\sphinxAtStartPar
This command is the equivalent of the {\hyperref[\detokenize{fix_nve_dotc_langevin::doc}]{\sphinxcrossref{\DUrole{doc}{fix nve/dotc/langevin}}}}
without damping and noise and can be used to determine the stability range
in a NVE ensemble prior to using the Langevin\sphinxhyphen{}type DOTC\sphinxhyphen{}integrator
(see also {\hyperref[\detokenize{fix_nve_dotc_langevin::doc}]{\sphinxcrossref{\DUrole{doc}{fix nve/dotc/langevin}}}}).
The command is equivalent to the {\hyperref[\detokenize{fix_nve::doc}]{\sphinxcrossref{\DUrole{doc}{fix nve}}}}.
The particles are always considered to have a finite size.

\sphinxAtStartPar
An example input file can be found in /examples/PACKAGES/cgdna/examples/duplex1/.
Further details of the implementation and stability of the integrator are contained in {\hyperref[\detokenize{fix_nve_dot:henrich4}]{\sphinxcrossref{\DUrole{std,std-ref}{(Henrich)}}}}.
The preprint version of the article can be found \sphinxhref{PDF/CG-DNA.pdf}{here}.


\bigskip\hrule\bigskip



\subsection{Restrictions}
\label{\detokenize{fix_nve_dot:restrictions}}
\sphinxAtStartPar
These pair styles can only be used if LAMMPS was built with the
{\hyperref[\detokenize{Packages_details:pkg-cg-dna}]{\sphinxcrossref{\DUrole{std,std-ref}{CG\sphinxhyphen{}DNA}}}} package and the MOLECULE and ASPHERE package.
See the {\hyperref[\detokenize{Build_package::doc}]{\sphinxcrossref{\DUrole{doc}{Build package}}}} page for more info.


\subsection{Related commands}
\label{\detokenize{fix_nve_dot:related-commands}}
\sphinxAtStartPar
{\hyperref[\detokenize{fix_nve_dotc_langevin::doc}]{\sphinxcrossref{\DUrole{doc}{fix nve/dotc/langevin}}}}, {\hyperref[\detokenize{fix_nve::doc}]{\sphinxcrossref{\DUrole{doc}{fix nve}}}}


\subsection{Default}
\label{\detokenize{fix_nve_dot:default}}
\sphinxAtStartPar
none


\bigskip\hrule\bigskip

\phantomsection\label{\detokenize{fix_nve_dot:davidchack4}}
\sphinxAtStartPar
\sphinxstylestrong{(Davidchack)} R.L Davidchack, T.E. Ouldridge, and M.V. Tretyakov. J. Chem. Phys. 142, 144114 (2015).

\phantomsection\label{\detokenize{fix_nve_dot:miller4}}
\sphinxAtStartPar
\sphinxstylestrong{(Miller)} T. F. Miller III, M. Eleftheriou, P. Pattnaik, A. Ndirango, G. J. Martyna, J. Chem. Phys., 116, 8649\sphinxhyphen{}8659 (2002).

\phantomsection\label{\detokenize{fix_nve_dot:henrich4}}
\sphinxAtStartPar
\sphinxstylestrong{(Henrich)} O. Henrich, Y. A. Gutierrez\sphinxhyphen{}Fosado, T. Curk, T. E. Ouldridge, Eur. Phys. J. E 41, 57 (2018).

\sphinxstepscope

\index{fix nve/dotc/langevin@\spxentry{fix nve/dotc/langevin}}\ignorespaces 

\section{fix nve/dotc/langevin command}
\label{\detokenize{fix_nve_dotc_langevin:fix-nve-dotc-langevin-command}}\label{\detokenize{fix_nve_dotc_langevin:index-0}}\label{\detokenize{fix_nve_dotc_langevin::doc}}

\subsection{Syntax}
\label{\detokenize{fix_nve_dotc_langevin:syntax}}
\begin{sphinxVerbatim}[commandchars=\\\{\}]
\PYG{k}{fix }\PYG{n+nv+nvIdentifier}{ID}\PYG{+w}{ }\PYG{n+nv+nvIdentifier}{group\PYGZhy{}ID}\PYG{+w}{ }\PYG{n}{nve}\PYG{o}{/}\PYG{n}{dotc}\PYG{o}{/}\PYG{n}{langevin}\PYG{+w}{ }\PYG{n}{Tstart}\PYG{+w}{ }\PYG{n}{Tstop}\PYG{+w}{ }\PYG{n}{damp}\PYG{+w}{ }\PYG{n}{seed}\PYG{+w}{ }\PYG{n}{keyword}\PYG{+w}{ }\PYG{n}{value}
\end{sphinxVerbatim}
\begin{itemize}
\item {} 
\sphinxAtStartPar
ID, group\sphinxhyphen{}ID are documented in {\hyperref[\detokenize{fix::doc}]{\sphinxcrossref{\DUrole{doc}{fix}}}} command

\item {} 
\sphinxAtStartPar
nve/dotc/langevin = style name of this fix command

\item {} 
\sphinxAtStartPar
Tstart,Tstop = desired temperature at start/end of run (temperature units)

\item {} 
\sphinxAtStartPar
damp = damping parameter (time units)

\item {} 
\sphinxAtStartPar
seed = random number seed to use for white noise (positive integer)

\item {} 
\sphinxAtStartPar
keyword = \sphinxstyleemphasis{angmom}
\begin{sphinxalltt}
\sphinxstyleemphasis{angmom} value = factor
  factor = do thermostat rotational degrees of freedom via the angular momentum and apply numeric scale factor as discussed below
\end{sphinxalltt}

\end{itemize}


\subsection{Examples}
\label{\detokenize{fix_nve_dotc_langevin:examples}}
\begin{sphinxVerbatim}[commandchars=\\\{\}]
\PYG{k}{fix }\PYG{n+nv+nvIdentifier}{1}\PYG{+w}{ }\PYG{n+nv+nvIdentifier}{all}\PYG{+w}{ }\PYG{n}{nve}\PYG{o}{/}\PYG{n}{dotc}\PYG{o}{/}\PYG{n}{langevin}\PYG{+w}{ }\PYG{l+m}{1.0}\PYG{+w}{ }\PYG{l+m}{1.0}\PYG{+w}{ }\PYG{l+m}{0.03}\PYG{+w}{ }\PYG{l+m}{457145}\PYG{+w}{ }\PYG{n}{angmom}\PYG{+w}{ }\PYG{l+m}{10}
\PYG{k}{fix }\PYG{n+nv+nvIdentifier}{1}\PYG{+w}{ }\PYG{n+nv+nvIdentifier}{all}\PYG{+w}{ }\PYG{n}{nve}\PYG{o}{/}\PYG{n}{dotc}\PYG{o}{/}\PYG{n}{langevin}\PYG{+w}{ }\PYG{l+m}{0.1}\PYG{+w}{ }\PYG{l+m}{0.1}\PYG{+w}{ }\PYG{l+m}{78.9375}\PYG{+w}{ }\PYG{l+m}{457145}\PYG{+w}{ }\PYG{n}{angmom}\PYG{+w}{ }\PYG{l+m}{10}
\end{sphinxVerbatim}


\subsection{Description}
\label{\detokenize{fix_nve_dotc_langevin:description}}
\sphinxAtStartPar
Apply a rigid\sphinxhyphen{}body Langevin\sphinxhyphen{}type integrator of the kind “Langevin C”
as described in {\hyperref[\detokenize{fix_nve_dotc_langevin:davidchack5}]{\sphinxcrossref{\DUrole{std,std-ref}{(Davidchack)}}}}
to a group of atoms, which models an interaction with an implicit background
solvent.  This command performs Brownian dynamics (BD)
via a technique that splits the integration into a deterministic Hamiltonian
part and the Ornstein\sphinxhyphen{}Uhlenbeck process for noise and damping.
The quaternion degrees of freedom are updated though an evolution
operator which performs a rotation in quaternion space, preserves
the quaternion norm and is akin to {\hyperref[\detokenize{fix_nve_dotc_langevin:miller5}]{\sphinxcrossref{\DUrole{std,std-ref}{(Miller)}}}}.

\sphinxAtStartPar
In terms of syntax this command has been closely modelled on the
{\hyperref[\detokenize{fix_langevin::doc}]{\sphinxcrossref{\DUrole{doc}{fix langevin}}}} and its \sphinxstyleemphasis{angmom} option. But it combines
the {\hyperref[\detokenize{fix_nve::doc}]{\sphinxcrossref{\DUrole{doc}{fix nve}}}} and the {\hyperref[\detokenize{fix_langevin::doc}]{\sphinxcrossref{\DUrole{doc}{fix langevin}}}} in
one single command. The main feature is improved stability
over the standard integrator, permitting slightly larger timestep sizes.

\begin{sphinxadmonition}{note}{Note:}
\sphinxAtStartPar
Unlike the {\hyperref[\detokenize{fix_langevin::doc}]{\sphinxcrossref{\DUrole{doc}{fix langevin}}}} this command performs
also time integration of the translational and quaternion degrees of freedom.
\end{sphinxadmonition}

\sphinxAtStartPar
The total force on each atom will have the form:
\begin{equation*}
\begin{split}F =   & F_c + F_f + F_r \\
F_f = & - \frac{m}{\mathrm{damp}} v \\
F_r \propto & \sqrt{\frac{k_B T m}{dt~\mathrm{damp}}}\end{split}
\end{equation*}
\sphinxAtStartPar
\(F_c\) is the conservative force computed via the usual
inter\sphinxhyphen{}particle interactions ({\hyperref[\detokenize{pair_style::doc}]{\sphinxcrossref{\DUrole{doc}{pair\_style}}}},
{\hyperref[\detokenize{bond_style::doc}]{\sphinxcrossref{\DUrole{doc}{bond\_style}}}}, etc). The \(F_f\) and \(F_r\)
terms are implicitly taken into account by this fix on a per\sphinxhyphen{}particle
basis.

\sphinxAtStartPar
\(F_f\) is a frictional drag or viscous damping term proportional to
the particle’s velocity.  The proportionality constant for each atom is
computed as \(\frac{m}{\mathrm{damp}}\), where \sphinxstyleemphasis{m} is the mass of
the particle and damp is the damping factor specified by the user.

\sphinxAtStartPar
\(F_r\) is a force due to solvent atoms at a temperature \(T\)
randomly bumping into the particle.  As derived from the
fluctuation/dissipation theorem, its magnitude as shown above is
proportional to \(\sqrt{\frac{k_B T m}{dt~\mathrm{damp}}}\), where
\(k_B\) is the Boltzmann constant, \(T\) is the desired temperature,
\sphinxstyleemphasis{m} is the mass of the particle, \sphinxstyleemphasis{dt} is the timestep size, and damp is
the damping factor.  Random numbers are used to randomize the direction
and magnitude of this force as described in {\hyperref[\detokenize{fix_nve_dotc_langevin:dunweg5}]{\sphinxcrossref{\DUrole{std,std-ref}{(Dunweg)}}}},
where a uniform random number is used (instead of a Gaussian random
number) for speed.


\bigskip\hrule\bigskip


\sphinxAtStartPar
\sphinxstyleemphasis{Tstart} and \sphinxstyleemphasis{Tstop} have to be constant values, i.e. they cannot
be variables. If used together with the oxDNA force field for
coarse\sphinxhyphen{}grained simulation of DNA please note that T = 0.1 in oxDNA units
corresponds to T = 300 K.

\sphinxAtStartPar
The \sphinxstyleemphasis{damp} parameter is specified in time units and determines how
rapidly the temperature is relaxed.  For example, a value of 0.03
means to relax the temperature in a timespan of (roughly) 0.03 time
units \(\tau\) (see the {\hyperref[\detokenize{units::doc}]{\sphinxcrossref{\DUrole{doc}{units}}}} command).
The damp factor can be thought of as inversely related to the
viscosity of the solvent, i.e. a small relaxation time implies a
high\sphinxhyphen{}viscosity solvent and vice versa.  See the discussion about gamma
and viscosity in the documentation for the {\hyperref[\detokenize{fix_viscous::doc}]{\sphinxcrossref{\DUrole{doc}{fix viscous}}}} command for more details.
Note that the value 78.9375 in the second example above corresponds
to a diffusion constant, which is about an order of magnitude larger
than realistic ones. This has been used to sample configurations faster
in Brownian dynamics simulations.

\sphinxAtStartPar
The random \# \sphinxstyleemphasis{seed} must be a positive integer. A Marsaglia random
number generator is used.  Each processor uses the input seed to
generate its own unique seed and its own stream of random numbers.
Thus the dynamics of the system will not be identical on two runs on
different numbers of processors.

\sphinxAtStartPar
The keyword/value option has to be used in the following way:

\sphinxAtStartPar
This fix has to be used together with the \sphinxstyleemphasis{angmom} keyword. The
particles are always considered to have a finite size.
The keyword \sphinxstyleemphasis{angmom} enables thermostatting of the rotational degrees of
freedom in addition to the usual translational degrees of freedom.

\sphinxAtStartPar
The scale factor after the \sphinxstyleemphasis{angmom} keyword gives the ratio of the
rotational to the translational friction coefficient.

\sphinxAtStartPar
An example input file can be found in examples/PACKAGES/cgdna/examples/duplex2/.
Further details of the implementation and stability of the integrators are contained in {\hyperref[\detokenize{fix_nve_dotc_langevin:henrich5}]{\sphinxcrossref{\DUrole{std,std-ref}{(Henrich)}}}}.
The preprint version of the article can be found \sphinxhref{PDF/CG-DNA.pdf}{here}.


\bigskip\hrule\bigskip



\subsection{Restrictions}
\label{\detokenize{fix_nve_dotc_langevin:restrictions}}
\sphinxAtStartPar
These pair styles can only be used if LAMMPS was built with the
{\hyperref[\detokenize{Packages_details:pkg-cg-dna}]{\sphinxcrossref{\DUrole{std,std-ref}{CG\sphinxhyphen{}DNA}}}} package and the MOLECULE and ASPHERE package.
See the {\hyperref[\detokenize{Build_package::doc}]{\sphinxcrossref{\DUrole{doc}{Build package}}}} page for more info.


\subsection{Related commands}
\label{\detokenize{fix_nve_dotc_langevin:related-commands}}
\sphinxAtStartPar
{\hyperref[\detokenize{fix_nve::doc}]{\sphinxcrossref{\DUrole{doc}{fix nve}}}}, {\hyperref[\detokenize{fix_langevin::doc}]{\sphinxcrossref{\DUrole{doc}{fix langevin}}}}, {\hyperref[\detokenize{fix_nve_dot::doc}]{\sphinxcrossref{\DUrole{doc}{fix nve/dot}}}}, {\hyperref[\detokenize{bond_oxdna::doc}]{\sphinxcrossref{\DUrole{doc}{bond\_style oxdna/fene}}}}, {\hyperref[\detokenize{bond_oxdna::doc}]{\sphinxcrossref{\DUrole{doc}{bond\_style oxdna2/fene}}}}, {\hyperref[\detokenize{pair_oxdna::doc}]{\sphinxcrossref{\DUrole{doc}{pair\_style oxdna/excv}}}}, {\hyperref[\detokenize{pair_oxdna2::doc}]{\sphinxcrossref{\DUrole{doc}{pair\_style oxdna2/excv}}}}


\subsection{Default}
\label{\detokenize{fix_nve_dotc_langevin:default}}
\sphinxAtStartPar
none


\bigskip\hrule\bigskip

\phantomsection\label{\detokenize{fix_nve_dotc_langevin:davidchack5}}
\sphinxAtStartPar
\sphinxstylestrong{(Davidchack)} R.L Davidchack, T.E. Ouldridge, M.V. Tretyakov. J. Chem. Phys. 142, 144114 (2015).

\phantomsection\label{\detokenize{fix_nve_dotc_langevin:miller5}}
\sphinxAtStartPar
\sphinxstylestrong{(Miller)} T. F. Miller III, M. Eleftheriou, P. Pattnaik, A. Ndirango, G. J. Martyna, J. Chem. Phys., 116, 8649\sphinxhyphen{}8659 (2002).

\phantomsection\label{\detokenize{fix_nve_dotc_langevin:dunweg5}}
\sphinxAtStartPar
\sphinxstylestrong{(Dunweg)} B. Dunweg, W. Paul, Int. J. Mod. Phys. C, 2, 817\sphinxhyphen{}27 (1991).

\phantomsection\label{\detokenize{fix_nve_dotc_langevin:henrich5}}
\sphinxAtStartPar
\sphinxstylestrong{(Henrich)} O. Henrich, Y. A. Gutierrez\sphinxhyphen{}Fosado, T. Curk, T. E. Ouldridge, Eur. Phys. J. E 41, 57 (2018).

\sphinxstepscope

\index{fix nve/eff@\spxentry{fix nve/eff}}\ignorespaces 

\section{fix nve/eff command}
\label{\detokenize{fix_nve_eff:fix-nve-eff-command}}\label{\detokenize{fix_nve_eff:index-0}}\label{\detokenize{fix_nve_eff::doc}}

\subsection{Syntax}
\label{\detokenize{fix_nve_eff:syntax}}
\begin{sphinxVerbatim}[commandchars=\\\{\}]
\PYG{k}{fix }\PYG{n+nv+nvIdentifier}{ID}\PYG{+w}{ }\PYG{n+nv+nvIdentifier}{group\PYGZhy{}ID}\PYG{+w}{ }\PYG{n}{nve}\PYG{o}{/}\PYG{n}{eff}
\end{sphinxVerbatim}
\begin{itemize}
\item {} 
\sphinxAtStartPar
ID, group\sphinxhyphen{}ID are documented in {\hyperref[\detokenize{fix::doc}]{\sphinxcrossref{\DUrole{doc}{fix}}}} command

\item {} 
\sphinxAtStartPar
nve/eff = style name of this fix command

\end{itemize}


\subsection{Examples}
\label{\detokenize{fix_nve_eff:examples}}
\begin{sphinxVerbatim}[commandchars=\\\{\}]
\PYG{k}{fix }\PYG{n+nv+nvIdentifier}{1}\PYG{+w}{ }\PYG{n+nv+nvIdentifier}{all}\PYG{+w}{ }\PYG{n}{nve}\PYG{o}{/}\PYG{n}{eff}
\end{sphinxVerbatim}


\subsection{Description}
\label{\detokenize{fix_nve_eff:description}}
\sphinxAtStartPar
Perform constant NVE integration to update position and velocity for
nuclei and electrons in the group for the {\hyperref[\detokenize{pair_eff::doc}]{\sphinxcrossref{\DUrole{doc}{electron force field}}}} model.  V is volume; E is energy.  This creates a
system trajectory consistent with the microcanonical ensemble.

\sphinxAtStartPar
The operation of this fix is exactly like that described by the {\hyperref[\detokenize{fix_nve::doc}]{\sphinxcrossref{\DUrole{doc}{fix nve}}}} command, except that the radius and radial velocity
of electrons are also updated.


\subsection{Restart, fix\_modify, output, run start/stop, minimize info}
\label{\detokenize{fix_nve_eff:restart-fix-modify-output-run-start-stop-minimize-info}}
\sphinxAtStartPar
No information about this fix is written to {\hyperref[\detokenize{restart::doc}]{\sphinxcrossref{\DUrole{doc}{binary restart files}}}}.  None of the {\hyperref[\detokenize{fix_modify::doc}]{\sphinxcrossref{\DUrole{doc}{fix\_modify}}}} options
are relevant to this fix.  No global or per\sphinxhyphen{}atom quantities are stored
by this fix for access by various {\hyperref[\detokenize{Howto_output::doc}]{\sphinxcrossref{\DUrole{doc}{output commands}}}}.
No parameter of this fix can be used with the \sphinxstyleemphasis{start/stop} keywords of
the {\hyperref[\detokenize{run::doc}]{\sphinxcrossref{\DUrole{doc}{run}}}} command.  This fix is not invoked during {\hyperref[\detokenize{minimize::doc}]{\sphinxcrossref{\DUrole{doc}{energy minimization}}}}.


\subsection{Restrictions}
\label{\detokenize{fix_nve_eff:restrictions}}
\sphinxAtStartPar
This fix is part of the EFF package.  It is only enabled if
LAMMPS was built with that package.  See the {\hyperref[\detokenize{Build_package::doc}]{\sphinxcrossref{\DUrole{doc}{Build package}}}} page for more info.


\subsection{Related commands}
\label{\detokenize{fix_nve_eff:related-commands}}
\sphinxAtStartPar
{\hyperref[\detokenize{fix_nve::doc}]{\sphinxcrossref{\DUrole{doc}{fix nve}}}}, {\hyperref[\detokenize{fix_nh_eff::doc}]{\sphinxcrossref{\DUrole{doc}{fix nvt/eff}}}}, {\hyperref[\detokenize{fix_nh_eff::doc}]{\sphinxcrossref{\DUrole{doc}{fix npt/eff}}}}


\subsection{Default}
\label{\detokenize{fix_nve_eff:default}}
\sphinxAtStartPar
none

\sphinxstepscope

\index{fix nve/limit@\spxentry{fix nve/limit}}\ignorespaces 

\section{fix nve/limit command}
\label{\detokenize{fix_nve_limit:fix-nve-limit-command}}\label{\detokenize{fix_nve_limit:index-0}}\label{\detokenize{fix_nve_limit::doc}}

\subsection{Syntax}
\label{\detokenize{fix_nve_limit:syntax}}
\begin{sphinxVerbatim}[commandchars=\\\{\}]
\PYG{k}{fix }\PYG{n+nv+nvIdentifier}{ID}\PYG{+w}{ }\PYG{n+nv+nvIdentifier}{group\PYGZhy{}ID}\PYG{+w}{ }\PYG{n}{nve}\PYG{o}{/}\PYG{n}{limit}\PYG{+w}{ }\PYG{n}{xmax}
\end{sphinxVerbatim}
\begin{itemize}
\item {} 
\sphinxAtStartPar
ID, group\sphinxhyphen{}ID are documented in {\hyperref[\detokenize{fix::doc}]{\sphinxcrossref{\DUrole{doc}{fix}}}} command

\item {} 
\sphinxAtStartPar
nve = style name of this fix command

\item {} 
\sphinxAtStartPar
xmax = maximum distance an atom can move in one timestep (distance units)

\end{itemize}


\subsection{Examples}
\label{\detokenize{fix_nve_limit:examples}}
\begin{sphinxVerbatim}[commandchars=\\\{\}]
\PYG{k}{fix }\PYG{n+nv+nvIdentifier}{1}\PYG{+w}{ }\PYG{n+nv+nvIdentifier}{all}\PYG{+w}{ }\PYG{n}{nve}\PYG{o}{/}\PYG{n}{limit}\PYG{+w}{ }\PYG{l+m}{0.1}
\end{sphinxVerbatim}


\subsection{Description}
\label{\detokenize{fix_nve_limit:description}}
\sphinxAtStartPar
Perform constant NVE updates of position and velocity for atoms in the
group each timestep.  A limit is imposed on the maximum distance an
atom can move in one timestep.  This is useful when starting a
simulation with a configuration containing highly overlapped atoms.
Normally this would generate huge forces which would blow atoms out of
the simulation box, causing LAMMPS to stop with an error.

\sphinxAtStartPar
Using this fix can overcome that problem.  Forces on atoms must still
be computable (which typically means two atoms must have a separation
distance \textgreater{} 0.0).  But large velocities generated by large forces are
reset to a value that corresponds to a displacement of length \sphinxstyleemphasis{xmax}
in a single timestep.  \sphinxstyleemphasis{Xmax} is specified in distance units; see the
{\hyperref[\detokenize{units::doc}]{\sphinxcrossref{\DUrole{doc}{units}}}} command for details.  The value of \sphinxstyleemphasis{xmax} should be
consistent with the neighbor skin distance and the frequency of
neighbor list re\sphinxhyphen{}building, so that pairwise interactions are not
missed on successive timesteps as atoms move.  See the
{\hyperref[\detokenize{neighbor::doc}]{\sphinxcrossref{\DUrole{doc}{neighbor}}}} and {\hyperref[\detokenize{neigh_modify::doc}]{\sphinxcrossref{\DUrole{doc}{neigh\_modify}}}} commands
for details.

\sphinxAtStartPar
Note that if a velocity reset occurs the integrator will not conserve
energy.  On steps where no velocity resets occur, this integrator is
exactly like the {\hyperref[\detokenize{fix_nve::doc}]{\sphinxcrossref{\DUrole{doc}{fix nve}}}} command.  Since forces are
unaltered, pressures computed by thermodynamic output will still be
very large for overlapped configurations.

\begin{sphinxadmonition}{note}{Note:}
\sphinxAtStartPar
You should not use {\hyperref[\detokenize{fix_shake::doc}]{\sphinxcrossref{\DUrole{doc}{fix shake}}}} in conjunction
with this fix.  That is because fix shake applies constraint forces
based on the predicted positions of atoms after the next timestep.
It has no way of knowing the timestep may change due to this fix,
which will cause the constraint forces to be invalid.  A better
strategy is to turn off fix shake when performing initial dynamics
that need this fix, then turn fix shake on when doing normal dynamics
with a fixed\sphinxhyphen{}size timestep.
\end{sphinxadmonition}


\subsection{Restart, fix\_modify, output, run start/stop, minimize info}
\label{\detokenize{fix_nve_limit:restart-fix-modify-output-run-start-stop-minimize-info}}
\sphinxAtStartPar
No information about this fix is written to {\hyperref[\detokenize{restart::doc}]{\sphinxcrossref{\DUrole{doc}{binary restart files}}}}.  None of the {\hyperref[\detokenize{fix_modify::doc}]{\sphinxcrossref{\DUrole{doc}{fix\_modify}}}} options
are relevant to this fix.

\sphinxAtStartPar
This fix computes a global scalar which can be accessed by various
{\hyperref[\detokenize{Howto_output::doc}]{\sphinxcrossref{\DUrole{doc}{output commands}}}}.  The scalar is the count of how
many updates of atom’s velocity/position were limited by the maximum
distance criterion.  This should be roughly the number of atoms so
affected, except that updates occur at both the beginning and end of a
timestep in a velocity Verlet timestepping algorithm.  This is a
cumulative quantity for the current run, but is re\sphinxhyphen{}initialized to zero
each time a run is performed.  The scalar value calculated by this fix
is “extensive”.

\sphinxAtStartPar
No parameter of this fix can be used with the \sphinxstyleemphasis{start/stop} keywords of
the {\hyperref[\detokenize{run::doc}]{\sphinxcrossref{\DUrole{doc}{run}}}} command.  This fix is not invoked during {\hyperref[\detokenize{minimize::doc}]{\sphinxcrossref{\DUrole{doc}{energy minimization}}}}.


\subsection{Restrictions}
\label{\detokenize{fix_nve_limit:restrictions}}\begin{quote}

\sphinxAtStartPar
none
\end{quote}


\subsection{Related commands}
\label{\detokenize{fix_nve_limit:related-commands}}
\sphinxAtStartPar
{\hyperref[\detokenize{fix_nve::doc}]{\sphinxcrossref{\DUrole{doc}{fix nve}}}}, {\hyperref[\detokenize{fix_nve_noforce::doc}]{\sphinxcrossref{\DUrole{doc}{fix nve/noforce}}}},
{\hyperref[\detokenize{pair_soft::doc}]{\sphinxcrossref{\DUrole{doc}{pair\_style soft}}}}


\subsection{Default}
\label{\detokenize{fix_nve_limit:default}}
\sphinxAtStartPar
none

\sphinxstepscope

\index{fix nve/line@\spxentry{fix nve/line}}\ignorespaces 

\section{fix nve/line command}
\label{\detokenize{fix_nve_line:fix-nve-line-command}}\label{\detokenize{fix_nve_line:index-0}}\label{\detokenize{fix_nve_line::doc}}

\subsection{Syntax}
\label{\detokenize{fix_nve_line:syntax}}
\begin{sphinxVerbatim}[commandchars=\\\{\}]
\PYG{k}{fix }\PYG{n+nv+nvIdentifier}{ID}\PYG{+w}{ }\PYG{n+nv+nvIdentifier}{group\PYGZhy{}ID}\PYG{+w}{ }\PYG{n}{nve}\PYG{o}{/}\PYG{n}{line}
\end{sphinxVerbatim}
\begin{itemize}
\item {} 
\sphinxAtStartPar
ID, group\sphinxhyphen{}ID are documented in {\hyperref[\detokenize{fix::doc}]{\sphinxcrossref{\DUrole{doc}{fix}}}} command

\item {} 
\sphinxAtStartPar
nve/line = style name of this fix command

\end{itemize}


\subsection{Examples}
\label{\detokenize{fix_nve_line:examples}}
\begin{sphinxVerbatim}[commandchars=\\\{\}]
\PYG{k}{fix }\PYG{n+nv+nvIdentifier}{1}\PYG{+w}{ }\PYG{n+nv+nvIdentifier}{all}\PYG{+w}{ }\PYG{n}{nve}\PYG{o}{/}\PYG{n}{line}
\end{sphinxVerbatim}


\subsection{Description}
\label{\detokenize{fix_nve_line:description}}
\sphinxAtStartPar
Perform constant NVE integration to update position, velocity,
orientation, and angular velocity for line segment particles in the
group each timestep.  V is volume; E is energy.  This creates a system
trajectory consistent with the microcanonical ensemble.  See {\hyperref[\detokenize{Howto_spherical::doc}]{\sphinxcrossref{\DUrole{doc}{Howto spherical}}}} page for an overview of using line
segment particles.

\sphinxAtStartPar
This fix differs from the {\hyperref[\detokenize{fix_nve::doc}]{\sphinxcrossref{\DUrole{doc}{fix nve}}}} command, which
assumes point particles and only updates their position and velocity.


\subsection{Restart, fix\_modify, output, run start/stop, minimize info}
\label{\detokenize{fix_nve_line:restart-fix-modify-output-run-start-stop-minimize-info}}
\sphinxAtStartPar
No information about this fix is written to {\hyperref[\detokenize{restart::doc}]{\sphinxcrossref{\DUrole{doc}{binary restart files}}}}.  None of the {\hyperref[\detokenize{fix_modify::doc}]{\sphinxcrossref{\DUrole{doc}{fix\_modify}}}} options
are relevant to this fix.  No global or per\sphinxhyphen{}atom quantities are stored
by this fix for access by various {\hyperref[\detokenize{Howto_output::doc}]{\sphinxcrossref{\DUrole{doc}{output commands}}}}.
No parameter of this fix can be used with the \sphinxstyleemphasis{start/stop} keywords of
the {\hyperref[\detokenize{run::doc}]{\sphinxcrossref{\DUrole{doc}{run}}}} command.  This fix is not invoked during {\hyperref[\detokenize{minimize::doc}]{\sphinxcrossref{\DUrole{doc}{energy minimization}}}}.


\subsection{Restrictions}
\label{\detokenize{fix_nve_line:restrictions}}
\sphinxAtStartPar
This fix is part of the ASPHERE package.  It is only enabled if LAMMPS
was built with that package.  See the {\hyperref[\detokenize{Build_package::doc}]{\sphinxcrossref{\DUrole{doc}{Build package}}}} page for more info.

\sphinxAtStartPar
This fix requires that particles be line segments as defined by the
{\hyperref[\detokenize{atom_style::doc}]{\sphinxcrossref{\DUrole{doc}{atom\_style line}}}} command.


\subsection{Related commands}
\label{\detokenize{fix_nve_line:related-commands}}
\sphinxAtStartPar
{\hyperref[\detokenize{fix_nve::doc}]{\sphinxcrossref{\DUrole{doc}{fix nve}}}}, {\hyperref[\detokenize{fix_nve_asphere::doc}]{\sphinxcrossref{\DUrole{doc}{fix nve/asphere}}}}


\subsection{Default}
\label{\detokenize{fix_nve_line:default}}
\sphinxAtStartPar
none

\sphinxstepscope

\index{fix nve/manifold/rattle@\spxentry{fix nve/manifold/rattle}}\ignorespaces 

\section{fix nve/manifold/rattle command}
\label{\detokenize{fix_nve_manifold_rattle:fix-nve-manifold-rattle-command}}\label{\detokenize{fix_nve_manifold_rattle:index-0}}\label{\detokenize{fix_nve_manifold_rattle::doc}}

\subsection{Syntax}
\label{\detokenize{fix_nve_manifold_rattle:syntax}}
\begin{sphinxVerbatim}[commandchars=\\\{\}]
\PYG{k}{fix }\PYG{n+nv+nvIdentifier}{ID}\PYG{+w}{ }\PYG{n+nv+nvIdentifier}{group\PYGZhy{}ID}\PYG{+w}{ }\PYG{n}{nve}\PYG{o}{/}\PYG{n}{manifold}\PYG{o}{/}\PYG{n}{rattle}\PYG{+w}{ }\PYG{n}{tol}\PYG{+w}{ }\PYG{n}{maxit}\PYG{+w}{ }\PYG{n}{manifold}\PYG{+w}{ }\PYG{n}{manifold}\PYG{o}{\PYGZhy{}}\PYG{n}{args}\PYG{+w}{ }\PYG{n}{keyword}\PYG{+w}{ }\PYG{n}{value}\PYG{+w}{ }\PYG{n}{...}
\end{sphinxVerbatim}
\begin{itemize}
\item {} 
\sphinxAtStartPar
ID, group\sphinxhyphen{}ID are documented in {\hyperref[\detokenize{fix::doc}]{\sphinxcrossref{\DUrole{doc}{fix}}}} command

\item {} 
\sphinxAtStartPar
nve/manifold/rattle = style name of this fix command

\item {} 
\sphinxAtStartPar
tol = tolerance to which Newton iteration must converge

\item {} 
\sphinxAtStartPar
maxit = maximum number of iterations to perform

\item {} 
\sphinxAtStartPar
manifold = name of the manifold

\item {} 
\sphinxAtStartPar
manifold\sphinxhyphen{}args = parameters for the manifold

\item {} 
\sphinxAtStartPar
one or more keyword/value pairs may be appended
\begin{sphinxalltt}
keyword = \sphinxstyleemphasis{every}
  \sphinxstyleemphasis{every} values = N
    N = print info about iteration every N steps. N = 0 means no output
\end{sphinxalltt}

\end{itemize}


\subsection{Examples}
\label{\detokenize{fix_nve_manifold_rattle:examples}}
\begin{sphinxVerbatim}[commandchars=\\\{\}]
\PYG{k}{fix }\PYG{n+nv+nvIdentifier}{1}\PYG{+w}{ }\PYG{n+nv+nvIdentifier}{all}\PYG{+w}{ }\PYG{n}{nve}\PYG{o}{/}\PYG{n}{manifold}\PYG{o}{/}\PYG{n}{rattle}\PYG{+w}{ }\PYG{l+m}{1e\PYGZhy{}4}\PYG{+w}{ }\PYG{l+m}{10}\PYG{+w}{ }\PYG{n}{sphere}\PYG{+w}{ }\PYG{l+m}{5.0}
\PYG{k}{fix }\PYG{n+nv+nvIdentifier}{step}\PYG{+w}{ }\PYG{n+nv+nvIdentifier}{all}\PYG{+w}{ }\PYG{n}{nve}\PYG{o}{/}\PYG{n}{manifold}\PYG{o}{/}\PYG{n}{rattle}\PYG{+w}{ }\PYG{l+m}{1e\PYGZhy{}8}\PYG{+w}{ }\PYG{l+m}{100}\PYG{+w}{ }\PYG{n}{ellipsoid}\PYG{+w}{ }\PYG{l+m}{2.5}\PYG{+w}{ }\PYG{l+m}{2.5}\PYG{+w}{ }\PYG{l+m}{5.0}\PYG{+w}{ }\PYG{n}{every}\PYG{+w}{ }\PYG{l+m}{25}
\end{sphinxVerbatim}


\subsection{Description}
\label{\detokenize{fix_nve_manifold_rattle:description}}
\sphinxAtStartPar
Perform constant NVE integration to update position and velocity for
atoms constrained to a curved surface (manifold) in the group each
timestep. The constraint is handled by RATTLE {\hyperref[\detokenize{fix_nve_manifold_rattle:andersen1}]{\sphinxcrossref{\DUrole{std,std-ref}{(Andersen)}}}}
written out for the special case of single\sphinxhyphen{}particle constraints as
explained in {\hyperref[\detokenize{fix_nve_manifold_rattle:paquay2}]{\sphinxcrossref{\DUrole{std,std-ref}{(Paquay)}}}}.  V is volume; E is energy. This way,
the dynamics of particles constrained to curved surfaces can be
studied. If combined with {\hyperref[\detokenize{fix_langevin::doc}]{\sphinxcrossref{\DUrole{doc}{fix langevin}}}}, this
generates Brownian motion of particles constrained to a curved
surface. For a list of currently supported manifolds and their
parameters, see the {\hyperref[\detokenize{Howto_manifold::doc}]{\sphinxcrossref{\DUrole{doc}{Howto manifold}}}} doc page.

\sphinxAtStartPar
Note that the particles must initially be close to the manifold in
question. If not, RATTLE will not be able to iterate until the
constraint is satisfied, and an error is generated. For simple
manifolds this can be achieved with \sphinxstyleemphasis{region} and \sphinxstyleemphasis{create\_atoms}
commands, but for more complex surfaces it might be more useful to
write a script.

\sphinxAtStartPar
The manifold args may be equal\sphinxhyphen{}style variables, like so:

\begin{sphinxVerbatim}[commandchars=\\\{\}]
\PYG{k}{variable }\PYG{n+nv+nvIdentifier}{R}\PYG{+w}{ }\PYG{n}{equal}\PYG{+w}{ }\PYG{l+s}{\PYGZdq{}}\PYG{l+s}{ramp(5.0,3.0)}\PYG{l+s}{\PYGZdq{}}
\PYG{k}{fix }\PYG{n+nv+nvIdentifier}{shrink\PYGZus{}sphere}\PYG{+w}{ }\PYG{n+nv+nvIdentifier}{all}\PYG{+w}{ }\PYG{n}{nve}\PYG{o}{/}\PYG{n}{manifold}\PYG{o}{/}\PYG{n}{rattle}\PYG{+w}{ }\PYG{l+m}{1e\PYGZhy{}4}\PYG{+w}{ }\PYG{l+m}{10}\PYG{+w}{ }\PYG{n}{sphere}\PYG{+w}{ }\PYG{n}{v\PYGZus{}R}
\end{sphinxVerbatim}

\sphinxAtStartPar
In this case, the manifold parameter will change in time according to
the variable.  This is not a problem for the time integrator as long
as the change of the manifold is slow with respect to the dynamics of
the particles.  Note that if the manifold has to exert work on the
particles because of these changes, the total energy might not be
conserved.


\bigskip\hrule\bigskip



\subsection{Restart, fix\_modify, output, run start/stop, minimize info}
\label{\detokenize{fix_nve_manifold_rattle:restart-fix-modify-output-run-start-stop-minimize-info}}
\sphinxAtStartPar
No information about this fix is written to {\hyperref[\detokenize{restart::doc}]{\sphinxcrossref{\DUrole{doc}{binary restart files}}}}.  None of the {\hyperref[\detokenize{fix_modify::doc}]{\sphinxcrossref{\DUrole{doc}{fix\_modify}}}} options
are relevant to this fix.  No global or per\sphinxhyphen{}atom quantities are stored
by this fix for access by various {\hyperref[\detokenize{Howto_output::doc}]{\sphinxcrossref{\DUrole{doc}{output commands}}}}.
No parameter of this fix can be used with the \sphinxstyleemphasis{start/stop} keywords of
the {\hyperref[\detokenize{run::doc}]{\sphinxcrossref{\DUrole{doc}{run}}}} command.  This fix is not invoked during {\hyperref[\detokenize{minimize::doc}]{\sphinxcrossref{\DUrole{doc}{energy minimization}}}}.


\bigskip\hrule\bigskip



\subsection{Restrictions}
\label{\detokenize{fix_nve_manifold_rattle:restrictions}}
\sphinxAtStartPar
This fix is part of the MANIFOLD package. It is only enabled if
LAMMPS was built with that package. See the {\hyperref[\detokenize{Build_package::doc}]{\sphinxcrossref{\DUrole{doc}{Build package}}}} page for more info.


\bigskip\hrule\bigskip



\subsection{Related commands}
\label{\detokenize{fix_nve_manifold_rattle:related-commands}}
\sphinxAtStartPar
{\hyperref[\detokenize{fix_nvt_manifold_rattle::doc}]{\sphinxcrossref{\DUrole{doc}{fix nvt/manifold/rattle}}}}, {\hyperref[\detokenize{fix_manifoldforce::doc}]{\sphinxcrossref{\DUrole{doc}{fix manifoldforce}}}}


\subsection{Default}
\label{\detokenize{fix_nve_manifold_rattle:default}}
\sphinxAtStartPar
every = 0, tchain = 3


\bigskip\hrule\bigskip

\phantomsection\label{\detokenize{fix_nve_manifold_rattle:andersen1}}
\sphinxAtStartPar
\sphinxstylestrong{(Andersen)} Andersen, J. Comp. Phys. 52, 24, (1983).

\phantomsection\label{\detokenize{fix_nve_manifold_rattle:paquay2}}
\sphinxAtStartPar
\sphinxstylestrong{(Paquay)} Paquay and Kusters, Biophys. J., 110, 6, (2016).
preprint available at \sphinxhref{https://arxiv.org/abs/1411.3019/}{arXiv:1411.3019}.

\sphinxstepscope

\index{fix nve/noforce@\spxentry{fix nve/noforce}}\ignorespaces 

\section{fix nve/noforce command}
\label{\detokenize{fix_nve_noforce:fix-nve-noforce-command}}\label{\detokenize{fix_nve_noforce:index-0}}\label{\detokenize{fix_nve_noforce::doc}}

\subsection{Syntax}
\label{\detokenize{fix_nve_noforce:syntax}}
\begin{sphinxVerbatim}[commandchars=\\\{\}]
\PYG{k}{fix }\PYG{n+nv+nvIdentifier}{ID}\PYG{+w}{ }\PYG{n+nv+nvIdentifier}{group\PYGZhy{}ID}\PYG{+w}{ }\PYG{n}{nve}
\end{sphinxVerbatim}
\begin{itemize}
\item {} 
\sphinxAtStartPar
ID, group\sphinxhyphen{}ID are documented in {\hyperref[\detokenize{fix::doc}]{\sphinxcrossref{\DUrole{doc}{fix}}}} command

\item {} 
\sphinxAtStartPar
nve/noforce = style name of this fix command

\end{itemize}


\subsection{Examples}
\label{\detokenize{fix_nve_noforce:examples}}
\begin{sphinxVerbatim}[commandchars=\\\{\}]
\PYG{k}{fix }\PYG{n+nv+nvIdentifier}{3}\PYG{+w}{ }\PYG{n+nv+nvIdentifier}{wall}\PYG{+w}{ }\PYG{n}{nve}\PYG{o}{/}\PYG{n}{noforce}
\end{sphinxVerbatim}


\subsection{Description}
\label{\detokenize{fix_nve_noforce:description}}
\sphinxAtStartPar
Perform updates of position, but not velocity for atoms in the group
each timestep.  In other words, the force on the atoms is ignored and
their velocity is not updated.  The atom velocities are used to update
their positions.

\sphinxAtStartPar
This can be useful for wall atoms, when you set their velocities, and
want the wall to move (or stay stationary) in a prescribed fashion.

\sphinxAtStartPar
This can also be accomplished via the {\hyperref[\detokenize{fix_setforce::doc}]{\sphinxcrossref{\DUrole{doc}{fix setforce}}}}
command, but with fix nve/noforce, the forces on the wall atoms are
unchanged, and can thus be printed by the {\hyperref[\detokenize{dump::doc}]{\sphinxcrossref{\DUrole{doc}{dump}}}} command or
queried with an equal\sphinxhyphen{}style {\hyperref[\detokenize{variable::doc}]{\sphinxcrossref{\DUrole{doc}{variable}}}} that uses the
fcm() group function to compute the total force on the group of atoms.


\subsection{Restart, fix\_modify, output, run start/stop, minimize info}
\label{\detokenize{fix_nve_noforce:restart-fix-modify-output-run-start-stop-minimize-info}}
\sphinxAtStartPar
No information about this fix is written to {\hyperref[\detokenize{restart::doc}]{\sphinxcrossref{\DUrole{doc}{binary restart files}}}}.  None of the {\hyperref[\detokenize{fix_modify::doc}]{\sphinxcrossref{\DUrole{doc}{fix\_modify}}}} options
are relevant to this fix.  No global or per\sphinxhyphen{}atom quantities are stored
by this fix for access by various {\hyperref[\detokenize{Howto_output::doc}]{\sphinxcrossref{\DUrole{doc}{output commands}}}}.
No parameter of this fix can be used with the \sphinxstyleemphasis{start/stop} keywords of
the {\hyperref[\detokenize{run::doc}]{\sphinxcrossref{\DUrole{doc}{run}}}} command.  This fix is not invoked during {\hyperref[\detokenize{minimize::doc}]{\sphinxcrossref{\DUrole{doc}{energy minimization}}}}.


\subsection{Restrictions}
\label{\detokenize{fix_nve_noforce:restrictions}}
\sphinxAtStartPar
none


\subsection{Related commands}
\label{\detokenize{fix_nve_noforce:related-commands}}
\sphinxAtStartPar
{\hyperref[\detokenize{fix_nve::doc}]{\sphinxcrossref{\DUrole{doc}{fix nve}}}}


\subsection{Default}
\label{\detokenize{fix_nve_noforce:default}}
\sphinxAtStartPar
none

\sphinxstepscope

\index{fix nve/sphere@\spxentry{fix nve/sphere}}\ignorespaces 
\index{fix nve/sphere/omp@\spxentry{fix nve/sphere/omp}}\ignorespaces 
\index{fix nve/sphere/kk@\spxentry{fix nve/sphere/kk}}\ignorespaces 

\section{fix nve/sphere command}
\label{\detokenize{fix_nve_sphere:fix-nve-sphere-command}}\label{\detokenize{fix_nve_sphere:index-2}}\label{\detokenize{fix_nve_sphere:index-1}}\label{\detokenize{fix_nve_sphere:index-0}}\label{\detokenize{fix_nve_sphere::doc}}
\sphinxAtStartPar
Accelerator Variants: \sphinxstyleemphasis{nve/sphere/omp}, \sphinxstyleemphasis{nve/sphere/kk}


\subsection{Syntax}
\label{\detokenize{fix_nve_sphere:syntax}}
\begin{sphinxVerbatim}[commandchars=\\\{\}]
\PYG{k}{fix }\PYG{n+nv+nvIdentifier}{ID}\PYG{+w}{ }\PYG{n+nv+nvIdentifier}{group\PYGZhy{}ID}\PYG{+w}{ }\PYG{n}{nve}\PYG{o}{/}\PYG{n}{sphere}
\end{sphinxVerbatim}
\begin{itemize}
\item {} 
\sphinxAtStartPar
ID, group\sphinxhyphen{}ID are documented in {\hyperref[\detokenize{fix::doc}]{\sphinxcrossref{\DUrole{doc}{fix}}}} command

\item {} 
\sphinxAtStartPar
nve/sphere = style name of this fix command

\item {} 
\sphinxAtStartPar
zero or more keyword/value pairs may be appended

\item {} 
\sphinxAtStartPar
keyword = \sphinxstyleemphasis{update} or \sphinxstyleemphasis{disc}
\begin{sphinxalltt}
\sphinxstyleemphasis{update} value = \sphinxstyleemphasis{dipole} or \sphinxstyleemphasis{dipole/dlm}
  dipole = update orientation of dipole moment during integration
  dipole/dlm = use DLM integrator to update dipole orientation
\sphinxstyleemphasis{disc} value = none = treat particles as 2d discs, not spheres
\end{sphinxalltt}

\end{itemize}


\subsection{Examples}
\label{\detokenize{fix_nve_sphere:examples}}
\begin{sphinxVerbatim}[commandchars=\\\{\}]
\PYG{k}{fix }\PYG{n+nv+nvIdentifier}{1}\PYG{+w}{ }\PYG{n+nv+nvIdentifier}{all}\PYG{+w}{ }\PYG{n}{nve}\PYG{o}{/}\PYG{n}{sphere}
\PYG{k}{fix }\PYG{n+nv+nvIdentifier}{1}\PYG{+w}{ }\PYG{n+nv+nvIdentifier}{all}\PYG{+w}{ }\PYG{n}{nve}\PYG{o}{/}\PYG{n}{sphere}\PYG{+w}{ }\PYG{n}{update}\PYG{+w}{ }\PYG{n}{dipole}
\PYG{k}{fix }\PYG{n+nv+nvIdentifier}{1}\PYG{+w}{ }\PYG{n+nv+nvIdentifier}{all}\PYG{+w}{ }\PYG{n}{nve}\PYG{o}{/}\PYG{n}{sphere}\PYG{+w}{ }\PYG{n}{disc}
\PYG{k}{fix }\PYG{n+nv+nvIdentifier}{1}\PYG{+w}{ }\PYG{n+nv+nvIdentifier}{all}\PYG{+w}{ }\PYG{n}{nve}\PYG{o}{/}\PYG{n}{sphere}\PYG{+w}{ }\PYG{n}{update}\PYG{+w}{ }\PYG{n}{dipole}\PYG{o}{/}\PYG{n}{dlm}
\end{sphinxVerbatim}


\subsection{Description}
\label{\detokenize{fix_nve_sphere:description}}
\sphinxAtStartPar
Perform constant NVE integration to update position, velocity, and
angular velocity for finite\sphinxhyphen{}size spherical particles in the group each
timestep.  V is volume; E is energy.  This creates a system trajectory
consistent with the microcanonical ensemble.

\sphinxAtStartPar
This fix differs from the {\hyperref[\detokenize{fix_nve::doc}]{\sphinxcrossref{\DUrole{doc}{fix nve}}}} command, which
assumes point particles and only updates their position and velocity.

\sphinxAtStartPar
If the \sphinxstyleemphasis{update} keyword is used with the \sphinxstyleemphasis{dipole} value, then the
orientation of the dipole moment of each particle is also updated
during the time integration.  This option should be used for models
where a dipole moment is assigned to finite\sphinxhyphen{}size particles,
e.g. spheroids via use of the {\hyperref[\detokenize{atom_style::doc}]{\sphinxcrossref{\DUrole{doc}{atom\_style hybrid sphere dipole}}}} command.

\sphinxAtStartPar
The default dipole orientation integrator can be changed to the
Dullweber\sphinxhyphen{}Leimkuhler\sphinxhyphen{}McLachlan integration scheme
{\hyperref[\detokenize{fix_nh:nh-dullweber}]{\sphinxcrossref{\DUrole{std,std-ref}{(Dullweber)}}}} when using \sphinxstyleemphasis{update} with the value
\sphinxstyleemphasis{dipole/dlm}. This integrator is symplectic and time\sphinxhyphen{}reversible,
giving better energy conservation and allows slightly longer timesteps
at only a small additional computational cost.

\sphinxAtStartPar
If the \sphinxstyleemphasis{disc} keyword is used, then each particle is treated as a 2d
disc (circle) instead of as a sphere.  This is only possible for 2d
simulations, as defined by the {\hyperref[\detokenize{dimension::doc}]{\sphinxcrossref{\DUrole{doc}{dimension}}}} keyword.
The only difference between discs and spheres in this context is their
moment of inertia, as used in the time integration.


\bigskip\hrule\bigskip


\sphinxAtStartPar
Styles with a \sphinxstyleemphasis{gpu}, \sphinxstyleemphasis{intel}, \sphinxstyleemphasis{kk}, \sphinxstyleemphasis{omp}, or \sphinxstyleemphasis{opt} suffix are
functionally the same as the corresponding style without the suffix.
They have been optimized to run faster, depending on your available
hardware, as discussed on the {\hyperref[\detokenize{Speed_packages::doc}]{\sphinxcrossref{\DUrole{doc}{Accelerator packages}}}}
page.  The accelerated styles take the same arguments and should
produce the same results, except for round\sphinxhyphen{}off and precision issues.

\sphinxAtStartPar
These accelerated styles are part of the GPU, INTEL, KOKKOS,
OPENMP, and OPT packages, respectively.  They are only enabled if
LAMMPS was built with those packages.  See the {\hyperref[\detokenize{Build_package::doc}]{\sphinxcrossref{\DUrole{doc}{Build package}}}} page for more info.

\sphinxAtStartPar
You can specify the accelerated styles explicitly in your input script
by including their suffix, or you can use the {\hyperref[\detokenize{Run_options::doc}]{\sphinxcrossref{\DUrole{doc}{\sphinxhyphen{}suffix command\sphinxhyphen{}line switch}}}} when you invoke LAMMPS, or you can use the
{\hyperref[\detokenize{suffix::doc}]{\sphinxcrossref{\DUrole{doc}{suffix}}}} command in your input script.

\sphinxAtStartPar
See the {\hyperref[\detokenize{Speed_packages::doc}]{\sphinxcrossref{\DUrole{doc}{Accelerator packages}}}} page for more
instructions on how to use the accelerated styles effectively.


\bigskip\hrule\bigskip



\subsection{Restart, fix\_modify, output, run start/stop, minimize info}
\label{\detokenize{fix_nve_sphere:restart-fix-modify-output-run-start-stop-minimize-info}}
\sphinxAtStartPar
No information about this fix is written to {\hyperref[\detokenize{restart::doc}]{\sphinxcrossref{\DUrole{doc}{binary restart files}}}}.  None of the {\hyperref[\detokenize{fix_modify::doc}]{\sphinxcrossref{\DUrole{doc}{fix\_modify}}}} options are
relevant to this fix.  No global or per\sphinxhyphen{}atom quantities are stored by
this fix for access by various {\hyperref[\detokenize{Howto_output::doc}]{\sphinxcrossref{\DUrole{doc}{output commands}}}}.
No parameter of this fix can be used with the \sphinxstyleemphasis{start/stop} keywords of
the {\hyperref[\detokenize{run::doc}]{\sphinxcrossref{\DUrole{doc}{run}}}} command.  This fix is not invoked during
{\hyperref[\detokenize{minimize::doc}]{\sphinxcrossref{\DUrole{doc}{energy minimization}}}}.


\subsection{Restrictions}
\label{\detokenize{fix_nve_sphere:restrictions}}
\sphinxAtStartPar
This fix requires that atoms store torque and angular velocity (omega)
and a radius as defined by the {\hyperref[\detokenize{atom_style::doc}]{\sphinxcrossref{\DUrole{doc}{atom\_style sphere}}}}
command.  If the \sphinxstyleemphasis{dipole} keyword is used, then they must also store a
dipole moment as defined by the {\hyperref[\detokenize{atom_style::doc}]{\sphinxcrossref{\DUrole{doc}{atom\_style dipole}}}}
command.

\sphinxAtStartPar
All particles in the group must be finite\sphinxhyphen{}size spheres.  They cannot
be point particles.

\sphinxAtStartPar
Use of the \sphinxstyleemphasis{disc} keyword is only allowed for 2d simulations, as
defined by the {\hyperref[\detokenize{dimension::doc}]{\sphinxcrossref{\DUrole{doc}{dimension}}}} keyword.


\subsection{Related commands}
\label{\detokenize{fix_nve_sphere:related-commands}}
\sphinxAtStartPar
{\hyperref[\detokenize{fix_nve::doc}]{\sphinxcrossref{\DUrole{doc}{fix nve}}}}, {\hyperref[\detokenize{fix_nve_asphere::doc}]{\sphinxcrossref{\DUrole{doc}{fix nve/asphere}}}}


\subsection{Default}
\label{\detokenize{fix_nve_sphere:default}}
\sphinxAtStartPar
none


\bigskip\hrule\bigskip

\phantomsection\label{\detokenize{fix_nve_sphere:nve-dullweber}}
\sphinxAtStartPar
\sphinxstylestrong{(Dullweber)} Dullweber, Leimkuhler and McLachlan, J Chem Phys, 107,
5840 (1997).

\sphinxstepscope

\index{fix nve/spin@\spxentry{fix nve/spin}}\ignorespaces 

\section{fix nve/spin command}
\label{\detokenize{fix_nve_spin:fix-nve-spin-command}}\label{\detokenize{fix_nve_spin:index-0}}\label{\detokenize{fix_nve_spin::doc}}

\subsection{Syntax}
\label{\detokenize{fix_nve_spin:syntax}}
\begin{sphinxVerbatim}[commandchars=\\\{\}]
\PYG{k}{fix }\PYG{n+nv+nvIdentifier}{ID}\PYG{+w}{ }\PYG{n+nv+nvIdentifier}{group\PYGZhy{}ID}\PYG{+w}{ }\PYG{n}{nve}\PYG{o}{/}\PYG{n}{spin}\PYG{+w}{ }\PYG{n}{keyword}\PYG{+w}{ }\PYG{n}{values}
\end{sphinxVerbatim}
\begin{itemize}
\item {} 
\sphinxAtStartPar
ID, group\sphinxhyphen{}ID are documented in {\hyperref[\detokenize{fix::doc}]{\sphinxcrossref{\DUrole{doc}{fix}}}} command

\item {} 
\sphinxAtStartPar
nve/spin = style name of this fix command

\item {} 
\sphinxAtStartPar
keyword = \sphinxstyleemphasis{lattice}
\begin{sphinxalltt}
\sphinxstyleemphasis{lattice} value = \sphinxstyleemphasis{moving} or \sphinxstyleemphasis{frozen}
  moving = integrate both spin and atomic degress of freedom
  frozen = integrate spins on a fixed lattice
\end{sphinxalltt}

\end{itemize}


\subsection{Examples}
\label{\detokenize{fix_nve_spin:examples}}
\begin{sphinxVerbatim}[commandchars=\\\{\}]
\PYG{k}{fix }\PYG{n+nv+nvIdentifier}{3}\PYG{+w}{ }\PYG{n+nv+nvIdentifier}{all}\PYG{+w}{ }\PYG{n}{nve}\PYG{o}{/}\PYG{n}{spin}\PYG{+w}{ }\PYG{n}{lattice}\PYG{+w}{ }\PYG{n}{moving}
\PYG{k}{fix }\PYG{n+nv+nvIdentifier}{1}\PYG{+w}{ }\PYG{n+nv+nvIdentifier}{all}\PYG{+w}{ }\PYG{n}{nve}\PYG{o}{/}\PYG{n}{spin}\PYG{+w}{ }\PYG{n}{lattice}\PYG{+w}{ }\PYG{n}{frozen}
\end{sphinxVerbatim}


\subsection{Description}
\label{\detokenize{fix_nve_spin:description}}
\sphinxAtStartPar
Perform a symplectic integration for the spin or spin\sphinxhyphen{}lattice system.

\sphinxAtStartPar
The \sphinxstyleemphasis{lattice} keyword defines if the spins are integrated on a lattice
of fixed atoms (lattice = frozen), or if atoms are moving
(lattice = moving).
The first case corresponds to a spin dynamics calculation, and
the second to a spin\sphinxhyphen{}lattice calculation.
By default a spin\sphinxhyphen{}lattice integration is performed (lattice = moving).

\sphinxAtStartPar
The \sphinxstyleemphasis{nve/spin} fix applies a Suzuki\sphinxhyphen{}Trotter decomposition to
the equations of motion of the spin lattice system, following the scheme:

\noindent{\hspace*{\fill}\sphinxincludegraphics{{fix_integration_spin_stdecomposition}.jpg}\hspace*{\fill}}

\sphinxAtStartPar
according to the implementation reported in {\hyperref[\detokenize{fix_nve_spin:omelyan1}]{\sphinxcrossref{\DUrole{std,std-ref}{(Omelyan)}}}}.

\sphinxAtStartPar
A sectoring method enables this scheme for parallel calculations.
The implementation of this sectoring algorithm is reported
in {\hyperref[\detokenize{fix_nve_spin:tranchida1}]{\sphinxcrossref{\DUrole{std,std-ref}{(Tranchida)}}}}.


\bigskip\hrule\bigskip



\subsection{Restrictions}
\label{\detokenize{fix_nve_spin:restrictions}}
\sphinxAtStartPar
This fix style can only be used if LAMMPS was built with the SPIN
package.  See the {\hyperref[\detokenize{Build_package::doc}]{\sphinxcrossref{\DUrole{doc}{Build package}}}} page for more
info.

\sphinxAtStartPar
To use the spin algorithm, it is necessary to define a map with
the atom\_modify command. Typically, by adding the command:

\begin{sphinxVerbatim}[commandchars=\\\{\}]
\PYG{k}{atom\PYGZus{}modify}\PYG{+w}{ }\PYG{n}{map}\PYG{+w}{ }\PYG{n}{array}
\end{sphinxVerbatim}

\sphinxAtStartPar
before you create the simulation box. Note that the keyword “hash”
instead of “array” is also valid.


\subsection{Related commands}
\label{\detokenize{fix_nve_spin:related-commands}}
\sphinxAtStartPar
{\hyperref[\detokenize{atom_style::doc}]{\sphinxcrossref{\DUrole{doc}{atom\_style spin}}}}, {\hyperref[\detokenize{fix_nve::doc}]{\sphinxcrossref{\DUrole{doc}{fix nve}}}}


\subsection{Default}
\label{\detokenize{fix_nve_spin:default}}
\sphinxAtStartPar
The option default is lattice = moving.


\bigskip\hrule\bigskip

\phantomsection\label{\detokenize{fix_nve_spin:omelyan1}}
\sphinxAtStartPar
\sphinxstylestrong{(Omelyan)} Omelyan, Mryglod, and Folk. Phys. Rev. Lett.
86(5), 898. (2001).

\phantomsection\label{\detokenize{fix_nve_spin:tranchida1}}
\sphinxAtStartPar
\sphinxstylestrong{(Tranchida)} Tranchida, Plimpton, Thibaudeau and Thompson,
Journal of Computational Physics, 372, 406\sphinxhyphen{}425, (2018).

\sphinxstepscope

\index{fix nve/tri@\spxentry{fix nve/tri}}\ignorespaces 

\section{fix nve/tri command}
\label{\detokenize{fix_nve_tri:fix-nve-tri-command}}\label{\detokenize{fix_nve_tri:index-0}}\label{\detokenize{fix_nve_tri::doc}}

\subsection{Syntax}
\label{\detokenize{fix_nve_tri:syntax}}
\begin{sphinxVerbatim}[commandchars=\\\{\}]
\PYG{k}{fix }\PYG{n+nv+nvIdentifier}{ID}\PYG{+w}{ }\PYG{n+nv+nvIdentifier}{group\PYGZhy{}ID}\PYG{+w}{ }\PYG{n}{nve}\PYG{o}{/}\PYG{n}{tri}
\end{sphinxVerbatim}
\begin{itemize}
\item {} 
\sphinxAtStartPar
ID, group\sphinxhyphen{}ID are documented in {\hyperref[\detokenize{fix::doc}]{\sphinxcrossref{\DUrole{doc}{fix}}}} command

\item {} 
\sphinxAtStartPar
nve/tri = style name of this fix command

\end{itemize}


\subsection{Examples}
\label{\detokenize{fix_nve_tri:examples}}
\begin{sphinxVerbatim}[commandchars=\\\{\}]
\PYG{k}{fix }\PYG{n+nv+nvIdentifier}{1}\PYG{+w}{ }\PYG{n+nv+nvIdentifier}{all}\PYG{+w}{ }\PYG{n}{nve}\PYG{o}{/}\PYG{n}{tri}
\end{sphinxVerbatim}


\subsection{Description}
\label{\detokenize{fix_nve_tri:description}}
\sphinxAtStartPar
Perform constant NVE integration to update position, velocity,
orientation, and angular momentum for triangular particles in the
group each timestep.  V is volume; E is energy.  This creates a system
trajectory consistent with the microcanonical ensemble.  See the
{\hyperref[\detokenize{Howto_spherical::doc}]{\sphinxcrossref{\DUrole{doc}{Howto spherical}}}} page for an overview of
using triangular particles.

\sphinxAtStartPar
This fix differs from the {\hyperref[\detokenize{fix_nve::doc}]{\sphinxcrossref{\DUrole{doc}{fix nve}}}} command, which
assumes point particles and only updates their position and velocity.


\subsection{Restart, fix\_modify, output, run start/stop, minimize info}
\label{\detokenize{fix_nve_tri:restart-fix-modify-output-run-start-stop-minimize-info}}
\sphinxAtStartPar
No information about this fix is written to {\hyperref[\detokenize{restart::doc}]{\sphinxcrossref{\DUrole{doc}{binary restart files}}}}.  None of the {\hyperref[\detokenize{fix_modify::doc}]{\sphinxcrossref{\DUrole{doc}{fix\_modify}}}} options
are relevant to this fix.  No global or per\sphinxhyphen{}atom quantities are stored
by this fix for access by various {\hyperref[\detokenize{Howto_output::doc}]{\sphinxcrossref{\DUrole{doc}{output commands}}}}.
No parameter of this fix can be used with the \sphinxstyleemphasis{start/stop} keywords of
the {\hyperref[\detokenize{run::doc}]{\sphinxcrossref{\DUrole{doc}{run}}}} command.  This fix is not invoked during {\hyperref[\detokenize{minimize::doc}]{\sphinxcrossref{\DUrole{doc}{energy minimization}}}}.


\subsection{Restrictions}
\label{\detokenize{fix_nve_tri:restrictions}}
\sphinxAtStartPar
This fix is part of the ASPHERE package.  It is only enabled if LAMMPS
was built with that package.  See the {\hyperref[\detokenize{Build_package::doc}]{\sphinxcrossref{\DUrole{doc}{Build package}}}} page for more info.

\sphinxAtStartPar
This fix requires that particles be triangles as defined by the
{\hyperref[\detokenize{atom_style::doc}]{\sphinxcrossref{\DUrole{doc}{atom\_style tri}}}} command.


\subsection{Related commands}
\label{\detokenize{fix_nve_tri:related-commands}}
\sphinxAtStartPar
{\hyperref[\detokenize{fix_nve::doc}]{\sphinxcrossref{\DUrole{doc}{fix nve}}}}, {\hyperref[\detokenize{fix_nve_asphere::doc}]{\sphinxcrossref{\DUrole{doc}{fix nve/asphere}}}}


\subsection{Default}
\label{\detokenize{fix_nve_tri:default}}
\sphinxAtStartPar
none

\sphinxstepscope

\index{fix nvk@\spxentry{fix nvk}}\ignorespaces 

\section{fix nvk command}
\label{\detokenize{fix_nvk:fix-nvk-command}}\label{\detokenize{fix_nvk:index-0}}\label{\detokenize{fix_nvk::doc}}

\subsection{Syntax}
\label{\detokenize{fix_nvk:syntax}}
\begin{sphinxVerbatim}[commandchars=\\\{\}]
\PYG{k}{fix }\PYG{n+nv+nvIdentifier}{ID}\PYG{+w}{ }\PYG{n+nv+nvIdentifier}{group\PYGZhy{}ID}\PYG{+w}{ }\PYG{n}{nvk}
\end{sphinxVerbatim}
\begin{itemize}
\item {} 
\sphinxAtStartPar
ID, group\sphinxhyphen{}ID are documented in {\hyperref[\detokenize{fix::doc}]{\sphinxcrossref{\DUrole{doc}{fix}}}} command

\item {} 
\sphinxAtStartPar
nvk = style name of this fix command

\end{itemize}


\subsection{Examples}
\label{\detokenize{fix_nvk:examples}}
\begin{sphinxVerbatim}[commandchars=\\\{\}]
\PYG{k}{fix }\PYG{n+nv+nvIdentifier}{1}\PYG{+w}{ }\PYG{n+nv+nvIdentifier}{all}\PYG{+w}{ }\PYG{n}{nvk}
\end{sphinxVerbatim}


\subsection{Description}
\label{\detokenize{fix_nvk:description}}
\sphinxAtStartPar
Perform constant kinetic energy integration using the Gaussian
thermostat to update position and velocity for atoms in the group each
timestep.  V is volume; K is kinetic energy. This creates a system
trajectory consistent with the isokinetic ensemble.

\sphinxAtStartPar
The equations of motion used are those of Minary et al in
{\hyperref[\detokenize{fix_nvk:nvk-minary}]{\sphinxcrossref{\DUrole{std,std-ref}{(Minary)}}}}, a variant of those initially given by Zhang in
{\hyperref[\detokenize{fix_nvk:nvk-zhang}]{\sphinxcrossref{\DUrole{std,std-ref}{(Zhang)}}}}.

\sphinxAtStartPar
The kinetic energy will be held constant at its value given when fix
nvk is initiated. If a different kinetic energy is desired, the
{\hyperref[\detokenize{velocity::doc}]{\sphinxcrossref{\DUrole{doc}{velocity}}}} command should be used to change the kinetic
energy prior to this fix.


\bigskip\hrule\bigskip



\subsection{Restart, fix\_modify, output, run start/stop, minimize info}
\label{\detokenize{fix_nvk:restart-fix-modify-output-run-start-stop-minimize-info}}
\sphinxAtStartPar
No information about this fix is written to {\hyperref[\detokenize{restart::doc}]{\sphinxcrossref{\DUrole{doc}{binary restart files}}}}.  None of the {\hyperref[\detokenize{fix_modify::doc}]{\sphinxcrossref{\DUrole{doc}{fix\_modify}}}} options
are relevant to this fix.  No global or per\sphinxhyphen{}atom quantities are stored
by this fix for access by various {\hyperref[\detokenize{Howto_output::doc}]{\sphinxcrossref{\DUrole{doc}{output commands}}}}.
No parameter of this fix can be used with the \sphinxstyleemphasis{start/stop} keywords of
the {\hyperref[\detokenize{run::doc}]{\sphinxcrossref{\DUrole{doc}{run}}}} command.  This fix is not invoked during {\hyperref[\detokenize{minimize::doc}]{\sphinxcrossref{\DUrole{doc}{energy minimization}}}}.


\subsection{Restrictions}
\label{\detokenize{fix_nvk:restrictions}}
\sphinxAtStartPar
The Gaussian thermostat only works when it is applied to all atoms in
the simulation box. Therefore, the group must be set to all.

\sphinxAtStartPar
This fix has not yet been implemented to work with the RESPA integrator.

\sphinxAtStartPar
This fix is part of the EXTRA\sphinxhyphen{}FIX package.  It is only enabled if
LAMMPS was built with that package.  See the
{\hyperref[\detokenize{Build_package::doc}]{\sphinxcrossref{\DUrole{doc}{Build package}}}} page for more info.


\subsection{Related commands}
\label{\detokenize{fix_nvk:related-commands}}
\sphinxAtStartPar
none


\subsection{Default}
\label{\detokenize{fix_nvk:default}}
\sphinxAtStartPar
none


\bigskip\hrule\bigskip

\phantomsection\label{\detokenize{fix_nvk:nvk-minary}}
\sphinxAtStartPar
\sphinxstylestrong{(Minary)} Minary, Martyna, and Tuckerman, J Chem Phys, 18, 2510 (2003).

\phantomsection\label{\detokenize{fix_nvk:nvk-zhang}}
\sphinxAtStartPar
\sphinxstylestrong{(Zhang)} Zhang, J Chem Phys, 106, 6102 (1997).

\sphinxstepscope

\index{fix nvt/asphere@\spxentry{fix nvt/asphere}}\ignorespaces 
\index{fix nvt/asphere/omp@\spxentry{fix nvt/asphere/omp}}\ignorespaces 

\section{fix nvt/asphere command}
\label{\detokenize{fix_nvt_asphere:fix-nvt-asphere-command}}\label{\detokenize{fix_nvt_asphere:index-1}}\label{\detokenize{fix_nvt_asphere:index-0}}\label{\detokenize{fix_nvt_asphere::doc}}
\sphinxAtStartPar
Accelerator Variants: \sphinxstyleemphasis{nvt/asphere/omp}


\subsection{Syntax}
\label{\detokenize{fix_nvt_asphere:syntax}}
\begin{sphinxVerbatim}[commandchars=\\\{\}]
\PYG{k}{fix }\PYG{n+nv+nvIdentifier}{ID}\PYG{+w}{ }\PYG{n+nv+nvIdentifier}{group\PYGZhy{}ID}\PYG{+w}{ }\PYG{n}{nvt}\PYG{o}{/}\PYG{n}{asphere}\PYG{+w}{ }\PYG{n}{keyword}\PYG{+w}{ }\PYG{n}{value}\PYG{+w}{ }\PYG{n}{...}
\end{sphinxVerbatim}
\begin{itemize}
\item {} 
\sphinxAtStartPar
ID, group\sphinxhyphen{}ID are documented in {\hyperref[\detokenize{fix::doc}]{\sphinxcrossref{\DUrole{doc}{fix}}}} command

\item {} 
\sphinxAtStartPar
nvt/asphere = style name of this fix command

\item {} 
\sphinxAtStartPar
additional thermostat related keyword/value pairs from the {\hyperref[\detokenize{fix_nh::doc}]{\sphinxcrossref{\DUrole{doc}{fix nvt}}}} command can be appended

\end{itemize}


\subsection{Examples}
\label{\detokenize{fix_nvt_asphere:examples}}
\begin{sphinxVerbatim}[commandchars=\\\{\}]
\PYG{k}{fix }\PYG{n+nv+nvIdentifier}{1}\PYG{+w}{ }\PYG{n+nv+nvIdentifier}{all}\PYG{+w}{ }\PYG{n}{nvt}\PYG{o}{/}\PYG{n}{asphere}\PYG{+w}{ }\PYG{n}{temp}\PYG{+w}{ }\PYG{l+m}{300.0}\PYG{+w}{ }\PYG{l+m}{300.0}\PYG{+w}{ }\PYG{l+m}{100.0}
\PYG{k}{fix }\PYG{n+nv+nvIdentifier}{1}\PYG{+w}{ }\PYG{n+nv+nvIdentifier}{all}\PYG{+w}{ }\PYG{n}{nvt}\PYG{o}{/}\PYG{n}{asphere}\PYG{+w}{ }\PYG{n}{temp}\PYG{+w}{ }\PYG{l+m}{300.0}\PYG{+w}{ }\PYG{l+m}{300.0}\PYG{+w}{ }\PYG{l+m}{100.0}\PYG{+w}{ }\PYG{n}{drag}\PYG{+w}{ }\PYG{l+m}{0.2}
\end{sphinxVerbatim}


\subsection{Description}
\label{\detokenize{fix_nvt_asphere:description}}
\sphinxAtStartPar
Perform constant NVT integration to update position, velocity,
orientation, and angular velocity each timestep for aspherical or
ellipsoidal particles in the group using a Nose/Hoover temperature
thermostat.  V is volume; T is temperature.  This creates a system
trajectory consistent with the canonical ensemble.

\sphinxAtStartPar
This fix differs from the {\hyperref[\detokenize{fix_nh::doc}]{\sphinxcrossref{\DUrole{doc}{fix nvt}}}} command, which
assumes point particles and only updates their position and velocity.

\sphinxAtStartPar
The thermostat is applied to both the translational and rotational
degrees of freedom for the aspherical particles, assuming a compute is
used which calculates a temperature that includes the rotational
degrees of freedom (see below).  The translational degrees of freedom
can also have a bias velocity removed from them before thermostatting
takes place; see the description below.

\sphinxAtStartPar
Additional parameters affecting the thermostat are specified by
keywords and values documented with the {\hyperref[\detokenize{fix_nh::doc}]{\sphinxcrossref{\DUrole{doc}{fix nvt}}}}
command.  See, for example, discussion of the \sphinxstyleemphasis{temp} and \sphinxstyleemphasis{drag}
keywords.

\sphinxAtStartPar
This fix computes a temperature each timestep.  To do this, the fix
creates its own compute of style “temp/asphere”, as if this command
had been issued:

\begin{sphinxVerbatim}[commandchars=\\\{\}]
\PYG{k}{compute }\PYG{n+nv+nvIdentifier}{fix\PYGZhy{}ID\PYGZus{}temp}\PYG{+w}{ }\PYG{n+nv+nvIdentifier}{group\PYGZhy{}ID}\PYG{+w}{ }\PYG{n}{temp}\PYG{o}{/}\PYG{n}{asphere}
\end{sphinxVerbatim}

\sphinxAtStartPar
See the {\hyperref[\detokenize{compute_temp_asphere::doc}]{\sphinxcrossref{\DUrole{doc}{compute temp/asphere}}}} command for
details.  Note that the ID of the new compute is the fix\sphinxhyphen{}ID +
underscore + “temp”, and the group for the new compute is the same as
the fix group.

\sphinxAtStartPar
Note that this is NOT the compute used by thermodynamic output (see
the {\hyperref[\detokenize{thermo_style::doc}]{\sphinxcrossref{\DUrole{doc}{thermo\_style}}}} command) with ID = \sphinxstyleemphasis{thermo\_temp}.
This means you can change the attributes of this fix’s temperature
(e.g. its degrees\sphinxhyphen{}of\sphinxhyphen{}freedom) via the
{\hyperref[\detokenize{compute_modify::doc}]{\sphinxcrossref{\DUrole{doc}{compute\_modify}}}} command or print this temperature
during thermodynamic output via the {\hyperref[\detokenize{thermo_style::doc}]{\sphinxcrossref{\DUrole{doc}{thermo\_style custom}}}} command using the appropriate compute\sphinxhyphen{}ID.
It also means that changing attributes of \sphinxstyleemphasis{thermo\_temp} will have no
effect on this fix.

\sphinxAtStartPar
Like other fixes that perform thermostatting, this fix can be used
with {\hyperref[\detokenize{compute::doc}]{\sphinxcrossref{\DUrole{doc}{compute commands}}}} that remove a “bias” from the
atom velocities.  E.g. to apply the thermostat only to atoms within a
spatial {\hyperref[\detokenize{region::doc}]{\sphinxcrossref{\DUrole{doc}{region}}}}, or to remove the center\sphinxhyphen{}of\sphinxhyphen{}mass
velocity from a group of atoms, or to remove the x\sphinxhyphen{}component of
velocity from the calculation.

\sphinxAtStartPar
This is not done by default, but only if the {\hyperref[\detokenize{fix_modify::doc}]{\sphinxcrossref{\DUrole{doc}{fix\_modify}}}} command is used to assign a temperature compute to this
fix that includes such a bias term.  See the doc pages for individual
{\hyperref[\detokenize{compute::doc}]{\sphinxcrossref{\DUrole{doc}{compute temp commands}}}} to determine which ones include
a bias.  In this case, the thermostat works in the following manner:
bias is removed from each atom, thermostatting is performed on the
remaining thermal degrees of freedom, and the bias is added back in.


\bigskip\hrule\bigskip


\sphinxAtStartPar
Styles with a \sphinxstyleemphasis{gpu}, \sphinxstyleemphasis{intel}, \sphinxstyleemphasis{kk}, \sphinxstyleemphasis{omp}, or \sphinxstyleemphasis{opt} suffix are
functionally the same as the corresponding style without the suffix.
They have been optimized to run faster, depending on your available
hardware, as discussed on the {\hyperref[\detokenize{Speed_packages::doc}]{\sphinxcrossref{\DUrole{doc}{Accelerator packages}}}}
page.  The accelerated styles take the same arguments and should
produce the same results, except for round\sphinxhyphen{}off and precision issues.

\sphinxAtStartPar
These accelerated styles are part of the GPU, INTEL, KOKKOS,
OPENMP, and OPT packages, respectively.  They are only enabled if
LAMMPS was built with those packages.  See the {\hyperref[\detokenize{Build_package::doc}]{\sphinxcrossref{\DUrole{doc}{Build package}}}} page for more info.

\sphinxAtStartPar
You can specify the accelerated styles explicitly in your input script
by including their suffix, or you can use the {\hyperref[\detokenize{Run_options::doc}]{\sphinxcrossref{\DUrole{doc}{\sphinxhyphen{}suffix command\sphinxhyphen{}line switch}}}} when you invoke LAMMPS, or you can use the
{\hyperref[\detokenize{suffix::doc}]{\sphinxcrossref{\DUrole{doc}{suffix}}}} command in your input script.

\sphinxAtStartPar
See the {\hyperref[\detokenize{Speed_packages::doc}]{\sphinxcrossref{\DUrole{doc}{Accelerator packages}}}} page for more
instructions on how to use the accelerated styles effectively.


\subsection{Restart, fix\_modify, output, run start/stop, minimize info}
\label{\detokenize{fix_nvt_asphere:restart-fix-modify-output-run-start-stop-minimize-info}}
\sphinxAtStartPar
This fix writes the state of the Nose/Hoover thermostat to
{\hyperref[\detokenize{restart::doc}]{\sphinxcrossref{\DUrole{doc}{binary restart files}}}}.  See the {\hyperref[\detokenize{read_restart::doc}]{\sphinxcrossref{\DUrole{doc}{read\_restart}}}} command for info on how to re\sphinxhyphen{}specify a fix in an
input script that reads a restart file, so that the operation of the
fix continues in an uninterrupted fashion.

\sphinxAtStartPar
The {\hyperref[\detokenize{fix_modify::doc}]{\sphinxcrossref{\DUrole{doc}{fix\_modify}}}} \sphinxstyleemphasis{temp} option is supported by this
fix.  You can use it to assign a {\hyperref[\detokenize{compute::doc}]{\sphinxcrossref{\DUrole{doc}{compute}}}} you have
defined to this fix which will be used in its thermostatting
procedure.

\sphinxAtStartPar
The cumulative energy change in the system imposed by this fix is
included in the {\hyperref[\detokenize{thermo_style::doc}]{\sphinxcrossref{\DUrole{doc}{thermodynamic output}}}} keywords
\sphinxstyleemphasis{ecouple} and \sphinxstyleemphasis{econserve}.  See the {\hyperref[\detokenize{thermo_style::doc}]{\sphinxcrossref{\DUrole{doc}{thermo\_style}}}}
doc page for details.

\sphinxAtStartPar
This fix computes the same global scalar and global vector of
quantities as does the {\hyperref[\detokenize{fix_nh::doc}]{\sphinxcrossref{\DUrole{doc}{fix nvt}}}} command.

\sphinxAtStartPar
This fix can ramp its target temperature over multiple runs, using the
\sphinxstyleemphasis{start} and \sphinxstyleemphasis{stop} keywords of the {\hyperref[\detokenize{run::doc}]{\sphinxcrossref{\DUrole{doc}{run}}}} command.  See the
{\hyperref[\detokenize{run::doc}]{\sphinxcrossref{\DUrole{doc}{run}}}} command for details of how to do this.

\sphinxAtStartPar
This fix is not invoked during {\hyperref[\detokenize{minimize::doc}]{\sphinxcrossref{\DUrole{doc}{energy minimization}}}}.


\subsection{Restrictions}
\label{\detokenize{fix_nvt_asphere:restrictions}}
\sphinxAtStartPar
This fix is part of the ASPHERE package.  It is only enabled if LAMMPS
was built with that package.  See the {\hyperref[\detokenize{Build_package::doc}]{\sphinxcrossref{\DUrole{doc}{Build package}}}} page for more info.

\sphinxAtStartPar
This fix requires that atoms store torque and angular momentum and a
quaternion as defined by the {\hyperref[\detokenize{atom_style::doc}]{\sphinxcrossref{\DUrole{doc}{atom\_style ellipsoid}}}}
command.

\sphinxAtStartPar
All particles in the group must be finite\sphinxhyphen{}size.  They cannot be point
particles, but they can be aspherical or spherical as defined by their
shape attribute.


\subsection{Related commands}
\label{\detokenize{fix_nvt_asphere:related-commands}}
\sphinxAtStartPar
{\hyperref[\detokenize{fix_nh::doc}]{\sphinxcrossref{\DUrole{doc}{fix nvt}}}}, {\hyperref[\detokenize{fix_nve_asphere::doc}]{\sphinxcrossref{\DUrole{doc}{fix nve\_asphere}}}}, {\hyperref[\detokenize{fix_npt_asphere::doc}]{\sphinxcrossref{\DUrole{doc}{fix npt\_asphere}}}}, {\hyperref[\detokenize{fix_modify::doc}]{\sphinxcrossref{\DUrole{doc}{fix\_modify}}}}


\subsection{Default}
\label{\detokenize{fix_nvt_asphere:default}}
\sphinxAtStartPar
none

\sphinxstepscope

\index{fix nvt/body@\spxentry{fix nvt/body}}\ignorespaces 

\section{fix nvt/body command}
\label{\detokenize{fix_nvt_body:fix-nvt-body-command}}\label{\detokenize{fix_nvt_body:index-0}}\label{\detokenize{fix_nvt_body::doc}}

\subsection{Syntax}
\label{\detokenize{fix_nvt_body:syntax}}
\begin{sphinxVerbatim}[commandchars=\\\{\}]
\PYG{k}{fix }\PYG{n+nv+nvIdentifier}{ID}\PYG{+w}{ }\PYG{n+nv+nvIdentifier}{group\PYGZhy{}ID}\PYG{+w}{ }\PYG{n}{nvt}\PYG{o}{/}\PYG{n}{body}\PYG{+w}{ }\PYG{n}{keyword}\PYG{+w}{ }\PYG{n}{value}\PYG{+w}{ }\PYG{n}{...}
\end{sphinxVerbatim}
\begin{itemize}
\item {} 
\sphinxAtStartPar
ID, group\sphinxhyphen{}ID are documented in {\hyperref[\detokenize{fix::doc}]{\sphinxcrossref{\DUrole{doc}{fix}}}} command

\item {} 
\sphinxAtStartPar
nvt/body = style name of this fix command

\item {} 
\sphinxAtStartPar
additional thermostat related keyword/value pairs from the {\hyperref[\detokenize{fix_nh::doc}]{\sphinxcrossref{\DUrole{doc}{fix nvt}}}} command can be appended

\end{itemize}


\subsection{Examples}
\label{\detokenize{fix_nvt_body:examples}}
\begin{sphinxVerbatim}[commandchars=\\\{\}]
\PYG{k}{fix }\PYG{n+nv+nvIdentifier}{1}\PYG{+w}{ }\PYG{n+nv+nvIdentifier}{all}\PYG{+w}{ }\PYG{n}{nvt}\PYG{o}{/}\PYG{n}{body}\PYG{+w}{ }\PYG{n}{temp}\PYG{+w}{ }\PYG{l+m}{300.0}\PYG{+w}{ }\PYG{l+m}{300.0}\PYG{+w}{ }\PYG{l+m}{100.0}
\PYG{k}{fix }\PYG{n+nv+nvIdentifier}{1}\PYG{+w}{ }\PYG{n+nv+nvIdentifier}{all}\PYG{+w}{ }\PYG{n}{nvt}\PYG{o}{/}\PYG{n}{body}\PYG{+w}{ }\PYG{n}{temp}\PYG{+w}{ }\PYG{l+m}{300.0}\PYG{+w}{ }\PYG{l+m}{300.0}\PYG{+w}{ }\PYG{l+m}{100.0}\PYG{+w}{ }\PYG{n}{drag}\PYG{+w}{ }\PYG{l+m}{0.2}
\end{sphinxVerbatim}


\subsection{Description}
\label{\detokenize{fix_nvt_body:description}}
\sphinxAtStartPar
Perform constant NVT integration to update position, velocity,
orientation, and angular velocity each timestep for body
particles in the group using a Nose/Hoover temperature
thermostat.  V is volume; T is temperature.  This creates a system
trajectory consistent with the canonical ensemble.

\sphinxAtStartPar
This fix differs from the {\hyperref[\detokenize{fix_nh::doc}]{\sphinxcrossref{\DUrole{doc}{fix nvt}}}} command, which
assumes point particles and only updates their position and velocity.

\sphinxAtStartPar
The thermostat is applied to both the translational and rotational
degrees of freedom for the body particles, assuming a compute is
used which calculates a temperature that includes the rotational
degrees of freedom (see below).  The translational degrees of freedom
can also have a bias velocity removed from them before thermostatting
takes place; see the description below.

\sphinxAtStartPar
Additional parameters affecting the thermostat are specified by
keywords and values documented with the {\hyperref[\detokenize{fix_nh::doc}]{\sphinxcrossref{\DUrole{doc}{fix nvt}}}}
command.  See, for example, discussion of the \sphinxstyleemphasis{temp} and \sphinxstyleemphasis{drag}
keywords.

\sphinxAtStartPar
This fix computes a temperature each timestep.  To do this, the fix
creates its own compute of style “temp/body”, as if this command
had been issued:

\begin{sphinxVerbatim}[commandchars=\\\{\}]
\PYG{k}{compute }\PYG{n+nv+nvIdentifier}{fix\PYGZhy{}ID\PYGZus{}temp}\PYG{+w}{ }\PYG{n+nv+nvIdentifier}{group\PYGZhy{}ID}\PYG{+w}{ }\PYG{n}{temp}\PYG{o}{/}\PYG{n}{body}
\end{sphinxVerbatim}

\sphinxAtStartPar
See the {\hyperref[\detokenize{compute_temp_body::doc}]{\sphinxcrossref{\DUrole{doc}{compute temp/body}}}} command for
details.  Note that the ID of the new compute is the fix\sphinxhyphen{}ID +
underscore + “temp”, and the group for the new compute is the same as
the fix group.

\sphinxAtStartPar
Note that this is NOT the compute used by thermodynamic output (see
the {\hyperref[\detokenize{thermo_style::doc}]{\sphinxcrossref{\DUrole{doc}{thermo\_style}}}} command) with ID = \sphinxstyleemphasis{thermo\_temp}.
This means you can change the attributes of this fix’s temperature
(e.g. its degrees\sphinxhyphen{}of\sphinxhyphen{}freedom) via the
{\hyperref[\detokenize{compute_modify::doc}]{\sphinxcrossref{\DUrole{doc}{compute\_modify}}}} command or print this temperature
during thermodynamic output via the {\hyperref[\detokenize{thermo_style::doc}]{\sphinxcrossref{\DUrole{doc}{thermo\_style custom}}}} command using the appropriate compute\sphinxhyphen{}ID.
It also means that changing attributes of \sphinxstyleemphasis{thermo\_temp} will have no
effect on this fix.

\sphinxAtStartPar
Like other fixes that perform thermostatting, this fix can be used
with {\hyperref[\detokenize{compute::doc}]{\sphinxcrossref{\DUrole{doc}{compute commands}}}} that remove a “bias” from the
atom velocities.  E.g. to apply the thermostat only to atoms within a
spatial {\hyperref[\detokenize{region::doc}]{\sphinxcrossref{\DUrole{doc}{region}}}}, or to remove the center\sphinxhyphen{}of\sphinxhyphen{}mass
velocity from a group of atoms, or to remove the x\sphinxhyphen{}component of
velocity from the calculation.

\sphinxAtStartPar
This is not done by default, but only if the {\hyperref[\detokenize{fix_modify::doc}]{\sphinxcrossref{\DUrole{doc}{fix\_modify}}}} command is used to assign a temperature compute to this
fix that includes such a bias term.  See the doc pages for individual
{\hyperref[\detokenize{compute::doc}]{\sphinxcrossref{\DUrole{doc}{compute temp commands}}}} to determine which ones include
a bias.  In this case, the thermostat works in the following manner:
bias is removed from each atom, thermostatting is performed on the
remaining thermal degrees of freedom, and the bias is added back in.


\bigskip\hrule\bigskip


\sphinxAtStartPar
Styles with a \sphinxstyleemphasis{gpu}, \sphinxstyleemphasis{intel}, \sphinxstyleemphasis{kk}, \sphinxstyleemphasis{omp}, or \sphinxstyleemphasis{opt} suffix are
functionally the same as the corresponding style without the suffix.
They have been optimized to run faster, depending on your available
hardware, as discussed on the {\hyperref[\detokenize{Speed_packages::doc}]{\sphinxcrossref{\DUrole{doc}{Accelerator packages}}}}
page.  The accelerated styles take the same arguments and should
produce the same results, except for round\sphinxhyphen{}off and precision issues.

\sphinxAtStartPar
These accelerated styles are part of the GPU, INTEL, KOKKOS,
OPENMP, and OPT packages, respectively.  They are only enabled if
LAMMPS was built with those packages.  See the {\hyperref[\detokenize{Build_package::doc}]{\sphinxcrossref{\DUrole{doc}{Build package}}}} page for more info.

\sphinxAtStartPar
You can specify the accelerated styles explicitly in your input script
by including their suffix, or you can use the {\hyperref[\detokenize{Run_options::doc}]{\sphinxcrossref{\DUrole{doc}{\sphinxhyphen{}suffix command\sphinxhyphen{}line switch}}}} when you invoke LAMMPS, or you can use the
{\hyperref[\detokenize{suffix::doc}]{\sphinxcrossref{\DUrole{doc}{suffix}}}} command in your input script.

\sphinxAtStartPar
See the {\hyperref[\detokenize{Speed_packages::doc}]{\sphinxcrossref{\DUrole{doc}{Accelerator packages}}}} page for more
instructions on how to use the accelerated styles effectively.


\subsection{Restart, fix\_modify, output, run start/stop, minimize info}
\label{\detokenize{fix_nvt_body:restart-fix-modify-output-run-start-stop-minimize-info}}
\sphinxAtStartPar
This fix writes the state of the Nose/Hoover thermostat to
{\hyperref[\detokenize{restart::doc}]{\sphinxcrossref{\DUrole{doc}{binary restart files}}}}.  See the {\hyperref[\detokenize{read_restart::doc}]{\sphinxcrossref{\DUrole{doc}{read\_restart}}}} command for info on how to re\sphinxhyphen{}specify a fix in an
input script that reads a restart file, so that the operation of the
fix continues in an uninterrupted fashion.

\sphinxAtStartPar
The {\hyperref[\detokenize{fix_modify::doc}]{\sphinxcrossref{\DUrole{doc}{fix\_modify}}}} \sphinxstyleemphasis{temp} option is supported by this
fix.  You can use it to assign a {\hyperref[\detokenize{compute::doc}]{\sphinxcrossref{\DUrole{doc}{compute}}}} you have
defined to this fix which will be used in its thermostatting
procedure.

\sphinxAtStartPar
The cumulative energy change in the system imposed by this fix is
included in the {\hyperref[\detokenize{thermo_style::doc}]{\sphinxcrossref{\DUrole{doc}{thermodynamic output}}}} keywords
\sphinxstyleemphasis{ecouple} and \sphinxstyleemphasis{econserve}.  See the {\hyperref[\detokenize{thermo_style::doc}]{\sphinxcrossref{\DUrole{doc}{thermo\_style}}}}
doc page for details.

\sphinxAtStartPar
This fix computes the same global scalar and global vector of
quantities as does the {\hyperref[\detokenize{fix_nh::doc}]{\sphinxcrossref{\DUrole{doc}{fix nvt}}}} command.

\sphinxAtStartPar
This fix can ramp its target temperature over multiple runs, using the
\sphinxstyleemphasis{start} and \sphinxstyleemphasis{stop} keywords of the {\hyperref[\detokenize{run::doc}]{\sphinxcrossref{\DUrole{doc}{run}}}} command.  See the
{\hyperref[\detokenize{run::doc}]{\sphinxcrossref{\DUrole{doc}{run}}}} command for details of how to do this.

\sphinxAtStartPar
This fix is not invoked during {\hyperref[\detokenize{minimize::doc}]{\sphinxcrossref{\DUrole{doc}{energy minimization}}}}.


\subsection{Restrictions}
\label{\detokenize{fix_nvt_body:restrictions}}
\sphinxAtStartPar
This fix is part of the BODY package.  It is only enabled if LAMMPS
was built with that package.  See the {\hyperref[\detokenize{Build_package::doc}]{\sphinxcrossref{\DUrole{doc}{Build package}}}} page for more info.

\sphinxAtStartPar
This fix requires that atoms store torque and angular momentum and a
quaternion as defined by the {\hyperref[\detokenize{atom_style::doc}]{\sphinxcrossref{\DUrole{doc}{atom\_style body}}}}
command.


\subsection{Related commands}
\label{\detokenize{fix_nvt_body:related-commands}}
\sphinxAtStartPar
{\hyperref[\detokenize{fix_nh::doc}]{\sphinxcrossref{\DUrole{doc}{fix nvt}}}}, {\hyperref[\detokenize{fix_nve_body::doc}]{\sphinxcrossref{\DUrole{doc}{fix nve\_body}}}}, {\hyperref[\detokenize{fix_npt_body::doc}]{\sphinxcrossref{\DUrole{doc}{fix npt\_body}}}}, {\hyperref[\detokenize{fix_modify::doc}]{\sphinxcrossref{\DUrole{doc}{fix\_modify}}}}


\subsection{Default}
\label{\detokenize{fix_nvt_body:default}}
\sphinxAtStartPar
none

\sphinxstepscope

\index{fix nvt/manifold/rattle@\spxentry{fix nvt/manifold/rattle}}\ignorespaces 

\section{fix nvt/manifold/rattle command}
\label{\detokenize{fix_nvt_manifold_rattle:fix-nvt-manifold-rattle-command}}\label{\detokenize{fix_nvt_manifold_rattle:index-0}}\label{\detokenize{fix_nvt_manifold_rattle::doc}}

\subsection{Syntax}
\label{\detokenize{fix_nvt_manifold_rattle:syntax}}
\begin{sphinxVerbatim}[commandchars=\\\{\}]
\PYG{k}{fix }\PYG{n+nv+nvIdentifier}{ID}\PYG{+w}{ }\PYG{n+nv+nvIdentifier}{group\PYGZhy{}ID}\PYG{+w}{ }\PYG{n}{nvt}\PYG{o}{/}\PYG{n}{manifold}\PYG{o}{/}\PYG{n}{rattle}\PYG{+w}{ }\PYG{n}{tol}\PYG{+w}{ }\PYG{n}{maxit}\PYG{+w}{ }\PYG{n}{manifold}\PYG{+w}{ }\PYG{n}{manifold}\PYG{o}{\PYGZhy{}}\PYG{n}{args}\PYG{+w}{ }\PYG{n}{keyword}\PYG{+w}{ }\PYG{n}{value}\PYG{+w}{ }\PYG{n}{...}
\end{sphinxVerbatim}
\begin{itemize}
\item {} 
\sphinxAtStartPar
ID, group\sphinxhyphen{}ID are documented in {\hyperref[\detokenize{fix::doc}]{\sphinxcrossref{\DUrole{doc}{fix}}}} command

\item {} 
\sphinxAtStartPar
nvt/manifold/rattle = style name of this fix command

\item {} 
\sphinxAtStartPar
tol = tolerance to which Newton iteration must converge

\item {} 
\sphinxAtStartPar
maxit = maximum number of iterations to perform

\item {} 
\sphinxAtStartPar
manifold = name of the manifold

\item {} 
\sphinxAtStartPar
manifold\sphinxhyphen{}args = parameters for the manifold

\item {} 
\sphinxAtStartPar
one or more keyword/value pairs may be appended
\begin{sphinxalltt}
keyword = \sphinxstyleemphasis{temp} or \sphinxstyleemphasis{tchain} or \sphinxstyleemphasis{every}
  \sphinxstyleemphasis{temp} values = Tstart Tstop Tdamp
    Tstart, Tstop = external temperature at start/end of run
    Tdamp = temperature damping parameter (time units)
  \sphinxstyleemphasis{tchain} value = N
    N = length of thermostat chain (1 = single thermostat)
  \sphinxstyleemphasis{every} value = N
    N = print info about iteration every N steps. N = 0 means no output
\end{sphinxalltt}

\end{itemize}


\subsection{Examples}
\label{\detokenize{fix_nvt_manifold_rattle:examples}}
\begin{sphinxVerbatim}[commandchars=\\\{\}]
\PYG{k}{fix }\PYG{n+nv+nvIdentifier}{1}\PYG{+w}{ }\PYG{n+nv+nvIdentifier}{all}\PYG{+w}{ }\PYG{n}{nvt}\PYG{o}{/}\PYG{n}{manifold}\PYG{o}{/}\PYG{n}{rattle}\PYG{+w}{ }\PYG{l+m}{1e\PYGZhy{}4}\PYG{+w}{ }\PYG{l+m}{10}\PYG{+w}{ }\PYG{n}{cylinder}\PYG{+w}{ }\PYG{l+m}{3.0}\PYG{+w}{ }\PYG{n}{temp}\PYG{+w}{ }\PYG{l+m}{1.0}\PYG{+w}{ }\PYG{l+m}{1.0}\PYG{+w}{ }\PYG{l+m}{10.0}
\end{sphinxVerbatim}


\subsection{Description}
\label{\detokenize{fix_nvt_manifold_rattle:description}}
\sphinxAtStartPar
This fix combines the RATTLE\sphinxhyphen{}based {\hyperref[\detokenize{fix_nvt_manifold_rattle:andersen2}]{\sphinxcrossref{\DUrole{std,std-ref}{(Andersen)}}}} time
integrator of {\hyperref[\detokenize{fix_nve_manifold_rattle::doc}]{\sphinxcrossref{\DUrole{doc}{fix nve/manifold/rattle}}}}
{\hyperref[\detokenize{fix_nvt_manifold_rattle:paquay3}]{\sphinxcrossref{\DUrole{std,std-ref}{(Paquay)}}}} with a Nose\sphinxhyphen{}Hoover\sphinxhyphen{}chain thermostat to sample the
canonical ensemble of particles constrained to a curved surface
(manifold). This sampling does suffer from discretization bias of
O(dt).  For a list of currently supported manifolds and their
parameters, see the {\hyperref[\detokenize{Howto_manifold::doc}]{\sphinxcrossref{\DUrole{doc}{Howto manifold}}}} doc page.


\bigskip\hrule\bigskip



\subsection{Restart, fix\_modify, output, run start/stop, minimize info}
\label{\detokenize{fix_nvt_manifold_rattle:restart-fix-modify-output-run-start-stop-minimize-info}}
\sphinxAtStartPar
No information about this fix is written to {\hyperref[\detokenize{restart::doc}]{\sphinxcrossref{\DUrole{doc}{binary restart files}}}}.  None of the {\hyperref[\detokenize{fix_modify::doc}]{\sphinxcrossref{\DUrole{doc}{fix\_modify}}}} options
are relevant to this fix.  No global or per\sphinxhyphen{}atom quantities are stored
by this fix for access by various {\hyperref[\detokenize{Howto_output::doc}]{\sphinxcrossref{\DUrole{doc}{output commands}}}}.
No parameter of this fix can be used with the \sphinxstyleemphasis{start/stop} keywords of
the {\hyperref[\detokenize{run::doc}]{\sphinxcrossref{\DUrole{doc}{run}}}} command.  This fix is not invoked during {\hyperref[\detokenize{minimize::doc}]{\sphinxcrossref{\DUrole{doc}{energy minimization}}}}.


\bigskip\hrule\bigskip



\subsection{Restrictions}
\label{\detokenize{fix_nvt_manifold_rattle:restrictions}}
\sphinxAtStartPar
This fix is part of the MANIFOLD package. It is only enabled if
LAMMPS was built with that package.  See the {\hyperref[\detokenize{Build_package::doc}]{\sphinxcrossref{\DUrole{doc}{Build package}}}} page for more info.


\bigskip\hrule\bigskip



\subsection{Related commands}
\label{\detokenize{fix_nvt_manifold_rattle:related-commands}}
\sphinxAtStartPar
{\hyperref[\detokenize{fix_nvt_manifold_rattle::doc}]{\sphinxcrossref{\DUrole{doc}{fix nve/manifold/rattle}}}}, {\hyperref[\detokenize{fix_manifoldforce::doc}]{\sphinxcrossref{\DUrole{doc}{fix manifoldforce}}}} \sphinxstylestrong{Default:} every = 0


\bigskip\hrule\bigskip

\phantomsection\label{\detokenize{fix_nvt_manifold_rattle:andersen2}}
\sphinxAtStartPar
\sphinxstylestrong{(Andersen)} Andersen, J. Comp. Phys. 52, 24, (1983).

\phantomsection\label{\detokenize{fix_nvt_manifold_rattle:paquay3}}
\sphinxAtStartPar
\sphinxstylestrong{(Paquay)} Paquay and Kusters, Biophys. J., 110, 6, (2016).
preprint available at \sphinxhref{https://arxiv.org/abs/1411.3019/}{arXiv:1411.3019}.

\sphinxstepscope

\index{fix nvt/sllod@\spxentry{fix nvt/sllod}}\ignorespaces 
\index{fix nvt/sllod/intel@\spxentry{fix nvt/sllod/intel}}\ignorespaces 
\index{fix nvt/sllod/omp@\spxentry{fix nvt/sllod/omp}}\ignorespaces 
\index{fix nvt/sllod/kk@\spxentry{fix nvt/sllod/kk}}\ignorespaces 

\section{fix nvt/sllod command}
\label{\detokenize{fix_nvt_sllod:fix-nvt-sllod-command}}\label{\detokenize{fix_nvt_sllod:index-3}}\label{\detokenize{fix_nvt_sllod:index-2}}\label{\detokenize{fix_nvt_sllod:index-1}}\label{\detokenize{fix_nvt_sllod:index-0}}\label{\detokenize{fix_nvt_sllod::doc}}
\sphinxAtStartPar
Accelerator Variants: \sphinxstyleemphasis{nvt/sllod/intel}, \sphinxstyleemphasis{nvt/sllod/omp}, \sphinxstyleemphasis{nvt/sllod/kk}


\subsection{Syntax}
\label{\detokenize{fix_nvt_sllod:syntax}}
\begin{sphinxVerbatim}[commandchars=\\\{\}]
\PYG{k}{fix }\PYG{n+nv+nvIdentifier}{ID}\PYG{+w}{ }\PYG{n+nv+nvIdentifier}{group\PYGZhy{}ID}\PYG{+w}{ }\PYG{n}{nvt}\PYG{o}{/}\PYG{n}{sllod}\PYG{+w}{ }\PYG{n}{keyword}\PYG{+w}{ }\PYG{n}{value}\PYG{+w}{ }\PYG{n}{...}
\end{sphinxVerbatim}
\begin{itemize}
\item {} 
\sphinxAtStartPar
ID, group\sphinxhyphen{}ID are documented in {\hyperref[\detokenize{fix::doc}]{\sphinxcrossref{\DUrole{doc}{fix}}}} command

\item {} 
\sphinxAtStartPar
nvt/sllod = style name of this fix command

\item {} 
\sphinxAtStartPar
zero or more keyword/value pairs may be appended
\begin{sphinxalltt}
keyword = \sphinxstyleemphasis{psllod}
  \sphinxstyleemphasis{psllod} value = \sphinxstyleemphasis{no} or \sphinxstyleemphasis{yes} = use SLLOD or p\sphinxhyphen{}SLLOD variant, respectively
\end{sphinxalltt}

\item {} 
\sphinxAtStartPar
additional thermostat related keyword/value pairs from the {\hyperref[\detokenize{fix_nh::doc}]{\sphinxcrossref{\DUrole{doc}{fix nvt}}}}
command can be appended

\end{itemize}


\subsection{Examples}
\label{\detokenize{fix_nvt_sllod:examples}}
\begin{sphinxVerbatim}[commandchars=\\\{\}]
\PYG{k}{fix }\PYG{n+nv+nvIdentifier}{1}\PYG{+w}{ }\PYG{n+nv+nvIdentifier}{all}\PYG{+w}{ }\PYG{n}{nvt}\PYG{o}{/}\PYG{n}{sllod}\PYG{+w}{ }\PYG{n}{temp}\PYG{+w}{ }\PYG{l+m}{300.0}\PYG{+w}{ }\PYG{l+m}{300.0}\PYG{+w}{ }\PYG{l+m}{100.0}
\PYG{k}{fix }\PYG{n+nv+nvIdentifier}{1}\PYG{+w}{ }\PYG{n+nv+nvIdentifier}{all}\PYG{+w}{ }\PYG{n}{nvt}\PYG{o}{/}\PYG{n}{sllod}\PYG{+w}{ }\PYG{n}{temp}\PYG{+w}{ }\PYG{l+m}{300.0}\PYG{+w}{ }\PYG{l+m}{300.0}\PYG{+w}{ }\PYG{l+m}{100.0}\PYG{+w}{ }\PYG{n}{drag}\PYG{+w}{ }\PYG{l+m}{0.2}
\end{sphinxVerbatim}


\subsection{Description}
\label{\detokenize{fix_nvt_sllod:description}}
\sphinxAtStartPar
Perform constant NVT integration to update positions and velocities
each timestep for atoms in the group using a Nose/Hoover temperature
thermostat.  V is volume; T is temperature.  This creates a system
trajectory consistent with the canonical ensemble.

\sphinxAtStartPar
This thermostat is used for a simulation box that is changing size
and/or shape, for example in a non\sphinxhyphen{}equilibrium MD (NEMD) simulation.
The size/shape change is induced by use of the {\hyperref[\detokenize{fix_deform::doc}]{\sphinxcrossref{\DUrole{doc}{fix deform}}}} command, so each point in the simulation box can be
thought of as having a “streaming” velocity.  This position\sphinxhyphen{}dependent
streaming velocity is subtracted from each atom’s actual velocity to
yield a thermal velocity which is used for temperature computation and
thermostatting.  For example, if the box is being sheared in x,
relative to y, then points at the bottom of the box (low y) have a
small x velocity, while points at the top of the box (hi y) have a
large x velocity.  These velocities do not contribute to the thermal
“temperature” of the atom.

\begin{sphinxadmonition}{note}{Note:}
\sphinxAtStartPar
{\hyperref[\detokenize{fix_deform::doc}]{\sphinxcrossref{\DUrole{doc}{Fix deform}}}} has an option for remapping either
atom coordinates or velocities to the changing simulation box.  To use
fix nvt/sllod, fix deform should NOT remap atom positions, because fix
nvt/sllod adjusts the atom positions and velocities to create a
velocity profile that matches the changing box size/shape.  Fix deform
SHOULD remap atom velocities when atoms cross periodic boundaries
since that is consistent with maintaining the velocity profile created
by fix nvt/sllod.  LAMMPS will give an error if this setting is not
consistent.
\end{sphinxadmonition}

\sphinxAtStartPar
The SLLOD equations of motion, originally proposed by Hoover and Ladd
(see {\hyperref[\detokenize{fix_nvt_sllod:evans3}]{\sphinxcrossref{\DUrole{std,std-ref}{(Evans and Morriss)}}}}), were proven to be equivalent
to Newton’s equations of motion for shear flow by {\hyperref[\detokenize{fix_nvt_sllod:evans3}]{\sphinxcrossref{\DUrole{std,std-ref}{(Evans and
Morriss)}}}}. They were later shown to generate the desired
velocity gradient and the correct production of work by stresses for all
forms of homogeneous flow by {\hyperref[\detokenize{fix_nvt_sllod:daivis}]{\sphinxcrossref{\DUrole{std,std-ref}{(Daivis and Todd)}}}}.

\sphinxAtStartPar
\DUrole{versionmodified,changed}{Changed in version 8Feb2023.}

\sphinxAtStartPar
For the default (\sphinxstyleemphasis{psllod} = \sphinxstyleemphasis{no}), the LAMMPS implementation adheres to
the standard SLLOD equations of motion, as defined by {\hyperref[\detokenize{fix_nvt_sllod:evans3}]{\sphinxcrossref{\DUrole{std,std-ref}{(Evans and
Morriss)}}}}.  The option \sphinxstyleemphasis{psllod} = \sphinxstyleemphasis{yes} invokes the slightly
different SLLOD variant first introduced by {\hyperref[\detokenize{fix_nvt_sllod:tuckerman}]{\sphinxcrossref{\DUrole{std,std-ref}{(Tuckerman et al.)}}}} as g\sphinxhyphen{}SLLOD and later by {\hyperref[\detokenize{fix_nvt_sllod:edwards}]{\sphinxcrossref{\DUrole{std,std-ref}{(Edwards)}}}} as
p\sphinxhyphen{}SLLOD.  In all cases, the equations of motion are coupled to a
Nose/Hoover chain thermostat in a velocity Verlet formulation, closely
following the implementation used for the {\hyperref[\detokenize{fix_nh::doc}]{\sphinxcrossref{\DUrole{doc}{fix nvt}}}}
command.

\begin{sphinxadmonition}{note}{Note:}
\sphinxAtStartPar
A recent (2017) book by {\hyperref[\detokenize{fix_nvt_sllod:todd-sllod}]{\sphinxcrossref{\DUrole{std,std-ref}{(Todd and Daivis)}}}}
discusses use of the SLLOD method and non\sphinxhyphen{}equilibrium MD (NEMD)
thermostatting generally, for both simple and complex fluids,
e.g. molecular systems.  The latter can be tricky to do correctly.
\end{sphinxadmonition}

\sphinxAtStartPar
Additional parameters affecting the thermostat are specified by
keywords and values documented with the {\hyperref[\detokenize{fix_nh::doc}]{\sphinxcrossref{\DUrole{doc}{fix nvt}}}}
command.  See, for example, discussion of the \sphinxstyleemphasis{temp} and \sphinxstyleemphasis{drag}
keywords.

\sphinxAtStartPar
This fix computes a temperature each timestep.  To do this, the fix
creates its own compute of style “temp/deform”, as if this command had
been issued:

\begin{sphinxVerbatim}[commandchars=\\\{\}]
\PYG{k}{compute }\PYG{n+nv+nvIdentifier}{fix\PYGZhy{}ID\PYGZus{}temp}\PYG{+w}{ }\PYG{n+nv+nvIdentifier}{group\PYGZhy{}ID}\PYG{+w}{ }\PYG{n}{temp}\PYG{o}{/}\PYG{n}{deform}
\end{sphinxVerbatim}

\sphinxAtStartPar
See the {\hyperref[\detokenize{compute_temp_deform::doc}]{\sphinxcrossref{\DUrole{doc}{compute temp/deform}}}} command for
details.  Note that the ID of the new compute is the fix\sphinxhyphen{}ID +
underscore + “temp”, and the group for the new compute is the same as
the fix group.

\sphinxAtStartPar
Note that this is NOT the compute used by thermodynamic output (see
the {\hyperref[\detokenize{thermo_style::doc}]{\sphinxcrossref{\DUrole{doc}{thermo\_style}}}} command) with ID =
\sphinxstyleemphasis{thermo\_temp}.  This means you can change the attributes of this fix’s
temperature (e.g. its degrees\sphinxhyphen{}of\sphinxhyphen{}freedom) via the {\hyperref[\detokenize{compute_modify::doc}]{\sphinxcrossref{\DUrole{doc}{compute\_modify}}}} command or print this temperature during
thermodynamic output via the {\hyperref[\detokenize{thermo_style::doc}]{\sphinxcrossref{\DUrole{doc}{thermo\_style custom}}}}
command using the appropriate compute\sphinxhyphen{}ID.  It also means that changing
attributes of \sphinxstyleemphasis{thermo\_temp} will have no effect on this fix.

\sphinxAtStartPar
Like other fixes that perform thermostatting, this fix can be used
with {\hyperref[\detokenize{compute::doc}]{\sphinxcrossref{\DUrole{doc}{compute commands}}}} that remove a “bias” from the
atom velocities.  E.g. to apply the thermostat only to atoms within a
spatial {\hyperref[\detokenize{region::doc}]{\sphinxcrossref{\DUrole{doc}{region}}}}, or to remove the center\sphinxhyphen{}of\sphinxhyphen{}mass
velocity from a group of atoms, or to remove the x\sphinxhyphen{}component of
velocity from the calculation.

\sphinxAtStartPar
This is not done by default, but only if the {\hyperref[\detokenize{fix_modify::doc}]{\sphinxcrossref{\DUrole{doc}{fix\_modify}}}} command is used to assign a temperature compute to this
fix that includes such a bias term.  See the doc pages for individual
{\hyperref[\detokenize{compute::doc}]{\sphinxcrossref{\DUrole{doc}{compute temp commands}}}} to determine which ones include
a bias.  In this case, the thermostat works in the following manner:
bias is removed from each atom, thermostatting is performed on the
remaining thermal degrees of freedom, and the bias is added back in.


\bigskip\hrule\bigskip


\sphinxAtStartPar
Styles with a \sphinxstyleemphasis{gpu}, \sphinxstyleemphasis{intel}, \sphinxstyleemphasis{kk}, \sphinxstyleemphasis{omp}, or \sphinxstyleemphasis{opt} suffix are
functionally the same as the corresponding style without the suffix.
They have been optimized to run faster, depending on your available
hardware, as discussed on the {\hyperref[\detokenize{Speed_packages::doc}]{\sphinxcrossref{\DUrole{doc}{Accelerator packages}}}}
page.  The accelerated styles take the same arguments and should
produce the same results, except for round\sphinxhyphen{}off and precision issues.

\sphinxAtStartPar
These accelerated styles are part of the GPU, INTEL, KOKKOS,
OPENMP, and OPT packages, respectively.  They are only enabled if
LAMMPS was built with those packages.  See the {\hyperref[\detokenize{Build_package::doc}]{\sphinxcrossref{\DUrole{doc}{Build package}}}} page for more info.

\sphinxAtStartPar
You can specify the accelerated styles explicitly in your input script
by including their suffix, or you can use the {\hyperref[\detokenize{Run_options::doc}]{\sphinxcrossref{\DUrole{doc}{\sphinxhyphen{}suffix command\sphinxhyphen{}line switch}}}} when you invoke LAMMPS, or you can use the
{\hyperref[\detokenize{suffix::doc}]{\sphinxcrossref{\DUrole{doc}{suffix}}}} command in your input script.

\sphinxAtStartPar
See the {\hyperref[\detokenize{Speed_packages::doc}]{\sphinxcrossref{\DUrole{doc}{Accelerator packages}}}} page for more
instructions on how to use the accelerated styles effectively.


\subsection{Restart, fix\_modify, output, run start/stop, minimize info}
\label{\detokenize{fix_nvt_sllod:restart-fix-modify-output-run-start-stop-minimize-info}}
\sphinxAtStartPar
This fix writes the state of the Nose/Hoover thermostat to
{\hyperref[\detokenize{restart::doc}]{\sphinxcrossref{\DUrole{doc}{binary restart files}}}}.  See the {\hyperref[\detokenize{read_restart::doc}]{\sphinxcrossref{\DUrole{doc}{read\_restart}}}} command for info on how to re\sphinxhyphen{}specify a fix in an
input script that reads a restart file, so that the operation of the
fix continues in an uninterrupted fashion.

\sphinxAtStartPar
The {\hyperref[\detokenize{fix_modify::doc}]{\sphinxcrossref{\DUrole{doc}{fix\_modify}}}} \sphinxstyleemphasis{temp} option is supported by this
fix.  You can use it to assign a {\hyperref[\detokenize{compute::doc}]{\sphinxcrossref{\DUrole{doc}{compute}}}} you have
defined to this fix which will be used in its thermostatting
procedure.

\sphinxAtStartPar
The cumulative energy change in the system imposed by this fix is
included in the {\hyperref[\detokenize{thermo_style::doc}]{\sphinxcrossref{\DUrole{doc}{thermodynamic output}}}} keywords
\sphinxstyleemphasis{ecouple} and \sphinxstyleemphasis{econserve}.  See the {\hyperref[\detokenize{thermo_style::doc}]{\sphinxcrossref{\DUrole{doc}{thermo\_style}}}}
doc page for details.

\sphinxAtStartPar
This fix computes the same global scalar and global vector of
quantities as does the {\hyperref[\detokenize{fix_nh::doc}]{\sphinxcrossref{\DUrole{doc}{fix nvt}}}} command.

\sphinxAtStartPar
This fix can ramp its target temperature over multiple runs, using the
\sphinxstyleemphasis{start} and \sphinxstyleemphasis{stop} keywords of the {\hyperref[\detokenize{run::doc}]{\sphinxcrossref{\DUrole{doc}{run}}}} command.  See the
{\hyperref[\detokenize{run::doc}]{\sphinxcrossref{\DUrole{doc}{run}}}} command for details of how to do this.

\sphinxAtStartPar
This fix is not invoked during {\hyperref[\detokenize{minimize::doc}]{\sphinxcrossref{\DUrole{doc}{energy minimization}}}}.


\subsection{Restrictions}
\label{\detokenize{fix_nvt_sllod:restrictions}}
\sphinxAtStartPar
This fix works best without Nose\sphinxhyphen{}Hoover chain thermostats, i.e. using
\sphinxstyleemphasis{tchain} = 1.  Setting \sphinxstyleemphasis{tchain} to larger values can result in poor
equilibration.


\subsection{Related commands}
\label{\detokenize{fix_nvt_sllod:related-commands}}
\sphinxAtStartPar
{\hyperref[\detokenize{fix_nve::doc}]{\sphinxcrossref{\DUrole{doc}{fix nve}}}}, {\hyperref[\detokenize{fix_nh::doc}]{\sphinxcrossref{\DUrole{doc}{fix nvt}}}}, {\hyperref[\detokenize{fix_temp_rescale::doc}]{\sphinxcrossref{\DUrole{doc}{fix temp/rescale}}}}, {\hyperref[\detokenize{fix_langevin::doc}]{\sphinxcrossref{\DUrole{doc}{fix langevin}}}},
{\hyperref[\detokenize{fix_modify::doc}]{\sphinxcrossref{\DUrole{doc}{fix\_modify}}}}, {\hyperref[\detokenize{compute_temp_deform::doc}]{\sphinxcrossref{\DUrole{doc}{compute temp/deform}}}}


\subsection{Default}
\label{\detokenize{fix_nvt_sllod:default}}
\sphinxAtStartPar
Same as {\hyperref[\detokenize{fix_nh::doc}]{\sphinxcrossref{\DUrole{doc}{fix nvt}}}}, except \sphinxstyleemphasis{tchain} = 1, psllod = \sphinxstyleemphasis{no}.


\bigskip\hrule\bigskip

\phantomsection\label{\detokenize{fix_nvt_sllod:evans3}}
\sphinxAtStartPar
\sphinxstylestrong{(Evans and Morriss)} Evans and Morriss, Phys Rev A, 30, 1528 (1984).

\phantomsection\label{\detokenize{fix_nvt_sllod:daivis}}
\sphinxAtStartPar
\sphinxstylestrong{(Daivis and Todd)} Daivis and Todd, J Chem Phys, 124, 194103 (2006).

\phantomsection\label{\detokenize{fix_nvt_sllod:todd-sllod}}
\sphinxAtStartPar
\sphinxstylestrong{(Todd and Daivis)} Todd and Daivis, Nonequilibrium Molecular Dynamics (book),
Cambridge University Press, (2017) \sphinxurl{https://doi.org/10.1017/9781139017848}.

\phantomsection\label{\detokenize{fix_nvt_sllod:tuckerman}}
\sphinxAtStartPar
\sphinxstylestrong{(Tuckerman et al.)} Tuckerman, Mundy, Balasubramanian, and Klein, J Chem Phys 106, 5615 (1997).

\phantomsection\label{\detokenize{fix_nvt_sllod:edwards}}
\sphinxAtStartPar
\sphinxstylestrong{(Edwards)} Edwards, Baig, and Keffer, J Chem Phys 124, 194104 (2006).

\sphinxstepscope

\index{fix nvt/sllod/eff@\spxentry{fix nvt/sllod/eff}}\ignorespaces 

\section{fix nvt/sllod/eff command}
\label{\detokenize{fix_nvt_sllod_eff:fix-nvt-sllod-eff-command}}\label{\detokenize{fix_nvt_sllod_eff:index-0}}\label{\detokenize{fix_nvt_sllod_eff::doc}}

\subsection{Syntax}
\label{\detokenize{fix_nvt_sllod_eff:syntax}}
\begin{sphinxVerbatim}[commandchars=\\\{\}]
\PYG{k}{fix }\PYG{n+nv+nvIdentifier}{ID}\PYG{+w}{ }\PYG{n+nv+nvIdentifier}{group\PYGZhy{}ID}\PYG{+w}{ }\PYG{n}{nvt}\PYG{o}{/}\PYG{n}{sllod}\PYG{o}{/}\PYG{n}{eff}\PYG{+w}{ }\PYG{n}{keyword}\PYG{+w}{ }\PYG{n}{value}\PYG{+w}{ }\PYG{n}{...}
\end{sphinxVerbatim}
\begin{itemize}
\item {} 
\sphinxAtStartPar
ID, group\sphinxhyphen{}ID are documented in {\hyperref[\detokenize{fix::doc}]{\sphinxcrossref{\DUrole{doc}{fix}}}} command

\item {} 
\sphinxAtStartPar
nvt/sllod/eff = style name of this fix command

\item {} 
\sphinxAtStartPar
zero or more keyword/value pairs may be appended
\begin{sphinxalltt}
keyword = \sphinxstyleemphasis{psllod}
  \sphinxstyleemphasis{psllod} value = \sphinxstyleemphasis{no} or \sphinxstyleemphasis{yes} = use SLLOD or p\sphinxhyphen{}SLLOD variant, respectively
\end{sphinxalltt}

\item {} 
\sphinxAtStartPar
additional thermostat related keyword/value pairs from the {\hyperref[\detokenize{fix_nh_eff::doc}]{\sphinxcrossref{\DUrole{doc}{fix
nvt/eff}}}} command may be appended, too.

\end{itemize}


\subsection{Examples}
\label{\detokenize{fix_nvt_sllod_eff:examples}}
\begin{sphinxVerbatim}[commandchars=\\\{\}]
\PYG{k}{fix }\PYG{n+nv+nvIdentifier}{1}\PYG{+w}{ }\PYG{n+nv+nvIdentifier}{all}\PYG{+w}{ }\PYG{n}{nvt}\PYG{o}{/}\PYG{n}{sllod}\PYG{o}{/}\PYG{n}{eff}\PYG{+w}{ }\PYG{n}{temp}\PYG{+w}{ }\PYG{l+m}{300.0}\PYG{+w}{ }\PYG{l+m}{300.0}\PYG{+w}{ }\PYG{l+m}{0.1}
\PYG{k}{fix }\PYG{n+nv+nvIdentifier}{1}\PYG{+w}{ }\PYG{n+nv+nvIdentifier}{all}\PYG{+w}{ }\PYG{n}{nvt}\PYG{o}{/}\PYG{n}{sllod}\PYG{o}{/}\PYG{n}{eff}\PYG{+w}{ }\PYG{n}{temp}\PYG{+w}{ }\PYG{l+m}{300.0}\PYG{+w}{ }\PYG{l+m}{300.0}\PYG{+w}{ }\PYG{l+m}{0.1}\PYG{+w}{ }\PYG{n}{drag}\PYG{+w}{ }\PYG{l+m}{0.2}
\end{sphinxVerbatim}


\subsection{Description}
\label{\detokenize{fix_nvt_sllod_eff:description}}
\sphinxAtStartPar
Perform constant NVT integration to update positions and velocities each
timestep for nuclei and electrons in the group for the {\hyperref[\detokenize{pair_eff::doc}]{\sphinxcrossref{\DUrole{doc}{electron
force field}}}} model, using a Nose/Hoover temperature
thermostat.  V is volume; T is temperature.  This creates a system
trajectory consistent with the canonical ensemble.

\sphinxAtStartPar
The operation of this fix is exactly like that described by the
{\hyperref[\detokenize{fix_nvt_sllod::doc}]{\sphinxcrossref{\DUrole{doc}{fix nvt/sllod}}}} command, except that the radius and
radial velocity of electrons are also updated and thermostatted.
Likewise the temperature calculated by the fix, using the compute it
creates (as discussed in the {\hyperref[\detokenize{fix_nh::doc}]{\sphinxcrossref{\DUrole{doc}{fix nvt, npt, and nph}}}} doc
page), is performed with a {\hyperref[\detokenize{compute_temp_deform_eff::doc}]{\sphinxcrossref{\DUrole{doc}{compute temp/deform/eff}}}} command that includes the eFF contribution to
the temperature from the electron radial velocity.


\subsection{Restart, fix\_modify, output, run start/stop, minimize info}
\label{\detokenize{fix_nvt_sllod_eff:restart-fix-modify-output-run-start-stop-minimize-info}}
\sphinxAtStartPar
This fix writes the state of the Nose/Hoover thermostat to
{\hyperref[\detokenize{restart::doc}]{\sphinxcrossref{\DUrole{doc}{binary restart files}}}}.  See the {\hyperref[\detokenize{read_restart::doc}]{\sphinxcrossref{\DUrole{doc}{read\_restart}}}} command for info on how to re\sphinxhyphen{}specify a fix in an
input script that reads a restart file, so that the operation of the
fix continues in an uninterrupted fashion.

\sphinxAtStartPar
The {\hyperref[\detokenize{fix_modify::doc}]{\sphinxcrossref{\DUrole{doc}{fix\_modify}}}} \sphinxstyleemphasis{temp} option is supported by this
fix.  You can use it to assign a {\hyperref[\detokenize{compute::doc}]{\sphinxcrossref{\DUrole{doc}{compute}}}} you have
defined to this fix which will be used in its thermostatting
procedure.

\sphinxAtStartPar
The cumulative energy change in the system imposed by this fix is
included in the {\hyperref[\detokenize{thermo_style::doc}]{\sphinxcrossref{\DUrole{doc}{thermodynamic output}}}} keywords
\sphinxstyleemphasis{ecouple} and \sphinxstyleemphasis{econserve}.  See the {\hyperref[\detokenize{thermo_style::doc}]{\sphinxcrossref{\DUrole{doc}{thermo\_style}}}}
doc page for details.

\sphinxAtStartPar
This fix computes the same global scalar and global vector of
quantities as does the {\hyperref[\detokenize{fix_nh_eff::doc}]{\sphinxcrossref{\DUrole{doc}{fix nvt/eff}}}} command.

\sphinxAtStartPar
This fix can ramp its target temperature over multiple runs, using the
\sphinxstyleemphasis{start} and \sphinxstyleemphasis{stop} keywords of the {\hyperref[\detokenize{run::doc}]{\sphinxcrossref{\DUrole{doc}{run}}}} command.  See the
{\hyperref[\detokenize{run::doc}]{\sphinxcrossref{\DUrole{doc}{run}}}} command for details of how to do this.

\sphinxAtStartPar
This fix is not invoked during {\hyperref[\detokenize{minimize::doc}]{\sphinxcrossref{\DUrole{doc}{energy minimization}}}}.


\subsection{Restrictions}
\label{\detokenize{fix_nvt_sllod_eff:restrictions}}
\sphinxAtStartPar
This fix is part of the EFF package.  It is only enabled if
LAMMPS was built with that package.  See the {\hyperref[\detokenize{Build_package::doc}]{\sphinxcrossref{\DUrole{doc}{Build package}}}} page for more info.

\sphinxAtStartPar
This fix works best without Nose\sphinxhyphen{}Hoover chain thermostats, i.e. using
tchain = 1.  Setting tchain to larger values can result in poor
equilibration.


\subsection{Related commands}
\label{\detokenize{fix_nvt_sllod_eff:related-commands}}
\sphinxAtStartPar
{\hyperref[\detokenize{fix_nve_eff::doc}]{\sphinxcrossref{\DUrole{doc}{fix nve/eff}}}}, {\hyperref[\detokenize{fix_nh_eff::doc}]{\sphinxcrossref{\DUrole{doc}{fix nvt/eff}}}}, {\hyperref[\detokenize{fix_langevin_eff::doc}]{\sphinxcrossref{\DUrole{doc}{fix langevin/eff}}}}, {\hyperref[\detokenize{fix_nvt_sllod::doc}]{\sphinxcrossref{\DUrole{doc}{fix nvt/sllod}}}}, {\hyperref[\detokenize{fix_modify::doc}]{\sphinxcrossref{\DUrole{doc}{fix\_modify}}}}, {\hyperref[\detokenize{compute_temp_deform_eff::doc}]{\sphinxcrossref{\DUrole{doc}{compute temp/deform/eff}}}}


\subsection{Default}
\label{\detokenize{fix_nvt_sllod_eff:default}}
\sphinxAtStartPar
Same as {\hyperref[\detokenize{fix_nh_eff::doc}]{\sphinxcrossref{\DUrole{doc}{fix nvt/eff}}}}, except tchain = 1.


\bigskip\hrule\bigskip

\phantomsection\label{\detokenize{fix_nvt_sllod_eff:tuckerman2}}
\sphinxAtStartPar
\sphinxstylestrong{(Tuckerman)} Tuckerman, Mundy, Balasubramanian, Klein, J Chem Phys,
106, 5615 (1997).

\sphinxstepscope

\index{fix nvt/sphere@\spxentry{fix nvt/sphere}}\ignorespaces 
\index{fix nvt/sphere/omp@\spxentry{fix nvt/sphere/omp}}\ignorespaces 

\section{fix nvt/sphere command}
\label{\detokenize{fix_nvt_sphere:fix-nvt-sphere-command}}\label{\detokenize{fix_nvt_sphere:index-1}}\label{\detokenize{fix_nvt_sphere:index-0}}\label{\detokenize{fix_nvt_sphere::doc}}
\sphinxAtStartPar
Accelerator Variants: \sphinxstyleemphasis{nvt/sphere/omp}


\subsection{Syntax}
\label{\detokenize{fix_nvt_sphere:syntax}}
\begin{sphinxVerbatim}[commandchars=\\\{\}]
\PYG{k}{fix }\PYG{n+nv+nvIdentifier}{ID}\PYG{+w}{ }\PYG{n+nv+nvIdentifier}{group\PYGZhy{}ID}\PYG{+w}{ }\PYG{n}{nvt}\PYG{o}{/}\PYG{n}{sphere}\PYG{+w}{ }\PYG{n}{keyword}\PYG{+w}{ }\PYG{n}{value}\PYG{+w}{ }\PYG{n}{...}
\end{sphinxVerbatim}
\begin{itemize}
\item {} 
\sphinxAtStartPar
ID, group\sphinxhyphen{}ID are documented in {\hyperref[\detokenize{fix::doc}]{\sphinxcrossref{\DUrole{doc}{fix}}}} command

\item {} 
\sphinxAtStartPar
nvt/sphere = style name of this fix command

\item {} 
\sphinxAtStartPar
zero or more keyword/value pairs may be appended

\item {} 
\sphinxAtStartPar
keyword = \sphinxstyleemphasis{disc}
\begin{sphinxalltt}
\sphinxstyleemphasis{disc} value = none = treat particles as 2d discs, not spheres
\end{sphinxalltt}

\item {} 
\sphinxAtStartPar
NOTE: additional thermostat and dipole related keyword/value pairs from the {\hyperref[\detokenize{fix_nh::doc}]{\sphinxcrossref{\DUrole{doc}{fix nvt}}}} command can be appended

\end{itemize}


\subsection{Examples}
\label{\detokenize{fix_nvt_sphere:examples}}
\begin{sphinxVerbatim}[commandchars=\\\{\}]
\PYG{k}{fix }\PYG{n+nv+nvIdentifier}{1}\PYG{+w}{ }\PYG{n+nv+nvIdentifier}{all}\PYG{+w}{ }\PYG{n}{nvt}\PYG{o}{/}\PYG{n}{sphere}\PYG{+w}{ }\PYG{n}{temp}\PYG{+w}{ }\PYG{l+m}{300.0}\PYG{+w}{ }\PYG{l+m}{300.0}\PYG{+w}{ }\PYG{l+m}{100.0}
\PYG{k}{fix }\PYG{n+nv+nvIdentifier}{1}\PYG{+w}{ }\PYG{n+nv+nvIdentifier}{all}\PYG{+w}{ }\PYG{n}{nvt}\PYG{o}{/}\PYG{n}{sphere}\PYG{+w}{ }\PYG{n}{temp}\PYG{+w}{ }\PYG{l+m}{300.0}\PYG{+w}{ }\PYG{l+m}{300.0}\PYG{+w}{ }\PYG{l+m}{100.0}\PYG{+w}{ }\PYG{n}{disc}
\PYG{k}{fix }\PYG{n+nv+nvIdentifier}{1}\PYG{+w}{ }\PYG{n+nv+nvIdentifier}{all}\PYG{+w}{ }\PYG{n}{nvt}\PYG{o}{/}\PYG{n}{sphere}\PYG{+w}{ }\PYG{n}{temp}\PYG{+w}{ }\PYG{l+m}{300.0}\PYG{+w}{ }\PYG{l+m}{300.0}\PYG{+w}{ }\PYG{l+m}{100.0}\PYG{+w}{ }\PYG{n}{drag}\PYG{+w}{ }\PYG{l+m}{0.2}
\PYG{k}{fix }\PYG{n+nv+nvIdentifier}{1}\PYG{+w}{ }\PYG{n+nv+nvIdentifier}{all}\PYG{+w}{ }\PYG{n}{nvt}\PYG{o}{/}\PYG{n}{sphere}\PYG{+w}{ }\PYG{n}{temp}\PYG{+w}{ }\PYG{l+m}{300.0}\PYG{+w}{ }\PYG{l+m}{300.0}\PYG{+w}{ }\PYG{l+m}{100.0}\PYG{+w}{ }\PYG{n}{update}\PYG{+w}{ }\PYG{n}{dipole}
\end{sphinxVerbatim}


\subsection{Description}
\label{\detokenize{fix_nvt_sphere:description}}
\sphinxAtStartPar
Perform constant NVT integration to update position, velocity, and
angular velocity each timestep for finite\sphinxhyphen{}size spherical particles in
the group using a Nose/Hoover temperature thermostat.  V is volume; T
is temperature.  This creates a system trajectory consistent with the
canonical ensemble.

\sphinxAtStartPar
This fix differs from the {\hyperref[\detokenize{fix_nh::doc}]{\sphinxcrossref{\DUrole{doc}{fix nvt}}}} command, which
assumes point particles and only updates their position and velocity.

\sphinxAtStartPar
The thermostat is applied to both the translational and rotational
degrees of freedom for the spherical particles, assuming a compute is
used which calculates a temperature that includes the rotational
degrees of freedom (see below).  The translational degrees of freedom
can also have a bias velocity removed from them before thermostatting
takes place; see the description below.

\sphinxAtStartPar
If the \sphinxstyleemphasis{disc} keyword is used, then each particle is treated as a 2d
disc (circle) instead of as a sphere.  This is only possible for 2d
simulations, as defined by the {\hyperref[\detokenize{dimension::doc}]{\sphinxcrossref{\DUrole{doc}{dimension}}}} keyword.
The only difference between discs and spheres in this context is their
moment of inertia, as used in the time integration.

\sphinxAtStartPar
Additional parameters affecting the thermostat are specified by
keywords and values documented with the {\hyperref[\detokenize{fix_nh::doc}]{\sphinxcrossref{\DUrole{doc}{fix nvt}}}}
command.  See, for example, discussion of the \sphinxstyleemphasis{temp} and \sphinxstyleemphasis{drag}
keywords.

\sphinxAtStartPar
This fix computes a temperature each timestep.  To do this, the fix
creates its own compute of style “temp/sphere”, as if this command
had been issued:

\begin{sphinxVerbatim}[commandchars=\\\{\}]
\PYG{k}{compute }\PYG{n+nv+nvIdentifier}{fix\PYGZhy{}ID\PYGZus{}temp}\PYG{+w}{ }\PYG{n+nv+nvIdentifier}{group\PYGZhy{}ID}\PYG{+w}{ }\PYG{n}{temp}\PYG{o}{/}\PYG{n}{sphere}
\end{sphinxVerbatim}

\sphinxAtStartPar
See the {\hyperref[\detokenize{compute_temp_sphere::doc}]{\sphinxcrossref{\DUrole{doc}{compute temp/sphere}}}} command for
details.  Note that the ID of the new compute is the fix\sphinxhyphen{}ID +
underscore + “temp”, and the group for the new compute is the same as
the fix group.

\sphinxAtStartPar
Note that this is NOT the compute used by thermodynamic output (see
the {\hyperref[\detokenize{thermo_style::doc}]{\sphinxcrossref{\DUrole{doc}{thermo\_style}}}} command) with ID =
\sphinxstyleemphasis{thermo\_temp}.  This means you can change the attributes of this fix’s
temperature (e.g. its degrees\sphinxhyphen{}of\sphinxhyphen{}freedom) via the {\hyperref[\detokenize{compute_modify::doc}]{\sphinxcrossref{\DUrole{doc}{compute\_modify}}}} command or print this temperature during
thermodynamic output via the {\hyperref[\detokenize{thermo_style::doc}]{\sphinxcrossref{\DUrole{doc}{thermo\_style custom}}}}
command using the appropriate compute\sphinxhyphen{}ID.  It also means that changing
attributes of \sphinxstyleemphasis{thermo\_temp} will have no effect on this fix.

\sphinxAtStartPar
Like other fixes that perform thermostatting, this fix can be used
with {\hyperref[\detokenize{compute::doc}]{\sphinxcrossref{\DUrole{doc}{compute commands}}}} that remove a “bias” from the
atom velocities.  E.g. to apply the thermostat only to atoms within a
spatial {\hyperref[\detokenize{region::doc}]{\sphinxcrossref{\DUrole{doc}{region}}}}, or to remove the center\sphinxhyphen{}of\sphinxhyphen{}mass
velocity from a group of atoms, or to remove the x\sphinxhyphen{}component of
velocity from the calculation.

\sphinxAtStartPar
This is not done by default, but only if the {\hyperref[\detokenize{fix_modify::doc}]{\sphinxcrossref{\DUrole{doc}{fix\_modify}}}} command is used to assign a temperature compute to this
fix that includes such a bias term.  See the doc pages for individual
{\hyperref[\detokenize{compute::doc}]{\sphinxcrossref{\DUrole{doc}{compute temp commands}}}} to determine which ones include
a bias.  In this case, the thermostat works in the following manner:
bias is removed from each atom, thermostatting is performed on the
remaining thermal degrees of freedom, and the bias is added back in.


\bigskip\hrule\bigskip


\sphinxAtStartPar
Styles with a \sphinxstyleemphasis{gpu}, \sphinxstyleemphasis{intel}, \sphinxstyleemphasis{kk}, \sphinxstyleemphasis{omp}, or \sphinxstyleemphasis{opt} suffix are
functionally the same as the corresponding style without the suffix.
They have been optimized to run faster, depending on your available
hardware, as discussed on the {\hyperref[\detokenize{Speed_packages::doc}]{\sphinxcrossref{\DUrole{doc}{Accelerator packages}}}}
page.  The accelerated styles take the same arguments and should
produce the same results, except for round\sphinxhyphen{}off and precision issues.

\sphinxAtStartPar
These accelerated styles are part of the GPU, INTEL, KOKKOS,
OPENMP, and OPT packages, respectively.  They are only enabled if
LAMMPS was built with those packages.  See the {\hyperref[\detokenize{Build_package::doc}]{\sphinxcrossref{\DUrole{doc}{Build package}}}} page for more info.

\sphinxAtStartPar
You can specify the accelerated styles explicitly in your input script
by including their suffix, or you can use the {\hyperref[\detokenize{Run_options::doc}]{\sphinxcrossref{\DUrole{doc}{\sphinxhyphen{}suffix command\sphinxhyphen{}line switch}}}} when you invoke LAMMPS, or you can use the
{\hyperref[\detokenize{suffix::doc}]{\sphinxcrossref{\DUrole{doc}{suffix}}}} command in your input script.

\sphinxAtStartPar
See the {\hyperref[\detokenize{Speed_packages::doc}]{\sphinxcrossref{\DUrole{doc}{Accelerator packages}}}} page for more
instructions on how to use the accelerated styles effectively.


\subsection{Restart, fix\_modify, output, run start/stop, minimize info}
\label{\detokenize{fix_nvt_sphere:restart-fix-modify-output-run-start-stop-minimize-info}}
\sphinxAtStartPar
This fix writes the state of the Nose/Hoover thermostat to
{\hyperref[\detokenize{restart::doc}]{\sphinxcrossref{\DUrole{doc}{binary restart files}}}}.  See the {\hyperref[\detokenize{read_restart::doc}]{\sphinxcrossref{\DUrole{doc}{read\_restart}}}} command for info on how to re\sphinxhyphen{}specify a fix in an
input script that reads a restart file, so that the operation of the
fix continues in an uninterrupted fashion.

\sphinxAtStartPar
The {\hyperref[\detokenize{fix_modify::doc}]{\sphinxcrossref{\DUrole{doc}{fix\_modify}}}} \sphinxstyleemphasis{temp} option is supported by this
fix.  You can use it to assign a {\hyperref[\detokenize{compute::doc}]{\sphinxcrossref{\DUrole{doc}{compute}}}} you have
defined to this fix which will be used in its thermostatting
procedure.

\sphinxAtStartPar
The cumulative energy change in the system imposed by this fix is
included in the {\hyperref[\detokenize{thermo_style::doc}]{\sphinxcrossref{\DUrole{doc}{thermodynamic output}}}} keywords
\sphinxstyleemphasis{ecouple} and \sphinxstyleemphasis{econserve}.  See the {\hyperref[\detokenize{thermo_style::doc}]{\sphinxcrossref{\DUrole{doc}{thermo\_style}}}}
doc page for details.

\sphinxAtStartPar
This fix computes the same global scalar and global vector of
quantities as does the {\hyperref[\detokenize{fix_nh::doc}]{\sphinxcrossref{\DUrole{doc}{fix nvt}}}} command.

\sphinxAtStartPar
This fix can ramp its target temperature over multiple runs, using the
\sphinxstyleemphasis{start} and \sphinxstyleemphasis{stop} keywords of the {\hyperref[\detokenize{run::doc}]{\sphinxcrossref{\DUrole{doc}{run}}}} command.  See the
{\hyperref[\detokenize{run::doc}]{\sphinxcrossref{\DUrole{doc}{run}}}} command for details of how to do this.

\sphinxAtStartPar
This fix is not invoked during {\hyperref[\detokenize{minimize::doc}]{\sphinxcrossref{\DUrole{doc}{energy minimization}}}}.


\subsection{Restrictions}
\label{\detokenize{fix_nvt_sphere:restrictions}}
\sphinxAtStartPar
This fix requires that atoms store torque and angular velocity (omega)
and a radius as defined by the {\hyperref[\detokenize{atom_style::doc}]{\sphinxcrossref{\DUrole{doc}{atom\_style sphere}}}}
command.

\sphinxAtStartPar
All particles in the group must be finite\sphinxhyphen{}size spheres.  They cannot
be point particles.

\sphinxAtStartPar
Use of the \sphinxstyleemphasis{disc} keyword is only allowed for 2d simulations, as
defined by the {\hyperref[\detokenize{dimension::doc}]{\sphinxcrossref{\DUrole{doc}{dimension}}}} keyword.


\subsection{Related commands}
\label{\detokenize{fix_nvt_sphere:related-commands}}\begin{description}
\sphinxlineitem{{\hyperref[\detokenize{fix_nh::doc}]{\sphinxcrossref{\DUrole{doc}{fix nvt}}}}, {\hyperref[\detokenize{fix_nve_sphere::doc}]{\sphinxcrossref{\DUrole{doc}{fix nve\_sphere}}}},}
\sphinxAtStartPar
{\hyperref[\detokenize{fix_nvt_asphere::doc}]{\sphinxcrossref{\DUrole{doc}{fix nvt\_asphere}}}}, {\hyperref[\detokenize{fix_npt_sphere::doc}]{\sphinxcrossref{\DUrole{doc}{fix npt\_sphere}}}}, {\hyperref[\detokenize{fix_modify::doc}]{\sphinxcrossref{\DUrole{doc}{fix\_modify}}}}

\end{description}


\subsection{Default}
\label{\detokenize{fix_nvt_sphere:default}}
\sphinxAtStartPar
none

\sphinxstepscope

\index{fix oneway@\spxentry{fix oneway}}\ignorespaces 

\section{fix oneway command}
\label{\detokenize{fix_oneway:fix-oneway-command}}\label{\detokenize{fix_oneway:index-0}}\label{\detokenize{fix_oneway::doc}}

\subsection{Syntax}
\label{\detokenize{fix_oneway:syntax}}
\begin{sphinxVerbatim}[commandchars=\\\{\}]
\PYG{k}{fix }\PYG{n+nv+nvIdentifier}{ID}\PYG{+w}{ }\PYG{n+nv+nvIdentifier}{group\PYGZhy{}ID}\PYG{+w}{ }\PYG{n}{oneway}\PYG{+w}{ }\PYG{n}{N}\PYG{+w}{ }\PYG{n}{region}\PYG{o}{\PYGZhy{}}\PYG{n}{ID}\PYG{+w}{ }\PYG{n}{direction}
\end{sphinxVerbatim}
\begin{itemize}
\item {} 
\sphinxAtStartPar
ID, group\sphinxhyphen{}ID are documented in {\hyperref[\detokenize{fix::doc}]{\sphinxcrossref{\DUrole{doc}{fix}}}} command

\item {} 
\sphinxAtStartPar
oneway = style name of this fix command

\item {} 
\sphinxAtStartPar
N = apply this fix every this many timesteps

\item {} 
\sphinxAtStartPar
region\sphinxhyphen{}ID = ID of region where fix is active

\item {} 
\sphinxAtStartPar
direction = \sphinxstyleemphasis{x} or \sphinxstyleemphasis{\sphinxhyphen{}x} or \sphinxstyleemphasis{y} or \sphinxstyleemphasis{\sphinxhyphen{}y} or \sphinxstyleemphasis{z} or \sphinxstyleemphasis{\sphinxhyphen{}z} = coordinate and direction of the oneway constraint

\end{itemize}


\subsection{Examples}
\label{\detokenize{fix_oneway:examples}}
\begin{sphinxVerbatim}[commandchars=\\\{\}]
\PYG{k}{fix }\PYG{n+nv+nvIdentifier}{1}\PYG{+w}{ }\PYG{n+nv+nvIdentifier}{ions}\PYG{+w}{ }\PYG{n}{oneway}\PYG{+w}{ }\PYG{l+m}{10}\PYG{+w}{ }\PYG{n}{semi}\PYG{+w}{ }\PYG{o}{\PYGZhy{}}\PYG{n}{x}
\PYG{k}{fix }\PYG{n+nv+nvIdentifier}{2}\PYG{+w}{ }\PYG{n+nv+nvIdentifier}{all}\PYG{+w}{ }\PYG{n}{oneway}\PYG{+w}{ }\PYG{l+m}{1}\PYG{+w}{ }\PYG{n}{left}\PYG{+w}{ }\PYG{o}{\PYGZhy{}}\PYG{n}{z}
\PYG{k}{fix }\PYG{n+nv+nvIdentifier}{3}\PYG{+w}{ }\PYG{n+nv+nvIdentifier}{all}\PYG{+w}{ }\PYG{n}{oneway}\PYG{+w}{ }\PYG{l+m}{1}\PYG{+w}{ }\PYG{n}{right}\PYG{+w}{ }\PYG{n}{z}
\end{sphinxVerbatim}


\subsection{Description}
\label{\detokenize{fix_oneway:description}}
\sphinxAtStartPar
Enforce that particles in the group and in a given region can only
move in one direction.  This is done by reversing a particle’s
velocity component, if it has the wrong sign in the specified
dimension.  The effect is that the particle moves in one direction
only.

\sphinxAtStartPar
This can be used, for example, as a simple model of a semi\sphinxhyphen{}permeable
membrane, or as an implementation of Maxwell’s demon.


\bigskip\hrule\bigskip



\subsection{Restart, fix\_modify, output, run start/stop, minimize info}
\label{\detokenize{fix_oneway:restart-fix-modify-output-run-start-stop-minimize-info}}
\sphinxAtStartPar
No information about this fix is written to {\hyperref[\detokenize{restart::doc}]{\sphinxcrossref{\DUrole{doc}{binary restart files}}}}.  None of the {\hyperref[\detokenize{fix_modify::doc}]{\sphinxcrossref{\DUrole{doc}{fix\_modify}}}} options
are relevant to this fix.  No global or per\sphinxhyphen{}atom quantities are stored
by this fix for access by various {\hyperref[\detokenize{Howto_output::doc}]{\sphinxcrossref{\DUrole{doc}{output commands}}}}.
No parameter of this fix can be used with the \sphinxstyleemphasis{start/stop} keywords of
the {\hyperref[\detokenize{run::doc}]{\sphinxcrossref{\DUrole{doc}{run}}}} command.  This fix is not invoked during {\hyperref[\detokenize{minimize::doc}]{\sphinxcrossref{\DUrole{doc}{energy minimization}}}}.


\subsection{Restrictions}
\label{\detokenize{fix_oneway:restrictions}}
\sphinxAtStartPar
This fix is part of the EXTRA\sphinxhyphen{}FIX package.  It is only enabled if LAMMPS
was built with that package.  See the {\hyperref[\detokenize{Build_package::doc}]{\sphinxcrossref{\DUrole{doc}{Build package}}}} page for more info.


\subsection{Related commands}
\label{\detokenize{fix_oneway:related-commands}}
\sphinxAtStartPar
{\hyperref[\detokenize{fix_wall_reflect::doc}]{\sphinxcrossref{\DUrole{doc}{fix wall/reflect}}}} command


\subsection{Default}
\label{\detokenize{fix_oneway:default}}
\sphinxAtStartPar
none

\sphinxstepscope

\index{fix orient/fcc@\spxentry{fix orient/fcc}}\ignorespaces 
\index{fix orient/bcc@\spxentry{fix orient/bcc}}\ignorespaces 

\section{fix orient/fcc command}
\label{\detokenize{fix_orient:fix-orient-fcc-command}}\label{\detokenize{fix_orient:index-1}}\label{\detokenize{fix_orient:index-0}}\label{\detokenize{fix_orient::doc}}

\section{fix orient/bcc command}
\label{\detokenize{fix_orient:fix-orient-bcc-command}}

\subsection{Syntax}
\label{\detokenize{fix_orient:syntax}}
\begin{sphinxVerbatim}[commandchars=\\\{\}]
\PYG{k}{fix }\PYG{n+nv+nvIdentifier}{ID}\PYG{+w}{ }\PYG{n+nv+nvIdentifier}{group\PYGZhy{}ID}\PYG{+w}{ }\PYG{n}{orient}\PYG{o}{/}\PYG{n}{fcc}\PYG{+w}{ }\PYG{n}{nstats}\PYG{+w}{ }\PYG{n}{dir}\PYG{+w}{ }\PYG{n}{alat}\PYG{+w}{ }\PYG{n}{dE}\PYG{+w}{ }\PYG{n}{cutlo}\PYG{+w}{ }\PYG{n}{cuthi}\PYG{+w}{ }\PYG{n}{file0}\PYG{+w}{ }\PYG{n}{file1}
\PYG{k}{fix }\PYG{n+nv+nvIdentifier}{ID}\PYG{+w}{ }\PYG{n+nv+nvIdentifier}{group\PYGZhy{}ID}\PYG{+w}{ }\PYG{n}{orient}\PYG{o}{/}\PYG{n}{bcc}\PYG{+w}{ }\PYG{n}{nstats}\PYG{+w}{ }\PYG{n}{dir}\PYG{+w}{ }\PYG{n}{alat}\PYG{+w}{ }\PYG{n}{dE}\PYG{+w}{ }\PYG{n}{cutlo}\PYG{+w}{ }\PYG{n}{cuthi}\PYG{+w}{ }\PYG{n}{file0}\PYG{+w}{ }\PYG{n}{file1}
\end{sphinxVerbatim}
\begin{itemize}
\item {} 
\sphinxAtStartPar
ID, group\sphinxhyphen{}ID are documented in {\hyperref[\detokenize{fix::doc}]{\sphinxcrossref{\DUrole{doc}{fix}}}} command

\item {} 
\sphinxAtStartPar
nstats = print stats every this many steps, 0 = never

\item {} 
\sphinxAtStartPar
dir = 0/1 for which crystal is used as reference

\item {} 
\sphinxAtStartPar
alat = fcc/bcc cubic lattice constant (distance units)

\item {} 
\sphinxAtStartPar
dE = energy added to each atom (energy units)

\item {} 
\sphinxAtStartPar
cutlo,cuthi = values between 0.0 and 1.0, cutlo \textless{} cuthi

\item {} 
\sphinxAtStartPar
file0,file1 = files that specify orientation of each grain

\end{itemize}


\subsection{Examples}
\label{\detokenize{fix_orient:examples}}
\begin{sphinxVerbatim}[commandchars=\\\{\}]
\PYG{k}{fix }\PYG{n+nv+nvIdentifier}{gb}\PYG{+w}{ }\PYG{n+nv+nvIdentifier}{all}\PYG{+w}{ }\PYG{n}{orient}\PYG{o}{/}\PYG{n}{fcc}\PYG{+w}{ }\PYG{l+m}{0}\PYG{+w}{ }\PYG{l+m}{1}\PYG{+w}{ }\PYG{l+m}{4.032008}\PYG{+w}{ }\PYG{l+m}{0.001}\PYG{+w}{ }\PYG{l+m}{0.25}\PYG{+w}{ }\PYG{l+m}{0.75}\PYG{+w}{ }\PYG{n}{xi.vec}\PYG{+w}{ }\PYG{n}{chi.vec}
\PYG{k}{fix }\PYG{n+nv+nvIdentifier}{gb}\PYG{+w}{ }\PYG{n+nv+nvIdentifier}{all}\PYG{+w}{ }\PYG{n}{orient}\PYG{o}{/}\PYG{n}{bcc}\PYG{+w}{ }\PYG{l+m}{0}\PYG{+w}{ }\PYG{l+m}{1}\PYG{+w}{ }\PYG{l+m}{2.882}\PYG{+w}{ }\PYG{l+m}{0.001}\PYG{+w}{ }\PYG{l+m}{0.25}\PYG{+w}{ }\PYG{l+m}{0.75}\PYG{+w}{ }\PYG{n}{ngb.left}\PYG{+w}{ }\PYG{n}{ngb.right}
\end{sphinxVerbatim}


\subsection{Description}
\label{\detokenize{fix_orient:description}}
\sphinxAtStartPar
The fix applies an orientation\sphinxhyphen{}dependent force to atoms near a planar
grain boundary which can be used to induce grain boundary migration
(in the direction perpendicular to the grain boundary plane).  The
motivation and explanation of this force and its application are
described in {\hyperref[\detokenize{fix_orient:janssens}]{\sphinxcrossref{\DUrole{std,std-ref}{(Janssens)}}}}. The adaptation to bcc crystals
is described in {\hyperref[\detokenize{fix_orient:wicaksono1}]{\sphinxcrossref{\DUrole{std,std-ref}{(Wicaksono1)}}}}. The computed force is only
applied to atoms in the fix group.

\sphinxAtStartPar
The basic idea is that atoms in one grain (on one side of the
boundary) have a potential energy dE added to them.  Atoms in the
other grain have 0.0 potential energy added.  Atoms near the boundary
(whose neighbor environment is intermediate between the two grain
orientations) have an energy between 0.0 and dE added.  This creates
an effective driving force to reduce the potential energy of atoms
near the boundary by pushing them towards one of the grain
orientations.  For dir = 1 and dE \textgreater{} 0, the boundary will thus move so
that the grain described by file0 grows and the grain described by
file1 shrinks.  Thus this fix is designed for simulations of two\sphinxhyphen{}grain
systems, either with one grain boundary and free surfaces parallel to
the boundary, or a system with periodic boundary conditions and two
equal and opposite grain boundaries.  In either case, the entire
system can displace during the simulation, and such motion should be
accounted for in measuring the grain boundary velocity.

\sphinxAtStartPar
The potential energy added to atom I is given by these formulas
\begin{equation*}
\begin{split} \xi_{i} = & \sum_{j=1}^{12} \left| \mathbf{r}_{j} - \mathbf{r}_{j}^{\rm I} \right| \qquad\qquad\left(1\right) \\
 \\
\xi_{\rm IJ} = & \sum_{j=1}^{12} \left| \mathbf{r}_{j}^{\rm J} - \mathbf{r}_{j}^{\rm I} \right| \qquad\qquad\left(2\right)\\
\\
\xi_{\rm low}  = & {\rm cutlo} \, \xi_{\rm IJ}  \qquad\qquad\qquad\left(3\right)\\
\xi_{\rm high}  = & {\rm cuthi} \, \xi_{\rm IJ} \qquad\qquad\qquad\left(4\right) \\
\\
\omega_{i} = & \frac{\pi}{2} \frac{\xi_{i} - \xi_{\rm low}}{\xi_{\rm high} - \xi_{\rm low}} \qquad\qquad\left(5\right)\\
\\
u_{i}  = & 0 \quad\quad\qquad\qquad\qquad \textrm{ for } \qquad \xi_{i} < \xi_{\rm low}\\
       = & {\rm dE}\,\frac{1 - \cos(2 \omega_{i})}{2}
 \qquad \mathrm{ for }\qquad \xi_{\rm low} < \xi_{i} < \xi_{\rm high}  \quad \left(6\right) \\
       = & {\rm dE} \quad\qquad\qquad\qquad\textrm{ for } \qquad \xi_{\rm high} < \xi_{i}\end{split}
\end{equation*}
\sphinxAtStartPar
which are fully explained in {\hyperref[\detokenize{fix_orient:janssens}]{\sphinxcrossref{\DUrole{std,std-ref}{(Janssens)}}}}.  For fcc crystals
this order parameter Xi for atom I in equation (1) is a sum over the
12 nearest neighbors of atom I. For bcc crystals it is the
corresponding sum of the 8 nearest neighbors. Rj is the vector from
atom I to its neighbor J, and RIj is a vector in the reference
(perfect) crystal.  That is, if dir = 0/1, then RIj is a vector to an
atom coord from file 0/1.  Equation (2) gives the expected value of
the order parameter XiIJ in the other grain.  Hi and lo cutoffs are
defined in equations (3) and (4), using the input parameters \sphinxstyleemphasis{cutlo}
and \sphinxstyleemphasis{cuthi} as thresholds to avoid adding grain boundary energy when
the deviation in the order parameter from 0 or 1 is small (e.g. due to
thermal fluctuations in a perfect crystal).  The added potential
energy Ui for atom I is given in equation (6) where it is interpolated
between 0 and dE using the two threshold Xi values and the Wi value of
equation (5).

\sphinxAtStartPar
The derivative of this energy expression gives the force on each atom
which thus depends on the orientation of its neighbors relative to the
2 grain orientations.  Only atoms near the grain boundary feel a net
force which tends to drive them to one of the two grain orientations.

\sphinxAtStartPar
In equation (1), the reference vector used for each neighbor is the
reference vector closest to the actual neighbor position.  This means
it is possible two different neighbors will use the same reference
vector.  In such cases, the atom in question is far from a perfect
orientation and will likely receive the full dE addition, so the
effect of duplicate reference vector usage is small.

\sphinxAtStartPar
The \sphinxstyleemphasis{dir} parameter determines which grain wants to grow at the
expense of the other.  A value of 0 means the first grain will shrink;
a value of 1 means it will grow.  This assumes that \sphinxstyleemphasis{dE} is positive.
The reverse will be true if \sphinxstyleemphasis{dE} is negative.

\sphinxAtStartPar
The \sphinxstyleemphasis{alat} parameter is the cubic lattice constant for the fcc or bcc
material and is only used to compute a cutoff distance of 1.57 * alat
/ sqrt(2) for finding the 12 or 8 nearest neighbors of each atom
(which should be valid for an fcc or bcc crystal).  A longer/shorter
cutoff can be imposed by adjusting \sphinxstyleemphasis{alat}.  If a particular atom has
less than 12 or 8 neighbors within the cutoff, the order parameter of
equation (1) is effectively multiplied by 12 or 8 divided by the
actual number of neighbors within the cutoff.

\sphinxAtStartPar
The \sphinxstyleemphasis{dE} parameter is the maximum amount of additional energy added to
each atom in the grain which wants to shrink.

\sphinxAtStartPar
The \sphinxstyleemphasis{cutlo} and \sphinxstyleemphasis{cuthi} parameters are used to reduce the force added
to bulk atoms in each grain far away from the boundary.  An atom in
the bulk surrounded by neighbors at the ideal grain orientation would
compute an order parameter of 0 or 1 and have no force added.
However, thermal vibrations in the solid will cause the order
parameters to be greater than 0 or less than 1.  The cutoff parameters
mask this effect, allowing forces to only be added to atoms with
order\sphinxhyphen{}parameters between the cutoff values.

\sphinxAtStartPar
\sphinxstyleemphasis{File0} and \sphinxstyleemphasis{file1} are filenames for the two grains which each
contain 6 vectors (6 lines with 3 values per line) which specify the
grain orientations.  Each vector is a displacement from a central atom
(0,0,0) to a nearest neighbor atom in an fcc lattice at the proper
orientation.  The vector lengths should all be identical since an fcc
lattice has a coordination number of 12.  Only 6 are listed due to
symmetry, so the list must include one from each pair of
equal\sphinxhyphen{}and\sphinxhyphen{}opposite neighbors.  A pair of orientation files for a
Sigma=5 tilt boundary are shown below. A tutorial that can help for
writing the orientation files is given in {\hyperref[\detokenize{fix_orient:wicaksono2}]{\sphinxcrossref{\DUrole{std,std-ref}{(Wicaksono2)}}}}


\subsection{Restart, fix\_modify, output, run start/stop, minimize info}
\label{\detokenize{fix_orient:restart-fix-modify-output-run-start-stop-minimize-info}}
\sphinxAtStartPar
No information about this fix is written to {\hyperref[\detokenize{restart::doc}]{\sphinxcrossref{\DUrole{doc}{binary restart files}}}}.

\sphinxAtStartPar
The {\hyperref[\detokenize{fix_modify::doc}]{\sphinxcrossref{\DUrole{doc}{fix\_modify}}}} \sphinxstyleemphasis{energy} option is supported by
this fix to add the potential energy of atom interactions with the
grain boundary driving force to the global potential energy of the
system as part of {\hyperref[\detokenize{thermo_style::doc}]{\sphinxcrossref{\DUrole{doc}{thermodynamic output}}}}.  The
default setting for this fix is {\hyperref[\detokenize{fix_modify::doc}]{\sphinxcrossref{\DUrole{doc}{fix\_modify energy no}}}}.

\sphinxAtStartPar
The {\hyperref[\detokenize{fix_modify::doc}]{\sphinxcrossref{\DUrole{doc}{fix\_modify}}}} \sphinxstyleemphasis{respa} option is supported by
these fixes. This allows to set at which level of the {\hyperref[\detokenize{run_style::doc}]{\sphinxcrossref{\DUrole{doc}{r\sphinxhyphen{}RESPA}}}} integrator a fix is adding its forces. Default is the
outermost level.

\sphinxAtStartPar
This fix calculates a global scalar which can be accessed by various
{\hyperref[\detokenize{Howto_output::doc}]{\sphinxcrossref{\DUrole{doc}{output commands}}}}.  The scalar is the potential
energy change due to this fix.  The scalar value calculated by this
fix is “extensive”.

\sphinxAtStartPar
This fix also calculates a per\sphinxhyphen{}atom array which can be accessed by
various {\hyperref[\detokenize{Howto_output::doc}]{\sphinxcrossref{\DUrole{doc}{output commands}}}}.  The array stores the
order parameter Xi and normalized order parameter (0 to 1) for each
atom.  The per\sphinxhyphen{}atom values can be accessed on any timestep.

\sphinxAtStartPar
No parameter of this fix can be used with the \sphinxstyleemphasis{start/stop} keywords of
the {\hyperref[\detokenize{run::doc}]{\sphinxcrossref{\DUrole{doc}{run}}}} command.

\sphinxAtStartPar
This fix is not invoked during {\hyperref[\detokenize{minimize::doc}]{\sphinxcrossref{\DUrole{doc}{energy minimization}}}}.


\subsection{Restrictions}
\label{\detokenize{fix_orient:restrictions}}
\sphinxAtStartPar
These fixes are part of the ORIENT package.  They are only enabled if LAMMPS
was built with that package.  See the {\hyperref[\detokenize{Build_package::doc}]{\sphinxcrossref{\DUrole{doc}{Build package}}}} page for more info.

\sphinxAtStartPar
These fixes should only be used with fcc or bcc lattices.


\subsection{Related commands}
\label{\detokenize{fix_orient:related-commands}}
\sphinxAtStartPar
{\hyperref[\detokenize{fix_modify::doc}]{\sphinxcrossref{\DUrole{doc}{fix\_modify}}}}


\subsection{Default}
\label{\detokenize{fix_orient:default}}
\sphinxAtStartPar
none


\bigskip\hrule\bigskip

\phantomsection\label{\detokenize{fix_orient:janssens}}
\sphinxAtStartPar
\sphinxstylestrong{(Janssens)} Janssens, Olmsted, Holm, Foiles, Plimpton, Derlet, Nature
Materials, 5, 124\sphinxhyphen{}127 (2006).

\phantomsection\label{\detokenize{fix_orient:wicaksono1}}
\sphinxAtStartPar
\sphinxstylestrong{(Wicaksono1)} Wicaksono, Sinclair, Militzer, Computational Materials
Science, 117, 397\sphinxhyphen{}405 (2016).

\phantomsection\label{\detokenize{fix_orient:wicaksono2}}
\sphinxAtStartPar
\sphinxstylestrong{(Wicaksono2)} Wicaksono, figshare,
\sphinxurl{https://doi.org/10.6084/m9.figshare.1488628.v1} (2015).


\bigskip\hrule\bigskip


\sphinxAtStartPar
For illustration purposes, here are example files that specify a
Sigma=5 \textless{}100\textgreater{} tilt boundary.  This is for a lattice constant of 3.5706
Angs.

\sphinxAtStartPar
file0:

\begin{sphinxVerbatim}[commandchars=\\\{\}]
 0.798410432046075    1.785300000000000    1.596820864092150
\PYGZhy{}0.798410432046075    1.785300000000000   \PYGZhy{}1.596820864092150
 2.395231296138225    0.000000000000000    0.798410432046075
 0.798410432046075    0.000000000000000   \PYGZhy{}2.395231296138225
 1.596820864092150    1.785300000000000   \PYGZhy{}0.798410432046075
 1.596820864092150   \PYGZhy{}1.785300000000000   \PYGZhy{}0.798410432046075
\end{sphinxVerbatim}

\sphinxAtStartPar
file1:

\begin{sphinxVerbatim}[commandchars=\\\{\}]
\PYGZhy{}0.798410432046075    1.785300000000000    1.596820864092150
 0.798410432046075    1.785300000000000   \PYGZhy{}1.596820864092150
 0.798410432046075    0.000000000000000    2.395231296138225
 2.395231296138225    0.000000000000000   \PYGZhy{}0.798410432046075
 1.596820864092150    1.785300000000000    0.798410432046075
 1.596820864092150   \PYGZhy{}1.785300000000000    0.798410432046075
\end{sphinxVerbatim}

\sphinxstepscope

\index{fix orient/eco@\spxentry{fix orient/eco}}\ignorespaces 

\section{fix orient/eco command}
\label{\detokenize{fix_orient_eco:fix-orient-eco-command}}\label{\detokenize{fix_orient_eco:index-0}}\label{\detokenize{fix_orient_eco::doc}}
\begin{sphinxVerbatim}[commandchars=\\\{\}]
fix ID group\PYGZhy{}ID orient/eco u0 eta cutoff orientationsFile
\end{sphinxVerbatim}
\begin{itemize}
\item {} 
\sphinxAtStartPar
ID, group\sphinxhyphen{}ID are documented in fix command

\item {} 
\sphinxAtStartPar
u0 = energy added to each atom (energy units)

\item {} 
\sphinxAtStartPar
eta = cutoff value (usually 0.25)

\item {} 
\sphinxAtStartPar
cutoff = cutoff radius for orientation parameter calculation

\item {} 
\sphinxAtStartPar
orientationsFile = file that specifies orientation of each grain

\end{itemize}


\subsection{Examples}
\label{\detokenize{fix_orient_eco:examples}}
\begin{sphinxVerbatim}[commandchars=\\\{\}]
\PYG{k}{fix }\PYG{n+nv+nvIdentifier}{gb}\PYG{+w}{ }\PYG{n+nv+nvIdentifier}{all}\PYG{+w}{ }\PYG{n}{orient}\PYG{o}{/}\PYG{n}{eco}\PYG{+w}{ }\PYG{l+m}{0.08}\PYG{+w}{ }\PYG{l+m}{0.25}\PYG{+w}{ }\PYG{l+m}{3.524}\PYG{+w}{ }\PYG{n}{sigma5.ori}
\end{sphinxVerbatim}


\subsection{Description}
\label{\detokenize{fix_orient_eco:description}}
\sphinxAtStartPar
The fix applies a synthetic driving force to a grain boundary which
can be used for the investigation of grain boundary motion. The
affiliation of atoms to either of the two grains forming the grain
boundary is determined from an orientation\sphinxhyphen{}dependent order parameter
as described in {\hyperref[\detokenize{fix_orient_eco:ulomek}]{\sphinxcrossref{\DUrole{std,std-ref}{(Ulomek)}}}}. The potential energy of
atoms is either increased by an amount of 0.5*\sphinxstyleemphasis{u0} or \sphinxhyphen{}0.5*\sphinxstyleemphasis{u0}
according to the orientation of the surrounding crystal. This creates
a potential energy gradient which pushes atoms near the grain boundary
to orient according to the energetically favorable grain
orientation. This fix is designed for applications in bicrystal system
with one grain boundary and open ends, or two opposite grain
boundaries in a periodic system. In either case, the entire system can
experience a displacement during the simulation which needs to be
accounted for in the evaluation of the grain boundary velocity. While
the basic method is described in {\hyperref[\detokenize{fix_orient_eco:ulomek}]{\sphinxcrossref{\DUrole{std,std-ref}{(Ulomek)}}}}, the
implementation follows the efficient implementation from
{\hyperref[\detokenize{fix_orient_eco:schratt}]{\sphinxcrossref{\DUrole{std,std-ref}{(Schratt \& Mohles)}}}}. The synthetic potential energy
added to an atom j is given by the following formulas
\begin{equation*}
\begin{split}w(|\vec{r}_{jk}|) = w_{jk} & = \left\{\begin{array}{lc} \frac{|\vec{r}_{jk}|^{4}}{r_{\mathrm{cut}}^{4}}
  -2\frac{|\vec{r}_{jk}|^{2}}{r_{\mathrm{cut}}^{2}}+1, & |\vec{r}_{jk}|<r_{\mathrm{cut}} \\
   0, & |\vec{r}_{jk}|\ge r_{\mathrm{cut}}
   \end{array}\right. \\
\chi_{j} & = \frac{1}{N}\sum_{l=1}^{3}\left\lbrack\left\vert\psi_{l}^{\mathrm{I}}(\vec{r}_{j})\right\vert^{2}-\left\vert\psi_{l}^{\mathrm{II}}(\vec{r}_{j})\right\vert^{2}\right\rbrack \\
\psi_{l}^{\mathrm{X}}(\vec{r}_{j}) & = \sum_{k\in\mathit{\Gamma}_{j}}w_{jk}\exp\left(\mathrm{i}\vec{r}_{jk}\cdot\vec{q}_{l}^{\mathrm{X}}\right) \\
u(\chi_{j}) & = \frac{u_{0}}{2}\left\{\begin{array}{lc}
1, & \chi_{j}\ge\eta\\
\sin\left(\frac{\pi\chi_{j}}{2\eta}\right), &  -\eta<\chi_{j}<\eta\\
-1, & \chi_{j}\le-\eta
\end{array}\right.\end{split}
\end{equation*}
\sphinxAtStartPar
which are fully explained in {\hyperref[\detokenize{fix_orient_eco:ulomek}]{\sphinxcrossref{\DUrole{std,std-ref}{(Ulomek)}}}}
and {\hyperref[\detokenize{fix_orient_eco:schratt}]{\sphinxcrossref{\DUrole{std,std-ref}{(Schratt \& Mohles)}}}}.

\sphinxAtStartPar
The force on each atom is the negative gradient of the synthetic
potential energy. It depends on the surrounding of this atom. An atom
far from the grain boundary does not experience a synthetic force as
its surrounding is that of an oriented single crystal and thermal
fluctuations are masked by the parameter \sphinxstyleemphasis{eta}. Near the grain
boundary however, the gradient is nonzero and synthetic force terms
are computed.  The orientationsFile specifies the perfect oriented
crystal basis vectors for the two adjoining crystals. The first three
lines (line=row vector) for the energetically penalized and the last
three lines for the energetically favored grain assuming \sphinxstyleemphasis{u0} is
positive. For negative \sphinxstyleemphasis{u0}, this is reversed. With the \sphinxstyleemphasis{cutoff}
parameter, the size of the region around each atom which is used in
the order parameter computation is defined. The cutoff must be smaller
than the interaction range of the MD potential. It should at least
include the nearest neighbor shell. For high temperatures or low angle
grain boundaries, it might be beneficial to increase the cutoff in
order to get a more precise identification of the atoms
surrounding. However, computation time will increase as more atoms are
considered in the order parameter and force computation.  It is also
worth noting that the cutoff radius must not exceed the communication
distance for ghost atoms in LAMMPS. With orientationsFile, the 6
oriented crystal basis vectors is specified. Each line of the input
file contains the three components of a primitive lattice vector
oriented according to the grain orientation in the simulation box. The
first (last) three lines correspond to the primitive lattice vectors
of the first (second) grain. An example for a
\(\Sigma\langle001\rangle\) mis\sphinxhyphen{}orientation is given at the end.

\sphinxAtStartPar
If no synthetic energy difference between the grains is created,
\(u0=0\), the force computation is omitted. In this case, still,
the order parameter of the driving force is computed and can be used
to track the grain boundary motion throughout the simulation.


\subsection{Restart, fix\_modify, output, run start/stop, minimize info}
\label{\detokenize{fix_orient_eco:restart-fix-modify-output-run-start-stop-minimize-info}}
\sphinxAtStartPar
No information about this fix is written to {\hyperref[\detokenize{restart::doc}]{\sphinxcrossref{\DUrole{doc}{binary restart
files}}}}.

\sphinxAtStartPar
The {\hyperref[\detokenize{fix_modify::doc}]{\sphinxcrossref{\DUrole{doc}{fix\_modify}}}} \sphinxstyleemphasis{energy} option is supported by
this fix to add the potential energy of atom interactions with the
grain boundary driving force to the global potential energy of the
system as part of {\hyperref[\detokenize{thermo_style::doc}]{\sphinxcrossref{\DUrole{doc}{thermodynamic output}}}}.  The
default setting for this fix is {\hyperref[\detokenize{fix_modify::doc}]{\sphinxcrossref{\DUrole{doc}{fix\_modify energy no}}}}.

\sphinxAtStartPar
This fix calculates a per\sphinxhyphen{}atom array with 2 columns, which can be
accessed by indices 1\sphinxhyphen{}1 by any command that uses per\sphinxhyphen{}atom values from
a fix as input.  See the {\hyperref[\detokenize{Howto_output::doc}]{\sphinxcrossref{\DUrole{doc}{Howto output}}}} doc page
for an overview of LAMMPS output options.

\sphinxAtStartPar
The first column is the order parameter for each atom; the second is
the thermal masking value for each atom.  Both are described above.

\sphinxAtStartPar
No parameter of this fix can be used with the start/stop keywords of
the run command. This fix is not invoked during energy minimization.


\subsection{Restrictions}
\label{\detokenize{fix_orient_eco:restrictions}}
\sphinxAtStartPar
This fix is part of the ORIENT package. It is only enabled if
LAMMPS was built with that package. See the {\hyperref[\detokenize{Build_package::doc}]{\sphinxcrossref{\DUrole{doc}{Build package}}}} page for more info.


\subsection{Related commands}
\label{\detokenize{fix_orient_eco:related-commands}}
\sphinxAtStartPar
{\hyperref[\detokenize{fix_modify::doc}]{\sphinxcrossref{\DUrole{doc}{fix\_modify}}}}

\sphinxAtStartPar
{\hyperref[\detokenize{fix_orient::doc}]{\sphinxcrossref{\DUrole{doc}{fix\_orient}}}}


\subsection{Default}
\label{\detokenize{fix_orient_eco:default}}
\sphinxAtStartPar
none


\bigskip\hrule\bigskip

\phantomsection\label{\detokenize{fix_orient_eco:ulomek}}
\sphinxAtStartPar
\sphinxstylestrong{(Ulomek)} Ulomek, Brien, Foiles, Mohles, Modelling Simul. Mater. Sci. Eng. 23 (2015) 025007

\phantomsection\label{\detokenize{fix_orient_eco:schratt}}
\sphinxAtStartPar
\sphinxstylestrong{(Schratt \& Mohles)} Schratt, Mohles. Comp. Mat. Sci. 182 (2020) 109774


\bigskip\hrule\bigskip


\sphinxAtStartPar
For illustration purposes, here is an example file that specifies a
\(\Sigma=5 \langle 001 \rangle\) tilt grain boundary. This is for a lattice constant of 3.52 Angstrom:

\begin{sphinxVerbatim}[commandchars=\\\{\}]
sigma5.ori:

1.671685  0.557228  1.76212
0.557228 \PYGZhy{}1.671685  1.76212
2.228913 \PYGZhy{}1.114456  0.00000
0.557228  1.671685  1.76212
1.671685 \PYGZhy{}0.557228  1.76212
2.228913  1.114456  0.00000
\end{sphinxVerbatim}

\sphinxstepscope

\index{fix pafi@\spxentry{fix pafi}}\ignorespaces 

\section{fix pafi command}
\label{\detokenize{fix_pafi:fix-pafi-command}}\label{\detokenize{fix_pafi:index-0}}\label{\detokenize{fix_pafi::doc}}

\subsection{Syntax}
\label{\detokenize{fix_pafi:syntax}}
\begin{sphinxVerbatim}[commandchars=\\\{\}]
\PYG{k}{fix }\PYG{n+nv+nvIdentifier}{ID}\PYG{+w}{ }\PYG{n+nv+nvIdentifier}{group\PYGZhy{}ID}\PYG{+w}{ }\PYG{n}{pafi}\PYG{+w}{ }\PYG{n}{compute}\PYG{o}{\PYGZhy{}}\PYG{n}{ID}\PYG{+w}{ }\PYG{n}{Temp}\PYG{+w}{ }\PYG{n}{Tdamp}\PYG{+w}{ }\PYG{n}{seed}\PYG{+w}{ }\PYG{n}{keyword}\PYG{+w}{ }\PYG{n}{values...}
\end{sphinxVerbatim}
\begin{itemize}
\item {} 
\sphinxAtStartPar
ID, group\sphinxhyphen{}ID are documented in {\hyperref[\detokenize{fix::doc}]{\sphinxcrossref{\DUrole{doc}{fix}}}} command

\item {} 
\sphinxAtStartPar
pafi = style name of this fix command

\item {} 
\sphinxAtStartPar
compute\sphinxhyphen{}ID = ID of a {\hyperref[\detokenize{compute_property_atom::doc}]{\sphinxcrossref{\DUrole{doc}{compute property/atom}}}} that holds data used by this fix

\item {} 
\sphinxAtStartPar
Temp = desired temperature (temperature units)

\item {} 
\sphinxAtStartPar
Tdamp = damping parameter (time units)

\item {} 
\sphinxAtStartPar
seed = random number seed to use for white noise (positive integer)

\item {} 
\sphinxAtStartPar
keyword = \sphinxstyleemphasis{overdamped} or \sphinxstyleemphasis{com}
\begin{sphinxalltt}
\sphinxstyleemphasis{overdamped} value = \sphinxstyleemphasis{yes} or \sphinxstyleemphasis{no} or 1 or 0
  \sphinxstyleemphasis{yes} or 1 = Brownian (overdamped) integration in hyperplane
  \sphinxstyleemphasis{no} or 0 = Langevin integration in hyperplane
\sphinxstyleemphasis{com} value = \sphinxstyleemphasis{yes} or \sphinxstyleemphasis{no} or 1 or 0
  \sphinxstyleemphasis{yes} or 1 = zero linear momentum, fixing center or mass (recommended)
  \sphinxstyleemphasis{no} or 0 = do not zero linear momentum, allowing center of mass drift
\end{sphinxalltt}

\end{itemize}


\subsection{Examples}
\label{\detokenize{fix_pafi:examples}}
\begin{sphinxVerbatim}[commandchars=\\\{\}]
\PYG{k}{compute }\PYG{n+nv+nvIdentifier}{pa}\PYG{+w}{ }\PYG{n+nv+nvIdentifier}{all}\PYG{+w}{ }\PYG{n}{property}\PYG{o}{/}\PYG{n}{atom}\PYG{+w}{ }\PYG{n}{d\PYGZus{}nx}\PYG{+w}{ }\PYG{n}{d\PYGZus{}ny}\PYG{+w}{ }\PYG{n}{d\PYGZus{}nz}\PYG{+w}{ }\PYG{n}{d\PYGZus{}dnx}\PYG{+w}{ }\PYG{n}{d\PYGZus{}dny}\PYG{+w}{ }\PYG{n}{d\PYGZus{}dnz}\PYG{+w}{ }\PYG{n}{d\PYGZus{}ddnx}\PYG{+w}{ }\PYG{n}{d\PYGZus{}ddny}\PYG{+w}{ }\PYG{n}{d\PYGZus{}ddnz}
\PYG{k}{run}\PYG{+w}{ }\PYG{l+m}{0}\PYG{+w}{ }\PYG{n}{post}\PYG{+w}{ }\PYG{n}{no}
\PYG{k}{fix }\PYG{n+nv+nvIdentifier}{hp}\PYG{+w}{ }\PYG{n+nv+nvIdentifier}{all}\PYG{+w}{ }\PYG{n}{pafi}\PYG{+w}{ }\PYG{n}{pa}\PYG{+w}{ }\PYG{l+m}{500.0}\PYG{+w}{ }\PYG{l+m}{0.01}\PYG{+w}{ }\PYG{l+m}{434}\PYG{+w}{ }\PYG{n}{overdamped}\PYG{+w}{ }\PYG{n}{yes}
\end{sphinxVerbatim}


\subsection{Description}
\label{\detokenize{fix_pafi:description}}
\sphinxAtStartPar
Perform Brownian or Langevin integration whilst constraining the system to lie
in some hyperplane, which is expected to be the tangent plane to some reference
pathway in a solid state system. The instantaneous value of a modified force
projection is also calculated, whose time integral can be shown to be equal to
the true free energy gradient along the minimum free energy path local to the reference pathway.
A detailed discussion of the projection technique can be found in {\hyperref[\detokenize{fix_pafi:swinburne}]{\sphinxcrossref{\DUrole{std,std-ref}{(Swinburne)}}}}.

\sphinxAtStartPar
This fix can be used with LAMMPS as demonstrated in examples/PACKAGES/pafi,
though it is primarily intended to be coupled with the PAFI C++ code, developed
at \sphinxurl{https://github.com/tomswinburne/pafi},
which distributes multiple LAMMPS workers in parallel to compute and collate
hyperplane\sphinxhyphen{}constrained averages, allowing the calculation of free energy barriers and pathways.

\sphinxAtStartPar
A {\hyperref[\detokenize{compute_property_atom::doc}]{\sphinxcrossref{\DUrole{doc}{compute property/atom}}}} must be provided with 9 fields per atom coordinate,
which in order are the x,y,z coordinates of a configuration on the reference path,
the x,y,z coordinates of the path tangent (derivative of path position with path coordinate)
and the x,y,z coordinates of the change in tangent (derivative of path tangent with path coordinate).

\sphinxAtStartPar
A 4\sphinxhyphen{}element vector is also calculated by this fix. The 4 components are the
modified projected force, its square, the expected projection of the minimum
free energy path tangent on the reference path tangent and the minimum image
distance between the current configuration and the reference configuration,
projected along the path tangent. This latter value should be essentially zero.

\begin{sphinxadmonition}{note}{Note:}
\sphinxAtStartPar
When com=yes/1, which is recommended, the provided tangent vector must also
have zero center of mass. This can be achieved by subtracting from each
coordinate of the path tangent the average x,y,z value. The PAFI C++ code
(see above) can generate these paths for use in LAMMPS.
\end{sphinxadmonition}

\begin{sphinxadmonition}{note}{Note:}
\sphinxAtStartPar
When overdamped=yes/1, the Tdamp parameter should be around 5\sphinxhyphen{}10 times smaller
than that used in typical Langevin integration.
See {\hyperref[\detokenize{fix_langevin::doc}]{\sphinxcrossref{\DUrole{doc}{fix langevin}}}} for typical values.
\end{sphinxadmonition}


\subsection{Restart, fix\_modify, output, run start/stop, minimize info}
\label{\detokenize{fix_pafi:restart-fix-modify-output-run-start-stop-minimize-info}}
\sphinxAtStartPar
No information about this fix is written to {\hyperref[\detokenize{restart::doc}]{\sphinxcrossref{\DUrole{doc}{binary restart files}}}}.
None of the {\hyperref[\detokenize{fix_modify::doc}]{\sphinxcrossref{\DUrole{doc}{fix\_modify}}}} options are relevant to this fix.

\sphinxAtStartPar
This fix produces a global vector each timestep which can be accessed by various {\hyperref[\detokenize{Howto_output::doc}]{\sphinxcrossref{\DUrole{doc}{output commands}}}}.


\subsection{Restrictions}
\label{\detokenize{fix_pafi:restrictions}}
\sphinxAtStartPar
This fix is part of the EXTRA\sphinxhyphen{}FIX package.  It is only enabled if
LAMMPS was built with that package.  See the {\hyperref[\detokenize{Build_package::doc}]{\sphinxcrossref{\DUrole{doc}{Build package}}}} page for more info.


\subsection{Default}
\label{\detokenize{fix_pafi:default}}
\sphinxAtStartPar
The option defaults are com = \sphinxstyleemphasis{yes}, overdamped = \sphinxstyleemphasis{no}


\bigskip\hrule\bigskip

\phantomsection\label{\detokenize{fix_pafi:swinburne}}
\sphinxAtStartPar
\sphinxstylestrong{(Swinburne)} Swinburne and Marinica, Physical Review Letters, 120, 1 (2018)

\sphinxstepscope

\index{fix pair@\spxentry{fix pair}}\ignorespaces 

\section{fix pair command}
\label{\detokenize{fix_pair:fix-pair-command}}\label{\detokenize{fix_pair:index-0}}\label{\detokenize{fix_pair::doc}}

\subsection{Syntax}
\label{\detokenize{fix_pair:syntax}}
\begin{sphinxVerbatim}[commandchars=\\\{\}]
\PYG{k}{fix }\PYG{n+nv+nvIdentifier}{ID}\PYG{+w}{ }\PYG{n+nv+nvIdentifier}{group\PYGZhy{}ID}\PYG{+w}{ }\PYG{n}{pair}\PYG{+w}{ }\PYG{n}{N}\PYG{+w}{ }\PYG{n}{pstyle}\PYG{+w}{ }\PYG{n}{name}\PYG{+w}{ }\PYG{n}{flag}\PYG{+w}{ }\PYG{n}{...}
\end{sphinxVerbatim}
\begin{itemize}
\item {} 
\sphinxAtStartPar
ID, group\sphinxhyphen{}ID are documented in {\hyperref[\detokenize{fix::doc}]{\sphinxcrossref{\DUrole{doc}{fix}}}} command

\item {} 
\sphinxAtStartPar
pair = style name of this fix command

\item {} 
\sphinxAtStartPar
N = invoke this fix once every N timesteps

\item {} 
\sphinxAtStartPar
pstyle = name of pair style to extract info from (e.g. eam)

\item {} 
\sphinxAtStartPar
one or more name/flag pairs can be listed

\item {} 
\sphinxAtStartPar
name = name of quantity the pair style allows extraction of

\item {} 
\sphinxAtStartPar
flag = 1 if pair style needs to be triggered to produce data for name, 0 if not

\end{itemize}


\subsection{Examples}
\label{\detokenize{fix_pair:examples}}
\begin{sphinxVerbatim}[commandchars=\\\{\}]
\PYG{k}{fix }\PYG{n+nv+nvIdentifier}{request}\PYG{+w}{ }\PYG{n+nv+nvIdentifier}{all}\PYG{+w}{ }\PYG{n}{pair}\PYG{+w}{ }\PYG{l+m}{100}\PYG{+w}{ }\PYG{n}{eam}\PYG{+w}{ }\PYG{n}{rho}\PYG{+w}{ }\PYG{l+m}{0}
\PYG{k}{fix }\PYG{n+nv+nvIdentifier}{request}\PYG{+w}{ }\PYG{n+nv+nvIdentifier}{all}\PYG{+w}{ }\PYG{n}{pair}\PYG{+w}{ }\PYG{l+m}{100}\PYG{+w}{ }\PYG{n}{amoeba}\PYG{+w}{ }\PYG{n}{uind}\PYG{+w}{ }\PYG{l+m}{0}\PYG{+w}{ }\PYG{n}{uinp}\PYG{+w}{ }\PYG{l+m}{0}
\end{sphinxVerbatim}


\subsection{Description}
\label{\detokenize{fix_pair:description}}
\sphinxAtStartPar
\DUrole{versionmodified,added}{Added in version 15Sep2022.}

\sphinxAtStartPar
Extract per\sphinxhyphen{}atom quantities from a pair style and store them in this
fix so they can be accessed by other LAMMPS commands, e.g. by a
{\hyperref[\detokenize{dump::doc}]{\sphinxcrossref{\DUrole{doc}{dump}}}} command or by another {\hyperref[\detokenize{fix::doc}]{\sphinxcrossref{\DUrole{doc}{fix}}}},
{\hyperref[\detokenize{compute::doc}]{\sphinxcrossref{\DUrole{doc}{compute}}}}, or {\hyperref[\detokenize{variable::doc}]{\sphinxcrossref{\DUrole{doc}{variable}}}} command.

\sphinxAtStartPar
These are example use cases:
\begin{itemize}
\item {} 
\sphinxAtStartPar
extract per\sphinxhyphen{}atom density from {\hyperref[\detokenize{pair_eam::doc}]{\sphinxcrossref{\DUrole{doc}{pair\_style eam}}}} to a dump file

\item {} 
\sphinxAtStartPar
extract induced dipoles from {\hyperref[\detokenize{pair_amoeba::doc}]{\sphinxcrossref{\DUrole{doc}{pair\_style amoeba}}}} to a dump file

\item {} 
\sphinxAtStartPar
extract accuracy metrics from a machine\sphinxhyphen{}learned potential to trigger output when
a condition is met (see the {\hyperref[\detokenize{dump_modify::doc}]{\sphinxcrossref{\DUrole{doc}{dump\_modify skip}}}} command)

\end{itemize}

\sphinxAtStartPar
The \sphinxstyleemphasis{N} argument determines how often the fix is invoked.

\sphinxAtStartPar
The \sphinxstyleemphasis{pstyle} argument is the name of the pair style.  It can be a
sub\sphinxhyphen{}style used in a {\hyperref[\detokenize{pair_hybrid::doc}]{\sphinxcrossref{\DUrole{doc}{pair\_style hybrid}}}} command.  If
there are multiple sub\sphinxhyphen{}styles using the same pair style, then \sphinxstyleemphasis{pstyle}
should be specified as “style:N”, where \sphinxstyleemphasis{N} is the number of the
instance of the pair style you wish monitor (e.g., the first or second).
For example, \sphinxstyleemphasis{pstyle} could be specified as “pace/extrapolation” or
“amoeba” or “eam:1” or “eam:2”.

\sphinxAtStartPar
One or more \sphinxstyleemphasis{name/flag} pairs of arguments follow.  Each \sphinxstyleemphasis{name} is a
per\sphinxhyphen{}atom quantity which the pair style must recognize as an extraction
request.  See the doc pages for individual {\hyperref[\detokenize{pair_style::doc}]{\sphinxcrossref{\DUrole{doc}{pair\_styles}}}} to see what fix pair requests (if any) they support.

\sphinxAtStartPar
The \sphinxstyleemphasis{flag} setting determines whether this fix will also trigger the
pair style to compute the named quantity so it can be extracted.  If the
quantity is always computed by the pair style, no trigger is needed;
specify \sphinxstyleemphasis{flag} = 0.  If the quantity is not always computed
(e.g. because it is expensive to calculate), then specify \sphinxstyleemphasis{flag} = 1.
This will trigger the quantity to be calculated only on timesteps it is
needed.  Again, see the doc pages for individual {\hyperref[\detokenize{pair_style::doc}]{\sphinxcrossref{\DUrole{doc}{pair\_styles}}}} to determine which fix pair requests (if any) need to be
triggered with a \sphinxstyleemphasis{flag} = 1 setting.

\sphinxAtStartPar
The per\sphinxhyphen{}atom data extracted from the pair style is stored by this fix
as either a per\sphinxhyphen{}atom vector or array.  If there is only one \sphinxstyleemphasis{name}
argument specified and the pair style computes a single value for each
atom, then this fix stores it as a per\sphinxhyphen{}atom vector.  Otherwise a
per\sphinxhyphen{}atom array is created, with its data in the order of the \sphinxstyleemphasis{name}
arguments.

\sphinxAtStartPar
For example, {\hyperref[\detokenize{pair_amoeba::doc}]{\sphinxcrossref{\DUrole{doc}{pair\_style amoeba}}}} allows extraction of
two named quantities: “uind” and “uinp”, both of which are 3\sphinxhyphen{}vectors for
each atom, i.e. dipole moments. In the example below a 6\sphinxhyphen{}column per\sphinxhyphen{}atom
array will be created.  Columns 1\sphinxhyphen{}3 will store the “uind” values;
columns 4\sphinxhyphen{}6 will store the “uinp” values.

\begin{sphinxVerbatim}[commandchars=\\\{\}]
\PYG{k}{pair\PYGZus{}style}\PYG{+w}{ }\PYG{n}{amoeba}
\PYG{k}{fix }\PYG{n+nv+nvIdentifier}{ex}\PYG{+w}{ }\PYG{n+nv+nvIdentifier}{all}\PYG{+w}{ }\PYG{n}{pair}\PYG{+w}{ }\PYG{l+m}{10}\PYG{+w}{ }\PYG{n}{amoeba}\PYG{+w}{ }\PYG{n}{uind}\PYG{+w}{ }\PYG{l+m}{0}\PYG{+w}{ }\PYG{n}{uinp}\PYG{+w}{ }\PYG{l+m}{0}
\end{sphinxVerbatim}


\subsection{Restart, fix\_modify, output, run start/stop, minimize info}
\label{\detokenize{fix_pair:restart-fix-modify-output-run-start-stop-minimize-info}}
\sphinxAtStartPar
No information about this fix is written to {\hyperref[\detokenize{restart::doc}]{\sphinxcrossref{\DUrole{doc}{binary restart files}}}}.  None of the {\hyperref[\detokenize{fix_modify::doc}]{\sphinxcrossref{\DUrole{doc}{fix\_modify}}}} options are
relevant to this fix.

\sphinxAtStartPar
As explained above, this fix produces a per\sphinxhyphen{}atom vector or array which
can be accessed by various {\hyperref[\detokenize{Howto_output::doc}]{\sphinxcrossref{\DUrole{doc}{output commands}}}}.  If
an array is produced, the number of columns is the sum of the number
of per\sphinxhyphen{}atom quantities produced by each \sphinxstyleemphasis{name} argument requested from
the pair style.


\subsection{Restrictions}
\label{\detokenize{fix_pair:restrictions}}
\sphinxAtStartPar
none


\subsection{Related commands}
\label{\detokenize{fix_pair:related-commands}}
\sphinxAtStartPar
{\hyperref[\detokenize{compute_pair::doc}]{\sphinxcrossref{\DUrole{doc}{compute pair}}}}


\subsection{Default}
\label{\detokenize{fix_pair:default}}
\sphinxAtStartPar
none

\sphinxstepscope

\index{fix phonon@\spxentry{fix phonon}}\ignorespaces 

\section{fix phonon command}
\label{\detokenize{fix_phonon:fix-phonon-command}}\label{\detokenize{fix_phonon:index-0}}\label{\detokenize{fix_phonon::doc}}

\subsection{Syntax}
\label{\detokenize{fix_phonon:syntax}}
\begin{sphinxVerbatim}[commandchars=\\\{\}]
\PYG{k}{fix }\PYG{n+nv+nvIdentifier}{ID}\PYG{+w}{ }\PYG{n+nv+nvIdentifier}{group\PYGZhy{}ID}\PYG{+w}{ }\PYG{n}{phonon}\PYG{+w}{ }\PYG{n}{N}\PYG{+w}{ }\PYG{n}{Noutput}\PYG{+w}{ }\PYG{n}{Nwait}\PYG{+w}{ }\PYG{n}{map\PYGZus{}file}\PYG{+w}{ }\PYG{n}{prefix}\PYG{+w}{ }\PYG{n}{keyword}\PYG{+w}{ }\PYG{n}{values}\PYG{+w}{ }\PYG{n}{...}
\end{sphinxVerbatim}
\begin{itemize}
\item {} 
\sphinxAtStartPar
ID, group\sphinxhyphen{}ID are documented in {\hyperref[\detokenize{fix::doc}]{\sphinxcrossref{\DUrole{doc}{fix}}}} command

\item {} 
\sphinxAtStartPar
phonon = style name of this fix command

\item {} 
\sphinxAtStartPar
N = measure the Green’s function every this many timesteps

\item {} 
\sphinxAtStartPar
Noutput = output the dynamical matrix every this many measurements

\item {} 
\sphinxAtStartPar
Nwait = wait this many timesteps before measuring

\item {} 
\sphinxAtStartPar
map\_file = \sphinxstyleemphasis{file} or \sphinxstyleemphasis{GAMMA}
\begin{sphinxalltt}
\sphinxstyleemphasis{file} is the file that contains the mapping info between atom ID and the lattice indices.
\end{sphinxalltt}
\begin{sphinxalltt}
\sphinxstyleemphasis{GAMMA} flags to treate the whole simulation box as a unit cell, so that the mapping
info can be generated internally. In this case, dynamical matrix at only the gamma\sphinxhyphen{}point
will/can be evaluated.
\end{sphinxalltt}

\item {} 
\sphinxAtStartPar
prefix = prefix for output files

\item {} 
\sphinxAtStartPar
one or none keyword/value pairs may be appended

\item {} 
\sphinxAtStartPar
keyword = \sphinxstyleemphasis{sysdim} or \sphinxstyleemphasis{nasr}
\begin{sphinxalltt}
\sphinxstyleemphasis{sysdim} value = d
  d = dimension of the system, usually the same as the MD model dimension
\sphinxstyleemphasis{nasr} value = n
  n = number of iterations to enforce the acoustic sum rule
\end{sphinxalltt}

\end{itemize}


\subsection{Examples}
\label{\detokenize{fix_phonon:examples}}
\begin{sphinxVerbatim}[commandchars=\\\{\}]
\PYG{k}{fix }\PYG{n+nv+nvIdentifier}{1}\PYG{+w}{ }\PYG{n+nv+nvIdentifier}{all}\PYG{+w}{ }\PYG{n}{phonon}\PYG{+w}{ }\PYG{l+m}{20}\PYG{+w}{ }\PYG{l+m}{5000}\PYG{+w}{ }\PYG{l+m}{200000}\PYG{+w}{ }\PYG{n}{map.in}\PYG{+w}{ }\PYG{n}{LJ1D}\PYG{+w}{ }\PYG{n}{sysdim}\PYG{+w}{ }\PYG{l+m}{1}
\PYG{k}{fix }\PYG{n+nv+nvIdentifier}{1}\PYG{+w}{ }\PYG{n+nv+nvIdentifier}{all}\PYG{+w}{ }\PYG{n}{phonon}\PYG{+w}{ }\PYG{l+m}{20}\PYG{+w}{ }\PYG{l+m}{5000}\PYG{+w}{ }\PYG{l+m}{200000}\PYG{+w}{ }\PYG{n}{map.in}\PYG{+w}{ }\PYG{n}{EAM3D}
\PYG{k}{fix }\PYG{n+nv+nvIdentifier}{1}\PYG{+w}{ }\PYG{n+nv+nvIdentifier}{all}\PYG{+w}{ }\PYG{n}{phonon}\PYG{+w}{ }\PYG{l+m}{10}\PYG{+w}{ }\PYG{l+m}{5000}\PYG{+w}{ }\PYG{l+m}{500000}\PYG{+w}{ }\PYG{n}{GAMMA}\PYG{+w}{ }\PYG{n}{EAM0D}\PYG{+w}{ }\PYG{n}{nasr}\PYG{+w}{ }\PYG{l+m}{100}
\end{sphinxVerbatim}


\subsection{Description}
\label{\detokenize{fix_phonon:description}}
\sphinxAtStartPar
Calculate the dynamical matrix from molecular dynamics simulations
based on fluctuation\sphinxhyphen{}dissipation theory for a group of atoms.

\sphinxAtStartPar
Consider a crystal with \(N\) unit cells in three dimensions labeled
\(l = (l_1, l_2, l_3)\) where \(l_i\) are integers.  Each unit cell is
defined by three linearly independent vectors \(\mathbf{a}_1\),
\(\mathbf{a}_2\), \(\mathbf{a}_3\) forming a parallelepiped,
containing \(K\) basis atoms labeled \(k\).

\sphinxAtStartPar
Based on fluctuation\sphinxhyphen{}dissipation theory, the force constant
coefficients of the system in reciprocal space are given by
({\hyperref[\detokenize{fix_phonon:campana}]{\sphinxcrossref{\DUrole{std,std-ref}{Campana}}}} , {\hyperref[\detokenize{fix_phonon:kong}]{\sphinxcrossref{\DUrole{std,std-ref}{Kong}}}})
\begin{equation*}
\begin{split}\mathbf{\Phi}_{k\alpha,k^\prime \beta}(\mathbf{q}) = k_B T \mathbf{G}^{-1}_{k\alpha,k^\prime \beta}(\mathbf{q})\end{split}
\end{equation*}
\sphinxAtStartPar
where \(\mathbf{G}\) is the Green’s functions coefficients given by
\begin{equation*}
\begin{split}\mathbf{G}_{k\alpha,k^\prime \beta}(\mathbf{q}) = \left< \mathbf{u}_{k\alpha}(\mathbf{q}) \bullet \mathbf{u}_{k^\prime \beta}^*(\mathbf{q}) \right>\end{split}
\end{equation*}
\sphinxAtStartPar
where \(\left< \ldots \right>\) denotes the ensemble average, and
\begin{equation*}
\begin{split}\mathbf{u}_{k\alpha}(\mathbf{q}) = \sum_l \mathbf{u}_{l k \alpha} \exp{(i\mathbf{qr}_l)}\end{split}
\end{equation*}
\sphinxAtStartPar
is the \(\alpha\) component of the atomic displacement for the \(k\)
th atom in the unit cell in reciprocal space at \(\mathbf{q}\). In
practice, the Green’s functions coefficients can also be measured
according to the following formula,
\begin{equation*}
\begin{split}\mathbf{G}_{k\alpha,k^\prime \beta}(\mathbf{q}) =
\left< \mathbf{R}_{k \alpha}(\mathbf{q}) \bullet \mathbf{R}^*_{k^\prime \beta}(\mathbf{q}) \right>
- \left<\mathbf{R}\right>_{k \alpha}(\mathbf{q}) \bullet \left<\mathbf{R}\right>^*_{k^\prime \beta}(\mathbf{q})\end{split}
\end{equation*}
\sphinxAtStartPar
where \(\mathbf{R}\) is the instantaneous positions of atoms, and
\(\left<\mathbf{R}\right>\) is the averaged atomic positions. It
gives essentially the same results as the displacement method and is
easier to implement in an MD code.

\sphinxAtStartPar
Once the force constant matrix is known, the dynamical matrix
\(\mathbf{D}\) can then be obtained by
\begin{equation*}
\begin{split}\mathbf{D}_{k\alpha, k^\prime\beta}(\mathbf{q}) =
(m_k m_{k^\prime})^{-\frac{1}{2}} \mathbf{\Phi}_{k \alpha, k^\prime \beta}(\mathbf{q})\end{split}
\end{equation*}
\sphinxAtStartPar
whose eigenvalues are exactly the phonon frequencies at \(\mathbf{q}\).

\sphinxAtStartPar
This fix uses positions of atoms in the specified group and calculates
two\sphinxhyphen{}point correlations.  To achieve this. the positions of the atoms
are examined every \sphinxstyleemphasis{Nevery} steps and are Fourier\sphinxhyphen{}transformed into
reciprocal space, where the averaging process and correlation
computation is then done.  After every \sphinxstyleemphasis{Noutput} measurements, the
matrix \(\mathbf{G}(\mathbf{q})\) is calculated and inverted to
obtain the elastic stiffness coefficients.  The dynamical matrices are
then constructed and written to \sphinxstyleemphasis{prefix}.bin.timestep files in binary
format and to the file \sphinxstyleemphasis{prefix}.log for each wave\sphinxhyphen{}vector
\(\mathbf{q}\).

\sphinxAtStartPar
A detailed description of this method can be found in
({\hyperref[\detokenize{fix_phonon:kong2011}]{\sphinxcrossref{\DUrole{std,std-ref}{Kong2011}}}}).

\sphinxAtStartPar
The \sphinxstyleemphasis{sysdim} keyword is optional.  If specified with a value smaller
than the dimensionality of the LAMMPS simulation, its value is used
for the dynamical matrix calculation.  For example, using LAMMPS to
model a 2D or 3D system, the phonon dispersion of a 1D atomic chain
can be computed using \sphinxstyleemphasis{sysdim} = 1.

\sphinxAtStartPar
The \sphinxstyleemphasis{nasr} keyword is optional.  An iterative procedure is employed to
enforce the acoustic sum rule on \(\Phi\) at \(\Gamma\), and the number
provided by keyword \sphinxstyleemphasis{nasr} gives the total number of iterations. For a
system whose unit cell has only one atom, \sphinxstyleemphasis{nasr} = 1 is sufficient;
for other systems, \sphinxstyleemphasis{nasr} = 10 is typically sufficient.

\sphinxAtStartPar
The \sphinxstyleemphasis{map\_file} contains the mapping information between the lattice
indices and the atom IDs, which tells the code which atom sits at
which lattice point; the lattice indices start from 0. An auxiliary
code, \sphinxhref{https://code.google.com/p/latgen}{latgen}, can be employed to
generate the compatible map file for various crystals.

\sphinxAtStartPar
In case one simulates a non\sphinxhyphen{}periodic system, where the whole simulation
box is treated as a unit cell, one can set \sphinxstyleemphasis{map\_file} as \sphinxstyleemphasis{GAMMA}, so
that the mapping info will be generated internally and a file is not
needed. In this case, the dynamical matrix at only the gamma\sphinxhyphen{}point
will/can be evaluated. Please keep in mind that fix\sphinxhyphen{}phonon is designed
for cyrstals, it will be inefficient and even degrade the performance
of LAMMPS in case the unit cell is too large.

\sphinxAtStartPar
The calculated dynamical matrix elements are written out in
{\hyperref[\detokenize{units::doc}]{\sphinxcrossref{\DUrole{doc}{energy/distance\textasciicircum{}2/mass}}}} units.  The coordinates for \sphinxstyleemphasis{q}
points in the log file is in the units of the basis vectors of the
corresponding reciprocal lattice.


\subsection{Restart, fix\_modify, output, run start/stop, minimize info}
\label{\detokenize{fix_phonon:restart-fix-modify-output-run-start-stop-minimize-info}}
\sphinxAtStartPar
No information about this fix is written to {\hyperref[\detokenize{restart::doc}]{\sphinxcrossref{\DUrole{doc}{binary restart files}}}}.

\sphinxAtStartPar
The {\hyperref[\detokenize{fix_modify::doc}]{\sphinxcrossref{\DUrole{doc}{fix\_modify}}}} \sphinxstyleemphasis{temp} option is supported by this
fix. You can use it to change the temperature compute from thermo\_temp
to the one that reflects the true temperature of atoms in the group.

\sphinxAtStartPar
No global scalar or vector or per\sphinxhyphen{}atom quantities are stored by this
fix for access by various {\hyperref[\detokenize{Howto_output::doc}]{\sphinxcrossref{\DUrole{doc}{output commands}}}}.

\sphinxAtStartPar
Instead, this fix outputs its initialization information (including
mapping information) and the calculated dynamical matrices to the file
\sphinxstyleemphasis{prefix}.log, with the specified \sphinxstyleemphasis{prefix}.  The dynamical matrices are
also written to files \sphinxstyleemphasis{prefix}.bin.timestep in binary format.  These
can be read by the post\sphinxhyphen{}processing tool in tools/phonon to compute the
phonon density of states and/or phonon dispersion curves.

\sphinxAtStartPar
No parameter of this fix can be used with the \sphinxstyleemphasis{start/stop} keywords
of the {\hyperref[\detokenize{run::doc}]{\sphinxcrossref{\DUrole{doc}{run}}}} command.

\sphinxAtStartPar
This fix is not invoked during {\hyperref[\detokenize{minimize::doc}]{\sphinxcrossref{\DUrole{doc}{energy minimization}}}}.


\subsection{Restrictions}
\label{\detokenize{fix_phonon:restrictions}}
\sphinxAtStartPar
This fix assumes a crystalline system with periodical lattice. The
temperature of the system should not exceed the melting temperature to
keep the system in its solid state.

\sphinxAtStartPar
This fix is part of the PHONON package.  It is only enabled if LAMMPS
was built with that package.  This fix also requires LAMMPS to be built
with 3d\sphinxhyphen{}FFT support which is included in the KSPACE package.  See the
{\hyperref[\detokenize{Build_package::doc}]{\sphinxcrossref{\DUrole{doc}{Build package}}}} page for more info.


\subsection{Related commands}
\label{\detokenize{fix_phonon:related-commands}}
\sphinxAtStartPar
{\hyperref[\detokenize{compute_msd::doc}]{\sphinxcrossref{\DUrole{doc}{compute msd}}}},
{\hyperref[\detokenize{dynamical_matrix::doc}]{\sphinxcrossref{\DUrole{doc}{dynamical\_matrix}}}}


\subsection{Default}
\label{\detokenize{fix_phonon:default}}
\sphinxAtStartPar
The option defaults are sysdim = the same dimension as specified by
the {\hyperref[\detokenize{dimension::doc}]{\sphinxcrossref{\DUrole{doc}{dimension}}}} command, and nasr = 20.


\bigskip\hrule\bigskip

\phantomsection\label{\detokenize{fix_phonon:campana}}
\sphinxAtStartPar
\sphinxstylestrong{(Campana)} C. Campana and
M. H. Muser, \sphinxstyleemphasis{Practical Green’s function approach to the
simulation of elastic semi\sphinxhyphen{}infinite solids}, \sphinxhref{https://doi.org/10.1103/PhysRevB.74.075420}{Phys. Rev. B {[}74{]}, 075420 (2006)}

\phantomsection\label{\detokenize{fix_phonon:kong}}
\sphinxAtStartPar
\sphinxstylestrong{(Kong)} L.T. Kong, G. Bartels, C. Campana,
C. Denniston, and Martin H. Muser, \sphinxstyleemphasis{Implementation of Green’s
function molecular dynamics: An extension to LAMMPS}, \sphinxhref{https://doi.org/10.1016/j.cpc.2008.12.035}{Computer Physics Communications {[}180{]}(6):1004\sphinxhyphen{}1010 (2009).}

\sphinxAtStartPar
L.T. Kong, C. Denniston, and Martin H. Muser,
\sphinxstyleemphasis{An improved version of the Green’s function molecular dynamics
method}, \sphinxhref{https://doi.org/10.1016/j.cpc.2010.10.006}{Computer Physics Communications {[}182{]}(2):540\sphinxhyphen{}541 (2011).}

\phantomsection\label{\detokenize{fix_phonon:kong2011}}
\sphinxAtStartPar
\sphinxstylestrong{(Kong2011)} L.T. Kong, \sphinxstyleemphasis{Phonon dispersion measured directly from
molecular dynamics simulations}, \sphinxhref{https://doi.org/10.1016/j.cpc.2011.04.019}{Computer Physics Communications {[}182{]}(10):2201\sphinxhyphen{}2207, (2011).}

\sphinxstepscope

\index{fix pimd/langevin@\spxentry{fix pimd/langevin}}\ignorespaces 
\index{fix pimd/nvt@\spxentry{fix pimd/nvt}}\ignorespaces 

\section{fix pimd/langevin command}
\label{\detokenize{fix_pimd:fix-pimd-langevin-command}}\label{\detokenize{fix_pimd:index-1}}\label{\detokenize{fix_pimd:index-0}}\label{\detokenize{fix_pimd::doc}}

\section{fix pimd/nvt command}
\label{\detokenize{fix_pimd:fix-pimd-nvt-command}}

\subsection{Syntax}
\label{\detokenize{fix_pimd:syntax}}
\begin{sphinxVerbatim}[commandchars=\\\{\}]
\PYG{k}{fix }\PYG{n+nv+nvIdentifier}{ID}\PYG{+w}{ }\PYG{n+nv+nvIdentifier}{group\PYGZhy{}ID}\PYG{+w}{ }\PYG{n}{style}\PYG{+w}{ }\PYG{n}{keyword}\PYG{+w}{ }\PYG{n}{value}\PYG{+w}{ }\PYG{n}{...}
\end{sphinxVerbatim}
\begin{itemize}
\item {} 
\sphinxAtStartPar
ID, group\sphinxhyphen{}ID are documented in {\hyperref[\detokenize{fix::doc}]{\sphinxcrossref{\DUrole{doc}{fix}}}} command

\item {} 
\sphinxAtStartPar
style = \sphinxstyleemphasis{pimd/langevin} or \sphinxstyleemphasis{pimd/nvt} = style name of this fix command

\item {} 
\sphinxAtStartPar
zero or more keyword/value pairs may be appended

\item {} 
\sphinxAtStartPar
keywords for style \sphinxstyleemphasis{pimd/nvt}
\begin{sphinxalltt}
\sphinxstyleemphasis{keywords} = \sphinxstyleemphasis{method} or \sphinxstyleemphasis{fmass} or \sphinxstyleemphasis{sp} or \sphinxstyleemphasis{temp} or \sphinxstyleemphasis{nhc}
\sphinxstyleemphasis{method} value = \sphinxstyleemphasis{pimd} or \sphinxstyleemphasis{nmpimd} or \sphinxstyleemphasis{cmd}
\sphinxstyleemphasis{fmass} value = scaling factor on mass
\sphinxstyleemphasis{sp} value = scaling factor on Planck constant
\sphinxstyleemphasis{temp} value = temperature (temperature units)
\sphinxstyleemphasis{nhc} value = Nc = number of chains in Nose\sphinxhyphen{}Hoover thermostat
\end{sphinxalltt}

\item {} 
\sphinxAtStartPar
keywords for style \sphinxstyleemphasis{pimd/langevin}
\begin{sphinxalltt}
\sphinxstyleemphasis{keywords} = \sphinxstyleemphasis{method} or \sphinxstyleemphasis{integrator} or \sphinxstyleemphasis{ensemble} or \sphinxstyleemphasis{fmmode} or \sphinxstyleemphasis{fmass} or \sphinxstyleemphasis{scale} or \sphinxstyleemphasis{temp} or \sphinxstyleemphasis{thermostat} or \sphinxstyleemphasis{tau} or \sphinxstyleemphasis{iso} or \sphinxstyleemphasis{aniso} or \sphinxstyleemphasis{barostat} or \sphinxstyleemphasis{taup} or \sphinxstyleemphasis{fixcom} or \sphinxstyleemphasis{lj}
\sphinxstyleemphasis{method} value = \sphinxstyleemphasis{nmpimd} (default) or \sphinxstyleemphasis{pimd}
\sphinxstyleemphasis{integrator} value = \sphinxstyleemphasis{obabo} or \sphinxstyleemphasis{baoab}
\sphinxstyleemphasis{fmmode} value = \sphinxstyleemphasis{physical} or \sphinxstyleemphasis{normal}
\sphinxstyleemphasis{fmass} value = scaling factor on mass
\sphinxstyleemphasis{temp} value = temperature (temperature unit)
   temperature = target temperature of the thermostat
\sphinxstyleemphasis{thermostat} values = style seed
   style value = \sphinxstyleemphasis{PILE\_L}
   seed = random number generator seed
\sphinxstyleemphasis{tau} value = thermostat damping parameter (time unit)
\sphinxstyleemphasis{scale} value = scaling factor of the damping times of non\sphinxhyphen{}centroid modes of PILE\_L thermostat
\sphinxstyleemphasis{iso} or \sphinxstyleemphasis{aniso} values = pressure (pressure unit)
  pressure = scalar external pressure of the barostat
\sphinxstyleemphasis{barostat} value = \sphinxstyleemphasis{BZP} or \sphinxstyleemphasis{MTTK}
\sphinxstyleemphasis{taup} value = barostat damping parameter (time unit)
\sphinxstyleemphasis{fixcom} value = \sphinxstyleemphasis{yes} or \sphinxstyleemphasis{no}
\sphinxstyleemphasis{lj} values = epsilon sigma mass planck mvv2e
   epsilon = energy scale for reduced units (energy units)
   sigma = length scale for reduced units (length units)
   mass = mass scale for reduced units (mass units)
   planck = Planck\textquotesingle{}s constant for other unit style
   mvv2e = mass * velocity\textasciicircum{}2 to energy conversion factor for other unit style
\end{sphinxalltt}

\end{itemize}


\subsection{Examples}
\label{\detokenize{fix_pimd:examples}}
\begin{sphinxVerbatim}[commandchars=\\\{\}]
\PYG{k}{fix }\PYG{n+nv+nvIdentifier}{1}\PYG{+w}{ }\PYG{n+nv+nvIdentifier}{all}\PYG{+w}{ }\PYG{n}{pimd}\PYG{o}{/}\PYG{n}{nvt}\PYG{+w}{ }\PYG{n}{method}\PYG{+w}{ }\PYG{n}{nmpimd}\PYG{+w}{ }\PYG{n}{fmass}\PYG{+w}{ }\PYG{l+m}{1.0}\PYG{+w}{ }\PYG{n}{sp}\PYG{+w}{ }\PYG{l+m}{2.0}\PYG{+w}{ }\PYG{n}{temp}\PYG{+w}{ }\PYG{l+m}{300.0}\PYG{+w}{ }\PYG{n}{nhc}\PYG{+w}{ }\PYG{l+m}{4}
\PYG{k}{fix }\PYG{n+nv+nvIdentifier}{1}\PYG{+w}{ }\PYG{n+nv+nvIdentifier}{all}\PYG{+w}{ }\PYG{n}{pimd}\PYG{o}{/}\PYG{n}{langevin}\PYG{+w}{ }\PYG{n}{ensemble}\PYG{+w}{ }\PYG{n}{npt}\PYG{+w}{ }\PYG{n}{integrator}\PYG{+w}{ }\PYG{n}{obabo}\PYG{+w}{ }\PYG{n}{temp}\PYG{+w}{ }\PYG{l+m}{113.15}\PYG{+w}{ }\PYG{n}{thermostat}\PYG{+w}{ }\PYG{n}{PILE\PYGZus{}L}\PYG{+w}{ }\PYG{l+m}{1234}\PYG{+w}{ }\PYG{n}{tau}\PYG{+w}{ }\PYG{l+m}{1.0}\PYG{+w}{ }\PYG{n}{iso}\PYG{+w}{ }\PYG{l+m}{1.0}\PYG{+w}{ }\PYG{n}{barostat}\PYG{+w}{ }\PYG{n}{BZP}\PYG{+w}{ }\PYG{n}{taup}\PYG{+w}{ }\PYG{l+m}{1.0}
\end{sphinxVerbatim}


\subsection{Description}
\label{\detokenize{fix_pimd:description}}
\sphinxAtStartPar
\DUrole{versionmodified,changed}{Changed in version 28Mar2023.}

\sphinxAtStartPar
Fix pimd was renamed to fix \sphinxstyleemphasis{pimd/nvt} and fix \sphinxstyleemphasis{pimd/langevin} was added.

\sphinxAtStartPar
These fix commands perform quantum molecular dynamics simulations based
on the Feynman path\sphinxhyphen{}integral to include effects of tunneling and
zero\sphinxhyphen{}point motion.  In this formalism, the isomorphism of a quantum
partition function for the original system to a classical partition
function for a ring\sphinxhyphen{}polymer system is exploited, to efficiently sample
configurations from the canonical ensemble {\hyperref[\detokenize{fix_pimd:feynman}]{\sphinxcrossref{\DUrole{std,std-ref}{(Feynman)}}}}.

\sphinxAtStartPar
The classical partition function and its components are given
by the following equations:
\begin{equation*}
\begin{split}Z = & \int d{\bf q} d{\bf p} \cdot \textrm{exp} [ -\beta H_{eff} ] \\
H_{eff} = & \bigg(\sum_{i=1}^P \frac{p_i^2}{2M_i}\bigg) + V_{eff} \\
V_{eff} = & \sum_{i=1}^P \bigg[ \frac{mP}{2\beta^2 \hbar^2} (q_i - q_{i+1})^2 + \frac{1}{P} V(q_i)\bigg]\end{split}
\end{equation*}
\sphinxAtStartPar
\(M_i\) is the fictitious mass of the \(i\)\sphinxhyphen{}th mode, and m is the actual mass of the atoms.

\sphinxAtStartPar
The interested user is referred to any of the numerous references on
this methodology, but briefly, each quantum particle in a path integral
simulation is represented by a ring\sphinxhyphen{}polymer of P quasi\sphinxhyphen{}beads, labeled
from 1 to P.  During the simulation, each quasi\sphinxhyphen{}bead interacts with
beads on the other ring\sphinxhyphen{}polymers with the same imaginary time index (the
second term in the effective potential above).  The quasi\sphinxhyphen{}beads also
interact with the two neighboring quasi\sphinxhyphen{}beads through the spring
potential in imaginary\sphinxhyphen{}time space (first term in effective potential).
To sample the canonical ensemble, any thermostat can be applied.

\sphinxAtStartPar
Fix \sphinxstyleemphasis{pimd/nvt} applies a Nose\sphinxhyphen{}Hoover massive chain thermostat
{\hyperref[\detokenize{fix_pimd:pimd-tuckerman}]{\sphinxcrossref{\DUrole{std,std-ref}{(Tuckerman)}}}}.  With the massive chain
algorithm, a chain of NH thermostats is coupled to each degree of
freedom for each quasi\sphinxhyphen{}bead.  The keyword \sphinxstyleemphasis{temp} sets the target
temperature for the system and the keyword \sphinxstyleemphasis{nhc} sets the number \sphinxstyleemphasis{Nc} of
thermostats in each chain.  For example, for a simulation of N particles
with P beads in each ring\sphinxhyphen{}polymer, the total number of NH thermostats
would be 3 x N x P x Nc.

\sphinxAtStartPar
Fix \sphinxstyleemphasis{pimd/langevin} implements a Langevin thermostat in the normal mode
representation, and also provides a barostat to sample the NPH/NPT ensembles.

\begin{sphinxadmonition}{note}{Note:}
\sphinxAtStartPar
Both these \sphinxstyleemphasis{fix} styles implement a complete velocity\sphinxhyphen{}verlet integrator
combined with a thermostat, so no other time integration fix should be used.
\end{sphinxadmonition}

\sphinxAtStartPar
The \sphinxstyleemphasis{method} keyword determines what style of PIMD is performed.  A
value of \sphinxstyleemphasis{pimd} is standard PIMD.  A value of \sphinxstyleemphasis{nmpimd} is for
normal\sphinxhyphen{}mode PIMD.  A value of \sphinxstyleemphasis{cmd} is for centroid molecular dynamics
(CMD).  The difference between the styles is as follows.
\begin{quote}

\sphinxAtStartPar
In standard PIMD, the value used for a bead’s fictitious mass is
arbitrary.  A common choice is to use \(M_i = m/P\), which results in the
mass of the entire ring\sphinxhyphen{}polymer being equal to the real quantum
particle.  But it can be difficult to efficiently integrate the
equations of motion for the stiff harmonic interactions in the ring
polymers.

\sphinxAtStartPar
A useful way to resolve this issue is to integrate the equations of
motion in a normal mode representation, using Normal Mode
Path\sphinxhyphen{}Integral Molecular Dynamics (NMPIMD) {\hyperref[\detokenize{fix_pimd:cao1}]{\sphinxcrossref{\DUrole{std,std-ref}{(Cao1)}}}}.  In
NMPIMD, the NH chains are attached to each normal mode of the
ring\sphinxhyphen{}polymer and the fictitious mass of each mode is chosen as Mk =
the eigenvalue of the Kth normal mode for k \textgreater{} 0. The k = 0 mode,
referred to as the zero\sphinxhyphen{}frequency mode or centroid, corresponds to
overall translation of the ring\sphinxhyphen{}polymer and is assigned the mass of
the real particle.
\end{quote}

\begin{sphinxadmonition}{note}{Note:}
\sphinxAtStartPar
Motion of the centroid can be effectively uncoupled from the other
normal modes by scaling the fictitious masses to achieve a partial
adiabatic separation.  This is called a Centroid Molecular Dynamics
(CMD) approximation {\hyperref[\detokenize{fix_pimd:cao2}]{\sphinxcrossref{\DUrole{std,std-ref}{(Cao2)}}}}.  The time\sphinxhyphen{}evolution (and
resulting dynamics) of the quantum particles can be used to obtain
centroid time correlation functions, which can be further used to
obtain the true quantum correlation function for the original system.
The CMD method also uses normal modes to evolve the system, except
only the k \textgreater{} 0 modes are thermostatted, not the centroid degrees of
freedom.
\end{sphinxadmonition}

\sphinxAtStartPar
\DUrole{versionmodified,added}{Added in version 21Nov2023: }Mode \sphinxstyleemphasis{pimd} added to fix pimd/langevin.

\sphinxAtStartPar
Fix pimd/langevin supports the \sphinxstyleemphasis{method} values \sphinxstyleemphasis{nmpimd} and \sphinxstyleemphasis{pimd}. The default value is \sphinxstyleemphasis{nmpimd}.
If \sphinxstyleemphasis{method} is \sphinxstyleemphasis{nmpimd}, the normal mode representation is used to integrate the equations of motion.
The exact solution of harmonic oscillator is used to propagate the free ring polymer part of the Hamiltonian.
If \sphinxstyleemphasis{method} is \sphinxstyleemphasis{pimd}, the Cartesian representation is used to integrate the equations of motion.
The harmonic force is added to the total force of the system, and the numerical integrator is used to propagate the Hamiltonian.

\sphinxAtStartPar
The keyword \sphinxstyleemphasis{integrator} specifies the Trotter splitting method used by \sphinxstyleemphasis{fix pimd/langevin}.
See {\hyperref[\detokenize{pair_hbond_dreiding:liu}]{\sphinxcrossref{\DUrole{std,std-ref}{(Liu)}}}} for a discussion on the OBABO and BAOAB splitting schemes. Typically
either of the two should work fine.

\sphinxAtStartPar
The keyword \sphinxstyleemphasis{fmass} sets a further scaling factor for the fictitious
masses of beads, which can be used for the Partial Adiabatic CMD
{\hyperref[\detokenize{fix_pimd:hone}]{\sphinxcrossref{\DUrole{std,std-ref}{(Hone)}}}}, or to be set as P, which results in the fictitious
masses to be equal to the real particle masses.

\sphinxAtStartPar
The keyword \sphinxstyleemphasis{fmmode} of \sphinxstyleemphasis{fix pimd/langevin} determines the mode of fictitious
mass preconditioning. There are two options: \sphinxstyleemphasis{physical} and \sphinxstyleemphasis{normal}. If \sphinxstyleemphasis{fmmode} is
\sphinxstyleemphasis{physical}, then the physical mass of the particles are used (and then multiplied by
\sphinxstyleemphasis{fmass}). If \sphinxstyleemphasis{fmmode} is \sphinxstyleemphasis{normal}, then the physical mass is first multiplied by the
eigenvalue of each normal mode, and then multiplied by \sphinxstyleemphasis{fmass}. More precisely, the
fictitious mass of \sphinxstyleemphasis{fix pimd/langevin} is determined by two factors: \sphinxstyleemphasis{fmmode} and \sphinxstyleemphasis{fmass}.
If \sphinxstyleemphasis{fmmode} is \sphinxstyleemphasis{physical}, then the fictitious mass is
\begin{equation*}
\begin{split}M_i = \mathrm{fmass} \times m\end{split}
\end{equation*}
\sphinxAtStartPar
If \sphinxstyleemphasis{fmmode} is \sphinxstyleemphasis{normal}, then the fictitious mass is
\begin{equation*}
\begin{split}M_i = \mathrm{fmass} \times \lambda_i \times m\end{split}
\end{equation*}
\sphinxAtStartPar
where \(\lambda_i\) is the eigenvalue of the \(i\)\sphinxhyphen{}th normal mode.

\begin{sphinxadmonition}{note}{Note:}
\sphinxAtStartPar
Fictitious mass is only used in the momentum of the equation of motion
(\(\mathbf{p}_i=M_i\mathbf{v}_i\)), and not used in the spring elastic energy
(\(\sum_{i=1}^P \frac{1}{2}m\omega_P^2(q_i - q_{i+1})^2\), \(m\) is always the
actual mass of the particles).
\end{sphinxadmonition}

\sphinxAtStartPar
The keyword \sphinxstyleemphasis{sp} is a scaling factor on Planck’s constant, which can
be useful for debugging or other purposes.  The default value of 1.0
is appropriate for most situations.

\sphinxAtStartPar
The keyword \sphinxstyleemphasis{ensemble} for fix style \sphinxstyleemphasis{pimd/langevin} determines which ensemble is it
going to sample. The value can be \sphinxstyleemphasis{nve} (microcanonical), \sphinxstyleemphasis{nvt} (canonical), \sphinxstyleemphasis{nph} (isoenthalpic),
and \sphinxstyleemphasis{npt} (isothermal\sphinxhyphen{}isobaric).

\sphinxAtStartPar
The keyword \sphinxstyleemphasis{temp} specifies temperature parameter for fix styles \sphinxstyleemphasis{pimd/nvt} and \sphinxstyleemphasis{pimd/langevin}. It should read
a positive floating\sphinxhyphen{}point number.

\begin{sphinxadmonition}{note}{Note:}
\sphinxAtStartPar
For pimd simulations, a temperature values should be specified even for nve ensemble. Temperature will make a difference
for nve pimd, since the spring elastic frequency between the beads will be affected by the temperature.
\end{sphinxadmonition}

\sphinxAtStartPar
The keyword \sphinxstyleemphasis{thermostat} reads \sphinxstyleemphasis{style} and \sphinxstyleemphasis{seed} of thermostat for fix style \sphinxstyleemphasis{pimd/langevin}. \sphinxstyleemphasis{style} can only
be \sphinxstyleemphasis{PILE\_L} (path integral Langevin equation local thermostat, as described in {\hyperref[\detokenize{fix_gle:ceriotti2}]{\sphinxcrossref{\DUrole{std,std-ref}{Ceriotti}}}}), and \sphinxstyleemphasis{seed} should a positive integer number, which serves as the seed of the pseudo random number generator.

\begin{sphinxadmonition}{note}{Note:}
\sphinxAtStartPar
The fix style \sphinxstyleemphasis{pimd/langevin} uses the stochastic PILE\_L thermostat to control temperature. This thermostat works on the normal modes
of the ring polymer. The \sphinxstyleemphasis{tau} parameter controls the centroid mode, and the \sphinxstyleemphasis{scale} parameter controls the non\sphinxhyphen{}centroid modes.
\end{sphinxadmonition}

\sphinxAtStartPar
The keyword \sphinxstyleemphasis{tau} specifies the thermostat damping time parameter for fix style \sphinxstyleemphasis{pimd/langevin}. It is in time unit. It only works on the centroid mode.

\sphinxAtStartPar
The keyword \sphinxstyleemphasis{scale} specifies a scaling parameter for the damping times of the non\sphinxhyphen{}centroid modes for fix style \sphinxstyleemphasis{pimd/langevin}. The default
damping time of the non\sphinxhyphen{}centroid mode \(i\) is \(\frac{P}{\beta\hbar}\sqrt{\lambda_i\times\mathrm{fmass}}\) (\sphinxstyleemphasis{fmmode} is \sphinxstyleemphasis{physical}) or  \(\frac{P}{\beta\hbar}\sqrt{\mathrm{fmass}}\) (\sphinxstyleemphasis{fmmode} is \sphinxstyleemphasis{normal}). The damping times of all non\sphinxhyphen{}centroid modes are the default values divided by \sphinxstyleemphasis{scale}.

\sphinxAtStartPar
The barostat parameters for fix style \sphinxstyleemphasis{pimd/langevin} with \sphinxstyleemphasis{npt} or \sphinxstyleemphasis{nph} ensemble is specified using one of \sphinxstyleemphasis{iso} and \sphinxstyleemphasis{aniso}
keywords. A \sphinxstyleemphasis{pressure} value should be given with pressure unit. The keyword \sphinxstyleemphasis{iso} means couple all 3 diagonal components together when pressure is computed (hydrostatic pressure), and dilate/contract the dimensions together. The keyword \sphinxstyleemphasis{aniso} means x, y, and z dimensions are controlled independently using the Pxx, Pyy, and Pzz components of the stress tensor as the driving forces, and the specified scalar external pressure.

\sphinxAtStartPar
The keyword \sphinxstyleemphasis{barostat} reads \sphinxstyleemphasis{style} of barostat for fix style \sphinxstyleemphasis{pimd/langevin}. \sphinxstyleemphasis{style} can
be \sphinxstyleemphasis{BZP} (Bussi\sphinxhyphen{}Zykova\sphinxhyphen{}Parrinello, as described in {\hyperref[\detokenize{fix_pimd:bussi}]{\sphinxcrossref{\DUrole{std,std-ref}{Bussi}}}}) or \sphinxstyleemphasis{MTTK} (Martyna\sphinxhyphen{}Tuckerman\sphinxhyphen{}Tobias\sphinxhyphen{}Klein, as described in {\hyperref[\detokenize{fix_pimd:martyna3}]{\sphinxcrossref{\DUrole{std,std-ref}{Martyna1}}}} and {\hyperref[\detokenize{fix_pimd:martyna4}]{\sphinxcrossref{\DUrole{std,std-ref}{Martyna2}}}}).

\sphinxAtStartPar
The keyword \sphinxstyleemphasis{taup} specifies the barostat damping time parameter for fix style \sphinxstyleemphasis{pimd/langevin}. It is in time unit.

\sphinxAtStartPar
The keyword \sphinxstyleemphasis{fixcom} specifies whether the center\sphinxhyphen{}of\sphinxhyphen{}mass of the extended ring\sphinxhyphen{}polymer system is fixed during the pimd simulation.
Once \sphinxstyleemphasis{fixcom} is set to be \sphinxstyleemphasis{yes}, the center\sphinxhyphen{}of\sphinxhyphen{}mass velocity will be distracted from the centroid\sphinxhyphen{}mode velocities in each step.

\sphinxAtStartPar
The keyword \sphinxstyleemphasis{lj} should be used if {\hyperref[\detokenize{units::doc}]{\sphinxcrossref{\DUrole{doc}{lj units}}}} is used for \sphinxstyleemphasis{fix pimd/langevin}. Typically one may want to use
reduced units to run the simulation, and then convert the results into some physical units (for example, {\hyperref[\detokenize{units::doc}]{\sphinxcrossref{\DUrole{doc}{metal units}}}}). In this case, the 5 quantities in the physical mass units are needed: epsilon (energy scale), sigma (length scale), mass, Planck’s constant, mvv2e (mass * velocity\textasciicircum{}2 to energy conversion factor). Planck’s constant and mvv2e can be found in src/update.cpp. If there is no need to convert reduced units to physical units, you can omit the keyword \sphinxstyleemphasis{lj} and these five values will be set to 1.

\sphinxAtStartPar
The PIMD algorithm in LAMMPS is implemented as a hyper\sphinxhyphen{}parallel scheme
as described in {\hyperref[\detokenize{fix_pimd:calhoun}]{\sphinxcrossref{\DUrole{std,std-ref}{Calhoun}}}}.  In LAMMPS this is done by using
{\hyperref[\detokenize{Howto_replica::doc}]{\sphinxcrossref{\DUrole{doc}{multi\sphinxhyphen{}replica feature}}}} in LAMMPS, where each
quasi\sphinxhyphen{}particle system is stored and simulated on a separate partition
of processors.  The following diagram illustrates this approach.  The
original system with 2 ring polymers is shown in red.  Since each ring
has 4 quasi\sphinxhyphen{}beads (imaginary time slices), there are 4 replicas of the
system, each running on one of the 4 partitions of processors.  Each
replica (shown in green) owns one quasi\sphinxhyphen{}bead in each ring.

\noindent{\hspace*{\fill}\sphinxincludegraphics{{pimd}.jpg}\hspace*{\fill}}

\sphinxAtStartPar
To run a PIMD simulation with M quasi\sphinxhyphen{}beads in each ring polymer using
N MPI tasks for each partition’s domain\sphinxhyphen{}decomposition, you would use P
= MxN processors (cores) and run the simulation as follows:

\begin{sphinxVerbatim}[commandchars=\\\{\}]
mpirun\PYG{+w}{ }\PYGZhy{}np\PYG{+w}{ }P\PYG{+w}{ }lmp\PYGZus{}mpi\PYG{+w}{ }\PYGZhy{}partition\PYG{+w}{ }MxN\PYG{+w}{ }\PYGZhy{}in\PYG{+w}{ }script
\end{sphinxVerbatim}

\sphinxAtStartPar
Note that in the LAMMPS input script for a multi\sphinxhyphen{}partition simulation,
it is often very useful to define a {\hyperref[\detokenize{variable::doc}]{\sphinxcrossref{\DUrole{doc}{uloop\sphinxhyphen{}style variable}}}} such as

\begin{sphinxVerbatim}[commandchars=\\\{\}]
\PYG{k}{variable }\PYG{n+nv+nvIdentifier}{ibead}\PYG{+w}{ }\PYG{n}{uloop}\PYG{+w}{ }\PYG{n}{M}\PYG{+w}{ }\PYG{n}{pad}
\end{sphinxVerbatim}

\sphinxAtStartPar
where M is the number of quasi\sphinxhyphen{}beads (partitions) used in the
calculation.  The uloop variable can then be used to manage I/O
related tasks for each of the partitions, e.g.

\begin{sphinxVerbatim}[commandchars=\\\{\}]
\PYG{k}{dump }\PYG{n+nv+nvIdentifier}{dcd}\PYG{+w}{ }\PYG{n+nv+nvIdentifier}{all}\PYG{+w}{ }\PYG{n}{dcd}\PYG{+w}{ }\PYG{l+m}{10}\PYG{+w}{ }\PYG{n}{system\PYGZus{}}\PYG{n+nv}{\PYGZdl{}\PYGZob{}}\PYG{n+nv}{ibead}\PYG{n+nv}{\PYGZcb{}}\PYG{n}{.dcd}
\PYG{k}{dump }\PYG{n+nv+nvIdentifier}{1}\PYG{+w}{ }\PYG{n+nv+nvIdentifier}{all}\PYG{+w}{ }\PYG{n}{custom}\PYG{+w}{ }\PYG{l+m}{100}\PYG{+w}{ }\PYG{n+nv}{\PYGZdl{}\PYGZob{}}\PYG{n+nv}{ibead}\PYG{n+nv}{\PYGZcb{}}\PYG{n}{.xyz}\PYG{+w}{ }\PYG{n}{id}\PYG{+w}{ }\PYG{n}{type}\PYG{+w}{ }\PYG{n}{x}\PYG{+w}{ }\PYG{n}{y}\PYG{+w}{ }\PYG{n}{z}\PYG{+w}{ }\PYG{n}{vx}\PYG{+w}{ }\PYG{n}{vy}\PYG{+w}{ }\PYG{n}{vz}\PYG{+w}{ }\PYG{n}{ix}\PYG{+w}{ }\PYG{n}{iy}\PYG{+w}{ }\PYG{n}{iz}\PYG{+w}{ }\PYG{n}{fx}\PYG{+w}{ }\PYG{n}{fy}\PYG{+w}{ }\PYG{n}{fz}
\PYG{k}{restart}\PYG{+w}{ }\PYG{l+m}{1000}\PYG{+w}{ }\PYG{n}{system\PYGZus{}}\PYG{n+nv}{\PYGZdl{}\PYGZob{}}\PYG{n+nv}{ibead}\PYG{n+nv}{\PYGZcb{}}\PYG{n}{.restart1}\PYG{+w}{ }\PYG{n}{system\PYGZus{}}\PYG{n+nv}{\PYGZdl{}\PYGZob{}}\PYG{n+nv}{ibead}\PYG{n+nv}{\PYGZcb{}}\PYG{n}{.restart2}
\PYG{k}{read\PYGZus{}restart}\PYG{+w}{ }\PYG{n}{system\PYGZus{}}\PYG{n+nv}{\PYGZdl{}\PYGZob{}}\PYG{n+nv}{ibead}\PYG{n+nv}{\PYGZcb{}}\PYG{n}{.restart2}
\end{sphinxVerbatim}

\begin{sphinxadmonition}{note}{Note:}
\sphinxAtStartPar
Fix \sphinxstyleemphasis{pimd/langevin} dumps the Cartesian coordinates, but dumps the velocities and
forces in the normal mode representation. If the Cartesian velocities and forces are
needed, it is easy to perform the transformation when doing post\sphinxhyphen{}processing.

\sphinxAtStartPar
It is recommended to dump the image flags (\sphinxstyleemphasis{ix iy iz}) for fix \sphinxstyleemphasis{pimd/langevin}. It
will be useful if you want to calculate some estimators during post\sphinxhyphen{}processing.
\end{sphinxadmonition}

\sphinxAtStartPar
Major differences of \sphinxstyleemphasis{fix pimd/nvt} and \sphinxstyleemphasis{fix pimd/langevin} are:
\begin{enumerate}
\sphinxsetlistlabels{\arabic}{enumi}{enumii}{}{.}%
\item {} 
\sphinxAtStartPar
\sphinxstyleemphasis{Fix pimd/nvt} includes Cartesian pimd, normal mode pimd, and centroid md. \sphinxstyleemphasis{Fix pimd/langevin} only intends to support normal mode pimd, as it is commonly enough for thermodynamic sampling.

\item {} 
\sphinxAtStartPar
\sphinxstyleemphasis{Fix pimd/nvt} uses Nose\sphinxhyphen{}Hoover chain thermostat. \sphinxstyleemphasis{Fix pimd/langevin} uses Langevin thermostat.

\item {} 
\sphinxAtStartPar
\sphinxstyleemphasis{Fix pimd/langevin} provides barostat, so the npt ensemble can be sampled. \sphinxstyleemphasis{Fix pimd/nvt} only support nvt ensemble.

\item {} 
\sphinxAtStartPar
\sphinxstyleemphasis{Fix pimd/langevin} provides several quantum estimators in output.

\item {} 
\sphinxAtStartPar
\sphinxstyleemphasis{Fix pimd/langevin} allows multiple processes for each bead. For \sphinxstyleemphasis{fix pimd/nvt}, there is a large chance that multi\sphinxhyphen{}process tasks for each bead may fail.

\item {} 
\sphinxAtStartPar
The dump of \sphinxstyleemphasis{fix pimd/nvt} are all Cartesian. \sphinxstyleemphasis{Fix pimd/langevin} dumps normal\sphinxhyphen{}mode velocities and forces, and Cartesian coordinates.

\end{enumerate}

\sphinxAtStartPar
Initially, the inter\sphinxhyphen{}replica communication and normal mode transformation parts of \sphinxstyleemphasis{fix pimd/langevin} are written based on
those of \sphinxstyleemphasis{fix pimd/nvt}, but are significantly revised.


\subsection{Restart, fix\_modify, output, run start/stop, minimize info}
\label{\detokenize{fix_pimd:restart-fix-modify-output-run-start-stop-minimize-info}}
\sphinxAtStartPar
Fix \sphinxstyleemphasis{pimd/nvt} writes the state of the Nose/Hoover thermostat over all
quasi\sphinxhyphen{}beads to {\hyperref[\detokenize{restart::doc}]{\sphinxcrossref{\DUrole{doc}{binary restart files}}}}.  See the
{\hyperref[\detokenize{read_restart::doc}]{\sphinxcrossref{\DUrole{doc}{read\_restart}}}} command for info on how to re\sphinxhyphen{}specify
a fix in an input script that reads a restart file, so that the
operation of the fix continues in an uninterrupted fashion.

\sphinxAtStartPar
Fix \sphinxstyleemphasis{pimd/langevin} writes the state of the barostat overall beads to
{\hyperref[\detokenize{restart::doc}]{\sphinxcrossref{\DUrole{doc}{binary restart files}}}}. Since it uses a stochastic thermostat,
the state of the thermostat is not written. However, the state of the system
can be restored by reading the restart file, except that it will re\sphinxhyphen{}initialize
the random number generator.

\sphinxAtStartPar
None of the {\hyperref[\detokenize{fix_modify::doc}]{\sphinxcrossref{\DUrole{doc}{fix\_modify}}}} options
are relevant to fix pimd/nvt.

\sphinxAtStartPar
Fix \sphinxstyleemphasis{pimd/nvt} computes a global 3\sphinxhyphen{}vector, which can be accessed by
various {\hyperref[\detokenize{Howto_output::doc}]{\sphinxcrossref{\DUrole{doc}{output commands}}}}.  The three quantities in
the global vector are:
\begin{enumerate}
\sphinxsetlistlabels{\arabic}{enumi}{enumii}{}{.}%
\item {} 
\sphinxAtStartPar
the total spring energy of the quasi\sphinxhyphen{}beads,

\item {} 
\sphinxAtStartPar
the current temperature of the classical system of ring polymers,

\item {} 
\sphinxAtStartPar
the current value of the scalar virial estimator for the kinetic
energy of the quantum system {\hyperref[\detokenize{fix_pimd:herman}]{\sphinxcrossref{\DUrole{std,std-ref}{(Herman)}}}}.

\end{enumerate}

\sphinxAtStartPar
The vector values calculated by fix \sphinxstyleemphasis{pimd/nvt} are “extensive”, except for the
temperature, which is “intensive”.

\sphinxAtStartPar
Fix \sphinxstyleemphasis{pimd/langevin} computes a global vector of quantities, which
can be accessed by various {\hyperref[\detokenize{Howto_output::doc}]{\sphinxcrossref{\DUrole{doc}{output commands}}}}. Note that
it outputs multiple log files, and different log files contain information
about different beads or modes (see detailed explanations below). If \sphinxstyleemphasis{ensemble}
is \sphinxstyleemphasis{nve} or \sphinxstyleemphasis{nvt}, the vector has 10 values:
\begin{enumerate}
\sphinxsetlistlabels{\arabic}{enumi}{enumii}{}{.}%
\item {} 
\sphinxAtStartPar
kinetic energy of the normal mode

\item {} 
\sphinxAtStartPar
spring elastic energy of the normal mode

\item {} 
\sphinxAtStartPar
potential energy of the bead

\item {} 
\sphinxAtStartPar
total energy of all beads (conserved if \sphinxstyleemphasis{ensemble} is \sphinxstyleemphasis{nve})

\item {} 
\sphinxAtStartPar
primitive kinetic energy estimator

\item {} 
\sphinxAtStartPar
virial energy estimator

\item {} 
\sphinxAtStartPar
centroid\sphinxhyphen{}virial energy estimator

\item {} 
\sphinxAtStartPar
primitive pressure estimator

\item {} 
\sphinxAtStartPar
thermodynamic pressure estimator

\item {} 
\sphinxAtStartPar
centroid\sphinxhyphen{}virial pressure estimator

\end{enumerate}

\sphinxAtStartPar
The first 3 are different for different log files, and the others are the same for different log files.

\sphinxAtStartPar
If \sphinxstyleemphasis{ensemble} is \sphinxstyleemphasis{nph} or \sphinxstyleemphasis{npt}, the vector stores internal variables of the barostat. If \sphinxstyleemphasis{iso} is used,
the vector has 15 values:
\begin{enumerate}
\sphinxsetlistlabels{\arabic}{enumi}{enumii}{}{.}%
\item {} 
\sphinxAtStartPar
kinetic energy of the normal mode

\item {} 
\sphinxAtStartPar
spring elastic energy of the normal mode

\item {} 
\sphinxAtStartPar
potential energy of the bead

\item {} 
\sphinxAtStartPar
total energy of all beads (conserved if \sphinxstyleemphasis{ensemble} is \sphinxstyleemphasis{nve})

\item {} 
\sphinxAtStartPar
primitive kinetic energy estimator

\item {} 
\sphinxAtStartPar
virial energy estimator

\item {} 
\sphinxAtStartPar
centroid\sphinxhyphen{}virial energy estimator

\item {} 
\sphinxAtStartPar
primitive pressure estimator

\item {} 
\sphinxAtStartPar
thermodynamic pressure estimator

\item {} 
\sphinxAtStartPar
centroid\sphinxhyphen{}virial pressure estimator

\item {} 
\sphinxAtStartPar
barostat velocity

\item {} 
\sphinxAtStartPar
barostat kinetic energy

\item {} 
\sphinxAtStartPar
barostat potential energy

\item {} 
\sphinxAtStartPar
barostat cell Jacobian

\item {} 
\sphinxAtStartPar
enthalpy of the extended system (sum of 4, 12, 13, and 14; conserved if \sphinxstyleemphasis{ensemble} is \sphinxstyleemphasis{nph})

\end{enumerate}

\sphinxAtStartPar
If \sphinxstyleemphasis{aniso} or \sphinxstyleemphasis{x} or \sphinxstyleemphasis{y} or \sphinxstyleemphasis{z} is used for the barostat, the vector has 17 values:
\begin{enumerate}
\sphinxsetlistlabels{\arabic}{enumi}{enumii}{}{.}%
\item {} 
\sphinxAtStartPar
kinetic energy of the normal mode

\item {} 
\sphinxAtStartPar
spring elastic energy of the normal mode

\item {} 
\sphinxAtStartPar
potential energy of the bead

\item {} 
\sphinxAtStartPar
total energy of all beads (conserved if \sphinxstyleemphasis{ensemble} is \sphinxstyleemphasis{nve})

\item {} 
\sphinxAtStartPar
primitive kinetic energy estimator

\item {} 
\sphinxAtStartPar
virial energy estimator

\item {} 
\sphinxAtStartPar
centroid\sphinxhyphen{}virial energy estimator

\item {} 
\sphinxAtStartPar
primitive pressure estimator

\item {} 
\sphinxAtStartPar
thermodynamic pressure estimator

\item {} 
\sphinxAtStartPar
centroid\sphinxhyphen{}virial pressure estimator

\item {} 
\sphinxAtStartPar
x component of barostat velocity

\item {} 
\sphinxAtStartPar
y component of barostat velocity

\item {} 
\sphinxAtStartPar
z component of barostat velocity

\item {} 
\sphinxAtStartPar
barostat kinetic energy

\item {} 
\sphinxAtStartPar
barostat potential energy

\item {} 
\sphinxAtStartPar
barostat cell Jacobian

\item {} 
\sphinxAtStartPar
enthalpy of the extended system (sum of 4, 14, 15, and 16; conserved if \sphinxstyleemphasis{ensemble} is \sphinxstyleemphasis{nph})

\end{enumerate}

\sphinxAtStartPar
No parameter of fix \sphinxstyleemphasis{pimd/nvt} or \sphinxstyleemphasis{pimd/langevin} can be used with the \sphinxstyleemphasis{start/stop} keywords
of the {\hyperref[\detokenize{run::doc}]{\sphinxcrossref{\DUrole{doc}{run}}}} command.  Fix \sphinxstyleemphasis{pimd/nvt} or \sphinxstyleemphasis{pimd/langevin} is not invoked during
{\hyperref[\detokenize{minimize::doc}]{\sphinxcrossref{\DUrole{doc}{energy minimization}}}}.


\subsection{Restrictions}
\label{\detokenize{fix_pimd:restrictions}}
\sphinxAtStartPar
These fixes are part of the REPLICA package.  They are only enabled if
LAMMPS was built with that package.  See the {\hyperref[\detokenize{Build_package::doc}]{\sphinxcrossref{\DUrole{doc}{Build package}}}} page for more info.

\sphinxAtStartPar
Fix \sphinxstyleemphasis{pimd/nvt} cannot be used with {\hyperref[\detokenize{units::doc}]{\sphinxcrossref{\DUrole{doc}{lj units}}}}.
Fix \sphinxstyleemphasis{pimd/langevin} can be used with {\hyperref[\detokenize{units::doc}]{\sphinxcrossref{\DUrole{doc}{lj units}}}}. See the above part for how to use it.

\sphinxAtStartPar
A PIMD simulation can be initialized with a single data file read via
the {\hyperref[\detokenize{read_data::doc}]{\sphinxcrossref{\DUrole{doc}{read\_data}}}} command.  However, this means all
quasi\sphinxhyphen{}beads in a ring polymer will have identical positions and
velocities, resulting in identical trajectories for all quasi\sphinxhyphen{}beads.  To
avoid this, users can simply initialize velocities with different random
number seeds assigned to each partition, as defined by the uloop
variable, e.g.

\begin{sphinxVerbatim}[commandchars=\\\{\}]
\PYG{k}{velocity }\PYG{n+nv+nvIdentifier}{all}\PYG{+w}{ }\PYG{n}{create}\PYG{+w}{ }\PYG{l+m}{300.0}\PYG{+w}{ }\PYG{l+m}{1234}\PYG{n+nv}{\PYGZdl{}\PYGZob{}}\PYG{n+nv}{ibead}\PYG{n+nv}{\PYGZcb{}}\PYG{+w}{ }\PYG{n}{rot}\PYG{+w}{ }\PYG{n}{yes}\PYG{+w}{ }\PYG{n}{dist}\PYG{+w}{ }\PYG{n}{gaussian}
\end{sphinxVerbatim}


\subsection{Default}
\label{\detokenize{fix_pimd:default}}
\sphinxAtStartPar
The keyword defaults for fix \sphinxstyleemphasis{pimd/nvt} are method = pimd, fmass = 1.0, sp
= 1.0, temp = 300.0, and nhc = 2.


\bigskip\hrule\bigskip

\phantomsection\label{\detokenize{fix_pimd:feynman}}
\sphinxAtStartPar
\sphinxstylestrong{(Feynman)} R. Feynman and A. Hibbs, Chapter 7, Quantum Mechanics and
Path Integrals, McGraw\sphinxhyphen{}Hill, New York (1965).

\phantomsection\label{\detokenize{fix_pimd:pimd-tuckerman}}
\sphinxAtStartPar
\sphinxstylestrong{(Tuckerman)} M. Tuckerman and B. Berne, J Chem Phys, 99, 2796 (1993).

\phantomsection\label{\detokenize{fix_pimd:cao1}}
\sphinxAtStartPar
\sphinxstylestrong{(Cao1)} J. Cao and B. Berne, J Chem Phys, 99, 2902 (1993).

\phantomsection\label{\detokenize{fix_pimd:cao2}}
\sphinxAtStartPar
\sphinxstylestrong{(Cao2)} J. Cao and G. Voth, J Chem Phys, 100, 5093 (1994).

\phantomsection\label{\detokenize{fix_pimd:hone}}
\sphinxAtStartPar
\sphinxstylestrong{(Hone)} T. Hone, P. Rossky, G. Voth, J Chem Phys, 124,
154103 (2006).

\phantomsection\label{\detokenize{fix_pimd:calhoun}}
\sphinxAtStartPar
\sphinxstylestrong{(Calhoun)} A. Calhoun, M. Pavese, G. Voth, Chem Phys Letters, 262,
415 (1996).

\phantomsection\label{\detokenize{fix_pimd:herman}}
\sphinxAtStartPar
\sphinxstylestrong{(Herman)} M. F. Herman, E. J. Bruskin, B. J. Berne, J Chem Phys, 76, 5150 (1982).

\phantomsection\label{\detokenize{fix_pimd:bussi}}
\sphinxAtStartPar
\sphinxstylestrong{(Bussi)} G. Bussi, T. Zykova\sphinxhyphen{}Timan, M. Parrinello, J Chem Phys, 130, 074101 (2009).

\phantomsection\label{\detokenize{fix_pimd:ceriotti3}}
\sphinxAtStartPar
\sphinxstylestrong{(Ceriotti)} M. Ceriotti, M. Parrinello, T. Markland, D. Manolopoulos, J. Chem. Phys. 133, 124104 (2010).

\phantomsection\label{\detokenize{fix_pimd:martyna3}}
\sphinxAtStartPar
\sphinxstylestrong{(Martyna1)} G. Martyna, D. Tobias, M. Klein, J. Chem. Phys. 101, 4177 (1994).

\phantomsection\label{\detokenize{fix_pimd:martyna4}}
\sphinxAtStartPar
\sphinxstylestrong{(Martyna2)} G. Martyna, A. Hughes, M. Tuckerman, J. Chem. Phys. 110, 3275 (1999).

\phantomsection\label{\detokenize{fix_pimd:liujian}}
\sphinxAtStartPar
\sphinxstylestrong{(Liu)} J. Liu, D. Li, X. Liu, J. Chem. Phys. 145, 024103 (2016).

\sphinxstepscope

\index{fix planeforce@\spxentry{fix planeforce}}\ignorespaces 

\section{fix planeforce command}
\label{\detokenize{fix_planeforce:fix-planeforce-command}}\label{\detokenize{fix_planeforce:index-0}}\label{\detokenize{fix_planeforce::doc}}

\subsection{Syntax}
\label{\detokenize{fix_planeforce:syntax}}
\begin{sphinxVerbatim}[commandchars=\\\{\}]
\PYG{k}{fix }\PYG{n+nv+nvIdentifier}{ID}\PYG{+w}{ }\PYG{n+nv+nvIdentifier}{group\PYGZhy{}ID}\PYG{+w}{ }\PYG{n}{planeforce}\PYG{+w}{ }\PYG{n}{x}\PYG{+w}{ }\PYG{n}{y}\PYG{+w}{ }\PYG{n}{z}
\end{sphinxVerbatim}
\begin{itemize}
\item {} 
\sphinxAtStartPar
ID, group\sphinxhyphen{}ID are documented in {\hyperref[\detokenize{fix::doc}]{\sphinxcrossref{\DUrole{doc}{fix}}}} command

\item {} 
\sphinxAtStartPar
planeforce = style name of this fix command

\item {} 
\sphinxAtStartPar
x y z = 3\sphinxhyphen{}vector that is normal to the plane

\end{itemize}


\subsection{Examples}
\label{\detokenize{fix_planeforce:examples}}
\begin{sphinxVerbatim}[commandchars=\\\{\}]
\PYG{k}{fix }\PYG{n+nv+nvIdentifier}{hold}\PYG{+w}{ }\PYG{n+nv+nvIdentifier}{boundary}\PYG{+w}{ }\PYG{n}{planeforce}\PYG{+w}{ }\PYG{l+m}{1.0}\PYG{+w}{ }\PYG{l+m}{0.0}\PYG{+w}{ }\PYG{l+m}{0.0}
\end{sphinxVerbatim}


\subsection{Description}
\label{\detokenize{fix_planeforce:description}}
\sphinxAtStartPar
Adjust the forces on each atom in the group so that only the
components of force in the plane specified by the normal vector
(x,y,z) remain.  This is done by subtracting out the component of
force perpendicular to the plane.

\sphinxAtStartPar
If the initial velocity of the atom is 0.0 (or in the plane), then it
should continue to move in the plane thereafter.


\subsection{Restart, fix\_modify, output, run start/stop, minimize info}
\label{\detokenize{fix_planeforce:restart-fix-modify-output-run-start-stop-minimize-info}}
\sphinxAtStartPar
No information about this fix is written to {\hyperref[\detokenize{restart::doc}]{\sphinxcrossref{\DUrole{doc}{binary restart files}}}}.  None of the {\hyperref[\detokenize{fix_modify::doc}]{\sphinxcrossref{\DUrole{doc}{fix\_modify}}}} options
are relevant to this fix.  No global or per\sphinxhyphen{}atom quantities are stored
by this fix for access by various {\hyperref[\detokenize{Howto_output::doc}]{\sphinxcrossref{\DUrole{doc}{output commands}}}}.
No parameter of this fix can be used with the \sphinxstyleemphasis{start/stop} keywords of
the {\hyperref[\detokenize{run::doc}]{\sphinxcrossref{\DUrole{doc}{run}}}} command.

\sphinxAtStartPar
The forces due to this fix are imposed during an energy minimization,
invoked by the {\hyperref[\detokenize{minimize::doc}]{\sphinxcrossref{\DUrole{doc}{minimize}}}} command.


\subsection{Restrictions}
\label{\detokenize{fix_planeforce:restrictions}}\begin{quote}

\sphinxAtStartPar
none
\end{quote}


\subsection{Related commands}
\label{\detokenize{fix_planeforce:related-commands}}
\sphinxAtStartPar
{\hyperref[\detokenize{fix_lineforce::doc}]{\sphinxcrossref{\DUrole{doc}{fix lineforce}}}}


\subsection{Default}
\label{\detokenize{fix_planeforce:default}}
\sphinxAtStartPar
none

\sphinxstepscope

\index{fix plumed@\spxentry{fix plumed}}\ignorespaces 

\section{fix plumed command}
\label{\detokenize{fix_plumed:fix-plumed-command}}\label{\detokenize{fix_plumed:index-0}}\label{\detokenize{fix_plumed::doc}}

\subsection{Syntax}
\label{\detokenize{fix_plumed:syntax}}
\begin{sphinxVerbatim}[commandchars=\\\{\}]
\PYG{k}{fix }\PYG{n+nv+nvIdentifier}{ID}\PYG{+w}{ }\PYG{n+nv+nvIdentifier}{group\PYGZhy{}ID}\PYG{+w}{ }\PYG{n}{plumed}\PYG{+w}{ }\PYG{n}{keyword}\PYG{+w}{ }\PYG{n}{value}\PYG{+w}{ }\PYG{n}{...}
\end{sphinxVerbatim}
\begin{itemize}
\item {} 
\sphinxAtStartPar
ID, group\sphinxhyphen{}ID are documented in {\hyperref[\detokenize{fix::doc}]{\sphinxcrossref{\DUrole{doc}{fix}}}} command

\item {} 
\sphinxAtStartPar
plumed = style name of this fix command

\item {} 
\sphinxAtStartPar
keyword = \sphinxstyleemphasis{plumedfile} or \sphinxstyleemphasis{outfile}
\begin{sphinxalltt}
\sphinxstyleemphasis{plumedfile} arg = name of PLUMED input file to use (default: NULL)
\sphinxstyleemphasis{outfile} arg = name of file on which to write the PLUMED log (default: NULL)
\end{sphinxalltt}

\end{itemize}


\subsection{Examples}
\label{\detokenize{fix_plumed:examples}}
\begin{sphinxVerbatim}[commandchars=\\\{\}]
\PYG{k}{fix }\PYG{n+nv+nvIdentifier}{pl}\PYG{+w}{ }\PYG{n+nv+nvIdentifier}{all}\PYG{+w}{ }\PYG{n}{plumed}\PYG{+w}{ }\PYG{n}{plumedfile}\PYG{+w}{ }\PYG{n}{plumed.dat}\PYG{+w}{ }\PYG{n}{outfile}\PYG{+w}{ }\PYG{n}{p.log}
\end{sphinxVerbatim}


\subsection{Description}
\label{\detokenize{fix_plumed:description}}
\sphinxAtStartPar
This fix instructs LAMMPS to call the \sphinxhref{https://www.plumed.org/}{PLUMED} library, which
allows one to perform various forms of trajectory analysis on the fly
and to also use methods such as umbrella sampling and metadynamics to
enhance the sampling of phase space.

\sphinxAtStartPar
The documentation included here only describes the fix plumed command
itself.  This command is LAMMPS specific, whereas most of the
functionality implemented in PLUMED will work with a range of MD codes,
and when PLUMED is used as a stand alone code for analysis.  The full
\sphinxhref{https://www.plumed.org/doc.html}{documentation for PLUMED} is available online and included
in the PLUMED source code.  The PLUMED library development is hosted at
\sphinxurl{https://github.com/plumed/plumed2}
A detailed discussion of the code can be found in {\hyperref[\detokenize{fix_plumed:tribello}]{\sphinxcrossref{\DUrole{std,std-ref}{(Tribello)}}}}.

\sphinxAtStartPar
There is an example input for using this package with LAMMPS in the
examples/PACKAGES/plumed directory.


\bigskip\hrule\bigskip


\sphinxAtStartPar
The command to make LAMMPS call PLUMED during a run requires two keyword
value pairs pointing to the PLUMED input file and an output file for the
PLUMED log. The user must specify these arguments every time PLUMED is
to be used.  Furthermore, the fix plumed command should appear in the
LAMMPS input file \sphinxstylestrong{after} relevant input parameters (e.g. the timestep)
have been set.

\sphinxAtStartPar
The \sphinxstyleemphasis{group\sphinxhyphen{}ID} entry is ignored. LAMMPS will always pass all the atoms
to PLUMED and there can only be one instance of the plumed fix at a
time. The way the plumed fix is implemented ensures that the minimum
amount of information required is communicated.  Furthermore, PLUMED
supports multiple, completely independent collective variables, multiple
independent biases and multiple independent forms of analysis.  There is
thus really no restriction in functionality by only allowing only one
plumed fix in the LAMMPS input.

\sphinxAtStartPar
The \sphinxstyleemphasis{plumedfile} keyword allows the user to specify the name of the
PLUMED input file.  Instructions as to what should be included in a
plumed input file can be found in the \sphinxhref{https://www.plumed.org/doc.html}{documentation for PLUMED}

\sphinxAtStartPar
The \sphinxstyleemphasis{outfile} keyword allows the user to specify the name of a file in
which to output the PLUMED log.  This log file normally just repeats the
information that is contained in the input file to confirm it was
correctly read and parsed.  The names of the files in which the results
are stored from the various analysis options performed by PLUMED will
be specified by the user in the PLUMED input file.


\subsection{Restart, fix\_modify, output, run start/stop, minimize info}
\label{\detokenize{fix_plumed:restart-fix-modify-output-run-start-stop-minimize-info}}
\sphinxAtStartPar
When performing a restart of a calculation that involves PLUMED you
must include a RESTART command in the PLUMED input file as detailed in
the \sphinxhref{https://www.plumed.org/doc.html}{PLUMED documentation}.  When the restart command
is found in the PLUMED input PLUMED will append to the files that were
generated in the run that was performed previously.  No part of the
PLUMED restart data is included in the LAMMPS restart files.
Furthermore, any history dependent bias potentials that were
accumulated in previous calculations will be read in when the RESTART
command is included in the PLUMED input.

\sphinxAtStartPar
The {\hyperref[\detokenize{fix_modify::doc}]{\sphinxcrossref{\DUrole{doc}{fix\_modify}}}} \sphinxstyleemphasis{energy} option is supported by
this fix to add the energy change from the biasing force added by
PLUMED to the global potential energy of the system as part of
{\hyperref[\detokenize{thermo_style::doc}]{\sphinxcrossref{\DUrole{doc}{thermodynamic output}}}}.  The default setting for
this fix is {\hyperref[\detokenize{fix_modify::doc}]{\sphinxcrossref{\DUrole{doc}{fix\_modify energy yes}}}}.

\sphinxAtStartPar
The {\hyperref[\detokenize{fix_modify::doc}]{\sphinxcrossref{\DUrole{doc}{fix\_modify}}}} \sphinxstyleemphasis{virial} option is supported by
this fix to add the contribution from the biasing force to the global
pressure of the system via the {\hyperref[\detokenize{compute_pressure::doc}]{\sphinxcrossref{\DUrole{doc}{compute pressure}}}} command.  This can be accessed by
{\hyperref[\detokenize{thermo_style::doc}]{\sphinxcrossref{\DUrole{doc}{thermodynamic output}}}}.  The default setting for
this fix is {\hyperref[\detokenize{fix_modify::doc}]{\sphinxcrossref{\DUrole{doc}{fix\_modify virial yes}}}}.

\sphinxAtStartPar
This fix computes a global scalar which can be accessed by various
{\hyperref[\detokenize{Howto_output::doc}]{\sphinxcrossref{\DUrole{doc}{output commands}}}}.  The scalar is the PLUMED
energy mentioned above.  The scalar value calculated by this fix is
“extensive”.

\sphinxAtStartPar
Note that other quantities of interest can be output by commands that
are native to PLUMED.


\subsection{Restrictions}
\label{\detokenize{fix_plumed:restrictions}}
\sphinxAtStartPar
This fix is part of the PLUMED package.  It is only enabled if
LAMMPS was built with that package.  See the {\hyperref[\detokenize{Build_package::doc}]{\sphinxcrossref{\DUrole{doc}{Build package}}}} page for more info.

\sphinxAtStartPar
There can only be one fix plumed command active at a time.


\subsection{Related commands}
\label{\detokenize{fix_plumed:related-commands}}
\sphinxAtStartPar
{\hyperref[\detokenize{fix_smd::doc}]{\sphinxcrossref{\DUrole{doc}{fix smd}}}}
{\hyperref[\detokenize{fix_colvars::doc}]{\sphinxcrossref{\DUrole{doc}{fix colvars}}}}


\subsection{Default}
\label{\detokenize{fix_plumed:default}}
\sphinxAtStartPar
The default options are plumedfile = NULL and outfile = NULL


\bigskip\hrule\bigskip

\phantomsection\label{\detokenize{fix_plumed:tribello}}
\sphinxAtStartPar
\sphinxstylestrong{(Tribello)} G.A. Tribello, M. Bonomi, D. Branduardi, C. Camilloni and G. Bussi, Comp. Phys. Comm 185, 604 (2014)

\sphinxstepscope

\index{fix poems@\spxentry{fix poems}}\ignorespaces 

\section{fix poems command}
\label{\detokenize{fix_poems:fix-poems-command}}\label{\detokenize{fix_poems:index-0}}\label{\detokenize{fix_poems::doc}}

\subsection{Syntax}
\label{\detokenize{fix_poems:syntax}}
\begin{sphinxVerbatim}[commandchars=\\\{\}]
\PYG{k}{fix }\PYG{n+nv+nvIdentifier}{ID}\PYG{+w}{ }\PYG{n+nv+nvIdentifier}{group\PYGZhy{}ID}\PYG{+w}{ }\PYG{n}{poems}\PYG{+w}{ }\PYG{n}{keyword}\PYG{+w}{ }\PYG{n}{values}
\end{sphinxVerbatim}
\begin{itemize}
\item {} 
\sphinxAtStartPar
ID, group\sphinxhyphen{}ID are documented in {\hyperref[\detokenize{fix::doc}]{\sphinxcrossref{\DUrole{doc}{fix}}}} command

\item {} 
\sphinxAtStartPar
poems = style name of this fix command

\item {} 
\sphinxAtStartPar
keyword = \sphinxstyleemphasis{group} or \sphinxstyleemphasis{file} or \sphinxstyleemphasis{molecule}
\begin{sphinxalltt}
\sphinxstyleemphasis{group} values = list of group IDs
\sphinxstyleemphasis{molecule} values = none
\sphinxstyleemphasis{file} values = filename
\end{sphinxalltt}

\end{itemize}


\subsection{Examples}
\label{\detokenize{fix_poems:examples}}
\begin{sphinxVerbatim}[commandchars=\\\{\}]
\PYG{k}{fix }\PYG{n+nv+nvIdentifier}{3}\PYG{+w}{ }\PYG{n+nv+nvIdentifier}{fluid}\PYG{+w}{ }\PYG{n}{poems}\PYG{+w}{ }\PYG{k}{group }\PYG{n+nv+nvIdentifier}{clump1}\PYG{+w}{ }\PYG{n}{clump2}\PYG{+w}{ }\PYG{n}{clump3}
\PYG{k}{fix }\PYG{n+nv+nvIdentifier}{3}\PYG{+w}{ }\PYG{n+nv+nvIdentifier}{fluid}\PYG{+w}{ }\PYG{n}{poems}\PYG{+w}{ }\PYG{n}{file}\PYG{+w}{ }\PYG{n}{cluster.list}
\end{sphinxVerbatim}


\subsection{Description}
\label{\detokenize{fix_poems:description}}
\sphinxAtStartPar
Treats one or more sets of atoms as coupled rigid bodies.  This means
that each timestep the total force and torque on each rigid body is
computed and the coordinates and velocities of the atoms are updated
so that the collection of bodies move as a coupled set.  This can be
useful for treating a large biomolecule as a collection of connected,
coarse\sphinxhyphen{}grained particles.

\sphinxAtStartPar
The coupling, associated motion constraints, and time integration is
performed by the software package \sphinxtitleref{Parallelizable Open source
Efficient Multibody Software (POEMS)} which computes the constrained
rigid\sphinxhyphen{}body motion of articulated (jointed) multibody systems
{\hyperref[\detokenize{fix_poems:anderson}]{\sphinxcrossref{\DUrole{std,std-ref}{(Anderson)}}}}.  POEMS was written and is distributed by
Prof Kurt Anderson, his graduate student Rudranarayan Mukherjee, and
other members of his group at Rensselaer Polytechnic Institute (RPI).
Rudranarayan developed the LAMMPS/POEMS interface.  For copyright
information on POEMS and other details, please refer to the documents
in the poems directory distributed with LAMMPS.

\sphinxAtStartPar
This fix updates the positions and velocities of the rigid atoms with
a constant\sphinxhyphen{}energy time integration, so you should not update the same
atoms via other fixes (e.g. nve, nvt, npt, temp/rescale, langevin).

\sphinxAtStartPar
Each body must have a non\sphinxhyphen{}degenerate inertia tensor, which means if
must contain at least 3 non\sphinxhyphen{}collinear atoms.  Which atoms are in which
bodies can be defined via several options.

\sphinxAtStartPar
For option \sphinxstyleemphasis{group}, each of the listed groups is treated as a rigid
body.  Note that only atoms that are also in the fix group are
included in each rigid body.

\sphinxAtStartPar
For option \sphinxstyleemphasis{molecule}, each set of atoms in the group with a different
molecule ID is treated as a rigid body.

\sphinxAtStartPar
For option \sphinxstyleemphasis{file}, sets of atoms are read from the specified file and
each set is treated as a rigid body.  Each line of the file specifies
a rigid body in the following format:

\sphinxAtStartPar
ID type atom1\sphinxhyphen{}ID atom2\sphinxhyphen{}ID atom3\sphinxhyphen{}ID …

\sphinxAtStartPar
ID as an integer from 1 to M (the number of rigid bodies).  Type is
any integer; it is not used by the fix poems command.  The remaining
arguments are IDs of atoms in the rigid body, each typically from 1 to
N (the number of atoms in the system).  Only atoms that are also in
the fix group are included in each rigid body.  Blank lines and lines
that begin with ‘\#’ are skipped.

\sphinxAtStartPar
A connection between a pair of rigid bodies is inferred if one atom is
common to both bodies.  The POEMS solver treats that atom as a
spherical joint with 3 degrees of freedom.  Currently, a collection of
bodies can only be connected by joints as a linear chain.  The entire
collection of rigid bodies can represent one or more chains.  Other
connection topologies (tree, ring) are not allowed, but will be added
later.  Note that if no joints exist, it is more efficient to use the
{\hyperref[\detokenize{fix_rigid::doc}]{\sphinxcrossref{\DUrole{doc}{fix rigid}}}} command to simulate the system.

\sphinxAtStartPar
When the poems fix is defined, it will print out statistics on the
total \# of clusters, bodies, joints, atoms involved.  A cluster in
this context means a set of rigid bodies connected by joints.

\sphinxAtStartPar
For computational efficiency, you should turn off pairwise and bond
interactions within each rigid body, as they no longer contribute to
the motion.  The “neigh\_modify exclude” and “delete\_bonds” commands
can be used to do this if each rigid body is a group.

\sphinxAtStartPar
For computational efficiency, you should only define one fix poems
which includes all the desired rigid bodies.  LAMMPS will allow
multiple poems fixes to be defined, but it is more expensive.

\sphinxAtStartPar
The degrees\sphinxhyphen{}of\sphinxhyphen{}freedom removed by coupled rigid bodies are accounted
for in temperature and pressure computations.  Similarly, the rigid
body contribution to the pressure virial is also accounted for.  The
latter is only correct if forces within the bodies have been turned
off, and there is only a single fix poems defined.


\subsection{Restart, fix\_modify, output, run start/stop, minimize info}
\label{\detokenize{fix_poems:restart-fix-modify-output-run-start-stop-minimize-info}}
\sphinxAtStartPar
No information about this fix is written to {\hyperref[\detokenize{restart::doc}]{\sphinxcrossref{\DUrole{doc}{binary restart files}}}}.

\sphinxAtStartPar
The {\hyperref[\detokenize{fix_modify::doc}]{\sphinxcrossref{\DUrole{doc}{fix\_modify}}}} \sphinxstyleemphasis{virial} option is supported by
this fix to add the contribution due to the added forces and torques
on atoms to both the global pressure and per\sphinxhyphen{}atom stress of the system
via the {\hyperref[\detokenize{compute_pressure::doc}]{\sphinxcrossref{\DUrole{doc}{compute pressure}}}} and {\hyperref[\detokenize{compute_stress_atom::doc}]{\sphinxcrossref{\DUrole{doc}{compute
stress/atom}}}} commands.  The former can be
accessed by {\hyperref[\detokenize{thermo_style::doc}]{\sphinxcrossref{\DUrole{doc}{thermodynamic output}}}}.  The default
setting for this fix is {\hyperref[\detokenize{fix_modify::doc}]{\sphinxcrossref{\DUrole{doc}{fix\_modify virial yes}}}}.

\sphinxAtStartPar
The {\hyperref[\detokenize{fix_modify::doc}]{\sphinxcrossref{\DUrole{doc}{fix\_modify}}}} \sphinxstyleemphasis{bodyforces} option is supported by
this fix style to set whether per\sphinxhyphen{}body forces and torques are computed
early or late in a timestep, i.e. at the post\sphinxhyphen{}force stage or at the
final\sphinxhyphen{}integrate stage, respectively.

\sphinxAtStartPar
No global or per\sphinxhyphen{}atom quantities are stored by this fix for access by
various {\hyperref[\detokenize{Howto_output::doc}]{\sphinxcrossref{\DUrole{doc}{output commands}}}}.  No parameter of this
fix can be used with the \sphinxstyleemphasis{start/stop} keywords of the {\hyperref[\detokenize{run::doc}]{\sphinxcrossref{\DUrole{doc}{run}}}}
command.  This fix is not invoked during {\hyperref[\detokenize{minimize::doc}]{\sphinxcrossref{\DUrole{doc}{energy minimization}}}}.


\subsection{Restrictions}
\label{\detokenize{fix_poems:restrictions}}
\sphinxAtStartPar
This fix is part of the {\hyperref[\detokenize{Packages_details:pkg-poems}]{\sphinxcrossref{\DUrole{std,std-ref}{POEMS}}}} package.  It is only
enabled if LAMMPS was built with that package, which also requires the
POEMS library be built and linked with LAMMPS.  See the {\hyperref[\detokenize{Build_package::doc}]{\sphinxcrossref{\DUrole{doc}{Build
package}}}} page for more info.


\subsection{Related commands}
\label{\detokenize{fix_poems:related-commands}}
\sphinxAtStartPar
{\hyperref[\detokenize{fix_rigid::doc}]{\sphinxcrossref{\DUrole{doc}{fix rigid}}}}, {\hyperref[\detokenize{delete_bonds::doc}]{\sphinxcrossref{\DUrole{doc}{delete\_bonds}}}},
{\hyperref[\detokenize{neigh_modify::doc}]{\sphinxcrossref{\DUrole{doc}{neigh\_modify}}}} exclude


\subsection{Default}
\label{\detokenize{fix_poems:default}}
\sphinxAtStartPar
none


\bigskip\hrule\bigskip

\phantomsection\label{\detokenize{fix_poems:anderson}}
\sphinxAtStartPar
\sphinxstylestrong{(Anderson)} Anderson, Mukherjee, Critchley, Ziegler, and Lipton
“POEMS: Parallelizable Open\sphinxhyphen{}source Efficient Multibody Software “,
Engineering With Computers (2006). (\sphinxhref{https://doi.org/10.1007/s00366-006-0026-x}{link to paper})

\sphinxstepscope

\index{fix polarize/bem/gmres@\spxentry{fix polarize/bem/gmres}}\ignorespaces 
\index{fix polarize/bem/icc@\spxentry{fix polarize/bem/icc}}\ignorespaces 
\index{fix polarize/functional@\spxentry{fix polarize/functional}}\ignorespaces 

\section{fix polarize/bem/gmres command}
\label{\detokenize{fix_polarize:fix-polarize-bem-gmres-command}}\label{\detokenize{fix_polarize:index-2}}\label{\detokenize{fix_polarize:index-1}}\label{\detokenize{fix_polarize:index-0}}\label{\detokenize{fix_polarize::doc}}

\section{fix polarize/bem/icc command}
\label{\detokenize{fix_polarize:fix-polarize-bem-icc-command}}

\section{fix polarize/functional command}
\label{\detokenize{fix_polarize:fix-polarize-functional-command}}

\subsection{Syntax}
\label{\detokenize{fix_polarize:syntax}}
\begin{sphinxVerbatim}[commandchars=\\\{\}]
\PYG{k}{fix }\PYG{n+nv+nvIdentifier}{ID}\PYG{+w}{ }\PYG{n+nv+nvIdentifier}{group\PYGZhy{}ID}\PYG{+w}{ }\PYG{n}{style}\PYG{+w}{ }\PYG{n}{nevery}\PYG{+w}{ }\PYG{n}{tolerance}
\end{sphinxVerbatim}
\begin{itemize}
\item {} 
\sphinxAtStartPar
ID, group\sphinxhyphen{}ID are documented in {\hyperref[\detokenize{fix::doc}]{\sphinxcrossref{\DUrole{doc}{fix}}}} command

\item {} 
\sphinxAtStartPar
style = \sphinxstyleemphasis{polarize/bem/gmres} or \sphinxstyleemphasis{polarize/bem/icc} or \sphinxstyleemphasis{polarize/functional}

\item {} 
\sphinxAtStartPar
nevery = this fixed is invoked every this many timesteps

\item {} 
\sphinxAtStartPar
tolerance = the relative tolerance for the iterative solver to stop

\end{itemize}


\subsection{Examples}
\label{\detokenize{fix_polarize:examples}}
\begin{sphinxVerbatim}[commandchars=\\\{\}]
\PYG{k}{fix }\PYG{n+nv+nvIdentifier}{2}\PYG{+w}{ }\PYG{n+nv+nvIdentifier}{interface}\PYG{+w}{ }\PYG{n}{polarize}\PYG{o}{/}\PYG{n}{bem}\PYG{o}{/}\PYG{n}{gmres}\PYG{+w}{ }\PYG{l+m}{5}\PYG{+w}{ }\PYG{l+m}{0.0001}
\PYG{k}{fix }\PYG{n+nv+nvIdentifier}{1}\PYG{+w}{ }\PYG{n+nv+nvIdentifier}{interface}\PYG{+w}{ }\PYG{n}{polarize}\PYG{o}{/}\PYG{n}{bem}\PYG{o}{/}\PYG{n}{icc}\PYG{+w}{ }\PYG{l+m}{1}\PYG{+w}{ }\PYG{l+m}{0.0001}
\PYG{k}{fix }\PYG{n+nv+nvIdentifier}{3}\PYG{+w}{ }\PYG{n+nv+nvIdentifier}{interface}\PYG{+w}{ }\PYG{n}{polarize}\PYG{o}{/}\PYG{n}{functional}\PYG{+w}{ }\PYG{l+m}{1}\PYG{+w}{ }\PYG{l+m}{0.0001}
\end{sphinxVerbatim}

\sphinxAtStartPar
Used in input scripts:
\begin{quote}

\begin{sphinxVerbatim}[commandchars=\\\{\}]
examples/PACKAGES/dielectric/in.confined
examples/PACKAGES/dielectric/in.nopbc
\end{sphinxVerbatim}
\end{quote}


\subsection{Description}
\label{\detokenize{fix_polarize:description}}
\sphinxAtStartPar
These fixes compute induced charges at the interface between two
impermeable media with different dielectric constants. The interfaces
need to be discretized into vertices, each representing a boundary
element.  The vertices are treated as if they were regular atoms or
particles.  {\hyperref[\detokenize{atom_style::doc}]{\sphinxcrossref{\DUrole{doc}{atom\_style dielectric}}}} should be used
since it defines the additional properties of each interface particle
such as interface normal vectors, element areas, and local dielectric
mismatch.  These fixes also require the use of {\hyperref[\detokenize{pair_style::doc}]{\sphinxcrossref{\DUrole{doc}{pair\_style}}}} and {\hyperref[\detokenize{kspace_style::doc}]{\sphinxcrossref{\DUrole{doc}{kspace\_style}}}} with the
\sphinxstyleemphasis{dielectric} suffix.  At every time step, given a configuration of the
physical charges in the system (such as atoms and charged particles)
these fixes compute and update the charge of the interface
particles. The interfaces are allowed to move during the simulation if
the appropriate time integrators are also set (for example, with
{\hyperref[\detokenize{fix_rigid::doc}]{\sphinxcrossref{\DUrole{doc}{fix\_rigid}}}}).

\sphinxAtStartPar
Consider an interface between two media: one with dielectric constant of
78 (water), the other of 4 (silica). The interface is discretized into
2000 boundary elements, each represented by an interface
particle. Suppose that each interface particle has a normal unit vector
pointing from the silica medium to water.  The dielectric difference
along the normal vector is then 78 \sphinxhyphen{} 4 = 74, the mean dielectric value
is (78 + 4) / 2 = 41. Each boundary element also has its area and the
local mean curvature, which is used by these fixes for computing a
correction term in the local electric field.  To model charged
interfaces, an interface particle will have a non\sphinxhyphen{}zero charge value,
coming from its area and surface charge density, and its local dielectric
constant set to the mean dielectric value.

\sphinxAtStartPar
For non\sphinxhyphen{}interface particles such as atoms and charged particles, the
interface normal vectors, element area, and dielectric mismatch are
irrelevant and unused. Their local dielectric value is used internally
to rescale their given charge when computing the Coulombic
interactions. For instance, to simulate a cation carrying a charge of +2
(in simulation charge units) in an implicit solvent with a dielectric
constant of 40, the cation’s charge should be set to +2 and its local
dielectric constant property (defined in the {\hyperref[\detokenize{atom_style::doc}]{\sphinxcrossref{\DUrole{doc}{atom\_style dielectric}}}}) should be set to 40; there is no need to manually rescale
charge. This will produce the proper force for any {\hyperref[\detokenize{pair_style::doc}]{\sphinxcrossref{\DUrole{doc}{pair\_style}}}} with the dielectric suffix.  It is assumed that the
particles cannot pass through the interface during the simulation
because the value of the local dielectric constant property does not
change.

\sphinxAtStartPar
There are some example scripts for using these fixes with LAMMPS in the
\sphinxcode{\sphinxupquote{examples/PACKAGES/dielectric}} directory. The README file therein
contains specific details on the system setup. Note that the example
data files show the additional fields (columns) needed for
{\hyperref[\detokenize{atom_style::doc}]{\sphinxcrossref{\DUrole{doc}{atom\_style dielectric}}}} beyond the conventional fields
\sphinxstyleemphasis{id}, \sphinxstyleemphasis{mol}, \sphinxstyleemphasis{type}, \sphinxstyleemphasis{q}, \sphinxstyleemphasis{x}, \sphinxstyleemphasis{y}, and \sphinxstyleemphasis{z}.


\bigskip\hrule\bigskip


\sphinxAtStartPar
For fix \sphinxstyleemphasis{polarize/bem/gmres} and fix \sphinxstyleemphasis{polarize/bem/icc} the induced
charges of the atoms in the specified group, which are the vertices on
the interface, are computed using the equation:
\begin{equation*}
\begin{split}\sigma_b(\mathbf{s}) = \dfrac{1 - \bar{\epsilon}}{\bar{\epsilon}}
   \sigma_f(\mathbf{s}) - \epsilon_0 \dfrac{\Delta \epsilon}{\bar{\epsilon}}
   \mathbf{E}(\mathbf{s}) \cdot \mathbf{n}(\mathbf{s})\end{split}
\end{equation*}\begin{itemize}
\item {} 
\sphinxAtStartPar
\(\sigma_b\) is the induced charge density at the interface vertex \(\mathbf{s}\).

\item {} 
\sphinxAtStartPar
\(\bar{\epsilon}\) is the mean dielectric constant at the interface vertex: \(\bar{\epsilon} = (\epsilon_1 + \epsilon_2)/2\).

\item {} 
\sphinxAtStartPar
\(\Delta \epsilon\) is the dielectric constant difference at the interface vertex: \(\Delta \epsilon = \epsilon_1 - \epsilon_2\)

\item {} 
\sphinxAtStartPar
\(\sigma_f\) is the free charge density at the interface vertex

\item {} 
\sphinxAtStartPar
\(\mathbf{E}(\mathbf{s})\) is the electrical field at the vertex

\item {} 
\sphinxAtStartPar
\(\mathbf{n}(\mathbf{s})\) is the unit normal vector at the vertex pointing from medium with \(\epsilon_2\) to that with \(\epsilon_1\)

\end{itemize}

\sphinxAtStartPar
Fix \sphinxstyleemphasis{polarize/bem/gmres} employs the Generalized Minimum Residual
(GMRES) as described in {\hyperref[\detokenize{fix_polarize:barros}]{\sphinxcrossref{\DUrole{std,std-ref}{(Barros)}}}} to solve
\(\sigma_b\).

\sphinxAtStartPar
Fix \sphinxstyleemphasis{polarize/bem/icc} employs the successive over\sphinxhyphen{}relaxation algorithm
as described in {\hyperref[\detokenize{fix_polarize:tyagi}]{\sphinxcrossref{\DUrole{std,std-ref}{(Tyagi)}}}} to solve \(\sigma_b\).

\sphinxAtStartPar
The iterative solvers would terminate either when the maximum relative
change in the induced charges in consecutive iterations is below the set
tolerance, or when the number of iterations reaches \sphinxstyleemphasis{iter\_max} (see
below).

\sphinxAtStartPar
Fix \sphinxstyleemphasis{polarize/functional} employs the energy functional variation
approach as described in {\hyperref[\detokenize{fix_polarize:jadhao}]{\sphinxcrossref{\DUrole{std,std-ref}{(Jadhao)}}}} to solve
\(\sigma_b\).

\sphinxAtStartPar
The induced charges computed by these fixes are stored in the \sphinxstyleemphasis{q\_scaled} field,
and can be accessed as in the following example:

\begin{sphinxVerbatim}[commandchars=\\\{\}]
\PYG{k}{compute }\PYG{n+nv+nvIdentifier}{qs}\PYG{+w}{ }\PYG{n+nv+nvIdentifier}{all}\PYG{+w}{ }\PYG{n}{property}\PYG{o}{/}\PYG{n}{atom}\PYG{+w}{ }\PYG{n}{q\PYGZus{}scaled}
\PYG{k}{dump }\PYG{n+nv+nvIdentifier}{1}\PYG{+w}{ }\PYG{n+nv+nvIdentifier}{all}\PYG{+w}{ }\PYG{n}{custom}\PYG{+w}{ }\PYG{l+m}{1000}\PYG{+w}{ }\PYG{n}{all.txt}\PYG{+w}{ }\PYG{n}{id}\PYG{+w}{ }\PYG{n}{type}\PYG{+w}{ }\PYG{n}{q}\PYG{+w}{ }\PYG{n}{x}\PYG{+w}{ }\PYG{n}{y}\PYG{+w}{ }\PYG{n}{z}\PYG{+w}{ }\PYG{n}{c\PYGZus{}qs}
\end{sphinxVerbatim}

\sphinxAtStartPar
Note that the \sphinxstyleemphasis{q} field is the regular atom charges, which do not change
during the simulation. For interface particles, \sphinxstyleemphasis{q\_scaled} is the sum
of the real charge, divided by the local dielectric constant \sphinxstyleemphasis{epsilon},
and their induced charges. For non\sphinxhyphen{}interface particles, \sphinxstyleemphasis{q\_scaled} is
the real charge, divided by the local dielectric constant \sphinxstyleemphasis{epsilon}.

\sphinxAtStartPar
More details on the implementation of these fixes and their recommended
use are described in {\hyperref[\detokenize{fix_polarize:nguyentd}]{\sphinxcrossref{\DUrole{std,std-ref}{(NguyenTD)}}}}.


\subsection{Restart, fix\_modify, output, run start/stop, minimize info}
\label{\detokenize{fix_polarize:restart-fix-modify-output-run-start-stop-minimize-info}}
\sphinxAtStartPar
No information about this fix is written to {\hyperref[\detokenize{restart::doc}]{\sphinxcrossref{\DUrole{doc}{binary restart files}}}}.

\sphinxAtStartPar
The {\hyperref[\detokenize{fix_modify::doc}]{\sphinxcrossref{\DUrole{doc}{fix\_modify}}}} command provides the ability to modify certain
settings:
\begin{quote}
\begin{sphinxalltt}
\sphinxstyleemphasis{itr\_max} arg
   arg = maximum number of iterations for convergence
\sphinxstyleemphasis{dielectrics} ediff emean epsilon area charge
   ediff = dielectric difference or NULL
   emean = dielectric mean or NULL
   epsilon = local dielectric value or NULL
   area = element area or NULL
   charge = real interface charge or NULL
\sphinxstyleemphasis{kspace} arg = yes or no
\sphinxstyleemphasis{rand} max seed
   max = range of random induced charges to be generated
   seed = random number seed to use when generating random charge
\sphinxstyleemphasis{mr} arg
   arg = maximum number of q\sphinxhyphen{}vectors to use when solving (GMRES only)
\sphinxstyleemphasis{omega} arg
   arg = relaxation parameter to use when iterating (ICC only)
\end{sphinxalltt}
\end{quote}

\sphinxAtStartPar
The \sphinxstyleemphasis{itr\_max} keyword sets the max number of iterations to be used for
solving each step.

\sphinxAtStartPar
The \sphinxstyleemphasis{dielectrics} keyword allows properties of the atoms in group
\sphinxstyleemphasis{group\sphinxhyphen{}ID} to be modified. Values passed to any of the arguments
(\sphinxstyleemphasis{ediff}, \sphinxstyleemphasis{emean}, \sphinxstyleemphasis{epsilon}, \sphinxstyleemphasis{area}, \sphinxstyleemphasis{charge}) will override existing
values for all atoms in the group \sphinxstyleemphasis{group\sphinxhyphen{}ID}. Passing NULL to any of
these arguments will preserve the existing value. Note that setting the
properties of the interface this way will change the properties of all
atoms associated with the fix (all atoms in \sphinxstyleemphasis{group\sphinxhyphen{}ID}), so multiple fix
and fix\_modify commands would be needed to change the properties of two
different interfaces to different values (one fix and fix\_modify for
each interface group).

\sphinxAtStartPar
The \sphinxstyleemphasis{kspace} keyword turns on long range interactions.

\sphinxAtStartPar
If the arguments of the \sphinxstyleemphasis{rand} keyword are set, then the atoms subject
to this fix will be assigned a random initial charge in a uniform
distribution from \sphinxhyphen{}\sphinxstyleemphasis{max}/2 to \sphinxstyleemphasis{max}/2, using random number seed \sphinxstyleemphasis{seed}.

\sphinxAtStartPar
The \sphinxstyleemphasis{mr} keyword only applies to \sphinxstyleemphasis{style} = \sphinxstyleemphasis{polarize/bem/gmres}. It is
the maximum number of q\sphinxhyphen{}vectors to use when solving for the surface
charge.

\sphinxAtStartPar
The \sphinxstyleemphasis{omega} keyword only applies when using \sphinxstyleemphasis{style} =
\sphinxstyleemphasis{polarize/bem/icc}. It is a relaxation parameter defined in
{\hyperref[\detokenize{fix_polarize:tyagi}]{\sphinxcrossref{\DUrole{std,std-ref}{(Tyagi)}}}} that should generally be set between 0 and 2.

\sphinxAtStartPar
Note that the local dielectric constant (epsilon) can also be set
independently using the {\hyperref[\detokenize{set::doc}]{\sphinxcrossref{\DUrole{doc}{set}}}} command.


\bigskip\hrule\bigskip


\sphinxAtStartPar
\sphinxstyleemphasis{polarize/bem/gmres} or \sphinxstyleemphasis{polarize/bem/icc} compute a global 2\sphinxhyphen{}element
vector which can be accessed by various {\hyperref[\detokenize{Howto_output::doc}]{\sphinxcrossref{\DUrole{doc}{output commands}}}}.  The first element is the number of iterations when the
solver terminates (of which the upper bound is set by \sphinxstyleemphasis{iter\_max}). The
second element is the RMS error.


\subsection{Restrictions}
\label{\detokenize{fix_polarize:restrictions}}
\sphinxAtStartPar
These fixes are part of the DIELECTRIC package.  They are only enabled
if LAMMPS was built with that package, which requires that also the
KSPACE package is installed.  See the {\hyperref[\detokenize{Build_package::doc}]{\sphinxcrossref{\DUrole{doc}{Build package}}}} page for more info.

\sphinxAtStartPar
Note that the \sphinxstyleemphasis{polarize/bem/gmres} and \sphinxstyleemphasis{polarize/bem/icc} fixes only
support {\hyperref[\detokenize{units::doc}]{\sphinxcrossref{\DUrole{doc}{units}}}} \sphinxstyleemphasis{lj}, \sphinxstyleemphasis{real}, \sphinxstyleemphasis{metal}, \sphinxstyleemphasis{si} and \sphinxstyleemphasis{nano} at
the moment.

\sphinxAtStartPar
Note that \sphinxstyleemphasis{polarize/functional} does not yet support charged interfaces.


\subsection{Related commands}
\label{\detokenize{fix_polarize:related-commands}}
\sphinxAtStartPar
{\hyperref[\detokenize{pair_coeff::doc}]{\sphinxcrossref{\DUrole{doc}{pair\_coeff}}}}, {\hyperref[\detokenize{fix_polarize::doc}]{\sphinxcrossref{\DUrole{doc}{fix polarize}}}}, {\hyperref[\detokenize{read_data::doc}]{\sphinxcrossref{\DUrole{doc}{read\_data}}}},
{\hyperref[\detokenize{pair_dielectric::doc}]{\sphinxcrossref{\DUrole{doc}{pair\_style lj/cut/coul/long/dielectric}}}},
{\hyperref[\detokenize{kspace_style::doc}]{\sphinxcrossref{\DUrole{doc}{kspace\_style pppm/dielectric}}}},
{\hyperref[\detokenize{compute_efield_atom::doc}]{\sphinxcrossref{\DUrole{doc}{compute efield/atom}}}}


\subsection{Default}
\label{\detokenize{fix_polarize:default}}
\sphinxAtStartPar
\sphinxstyleemphasis{iter\_max} = 50

\sphinxAtStartPar
\sphinxstyleemphasis{kspace} = yes

\sphinxAtStartPar
\sphinxstyleemphasis{omega} = 0.7 (ICC only)

\sphinxAtStartPar
\sphinxstyleemphasis{mr} = \# atoms in group \sphinxstyleemphasis{group\sphinxhyphen{}ID} minus 1 (GMRES only)

\sphinxAtStartPar
No random charge initialization happens by default.


\bigskip\hrule\bigskip

\phantomsection\label{\detokenize{fix_polarize:barros}}
\sphinxAtStartPar
\sphinxstylestrong{(Barros)} Barros, Sinkovits, Luijten, J. Chem. Phys, 140, 064903 (2014)

\phantomsection\label{\detokenize{fix_polarize:tyagi}}
\sphinxAtStartPar
\sphinxstylestrong{(Tyagi)} Tyagi, Suzen, Sega, Barbosa, Kantorovich, Holm, J Chem Phys, 132, 154112 (2010)

\phantomsection\label{\detokenize{fix_polarize:jadhao}}
\sphinxAtStartPar
\sphinxstylestrong{(Jadhao)} Jadhao, Solis, Olvera de la Cruz, J Chem Phys, 138, 054119 (2013)

\phantomsection\label{\detokenize{fix_polarize:nguyentd}}
\sphinxAtStartPar
\sphinxstylestrong{(NguyenTD)} Nguyen, Li, Bagchi, Solis, Olvera de la Cruz, Comput Phys Commun 241, 80\sphinxhyphen{}19 (2019)

\sphinxstepscope

\index{fix pour@\spxentry{fix pour}}\ignorespaces 

\section{fix pour command}
\label{\detokenize{fix_pour:fix-pour-command}}\label{\detokenize{fix_pour:index-0}}\label{\detokenize{fix_pour::doc}}

\subsection{Syntax}
\label{\detokenize{fix_pour:syntax}}
\begin{sphinxVerbatim}[commandchars=\\\{\}]
\PYG{k}{fix }\PYG{n+nv+nvIdentifier}{ID}\PYG{+w}{ }\PYG{n+nv+nvIdentifier}{group\PYGZhy{}ID}\PYG{+w}{ }\PYG{n}{pour}\PYG{+w}{ }\PYG{n}{N}\PYG{+w}{ }\PYG{n}{type}\PYG{+w}{ }\PYG{n}{seed}\PYG{+w}{ }\PYG{n}{keyword}\PYG{+w}{ }\PYG{n}{values}\PYG{+w}{ }\PYG{n}{...}
\end{sphinxVerbatim}
\begin{itemize}
\item {} 
\sphinxAtStartPar
ID, group\sphinxhyphen{}ID are documented in {\hyperref[\detokenize{fix::doc}]{\sphinxcrossref{\DUrole{doc}{fix}}}} command

\item {} 
\sphinxAtStartPar
pour = style name of this fix command

\item {} 
\sphinxAtStartPar
N = \# of particles to insert

\item {} 
\sphinxAtStartPar
type = atom type to assign to inserted particles (offset for molecule insertion)

\item {} 
\sphinxAtStartPar
seed = random \# seed (positive integer)

\item {} 
\sphinxAtStartPar
one or more keyword/value pairs may be appended to args

\item {} 
\sphinxAtStartPar
keyword = \sphinxstyleemphasis{region} or \sphinxstyleemphasis{diam} or \sphinxstyleemphasis{id} or \sphinxstyleemphasis{vol} or \sphinxstyleemphasis{rate} or \sphinxstyleemphasis{dens} or \sphinxstyleemphasis{vel} or \sphinxstyleemphasis{mol} or \sphinxstyleemphasis{molfrac} or \sphinxstyleemphasis{rigid} or \sphinxstyleemphasis{shake} or \sphinxstyleemphasis{ignore}
\begin{sphinxalltt}
\sphinxstyleemphasis{region} value = region\sphinxhyphen{}ID
  region\sphinxhyphen{}ID = ID of region to use as insertion volume
\sphinxstyleemphasis{diam} values = dstyle args
  dstyle = \sphinxstyleemphasis{one} or \sphinxstyleemphasis{range} or \sphinxstyleemphasis{poly}
    \sphinxstyleemphasis{one} args = D
      D = single diameter for inserted particles (distance units)
    \sphinxstyleemphasis{range} args = Dlo Dhi
      Dlo,Dhi = range of diameters for inserted particles (distance units)
    \sphinxstyleemphasis{poly} args = Npoly D1 P1 D2 P2 ...
      Npoly = \# of (D,P) pairs
      D1,D2,... = diameter for subset of inserted particles (distance units)
      P1,P2,... = percentage of inserted particles with this diameter (0\sphinxhyphen{}1)
\sphinxstyleemphasis{id} values = idflag
  idflag = \sphinxstyleemphasis{max} or \sphinxstyleemphasis{next} = how to choose IDs for inserted particles and molecules
\sphinxstyleemphasis{vol} values = fraction Nattempt
  fraction = desired volume fraction for filling insertion volume
  Nattempt = max \# of insertion attempts per particle
\sphinxstyleemphasis{rate} value = V
  V = z velocity (3d) or y velocity (2d) at which
      insertion volume moves (velocity units)
\sphinxstyleemphasis{dens} values = Rholo Rhohi
  Rholo,Rhohi = range of densities for inserted particles (mass/volume units)
\sphinxstyleemphasis{vel} values (3d) = vxlo vxhi vylo vyhi vz
\sphinxstyleemphasis{vel} values (2d) = vxlo vxhi vy
  vxlo,vxhi = range of x velocities for inserted particles (velocity units)
  vylo,vyhi = range of y velocities for inserted particles (velocity units)
  vz = z velocity (3d) assigned to inserted particles (velocity units)
  vy = y velocity (2d) assigned to inserted particles (velocity units)
\sphinxstyleemphasis{mol} value = template\sphinxhyphen{}ID
  template\sphinxhyphen{}ID = ID of molecule template specified in a separate {\hyperref[\detokenize{molecule::doc}]{\sphinxcrossref{\DUrole{doc}{molecule}}}} command
\sphinxstyleemphasis{molfrac} values = f1 f2 ... fN
  f1 to fN = relative probability of creating each of N molecules in template\sphinxhyphen{}ID
\sphinxstyleemphasis{rigid} value = fix\sphinxhyphen{}ID
  fix\sphinxhyphen{}ID = ID of {\hyperref[\detokenize{fix_rigid::doc}]{\sphinxcrossref{\DUrole{doc}{fix rigid/small}}}} command
\sphinxstyleemphasis{shake} value = fix\sphinxhyphen{}ID
  fix\sphinxhyphen{}ID = ID of {\hyperref[\detokenize{fix_shake::doc}]{\sphinxcrossref{\DUrole{doc}{fix shake}}}} command
\sphinxstyleemphasis{ignore} value = none
  skip any line or triangle particles when detecting possible
    overlaps with inserted particles
\end{sphinxalltt}

\end{itemize}


\subsection{Examples}
\label{\detokenize{fix_pour:examples}}
\begin{sphinxVerbatim}[commandchars=\\\{\}]
\PYG{k}{fix }\PYG{n+nv+nvIdentifier}{3}\PYG{+w}{ }\PYG{n+nv+nvIdentifier}{all}\PYG{+w}{ }\PYG{n}{pour}\PYG{+w}{ }\PYG{l+m}{1000}\PYG{+w}{ }\PYG{l+m}{2}\PYG{+w}{ }\PYG{l+m}{29494}\PYG{+w}{ }\PYG{k}{region }\PYG{n+nv+nvIdentifier}{myblock}
\PYG{k}{fix }\PYG{n+nv+nvIdentifier}{2}\PYG{+w}{ }\PYG{n+nv+nvIdentifier}{all}\PYG{+w}{ }\PYG{n}{pour}\PYG{+w}{ }\PYG{l+m}{10000}\PYG{+w}{ }\PYG{l+m}{1}\PYG{+w}{ }\PYG{l+m}{19985583}\PYG{+w}{ }\PYG{k}{region }\PYG{n+nv+nvIdentifier}{disk}\PYG{+w}{ }\PYG{n}{vol}\PYG{+w}{ }\PYG{l+m}{0.33}\PYG{+w}{ }\PYG{l+m}{100}\PYG{+w}{ }\PYG{n}{rate}\PYG{+w}{ }\PYG{l+m}{1.0}\PYG{+w}{ }\PYG{n}{diam}\PYG{+w}{ }\PYG{n}{range}\PYG{+w}{ }\PYG{l+m}{0.9}\PYG{+w}{ }\PYG{l+m}{1.1}
\PYG{k}{fix }\PYG{n+nv+nvIdentifier}{2}\PYG{+w}{ }\PYG{n+nv+nvIdentifier}{all}\PYG{+w}{ }\PYG{n}{pour}\PYG{+w}{ }\PYG{l+m}{10000}\PYG{+w}{ }\PYG{l+m}{1}\PYG{+w}{ }\PYG{l+m}{19985583}\PYG{+w}{ }\PYG{k}{region }\PYG{n+nv+nvIdentifier}{disk}\PYG{+w}{ }\PYG{n}{diam}\PYG{+w}{ }\PYG{n}{poly}\PYG{+w}{ }\PYG{l+m}{2}\PYG{+w}{ }\PYG{l+m}{0.7}\PYG{+w}{ }\PYG{l+m}{0.4}\PYG{+w}{ }\PYG{l+m}{1.5}\PYG{+w}{ }\PYG{l+m}{0.6}
\PYG{k}{fix }\PYG{n+nv+nvIdentifier}{ins}\PYG{+w}{ }\PYG{n+nv+nvIdentifier}{all}\PYG{+w}{ }\PYG{n}{pour}\PYG{+w}{ }\PYG{l+m}{500}\PYG{+w}{ }\PYG{l+m}{1}\PYG{+w}{ }\PYG{l+m}{4767548}\PYG{+w}{ }\PYG{n}{vol}\PYG{+w}{ }\PYG{l+m}{0.8}\PYG{+w}{ }\PYG{l+m}{10}\PYG{+w}{ }\PYG{k}{region }\PYG{n+nv+nvIdentifier}{slab}\PYG{+w}{ }\PYG{n}{mol}\PYG{+w}{ }\PYG{n}{object}\PYG{+w}{ }\PYG{n}{rigid}\PYG{+w}{ }\PYG{n}{myRigid}
\end{sphinxVerbatim}


\subsection{Description}
\label{\detokenize{fix_pour:description}}
\sphinxAtStartPar
Insert finite\sphinxhyphen{}size particles or molecules into the simulation box
every few timesteps within a specified region until N particles or
molecules have been inserted.  This is typically used to model the
pouring of granular particles into a container under the influence of
gravity.  For the remainder of this doc page, a single inserted atom
or molecule is referred to as a “particle”.

\sphinxAtStartPar
If inserted particles are individual atoms, they are assigned the
specified atom type.  If they are molecules, the type of each atom in
the inserted molecule is specified in the file read by the
{\hyperref[\detokenize{molecule::doc}]{\sphinxcrossref{\DUrole{doc}{molecule}}}} command, and those values are added to the
specified atom type.  E.g. if the file specifies atom types 1,2,3, and
those are the atom types you want for inserted molecules, then specify
\sphinxstyleemphasis{type} = 0.  If you specify \sphinxstyleemphasis{type} = 2, the in the inserted molecule
will have atom types 3,4,5.

\sphinxAtStartPar
All atoms in the inserted particle are assigned to two groups: the
default group “all” and the group specified in the fix pour command
(which can also be “all”).

\sphinxAtStartPar
This command must use the \sphinxstyleemphasis{region} keyword to define an insertion
volume.  The specified region must have been previously defined with a
{\hyperref[\detokenize{region::doc}]{\sphinxcrossref{\DUrole{doc}{region}}}} command.  It must be of type \sphinxstyleemphasis{block} or a z\sphinxhyphen{}axis
\sphinxstyleemphasis{cylinder} and must be defined with side = \sphinxstyleemphasis{in}.  The cylinder style
of region can only be used with 3d simulations.

\sphinxAtStartPar
Individual atoms are inserted, unless the \sphinxstyleemphasis{mol} keyword is used.  It
specifies a \sphinxstyleemphasis{template\sphinxhyphen{}ID} previously defined using the
{\hyperref[\detokenize{molecule::doc}]{\sphinxcrossref{\DUrole{doc}{molecule}}}} command, which reads a file that defines the
molecule.  The coordinates, atom types, center\sphinxhyphen{}of\sphinxhyphen{}mass, moments of
inertia, etc, as well as any bond/angle/etc and special neighbor
information for the molecule can be specified in the molecule file.
See the {\hyperref[\detokenize{molecule::doc}]{\sphinxcrossref{\DUrole{doc}{molecule}}}} command for details.  The only
settings required to be in this file are the coordinates and types of
atoms in the molecule.

\sphinxAtStartPar
If the molecule template contains more than one molecule, the relative
probability of depositing each molecule can be specified by the
\sphinxstyleemphasis{molfrac} keyword.  N relative probabilities, each from 0.0 to 1.0, are
specified, where N is the number of molecules in the template.  Each
time a molecule is inserted, a random number is used to sample from
the list of relative probabilities.  The N values must sum to 1.0.

\sphinxAtStartPar
If you wish to insert molecules via the \sphinxstyleemphasis{mol} keyword, that will be
treated as rigid bodies, use the \sphinxstyleemphasis{rigid} keyword, specifying as its
value the ID of a separate {\hyperref[\detokenize{fix_rigid::doc}]{\sphinxcrossref{\DUrole{doc}{fix rigid/small}}}}
command which also appears in your input script.

\begin{sphinxadmonition}{note}{Note:}
\sphinxAtStartPar
If you wish the new rigid molecules (and other rigid molecules) to be
thermostatted correctly via {\hyperref[\detokenize{fix_rigid::doc}]{\sphinxcrossref{\DUrole{doc}{fix rigid/small/nvt}}}} or
{\hyperref[\detokenize{fix_rigid::doc}]{\sphinxcrossref{\DUrole{doc}{fix rigid/small/npt}}}}, then you need to use the
{\hyperref[\detokenize{fix_modify::doc}]{\sphinxcrossref{\DUrole{doc}{fix\_modify dynamic/dof yes}}}} command for the rigid
fix.  This is to inform that fix that the molecule count will vary
dynamically.
\end{sphinxadmonition}

\sphinxAtStartPar
If you wish to insert molecules via the \sphinxstyleemphasis{mol} keyword, that will have
their bonds or angles constrained via SHAKE, use the \sphinxstyleemphasis{shake} keyword,
specifying as its value the ID of a separate {\hyperref[\detokenize{fix_shake::doc}]{\sphinxcrossref{\DUrole{doc}{fix shake}}}} command which also appears in your input script.

\sphinxAtStartPar
Each timestep particles are inserted, they are placed randomly inside
the insertion volume so as to mimic a stream of poured particles.  If
they are molecules they are also oriented randomly.  Each atom in the
particle is tested for overlaps with existing particles, including
effects due to periodic boundary conditions if applicable.  If an
overlap is detected, another random insertion attempt is made; see the
\sphinxstyleemphasis{vol} keyword discussion below.  The larger the volume of the
insertion region, the more particles that can be inserted at any one
timestep.  Particles are inserted again after enough time has elapsed
that the previously inserted particles fall out of the insertion
volume under the influence of gravity.  Insertions continue every so
many timesteps until the desired \# of particles has been inserted.

\begin{sphinxadmonition}{note}{Note:}
\sphinxAtStartPar
If you are monitoring the temperature of a system where the particle
count is changing due to adding particles, you typically should use
the {\hyperref[\detokenize{compute_modify::doc}]{\sphinxcrossref{\DUrole{doc}{compute\_modify dynamic/dof yes}}}} command
for the temperature compute you are using.
\end{sphinxadmonition}


\bigskip\hrule\bigskip


\sphinxAtStartPar
All other keywords are optional with defaults as shown below.

\sphinxAtStartPar
The \sphinxstyleemphasis{diam} option is only used when inserting atoms and specifies the
diameters of inserted particles.  There are 3 styles: \sphinxstyleemphasis{one}, \sphinxstyleemphasis{range},
or \sphinxstyleemphasis{poly}.  For \sphinxstyleemphasis{one}, all particles will have diameter \sphinxstyleemphasis{D}.  For
\sphinxstyleemphasis{range}, the diameter of each particle will be chosen randomly and
uniformly between the specified \sphinxstyleemphasis{Dlo} and \sphinxstyleemphasis{Dhi} bounds.  For \sphinxstyleemphasis{poly}, a
series of \sphinxstyleemphasis{Npoly} diameters is specified.  For each diameter a
percentage value from 0.0 to 1.0 is also specified.  The \sphinxstyleemphasis{Npoly}
percentages must sum to 1.0.  For the example shown above with “diam 2
0.7 0.4 1.5 0.6”, all inserted particles will have a diameter of 0.7
or 1.5.  40\% of the particles will be small; 60\% will be large.

\sphinxAtStartPar
Note that for molecule insertion, the diameters of individual atoms in
the molecule can be specified in the file read by the
{\hyperref[\detokenize{molecule::doc}]{\sphinxcrossref{\DUrole{doc}{molecule}}}} command.  If not specified, the diameter of
each atom in the molecule has a default diameter of 1.0.

\sphinxAtStartPar
The \sphinxstyleemphasis{id} option has two settings which are used to determine the atom
or molecule IDs to assign to inserted particles/molecules.  In both
cases a check is done of the current system to find the maximum
current atom and molecule ID of any existing particle.  Newly inserted
particles and molecules are assigned IDs that increment those max
values.  For the \sphinxstyleemphasis{max} setting, which is the default, this check is
done at every insertion step, which allows for particles to leave the
system, and their IDs to potentially be re\sphinxhyphen{}used.  For the \sphinxstyleemphasis{next}
setting this check is done only once when the fix is specified, which
can be more efficient if you are sure particles will not be added in
some other way.

\sphinxAtStartPar
The \sphinxstyleemphasis{vol} option specifies what volume fraction of the insertion
volume will be filled with particles.  For particles with a size
specified by the \sphinxstyleemphasis{diam range} keyword, they are assumed to all be of
maximum diameter \sphinxstyleemphasis{Dhi} for purposes of computing their contribution to
the volume fraction.

\sphinxAtStartPar
The higher the volume fraction value, the more particles are inserted
each timestep.  Since inserted particles cannot overlap, the maximum
volume fraction should be no higher than about 0.6.  Each timestep
particles are inserted, LAMMPS will make up to a total of M tries to
insert the new particles without overlaps, where M = \# of inserted
particles * Nattempt.  If LAMMPS is unsuccessful at completing all
insertions, it prints a warning.

\sphinxAtStartPar
The \sphinxstyleemphasis{dens} and \sphinxstyleemphasis{vel} options enable inserted particles to have a range
of densities or xy velocities.  The specific values for a particular
inserted particle will be chosen randomly and uniformly between the
specified bounds.  Internally, the density value for a particle is
converted to a mass, based on the radius (volume) of the particle.
The \sphinxstyleemphasis{vz} or \sphinxstyleemphasis{vy} value for option \sphinxstyleemphasis{vel} assigns a z\sphinxhyphen{}velocity (3d) or
y\sphinxhyphen{}velocity (2d) to each inserted particle.

\sphinxAtStartPar
The \sphinxstyleemphasis{rate} option moves the insertion volume in the z direction (3d)
or y direction (2d).  This enables pouring particles from a
successively higher height over time.

\sphinxAtStartPar
The \sphinxstyleemphasis{ignore} option is useful when running a simulation that used line
segment (2d) or triangle (3d) particles, typically to define
boundaries for spherical granular particles to interact with.  See the
{\hyperref[\detokenize{atom_style::doc}]{\sphinxcrossref{\DUrole{doc}{atom\_style line or tri}}}} command for details.  Lines
and triangles store their size, and if the size is large it may
overlap (in a spherical sense) with the insertion region, even if the
line/triangle is oriented such that there is no actual overlap.  This
can prevent particles from being inserted.  The \sphinxstyleemphasis{ignore} keyword
causes the overlap check to skip any line or triangle particles.
Obviously you should only use it if there is in fact no overlap of the
line or triangle particles with the insertion region.


\bigskip\hrule\bigskip



\subsection{Restart, fix\_modify, output, run start/stop, minimize info}
\label{\detokenize{fix_pour:restart-fix-modify-output-run-start-stop-minimize-info}}
\sphinxAtStartPar
No information about this fix is written to {\hyperref[\detokenize{restart::doc}]{\sphinxcrossref{\DUrole{doc}{binary restart files}}}}.  This means you must be careful when restarting a
pouring simulation, when the restart file was written in the middle of
the pouring operation.  Specifically, you should use a new fix pour
command in the input script for the restarted simulation that
continues the operation.  You will need to adjust the arguments of the
original fix pour command to do this.

\sphinxAtStartPar
Also note that because the state of the random number generator is not
saved in restart files, you cannot do “exact” restarts with this fix,
where the simulation continues on the same as if no restart had taken
place.  However, in a statistical sense, a restarted simulation should
produce the same behavior if you adjust the fix pour parameters
appropriately.

\sphinxAtStartPar
None of the {\hyperref[\detokenize{fix_modify::doc}]{\sphinxcrossref{\DUrole{doc}{fix\_modify}}}} options are relevant to this
fix.  This fix computes a global scalar, which can be accessed by various
output commands.  The scalar is the cumulative number of insertions.  The
scalar value calculated by this fix is “intensive”.  No parameter
of this fix can be used with the \sphinxstyleemphasis{start/stop} keywords of the
{\hyperref[\detokenize{run::doc}]{\sphinxcrossref{\DUrole{doc}{run}}}} command.  This fix is not invoked during {\hyperref[\detokenize{minimize::doc}]{\sphinxcrossref{\DUrole{doc}{energy minimization}}}}.


\subsection{Restrictions}
\label{\detokenize{fix_pour:restrictions}}
\sphinxAtStartPar
This fix is part of the GRANULAR package.  It is only enabled if
LAMMPS was built with that package.  See the {\hyperref[\detokenize{Build_package::doc}]{\sphinxcrossref{\DUrole{doc}{Build package}}}} page for more info.

\sphinxAtStartPar
For 3d simulations, a gravity fix in the \sphinxhyphen{}z direction must be defined
for use in conjunction with this fix.  For 2d simulations, gravity
must be defined in the \sphinxhyphen{}y direction.

\sphinxAtStartPar
The specified insertion region cannot be a “dynamic” region, as
defined by the {\hyperref[\detokenize{region::doc}]{\sphinxcrossref{\DUrole{doc}{region}}}} command.


\subsection{Related commands}
\label{\detokenize{fix_pour:related-commands}}
\sphinxAtStartPar
{\hyperref[\detokenize{fix_deposit::doc}]{\sphinxcrossref{\DUrole{doc}{fix deposit}}}}, {\hyperref[\detokenize{fix_gravity::doc}]{\sphinxcrossref{\DUrole{doc}{fix gravity}}}},
{\hyperref[\detokenize{region::doc}]{\sphinxcrossref{\DUrole{doc}{region}}}}


\subsection{Default}
\label{\detokenize{fix_pour:default}}
\sphinxAtStartPar
Insertions are performed for individual particles, i.e. no \sphinxstyleemphasis{mol}
setting is defined.  If the \sphinxstyleemphasis{mol} keyword is used, the default for
\sphinxstyleemphasis{molfrac} is an equal probabilities for all molecules in the template.
Additional option defaults are diam = one 1.0, dens = 1.0 1.0, vol =
0.25 50, rate = 0.0, vel = 0.0 0.0 0.0 0.0 0.0 (for 3d), vel = 0.0 0.0 0.0
(for 2d), and id = max.

\sphinxstepscope

\index{fix precession/spin@\spxentry{fix precession/spin}}\ignorespaces 

\section{fix precession/spin command}
\label{\detokenize{fix_precession_spin:fix-precession-spin-command}}\label{\detokenize{fix_precession_spin:index-0}}\label{\detokenize{fix_precession_spin::doc}}

\subsection{Syntax}
\label{\detokenize{fix_precession_spin:syntax}}
\begin{sphinxVerbatim}[commandchars=\\\{\}]
\PYG{k}{fix }\PYG{n+nv+nvIdentifier}{ID}\PYG{+w}{ }\PYG{n+nv+nvIdentifier}{group}\PYG{+w}{ }\PYG{n}{precession}\PYG{o}{/}\PYG{n}{spin}\PYG{+w}{ }\PYG{n}{style}\PYG{+w}{ }\PYG{n}{args}
\end{sphinxVerbatim}
\begin{itemize}
\item {} 
\sphinxAtStartPar
ID, group are documented in {\hyperref[\detokenize{fix::doc}]{\sphinxcrossref{\DUrole{doc}{fix}}}} command

\item {} 
\sphinxAtStartPar
precession/spin = style name of this fix command

\item {} 
\sphinxAtStartPar
style = \sphinxstyleemphasis{zeeman} or \sphinxstyleemphasis{anisotropy} or \sphinxstyleemphasis{cubic} or \sphinxstyleemphasis{stt}
\begin{sphinxalltt}
\sphinxstyleemphasis{zeeman} args = H x y z
  H = intensity of the magnetic field (in Tesla)
  x y z = vector direction of the field
\sphinxstyleemphasis{anisotropy} args = K x y z
  K = intensity of the magnetic anisotropy (in eV)
  x y z = vector direction of the anisotropy
\sphinxstyleemphasis{cubic} args = K1 K2c n1x n1y n1x n2x n2y n2z n3x n3y n3z
  K1 and K2c = intensity of the magnetic anisotropy (in eV)
  n1x to n3z = three direction vectors of the cubic anisotropy
\sphinxstyleemphasis{stt} args = J x y z
  J = intensity of the spin\sphinxhyphen{}transfer torque field
  x y z = vector direction of the field
\end{sphinxalltt}

\end{itemize}


\subsection{Examples}
\label{\detokenize{fix_precession_spin:examples}}
\begin{sphinxVerbatim}[commandchars=\\\{\}]
\PYG{k}{fix }\PYG{n+nv+nvIdentifier}{1}\PYG{+w}{ }\PYG{n+nv+nvIdentifier}{all}\PYG{+w}{ }\PYG{n}{precession}\PYG{o}{/}\PYG{n}{spin}\PYG{+w}{ }\PYG{n}{zeeman}\PYG{+w}{ }\PYG{l+m}{0.1}\PYG{+w}{ }\PYG{l+m}{0.0}\PYG{+w}{ }\PYG{l+m}{0.0}\PYG{+w}{ }\PYG{l+m}{1.0}
\PYG{k}{fix }\PYG{n+nv+nvIdentifier}{1}\PYG{+w}{ }\PYG{n+nv+nvIdentifier}{3}\PYG{+w}{ }\PYG{n}{precession}\PYG{o}{/}\PYG{n}{spin}\PYG{+w}{ }\PYG{n}{anisotropy}\PYG{+w}{ }\PYG{l+m}{0.001}\PYG{+w}{ }\PYG{l+m}{0.0}\PYG{+w}{ }\PYG{l+m}{0.0}\PYG{+w}{ }\PYG{l+m}{1.0}
\PYG{k}{fix }\PYG{n+nv+nvIdentifier}{1}\PYG{+w}{ }\PYG{n+nv+nvIdentifier}{iron}\PYG{+w}{ }\PYG{n}{precession}\PYG{o}{/}\PYG{n}{spin}\PYG{+w}{ }\PYG{n}{cubic}\PYG{+w}{ }\PYG{l+m}{0.001}\PYG{+w}{ }\PYG{l+m}{0.0005}\PYG{+w}{ }\PYG{l+m}{1.0}\PYG{+w}{ }\PYG{l+m}{0.0}\PYG{+w}{ }\PYG{l+m}{0.0}\PYG{+w}{ }\PYG{l+m}{0.0}\PYG{+w}{ }\PYG{l+m}{1.0}\PYG{+w}{ }\PYG{l+m}{0.0}\PYG{+w}{ }\PYG{l+m}{0.0}\PYG{+w}{ }\PYG{l+m}{0.0}\PYG{+w}{ }\PYG{l+m}{1.0}
\PYG{k}{fix }\PYG{n+nv+nvIdentifier}{1}\PYG{+w}{ }\PYG{n+nv+nvIdentifier}{all}\PYG{+w}{ }\PYG{n}{precession}\PYG{o}{/}\PYG{n}{spin}\PYG{+w}{ }\PYG{n}{zeeman}\PYG{+w}{ }\PYG{l+m}{0.1}\PYG{+w}{ }\PYG{l+m}{0.0}\PYG{+w}{ }\PYG{l+m}{0.0}\PYG{+w}{ }\PYG{l+m}{1.0}\PYG{+w}{ }\PYG{n}{anisotropy}\PYG{+w}{ }\PYG{l+m}{0.001}\PYG{+w}{ }\PYG{l+m}{0.0}\PYG{+w}{ }\PYG{l+m}{0.0}\PYG{+w}{ }\PYG{l+m}{1.0}
\end{sphinxVerbatim}


\subsection{Description}
\label{\detokenize{fix_precession_spin:description}}
\sphinxAtStartPar
This fix applies a precession torque to each magnetic spin in the
group.

\sphinxAtStartPar
Style \sphinxstyleemphasis{zeeman} is used for the simulation of the interaction between
the magnetic spins in the defined group and an external magnetic
field:
\begin{equation*}
\begin{split}H_{Zeeman} = -g \sum_{i=0}^{N}\mu_{i}\, \vec{s}_{i} \cdot\vec{B}_{ext}\end{split}
\end{equation*}
\sphinxAtStartPar
with:
\begin{itemize}
\item {} 
\sphinxAtStartPar
\(\vec{B}_{ext}\) the external magnetic field (in T)

\item {} 
\sphinxAtStartPar
\(g\) the Lande factor (hard\sphinxhyphen{}coded as \(g=2.0\))

\item {} 
\sphinxAtStartPar
\(\vec{s}_i\) the unitary vector describing the orientation of spin \(i\)

\item {} 
\sphinxAtStartPar
\(\mu_i\) the atomic moment of spin \(i\) given as a multiple of the
Bohr magneton \(\mu_B\) (for example, \(\mu_i \approx 2.2\) in bulk iron).

\end{itemize}

\sphinxAtStartPar
The field value in Tesla is multiplied by the gyromagnetic
ratio, \(g \cdot \mu_B/\hbar\), converting it into a precession frequency in
rad.THz (in metal units and with \(\mu_B = 5.788\cdot 10^{-5}\)
eV/T).

\sphinxAtStartPar
As a comparison, the figure below displays the simulation of a single
spin (of norm \(\mu_i = 1.0\)) submitted to an external magnetic
field of \(\vert B_{ext}\vert = 10.0\; \mathrm{Tesla}\) (and
oriented along the z axis).  The upper plot shows the average
magnetization along the external magnetic field axis and the lower
plot the Zeeman energy, both as a function of temperature.  The
reference result is provided by the plot of the Langevin function for
the same parameters.

\noindent{\hspace*{\fill}\sphinxincludegraphics[width=600\sphinxpxdimen]{{zeeman_langevin}.jpg}\hspace*{\fill}}

\sphinxAtStartPar
The temperature effects are accounted for by connecting the spin
\(i\) to a thermal bath using a Langevin thermostat (see
{\hyperref[\detokenize{fix_langevin_spin::doc}]{\sphinxcrossref{\DUrole{doc}{fix langevin/spin}}}} for the definition of
this thermostat).

\sphinxAtStartPar
Style \sphinxstyleemphasis{anisotropy} is used to simulate an easy axis or an easy plane
for the magnetic spins in the defined group:
\begin{equation*}
\begin{split}H_{aniso} = -\sum_{{ i}=1}^{N} K_{an}(\mathbf{r}_{i})\, \left(
\vec{s}_{i} \cdot \vec{n}_{i} \right)^2\end{split}
\end{equation*}
\sphinxAtStartPar
with \(n\) defining the direction of the anisotropy, and \(K\)
(in eV) its intensity.  If \(K > 0\), an easy axis is defined, and
if \(K < 0\), an easy plane is defined.

\sphinxAtStartPar
Style \sphinxstyleemphasis{cubic} is used to simulate a cubic anisotropy, with three
possible easy axis for the magnetic spins in the defined group:
\begin{equation*}
\begin{split}H_{cubic} = -\sum_{{ i}=1}^{N} K_{1}
\Big[
\left(\vec{s}_{i} \cdot \vec{n}_1 \right)^2
\left(\vec{s}_{i} \cdot \vec{n}_2 \right)^2 +
\left(\vec{s}_{i} \cdot \vec{n}_2 \right)^2
\left(\vec{s}_{i} \cdot \vec{n}_3 \right)^2 +
\left(\vec{s}_{i} \cdot \vec{n}_1 \right)^2
\left(\vec{s}_{i} \cdot \vec{n}_3 \right)^2 \Big]
+K_{2}^{(c)} \left(\vec{s}_{i} \cdot \vec{n}_1 \right)^2
\left(\vec{s}_{i} \cdot \vec{n}_2 \right)^2
\left(\vec{s}_{i} \cdot \vec{n}_3 \right)^2\end{split}
\end{equation*}
\sphinxAtStartPar
with \(K_1\) and \(K_{2c}\) (in eV) the intensity coefficients
and \(\vec{n}_1\), \(\vec{n}_2\) and \(\vec{n}_3\)
defining the three anisotropic directions defined by the command (from
\sphinxstyleemphasis{n1x} to \sphinxstyleemphasis{n3z}).  For \(\vec{n}_1 = (1 0 0)\), \(\vec{n}_2 =
(0 1 0)\), and \(\vec{n}_3 = (0 0 1)\), \(K_1 < 0\) defines an
iron type anisotropy (easy axis along the \((0 0 1)\)\sphinxhyphen{}type cube
edges), and \(K_1 > 0\) defines a nickel type anisotropy (easy
axis along the \((1 1 1)\)\sphinxhyphen{}type cube diagonals).  \(K_2^c >
0\) also defines easy axis along the \((1 1 1)\)\sphinxhyphen{}type cube
diagonals.  See chapter 2 of {\hyperref[\detokenize{fix_precession_spin:skomski1}]{\sphinxcrossref{\DUrole{std,std-ref}{(Skomski)}}}} for more
details on cubic anisotropies.

\sphinxAtStartPar
Style \sphinxstyleemphasis{stt} is used to simulate the interaction between the spins and
a spin\sphinxhyphen{}transfer torque.
See equation (7) of {\hyperref[\detokenize{fix_precession_spin:chirac1}]{\sphinxcrossref{\DUrole{std,std-ref}{(Chirac)}}}} for more details about the
implemented spin\sphinxhyphen{}transfer torque term.

\sphinxAtStartPar
In all cases, the choice of \((x y z)\) only imposes the vector
directions for the forces. Only the direction of the vector is
important; its length is ignored (the entered vectors are
normalized).

\sphinxAtStartPar
Those styles can be combined within one single command line.

\begin{sphinxadmonition}{note}{Note:}
\sphinxAtStartPar
The norm of all vectors defined with the precession/spin command
have to be non\sphinxhyphen{}zero. For example, defining
“fix 1 all precession/spin zeeman 0.1 0.0 0.0 0.0” would result
in an error message.
Since those vector components are used to compute the inverse of the
field (or anisotropy) vector norm, setting a zero\sphinxhyphen{}vector would result
in a division by zero.
\end{sphinxadmonition}


\bigskip\hrule\bigskip



\subsection{Restart, fix\_modify, output, run start/stop, minimize info}
\label{\detokenize{fix_precession_spin:restart-fix-modify-output-run-start-stop-minimize-info}}
\sphinxAtStartPar
No information about this fix is written to {\hyperref[\detokenize{restart::doc}]{\sphinxcrossref{\DUrole{doc}{binary restart files}}}}.

\sphinxAtStartPar
The {\hyperref[\detokenize{fix_modify::doc}]{\sphinxcrossref{\DUrole{doc}{fix\_modify}}}} \sphinxstyleemphasis{energy} option is supported by
this fix to add the energy associated with the spin precession
torque to the global potential energy of the system as part of
{\hyperref[\detokenize{thermo_style::doc}]{\sphinxcrossref{\DUrole{doc}{thermodynamic output}}}}.  The default setting for
this fix is {\hyperref[\detokenize{fix_modify::doc}]{\sphinxcrossref{\DUrole{doc}{fix\_modify energy no}}}}.

\sphinxAtStartPar
This fix computes a global scalar which can be accessed by various
{\hyperref[\detokenize{Howto_output::doc}]{\sphinxcrossref{\DUrole{doc}{output commands}}}}.  The scalar is the potential
energy (in energy units) discussed in the previous paragraph.  The
scalar value is an “extensive” quantity.

\sphinxAtStartPar
No information about this fix is written to {\hyperref[\detokenize{restart::doc}]{\sphinxcrossref{\DUrole{doc}{binary restart files}}}}.


\subsection{Restrictions}
\label{\detokenize{fix_precession_spin:restrictions}}
\sphinxAtStartPar
The \sphinxstyleemphasis{precession/spin} style is part of the SPIN package.  This style
is only enabled if LAMMPS was built with this package, and if the
atom\_style “spin” was declared.  See the {\hyperref[\detokenize{Build_package::doc}]{\sphinxcrossref{\DUrole{doc}{Build package}}}} page for more info.


\subsection{Related commands}
\label{\detokenize{fix_precession_spin:related-commands}}
\sphinxAtStartPar
{\hyperref[\detokenize{atom_style::doc}]{\sphinxcrossref{\DUrole{doc}{atom\_style spin}}}}


\subsection{Default}
\label{\detokenize{fix_precession_spin:default}}
\sphinxAtStartPar
none


\bigskip\hrule\bigskip

\phantomsection\label{\detokenize{fix_precession_spin:skomski1}}
\sphinxAtStartPar
\sphinxstylestrong{(Skomski)} Skomski, R. (2008). Simple models of magnetism.
Oxford University Press.

\phantomsection\label{\detokenize{fix_precession_spin:chirac1}}
\sphinxAtStartPar
\sphinxstylestrong{(Chirac)} Chirac, Theophile, et al.  Ultrafast antiferromagnetic
switching in NiO induced by spin transfer torques.
Physical Review B 102.13 (2020): 134415.

\sphinxstepscope

\index{fix press/berendsen@\spxentry{fix press/berendsen}}\ignorespaces 

\section{fix press/berendsen command}
\label{\detokenize{fix_press_berendsen:fix-press-berendsen-command}}\label{\detokenize{fix_press_berendsen:index-0}}\label{\detokenize{fix_press_berendsen::doc}}

\subsection{Syntax}
\label{\detokenize{fix_press_berendsen:syntax}}
\begin{sphinxVerbatim}[commandchars=\\\{\}]
\PYG{k}{fix }\PYG{n+nv+nvIdentifier}{ID}\PYG{+w}{ }\PYG{n+nv+nvIdentifier}{group\PYGZhy{}ID}\PYG{+w}{ }\PYG{n}{press}\PYG{o}{/}\PYG{n}{berendsen}\PYG{+w}{ }\PYG{n}{keyword}\PYG{+w}{ }\PYG{n}{value}\PYG{+w}{ }\PYG{n}{...}
\end{sphinxVerbatim}
\begin{itemize}
\item {} 
\sphinxAtStartPar
ID, group\sphinxhyphen{}ID are documented in {\hyperref[\detokenize{fix::doc}]{\sphinxcrossref{\DUrole{doc}{fix}}}} command

\item {} 
\sphinxAtStartPar
press/berendsen = style name of this fix command
\begin{sphinxalltt}
one or more keyword value pairs may be appended
keyword = \sphinxstyleemphasis{iso} or \sphinxstyleemphasis{aniso} or \sphinxstyleemphasis{x} or \sphinxstyleemphasis{y} or \sphinxstyleemphasis{z} or \sphinxstyleemphasis{couple} or \sphinxstyleemphasis{dilate} or \sphinxstyleemphasis{modulus}
  \sphinxstyleemphasis{iso} or \sphinxstyleemphasis{aniso} values = Pstart Pstop Pdamp
    Pstart,Pstop = scalar external pressure at start/end of run (pressure units)
    Pdamp = pressure damping parameter (time units)
  \sphinxstyleemphasis{x} or \sphinxstyleemphasis{y} or \sphinxstyleemphasis{z} values = Pstart Pstop Pdamp
    Pstart,Pstop = external stress tensor component at start/end of run (pressure units)
    Pdamp = stress damping parameter (time units)
  \sphinxstyleemphasis{couple} = \sphinxstyleemphasis{none} or \sphinxstyleemphasis{xyz} or \sphinxstyleemphasis{xy} or \sphinxstyleemphasis{yz} or \sphinxstyleemphasis{xz}
  \sphinxstyleemphasis{modulus} value = bulk modulus of system (pressure units)
  \sphinxstyleemphasis{dilate} value = \sphinxstyleemphasis{all} or \sphinxstyleemphasis{partial}
\end{sphinxalltt}

\end{itemize}


\subsection{Examples}
\label{\detokenize{fix_press_berendsen:examples}}
\begin{sphinxVerbatim}[commandchars=\\\{\}]
\PYG{k}{fix }\PYG{n+nv+nvIdentifier}{1}\PYG{+w}{ }\PYG{n+nv+nvIdentifier}{all}\PYG{+w}{ }\PYG{n}{press}\PYG{o}{/}\PYG{n}{berendsen}\PYG{+w}{ }\PYG{n}{iso}\PYG{+w}{ }\PYG{l+m}{0.0}\PYG{+w}{ }\PYG{l+m}{0.0}\PYG{+w}{ }\PYG{l+m}{1000.0}
\PYG{k}{fix }\PYG{n+nv+nvIdentifier}{2}\PYG{+w}{ }\PYG{n+nv+nvIdentifier}{all}\PYG{+w}{ }\PYG{n}{press}\PYG{o}{/}\PYG{n}{berendsen}\PYG{+w}{ }\PYG{n}{aniso}\PYG{+w}{ }\PYG{l+m}{0.0}\PYG{+w}{ }\PYG{l+m}{0.0}\PYG{+w}{ }\PYG{l+m}{1000.0}\PYG{+w}{ }\PYG{n}{dilate}\PYG{+w}{ }\PYG{n}{partial}
\end{sphinxVerbatim}


\subsection{Description}
\label{\detokenize{fix_press_berendsen:description}}
\sphinxAtStartPar
Reset the pressure of the system by using a Berendsen barostat
{\hyperref[\detokenize{fix_press_berendsen:berendsen1}]{\sphinxcrossref{\DUrole{std,std-ref}{(Berendsen)}}}}, which rescales the system volume and
(optionally) the atoms coordinates within the simulation box every
timestep.

\sphinxAtStartPar
Regardless of what atoms are in the fix group, a global pressure is
computed for all atoms.  Similarly, when the size of the simulation
box is changed, all atoms are re\sphinxhyphen{}scaled to new positions, unless the
keyword \sphinxstyleemphasis{dilate} is specified with a value of \sphinxstyleemphasis{partial}, in which case
only the atoms in the fix group are re\sphinxhyphen{}scaled.  The latter can be
useful for leaving the coordinates of atoms in a solid substrate
unchanged and controlling the pressure of a surrounding fluid.

\begin{sphinxadmonition}{note}{Note:}
\sphinxAtStartPar
Unlike the {\hyperref[\detokenize{fix_nh::doc}]{\sphinxcrossref{\DUrole{doc}{fix npt}}}} or {\hyperref[\detokenize{fix_nh::doc}]{\sphinxcrossref{\DUrole{doc}{fix nph}}}}
commands which perform Nose/Hoover barostatting AND time integration,
this fix does NOT perform time integration.  It only modifies the box
size and atom coordinates to effect barostatting.  Thus you must use a
separate time integration fix, like {\hyperref[\detokenize{fix_nve::doc}]{\sphinxcrossref{\DUrole{doc}{fix nve}}}} or {\hyperref[\detokenize{fix_nh::doc}]{\sphinxcrossref{\DUrole{doc}{fix nvt}}}} to actually update the positions and velocities of
atoms.  This fix can be used in conjunction with thermostatting fixes
to control the temperature, such as {\hyperref[\detokenize{fix_nh::doc}]{\sphinxcrossref{\DUrole{doc}{fix nvt}}}} or {\hyperref[\detokenize{fix_langevin::doc}]{\sphinxcrossref{\DUrole{doc}{fix langevin}}}} or {\hyperref[\detokenize{fix_temp_berendsen::doc}]{\sphinxcrossref{\DUrole{doc}{fix temp/berendsen}}}}.
\end{sphinxadmonition}

\sphinxAtStartPar
See the {\hyperref[\detokenize{Howto_barostat::doc}]{\sphinxcrossref{\DUrole{doc}{Howto barostat}}}} page for a
discussion of different ways to perform barostatting.


\bigskip\hrule\bigskip


\sphinxAtStartPar
The barostat is specified using one or more of the \sphinxstyleemphasis{iso}, \sphinxstyleemphasis{aniso},
\sphinxstyleemphasis{x}, \sphinxstyleemphasis{y}, \sphinxstyleemphasis{z}, and \sphinxstyleemphasis{couple} keywords.  These keywords give you the
ability to specify the 3 diagonal components of an external stress
tensor, and to couple various of these components together so that the
dimensions they represent are varied together during a
constant\sphinxhyphen{}pressure simulation.  Unlike the {\hyperref[\detokenize{fix_nh::doc}]{\sphinxcrossref{\DUrole{doc}{fix npt}}}} and
{\hyperref[\detokenize{fix_nh::doc}]{\sphinxcrossref{\DUrole{doc}{fix nph}}}} commands, this fix cannot be used with triclinic
(non\sphinxhyphen{}orthogonal) simulation boxes to control all 6 components of the
general pressure tensor.

\sphinxAtStartPar
The target pressures for each of the 3 diagonal components of the
stress tensor can be specified independently via the \sphinxstyleemphasis{x}, \sphinxstyleemphasis{y}, \sphinxstyleemphasis{z},
keywords, which correspond to the 3 simulation box dimensions.  For
each component, the external pressure or tensor component at each
timestep is a ramped value during the run from \sphinxstyleemphasis{Pstart} to \sphinxstyleemphasis{Pstop}.
If a target pressure is specified for a component, then the
corresponding box dimension will change during a simulation.  For
example, if the \sphinxstyleemphasis{y} keyword is used, the y\sphinxhyphen{}box length will change.  A
box dimension will not change if that component is not specified,
although you have the option to change that dimension via the {\hyperref[\detokenize{fix_deform::doc}]{\sphinxcrossref{\DUrole{doc}{fix deform}}}} command.

\sphinxAtStartPar
For all barostat keywords, the \sphinxstyleemphasis{Pdamp} parameter determines the time
scale on which pressure is relaxed.  For example, a value of 10.0
means to relax the pressure in a timespan of (roughly) 10 time units
(tau or fs or ps \sphinxhyphen{} see the {\hyperref[\detokenize{units::doc}]{\sphinxcrossref{\DUrole{doc}{units}}}} command).

\begin{sphinxadmonition}{note}{Note:}
\sphinxAtStartPar
A Berendsen barostat will not work well for arbitrary values of
\sphinxstyleemphasis{Pdamp}.  If \sphinxstyleemphasis{Pdamp} is too small, the pressure and volume can
fluctuate wildly; if it is too large, the pressure will take a very
long time to equilibrate.  A good choice for many models is a \sphinxstyleemphasis{Pdamp}
of around 1000 timesteps.  However, note that \sphinxstyleemphasis{Pdamp} is specified in
time units, and that timesteps are NOT the same as time units for most
{\hyperref[\detokenize{units::doc}]{\sphinxcrossref{\DUrole{doc}{units}}}} settings.
\end{sphinxadmonition}

\begin{sphinxadmonition}{note}{Note:}
\sphinxAtStartPar
The relaxation time is actually also a function of the bulk
modulus of the system (inverse of isothermal compressibility).  The
bulk modulus has units of pressure and is the amount of pressure that
would need to be applied (isotropically) to reduce the volume of the
system by a factor of 2 (assuming the bulk modulus was a constant,
independent of density, which it’s not).  The bulk modulus can be set
via the keyword \sphinxstyleemphasis{modulus}.  The \sphinxstyleemphasis{Pdamp} parameter is effectively
multiplied by the bulk modulus, so if the pressure is relaxing faster
than expected or desired, increasing the bulk modulus has the same
effect as increasing \sphinxstyleemphasis{Pdamp}.  The converse is also true.  LAMMPS does
not attempt to guess a correct value of the bulk modulus; it just uses
10.0 as a default value which gives reasonable relaxation for a
Lennard\sphinxhyphen{}Jones liquid, but will be way off for other materials and way
too small for solids.  Thus you should experiment to find appropriate
values of \sphinxstyleemphasis{Pdamp} and/or the \sphinxstyleemphasis{modulus} when using this fix.
\end{sphinxadmonition}


\bigskip\hrule\bigskip


\sphinxAtStartPar
The \sphinxstyleemphasis{couple} keyword allows two or three of the diagonal components of
the pressure tensor to be “coupled” together.  The value specified
with the keyword determines which are coupled.  For example, \sphinxstyleemphasis{xz}
means the \sphinxstyleemphasis{Pxx} and \sphinxstyleemphasis{Pzz} components of the stress tensor are coupled.
\sphinxstyleemphasis{Xyz} means all 3 diagonal components are coupled.  Coupling means two
things: the instantaneous stress will be computed as an average of the
corresponding diagonal components, and the coupled box dimensions will
be changed together in lockstep, meaning coupled dimensions will be
dilated or contracted by the same percentage every timestep.  The
\sphinxstyleemphasis{Pstart}, \sphinxstyleemphasis{Pstop}, \sphinxstyleemphasis{Pdamp} parameters for any coupled dimensions must
be identical.  \sphinxstyleemphasis{Couple xyz} can be used for a 2d simulation; the \sphinxstyleemphasis{z}
dimension is simply ignored.


\bigskip\hrule\bigskip


\sphinxAtStartPar
The \sphinxstyleemphasis{iso} and \sphinxstyleemphasis{aniso} keywords are simply shortcuts that are
equivalent to specifying several other keywords together.

\sphinxAtStartPar
The keyword \sphinxstyleemphasis{iso} means couple all 3 diagonal components together when
pressure is computed (hydrostatic pressure), and dilate/contract the
dimensions together.  Using “iso Pstart Pstop Pdamp” is the same as
specifying these 4 keywords:

\begin{sphinxVerbatim}[commandchars=\\\{\}]
x Pstart Pstop Pdamp
y Pstart Pstop Pdamp
z Pstart Pstop Pdamp
couple xyz
\end{sphinxVerbatim}

\sphinxAtStartPar
The keyword \sphinxstyleemphasis{aniso} means \sphinxstyleemphasis{x}, \sphinxstyleemphasis{y}, and \sphinxstyleemphasis{z} dimensions are controlled
independently using the \sphinxstyleemphasis{Pxx}, \sphinxstyleemphasis{Pyy}, and \sphinxstyleemphasis{Pzz} components of the
stress tensor as the driving forces, and the specified scalar external
pressure.  Using “aniso Pstart Pstop Pdamp” is the same as specifying
these 4 keywords:

\begin{sphinxVerbatim}[commandchars=\\\{\}]
x Pstart Pstop Pdamp
y Pstart Pstop Pdamp
z Pstart Pstop Pdamp
couple none
\end{sphinxVerbatim}


\bigskip\hrule\bigskip


\sphinxAtStartPar
This fix computes a temperature and pressure each timestep.  To do
this, the fix creates its own computes of style “temp” and “pressure”,
as if these commands had been issued:

\begin{sphinxVerbatim}[commandchars=\\\{\}]
\PYG{k}{compute }\PYG{n+nv+nvIdentifier}{fix\PYGZhy{}ID\PYGZus{}temp}\PYG{+w}{ }\PYG{n+nv+nvIdentifier}{group\PYGZhy{}ID}\PYG{+w}{ }\PYG{n}{temp}
\PYG{k}{compute }\PYG{n+nv+nvIdentifier}{fix\PYGZhy{}ID\PYGZus{}press}\PYG{+w}{ }\PYG{n+nv+nvIdentifier}{group\PYGZhy{}ID}\PYG{+w}{ }\PYG{n}{pressure}\PYG{+w}{ }\PYG{n}{fix}\PYG{o}{\PYGZhy{}}\PYG{n}{ID\PYGZus{}temp}
\end{sphinxVerbatim}

\sphinxAtStartPar
See the {\hyperref[\detokenize{compute_temp::doc}]{\sphinxcrossref{\DUrole{doc}{compute temp}}}} and {\hyperref[\detokenize{compute_pressure::doc}]{\sphinxcrossref{\DUrole{doc}{compute pressure}}}} commands for details.  Note that the
IDs of the new computes are the fix\sphinxhyphen{}ID + underscore + “temp” or fix\_ID
+ underscore + “press”, and the group for the new computes is the same
as the fix group.

\sphinxAtStartPar
Note that these are NOT the computes used by thermodynamic output (see
the {\hyperref[\detokenize{thermo_style::doc}]{\sphinxcrossref{\DUrole{doc}{thermo\_style}}}} command) with ID = \sphinxstyleemphasis{thermo\_temp}
and \sphinxstyleemphasis{thermo\_press}.  This means you can change the attributes of this
fix’s temperature or pressure via the
{\hyperref[\detokenize{compute_modify::doc}]{\sphinxcrossref{\DUrole{doc}{compute\_modify}}}} command or print this temperature
or pressure during thermodynamic output via the {\hyperref[\detokenize{thermo_style::doc}]{\sphinxcrossref{\DUrole{doc}{thermo\_style custom}}}} command using the appropriate compute\sphinxhyphen{}ID.
It also means that changing attributes of \sphinxstyleemphasis{thermo\_temp} or
\sphinxstyleemphasis{thermo\_press} will have no effect on this fix.


\subsection{Restart, fix\_modify, output, run start/stop, minimize info}
\label{\detokenize{fix_press_berendsen:restart-fix-modify-output-run-start-stop-minimize-info}}
\sphinxAtStartPar
No information about this fix is written to {\hyperref[\detokenize{restart::doc}]{\sphinxcrossref{\DUrole{doc}{binary restart files}}}}.

\sphinxAtStartPar
The {\hyperref[\detokenize{fix_modify::doc}]{\sphinxcrossref{\DUrole{doc}{fix\_modify}}}} \sphinxstyleemphasis{temp} and \sphinxstyleemphasis{press} options are
supported by this fix.  You can use them to assign a
{\hyperref[\detokenize{compute::doc}]{\sphinxcrossref{\DUrole{doc}{compute}}}} you have defined to this fix which will be used
in its temperature and pressure calculations.  If you do this, note
that the kinetic energy derived from the compute temperature should be
consistent with the virial term computed using all atoms for the
pressure.  LAMMPS will warn you if you choose to compute temperature
on a subset of atoms.

\sphinxAtStartPar
No global or per\sphinxhyphen{}atom quantities are stored by this fix for access by
various {\hyperref[\detokenize{Howto_output::doc}]{\sphinxcrossref{\DUrole{doc}{output commands}}}}.

\sphinxAtStartPar
This fix can ramp its target pressure over multiple runs, using the
\sphinxstyleemphasis{start} and \sphinxstyleemphasis{stop} keywords of the {\hyperref[\detokenize{run::doc}]{\sphinxcrossref{\DUrole{doc}{run}}}} command.  See the
{\hyperref[\detokenize{run::doc}]{\sphinxcrossref{\DUrole{doc}{run}}}} command for details of how to do this.

\sphinxAtStartPar
This fix is not invoked during {\hyperref[\detokenize{minimize::doc}]{\sphinxcrossref{\DUrole{doc}{energy minimization}}}}.


\subsection{Restrictions}
\label{\detokenize{fix_press_berendsen:restrictions}}
\sphinxAtStartPar
Any dimension being adjusted by this fix must be periodic.


\subsection{Related commands}
\label{\detokenize{fix_press_berendsen:related-commands}}
\sphinxAtStartPar
{\hyperref[\detokenize{fix_nve::doc}]{\sphinxcrossref{\DUrole{doc}{fix nve}}}}, {\hyperref[\detokenize{fix_nh::doc}]{\sphinxcrossref{\DUrole{doc}{fix nph}}}}, {\hyperref[\detokenize{fix_nh::doc}]{\sphinxcrossref{\DUrole{doc}{fix npt}}}}, {\hyperref[\detokenize{fix_temp_berendsen::doc}]{\sphinxcrossref{\DUrole{doc}{fix temp/berendsen}}}},
{\hyperref[\detokenize{fix_modify::doc}]{\sphinxcrossref{\DUrole{doc}{fix\_modify}}}}


\subsection{Default}
\label{\detokenize{fix_press_berendsen:default}}
\sphinxAtStartPar
The keyword defaults are dilate = all, modulus = 10.0 in units of
pressure for whatever {\hyperref[\detokenize{units::doc}]{\sphinxcrossref{\DUrole{doc}{units}}}} are defined.


\bigskip\hrule\bigskip

\phantomsection\label{\detokenize{fix_press_berendsen:berendsen1}}
\sphinxAtStartPar
\sphinxstylestrong{(Berendsen)} Berendsen, Postma, van Gunsteren, DiNola, Haak, J Chem
Phys, 81, 3684 (1984).

\sphinxstepscope

\index{fix press/langevin@\spxentry{fix press/langevin}}\ignorespaces 

\section{fix press/langevin command}
\label{\detokenize{fix_press_langevin:fix-press-langevin-command}}\label{\detokenize{fix_press_langevin:index-0}}\label{\detokenize{fix_press_langevin::doc}}

\subsection{Syntax}
\label{\detokenize{fix_press_langevin:syntax}}
\begin{sphinxVerbatim}[commandchars=\\\{\}]
fix ID group\PYGZhy{}ID press/langevin keyword value ...
\end{sphinxVerbatim}
\begin{itemize}
\item {} 
\sphinxAtStartPar
ID, group\sphinxhyphen{}ID are documented in {\hyperref[\detokenize{fix::doc}]{\sphinxcrossref{\DUrole{doc}{fix}}}} command

\item {} 
\sphinxAtStartPar
press/langevin = style name of this fix command
\begin{sphinxalltt}
one or more keyword value pairs may be appended
keyword = \sphinxstyleemphasis{iso} or \sphinxstyleemphasis{aniso} or \sphinxstyleemphasis{tri} or \sphinxstyleemphasis{x} or \sphinxstyleemphasis{y} or \sphinxstyleemphasis{z} or \sphinxstyleemphasis{xy} or \sphinxstyleemphasis{xz} or \sphinxstyleemphasis{yz} or \sphinxstyleemphasis{couple} or \sphinxstyleemphasis{dilate} or \sphinxstyleemphasis{modulus} or \sphinxstyleemphasis{temp} or \sphinxstyleemphasis{flip}
  \sphinxstyleemphasis{iso} or \sphinxstyleemphasis{aniso} or \sphinxstyleemphasis{tri} values = Pstart Pstop Pdamp
    Pstart,Pstop = scalar external pressure at start/end of run (pressure units)
    Pdamp = pressure damping parameter (time units)
  \sphinxstyleemphasis{x} or \sphinxstyleemphasis{y} or \sphinxstyleemphasis{z} or \sphinxstyleemphasis{xy} or \sphinxstyleemphasis{xz} or \sphinxstyleemphasis{yz} values = Pstart Pstop Pdamp
    Pstart,Pstop = external stress tensor component at start/end of run (pressure units)
    Pdamp = pressure damping parameter
  \sphinxstyleemphasis{flip} value = \sphinxstyleemphasis{yes} or \sphinxstyleemphasis{no} = allow or disallow box flips when it becomes highly skewed
  \sphinxstyleemphasis{couple} = \sphinxstyleemphasis{none} or \sphinxstyleemphasis{xyz} or \sphinxstyleemphasis{xy} or \sphinxstyleemphasis{yz} or \sphinxstyleemphasis{xz}
  \sphinxstyleemphasis{friction} value = Friction coefficient for the barostat (time units)
  \sphinxstyleemphasis{temp} values = Tstart, Tstop, seed
  Tstart, Tstop = target temperature used for the barostat at start/end of run
  seed = seed of the random number generator
  \sphinxstyleemphasis{dilate} value = \sphinxstyleemphasis{all} or \sphinxstyleemphasis{partial}
\end{sphinxalltt}

\end{itemize}


\subsection{Examples}
\label{\detokenize{fix_press_langevin:examples}}
\begin{sphinxVerbatim}[commandchars=\\\{\}]
\PYG{k}{fix }\PYG{n+nv+nvIdentifier}{1}\PYG{+w}{ }\PYG{n+nv+nvIdentifier}{all}\PYG{+w}{ }\PYG{n}{press}\PYG{o}{/}\PYG{n}{langevin}\PYG{+w}{ }\PYG{n}{iso}\PYG{+w}{ }\PYG{l+m}{0.0}\PYG{+w}{ }\PYG{l+m}{0.0}\PYG{+w}{ }\PYG{l+m}{1000.0}\PYG{+w}{ }\PYG{n}{temp}\PYG{+w}{ }\PYG{l+m}{300}\PYG{+w}{ }\PYG{l+m}{300}\PYG{+w}{ }\PYG{l+m}{487374}
\PYG{k}{fix }\PYG{n+nv+nvIdentifier}{2}\PYG{+w}{ }\PYG{n+nv+nvIdentifier}{all}\PYG{+w}{ }\PYG{n}{press}\PYG{o}{/}\PYG{n}{langevin}\PYG{+w}{ }\PYG{n}{aniso}\PYG{+w}{ }\PYG{l+m}{0.0}\PYG{+w}{ }\PYG{l+m}{0.0}\PYG{+w}{ }\PYG{l+m}{1000.0}\PYG{+w}{ }\PYG{n}{temp}\PYG{+w}{ }\PYG{l+m}{100}\PYG{+w}{ }\PYG{l+m}{300}\PYG{+w}{ }\PYG{l+m}{238}\PYG{+w}{ }\PYG{n}{dilate}\PYG{+w}{ }\PYG{n}{partial}
\end{sphinxVerbatim}


\subsection{Description}
\label{\detokenize{fix_press_langevin:description}}
\sphinxAtStartPar
Adjust the pressure of the system by using a Langevin stochastic barostat
{\hyperref[\detokenize{fix_press_langevin:gronbech}]{\sphinxcrossref{\DUrole{std,std-ref}{(Gronbech)}}}}, which rescales the system volume and
(optionally) the atoms coordinates within the simulation box every
timestep.

\sphinxAtStartPar
The Langevin barostat couple each direction \sphinxstyleemphasis{L} with a pseudo\sphinxhyphen{}particle that obeys
the Langevin equation such as:
\begin{equation*}
\begin{split}f_P = & \frac{N k_B T_{target}}{V} + \frac{1}{V d}\sum_{i=1}^{N} \vec r_i \cdot \vec f_i - P_{target} \\
Q\ddot{L} + \alpha{}\dot{L} = & f_P + \beta(t)\\
L^{n+1} = & L^{n} + bdt\dot{L}^{n} \frac{bdt^{2}}{2Q} \\
\dot{L}^{n+1} = & \alpha\dot{L}^{n} + \frac{dt}{2Q}\left(a f^{n}_{P} + f^{n+1}_{P}\right) + \frac{b}{Q}\beta^{n+1} \\
a = & \frac{1-\frac{\alpha{}dt}{2Q}}{1+\frac{\alpha{}dt}{2Q}} \\
b = & \frac{1}{1+\frac{\alpha{}dt}{2Q}} \\
\left< \beta(t)\beta(t') \right> = & 2\alpha k_B Tdt\end{split}
\end{equation*}
\sphinxAtStartPar
Where \(dt\) is the timestep \(\dot{L}\) and \(\ddot{L}\) the first
and second derivatives of the coupled direction with regard to time,
\(\alpha\) is a friction coefficient, \(\beta\) is a random gaussian
variable and \(Q\) the effective mass of the coupled pseudoparticle. The
two first terms on the right\sphinxhyphen{}hand side of the first equation are the virial
expression of the canonical pressure. It is to be noted that the temperature
used to compute the pressure is not based on the atom velocities but rather on
the canonical
target temperature directly. This temperature is specified using the \sphinxstyleemphasis{temp}
keyword parameter and should be close to the expected target temperature of the
system.

\sphinxAtStartPar
Regardless of what atoms are in the fix group, a global pressure is
computed for all atoms. Similarly, when the size of the simulation
box is changed, all atoms are re\sphinxhyphen{}scaled to new positions, unless the
keyword \sphinxstyleemphasis{dilate} is specified with a value of \sphinxstyleemphasis{partial}, in which case
only the atoms in the fix group are re\sphinxhyphen{}scaled. The latter can be
useful for leaving the coordinates of atoms in a solid substrate
unchanged and controlling the pressure of a surrounding fluid.

\begin{sphinxadmonition}{note}{Note:}
\sphinxAtStartPar
Unlike the {\hyperref[\detokenize{fix_nh::doc}]{\sphinxcrossref{\DUrole{doc}{fix npt}}}} or {\hyperref[\detokenize{fix_nh::doc}]{\sphinxcrossref{\DUrole{doc}{fix nph}}}} commands which
perform Nose\sphinxhyphen{}Hoover barostatting AND time integration, this fix does NOT
perform time integration of the atoms but only of the barostat coupled
coordinate. It then only modifies the box size and atom coordinates to
effect barostatting. Thus you must use a separate time integration fix,
like {\hyperref[\detokenize{fix_nve::doc}]{\sphinxcrossref{\DUrole{doc}{fix nve}}}} or {\hyperref[\detokenize{fix_nh::doc}]{\sphinxcrossref{\DUrole{doc}{fix nvt}}}} to actually update
the positions and velocities of atoms.  This fix can be used in conjunction
with thermostatting fixes to control the temperature, such as {\hyperref[\detokenize{fix_nh::doc}]{\sphinxcrossref{\DUrole{doc}{fix nvt}}}} or {\hyperref[\detokenize{fix_langevin::doc}]{\sphinxcrossref{\DUrole{doc}{fix langevin}}}} or {\hyperref[\detokenize{fix_temp_berendsen::doc}]{\sphinxcrossref{\DUrole{doc}{fix temp/berendsen}}}}.
\end{sphinxadmonition}

\sphinxAtStartPar
See the {\hyperref[\detokenize{Howto_barostat::doc}]{\sphinxcrossref{\DUrole{doc}{Howto barostat}}}} page for a
discussion of different ways to perform barostatting.


\bigskip\hrule\bigskip


\sphinxAtStartPar
The barostat is specified using one or more of the \sphinxstyleemphasis{iso}, \sphinxstyleemphasis{aniso}, \sphinxstyleemphasis{tri} \sphinxstyleemphasis{x},
\sphinxstyleemphasis{y}, \sphinxstyleemphasis{z}, \sphinxstyleemphasis{xy}, \sphinxstyleemphasis{xz}, \sphinxstyleemphasis{yz}, and \sphinxstyleemphasis{couple} keywords.  These keywords give you the
ability to specify the 3 diagonal components of an external stress tensor, and
to couple various of these components together so that the dimensions they
represent are varied together during a constant\sphinxhyphen{}pressure simulation.

\sphinxAtStartPar
The target pressures for each of the 6 diagonal components of the stress tensor
can be specified independently via the \sphinxstyleemphasis{x}, \sphinxstyleemphasis{y}, \sphinxstyleemphasis{z}, keywords, which
correspond to the 3 simulation box dimensions, and the \sphinxstyleemphasis{xy}, \sphinxstyleemphasis{xz} and \sphinxstyleemphasis{yz}
keywords which corresponds to the 3 simulation box tilt factors. For each
component, the external pressure or tensor component at each timestep is a
ramped value during the run from \sphinxstyleemphasis{Pstart} to \sphinxstyleemphasis{Pstop}. If a target pressure is
specified for a component, then the corresponding box dimension will change
during a simulation.  For example, if the \sphinxstyleemphasis{y} keyword is used, the y\sphinxhyphen{}box length
will change.  A box dimension will not change if that component is not
specified, although you have the option to change that dimension via the
{\hyperref[\detokenize{fix_deform::doc}]{\sphinxcrossref{\DUrole{doc}{fix deform}}}} command.

\sphinxAtStartPar
The \sphinxstyleemphasis{Pdamp} parameter can be seen in the same way as a Nose\sphinxhyphen{}Hoover parameter as
it is used to compute the mass of the fictitious particle. Without friction,
the barostat can be compared to a single particle Nose\sphinxhyphen{}Hoover barostat and
should follow a similar decay in time. The mass of the barostat is
linked to \sphinxstyleemphasis{Pdamp} by the relation
\(Q=(N_{at}+1)\cdot{}k_BT_{target}\cdot{}P_{damp}^2\). Note that \sphinxstyleemphasis{Pdamp}
should be expressed in time units.

\begin{sphinxadmonition}{note}{Note:}
\sphinxAtStartPar
As for Berendsen barostat, a Langevin barostat will not work well for
arbitrary values of \sphinxstyleemphasis{Pdamp}.  If \sphinxstyleemphasis{Pdamp} is too small, the pressure and
volume can fluctuate wildly; if it is too large, the pressure will take a
very long time to equilibrate.  A good choice for many models is a \sphinxstyleemphasis{Pdamp}
of around 1000 timesteps.  However, note that \sphinxstyleemphasis{Pdamp} is specified in time
units, and that timesteps are NOT the same as time units for most
{\hyperref[\detokenize{units::doc}]{\sphinxcrossref{\DUrole{doc}{units}}}} settings.
\end{sphinxadmonition}


\bigskip\hrule\bigskip


\sphinxAtStartPar
The \sphinxstyleemphasis{temp} keyword sets the temperature to use in the equation of motion of the
barostat. This value is used to compute the value of the force \(f_P\) in
the equation of motion. It is important to note that this value is not the
instantaneous temperature but a target temperature that ramps from \sphinxstyleemphasis{Tstart} to
\sphinxstyleemphasis{Tstop}. Also the required argument \sphinxstyleemphasis{seed} sets the seed for the random
number generator used in the generation of the random forces.


\bigskip\hrule\bigskip


\sphinxAtStartPar
The \sphinxstyleemphasis{couple} keyword allows two or three of the diagonal components of
the pressure tensor to be “coupled” together.  The value specified
with the keyword determines which are coupled.  For example, \sphinxstyleemphasis{xz}
means the \sphinxstyleemphasis{Pxx} and \sphinxstyleemphasis{Pzz} components of the stress tensor are coupled.
\sphinxstyleemphasis{Xyz} means all 3 diagonal components are coupled.  Coupling means two
things: the instantaneous stress will be computed as an average of the
corresponding diagonal components, and the coupled box dimensions will
be changed together in lockstep, meaning coupled dimensions will be
dilated or contracted by the same percentage every timestep.  The
\sphinxstyleemphasis{Pstart}, \sphinxstyleemphasis{Pstop}, \sphinxstyleemphasis{Pdamp} parameters for any coupled dimensions must
be identical.  \sphinxstyleemphasis{Couple xyz} can be used for a 2d simulation; the \sphinxstyleemphasis{z}
dimension is simply ignored.


\bigskip\hrule\bigskip


\sphinxAtStartPar
The \sphinxstyleemphasis{iso}, \sphinxstyleemphasis{aniso} and \sphinxstyleemphasis{tri} keywords are simply shortcuts that are
equivalent to specifying several other keywords together.

\sphinxAtStartPar
The keyword \sphinxstyleemphasis{iso} means couple all 3 diagonal components together when
pressure is computed (hydrostatic pressure), and dilate/contract the
dimensions together.  Using “iso Pstart Pstop Pdamp” is the same as
specifying these 4 keywords:

\begin{sphinxVerbatim}[commandchars=\\\{\}]
x Pstart Pstop Pdamp
y Pstart Pstop Pdamp
z Pstart Pstop Pdamp
couple xyz
\end{sphinxVerbatim}

\sphinxAtStartPar
The keyword \sphinxstyleemphasis{aniso} means \sphinxstyleemphasis{x}, \sphinxstyleemphasis{y}, and \sphinxstyleemphasis{z} dimensions are controlled
independently using the \sphinxstyleemphasis{Pxx}, \sphinxstyleemphasis{Pyy}, and \sphinxstyleemphasis{Pzz} components of the
stress tensor as the driving forces, and the specified scalar external
pressure.  Using “aniso Pstart Pstop Pdamp” is the same as specifying
these 4 keywords:

\begin{sphinxVerbatim}[commandchars=\\\{\}]
x Pstart Pstop Pdamp
y Pstart Pstop Pdamp
z Pstart Pstop Pdamp
couple none
\end{sphinxVerbatim}

\sphinxAtStartPar
The keyword \sphinxstyleemphasis{tri} is the same as \sphinxstyleemphasis{aniso} but also adds the control on the
shear pressure coupled with the tilt factors.

\begin{sphinxVerbatim}[commandchars=\\\{\}]
x Pstart Pstop Pdamp
y Pstart Pstop Pdamp
z Pstart Pstop Pdamp
xy Pstart Pstop Pdamp
xz Pstart Pstop Pdamp
yz Pstart Pstop Pdamp
couple none
\end{sphinxVerbatim}


\bigskip\hrule\bigskip


\sphinxAtStartPar
The \sphinxstyleemphasis{flip} keyword allows the tilt factors for a triclinic box to
exceed half the distance of the parallel box length, as discussed
below.  If the \sphinxstyleemphasis{flip} value is set to \sphinxstyleemphasis{yes}, the bound is enforced by
flipping the box when it is exceeded.  If the \sphinxstyleemphasis{flip} value is set to
\sphinxstyleemphasis{no}, the tilt will continue to change without flipping.  Note that if
applied stress induces large deformations (e.g. in a liquid), this
means the box shape can tilt dramatically and LAMMPS will run less
efficiently, due to the large volume of communication needed to
acquire ghost atoms around a processor’s irregular\sphinxhyphen{}shaped subdomain.
For extreme values of tilt, LAMMPS may also lose atoms and generate an
error.


\bigskip\hrule\bigskip


\sphinxAtStartPar
The \sphinxstyleemphasis{friction} keyword sets the friction parameter \(\alpha\) in the
equations of motion of the barostat. For each barostat direction, the value of
\(\alpha\) depends on both \sphinxstyleemphasis{Pdamp} and \sphinxstyleemphasis{friction}. The value given as a
parameter is the Langevin characteristic time
\(\tau_{L}=\frac{Q}{\alpha}\) in time units. The langevin time can be understood as a
decorrelation time for the pressure. A long Langevin time value will make the
barostat act as an underdamped oscillator while a short value will make it
act as an overdamped oscillator. The ideal configuration would be to find
the critical parameter of the barostat. Empirically this is observed to
occur for \(\tau_{L}\approx{}P_{damp}\).  For this reason, if the \sphinxstyleemphasis{friction}
keyword is not used, the default value \sphinxstyleemphasis{Pdamp} is used for each barostat direction.


\bigskip\hrule\bigskip


\sphinxAtStartPar
This fix computes pressure each timestep. To do
this, the fix creates its own computes of style “pressure”,
as if this command had been issued:

\begin{sphinxVerbatim}[commandchars=\\\{\}]
\PYG{k}{compute }\PYG{n+nv+nvIdentifier}{fix\PYGZhy{}ID\PYGZus{}press}\PYG{+w}{ }\PYG{n+nv+nvIdentifier}{group\PYGZhy{}ID}\PYG{+w}{ }\PYG{n}{pressure}\PYG{+w}{ }\PYG{n}{NULL}\PYG{+w}{ }\PYG{n}{virial}
\end{sphinxVerbatim}

\sphinxAtStartPar
The kinetic contribution to the pressure is taken as the ensemble value
\(\frac{Nk_bT}{V}\) and computed by the fix itself.

\sphinxAtStartPar
See the {\hyperref[\detokenize{compute_pressure::doc}]{\sphinxcrossref{\DUrole{doc}{compute pressure}}}} command for details.  Note
that the IDs of the new compute is the fix\sphinxhyphen{}ID + underscore + “press” and the
group for the new computes is the same as the fix group.

\sphinxAtStartPar
Note that this is NOT the compute used by thermodynamic output (see the
{\hyperref[\detokenize{thermo_style::doc}]{\sphinxcrossref{\DUrole{doc}{thermo\_style}}}} command) with ID = \sphinxstyleemphasis{thermo\_press}. This
means you can change the attributes of this fix’s pressure via the
{\hyperref[\detokenize{compute_modify::doc}]{\sphinxcrossref{\DUrole{doc}{compute\_modify}}}} command or print this temperature or
pressure during thermodynamic output via the {\hyperref[\detokenize{thermo_style::doc}]{\sphinxcrossref{\DUrole{doc}{thermo\_style custom}}}} command using the appropriate compute\sphinxhyphen{}ID. It also means that
changing attributes of \sphinxstyleemphasis{thermo\_temp} or \sphinxstyleemphasis{thermo\_press} will have no effect on
this fix.


\subsection{Restart, fix\_modify, output, run start/stop, minimize info}
\label{\detokenize{fix_press_langevin:restart-fix-modify-output-run-start-stop-minimize-info}}
\sphinxAtStartPar
No information about this fix is written to {\hyperref[\detokenize{restart::doc}]{\sphinxcrossref{\DUrole{doc}{binary restart files}}}}.

\sphinxAtStartPar
The {\hyperref[\detokenize{fix_modify::doc}]{\sphinxcrossref{\DUrole{doc}{fix\_modify}}}} \sphinxstyleemphasis{press} option is
supported by this fix.  You can use it to assign a
{\hyperref[\detokenize{compute::doc}]{\sphinxcrossref{\DUrole{doc}{compute}}}} you have defined to this fix which will be used
in its pressure calculations.

\sphinxAtStartPar
No global or per\sphinxhyphen{}atom quantities are stored by this fix for access by
various {\hyperref[\detokenize{Howto_output::doc}]{\sphinxcrossref{\DUrole{doc}{output commands}}}}.

\sphinxAtStartPar
This fix can ramp its target pressure and temperature over multiple runs, using
the \sphinxstyleemphasis{start} and \sphinxstyleemphasis{stop} keywords of the {\hyperref[\detokenize{run::doc}]{\sphinxcrossref{\DUrole{doc}{run}}}} command.  See the
{\hyperref[\detokenize{run::doc}]{\sphinxcrossref{\DUrole{doc}{run}}}} command for details of how to do this. It is recommended that
the ramped temperature is the same as the effective temperature of the
thermostatted system. That is, if the system’s temperature is ramped by other
commands, it is recommended to do the same with this pressure control.

\sphinxAtStartPar
This fix is not invoked during {\hyperref[\detokenize{minimize::doc}]{\sphinxcrossref{\DUrole{doc}{energy minimization}}}}.


\subsection{Restrictions}
\label{\detokenize{fix_press_langevin:restrictions}}
\sphinxAtStartPar
Any dimension being adjusted by this fix must be periodic.


\subsection{Related commands}
\label{\detokenize{fix_press_langevin:related-commands}}
\sphinxAtStartPar
{\hyperref[\detokenize{fix_press_berendsen::doc}]{\sphinxcrossref{\DUrole{doc}{fix press/berendsen}}}},
{\hyperref[\detokenize{fix_nve::doc}]{\sphinxcrossref{\DUrole{doc}{fix nve}}}}, {\hyperref[\detokenize{fix_nh::doc}]{\sphinxcrossref{\DUrole{doc}{fix nph}}}}, {\hyperref[\detokenize{fix_nh::doc}]{\sphinxcrossref{\DUrole{doc}{fix npt}}}}, {\hyperref[\detokenize{fix_langevin::doc}]{\sphinxcrossref{\DUrole{doc}{fix langevin}}}},
{\hyperref[\detokenize{fix_modify::doc}]{\sphinxcrossref{\DUrole{doc}{fix\_modify}}}}


\subsection{Default}
\label{\detokenize{fix_press_langevin:default}}
\sphinxAtStartPar
The keyword defaults are \sphinxstyleemphasis{dilate} = all, \sphinxstyleemphasis{flip} = yes, and \sphinxstyleemphasis{friction} = \sphinxstyleemphasis{Pdamp}.


\bigskip\hrule\bigskip

\phantomsection\label{\detokenize{fix_press_langevin:gronbech}}
\sphinxAtStartPar
\sphinxstylestrong{(Gronbech)} Gronbech\sphinxhyphen{}Jensen, Farago, J Chem Phys, 141, 194108 (2014).

\sphinxstepscope

\index{fix print@\spxentry{fix print}}\ignorespaces 

\section{fix print command}
\label{\detokenize{fix_print:fix-print-command}}\label{\detokenize{fix_print:index-0}}\label{\detokenize{fix_print::doc}}

\subsection{Syntax}
\label{\detokenize{fix_print:syntax}}
\begin{sphinxVerbatim}[commandchars=\\\{\}]
\PYG{k}{fix }\PYG{n+nv+nvIdentifier}{ID}\PYG{+w}{ }\PYG{n+nv+nvIdentifier}{group\PYGZhy{}ID}\PYG{+w}{ }\PYG{n}{print}\PYG{+w}{ }\PYG{n}{N}\PYG{+w}{ }\PYG{n}{string}\PYG{+w}{ }\PYG{n}{keyword}\PYG{+w}{ }\PYG{n}{value}\PYG{+w}{ }\PYG{n}{...}
\end{sphinxVerbatim}
\begin{itemize}
\item {} 
\sphinxAtStartPar
ID, group\sphinxhyphen{}ID are documented in {\hyperref[\detokenize{fix::doc}]{\sphinxcrossref{\DUrole{doc}{fix}}}} command

\item {} 
\sphinxAtStartPar
print = style name of this fix command

\item {} 
\sphinxAtStartPar
N = print every N steps; N can be a variable (see below)

\item {} 
\sphinxAtStartPar
string = text string to print with optional variable names

\item {} 
\sphinxAtStartPar
zero or more keyword/value pairs may be appended

\item {} 
\sphinxAtStartPar
keyword = \sphinxstyleemphasis{file} or \sphinxstyleemphasis{append} or \sphinxstyleemphasis{screen} or \sphinxstyleemphasis{title}
\begin{sphinxalltt}
\sphinxstyleemphasis{file} value = filename
\sphinxstyleemphasis{append} value = filename
\sphinxstyleemphasis{screen} value = \sphinxstyleemphasis{yes} or \sphinxstyleemphasis{no}
\sphinxstyleemphasis{title} value = string
  string =  text to print as 1st line of output file
\end{sphinxalltt}

\end{itemize}


\subsection{Examples}
\label{\detokenize{fix_print:examples}}
\begin{sphinxVerbatim}[commandchars=\\\{\}]
\PYG{k}{fix }\PYG{n+nv+nvIdentifier}{extra}\PYG{+w}{ }\PYG{n+nv+nvIdentifier}{all}\PYG{+w}{ }\PYG{n}{print}\PYG{+w}{ }\PYG{l+m}{100}\PYG{+w}{ }\PYG{l+s}{\PYGZdq{}}\PYG{l+s}{Coords of marker atom = \PYGZdl{}x \PYGZdl{}y \PYGZdl{}z}\PYG{l+s}{\PYGZdq{}}
\PYG{k}{fix }\PYG{n+nv+nvIdentifier}{extra}\PYG{+w}{ }\PYG{n+nv+nvIdentifier}{all}\PYG{+w}{ }\PYG{n}{print}\PYG{+w}{ }\PYG{l+m}{100}\PYG{+w}{ }\PYG{l+s}{\PYGZdq{}}\PYG{l+s}{Coords of marker atom = \PYGZdl{}x \PYGZdl{}y \PYGZdl{}z}\PYG{l+s}{\PYGZdq{}}\PYG{+w}{ }\PYG{n}{file}\PYG{+w}{ }\PYG{n}{coord.txt}
\end{sphinxVerbatim}


\subsection{Description}
\label{\detokenize{fix_print:description}}
\sphinxAtStartPar
Print a text string every N steps during a simulation run.  This can
be used for diagnostic purposes or as a debugging tool to monitor some
quantity during a run.  The text string must be a single argument, so
it should be enclosed in single or double quotes if it is more than
one word.  If it contains variables it must be enclosed in double
quotes to ensure they are not evaluated when the input script line is
read, but will instead be evaluated each time the string is printed.

\sphinxAtStartPar
\DUrole{versionmodified,added}{Added in version 15Jun2023: }support for vector style variables

\sphinxAtStartPar
See the {\hyperref[\detokenize{variable::doc}]{\sphinxcrossref{\DUrole{doc}{variable}}}} command for a description of
\sphinxstyleemphasis{equal} and \sphinxstyleemphasis{vector} style variables which are typically the most
useful ones to use with the print command.  Equal\sphinxhyphen{} and vector\sphinxhyphen{}style
variables can calculate formulas involving mathematical operations,
atom properties, group properties, thermodynamic properties, global
values calculated by a {\hyperref[\detokenize{compute::doc}]{\sphinxcrossref{\DUrole{doc}{compute}}}} or {\hyperref[\detokenize{fix::doc}]{\sphinxcrossref{\DUrole{doc}{fix}}}},
or references to other {\hyperref[\detokenize{variable::doc}]{\sphinxcrossref{\DUrole{doc}{variables}}}}.  Vector\sphinxhyphen{}style
variables are printed in a bracketed, comma\sphinxhyphen{}separated format,
e.g. {[}1,2,3,4{]} or {[}12.5,2,4.6,10.1{]}.

\begin{sphinxadmonition}{note}{Note:}
\sphinxAtStartPar
As discussed on the {\hyperref[\detokenize{Commands_parse::doc}]{\sphinxcrossref{\DUrole{doc}{Commands parse}}}} doc
page, the text string can use “immediate” variables, specified as
\$(formula) with parenthesis, where the numeric formula has the same
syntax as equal\sphinxhyphen{}style variables described on the {\hyperref[\detokenize{variable::doc}]{\sphinxcrossref{\DUrole{doc}{variable}}}} doc page.  This is a convenient way to evaluate a
formula immediately without using the variable command to define a
named variable and then use that variable in the text string.  The
formula can include a trailing colon and format string which
determines the precision with which the numeric value is output.
This is also explained on the {\hyperref[\detokenize{Commands_parse::doc}]{\sphinxcrossref{\DUrole{doc}{Commands parse}}}} doc page.
\end{sphinxadmonition}

\sphinxAtStartPar
Instead of a numeric value, N can be specified as an {\hyperref[\detokenize{variable::doc}]{\sphinxcrossref{\DUrole{doc}{equal\sphinxhyphen{}style
variable}}}}, which should be specified as v\_name, where name
is the variable name.  In this case, the variable is evaluated at the
beginning of a run to determine the \sphinxstylestrong{next} timestep at which the
string will be written out.  On that timestep, the variable will be
evaluated again to determine the next timestep, etc.  Thus the
variable should return timestep values.  See the stagger() and
logfreq() and stride() math functions for {\hyperref[\detokenize{variable::doc}]{\sphinxcrossref{\DUrole{doc}{equal\sphinxhyphen{}style variables}}}}, as examples of useful functions to use in this
context. For example, the following commands will print output at
timesteps 10,20,30,100,200,300,1000,2000,etc:

\begin{sphinxVerbatim}[commandchars=\\\{\}]
\PYG{k}{variable        }\PYG{n+nv+nvIdentifier}{s}\PYG{+w}{ }\PYG{n}{equal}\PYG{+w}{ }\PYG{n}{logfreq}\PYG{n+nv}{(}\PYG{n+nv}{10,3,10}\PYG{n+nv}{)}
\PYG{k}{fix }\PYG{n+nv+nvIdentifier}{extra}\PYG{+w}{ }\PYG{n+nv+nvIdentifier}{all}\PYG{+w}{ }\PYG{n}{print}\PYG{+w}{ }\PYG{n}{v\PYGZus{}s}\PYG{+w}{ }\PYG{l+s}{\PYGZdq{}}\PYG{l+s}{Coords of marker atom = \PYGZdl{}x \PYGZdl{}y \PYGZdl{}z}\PYG{l+s}{\PYGZdq{}}
\end{sphinxVerbatim}

\sphinxAtStartPar
The specified group\sphinxhyphen{}ID is ignored by this fix.

\sphinxAtStartPar
If the \sphinxstyleemphasis{file} or \sphinxstyleemphasis{append} keyword is used, a filename is specified to
which the output generated by this fix will be written.  If \sphinxstyleemphasis{file} is
used, then the filename is overwritten if it already exists.  If
\sphinxstyleemphasis{append} is used, then the filename is appended to if it already
exists, or created if it does not exist.

\sphinxAtStartPar
If the \sphinxstyleemphasis{screen} keyword is used, output by this fix to the screen and
logfile can be turned on or off as desired.

\sphinxAtStartPar
The \sphinxstyleemphasis{title} keyword allow specification of the string that will be
printed as the first line of the output file, assuming the \sphinxstyleemphasis{file}
keyword was used.  By default, the title line is as follows:

\begin{sphinxVerbatim}[commandchars=\\\{\}]
\PYGZsh{} Fix print output for fix ID
\end{sphinxVerbatim}

\sphinxAtStartPar
where ID is replaced with the fix\sphinxhyphen{}ID.


\subsection{Restart, fix\_modify, output, run start/stop, minimize info}
\label{\detokenize{fix_print:restart-fix-modify-output-run-start-stop-minimize-info}}
\sphinxAtStartPar
No information about this fix is written to {\hyperref[\detokenize{restart::doc}]{\sphinxcrossref{\DUrole{doc}{binary restart files}}}}.  None of the {\hyperref[\detokenize{fix_modify::doc}]{\sphinxcrossref{\DUrole{doc}{fix\_modify}}}} options are
relevant to this fix.  No global or per\sphinxhyphen{}atom quantities are stored by
this fix for access by various {\hyperref[\detokenize{Howto_output::doc}]{\sphinxcrossref{\DUrole{doc}{output commands}}}}.
No parameter of this fix can be used with the \sphinxstyleemphasis{start/stop} keywords of
the {\hyperref[\detokenize{run::doc}]{\sphinxcrossref{\DUrole{doc}{run}}}} command.  This fix is not invoked during
{\hyperref[\detokenize{minimize::doc}]{\sphinxcrossref{\DUrole{doc}{energy minimization}}}}.


\subsection{Restrictions}
\label{\detokenize{fix_print:restrictions}}\begin{quote}

\sphinxAtStartPar
none
\end{quote}


\subsection{Related commands}
\label{\detokenize{fix_print:related-commands}}
\sphinxAtStartPar
{\hyperref[\detokenize{variable::doc}]{\sphinxcrossref{\DUrole{doc}{variable}}}}, {\hyperref[\detokenize{print::doc}]{\sphinxcrossref{\DUrole{doc}{print}}}}


\subsection{Default}
\label{\detokenize{fix_print:default}}
\sphinxAtStartPar
The option defaults are no file output, screen = yes, and title string
as described above.

\sphinxstepscope

\index{fix propel/self@\spxentry{fix propel/self}}\ignorespaces 

\section{fix propel/self command}
\label{\detokenize{fix_propel_self:fix-propel-self-command}}\label{\detokenize{fix_propel_self:index-0}}\label{\detokenize{fix_propel_self::doc}}

\subsection{Syntax}
\label{\detokenize{fix_propel_self:syntax}}
\begin{sphinxVerbatim}[commandchars=\\\{\}]
\PYG{k}{fix }\PYG{n+nv+nvIdentifier}{ID}\PYG{+w}{ }\PYG{n+nv+nvIdentifier}{group\PYGZhy{}ID}\PYG{+w}{ }\PYG{n}{propel}\PYG{o}{/}\PYG{n}{self}\PYG{+w}{ }\PYG{n}{mode}\PYG{+w}{ }\PYG{n}{magnitude}\PYG{+w}{ }\PYG{n}{keyword}\PYG{+w}{ }\PYG{n}{values}
\end{sphinxVerbatim}
\begin{itemize}
\item {} 
\sphinxAtStartPar
ID, group\sphinxhyphen{}ID are documented in {\hyperref[\detokenize{fix::doc}]{\sphinxcrossref{\DUrole{doc}{fix}}}} command

\item {} 
\sphinxAtStartPar
propel/self = style name of this fix command

\item {} 
\sphinxAtStartPar
mode = \sphinxstyleemphasis{dipole} or \sphinxstyleemphasis{velocity} or \sphinxstyleemphasis{quat}

\item {} 
\sphinxAtStartPar
magnitude = magnitude of self\sphinxhyphen{}propulsion force

\item {} 
\sphinxAtStartPar
zero or one keyword/value pairs may be appended

\item {} 
\sphinxAtStartPar
keyword = \sphinxstyleemphasis{qvector}
\begin{sphinxalltt}
\sphinxstyleemphasis{qvector} value = direction of force in ellipsoid frame
 \sphinxstyleemphasis{sx}, \sphinxstyleemphasis{sy}, \sphinxstyleemphasis{sz} = components of \sphinxstyleemphasis{qvector}
\end{sphinxalltt}

\end{itemize}


\subsection{Examples}
\label{\detokenize{fix_propel_self:examples}}
\begin{sphinxVerbatim}[commandchars=\\\{\}]
\PYG{k}{fix }\PYG{n+nv+nvIdentifier}{active}\PYG{+w}{ }\PYG{n+nv+nvIdentifier}{all}\PYG{+w}{ }\PYG{n}{propel}\PYG{o}{/}\PYG{n}{self}\PYG{+w}{ }\PYG{n}{dipole}\PYG{+w}{ }\PYG{l+m}{40.0}
\PYG{k}{fix }\PYG{n+nv+nvIdentifier}{active}\PYG{+w}{ }\PYG{n+nv+nvIdentifier}{all}\PYG{+w}{ }\PYG{n}{propel}\PYG{o}{/}\PYG{n}{self}\PYG{+w}{ }\PYG{k}{velocity }\PYG{n+nv+nvIdentifier}{10.0}
\PYG{k}{fix }\PYG{n+nv+nvIdentifier}{active}\PYG{+w}{ }\PYG{n+nv+nvIdentifier}{all}\PYG{+w}{ }\PYG{n}{propel}\PYG{o}{/}\PYG{n}{self}\PYG{+w}{ }\PYG{n}{quat}\PYG{+w}{ }\PYG{l+m}{15.7}\PYG{+w}{ }\PYG{n}{qvector}\PYG{+w}{ }\PYG{l+m}{1.0}\PYG{+w}{ }\PYG{l+m}{0.0}\PYG{+w}{ }\PYG{l+m}{0.0}
\end{sphinxVerbatim}


\subsection{Description}
\label{\detokenize{fix_propel_self:description}}
\sphinxAtStartPar
Add a force to each atom in the group due to a self\sphinxhyphen{}propulsion force. The
force is given by
\begin{equation*}
\begin{split}F_i = f_P e_i\end{split}
\end{equation*}
\sphinxAtStartPar
where \sphinxstyleemphasis{i} is the particle the force is being applied to, \(f_P\)
is the magnitude of the force, and \(e_i\) is the vector direction
of the force. The specification of \(e_i\) is based on which of the
three keywords (\sphinxstyleemphasis{dipole} or \sphinxstyleemphasis{velocity} or \sphinxstyleemphasis{quat}) one selects.

\sphinxAtStartPar
For mode \sphinxstyleemphasis{dipole}, \(e_i\) is just equal to
the dipole vectors of the atoms in the group. Therefore, if the dipoles
are not unit vectors, the \(e_i\) will not be unit vectors.

\begin{sphinxadmonition}{note}{Note:}
\sphinxAtStartPar
If another command changes the magnitude of the dipole, this force will
change accordingly (since \(|e_i|\) will change, which is physically
equivalent to re\sphinxhyphen{}scaling \(f_P\) while keeping \(|e_i|\) constant),
and no warning will be provided by LAMMPS. This is almost never what you
want, so ensure you are not changing dipole magnitudes with another LAMMPS
fix or pair style. Furthermore, self\sphinxhyphen{}propulsion forces (almost) always
set \(e_i\)  to be a unit vector for all times, so it’s best to set
all the dipole magnitudes to 1.0 unless you have a good reason not to
(see the {\hyperref[\detokenize{set::doc}]{\sphinxcrossref{\DUrole{doc}{set}}}} command on how to do this).
\end{sphinxadmonition}

\sphinxAtStartPar
For mode \sphinxstyleemphasis{velocity}, \(e_i\) points in the direction
of the current velocity (a unit\sphinxhyphen{}vector). This can be interpreted as a
velocity\sphinxhyphen{}dependent friction, as proposed by e.g. {\hyperref[\detokenize{fix_propel_self:erdmann1}]{\sphinxcrossref{\DUrole{std,std-ref}{(Erdmann)}}}}.

\sphinxAtStartPar
For mode \sphinxstyleemphasis{quat}, \(e_i\) points in the direction of a unit
vector, oriented in the coordinate frame of the ellipsoidal particles,
which defaults to point along the x\sphinxhyphen{}direction. This default behavior
can be changed by via the \sphinxstyleemphasis{quatvec} keyword.

\sphinxAtStartPar
The optional \sphinxstyleemphasis{quatvec} keyword specifies the direction of self\sphinxhyphen{}propulsion
via a unit vector (sx,sy,sz). The arguments \sphinxstyleemphasis{sx}, \sphinxstyleemphasis{sy}, and \sphinxstyleemphasis{sz}, are
defined within the coordinate frame of the atom’s
ellipsoid. For instance, for an ellipsoid with long axis along
its x\sphinxhyphen{}direction, if one wanted the self\sphinxhyphen{}propulsion force to also
be along this axis, set \sphinxstyleemphasis{sx} equal to 1 and \sphinxstyleemphasis{sy}, \sphinxstyleemphasis{sz} both equal
to zero. This keyword may only be specified for mode \sphinxstyleemphasis{quat}.

\begin{sphinxadmonition}{note}{Note:}
\sphinxAtStartPar
In using keyword \sphinxstyleemphasis{quatvec}, the three arguments \sphinxstyleemphasis{sx},
\sphinxstyleemphasis{sy}, and \sphinxstyleemphasis{sz} will be automatically normalized to components
of a unit vector internally to avoid users having to explicitly
do so themselves. Therefore, in mode \sphinxstyleemphasis{quat}, the vectors \(e_i\)
will always be of unit length.
\end{sphinxadmonition}

\sphinxAtStartPar
Along with adding a force contribution, this fix can also
contribute to the virial (pressure) of the system, defined as
\(f_P \sum_i <e_i . r_i>/(d V)\), where \(r_i\) is the
\sphinxstyleemphasis{unwrapped} coordinate of particle i in the case of periodic
boundary conditions. See {\hyperref[\detokenize{fix_propel_self:winkler1}]{\sphinxcrossref{\DUrole{std,std-ref}{(Winkler)}}}} for a
discussion of this active pressure contribution.

\sphinxAtStartPar
For modes \sphinxstyleemphasis{dipole} and \sphinxstyleemphasis{quat}, this fix is by default
included in pressure computations.

\sphinxAtStartPar
For mode \sphinxstyleemphasis{velocity}, this fix is by default not included
in pressure computations.

\begin{sphinxadmonition}{note}{Note:}
\sphinxAtStartPar
In contrast to equilibrium systems, pressure of active systems
in general depends on the geometry of the container.
The active pressure contribution as calculated in this fix
is only valid for certain boundary conditions (spherical
walls, rectangular walls, or periodic boundary conditions).
For other geometries, the pressure must be measured via
explicit calculation of the force per unit area on a wall,
and so one must not calculate it using this fix.
(Use {\hyperref[\detokenize{fix_modify::doc}]{\sphinxcrossref{\DUrole{doc}{fix\_modify}}}} as described below
to turn off the virial contribution of this fix). Again,
see {\hyperref[\detokenize{fix_propel_self:winkler1}]{\sphinxcrossref{\DUrole{std,std-ref}{(Winkler)}}}} for discussion of why this
is the case.

\sphinxAtStartPar
Furthermore, when dealing with active systems, the temperature
is no longer well defined. Therefore, one should ensure that
the \sphinxstyleemphasis{virial} flag is used in the
{\hyperref[\detokenize{compute_pressure::doc}]{\sphinxcrossref{\DUrole{doc}{compute pressure}}}} command (turning
off temperature contributions).
\end{sphinxadmonition}


\bigskip\hrule\bigskip



\subsection{Restart, fix\_modify, output, run start/stop, minimize info}
\label{\detokenize{fix_propel_self:restart-fix-modify-output-run-start-stop-minimize-info}}
\sphinxAtStartPar
No information about this fix is written to {\hyperref[\detokenize{restart::doc}]{\sphinxcrossref{\DUrole{doc}{binary restart files}}}}.

\sphinxAtStartPar
The {\hyperref[\detokenize{fix_modify::doc}]{\sphinxcrossref{\DUrole{doc}{fix\_modify}}}} \sphinxstyleemphasis{virial} option is supported by this
fix to add the contribution due to the added forces on atoms to the
system’s virial as part of {\hyperref[\detokenize{thermo_style::doc}]{\sphinxcrossref{\DUrole{doc}{thermodynamic output}}}}.
The default is \sphinxstyleemphasis{virial yes} for keywords \sphinxstyleemphasis{dipole} and \sphinxstyleemphasis{quat}. The
default is \sphinxstyleemphasis{virial no} for keyword \sphinxstyleemphasis{velocity}.

\sphinxAtStartPar
No parameter of this fix can be used with the \sphinxstyleemphasis{start/stop} keywords of
the {\hyperref[\detokenize{run::doc}]{\sphinxcrossref{\DUrole{doc}{run}}}} command.


\subsection{Restrictions}
\label{\detokenize{fix_propel_self:restrictions}}
\sphinxAtStartPar
With keyword \sphinxstyleemphasis{dipole}, this fix only works when the DIPOLE package is enabled.
See the {\hyperref[\detokenize{Build_package::doc}]{\sphinxcrossref{\DUrole{doc}{Build package}}}} page for more info.

\sphinxAtStartPar
This fix is part of the BROWNIAN package.  It is only enabled if
LAMMPS was built with that package.  See the {\hyperref[\detokenize{Build_package::doc}]{\sphinxcrossref{\DUrole{doc}{Build package}}}}
doc page for more info.


\subsection{Related commands}
\label{\detokenize{fix_propel_self:related-commands}}
\sphinxAtStartPar
{\hyperref[\detokenize{fix_efield::doc}]{\sphinxcrossref{\DUrole{doc}{fix efield}}}} , {\hyperref[\detokenize{fix_setforce::doc}]{\sphinxcrossref{\DUrole{doc}{fix setforce}}}},
{\hyperref[\detokenize{fix_addforce::doc}]{\sphinxcrossref{\DUrole{doc}{fix addforce}}}}


\subsection{Default}
\label{\detokenize{fix_propel_self:default}}
\sphinxAtStartPar
none


\bigskip\hrule\bigskip

\phantomsection\label{\detokenize{fix_propel_self:erdmann1}}
\sphinxAtStartPar
\sphinxstylestrong{(Erdmann)} U. Erdmann , W. Ebeling, L. Schimansky\sphinxhyphen{}Geier, and F. Schweitzer,
Eur. Phys. J. B 15, 105\sphinxhyphen{}113, 2000.

\phantomsection\label{\detokenize{fix_propel_self:winkler1}}
\sphinxAtStartPar
\sphinxstylestrong{(Winkler)} Winkler, Wysocki, and Gompper, Soft Matter, 11, 6680 (2015).

\sphinxstepscope

\index{fix property/atom@\spxentry{fix property/atom}}\ignorespaces 
\index{fix property/atom/kk@\spxentry{fix property/atom/kk}}\ignorespaces 

\section{fix property/atom command}
\label{\detokenize{fix_property_atom:fix-property-atom-command}}\label{\detokenize{fix_property_atom:index-1}}\label{\detokenize{fix_property_atom:index-0}}\label{\detokenize{fix_property_atom::doc}}
\sphinxAtStartPar
Accelerator Variants: \sphinxstyleemphasis{property/atom/kk}


\subsection{Syntax}
\label{\detokenize{fix_property_atom:syntax}}
\begin{sphinxVerbatim}[commandchars=\\\{\}]
\PYG{k}{fix }\PYG{n+nv+nvIdentifier}{ID}\PYG{+w}{ }\PYG{n+nv+nvIdentifier}{group\PYGZhy{}ID}\PYG{+w}{ }\PYG{n}{property}\PYG{o}{/}\PYG{n}{atom}\PYG{+w}{ }\PYG{n}{name1}\PYG{+w}{ }\PYG{n}{name2}\PYG{+w}{ }\PYG{n}{...}\PYG{+w}{ }\PYG{n}{keyword}\PYG{+w}{ }\PYG{n}{value}\PYG{+w}{ }\PYG{n}{...}
\end{sphinxVerbatim}
\begin{itemize}
\item {} 
\sphinxAtStartPar
ID, group\sphinxhyphen{}ID are documented in {\hyperref[\detokenize{fix::doc}]{\sphinxcrossref{\DUrole{doc}{fix}}}} command

\item {} 
\sphinxAtStartPar
property/atom = style name of this fix command

\item {} 
\sphinxAtStartPar
name1,name2,… = \sphinxstyleemphasis{mol} or \sphinxstyleemphasis{q} or \sphinxstyleemphasis{rmass} or i\_name or d\_name or i2\_name or d2\_name
\begin{sphinxalltt}
\sphinxstyleemphasis{mol} = molecule IDs
\sphinxstyleemphasis{q} = charge
\sphinxstyleemphasis{rmass} = per\sphinxhyphen{}atom mass
\sphinxstyleemphasis{temperature} = internal temperature of atom
\sphinxstyleemphasis{heatflow} = internal heat flow of atom
i\_name = new integer vector referenced by name
d\_name = new floating\sphinxhyphen{}point vector referenced by name
i2\_name = new integer array referenced by name
   i2\_name arg = N = number of columns in the array
d2\_name = new floating\sphinxhyphen{}point array referenced by name
   d2\_name arg = N = number of columns in the array
\end{sphinxalltt}

\item {} 
\sphinxAtStartPar
zero of more keyword/value pairs may be appended

\item {} 
\sphinxAtStartPar
keyword = \sphinxstyleemphasis{ghost}
\begin{sphinxalltt}
\sphinxstyleemphasis{ghost} value = \sphinxstyleemphasis{no} or \sphinxstyleemphasis{yes} for whether ghost atom info in communicated
\end{sphinxalltt}

\end{itemize}


\subsection{Examples}
\label{\detokenize{fix_property_atom:examples}}
\begin{sphinxVerbatim}[commandchars=\\\{\}]
\PYG{k}{fix }\PYG{n+nv+nvIdentifier}{1}\PYG{+w}{ }\PYG{n+nv+nvIdentifier}{all}\PYG{+w}{ }\PYG{n}{property}\PYG{o}{/}\PYG{n}{atom}\PYG{+w}{ }\PYG{n}{mol}
\PYG{k}{fix }\PYG{n+nv+nvIdentifier}{1}\PYG{+w}{ }\PYG{n+nv+nvIdentifier}{all}\PYG{+w}{ }\PYG{n}{property}\PYG{o}{/}\PYG{n}{atom}\PYG{+w}{ }\PYG{n}{i\PYGZus{}myflag1}\PYG{+w}{ }\PYG{n}{i\PYGZus{}myflag2}
\PYG{k}{fix }\PYG{n+nv+nvIdentifier}{1}\PYG{+w}{ }\PYG{n+nv+nvIdentifier}{all}\PYG{+w}{ }\PYG{n}{property}\PYG{o}{/}\PYG{n}{atom}\PYG{+w}{ }\PYG{n}{d2\PYGZus{}sxyz}\PYG{+w}{ }\PYG{l+m}{3}\PYG{+w}{ }\PYG{n}{ghost}\PYG{+w}{ }\PYG{n}{yes}
\end{sphinxVerbatim}


\subsection{Description}
\label{\detokenize{fix_property_atom:description}}
\sphinxAtStartPar
Create one or more additional per\sphinxhyphen{}atom vectors or arrays to store
information about atoms and to use during a simulation.  The specified
\sphinxstyleemphasis{group\sphinxhyphen{}ID} is ignored by this fix.

\sphinxAtStartPar
The atom style used for a simulation defines a set of per\sphinxhyphen{}atom
properties, as explained on the {\hyperref[\detokenize{atom_style::doc}]{\sphinxcrossref{\DUrole{doc}{atom\_style}}}} and
{\hyperref[\detokenize{read_data::doc}]{\sphinxcrossref{\DUrole{doc}{read\_data}}}} doc pages.  The latter command defines
these properties for each atom in the system when a data file is read.
This fix augments the set of per\sphinxhyphen{}atom properties with new custom
ones. This can be useful in several scenarios.

\sphinxAtStartPar
If the atom style does not define molecule IDs, per\sphinxhyphen{}atom charge,
per\sphinxhyphen{}atom mass, internal temperature, or internal heat flow, they can
be added using the \sphinxstyleemphasis{mol}, \sphinxstyleemphasis{q}, \sphinxstyleemphasis{rmass}, \sphinxstyleemphasis{temperature}, or \sphinxstyleemphasis{heatflow}
keywords.  This could be useful to define “molecules” to use as rigid
bodies with the {\hyperref[\detokenize{fix_rigid::doc}]{\sphinxcrossref{\DUrole{doc}{fix rigid}}}} command, or to carry
around an extra flag with atoms (stored as a molecule ID) that can be
used by various commands like {\hyperref[\detokenize{compute_chunk_atom::doc}]{\sphinxcrossref{\DUrole{doc}{compute chunk/atom}}}} to group atoms without having to use the group
command (which is limited to a total of 32 groups including \sphinxstyleemphasis{all}).
For finite\sphinxhyphen{}size particles, an internal temperature and heat flow can
be used to model heat conduction as in the
{\hyperref[\detokenize{Howto_granular::doc}]{\sphinxcrossref{\DUrole{doc}{GRANULAR package}}}}.

\sphinxAtStartPar
Another application is to use the \sphinxstyleemphasis{rmass} flag in order to have
per\sphinxhyphen{}atom masses instead of per\sphinxhyphen{}type masses.  This could be used to
study isotope effects with partial isotope substitution. {\hyperref[\detokenize{fix_property_atom:isotopes}]{\sphinxcrossref{\DUrole{std,std-ref}{See
below}}}} for an example of simulating a mixture of light and
heavy water with the TIP4P water potential.

\sphinxAtStartPar
An alternative to using fix \sphinxstyleemphasis{property/atom} for these examples is to
use an atom style that does define molecule IDs or charge or per\sphinxhyphen{}atom
mass (indirectly via diameter and density) or to use a hybrid atom
style that combines two or more atom styles to provide the union of
all their atom properties. However, this has two practical drawbacks:
first, it typically necessitates changing the format of the Atoms
section in the data file and second, it may define additional
properties that are not needed such as bond lists, which incurs some
overhead when there are no bonds.

\sphinxAtStartPar
In the future, we may add additional existing per\sphinxhyphen{}atom properties to
fix property/atom, similar to \sphinxstyleemphasis{mol}, \sphinxstyleemphasis{q}, \sphinxstyleemphasis{rmass}, \sphinxstyleemphasis{temperature},
or \sphinxstyleemphasis{heatflow} which “turn\sphinxhyphen{}on” specific properties defined by some atom
styles, so they can be easily used by atom styles that do not define
them.

\sphinxAtStartPar
More generally, the \sphinxstyleemphasis{i\_name} and \sphinxstyleemphasis{d\_name} options allow one or more
new custom per\sphinxhyphen{}atom vectors to be defined.  Likewise the \sphinxstyleemphasis{i2\_name} and
\sphinxstyleemphasis{d2\_name} options allow one or more custom per\sphinxhyphen{}atom arrays to be
defined.  The \sphinxstyleemphasis{i2\_name} and \sphinxstyleemphasis{d2\_name} options take an argument \sphinxstyleemphasis{N}
which specifies the number of columns in the per\sphinxhyphen{}atom array, i.e. the
number of attributes associated with each atom.  \sphinxstyleemphasis{N} \textgreater{}= 1 is required.

\sphinxAtStartPar
Each name must be unique and can use alphanumeric or underscore
characters.  These vectors and arrays can store whatever values you
decide are useful in your simulation.  As explained below there are
several ways to initialize, access, and output these values, via input
script commands, data files, and in new code you add to LAMMPS.

\sphinxAtStartPar
This is effectively a simple way to add per\sphinxhyphen{}atom properties to a model
without needing to write code for a new {\hyperref[\detokenize{atom_style::doc}]{\sphinxcrossref{\DUrole{doc}{atom style}}}}
that defines the properties.  Note however that implementing a new
atom style allows new atom properties to be more tightly and
seamlessly integrated with the rest of the code.

\sphinxAtStartPar
The new atom properties encode values that migrate with atoms to new
processors and are written to restart files.  If you want the new
properties to also be defined for ghost atoms, then use the \sphinxstyleemphasis{ghost}
keyword with a value of \sphinxstyleemphasis{yes}.  This will invoke extra communication
when ghost atoms are created (at every re\sphinxhyphen{}neighboring) to ensure the
new properties are also defined for the ghost atoms.

\begin{sphinxadmonition}{note}{Properties on ghost atoms}

\sphinxAtStartPar
If you use the \sphinxstyleemphasis{mol}, \sphinxstyleemphasis{q} or \sphinxstyleemphasis{rmass} names, you most likely want
to set \sphinxstyleemphasis{ghost} yes, since these properties are stored with ghost
atoms if you use an {\hyperref[\detokenize{atom_style::doc}]{\sphinxcrossref{\DUrole{doc}{atom\_style}}}} that defines
them.  Many LAMMPS operations that use molecule IDs or charge, such
as neighbor lists and pair styles, will expect ghost atoms to have
these values.  LAMMPS will issue a warning it you define those
vectors but do not set \sphinxstyleemphasis{ghost} yes.
\end{sphinxadmonition}

\begin{sphinxadmonition}{note}{Limitations on ghost atom properties}

\sphinxAtStartPar
The specified properties for ghost atoms are not updated every
timestep, but only once every few steps when neighbor lists are
re\sphinxhyphen{}built.  Thus the \sphinxstyleemphasis{ghost} keyword is suitable for static
properties, like molecule IDs, but not for dynamic properties that
change every step.  For the latter, the code you add to LAMMPS to
change the properties will also need to communicate their new
values to/from ghost atoms, an operation that can be invoked from
within a {\hyperref[\detokenize{pair_style::doc}]{\sphinxcrossref{\DUrole{doc}{pair style}}}} or {\hyperref[\detokenize{fix::doc}]{\sphinxcrossref{\DUrole{doc}{fix}}}} or
{\hyperref[\detokenize{compute::doc}]{\sphinxcrossref{\DUrole{doc}{compute}}}} that you write.
\end{sphinxadmonition}


\bigskip\hrule\bigskip


\sphinxAtStartPar
This fix is one of a small number that can be defined in an input
script before the simulation box is created or atoms are defined.
This is so it can be used with the {\hyperref[\detokenize{read_data::doc}]{\sphinxcrossref{\DUrole{doc}{read\_data}}}}
command as described next.

\sphinxAtStartPar
Per\sphinxhyphen{}atom properties that are defined by the {\hyperref[\detokenize{atom_style::doc}]{\sphinxcrossref{\DUrole{doc}{atom style}}}} are initialized when atoms are created, e.g. by the
{\hyperref[\detokenize{read_data::doc}]{\sphinxcrossref{\DUrole{doc}{read\_data}}}} or {\hyperref[\detokenize{create_atoms::doc}]{\sphinxcrossref{\DUrole{doc}{create\_atoms}}}}
commands.  The per\sphinxhyphen{}atom properties defined by this fix are not.  So
you need to initialize them explicitly.  One way to do this is
{\hyperref[\detokenize{read_data::doc}]{\sphinxcrossref{\DUrole{doc}{read\_data}}}} command, using its \sphinxstyleemphasis{fix} keyword and
passing it the fix\sphinxhyphen{}ID of this fix.

\sphinxAtStartPar
Thus these commands:

\begin{sphinxVerbatim}[commandchars=\\\{\}]
\PYG{k}{fix }\PYG{n+nv+nvIdentifier}{prop}\PYG{+w}{ }\PYG{n+nv+nvIdentifier}{all}\PYG{+w}{ }\PYG{n}{property}\PYG{o}{/}\PYG{n}{atom}\PYG{+w}{ }\PYG{n}{mol}\PYG{+w}{ }\PYG{n}{d\PYGZus{}flag}
\PYG{k}{read\PYGZus{}data}\PYG{+w}{ }\PYG{n}{data.txt}\PYG{+w}{ }\PYG{k}{fix }\PYG{n+nv+nvIdentifier}{prop}\PYG{+w}{ }\PYG{n+nv+nvIdentifier}{NULL}\PYG{+w}{ }\PYG{n}{Molecules}
\end{sphinxVerbatim}

\sphinxAtStartPar
would allow a data file to have a section like this:

\begin{sphinxVerbatim}[commandchars=\\\{\}]
Molecules

1 4 1.5
2 4 3.0
3 10 1.0
4 10 1.0
5 10 1.0
...
N 763 4.5
\end{sphinxVerbatim}

\sphinxAtStartPar
where N is the number of atoms, the first field on each line is the
atom\sphinxhyphen{}ID, the next two are a molecule\sphinxhyphen{}ID and a floating point value
that will be stored in a new property called “flag”.  If a per\sphinxhyphen{}atom
array was specified in the fix property/atom command then the \sphinxstyleemphasis{N}
values for that array must be specified consecutively for that
property on each line.  Note that the order of values on each line
corresponds to the order of custom names in the fix property/atom
command.

\sphinxAtStartPar
Note that the the lines of per\sphinxhyphen{}atom properties can be listed in any
order.  Also note that all the per\sphinxhyphen{}atom properties specified by the
fix ID (prop in this case) must be included on each line in the
specified data file section (Molecules in this case).

\sphinxAtStartPar
Another way of initializing the new properties is via the {\hyperref[\detokenize{set::doc}]{\sphinxcrossref{\DUrole{doc}{set}}}} command.  For example, if you wanted molecules defined for
every set of 10 atoms, based on their atom\sphinxhyphen{}IDs, these commands could
be used:

\begin{sphinxVerbatim}[commandchars=\\\{\}]
\PYG{k}{fix }\PYG{n+nv+nvIdentifier}{prop}\PYG{+w}{ }\PYG{n+nv+nvIdentifier}{all}\PYG{+w}{ }\PYG{n}{property}\PYG{o}{/}\PYG{n}{atom}\PYG{+w}{ }\PYG{n}{mol}
\PYG{k}{variable }\PYG{n+nv+nvIdentifier}{cluster}\PYG{+w}{ }\PYG{n}{atom}\PYG{+w}{ }\PYG{n+nv}{(}\PYG{n+nv}{(}\PYG{n+nv}{id\PYGZhy{}1}\PYG{n+nv}{)}\PYG{n+nv}{/10}\PYG{n+nv}{)}\PYG{o}{+}\PYG{l+m}{1}
\PYG{k}{set}\PYG{+w}{ }\PYG{n}{atom}\PYG{+w}{ }\PYG{o}{*}\PYG{+w}{ }\PYG{n}{mol}\PYG{+w}{ }\PYG{n}{v\PYGZus{}cluster}
\end{sphinxVerbatim}

\sphinxAtStartPar
The {\hyperref[\detokenize{variable::doc}]{\sphinxcrossref{\DUrole{doc}{atom\sphinxhyphen{}style variable}}}} will create values for atoms
with IDs 31,32,33,…40 that are 4.0,4.1,4.2,…,4.9.  When the
{\hyperref[\detokenize{set::doc}]{\sphinxcrossref{\DUrole{doc}{set}}}} commands assigns them to the molecule ID for each
atom, they will be truncated to an integer value, so atoms 31\sphinxhyphen{}40 will
all be assigned a molecule ID of 4.

\sphinxAtStartPar
Note that {\hyperref[\detokenize{variable::doc}]{\sphinxcrossref{\DUrole{doc}{atomfile\sphinxhyphen{}style variables}}}} can also be used
in place of atom\sphinxhyphen{}style variables, which means in this case that the
molecule IDs could be read\sphinxhyphen{}in from a separate file and assigned by the
{\hyperref[\detokenize{set::doc}]{\sphinxcrossref{\DUrole{doc}{set}}}} command.  This allows you to initialize new per\sphinxhyphen{}atom
properties in a completely general fashion.


\bigskip\hrule\bigskip


\sphinxAtStartPar
For new atom properties specified as \sphinxstyleemphasis{i\_name}, \sphinxstyleemphasis{d\_name}, \sphinxstyleemphasis{i2\_name}, or
\sphinxstyleemphasis{d2\_name}, the {\hyperref[\detokenize{dump::doc}]{\sphinxcrossref{\DUrole{doc}{dump custom}}}} and {\hyperref[\detokenize{compute_property_atom::doc}]{\sphinxcrossref{\DUrole{doc}{compute
property/atom}}}} commands can access their
values.  This means that the values can be used accessed by fixes like
{\hyperref[\detokenize{fix_ave_atom::doc}]{\sphinxcrossref{\DUrole{doc}{fix ave/atom}}}}, accessed by other computes like
{\hyperref[\detokenize{compute_reduce::doc}]{\sphinxcrossref{\DUrole{doc}{compute reduce}}}}, or used in {\hyperref[\detokenize{variable::doc}]{\sphinxcrossref{\DUrole{doc}{atom\sphinxhyphen{}style
variables}}}}.

\sphinxAtStartPar
For example, these commands will output both the instantaneous and
time\sphinxhyphen{}averaged values of two new properties to a custom dump file:

\begin{sphinxVerbatim}[commandchars=\\\{\}]
\PYG{k}{fix }\PYG{n+nv+nvIdentifier}{myprops}\PYG{+w}{ }\PYG{n+nv+nvIdentifier}{all}\PYG{+w}{ }\PYG{n}{property}\PYG{o}{/}\PYG{n}{atom}\PYG{+w}{ }\PYG{n}{i\PYGZus{}flag1}\PYG{+w}{ }\PYG{n}{d\PYGZus{}flag2}
\PYG{k}{compute }\PYG{n+nv+nvIdentifier}{1}\PYG{+w}{ }\PYG{n+nv+nvIdentifier}{all}\PYG{+w}{ }\PYG{n}{property}\PYG{o}{/}\PYG{n}{atom}\PYG{+w}{ }\PYG{n}{i\PYGZus{}flag1}\PYG{+w}{ }\PYG{n}{d\PYGZus{}flag2}
\PYG{k}{fix }\PYG{n+nv+nvIdentifier}{1}\PYG{+w}{ }\PYG{n+nv+nvIdentifier}{all}\PYG{+w}{ }\PYG{n}{ave}\PYG{o}{/}\PYG{n}{atom}\PYG{+w}{ }\PYG{l+m}{10}\PYG{+w}{ }\PYG{l+m}{10}\PYG{+w}{ }\PYG{l+m}{100}\PYG{+w}{ }\PYG{n}{c\PYGZus{}1[1]}\PYG{+w}{ }\PYG{n}{c\PYGZus{}1[2]}
\PYG{k}{dump }\PYG{n+nv+nvIdentifier}{1}\PYG{+w}{ }\PYG{n+nv+nvIdentifier}{all}\PYG{+w}{ }\PYG{n}{custom}\PYG{+w}{ }\PYG{l+m}{100}\PYG{+w}{ }\PYG{n}{tmp.dump}\PYG{+w}{ }\PYG{n}{id}\PYG{+w}{ }\PYG{n}{x}\PYG{+w}{ }\PYG{n}{y}\PYG{+w}{ }\PYG{n}{z}\PYG{+w}{ }\PYG{n}{i\PYGZus{}flag1}\PYG{+w}{ }\PYG{n}{d\PYGZus{}flag2}\PYG{+w}{ }\PYG{n}{f\PYGZus{}1[1]}\PYG{+w}{ }\PYG{n}{f\PYGZus{}1[2]}
\end{sphinxVerbatim}


\bigskip\hrule\bigskip


\sphinxAtStartPar
If you wish to add new {\hyperref[\detokenize{pair_style::doc}]{\sphinxcrossref{\DUrole{doc}{pair styles}}}}, {\hyperref[\detokenize{fix::doc}]{\sphinxcrossref{\DUrole{doc}{fixes}}}}, or {\hyperref[\detokenize{compute::doc}]{\sphinxcrossref{\DUrole{doc}{computes}}}} that use the per\sphinxhyphen{}atom properties
defined by this fix, see the {\hyperref[\detokenize{Modify_atom::doc}]{\sphinxcrossref{\DUrole{doc}{Modify atom}}}} doc page
which has details on how the custom properties of this fix can be
accessed from added classes.


\bigskip\hrule\bigskip

\phantomsection\label{\detokenize{fix_property_atom:isotopes}}
\sphinxAtStartPar
Here is an example of using per\sphinxhyphen{}atom masses with TIP4P water to study
isotope effects. When setting up simulations with the {\hyperref[\detokenize{Howto_tip4p::doc}]{\sphinxcrossref{\DUrole{doc}{TIP4P pair
styles}}}} for water, you have to provide exactly one atom
type each to identify the water oxygen and hydrogen atoms. Since the
atom mass is normally tied to the atom type, this makes it impossible
to study multiple isotopes in the same simulation.  With \sphinxstyleemphasis{fix
property/atom rmass} however, the per\sphinxhyphen{}type masses are replaced by
per\sphinxhyphen{}atom masses. Asumming you have a working input deck for regular
TIP4P water, where water oxygen is atom type 1 and water hydrogen is
atom type 2, the following lines of input script convert this to using
per\sphinxhyphen{}atom masses:

\begin{sphinxVerbatim}[commandchars=\\\{\}]
\PYG{k}{fix }\PYG{n+nv+nvIdentifier}{Isotopes}\PYG{+w}{ }\PYG{n+nv+nvIdentifier}{all}\PYG{+w}{ }\PYG{n}{property}\PYG{o}{/}\PYG{n}{atom}\PYG{+w}{ }\PYG{n}{rmass}\PYG{+w}{ }\PYG{n}{ghost}\PYG{+w}{ }\PYG{n}{yes}
\PYG{k}{set}\PYG{+w}{ }\PYG{n}{type}\PYG{+w}{ }\PYG{l+m}{1}\PYG{+w}{ }\PYG{n}{mass}\PYG{+w}{ }\PYG{l+m}{15.9994}
\PYG{k}{set}\PYG{+w}{ }\PYG{n}{type}\PYG{+w}{ }\PYG{l+m}{2}\PYG{+w}{ }\PYG{n}{mass}\PYG{+w}{ }\PYG{l+m}{1.008}
\end{sphinxVerbatim}

\sphinxAtStartPar
When writing out the system data with the {\hyperref[\detokenize{write_data::doc}]{\sphinxcrossref{\DUrole{doc}{write\_data}}}} command, there will be a new section named with the
fix\sphinxhyphen{}ID (i.e. \sphinxstyleemphasis{Isotopes} in this case). Alternatively, you can take an
existing data file and just add this \sphinxstyleemphasis{Isotopes} section with one line
per atom containing atom\sphinxhyphen{}ID and mass. Either way, the extended data
file can be read back with:

\begin{sphinxVerbatim}[commandchars=\\\{\}]
\PYG{k}{fix }\PYG{n+nv+nvIdentifier}{Isotopes}\PYG{+w}{ }\PYG{n+nv+nvIdentifier}{all}\PYG{+w}{ }\PYG{n}{property}\PYG{o}{/}\PYG{n}{atom}\PYG{+w}{ }\PYG{n}{rmass}\PYG{+w}{ }\PYG{n}{ghost}\PYG{+w}{ }\PYG{n}{yes}
\PYG{k}{read\PYGZus{}data}\PYG{+w}{ }\PYG{n}{tip4p}\PYG{o}{\PYGZhy{}}\PYG{n}{isotopes.data}\PYG{+w}{ }\PYG{k}{fix }\PYG{n+nv+nvIdentifier}{Isotopes}\PYG{+w}{ }\PYG{n+nv+nvIdentifier}{NULL}\PYG{+w}{ }\PYG{n}{Isotopes}
\end{sphinxVerbatim}

\sphinxAtStartPar
Please note that the first \sphinxstyleemphasis{Isotopes} refers to the fix\sphinxhyphen{}ID and the
second to the name of the section. The following input script code
will now change the first 100 water molecules in this example to heavy
water:

\begin{sphinxVerbatim}[commandchars=\\\{\}]
\PYG{k}{group }\PYG{n+nv+nvIdentifier}{hwat}\PYG{+w}{ }\PYG{n}{id}\PYG{+w}{ }\PYG{l+m}{2:300:3}
\PYG{k}{group }\PYG{n+nv+nvIdentifier}{hwat}\PYG{+w}{ }\PYG{n}{id}\PYG{+w}{ }\PYG{l+m}{3:300:3}
\PYG{k}{set}\PYG{+w}{ }\PYG{k}{group }\PYG{n+nv+nvIdentifier}{hwat}\PYG{+w}{ }\PYG{n}{mass}\PYG{+w}{ }\PYG{l+m}{2.0141018}
\end{sphinxVerbatim}


\bigskip\hrule\bigskip


\sphinxAtStartPar
Styles with a \sphinxstyleemphasis{gpu}, \sphinxstyleemphasis{intel}, \sphinxstyleemphasis{kk}, \sphinxstyleemphasis{omp}, or \sphinxstyleemphasis{opt} suffix are
functionally the same as the corresponding style without the suffix.
They have been optimized to run faster, depending on your available
hardware, as discussed on the {\hyperref[\detokenize{Speed_packages::doc}]{\sphinxcrossref{\DUrole{doc}{Accelerator packages}}}}
page.  The accelerated styles take the same arguments and should
produce the same results, except for round\sphinxhyphen{}off and precision issues.

\sphinxAtStartPar
These accelerated styles are part of the GPU, INTEL, KOKKOS,
OPENMP, and OPT packages, respectively.  They are only enabled if
LAMMPS was built with those packages.  See the {\hyperref[\detokenize{Build_package::doc}]{\sphinxcrossref{\DUrole{doc}{Build package}}}} page for more info.

\sphinxAtStartPar
You can specify the accelerated styles explicitly in your input script
by including their suffix, or you can use the {\hyperref[\detokenize{Run_options::doc}]{\sphinxcrossref{\DUrole{doc}{\sphinxhyphen{}suffix command\sphinxhyphen{}line switch}}}} when you invoke LAMMPS, or you can use the
{\hyperref[\detokenize{suffix::doc}]{\sphinxcrossref{\DUrole{doc}{suffix}}}} command in your input script.

\sphinxAtStartPar
See the {\hyperref[\detokenize{Speed_packages::doc}]{\sphinxcrossref{\DUrole{doc}{Accelerator packages}}}} page for more
instructions on how to use the accelerated styles effectively.


\bigskip\hrule\bigskip



\subsection{Restart, fix\_modify, output, run start/stop, minimize info}
\label{\detokenize{fix_property_atom:restart-fix-modify-output-run-start-stop-minimize-info}}
\sphinxAtStartPar
This fix writes the per\sphinxhyphen{}atom values it stores to {\hyperref[\detokenize{restart::doc}]{\sphinxcrossref{\DUrole{doc}{binary restart
files}}}}, so that the values can be restored when a simulation
is restarted.  See the {\hyperref[\detokenize{read_restart::doc}]{\sphinxcrossref{\DUrole{doc}{read\_restart}}}} command for
info on how to re\sphinxhyphen{}specify a fix in an input script that reads a
restart file, so that the operation of the fix continues in an
uninterrupted fashion.

\begin{sphinxadmonition}{warning}{Warning:}
\sphinxAtStartPar
When reading data from a restart file, this fix command has to be
specified \sphinxstylestrong{after} the \sphinxstyleemphasis{read\_restart} command and \sphinxstylestrong{exactly} the
same was in the input script that created the restart file.  LAMMPS
will only check whether a fix is of the same style and has the same
fix ID and in case of a match will then try to initialize the fix
with the data stored in the binary restart file.  If the names and
associated date types in the new fix property/atom command do not
match the old one exactly, data can be corrupted or LAMMPS may crash.
If the fix is specified \sphinxstylestrong{before} the \sphinxstyleemphasis{read\_restart} command its
data will not be restored.
\end{sphinxadmonition}

\sphinxAtStartPar
None of the {\hyperref[\detokenize{fix_modify::doc}]{\sphinxcrossref{\DUrole{doc}{fix\_modify}}}} options are relevant to
this fix.  No global or per\sphinxhyphen{}atom quantities are stored by this fix for
access by various {\hyperref[\detokenize{Howto_output::doc}]{\sphinxcrossref{\DUrole{doc}{output commands}}}}.  No parameter
of this fix can be used with the \sphinxstyleemphasis{start/stop} keywords of the
{\hyperref[\detokenize{run::doc}]{\sphinxcrossref{\DUrole{doc}{run}}}} command.  This fix is not invoked during {\hyperref[\detokenize{minimize::doc}]{\sphinxcrossref{\DUrole{doc}{energy
minimization}}}}.


\subsection{Restrictions}
\label{\detokenize{fix_property_atom:restrictions}}\begin{quote}

\sphinxAtStartPar
none
\end{quote}


\subsection{Related commands}
\label{\detokenize{fix_property_atom:related-commands}}
\sphinxAtStartPar
{\hyperref[\detokenize{read_data::doc}]{\sphinxcrossref{\DUrole{doc}{read\_data}}}}, {\hyperref[\detokenize{set::doc}]{\sphinxcrossref{\DUrole{doc}{set}}}},
{\hyperref[\detokenize{compute_property_atom::doc}]{\sphinxcrossref{\DUrole{doc}{compute property/atom}}}}


\subsection{Default}
\label{\detokenize{fix_property_atom:default}}
\sphinxAtStartPar
The default keyword value is ghost = no.

\sphinxstepscope

\index{fix python/invoke@\spxentry{fix python/invoke}}\ignorespaces 

\section{fix python/invoke command}
\label{\detokenize{fix_python_invoke:fix-python-invoke-command}}\label{\detokenize{fix_python_invoke:index-0}}\label{\detokenize{fix_python_invoke::doc}}

\subsection{Syntax}
\label{\detokenize{fix_python_invoke:syntax}}
\begin{sphinxVerbatim}[commandchars=\\\{\}]
\PYG{k}{fix }\PYG{n+nv+nvIdentifier}{ID}\PYG{+w}{ }\PYG{n+nv+nvIdentifier}{group\PYGZhy{}ID}\PYG{+w}{ }\PYG{n}{python}\PYG{o}{/}\PYG{n}{invoke}\PYG{+w}{ }\PYG{n}{N}\PYG{+w}{ }\PYG{n}{callback}\PYG{+w}{ }\PYG{n}{function\PYGZus{}name}
\end{sphinxVerbatim}
\begin{itemize}
\item {} 
\sphinxAtStartPar
ID, group\sphinxhyphen{}ID are ignored by this fix

\item {} 
\sphinxAtStartPar
python/invoke = style name of this fix command

\item {} 
\sphinxAtStartPar
N = execute every N steps

\item {} 
\sphinxAtStartPar
callback = \sphinxstyleemphasis{post\_force} or \sphinxstyleemphasis{end\_of\_step}
\begin{sphinxalltt}
\sphinxstyleemphasis{post\_force} = callback after force computations on atoms every N time steps
\sphinxstyleemphasis{end\_of\_step} = callback after every N time steps
\end{sphinxalltt}

\end{itemize}


\subsection{Examples}
\label{\detokenize{fix_python_invoke:examples}}
\begin{sphinxVerbatim}[commandchars=\\\{\}]
\PYG{k}{python}\PYG{+w}{ }\PYG{n}{post\PYGZus{}force\PYGZus{}callback}\PYG{+w}{ }\PYG{n}{here}\PYG{+w}{ }\PYG{l+s}{\PYGZdq{}}\PYG{l+s}{\PYGZdq{}}\PYG{l+s}{\PYGZdq{}}
\PYG{l+s}{from lammps import lammps}

\PYG{l+s}{def post\PYGZus{}force\PYGZus{}callback(lammps\PYGZus{}ptr, vflag):}
\PYG{l+s}{    lmp = lammps(ptr=lammps\PYGZus{}ptr)}
\PYG{l+s}{    \PYGZsh{} access LAMMPS state using Python interface}
\PYG{l+s}{\PYGZdq{}}\PYG{l+s}{\PYGZdq{}}\PYG{l+s}{\PYGZdq{}}

\PYG{k}{python}\PYG{+w}{ }\PYG{n}{end\PYGZus{}of\PYGZus{}step\PYGZus{}callback}\PYG{+w}{ }\PYG{n}{here}\PYG{+w}{ }\PYG{l+s}{\PYGZdq{}}\PYG{l+s}{\PYGZdq{}}\PYG{l+s}{\PYGZdq{}}
\PYG{l+s}{def end\PYGZus{}of\PYGZus{}step\PYGZus{}callback(lammps\PYGZus{}ptr):}
\PYG{l+s}{    lmp = lammps(ptr=lammps\PYGZus{}ptr)}
\PYG{l+s}{    \PYGZsh{} access LAMMPS state using Python interface}
\PYG{l+s}{\PYGZdq{}}\PYG{l+s}{\PYGZdq{}}\PYG{l+s}{\PYGZdq{}}

\PYG{k}{fix }\PYG{n+nv+nvIdentifier}{pf}\PYG{+w}{  }\PYG{n+nv+nvIdentifier}{all}\PYG{+w}{ }\PYG{n}{python}\PYG{o}{/}\PYG{n}{invoke}\PYG{+w}{ }\PYG{l+m}{50}\PYG{+w}{ }\PYG{n}{post\PYGZus{}force}\PYG{+w}{ }\PYG{n}{post\PYGZus{}force\PYGZus{}callback}
\PYG{k}{fix }\PYG{n+nv+nvIdentifier}{eos}\PYG{+w}{ }\PYG{n+nv+nvIdentifier}{all}\PYG{+w}{ }\PYG{n}{python}\PYG{o}{/}\PYG{n}{invoke}\PYG{+w}{ }\PYG{l+m}{50}\PYG{+w}{ }\PYG{n}{end\PYGZus{}of\PYGZus{}step}\PYG{+w}{ }\PYG{n}{end\PYGZus{}of\PYGZus{}step\PYGZus{}callback}
\end{sphinxVerbatim}


\subsection{Description}
\label{\detokenize{fix_python_invoke:description}}
\sphinxAtStartPar
This fix allows you to call a Python function during a simulation run.
The callback is either executed after forces have been applied to atoms
or at the end of every N time steps.

\sphinxAtStartPar
Callback functions must be declared in the global scope of the
active Python interpreter. This can either be done by defining it
inline using the python command or by importing functions from other
Python modules. If LAMMPS is driven using the library interface from
Python, functions defined in the driving Python interpreter can also
be executed.

\sphinxAtStartPar
Each callback is given a pointer object as first argument. This can be
used to initialize an instance of the lammps Python interface, which
gives access to the LAMMPS state from Python.

\begin{sphinxadmonition}{warning}{Warning:}
\sphinxAtStartPar
While you can access the state of LAMMPS via library functions
from these callbacks, trying to execute input script commands will in the best
case not work or in the worst case result in undefined behavior.
\end{sphinxadmonition}


\subsection{Restrictions}
\label{\detokenize{fix_python_invoke:restrictions}}
\sphinxAtStartPar
This fix is part of the PYTHON package.  It is only enabled if LAMMPS
was built with that package.  See the {\hyperref[\detokenize{Build_package::doc}]{\sphinxcrossref{\DUrole{doc}{Build package}}}} page for more info.

\sphinxAtStartPar
Building LAMMPS with the PYTHON package will link LAMMPS with the
Python library on your system.  Settings to enable this are in the
lib/python/Makefile.lammps file.  See the lib/python/README file for
information on those settings.


\subsection{Related commands}
\label{\detokenize{fix_python_invoke:related-commands}}
\sphinxAtStartPar
{\hyperref[\detokenize{python::doc}]{\sphinxcrossref{\DUrole{doc}{python command}}}}

\sphinxstepscope

\index{fix python/move@\spxentry{fix python/move}}\ignorespaces 

\section{fix python/move command}
\label{\detokenize{fix_python_move:fix-python-move-command}}\label{\detokenize{fix_python_move:index-0}}\label{\detokenize{fix_python_move::doc}}

\subsection{Syntax}
\label{\detokenize{fix_python_move:syntax}}
\begin{sphinxVerbatim}[commandchars=\\\{\}]
\PYG{k}{fix }\PYG{n+nv+nvIdentifier}{python}\PYG{o}{/}\PYG{n}{move}\PYG{+w}{ }\PYG{n}{pymodule.CLASS}
\end{sphinxVerbatim}

\sphinxAtStartPar
pymodule.CLASS = use class \sphinxstylestrong{CLASS} in module/file \sphinxstylestrong{pymodule} to compute how to move atoms


\subsection{Examples}
\label{\detokenize{fix_python_move:examples}}
\begin{sphinxVerbatim}[commandchars=\\\{\}]
\PYG{k}{fix  }\PYG{n+nv+nvIdentifier}{1}\PYG{+w}{ }\PYG{n+nv+nvIdentifier}{all}\PYG{+w}{ }\PYG{n}{python}\PYG{o}{/}\PYG{n}{move}\PYG{+w}{ }\PYG{n}{py\PYGZus{}nve.NVE}
\PYG{k}{fix  }\PYG{n+nv+nvIdentifier}{1}\PYG{+w}{ }\PYG{n+nv+nvIdentifier}{all}\PYG{+w}{ }\PYG{n}{python}\PYG{o}{/}\PYG{n}{move}\PYG{+w}{ }\PYG{n}{py\PYGZus{}nve.NVE\PYGZus{}OPT}
\end{sphinxVerbatim}


\subsection{Description}
\label{\detokenize{fix_python_move:description}}
\sphinxAtStartPar
The \sphinxstyleemphasis{python/move} fix style provides a way to define ways how particles
are moved during an MD run from python script code, that is loaded from
a file into LAMMPS and executed at the various steps where other fixes
can be executed. This python script must contain specific python class
definitions.

\sphinxAtStartPar
This allows to implement complex position updates and also modified
time integration methods. Due to python being an interpreted language,
however, the performance of this fix can be moderately to significantly
slower than the corresponding C++ code. For specific cases, this
performance penalty can be limited through effective use of NumPy.


\bigskip\hrule\bigskip


\sphinxAtStartPar
The python module file has to start with the following code:

\begin{sphinxVerbatim}[commandchars=\\\{\}]
\PYG{k+kn}{from}\PYG{+w}{ }\PYG{n+nn}{\PYGZus{}\PYGZus{}future\PYGZus{}\PYGZus{}}\PYG{+w}{ }\PYG{k+kn}{import} \PYG{n}{print\PYGZus{}function}
\PYG{k+kn}{import}\PYG{+w}{ }\PYG{n+nn}{lammps}
\PYG{k+kn}{import}\PYG{+w}{ }\PYG{n+nn}{ctypes}
\PYG{k+kn}{import}\PYG{+w}{ }\PYG{n+nn}{traceback}
\PYG{k+kn}{import}\PYG{+w}{ }\PYG{n+nn}{numpy}\PYG{+w}{ }\PYG{k}{as}\PYG{+w}{ }\PYG{n+nn}{np}
\PYG{c+c1}{\PYGZsh{}}
\PYG{k}{class}\PYG{+w}{ }\PYG{n+nc}{LAMMPSFix}\PYG{p}{(}\PYG{n+nb}{object}\PYG{p}{)}\PYG{p}{:}
    \PYG{k}{def}\PYG{+w}{ }\PYG{n+nf+fm}{\PYGZus{}\PYGZus{}init\PYGZus{}\PYGZus{}}\PYG{p}{(}\PYG{n+nb+bp}{self}\PYG{p}{,} \PYG{n}{ptr}\PYG{p}{,} \PYG{n}{group\PYGZus{}name}\PYG{o}{=}\PYG{l+s+s2}{\PYGZdq{}}\PYG{l+s+s2}{all}\PYG{l+s+s2}{\PYGZdq{}}\PYG{p}{)}\PYG{p}{:}
        \PYG{n+nb+bp}{self}\PYG{o}{.}\PYG{n}{lmp} \PYG{o}{=} \PYG{n}{lammps}\PYG{o}{.}\PYG{n}{lammps}\PYG{p}{(}\PYG{n}{ptr}\PYG{o}{=}\PYG{n}{ptr}\PYG{p}{)}
        \PYG{n+nb+bp}{self}\PYG{o}{.}\PYG{n}{group\PYGZus{}name} \PYG{o}{=} \PYG{n}{group\PYGZus{}name}
\PYG{c+c1}{\PYGZsh{}}
\PYG{k}{class}\PYG{+w}{ }\PYG{n+nc}{LAMMPSFixMove}\PYG{p}{(}\PYG{n}{LAMMPSFix}\PYG{p}{)}\PYG{p}{:}
    \PYG{k}{def}\PYG{+w}{ }\PYG{n+nf+fm}{\PYGZus{}\PYGZus{}init\PYGZus{}\PYGZus{}}\PYG{p}{(}\PYG{n+nb+bp}{self}\PYG{p}{,} \PYG{n}{ptr}\PYG{p}{,} \PYG{n}{group\PYGZus{}name}\PYG{o}{=}\PYG{l+s+s2}{\PYGZdq{}}\PYG{l+s+s2}{all}\PYG{l+s+s2}{\PYGZdq{}}\PYG{p}{)}\PYG{p}{:}
        \PYG{n+nb}{super}\PYG{p}{(}\PYG{n}{LAMMPSFixMove}\PYG{p}{,} \PYG{n+nb+bp}{self}\PYG{p}{)}\PYG{o}{.}\PYG{n+nf+fm}{\PYGZus{}\PYGZus{}init\PYGZus{}\PYGZus{}}\PYG{p}{(}\PYG{n}{ptr}\PYG{p}{,} \PYG{n}{group\PYGZus{}name}\PYG{p}{)}
\PYG{c+c1}{\PYGZsh{}}
    \PYG{k}{def}\PYG{+w}{ }\PYG{n+nf}{init}\PYG{p}{(}\PYG{n+nb+bp}{self}\PYG{p}{)}\PYG{p}{:}
        \PYG{k}{pass}
\PYG{c+c1}{\PYGZsh{}}
    \PYG{k}{def}\PYG{+w}{ }\PYG{n+nf}{initial\PYGZus{}integrate}\PYG{p}{(}\PYG{n+nb+bp}{self}\PYG{p}{,} \PYG{n}{vflag}\PYG{p}{)}\PYG{p}{:}
        \PYG{k}{pass}
\PYG{c+c1}{\PYGZsh{}}
    \PYG{k}{def}\PYG{+w}{ }\PYG{n+nf}{final\PYGZus{}integrate}\PYG{p}{(}\PYG{n+nb+bp}{self}\PYG{p}{)}\PYG{p}{:}
        \PYG{k}{pass}
\PYG{c+c1}{\PYGZsh{}}
    \PYG{k}{def}\PYG{+w}{ }\PYG{n+nf}{initial\PYGZus{}integrate\PYGZus{}respa}\PYG{p}{(}\PYG{n+nb+bp}{self}\PYG{p}{,} \PYG{n}{vflag}\PYG{p}{,} \PYG{n}{ilevel}\PYG{p}{,} \PYG{n}{iloop}\PYG{p}{)}\PYG{p}{:}
        \PYG{k}{pass}
\PYG{c+c1}{\PYGZsh{}}
    \PYG{k}{def}\PYG{+w}{ }\PYG{n+nf}{final\PYGZus{}integrate\PYGZus{}respa}\PYG{p}{(}\PYG{n+nb+bp}{self}\PYG{p}{,} \PYG{n}{ilevel}\PYG{p}{,} \PYG{n}{iloop}\PYG{p}{)}\PYG{p}{:}
        \PYG{k}{pass}
\PYG{c+c1}{\PYGZsh{}}
    \PYG{k}{def}\PYG{+w}{ }\PYG{n+nf}{reset\PYGZus{}dt}\PYG{p}{(}\PYG{n+nb+bp}{self}\PYG{p}{)}\PYG{p}{:}
        \PYG{k}{pass}
\end{sphinxVerbatim}

\sphinxAtStartPar
Any classes implementing new atom motion functionality have to be
derived from the \sphinxstylestrong{LAMMPSFixMove} class, overriding the available
methods as needed.

\sphinxAtStartPar
Examples for how to do this are in the \sphinxstyleemphasis{examples/python} folder.


\bigskip\hrule\bigskip



\subsection{Restart, fix\_modify, output, run start/stop, minimize info}
\label{\detokenize{fix_python_move:restart-fix-modify-output-run-start-stop-minimize-info}}
\sphinxAtStartPar
No information about this fix is written to {\hyperref[\detokenize{restart::doc}]{\sphinxcrossref{\DUrole{doc}{binary restart files}}}}.  None of the {\hyperref[\detokenize{fix_modify::doc}]{\sphinxcrossref{\DUrole{doc}{fix\_modify}}}} options
are relevant to this fix.  No global or per\sphinxhyphen{}atom quantities are stored
by this fix for access by various {\hyperref[\detokenize{Howto_output::doc}]{\sphinxcrossref{\DUrole{doc}{output commands}}}}.
No parameter of this fix can be used with the \sphinxstyleemphasis{start/stop} keywords of
the {\hyperref[\detokenize{run::doc}]{\sphinxcrossref{\DUrole{doc}{run}}}} command.  This fix is not invoked during {\hyperref[\detokenize{minimize::doc}]{\sphinxcrossref{\DUrole{doc}{energy minimization}}}}.


\subsection{Restrictions}
\label{\detokenize{fix_python_move:restrictions}}
\sphinxAtStartPar
This pair style is part of the PYTHON package.  It is only enabled if
LAMMPS was built with that package.  See the {\hyperref[\detokenize{Build_package::doc}]{\sphinxcrossref{\DUrole{doc}{Build package}}}} page for more info.


\subsection{Related commands}
\label{\detokenize{fix_python_move:related-commands}}
\sphinxAtStartPar
{\hyperref[\detokenize{fix_nve::doc}]{\sphinxcrossref{\DUrole{doc}{fix nve}}}}, {\hyperref[\detokenize{fix_python_invoke::doc}]{\sphinxcrossref{\DUrole{doc}{fix python/invoke}}}}


\subsection{Default}
\label{\detokenize{fix_python_move:default}}
\sphinxAtStartPar
none

\sphinxstepscope

\index{fix qbmsst@\spxentry{fix qbmsst}}\ignorespaces 

\section{fix qbmsst command}
\label{\detokenize{fix_qbmsst:fix-qbmsst-command}}\label{\detokenize{fix_qbmsst:index-0}}\label{\detokenize{fix_qbmsst::doc}}

\subsection{Syntax}
\label{\detokenize{fix_qbmsst:syntax}}
\begin{sphinxVerbatim}[commandchars=\\\{\}]
\PYG{k}{fix }\PYG{n+nv+nvIdentifier}{ID}\PYG{+w}{ }\PYG{n+nv+nvIdentifier}{group\PYGZhy{}ID}\PYG{+w}{ }\PYG{n}{qbmsst}\PYG{+w}{ }\PYG{n}{dir}\PYG{+w}{ }\PYG{n}{shockvel}\PYG{+w}{ }\PYG{n}{keyword}\PYG{+w}{ }\PYG{n}{value}\PYG{+w}{ }\PYG{n}{...}
\end{sphinxVerbatim}
\begin{itemize}
\item {} 
\sphinxAtStartPar
ID, group\sphinxhyphen{}ID are documented in {\hyperref[\detokenize{fix::doc}]{\sphinxcrossref{\DUrole{doc}{fix}}}} command

\item {} 
\sphinxAtStartPar
qbmsst = style name of this fix

\item {} 
\sphinxAtStartPar
dir = \sphinxstyleemphasis{x} or \sphinxstyleemphasis{y} or \sphinxstyleemphasis{z}

\item {} 
\sphinxAtStartPar
shockvel = shock velocity (strictly positive, velocity units)

\item {} 
\sphinxAtStartPar
zero or more keyword/value pairs may be appended

\item {} 
\sphinxAtStartPar
keyword = \sphinxstyleemphasis{q} or \sphinxstyleemphasis{mu} or \sphinxstyleemphasis{p0} or \sphinxstyleemphasis{v0} or \sphinxstyleemphasis{e0} or \sphinxstyleemphasis{tscale} or \sphinxstyleemphasis{damp} or \sphinxstyleemphasis{seed} or \sphinxstyleemphasis{f\_max} or \sphinxstyleemphasis{N\_f} or \sphinxstyleemphasis{eta} or \sphinxstyleemphasis{beta} or \sphinxstyleemphasis{T\_init}
\begin{sphinxalltt}
\sphinxstyleemphasis{q} value = cell mass\sphinxhyphen{}like parameter (mass\textasciicircum{}2/distance\textasciicircum{}4 units)
\sphinxstyleemphasis{mu} value = artificial viscosity (mass/distance/time units)
\sphinxstyleemphasis{p0} value = initial pressure in the shock equations (pressure units)
\sphinxstyleemphasis{v0} value = initial simulation cell volume in the shock equations (distance\textasciicircum{}3 units)
\sphinxstyleemphasis{e0} value = initial total energy (energy units)
\sphinxstyleemphasis{tscale} value = reduction in initial temperature (unitless fraction between 0.0 and 1.0)
\sphinxstyleemphasis{damp} value = damping parameter (time units) inverse of friction \sphinxstyleemphasis{gamma}
\sphinxstyleemphasis{seed} value = random number seed (positive integer)
\sphinxstyleemphasis{f\_max} value = upper cutoff frequency of the vibration spectrum (1/time units)
\sphinxstyleemphasis{N\_f} value = number of frequency bins (positive integer)
\sphinxstyleemphasis{eta} value = coupling constant between the shock system and the quantum thermal bath (positive unitless)
\sphinxstyleemphasis{beta} value = the quantum temperature is updated every beta time steps (positive integer)
\sphinxstyleemphasis{T\_init} value = quantum temperature for the initial state (temperature units)
\end{sphinxalltt}

\end{itemize}


\subsection{Examples}
\label{\detokenize{fix_qbmsst:examples}}
\begin{sphinxVerbatim}[commandchars=\\\{\}]
\PYG{c}{\PYGZsh{} (liquid methane modeled with the REAX force field, real units)}
\PYG{k}{fix }\PYG{n+nv+nvIdentifier}{1}\PYG{+w}{ }\PYG{n+nv+nvIdentifier}{all}\PYG{+w}{ }\PYG{n}{qbmsst}\PYG{+w}{ }\PYG{n}{z}\PYG{+w}{ }\PYG{l+m}{0.122}\PYG{+w}{ }\PYG{n}{q}\PYG{+w}{ }\PYG{l+m}{25}\PYG{+w}{ }\PYG{n}{mu}\PYG{+w}{ }\PYG{l+m}{0.9}\PYG{+w}{ }\PYG{n}{tscale}\PYG{+w}{ }\PYG{l+m}{0.01}\PYG{+w}{ }\PYG{n}{damp}\PYG{+w}{ }\PYG{l+m}{200}\PYG{+w}{ }\PYG{n}{seed}\PYG{+w}{ }\PYG{l+m}{35082}\PYG{+w}{ }\PYG{n}{f\PYGZus{}max}\PYG{+w}{ }\PYG{l+m}{0.3}\PYG{+w}{ }\PYG{n}{N\PYGZus{}f}\PYG{+w}{ }\PYG{l+m}{100}\PYG{+w}{ }\PYG{n}{eta}\PYG{+w}{ }\PYG{l+m}{1}\PYG{+w}{ }\PYG{n}{beta}\PYG{+w}{ }\PYG{l+m}{400}\PYG{+w}{ }\PYG{n}{T\PYGZus{}init}\PYG{+w}{ }\PYG{l+m}{110}
\PYG{c}{\PYGZsh{} (quartz modeled with the BKS force field, metal units)}
\PYG{k}{fix }\PYG{n+nv+nvIdentifier}{2}\PYG{+w}{ }\PYG{n+nv+nvIdentifier}{all}\PYG{+w}{ }\PYG{n}{qbmsst}\PYG{+w}{ }\PYG{n}{z}\PYG{+w}{ }\PYG{l+m}{72}\PYG{+w}{ }\PYG{n}{q}\PYG{+w}{ }\PYG{l+m}{40}\PYG{+w}{ }\PYG{n}{tscale}\PYG{+w}{ }\PYG{l+m}{0.05}\PYG{+w}{ }\PYG{n}{damp}\PYG{+w}{ }\PYG{l+m}{1}\PYG{+w}{ }\PYG{n}{seed}\PYG{+w}{ }\PYG{l+m}{47508}\PYG{+w}{ }\PYG{n}{f\PYGZus{}max}\PYG{+w}{ }\PYG{l+m}{120.0}\PYG{+w}{ }\PYG{n}{N\PYGZus{}f}\PYG{+w}{ }\PYG{l+m}{100}\PYG{+w}{ }\PYG{n}{eta}\PYG{+w}{ }\PYG{l+m}{1.0}\PYG{+w}{ }\PYG{n}{beta}\PYG{+w}{ }\PYG{l+m}{500}\PYG{+w}{ }\PYG{n}{T\PYGZus{}init}\PYG{+w}{ }\PYG{l+m}{300}
\end{sphinxVerbatim}

\sphinxAtStartPar
Two example input scripts are given, including shocked
\(\alpha\textrm{-quartz}\) and shocked liquid methane.  The input
script first equilibrates an initial state with the quantum thermal
bath at the target temperature and then applies \sphinxstyleemphasis{fix qbmsst} to simulate
shock compression with quantum nuclear correction.  The following two
figures plot relevant quantities for shocked
\(\alpha\textrm{-quartz}\).

\noindent{\hspace*{\fill}\sphinxincludegraphics{{qbmsst_init}.jpg}\hspace*{\fill}}

\sphinxAtStartPar
Figure 1. Classical temperature
\(T_{cl} = \sum \frac{m_iv_i^2}{3Nk_B}\) vs. time for coupling the
\(\alpha\textrm{-quartz}\) initial state with the quantum thermal
bath at target quantum temperature \(T^{qm} = 300 K\). The NpH
ensemble is used for time integration while QTB provides the colored
random force. \(T^{cl}\) converges at the timescale of \sphinxstyleemphasis{damp}
which is set to be 1 ps.

\noindent{\hspace*{\fill}\sphinxincludegraphics{{qbmsst_shock}.jpg}\hspace*{\fill}}

\sphinxAtStartPar
Figure 2. Quantum temperature and pressure vs. time for simulating
shocked \(\alpha\textrm{-quartz}\) with \sphinxstyleemphasis{fix qbmsst}. The shock
propagates along the z direction. Restart of the \sphinxstyleemphasis{fix qbmsst} command
is demonstrated in the example input script. Thermodynamic quantities
stay continuous before and after the restart.


\subsection{Description}
\label{\detokenize{fix_qbmsst:description}}
\sphinxAtStartPar
This command performs the Quantum\sphinxhyphen{}Bath coupled Multi\sphinxhyphen{}Scale Shock
Technique (QBMSST) integration. See {\hyperref[\detokenize{fix_qbmsst:qi}]{\sphinxcrossref{\DUrole{std,std-ref}{(Qi)}}}} for a detailed
description of this method.  QBMSST provides description of the
thermodynamics and kinetics of shock processes while incorporating
quantum nuclear effects.  The \sphinxstyleemphasis{shockvel} setting determines the steady
shock velocity that will be simulated along direction \sphinxstyleemphasis{dir}.

\sphinxAtStartPar
Quantum nuclear effects {\hyperref[\detokenize{fix_qtb::doc}]{\sphinxcrossref{\DUrole{doc}{(fix qtb)}}}} can be crucial
especially when the temperature of the initial state is below the
classical limit or there is a great change in the zero point energies
between the initial and final states. Theoretical post processing
quantum corrections of shock compressed water and methane have been
reported as much as 30\% of the temperatures {\hyperref[\detokenize{fix_qbmsst:goldman1}]{\sphinxcrossref{\DUrole{std,std-ref}{(Goldman)}}}}.  A
self\sphinxhyphen{}consistent method that couples the shock to a quantum thermal
bath described by a colored noise Langevin thermostat has been
developed by Qi et al {\hyperref[\detokenize{fix_qbmsst:qi}]{\sphinxcrossref{\DUrole{std,std-ref}{(Qi)}}}} and applied to shocked methane.  The
onset of chemistry is reported to be at a pressure on the shock
Hugoniot that is 40\% lower than observed with classical molecular
dynamics.

\sphinxAtStartPar
It is highly recommended that the system be already in an equilibrium
state with a quantum thermal bath at temperature of \sphinxstyleemphasis{T\_init}.  The fix
command {\hyperref[\detokenize{fix_qtb::doc}]{\sphinxcrossref{\DUrole{doc}{fix qtb}}}} at constant temperature \sphinxstyleemphasis{T\_init} could
be used before applying this command to introduce self\sphinxhyphen{}consistent
quantum nuclear effects into the initial state.

\sphinxAtStartPar
The parameters \sphinxstyleemphasis{q}, \sphinxstyleemphasis{mu}, \sphinxstyleemphasis{e0}, \sphinxstyleemphasis{p0}, \sphinxstyleemphasis{v0} and \sphinxstyleemphasis{tscale} are described
in the command {\hyperref[\detokenize{fix_msst::doc}]{\sphinxcrossref{\DUrole{doc}{fix msst}}}}. The values of \sphinxstyleemphasis{e0}, \sphinxstyleemphasis{p0}, or
\sphinxstyleemphasis{v0} will be calculated on the first step if not specified.  The
parameter of \sphinxstyleemphasis{damp}, \sphinxstyleemphasis{f\_max}, and \sphinxstyleemphasis{N\_f} are described in the command
{\hyperref[\detokenize{fix_qtb::doc}]{\sphinxcrossref{\DUrole{doc}{fix qtb}}}}.

\sphinxAtStartPar
The \sphinxstyleemphasis{fix qbmsst} command couples the shock system to a quantum thermal
bath with a rate that is proportional to the change of the total
energy of the shock system, \(E^{tot} - E^{tot}_0\).
Here \(E^{etot}\) consists of both the system energy and a thermal
term, see {\hyperref[\detokenize{fix_qbmsst:qi}]{\sphinxcrossref{\DUrole{std,std-ref}{(Qi)}}}}, and \(E^{tot}_0 = e0\) is the
initial total energy.

\sphinxAtStartPar
The \sphinxstyleemphasis{eta} (\(\eta\)) parameter is a unitless coupling constant
between the shock system and the quantum thermal bath. A small \(\eta\)
value cannot adjust the quantum temperature fast enough during the
temperature ramping period of shock compression while large \(\eta\)
leads to big temperature oscillation. A value of \(\eta\) between 0.3 and
1 is usually appropriate for simulating most systems under shock
compression. We observe that different values of \(\eta\) lead to almost
the same final thermodynamic state behind the shock, as expected.

\sphinxAtStartPar
The quantum temperature is updated every \sphinxstyleemphasis{beta} (\(\beta\)) steps
with an integration time interval \(\beta\) times longer than the
simulation time step. In that case, \(E^{tot}\) is taken as its
average over the past \(\beta\) steps. The temperature of the quantum
thermal bath \(T^{qm}\) changes dynamically according to
the following equation where \(\Delta_t\) is the MD time step and
\(\gamma\) is the friction constant which is equal to the inverse
of the \sphinxstyleemphasis{damp} parameter.
\begin{equation*}
\begin{split}\frac{dT^{qm}}{dt} = \gamma\eta\sum^\beta_{l=1}\frac{E^{tot}(t-l\Delta t) - E^{tot}_0}{3\beta N k_B}\end{split}
\end{equation*}
\sphinxAtStartPar
The parameter \sphinxstyleemphasis{T\_init} is the initial temperature of the quantum
thermal bath and the system before shock loading.

\sphinxAtStartPar
For all pressure styles, the simulation box stays orthorhombic in
shape. Parrinello\sphinxhyphen{}Rahman boundary conditions (tilted box) are
supported by LAMMPS, but are not implemented for QBMSST.


\bigskip\hrule\bigskip



\subsection{Restart, fix\_modify, output, run start/stop, minimize info}
\label{\detokenize{fix_qbmsst:restart-fix-modify-output-run-start-stop-minimize-info}}
\sphinxAtStartPar
Because the state of the random number generator is not written to
{\hyperref[\detokenize{restart::doc}]{\sphinxcrossref{\DUrole{doc}{binary restart files}}}}, this fix cannot be restarted
“exactly” in an uninterrupted fashion. However, in a statistical
sense, a restarted simulation should produce similar behaviors of the
system as if it is not interrupted.  To achieve such a restart, one
should write explicitly the same value for \sphinxstyleemphasis{q}, \sphinxstyleemphasis{mu}, \sphinxstyleemphasis{damp},
\sphinxstyleemphasis{f\_max}, \sphinxstyleemphasis{N\_f}, \sphinxstyleemphasis{eta}, and \sphinxstyleemphasis{beta} and set \sphinxstyleemphasis{tscale} = 0 if the system
is compressed during the first run.

\sphinxAtStartPar
The cumulative energy change in the system imposed by this fix is
included in the {\hyperref[\detokenize{thermo_style::doc}]{\sphinxcrossref{\DUrole{doc}{thermodynamic output}}}} keywords
\sphinxstyleemphasis{ecouple} and \sphinxstyleemphasis{econserve}.  See the {\hyperref[\detokenize{thermo_style::doc}]{\sphinxcrossref{\DUrole{doc}{thermo\_style}}}}
doc page for details.

\sphinxAtStartPar
This fix computes a global scalar which can be accessed by various
{\hyperref[\detokenize{Howto_output::doc}]{\sphinxcrossref{\DUrole{doc}{output commands}}}}.  The scalar is the same
cumulative energy change due to this fix described in the previous
paragraph.  The scalar value calculated by this fix is “extensive”.

\sphinxAtStartPar
The progress of the QBMSST can be monitored by printing the global
scalar and global vector quantities computed by the fix.

\sphinxAtStartPar
As mentioned above, the scalar is the cumulative energy change due to
the fix.  By monitoring the thermodynamic \sphinxstyleemphasis{econserve} output, this can
be used to test if the MD timestep is sufficiently small for accurate
integration of the dynamic equations.

\sphinxAtStartPar
The global vector contains five values in the following order.  The
vector values output by this fix are “intensive”.

\sphinxAtStartPar
{[}\sphinxstyleemphasis{dhugoniot}, \sphinxstyleemphasis{drayleigh}, \sphinxstyleemphasis{lagrangian\_speed}, \sphinxstyleemphasis{lagrangian\_position},
\sphinxstyleemphasis{quantum\_temperature}{]}
\begin{enumerate}
\sphinxsetlistlabels{\arabic}{enumi}{enumii}{}{.}%
\item {} 
\sphinxAtStartPar
\sphinxstyleemphasis{dhugoniot} is the departure from the Hugoniot (temperature units).

\item {} 
\sphinxAtStartPar
\sphinxstyleemphasis{drayleigh} is the departure from the Rayleigh line (pressure units).

\item {} 
\sphinxAtStartPar
\sphinxstyleemphasis{lagrangian\_speed} is the laboratory\sphinxhyphen{}frame Lagrangian speed (particle velocity) of the computational cell (velocity units).

\item {} 
\sphinxAtStartPar
\sphinxstyleemphasis{lagrangian\_position} is the computational cell position in the reference frame moving at the shock speed. This is the distance of the computational cell behind the shock front.

\item {} 
\sphinxAtStartPar
\sphinxstyleemphasis{quantum\_temperature} is the temperature of the quantum thermal bath \(T^{qm}\).

\end{enumerate}

\sphinxAtStartPar
To print these quantities to the log file with descriptive column
headers, the following LAMMPS commands are suggested.

\begin{sphinxVerbatim}[commandchars=\\\{\}]
fix             fix\PYGZus{}id all msst z
variable        dhug    equal f\PYGZus{}fix\PYGZus{}id[1]
variable        dray    equal f\PYGZus{}fix\PYGZus{}id[2]
variable        lgr\PYGZus{}vel equal f\PYGZus{}fix\PYGZus{}id[3]
variable        lgr\PYGZus{}pos equal f\PYGZus{}fix\PYGZus{}id[4]
variable        T\PYGZus{}qm    equal f\PYGZus{}fix\PYGZus{}id[5]
thermo\PYGZus{}style    custom  step temp ke pe lz pzz econserve v\PYGZus{}dhug v\PYGZus{}dray v\PYGZus{}lgr\PYGZus{}vel v\PYGZus{}lgr\PYGZus{}pos v\PYGZus{}T\PYGZus{}qm f\PYGZus{}fix\PYGZus{}id
\end{sphinxVerbatim}

\sphinxAtStartPar
It is worth noting that the temp keyword for the {\hyperref[\detokenize{thermo_style::doc}]{\sphinxcrossref{\DUrole{doc}{thermo\_style}}}} command prints the instantaneous classical temperature
\(T^{cl}\) as described by the {\hyperref[\detokenize{fix_qtb::doc}]{\sphinxcrossref{\DUrole{doc}{fix qtb}}}} command.


\bigskip\hrule\bigskip



\subsection{Restrictions}
\label{\detokenize{fix_qbmsst:restrictions}}
\sphinxAtStartPar
This fix style is part of the QTB package.  It is only enabled if
LAMMPS was built with that package. See the {\hyperref[\detokenize{Build_package::doc}]{\sphinxcrossref{\DUrole{doc}{Build package}}}} page for more info.

\sphinxAtStartPar
All cell dimensions must be periodic. This fix can not be used with a
triclinic cell.  The QBMSST fix has been tested only for the group\sphinxhyphen{}ID
all.


\bigskip\hrule\bigskip



\subsection{Related commands}
\label{\detokenize{fix_qbmsst:related-commands}}
\sphinxAtStartPar
{\hyperref[\detokenize{fix_qtb::doc}]{\sphinxcrossref{\DUrole{doc}{fix qtb}}}}, {\hyperref[\detokenize{fix_msst::doc}]{\sphinxcrossref{\DUrole{doc}{fix msst}}}}


\bigskip\hrule\bigskip



\subsection{Default}
\label{\detokenize{fix_qbmsst:default}}
\sphinxAtStartPar
The keyword defaults are q = 10, mu = 0, tscale = 0.01, damp = 1, seed
= 880302, f\_max = 200.0, N\_f = 100, eta = 1.0, beta = 100, and
T\_init=300.0. e0, p0, and v0 are calculated on the first step.


\bigskip\hrule\bigskip

\phantomsection\label{\detokenize{fix_qbmsst:goldman1}}
\sphinxAtStartPar
\sphinxstylestrong{(Goldman)} Goldman, Reed and Fried, J. Chem. Phys. 131, 204103 (2009)

\phantomsection\label{\detokenize{fix_qbmsst:qi}}
\sphinxAtStartPar
\sphinxstylestrong{(Qi)} Qi and Reed, J. Phys. Chem. A 116, 10451 (2012).

\sphinxstepscope

\index{fix qeq/point@\spxentry{fix qeq/point}}\ignorespaces 
\index{fix qeq/shielded@\spxentry{fix qeq/shielded}}\ignorespaces 
\index{fix qeq/slater@\spxentry{fix qeq/slater}}\ignorespaces 
\index{fix qeq/dynamic@\spxentry{fix qeq/dynamic}}\ignorespaces 
\index{fix qeq/fire@\spxentry{fix qeq/fire}}\ignorespaces 

\section{fix qeq/point command}
\label{\detokenize{fix_qeq:fix-qeq-point-command}}\label{\detokenize{fix_qeq:index-4}}\label{\detokenize{fix_qeq:index-3}}\label{\detokenize{fix_qeq:index-2}}\label{\detokenize{fix_qeq:index-1}}\label{\detokenize{fix_qeq:index-0}}\label{\detokenize{fix_qeq::doc}}

\section{fix qeq/shielded command}
\label{\detokenize{fix_qeq:fix-qeq-shielded-command}}

\section{fix qeq/slater command}
\label{\detokenize{fix_qeq:fix-qeq-slater-command}}

\section{fix qeq/dynamic command}
\label{\detokenize{fix_qeq:fix-qeq-dynamic-command}}

\section{fix qeq/fire command}
\label{\detokenize{fix_qeq:fix-qeq-fire-command}}

\subsection{Syntax}
\label{\detokenize{fix_qeq:syntax}}
\begin{sphinxVerbatim}[commandchars=\\\{\}]
\PYG{k}{fix }\PYG{n+nv+nvIdentifier}{ID}\PYG{+w}{ }\PYG{n+nv+nvIdentifier}{group\PYGZhy{}ID}\PYG{+w}{ }\PYG{n}{style}\PYG{+w}{ }\PYG{n}{Nevery}\PYG{+w}{ }\PYG{n}{cutoff}\PYG{+w}{ }\PYG{n}{tolerance}\PYG{+w}{ }\PYG{n}{maxiter}\PYG{+w}{ }\PYG{n}{qfile}\PYG{+w}{ }\PYG{n}{keyword}\PYG{+w}{ }\PYG{n}{...}
\end{sphinxVerbatim}
\begin{itemize}
\item {} 
\sphinxAtStartPar
ID, group\sphinxhyphen{}ID are documented in {\hyperref[\detokenize{fix::doc}]{\sphinxcrossref{\DUrole{doc}{fix}}}} command

\item {} 
\sphinxAtStartPar
style = \sphinxstyleemphasis{qeq/point} or \sphinxstyleemphasis{qeq/shielded} or \sphinxstyleemphasis{qeq/slater} or \sphinxstyleemphasis{qeq/dynamic} or \sphinxstyleemphasis{qeq/fire}

\item {} 
\sphinxAtStartPar
Nevery = perform charge equilibration every this many steps

\item {} 
\sphinxAtStartPar
cutoff = global cutoff for charge\sphinxhyphen{}charge interactions (distance unit)

\item {} 
\sphinxAtStartPar
tolerance = precision to which charges will be equilibrated

\item {} 
\sphinxAtStartPar
maxiter = maximum iterations to perform charge equilibration

\item {} 
\sphinxAtStartPar
qfile = a filename with QEq parameters or \sphinxstyleemphasis{coul/streitz} or \sphinxstyleemphasis{reaxff}

\item {} 
\sphinxAtStartPar
zero or more keyword/value pairs may be appended

\item {} 
\sphinxAtStartPar
keyword = \sphinxstyleemphasis{alpha} or \sphinxstyleemphasis{qdamp} or \sphinxstyleemphasis{qstep} or \sphinxstyleemphasis{warn}
\begin{sphinxalltt}
\sphinxstyleemphasis{alpha} value = Slater type orbital exponent (qeq/slater only)
\sphinxstyleemphasis{qdamp} value = damping factor for damped dynamics charge solver (qeq/dynamic and qeq/fire only)
\sphinxstyleemphasis{qstep} value = time step size for damped dynamics charge solver (qeq/dynamic and qeq/fire only)
\sphinxstyleemphasis{warn} value = do (=yes) or do not (=no) print a warning when the maximum number of iterations is reached
\end{sphinxalltt}

\end{itemize}


\subsection{Examples}
\label{\detokenize{fix_qeq:examples}}
\begin{sphinxVerbatim}[commandchars=\\\{\}]
\PYG{k}{fix }\PYG{n+nv+nvIdentifier}{1}\PYG{+w}{ }\PYG{n+nv+nvIdentifier}{all}\PYG{+w}{ }\PYG{n}{qeq}\PYG{o}{/}\PYG{n}{point}\PYG{+w}{ }\PYG{l+m}{1}\PYG{+w}{ }\PYG{l+m}{10}\PYG{+w}{ }\PYG{l+m}{1.0e\PYGZhy{}6}\PYG{+w}{ }\PYG{l+m}{200}\PYG{+w}{ }\PYG{n}{param.qeq1}
\PYG{k}{fix }\PYG{n+nv+nvIdentifier}{1}\PYG{+w}{ }\PYG{n+nv+nvIdentifier}{qeq}\PYG{+w}{ }\PYG{n}{qeq}\PYG{o}{/}\PYG{n}{shielded}\PYG{+w}{ }\PYG{l+m}{1}\PYG{+w}{ }\PYG{l+m}{8}\PYG{+w}{ }\PYG{l+m}{1.0e\PYGZhy{}6}\PYG{+w}{ }\PYG{l+m}{100}\PYG{+w}{ }\PYG{n}{param.qeq2}
\PYG{k}{fix }\PYG{n+nv+nvIdentifier}{1}\PYG{+w}{ }\PYG{n+nv+nvIdentifier}{all}\PYG{+w}{ }\PYG{n}{qeq}\PYG{o}{/}\PYG{n}{slater}\PYG{+w}{ }\PYG{l+m}{5}\PYG{+w}{ }\PYG{l+m}{10}\PYG{+w}{ }\PYG{l+m}{1.0e\PYGZhy{}6}\PYG{+w}{ }\PYG{l+m}{100}\PYG{+w}{ }\PYG{n}{params}\PYG{+w}{ }\PYG{n}{alpha}\PYG{+w}{ }\PYG{l+m}{0.2}
\PYG{k}{fix }\PYG{n+nv+nvIdentifier}{1}\PYG{+w}{ }\PYG{n+nv+nvIdentifier}{qeq}\PYG{+w}{ }\PYG{n}{qeq}\PYG{o}{/}\PYG{n}{dynamic}\PYG{+w}{ }\PYG{l+m}{1}\PYG{+w}{ }\PYG{l+m}{12}\PYG{+w}{ }\PYG{l+m}{1.0e\PYGZhy{}3}\PYG{+w}{ }\PYG{l+m}{100}\PYG{+w}{ }\PYG{n}{my\PYGZus{}qeq}
\PYG{k}{fix }\PYG{n+nv+nvIdentifier}{1}\PYG{+w}{ }\PYG{n+nv+nvIdentifier}{all}\PYG{+w}{ }\PYG{n}{qeq}\PYG{o}{/}\PYG{n}{fire}\PYG{+w}{ }\PYG{l+m}{1}\PYG{+w}{ }\PYG{l+m}{10}\PYG{+w}{ }\PYG{l+m}{1.0e\PYGZhy{}3}\PYG{+w}{ }\PYG{l+m}{100}\PYG{+w}{ }\PYG{n}{my\PYGZus{}qeq}\PYG{+w}{ }\PYG{n}{qdamp}\PYG{+w}{ }\PYG{l+m}{0.2}\PYG{+w}{ }\PYG{n}{qstep}\PYG{+w}{ }\PYG{l+m}{0.1}
\end{sphinxVerbatim}


\subsection{Description}
\label{\detokenize{fix_qeq:description}}
\sphinxAtStartPar
Perform the charge equilibration (QEq) method as described in
{\hyperref[\detokenize{fix_qeq:rappe1}]{\sphinxcrossref{\DUrole{std,std-ref}{(Rappe and Goddard)}}}} and formulated in {\hyperref[\detokenize{fix_qeq:nakano1}]{\sphinxcrossref{\DUrole{std,std-ref}{(Nakano)}}}} (also known as the matrix inversion method) and in
{\hyperref[\detokenize{fix_qeq:rick1}]{\sphinxcrossref{\DUrole{std,std-ref}{(Rick and Stuart)}}}} (also known as the extended Lagrangian
method) based on the electronegativity equilization principle.

\sphinxAtStartPar
These fixes can be used with any {\hyperref[\detokenize{pair_style::doc}]{\sphinxcrossref{\DUrole{doc}{pair style}}}} in
LAMMPS, so long as per\sphinxhyphen{}atom charges are defined.  The most typical
use\sphinxhyphen{}case is in conjunction with a {\hyperref[\detokenize{pair_style::doc}]{\sphinxcrossref{\DUrole{doc}{pair style}}}} that
performs charge equilibration periodically (e.g. every timestep), such
as the ReaxFF or Streitz\sphinxhyphen{}Mintmire potential.  But these fixes can also
be used with potentials that normally assume per\sphinxhyphen{}atom charges are fixed,
e.g. a {\hyperref[\detokenize{pair_buck::doc}]{\sphinxcrossref{\DUrole{doc}{Buckingham}}}} or {\hyperref[\detokenize{pair_lj::doc}]{\sphinxcrossref{\DUrole{doc}{LJ/Coulombic}}}}
potential.

\sphinxAtStartPar
Because the charge equilibration calculation is effectively independent
of the pair style, these fixes can also be used to perform a one\sphinxhyphen{}time
assignment of charges to atoms.  For example, you could define the QEq
fix, perform a zero\sphinxhyphen{}timestep run via the {\hyperref[\detokenize{run::doc}]{\sphinxcrossref{\DUrole{doc}{run}}}} command
without any pair style defined which would set per\sphinxhyphen{}atom charges (based
on the current atom configuration), then remove the fix via the
{\hyperref[\detokenize{unfix::doc}]{\sphinxcrossref{\DUrole{doc}{unfix}}}} command before performing further dynamics.

\begin{sphinxadmonition}{note}{Note:}
\sphinxAtStartPar
Computing and using charge values different from published
values defined for a fixed\sphinxhyphen{}charge potential like Buckingham or CHARMM
or AMBER, can have a strong effect on energies and forces, and
produces a different model than the published versions.
\end{sphinxadmonition}

\begin{sphinxadmonition}{note}{Note:}
\sphinxAtStartPar
The {\hyperref[\detokenize{fix_qeq_comb::doc}]{\sphinxcrossref{\DUrole{doc}{fix qeq/comb}}}} command must still be used to
perform charge equilibration with the {\hyperref[\detokenize{pair_comb::doc}]{\sphinxcrossref{\DUrole{doc}{COMB potential}}}}.  The {\hyperref[\detokenize{fix_qeq_reaxff::doc}]{\sphinxcrossref{\DUrole{doc}{fix qeq/reaxff}}}} command can be
used to perform charge equilibration with the {\hyperref[\detokenize{pair_reaxff::doc}]{\sphinxcrossref{\DUrole{doc}{ReaxFF force
field}}}}, although fix qeq/shielded yields the same
results as fix qeq/reaxff if \sphinxstyleemphasis{Nevery}, \sphinxstyleemphasis{cutoff}, and \sphinxstyleemphasis{tolerance}
are the same.  Eventually the fix qeq/reaxff command will be
deprecated.
\end{sphinxadmonition}

\sphinxAtStartPar
The QEq method minimizes the electrostatic energy of the system (or
equalizes the derivative of energy with respect to charge of all the
atoms) by adjusting the partial charge on individual atoms based on
interactions with their neighbors within \sphinxstyleemphasis{cutoff}.  It requires a few
parameters in the appropriate units for each atom type which are read
from a file specified by \sphinxstyleemphasis{qfile}.  The file has the following format

\begin{sphinxVerbatim}[commandchars=\\\{\}]
1 chi eta gamma zeta qcore
2 chi eta gamma zeta qcore
...
Ntype chi eta gamma zeta qcore
\end{sphinxVerbatim}

\sphinxAtStartPar
There have to be parameters given for every atom type. Wildcard entries
are possible using the same type range syntax as for “coeff” commands
(i.e., n*m, n*, *m, *). Later entries will overwrite previous ones.
Empty lines or any text following the pound sign (\#) are ignored.
Each line starts with the atom type followed by five parameters.
Only a subset of the parameters is used by each QEq style as described
below, thus the others can be set to 0.0 if desired, but all five
entries per line are required.
\begin{itemize}
\item {} 
\sphinxAtStartPar
\sphinxstyleemphasis{chi} = electronegativity in energy units

\item {} 
\sphinxAtStartPar
\sphinxstyleemphasis{eta} = self\sphinxhyphen{}Coulomb potential in energy units

\item {} 
\sphinxAtStartPar
\sphinxstyleemphasis{gamma} = shielded Coulomb constant defined by {\hyperref[\detokenize{fix_qeq:vanduin}]{\sphinxcrossref{\DUrole{std,std-ref}{ReaxFF force field}}}} in distance units

\item {} 
\sphinxAtStartPar
\sphinxstyleemphasis{zeta} = Slater type orbital exponent defined by the {\hyperref[\detokenize{fix_qeq:streitz1}]{\sphinxcrossref{\DUrole{std,std-ref}{Streitz\sphinxhyphen{}Mintmire}}}} potential in reverse distance units

\item {} 
\sphinxAtStartPar
\sphinxstyleemphasis{qcore} = charge of the nucleus defined by the {\hyperref[\detokenize{fix_qeq:streitz1}]{\sphinxcrossref{\DUrole{std,std-ref}{Streitz\sphinxhyphen{}Mintmire potential}}}} potential in charge units

\end{itemize}

\sphinxAtStartPar
The fix qeq styles will print a warning if the charges are not
equilibrated within \sphinxstyleemphasis{tolerance} by \sphinxstyleemphasis{maxiter} steps, unless the
\sphinxstyleemphasis{warn} keyword is used with “no” as argument.  This latter option
may be useful for testing and benchmarking purposes, as it allows
to use a fixed number of QEq iterations when \sphinxstyleemphasis{tolerance} is set
to a small enough value to always reach the \sphinxstyleemphasis{maxiter} limit.  Turning
off warnings will avoid the excessive output in that case.

\sphinxAtStartPar
The \sphinxstyleemphasis{qeq/point} style describes partial charges on atoms as point
charges.  Interaction between a pair of charged particles is 1/r,
which is the simplest description of the interaction between charges.
Only the \sphinxstyleemphasis{chi} and \sphinxstyleemphasis{eta} parameters from the \sphinxstyleemphasis{qfile} file are used.
Note that Coulomb catastrophe can occur if repulsion between the pair
of charged particles is too weak.  This style solves partial charges
on atoms via the matrix inversion method.  A tolerance of 1.0e\sphinxhyphen{}6 is
usually a good number.

\sphinxAtStartPar
The \sphinxstyleemphasis{qeq/shielded} style describes partial charges on atoms also as
point charges, but uses a shielded Coulomb potential to describe the
interaction between a pair of charged particles.  Interaction through
the shielded Coulomb is given by equation (13) of the {\hyperref[\detokenize{fix_qeq:vanduin}]{\sphinxcrossref{\DUrole{std,std-ref}{ReaxFF force
field}}}} paper.  The shielding accounts for charge overlap
between charged particles at small separation.  This style is the same
as {\hyperref[\detokenize{fix_qeq_reaxff::doc}]{\sphinxcrossref{\DUrole{doc}{fix qeq/reaxff}}}}, and can be used with
{\hyperref[\detokenize{pair_reaxff::doc}]{\sphinxcrossref{\DUrole{doc}{pair\_style reaxff}}}}.  Only the \sphinxstyleemphasis{chi}, \sphinxstyleemphasis{eta}, and
\sphinxstyleemphasis{gamma} parameters from the \sphinxstyleemphasis{qfile} file are used. When using the string
\sphinxstyleemphasis{reaxff} as filename, these parameters are extracted directly from an
active \sphinxstyleemphasis{reaxff} pair style.  This style solves partial charges on atoms
via the matrix inversion method.  A tolerance of 1.0e\sphinxhyphen{}6 is usually a
good number.

\sphinxAtStartPar
The \sphinxstyleemphasis{qeq/slater} style describes partial charges on atoms as spherical
charge densities centered around atoms via the Slater 1\sphinxstyleemphasis{s} orbital, so
that the interaction between a pair of charged particles is the product
of two Slater 1\sphinxstyleemphasis{s} orbitals.  The expression for the Slater 1\sphinxstyleemphasis{s}
orbital is given under equation (6) of the {\hyperref[\detokenize{fix_qeq:streitz1}]{\sphinxcrossref{\DUrole{std,std-ref}{Streitz\sphinxhyphen{}Mintmire}}}} paper.  Only the \sphinxstyleemphasis{chi}, \sphinxstyleemphasis{eta}, \sphinxstyleemphasis{zeta}, and \sphinxstyleemphasis{qcore}
parameters from the \sphinxstyleemphasis{qfile} file are used. When using the string
\sphinxstyleemphasis{coul/streitz} as filename, these parameters are extracted directly from
an active \sphinxstyleemphasis{coul/streitz} pair style.  This style solves partial charges
on atoms via the matrix inversion method.  A tolerance of 1.0e\sphinxhyphen{}6 is
usually a good number.  Keyword \sphinxstyleemphasis{alpha} can be used to change the Slater
type orbital exponent.

\sphinxAtStartPar
The \sphinxstyleemphasis{qeq/dynamic} style describes partial charges on atoms as point
charges that interact through 1/r, but the extended Lagrangian method is
used to solve partial charges on atoms.  Only the \sphinxstyleemphasis{chi} and \sphinxstyleemphasis{eta}
parameters from the \sphinxstyleemphasis{qfile} file are used.  Note that Coulomb
catastrophe can occur if repulsion between the pair of charged particles
is too weak.  A tolerance of 1.0e\sphinxhyphen{}3 is usually a good number.  Keyword
\sphinxstyleemphasis{qdamp} can be used to change the damping factor, while keyword \sphinxstyleemphasis{qstep}
can be used to change the time step size.

\sphinxAtStartPar
The {\hyperref[\detokenize{fix_qeq:shan}]{\sphinxcrossref{\DUrole{std,std-ref}{*qeq/fire*}}}} style describes the same charge model
and charge solver as the \sphinxstyleemphasis{qeq/dynamic} style, but employs a FIRE
minimization algorithm to solve for equilibrium charges.  Keyword
\sphinxstyleemphasis{qdamp} can be used to change the damping factor, while keyword \sphinxstyleemphasis{qstep}
can be used to change the time step size.

\sphinxAtStartPar
Note that \sphinxstyleemphasis{qeq/point}, \sphinxstyleemphasis{qeq/shielded}, and \sphinxstyleemphasis{qeq/slater} describe
different charge models, whereas the matrix inversion method and the
extended Lagrangian method (\sphinxstyleemphasis{qeq/dynamic} and \sphinxstyleemphasis{qeq/fire}) are
different solvers.

\sphinxAtStartPar
Note that \sphinxstyleemphasis{qeq/point}, \sphinxstyleemphasis{qeq/dynamic} and \sphinxstyleemphasis{qeq/fire} styles all
describe charges as point charges that interact through 1/r
relationship, but solve partial charges on atoms using different
solvers.  These three styles should yield comparable results if the QEq
parameters and \sphinxstyleemphasis{Nevery}, \sphinxstyleemphasis{cutoff}, and \sphinxstyleemphasis{tolerance} are the same.
Style \sphinxstyleemphasis{qeq/point} is typically faster, \sphinxstyleemphasis{qeq/dynamic} scales better on
larger sizes, and \sphinxstyleemphasis{qeq/fire} is faster than \sphinxstyleemphasis{qeq/dynamic}.

\begin{sphinxadmonition}{note}{Note:}
\sphinxAtStartPar
In order to solve the self\sphinxhyphen{}consistent equations for electronegativity
equalization, LAMMPS imposes the additional constraint that all the
charges in the fix group must add up to zero.  The initial charge
assignments should also satisfy this constraint.  LAMMPS will print a
warning if that is not the case.
\end{sphinxadmonition}

\begin{sphinxadmonition}{note}{Note:}
\sphinxAtStartPar
Developing QEq parameters (chi, eta, gamma, zeta, and qcore) is
non\sphinxhyphen{}trivial.  Charges on atoms are not guaranteed to equilibrate with
arbitrary choices of these parameters.  We do not develop these QEq
parameters.  See the examples/qeq directory for some examples.
\end{sphinxadmonition}


\subsection{Restart, fix\_modify, output, run start/stop, minimize info}
\label{\detokenize{fix_qeq:restart-fix-modify-output-run-start-stop-minimize-info}}
\sphinxAtStartPar
No information about these fixes is written to {\hyperref[\detokenize{restart::doc}]{\sphinxcrossref{\DUrole{doc}{binary restart
files}}}}.  No global scalar or vector or per\sphinxhyphen{}atom quantities are
stored by these fixes for access by various {\hyperref[\detokenize{Howto_output::doc}]{\sphinxcrossref{\DUrole{doc}{output commands}}}}.  No parameter of these fixes can be used with the
\sphinxstyleemphasis{start/stop} keywords of the {\hyperref[\detokenize{run::doc}]{\sphinxcrossref{\DUrole{doc}{run}}}} command.

\sphinxAtStartPar
Thexe fixes are invoked during {\hyperref[\detokenize{minimize::doc}]{\sphinxcrossref{\DUrole{doc}{energy minimization}}}}.


\subsection{Restrictions}
\label{\detokenize{fix_qeq:restrictions}}
\sphinxAtStartPar
These fixes are part of the QEQ package.  They are only enabled if
LAMMPS was built with that package.  See the {\hyperref[\detokenize{Build_package::doc}]{\sphinxcrossref{\DUrole{doc}{Build package}}}} page for more info.

\sphinxAtStartPar
These qeq fixes will ignore electric field contributions from
{\hyperref[\detokenize{fix_efield::doc}]{\sphinxcrossref{\DUrole{doc}{fix efield}}}}.


\subsection{Related commands}
\label{\detokenize{fix_qeq:related-commands}}
\sphinxAtStartPar
{\hyperref[\detokenize{fix_qeq_reaxff::doc}]{\sphinxcrossref{\DUrole{doc}{fix qeq/reaxff}}}}, {\hyperref[\detokenize{fix_qeq_comb::doc}]{\sphinxcrossref{\DUrole{doc}{fix qeq/comb}}}}


\subsection{Default}
\label{\detokenize{fix_qeq:default}}
\sphinxAtStartPar
warn yes


\bigskip\hrule\bigskip

\phantomsection\label{\detokenize{fix_qeq:rappe1}}
\sphinxAtStartPar
\sphinxstylestrong{(Rappe and Goddard)} A. K. Rappe and W. A. Goddard III, J Physical
Chemistry, 95, 3358\sphinxhyphen{}3363 (1991).

\phantomsection\label{\detokenize{fix_qeq:nakano1}}
\sphinxAtStartPar
\sphinxstylestrong{(Nakano)} A. Nakano, Computer Physics Communications, 104, 59\sphinxhyphen{}69 (1997).

\phantomsection\label{\detokenize{fix_qeq:rick1}}
\sphinxAtStartPar
\sphinxstylestrong{(Rick and Stuart)} S. W. Rick, S. J. Stuart, B. J. Berne, J Chemical Physics
101, 16141 (1994).

\phantomsection\label{\detokenize{fix_qeq:streitz1}}
\sphinxAtStartPar
\sphinxstylestrong{(Streitz\sphinxhyphen{}Mintmire)} F. H. Streitz, J. W. Mintmire, Physical Review B, 50,
16, 11996 (1994)

\phantomsection\label{\detokenize{fix_qeq:vanduin}}
\sphinxAtStartPar
\sphinxstylestrong{(ReaxFF)} A. C. T. van Duin, S. Dasgupta, F. Lorant, W. A. Goddard III, J
Physical Chemistry, 105, 9396\sphinxhyphen{}9049 (2001)

\phantomsection\label{\detokenize{fix_qeq:shan}}
\sphinxAtStartPar
\sphinxstylestrong{(QEq/Fire)} T.\sphinxhyphen{}R. Shan, A. P. Thompson, S. J. Plimpton, in preparation

\sphinxstepscope

\index{fix qeq/comb@\spxentry{fix qeq/comb}}\ignorespaces 
\index{fix qeq/comb/omp@\spxentry{fix qeq/comb/omp}}\ignorespaces 

\section{fix qeq/comb command}
\label{\detokenize{fix_qeq_comb:fix-qeq-comb-command}}\label{\detokenize{fix_qeq_comb:index-1}}\label{\detokenize{fix_qeq_comb:index-0}}\label{\detokenize{fix_qeq_comb::doc}}
\sphinxAtStartPar
Accelerator Variants: \sphinxstyleemphasis{qeq/comb/omp}


\subsection{Syntax}
\label{\detokenize{fix_qeq_comb:syntax}}
\begin{sphinxVerbatim}[commandchars=\\\{\}]
\PYG{k}{fix }\PYG{n+nv+nvIdentifier}{ID}\PYG{+w}{ }\PYG{n+nv+nvIdentifier}{group\PYGZhy{}ID}\PYG{+w}{ }\PYG{n}{qeq}\PYG{o}{/}\PYG{n}{comb}\PYG{+w}{ }\PYG{n}{Nevery}\PYG{+w}{ }\PYG{n}{precision}\PYG{+w}{ }\PYG{n}{keyword}\PYG{+w}{ }\PYG{n}{value}\PYG{+w}{ }\PYG{n}{...}
\end{sphinxVerbatim}
\begin{itemize}
\item {} 
\sphinxAtStartPar
ID, group\sphinxhyphen{}ID are documented in {\hyperref[\detokenize{fix::doc}]{\sphinxcrossref{\DUrole{doc}{fix}}}} command

\item {} 
\sphinxAtStartPar
qeq/comb = style name of this fix command

\item {} 
\sphinxAtStartPar
Nevery = perform charge equilibration every this many steps

\item {} 
\sphinxAtStartPar
precision = convergence criterion for charge equilibration

\item {} 
\sphinxAtStartPar
zero or more keyword/value pairs may be appended

\item {} 
\sphinxAtStartPar
keyword = \sphinxstyleemphasis{file}
\begin{sphinxalltt}
\sphinxstyleemphasis{file} value = filename
  filename = name of file to write QEQ equilibration info to
\end{sphinxalltt}

\end{itemize}


\subsection{Examples}
\label{\detokenize{fix_qeq_comb:examples}}
\begin{sphinxVerbatim}[commandchars=\\\{\}]
\PYG{k}{fix }\PYG{n+nv+nvIdentifier}{1}\PYG{+w}{ }\PYG{n+nv+nvIdentifier}{surface}\PYG{+w}{ }\PYG{n}{qeq}\PYG{o}{/}\PYG{n}{comb}\PYG{+w}{ }\PYG{l+m}{10}\PYG{+w}{ }\PYG{l+m}{0.0001}
\end{sphinxVerbatim}


\subsection{Description}
\label{\detokenize{fix_qeq_comb:description}}
\sphinxAtStartPar
Perform charge equilibration (QeQ) in conjunction with the COMB
(Charge\sphinxhyphen{}Optimized Many\sphinxhyphen{}Body) potential as described in
{\hyperref[\detokenize{fix_qeq_comb:comb-1}]{\sphinxcrossref{\DUrole{std,std-ref}{(COMB\_1)}}}} and {\hyperref[\detokenize{fix_qeq_comb:comb-2}]{\sphinxcrossref{\DUrole{std,std-ref}{(COMB\_2)}}}}.  It performs the charge
equilibration portion of the calculation using the so\sphinxhyphen{}called QEq
method, whereby the charge on each atom is adjusted to minimize the
energy of the system.  This fix can only be used with the COMB
potential; see the {\hyperref[\detokenize{fix_qeq_reaxff::doc}]{\sphinxcrossref{\DUrole{doc}{fix qeq/reaxff}}}} command for a QeQ
calculation that can be used with any potential.

\sphinxAtStartPar
Only charges on the atoms in the specified group are equilibrated.
The fix relies on the pair style (COMB in this case) to calculate the
per\sphinxhyphen{}atom electronegativity (effective force on the charges).  An
electronegativity equalization calculation (or QEq) is performed in an
iterative fashion, which in parallel requires communication at each
iteration for processors to exchange charge information about nearby
atoms with each other.  See {\hyperref[\detokenize{fix_qeq_comb:rappe-and-goddard}]{\sphinxcrossref{\DUrole{std,std-ref}{Rappe\_and\_Goddard}}}} and
{\hyperref[\detokenize{fix_qeq_comb:rick-and-stuart}]{\sphinxcrossref{\DUrole{std,std-ref}{Rick\_and\_Stuart}}}} for details.

\sphinxAtStartPar
During a run, charge equilibration is performed every \sphinxstyleemphasis{Nevery} time
steps.  Charge equilibration is also always enforced on the first step
of each run.  The \sphinxstyleemphasis{precision} argument controls the tolerance for the
difference in electronegativity for all atoms during charge
equilibration.  \sphinxstyleemphasis{Precision} is a trade\sphinxhyphen{}off between the cost of
performing charge equilibration (more iterations) and accuracy.

\sphinxAtStartPar
If the \sphinxstyleemphasis{file} keyword is used, then information about each
equilibration calculation is written to the specified file.

\begin{sphinxadmonition}{note}{Note:}
\sphinxAtStartPar
In order to solve the self\sphinxhyphen{}consistent equations for electronegativity
equalization, LAMMPS imposes the additional constraint that all the
charges in the fix group must add up to zero.  The initial charge
assignments should also satisfy this constraint.  LAMMPS will print a
warning if that is not the case.
\end{sphinxadmonition}


\bigskip\hrule\bigskip


\sphinxAtStartPar
Styles with a \sphinxstyleemphasis{gpu}, \sphinxstyleemphasis{intel}, \sphinxstyleemphasis{kk}, \sphinxstyleemphasis{omp}, or \sphinxstyleemphasis{opt} suffix are
functionally the same as the corresponding style without the suffix.
They have been optimized to run faster, depending on your available
hardware, as discussed on the {\hyperref[\detokenize{Speed_packages::doc}]{\sphinxcrossref{\DUrole{doc}{Accelerator packages}}}}
page.  The accelerated styles take the same arguments and should
produce the same results, except for round\sphinxhyphen{}off and precision issues.

\sphinxAtStartPar
These accelerated styles are part of the GPU, INTEL, KOKKOS,
OPENMP, and OPT packages, respectively.  They are only enabled if
LAMMPS was built with those packages.  See the {\hyperref[\detokenize{Build_package::doc}]{\sphinxcrossref{\DUrole{doc}{Build package}}}} page for more info.

\sphinxAtStartPar
You can specify the accelerated styles explicitly in your input script
by including their suffix, or you can use the {\hyperref[\detokenize{Run_options::doc}]{\sphinxcrossref{\DUrole{doc}{\sphinxhyphen{}suffix command\sphinxhyphen{}line switch}}}} when you invoke LAMMPS, or you can use the
{\hyperref[\detokenize{suffix::doc}]{\sphinxcrossref{\DUrole{doc}{suffix}}}} command in your input script.

\sphinxAtStartPar
See the {\hyperref[\detokenize{Speed_packages::doc}]{\sphinxcrossref{\DUrole{doc}{Accelerator packages}}}} page for more
instructions on how to use the accelerated styles effectively.


\bigskip\hrule\bigskip



\subsection{Restart, fix\_modify, output, run start/stop, minimize info}
\label{\detokenize{fix_qeq_comb:restart-fix-modify-output-run-start-stop-minimize-info}}
\sphinxAtStartPar
No information about this fix is written to {\hyperref[\detokenize{restart::doc}]{\sphinxcrossref{\DUrole{doc}{binary restart files}}}}.

\sphinxAtStartPar
The {\hyperref[\detokenize{fix_modify::doc}]{\sphinxcrossref{\DUrole{doc}{fix\_modify}}}} \sphinxstyleemphasis{respa} option is supported by this
fix. This allows to set at which level of the {\hyperref[\detokenize{run_style::doc}]{\sphinxcrossref{\DUrole{doc}{r\sphinxhyphen{}RESPA}}}}
integrator the fix is performing charge equilibration. Default is
the outermost level.

\sphinxAtStartPar
This fix produces a per\sphinxhyphen{}atom vector which can be accessed by various
{\hyperref[\detokenize{Howto_output::doc}]{\sphinxcrossref{\DUrole{doc}{output commands}}}}.  The vector stores the gradient
of the charge on each atom.  The per\sphinxhyphen{}atom values be accessed on any
timestep.

\sphinxAtStartPar
No parameter of this fix can be used with the \sphinxstyleemphasis{start/stop} keywords of
the {\hyperref[\detokenize{run::doc}]{\sphinxcrossref{\DUrole{doc}{run}}}} command.

\sphinxAtStartPar
This fix can be invoked during {\hyperref[\detokenize{minimize::doc}]{\sphinxcrossref{\DUrole{doc}{energy minimization}}}}.


\subsection{Restrictions}
\label{\detokenize{fix_qeq_comb:restrictions}}
\sphinxAtStartPar
This fix command currently only supports {\hyperref[\detokenize{pair_comb::doc}]{\sphinxcrossref{\DUrole{doc}{pair style *comb*}}}}.


\subsection{Related commands}
\label{\detokenize{fix_qeq_comb:related-commands}}
\sphinxAtStartPar
{\hyperref[\detokenize{pair_comb::doc}]{\sphinxcrossref{\DUrole{doc}{pair\_style comb}}}}


\subsection{Default}
\label{\detokenize{fix_qeq_comb:default}}
\sphinxAtStartPar
No file output is performed.


\bigskip\hrule\bigskip

\phantomsection\label{\detokenize{fix_qeq_comb:comb-1}}
\sphinxAtStartPar
\sphinxstylestrong{(COMB\_1)} J. Yu, S. B. Sinnott, S. R. Phillpot, Phys Rev B, 75, 085311 (2007),

\phantomsection\label{\detokenize{fix_qeq_comb:comb-2}}
\sphinxAtStartPar
\sphinxstylestrong{(COMB\_2)} T.\sphinxhyphen{}R. Shan, B. D. Devine, T. W. Kemper, S. B. Sinnott, S. R.
Phillpot, Phys Rev B, 81, 125328 (2010).

\phantomsection\label{\detokenize{fix_qeq_comb:rappe-and-goddard}}
\sphinxAtStartPar
\sphinxstylestrong{(Rappe\_and\_Goddard)} A. K. Rappe, W. A. Goddard, J Phys Chem 95, 3358
(1991).

\phantomsection\label{\detokenize{fix_qeq_comb:rick-and-stuart}}
\sphinxAtStartPar
\sphinxstylestrong{(Rick\_and\_Stuart)} S. W. Rick, S. J. Stuart, B. J. Berne, J Chem Phys
101, 16141 (1994).

\sphinxstepscope

\index{fix qeq/reaxff@\spxentry{fix qeq/reaxff}}\ignorespaces 
\index{fix qeq/reaxff/kk@\spxentry{fix qeq/reaxff/kk}}\ignorespaces 
\index{fix qeq/reaxff/omp@\spxentry{fix qeq/reaxff/omp}}\ignorespaces 

\section{fix qeq/reaxff command}
\label{\detokenize{fix_qeq_reaxff:fix-qeq-reaxff-command}}\label{\detokenize{fix_qeq_reaxff:index-2}}\label{\detokenize{fix_qeq_reaxff:index-1}}\label{\detokenize{fix_qeq_reaxff:index-0}}\label{\detokenize{fix_qeq_reaxff::doc}}
\sphinxAtStartPar
Accelerator Variants: \sphinxstyleemphasis{qeq/reaxff/kk}, \sphinxstyleemphasis{qeq/reaxff/omp}


\subsection{Syntax}
\label{\detokenize{fix_qeq_reaxff:syntax}}
\begin{sphinxVerbatim}[commandchars=\\\{\}]
\PYG{k}{fix }\PYG{n+nv+nvIdentifier}{ID}\PYG{+w}{ }\PYG{n+nv+nvIdentifier}{group\PYGZhy{}ID}\PYG{+w}{ }\PYG{n}{qeq}\PYG{o}{/}\PYG{n}{reaxff}\PYG{+w}{ }\PYG{n}{Nevery}\PYG{+w}{ }\PYG{n}{cutlo}\PYG{+w}{ }\PYG{n}{cuthi}\PYG{+w}{ }\PYG{n}{tolerance}\PYG{+w}{ }\PYG{n}{params}\PYG{+w}{ }\PYG{n}{args}
\end{sphinxVerbatim}
\begin{itemize}
\item {} 
\sphinxAtStartPar
ID, group\sphinxhyphen{}ID are documented in {\hyperref[\detokenize{fix::doc}]{\sphinxcrossref{\DUrole{doc}{fix}}}} command

\item {} 
\sphinxAtStartPar
qeq/reaxff = style name of this fix command

\item {} 
\sphinxAtStartPar
Nevery = perform QEq every this many steps

\item {} 
\sphinxAtStartPar
cutlo,cuthi = lo and hi cutoff for Taper radius

\item {} 
\sphinxAtStartPar
tolerance = precision to which charges will be equilibrated

\item {} 
\sphinxAtStartPar
params = reaxff or a filename

\item {} 
\sphinxAtStartPar
one or more keywords or keyword/value pairs may be appended
\begin{sphinxalltt}
keyword = \sphinxstyleemphasis{dual} or \sphinxstyleemphasis{maxiter} or \sphinxstyleemphasis{nowarn}
  \sphinxstyleemphasis{dual} = process S and T matrix in parallel (only for qeq/reaxff/omp)
  \sphinxstyleemphasis{maxiter} N = limit the number of iterations to \sphinxstyleemphasis{N}
  \sphinxstyleemphasis{nowarn} = do not print a warning message if the maximum number of iterations was reached
\end{sphinxalltt}

\end{itemize}


\subsection{Examples}
\label{\detokenize{fix_qeq_reaxff:examples}}
\begin{sphinxVerbatim}[commandchars=\\\{\}]
\PYG{k}{fix }\PYG{n+nv+nvIdentifier}{1}\PYG{+w}{ }\PYG{n+nv+nvIdentifier}{all}\PYG{+w}{ }\PYG{n}{qeq}\PYG{o}{/}\PYG{n}{reaxff}\PYG{+w}{ }\PYG{l+m}{1}\PYG{+w}{ }\PYG{l+m}{0.0}\PYG{+w}{ }\PYG{l+m}{10.0}\PYG{+w}{ }\PYG{l+m}{1.0e\PYGZhy{}6}\PYG{+w}{ }\PYG{n}{reaxff}
\PYG{k}{fix }\PYG{n+nv+nvIdentifier}{1}\PYG{+w}{ }\PYG{n+nv+nvIdentifier}{all}\PYG{+w}{ }\PYG{n}{qeq}\PYG{o}{/}\PYG{n}{reaxff}\PYG{+w}{ }\PYG{l+m}{1}\PYG{+w}{ }\PYG{l+m}{0.0}\PYG{+w}{ }\PYG{l+m}{10.0}\PYG{+w}{ }\PYG{l+m}{1.0e\PYGZhy{}6}\PYG{+w}{ }\PYG{n}{param.qeq}\PYG{+w}{ }\PYG{n}{maxiter}\PYG{+w}{ }\PYG{l+m}{500}
\end{sphinxVerbatim}


\subsection{Description}
\label{\detokenize{fix_qeq_reaxff:description}}
\sphinxAtStartPar
Perform the charge equilibration (QEq) method as described in
{\hyperref[\detokenize{fix_qeq_reaxff:rappe2}]{\sphinxcrossref{\DUrole{std,std-ref}{(Rappe and Goddard)}}}} and formulated in {\hyperref[\detokenize{fix_qeq_reaxff:nakano2}]{\sphinxcrossref{\DUrole{std,std-ref}{(Nakano)}}}}.  It is typically used in conjunction with the ReaxFF force
field model as implemented in the {\hyperref[\detokenize{pair_reaxff::doc}]{\sphinxcrossref{\DUrole{doc}{pair\_style reaxff}}}}
command, but it can be used with any potential in LAMMPS, so long as it
defines and uses charges on each atom.  The {\hyperref[\detokenize{fix_qeq_comb::doc}]{\sphinxcrossref{\DUrole{doc}{fix qeq/comb}}}} command should be used to perform charge equilibration
with the {\hyperref[\detokenize{pair_comb::doc}]{\sphinxcrossref{\DUrole{doc}{COMB potential}}}}.  For more technical details
about the charge equilibration performed by fix qeq/reaxff, see the
{\hyperref[\detokenize{fix_qeq_reaxff:qeq-aktulga}]{\sphinxcrossref{\DUrole{std,std-ref}{(Aktulga)}}}} paper.

\sphinxAtStartPar
The QEq method minimizes the electrostatic energy of the system by
adjusting the partial charge on individual atoms based on interactions
with their neighbors.  It requires some parameters for each atom type.
If the \sphinxstyleemphasis{params} setting above is the word “reaxff”, then these are
extracted from the {\hyperref[\detokenize{pair_reaxff::doc}]{\sphinxcrossref{\DUrole{doc}{pair\_style reaxff}}}} command and
the ReaxFF force field file it reads in.  If a file name is specified
for \sphinxstyleemphasis{params}, then the parameters are taken from the specified file
and the file must contain one line for each atom type.  The latter
form must be used when performing QeQ with a non\sphinxhyphen{}ReaxFF potential.
Each line should be formatted as follows:

\begin{sphinxVerbatim}[commandchars=\\\{\}]
itype chi eta gamma
\end{sphinxVerbatim}

\sphinxAtStartPar
where \sphinxstyleemphasis{itype} is the atom type from 1 to Ntypes, \sphinxstyleemphasis{chi} denotes the
electronegativity in eV, \sphinxstyleemphasis{eta} denotes the self\sphinxhyphen{}Coulomb
potential in eV, and \sphinxstyleemphasis{gamma} denotes the valence orbital
exponent.  Note that these 3 quantities are also in the ReaxFF
potential file, except that eta is defined here as twice the eta value
in the ReaxFF file. Note that unlike the rest of LAMMPS, the units
of this fix are hard\sphinxhyphen{}coded to be A, eV, and electronic charge.

\sphinxAtStartPar
The optional \sphinxstyleemphasis{dual} keyword allows to perform the optimization
of the S and T matrices in parallel. This is only supported for
the \sphinxstyleemphasis{qeq/reaxff/omp} style. Otherwise they are processed separately.
The \sphinxstyleemphasis{qeq/reaxff/kk} style always solves the S and T matrices in
parallel.

\sphinxAtStartPar
The optional \sphinxstyleemphasis{maxiter} keyword allows changing the max number
of iterations in the linear solver. The default value is 200.

\sphinxAtStartPar
The optional \sphinxstyleemphasis{nowarn} keyword silences the warning message printed
when the maximum number of iterations was reached.  This can be
useful for comparing serial and parallel results where having the
same fixed number of QEq iterations is desired, which can be achieved
by using a very small tolerance and setting \sphinxstyleemphasis{maxiter} to the desired
number of iterations.

\begin{sphinxadmonition}{note}{Note:}
\sphinxAtStartPar
In order to solve the self\sphinxhyphen{}consistent equations for electronegativity
equalization, LAMMPS imposes the additional constraint that all the
charges in the fix group must add up to zero.  The initial charge
assignments should also satisfy this constraint.  LAMMPS will print a
warning if that is not the case.
\end{sphinxadmonition}


\subsection{Restart, fix\_modify, output, run start/stop, minimize info}
\label{\detokenize{fix_qeq_reaxff:restart-fix-modify-output-run-start-stop-minimize-info}}
\sphinxAtStartPar
No information about this fix is written to {\hyperref[\detokenize{restart::doc}]{\sphinxcrossref{\DUrole{doc}{binary restart files}}}}.  This fix computes a global scalar (the number of
iterations) for access by various {\hyperref[\detokenize{Howto_output::doc}]{\sphinxcrossref{\DUrole{doc}{output commands}}}}.
No parameter of this fix can be used with the \sphinxstyleemphasis{start/stop} keywords of
the {\hyperref[\detokenize{run::doc}]{\sphinxcrossref{\DUrole{doc}{run}}}} command.

\sphinxAtStartPar
This fix is invoked during {\hyperref[\detokenize{minimize::doc}]{\sphinxcrossref{\DUrole{doc}{energy minimization}}}}.


\bigskip\hrule\bigskip


\sphinxAtStartPar
Styles with a \sphinxstyleemphasis{gpu}, \sphinxstyleemphasis{intel}, \sphinxstyleemphasis{kk}, \sphinxstyleemphasis{omp}, or \sphinxstyleemphasis{opt} suffix are
functionally the same as the corresponding style without the suffix.
They have been optimized to run faster, depending on your available
hardware, as discussed on the {\hyperref[\detokenize{Speed_packages::doc}]{\sphinxcrossref{\DUrole{doc}{Accelerator packages}}}}
page.  The accelerated styles take the same arguments and should
produce the same results, except for round\sphinxhyphen{}off and precision issues.

\sphinxAtStartPar
These accelerated styles are part of the GPU, INTEL, KOKKOS,
OPENMP, and OPT packages, respectively.  They are only enabled if
LAMMPS was built with those packages.  See the {\hyperref[\detokenize{Build_package::doc}]{\sphinxcrossref{\DUrole{doc}{Build package}}}} page for more info.

\sphinxAtStartPar
You can specify the accelerated styles explicitly in your input script
by including their suffix, or you can use the {\hyperref[\detokenize{Run_options::doc}]{\sphinxcrossref{\DUrole{doc}{\sphinxhyphen{}suffix command\sphinxhyphen{}line switch}}}} when you invoke LAMMPS, or you can use the
{\hyperref[\detokenize{suffix::doc}]{\sphinxcrossref{\DUrole{doc}{suffix}}}} command in your input script.

\sphinxAtStartPar
See the {\hyperref[\detokenize{Speed_packages::doc}]{\sphinxcrossref{\DUrole{doc}{Accelerator packages}}}} page for more
instructions on how to use the accelerated styles effectively.


\bigskip\hrule\bigskip



\subsection{Restrictions}
\label{\detokenize{fix_qeq_reaxff:restrictions}}
\sphinxAtStartPar
This fix is part of the REAXFF package.  It is only enabled if
LAMMPS was built with that package. See the {\hyperref[\detokenize{Build_package::doc}]{\sphinxcrossref{\DUrole{doc}{Build package}}}} page for more info.

\sphinxAtStartPar
This fix does not correctly handle interactions involving multiple
periodic images of the same atom.  Hence, it should not be used for
periodic cell dimensions less than 10 Angstroms.

\sphinxAtStartPar
This fix may be used in combination with {\hyperref[\detokenize{fix_efield::doc}]{\sphinxcrossref{\DUrole{doc}{fix efield}}}}
and will apply the external electric field during charge equilibration,
but there may be only one fix efield instance used and the electric field
vector may only have components in non\sphinxhyphen{}periodic directions. Equal\sphinxhyphen{}style
variables can be used for electric field vector components without any further
settings. Atom\sphinxhyphen{}style variables can be used for spatially\sphinxhyphen{}varying electric field
vector components, but the resulting electric potential must be specified
as an atom\sphinxhyphen{}style variable using the \sphinxstyleemphasis{potential} keyword for \sphinxtitleref{fix efield}.


\subsection{Related commands}
\label{\detokenize{fix_qeq_reaxff:related-commands}}
\sphinxAtStartPar
{\hyperref[\detokenize{pair_reaxff::doc}]{\sphinxcrossref{\DUrole{doc}{pair\_style reaxff}}}}, {\hyperref[\detokenize{fix_qeq::doc}]{\sphinxcrossref{\DUrole{doc}{fix qeq/shielded}}}}


\subsection{Default}
\label{\detokenize{fix_qeq_reaxff:default}}
\sphinxAtStartPar
maxiter 200


\bigskip\hrule\bigskip

\phantomsection\label{\detokenize{fix_qeq_reaxff:rappe2}}
\sphinxAtStartPar
\sphinxstylestrong{(Rappe)} Rappe and Goddard III, Journal of Physical Chemistry, 95,
3358\sphinxhyphen{}3363 (1991).

\phantomsection\label{\detokenize{fix_qeq_reaxff:nakano2}}
\sphinxAtStartPar
\sphinxstylestrong{(Nakano)} Nakano, Computer Physics Communications, 104, 59\sphinxhyphen{}69 (1997).

\phantomsection\label{\detokenize{fix_qeq_reaxff:qeq-aktulga}}
\sphinxAtStartPar
\sphinxstylestrong{(Aktulga)} Aktulga, Fogarty, Pandit, Grama, Parallel Computing, 38,
245\sphinxhyphen{}259 (2012).

\sphinxstepscope

\index{fix qmmm@\spxentry{fix qmmm}}\ignorespaces 

\section{fix qmmm command}
\label{\detokenize{fix_qmmm:fix-qmmm-command}}\label{\detokenize{fix_qmmm:index-0}}\label{\detokenize{fix_qmmm::doc}}

\subsection{Syntax}
\label{\detokenize{fix_qmmm:syntax}}
\begin{sphinxVerbatim}[commandchars=\\\{\}]
\PYG{k}{fix }\PYG{n+nv+nvIdentifier}{ID}\PYG{+w}{ }\PYG{n+nv+nvIdentifier}{group\PYGZhy{}ID}\PYG{+w}{ }\PYG{n}{qmmm}
\end{sphinxVerbatim}
\begin{itemize}
\item {} 
\sphinxAtStartPar
ID, group\sphinxhyphen{}ID are documented in {\hyperref[\detokenize{fix::doc}]{\sphinxcrossref{\DUrole{doc}{fix}}}} command

\item {} 
\sphinxAtStartPar
qmmm = style name of this fix command

\end{itemize}


\subsection{Examples}
\label{\detokenize{fix_qmmm:examples}}
\begin{sphinxVerbatim}[commandchars=\\\{\}]
\PYG{k}{fix }\PYG{n+nv+nvIdentifier}{1}\PYG{+w}{ }\PYG{n+nv+nvIdentifier}{qmol}\PYG{+w}{ }\PYG{n}{qmmm}
\end{sphinxVerbatim}


\subsection{Description}
\label{\detokenize{fix_qmmm:description}}
\sphinxAtStartPar
This fix provides functionality to enable a quantum
mechanics/molecular mechanics (QM/MM) coupling of LAMMPS to a quantum
mechanical code.  The current implementation only supports an ONIOM
style mechanical coupling to the \sphinxhref{https://www.quantum-espresso.org}{Quantum ESPRESSO} plane
wave DFT package.  Electrostatic coupling is in preparation and the
interface has been written in a manner that coupling to other QM codes
should be possible without changes to LAMMPS itself.

\sphinxAtStartPar
The interface code for this is in the lib/qmmm directory of the LAMMPS
distribution and is being made available at this early stage of
development in order to encourage contributions for interfaces to
other QM codes.  This will allow the LAMMPS side of the implementation
to be adapted if necessary before being finalized.

\sphinxAtStartPar
Details about how to use this fix are currently documented in the
description of the QM/MM interface code itself in lib/qmmm/README.


\subsection{Restart, fix\_modify, output, run start/stop, minimize info}
\label{\detokenize{fix_qmmm:restart-fix-modify-output-run-start-stop-minimize-info}}
\sphinxAtStartPar
No information about this fix is written to {\hyperref[\detokenize{restart::doc}]{\sphinxcrossref{\DUrole{doc}{binary restart files}}}}.  None of the {\hyperref[\detokenize{fix_modify::doc}]{\sphinxcrossref{\DUrole{doc}{fix\_modify}}}} options
are relevant to this fix.  No global scalar or vector or per\sphinxhyphen{}atom
quantities are stored by this fix for access by various {\hyperref[\detokenize{Howto_output::doc}]{\sphinxcrossref{\DUrole{doc}{output commands}}}}.  No parameter of this fix can be used
with the \sphinxstyleemphasis{start/stop} keywords of the {\hyperref[\detokenize{run::doc}]{\sphinxcrossref{\DUrole{doc}{run}}}} command.  This
fix is not invoked during {\hyperref[\detokenize{minimize::doc}]{\sphinxcrossref{\DUrole{doc}{energy minimization}}}}.


\subsection{Restrictions}
\label{\detokenize{fix_qmmm:restrictions}}
\sphinxAtStartPar
This fix is part of the QMMM package.  It is only enabled if
LAMMPS was built with that package. It also requires building a
library provided with LAMMPS.  See the {\hyperref[\detokenize{Build_package::doc}]{\sphinxcrossref{\DUrole{doc}{Build package}}}} page for more info.

\sphinxAtStartPar
The fix is only functional when LAMMPS is built as a library and
linked with a compatible QM program and a QM/MM front end into a QM/MM
executable.  See the lib/qmmm/README file for details.


\subsection{Related commands}
\label{\detokenize{fix_qmmm:related-commands}}
\sphinxAtStartPar
none


\subsection{Default}
\label{\detokenize{fix_qmmm:default}}
\sphinxAtStartPar
none

\sphinxstepscope

\index{fix qtb@\spxentry{fix qtb}}\ignorespaces 

\section{fix qtb command}
\label{\detokenize{fix_qtb:fix-qtb-command}}\label{\detokenize{fix_qtb:index-0}}\label{\detokenize{fix_qtb::doc}}

\subsection{Syntax}
\label{\detokenize{fix_qtb:syntax}}
\begin{sphinxVerbatim}[commandchars=\\\{\}]
\PYG{k}{fix }\PYG{n+nv+nvIdentifier}{ID}\PYG{+w}{ }\PYG{n+nv+nvIdentifier}{group\PYGZhy{}ID}\PYG{+w}{ }\PYG{n}{qtb}\PYG{+w}{ }\PYG{n}{keyword}\PYG{+w}{ }\PYG{n}{value}\PYG{+w}{ }\PYG{n}{...}
\end{sphinxVerbatim}
\begin{itemize}
\item {} 
\sphinxAtStartPar
ID, group\sphinxhyphen{}ID are documented in {\hyperref[\detokenize{fix::doc}]{\sphinxcrossref{\DUrole{doc}{fix}}}} command

\item {} 
\sphinxAtStartPar
qtb = style name of this fix

\item {} 
\sphinxAtStartPar
zero or more keyword/value pairs may be appended

\item {} 
\sphinxAtStartPar
keyword = \sphinxstyleemphasis{temp} or \sphinxstyleemphasis{damp} or \sphinxstyleemphasis{seed} or \sphinxstyleemphasis{f\_max} or \sphinxstyleemphasis{N\_f}
\begin{sphinxalltt}
\sphinxstyleemphasis{temp} value = target quantum temperature (temperature units)
\sphinxstyleemphasis{damp} value = damping parameter (time units) inverse of friction \sphinxstyleemphasis{gamma}
\sphinxstyleemphasis{seed} value = random number seed (positive integer)
\sphinxstyleemphasis{f\_max} value = upper cutoff frequency of the vibration spectrum (1/time units)
\sphinxstyleemphasis{N\_f} value = number of frequency bins (positive integer)
\end{sphinxalltt}

\end{itemize}


\subsection{Examples}
\label{\detokenize{fix_qtb:examples}}
\begin{sphinxVerbatim}[commandchars=\\\{\}]
\PYG{c}{\PYGZsh{} (liquid methane modeled with the REAX force field, real units)}
\PYG{k}{fix }\PYG{n+nv+nvIdentifier}{1}\PYG{+w}{ }\PYG{n+nv+nvIdentifier}{all}\PYG{+w}{ }\PYG{n}{nve}
\PYG{k}{fix }\PYG{n+nv+nvIdentifier}{1}\PYG{+w}{ }\PYG{n+nv+nvIdentifier}{all}\PYG{+w}{ }\PYG{n}{qtb}\PYG{+w}{ }\PYG{n}{temp}\PYG{+w}{ }\PYG{l+m}{110}\PYG{+w}{ }\PYG{n}{damp}\PYG{+w}{ }\PYG{l+m}{200}\PYG{+w}{ }\PYG{n}{seed}\PYG{+w}{ }\PYG{l+m}{35082}\PYG{+w}{ }\PYG{n}{f\PYGZus{}max}\PYG{+w}{ }\PYG{l+m}{0.3}\PYG{+w}{ }\PYG{n}{N\PYGZus{}f}\PYG{+w}{ }\PYG{l+m}{100}
\PYG{c}{\PYGZsh{} (quartz modeled with the BKS force field, metal units)}
\PYG{k}{fix }\PYG{n+nv+nvIdentifier}{2}\PYG{+w}{ }\PYG{n+nv+nvIdentifier}{all}\PYG{+w}{ }\PYG{n}{nph}\PYG{+w}{ }\PYG{n}{iso}\PYG{+w}{ }\PYG{l+m}{1.01325}\PYG{+w}{ }\PYG{l+m}{1.01325}\PYG{+w}{ }\PYG{l+m}{1}
\PYG{k}{fix }\PYG{n+nv+nvIdentifier}{2}\PYG{+w}{ }\PYG{n+nv+nvIdentifier}{all}\PYG{+w}{ }\PYG{n}{qtb}\PYG{+w}{ }\PYG{n}{temp}\PYG{+w}{ }\PYG{l+m}{300}\PYG{+w}{ }\PYG{n}{damp}\PYG{+w}{ }\PYG{l+m}{1}\PYG{+w}{ }\PYG{n}{seed}\PYG{+w}{ }\PYG{l+m}{47508}\PYG{+w}{ }\PYG{n}{f\PYGZus{}max}\PYG{+w}{ }\PYG{l+m}{120.0}\PYG{+w}{ }\PYG{n}{N\PYGZus{}f}\PYG{+w}{ }\PYG{l+m}{100}
\end{sphinxVerbatim}


\subsection{Description}
\label{\detokenize{fix_qtb:description}}
\sphinxAtStartPar
This command performs the quantum thermal bath scheme proposed by
{\hyperref[\detokenize{fix_qtb:dammak}]{\sphinxcrossref{\DUrole{std,std-ref}{(Dammak)}}}} to include self\sphinxhyphen{}consistent quantum nuclear effects,
when used in conjunction with the {\hyperref[\detokenize{fix_nve::doc}]{\sphinxcrossref{\DUrole{doc}{fix nve}}}} or {\hyperref[\detokenize{fix_nh::doc}]{\sphinxcrossref{\DUrole{doc}{fix nph}}}} commands.

\sphinxAtStartPar
Classical molecular dynamics simulation does not include any quantum
nuclear effect. Quantum treatment of the vibrational modes will
introduce zero point energy into the system, alter the energy power
spectrum and bias the heat capacity from the classical limit. Missing
all the quantum nuclear effects, classical MD cannot model systems at
temperatures lower than their classical limits. This effect is
especially important for materials with a large population of hydrogen
atoms and thus higher classical limits.

\sphinxAtStartPar
The equation of motion implemented by this command follows a Langevin
form:
\begin{equation*}
\begin{split}m_i a_i = f_i + R_i - m_i\gamma v_i\end{split}
\end{equation*}
\sphinxAtStartPar
Here \(m_i, a_i, f_i, R_i, \gamma, \textrm{and} v_i\)
represent in this order mass, acceleration, force exerted by all other atoms, random
force, frictional coefficient (the inverse of damping parameter damp),
and velocity. The random force \(R_i\) is “colored” so
that any vibrational mode with frequency \(\omega\) will have a
temperature\sphinxhyphen{}sensitive energy \(\theta(\omega,T)\) which
resembles the energy expectation for a quantum harmonic oscillator
with the same natural frequency:
\begin{equation*}
\begin{split}\theta(\omega T) = \frac{\hbar}{2} + \hbar\omega \left[\exp(\frac{\hbar\omega}{k_B T})-1 \right]^{-1}\end{split}
\end{equation*}
\sphinxAtStartPar
To efficiently generate the random forces, we employ the method
of {\hyperref[\detokenize{fix_qtb:barrat}]{\sphinxcrossref{\DUrole{std,std-ref}{(Barrat)}}}}, that circumvents the need to generate all
random forces for all times before the simulation. The memory
requirement of this approach is less demanding and independent
of the simulation duration. Since the total random force \(R_{tot}\)
does not necessarily vanish for a finite number of atoms,
\(R_i\) is replaced by \(R_i - \frac{R_{tot}}{N_{tot}}\)
to avoid collective motion of the system.

\sphinxAtStartPar
The \sphinxstyleemphasis{temp} parameter sets the target quantum temperature. LAMMPS will
still have an output temperature in its thermo style. That is the
instantaneous classical temperature \(T^{cl}\) derived from
the atom velocities at thermal equilibrium. A non\sphinxhyphen{}zero
\(T^{cl}\) will be present even when the quantum
temperature approaches zero. This is associated with zero\sphinxhyphen{}point energy
at low temperatures.
\begin{equation*}
\begin{split}T^{cl} = \sum \frac{m_i v_i^2}{3 N k_B}\end{split}
\end{equation*}
\sphinxAtStartPar
The \sphinxstyleemphasis{damp} parameter is specified in time units, and it equals the
inverse of the frictional coefficient \(\gamma\). \(\gamma\)
should be as small as possible but slightly larger than the timescale
of anharmonic coupling in the system which is about 10 ps to 100
ps. When \(\gamma\) is too large, it gives an energy spectrum that
differs from the desired Bose\sphinxhyphen{}Einstein spectrum. When \(\gamma\)
is too small, the quantum thermal bath coupling to the system will be
less significant than anharmonic effects, reducing to a classical
limit. We find that setting \(\gamma\) between 5 THz and 1 THz
could be appropriate depending on the system.

\sphinxAtStartPar
The random number \sphinxstyleemphasis{seed} is a positive integer used to initiate a
Marsaglia random number generator. Each processor uses the input seed
to generate its own unique seed and its own stream of random
numbers. Thus the dynamics of the system will not be identical on two
runs on different numbers of processors.

\sphinxAtStartPar
The \sphinxstyleemphasis{f\_max} parameter truncate the noise frequency domain so that
vibrational modes with frequencies higher than \sphinxstyleemphasis{f\_max} will not be
modulated. If we denote \(\Delta t\) as the time interval for the
MD integration, \sphinxstyleemphasis{f\_max} is always reset by the code to make
\(\alpha = (int)(2\) \sphinxstyleemphasis{f\_max} \(\Delta t)^{-1}\) a
positive integer and print out relative information. An appropriate
value for the cutoff frequency \sphinxstyleemphasis{f\_max} would be around 2\textasciitilde{}3 \(f_D\),
where \(f_D\) is the Debye frequency.

\sphinxAtStartPar
The \sphinxstyleemphasis{N\_f} parameter is the frequency grid size, the number of points
from 0 to \sphinxstyleemphasis{f\_max} in the frequency domain that will be
sampled. 3*2\sphinxstyleemphasis{N\_f} per\sphinxhyphen{}atom random numbers are required
in the random force generation and there could be as many atoms as in
the whole simulation that can migrate into every individual
processor. A larger \sphinxstyleemphasis{N\_f} provides a more accurate sampling of the
spectrum while consumes more memory.  With fixed \sphinxstyleemphasis{f\_max} and
\(\gamma\), \sphinxstyleemphasis{N\_f} should be big enough to converge the classical
temperature \(T^{cl}\) as a function of target quantum bath
temperature. Memory usage per processor could be from 10 to 100
MBytes.

\begin{sphinxadmonition}{note}{Note:}
\sphinxAtStartPar
Unlike the {\hyperref[\detokenize{fix_nh::doc}]{\sphinxcrossref{\DUrole{doc}{fix nvt}}}} command which performs
Nose/Hoover thermostatting AND time integration, this fix does NOT
perform time integration. It only modifies forces to a colored
thermostat. Thus you must use a separate time integration fix, like
{\hyperref[\detokenize{fix_nve::doc}]{\sphinxcrossref{\DUrole{doc}{fix nve}}}} or {\hyperref[\detokenize{fix_nh::doc}]{\sphinxcrossref{\DUrole{doc}{fix nph}}}} to actually
update the velocities and positions of atoms (as shown in the
examples). Likewise, this fix should not normally be used with
other fixes or commands that also specify system temperatures ,
e.g. {\hyperref[\detokenize{fix_nh::doc}]{\sphinxcrossref{\DUrole{doc}{fix nvt}}}} and {\hyperref[\detokenize{fix_temp_rescale::doc}]{\sphinxcrossref{\DUrole{doc}{fix temp/rescale}}}}.
\end{sphinxadmonition}


\bigskip\hrule\bigskip



\subsection{Restart, fix\_modify, output, run start/stop, minimize info}
\label{\detokenize{fix_qtb:restart-fix-modify-output-run-start-stop-minimize-info}}
\sphinxAtStartPar
No information about this fix is written to {\hyperref[\detokenize{restart::doc}]{\sphinxcrossref{\DUrole{doc}{binary restart files}}}}.  Because the state of the random number generator is not
saved in restart files, this means you cannot do “exact” restarts with
this fix. However, in a statistical sense, a restarted simulation
should produce similar behaviors of the system.

\sphinxAtStartPar
This fix is not invoked during {\hyperref[\detokenize{minimize::doc}]{\sphinxcrossref{\DUrole{doc}{energy minimization}}}}.


\bigskip\hrule\bigskip



\subsection{Restrictions}
\label{\detokenize{fix_qtb:restrictions}}
\sphinxAtStartPar
This fix style is part of the QTB package.  It is only enabled if
LAMMPS was built with that package. See the {\hyperref[\detokenize{Build_package::doc}]{\sphinxcrossref{\DUrole{doc}{Build package}}}} page for more info.


\bigskip\hrule\bigskip



\subsection{Related commands}
\label{\detokenize{fix_qtb:related-commands}}
\sphinxAtStartPar
{\hyperref[\detokenize{fix_nve::doc}]{\sphinxcrossref{\DUrole{doc}{fix nve}}}}, {\hyperref[\detokenize{fix_nh::doc}]{\sphinxcrossref{\DUrole{doc}{fix nph}}}},
{\hyperref[\detokenize{fix_langevin::doc}]{\sphinxcrossref{\DUrole{doc}{fix langevin}}}}, {\hyperref[\detokenize{fix_qbmsst::doc}]{\sphinxcrossref{\DUrole{doc}{fix qbmsst}}}}


\bigskip\hrule\bigskip



\subsection{Default}
\label{\detokenize{fix_qtb:default}}
\sphinxAtStartPar
The keyword defaults are temp = 300, damp = 1, seed = 880302,
f\_max=200.0 and N\_f = 100.


\bigskip\hrule\bigskip

\phantomsection\label{\detokenize{fix_qtb:dammak}}
\sphinxAtStartPar
\sphinxstylestrong{(Dammak)} Dammak, Chalopin, Laroche, Hayoun, and Greffet, Phys Rev
Lett, 103, 190601 (2009).

\phantomsection\label{\detokenize{fix_qtb:barrat}}
\sphinxAtStartPar
\sphinxstylestrong{(Barrat)} Barrat and Rodney, J. Stat. Phys, 144, 679 (2011).

\sphinxstepscope

\index{fix reaxff/bonds@\spxentry{fix reaxff/bonds}}\ignorespaces 
\index{fix reaxff/bonds/kk@\spxentry{fix reaxff/bonds/kk}}\ignorespaces 

\section{fix reaxff/bonds command}
\label{\detokenize{fix_reaxff_bonds:fix-reaxff-bonds-command}}\label{\detokenize{fix_reaxff_bonds:index-1}}\label{\detokenize{fix_reaxff_bonds:index-0}}\label{\detokenize{fix_reaxff_bonds::doc}}
\sphinxAtStartPar
Accelerator Variants: \sphinxstyleemphasis{reaxff/bonds/kk}


\subsection{Syntax}
\label{\detokenize{fix_reaxff_bonds:syntax}}
\begin{sphinxVerbatim}[commandchars=\\\{\}]
\PYG{k}{fix }\PYG{n+nv+nvIdentifier}{ID}\PYG{+w}{ }\PYG{n+nv+nvIdentifier}{group\PYGZhy{}ID}\PYG{+w}{ }\PYG{n}{reaxff}\PYG{o}{/}\PYG{n}{bonds}\PYG{+w}{ }\PYG{n}{Nevery}\PYG{+w}{ }\PYG{n}{filename}
\end{sphinxVerbatim}
\begin{itemize}
\item {} 
\sphinxAtStartPar
ID, group\sphinxhyphen{}ID are documented in {\hyperref[\detokenize{fix::doc}]{\sphinxcrossref{\DUrole{doc}{fix}}}} command

\item {} 
\sphinxAtStartPar
reax/bonds = style name of this fix command

\item {} 
\sphinxAtStartPar
Nevery = output interval in timesteps

\item {} 
\sphinxAtStartPar
filename = name of output file

\end{itemize}


\subsection{Examples}
\label{\detokenize{fix_reaxff_bonds:examples}}
\begin{sphinxVerbatim}[commandchars=\\\{\}]
\PYG{k}{fix }\PYG{n+nv+nvIdentifier}{1}\PYG{+w}{ }\PYG{n+nv+nvIdentifier}{all}\PYG{+w}{ }\PYG{n}{reaxff}\PYG{o}{/}\PYG{n}{bonds}\PYG{+w}{ }\PYG{l+m}{100}\PYG{+w}{ }\PYG{n}{bonds.reaxff}
\end{sphinxVerbatim}


\subsection{Description}
\label{\detokenize{fix_reaxff_bonds:description}}
\sphinxAtStartPar
Write out the bond information computed by the ReaxFF potential specified
by {\hyperref[\detokenize{pair_reaxff::doc}]{\sphinxcrossref{\DUrole{doc}{pair\_style reaxff}}}} in the exact same format as the
original stand\sphinxhyphen{}alone ReaxFF code of Adri van Duin.  The bond information
is written to \sphinxstyleemphasis{filename} on timesteps that are multiples of \sphinxstyleemphasis{Nevery},
including timestep 0.  For time\sphinxhyphen{}averaged chemical species analysis,
please see the {\hyperref[\detokenize{fix_reaxff_species::doc}]{\sphinxcrossref{\DUrole{doc}{fix reaxff/species}}}} command.

\sphinxAtStartPar
The specified group\sphinxhyphen{}ID is ignored by this fix.

\sphinxAtStartPar
The format of the output file should be reasonably self\sphinxhyphen{}explanatory.
The meaning of the column header abbreviations is as follows:
\begin{itemize}
\item {} 
\sphinxAtStartPar
id = atom id

\item {} 
\sphinxAtStartPar
type = atom type

\item {} 
\sphinxAtStartPar
nb = number of bonds

\item {} 
\sphinxAtStartPar
id\_1 = atom id of first bond

\item {} 
\sphinxAtStartPar
id\_nb = atom id of Nth bond

\item {} 
\sphinxAtStartPar
mol = molecule id

\item {} 
\sphinxAtStartPar
bo\_1 = bond order of first bond

\item {} 
\sphinxAtStartPar
bo\_nb = bond order of Nth bond

\item {} 
\sphinxAtStartPar
abo = atom bond order (sum of all bonds)

\item {} 
\sphinxAtStartPar
nlp = number of lone pairs

\item {} 
\sphinxAtStartPar
q = atomic charge

\end{itemize}

\sphinxAtStartPar
If the filename ends with “.gz”, the output file is written in gzipped
format.  A gzipped dump file will be about 3x smaller than the text
version, but will also take longer to write.


\bigskip\hrule\bigskip



\subsection{Restart, fix\_modify, output, run start/stop, minimize info}
\label{\detokenize{fix_reaxff_bonds:restart-fix-modify-output-run-start-stop-minimize-info}}
\sphinxAtStartPar
No information about this fix is written to {\hyperref[\detokenize{restart::doc}]{\sphinxcrossref{\DUrole{doc}{binary restart files}}}}.  None of the {\hyperref[\detokenize{fix_modify::doc}]{\sphinxcrossref{\DUrole{doc}{fix\_modify}}}} options
are relevant to this fix.  No global or per\sphinxhyphen{}atom quantities are stored
by this fix for access by various {\hyperref[\detokenize{Howto_output::doc}]{\sphinxcrossref{\DUrole{doc}{output commands}}}}.
No parameter of this fix can be used with the \sphinxstyleemphasis{start/stop} keywords of
the {\hyperref[\detokenize{run::doc}]{\sphinxcrossref{\DUrole{doc}{run}}}} command.  This fix is not invoked during {\hyperref[\detokenize{minimize::doc}]{\sphinxcrossref{\DUrole{doc}{energy minimization}}}}.


\bigskip\hrule\bigskip


\sphinxAtStartPar
Styles with a \sphinxstyleemphasis{gpu}, \sphinxstyleemphasis{intel}, \sphinxstyleemphasis{kk}, \sphinxstyleemphasis{omp}, or \sphinxstyleemphasis{opt} suffix are
functionally the same as the corresponding style without the suffix.
They have been optimized to run faster, depending on your available
hardware, as discussed on the {\hyperref[\detokenize{Speed_packages::doc}]{\sphinxcrossref{\DUrole{doc}{Accelerator packages}}}}
page.  The accelerated styles take the same arguments and should
produce the same results, except for round\sphinxhyphen{}off and precision issues.

\sphinxAtStartPar
These accelerated styles are part of the GPU, INTEL, KOKKOS,
OPENMP, and OPT packages, respectively.  They are only enabled if
LAMMPS was built with those packages.  See the {\hyperref[\detokenize{Build_package::doc}]{\sphinxcrossref{\DUrole{doc}{Build package}}}} page for more info.

\sphinxAtStartPar
You can specify the accelerated styles explicitly in your input script
by including their suffix, or you can use the {\hyperref[\detokenize{Run_options::doc}]{\sphinxcrossref{\DUrole{doc}{\sphinxhyphen{}suffix command\sphinxhyphen{}line switch}}}} when you invoke LAMMPS, or you can use the
{\hyperref[\detokenize{suffix::doc}]{\sphinxcrossref{\DUrole{doc}{suffix}}}} command in your input script.

\sphinxAtStartPar
See the {\hyperref[\detokenize{Speed_packages::doc}]{\sphinxcrossref{\DUrole{doc}{Accelerator packages}}}} page for more
instructions on how to use the accelerated styles effectively.


\bigskip\hrule\bigskip



\subsection{Restrictions}
\label{\detokenize{fix_reaxff_bonds:restrictions}}
\sphinxAtStartPar
The fix reaxff/bonds command requires that the {\hyperref[\detokenize{pair_reaxff::doc}]{\sphinxcrossref{\DUrole{doc}{pair\_style reaxff}}}} is invoked.  This fix is part of the
REAXFF package.  It is only enabled if LAMMPS was built with that
package.  See the {\hyperref[\detokenize{Build_package::doc}]{\sphinxcrossref{\DUrole{doc}{Build package}}}} page for more
info.

\sphinxAtStartPar
To write gzipped bond files, you must compile LAMMPS with the
\sphinxhyphen{}DLAMMPS\_GZIP option.


\subsection{Related commands}
\label{\detokenize{fix_reaxff_bonds:related-commands}}
\sphinxAtStartPar
{\hyperref[\detokenize{pair_reaxff::doc}]{\sphinxcrossref{\DUrole{doc}{pair\_style reaxff}}}}, {\hyperref[\detokenize{fix_reaxff_species::doc}]{\sphinxcrossref{\DUrole{doc}{fix reaxff/species}}}}


\subsection{Default}
\label{\detokenize{fix_reaxff_bonds:default}}
\sphinxAtStartPar
none

\sphinxstepscope

\index{fix reaxff/species@\spxentry{fix reaxff/species}}\ignorespaces 
\index{fix reaxff/species/kk@\spxentry{fix reaxff/species/kk}}\ignorespaces 

\section{fix reaxff/species command}
\label{\detokenize{fix_reaxff_species:fix-reaxff-species-command}}\label{\detokenize{fix_reaxff_species:index-1}}\label{\detokenize{fix_reaxff_species:index-0}}\label{\detokenize{fix_reaxff_species::doc}}
\sphinxAtStartPar
Accelerator Variants: \sphinxstyleemphasis{reaxff/species/kk}


\subsection{Syntax}
\label{\detokenize{fix_reaxff_species:syntax}}
\begin{sphinxVerbatim}[commandchars=\\\{\}]
\PYG{k}{fix }\PYG{n+nv+nvIdentifier}{ID}\PYG{+w}{ }\PYG{n+nv+nvIdentifier}{group\PYGZhy{}ID}\PYG{+w}{ }\PYG{n}{reaxff}\PYG{o}{/}\PYG{n}{species}\PYG{+w}{ }\PYG{n}{Nevery}\PYG{+w}{ }\PYG{n}{Nrepeat}\PYG{+w}{ }\PYG{n}{Nfreq}\PYG{+w}{ }\PYG{n}{filename}\PYG{+w}{ }\PYG{n}{keyword}\PYG{+w}{ }\PYG{n}{value}\PYG{+w}{ }\PYG{n}{...}
\end{sphinxVerbatim}
\begin{itemize}
\item {} 
\sphinxAtStartPar
ID, group\sphinxhyphen{}ID are documented in {\hyperref[\detokenize{fix::doc}]{\sphinxcrossref{\DUrole{doc}{fix}}}} command

\item {} 
\sphinxAtStartPar
reaxff/species = style name of this command

\item {} 
\sphinxAtStartPar
Nevery = sample bond\sphinxhyphen{}order every this many timesteps

\item {} 
\sphinxAtStartPar
Nrepeat = \# of bond\sphinxhyphen{}order samples used for calculating averages

\item {} 
\sphinxAtStartPar
Nfreq = calculate average bond\sphinxhyphen{}order every this many timesteps

\item {} 
\sphinxAtStartPar
filename = name of output file

\item {} 
\sphinxAtStartPar
zero or more keyword/value pairs may be appended

\item {} 
\sphinxAtStartPar
keyword = \sphinxstyleemphasis{cutoff} or \sphinxstyleemphasis{element} or \sphinxstyleemphasis{position} or \sphinxstyleemphasis{delete} or \sphinxstyleemphasis{delete\_rate\_limit}
\begin{sphinxalltt}
\sphinxstyleemphasis{cutoff} value = I J Cutoff
  I, J = atom types (see asterisk form below)
  Cutoff = Bond\sphinxhyphen{}order cutoff value for this pair of atom types
\sphinxstyleemphasis{element} value = Element1, Element2, ...
\sphinxstyleemphasis{position} value = posfreq filepos
  posfreq = write position files every this many timestep
  filepos = name of position output file
\sphinxstyleemphasis{delete} value = filedel keyword value
  filedel = name of delete species output file
  keyword = \sphinxstyleemphasis{specieslist} or \sphinxstyleemphasis{masslimit}
    \sphinxstyleemphasis{specieslist} value = Nspecies Species1 Species2 ...
      Nspecies = number of species in list
    \sphinxstyleemphasis{masslimit} value = massmin massmax
      massmin = minimum molecular weight of species to delete
      massmax = maximum molecular weight of species to delete
\sphinxstyleemphasis{delete\_rate\_limit} value = Nlimit Nsteps
      Nlimit = maximum number of deletions allowed to occur within interval
      Nsteps = the interval (number of timesteps) over which to count deletions
\end{sphinxalltt}

\end{itemize}


\subsection{Examples}
\label{\detokenize{fix_reaxff_species:examples}}
\begin{sphinxVerbatim}[commandchars=\\\{\}]
\PYG{k}{fix }\PYG{n+nv+nvIdentifier}{1}\PYG{+w}{ }\PYG{n+nv+nvIdentifier}{all}\PYG{+w}{ }\PYG{n}{reaxff}\PYG{o}{/}\PYG{n}{species}\PYG{+w}{ }\PYG{l+m}{10}\PYG{+w}{ }\PYG{l+m}{10}\PYG{+w}{ }\PYG{l+m}{100}\PYG{+w}{ }\PYG{n}{species.out}
\PYG{k}{fix }\PYG{n+nv+nvIdentifier}{1}\PYG{+w}{ }\PYG{n+nv+nvIdentifier}{all}\PYG{+w}{ }\PYG{n}{reaxff}\PYG{o}{/}\PYG{n}{species}\PYG{+w}{ }\PYG{l+m}{1}\PYG{+w}{ }\PYG{l+m}{2}\PYG{+w}{ }\PYG{l+m}{20}\PYG{+w}{ }\PYG{n}{species.out}\PYG{+w}{ }\PYG{n}{cutoff}\PYG{+w}{ }\PYG{l+m}{1}\PYG{+w}{ }\PYG{l+m}{1}\PYG{+w}{ }\PYG{l+m}{0.40}\PYG{+w}{ }\PYG{n}{cutoff}\PYG{+w}{ }\PYG{l+m}{1}\PYG{+w}{ }\PYG{l+m}{2}\PYG{o}{*}\PYG{l+m}{3}\PYG{+w}{ }\PYG{l+m}{0.55}
\PYG{k}{fix }\PYG{n+nv+nvIdentifier}{1}\PYG{+w}{ }\PYG{n+nv+nvIdentifier}{all}\PYG{+w}{ }\PYG{n}{reaxff}\PYG{o}{/}\PYG{n}{species}\PYG{+w}{ }\PYG{l+m}{1}\PYG{+w}{ }\PYG{l+m}{100}\PYG{+w}{ }\PYG{l+m}{100}\PYG{+w}{ }\PYG{n}{species.out}\PYG{+w}{ }\PYG{n}{element}\PYG{+w}{ }\PYG{n}{Au}\PYG{+w}{ }\PYG{n}{O}\PYG{+w}{ }\PYG{n}{H}\PYG{+w}{ }\PYG{n}{position}\PYG{+w}{ }\PYG{l+m}{1000}\PYG{+w}{ }\PYG{n}{AuOH.pos}
\PYG{k}{fix }\PYG{n+nv+nvIdentifier}{1}\PYG{+w}{ }\PYG{n+nv+nvIdentifier}{all}\PYG{+w}{ }\PYG{n}{reaxff}\PYG{o}{/}\PYG{n}{species}\PYG{+w}{ }\PYG{l+m}{1}\PYG{+w}{ }\PYG{l+m}{100}\PYG{+w}{ }\PYG{l+m}{100}\PYG{+w}{ }\PYG{n}{species.out}\PYG{+w}{ }\PYG{n}{delete}\PYG{+w}{ }\PYG{n}{species.del}\PYG{+w}{ }\PYG{n}{masslimit}\PYG{+w}{ }\PYG{l+m}{0}\PYG{+w}{ }\PYG{l+m}{50}
\end{sphinxVerbatim}


\subsection{Description}
\label{\detokenize{fix_reaxff_species:description}}
\sphinxAtStartPar
Write out the chemical species information computed by the ReaxFF
potential specified by {\hyperref[\detokenize{pair_reaxff::doc}]{\sphinxcrossref{\DUrole{doc}{pair\_style reaxff}}}}.
Bond\sphinxhyphen{}order values (either averaged or instantaneous, depending on
value of \sphinxstyleemphasis{Nrepeat}) are used to determine chemical bonds.  Every
\sphinxstyleemphasis{Nfreq} timesteps, chemical species information is written to
\sphinxstyleemphasis{filename} as a two line output.  The first line is a header
containing labels. The second line consists of the following:
timestep, total number of molecules, total number of distinct species,
number of molecules of each species.  In this context, “species” means
a unique molecule.  The chemical formula of each species is given in
the first line.

\begin{sphinxadmonition}{warning}{Warning:}
\sphinxAtStartPar
In order to compute averaged data, it is required that there are no
neighbor list rebuilds for at least Nrepeat*Nevery steps preceding
each \sphinxstyleemphasis{Nfreq} step.  For that reason, fix \sphinxstyleemphasis{reaxff/species} may
change your neighbor list settings.  Reneighboring will occur no
more frequently than every Nrepeat*Nevery timesteps, and will
occur less frequently if \sphinxstyleemphasis{Nfreq} is not a multiple of
Nrepeat*Nevery.  There will be a warning message showing the new
settings.  Having a \sphinxstyleemphasis{Nfreq} setting that is larger than what is
required for correct computation of the ReaxFF force field
interactions, in combination with certain \sphinxstyleemphasis{Nrepeat} and \sphinxstyleemphasis{Nevery}
settings, can thus lead to incorrect results.  For typical ReaxFF
calculations, reneighboring only every 100 steps is already quite a
low frequency.
\end{sphinxadmonition}

\sphinxAtStartPar
If the filename ends with “.gz”, the output file is written in gzipped
format.  A gzipped dump file will be about 3x smaller than the text version,
but will also take longer to write.

\sphinxAtStartPar
\DUrole{versionmodified,added}{Added in version 15Jun2023: }Support for wildcards added

\sphinxAtStartPar
Optional keyword \sphinxstyleemphasis{cutoff} can be assigned to change the minimum
bond\sphinxhyphen{}order values used in identifying chemical bonds between pairs of
atoms.  Bond\sphinxhyphen{}order cutoffs should be carefully chosen, as bond\sphinxhyphen{}order
cutoffs that are too small may include too many bonds (which will result
in an error), while cutoffs that are too large will result in fragmented
molecules.  The default cutoff of 0.3 usually gives good results.  A
wildcard asterisk can be used in place of or in conjunction with the I,J
arguments to set the bond\sphinxhyphen{}order cutoff for multiple pairs of atom types.
This takes the form “*” or “*n” or “n*” or “m*n”.  If \(N\) is
the number of atom types, then an asterisk with no numeric values means
all types from 1 to \(N\).  A leading asterisk means all types from
1 to n (inclusive).  A trailing asterisk means all types from n to
\(N\) (inclusive).  A middle asterisk means all types from m to n
(inclusive).

\sphinxAtStartPar
The optional keyword \sphinxstyleemphasis{element} can be used to specify the chemical
symbol printed for each LAMMPS atom type. The number of symbols must
match the number of LAMMPS atom types and each symbol must consist of
1 or 2 alphanumeric characters. By default, these symbols are the same
as the chemical identity of each LAMMPS atom type, as specified by the
{\hyperref[\detokenize{pair_reaxff::doc}]{\sphinxcrossref{\DUrole{doc}{ReaxFF pair\_coeff}}}} command and the ReaxFF force
field file.

\sphinxAtStartPar
The optional keyword \sphinxstyleemphasis{position} writes center\sphinxhyphen{}of\sphinxhyphen{}mass positions of
each identified molecules to file \sphinxstyleemphasis{filepos} every \sphinxstyleemphasis{posfreq} timesteps.
The first line contains information on timestep, total number of
molecules, total number of distinct species, and box dimensions.  The
second line is a header containing labels.  From the third line
downward, each molecule writes a line of output containing the
following information: molecule ID, number of atoms in this molecule,
chemical formula, total charge, and center\sphinxhyphen{}of\sphinxhyphen{}mass xyz positions of
this molecule.  The xyz positions are in fractional coordinates
relative to the box dimensions.

\sphinxAtStartPar
For the keyword \sphinxstyleemphasis{position}, the \sphinxstyleemphasis{filepos} is the name of the output
file.  It can contain the wildcard character “*”.  If the “*”
character appears in \sphinxstyleemphasis{filepos}, then one file per snapshot is written
at \sphinxstyleemphasis{posfreq} and the “*” character is replaced with the timestep
value.  For example, AuO.pos.* becomes AuO.pos.0, AuO.pos.1000, etc.

\sphinxAtStartPar
\DUrole{versionmodified,added}{Added in version 3Aug2022.}

\sphinxAtStartPar
The optional keyword \sphinxstyleemphasis{delete} enables the periodic removal of molecules
from the system {\hyperref[\detokenize{fix_reaxff_species:delete}]{\sphinxcrossref{\DUrole{std,std-ref}{(Gissinger)}}}}.  Criteria for deletion can
be either a list of specific chemical formulae or a range of molecular
weights.  Molecules are deleted every \sphinxstyleemphasis{Nfreq} timesteps, and bond
connectivity is determined using the \sphinxstyleemphasis{Nevery} and \sphinxstyleemphasis{Nrepeat} keywords.  The
\sphinxstyleemphasis{filedel} argument is the name of the output file that records the species
that are removed from the system.  The \sphinxstyleemphasis{specieslist} keyword permits
specific chemical species to be deleted.  The \sphinxstyleemphasis{Nspecies} argument specifies
how many species are eligible for deletion and is followed by a list of
chemical formulae, whose strings are compared to species identified by this
fix.  For example, “specieslist 2 CO CO2” deletes molecules that are
identified as “CO” and “CO2” in the species output file.  When using the
\sphinxstyleemphasis{specieslist} keyword, the \sphinxstyleemphasis{filedel} file has the following format: the
first line lists the chemical formulae eligible for deletion, and each
additional line contains the timestep on which a molecule deletion occurs
and the number of each species deleted on that timestep.  The \sphinxstyleemphasis{masslimit}
keyword permits deletion of molecules with molecular weights between
\sphinxstyleemphasis{massmin} and \sphinxstyleemphasis{massmax}.  When using the \sphinxstyleemphasis{masslimit} keyword, each line of
the \sphinxstyleemphasis{filedel} file contains the timestep on which deletions occurs,
followed by how many of each species are deleted (with quantities preceding
chemical formulae).  The \sphinxstyleemphasis{specieslist} and \sphinxstyleemphasis{masslimit} keywords cannot both
be used in the same \sphinxstyleemphasis{reaxff/species} fix.  The \sphinxstyleemphasis{delete\_rate\_limit} keyword
can enforce an upper limit on the overall rate of molecule deletion.  The
number of deletion occurrences is limited to Nlimit within an interval of
Nsteps timesteps.  Nlimit can be specified with an equal\sphinxhyphen{}style
{\hyperref[\detokenize{variable::doc}]{\sphinxcrossref{\DUrole{doc}{variable}}}}.  When using the \sphinxstyleemphasis{delete\_rate\_limit} keyword, no
deletions are permitted to occur within the first Nsteps timesteps of the
first run (after reading a either a data or restart file).


\bigskip\hrule\bigskip


\sphinxAtStartPar
The \sphinxstyleemphasis{Nevery}, \sphinxstyleemphasis{Nrepeat}, and \sphinxstyleemphasis{Nfreq} arguments specify on what
timesteps the bond\sphinxhyphen{}order values are sampled to get the average bond
order.  The species analysis is performed using the average bond\sphinxhyphen{}order
on timesteps that are a multiple of \sphinxstyleemphasis{Nfreq}.  The average is over
\sphinxstyleemphasis{Nrepeat} bond\sphinxhyphen{}order samples, computed in the preceding portion of the
simulation every \sphinxstyleemphasis{Nevery} timesteps.  \sphinxstyleemphasis{Nfreq} must be a multiple of
\sphinxstyleemphasis{Nevery} and \sphinxstyleemphasis{Nevery} must be non\sphinxhyphen{}zero even if \sphinxstyleemphasis{Nrepeat} is 1.
Also, the timesteps
contributing to the average bond\sphinxhyphen{}order cannot overlap,
i.e. Nrepeat*Nevery can not exceed Nfreq.

\sphinxAtStartPar
For example, if Nevery=2, Nrepeat=6, and Nfreq=100, then bond\sphinxhyphen{}order
values on timesteps 90,92,94,96,98,100 will be used to compute the
average bond\sphinxhyphen{}order for the species analysis output on timestep 100.


\bigskip\hrule\bigskip



\subsection{Restart, fix\_modify, output, run start/stop, minimize info}
\label{\detokenize{fix_reaxff_species:restart-fix-modify-output-run-start-stop-minimize-info}}
\sphinxAtStartPar
No information about this fix is written to {\hyperref[\detokenize{restart::doc}]{\sphinxcrossref{\DUrole{doc}{binary restart files}}}}.
None of the {\hyperref[\detokenize{fix_modify::doc}]{\sphinxcrossref{\DUrole{doc}{fix\_modify}}}} options are relevant to this fix.

\sphinxAtStartPar
This fix computes both a global vector of length 2 and a per\sphinxhyphen{}atom vector,
either of which can be accessed by various {\hyperref[\detokenize{Howto_output::doc}]{\sphinxcrossref{\DUrole{doc}{output commands}}}}.
The values in the global vector are “intensive”.

\sphinxAtStartPar
The 2 values in the global vector are as follows:
\begin{enumerate}
\sphinxsetlistlabels{\arabic}{enumi}{enumii}{}{.}%
\item {} 
\sphinxAtStartPar
total number of molecules

\item {} 
\sphinxAtStartPar
total number of distinct species

\end{enumerate}

\sphinxAtStartPar
The per\sphinxhyphen{}atom vector stores the molecule ID for each atom as identified
by the fix.  If an atom is not in a molecule, its ID will be 0.
For atoms in the same molecule, the molecule ID for all of them
will be the same and will be equal to the smallest atom ID of
any atom in the molecule.

\sphinxAtStartPar
No parameter of this fix can be used with the \sphinxstyleemphasis{start/stop} keywords of
the {\hyperref[\detokenize{run::doc}]{\sphinxcrossref{\DUrole{doc}{run}}}} command.
This fix is not invoked during {\hyperref[\detokenize{minimize::doc}]{\sphinxcrossref{\DUrole{doc}{energy minimization}}}}.


\bigskip\hrule\bigskip


\sphinxAtStartPar
Styles with a \sphinxstyleemphasis{gpu}, \sphinxstyleemphasis{intel}, \sphinxstyleemphasis{kk}, \sphinxstyleemphasis{omp}, or \sphinxstyleemphasis{opt} suffix are
functionally the same as the corresponding style without the suffix.
They have been optimized to run faster, depending on your available
hardware, as discussed on the {\hyperref[\detokenize{Speed_packages::doc}]{\sphinxcrossref{\DUrole{doc}{Accelerator packages}}}}
page.  The accelerated styles take the same arguments and should
produce the same results, except for round\sphinxhyphen{}off and precision issues.

\sphinxAtStartPar
These accelerated styles are part of the GPU, INTEL, KOKKOS,
OPENMP, and OPT packages, respectively.  They are only enabled if
LAMMPS was built with those packages.  See the {\hyperref[\detokenize{Build_package::doc}]{\sphinxcrossref{\DUrole{doc}{Build package}}}} page for more info.

\sphinxAtStartPar
You can specify the accelerated styles explicitly in your input script
by including their suffix, or you can use the {\hyperref[\detokenize{Run_options::doc}]{\sphinxcrossref{\DUrole{doc}{\sphinxhyphen{}suffix command\sphinxhyphen{}line switch}}}} when you invoke LAMMPS, or you can use the
{\hyperref[\detokenize{suffix::doc}]{\sphinxcrossref{\DUrole{doc}{suffix}}}} command in your input script.

\sphinxAtStartPar
See the {\hyperref[\detokenize{Speed_packages::doc}]{\sphinxcrossref{\DUrole{doc}{Accelerator packages}}}} page for more
instructions on how to use the accelerated styles effectively.


\bigskip\hrule\bigskip



\subsection{Restrictions}
\label{\detokenize{fix_reaxff_species:restrictions}}
\sphinxAtStartPar
The “fix reaxff/species” requires that {\hyperref[\detokenize{pair_reaxff::doc}]{\sphinxcrossref{\DUrole{doc}{pair\_style reaxff}}}} is used.
This fix is part of the REAXFF package.  It is only enabled if LAMMPS was built with that
package.  See the {\hyperref[\detokenize{Build_package::doc}]{\sphinxcrossref{\DUrole{doc}{Build package}}}} page for more info.

\sphinxAtStartPar
To write gzipped species files, you must compile LAMMPS with the \sphinxhyphen{}DLAMMPS\_GZIP option.


\subsection{Related commands}
\label{\detokenize{fix_reaxff_species:related-commands}}
\sphinxAtStartPar
{\hyperref[\detokenize{pair_reaxff::doc}]{\sphinxcrossref{\DUrole{doc}{pair\_style reaxff}}}}, {\hyperref[\detokenize{fix_reaxff_bonds::doc}]{\sphinxcrossref{\DUrole{doc}{fix reaxff/bonds}}}}


\subsection{Default}
\label{\detokenize{fix_reaxff_species:default}}
\sphinxAtStartPar
The default values for bond\sphinxhyphen{}order cutoffs are 0.3 for all I\sphinxhyphen{}J pairs.
The default element symbols are taken from the ReaxFF pair\_coeff command.
Position files are not written by default.

\phantomsection\label{\detokenize{fix_reaxff_species:delete}}
\sphinxAtStartPar
\sphinxstylestrong{(Gissinger)} Jacob R. Gissinger, Scott R. Zavada, Joseph G. Smith, Josh Kemppainen, Ivan Gallegos, Gregory M. Odegard, Emilie J. Siochi, and Kristopher E. Wise, Carbon, 202, 336\sphinxhyphen{}347 (2023).

\sphinxstepscope

\index{fix recenter@\spxentry{fix recenter}}\ignorespaces 

\section{fix recenter command}
\label{\detokenize{fix_recenter:fix-recenter-command}}\label{\detokenize{fix_recenter:index-0}}\label{\detokenize{fix_recenter::doc}}

\subsection{Syntax}
\label{\detokenize{fix_recenter:syntax}}
\begin{sphinxVerbatim}[commandchars=\\\{\}]
\PYG{k}{fix }\PYG{n+nv+nvIdentifier}{ID}\PYG{+w}{ }\PYG{n+nv+nvIdentifier}{group\PYGZhy{}ID}\PYG{+w}{ }\PYG{n}{recenter}\PYG{+w}{ }\PYG{n}{x}\PYG{+w}{ }\PYG{n}{y}\PYG{+w}{ }\PYG{n}{z}\PYG{+w}{ }\PYG{n}{keyword}\PYG{+w}{ }\PYG{n}{value}\PYG{+w}{ }\PYG{n}{...}
\end{sphinxVerbatim}
\begin{itemize}
\item {} 
\sphinxAtStartPar
ID, group\sphinxhyphen{}ID are documented in {\hyperref[\detokenize{fix::doc}]{\sphinxcrossref{\DUrole{doc}{fix}}}} command

\item {} 
\sphinxAtStartPar
recenter = style name of this fix command

\item {} 
\sphinxAtStartPar
x,y,z = constrain center\sphinxhyphen{}of\sphinxhyphen{}mass to these coords (distance units),         any coord can also be NULL or INIT (see below)

\item {} 
\sphinxAtStartPar
zero or more keyword/value pairs may be appended

\item {} 
\sphinxAtStartPar
keyword = \sphinxstyleemphasis{shift} or \sphinxstyleemphasis{units}
\begin{sphinxalltt}
\sphinxstyleemphasis{shift} value = group\sphinxhyphen{}ID
  group\sphinxhyphen{}ID = group of atoms whose coords are shifted
\sphinxstyleemphasis{units} value = \sphinxstyleemphasis{box} or \sphinxstyleemphasis{lattice} or \sphinxstyleemphasis{fraction}
\end{sphinxalltt}

\end{itemize}


\subsection{Examples}
\label{\detokenize{fix_recenter:examples}}
\begin{sphinxVerbatim}[commandchars=\\\{\}]
\PYG{k}{fix }\PYG{n+nv+nvIdentifier}{1}\PYG{+w}{ }\PYG{n+nv+nvIdentifier}{all}\PYG{+w}{ }\PYG{n}{recenter}\PYG{+w}{ }\PYG{l+m}{0.0}\PYG{+w}{ }\PYG{l+m}{0.5}\PYG{+w}{ }\PYG{l+m}{0.0}
\PYG{k}{fix }\PYG{n+nv+nvIdentifier}{1}\PYG{+w}{ }\PYG{n+nv+nvIdentifier}{all}\PYG{+w}{ }\PYG{n}{recenter}\PYG{+w}{ }\PYG{n}{INIT}\PYG{+w}{ }\PYG{n}{INIT}\PYG{+w}{ }\PYG{n}{NULL}
\PYG{k}{fix }\PYG{n+nv+nvIdentifier}{1}\PYG{+w}{ }\PYG{n+nv+nvIdentifier}{all}\PYG{+w}{ }\PYG{n}{recenter}\PYG{+w}{ }\PYG{n}{INIT}\PYG{+w}{ }\PYG{l+m}{0.0}\PYG{+w}{ }\PYG{l+m}{0.0}\PYG{+w}{ }\PYG{n}{units}\PYG{+w}{ }\PYG{n}{box}
\end{sphinxVerbatim}


\subsection{Description}
\label{\detokenize{fix_recenter:description}}
\sphinxAtStartPar
Constrain the center\sphinxhyphen{}of\sphinxhyphen{}mass position of a group of atoms by adjusting
the coordinates of the atoms every timestep.  This is simply a small
shift that does not alter the dynamics of the system or change the
relative coordinates of any pair of atoms in the group.  This can be
used to ensure the entire collection of atoms (or a portion of them)
do not drift during the simulation due to random perturbations
(e.g. {\hyperref[\detokenize{fix_langevin::doc}]{\sphinxcrossref{\DUrole{doc}{fix langevin}}}} thermostatting).

\sphinxAtStartPar
Distance units for the x,y,z values are determined by the setting of
the \sphinxstyleemphasis{units} keyword, as discussed below.  One or more x,y,z values can
also be specified as NULL, which means exclude that dimension from
this operation.  Or it can be specified as INIT which means to
constrain the center\sphinxhyphen{}of\sphinxhyphen{}mass to its initial value at the beginning of
the run.

\sphinxAtStartPar
The center\sphinxhyphen{}of\sphinxhyphen{}mass (COM) is computed for the group specified by the
fix.  If the current COM is different than the specified x,y,z, then a
group of atoms has their coordinates shifted by the difference.  By
default the shifted group is also the group specified by the fix.  A
different group can be shifted by using the \sphinxstyleemphasis{shift} keyword.  For
example, the COM could be computed on a protein to keep it in the
center of the simulation box.  But the entire system (protein + water)
could be shifted.

\sphinxAtStartPar
If the \sphinxstyleemphasis{units} keyword is set to \sphinxstyleemphasis{box}, then the distance units of
x,y,z are defined by the {\hyperref[\detokenize{units::doc}]{\sphinxcrossref{\DUrole{doc}{units}}}} command \sphinxhyphen{} e.g. Angstroms
for \sphinxstyleemphasis{real} units.  A \sphinxstyleemphasis{lattice} value means the distance units are in
lattice spacings.  The {\hyperref[\detokenize{lattice::doc}]{\sphinxcrossref{\DUrole{doc}{lattice}}}} command must have been
previously used to define the lattice spacing.  A \sphinxstyleemphasis{fraction} value
means a fractional distance between the lo/hi box boundaries, e.g. 0.5
= middle of the box.  The default is to use lattice units.

\sphinxAtStartPar
Note that the {\hyperref[\detokenize{velocity::doc}]{\sphinxcrossref{\DUrole{doc}{velocity}}}} command can be used to create
velocities with zero aggregate linear and/or angular momentum.

\begin{sphinxadmonition}{note}{Note:}
\sphinxAtStartPar
This fix performs its operations at the same point in the
timestep as other time integration fixes, such as {\hyperref[\detokenize{fix_nve::doc}]{\sphinxcrossref{\DUrole{doc}{fix nve}}}}, {\hyperref[\detokenize{fix_nh::doc}]{\sphinxcrossref{\DUrole{doc}{fix nvt}}}}, or {\hyperref[\detokenize{fix_nh::doc}]{\sphinxcrossref{\DUrole{doc}{fix npt}}}}.
Thus fix recenter should normally be the last such fix specified in
the input script, since the adjustments it makes to atom coordinates
should come after the changes made by time integration.  LAMMPS will
warn you if your fixes are not ordered this way.
\end{sphinxadmonition}

\begin{sphinxadmonition}{note}{Note:}
\sphinxAtStartPar
If you use this fix on a small group of atoms (e.g. a molecule
in solvent) without using the \sphinxstyleemphasis{shift} keyword to adjust the positions
of all atoms in the system, then the results can be unpredictable.
For example, if the molecule is pushed consistently in one direction
by a flowing solvent, its velocity will increase.  But its coordinates
will be re\sphinxhyphen{}centered, meaning it is moved back towards the force.  Thus
over time, the velocity and effective temperature of the molecule
could become very large, though it won’t actually be moving due to the
re\sphinxhyphen{}centering.  If you are thermostatting the entire system, then the
solvent would be cooled to compensate.  A better solution for this
simulation scenario is to use the {\hyperref[\detokenize{fix_spring::doc}]{\sphinxcrossref{\DUrole{doc}{fix spring}}}} command
to tether the molecule in place.
\end{sphinxadmonition}


\subsection{Restart, fix\_modify, output, run start/stop, minimize info}
\label{\detokenize{fix_recenter:restart-fix-modify-output-run-start-stop-minimize-info}}
\sphinxAtStartPar
No information about this fix is written to {\hyperref[\detokenize{restart::doc}]{\sphinxcrossref{\DUrole{doc}{binary restart files}}}}.  None of the {\hyperref[\detokenize{fix_modify::doc}]{\sphinxcrossref{\DUrole{doc}{fix\_modify}}}} options
are relevant to this fix.

\sphinxAtStartPar
This fix computes a global scalar which can be accessed by various
{\hyperref[\detokenize{Howto_output::doc}]{\sphinxcrossref{\DUrole{doc}{output commands}}}}.  The scalar is the distance the
group is moved by fix recenter.

\sphinxAtStartPar
This fix also computes global 3\sphinxhyphen{}vector which can be accessed by
various {\hyperref[\detokenize{Howto_output::doc}]{\sphinxcrossref{\DUrole{doc}{output commands}}}}.  The 3 quantities in the
vector are xyz components of displacement applied to the group of
atoms by the fix.

\sphinxAtStartPar
The scalar and vector values calculated by this fix are “extensive”.

\sphinxAtStartPar
No parameter of this fix can be used with the \sphinxstyleemphasis{start/stop} keywords of
the {\hyperref[\detokenize{run::doc}]{\sphinxcrossref{\DUrole{doc}{run}}}} command.  This fix is not invoked during {\hyperref[\detokenize{minimize::doc}]{\sphinxcrossref{\DUrole{doc}{energy minimization}}}}.


\subsection{Restrictions}
\label{\detokenize{fix_recenter:restrictions}}
\sphinxAtStartPar
This fix should not be used with an x,y,z setting that causes a large
shift in the system on the first timestep, due to the requested COM
being very different from the initial COM.  This could cause atoms to
be lost, especially in parallel.  Instead, use the
{\hyperref[\detokenize{displace_atoms::doc}]{\sphinxcrossref{\DUrole{doc}{displace\_atoms}}}} command, which can be used to
move atoms a large distance.


\subsection{Related commands}
\label{\detokenize{fix_recenter:related-commands}}
\sphinxAtStartPar
{\hyperref[\detokenize{fix_momentum::doc}]{\sphinxcrossref{\DUrole{doc}{fix momentum}}}}, {\hyperref[\detokenize{velocity::doc}]{\sphinxcrossref{\DUrole{doc}{velocity}}}}


\subsection{Default}
\label{\detokenize{fix_recenter:default}}
\sphinxAtStartPar
The option defaults are shift = fix group\sphinxhyphen{}ID, and units = lattice.

\sphinxstepscope

\index{fix restrain@\spxentry{fix restrain}}\ignorespaces 

\section{fix restrain command}
\label{\detokenize{fix_restrain:fix-restrain-command}}\label{\detokenize{fix_restrain:index-0}}\label{\detokenize{fix_restrain::doc}}

\subsection{Syntax}
\label{\detokenize{fix_restrain:syntax}}
\begin{sphinxVerbatim}[commandchars=\\\{\}]
\PYG{k}{fix }\PYG{n+nv+nvIdentifier}{ID}\PYG{+w}{ }\PYG{n+nv+nvIdentifier}{group\PYGZhy{}ID}\PYG{+w}{ }\PYG{n}{restrain}\PYG{+w}{ }\PYG{n}{keyword}\PYG{+w}{ }\PYG{n}{args}\PYG{+w}{ }\PYG{n}{...}
\end{sphinxVerbatim}
\begin{itemize}
\item {} 
\sphinxAtStartPar
ID, group\sphinxhyphen{}ID are documented in {\hyperref[\detokenize{fix::doc}]{\sphinxcrossref{\DUrole{doc}{fix}}}} command

\item {} 
\sphinxAtStartPar
restrain = style name of this fix command

\item {} 
\sphinxAtStartPar
one or more keyword/arg pairs may be appended

\item {} 
\sphinxAtStartPar
keyword = \sphinxstyleemphasis{bond} or \sphinxstyleemphasis{lbound} or \sphinxstyleemphasis{angle} or \sphinxstyleemphasis{dihedral}
\begin{sphinxalltt}
\sphinxstyleemphasis{bond} args = atom1 atom2 Kstart Kstop r0start (r0stop)
  atom1,atom2 = IDs of two atoms in bond
  Kstart,Kstop = restraint coefficients at start/end of run (energy units)
  r0start = equilibrium bond distance at start of run (distance units)
  r0stop = equilibrium bond distance at end of run (optional) (distance units). If not
    specified it is assumed to be equal to r0start
\sphinxstyleemphasis{lbound} args = atom1 atom2 Kstart Kstop r0start (r0stop)
  atom1,atom2 = IDs of two atoms in bond
  Kstart,Kstop = restraint coefficients at start/end of run (energy units)
  r0start = equilibrium bond distance at start of run (distance units)
  r0stop = equilibrium bond distance at end of run (optional) (distance units). If not
    specified it is assumed to be equal to r0start
\sphinxstyleemphasis{angle} args = atom1 atom2 atom3 Kstart Kstop theta0
  atom1,atom2,atom3 = IDs of three atoms in angle, atom2 = middle atom
  Kstart,Kstop = restraint coefficients at start/end of run (energy units)
  theta0 = equilibrium angle theta (degrees)
\sphinxstyleemphasis{dihedral} args = atom1 atom2 atom3 atom4 Kstart Kstop phi0 keyword/value
  atom1,atom2,atom3,atom4 = IDs of 4 atoms in dihedral in linear order
  Kstart,Kstop = restraint coefficients at start/end of run (energy units)
  phi0 = equilibrium dihedral angle phi (degrees)
  keyword/value = optional keyword value pairs. supported keyword/value pairs:
    \sphinxstyleemphasis{mult} n = dihedral multiplicity n (integer \textgreater{}= 0, default = 1)
\end{sphinxalltt}

\end{itemize}


\subsection{Examples}
\label{\detokenize{fix_restrain:examples}}
\begin{sphinxVerbatim}[commandchars=\\\{\}]
\PYG{k}{fix }\PYG{n+nv+nvIdentifier}{holdem}\PYG{+w}{ }\PYG{n+nv+nvIdentifier}{all}\PYG{+w}{ }\PYG{n}{restrain}\PYG{+w}{ }\PYG{n}{bond}\PYG{+w}{ }\PYG{l+m}{45}\PYG{+w}{ }\PYG{l+m}{48}\PYG{+w}{ }\PYG{l+m}{2000.0}\PYG{+w}{ }\PYG{l+m}{2000.0}\PYG{+w}{ }\PYG{l+m}{2.75}
\PYG{k}{fix }\PYG{n+nv+nvIdentifier}{holdem}\PYG{+w}{ }\PYG{n+nv+nvIdentifier}{all}\PYG{+w}{ }\PYG{n}{restrain}\PYG{+w}{ }\PYG{n}{lbound}\PYG{+w}{ }\PYG{l+m}{45}\PYG{+w}{ }\PYG{l+m}{48}\PYG{+w}{ }\PYG{l+m}{2000.0}\PYG{+w}{ }\PYG{l+m}{2000.0}\PYG{+w}{ }\PYG{l+m}{2.75}
\PYG{k}{fix }\PYG{n+nv+nvIdentifier}{holdem}\PYG{+w}{ }\PYG{n+nv+nvIdentifier}{all}\PYG{+w}{ }\PYG{n}{restrain}\PYG{+w}{ }\PYG{n}{dihedral}\PYG{+w}{ }\PYG{l+m}{1}\PYG{+w}{ }\PYG{l+m}{2}\PYG{+w}{ }\PYG{l+m}{3}\PYG{+w}{ }\PYG{l+m}{4}\PYG{+w}{ }\PYG{l+m}{2000.0}\PYG{+w}{ }\PYG{l+m}{2000.0}\PYG{+w}{ }\PYG{l+m}{120.0}
\PYG{k}{fix }\PYG{n+nv+nvIdentifier}{holdem}\PYG{+w}{ }\PYG{n+nv+nvIdentifier}{all}\PYG{+w}{ }\PYG{n}{restrain}\PYG{+w}{ }\PYG{n}{bond}\PYG{+w}{ }\PYG{l+m}{45}\PYG{+w}{ }\PYG{l+m}{48}\PYG{+w}{ }\PYG{l+m}{2000.0}\PYG{+w}{ }\PYG{l+m}{2000.0}\PYG{+w}{ }\PYG{l+m}{2.75}\PYG{+w}{ }\PYG{n}{dihedral}\PYG{+w}{ }\PYG{l+m}{1}\PYG{+w}{ }\PYG{l+m}{2}\PYG{+w}{ }\PYG{l+m}{3}\PYG{+w}{ }\PYG{l+m}{4}\PYG{+w}{ }\PYG{l+m}{2000.0}\PYG{+w}{ }\PYG{l+m}{2000.0}\PYG{+w}{ }\PYG{l+m}{120.0}
\PYG{k}{fix }\PYG{n+nv+nvIdentifier}{texas\PYGZus{}holdem}\PYG{+w}{ }\PYG{n+nv+nvIdentifier}{all}\PYG{+w}{ }\PYG{n}{restrain}\PYG{+w}{ }\PYG{n}{dihedral}\PYG{+w}{ }\PYG{l+m}{1}\PYG{+w}{ }\PYG{l+m}{2}\PYG{+w}{ }\PYG{l+m}{3}\PYG{+w}{ }\PYG{l+m}{4}\PYG{+w}{ }\PYG{l+m}{0.0}\PYG{+w}{ }\PYG{l+m}{2000.0}\PYG{+w}{ }\PYG{l+m}{120.0}\PYG{+w}{ }\PYG{n}{dihedral}\PYG{+w}{ }\PYG{l+m}{1}\PYG{+w}{ }\PYG{l+m}{2}\PYG{+w}{ }\PYG{l+m}{3}\PYG{+w}{ }\PYG{l+m}{5}\PYG{+w}{ }\PYG{l+m}{0.0}\PYG{+w}{ }\PYG{l+m}{2000.0}\PYG{+w}{ }\PYG{o}{\PYGZhy{}}\PYG{l+m}{120.0}\PYG{+w}{ }\PYG{n}{dihedral}\PYG{+w}{ }\PYG{l+m}{1}\PYG{+w}{ }\PYG{l+m}{2}\PYG{+w}{ }\PYG{l+m}{3}\PYG{+w}{ }\PYG{l+m}{6}\PYG{+w}{ }\PYG{l+m}{0.0}\PYG{+w}{ }\PYG{l+m}{2000.0}\PYG{+w}{ }\PYG{l+m}{0.0}
\end{sphinxVerbatim}


\subsection{Description}
\label{\detokenize{fix_restrain:description}}
\sphinxAtStartPar
Restrain the motion of the specified sets of atoms by making them part
of a bond or angle or dihedral interaction whose strength can vary
over time during a simulation.  This is functionally similar to
creating a bond or angle or dihedral for the same atoms in a data
file, as specified by the {\hyperref[\detokenize{read_data::doc}]{\sphinxcrossref{\DUrole{doc}{read\_data}}}} command, albeit
with a time\sphinxhyphen{}varying prefactor coefficient, and except for exclusion
rules, as explained below.

\sphinxAtStartPar
For the purpose of force field parameter\sphinxhyphen{}fitting or mapping a molecular
potential energy surface, this fix reduces the hassle and risk
associated with modifying data files.  In other words, use this fix to
temporarily force a molecule to adopt a particular conformation.  To
create a permanent bond or angle or dihedral, you should modify the
data file.

\begin{sphinxadmonition}{note}{Note:}
\sphinxAtStartPar
Adding a bond/angle/dihedral with this command does not apply
the exclusion rules and weighting factors specified by the
{\hyperref[\detokenize{special_bonds::doc}]{\sphinxcrossref{\DUrole{doc}{special\_bonds}}}} command to atoms in the restraint
that are now bonded (1\sphinxhyphen{}2,1\sphinxhyphen{}3,1\sphinxhyphen{}4 neighbors) as a result.  If they are
close enough to interact in a {\hyperref[\detokenize{pair_style::doc}]{\sphinxcrossref{\DUrole{doc}{pair\_style}}}} sense
(non\sphinxhyphen{}bonded interaction), then the bond/angle/dihedral restraint
interaction will simply be superposed on top of that interaction.
\end{sphinxadmonition}

\sphinxAtStartPar
The group\sphinxhyphen{}ID specified by this fix is ignored.

\sphinxAtStartPar
The second example above applies a restraint to hold the dihedral
angle formed by atoms 1, 2, 3, and 4 near 120 degrees using a constant
restraint coefficient.  The fourth example applies similar restraints
to multiple dihedral angles using a restraint coefficient that
increases from 0.0 to 2000.0 over the course of the run.

\begin{sphinxadmonition}{note}{Note:}
\sphinxAtStartPar
Adding a force to atoms implies a change in their potential
energy as they move due to the applied force field.  For dynamics via
the {\hyperref[\detokenize{run::doc}]{\sphinxcrossref{\DUrole{doc}{run}}}} command, this energy can be added to the system’s
potential energy for thermodynamic output (see below).  For energy
minimization via the {\hyperref[\detokenize{minimize::doc}]{\sphinxcrossref{\DUrole{doc}{minimize}}}} command, this energy
must be added to the system’s potential energy to formulate a
self\sphinxhyphen{}consistent minimization problem (see below).
\end{sphinxadmonition}

\sphinxAtStartPar
In order for a restraint to be effective, the restraint force must
typically be significantly larger than the forces associated with
conventional force field terms.  If the restraint is applied during a
dynamics run (as opposed to during an energy minimization), a large
restraint coefficient can significantly reduce the stable timestep
size, especially if the atoms are initially far from the preferred
conformation.  You may need to experiment to determine what value of \(K\)
works best for a given application.

\sphinxAtStartPar
For the case of finding a minimum energy structure for a single
molecule with particular restraints (e.g. for fitting force field
parameters or constructing a potential energy surface), commands such
as the following may be useful:

\begin{sphinxVerbatim}[commandchars=\\\{\}]
\PYG{c}{\PYGZsh{} minimize molecule energy with restraints}
\PYG{k}{velocity }\PYG{n+nv+nvIdentifier}{all}\PYG{+w}{ }\PYG{n}{create}\PYG{+w}{ }\PYG{l+m}{600.0}\PYG{+w}{ }\PYG{l+m}{8675309}\PYG{+w}{ }\PYG{n}{mom}\PYG{+w}{ }\PYG{n}{yes}\PYG{+w}{ }\PYG{n}{rot}\PYG{+w}{ }\PYG{n}{yes}\PYG{+w}{ }\PYG{n}{dist}\PYG{+w}{ }\PYG{n}{gaussian}
\PYG{k}{fix }\PYG{n+nv+nvIdentifier}{NVE}\PYG{+w}{ }\PYG{n+nv+nvIdentifier}{all}\PYG{+w}{ }\PYG{n}{nve}
\PYG{k}{fix }\PYG{n+nv+nvIdentifier}{TFIX}\PYG{+w}{ }\PYG{n+nv+nvIdentifier}{all}\PYG{+w}{ }\PYG{n}{langevin}\PYG{+w}{ }\PYG{l+m}{600.0}\PYG{+w}{ }\PYG{l+m}{0.0}\PYG{+w}{ }\PYG{l+m}{100}\PYG{+w}{ }\PYG{l+m}{24601}
\PYG{k}{fix }\PYG{n+nv+nvIdentifier}{REST}\PYG{+w}{ }\PYG{n+nv+nvIdentifier}{all}\PYG{+w}{ }\PYG{n}{restrain}\PYG{+w}{ }\PYG{n}{dihedral}\PYG{+w}{ }\PYG{l+m}{2}\PYG{+w}{ }\PYG{l+m}{1}\PYG{+w}{ }\PYG{l+m}{3}\PYG{+w}{ }\PYG{l+m}{8}\PYG{+w}{ }\PYG{l+m}{0.0}\PYG{+w}{ }\PYG{l+m}{5000.0}\PYG{+w}{ }\PYG{n+nv}{\PYGZdl{}\PYGZob{}}\PYG{n+nv}{angle1}\PYG{n+nv}{\PYGZcb{}}\PYG{+w}{ }\PYG{n}{dihedral}\PYG{+w}{ }\PYG{l+m}{3}\PYG{+w}{ }\PYG{l+m}{1}\PYG{+w}{ }\PYG{l+m}{2}\PYG{+w}{ }\PYG{l+m}{9}\PYG{+w}{ }\PYG{l+m}{0.0}\PYG{+w}{ }\PYG{l+m}{5000.0}\PYG{+w}{ }\PYG{n+nv}{\PYGZdl{}\PYGZob{}}\PYG{n+nv}{angle2}\PYG{n+nv}{\PYGZcb{}}
\PYG{k}{fix\PYGZus{}modify }\PYG{n+nv+nvIdentifier}{REST}\PYG{+w}{ }\PYG{n}{energy}\PYG{+w}{ }\PYG{n}{yes}
\PYG{k}{run}\PYG{+w}{ }\PYG{l+m}{10000}
\PYG{k}{fix }\PYG{n+nv+nvIdentifier}{TFIX}\PYG{+w}{ }\PYG{n+nv+nvIdentifier}{all}\PYG{+w}{ }\PYG{n}{langevin}\PYG{+w}{ }\PYG{l+m}{0.0}\PYG{+w}{ }\PYG{l+m}{0.0}\PYG{+w}{ }\PYG{l+m}{100}\PYG{+w}{ }\PYG{l+m}{24601}
\PYG{k}{fix }\PYG{n+nv+nvIdentifier}{REST}\PYG{+w}{ }\PYG{n+nv+nvIdentifier}{all}\PYG{+w}{ }\PYG{n}{restrain}\PYG{+w}{ }\PYG{n}{dihedral}\PYG{+w}{ }\PYG{l+m}{2}\PYG{+w}{ }\PYG{l+m}{1}\PYG{+w}{ }\PYG{l+m}{3}\PYG{+w}{ }\PYG{l+m}{8}\PYG{+w}{ }\PYG{l+m}{5000.0}\PYG{+w}{ }\PYG{l+m}{5000.0}\PYG{+w}{ }\PYG{n+nv}{\PYGZdl{}\PYGZob{}}\PYG{n+nv}{angle1}\PYG{n+nv}{\PYGZcb{}}\PYG{+w}{ }\PYG{n}{dihedral}\PYG{+w}{ }\PYG{l+m}{3}\PYG{+w}{ }\PYG{l+m}{1}\PYG{+w}{ }\PYG{l+m}{2}\PYG{+w}{ }\PYG{l+m}{9}\PYG{+w}{ }\PYG{l+m}{5000.0}\PYG{+w}{ }\PYG{l+m}{5000.0}\PYG{+w}{ }\PYG{n+nv}{\PYGZdl{}\PYGZob{}}\PYG{n+nv}{angle2}\PYG{n+nv}{\PYGZcb{}}
\PYG{k}{fix\PYGZus{}modify }\PYG{n+nv+nvIdentifier}{REST}\PYG{+w}{ }\PYG{n}{energy}\PYG{+w}{ }\PYG{n}{yes}
\PYG{k}{run}\PYG{+w}{ }\PYG{l+m}{10000}
\PYG{c}{\PYGZsh{} sanity check for convergence}
\PYG{k}{minimize}\PYG{+w}{ }\PYG{l+m}{1e\PYGZhy{}6}\PYG{+w}{ }\PYG{l+m}{1e\PYGZhy{}9}\PYG{+w}{ }\PYG{l+m}{1000}\PYG{+w}{ }\PYG{l+m}{100000}
\PYG{c}{\PYGZsh{} report unrestrained energies}
\PYG{k}{unfix }\PYG{n+nv+nvIdentifier}{REST}
\PYG{k}{run}\PYG{+w}{ }\PYG{l+m}{0}
\end{sphinxVerbatim}


\bigskip\hrule\bigskip


\sphinxAtStartPar
The \sphinxstyleemphasis{bond} keyword applies a bond restraint to the specified atoms
using the same functional form used by the {\hyperref[\detokenize{bond_harmonic::doc}]{\sphinxcrossref{\DUrole{doc}{bond\_style harmonic}}}} command.  The potential associated with
the restraint is
\begin{equation*}
\begin{split}E = K (r - r_0)^2\end{split}
\end{equation*}
\sphinxAtStartPar
with the following coefficients:
\begin{itemize}
\item {} 
\sphinxAtStartPar
\(K\) (energy/distance\textasciicircum{}2)

\item {} 
\sphinxAtStartPar
\(r_0\) (distance)

\end{itemize}

\sphinxAtStartPar
\(K\) and \(r_0\) are specified with the fix.  Note that the usual 1/2 factor
is included in \(K\).


\bigskip\hrule\bigskip


\sphinxAtStartPar
The \sphinxstyleemphasis{lbound} keyword applies a lower bound bond restraint to the specified atoms
using the same functional form used by the {\hyperref[\detokenize{bond_harmonic::doc}]{\sphinxcrossref{\DUrole{doc}{bond\_style harmonic}}}} command if the distance between
the atoms is smaller than the equilibrium bond distance and 0 otherwise. The potential associated with
the restraint is
\begin{equation*}
\begin{split}E = K (r - r_0)^2 ,if\;r < r_0\end{split}
\end{equation*}\begin{equation*}
\begin{split}E = 0 \qquad\quad\quad ,if\;r \ge r_0\end{split}
\end{equation*}
\sphinxAtStartPar
with the following coefficients:
\begin{itemize}
\item {} 
\sphinxAtStartPar
\(K\) (energy/distance\textasciicircum{}2)

\item {} 
\sphinxAtStartPar
\(r_0\) (distance)

\end{itemize}

\sphinxAtStartPar
\(K\) and \(r_0\) are specified with the fix.  Note that the usual 1/2 factor
is included in \(K\).


\bigskip\hrule\bigskip


\sphinxAtStartPar
The \sphinxstyleemphasis{angle} keyword applies an angle restraint to the specified atoms
using the same functional form used by the {\hyperref[\detokenize{angle_harmonic::doc}]{\sphinxcrossref{\DUrole{doc}{angle\_style harmonic}}}} command.  The potential associated with
the restraint is
\begin{equation*}
\begin{split}E = K (\theta - \theta_0)^2\end{split}
\end{equation*}
\sphinxAtStartPar
with the following coefficients:
\begin{itemize}
\item {} 
\sphinxAtStartPar
\(K\) (energy)

\item {} 
\sphinxAtStartPar
\(\theta_0\) (degrees)

\end{itemize}

\sphinxAtStartPar
\(K\) and \(\theta_0\) are specified with the fix.
\(\theta_0\) is specified in degrees, but LAMMPS converts it to
radians internally; hence \(K\) is effectively energy per
radian\textasciicircum{}2.  Note that the usual 1/2 factor is included in \(K\).


\bigskip\hrule\bigskip


\sphinxAtStartPar
The \sphinxstyleemphasis{dihedral} keyword applies a dihedral restraint to the specified
atoms using a simplified form of the function used by the
{\hyperref[\detokenize{dihedral_charmm::doc}]{\sphinxcrossref{\DUrole{doc}{dihedral\_style charmm}}}} command.  The potential
associated with the restraint is
\begin{equation*}
\begin{split}E = K [ 1 + \cos (n \phi - d) ]\end{split}
\end{equation*}
\sphinxAtStartPar
with the following coefficients:
\begin{itemize}
\item {} 
\sphinxAtStartPar
\(K\) (energy)

\item {} 
\sphinxAtStartPar
\(n\) (multiplicity, \textgreater{}= 0)

\item {} 
\sphinxAtStartPar
\(d\) (degrees) = \(\phi_0 + 180\)

\end{itemize}

\sphinxAtStartPar
\(K\) and \(\phi_0\) are specified with the fix.  Note that the value of the
dihedral multiplicity \(n\) is set by default to 1. You can use the
optional \sphinxstyleemphasis{mult} keyword to set it to a different positive integer.
Also note that the energy will be a minimum when the
current dihedral angle \(\phi\) is equal to \(\phi_0\).


\bigskip\hrule\bigskip



\subsection{Restart, fix\_modify, output, run start/stop, minimize info}
\label{\detokenize{fix_restrain:restart-fix-modify-output-run-start-stop-minimize-info}}
\sphinxAtStartPar
No information about this fix is written to {\hyperref[\detokenize{restart::doc}]{\sphinxcrossref{\DUrole{doc}{binary restart files}}}}.

\sphinxAtStartPar
The {\hyperref[\detokenize{fix_modify::doc}]{\sphinxcrossref{\DUrole{doc}{fix\_modify}}}} \sphinxstyleemphasis{energy} option is supported by
this fix to add the potential energy associated with this fix to the
global potential energy of the system as part of {\hyperref[\detokenize{thermo_style::doc}]{\sphinxcrossref{\DUrole{doc}{thermodynamic
output}}}} The default setting for this fix is
{\hyperref[\detokenize{fix_modify::doc}]{\sphinxcrossref{\DUrole{doc}{fix\_modify energy no}}}}.

\sphinxAtStartPar
The {\hyperref[\detokenize{fix_modify::doc}]{\sphinxcrossref{\DUrole{doc}{fix\_modify}}}} \sphinxstyleemphasis{respa} option is supported by this
fix. This allows to set at which level of the {\hyperref[\detokenize{run_style::doc}]{\sphinxcrossref{\DUrole{doc}{r\sphinxhyphen{}RESPA}}}} integrator the fix is adding its forces. Default is the
outermost level.

\begin{sphinxadmonition}{note}{Note:}
\sphinxAtStartPar
If you want the fictitious potential energy associated with the
added forces to be included in the total potential energy of the
system (the quantity being minimized), you MUST enable the
{\hyperref[\detokenize{fix_modify::doc}]{\sphinxcrossref{\DUrole{doc}{fix\_modify}}}} \sphinxstyleemphasis{energy} option for this fix.
\end{sphinxadmonition}

\sphinxAtStartPar
This fix computes a global scalar and a global vector of length 3,
which can be accessed by various {\hyperref[\detokenize{Howto_output::doc}]{\sphinxcrossref{\DUrole{doc}{output commands}}}}.
The scalar is the total potential energy for \sphinxstyleemphasis{all} the restraints as
discussed above. The vector values are the sum of contributions to the
following individual categories:
\begin{enumerate}
\sphinxsetlistlabels{\arabic}{enumi}{enumii}{}{.}%
\item {} 
\sphinxAtStartPar
bond energy

\item {} 
\sphinxAtStartPar
angle energy

\item {} 
\sphinxAtStartPar
dihedral energy

\end{enumerate}

\sphinxAtStartPar
The scalar and vector values calculated by this fix are “extensive”.

\sphinxAtStartPar
No parameter of this fix can be used with the \sphinxstyleemphasis{start/stop} keywords of
the {\hyperref[\detokenize{run::doc}]{\sphinxcrossref{\DUrole{doc}{run}}}} command.


\subsection{Restrictions}
\label{\detokenize{fix_restrain:restrictions}}\begin{quote}

\sphinxAtStartPar
none
\end{quote}


\subsection{Related commands}
\label{\detokenize{fix_restrain:related-commands}}
\sphinxAtStartPar
none


\subsection{Default}
\label{\detokenize{fix_restrain:default}}
\sphinxAtStartPar
none

\sphinxstepscope

\index{fix rheo@\spxentry{fix rheo}}\ignorespaces 

\section{fix rheo command}
\label{\detokenize{fix_rheo:fix-rheo-command}}\label{\detokenize{fix_rheo:index-0}}\label{\detokenize{fix_rheo::doc}}

\subsection{Syntax}
\label{\detokenize{fix_rheo:syntax}}
\begin{sphinxVerbatim}[commandchars=\\\{\}]
fix ID group\PYGZhy{}ID rheo cut kstyle zmin keyword values...
\end{sphinxVerbatim}
\begin{itemize}
\item {} 
\sphinxAtStartPar
ID, group\sphinxhyphen{}ID are documented in {\hyperref[\detokenize{fix::doc}]{\sphinxcrossref{\DUrole{doc}{fix}}}} command

\item {} 
\sphinxAtStartPar
rheo = style name of this fix command

\item {} 
\sphinxAtStartPar
cut = cutoff for the kernel (distance)

\item {} 
\sphinxAtStartPar
kstyle = \sphinxstyleemphasis{quintic} or \sphinxstyleemphasis{RK0} or \sphinxstyleemphasis{RK1} or \sphinxstyleemphasis{RK2}

\item {} 
\sphinxAtStartPar
zmin = minimal number of neighbors for reproducing kernels

\item {} 
\sphinxAtStartPar
zero or more keyword/value pairs may be appended to args

\item {} \begin{description}
\sphinxlineitem{keyword = \sphinxstyleemphasis{thermal} or \sphinxstyleemphasis{interface/reconstruct} or \sphinxstyleemphasis{surface/detection} or}
\sphinxAtStartPar
\sphinxstyleemphasis{shift} or \sphinxstyleemphasis{rho/sum} or \sphinxstyleemphasis{density} or \sphinxstyleemphasis{self/mass} or \sphinxstyleemphasis{speed/sound}

\end{description}
\begin{sphinxalltt}
\sphinxstyleemphasis{thermal} values = none, turns on thermal evolution
\sphinxstyleemphasis{interface/reconstruct} values = none, reconstructs interfaces with solid particles
\sphinxstyleemphasis{surface/detection} values = \sphinxstyleemphasis{sdstyle} \sphinxstyleemphasis{limit} \sphinxstyleemphasis{limit/splash}
  \sphinxstyleemphasis{sdstyle} = \sphinxstyleemphasis{coordination} or \sphinxstyleemphasis{divergence}
  \sphinxstyleemphasis{limit} = threshold for surface particles
  \sphinxstyleemphasis{limit/splash} = threshold for splash particles
\sphinxstyleemphasis{shift} values = none, turns on velocity shifting
\sphinxstyleemphasis{rho/sum} values = none, uses the kernel to compute the density of particles
\sphinxstyleemphasis{self/mass} values = none, a particle uses its own mass in a rho summation
\sphinxstyleemphasis{density} values = \sphinxstyleemphasis{rho01}, ... \sphinxstyleemphasis{rho0N} (density)
\sphinxstyleemphasis{speed/sound} values = \sphinxstyleemphasis{cs0}, ... \sphinxstyleemphasis{csN} (velocity)
\end{sphinxalltt}

\end{itemize}


\subsection{Examples}
\label{\detokenize{fix_rheo:examples}}
\begin{sphinxVerbatim}[commandchars=\\\{\}]
\PYG{k}{fix }\PYG{n+nv+nvIdentifier}{1}\PYG{+w}{ }\PYG{n+nv+nvIdentifier}{all}\PYG{+w}{ }\PYG{n}{rheo}\PYG{+w}{ }\PYG{l+m}{3.0}\PYG{+w}{ }\PYG{n}{quintic}\PYG{+w}{ }\PYG{l+m}{0}\PYG{+w}{ }\PYG{n}{thermal}\PYG{+w}{ }\PYG{n}{density}\PYG{+w}{ }\PYG{l+m}{0.1}\PYG{+w}{ }\PYG{l+m}{0.1}\PYG{+w}{ }\PYG{n}{speed}\PYG{o}{/}\PYG{n}{sound}\PYG{+w}{ }\PYG{l+m}{10.0}\PYG{+w}{ }\PYG{l+m}{1.0}
\PYG{k}{fix }\PYG{n+nv+nvIdentifier}{1}\PYG{+w}{ }\PYG{n+nv+nvIdentifier}{all}\PYG{+w}{ }\PYG{n}{rheo}\PYG{+w}{ }\PYG{l+m}{3.0}\PYG{+w}{ }\PYG{n}{RK1}\PYG{+w}{ }\PYG{l+m}{10}\PYG{+w}{ }\PYG{n}{shift}\PYG{+w}{ }\PYG{n}{surface}\PYG{o}{/}\PYG{n}{detection}\PYG{+w}{ }\PYG{n}{coordination}\PYG{+w}{ }\PYG{l+m}{40}
\end{sphinxVerbatim}


\subsection{Description}
\label{\detokenize{fix_rheo:description}}
\sphinxAtStartPar
\DUrole{versionmodified,added}{Added in version 29Aug2024.}

\sphinxAtStartPar
Perform time integration for RHEO particles, updating positions, velocities,
and densities. For an overview of other features available in the RHEO package,
see {\hyperref[\detokenize{Howto_rheo::doc}]{\sphinxcrossref{\DUrole{doc}{the RHEO howto}}}}.

\sphinxAtStartPar
The type of kernel is specified using \sphinxstyleemphasis{kstyle} and the cutoff is \sphinxstyleemphasis{cut}. Four
kernels are currently available. The \sphinxstyleemphasis{quintic} kernel is a standard quintic
spline function commonly used in SPH. The other options, \sphinxstyleemphasis{RK0}, \sphinxstyleemphasis{RK1}, and
\sphinxstyleemphasis{RK2}, are zeroth, first, and second order reproducing. To generate a
reproducing kernel, a particle must have sufficient neighbors inside the
kernel cutoff distance (a coordination number) to accurately calculate
moments. This threshold is set by \sphinxstyleemphasis{zmin}. If reproducing kernels are
requested but a particle has fewer neighbors, then it will revert to a
non\sphinxhyphen{}reproducing quintic kernel until it gains more neighbors.

\sphinxAtStartPar
To model temperature evolution, one must specify the \sphinxstyleemphasis{thermal} keyword,
define a separate instance of {\hyperref[\detokenize{fix_rheo_thermal::doc}]{\sphinxcrossref{\DUrole{doc}{fix rheo/thermal}}}},
and use atom style rheo/thermal.

\sphinxAtStartPar
By default, the density of solid RHEO particles does not evolve and forces
with fluid particles are calculated using the current velocity of the solid
particle. If the \sphinxstyleemphasis{interface/reconstruct} keyword is used, then the density
and velocity of solid particles are alternatively reconstructed for every
fluid\sphinxhyphen{}solid interaction to ensure no\sphinxhyphen{}slip and pressure\sphinxhyphen{}balanced boundaries.
This is done by estimating the location of the fluid\sphinxhyphen{}solid interface and
extrapolating fluid particle properties across the interface to calculate a
temporary apparent density and velocity for a solid particle.

\sphinxAtStartPar
A modified form of Fickian particle shifting can be enabled with the
\sphinxstyleemphasis{shift} keyword. This effectively shifts particle positions to generate a
more uniform spatial distribution. Shifting currently does not consider the
type of a particle and therefore may be inappropriate in systems consisting
of multiple fluid phases.

\sphinxAtStartPar
In systems with free surfaces, the \sphinxstyleemphasis{surface/detection} keyword can be used
to classify the location of particles as being within the bulk fluid, on a
free surface, or isolated from other particles in a splash or droplet.
Shifting is then disabled in the normal direction away from the free surface
to prevent particles from diffusing away. Surface detection can also be used
to control surface\sphinxhyphen{}nucleated effects like oxidation when used in combination
with {\hyperref[\detokenize{fix_rheo_oxidation::doc}]{\sphinxcrossref{\DUrole{doc}{fix rheo/oxidation}}}}. Surface detection is not
performed on solid bodies.

\sphinxAtStartPar
The \sphinxstyleemphasis{surface/detection} keyword takes three arguments: \sphinxstyleemphasis{sdstyle}, \sphinxstyleemphasis{limit},
and \sphinxstyleemphasis{limit/splash}. The first, \sphinxstyleemphasis{sdstyle}, specifies whether surface particles
are identified using a coordination number (\sphinxstyleemphasis{coordination}) or the divergence
of the local particle positions (\sphinxstyleemphasis{divergence}). The threshold value for a
surface particle for either of these criteria is set by the numerical value
of \sphinxstyleemphasis{limit}. Additionally, if a particle’s coordination number is too low,
i.e. if it has separated off from the bulk in a droplet, it is not possible
to define surfaces and the particle is classified as a splash. The coordination
threshold for this classification is set by the numerical value of
\sphinxstyleemphasis{limit/splash}.

\sphinxAtStartPar
By default, RHEO integrates particles’ densities using a mass diffusion
equation. Alternatively, one can update densities every timestep by performing
a kernel summation of the masses of neighboring particles by specifying the \sphinxstyleemphasis{rho/sum}
keyword.

\sphinxAtStartPar
The \sphinxstyleemphasis{self/mass} keyword modifies the behavior of the density summation in \sphinxstyleemphasis{rho/sum}.
Typically, the density \(\rho\) of a particle is calculated as the sum over neighbors
\begin{equation*}
\begin{split}\rho_i = \sum_{j} W_{ij} M_j\end{split}
\end{equation*}
\sphinxAtStartPar
where \(W_{ij}\) is the kernel, and \(M_j\) is the mass of particle \(j\).
The \sphinxstyleemphasis{self/mass} keyword augments this expression by replacing \(M_j\) with
\(M_i\). This may be useful in simulations of multiple fluid phases with large
differences in density, {\hyperref[\detokenize{fix_rheo:fix-rheo-hu}]{\sphinxcrossref{\DUrole{std,std-ref}{(Hu)}}}}.

\sphinxAtStartPar
The \sphinxstyleemphasis{density} keyword is used to specify the equilibrium density of each of the N
particle types. It must be followed by N numerical values specifying each type’s
equilibrium density \sphinxstyleemphasis{rho0}.

\sphinxAtStartPar
The \sphinxstyleemphasis{speed/sound} keyword is used to specify the speed of sound of each of the
N particle types. It must be followed by N numerical values specifying each type’s
speed of sound \sphinxstyleemphasis{cs}.


\subsection{Restart, fix\_modify, output, run start/stop, minimize info}
\label{\detokenize{fix_rheo:restart-fix-modify-output-run-start-stop-minimize-info}}
\sphinxAtStartPar
No information about this fix is written to {\hyperref[\detokenize{restart::doc}]{\sphinxcrossref{\DUrole{doc}{binary restart files}}}}.
None of the {\hyperref[\detokenize{fix_modify::doc}]{\sphinxcrossref{\DUrole{doc}{fix\_modify}}}} options
are relevant to this fix.  No global or per\sphinxhyphen{}atom quantities are stored
by this fix for access by various {\hyperref[\detokenize{Howto_output::doc}]{\sphinxcrossref{\DUrole{doc}{output commands}}}}.
No parameter of this fix can be used with the \sphinxstyleemphasis{start/stop} keywords of
the {\hyperref[\detokenize{run::doc}]{\sphinxcrossref{\DUrole{doc}{run}}}} command.  This fix is not invoked during
{\hyperref[\detokenize{minimize::doc}]{\sphinxcrossref{\DUrole{doc}{energy minimization}}}}.


\subsection{Restrictions}
\label{\detokenize{fix_rheo:restrictions}}
\sphinxAtStartPar
This fix must be used with atom style rheo or rheo/thermal. This fix must
be used in conjunction with {\hyperref[\detokenize{fix_rheo_pressure::doc}]{\sphinxcrossref{\DUrole{doc}{fix rheo/pressure}}}}.
and {\hyperref[\detokenize{fix_rheo_viscosity::doc}]{\sphinxcrossref{\DUrole{doc}{fix rheo/viscosity}}}}. If the \sphinxstyleemphasis{thermal} setting
is used, there must also be an instance of
{\hyperref[\detokenize{fix_rheo_thermal::doc}]{\sphinxcrossref{\DUrole{doc}{fix rheo/thermal}}}}. The fix group must be set to all.
Only one instance of fix rheo may be defined and it  must be defined prior
to all other RHEO fixes in the input script.

\sphinxAtStartPar
This fix is part of the RHEO package.  It is only enabled if LAMMPS was built
with that package. See the {\hyperref[\detokenize{Build_package::doc}]{\sphinxcrossref{\DUrole{doc}{Build package}}}} page for more info.


\subsection{Related commands}
\label{\detokenize{fix_rheo:related-commands}}
\sphinxAtStartPar
{\hyperref[\detokenize{fix_rheo_viscosity::doc}]{\sphinxcrossref{\DUrole{doc}{fix rheo/viscosity}}}},
{\hyperref[\detokenize{fix_rheo_pressure::doc}]{\sphinxcrossref{\DUrole{doc}{fix rheo/pressure}}}},
{\hyperref[\detokenize{fix_rheo_thermal::doc}]{\sphinxcrossref{\DUrole{doc}{fix rheo/thermal}}}},
{\hyperref[\detokenize{pair_rheo::doc}]{\sphinxcrossref{\DUrole{doc}{pair rheo}}}},
{\hyperref[\detokenize{compute_rheo_property_atom::doc}]{\sphinxcrossref{\DUrole{doc}{compute rheo/property/atom}}}}


\subsection{Default}
\label{\detokenize{fix_rheo:default}}
\sphinxAtStartPar
\sphinxstyleemphasis{rho0} and \sphinxstyleemphasis{cs} are set to 1.0 for all atom types.


\bigskip\hrule\bigskip

\phantomsection\label{\detokenize{fix_rheo:fix-rheo-hu}}
\sphinxAtStartPar
\sphinxstylestrong{(Hu)} Hu, and Adams J. Comp. Physics, 213, 844\sphinxhyphen{}861 (2006).

\sphinxstepscope

\index{fix rheo/oxidation@\spxentry{fix rheo/oxidation}}\ignorespaces 

\section{fix rheo/oxidation command}
\label{\detokenize{fix_rheo_oxidation:fix-rheo-oxidation-command}}\label{\detokenize{fix_rheo_oxidation:index-0}}\label{\detokenize{fix_rheo_oxidation::doc}}

\subsection{Syntax}
\label{\detokenize{fix_rheo_oxidation:syntax}}
\begin{sphinxVerbatim}[commandchars=\\\{\}]
fix ID group\PYGZhy{}ID rheo/oxidation cut btype rsurf
\end{sphinxVerbatim}
\begin{itemize}
\item {} 
\sphinxAtStartPar
ID, group\sphinxhyphen{}ID are documented in {\hyperref[\detokenize{fix::doc}]{\sphinxcrossref{\DUrole{doc}{fix}}}} command

\item {} 
\sphinxAtStartPar
rheo/oxidation = style name of this fix command

\item {} 
\sphinxAtStartPar
cut = maximum bond length (distance units)

\item {} 
\sphinxAtStartPar
btype = type of bonds created

\item {} 
\sphinxAtStartPar
rsurf = distance from surface to create bonds (distance units)

\end{itemize}


\subsection{Examples}
\label{\detokenize{fix_rheo_oxidation:examples}}
\begin{sphinxVerbatim}[commandchars=\\\{\}]
\PYG{k}{fix }\PYG{n+nv+nvIdentifier}{1}\PYG{+w}{ }\PYG{n+nv+nvIdentifier}{all}\PYG{+w}{ }\PYG{n}{rheo}\PYG{o}{/}\PYG{n}{oxidation}\PYG{+w}{ }\PYG{l+m}{1.5}\PYG{+w}{ }\PYG{l+m}{2}\PYG{+w}{ }\PYG{l+m}{0.0}
\PYG{k}{fix }\PYG{n+nv+nvIdentifier}{1}\PYG{+w}{ }\PYG{n+nv+nvIdentifier}{all}\PYG{+w}{ }\PYG{n}{rheo}\PYG{o}{/}\PYG{n}{oxidation}\PYG{+w}{ }\PYG{l+m}{1.0}\PYG{+w}{ }\PYG{l+m}{1}\PYG{+w}{ }\PYG{l+m}{2.0}
\end{sphinxVerbatim}


\subsection{Description}
\label{\detokenize{fix_rheo_oxidation:description}}
\sphinxAtStartPar
\DUrole{versionmodified,added}{Added in version 29Aug2024.}

\sphinxAtStartPar
This fix dynamically creates bonds on the surface of fluids to
represent physical processes such as oxidation. It is intended
for use with bond style {\hyperref[\detokenize{bond_rheo_shell::doc}]{\sphinxcrossref{\DUrole{doc}{bond rheo/shell}}}}.

\sphinxAtStartPar
Every timestep, particles check neighbors within a distance of \sphinxstyleemphasis{cut}.
This distance must be smaller than the kernel length defined in
{\hyperref[\detokenize{fix_rheo::doc}]{\sphinxcrossref{\DUrole{doc}{fix rheo}}}}. Bonds of type \sphinxstyleemphasis{btype} are created between
a fluid particle and either a fluid or solid neighbor. The fluid particles
must also be on the fluid surface, or within a distance of \sphinxstyleemphasis{rsurf} from
the surface. This process is further described in
{\hyperref[\detokenize{fix_rheo_oxidation:howto-rheo-clemmer2}]{\sphinxcrossref{\DUrole{std,std-ref}{(Clemmer)}}}}.

\sphinxAtStartPar
If used in conjunction with solid bodies, such as those generated
by the \sphinxstyleemphasis{react} option of {\hyperref[\detokenize{fix_rheo_thermal::doc}]{\sphinxcrossref{\DUrole{doc}{fix rheo/thermal}}}},
it is recommended to use a {\hyperref[\detokenize{bond_hybrid::doc}]{\sphinxcrossref{\DUrole{doc}{hybrid bond style}}}}
with different bond types for solid and oxide bonds.


\subsection{Restart, fix\_modify, output, run start/stop, minimize info}
\label{\detokenize{fix_rheo_oxidation:restart-fix-modify-output-run-start-stop-minimize-info}}
\sphinxAtStartPar
No information about this fix is written to {\hyperref[\detokenize{restart::doc}]{\sphinxcrossref{\DUrole{doc}{binary restart files}}}}.
None of the {\hyperref[\detokenize{fix_modify::doc}]{\sphinxcrossref{\DUrole{doc}{fix\_modify}}}} options
are relevant to this fix.  No global or per\sphinxhyphen{}atom quantities are stored
by this fix for access by various {\hyperref[\detokenize{Howto_output::doc}]{\sphinxcrossref{\DUrole{doc}{output commands}}}}.
No parameter of this fix can be used with the \sphinxstyleemphasis{start/stop} keywords of
the {\hyperref[\detokenize{run::doc}]{\sphinxcrossref{\DUrole{doc}{run}}}} command.  This fix is not invoked during {\hyperref[\detokenize{minimize::doc}]{\sphinxcrossref{\DUrole{doc}{energy minimization}}}}.


\subsection{Restrictions}
\label{\detokenize{fix_rheo_oxidation:restrictions}}
\sphinxAtStartPar
This fix must be used with the bond style {\hyperref[\detokenize{bond_rheo_shell::doc}]{\sphinxcrossref{\DUrole{doc}{rheo/shell}}}}
and {\hyperref[\detokenize{fix_rheo::doc}]{\sphinxcrossref{\DUrole{doc}{fix rheo}}}} with surface detection enabled.

\sphinxAtStartPar
This fix is part of the RHEO package.  It is only enabled if
LAMMPS was built with that package.  See the {\hyperref[\detokenize{Build_package::doc}]{\sphinxcrossref{\DUrole{doc}{Build package}}}}
page for more info.


\subsection{Related commands}
\label{\detokenize{fix_rheo_oxidation:related-commands}}
\sphinxAtStartPar
{\hyperref[\detokenize{fix_rheo::doc}]{\sphinxcrossref{\DUrole{doc}{fix rheo}}}},
{\hyperref[\detokenize{bond_rheo_shell::doc}]{\sphinxcrossref{\DUrole{doc}{bond rheo/shell}}}},
{\hyperref[\detokenize{compute_rheo_property_atom::doc}]{\sphinxcrossref{\DUrole{doc}{compute rheo/property/atom}}}}


\subsection{Default}
\label{\detokenize{fix_rheo_oxidation:default}}
\sphinxAtStartPar
none


\bigskip\hrule\bigskip

\phantomsection\label{\detokenize{fix_rheo_oxidation:howto-rheo-clemmer2}}
\sphinxAtStartPar
\sphinxstylestrong{(Clemmer)} Clemmer, Pierce, O’Connor, Nevins, Jones, Lechman, Tencer, Appl. Math. Model., 130, 310\sphinxhyphen{}326 (2024).

\sphinxstepscope

\index{fix rheo/pressure@\spxentry{fix rheo/pressure}}\ignorespaces 

\section{fix rheo/pressure command}
\label{\detokenize{fix_rheo_pressure:fix-rheo-pressure-command}}\label{\detokenize{fix_rheo_pressure:index-0}}\label{\detokenize{fix_rheo_pressure::doc}}

\subsection{Syntax}
\label{\detokenize{fix_rheo_pressure:syntax}}
\begin{sphinxVerbatim}[commandchars=\\\{\}]
fix ID group\PYGZhy{}ID rheo/pressure type1 pstyle1 args1 ... typeN pstyleN argsN
\end{sphinxVerbatim}
\begin{itemize}
\item {} 
\sphinxAtStartPar
ID, group\sphinxhyphen{}ID are documented in {\hyperref[\detokenize{fix::doc}]{\sphinxcrossref{\DUrole{doc}{fix}}}} command

\item {} 
\sphinxAtStartPar
rheo/pressure = style name of this fix command

\item {} 
\sphinxAtStartPar
one or more types and pressure styles must be appended

\item {} 
\sphinxAtStartPar
types = lists of types (see below)

\item {} 
\sphinxAtStartPar
pstyle = \sphinxstyleemphasis{linear} or \sphinxstyleemphasis{taitwater} or \sphinxstyleemphasis{cubic}
\begin{sphinxalltt}
\sphinxstyleemphasis{linear} args = none
\sphinxstyleemphasis{taitwater} args = none
\sphinxstyleemphasis{cubic} args = cubic prefactor \(A_3\) (pressure/density\textasciicircum{}2)
\end{sphinxalltt}

\end{itemize}


\subsection{Examples}
\label{\detokenize{fix_rheo_pressure:examples}}
\begin{sphinxVerbatim}[commandchars=\\\{\}]
\PYG{k}{fix }\PYG{n+nv+nvIdentifier}{1}\PYG{+w}{ }\PYG{n+nv+nvIdentifier}{all}\PYG{+w}{ }\PYG{n}{rheo}\PYG{o}{/}\PYG{n}{pressure}\PYG{+w}{ }\PYG{o}{*}\PYG{+w}{ }\PYG{n}{linear}
\PYG{k}{fix }\PYG{n+nv+nvIdentifier}{1}\PYG{+w}{ }\PYG{n+nv+nvIdentifier}{all}\PYG{+w}{ }\PYG{n}{rheo}\PYG{o}{/}\PYG{n}{pressure}\PYG{+w}{ }\PYG{l+m}{1}\PYG{+w}{ }\PYG{n}{linear}\PYG{+w}{ }\PYG{l+m}{2}\PYG{+w}{ }\PYG{n}{cubic}\PYG{+w}{ }\PYG{l+m}{10.0}
\end{sphinxVerbatim}


\subsection{Description}
\label{\detokenize{fix_rheo_pressure:description}}
\sphinxAtStartPar
\DUrole{versionmodified,added}{Added in version 29Aug2024.}

\sphinxAtStartPar
This fix defines a pressure equation of state for RHEO particles. One can
define different equations of state for different atom types. An equation
must be specified for every atom type.

\sphinxAtStartPar
One first defines the atom \sphinxstyleemphasis{types}. A wild\sphinxhyphen{}card asterisk can be used in place
of or in conjunction with the \sphinxstyleemphasis{types} argument to set the coefficients for
multiple pairs of atom types.  This takes the form “*” or “*n” or “m*”
or “m*n”.  If \(N\) is the number of atom types, then an asterisk with
no numeric values means all types from 1 to \(N\).  A leading asterisk
means all types from 1 to n (inclusive).  A trailing asterisk means all types
from m to \(N\) (inclusive).  A middle asterisk means all types from m to n
(inclusive).

\sphinxAtStartPar
The \sphinxstyleemphasis{types} definition is followed by the pressure style, \sphinxstyleemphasis{pstyle}. Current
options \sphinxstyleemphasis{linear}, \sphinxstyleemphasis{taitwater}, and \sphinxstyleemphasis{cubic}. Style \sphinxstyleemphasis{linear} is a linear
equation of state with a particle pressure \(P\) calculated as
\begin{equation*}
\begin{split}P = c (\rho - \rho_0)\end{split}
\end{equation*}
\sphinxAtStartPar
where \(c\) is the speed of sound, \(\rho_0\) is the equilibrium density,
and \(\rho\) is the current density of a particle. The numerical values of
\(c\) and \(\rho_0\) are set in {\hyperref[\detokenize{fix_rheo::doc}]{\sphinxcrossref{\DUrole{doc}{fix rheo}}}}. Style \sphinxstyleemphasis{cubic}
is a cubic equation of state which has an extra argument \(A_3\),
\begin{equation*}
\begin{split}P = c ((\rho - \rho_0) + A_3 (\rho - \rho_0)^3) .\end{split}
\end{equation*}
\sphinxAtStartPar
Style \sphinxstyleemphasis{taitwater} is Tait’s equation of state:
\begin{equation*}
\begin{split}P = \frac{c^2 \rho_0}{7} \biggl[\left(\frac{\rho}{\rho_0}\right)^{7} - 1\biggr].\end{split}
\end{equation*}

\subsection{Restart, fix\_modify, output, run start/stop, minimize info}
\label{\detokenize{fix_rheo_pressure:restart-fix-modify-output-run-start-stop-minimize-info}}
\sphinxAtStartPar
No information about this fix is written to {\hyperref[\detokenize{restart::doc}]{\sphinxcrossref{\DUrole{doc}{binary restart files}}}}.
None of the {\hyperref[\detokenize{fix_modify::doc}]{\sphinxcrossref{\DUrole{doc}{fix\_modify}}}} options
are relevant to this fix.  No global or per\sphinxhyphen{}atom quantities are stored
by this fix for access by various {\hyperref[\detokenize{Howto_output::doc}]{\sphinxcrossref{\DUrole{doc}{output commands}}}}.
No parameter of this fix can be used with the \sphinxstyleemphasis{start/stop} keywords of
the {\hyperref[\detokenize{run::doc}]{\sphinxcrossref{\DUrole{doc}{run}}}} command.  This fix is not invoked during {\hyperref[\detokenize{minimize::doc}]{\sphinxcrossref{\DUrole{doc}{energy minimization}}}}.


\subsection{Restrictions}
\label{\detokenize{fix_rheo_pressure:restrictions}}
\sphinxAtStartPar
This fix must be used with an atom style that includes density
such as atom\_style rheo or rheo/thermal. This fix must be used in
conjunction with {\hyperref[\detokenize{fix_rheo::doc}]{\sphinxcrossref{\DUrole{doc}{fix rheo}}}}. The fix group must be
set to all. Only one instance of fix rheo/pressure can be defined.

\sphinxAtStartPar
This fix is part of the RHEO package.  It is only enabled if
LAMMPS was built with that package.  See the {\hyperref[\detokenize{Build_package::doc}]{\sphinxcrossref{\DUrole{doc}{Build package}}}}
page for more info.


\subsection{Related commands}
\label{\detokenize{fix_rheo_pressure:related-commands}}
\sphinxAtStartPar
{\hyperref[\detokenize{fix_rheo::doc}]{\sphinxcrossref{\DUrole{doc}{fix rheo}}}},
{\hyperref[\detokenize{pair_rheo::doc}]{\sphinxcrossref{\DUrole{doc}{pair rheo}}}},
{\hyperref[\detokenize{compute_rheo_property_atom::doc}]{\sphinxcrossref{\DUrole{doc}{compute rheo/property/atom}}}}


\subsection{Default}
\label{\detokenize{fix_rheo_pressure:default}}
\sphinxAtStartPar
none

\sphinxstepscope

\index{fix rheo/thermal@\spxentry{fix rheo/thermal}}\ignorespaces 

\section{fix rheo/thermal command}
\label{\detokenize{fix_rheo_thermal:fix-rheo-thermal-command}}\label{\detokenize{fix_rheo_thermal:index-0}}\label{\detokenize{fix_rheo_thermal::doc}}

\subsection{Syntax}
\label{\detokenize{fix_rheo_thermal:syntax}}
\begin{sphinxVerbatim}[commandchars=\\\{\}]
fix ID group\PYGZhy{}ID rheo/thermal attribute values ...
\end{sphinxVerbatim}
\begin{itemize}
\item {} 
\sphinxAtStartPar
ID, group\sphinxhyphen{}ID are documented in {\hyperref[\detokenize{fix::doc}]{\sphinxcrossref{\DUrole{doc}{fix}}}} command

\item {} 
\sphinxAtStartPar
rheo/thermal = style name of this fix command

\item {} 
\sphinxAtStartPar
one or more attributes may be appended

\item {} 
\sphinxAtStartPar
attribute = \sphinxstyleemphasis{conductivity} or \sphinxstyleemphasis{specific/heat} or \sphinxstyleemphasis{latent/heat} or \sphinxstyleemphasis{Tfreeze} or \sphinxstyleemphasis{react}
\begin{sphinxalltt}
\sphinxstyleemphasis{conductivity} args = types style args
  types = lists of types (see below)
  style = \sphinxstyleemphasis{constant}
    \sphinxstyleemphasis{constant} arg = conductivity (power/temperature)
\sphinxstyleemphasis{specific/heat} args = types style args
  types = lists of types (see below)
  style = \sphinxstyleemphasis{constant}
    \sphinxstyleemphasis{constant} arg = specific heat (energy/(mass*temperature))
\sphinxstyleemphasis{latent/heat} args = types style args
  types = lists of types (see below)
  style = \sphinxstyleemphasis{constant}
    \sphinxstyleemphasis{constant} arg = latent heat (energy/mass)
\sphinxstyleemphasis{Tfreeze} args = types style args
  types = lists of types (see below)
  style = \sphinxstyleemphasis{constant}
    \sphinxstyleemphasis{constant} arg = freezing temperature (temperature)
\sphinxstyleemphasis{react} args = cut type
  cut = maximum bond distance
  type = bond type
\end{sphinxalltt}

\end{itemize}


\subsection{Examples}
\label{\detokenize{fix_rheo_thermal:examples}}
\begin{sphinxVerbatim}[commandchars=\\\{\}]
\PYG{k}{fix }\PYG{n+nv+nvIdentifier}{1}\PYG{+w}{ }\PYG{n+nv+nvIdentifier}{all}\PYG{+w}{ }\PYG{n}{rheo}\PYG{o}{/}\PYG{n}{thermal}\PYG{+w}{ }\PYG{n}{conductivity}\PYG{+w}{ }\PYG{o}{*}\PYG{+w}{ }\PYG{n}{constant}\PYG{+w}{ }\PYG{l+m}{1.0}\PYG{+w}{ }\PYG{n}{specific}\PYG{o}{/}\PYG{n}{heat}\PYG{+w}{ }\PYG{o}{*}\PYG{+w}{ }\PYG{n}{constant}\PYG{+w}{ }\PYG{l+m}{1.0}\PYG{+w}{ }\PYG{n}{Tfreeze}\PYG{+w}{ }\PYG{o}{*}\PYG{+w}{ }\PYG{n}{constant}\PYG{+w}{ }\PYG{l+m}{1.0}
\PYG{k}{fix }\PYG{n+nv+nvIdentifier}{1}\PYG{+w}{ }\PYG{n+nv+nvIdentifier}{all}\PYG{+w}{ }\PYG{n}{rheo}\PYG{o}{/}\PYG{n}{pressure}\PYG{+w}{ }\PYG{n}{conductivity}\PYG{+w}{ }\PYG{l+m}{1}\PYG{o}{*}\PYG{l+m}{2}\PYG{+w}{ }\PYG{n}{constant}\PYG{+w}{ }\PYG{l+m}{1.0}\PYG{+w}{ }\PYG{n}{conductivity}\PYG{+w}{ }\PYG{l+m}{3}\PYG{o}{*}\PYG{l+m}{4}\PYG{+w}{ }\PYG{n}{constant}\PYG{+w}{ }\PYG{l+m}{2.0}\PYG{+w}{ }\PYG{n}{specific}\PYG{o}{/}\PYG{n}{heat}\PYG{+w}{ }\PYG{o}{*}\PYG{+w}{ }\PYG{n}{constant}\PYG{+w}{ }\PYG{l+m}{1.0}
\end{sphinxVerbatim}


\subsection{Description}
\label{\detokenize{fix_rheo_thermal:description}}
\sphinxAtStartPar
\DUrole{versionmodified,added}{Added in version 29Aug2024.}

\sphinxAtStartPar
This fix performs time integration of temperature for atom style rheo/thermal.
In addition, it defines multiple thermal properties of particles and handles
melting/solidification, if applicable. For more details on phase transitions
in RHEO, see {\hyperref[\detokenize{Howto_rheo::doc}]{\sphinxcrossref{\DUrole{doc}{the RHEO howto}}}}.

\sphinxAtStartPar
Note that the temperature of a particle is always derived from the energy.
This implies the \sphinxstyleemphasis{temperature} attribute of {\hyperref[\detokenize{set::doc}]{\sphinxcrossref{\DUrole{doc}{the set command}}}} does
not affect particles. Instead, one should use the \sphinxstyleemphasis{sph/e} attribute.

\sphinxAtStartPar
For each atom type, one can define expressions for the \sphinxstyleemphasis{conductivity},
\sphinxstyleemphasis{specific/heat}, \sphinxstyleemphasis{latent/heat}, and critical temperature (\sphinxstyleemphasis{Tfreeze}).
The conductivity and specific heat must be defined for all atom types.
The latent heat and critical temperature are optional. However, a
critical temperature must be defined to specify a latent heat.

\sphinxAtStartPar
Note, if shifting is turned on in {\hyperref[\detokenize{fix_rheo::doc}]{\sphinxcrossref{\DUrole{doc}{fix rheo}}}}, the gradient
of the energy is used to shift energies. This may be inappropriate in systems
with multiple atom types with different specific heats.

\sphinxAtStartPar
For each property, one must first define a list of atom types. A wild\sphinxhyphen{}card
asterisk can be used in place of or in conjunction with the \sphinxstyleemphasis{types} argument
to set the coefficients for multiple pairs of atom types.  This takes the
form “*” or “*n” or “m*” or “m*n”.  If \(N\) is the number of atom
types, then an asterisk with no numeric values means all types from 1 to
\(N\).  A leading asterisk means all types from 1 to n (inclusive).
A trailing asterisk means all types from m to \(N\) (inclusive).  A
middle asterisk means all types from m to n (inclusive).

\sphinxAtStartPar
The \sphinxstyleemphasis{types} definition for each property is followed by the style. Currently,
the only option is \sphinxstyleemphasis{constant}. Style \sphinxstyleemphasis{constant} simply applies a constant value
of respective property to each particle of the assigned type.

\sphinxAtStartPar
The \sphinxstyleemphasis{react} keyword controls whether bonds are created/deleted when particles
transition between a fluid and solid state. This option only applies to atom
types that have a defined value of \sphinxstyleemphasis{Tfreeze}. When a fluid particle’s
temperature drops below \sphinxstyleemphasis{Tfreeze}, bonds of type \sphinxstyleemphasis{btype} are created between
nearby solid particles within a distance of \sphinxstyleemphasis{cut}. The particle’s status also
swaps to a solid state. When a solid particle’s temperature rises above
\sphinxstyleemphasis{Tfreeze}, all bonds of type \sphinxstyleemphasis{btype} are broken and the particle’s status swaps
to a fluid state.


\subsection{Restart, fix\_modify, output, run start/stop, minimize info}
\label{\detokenize{fix_rheo_thermal:restart-fix-modify-output-run-start-stop-minimize-info}}
\sphinxAtStartPar
No information about this fix is written to {\hyperref[\detokenize{restart::doc}]{\sphinxcrossref{\DUrole{doc}{binary restart files}}}}.
None of the {\hyperref[\detokenize{fix_modify::doc}]{\sphinxcrossref{\DUrole{doc}{fix\_modify}}}} options
are relevant to this fix.  No global or per\sphinxhyphen{}atom quantities are stored
by this fix for access by various {\hyperref[\detokenize{Howto_output::doc}]{\sphinxcrossref{\DUrole{doc}{output commands}}}}.
No parameter of this fix can be used with the \sphinxstyleemphasis{start/stop} keywords of
the {\hyperref[\detokenize{run::doc}]{\sphinxcrossref{\DUrole{doc}{run}}}} command.  This fix is not invoked during {\hyperref[\detokenize{minimize::doc}]{\sphinxcrossref{\DUrole{doc}{energy minimization}}}}.


\subsection{Restrictions}
\label{\detokenize{fix_rheo_thermal:restrictions}}
\sphinxAtStartPar
This fix must be used with an atom style that includes temperature,
heatflow, and conductivity such as atom\_style rheo/thermal This fix
must be used in conjunction with {\hyperref[\detokenize{fix_rheo::doc}]{\sphinxcrossref{\DUrole{doc}{fix rheo}}}} with the
\sphinxstyleemphasis{thermal} setting. The fix group must be set to all. Only one
instance of fix rheo/pressure can be defined.

\sphinxAtStartPar
This fix is part of the RHEO package.  It is only enabled if
LAMMPS was built with that package.  See the {\hyperref[\detokenize{Build_package::doc}]{\sphinxcrossref{\DUrole{doc}{Build package}}}}
page for more info.


\subsection{Related commands}
\label{\detokenize{fix_rheo_thermal:related-commands}}
\sphinxAtStartPar
{\hyperref[\detokenize{fix_rheo::doc}]{\sphinxcrossref{\DUrole{doc}{fix rheo}}}},
{\hyperref[\detokenize{pair_rheo::doc}]{\sphinxcrossref{\DUrole{doc}{pair rheo}}}},
{\hyperref[\detokenize{compute_rheo_property_atom::doc}]{\sphinxcrossref{\DUrole{doc}{compute rheo/property/atom}}}},
{\hyperref[\detokenize{fix_add_heat::doc}]{\sphinxcrossref{\DUrole{doc}{fix add/heat}}}}


\subsection{Default}
\label{\detokenize{fix_rheo_thermal:default}}
\sphinxAtStartPar
none

\sphinxstepscope

\index{fix rheo/viscosity@\spxentry{fix rheo/viscosity}}\ignorespaces 

\section{fix rheo/viscosity command}
\label{\detokenize{fix_rheo_viscosity:fix-rheo-viscosity-command}}\label{\detokenize{fix_rheo_viscosity:index-0}}\label{\detokenize{fix_rheo_viscosity::doc}}

\subsection{Syntax}
\label{\detokenize{fix_rheo_viscosity:syntax}}
\begin{sphinxVerbatim}[commandchars=\\\{\}]
fix ID group\PYGZhy{}ID rheo/viscosity type1 pstyle1 args1 ... typeN pstyleN argsN
\end{sphinxVerbatim}
\begin{itemize}
\item {} 
\sphinxAtStartPar
ID, group\sphinxhyphen{}ID are documented in {\hyperref[\detokenize{fix::doc}]{\sphinxcrossref{\DUrole{doc}{fix}}}} command

\item {} 
\sphinxAtStartPar
rheo/viscosity = style name of this fix command

\item {} 
\sphinxAtStartPar
one or more types and viscosity styles must be appended

\item {} 
\sphinxAtStartPar
types = lists of types (see below)

\item {} 
\sphinxAtStartPar
vstyle = \sphinxstyleemphasis{constant} or \sphinxstyleemphasis{power}
\begin{sphinxalltt}
\sphinxstyleemphasis{constant} args = \sphinxstyleemphasis{eta}
  \sphinxstyleemphasis{eta} = viscosity

\sphinxstyleemphasis{power} args = \sphinxstyleemphasis{eta}, \sphinxstyleemphasis{gd0}, \sphinxstyleemphasis{K}, \sphinxstyleemphasis{n}
  \sphinxstyleemphasis{eta} = viscosity
  \sphinxstyleemphasis{gd0} = critical strain rate
  \sphinxstyleemphasis{K} = consistency index
  \sphinxstyleemphasis{n} = power\sphinxhyphen{}law exponent
\end{sphinxalltt}

\end{itemize}


\subsection{Examples}
\label{\detokenize{fix_rheo_viscosity:examples}}
\begin{sphinxVerbatim}[commandchars=\\\{\}]
\PYG{k}{fix }\PYG{n+nv+nvIdentifier}{1}\PYG{+w}{ }\PYG{n+nv+nvIdentifier}{all}\PYG{+w}{ }\PYG{n}{rheo}\PYG{o}{/}\PYG{n}{viscosity}\PYG{+w}{ }\PYG{o}{*}\PYG{+w}{ }\PYG{n}{constant}\PYG{+w}{ }\PYG{l+m}{1.0}
\PYG{k}{fix }\PYG{n+nv+nvIdentifier}{1}\PYG{+w}{ }\PYG{n+nv+nvIdentifier}{all}\PYG{+w}{ }\PYG{n}{rheo}\PYG{o}{/}\PYG{n}{viscosity}\PYG{+w}{ }\PYG{l+m}{1}\PYG{+w}{ }\PYG{n}{constant}\PYG{+w}{ }\PYG{l+m}{1.0}\PYG{+w}{ }\PYG{l+m}{2}\PYG{+w}{ }\PYG{n}{power}\PYG{+w}{ }\PYG{l+m}{0.1}\PYG{+w}{ }\PYG{l+m}{5e\PYGZhy{}4}\PYG{+w}{ }\PYG{l+m}{0.001}\PYG{+w}{ }\PYG{l+m}{0.5}
\end{sphinxVerbatim}


\subsection{Description}
\label{\detokenize{fix_rheo_viscosity:description}}
\sphinxAtStartPar
\DUrole{versionmodified,added}{Added in version 29Aug2024.}

\sphinxAtStartPar
This fix defines a viscosity for RHEO particles. One can define different
viscosities for different atom types, but a viscosity must be specified for
every atom type.

\sphinxAtStartPar
One first defines the atom \sphinxstyleemphasis{types}. A wild\sphinxhyphen{}card asterisk can be used in place
of or in conjunction with the \sphinxstyleemphasis{types} argument to set the coefficients for
multiple pairs of atom types.  This takes the form “*” or “*n” or “m*”
or “m*n”.  If \(N\) is the number of atom types, then an asterisk with
no numeric values means all types from 1 to \(N\).  A leading asterisk
means all types from 1 to n (inclusive).  A trailing asterisk means all types
from m to \(N\) (inclusive).  A middle asterisk means all types from m to n
(inclusive).

\sphinxAtStartPar
The \sphinxstyleemphasis{types} definition is followed by the viscosity style, \sphinxstyleemphasis{vstyle}. Two
options are available, \sphinxstyleemphasis{constant} and \sphinxstyleemphasis{power}. Style \sphinxstyleemphasis{constant} simply
applies a constant value of the viscosity \sphinxstyleemphasis{eta} to each particle of the
assigned type. Style \sphinxstyleemphasis{power} is a Hershchel\sphinxhyphen{}Bulkley constitutive equation
for the stress \(\tau\)
\begin{equation*}
\begin{split}\tau = \left(\frac{\tau_0}{\dot{\gamma}} + K \dot{\gamma}^{n - 1}\right) \dot{\gamma}, \tau \ge \tau_0\end{split}
\end{equation*}
\sphinxAtStartPar
where \(\dot{\gamma}\) is the strain rate and \(\tau_0\) is the critical
yield stress, below which \(\dot{\gamma} = 0.0\). To avoid divergences, this
expression is regularized by defining a critical strain rate \sphinxstyleemphasis{gd0}. If the local
strain rate on a particle falls below this limit, a constant viscosity of \sphinxstyleemphasis{eta}
is assigned. This implies a value of
\begin{equation*}
\begin{split}\tau_0 = \eta \dot{\gamma}_0 - K \dot{\gamma}_0^N\end{split}
\end{equation*}

\subsection{Restart, fix\_modify, output, run start/stop, minimize info}
\label{\detokenize{fix_rheo_viscosity:restart-fix-modify-output-run-start-stop-minimize-info}}
\sphinxAtStartPar
No information about this fix is written to {\hyperref[\detokenize{restart::doc}]{\sphinxcrossref{\DUrole{doc}{binary restart files}}}}.
None of the {\hyperref[\detokenize{fix_modify::doc}]{\sphinxcrossref{\DUrole{doc}{fix\_modify}}}} options
are relevant to this fix.  No global or per\sphinxhyphen{}atom quantities are stored
by this fix for access by various {\hyperref[\detokenize{Howto_output::doc}]{\sphinxcrossref{\DUrole{doc}{output commands}}}}.
No parameter of this fix can be used with the \sphinxstyleemphasis{start/stop} keywords of
the {\hyperref[\detokenize{run::doc}]{\sphinxcrossref{\DUrole{doc}{run}}}} command.  This fix is not invoked during
{\hyperref[\detokenize{minimize::doc}]{\sphinxcrossref{\DUrole{doc}{energy minimization}}}}.


\subsection{Restrictions}
\label{\detokenize{fix_rheo_viscosity:restrictions}}
\sphinxAtStartPar
This fix must be used with an atom style that includes viscosity
such as atom\_style rheo or rheo/thermal. This fix must be used in
conjunction with {\hyperref[\detokenize{fix_rheo::doc}]{\sphinxcrossref{\DUrole{doc}{fix rheo}}}}. The fix group must be
set to all. Only one instance of fix rheo/viscosity can be defined.

\sphinxAtStartPar
This fix is part of the RHEO package.  It is only enabled if
LAMMPS was built with that package.  See the
{\hyperref[\detokenize{Build_package::doc}]{\sphinxcrossref{\DUrole{doc}{Build package}}}} page for more info.


\subsection{Related commands}
\label{\detokenize{fix_rheo_viscosity:related-commands}}
\sphinxAtStartPar
{\hyperref[\detokenize{fix_rheo::doc}]{\sphinxcrossref{\DUrole{doc}{fix rheo}}}},
{\hyperref[\detokenize{pair_rheo::doc}]{\sphinxcrossref{\DUrole{doc}{pair rheo}}}},
{\hyperref[\detokenize{compute_rheo_property_atom::doc}]{\sphinxcrossref{\DUrole{doc}{compute rheo/property/atom}}}}


\subsection{Default}
\label{\detokenize{fix_rheo_viscosity:default}}
\sphinxAtStartPar
none

\sphinxstepscope

\index{fix rhok@\spxentry{fix rhok}}\ignorespaces 

\section{fix rhok command}
\label{\detokenize{fix_rhok:fix-rhok-command}}\label{\detokenize{fix_rhok:index-0}}\label{\detokenize{fix_rhok::doc}}

\subsection{Syntax}
\label{\detokenize{fix_rhok:syntax}}
\begin{sphinxVerbatim}[commandchars=\\\{\}]
\PYG{k}{fix }\PYG{n+nv+nvIdentifier}{ID}\PYG{+w}{ }\PYG{n+nv+nvIdentifier}{group\PYGZhy{}ID}\PYG{+w}{ }\PYG{n}{rhok}\PYG{+w}{ }\PYG{n}{nx}\PYG{+w}{ }\PYG{n}{ny}\PYG{+w}{ }\PYG{n}{nz}\PYG{+w}{ }\PYG{n}{K}\PYG{+w}{ }\PYG{n}{a}
\end{sphinxVerbatim}
\begin{itemize}
\item {} 
\sphinxAtStartPar
ID, group\sphinxhyphen{}ID are documented in {\hyperref[\detokenize{fix::doc}]{\sphinxcrossref{\DUrole{doc}{fix}}}} command

\item {} 
\sphinxAtStartPar
nx, ny, nz = k\sphinxhyphen{}vector of collective density field

\item {} 
\sphinxAtStartPar
K = spring constant of bias potential

\item {} 
\sphinxAtStartPar
a = anchor point of bias potential

\end{itemize}


\subsection{Examples}
\label{\detokenize{fix_rhok:examples}}
\begin{sphinxVerbatim}[commandchars=\\\{\}]
\PYG{k}{fix }\PYG{n+nv+nvIdentifier}{bias}\PYG{+w}{ }\PYG{n+nv+nvIdentifier}{all}\PYG{+w}{ }\PYG{n}{rhok}\PYG{+w}{ }\PYG{l+m}{16}\PYG{+w}{ }\PYG{l+m}{0}\PYG{+w}{ }\PYG{l+m}{0}\PYG{+w}{ }\PYG{l+m}{4.0}\PYG{+w}{ }\PYG{l+m}{16.0}
\PYG{k}{fix }\PYG{n+nv+nvIdentifier}{1}\PYG{+w}{ }\PYG{n+nv+nvIdentifier}{all}\PYG{+w}{ }\PYG{n}{npt}\PYG{+w}{ }\PYG{n}{temp}\PYG{+w}{ }\PYG{l+m}{0.8}\PYG{+w}{ }\PYG{l+m}{0.8}\PYG{+w}{ }\PYG{l+m}{4.0}\PYG{+w}{ }\PYG{n}{z}\PYG{+w}{ }\PYG{l+m}{2.2}\PYG{+w}{ }\PYG{l+m}{2.2}\PYG{+w}{ }\PYG{l+m}{8.0}
\PYG{c}{\PYGZsh{} output of 4 values from fix rhok: U\PYGZus{}bias rho\PYGZus{}k\PYGZus{}RE  rho\PYGZus{}k\PYGZus{}IM  \PYGZbs{}|rho\PYGZus{}k\PYGZbs{}|}
\PYG{k}{thermo\PYGZus{}style}\PYG{+w}{ }\PYG{n}{custom}\PYG{+w}{ }\PYG{n}{step}\PYG{+w}{ }\PYG{n}{temp}\PYG{+w}{ }\PYG{n}{pzz}\PYG{+w}{ }\PYG{n}{lz}\PYG{+w}{ }\PYG{n}{f\PYGZus{}bias}\PYG{+w}{ }\PYG{n}{f\PYGZus{}bias[1]}\PYG{+w}{ }\PYG{n}{f\PYGZus{}bias[2]}\PYG{+w}{ }\PYG{n}{f\PYGZus{}bias[3]}
\end{sphinxVerbatim}


\subsection{Description}
\label{\detokenize{fix_rhok:description}}
\sphinxAtStartPar
The fix applies a force to atoms given by the potential
\begin{equation*}
\begin{split}U  = &  \frac{1}{2} K (|\rho_{\vec{k}}| - a)^2 \\
\rho_{\vec{k}}  = & \sum_j^N \exp(-i\vec{k} \cdot \vec{r}_j )/\sqrt{N} \\
\vec{k}  = & (2\pi n_x /L_x , 2\pi n_y  /L_y , 2\pi n_z/L_z )\end{split}
\end{equation*}
\sphinxAtStartPar
as described in {\hyperref[\detokenize{fix_rhok:pedersen}]{\sphinxcrossref{\DUrole{std,std-ref}{(Pedersen)}}}}.

\sphinxAtStartPar
This field, which biases configurations with long\sphinxhyphen{}range order, can be
used to study crystal\sphinxhyphen{}liquid interfaces and determine melting
temperatures {\hyperref[\detokenize{fix_rhok:pedersen}]{\sphinxcrossref{\DUrole{std,std-ref}{(Pedersen)}}}}.

\sphinxAtStartPar
An example of using the interface pinning method is located in the
\sphinxstyleemphasis{examples/PACKAGES/rhok} directory.


\subsection{Restart, fix\_modify, output, run start/stop, minimize info}
\label{\detokenize{fix_rhok:restart-fix-modify-output-run-start-stop-minimize-info}}
\sphinxAtStartPar
No information about this fix is written to {\hyperref[\detokenize{restart::doc}]{\sphinxcrossref{\DUrole{doc}{binary restart files}}}}.

\sphinxAtStartPar
The {\hyperref[\detokenize{fix_modify::doc}]{\sphinxcrossref{\DUrole{doc}{fix\_modify}}}} \sphinxstyleemphasis{energy} option is supported by
this fix to add the potential energy calculated by the fix to the
global potential energy of the system as part of {\hyperref[\detokenize{thermo_style::doc}]{\sphinxcrossref{\DUrole{doc}{thermodynamic
output}}}}.  The default setting for this fix is
{\hyperref[\detokenize{fix_modify::doc}]{\sphinxcrossref{\DUrole{doc}{fix\_modify energy no}}}}.

\sphinxAtStartPar
This fix computes a global scalar which can be accessed by various
{\hyperref[\detokenize{Howto_output::doc}]{\sphinxcrossref{\DUrole{doc}{output commands}}}}.  The scalar is the potential
energy discussed in the preceding paragraph.  The scalar stored by
this fix is “extensive”.

\sphinxAtStartPar
No parameter of this fix can be used with the \sphinxstyleemphasis{start/stop} keywords of
the {\hyperref[\detokenize{run::doc}]{\sphinxcrossref{\DUrole{doc}{run}}}} command.

\sphinxAtStartPar
This fix is not invoked during {\hyperref[\detokenize{minimize::doc}]{\sphinxcrossref{\DUrole{doc}{energy minimization}}}}.


\subsection{Restrictions}
\label{\detokenize{fix_rhok:restrictions}}
\sphinxAtStartPar
This fix is part of the EXTRA\sphinxhyphen{}FIX package.  It is only enabled if
LAMMPS was built with that package.  See the {\hyperref[\detokenize{Build_package::doc}]{\sphinxcrossref{\DUrole{doc}{Build package}}}} page for more info.


\subsection{Related commands}
\label{\detokenize{fix_rhok:related-commands}}
\sphinxAtStartPar
{\hyperref[\detokenize{thermo_style::doc}]{\sphinxcrossref{\DUrole{doc}{thermo\_style}}}}


\subsection{Default}
\label{\detokenize{fix_rhok:default}}
\sphinxAtStartPar
none


\bigskip\hrule\bigskip

\phantomsection\label{\detokenize{fix_rhok:pedersen}}
\sphinxAtStartPar
\sphinxstylestrong{(Pedersen)} Pedersen, J. Chem. Phys., 139, 104102 (2013).

\sphinxstepscope

\index{fix rigid@\spxentry{fix rigid}}\ignorespaces 
\index{fix rigid/omp@\spxentry{fix rigid/omp}}\ignorespaces 
\index{fix rigid/nve@\spxentry{fix rigid/nve}}\ignorespaces 
\index{fix rigid/nve/omp@\spxentry{fix rigid/nve/omp}}\ignorespaces 
\index{fix rigid/nvt@\spxentry{fix rigid/nvt}}\ignorespaces 
\index{fix rigid/nvt/omp@\spxentry{fix rigid/nvt/omp}}\ignorespaces 
\index{fix rigid/npt@\spxentry{fix rigid/npt}}\ignorespaces 
\index{fix rigid/npt/omp@\spxentry{fix rigid/npt/omp}}\ignorespaces 
\index{fix rigid/nph@\spxentry{fix rigid/nph}}\ignorespaces 
\index{fix rigid/nph/omp@\spxentry{fix rigid/nph/omp}}\ignorespaces 
\index{fix rigid/small@\spxentry{fix rigid/small}}\ignorespaces 
\index{fix rigid/small/omp@\spxentry{fix rigid/small/omp}}\ignorespaces 
\index{fix rigid/nve/small@\spxentry{fix rigid/nve/small}}\ignorespaces 
\index{fix rigid/nvt/small@\spxentry{fix rigid/nvt/small}}\ignorespaces 
\index{fix rigid/npt/small@\spxentry{fix rigid/npt/small}}\ignorespaces 
\index{fix rigid/nph/small@\spxentry{fix rigid/nph/small}}\ignorespaces 

\section{fix rigid command}
\label{\detokenize{fix_rigid:fix-rigid-command}}\label{\detokenize{fix_rigid:index-15}}\label{\detokenize{fix_rigid:index-14}}\label{\detokenize{fix_rigid:index-13}}\label{\detokenize{fix_rigid:index-12}}\label{\detokenize{fix_rigid:index-11}}\label{\detokenize{fix_rigid:index-10}}\label{\detokenize{fix_rigid:index-9}}\label{\detokenize{fix_rigid:index-8}}\label{\detokenize{fix_rigid:index-7}}\label{\detokenize{fix_rigid:index-6}}\label{\detokenize{fix_rigid:index-5}}\label{\detokenize{fix_rigid:index-4}}\label{\detokenize{fix_rigid:index-3}}\label{\detokenize{fix_rigid:index-2}}\label{\detokenize{fix_rigid:index-1}}\label{\detokenize{fix_rigid:index-0}}\label{\detokenize{fix_rigid::doc}}
\sphinxAtStartPar
Accelerator Variants: \sphinxstyleemphasis{rigid/omp}


\section{fix rigid/nve command}
\label{\detokenize{fix_rigid:fix-rigid-nve-command}}
\sphinxAtStartPar
Accelerator Variants: \sphinxstyleemphasis{rigid/nve/omp}


\section{fix rigid/nvt command}
\label{\detokenize{fix_rigid:fix-rigid-nvt-command}}
\sphinxAtStartPar
Accelerator Variants: \sphinxstyleemphasis{rigid/nvt/omp}


\section{fix rigid/npt command}
\label{\detokenize{fix_rigid:fix-rigid-npt-command}}
\sphinxAtStartPar
Accelerator Variants: \sphinxstyleemphasis{rigid/npt/omp}


\section{fix rigid/nph command}
\label{\detokenize{fix_rigid:fix-rigid-nph-command}}
\sphinxAtStartPar
Accelerator Variants: \sphinxstyleemphasis{rigid/nph/omp}


\section{fix rigid/small command}
\label{\detokenize{fix_rigid:fix-rigid-small-command}}
\sphinxAtStartPar
Accelerator Variants: \sphinxstyleemphasis{rigid/small/omp}


\section{fix rigid/nve/small command}
\label{\detokenize{fix_rigid:fix-rigid-nve-small-command}}

\section{fix rigid/nvt/small command}
\label{\detokenize{fix_rigid:fix-rigid-nvt-small-command}}

\section{fix rigid/npt/small command}
\label{\detokenize{fix_rigid:fix-rigid-npt-small-command}}

\section{fix rigid/nph/small command}
\label{\detokenize{fix_rigid:fix-rigid-nph-small-command}}

\subsection{Syntax}
\label{\detokenize{fix_rigid:syntax}}
\begin{sphinxVerbatim}[commandchars=\\\{\}]
\PYG{k}{fix }\PYG{n+nv+nvIdentifier}{ID}\PYG{+w}{ }\PYG{n+nv+nvIdentifier}{group\PYGZhy{}ID}\PYG{+w}{ }\PYG{n}{style}\PYG{+w}{ }\PYG{n}{bodystyle}\PYG{+w}{ }\PYG{n}{args}\PYG{+w}{ }\PYG{n}{keyword}\PYG{+w}{ }\PYG{n}{values}\PYG{+w}{ }\PYG{n}{...}
\end{sphinxVerbatim}
\begin{itemize}
\item {} 
\sphinxAtStartPar
ID, group\sphinxhyphen{}ID are documented in {\hyperref[\detokenize{fix::doc}]{\sphinxcrossref{\DUrole{doc}{fix}}}} command

\item {} 
\sphinxAtStartPar
style = \sphinxstyleemphasis{rigid} or \sphinxstyleemphasis{rigid/nve} or \sphinxstyleemphasis{rigid/nvt} or \sphinxstyleemphasis{rigid/npt} or \sphinxstyleemphasis{rigid/nph} or \sphinxstyleemphasis{rigid/small} or \sphinxstyleemphasis{rigid/nve/small} or \sphinxstyleemphasis{rigid/nvt/small} or \sphinxstyleemphasis{rigid/npt/small} or \sphinxstyleemphasis{rigid/nph/small}

\item {} 
\sphinxAtStartPar
bodystyle = \sphinxstyleemphasis{single} or \sphinxstyleemphasis{molecule} or \sphinxstyleemphasis{group}
\begin{sphinxalltt}
\sphinxstyleemphasis{single} args = none
\sphinxstyleemphasis{molecule} args = none
\sphinxstyleemphasis{custom} args = \sphinxstyleemphasis{i\_propname} or \sphinxstyleemphasis{v\_varname}
  i\_propname = a custom integer vector defined via fix property/atom
  v\_varname  = an atom\sphinxhyphen{}style or atomfile\sphinxhyphen{}style variable
\sphinxstyleemphasis{group} args = N groupID1 groupID2 ...
  N = \# of groups
  groupID1, groupID2, ... = list of N group IDs
\end{sphinxalltt}

\item {} 
\sphinxAtStartPar
zero or more keyword/value pairs may be appended

\item {} 
\sphinxAtStartPar
keyword = \sphinxstyleemphasis{langevin} or \sphinxstyleemphasis{reinit} or \sphinxstyleemphasis{temp}  or \sphinxstyleemphasis{mol} or \sphinxstyleemphasis{iso} or \sphinxstyleemphasis{aniso} or \sphinxstyleemphasis{x} or \sphinxstyleemphasis{y} or \sphinxstyleemphasis{z} or \sphinxstyleemphasis{couple} or \sphinxstyleemphasis{tparam} or \sphinxstyleemphasis{pchain} or \sphinxstyleemphasis{dilate} or \sphinxstyleemphasis{force} or \sphinxstyleemphasis{torque} or \sphinxstyleemphasis{infile} or \sphinxstyleemphasis{gravity}
\begin{sphinxalltt}
\sphinxstyleemphasis{langevin} values = Tstart Tstop Tperiod seed
  Tstart,Tstop = desired temperature at start/stop of run (temperature units)
  Tdamp = temperature damping parameter (time units)
  seed = random number seed to use for white noise (positive integer)
\sphinxstyleemphasis{reinit} value = \sphinxstyleemphasis{yes} or \sphinxstyleemphasis{no}
\sphinxstyleemphasis{temp} values = Tstart Tstop Tdamp
  Tstart,Tstop = desired temperature at start/stop of run (temperature units)
  Tdamp = temperature damping parameter (time units)
\sphinxstyleemphasis{mol} value = template\sphinxhyphen{}ID
  template\sphinxhyphen{}ID = ID of molecule template specified in a separate {\hyperref[\detokenize{molecule::doc}]{\sphinxcrossref{\DUrole{doc}{molecule}}}} command
\sphinxstyleemphasis{iso} or \sphinxstyleemphasis{aniso} values = Pstart Pstop Pdamp
  Pstart,Pstop = scalar external pressure at start/end of run (pressure units)
  Pdamp = pressure damping parameter (time units)
\sphinxstyleemphasis{x} or \sphinxstyleemphasis{y} or \sphinxstyleemphasis{z} values = Pstart Pstop Pdamp
  Pstart,Pstop = external stress tensor component at start/end of run (pressure units)
  Pdamp = stress damping parameter (time units)
\sphinxstyleemphasis{couple} value = \sphinxstyleemphasis{none} or \sphinxstyleemphasis{xyz} or \sphinxstyleemphasis{xy} or \sphinxstyleemphasis{yz} or \sphinxstyleemphasis{xz}
\sphinxstyleemphasis{tparam} values = Tchain Titer Torder
  Tchain = length of Nose/Hoover thermostat chain
  Titer = number of thermostat iterations performed
  Torder = 3 or 5 = Yoshida\sphinxhyphen{}Suzuki integration parameters
\sphinxstyleemphasis{pchain} values = Pchain
  Pchain = length of the Nose/Hoover thermostat chain coupled with the barostat
\sphinxstyleemphasis{dilate} value = dilate\sphinxhyphen{}group\sphinxhyphen{}ID
  dilate\sphinxhyphen{}group\sphinxhyphen{}ID = only dilate atoms in this group due to barostat volume changes
\sphinxstyleemphasis{force} values = M xflag yflag zflag
  M = which rigid body from 1\sphinxhyphen{}Nbody (see asterisk form below)
  xflag,yflag,zflag = off/on if component of center\sphinxhyphen{}of\sphinxhyphen{}mass force is active
\sphinxstyleemphasis{torque} values = M xflag yflag zflag
  M = which rigid body from 1\sphinxhyphen{}Nbody (see asterisk form below)
  xflag,yflag,zflag = off/on if component of center\sphinxhyphen{}of\sphinxhyphen{}mass torque is active
\sphinxstyleemphasis{infile} filename
  filename = file with per\sphinxhyphen{}body values of mass, center\sphinxhyphen{}of\sphinxhyphen{}mass, moments of inertia
\sphinxstyleemphasis{gravity} values = gravity\sphinxhyphen{}ID
  gravity\sphinxhyphen{}ID = ID of fix gravity command to add gravitational forces
\end{sphinxalltt}

\end{itemize}


\subsection{Examples}
\label{\detokenize{fix_rigid:examples}}
\begin{sphinxVerbatim}[commandchars=\\\{\}]
\PYG{k}{fix }\PYG{n+nv+nvIdentifier}{1}\PYG{+w}{ }\PYG{n+nv+nvIdentifier}{clump}\PYG{+w}{ }\PYG{n}{rigid}\PYG{+w}{ }\PYG{n}{single}\PYG{+w}{ }\PYG{n}{reinit}\PYG{+w}{ }\PYG{n}{yes}
\PYG{k}{fix }\PYG{n+nv+nvIdentifier}{1}\PYG{+w}{ }\PYG{n+nv+nvIdentifier}{clump}\PYG{+w}{ }\PYG{n}{rigid}\PYG{o}{/}\PYG{n}{small}\PYG{+w}{ }\PYG{n}{molecule}
\PYG{k}{fix }\PYG{n+nv+nvIdentifier}{1}\PYG{+w}{ }\PYG{n+nv+nvIdentifier}{clump}\PYG{+w}{ }\PYG{n}{rigid}\PYG{+w}{ }\PYG{n}{single}\PYG{+w}{ }\PYG{n}{force}\PYG{+w}{ }\PYG{l+m}{1}\PYG{+w}{ }\PYG{n}{off}\PYG{+w}{ }\PYG{n}{off}\PYG{+w}{ }\PYG{n}{on}\PYG{+w}{ }\PYG{n}{langevin}\PYG{+w}{ }\PYG{l+m}{1.0}\PYG{+w}{ }\PYG{l+m}{1.0}\PYG{+w}{ }\PYG{l+m}{1.0}\PYG{+w}{ }\PYG{l+m}{428984}
\PYG{k}{fix }\PYG{n+nv+nvIdentifier}{1}\PYG{+w}{ }\PYG{n+nv+nvIdentifier}{polychains}\PYG{+w}{ }\PYG{n}{rigid}\PYG{o}{/}\PYG{n}{nvt}\PYG{+w}{ }\PYG{n}{molecule}\PYG{+w}{ }\PYG{n}{temp}\PYG{+w}{ }\PYG{l+m}{1.0}\PYG{+w}{ }\PYG{l+m}{1.0}\PYG{+w}{ }\PYG{l+m}{5.0}\PYG{+w}{ }\PYG{n}{reinit}\PYG{+w}{ }\PYG{n}{no}
\PYG{k}{fix }\PYG{n+nv+nvIdentifier}{1}\PYG{+w}{ }\PYG{n+nv+nvIdentifier}{polychains}\PYG{+w}{ }\PYG{n}{rigid}\PYG{+w}{ }\PYG{n}{molecule}\PYG{+w}{ }\PYG{n}{force}\PYG{+w}{ }\PYG{l+m}{1}\PYG{o}{*}\PYG{l+m}{5}\PYG{+w}{ }\PYG{n}{off}\PYG{+w}{ }\PYG{n}{off}\PYG{+w}{ }\PYG{n}{off}\PYG{+w}{ }\PYG{n}{force}\PYG{+w}{ }\PYG{l+m}{6}\PYG{o}{*}\PYG{l+m}{10}\PYG{+w}{ }\PYG{n}{off}\PYG{+w}{ }\PYG{n}{off}\PYG{+w}{ }\PYG{n}{on}
\PYG{k}{fix }\PYG{n+nv+nvIdentifier}{1}\PYG{+w}{ }\PYG{n+nv+nvIdentifier}{polychains}\PYG{+w}{ }\PYG{n}{rigid}\PYG{o}{/}\PYG{n}{small}\PYG{+w}{ }\PYG{n}{molecule}\PYG{+w}{ }\PYG{n}{langevin}\PYG{+w}{ }\PYG{l+m}{1.0}\PYG{+w}{ }\PYG{l+m}{1.0}\PYG{+w}{ }\PYG{l+m}{1.0}\PYG{+w}{ }\PYG{l+m}{428984}
\PYG{k}{fix }\PYG{n+nv+nvIdentifier}{2}\PYG{+w}{ }\PYG{n+nv+nvIdentifier}{fluid}\PYG{+w}{ }\PYG{n}{rigid}\PYG{+w}{ }\PYG{k}{group }\PYG{n+nv+nvIdentifier}{3}\PYG{+w}{ }\PYG{n}{clump1}\PYG{+w}{ }\PYG{n}{clump2}\PYG{+w}{ }\PYG{n}{clump3}\PYG{+w}{ }\PYG{n}{torque}\PYG{+w}{ }\PYG{o}{*}\PYG{+w}{ }\PYG{n}{off}\PYG{+w}{ }\PYG{n}{off}\PYG{+w}{ }\PYG{n}{off}
\PYG{k}{fix }\PYG{n+nv+nvIdentifier}{1}\PYG{+w}{ }\PYG{n+nv+nvIdentifier}{rods}\PYG{+w}{ }\PYG{n}{rigid}\PYG{o}{/}\PYG{n}{npt}\PYG{+w}{ }\PYG{n}{molecule}\PYG{+w}{ }\PYG{n}{temp}\PYG{+w}{ }\PYG{l+m}{300.0}\PYG{+w}{ }\PYG{l+m}{300.0}\PYG{+w}{ }\PYG{l+m}{100.0}\PYG{+w}{ }\PYG{n}{iso}\PYG{+w}{ }\PYG{l+m}{0.5}\PYG{+w}{ }\PYG{l+m}{0.5}\PYG{+w}{ }\PYG{l+m}{10.0}
\PYG{k}{fix }\PYG{n+nv+nvIdentifier}{1}\PYG{+w}{ }\PYG{n+nv+nvIdentifier}{particles}\PYG{+w}{ }\PYG{n}{rigid}\PYG{o}{/}\PYG{n}{npt}\PYG{+w}{ }\PYG{n}{molecule}\PYG{+w}{ }\PYG{n}{temp}\PYG{+w}{ }\PYG{l+m}{1.0}\PYG{+w}{ }\PYG{l+m}{1.0}\PYG{+w}{ }\PYG{l+m}{5.0}\PYG{+w}{ }\PYG{n}{x}\PYG{+w}{ }\PYG{l+m}{0.5}\PYG{+w}{ }\PYG{l+m}{0.5}\PYG{+w}{ }\PYG{l+m}{1.0}\PYG{+w}{ }\PYG{n}{z}\PYG{+w}{ }\PYG{l+m}{0.5}\PYG{+w}{ }\PYG{l+m}{0.5}\PYG{+w}{ }\PYG{l+m}{1.0}\PYG{+w}{ }\PYG{n}{couple}\PYG{+w}{ }\PYG{n}{xz}
\PYG{k}{fix }\PYG{n+nv+nvIdentifier}{1}\PYG{+w}{ }\PYG{n+nv+nvIdentifier}{water}\PYG{+w}{ }\PYG{n}{rigid}\PYG{o}{/}\PYG{n}{nph}\PYG{+w}{ }\PYG{n}{molecule}\PYG{+w}{ }\PYG{n}{iso}\PYG{+w}{ }\PYG{l+m}{0.5}\PYG{+w}{ }\PYG{l+m}{0.5}\PYG{+w}{ }\PYG{l+m}{1.0}
\PYG{k}{fix }\PYG{n+nv+nvIdentifier}{1}\PYG{+w}{ }\PYG{n+nv+nvIdentifier}{particles}\PYG{+w}{ }\PYG{n}{rigid}\PYG{o}{/}\PYG{n}{npt}\PYG{o}{/}\PYG{n}{small}\PYG{+w}{ }\PYG{n}{molecule}\PYG{+w}{ }\PYG{n}{temp}\PYG{+w}{ }\PYG{l+m}{1.0}\PYG{+w}{ }\PYG{l+m}{1.0}\PYG{+w}{ }\PYG{l+m}{1.0}\PYG{+w}{ }\PYG{n}{iso}\PYG{+w}{ }\PYG{l+m}{0.5}\PYG{+w}{ }\PYG{l+m}{0.5}\PYG{+w}{ }\PYG{l+m}{1.0}

\PYG{k}{variable }\PYG{n+nv+nvIdentifier}{bodyid}\PYG{+w}{ }\PYG{n}{atom}\PYG{+w}{ }\PYG{l+m}{1.0}\PYG{o}{*}\PYG{n}{gmask}\PYG{n+nv}{(}\PYG{n+nv}{clump1}\PYG{n+nv}{)}\PYG{o}{+}\PYG{l+m}{2.0}\PYG{o}{*}\PYG{n}{gmask}\PYG{n+nv}{(}\PYG{n+nv}{clump2}\PYG{n+nv}{)}\PYG{o}{+}\PYG{l+m}{3.0}\PYG{o}{*}\PYG{n}{gmask}\PYG{n+nv}{(}\PYG{n+nv}{clump3}\PYG{n+nv}{)}
\PYG{k}{fix }\PYG{n+nv+nvIdentifier}{1}\PYG{+w}{ }\PYG{n+nv+nvIdentifier}{clump}\PYG{+w}{ }\PYG{n}{rigid}\PYG{+w}{ }\PYG{n}{custom}\PYG{+w}{ }\PYG{n}{v\PYGZus{}bodyid}

\PYG{k}{variable }\PYG{n+nv+nvIdentifier}{bodyid}\PYG{+w}{ }\PYG{n}{atomfile}\PYG{+w}{ }\PYG{n}{bodies.txt}
\PYG{k}{fix }\PYG{n+nv+nvIdentifier}{1}\PYG{+w}{ }\PYG{n+nv+nvIdentifier}{clump}\PYG{+w}{ }\PYG{n}{rigid}\PYG{+w}{ }\PYG{n}{custom}\PYG{+w}{ }\PYG{n}{v\PYGZus{}bodyid}

\PYG{k}{fix }\PYG{n+nv+nvIdentifier}{0}\PYG{+w}{ }\PYG{n+nv+nvIdentifier}{all}\PYG{+w}{ }\PYG{n}{property}\PYG{o}{/}\PYG{n}{atom}\PYG{+w}{ }\PYG{n}{i\PYGZus{}bodyid}
\PYG{k}{read\PYGZus{}restart}\PYG{+w}{ }\PYG{n}{data.rigid}\PYG{+w}{ }\PYG{k}{fix }\PYG{n+nv+nvIdentifier}{0}\PYG{+w}{ }\PYG{n+nv+nvIdentifier}{NULL}\PYG{+w}{ }\PYG{n}{Bodies}
\PYG{k}{fix }\PYG{n+nv+nvIdentifier}{1}\PYG{+w}{ }\PYG{n+nv+nvIdentifier}{clump}\PYG{+w}{ }\PYG{n}{rigid}\PYG{o}{/}\PYG{n}{small}\PYG{+w}{ }\PYG{n}{custom}\PYG{+w}{ }\PYG{n}{i\PYGZus{}bodyid}
\end{sphinxVerbatim}


\subsection{Description}
\label{\detokenize{fix_rigid:description}}
\sphinxAtStartPar
Treat one or more sets of atoms as independent rigid bodies.  This
means that each timestep the total force and torque on each rigid body
is computed as the sum of the forces and torques on its constituent
particles.  The coordinates, velocities, and orientations of the atoms
in each body are then updated so that the body moves and rotates as a
single entity.  This is implemented by creating internal data structures
for each rigid body and performing time integration on these data
structures.  Positions, velocities, and orientations of the constituent
particles are regenerated from the rigid body data structures in every
time step. This restricts which operations and fixes can be applied to
rigid bodies. See below for a detailed discussion.

\sphinxAtStartPar
Examples of large rigid bodies are a colloidal particle, or portions
of a biomolecule such as a protein.

\sphinxAtStartPar
Example of small rigid bodies are patchy nanoparticles, such as those
modeled in {\hyperref[\detokenize{fix_rigid:zhang1}]{\sphinxcrossref{\DUrole{std,std-ref}{this paper}}}} by Sharon Glotzer’s group,
clumps of granular particles, lipid molecules consisting of one or
more point dipoles connected to other spheroids or ellipsoids,
irregular particles built from line segments (2d) or triangles (3d),
and coarse\sphinxhyphen{}grain models of nano or colloidal particles consisting of a
small number of constituent particles.  Note that the {\hyperref[\detokenize{fix_shake::doc}]{\sphinxcrossref{\DUrole{doc}{fix shake}}}} command can also be used to rigidify small molecules of
2, 3, or 4 atoms, e.g. water molecules.  That fix treats the
constituent atoms as point masses.

\sphinxAtStartPar
These fixes also update the positions and velocities of the atoms in
each rigid body via time integration, in the NVE, NVT, NPT, or NPH
ensemble, as described below.

\sphinxAtStartPar
There are two main variants of this fix, fix rigid and fix
rigid/small.  The NVE/NVT/NPT/NHT versions belong to one of the two
variants, as their style names indicate.

\begin{sphinxadmonition}{note}{Note:}
\sphinxAtStartPar
Not all of the \sphinxstyleemphasis{bodystyle} options and keyword/value options are
available for both the \sphinxstyleemphasis{rigid} and \sphinxstyleemphasis{rigid/small} variants.  See
details below.
\end{sphinxadmonition}

\sphinxAtStartPar
The \sphinxstyleemphasis{rigid} styles are typically the best choice for a system with a
small number of large rigid bodies, each of which can extend across
the domain of many processors.  It operates by creating a single
global list of rigid bodies, which all processors contribute to.
MPI\_Allreduce operations are performed each timestep to sum the
contributions from each processor to the force and torque on all the
bodies.  This operation will not scale well in parallel if large
numbers of rigid bodies are simulated.

\sphinxAtStartPar
The \sphinxstyleemphasis{rigid/small} styles are typically best for a system with a large
number of small rigid bodies.  Each body is assigned to the atom
closest to the geometrical center of the body.  The fix operates using
local lists of rigid bodies owned by each processor and information is
exchanged and summed via local communication between neighboring
processors when ghost atom info is accumulated.

\begin{sphinxadmonition}{note}{Note:}
\sphinxAtStartPar
To use the \sphinxstyleemphasis{rigid/small} styles the ghost atom cutoff must be large
enough to span the distance between the atom that owns the body and
every other atom in the body.  This distance value is printed out
when the rigid bodies are defined.  If the {\hyperref[\detokenize{pair_style::doc}]{\sphinxcrossref{\DUrole{doc}{pair\_style}}}} cutoff plus neighbor skin does not span this
distance, then you should use the {\hyperref[\detokenize{comm_modify::doc}]{\sphinxcrossref{\DUrole{doc}{comm\_modify cutoff}}}} command with a setting epsilon larger than the
distance.
\end{sphinxadmonition}

\sphinxAtStartPar
Which of the two variants is faster for a particular problem is hard
to predict.  The best way to decide is to perform a short test run.
Both variants should give identical numerical answers for short runs.
Long runs should give statistically similar results, but round\sphinxhyphen{}off
differences may accumulate to produce divergent trajectories.

\begin{sphinxadmonition}{note}{Note:}
\sphinxAtStartPar
You should not update the atoms in rigid bodies via other
time\sphinxhyphen{}integration fixes (e.g. {\hyperref[\detokenize{fix_nve::doc}]{\sphinxcrossref{\DUrole{doc}{fix nve}}}}, {\hyperref[\detokenize{fix_nh::doc}]{\sphinxcrossref{\DUrole{doc}{fix
nvt}}}}, {\hyperref[\detokenize{fix_nh::doc}]{\sphinxcrossref{\DUrole{doc}{fix npt}}}}, {\hyperref[\detokenize{fix_move::doc}]{\sphinxcrossref{\DUrole{doc}{fix move}}}}), or you will have conflicting updates to positions and
velocities resulting in unphysical behavior in most cases. When
performing a hybrid simulation with some atoms in rigid bodies, and
some not, a separate time integration fix like {\hyperref[\detokenize{fix_nve::doc}]{\sphinxcrossref{\DUrole{doc}{fix nve}}}} or {\hyperref[\detokenize{fix_nh::doc}]{\sphinxcrossref{\DUrole{doc}{fix nvt}}}} should be used for the
non\sphinxhyphen{}rigid particles.
\end{sphinxadmonition}

\begin{sphinxadmonition}{note}{Note:}
\sphinxAtStartPar
These fixes are overkill if you simply want to hold a collection of
atoms stationary or have them move with a constant velocity.  A
simpler way to hold atoms stationary is to not include those atoms
in your time integration fix.  E.g. use “fix 1 mobile nve” instead
of “fix 1 all nve”, where “mobile” is the group of atoms that you
want to move.  You can move atoms with a constant velocity by
assigning them an initial velocity (via the {\hyperref[\detokenize{velocity::doc}]{\sphinxcrossref{\DUrole{doc}{velocity}}}} command), setting the force on them to 0.0 (via the
{\hyperref[\detokenize{fix_setforce::doc}]{\sphinxcrossref{\DUrole{doc}{fix setforce}}}} command), and integrating them
as usual (e.g. via the {\hyperref[\detokenize{fix_nve::doc}]{\sphinxcrossref{\DUrole{doc}{fix nve}}}} command).
\end{sphinxadmonition}

\begin{sphinxadmonition}{warning}{Warning:}
\sphinxAtStartPar
The aggregate properties of each rigid body are calculated at the
start of a simulation run and are maintained in internal data
structures. The properties include the position and velocity of the
center\sphinxhyphen{}of\sphinxhyphen{}mass of the body, its moments of inertia, and its angular
momentum.  This is done using the properties of the constituent
atoms of the body at that point in time (or see the \sphinxstyleemphasis{infile}
keyword option).  Thereafter, changing these properties of
individual atoms in the body will have no effect on a rigid body’s
dynamics, unless they effect any computation of per\sphinxhyphen{}atom forces or
torques. If the keyword \sphinxstyleemphasis{reinit} is set to \sphinxstyleemphasis{yes} (the default), the
rigid body data structures will be recreated at the beginning of
each \sphinxstyleemphasis{run} command; if the keyword \sphinxstyleemphasis{reinit} is set to \sphinxstyleemphasis{no}, the
rigid body data structures will be built only at the very first
\sphinxstyleemphasis{run} command and maintained for as long as the rigid fix is
defined. For example, you might think you could displace the atoms
in a body or add a large velocity to each atom in a body to make it
move in a desired direction before a second run is performed, using
the {\hyperref[\detokenize{set::doc}]{\sphinxcrossref{\DUrole{doc}{set}}}} or {\hyperref[\detokenize{displace_atoms::doc}]{\sphinxcrossref{\DUrole{doc}{displace\_atoms}}}} or
{\hyperref[\detokenize{velocity::doc}]{\sphinxcrossref{\DUrole{doc}{velocity}}}} commands.  But these commands will not
affect the internal attributes of the body unless \sphinxstyleemphasis{reinit} is set
to \sphinxstyleemphasis{yes}. With \sphinxstyleemphasis{reinit} set to \sphinxstyleemphasis{no} (or using the \sphinxstyleemphasis{infile}
option, which implies \sphinxstyleemphasis{reinit} \sphinxstyleemphasis{no}) the position and velocity of
individual atoms in the body will be reset when time integration
starts again.
\end{sphinxadmonition}


\bigskip\hrule\bigskip


\sphinxAtStartPar
Each rigid body must have two or more atoms.  An atom can belong to at
most one rigid body.  Which atoms are in which bodies can be defined
via several options.

\begin{sphinxadmonition}{note}{Note:}
\sphinxAtStartPar
With the \sphinxstyleemphasis{rigid/small} styles, which require that \sphinxstyleemphasis{bodystyle} be
specified as \sphinxstyleemphasis{molecule} or \sphinxstyleemphasis{custom}, you can define a system that has
no rigid bodies initially.  This is useful when you are using the
\sphinxstyleemphasis{mol} keyword in conjunction with another fix that is adding rigid
bodies on\sphinxhyphen{}the\sphinxhyphen{}fly as molecules, such as {\hyperref[\detokenize{fix_deposit::doc}]{\sphinxcrossref{\DUrole{doc}{fix deposit}}}}
or {\hyperref[\detokenize{fix_pour::doc}]{\sphinxcrossref{\DUrole{doc}{fix pour}}}}.
\end{sphinxadmonition}

\sphinxAtStartPar
For bodystyle \sphinxstyleemphasis{single} the entire fix group of atoms is treated as one
rigid body.  This option is only allowed for the \sphinxstyleemphasis{rigid} styles.

\sphinxAtStartPar
For bodystyle \sphinxstyleemphasis{molecule}, atoms are grouped into rigid bodies by their
respective molecule IDs: each set of atoms in the fix group with the
same molecule ID is treated as a different rigid body.  This option is
allowed for both the \sphinxstyleemphasis{rigid} and \sphinxstyleemphasis{rigid/small} styles.  Note that
atoms with a molecule ID = 0 will be treated as a single rigid body.
For a system with atomic solvent (typically this is atoms with
molecule ID = 0) surrounding rigid bodies, this may not be what you
want.  Thus you should be careful to use a fix group that only
includes atoms you want to be part of rigid bodies.

\sphinxAtStartPar
Bodystyle \sphinxstyleemphasis{custom} is similar to bodystyle \sphinxstyleemphasis{molecule} except that it
is more flexible in using other per\sphinxhyphen{}atom properties to define the sets
of atoms that form rigid bodies.  A custom per\sphinxhyphen{}atom integer vector
defined by the {\hyperref[\detokenize{fix_property_atom::doc}]{\sphinxcrossref{\DUrole{doc}{fix property/atom}}}} command
can be used.  Or an {\hyperref[\detokenize{variable::doc}]{\sphinxcrossref{\DUrole{doc}{atom\sphinxhyphen{}style or atomfile\sphinxhyphen{}style variable}}}} can be used; the floating\sphinxhyphen{}point value produced by the
variable is rounded to an integer.  As with bodystyle \sphinxstyleemphasis{molecule},
each set of atoms in the fix groups with the same integer value is
treated as a different rigid body.  Since fix property/atom custom
vectors and atom\sphinxhyphen{}style variables produce values for all atoms, you
should be careful to use a fix group that only includes atoms you want
to be part of rigid bodies.

\begin{sphinxadmonition}{note}{Note:}
\sphinxAtStartPar
To compute the initial center\sphinxhyphen{}of\sphinxhyphen{}mass position and other properties
of each rigid body, the image flags for each atom in the body are
used to “unwrap” the atom coordinates.  Thus you must ensure that
these image flags are consistent so that the unwrapping creates a
valid rigid body (one where the atoms are close together),
particularly if the atoms in a single rigid body straddle a
periodic boundary.  This means the input data file or restart file
must define the image flags for each atom consistently or that you
have used the {\hyperref[\detokenize{set::doc}]{\sphinxcrossref{\DUrole{doc}{set}}}} command to specify them correctly.
If a dimension is non\sphinxhyphen{}periodic then the image flag of each atom
must be 0 in that dimension, else an error is generated.
\end{sphinxadmonition}

\sphinxAtStartPar
The \sphinxstyleemphasis{force} and \sphinxstyleemphasis{torque} keywords discussed next are only allowed for
the \sphinxstyleemphasis{rigid} styles.

\sphinxAtStartPar
By default, each rigid body is acted on by other atoms which induce an
external force and torque on its center of mass, causing it to
translate and rotate.  Components of the external center\sphinxhyphen{}of\sphinxhyphen{}mass force
and torque can be turned off by the \sphinxstyleemphasis{force} and \sphinxstyleemphasis{torque} keywords.
This may be useful if you wish a body to rotate but not translate, or
vice versa, or if you wish it to rotate or translate continuously
unaffected by interactions with other particles.  Note that if you
expect a rigid body not to move or rotate by using these keywords, you
must ensure its initial center\sphinxhyphen{}of\sphinxhyphen{}mass translational or angular
velocity is 0.0.  Otherwise the initial translational or angular
momentum the body has will persist.

\sphinxAtStartPar
An xflag, yflag, or zflag set to \sphinxstyleemphasis{off} means turn off the component of
force of torque in that dimension.  A setting of \sphinxstyleemphasis{on} means turn on
the component, which is the default.  Which rigid body(s) the settings
apply to is determined by the first argument of the \sphinxstyleemphasis{force} and
\sphinxstyleemphasis{torque} keywords.  It can be an integer M from 1 to Nbody, where
Nbody is the number of rigid bodies defined.  A wild\sphinxhyphen{}card asterisk can
be used in place of, or in conjunction with, the M argument to set the
flags for multiple rigid bodies.  This takes the form “*” or “*n” or
“n*” or “m*n”.  If N = the number of rigid bodies, then an asterisk
with no numeric values means all bodies from 1 to N.  A leading
asterisk means all bodies from 1 to n (inclusive).  A trailing
asterisk means all bodies from n to N (inclusive).  A middle asterisk
means all types from m to n (inclusive).  Note that you can use the
\sphinxstyleemphasis{force} or \sphinxstyleemphasis{torque} keywords as many times as you like.  If a
particular rigid body has its component flags set multiple times, the
settings from the final keyword are used.

\begin{sphinxadmonition}{note}{Note:}
\sphinxAtStartPar
For computational efficiency, you may wish to turn off pairwise and
bond interactions within each rigid body, as they no longer
contribute to the motion.  The {\hyperref[\detokenize{neigh_modify::doc}]{\sphinxcrossref{\DUrole{doc}{neigh\_modify exclude}}}} and {\hyperref[\detokenize{delete_bonds::doc}]{\sphinxcrossref{\DUrole{doc}{delete\_bonds}}}} commands are
used to do this.  If the rigid bodies have strongly overlapping
atoms, you may need to turn off these interactions to avoid
numerical problems due to large equal/opposite intra\sphinxhyphen{}body forces
swamping the contribution of small inter\sphinxhyphen{}body forces.
\end{sphinxadmonition}

\sphinxAtStartPar
For computational efficiency, you should typically define one fix
rigid or fix rigid/small command which includes all the desired rigid
bodies.  LAMMPS will allow multiple rigid fixes to be defined, but it
is more expensive.


\bigskip\hrule\bigskip


\sphinxAtStartPar
The constituent particles within a rigid body can be point particles
(the default in LAMMPS) or finite\sphinxhyphen{}size particles, such as spheres or
ellipsoids or line segments or triangles.  See the {\hyperref[\detokenize{atom_style::doc}]{\sphinxcrossref{\DUrole{doc}{atom\_style
sphere and ellipsoid and line and tri}}}} commands for more
details on these kinds of particles.  Finite\sphinxhyphen{}size particles contribute
differently to the moment of inertia of a rigid body than do point
particles.  Finite\sphinxhyphen{}size particles can also experience torque (e.g. due
to {\hyperref[\detokenize{pair_gran::doc}]{\sphinxcrossref{\DUrole{doc}{frictional granular interactions}}}}) and have an
orientation.  These contributions are accounted for by these fixes.

\sphinxAtStartPar
Forces between particles within a body do not contribute to the
external force or torque on the body.  Thus for computational
efficiency, you may wish to turn off pairwise and bond interactions
between particles within each rigid body.  The {\hyperref[\detokenize{neigh_modify::doc}]{\sphinxcrossref{\DUrole{doc}{neigh\_modify
exclude}}}} and {\hyperref[\detokenize{delete_bonds::doc}]{\sphinxcrossref{\DUrole{doc}{delete\_bonds}}}}
commands are used to do this.  For finite\sphinxhyphen{}size particles this also
means the particles can be highly overlapped when creating the rigid
body.


\bigskip\hrule\bigskip


\sphinxAtStartPar
The \sphinxstyleemphasis{rigid}, \sphinxstyleemphasis{rigid/nve}, \sphinxstyleemphasis{rigid/small}, and \sphinxstyleemphasis{rigid/small/nve} styles
perform constant NVE time integration.  They are referred to below as
the 4 NVE rigid styles.  The only difference is that the \sphinxstyleemphasis{rigid} and
\sphinxstyleemphasis{rigid/small} styles use an integration technique based on Richardson
iterations.  The \sphinxstyleemphasis{rigid/nve} and \sphinxstyleemphasis{rigid/small/nve} styles uses the
methods described in the paper by {\hyperref[\detokenize{fix_rigid:miller3}]{\sphinxcrossref{\DUrole{std,std-ref}{Miller}}}}, which are
thought to provide better energy conservation than an iterative
approach.

\sphinxAtStartPar
The \sphinxstyleemphasis{rigid/nvt} and \sphinxstyleemphasis{rigid/nvt/small} styles performs constant NVT
integration using a Nose/Hoover thermostat with chains as described
originally in {\hyperref[\detokenize{fix_rigid:hoover}]{\sphinxcrossref{\DUrole{std,std-ref}{(Hoover)}}}} and {\hyperref[\detokenize{fix_rigid:martyna2}]{\sphinxcrossref{\DUrole{std,std-ref}{(Martyna)}}}}, which thermostats both the translational and rotational
degrees of freedom of the rigid bodies.  They are referred to below as
the 2 NVT rigid styles.  The rigid\sphinxhyphen{}body algorithm used by \sphinxstyleemphasis{rigid/nvt}
is described in the paper by {\hyperref[\detokenize{fix_rigid:kamberaj}]{\sphinxcrossref{\DUrole{std,std-ref}{Kamberaj}}}}.

\sphinxAtStartPar
The \sphinxstyleemphasis{rigid/npt}, \sphinxstyleemphasis{rigid/nph}, \sphinxstyleemphasis{rigid/npt/small}, and \sphinxstyleemphasis{rigid/nph/small}
styles perform constant NPT or NPH integration using a Nose/Hoover
barostat with chains.  They are referred to below as the 4 NPT and NPH
rigid styles.  For the NPT case, the same Nose/Hoover thermostat is
also used as with \sphinxstyleemphasis{rigid/nvt} and \sphinxstyleemphasis{rigid/nvt/small}.

\sphinxAtStartPar
The barostat parameters are specified using one or more of the \sphinxstyleemphasis{iso},
\sphinxstyleemphasis{aniso}, \sphinxstyleemphasis{x}, \sphinxstyleemphasis{y}, \sphinxstyleemphasis{z} and \sphinxstyleemphasis{couple} keywords.  These keywords give you
the ability to specify 3 diagonal components of the external stress
tensor, and to couple these components together so that the dimensions
they represent are varied together during a constant\sphinxhyphen{}pressure
simulation.  The effects of these keywords are similar to those
defined in {\hyperref[\detokenize{fix_nh::doc}]{\sphinxcrossref{\DUrole{doc}{fix npt/nph}}}}

\begin{sphinxadmonition}{note}{Note:}
\sphinxAtStartPar
Currently the \sphinxstyleemphasis{rigid/npt}, \sphinxstyleemphasis{rigid/nph}, \sphinxstyleemphasis{rigid/npt/small}, and
\sphinxstyleemphasis{rigid/nph/small} styles do not support triclinic (non\sphinxhyphen{}orthogonal)
boxes.
\end{sphinxadmonition}

\sphinxAtStartPar
The target pressures for each of the 6 components of the stress tensor
can be specified independently via the \sphinxstyleemphasis{x}, \sphinxstyleemphasis{y}, \sphinxstyleemphasis{z} keywords, which
correspond to the 3 simulation box dimensions.  For each component,
the external pressure or tensor component at each timestep is a ramped
value during the run from \sphinxstyleemphasis{Pstart} to \sphinxstyleemphasis{Pstop}. If a target pressure
is specified for a component, then the corresponding box dimension
will change during a simulation.  For example, if the \sphinxstyleemphasis{y} keyword is
used, the y\sphinxhyphen{}box length will change.  A box dimension will not change
if that component is not specified, although you have the option to
change that dimension via the {\hyperref[\detokenize{fix_deform::doc}]{\sphinxcrossref{\DUrole{doc}{fix deform}}}} command.

\sphinxAtStartPar
For all barostat keywords, the \sphinxstyleemphasis{Pdamp} parameter operates like the
\sphinxstyleemphasis{Tdamp} parameter, determining the time scale on which pressure is
relaxed.  For example, a value of 10.0 means to relax the pressure in
a timespan of (roughly) 10 time units (e.g. \(\tau\) or fs or ps
\sphinxhyphen{} see the {\hyperref[\detokenize{units::doc}]{\sphinxcrossref{\DUrole{doc}{units}}}} command).

\sphinxAtStartPar
Regardless of what atoms are in the fix group (the only atoms which
are time integrated), a global pressure or stress tensor is computed
for all atoms.  Similarly, when the size of the simulation box is
changed, all atoms are re\sphinxhyphen{}scaled to new positions, unless the keyword
\sphinxstyleemphasis{dilate} is specified with a \sphinxstyleemphasis{dilate\sphinxhyphen{}group\sphinxhyphen{}ID} for a group that
represents a subset of the atoms.  This can be useful, for example, to
leave the coordinates of atoms in a solid substrate unchanged and
controlling the pressure of a surrounding fluid.  Another example is a
system consisting of rigid bodies and point particles where the
barostat is only coupled with the rigid bodies.  This option should be
used with care, since it can be unphysical to dilate some atoms and
not others, because it can introduce large, instantaneous
displacements between a pair of atoms (one dilated, one not) that are
far from the dilation origin.

\sphinxAtStartPar
The \sphinxstyleemphasis{couple} keyword allows two or three of the diagonal components of
the pressure tensor to be “coupled” together.  The value specified
with the keyword determines which are coupled.  For example, \sphinxstyleemphasis{xz}
means the \sphinxstyleemphasis{Pxx} and \sphinxstyleemphasis{Pzz} components of the stress tensor are coupled.
\sphinxstyleemphasis{Xyz} means all 3 diagonal components are coupled.  Coupling means two
things: the instantaneous stress will be computed as an average of the
corresponding diagonal components, and the coupled box dimensions will
be changed together in lockstep, meaning coupled dimensions will be
dilated or contracted by the same percentage every timestep.  The
\sphinxstyleemphasis{Pstart}, \sphinxstyleemphasis{Pstop}, \sphinxstyleemphasis{Pdamp} parameters for any coupled dimensions must
be identical.  \sphinxstyleemphasis{Couple xyz} can be used for a 2d simulation; the \sphinxstyleemphasis{z}
dimension is simply ignored.

\sphinxAtStartPar
The \sphinxstyleemphasis{iso} and \sphinxstyleemphasis{aniso} keywords are simply shortcuts that are
equivalent to specifying several other keywords together.

\sphinxAtStartPar
The keyword \sphinxstyleemphasis{iso} means couple all 3 diagonal components together when
pressure is computed (hydrostatic pressure), and dilate/contract the
dimensions together.  Using “iso Pstart Pstop Pdamp” is the same as
specifying these 4 keywords:

\begin{sphinxVerbatim}[commandchars=\\\{\}]
x Pstart Pstop Pdamp
y Pstart Pstop Pdamp
z Pstart Pstop Pdamp
couple xyz
\end{sphinxVerbatim}

\sphinxAtStartPar
The keyword \sphinxstyleemphasis{aniso} means \sphinxstyleemphasis{x}, \sphinxstyleemphasis{y}, and \sphinxstyleemphasis{z} dimensions are controlled
independently using the \sphinxstyleemphasis{Pxx}, \sphinxstyleemphasis{Pyy}, and \sphinxstyleemphasis{Pzz} components of the
stress tensor as the driving forces, and the specified scalar external
pressure.  Using “aniso Pstart Pstop Pdamp” is the same as specifying
these 4 keywords:

\begin{sphinxVerbatim}[commandchars=\\\{\}]
x Pstart Pstop Pdamp
y Pstart Pstop Pdamp
z Pstart Pstop Pdamp
couple none
\end{sphinxVerbatim}


\bigskip\hrule\bigskip


\sphinxAtStartPar
The keyword/value option pairs are used in the following ways.

\sphinxAtStartPar
The \sphinxstyleemphasis{reinit} keyword determines, whether the rigid body properties
are re\sphinxhyphen{}initialized between run commands. With the option \sphinxstyleemphasis{yes} (the
default) this is done, with the option \sphinxstyleemphasis{no} this is not done. Turning
off the re\sphinxhyphen{}initialization can be helpful to protect rigid bodies against
unphysical manipulations between runs or when properties cannot be
easily re\sphinxhyphen{}computed (e.g. when read from a file). When using the \sphinxstyleemphasis{infile}
keyword, the \sphinxstyleemphasis{reinit} option is automatically set to \sphinxstyleemphasis{no}.

\sphinxAtStartPar
The \sphinxstyleemphasis{langevin} and \sphinxstyleemphasis{temp} and \sphinxstyleemphasis{tparam} keywords perform thermostatting
of the rigid bodies, altering both their translational and rotational
degrees of freedom.  What is meant by “temperature” of a collection of
rigid bodies and how it can be monitored via the fix output is
discussed below.

\sphinxAtStartPar
The \sphinxstyleemphasis{langevin} keyword applies a Langevin thermostat to the constant
NVE time integration performed by any of the 4 NVE rigid styles:
\sphinxstyleemphasis{rigid}, \sphinxstyleemphasis{rigid/nve}, \sphinxstyleemphasis{rigid/small}, \sphinxstyleemphasis{rigid/small/nve}.  It cannot
be used with the 2 NVT rigid styles: \sphinxstyleemphasis{rigid/nvt}, \sphinxstyleemphasis{rigid/small/nvt}.
The desired temperature at each timestep is a ramped value during the
run from \sphinxstyleemphasis{Tstart} to \sphinxstyleemphasis{Tstop}.  The \sphinxstyleemphasis{Tdamp} parameter is specified in
time units and determines how rapidly the temperature is relaxed.  For
example, a value of 100.0 means to relax the temperature in a timespan
of (roughly) 100 time units (\(\tau\) or fs or ps \sphinxhyphen{} see the
{\hyperref[\detokenize{units::doc}]{\sphinxcrossref{\DUrole{doc}{units}}}} command).  The random \# \sphinxstyleemphasis{seed} must be a positive
integer.

\sphinxAtStartPar
The way that Langevin thermostatting operates is explained on the {\hyperref[\detokenize{fix_langevin::doc}]{\sphinxcrossref{\DUrole{doc}{fix langevin}}}} doc page.  If you wish to simply viscously
damp the rotational motion without thermostatting, you can set
\sphinxstyleemphasis{Tstart} and \sphinxstyleemphasis{Tstop} to 0.0, which means only the viscous drag term in
the Langevin thermostat will be applied.  See the discussion on the
{\hyperref[\detokenize{fix_viscous::doc}]{\sphinxcrossref{\DUrole{doc}{fix viscous}}}} page for details.

\begin{sphinxadmonition}{note}{Note:}
\sphinxAtStartPar
When the \sphinxstyleemphasis{langevin} keyword is used with fix rigid versus fix
rigid/small, different dynamics will result for parallel runs.  This
is because of the way random numbers are used in the two cases.  The
dynamics for the two cases should be statistically similar, but will
not be identical, even for a single timestep.
\end{sphinxadmonition}

\sphinxAtStartPar
The \sphinxstyleemphasis{temp} and \sphinxstyleemphasis{tparam} keywords apply a Nose/Hoover thermostat to the
NVT time integration performed by the 2 NVT rigid styles.  They cannot
be used with the 4 NVE rigid styles.  The desired temperature at each
timestep is a ramped value during the run from \sphinxstyleemphasis{Tstart} to \sphinxstyleemphasis{Tstop}.
The \sphinxstyleemphasis{Tdamp} parameter is specified in time units and determines how
rapidly the temperature is relaxed.  For example, a value of 100.0
means to relax the temperature in a timespan of (roughly) 100 time
units (tau or fs or ps \sphinxhyphen{} see the {\hyperref[\detokenize{units::doc}]{\sphinxcrossref{\DUrole{doc}{units}}}} command).

\sphinxAtStartPar
Nose/Hoover chains are used in conjunction with this thermostat.  The
\sphinxstyleemphasis{tparam} keyword can optionally be used to change the chain settings
used.  \sphinxstyleemphasis{Tchain} is the number of thermostats in the Nose Hoover chain.
This value, along with \sphinxstyleemphasis{Tdamp} can be varied to dampen undesirable
oscillations in temperature that can occur in a simulation.  As a rule
of thumb, increasing the chain length should lead to smaller
oscillations. The keyword \sphinxstyleemphasis{pchain} specifies the number of thermostats
in the chain thermostatting the barostat degrees of freedom.

\begin{sphinxadmonition}{note}{Note:}
\sphinxAtStartPar
There are alternate ways to thermostat a system of rigid bodies.
You can use {\hyperref[\detokenize{fix_langevin::doc}]{\sphinxcrossref{\DUrole{doc}{fix langevin}}}} to treat the
individual particles in the rigid bodies as effectively immersed in
an implicit solvent, e.g. a Brownian dynamics model.  For hybrid
systems with both rigid bodies and solvent particles, you can
thermostat only the solvent particles that surround one or more
rigid bodies by appropriate choice of groups in the compute and fix
commands for temperature and thermostatting.  The solvent
interactions with the rigid bodies should then effectively
thermostat the rigid body temperature as well without use of the
Langevin or Nose/Hoover options associated with the fix rigid
commands.
\end{sphinxadmonition}


\bigskip\hrule\bigskip


\sphinxAtStartPar
The \sphinxstyleemphasis{mol} keyword can only be used with the \sphinxstyleemphasis{rigid/small} styles.  It
must be used when other commands, such as {\hyperref[\detokenize{fix_deposit::doc}]{\sphinxcrossref{\DUrole{doc}{fix deposit}}}} or {\hyperref[\detokenize{fix_pour::doc}]{\sphinxcrossref{\DUrole{doc}{fix pour}}}}, add rigid bodies
on\sphinxhyphen{}the\sphinxhyphen{}fly during a simulation.  You specify a \sphinxstyleemphasis{template\sphinxhyphen{}ID}
previously defined using the {\hyperref[\detokenize{molecule::doc}]{\sphinxcrossref{\DUrole{doc}{molecule}}}} command, which
reads a file that defines the molecule.  You must use the same
\sphinxstyleemphasis{template\sphinxhyphen{}ID} that the other fix which is adding rigid bodies uses.
The coordinates, atom types, atom diameters, center\sphinxhyphen{}of\sphinxhyphen{}mass, and
moments of inertia can be specified in the molecule file.  See the
{\hyperref[\detokenize{molecule::doc}]{\sphinxcrossref{\DUrole{doc}{molecule}}}} command for details.  The only settings
required to be in this file are the coordinates and types of atoms in
the molecule, in which case the molecule command calculates the other
quantities itself.

\sphinxAtStartPar
Note that these other fixes create new rigid bodies, in addition to
those defined initially by this fix via the \sphinxstyleemphasis{bodystyle} setting.

\sphinxAtStartPar
Also note that when using the \sphinxstyleemphasis{mol} keyword, extra restart information
about all rigid bodies is written out whenever a restart file is
written out.  See the NOTE in the next section for details.


\bigskip\hrule\bigskip


\sphinxAtStartPar
The \sphinxstyleemphasis{infile} keyword allows a file of rigid body attributes to be read
in from a file, rather then having LAMMPS compute them.  There are 5
such attributes: the total mass of the rigid body, its center\sphinxhyphen{}of\sphinxhyphen{}mass
position, its 6 moments of inertia, its center\sphinxhyphen{}of\sphinxhyphen{}mass velocity, and
the 3 image flags of the center\sphinxhyphen{}of\sphinxhyphen{}mass position.  For rigid bodies
consisting of point particles or non\sphinxhyphen{}overlapping finite\sphinxhyphen{}size
particles, LAMMPS can compute these values accurately.

\sphinxAtStartPar
However, for rigid bodies consisting of finite\sphinxhyphen{}size particles which
overlap each other, LAMMPS will ignore the overlaps when computing
these 4 attributes, which means the dynamics of the bodies will be
incorrect.  The amount of error this induces depends on the amount of
overlap.  To avoid this issue, the values can be pre\sphinxhyphen{}computed
(e.g. using Monte Carlo integration).

\sphinxAtStartPar
The format of the file is as follows.  Note that the file does not
have to list attributes for every rigid body integrated by fix rigid.
Only bodies which the file specifies will have their computed
attributes overridden.  The file can contain initial blank lines or
comment lines starting with “\#” which are ignored.  The first
non\sphinxhyphen{}blank, non\sphinxhyphen{}comment line should list N = the number of lines to
follow.  The N successive lines contain the following information:

\begin{sphinxVerbatim}[commandchars=\\\{\}]
ID1 masstotal xcm ycm zcm ixx iyy izz ixy ixz iyz vxcm vycm vzcm lx ly lz ixcm iycm izcm
ID2 masstotal xcm ycm zcm ixx iyy izz ixy ixz iyz vxcm vycm vzcm lx ly lz ixcm iycm izcm
...
IDN masstotal xcm ycm zcm ixx iyy izz ixy ixz iyz vxcm vycm vzcm lx ly lz ixcm iycm izcm
\end{sphinxVerbatim}

\sphinxAtStartPar
The rigid body IDs are all positive integers.  For the \sphinxstyleemphasis{single}
bodystyle, only an ID of 1 can be used.  For the \sphinxstyleemphasis{group} bodystyle,
IDs from 1 to Ng can be used where Ng is the number of specified
groups.  For the \sphinxstyleemphasis{molecule} bodystyle, use the molecule ID for the
atoms in a specific rigid body as the rigid body ID.

\sphinxAtStartPar
The masstotal and center\sphinxhyphen{}of\sphinxhyphen{}mass coordinates (xcm,ycm,zcm) are
self\sphinxhyphen{}explanatory.  The center\sphinxhyphen{}of\sphinxhyphen{}mass should be consistent with what
is calculated for the position of the rigid body with all its atoms
unwrapped by their respective image flags.  If this produces a
center\sphinxhyphen{}of\sphinxhyphen{}mass that is outside the simulation box, LAMMPS wraps it
back into the box.

\sphinxAtStartPar
The 6 moments of inertia (ixx,iyy,izz,ixy,ixz,iyz) should be the
values consistent with the current orientation of the rigid body
around its center of mass.  The values are with respect to the
simulation box XYZ axes, not with respect to the principal axes of the
rigid body itself.  LAMMPS performs the latter calculation internally.

\sphinxAtStartPar
The (vxcm,vycm,vzcm) values are the velocity of the center of mass.
The (lx,ly,lz) values are the angular momentum of the body.  The
(vxcm,vycm,vzcm) and (lx,ly,lz) values can simply be set to 0 if you
wish the body to have no initial motion.

\sphinxAtStartPar
The (ixcm,iycm,izcm) values are the image flags of the center of mass
of the body.  For periodic dimensions, they specify which image of the
simulation box the body is considered to be in.  An image of 0 means
it is inside the box as defined.  A value of 2 means add 2 box lengths
to get the true value.  A value of \sphinxhyphen{}1 means subtract 1 box length to
get the true value.  LAMMPS updates these flags as the rigid bodies
cross periodic boundaries during the simulation.

\begin{sphinxadmonition}{note}{Note:}
\sphinxAtStartPar
If you use the \sphinxstyleemphasis{infile} or \sphinxstyleemphasis{mol} keywords and write restart files
during a simulation, then each time a restart file is written, the
fix also write an auxiliary restart file with the name rfile.rigid,
where “rfile” is the name of the restart file,
e.g. tmp.restart.10000 and tmp.restart.10000.rigid.  This auxiliary
file is in the same format described above.  Thus it can be used in
a new input script that restarts the run and re\sphinxhyphen{}specifies a rigid
fix using an \sphinxstyleemphasis{infile} keyword and the appropriate filename.  Note
that the auxiliary file will contain one line for every rigid body,
even if the original file only listed a subset of the rigid bodies.
\end{sphinxadmonition}

\sphinxAtStartPar
If the system has rigid bodies with finite\sphinxhyphen{}size overlapping particles
and the model uses the {\hyperref[\detokenize{fix_gravity::doc}]{\sphinxcrossref{\DUrole{doc}{fix gravity}}}} command to
apply a gravitational force to the rigid bodies, then the \sphinxstyleemphasis{gravity}
keyword should be used in the following manner.

\sphinxAtStartPar
First, the group specified for the {\hyperref[\detokenize{fix_gravity::doc}]{\sphinxcrossref{\DUrole{doc}{fix gravity}}}}
command should not include any atoms in rigid bodies which have
overlapping particles.  It can be empty (see the {\hyperref[\detokenize{group::doc}]{\sphinxcrossref{\DUrole{doc}{group empty}}}} command) or only contain single particles not in rigid
bodies, e.g. background particles.

\sphinxAtStartPar
Second, the \sphinxstyleemphasis{infile} keyword should be used to specify the total mass
and other properties of the rigid bodies with overlaps, so that their
dynamics will be modeled correctly, as explained above.

\sphinxAtStartPar
Third, the \sphinxstyleemphasis{gravity} keyword should be used the with the ID of the
{\hyperref[\detokenize{fix_gravity::doc}]{\sphinxcrossref{\DUrole{doc}{fix gravity}}}} command as its argument.  The rigid
fixes will access the gravity fix to extract the current direction of
the gravity vector at each timestep (which can be static or dynamic).
A gravity force will then be applied to each rigid body at its
center\sphinxhyphen{}of\sphinxhyphen{}mass position using its total mass.


\bigskip\hrule\bigskip


\sphinxAtStartPar
If you use a {\hyperref[\detokenize{compute::doc}]{\sphinxcrossref{\DUrole{doc}{temperature compute}}}} with a group that
includes particles in rigid bodies, the degrees\sphinxhyphen{}of\sphinxhyphen{}freedom removed by
each rigid body are accounted for in the temperature (and pressure)
computation, but only if the temperature group includes all the
particles in a particular rigid body.

\sphinxAtStartPar
A 3d rigid body has 6 degrees of freedom (3 translational, 3
rotational), except for a collection of point particles lying on a
straight line, which has only 5, e.g a dimer.  A 2d rigid body has 3
degrees of freedom (2 translational, 1 rotational).

\begin{sphinxadmonition}{note}{Note:}
\sphinxAtStartPar
You may wish to explicitly subtract additional
degrees\sphinxhyphen{}of\sphinxhyphen{}freedom if you use the \sphinxstyleemphasis{force} and \sphinxstyleemphasis{torque} keywords to
eliminate certain motions of one or more rigid bodies.  LAMMPS does
not do this automatically.
\end{sphinxadmonition}

\sphinxAtStartPar
The rigid body contribution to the pressure of the system (virial) is
also accounted for by this fix.


\bigskip\hrule\bigskip


\sphinxAtStartPar
If your simulation is a hybrid model with a mixture of rigid bodies
and non\sphinxhyphen{}rigid particles (e.g. solvent) there are several ways these
rigid fixes can be used in tandem with {\hyperref[\detokenize{fix_nve::doc}]{\sphinxcrossref{\DUrole{doc}{fix nve}}}},
{\hyperref[\detokenize{fix_nh::doc}]{\sphinxcrossref{\DUrole{doc}{fix nvt}}}}, {\hyperref[\detokenize{fix_nh::doc}]{\sphinxcrossref{\DUrole{doc}{fix npt}}}}, and {\hyperref[\detokenize{fix_nh::doc}]{\sphinxcrossref{\DUrole{doc}{fix nph}}}}.

\sphinxAtStartPar
If you wish to perform NVE dynamics (no thermostatting or
barostatting), use one of 4 NVE rigid styles to integrate the rigid
bodies, and {\hyperref[\detokenize{fix_nve::doc}]{\sphinxcrossref{\DUrole{doc}{fix nve}}}} to integrate the non\sphinxhyphen{}rigid
particles.

\sphinxAtStartPar
If you wish to perform NVT dynamics (thermostatting, but no
barostatting), you can use one of the 2 NVT rigid styles for the rigid
bodies, and any thermostatting fix for the non\sphinxhyphen{}rigid particles
({\hyperref[\detokenize{fix_nh::doc}]{\sphinxcrossref{\DUrole{doc}{fix nvt}}}}, {\hyperref[\detokenize{fix_langevin::doc}]{\sphinxcrossref{\DUrole{doc}{fix langevin}}}},
{\hyperref[\detokenize{fix_temp_berendsen::doc}]{\sphinxcrossref{\DUrole{doc}{fix temp/berendsen}}}}).  You can also use one
of the 4 NVE rigid styles for the rigid bodies and thermostat them
using {\hyperref[\detokenize{fix_langevin::doc}]{\sphinxcrossref{\DUrole{doc}{fix langevin}}}} on the group that contains
all the particles in the rigid bodies.  The net force added by
{\hyperref[\detokenize{fix_langevin::doc}]{\sphinxcrossref{\DUrole{doc}{fix langevin}}}} to each rigid body effectively
thermostats its translational center\sphinxhyphen{}of\sphinxhyphen{}mass motion.  Not sure how
well it does at thermostatting its rotational motion.

\sphinxAtStartPar
If you wish to perform NPT or NPH dynamics (barostatting), you cannot
use both {\hyperref[\detokenize{fix_nh::doc}]{\sphinxcrossref{\DUrole{doc}{fix npt}}}} and the NPT or NPH rigid styles.  This
is because there can only be one fix which monitors the global
pressure and changes the simulation box dimensions.  So you have 3
choices:
\begin{enumerate}
\sphinxsetlistlabels{\arabic}{enumi}{enumii}{}{.}%
\item {} 
\sphinxAtStartPar
Use one of the 4 NPT or NPH styles for the rigid bodies.  Use the
\sphinxstyleemphasis{dilate} all option so that it will dilate the positions of the
non\sphinxhyphen{}rigid particles as well.  Use {\hyperref[\detokenize{fix_nh::doc}]{\sphinxcrossref{\DUrole{doc}{fix nvt}}}} (or any
other thermostat) for the non\sphinxhyphen{}rigid particles.

\item {} 
\sphinxAtStartPar
Use {\hyperref[\detokenize{fix_nh::doc}]{\sphinxcrossref{\DUrole{doc}{fix npt}}}} for the group of non\sphinxhyphen{}rigid particles.  Use
the \sphinxstyleemphasis{dilate} all option so that it will dilate the center\sphinxhyphen{}of\sphinxhyphen{}mass
positions of the rigid bodies as well.  Use one of the 4 NVE or 2 NVT
rigid styles for the rigid bodies.

\item {} 
\sphinxAtStartPar
Use {\hyperref[\detokenize{fix_press_berendsen::doc}]{\sphinxcrossref{\DUrole{doc}{fix press/berendsen}}}} to compute the
pressure and change the box dimensions.  Use one of the 4 NVE or 2 NVT
rigid styles for the rigid bodies.  Use {\hyperref[\detokenize{fix_nh::doc}]{\sphinxcrossref{\DUrole{doc}{fix nvt}}}} (or
any other thermostat) for the non\sphinxhyphen{}rigid particles.

\end{enumerate}

\sphinxAtStartPar
In all case, the rigid bodies and non\sphinxhyphen{}rigid particles both contribute
to the global pressure and the box is scaled the same by any of the
barostatting fixes.

\sphinxAtStartPar
You could even use the second and third options for a non\sphinxhyphen{}hybrid
simulation consisting of only rigid bodies, assuming you give
{\hyperref[\detokenize{fix_nh::doc}]{\sphinxcrossref{\DUrole{doc}{fix npt}}}} an empty group, though it’s an odd thing to
do.  The barostatting fixes ({\hyperref[\detokenize{fix_nh::doc}]{\sphinxcrossref{\DUrole{doc}{fix npt}}}} and {\hyperref[\detokenize{fix_press_berendsen::doc}]{\sphinxcrossref{\DUrole{doc}{fix
press/berensen}}}}) will monitor the pressure and
change the box dimensions, but not time integrate any particles.  The
integration of the rigid bodies will be performed by fix rigid/nvt.


\bigskip\hrule\bigskip


\sphinxAtStartPar
Styles with a \sphinxstyleemphasis{gpu}, \sphinxstyleemphasis{intel}, \sphinxstyleemphasis{kk}, \sphinxstyleemphasis{omp}, or \sphinxstyleemphasis{opt} suffix are
functionally the same as the corresponding style without the suffix.
They have been optimized to run faster, depending on your available
hardware, as discussed on the {\hyperref[\detokenize{Speed_packages::doc}]{\sphinxcrossref{\DUrole{doc}{Accelerator packages}}}}
page.  The accelerated styles take the same arguments and should
produce the same results, except for round\sphinxhyphen{}off and precision issues.

\sphinxAtStartPar
These accelerated styles are part of the GPU, INTEL, KOKKOS,
OPENMP, and OPT packages, respectively.  They are only enabled if
LAMMPS was built with those packages.  See the {\hyperref[\detokenize{Build_package::doc}]{\sphinxcrossref{\DUrole{doc}{Build package}}}} page for more info.

\sphinxAtStartPar
You can specify the accelerated styles explicitly in your input script
by including their suffix, or you can use the {\hyperref[\detokenize{Run_options::doc}]{\sphinxcrossref{\DUrole{doc}{\sphinxhyphen{}suffix command\sphinxhyphen{}line switch}}}} when you invoke LAMMPS, or you can use the
{\hyperref[\detokenize{suffix::doc}]{\sphinxcrossref{\DUrole{doc}{suffix}}}} command in your input script.

\sphinxAtStartPar
See the {\hyperref[\detokenize{Speed_packages::doc}]{\sphinxcrossref{\DUrole{doc}{Accelerator packages}}}} page for more
instructions on how to use the accelerated styles effectively.


\bigskip\hrule\bigskip



\subsection{Restart, fix\_modify, output, run start/stop, minimize info}
\label{\detokenize{fix_rigid:restart-fix-modify-output-run-start-stop-minimize-info}}
\sphinxAtStartPar
No information about the 4 NVE rigid styles is written to {\hyperref[\detokenize{restart::doc}]{\sphinxcrossref{\DUrole{doc}{binary
restart files}}}}.  The exception is if the \sphinxstyleemphasis{infile} or \sphinxstyleemphasis{mol}
keyword is used, in which case an auxiliary file is written out with
rigid body information each time a restart file is written, as
explained above for the \sphinxstyleemphasis{infile} keyword.  For the 2 NVT rigid styles,
the state of the Nose/Hoover thermostat is written to {\hyperref[\detokenize{restart::doc}]{\sphinxcrossref{\DUrole{doc}{binary
restart files}}}}.  Ditto for the 4 NPT and NPH rigid styles,
and the state of the Nose/Hoover barostat.  See the {\hyperref[\detokenize{read_restart::doc}]{\sphinxcrossref{\DUrole{doc}{read\_restart}}}} command for info on how to re\sphinxhyphen{}specify a fix in an
input script that reads a restart file, so that the operation of the
fix continues in an uninterrupted fashion.

\sphinxAtStartPar
The {\hyperref[\detokenize{fix_modify::doc}]{\sphinxcrossref{\DUrole{doc}{fix\_modify}}}} \sphinxstyleemphasis{temp} and \sphinxstyleemphasis{press} options are
supported by the 4 NPT and NPH rigid styles to change the computes
used to calculate the instantaneous pressure tensor.  Note that the 2
NVT rigid fixes do not use any external compute to compute
instantaneous temperature.

\sphinxAtStartPar
The {\hyperref[\detokenize{fix_modify::doc}]{\sphinxcrossref{\DUrole{doc}{fix\_modify}}}} \sphinxstyleemphasis{bodyforces} option is supported by
all rigid styles to set whether per\sphinxhyphen{}body forces and torques are
computed early or late in a timestep, i.e. at the post\sphinxhyphen{}force stage or
at the final\sphinxhyphen{}integrate stage or the timestep, respectively.

\sphinxAtStartPar
The cumulative energy change in the system imposed by the 6 NVT, NPT,
NPH rigid fixes, via either thermostatting and/or barostatting, is
included in the {\hyperref[\detokenize{thermo_style::doc}]{\sphinxcrossref{\DUrole{doc}{thermodynamic output}}}} keywords
\sphinxstyleemphasis{ecouple} and \sphinxstyleemphasis{econserve}.  See the {\hyperref[\detokenize{thermo_style::doc}]{\sphinxcrossref{\DUrole{doc}{thermo\_style}}}}
doc page for details.

\sphinxAtStartPar
The 2 NVE rigid fixes compute a global scalar which can be accessed by
various {\hyperref[\detokenize{Howto_output::doc}]{\sphinxcrossref{\DUrole{doc}{output commands}}}}.  The scalar value
calculated by these fixes is “intensive”.  The scalar is the current
temperature of the collection of rigid bodies.  This is averaged over
all rigid bodies and their translational and rotational degrees of
freedom.  The translational energy of a rigid body is 1/2 m v\textasciicircum{}2,
where m = total mass of the body and v = the velocity of its center of
mass.  The rotational energy of a rigid body is 1/2 I w\textasciicircum{}2, where I =
the moment of inertia tensor of the body and w = its angular velocity.
Degrees of freedom constrained by the \sphinxstyleemphasis{force} and \sphinxstyleemphasis{torque} keywords
are removed from this calculation, but only for the \sphinxstyleemphasis{rigid} and
\sphinxstyleemphasis{rigid/nve} fixes.

\sphinxAtStartPar
The 6 NVT, NPT, NPH rigid fixes compute a global scalar which can be
accessed by various {\hyperref[\detokenize{Howto_output::doc}]{\sphinxcrossref{\DUrole{doc}{output commands}}}}.  The scalar
is the same cumulative energy change due to these fixes described
above.  The scalar value calculated by this fix is “extensive”.

\sphinxAtStartPar
The {\hyperref[\detokenize{fix_modify::doc}]{\sphinxcrossref{\DUrole{doc}{fix\_modify}}}} \sphinxstyleemphasis{virial} option is supported by
these fixes to add the contribution due to the added forces on atoms
to both the global pressure and per\sphinxhyphen{}atom stress of the system via the
{\hyperref[\detokenize{compute_pressure::doc}]{\sphinxcrossref{\DUrole{doc}{compute pressure}}}} and {\hyperref[\detokenize{compute_stress_atom::doc}]{\sphinxcrossref{\DUrole{doc}{compute
stress/atom}}}} commands.  The former can be
accessed by {\hyperref[\detokenize{thermo_style::doc}]{\sphinxcrossref{\DUrole{doc}{thermodynamic output}}}}.  The default
setting for this fix is {\hyperref[\detokenize{fix_modify::doc}]{\sphinxcrossref{\DUrole{doc}{fix\_modify virial yes}}}}.

\sphinxAtStartPar
All of the \sphinxstyleemphasis{rigid} styles (but not the \sphinxstyleemphasis{rigid/small} styles) compute a
global array of values which can be accessed by various {\hyperref[\detokenize{Howto_output::doc}]{\sphinxcrossref{\DUrole{doc}{output
commands}}}}.  Similar information about the bodies
defined by the \sphinxstyleemphasis{rigid/small} styles can be accessed via the
{\hyperref[\detokenize{compute_rigid_local::doc}]{\sphinxcrossref{\DUrole{doc}{compute rigid/local}}}} command.

\sphinxAtStartPar
The number of rows in the array is equal to the number of rigid
bodies.  The number of columns is 15.  Thus for each rigid body, 15
values are stored: the xyz coords of the center of mass (COM), the xyz
components of the COM velocity, the xyz components of the force acting
on the COM, the xyz components of the torque acting on the COM, and
the xyz image flags of the COM.

\sphinxAtStartPar
The center of mass (COM) for each body is similar to unwrapped
coordinates written to a dump file.  It will always be inside (or
slightly outside) the simulation box.  The image flags have the same
meaning as image flags for atom positions (see the “dump” command).
This means you can calculate the unwrapped COM by applying the image
flags to the COM, the same as when unwrapped coordinates are written
to a dump file.

\sphinxAtStartPar
The force and torque values in the array are not affected by the
\sphinxstyleemphasis{force} and \sphinxstyleemphasis{torque} keywords in the fix rigid command; they reflect
values before any changes are made by those keywords.

\sphinxAtStartPar
The ordering of the rigid bodies (by row in the array) is as follows.
For the \sphinxstyleemphasis{single} keyword there is just one rigid body.  For the
\sphinxstyleemphasis{molecule} keyword, the bodies are ordered by ascending molecule ID.
For the \sphinxstyleemphasis{group} keyword, the list of group IDs determines the ordering
of bodies.

\sphinxAtStartPar
The array values calculated by these fixes are “intensive”, meaning
they are independent of the number of atoms in the simulation.

\sphinxAtStartPar
No parameter of these fixes can be used with the \sphinxstyleemphasis{start/stop} keywords
of the {\hyperref[\detokenize{run::doc}]{\sphinxcrossref{\DUrole{doc}{run}}}} command.  These fixes are not invoked during
{\hyperref[\detokenize{minimize::doc}]{\sphinxcrossref{\DUrole{doc}{energy minimization}}}}.


\bigskip\hrule\bigskip



\subsection{Restrictions}
\label{\detokenize{fix_rigid:restrictions}}
\sphinxAtStartPar
These fixes are all part of the RIGID package.  It is only enabled if
LAMMPS was built with that package.  See the {\hyperref[\detokenize{Build_package::doc}]{\sphinxcrossref{\DUrole{doc}{Build package}}}} page for more info.

\sphinxAtStartPar
Assigning a temperature via the {\hyperref[\detokenize{velocity::doc}]{\sphinxcrossref{\DUrole{doc}{velocity create}}}}
command to a system with {\hyperref[\detokenize{fix_rigid::doc}]{\sphinxcrossref{\DUrole{doc}{rigid bodies}}}} may not have
the desired outcome for two reasons.  First, the velocity command can
be invoked before the rigid\sphinxhyphen{}body fix is invoked or initialized and the
number of adjusted degrees of freedom (DOFs) is known.  Thus it is not
possible to compute the target temperature correctly.  Second, the
assigned velocities may be partially canceled when constraints are
first enforced, leading to a different temperature than desired.  A
workaround for this is to perform a {\hyperref[\detokenize{run::doc}]{\sphinxcrossref{\DUrole{doc}{run 0}}}} command, which
ensures all DOFs are accounted for properly, and then rescale the
temperature to the desired value before performing a simulation.  For
example:

\begin{sphinxVerbatim}[commandchars=\\\{\}]
\PYG{k}{velocity }\PYG{n+nv+nvIdentifier}{all}\PYG{+w}{ }\PYG{n}{create}\PYG{+w}{ }\PYG{l+m}{300.0}\PYG{+w}{ }\PYG{l+m}{12345}
\PYG{k}{run}\PYG{+w}{ }\PYG{l+m}{0}\PYG{+w}{                             }\PYG{c}{\PYGZsh{} temperature may not be 300K}
\PYG{k}{velocity }\PYG{n+nv+nvIdentifier}{all}\PYG{+w}{ }\PYG{n}{scale}\PYG{+w}{ }\PYG{l+m}{300.0}\PYG{+w}{          }\PYG{c}{\PYGZsh{} now it should be}
\end{sphinxVerbatim}


\subsection{Related commands}
\label{\detokenize{fix_rigid:related-commands}}
\sphinxAtStartPar
{\hyperref[\detokenize{delete_bonds::doc}]{\sphinxcrossref{\DUrole{doc}{delete\_bonds}}}}, {\hyperref[\detokenize{neigh_modify::doc}]{\sphinxcrossref{\DUrole{doc}{neigh\_modify}}}}
exclude, {\hyperref[\detokenize{fix_shake::doc}]{\sphinxcrossref{\DUrole{doc}{fix shake}}}}


\subsection{Default}
\label{\detokenize{fix_rigid:default}}
\sphinxAtStartPar
The option defaults are force * on on on and torque * on on on,
meaning all rigid bodies are acted on by center\sphinxhyphen{}of\sphinxhyphen{}mass force and
torque.  Also Tchain = Pchain = 10, Titer = 1, Torder = 3, reinit = yes.


\bigskip\hrule\bigskip

\phantomsection\label{\detokenize{fix_rigid:hoover}}
\sphinxAtStartPar
\sphinxstylestrong{(Hoover)} Hoover, Phys Rev A, 31, 1695 (1985).

\phantomsection\label{\detokenize{fix_rigid:kamberaj}}
\sphinxAtStartPar
\sphinxstylestrong{(Kamberaj)} Kamberaj, Low, Neal, J Chem Phys, 122, 224114 (2005).

\phantomsection\label{\detokenize{fix_rigid:martyna2}}
\sphinxAtStartPar
\sphinxstylestrong{(Martyna)} Martyna, Klein, Tuckerman, J Chem Phys, 97, 2635 (1992);
Martyna, Tuckerman, Tobias, Klein, Mol Phys, 87, 1117.

\phantomsection\label{\detokenize{fix_rigid:miller3}}
\sphinxAtStartPar
\sphinxstylestrong{(Miller)} Miller, Eleftheriou, Pattnaik, Ndirango, and Newns,
J Chem Phys, 116, 8649 (2002).

\phantomsection\label{\detokenize{fix_rigid:zhang1}}
\sphinxAtStartPar
\sphinxstylestrong{(Zhang)} Zhang, Glotzer, Nanoletters, 4, 1407\sphinxhyphen{}1413 (2004).

\sphinxstepscope

\index{fix rigid/meso@\spxentry{fix rigid/meso}}\ignorespaces 

\section{fix rigid/meso command}
\label{\detokenize{fix_rigid_meso:fix-rigid-meso-command}}\label{\detokenize{fix_rigid_meso:index-0}}\label{\detokenize{fix_rigid_meso::doc}}

\subsection{Syntax}
\label{\detokenize{fix_rigid_meso:syntax}}
\begin{sphinxVerbatim}[commandchars=\\\{\}]
\PYG{k}{fix }\PYG{n+nv+nvIdentifier}{ID}\PYG{+w}{ }\PYG{n+nv+nvIdentifier}{group\PYGZhy{}ID}\PYG{+w}{ }\PYG{n}{rigid}\PYG{o}{/}\PYG{n}{meso}\PYG{+w}{ }\PYG{n}{bodystyle}\PYG{+w}{ }\PYG{n}{args}\PYG{+w}{ }\PYG{n}{keyword}\PYG{+w}{ }\PYG{n}{values}\PYG{+w}{ }\PYG{n}{...}
\end{sphinxVerbatim}
\begin{itemize}
\item {} 
\sphinxAtStartPar
ID, group\sphinxhyphen{}ID are documented in {\hyperref[\detokenize{fix::doc}]{\sphinxcrossref{\DUrole{doc}{fix}}}} command

\item {} 
\sphinxAtStartPar
rigid/meso = style name of this fix command

\item {} 
\sphinxAtStartPar
bodystyle = \sphinxstyleemphasis{single} or \sphinxstyleemphasis{molecule} or \sphinxstyleemphasis{group}
\begin{sphinxalltt}
\sphinxstyleemphasis{single} args = none
\sphinxstyleemphasis{molecule} args = none
\sphinxstyleemphasis{custom} args = \sphinxstyleemphasis{i\_propname} or \sphinxstyleemphasis{v\_varname}
  i\_propname = an integer property defined via fix property/atom
  v\_varname  = an atom\sphinxhyphen{}style or atomfile\sphinxhyphen{}style variable
\sphinxstyleemphasis{group} args = N groupID1 groupID2 ...
  N = \# of groups
  groupID1, groupID2, ... = list of N group IDs
\end{sphinxalltt}

\item {} 
\sphinxAtStartPar
zero or more keyword/value pairs may be appended

\item {} 
\sphinxAtStartPar
keyword = \sphinxstyleemphasis{reinit} or \sphinxstyleemphasis{force} or \sphinxstyleemphasis{torque} or \sphinxstyleemphasis{infile}
\begin{sphinxalltt}
\sphinxstyleemphasis{reinit} = \sphinxstyleemphasis{yes} or \sphinxstyleemphasis{no}
\sphinxstyleemphasis{force} values = M xflag yflag zflag
  M = which rigid body from 1\sphinxhyphen{}Nbody (see asterisk form below)
  xflag,yflag,zflag = off/on if component of center\sphinxhyphen{}of\sphinxhyphen{}mass force is active
\sphinxstyleemphasis{torque} values = M xflag yflag zflag
  M = which rigid body from 1\sphinxhyphen{}Nbody (see asterisk form below)
  xflag,yflag,zflag = off/on if component of center\sphinxhyphen{}of\sphinxhyphen{}mass torque is active
\sphinxstyleemphasis{infile} filename
  filename = file with per\sphinxhyphen{}body values of mass, center\sphinxhyphen{}of\sphinxhyphen{}mass, moments of inertia
\end{sphinxalltt}

\end{itemize}


\subsection{Examples}
\label{\detokenize{fix_rigid_meso:examples}}
\begin{sphinxVerbatim}[commandchars=\\\{\}]
\PYG{k}{fix }\PYG{n+nv+nvIdentifier}{1}\PYG{+w}{ }\PYG{n+nv+nvIdentifier}{ellipsoid}\PYG{+w}{ }\PYG{n}{rigid}\PYG{o}{/}\PYG{n}{meso}\PYG{+w}{ }\PYG{n}{single}
\PYG{k}{fix }\PYG{n+nv+nvIdentifier}{1}\PYG{+w}{ }\PYG{n+nv+nvIdentifier}{rods}\PYG{+w}{      }\PYG{n}{rigid}\PYG{o}{/}\PYG{n}{meso}\PYG{+w}{ }\PYG{n}{molecule}
\PYG{k}{fix }\PYG{n+nv+nvIdentifier}{1}\PYG{+w}{ }\PYG{n+nv+nvIdentifier}{spheres}\PYG{+w}{   }\PYG{n}{rigid}\PYG{o}{/}\PYG{n}{meso}\PYG{+w}{ }\PYG{n}{single}\PYG{+w}{ }\PYG{n}{force}\PYG{+w}{ }\PYG{l+m}{1}\PYG{+w}{ }\PYG{n}{off}\PYG{+w}{ }\PYG{n}{off}\PYG{+w}{ }\PYG{n}{on}
\PYG{k}{fix }\PYG{n+nv+nvIdentifier}{1}\PYG{+w}{ }\PYG{n+nv+nvIdentifier}{particles}\PYG{+w}{ }\PYG{n}{rigid}\PYG{o}{/}\PYG{n}{meso}\PYG{+w}{ }\PYG{n}{molecule}\PYG{+w}{ }\PYG{n}{force}\PYG{+w}{ }\PYG{l+m}{1}\PYG{o}{*}\PYG{l+m}{5}\PYG{+w}{ }\PYG{n}{off}\PYG{+w}{ }\PYG{n}{off}\PYG{+w}{ }\PYG{n}{off}\PYG{+w}{ }\PYG{n}{force}\PYG{+w}{ }\PYG{l+m}{6}\PYG{o}{*}\PYG{l+m}{10}\PYG{+w}{ }\PYG{n}{off}\PYG{+w}{ }\PYG{n}{off}\PYG{+w}{ }\PYG{n}{on}
\PYG{k}{fix }\PYG{n+nv+nvIdentifier}{2}\PYG{+w}{ }\PYG{n+nv+nvIdentifier}{spheres}\PYG{+w}{   }\PYG{n}{rigid}\PYG{o}{/}\PYG{n}{meso}\PYG{+w}{ }\PYG{k}{group }\PYG{n+nv+nvIdentifier}{3}\PYG{+w}{ }\PYG{n}{sphere1}\PYG{+w}{ }\PYG{n}{sphere2}\PYG{+w}{ }\PYG{n}{sphere3}\PYG{+w}{ }\PYG{n}{torque}\PYG{+w}{ }\PYG{o}{*}\PYG{+w}{ }\PYG{n}{off}\PYG{+w}{ }\PYG{n}{off}\PYG{+w}{ }\PYG{n}{off}
\end{sphinxVerbatim}


\subsection{Description}
\label{\detokenize{fix_rigid_meso:description}}
\sphinxAtStartPar
Treat one or more sets of mesoscopic SPH/SDPD particles as independent
rigid bodies.  This means that each timestep the total force and torque
on each rigid body is computed as the sum of the forces and torques on
its constituent particles.  The coordinates and velocities of the
particles in each body are then updated so that the body moves and
rotates as a single entity using the methods described in the paper by
{\hyperref[\detokenize{fix_rigid_meso:miller}]{\sphinxcrossref{\DUrole{std,std-ref}{(Miller)}}}}. Density and internal energy of the particles will
also be updated. This is implemented by creating internal data structures
for each rigid body and performing time integration on these data
structures.  Positions and velocities of the constituent particles are
regenerated from the rigid body data structures in every time step. This
restricts which operations and fixes can be applied to rigid bodies. See
below for a detailed discussion.

\sphinxAtStartPar
The operation of this fix is exactly like that described by the
{\hyperref[\detokenize{fix_rigid::doc}]{\sphinxcrossref{\DUrole{doc}{fix rigid/nve}}}} command, except that particles’ density,
internal energy and extrapolated velocity are also updated.

\begin{sphinxadmonition}{note}{Note:}
\sphinxAtStartPar
You should not update the particles in rigid bodies via other
time\sphinxhyphen{}integration fixes (e.g. {\hyperref[\detokenize{fix_sph::doc}]{\sphinxcrossref{\DUrole{doc}{fix sph}}}}, {\hyperref[\detokenize{fix_sph_stationary::doc}]{\sphinxcrossref{\DUrole{doc}{fix
sph/stationary}}}}), or you will have conflicting
updates to positions and velocities resulting in unphysical
behavior in most cases. When performing a hybrid simulation with
some atoms in rigid bodies, and some not, a separate time
integration fix like {\hyperref[\detokenize{fix_sph::doc}]{\sphinxcrossref{\DUrole{doc}{fix sph}}}} should be used for
the non\sphinxhyphen{}rigid particles.
\end{sphinxadmonition}

\begin{sphinxadmonition}{note}{Note:}
\sphinxAtStartPar
These fixes are overkill if you simply want to hold a collection of
particles stationary or have them move with a constant velocity. To
hold particles stationary use {\hyperref[\detokenize{fix_sph_stationary::doc}]{\sphinxcrossref{\DUrole{doc}{fix sph/stationary}}}} instead. If you would like to move particles
with a constant velocity use {\hyperref[\detokenize{fix_meso_move::doc}]{\sphinxcrossref{\DUrole{doc}{fix meso/move}}}}.
\end{sphinxadmonition}

\begin{sphinxadmonition}{warning}{Warning:}
\sphinxAtStartPar
The aggregate properties of each rigid body are
calculated at the start of a simulation run and are maintained in
internal data structures. The properties include the position and
velocity of the center\sphinxhyphen{}of\sphinxhyphen{}mass of the body, its moments of inertia, and
its angular momentum.  This is done using the properties of the
constituent particles of the body at that point in time (or see the \sphinxstyleemphasis{infile}
keyword option).  Thereafter, changing these properties of individual
particles in the body will have no effect on a rigid body’s dynamics, unless
they effect any computation of per\sphinxhyphen{}particle forces or torques. If the
keyword \sphinxstyleemphasis{reinit} is set to \sphinxstyleemphasis{yes} (the default), the rigid body data
structures will be recreated at the beginning of each \sphinxstyleemphasis{run} command;
if the keyword \sphinxstyleemphasis{reinit} is set to \sphinxstyleemphasis{no}, the rigid body data structures
will be built only at the very first \sphinxstyleemphasis{run} command and maintained for
as long as the rigid fix is defined. For example, you might think you
could displace the particles in a body or add a large velocity to each particle
in a body to make it move in a desired direction before a second run is
performed, using the {\hyperref[\detokenize{set::doc}]{\sphinxcrossref{\DUrole{doc}{set}}}} or
{\hyperref[\detokenize{displace_atoms::doc}]{\sphinxcrossref{\DUrole{doc}{displace\_atoms}}}} or {\hyperref[\detokenize{velocity::doc}]{\sphinxcrossref{\DUrole{doc}{velocity}}}}
commands.  But these commands will not affect the internal attributes
of the body unless \sphinxstyleemphasis{reinit} is set to \sphinxstyleemphasis{yes}. With \sphinxstyleemphasis{reinit} set to \sphinxstyleemphasis{no}
(or using the \sphinxstyleemphasis{infile} option, which implies \sphinxstyleemphasis{reinit} \sphinxstyleemphasis{no}) the position
and velocity of individual particles in the body will be reset when time
integration starts again.
\end{sphinxadmonition}


\bigskip\hrule\bigskip


\sphinxAtStartPar
Each rigid body must have two or more particles.  A particle can belong
to at most one rigid body.  Which particles are in which bodies can be
defined via several options.

\sphinxAtStartPar
For bodystyle \sphinxstyleemphasis{single} the entire fix group of particles is treated as
one rigid body.

\sphinxAtStartPar
For bodystyle \sphinxstyleemphasis{molecule}, particles are grouped into rigid bodies by their
respective molecule IDs: each set of particles in the fix group with the
same molecule ID is treated as a different rigid body.  Note that particles
with a molecule ID = 0 will be treated as a single rigid body. For a
system with solvent (typically this is particles with molecule ID = 0)
surrounding rigid bodies, this may not be what you want.  Thus you
should be careful to use a fix group that only includes particles you
want to be part of rigid bodies.

\sphinxAtStartPar
Bodystyle \sphinxstyleemphasis{custom} is similar to bodystyle \sphinxstyleemphasis{molecule} except that it
is more flexible in using other per\sphinxhyphen{}atom properties to define the sets
of particles that form rigid bodies.  An integer vector defined by the
{\hyperref[\detokenize{fix_property_atom::doc}]{\sphinxcrossref{\DUrole{doc}{fix property/atom}}}} command can be used.  Or an
{\hyperref[\detokenize{variable::doc}]{\sphinxcrossref{\DUrole{doc}{atom\sphinxhyphen{}style or atomfile\sphinxhyphen{}style variable}}}} can be used; the
floating\sphinxhyphen{}point value produced by the variable is rounded to an
integer.  As with bodystyle \sphinxstyleemphasis{molecule}, each set of particles in the fix
groups with the same integer value is treated as a different rigid
body.  Since fix property/atom vectors and atom\sphinxhyphen{}style variables
produce values for all particles, you should be careful to use a fix group
that only includes particles you want to be part of rigid bodies.

\sphinxAtStartPar
For bodystyle \sphinxstyleemphasis{group}, each of the listed groups is treated as a
separate rigid body.  Only particles that are also in the fix group are
included in each rigid body.

\begin{sphinxadmonition}{note}{Note:}
\sphinxAtStartPar
To compute the initial center\sphinxhyphen{}of\sphinxhyphen{}mass position and other
properties of each rigid body, the image flags for each particle in the
body are used to “unwrap” the particle coordinates.  Thus you must
ensure that these image flags are consistent so that the unwrapping
creates a valid rigid body (one where the particles are close together)
, particularly if the particles in a single rigid body straddle a
periodic boundary.  This means the input data file or restart file must
define the image flags for each particle consistently or that you have
used the {\hyperref[\detokenize{set::doc}]{\sphinxcrossref{\DUrole{doc}{set}}}} command to specify them correctly.  If a
dimension is non\sphinxhyphen{}periodic then the image flag of each particle must be
0 in that dimension, else an error is generated.
\end{sphinxadmonition}

\sphinxAtStartPar
By default, each rigid body is acted on by other particles which induce
an external force and torque on its center of mass, causing it to
translate and rotate.  Components of the external center\sphinxhyphen{}of\sphinxhyphen{}mass force
and torque can be turned off by the \sphinxstyleemphasis{force} and \sphinxstyleemphasis{torque} keywords.
This may be useful if you wish a body to rotate but not translate, or
vice versa, or if you wish it to rotate or translate continuously
unaffected by interactions with other particles.  Note that if you
expect a rigid body not to move or rotate by using these keywords, you
must ensure its initial center\sphinxhyphen{}of\sphinxhyphen{}mass translational or angular
velocity is 0.0. Otherwise the initial translational or angular
momentum, the body has, will persist.

\sphinxAtStartPar
An xflag, yflag, or zflag set to \sphinxstyleemphasis{off} means turn off the component of
force or torque in that dimension.  A setting of \sphinxstyleemphasis{on} means turn on
the component, which is the default.  Which rigid body(s) the settings
apply to is determined by the first argument of the \sphinxstyleemphasis{force} and
\sphinxstyleemphasis{torque} keywords.  It can be an integer M from 1 to Nbody, where
Nbody is the number of rigid bodies defined.  A wild\sphinxhyphen{}card asterisk can
be used in place of, or in conjunction with, the M argument to set the
flags for multiple rigid bodies.  This takes the form “*” or “*n” or
“n*” or “m*n”.  If N = the number of rigid bodies, then an asterisk
with no numeric values means all bodies from 1 to N.  A leading
asterisk means all bodies from 1 to n (inclusive).  A trailing
asterisk means all bodies from n to N (inclusive).  A middle asterisk
means all bodies from m to n (inclusive).  Note that you can use the
\sphinxstyleemphasis{force} or \sphinxstyleemphasis{torque} keywords as many times as you like.  If a
particular rigid body has its component flags set multiple times, the
settings from the final keyword are used.

\sphinxAtStartPar
For computational efficiency, you should typically define one fix
rigid/meso command which includes all the desired rigid bodies. LAMMPS
will allow multiple rigid/meso fixes to be defined, but it is more
expensive.


\bigskip\hrule\bigskip


\sphinxAtStartPar
The keyword/value option pairs are used in the following ways.

\sphinxAtStartPar
The \sphinxstyleemphasis{reinit} keyword determines, whether the rigid body properties
are re\sphinxhyphen{}initialized between run commands. With the option \sphinxstyleemphasis{yes} (the
default) this is done, with the option \sphinxstyleemphasis{no} this is not done. Turning
off the re\sphinxhyphen{}initialization can be helpful to protect rigid bodies against
unphysical manipulations between runs or when properties cannot be
easily re\sphinxhyphen{}computed (e.g. when read from a file). When using the \sphinxstyleemphasis{infile}
keyword, the \sphinxstyleemphasis{reinit} option is automatically set to \sphinxstyleemphasis{no}.


\bigskip\hrule\bigskip


\sphinxAtStartPar
The \sphinxstyleemphasis{infile} keyword allows a file of rigid body attributes to be read
in from a file, rather then having LAMMPS compute them.  There are 5
such attributes: the total mass of the rigid body, its center\sphinxhyphen{}of\sphinxhyphen{}mass
position, its 6 moments of inertia, its center\sphinxhyphen{}of\sphinxhyphen{}mass velocity, and
the 3 image flags of the center\sphinxhyphen{}of\sphinxhyphen{}mass position.  For rigid bodies
consisting of point particles or non\sphinxhyphen{}overlapping finite\sphinxhyphen{}size
particles, LAMMPS can compute these values accurately.  However, for
rigid bodies consisting of finite\sphinxhyphen{}size particles which overlap each
other, LAMMPS will ignore the overlaps when computing these 4
attributes.  The amount of error this induces depends on the amount of
overlap.  To avoid this issue, the values can be pre\sphinxhyphen{}computed
(e.g. using Monte Carlo integration).

\sphinxAtStartPar
The format of the file is as follows.  Note that the file does not
have to list attributes for every rigid body integrated by fix rigid.
Only bodies which the file specifies will have their computed
attributes overridden.  The file can contain initial blank lines or
comment lines starting with “\#” which are ignored.  The first
non\sphinxhyphen{}blank, non\sphinxhyphen{}comment line should list N = the number of lines to
follow.  The N successive lines contain the following information:

\begin{sphinxVerbatim}[commandchars=\\\{\}]
ID1 masstotal xcm ycm zcm ixx iyy izz ixy ixz iyz vxcm vycm vzcm lx ly lz ixcm iycm izcm
ID2 masstotal xcm ycm zcm ixx iyy izz ixy ixz iyz vxcm vycm vzcm lx ly lz ixcm iycm izcm
...
IDN masstotal xcm ycm zcm ixx iyy izz ixy ixz iyz vxcm vycm vzcm lx ly lz ixcm iycm izcm
\end{sphinxVerbatim}

\sphinxAtStartPar
The rigid body IDs are all positive integers.  For the \sphinxstyleemphasis{single}
bodystyle, only an ID of 1 can be used.  For the \sphinxstyleemphasis{group} bodystyle,
IDs from 1 to Ng can be used where Ng is the number of specified
groups.  For the \sphinxstyleemphasis{molecule} bodystyle, use the molecule ID for the
atoms in a specific rigid body as the rigid body ID.

\sphinxAtStartPar
The masstotal and center\sphinxhyphen{}of\sphinxhyphen{}mass coordinates (xcm,ycm,zcm) are
self\sphinxhyphen{}explanatory.  The center\sphinxhyphen{}of\sphinxhyphen{}mass should be consistent with what
is calculated for the position of the rigid body with all its atoms
unwrapped by their respective image flags.  If this produces a
center\sphinxhyphen{}of\sphinxhyphen{}mass that is outside the simulation box, LAMMPS wraps it
back into the box.

\sphinxAtStartPar
The 6 moments of inertia (ixx,iyy,izz,ixy,ixz,iyz) should be the
values consistent with the current orientation of the rigid body
around its center of mass.  The values are with respect to the
simulation box XYZ axes, not with respect to the principal axes of the
rigid body itself.  LAMMPS performs the latter calculation internally.

\sphinxAtStartPar
The (vxcm,vycm,vzcm) values are the velocity of the center of mass.
The (lx,ly,lz) values are the angular momentum of the body.  The
(vxcm,vycm,vzcm) and (lx,ly,lz) values can simply be set to 0 if you
wish the body to have no initial motion.

\sphinxAtStartPar
The (ixcm,iycm,izcm) values are the image flags of the center of mass
of the body.  For periodic dimensions, they specify which image of the
simulation box the body is considered to be in.  An image of 0 means
it is inside the box as defined.  A value of 2 means add 2 box lengths
to get the true value.  A value of \sphinxhyphen{}1 means subtract 1 box length to
get the true value.  LAMMPS updates these flags as the rigid bodies
cross periodic boundaries during the simulation.

\begin{sphinxadmonition}{note}{Note:}
\sphinxAtStartPar
If you use the \sphinxstyleemphasis{infile} keyword and write restart
files during a simulation, then each time a restart file is written,
the fix also write an auxiliary restart file with the name
rfile.rigid, where “rfile” is the name of the restart file,
e.g. tmp.restart.10000 and tmp.restart.10000.rigid.  This auxiliary
file is in the same format described above.  Thus it can be used in a
new input script that restarts the run and re\sphinxhyphen{}specifies a rigid fix
using an \sphinxstyleemphasis{infile} keyword and the appropriate filename.  Note that the
auxiliary file will contain one line for every rigid body, even if the
original file only listed a subset of the rigid bodies.
\end{sphinxadmonition}


\bigskip\hrule\bigskip



\subsection{Restart, fix\_modify, output, run start/stop, minimize info}
\label{\detokenize{fix_rigid_meso:restart-fix-modify-output-run-start-stop-minimize-info}}
\sphinxAtStartPar
No information is written to {\hyperref[\detokenize{restart::doc}]{\sphinxcrossref{\DUrole{doc}{binary restart files}}}}.
If the \sphinxstyleemphasis{infile} keyword is used, an auxiliary file is written out
with rigid body information each time a restart file is written, as
explained above for the \sphinxstyleemphasis{infile} keyword.

\sphinxAtStartPar
None of the {\hyperref[\detokenize{fix_modify::doc}]{\sphinxcrossref{\DUrole{doc}{fix\_modify}}}} options are relevant to this
fix.

\sphinxAtStartPar
This fix computes a global array of values which can be accessed by
various {\hyperref[\detokenize{Howto_output::doc}]{\sphinxcrossref{\DUrole{doc}{output commands}}}}.

\sphinxAtStartPar
The number of rows in the array is equal to the number of rigid
bodies.  The number of columns is 28.  Thus for each rigid body, 28
values are stored: the xyz coords of the center of mass (COM), the xyz
components of the COM velocity, the xyz components of the force acting
on the COM, the components of the 4\sphinxhyphen{}vector quaternion representing the
orientation of the rigid body, the xyz components of the angular velocity
of the body around its COM, the xyz components of the torque acting on the
COM, the 3 principal components of the moment of inertia, the xyz components
of the angular momentum of the body around its COM, and the xyz image
flags of the COM.

\sphinxAtStartPar
The center of mass (COM) for each body is similar to unwrapped
coordinates written to a dump file.  It will always be inside (or
slightly outside) the simulation box.  The image flags have the same
meaning as image flags for particle positions (see the “dump” command).
This means you can calculate the unwrapped COM by applying the image
flags to the COM, the same as when unwrapped coordinates are written
to a dump file.

\sphinxAtStartPar
The force and torque values in the array are not affected by the
\sphinxstyleemphasis{force} and \sphinxstyleemphasis{torque} keywords in the fix rigid command; they reflect
values before any changes are made by those keywords.

\sphinxAtStartPar
The ordering of the rigid bodies (by row in the array) is as follows.
For the \sphinxstyleemphasis{single} keyword there is just one rigid body.  For the
\sphinxstyleemphasis{molecule} keyword, the bodies are ordered by ascending molecule ID.
For the \sphinxstyleemphasis{group} keyword, the list of group IDs determines the ordering
of bodies.

\sphinxAtStartPar
The array values calculated by this fix are “intensive”, meaning they
are independent of the number of particles in the simulation.

\sphinxAtStartPar
No parameter of this fix can be used with the \sphinxstyleemphasis{start/stop} keywords of
the {\hyperref[\detokenize{run::doc}]{\sphinxcrossref{\DUrole{doc}{run}}}} command.

\sphinxAtStartPar
This fix is not invoked during {\hyperref[\detokenize{minimize::doc}]{\sphinxcrossref{\DUrole{doc}{energy minimization}}}}.


\bigskip\hrule\bigskip



\subsection{Restrictions}
\label{\detokenize{fix_rigid_meso:restrictions}}
\sphinxAtStartPar
This fix is part of the DPD\sphinxhyphen{}SMOOTH package and also depends on the RIGID
package.  It is only enabled if LAMMPS was built with both packages. See
the {\hyperref[\detokenize{Build_package::doc}]{\sphinxcrossref{\DUrole{doc}{Build package}}}} page for more info.

\sphinxAtStartPar
This fix requires that atoms store density and internal energy as
defined by the {\hyperref[\detokenize{atom_style::doc}]{\sphinxcrossref{\DUrole{doc}{atom\_style sph}}}} command.

\sphinxAtStartPar
All particles in the group must be mesoscopic SPH/SDPD particles.

\sphinxAtStartPar
\DUrole{versionmodified,changed}{Changed in version 29Aug2024.}

\sphinxAtStartPar
This fix is incompatible with deformation controls that remap velocity,
for instance the \sphinxstyleemphasis{remap v} option of {\hyperref[\detokenize{fix_deform::doc}]{\sphinxcrossref{\DUrole{doc}{fix deform}}}}.


\subsection{Related commands}
\label{\detokenize{fix_rigid_meso:related-commands}}
\sphinxAtStartPar
{\hyperref[\detokenize{fix_meso_move::doc}]{\sphinxcrossref{\DUrole{doc}{fix meso/move}}}}, {\hyperref[\detokenize{fix_rigid::doc}]{\sphinxcrossref{\DUrole{doc}{fix rigid}}}},
{\hyperref[\detokenize{neigh_modify::doc}]{\sphinxcrossref{\DUrole{doc}{neigh\_modify exclude}}}}


\subsection{Default}
\label{\detokenize{fix_rigid_meso:default}}
\sphinxAtStartPar
The option defaults are force * on on on and torque * on on on,
meaning all rigid bodies are acted on by center\sphinxhyphen{}of\sphinxhyphen{}mass force and
torque. Also reinit = yes.


\bigskip\hrule\bigskip

\phantomsection\label{\detokenize{fix_rigid_meso:miller}}
\sphinxAtStartPar
\sphinxstylestrong{(Miller)} Miller, Eleftheriou, Pattnaik, Ndirango, and Newns,
J Chem Phys, 116, 8649 (2002).

\sphinxstepscope

\index{fix rx@\spxentry{fix rx}}\ignorespaces 
\index{fix rx/kk@\spxentry{fix rx/kk}}\ignorespaces 

\section{fix rx command}
\label{\detokenize{fix_rx:fix-rx-command}}\label{\detokenize{fix_rx:index-1}}\label{\detokenize{fix_rx:index-0}}\label{\detokenize{fix_rx::doc}}
\sphinxAtStartPar
Accelerator Variants: \sphinxstyleemphasis{rx/kk}


\subsection{Syntax}
\label{\detokenize{fix_rx:syntax}}
\begin{sphinxVerbatim}[commandchars=\\\{\}]
\PYG{k}{fix }\PYG{n+nv+nvIdentifier}{ID}\PYG{+w}{ }\PYG{n+nv+nvIdentifier}{group\PYGZhy{}ID}\PYG{+w}{ }\PYG{n}{rx}\PYG{+w}{ }\PYG{n}{file}\PYG{+w}{ }\PYG{n}{localTemp}\PYG{+w}{ }\PYG{n}{matrix}\PYG{+w}{ }\PYG{n}{solver}\PYG{+w}{ }\PYG{n}{minSteps}\PYG{+w}{ }\PYG{n}{...}
\end{sphinxVerbatim}
\begin{itemize}
\item {} 
\sphinxAtStartPar
ID, group\sphinxhyphen{}ID are documented in {\hyperref[\detokenize{fix::doc}]{\sphinxcrossref{\DUrole{doc}{fix}}}} command

\item {} 
\sphinxAtStartPar
rx = style name of this fix command

\item {} 
\sphinxAtStartPar
file = filename containing the reaction kinetic equations and Arrhenius parameters

\item {} 
\sphinxAtStartPar
localTemp = \sphinxstyleemphasis{none,lucy} = no local temperature averaging or local temperature defined through Lucy weighting function

\item {} 
\sphinxAtStartPar
matrix = \sphinxstyleemphasis{sparse, dense} format for the stoichiometric matrix

\item {} 
\sphinxAtStartPar
solver = \sphinxstyleemphasis{lammps\_rk4,rkf45} = rk4 is an explicit fourth order Runge\sphinxhyphen{}Kutta method; rkf45 is an adaptive fourth\sphinxhyphen{}order Runge\sphinxhyphen{}Kutta\sphinxhyphen{}Fehlberg method

\item {} 
\sphinxAtStartPar
minSteps = \# of steps for rk4 solver or minimum \# of steps for rkf45 (rk4 or rkf45)

\item {} 
\sphinxAtStartPar
maxSteps = maximum number of steps for the rkf45 solver (rkf45 only)

\item {} 
\sphinxAtStartPar
relTol = relative tolerance for the rkf45 solver (rkf45 only)

\item {} 
\sphinxAtStartPar
absTol = absolute tolerance for the rkf45 solver (rkf45 only)

\item {} 
\sphinxAtStartPar
diag   = Diagnostics frequency for the rkf45 solver (optional, rkf45 only)

\end{itemize}


\subsection{Examples}
\label{\detokenize{fix_rx:examples}}
\begin{sphinxVerbatim}[commandchars=\\\{\}]
\PYG{k}{fix }\PYG{n+nv+nvIdentifier}{1}\PYG{+w}{ }\PYG{n+nv+nvIdentifier}{all}\PYG{+w}{ }\PYG{n}{rx}\PYG{+w}{ }\PYG{n}{kinetics.rx}\PYG{+w}{ }\PYG{n}{none}\PYG{+w}{ }\PYG{n}{dense}\PYG{+w}{ }\PYG{n}{lammps\PYGZus{}rk4}
\PYG{k}{fix }\PYG{n+nv+nvIdentifier}{1}\PYG{+w}{ }\PYG{n+nv+nvIdentifier}{all}\PYG{+w}{ }\PYG{n}{rx}\PYG{+w}{ }\PYG{n}{kinetics.rx}\PYG{+w}{ }\PYG{n}{none}\PYG{+w}{ }\PYG{n}{sparse}\PYG{+w}{ }\PYG{n}{lammps\PYGZus{}rk4}\PYG{+w}{ }\PYG{l+m}{1}
\PYG{k}{fix }\PYG{n+nv+nvIdentifier}{1}\PYG{+w}{ }\PYG{n+nv+nvIdentifier}{all}\PYG{+w}{ }\PYG{n}{rx}\PYG{+w}{ }\PYG{n}{kinetics.rx}\PYG{+w}{ }\PYG{n}{lucy}\PYG{+w}{ }\PYG{n}{sparse}\PYG{+w}{ }\PYG{n}{lammps\PYGZus{}rk4}\PYG{+w}{ }\PYG{l+m}{10}
\PYG{k}{fix }\PYG{n+nv+nvIdentifier}{1}\PYG{+w}{ }\PYG{n+nv+nvIdentifier}{all}\PYG{+w}{ }\PYG{n}{rx}\PYG{+w}{ }\PYG{n}{kinetics.rx}\PYG{+w}{ }\PYG{n}{none}\PYG{+w}{ }\PYG{n}{dense}\PYG{+w}{ }\PYG{n}{rkf45}\PYG{+w}{ }\PYG{l+m}{1}\PYG{+w}{ }\PYG{l+m}{100}\PYG{+w}{ }\PYG{l+m}{1e\PYGZhy{}6}\PYG{+w}{ }\PYG{l+m}{1e\PYGZhy{}8}
\PYG{k}{fix }\PYG{n+nv+nvIdentifier}{1}\PYG{+w}{ }\PYG{n+nv+nvIdentifier}{all}\PYG{+w}{ }\PYG{n}{rx}\PYG{+w}{ }\PYG{n}{kinetics.rx}\PYG{+w}{ }\PYG{n}{none}\PYG{+w}{ }\PYG{n}{dense}\PYG{+w}{ }\PYG{n}{rkf45}\PYG{+w}{ }\PYG{l+m}{1}\PYG{+w}{ }\PYG{l+m}{100}\PYG{+w}{ }\PYG{l+m}{1e\PYGZhy{}6}\PYG{+w}{ }\PYG{l+m}{1e\PYGZhy{}8}\PYG{+w}{ }\PYG{o}{\PYGZhy{}}\PYG{l+m}{1}
\end{sphinxVerbatim}


\subsection{Description}
\label{\detokenize{fix_rx:description}}
\sphinxAtStartPar
Fix \sphinxstyleemphasis{rx} solves the reaction kinetic ODEs for a given reaction set that is
defined within the file associated with this command.

\sphinxAtStartPar
For a general reaction such that
\begin{equation*}
\begin{split}\nu_{A}A + \nu_{B}B \rightarrow \nu_{C}C\end{split}
\end{equation*}
\sphinxAtStartPar
the reaction rate equation is defined to be of the form
\begin{equation*}
\begin{split}r = k(T)[A]^{\nu_{A}}[B]^{\nu_{B}}\end{split}
\end{equation*}
\sphinxAtStartPar
In the current implementation, the exponents are defined to be equal
to the stoichiometric coefficients.  A given reaction set consisting
of \sphinxstyleemphasis{n} reaction equations will contain a total of \sphinxstyleemphasis{m} species.  A set
of \sphinxstyleemphasis{m} ordinary differential equations (ODEs) that describe the change
in concentration of a given species as a function of time are then
constructed based on the \sphinxstyleemphasis{n} reaction rate equations.

\sphinxAtStartPar
The ODE systems are solved over the full DPD timestep \sphinxstyleemphasis{dt} using either a fourth
order Runge\sphinxhyphen{}Kutta \sphinxstyleemphasis{rk4} method with a fixed step\sphinxhyphen{}size \sphinxstyleemphasis{h}, specified
by the \sphinxstyleemphasis{lammps\_rk4} keyword, or a fourth order Runge\sphinxhyphen{}Kutta\sphinxhyphen{}Fehlberg (rkf45) method
with an adaptive step\sphinxhyphen{}size for \sphinxstyleemphasis{h}. The number of ODE steps per DPD timestep
for the rk4 method is optionally specified immediately after the rk4
keyword. The ODE step\sphinxhyphen{}size is set as \sphinxstyleemphasis{dt/num\_steps}. Smaller
step\sphinxhyphen{}sizes tend to yield more accurate results but there is not
control on the error. For error control, use the rkf45 ODE solver.

\sphinxAtStartPar
The rkf45 method adjusts the step\sphinxhyphen{}size so that the local truncation error is held
within the specified absolute and relative tolerances. The initial step\sphinxhyphen{}size \sphinxstyleemphasis{h0}
can be specified by the user or estimated internally. It is recommended that the user
specify \sphinxstyleemphasis{h0} since this will generally reduced the number of ODE integration steps
required. \sphinxstyleemphasis{h0} is defined as \sphinxstyleemphasis{dt / min\_steps} if min\_steps \textgreater{}= 1. If min\_steps == 0,
\sphinxstyleemphasis{h0} is estimated such that an explicit Euler method would likely produce
an acceptable solution. This is generally overly conservative for the fourth\sphinxhyphen{}order
method and users are advised to specify \sphinxstyleemphasis{h0} as some fraction of the DPD timestep.
For small DPD timesteps, only one step may be necessary depending upon the tolerances.
Note that more than min\_steps ODE steps may be taken depending upon the ODE stiffness
but no more than max\_steps will be taken. If max\_steps is reached, an error warning
is printed and the simulation is stopped.

\sphinxAtStartPar
After each ODE step, the solution error \sphinxstyleemphasis{e} is tested and weighted using the absTol
and relTol values. The error vector is weighted as \sphinxstyleemphasis{e} / (relTol * |\sphinxstyleemphasis{u}| + absTol)
where \sphinxstyleemphasis{u} is the solution vector. If the norm of the error is \textless{}= 1, the solution is
accepted, \sphinxstyleemphasis{h} is increased by a proportional amount, and the next ODE step is begun.
Otherwise, \sphinxstyleemphasis{h} is shrunk and the ODE step is repeated.

\sphinxAtStartPar
Run\sphinxhyphen{}time diagnostics are available for the rkf45 ODE solver. The frequency
(in timesteps) that diagnostics are reported is controlled by the last (optional)
12th argument. A negative frequency means that diagnostics are reported once at the
end of each run. A positive value N means that the diagnostics are reported once
per N timesteps.

\sphinxAtStartPar
The diagnostics report the average \# of integrator steps and RHS function evaluations
and run\sphinxhyphen{}time per ODE as well as the average/RMS/min/max per process. If the
reporting frequency is 1, the RMS/min/max per ODE are also reported. The per ODE
statistics can be used to adjust the tolerance and min/max step parameters. The
statistics per MPI process can be useful to examine any load imbalance caused by the
adaptive ODE solver. (Some DPD particles can take longer to solve than others. This
can lead to an imbalance across the MPI processes.)


\bigskip\hrule\bigskip


\sphinxAtStartPar
The filename specifies a file that contains the entire set of reaction
kinetic equations and corresponding Arrhenius parameters.  The format of
this file is described below.

\sphinxAtStartPar
There is no restriction on the total number or reaction equations that
are specified.  The species names are arbitrary string names that are
associated with the species concentrations.  Each species in a given
reaction must be preceded by it’s stoichiometric coefficient.  The
only delimiters that are recognized between the species are either a
\sphinxstyleemphasis{+} or \sphinxstyleemphasis{=} character.  The \sphinxstyleemphasis{=} character corresponds to an
irreversible reaction.  After specifying the reaction, the reaction
rate constant is determined through the temperature dependent
Arrhenius equation:
\begin{equation*}
\begin{split}k = AT^{n}e^{\frac{-E_{a}}{k_{B}T}}\end{split}
\end{equation*}
\sphinxAtStartPar
where \sphinxstyleemphasis{A} is the Arrhenius factor in time units or concentration/time
units, \sphinxstyleemphasis{n} is the unitless exponent of the temperature dependence, and
\(E_a\) is the activation energy in energy units.  The temperature
dependence can be removed by specifying the exponent as zero.

\sphinxAtStartPar
The internal temperature of the coarse\sphinxhyphen{}grained particles can be used
in constructing the reaction rate constants at every DPD timestep by
specifying the keyword \sphinxstyleemphasis{none}.  Alternatively, the keyword \sphinxstyleemphasis{lucy} can
be specified to compute a local\sphinxhyphen{}average particle internal temperature
for use in the reaction rate constant expressions.  The local\sphinxhyphen{}average
particle internal temperature is defined as:
\begin{equation*}
\begin{split}\theta_i^{-1} = \frac{\sum_{j=1}\omega_{Lucy}\left(r_{ij}\right)\theta_j^{-1}}{\sum_{j=1}\omega_{Lucy}\left(r_{ij}\right)}\end{split}
\end{equation*}
\sphinxAtStartPar
where the Lucy function is expressed as:
\begin{equation*}
\begin{split}\omega_{Lucy}\left(r_{ij}\right) = \left( 1 + \frac{3r_{ij}}{r_c} \right) \left( 1 - \frac{r_{ij}}{r_c} \right)^3\end{split}
\end{equation*}
\sphinxAtStartPar
The self\sphinxhyphen{}particle interaction is included in the above equation.

\sphinxAtStartPar
The stoichiometric coefficients for the reaction mechanism are stored
in either a sparse or dense matrix format. The dense matrix should only be
used for small reaction mechanisms. The sparse matrix should be used when there
are many reactions (e.g., more than 5). This allows the number of reactions and
species to grow while keeping the computational cost tractable. The matrix
format can be specified as using either the \sphinxstyleemphasis{sparse} or \sphinxstyleemphasis{dense} keywords.
If all stoichiometric coefficients for a reaction are small integers (whole
numbers \textless{}= 3), a fast exponential function is used. This can save significant
computational time so users are encouraged to use integer coefficients
where possible.


\bigskip\hrule\bigskip


\sphinxAtStartPar
The format of a tabulated file is as follows (without the
parenthesized comments):

\begin{sphinxVerbatim}[commandchars=\\\{\}]
\PYGZsh{} Rxn equations and parameters                                               (one or more comment or blank lines)

1.0 hcn + 1.0 no2 = 1.0  no + 0.5 n2  + 0.5 h2 + 1.0 co   2.49E+01 0.0 1.34  (rxn equation, A, n, Ea)
1.0 hcn + 1.0  no = 1.0  co + 1.0 n2  + 0.5 h2            2.16E+00 0.0 1.52
...
1.0  no + 1.0  co = 0.5  n2 + 1.0 co2                     1.66E+06 0.0 0.69
\end{sphinxVerbatim}

\sphinxAtStartPar
A section begins with a non\sphinxhyphen{}blank line whose first character is not a
“\#”; blank lines or lines starting with “\#” can be used as comments
between sections.

\sphinxAtStartPar
Following a blank line, the next N lines list the N reaction
equations.  Each species within the reaction equation is specified
through its stoichiometric coefficient and a species tag.  Reactant
species are specified on the left\sphinxhyphen{}hand side of the equation and
product species are specified on the right\sphinxhyphen{}hand side of the equation.
After specifying the reactant and product species, the final three
arguments of each line represent the Arrhenius parameter \sphinxstyleemphasis{A}, the
temperature exponent \sphinxstyleemphasis{n}, and the activation energy \sphinxstyleemphasis{Ea}.

\sphinxAtStartPar
Note that the species tags that are defined in the reaction equations
are used by the {\hyperref[\detokenize{fix_eos_table_rx::doc}]{\sphinxcrossref{\DUrole{doc}{fix eos/table/rx}}}} command to
define the thermodynamic properties of each species.  Furthermore, the
number of species molecules (i.e., concentration) can be specified
either with the {\hyperref[\detokenize{set::doc}]{\sphinxcrossref{\DUrole{doc}{set}}}} command using the “d\_” prefix or by
reading directly the concentrations from a data file.  For the latter
case, the {\hyperref[\detokenize{read_data::doc}]{\sphinxcrossref{\DUrole{doc}{read\_data}}}} command with the fix keyword
should be specified, where the fix\sphinxhyphen{}ID will be the “fix rx\textasciigrave{}ID with a \textless{}SPECIES”\textgreater{}\textasciigrave{}\_ suffix, e.g.

\sphinxAtStartPar
fix          foo all rx reaction.file …
read\_data    data.dpd fix foo\_SPECIES NULL Species


\bigskip\hrule\bigskip


\sphinxAtStartPar
Styles with a \sphinxstyleemphasis{gpu}, \sphinxstyleemphasis{intel}, \sphinxstyleemphasis{kk}, \sphinxstyleemphasis{omp}, or \sphinxstyleemphasis{opt} suffix are
functionally the same as the corresponding style without the suffix.
They have been optimized to run faster, depending on your available
hardware, as discussed on the {\hyperref[\detokenize{Speed_packages::doc}]{\sphinxcrossref{\DUrole{doc}{Accelerator packages}}}}
page.  The accelerated styles take the same arguments and should
produce the same results, except for round\sphinxhyphen{}off and precision issues.

\sphinxAtStartPar
These accelerated styles are part of the GPU, INTEL, KOKKOS,
OPENMP, and OPT packages, respectively.  They are only enabled if
LAMMPS was built with those packages.  See the {\hyperref[\detokenize{Build_package::doc}]{\sphinxcrossref{\DUrole{doc}{Build package}}}} page for more info.

\sphinxAtStartPar
You can specify the accelerated styles explicitly in your input script
by including their suffix, or you can use the {\hyperref[\detokenize{Run_options::doc}]{\sphinxcrossref{\DUrole{doc}{\sphinxhyphen{}suffix command\sphinxhyphen{}line switch}}}} when you invoke LAMMPS, or you can use the
{\hyperref[\detokenize{suffix::doc}]{\sphinxcrossref{\DUrole{doc}{suffix}}}} command in your input script.

\sphinxAtStartPar
See the {\hyperref[\detokenize{Speed_packages::doc}]{\sphinxcrossref{\DUrole{doc}{Accelerator packages}}}} page for more
instructions on how to use the accelerated styles effectively.


\bigskip\hrule\bigskip



\subsection{Restrictions}
\label{\detokenize{fix_rx:restrictions}}
\sphinxAtStartPar
This command is part of the DPD\sphinxhyphen{}REACT package.  It is only enabled if
LAMMPS was built with that package.  See the {\hyperref[\detokenize{Build_package::doc}]{\sphinxcrossref{\DUrole{doc}{Build package}}}} page for more info.

\sphinxAtStartPar
This command also requires use of the {\hyperref[\detokenize{atom_style::doc}]{\sphinxcrossref{\DUrole{doc}{atom\_style dpd}}}}
command.

\sphinxAtStartPar
This command can only be used with a constant energy or constant
enthalpy DPD simulation.


\subsection{Related commands}
\label{\detokenize{fix_rx:related-commands}}
\sphinxAtStartPar
{\hyperref[\detokenize{fix_eos_table_rx::doc}]{\sphinxcrossref{\DUrole{doc}{fix eos/table/rx}}}},
{\hyperref[\detokenize{fix_shardlow::doc}]{\sphinxcrossref{\DUrole{doc}{fix shardlow}}}},
{\hyperref[\detokenize{pair_dpd_fdt::doc}]{\sphinxcrossref{\DUrole{doc}{pair dpd/fdt/energy}}}}


\subsection{Default}
\label{\detokenize{fix_rx:default}}
\sphinxAtStartPar
none

\sphinxstepscope

\index{fix saed/vtk@\spxentry{fix saed/vtk}}\ignorespaces 

\section{fix saed/vtk command}
\label{\detokenize{fix_saed_vtk:fix-saed-vtk-command}}\label{\detokenize{fix_saed_vtk:index-0}}\label{\detokenize{fix_saed_vtk::doc}}

\subsection{Syntax}
\label{\detokenize{fix_saed_vtk:syntax}}
\begin{sphinxVerbatim}[commandchars=\\\{\}]
\PYG{k}{fix }\PYG{n+nv+nvIdentifier}{ID}\PYG{+w}{ }\PYG{n+nv+nvIdentifier}{group\PYGZhy{}ID}\PYG{+w}{ }\PYG{n}{saed}\PYG{o}{/}\PYG{n}{vtk}\PYG{+w}{ }\PYG{n}{Nevery}\PYG{+w}{ }\PYG{n}{Nrepeat}\PYG{+w}{ }\PYG{n}{Nfreak}\PYG{+w}{ }\PYG{n}{c\PYGZus{}ID}\PYG{+w}{ }\PYG{n}{attribute}\PYG{+w}{ }\PYG{n}{args}\PYG{+w}{ }\PYG{n}{...}\PYG{+w}{ }\PYG{n}{keyword}\PYG{+w}{ }\PYG{n}{args}\PYG{+w}{ }\PYG{n}{...}
\end{sphinxVerbatim}
\begin{itemize}
\item {} 
\sphinxAtStartPar
ID, group\sphinxhyphen{}ID are documented in {\hyperref[\detokenize{fix::doc}]{\sphinxcrossref{\DUrole{doc}{fix}}}} command

\item {} 
\sphinxAtStartPar
saed/vtk = style name of this fix command

\item {} 
\sphinxAtStartPar
Nevery = use input values every this many timesteps

\item {} 
\sphinxAtStartPar
Nrepeat = \# of times to use input values for calculating averages

\item {} 
\sphinxAtStartPar
Nfreq = calculate averages every this many timesteps

\item {} 
\sphinxAtStartPar
c\_ID = saed compute ID
\begin{sphinxalltt}
keyword = \sphinxstyleemphasis{file} or \sphinxstyleemphasis{ave} or \sphinxstyleemphasis{start} or \sphinxstyleemphasis{file} or \sphinxstyleemphasis{overwrite}:l
  \sphinxstyleemphasis{ave} args = \sphinxstyleemphasis{one} or \sphinxstyleemphasis{running} or \sphinxstyleemphasis{window M}
    one = output a new average value every Nfreq steps
    running = output cumulative average of all previous Nfreq steps
    window M = output average of M most recent Nfreq steps
  \sphinxstyleemphasis{start} args = Nstart
    Nstart = start averaging on this timestep
  \sphinxstyleemphasis{file} arg = filename
    filename = name of file to output time averages to
\end{sphinxalltt}

\end{itemize}


\subsection{Examples}
\label{\detokenize{fix_saed_vtk:examples}}
\begin{sphinxVerbatim}[commandchars=\\\{\}]
\PYG{k}{compute }\PYG{n+nv+nvIdentifier}{1}\PYG{+w}{ }\PYG{n+nv+nvIdentifier}{all}\PYG{+w}{ }\PYG{n}{saed}\PYG{+w}{ }\PYG{l+m}{0.0251}\PYG{+w}{ }\PYG{n}{Al}\PYG{+w}{ }\PYG{n}{O}\PYG{+w}{ }\PYG{n}{Kmax}\PYG{+w}{ }\PYG{l+m}{1.70}\PYG{+w}{ }\PYG{n}{Zone}\PYG{+w}{ }\PYG{l+m}{0}\PYG{+w}{ }\PYG{l+m}{0}\PYG{+w}{ }\PYG{l+m}{1}\PYG{+w}{ }\PYG{n}{dR\PYGZus{}Ewald}\PYG{+w}{ }\PYG{l+m}{0.01}\PYG{+w}{ }\PYG{n}{c}\PYG{+w}{ }\PYG{l+m}{0.5}\PYG{+w}{ }\PYG{l+m}{0.5}\PYG{+w}{ }\PYG{l+m}{0.5}
\PYG{k}{compute }\PYG{n+nv+nvIdentifier}{2}\PYG{+w}{ }\PYG{n+nv+nvIdentifier}{all}\PYG{+w}{ }\PYG{n}{saed}\PYG{+w}{ }\PYG{l+m}{0.0251}\PYG{+w}{ }\PYG{n}{Ni}\PYG{+w}{ }\PYG{n}{Kmax}\PYG{+w}{ }\PYG{l+m}{1.70}\PYG{+w}{ }\PYG{n}{Zone}\PYG{+w}{ }\PYG{l+m}{0}\PYG{+w}{ }\PYG{l+m}{0}\PYG{+w}{ }\PYG{l+m}{0}\PYG{+w}{ }\PYG{n}{c}\PYG{+w}{ }\PYG{l+m}{0.05}\PYG{+w}{ }\PYG{l+m}{0.05}\PYG{+w}{ }\PYG{l+m}{0.05}\PYG{+w}{ }\PYG{n}{manual}\PYG{+w}{ }\PYG{n}{echo}

\PYG{k}{fix }\PYG{n+nv+nvIdentifier}{1}\PYG{+w}{ }\PYG{n+nv+nvIdentifier}{all}\PYG{+w}{ }\PYG{n}{saed}\PYG{o}{/}\PYG{n}{vtk}\PYG{+w}{ }\PYG{l+m}{1}\PYG{+w}{ }\PYG{l+m}{1}\PYG{+w}{ }\PYG{l+m}{1}\PYG{+w}{ }\PYG{n}{c\PYGZus{}1}\PYG{+w}{ }\PYG{n}{file}\PYG{+w}{ }\PYG{n}{Al2O3\PYGZus{}001.saed}
\PYG{k}{fix }\PYG{n+nv+nvIdentifier}{2}\PYG{+w}{ }\PYG{n+nv+nvIdentifier}{all}\PYG{+w}{ }\PYG{n}{saed}\PYG{o}{/}\PYG{n}{vtk}\PYG{+w}{ }\PYG{l+m}{1}\PYG{+w}{ }\PYG{l+m}{1}\PYG{+w}{ }\PYG{l+m}{1}\PYG{+w}{ }\PYG{n}{c\PYGZus{}2}\PYG{+w}{ }\PYG{n}{file}\PYG{+w}{ }\PYG{n}{Ni\PYGZus{}000.saed}
\end{sphinxVerbatim}


\subsection{Description}
\label{\detokenize{fix_saed_vtk:description}}
\sphinxAtStartPar
Time average computed intensities from {\hyperref[\detokenize{compute_saed::doc}]{\sphinxcrossref{\DUrole{doc}{compute saed}}}} and
write output to a file in the third generation vtk image data format for
visualization directly in parallelized visualization software packages
like ParaView and VisIt. Note that if no time averaging is done, this
command can be used as a convenient way to simply output diffraction
intensities at a single snapshot.

\sphinxAtStartPar
To produce output in the image data vtk format ghost data is added
outside the \sphinxstyleemphasis{Kmax} range assigned in the compute saed. The ghost data is
assigned a value of \sphinxhyphen{}1 and can be removed setting a minimum isovolume
of 0 within the visualization software. SAED images can be created by
visualizing a spherical slice of the data that is centered at
R\_Ewald*{[}h k l{]}/norm({[}h k l{]}), where R\_Ewald=1/lambda.

\sphinxAtStartPar
The group specified within this command is ignored. However, note that
specified values may represent calculations performed by saed computes
which store their own “group” definitions.

\sphinxAtStartPar
Fix saed/vtk is designed to work only with {\hyperref[\detokenize{compute_saed::doc}]{\sphinxcrossref{\DUrole{doc}{compute saed}}}}
values, e.g.

\begin{sphinxVerbatim}[commandchars=\\\{\}]
\PYG{k}{compute }\PYG{n+nv+nvIdentifier}{3}\PYG{+w}{ }\PYG{n+nv+nvIdentifier}{top}\PYG{+w}{ }\PYG{n}{saed}\PYG{+w}{ }\PYG{l+m}{0.0251}\PYG{+w}{ }\PYG{n}{Al}\PYG{+w}{ }\PYG{n}{O}
\PYG{k}{fix }\PYG{n+nv+nvIdentifier}{saed}\PYG{o}{/}\PYG{n}{vtk}\PYG{+w}{ }\PYG{l+m}{1}\PYG{+w}{ }\PYG{l+m}{1}\PYG{+w}{ }\PYG{l+m}{1}\PYG{+w}{ }\PYG{n}{c\PYGZus{}3}\PYG{+w}{ }\PYG{n}{file}\PYG{+w}{ }\PYG{n}{Al2O3\PYGZus{}001.saed}
\end{sphinxVerbatim}


\bigskip\hrule\bigskip


\sphinxAtStartPar
The \sphinxstyleemphasis{Nevery}, \sphinxstyleemphasis{Nrepeat}, and \sphinxstyleemphasis{Nfreq} arguments specify on what
timesteps the input values will be used in order to contribute to the
average.  The final averaged quantities are generated on timesteps
that are a multiple of \sphinxstyleemphasis{Nfreq}.  The average is over \sphinxstyleemphasis{Nrepeat}
quantities, computed in the preceding portion of the simulation every
\sphinxstyleemphasis{Nevery} timesteps.  \sphinxstyleemphasis{Nfreq} must be a multiple of \sphinxstyleemphasis{Nevery} and
\sphinxstyleemphasis{Nevery} must be non\sphinxhyphen{}zero even if \sphinxstyleemphasis{Nrepeat} is 1.
Also, the timesteps
contributing to the average value cannot overlap,
i.e. Nrepeat*Nevery can not exceed Nfreq.

\sphinxAtStartPar
For example, if Nevery=2, Nrepeat=6, and Nfreq=100, then values on
timesteps 90,92,94,96,98,100 will be used to compute the final average
on timestep 100.  Similarly for timesteps 190,192,194,196,198,200 on
timestep 200, etc.  If Nrepeat=1 and Nfreq = 100, then no time
averaging is done; values are simply generated on timesteps
100,200,etc.


\bigskip\hrule\bigskip


\sphinxAtStartPar
The output for fix ave/time/saed is a file written with the third generation
vtk image data formatting.  The filename assigned by the \sphinxstyleemphasis{file} keyword is
appended with \_N.vtk where N is an index (0,1,2…) to account for multiple
diffraction intensity outputs.

\sphinxAtStartPar
By default the header contains the following information (with example data):

\begin{sphinxVerbatim}[commandchars=\\\{\}]
\PYGZsh{} vtk DataFile Version 3.0 c\PYGZus{}SAED
Image data set
ASCII
DATASET STRUCTURED\PYGZus{}POINTS
DIMENSIONS 337 219 209
ASPECT\PYGZus{}RATIO 0.00507953 0.00785161 0.00821458
ORIGIN \PYGZhy{}0.853361 \PYGZhy{}0.855826 \PYGZhy{}0.854316
POINT\PYGZus{}DATA 15424827
SCALARS intensity float
LOOKUP\PYGZus{}TABLE default
...data
\end{sphinxVerbatim}

\sphinxAtStartPar
In this example, kspace is sampled across a 337 x 219 x 209 point mesh
where the mesh spacing is approximately 0.005, 0.007, and 0.008
inv(length) units in the k1, k2, and k3 directions, respectively.
The data is shifted by \sphinxhyphen{}0.85, \sphinxhyphen{}0.85, \sphinxhyphen{}0.85 inv(length) units so that
the origin will lie at 0, 0, 0.   Here, 15,424,827 kspace points are
sampled in total.


\bigskip\hrule\bigskip


\sphinxAtStartPar
Additional optional keywords also affect the operation of this fix.

\sphinxAtStartPar
The \sphinxstyleemphasis{ave} keyword determines how the values produced every \sphinxstyleemphasis{Nfreq}
steps are averaged with values produced on previous steps that were
multiples of \sphinxstyleemphasis{Nfreq}, before they are accessed by another output
command or written to a file.

\sphinxAtStartPar
If the \sphinxstyleemphasis{ave} setting is \sphinxstyleemphasis{one}, then the values produced on timesteps
that are multiples of \sphinxstyleemphasis{Nfreq} are independent of each other; they are
output as\sphinxhyphen{}is without further averaging.

\sphinxAtStartPar
If the \sphinxstyleemphasis{ave} setting is \sphinxstyleemphasis{running}, then the values produced on
timesteps that are multiples of \sphinxstyleemphasis{Nfreq} are summed and averaged in a
cumulative sense before being output.  Each output value is thus the
average of the value produced on that timestep with all preceding
values.  This running average begins when the fix is defined; it can
only be restarted by deleting the fix via the {\hyperref[\detokenize{unfix::doc}]{\sphinxcrossref{\DUrole{doc}{unfix}}}}
command, or by re\sphinxhyphen{}defining the fix by re\sphinxhyphen{}specifying it.

\sphinxAtStartPar
If the \sphinxstyleemphasis{ave} setting is \sphinxstyleemphasis{window}, then the values produced on
timesteps that are multiples of \sphinxstyleemphasis{Nfreq} are summed and averaged within
a moving “window” of time, so that the last M values are used to
produce the output.  E.g. if M = 3 and Nfreq = 1000, then the output
on step 10000 will be the average of the individual values on steps
8000,9000,10000.  Outputs on early steps will average over less than M
values if they are not available.

\sphinxAtStartPar
The \sphinxstyleemphasis{start} keyword specifies what timestep averaging will begin on.
The default is step 0.  Often input values can be 0.0 at time 0, so
setting \sphinxstyleemphasis{start} to a larger value can avoid including a 0.0 in a
running or windowed average.

\sphinxAtStartPar
The \sphinxstyleemphasis{file} keyword allows a filename to be specified.  Every \sphinxstyleemphasis{Nfreq}
steps, the vector of saed intensity data is written to a new file using
the third generation vtk format.  The base of each file is assigned by
the \sphinxstyleemphasis{file} keyword and this string is appended with \_N.vtk where N is
an index (0,1,2…) to account for situations with multiple diffraction
intensity outputs.


\subsection{Restart, fix\_modify, output, run start/stop, minimize info}
\label{\detokenize{fix_saed_vtk:restart-fix-modify-output-run-start-stop-minimize-info}}
\sphinxAtStartPar
This fix is part of the DIFFRACTION package.  It is only enabled if
LAMMPS was built with that package.  See the {\hyperref[\detokenize{Build_package::doc}]{\sphinxcrossref{\DUrole{doc}{Build package}}}} page for more info.

\sphinxAtStartPar
No information about this fix is written to {\hyperref[\detokenize{restart::doc}]{\sphinxcrossref{\DUrole{doc}{binary restart files}}}}.  None of the {\hyperref[\detokenize{fix_modify::doc}]{\sphinxcrossref{\DUrole{doc}{fix\_modify}}}} options are
relevant to this fix.

\sphinxAtStartPar
No parameter of this fix can be used with the \sphinxstyleemphasis{start/stop} keywords of
the {\hyperref[\detokenize{run::doc}]{\sphinxcrossref{\DUrole{doc}{run}}}} command.  This fix is not invoked during
{\hyperref[\detokenize{minimize::doc}]{\sphinxcrossref{\DUrole{doc}{energy minimization}}}}.


\subsection{Restrictions}
\label{\detokenize{fix_saed_vtk:restrictions}}
\sphinxAtStartPar
The attributes for fix\_saed\_vtk must match the values assigned in the
associated {\hyperref[\detokenize{compute_saed::doc}]{\sphinxcrossref{\DUrole{doc}{compute\_saed}}}} command.


\subsection{Related commands}
\label{\detokenize{fix_saed_vtk:related-commands}}
\sphinxAtStartPar
{\hyperref[\detokenize{compute_saed::doc}]{\sphinxcrossref{\DUrole{doc}{compute\_saed}}}}


\subsection{Default}
\label{\detokenize{fix_saed_vtk:default}}
\sphinxAtStartPar
The option defaults are ave = one, start = 0, no file output.


\bigskip\hrule\bigskip

\phantomsection\label{\detokenize{fix_saed_vtk:coleman}}
\sphinxAtStartPar
\sphinxstylestrong{(Coleman)} Coleman, Spearot, Capolungo, MSMSE, 21, 055020
(2013).

\sphinxstepscope

\index{fix setforce@\spxentry{fix setforce}}\ignorespaces 
\index{fix setforce/kk@\spxentry{fix setforce/kk}}\ignorespaces 
\index{fix setforce/spin@\spxentry{fix setforce/spin}}\ignorespaces 

\section{fix setforce command}
\label{\detokenize{fix_setforce:fix-setforce-command}}\label{\detokenize{fix_setforce:index-2}}\label{\detokenize{fix_setforce:index-1}}\label{\detokenize{fix_setforce:index-0}}\label{\detokenize{fix_setforce::doc}}
\sphinxAtStartPar
Accelerator Variants: \sphinxstyleemphasis{setforce/kk}


\section{fix setforce/spin command}
\label{\detokenize{fix_setforce:fix-setforce-spin-command}}

\subsection{Syntax}
\label{\detokenize{fix_setforce:syntax}}
\begin{sphinxVerbatim}[commandchars=\\\{\}]
\PYG{k}{fix }\PYG{n+nv+nvIdentifier}{ID}\PYG{+w}{ }\PYG{n+nv+nvIdentifier}{group\PYGZhy{}ID}\PYG{+w}{ }\PYG{n}{setforce}\PYG{+w}{ }\PYG{n}{fx}\PYG{+w}{ }\PYG{n}{fy}\PYG{+w}{ }\PYG{n}{fz}\PYG{+w}{ }\PYG{n}{keyword}\PYG{+w}{ }\PYG{n}{value}\PYG{+w}{ }\PYG{n}{...}
\end{sphinxVerbatim}
\begin{itemize}
\item {} 
\sphinxAtStartPar
ID, group\sphinxhyphen{}ID are documented in {\hyperref[\detokenize{fix::doc}]{\sphinxcrossref{\DUrole{doc}{fix}}}} command

\item {} 
\sphinxAtStartPar
setforce = style name of this fix command

\item {} 
\sphinxAtStartPar
fx,fy,fz = force component values

\item {} 
\sphinxAtStartPar
any of fx,fy,fz can be a variable (see below)

\item {} 
\sphinxAtStartPar
zero or more keyword/value pairs may be appended to args

\item {} 
\sphinxAtStartPar
keyword = \sphinxstyleemphasis{region}
\begin{sphinxalltt}
\sphinxstyleemphasis{region} value = region\sphinxhyphen{}ID
  region\sphinxhyphen{}ID = ID of region atoms must be in to have added force
\end{sphinxalltt}

\end{itemize}


\subsection{Examples}
\label{\detokenize{fix_setforce:examples}}
\begin{sphinxVerbatim}[commandchars=\\\{\}]
\PYG{k}{fix }\PYG{n+nv+nvIdentifier}{freeze}\PYG{+w}{ }\PYG{n+nv+nvIdentifier}{indenter}\PYG{+w}{ }\PYG{n}{setforce}\PYG{+w}{ }\PYG{l+m}{0.0}\PYG{+w}{ }\PYG{l+m}{0.0}\PYG{+w}{ }\PYG{l+m}{0.0}
\PYG{k}{fix }\PYG{n+nv+nvIdentifier}{2}\PYG{+w}{ }\PYG{n+nv+nvIdentifier}{edge}\PYG{+w}{ }\PYG{n}{setforce}\PYG{+w}{ }\PYG{n}{NULL}\PYG{+w}{ }\PYG{l+m}{0.0}\PYG{+w}{ }\PYG{l+m}{0.0}
\PYG{k}{fix }\PYG{n+nv+nvIdentifier}{1}\PYG{+w}{ }\PYG{n+nv+nvIdentifier}{edge}\PYG{+w}{ }\PYG{n}{setforce}\PYG{o}{/}\PYG{n}{spin}\PYG{+w}{ }\PYG{l+m}{0.0}\PYG{+w}{ }\PYG{l+m}{0.0}\PYG{+w}{ }\PYG{l+m}{0.0}
\PYG{k}{fix }\PYG{n+nv+nvIdentifier}{2}\PYG{+w}{ }\PYG{n+nv+nvIdentifier}{edge}\PYG{+w}{ }\PYG{n}{setforce}\PYG{+w}{ }\PYG{n}{NULL}\PYG{+w}{ }\PYG{l+m}{0.0}\PYG{+w}{ }\PYG{n}{v\PYGZus{}oscillate}
\end{sphinxVerbatim}


\subsection{Description}
\label{\detokenize{fix_setforce:description}}
\sphinxAtStartPar
Set each component of force on each atom in the group to the specified
values fx,fy,fz.  This erases all previously computed forces on the
atom, though additional fixes could add new forces.  This command can
be used to freeze certain atoms in the simulation by zeroing their
force, either for running dynamics or performing an energy
minimization.  For dynamics, this assumes their initial velocity is
also zero.

\sphinxAtStartPar
Any of the fx,fy,fz values can be specified as NULL which means do not
alter the force component in that dimension.

\sphinxAtStartPar
Any of the 3 quantities defining the force components can be specified
as an equal\sphinxhyphen{}style or atom\sphinxhyphen{}style {\hyperref[\detokenize{variable::doc}]{\sphinxcrossref{\DUrole{doc}{variable}}}}, namely \sphinxstyleemphasis{fx},
\sphinxstyleemphasis{fy}, \sphinxstyleemphasis{fz}.  If the value is a variable, it should be specified as
v\_name, where name is the variable name.  In this case, the variable
will be evaluated each timestep, and its value used to determine the
force component.

\sphinxAtStartPar
Equal\sphinxhyphen{}style variables can specify formulas with various mathematical
functions, and include {\hyperref[\detokenize{thermo_style::doc}]{\sphinxcrossref{\DUrole{doc}{thermo\_style}}}} command
keywords for the simulation box parameters and timestep and elapsed
time.  Thus it is easy to specify a time\sphinxhyphen{}dependent force field.

\sphinxAtStartPar
Atom\sphinxhyphen{}style variables can specify the same formulas as equal\sphinxhyphen{}style
variables but can also include per\sphinxhyphen{}atom values, such as atom
coordinates.  Thus it is easy to specify a spatially\sphinxhyphen{}dependent force
field with optional time\sphinxhyphen{}dependence as well.

\sphinxAtStartPar
If the \sphinxstyleemphasis{region} keyword is used, the atom must also be in the
specified geometric {\hyperref[\detokenize{region::doc}]{\sphinxcrossref{\DUrole{doc}{region}}}} in order to have force added
to it.


\bigskip\hrule\bigskip


\sphinxAtStartPar
Style \sphinxstyleemphasis{spin} suffix sets the components of the magnetic precession
vectors instead of the mechanical forces. This also erases all
previously computed magnetic precession vectors on the atom, though
additional magnetic fixes could add new forces.

\sphinxAtStartPar
This command can be used to freeze the magnetic moment of certain
atoms in the simulation by zeroing their precession vector.

\sphinxAtStartPar
All options defined above remain valid, they just apply to the magnetic
precession vectors instead of the forces.


\bigskip\hrule\bigskip


\sphinxAtStartPar
Styles with a \sphinxstyleemphasis{gpu}, \sphinxstyleemphasis{intel}, \sphinxstyleemphasis{kk}, \sphinxstyleemphasis{omp}, or \sphinxstyleemphasis{opt} suffix are
functionally the same as the corresponding style without the suffix.
They have been optimized to run faster, depending on your available
hardware, as discussed on the {\hyperref[\detokenize{Speed_packages::doc}]{\sphinxcrossref{\DUrole{doc}{Accelerator packages}}}}
page.  The accelerated styles take the same arguments and should
produce the same results, except for round\sphinxhyphen{}off and precision issues.

\sphinxAtStartPar
These accelerated styles are part of the GPU, INTEL, KOKKOS,
OPENMP, and OPT packages, respectively.  They are only enabled if
LAMMPS was built with those packages.  See the {\hyperref[\detokenize{Build_package::doc}]{\sphinxcrossref{\DUrole{doc}{Build package}}}} page for more info.

\sphinxAtStartPar
You can specify the accelerated styles explicitly in your input script
by including their suffix, or you can use the {\hyperref[\detokenize{Run_options::doc}]{\sphinxcrossref{\DUrole{doc}{\sphinxhyphen{}suffix command\sphinxhyphen{}line switch}}}} when you invoke LAMMPS, or you can use the
{\hyperref[\detokenize{suffix::doc}]{\sphinxcrossref{\DUrole{doc}{suffix}}}} command in your input script.

\sphinxAtStartPar
See the {\hyperref[\detokenize{Speed_packages::doc}]{\sphinxcrossref{\DUrole{doc}{Accelerator packages}}}} page for more
instructions on how to use the accelerated styles effectively.

\begin{sphinxadmonition}{note}{Note:}
\sphinxAtStartPar
The region keyword is supported by Kokkos, but a Kokkos\sphinxhyphen{}enabled
region must be used. See the region {\hyperref[\detokenize{region::doc}]{\sphinxcrossref{\DUrole{doc}{region}}}} command for
more information.
\end{sphinxadmonition}


\bigskip\hrule\bigskip



\subsection{Restart, fix\_modify, output, run start/stop, minimize info}
\label{\detokenize{fix_setforce:restart-fix-modify-output-run-start-stop-minimize-info}}
\sphinxAtStartPar
No information about this fix is written to {\hyperref[\detokenize{restart::doc}]{\sphinxcrossref{\DUrole{doc}{binary restart files}}}}.

\sphinxAtStartPar
The {\hyperref[\detokenize{fix_modify::doc}]{\sphinxcrossref{\DUrole{doc}{fix\_modify}}}} \sphinxstyleemphasis{respa} option is supported by
this fix. This allows to set at which level of the {\hyperref[\detokenize{run_style::doc}]{\sphinxcrossref{\DUrole{doc}{r\sphinxhyphen{}RESPA}}}}
integrator the fix is setting the forces to the desired values; on all
other levels, the force is set to 0.0 for the atoms in the fix group,
so that setforce values are not counted multiple times. Default is to
to override forces at the outermost level.

\sphinxAtStartPar
This fix computes a global 3\sphinxhyphen{}vector of forces, which can be accessed
by various {\hyperref[\detokenize{Howto_output::doc}]{\sphinxcrossref{\DUrole{doc}{output commands}}}}.  This is the total
force on the group of atoms before the forces on individual atoms are
changed by the fix.  The vector values calculated by this fix are
“extensive”.

\sphinxAtStartPar
No parameter of this fix can be used with the \sphinxstyleemphasis{start/stop} keywords of
the {\hyperref[\detokenize{run::doc}]{\sphinxcrossref{\DUrole{doc}{run}}}} command.

\sphinxAtStartPar
The forces due to this fix are imposed during an energy minimization,
invoked by the {\hyperref[\detokenize{minimize::doc}]{\sphinxcrossref{\DUrole{doc}{minimize}}}} command, but you cannot set
forces to any value besides zero when performing a minimization.  Use
the {\hyperref[\detokenize{fix_addforce::doc}]{\sphinxcrossref{\DUrole{doc}{fix addforce}}}} command if you want to apply a
non\sphinxhyphen{}zero force to atoms during a minimization.


\subsection{Restrictions}
\label{\detokenize{fix_setforce:restrictions}}
\sphinxAtStartPar
Fix \sphinxstyleemphasis{setforce/spin} is part of the SPIN package.  It is only
enabled if LAMMPS was built with that package.  See the {\hyperref[\detokenize{Build_package::doc}]{\sphinxcrossref{\DUrole{doc}{Build
package}}}} page for more info.


\subsection{Related commands}
\label{\detokenize{fix_setforce:related-commands}}
\sphinxAtStartPar
{\hyperref[\detokenize{fix_addforce::doc}]{\sphinxcrossref{\DUrole{doc}{fix addforce}}}}, {\hyperref[\detokenize{fix_aveforce::doc}]{\sphinxcrossref{\DUrole{doc}{fix aveforce}}}}


\subsection{Default}
\label{\detokenize{fix_setforce:default}}
\sphinxAtStartPar
none

\sphinxstepscope

\index{fix sgcmc@\spxentry{fix sgcmc}}\ignorespaces 

\section{fix sgcmc command}
\label{\detokenize{fix_sgcmc:fix-sgcmc-command}}\label{\detokenize{fix_sgcmc:index-0}}\label{\detokenize{fix_sgcmc::doc}}

\subsection{Syntax}
\label{\detokenize{fix_sgcmc:syntax}}
\begin{sphinxVerbatim}[commandchars=\\\{\}]
\PYG{k}{fix }\PYG{n+nv+nvIdentifier}{ID}\PYG{+w}{ }\PYG{n+nv+nvIdentifier}{group\PYGZhy{}ID}\PYG{+w}{ }\PYG{n}{sgcmc}\PYG{+w}{ }\PYG{n}{every\PYGZus{}nsteps}\PYG{+w}{ }\PYG{n}{swap\PYGZus{}fraction}\PYG{+w}{ }\PYG{n}{temperature}\PYG{+w}{ }\PYG{n}{deltamu}\PYG{+w}{ }\PYG{n}{...}
\end{sphinxVerbatim}
\begin{itemize}
\item {} 
\sphinxAtStartPar
ID, group\sphinxhyphen{}ID are documented in {\hyperref[\detokenize{fix::doc}]{\sphinxcrossref{\DUrole{doc}{fix}}}} command

\item {} 
\sphinxAtStartPar
sgcmc = style name of this fix command

\item {} 
\sphinxAtStartPar
every\_nsteps = number of MD steps between MC cycles

\item {} 
\sphinxAtStartPar
swap\_fraction = fraction of a full MC cycle carried out at each call (a value of 1.0 will perform as many trial moves as there are atoms)

\item {} 
\sphinxAtStartPar
temperature = temperature that enters Boltzmann factor in Metropolis criterion (usually the same as MD temperature)

\item {} 
\sphinxAtStartPar
deltamu = \sphinxtitleref{N\sphinxhyphen{}1} chemical potential differences \(\mu_1-\mu_2, \ldots, \mu_1-\mu_N\) (\sphinxtitleref{N} is the number of atom types)

\item {} 
\sphinxAtStartPar
Zero or more keyword/value pairs may be appended to fix definition line:
\begin{sphinxalltt}
keyword = \sphinxstyleemphasis{variance} or \sphinxstyleemphasis{randseed} or \sphinxstyleemphasis{window\_moves} or \sphinxstyleemphasis{window\_size}
  \sphinxstyleemphasis{variance} kappa conc1 {[}conc2{]} ... {[}concN{]}
    kappa = variance constraint parameter
    \sphinxtitleref{c\_2}, \sphinxtitleref{c\_3},..., \sphinxtitleref{c\_N} = \sphinxtitleref{N\sphinxhyphen{}1} target concentration fractions
  \sphinxstyleemphasis{randseed} N
    N = seed for pseudo random number generator
  \sphinxstyleemphasis{window\_moves} N
    N = number of times sampling window is moved during one MC cycle
  \sphinxstyleemphasis{window\_size} frac
    frac = size of sampling window (must be between 0.5 and 1.0)
\end{sphinxalltt}

\end{itemize}


\subsection{Examples}
\label{\detokenize{fix_sgcmc:examples}}
\begin{sphinxVerbatim}[commandchars=\\\{\}]
\PYG{k}{fix }\PYG{n+nv+nvIdentifier}{mc}\PYG{+w}{ }\PYG{n+nv+nvIdentifier}{all}\PYG{+w}{ }\PYG{n}{sgcmc}\PYG{+w}{ }\PYG{l+m}{50}\PYG{+w}{ }\PYG{l+m}{0.1}\PYG{+w}{ }\PYG{l+m}{400.0}\PYG{+w}{ }\PYG{o}{\PYGZhy{}}\PYG{l+m}{0.55}
\PYG{k}{fix }\PYG{n+nv+nvIdentifier}{vc}\PYG{+w}{ }\PYG{n+nv+nvIdentifier}{all}\PYG{+w}{ }\PYG{n}{sgcmc}\PYG{+w}{ }\PYG{l+m}{20}\PYG{+w}{ }\PYG{l+m}{0.2}\PYG{+w}{ }\PYG{l+m}{700.0}\PYG{+w}{ }\PYG{o}{\PYGZhy{}}\PYG{l+m}{0.7}\PYG{+w}{ }\PYG{n}{randseed}\PYG{+w}{ }\PYG{l+m}{324234}\PYG{+w}{ }\PYG{n}{variance}\PYG{+w}{ }\PYG{l+m}{2000.0}\PYG{+w}{ }\PYG{l+m}{0.05}
\PYG{k}{fix }\PYG{n+nv+nvIdentifier}{2}\PYG{+w}{  }\PYG{n+nv+nvIdentifier}{all}\PYG{+w}{ }\PYG{n}{sgcmc}\PYG{+w}{ }\PYG{l+m}{20}\PYG{+w}{ }\PYG{l+m}{0.1}\PYG{+w}{ }\PYG{l+m}{700.0}\PYG{+w}{ }\PYG{o}{\PYGZhy{}}\PYG{l+m}{0.7}\PYG{+w}{ }\PYG{n}{window\PYGZus{}moves}\PYG{+w}{ }\PYG{l+m}{20}
\end{sphinxVerbatim}


\subsection{Description}
\label{\detokenize{fix_sgcmc:description}}
\sphinxAtStartPar
\DUrole{versionmodified,added}{Added in version 22Dec2022.}

\sphinxAtStartPar
This command allows to carry out parallel hybrid molecular
dynamics/Monte Carlo (MD/MC) simulations using the algorithms described
in {\hyperref[\detokenize{fix_sgcmc:sadigh1}]{\sphinxcrossref{\DUrole{std,std-ref}{(Sadigh1)}}}}.  Simulations can be carried out in either
the semi\sphinxhyphen{}grand canonical (SGC) or variance constrained semi\sphinxhyphen{}grand
canonical (VC\sphinxhyphen{}SGC) ensemble {\hyperref[\detokenize{fix_sgcmc:sadigh2}]{\sphinxcrossref{\DUrole{std,std-ref}{(Sadigh2)}}}}. Only atom type
swaps are performed by the SGCMC fix. Relaxations are accounted for by
the molecular dynamics integration steps.

\sphinxAtStartPar
This fix can be used with standard multi\sphinxhyphen{}element EAM potentials
({\hyperref[\detokenize{pair_eam::doc}]{\sphinxcrossref{\DUrole{doc}{pair styles eam/alloy or eam/fs}}}})

\sphinxAtStartPar
The SGCMC fix can handle Finnis/Sinclair type EAM potentials where
\(\rho(r)\) is atom\sphinxhyphen{}type specific, such that different elements can
contribute differently to the total electron density at an atomic site
depending on the identity of the element at that atomic site.


\bigskip\hrule\bigskip


\sphinxAtStartPar
If this fix is applied, the regular MD simulation will be interrupted in
defined intervals to carry out a fraction of a Monte Carlo (MC)
cycle. The interval is set using the parameter \sphinxstyleemphasis{every\_nsteps} which
determines how many MD integrator steps are taken between subsequent
calls to the MC routine.

\sphinxAtStartPar
It is possible to carry out pure lattice MC simulations by setting
\sphinxstyleemphasis{every\_nsteps} to 1 and not defining an integration fix such as NVE,
NPT etc.  In that case, the particles will not move and only the MC
routine will be called to perform atom type swaps.

\sphinxAtStartPar
The parameter \sphinxstyleemphasis{swap\_fraction} determines how many MC trial steps are carried
out every time the MC routine is entered. It is measured in units of full MC
cycles where one full cycle, \sphinxstyleemphasis{swap\_fraction=1}, corresponds to as many MC
trial steps as there are atoms.


\bigskip\hrule\bigskip


\sphinxAtStartPar
The parameter \sphinxstyleemphasis{temperature} specifies the temperature that is used
to evaluate the Metropolis acceptance criterion. While it usually
should be set to the same value as the MD temperature there are cases
when it can be useful to use two different values for at least part of
the simulation, e.g., to speed up equilibration at low temperatures.


\bigskip\hrule\bigskip


\sphinxAtStartPar
The parameter \sphinxstyleemphasis{deltamu} is used to set the chemical potential differences
in the SGC MC algorithm (see Eq. 16 in {\hyperref[\detokenize{fix_sgcmc:sadigh1}]{\sphinxcrossref{\DUrole{std,std-ref}{Sadigh1}}}}).
The \sphinxtitleref{N\sphinxhyphen{}1} differences are defined as \(\mu_1-\mu_2, \ldots, \mu_1-\mu_N\),
where \sphinxtitleref{N} is the number of atom types.


\bigskip\hrule\bigskip


\sphinxAtStartPar
The variance\sphinxhyphen{}constrained SGC MC algorithm is activated if the keyword
\sphinxstyleemphasis{variance} is used. In that case the fix parameter \sphinxstyleemphasis{deltamu} determines
the effective average constraint in the parallel VC\sphinxhyphen{}SGC MC algorithm
(parameter \(\delta\mu_0\) in Eq. (20) of {\hyperref[\detokenize{fix_sgcmc:sadigh1}]{\sphinxcrossref{\DUrole{std,std-ref}{Sadigh1}}}}). The parameter \sphinxstyleemphasis{kappa} specifies the variance constraint
(see Eqs. (20\sphinxhyphen{}21) in {\hyperref[\detokenize{fix_sgcmc:sadigh1}]{\sphinxcrossref{\DUrole{std,std-ref}{Sadigh1}}}}).
The parameter \sphinxstyleemphasis{conc} sets the \sphinxtitleref{N\sphinxhyphen{}1} target atomic concentration
fractions (parameter \(c_0\) in Eqs.  (20\sphinxhyphen{}21) of {\hyperref[\detokenize{fix_sgcmc:sadigh1}]{\sphinxcrossref{\DUrole{std,std-ref}{Sadigh1}}}})
\(0 \le c_2, \ldots, c_N \le 1\), with
\(c_1 = 1 - \Sigma_{i=2}^N c_i\).
When the simulation includes \sphinxtitleref{N} atom types (elements),
\sphinxtitleref{N\sphinxhyphen{}1} concentration values must be specified.


\bigskip\hrule\bigskip


\sphinxAtStartPar
There are several technical parameters that can be set via optional flags.

\sphinxAtStartPar
\sphinxstyleemphasis{randseed} is expected to be a positive integer number and is used
to initialize the random number generator on each processor.

\sphinxAtStartPar
\sphinxstyleemphasis{window\_size} controls the size of the sampling window in a parallel MC
simulation. The size has to lie between 0.5 and 1.0. Normally, this
parameter should be left unspecified which instructs the code to choose
the optimal window size automatically (see Sect. III.B and Figure 6 in
{\hyperref[\detokenize{fix_sgcmc:sadigh1}]{\sphinxcrossref{\DUrole{std,std-ref}{Sadigh1}}}} for details).

\sphinxAtStartPar
The number of times the window is moved during a MC cycle is set using
the parameter \sphinxstyleemphasis{window\_moves} (see Sect. III.B in {\hyperref[\detokenize{fix_sgcmc:sadigh1}]{\sphinxcrossref{\DUrole{std,std-ref}{Sadigh1}}}} for details).


\bigskip\hrule\bigskip



\subsection{Restart, fix\_modify, output, run start/stop, minimize info}
\label{\detokenize{fix_sgcmc:restart-fix-modify-output-run-start-stop-minimize-info}}
\sphinxAtStartPar
No information about this fix is written to restart files.

\sphinxAtStartPar
The MC routine keeps track of the global concentration(s) as well as the
number of accepted and rejected trial swaps during each MC step. These
values are provided by the sgcmc fix in the form of a global vector that
can be accessed by various {\hyperref[\detokenize{Howto_output::doc}]{\sphinxcrossref{\DUrole{doc}{output commands}}}}
components of the vector represent the following quantities:
\begin{itemize}
\item {} 
\sphinxAtStartPar
1 = The absolute number of accepted trial swaps during the last MC step

\item {} 
\sphinxAtStartPar
2 = The absolute number of rejected trial swaps during the last MC step

\item {} 
\sphinxAtStartPar
3 = Current global concentration \sphinxtitleref{c\_1} (= number of atoms of type 1 / total number of atoms)

\item {} 
\sphinxAtStartPar
4 = Current global concentration \sphinxtitleref{c\_2} (= number of atoms of type 2 / total number of atoms)

\item {} 
\sphinxAtStartPar
…

\item {} 
\sphinxAtStartPar
N+2 = Current global concentration \sphinxtitleref{c\_N} (= number of atoms of type \sphinxstyleemphasis{N} / total number of atoms)

\end{itemize}

\sphinxAtStartPar
The vector values calculated by this fix are “intensive”.


\subsection{Restrictions}
\label{\detokenize{fix_sgcmc:restrictions}}
\sphinxAtStartPar
This fix is part of the MC package. It is only enabled if LAMMPS was
built with that package.  See the {\hyperref[\detokenize{Build_package::doc}]{\sphinxcrossref{\DUrole{doc}{Build package}}}}
page for more info.

\sphinxAtStartPar
This fix style requires an {\hyperref[\detokenize{atom_style::doc}]{\sphinxcrossref{\DUrole{doc}{atom style}}}} with per atom
type masses.

\sphinxAtStartPar
At present the fix provides optimized subroutines for EAM type
potentials (see above) that calculate potential energy changes due to
\sphinxstyleemphasis{local} atom type swaps very efficiently.  Other potentials are
supported by using the generic potential functions. This, however, will
lead to exceedingly slow simulations since it implies that the
energy of the \sphinxstyleemphasis{entire} system is recomputed at each MC trial step.  If
other potentials are to be used it is strongly recommended to modify and
optimize the existing generic potential functions for this purpose.
Also, the generic energy calculation can not be used for parallel
execution i.e. it only works with a single MPI process.


\bigskip\hrule\bigskip



\subsection{Default}
\label{\detokenize{fix_sgcmc:default}}
\sphinxAtStartPar
The optional parameters default to the following values:
\begin{itemize}
\item {} 
\sphinxAtStartPar
\sphinxstyleemphasis{randseed} = 324234

\item {} 
\sphinxAtStartPar
\sphinxstyleemphasis{window\_moves} = 8

\item {} 
\sphinxAtStartPar
\sphinxstyleemphasis{window\_size} = automatic

\end{itemize}


\bigskip\hrule\bigskip

\phantomsection\label{\detokenize{fix_sgcmc:sadigh1}}
\sphinxAtStartPar
\sphinxstylestrong{(Sadigh1)} B. Sadigh, P. Erhart, A. Stukowski, A. Caro, E. Martinez, and L. Zepeda\sphinxhyphen{}Ruiz, Phys. Rev. B \sphinxstylestrong{85}, 184203 (2012)

\phantomsection\label{\detokenize{fix_sgcmc:sadigh2}}
\sphinxAtStartPar
\sphinxstylestrong{(Sadigh2)} B. Sadigh and P. Erhart, Phys. Rev. B \sphinxstylestrong{86}, 134204 (2012)

\sphinxstepscope

\index{fix shake@\spxentry{fix shake}}\ignorespaces 
\index{fix shake/kk@\spxentry{fix shake/kk}}\ignorespaces 
\index{fix rattle@\spxentry{fix rattle}}\ignorespaces 

\section{fix shake command}
\label{\detokenize{fix_shake:fix-shake-command}}\label{\detokenize{fix_shake:index-2}}\label{\detokenize{fix_shake:index-1}}\label{\detokenize{fix_shake:index-0}}\label{\detokenize{fix_shake::doc}}
\sphinxAtStartPar
Accelerator Variants: \sphinxstyleemphasis{shake/kk}


\section{fix rattle command}
\label{\detokenize{fix_shake:fix-rattle-command}}

\subsection{Syntax}
\label{\detokenize{fix_shake:syntax}}
\begin{sphinxVerbatim}[commandchars=\\\{\}]
\PYG{k}{fix }\PYG{n+nv+nvIdentifier}{ID}\PYG{+w}{ }\PYG{n+nv+nvIdentifier}{group\PYGZhy{}ID}\PYG{+w}{ }\PYG{n}{style}\PYG{+w}{ }\PYG{n}{tol}\PYG{+w}{ }\PYG{n}{iter}\PYG{+w}{ }\PYG{n}{N}\PYG{+w}{ }\PYG{n}{constraint}\PYG{+w}{ }\PYG{n}{values}\PYG{+w}{ }\PYG{n}{...}\PYG{+w}{ }\PYG{n}{keyword}\PYG{+w}{ }\PYG{n}{value}\PYG{+w}{ }\PYG{n}{...}
\end{sphinxVerbatim}
\begin{itemize}
\item {} 
\sphinxAtStartPar
ID, group\sphinxhyphen{}ID are documented in {\hyperref[\detokenize{fix::doc}]{\sphinxcrossref{\DUrole{doc}{fix}}}} command

\item {} 
\sphinxAtStartPar
style = shake or rattle = style name of this fix command

\item {} 
\sphinxAtStartPar
tol = accuracy tolerance of SHAKE solution

\item {} 
\sphinxAtStartPar
iter = max \# of iterations in each SHAKE solution

\item {} 
\sphinxAtStartPar
N = print SHAKE statistics every this many timesteps (0 = never)

\item {} 
\sphinxAtStartPar
one or more constraint/value pairs are appended

\item {} 
\sphinxAtStartPar
constraint = \sphinxstyleemphasis{b} or \sphinxstyleemphasis{a} or \sphinxstyleemphasis{t} or \sphinxstyleemphasis{m}
\begin{sphinxalltt}
\sphinxstyleemphasis{b} values = one or more bond types (may use type labels)
\sphinxstyleemphasis{a} values = one or more angle types  (may use type labels)
\sphinxstyleemphasis{t} values = one or more atom types (may use type labels)
\sphinxstyleemphasis{m} value = one or more mass values
\end{sphinxalltt}

\item {} 
\sphinxAtStartPar
zero or more keyword/value pairs may be appended

\item {} 
\sphinxAtStartPar
keyword = \sphinxstyleemphasis{mol} or \sphinxstyleemphasis{kbond}
\begin{sphinxalltt}
\sphinxstyleemphasis{mol} value = template\sphinxhyphen{}ID
  template\sphinxhyphen{}ID = ID of molecule template specified in a separate {\hyperref[\detokenize{molecule::doc}]{\sphinxcrossref{\DUrole{doc}{molecule}}}} command
\sphinxstyleemphasis{kbond} value = force constant
  force constant = force constant used to apply a restraint force when used during minimization
\end{sphinxalltt}

\end{itemize}


\subsection{Examples}
\label{\detokenize{fix_shake:examples}}
\begin{sphinxVerbatim}[commandchars=\\\{\}]
\PYG{k}{fix }\PYG{n+nv+nvIdentifier}{1}\PYG{+w}{ }\PYG{n+nv+nvIdentifier}{sub}\PYG{+w}{ }\PYG{n}{shake}\PYG{+w}{ }\PYG{l+m}{0.0001}\PYG{+w}{ }\PYG{l+m}{20}\PYG{+w}{ }\PYG{l+m}{10}\PYG{+w}{ }\PYG{n}{b}\PYG{+w}{ }\PYG{l+m}{4}\PYG{+w}{ }\PYG{l+m}{19}\PYG{+w}{ }\PYG{n}{a}\PYG{+w}{ }\PYG{l+m}{3}\PYG{+w}{ }\PYG{l+m}{5}\PYG{+w}{ }\PYG{l+m}{2}
\PYG{k}{fix }\PYG{n+nv+nvIdentifier}{1}\PYG{+w}{ }\PYG{n+nv+nvIdentifier}{sub}\PYG{+w}{ }\PYG{n}{shake}\PYG{+w}{ }\PYG{l+m}{0.0001}\PYG{+w}{ }\PYG{l+m}{20}\PYG{+w}{ }\PYG{l+m}{10}\PYG{+w}{ }\PYG{n}{t}\PYG{+w}{ }\PYG{l+m}{5}\PYG{+w}{ }\PYG{l+m}{6}\PYG{+w}{ }\PYG{n}{m}\PYG{+w}{ }\PYG{l+m}{1.0}\PYG{+w}{ }\PYG{n}{a}\PYG{+w}{ }\PYG{l+m}{31}
\PYG{k}{fix }\PYG{n+nv+nvIdentifier}{1}\PYG{+w}{ }\PYG{n+nv+nvIdentifier}{sub}\PYG{+w}{ }\PYG{n}{shake}\PYG{+w}{ }\PYG{l+m}{0.0001}\PYG{+w}{ }\PYG{l+m}{20}\PYG{+w}{ }\PYG{l+m}{10}\PYG{+w}{ }\PYG{n}{t}\PYG{+w}{ }\PYG{l+m}{5}\PYG{+w}{ }\PYG{l+m}{6}\PYG{+w}{ }\PYG{n}{m}\PYG{+w}{ }\PYG{l+m}{1.0}\PYG{+w}{ }\PYG{n}{a}\PYG{+w}{ }\PYG{l+m}{31}\PYG{+w}{ }\PYG{n}{mol}\PYG{+w}{ }\PYG{n}{myMol}
\PYG{k}{fix }\PYG{n+nv+nvIdentifier}{1}\PYG{+w}{ }\PYG{n+nv+nvIdentifier}{sub}\PYG{+w}{ }\PYG{n}{rattle}\PYG{+w}{ }\PYG{l+m}{0.0001}\PYG{+w}{ }\PYG{l+m}{20}\PYG{+w}{ }\PYG{l+m}{10}\PYG{+w}{ }\PYG{n}{t}\PYG{+w}{ }\PYG{l+m}{5}\PYG{+w}{ }\PYG{l+m}{6}\PYG{+w}{ }\PYG{n}{m}\PYG{+w}{ }\PYG{l+m}{1.0}\PYG{+w}{ }\PYG{n}{a}\PYG{+w}{ }\PYG{l+m}{31}
\PYG{k}{fix }\PYG{n+nv+nvIdentifier}{1}\PYG{+w}{ }\PYG{n+nv+nvIdentifier}{sub}\PYG{+w}{ }\PYG{n}{rattle}\PYG{+w}{ }\PYG{l+m}{0.0001}\PYG{+w}{ }\PYG{l+m}{20}\PYG{+w}{ }\PYG{l+m}{10}\PYG{+w}{ }\PYG{n}{t}\PYG{+w}{ }\PYG{l+m}{5}\PYG{+w}{ }\PYG{l+m}{6}\PYG{+w}{ }\PYG{n}{m}\PYG{+w}{ }\PYG{l+m}{1.0}\PYG{+w}{ }\PYG{n}{a}\PYG{+w}{ }\PYG{l+m}{31}\PYG{+w}{ }\PYG{n}{mol}\PYG{+w}{ }\PYG{n}{myMol}
\end{sphinxVerbatim}


\subsection{Description}
\label{\detokenize{fix_shake:description}}
\sphinxAtStartPar
Apply bond and angle constraints to specified bonds and angles in the
simulation by either the SHAKE or RATTLE algorithms.  This typically
enables a longer timestep.  The SHAKE or RATTLE constraint algorithms,
however, can \sphinxstyleemphasis{only} be applied during molecular dynamics runs.

\sphinxAtStartPar
\DUrole{versionmodified,changed}{Changed in version 15Sep2022.}

\sphinxAtStartPar
These fixes may now also be used during minimization.  In that case the
constraints are \sphinxstyleemphasis{approximated} by strong harmonic restraints.

\sphinxAtStartPar
\sphinxstylestrong{SHAKE vs RATTLE:}

\sphinxAtStartPar
The SHAKE algorithm was invented for schemes such as standard Verlet
timestepping, where only the coordinates are integrated and the
velocities are approximated as finite differences to the trajectories
({\hyperref[\detokenize{fix_shake:ryckaert}]{\sphinxcrossref{\DUrole{std,std-ref}{Ryckaert et al. (1977)}}}}).  If the velocities are
integrated explicitly, as with velocity Verlet which is what LAMMPS
uses as an integration method, a second set of constraining forces is
required in order to eliminate velocity components along the bonds
({\hyperref[\detokenize{fix_shake:andersen3}]{\sphinxcrossref{\DUrole{std,std-ref}{Andersen (1983)}}}}).

\sphinxAtStartPar
In order to formulate individual constraints for SHAKE and RATTLE,
focus on a single molecule whose bonds are constrained.  Let Ri and Vi
be the position and velocity of atom \sphinxstyleemphasis{i} at time \sphinxstyleemphasis{n}, for
\sphinxstyleemphasis{i} =1,…,\sphinxstyleemphasis{N}, where \sphinxstyleemphasis{N} is the number of sites of our reference
molecule. The distance vector between sites \sphinxstyleemphasis{i} and \sphinxstyleemphasis{j} is given by
\begin{equation*}
\begin{split}\mathbf r^{n+1}_{ij} = \mathbf r^n_j - \mathbf r^n_i\end{split}
\end{equation*}
\sphinxAtStartPar
The constraints can then be formulated as
\begin{equation*}
\begin{split}\mathbf r^{n+1}_{ij} \cdot \mathbf r^{n+1}_{ij} &= d^2_{ij} \quad \text{and} \\
\mathbf v^{n+1}_{ij} \cdot \mathbf r^{n+1}_{ij} &= 0\end{split}
\end{equation*}
\sphinxAtStartPar
The SHAKE algorithm satisfies the first condition, i.e. the sites at
time \sphinxstyleemphasis{n+1} will have the desired separations Dij immediately after the
coordinates are integrated.  If we also enforce the second condition,
the velocity components along the bonds will vanish.  RATTLE satisfies
both conditions.  As implemented in LAMMPS, \sphinxstyleemphasis{fix rattle} uses fix shake
for satisfying the coordinate constraints. Therefore the settings and
optional keywords are the same for both fixes, and all the information
below about SHAKE is also relevant for RATTLE.

\sphinxAtStartPar
\sphinxstylestrong{SHAKE:}

\sphinxAtStartPar
Each timestep the specified bonds and angles are reset to their
equilibrium lengths and angular values via the SHAKE algorithm
({\hyperref[\detokenize{fix_shake:ryckaert}]{\sphinxcrossref{\DUrole{std,std-ref}{Ryckaert et al. (1977)}}}}).  This is done by applying an
additional constraint force so that the new positions preserve the
desired atom separations.  The equations for the additional force are
solved via an iterative method that typically converges to an accurate
solution in a few iterations.  The desired tolerance (e.g. 1.0e\sphinxhyphen{}4 = 1
part in 10000) and maximum \# of iterations are specified as arguments.
Setting the N argument will print statistics to the screen and log
file about regarding the lengths of bonds and angles that are being
constrained.  Small delta values mean SHAKE is doing a good job.

\sphinxAtStartPar
In LAMMPS, only small clusters of atoms can be constrained.  This is
so the constraint calculation for a cluster can be performed by a
single processor, to enable good parallel performance.  A cluster is
defined as a central atom connected to others in the cluster by
constrained bonds.  LAMMPS allows for the following kinds of clusters
to be constrained: one central atom bonded to 1 or 2 or 3 atoms, or
one central atom bonded to 2 others and the angle between the three atoms
also constrained.  This means water molecules or CH2 or CH3 groups may
be constrained, but not all the C\sphinxhyphen{}C backbone bonds of a long polymer
chain.

\sphinxAtStartPar
The \sphinxstyleemphasis{b} constraint lists bond types that will be constrained.  The \sphinxstyleemphasis{t}
constraint lists atom types.  All bonds connected to an atom of the
specified type will be constrained.  The \sphinxstyleemphasis{m} constraint lists atom
masses.  All bonds connected to atoms of the specified masses will be
constrained (within a fudge factor of MASSDELTA specified in
\sphinxcode{\sphinxupquote{src/RIGID/fix\_shake.cpp}}).  The \sphinxstyleemphasis{a} constraint lists angle types.  If
both bonds in the angle are constrained then the angle will also be
constrained if its type is in the list.

\sphinxAtStartPar
\DUrole{versionmodified,changed}{Changed in version 29Aug2024.}

\sphinxAtStartPar
The types may be given as type labels \sphinxstyleemphasis{only} if there is no atom, bond,
or angle type label named \sphinxstyleemphasis{b}, \sphinxstyleemphasis{a}, \sphinxstyleemphasis{t}, or \sphinxstyleemphasis{m} defined in the
simulation.  If that is the case, type labels cannot be used as
constraint type index with these two fixes, because the type labels
would be incorrectly treated as a new type of constraint instead.
Thus, LAMMPS will print a warning and type label handling is disabled
and numeric types must be used.

\sphinxAtStartPar
For all constraints, a particular bond is only constrained if \sphinxstyleemphasis{both}
atoms in the bond are in the group specified with the SHAKE fix.

\sphinxAtStartPar
The degrees\sphinxhyphen{}of\sphinxhyphen{}freedom removed by SHAKE bonds and angles are accounted
for in temperature and pressure computations.  Similarly, the SHAKE
contribution to the pressure of the system (virial) is also accounted
for.

\begin{sphinxadmonition}{note}{Note:}
\sphinxAtStartPar
This command works by using the current forces on atoms to calculate
an additional constraint force which when added will leave the atoms
in positions that satisfy the SHAKE constraints (e.g. bond length)
after the next time integration step.  If you define fixes
(e.g. {\hyperref[\detokenize{fix_efield::doc}]{\sphinxcrossref{\DUrole{doc}{fix efield}}}}) that add additional force to
the atoms after \sphinxstyleemphasis{fix shake} operates, then this fix will not take them
into account and the time integration will typically not satisfy the
SHAKE constraints.  The solution for this is to make sure that fix
shake is defined in your input script after any other fixes which add
or change forces (to atoms that \sphinxstyleemphasis{fix shake} operates on).
\end{sphinxadmonition}


\bigskip\hrule\bigskip


\sphinxAtStartPar
The \sphinxstyleemphasis{mol} keyword should be used when other commands, such as {\hyperref[\detokenize{fix_deposit::doc}]{\sphinxcrossref{\DUrole{doc}{fix
deposit}}}} or {\hyperref[\detokenize{fix_pour::doc}]{\sphinxcrossref{\DUrole{doc}{fix pour}}}}, add molecules
on\sphinxhyphen{}the\sphinxhyphen{}fly during a simulation, and you wish to constrain the new
molecules via SHAKE.  You specify a \sphinxstyleemphasis{template\sphinxhyphen{}ID} previously defined
using the {\hyperref[\detokenize{molecule::doc}]{\sphinxcrossref{\DUrole{doc}{molecule}}}} command, which reads a file that
defines the molecule.  You must use the same \sphinxstyleemphasis{template\sphinxhyphen{}ID} that the
command adding molecules uses.  The coordinates, atom types, special
bond restrictions, and SHAKE info can be specified in the molecule file.
See the {\hyperref[\detokenize{molecule::doc}]{\sphinxcrossref{\DUrole{doc}{molecule}}}} command for details.  The only
settings required to be in this file (by this command) are the SHAKE
info of atoms in the molecule.

\sphinxAtStartPar
The \sphinxstyleemphasis{kbond} keyword sets the restraint force constant when \sphinxstyleemphasis{fix shake}
or fix rattle are used during minimization.  In that case the constraint
algorithms are \sphinxstyleemphasis{not} applied and restraint forces are used instead to
maintain the geometries similar to the constraints.  How well the
geometries are maintained and how quickly a minimization converges,
depends largely on the force constant \sphinxstyleemphasis{kbond}: larger values will reduce
the deviation from the desired geometry, but can also lead to slower
convergence of the minimization or lead to instabilities depending on
the minimization algorithm requiring to reduce the value of
{\hyperref[\detokenize{timestep::doc}]{\sphinxcrossref{\DUrole{doc}{timestep}}}}.  Even though the restraints will not
preserve the bond lengths and angles as closely as the constraints
during the MD, they are generally close enough so that the constraints
will be fulfilled to the desired accuracy within a few MD steps
following the minimization. The default value for \sphinxstyleemphasis{kbond} depends on the
{\hyperref[\detokenize{units::doc}]{\sphinxcrossref{\DUrole{doc}{units}}}} setting and is 1.0e6*k\_B.


\bigskip\hrule\bigskip


\sphinxAtStartPar
Styles with a \sphinxstyleemphasis{gpu}, \sphinxstyleemphasis{intel}, \sphinxstyleemphasis{kk}, \sphinxstyleemphasis{omp}, or \sphinxstyleemphasis{opt} suffix are
functionally the same as the corresponding style without the suffix.
They have been optimized to run faster, depending on your available
hardware, as discussed on the {\hyperref[\detokenize{Speed_packages::doc}]{\sphinxcrossref{\DUrole{doc}{Accelerator packages}}}}
page.  The accelerated styles take the same arguments and should
produce the same results, except for round\sphinxhyphen{}off and precision issues.

\sphinxAtStartPar
These accelerated styles are part of the GPU, INTEL, KOKKOS,
OPENMP, and OPT packages, respectively.  They are only enabled if
LAMMPS was built with those packages.  See the {\hyperref[\detokenize{Build_package::doc}]{\sphinxcrossref{\DUrole{doc}{Build package}}}} page for more info.

\sphinxAtStartPar
You can specify the accelerated styles explicitly in your input script
by including their suffix, or you can use the {\hyperref[\detokenize{Run_options::doc}]{\sphinxcrossref{\DUrole{doc}{\sphinxhyphen{}suffix command\sphinxhyphen{}line switch}}}} when you invoke LAMMPS, or you can use the
{\hyperref[\detokenize{suffix::doc}]{\sphinxcrossref{\DUrole{doc}{suffix}}}} command in your input script.

\sphinxAtStartPar
See the {\hyperref[\detokenize{Speed_packages::doc}]{\sphinxcrossref{\DUrole{doc}{Accelerator packages}}}} page for more
instructions on how to use the accelerated styles effectively.


\bigskip\hrule\bigskip


\sphinxAtStartPar
\sphinxstylestrong{RATTLE:}

\sphinxAtStartPar
The velocity constraints lead to a linear system of equations which
can be solved analytically.  The implementation of the algorithm in
LAMMPS closely follows ({\hyperref[\detokenize{fix_shake:andersen3}]{\sphinxcrossref{\DUrole{std,std-ref}{Andersen (1983)}}}}).

\begin{sphinxadmonition}{note}{Note:}
\sphinxAtStartPar
The \sphinxstyleemphasis{fix rattle} command modifies forces and velocities and thus
should be defined after all other integration fixes in your input
script.  If you define other fixes that modify velocities or forces
after \sphinxstyleemphasis{fix rattle} operates, then \sphinxstyleemphasis{fix rattle} will not take them
into account and the overall time integration will typically not
satisfy the RATTLE constraints.  You can check whether the
constraints work correctly by setting the value of RATTLE\_DEBUG in
\sphinxcode{\sphinxupquote{src/RIGID/fix\_rattle.cpp}} to 1 and recompiling LAMMPS.
\end{sphinxadmonition}


\bigskip\hrule\bigskip



\subsection{Restart, fix\_modify, output, run start/stop, minimize info}
\label{\detokenize{fix_shake:restart-fix-modify-output-run-start-stop-minimize-info}}
\sphinxAtStartPar
No information about these fixes is written to {\hyperref[\detokenize{restart::doc}]{\sphinxcrossref{\DUrole{doc}{binary restart
files}}}}.

\sphinxAtStartPar
Both fix \sphinxstyleemphasis{shake} and fix \sphinxstyleemphasis{rattle} behave differently during a minimization
in comparison to a molecular dynamics run:
\begin{itemize}
\item {} 
\sphinxAtStartPar
When used during a minimization, the SHAKE or RATTLE constraint
algorithms themselves are \sphinxstylestrong{not} applied.  Instead the constraints
are replaced by harmonic restraint forces.  The energy and virial
contributions due to the restraint forces are tallied into global and
per\sphinxhyphen{}atom accumulators. The total restraint energy is also accessible
as a global scalar property of the fix.

\item {} 
\sphinxAtStartPar
During molecular dynamics runs, however, the fixes do apply the
requested SHAKE or RATTLE constraint algorithms.

\sphinxAtStartPar
The {\hyperref[\detokenize{fix_modify::doc}]{\sphinxcrossref{\DUrole{doc}{fix\_modify}}}} \sphinxstyleemphasis{virial} option is supported by
these fixes to add the contribution due to the added constraint forces
on atoms to both the global pressure and per\sphinxhyphen{}atom stress of the system
via the {\hyperref[\detokenize{compute_pressure::doc}]{\sphinxcrossref{\DUrole{doc}{compute pressure}}}} and {\hyperref[\detokenize{compute_stress_atom::doc}]{\sphinxcrossref{\DUrole{doc}{compute
stress/atom}}}} commands.  The former can be
accessed by {\hyperref[\detokenize{thermo_style::doc}]{\sphinxcrossref{\DUrole{doc}{thermodynamic output}}}}.

\sphinxAtStartPar
The  default  setting  for  this  fix  is  {\hyperref[\detokenize{fix_modify::doc}]{\sphinxcrossref{\DUrole{doc}{fix\_modify  virial  yes}}}}.  No global  or per\sphinxhyphen{}atom  quantities are  stored by  these
fixes for access by various {\hyperref[\detokenize{Howto_output::doc}]{\sphinxcrossref{\DUrole{doc}{output commands}}}} during
an  MD  run.   No  parameter  of  these  fixes  can  be  used  with  the
\sphinxstyleemphasis{start/stop} keywords of the {\hyperref[\detokenize{run::doc}]{\sphinxcrossref{\DUrole{doc}{run}}}} command.

\end{itemize}


\subsection{Restrictions}
\label{\detokenize{fix_shake:restrictions}}
\sphinxAtStartPar
These fixes are part of the RIGID package.  They are only enabled if
LAMMPS was built with that package.  See the {\hyperref[\detokenize{Build_package::doc}]{\sphinxcrossref{\DUrole{doc}{Build package}}}} page for more info.

\sphinxAtStartPar
For computational efficiency, there can only be one shake or rattle
fix defined in a simulation.

\sphinxAtStartPar
If you use a tolerance that is too large or a max\sphinxhyphen{}iteration count that
is too small, the constraints will not be enforced very strongly,
which can lead to poor energy conservation.  You can test for this in
your system by running a constant NVE simulation with a particular set
of SHAKE parameters and monitoring the energy versus time.

\sphinxAtStartPar
SHAKE or RATTLE should not be used to constrain an angle at 180 degrees
(e.g. a linear CO2 molecule).  This causes a divergence when solving the
constraint equations numerically.  You can use {\hyperref[\detokenize{fix_rigid::doc}]{\sphinxcrossref{\DUrole{doc}{fix rigid or fix
rigid/small}}}} instead to make a linear molecule rigid.

\sphinxAtStartPar
When used during minimization choosing a too large value of the \sphinxstyleemphasis{kbond}
can make minimization very inefficient and also cause stability problems
with some minimization algorithms.  Sometimes those can be avoided by
reducing the {\hyperref[\detokenize{timestep::doc}]{\sphinxcrossref{\DUrole{doc}{timestep}}}}.


\subsection{Related commands}
\label{\detokenize{fix_shake:related-commands}}
\sphinxAtStartPar
{\hyperref[\detokenize{fix_rigid::doc}]{\sphinxcrossref{\DUrole{doc}{fix rigid}}}}, {\hyperref[\detokenize{fix_ehex::doc}]{\sphinxcrossref{\DUrole{doc}{fix ehex}}}},
{\hyperref[\detokenize{fix_nve_manifold_rattle::doc}]{\sphinxcrossref{\DUrole{doc}{fix nve/manifold/rattle}}}}


\subsection{Default}
\label{\detokenize{fix_shake:default}}
\sphinxAtStartPar
kbond = 1.0e9*k\_B


\bigskip\hrule\bigskip

\phantomsection\label{\detokenize{fix_shake:ryckaert}}
\sphinxAtStartPar
\sphinxstylestrong{(Ryckaert)} J.\sphinxhyphen{}P. Ryckaert, G. Ciccotti and H. J. C. Berendsen,
J of Comp Phys, 23, 327\sphinxhyphen{}341 (1977).

\phantomsection\label{\detokenize{fix_shake:andersen3}}
\sphinxAtStartPar
\sphinxstylestrong{(Andersen)} H. Andersen, J of Comp Phys, 52, 24\sphinxhyphen{}34 (1983).

\sphinxstepscope

\index{fix shardlow@\spxentry{fix shardlow}}\ignorespaces 
\index{fix shardlow/kk@\spxentry{fix shardlow/kk}}\ignorespaces 

\section{fix shardlow command}
\label{\detokenize{fix_shardlow:fix-shardlow-command}}\label{\detokenize{fix_shardlow:index-1}}\label{\detokenize{fix_shardlow:index-0}}\label{\detokenize{fix_shardlow::doc}}
\sphinxAtStartPar
Accelerator Variants: \sphinxstyleemphasis{shardlow/kk}


\subsection{Syntax}
\label{\detokenize{fix_shardlow:syntax}}
\begin{sphinxVerbatim}[commandchars=\\\{\}]
\PYG{k}{fix }\PYG{n+nv+nvIdentifier}{ID}\PYG{+w}{ }\PYG{n+nv+nvIdentifier}{group\PYGZhy{}ID}\PYG{+w}{ }\PYG{n}{shardlow}
\end{sphinxVerbatim}
\begin{itemize}
\item {} 
\sphinxAtStartPar
ID, group\sphinxhyphen{}ID are documented in {\hyperref[\detokenize{fix::doc}]{\sphinxcrossref{\DUrole{doc}{fix}}}} command

\item {} 
\sphinxAtStartPar
shardlow = style name of this fix command

\end{itemize}


\subsection{Examples}
\label{\detokenize{fix_shardlow:examples}}
\begin{sphinxVerbatim}[commandchars=\\\{\}]
\PYG{k}{fix }\PYG{n+nv+nvIdentifier}{1}\PYG{+w}{ }\PYG{n+nv+nvIdentifier}{all}\PYG{+w}{ }\PYG{n}{shardlow}
\end{sphinxVerbatim}


\subsection{Description}
\label{\detokenize{fix_shardlow:description}}
\sphinxAtStartPar
Specifies that the Shardlow splitting algorithm (SSA) is to be used to
integrate the DPD equations of motion.  The SSA splits the integration
into a stochastic and deterministic integration step.  The fix
\sphinxstyleemphasis{shardlow} performs the stochastic integration step and must be used in
conjunction with a deterministic integrator (e.g. {\hyperref[\detokenize{fix_nve::doc}]{\sphinxcrossref{\DUrole{doc}{fix nve}}}} or {\hyperref[\detokenize{fix_nh::doc}]{\sphinxcrossref{\DUrole{doc}{fix nph}}}}).  The stochastic integration of
the dissipative and random forces is performed prior to the
deterministic integration of the conservative force. Further details
regarding the method are provided in {\hyperref[\detokenize{fix_shardlow:lisal}]{\sphinxcrossref{\DUrole{std,std-ref}{(Lisal)}}}} and
{\hyperref[\detokenize{fix_shardlow:larentzos1sh}]{\sphinxcrossref{\DUrole{std,std-ref}{(Larentzos1)}}}}.

\sphinxAtStartPar
The fix \sphinxstyleemphasis{shardlow} must be used with the {\hyperref[\detokenize{pair_style::doc}]{\sphinxcrossref{\DUrole{doc}{pair\_style dpd/fdt}}}} or {\hyperref[\detokenize{pair_style::doc}]{\sphinxcrossref{\DUrole{doc}{pair\_style dpd/fdt/energy}}}} command
to properly initialize the fluctuation\sphinxhyphen{}dissipation theorem parameter(s)
sigma (and kappa, if necessary).

\sphinxAtStartPar
Note that numerous variants of DPD can be specified by choosing an
appropriate combination of the integrator and {\hyperref[\detokenize{pair_style::doc}]{\sphinxcrossref{\DUrole{doc}{pair\_style dpd/fdt}}}} command.  DPD under isothermal conditions can be specified
by using fix \sphinxstyleemphasis{shardlow}, fix \sphinxstyleemphasis{nve} and pair\_style \sphinxstyleemphasis{dpd/fdt}.  DPD
under isoenergetic conditions can be specified by using fix \sphinxstyleemphasis{shardlow},
fix \sphinxstyleemphasis{nve} and pair\_style \sphinxstyleemphasis{dpd/fdt/energy}.  DPD under isobaric
conditions can be specified by using fix shardlow, fix \sphinxstyleemphasis{nph} and
pair\_style \sphinxstyleemphasis{dpd/fdt}.  DPD under isoenthalpic conditions can be
specified by using fix shardlow, fix \sphinxstyleemphasis{nph} and pair\_style
\sphinxstyleemphasis{dpd/fdt/energy}.  Examples of each DPD variant are provided in the
examples/PACKAGES/dpd\sphinxhyphen{}react directory.


\bigskip\hrule\bigskip


\sphinxAtStartPar
Styles with a \sphinxstyleemphasis{gpu}, \sphinxstyleemphasis{intel}, \sphinxstyleemphasis{kk}, \sphinxstyleemphasis{omp}, or \sphinxstyleemphasis{opt} suffix are
functionally the same as the corresponding style without the suffix.
They have been optimized to run faster, depending on your available
hardware, as discussed on the {\hyperref[\detokenize{Speed_packages::doc}]{\sphinxcrossref{\DUrole{doc}{Accelerator packages}}}}
page.  The accelerated styles take the same arguments and should
produce the same results, except for round\sphinxhyphen{}off and precision issues.

\sphinxAtStartPar
These accelerated styles are part of the GPU, INTEL, KOKKOS,
OPENMP, and OPT packages, respectively.  They are only enabled if
LAMMPS was built with those packages.  See the {\hyperref[\detokenize{Build_package::doc}]{\sphinxcrossref{\DUrole{doc}{Build package}}}} page for more info.

\sphinxAtStartPar
You can specify the accelerated styles explicitly in your input script
by including their suffix, or you can use the {\hyperref[\detokenize{Run_options::doc}]{\sphinxcrossref{\DUrole{doc}{\sphinxhyphen{}suffix command\sphinxhyphen{}line switch}}}} when you invoke LAMMPS, or you can use the
{\hyperref[\detokenize{suffix::doc}]{\sphinxcrossref{\DUrole{doc}{suffix}}}} command in your input script.

\sphinxAtStartPar
See the {\hyperref[\detokenize{Speed_packages::doc}]{\sphinxcrossref{\DUrole{doc}{Accelerator packages}}}} page for more
instructions on how to use the accelerated styles effectively.


\bigskip\hrule\bigskip



\subsection{Restrictions}
\label{\detokenize{fix_shardlow:restrictions}}
\sphinxAtStartPar
This command is part of the DPD\sphinxhyphen{}REACT package.  It is only enabled if
LAMMPS was built with that package.  See the {\hyperref[\detokenize{Build_package::doc}]{\sphinxcrossref{\DUrole{doc}{Build package}}}} page for more info.

\sphinxAtStartPar
This fix is currently limited to orthogonal simulation cell
geometries.

\sphinxAtStartPar
This fix must be used with an additional fix that specifies time
integration, e.g. {\hyperref[\detokenize{fix_nve::doc}]{\sphinxcrossref{\DUrole{doc}{fix nve}}}} or {\hyperref[\detokenize{fix_nh::doc}]{\sphinxcrossref{\DUrole{doc}{fix nph}}}}.

\sphinxAtStartPar
The Shardlow splitting algorithm requires the sizes of the subdomain
lengths to be larger than twice the cutoff+skin.  Generally, the
domain decomposition is dependent on the number of processors
requested.


\subsection{Related commands}
\label{\detokenize{fix_shardlow:related-commands}}
\sphinxAtStartPar
{\hyperref[\detokenize{pair_dpd_fdt::doc}]{\sphinxcrossref{\DUrole{doc}{pair\_style dpd/fdt}}}}, {\hyperref[\detokenize{fix_eos_cv::doc}]{\sphinxcrossref{\DUrole{doc}{fix eos/cv}}}}


\subsection{Default}
\label{\detokenize{fix_shardlow:default}}
\sphinxAtStartPar
none


\bigskip\hrule\bigskip

\phantomsection\label{\detokenize{fix_shardlow:lisal}}
\sphinxAtStartPar
\sphinxstylestrong{(Lisal)} M. Lisal, J.K. Brennan, J. Bonet Avalos, J. Chem. Phys., 135, 204105 (2011).

\phantomsection\label{\detokenize{fix_shardlow:larentzos1sh}}
\sphinxAtStartPar
\sphinxstylestrong{(Larentzos1)} J.P. Larentzos, J.K. Brennan, J.D. Moore, M. Lisal and W.D. Mattson, Comput. Phys. Commun., 185, 1987\sphinxhyphen{}1998 (2014).

\phantomsection\label{\detokenize{fix_shardlow:larentzos2sh}}
\sphinxAtStartPar
\sphinxstylestrong{(Larentzos2)} J.P. Larentzos, J.K. Brennan, J.D. Moore, and W.D. Mattson, ARL\sphinxhyphen{}TR\sphinxhyphen{}6863, U.S. Army Research Laboratory, Aberdeen Proving Ground, MD (2014).

\sphinxstepscope

\index{fix smd@\spxentry{fix smd}}\ignorespaces 

\section{fix smd command}
\label{\detokenize{fix_smd:fix-smd-command}}\label{\detokenize{fix_smd:index-0}}\label{\detokenize{fix_smd::doc}}

\subsection{Syntax}
\label{\detokenize{fix_smd:syntax}}
\begin{sphinxVerbatim}[commandchars=\\\{\}]
\PYG{k}{fix }\PYG{n+nv+nvIdentifier}{ID}\PYG{+w}{ }\PYG{n+nv+nvIdentifier}{group\PYGZhy{}ID}\PYG{+w}{ }\PYG{n}{smd}\PYG{+w}{ }\PYG{n}{type}\PYG{+w}{ }\PYG{n}{values}\PYG{+w}{ }\PYG{n}{keyword}\PYG{+w}{ }\PYG{n}{values}
\end{sphinxVerbatim}
\begin{itemize}
\item {} 
\sphinxAtStartPar
ID, group\sphinxhyphen{}ID are documented in {\hyperref[\detokenize{fix::doc}]{\sphinxcrossref{\DUrole{doc}{fix}}}} command

\item {} 
\sphinxAtStartPar
smd  = style name of this fix command

\item {} 
\sphinxAtStartPar
mode = \sphinxstyleemphasis{cvel} or \sphinxstyleemphasis{cfor} to select constant velocity or constant force SMD
\begin{sphinxalltt}
\sphinxstyleemphasis{cvel} values = K vel
  K = spring constant (force/distance units)
  vel = velocity of pulling (distance/time units)
\sphinxstyleemphasis{cfor} values = force
  force = pulling force (force units)
\end{sphinxalltt}

\item {} 
\sphinxAtStartPar
keyword = \sphinxstyleemphasis{tether} or \sphinxstyleemphasis{couple}
\begin{sphinxalltt}
\sphinxstyleemphasis{tether} values = x y z R0
  x,y,z = point to which spring is tethered
  R0 = distance of end of spring from tether point (distance units)
\sphinxstyleemphasis{couple} values = group\sphinxhyphen{}ID2 x y z R0
  group\sphinxhyphen{}ID2 = 2nd group to couple to fix group with a spring
  x,y,z = direction of spring, automatically computed with \textquotesingle{}auto\textquotesingle{}
  R0 = distance of end of spring (distance units)
\end{sphinxalltt}

\end{itemize}


\subsection{Examples}
\label{\detokenize{fix_smd:examples}}
\begin{sphinxVerbatim}[commandchars=\\\{\}]
\PYG{k}{fix  }\PYG{n+nv+nvIdentifier}{pull}\PYG{+w}{    }\PYG{n+nv+nvIdentifier}{cterm}\PYG{+w}{ }\PYG{n}{smd}\PYG{+w}{ }\PYG{n}{cvel}\PYG{+w}{ }\PYG{l+m}{20.0}\PYG{+w}{ }\PYG{o}{\PYGZhy{}}\PYG{l+m}{0.00005}\PYG{+w}{ }\PYG{n}{tether}\PYG{+w}{ }\PYG{n}{NULL}\PYG{+w}{ }\PYG{n}{NULL}\PYG{+w}{ }\PYG{l+m}{100.0}\PYG{+w}{ }\PYG{l+m}{0.0}
\PYG{k}{fix  }\PYG{n+nv+nvIdentifier}{pull}\PYG{+w}{    }\PYG{n+nv+nvIdentifier}{cterm}\PYG{+w}{ }\PYG{n}{smd}\PYG{+w}{ }\PYG{n}{cvel}\PYG{+w}{ }\PYG{l+m}{20.0}\PYG{+w}{ }\PYG{o}{\PYGZhy{}}\PYG{l+m}{0.0001}\PYG{+w}{ }\PYG{n}{tether}\PYG{+w}{ }\PYG{l+m}{25.0}\PYG{+w}{ }\PYG{l+m}{25}\PYG{+w}{ }\PYG{l+m}{25.0}\PYG{+w}{ }\PYG{l+m}{0.0}
\PYG{k}{fix  }\PYG{n+nv+nvIdentifier}{stretch}\PYG{+w}{ }\PYG{n+nv+nvIdentifier}{cterm}\PYG{+w}{ }\PYG{n}{smd}\PYG{+w}{ }\PYG{n}{cvel}\PYG{+w}{ }\PYG{l+m}{20.0}\PYG{+w}{  }\PYG{l+m}{0.0001}\PYG{+w}{ }\PYG{n}{couple}\PYG{+w}{ }\PYG{n}{nterm}\PYG{+w}{ }\PYG{n}{auto}\PYG{+w}{ }\PYG{n}{auto}\PYG{+w}{ }\PYG{n}{auto}\PYG{+w}{ }\PYG{l+m}{0.0}
\PYG{k}{fix  }\PYG{n+nv+nvIdentifier}{pull}\PYG{+w}{    }\PYG{n+nv+nvIdentifier}{cterm}\PYG{+w}{ }\PYG{n}{smd}\PYG{+w}{ }\PYG{n}{cfor}\PYG{+w}{  }\PYG{l+m}{5.0}\PYG{+w}{ }\PYG{n}{tether}\PYG{+w}{ }\PYG{l+m}{25.0}\PYG{+w}{ }\PYG{l+m}{25.0}\PYG{+w}{ }\PYG{l+m}{25.0}\PYG{+w}{ }\PYG{l+m}{0.0}
\end{sphinxVerbatim}


\subsection{Description}
\label{\detokenize{fix_smd:description}}
\sphinxAtStartPar
This fix implements several options of steered MD (SMD) as reviewed in
{\hyperref[\detokenize{fix_smd:izrailev}]{\sphinxcrossref{\DUrole{std,std-ref}{(Izrailev)}}}}, which allows to induce conformational
changes in systems and to compute the potential of mean force (PMF)
along the assumed reaction coordinate {\hyperref[\detokenize{fix_smd:park}]{\sphinxcrossref{\DUrole{std,std-ref}{(Park)}}}} based on
Jarzynski’s equality {\hyperref[\detokenize{fix_smd:jarzynski}]{\sphinxcrossref{\DUrole{std,std-ref}{(Jarzynski)}}}}.  This fix borrows
a lot from {\hyperref[\detokenize{fix_spring::doc}]{\sphinxcrossref{\DUrole{doc}{fix spring}}}} and {\hyperref[\detokenize{fix_setforce::doc}]{\sphinxcrossref{\DUrole{doc}{fix setforce}}}}.

\sphinxAtStartPar
You can apply a moving spring force to a group of atoms (\sphinxstyleemphasis{tether}
style) or between two groups of atoms (\sphinxstyleemphasis{couple} style).  The spring
can then be used in either constant velocity (\sphinxstyleemphasis{cvel}) mode or in
constant force (\sphinxstyleemphasis{cfor}) mode to induce transitions in your systems.
When running in \sphinxstyleemphasis{tether} style, you may need some way to fix some
other part of the system (e.g. via {\hyperref[\detokenize{fix_spring_self::doc}]{\sphinxcrossref{\DUrole{doc}{fix spring/self}}}})

\sphinxAtStartPar
The \sphinxstyleemphasis{tether} style attaches a spring between a point at a distance of
R0 away from a fixed point \sphinxstyleemphasis{x,y,z} and the center of mass of the fix
group of atoms.  A restoring force of magnitude K (R \sphinxhyphen{} R0) Mi / M is
applied to each atom in the group where \sphinxstyleemphasis{K} is the spring constant, Mi
is the mass of the atom, and M is the total mass of all atoms in the
group.  Note that \sphinxstyleemphasis{K} thus represents the total force on the group of
atoms, not a per\sphinxhyphen{}atom force.

\sphinxAtStartPar
In \sphinxstyleemphasis{cvel} mode the distance R is incremented or decremented
monotonously according to the pulling (or pushing) velocity.
In \sphinxstyleemphasis{cfor} mode a constant force is added and the actual distance
in direction of the spring is recorded.

\sphinxAtStartPar
The \sphinxstyleemphasis{couple} style links two groups of atoms together.  The first
group is the fix group; the second is specified by group\sphinxhyphen{}ID2.  The
groups are coupled together by a spring that is at equilibrium when
the two groups are displaced by a vector in direction \sphinxstyleemphasis{x,y,z} with
respect to each other and at a distance R0 from that displacement.
Note that \sphinxstyleemphasis{x,y,z} only provides a direction and will be internally
normalized. But since it represents the \sphinxstyleemphasis{absolute} displacement of
group\sphinxhyphen{}ID2 relative to the fix group, (1,1,0) is a different spring
than (\sphinxhyphen{}1,\sphinxhyphen{}1,0).  For each vector component, the displacement can be
described with the \sphinxstyleemphasis{auto} parameter. In this case the direction is
re\sphinxhyphen{}computed in every step, which can be useful for steering a local
process where the whole object undergoes some other change.  When the
relative positions and distance between the two groups are not in
equilibrium, the same spring force described above is applied to atoms
in each of the two groups.

\sphinxAtStartPar
For both the \sphinxstyleemphasis{tether} and \sphinxstyleemphasis{couple} styles, any of the x,y,z values can
be specified as NULL which means do not include that dimension in the
distance calculation or force application.

\sphinxAtStartPar
For constant velocity pulling (\sphinxstyleemphasis{cvel} mode), the running integral
over the pulling force in direction of the spring is recorded and
can then later be used to compute the potential of mean force (PMF)
by averaging over multiple independent trajectories along the same
pulling path.


\subsection{Restart, fix\_modify, output, run start/stop, minimize info}
\label{\detokenize{fix_smd:restart-fix-modify-output-run-start-stop-minimize-info}}
\sphinxAtStartPar
The fix stores the direction of the spring, current pulling target
distance and the running PMF to {\hyperref[\detokenize{restart::doc}]{\sphinxcrossref{\DUrole{doc}{binary restart files}}}}.
See the {\hyperref[\detokenize{read_restart::doc}]{\sphinxcrossref{\DUrole{doc}{read\_restart}}}} command for info on how to
re\sphinxhyphen{}specify a fix in an input script that reads a restart file, so that
the operation of the fix continues in an uninterrupted fashion.

\sphinxAtStartPar
The {\hyperref[\detokenize{fix_modify::doc}]{\sphinxcrossref{\DUrole{doc}{fix\_modify}}}} \sphinxstyleemphasis{virial} option is supported by
this fix to add the contribution due to the added forces on atoms to
both the global pressure and per\sphinxhyphen{}atom stress of the system via the
{\hyperref[\detokenize{compute_pressure::doc}]{\sphinxcrossref{\DUrole{doc}{compute pressure}}}} and {\hyperref[\detokenize{compute_stress_atom::doc}]{\sphinxcrossref{\DUrole{doc}{compute
stress/atom}}}} commands.  The former can be
accessed by {\hyperref[\detokenize{thermo_style::doc}]{\sphinxcrossref{\DUrole{doc}{thermodynamic output}}}}.  The default
setting for this fix is {\hyperref[\detokenize{fix_modify::doc}]{\sphinxcrossref{\DUrole{doc}{fix\_modify virial no}}}}.

\sphinxAtStartPar
The {\hyperref[\detokenize{fix_modify::doc}]{\sphinxcrossref{\DUrole{doc}{fix\_modify}}}} \sphinxstyleemphasis{respa} option is supported by this
fix. This allows to set at which level of the {\hyperref[\detokenize{run_style::doc}]{\sphinxcrossref{\DUrole{doc}{r\sphinxhyphen{}RESPA}}}} integrator the fix is adding its forces. Default is the
outermost level.

\sphinxAtStartPar
This fix computes a vector list of 7 quantities, which can be accessed
by various {\hyperref[\detokenize{Howto_output::doc}]{\sphinxcrossref{\DUrole{doc}{output commands}}}}.  The quantities in
the vector are in this order: the x\sphinxhyphen{}, y\sphinxhyphen{}, and z\sphinxhyphen{}component of the
pulling force, the total force in direction of the pull, the
equilibrium distance of the spring, the distance between the two
reference points, and finally the accumulated PMF (the sum of pulling
forces times displacement).

\sphinxAtStartPar
The force is the total force on the group of atoms by the spring.  In
the case of the \sphinxstyleemphasis{couple} style, it is the force on the fix group
(group\sphinxhyphen{}ID) or the negative of the force on the second group
(group\sphinxhyphen{}ID2).  The vector values calculated by this fix are
“extensive”.

\sphinxAtStartPar
No parameter of this fix can be used with the \sphinxstyleemphasis{start/stop} keywords of
the {\hyperref[\detokenize{run::doc}]{\sphinxcrossref{\DUrole{doc}{run}}}} command.  This fix is not invoked during
{\hyperref[\detokenize{minimize::doc}]{\sphinxcrossref{\DUrole{doc}{energy minimization}}}}.


\subsection{Restrictions}
\label{\detokenize{fix_smd:restrictions}}
\sphinxAtStartPar
This fix is part of the EXTRA\sphinxhyphen{}FIX package.  It is only enabled if
LAMMPS was built with that package.  See the {\hyperref[\detokenize{Build_package::doc}]{\sphinxcrossref{\DUrole{doc}{Build package}}}} page for more info.


\subsection{Related commands}
\label{\detokenize{fix_smd:related-commands}}
\sphinxAtStartPar
{\hyperref[\detokenize{fix_drag::doc}]{\sphinxcrossref{\DUrole{doc}{fix drag}}}}, {\hyperref[\detokenize{fix_spring::doc}]{\sphinxcrossref{\DUrole{doc}{fix spring}}}},
{\hyperref[\detokenize{fix_spring_self::doc}]{\sphinxcrossref{\DUrole{doc}{fix spring/self}}}},
{\hyperref[\detokenize{fix_spring_rg::doc}]{\sphinxcrossref{\DUrole{doc}{fix spring/rg}}}},
{\hyperref[\detokenize{fix_colvars::doc}]{\sphinxcrossref{\DUrole{doc}{fix colvars}}}}, {\hyperref[\detokenize{fix_plumed::doc}]{\sphinxcrossref{\DUrole{doc}{fix plumed}}}}


\subsection{Default}
\label{\detokenize{fix_smd:default}}
\sphinxAtStartPar
none


\bigskip\hrule\bigskip

\phantomsection\label{\detokenize{fix_smd:izrailev}}
\sphinxAtStartPar
\sphinxstylestrong{(Izrailev)} Izrailev, Stepaniants, Isralewitz, Kosztin, Lu, Molnar,
Wriggers, Schulten. Computational Molecular Dynamics: Challenges,
Methods, Ideas, volume 4 of Lecture Notes in Computational Science and
Engineering, pp. 39\sphinxhyphen{}65. Springer\sphinxhyphen{}Verlag, Berlin, 1998.

\phantomsection\label{\detokenize{fix_smd:park}}
\sphinxAtStartPar
\sphinxstylestrong{(Park)} Park, Schulten, J. Chem. Phys. 120 (13), 5946 (2004)

\phantomsection\label{\detokenize{fix_smd:jarzynski}}
\sphinxAtStartPar
\sphinxstylestrong{(Jarzynski)} Jarzynski, Phys. Rev. Lett. 78, 2690 (1997)

\sphinxstepscope

\index{fix smd/adjust\_dt@\spxentry{fix smd/adjust\_dt}}\ignorespaces 

\section{fix smd/adjust\_dt command}
\label{\detokenize{fix_smd_adjust_dt:fix-smd-adjust-dt-command}}\label{\detokenize{fix_smd_adjust_dt:index-0}}\label{\detokenize{fix_smd_adjust_dt::doc}}

\subsection{Syntax}
\label{\detokenize{fix_smd_adjust_dt:syntax}}
\begin{sphinxVerbatim}[commandchars=\\\{\}]
\PYG{k}{fix }\PYG{n+nv+nvIdentifier}{ID}\PYG{+w}{ }\PYG{n+nv+nvIdentifier}{group\PYGZhy{}ID}\PYG{+w}{ }\PYG{n}{smd}\PYG{o}{/}\PYG{n}{adjust\PYGZus{}dt}\PYG{+w}{ }\PYG{n}{arg}
\end{sphinxVerbatim}
\begin{itemize}
\item {} 
\sphinxAtStartPar
ID, group\sphinxhyphen{}ID are documented in {\hyperref[\detokenize{fix::doc}]{\sphinxcrossref{\DUrole{doc}{fix}}}} command

\item {} 
\sphinxAtStartPar
smd/adjust\_dt = style name of this fix command

\item {} 
\sphinxAtStartPar
arg = \sphinxstyleemphasis{s\_fact}
\begin{sphinxalltt}
\sphinxstyleemphasis{s\_fact} = safety factor
\end{sphinxalltt}

\end{itemize}


\subsection{Examples}
\label{\detokenize{fix_smd_adjust_dt:examples}}
\begin{sphinxVerbatim}[commandchars=\\\{\}]
\PYG{k}{fix }\PYG{n+nv+nvIdentifier}{1}\PYG{+w}{ }\PYG{n+nv+nvIdentifier}{all}\PYG{+w}{ }\PYG{n}{smd}\PYG{o}{/}\PYG{n}{adjust\PYGZus{}dt}\PYG{+w}{ }\PYG{l+m}{0.1}
\end{sphinxVerbatim}


\subsection{Description}
\label{\detokenize{fix_smd_adjust_dt:description}}
\sphinxAtStartPar
The fix calculates a new stable time increment for use with the SMD
time integrators.

\sphinxAtStartPar
The stable time increment is based on multiple conditions. For the SPH
pair styles, a CFL criterion (Courant, Friedrichs \& Lewy, 1928) is
evaluated, which determines the speed of sound cannot propagate
further than a typical spacing between particles within a single time
step to ensure no information is lost. For the contact pair styles, a
linear analysis of the pair potential determines a stable maximum time
step.

\sphinxAtStartPar
This fix inquires the minimum stable time increment across all
particles contained in the group for which this fix is defined. An
additional safety factor \sphinxstyleemphasis{s\_fact} is applied to the time increment.

\sphinxAtStartPar
See \sphinxhref{PDF/MACHDYN\_LAMMPS\_userguide.pdf}{this PDF guide} to use Smooth Mach
Dynamics in LAMMPS.


\subsection{Restart, fix\_modify, output, run start/stop, minimize info}
\label{\detokenize{fix_smd_adjust_dt:restart-fix-modify-output-run-start-stop-minimize-info}}
\sphinxAtStartPar
Currently, no part of MACHDYN supports restarting nor minimization.


\subsection{Restrictions}
\label{\detokenize{fix_smd_adjust_dt:restrictions}}
\sphinxAtStartPar
This fix is part of the MACHDYN package.  It is only enabled if
LAMMPS was built with that package.  See the {\hyperref[\detokenize{Build_package::doc}]{\sphinxcrossref{\DUrole{doc}{Build package}}}} page for more info.


\subsection{Related commands}
\label{\detokenize{fix_smd_adjust_dt:related-commands}}
\sphinxAtStartPar
{\hyperref[\detokenize{compute_smd_tlsph_dt::doc}]{\sphinxcrossref{\DUrole{doc}{smd/tlsph\_dt}}}}


\subsection{Default}
\label{\detokenize{fix_smd_adjust_dt:default}}
\sphinxAtStartPar
none

\sphinxstepscope

\index{fix smd/integrate\_tlsph@\spxentry{fix smd/integrate\_tlsph}}\ignorespaces 

\section{fix smd/integrate\_tlsph command}
\label{\detokenize{fix_smd_integrate_tlsph:fix-smd-integrate-tlsph-command}}\label{\detokenize{fix_smd_integrate_tlsph:index-0}}\label{\detokenize{fix_smd_integrate_tlsph::doc}}

\subsection{Syntax}
\label{\detokenize{fix_smd_integrate_tlsph:syntax}}
\begin{sphinxVerbatim}[commandchars=\\\{\}]
\PYG{k}{fix }\PYG{n+nv+nvIdentifier}{ID}\PYG{+w}{ }\PYG{n+nv+nvIdentifier}{group\PYGZhy{}ID}\PYG{+w}{ }\PYG{n}{smd}\PYG{o}{/}\PYG{n}{integrate\PYGZus{}tlsph}\PYG{+w}{ }\PYG{n}{keyword}\PYG{+w}{ }\PYG{n}{values}
\end{sphinxVerbatim}
\begin{itemize}
\item {} 
\sphinxAtStartPar
ID, group\sphinxhyphen{}ID are documented in {\hyperref[\detokenize{fix::doc}]{\sphinxcrossref{\DUrole{doc}{fix}}}} command

\item {} 
\sphinxAtStartPar
smd/integrate\_tlsph = style name of this fix command

\item {} 
\sphinxAtStartPar
zero or more keyword/value pairs may be appended

\item {} 
\sphinxAtStartPar
keyword = \sphinxstyleemphasis{limit\_velocity}

\end{itemize}
\begin{sphinxalltt}
\sphinxstyleemphasis{limit\_velocity} value = max\_vel
  max\_vel = maximum allowed velocity
\end{sphinxalltt}


\subsection{Examples}
\label{\detokenize{fix_smd_integrate_tlsph:examples}}
\begin{sphinxVerbatim}[commandchars=\\\{\}]
\PYG{k}{fix }\PYG{n+nv+nvIdentifier}{1}\PYG{+w}{ }\PYG{n+nv+nvIdentifier}{all}\PYG{+w}{ }\PYG{n}{smd}\PYG{o}{/}\PYG{n}{integrate\PYGZus{}tlsph}
\PYG{k}{fix }\PYG{n+nv+nvIdentifier}{1}\PYG{+w}{ }\PYG{n+nv+nvIdentifier}{all}\PYG{+w}{ }\PYG{n}{smd}\PYG{o}{/}\PYG{n}{integrate\PYGZus{}tlsph}\PYG{+w}{ }\PYG{n}{limit\PYGZus{}velocity}\PYG{+w}{ }\PYG{l+m}{1000}
\end{sphinxVerbatim}


\subsection{Description}
\label{\detokenize{fix_smd_integrate_tlsph:description}}
\sphinxAtStartPar
The fix performs explicit time integration for particles which
interact according with the Total\sphinxhyphen{}Lagrangian SPH pair style.

\sphinxAtStartPar
See \sphinxhref{PDF/MACHDYN\_LAMMPS\_userguide.pdf}{this PDF guide} to using Smooth Mach
Dynamics in LAMMPS.

\sphinxAtStartPar
The \sphinxstyleemphasis{limit\_velocity} keyword will control the velocity, scaling the
norm of the velocity vector to max\_vel in case it exceeds this
velocity limit.


\subsection{Restart, fix\_modify, output, run start/stop, minimize info}
\label{\detokenize{fix_smd_integrate_tlsph:restart-fix-modify-output-run-start-stop-minimize-info}}
\sphinxAtStartPar
Currently, no part of MACHDYN supports restarting nor
minimization. This fix has no outputs.


\subsection{Restrictions}
\label{\detokenize{fix_smd_integrate_tlsph:restrictions}}
\sphinxAtStartPar
This fix is part of the MACHDYN package.  It is only enabled if
LAMMPS was built with that package.  See the {\hyperref[\detokenize{Build_package::doc}]{\sphinxcrossref{\DUrole{doc}{Build package}}}} page for more info.

\sphinxAtStartPar
\DUrole{versionmodified,changed}{Changed in version 29Aug2024.}

\sphinxAtStartPar
This fix is incompatible with deformation controls that remap velocity,
for instance the \sphinxstyleemphasis{remap v} option of {\hyperref[\detokenize{fix_deform::doc}]{\sphinxcrossref{\DUrole{doc}{fix deform}}}}.


\subsection{Related commands}
\label{\detokenize{fix_smd_integrate_tlsph:related-commands}}
\sphinxAtStartPar
{\hyperref[\detokenize{fix_smd_integrate_ulsph::doc}]{\sphinxcrossref{\DUrole{doc}{smd/integrate\_ulsph}}}}


\subsection{Default}
\label{\detokenize{fix_smd_integrate_tlsph:default}}
\sphinxAtStartPar
none

\sphinxstepscope

\index{fix smd/integrate\_ulsph@\spxentry{fix smd/integrate\_ulsph}}\ignorespaces 

\section{fix smd/integrate\_ulsph command}
\label{\detokenize{fix_smd_integrate_ulsph:fix-smd-integrate-ulsph-command}}\label{\detokenize{fix_smd_integrate_ulsph:index-0}}\label{\detokenize{fix_smd_integrate_ulsph::doc}}

\subsection{Syntax}
\label{\detokenize{fix_smd_integrate_ulsph:syntax}}
\begin{sphinxVerbatim}[commandchars=\\\{\}]
\PYG{k}{fix }\PYG{n+nv+nvIdentifier}{ID}\PYG{+w}{ }\PYG{n+nv+nvIdentifier}{group\PYGZhy{}ID}\PYG{+w}{ }\PYG{n}{smd}\PYG{o}{/}\PYG{n}{integrate\PYGZus{}ulsph}\PYG{+w}{ }\PYG{n}{keyword}
\end{sphinxVerbatim}
\begin{itemize}
\item {} 
\sphinxAtStartPar
ID, group\sphinxhyphen{}ID are documented in {\hyperref[\detokenize{fix::doc}]{\sphinxcrossref{\DUrole{doc}{fix}}}} command

\item {} 
\sphinxAtStartPar
smd/integrate\_ulsph = style name of this fix command

\item {} 
\sphinxAtStartPar
zero or more keyword/value pairs may be appended

\item {} 
\sphinxAtStartPar
keyword = adjust\_radius or limit\_velocity

\end{itemize}
\begin{description}
\sphinxlineitem{adjust\_radius values = adjust\_radius\_factor min\_nn max\_nn}
\sphinxAtStartPar
adjust\_radius\_factor = factor which scale the smooth/kernel radius
min\_nn = minimum number of neighbors
max\_nn = maximum number of neighbors

\sphinxlineitem{limit\_velocity values = max\_velocity}
\sphinxAtStartPar
max\_velocity = maximum allowed velocity.

\end{description}


\subsection{Examples}
\label{\detokenize{fix_smd_integrate_ulsph:examples}}
\begin{sphinxVerbatim}[commandchars=\\\{\}]
\PYG{k}{fix }\PYG{n+nv+nvIdentifier}{1}\PYG{+w}{ }\PYG{n+nv+nvIdentifier}{all}\PYG{+w}{ }\PYG{n}{smd}\PYG{o}{/}\PYG{n}{integrate\PYGZus{}ulsph}\PYG{+w}{ }\PYG{n}{adjust\PYGZus{}radius}\PYG{+w}{ }\PYG{l+m}{1.02}\PYG{+w}{ }\PYG{l+m}{25}\PYG{+w}{ }\PYG{l+m}{50}
\PYG{k}{fix }\PYG{n+nv+nvIdentifier}{1}\PYG{+w}{ }\PYG{n+nv+nvIdentifier}{all}\PYG{+w}{ }\PYG{n}{smd}\PYG{o}{/}\PYG{n}{integrate\PYGZus{}ulsph}\PYG{+w}{ }\PYG{n}{limit\PYGZus{}velocity}\PYG{+w}{ }\PYG{l+m}{1000}
\end{sphinxVerbatim}


\subsection{Description}
\label{\detokenize{fix_smd_integrate_ulsph:description}}
\sphinxAtStartPar
The fix performs explicit time integration for particles which
interact with the updated Lagrangian SPH pair style.

\sphinxAtStartPar
See \sphinxhref{PDF/MACHDYN\_LAMMPS\_userguide.pdf}{this PDF guide} to using Smooth Mach
Dynamics in LAMMPS.

\sphinxAtStartPar
The \sphinxstyleemphasis{adjust\_radius} keyword activates dynamic adjustment of the
per\sphinxhyphen{}particle SPH smoothing kernel radius such that the number of
neighbors per particles remains within the interval \sphinxstyleemphasis{min\_nn} to
\sphinxstyleemphasis{max\_nn}. The parameter \sphinxstyleemphasis{adjust\_radius\_factor} determines the amount
of adjustment per timestep. Typical values are \sphinxstyleemphasis{adjust\_radius\_factor}
=1.02, \sphinxstyleemphasis{min\_nn} =15, and \sphinxstyleemphasis{max\_nn} =20.

\sphinxAtStartPar
The \sphinxstyleemphasis{limit\_velocity} keyword will control the velocity, scaling the norm of
the velocity vector to max\_vel in case it exceeds this velocity limit.


\subsection{Restart, fix\_modify, output, run start/stop, minimize info}
\label{\detokenize{fix_smd_integrate_ulsph:restart-fix-modify-output-run-start-stop-minimize-info}}
\sphinxAtStartPar
Currently, no part of MACHDYN supports restarting nor
minimization. This fix has no outputs.


\subsection{Restrictions}
\label{\detokenize{fix_smd_integrate_ulsph:restrictions}}
\sphinxAtStartPar
This fix is part of the MACHDYN package.  It is only enabled if
LAMMPS was built with that package.  See the {\hyperref[\detokenize{Build_package::doc}]{\sphinxcrossref{\DUrole{doc}{Build package}}}} page for more info.

\sphinxAtStartPar
\DUrole{versionmodified,changed}{Changed in version 29Aug2024.}

\sphinxAtStartPar
This fix is incompatible with deformation controls that remap velocity,
for instance the \sphinxstyleemphasis{remap v} option of {\hyperref[\detokenize{fix_deform::doc}]{\sphinxcrossref{\DUrole{doc}{fix deform}}}}.


\subsection{Related commands}
\label{\detokenize{fix_smd_integrate_ulsph:related-commands}}
\sphinxAtStartPar
none


\subsection{Default}
\label{\detokenize{fix_smd_integrate_ulsph:default}}
\sphinxAtStartPar
none

\sphinxstepscope

\index{fix smd/move\_tri\_surf@\spxentry{fix smd/move\_tri\_surf}}\ignorespaces 

\section{fix smd/move\_tri\_surf command}
\label{\detokenize{fix_smd_move_triangulated_surface:fix-smd-move-tri-surf-command}}\label{\detokenize{fix_smd_move_triangulated_surface:index-0}}\label{\detokenize{fix_smd_move_triangulated_surface::doc}}

\subsection{Syntax}
\label{\detokenize{fix_smd_move_triangulated_surface:syntax}}
\begin{sphinxVerbatim}[commandchars=\\\{\}]
\PYG{k}{fix }\PYG{n+nv+nvIdentifier}{ID}\PYG{+w}{ }\PYG{n+nv+nvIdentifier}{group\PYGZhy{}ID}\PYG{+w}{ }\PYG{n}{smd}\PYG{o}{/}\PYG{n}{move\PYGZus{}tri\PYGZus{}surf}\PYG{+w}{ }\PYG{n}{keyword}
\end{sphinxVerbatim}
\begin{itemize}
\item {} 
\sphinxAtStartPar
ID, group\sphinxhyphen{}ID are documented in {\hyperref[\detokenize{fix::doc}]{\sphinxcrossref{\DUrole{doc}{fix}}}} command

\item {} 
\sphinxAtStartPar
smd/move\_tri\_surf keyword = style name of this fix command

\item {} 
\sphinxAtStartPar
keyword = \sphinxstyleemphasis{*LINEAR} or \sphinxstyleemphasis{*WIGGLE} or \sphinxstyleemphasis{*ROTATE}
\begin{sphinxalltt}
\sphinxstyleemphasis{*LINEAR} args = Vx Vy Vz
   Vx,Vy,Vz = components of velocity vector (velocity units), any component can be specified as NULL
\sphinxstyleemphasis{*WIGGLE} args = Vx Vy Vz max\_travel
   vx,vy,vz = components of velocity vector (velocity units), any component can be specified as NULL
   max\_travel = wiggle amplitude
\sphinxstyleemphasis{*ROTATE} args = Px Py Pz Rx Ry Rz period
   Px,Py,Pz = origin point of axis of rotation (distance units)
   Rx,Ry,Rz = axis of rotation vector
   period = period of rotation (time units)
\end{sphinxalltt}

\end{itemize}


\subsection{Examples}
\label{\detokenize{fix_smd_move_triangulated_surface:examples}}
\begin{sphinxVerbatim}[commandchars=\\\{\}]
\PYG{k}{fix }\PYG{n+nv+nvIdentifier}{1}\PYG{+w}{ }\PYG{n+nv+nvIdentifier}{tool}\PYG{+w}{ }\PYG{n}{smd}\PYG{o}{/}\PYG{n}{move\PYGZus{}tri\PYGZus{}surf}\PYG{+w}{ }\PYG{o}{*}\PYG{n}{LINEAR}\PYG{+w}{ }\PYG{l+m}{20}\PYG{+w}{ }\PYG{l+m}{20}\PYG{+w}{ }\PYG{l+m}{10}
\PYG{k}{fix }\PYG{n+nv+nvIdentifier}{2}\PYG{+w}{ }\PYG{n+nv+nvIdentifier}{tool}\PYG{+w}{ }\PYG{n}{smd}\PYG{o}{/}\PYG{n}{move\PYGZus{}tri\PYGZus{}surf}\PYG{+w}{ }\PYG{o}{*}\PYG{n}{WIGGLE}\PYG{+w}{ }\PYG{l+m}{20}\PYG{+w}{ }\PYG{l+m}{20}\PYG{+w}{ }\PYG{l+m}{10}
\PYG{k}{fix }\PYG{n+nv+nvIdentifier}{2}\PYG{+w}{ }\PYG{n+nv+nvIdentifier}{tool}\PYG{+w}{ }\PYG{n}{smd}\PYG{o}{/}\PYG{n}{move\PYGZus{}tri\PYGZus{}surf}\PYG{+w}{ }\PYG{o}{*}\PYG{n}{ROTATE}\PYG{+w}{ }\PYG{l+m}{0}\PYG{+w}{ }\PYG{l+m}{0}\PYG{+w}{ }\PYG{l+m}{0}\PYG{+w}{ }\PYG{l+m}{5}\PYG{+w}{ }\PYG{l+m}{2}\PYG{+w}{ }\PYG{l+m}{1}
\end{sphinxVerbatim}


\subsection{Description}
\label{\detokenize{fix_smd_move_triangulated_surface:description}}
\sphinxAtStartPar
This fix applies only to rigid surfaces read from .STL files via fix
{\hyperref[\detokenize{fix_smd_wall_surface::doc}]{\sphinxcrossref{\DUrole{doc}{smd/wall\_surface}}}} .  It updates position
and velocity for the particles in the group each timestep without
regard to forces on the particles.  The rigid surfaces can thus be
moved along simple trajectories during the simulation.

\sphinxAtStartPar
The \sphinxstyleemphasis{*LINEAR} style moves particles with the specified constant velocity
vector V = (Vx,Vy,Vz). This style also sets the velocity of each particle
to V = (Vx,Vy,Vz).

\sphinxAtStartPar
The \sphinxstyleemphasis{*WIGGLE} style moves particles in an oscillatory fashion.
Particles are moved along (vx, vy, vz) with constant velocity until a
displacement of max\_travel is reached. Then, the velocity vector is
reversed. This process is repeated.

\sphinxAtStartPar
The \sphinxstyleemphasis{*ROTATE} style rotates particles around a rotation axis R =
(Rx,Ry,Rz) that goes through a point P = (Px,Py,Pz). The period of the
rotation is also specified. This style also sets the velocity of each
particle to (omega cross Rperp) where omega is its angular velocity
around the rotation axis and Rperp is a perpendicular vector from the
rotation axis to the particle.

\sphinxAtStartPar
See \sphinxhref{PDF/MACHDYN\_LAMMPS\_userguide.pdf}{this PDF guide} to using Smooth Mach
Dynamics in LAMMPS.


\subsection{Restart, fix\_modify, output, run start/stop, minimize info}
\label{\detokenize{fix_smd_move_triangulated_surface:restart-fix-modify-output-run-start-stop-minimize-info}}
\sphinxAtStartPar
Currently, no part of MACHDYN supports restarting nor
minimization. This fix has no outputs.


\subsection{Restrictions}
\label{\detokenize{fix_smd_move_triangulated_surface:restrictions}}
\sphinxAtStartPar
This fix is part of the MACHDYN package.  It is only enabled if
LAMMPS was built with that package.  See the {\hyperref[\detokenize{Build_package::doc}]{\sphinxcrossref{\DUrole{doc}{Build package}}}} page for more info.


\subsection{Related commands}
\label{\detokenize{fix_smd_move_triangulated_surface:related-commands}}
\sphinxAtStartPar
{\hyperref[\detokenize{compute_smd_triangle_vertices::doc}]{\sphinxcrossref{\DUrole{doc}{smd/triangle\_mesh\_vertices}}}},
{\hyperref[\detokenize{fix_smd_wall_surface::doc}]{\sphinxcrossref{\DUrole{doc}{smd/wall\_surface}}}}


\subsection{Default}
\label{\detokenize{fix_smd_move_triangulated_surface:default}}
\sphinxAtStartPar
none

\sphinxstepscope

\index{fix smd/setvel@\spxentry{fix smd/setvel}}\ignorespaces 

\section{fix smd/setvel command}
\label{\detokenize{fix_smd_setvel:fix-smd-setvel-command}}\label{\detokenize{fix_smd_setvel:index-0}}\label{\detokenize{fix_smd_setvel::doc}}

\subsection{Syntax}
\label{\detokenize{fix_smd_setvel:syntax}}
\begin{sphinxVerbatim}[commandchars=\\\{\}]
\PYG{k}{fix }\PYG{n+nv+nvIdentifier}{ID}\PYG{+w}{ }\PYG{n+nv+nvIdentifier}{group\PYGZhy{}ID}\PYG{+w}{ }\PYG{n}{smd}\PYG{o}{/}\PYG{n}{setvel}\PYG{+w}{ }\PYG{n}{vx}\PYG{+w}{ }\PYG{n}{vy}\PYG{+w}{ }\PYG{n}{vz}\PYG{+w}{ }\PYG{n}{keyword}\PYG{+w}{ }\PYG{n}{value}\PYG{+w}{ }\PYG{n}{...}
\end{sphinxVerbatim}
\begin{itemize}
\item {} 
\sphinxAtStartPar
ID, group\sphinxhyphen{}ID are documented in {\hyperref[\detokenize{fix::doc}]{\sphinxcrossref{\DUrole{doc}{fix}}}} command

\item {} 
\sphinxAtStartPar
smd/setvel = style name of this fix command

\item {} 
\sphinxAtStartPar
vx,vy,vz = velocity component values

\item {} 
\sphinxAtStartPar
any of vx,vy,vz can be a variable (see below)

\item {} 
\sphinxAtStartPar
zero or more keyword/value pairs may be appended to args

\item {} 
\sphinxAtStartPar
keyword = \sphinxstyleemphasis{region}
\begin{sphinxalltt}
\sphinxstyleemphasis{region} value = region\sphinxhyphen{}ID
  region\sphinxhyphen{}ID = ID of region particles must be in to have their velocities set
\end{sphinxalltt}

\end{itemize}


\subsection{Examples}
\label{\detokenize{fix_smd_setvel:examples}}
\begin{sphinxVerbatim}[commandchars=\\\{\}]
\PYG{k}{fix }\PYG{n+nv+nvIdentifier}{top\PYGZus{}velocity}\PYG{+w}{ }\PYG{n+nv+nvIdentifier}{top\PYGZus{}group}\PYG{+w}{ }\PYG{n}{smd}\PYG{o}{/}\PYG{n}{setvel}\PYG{+w}{ }\PYG{l+m}{1.0}\PYG{+w}{ }\PYG{l+m}{0.0}\PYG{+w}{ }\PYG{l+m}{0.0}
\end{sphinxVerbatim}


\subsection{Description}
\label{\detokenize{fix_smd_setvel:description}}
\sphinxAtStartPar
Set each component of velocity on each particle in the group to the specified
values vx,vy,vz, regardless of the forces acting on the particle.  This command can
be used to impose velocity boundary conditions.

\sphinxAtStartPar
Any of the vx,vy,vz values can be specified as NULL which means do not
alter the velocity component in that dimension.

\sphinxAtStartPar
This fix is indented to be used together with a time integration fix.

\sphinxAtStartPar
Any of the 3 quantities defining the velocity components can be specified
as an equal\sphinxhyphen{}style or atom\sphinxhyphen{}style {\hyperref[\detokenize{variable::doc}]{\sphinxcrossref{\DUrole{doc}{variable}}}}, namely \sphinxstyleemphasis{vx},
\sphinxstyleemphasis{vy}, \sphinxstyleemphasis{vz}.  If the value is a variable, it should be specified as
v\_name, where name is the variable name.  In this case, the variable
will be evaluated each timestep, and its value used to determine the
force component.

\sphinxAtStartPar
Equal\sphinxhyphen{}style variables can specify formulas with various mathematical
functions, and include {\hyperref[\detokenize{thermo_style::doc}]{\sphinxcrossref{\DUrole{doc}{thermo\_style}}}} command
keywords for the simulation box parameters and timestep and elapsed
time.  Thus it is easy to specify a time\sphinxhyphen{}dependent velocity field.

\sphinxAtStartPar
Atom\sphinxhyphen{}style variables can specify the same formulas as equal\sphinxhyphen{}style
variables but can also include per\sphinxhyphen{}atom values, such as atom
coordinates.  Thus it is easy to specify a spatially\sphinxhyphen{}dependent velocity
field with optional time\sphinxhyphen{}dependence as well.

\sphinxAtStartPar
If the \sphinxstyleemphasis{region} keyword is used, the particle must also be in the
specified geometric {\hyperref[\detokenize{region::doc}]{\sphinxcrossref{\DUrole{doc}{region}}}} in order to have its velocity set by this command.


\bigskip\hrule\bigskip



\subsection{Restart, fix\_modify, output, run start/stop, minimize info}
\label{\detokenize{fix_smd_setvel:restart-fix-modify-output-run-start-stop-minimize-info}}
\sphinxAtStartPar
Currently, no part of MACHDYN supports restarting nor minimization
None of the {\hyperref[\detokenize{fix_modify::doc}]{\sphinxcrossref{\DUrole{doc}{fix\_modify}}}} options
are relevant to this fix.

\sphinxAtStartPar
This fix computes a global 3\sphinxhyphen{}vector of forces, which can be accessed
by various {\hyperref[\detokenize{Howto_output::doc}]{\sphinxcrossref{\DUrole{doc}{output commands}}}}.  This is the total
force on the group of atoms.  The vector values calculated by this fix
are “extensive”.

\sphinxAtStartPar
No parameter of this fix can be used with the \sphinxstyleemphasis{start/stop} keywords of
the {\hyperref[\detokenize{run::doc}]{\sphinxcrossref{\DUrole{doc}{run}}}} command.


\subsection{Restrictions}
\label{\detokenize{fix_smd_setvel:restrictions}}
\sphinxAtStartPar
This fix is part of the MACHDYN package.  It is only enabled if
LAMMPS was built with that package.  See the {\hyperref[\detokenize{Build_package::doc}]{\sphinxcrossref{\DUrole{doc}{Build package}}}} page for more info.


\subsection{Related commands}
\label{\detokenize{fix_smd_setvel:related-commands}}
\sphinxAtStartPar
none


\subsection{Default}
\label{\detokenize{fix_smd_setvel:default}}
\sphinxAtStartPar
none

\sphinxstepscope

\index{fix smd/wall\_surface@\spxentry{fix smd/wall\_surface}}\ignorespaces 

\section{fix smd/wall\_surface command}
\label{\detokenize{fix_smd_wall_surface:fix-smd-wall-surface-command}}\label{\detokenize{fix_smd_wall_surface:index-0}}\label{\detokenize{fix_smd_wall_surface::doc}}

\subsection{Syntax}
\label{\detokenize{fix_smd_wall_surface:syntax}}
\begin{sphinxVerbatim}[commandchars=\\\{\}]
\PYG{k}{fix }\PYG{n+nv+nvIdentifier}{ID}\PYG{+w}{ }\PYG{n+nv+nvIdentifier}{group\PYGZhy{}ID}\PYG{+w}{ }\PYG{n}{smd}\PYG{o}{/}\PYG{n}{wall\PYGZus{}surface}\PYG{+w}{ }\PYG{n}{arg}\PYG{+w}{ }\PYG{n}{type}\PYG{+w}{ }\PYG{n}{mol}\PYG{o}{\PYGZhy{}}\PYG{n}{ID}
\end{sphinxVerbatim}
\begin{itemize}
\item {} 
\sphinxAtStartPar
ID, group\sphinxhyphen{}ID are documented in {\hyperref[\detokenize{fix::doc}]{\sphinxcrossref{\DUrole{doc}{fix}}}} command

\item {} 
\sphinxAtStartPar
smd/wall\_surface = style name of this fix command

\item {} 
\sphinxAtStartPar
arg = \sphinxstyleemphasis{file}
\begin{sphinxalltt}
\sphinxstyleemphasis{file} = file name of a triangular mesh in stl format
\end{sphinxalltt}

\item {} 
\sphinxAtStartPar
type = particle type to be given to the new particles created by this fix

\item {} 
\sphinxAtStartPar
mol\sphinxhyphen{}ID = molecule\sphinxhyphen{}ID to be given to the new particles created by this fix (must be \textgreater{}= 65535)

\end{itemize}


\subsection{Examples}
\label{\detokenize{fix_smd_wall_surface:examples}}
\begin{sphinxVerbatim}[commandchars=\\\{\}]
\PYG{k}{fix }\PYG{n+nv+nvIdentifier}{stl\PYGZus{}surf}\PYG{+w}{ }\PYG{n+nv+nvIdentifier}{all}\PYG{+w}{ }\PYG{n}{smd}\PYG{o}{/}\PYG{n}{wall\PYGZus{}surface}\PYG{+w}{ }\PYG{n}{tool.stl}\PYG{+w}{ }\PYG{l+m}{2}\PYG{+w}{ }\PYG{l+m}{65535}
\end{sphinxVerbatim}


\subsection{Description}
\label{\detokenize{fix_smd_wall_surface:description}}
\sphinxAtStartPar
This fix creates reads a triangulated surface from a file in .STL
format.  For each triangle, a new particle is created which stores the
barycenter of the triangle and the vertex positions.  The radius of
the new particle is that of the minimum circle which encompasses the
triangle vertices.

\sphinxAtStartPar
The triangulated surface can be used as a complex rigid wall via the
{\hyperref[\detokenize{pair_smd_triangulated_surface::doc}]{\sphinxcrossref{\DUrole{doc}{smd/tri\_surface}}}} pair style.  It
is possible to move the triangulated surface via the
{\hyperref[\detokenize{fix_smd_move_triangulated_surface::doc}]{\sphinxcrossref{\DUrole{doc}{smd/move\_tri\_surf}}}} fix style.

\sphinxAtStartPar
Immediately after a .STL file has been read, the simulation needs to
be run for 0 timesteps in order to properly register the new particles
in the system. See the “funnel\_flow” example in the MACHDYN examples
directory.

\sphinxAtStartPar
See \sphinxhref{PDF/MACHDYN\_LAMMPS\_userguide.pdf}{this PDF guide} to use Smooth Mach
Dynamics in LAMMPS.


\subsection{Restart, fix\_modify, output, run start/stop, minimize info}
\label{\detokenize{fix_smd_wall_surface:restart-fix-modify-output-run-start-stop-minimize-info}}
\sphinxAtStartPar
Currently, no part of MACHDYN supports restarting nor
minimization. This fix has no outputs.


\subsection{Restrictions}
\label{\detokenize{fix_smd_wall_surface:restrictions}}
\sphinxAtStartPar
This fix is part of the MACHDYN package.  It is only enabled if
LAMMPS was built with that package.  See the {\hyperref[\detokenize{Build_package::doc}]{\sphinxcrossref{\DUrole{doc}{Build package}}}} page for more info.

\sphinxAtStartPar
The molecule ID given to the particles created by this fix have to be
equal to or larger than 65535.

\sphinxAtStartPar
Within each .STL file, only a single triangulated object must be
present, even though the STL format allows for the possibility of
multiple objects in one file.


\subsection{Related commands}
\label{\detokenize{fix_smd_wall_surface:related-commands}}
\sphinxAtStartPar
{\hyperref[\detokenize{compute_smd_triangle_vertices::doc}]{\sphinxcrossref{\DUrole{doc}{smd/triangle\_mesh\_vertices}}}},
{\hyperref[\detokenize{fix_smd_move_triangulated_surface::doc}]{\sphinxcrossref{\DUrole{doc}{smd/move\_tri\_surf}}}},
{\hyperref[\detokenize{pair_smd_triangulated_surface::doc}]{\sphinxcrossref{\DUrole{doc}{smd/tri\_surface}}}}


\subsection{Default}
\label{\detokenize{fix_smd_wall_surface:default}}
\sphinxAtStartPar
none

\sphinxstepscope

\index{fix sph@\spxentry{fix sph}}\ignorespaces 

\section{fix sph command}
\label{\detokenize{fix_sph:fix-sph-command}}\label{\detokenize{fix_sph:index-0}}\label{\detokenize{fix_sph::doc}}

\subsection{Syntax}
\label{\detokenize{fix_sph:syntax}}
\begin{sphinxVerbatim}[commandchars=\\\{\}]
\PYG{k}{fix }\PYG{n+nv+nvIdentifier}{ID}\PYG{+w}{ }\PYG{n+nv+nvIdentifier}{group\PYGZhy{}ID}\PYG{+w}{ }\PYG{n}{sph}
\end{sphinxVerbatim}
\begin{itemize}
\item {} 
\sphinxAtStartPar
ID, group\sphinxhyphen{}ID are documented in {\hyperref[\detokenize{fix::doc}]{\sphinxcrossref{\DUrole{doc}{fix}}}} command

\item {} 
\sphinxAtStartPar
sph = style name of this fix command

\end{itemize}


\subsection{Examples}
\label{\detokenize{fix_sph:examples}}
\begin{sphinxVerbatim}[commandchars=\\\{\}]
\PYG{k}{fix }\PYG{n+nv+nvIdentifier}{1}\PYG{+w}{ }\PYG{n+nv+nvIdentifier}{all}\PYG{+w}{ }\PYG{n}{sph}
\end{sphinxVerbatim}


\subsection{Description}
\label{\detokenize{fix_sph:description}}
\sphinxAtStartPar
Perform time integration to update position, velocity, internal energy
and local density for atoms in the group each timestep. This fix is
needed to time\sphinxhyphen{}integrate SPH systems where particles carry internal
variables such as internal energy.  SPH stands for Smoothed Particle
Hydrodynamics.

\sphinxAtStartPar
See \sphinxhref{PDF/SPH\_LAMMPS\_userguide.pdf}{this PDF guide} to using SPH in
LAMMPS.


\subsection{Restart, fix\_modify, output, run start/stop, minimize info}
\label{\detokenize{fix_sph:restart-fix-modify-output-run-start-stop-minimize-info}}
\sphinxAtStartPar
No information about this fix is written to {\hyperref[\detokenize{restart::doc}]{\sphinxcrossref{\DUrole{doc}{binary restart files}}}}.  None of the {\hyperref[\detokenize{fix_modify::doc}]{\sphinxcrossref{\DUrole{doc}{fix\_modify}}}} options
are relevant to this fix.  No global or per\sphinxhyphen{}atom quantities are stored
by this fix for access by various {\hyperref[\detokenize{Howto_output::doc}]{\sphinxcrossref{\DUrole{doc}{output commands}}}}.
No parameter of this fix can be used with the \sphinxstyleemphasis{start/stop} keywords of
the {\hyperref[\detokenize{run::doc}]{\sphinxcrossref{\DUrole{doc}{run}}}} command.  This fix is not invoked during {\hyperref[\detokenize{minimize::doc}]{\sphinxcrossref{\DUrole{doc}{energy minimization}}}}.


\subsection{Restrictions}
\label{\detokenize{fix_sph:restrictions}}
\sphinxAtStartPar
This fix is part of the SPH package.  It is only enabled if
LAMMPS was built with that package.  See the {\hyperref[\detokenize{Build_package::doc}]{\sphinxcrossref{\DUrole{doc}{Build package}}}} page for more info.


\subsection{Related commands}
\label{\detokenize{fix_sph:related-commands}}
\sphinxAtStartPar
{\hyperref[\detokenize{fix_sph_stationary::doc}]{\sphinxcrossref{\DUrole{doc}{fix sph/stationary}}}}


\subsection{Default}
\label{\detokenize{fix_sph:default}}
\sphinxAtStartPar
none

\sphinxstepscope

\index{fix sph/stationary@\spxentry{fix sph/stationary}}\ignorespaces 

\section{fix sph/stationary command}
\label{\detokenize{fix_sph_stationary:fix-sph-stationary-command}}\label{\detokenize{fix_sph_stationary:index-0}}\label{\detokenize{fix_sph_stationary::doc}}

\subsection{Syntax}
\label{\detokenize{fix_sph_stationary:syntax}}
\begin{sphinxVerbatim}[commandchars=\\\{\}]
\PYG{k}{fix }\PYG{n+nv+nvIdentifier}{ID}\PYG{+w}{ }\PYG{n+nv+nvIdentifier}{group\PYGZhy{}ID}\PYG{+w}{ }\PYG{n}{sph}\PYG{o}{/}\PYG{n}{stationary}
\end{sphinxVerbatim}
\begin{itemize}
\item {} 
\sphinxAtStartPar
ID, group\sphinxhyphen{}ID are documented in {\hyperref[\detokenize{fix::doc}]{\sphinxcrossref{\DUrole{doc}{fix}}}} command

\item {} 
\sphinxAtStartPar
sph = style name of this fix command

\end{itemize}


\subsection{Examples}
\label{\detokenize{fix_sph_stationary:examples}}
\begin{sphinxVerbatim}[commandchars=\\\{\}]
\PYG{k}{fix }\PYG{n+nv+nvIdentifier}{1}\PYG{+w}{ }\PYG{n+nv+nvIdentifier}{boundary}\PYG{+w}{ }\PYG{n}{sph}\PYG{o}{/}\PYG{n}{stationary}
\end{sphinxVerbatim}


\subsection{Description}
\label{\detokenize{fix_sph_stationary:description}}
\sphinxAtStartPar
Perform time integration to update internal energy and local density,
but not position or velocity for atoms in the group each timestep.
This fix is needed for SPH simulations to correctly time\sphinxhyphen{}integrate
fixed boundary particles which constrain a fluid to a given region in
space.  SPH stands for Smoothed Particle Hydrodynamics.

\sphinxAtStartPar
See \sphinxhref{PDF/SPH\_LAMMPS\_userguide.pdf}{this PDF guide} to using SPH in
LAMMPS.


\subsection{Restart, fix\_modify, output, run start/stop, minimize info}
\label{\detokenize{fix_sph_stationary:restart-fix-modify-output-run-start-stop-minimize-info}}
\sphinxAtStartPar
No information about this fix is written to {\hyperref[\detokenize{restart::doc}]{\sphinxcrossref{\DUrole{doc}{binary restart files}}}}.  None of the {\hyperref[\detokenize{fix_modify::doc}]{\sphinxcrossref{\DUrole{doc}{fix\_modify}}}} options
are relevant to this fix.  No global or per\sphinxhyphen{}atom quantities are stored
by this fix for access by various {\hyperref[\detokenize{Howto_output::doc}]{\sphinxcrossref{\DUrole{doc}{output commands}}}}.
No parameter of this fix can be used with the \sphinxstyleemphasis{start/stop} keywords of
the {\hyperref[\detokenize{run::doc}]{\sphinxcrossref{\DUrole{doc}{run}}}} command.  This fix is not invoked during {\hyperref[\detokenize{minimize::doc}]{\sphinxcrossref{\DUrole{doc}{energy minimization}}}}.


\subsection{Restrictions}
\label{\detokenize{fix_sph_stationary:restrictions}}
\sphinxAtStartPar
This fix is part of the SPH package.  It is only enabled if
LAMMPS was built with that package.  See the {\hyperref[\detokenize{Build_package::doc}]{\sphinxcrossref{\DUrole{doc}{Build package}}}} page for more info.


\subsection{Related commands}
\label{\detokenize{fix_sph_stationary:related-commands}}
\sphinxAtStartPar
{\hyperref[\detokenize{fix_sph::doc}]{\sphinxcrossref{\DUrole{doc}{fix sph}}}}


\subsection{Default}
\label{\detokenize{fix_sph_stationary:default}}
\sphinxAtStartPar
none

\sphinxstepscope

\index{fix spring@\spxentry{fix spring}}\ignorespaces 

\section{fix spring command}
\label{\detokenize{fix_spring:fix-spring-command}}\label{\detokenize{fix_spring:index-0}}\label{\detokenize{fix_spring::doc}}

\subsection{Syntax}
\label{\detokenize{fix_spring:syntax}}
\begin{sphinxVerbatim}[commandchars=\\\{\}]
\PYG{k}{fix }\PYG{n+nv+nvIdentifier}{ID}\PYG{+w}{ }\PYG{n+nv+nvIdentifier}{group\PYGZhy{}ID}\PYG{+w}{ }\PYG{n}{spring}\PYG{+w}{ }\PYG{n}{keyword}\PYG{+w}{ }\PYG{n}{values}
\end{sphinxVerbatim}
\begin{itemize}
\item {} 
\sphinxAtStartPar
ID, group\sphinxhyphen{}ID are documented in {\hyperref[\detokenize{fix::doc}]{\sphinxcrossref{\DUrole{doc}{fix}}}} command

\item {} 
\sphinxAtStartPar
spring = style name of this fix command

\item {} 
\sphinxAtStartPar
keyword = \sphinxstyleemphasis{tether} or \sphinxstyleemphasis{couple}
\begin{sphinxalltt}
\sphinxstyleemphasis{tether} values = K x y z R0
  K = spring constant (force/distance units)
  x,y,z = point to which spring is tethered
  R0 = equilibrium distance from tether point (distance units)
\sphinxstyleemphasis{couple} values = group\sphinxhyphen{}ID2 K x y z R0
  group\sphinxhyphen{}ID2 = 2nd group to couple to fix group with a spring
  K = spring constant (force/distance units)
  x,y,z = direction of spring
  R0 = equilibrium distance of spring (distance units)
\end{sphinxalltt}

\end{itemize}


\subsection{Examples}
\label{\detokenize{fix_spring:examples}}
\begin{sphinxVerbatim}[commandchars=\\\{\}]
\PYG{k}{fix }\PYG{n+nv+nvIdentifier}{pull}\PYG{+w}{ }\PYG{n+nv+nvIdentifier}{ligand}\PYG{+w}{ }\PYG{n}{spring}\PYG{+w}{ }\PYG{n}{tether}\PYG{+w}{ }\PYG{l+m}{50.0}\PYG{+w}{ }\PYG{l+m}{0.0}\PYG{+w}{ }\PYG{l+m}{0.0}\PYG{+w}{ }\PYG{l+m}{0.0}\PYG{+w}{ }\PYG{l+m}{0.0}
\PYG{k}{fix }\PYG{n+nv+nvIdentifier}{pull}\PYG{+w}{ }\PYG{n+nv+nvIdentifier}{ligand}\PYG{+w}{ }\PYG{n}{spring}\PYG{+w}{ }\PYG{n}{tether}\PYG{+w}{ }\PYG{l+m}{50.0}\PYG{+w}{ }\PYG{l+m}{0.0}\PYG{+w}{ }\PYG{l+m}{0.0}\PYG{+w}{ }\PYG{l+m}{0.0}\PYG{+w}{ }\PYG{l+m}{5.0}
\PYG{k}{fix }\PYG{n+nv+nvIdentifier}{pull}\PYG{+w}{ }\PYG{n+nv+nvIdentifier}{ligand}\PYG{+w}{ }\PYG{n}{spring}\PYG{+w}{ }\PYG{n}{tether}\PYG{+w}{ }\PYG{l+m}{50.0}\PYG{+w}{ }\PYG{n}{NULL}\PYG{+w}{ }\PYG{n}{NULL}\PYG{+w}{ }\PYG{l+m}{2.0}\PYG{+w}{ }\PYG{l+m}{3.0}
\PYG{k}{fix }\PYG{n+nv+nvIdentifier}{5}\PYG{+w}{ }\PYG{n+nv+nvIdentifier}{bilayer1}\PYG{+w}{ }\PYG{n}{spring}\PYG{+w}{ }\PYG{n}{couple}\PYG{+w}{ }\PYG{n}{bilayer2}\PYG{+w}{ }\PYG{l+m}{100.0}\PYG{+w}{ }\PYG{n}{NULL}\PYG{+w}{ }\PYG{n}{NULL}\PYG{+w}{ }\PYG{l+m}{10.0}\PYG{+w}{ }\PYG{l+m}{0.0}
\PYG{k}{fix }\PYG{n+nv+nvIdentifier}{longitudinal}\PYG{+w}{ }\PYG{n+nv+nvIdentifier}{pore}\PYG{+w}{ }\PYG{n}{spring}\PYG{+w}{ }\PYG{n}{couple}\PYG{+w}{ }\PYG{n}{ion}\PYG{+w}{ }\PYG{l+m}{100.0}\PYG{+w}{ }\PYG{n}{NULL}\PYG{+w}{ }\PYG{n}{NULL}\PYG{+w}{ }\PYG{o}{\PYGZhy{}}\PYG{l+m}{20.0}\PYG{+w}{ }\PYG{l+m}{0.0}
\PYG{k}{fix }\PYG{n+nv+nvIdentifier}{radial}\PYG{+w}{ }\PYG{n+nv+nvIdentifier}{pore}\PYG{+w}{ }\PYG{n}{spring}\PYG{+w}{ }\PYG{n}{couple}\PYG{+w}{ }\PYG{n}{ion}\PYG{+w}{ }\PYG{l+m}{100.0}\PYG{+w}{ }\PYG{l+m}{0.0}\PYG{+w}{ }\PYG{l+m}{0.0}\PYG{+w}{ }\PYG{n}{NULL}\PYG{+w}{ }\PYG{l+m}{5.0}
\end{sphinxVerbatim}


\subsection{Description}
\label{\detokenize{fix_spring:description}}
\sphinxAtStartPar
Apply a spring force to a group of atoms or between two groups of
atoms.  This is useful for applying an umbrella force to a small
molecule or lightly tethering a large group of atoms (e.g. all the
solvent or a large molecule) to the center of the simulation box so
that it does not wander away over the course of a long simulation.  It
can also be used to hold the centers of mass of two groups of atoms at
a given distance or orientation with respect to each other.

\sphinxAtStartPar
The \sphinxstyleemphasis{tether} style attaches a spring between a fixed point \sphinxstyleemphasis{x,y,z} and
the center of mass of the fix group of atoms.  The equilibrium
position of the spring is R0.  At each timestep the distance R from
the center of mass of the group of atoms to the tethering point is
computed, taking account of wrap\sphinxhyphen{}around in a periodic simulation box.
A restoring force of magnitude K (R \sphinxhyphen{} R0) Mi / M is applied to each
atom in the group where \sphinxstyleemphasis{K} is the spring constant, Mi is the mass of
the atom, and M is the total mass of all atoms in the group.  Note
that \sphinxstyleemphasis{K} thus represents the spring constant for the total force on
the group of atoms, not for a spring applied to each atom.

\sphinxAtStartPar
The \sphinxstyleemphasis{couple} style links two groups of atoms together.  The first
group is the fix group; the second is specified by group\sphinxhyphen{}ID2.  The
groups are coupled together by a spring that is at equilibrium when
the two groups are displaced by a vector \sphinxstyleemphasis{x,y,z} with respect to each
other and at a distance R0 from that displacement.  Note that \sphinxstyleemphasis{x,y,z}
is the equilibrium displacement of group\sphinxhyphen{}ID2 relative to the fix
group.  Thus (1,1,0) is a different spring than (\sphinxhyphen{}1,\sphinxhyphen{}1,0).  When the
relative positions and distance between the two groups are not in
equilibrium, the same spring force described above is applied to atoms
in each of the two groups.

\sphinxAtStartPar
For both the \sphinxstyleemphasis{tether} and \sphinxstyleemphasis{couple} styles, any of the x,y,z values can
be specified as NULL which means do not include that dimension in the
distance calculation or force application.

\sphinxAtStartPar
The first example above pulls the ligand towards the point (0,0,0).
The second example holds the ligand near the surface of a sphere of
radius 5 around the point (0,0,0).  The third example holds the ligand
a distance 3 away from the z=2 plane (on either side).

\sphinxAtStartPar
The fourth example holds 2 bilayers a distance 10 apart in z.  For the
last two examples, imagine a pore (a slab of atoms with a cylindrical
hole cut out) oriented with the pore axis along z, and an ion moving
within the pore.  The fifth example holds the ion a distance of \sphinxhyphen{}20
below the z = 0 center plane of the pore (umbrella sampling).  The
last example holds the ion a distance 5 away from the pore axis
(assuming the center\sphinxhyphen{}of\sphinxhyphen{}mass of the pore in x,y is the pore axis).

\begin{sphinxadmonition}{note}{Note:}
\sphinxAtStartPar
The center of mass of a group of atoms is calculated in
“unwrapped” coordinates using atom image flags, which means that the
group can straddle a periodic boundary.  See the {\hyperref[\detokenize{dump::doc}]{\sphinxcrossref{\DUrole{doc}{dump}}}} doc
page for a discussion of unwrapped coordinates.  It also means that a
spring connecting two groups or a group and the tether point can cross
a periodic boundary and its length be calculated correctly.
\end{sphinxadmonition}


\subsection{Restart, fix\_modify, output, run start/stop, minimize info}
\label{\detokenize{fix_spring:restart-fix-modify-output-run-start-stop-minimize-info}}
\sphinxAtStartPar
No information about this fix is written to {\hyperref[\detokenize{restart::doc}]{\sphinxcrossref{\DUrole{doc}{binary restart files}}}}.

\sphinxAtStartPar
The {\hyperref[\detokenize{fix_modify::doc}]{\sphinxcrossref{\DUrole{doc}{fix\_modify}}}} \sphinxstyleemphasis{energy} option is supported by
this fix to add the energy stored in the spring to the global
potential energy of the system as part of {\hyperref[\detokenize{thermo_style::doc}]{\sphinxcrossref{\DUrole{doc}{thermodynamic output}}}}. The default setting for this fix is {\hyperref[\detokenize{fix_modify::doc}]{\sphinxcrossref{\DUrole{doc}{fix\_modify
energy no}}}}.

\sphinxAtStartPar
The {\hyperref[\detokenize{fix_modify::doc}]{\sphinxcrossref{\DUrole{doc}{fix\_modify}}}} \sphinxstyleemphasis{respa} option is supported by this
fix. This allows to set at which level of the {\hyperref[\detokenize{run_style::doc}]{\sphinxcrossref{\DUrole{doc}{r\sphinxhyphen{}RESPA}}}}
integrator the fix is adding its forces. Default is the outermost level.

\sphinxAtStartPar
This fix computes a global scalar which can be accessed by various
{\hyperref[\detokenize{Howto_output::doc}]{\sphinxcrossref{\DUrole{doc}{output commands}}}}.  The scalar is the spring energy
= 0.5 * K * r\textasciicircum{}2.

\sphinxAtStartPar
This fix also computes global 4\sphinxhyphen{}vector which can be accessed by
various {\hyperref[\detokenize{Howto_output::doc}]{\sphinxcrossref{\DUrole{doc}{output commands}}}}.  The first 3 quantities
in the vector are xyz components of the total force added to the group
of atoms by the spring.  In the case of the \sphinxstyleemphasis{couple} style, it is the
force on the fix group (group\sphinxhyphen{}ID) or the negative of the force on the
second group (group\sphinxhyphen{}ID2).  The fourth quantity in the vector is the
magnitude of the force added by the spring, as a positive value if
(r\sphinxhyphen{}R0) \textgreater{} 0 and a negative value if (r\sphinxhyphen{}R0) \textless{} 0.  This sign convention
can be useful when using the spring force to compute a potential of
mean force (PMF).

\sphinxAtStartPar
The scalar and vector values calculated by this fix are “extensive”.

\sphinxAtStartPar
No parameter of this fix can be used with the \sphinxstyleemphasis{start/stop} keywords of
the {\hyperref[\detokenize{run::doc}]{\sphinxcrossref{\DUrole{doc}{run}}}} command.

\sphinxAtStartPar
The forces due to this fix are imposed during an energy minimization,
invoked by the {\hyperref[\detokenize{minimize::doc}]{\sphinxcrossref{\DUrole{doc}{minimize}}}} command.

\begin{sphinxadmonition}{note}{Note:}
\sphinxAtStartPar
If you want the spring energy to be included in the total
potential energy of the system (the quantity being minimized), you
MUST enable the {\hyperref[\detokenize{fix_modify::doc}]{\sphinxcrossref{\DUrole{doc}{fix\_modify}}}} \sphinxstyleemphasis{energy} option for this
fix.
\end{sphinxadmonition}


\subsection{Restrictions}
\label{\detokenize{fix_spring:restrictions}}\begin{quote}

\sphinxAtStartPar
none
\end{quote}


\subsection{Related commands}
\label{\detokenize{fix_spring:related-commands}}
\sphinxAtStartPar
{\hyperref[\detokenize{fix_drag::doc}]{\sphinxcrossref{\DUrole{doc}{fix drag}}}}, {\hyperref[\detokenize{fix_spring_self::doc}]{\sphinxcrossref{\DUrole{doc}{fix spring/self}}}},
{\hyperref[\detokenize{fix_spring_rg::doc}]{\sphinxcrossref{\DUrole{doc}{fix spring/rg}}}}, {\hyperref[\detokenize{fix_smd::doc}]{\sphinxcrossref{\DUrole{doc}{fix smd}}}}


\subsection{Default}
\label{\detokenize{fix_spring:default}}
\sphinxAtStartPar
none

\sphinxstepscope

\index{fix spring/chunk@\spxentry{fix spring/chunk}}\ignorespaces 

\section{fix spring/chunk command}
\label{\detokenize{fix_spring_chunk:fix-spring-chunk-command}}\label{\detokenize{fix_spring_chunk:index-0}}\label{\detokenize{fix_spring_chunk::doc}}

\subsection{Syntax}
\label{\detokenize{fix_spring_chunk:syntax}}
\begin{sphinxVerbatim}[commandchars=\\\{\}]
\PYG{k}{fix }\PYG{n+nv+nvIdentifier}{ID}\PYG{+w}{ }\PYG{n+nv+nvIdentifier}{group\PYGZhy{}ID}\PYG{+w}{ }\PYG{n}{spring}\PYG{o}{/}\PYG{n}{chunk}\PYG{+w}{ }\PYG{n}{K}\PYG{+w}{ }\PYG{n}{chunkID}\PYG{+w}{ }\PYG{n}{comID}
\end{sphinxVerbatim}
\begin{itemize}
\item {} 
\sphinxAtStartPar
ID, group\sphinxhyphen{}ID are documented in {\hyperref[\detokenize{fix::doc}]{\sphinxcrossref{\DUrole{doc}{fix}}}} command

\item {} 
\sphinxAtStartPar
spring/chunk = style name of this fix command

\item {} 
\sphinxAtStartPar
K = spring constant for each chunk (force/distance units)

\item {} 
\sphinxAtStartPar
chunkID = ID of {\hyperref[\detokenize{compute_chunk_atom::doc}]{\sphinxcrossref{\DUrole{doc}{compute chunk/atom}}}} command

\item {} 
\sphinxAtStartPar
comID = ID of {\hyperref[\detokenize{compute_com_chunk::doc}]{\sphinxcrossref{\DUrole{doc}{compute com/chunk}}}} command

\end{itemize}


\subsection{Examples}
\label{\detokenize{fix_spring_chunk:examples}}
\begin{sphinxVerbatim}[commandchars=\\\{\}]
\PYG{k}{fix }\PYG{n+nv+nvIdentifier}{restrain}\PYG{+w}{ }\PYG{n+nv+nvIdentifier}{all}\PYG{+w}{ }\PYG{n}{spring}\PYG{o}{/}\PYG{n}{chunk}\PYG{+w}{ }\PYG{l+m}{100}\PYG{+w}{ }\PYG{n}{chunkID}\PYG{+w}{ }\PYG{n}{comID}
\end{sphinxVerbatim}


\subsection{Description}
\label{\detokenize{fix_spring_chunk:description}}
\sphinxAtStartPar
Apply a spring force to the center\sphinxhyphen{}of\sphinxhyphen{}mass (COM) of chunks of atoms as
defined by the {\hyperref[\detokenize{compute_chunk_atom::doc}]{\sphinxcrossref{\DUrole{doc}{compute chunk/atom}}}} command.
Chunks can be molecules or spatial bins or other groupings of atoms.
This is a way of tethering each chunk to its initial COM coordinates.

\sphinxAtStartPar
The \sphinxstyleemphasis{chunkID} is the ID of a compute chunk/atom command defined in the
input script.  It is used to define the chunks.  The \sphinxstyleemphasis{comID} is the ID
of a compute com/chunk command defined in the input script.  It is
used to compute the COMs of each chunk.

\sphinxAtStartPar
At the beginning of the first {\hyperref[\detokenize{run::doc}]{\sphinxcrossref{\DUrole{doc}{run}}}} or
{\hyperref[\detokenize{minimize::doc}]{\sphinxcrossref{\DUrole{doc}{minimize}}}} command after this fix is defined, the
initial COM of each chunk is calculated and stored as R0m, where M is
the chunk number.  Thereafter, at every timestep (or minimization
iteration), the current COM of each chunk is calculated as Rm.  A
restoring force of magnitude K (Rm \sphinxhyphen{} R0m) Mi / Mm is applied to each
atom in each chunk where \sphinxstyleemphasis{K} is the specified spring constant, Mi is
the mass of the atom, and Mm is the total mass of all atoms in the
chunk.  Note that \sphinxstyleemphasis{K} thus represents the spring constant for the
total force on each chunk of atoms, not for a spring applied to each
atom.


\subsection{Restart, fix\_modify, output, run start/stop, minimize info}
\label{\detokenize{fix_spring_chunk:restart-fix-modify-output-run-start-stop-minimize-info}}
\sphinxAtStartPar
This fix writes the locations of the initial per\sphinxhyphen{}chunk center of mass
coordinates to {\hyperref[\detokenize{restart::doc}]{\sphinxcrossref{\DUrole{doc}{binary restart files}}}}.  See the
{\hyperref[\detokenize{read_restart::doc}]{\sphinxcrossref{\DUrole{doc}{read\_restart}}}} command for info on how to
re\sphinxhyphen{}specify a fix in an input script that reads a restart file, so that
the fix continues in an uninterrupted fashion.  Since this fix depends
on an instance of {\hyperref[\detokenize{compute_chunk_atom::doc}]{\sphinxcrossref{\DUrole{doc}{compute chunk/atom}}}}
it will check when reading the restart if the chunk still exists and
will define the same number of chunks. The restart data is only applied
when the number of chunks matches. Otherwise the center of mass
coordinates are recomputed.

\sphinxAtStartPar
The {\hyperref[\detokenize{fix_modify::doc}]{\sphinxcrossref{\DUrole{doc}{fix\_modify}}}} \sphinxstyleemphasis{energy} option is supported by
this fix to add the energy stored in all the springs to the global
potential energy of the system as part of {\hyperref[\detokenize{thermo_style::doc}]{\sphinxcrossref{\DUrole{doc}{thermodynamic output}}}}.  The default setting for this fix is {\hyperref[\detokenize{fix_modify::doc}]{\sphinxcrossref{\DUrole{doc}{fix\_modify
energy no}}}}.

\sphinxAtStartPar
The {\hyperref[\detokenize{fix_modify::doc}]{\sphinxcrossref{\DUrole{doc}{fix\_modify}}}} \sphinxstyleemphasis{respa} option is supported by this
fix. This allows to set at which level of the {\hyperref[\detokenize{run_style::doc}]{\sphinxcrossref{\DUrole{doc}{r\sphinxhyphen{}RESPA}}}}
integrator the fix is adding its forces. Default is the outermost level.

\sphinxAtStartPar
This fix computes a global scalar which can be accessed by various
{\hyperref[\detokenize{Howto_output::doc}]{\sphinxcrossref{\DUrole{doc}{output commands}}}}.  The scalar is the energy of all
the springs, i.e. 0.5 * K * r\textasciicircum{}2 per\sphinxhyphen{}spring.

\sphinxAtStartPar
The scalar value calculated by this fix is “extensive”.

\sphinxAtStartPar
No parameter of this fix can be used with the \sphinxstyleemphasis{start/stop} keywords of
the {\hyperref[\detokenize{run::doc}]{\sphinxcrossref{\DUrole{doc}{run}}}} command.

\sphinxAtStartPar
The forces due to this fix are imposed during an energy minimization,
invoked by the {\hyperref[\detokenize{minimize::doc}]{\sphinxcrossref{\DUrole{doc}{minimize}}}} command.

\begin{sphinxadmonition}{note}{Note:}
\sphinxAtStartPar
If you want the spring energies to be included in the total
potential energy of the system (the quantity being minimized), you
MUST enable the {\hyperref[\detokenize{fix_modify::doc}]{\sphinxcrossref{\DUrole{doc}{fix\_modify}}}} \sphinxstyleemphasis{energy} option for this
fix.
\end{sphinxadmonition}


\subsection{Restrictions}
\label{\detokenize{fix_spring_chunk:restrictions}}\begin{quote}

\sphinxAtStartPar
none
\end{quote}


\subsection{Related commands}
\label{\detokenize{fix_spring_chunk:related-commands}}
\sphinxAtStartPar
{\hyperref[\detokenize{fix_spring::doc}]{\sphinxcrossref{\DUrole{doc}{fix spring}}}}, {\hyperref[\detokenize{fix_spring_self::doc}]{\sphinxcrossref{\DUrole{doc}{fix spring/self}}}},
{\hyperref[\detokenize{fix_spring_rg::doc}]{\sphinxcrossref{\DUrole{doc}{fix spring/rg}}}}


\subsection{Default}
\label{\detokenize{fix_spring_chunk:default}}
\sphinxAtStartPar
none

\sphinxstepscope

\index{fix spring/rg@\spxentry{fix spring/rg}}\ignorespaces 

\section{fix spring/rg command}
\label{\detokenize{fix_spring_rg:fix-spring-rg-command}}\label{\detokenize{fix_spring_rg:index-0}}\label{\detokenize{fix_spring_rg::doc}}

\subsection{Syntax}
\label{\detokenize{fix_spring_rg:syntax}}
\begin{sphinxVerbatim}[commandchars=\\\{\}]
\PYG{k}{fix }\PYG{n+nv+nvIdentifier}{ID}\PYG{+w}{ }\PYG{n+nv+nvIdentifier}{group\PYGZhy{}ID}\PYG{+w}{ }\PYG{n}{spring}\PYG{o}{/}\PYG{n}{rg}\PYG{+w}{ }\PYG{n}{K}\PYG{+w}{ }\PYG{n}{RG0}
\end{sphinxVerbatim}
\begin{itemize}
\item {} 
\sphinxAtStartPar
ID, group\sphinxhyphen{}ID are documented in {\hyperref[\detokenize{fix::doc}]{\sphinxcrossref{\DUrole{doc}{fix}}}} command

\item {} 
\sphinxAtStartPar
spring/rg = style name of this fix command

\item {} 
\sphinxAtStartPar
K = harmonic force constant (force/distance units)

\item {} 
\sphinxAtStartPar
RG0 = target radius of gyration to constrain to (distance units)

\end{itemize}

\begin{sphinxVerbatim}[commandchars=\\\{\}]
if RG0 = NULL, use the current RG as the target value
\end{sphinxVerbatim}


\subsection{Examples}
\label{\detokenize{fix_spring_rg:examples}}
\begin{sphinxVerbatim}[commandchars=\\\{\}]
\PYG{k}{fix }\PYG{n+nv+nvIdentifier}{1}\PYG{+w}{ }\PYG{n+nv+nvIdentifier}{protein}\PYG{+w}{ }\PYG{n}{spring}\PYG{o}{/}\PYG{n}{rg}\PYG{+w}{ }\PYG{l+m}{5.0}\PYG{+w}{ }\PYG{l+m}{10.0}
\PYG{k}{fix }\PYG{n+nv+nvIdentifier}{2}\PYG{+w}{ }\PYG{n+nv+nvIdentifier}{micelle}\PYG{+w}{ }\PYG{n}{spring}\PYG{o}{/}\PYG{n}{rg}\PYG{+w}{ }\PYG{l+m}{5.0}\PYG{+w}{ }\PYG{n}{NULL}
\end{sphinxVerbatim}


\subsection{Description}
\label{\detokenize{fix_spring_rg:description}}
\sphinxAtStartPar
Apply a harmonic restraining force to atoms in the group to affect
their central moment about the center of mass (radius of gyration).
This fix is useful to encourage a protein or polymer to fold/unfold
and also when sampling along the radius of gyration as a reaction
coordinate (i.e. for protein folding).

\sphinxAtStartPar
The radius of gyration is defined as RG in the first formula.  The
energy of the constraint and associated force on each atom is given by
the second and third formulas, when the group is at a different RG
than the target value RG0.
\begin{equation*}
\begin{split}{R_G}^2 & = \frac{1}{M}\sum_{i}^{N}{m_{i}\left( x_{i} -
\frac{1}{M}\sum_{j}^{N}{m_{j}x_{j}} \right)^{2}} \\
E & = K\left( R_G - R_{G0} \right)^{2} \\
F_{i} & = 2K\frac{m_{i}}{M}\left( 1-\frac{R_{G0}}{R_G}
\right)\left( x_{i} - \frac{1}{M}\sum_{j}^{N}{m_{j}x_{j}} \right)\end{split}
\end{equation*}
\sphinxAtStartPar
The (\(x_i\) \sphinxhyphen{} center\sphinxhyphen{}of\sphinxhyphen{}mass) term is computed taking into account
periodic boundary conditions, \(m_i\) is the mass of the atom, and
\sphinxstyleemphasis{M} is the mass of the entire group.  Note that K is thus a force constant
for the aggregate force on the group of atoms, not a per\sphinxhyphen{}atom force.

\sphinxAtStartPar
If \(R_{G0}\) is specified as NULL, then the RG of the group is computed at
the time the fix is specified, and that value is used as the target.


\subsection{Restart, fix\_modify, output, run start/stop, minimize info}
\label{\detokenize{fix_spring_rg:restart-fix-modify-output-run-start-stop-minimize-info}}
\sphinxAtStartPar
This fix writes the currently used reference RG (\(R_{G0}\)) to
{\hyperref[\detokenize{restart::doc}]{\sphinxcrossref{\DUrole{doc}{binary restart files}}}}.  See the {\hyperref[\detokenize{read_restart::doc}]{\sphinxcrossref{\DUrole{doc}{read\_restart}}}} command for info on how to re\sphinxhyphen{}specify a fix in an input
script that reads a restart file, so that the fix continues in an
uninterrupted fashion.

\sphinxAtStartPar
None of the {\hyperref[\detokenize{fix_modify::doc}]{\sphinxcrossref{\DUrole{doc}{fix\_modify}}}} options
are relevant to this fix.

\sphinxAtStartPar
This fix computes a global scalar which can be accessed by various
{\hyperref[\detokenize{Howto_output::doc}]{\sphinxcrossref{\DUrole{doc}{output commands}}}}.  The scalar is the reference
radius of gyration \(R_{G0}\) used by the fix.  energy change due to
this fix.  The scalar value calculated by this fix is “intensive”.

\sphinxAtStartPar
No parameter of this fix can be used with the \sphinxstyleemphasis{start/stop} keywords of
the {\hyperref[\detokenize{run::doc}]{\sphinxcrossref{\DUrole{doc}{run}}}} command.  This fix is not invoked during {\hyperref[\detokenize{minimize::doc}]{\sphinxcrossref{\DUrole{doc}{energy minimization}}}}.

\sphinxAtStartPar
The {\hyperref[\detokenize{fix_modify::doc}]{\sphinxcrossref{\DUrole{doc}{fix\_modify}}}} \sphinxstyleemphasis{respa} option is supported by this
fix. This allows to set at which level of the {\hyperref[\detokenize{run_style::doc}]{\sphinxcrossref{\DUrole{doc}{r\sphinxhyphen{}RESPA}}}}
integrator the fix is adding its forces. Default is the outermost level.


\subsection{Restrictions}
\label{\detokenize{fix_spring_rg:restrictions}}
\sphinxAtStartPar
This fix is part of the EXTRA\sphinxhyphen{}FIX package.  It is only enabled if LAMMPS
was built with that package.  See the {\hyperref[\detokenize{Build_package::doc}]{\sphinxcrossref{\DUrole{doc}{Build package}}}} page for more info.


\subsection{Related commands}
\label{\detokenize{fix_spring_rg:related-commands}}
\sphinxAtStartPar
{\hyperref[\detokenize{fix_spring::doc}]{\sphinxcrossref{\DUrole{doc}{fix spring}}}}, {\hyperref[\detokenize{fix_spring_self::doc}]{\sphinxcrossref{\DUrole{doc}{fix spring/self}}}}
{\hyperref[\detokenize{fix_drag::doc}]{\sphinxcrossref{\DUrole{doc}{fix drag}}}}, {\hyperref[\detokenize{fix_smd::doc}]{\sphinxcrossref{\DUrole{doc}{fix smd}}}}


\subsection{Default}
\label{\detokenize{fix_spring_rg:default}}
\sphinxAtStartPar
none

\sphinxstepscope

\index{fix spring/self@\spxentry{fix spring/self}}\ignorespaces 
\index{fix spring/self/kk@\spxentry{fix spring/self/kk}}\ignorespaces 

\section{fix spring/self command}
\label{\detokenize{fix_spring_self:fix-spring-self-command}}\label{\detokenize{fix_spring_self:index-1}}\label{\detokenize{fix_spring_self:index-0}}\label{\detokenize{fix_spring_self::doc}}

\subsection{Syntax}
\label{\detokenize{fix_spring_self:syntax}}
\begin{sphinxVerbatim}[commandchars=\\\{\}]
\PYG{k}{fix }\PYG{n+nv+nvIdentifier}{ID}\PYG{+w}{ }\PYG{n+nv+nvIdentifier}{group\PYGZhy{}ID}\PYG{+w}{ }\PYG{n}{spring}\PYG{o}{/}\PYG{n}{self}\PYG{+w}{ }\PYG{n}{K}\PYG{+w}{ }\PYG{n}{dir}
\end{sphinxVerbatim}
\begin{itemize}
\item {} 
\sphinxAtStartPar
ID, group\sphinxhyphen{}ID are documented in {\hyperref[\detokenize{fix::doc}]{\sphinxcrossref{\DUrole{doc}{fix}}}} command

\item {} 
\sphinxAtStartPar
spring/self = style name of this fix command

\item {} 
\sphinxAtStartPar
K = spring constant (force/distance units)

\item {} 
\sphinxAtStartPar
dir = xyz, xy, xz, yz, x, y, or z (optional, default: xyz)

\end{itemize}


\subsection{Examples}
\label{\detokenize{fix_spring_self:examples}}
\begin{sphinxVerbatim}[commandchars=\\\{\}]
\PYG{k}{fix }\PYG{n+nv+nvIdentifier}{tether}\PYG{+w}{ }\PYG{n+nv+nvIdentifier}{boundary\PYGZhy{}atoms}\PYG{+w}{ }\PYG{n}{spring}\PYG{o}{/}\PYG{n}{self}\PYG{+w}{ }\PYG{l+m}{10.0}
\PYG{k}{fix }\PYG{n+nv+nvIdentifier}{zrest}\PYG{+w}{  }\PYG{n+nv+nvIdentifier}{move}\PYG{+w}{ }\PYG{n}{spring}\PYG{o}{/}\PYG{n}{self}\PYG{+w}{ }\PYG{l+m}{10.0}\PYG{+w}{ }\PYG{n}{z}
\end{sphinxVerbatim}


\subsection{Description}
\label{\detokenize{fix_spring_self:description}}
\sphinxAtStartPar
Apply a spring force independently to each atom in the group to tether
it to its initial position.  The initial position for each atom is its
location at the time the fix command was issued.  At each timestep,
the magnitude of the force on each atom is \sphinxhyphen{}Kr, where r is the
displacement of the atom from its current position to its initial
position.  The distance r correctly takes into account any crossings
of periodic boundary by the atom since it was in its initial
position.

\sphinxAtStartPar
With the (optional) dir flag, one can select in which direction the
spring force is applied. By default, the restraint is applied in all
directions, but it can be limited to the xy\sphinxhyphen{}, xz\sphinxhyphen{}, yz\sphinxhyphen{}plane and the
x\sphinxhyphen{}, y\sphinxhyphen{}, or z\sphinxhyphen{}direction, thus restraining the atoms to a line or a
plane, respectively.


\subsection{Restart, fix\_modify, output, run start/stop, minimize info}
\label{\detokenize{fix_spring_self:restart-fix-modify-output-run-start-stop-minimize-info}}
\sphinxAtStartPar
This fix writes the original coordinates of tethered atoms to
{\hyperref[\detokenize{restart::doc}]{\sphinxcrossref{\DUrole{doc}{binary restart files}}}}, so that the spring effect will
be the same in a restarted simulation.  See the {\hyperref[\detokenize{read_restart::doc}]{\sphinxcrossref{\DUrole{doc}{read\_restart}}}} command for info on how to re\sphinxhyphen{}specify a fix in an
input script that reads a restart file, so that the operation of the
fix continues in an uninterrupted fashion.

\sphinxAtStartPar
The {\hyperref[\detokenize{fix_modify::doc}]{\sphinxcrossref{\DUrole{doc}{fix\_modify}}}} \sphinxstyleemphasis{energy} option is supported by
this fix to add the energy stored in the per\sphinxhyphen{}atom springs to the
global potential energy of the system as part of {\hyperref[\detokenize{thermo_style::doc}]{\sphinxcrossref{\DUrole{doc}{thermodynamic
output}}}}.  The default setting for this fix is
{\hyperref[\detokenize{fix_modify::doc}]{\sphinxcrossref{\DUrole{doc}{fix\_modify energy no}}}}.

\sphinxAtStartPar
The {\hyperref[\detokenize{fix_modify::doc}]{\sphinxcrossref{\DUrole{doc}{fix\_modify}}}} \sphinxstyleemphasis{respa} option is supported by
this fix. This allows to set at which level of the {\hyperref[\detokenize{run_style::doc}]{\sphinxcrossref{\DUrole{doc}{r\sphinxhyphen{}RESPA}}}}
integrator the fix is adding its forces. Default is the outermost level.

\sphinxAtStartPar
This fix computes a global scalar which can be accessed by various
{\hyperref[\detokenize{Howto_output::doc}]{\sphinxcrossref{\DUrole{doc}{output commands}}}}.  The scalar is an energy which is
the sum of the spring energy for each atom, where the per\sphinxhyphen{}atom energy
is 0.5 * K * r\textasciicircum{}2.  The scalar value calculated by this fix is
“extensive”.

\sphinxAtStartPar
No parameter of this fix can be used with the \sphinxstyleemphasis{start/stop} keywords of
the {\hyperref[\detokenize{run::doc}]{\sphinxcrossref{\DUrole{doc}{run}}}} command.

\sphinxAtStartPar
The forces due to this fix are imposed during an energy minimization,
invoked by the {\hyperref[\detokenize{minimize::doc}]{\sphinxcrossref{\DUrole{doc}{minimize}}}} command.

\begin{sphinxadmonition}{note}{Note:}
\sphinxAtStartPar
If you want the per\sphinxhyphen{}atom spring energy to be included in the
total potential energy of the system (the quantity being minimized),
you MUST enable the {\hyperref[\detokenize{fix_modify::doc}]{\sphinxcrossref{\DUrole{doc}{fix\_modify}}}} \sphinxstyleemphasis{energy} option for
this fix.
\end{sphinxadmonition}


\bigskip\hrule\bigskip


\sphinxAtStartPar
Styles with a \sphinxstyleemphasis{gpu}, \sphinxstyleemphasis{intel}, \sphinxstyleemphasis{kk}, \sphinxstyleemphasis{omp}, or \sphinxstyleemphasis{opt} suffix are
functionally the same as the corresponding style without the suffix.
They have been optimized to run faster, depending on your available
hardware, as discussed on the {\hyperref[\detokenize{Speed_packages::doc}]{\sphinxcrossref{\DUrole{doc}{Accelerator packages}}}}
page.  The accelerated styles take the same arguments and should
produce the same results, except for round\sphinxhyphen{}off and precision issues.

\sphinxAtStartPar
These accelerated styles are part of the GPU, INTEL, KOKKOS,
OPENMP, and OPT packages, respectively.  They are only enabled if
LAMMPS was built with those packages.  See the {\hyperref[\detokenize{Build_package::doc}]{\sphinxcrossref{\DUrole{doc}{Build package}}}} page for more info.

\sphinxAtStartPar
You can specify the accelerated styles explicitly in your input script
by including their suffix, or you can use the {\hyperref[\detokenize{Run_options::doc}]{\sphinxcrossref{\DUrole{doc}{\sphinxhyphen{}suffix command\sphinxhyphen{}line switch}}}} when you invoke LAMMPS, or you can use the
{\hyperref[\detokenize{suffix::doc}]{\sphinxcrossref{\DUrole{doc}{suffix}}}} command in your input script.

\sphinxAtStartPar
See the {\hyperref[\detokenize{Speed_packages::doc}]{\sphinxcrossref{\DUrole{doc}{Accelerator packages}}}} page for more
instructions on how to use the accelerated styles effectively.


\bigskip\hrule\bigskip



\subsection{Restrictions}
\label{\detokenize{fix_spring_self:restrictions}}\begin{quote}

\sphinxAtStartPar
none
\end{quote}


\subsection{Related commands}
\label{\detokenize{fix_spring_self:related-commands}}
\sphinxAtStartPar
{\hyperref[\detokenize{fix_drag::doc}]{\sphinxcrossref{\DUrole{doc}{fix drag}}}}, {\hyperref[\detokenize{fix_spring::doc}]{\sphinxcrossref{\DUrole{doc}{fix spring}}}},
{\hyperref[\detokenize{fix_smd::doc}]{\sphinxcrossref{\DUrole{doc}{fix smd}}}}, {\hyperref[\detokenize{fix_spring_rg::doc}]{\sphinxcrossref{\DUrole{doc}{fix spring/rg}}}}


\subsection{Default}
\label{\detokenize{fix_spring_self:default}}
\sphinxAtStartPar
none

\sphinxstepscope

\index{fix srd@\spxentry{fix srd}}\ignorespaces 

\section{fix srd command}
\label{\detokenize{fix_srd:fix-srd-command}}\label{\detokenize{fix_srd:index-0}}\label{\detokenize{fix_srd::doc}}

\subsection{Syntax}
\label{\detokenize{fix_srd:syntax}}
\begin{sphinxVerbatim}[commandchars=\\\{\}]
\PYG{k}{fix }\PYG{n+nv+nvIdentifier}{ID}\PYG{+w}{ }\PYG{n+nv+nvIdentifier}{group\PYGZhy{}ID}\PYG{+w}{ }\PYG{n}{srd}\PYG{+w}{ }\PYG{n}{N}\PYG{+w}{ }\PYG{n}{groupbig}\PYG{o}{\PYGZhy{}}\PYG{n}{ID}\PYG{+w}{ }\PYG{n}{Tsrd}\PYG{+w}{ }\PYG{n}{hgrid}\PYG{+w}{ }\PYG{n}{seed}\PYG{+w}{ }\PYG{n}{keyword}\PYG{+w}{ }\PYG{n}{value}\PYG{+w}{ }\PYG{n}{...}
\end{sphinxVerbatim}
\begin{itemize}
\item {} 
\sphinxAtStartPar
ID, group\sphinxhyphen{}ID are documented in {\hyperref[\detokenize{fix::doc}]{\sphinxcrossref{\DUrole{doc}{fix}}}} command

\item {} 
\sphinxAtStartPar
srd = style name of this fix command

\item {} 
\sphinxAtStartPar
N = reset SRD particle velocities every this many timesteps

\item {} 
\sphinxAtStartPar
groupbig\sphinxhyphen{}ID = ID of group of large particles that SRDs interact with

\item {} 
\sphinxAtStartPar
Tsrd = temperature of SRD particles (temperature units)

\item {} 
\sphinxAtStartPar
hgrid = grid spacing for SRD grouping (distance units)

\item {} 
\sphinxAtStartPar
seed = random \# seed (positive integer)

\item {} 
\sphinxAtStartPar
zero or more keyword/value pairs may be appended

\item {} 
\sphinxAtStartPar
keyword = \sphinxstyleemphasis{lamda} or \sphinxstyleemphasis{collision} or \sphinxstyleemphasis{overlap} or \sphinxstyleemphasis{inside} or \sphinxstyleemphasis{exact} or \sphinxstyleemphasis{radius} or \sphinxstyleemphasis{bounce} or \sphinxstyleemphasis{search} or \sphinxstyleemphasis{cubic} or \sphinxstyleemphasis{shift} or \sphinxstyleemphasis{tstat} or \sphinxstyleemphasis{rescale}
\begin{sphinxalltt}
\sphinxstyleemphasis{lamda} value = mean free path of SRD particles (distance units)
\sphinxstyleemphasis{collision} value = \sphinxstyleemphasis{noslip} or \sphinxstyleemphasis{slip} = collision model
\sphinxstyleemphasis{overlap} value = \sphinxstyleemphasis{yes} or \sphinxstyleemphasis{no} = whether big particles may overlap
\sphinxstyleemphasis{inside} value = \sphinxstyleemphasis{error} or \sphinxstyleemphasis{warn} or \sphinxstyleemphasis{ignore} = how SRD particles which end up inside a big particle are treated
\sphinxstyleemphasis{exact} value = \sphinxstyleemphasis{yes} or \sphinxstyleemphasis{no}
\sphinxstyleemphasis{radius} value = rfactor = scale collision radius by this factor
\sphinxstyleemphasis{bounce} value = Nbounce = max \# of collisions an SRD particle can undergo in one timestep
\sphinxstyleemphasis{search} value = sgrid = grid spacing for collision partner searching (distance units)
\sphinxstyleemphasis{cubic} values = style tolerance
  style = \sphinxstyleemphasis{error} or \sphinxstyleemphasis{warn}
  tolerance = fractional difference allowed (0 \textless{}= tol \textless{}= 1)
\sphinxstyleemphasis{shift} values = flag shiftseed
  flag = \sphinxstyleemphasis{yes} or \sphinxstyleemphasis{no} or \sphinxstyleemphasis{possible} = SRD bin shifting for better statistics
    \sphinxstyleemphasis{yes} = perform bin shifting each time SRD velocities are rescaled
    \sphinxstyleemphasis{no} = no shifting
    \sphinxstyleemphasis{possible} = shift depending on mean free path and bin size
  shiftseed = random \# seed (positive integer)
\sphinxstyleemphasis{tstat} value = \sphinxstyleemphasis{yes} or \sphinxstyleemphasis{no} = thermostat SRD particles or not
\sphinxstyleemphasis{rescale} value = \sphinxstyleemphasis{yes} or \sphinxstyleemphasis{no} or \sphinxstyleemphasis{rotate} or \sphinxstyleemphasis{collide} = rescaling of SRD velocities
  \sphinxstyleemphasis{yes} = rescale during velocity rotation and collisions
  \sphinxstyleemphasis{no} = no rescaling
  \sphinxstyleemphasis{rotate} = rescale during velocity rotation, but not collisions
  \sphinxstyleemphasis{collide} = rescale during collisions, but not velocity rotation
\end{sphinxalltt}

\end{itemize}


\subsection{Examples}
\label{\detokenize{fix_srd:examples}}
\begin{sphinxVerbatim}[commandchars=\\\{\}]
\PYG{k}{fix }\PYG{n+nv+nvIdentifier}{1}\PYG{+w}{ }\PYG{n+nv+nvIdentifier}{srd}\PYG{+w}{ }\PYG{n}{srd}\PYG{+w}{ }\PYG{l+m}{10}\PYG{+w}{ }\PYG{n}{big}\PYG{+w}{ }\PYG{l+m}{1.0}\PYG{+w}{ }\PYG{l+m}{0.25}\PYG{+w}{ }\PYG{l+m}{482984}
\PYG{k}{fix }\PYG{n+nv+nvIdentifier}{1}\PYG{+w}{ }\PYG{n+nv+nvIdentifier}{srd}\PYG{+w}{ }\PYG{n}{srd}\PYG{+w}{ }\PYG{l+m}{10}\PYG{+w}{ }\PYG{n}{big}\PYG{+w}{ }\PYG{l+m}{0.5}\PYG{+w}{ }\PYG{l+m}{0.25}\PYG{+w}{ }\PYG{l+m}{482984}\PYG{+w}{ }\PYG{n}{collision}\PYG{+w}{ }\PYG{n}{slip}\PYG{+w}{ }\PYG{n}{search}\PYG{+w}{ }\PYG{l+m}{0.5}
\end{sphinxVerbatim}


\subsection{Description}
\label{\detokenize{fix_srd:description}}
\sphinxAtStartPar
Treat a group of particles as stochastic rotation dynamics (SRD)
particles that serve as a background solvent when interacting with big
(colloidal) particles in groupbig\sphinxhyphen{}ID.  The SRD formalism is described
in {\hyperref[\detokenize{fix_srd:hecht}]{\sphinxcrossref{\DUrole{std,std-ref}{(Hecht)}}}}.  The same methodology is also called
multi\sphinxhyphen{}particle collision dynamics (MPCD) in the literature.

\sphinxAtStartPar
The key idea behind using SRD particles as a cheap coarse\sphinxhyphen{}grained
solvent is that SRD particles do not interact with each other, but
only with the solute particles, which in LAMMPS can be spheroids,
ellipsoids, or line segments, or triangles, or rigid bodies containing
multiple spheroids or ellipsoids or line segments or triangles.  The
collision and rotation properties of the model imbue the SRD particles
with fluid\sphinxhyphen{}like properties, including an effective viscosity.  Thus
simulations with large solute particles can be run more quickly, to
measure solute properties like diffusivity and viscosity in a
background fluid.  The usual LAMMPS fixes for such simulations, such
as {\hyperref[\detokenize{fix_deform::doc}]{\sphinxcrossref{\DUrole{doc}{fix deform}}}}, {\hyperref[\detokenize{fix_viscosity::doc}]{\sphinxcrossref{\DUrole{doc}{fix viscosity}}}}, and {\hyperref[\detokenize{fix_nvt_sllod::doc}]{\sphinxcrossref{\DUrole{doc}{fix nvt/sllod}}}}, can be
used in conjunction with the SRD model.

\sphinxAtStartPar
These 3 papers give more details on how the SRD model is implemented
in LAMMPS.  {\hyperref[\detokenize{fix_srd:petersen1}]{\sphinxcrossref{\DUrole{std,std-ref}{(Petersen)}}}} describes pure SRD fluid
systems.  {\hyperref[\detokenize{fix_srd:bolintineanu1}]{\sphinxcrossref{\DUrole{std,std-ref}{(Bolintineanu1)}}}} describes models
where pure SRD fluids interact with boundary walls.
{\hyperref[\detokenize{fix_srd:bolintineanu2}]{\sphinxcrossref{\DUrole{std,std-ref}{(Bolintineanu2)}}}} describes mixture models where
large colloidal particles are solvated by an SRD fluid.  See the
\sphinxcode{\sphinxupquote{examples/srd}} directory for sample input scripts.

\sphinxAtStartPar
This fix does two things:
\begin{enumerate}
\sphinxsetlistlabels{\arabic}{enumi}{enumii}{}{.}%
\item {} 
\sphinxAtStartPar
It advects the SRD particles, performing collisions between SRD
and big particles or walls every timestep, imparting force and torque
to the big particles.  Collisions also change the position and
velocity of SRD particles.

\item {} 
\sphinxAtStartPar
It resets the velocity distribution of SRD particles via random
rotations every N timesteps.

\end{enumerate}

\sphinxAtStartPar
SRD particles have a mass, temperature, characteristic timestep
\(dt_{SRD}\), and mean free path between collisions
(\(\lambda\)).  The fundamental equation relating these 4 quantities
is
\begin{equation*}
\begin{split}\lambda = dt_{SRD} \sqrt{\frac{k_B T_{SRD}}{m}}\end{split}
\end{equation*}
\sphinxAtStartPar
The mass \sphinxstyleemphasis{m} of SRD particles is set by the {\hyperref[\detokenize{mass::doc}]{\sphinxcrossref{\DUrole{doc}{mass}}}} command
elsewhere in the input script.  The SRD timestep \(dt_{SRD}\) is N
times the step \sphinxstyleemphasis{dt} defined by the {\hyperref[\detokenize{timestep::doc}]{\sphinxcrossref{\DUrole{doc}{timestep}}}} command.
Big particles move in the normal way via a time integration {\hyperref[\detokenize{fix::doc}]{\sphinxcrossref{\DUrole{doc}{fix}}}} with a short timestep dt.  SRD particles advect with a large
timestep \(dt_{SRD} \ge dt\).

\sphinxAtStartPar
If the \sphinxstyleemphasis{lamda} keyword is not specified, the SRD temperature
\(T_{SRD}\) is used in the above formula to compute \(\lambda\).
If the \sphinxstyleemphasis{lamda} keyword is specified, then the \sphinxstyleemphasis{Tsrd} setting is ignored
and the above equation is used to compute the SRD temperature.

\sphinxAtStartPar
The characteristic length scale for the SRD fluid is set by \sphinxstyleemphasis{hgrid}
which is used to bin SRD particles for purposes of resetting their
velocities.  Normally hgrid is set to be 1/4 of the big particle
diameter or smaller, to adequately resolve fluid properties around the
big particles.

\sphinxAtStartPar
\(\lambda\) cannot be smaller than 0.6 * hgrid, else an error is
generated (unless the \sphinxstyleemphasis{shift} keyword is used, see below).  The
velocities of SRD particles are bounded by Vmax, which is set so that an
SRD particle will not advect further than \(D_{max} = 4 \lambda\) in
\(dt_{SRD}\).  This means that roughly speaking, \(D_{max}\)
should not be larger than a big particle diameter, else SRDs may pass
through big particles without colliding.  A warning is generated if this
is the case.

\sphinxAtStartPar
Collisions between SRD particles and big particles or walls are
modeled as a lightweight SRD point particle hitting a heavy big
particle of given diameter or a wall at a point on its surface and
bouncing off with a new velocity.  The collision changes the momentum
of the SRD particle.  It imparts a force and torque to the big
particle.  It imparts a force to a wall.  Static or moving SRD walls
are setup via the {\hyperref[\detokenize{fix_wall_srd::doc}]{\sphinxcrossref{\DUrole{doc}{fix wall/srd}}}} command.  For the
remainder of this doc page, a collision of an SRD particle with a wall
can be viewed as a collision with a big particle of infinite radius
and mass.

\sphinxAtStartPar
The \sphinxstyleemphasis{collision} keyword sets the style of collisions.  The \sphinxstyleemphasis{slip}
style means that the tangential component of the SRD particle momentum
is preserved.  Thus a force is imparted to a big particle, but no
torque.  The normal component of the new SRD velocity is sampled from
a Gaussian distribution at temperature \sphinxstyleemphasis{Tsrd}.

\sphinxAtStartPar
For the \sphinxstyleemphasis{noslip} style, both the normal and tangential components of
the new SRD velocity are sampled from a Gaussian distribution at
temperature \sphinxstyleemphasis{Tsrd}.  Additionally, a new tangential direction for the
SRD velocity is chosen randomly.  This collision style imparts torque
to a big particle.  Thus a time integrator {\hyperref[\detokenize{fix::doc}]{\sphinxcrossref{\DUrole{doc}{fix}}}} that rotates
the big particles appropriately should be used.


\bigskip\hrule\bigskip


\sphinxAtStartPar
The \sphinxstyleemphasis{overlap} keyword should be set to \sphinxstyleemphasis{yes} if two (or more) big
particles can ever overlap.  This depends on the pair potential
interaction used for big\sphinxhyphen{}big interactions, or could be the case if
multiple big particles are held together as rigid bodies via the
{\hyperref[\detokenize{fix_rigid::doc}]{\sphinxcrossref{\DUrole{doc}{fix rigid}}}} command.  If the \sphinxstyleemphasis{overlap} keyword is \sphinxstyleemphasis{no}
and big particles do in fact overlap, then SRD/big collisions can
generate an error if an SRD ends up inside two (or more) big particles
at once.  How this error is treated is determined by the \sphinxstyleemphasis{inside}
keyword.  Running with \sphinxstyleemphasis{overlap} set to \sphinxstyleemphasis{no} allows for faster collision
checking, so it should only be set to \sphinxstyleemphasis{yes} if needed.

\sphinxAtStartPar
The \sphinxstyleemphasis{inside} keyword determines how a collision is treated if the
computation determines that the timestep started with the SRD particle
already inside a big particle.  If the setting is \sphinxstyleemphasis{error} then this
generates an error message and LAMMPS stops.  If the setting is \sphinxstyleemphasis{warn}
then this generates a warning message and the code continues.  If the
setting is \sphinxstyleemphasis{ignore} then no message is generated.  One of the output
quantities logged by the fix (see below) tallies the number of such
events, so it can be monitored.  Note that once an SRD particle is
inside a big particle, it may remain there for several steps until it
drifts outside the big particle.

\sphinxAtStartPar
The \sphinxstyleemphasis{exact} keyword determines how accurately collisions are computed.
A setting of \sphinxstyleemphasis{yes} computes the time and position of each collision as
SRD and big particles move together.  A setting of \sphinxstyleemphasis{no} estimates the
position of each collision based on the end\sphinxhyphen{}of\sphinxhyphen{}timestep positions of
the SRD and big particle.  If \sphinxstyleemphasis{overlap} is set to yes, the setting of
the \sphinxstyleemphasis{exact} keyword is ignored since time\sphinxhyphen{}accurate collisions are
needed.

\sphinxAtStartPar
The \sphinxstyleemphasis{radius} keyword scales the effective size of big particles.  If
big particles will overlap as they undergo dynamics, then this keyword
can be used to scale down their effective collision radius by an
amount \sphinxstyleemphasis{rfactor}, so that SRD particle will only collide with one big
particle at a time.  For example, in a Lennard\sphinxhyphen{}Jones system at a
temperature of 1.0 (in reduced LJ units), the minimum separation
between two big particles is as small as about 0.88 sigma.  Thus an
\sphinxstyleemphasis{rfactor} value of 0.85 should prevent dual collisions.

\sphinxAtStartPar
The \sphinxstyleemphasis{bounce} keyword can be used to limit the maximum number of
collisions an SRD particle undergoes in a single timestep as it
bounces between nearby big particles.  Note that if the limit is
reached, the SRD can be left inside a big particle.  A setting of 0 is
the same as no limit.


\bigskip\hrule\bigskip


\sphinxAtStartPar
There are 2 kinds of bins created and maintained when running an SRD
simulation.  The first are “SRD bins” which are used to bin SRD
particles and reset their velocities, as discussed above.  The second
are “search bins” which are used to identify SRD/big particle
collisions.

\sphinxAtStartPar
The \sphinxstyleemphasis{search} keyword can be used to choose a search bin size for
identifying SRD/big particle collisions.  The default is to use the
\sphinxstyleemphasis{hgrid} parameter for SRD bins as the search bin size.  Choosing a
smaller or large value may be more efficient, depending on the
problem.  But, in a statistical sense, it should not change the
simulation results.

\sphinxAtStartPar
The \sphinxstyleemphasis{cubic} keyword can be used to generate an error or warning when
the bin size chosen by LAMMPS creates SRD bins that are non\sphinxhyphen{}cubic or
different than the requested value of \sphinxstyleemphasis{hgrid} by a specified
\sphinxstyleemphasis{tolerance}.  Note that using non\sphinxhyphen{}cubic SRD bins can lead to
undetermined behavior when rotating the velocities of SRD particles,
hence LAMMPS tries to protect you from this problem.

\sphinxAtStartPar
LAMMPS attempts to set the SRD bin size to exactly \sphinxstyleemphasis{hgrid}.  However,
there must be an integer number of bins in each dimension of the
simulation box.  Thus the actual bin size will depend on the size and
shape of the overall simulation box.  The actual bin size is printed
as part of the SRD output when a simulation begins.

\sphinxAtStartPar
If the actual bin size in non\sphinxhyphen{}cubic by an amount exceeding the
tolerance, an error or warning is printed, depending on the style of
the \sphinxstyleemphasis{cubic} keyword.  Likewise, if the actual bin size differs from
the requested \sphinxstyleemphasis{hgrid} value by an amount exceeding the tolerance, then
an error or warning is printed.  The \sphinxstyleemphasis{tolerance} is a fractional
difference.  E.g. a tolerance setting of 0.01 on the shape means that
if the ratio of any 2 bin dimensions exceeds (1 +/\sphinxhyphen{} tolerance) then an
error or warning is generated.  Similarly, if the ratio of any bin
dimension with \sphinxstyleemphasis{hgrid} exceeds (1 +/\sphinxhyphen{} tolerance), then an error or
warning is generated.

\begin{sphinxadmonition}{note}{Note:}
\sphinxAtStartPar
The fix srd command can be used with simulations where the size and/or
shape of the simulation box changes.  This can be due to non\sphinxhyphen{}periodic
boundary conditions or the use of fixes such as the {\hyperref[\detokenize{fix_deform::doc}]{\sphinxcrossref{\DUrole{doc}{fix deform}}}} or {\hyperref[\detokenize{fix_wall_srd::doc}]{\sphinxcrossref{\DUrole{doc}{fix wall/srd}}}} commands
to impose a shear on an SRD fluid or an interaction with an external
wall.  If the box size changes then the size of SRD bins must be
recalculated every reneighboring.  This is not necessary if only the
box shape changes.  This re\sphinxhyphen{}binning is always done so as to fit an
integer number of bins in the current box dimension, whether it be a
fixed, shrink\sphinxhyphen{}wrapped, or periodic boundary, as set by the
{\hyperref[\detokenize{boundary::doc}]{\sphinxcrossref{\DUrole{doc}{boundary}}}} command.  If the box size or shape changes,
then the size of the search bins must be recalculated every
reneighboring.  Note that changing the SRD bin size may alter the
properties of the SRD fluid, such as its viscosity.
\end{sphinxadmonition}

\sphinxAtStartPar
The \sphinxstyleemphasis{shift} keyword determines whether the coordinates of SRD particles
are randomly shifted when binned for purposes of rotating their
velocities.  When no shifting is performed, SRD particles are binned and
the velocity distribution of the set of SRD particles in each bin is
adjusted via a rotation operator.  This is a statistically valid
operation if SRD particles move sufficiently far between successive
rotations.  This is determined by their mean\sphinxhyphen{}free path \(\lambda\).
If \(\lambda\) is less than 0.6 of the SRD bin size, then shifting
is required.  A shift means that all of the SRD particles are shifted by
a vector whose coordinates are chosen randomly in the range {[}\sphinxhyphen{}1/2 bin
size, 1/2 bin size{]}.  Note that all particles are shifted by the same
vector.  The specified random number \sphinxstyleemphasis{shiftseed} is used to generate
these vectors.  This operation sufficiently randomizes which SRD
particles are in the same bin, even if \(lambda\) is small.

\sphinxAtStartPar
If the \sphinxstyleemphasis{shift} flag is set to \sphinxstyleemphasis{no}, then no shifting is performed, but
bin data will be communicated if bins overlap processor boundaries.  An
error will be generated if \(\lambda < 0.6\) of the SRD bin size.
If the \sphinxstyleemphasis{shift} flag is set to \sphinxstyleemphasis{possible}, then shifting is performed
only if \(\lambda < 0.6\) of the SRD bin size.  A warning is
generated to let you know this is occurring.  If the \sphinxstyleemphasis{shift} flag is set
to \sphinxstyleemphasis{yes} then shifting is performed regardless of the magnitude of
\(\lambda\).  Note that the \sphinxstyleemphasis{shiftseed} is not used if the \sphinxstyleemphasis{shift}
flag is set to \sphinxstyleemphasis{no}, but must still be specified.

\sphinxAtStartPar
Note that shifting of SRD coordinates requires extra communication,
hence it should not normally be enabled unless required.

\sphinxAtStartPar
The \sphinxstyleemphasis{tstat} keyword will thermostat the SRD particles to the specified
\sphinxstyleemphasis{Tsrd}.  This is done every N timesteps, during the velocity rotation
operation, by rescaling the thermal velocity of particles in each SRD
bin to the desired temperature.  If there is a streaming velocity
associated with the system, e.g. due to use of the {\hyperref[\detokenize{fix_deform::doc}]{\sphinxcrossref{\DUrole{doc}{fix deform}}}} command to perform a simulation undergoing
shear, then that is also accounted for.  The mean velocity of each bin
of SRD particles is set to the position\sphinxhyphen{}dependent streaming velocity,
based on the coordinates of the center of the SRD bin.  Note that
collisions of SRD particles with big particles or walls has a
thermostatting effect on the colliding particles, so it may not be
necessary to thermostat the SRD particles on a bin by bin basis in
that case.  Also note that for streaming simulations, if no
thermostatting is performed (the default), then it may take a long
time for the SRD fluid to come to equilibrium with a velocity profile
that matches the simulation box deformation.

\sphinxAtStartPar
The \sphinxstyleemphasis{rescale} keyword enables rescaling of an SRD particle’s velocity
if it would travel more than 4 mean\sphinxhyphen{}free paths in an SRD timestep.  If
an SRD particle exceeds this velocity it is possible it will be lost
when migrating to other processors or that collisions with big
particles will be missed, either of which will generate errors.  Thus
the safest mode is to run with rescaling enabled.  However rescaling
removes kinetic energy from the system (the particle’s velocity is
reduced).  The latter will not typically be a problem if
thermostatting is enabled via the \sphinxstyleemphasis{tstat} keyword or if SRD collisions
with big particles or walls effectively thermostat the system.  If you
wish to turn off rescaling (on is the default), e.g. for a pure SRD
system with no thermostatting so that the temperature does not decline
over time, the \sphinxstyleemphasis{rescale} keyword can be used.  The \sphinxstyleemphasis{no} value turns
rescaling off during collisions and the per\sphinxhyphen{}bin velocity rotation
operation.  The \sphinxstyleemphasis{collide} and \sphinxstyleemphasis{rotate} values turn it on for
one of the operations and off for the other.


\bigskip\hrule\bigskip


\begin{sphinxadmonition}{note}{Note:}
\sphinxAtStartPar
This fix is normally used for simulations with a huge number of
SRD particles relative to the number of big particles, e.g. 100 to 1.
In this scenario, computations that involve only big particles
(neighbor list creation, communication, time integration) can slow
down dramatically due to the large number of background SRD particles.
\end{sphinxadmonition}

\sphinxAtStartPar
Three other input script commands will largely overcome this effect,
speeding up an SRD simulation by a significant amount.  These are the
{\hyperref[\detokenize{atom_modify::doc}]{\sphinxcrossref{\DUrole{doc}{atom\_modify first}}}}, {\hyperref[\detokenize{neigh_modify::doc}]{\sphinxcrossref{\DUrole{doc}{neigh\_modify include}}}}, and {\hyperref[\detokenize{comm_modify::doc}]{\sphinxcrossref{\DUrole{doc}{comm\_modify group}}}}
commands.  Each takes a group\sphinxhyphen{}ID as an argument, which in this case
should be the group\sphinxhyphen{}ID of the big solute particles.

\sphinxAtStartPar
Additionally, when a {\hyperref[\detokenize{pair_style::doc}]{\sphinxcrossref{\DUrole{doc}{pair\_style}}}} for big/big particle
interactions is specified, the {\hyperref[\detokenize{pair_coeff::doc}]{\sphinxcrossref{\DUrole{doc}{pair\_coeff}}}} command
should be used to turn off big/SRD interactions, e.g. by setting their
epsilon or cutoff length to 0.0.

\sphinxAtStartPar
The “delete\_atoms overlap” command may be useful in setting up an SRD
simulation to ensure there are no initial overlaps between big and SRD
particles.


\bigskip\hrule\bigskip



\subsection{Restart, fix\_modify, output, run start/stop, minimize info}
\label{\detokenize{fix_srd:restart-fix-modify-output-run-start-stop-minimize-info}}
\sphinxAtStartPar
No information about this fix is written to {\hyperref[\detokenize{restart::doc}]{\sphinxcrossref{\DUrole{doc}{binary restart files}}}}.  None of the {\hyperref[\detokenize{fix_modify::doc}]{\sphinxcrossref{\DUrole{doc}{fix\_modify}}}} options
are relevant to this fix.

\sphinxAtStartPar
This fix tabulates several SRD statistics which are stored in a vector
of length 12, which can be accessed by various {\hyperref[\detokenize{Howto_output::doc}]{\sphinxcrossref{\DUrole{doc}{output commands}}}}.  The vector values calculated by this fix
are “intensive”, meaning they do not scale with the size of the
simulation.  Technically, the first 8 do scale with the size of the
simulation, but treating them as intensive means they are not scaled
when printed as part of thermodynamic output.

\sphinxAtStartPar
These are the 12 quantities.  All are values for the current timestep,
except for quantity 5 and the last three, each of which are
cumulative quantities since the beginning of the run.
\begin{enumerate}
\sphinxsetlistlabels{\arabic}{enumi}{enumii}{(}{)}%
\item {} 
\sphinxAtStartPar
\# of SRD/big collision checks performed

\item {} 
\sphinxAtStartPar
\# of SRDs which had a collision

\item {} 
\sphinxAtStartPar
\# of SRD/big collisions (including multiple bounces)

\item {} 
\sphinxAtStartPar
\# of SRD particles inside a big particle

\item {} 
\sphinxAtStartPar
\# of SRD particles whose velocity was rescaled to be \textless{} Vmax

\item {} 
\sphinxAtStartPar
\# of bins for collision searching

\item {} 
\sphinxAtStartPar
\# of bins for SRD velocity rotation

\item {} 
\sphinxAtStartPar
\# of bins in which SRD temperature was computed

\item {} 
\sphinxAtStartPar
SRD temperature

\item {} 
\sphinxAtStartPar
\# of SRD particles which have undergone max \# of bounces

\item {} 
\sphinxAtStartPar
max \# of bounces any SRD particle has had in a single step

\item {} 
\sphinxAtStartPar
\# of reneighborings due to SRD particles moving too far

\end{enumerate}

\sphinxAtStartPar
No parameter of this fix can be used with the \sphinxstyleemphasis{start/stop} keywords of
the {\hyperref[\detokenize{run::doc}]{\sphinxcrossref{\DUrole{doc}{run}}}} command.  This fix is not invoked during
{\hyperref[\detokenize{minimize::doc}]{\sphinxcrossref{\DUrole{doc}{energy minimization}}}}.


\subsection{Restrictions}
\label{\detokenize{fix_srd:restrictions}}
\sphinxAtStartPar
This command can only be used if LAMMPS was built with the SRD package.
See the {\hyperref[\detokenize{Build_package::doc}]{\sphinxcrossref{\DUrole{doc}{Build package}}}} doc page for more info.


\subsection{Related commands}
\label{\detokenize{fix_srd:related-commands}}
\sphinxAtStartPar
{\hyperref[\detokenize{fix_wall_srd::doc}]{\sphinxcrossref{\DUrole{doc}{fix wall/srd}}}}


\subsection{Default}
\label{\detokenize{fix_srd:default}}
\sphinxAtStartPar
The option defaults are: \sphinxstyleemphasis{lamda} (\(\lambda\)) is inferred from \sphinxstyleemphasis{Tsrd},
collision = noslip, overlap = no, inside = error, exact = yes, radius =
1.0, bounce = 0, search = hgrid, cubic = error 0.01, shift = no, tstat =
no, and rescale = yes.


\bigskip\hrule\bigskip

\phantomsection\label{\detokenize{fix_srd:hecht}}
\sphinxAtStartPar
\sphinxstylestrong{(Hecht)} Hecht, Harting, Ihle, Herrmann, Phys Rev E, 72, 011408 (2005).

\phantomsection\label{\detokenize{fix_srd:petersen1}}
\sphinxAtStartPar
\sphinxstylestrong{(Petersen)} Petersen, Lechman, Plimpton, Grest, in’ t Veld, Schunk, J
Chem Phys, 132, 174106 (2010).

\phantomsection\label{\detokenize{fix_srd:bolintineanu1}}
\sphinxAtStartPar
\sphinxstylestrong{(Bolintineanu1)} Bolintineanu, Lechman, Plimpton, Grest, Phys Rev E, 86, 066703 (2012).

\phantomsection\label{\detokenize{fix_srd:bolintineanu2}}
\sphinxAtStartPar
\sphinxstylestrong{(Bolintineanu2)} Bolintineanu, Grest, Lechman, Pierce, Plimpton, Schunk, Comp Particle Mechanics, 1, 321\sphinxhyphen{}356 (2014).

\sphinxstepscope

\index{fix store/force@\spxentry{fix store/force}}\ignorespaces 

\section{fix store/force command}
\label{\detokenize{fix_store_force:fix-store-force-command}}\label{\detokenize{fix_store_force:index-0}}\label{\detokenize{fix_store_force::doc}}

\subsection{Syntax}
\label{\detokenize{fix_store_force:syntax}}
\begin{sphinxVerbatim}[commandchars=\\\{\}]
\PYG{k}{fix }\PYG{n+nv+nvIdentifier}{ID}\PYG{+w}{ }\PYG{n+nv+nvIdentifier}{group\PYGZhy{}ID}\PYG{+w}{ }\PYG{n}{store}\PYG{o}{/}\PYG{n}{force}
\end{sphinxVerbatim}
\begin{itemize}
\item {} 
\sphinxAtStartPar
ID, group\sphinxhyphen{}ID are documented in {\hyperref[\detokenize{fix::doc}]{\sphinxcrossref{\DUrole{doc}{fix}}}} command

\item {} 
\sphinxAtStartPar
store/force = style name of this fix command

\end{itemize}


\subsection{Examples}
\label{\detokenize{fix_store_force:examples}}
\begin{sphinxVerbatim}[commandchars=\\\{\}]
\PYG{k}{fix }\PYG{n+nv+nvIdentifier}{1}\PYG{+w}{ }\PYG{n+nv+nvIdentifier}{all}\PYG{+w}{ }\PYG{n}{store}\PYG{o}{/}\PYG{n}{force}
\end{sphinxVerbatim}


\subsection{Description}
\label{\detokenize{fix_store_force:description}}
\sphinxAtStartPar
Store the forces on atoms in the group at the point during each timestep
when the fix is invoked, as described below.  This is useful for storing
forces before constraints or other boundary conditions are computed
which modify the forces, so that unmodified forces can be {\hyperref[\detokenize{dump::doc}]{\sphinxcrossref{\DUrole{doc}{written
to a dump file}}}} or accessed by other {\hyperref[\detokenize{Howto_output::doc}]{\sphinxcrossref{\DUrole{doc}{output commands}}}} that use per\sphinxhyphen{}atom quantities.

\sphinxAtStartPar
This fix is invoked at the point in the velocity\sphinxhyphen{}Verlet timestepping
immediately after {\hyperref[\detokenize{pair_style::doc}]{\sphinxcrossref{\DUrole{doc}{pair}}}}, {\hyperref[\detokenize{bond_style::doc}]{\sphinxcrossref{\DUrole{doc}{bond}}}},
{\hyperref[\detokenize{angle_style::doc}]{\sphinxcrossref{\DUrole{doc}{angle}}}}, {\hyperref[\detokenize{dihedral_style::doc}]{\sphinxcrossref{\DUrole{doc}{dihedral}}}},
{\hyperref[\detokenize{improper_style::doc}]{\sphinxcrossref{\DUrole{doc}{improper}}}}, and {\hyperref[\detokenize{kspace_style::doc}]{\sphinxcrossref{\DUrole{doc}{long\sphinxhyphen{}range}}}}
forces have been calculated.  It is the point in the timestep when
various fixes that compute constraint forces are calculated and
potentially modify the force on each atom.  Examples of such fixes are
{\hyperref[\detokenize{fix_shake::doc}]{\sphinxcrossref{\DUrole{doc}{fix shake}}}}, {\hyperref[\detokenize{fix_wall::doc}]{\sphinxcrossref{\DUrole{doc}{fix wall}}}}, and {\hyperref[\detokenize{fix_indent::doc}]{\sphinxcrossref{\DUrole{doc}{fix
indent}}}}.

\begin{sphinxadmonition}{note}{Note:}
\sphinxAtStartPar
The order in which various fixes are applied which operate at the
same point during the timestep, is the same as the order they are
specified in the input script.  Thus normally, if you want to store
per\sphinxhyphen{}atom forces due to force field interactions, before constraints
are applied, you should list this fix first within that set of fixes,
i.e. before other fixes that apply constraints.  However, if you wish
to include certain constraints (e.g. fix shake) in the stored force,
then it could be specified after some fixes and before others.
\end{sphinxadmonition}


\subsection{Restart, fix\_modify, output, run start/stop, minimize info}
\label{\detokenize{fix_store_force:restart-fix-modify-output-run-start-stop-minimize-info}}
\sphinxAtStartPar
No information about this fix is written to {\hyperref[\detokenize{restart::doc}]{\sphinxcrossref{\DUrole{doc}{binary restart files}}}}.  None of the {\hyperref[\detokenize{fix_modify::doc}]{\sphinxcrossref{\DUrole{doc}{fix\_modify}}}} options are
relevant to this fix.

\sphinxAtStartPar
This fix produces a per\sphinxhyphen{}atom array which can be accessed by various
{\hyperref[\detokenize{Howto_output::doc}]{\sphinxcrossref{\DUrole{doc}{output commands}}}}.  The number of columns for each
atom is 3, and the columns store the x,y,z forces on each atom.  The
per\sphinxhyphen{}atom values be accessed on any timestep.

\sphinxAtStartPar
No parameter of this fix can be used with the \sphinxstyleemphasis{start/stop} keywords of
the {\hyperref[\detokenize{run::doc}]{\sphinxcrossref{\DUrole{doc}{run}}}} command.  This fix is not invoked during
{\hyperref[\detokenize{minimize::doc}]{\sphinxcrossref{\DUrole{doc}{energy minimization}}}}.


\subsection{Restrictions}
\label{\detokenize{fix_store_force:restrictions}}\begin{quote}

\sphinxAtStartPar
none
\end{quote}


\subsection{Related commands}
\label{\detokenize{fix_store_force:related-commands}}
\sphinxAtStartPar
{\hyperref[\detokenize{fix_store_state::doc}]{\sphinxcrossref{\DUrole{doc}{fix store\_state}}}}


\subsection{Default}
\label{\detokenize{fix_store_force:default}}
\sphinxAtStartPar
none

\sphinxstepscope

\index{fix store/state@\spxentry{fix store/state}}\ignorespaces 

\section{fix store/state command}
\label{\detokenize{fix_store_state:fix-store-state-command}}\label{\detokenize{fix_store_state:index-0}}\label{\detokenize{fix_store_state::doc}}

\subsection{Syntax}
\label{\detokenize{fix_store_state:syntax}}
\begin{sphinxVerbatim}[commandchars=\\\{\}]
\PYG{k}{fix }\PYG{n+nv+nvIdentifier}{ID}\PYG{+w}{ }\PYG{n+nv+nvIdentifier}{group\PYGZhy{}ID}\PYG{+w}{ }\PYG{n}{store}\PYG{o}{/}\PYG{n}{state}\PYG{+w}{ }\PYG{n}{N}\PYG{+w}{ }\PYG{n}{input1}\PYG{+w}{ }\PYG{n}{input2}\PYG{+w}{ }\PYG{n}{...}\PYG{+w}{ }\PYG{n}{keyword}\PYG{+w}{ }\PYG{n}{value}\PYG{+w}{ }\PYG{n}{...}
\end{sphinxVerbatim}
\begin{itemize}
\item {} 
\sphinxAtStartPar
ID, group\sphinxhyphen{}ID are documented in {\hyperref[\detokenize{fix::doc}]{\sphinxcrossref{\DUrole{doc}{fix}}}} command

\item {} 
\sphinxAtStartPar
store/state = style name of this fix command

\item {} 
\sphinxAtStartPar
N = store atom attributes every N steps, N = 0 for initial store only

\item {} 
\sphinxAtStartPar
input = one or more atom attributes

\begin{sphinxVerbatim}[commandchars=\\\{\}]
possible attributes = id, mol, type, mass,
                      x, y, z, xs, ys, zs, xu, yu, zu, xsu, ysu, zsu, ix, iy, iz,
                      vx, vy, vz, fx, fy, fz,
                      q, mux, muy, muz, mu,
                      radius, diameter, omegax, omegay, omegaz,
                      angmomx, angmomy, angmomz, tqx, tqy, tqz,
                      c\PYGZus{}ID, c\PYGZus{}ID[I], f\PYGZus{}ID, f\PYGZus{}ID[I], v\PYGZus{}name,
                      d\PYGZus{}name, i\PYGZus{}name, i2\PYGZus{}name[I], d2\PYGZus{}name[I],
\end{sphinxVerbatim}
\begin{sphinxalltt}
id = atom ID
mol = molecule ID
type = atom type
mass = atom mass
x,y,z = unscaled atom coordinates
xs,ys,zs = scaled atom coordinates
xu,yu,zu = unwrapped atom coordinates
xsu,ysu,zsu = scaled unwrapped atom coordinates
ix,iy,iz = box image that the atom is in
vx,vy,vz = atom velocities
fx,fy,fz = forces on atoms
q = atom charge
mux,muy,muz = orientation of dipolar atom
mu = magnitued of dipole moment of atom
radius,diameter = radius.diameter of spherical particle
omegax,omegay,omegaz = angular velocity of spherical particle
angmomx,angmomy,angmomz = angular momentum of aspherical particle
tqx,tqy,tqz = torque on finite\sphinxhyphen{}size particles
\sphinxstyleemphasis{c\_ID} = per\sphinxhyphen{}atom vector calculated by a compute with ID
\sphinxstyleemphasis{c\_ID{[}I{]}} = Ith column of per\sphinxhyphen{}atom array calculated by a compute with ID
\sphinxstyleemphasis{f\_ID} = per\sphinxhyphen{}atom vector calculated by a fix with ID
\sphinxstyleemphasis{f\_ID{[}I{]}} = Ith column of per\sphinxhyphen{}atom array calculated by a fix with ID
\sphinxstyleemphasis{v\_name} = per\sphinxhyphen{}atom vector calculated by an atom\sphinxhyphen{}style variable with name
\sphinxstyleemphasis{i\_name} = custom integer vector with name
\sphinxstyleemphasis{d\_name} = custom floating point vector with name
\sphinxstyleemphasis{i2\_name{[}I{]}} = Ith column of custom integer array with name
\sphinxstyleemphasis{d2\_name{[}I{]}} = Ith column of custom floating\sphinxhyphen{}point array with name
\end{sphinxalltt}

\item {} 
\sphinxAtStartPar
zero or more keyword/value pairs may be appended

\item {} 
\sphinxAtStartPar
keyword = \sphinxstyleemphasis{com}
\begin{sphinxalltt}
\sphinxstyleemphasis{com} value = \sphinxstyleemphasis{yes} or \sphinxstyleemphasis{no}
\end{sphinxalltt}

\end{itemize}


\subsection{Examples}
\label{\detokenize{fix_store_state:examples}}
\begin{sphinxVerbatim}[commandchars=\\\{\}]
\PYG{k}{fix }\PYG{n+nv+nvIdentifier}{1}\PYG{+w}{ }\PYG{n+nv+nvIdentifier}{all}\PYG{+w}{ }\PYG{n}{store}\PYG{o}{/}\PYG{n}{state}\PYG{+w}{ }\PYG{l+m}{0}\PYG{+w}{ }\PYG{n}{x}\PYG{+w}{ }\PYG{n}{y}\PYG{+w}{ }\PYG{n}{z}
\PYG{k}{fix }\PYG{n+nv+nvIdentifier}{1}\PYG{+w}{ }\PYG{n+nv+nvIdentifier}{all}\PYG{+w}{ }\PYG{n}{store}\PYG{o}{/}\PYG{n}{state}\PYG{+w}{ }\PYG{l+m}{0}\PYG{+w}{ }\PYG{n}{xu}\PYG{+w}{ }\PYG{n}{yu}\PYG{+w}{ }\PYG{n}{zu}\PYG{+w}{ }\PYG{n}{com}\PYG{+w}{ }\PYG{n}{yes}
\PYG{k}{fix }\PYG{n+nv+nvIdentifier}{2}\PYG{+w}{ }\PYG{n+nv+nvIdentifier}{all}\PYG{+w}{ }\PYG{n}{store}\PYG{o}{/}\PYG{n}{state}\PYG{+w}{ }\PYG{l+m}{1000}\PYG{+w}{ }\PYG{n}{vx}\PYG{+w}{ }\PYG{n}{vy}\PYG{+w}{ }\PYG{n}{vz}
\end{sphinxVerbatim}


\subsection{Description}
\label{\detokenize{fix_store_state:description}}
\sphinxAtStartPar
Define a fix that stores attributes for each atom in the group at the
time the fix is defined.  If \sphinxstyleemphasis{N} is 0, then the values are never
updated, so this is a way of archiving an atom attribute at a given
time for future use in a calculation or output.  See the discussion of
{\hyperref[\detokenize{Howto_output::doc}]{\sphinxcrossref{\DUrole{doc}{output commands}}}} that take fixes as inputs.

\sphinxAtStartPar
If \sphinxstyleemphasis{N} is not zero, then the attributes will be updated every \sphinxstyleemphasis{N}
steps.

\begin{sphinxadmonition}{note}{Note:}
\sphinxAtStartPar
Actually, only atom attributes specified by keywords like \sphinxstyleemphasis{xu}
or \sphinxstyleemphasis{vy} or \sphinxstyleemphasis{radius} are initially stored immediately at the point in
your input script when the fix is defined.  Attributes specified by a
compute, fix, or variable are not initially stored until the first run
following the fix definition begins.  This is because calculating
those attributes may require quantities that are not defined in
between runs.
\end{sphinxadmonition}

\sphinxAtStartPar
The list of possible attributes is the same as that used by the
{\hyperref[\detokenize{dump::doc}]{\sphinxcrossref{\DUrole{doc}{dump custom}}}} command, which describes their meaning.

\sphinxAtStartPar
If the \sphinxstyleemphasis{com} keyword is set to \sphinxstyleemphasis{yes} then the \sphinxstyleemphasis{xu}, \sphinxstyleemphasis{yu}, and \sphinxstyleemphasis{zu}
inputs store the position of each atom relative to the center\sphinxhyphen{}of\sphinxhyphen{}mass
of the group of atoms, instead of storing the absolute position.

\sphinxAtStartPar
The requested values are stored in a per\sphinxhyphen{}atom vector or array as
discussed below.  Zeroes are stored for atoms not in the specified
group.


\subsection{Restart, fix\_modify, output, run start/stop, minimize info}
\label{\detokenize{fix_store_state:restart-fix-modify-output-run-start-stop-minimize-info}}
\sphinxAtStartPar
This fix writes the per\sphinxhyphen{}atom values it stores to {\hyperref[\detokenize{restart::doc}]{\sphinxcrossref{\DUrole{doc}{binary restart
files}}}}, so that the values can be restored when a simulation
is restarted.  See the {\hyperref[\detokenize{read_restart::doc}]{\sphinxcrossref{\DUrole{doc}{read\_restart}}}} command for
info on how to re\sphinxhyphen{}specify a fix in an input script that reads a
restart file, so that the operation of the fix continues in an
uninterrupted fashion.

\begin{sphinxadmonition}{warning}{Warning:}
\sphinxAtStartPar
When reading data from a restart file, this fix command has to be specified
\sphinxstylestrong{exactly} the same way as before. LAMMPS will only check whether a
fix is of the same style and has the same fix ID and in case of a match
will then try to initialize the fix with the data stored in the binary
restart file.  If the fix store/state command does not match exactly,
data can be corrupted or LAMMPS may crash.
\end{sphinxadmonition}

\sphinxAtStartPar
None of the {\hyperref[\detokenize{fix_modify::doc}]{\sphinxcrossref{\DUrole{doc}{fix\_modify}}}} options are relevant to this
fix.

\sphinxAtStartPar
If a single input is specified, this fix produces a per\sphinxhyphen{}atom vector.
If multiple inputs are specified, a per\sphinxhyphen{}atom array is produced where
the number of columns for each atom is the number of inputs.  These
can be accessed by various {\hyperref[\detokenize{Howto_output::doc}]{\sphinxcrossref{\DUrole{doc}{output commands}}}}.  The
per\sphinxhyphen{}atom values be accessed on any timestep.

\sphinxAtStartPar
No parameter of this fix can be used with the \sphinxstyleemphasis{start/stop} keywords of
the {\hyperref[\detokenize{run::doc}]{\sphinxcrossref{\DUrole{doc}{run}}}} command.  This fix is not invoked during
{\hyperref[\detokenize{minimize::doc}]{\sphinxcrossref{\DUrole{doc}{energy minimization}}}}.


\subsection{Restrictions}
\label{\detokenize{fix_store_state:restrictions}}\begin{quote}

\sphinxAtStartPar
none
\end{quote}


\subsection{Related commands}
\label{\detokenize{fix_store_state:related-commands}}
\sphinxAtStartPar
{\hyperref[\detokenize{dump::doc}]{\sphinxcrossref{\DUrole{doc}{dump custom}}}}, {\hyperref[\detokenize{compute_property_atom::doc}]{\sphinxcrossref{\DUrole{doc}{compute property/atom}}}},
{\hyperref[\detokenize{fix_property_atom::doc}]{\sphinxcrossref{\DUrole{doc}{fix property/atom}}}}, {\hyperref[\detokenize{variable::doc}]{\sphinxcrossref{\DUrole{doc}{variable}}}}


\subsection{Default}
\label{\detokenize{fix_store_state:default}}
\sphinxAtStartPar
The option default is com = no.

\sphinxstepscope

\index{fix temp/berendsen@\spxentry{fix temp/berendsen}}\ignorespaces 
\index{fix temp/berendsen/kk@\spxentry{fix temp/berendsen/kk}}\ignorespaces 

\section{fix temp/berendsen command}
\label{\detokenize{fix_temp_berendsen:fix-temp-berendsen-command}}\label{\detokenize{fix_temp_berendsen:index-1}}\label{\detokenize{fix_temp_berendsen:index-0}}\label{\detokenize{fix_temp_berendsen::doc}}
\sphinxAtStartPar
Accelerator Variants: \sphinxstyleemphasis{temp/berendsen/kk}


\subsection{Syntax}
\label{\detokenize{fix_temp_berendsen:syntax}}
\begin{sphinxVerbatim}[commandchars=\\\{\}]
\PYG{k}{fix }\PYG{n+nv+nvIdentifier}{ID}\PYG{+w}{ }\PYG{n+nv+nvIdentifier}{group\PYGZhy{}ID}\PYG{+w}{ }\PYG{n}{temp}\PYG{o}{/}\PYG{n}{berendsen}\PYG{+w}{ }\PYG{n}{Tstart}\PYG{+w}{ }\PYG{n}{Tstop}\PYG{+w}{ }\PYG{n}{Tdamp}
\end{sphinxVerbatim}
\begin{itemize}
\item {} 
\sphinxAtStartPar
ID, group\sphinxhyphen{}ID are documented in {\hyperref[\detokenize{fix::doc}]{\sphinxcrossref{\DUrole{doc}{fix}}}} command

\item {} 
\sphinxAtStartPar
temp/berendsen = style name of this fix command

\item {} 
\sphinxAtStartPar
Tstart,Tstop = desired temperature at start/end of run

\begin{sphinxVerbatim}[commandchars=\\\{\}]
Tstart can be a variable (see below)
\end{sphinxVerbatim}

\item {} 
\sphinxAtStartPar
Tdamp = temperature damping parameter (time units)

\end{itemize}


\subsection{Examples}
\label{\detokenize{fix_temp_berendsen:examples}}
\begin{sphinxVerbatim}[commandchars=\\\{\}]
\PYG{k}{fix }\PYG{n+nv+nvIdentifier}{1}\PYG{+w}{ }\PYG{n+nv+nvIdentifier}{all}\PYG{+w}{ }\PYG{n}{temp}\PYG{o}{/}\PYG{n}{berendsen}\PYG{+w}{ }\PYG{l+m}{300.0}\PYG{+w}{ }\PYG{l+m}{300.0}\PYG{+w}{ }\PYG{l+m}{100.0}
\end{sphinxVerbatim}


\subsection{Description}
\label{\detokenize{fix_temp_berendsen:description}}
\sphinxAtStartPar
Reset the temperature of a group of atoms by using a Berendsen
thermostat {\hyperref[\detokenize{fix_temp_berendsen:berendsen2}]{\sphinxcrossref{\DUrole{std,std-ref}{(Berendsen)}}}}, which rescales their velocities
every timestep.

\sphinxAtStartPar
The thermostat is applied to only the translational degrees of freedom
for the particles, which is an important consideration for finite\sphinxhyphen{}size
particles which have rotational degrees of freedom are being
thermostatted with this fix.  The translational degrees of freedom can
also have a bias velocity removed from them before thermostatting
takes place; see the description below.

\sphinxAtStartPar
The desired temperature at each timestep is a ramped value during the
run from \sphinxstyleemphasis{Tstart} to \sphinxstyleemphasis{Tstop}.  The \sphinxstyleemphasis{Tdamp} parameter is specified in
time units and determines how rapidly the temperature is relaxed.  For
example, a value of 100.0 means to relax the temperature in a timespan
of (roughly) 100 time units (tau or fs or ps \sphinxhyphen{} see the
{\hyperref[\detokenize{units::doc}]{\sphinxcrossref{\DUrole{doc}{units}}}} command).

\sphinxAtStartPar
\sphinxstyleemphasis{Tstart} can be specified as an equal\sphinxhyphen{}style {\hyperref[\detokenize{variable::doc}]{\sphinxcrossref{\DUrole{doc}{variable}}}}.
In this case, the \sphinxstyleemphasis{Tstop} setting is ignored.  If the value is a
variable, it should be specified as v\_name, where name is the variable
name.  In this case, the variable will be evaluated each timestep, and
its value used to determine the target temperature.

\begin{sphinxadmonition}{note}{Note:}
\sphinxAtStartPar
This thermostat will generate an error if the current
temperature is zero at the end of a timestep.  It cannot rescale a
zero temperature.
\end{sphinxadmonition}

\sphinxAtStartPar
Equal\sphinxhyphen{}style variables can specify formulas with various mathematical
functions, and include {\hyperref[\detokenize{thermo_style::doc}]{\sphinxcrossref{\DUrole{doc}{thermo\_style}}}} command
keywords for the simulation box parameters and timestep and elapsed
time.  Thus it is easy to specify a time\sphinxhyphen{}dependent temperature.

\begin{sphinxadmonition}{note}{Note:}
\sphinxAtStartPar
Unlike the {\hyperref[\detokenize{fix_nh::doc}]{\sphinxcrossref{\DUrole{doc}{fix nvt}}}} command which performs
Nose/Hoover thermostatting AND time integration, this fix does NOT
perform time integration.  It only modifies velocities to effect
thermostatting.  Thus you must use a separate time integration fix,
like {\hyperref[\detokenize{fix_nve::doc}]{\sphinxcrossref{\DUrole{doc}{fix nve}}}} to actually update the positions of atoms
using the modified velocities.  Likewise, this fix should not normally
be used on atoms that also have their temperature controlled by
another fix \sphinxhyphen{} e.g. by {\hyperref[\detokenize{fix_nh::doc}]{\sphinxcrossref{\DUrole{doc}{fix nvt}}}} or {\hyperref[\detokenize{fix_langevin::doc}]{\sphinxcrossref{\DUrole{doc}{fix langevin}}}} commands.
\end{sphinxadmonition}

\sphinxAtStartPar
See the {\hyperref[\detokenize{Howto_thermostat::doc}]{\sphinxcrossref{\DUrole{doc}{Howto thermostat}}}} page for a
discussion of different ways to compute temperature and perform
thermostatting.

\sphinxAtStartPar
This fix computes a temperature each timestep.  To do this, the fix
creates its own compute of style “temp”, as if this command had been
issued:

\begin{sphinxVerbatim}[commandchars=\\\{\}]
\PYG{k}{compute }\PYG{n+nv+nvIdentifier}{fix\PYGZhy{}ID\PYGZus{}temp}\PYG{+w}{ }\PYG{n+nv+nvIdentifier}{group\PYGZhy{}ID}\PYG{+w}{ }\PYG{n}{temp}
\end{sphinxVerbatim}

\sphinxAtStartPar
See the {\hyperref[\detokenize{compute_temp::doc}]{\sphinxcrossref{\DUrole{doc}{compute temp}}}} command for details.  Note
that the ID of the new compute is the fix\sphinxhyphen{}ID + underscore + “temp”,
and the group for the new compute is the same as the fix group.

\sphinxAtStartPar
Note that this is NOT the compute used by thermodynamic output (see
the {\hyperref[\detokenize{thermo_style::doc}]{\sphinxcrossref{\DUrole{doc}{thermo\_style}}}} command) with ID = \sphinxstyleemphasis{thermo\_temp}.
This means you can change the attributes of this fix’s temperature
(e.g. its degrees\sphinxhyphen{}of\sphinxhyphen{}freedom) via the
{\hyperref[\detokenize{compute_modify::doc}]{\sphinxcrossref{\DUrole{doc}{compute\_modify}}}} command or print this temperature
during thermodynamic output via the {\hyperref[\detokenize{thermo_style::doc}]{\sphinxcrossref{\DUrole{doc}{thermo\_style custom}}}} command using the appropriate compute\sphinxhyphen{}ID.
It also means that changing attributes of \sphinxstyleemphasis{thermo\_temp} will have no
effect on this fix.

\sphinxAtStartPar
Like other fixes that perform thermostatting, this fix can be used
with {\hyperref[\detokenize{compute::doc}]{\sphinxcrossref{\DUrole{doc}{compute commands}}}} that remove a “bias” from the
atom velocities.  E.g. to apply the thermostat only to atoms within a
spatial {\hyperref[\detokenize{region::doc}]{\sphinxcrossref{\DUrole{doc}{region}}}}, or to remove the center\sphinxhyphen{}of\sphinxhyphen{}mass
velocity from a group of atoms, or to remove the x\sphinxhyphen{}component of
velocity from the calculation.

\sphinxAtStartPar
This is not done by default, but only if the {\hyperref[\detokenize{fix_modify::doc}]{\sphinxcrossref{\DUrole{doc}{fix\_modify}}}} command is used to assign a temperature compute to this
fix that includes such a bias term.  See the doc pages for individual
{\hyperref[\detokenize{compute::doc}]{\sphinxcrossref{\DUrole{doc}{compute temp commands}}}} to determine which ones include
a bias.  In this case, the thermostat works in the following manner:
bias is removed from each atom, thermostatting is performed on the
remaining thermal degrees of freedom, and the bias is added back in.


\bigskip\hrule\bigskip


\sphinxAtStartPar
Styles with a \sphinxstyleemphasis{gpu}, \sphinxstyleemphasis{intel}, \sphinxstyleemphasis{kk}, \sphinxstyleemphasis{omp}, or \sphinxstyleemphasis{opt} suffix are
functionally the same as the corresponding style without the suffix.
They have been optimized to run faster, depending on your available
hardware, as discussed on the {\hyperref[\detokenize{Speed_packages::doc}]{\sphinxcrossref{\DUrole{doc}{Accelerator packages}}}}
page.  The accelerated styles take the same arguments and should
produce the same results, except for round\sphinxhyphen{}off and precision issues.

\sphinxAtStartPar
These accelerated styles are part of the GPU, INTEL, KOKKOS,
OPENMP, and OPT packages, respectively.  They are only enabled if
LAMMPS was built with those packages.  See the {\hyperref[\detokenize{Build_package::doc}]{\sphinxcrossref{\DUrole{doc}{Build package}}}} page for more info.

\sphinxAtStartPar
You can specify the accelerated styles explicitly in your input script
by including their suffix, or you can use the {\hyperref[\detokenize{Run_options::doc}]{\sphinxcrossref{\DUrole{doc}{\sphinxhyphen{}suffix command\sphinxhyphen{}line switch}}}} when you invoke LAMMPS, or you can use the
{\hyperref[\detokenize{suffix::doc}]{\sphinxcrossref{\DUrole{doc}{suffix}}}} command in your input script.

\sphinxAtStartPar
See the {\hyperref[\detokenize{Speed_packages::doc}]{\sphinxcrossref{\DUrole{doc}{Accelerator packages}}}} page for more
instructions on how to use the accelerated styles effectively.


\bigskip\hrule\bigskip



\subsection{Restart, fix\_modify, output, run start/stop, minimize info}
\label{\detokenize{fix_temp_berendsen:restart-fix-modify-output-run-start-stop-minimize-info}}
\sphinxAtStartPar
This fix writes the cumulative global energy change to
{\hyperref[\detokenize{restart::doc}]{\sphinxcrossref{\DUrole{doc}{binary restart files}}}}.  See the
{\hyperref[\detokenize{read_restart::doc}]{\sphinxcrossref{\DUrole{doc}{read\_restart}}}} command for info on how to
re\sphinxhyphen{}specify a fix in an input script that reads a restart file,
so that the fix continues in an uninterrupted fashion.

\sphinxAtStartPar
The {\hyperref[\detokenize{fix_modify::doc}]{\sphinxcrossref{\DUrole{doc}{fix\_modify}}}} \sphinxstyleemphasis{temp} option is supported by this
fix.  You can use it to assign a temperature {\hyperref[\detokenize{compute::doc}]{\sphinxcrossref{\DUrole{doc}{compute}}}}
you have defined to this fix which will be used in its thermostatting
procedure, as described above.  For consistency, the group used by
this fix and by the compute should be the same.

\sphinxAtStartPar
The cumulative energy change in the system imposed by this fix is
included in the {\hyperref[\detokenize{thermo_style::doc}]{\sphinxcrossref{\DUrole{doc}{thermodynamic output}}}} keywords
\sphinxstyleemphasis{ecouple} and \sphinxstyleemphasis{econserve}.  See the {\hyperref[\detokenize{thermo_style::doc}]{\sphinxcrossref{\DUrole{doc}{thermo\_style}}}}
doc page for details.

\sphinxAtStartPar
This fix computes a global scalar which can be accessed by various
{\hyperref[\detokenize{Howto_output::doc}]{\sphinxcrossref{\DUrole{doc}{output commands}}}}.  The scalar is the same
cumulative energy change due to this fix described in the previous
paragraph.  The scalar value calculated by this fix is “extensive”.

\sphinxAtStartPar
This fix can ramp its target temperature over multiple runs, using the
\sphinxstyleemphasis{start} and \sphinxstyleemphasis{stop} keywords of the {\hyperref[\detokenize{run::doc}]{\sphinxcrossref{\DUrole{doc}{run}}}} command.  See the
{\hyperref[\detokenize{run::doc}]{\sphinxcrossref{\DUrole{doc}{run}}}} command for details of how to do this.

\sphinxAtStartPar
This fix is not invoked during {\hyperref[\detokenize{minimize::doc}]{\sphinxcrossref{\DUrole{doc}{energy minimization}}}}.


\subsection{Restrictions}
\label{\detokenize{fix_temp_berendsen:restrictions}}
\sphinxAtStartPar
This fix can be used with dynamic groups as defined by the
{\hyperref[\detokenize{group::doc}]{\sphinxcrossref{\DUrole{doc}{group}}}} command.  Likewise it can be used with groups to
which atoms are added or deleted over time, e.g. a deposition
simulation.  However, the conservation properties of the thermostat
and barostat are defined for systems with a static set of atoms.  You
may observe odd behavior if the atoms in a group vary dramatically
over time or the atom count becomes very small.


\subsection{Related commands}
\label{\detokenize{fix_temp_berendsen:related-commands}}
\sphinxAtStartPar
{\hyperref[\detokenize{fix_nve::doc}]{\sphinxcrossref{\DUrole{doc}{fix nve}}}}, {\hyperref[\detokenize{fix_nh::doc}]{\sphinxcrossref{\DUrole{doc}{fix nvt}}}}, {\hyperref[\detokenize{fix_temp_rescale::doc}]{\sphinxcrossref{\DUrole{doc}{fix temp/rescale}}}}, {\hyperref[\detokenize{fix_langevin::doc}]{\sphinxcrossref{\DUrole{doc}{fix langevin}}}},
{\hyperref[\detokenize{fix_modify::doc}]{\sphinxcrossref{\DUrole{doc}{fix\_modify}}}}, {\hyperref[\detokenize{compute_temp::doc}]{\sphinxcrossref{\DUrole{doc}{compute temp}}}},
{\hyperref[\detokenize{fix_press_berendsen::doc}]{\sphinxcrossref{\DUrole{doc}{fix press/berendsen}}}}


\subsection{Default}
\label{\detokenize{fix_temp_berendsen:default}}
\sphinxAtStartPar
none


\bigskip\hrule\bigskip

\phantomsection\label{\detokenize{fix_temp_berendsen:berendsen2}}
\sphinxAtStartPar
\sphinxstylestrong{(Berendsen)} Berendsen, Postma, van Gunsteren, DiNola, Haak, J Chem
Phys, 81, 3684 (1984).

\sphinxstepscope

\index{fix temp/csvr@\spxentry{fix temp/csvr}}\ignorespaces 
\index{fix temp/csld@\spxentry{fix temp/csld}}\ignorespaces 

\section{fix temp/csvr command}
\label{\detokenize{fix_temp_csvr:fix-temp-csvr-command}}\label{\detokenize{fix_temp_csvr:index-1}}\label{\detokenize{fix_temp_csvr:index-0}}\label{\detokenize{fix_temp_csvr::doc}}

\section{fix temp/csld command}
\label{\detokenize{fix_temp_csvr:fix-temp-csld-command}}

\subsection{Syntax}
\label{\detokenize{fix_temp_csvr:syntax}}
\begin{sphinxVerbatim}[commandchars=\\\{\}]
\PYG{k}{fix }\PYG{n+nv+nvIdentifier}{ID}\PYG{+w}{ }\PYG{n+nv+nvIdentifier}{group\PYGZhy{}ID}\PYG{+w}{ }\PYG{n}{temp}\PYG{o}{/}\PYG{n}{csvr}\PYG{+w}{ }\PYG{n}{Tstart}\PYG{+w}{ }\PYG{n}{Tstop}\PYG{+w}{ }\PYG{n}{Tdamp}\PYG{+w}{ }\PYG{n}{seed}
\PYG{k}{fix }\PYG{n+nv+nvIdentifier}{ID}\PYG{+w}{ }\PYG{n+nv+nvIdentifier}{group\PYGZhy{}ID}\PYG{+w}{ }\PYG{n}{temp}\PYG{o}{/}\PYG{n}{csld}\PYG{+w}{ }\PYG{n}{Tstart}\PYG{+w}{ }\PYG{n}{Tstop}\PYG{+w}{ }\PYG{n}{Tdamp}\PYG{+w}{ }\PYG{n}{seed}
\end{sphinxVerbatim}
\begin{itemize}
\item {} 
\sphinxAtStartPar
ID, group\sphinxhyphen{}ID are documented in {\hyperref[\detokenize{fix::doc}]{\sphinxcrossref{\DUrole{doc}{fix}}}} command

\item {} 
\sphinxAtStartPar
temp/csvr or temp/csld = style name of this fix command

\item {} 
\sphinxAtStartPar
Tstart,Tstop = desired temperature at start/end of run

\begin{sphinxVerbatim}[commandchars=\\\{\}]
Tstart can be a variable (see below)
\end{sphinxVerbatim}

\item {} 
\sphinxAtStartPar
Tdamp = temperature damping parameter (time units)

\item {} 
\sphinxAtStartPar
seed = random number seed to use for white noise (positive integer)

\end{itemize}


\subsection{Examples}
\label{\detokenize{fix_temp_csvr:examples}}
\begin{sphinxVerbatim}[commandchars=\\\{\}]
\PYG{k}{fix }\PYG{n+nv+nvIdentifier}{1}\PYG{+w}{ }\PYG{n+nv+nvIdentifier}{all}\PYG{+w}{ }\PYG{n}{temp}\PYG{o}{/}\PYG{n}{csvr}\PYG{+w}{ }\PYG{l+m}{300.0}\PYG{+w}{ }\PYG{l+m}{300.0}\PYG{+w}{ }\PYG{l+m}{100.0}\PYG{+w}{ }\PYG{l+m}{54324}
\PYG{k}{fix }\PYG{n+nv+nvIdentifier}{1}\PYG{+w}{ }\PYG{n+nv+nvIdentifier}{all}\PYG{+w}{ }\PYG{n}{temp}\PYG{o}{/}\PYG{n}{csld}\PYG{+w}{ }\PYG{l+m}{100.0}\PYG{+w}{ }\PYG{l+m}{300.0}\PYG{+w}{ }\PYG{l+m}{10.0}\PYG{+w}{ }\PYG{l+m}{123321}
\end{sphinxVerbatim}


\subsection{Description}
\label{\detokenize{fix_temp_csvr:description}}
\sphinxAtStartPar
Adjust the temperature with a canonical sampling thermostat that uses
global velocity rescaling with Hamiltonian dynamics (\sphinxstyleemphasis{temp/csvr})
{\hyperref[\detokenize{fix_temp_csvr:bussi1}]{\sphinxcrossref{\DUrole{std,std-ref}{(Bussi1)}}}}, or Langevin dynamics (\sphinxstyleemphasis{temp/csld})
{\hyperref[\detokenize{fix_temp_csvr:bussi2}]{\sphinxcrossref{\DUrole{std,std-ref}{(Bussi2)}}}}.  In the case of \sphinxstyleemphasis{temp/csvr} the thermostat is
similar to the empirical Berendsen thermostat in
{\hyperref[\detokenize{fix_temp_berendsen::doc}]{\sphinxcrossref{\DUrole{doc}{temp/berendsen}}}}, but chooses the actual
scaling factor from a suitably chosen (gaussian) distribution rather
than having it determined from the time constant directly. In the case
of \sphinxstyleemphasis{temp/csld} the velocities are updated to a linear combination of
the current velocities with a gaussian distribution of velocities at
the desired temperature.  Both thermostats are applied every timestep.

\sphinxAtStartPar
The thermostat is applied to only the translational degrees of freedom
for the particles, which is an important consideration for finite\sphinxhyphen{}size
particles which have rotational degrees of freedom are being
thermostatted with these fixes.  The translational degrees of freedom
can also have a bias velocity removed from them before thermostatting
takes place; see the description below.

\sphinxAtStartPar
The desired temperature at each timestep is a ramped value during the
run from \sphinxstyleemphasis{Tstart} to \sphinxstyleemphasis{Tstop}.  The \sphinxstyleemphasis{Tdamp} parameter is specified in
time units and determines how rapidly the temperature is relaxed.  For
example, a value of 100.0 means to relax the temperature in a timespan
of (roughly) 100 time units (tau or fs or ps \sphinxhyphen{} see the
{\hyperref[\detokenize{units::doc}]{\sphinxcrossref{\DUrole{doc}{units}}}} command).

\sphinxAtStartPar
\sphinxstyleemphasis{Tstart} can be specified as an equal\sphinxhyphen{}style {\hyperref[\detokenize{variable::doc}]{\sphinxcrossref{\DUrole{doc}{variable}}}}.
In this case, the \sphinxstyleemphasis{Tstop} setting is ignored.  If the value is a
variable, it should be specified as v\_name, where name is the variable
name.  In this case, the variable will be evaluated each timestep, and
its value used to determine the target temperature.

\sphinxAtStartPar
Equal\sphinxhyphen{}style variables can specify formulas with various mathematical
functions, and include {\hyperref[\detokenize{thermo_style::doc}]{\sphinxcrossref{\DUrole{doc}{thermo\_style}}}} command
keywords for the simulation box parameters and timestep and elapsed
time.  Thus it is easy to specify a time\sphinxhyphen{}dependent temperature.

\begin{sphinxadmonition}{note}{Note:}
\sphinxAtStartPar
Unlike the {\hyperref[\detokenize{fix_nh::doc}]{\sphinxcrossref{\DUrole{doc}{fix nvt}}}} command which performs
Nose/Hoover thermostatting AND time integration, these fixes do NOT
perform time integration. They only modify velocities to effect
thermostatting.  Thus you must use a separate time integration fix,
like {\hyperref[\detokenize{fix_nve::doc}]{\sphinxcrossref{\DUrole{doc}{fix nve}}}} to actually update the positions of atoms
using the modified velocities.  Likewise, these fixes should not
normally be used on atoms that also have their temperature controlled
by another fix \sphinxhyphen{} e.g. by {\hyperref[\detokenize{fix_nh::doc}]{\sphinxcrossref{\DUrole{doc}{fix nvt}}}} or {\hyperref[\detokenize{fix_langevin::doc}]{\sphinxcrossref{\DUrole{doc}{fix langevin}}}} commands.
\end{sphinxadmonition}

\sphinxAtStartPar
See the {\hyperref[\detokenize{Howto_thermostat::doc}]{\sphinxcrossref{\DUrole{doc}{Howto thermostat}}}} page for a
discussion of different ways to compute temperature and perform
thermostatting.

\sphinxAtStartPar
These fixes compute a temperature each timestep.  To do this, the fix
creates its own compute of style “temp”, as if this command had been
issued:

\begin{sphinxVerbatim}[commandchars=\\\{\}]
\PYG{k}{compute }\PYG{n+nv+nvIdentifier}{fix\PYGZhy{}ID\PYGZus{}temp}\PYG{+w}{ }\PYG{n+nv+nvIdentifier}{group\PYGZhy{}ID}\PYG{+w}{ }\PYG{n}{temp}
\end{sphinxVerbatim}

\sphinxAtStartPar
See the {\hyperref[\detokenize{compute_temp::doc}]{\sphinxcrossref{\DUrole{doc}{compute temp}}}} command for details.  Note
that the ID of the new compute is the fix\sphinxhyphen{}ID + underscore + “temp”,
and the group for the new compute is the same as the fix group.

\sphinxAtStartPar
Note that this is NOT the compute used by thermodynamic output (see
the {\hyperref[\detokenize{thermo_style::doc}]{\sphinxcrossref{\DUrole{doc}{thermo\_style}}}} command) with ID = \sphinxstyleemphasis{thermo\_temp}.
This means you can change the attributes of this fix’s temperature
(e.g. its degrees\sphinxhyphen{}of\sphinxhyphen{}freedom) via the
{\hyperref[\detokenize{compute_modify::doc}]{\sphinxcrossref{\DUrole{doc}{compute\_modify}}}} command or print this temperature
during thermodynamic output via the {\hyperref[\detokenize{thermo_style::doc}]{\sphinxcrossref{\DUrole{doc}{thermo\_style custom}}}} command using the appropriate compute\sphinxhyphen{}ID.
It also means that changing attributes of \sphinxstyleemphasis{thermo\_temp} will have no
effect on this fix.

\sphinxAtStartPar
Like other fixes that perform thermostatting, this fix can be used
with {\hyperref[\detokenize{compute::doc}]{\sphinxcrossref{\DUrole{doc}{compute commands}}}} that remove a “bias” from the
atom velocities.  E.g. to apply the thermostat only to atoms within a
spatial {\hyperref[\detokenize{region::doc}]{\sphinxcrossref{\DUrole{doc}{region}}}}, or to remove the center\sphinxhyphen{}of\sphinxhyphen{}mass
velocity from a group of atoms, or to remove the x\sphinxhyphen{}component of
velocity from the calculation.

\sphinxAtStartPar
This is not done by default, but only if the {\hyperref[\detokenize{fix_modify::doc}]{\sphinxcrossref{\DUrole{doc}{fix\_modify}}}} command is used to assign a temperature compute to this
fix that includes such a bias term.  See the doc pages for individual
{\hyperref[\detokenize{compute::doc}]{\sphinxcrossref{\DUrole{doc}{compute temp commands}}}} to determine which ones include
a bias.  In this case, the thermostat works in the following manner:
bias is removed from each atom, thermostatting is performed on the
remaining thermal degrees of freedom, and the bias is added back in.

\sphinxAtStartPar
An important feature of these thermostats is that they have an
associated effective energy that is a constant of motion.  The
effective energy is the total energy (kinetic + potential) plus the
accumulated kinetic energy changes due to the thermostat. The latter
quantity is the global scalar computed by these fixes. This feature is
useful to check the integration of the equations of motion against
discretization errors. In other words, the conservation of the
effective energy can be used to choose an appropriate integration
{\hyperref[\detokenize{timestep::doc}]{\sphinxcrossref{\DUrole{doc}{timestep}}}}. This is similar to the usual paradigm of
checking the conservation of the total energy in the microcanonical
ensemble.


\bigskip\hrule\bigskip



\subsection{Restart, fix\_modify, output, run start/stop, minimize info}
\label{\detokenize{fix_temp_csvr:restart-fix-modify-output-run-start-stop-minimize-info}}
\sphinxAtStartPar
These fixes write the cumulative global energy change and the random
number generator states to {\hyperref[\detokenize{restart::doc}]{\sphinxcrossref{\DUrole{doc}{binary restart files}}}}.  See
the {\hyperref[\detokenize{read_restart::doc}]{\sphinxcrossref{\DUrole{doc}{read\_restart}}}} command for info on how to
re\sphinxhyphen{}specify a fix in an input script that reads a restart file, so that
the selected fix continues in an uninterrupted fashion.  The random
number generator state can only be restored when the number of
processors remains unchanged from what is recorded in the restart
file.

\sphinxAtStartPar
The {\hyperref[\detokenize{fix_modify::doc}]{\sphinxcrossref{\DUrole{doc}{fix\_modify}}}} \sphinxstyleemphasis{temp} option is supported by these
fixes.  You can use it to assign a temperature {\hyperref[\detokenize{compute::doc}]{\sphinxcrossref{\DUrole{doc}{compute}}}} you have defined to these fixes which will be used in its
thermostatting procedure, as described above.  For consistency, the
group used by these fixes and by the compute should be the same.

\sphinxAtStartPar
The cumulative energy change in the system imposed by these fixes is
included in the {\hyperref[\detokenize{thermo_style::doc}]{\sphinxcrossref{\DUrole{doc}{thermodynamic output}}}} keywords
\sphinxstyleemphasis{ecouple} and \sphinxstyleemphasis{econserve}.  See the {\hyperref[\detokenize{thermo_style::doc}]{\sphinxcrossref{\DUrole{doc}{thermo\_style}}}}
doc page for details.

\sphinxAtStartPar
These fixes compute a global scalar which can be accessed by various
{\hyperref[\detokenize{Howto_output::doc}]{\sphinxcrossref{\DUrole{doc}{output commands}}}}.  The scalar is the same
cumulative energy change due to this fix described in the previous
paragraph.  The scalar value calculated by this fix is “extensive”.

\sphinxAtStartPar
These fixes can ramp their target temperature over multiple runs,
using the \sphinxstyleemphasis{start} and \sphinxstyleemphasis{stop} keywords of the {\hyperref[\detokenize{run::doc}]{\sphinxcrossref{\DUrole{doc}{run}}}} command.
See the {\hyperref[\detokenize{run::doc}]{\sphinxcrossref{\DUrole{doc}{run}}}} command for details of how to do this.

\sphinxAtStartPar
These fixes are not invoked during {\hyperref[\detokenize{minimize::doc}]{\sphinxcrossref{\DUrole{doc}{energy minimization}}}}.


\subsection{Restrictions}
\label{\detokenize{fix_temp_csvr:restrictions}}
\sphinxAtStartPar
Fix \sphinxstyleemphasis{temp/csld} is not compatible with {\hyperref[\detokenize{fix_shake::doc}]{\sphinxcrossref{\DUrole{doc}{fix shake}}}}.

\sphinxAtStartPar
These fixes are part of the EXTRA\sphinxhyphen{}FIX package.  They are only enabled if LAMMPS
was built with that package.  See the {\hyperref[\detokenize{Build_package::doc}]{\sphinxcrossref{\DUrole{doc}{Build package}}}} page for more info.

\sphinxAtStartPar
These fixes can be used with dynamic groups as defined by the
{\hyperref[\detokenize{group::doc}]{\sphinxcrossref{\DUrole{doc}{group}}}} command.  Likewise it can be used with groups to
which atoms are added or deleted over time, e.g. a deposition
simulation.  However, the conservation properties of the thermostat
and barostat are defined for systems with a static set of atoms.  You
may observe odd behavior if the atoms in a group vary dramatically
over time or the atom count becomes very small.


\subsection{Related commands}
\label{\detokenize{fix_temp_csvr:related-commands}}
\sphinxAtStartPar
{\hyperref[\detokenize{fix_nve::doc}]{\sphinxcrossref{\DUrole{doc}{fix nve}}}}, {\hyperref[\detokenize{fix_nh::doc}]{\sphinxcrossref{\DUrole{doc}{fix nvt}}}}, {\hyperref[\detokenize{fix_temp_rescale::doc}]{\sphinxcrossref{\DUrole{doc}{fix temp/rescale}}}}, {\hyperref[\detokenize{fix_langevin::doc}]{\sphinxcrossref{\DUrole{doc}{fix langevin}}}},
{\hyperref[\detokenize{fix_modify::doc}]{\sphinxcrossref{\DUrole{doc}{fix\_modify}}}}, {\hyperref[\detokenize{compute_temp::doc}]{\sphinxcrossref{\DUrole{doc}{compute temp}}}},
{\hyperref[\detokenize{fix_temp_berendsen::doc}]{\sphinxcrossref{\DUrole{doc}{fix temp/berendsen}}}}


\subsection{Default}
\label{\detokenize{fix_temp_csvr:default}}
\sphinxAtStartPar
none


\bigskip\hrule\bigskip

\phantomsection\label{\detokenize{fix_temp_csvr:bussi1}}\phantomsection\label{\detokenize{fix_temp_csvr:bussi2}}
\sphinxAtStartPar
\sphinxstylestrong{(Bussi1)} Bussi, Donadio and Parrinello, J. Chem. Phys. 126, 014101(2007)

\sphinxAtStartPar
\sphinxstylestrong{(Bussi2)} Bussi and Parrinello, Phys. Rev. E 75, 056707 (2007)

\sphinxstepscope

\index{fix temp/rescale@\spxentry{fix temp/rescale}}\ignorespaces 
\index{fix temp/rescale/kk@\spxentry{fix temp/rescale/kk}}\ignorespaces 

\section{fix temp/rescale command}
\label{\detokenize{fix_temp_rescale:fix-temp-rescale-command}}\label{\detokenize{fix_temp_rescale:index-1}}\label{\detokenize{fix_temp_rescale:index-0}}\label{\detokenize{fix_temp_rescale::doc}}
\sphinxAtStartPar
Accelerator Variants: \sphinxstyleemphasis{temp/rescale/kk}


\subsection{Syntax}
\label{\detokenize{fix_temp_rescale:syntax}}
\begin{sphinxVerbatim}[commandchars=\\\{\}]
\PYG{k}{fix }\PYG{n+nv+nvIdentifier}{ID}\PYG{+w}{ }\PYG{n+nv+nvIdentifier}{group\PYGZhy{}ID}\PYG{+w}{ }\PYG{n}{temp}\PYG{o}{/}\PYG{n}{rescale}\PYG{+w}{ }\PYG{n}{N}\PYG{+w}{ }\PYG{n}{Tstart}\PYG{+w}{ }\PYG{n}{Tstop}\PYG{+w}{ }\PYG{n}{window}\PYG{+w}{ }\PYG{n}{fraction}
\end{sphinxVerbatim}
\begin{itemize}
\item {} 
\sphinxAtStartPar
ID, group\sphinxhyphen{}ID are documented in {\hyperref[\detokenize{fix::doc}]{\sphinxcrossref{\DUrole{doc}{fix}}}} command

\item {} 
\sphinxAtStartPar
temp/rescale = style name of this fix command

\item {} 
\sphinxAtStartPar
N = perform rescaling every N steps

\item {} 
\sphinxAtStartPar
Tstart,Tstop = desired temperature at start/end of run (temperature units)

\begin{sphinxVerbatim}[commandchars=\\\{\}]
Tstart can be a variable (see below)
\end{sphinxVerbatim}

\item {} 
\sphinxAtStartPar
window = only rescale if temperature is outside this window (temperature units)

\item {} 
\sphinxAtStartPar
fraction = rescale to target temperature by this fraction

\end{itemize}


\subsection{Examples}
\label{\detokenize{fix_temp_rescale:examples}}
\begin{sphinxVerbatim}[commandchars=\\\{\}]
\PYG{k}{fix }\PYG{n+nv+nvIdentifier}{3}\PYG{+w}{ }\PYG{n+nv+nvIdentifier}{flow}\PYG{+w}{ }\PYG{n}{temp}\PYG{o}{/}\PYG{n}{rescale}\PYG{+w}{ }\PYG{l+m}{100}\PYG{+w}{ }\PYG{l+m}{1.0}\PYG{+w}{ }\PYG{l+m}{1.1}\PYG{+w}{ }\PYG{l+m}{0.02}\PYG{+w}{ }\PYG{l+m}{0.5}
\PYG{k}{fix }\PYG{n+nv+nvIdentifier}{3}\PYG{+w}{ }\PYG{n+nv+nvIdentifier}{boundary}\PYG{+w}{ }\PYG{n}{temp}\PYG{o}{/}\PYG{n}{rescale}\PYG{+w}{ }\PYG{l+m}{1}\PYG{+w}{ }\PYG{l+m}{1.0}\PYG{+w}{ }\PYG{l+m}{1.5}\PYG{+w}{ }\PYG{l+m}{0.05}\PYG{+w}{ }\PYG{l+m}{1.0}
\PYG{k}{fix }\PYG{n+nv+nvIdentifier}{3}\PYG{+w}{ }\PYG{n+nv+nvIdentifier}{boundary}\PYG{+w}{ }\PYG{n}{temp}\PYG{o}{/}\PYG{n}{rescale}\PYG{+w}{ }\PYG{l+m}{1}\PYG{+w}{ }\PYG{l+m}{1.0}\PYG{+w}{ }\PYG{l+m}{1.5}\PYG{+w}{ }\PYG{l+m}{0.05}\PYG{+w}{ }\PYG{l+m}{1.0}
\end{sphinxVerbatim}


\subsection{Description}
\label{\detokenize{fix_temp_rescale:description}}
\sphinxAtStartPar
Reset the temperature of a group of atoms by explicitly rescaling
their velocities.

\sphinxAtStartPar
The rescaling is applied to only the translational degrees of freedom
for the particles, which is an important consideration if finite\sphinxhyphen{}size
particles which have rotational degrees of freedom are being
thermostatted with this fix.  The translational degrees of freedom can
also have a bias velocity removed from them before thermostatting
takes place; see the description below.

\sphinxAtStartPar
Rescaling is performed every N timesteps.  The target temperature is a
ramped value between the \sphinxstyleemphasis{Tstart} and \sphinxstyleemphasis{Tstop} temperatures at the
beginning and end of the run.

\begin{sphinxadmonition}{note}{Note:}
\sphinxAtStartPar
This thermostat will generate an error if the current
temperature is zero at the end of a timestep it is invoked on.  It
cannot rescale a zero temperature.
\end{sphinxadmonition}

\sphinxAtStartPar
\sphinxstyleemphasis{Tstart} can be specified as an equal\sphinxhyphen{}style {\hyperref[\detokenize{variable::doc}]{\sphinxcrossref{\DUrole{doc}{variable}}}}.
In this case, the \sphinxstyleemphasis{Tstop} setting is ignored.  If the value is a
variable, it should be specified as v\_name, where name is the variable
name.  In this case, the variable will be evaluated each timestep, and
its value used to determine the target temperature.

\sphinxAtStartPar
Equal\sphinxhyphen{}style variables can specify formulas with various mathematical
functions, and include {\hyperref[\detokenize{thermo_style::doc}]{\sphinxcrossref{\DUrole{doc}{thermo\_style}}}} command
keywords for the simulation box parameters and timestep and elapsed
time.  Thus it is easy to specify a time\sphinxhyphen{}dependent temperature.

\sphinxAtStartPar
Rescaling is only performed if the difference between the current and
desired temperatures is greater than the \sphinxstyleemphasis{window} value.  The amount
of rescaling that is applied is a \sphinxstyleemphasis{fraction} (from 0.0 to 1.0) of the
difference between the actual and desired temperature.  E.g. if
\sphinxstyleemphasis{fraction} = 1.0, the temperature is reset to exactly the desired
value.

\begin{sphinxadmonition}{note}{Note:}
\sphinxAtStartPar
Unlike the {\hyperref[\detokenize{fix_nh::doc}]{\sphinxcrossref{\DUrole{doc}{fix nvt}}}} command which performs
Nose/Hoover thermostatting AND time integration, this fix does NOT
perform time integration.  It only modifies velocities to effect
thermostatting.  Thus you must use a separate time integration fix,
like {\hyperref[\detokenize{fix_nve::doc}]{\sphinxcrossref{\DUrole{doc}{fix nve}}}} to actually update the positions of atoms
using the modified velocities.  Likewise, this fix should not normally
be used on atoms that also have their temperature controlled by
another fix \sphinxhyphen{} e.g. by {\hyperref[\detokenize{fix_nh::doc}]{\sphinxcrossref{\DUrole{doc}{fix nvt}}}} or {\hyperref[\detokenize{fix_langevin::doc}]{\sphinxcrossref{\DUrole{doc}{fix langevin}}}} commands.
\end{sphinxadmonition}

\sphinxAtStartPar
See the {\hyperref[\detokenize{Howto_thermostat::doc}]{\sphinxcrossref{\DUrole{doc}{Howto thermostat}}}} page for a
discussion of different ways to compute temperature and perform
thermostatting.

\sphinxAtStartPar
This fix computes a temperature each timestep.  To do this, the fix
creates its own compute of style “temp”, as if one of this command had
been issued:

\begin{sphinxVerbatim}[commandchars=\\\{\}]
\PYG{k}{compute }\PYG{n+nv+nvIdentifier}{fix\PYGZhy{}ID\PYGZus{}temp}\PYG{+w}{ }\PYG{n+nv+nvIdentifier}{group\PYGZhy{}ID}\PYG{+w}{ }\PYG{n}{temp}
\end{sphinxVerbatim}

\sphinxAtStartPar
See the {\hyperref[\detokenize{compute_temp::doc}]{\sphinxcrossref{\DUrole{doc}{compute temp}}}} for details.  Note that the
ID of the new compute is the fix\sphinxhyphen{}ID + underscore + “temp”, and the
group for the new compute is the same as the fix group.

\sphinxAtStartPar
Note that this is NOT the compute used by thermodynamic output (see
the {\hyperref[\detokenize{thermo_style::doc}]{\sphinxcrossref{\DUrole{doc}{thermo\_style}}}} command) with ID =
\sphinxstyleemphasis{thermo\_temp}.  This means you can change the attributes of this fix’s
temperature (e.g. its degrees\sphinxhyphen{}of\sphinxhyphen{}freedom) via the {\hyperref[\detokenize{compute_modify::doc}]{\sphinxcrossref{\DUrole{doc}{compute\_modify}}}} command or print this temperature during
thermodynamic output via the {\hyperref[\detokenize{thermo_style::doc}]{\sphinxcrossref{\DUrole{doc}{thermo\_style custom}}}}
command using the appropriate compute\sphinxhyphen{}ID.  It also means that changing
attributes of \sphinxstyleemphasis{thermo\_temp} will have no effect on this fix.

\sphinxAtStartPar
Like other fixes that perform thermostatting, this fix can be used
with {\hyperref[\detokenize{compute::doc}]{\sphinxcrossref{\DUrole{doc}{compute commands}}}} that remove a “bias” from the
atom velocities.  E.g. to apply the thermostat only to atoms within a
spatial {\hyperref[\detokenize{region::doc}]{\sphinxcrossref{\DUrole{doc}{region}}}}, or to remove the center\sphinxhyphen{}of\sphinxhyphen{}mass
velocity from a group of atoms, or to remove the x\sphinxhyphen{}component of
velocity from the calculation.

\sphinxAtStartPar
This is not done by default, but only if the {\hyperref[\detokenize{fix_modify::doc}]{\sphinxcrossref{\DUrole{doc}{fix\_modify}}}} command is used to assign a temperature compute to this
fix that includes such a bias term.  See the doc pages for individual
{\hyperref[\detokenize{compute::doc}]{\sphinxcrossref{\DUrole{doc}{compute temp commands}}}} to determine which ones include
a bias.  In this case, the thermostat works in the following manner:
bias is removed from each atom, thermostatting is performed on the
remaining thermal degrees of freedom, and the bias is added back in.


\bigskip\hrule\bigskip


\sphinxAtStartPar
Styles with a \sphinxstyleemphasis{gpu}, \sphinxstyleemphasis{intel}, \sphinxstyleemphasis{kk}, \sphinxstyleemphasis{omp}, or \sphinxstyleemphasis{opt} suffix are
functionally the same as the corresponding style without the suffix.
They have been optimized to run faster, depending on your available
hardware, as discussed on the {\hyperref[\detokenize{Speed_packages::doc}]{\sphinxcrossref{\DUrole{doc}{Accelerator packages}}}}
page.  The accelerated styles take the same arguments and should
produce the same results, except for round\sphinxhyphen{}off and precision issues.

\sphinxAtStartPar
These accelerated styles are part of the GPU, INTEL, KOKKOS,
OPENMP, and OPT packages, respectively.  They are only enabled if
LAMMPS was built with those packages.  See the {\hyperref[\detokenize{Build_package::doc}]{\sphinxcrossref{\DUrole{doc}{Build package}}}} page for more info.

\sphinxAtStartPar
You can specify the accelerated styles explicitly in your input script
by including their suffix, or you can use the {\hyperref[\detokenize{Run_options::doc}]{\sphinxcrossref{\DUrole{doc}{\sphinxhyphen{}suffix command\sphinxhyphen{}line switch}}}} when you invoke LAMMPS, or you can use the
{\hyperref[\detokenize{suffix::doc}]{\sphinxcrossref{\DUrole{doc}{suffix}}}} command in your input script.

\sphinxAtStartPar
See the {\hyperref[\detokenize{Speed_packages::doc}]{\sphinxcrossref{\DUrole{doc}{Accelerator packages}}}} page for more
instructions on how to use the accelerated styles effectively.


\bigskip\hrule\bigskip



\subsection{Restart, fix\_modify, output, run start/stop, minimize info}
\label{\detokenize{fix_temp_rescale:restart-fix-modify-output-run-start-stop-minimize-info}}
\sphinxAtStartPar
This fix writes the cumulative global energy change to {\hyperref[\detokenize{restart::doc}]{\sphinxcrossref{\DUrole{doc}{binary
restart files}}}}.  See the {\hyperref[\detokenize{read_restart::doc}]{\sphinxcrossref{\DUrole{doc}{read\_restart}}}}
command for info on how to re\sphinxhyphen{}specify a fix in an input script that
reads a restart file, so that the fix continues in an uninterrupted
fashion.

\sphinxAtStartPar
The {\hyperref[\detokenize{fix_modify::doc}]{\sphinxcrossref{\DUrole{doc}{fix\_modify}}}} \sphinxstyleemphasis{temp} option is supported by this
fix.  You can use it to assign a temperature {\hyperref[\detokenize{compute::doc}]{\sphinxcrossref{\DUrole{doc}{compute}}}}
you have defined to this fix which will be used in its thermostatting
procedure, as described above.  For consistency, the group used by
this fix and by the compute should be the same.

\sphinxAtStartPar
The cumulative energy change in the system imposed by this fix is
included in the {\hyperref[\detokenize{thermo_style::doc}]{\sphinxcrossref{\DUrole{doc}{thermodynamic output}}}} keywords
\sphinxstyleemphasis{ecouple} and \sphinxstyleemphasis{econserve}.  See the {\hyperref[\detokenize{thermo_style::doc}]{\sphinxcrossref{\DUrole{doc}{thermo\_style}}}}
doc page for details.

\sphinxAtStartPar
This fix computes a global scalar which can be accessed by various
{\hyperref[\detokenize{Howto_output::doc}]{\sphinxcrossref{\DUrole{doc}{output commands}}}}.  The scalar is the same
cumulative energy change due to this fix described in the previous
paragraph.  The scalar value calculated by this fix is “extensive”.

\sphinxAtStartPar
This fix can ramp its target temperature over multiple runs, using the
\sphinxstyleemphasis{start} and \sphinxstyleemphasis{stop} keywords of the {\hyperref[\detokenize{run::doc}]{\sphinxcrossref{\DUrole{doc}{run}}}} command.  See the
{\hyperref[\detokenize{run::doc}]{\sphinxcrossref{\DUrole{doc}{run}}}} command for details of how to do this.

\sphinxAtStartPar
This fix is not invoked during {\hyperref[\detokenize{minimize::doc}]{\sphinxcrossref{\DUrole{doc}{energy minimization}}}}.


\subsection{Restrictions}
\label{\detokenize{fix_temp_rescale:restrictions}}\begin{quote}

\sphinxAtStartPar
none
\end{quote}


\subsection{Related commands}
\label{\detokenize{fix_temp_rescale:related-commands}}
\sphinxAtStartPar
{\hyperref[\detokenize{fix_langevin::doc}]{\sphinxcrossref{\DUrole{doc}{fix langevin}}}}, {\hyperref[\detokenize{fix_nh::doc}]{\sphinxcrossref{\DUrole{doc}{fix nvt}}}},
{\hyperref[\detokenize{fix_modify::doc}]{\sphinxcrossref{\DUrole{doc}{fix\_modify}}}}


\subsection{Default}
\label{\detokenize{fix_temp_rescale:default}}
\sphinxAtStartPar
none

\sphinxstepscope

\index{fix temp/rescale/eff@\spxentry{fix temp/rescale/eff}}\ignorespaces 

\section{fix temp/rescale/eff command}
\label{\detokenize{fix_temp_rescale_eff:fix-temp-rescale-eff-command}}\label{\detokenize{fix_temp_rescale_eff:index-0}}\label{\detokenize{fix_temp_rescale_eff::doc}}

\subsection{Syntax}
\label{\detokenize{fix_temp_rescale_eff:syntax}}
\begin{sphinxVerbatim}[commandchars=\\\{\}]
\PYG{k}{fix }\PYG{n+nv+nvIdentifier}{ID}\PYG{+w}{ }\PYG{n+nv+nvIdentifier}{group\PYGZhy{}ID}\PYG{+w}{ }\PYG{n}{temp}\PYG{o}{/}\PYG{n}{rescale}\PYG{o}{/}\PYG{n}{eff}\PYG{+w}{ }\PYG{n}{N}\PYG{+w}{ }\PYG{n}{Tstart}\PYG{+w}{ }\PYG{n}{Tstop}\PYG{+w}{ }\PYG{n}{window}\PYG{+w}{ }\PYG{n}{fraction}
\end{sphinxVerbatim}
\begin{itemize}
\item {} 
\sphinxAtStartPar
ID, group\sphinxhyphen{}ID are documented in {\hyperref[\detokenize{fix::doc}]{\sphinxcrossref{\DUrole{doc}{fix}}}} command

\item {} 
\sphinxAtStartPar
temp/rescale/eff = style name of this fix command

\item {} 
\sphinxAtStartPar
N = perform rescaling every N steps

\item {} 
\sphinxAtStartPar
Tstart,Tstop = desired temperature at start/end of run (temperature units)

\item {} 
\sphinxAtStartPar
window = only rescale if temperature is outside this window (temperature units)

\item {} 
\sphinxAtStartPar
fraction = rescale to target temperature by this fraction

\end{itemize}


\subsection{Examples}
\label{\detokenize{fix_temp_rescale_eff:examples}}
\begin{sphinxVerbatim}[commandchars=\\\{\}]
\PYG{k}{fix }\PYG{n+nv+nvIdentifier}{3}\PYG{+w}{ }\PYG{n+nv+nvIdentifier}{flow}\PYG{+w}{ }\PYG{n}{temp}\PYG{o}{/}\PYG{n}{rescale}\PYG{o}{/}\PYG{n}{eff}\PYG{+w}{ }\PYG{l+m}{10}\PYG{+w}{ }\PYG{l+m}{1.0}\PYG{+w}{ }\PYG{l+m}{100.0}\PYG{+w}{ }\PYG{l+m}{0.02}\PYG{+w}{ }\PYG{l+m}{1.0}
\end{sphinxVerbatim}


\subsection{Description}
\label{\detokenize{fix_temp_rescale_eff:description}}
\sphinxAtStartPar
Reset the temperature of a group of nuclei and electrons in the
{\hyperref[\detokenize{pair_eff::doc}]{\sphinxcrossref{\DUrole{doc}{electron force field}}}} model by explicitly rescaling
their velocities.

\sphinxAtStartPar
The operation of this fix is exactly like that described by the {\hyperref[\detokenize{fix_temp_rescale::doc}]{\sphinxcrossref{\DUrole{doc}{fix temp/rescale}}}} command, except that the rescaling
is also applied to the radial electron velocity for electron
particles.


\subsection{Restart, fix\_modify, output, run start/stop, minimize info}
\label{\detokenize{fix_temp_rescale_eff:restart-fix-modify-output-run-start-stop-minimize-info}}
\sphinxAtStartPar
No information about this fix is written to {\hyperref[\detokenize{restart::doc}]{\sphinxcrossref{\DUrole{doc}{binary restart files}}}}.

\sphinxAtStartPar
The {\hyperref[\detokenize{fix_modify::doc}]{\sphinxcrossref{\DUrole{doc}{fix\_modify}}}} \sphinxstyleemphasis{temp} option is supported by this
fix.  You can use it to assign a temperature {\hyperref[\detokenize{compute::doc}]{\sphinxcrossref{\DUrole{doc}{compute}}}}
you have defined to this fix which will be used in its thermostatting
procedure, as described above.  For consistency, the group used by
this fix and by the compute should be the same.

\sphinxAtStartPar
The cumulative energy change in the system imposed by this fix is
included in the {\hyperref[\detokenize{thermo_style::doc}]{\sphinxcrossref{\DUrole{doc}{thermodynamic output}}}} keywords
\sphinxstyleemphasis{ecouple} and \sphinxstyleemphasis{econserve}.  See the {\hyperref[\detokenize{thermo_style::doc}]{\sphinxcrossref{\DUrole{doc}{thermo\_style}}}}
doc page for details.

\sphinxAtStartPar
This fix computes a global scalar which can be accessed by various
{\hyperref[\detokenize{Howto_output::doc}]{\sphinxcrossref{\DUrole{doc}{output commands}}}}.  The scalar is the same
cumulative energy change due to this fix described in the previous
paragraph.  The scalar value calculated by this fix is “extensive”.

\sphinxAtStartPar
This fix can ramp its target temperature over multiple runs, using the
\sphinxstyleemphasis{start} and \sphinxstyleemphasis{stop} keywords of the {\hyperref[\detokenize{run::doc}]{\sphinxcrossref{\DUrole{doc}{run}}}} command.  See the
{\hyperref[\detokenize{run::doc}]{\sphinxcrossref{\DUrole{doc}{run}}}} command for details of how to do this.

\sphinxAtStartPar
This fix is not invoked during {\hyperref[\detokenize{minimize::doc}]{\sphinxcrossref{\DUrole{doc}{energy minimization}}}}.


\subsection{Restrictions}
\label{\detokenize{fix_temp_rescale_eff:restrictions}}
\sphinxAtStartPar
This fix is part of the EFF package.  It is only enabled if
LAMMPS was built with that package.  See the {\hyperref[\detokenize{Build_package::doc}]{\sphinxcrossref{\DUrole{doc}{Build package}}}} page for more info.


\subsection{Related commands}
\label{\detokenize{fix_temp_rescale_eff:related-commands}}
\sphinxAtStartPar
{\hyperref[\detokenize{fix_langevin_eff::doc}]{\sphinxcrossref{\DUrole{doc}{fix langevin/eff}}}}, {\hyperref[\detokenize{fix_nh_eff::doc}]{\sphinxcrossref{\DUrole{doc}{fix nvt/eff}}}}, {\hyperref[\detokenize{fix_modify::doc}]{\sphinxcrossref{\DUrole{doc}{fix\_modify}}}},
{\hyperref[\detokenize{fix_temp_rescale::doc}]{\sphinxcrossref{\DUrole{doc}{fix temp rescale}}}},


\subsection{Default}
\label{\detokenize{fix_temp_rescale_eff:default}}
\sphinxAtStartPar
none

\sphinxstepscope

\index{fix tfmc@\spxentry{fix tfmc}}\ignorespaces 

\section{fix tfmc command}
\label{\detokenize{fix_tfmc:fix-tfmc-command}}\label{\detokenize{fix_tfmc:index-0}}\label{\detokenize{fix_tfmc::doc}}

\subsection{Syntax}
\label{\detokenize{fix_tfmc:syntax}}
\begin{sphinxVerbatim}[commandchars=\\\{\}]
\PYG{k}{fix }\PYG{n+nv+nvIdentifier}{ID}\PYG{+w}{ }\PYG{n+nv+nvIdentifier}{group\PYGZhy{}ID}\PYG{+w}{ }\PYG{n}{tfmc}\PYG{+w}{ }\PYG{n}{Delta}\PYG{+w}{ }\PYG{n}{Temp}\PYG{+w}{ }\PYG{n}{seed}\PYG{+w}{ }\PYG{n}{keyword}\PYG{+w}{ }\PYG{n}{value}
\end{sphinxVerbatim}
\begin{itemize}
\item {} 
\sphinxAtStartPar
ID, group\sphinxhyphen{}ID are documented in {\hyperref[\detokenize{fix::doc}]{\sphinxcrossref{\DUrole{doc}{fix}}}} command

\item {} 
\sphinxAtStartPar
tfmc = style name of this fix command

\item {} 
\sphinxAtStartPar
Delta = maximal displacement length (distance units)

\item {} 
\sphinxAtStartPar
Temp = imposed temperature of the system

\item {} 
\sphinxAtStartPar
seed = random number seed (positive integer)

\item {} 
\sphinxAtStartPar
zero or more keyword/arg pairs may be appended

\item {} 
\sphinxAtStartPar
keyword = \sphinxstyleemphasis{com} or \sphinxstyleemphasis{rot}
\begin{sphinxalltt}
\sphinxstyleemphasis{com} args = xflag yflag zflag
  xflag,yflag,zflag = 0/1 to exclude/include each dimension
\sphinxstyleemphasis{rot} args = none
\end{sphinxalltt}

\end{itemize}


\subsection{Examples}
\label{\detokenize{fix_tfmc:examples}}
\begin{sphinxVerbatim}[commandchars=\\\{\}]
\PYG{k}{fix }\PYG{n+nv+nvIdentifier}{1}\PYG{+w}{ }\PYG{n+nv+nvIdentifier}{all}\PYG{+w}{ }\PYG{n}{tfmc}\PYG{+w}{ }\PYG{l+m}{0.1}\PYG{+w}{ }\PYG{l+m}{1000.0}\PYG{+w}{ }\PYG{l+m}{159345}
\PYG{k}{fix }\PYG{n+nv+nvIdentifier}{1}\PYG{+w}{ }\PYG{n+nv+nvIdentifier}{all}\PYG{+w}{ }\PYG{n}{tfmc}\PYG{+w}{ }\PYG{l+m}{0.05}\PYG{+w}{ }\PYG{l+m}{600.0}\PYG{+w}{ }\PYG{l+m}{658943}\PYG{+w}{ }\PYG{n}{com}\PYG{+w}{ }\PYG{l+m}{1}\PYG{+w}{ }\PYG{l+m}{1}\PYG{+w}{ }\PYG{l+m}{0}
\PYG{k}{fix }\PYG{n+nv+nvIdentifier}{1}\PYG{+w}{ }\PYG{n+nv+nvIdentifier}{all}\PYG{+w}{ }\PYG{n}{tfmc}\PYG{+w}{ }\PYG{l+m}{0.1}\PYG{+w}{ }\PYG{l+m}{750.0}\PYG{+w}{ }\PYG{l+m}{387068}\PYG{+w}{ }\PYG{n}{com}\PYG{+w}{ }\PYG{l+m}{1}\PYG{+w}{ }\PYG{l+m}{1}\PYG{+w}{ }\PYG{l+m}{1}\PYG{+w}{ }\PYG{n}{rot}
\end{sphinxVerbatim}


\subsection{Description}
\label{\detokenize{fix_tfmc:description}}
\sphinxAtStartPar
Perform uniform\sphinxhyphen{}acceptance force\sphinxhyphen{}bias Monte Carlo (fbMC) simulations,
using the time\sphinxhyphen{}stamped force\sphinxhyphen{}bias Monte Carlo (tfMC) algorithm
described in {\hyperref[\detokenize{fix_tfmc:mees}]{\sphinxcrossref{\DUrole{std,std-ref}{(Mees)}}}} and {\hyperref[\detokenize{fix_tfmc:bal}]{\sphinxcrossref{\DUrole{std,std-ref}{(Bal)}}}}.

\sphinxAtStartPar
One successful use case of force\sphinxhyphen{}bias Monte Carlo methods is that they
can be used to extend the time scale of atomistic simulations, in
particular when long time scale relaxation effects must be considered;
some interesting examples are given in the review by {\hyperref[\detokenize{fix_tfmc:neyts}]{\sphinxcrossref{\DUrole{std,std-ref}{(Neyts)}}}}.
An example of a typical use case would be the modelling of chemical
vapor deposition (CVD) processes on a surface, in which impacts by
gas\sphinxhyphen{}phase species can be performed using MD, but subsequent relaxation
of the surface is too slow to be done using MD only. Using tfMC can
allow for a much faster relaxation of the surface, so that higher
fluxes can be used, effectively extending the time scale of the
simulation. (Such an alternating simulation approach could be set up
using a {\hyperref[\detokenize{jump::doc}]{\sphinxcrossref{\DUrole{doc}{loop}}}}.)

\sphinxAtStartPar
The initial version of tfMC algorithm in {\hyperref[\detokenize{fix_tfmc:mees}]{\sphinxcrossref{\DUrole{std,std-ref}{(Mees)}}}} contained an
estimation of the effective time scale of such a simulation, but it
was later shown that the speed\sphinxhyphen{}up one can gain from a tfMC simulation
is system\sphinxhyphen{} and process\sphinxhyphen{}dependent, ranging from none to several orders
of magnitude. In general, solid\sphinxhyphen{}state processes such as
(re)crystallization or growth can be accelerated by up to two or three
orders of magnitude, whereas diffusion in the liquid phase is not
accelerated at all. The observed pseudodynamics when using the tfMC
method is not the actual dynamics one would obtain using MD, but the
relative importance of processes can match the actual relative
dynamics of the system quite well, provided \sphinxstyleemphasis{Delta} is chosen with
care. Thus, the system’s equilibrium is reached faster than in MD,
along a path that is generally roughly similar to a typical MD
simulation (but not necessarily so). See {\hyperref[\detokenize{fix_tfmc:bal}]{\sphinxcrossref{\DUrole{std,std-ref}{(Bal)}}}} for details.

\sphinxAtStartPar
Each step, all atoms in the selected group are displaced using the
stochastic tfMC algorithm, which is designed to sample the canonical
(NVT) ensemble at the temperature \sphinxstyleemphasis{Temp}. Although tfMC is a Monte
Carlo algorithm and thus strictly speaking does not perform time
integration, it is similar in the sense that it uses the forces on all
atoms in order to update their positions. Therefore, it is implemented
as a time integration fix, and no other fixes of this type (such as
{\hyperref[\detokenize{fix_nve::doc}]{\sphinxcrossref{\DUrole{doc}{fix nve}}}}) should be used at the same time. Because
velocities do not play a role in this kind of Monte Carlo simulations,
instantaneous temperatures as calculated by {\hyperref[\detokenize{compute_temp::doc}]{\sphinxcrossref{\DUrole{doc}{temperature computes}}}} or {\hyperref[\detokenize{thermo_style::doc}]{\sphinxcrossref{\DUrole{doc}{thermodynamic output}}}} have no meaning: the only relevant
temperature is the sampling temperature \sphinxstyleemphasis{Temp}.  Similarly, performing
tfMC simulations does not require setting a {\hyperref[\detokenize{timestep::doc}]{\sphinxcrossref{\DUrole{doc}{timestep}}}}
and the {\hyperref[\detokenize{thermo_style::doc}]{\sphinxcrossref{\DUrole{doc}{simulated time}}}} as calculated by LAMMPS is
meaningless.

\sphinxAtStartPar
The critical parameter determining the success of a tfMC simulation is
\sphinxstyleemphasis{Delta}, the maximal displacement length of the lightest element in
the system: the larger it is, the longer the effective time scale of
the simulation will be (there is an approximately quadratic
dependence). However, \sphinxstyleemphasis{Delta} must also be chosen sufficiently small
in order to comply with detailed balance; in general values between 5
and 10 \% of the nearest neighbor distance are found to be a good
choice. For a more extensive discussion with specific examples, please
refer to {\hyperref[\detokenize{fix_tfmc:bal}]{\sphinxcrossref{\DUrole{std,std-ref}{(Bal)}}}}, which also describes how the code calculates
element\sphinxhyphen{}specific maximal displacements from \sphinxstyleemphasis{Delta}, based on the
fourth root of their mass.

\sphinxAtStartPar
Because of the uncorrelated movements of the atoms, the center\sphinxhyphen{}of\sphinxhyphen{}mass
of the fix group will not necessarily be stationary, just like its
orientation. When the \sphinxstyleemphasis{com} keyword is used, all atom positions will
be shifted (after every tfMC iteration) in order to fix the position
of the center\sphinxhyphen{}of\sphinxhyphen{}mass along the included directions, by setting the
corresponding flag to 1. The \sphinxstyleemphasis{rot} keyword does the same for the
rotational component of the tfMC displacements after every iteration.

\begin{sphinxadmonition}{note}{Note:}
\sphinxAtStartPar
the \sphinxstyleemphasis{com} and \sphinxstyleemphasis{rot} keywords should not be used if an external
force is acting on the specified fix group, along the included
directions. This can be either a true external force (e.g.  through
{\hyperref[\detokenize{fix_wall::doc}]{\sphinxcrossref{\DUrole{doc}{fix wall}}}}) or forces due to the interaction with atoms
not included in the fix group. This is because in such cases,
translations or rotations of the fix group could be induced by these
external forces, and removing them will lead to a violation of
detailed balance.
\end{sphinxadmonition}


\bigskip\hrule\bigskip



\subsection{Restart, fix\_modify, output, run start/stop, minimize info}
\label{\detokenize{fix_tfmc:restart-fix-modify-output-run-start-stop-minimize-info}}
\sphinxAtStartPar
No information about this fix is written to {\hyperref[\detokenize{restart::doc}]{\sphinxcrossref{\DUrole{doc}{binary restart files}}}}.

\sphinxAtStartPar
None of the {\hyperref[\detokenize{fix_modify::doc}]{\sphinxcrossref{\DUrole{doc}{fix\_modify}}}} options are relevant to this
fix.

\sphinxAtStartPar
This fix is not invoked during {\hyperref[\detokenize{minimize::doc}]{\sphinxcrossref{\DUrole{doc}{energy minimization}}}}.


\subsection{Restrictions}
\label{\detokenize{fix_tfmc:restrictions}}
\sphinxAtStartPar
This fix is part of the MC package.  It is only enabled if LAMMPS was
built with that package.  See the {\hyperref[\detokenize{Build_package::doc}]{\sphinxcrossref{\DUrole{doc}{Build package}}}}
doc page for more info.

\sphinxAtStartPar
This fix is not compatible with {\hyperref[\detokenize{fix_shake::doc}]{\sphinxcrossref{\DUrole{doc}{fix shake}}}}.


\subsection{Related commands}
\label{\detokenize{fix_tfmc:related-commands}}
\sphinxAtStartPar
{\hyperref[\detokenize{fix_gcmc::doc}]{\sphinxcrossref{\DUrole{doc}{fix gcmc}}}}, {\hyperref[\detokenize{fix_nh::doc}]{\sphinxcrossref{\DUrole{doc}{fix nvt}}}}


\subsection{Default}
\label{\detokenize{fix_tfmc:default}}
\sphinxAtStartPar
The option default is com = 0 0 0


\bigskip\hrule\bigskip

\phantomsection\label{\detokenize{fix_tfmc:bal}}
\sphinxAtStartPar
\sphinxstylestrong{(Bal)} K. M Bal and E. C. Neyts, J. Chem. Phys. 141, 204104 (2014).

\phantomsection\label{\detokenize{fix_tfmc:mees}}
\sphinxAtStartPar
\sphinxstylestrong{(Mees)} M. J. Mees, G. Pourtois, E. C. Neyts, B. J. Thijsse, and
A. Stesmans, Phys. Rev. B 85, 134301 (2012).

\phantomsection\label{\detokenize{fix_tfmc:neyts}}
\sphinxAtStartPar
\sphinxstylestrong{(Neyts)} E. C. Neyts and A. Bogaerts, Theor. Chem. Acc. 132, 1320
(2013).

\sphinxstepscope

\index{fix tgnvt/drude@\spxentry{fix tgnvt/drude}}\ignorespaces 
\index{fix tgnpt/drude@\spxentry{fix tgnpt/drude}}\ignorespaces 

\section{fix tgnvt/drude command}
\label{\detokenize{fix_tgnh_drude:fix-tgnvt-drude-command}}\label{\detokenize{fix_tgnh_drude:index-1}}\label{\detokenize{fix_tgnh_drude:index-0}}\label{\detokenize{fix_tgnh_drude::doc}}

\section{fix tgnpt/drude command}
\label{\detokenize{fix_tgnh_drude:fix-tgnpt-drude-command}}

\subsection{Syntax}
\label{\detokenize{fix_tgnh_drude:syntax}}
\begin{sphinxVerbatim}[commandchars=\\\{\}]
\PYG{k}{fix }\PYG{n+nv+nvIdentifier}{ID}\PYG{+w}{ }\PYG{n+nv+nvIdentifier}{group\PYGZhy{}ID}\PYG{+w}{ }\PYG{n}{style\PYGZus{}name}\PYG{+w}{ }\PYG{n}{keyword}\PYG{+w}{ }\PYG{n}{values}\PYG{+w}{ }\PYG{n}{...}
\end{sphinxVerbatim}
\begin{itemize}
\item {} 
\sphinxAtStartPar
ID, group\sphinxhyphen{}ID are documented in {\hyperref[\detokenize{fix::doc}]{\sphinxcrossref{\DUrole{doc}{fix}}}} command

\item {} 
\sphinxAtStartPar
style\_name = \sphinxstyleemphasis{tgnvt/drude} or \sphinxstyleemphasis{tgnpt/drude}

\item {} 
\sphinxAtStartPar
one or more keyword/values pairs may be appended
\begin{sphinxalltt}
keyword = \sphinxstyleemphasis{temp} \sphinxstyleemphasis{iso} or \sphinxstyleemphasis{aniso} or \sphinxstyleemphasis{tri} or \sphinxstyleemphasis{x} or \sphinxstyleemphasis{y} or \sphinxstyleemphasis{z} or \sphinxstyleemphasis{xy} or \sphinxstyleemphasis{yz} or \sphinxstyleemphasis{xz} or \sphinxstyleemphasis{couple} or \sphinxstyleemphasis{tchain} or \sphinxstyleemphasis{pchain} or \sphinxstyleemphasis{mtk} or \sphinxstyleemphasis{tloop} or \sphinxstyleemphasis{ploop} or \sphinxstyleemphasis{nreset} or \sphinxstyleemphasis{scalexy} or \sphinxstyleemphasis{scaleyz} or \sphinxstyleemphasis{scalexz} or \sphinxstyleemphasis{flip} or \sphinxstyleemphasis{fixedpoint}
  \sphinxstyleemphasis{temp} values = Tstart Tstop Tdamp Tdrude Tdamp\_drude
    Tstart, Tstop = external temperature at start/end of run (temperature units)
    Tdamp = temperature damping parameter (time units)
    Tdrude = desired temperature of Drude oscillators (temperature units)
    Tdamp\_drude = temperature damping parameter for Drude oscillators (time units)
  \sphinxstyleemphasis{iso} or \sphinxstyleemphasis{aniso} or \sphinxstyleemphasis{tri} values = Pstart Pstop Pdamp
    Pstart,Pstop = scalar external pressure at start/end of run (pressure units)
    Pdamp = pressure damping parameter (time units)
  \sphinxstyleemphasis{x} or \sphinxstyleemphasis{y} or \sphinxstyleemphasis{z} or \sphinxstyleemphasis{xy} or \sphinxstyleemphasis{yz} or \sphinxstyleemphasis{xz} values = Pstart Pstop Pdamp
    Pstart,Pstop = external stress tensor component at start/end of run (pressure units)
    Pdamp = stress damping parameter (time units)
  \sphinxstyleemphasis{couple} = \sphinxstyleemphasis{none} or \sphinxstyleemphasis{xyz} or \sphinxstyleemphasis{xy} or \sphinxstyleemphasis{yz} or \sphinxstyleemphasis{xz}
  \sphinxstyleemphasis{tchain} value = N
    N = length of thermostat chain (1 = single thermostat)
  \sphinxstyleemphasis{pchain} value = N
    N length of thermostat chain on barostat (0 = no thermostat)
  \sphinxstyleemphasis{mtk} value = \sphinxstyleemphasis{yes} or \sphinxstyleemphasis{no} = add in MTK adjustment term or not
  \sphinxstyleemphasis{tloop} value = M
    M = number of sub\sphinxhyphen{}cycles to perform on thermostat
  \sphinxstyleemphasis{ploop} value = M
    M = number of sub\sphinxhyphen{}cycles to perform on barostat thermostat
  \sphinxstyleemphasis{nreset} value = reset reference cell every this many timesteps
  \sphinxstyleemphasis{scalexy} value = \sphinxstyleemphasis{yes} or \sphinxstyleemphasis{no} = scale xy with ly
  \sphinxstyleemphasis{scaleyz} value = \sphinxstyleemphasis{yes} or \sphinxstyleemphasis{no} = scale yz with lz
  \sphinxstyleemphasis{scalexz} value = \sphinxstyleemphasis{yes} or \sphinxstyleemphasis{no} = scale xz with lz
  \sphinxstyleemphasis{flip} value = \sphinxstyleemphasis{yes} or \sphinxstyleemphasis{no} = allow or disallow box flips when it becomes highly skewed
  \sphinxstyleemphasis{fixedpoint} values = x y z
    x,y,z = perform barostat dilation/contraction around this point (distance units)
\end{sphinxalltt}

\end{itemize}


\subsection{Examples}
\label{\detokenize{fix_tgnh_drude:examples}}
\begin{sphinxVerbatim}[commandchars=\\\{\}]
\PYG{k}{comm\PYGZus{}modify}\PYG{+w}{ }\PYG{n}{vel}\PYG{+w}{ }\PYG{n}{yes}
\PYG{k}{fix }\PYG{n+nv+nvIdentifier}{1}\PYG{+w}{ }\PYG{n+nv+nvIdentifier}{all}\PYG{+w}{ }\PYG{n}{tgnvt}\PYG{o}{/}\PYG{n}{drude}\PYG{+w}{ }\PYG{n}{temp}\PYG{+w}{ }\PYG{l+m}{300.0}\PYG{+w}{ }\PYG{l+m}{300.0}\PYG{+w}{ }\PYG{l+m}{100.0}\PYG{+w}{ }\PYG{l+m}{1.0}\PYG{+w}{ }\PYG{l+m}{20.0}
\PYG{k}{fix }\PYG{n+nv+nvIdentifier}{1}\PYG{+w}{ }\PYG{n+nv+nvIdentifier}{water}\PYG{+w}{ }\PYG{n}{tgnpt}\PYG{o}{/}\PYG{n}{drude}\PYG{+w}{ }\PYG{n}{temp}\PYG{+w}{ }\PYG{l+m}{300.0}\PYG{+w}{ }\PYG{l+m}{300.0}\PYG{+w}{ }\PYG{l+m}{100.0}\PYG{+w}{ }\PYG{l+m}{1.0}\PYG{+w}{ }\PYG{l+m}{20.0}\PYG{+w}{ }\PYG{n}{iso}\PYG{+w}{ }\PYG{l+m}{0.0}\PYG{+w}{ }\PYG{l+m}{0.0}\PYG{+w}{ }\PYG{l+m}{1000.0}
\PYG{k}{fix }\PYG{n+nv+nvIdentifier}{2}\PYG{+w}{ }\PYG{n+nv+nvIdentifier}{jello}\PYG{+w}{ }\PYG{n}{tgnpt}\PYG{o}{/}\PYG{n}{drude}\PYG{+w}{ }\PYG{n}{temp}\PYG{+w}{ }\PYG{l+m}{300.0}\PYG{+w}{ }\PYG{l+m}{300.0}\PYG{+w}{ }\PYG{l+m}{100.0}\PYG{+w}{ }\PYG{l+m}{1.0}\PYG{+w}{ }\PYG{l+m}{20.0}\PYG{+w}{ }\PYG{n}{tri}\PYG{+w}{ }\PYG{l+m}{5.0}\PYG{+w}{ }\PYG{l+m}{5.0}\PYG{+w}{ }\PYG{l+m}{1000.0}
\PYG{k}{fix }\PYG{n+nv+nvIdentifier}{2}\PYG{+w}{ }\PYG{n+nv+nvIdentifier}{ice}\PYG{+w}{ }\PYG{n}{tgnpt}\PYG{o}{/}\PYG{n}{drude}\PYG{+w}{ }\PYG{n}{temp}\PYG{+w}{ }\PYG{l+m}{250.0}\PYG{+w}{ }\PYG{l+m}{250.0}\PYG{+w}{ }\PYG{l+m}{100.0}\PYG{+w}{ }\PYG{l+m}{1.0}\PYG{+w}{ }\PYG{l+m}{20.0}\PYG{+w}{ }\PYG{n}{x}\PYG{+w}{ }\PYG{l+m}{1.0}\PYG{+w}{ }\PYG{l+m}{1.0}\PYG{+w}{ }\PYG{l+m}{0.5}\PYG{+w}{ }\PYG{n}{y}\PYG{+w}{ }\PYG{l+m}{2.0}\PYG{+w}{ }\PYG{l+m}{2.0}\PYG{+w}{ }\PYG{l+m}{0.5}\PYG{+w}{ }\PYG{n}{z}\PYG{+w}{ }\PYG{l+m}{3.0}\PYG{+w}{ }\PYG{l+m}{3.0}\PYG{+w}{ }\PYG{l+m}{0.5}\PYG{+w}{ }\PYG{n}{yz}\PYG{+w}{ }\PYG{l+m}{0.1}\PYG{+w}{ }\PYG{l+m}{0.1}\PYG{+w}{ }\PYG{l+m}{0.5}\PYG{+w}{ }\PYG{n}{xz}\PYG{+w}{ }\PYG{l+m}{0.2}\PYG{+w}{ }\PYG{l+m}{0.2}\PYG{+w}{ }\PYG{l+m}{0.5}\PYG{+w}{ }\PYG{n}{xy}\PYG{+w}{ }\PYG{l+m}{0.3}\PYG{+w}{ }\PYG{l+m}{0.3}\PYG{+w}{ }\PYG{l+m}{0.5}\PYG{+w}{ }\PYG{n}{nreset}\PYG{+w}{ }\PYG{l+m}{1000}
\end{sphinxVerbatim}

\sphinxAtStartPar
Example input scripts available: examples/PACKAGES/drude


\subsection{Description}
\label{\detokenize{fix_tgnh_drude:description}}
\sphinxAtStartPar
These commands are variants of the Nose\sphinxhyphen{}Hoover fix styles {\hyperref[\detokenize{fix_nh::doc}]{\sphinxcrossref{\DUrole{doc}{fix nvt}}}} and {\hyperref[\detokenize{fix_nh::doc}]{\sphinxcrossref{\DUrole{doc}{fix npt}}}} for thermalized Drude polarizable
models.  They apply temperature\sphinxhyphen{}grouped Nose\sphinxhyphen{}Hoover thermostat (TGNH)
proposed by {\hyperref[\detokenize{Howto_drude2:tgnh-son}]{\sphinxcrossref{\DUrole{std,std-ref}{(Son)}}}}.  When there are fast vibrational
modes with frequencies close to Drude oscillators (e.g. double bonds or
out\sphinxhyphen{}of\sphinxhyphen{}plane torsions), this thermostat can provide better kinetic
energy equipartitioning.

\sphinxAtStartPar
The difference between TGNH and the original Nose\sphinxhyphen{}Hoover thermostat is that,
TGNH separates the kinetic energy of the group into three contributions:
molecular center of mass (COM) motion,
motion of COM of atom\sphinxhyphen{}Drude pairs or non\sphinxhyphen{}polarizable atoms relative to molecular COM,
and relative motion of atom\sphinxhyphen{}Drude pairs.
An independent Nose\sphinxhyphen{}Hoover chain is applied to each type of motion.
The temperatures for these three types of motion are denoted as
molecular translational temperature (\(T_\mathrm{M}\)), real atomic temperature (\(T_\mathrm{R}\)) and Drude temperature (\(T_\mathrm{D}\)),
which are defined in terms of their associated degrees of freedom (DOF):
\begin{equation*}
\begin{split}T_\mathrm{M}=\frac{\Sigma_{i}^{N_\mathrm{mol}} M_i V_i^2}{3 \left ( N_\mathrm{mol} - \frac{N_\mathrm{mol}}{N_\mathrm{mol,sys}} \right ) k_\mathrm{B}}\end{split}
\end{equation*}\begin{equation*}
\begin{split}T_\mathrm{R}=\frac{\Sigma_{i}^{N_\mathrm{real}} m_i (v_i-v_{M,i})^2}{(N_\mathrm{DOF} - 3 N_\mathrm{mol} + 3 \frac{N_\mathrm{mol}}{N_\mathrm{mol,sys}} - 3 N_\mathrm{drude}) k_\mathrm{B}}\end{split}
\end{equation*}\begin{equation*}
\begin{split}T_\mathrm{D}=\frac{\Sigma_{i}^{N_\mathrm{drude}} m_i^{\prime} v_i^{\prime 2}}{3 N_\mathrm{drude} k_\mathrm{B}}\end{split}
\end{equation*}
\sphinxAtStartPar
Here \(N_\mathrm{mol}\) and \(N_\mathrm{mol,sys}\) are the numbers of molecules in the group and in the whole system, respectively.
\(N_\mathrm{real}\) is the number of atom\sphinxhyphen{}Drude pairs and non\sphinxhyphen{}polarizable atoms in the group.
\(N_\mathrm{drude}\) is the number of Drude particles in the group.
\(N_\mathrm{DOF}\) is the DOF of the group.
\(M_i\) and \(V_i\) are the mass and the COM velocity of the i\sphinxhyphen{}th molecule.
\(m_i\) is the mass of the i\sphinxhyphen{}th atom\sphinxhyphen{}Drude pair or non\sphinxhyphen{}polarizable atom.
\(v_i\) is the velocity of COM of i\sphinxhyphen{}th atom\sphinxhyphen{}Drude pair or non\sphinxhyphen{}polarizable atom.
\(v_{M,i}\) is the COM velocity of the molecule the i\sphinxhyphen{}th atom\sphinxhyphen{}Drude pair or non\sphinxhyphen{}polarizable atom belongs to.
\(m_i^\prime\) and \(v_i^\prime\) are the reduced mass and the relative velocity of the i\sphinxhyphen{}th atom\sphinxhyphen{}Drude pair.

\begin{sphinxadmonition}{note}{Note:}
\sphinxAtStartPar
These fixes require that each atom knows whether it is a Drude particle or
not.  You must therefore use the {\hyperref[\detokenize{fix_drude::doc}]{\sphinxcrossref{\DUrole{doc}{fix drude}}}} command to
specify the Drude status of each atom type.

\sphinxAtStartPar
Because the TGNH thermostat thermostats the molecular COM motion,
all atoms belonging to the same molecule must be in the same group.
That is, these fixes can not be applied to a subset of a molecule.

\sphinxAtStartPar
For this fix to act correctly, ghost atoms need to know their velocity.
You must use the {\hyperref[\detokenize{comm_modify::doc}]{\sphinxcrossref{\DUrole{doc}{comm\_modify}}}} command to enable this.

\sphinxAtStartPar
These fixes assume that the translational DOF of the whole system is removed.
It is therefore recommended to invoke {\hyperref[\detokenize{fix_momentum::doc}]{\sphinxcrossref{\DUrole{doc}{fix momentum}}}} command so that the \(T_\mathrm{M}\) is calculated correctly.
\end{sphinxadmonition}


\bigskip\hrule\bigskip


\sphinxAtStartPar
The thermostat parameters are specified using the \sphinxstyleemphasis{temp} keyword.
The thermostat is applied to only the translational DOF
for the particles.  The translational DOF can also have
a bias velocity removed before thermostatting takes place; see the
description below.  The desired temperature for molecular and real atomic motion is a
ramped value during the run from \sphinxstyleemphasis{Tstart} to \sphinxstyleemphasis{Tstop}.  The \sphinxstyleemphasis{Tdamp}
parameter is specified in time units and determines how rapidly the
temperature is relaxed.  For example, a value of 10.0 means to relax
the temperature in a timespan of (roughly) 10 time units (e.g. \(\tau\)
or fs or ps \sphinxhyphen{} see the {\hyperref[\detokenize{units::doc}]{\sphinxcrossref{\DUrole{doc}{units}}}} command).
The parameter \sphinxstyleemphasis{Tdrude} is the desired temperature for Drude motion at each timestep.
Similar to \sphinxstyleemphasis{Tdamp}, the \sphinxstyleemphasis{Tdamp\_drude} parameter determines the relaxation speed for Drude motion.
Fix group are the only ones whose velocities and positions are updated
by the velocity/position update portion of the integration.
Other thermostat\sphinxhyphen{}related keywords are \sphinxstyleemphasis{tchain} and \sphinxstyleemphasis{tloop},
which are detailed in {\hyperref[\detokenize{fix_nh::doc}]{\sphinxcrossref{\DUrole{doc}{fix nvt}}}}.

\begin{sphinxadmonition}{note}{Note:}
\sphinxAtStartPar
A Nose\sphinxhyphen{}Hoover thermostat will not work well for arbitrary values
of \sphinxstyleemphasis{Tdamp}.  If \sphinxstyleemphasis{Tdamp} is too small, the temperature can fluctuate
wildly; if it is too large, the temperature will take a very long time
to equilibrate.  A good choice for many models is a \sphinxstyleemphasis{Tdamp} of around
100 timesteps.  A smaller \sphinxstyleemphasis{Tdamp\_drude} value would be required
to maintain Drude motion at low temperature.
\end{sphinxadmonition}

\begin{sphinxVerbatim}[commandchars=\\\{\}]
\PYG{k}{fix }\PYG{n+nv+nvIdentifier}{1}\PYG{+w}{ }\PYG{n+nv+nvIdentifier}{all}\PYG{+w}{ }\PYG{n}{nvt}\PYG{+w}{ }\PYG{n}{temp}\PYG{+w}{ }\PYG{l+m}{300.0}\PYG{+w}{ }\PYG{l+m}{300.0}\PYG{+w}{ }\PYG{n+nv}{\PYGZdl{}(}\PYG{n+nv}{100.0*dt}\PYG{n+nv}{)}\PYG{+w}{ }\PYG{l+m}{1.0}\PYG{+w}{ }\PYG{n+nv}{\PYGZdl{}(}\PYG{n+nv}{20.0*dt}\PYG{n+nv}{)}
\end{sphinxVerbatim}


\bigskip\hrule\bigskip


\sphinxAtStartPar
The barostat parameters for fix style \sphinxstyleemphasis{tgnpt/drude} is specified
using one or more of the \sphinxstyleemphasis{iso}, \sphinxstyleemphasis{aniso}, \sphinxstyleemphasis{tri}, \sphinxstyleemphasis{x}, \sphinxstyleemphasis{y}, \sphinxstyleemphasis{z}, \sphinxstyleemphasis{xy},
\sphinxstyleemphasis{xz}, \sphinxstyleemphasis{yz}, and \sphinxstyleemphasis{couple} keywords.  These keywords give you the
ability to specify all 6 components of an external stress tensor, and
to couple various of these components together so that the dimensions
they represent are varied together during a constant\sphinxhyphen{}pressure
simulation. Other barostat\sphinxhyphen{}related keywords are \sphinxstyleemphasis{pchain}, \sphinxstyleemphasis{mtk}, \sphinxstyleemphasis{ploop},
\sphinxstyleemphasis{nreset}, \sphinxstyleemphasis{scalexy}, \sphinxstyleemphasis{scaleyz}, \sphinxstyleemphasis{scalexz}, \sphinxstyleemphasis{flip}and \sphinxstyleemphasis{fixedpoint}.
The meaning of barostat parameters are detailed in {\hyperref[\detokenize{fix_nh::doc}]{\sphinxcrossref{\DUrole{doc}{fix npt}}}}.

\sphinxAtStartPar
Regardless of what atoms are in the fix group (the only atoms which
are time integrated), a global pressure or stress tensor is computed
for all atoms.  Similarly, when the size of the simulation box is
changed, all atoms are re\sphinxhyphen{}scaled to new positions.

\begin{sphinxadmonition}{note}{Note:}
\sphinxAtStartPar
Unlike the {\hyperref[\detokenize{fix_temp_berendsen::doc}]{\sphinxcrossref{\DUrole{doc}{fix temp/berendsen}}}} command
which performs thermostatting but NO time integration, these fixes
perform thermostatting/barostatting AND time integration.  Thus you
should not use any other time integration fix, such as {\hyperref[\detokenize{fix_nve::doc}]{\sphinxcrossref{\DUrole{doc}{fix nve}}}} on atoms to which this fix is applied.
Likewise, these fixes should not be used on atoms that also
have their temperature controlled by another fix \sphinxhyphen{} e.g. by {\hyperref[\detokenize{fix_langevin_drude::doc}]{\sphinxcrossref{\DUrole{doc}{fix langevin/drude}}}} command.
\end{sphinxadmonition}

\sphinxAtStartPar
See the {\hyperref[\detokenize{Howto_thermostat::doc}]{\sphinxcrossref{\DUrole{doc}{Howto thermostat}}}} and {\hyperref[\detokenize{Howto_barostat::doc}]{\sphinxcrossref{\DUrole{doc}{Howto barostat}}}} doc pages for a discussion of different
ways to compute temperature and perform thermostatting and
barostatting.


\bigskip\hrule\bigskip


\sphinxAtStartPar
Like other fixes that perform thermostatting, this fix can be used
with {\hyperref[\detokenize{compute::doc}]{\sphinxcrossref{\DUrole{doc}{compute commands}}}} that remove a “bias” from the
atom velocities.  E.g. to apply the thermostat only to atoms within a
spatial {\hyperref[\detokenize{region::doc}]{\sphinxcrossref{\DUrole{doc}{region}}}}, or to remove the center\sphinxhyphen{}of\sphinxhyphen{}mass
velocity from a group of atoms, or to remove the x\sphinxhyphen{}component of
velocity from the calculation.

\sphinxAtStartPar
This is not done by default, but only if the {\hyperref[\detokenize{fix_modify::doc}]{\sphinxcrossref{\DUrole{doc}{fix\_modify}}}} command is used to assign a temperature compute to this
fix that includes such a bias term.  See the doc pages for individual
{\hyperref[\detokenize{compute::doc}]{\sphinxcrossref{\DUrole{doc}{compute temp commands}}}} to determine which ones include
a bias.  In this case, the thermostat works in the following manner:
bias is removed from each atom, thermostatting is performed on the
remaining thermal degrees of freedom, and the bias is added back in.

\begin{sphinxadmonition}{note}{Note:}
\sphinxAtStartPar
However, not all temperature compute commands are valid to be used
with these fixes.  Precisely, only temperature compute that does
not modify the DOF of the group can be used.  E.g. {\hyperref[\detokenize{compute_temp_ramp::doc}]{\sphinxcrossref{\DUrole{doc}{compute
temp/ramp}}}} and {\hyperref[\detokenize{compute_viscosity_cos::doc}]{\sphinxcrossref{\DUrole{doc}{compute viscosity/cos}}}} compute the kinetic energy after remove a
velocity gradient without affecting the DOF of the group, then they
can be invoked in this way.  In contrast, {\hyperref[\detokenize{compute_temp_partial::doc}]{\sphinxcrossref{\DUrole{doc}{compute
temp/partial}}}} may remove the DOF at one or
more dimensions, therefore it cannot be used with these fixes.
\end{sphinxadmonition}


\bigskip\hrule\bigskip



\subsection{Restart, fix\_modify, output, run start/stop, minimize info}
\label{\detokenize{fix_tgnh_drude:restart-fix-modify-output-run-start-stop-minimize-info}}
\sphinxAtStartPar
These fixes writes the state of all the thermostat and barostat
variables to {\hyperref[\detokenize{restart::doc}]{\sphinxcrossref{\DUrole{doc}{binary restart files}}}}.  See the
{\hyperref[\detokenize{read_restart::doc}]{\sphinxcrossref{\DUrole{doc}{read\_restart}}}} command for info on how to re\sphinxhyphen{}specify
a fix in an input script that reads a restart file, so that the
operation of the fix continues in an uninterrupted fashion.

\sphinxAtStartPar
The {\hyperref[\detokenize{fix_modify::doc}]{\sphinxcrossref{\DUrole{doc}{fix\_modify}}}} \sphinxstyleemphasis{temp} and \sphinxstyleemphasis{press} options are
supported by these fixes.  You can use them to assign a {\hyperref[\detokenize{compute::doc}]{\sphinxcrossref{\DUrole{doc}{compute}}}} you have defined to this fix which will be used in its
thermostatting or barostatting procedure, as described above.  If you
do this, note that the kinetic energy derived from the compute
temperature should be consistent with the virial term computed using
all atoms for the pressure.  LAMMPS will warn you if you choose to
compute temperature on a subset of atoms.

\begin{sphinxadmonition}{note}{Note:}
\sphinxAtStartPar
If both the \sphinxstyleemphasis{temp} and \sphinxstyleemphasis{press} keywords are used in a single
thermo\_modify command (or in two separate commands), then the order
in which the keywords are specified is important.  Note that a
{\hyperref[\detokenize{compute_pressure::doc}]{\sphinxcrossref{\DUrole{doc}{pressure compute}}}} defines its own
temperature compute as an argument when it is specified.  The
\sphinxstyleemphasis{temp} keyword will override this (for the pressure compute being
used by fix npt), but only if the \sphinxstyleemphasis{temp} keyword comes after the
\sphinxstyleemphasis{press} keyword.  If the \sphinxstyleemphasis{temp} keyword comes before the \sphinxstyleemphasis{press}
keyword, then the new pressure compute specified by the \sphinxstyleemphasis{press}
keyword will be unaffected by the \sphinxstyleemphasis{temp} setting.
\end{sphinxadmonition}

\sphinxAtStartPar
The cumulative energy change in the system imposed by these fixes, due
to thermostatting and/or barostatting, are included in the
{\hyperref[\detokenize{thermo_style::doc}]{\sphinxcrossref{\DUrole{doc}{thermodynamic output}}}} keywords \sphinxstyleemphasis{ecouple} and
\sphinxstyleemphasis{econserve}.  See the {\hyperref[\detokenize{thermo_style::doc}]{\sphinxcrossref{\DUrole{doc}{thermo\_style}}}} page for
details.

\sphinxAtStartPar
These fixes compute a global scalar which can be accessed by various
{\hyperref[\detokenize{Howto_output::doc}]{\sphinxcrossref{\DUrole{doc}{output commands}}}}.  The scalar is the same
cumulative energy change due to this fix described in the previous
paragraph.  The scalar value calculated by this fix is “extensive”.

\sphinxAtStartPar
These fixes also compute a global vector of quantities, which can be
accessed by various {\hyperref[\detokenize{Howto_output::doc}]{\sphinxcrossref{\DUrole{doc}{output commands}}}}.  The vector
values are “intensive”.  The vector stores the three temperatures
\(T_\mathrm{M}\), \(T_\mathrm{R}\) and \(T_\mathrm{D}\).

\sphinxAtStartPar
These fixes can ramp their external temperature and pressure over
multiple runs, using the \sphinxstyleemphasis{start} and \sphinxstyleemphasis{stop} keywords of the {\hyperref[\detokenize{run::doc}]{\sphinxcrossref{\DUrole{doc}{run}}}} command.  See the {\hyperref[\detokenize{run::doc}]{\sphinxcrossref{\DUrole{doc}{run}}}} command for details of how
to do this.

\sphinxAtStartPar
These fixes are not invoked during {\hyperref[\detokenize{minimize::doc}]{\sphinxcrossref{\DUrole{doc}{energy minimization}}}}.


\bigskip\hrule\bigskip



\subsection{Restrictions}
\label{\detokenize{fix_tgnh_drude:restrictions}}
\sphinxAtStartPar
These fixes are only available when LAMMPS was built with the
DRUDE package.  These fixes cannot be used with dynamic groups as
defined by the {\hyperref[\detokenize{group::doc}]{\sphinxcrossref{\DUrole{doc}{group}}}} command.  These fixes cannot be
used in 2D simulations.

\sphinxAtStartPar
\sphinxstyleemphasis{X}, \sphinxstyleemphasis{y}, \sphinxstyleemphasis{z} cannot be barostatted if the associated dimension is not
periodic.  \sphinxstyleemphasis{Xy}, \sphinxstyleemphasis{xz}, and \sphinxstyleemphasis{yz} can only be barostatted if the
simulation domain is triclinic and the second dimension in the keyword
(\sphinxstyleemphasis{y} dimension in \sphinxstyleemphasis{xy}) is periodic.  The {\hyperref[\detokenize{create_box::doc}]{\sphinxcrossref{\DUrole{doc}{create\_box}}}},
{\hyperref[\detokenize{read_data::doc}]{\sphinxcrossref{\DUrole{doc}{read data}}}}, and {\hyperref[\detokenize{read_restart::doc}]{\sphinxcrossref{\DUrole{doc}{read\_restart}}}}
commands specify whether the simulation box is orthogonal or
non\sphinxhyphen{}orthogonal (triclinic) and explain the meaning of the xy,xz,yz
tilt factors.

\sphinxAtStartPar
For the \sphinxstyleemphasis{temp} keyword, the final \sphinxstyleemphasis{Tstop} cannot be 0.0 since it would
make the external T = 0.0 at some timestep during the simulation which
is not allowed in the Nose/Hoover formulation.

\sphinxAtStartPar
The \sphinxstyleemphasis{scaleyz yes}, \sphinxstyleemphasis{scalexz yes}, and \sphinxstyleemphasis{scalexy yes} options
can only be used if the second dimension in the keyword is periodic,
and if the tilt factor is not coupled to the barostat via keywords
\sphinxstyleemphasis{tri}, \sphinxstyleemphasis{yz}, \sphinxstyleemphasis{xz}, and \sphinxstyleemphasis{xy}.


\subsection{Related commands}
\label{\detokenize{fix_tgnh_drude:related-commands}}
\sphinxAtStartPar
{\hyperref[\detokenize{fix_drude::doc}]{\sphinxcrossref{\DUrole{doc}{fix drude}}}}, {\hyperref[\detokenize{fix_nh::doc}]{\sphinxcrossref{\DUrole{doc}{fix nvt}}}}, {\hyperref[\detokenize{fix_nh::doc}]{\sphinxcrossref{\DUrole{doc}{fix\_npt}}}},
{\hyperref[\detokenize{fix_modify::doc}]{\sphinxcrossref{\DUrole{doc}{fix\_modify}}}}


\subsection{Default}
\label{\detokenize{fix_tgnh_drude:default}}
\sphinxAtStartPar
The keyword defaults are tchain = 3, pchain = 3, mtk = yes, tloop = 1,
ploop = 1, nreset = 0, couple = none,
flip = yes, scaleyz = scalexz = scalexy = yes if periodic in second
dimension and not coupled to barostat, otherwise no.


\bigskip\hrule\bigskip

\phantomsection\label{\detokenize{fix_tgnh_drude:tgnh-son}}
\sphinxAtStartPar
\sphinxstylestrong{(Son)} Son, McDaniel, Cui and Yethiraj, J Phys Chem Lett, 10, 7523 (2019).

\sphinxstepscope

\index{fix thermal/conductivity@\spxentry{fix thermal/conductivity}}\ignorespaces 

\section{fix thermal/conductivity command}
\label{\detokenize{fix_thermal_conductivity:fix-thermal-conductivity-command}}\label{\detokenize{fix_thermal_conductivity:index-0}}\label{\detokenize{fix_thermal_conductivity::doc}}

\subsection{Syntax}
\label{\detokenize{fix_thermal_conductivity:syntax}}
\begin{sphinxVerbatim}[commandchars=\\\{\}]
\PYG{k}{fix }\PYG{n+nv+nvIdentifier}{ID}\PYG{+w}{ }\PYG{n+nv+nvIdentifier}{group\PYGZhy{}ID}\PYG{+w}{ }\PYG{n}{thermal}\PYG{o}{/}\PYG{n}{conductivity}\PYG{+w}{ }\PYG{n}{N}\PYG{+w}{ }\PYG{n}{edim}\PYG{+w}{ }\PYG{n}{Nbin}\PYG{+w}{ }\PYG{n}{keyword}\PYG{+w}{ }\PYG{n}{value}\PYG{+w}{ }\PYG{n}{...}
\end{sphinxVerbatim}
\begin{itemize}
\item {} 
\sphinxAtStartPar
ID, group\sphinxhyphen{}ID are documented in {\hyperref[\detokenize{fix::doc}]{\sphinxcrossref{\DUrole{doc}{fix}}}} command

\item {} 
\sphinxAtStartPar
thermal/conductivity = style name of this fix command

\item {} 
\sphinxAtStartPar
N = perform kinetic energy exchange every N steps

\item {} 
\sphinxAtStartPar
edim = \sphinxstyleemphasis{x} or \sphinxstyleemphasis{y} or \sphinxstyleemphasis{z} = direction of kinetic energy transfer

\item {} 
\sphinxAtStartPar
Nbin = \# of layers in edim direction (must be even number)

\item {} 
\sphinxAtStartPar
zero or more keyword/value pairs may be appended

\item {} 
\sphinxAtStartPar
keyword = \sphinxstyleemphasis{swap}
\begin{sphinxalltt}
\sphinxstyleemphasis{swap} value = Nswap = number of swaps to perform every N steps
\end{sphinxalltt}

\end{itemize}


\subsection{Examples}
\label{\detokenize{fix_thermal_conductivity:examples}}
\begin{sphinxVerbatim}[commandchars=\\\{\}]
\PYG{k}{fix }\PYG{n+nv+nvIdentifier}{1}\PYG{+w}{ }\PYG{n+nv+nvIdentifier}{all}\PYG{+w}{ }\PYG{n}{thermal}\PYG{o}{/}\PYG{n}{conductivity}\PYG{+w}{ }\PYG{l+m}{100}\PYG{+w}{ }\PYG{n}{z}\PYG{+w}{ }\PYG{l+m}{20}
\PYG{k}{fix }\PYG{n+nv+nvIdentifier}{1}\PYG{+w}{ }\PYG{n+nv+nvIdentifier}{all}\PYG{+w}{ }\PYG{n}{thermal}\PYG{o}{/}\PYG{n}{conductivity}\PYG{+w}{ }\PYG{l+m}{50}\PYG{+w}{ }\PYG{n}{z}\PYG{+w}{ }\PYG{l+m}{20}\PYG{+w}{ }\PYG{n}{swap}\PYG{+w}{ }\PYG{l+m}{2}
\end{sphinxVerbatim}


\subsection{Description}
\label{\detokenize{fix_thermal_conductivity:description}}
\sphinxAtStartPar
Use the Muller\sphinxhyphen{}Plathe algorithm described in {\hyperref[\detokenize{fix_thermal_conductivity:muller-plathe1}]{\sphinxcrossref{\DUrole{std,std-ref}{this paper}}}} to exchange kinetic energy between two particles
in different regions of the simulation box every N steps.  This
induces a temperature gradient in the system.  As described below this
enables the thermal conductivity of a material to be calculated.  This
algorithm is sometimes called a reverse non\sphinxhyphen{}equilibrium MD (reverse
NEMD) approach to computing thermal conductivity.  This is because the
usual NEMD approach is to impose a temperature gradient on the system
and measure the response as the resulting heat flux.  In the
Muller\sphinxhyphen{}Plathe method, the heat flux is imposed, and the temperature
gradient is the system’s response.

\sphinxAtStartPar
See the {\hyperref[\detokenize{compute_heat_flux::doc}]{\sphinxcrossref{\DUrole{doc}{compute heat/flux}}}} command for details
on how to compute thermal conductivity in an alternate way, via the
Green\sphinxhyphen{}Kubo formalism.

\sphinxAtStartPar
The simulation box is divided into \sphinxstyleemphasis{Nbin} layers in the \sphinxstyleemphasis{edim}
direction, where the layer 1 is at the low end of that dimension and
the layer \sphinxstyleemphasis{Nbin} is at the high end.  Every N steps, Nswap pairs of
atoms are chosen in the following manner.  Only atoms in the fix group
are considered.  The hottest Nswap atoms in layer 1 are selected.
Similarly, the coldest Nswap atoms in the “middle” layer (see below)
are selected.  The two sets of Nswap atoms are paired up and their
velocities are exchanged.  This effectively swaps their kinetic
energies, assuming their masses are the same.  If the masses are
different, an exchange of velocities relative to center of mass motion
of the two atoms is performed, to conserve kinetic energy.  Over time,
this induces a temperature gradient in the system which can be
measured using commands such as the following, which writes the
temperature profile (assuming z = edim) to the file tmp.profile:

\begin{sphinxVerbatim}[commandchars=\\\{\}]
\PYG{k}{compute   }\PYG{n+nv+nvIdentifier}{ke}\PYG{+w}{ }\PYG{n+nv+nvIdentifier}{all}\PYG{+w}{ }\PYG{n}{ke}\PYG{o}{/}\PYG{n}{atom}
\PYG{k}{variable  }\PYG{n+nv+nvIdentifier}{temp}\PYG{+w}{ }\PYG{n}{atom}\PYG{+w}{ }\PYG{n}{c\PYGZus{}ke}\PYG{o}{/}\PYG{l+m}{1.5}
\PYG{k}{compute   }\PYG{n+nv+nvIdentifier}{layers}\PYG{+w}{ }\PYG{n+nv+nvIdentifier}{all}\PYG{+w}{ }\PYG{n}{chunk}\PYG{o}{/}\PYG{n}{atom}\PYG{+w}{ }\PYG{n}{bin}\PYG{o}{/}\PYG{l+m}{1}\PYG{n}{d}\PYG{+w}{ }\PYG{n}{z}\PYG{+w}{ }\PYG{n}{lower}\PYG{+w}{ }\PYG{l+m}{0.05}\PYG{+w}{ }\PYG{n}{units}\PYG{+w}{ }\PYG{n}{reduced}
\PYG{k}{fix       }\PYG{n+nv+nvIdentifier}{3}\PYG{+w}{ }\PYG{n+nv+nvIdentifier}{all}\PYG{+w}{ }\PYG{n}{ave}\PYG{o}{/}\PYG{n}{chunk}\PYG{+w}{ }\PYG{l+m}{10}\PYG{+w}{ }\PYG{l+m}{100}\PYG{+w}{ }\PYG{l+m}{1000}\PYG{+w}{ }\PYG{n}{layers}\PYG{+w}{ }\PYG{n}{v\PYGZus{}temp}\PYG{+w}{ }\PYG{n}{file}\PYG{+w}{ }\PYG{n}{tmp.profile}
\end{sphinxVerbatim}

\sphinxAtStartPar
Note that by default, Nswap = 1, though this can be changed by the
optional \sphinxstyleemphasis{swap} keyword.  Setting this parameter appropriately, in
conjunction with the swap rate N, allows the heat flux to be adjusted
across a wide range of values, and the kinetic energy to be exchanged
in large chunks or more smoothly.

\sphinxAtStartPar
The “middle” layer for velocity swapping is defined as the \sphinxstyleemphasis{Nbin}/2 +
1 layer.  Thus if \sphinxstyleemphasis{Nbin} = 20, the two swapping layers are 1 and 11.
This should lead to a symmetric temperature profile since the two
layers are separated by the same distance in both directions in a
periodic sense.  This is why \sphinxstyleemphasis{Nbin} is restricted to being an even
number.

\sphinxAtStartPar
As described below, the total kinetic energy transferred by these
swaps is computed by the fix and can be output.  Dividing this
quantity by time and the cross\sphinxhyphen{}sectional area of the simulation box
yields a heat flux.  The ratio of heat flux to the slope of the
temperature profile is proportional to the thermal conductivity of the
fluid, in appropriate units.  See the {\hyperref[\detokenize{fix_thermal_conductivity:muller-plathe1}]{\sphinxcrossref{\DUrole{std,std-ref}{Muller\sphinxhyphen{}Plathe paper}}}} for details.

\begin{sphinxadmonition}{note}{Note:}
\sphinxAtStartPar
If your system is periodic in the direction of the heat flux,
then the flux is going in 2 directions.  This means the effective heat
flux in one direction is reduced by a factor of 2.  You will see this
in the equations for thermal conductivity (kappa) in the Muller\sphinxhyphen{}Plathe
paper.  LAMMPS is simply tallying kinetic energy which does not
account for whether or not your system is periodic; you must use the
value appropriately to yield a kappa for your system.
\end{sphinxadmonition}

\begin{sphinxadmonition}{note}{Note:}
\sphinxAtStartPar
After equilibration, if the temperature gradient you observe is
not linear, then you are likely swapping energy too frequently and are
not in a regime of linear response.  In this case you cannot
accurately infer a thermal conductivity and should try increasing the
Nevery parameter.
\end{sphinxadmonition}


\subsection{Restart, fix\_modify, output, run start/stop, minimize info}
\label{\detokenize{fix_thermal_conductivity:restart-fix-modify-output-run-start-stop-minimize-info}}
\sphinxAtStartPar
No information about this fix is written to {\hyperref[\detokenize{restart::doc}]{\sphinxcrossref{\DUrole{doc}{binary restart files}}}}.  None of the {\hyperref[\detokenize{fix_modify::doc}]{\sphinxcrossref{\DUrole{doc}{fix\_modify}}}} options
are relevant to this fix.

\sphinxAtStartPar
This fix computes a global scalar which can be accessed by various
{\hyperref[\detokenize{Howto_output::doc}]{\sphinxcrossref{\DUrole{doc}{output commands}}}}.  The scalar is the cumulative
kinetic energy transferred between the bottom and middle of the
simulation box (in the \sphinxstyleemphasis{edim} direction) is stored as a scalar
quantity by this fix.  This quantity is zeroed when the fix is defined
and accumulates thereafter, once every N steps.  The units of the
quantity are energy; see the {\hyperref[\detokenize{units::doc}]{\sphinxcrossref{\DUrole{doc}{units}}}} command for details.
The scalar value calculated by this fix is “intensive”.

\sphinxAtStartPar
No parameter of this fix can be used with the \sphinxstyleemphasis{start/stop} keywords of
the {\hyperref[\detokenize{run::doc}]{\sphinxcrossref{\DUrole{doc}{run}}}} command.  This fix is not invoked during {\hyperref[\detokenize{minimize::doc}]{\sphinxcrossref{\DUrole{doc}{energy minimization}}}}.


\subsection{Restrictions}
\label{\detokenize{fix_thermal_conductivity:restrictions}}
\sphinxAtStartPar
Swaps conserve both momentum and kinetic energy, even if the masses of
the swapped atoms are not equal.  Thus you should not need to
thermostat the system.  If you do use a thermostat, you may want to
apply it only to the non\sphinxhyphen{}swapped dimensions (other than \sphinxstyleemphasis{vdim}).

\sphinxAtStartPar
LAMMPS does not check, but you should not use this fix to swap the
kinetic energy of atoms that are in constrained molecules, e.g. via
{\hyperref[\detokenize{fix_shake::doc}]{\sphinxcrossref{\DUrole{doc}{fix shake}}}} or {\hyperref[\detokenize{fix_rigid::doc}]{\sphinxcrossref{\DUrole{doc}{fix rigid}}}}.  This is
because application of the constraints will alter the amount of
transferred momentum.  You should, however, be able to use flexible
molecules.  See the {\hyperref[\detokenize{fix_thermal_conductivity:zhang2}]{\sphinxcrossref{\DUrole{std,std-ref}{Zhang paper}}}} for a discussion and results
of this idea.

\sphinxAtStartPar
When running a simulation with large, massive particles or molecules
in a background solvent, you may want to only exchange kinetic energy
between solvent particles.


\subsection{Related commands}
\label{\detokenize{fix_thermal_conductivity:related-commands}}
\sphinxAtStartPar
{\hyperref[\detokenize{fix_ehex::doc}]{\sphinxcrossref{\DUrole{doc}{fix ehex}}}}, {\hyperref[\detokenize{fix_heat::doc}]{\sphinxcrossref{\DUrole{doc}{fix heat}}}}, {\hyperref[\detokenize{fix_ave_chunk::doc}]{\sphinxcrossref{\DUrole{doc}{fix ave/chunk}}}}, {\hyperref[\detokenize{fix_viscosity::doc}]{\sphinxcrossref{\DUrole{doc}{fix viscosity}}}},
{\hyperref[\detokenize{compute_heat_flux::doc}]{\sphinxcrossref{\DUrole{doc}{compute heat/flux}}}}


\subsection{Default}
\label{\detokenize{fix_thermal_conductivity:default}}
\sphinxAtStartPar
The option defaults are swap = 1.


\bigskip\hrule\bigskip

\phantomsection\label{\detokenize{fix_thermal_conductivity:muller-plathe1}}
\sphinxAtStartPar
\sphinxstylestrong{(Muller\sphinxhyphen{}Plathe)} Muller\sphinxhyphen{}Plathe, J Chem Phys, 106, 6082 (1997).

\phantomsection\label{\detokenize{fix_thermal_conductivity:zhang2}}
\sphinxAtStartPar
\sphinxstylestrong{(Zhang)} Zhang, Lussetti, de Souza, Muller\sphinxhyphen{}Plathe, J Phys Chem B,
109, 15060\sphinxhyphen{}15067 (2005).

\sphinxstepscope

\index{fix ti/spring@\spxentry{fix ti/spring}}\ignorespaces 

\section{fix ti/spring command}
\label{\detokenize{fix_ti_spring:fix-ti-spring-command}}\label{\detokenize{fix_ti_spring:index-0}}\label{\detokenize{fix_ti_spring::doc}}

\subsection{Syntax}
\label{\detokenize{fix_ti_spring:syntax}}
\begin{sphinxVerbatim}[commandchars=\\\{\}]
\PYG{k}{fix }\PYG{n+nv+nvIdentifier}{ID}\PYG{+w}{ }\PYG{n+nv+nvIdentifier}{group\PYGZhy{}ID}\PYG{+w}{ }\PYG{n}{ti}\PYG{o}{/}\PYG{n}{spring}\PYG{+w}{ }\PYG{n}{k}\PYG{+w}{ }\PYG{n}{t\PYGZus{}s}\PYG{+w}{ }\PYG{n}{t\PYGZus{}eq}\PYG{+w}{ }\PYG{n}{keyword}\PYG{+w}{ }\PYG{n}{value}\PYG{+w}{ }\PYG{n}{...}
\end{sphinxVerbatim}
\begin{itemize}
\item {} 
\sphinxAtStartPar
ID, group\sphinxhyphen{}ID are documented in {\hyperref[\detokenize{fix::doc}]{\sphinxcrossref{\DUrole{doc}{fix}}}} command

\item {} 
\sphinxAtStartPar
ti/spring = style name of this fix command

\item {} 
\sphinxAtStartPar
k = spring constant (force/distance units)

\item {} 
\sphinxAtStartPar
t\_eq = number of steps for the equilibration procedure

\item {} 
\sphinxAtStartPar
t\_s = number of steps for the switching procedure

\item {} 
\sphinxAtStartPar
zero or more keyword/value pairs may be appended to args

\item {} 
\sphinxAtStartPar
keyword = \sphinxstyleemphasis{function}
\begin{sphinxalltt}
\sphinxstyleemphasis{function} value = function\sphinxhyphen{}ID
  function\sphinxhyphen{}ID = ID of the switching function (1 or 2)
\end{sphinxalltt}

\end{itemize}


\subsection{Example}
\label{\detokenize{fix_ti_spring:example}}
\begin{sphinxVerbatim}[commandchars=\\\{\}]
\PYG{k}{fix }\PYG{n+nv+nvIdentifier}{1}\PYG{+w}{ }\PYG{n+nv+nvIdentifier}{all}\PYG{+w}{ }\PYG{n}{ti}\PYG{o}{/}\PYG{n}{spring}\PYG{+w}{ }\PYG{l+m}{50.0}\PYG{+w}{ }\PYG{l+m}{2000}\PYG{+w}{ }\PYG{l+m}{1000}\PYG{+w}{ }\PYG{n}{function}\PYG{+w}{ }\PYG{l+m}{2}
\end{sphinxVerbatim}


\subsection{Description}
\label{\detokenize{fix_ti_spring:description}}
\sphinxAtStartPar
This fix allows you to compute the free energy of crystalline solids
by performing a nonequilibrium thermodynamic integration between the
solid of interest and an Einstein crystal. A detailed explanation of
how to use this command and choose its parameters for optimal
performance and accuracy is given in the paper by
{\hyperref[\detokenize{fix_ti_spring:freitas1}]{\sphinxcrossref{\DUrole{std,std-ref}{Freitas}}}}. The paper also presents a short summary of the
theory of nonequilibrium thermodynamic integration.

\sphinxAtStartPar
The thermodynamic integration procedure is performed by rescaling the
force on each atom. Given an atomic configuration the force (F) on
each atom is given by
\begin{equation*}
\begin{split}F = \left( 1-\lambda \right) F_{\text{solid}} + \lambda F_{\text{harm}}\end{split}
\end{equation*}
\sphinxAtStartPar
where F\_solid is the force that acts on an atom due to an interatomic
potential (\sphinxstyleemphasis{e.g.} EAM potential), F\_harm is the force due to the
Einstein crystal harmonic spring, and lambda is the coupling parameter
of the thermodynamic integration. An Einstein crystal is a solid where
each atom is attached to its equilibrium position by a harmonic spring
with spring constant \sphinxstyleemphasis{k}. With this fix a spring force is applied
independently to each atom in the group defined by the fix to tether
it to its initial position. The initial position of each atom is its
position at the time the fix command was issued.

\sphinxAtStartPar
The fix acts as follows: during the first \sphinxstyleemphasis{t\_eq} steps after the fix
is defined the value of lambda is zero. This is the period to
equilibrate the system in the lambda = 0 state. After this the value
of lambda changes dynamically during the simulation from 0 to 1
according to the function defined using the keyword \sphinxstyleemphasis{function}
(described below), this switching from lambda from 0 to 1 is done in
\sphinxstyleemphasis{t\_s} steps. Then comes the second equilibration period of \sphinxstyleemphasis{t\_eq} to
equilibrate the system in the lambda = 1 state. After that, the
switching back to the lambda = 0 state is made using \sphinxstyleemphasis{t\_s} timesteps
and following the same switching function. After this period the value
of lambda is kept equal to zero and the fix has no other effect on the
dynamics of the system.

\sphinxAtStartPar
The processes described above is known as nonequilibrium thermodynamic
integration and is has been shown ({\hyperref[\detokenize{fix_ti_spring:freitas1}]{\sphinxcrossref{\DUrole{std,std-ref}{Freitas}}}}) to present a
much superior efficiency when compared to standard equilibrium
methods. The reason why the switching it is made in both directions
(potential to Einstein crystal and back) is to eliminate the
dissipated heat due to the nonequilibrium process. Further details
about nonequilibrium thermodynamic integration and its implementation
in LAMMPS is available in {\hyperref[\detokenize{fix_ti_spring:freitas1}]{\sphinxcrossref{\DUrole{std,std-ref}{Freitas}}}}.

\sphinxAtStartPar
The \sphinxstyleemphasis{function} keyword allows the use of two different lambda
paths. Option \sphinxstyleemphasis{1} results in a constant rate of change of lambda with
time:
\begin{equation*}
\begin{split}\lambda(\tau) = \tau\end{split}
\end{equation*}
\sphinxAtStartPar
where \(\tau\) is the scaled time variable \sphinxstyleemphasis{t/t\_s}. The option \sphinxstyleemphasis{2}
performs the lambda switching at a rate defined by the following
switching function
\begin{equation*}
\begin{split}\lambda(\tau) = \tau^5 \left( 70 \tau^4 - 315 \tau^3 + 540 \tau^2 -
420 \tau + 126 \right)\end{split}
\end{equation*}
\sphinxAtStartPar
This function has zero slope as lambda approaches its extreme values
(0 and 1), according to {\hyperref[\detokenize{fix_ti_spring:dekoning96}]{\sphinxcrossref{\DUrole{std,std-ref}{de Koning}}}} this results in
smaller fluctuations on the integral to be computed on the
thermodynamic integration. The use of option \sphinxstyleemphasis{2} is recommended since
it results in better accuracy and less dissipation without any
increase in computational resources cost.

\begin{sphinxadmonition}{note}{Note:}
\sphinxAtStartPar
As described in {\hyperref[\detokenize{fix_ti_spring:freitas1}]{\sphinxcrossref{\DUrole{std,std-ref}{Freitas}}}}, it is important to keep
the center\sphinxhyphen{}of\sphinxhyphen{}mass fixed during the thermodynamic integration. A
nonzero total velocity will result in divergences during the
integration due to the fact that the atoms are ‘attached’ to their
equilibrium positions by the Einstein crystal. Check the option
\sphinxstyleemphasis{zero} of {\hyperref[\detokenize{fix_langevin::doc}]{\sphinxcrossref{\DUrole{doc}{fix langevin}}}} and {\hyperref[\detokenize{velocity::doc}]{\sphinxcrossref{\DUrole{doc}{velocity}}}}. The use of the Nose\sphinxhyphen{}Hoover thermostat ({\hyperref[\detokenize{fix_nh::doc}]{\sphinxcrossref{\DUrole{doc}{fix nvt}}}}) is \sphinxstyleemphasis{NOT} recommended due to its well documented issues
with the canonical sampling of harmonic degrees of freedom (notice
that the \sphinxstyleemphasis{chain} option will \sphinxstyleemphasis{NOT} solve this problem). The
Langevin thermostat ({\hyperref[\detokenize{fix_langevin::doc}]{\sphinxcrossref{\DUrole{doc}{fix langevin}}}}) correctly
thermostats the system and we advise its usage with ti/spring
command.
\end{sphinxadmonition}


\subsection{Restart, fix\_modify, output, run start/stop, minimize info}
\label{\detokenize{fix_ti_spring:restart-fix-modify-output-run-start-stop-minimize-info}}
\sphinxAtStartPar
This fix writes the original coordinates of tethered atoms to
{\hyperref[\detokenize{restart::doc}]{\sphinxcrossref{\DUrole{doc}{binary restart files}}}}, so that the spring effect will
be the same in a restarted simulation. See the {\hyperref[\detokenize{read_restart::doc}]{\sphinxcrossref{\DUrole{doc}{read restart}}}} command for info on how to re\sphinxhyphen{}specify a fix in an
input script that reads a restart file, so that the operation of the
fix continues in an uninterrupted fashion.

\sphinxAtStartPar
The {\hyperref[\detokenize{fix_modify::doc}]{\sphinxcrossref{\DUrole{doc}{fix\_modify}}}} \sphinxstyleemphasis{energy} option is supported by
this fix to add the energy stored in the per\sphinxhyphen{}atom springs to the
global potential energy of the system as part of {\hyperref[\detokenize{thermo_style::doc}]{\sphinxcrossref{\DUrole{doc}{thermodynamic
output}}}}. The default setting for this fix is
{\hyperref[\detokenize{fix_modify::doc}]{\sphinxcrossref{\DUrole{doc}{fix\_modify energy no}}}}.

\sphinxAtStartPar
This fix computes a global scalar and a global vector quantities which
can be accessed by various {\hyperref[\detokenize{Howto_output::doc}]{\sphinxcrossref{\DUrole{doc}{output commands}}}}. The
scalar is an energy which is the sum of the spring energy for each
atom, where the per\sphinxhyphen{}atom energy is \(0.5 \cdot k \cdot r^2\).
The vector stores 2 values.  The first value is the coupling parameter lambda.
The second value is the derivative of lambda with respect to the integer
timestep \sphinxstyleemphasis{s}, i.e. \(\frac{d \lambda}{d s}\).  In order to obtain
\(\frac{d \lambda}{d t}\),
where t is simulation time, this 2nd value needs to be divided by the
timestep size (e.g. 0.5 fs).  The scalar and vector values calculated
by this fix are “extensive”.

\sphinxAtStartPar
No parameter of this fix can be used with the \sphinxstyleemphasis{start/stop} keywords of
the {\hyperref[\detokenize{run::doc}]{\sphinxcrossref{\DUrole{doc}{run}}}} command.

\sphinxAtStartPar
The forces due to this fix are imposed during an energy minimization,
invoked by the {\hyperref[\detokenize{minimize::doc}]{\sphinxcrossref{\DUrole{doc}{minimize}}}} command.

\begin{sphinxadmonition}{note}{Note:}
\sphinxAtStartPar
If you want the per\sphinxhyphen{}atom spring energy to be included in the
total potential energy of the system (the quantity being minimized),
you MUST enable the {\hyperref[\detokenize{fix_modify::doc}]{\sphinxcrossref{\DUrole{doc}{fix modify}}}} \sphinxstyleemphasis{energy} option for
this fix.
\end{sphinxadmonition}


\subsection{Related commands}
\label{\detokenize{fix_ti_spring:related-commands}}
\sphinxAtStartPar
{\hyperref[\detokenize{fix_spring::doc}]{\sphinxcrossref{\DUrole{doc}{fix spring}}}}, {\hyperref[\detokenize{fix_adapt::doc}]{\sphinxcrossref{\DUrole{doc}{fix adapt}}}}


\subsection{Restrictions}
\label{\detokenize{fix_ti_spring:restrictions}}
\sphinxAtStartPar
This fix is part of the EXTRA\sphinxhyphen{}FIX package. It is only enabled if
LAMMPS was built with that package. See the
{\hyperref[\detokenize{Build_package::doc}]{\sphinxcrossref{\DUrole{doc}{Build package}}}} page for more info.


\subsection{Default}
\label{\detokenize{fix_ti_spring:default}}
\sphinxAtStartPar
The keyword default is function = 1.


\bigskip\hrule\bigskip

\phantomsection\label{\detokenize{fix_ti_spring:freitas1}}
\sphinxAtStartPar
\sphinxstylestrong{(Freitas)} Freitas, Asta, and de Koning, Computational Materials
Science, 112, 333 (2016).

\phantomsection\label{\detokenize{fix_ti_spring:dekoning96}}
\sphinxAtStartPar
\sphinxstylestrong{(de Koning)} de Koning and Antonelli, Phys Rev E, 53, 465 (1996).

\sphinxstepscope

\index{fix tmd@\spxentry{fix tmd}}\ignorespaces 

\section{fix tmd command}
\label{\detokenize{fix_tmd:fix-tmd-command}}\label{\detokenize{fix_tmd:index-0}}\label{\detokenize{fix_tmd::doc}}

\subsection{Syntax}
\label{\detokenize{fix_tmd:syntax}}
\begin{sphinxVerbatim}[commandchars=\\\{\}]
\PYG{k}{fix }\PYG{n+nv+nvIdentifier}{ID}\PYG{+w}{ }\PYG{n+nv+nvIdentifier}{group\PYGZhy{}ID}\PYG{+w}{ }\PYG{n}{tmd}\PYG{+w}{ }\PYG{n}{rho\PYGZus{}final}\PYG{+w}{ }\PYG{n}{file1}\PYG{+w}{ }\PYG{n}{N}\PYG{+w}{ }\PYG{n}{file2}
\end{sphinxVerbatim}
\begin{itemize}
\item {} 
\sphinxAtStartPar
ID, group\sphinxhyphen{}ID are documented in {\hyperref[\detokenize{fix::doc}]{\sphinxcrossref{\DUrole{doc}{fix}}}} command

\item {} 
\sphinxAtStartPar
tmd = style name of this fix command

\item {} 
\sphinxAtStartPar
rho\_final = desired value of rho at the end of the run (distance units)

\item {} 
\sphinxAtStartPar
file1 = filename to read target structure from

\item {} 
\sphinxAtStartPar
N = dump TMD statistics every this many timesteps, 0 = no dump

\item {} 
\sphinxAtStartPar
file2 = filename to write TMD statistics to (only needed if N \textgreater{} 0)

\end{itemize}


\subsection{Examples}
\label{\detokenize{fix_tmd:examples}}
\begin{sphinxVerbatim}[commandchars=\\\{\}]
\PYG{k}{fix }\PYG{n+nv+nvIdentifier}{1}\PYG{+w}{ }\PYG{n+nv+nvIdentifier}{all}\PYG{+w}{ }\PYG{n}{nve}
\PYG{k}{fix }\PYG{n+nv+nvIdentifier}{2}\PYG{+w}{ }\PYG{n+nv+nvIdentifier}{tmdatoms}\PYG{+w}{ }\PYG{n}{tmd}\PYG{+w}{ }\PYG{l+m}{1.0}\PYG{+w}{ }\PYG{n}{target\PYGZus{}file}\PYG{+w}{ }\PYG{l+m}{100}\PYG{+w}{ }\PYG{n}{tmd\PYGZus{}dump\PYGZus{}file}
\end{sphinxVerbatim}


\subsection{Description}
\label{\detokenize{fix_tmd:description}}
\sphinxAtStartPar
Perform targeted molecular dynamics (TMD) on a group of atoms.  A
holonomic constraint is used to force the atoms to move towards (or away
from) the target configuration.  The parameter “rho” is monotonically
decreased (or increased) from its initial value to rho\_final at the end
of the run.

\sphinxAtStartPar
Rho has distance units and is a measure of the root\sphinxhyphen{}mean\sphinxhyphen{}squared
distance (RMSD) between the current configuration of the atoms in the
group and the target coordinates listed in file1.  Thus a value of
rho\_final = 0.0 means move the atoms all the way to the final
structure during the course of the run.

\sphinxAtStartPar
The target file1 can be ASCII text or a gzipped text file (detected by
a .gz suffix).  The format of the target file1 is as follows:

\begin{sphinxVerbatim}[commandchars=\\\{\}]
0.0 25.0 xlo xhi
0.0 25.0 ylo yhi
0.0 25.0 zlo zhi
125     24.97311   1.69005     23.46956 0 0 \PYGZhy{}1
126     1.94691    2.79640     1.92799  1 0 0
127     0.15906    3.46099     0.79121  1 0 0
...
\end{sphinxVerbatim}

\sphinxAtStartPar
The first 3 lines may or may not be needed, depending on the format of
the atoms to follow.  If image flags are included with the atoms, the
first 3 lo/hi lines \sphinxstylestrong{must} appear in the file.  If image flags are not
included, the first 3 lines \sphinxstylestrong{must not} appear.  The 3 lines contain the
simulation box dimensions for the atom coordinates, in the same format
as in a LAMMPS data file (see the {\hyperref[\detokenize{read_data::doc}]{\sphinxcrossref{\DUrole{doc}{read\_data}}}} command).

\sphinxAtStartPar
The remaining lines each contain an atom ID and its target x,y,z
coordinates.  The atom lines (all or none of them) can optionally be
followed by 3 integer values: nx,ny,nz.For periodic dimensions, they
specify which image of the box the atom is considered to be in, i.e. a
value of N (positive or negative) means add N times the box length to
the coordinate to get the true value.  Those 3 integers either must
be given for all atoms or none.

\sphinxAtStartPar
The atom lines can be listed in any order, but every atom in the group
must be listed in the file.  Atoms not in the fix group may also be
listed; they will be ignored.

\sphinxAtStartPar
Comments starting with ‘\#’ and empty lines may be included as well.

\sphinxAtStartPar
TMD statistics are written to file2 every N timesteps, unless N is
specified as 0, which means no statistics.

\sphinxAtStartPar
The atoms in the fix tmd group should be integrated (via a fix nve,
nvt, npt) along with other atoms in the system.

\sphinxAtStartPar
Restarts can be used with a fix tmd command.  For example, imagine a
10000 timestep run with a rho\_initial = 11 and a rho\_final = 1.  If a
restart file was written after 2000 time steps, then the configuration
in the file would have a rho value of 9.  A new 8000 time step run
could be performed with the same rho\_final = 1 to complete the
conformational change at the same transition rate.  Note that for
restarted runs, the name of the TMD statistics file should be changed
to prevent it being overwritten.

\sphinxAtStartPar
For more information about TMD, see {\hyperref[\detokenize{fix_tmd:schlitter1}]{\sphinxcrossref{\DUrole{std,std-ref}{(Schlitter1)}}}} and
{\hyperref[\detokenize{fix_tmd:schlitter2}]{\sphinxcrossref{\DUrole{std,std-ref}{(Schlitter2)}}}}.


\subsection{Restart, fix\_modify, output, run start/stop, minimize info}
\label{\detokenize{fix_tmd:restart-fix-modify-output-run-start-stop-minimize-info}}
\sphinxAtStartPar
No information about this fix is written to {\hyperref[\detokenize{restart::doc}]{\sphinxcrossref{\DUrole{doc}{binary restart files}}}}.  None of the {\hyperref[\detokenize{fix_modify::doc}]{\sphinxcrossref{\DUrole{doc}{fix\_modify}}}} options
are relevant to this fix.  No global or per\sphinxhyphen{}atom quantities are stored
by this fix for access by various {\hyperref[\detokenize{Howto_output::doc}]{\sphinxcrossref{\DUrole{doc}{output commands}}}}.

\sphinxAtStartPar
This fix can ramp its rho parameter over multiple runs, using the
\sphinxstyleemphasis{start} and \sphinxstyleemphasis{stop} keywords of the {\hyperref[\detokenize{run::doc}]{\sphinxcrossref{\DUrole{doc}{run}}}} command.  See the
{\hyperref[\detokenize{run::doc}]{\sphinxcrossref{\DUrole{doc}{run}}}} command for details of how to do this.

\sphinxAtStartPar
This fix is not invoked during {\hyperref[\detokenize{minimize::doc}]{\sphinxcrossref{\DUrole{doc}{energy minimization}}}}.


\subsection{Restrictions}
\label{\detokenize{fix_tmd:restrictions}}
\sphinxAtStartPar
This fix is part of the EXTRA\sphinxhyphen{}FIX package.  It is only enabled if LAMMPS
was built with that package.  See the {\hyperref[\detokenize{Build_package::doc}]{\sphinxcrossref{\DUrole{doc}{Build package}}}} page for more info.

\sphinxAtStartPar
All TMD fixes must be listed in the input script after all integrator
fixes (nve, nvt, npt) are applied.  This ensures that atoms are moved
before their positions are corrected to comply with the constraint.

\sphinxAtStartPar
Atoms that have a TMD fix applied should not be part of a group to
which a SHAKE fix is applied.  This is because LAMMPS assumes there
are not multiple competing holonomic constraints applied to the same
atoms.

\sphinxAtStartPar
To read gzipped target files, you must compile LAMMPS with the
\sphinxhyphen{}DLAMMPS\_GZIP option.  See the {\hyperref[\detokenize{Build_settings::doc}]{\sphinxcrossref{\DUrole{doc}{Build settings}}}}
doc page for details.


\subsection{Related commands}
\label{\detokenize{fix_tmd:related-commands}}
\sphinxAtStartPar
none


\subsection{Default}
\label{\detokenize{fix_tmd:default}}
\sphinxAtStartPar
none


\bigskip\hrule\bigskip

\phantomsection\label{\detokenize{fix_tmd:schlitter1}}
\sphinxAtStartPar
\sphinxstylestrong{(Schlitter1)} Schlitter, Swegat, Mulders, “Distance\sphinxhyphen{}type reaction
coordinates for modelling activated processes”, J Molecular Modeling,
7, 171\sphinxhyphen{}177 (2001).

\phantomsection\label{\detokenize{fix_tmd:schlitter2}}
\sphinxAtStartPar
\sphinxstylestrong{(Schlitter2)} Schlitter and Klahn, “The free energy of a reaction
coordinate at multiple constraints: a concise formulation”, Molecular
Physics, 101, 3439\sphinxhyphen{}3443 (2003).

\sphinxstepscope

\index{fix ttm@\spxentry{fix ttm}}\ignorespaces 
\index{fix ttm/grid@\spxentry{fix ttm/grid}}\ignorespaces 
\index{fix ttm/mod@\spxentry{fix ttm/mod}}\ignorespaces 

\section{fix ttm command}
\label{\detokenize{fix_ttm:fix-ttm-command}}\label{\detokenize{fix_ttm:index-2}}\label{\detokenize{fix_ttm:index-1}}\label{\detokenize{fix_ttm:index-0}}\label{\detokenize{fix_ttm::doc}}

\section{fix ttm/grid command}
\label{\detokenize{fix_ttm:fix-ttm-grid-command}}

\section{fix ttm/mod command}
\label{\detokenize{fix_ttm:fix-ttm-mod-command}}

\subsection{Syntax}
\label{\detokenize{fix_ttm:syntax}}
\begin{sphinxVerbatim}[commandchars=\\\{\}]
\PYG{k}{fix }\PYG{n+nv+nvIdentifier}{ID}\PYG{+w}{ }\PYG{n+nv+nvIdentifier}{group\PYGZhy{}ID}\PYG{+w}{ }\PYG{n}{ttm}\PYG{+w}{ }\PYG{n}{seed}\PYG{+w}{ }\PYG{n}{C\PYGZus{}e}\PYG{+w}{ }\PYG{n}{rho\PYGZus{}e}\PYG{+w}{ }\PYG{n}{kappa\PYGZus{}e}\PYG{+w}{ }\PYG{n}{gamma\PYGZus{}p}\PYG{+w}{ }\PYG{n}{gamma\PYGZus{}s}\PYG{+w}{ }\PYG{n}{v\PYGZus{}0}\PYG{+w}{ }\PYG{n}{Nx}\PYG{+w}{ }\PYG{n}{Ny}\PYG{+w}{ }\PYG{n}{Nz}\PYG{+w}{ }\PYG{n}{keyword}\PYG{+w}{ }\PYG{n}{value}\PYG{+w}{ }\PYG{n}{...}
\PYG{k}{fix }\PYG{n+nv+nvIdentifier}{ID}\PYG{+w}{ }\PYG{n+nv+nvIdentifier}{group\PYGZhy{}ID}\PYG{+w}{ }\PYG{n}{ttm}\PYG{o}{/}\PYG{n}{mod}\PYG{+w}{ }\PYG{n}{seed}\PYG{+w}{ }\PYG{n}{init\PYGZus{}file}\PYG{+w}{ }\PYG{n}{Nx}\PYG{+w}{ }\PYG{n}{Ny}\PYG{+w}{ }\PYG{n}{Nz}\PYG{+w}{ }\PYG{n}{keyword}\PYG{+w}{ }\PYG{n}{value}\PYG{+w}{ }\PYG{n}{...}
\end{sphinxVerbatim}
\begin{itemize}
\item {} 
\sphinxAtStartPar
ID, group\sphinxhyphen{}ID are documented in {\hyperref[\detokenize{fix::doc}]{\sphinxcrossref{\DUrole{doc}{fix}}}} command

\item {} 
\sphinxAtStartPar
style = \sphinxstyleemphasis{ttm} or \sphinxstyleemphasis{ttm/grid} or \sphinxstyleemphasis{ttm/mod}

\item {} 
\sphinxAtStartPar
seed = random number seed to use for white noise (positive integer)

\item {} 
\sphinxAtStartPar
remaining arguments for fix ttm or fix ttm/grid
\begin{sphinxalltt}
C\_e  = electronic specific heat (energy/(electron*temperature) units)
rho\_e = electronic density (electrons/volume units)
kappa\_e = electronic thermal conductivity (energy/(time*distance*temperature) units)
gamma\_p = friction coefficient due to electron\sphinxhyphen{}ion interactions (mass/time units)
gamma\_s = friction coefficient due to electronic stopping (mass/time units)
v\_0 = electronic stopping critical velocity (velocity units)
Nx = number of thermal solve grid points in the x\sphinxhyphen{}direction (positive integer)
Ny = number of thermal solve grid points in the y\sphinxhyphen{}direction (positive integer)
Nz = number of thermal solve grid points in the z\sphinxhyphen{}direction (positive integer)
\end{sphinxalltt}

\item {} 
\sphinxAtStartPar
remaining arguments for fix ttm/mod:

\begin{sphinxVerbatim}[commandchars=\\\{\}]
init\PYGZus{}file = file with the parameters to TTM
Nx = number of thermal solve grid points in the x\PYGZhy{}direction (positive integer)
Ny = number of thermal solve grid points in the y\PYGZhy{}direction (positive integer)
Nz = number of thermal solve grid points in the z\PYGZhy{}direction (positive integer)
\end{sphinxVerbatim}

\item {} 
\sphinxAtStartPar
zero or more keyword/value(s) pairs may be appended

\item {} 
\sphinxAtStartPar
keyword = \sphinxstyleemphasis{set} or \sphinxstyleemphasis{infile} or \sphinxstyleemphasis{outfile}
\begin{sphinxalltt}
\sphinxstyleemphasis{set} value = Tinit
  Tinit = initial electronic temperature at all grid points (temperature units)
\sphinxstyleemphasis{infile} value = file.in with grid values for electronic temperatures
\sphinxstyleemphasis{outfile} values = Nout file.out
  Nout = dump grid temperatures every this many timesteps
  file.out = filename to write grid temperatures to
\end{sphinxalltt}

\end{itemize}


\subsection{Examples}
\label{\detokenize{fix_ttm:examples}}
\begin{sphinxVerbatim}[commandchars=\\\{\}]
\PYG{k}{fix }\PYG{n+nv+nvIdentifier}{2}\PYG{+w}{ }\PYG{n+nv+nvIdentifier}{all}\PYG{+w}{ }\PYG{n}{ttm}\PYG{+w}{ }\PYG{l+m}{699489}\PYG{+w}{ }\PYG{l+m}{1.0}\PYG{+w}{ }\PYG{l+m}{1.0}\PYG{+w}{ }\PYG{l+m}{10}\PYG{+w}{ }\PYG{l+m}{0.1}\PYG{+w}{ }\PYG{l+m}{0.0}\PYG{+w}{ }\PYG{l+m}{2.0}\PYG{+w}{ }\PYG{l+m}{1}\PYG{+w}{ }\PYG{l+m}{12}\PYG{+w}{ }\PYG{l+m}{1}\PYG{+w}{ }\PYG{n}{infile}\PYG{+w}{ }\PYG{n}{initial}\PYG{+w}{ }\PYG{n}{outfile}\PYG{+w}{ }\PYG{l+m}{1000}\PYG{+w}{ }\PYG{n}{T.out}
\PYG{k}{fix }\PYG{n+nv+nvIdentifier}{3}\PYG{+w}{ }\PYG{n+nv+nvIdentifier}{all}\PYG{+w}{ }\PYG{n}{ttm}\PYG{o}{/}\PYG{n}{grid}\PYG{+w}{ }\PYG{l+m}{123456}\PYG{+w}{ }\PYG{l+m}{1.0}\PYG{+w}{ }\PYG{l+m}{1.0}\PYG{+w}{ }\PYG{l+m}{1.0}\PYG{+w}{ }\PYG{l+m}{1.0}\PYG{+w}{ }\PYG{l+m}{1.0}\PYG{+w}{ }\PYG{l+m}{5.0}\PYG{+w}{ }\PYG{l+m}{5}\PYG{+w}{ }\PYG{l+m}{5}\PYG{+w}{ }\PYG{l+m}{5}\PYG{+w}{ }\PYG{n}{infile}\PYG{+w}{ }\PYG{n}{Te.in}
\PYG{k}{fix }\PYG{n+nv+nvIdentifier}{4}\PYG{+w}{ }\PYG{n+nv+nvIdentifier}{all}\PYG{+w}{ }\PYG{n}{ttm}\PYG{o}{/}\PYG{n}{mod}\PYG{+w}{ }\PYG{l+m}{34277}\PYG{+w}{ }\PYG{n}{parameters.txt}\PYG{+w}{ }\PYG{l+m}{5}\PYG{+w}{ }\PYG{l+m}{5}\PYG{+w}{ }\PYG{l+m}{5}\PYG{+w}{ }\PYG{n}{infile}\PYG{+w}{ }\PYG{n}{T\PYGZus{}init}\PYG{+w}{ }\PYG{n}{outfile}\PYG{+w}{ }\PYG{l+m}{10}\PYG{+w}{ }\PYG{n}{T\PYGZus{}out}
\end{sphinxVerbatim}

\sphinxAtStartPar
Example input scripts using these commands can be found in examples/ttm.


\subsection{Description}
\label{\detokenize{fix_ttm:description}}
\sphinxAtStartPar
Use a two\sphinxhyphen{}temperature model (TTM) to represent heat transfer through
and between electronic and atomic subsystems.  LAMMPS models the
atomic subsystem as usual with a molecular dynamics model and the
classical force field specified by the user.  The electronic subsystem
is modeled as a continuum, or a background “gas”, on a regular grid
which overlays the simulation domain.  Energy can be transferred
spatially within the grid representing the electrons.  Energy can also
be transferred between the electronic and atomic subsystems.  The
algorithm underlying this fix was derived by D. M.  Duffy
and A. M. Rutherford and is discussed in two J Physics: Condensed
Matter papers: {\hyperref[\detokenize{fix_ttm:duffy}]{\sphinxcrossref{\DUrole{std,std-ref}{(Duffy)}}}} and {\hyperref[\detokenize{fix_ttm:rutherford}]{\sphinxcrossref{\DUrole{std,std-ref}{(Rutherford)}}}}.  They used this algorithm in cascade simulations where
a primary knock\sphinxhyphen{}on atom (PKA) was initialized with a high velocity to
simulate a radiation event.

\sphinxAtStartPar
The description in this subsection applies to all 3 fix styles:
\sphinxstyleemphasis{ttm}, \sphinxstyleemphasis{ttm/grid}, and \sphinxstyleemphasis{ttm/mod}.

\sphinxAtStartPar
Fix \sphinxstyleemphasis{ttm/grid} distributes the regular grid across processors consistent
with the subdomains of atoms owned by each processor, but is otherwise
identical to fix ttm.  Note that fix \sphinxstyleemphasis{ttm} stores a copy of the grid on
each processor, which is acceptable when the overall grid is reasonably
small.  For larger grids you should use fix \sphinxstyleemphasis{ttm/grid} instead.

\sphinxAtStartPar
Fix \sphinxstyleemphasis{ttm/mod} adds options to account for external heat sources (e.g. at
a surface) and for specifying parameters that allow the electronic heat
capacity to depend strongly on electronic temperature.  It is more
expensive computationally than fix \sphinxstyleemphasis{ttm} because it treats the thermal
diffusion equation as non\sphinxhyphen{}linear.  More details on fix \sphinxstyleemphasis{ttm/mod} are
given below.

\sphinxAtStartPar
Heat transfer between the electronic and atomic subsystems is carried
out via an inhomogeneous Langevin thermostat.  Only atoms in the fix
group contribute to and are affected by this heat transfer.

\sphinxAtStartPar
This thermostatting differs from the regular Langevin thermostat
({\hyperref[\detokenize{fix_langevin::doc}]{\sphinxcrossref{\DUrole{doc}{fix langevin}}}}) in three important ways.  First,
the Langevin thermostat is applied uniformly to all atoms in the
user\sphinxhyphen{}specified group for a single target temperature, whereas the TTM
fixes apply Langevin thermostatting locally to atoms within the
volumes represented by the user\sphinxhyphen{}specified grid points with a target
temperature specific to that grid point.  Second, the Langevin
thermostat couples the temperature of the atoms to an infinite heat
reservoir, whereas the heat reservoir for the TTM fixes is finite and
represents the local electrons.  Third, the TTM fixes allow users to
specify not just one friction coefficient, but rather two independent
friction coefficients: one for the electron\sphinxhyphen{}ion interactions
(\sphinxstyleemphasis{gamma\_p}), and one for electron stopping (\sphinxstyleemphasis{gamma\_s}).

\sphinxAtStartPar
When the friction coefficient due to electron stopping, \sphinxstyleemphasis{gamma\_s}, is
non\sphinxhyphen{}zero, electron stopping effects are included for atoms moving
faster than the electron stopping critical velocity, \sphinxstyleemphasis{v\_0}.  For
further details about this algorithm, see {\hyperref[\detokenize{fix_ttm:duffy}]{\sphinxcrossref{\DUrole{std,std-ref}{(Duffy)}}}} and
{\hyperref[\detokenize{fix_ttm:rutherford}]{\sphinxcrossref{\DUrole{std,std-ref}{(Rutherford)}}}}.

\sphinxAtStartPar
Energy transport within the electronic subsystem is solved according
to the heat diffusion equation with added source terms for heat
transfer between the subsystems:
\begin{equation*}
\begin{split}C_e \rho_e \frac{\partial T_e}{\partial t} =
\bigtriangledown (\kappa_e \bigtriangledown T_e) -
g_p (T_e - T_a) + g_s T_a'\end{split}
\end{equation*}
\sphinxAtStartPar
where \(C_e\) is the specific heat, \(\rho_e\) is the density,
\(\kappa_e\) is the thermal conductivity, \sphinxstyleemphasis{T} is temperature, the
“e” and “a” subscripts represent electronic and atomic subsystems
respectively, \(g_p\) is the coupling constant for the electron\sphinxhyphen{}ion
interaction, and \(g_s\) is the electron stopping coupling
parameter.  \(C_e\), \(\rho_e\), and \(\kappa_e\) are
specified as parameters to the fix \sphinxstyleemphasis{ttm} or \sphinxstyleemphasis{ttm/grid}.  The other
quantities are derived.  The form of the heat diffusion equation used
here is almost the same as that in equation 6 of {\hyperref[\detokenize{fix_ttm:duffy}]{\sphinxcrossref{\DUrole{std,std-ref}{(Duffy)}}}},
with the exception that the electronic density is explicitly
represented, rather than being part of the specific heat parameter.

\sphinxAtStartPar
Currently, the TTM fixes assume that none of the user\sphinxhyphen{}supplied
parameters will vary with temperature. Note that {\hyperref[\detokenize{fix_ttm:duffy}]{\sphinxcrossref{\DUrole{std,std-ref}{(Duffy)}}}}
used a tanh() functional form for the temperature dependence of the
electronic specific heat, but ignored temperature dependencies of any of
the other parameters.  See more discussion below for fix \sphinxstyleemphasis{ttm/mod}.

\begin{sphinxadmonition}{note}{Note:}
\sphinxAtStartPar
These fixes do not perform time integration of the atoms in the fix
group, they only rescale their velocities.  Thus a time integration
fix such as {\hyperref[\detokenize{fix_nve::doc}]{\sphinxcrossref{\DUrole{doc}{fix nve}}}} should be used in conjunction
with these fixes.  These fixes should not normally be used on atoms
that have their temperature controlled by another thermostatting
fix, e.g. {\hyperref[\detokenize{fix_nh::doc}]{\sphinxcrossref{\DUrole{doc}{fix nvt}}}} or {\hyperref[\detokenize{fix_langevin::doc}]{\sphinxcrossref{\DUrole{doc}{fix langevin}}}}.
\end{sphinxadmonition}

\begin{sphinxadmonition}{note}{Note:}
\sphinxAtStartPar
These fixes require use of an orthogonal 3d simulation box with
periodic boundary conditions in all dimensions.  They also require
that the size and shape of the simulation box do not vary
dynamically, e.g. due to use of the {\hyperref[\detokenize{fix_nh::doc}]{\sphinxcrossref{\DUrole{doc}{fix npt}}}} command.
Likewise, the size/shape of processor subdomains cannot vary due to
dynamic load\sphinxhyphen{}balancing via use of the {\hyperref[\detokenize{fix_balance::doc}]{\sphinxcrossref{\DUrole{doc}{fix balance}}}} command.  It is possible however to load balance
before the simulation starts using the {\hyperref[\detokenize{balance::doc}]{\sphinxcrossref{\DUrole{doc}{balance}}}}
command, so that each processor has a different size subdomain.
\end{sphinxadmonition}

\sphinxAtStartPar
Periodic boundary conditions are also used in the heat equation solve
for the electronic subsystem.  This varies from the approach of
{\hyperref[\detokenize{fix_ttm:rutherford}]{\sphinxcrossref{\DUrole{std,std-ref}{(Rutherford)}}}} where the atomic subsystem was
embedded within a larger continuum representation of the electronic
subsystem.

\sphinxAtStartPar
The \sphinxstyleemphasis{set} keyword specifies a \sphinxstyleemphasis{Tinit} temperature value to initialize
the value stored on all grid points.  By default the temperatures
are all zero when the grid is created.

\sphinxAtStartPar
The \sphinxstyleemphasis{infile} keyword specifies an input file of electronic temperatures
for each grid point to be read in to initialize the grid, as an alternative
to using the \sphinxstyleemphasis{set} keyword.

\sphinxAtStartPar
The input file is a text file which may have comments starting with
the ‘\#’ character.  Each line contains four numeric columns:
ix,iy,iz,Temperature.  Empty or comment\sphinxhyphen{}only lines will be
ignored. The number of lines must be equal to the number of
user\sphinxhyphen{}specified grid points (Nx by Ny by Nz).  The ix,iy,iz are grid
point indices ranging from 1 to Nxyz inclusive in each dimension.  The
lines can appear in any order.  For example, the initial electronic
temperatures on a 1 by 2 by 3 grid could be specified in the file as
follows:

\begin{sphinxVerbatim}[commandchars=\\\{\}]
\PYGZsh{} UNITS: metal COMMENT: initial electron temperature
1 1 1 1.0
1 1 2 1.0
1 1 3 1.0
1 2 1 2.0
1 2 2 2.0
1 2 3 2.0
\end{sphinxVerbatim}

\sphinxAtStartPar
where the electronic temperatures along the y=0 plane have been set to
1.0, and the electronic temperatures along the y=1 plane have been set
to 2.0.  If all the grid point values are not specified, LAMMPS will
generate an error. LAMMPS will check if a “UNITS:” tag is in the first
line and stop with an error, if there is a mismatch with the current
units used.

\begin{sphinxadmonition}{note}{Note:}
\sphinxAtStartPar
The electronic temperature at each grid point must be a non\sphinxhyphen{}zero
positive value, both initially, and as the temperature evolves over
time.  Thus you must use either the \sphinxstyleemphasis{set} or \sphinxstyleemphasis{infile} keyword or be
restarting a simulation that used this fix previously.
\end{sphinxadmonition}

\sphinxAtStartPar
The \sphinxstyleemphasis{outfile} keyword has 2 values.  The first value \sphinxstyleemphasis{Nout} triggers
output of the electronic temperatures for each grid point every Nout
timesteps.  The second value is the filename for output, which will be
suffixed by the timestep.  The format of each output file is exactly
the same as the input temperature file. It will contain a comment in
the first line reporting the date the file was created, the LAMMPS
units setting in use, grid size and the current timestep.

\begin{sphinxadmonition}{note}{Note:}
\sphinxAtStartPar
The fix ttm/grid command does not support the \sphinxstyleemphasis{outfile} keyword.
Instead you can use the {\hyperref[\detokenize{dump::doc}]{\sphinxcrossref{\DUrole{doc}{dump grid}}}} command to output
the electronic temperature on the distributed grid to a dump file or
the {\hyperref[\detokenize{restart::doc}]{\sphinxcrossref{\DUrole{doc}{restart}}}} command which creates a file specific
to this fix which the {\hyperref[\detokenize{read_restart::doc}]{\sphinxcrossref{\DUrole{doc}{read restart}}}} command
reads.  The file has the same format as the file the \sphinxstyleemphasis{infile} option
reads.
\end{sphinxadmonition}

\sphinxAtStartPar
For the fix ttm and fix ttm/mod commands, the corresponding atomic
temperature for atoms in each grid cell can be computed and output by
the {\hyperref[\detokenize{fix_ave_chunk::doc}]{\sphinxcrossref{\DUrole{doc}{fix ave/chunk}}}} command using the
{\hyperref[\detokenize{compute_chunk_atom::doc}]{\sphinxcrossref{\DUrole{doc}{compute chunk/atom}}}} command to create a 3d
array of chunks consistent with the grid used by this fix.

\sphinxAtStartPar
For the fix ttm/grid command the same thing can be done using the
{\hyperref[\detokenize{fix_ave_grid::doc}]{\sphinxcrossref{\DUrole{doc}{fix ave/grid}}}} command and its per\sphinxhyphen{}grid values can
be output via the {\hyperref[\detokenize{dump::doc}]{\sphinxcrossref{\DUrole{doc}{dump grid}}}} command.


\bigskip\hrule\bigskip


\sphinxAtStartPar
\sphinxstylestrong{Additional details for fix ttm/mod}

\sphinxAtStartPar
Fix \sphinxstyleemphasis{ttm/mod} uses the heat diffusion equation with possible external
heat sources (e.g. laser heating in ablation simulations):
\begin{equation*}
\begin{split}C_e \rho_e \frac{\partial T_e}{\partial t} =
\bigtriangledown (\kappa_e \bigtriangledown T_e) -
g_p (T_e - T_a) + g_s T_a' + \theta (x-x_{surface})I_0 \exp(-x/l_{skin})\end{split}
\end{equation*}
\sphinxAtStartPar
where \(\theta\) is the Heaviside step function, \(I_0\) is the
(absorbed) laser pulse intensity for ablation simulations,
\(l_{skin}\) is the depth of the skin\sphinxhyphen{}layer, and all other
designations have the same meaning as in the former equation. The
duration of the pulse is set by the parameter \sphinxstyleemphasis{tau} in the \sphinxstyleemphasis{init\_file}.

\sphinxAtStartPar
Fix \sphinxstyleemphasis{ttm/mod} also allows users to specify the dependencies of
\(C_e\) and \(\kappa_e\) on the electronic temperature. The
specific heat is expressed as
\begin{equation*}
\begin{split}C_e = C_0 + (a_0 + a_1 X + a_2 X^2 + a_3 X^3 + a_4 X^4) \exp (-(AX)^2)\end{split}
\end{equation*}
\sphinxAtStartPar
where \(X = \frac{T_e}{1000}\), and the thermal conductivity is
defined as \(\kappa_e = D_e \cdot rho_e \cdot C_e\), where
\(D_e\) is the thermal diffusion coefficient.

\sphinxAtStartPar
Electronic pressure effects are included in the TTM model to account for
the blast force acting on ions because of electronic pressure gradient
(see {\hyperref[\detokenize{fix_ttm:chen}]{\sphinxcrossref{\DUrole{std,std-ref}{(Chen)}}}}, {\hyperref[\detokenize{fix_ttm:norman}]{\sphinxcrossref{\DUrole{std,std-ref}{(Norman)}}}}).  The total force
acting on an ion is:
\begin{equation*}
\begin{split}{\vec F}_i = - \partial U / \partial {\vec r}_i + {\vec
F}_{langevin} - \nabla P_e/n_{ion}\end{split}
\end{equation*}
\sphinxAtStartPar
where \(F_{langevin}\) is a force from Langevin thermostat
simulating electron\sphinxhyphen{}phonon coupling, and \(\nabla P_e/n_{ion}\) is
the electron blast force.

\sphinxAtStartPar
The electronic pressure is taken to be \(P_e = B \cdot rho_e \cdot
C_e \cdot T_e\)

\sphinxAtStartPar
The current fix \sphinxstyleemphasis{ttm/mod} implementation allows TTM simulations with a
vacuum. The vacuum region is defined as the grid cells with zero
electronic temperature. The numerical scheme does not allow energy
exchange with such cells. Since the material can expand to previously
unoccupied region in some simulations, the vacuum border can be allowed
to move. It is controlled by the \sphinxstyleemphasis{surface\_movement} parameter in the
\sphinxstyleemphasis{init\_file}. If it is set to 1, then “vacuum” cells can be changed to
“electron\sphinxhyphen{}filled” cells with the temperature \sphinxstyleemphasis{T\_e\_min} if atoms move
into them (currently only implemented for the case of 1\sphinxhyphen{}dimensional
motion of a flat surface normal to the X axis). The initial locations of
the interfaces of the electron density to the vacuum can be set in the
\sphinxstyleemphasis{init\_file} via \sphinxstyleemphasis{lsurface} and \sphinxstyleemphasis{rsurface} parameters. In this case,
electronic pressure gradient is calculated as
\begin{equation*}
\begin{split}\nabla_x P_e = \left[\frac{C_e{}T_e(x)\lambda}{(x+\lambda)^2} +
\frac{x}{x+\lambda}\frac{(C_e{}T_e)_{x+\Delta
x}-(C_e{}T_e)_{x}}{\Delta x} \right]\end{split}
\end{equation*}
\sphinxAtStartPar
where \(\lambda\) is the electron mean free path (see {\hyperref[\detokenize{fix_ttm:norman}]{\sphinxcrossref{\DUrole{std,std-ref}{(Norman)}}}}, {\hyperref[\detokenize{fix_ttm:pisarev}]{\sphinxcrossref{\DUrole{std,std-ref}{(Pisarev)}}}})

\sphinxAtStartPar
The fix \sphinxstyleemphasis{ttm/mod} parameter file \sphinxstyleemphasis{init\_file} has the following syntax.
Every line with an odd number is considered as a comment and
ignored. The lines with the even numbers are treated as follows:
\begin{sphinxalltt}
a\_0, energy/(temperature*electron) units
a\_1, energy/(temperature\textasciicircum{}2*electron) units
a\_2, energy/(temperature\textasciicircum{}3*electron) units
a\_3, energy/(temperature\textasciicircum{}4*electron) units
a\_4, energy/(temperature\textasciicircum{}5*electron) units
C\_0, energy/(temperature*electron) units
A, 1/temperature units
rho\_e, electrons/volume units
D\_e, length\textasciicircum{}2/time units
gamma\_p, mass/time units
gamma\_s, mass/time units
v\_0, length/time units
I\_0, energy/(time*length\textasciicircum{}2) units
lsurface, electron grid units (positive integer)
rsurface, electron grid units (positive integer)
l\_skin, length units
tau, time units
B, dimensionless
lambda, length units
n\_ion, ions/volume units
surface\_movement: 0 to disable tracking of surface motion, 1 to enable
T\_e\_min, temperature units
\end{sphinxalltt}


\bigskip\hrule\bigskip



\subsection{Restart, fix\_modify, output, run start/stop, minimize info}
\label{\detokenize{fix_ttm:restart-fix-modify-output-run-start-stop-minimize-info}}
\sphinxAtStartPar
The fix ttm and fix ttm/mod commands write the state of the electronic
subsystem and the energy exchange between the subsystems to
{\hyperref[\detokenize{restart::doc}]{\sphinxcrossref{\DUrole{doc}{binary restart files}}}}.  The fix ttm/grid command does
not yet support writing of its distributed grid to a restart file.

\sphinxAtStartPar
See the {\hyperref[\detokenize{read_restart::doc}]{\sphinxcrossref{\DUrole{doc}{read\_restart}}}} command for info on how to
re\sphinxhyphen{}specify a fix in an input script that reads a restart file, so that
the operation of the fix continues in an uninterrupted fashion.  Note
that the restart script must define the same size grid as the original
script.

\sphinxAtStartPar
The fix ttm/grid command also outputs an auxiliary file each time a
restart file is written, with the electron temperatures for each grid
cell.  The format of this file is the same as that read by the
\sphinxstyleemphasis{infile} option explained above.  The filename is the same as the
restart filename with “.ttm” appended.  This auxiliary file can be
read in for a restarted run by using the \sphinxstyleemphasis{infile} option for the fix
ttm/grid command, following the {\hyperref[\detokenize{read_restart::doc}]{\sphinxcrossref{\DUrole{doc}{read\_restart}}}}
command.

\sphinxAtStartPar
None of the {\hyperref[\detokenize{fix_modify::doc}]{\sphinxcrossref{\DUrole{doc}{fix\_modify}}}} options are relevant to
these fixes.

\sphinxAtStartPar
These fixes compute 2 output quantities stored in a vector of length
2, which can be accessed by various {\hyperref[\detokenize{Howto_output::doc}]{\sphinxcrossref{\DUrole{doc}{output commands}}}}.  The first quantity is the total energy of the
electronic subsystem.  The second quantity is the energy transferred
from the electronic to the atomic subsystem on that timestep. Note
that the velocity verlet integrator applies the fix ttm forces to the
atomic subsystem as two half\sphinxhyphen{}step velocity updates: one on the current
timestep and one on the subsequent timestep.  Consequently, the change
in the atomic subsystem energy is lagged by half a timestep relative
to the change in the electronic subsystem energy. As a result of this,
users may notice slight fluctuations in the sum of the atomic and
electronic subsystem energies reported at the end of the timestep.

\sphinxAtStartPar
The vector values calculated are “extensive”.

\sphinxAtStartPar
The fix ttm/grid command also outputs a per\sphinxhyphen{}grid vector which stores
the electron temperature for each grid cell in temperature {\hyperref[\detokenize{units::doc}]{\sphinxcrossref{\DUrole{doc}{units}}}}. which can be accessed by various {\hyperref[\detokenize{Howto_output::doc}]{\sphinxcrossref{\DUrole{doc}{output commands}}}}.  The length of the vector (distributed across all
processors) is Nx * Ny * Nz.  For access by other commands, the name
of the single grid produced by fix ttm/grid is “grid”.  The name of
its per\sphinxhyphen{}grid data is “data”.

\sphinxAtStartPar
No parameter of the fixes can be used with the \sphinxstyleemphasis{start/stop} keywords
of the {\hyperref[\detokenize{run::doc}]{\sphinxcrossref{\DUrole{doc}{run}}}} command.  The fixes are not invoked during
{\hyperref[\detokenize{minimize::doc}]{\sphinxcrossref{\DUrole{doc}{energy minimization}}}}.


\subsection{Restrictions}
\label{\detokenize{fix_ttm:restrictions}}
\sphinxAtStartPar
All these fixes are part of the EXTRA\sphinxhyphen{}FIX package. They are only
enabled if LAMMPS was built with that package.  See the {\hyperref[\detokenize{Build_package::doc}]{\sphinxcrossref{\DUrole{doc}{Build
package}}}} page for more info.

\sphinxAtStartPar
As mentioned above, these fixes require 3d simulations and orthogonal
simulation boxes periodic in all 3 dimensions.

\sphinxAtStartPar
These fixes used a random number generator to Langevin thermostat the
electron temperature.  This means you will not get identical answers
when running on different numbers of processors or when restarting a
simulation (even on the same number of processors).  However, in a
statistical sense, simulations on different processor counts and
restarted simulation should produce results which are statistically
the same.


\subsection{Related commands}
\label{\detokenize{fix_ttm:related-commands}}
\sphinxAtStartPar
{\hyperref[\detokenize{fix_langevin::doc}]{\sphinxcrossref{\DUrole{doc}{fix langevin}}}}, {\hyperref[\detokenize{fix_dt_reset::doc}]{\sphinxcrossref{\DUrole{doc}{fix dt/reset}}}}


\subsection{Default}
\label{\detokenize{fix_ttm:default}}
\sphinxAtStartPar
none


\bigskip\hrule\bigskip

\phantomsection\label{\detokenize{fix_ttm:duffy}}
\sphinxAtStartPar
\sphinxstylestrong{(Duffy)} D M Duffy and A M Rutherford, J. Phys.: Condens. Matter, 19,
016207\sphinxhyphen{}016218 (2007).

\phantomsection\label{\detokenize{fix_ttm:rutherford}}
\sphinxAtStartPar
\sphinxstylestrong{(Rutherford)} A M Rutherford and D M Duffy, J. Phys.:
Condens. Matter, 19, 496201\sphinxhyphen{}496210 (2007).

\phantomsection\label{\detokenize{fix_ttm:chen}}
\sphinxAtStartPar
\sphinxstylestrong{(Chen)} J Chen, D Tzou and J Beraun, Int. J. Heat
Mass Transfer, 49, 307\sphinxhyphen{}316 (2006).

\phantomsection\label{\detokenize{fix_ttm:norman}}
\sphinxAtStartPar
\sphinxstylestrong{(Norman)} G E Norman, S V Starikov, V V Stegailov et al., Contrib.
Plasma Phys., 53, 129\sphinxhyphen{}139 (2013).

\phantomsection\label{\detokenize{fix_ttm:pisarev}}
\sphinxAtStartPar
\sphinxstylestrong{(Pisarev)} V V Pisarev and S V Starikov, J. Phys.: Condens. Matter, 26,
475401 (2014).

\sphinxstepscope

\index{fix tune/kspace@\spxentry{fix tune/kspace}}\ignorespaces 

\section{fix tune/kspace command}
\label{\detokenize{fix_tune_kspace:fix-tune-kspace-command}}\label{\detokenize{fix_tune_kspace:index-0}}\label{\detokenize{fix_tune_kspace::doc}}

\subsection{Syntax}
\label{\detokenize{fix_tune_kspace:syntax}}
\begin{sphinxVerbatim}[commandchars=\\\{\}]
\PYG{k}{fix }\PYG{n+nv+nvIdentifier}{ID}\PYG{+w}{ }\PYG{n+nv+nvIdentifier}{group\PYGZhy{}ID}\PYG{+w}{ }\PYG{n}{tune}\PYG{o}{/}\PYG{n}{kspace}\PYG{+w}{ }\PYG{n}{N}
\end{sphinxVerbatim}
\begin{itemize}
\item {} 
\sphinxAtStartPar
ID, group\sphinxhyphen{}ID are documented in {\hyperref[\detokenize{fix::doc}]{\sphinxcrossref{\DUrole{doc}{fix}}}} command

\item {} 
\sphinxAtStartPar
tune/kspace = style name of this fix command

\item {} 
\sphinxAtStartPar
N = invoke this fix every N steps

\end{itemize}


\subsection{Examples}
\label{\detokenize{fix_tune_kspace:examples}}
\begin{sphinxVerbatim}[commandchars=\\\{\}]
\PYG{k}{fix }\PYG{n+nv+nvIdentifier}{2}\PYG{+w}{ }\PYG{n+nv+nvIdentifier}{all}\PYG{+w}{ }\PYG{n}{tune}\PYG{o}{/}\PYG{n}{kspace}\PYG{+w}{ }\PYG{l+m}{100}
\end{sphinxVerbatim}


\subsection{Description}
\label{\detokenize{fix_tune_kspace:description}}
\sphinxAtStartPar
This fix tests each kspace style (Ewald, PPPM, and MSM), and
automatically selects the fastest style to use for the remainder
of the run. If the fastest style is Ewald or PPPM, the fix also
adjusts the Coulombic cutoff towards optimal speed. Future versions
of this fix will automatically select other kspace parameters
to use for maximum simulation speed. The kspace parameters may
include the style, cutoff, grid points in each direction, order,
Ewald parameter, MSM parallelization cut\sphinxhyphen{}point, MPI tasks to use, etc.

\sphinxAtStartPar
The rationale for this fix is to provide the user with
as\sphinxhyphen{}fast\sphinxhyphen{}as\sphinxhyphen{}possible simulations that include long\sphinxhyphen{}range electrostatics
(kspace) while meeting the user\sphinxhyphen{}prescribed accuracy requirement. A
simple heuristic could never capture the optimal combination of
parameters for every possible run\sphinxhyphen{}time scenario. But by performing
short tests of various kspace parameter sets, this fix allows
parameters to be tailored specifically to the user’s machine, MPI
ranks, use of threading or accelerators, the simulated system, and the
simulation details. In addition, it is possible that parameters could
be evolved with the simulation on\sphinxhyphen{}the\sphinxhyphen{}fly, which is useful for systems
that are dynamically evolving (e.g. changes in box size/shape or
number of particles).

\sphinxAtStartPar
When this fix is invoked, LAMMPS will perform short timed tests of
various parameter sets to determine the optimal parameters. Tests are
performed on\sphinxhyphen{}the\sphinxhyphen{}fly, with a new test initialized every N steps. N should
be chosen large enough so that adequate CPU time lapses between tests,
thereby providing statistically significant timings. But N should not be
chosen to be so large that an unfortunate parameter set test takes an
inordinate amount of wall time to complete. An N of 100 for most problems
seems reasonable. Once an optimal parameter set is found, that set is
used for the remainder of the run.

\sphinxAtStartPar
This fix uses heuristics to guide it’s selection of parameter sets to test,
but the actual timed results will be used to decide which set to use in the
simulation.

\sphinxAtStartPar
It is not necessary to discard trajectories produced using sub\sphinxhyphen{}optimal
parameter sets, or a mix of various parameter sets, since the user\sphinxhyphen{}prescribed
accuracy will have been maintained throughout. However, some users may prefer
to use this fix only to discover the optimal parameter set for a given setup
that can then be used on subsequent production runs.

\sphinxAtStartPar
This fix starts with kspace parameters that are set by the user with the
{\hyperref[\detokenize{kspace_style::doc}]{\sphinxcrossref{\DUrole{doc}{kspace\_style}}}} and {\hyperref[\detokenize{kspace_modify::doc}]{\sphinxcrossref{\DUrole{doc}{kspace\_modify}}}}
commands. The prescribed accuracy will be maintained by this fix throughout
the simulation.

\sphinxAtStartPar
None of the {\hyperref[\detokenize{fix_modify::doc}]{\sphinxcrossref{\DUrole{doc}{fix\_modify}}}} options are relevant to this
fix.

\sphinxAtStartPar
No parameter of this fix can be used with the \sphinxstyleemphasis{start/stop} keywords of
the {\hyperref[\detokenize{run::doc}]{\sphinxcrossref{\DUrole{doc}{run}}}} command.  This fix is not invoked during {\hyperref[\detokenize{minimize::doc}]{\sphinxcrossref{\DUrole{doc}{energy minimization}}}}.


\subsection{Restrictions}
\label{\detokenize{fix_tune_kspace:restrictions}}
\sphinxAtStartPar
This fix is part of the KSPACE package.  It is only enabled if LAMMPS
was built with that package.  See the {\hyperref[\detokenize{Build_package::doc}]{\sphinxcrossref{\DUrole{doc}{Build package}}}} page for more info.

\sphinxAtStartPar
Do not set “neigh\_modify once yes” or else this fix will never be
called.  Reneighboring is required.

\sphinxAtStartPar
This fix is not compatible with a hybrid pair style, long\sphinxhyphen{}range dispersion,
TIP4P water support, or long\sphinxhyphen{}range point dipole support.


\subsection{Related commands}
\label{\detokenize{fix_tune_kspace:related-commands}}
\sphinxAtStartPar
{\hyperref[\detokenize{kspace_style::doc}]{\sphinxcrossref{\DUrole{doc}{kspace\_style}}}}, {\hyperref[\detokenize{boundary::doc}]{\sphinxcrossref{\DUrole{doc}{boundary}}}}
{\hyperref[\detokenize{kspace_modify::doc}]{\sphinxcrossref{\DUrole{doc}{kspace\_modify}}}}, {\hyperref[\detokenize{pair_lj_cut_coul::doc}]{\sphinxcrossref{\DUrole{doc}{pair\_style lj/cut/coul/long}}}}, {\hyperref[\detokenize{pair_charmm::doc}]{\sphinxcrossref{\DUrole{doc}{pair\_style lj/charmm/coul/long}}}}, {\hyperref[\detokenize{pair_lj_long::doc}]{\sphinxcrossref{\DUrole{doc}{pair\_style lj/long}}}}, {\hyperref[\detokenize{pair_lj_long::doc}]{\sphinxcrossref{\DUrole{doc}{pair\_style lj/long/coul/long}}}},
{\hyperref[\detokenize{pair_buck::doc}]{\sphinxcrossref{\DUrole{doc}{pair\_style buck/coul/long}}}}


\subsection{Default}
\label{\detokenize{fix_tune_kspace:default}}
\sphinxstepscope

\index{fix vector@\spxentry{fix vector}}\ignorespaces 

\section{fix vector command}
\label{\detokenize{fix_vector:fix-vector-command}}\label{\detokenize{fix_vector:index-0}}\label{\detokenize{fix_vector::doc}}

\subsection{Syntax}
\label{\detokenize{fix_vector:syntax}}
\begin{sphinxVerbatim}[commandchars=\\\{\}]
\PYG{k}{fix }\PYG{n+nv+nvIdentifier}{ID}\PYG{+w}{ }\PYG{n+nv+nvIdentifier}{group\PYGZhy{}ID}\PYG{+w}{ }\PYG{n}{vector}\PYG{+w}{ }\PYG{n}{Nevery}\PYG{+w}{ }\PYG{n}{value1}\PYG{+w}{ }\PYG{n}{value2}\PYG{+w}{ }\PYG{n}{...}\PYG{+w}{ }\PYG{n}{keyword}\PYG{+w}{ }\PYG{n}{args}\PYG{+w}{ }\PYG{n}{...}
\end{sphinxVerbatim}
\begin{itemize}
\item {} 
\sphinxAtStartPar
ID, group\sphinxhyphen{}ID are documented in {\hyperref[\detokenize{fix::doc}]{\sphinxcrossref{\DUrole{doc}{fix}}}} command

\item {} 
\sphinxAtStartPar
vector = style name of this fix command

\item {} 
\sphinxAtStartPar
Nevery = use input values every this many timesteps

\item {} 
\sphinxAtStartPar
one or more input values can be listed

\item {} 
\sphinxAtStartPar
value = c\_ID, c\_ID{[}N{]}, f\_ID, f\_ID{[}N{]}, v\_name

\begin{sphinxVerbatim}[commandchars=\\\{\}]
c\PYGZus{}ID = global scalar calculated by a compute with ID
c\PYGZus{}ID[I] = Ith component of global vector calculated by a compute with ID
f\PYGZus{}ID = global scalar calculated by a fix with ID
f\PYGZus{}ID[I] = Ith component of global vector calculated by a fix with ID
v\PYGZus{}name = value calculated by an equal\PYGZhy{}style variable with name
v\PYGZus{}name[I] = Ith component of vector\PYGZhy{}style variable with name
\end{sphinxVerbatim}

\item {} 
\sphinxAtStartPar
zero or more keyword/args pairs may be appended

\item {} 
\sphinxAtStartPar
keyword = \sphinxstyleemphasis{nmax}
\begin{sphinxalltt}
\sphinxstyleemphasis{nmax} length = set maximal length of vector to \textless{}length\textgreater{}
\end{sphinxalltt}

\end{itemize}


\subsection{Examples}
\label{\detokenize{fix_vector:examples}}
\begin{sphinxVerbatim}[commandchars=\\\{\}]
\PYG{k}{fix }\PYG{n+nv+nvIdentifier}{1}\PYG{+w}{ }\PYG{n+nv+nvIdentifier}{all}\PYG{+w}{ }\PYG{n}{vector}\PYG{+w}{ }\PYG{l+m}{100}\PYG{+w}{ }\PYG{n}{c\PYGZus{}myTemp}
\PYG{k}{fix }\PYG{n+nv+nvIdentifier}{1}\PYG{+w}{ }\PYG{n+nv+nvIdentifier}{all}\PYG{+w}{ }\PYG{n}{vector}\PYG{+w}{ }\PYG{l+m}{5}\PYG{+w}{ }\PYG{n}{c\PYGZus{}myTemp}\PYG{+w}{ }\PYG{n}{v\PYGZus{}integral}
\PYG{k}{fix }\PYG{n+nv+nvIdentifier}{1}\PYG{+w}{ }\PYG{n+nv+nvIdentifier}{all}\PYG{+w}{ }\PYG{n}{vector}\PYG{+w}{ }\PYG{l+m}{50}\PYG{+w}{ }\PYG{n}{c\PYGZus{}myTemp}\PYG{+w}{ }\PYG{n}{nmax}\PYG{+w}{ }\PYG{l+m}{200}
\end{sphinxVerbatim}


\subsection{Description}
\label{\detokenize{fix_vector:description}}
\sphinxAtStartPar
Use one or more global values as inputs every few timesteps, and simply
store them as a sequence.  For a single specified value, the values are
stored as a global vector of growing length.  For multiple specified
values, they are stored as rows in a global array, whose number of rows
is growing.  The resulting vector or array can be used by other
{\hyperref[\detokenize{Howto_output::doc}]{\sphinxcrossref{\DUrole{doc}{output commands}}}}.

\sphinxAtStartPar
The optional \sphinxstyleemphasis{nmax} keyword can be used to restrict the length of the
vector to the given \sphinxstyleemphasis{length} value.  Once the restricted vector is
filled, the oldest entry will be discarded when a entry is added.

\sphinxAtStartPar
One way to to use this command is to accumulate a vector that is
numerically integrated using the {\hyperref[\detokenize{variable::doc}]{\sphinxcrossref{\DUrole{doc}{variable trap()}}}}
function. For example, the velocity auto\sphinxhyphen{}correlation function (VACF)
can be integrated, to yield a diffusion coefficient, as follows:

\begin{sphinxVerbatim}[commandchars=\\\{\}]
\PYG{k}{compute         }\PYG{n+nv+nvIdentifier}{2}\PYG{+w}{ }\PYG{n+nv+nvIdentifier}{all}\PYG{+w}{ }\PYG{n}{vacf}
\PYG{k}{fix             }\PYG{n+nv+nvIdentifier}{5}\PYG{+w}{ }\PYG{n+nv+nvIdentifier}{all}\PYG{+w}{ }\PYG{n}{vector}\PYG{+w}{ }\PYG{l+m}{1}\PYG{+w}{ }\PYG{n}{c\PYGZus{}2[4]}
\PYG{k}{variable        }\PYG{n+nv+nvIdentifier}{diff}\PYG{+w}{ }\PYG{n}{equal}\PYG{+w}{ }\PYG{n}{dt}\PYG{o}{*}\PYG{n}{trap}\PYG{n+nv}{(}\PYG{n+nv}{f\PYGZus{}5}\PYG{n+nv}{)}
\PYG{k}{thermo\PYGZus{}style}\PYG{+w}{    }\PYG{n}{custom}\PYG{+w}{ }\PYG{n}{step}\PYG{+w}{ }\PYG{n}{v\PYGZus{}diff}
\end{sphinxVerbatim}

\sphinxAtStartPar
The group specified with this command is ignored.  However, note that
specified values may represent calculations performed by computes and
fixes which store their own “group” definitions.

\sphinxAtStartPar
Each listed value can be the result of a {\hyperref[\detokenize{compute::doc}]{\sphinxcrossref{\DUrole{doc}{compute}}}} or
{\hyperref[\detokenize{fix::doc}]{\sphinxcrossref{\DUrole{doc}{fix}}}} or the evaluation of an equal\sphinxhyphen{}style or vector\sphinxhyphen{}style
{\hyperref[\detokenize{variable::doc}]{\sphinxcrossref{\DUrole{doc}{variable}}}}.  In each case, the compute, fix, or variable
must produce a global quantity, not a per\sphinxhyphen{}atom or local quantity.  And
the global quantity must be a scalar, not a vector or array.

\sphinxAtStartPar
{\hyperref[\detokenize{compute::doc}]{\sphinxcrossref{\DUrole{doc}{Computes}}}} that produce global quantities are those which
do not have the word \sphinxstyleemphasis{atom} in their style name.  Only a few
{\hyperref[\detokenize{fix::doc}]{\sphinxcrossref{\DUrole{doc}{fixes}}}} produce global quantities.  See the doc pages for
individual fixes for info on which ones produce such values.
{\hyperref[\detokenize{variable::doc}]{\sphinxcrossref{\DUrole{doc}{Variables}}}} of style \sphinxstyleemphasis{equal} or \sphinxstyleemphasis{vector} are the only
ones that can be used with this fix.  Variables of style \sphinxstyleemphasis{atom} cannot
be used, since they produce per\sphinxhyphen{}atom values.

\sphinxAtStartPar
The \sphinxstyleemphasis{Nevery} argument specifies on what timesteps the input values
will be used in order to be stored.  Only timesteps that are a
multiple of \sphinxstyleemphasis{Nevery}, including timestep 0, will contribute values.

\begin{sphinxadmonition}{note}{Note:}\begin{quote}

\sphinxAtStartPar
If \sphinxstyleemphasis{Nevery} is a small number and the simulation runs for many
steps, the accumulated vector or array can become very large and
thus consume a lot of memory. The implementation limit is about
2 billion entries. Using the \sphinxstyleemphasis{nmax} keyword mentioned above can
avoid that by limiting the size of the vector.
\end{quote}
\end{sphinxadmonition}

\sphinxAtStartPar
Note that if you perform multiple runs, using the “pre no” option of
the {\hyperref[\detokenize{run::doc}]{\sphinxcrossref{\DUrole{doc}{run}}}} command to avoid initialization on subsequent runs,
then you need to use the \sphinxstyleemphasis{stop} keyword with the first {\hyperref[\detokenize{run::doc}]{\sphinxcrossref{\DUrole{doc}{run}}}}
command with a timestep value that encompasses all the runs.  This is
so that the vector or array stored by this fix can be allocated to a
sufficient size.


\bigskip\hrule\bigskip


\sphinxAtStartPar
If a value begins with “c\_”, a compute ID must follow which has been
previously defined in the input script.  If no bracketed term is
appended, the global scalar calculated by the compute is used.  If a
bracketed term is appended, the Ith element of the global vector
calculated by the compute is used.

\sphinxAtStartPar
Note that there is a {\hyperref[\detokenize{compute_reduce::doc}]{\sphinxcrossref{\DUrole{doc}{compute reduce}}}} command
which can sum per\sphinxhyphen{}atom quantities into a global scalar or vector which
can thus be accessed by fix vector.  Or it can be a compute defined not
in your input script, but by {\hyperref[\detokenize{thermo_style::doc}]{\sphinxcrossref{\DUrole{doc}{thermodynamic output}}}}
or other fixes such as {\hyperref[\detokenize{fix_nh::doc}]{\sphinxcrossref{\DUrole{doc}{fix nvt}}}} or {\hyperref[\detokenize{fix_temp_rescale::doc}]{\sphinxcrossref{\DUrole{doc}{fix temp/rescale}}}}.  See the doc pages for these commands which give
the IDs of these computes.  Users can also write code for their own
compute styles and {\hyperref[\detokenize{Modify::doc}]{\sphinxcrossref{\DUrole{doc}{add them to LAMMPS}}}}.

\sphinxAtStartPar
If a value begins with “f\_”, a fix ID must follow which has been
previously defined in the input script.  If no bracketed term is
appended, the global scalar calculated by the fix is used.  If a
bracketed term is appended, the Ith element of the global vector
calculated by the fix is used.

\sphinxAtStartPar
Note that some fixes only produce their values on certain timesteps,
which must be compatible with \sphinxstyleemphasis{Nevery}, else an error will result.
Users can also write code for their own fix styles and {\hyperref[\detokenize{Modify::doc}]{\sphinxcrossref{\DUrole{doc}{add them to
LAMMPS}}}}.

\sphinxAtStartPar
If a value begins with “v\_”, a variable name must follow which has
been previously defined in the input script.  An equal\sphinxhyphen{}style or
vector\sphinxhyphen{}style variable can be referenced; the latter requires a
bracketed term to specify the Ith element of the vector calculated by
the variable.  See the {\hyperref[\detokenize{variable::doc}]{\sphinxcrossref{\DUrole{doc}{variable}}}} command for details.
Note that variables of style \sphinxstyleemphasis{equal} and \sphinxstyleemphasis{vector} define a formula
which can reference individual atom properties or thermodynamic
keywords, or they can invoke other computes, fixes, or variables when
they are evaluated, so this is a very general means of specifying
quantities to be stored by fix vector.


\bigskip\hrule\bigskip



\subsection{Restart, fix\_modify, output, run start/stop, minimize info}
\label{\detokenize{fix_vector:restart-fix-modify-output-run-start-stop-minimize-info}}
\sphinxAtStartPar
No information about this fix is written to {\hyperref[\detokenize{restart::doc}]{\sphinxcrossref{\DUrole{doc}{binary restart files}}}}.  None of the {\hyperref[\detokenize{fix_modify::doc}]{\sphinxcrossref{\DUrole{doc}{fix\_modify}}}} options are
relevant to this fix.

\sphinxAtStartPar
This fix produces a global vector or global array which can be
accessed by various {\hyperref[\detokenize{Howto_output::doc}]{\sphinxcrossref{\DUrole{doc}{output commands}}}}.  The values
can only be accessed on timesteps that are multiples of \sphinxstyleemphasis{Nevery}.

\sphinxAtStartPar
A vector is produced if only a single input value is specified.
An array is produced if multiple input values are specified.
The length of the vector or the number of rows in the array grows
by 1 every \sphinxstyleemphasis{Nevery} timesteps.

\sphinxAtStartPar
If the fix produces a vector, then the entire vector will be either
“intensive” or “extensive”, depending on whether the values stored in
the vector are “intensive” or “extensive”.  If the fix produces an
array, then all elements in the array must be the same, either
“intensive” or “extensive”.  If a compute or fix provides the value
stored, then the compute or fix determines whether the value is
intensive or extensive; see the page for that compute or fix for further
info.  Values produced by a variable are treated as intensive.

\sphinxAtStartPar
This fix can allocate storage for stored values accumulated over
multiple runs, using the \sphinxstyleemphasis{start} and \sphinxstyleemphasis{stop} keywords of the {\hyperref[\detokenize{run::doc}]{\sphinxcrossref{\DUrole{doc}{run}}}} command.  See the {\hyperref[\detokenize{run::doc}]{\sphinxcrossref{\DUrole{doc}{run}}}} command for details of how to
do this.  If using the {\hyperref[\detokenize{run::doc}]{\sphinxcrossref{\DUrole{doc}{run pre no}}}} command option, this is
required to allow the fix to allocate sufficient storage for stored
values.

\sphinxAtStartPar
This fix is not invoked during {\hyperref[\detokenize{minimize::doc}]{\sphinxcrossref{\DUrole{doc}{energy minimization}}}}.


\subsection{Restrictions}
\label{\detokenize{fix_vector:restrictions}}\begin{quote}

\sphinxAtStartPar
none
\end{quote}


\subsection{Related commands}
\label{\detokenize{fix_vector:related-commands}}
\sphinxAtStartPar
{\hyperref[\detokenize{compute::doc}]{\sphinxcrossref{\DUrole{doc}{compute}}}}, {\hyperref[\detokenize{variable::doc}]{\sphinxcrossref{\DUrole{doc}{variable}}}}


\subsection{Defaults}
\label{\detokenize{fix_vector:defaults}}
\sphinxAtStartPar
The default value of \sphinxstyleemphasis{nmax} is deduced from the number of steps
in a run (or multiple runs when using the \sphinxstyleemphasis{start} and \sphinxstyleemphasis{stop}
keywords of the {\hyperref[\detokenize{run::doc}]{\sphinxcrossref{\DUrole{doc}{run command}}}}) divided by the choice
of \sphinxstyleemphasis{Nevery} plus 1.

\sphinxstepscope

\index{fix viscosity@\spxentry{fix viscosity}}\ignorespaces 

\section{fix viscosity command}
\label{\detokenize{fix_viscosity:fix-viscosity-command}}\label{\detokenize{fix_viscosity:index-0}}\label{\detokenize{fix_viscosity::doc}}

\subsection{Syntax}
\label{\detokenize{fix_viscosity:syntax}}
\begin{sphinxVerbatim}[commandchars=\\\{\}]
\PYG{k}{fix }\PYG{n+nv+nvIdentifier}{ID}\PYG{+w}{ }\PYG{n+nv+nvIdentifier}{group\PYGZhy{}ID}\PYG{+w}{ }\PYG{n}{viscosity}\PYG{+w}{ }\PYG{n}{N}\PYG{+w}{ }\PYG{n}{vdim}\PYG{+w}{ }\PYG{n}{pdim}\PYG{+w}{ }\PYG{n}{Nbin}\PYG{+w}{ }\PYG{n}{keyword}\PYG{+w}{ }\PYG{n}{value}\PYG{+w}{ }\PYG{n}{...}
\end{sphinxVerbatim}
\begin{itemize}
\item {} 
\sphinxAtStartPar
ID, group\sphinxhyphen{}ID are documented in {\hyperref[\detokenize{fix::doc}]{\sphinxcrossref{\DUrole{doc}{fix}}}} command

\item {} 
\sphinxAtStartPar
viscosity = style name of this fix command

\item {} 
\sphinxAtStartPar
N = perform momentum exchange every N steps

\item {} 
\sphinxAtStartPar
vdim = \sphinxstyleemphasis{x} or \sphinxstyleemphasis{y} or \sphinxstyleemphasis{z} = which momentum component to exchange

\item {} 
\sphinxAtStartPar
pdim = \sphinxstyleemphasis{x} or \sphinxstyleemphasis{y} or \sphinxstyleemphasis{z} = direction of momentum transfer

\item {} 
\sphinxAtStartPar
Nbin = \# of layers in pdim direction (must be even number)

\item {} 
\sphinxAtStartPar
zero or more keyword/value pairs may be appended

\item {} 
\sphinxAtStartPar
keyword = \sphinxstyleemphasis{swap} or \sphinxstyleemphasis{vtarget}
\begin{sphinxalltt}
\sphinxstyleemphasis{swap} value = Nswap = number of swaps to perform every N steps
\sphinxstyleemphasis{vtarget} value = V or INF = target velocity of swap partners (velocity units)
\end{sphinxalltt}

\end{itemize}


\subsection{Examples}
\label{\detokenize{fix_viscosity:examples}}
\begin{sphinxVerbatim}[commandchars=\\\{\}]
\PYG{k}{fix }\PYG{n+nv+nvIdentifier}{1}\PYG{+w}{ }\PYG{n+nv+nvIdentifier}{all}\PYG{+w}{ }\PYG{n}{viscosity}\PYG{+w}{ }\PYG{l+m}{100}\PYG{+w}{ }\PYG{n}{x}\PYG{+w}{ }\PYG{n}{z}\PYG{+w}{ }\PYG{l+m}{20}
\PYG{k}{fix }\PYG{n+nv+nvIdentifier}{1}\PYG{+w}{ }\PYG{n+nv+nvIdentifier}{all}\PYG{+w}{ }\PYG{n}{viscosity}\PYG{+w}{ }\PYG{l+m}{50}\PYG{+w}{ }\PYG{n}{x}\PYG{+w}{ }\PYG{n}{z}\PYG{+w}{ }\PYG{l+m}{20}\PYG{+w}{ }\PYG{n}{swap}\PYG{+w}{ }\PYG{l+m}{2}\PYG{+w}{ }\PYG{n}{vtarget}\PYG{+w}{ }\PYG{l+m}{1.5}
\end{sphinxVerbatim}


\subsection{Description}
\label{\detokenize{fix_viscosity:description}}
\sphinxAtStartPar
Use the Muller\sphinxhyphen{}Plathe algorithm described in {\hyperref[\detokenize{fix_viscosity:muller-plathe2}]{\sphinxcrossref{\DUrole{std,std-ref}{this paper}}}} to exchange momenta between two particles in
different regions of the simulation box every N steps.  This induces a
shear velocity profile in the system.  As described below this enables
a viscosity of the fluid to be calculated.  This algorithm is
sometimes called a reverse non\sphinxhyphen{}equilibrium MD (reverse NEMD) approach
to computing viscosity.  This is because the usual NEMD approach is to
impose a shear velocity profile on the system and measure the response
via an off\sphinxhyphen{}diagonal component of the stress tensor, which is
proportional to the momentum flux.  In the Muller\sphinxhyphen{}Plathe method, the
momentum flux is imposed, and the shear velocity profile is the
system’s response.

\sphinxAtStartPar
The simulation box is divided into \sphinxstyleemphasis{Nbin} layers in the \sphinxstyleemphasis{pdim}
direction, where the layer 1 is at the low end of that dimension and
the layer \sphinxstyleemphasis{Nbin} is at the high end.  Every N steps, Nswap pairs of
atoms are chosen in the following manner.  Only atoms in the fix group
are considered.  Nswap atoms in layer 1 with positive velocity
components in the \sphinxstyleemphasis{vdim} direction closest to the target value \sphinxstyleemphasis{V} are
selected.  Similarly, Nswap atoms in the “middle” layer (see below) with
negative velocity components in the \sphinxstyleemphasis{vdim} direction closest to the
negative of the target value \sphinxstyleemphasis{V} are selected.  The two sets of Nswap
atoms are paired up and their \sphinxstyleemphasis{vdim} momenta components are swapped
within each pair.  This resets their velocities, typically in opposite
directions.  Over time, this induces a shear velocity profile in the
system which can be measured using commands such as the following,
which writes the profile to the file tmp.profile:

\begin{sphinxVerbatim}[commandchars=\\\{\}]
\PYG{k}{compute }\PYG{n+nv+nvIdentifier}{layers}\PYG{+w}{ }\PYG{n+nv+nvIdentifier}{all}\PYG{+w}{ }\PYG{n}{chunk}\PYG{o}{/}\PYG{n}{atom}\PYG{+w}{ }\PYG{n}{bin}\PYG{o}{/}\PYG{l+m}{1}\PYG{n}{d}\PYG{+w}{ }\PYG{n}{z}\PYG{+w}{ }\PYG{n}{lower}\PYG{+w}{ }\PYG{l+m}{0.05}\PYG{+w}{ }\PYG{n}{units}\PYG{+w}{ }\PYG{n}{reduced}
\PYG{k}{fix }\PYG{n+nv+nvIdentifier}{f1}\PYG{+w}{ }\PYG{n+nv+nvIdentifier}{all}\PYG{+w}{ }\PYG{n}{ave}\PYG{o}{/}\PYG{n}{chunk}\PYG{+w}{ }\PYG{l+m}{100}\PYG{+w}{ }\PYG{l+m}{10}\PYG{+w}{ }\PYG{l+m}{1000}\PYG{+w}{ }\PYG{n}{layers}\PYG{+w}{ }\PYG{n}{vx}\PYG{+w}{ }\PYG{n}{file}\PYG{+w}{ }\PYG{n}{tmp.profile}
\end{sphinxVerbatim}

\sphinxAtStartPar
Note that by default, Nswap = 1 and vtarget = INF, though this can be
changed by the optional \sphinxstyleemphasis{swap} and \sphinxstyleemphasis{vtarget} keywords.  When vtarget =
INF, one or more atoms with the most positive and negative velocity
components are selected.  Setting these parameters appropriately, in
conjunction with the swap rate N, allows the momentum flux rate to be
adjusted across a wide range of values, and the momenta to be
exchanged in large chunks or more smoothly.

\sphinxAtStartPar
The “middle” layer for momenta swapping is defined as the \sphinxstyleemphasis{Nbin}/2 + 1
layer.  Thus if \sphinxstyleemphasis{Nbin} = 20, the two swapping layers are 1 and 11.
This should lead to a symmetric velocity profile since the two layers
are separated by the same distance in both directions in a periodic
sense.  This is why \sphinxstyleemphasis{Nbin} is restricted to being an even number.

\sphinxAtStartPar
As described below, the total momentum transferred by these velocity
swaps is computed by the fix and can be output.  Dividing this
quantity by time and the cross\sphinxhyphen{}sectional area of the simulation box
yields a momentum flux.  The ratio of momentum flux to the slope of
the shear velocity profile is proportional to the viscosity of the
fluid, in appropriate units.  See the {\hyperref[\detokenize{fix_viscosity:muller-plathe2}]{\sphinxcrossref{\DUrole{std,std-ref}{Muller\sphinxhyphen{}Plathe paper}}}} for details.

\begin{sphinxadmonition}{note}{Note:}
\sphinxAtStartPar
If your system is periodic in the direction of the momentum
flux, then the flux is going in 2 directions.  This means the
effective momentum flux in one direction is reduced by a factor of 2.
You will see this in the equations for viscosity in the Muller\sphinxhyphen{}Plathe
paper.  LAMMPS is simply tallying momentum which does not account for
whether or not your system is periodic; you must use the value
appropriately to yield a viscosity for your system.
\end{sphinxadmonition}

\begin{sphinxadmonition}{note}{Note:}
\sphinxAtStartPar
After equilibration, if the velocity profile you observe is not
linear, then you are likely swapping momentum too frequently and are
not in a regime of linear response.  In this case you cannot
accurately infer a viscosity and should try increasing the Nevery
parameter.
\end{sphinxadmonition}

\sphinxAtStartPar
An alternative method for calculating a viscosity is to run a NEMD
simulation, as described on the {\hyperref[\detokenize{Howto_nemd::doc}]{\sphinxcrossref{\DUrole{doc}{Howto nemd}}}} doc page.
NEMD simulations deform the simulation box via the {\hyperref[\detokenize{fix_deform::doc}]{\sphinxcrossref{\DUrole{doc}{fix deform}}}} command.

\sphinxAtStartPar
Some features or combination of settings in LAMMPS do not support
non\sphinxhyphen{}orthogonal boxes.  Using fix viscosity keeps the box orthogonal;
thus it does not suffer from these limitations.


\subsection{Restart, fix\_modify, output, run start/stop, minimize info}
\label{\detokenize{fix_viscosity:restart-fix-modify-output-run-start-stop-minimize-info}}
\sphinxAtStartPar
No information about this fix is written to {\hyperref[\detokenize{restart::doc}]{\sphinxcrossref{\DUrole{doc}{binary restart files}}}}.  None of the {\hyperref[\detokenize{fix_modify::doc}]{\sphinxcrossref{\DUrole{doc}{fix\_modify}}}} options
are relevant to this fix.

\sphinxAtStartPar
This fix computes a global scalar which can be accessed by various
{\hyperref[\detokenize{Howto_output::doc}]{\sphinxcrossref{\DUrole{doc}{output commands}}}}.  The scalar is the cumulative
momentum transferred between the bottom and middle of the simulation
box (in the \sphinxstyleemphasis{pdim} direction) is stored as a scalar quantity by this
fix.  This quantity is zeroed when the fix is defined and accumulates
thereafter, once every N steps.  The units of the quantity are
momentum = mass*velocity.  The scalar value calculated by this fix is
“intensive”.

\sphinxAtStartPar
No parameter of this fix can be used with the \sphinxstyleemphasis{start/stop} keywords of
the {\hyperref[\detokenize{run::doc}]{\sphinxcrossref{\DUrole{doc}{run}}}} command.  This fix is not invoked during {\hyperref[\detokenize{minimize::doc}]{\sphinxcrossref{\DUrole{doc}{energy minimization}}}}.


\subsection{Restrictions}
\label{\detokenize{fix_viscosity:restrictions}}
\sphinxAtStartPar
This fix is part of the EXTRA\sphinxhyphen{}FIX package.  It is only enabled if LAMMPS
was built with that package.  See the {\hyperref[\detokenize{Build_package::doc}]{\sphinxcrossref{\DUrole{doc}{Build package}}}} page for more info.

\sphinxAtStartPar
Swaps conserve both momentum and kinetic energy, even if the masses of
the swapped atoms are not equal.  Thus you should not need to
thermostat the system.  If you do use a thermostat, you may want to
apply it only to the non\sphinxhyphen{}swapped dimensions (other than \sphinxstyleemphasis{vdim}).

\sphinxAtStartPar
LAMMPS does not check, but you should not use this fix to swap
velocities of atoms that are in constrained molecules, e.g. via {\hyperref[\detokenize{fix_shake::doc}]{\sphinxcrossref{\DUrole{doc}{fix shake}}}} or {\hyperref[\detokenize{fix_rigid::doc}]{\sphinxcrossref{\DUrole{doc}{fix rigid}}}}.  This is because
application of the constraints will alter the amount of transferred
momentum.  You should, however, be able to use flexible molecules.
See the {\hyperref[\detokenize{fix_viscosity:maginn}]{\sphinxcrossref{\DUrole{std,std-ref}{Maginn paper}}}} for an example of using this algorithm
in a computation of alcohol molecule properties.

\sphinxAtStartPar
When running a simulation with large, massive particles or molecules
in a background solvent, you may want to only exchange momenta between
solvent particles.


\subsection{Related commands}
\label{\detokenize{fix_viscosity:related-commands}}
\sphinxAtStartPar
{\hyperref[\detokenize{fix_ave_chunk::doc}]{\sphinxcrossref{\DUrole{doc}{fix ave/chunk}}}}, {\hyperref[\detokenize{fix_thermal_conductivity::doc}]{\sphinxcrossref{\DUrole{doc}{fix thermal/conductivity}}}}


\subsection{Default}
\label{\detokenize{fix_viscosity:default}}
\sphinxAtStartPar
The option defaults are swap = 1 and vtarget = INF.


\bigskip\hrule\bigskip

\phantomsection\label{\detokenize{fix_viscosity:muller-plathe2}}
\sphinxAtStartPar
\sphinxstylestrong{(Muller\sphinxhyphen{}Plathe)} Muller\sphinxhyphen{}Plathe, Phys Rev E, 59, 4894\sphinxhyphen{}4898 (1999).

\phantomsection\label{\detokenize{fix_viscosity:maginn}}
\sphinxAtStartPar
\sphinxstylestrong{(Maginn)} Kelkar, Rafferty, Maginn, Siepmann, Fluid Phase Equilibria,
260, 218\sphinxhyphen{}231 (2007).

\sphinxstepscope

\index{fix viscous@\spxentry{fix viscous}}\ignorespaces 
\index{fix viscous/kk@\spxentry{fix viscous/kk}}\ignorespaces 

\section{fix viscous command}
\label{\detokenize{fix_viscous:fix-viscous-command}}\label{\detokenize{fix_viscous:index-1}}\label{\detokenize{fix_viscous:index-0}}\label{\detokenize{fix_viscous::doc}}
\sphinxAtStartPar
Accelerator Variants: \sphinxstyleemphasis{viscous/kk}


\subsection{Syntax}
\label{\detokenize{fix_viscous:syntax}}
\begin{sphinxVerbatim}[commandchars=\\\{\}]
\PYG{k}{fix }\PYG{n+nv+nvIdentifier}{ID}\PYG{+w}{ }\PYG{n+nv+nvIdentifier}{group\PYGZhy{}ID}\PYG{+w}{ }\PYG{n}{viscous}\PYG{+w}{ }\PYG{n}{gamma}\PYG{+w}{ }\PYG{n}{keyword}\PYG{+w}{ }\PYG{n}{values}\PYG{+w}{ }\PYG{n}{...}
\end{sphinxVerbatim}
\begin{itemize}
\item {} 
\sphinxAtStartPar
ID, group\sphinxhyphen{}ID are documented in {\hyperref[\detokenize{fix::doc}]{\sphinxcrossref{\DUrole{doc}{fix}}}} command

\item {} 
\sphinxAtStartPar
viscous = style name of this fix command

\item {} 
\sphinxAtStartPar
gamma = damping coefficient (force/velocity units)

\item {} 
\sphinxAtStartPar
zero or more keyword/value pairs may be appended
\begin{sphinxalltt}
keyword = \sphinxstyleemphasis{scale}
  \sphinxstyleemphasis{scale} values = type ratio
    type = atom type (1\sphinxhyphen{}N)
    ratio = factor to scale the damping coefficient by
\end{sphinxalltt}

\end{itemize}


\subsection{Examples}
\label{\detokenize{fix_viscous:examples}}
\begin{sphinxVerbatim}[commandchars=\\\{\}]
\PYG{k}{fix }\PYG{n+nv+nvIdentifier}{1}\PYG{+w}{ }\PYG{n+nv+nvIdentifier}{flow}\PYG{+w}{ }\PYG{n}{viscous}\PYG{+w}{ }\PYG{l+m}{0.1}
\PYG{k}{fix }\PYG{n+nv+nvIdentifier}{1}\PYG{+w}{ }\PYG{n+nv+nvIdentifier}{damp}\PYG{+w}{ }\PYG{n}{viscous}\PYG{+w}{ }\PYG{l+m}{0.5}\PYG{+w}{ }\PYG{n}{scale}\PYG{+w}{ }\PYG{l+m}{3}\PYG{+w}{ }\PYG{l+m}{2.5}
\end{sphinxVerbatim}


\subsection{Description}
\label{\detokenize{fix_viscous:description}}
\sphinxAtStartPar
Add a viscous damping force to atoms in the group that is proportional
to the velocity of the atom.  The added force can be thought of as a
frictional interaction with implicit solvent, i.e. the no\sphinxhyphen{}slip Stokes
drag on a spherical particle.  In granular simulations this can be
useful for draining the kinetic energy from the system in a controlled
fashion.  If used without additional thermostatting (to add kinetic
energy to the system), it has the effect of slowly (or rapidly)
freezing the system; hence it can also be used as a simple energy
minimization technique.

\sphinxAtStartPar
The damping force \(F_i\) is given by \(F_i = - \gamma v_i\).
The larger the coefficient, the faster the kinetic energy is reduced.
If the optional keyword \sphinxstyleemphasis{scale} is used, \(\gamma\) can scaled up or
down by the specified factor for atoms of that type.  It can be used
multiple times to adjust \(\gamma\) for several atom types.

\begin{sphinxadmonition}{note}{Note:}
\sphinxAtStartPar
You should specify gamma in force/velocity units.  This is not
the same as mass/time units, at least for some of the LAMMPS
{\hyperref[\detokenize{units::doc}]{\sphinxcrossref{\DUrole{doc}{units}}}} options like “real” or “metal” that are not
self\sphinxhyphen{}consistent.
\end{sphinxadmonition}

\sphinxAtStartPar
In a Brownian dynamics context, \(\gamma = \frac{k_B T}{D}\), where
\(k_B =\) Boltzmann’s constant, \(T\) = temperature, and \sphinxstyleemphasis{D} =
particle diffusion coefficient.  \sphinxstyleemphasis{D} can be written as \(\frac{k_B
T}{3 \pi \eta d}\), where \(\eta =\) dynamic viscosity of the
frictional fluid and d = diameter of particle.  This means \(\gamma
= 3 \pi \eta d\), and thus is proportional to the viscosity of the fluid
and the particle diameter.

\sphinxAtStartPar
In the current implementation, rather than have the user specify a
viscosity, \(\gamma\) is specified directly in force/velocity units.
If needed, \(\gamma\) can be adjusted for atoms of different sizes
(i.e. \(\sigma\)) by using the \sphinxstyleemphasis{scale} keyword.

\sphinxAtStartPar
Note that Brownian dynamics models also typically include a randomized
force term to thermostat the system at a chosen temperature.  The
{\hyperref[\detokenize{fix_langevin::doc}]{\sphinxcrossref{\DUrole{doc}{fix langevin}}}} command does this.  It has the same
viscous damping term as fix viscous and adds a random force to each
atom.  The random force term is proportional to the square root of the
chosen thermostatting temperature.  Thus if you use fix langevin with a
target \(T = 0\), its random force term is zero, and you are
essentially performing the same operation as fix viscous.  Also note
that the gamma of fix viscous is related to the damping parameter of
{\hyperref[\detokenize{fix_langevin::doc}]{\sphinxcrossref{\DUrole{doc}{fix langevin}}}}, however the former is specified in
units of force/velocity and the latter in units of time, so that it can
more easily be used as a thermostat.


\bigskip\hrule\bigskip


\sphinxAtStartPar
Styles with a \sphinxstyleemphasis{gpu}, \sphinxstyleemphasis{intel}, \sphinxstyleemphasis{kk}, \sphinxstyleemphasis{omp}, or \sphinxstyleemphasis{opt} suffix are
functionally the same as the corresponding style without the suffix.
They have been optimized to run faster, depending on your available
hardware, as discussed on the {\hyperref[\detokenize{Speed_packages::doc}]{\sphinxcrossref{\DUrole{doc}{Accelerator packages}}}}
page.  The accelerated styles take the same arguments and should
produce the same results, except for round\sphinxhyphen{}off and precision issues.

\sphinxAtStartPar
These accelerated styles are part of the GPU, INTEL, KOKKOS,
OPENMP, and OPT packages, respectively.  They are only enabled if
LAMMPS was built with those packages.  See the {\hyperref[\detokenize{Build_package::doc}]{\sphinxcrossref{\DUrole{doc}{Build package}}}} page for more info.

\sphinxAtStartPar
You can specify the accelerated styles explicitly in your input script
by including their suffix, or you can use the {\hyperref[\detokenize{Run_options::doc}]{\sphinxcrossref{\DUrole{doc}{\sphinxhyphen{}suffix command\sphinxhyphen{}line switch}}}} when you invoke LAMMPS, or you can use the
{\hyperref[\detokenize{suffix::doc}]{\sphinxcrossref{\DUrole{doc}{suffix}}}} command in your input script.

\sphinxAtStartPar
See the {\hyperref[\detokenize{Speed_packages::doc}]{\sphinxcrossref{\DUrole{doc}{Accelerator packages}}}} page for more
instructions on how to use the accelerated styles effectively.


\bigskip\hrule\bigskip



\subsection{Restart, fix\_modify, output, run start/stop, minimize info}
\label{\detokenize{fix_viscous:restart-fix-modify-output-run-start-stop-minimize-info}}
\sphinxAtStartPar
No information about this fix is written to {\hyperref[\detokenize{restart::doc}]{\sphinxcrossref{\DUrole{doc}{binary restart files}}}}.  None of the {\hyperref[\detokenize{fix_modify::doc}]{\sphinxcrossref{\DUrole{doc}{fix\_modify}}}} options are
relevant to this fix.  No global or per\sphinxhyphen{}atom quantities are stored by
this fix for access by various {\hyperref[\detokenize{Howto_output::doc}]{\sphinxcrossref{\DUrole{doc}{output commands}}}}.
No parameter of this fix can be used with the \sphinxstyleemphasis{start/stop} keywords of
the {\hyperref[\detokenize{run::doc}]{\sphinxcrossref{\DUrole{doc}{run}}}} command.

\sphinxAtStartPar
The {\hyperref[\detokenize{fix_modify::doc}]{\sphinxcrossref{\DUrole{doc}{fix\_modify}}}} \sphinxstyleemphasis{respa} option is supported by this
fix. This allows to set at which level of the {\hyperref[\detokenize{run_style::doc}]{\sphinxcrossref{\DUrole{doc}{r\sphinxhyphen{}RESPA}}}}
integrator the fix is modifying forces. Default is the outermost level.

\sphinxAtStartPar
The forces due to this fix are imposed during an energy minimization,
invoked by the {\hyperref[\detokenize{minimize::doc}]{\sphinxcrossref{\DUrole{doc}{minimize}}}} command.  This fix should only
be used with damped dynamics minimizers that allow for
non\sphinxhyphen{}conservative forces.  See the {\hyperref[\detokenize{min_style::doc}]{\sphinxcrossref{\DUrole{doc}{min\_style}}}} command
for details.


\subsection{Restrictions}
\label{\detokenize{fix_viscous:restrictions}}
\sphinxAtStartPar
none


\subsection{Related commands}
\label{\detokenize{fix_viscous:related-commands}}
\sphinxAtStartPar
{\hyperref[\detokenize{fix_langevin::doc}]{\sphinxcrossref{\DUrole{doc}{fix langevin}}}},
{\hyperref[\detokenize{fix_viscous_sphere::doc}]{\sphinxcrossref{\DUrole{doc}{fix viscous/sphere}}}},
{\hyperref[\detokenize{fix_damping_cundall::doc}]{\sphinxcrossref{\DUrole{doc}{fix damping/cundall}}}}


\subsection{Default}
\label{\detokenize{fix_viscous:default}}
\sphinxAtStartPar
none

\sphinxstepscope

\index{fix viscous/sphere@\spxentry{fix viscous/sphere}}\ignorespaces 

\section{fix viscous/sphere command}
\label{\detokenize{fix_viscous_sphere:fix-viscous-sphere-command}}\label{\detokenize{fix_viscous_sphere:index-0}}\label{\detokenize{fix_viscous_sphere::doc}}

\subsection{Syntax}
\label{\detokenize{fix_viscous_sphere:syntax}}
\begin{sphinxVerbatim}[commandchars=\\\{\}]
\PYG{k}{fix }\PYG{n+nv+nvIdentifier}{ID}\PYG{+w}{ }\PYG{n+nv+nvIdentifier}{group\PYGZhy{}ID}\PYG{+w}{ }\PYG{n}{viscous}\PYG{o}{/}\PYG{n}{sphere}\PYG{+w}{ }\PYG{n}{gamma}\PYG{+w}{ }\PYG{n}{keyword}\PYG{+w}{ }\PYG{n}{values}\PYG{+w}{ }\PYG{n}{...}
\end{sphinxVerbatim}
\begin{itemize}
\item {} 
\sphinxAtStartPar
ID, group\sphinxhyphen{}ID are documented in {\hyperref[\detokenize{fix::doc}]{\sphinxcrossref{\DUrole{doc}{fix}}}} command

\item {} 
\sphinxAtStartPar
viscous/sphere = style name of this fix command

\item {} 
\sphinxAtStartPar
gamma = damping coefficient (torque/angular velocity units)

\item {} 
\sphinxAtStartPar
zero or more keyword/value pairs may be appended
\begin{sphinxalltt}
keyword = \sphinxstyleemphasis{scale}
  \sphinxstyleemphasis{scale} values = \sphinxstyleemphasis{type ratio} or \sphinxstyleemphasis{v\_name}
    type = atom type (1\sphinxhyphen{}N)
    ratio = factor to scale the damping coefficients by
    v\_name = reference to atom style variable \sphinxstyleemphasis{name}
\end{sphinxalltt}

\end{itemize}


\subsection{Examples}
\label{\detokenize{fix_viscous_sphere:examples}}
\begin{sphinxVerbatim}[commandchars=\\\{\}]
\PYG{k}{fix }\PYG{n+nv+nvIdentifier}{1}\PYG{+w}{ }\PYG{n+nv+nvIdentifier}{flow}\PYG{+w}{ }\PYG{n}{viscous}\PYG{o}{/}\PYG{n}{sphere}\PYG{+w}{ }\PYG{l+m}{0.1}
\PYG{k}{fix }\PYG{n+nv+nvIdentifier}{1}\PYG{+w}{ }\PYG{n+nv+nvIdentifier}{damp}\PYG{+w}{ }\PYG{n}{viscous}\PYG{o}{/}\PYG{n}{sphere}\PYG{+w}{ }\PYG{l+m}{0.5}\PYG{+w}{ }\PYG{n}{scale}\PYG{+w}{ }\PYG{l+m}{3}\PYG{+w}{ }\PYG{l+m}{2.5}
\PYG{k}{fix }\PYG{n+nv+nvIdentifier}{1}\PYG{+w}{ }\PYG{n+nv+nvIdentifier}{damp}\PYG{+w}{ }\PYG{n}{viscous}\PYG{o}{/}\PYG{n}{sphere}\PYG{+w}{ }\PYG{l+m}{0.5}\PYG{+w}{ }\PYG{n}{scale}\PYG{+w}{ }\PYG{n}{v\PYGZus{}radscale}
\end{sphinxVerbatim}


\subsection{Description}
\label{\detokenize{fix_viscous_sphere:description}}
\sphinxAtStartPar
Add a viscous damping torque to finite\sphinxhyphen{}size spherical particles in the group
that is proportional to the angular velocity of the atom.  In granular
simulations this can be useful for draining the rotational kinetic energy from
the system in a controlled fashion.  If used without additional thermostatting
(to add kinetic energy to the system), it has the effect of slowly (or rapidly)
freezing the system; hence it can also be used as a simple energy minimization
technique.

\sphinxAtStartPar
The damping torque \(T_i\) is given by \(T_i = - \gamma \omega_i\).
The larger the coefficient, the faster the rotational kinetic energy is reduced.

\sphinxAtStartPar
If the optional keyword \sphinxstyleemphasis{scale} is used, \(\gamma\) can be scaled up
or down by the specified factor for atoms.  This factor can be set for
different atom types and thus the \sphinxstyleemphasis{scale} keyword used multiple times
followed by the atom type and the associated scale factor.  Alternately
the scaling factor can be computed for each atom (e.g. based on its
radius) by using an {\hyperref[\detokenize{variable::doc}]{\sphinxcrossref{\DUrole{doc}{atom\sphinxhyphen{}style variable}}}}.

\begin{sphinxadmonition}{note}{Note:}
\sphinxAtStartPar
You should specify gamma in torque/angular velocity units.  This is not
the same as mass/time units, at least for some of the LAMMPS
{\hyperref[\detokenize{units::doc}]{\sphinxcrossref{\DUrole{doc}{units}}}} options like “real” or “metal” that are not
self\sphinxhyphen{}consistent.
\end{sphinxadmonition}

\sphinxAtStartPar
In the current implementation, rather than have the user specify a viscosity,
\(\gamma\) is specified directly in torque/angular velocity units.
If needed, \(\gamma\) can be adjusted for atoms of different sizes
(i.e. \(\sigma\)) by using the \sphinxstyleemphasis{scale} keyword.


\bigskip\hrule\bigskip



\subsection{Restart, fix\_modify, output, run start/stop, minimize info}
\label{\detokenize{fix_viscous_sphere:restart-fix-modify-output-run-start-stop-minimize-info}}
\sphinxAtStartPar
No information about this fix is written to {\hyperref[\detokenize{restart::doc}]{\sphinxcrossref{\DUrole{doc}{binary restart files}}}}.  None of the {\hyperref[\detokenize{fix_modify::doc}]{\sphinxcrossref{\DUrole{doc}{fix\_modify}}}} options are
relevant to this fix.  No global or per\sphinxhyphen{}atom quantities are stored by
this fix for access by various {\hyperref[\detokenize{Howto_output::doc}]{\sphinxcrossref{\DUrole{doc}{output commands}}}}.
No parameter of this fix can be used with the \sphinxstyleemphasis{start/stop} keywords of
the {\hyperref[\detokenize{run::doc}]{\sphinxcrossref{\DUrole{doc}{run}}}} command.

\sphinxAtStartPar
The {\hyperref[\detokenize{fix_modify::doc}]{\sphinxcrossref{\DUrole{doc}{fix\_modify}}}} \sphinxstyleemphasis{respa} option is supported by this
fix. This allows to set at which level of the {\hyperref[\detokenize{run_style::doc}]{\sphinxcrossref{\DUrole{doc}{r\sphinxhyphen{}RESPA}}}}
integrator the fix is modifying torques. Default is the outermost level.

\sphinxAtStartPar
The torques due to this fix are imposed during an energy minimization,
invoked by the {\hyperref[\detokenize{minimize::doc}]{\sphinxcrossref{\DUrole{doc}{minimize}}}} command.  This fix should only
be used with damped dynamics minimizers that allow for
non\sphinxhyphen{}conservative forces.  See the {\hyperref[\detokenize{min_style::doc}]{\sphinxcrossref{\DUrole{doc}{min\_style}}}} command
for details.


\subsection{Restrictions}
\label{\detokenize{fix_viscous_sphere:restrictions}}
\sphinxAtStartPar
This fix is part of the EXTRA\sphinxhyphen{}FIX package.  It is only enabled if
LAMMPS was built with that package.  See the {\hyperref[\detokenize{Build_package::doc}]{\sphinxcrossref{\DUrole{doc}{Build package}}}} page for more info.

\sphinxAtStartPar
This fix requires that atoms store torque and angular velocity (omega)
and a radius as defined by the {\hyperref[\detokenize{atom_style::doc}]{\sphinxcrossref{\DUrole{doc}{atom\_style sphere}}}}
command.

\sphinxAtStartPar
All particles in the group must be finite\sphinxhyphen{}size spheres.  They cannot
be point particles.


\subsection{Related commands}
\label{\detokenize{fix_viscous_sphere:related-commands}}
\sphinxAtStartPar
{\hyperref[\detokenize{fix_viscous::doc}]{\sphinxcrossref{\DUrole{doc}{fix viscous}}}}, {\hyperref[\detokenize{fix_damping_cundall::doc}]{\sphinxcrossref{\DUrole{doc}{fix damping/cundall}}}}


\subsection{Default}
\label{\detokenize{fix_viscous_sphere:default}}
\sphinxAtStartPar
none

\sphinxstepscope

\index{fix wall/lj93@\spxentry{fix wall/lj93}}\ignorespaces 
\index{fix wall/lj93/kk@\spxentry{fix wall/lj93/kk}}\ignorespaces 
\index{fix wall/lj126@\spxentry{fix wall/lj126}}\ignorespaces 
\index{fix wall/lj1043@\spxentry{fix wall/lj1043}}\ignorespaces 
\index{fix wall/colloid@\spxentry{fix wall/colloid}}\ignorespaces 
\index{fix wall/harmonic@\spxentry{fix wall/harmonic}}\ignorespaces 
\index{fix wall/lepton@\spxentry{fix wall/lepton}}\ignorespaces 
\index{fix wall/morse@\spxentry{fix wall/morse}}\ignorespaces 
\index{fix wall/table@\spxentry{fix wall/table}}\ignorespaces 

\section{fix wall/lj93 command}
\label{\detokenize{fix_wall:fix-wall-lj93-command}}\label{\detokenize{fix_wall:index-8}}\label{\detokenize{fix_wall:index-7}}\label{\detokenize{fix_wall:index-6}}\label{\detokenize{fix_wall:index-5}}\label{\detokenize{fix_wall:index-4}}\label{\detokenize{fix_wall:index-3}}\label{\detokenize{fix_wall:index-2}}\label{\detokenize{fix_wall:index-1}}\label{\detokenize{fix_wall:index-0}}\label{\detokenize{fix_wall::doc}}
\sphinxAtStartPar
Accelerator Variants: \sphinxstyleemphasis{wall/lj93/kk}


\section{fix wall/lj126 command}
\label{\detokenize{fix_wall:fix-wall-lj126-command}}

\section{fix wall/lj1043 command}
\label{\detokenize{fix_wall:fix-wall-lj1043-command}}

\section{fix wall/colloid command}
\label{\detokenize{fix_wall:fix-wall-colloid-command}}

\section{fix wall/harmonic command}
\label{\detokenize{fix_wall:fix-wall-harmonic-command}}

\section{fix wall/lepton command}
\label{\detokenize{fix_wall:fix-wall-lepton-command}}

\section{fix wall/morse command}
\label{\detokenize{fix_wall:fix-wall-morse-command}}

\section{fix wall/table command}
\label{\detokenize{fix_wall:fix-wall-table-command}}

\subsection{Syntax}
\label{\detokenize{fix_wall:syntax}}
\begin{sphinxVerbatim}[commandchars=\\\{\}]
\PYG{k}{fix }\PYG{n+nv+nvIdentifier}{ID}\PYG{+w}{ }\PYG{n+nv+nvIdentifier}{group\PYGZhy{}ID}\PYG{+w}{ }\PYG{n}{style}\PYG{+w}{ }\PYG{n}{[tabstyle]}\PYG{+w}{ }\PYG{n}{[N]}\PYG{+w}{ }\PYG{n}{face}\PYG{+w}{ }\PYG{n}{args}\PYG{+w}{ }\PYG{n}{...}\PYG{+w}{ }\PYG{n}{keyword}\PYG{+w}{ }\PYG{n}{value}\PYG{+w}{ }\PYG{n}{...}
\end{sphinxVerbatim}
\begin{itemize}
\item {} 
\sphinxAtStartPar
ID, group\sphinxhyphen{}ID are documented in {\hyperref[\detokenize{fix::doc}]{\sphinxcrossref{\DUrole{doc}{fix}}}} command

\item {} 
\sphinxAtStartPar
style = \sphinxstyleemphasis{wall/lj93} or \sphinxstyleemphasis{wall/lj126} or \sphinxstyleemphasis{wall/lj1043} or \sphinxstyleemphasis{wall/colloid} or \sphinxstyleemphasis{wall/harmonic} or \sphinxstyleemphasis{wall/lepton} or \sphinxstyleemphasis{wall/morse} or \sphinxstyleemphasis{wall/table}

\item {} 
\sphinxAtStartPar
tabstyle = \sphinxstyleemphasis{linear} or \sphinxstyleemphasis{spline} = method of table interpolation (only applies to \sphinxstyleemphasis{wall/table})

\item {} 
\sphinxAtStartPar
N = use N values in \sphinxstyleemphasis{linear} or \sphinxstyleemphasis{spline} interpolation (only applies to \sphinxstyleemphasis{wall/table})

\item {} 
\sphinxAtStartPar
one or more face/arg pairs may be appended

\item {} 
\sphinxAtStartPar
face = \sphinxstyleemphasis{xlo} or \sphinxstyleemphasis{xhi} or \sphinxstyleemphasis{ylo} or \sphinxstyleemphasis{yhi} or \sphinxstyleemphasis{zlo} or \sphinxstyleemphasis{zhi}

\end{itemize}
\begin{itemize}
\item {} 
\sphinxAtStartPar
args for styles \sphinxstyleemphasis{lj93} or \sphinxstyleemphasis{lj126} or \sphinxstyleemphasis{lj1043} or \sphinxstyleemphasis{colloid} or \sphinxstyleemphasis{harmonic}
\begin{sphinxalltt}
args = coord epsilon sigma cutoff
coord = position of wall = EDGE or constant or variable
  EDGE = current lo or hi edge of simulation box
  constant = number like 0.0 or \sphinxhyphen{}30.0 (distance units)
  variable = {\hyperref[\detokenize{variable::doc}]{\sphinxcrossref{\DUrole{doc}{equal\sphinxhyphen{}style variable}}}} like v\_x or v\_wiggle
epsilon = strength factor for wall\sphinxhyphen{}particle interaction (energy or energy/distance\textasciicircum{}2 units)
  epsilon can be a variable (see below)
sigma = size factor for wall\sphinxhyphen{}particle interaction (distance units)
  sigma can be a variable (see below)
cutoff = distance from wall at which wall\sphinxhyphen{}particle interactions are cut off (distance units)
\end{sphinxalltt}

\item {} 
\sphinxAtStartPar
args for style \sphinxstyleemphasis{lepton}
\begin{sphinxalltt}
args = coord expression cutoff
coord = position of wall = EDGE or constant or variable
  EDGE = current lo or hi edge of simulation box
  constant = number like 0.0 or \sphinxhyphen{}30.0 (distance units)
  variable = {\hyperref[\detokenize{variable::doc}]{\sphinxcrossref{\DUrole{doc}{equal\sphinxhyphen{}style variable}}}} like v\_x or v\_wiggle
expression = Lepton expression for the potential  (energy units)
cutoff = distance from wall at which wall\sphinxhyphen{}particle interactions are cut off (distance units)
\end{sphinxalltt}

\item {} 
\sphinxAtStartPar
args for style \sphinxstyleemphasis{morse}
\begin{sphinxalltt}
args = coord D\_0 alpha r\_0 cutoff
coord = position of wall = EDGE or constant or variable
  EDGE = current lo or hi edge of simulation box
  constant = number like 0.0 or \sphinxhyphen{}30.0 (distance units)
  variable = {\hyperref[\detokenize{variable::doc}]{\sphinxcrossref{\DUrole{doc}{equal\sphinxhyphen{}style variable}}}} like v\_x or v\_wiggle
D\_0 = depth of the potential (energy units)
  D\_0 can be a variable (see below)
alpha = width factor for wall\sphinxhyphen{}particle interaction (1/distance units)
  alpha can be a variable (see below)
r\_0 = distance of the potential minimum from the face of region (distance units)
  r\_0 can be a variable (see below)
cutoff = distance from wall at which wall\sphinxhyphen{}particle interactions are cut off (distance units)
\end{sphinxalltt}

\item {} 
\sphinxAtStartPar
args for style \sphinxstyleemphasis{table}
\begin{sphinxalltt}
args = coord filename keyword cutoff
coord = position of wall = EDGE or constant or variable
  EDGE = current lo or hi edge of simulation box
  constant = number like 0.0 or \sphinxhyphen{}30.0 (distance units)
  variable = {\hyperref[\detokenize{variable::doc}]{\sphinxcrossref{\DUrole{doc}{equal\sphinxhyphen{}style variable}}}} like v\_x or v\_wiggle
filename = file containing tabulated energy and force values
keyword = section identifier to select a specific table in table file
cutoff = distance from wall at which wall\sphinxhyphen{}particle interactions are cut off (distance units)
\end{sphinxalltt}

\item {} 
\sphinxAtStartPar
zero or more keyword/value pairs may be appended

\item {} 
\sphinxAtStartPar
keyword = \sphinxstyleemphasis{units} or \sphinxstyleemphasis{fld} or \sphinxstyleemphasis{pbc}
\begin{sphinxalltt}
\sphinxstyleemphasis{units} value = \sphinxstyleemphasis{lattice} or \sphinxstyleemphasis{box}
  \sphinxstyleemphasis{lattice} = the wall position is defined in lattice units
  \sphinxstyleemphasis{box} = the wall position is defined in simulation box units
\sphinxstyleemphasis{fld} value = \sphinxstyleemphasis{yes} or \sphinxstyleemphasis{no}
  \sphinxstyleemphasis{yes} = invoke the wall constraint to be compatible with implicit FLD
  \sphinxstyleemphasis{no} = invoke the wall constraint in the normal way
\sphinxstyleemphasis{pbc} value = \sphinxstyleemphasis{yes} or \sphinxstyleemphasis{no}
  \sphinxstyleemphasis{yes} = allow periodic boundary in a wall dimension
  \sphinxstyleemphasis{no} = require non\sphinxhyphen{}perioidic boundaries in any wall dimension
\end{sphinxalltt}

\end{itemize}


\subsection{Examples}
\label{\detokenize{fix_wall:examples}}
\begin{sphinxVerbatim}[commandchars=\\\{\}]
\PYG{k}{fix }\PYG{n+nv+nvIdentifier}{wallhi}\PYG{+w}{ }\PYG{n+nv+nvIdentifier}{all}\PYG{+w}{ }\PYG{n}{wall}\PYG{o}{/}\PYG{n}{lj93}\PYG{+w}{ }\PYG{n}{xlo}\PYG{+w}{ }\PYG{o}{\PYGZhy{}}\PYG{l+m}{1.0}\PYG{+w}{ }\PYG{l+m}{1.0}\PYG{+w}{ }\PYG{l+m}{1.0}\PYG{+w}{ }\PYG{l+m}{2.5}\PYG{+w}{ }\PYG{n}{units}\PYG{+w}{ }\PYG{n}{box}
\PYG{k}{fix }\PYG{n+nv+nvIdentifier}{wallhi}\PYG{+w}{ }\PYG{n+nv+nvIdentifier}{all}\PYG{+w}{ }\PYG{n}{wall}\PYG{o}{/}\PYG{n}{lj93}\PYG{+w}{ }\PYG{n}{xhi}\PYG{+w}{ }\PYG{n}{EDGE}\PYG{+w}{ }\PYG{l+m}{1.0}\PYG{+w}{ }\PYG{l+m}{1.0}\PYG{+w}{ }\PYG{l+m}{2.5}
\PYG{k}{fix }\PYG{n+nv+nvIdentifier}{wallhi}\PYG{+w}{ }\PYG{n+nv+nvIdentifier}{all}\PYG{+w}{ }\PYG{n}{wall}\PYG{o}{/}\PYG{n}{harmonic}\PYG{+w}{ }\PYG{n}{xhi}\PYG{+w}{ }\PYG{n}{EDGE}\PYG{+w}{ }\PYG{l+m}{100.0}\PYG{+w}{ }\PYG{l+m}{0.0}\PYG{+w}{ }\PYG{l+m}{4.0}\PYG{+w}{ }\PYG{n}{units}\PYG{+w}{ }\PYG{n}{box}
\PYG{k}{fix }\PYG{n+nv+nvIdentifier}{wallhi}\PYG{+w}{ }\PYG{n+nv+nvIdentifier}{all}\PYG{+w}{ }\PYG{n}{wall}\PYG{o}{/}\PYG{n}{morse}\PYG{+w}{ }\PYG{n}{xhi}\PYG{+w}{ }\PYG{n}{EDGE}\PYG{+w}{ }\PYG{l+m}{1.0}\PYG{+w}{ }\PYG{l+m}{1.0}\PYG{+w}{ }\PYG{l+m}{1.0}\PYG{+w}{ }\PYG{l+m}{2.5}\PYG{+w}{ }\PYG{n}{units}\PYG{+w}{ }\PYG{n}{box}
\PYG{k}{fix }\PYG{n+nv+nvIdentifier}{wallhi}\PYG{+w}{ }\PYG{n+nv+nvIdentifier}{all}\PYG{+w}{ }\PYG{n}{wall}\PYG{o}{/}\PYG{n}{lj126}\PYG{+w}{ }\PYG{n}{v\PYGZus{}wiggle}\PYG{+w}{ }\PYG{l+m}{23.2}\PYG{+w}{ }\PYG{l+m}{1.0}\PYG{+w}{ }\PYG{l+m}{1.0}\PYG{+w}{ }\PYG{l+m}{2.5}
\PYG{k}{fix }\PYG{n+nv+nvIdentifier}{zwalls}\PYG{+w}{ }\PYG{n+nv+nvIdentifier}{all}\PYG{+w}{ }\PYG{n}{wall}\PYG{o}{/}\PYG{n}{colloid}\PYG{+w}{ }\PYG{n}{zlo}\PYG{+w}{ }\PYG{l+m}{0.0}\PYG{+w}{ }\PYG{l+m}{1.0}\PYG{+w}{ }\PYG{l+m}{1.0}\PYG{+w}{ }\PYG{l+m}{0.858}\PYG{+w}{ }\PYG{n}{zhi}\PYG{+w}{ }\PYG{l+m}{40.0}\PYG{+w}{ }\PYG{l+m}{1.0}\PYG{+w}{ }\PYG{l+m}{1.0}\PYG{+w}{ }\PYG{l+m}{0.858}
\PYG{k}{fix }\PYG{n+nv+nvIdentifier}{xwall}\PYG{+w}{ }\PYG{n+nv+nvIdentifier}{mobile}\PYG{+w}{ }\PYG{n}{wall}\PYG{o}{/}\PYG{n}{table}\PYG{+w}{ }\PYG{n}{spline}\PYG{+w}{ }\PYG{l+m}{200}\PYG{+w}{ }\PYG{n}{EDGE}\PYG{+w}{ }\PYG{o}{\PYGZhy{}}\PYG{l+m}{5.0}\PYG{+w}{ }\PYG{n}{walltab.dat}\PYG{+w}{ }\PYG{n}{HARMONIC}\PYG{+w}{ }\PYG{l+m}{4.0}
\PYG{k}{fix }\PYG{n+nv+nvIdentifier}{xwalls}\PYG{+w}{ }\PYG{n+nv+nvIdentifier}{mobile}\PYG{+w}{ }\PYG{n}{wall}\PYG{o}{/}\PYG{n}{lepton}\PYG{+w}{ }\PYG{n}{xlo}\PYG{+w}{ }\PYG{o}{\PYGZhy{}}\PYG{l+m}{5.0}\PYG{+w}{ }\PYG{l+s}{\PYGZdq{}}\PYG{l+s}{k*(r\PYGZhy{}rc)\PYGZca{}2;k=100.0}\PYG{l+s}{\PYGZdq{}}\PYG{+w}{ }\PYG{l+m}{4.0}\PYG{+w}{ }\PYG{n}{xhi}\PYG{+w}{ }\PYG{l+m}{5.0}\PYG{+w}{ }\PYG{l+s}{\PYGZdq{}}\PYG{l+s}{k*(r\PYGZhy{}rc)\PYGZca{}2;k=100.0}\PYG{l+s}{\PYGZdq{}}\PYG{+w}{ }\PYG{l+m}{4.0}
\end{sphinxVerbatim}


\subsection{Description}
\label{\detokenize{fix_wall:description}}
\sphinxAtStartPar
Bound the simulation domain on one or more of its faces with a flat
wall that interacts with the atoms in the group by generating a force
on the atom in a direction perpendicular to the wall.  The energy of
wall\sphinxhyphen{}particle interactions depends on the style.

\sphinxAtStartPar
For style \sphinxstyleemphasis{wall/lj93}, the energy E is given by the 9\sphinxhyphen{}3 Lennard\sphinxhyphen{}Jones potential:
\begin{equation*}
\begin{split}E = \epsilon \left[ \frac{2}{15} \left(\frac{\sigma}{r}\right)^{9} -
                      \left(\frac{\sigma}{r}\right)^3 \right]
                      \qquad r < r_c\end{split}
\end{equation*}
\sphinxAtStartPar
For style \sphinxstyleemphasis{wall/lj126}, the energy E is given by the 12\sphinxhyphen{}6 Lennard\sphinxhyphen{}Jones potential:
\begin{equation*}
\begin{split}E = 4 \epsilon \left[ \left(\frac{\sigma}{r}\right)^{12} -
                      \left(\frac{\sigma}{r}\right)^6 \right]
                      \qquad r < r_c\end{split}
\end{equation*}
\sphinxAtStartPar
For style \sphinxstyleemphasis{wall/lj1043}, the energy E is given by the 10\sphinxhyphen{}4\sphinxhyphen{}3 Lennard\sphinxhyphen{}Jones potential:
\begin{equation*}
\begin{split}E = 2 \pi \epsilon \left[ \frac{2}{5} \left(\frac{\sigma}{r}\right)^{10} -
                      \left(\frac{\sigma}{r}\right)^4 -
                      \frac{\sqrt(2)\sigma^3}{3\left(r+\left(0.61/\sqrt(2)\right)\sigma\right)^3}\right]
                      \qquad r < r_c\end{split}
\end{equation*}
\sphinxAtStartPar
For style \sphinxstyleemphasis{wall/colloid}, the energy E is given by an integrated form
of the {\hyperref[\detokenize{pair_colloid::doc}]{\sphinxcrossref{\DUrole{doc}{pair\_style colloid}}}} potential:
\begin{equation*}
\begin{split}E = & \epsilon \left[ \frac{\sigma^{6}}{7560}
\left(\frac{6R-D}{D^{7}} + \frac{D+8R}{(D+2R)^{7}} \right) \right. \\
 & \left. - \frac{1}{6} \left(\frac{2R(D+R) + D(D+2R)
 \left[ \ln D - \ln (D+2R) \right]}{D(D+2R)} \right) \right] \qquad r < r_c\end{split}
\end{equation*}
\sphinxAtStartPar
For style \sphinxstyleemphasis{wall/harmonic}, the energy E is given by a repulsive\sphinxhyphen{}only harmonic
spring potential:
\begin{equation*}
\begin{split}E = \epsilon \quad (r - r_c)^2 \qquad r < r_c\end{split}
\end{equation*}
\sphinxAtStartPar
For style \sphinxstyleemphasis{wall/morse}, the energy E is given by a Morse potential:
\begin{equation*}
\begin{split}E = D_0 \left[ e^{- 2 \alpha (r - r_0)} - 2 e^{- \alpha (r - r_0)} \right]
    \qquad r < r_c\end{split}
\end{equation*}
\sphinxAtStartPar
\DUrole{versionmodified,added}{Added in version 28Mar2023.}

\sphinxAtStartPar
For style \sphinxstyleemphasis{wall/lepton}, the energy E is provided as an Lepton
expression string using “r” as the distance variable.  The \sphinxhref{https://simtk.org/projects/lepton}{Lepton
library}, that the \sphinxstyleemphasis{wall/lepton}
style interfaces with, evaluates this expression string at run time to
compute the wall\sphinxhyphen{}particle energy.  It also creates an analytical
representation of the first derivative of this expression with respect
to “r” and then uses that to compute the force between the wall and
atoms in the fix group.  The Lepton expression must be either enclosed
in quotes or must not contain any whitespace so that LAMMPS recognizes
it as a single keyword.

\sphinxAtStartPar
Optionally, the expression may use “rc” to refer to the cutoff distance
for the given wall.  Further constants in the expression can be defined
in the same string as additional expressions separated by semicolons.
The expression “k*(r\sphinxhyphen{}rc)\textasciicircum{}2;k=100.0” represents a repulsive\sphinxhyphen{}only harmonic
spring as in fix \sphinxstyleemphasis{wall/harmonic} with a force constant \sphinxstyleemphasis{K} (same as
\(\epsilon\) above) of 100 energy units.  More details on the Lepton
expression strings are given below.

\sphinxAtStartPar
\DUrole{versionmodified,added}{Added in version 28Mar2023.}

\sphinxAtStartPar
For style \sphinxstyleemphasis{wall/table}, the energy E and forces are determined from
interpolation tables listed in one or more files as a function of
distance.  The interpolation tables are used to evaluate energy and
forces between particles and the wall similar to how analytic formulas
are used for the other wall styles.

\sphinxAtStartPar
The interpolation tables are created as a pre\sphinxhyphen{}computation by fitting
cubic splines to the file values and interpolating energy and force
values at each of \sphinxstyleemphasis{N} distances.  During a simulation, the tables are
used to interpolate energy and force values as needed for each wall and
particle separated by a distance \sphinxstyleemphasis{R}.  The interpolation is done in
one of two styles: \sphinxstyleemphasis{linear} or \sphinxstyleemphasis{spline}.

\sphinxAtStartPar
For the \sphinxstyleemphasis{linear} style, the distance \sphinxstyleemphasis{R} is used to find the 2
surrounding table values from which an energy or force is computed by
linear interpolation.

\sphinxAtStartPar
For the \sphinxstyleemphasis{spline} style, cubic spline coefficients are computed and
stored for each of the \sphinxstyleemphasis{N} values in the table, one set of splines for
energy, another for force.  Note that these splines are different than
the ones used to pre\sphinxhyphen{}compute the \sphinxstyleemphasis{N} values.  Those splines were fit
to the \sphinxstyleemphasis{Nfile} values in the tabulated file, where often \sphinxstyleemphasis{Nfile} \textless{}
\sphinxstyleemphasis{N}.  The distance \sphinxstyleemphasis{R} is used to find the appropriate set of spline
coefficients which are used to evaluate a cubic polynomial which
computes the energy or force.

\sphinxAtStartPar
For each wall a filename and a keyword must be provided as in the
examples above.  The filename specifies a file containing tabulated
energy and force values.  The keyword specifies a section of the file.
The format of this file is described below.

\sphinxAtStartPar
In all cases, \sphinxstyleemphasis{r} is the distance from the particle to the wall at
position \sphinxstyleemphasis{coord}, and \(r_c\) is the \sphinxstyleemphasis{cutoff} distance at which the
particle and wall no longer interact.  The energy of the wall
potential is shifted so that the wall\sphinxhyphen{}particle interaction energy is
0.0 at the cutoff distance.

\sphinxAtStartPar
Up to 6 walls or faces can be specified in a single command: \sphinxstyleemphasis{xlo},
\sphinxstyleemphasis{xhi}, \sphinxstyleemphasis{ylo}, \sphinxstyleemphasis{yhi}, \sphinxstyleemphasis{zlo}, \sphinxstyleemphasis{zhi}.  A \sphinxstyleemphasis{lo} face interacts with
particles near the lower side of the simulation box in that dimension.
A \sphinxstyleemphasis{hi} face interacts with particles near the upper side of the
simulation box in that dimension.

\sphinxAtStartPar
The position of each wall can be specified in one of 3 ways: as the
EDGE of the simulation box, as a constant value, or as a variable.  If
EDGE is used, then the corresponding boundary of the current
simulation box is used.  If a numeric constant is specified then the
wall is placed at that position in the appropriate dimension (x, y, or
z).  In both the EDGE and constant cases, the wall will never move.
If the wall position is a variable, it should be specified as v\_name,
where name is an {\hyperref[\detokenize{variable::doc}]{\sphinxcrossref{\DUrole{doc}{equal\sphinxhyphen{}style variable}}}} name.  In this
case the variable is evaluated each timestep and the result becomes
the current position of the reflecting wall.  Equal\sphinxhyphen{}style variables
can specify formulas with various mathematical functions, and include
{\hyperref[\detokenize{thermo_style::doc}]{\sphinxcrossref{\DUrole{doc}{thermo\_style}}}} command keywords for the simulation
box parameters and timestep and elapsed time.  Thus it is easy to
specify a time\sphinxhyphen{}dependent wall position.  See examples below.

\sphinxAtStartPar
For the \sphinxstyleemphasis{wall/lj93} and \sphinxstyleemphasis{wall/lj126} and \sphinxstyleemphasis{wall/lj1043} styles,
\(\epsilon\) and \(\sigma\) are the usual Lennard\sphinxhyphen{}Jones
parameters, which determine the strength and size of the particle as it
interacts with the wall.  Epsilon has energy units.  Note that this
\(\epsilon\) and \(\sigma\) may be different than any
\(\epsilon\) or \(\sigma\) values defined for a pair style that
computes particle\sphinxhyphen{}particle interactions.

\sphinxAtStartPar
The \sphinxstyleemphasis{wall/lj93} interaction is derived by integrating over a 3d
half\sphinxhyphen{}lattice of Lennard\sphinxhyphen{}Jones 12/6 particles.  The \sphinxstyleemphasis{wall/lj126}
interaction is effectively a harder, more repulsive wall interaction.
The \sphinxstyleemphasis{wall/lj1043} interaction is yet a different form of wall
interaction, described in Magda et al in {\hyperref[\detokenize{fix_wall:magda}]{\sphinxcrossref{\DUrole{std,std-ref}{(Magda)}}}}.

\sphinxAtStartPar
For the \sphinxstyleemphasis{wall/colloid} style, \sphinxstyleemphasis{R} is the radius of the colloid particle,
\sphinxstyleemphasis{D} is the distance from the surface of the colloid particle to the wall
(r\sphinxhyphen{}R), and \(\sigma\) is the size of a constituent LJ particle
inside the colloid particle and wall.  Note that the cutoff distance Rc
in this case is the distance from the colloid particle center to the
wall.  The prefactor \(\epsilon\) can be thought of as an effective
Hamaker constant with energy units for the strength of the colloid\sphinxhyphen{}wall
interaction.  More specifically, the \(\epsilon\) prefactor is
\(4\pi^2 \rho_{wall} \rho_{colloid} \epsilon \sigma^6\), where
\(\epsilon\) and \(\sigma\) are the LJ parameters for the
constituent LJ particles. \(\rho_{wall}\) and \(\rho_{colloid}\)
are the number density of the constituent particles, in the wall and
colloid respectively, in units of 1/volume.

\sphinxAtStartPar
The \sphinxstyleemphasis{wall/colloid} interaction is derived by integrating over
constituent LJ particles of size \(\sigma\) within the colloid
particle and a 3d half\sphinxhyphen{}lattice of Lennard\sphinxhyphen{}Jones 12/6 particles of size
\(\sigma\) in the wall.  As mentioned in the preceding paragraph,
the density of particles in the wall and colloid can be different, as
specified by the \(\epsilon\) prefactor.

\sphinxAtStartPar
For the \sphinxstyleemphasis{wall/harmonic} style, \(\epsilon\) is effectively the spring
constant K, and has units (energy/distance\textasciicircum{}2).  The input parameter
\(\sigma\) is ignored.  The minimum energy position of the harmonic
spring is at the \sphinxstyleemphasis{cutoff}.  This is a repulsive\sphinxhyphen{}only spring since the
interaction is truncated at the \sphinxstyleemphasis{cutoff}

\sphinxAtStartPar
For the \sphinxstyleemphasis{wall/morse} style, the three parameters are in this order:
\(D_0\) the depth of the potential, \(\alpha\) the width
parameter, and \(r_0\) the location of the minimum.  \(D_0\) has
energy units, \(\alpha\) inverse distance units, and \(r_0\)
distance units.

\sphinxAtStartPar
For any wall that supports them, the \(\epsilon\) and/or
\(\sigma\) and/or \(\alpha\) parameter can be specified as an
{\hyperref[\detokenize{variable::doc}]{\sphinxcrossref{\DUrole{doc}{equal\sphinxhyphen{}style variable}}}}, in which case it should be
specified as v\_name, where name is the variable name.  As with a
variable wall position, the variable is evaluated each timestep and the
result becomes the current epsilon or sigma of the wall.  Equal\sphinxhyphen{}style
variables can specify formulas with various mathematical functions, and
include {\hyperref[\detokenize{thermo_style::doc}]{\sphinxcrossref{\DUrole{doc}{thermo\_style}}}} command keywords for the
simulation box parameters and timestep and elapsed time.  Thus it is
easy to specify a time\sphinxhyphen{}dependent wall interaction.

\begin{sphinxadmonition}{note}{Note:}
\sphinxAtStartPar
For all of the styles, you must ensure that r is always \textgreater{} 0 for
all particles in the group, or LAMMPS will generate an error.  This
means you cannot start your simulation with particles at the wall
position \sphinxstyleemphasis{coord} (r = 0) or with particles on the wrong side of the
wall (r \textless{} 0).  For the \sphinxstyleemphasis{wall/lj93} and \sphinxstyleemphasis{wall/lj126} styles, the energy
of the wall/particle interaction (and hence the force on the particle)
blows up as r \sphinxhyphen{}\textgreater{} 0.  The \sphinxstyleemphasis{wall/colloid} style is even more
restrictive, since the energy blows up as D = r\sphinxhyphen{}R \sphinxhyphen{}\textgreater{} 0.  This means
the finite\sphinxhyphen{}size particles of radius R must be a distance larger than R
from the wall position \sphinxstyleemphasis{coord}.  The \sphinxstyleemphasis{harmonic} style is a softer
potential and does not blow up as r \sphinxhyphen{}\textgreater{} 0, but you must use a large
enough \(\epsilon\) that particles always reamin on the correct side of
the wall (r \textgreater{} 0).
\end{sphinxadmonition}

\sphinxAtStartPar
The \sphinxstyleemphasis{units} keyword determines the meaning of the distance units used
to define a wall position, but only when a numeric constant or
variable is used.  It is not relevant when EDGE is used to specify a
face position.  In the variable case, the variable is assumed to
produce a value compatible with the \sphinxstyleemphasis{units} setting you specify.

\sphinxAtStartPar
A \sphinxstyleemphasis{box} value selects standard distance units as defined by the
{\hyperref[\detokenize{units::doc}]{\sphinxcrossref{\DUrole{doc}{units}}}} command, e.g. Angstroms for units = real or metal.
A \sphinxstyleemphasis{lattice} value means the distance units are in lattice spacings.
The {\hyperref[\detokenize{lattice::doc}]{\sphinxcrossref{\DUrole{doc}{lattice}}}} command must have been previously used to
define the lattice spacings.

\sphinxAtStartPar
The \sphinxstyleemphasis{fld} keyword can be used with a \sphinxstyleemphasis{yes} setting to invoke the wall
constraint before pairwise interactions are computed.  This allows an
implicit FLD model using {\hyperref[\detokenize{pair_lubricateU::doc}]{\sphinxcrossref{\DUrole{doc}{pair\_style lubricateU}}}}
to include the wall force in its calculations.  If the setting is \sphinxstyleemphasis{no},
wall forces are imposed after pairwise interactions, in the usual
manner.

\sphinxAtStartPar
The \sphinxstyleemphasis{pbc} keyword can be used with a \sphinxstyleemphasis{yes} setting to allow walls to be
specified in a periodic dimension.  See the {\hyperref[\detokenize{boundary::doc}]{\sphinxcrossref{\DUrole{doc}{boundary}}}}
command for options on simulation box boundaries.  The default for \sphinxstyleemphasis{pbc}
is \sphinxstyleemphasis{no}, which means the system must be non\sphinxhyphen{}periodic when using a wall.
But you may wish to use a periodic box.  E.g. to allow some particles to
interact with the wall via the fix group\sphinxhyphen{}ID, and others to pass through
it and wrap around a periodic box.  In this case you should ensure that
the wall is sufficiently far enough away from the box boundary.  If you
do not, then particles may interact with both the wall and with periodic
images on the other side of the box, which is probably not what you
want.


\bigskip\hrule\bigskip


\sphinxAtStartPar
Here are examples of variable definitions that move the wall position
in a time\sphinxhyphen{}dependent fashion using equal\sphinxhyphen{}style
{\hyperref[\detokenize{variable::doc}]{\sphinxcrossref{\DUrole{doc}{variables}}}}.  The wall interaction parameters (epsilon,
sigma) could be varied with additional variable definitions.

\begin{sphinxVerbatim}[commandchars=\\\{\}]
\PYG{k}{variable }\PYG{n+nv+nvIdentifier}{ramp}\PYG{+w}{ }\PYG{n}{equal}\PYG{+w}{ }\PYG{n}{ramp}\PYG{n+nv}{(}\PYG{n+nv}{0,10}\PYG{n+nv}{)}
\PYG{k}{fix }\PYG{n+nv+nvIdentifier}{1}\PYG{+w}{ }\PYG{n+nv+nvIdentifier}{all}\PYG{+w}{ }\PYG{n}{wall}\PYG{+w}{ }\PYG{n}{xlo}\PYG{+w}{ }\PYG{n}{v\PYGZus{}ramp}\PYG{+w}{ }\PYG{l+m}{1.0}\PYG{+w}{ }\PYG{l+m}{1.0}\PYG{+w}{ }\PYG{l+m}{2.5}

\PYG{k}{variable }\PYG{n+nv+nvIdentifier}{linear}\PYG{+w}{ }\PYG{n}{equal}\PYG{+w}{ }\PYG{n}{vdisplace}\PYG{n+nv}{(}\PYG{n+nv}{0,20}\PYG{n+nv}{)}
\PYG{k}{fix }\PYG{n+nv+nvIdentifier}{1}\PYG{+w}{ }\PYG{n+nv+nvIdentifier}{all}\PYG{+w}{ }\PYG{n}{wall}\PYG{+w}{ }\PYG{n}{xlo}\PYG{+w}{ }\PYG{n}{v\PYGZus{}linear}\PYG{+w}{ }\PYG{l+m}{1.0}\PYG{+w}{ }\PYG{l+m}{1.0}\PYG{+w}{ }\PYG{l+m}{2.5}

\PYG{k}{variable }\PYG{n+nv+nvIdentifier}{wiggle}\PYG{+w}{ }\PYG{n}{equal}\PYG{+w}{ }\PYG{n}{swiggle}\PYG{n+nv}{(}\PYG{n+nv}{0.0,5.0,3.0}\PYG{n+nv}{)}
\PYG{k}{fix }\PYG{n+nv+nvIdentifier}{1}\PYG{+w}{ }\PYG{n+nv+nvIdentifier}{all}\PYG{+w}{ }\PYG{n}{wall}\PYG{+w}{ }\PYG{n}{xlo}\PYG{+w}{ }\PYG{n}{v\PYGZus{}wiggle}\PYG{+w}{ }\PYG{l+m}{1.0}\PYG{+w}{ }\PYG{l+m}{1.0}\PYG{+w}{ }\PYG{l+m}{2.5}

\PYG{k}{variable }\PYG{n+nv+nvIdentifier}{wiggle}\PYG{+w}{ }\PYG{n}{equal}\PYG{+w}{ }\PYG{n}{cwiggle}\PYG{n+nv}{(}\PYG{n+nv}{0.0,5.0,3.0}\PYG{n+nv}{)}
\PYG{k}{fix }\PYG{n+nv+nvIdentifier}{1}\PYG{+w}{ }\PYG{n+nv+nvIdentifier}{all}\PYG{+w}{ }\PYG{n}{wall}\PYG{+w}{ }\PYG{n}{xlo}\PYG{+w}{ }\PYG{n}{v\PYGZus{}wiggle}\PYG{+w}{ }\PYG{l+m}{1.0}\PYG{+w}{ }\PYG{l+m}{1.0}\PYG{+w}{ }\PYG{l+m}{2.5}
\end{sphinxVerbatim}

\sphinxAtStartPar
The \sphinxstyleemphasis{ramp(lo,hi)} function adjusts the wall position linearly from \sphinxstyleemphasis{lo} to
\sphinxstyleemphasis{hi} over the course of a run.  The \sphinxstyleemphasis{vdisplace(c0,velocity)} function does
something similar using the equation \sphinxstyleemphasis{position = c0 + velocity*delta},
where \sphinxstyleemphasis{delta} is the elapsed time.

\sphinxAtStartPar
The \sphinxstyleemphasis{swiggle(c0,A,period)} function causes the wall position to
oscillate sinusoidally according to this equation, where \sphinxstyleemphasis{omega = 2 PI
/ period}:
\begin{sphinxalltt}
position = c0 + A sin(omega*delta)
\end{sphinxalltt}

\sphinxAtStartPar
The \sphinxstyleemphasis{cwiggle(c0,A,period)} function causes the wall position to
oscillate sinusoidally according to this equation, which will have an
initial wall velocity of 0.0, and thus may impose a gentler
perturbation on the particles:
\begin{sphinxalltt}
position = c0 + A (1 \sphinxhyphen{} cos(omega*delta))
\end{sphinxalltt}


\bigskip\hrule\bigskip



\subsection{Lepton expression syntax and features}
\label{\detokenize{fix_wall:lepton-expression-syntax-and-features}}
\sphinxAtStartPar
Lepton supports the following operators in expressions:


\begin{savenotes}\sphinxattablestart
\sphinxthistablewithglobalstyle
\centering
\begin{tabulary}{\linewidth}[t]{TTTTTTTTTTTTTT}
\sphinxtoprule
\sphinxtableatstartofbodyhook
\sphinxAtStartPar
+
&
\sphinxAtStartPar
Add
&&
\sphinxAtStartPar
\sphinxhyphen{}
&
\sphinxAtStartPar
Subtract
&&
\sphinxAtStartPar
*
&
\sphinxAtStartPar
Multiply
&&
\sphinxAtStartPar
/
&
\sphinxAtStartPar
Divide
&&
\sphinxAtStartPar
\textasciicircum{}
&
\sphinxAtStartPar
Power
\\
\sphinxbottomrule
\end{tabulary}
\sphinxtableafterendhook\par
\sphinxattableend\end{savenotes}

\sphinxAtStartPar
The following mathematical functions are available:


\begin{savenotes}\sphinxattablestart
\sphinxthistablewithglobalstyle
\centering
\begin{tabulary}{\linewidth}[t]{TTTT}
\sphinxtoprule
\sphinxtableatstartofbodyhook
\sphinxAtStartPar
sqrt(x)
&
\sphinxAtStartPar
Square root
&
\sphinxAtStartPar
exp(x)
&
\sphinxAtStartPar
Exponential
\\
\sphinxhline
\sphinxAtStartPar
log(x)
&
\sphinxAtStartPar
Natural logarithm
&
\sphinxAtStartPar
sin(x)
&
\sphinxAtStartPar
Sine (angle in radians)
\\
\sphinxhline
\sphinxAtStartPar
cos(x)
&
\sphinxAtStartPar
Cosine (angle in radians)
&
\sphinxAtStartPar
sec(x)
&
\sphinxAtStartPar
Secant (angle in radians)
\\
\sphinxhline
\sphinxAtStartPar
csc(x)
&
\sphinxAtStartPar
Cosecant (angle in radians)
&
\sphinxAtStartPar
tan(x)
&
\sphinxAtStartPar
Tangent (angle in radians)
\\
\sphinxhline
\sphinxAtStartPar
cot(x)
&
\sphinxAtStartPar
Cotangent (angle in radians)
&
\sphinxAtStartPar
asin(x)
&
\sphinxAtStartPar
Inverse sine (in radians)
\\
\sphinxhline
\sphinxAtStartPar
acos(x)
&
\sphinxAtStartPar
Inverse cosine (in radians)
&
\sphinxAtStartPar
atan(x)
&
\sphinxAtStartPar
Inverse tangent (in radians)
\\
\sphinxhline
\sphinxAtStartPar
sinh(x)
&
\sphinxAtStartPar
Hyperbolic sine
&
\sphinxAtStartPar
cosh(x)
&
\sphinxAtStartPar
Hyperbolic cosine
\\
\sphinxhline
\sphinxAtStartPar
tanh(x)
&
\sphinxAtStartPar
Hyperbolic tangent
&
\sphinxAtStartPar
erf(x)
&
\sphinxAtStartPar
Error function
\\
\sphinxhline
\sphinxAtStartPar
erfc(x)
&
\sphinxAtStartPar
Complementary Error function
&
\sphinxAtStartPar
abs(x)
&
\sphinxAtStartPar
Absolute value
\\
\sphinxhline
\sphinxAtStartPar
min(x,y)
&
\sphinxAtStartPar
Minimum of two values
&
\sphinxAtStartPar
max(x,y)
&
\sphinxAtStartPar
Maximum of two values
\\
\sphinxhline
\sphinxAtStartPar
delta(x)
&
\sphinxAtStartPar
delta(x) is 1 for \sphinxtitleref{x = 0}, otherwise 0
&
\sphinxAtStartPar
step(x)
&
\sphinxAtStartPar
step(x) is 0 for \sphinxtitleref{x \textless{} 0}, otherwise 1
\\
\sphinxbottomrule
\end{tabulary}
\sphinxtableafterendhook\par
\sphinxattableend\end{savenotes}

\sphinxAtStartPar
Numbers may be given in either decimal or exponential form.  All of the
following are valid numbers: \sphinxtitleref{5}, \sphinxtitleref{\sphinxhyphen{}3.1}, \sphinxtitleref{1e6}, and \sphinxtitleref{3.12e\sphinxhyphen{}2}.

\sphinxAtStartPar
As an extension to the standard Lepton syntax, it is also possible to
use LAMMPS {\hyperref[\detokenize{variable::doc}]{\sphinxcrossref{\DUrole{doc}{variables}}}} in the format “v\_name”.  Before
evaluating the expression, “v\_name” will be replaced with the value of
the variable “name”.  This is compatible with all kinds of scalar
variables, but not with vectors, arrays, local, or per\sphinxhyphen{}atom
variables.  If necessary, a custom scalar variable needs to be defined
that can access the desired (single) item from a non\sphinxhyphen{}scalar variable.
As an example, the following lines will instruct LAMMPS to ramp
the force constant for a harmonic bond from 100.0 to 200.0 during the
next run:

\begin{sphinxVerbatim}[commandchars=\\\{\}]
\PYG{k}{variable }\PYG{n+nv+nvIdentifier}{fconst}\PYG{+w}{ }\PYG{n}{equal}\PYG{+w}{ }\PYG{n}{ramp}\PYG{n+nv}{(}\PYG{n+nv}{100.0, 200}\PYG{n+nv}{)}
\PYG{k}{bond\PYGZus{}style}\PYG{+w}{ }\PYG{n}{lepton}
\PYG{k}{bond\PYGZus{}coeff}\PYG{+w}{ }\PYG{l+m}{1}\PYG{+w}{ }\PYG{l+m}{1.5}\PYG{+w}{ }\PYG{l+s}{\PYGZdq{}}\PYG{l+s}{v\PYGZus{}fconst * (r\PYGZca{}2)}\PYG{l+s}{\PYGZdq{}}
\end{sphinxVerbatim}

\sphinxAtStartPar
An expression may be followed by definitions for intermediate values that appear in the
expression. A semicolon “;” is used as a delimiter between value definitions. For example,
the expression:

\begin{sphinxVerbatim}[commandchars=\\\{\}]
\PYG{n}{a}\PYG{o}{\PYGZca{}}\PYG{l+m+mi}{2}\PYG{o}{+}\PYG{n}{a}\PYG{o}{*}\PYG{n}{b}\PYG{o}{+}\PYG{n}{b}\PYG{o}{\PYGZca{}}\PYG{l+m+mi}{2}\PYG{p}{;}\PYG{+w}{ }\PYG{n}{a}\PYG{o}{=}\PYG{n}{a1}\PYG{o}{+}\PYG{n}{a2}\PYG{p}{;}\PYG{+w}{ }\PYG{n}{b}\PYG{o}{=}\PYG{n}{b1}\PYG{o}{+}\PYG{n}{b2}
\end{sphinxVerbatim}

\sphinxAtStartPar
is exactly equivalent to

\begin{sphinxVerbatim}[commandchars=\\\{\}]
\PYG{p}{(}\PYG{n}{a1}\PYG{o}{+}\PYG{n}{a2}\PYG{p}{)}\PYG{o}{\PYGZca{}}\PYG{l+m+mi}{2}\PYG{o}{+}\PYG{p}{(}\PYG{n}{a1}\PYG{o}{+}\PYG{n}{a2}\PYG{p}{)}\PYG{o}{*}\PYG{p}{(}\PYG{n}{b1}\PYG{o}{+}\PYG{n}{b2}\PYG{p}{)}\PYG{o}{+}\PYG{p}{(}\PYG{n}{b1}\PYG{o}{+}\PYG{n}{b2}\PYG{p}{)}\PYG{o}{\PYGZca{}}\PYG{l+m+mi}{2}
\end{sphinxVerbatim}

\sphinxAtStartPar
The definition of an intermediate value may itself involve other
intermediate values. Whitespace and quotation characters (’\textquotesingle{}’ and ‘”’)
are ignored.  All uses of a value must appear \sphinxstyleemphasis{before} that value’s
definition.  For efficiency reasons, the expression string is parsed,
optimized, and then stored in an internal, pre\sphinxhyphen{}parsed representation for
evaluation.

\sphinxAtStartPar
Evaluating a Lepton expression is typically between 2.5 and 5 times
slower than the corresponding compiled and optimized C++ code.  If
additional speed or GPU acceleration (via GPU or KOKKOS) is required,
the interaction can be represented as a table.  Suitable table files
can be created either internally using the {\hyperref[\detokenize{pair_write::doc}]{\sphinxcrossref{\DUrole{doc}{pair\_write}}}}
or {\hyperref[\detokenize{bond_write::doc}]{\sphinxcrossref{\DUrole{doc}{bond\_write}}}} command or through the Python scripts
in the {\hyperref[\detokenize{Tools:tabulate}]{\sphinxcrossref{\DUrole{std,std-ref}{tools/tabulate}}}} folder.


\bigskip\hrule\bigskip



\subsection{Table file format}
\label{\detokenize{fix_wall:table-file-format}}
\sphinxAtStartPar
Suitable tables for use with fix \sphinxstyleemphasis{wall/table} can be created by the
Python code in the \sphinxcode{\sphinxupquote{tools/tabulate}} folder of the LAMMPS source code
distribution.

\sphinxAtStartPar
The format of a tabulated file is as follows (without the parenthesized
comments):

\begin{sphinxVerbatim}[commandchars=\\\{\}]
\PYGZsh{} Tabulated wall potential UNITS: real

HARMONIC                       (keyword is the first text on a line)
N 100 FP 200 200
                               (blank line)
1  0.04    1568.16    792.00   (index, distance to wall, energy, force)
2  0.08    1536.64    784.00
3  0.12    1505.44    776.00
...
99  3.96      0.16      8.00
100  4.00     0         0
\end{sphinxVerbatim}

\sphinxAtStartPar
A section begins with a non\sphinxhyphen{}blank line whose first character is not a
“\#”; blank lines or lines starting with “\#” can be used as comments
between sections.  The first line begins with a keyword which identifies
the section.  The line can contain additional text, but the initial text
must match the argument specified in the fix \sphinxstyleemphasis{wall/table} command.  The
next line lists (in any order) one or more parameters for the table.
Each parameter is a keyword followed by one or more numeric values.

\sphinxAtStartPar
The parameter “N” is required and its value is the number of table
entries that follow.  Note that this may be different than the \sphinxstyleemphasis{N}
specified in the fix \sphinxstyleemphasis{wall/table} command.  Let Ntable = \sphinxstyleemphasis{N} in the fix
command, and Nfile = “N” in the tabulated file.  What LAMMPS does is a
preliminary interpolation by creating splines using the Nfile tabulated
values as nodal points.  It uses these to interpolate as needed to
generate energy and force values at Ntable different points.  The
resulting tables of length Ntable are then used as described above, when
computing energy and force for wall\sphinxhyphen{}particle interactions.  This means that
if you want the interpolation tables of length Ntable to match exactly
what is in the tabulated file (with effectively no preliminary
interpolation), you should set Ntable = Nfile.


\bigskip\hrule\bigskip



\subsection{Restart, fix\_modify, output, run start/stop, minimize info}
\label{\detokenize{fix_wall:restart-fix-modify-output-run-start-stop-minimize-info}}
\sphinxAtStartPar
No information about this fix is written to {\hyperref[\detokenize{restart::doc}]{\sphinxcrossref{\DUrole{doc}{binary restart files}}}}.

\sphinxAtStartPar
The {\hyperref[\detokenize{fix_modify::doc}]{\sphinxcrossref{\DUrole{doc}{fix\_modify}}}} \sphinxstyleemphasis{energy} option is supported by
this fix to add the energy of interaction between atoms and all the
specified walls to the global potential energy of the system as part
of {\hyperref[\detokenize{thermo_style::doc}]{\sphinxcrossref{\DUrole{doc}{thermodynamic output}}}}.  The default setting
for this fix is {\hyperref[\detokenize{fix_modify::doc}]{\sphinxcrossref{\DUrole{doc}{fix\_modify energy no}}}}.

\sphinxAtStartPar
The {\hyperref[\detokenize{fix_modify::doc}]{\sphinxcrossref{\DUrole{doc}{fix\_modify}}}} \sphinxstyleemphasis{virial} option is supported by
this fix to add the contribution due to the interaction between atoms
and all the specified walls to both the global pressure and per\sphinxhyphen{}atom
stress of the system via the {\hyperref[\detokenize{compute_pressure::doc}]{\sphinxcrossref{\DUrole{doc}{compute pressure}}}} and {\hyperref[\detokenize{compute_stress_atom::doc}]{\sphinxcrossref{\DUrole{doc}{compute stress/atom}}}} commands.  The former can be accessed by
{\hyperref[\detokenize{thermo_style::doc}]{\sphinxcrossref{\DUrole{doc}{thermodynamic output}}}}.  The default setting for
this fix is {\hyperref[\detokenize{fix_modify::doc}]{\sphinxcrossref{\DUrole{doc}{fix\_modify virial no}}}}.

\sphinxAtStartPar
The {\hyperref[\detokenize{fix_modify::doc}]{\sphinxcrossref{\DUrole{doc}{fix\_modify}}}} \sphinxstyleemphasis{respa} option is supported by this
fix. This allows to set at which level of the {\hyperref[\detokenize{run_style::doc}]{\sphinxcrossref{\DUrole{doc}{r\sphinxhyphen{}RESPA}}}} integrator the fix is adding its forces. Default is the
outermost level.

\sphinxAtStartPar
This fix computes a global scalar energy and a global vector of forces,
which can be accessed by various {\hyperref[\detokenize{Howto_output::doc}]{\sphinxcrossref{\DUrole{doc}{output commands}}}}.
Note that the scalar energy is the sum of interactions with all defined
walls.  If you want the energy on a per\sphinxhyphen{}wall basis, you need to use
multiple fix wall commands.  The length of the vector is equal to the
number of walls defined by the fix.  Each vector value is the normal
force on a specific wall.  Note that an outward force on a wall will be
a negative value for \sphinxstyleemphasis{lo} walls and a positive value for \sphinxstyleemphasis{hi} walls.
The scalar and vector values calculated by this fix are “extensive”.

\sphinxAtStartPar
No parameter of this fix can be used with the \sphinxstyleemphasis{start/stop} keywords of
the {\hyperref[\detokenize{run::doc}]{\sphinxcrossref{\DUrole{doc}{run}}}} command.

\sphinxAtStartPar
The forces due to this fix are imposed during an energy minimization,
invoked by the {\hyperref[\detokenize{minimize::doc}]{\sphinxcrossref{\DUrole{doc}{minimize}}}} command.

\begin{sphinxadmonition}{note}{Note:}
\sphinxAtStartPar
If you want the atom/wall interaction energy to be included in
the total potential energy of the system (the quantity being
minimized), you MUST enable the {\hyperref[\detokenize{fix_modify::doc}]{\sphinxcrossref{\DUrole{doc}{fix\_modify}}}} \sphinxstyleemphasis{energy}
option for this fix.
\end{sphinxadmonition}


\bigskip\hrule\bigskip


\sphinxAtStartPar
Styles with a \sphinxstyleemphasis{gpu}, \sphinxstyleemphasis{intel}, \sphinxstyleemphasis{kk}, \sphinxstyleemphasis{omp}, or \sphinxstyleemphasis{opt} suffix are
functionally the same as the corresponding style without the suffix.
They have been optimized to run faster, depending on your available
hardware, as discussed on the {\hyperref[\detokenize{Speed_packages::doc}]{\sphinxcrossref{\DUrole{doc}{Accelerator packages}}}}
page.  The accelerated styles take the same arguments and should
produce the same results, except for round\sphinxhyphen{}off and precision issues.

\sphinxAtStartPar
These accelerated styles are part of the GPU, INTEL, KOKKOS,
OPENMP, and OPT packages, respectively.  They are only enabled if
LAMMPS was built with those packages.  See the {\hyperref[\detokenize{Build_package::doc}]{\sphinxcrossref{\DUrole{doc}{Build package}}}} page for more info.

\sphinxAtStartPar
You can specify the accelerated styles explicitly in your input script
by including their suffix, or you can use the {\hyperref[\detokenize{Run_options::doc}]{\sphinxcrossref{\DUrole{doc}{\sphinxhyphen{}suffix command\sphinxhyphen{}line switch}}}} when you invoke LAMMPS, or you can use the
{\hyperref[\detokenize{suffix::doc}]{\sphinxcrossref{\DUrole{doc}{suffix}}}} command in your input script.

\sphinxAtStartPar
See the {\hyperref[\detokenize{Speed_packages::doc}]{\sphinxcrossref{\DUrole{doc}{Accelerator packages}}}} page for more
instructions on how to use the accelerated styles effectively.


\bigskip\hrule\bigskip



\subsection{Restrictions}
\label{\detokenize{fix_wall:restrictions}}
\sphinxAtStartPar
Fix \sphinxstyleemphasis{wall/lepton} is part of the LEPTON package and only enabled if
LAMMPS was built with this package.  See the {\hyperref[\detokenize{Build_package::doc}]{\sphinxcrossref{\DUrole{doc}{Build package}}}} page for more info.


\subsection{Related commands}
\label{\detokenize{fix_wall:related-commands}}
\sphinxAtStartPar
{\hyperref[\detokenize{fix_wall_reflect::doc}]{\sphinxcrossref{\DUrole{doc}{fix wall/reflect}}}},
{\hyperref[\detokenize{fix_wall_gran::doc}]{\sphinxcrossref{\DUrole{doc}{fix wall/gran}}}},
{\hyperref[\detokenize{fix_wall_region::doc}]{\sphinxcrossref{\DUrole{doc}{fix wall/region}}}}


\subsection{Default}
\label{\detokenize{fix_wall:default}}
\sphinxAtStartPar
The option defaults units = lattice, fld = no, and pbc = no.


\bigskip\hrule\bigskip

\phantomsection\label{\detokenize{fix_wall:magda}}
\sphinxAtStartPar
\sphinxstylestrong{(Magda)} Magda, Tirrell, Davis, J Chem Phys, 83, 1888\sphinxhyphen{}1901 (1985);
erratum in JCP 84, 2901 (1986).

\sphinxstepscope

\index{fix wall/body/polygon@\spxentry{fix wall/body/polygon}}\ignorespaces 

\section{fix wall/body/polygon command}
\label{\detokenize{fix_wall_body_polygon:fix-wall-body-polygon-command}}\label{\detokenize{fix_wall_body_polygon:index-0}}\label{\detokenize{fix_wall_body_polygon::doc}}

\subsection{Syntax}
\label{\detokenize{fix_wall_body_polygon:syntax}}
\begin{sphinxVerbatim}[commandchars=\\\{\}]
\PYG{k}{fix }\PYG{n+nv+nvIdentifier}{ID}\PYG{+w}{ }\PYG{n+nv+nvIdentifier}{group\PYGZhy{}ID}\PYG{+w}{ }\PYG{n}{wall}\PYG{o}{/}\PYG{n}{body}\PYG{o}{/}\PYG{n}{polygon}\PYG{+w}{ }\PYG{n}{k\PYGZus{}n}\PYG{+w}{ }\PYG{n}{c\PYGZus{}n}\PYG{+w}{ }\PYG{n}{c\PYGZus{}t}\PYG{+w}{ }\PYG{n}{wallstyle}\PYG{+w}{ }\PYG{n}{args}\PYG{+w}{ }\PYG{n}{keyword}\PYG{+w}{ }\PYG{n}{values}\PYG{+w}{ }\PYG{n}{...}
\end{sphinxVerbatim}
\begin{itemize}
\item {} 
\sphinxAtStartPar
ID, group\sphinxhyphen{}ID are documented in {\hyperref[\detokenize{fix::doc}]{\sphinxcrossref{\DUrole{doc}{fix}}}} command

\item {} 
\sphinxAtStartPar
wall/body/polygon = style name of this fix command

\item {} 
\sphinxAtStartPar
k\_n = normal repulsion strength (force/distance or pressure units)

\item {} 
\sphinxAtStartPar
c\_n = normal damping coefficient (force/distance or pressure units)

\item {} 
\sphinxAtStartPar
c\_t = tangential damping coefficient (force/distance or pressure units)

\item {} 
\sphinxAtStartPar
wallstyle = \sphinxstyleemphasis{xplane} or \sphinxstyleemphasis{yplane} or \sphinxstyleemphasis{zcylinder}

\item {} 
\sphinxAtStartPar
args = list of arguments for a particular style
\begin{sphinxalltt}
\sphinxstyleemphasis{xplane} or \sphinxstyleemphasis{yplane} args = lo hi
  lo,hi = position of lower and upper plane (distance units), either can be NULL)
\sphinxstyleemphasis{zcylinder} args = radius
  radius = cylinder radius (distance units)
\end{sphinxalltt}

\item {} 
\sphinxAtStartPar
zero or more keyword/value pairs may be appended to args

\item {} 
\sphinxAtStartPar
keyword = \sphinxstyleemphasis{wiggle}
\begin{sphinxalltt}
\sphinxstyleemphasis{wiggle} values = dim amplitude period
  dim = \sphinxstyleemphasis{x} or \sphinxstyleemphasis{y} or \sphinxstyleemphasis{z}
  amplitude = size of oscillation (distance units)
  period = time of oscillation (time units)
\end{sphinxalltt}

\end{itemize}


\subsection{Examples}
\label{\detokenize{fix_wall_body_polygon:examples}}
\begin{sphinxVerbatim}[commandchars=\\\{\}]
\PYG{k}{fix }\PYG{n+nv+nvIdentifier}{1}\PYG{+w}{ }\PYG{n+nv+nvIdentifier}{all}\PYG{+w}{ }\PYG{n}{wall}\PYG{o}{/}\PYG{n}{body}\PYG{o}{/}\PYG{n}{polygon}\PYG{+w}{ }\PYG{l+m}{1000.0}\PYG{+w}{ }\PYG{l+m}{20.0}\PYG{+w}{ }\PYG{l+m}{5.0}\PYG{+w}{ }\PYG{n}{xplane}\PYG{+w}{ }\PYG{o}{\PYGZhy{}}\PYG{l+m}{10.0}\PYG{+w}{ }\PYG{l+m}{10.0}
\end{sphinxVerbatim}


\subsection{Description}
\label{\detokenize{fix_wall_body_polygon:description}}
\sphinxAtStartPar
This fix is for use with 2d models of body particles of style
\sphinxstyleemphasis{rounded/polygon}.  It bounds the simulation domain with wall(s).  All
particles in the group interact with the wall when they are close
enough to touch it.  The nature of the interaction between the wall
and the polygon particles is the same as that between the polygon
particles themselves, which is similar to a Hookean potential.  See
the {\hyperref[\detokenize{Howto_body::doc}]{\sphinxcrossref{\DUrole{doc}{Howto body}}}} page for more details on using
body particles.

\sphinxAtStartPar
The parameters \sphinxstyleemphasis{k\_n}, \sphinxstyleemphasis{c\_n}, \sphinxstyleemphasis{c\_t} have the same meaning and units as
those specified with the {\hyperref[\detokenize{pair_body_rounded_polygon::doc}]{\sphinxcrossref{\DUrole{doc}{pair\_style body/rounded/polygon}}}} command.

\sphinxAtStartPar
The \sphinxstyleemphasis{wallstyle} can be planar or cylindrical.  The 2 planar options
specify a pair of walls in a dimension.  Wall positions are given by
\sphinxstyleemphasis{lo} and \sphinxstyleemphasis{hi}.  Either of the values can be specified as NULL if a
single wall is desired.  For a \sphinxstyleemphasis{zcylinder} wallstyle, the cylinder’s
axis is at x = y = 0.0, and the radius of the cylinder is specified.

\sphinxAtStartPar
Optionally, the wall can be moving, if the \sphinxstyleemphasis{wiggle} keyword is
appended.

\sphinxAtStartPar
For the \sphinxstyleemphasis{wiggle} keyword, the wall oscillates sinusoidally, similar to
the oscillations of particles which can be specified by the {\hyperref[\detokenize{fix_move::doc}]{\sphinxcrossref{\DUrole{doc}{fix move}}}} command.  This is useful in packing simulations of
particles.  The arguments to the \sphinxstyleemphasis{wiggle} keyword specify a dimension
for the motion, as well as it’s \sphinxstyleemphasis{amplitude} and \sphinxstyleemphasis{period}.  Note that
if the dimension is in the plane of the wall, this is effectively a
shearing motion.  If the dimension is perpendicular to the wall, it is
more of a shaking motion.  A \sphinxstyleemphasis{zcylinder} wall can only be wiggled in
the z dimension.

\sphinxAtStartPar
Each timestep, the position of a wiggled wall in the appropriate \sphinxstyleemphasis{dim}
is set according to this equation:
\begin{sphinxalltt}
position = coord + A \sphinxhyphen{} A cos (omega * delta)
\end{sphinxalltt}

\sphinxAtStartPar
where \sphinxstyleemphasis{coord} is the specified initial position of the wall, \sphinxstyleemphasis{A} is
the \sphinxstyleemphasis{amplitude}, \sphinxstyleemphasis{omega} is 2 PI / \sphinxstyleemphasis{period}, and \sphinxstyleemphasis{delta} is the time
elapsed since the fix was specified.  The velocity of the wall is set
to the derivative of this expression.


\subsection{Restart, fix\_modify, output, run start/stop, minimize info}
\label{\detokenize{fix_wall_body_polygon:restart-fix-modify-output-run-start-stop-minimize-info}}
\sphinxAtStartPar
None of the {\hyperref[\detokenize{fix_modify::doc}]{\sphinxcrossref{\DUrole{doc}{fix\_modify}}}} options are relevant to this
fix.  No global or per\sphinxhyphen{}atom quantities are stored by this fix for
access by various {\hyperref[\detokenize{Howto_output::doc}]{\sphinxcrossref{\DUrole{doc}{output commands}}}}.  No parameter
of this fix can be used with the \sphinxstyleemphasis{start/stop} keywords of the
{\hyperref[\detokenize{run::doc}]{\sphinxcrossref{\DUrole{doc}{run}}}} command.  This fix is not invoked during {\hyperref[\detokenize{minimize::doc}]{\sphinxcrossref{\DUrole{doc}{energy minimization}}}}.


\subsection{Restrictions}
\label{\detokenize{fix_wall_body_polygon:restrictions}}
\sphinxAtStartPar
This fix is part of the BODY package.  It is only enabled if LAMMPS
was built with that package.  See the {\hyperref[\detokenize{Build_package::doc}]{\sphinxcrossref{\DUrole{doc}{Build package}}}} page for more info.

\sphinxAtStartPar
Any dimension (xy) that has a wall must be non\sphinxhyphen{}periodic.


\subsection{Related commands}
\label{\detokenize{fix_wall_body_polygon:related-commands}}
\sphinxAtStartPar
{\hyperref[\detokenize{atom_style::doc}]{\sphinxcrossref{\DUrole{doc}{atom\_style body}}}}, {\hyperref[\detokenize{pair_body_rounded_polygon::doc}]{\sphinxcrossref{\DUrole{doc}{pair\_style body/rounded/polygon}}}}


\subsection{Default}
\label{\detokenize{fix_wall_body_polygon:default}}
\sphinxAtStartPar
none

\sphinxstepscope

\index{fix wall/body/polyhedron@\spxentry{fix wall/body/polyhedron}}\ignorespaces 

\section{fix wall/body/polyhedron command}
\label{\detokenize{fix_wall_body_polyhedron:fix-wall-body-polyhedron-command}}\label{\detokenize{fix_wall_body_polyhedron:index-0}}\label{\detokenize{fix_wall_body_polyhedron::doc}}

\subsection{Syntax}
\label{\detokenize{fix_wall_body_polyhedron:syntax}}
\begin{sphinxVerbatim}[commandchars=\\\{\}]
\PYG{k}{fix }\PYG{n+nv+nvIdentifier}{ID}\PYG{+w}{ }\PYG{n+nv+nvIdentifier}{group\PYGZhy{}ID}\PYG{+w}{ }\PYG{n}{wall}\PYG{o}{/}\PYG{n}{body}\PYG{o}{/}\PYG{n}{polyhedron}\PYG{+w}{ }\PYG{n}{k\PYGZus{}n}\PYG{+w}{ }\PYG{n}{c\PYGZus{}n}\PYG{+w}{ }\PYG{n}{c\PYGZus{}t}\PYG{+w}{ }\PYG{n}{wallstyle}\PYG{+w}{ }\PYG{n}{args}\PYG{+w}{ }\PYG{n}{keyword}\PYG{+w}{ }\PYG{n}{values}\PYG{+w}{ }\PYG{n}{...}
\end{sphinxVerbatim}
\begin{itemize}
\item {} 
\sphinxAtStartPar
ID, group\sphinxhyphen{}ID are documented in {\hyperref[\detokenize{fix::doc}]{\sphinxcrossref{\DUrole{doc}{fix}}}} command

\item {} 
\sphinxAtStartPar
wall/body/polyhedron = style name of this fix command

\item {} 
\sphinxAtStartPar
k\_n = normal repulsion strength (force/distance units or pressure units \sphinxhyphen{} see discussion below)

\item {} 
\sphinxAtStartPar
c\_n = normal damping coefficient (force/distance units or pressure units \sphinxhyphen{} see discussion below)

\item {} 
\sphinxAtStartPar
c\_t = tangential damping coefficient (force/distance units or pressure units \sphinxhyphen{} see discussion below)

\item {} 
\sphinxAtStartPar
wallstyle = \sphinxstyleemphasis{xplane} or \sphinxstyleemphasis{yplane} or \sphinxstyleemphasis{zplane}

\item {} 
\sphinxAtStartPar
args = list of arguments for a particular style
\begin{sphinxalltt}
\sphinxstyleemphasis{xplane} or \sphinxstyleemphasis{yplane} or \sphinxstyleemphasis{zplane} args = lo hi
  lo,hi = position of lower and upper plane (distance units), either can be NULL)
\end{sphinxalltt}

\item {} 
\sphinxAtStartPar
zero or more keyword/value pairs may be appended to args

\item {} 
\sphinxAtStartPar
keyword = \sphinxstyleemphasis{wiggle}
\begin{sphinxalltt}
\sphinxstyleemphasis{wiggle} values = dim amplitude period
  dim = \sphinxstyleemphasis{x} or \sphinxstyleemphasis{y} or \sphinxstyleemphasis{z}
  amplitude = size of oscillation (distance units)
  period = time of oscillation (time units)
\end{sphinxalltt}

\end{itemize}


\subsection{Examples}
\label{\detokenize{fix_wall_body_polyhedron:examples}}
\begin{sphinxVerbatim}[commandchars=\\\{\}]
\PYG{k}{fix }\PYG{n+nv+nvIdentifier}{1}\PYG{+w}{ }\PYG{n+nv+nvIdentifier}{all}\PYG{+w}{ }\PYG{n}{wall}\PYG{o}{/}\PYG{n}{body}\PYG{o}{/}\PYG{n}{polyhedron}\PYG{+w}{ }\PYG{l+m}{1000.0}\PYG{+w}{ }\PYG{l+m}{20.0}\PYG{+w}{ }\PYG{l+m}{5.0}\PYG{+w}{ }\PYG{n}{xplane}\PYG{+w}{ }\PYG{o}{\PYGZhy{}}\PYG{l+m}{10.0}\PYG{+w}{ }\PYG{l+m}{10.0}
\end{sphinxVerbatim}


\subsection{Description}
\label{\detokenize{fix_wall_body_polyhedron:description}}
\sphinxAtStartPar
This fix is for use with 3d models of body particles of style
\sphinxstyleemphasis{rounded/polyhedron}.  It bounds the simulation domain with wall(s).
All particles in the group interact with the wall when they are close
enough to touch it.  The nature of the interaction between the wall
and the polygon particles is the same as that between the polygon
particles themselves, which is similar to a Hookean potential.  See
the {\hyperref[\detokenize{Howto_body::doc}]{\sphinxcrossref{\DUrole{doc}{Howto body}}}} page for more details on using
body particles.

\sphinxAtStartPar
The parameters \sphinxstyleemphasis{k\_n}, \sphinxstyleemphasis{c\_n}, \sphinxstyleemphasis{c\_t} have the same meaning and units as
those specified with the {\hyperref[\detokenize{pair_body_rounded_polyhedron::doc}]{\sphinxcrossref{\DUrole{doc}{pair\_style body/rounded/polyhedron}}}} command.

\sphinxAtStartPar
The \sphinxstyleemphasis{wallstyle} can be planar or cylindrical.  The 3 planar options
specify a pair of walls in a dimension.  Wall positions are given by
\sphinxstyleemphasis{lo} and \sphinxstyleemphasis{hi}.  Either of the values can be specified as NULL if a
single wall is desired.

\sphinxAtStartPar
Optionally, the wall can be moving, if the \sphinxstyleemphasis{wiggle} keyword is appended.

\sphinxAtStartPar
For the \sphinxstyleemphasis{wiggle} keyword, the wall oscillates sinusoidally, similar to
the oscillations of particles which can be specified by the {\hyperref[\detokenize{fix_move::doc}]{\sphinxcrossref{\DUrole{doc}{fix move}}}} command.  This is useful in packing simulations of
particles.  The arguments to the \sphinxstyleemphasis{wiggle} keyword specify a dimension
for the motion, as well as it’s \sphinxstyleemphasis{amplitude} and \sphinxstyleemphasis{period}.  Note that
if the dimension is in the plane of the wall, this is effectively a
shearing motion.  If the dimension is perpendicular to the wall, it is
more of a shaking motion.

\sphinxAtStartPar
Each timestep, the position of a wiggled wall in the appropriate \sphinxstyleemphasis{dim}
is set according to this equation:
\begin{sphinxalltt}
position = coord + A \sphinxhyphen{} A cos (omega * delta)
\end{sphinxalltt}

\sphinxAtStartPar
where \sphinxstyleemphasis{coord} is the specified initial position of the wall, \sphinxstyleemphasis{A} is
the \sphinxstyleemphasis{amplitude}, \sphinxstyleemphasis{omega} is 2 PI / \sphinxstyleemphasis{period}, and \sphinxstyleemphasis{delta} is the time
elapsed since the fix was specified.  The velocity of the wall is set
to the derivative of this expression.


\subsection{Restart, fix\_modify, output, run start/stop, minimize info}
\label{\detokenize{fix_wall_body_polyhedron:restart-fix-modify-output-run-start-stop-minimize-info}}
\sphinxAtStartPar
None of the {\hyperref[\detokenize{fix_modify::doc}]{\sphinxcrossref{\DUrole{doc}{fix\_modify}}}} options are relevant to this
fix.  No global or per\sphinxhyphen{}atom quantities are stored by this fix for
access by various {\hyperref[\detokenize{Howto_output::doc}]{\sphinxcrossref{\DUrole{doc}{output commands}}}}.  No parameter
of this fix can be used with the \sphinxstyleemphasis{start/stop} keywords of the
{\hyperref[\detokenize{run::doc}]{\sphinxcrossref{\DUrole{doc}{run}}}} command.  This fix is not invoked during {\hyperref[\detokenize{minimize::doc}]{\sphinxcrossref{\DUrole{doc}{energy minimization}}}}.


\subsection{Restrictions}
\label{\detokenize{fix_wall_body_polyhedron:restrictions}}
\sphinxAtStartPar
This fix is part of the BODY package.  It is only enabled if LAMMPS
was built with that package.  See the {\hyperref[\detokenize{Build_package::doc}]{\sphinxcrossref{\DUrole{doc}{Build package}}}} page for more info.

\sphinxAtStartPar
Any dimension (xyz) that has a wall must be non\sphinxhyphen{}periodic.


\subsection{Related commands}
\label{\detokenize{fix_wall_body_polyhedron:related-commands}}
\sphinxAtStartPar
{\hyperref[\detokenize{atom_style::doc}]{\sphinxcrossref{\DUrole{doc}{atom\_style body}}}}, {\hyperref[\detokenize{pair_body_rounded_polyhedron::doc}]{\sphinxcrossref{\DUrole{doc}{pair\_style body/rounded/polyhedron}}}}


\subsection{Default}
\label{\detokenize{fix_wall_body_polyhedron:default}}
\sphinxAtStartPar
none

\sphinxstepscope

\index{fix wall/ees@\spxentry{fix wall/ees}}\ignorespaces 
\index{fix wall/region/ees@\spxentry{fix wall/region/ees}}\ignorespaces 

\section{fix wall/ees command}
\label{\detokenize{fix_wall_ees:fix-wall-ees-command}}\label{\detokenize{fix_wall_ees:index-1}}\label{\detokenize{fix_wall_ees:index-0}}\label{\detokenize{fix_wall_ees::doc}}

\section{fix wall/region/ees command}
\label{\detokenize{fix_wall_ees:fix-wall-region-ees-command}}

\subsection{Syntax}
\label{\detokenize{fix_wall_ees:syntax}}
\begin{sphinxVerbatim}[commandchars=\\\{\}]
\PYG{k}{fix }\PYG{n+nv+nvIdentifier}{ID}\PYG{+w}{ }\PYG{n+nv+nvIdentifier}{group\PYGZhy{}ID}\PYG{+w}{ }\PYG{n}{style}\PYG{+w}{ }\PYG{n}{args}
\end{sphinxVerbatim}
\begin{itemize}
\item {} 
\sphinxAtStartPar
ID, group\sphinxhyphen{}ID are documented in {\hyperref[\detokenize{fix::doc}]{\sphinxcrossref{\DUrole{doc}{fix}}}} command

\item {} 
\sphinxAtStartPar
style = \sphinxstyleemphasis{wall/ees} or \sphinxstyleemphasis{wall/region/ees}
\begin{sphinxalltt}
args for style \sphinxstyleemphasis{wall/ees}: one or more \sphinxstyleemphasis{face parameters} groups may be appended
face = \sphinxstyleemphasis{xlo} or \sphinxstyleemphasis{xhi} or \sphinxstyleemphasis{ylo} or \sphinxstyleemphasis{yhi} or \sphinxstyleemphasis{zlo} or \sphinxstyleemphasis{zhi}
parameters = coord epsilon sigma cutoff
  coord = position of wall = EDGE or constant or variable
    EDGE = current lo or hi edge of simulation box
    constant = number like 0.0 or \sphinxhyphen{}30.0 (distance units)
    variable = {\hyperref[\detokenize{variable::doc}]{\sphinxcrossref{\DUrole{doc}{equal\sphinxhyphen{}style variable}}}} like v\_x or v\_wiggle
  epsilon = strength factor for wall\sphinxhyphen{}particle interaction (energy or energy/distance\textasciicircum{}2 units)
    epsilon can be a variable (see below)
  sigma = size factor for wall\sphinxhyphen{}particle interaction (distance units)
    sigma can be a variable (see below)
  cutoff = distance from wall at which wall\sphinxhyphen{}particle interaction is cut off (distance units)
\end{sphinxalltt}
\begin{sphinxalltt}
args for style \sphinxstyleemphasis{wall/region/ees}: \sphinxstyleemphasis{region\sphinxhyphen{}ID} \sphinxstyleemphasis{epsilon} \sphinxstyleemphasis{sigma} \sphinxstyleemphasis{cutoff}
  region\sphinxhyphen{}ID = region whose boundary will act as wall
  epsilon = strength factor for wall\sphinxhyphen{}particle interaction (energy or energy/distance\textasciicircum{}2 units)
  sigma = size factor for wall\sphinxhyphen{}particle interaction (distance units)
  cutoff = distance from wall at which wall\sphinxhyphen{}particle interaction is cut off (distance units)
\end{sphinxalltt}

\end{itemize}


\subsection{Examples}
\label{\detokenize{fix_wall_ees:examples}}
\begin{sphinxVerbatim}[commandchars=\\\{\}]
\PYG{k}{fix }\PYG{n+nv+nvIdentifier}{wallhi}\PYG{+w}{ }\PYG{n+nv+nvIdentifier}{all}\PYG{+w}{ }\PYG{n}{wall}\PYG{o}{/}\PYG{n}{ees}\PYG{+w}{ }\PYG{n}{xlo}\PYG{+w}{ }\PYG{o}{\PYGZhy{}}\PYG{l+m}{1.0}\PYG{+w}{ }\PYG{l+m}{1.0}\PYG{+w}{ }\PYG{l+m}{1.0}\PYG{+w}{ }\PYG{l+m}{2.5}\PYG{+w}{ }\PYG{n}{units}\PYG{+w}{ }\PYG{n}{box}
\PYG{k}{fix }\PYG{n+nv+nvIdentifier}{wallhi}\PYG{+w}{ }\PYG{n+nv+nvIdentifier}{all}\PYG{+w}{ }\PYG{n}{wall}\PYG{o}{/}\PYG{n}{ees}\PYG{+w}{ }\PYG{n}{xhi}\PYG{+w}{ }\PYG{n}{EDGE}\PYG{+w}{ }\PYG{l+m}{1.0}\PYG{+w}{ }\PYG{l+m}{1.0}\PYG{+w}{ }\PYG{l+m}{2.5}
\PYG{k}{fix }\PYG{n+nv+nvIdentifier}{wallhi}\PYG{+w}{ }\PYG{n+nv+nvIdentifier}{all}\PYG{+w}{ }\PYG{n}{wall}\PYG{o}{/}\PYG{n}{ees}\PYG{+w}{ }\PYG{n}{v\PYGZus{}wiggle}\PYG{+w}{ }\PYG{l+m}{23.2}\PYG{+w}{ }\PYG{l+m}{1.0}\PYG{+w}{ }\PYG{l+m}{1.0}\PYG{+w}{ }\PYG{l+m}{2.5}
\PYG{k}{fix }\PYG{n+nv+nvIdentifier}{zwalls}\PYG{+w}{ }\PYG{n+nv+nvIdentifier}{all}\PYG{+w}{ }\PYG{n}{wall}\PYG{o}{/}\PYG{n}{ees}\PYG{+w}{ }\PYG{n}{zlo}\PYG{+w}{ }\PYG{l+m}{0.0}\PYG{+w}{ }\PYG{l+m}{1.0}\PYG{+w}{ }\PYG{l+m}{1.0}\PYG{+w}{ }\PYG{l+m}{0.858}\PYG{+w}{ }\PYG{n}{zhi}\PYG{+w}{ }\PYG{l+m}{40.0}\PYG{+w}{ }\PYG{l+m}{1.0}\PYG{+w}{ }\PYG{l+m}{1.0}\PYG{+w}{ }\PYG{l+m}{0.858}

\PYG{k}{fix }\PYG{n+nv+nvIdentifier}{ees\PYGZus{}cube}\PYG{+w}{ }\PYG{n+nv+nvIdentifier}{all}\PYG{+w}{ }\PYG{n}{wall}\PYG{o}{/}\PYG{n}{region}\PYG{o}{/}\PYG{n}{ees}\PYG{+w}{ }\PYG{n}{myCube}\PYG{+w}{ }\PYG{l+m}{1.0}\PYG{+w}{ }\PYG{l+m}{1.0}\PYG{+w}{ }\PYG{l+m}{2.5}
\end{sphinxVerbatim}


\subsection{Description}
\label{\detokenize{fix_wall_ees:description}}
\sphinxAtStartPar
Fix \sphinxstyleemphasis{wall/ees} bounds the simulation domain on one or more of its
faces with a flat wall that interacts with the ellipsoidal atoms in
the group by generating a force on the atom in a direction
perpendicular to the wall and a torque parallel with the wall.  The
energy of wall\sphinxhyphen{}particle interactions E is given by:
\begin{equation*}
\begin{split}E = \epsilon \left[ \frac{2 \sigma_{LJ}^{12} \left(7 r^5+14 r^3
\sigma_{n}^2+3 r \sigma_{n}^4\right) }{945
\left(r^2-\sigma_{n}^2\right)^7} -\frac{ \sigma_{LJ}^6 \left(2 r
\sigma_{n}^3+\sigma_{n}^2 \left(r^2-\sigma_{n}^2\right)\log{
\left[\frac{r-\sigma_{n}}{r+\sigma_{n}}\right]}\right) }{12
\sigma_{n}^5 \left(r^2-\sigma_{n}^2\right)} \right]\qquad \sigma_n
< r < r_c\end{split}
\end{equation*}
\sphinxAtStartPar
Introduced by Babadi and Ejtehadi in {\hyperref[\detokenize{fix_wall_ees:babadiejtehadi}]{\sphinxcrossref{\DUrole{std,std-ref}{(Babadi2)}}}}. Here, \sphinxstyleemphasis{r} is the distance from the particle to the
wall at position \sphinxstyleemphasis{coord}, and Rc is the \sphinxstyleemphasis{cutoff} distance at which
the particle and wall no longer interact. Also, \(\sigma_n\) is
the distance between center of ellipsoid and the nearest point of its
surface to the wall as shown below.

\noindent{\hspace*{\fill}\sphinxincludegraphics{{fix_wall_ees_image}.jpg}\hspace*{\fill}}

\sphinxAtStartPar
Details of using this command and specifications are the same as
fix/wall command. You can also find an example in USER/ees/ under
examples/ directory.

\sphinxAtStartPar
The prefactor \(\epsilon\) can be thought of as an
effective Hamaker constant with energy units for the strength of the
ellipsoid\sphinxhyphen{}wall interaction.  More specifically, the \(\epsilon\)
prefactor is
\begin{equation*}
\begin{split}8 \pi^2 \quad \rho_{wall} \quad \rho_{ellipsoid} \quad \epsilon
\quad \sigma_a \quad \sigma_b \quad \sigma_c\end{split}
\end{equation*}
\sphinxAtStartPar
where \(\epsilon\) is the LJ energy parameter for the constituent
LJ particles and \(\sigma_a\), \(\sigma_b\), and
\(\sigma_c\) are the radii of the ellipsoidal
particles. \(\rho_{wall}\) and \(\rho_{ellipsoid}\) are the
number density of the constituent particles, in the wall and ellipsoid
respectively, in units of 1/volume.

\begin{sphinxadmonition}{note}{Note:}
\sphinxAtStartPar
You must ensure that r is always bigger than \(\sigma_n\) for
all particles in the group, or LAMMPS will generate an error.  This
means you cannot start your simulation with particles touching the wall
position \sphinxstyleemphasis{coord} (\(r = \sigma_n\)) or with particles penetrating
the wall (\(0 =< r < \sigma_n\)) or with particles on the wrong
side of the wall (\(r < 0\)).
\end{sphinxadmonition}

\sphinxAtStartPar
Fix \sphinxstyleemphasis{wall/region/ees} treats the surface of the geometric region defined
by the \sphinxstyleemphasis{region\sphinxhyphen{}ID} as a bounding wall which interacts with nearby
ellipsoidal particles according to the EES potential introduced above.

\sphinxAtStartPar
Other details of this command are the same as for the {\hyperref[\detokenize{fix_wall_region::doc}]{\sphinxcrossref{\DUrole{doc}{fix
wall/region}}}} command.  One may also find an example
of using this fix in the examples/PACKAGES/ees/ directory.


\bigskip\hrule\bigskip



\subsection{Restart, fix\_modify, output, run start/stop, minimize info}
\label{\detokenize{fix_wall_ees:restart-fix-modify-output-run-start-stop-minimize-info}}
\sphinxAtStartPar
No information about these fixes are written to {\hyperref[\detokenize{restart::doc}]{\sphinxcrossref{\DUrole{doc}{binary restart
files}}}}.

\sphinxAtStartPar
The {\hyperref[\detokenize{fix_modify::doc}]{\sphinxcrossref{\DUrole{doc}{fix\_modify}}}} \sphinxstyleemphasis{energy} option is supported by
these fixes to add the energy of interaction between atoms and all the
specified walls or region wall to the global potential energy of the
system as part of {\hyperref[\detokenize{thermo_style::doc}]{\sphinxcrossref{\DUrole{doc}{thermodynamic output}}}}.  The
default settings for these fixes are {\hyperref[\detokenize{fix_modify::doc}]{\sphinxcrossref{\DUrole{doc}{fix\_modify energy no}}}}.

\sphinxAtStartPar
The {\hyperref[\detokenize{fix_modify::doc}]{\sphinxcrossref{\DUrole{doc}{fix\_modify}}}} \sphinxstyleemphasis{respa} option is supported by
these fixes. This allows to set at which level of the {\hyperref[\detokenize{run_style::doc}]{\sphinxcrossref{\DUrole{doc}{r\sphinxhyphen{}RESPA}}}} integrator the fix is adding its forces. Default is the
outermost level.

\sphinxAtStartPar
These fixes computes a global scalar and a global vector of forces,
which can be accessed by various {\hyperref[\detokenize{Howto_output::doc}]{\sphinxcrossref{\DUrole{doc}{output commands}}}}.  See the {\hyperref[\detokenize{fix_wall::doc}]{\sphinxcrossref{\DUrole{doc}{fix wall}}}} command for a
description of the scalar and vector.

\sphinxAtStartPar
No parameter of these fixes can be used with the \sphinxstyleemphasis{start/stop} keywords of
the {\hyperref[\detokenize{run::doc}]{\sphinxcrossref{\DUrole{doc}{run}}}} command.

\sphinxAtStartPar
The forces due to these fixes are imposed during an energy
minimization, invoked by the {\hyperref[\detokenize{minimize::doc}]{\sphinxcrossref{\DUrole{doc}{minimize}}}} command.

\begin{sphinxadmonition}{note}{Note:}
\sphinxAtStartPar
If you want the atom/wall interaction energy to be included in
the total potential energy of the system (the quantity being
minimized), you MUST enable the {\hyperref[\detokenize{fix_modify::doc}]{\sphinxcrossref{\DUrole{doc}{fix\_modify}}}} \sphinxstyleemphasis{energy}
option for this fix.
\end{sphinxadmonition}


\subsection{Restrictions}
\label{\detokenize{fix_wall_ees:restrictions}}
\sphinxAtStartPar
These fixes are part of the EXTRA\sphinxhyphen{}FIX package.  They are only enabled
if LAMMPS was built with that package.  See the {\hyperref[\detokenize{Build_package::doc}]{\sphinxcrossref{\DUrole{doc}{Build package}}}} page for more info.

\sphinxAtStartPar
These fixes requires that atoms be ellipsoids as defined by the
{\hyperref[\detokenize{atom_style::doc}]{\sphinxcrossref{\DUrole{doc}{atom\_style ellipsoid}}}} command.


\subsection{Related commands}
\label{\detokenize{fix_wall_ees:related-commands}}
\sphinxAtStartPar
{\hyperref[\detokenize{fix_wall::doc}]{\sphinxcrossref{\DUrole{doc}{fix wall}}}},
{\hyperref[\detokenize{pair_resquared::doc}]{\sphinxcrossref{\DUrole{doc}{pair resquared}}}}


\subsection{Default}
\label{\detokenize{fix_wall_ees:default}}
\sphinxAtStartPar
none


\bigskip\hrule\bigskip

\phantomsection\label{\detokenize{fix_wall_ees:babadiejtehadi}}
\sphinxAtStartPar
\sphinxstylestrong{(Babadi2)} Babadi and Ejtehadi, EPL, 77 (2007) 23002.

\sphinxstepscope

\index{fix wall/flow@\spxentry{fix wall/flow}}\ignorespaces 
\index{fix wall/flow/kk@\spxentry{fix wall/flow/kk}}\ignorespaces 

\section{fix wall/flow command}
\label{\detokenize{fix_wall_flow:fix-wall-flow-command}}\label{\detokenize{fix_wall_flow:index-1}}\label{\detokenize{fix_wall_flow:index-0}}\label{\detokenize{fix_wall_flow::doc}}
\sphinxAtStartPar
Accelerator Variants: \sphinxstyleemphasis{wall/flow/kk}


\subsection{Syntax}
\label{\detokenize{fix_wall_flow:syntax}}
\begin{sphinxVerbatim}[commandchars=\\\{\}]
\PYG{k}{fix }\PYG{n+nv+nvIdentifier}{ID}\PYG{+w}{ }\PYG{n+nv+nvIdentifier}{group\PYGZhy{}ID}\PYG{+w}{ }\PYG{n}{wall}\PYG{o}{/}\PYG{n}{flow}\PYG{+w}{ }\PYG{n}{axis}\PYG{+w}{ }\PYG{n}{vflow}\PYG{+w}{ }\PYG{n}{T}\PYG{+w}{ }\PYG{n}{seed}\PYG{+w}{ }\PYG{n}{N}\PYG{+w}{ }\PYG{n}{coords}\PYG{+w}{ }\PYG{n}{...}\PYG{+w}{ }\PYG{n}{keyword}\PYG{+w}{ }\PYG{n}{value}
\end{sphinxVerbatim}
\begin{itemize}
\item {} 
\sphinxAtStartPar
ID, group\sphinxhyphen{}ID are documented in {\hyperref[\detokenize{fix::doc}]{\sphinxcrossref{\DUrole{doc}{fix}}}} command

\item {} 
\sphinxAtStartPar
wall/flow = style name of this fix command

\item {} 
\sphinxAtStartPar
axis = flow axis (\sphinxstyleemphasis{x}, \sphinxstyleemphasis{y}, or \sphinxstyleemphasis{z})

\item {} 
\sphinxAtStartPar
vflow = generated flow velocity in \sphinxstyleemphasis{axis} direction (velocity units)

\item {} 
\sphinxAtStartPar
T = flow temperature (temperature units)

\item {} 
\sphinxAtStartPar
seed = random seed for stochasticity (positive integer)

\item {} 
\sphinxAtStartPar
N = number of walls

\item {} 
\sphinxAtStartPar
coords = list of N wall positions along the \sphinxstyleemphasis{axis} direction in ascending order (distance units)

\item {} 
\sphinxAtStartPar
zero or more keyword/value pairs may be appended

\item {} 
\sphinxAtStartPar
keyword = \sphinxstyleemphasis{units}
\begin{sphinxalltt}
\sphinxstyleemphasis{units} value = \sphinxstyleemphasis{lattice} or \sphinxstyleemphasis{box}
  \sphinxstyleemphasis{lattice} = wall positions are defined in lattice units
  \sphinxstyleemphasis{box} = the wall positions are defined in simulation box units
\end{sphinxalltt}

\end{itemize}


\subsection{Examples}
\label{\detokenize{fix_wall_flow:examples}}
\begin{sphinxVerbatim}[commandchars=\\\{\}]
\PYG{k}{fix }\PYG{n+nv+nvIdentifier}{1}\PYG{+w}{ }\PYG{n+nv+nvIdentifier}{all}\PYG{+w}{ }\PYG{n}{wall}\PYG{o}{/}\PYG{n}{flow}\PYG{+w}{ }\PYG{n}{x}\PYG{+w}{ }\PYG{l+m}{0.4}\PYG{+w}{ }\PYG{l+m}{1.5}\PYG{+w}{ }\PYG{l+m}{593894}\PYG{+w}{ }\PYG{l+m}{4}\PYG{+w}{ }\PYG{l+m}{2.0}\PYG{+w}{ }\PYG{l+m}{4.0}\PYG{+w}{ }\PYG{l+m}{6.0}\PYG{+w}{ }\PYG{l+m}{8.0}
\end{sphinxVerbatim}


\subsection{Description}
\label{\detokenize{fix_wall_flow:description}}
\sphinxAtStartPar
\DUrole{versionmodified,added}{Added in version 17Apr2024.}

\sphinxAtStartPar
This fix implements flow boundary conditions (FBC) introduced in
{\hyperref[\detokenize{fix_wall_flow:fbc-pavlov1}]{\sphinxcrossref{\DUrole{std,std-ref}{(Pavlov1)}}}} and {\hyperref[\detokenize{fix_wall_flow:fbc-pavlov2}]{\sphinxcrossref{\DUrole{std,std-ref}{(Pavlov2)}}}}.
The goal is to generate a stationary flow with a shifted Maxwell
velocity distribution:
\begin{equation*}
\begin{split}f_a(v_a) \propto \exp{\left(-\frac{m (v_a-v_{\text{flow}})^2}{2 kB T}\right)}\end{split}
\end{equation*}
\sphinxAtStartPar
where \(v_a\) is the component of velocity along the specified
\sphinxstyleemphasis{axis} argument (a = x,y,z), \(v_{\text{flow}}\) is the flow
velocity specified as the \sphinxstyleemphasis{vflow} argument, \sphinxstyleemphasis{T} is the specified flow
temperature, \sphinxstyleemphasis{m} is the particle mass, and \sphinxstyleemphasis{kB} is the Boltzmann
constant.

\sphinxAtStartPar
This is achieved by defining a series of \sphinxstyleemphasis{N} transparent walls along
the flow \sphinxstyleemphasis{axis} direction.  Each wall is at the specified position
listed in the \sphinxstyleemphasis{coords} argument.  Note that an additional transparent
wall is defined by the code at the boundary of the (periodic)
simulation domain in the \sphinxstyleemphasis{axis} direction.  So there are effectively
N+1 walls.

\sphinxAtStartPar
Each time a particle in the specified group passes through one of the
transparent walls, its velocity is re\sphinxhyphen{}assigned.  Particles not in the
group do not interact with the wall. This can be used, for example, to
add obstacles composed of atoms, or to simulate a solution of complex
molecules in a one\sphinxhyphen{}atom liquid (note that the fix has been tested for
one\sphinxhyphen{}atom systems only).

\sphinxAtStartPar
Conceptually, the velocity re\sphinxhyphen{}assignment represents creation of a new
particle within the system with simultaneous removal of the particle
which passed through the wall.  The velocity components in directions
parallel to the wall are re\sphinxhyphen{}assigned according to the standard Maxwell
velocity distribution for the specified temperature \sphinxstyleemphasis{T}.  The velocity
component perpendicular to the wall is re\sphinxhyphen{}assigned according to the
shifted Maxwell distribution defined above:
\begin{equation*}
\begin{split}f_{\text{a generated}}(v_a) \propto v_a f_a(v_a)\end{split}
\end{equation*}
\sphinxAtStartPar
It can be shown that for an ideal\sphinxhyphen{}gas scenario this procedure makes
the velocity distribution of particles between walls exactly as
desired.

\sphinxAtStartPar
Since in most cases simulated systems are not an ideal gas, multiple
walls can be defined, since a single wall may not be sufficient for
maintaining a stationary flow without “congestion” which can manifest
itself as regions in the flow with increased particle density located
upstream from static obstacles.

\sphinxAtStartPar
For the same reason, the actual temperature and velocity of the
generated flow may differ from what is requested.  The degree of
discrepancy is determined by how different from an ideal gas the
simulated system is.  Therefore, a calibration procedure may be
required for such a system as described in {\hyperref[\detokenize{fix_wall_flow:fbc-pavlov2}]{\sphinxcrossref{\DUrole{std,std-ref}{(Pavlov)}}}}.

\sphinxAtStartPar
Note that the interactions between particles on different sides of a
transparent wall are not disabled or neglected.  Likewise particle
positions are not altered by the velocity reassignment.  This removes
the need to modify the force field to work correctly in cases when a
particle is close to a wall.

\sphinxAtStartPar
For example, if particle positions were uniformly redistributed across
the surface of a wall, two particles could end up too close to each
other, potentially causing the simulation to explode.  However due to
this compromise, some collective phenomena such as regions with
increased/decreased density or collective movements are not fully
removed when particles cross a wall.  This unwanted consequence can
also be potentially mitigated by using more multiple walls.

\begin{sphinxadmonition}{note}{Note:}
\sphinxAtStartPar
When the specified flow has a high velocity, a lost atoms error can
occur (see {\hyperref[\detokenize{Errors_messages::doc}]{\sphinxcrossref{\DUrole{doc}{error messages}}}}).  If this
happens, you should ensure the checks for neighbor list rebuilds,
set via the {\hyperref[\detokenize{neigh_modify::doc}]{\sphinxcrossref{\DUrole{doc}{neigh\_modify}}}} command, are as
conservative as possible (every timestep if needed).  Those are the
default settings.
\end{sphinxadmonition}


\subsection{Restart, fix\_modify, output, run start/stop, minimize info}
\label{\detokenize{fix_wall_flow:restart-fix-modify-output-run-start-stop-minimize-info}}
\sphinxAtStartPar
No information about this fix is written to {\hyperref[\detokenize{restart::doc}]{\sphinxcrossref{\DUrole{doc}{binary restart files}}}}.

\sphinxAtStartPar
None of the {\hyperref[\detokenize{fix_modify::doc}]{\sphinxcrossref{\DUrole{doc}{fix\_modify}}}} options are relevant to
this fix.

\sphinxAtStartPar
No global or per\sphinxhyphen{}atom quantities are stored by this fix for access by
various {\hyperref[\detokenize{Howto_output::doc}]{\sphinxcrossref{\DUrole{doc}{output commands}}}}.

\sphinxAtStartPar
No parameter of this fix can be used with the \sphinxstyleemphasis{start/stop} keywords of
the {\hyperref[\detokenize{run::doc}]{\sphinxcrossref{\DUrole{doc}{run}}}} command.

\sphinxAtStartPar
This fix is not invoked during {\hyperref[\detokenize{minimize::doc}]{\sphinxcrossref{\DUrole{doc}{energy minimization}}}}.


\subsection{Restrictions}
\label{\detokenize{fix_wall_flow:restrictions}}
\sphinxAtStartPar
Fix \sphinxstyleemphasis{wall\_flow} is part of the EXTRA\sphinxhyphen{}FIX package.  It is only enabled
if LAMMPS was built with that package.  See the {\hyperref[\detokenize{Build_package::doc}]{\sphinxcrossref{\DUrole{doc}{Build package}}}} page for more info.

\sphinxAtStartPar
Flow boundary conditions should not be used with rigid bodies such as
those defined by a “fix rigid” command.

\sphinxAtStartPar
This fix can only be used with periodic boundary conditions along the
flow axis. The size of the box in this direction must not change. Also,
the fix is designed to work only in an orthogonal simulation box.


\subsection{Related commands}
\label{\detokenize{fix_wall_flow:related-commands}}
\sphinxAtStartPar
{\hyperref[\detokenize{fix_wall::doc}]{\sphinxcrossref{\DUrole{doc}{fix wall/reflect}}}} command


\subsection{Default}
\label{\detokenize{fix_wall_flow:default}}
\sphinxAtStartPar
The default for the units keyword is lattice.


\bigskip\hrule\bigskip

\phantomsection\label{\detokenize{fix_wall_flow:fbc-pavlov1}}
\sphinxAtStartPar
\sphinxstylestrong{(Pavlov1)} Pavlov, Kolotinskii, Stegailov, “GPU\sphinxhyphen{}Based Molecular Dynamics of Turbulent Liquid Flows with OpenMM”, Proceedings of PPAM\sphinxhyphen{}2022, LNCS (Springer), vol. 13826, pp. 346\sphinxhyphen{}358 (2023)

\phantomsection\label{\detokenize{fix_wall_flow:fbc-pavlov2}}
\sphinxAtStartPar
\sphinxstylestrong{(Pavlov2)} Pavlov, Galigerov, Kolotinskii, Nikolskiy, Stegailov, “GPU\sphinxhyphen{}based Molecular Dynamics of Fluid Flows: Reaching for Turbulence”, Int. J. High Perf. Comp. Appl., (2024)

\sphinxstepscope

\index{fix wall/gran@\spxentry{fix wall/gran}}\ignorespaces 
\index{fix wall/gran/kk@\spxentry{fix wall/gran/kk}}\ignorespaces 

\section{fix wall/gran command}
\label{\detokenize{fix_wall_gran:fix-wall-gran-command}}\label{\detokenize{fix_wall_gran:index-1}}\label{\detokenize{fix_wall_gran:index-0}}\label{\detokenize{fix_wall_gran::doc}}
\sphinxAtStartPar
Accelerator Variants: \sphinxstyleemphasis{wall/gran/kk}


\subsection{Syntax}
\label{\detokenize{fix_wall_gran:syntax}}
\begin{sphinxVerbatim}[commandchars=\\\{\}]
\PYG{k}{fix }\PYG{n+nv+nvIdentifier}{ID}\PYG{+w}{ }\PYG{n+nv+nvIdentifier}{group\PYGZhy{}ID}\PYG{+w}{ }\PYG{n}{wall}\PYG{o}{/}\PYG{n}{gran}\PYG{+w}{ }\PYG{n}{fstyle}\PYG{+w}{ }\PYG{n}{fstyle\PYGZus{}params}\PYG{+w}{ }\PYG{n}{wallstyle}\PYG{+w}{ }\PYG{n}{args}\PYG{+w}{ }\PYG{n}{keyword}\PYG{+w}{ }\PYG{n}{values}\PYG{+w}{ }\PYG{n}{...}
\end{sphinxVerbatim}
\begin{itemize}
\item {} 
\sphinxAtStartPar
ID, group\sphinxhyphen{}ID are documented in {\hyperref[\detokenize{fix::doc}]{\sphinxcrossref{\DUrole{doc}{fix}}}} command

\item {} 
\sphinxAtStartPar
wall/gran = style name of this fix command

\item {} 
\sphinxAtStartPar
fstyle = style of force interactions between particles and wall

\begin{sphinxVerbatim}[commandchars=\\\{\}]
possible choices: hooke, hooke/history, hertz/history, granular
\end{sphinxVerbatim}

\item {} 
\sphinxAtStartPar
fstyle\_params = parameters associated with force interaction style
\begin{sphinxalltt}
For \sphinxstyleemphasis{hooke}, \sphinxstyleemphasis{hooke/history}, and \sphinxstyleemphasis{hertz/history}, \sphinxstyleemphasis{fstyle\_params} are:
      Kn = elastic constant for normal particle repulsion (force/distance units or pressure units \sphinxhyphen{} see discussion below)
      Kt = elastic constant for tangential contact (force/distance units or pressure units \sphinxhyphen{} see discussion below)
      gamma\_n = damping coefficient for collisions in normal direction (1/time units or 1/time\sphinxhyphen{}distance units \sphinxhyphen{} see discussion below)
      gamma\_t = damping coefficient for collisions in tangential direction (1/time units or 1/time\sphinxhyphen{}distance units \sphinxhyphen{} see discussion below)
      xmu = static yield criterion (unitless value between 0.0 and 1.0e4)
      dampflag = 0 or 1 if tangential damping force is excluded or included
      optional keyword = \sphinxstyleemphasis{limit\_damping}, limit damping to prevent attractive interaction
\end{sphinxalltt}
\begin{sphinxalltt}
For \sphinxstyleemphasis{granular}, \sphinxstyleemphasis{fstyle\_params} are set using the same syntax as for the \sphinxstyleemphasis{pair\_coeff} command of {\hyperref[\detokenize{pair_granular::doc}]{\sphinxcrossref{\DUrole{doc}{pair\_style granular}}}}
\end{sphinxalltt}

\item {} 
\sphinxAtStartPar
wallstyle = \sphinxstyleemphasis{xplane} or \sphinxstyleemphasis{yplane} or \sphinxstyleemphasis{zplane} or \sphinxstyleemphasis{zcylinder}

\item {} 
\sphinxAtStartPar
args = list of arguments for a particular style
\begin{sphinxalltt}
\sphinxstyleemphasis{xplane} or \sphinxstyleemphasis{yplane} or \sphinxstyleemphasis{zplane} args = lo hi
  lo,hi = position of lower and upper plane (distance units), either can be NULL)
\sphinxstyleemphasis{zcylinder} args = radius
  radius = cylinder radius (distance units)
\end{sphinxalltt}

\item {} 
\sphinxAtStartPar
zero or more keyword/value pairs may be appended to args

\item {} 
\sphinxAtStartPar
keyword = \sphinxstyleemphasis{wiggle} or \sphinxstyleemphasis{shear} or \sphinxstyleemphasis{contacts} or \sphinxstyleemphasis{temperature}
\begin{sphinxalltt}
 \sphinxstyleemphasis{wiggle} values = dim amplitude period
   dim = \sphinxstyleemphasis{x} or \sphinxstyleemphasis{y} or \sphinxstyleemphasis{z}
   amplitude = size of oscillation (distance units)
   period = time of oscillation (time units)
 \sphinxstyleemphasis{shear} values = dim vshear
   dim = \sphinxstyleemphasis{x} or \sphinxstyleemphasis{y} or \sphinxstyleemphasis{z}
   vshear = magnitude of shear velocity (velocity units)
\sphinxstyleemphasis{contacts} value = none
   generate contact information for each particle
\sphinxstyleemphasis{temperature} value = temperature
   specify temperature of wall
\end{sphinxalltt}

\end{itemize}


\subsection{Examples}
\label{\detokenize{fix_wall_gran:examples}}
\begin{sphinxVerbatim}[commandchars=\\\{\}]
\PYG{k}{fix }\PYG{n+nv+nvIdentifier}{1}\PYG{+w}{ }\PYG{n+nv+nvIdentifier}{all}\PYG{+w}{ }\PYG{n}{wall}\PYG{o}{/}\PYG{n}{gran}\PYG{+w}{ }\PYG{n}{hooke}\PYG{+w}{  }\PYG{l+m}{200000.0}\PYG{+w}{ }\PYG{n}{NULL}\PYG{+w}{ }\PYG{l+m}{50.0}\PYG{+w}{ }\PYG{n}{NULL}\PYG{+w}{ }\PYG{l+m}{0.5}\PYG{+w}{ }\PYG{l+m}{0}\PYG{+w}{ }\PYG{n}{xplane}\PYG{+w}{ }\PYG{o}{\PYGZhy{}}\PYG{l+m}{10.0}\PYG{+w}{ }\PYG{l+m}{10.0}
\PYG{k}{fix }\PYG{n+nv+nvIdentifier}{1}\PYG{+w}{ }\PYG{n+nv+nvIdentifier}{all}\PYG{+w}{ }\PYG{n}{wall}\PYG{o}{/}\PYG{n}{gran}\PYG{+w}{ }\PYG{n}{hooke}\PYG{o}{/}\PYG{n}{history}\PYG{+w}{ }\PYG{l+m}{200000.0}\PYG{+w}{ }\PYG{n}{NULL}\PYG{+w}{ }\PYG{l+m}{50.0}\PYG{+w}{ }\PYG{n}{NULL}\PYG{+w}{ }\PYG{l+m}{0.5}\PYG{+w}{ }\PYG{l+m}{0}\PYG{+w}{ }\PYG{n}{zplane}\PYG{+w}{ }\PYG{l+m}{0.0}\PYG{+w}{ }\PYG{n}{NULL}
\PYG{k}{fix }\PYG{n+nv+nvIdentifier}{2}\PYG{+w}{ }\PYG{n+nv+nvIdentifier}{all}\PYG{+w}{ }\PYG{n}{wall}\PYG{o}{/}\PYG{n}{gran}\PYG{+w}{ }\PYG{n}{hooke}\PYG{+w}{ }\PYG{l+m}{100000.0}\PYG{+w}{ }\PYG{l+m}{20000.0}\PYG{+w}{ }\PYG{l+m}{50.0}\PYG{+w}{ }\PYG{l+m}{30.0}\PYG{+w}{ }\PYG{l+m}{0.5}\PYG{+w}{ }\PYG{l+m}{1}\PYG{+w}{ }\PYG{n}{zcylinder}\PYG{+w}{ }\PYG{l+m}{15.0}\PYG{+w}{ }\PYG{n}{wiggle}\PYG{+w}{ }\PYG{n}{z}\PYG{+w}{ }\PYG{l+m}{3.0}\PYG{+w}{ }\PYG{l+m}{2.0}
\PYG{k}{fix }\PYG{n+nv+nvIdentifier}{3}\PYG{+w}{ }\PYG{n+nv+nvIdentifier}{all}\PYG{+w}{ }\PYG{n}{wall}\PYG{o}{/}\PYG{n}{gran}\PYG{+w}{ }\PYG{n}{granular}\PYG{+w}{ }\PYG{n}{hooke}\PYG{+w}{ }\PYG{l+m}{1000.0}\PYG{+w}{ }\PYG{l+m}{50.0}\PYG{+w}{ }\PYG{n}{tangential}\PYG{+w}{ }\PYG{n}{linear\PYGZus{}nohistory}\PYG{+w}{ }\PYG{l+m}{1.0}\PYG{+w}{ }\PYG{l+m}{0.4}\PYG{+w}{ }\PYG{n}{damping}\PYG{+w}{ }\PYG{k}{velocity }\PYG{n+nv+nvIdentifier}{region}\PYG{+w}{ }\PYG{n}{myBox}
\PYG{k}{fix }\PYG{n+nv+nvIdentifier}{4}\PYG{+w}{ }\PYG{n+nv+nvIdentifier}{all}\PYG{+w}{ }\PYG{n}{wall}\PYG{o}{/}\PYG{n}{gran}\PYG{+w}{ }\PYG{n}{granular}\PYG{+w}{ }\PYG{n}{jkr}\PYG{+w}{ }\PYG{l+m}{1e5}\PYG{+w}{ }\PYG{l+m}{1500.0}\PYG{+w}{ }\PYG{l+m}{0.3}\PYG{+w}{ }\PYG{l+m}{10.0}\PYG{+w}{ }\PYG{n}{tangential}\PYG{+w}{ }\PYG{n}{mindlin}\PYG{+w}{ }\PYG{n}{NULL}\PYG{+w}{ }\PYG{l+m}{1.0}\PYG{+w}{ }\PYG{l+m}{0.5}\PYG{+w}{ }\PYG{n}{rolling}\PYG{+w}{ }\PYG{n}{sds}\PYG{+w}{ }\PYG{l+m}{500.0}\PYG{+w}{ }\PYG{l+m}{200.0}\PYG{+w}{ }\PYG{l+m}{0.5}\PYG{+w}{ }\PYG{n}{twisting}\PYG{+w}{ }\PYG{n}{marshall}\PYG{+w}{ }\PYG{k}{region }\PYG{n+nv+nvIdentifier}{myCone}
\PYG{k}{fix }\PYG{n+nv+nvIdentifier}{5}\PYG{+w}{ }\PYG{n+nv+nvIdentifier}{all}\PYG{+w}{ }\PYG{n}{wall}\PYG{o}{/}\PYG{n}{gran}\PYG{+w}{ }\PYG{n}{granular}\PYG{+w}{ }\PYG{n}{dmt}\PYG{+w}{ }\PYG{l+m}{1e5}\PYG{+w}{ }\PYG{l+m}{0.2}\PYG{+w}{ }\PYG{l+m}{0.3}\PYG{+w}{ }\PYG{l+m}{10.0}\PYG{+w}{ }\PYG{n}{tangential}\PYG{+w}{ }\PYG{n}{mindlin}\PYG{+w}{ }\PYG{n}{NULL}\PYG{+w}{ }\PYG{l+m}{1.0}\PYG{+w}{ }\PYG{l+m}{0.5}\PYG{+w}{ }\PYG{n}{rolling}\PYG{+w}{ }\PYG{n}{sds}\PYG{+w}{ }\PYG{l+m}{500.0}\PYG{+w}{ }\PYG{l+m}{200.0}\PYG{+w}{ }\PYG{l+m}{0.5}\PYG{+w}{ }\PYG{n}{twisting}\PYG{+w}{ }\PYG{n}{marshall}\PYG{+w}{ }\PYG{n}{damping}\PYG{+w}{ }\PYG{n}{tsuji}\PYG{+w}{ }\PYG{n}{heat}\PYG{+w}{ }\PYG{l+m}{10}\PYG{+w}{ }\PYG{k}{region }\PYG{n+nv+nvIdentifier}{myCone}\PYG{+w}{ }\PYG{n}{temperature}\PYG{+w}{ }\PYG{l+m}{1.0}
\PYG{k}{fix }\PYG{n+nv+nvIdentifier}{6}\PYG{+w}{ }\PYG{n+nv+nvIdentifier}{all}\PYG{+w}{ }\PYG{n}{wall}\PYG{o}{/}\PYG{n}{gran}\PYG{+w}{ }\PYG{n}{hooke}\PYG{+w}{  }\PYG{l+m}{200000.0}\PYG{+w}{ }\PYG{n}{NULL}\PYG{+w}{ }\PYG{l+m}{50.0}\PYG{+w}{ }\PYG{n}{NULL}\PYG{+w}{ }\PYG{l+m}{0.5}\PYG{+w}{ }\PYG{l+m}{0}\PYG{+w}{ }\PYG{n}{xplane}\PYG{+w}{ }\PYG{o}{\PYGZhy{}}\PYG{l+m}{10.0}\PYG{+w}{ }\PYG{l+m}{10.0}\PYG{+w}{ }\PYG{n}{contacts}
\end{sphinxVerbatim}


\subsection{Description}
\label{\detokenize{fix_wall_gran:description}}
\sphinxAtStartPar
Bound the simulation domain of a granular system with a frictional
wall.  All particles in the group interact with the wall when they are
close enough to touch it.

\sphinxAtStartPar
The nature of the wall/particle interactions are determined by the
\sphinxstyleemphasis{fstyle} setting.  It can be any of the styles defined by the
{\hyperref[\detokenize{pair_gran::doc}]{\sphinxcrossref{\DUrole{doc}{pair\_style gran/*}}}} or the more general
{\hyperref[\detokenize{pair_granular::doc}]{\sphinxcrossref{\DUrole{doc}{pair\_style granular}}}} commands.  Currently the
options are \sphinxstyleemphasis{hooke}, \sphinxstyleemphasis{hooke/history}, or \sphinxstyleemphasis{hertz/history} for the
former, and \sphinxstyleemphasis{granular} with all the possible options of the associated
\sphinxstyleemphasis{pair\_coeff} command for the latter.  The equation for the force
between the wall and particles touching it is the same as the
corresponding equation on the {\hyperref[\detokenize{pair_gran::doc}]{\sphinxcrossref{\DUrole{doc}{pair\_style gran/*}}}} and
{\hyperref[\detokenize{pair_granular::doc}]{\sphinxcrossref{\DUrole{doc}{pair\_style granular}}}} doc pages, in the limit of
one of the two particles going to infinite radius and mass (flat wall).
Specifically, delta = radius \sphinxhyphen{} r = overlap of particle with wall, m\_eff
= mass of particle, and the effective radius of contact = RiRj/Ri+Rj is
set to the radius of the particle.

\sphinxAtStartPar
The parameters \sphinxstyleemphasis{Kn}, \sphinxstyleemphasis{Kt}, \sphinxstyleemphasis{gamma\_n}, \sphinxstyleemphasis{gamma\_t}, \sphinxstyleemphasis{xmu}, \sphinxstyleemphasis{dampflag},
and the optional keyword \sphinxstyleemphasis{limit\_damping}
have the same meaning and units as those specified with the
{\hyperref[\detokenize{pair_gran::doc}]{\sphinxcrossref{\DUrole{doc}{pair\_style gran/*}}}} commands.  This means a NULL can be
used for either \sphinxstyleemphasis{Kt} or \sphinxstyleemphasis{gamma\_t} as described on that page.  If a
NULL is used for \sphinxstyleemphasis{Kt}, then a default value is used where \sphinxstyleemphasis{Kt} = 2/7
\sphinxstyleemphasis{Kn}.  If a NULL is used for \sphinxstyleemphasis{gamma\_t}, then a default value is used
where \sphinxstyleemphasis{gamma\_t} = 1/2 \sphinxstyleemphasis{gamma\_n}.

\sphinxAtStartPar
All the model choices for cohesion, tangential friction, rolling
friction and twisting friction supported by the {\hyperref[\detokenize{pair_granular::doc}]{\sphinxcrossref{\DUrole{doc}{pair\_style granular}}}} through its \sphinxstyleemphasis{pair\_coeff} command are also
supported for walls. These are discussed in greater detail on the doc
page for {\hyperref[\detokenize{pair_granular::doc}]{\sphinxcrossref{\DUrole{doc}{pair\_style granular}}}}.

\begin{sphinxadmonition}{note}{Note:}
\sphinxAtStartPar
When \sphinxstyleemphasis{fstyle} \sphinxstyleemphasis{granular} is specified, the associated \sphinxstyleemphasis{fstyle\_params} are taken as
those for a wall/particle interaction. For example, for the \sphinxstyleemphasis{hertz/material} normal
contact model with \(E = 960\) and \(\nu = 0.2\), the effective Young’s
modulus for a wall/particle interaction is computed as
\(E_{eff} = \frac{960}{2(1-0.2^2)} = 500\). Any pair coefficients defined by
{\hyperref[\detokenize{pair_granular::doc}]{\sphinxcrossref{\DUrole{doc}{pair\_style granular}}}} are not taken into consideration. To
model different wall/particle interactions for particles of different material
types, the user may define multiple fix wall/gran commands operating on separate
groups (e.g. based on particle type) each with a different wall/particle effective
Young’s modulus.
\end{sphinxadmonition}

\sphinxAtStartPar
Note that you can choose a different force styles and/or different
values for the wall/particle coefficients than for particle/particle
interactions.  E.g. if you wish to model the wall as a different
material.

\begin{sphinxadmonition}{note}{Note:}
\sphinxAtStartPar
As discussed on the page for {\hyperref[\detokenize{pair_gran::doc}]{\sphinxcrossref{\DUrole{doc}{pair\_style gran/*}}}},
versions of LAMMPS before 9Jan09 used a different equation for
Hertzian interactions.  This means Hertizian wall/particle
interactions have also changed.  They now include a sqrt(radius) term
which was not present before.  Also the previous versions used Kn and
Kt from the pairwise interaction and hardwired dampflag to 1, rather
than letting them be specified directly.  This means you can set the
values of the wall/particle coefficients appropriately in the current
code to reproduce the results of a previous Hertzian monodisperse
calculation.  For example, for the common case of a monodisperse
system with particles of diameter 1, Kn, Kt, gamma\_n, and gamma\_s
should be set sqrt(2.0) larger than they were previously.
\end{sphinxadmonition}

\sphinxAtStartPar
The effective mass \sphinxstyleemphasis{m\_eff} in the formulas listed on the {\hyperref[\detokenize{pair_gran::doc}]{\sphinxcrossref{\DUrole{doc}{pair\_style granular}}}} page is the mass of the particle for
particle/wall interactions (mass of wall is infinite).  If the
particle is part of a rigid body, its mass is replaced by the mass of
the rigid body in those formulas.  This is determined by searching for
a {\hyperref[\detokenize{fix_rigid::doc}]{\sphinxcrossref{\DUrole{doc}{fix rigid}}}} command (or its variants).

\sphinxAtStartPar
The \sphinxstyleemphasis{wallstyle} can be planar or cylindrical.  The 3 planar options
specify a pair of walls in a dimension.  Wall positions are given by
\sphinxstyleemphasis{lo} and \sphinxstyleemphasis{hi}.  Either of the values can be specified as NULL if a
single wall is desired.  For a \sphinxstyleemphasis{zcylinder} wallstyle, the cylinder’s
axis is at x = y = 0.0, and the radius of the cylinder is specified.

\sphinxAtStartPar
Optionally, the wall can be moving, if the \sphinxstyleemphasis{wiggle} or \sphinxstyleemphasis{shear}
keywords are appended.  Both keywords cannot be used together.

\sphinxAtStartPar
For the \sphinxstyleemphasis{wiggle} keyword, the wall oscillates sinusoidally, similar to
the oscillations of particles which can be specified by the {\hyperref[\detokenize{fix_move::doc}]{\sphinxcrossref{\DUrole{doc}{fix move}}}} command.  This is useful in packing simulations of
granular particles.  The arguments to the \sphinxstyleemphasis{wiggle} keyword specify a
dimension for the motion, as well as it’s \sphinxstyleemphasis{amplitude} and \sphinxstyleemphasis{period}.
Note that if the dimension is in the plane of the wall, this is
effectively a shearing motion.  If the dimension is perpendicular to
the wall, it is more of a shaking motion.  A \sphinxstyleemphasis{zcylinder} wall can only
be wiggled in the z dimension.

\sphinxAtStartPar
Each timestep, the position of a wiggled wall in the appropriate \sphinxstyleemphasis{dim}
is set according to this equation:
\begin{sphinxalltt}
position = coord + A \sphinxhyphen{} A cos (omega * delta)
\end{sphinxalltt}

\sphinxAtStartPar
where \sphinxstyleemphasis{coord} is the specified initial position of the wall, \sphinxstyleemphasis{A} is
the \sphinxstyleemphasis{amplitude}, \sphinxstyleemphasis{omega} is 2 PI / \sphinxstyleemphasis{period}, and \sphinxstyleemphasis{delta} is the time
elapsed since the fix was specified.  The velocity of the wall is set
to the derivative of this expression.

\sphinxAtStartPar
For the \sphinxstyleemphasis{shear} keyword, the wall moves continuously in the specified
dimension with velocity \sphinxstyleemphasis{vshear}.  The dimension must be tangential to
walls with a planar \sphinxstyleemphasis{wallstyle}, e.g. in the \sphinxstyleemphasis{y} or \sphinxstyleemphasis{z} directions for
an \sphinxstyleemphasis{xplane} wall.  For \sphinxstyleemphasis{zcylinder} walls, a dimension of \sphinxstyleemphasis{z} means the
cylinder is moving in the z\sphinxhyphen{}direction along it’s axis.  A dimension of
\sphinxstyleemphasis{x} or \sphinxstyleemphasis{y} means the cylinder is spinning around the z\sphinxhyphen{}axis, either in
the clockwise direction for \sphinxstyleemphasis{vshear} \textgreater{} 0 or counter\sphinxhyphen{}clockwise for
\sphinxstyleemphasis{vshear} \textless{} 0.  In this case, \sphinxstyleemphasis{vshear} is the tangential velocity of
the wall at whatever \sphinxstyleemphasis{radius} has been defined.

\sphinxAtStartPar
The \sphinxstyleemphasis{temperature} keyword is used to assign a temperature to the wall.
The following value can either be a numeric value or an equal\sphinxhyphen{}style
{\hyperref[\detokenize{variable::doc}]{\sphinxcrossref{\DUrole{doc}{variable}}}}.  If the value is a variable, it should be
specified as v\_name, where name is the variable name.  In this case, the
variable will be evaluated each timestep, and its value used to determine
the temperature. This option must be used in conjunction with a heat
conduction model defined in {\hyperref[\detokenize{pair_granular::doc}]{\sphinxcrossref{\DUrole{doc}{pair\_style granular}}}},
{\hyperref[\detokenize{fix_property_atom::doc}]{\sphinxcrossref{\DUrole{doc}{fix property/atom}}}} to store temperature and a
heat flow, and {\hyperref[\detokenize{fix_heat_flow::doc}]{\sphinxcrossref{\DUrole{doc}{fix heat/flow}}}} to integrate heat
flow.


\bigskip\hrule\bigskip


\sphinxAtStartPar
Styles with a \sphinxstyleemphasis{gpu}, \sphinxstyleemphasis{intel}, \sphinxstyleemphasis{kk}, \sphinxstyleemphasis{omp}, or \sphinxstyleemphasis{opt} suffix are
functionally the same as the corresponding style without the suffix.
They have been optimized to run faster, depending on your available
hardware, as discussed on the {\hyperref[\detokenize{Speed_packages::doc}]{\sphinxcrossref{\DUrole{doc}{Accelerator packages}}}}
page.  The accelerated styles take the same arguments and should
produce the same results, except for round\sphinxhyphen{}off and precision issues.

\sphinxAtStartPar
These accelerated styles are part of the GPU, INTEL, KOKKOS,
OPENMP, and OPT packages, respectively.  They are only enabled if
LAMMPS was built with those packages.  See the {\hyperref[\detokenize{Build_package::doc}]{\sphinxcrossref{\DUrole{doc}{Build package}}}} page for more info.

\sphinxAtStartPar
You can specify the accelerated styles explicitly in your input script
by including their suffix, or you can use the {\hyperref[\detokenize{Run_options::doc}]{\sphinxcrossref{\DUrole{doc}{\sphinxhyphen{}suffix command\sphinxhyphen{}line switch}}}} when you invoke LAMMPS, or you can use the
{\hyperref[\detokenize{suffix::doc}]{\sphinxcrossref{\DUrole{doc}{suffix}}}} command in your input script.

\sphinxAtStartPar
See the {\hyperref[\detokenize{Speed_packages::doc}]{\sphinxcrossref{\DUrole{doc}{Accelerator packages}}}} page for more
instructions on how to use the accelerated styles effectively.


\bigskip\hrule\bigskip



\subsection{Restart, fix\_modify, output, run start/stop, minimize info}
\label{\detokenize{fix_wall_gran:restart-fix-modify-output-run-start-stop-minimize-info}}
\sphinxAtStartPar
This fix writes the shear friction state of atoms interacting with the
wall to {\hyperref[\detokenize{restart::doc}]{\sphinxcrossref{\DUrole{doc}{binary restart files}}}}, so that a simulation can
continue correctly if granular potentials with shear “history” effects
are being used.  See the {\hyperref[\detokenize{read_restart::doc}]{\sphinxcrossref{\DUrole{doc}{read\_restart}}}} command for
info on how to re\sphinxhyphen{}specify a fix in an input script that reads a
restart file, so that the operation of the fix continues in an
uninterrupted fashion.

\sphinxAtStartPar
If the \sphinxcode{\sphinxupquote{contacts}} option is used, this fix generates a per\sphinxhyphen{}atom array
with 8 columns as output, containing the contact information for owned
particles (nlocal on each processor). All columns in this per\sphinxhyphen{}atom array will
be zero if no contact has occurred.  The values of these columns are listed in
the following table:


\begin{savenotes}\sphinxattablestart
\sphinxthistablewithglobalstyle
\centering
\begin{tabulary}{\linewidth}[t]{TTT}
\sphinxtoprule
\sphinxstyletheadfamily 
\sphinxAtStartPar
Index
&\sphinxstyletheadfamily 
\sphinxAtStartPar
Value
&\sphinxstyletheadfamily 
\sphinxAtStartPar
Units
\\
\sphinxmidrule
\sphinxtableatstartofbodyhook
\sphinxAtStartPar
1
&
\sphinxAtStartPar
1.0 if particle is in contact with wall,
0.0 otherwise
&\\
\sphinxhline
\sphinxAtStartPar
2
&
\sphinxAtStartPar
Force \(f_x\) exerted by the wall
&
\sphinxAtStartPar
force units
\\
\sphinxhline
\sphinxAtStartPar
3
&
\sphinxAtStartPar
Force \(f_y\) exerted by the wall
&
\sphinxAtStartPar
force units
\\
\sphinxhline
\sphinxAtStartPar
4
&
\sphinxAtStartPar
Force \(f_z\) exerted by the wall
&
\sphinxAtStartPar
force units
\\
\sphinxhline
\sphinxAtStartPar
5
&
\sphinxAtStartPar
\(x\)\sphinxhyphen{}coordinate of contact point on wall
&
\sphinxAtStartPar
distance units
\\
\sphinxhline
\sphinxAtStartPar
6
&
\sphinxAtStartPar
\(y\)\sphinxhyphen{}coordinate of contact point on wall
&
\sphinxAtStartPar
distance units
\\
\sphinxhline
\sphinxAtStartPar
7
&
\sphinxAtStartPar
\(z\)\sphinxhyphen{}coordinate of contact point on wall
&
\sphinxAtStartPar
distance units
\\
\sphinxhline
\sphinxAtStartPar
8
&
\sphinxAtStartPar
Radius \(r\) of atom
&
\sphinxAtStartPar
distance units
\\
\sphinxbottomrule
\end{tabulary}
\sphinxtableafterendhook\par
\sphinxattableend\end{savenotes}

\sphinxAtStartPar
None of the {\hyperref[\detokenize{fix_modify::doc}]{\sphinxcrossref{\DUrole{doc}{fix\_modify}}}} options are relevant to this fix.
No parameter of this fix can be used with the \sphinxstyleemphasis{start/stop} keywords of the
{\hyperref[\detokenize{run::doc}]{\sphinxcrossref{\DUrole{doc}{run}}}} command. This fix is not invoked during {\hyperref[\detokenize{minimize::doc}]{\sphinxcrossref{\DUrole{doc}{energy
minimization}}}}.


\subsection{Restrictions}
\label{\detokenize{fix_wall_gran:restrictions}}
\sphinxAtStartPar
This fix is part of the GRANULAR package.  It is only enabled if
LAMMPS was built with that package.  See the {\hyperref[\detokenize{Build_package::doc}]{\sphinxcrossref{\DUrole{doc}{Build package}}}} page for more info.

\sphinxAtStartPar
Any dimension (xyz) that has a granular wall must be non\sphinxhyphen{}periodic.


\subsection{Related commands}
\label{\detokenize{fix_wall_gran:related-commands}}
\sphinxAtStartPar
{\hyperref[\detokenize{fix_move::doc}]{\sphinxcrossref{\DUrole{doc}{fix move}}}},
{\hyperref[\detokenize{fix_wall_gran_region::doc}]{\sphinxcrossref{\DUrole{doc}{fix wall/gran/region}}}},
{\hyperref[\detokenize{pair_gran::doc}]{\sphinxcrossref{\DUrole{doc}{pair\_style gran/*}}}}
{\hyperref[\detokenize{pair_granular::doc}]{\sphinxcrossref{\DUrole{doc}{pair\_style granular}}}}


\subsection{Default}
\label{\detokenize{fix_wall_gran:default}}
\sphinxAtStartPar
none

\sphinxstepscope

\index{fix wall/gran/region@\spxentry{fix wall/gran/region}}\ignorespaces 

\section{fix wall/gran/region command}
\label{\detokenize{fix_wall_gran_region:fix-wall-gran-region-command}}\label{\detokenize{fix_wall_gran_region:index-0}}\label{\detokenize{fix_wall_gran_region::doc}}

\subsection{Syntax}
\label{\detokenize{fix_wall_gran_region:syntax}}
\begin{sphinxVerbatim}[commandchars=\\\{\}]
\PYG{k}{fix }\PYG{n+nv+nvIdentifier}{ID}\PYG{+w}{ }\PYG{n+nv+nvIdentifier}{group\PYGZhy{}ID}\PYG{+w}{ }\PYG{n}{wall}\PYG{o}{/}\PYG{n}{gran}\PYG{o}{/}\PYG{k}{region }\PYG{n+nv+nvIdentifier}{fstyle}\PYG{+w}{ }\PYG{n}{fstyle\PYGZus{}params}\PYG{+w}{ }\PYG{n}{wallstyle}\PYG{+w}{ }\PYG{n}{regionID}\PYG{+w}{ }\PYG{n}{keyword}\PYG{+w}{ }\PYG{n}{values}\PYG{+w}{ }\PYG{n}{...}
\end{sphinxVerbatim}
\begin{itemize}
\item {} 
\sphinxAtStartPar
ID, group\sphinxhyphen{}ID are documented in {\hyperref[\detokenize{fix::doc}]{\sphinxcrossref{\DUrole{doc}{fix}}}} command

\item {} 
\sphinxAtStartPar
wall/region = style name of this fix command

\item {} 
\sphinxAtStartPar
fstyle = style of force interactions between particles and wall

\begin{sphinxVerbatim}[commandchars=\\\{\}]
possible choices: hooke, hooke/history, hertz/history, granular
\end{sphinxVerbatim}

\item {} 
\sphinxAtStartPar
fstyle\_params = parameters associated with force interaction style
\begin{sphinxalltt}
For \sphinxstyleemphasis{hooke}, \sphinxstyleemphasis{hooke/history}, and \sphinxstyleemphasis{hertz/history}, \sphinxstyleemphasis{fstyle\_params} are:
      Kn = elastic constant for normal particle repulsion (force/distance units or pressure units \sphinxhyphen{} see discussion below)
      Kt = elastic constant for tangential contact (force/distance units or pressure units \sphinxhyphen{} see discussion below)
      gamma\_n = damping coefficient for collisions in normal direction (1/time units or 1/time\sphinxhyphen{}distance units \sphinxhyphen{} see discussion below)
      gamma\_t = damping coefficient for collisions in tangential direction (1/time units or 1/time\sphinxhyphen{}distance units \sphinxhyphen{} see discussion below)
      xmu = static yield criterion (unitless value between 0.0 and 1.0e4)
      dampflag = 0 or 1 if tangential damping force is excluded or included
\end{sphinxalltt}
\begin{sphinxalltt}
For \sphinxstyleemphasis{granular}, \sphinxstyleemphasis{fstyle\_params} are set using the same syntax as for the \sphinxstyleemphasis{pair\_coeff} command of {\hyperref[\detokenize{pair_granular::doc}]{\sphinxcrossref{\DUrole{doc}{pair\_style granular}}}}
\end{sphinxalltt}

\item {} 
\sphinxAtStartPar
wallstyle = region (see {\hyperref[\detokenize{fix_wall_gran::doc}]{\sphinxcrossref{\DUrole{doc}{fix wall/gran}}}} for options for other kinds of walls)

\item {} 
\sphinxAtStartPar
region\sphinxhyphen{}ID = region whose boundary will act as wall

\item {} 
\sphinxAtStartPar
keyword = \sphinxstyleemphasis{contacts} or \sphinxstyleemphasis{temperature}
\begin{sphinxalltt}
\sphinxstyleemphasis{contacts} value = none
   generate contact information for each particle
\sphinxstyleemphasis{temperature} value = temperature
   specify temperature of wall
\end{sphinxalltt}

\end{itemize}


\subsection{Examples}
\label{\detokenize{fix_wall_gran_region:examples}}
\begin{sphinxVerbatim}[commandchars=\\\{\}]
\PYG{k}{fix }\PYG{n+nv+nvIdentifier}{wall}\PYG{+w}{ }\PYG{n+nv+nvIdentifier}{all}\PYG{+w}{ }\PYG{n}{wall}\PYG{o}{/}\PYG{n}{gran}\PYG{o}{/}\PYG{k}{region }\PYG{n+nv+nvIdentifier}{hooke}\PYG{o}{/}\PYG{n}{history}\PYG{+w}{ }\PYG{l+m}{1000.0}\PYG{+w}{ }\PYG{l+m}{200.0}\PYG{+w}{ }\PYG{l+m}{200.0}\PYG{+w}{ }\PYG{l+m}{100.0}\PYG{+w}{ }\PYG{l+m}{0.5}\PYG{+w}{ }\PYG{l+m}{1}\PYG{+w}{ }\PYG{k}{region }\PYG{n+nv+nvIdentifier}{myCone}
\PYG{k}{fix }\PYG{n+nv+nvIdentifier}{3}\PYG{+w}{ }\PYG{n+nv+nvIdentifier}{all}\PYG{+w}{ }\PYG{n}{wall}\PYG{o}{/}\PYG{n}{gran}\PYG{o}{/}\PYG{k}{region }\PYG{n+nv+nvIdentifier}{granular}\PYG{+w}{ }\PYG{n}{hooke}\PYG{+w}{ }\PYG{l+m}{1000.0}\PYG{+w}{ }\PYG{l+m}{50.0}\PYG{+w}{ }\PYG{n}{tangential}\PYG{+w}{ }\PYG{n}{linear\PYGZus{}nohistory}\PYG{+w}{ }\PYG{l+m}{1.0}\PYG{+w}{ }\PYG{l+m}{0.4}\PYG{+w}{ }\PYG{n}{damping}\PYG{+w}{ }\PYG{k}{velocity }\PYG{n+nv+nvIdentifier}{region}\PYG{+w}{ }\PYG{n}{myBox}
\PYG{k}{fix }\PYG{n+nv+nvIdentifier}{4}\PYG{+w}{ }\PYG{n+nv+nvIdentifier}{all}\PYG{+w}{ }\PYG{n}{wall}\PYG{o}{/}\PYG{n}{gran}\PYG{o}{/}\PYG{k}{region }\PYG{n+nv+nvIdentifier}{granular}\PYG{+w}{ }\PYG{n}{jkr}\PYG{+w}{ }\PYG{l+m}{1e5}\PYG{+w}{ }\PYG{l+m}{1500.0}\PYG{+w}{ }\PYG{l+m}{0.3}\PYG{+w}{ }\PYG{l+m}{10.0}\PYG{+w}{ }\PYG{n}{tangential}\PYG{+w}{ }\PYG{n}{mindlin}\PYG{+w}{ }\PYG{n}{NULL}\PYG{+w}{ }\PYG{l+m}{1.0}\PYG{+w}{ }\PYG{l+m}{0.5}\PYG{+w}{ }\PYG{n}{rolling}\PYG{+w}{ }\PYG{n}{sds}\PYG{+w}{ }\PYG{l+m}{500.0}\PYG{+w}{ }\PYG{l+m}{200.0}\PYG{+w}{ }\PYG{l+m}{0.5}\PYG{+w}{ }\PYG{n}{twisting}\PYG{+w}{ }\PYG{n}{marshall}\PYG{+w}{ }\PYG{k}{region }\PYG{n+nv+nvIdentifier}{myCone}
\PYG{k}{fix }\PYG{n+nv+nvIdentifier}{5}\PYG{+w}{ }\PYG{n+nv+nvIdentifier}{all}\PYG{+w}{ }\PYG{n}{wall}\PYG{o}{/}\PYG{n}{gran}\PYG{o}{/}\PYG{k}{region }\PYG{n+nv+nvIdentifier}{granular}\PYG{+w}{ }\PYG{n}{dmt}\PYG{+w}{ }\PYG{l+m}{1e5}\PYG{+w}{ }\PYG{l+m}{0.2}\PYG{+w}{ }\PYG{l+m}{0.3}\PYG{+w}{ }\PYG{l+m}{10.0}\PYG{+w}{ }\PYG{n}{tangential}\PYG{+w}{ }\PYG{n}{mindlin}\PYG{+w}{ }\PYG{n}{NULL}\PYG{+w}{ }\PYG{l+m}{1.0}\PYG{+w}{ }\PYG{l+m}{0.5}\PYG{+w}{ }\PYG{n}{rolling}\PYG{+w}{ }\PYG{n}{sds}\PYG{+w}{ }\PYG{l+m}{500.0}\PYG{+w}{ }\PYG{l+m}{200.0}\PYG{+w}{ }\PYG{l+m}{0.5}\PYG{+w}{ }\PYG{n}{twisting}\PYG{+w}{ }\PYG{n}{marshall}\PYG{+w}{ }\PYG{n}{damping}\PYG{+w}{ }\PYG{n}{tsuji}\PYG{+w}{ }\PYG{k}{region }\PYG{n+nv+nvIdentifier}{myCone}
\PYG{k}{fix }\PYG{n+nv+nvIdentifier}{wall}\PYG{+w}{ }\PYG{n+nv+nvIdentifier}{all}\PYG{+w}{ }\PYG{n}{wall}\PYG{o}{/}\PYG{n}{gran}\PYG{o}{/}\PYG{k}{region }\PYG{n+nv+nvIdentifier}{hooke}\PYG{o}{/}\PYG{n}{history}\PYG{+w}{ }\PYG{l+m}{1000.0}\PYG{+w}{ }\PYG{l+m}{200.0}\PYG{+w}{ }\PYG{l+m}{200.0}\PYG{+w}{ }\PYG{l+m}{100.0}\PYG{+w}{ }\PYG{l+m}{0.5}\PYG{+w}{ }\PYG{l+m}{1}\PYG{+w}{ }\PYG{k}{region }\PYG{n+nv+nvIdentifier}{myCone}\PYG{+w}{ }\PYG{n}{contacts}
\end{sphinxVerbatim}


\subsection{Description}
\label{\detokenize{fix_wall_gran_region:description}}
\sphinxAtStartPar
Treat the surface of the geometric region defined by the \sphinxstyleemphasis{region\sphinxhyphen{}ID}
as a bounding frictional wall which interacts with nearby finite\sphinxhyphen{}size
granular particles when they are close enough to touch the wall.  See
the {\hyperref[\detokenize{fix_wall_region::doc}]{\sphinxcrossref{\DUrole{doc}{fix wall/region}}}} and {\hyperref[\detokenize{fix_wall_gran::doc}]{\sphinxcrossref{\DUrole{doc}{fix wall/gran}}}} commands for related kinds of walls for
non\sphinxhyphen{}granular particles and simpler wall geometries, respectively.

\sphinxAtStartPar
Here are snapshots of example models using this command.  Corresponding
input scripts can be found in examples/granregion.  Movies of these
simulations are \sphinxhref{https://www.lammps.org/movies.html\#granregion}{here on the Movies page}
of the LAMMPS website.

\sphinxAtStartPar
\sphinxincludegraphics[width=0.480\linewidth]{{gran_funnel}.png}  \sphinxincludegraphics[width=0.480\linewidth]{{gran_mixer}.png}




\bigskip\hrule\bigskip


\sphinxAtStartPar
The distance between a particle and the region boundary is the
distance to the nearest point on the region surface.  The force the
wall exerts on the particle is along the direction between that point
and the particle center, which is the direction normal to the surface
at that point.  Note that if the region surface is comprised of
multiple “faces”, then each face can exert a force on the particle if
it is close enough.  E.g. for {\hyperref[\detokenize{region::doc}]{\sphinxcrossref{\DUrole{doc}{region\_style block}}}}, a
particle in the interior, near a corner of the block, could feel wall
forces from 1, 2, or 3 faces of the block.

\sphinxAtStartPar
Regions are defined using the {\hyperref[\detokenize{region::doc}]{\sphinxcrossref{\DUrole{doc}{region}}}} command.  Note that
the region volume can be interior or exterior to the bounding surface,
which will determine in which direction the surface interacts with
particles, i.e. the direction of the surface normal. The exception to
this is if one or more \sphinxstyleemphasis{open} options are specified for the region
command, in which case particles interact with both the interior and
exterior surfaces of regions.

\sphinxAtStartPar
Regions can either be primitive shapes (block, sphere, cylinder, etc)
or combinations of primitive shapes specified via the \sphinxstyleemphasis{union} or
\sphinxstyleemphasis{intersect} region styles.  These latter styles can be used to
construct particle containers with complex shapes.

\sphinxAtStartPar
Regions can also move dynamically via the {\hyperref[\detokenize{region::doc}]{\sphinxcrossref{\DUrole{doc}{region}}}} command
keywords (move) and \sphinxstyleemphasis{rotate}, or change their shape by use of variables
as inputs to the {\hyperref[\detokenize{region::doc}]{\sphinxcrossref{\DUrole{doc}{region}}}} command.  If such a region is used
with this fix, then the region surface will move in time in the
corresponding manner.

\begin{sphinxadmonition}{note}{Note:}
\sphinxAtStartPar
As discussed on the {\hyperref[\detokenize{region::doc}]{\sphinxcrossref{\DUrole{doc}{region}}}} command doc page,
regions in LAMMPS do not get wrapped across periodic boundaries.  It
is up to you to ensure that the region location with respect to
periodic or non\sphinxhyphen{}periodic boundaries is specified appropriately via the
{\hyperref[\detokenize{region::doc}]{\sphinxcrossref{\DUrole{doc}{region}}}} and {\hyperref[\detokenize{boundary::doc}]{\sphinxcrossref{\DUrole{doc}{boundary}}}} commands when using
a region as a wall that bounds particle motion.
\end{sphinxadmonition}

\begin{sphinxadmonition}{note}{Note:}
\sphinxAtStartPar
For primitive regions with sharp corners and/or edges (e.g. a
block or cylinder), wall/particle forces are computed accurately for
both interior and exterior regions.  For \sphinxstyleemphasis{union} and \sphinxstyleemphasis{intersect}
regions, additional sharp corners and edges may be present due to the
intersection of the surfaces of 2 or more primitive volumes.  These
corners and edges can be of two types: concave or convex.  Concave
points/edges are like the corners of a cube as seen by particles in
the interior of a cube.  Wall/particle forces around these features
are computed correctly.  Convex points/edges are like the corners of a
cube as seen by particles exterior to the cube, i.e. the points jut
into the volume where particles are present.  LAMMPS does NOT compute
the location of these convex points directly, and hence wall/particle
forces in the cutoff volume around these points suffer from
inaccuracies.  The basic problem is that the outward normal of the
surface is not continuous at these points.  This can cause particles
to feel no force (they don’t “see” the wall) when in one location,
then move a distance epsilon, and suddenly feel a large force because
they now “see” the wall.  In a worst\sphinxhyphen{}case scenario, this can blow
particles out of the simulation box.  Thus, as a general rule you
should not use the fix wall/gran/region command with \sphinxstyleemphasis{union} or
\sphinxstyleemphasis{interesect} regions that have convex points or edges resulting from
the union/intersection (convex points/edges in the union/intersection
due to a single sub\sphinxhyphen{}region are still OK).
\end{sphinxadmonition}

\begin{sphinxadmonition}{note}{Note:}
\sphinxAtStartPar
Similarly, you should not define \sphinxstyleemphasis{union} or \sphinxstyleemphasis{intersert} regions
for use with this command that share an overlapping common face that
is part of the overall outer boundary (interior boundary is OK), even
if the face is smooth.  E.g. two regions of style block in a \sphinxstyleemphasis{union}
region, where the two blocks overlap on one or more of their faces.
This is because LAMMPS discards points that are part of multiple
sub\sphinxhyphen{}regions when calculating wall/particle interactions, to avoid
double\sphinxhyphen{}counting the interaction.  Having two coincident faces could
cause the face to become invisible to the particles.  The solution is
to make the two faces differ by epsilon in their position.
\end{sphinxadmonition}

\sphinxAtStartPar
The nature of the wall/particle interactions are determined by the
\sphinxstyleemphasis{fstyle} setting.  It can be any of the styles defined by the
{\hyperref[\detokenize{pair_gran::doc}]{\sphinxcrossref{\DUrole{doc}{pair\_style gran/*}}}} or the more general
{\hyperref[\detokenize{pair_granular::doc}]{\sphinxcrossref{\DUrole{doc}{pair\_style granular}}}} commands.  Currently the
options are \sphinxstyleemphasis{hooke}, \sphinxstyleemphasis{hooke/history}, or \sphinxstyleemphasis{hertz/history} for the
former, and \sphinxstyleemphasis{granular} with all the possible options of the associated
\sphinxstyleemphasis{pair\_coeff} command for the latter.  The equation for the force
between the wall and particles touching it is the same as the
corresponding equation on the {\hyperref[\detokenize{pair_gran::doc}]{\sphinxcrossref{\DUrole{doc}{pair\_style gran/*}}}} and
{\hyperref[\detokenize{pair_granular::doc}]{\sphinxcrossref{\DUrole{doc}{pair\_style granular}}}} doc pages, but the effective
radius is calculated using the radius of the particle and the radius of
curvature of the wall at the contact point.

\sphinxAtStartPar
Specifically, delta = radius \sphinxhyphen{} r = overlap of particle with wall,
m\_eff = mass of particle, and RiRj/Ri+Rj is the effective radius, with
Rj replaced by the radius of curvature of the wall at the contact
point.  The radius of curvature can be negative for a concave wall
section, e.g. the interior of cylinder.  For a flat wall, delta =
radius \sphinxhyphen{} r = overlap of particle with wall, m\_eff = mass of particle,
and the effective radius of contact is just the radius of the
particle.

\sphinxAtStartPar
The parameters \sphinxstyleemphasis{Kn}, \sphinxstyleemphasis{Kt}, \sphinxstyleemphasis{gamma\_n}, \sphinxstyleemphasis{gamma\_t}, \sphinxstyleemphasis{xmu}, \sphinxstyleemphasis{dampflag},
and the optional keyword \sphinxstyleemphasis{limit\_damping}
have the same meaning and units as those specified with the
{\hyperref[\detokenize{pair_gran::doc}]{\sphinxcrossref{\DUrole{doc}{pair\_style gran/*}}}} commands.  This means a NULL can be
used for either \sphinxstyleemphasis{Kt} or \sphinxstyleemphasis{gamma\_t} as described on that page.  If a
NULL is used for \sphinxstyleemphasis{Kt}, then a default value is used where \sphinxstyleemphasis{Kt} = 2/7
\sphinxstyleemphasis{Kn}.  If a NULL is used for \sphinxstyleemphasis{gamma\_t}, then a default value is used
where \sphinxstyleemphasis{gamma\_t} = 1/2 \sphinxstyleemphasis{gamma\_n}.

\sphinxAtStartPar
All the model choices for cohesion, tangential friction, rolling
friction and twisting friction supported by the {\hyperref[\detokenize{pair_granular::doc}]{\sphinxcrossref{\DUrole{doc}{pair\_style granular}}}} through its \sphinxstyleemphasis{pair\_coeff} command are also
supported for walls. These are discussed in greater detail on the doc
page for {\hyperref[\detokenize{pair_granular::doc}]{\sphinxcrossref{\DUrole{doc}{pair\_style granular}}}}.

\sphinxAtStartPar
Note that you can choose a different force styles and/or different
values for the 6 wall/particle coefficients than for particle/particle
interactions.  E.g. if you wish to model the wall as a different
material.

\sphinxAtStartPar
The \sphinxstyleemphasis{temperature} keyword is used to assign a temperature to the wall.
The following value can either be a numeric value or an equal\sphinxhyphen{}style
{\hyperref[\detokenize{variable::doc}]{\sphinxcrossref{\DUrole{doc}{variable}}}}.  If the value is a variable, it should be
specified as v\_name, where name is the variable name.  In this case, the
variable will be evaluated each timestep, and its value used to determine
the temperature. This option must be used in conjunction with a heat
conduction model defined in {\hyperref[\detokenize{pair_granular::doc}]{\sphinxcrossref{\DUrole{doc}{pair\_style granular}}}},
{\hyperref[\detokenize{fix_property_atom::doc}]{\sphinxcrossref{\DUrole{doc}{fix property/atom}}}} to store temperature and a
heat flow, and {\hyperref[\detokenize{fix_heat_flow::doc}]{\sphinxcrossref{\DUrole{doc}{fix heat/flow}}}} to integrate heat
flow.


\subsection{Restart, fix\_modify, output, run start/stop, minimize info}
\label{\detokenize{fix_wall_gran_region:restart-fix-modify-output-run-start-stop-minimize-info}}
\sphinxAtStartPar
Similar to {\hyperref[\detokenize{fix_wall_gran::doc}]{\sphinxcrossref{\DUrole{doc}{fix wall/gran}}}} command, this fix writes
the shear friction state of atoms interacting with the wall to {\hyperref[\detokenize{restart::doc}]{\sphinxcrossref{\DUrole{doc}{binary restart files}}}}, so that a simulation can continue
correctly if granular potentials with shear “history” effects are
being used.  This fix also includes info about a moving region in the
restart file.  See the {\hyperref[\detokenize{read_restart::doc}]{\sphinxcrossref{\DUrole{doc}{read\_restart}}}} command for
info on how to re\sphinxhyphen{}specify a fix in an input script that reads a
restart file, so that the operation of the fix continues in an
uninterrupted fashion.

\begin{sphinxadmonition}{note}{Note:}
\sphinxAtStartPar
Information about region definitions is NOT included in restart
files, as discussed on the {\hyperref[\detokenize{read_restart::doc}]{\sphinxcrossref{\DUrole{doc}{read\_restart}}}} doc page.
So you must re\sphinxhyphen{}define your region and if it is a moving region, define
its motion attributes in a way that is consistent with the simulation
that wrote the restart file.  In particular, if you want to change the
region motion attributes (e.g. its velocity), then you should ensure
the position/orientation of the region at the initial restart timestep
is the same as it was on the timestep the restart file was written.
If this is not possible, you may need to ignore info in the restart
file by defining a new fix wall/gran/region command in your restart
script, e.g. with a different fix ID.  Or if you want to keep the
shear history info but discard the region motion information, you can
use the same fix ID for fix wall/gran/region, but assign it a region
with a different region ID.
\end{sphinxadmonition}

\sphinxAtStartPar
If the \sphinxcode{\sphinxupquote{contacts}} option is used, this fix generates a per\sphinxhyphen{}atom array
with 8 columns as output, containing the contact information for owned
particles (nlocal on each processor). All columns in this per\sphinxhyphen{}atom array will
be zero if no contact has occurred. The values of these columns are listed in
the following table:


\begin{savenotes}\sphinxattablestart
\sphinxthistablewithglobalstyle
\centering
\begin{tabulary}{\linewidth}[t]{TTT}
\sphinxtoprule
\sphinxstyletheadfamily 
\sphinxAtStartPar
Index
&\sphinxstyletheadfamily 
\sphinxAtStartPar
Value
&\sphinxstyletheadfamily 
\sphinxAtStartPar
Units
\\
\sphinxmidrule
\sphinxtableatstartofbodyhook
\sphinxAtStartPar
1
&
\sphinxAtStartPar
1.0 if particle is in contact with wall,
0.0 otherwise
&\\
\sphinxhline
\sphinxAtStartPar
2
&
\sphinxAtStartPar
Force \(f_x\) exerted by the wall
&
\sphinxAtStartPar
force units
\\
\sphinxhline
\sphinxAtStartPar
3
&
\sphinxAtStartPar
Force \(f_y\) exerted by the wall
&
\sphinxAtStartPar
force units
\\
\sphinxhline
\sphinxAtStartPar
4
&
\sphinxAtStartPar
Force \(f_z\) exerted by the wall
&
\sphinxAtStartPar
force units
\\
\sphinxhline
\sphinxAtStartPar
5
&
\sphinxAtStartPar
\(x\)\sphinxhyphen{}coordinate of contact point on wall
&
\sphinxAtStartPar
distance units
\\
\sphinxhline
\sphinxAtStartPar
6
&
\sphinxAtStartPar
\(y\)\sphinxhyphen{}coordinate of contact point on wall
&
\sphinxAtStartPar
distance units
\\
\sphinxhline
\sphinxAtStartPar
7
&
\sphinxAtStartPar
\(z\)\sphinxhyphen{}coordinate of contact point on wall
&
\sphinxAtStartPar
distance units
\\
\sphinxhline
\sphinxAtStartPar
8
&
\sphinxAtStartPar
Radius \(r\) of atom
&
\sphinxAtStartPar
distance units
\\
\sphinxbottomrule
\end{tabulary}
\sphinxtableafterendhook\par
\sphinxattableend\end{savenotes}

\sphinxAtStartPar
None of the {\hyperref[\detokenize{fix_modify::doc}]{\sphinxcrossref{\DUrole{doc}{fix\_modify}}}} options are relevant to this fix.
No parameter of this fix can be used with the \sphinxstyleemphasis{start/stop} keywords of the
{\hyperref[\detokenize{run::doc}]{\sphinxcrossref{\DUrole{doc}{run}}}} command. This fix is not invoked during {\hyperref[\detokenize{minimize::doc}]{\sphinxcrossref{\DUrole{doc}{energy
minimization}}}}.


\subsection{Restrictions}
\label{\detokenize{fix_wall_gran_region:restrictions}}
\sphinxAtStartPar
This fix is part of the GRANULAR package.  It is only enabled if
LAMMPS was built with that package.  See the {\hyperref[\detokenize{Build_package::doc}]{\sphinxcrossref{\DUrole{doc}{Build package}}}} page for more info.


\subsection{Related commands}
\label{\detokenize{fix_wall_gran_region:related-commands}}
\sphinxAtStartPar
{\hyperref[\detokenize{fix_move::doc}]{\sphinxcrossref{\DUrole{doc}{fix\_move}}}},
{\hyperref[\detokenize{fix_wall_gran::doc}]{\sphinxcrossref{\DUrole{doc}{fix wall/gran}}}},
{\hyperref[\detokenize{fix_wall_region::doc}]{\sphinxcrossref{\DUrole{doc}{fix wall/region}}}},
{\hyperref[\detokenize{pair_gran::doc}]{\sphinxcrossref{\DUrole{doc}{pair\_style granular}}}},
{\hyperref[\detokenize{region::doc}]{\sphinxcrossref{\DUrole{doc}{region}}}}


\subsection{Default}
\label{\detokenize{fix_wall_gran_region:default}}
\sphinxAtStartPar
none

\sphinxstepscope

\index{fix wall/piston@\spxentry{fix wall/piston}}\ignorespaces 

\section{fix wall/piston command}
\label{\detokenize{fix_wall_piston:fix-wall-piston-command}}\label{\detokenize{fix_wall_piston:index-0}}\label{\detokenize{fix_wall_piston::doc}}

\subsection{Syntax}
\label{\detokenize{fix_wall_piston:syntax}}
\begin{sphinxVerbatim}[commandchars=\\\{\}]
\PYG{k}{fix }\PYG{n+nv+nvIdentifier}{ID}\PYG{+w}{ }\PYG{n+nv+nvIdentifier}{group\PYGZhy{}ID}\PYG{+w}{ }\PYG{n}{wall}\PYG{o}{/}\PYG{n}{piston}\PYG{+w}{ }\PYG{n}{face}\PYG{+w}{ }\PYG{n}{...}\PYG{+w}{ }\PYG{n}{keyword}\PYG{+w}{ }\PYG{n}{value}\PYG{+w}{ }\PYG{n}{...}
\end{sphinxVerbatim}
\begin{itemize}
\item {} 
\sphinxAtStartPar
ID, group\sphinxhyphen{}ID are documented in {\hyperref[\detokenize{fix::doc}]{\sphinxcrossref{\DUrole{doc}{fix}}}} command

\item {} 
\sphinxAtStartPar
wall/piston = style name of this fix command

\item {} 
\sphinxAtStartPar
face = \sphinxstyleemphasis{zlo}

\item {} 
\sphinxAtStartPar
zero or more keyword/value pairs may be appended

\item {} 
\sphinxAtStartPar
keyword = \sphinxstyleemphasis{pos} or \sphinxstyleemphasis{vel} or \sphinxstyleemphasis{ramp} or \sphinxstyleemphasis{temp} or \sphinxstyleemphasis{units}
\begin{sphinxalltt}
\sphinxstyleemphasis{pos} args = z
  z = z coordinate at which the piston begins (distance units)
\sphinxstyleemphasis{vel} args = vz
  vz = final velocity of the piston (velocity units)
\sphinxstyleemphasis{ramp} = use a linear velocity ramp from 0 to vz
\sphinxstyleemphasis{temp} args = target damp seed extent
  target = target velocity for region immediately ahead of the piston
  damp = damping parameter (time units)
  seed = random number seed for langevin kicks
  extent = extent of thermostatted region (distance units)
\sphinxstyleemphasis{units} value = \sphinxstyleemphasis{lattice} or \sphinxstyleemphasis{box}
  \sphinxstyleemphasis{lattice} = the wall position is defined in lattice units
  \sphinxstyleemphasis{box} = the wall position is defined in simulation box units
\end{sphinxalltt}

\end{itemize}


\subsection{Examples}
\label{\detokenize{fix_wall_piston:examples}}
\begin{sphinxVerbatim}[commandchars=\\\{\}]
\PYG{k}{fix }\PYG{n+nv+nvIdentifier}{xwalls}\PYG{+w}{ }\PYG{n+nv+nvIdentifier}{all}\PYG{+w}{ }\PYG{n}{wall}\PYG{o}{/}\PYG{n}{piston}\PYG{+w}{ }\PYG{n}{zlo}
\PYG{k}{fix }\PYG{n+nv+nvIdentifier}{walls}\PYG{+w}{ }\PYG{n+nv+nvIdentifier}{all}\PYG{+w}{ }\PYG{n}{wall}\PYG{o}{/}\PYG{n}{piston}\PYG{+w}{ }\PYG{n}{zlo}\PYG{+w}{ }\PYG{n}{pos}\PYG{+w}{ }\PYG{l+m}{1.0}\PYG{+w}{ }\PYG{n}{vel}\PYG{+w}{ }\PYG{l+m}{10.0}\PYG{+w}{ }\PYG{n}{units}\PYG{+w}{ }\PYG{n}{box}
\PYG{k}{fix }\PYG{n+nv+nvIdentifier}{top}\PYG{+w}{ }\PYG{n+nv+nvIdentifier}{all}\PYG{+w}{ }\PYG{n}{wall}\PYG{o}{/}\PYG{n}{piston}\PYG{+w}{ }\PYG{n}{zlo}\PYG{+w}{ }\PYG{n}{vel}\PYG{+w}{ }\PYG{l+m}{10.0}\PYG{+w}{ }\PYG{n}{ramp}
\end{sphinxVerbatim}


\subsection{Description}
\label{\detokenize{fix_wall_piston:description}}
\sphinxAtStartPar
Bound the simulation with a moving wall which reflect particles in the
specified group and drive the system with an effective infinite\sphinxhyphen{}mass
piston capable of driving shock waves.

\sphinxAtStartPar
A momentum mirror technique is used, which means that if an atom (or
the wall) moves such that an atom is outside the wall on a timestep by
a distance delta (e.g. due to {\hyperref[\detokenize{fix_nve::doc}]{\sphinxcrossref{\DUrole{doc}{fix nve}}}}), then it is put
back inside the face by the same delta, and the velocity relative to
the moving wall is flipped in z.  For instance, a stationary particle
hit with a piston wall with velocity vz, will end the timestep with a
velocity of 2*vz.

\sphinxAtStartPar
Currently the \sphinxstyleemphasis{face} keyword can only be \sphinxstyleemphasis{zlo}.  This creates a piston
moving in the positive z direction.  Particles with z coordinate less
than the wall position are reflected to a z coordinate greater than
the wall position.  If the piston velocity is vpz and the particle
velocity before reflection is vzi, the particle velocity after
reflection is \sphinxhyphen{}vzi + 2*vpz.

\sphinxAtStartPar
The initial position of the wall can be specified by the \sphinxstyleemphasis{pos} keyword.

\sphinxAtStartPar
The final velocity of the wall can be specified by the \sphinxstyleemphasis{vel} keyword

\sphinxAtStartPar
The \sphinxstyleemphasis{ramp} keyword will cause the wall/piston to adjust the velocity
linearly from zero velocity to \sphinxstyleemphasis{vel} over the course of the run. If
the \sphinxstyleemphasis{ramp} keyword is omitted then the wall/piston moves at a constant
velocity defined by \sphinxstyleemphasis{vel}.

\sphinxAtStartPar
The \sphinxstyleemphasis{temp} keyword will cause the region immediately in front of the
wall/piston to be thermostatted with a Langevin thermostat.  This
region moves with the piston.  The damping and kicking are measured in
the reference frame of the piston.  So, a temperature of zero would
mean all particles were moving at exactly the speed of the
wall/piston.

\sphinxAtStartPar
The \sphinxstyleemphasis{units} keyword determines the meaning of the distance units used
to define a wall position, but only when a numeric constant is used.

\sphinxAtStartPar
A \sphinxstyleemphasis{box} value selects standard distance units as defined by the
{\hyperref[\detokenize{units::doc}]{\sphinxcrossref{\DUrole{doc}{units}}}} command, e.g. Angstroms for units = real or metal.
A \sphinxstyleemphasis{lattice} value means the distance units are in lattice spacings.
The {\hyperref[\detokenize{lattice::doc}]{\sphinxcrossref{\DUrole{doc}{lattice}}}} command must have been previously used to
define the lattice spacings.


\bigskip\hrule\bigskip



\subsection{Restart, fix\_modify, output, run start/stop, minimize info}
\label{\detokenize{fix_wall_piston:restart-fix-modify-output-run-start-stop-minimize-info}}
\sphinxAtStartPar
No information about this fix is written to {\hyperref[\detokenize{restart::doc}]{\sphinxcrossref{\DUrole{doc}{binary restart files}}}}.  None of the {\hyperref[\detokenize{fix_modify::doc}]{\sphinxcrossref{\DUrole{doc}{fix\_modify}}}} options
are relevant to this fix.  No global or per\sphinxhyphen{}atom quantities are stored
by this fix for access by various {\hyperref[\detokenize{Howto_output::doc}]{\sphinxcrossref{\DUrole{doc}{output commands}}}}.
No parameter of this fix can be used with the \sphinxstyleemphasis{start/stop} keywords of
the {\hyperref[\detokenize{run::doc}]{\sphinxcrossref{\DUrole{doc}{run}}}} command.  This fix is not invoked during {\hyperref[\detokenize{minimize::doc}]{\sphinxcrossref{\DUrole{doc}{energy minimization}}}}.


\subsection{Restrictions}
\label{\detokenize{fix_wall_piston:restrictions}}
\sphinxAtStartPar
This fix style is part of the SHOCK package.  It is only enabled if
LAMMPS was built with that package. See the {\hyperref[\detokenize{Build_package::doc}]{\sphinxcrossref{\DUrole{doc}{Build package}}}} page for more info.

\sphinxAtStartPar
The face that has the wall/piston must be boundary type ‘s’
(shrink\sphinxhyphen{}wrapped). The opposing face can be
any boundary type other than periodic.

\sphinxAtStartPar
A wall/piston should not be used with rigid bodies such as those
defined by a “fix rigid” command.  This is because the wall/piston
displaces atoms directly rather than exerting a force on them.


\subsection{Related commands}
\label{\detokenize{fix_wall_piston:related-commands}}
\sphinxAtStartPar
{\hyperref[\detokenize{fix_wall::doc}]{\sphinxcrossref{\DUrole{doc}{fix wall/reflect}}}} command, {\hyperref[\detokenize{fix_append_atoms::doc}]{\sphinxcrossref{\DUrole{doc}{fix append/atoms}}}} command


\subsection{Default}
\label{\detokenize{fix_wall_piston:default}}
\sphinxAtStartPar
The keyword defaults are pos = 0, vel = 0, units = lattice.

\sphinxstepscope

\index{fix wall/reflect@\spxentry{fix wall/reflect}}\ignorespaces 
\index{fix wall/reflect/kk@\spxentry{fix wall/reflect/kk}}\ignorespaces 

\section{fix wall/reflect command}
\label{\detokenize{fix_wall_reflect:fix-wall-reflect-command}}\label{\detokenize{fix_wall_reflect:index-1}}\label{\detokenize{fix_wall_reflect:index-0}}\label{\detokenize{fix_wall_reflect::doc}}
\sphinxAtStartPar
Accelerator Variants: \sphinxstyleemphasis{wall/reflect/kk}


\subsection{Syntax}
\label{\detokenize{fix_wall_reflect:syntax}}
\begin{sphinxVerbatim}[commandchars=\\\{\}]
\PYG{k}{fix }\PYG{n+nv+nvIdentifier}{ID}\PYG{+w}{ }\PYG{n+nv+nvIdentifier}{group\PYGZhy{}ID}\PYG{+w}{ }\PYG{n}{wall}\PYG{o}{/}\PYG{n}{reflect}\PYG{+w}{ }\PYG{n}{face}\PYG{+w}{ }\PYG{n}{arg}\PYG{+w}{ }\PYG{n}{...}\PYG{+w}{ }\PYG{n}{keyword}\PYG{+w}{ }\PYG{n}{value}\PYG{+w}{ }\PYG{n}{...}
\end{sphinxVerbatim}
\begin{itemize}
\item {} 
\sphinxAtStartPar
ID, group\sphinxhyphen{}ID are documented in {\hyperref[\detokenize{fix::doc}]{\sphinxcrossref{\DUrole{doc}{fix}}}} command

\item {} 
\sphinxAtStartPar
wall/reflect = style name of this fix command

\item {} 
\sphinxAtStartPar
one or more face/arg pairs may be appended

\item {} 
\sphinxAtStartPar
face = \sphinxstyleemphasis{xlo} or \sphinxstyleemphasis{xhi} or \sphinxstyleemphasis{ylo} or \sphinxstyleemphasis{yhi} or \sphinxstyleemphasis{zlo} or \sphinxstyleemphasis{zhi}
\begin{sphinxalltt}
arg = EDGE or constant or variable
  EDGE = current lo edge of simulation box
  constant = number like 0.0 or 30.0 (distance units)
  variable = {\hyperref[\detokenize{variable::doc}]{\sphinxcrossref{\DUrole{doc}{equal\sphinxhyphen{}style variable}}}} like v\_x or v\_wiggle
\end{sphinxalltt}

\item {} 
\sphinxAtStartPar
zero or more keyword/value pairs may be appended

\item {} 
\sphinxAtStartPar
keyword = \sphinxstyleemphasis{units}
\begin{sphinxalltt}
\sphinxstyleemphasis{units} value = \sphinxstyleemphasis{lattice} or \sphinxstyleemphasis{box}
  \sphinxstyleemphasis{lattice} = the wall position is defined in lattice units
  \sphinxstyleemphasis{box} = the wall position is defined in simulation box units
\end{sphinxalltt}

\end{itemize}


\subsection{Examples}
\label{\detokenize{fix_wall_reflect:examples}}
\begin{sphinxVerbatim}[commandchars=\\\{\}]
\PYG{k}{fix }\PYG{n+nv+nvIdentifier}{xwalls}\PYG{+w}{ }\PYG{n+nv+nvIdentifier}{all}\PYG{+w}{ }\PYG{n}{wall}\PYG{o}{/}\PYG{n}{reflect}\PYG{+w}{ }\PYG{n}{xlo}\PYG{+w}{ }\PYG{n}{EDGE}\PYG{+w}{ }\PYG{n}{xhi}\PYG{+w}{ }\PYG{n}{EDGE}
\PYG{k}{fix }\PYG{n+nv+nvIdentifier}{walls}\PYG{+w}{ }\PYG{n+nv+nvIdentifier}{all}\PYG{+w}{ }\PYG{n}{wall}\PYG{o}{/}\PYG{n}{reflect}\PYG{+w}{ }\PYG{n}{xlo}\PYG{+w}{ }\PYG{l+m}{0.0}\PYG{+w}{ }\PYG{n}{ylo}\PYG{+w}{ }\PYG{l+m}{10.0}\PYG{+w}{ }\PYG{n}{units}\PYG{+w}{ }\PYG{n}{box}
\PYG{k}{fix }\PYG{n+nv+nvIdentifier}{top}\PYG{+w}{ }\PYG{n+nv+nvIdentifier}{all}\PYG{+w}{ }\PYG{n}{wall}\PYG{o}{/}\PYG{n}{reflect}\PYG{+w}{ }\PYG{n}{zhi}\PYG{+w}{ }\PYG{n}{v\PYGZus{}pressdown}
\end{sphinxVerbatim}


\subsection{Description}
\label{\detokenize{fix_wall_reflect:description}}
\sphinxAtStartPar
Bound the simulation with one or more walls which reflect particles
in the specified group when they attempt to move through them.

\sphinxAtStartPar
Reflection means that if an atom moves outside the wall on a timestep
by a distance delta (e.g. due to {\hyperref[\detokenize{fix_nve::doc}]{\sphinxcrossref{\DUrole{doc}{fix nve}}}}), then it is
put back inside the face by the same delta, and the sign of the
corresponding component of its velocity is flipped.

\sphinxAtStartPar
When used in conjunction with {\hyperref[\detokenize{fix_nve::doc}]{\sphinxcrossref{\DUrole{doc}{fix nve}}}} and
{\hyperref[\detokenize{run_style::doc}]{\sphinxcrossref{\DUrole{doc}{run\_style verlet}}}}, the resultant time\sphinxhyphen{}integration
algorithm is equivalent to the primitive splitting algorithm (PSA)
described by {\hyperref[\detokenize{fix_wall_reflect:bond1}]{\sphinxcrossref{\DUrole{std,std-ref}{Bond}}}}.  Because each reflection event
divides the corresponding timestep asymmetrically, energy conservation
is only satisfied to O(dt), rather than to O(dt\textasciicircum{}2) as it would be for
velocity\sphinxhyphen{}Verlet integration without reflective walls.

\sphinxAtStartPar
Up to 6 walls or faces can be specified in a single command: \sphinxstyleemphasis{xlo},
\sphinxstyleemphasis{xhi}, \sphinxstyleemphasis{ylo}, \sphinxstyleemphasis{yhi}, \sphinxstyleemphasis{zlo}, \sphinxstyleemphasis{zhi}.  A \sphinxstyleemphasis{lo} face reflects particles
that move to a coordinate less than the wall position, back in the
\sphinxstyleemphasis{hi} direction.  A \sphinxstyleemphasis{hi} face reflects particles that move to a
coordinate higher than the wall position, back in the \sphinxstyleemphasis{lo} direction.

\sphinxAtStartPar
The position of each wall can be specified in one of 3 ways: as the
EDGE of the simulation box, as a constant value, or as a variable.  If
EDGE is used, then the corresponding boundary of the current
simulation box is used.  If a numeric constant is specified then the
wall is placed at that position in the appropriate dimension (x, y, or
z).  In both the EDGE and constant cases, the wall will never move.
If the wall position is a variable, it should be specified as v\_name,
where name is an {\hyperref[\detokenize{variable::doc}]{\sphinxcrossref{\DUrole{doc}{equal\sphinxhyphen{}style variable}}}} name.  In this
case the variable is evaluated each timestep and the result becomes
the current position of the reflecting wall.  Equal\sphinxhyphen{}style variables
can specify formulas with various mathematical functions, and include
{\hyperref[\detokenize{thermo_style::doc}]{\sphinxcrossref{\DUrole{doc}{thermo\_style}}}} command keywords for the simulation
box parameters and timestep and elapsed time.  Thus it is easy to
specify a time\sphinxhyphen{}dependent wall position.

\sphinxAtStartPar
The \sphinxstyleemphasis{units} keyword determines the meaning of the distance units used
to define a wall position, but only when a numeric constant or
variable is used.  It is not relevant when EDGE is used to specify a
face position.  In the variable case, the variable is assumed to
produce a value compatible with the \sphinxstyleemphasis{units} setting you specify.

\sphinxAtStartPar
A \sphinxstyleemphasis{box} value selects standard distance units as defined by the
{\hyperref[\detokenize{units::doc}]{\sphinxcrossref{\DUrole{doc}{units}}}} command, e.g. Angstroms for units = real or metal.
A \sphinxstyleemphasis{lattice} value means the distance units are in lattice spacings.
The {\hyperref[\detokenize{lattice::doc}]{\sphinxcrossref{\DUrole{doc}{lattice}}}} command must have been previously used to
define the lattice spacings.


\bigskip\hrule\bigskip


\sphinxAtStartPar
Here are examples of variable definitions that move the wall position
in a time\sphinxhyphen{}dependent fashion using equal\sphinxhyphen{}style
{\hyperref[\detokenize{variable::doc}]{\sphinxcrossref{\DUrole{doc}{variables}}}}.

\begin{sphinxVerbatim}[commandchars=\\\{\}]
\PYG{k}{variable }\PYG{n+nv+nvIdentifier}{ramp}\PYG{+w}{ }\PYG{n}{equal}\PYG{+w}{ }\PYG{n}{ramp}\PYG{n+nv}{(}\PYG{n+nv}{0,10}\PYG{n+nv}{)}
\PYG{k}{fix }\PYG{n+nv+nvIdentifier}{1}\PYG{+w}{ }\PYG{n+nv+nvIdentifier}{all}\PYG{+w}{ }\PYG{n}{wall}\PYG{o}{/}\PYG{n}{reflect}\PYG{+w}{ }\PYG{n}{xlo}\PYG{+w}{ }\PYG{n}{v\PYGZus{}ramp}

\PYG{k}{variable }\PYG{n+nv+nvIdentifier}{linear}\PYG{+w}{ }\PYG{n}{equal}\PYG{+w}{ }\PYG{n}{vdisplace}\PYG{n+nv}{(}\PYG{n+nv}{0,20}\PYG{n+nv}{)}
\PYG{k}{fix }\PYG{n+nv+nvIdentifier}{1}\PYG{+w}{ }\PYG{n+nv+nvIdentifier}{all}\PYG{+w}{ }\PYG{n}{wall}\PYG{o}{/}\PYG{n}{reflect}\PYG{+w}{ }\PYG{n}{xlo}\PYG{+w}{ }\PYG{n}{v\PYGZus{}linear}

\PYG{k}{variable }\PYG{n+nv+nvIdentifier}{wiggle}\PYG{+w}{ }\PYG{n}{equal}\PYG{+w}{ }\PYG{n}{swiggle}\PYG{n+nv}{(}\PYG{n+nv}{0.0,5.0,3.0}\PYG{n+nv}{)}
\PYG{k}{fix }\PYG{n+nv+nvIdentifier}{1}\PYG{+w}{ }\PYG{n+nv+nvIdentifier}{all}\PYG{+w}{ }\PYG{n}{wall}\PYG{o}{/}\PYG{n}{reflect}\PYG{+w}{ }\PYG{n}{xlo}\PYG{+w}{ }\PYG{n}{v\PYGZus{}wiggle}

\PYG{k}{variable }\PYG{n+nv+nvIdentifier}{wiggle}\PYG{+w}{ }\PYG{n}{equal}\PYG{+w}{ }\PYG{n}{cwiggle}\PYG{n+nv}{(}\PYG{n+nv}{0.0,5.0,3.0}\PYG{n+nv}{)}
\PYG{k}{fix }\PYG{n+nv+nvIdentifier}{1}\PYG{+w}{ }\PYG{n+nv+nvIdentifier}{all}\PYG{+w}{ }\PYG{n}{wall}\PYG{o}{/}\PYG{n}{reflect}\PYG{+w}{ }\PYG{n}{xlo}\PYG{+w}{ }\PYG{n}{v\PYGZus{}wiggle}
\end{sphinxVerbatim}

\sphinxAtStartPar
The \sphinxstyleemphasis{ramp(lo,hi)} function adjusts the wall position linearly from \sphinxstyleemphasis{lo} to
\sphinxstyleemphasis{hi} over the course of a run.  The \sphinxstyleemphasis{vdisplace(c0,velocity)} function does
something similar using the equation \sphinxstyleemphasis{position = c0 + velocity*delta},
where \sphinxstyleemphasis{delta} is the elapsed time.

\sphinxAtStartPar
The \sphinxstyleemphasis{swiggle(c0,A,period)} function causes the wall position to
oscillate sinusoidally according to this equation, where \sphinxstyleemphasis{omega = 2 PI
/ period}:
\begin{sphinxalltt}
position = c0 + A sin(omega*delta)
\end{sphinxalltt}

\sphinxAtStartPar
The \sphinxstyleemphasis{cwiggle(c0,A,period)} function causes the wall position to
oscillate sinusoidally according to this equation, which will have an
initial wall velocity of 0.0, and thus may impose a gentler
perturbation on the particles:
\begin{sphinxalltt}
position = c0 + A (1 \sphinxhyphen{} cos(omega*delta))
\end{sphinxalltt}


\bigskip\hrule\bigskip


\sphinxAtStartPar
Styles with a \sphinxstyleemphasis{gpu}, \sphinxstyleemphasis{intel}, \sphinxstyleemphasis{kk}, \sphinxstyleemphasis{omp}, or \sphinxstyleemphasis{opt} suffix are
functionally the same as the corresponding style without the suffix.
They have been optimized to run faster, depending on your available
hardware, as discussed on the {\hyperref[\detokenize{Speed_packages::doc}]{\sphinxcrossref{\DUrole{doc}{Accelerator packages}}}}
page.  The accelerated styles take the same arguments and should
produce the same results, except for round\sphinxhyphen{}off and precision issues.

\sphinxAtStartPar
These accelerated styles are part of the GPU, INTEL, KOKKOS,
OPENMP, and OPT packages, respectively.  They are only enabled if
LAMMPS was built with those packages.  See the {\hyperref[\detokenize{Build_package::doc}]{\sphinxcrossref{\DUrole{doc}{Build package}}}} page for more info.

\sphinxAtStartPar
You can specify the accelerated styles explicitly in your input script
by including their suffix, or you can use the {\hyperref[\detokenize{Run_options::doc}]{\sphinxcrossref{\DUrole{doc}{\sphinxhyphen{}suffix command\sphinxhyphen{}line switch}}}} when you invoke LAMMPS, or you can use the
{\hyperref[\detokenize{suffix::doc}]{\sphinxcrossref{\DUrole{doc}{suffix}}}} command in your input script.

\sphinxAtStartPar
See the {\hyperref[\detokenize{Speed_packages::doc}]{\sphinxcrossref{\DUrole{doc}{Accelerator packages}}}} page for more
instructions on how to use the accelerated styles effectively.


\bigskip\hrule\bigskip



\subsection{Restart, fix\_modify, output, run start/stop, minimize info}
\label{\detokenize{fix_wall_reflect:restart-fix-modify-output-run-start-stop-minimize-info}}
\sphinxAtStartPar
No information about this fix is written to {\hyperref[\detokenize{restart::doc}]{\sphinxcrossref{\DUrole{doc}{binary restart files}}}}.  None of the {\hyperref[\detokenize{fix_modify::doc}]{\sphinxcrossref{\DUrole{doc}{fix\_modify}}}} options are
relevant to this fix.  No global or per\sphinxhyphen{}atom quantities are stored by
this fix for access by various {\hyperref[\detokenize{Howto_output::doc}]{\sphinxcrossref{\DUrole{doc}{output commands}}}}.
No parameter of this fix can be used with the \sphinxstyleemphasis{start/stop} keywords of
the {\hyperref[\detokenize{run::doc}]{\sphinxcrossref{\DUrole{doc}{run}}}} command.  This fix is not invoked during
{\hyperref[\detokenize{minimize::doc}]{\sphinxcrossref{\DUrole{doc}{energy minimization}}}}.


\subsection{Restrictions}
\label{\detokenize{fix_wall_reflect:restrictions}}
\sphinxAtStartPar
Any dimension (xyz) that has a reflecting wall must be non\sphinxhyphen{}periodic.

\sphinxAtStartPar
A reflecting wall should not be used with rigid bodies such as those
defined by a “fix rigid” command.  This is because the wall/reflect
displaces atoms directly rather than exerts a force on them.  For
rigid bodies, use a soft wall instead, such as {\hyperref[\detokenize{fix_wall::doc}]{\sphinxcrossref{\DUrole{doc}{fix wall/lj93}}}}.  LAMMPS will flag the use of a rigid fix with fix
wall/reflect with a warning, but will not generate an error.


\subsection{Related commands}
\label{\detokenize{fix_wall_reflect:related-commands}}
\sphinxAtStartPar
{\hyperref[\detokenize{fix_wall::doc}]{\sphinxcrossref{\DUrole{doc}{fix wall/lj93}}}}, {\hyperref[\detokenize{fix_oneway::doc}]{\sphinxcrossref{\DUrole{doc}{fix oneway}}}}


\subsection{Default}
\label{\detokenize{fix_wall_reflect:default}}
\sphinxAtStartPar
The default for the units keyword is lattice.


\bigskip\hrule\bigskip

\phantomsection\label{\detokenize{fix_wall_reflect:bond1}}
\sphinxAtStartPar
\sphinxstylestrong{(Bond)} Bond and Leimkuhler, SIAM J Sci Comput, 30, p 134 (2007).

\sphinxstepscope

\index{fix wall/reflect/stochastic@\spxentry{fix wall/reflect/stochastic}}\ignorespaces 

\section{fix wall/reflect/stochastic command}
\label{\detokenize{fix_wall_reflect_stochastic:fix-wall-reflect-stochastic-command}}\label{\detokenize{fix_wall_reflect_stochastic:index-0}}\label{\detokenize{fix_wall_reflect_stochastic::doc}}

\subsection{Syntax}
\label{\detokenize{fix_wall_reflect_stochastic:syntax}}
\begin{sphinxVerbatim}[commandchars=\\\{\}]
\PYG{k}{fix }\PYG{n+nv+nvIdentifier}{ID}\PYG{+w}{ }\PYG{n+nv+nvIdentifier}{group\PYGZhy{}ID}\PYG{+w}{ }\PYG{n}{wall}\PYG{o}{/}\PYG{n}{reflect}\PYG{o}{/}\PYG{n}{stochastic}\PYG{+w}{ }\PYG{n}{rstyle}\PYG{+w}{ }\PYG{n}{seed}\PYG{+w}{ }\PYG{n}{face}\PYG{+w}{ }\PYG{n}{args}\PYG{+w}{ }\PYG{n}{...}\PYG{+w}{ }\PYG{n}{keyword}\PYG{+w}{ }\PYG{n}{value}\PYG{+w}{ }\PYG{n}{...}
\end{sphinxVerbatim}
\begin{itemize}
\item {} 
\sphinxAtStartPar
ID, group\sphinxhyphen{}ID are documented in {\hyperref[\detokenize{fix::doc}]{\sphinxcrossref{\DUrole{doc}{fix}}}} command

\item {} 
\sphinxAtStartPar
wall/reflect/stochastic = style name of this fix command

\item {} 
\sphinxAtStartPar
rstyle = diffusive or maxwell or ccl

\item {} 
\sphinxAtStartPar
seed = random seed for stochasticity (positive integer)

\item {} 
\sphinxAtStartPar
one or more face/args pairs may be appended

\item {} 
\sphinxAtStartPar
face = \sphinxstyleemphasis{xlo} or \sphinxstyleemphasis{xhi} or \sphinxstyleemphasis{ylo} or \sphinxstyleemphasis{yhi} or \sphinxstyleemphasis{zlo} or \sphinxstyleemphasis{zhi}

\begin{sphinxVerbatim}[commandchars=\\\{\}]
args = pos temp velx vely velz accomx accomy accomz
  pos = EDGE or constant
    EDGE = current lo or hi edge of simulation box
    constant = number like 0.0 or 30.0 (distance units)
  temp = wall temperature (temperature units)
  velx,vely,velz = wall velocity in x,y,z directions (velocity units)
  accomx,accomy,accomz = accommodation coeffs in x,y,z directions (unitless)
    not specified for rstyle = diffusive
    single accom coeff specified for rstyle maxwell
    all 3 coeffs specified for rstyle cll
\end{sphinxVerbatim}

\item {} 
\sphinxAtStartPar
zero or more keyword/value pairs may be appended

\item {} 
\sphinxAtStartPar
keyword = \sphinxstyleemphasis{units}
\begin{sphinxalltt}
\sphinxstyleemphasis{units} value = \sphinxstyleemphasis{lattice} or \sphinxstyleemphasis{box}
  \sphinxstyleemphasis{lattice} = the wall position is defined in lattice units
  \sphinxstyleemphasis{box} = the wall position is defined in simulation box units
\end{sphinxalltt}

\end{itemize}


\subsection{Examples}
\label{\detokenize{fix_wall_reflect_stochastic:examples}}
\begin{sphinxVerbatim}[commandchars=\\\{\}]
\PYG{k}{fix }\PYG{n+nv+nvIdentifier}{zwalls}\PYG{+w}{ }\PYG{n+nv+nvIdentifier}{all}\PYG{+w}{ }\PYG{n}{wall}\PYG{o}{/}\PYG{n}{reflect}\PYG{o}{/}\PYG{n}{stochastic}\PYG{+w}{ }\PYG{n}{diffusive}\PYG{+w}{ }\PYG{l+m}{23424}\PYG{+w}{ }\PYG{n}{zlo}\PYG{+w}{ }\PYG{n}{EDGE}\PYG{+w}{ }\PYG{l+m}{300}\PYG{+w}{ }\PYG{l+m}{0.1}\PYG{+w}{ }\PYG{l+m}{0.1}\PYG{+w}{ }\PYG{l+m}{0}\PYG{+w}{ }\PYG{n}{zhi}\PYG{+w}{ }\PYG{n}{EDGE}\PYG{+w}{ }\PYG{l+m}{200}\PYG{+w}{ }\PYG{l+m}{0.1}\PYG{+w}{ }\PYG{l+m}{0.1}\PYG{+w}{ }\PYG{l+m}{0}
\PYG{k}{fix }\PYG{n+nv+nvIdentifier}{ywalls}\PYG{+w}{ }\PYG{n+nv+nvIdentifier}{all}\PYG{+w}{ }\PYG{n}{wall}\PYG{o}{/}\PYG{n}{reflect}\PYG{o}{/}\PYG{n}{stochastic}\PYG{+w}{ }\PYG{n}{maxwell}\PYG{+w}{ }\PYG{l+m}{345533}\PYG{+w}{ }\PYG{n}{ylo}\PYG{+w}{ }\PYG{l+m}{5.0}\PYG{+w}{ }\PYG{l+m}{300}\PYG{+w}{ }\PYG{l+m}{0.1}\PYG{+w}{ }\PYG{l+m}{0.0}\PYG{+w}{ }\PYG{l+m}{0.0}\PYG{+w}{ }\PYG{l+m}{0.8}\PYG{+w}{ }\PYG{n}{yhi}\PYG{+w}{ }\PYG{l+m}{10.0}\PYG{+w}{ }\PYG{l+m}{300}\PYG{+w}{ }\PYG{l+m}{0.1}\PYG{+w}{ }\PYG{l+m}{0.0}\PYG{+w}{ }\PYG{l+m}{0.0}\PYG{+w}{ }\PYG{l+m}{0.8}
\PYG{k}{fix }\PYG{n+nv+nvIdentifier}{xwalls}\PYG{+w}{ }\PYG{n+nv+nvIdentifier}{all}\PYG{+w}{ }\PYG{n}{wall}\PYG{o}{/}\PYG{n}{reflect}\PYG{o}{/}\PYG{n}{stochastic}\PYG{+w}{ }\PYG{n}{cercignanilampis}\PYG{+w}{ }\PYG{l+m}{2308}\PYG{+w}{ }\PYG{n}{xlo}\PYG{+w}{ }\PYG{l+m}{0.0}\PYG{+w}{ }\PYG{l+m}{300}\PYG{+w}{ }\PYG{l+m}{0.0}\PYG{+w}{ }\PYG{l+m}{0.1}\PYG{+w}{ }\PYG{l+m}{0.9}\PYG{+w}{ }\PYG{l+m}{0.8}\PYG{+w}{ }\PYG{l+m}{0.7}\PYG{+w}{ }\PYG{n}{xhi}\PYG{+w}{ }\PYG{n}{EDGE}\PYG{+w}{ }\PYG{l+m}{300}\PYG{+w}{ }\PYG{l+m}{0.0}\PYG{+w}{ }\PYG{l+m}{0.1}\PYG{+w}{ }\PYG{l+m}{0}\PYG{+w}{ }\PYG{l+m}{0.9}\PYG{+w}{ }\PYG{l+m}{0.8}\PYG{+w}{ }\PYG{l+m}{0.7}\PYG{+w}{ }\PYG{n}{units}\PYG{+w}{ }\PYG{n}{box}
\end{sphinxVerbatim}


\subsection{Description}
\label{\detokenize{fix_wall_reflect_stochastic:description}}
\sphinxAtStartPar
Bound the simulation with one or more walls which reflect particles
in the specified group when they attempt to move through them.

\sphinxAtStartPar
Reflection means that if an atom moves outside the wall on a timestep
(e.g. due to the {\hyperref[\detokenize{fix_nve::doc}]{\sphinxcrossref{\DUrole{doc}{fix nve}}}} command), then it is put back
inside the wall with a changed velocity.

\sphinxAtStartPar
This fix models treats the wall as a moving solid boundary with a
finite temperature, which can exchange energy with particles that
collide with it.  This is different than the simpler {\hyperref[\detokenize{fix_wall_reflect::doc}]{\sphinxcrossref{\DUrole{doc}{fix wall/reflect}}}} command which models mirror
reflection.  For this fix, the post collision velocity of each
particle is treated stochastically.  The randomness can come from many
sources: thermal motion of the wall atoms, surface roughness, etc.
Three stochastic reflection models are currently implemented.

\sphinxAtStartPar
For rstyle \sphinxstyleemphasis{diffusive}, particles are reflected diffusively. Their
velocity distribution corresponds to an equilibrium distribution of
particles at the wall temperature.  No accommodation coefficients
are specified.

\sphinxAtStartPar
For rstyle \sphinxstyleemphasis{maxwell}, particle reflection is Maxwellian which means
partially diffusive and partially specular ({\hyperref[\detokenize{fix_wall_reflect_stochastic:maxwell}]{\sphinxcrossref{\DUrole{std,std-ref}{Maxwell}}}}).  A
single accommodation coeff is specified which must be between 0.0 and
1.0 inclusive.  It determines the fraction of the collision which is
diffusive versus specular.  An accommodation coefficient of 1.0 is fully
diffusive; a coefficient of 0.0 is fully specular.

\sphinxAtStartPar
For rstyle \sphinxstyleemphasis{cll}, particle collisions are computed by the
Cercignani/Lampis model.  See {\hyperref[\detokenize{fix_wall_reflect_stochastic:cl}]{\sphinxcrossref{\DUrole{std,std-ref}{CL}}}} and {\hyperref[\detokenize{fix_wall_reflect_stochastic:to}]{\sphinxcrossref{\DUrole{std,std-ref}{To}}}} for details.
Three accommodations coefficient are specified.  Each must be between
0.0 and 1.0 inclusive.  Two are velocity accommodation coefficients;
one is a normal kinetic energy accommodation.  The normal coeff is the
one corresponding to the normal of the wall itself.  For example if
the wall is \sphinxstyleemphasis{ylo} or \sphinxstyleemphasis{yhi}, \sphinxstyleemphasis{accomx} and \sphinxstyleemphasis{accomz} are the tangential
velocity accommodation coefficients, and \sphinxstyleemphasis{accomy} is the normal
kinetic energy accommodation coefficient.

\sphinxAtStartPar
The optional \sphinxstyleemphasis{units} keyword determines the distance units used to
define a wall position.  A \sphinxstyleemphasis{box} value selects standard distance units
as defined by the {\hyperref[\detokenize{units::doc}]{\sphinxcrossref{\DUrole{doc}{units}}}} command, e.g. Angstroms for units
= real or metal.  A \sphinxstyleemphasis{lattice} value means the distance units are in
lattice spacings. The {\hyperref[\detokenize{lattice::doc}]{\sphinxcrossref{\DUrole{doc}{lattice}}}} command must have been
previously used to define the lattice spacings.


\bigskip\hrule\bigskip



\subsection{Restrictions}
\label{\detokenize{fix_wall_reflect_stochastic:restrictions}}
\sphinxAtStartPar
This fix has the same limitations as the {\hyperref[\detokenize{fix_wall_reflect::doc}]{\sphinxcrossref{\DUrole{doc}{fix wall/reflect}}}} command.  Any dimension (xyz) that
has a wall must be non\sphinxhyphen{}periodic.  It should not be used with rigid
bodies such as those defined by the {\hyperref[\detokenize{fix_rigid::doc}]{\sphinxcrossref{\DUrole{doc}{fix rigid}}}}
command.  The wall velocity must lie on the same plane as the wall
itself.

\sphinxAtStartPar
This fix is part of the EXTRA\sphinxhyphen{}FIX package.  It is only enabled if
LAMMPS was built with that package.  See the {\hyperref[\detokenize{Build_package::doc}]{\sphinxcrossref{\DUrole{doc}{Build package}}}} page for more info.


\subsection{Related commands}
\label{\detokenize{fix_wall_reflect_stochastic:related-commands}}
\sphinxAtStartPar
{\hyperref[\detokenize{fix_wall_reflect::doc}]{\sphinxcrossref{\DUrole{doc}{fix wall/reflect}}}}


\subsection{Default}
\label{\detokenize{fix_wall_reflect_stochastic:default}}
\sphinxAtStartPar
The default for the units keyword is lattice.


\bigskip\hrule\bigskip

\phantomsection\label{\detokenize{fix_wall_reflect_stochastic:maxwell}}
\sphinxAtStartPar
\sphinxstylestrong{(Maxwell)} J.C. Maxwell, Philos. Tans. Royal Soc. London, 157: 49\sphinxhyphen{}88
(1867).

\phantomsection\label{\detokenize{fix_wall_reflect_stochastic:cl}}
\sphinxAtStartPar
\sphinxstylestrong{(Cercignani)} C. Cercignani and M. Lampis. Trans. Theory
Stat. Phys. 1, 2, 101 (1971).

\phantomsection\label{\detokenize{fix_wall_reflect_stochastic:to}}
\sphinxAtStartPar
\sphinxstylestrong{(To)} Q.D. To, V.H. Vu, G. Lauriat, and
C. Leonard. J. Math. Phys. 56, 103101 (2015).

\sphinxstepscope

\index{fix wall/region@\spxentry{fix wall/region}}\ignorespaces 

\section{fix wall/region command}
\label{\detokenize{fix_wall_region:fix-wall-region-command}}\label{\detokenize{fix_wall_region:index-0}}\label{\detokenize{fix_wall_region::doc}}

\subsection{Syntax}
\label{\detokenize{fix_wall_region:syntax}}
\begin{sphinxVerbatim}[commandchars=\\\{\}]
\PYG{k}{fix }\PYG{n+nv+nvIdentifier}{ID}\PYG{+w}{ }\PYG{n+nv+nvIdentifier}{group\PYGZhy{}ID}\PYG{+w}{ }\PYG{n}{wall}\PYG{o}{/}\PYG{k}{region }\PYG{n+nv+nvIdentifier}{region\PYGZhy{}ID}\PYG{+w}{ }\PYG{n}{style}\PYG{+w}{ }\PYG{n}{args}\PYG{+w}{ }\PYG{n}{...}\PYG{+w}{ }\PYG{n}{cutoff}
\end{sphinxVerbatim}
\begin{itemize}
\item {} 
\sphinxAtStartPar
ID, group\sphinxhyphen{}ID are documented in {\hyperref[\detokenize{fix::doc}]{\sphinxcrossref{\DUrole{doc}{fix}}}} command

\item {} 
\sphinxAtStartPar
wall/region = style name of this fix command

\item {} 
\sphinxAtStartPar
region\sphinxhyphen{}ID = region whose boundary will act as wall

\item {} 
\sphinxAtStartPar
style = \sphinxstyleemphasis{lj93} or \sphinxstyleemphasis{lj126} or \sphinxstyleemphasis{lj1043} or \sphinxstyleemphasis{colloid} or \sphinxstyleemphasis{harmonic} or \sphinxstyleemphasis{morse}

\item {} 
\sphinxAtStartPar
args for styles \sphinxstyleemphasis{lj93} or \sphinxstyleemphasis{lj126} or \sphinxstyleemphasis{lj1043} or \sphinxstyleemphasis{colloid} or \sphinxstyleemphasis{harmonic} =
\begin{sphinxalltt}
epsilon = strength factor for wall\sphinxhyphen{}particle interaction (energy or energy/distance\textasciicircum{}2 units)
sigma = size factor for wall\sphinxhyphen{}particle interaction (distance units)
\end{sphinxalltt}

\item {} 
\sphinxAtStartPar
args for style \sphinxstyleemphasis{morse} =

\begin{sphinxVerbatim}[commandchars=\\\{\}]
D\PYGZus{}0 = depth of the potential (energy units)
alpha = width parameter (1/distance units)
r\PYGZus{}0 = distance of the potential minimum from wall position (distance units)
\end{sphinxVerbatim}

\item {} 
\sphinxAtStartPar
cutoff = distance from wall at which wall\sphinxhyphen{}particle interaction is cut off (distance units)

\end{itemize}


\subsection{Examples}
\label{\detokenize{fix_wall_region:examples}}
\begin{sphinxVerbatim}[commandchars=\\\{\}]
\PYG{k}{fix }\PYG{n+nv+nvIdentifier}{wall}\PYG{+w}{ }\PYG{n+nv+nvIdentifier}{all}\PYG{+w}{ }\PYG{n}{wall}\PYG{o}{/}\PYG{k}{region }\PYG{n+nv+nvIdentifier}{mySphere}\PYG{+w}{ }\PYG{n}{lj93}\PYG{+w}{ }\PYG{l+m}{1.0}\PYG{+w}{ }\PYG{l+m}{1.0}\PYG{+w}{ }\PYG{l+m}{2.5}
\PYG{k}{fix }\PYG{n+nv+nvIdentifier}{wall}\PYG{+w}{ }\PYG{n+nv+nvIdentifier}{all}\PYG{+w}{ }\PYG{n}{wall}\PYG{o}{/}\PYG{k}{region }\PYG{n+nv+nvIdentifier}{mySphere}\PYG{+w}{ }\PYG{n}{harmonic}\PYG{+w}{ }\PYG{l+m}{1.0}\PYG{+w}{ }\PYG{l+m}{0.0}\PYG{+w}{ }\PYG{l+m}{2.5}
\PYG{k}{fix }\PYG{n+nv+nvIdentifier}{wall}\PYG{+w}{ }\PYG{n+nv+nvIdentifier}{all}\PYG{+w}{ }\PYG{n}{wall}\PYG{o}{/}\PYG{k}{region }\PYG{n+nv+nvIdentifier}{box\PYGZus{}top}\PYG{+w}{ }\PYG{n}{morse}\PYG{+w}{ }\PYG{l+m}{1.0}\PYG{+w}{ }\PYG{l+m}{1.0}\PYG{+w}{ }\PYG{l+m}{1.5}\PYG{+w}{ }\PYG{l+m}{3.0}
\end{sphinxVerbatim}


\subsection{Description}
\label{\detokenize{fix_wall_region:description}}
\sphinxAtStartPar
Treat the surface of the geometric region defined by the \sphinxstyleemphasis{region\sphinxhyphen{}ID}
as a bounding wall which interacts with nearby particles according to
the specified style.

\sphinxAtStartPar
The distance between a particle and the surface is the distance to the
nearest point on the surface and the force the wall exerts on the
particle is along the direction between that point and the particle,
which is the direction normal to the surface at that point.  Note that
if the region surface is comprised of multiple “faces”, then each face
can exert a force on the particle if it is close enough.  E.g. for
{\hyperref[\detokenize{region::doc}]{\sphinxcrossref{\DUrole{doc}{region\_style block}}}}, a particle in the interior, near a
corner of the block, could feel wall forces from 1, 2, or 3 faces of
the block.

\sphinxAtStartPar
Regions are defined using the {\hyperref[\detokenize{region::doc}]{\sphinxcrossref{\DUrole{doc}{region}}}} command.  Note that
the region volume can be interior or exterior to the bounding surface,
which will determine in which direction the surface interacts with
particles, i.e. the direction of the surface normal.  The surface of
the region only exerts forces on particles “inside” the region; if a
particle is “outside” the region it will generate an error, because it
has moved through the wall.

\sphinxAtStartPar
Regions can either be primitive shapes (block, sphere, cylinder, etc)
or combinations of primitive shapes specified via the \sphinxstyleemphasis{union} or
\sphinxstyleemphasis{intersect} region styles.  These latter styles can be used to
construct particle containers with complex shapes.  Regions can also
change over time via the {\hyperref[\detokenize{region::doc}]{\sphinxcrossref{\DUrole{doc}{region}}}} command keywords (move)
and \sphinxstyleemphasis{rotate}.  If such a region is used with this fix, then the of
region surface will move over time in the corresponding manner.

\begin{sphinxadmonition}{note}{Note:}
\sphinxAtStartPar
As discussed on the {\hyperref[\detokenize{region::doc}]{\sphinxcrossref{\DUrole{doc}{region}}}} command doc page,
regions in LAMMPS do not get wrapped across periodic boundaries.  It
is up to you to ensure that periodic or non\sphinxhyphen{}periodic boundaries are
specified appropriately via the {\hyperref[\detokenize{boundary::doc}]{\sphinxcrossref{\DUrole{doc}{boundary}}}} command when
using a region as a wall that bounds particle motion.  This also means
that if you embed a region in your simulation box and want it to
repulse particles from its surface (using the “side out” option in the
{\hyperref[\detokenize{region::doc}]{\sphinxcrossref{\DUrole{doc}{region}}}} command), that its repulsive force will not be
felt across a periodic boundary.
\end{sphinxadmonition}

\begin{sphinxadmonition}{note}{Note:}
\sphinxAtStartPar
For primitive regions with sharp corners and/or edges (e.g. a
block or cylinder), wall/particle forces are computed accurately for
both interior and exterior regions.  For \sphinxstyleemphasis{union} and \sphinxstyleemphasis{intersect}
regions, additional sharp corners and edges may be present due to the
intersection of the surfaces of 2 or more primitive volumes.  These
corners and edges can be of two types: concave or convex.  Concave
points/edges are like the corners of a cube as seen by particles in
the interior of a cube.  Wall/particle forces around these features
are computed correctly.  Convex points/edges are like the corners of a
cube as seen by particles exterior to the cube, i.e. the points jut
into the volume where particles are present.  LAMMPS does NOT compute
the location of these convex points directly, and hence wall/particle
forces in the cutoff volume around these points suffer from
inaccuracies.  The basic problem is that the outward normal of the
surface is not continuous at these points.  This can cause particles
to feel no force (they don’t “see” the wall) when in one location,
then move a distance epsilon, and suddenly feel a large force because
they now “see” the wall.  In a worst\sphinxhyphen{}case scenario, this can blow
particles out of the simulation box.  Thus, as a general rule you
should not use the fix wall/gran/region command with \sphinxstyleemphasis{union} or
\sphinxstyleemphasis{interesect} regions that have convex points or edges resulting from
the union/intersection (convex points/edges in the union/intersection
due to a single sub\sphinxhyphen{}region are still OK).
\end{sphinxadmonition}

\begin{sphinxadmonition}{note}{Note:}
\sphinxAtStartPar
Similarly, you should not define \sphinxstyleemphasis{union} or \sphinxstyleemphasis{intersert} regions
for use with this command that share an overlapping common face that
is part of the overall outer boundary (interior boundary is OK), even
if the face is smooth.  E.g. two regions of style block in a \sphinxstyleemphasis{union}
region, where the two blocks overlap on one or more of their faces.
This is because LAMMPS discards points that are part of multiple
sub\sphinxhyphen{}regions when calculating wall/particle interactions, to avoid
double\sphinxhyphen{}counting the interaction.  Having two coincident faces could
cause the face to become invisible to the particles.  The solution is
to make the two faces differ by epsilon in their position.
\end{sphinxadmonition}

\sphinxAtStartPar
The energy of wall\sphinxhyphen{}particle interactions depends on the specified
style.

\sphinxAtStartPar
For style \sphinxstyleemphasis{lj93}, the energy E is given by the 9/3 potential:
\begin{equation*}
\begin{split}E = \epsilon \left[ \frac{2}{15} \left(\frac{\sigma}{r}\right)^{9} -
                      \left(\frac{\sigma}{r}\right)^3 \right]
                      \qquad r < r_c\end{split}
\end{equation*}
\sphinxAtStartPar
For style \sphinxstyleemphasis{lj126}, the energy E is given by the 12/6 potential:
\begin{equation*}
\begin{split}E = 4 \epsilon \left[ \left(\frac{\sigma}{r}\right)^{12} -
                      \left(\frac{\sigma}{r}\right)^6 \right]
                      \qquad r < r_c\end{split}
\end{equation*}
\sphinxAtStartPar
For style \sphinxstyleemphasis{wall/lj1043}, the energy E is given by the 10/4/3 potential:
\begin{equation*}
\begin{split}E = 2 \pi \epsilon \left[ \frac{2}{5} \left(\frac{\sigma}{r}\right)^{10} -
                      \left(\frac{\sigma}{r}\right)^4 -
                      \frac{\sqrt(2)\sigma^3}{3\left(r+\left(0.61/\sqrt(2)\right)\sigma\right)^3}\right]
                      \qquad r < r_c\end{split}
\end{equation*}
\sphinxAtStartPar
For style \sphinxstyleemphasis{colloid}, the energy E is given by an integrated form of
the {\hyperref[\detokenize{pair_colloid::doc}]{\sphinxcrossref{\DUrole{doc}{pair\_style colloid}}}} potential:
\begin{equation*}
\begin{split}E = & \epsilon \left[ \frac{\sigma^{6}}{7560}
\left(\frac{6R-D}{D^{7}} + \frac{D+8R}{(D+2R)^{7}} \right) \right. \\
 & \left. - \frac{1}{6} \left(\frac{2R(D+R) + D(D+2R)
 \left[ \ln D - \ln (D+2R) \right]}{D(D+2R)} \right) \right] \qquad r < r_c\end{split}
\end{equation*}
\sphinxAtStartPar
For style \sphinxstyleemphasis{wall/harmonic}, the energy E is given by a harmonic spring
potential (the distance parameter is ignored):
\begin{equation*}
\begin{split}E = \epsilon \quad (r - r_c)^2 \qquad r < r_c\end{split}
\end{equation*}
\sphinxAtStartPar
For style \sphinxstyleemphasis{wall/morse}, the energy E is given by the Morse potential:
\begin{equation*}
\begin{split}E = D_0 \left[ e^{- 2 \alpha (r - r_0)} - 2 e^{- \alpha (r - r_0)} \right]
    \qquad r < r_c\end{split}
\end{equation*}
\sphinxAtStartPar
Unlike other styles, this requires three parameters (\(D_0\),
\(\alpha\), and \(r_0\) in this order) instead of two like
for the other wall styles.

\sphinxAtStartPar
In all cases, \sphinxstyleemphasis{r} is the distance from the particle to the region
surface, and Rc is the \sphinxstyleemphasis{cutoff} distance at which the particle and
surface no longer interact.  The cutoff is always the last argument.
The energy of the wall potential is shifted so that the wall\sphinxhyphen{}particle
interaction energy is 0.0 at the cutoff distance.

\sphinxAtStartPar
For a full description of these wall styles, see fix\_style
{\hyperref[\detokenize{fix_wall::doc}]{\sphinxcrossref{\DUrole{doc}{wall}}}}


\subsection{Restart, fix\_modify, output, run start/stop, minimize info}
\label{\detokenize{fix_wall_region:restart-fix-modify-output-run-start-stop-minimize-info}}
\sphinxAtStartPar
No information about this fix is written to {\hyperref[\detokenize{restart::doc}]{\sphinxcrossref{\DUrole{doc}{binary restart files}}}}.

\sphinxAtStartPar
The {\hyperref[\detokenize{fix_modify::doc}]{\sphinxcrossref{\DUrole{doc}{fix\_modify}}}} \sphinxstyleemphasis{energy} option is supported by
this fix to add the energy of interaction between atoms and the region
wall to the global potential energy of the system as part of
{\hyperref[\detokenize{thermo_style::doc}]{\sphinxcrossref{\DUrole{doc}{thermodynamic output}}}}.  The default setting for
this fix is {\hyperref[\detokenize{fix_modify::doc}]{\sphinxcrossref{\DUrole{doc}{fix\_modify energy no}}}}.

\sphinxAtStartPar
The {\hyperref[\detokenize{fix_modify::doc}]{\sphinxcrossref{\DUrole{doc}{fix\_modify}}}} \sphinxstyleemphasis{virial} option is supported by
this fix to add the contribution due to the interaction between atoms
and the region wall to both the global pressure and per\sphinxhyphen{}atom stress of
the system via the {\hyperref[\detokenize{compute_pressure::doc}]{\sphinxcrossref{\DUrole{doc}{compute pressure}}}} and
{\hyperref[\detokenize{compute_stress_atom::doc}]{\sphinxcrossref{\DUrole{doc}{compute stress/atom}}}} commands.  The former
can be accessed by {\hyperref[\detokenize{thermo_style::doc}]{\sphinxcrossref{\DUrole{doc}{thermodynamic output}}}}.  The
default setting for this fix is {\hyperref[\detokenize{fix_modify::doc}]{\sphinxcrossref{\DUrole{doc}{fix\_modify virial no}}}}.

\sphinxAtStartPar
The {\hyperref[\detokenize{fix_modify::doc}]{\sphinxcrossref{\DUrole{doc}{fix\_modify}}}} \sphinxstyleemphasis{respa} option is supported by this
fix. This allows to set at which level of the {\hyperref[\detokenize{run_style::doc}]{\sphinxcrossref{\DUrole{doc}{r\sphinxhyphen{}RESPA}}}}
integrator the fix is adding its forces. Default is the outermost level.

\sphinxAtStartPar
This fix computes a global scalar energy and a global 3\sphinxhyphen{}length vector
of forces, which can be accessed by various {\hyperref[\detokenize{Howto_output::doc}]{\sphinxcrossref{\DUrole{doc}{output commands}}}}.  The scalar energy is the sum of energy interactions
for all particles interacting with the wall represented by the region
surface.  The 3 vector quantities are the x,y,z components of the
total force acting on the wall due to the particles.  The scalar and
vector values calculated by this fix are “extensive”.

\sphinxAtStartPar
No parameter of this fix can be used with the \sphinxstyleemphasis{start/stop} keywords of
the {\hyperref[\detokenize{run::doc}]{\sphinxcrossref{\DUrole{doc}{run}}}} command.

\sphinxAtStartPar
The forces due to this fix are imposed during an energy minimization,
invoked by the {\hyperref[\detokenize{minimize::doc}]{\sphinxcrossref{\DUrole{doc}{minimize}}}} command.

\begin{sphinxadmonition}{note}{Note:}
\sphinxAtStartPar
If you want the atom/wall interaction energy to be included in the
total potential energy of the system (the quantity being
minimized), you MUST enable the {\hyperref[\detokenize{fix_modify::doc}]{\sphinxcrossref{\DUrole{doc}{fix\_modify}}}}
\sphinxstyleemphasis{energy} option for this fix.
\end{sphinxadmonition}


\subsection{Restrictions}
\label{\detokenize{fix_wall_region:restrictions}}\begin{quote}

\sphinxAtStartPar
none
\end{quote}


\subsection{Related commands}
\label{\detokenize{fix_wall_region:related-commands}}
\sphinxAtStartPar
{\hyperref[\detokenize{fix_wall::doc}]{\sphinxcrossref{\DUrole{doc}{fix wall/lj93}}}},
{\hyperref[\detokenize{fix_wall::doc}]{\sphinxcrossref{\DUrole{doc}{fix wall/lj126}}}},
{\hyperref[\detokenize{fix_wall::doc}]{\sphinxcrossref{\DUrole{doc}{fix wall/lj1043}}}},
{\hyperref[\detokenize{fix_wall::doc}]{\sphinxcrossref{\DUrole{doc}{fix wall/colloid}}}},
{\hyperref[\detokenize{fix_wall::doc}]{\sphinxcrossref{\DUrole{doc}{fix wall/harmonic}}}},
{\hyperref[\detokenize{fix_wall_gran::doc}]{\sphinxcrossref{\DUrole{doc}{fix wall/gran}}}}


\subsection{Default}
\label{\detokenize{fix_wall_region:default}}
\sphinxAtStartPar
none

\sphinxstepscope

\index{fix wall/srd@\spxentry{fix wall/srd}}\ignorespaces 

\section{fix wall/srd command}
\label{\detokenize{fix_wall_srd:fix-wall-srd-command}}\label{\detokenize{fix_wall_srd:index-0}}\label{\detokenize{fix_wall_srd::doc}}

\subsection{Syntax}
\label{\detokenize{fix_wall_srd:syntax}}
\begin{sphinxVerbatim}[commandchars=\\\{\}]
\PYG{k}{fix }\PYG{n+nv+nvIdentifier}{ID}\PYG{+w}{ }\PYG{n+nv+nvIdentifier}{group\PYGZhy{}ID}\PYG{+w}{ }\PYG{n}{wall}\PYG{o}{/}\PYG{n}{srd}\PYG{+w}{ }\PYG{n}{face}\PYG{+w}{ }\PYG{n}{arg}\PYG{+w}{ }\PYG{n}{...}\PYG{+w}{ }\PYG{n}{keyword}\PYG{+w}{ }\PYG{n}{value}\PYG{+w}{ }\PYG{n}{...}
\end{sphinxVerbatim}
\begin{itemize}
\item {} 
\sphinxAtStartPar
ID, group\sphinxhyphen{}ID are documented in {\hyperref[\detokenize{fix::doc}]{\sphinxcrossref{\DUrole{doc}{fix}}}} command

\item {} 
\sphinxAtStartPar
wall/srd = style name of this fix command

\item {} 
\sphinxAtStartPar
one or more face/arg pairs may be appended

\item {} 
\sphinxAtStartPar
face = \sphinxstyleemphasis{xlo} or \sphinxstyleemphasis{xhi} or \sphinxstyleemphasis{ylo} or \sphinxstyleemphasis{yhi} or \sphinxstyleemphasis{zlo} or \sphinxstyleemphasis{zhi}
\begin{sphinxalltt}
\sphinxstyleemphasis{xlo},\sphinxstyleemphasis{ylo},\sphinxstyleemphasis{zlo} arg = EDGE or constant or variable
  EDGE = current lo edge of simulation box
  constant = number like 0.0 or \sphinxhyphen{}30.0 (distance units)
  variable = {\hyperref[\detokenize{variable::doc}]{\sphinxcrossref{\DUrole{doc}{equal\sphinxhyphen{}style variable}}}} like v\_x or v\_wiggle
\sphinxstyleemphasis{xhi},\sphinxstyleemphasis{yhi},\sphinxstyleemphasis{zhi} arg = EDGE or constant or variable
  EDGE = current hi edge of simulation box
  constant = number like 50.0 or 100.3 (distance units)
  variable = {\hyperref[\detokenize{variable::doc}]{\sphinxcrossref{\DUrole{doc}{equal\sphinxhyphen{}style variable}}}} like v\_x or v\_wiggle
\end{sphinxalltt}

\item {} 
\sphinxAtStartPar
zero or more keyword/value pairs may be appended

\item {} 
\sphinxAtStartPar
keyword = \sphinxstyleemphasis{units}
\begin{sphinxalltt}
\sphinxstyleemphasis{units} value = \sphinxstyleemphasis{lattice} or \sphinxstyleemphasis{box}
  \sphinxstyleemphasis{lattice} = the wall position is defined in lattice units
  \sphinxstyleemphasis{box} = the wall position is defined in simulation box units
\end{sphinxalltt}

\end{itemize}


\subsection{Examples}
\label{\detokenize{fix_wall_srd:examples}}
\begin{sphinxVerbatim}[commandchars=\\\{\}]
\PYG{k}{fix }\PYG{n+nv+nvIdentifier}{xwalls}\PYG{+w}{ }\PYG{n+nv+nvIdentifier}{all}\PYG{+w}{ }\PYG{n}{wall}\PYG{o}{/}\PYG{n}{srd}\PYG{+w}{ }\PYG{n}{xlo}\PYG{+w}{ }\PYG{n}{EDGE}\PYG{+w}{ }\PYG{n}{xhi}\PYG{+w}{ }\PYG{n}{EDGE}
\PYG{k}{fix }\PYG{n+nv+nvIdentifier}{walls}\PYG{+w}{ }\PYG{n+nv+nvIdentifier}{all}\PYG{+w}{ }\PYG{n}{wall}\PYG{o}{/}\PYG{n}{srd}\PYG{+w}{ }\PYG{n}{xlo}\PYG{+w}{ }\PYG{l+m}{0.0}\PYG{+w}{ }\PYG{n}{ylo}\PYG{+w}{ }\PYG{l+m}{10.0}\PYG{+w}{ }\PYG{n}{units}\PYG{+w}{ }\PYG{n}{box}
\PYG{k}{fix }\PYG{n+nv+nvIdentifier}{top}\PYG{+w}{ }\PYG{n+nv+nvIdentifier}{all}\PYG{+w}{ }\PYG{n}{wall}\PYG{o}{/}\PYG{n}{srd}\PYG{+w}{ }\PYG{n}{zhi}\PYG{+w}{ }\PYG{n}{v\PYGZus{}pressdown}
\end{sphinxVerbatim}


\subsection{Description}
\label{\detokenize{fix_wall_srd:description}}
\sphinxAtStartPar
Bound the simulation with one or more walls which interact with
stochastic reaction dynamics (SRD) particles as slip (smooth) or
no\sphinxhyphen{}slip (rough) flat surfaces.  The wall interaction is actually
invoked via the {\hyperref[\detokenize{fix_srd::doc}]{\sphinxcrossref{\DUrole{doc}{fix srd}}}} command, only on the group of
SRD particles it defines, so the group setting for the fix wall/srd
command is ignored.

\sphinxAtStartPar
A particle/wall collision occurs if an SRD particle moves outside the
wall on a timestep.  This alters the position and velocity of the SRD
particle and imparts a force to the wall.

\sphinxAtStartPar
The \sphinxstyleemphasis{collision} and \sphinxstyleemphasis{Tsrd} settings specified via the {\hyperref[\detokenize{fix_srd::doc}]{\sphinxcrossref{\DUrole{doc}{fix srd}}}} command affect the SRD/wall collisions.  A \sphinxstyleemphasis{slip}
setting for the \sphinxstyleemphasis{collision} keyword means that the tangential
component of the SRD particle momentum is preserved.  Thus only a
normal force is imparted to the wall.  The normal component of the new
SRD velocity is sampled from a Gaussian distribution at temperature
\sphinxstyleemphasis{Tsrd}.

\sphinxAtStartPar
For a \sphinxstyleemphasis{noslip} setting of the \sphinxstyleemphasis{collision} keyword, both the normal and
tangential components of the new SRD velocity are sampled from a
Gaussian distribution at temperature \sphinxstyleemphasis{Tsrd}.  Additionally, a new
tangential direction for the SRD velocity is chosen randomly.  This
collision style imparts both a normal and tangential force to the
wall.

\sphinxAtStartPar
Up to 6 walls or faces can be specified in a single command: \sphinxstyleemphasis{xlo},
\sphinxstyleemphasis{xhi}, \sphinxstyleemphasis{ylo}, \sphinxstyleemphasis{yhi}, \sphinxstyleemphasis{zlo}, \sphinxstyleemphasis{zhi}.  A \sphinxstyleemphasis{lo} face reflects particles
that move to a coordinate less than the wall position, back in the
\sphinxstyleemphasis{hi} direction.  A \sphinxstyleemphasis{hi} face reflects particles that move to a
coordinate higher than the wall position, back in the \sphinxstyleemphasis{lo} direction.

\sphinxAtStartPar
The position of each wall can be specified in one of 3 ways: as the
EDGE of the simulation box, as a constant value, or as a variable.  If
EDGE is used, then the corresponding boundary of the current
simulation box is used.  If a numeric constant is specified then the
wall is placed at that position in the appropriate dimension (x, y, or
z).  In both the EDGE and constant cases, the wall will never move.
If the wall position is a variable, it should be specified as v\_name,
where name is an {\hyperref[\detokenize{variable::doc}]{\sphinxcrossref{\DUrole{doc}{equal\sphinxhyphen{}style variable}}}} name.  In this
case the variable is evaluated each timestep and the result becomes
the current position of the reflecting wall.  Equal\sphinxhyphen{}style variables
can specify formulas with various mathematical functions, and include
{\hyperref[\detokenize{thermo_style::doc}]{\sphinxcrossref{\DUrole{doc}{thermo\_style}}}} command keywords for the simulation
box parameters and timestep and elapsed time.  Thus it is easy to
specify a time\sphinxhyphen{}dependent wall position.

\begin{sphinxadmonition}{note}{Note:}
\sphinxAtStartPar
Because the trajectory of the SRD particle is tracked as it
collides with the wall, you must ensure that r = distance of the
particle from the wall, is always \textgreater{} 0 for SRD particles, or LAMMPS
will generate an error.  This means you cannot start your simulation
with SRD particles at the wall position \sphinxstyleemphasis{coord} (r = 0) or with
particles on the wrong side of the wall (r \textless{} 0).
\end{sphinxadmonition}

\begin{sphinxadmonition}{note}{Note:}
\sphinxAtStartPar
If you have 2 or more walls that come together at an edge or
corner (e.g. walls in the x and y dimensions), then be sure to set the
\sphinxstyleemphasis{overlap} keyword to \sphinxstyleemphasis{yes} in the {\hyperref[\detokenize{fix_srd::doc}]{\sphinxcrossref{\DUrole{doc}{fix srd}}}} command,
since the walls effectively overlap when SRD particles collide with
them.  LAMMPS will issue a warning if you do not do this.
\end{sphinxadmonition}

\begin{sphinxadmonition}{note}{Note:}
\sphinxAtStartPar
The walls of this fix only interact with SRD particles, as
defined by the {\hyperref[\detokenize{fix_srd::doc}]{\sphinxcrossref{\DUrole{doc}{fix srd}}}} command.  If you are simulating
a mixture containing other kinds of particles, then you should
typically use {\hyperref[\detokenize{fix_wall::doc}]{\sphinxcrossref{\DUrole{doc}{another wall command}}}} to act on the other
particles.  Since SRD particles will be colliding both with the walls
and the other particles, it is important to ensure that the other
particle’s finite extent does not overlap an SRD wall.  If you do not
do this, you may generate errors when SRD particles end up “inside”
another particle or a wall at the beginning of a collision step.
\end{sphinxadmonition}

\sphinxAtStartPar
The \sphinxstyleemphasis{units} keyword determines the meaning of the distance units used
to define a wall position, but only when a numeric constant is used.
It is not relevant when EDGE or a variable is used to specify a face
position.

\sphinxAtStartPar
A \sphinxstyleemphasis{box} value selects standard distance units as defined by the
{\hyperref[\detokenize{units::doc}]{\sphinxcrossref{\DUrole{doc}{units}}}} command, e.g. Angstroms for units = real or metal.
A \sphinxstyleemphasis{lattice} value means the distance units are in lattice spacings.
The {\hyperref[\detokenize{lattice::doc}]{\sphinxcrossref{\DUrole{doc}{lattice}}}} command must have been previously used to
define the lattice spacings.


\bigskip\hrule\bigskip


\sphinxAtStartPar
Here are examples of variable definitions that move the wall position
in a time\sphinxhyphen{}dependent fashion using equal\sphinxhyphen{}style
{\hyperref[\detokenize{variable::doc}]{\sphinxcrossref{\DUrole{doc}{variables}}}}.

\begin{sphinxVerbatim}[commandchars=\\\{\}]
\PYG{k}{variable }\PYG{n+nv+nvIdentifier}{ramp}\PYG{+w}{ }\PYG{n}{equal}\PYG{+w}{ }\PYG{n}{ramp}\PYG{n+nv}{(}\PYG{n+nv}{0,10}\PYG{n+nv}{)}
\PYG{k}{fix }\PYG{n+nv+nvIdentifier}{1}\PYG{+w}{ }\PYG{n+nv+nvIdentifier}{all}\PYG{+w}{ }\PYG{n}{wall}\PYG{o}{/}\PYG{n}{srd}\PYG{+w}{ }\PYG{n}{xlo}\PYG{+w}{ }\PYG{n}{v\PYGZus{}ramp}

\PYG{k}{variable }\PYG{n+nv+nvIdentifier}{linear}\PYG{+w}{ }\PYG{n}{equal}\PYG{+w}{ }\PYG{n}{vdisplace}\PYG{n+nv}{(}\PYG{n+nv}{0,20}\PYG{n+nv}{)}
\PYG{k}{fix }\PYG{n+nv+nvIdentifier}{1}\PYG{+w}{ }\PYG{n+nv+nvIdentifier}{all}\PYG{+w}{ }\PYG{n}{wall}\PYG{o}{/}\PYG{n}{srd}\PYG{+w}{ }\PYG{n}{xlo}\PYG{+w}{ }\PYG{n}{v\PYGZus{}linear}

\PYG{k}{variable }\PYG{n+nv+nvIdentifier}{wiggle}\PYG{+w}{ }\PYG{n}{equal}\PYG{+w}{ }\PYG{n}{swiggle}\PYG{n+nv}{(}\PYG{n+nv}{0.0,5.0,3.0}\PYG{n+nv}{)}
\PYG{k}{fix }\PYG{n+nv+nvIdentifier}{1}\PYG{+w}{ }\PYG{n+nv+nvIdentifier}{all}\PYG{+w}{ }\PYG{n}{wall}\PYG{o}{/}\PYG{n}{srd}\PYG{+w}{ }\PYG{n}{xlo}\PYG{+w}{ }\PYG{n}{v\PYGZus{}wiggle}

\PYG{k}{variable }\PYG{n+nv+nvIdentifier}{wiggle}\PYG{+w}{ }\PYG{n}{equal}\PYG{+w}{ }\PYG{n}{cwiggle}\PYG{n+nv}{(}\PYG{n+nv}{0.0,5.0,3.0}\PYG{n+nv}{)}
\PYG{k}{fix }\PYG{n+nv+nvIdentifier}{1}\PYG{+w}{ }\PYG{n+nv+nvIdentifier}{all}\PYG{+w}{ }\PYG{n}{wall}\PYG{o}{/}\PYG{n}{srd}\PYG{+w}{ }\PYG{n}{xlo}\PYG{+w}{ }\PYG{n}{v\PYGZus{}wiggle}
\end{sphinxVerbatim}

\sphinxAtStartPar
The \sphinxstyleemphasis{ramp(lo,hi)} function adjusts the wall position linearly from \sphinxstyleemphasis{lo}
to \sphinxstyleemphasis{hi} over the course of a run.  The \sphinxstyleemphasis{vdisplace(c0,velocity)} function
does something similar using the equation \sphinxstyleemphasis{position = c0 +
velocity*delta}, where \sphinxstyleemphasis{delta} is the elapsed time.

\sphinxAtStartPar
The \sphinxstyleemphasis{swiggle(c0,A,period)} function causes the wall position to
oscillate sinusoidally according to this equation, where \sphinxstyleemphasis{omega = 2 PI
/ period}:
\begin{sphinxalltt}
position = c0 + A sin(omega*delta)
\end{sphinxalltt}

\sphinxAtStartPar
The \sphinxstyleemphasis{cwiggle(c0,A,period)} function causes the wall position to
oscillate sinusoidally according to this equation, which will have an
initial wall velocity of 0.0, and thus may impose a gentler
perturbation on the particles:
\begin{sphinxalltt}
position = c0 + A (1 \sphinxhyphen{} cos(omega*delta))
\end{sphinxalltt}


\bigskip\hrule\bigskip



\subsection{Restart, fix\_modify, output, run start/stop, minimize info}
\label{\detokenize{fix_wall_srd:restart-fix-modify-output-run-start-stop-minimize-info}}
\sphinxAtStartPar
No information about this fix is written to {\hyperref[\detokenize{restart::doc}]{\sphinxcrossref{\DUrole{doc}{binary restart files}}}}.  None of the {\hyperref[\detokenize{fix_modify::doc}]{\sphinxcrossref{\DUrole{doc}{fix\_modify}}}} options
are relevant to this fix.

\sphinxAtStartPar
This fix computes a global array of values which can be accessed by
various {\hyperref[\detokenize{Howto_output::doc}]{\sphinxcrossref{\DUrole{doc}{output commands}}}}.  The number of rows in
the array is equal to the number of walls defined by the fix.  The
number of columns is 3, for the x,y,z components of force on each
wall.

\sphinxAtStartPar
Note that an outward normal force on a wall will be a negative value
for \sphinxstyleemphasis{lo} walls and a positive value for \sphinxstyleemphasis{hi} walls.  The array values
calculated by this fix are “extensive”.

\sphinxAtStartPar
No parameter of this fix can be used with the \sphinxstyleemphasis{start/stop} keywords of
the {\hyperref[\detokenize{run::doc}]{\sphinxcrossref{\DUrole{doc}{run}}}} command.  This fix is not invoked during {\hyperref[\detokenize{minimize::doc}]{\sphinxcrossref{\DUrole{doc}{energy minimization}}}}.


\subsection{Restrictions}
\label{\detokenize{fix_wall_srd:restrictions}}
\sphinxAtStartPar
Any dimension (xyz) that has an SRD wall must be non\sphinxhyphen{}periodic.


\subsection{Related commands}
\label{\detokenize{fix_wall_srd:related-commands}}
\sphinxAtStartPar
{\hyperref[\detokenize{fix_srd::doc}]{\sphinxcrossref{\DUrole{doc}{fix srd}}}}


\subsection{Default}
\label{\detokenize{fix_wall_srd:default}}
\sphinxAtStartPar
none

\sphinxstepscope

\index{fix widom@\spxentry{fix widom}}\ignorespaces 

\section{fix widom command}
\label{\detokenize{fix_widom:fix-widom-command}}\label{\detokenize{fix_widom:index-0}}\label{\detokenize{fix_widom::doc}}

\subsection{Syntax}
\label{\detokenize{fix_widom:syntax}}
\begin{sphinxVerbatim}[commandchars=\\\{\}]
\PYG{k}{fix }\PYG{n+nv+nvIdentifier}{ID}\PYG{+w}{ }\PYG{n+nv+nvIdentifier}{group\PYGZhy{}ID}\PYG{+w}{ }\PYG{n}{widom}\PYG{+w}{ }\PYG{n}{N}\PYG{+w}{ }\PYG{n}{M}\PYG{+w}{ }\PYG{n}{type}\PYG{+w}{ }\PYG{n}{seed}\PYG{+w}{ }\PYG{n}{T}\PYG{+w}{ }\PYG{n}{keyword}\PYG{+w}{ }\PYG{n}{values}\PYG{+w}{ }\PYG{n}{...}
\end{sphinxVerbatim}
\begin{itemize}
\item {} 
\sphinxAtStartPar
ID, group\sphinxhyphen{}ID are documented in {\hyperref[\detokenize{fix::doc}]{\sphinxcrossref{\DUrole{doc}{fix}}}} command

\item {} 
\sphinxAtStartPar
widom = style name of this fix command

\item {} 
\sphinxAtStartPar
N = invoke this fix every N steps

\item {} 
\sphinxAtStartPar
M = number of Widom insertions to attempt every N steps

\item {} 
\sphinxAtStartPar
type = atom type (1\sphinxhyphen{}Ntypes or type label) for inserted atoms (must be 0 if mol keyword used)

\item {} 
\sphinxAtStartPar
seed = random \# seed (positive integer)

\item {} 
\sphinxAtStartPar
T = temperature of the system (temperature units)

\item {} 
\sphinxAtStartPar
zero or more keyword/value pairs may be appended to args
\begin{sphinxalltt}
keyword = \sphinxstyleemphasis{mol}, \sphinxstyleemphasis{region}, \sphinxstyleemphasis{full\_energy}, \sphinxstyleemphasis{charge}, \sphinxstyleemphasis{intra\_energy}
  \sphinxstyleemphasis{mol} value = template\sphinxhyphen{}ID
    template\sphinxhyphen{}ID = ID of molecule template specified in a separate {\hyperref[\detokenize{molecule::doc}]{\sphinxcrossref{\DUrole{doc}{molecule}}}} command
  \sphinxstyleemphasis{region} value = region\sphinxhyphen{}ID
    region\sphinxhyphen{}ID = ID of region where Widom insertions are allowed
  \sphinxstyleemphasis{full\_energy} = compute the entire system energy when performing Widom insertions
  \sphinxstyleemphasis{charge} value = charge of inserted atoms (charge units)
  \sphinxstyleemphasis{intra\_energy} value = intramolecular energy (energy units)
\end{sphinxalltt}

\end{itemize}


\subsection{Examples}
\label{\detokenize{fix_widom:examples}}
\begin{sphinxVerbatim}[commandchars=\\\{\}]
\PYG{k}{fix }\PYG{n+nv+nvIdentifier}{2}\PYG{+w}{ }\PYG{n+nv+nvIdentifier}{gas}\PYG{+w}{ }\PYG{n}{widom}\PYG{+w}{ }\PYG{l+m}{1}\PYG{+w}{ }\PYG{l+m}{50000}\PYG{+w}{ }\PYG{l+m}{1}\PYG{+w}{ }\PYG{l+m}{19494}\PYG{+w}{ }\PYG{l+m}{2.0}
\PYG{k}{fix }\PYG{n+nv+nvIdentifier}{3}\PYG{+w}{ }\PYG{n+nv+nvIdentifier}{water}\PYG{+w}{ }\PYG{n}{widom}\PYG{+w}{ }\PYG{l+m}{1000}\PYG{+w}{ }\PYG{l+m}{100}\PYG{+w}{ }\PYG{l+m}{0}\PYG{+w}{ }\PYG{l+m}{29494}\PYG{+w}{ }\PYG{l+m}{300.0}\PYG{+w}{ }\PYG{n}{mol}\PYG{+w}{ }\PYG{n}{h2omol}\PYG{+w}{ }\PYG{n}{full\PYGZus{}energy}

\PYG{k}{labelmap}\PYG{+w}{ }\PYG{n}{atom}\PYG{+w}{ }\PYG{l+m}{1}\PYG{+w}{ }\PYG{n}{Li}
\PYG{k}{fix }\PYG{n+nv+nvIdentifier}{2}\PYG{+w}{ }\PYG{n+nv+nvIdentifier}{ion}\PYG{+w}{ }\PYG{n}{widom}\PYG{+w}{ }\PYG{l+m}{1}\PYG{+w}{ }\PYG{l+m}{50000}\PYG{+w}{ }\PYG{n}{Li}\PYG{+w}{ }\PYG{l+m}{19494}\PYG{+w}{ }\PYG{l+m}{2.0}
\end{sphinxVerbatim}


\subsection{Description}
\label{\detokenize{fix_widom:description}}
\sphinxAtStartPar
This fix performs Widom insertions of atoms or molecules at the given
temperature as discussed in {\hyperref[\detokenize{fix_widom:frenkel1}]{\sphinxcrossref{\DUrole{std,std-ref}{(Frenkel)}}}}. Specific uses
include computation of Henry constants of small molecules in microporous
materials or amorphous systems.

\sphinxAtStartPar
Every N timesteps the fix attempts M number of Widom insertions of atoms
or molecules.

\sphinxAtStartPar
If the \sphinxstyleemphasis{mol} keyword is used, only molecule insertions are performed.
Conversely, if the \sphinxstyleemphasis{mol} keyword is not used, only atom insertions are
performed.

\sphinxAtStartPar
This command may optionally use the \sphinxstyleemphasis{region} keyword to define an
insertion volume.  The specified region must have been previously
defined with a {\hyperref[\detokenize{region::doc}]{\sphinxcrossref{\DUrole{doc}{region}}}} command.  It must be defined with
side = \sphinxstyleemphasis{in}.  Insertion attempts occur only within the specified
region. For non\sphinxhyphen{}rectangular regions, random trial points are generated
within the rectangular bounding box until a point is found that lies
inside the region. If no valid point is generated after 1000 trials, no
insertion is performed. If an attempted insertion places the atom or
molecule center\sphinxhyphen{}of\sphinxhyphen{}mass outside the specified region, a new attempted
insertion is generated. This process is repeated until the atom or
molecule center\sphinxhyphen{}of\sphinxhyphen{}mass is inside the specified region.

\sphinxAtStartPar
Note that neighbor lists are re\sphinxhyphen{}built every timestep that this fix is
invoked, so you should not set N to be too small. See the {\hyperref[\detokenize{neighbor::doc}]{\sphinxcrossref{\DUrole{doc}{neighbor}}}} command for details.

\sphinxAtStartPar
When an atom or molecule is to be inserted, its coordinates are chosen
at a random position within the current simulation cell or region.
Relative coordinates for atoms in a molecule are taken from the
template molecule provided by the user. The center of mass of the
molecule is placed at the insertion point. The orientation of the
molecule is chosen at random by rotating about this point.

\sphinxAtStartPar
Individual atoms are inserted, unless the \sphinxstyleemphasis{mol} keyword is used.  It
specifies a \sphinxstyleemphasis{template\sphinxhyphen{}ID} previously defined using the {\hyperref[\detokenize{molecule::doc}]{\sphinxcrossref{\DUrole{doc}{molecule}}}} command, which reads a file that defines the molecule.  The
coordinates, atom types, charges, etc., as well as any bonding and
special neighbor information for the molecule can be specified in the
molecule file.  See the {\hyperref[\detokenize{molecule::doc}]{\sphinxcrossref{\DUrole{doc}{molecule}}}} command for details.
The only settings required to be in this file are the coordinates and
types of atoms in the molecule.

\sphinxAtStartPar
Note that fix widom does not use configurational bias MC or any other
kind of sampling of intramolecular degrees of freedom.  Inserted
molecules can have different orientations, but they will all have the
same intramolecular configuration, which was specified in the molecule
command input.

\sphinxAtStartPar
For atoms, inserted particles have the specified atom type.  For
molecules, they use the same atom types as in the template molecule
supplied by the user.

\sphinxAtStartPar
The excess chemical potential mu\_ex is defined as:
\begin{equation*}
\begin{split}\mu_{ex} = -kT \ln(<\exp(-(U_{N+1}-U_{N})/{k_B T})>)\end{split}
\end{equation*}
\sphinxAtStartPar
where \(k_B\) is the Boltzmann constant, \(T\) is the
user\sphinxhyphen{}specified temperature, \(U_N\) and \(U_{N+1}\) is the
potential energy of the system with \(N\) and \(N+1\) particles.

\sphinxAtStartPar
The \sphinxstyleemphasis{full\_energy} option means that the fix calculates the total
potential energy of the entire simulated system, instead of just the
energy of the part that is changed. By default, this option is off, in
which case only partial energies are computed to determine the energy
difference due to the proposed change.

\sphinxAtStartPar
The \sphinxstyleemphasis{full\_energy} option is needed for systems with complicated
potential energy calculations, including the following:
\begin{itemize}
\item {} 
\sphinxAtStartPar
long\sphinxhyphen{}range electrostatics (kspace)

\item {} 
\sphinxAtStartPar
many\sphinxhyphen{}body pair styles

\item {} 
\sphinxAtStartPar
hybrid pair styles

\item {} 
\sphinxAtStartPar
eam pair styles

\item {} 
\sphinxAtStartPar
tail corrections

\item {} 
\sphinxAtStartPar
need to include potential energy contributions from other fixes

\end{itemize}

\sphinxAtStartPar
In these cases, LAMMPS will automatically apply the \sphinxstyleemphasis{full\_energy}
keyword and issue a warning message.

\sphinxAtStartPar
When the \sphinxstyleemphasis{mol} keyword is used, the \sphinxstyleemphasis{full\_energy} option also includes
the intramolecular energy of inserted and deleted molecules, whereas
this energy is not included when \sphinxstyleemphasis{full\_energy} is not used. If this is
not desired, the \sphinxstyleemphasis{intra\_energy} keyword can be used to define an amount
of energy that is subtracted from the final energy when a molecule is
inserted, and subtracted from the initial energy when a molecule is
deleted. For molecules that have a non\sphinxhyphen{}zero intramolecular energy, this
will ensure roughly the same behavior whether or not the \sphinxstyleemphasis{full\_energy}
option is used.

\sphinxAtStartPar
Some fixes have an associated potential energy. Examples of such fixes
include: {\hyperref[\detokenize{fix_efield::doc}]{\sphinxcrossref{\DUrole{doc}{efield}}}}, {\hyperref[\detokenize{fix_gravity::doc}]{\sphinxcrossref{\DUrole{doc}{gravity}}}},
{\hyperref[\detokenize{fix_addforce::doc}]{\sphinxcrossref{\DUrole{doc}{addforce}}}}, {\hyperref[\detokenize{fix_restrain::doc}]{\sphinxcrossref{\DUrole{doc}{restrain}}}}, and
{\hyperref[\detokenize{fix_wall::doc}]{\sphinxcrossref{\DUrole{doc}{wall fixes}}}}.  For that energy to be included in the
total potential energy of the system (the quantity used when performing
Widom insertions), you MUST enable the {\hyperref[\detokenize{fix_modify::doc}]{\sphinxcrossref{\DUrole{doc}{fix\_modify}}}}
\sphinxstyleemphasis{energy} option for that fix.  The doc pages for individual {\hyperref[\detokenize{fix::doc}]{\sphinxcrossref{\DUrole{doc}{fix}}}} commands specify if this should be done.

\sphinxAtStartPar
Use the \sphinxstyleemphasis{charge} option to insert atoms with a user\sphinxhyphen{}specified point
charge. Note that doing so will cause the system to become non\sphinxhyphen{}neutral.
LAMMPS issues a warning when using long\sphinxhyphen{}range electrostatics (kspace)
with non\sphinxhyphen{}neutral systems. See the {\hyperref[\detokenize{compute_group_group::doc}]{\sphinxcrossref{\DUrole{doc}{compute group/group}}}} documentation for more details about simulating
non\sphinxhyphen{}neutral systems with kspace on.


\subsection{Restart, fix\_modify, output, run start/stop, minimize info}
\label{\detokenize{fix_widom:restart-fix-modify-output-run-start-stop-minimize-info}}
\sphinxAtStartPar
This fix writes the state of the fix to {\hyperref[\detokenize{restart::doc}]{\sphinxcrossref{\DUrole{doc}{binary restart files}}}}.  This includes information about the random number
generator seed, the next timestep for Widom insertions etc.  See the
{\hyperref[\detokenize{read_restart::doc}]{\sphinxcrossref{\DUrole{doc}{read\_restart}}}} command for info on how to
re\sphinxhyphen{}specify a fix in an input script that reads a restart file, so that
the operation of the fix continues in an uninterrupted fashion.

\begin{sphinxadmonition}{note}{Note:}
\sphinxAtStartPar
For this to work correctly, the timestep must \sphinxstylestrong{not} be changed
after reading the restart with {\hyperref[\detokenize{reset_timestep::doc}]{\sphinxcrossref{\DUrole{doc}{reset\_timestep}}}}.  The fix will try to detect it and stop with an
error.
\end{sphinxadmonition}

\sphinxAtStartPar
None of the {\hyperref[\detokenize{fix_modify::doc}]{\sphinxcrossref{\DUrole{doc}{fix\_modify}}}} options are relevant to this
fix.

\sphinxAtStartPar
This fix computes a global vector of length 3, which can be accessed by
various {\hyperref[\detokenize{Howto_output::doc}]{\sphinxcrossref{\DUrole{doc}{output commands}}}}.  The vector values are
the following global cumulative quantities:
\begin{enumerate}
\sphinxsetlistlabels{\arabic}{enumi}{enumii}{}{.}%
\item {} 
\sphinxAtStartPar
average excess chemical potential on each timestep

\item {} 
\sphinxAtStartPar
average difference in potential energy on each timestep

\item {} 
\sphinxAtStartPar
volume of the insertion region

\end{enumerate}

\sphinxAtStartPar
The vector values calculated by this fix are “intensive”.

\sphinxAtStartPar
No parameter of this fix can be used with the \sphinxstyleemphasis{start/stop} keywords of
the {\hyperref[\detokenize{run::doc}]{\sphinxcrossref{\DUrole{doc}{run}}}} command.  This fix is not invoked during
{\hyperref[\detokenize{minimize::doc}]{\sphinxcrossref{\DUrole{doc}{energy minimization}}}}.


\subsection{Restrictions}
\label{\detokenize{fix_widom:restrictions}}
\sphinxAtStartPar
This fix is part of the MC package.  It is only enabled if LAMMPS was
built with that package.  See the {\hyperref[\detokenize{Build_package::doc}]{\sphinxcrossref{\DUrole{doc}{Build package}}}}
doc page for more info.

\sphinxAtStartPar
Do not set “neigh\_modify once yes” or else this fix will never be
called.  Reneighboring is \sphinxstylestrong{required}.

\sphinxAtStartPar
This fix style requires an {\hyperref[\detokenize{atom_style::doc}]{\sphinxcrossref{\DUrole{doc}{atom style}}}} with per atom
type masses.

\sphinxAtStartPar
Can be run in parallel, but some aspects of the insertion procedure
will not scale well in parallel. Only usable for 3D simulations.


\subsection{Related commands}
\label{\detokenize{fix_widom:related-commands}}
\sphinxAtStartPar
{\hyperref[\detokenize{fix_gcmc::doc}]{\sphinxcrossref{\DUrole{doc}{fix gcmc}}}}
{\hyperref[\detokenize{fix_atom_swap::doc}]{\sphinxcrossref{\DUrole{doc}{fix atom/swap}}}},
{\hyperref[\detokenize{neighbor::doc}]{\sphinxcrossref{\DUrole{doc}{neighbor}}}},
{\hyperref[\detokenize{fix_deposit::doc}]{\sphinxcrossref{\DUrole{doc}{fix deposit}}}}, {\hyperref[\detokenize{fix_evaporate::doc}]{\sphinxcrossref{\DUrole{doc}{fix evaporate}}}},


\subsection{Default}
\label{\detokenize{fix_widom:default}}
\sphinxAtStartPar
The option defaults are mol = no, intra\_energy = 0.0 and full\_energy =
no, except for the situations where full\_energy is required, as listed
above.


\bigskip\hrule\bigskip

\phantomsection\label{\detokenize{fix_widom:frenkel1}}
\sphinxAtStartPar
\sphinxstylestrong{(Frenkel)} Frenkel and Smit, Understanding Molecular Simulation,
Academic Press, London, 2002.

\sphinxstepscope


\chapter{Computes}
\label{\detokenize{computes:computes}}\label{\detokenize{computes::doc}}
\sphinxstepscope

\index{compute ackland/atom@\spxentry{compute ackland/atom}}\ignorespaces 

\section{compute ackland/atom command}
\label{\detokenize{compute_ackland_atom:compute-ackland-atom-command}}\label{\detokenize{compute_ackland_atom:index-0}}\label{\detokenize{compute_ackland_atom::doc}}

\subsection{Syntax}
\label{\detokenize{compute_ackland_atom:syntax}}
\begin{sphinxVerbatim}[commandchars=\\\{\}]
\PYG{k}{compute }\PYG{n+nv+nvIdentifier}{ID}\PYG{+w}{ }\PYG{n+nv+nvIdentifier}{group\PYGZhy{}ID}\PYG{+w}{ }\PYG{n}{ackland}\PYG{o}{/}\PYG{n}{atom}\PYG{+w}{ }\PYG{n}{keyword}\PYG{o}{/}\PYG{n}{value}
\end{sphinxVerbatim}
\begin{itemize}
\item {} 
\sphinxAtStartPar
ID, group\sphinxhyphen{}ID are documented in {\hyperref[\detokenize{compute::doc}]{\sphinxcrossref{\DUrole{doc}{compute}}}} command

\item {} 
\sphinxAtStartPar
ackland/atom = style name of this compute command

\item {} 
\sphinxAtStartPar
zero or more keyword/value pairs may be appended

\item {} 
\sphinxAtStartPar
keyword = \sphinxstyleemphasis{legacy}
\begin{sphinxalltt}
\sphinxstyleemphasis{legacy} args = \sphinxstyleemphasis{yes} or \sphinxstyleemphasis{no} = use (\sphinxstyleemphasis{yes}) or do not use (\sphinxstyleemphasis{no}) legacy Ackland algorithm implementation
\end{sphinxalltt}

\end{itemize}


\subsection{Examples}
\label{\detokenize{compute_ackland_atom:examples}}
\begin{sphinxVerbatim}[commandchars=\\\{\}]
\PYG{k}{compute }\PYG{n+nv+nvIdentifier}{1}\PYG{+w}{ }\PYG{n+nv+nvIdentifier}{all}\PYG{+w}{ }\PYG{n}{ackland}\PYG{o}{/}\PYG{n}{atom}
\PYG{k}{compute }\PYG{n+nv+nvIdentifier}{1}\PYG{+w}{ }\PYG{n+nv+nvIdentifier}{all}\PYG{+w}{ }\PYG{n}{ackland}\PYG{o}{/}\PYG{n}{atom}\PYG{+w}{ }\PYG{n}{legacy}\PYG{+w}{ }\PYG{n}{yes}
\end{sphinxVerbatim}


\subsection{Description}
\label{\detokenize{compute_ackland_atom:description}}
\sphinxAtStartPar
Defines a computation that calculates the local lattice structure
according to the formulation given in {\hyperref[\detokenize{compute_ackland_atom:ackland}]{\sphinxcrossref{\DUrole{std,std-ref}{(Ackland)}}}}.
Historically, LAMMPS had two, slightly different implementations of
the algorithm from the paper. With the \sphinxstyleemphasis{legacy} keyword, it is
possible to switch between the pre\sphinxhyphen{}2015 (\sphinxstyleemphasis{legacy yes}) and post\sphinxhyphen{}2015
implementation (\sphinxstyleemphasis{legacy no}). The post\sphinxhyphen{}2015 variant is the default.

\sphinxAtStartPar
In contrast to the {\hyperref[\detokenize{compute_centro_atom::doc}]{\sphinxcrossref{\DUrole{doc}{centro\sphinxhyphen{}symmetry parameter}}}} this method is stable against
temperature boost, because it is based not on the distance between
particles but the angles.  Therefore statistical fluctuations are
averaged out a little more.  A comparison with the Common Neighbor
Analysis metric is made in the paper.

\sphinxAtStartPar
The result is a number which is mapped to the following different
lattice structures:
\begin{itemize}
\item {} 
\sphinxAtStartPar
0 = UNKNOWN

\item {} 
\sphinxAtStartPar
1 = BCC

\item {} 
\sphinxAtStartPar
2 = FCC

\item {} 
\sphinxAtStartPar
3 = HCP

\item {} 
\sphinxAtStartPar
4 = ICO

\end{itemize}

\sphinxAtStartPar
The neighbor list needed to compute this quantity is constructed each
time the calculation is performed (i.e. each time a snapshot of atoms
is dumped).  Thus it can be inefficient to compute/dump this quantity
too frequently or to have multiple compute/dump commands, each of
which computes this quantity.\sphinxhyphen{}


\subsection{Output info}
\label{\detokenize{compute_ackland_atom:output-info}}
\sphinxAtStartPar
This compute calculates a per\sphinxhyphen{}atom vector, which can be accessed by
any command that uses per\sphinxhyphen{}atom values from a compute as input.  See
the {\hyperref[\detokenize{Howto_output::doc}]{\sphinxcrossref{\DUrole{doc}{Howto output}}}} page for an overview of
LAMMPS output options.


\subsection{Restrictions}
\label{\detokenize{compute_ackland_atom:restrictions}}
\sphinxAtStartPar
This compute is part of the EXTRA\sphinxhyphen{}COMPUTE package.  It is only enabled if
LAMMPS was built with that package.  See the {\hyperref[\detokenize{Build_package::doc}]{\sphinxcrossref{\DUrole{doc}{Build package}}}} page for more info.

\sphinxAtStartPar
The per\sphinxhyphen{}atom vector values will be unitless since they are the
integers defined above.


\subsection{Related commands}
\label{\detokenize{compute_ackland_atom:related-commands}}
\sphinxAtStartPar
{\hyperref[\detokenize{compute_centro_atom::doc}]{\sphinxcrossref{\DUrole{doc}{compute centro/atom}}}}


\subsection{Default}
\label{\detokenize{compute_ackland_atom:default}}
\sphinxAtStartPar
The keyword \sphinxstyleemphasis{legacy} defaults to \sphinxstyleemphasis{no}.


\bigskip\hrule\bigskip

\phantomsection\label{\detokenize{compute_ackland_atom:ackland}}
\sphinxAtStartPar
\sphinxstylestrong{(Ackland)} Ackland, Jones, Phys Rev B, 73, 054104 (2006).

\sphinxstepscope

\index{compute adf@\spxentry{compute adf}}\ignorespaces 

\section{compute adf command}
\label{\detokenize{compute_adf:compute-adf-command}}\label{\detokenize{compute_adf:index-0}}\label{\detokenize{compute_adf::doc}}

\subsection{Syntax}
\label{\detokenize{compute_adf:syntax}}
\begin{sphinxVerbatim}[commandchars=\\\{\}]
\PYG{k}{compute }\PYG{n+nv+nvIdentifier}{ID}\PYG{+w}{ }\PYG{n+nv+nvIdentifier}{group\PYGZhy{}ID}\PYG{+w}{ }\PYG{n}{adf}\PYG{+w}{ }\PYG{n}{Nbin}\PYG{+w}{ }\PYG{n}{itype1}\PYG{+w}{ }\PYG{n}{jtype1}\PYG{+w}{ }\PYG{n}{ktype1}\PYG{+w}{ }\PYG{n}{Rjinner1}\PYG{+w}{ }\PYG{n}{Rjouter1}\PYG{+w}{ }\PYG{n}{Rkinner1}\PYG{+w}{ }\PYG{n}{Rkouter1}\PYG{+w}{ }\PYG{n}{...}
\end{sphinxVerbatim}
\begin{itemize}
\item {} 
\sphinxAtStartPar
ID, group\sphinxhyphen{}ID are documented in {\hyperref[\detokenize{compute::doc}]{\sphinxcrossref{\DUrole{doc}{compute}}}} command

\item {} 
\sphinxAtStartPar
adf = style name of this compute command

\item {} 
\sphinxAtStartPar
Nbin = number of ADF bins

\item {} 
\sphinxAtStartPar
itypeN = central atom type for Nth ADF histogram (see asterisk form below)

\item {} 
\sphinxAtStartPar
jtypeN = J atom type for Nth ADF histogram (see asterisk form below)

\item {} 
\sphinxAtStartPar
ktypeN = K atom type for Nth ADF histogram (see asterisk form below)

\item {} 
\sphinxAtStartPar
RjinnerN = inner radius of J atom shell for Nth ADF histogram (distance units)

\item {} 
\sphinxAtStartPar
RjouterN = outer radius of J atom shell for Nth ADF histogram (distance units)

\item {} 
\sphinxAtStartPar
RkinnerN = inner radius of K atom shell for Nth ADF histogram (distance units)

\item {} 
\sphinxAtStartPar
RkouterN = outer radius of K atom shell for Nth ADF histogram (distance units)

\item {} 
\sphinxAtStartPar
zero or one keyword/value pairs may be appended

\item {} 
\sphinxAtStartPar
keyword = \sphinxstyleemphasis{ordinate}
\begin{sphinxalltt}
\sphinxstyleemphasis{ordinate} value = \sphinxstyleemphasis{degree} or \sphinxstyleemphasis{radian} or \sphinxstyleemphasis{cosine}
  Choose the ordinate parameter for the histogram
\end{sphinxalltt}

\end{itemize}


\subsection{Examples}
\label{\detokenize{compute_adf:examples}}
\begin{sphinxVerbatim}[commandchars=\\\{\}]
\PYG{k}{compute }\PYG{n+nv+nvIdentifier}{1}\PYG{+w}{ }\PYG{n+nv+nvIdentifier}{fluid}\PYG{+w}{ }\PYG{n}{adf}\PYG{+w}{ }\PYG{l+m}{32}\PYG{+w}{ }\PYG{l+m}{1}\PYG{+w}{ }\PYG{l+m}{1}\PYG{+w}{ }\PYG{l+m}{1}\PYG{+w}{ }\PYG{l+m}{0.0}\PYG{+w}{ }\PYG{l+m}{1.2}\PYG{+w}{ }\PYG{l+m}{0.0}\PYG{+w}{ }\PYG{l+m}{1.2}\PYG{+w}{ }\PYG{o}{\PYGZam{}}
\PYG{+w}{                       }\PYG{l+m}{1}\PYG{+w}{ }\PYG{l+m}{1}\PYG{+w}{ }\PYG{l+m}{2}\PYG{+w}{ }\PYG{l+m}{0.0}\PYG{+w}{ }\PYG{l+m}{1.2}\PYG{+w}{ }\PYG{l+m}{0.0}\PYG{+w}{ }\PYG{l+m}{1.5}\PYG{+w}{ }\PYG{o}{\PYGZam{}}
\PYG{+w}{                       }\PYG{l+m}{1}\PYG{+w}{ }\PYG{l+m}{2}\PYG{+w}{ }\PYG{l+m}{2}\PYG{+w}{ }\PYG{l+m}{0.0}\PYG{+w}{ }\PYG{l+m}{1.5}\PYG{+w}{ }\PYG{l+m}{0.0}\PYG{+w}{ }\PYG{l+m}{1.5}\PYG{+w}{ }\PYG{o}{\PYGZam{}}
\PYG{+w}{                       }\PYG{l+m}{2}\PYG{+w}{ }\PYG{l+m}{1}\PYG{+w}{ }\PYG{l+m}{1}\PYG{+w}{ }\PYG{l+m}{0.0}\PYG{+w}{ }\PYG{l+m}{1.2}\PYG{+w}{ }\PYG{l+m}{0.0}\PYG{+w}{ }\PYG{l+m}{1.2}\PYG{+w}{ }\PYG{o}{\PYGZam{}}
\PYG{+w}{                       }\PYG{l+m}{2}\PYG{+w}{ }\PYG{l+m}{1}\PYG{+w}{ }\PYG{l+m}{2}\PYG{+w}{ }\PYG{l+m}{0.0}\PYG{+w}{ }\PYG{l+m}{1.5}\PYG{+w}{ }\PYG{l+m}{2.0}\PYG{+w}{ }\PYG{l+m}{3.5}\PYG{+w}{ }\PYG{o}{\PYGZam{}}
\PYG{+w}{                       }\PYG{l+m}{2}\PYG{+w}{ }\PYG{l+m}{2}\PYG{+w}{ }\PYG{l+m}{2}\PYG{+w}{ }\PYG{l+m}{2.0}\PYG{+w}{ }\PYG{l+m}{3.5}\PYG{+w}{ }\PYG{l+m}{2.0}\PYG{+w}{ }\PYG{l+m}{3.5}
\PYG{k}{compute }\PYG{n+nv+nvIdentifier}{1}\PYG{+w}{ }\PYG{n+nv+nvIdentifier}{fluid}\PYG{+w}{ }\PYG{n}{adf}\PYG{+w}{ }\PYG{l+m}{32}\PYG{+w}{ }\PYG{l+m}{1}\PYG{o}{*}\PYG{l+m}{2}\PYG{+w}{ }\PYG{l+m}{1}\PYG{o}{*}\PYG{l+m}{2}\PYG{+w}{ }\PYG{l+m}{1}\PYG{o}{*}\PYG{l+m}{2}\PYG{+w}{ }\PYG{l+m}{0.5}\PYG{+w}{ }\PYG{l+m}{3.5}
\PYG{k}{compute }\PYG{n+nv+nvIdentifier}{1}\PYG{+w}{ }\PYG{n+nv+nvIdentifier}{fluid}\PYG{+w}{ }\PYG{n}{adf}\PYG{+w}{ }\PYG{l+m}{32}
\end{sphinxVerbatim}


\subsection{Description}
\label{\detokenize{compute_adf:description}}
\sphinxAtStartPar
Define a computation that calculates one or more angular distribution functions
(ADF) for a group of particles.  Each ADF is calculated in histogram form
by measuring the angle formed by a central atom and two neighbor atoms and
binning these angles into \sphinxstyleemphasis{Nbin} bins.
Only neighbors for which \sphinxstyleemphasis{Rinner} \textless{} \sphinxstyleemphasis{R} \textless{} \sphinxstyleemphasis{Router} are counted, where
\sphinxstyleemphasis{Rinner} and \sphinxstyleemphasis{Router} are specified separately for the first and second
neighbor atom in each requested ADF.

\begin{sphinxadmonition}{note}{Note:}
\sphinxAtStartPar
If you have a bonded system, then the settings of
{\hyperref[\detokenize{special_bonds::doc}]{\sphinxcrossref{\DUrole{doc}{special\_bonds}}}} command can remove pairwise
interactions between atoms in the same bond, angle, or dihedral.  This
is the default setting for the {\hyperref[\detokenize{special_bonds::doc}]{\sphinxcrossref{\DUrole{doc}{special\_bonds}}}}
command, and means those pairwise interactions do not appear in the
neighbor list.  Because this fix uses a neighbor list, it also means
those pairs will not be included in the ADF. This does not apply when
using long\sphinxhyphen{}range coulomb interactions (\sphinxstyleemphasis{coul/long}, \sphinxstyleemphasis{coul/msm},
\sphinxstyleemphasis{coul/wolf} or similar.  One way to get around this would be to set
special\_bond scaling factors to very tiny numbers that are not exactly
zero (e.g. 1.0e\sphinxhyphen{}50). Another workaround is to write a dump file, and
use the {\hyperref[\detokenize{rerun::doc}]{\sphinxcrossref{\DUrole{doc}{rerun}}}} command to compute the ADF for snapshots in
the dump file.  The rerun script can use a
{\hyperref[\detokenize{special_bonds::doc}]{\sphinxcrossref{\DUrole{doc}{special\_bonds}}}} command that includes all pairs in
the neighbor list.
\end{sphinxadmonition}

\begin{sphinxadmonition}{note}{Note:}
\sphinxAtStartPar
If you request any outer cutoff \sphinxstyleemphasis{Router} \textgreater{} force cutoff, or if no
pair style is defined,  e.g. the {\hyperref[\detokenize{rerun::doc}]{\sphinxcrossref{\DUrole{doc}{rerun}}}} command is being used to
post\sphinxhyphen{}process a dump file of snapshots you must ensure ghost atom information
out to the largest value of \sphinxstyleemphasis{Router} + \sphinxstyleemphasis{skin} is communicated, via the
{\hyperref[\detokenize{comm_modify::doc}]{\sphinxcrossref{\DUrole{doc}{comm\_modify cutoff}}}} command, else the ADF computation
cannot be performed, and LAMMPS will give an error message.  The \sphinxstyleemphasis{skin} value
is what is specified with the {\hyperref[\detokenize{neighbor::doc}]{\sphinxcrossref{\DUrole{doc}{neighbor}}}} command.
\end{sphinxadmonition}

\sphinxAtStartPar
The \sphinxstyleemphasis{itypeN},\sphinxstyleemphasis{jtypeN},\sphinxstyleemphasis{ktypeN} settings can be specified in one of two
ways.  An explicit numeric value can be used, as in the first example
above.  Or a wild\sphinxhyphen{}card asterisk can be used to specify a range of atom
types as in the second example above.
This takes the form “*” or “*n” or “n*” or “m*n”.  If N = the
number of atom types, then an asterisk with no numeric values means
all types from 1 to N.  A leading asterisk means all types from 1 to n
(inclusive).  A trailing asterisk means all types from n to N
(inclusive).  A middle asterisk means all types from m to n
(inclusive).

\sphinxAtStartPar
If \sphinxstyleemphasis{itypeN}, \sphinxstyleemphasis{jtypeN}, and \sphinxstyleemphasis{ktypeN} are single values, as in the first example
above, this means that the ADF is computed where atoms of type \sphinxstyleemphasis{itypeN}
are the central atom, and neighbor atoms of type \sphinxstyleemphasis{jtypeN} and \sphinxstyleemphasis{ktypeN}
are forming the angle.  If any of \sphinxstyleemphasis{itypeN}, \sphinxstyleemphasis{jtypeN}, or \sphinxstyleemphasis{ktypeN}
represent a range of values via
the wild\sphinxhyphen{}card asterisk, as in the second example above, this means that the
ADF is computed where atoms of any of the range of types represented
by \sphinxstyleemphasis{itypeN} are the central atom, and the angle is formed by two neighbors,
one neighbor in the range of types represented by \sphinxstyleemphasis{jtypeN} and another neighbor
in the range of types represented by \sphinxstyleemphasis{ktypeN}.

\sphinxAtStartPar
If no \sphinxstyleemphasis{itypeN}, \sphinxstyleemphasis{jtypeN}, \sphinxstyleemphasis{ktypeN} settings are specified, then
LAMMPS will generate a single ADF for all atoms in the group.
The inner cutoff is set to zero and the outer cutoff is set
to the force cutoff. If no pair\_style is specified, there is no
force cutoff and LAMMPS will give an error message. Note that
in most cases, generating an ADF for all atoms is not a good thing.
Such an ADF is both uninformative and
extremely expensive to compute.  For example, with liquid water
with a 10 A force cutoff, there are 80,000 angles per atom.
In addition, most of the interesting angular structure occurs for
neighbors that are the closest to the central atom, involving
just a few dozen angles.

\sphinxAtStartPar
Angles for each ADF are generated by double\sphinxhyphen{}looping over the list of
neighbors of each central atom I,
just as they would be in the force calculation for
a three\sphinxhyphen{}body potential such as {\hyperref[\detokenize{pair_sw::doc}]{\sphinxcrossref{\DUrole{doc}{Stillinger\sphinxhyphen{}Weber}}}}.
The angle formed by central atom I and neighbor atoms J and K is included in an
ADF if the following criteria are met:
\begin{itemize}
\item {} 
\sphinxAtStartPar
atoms I,J,K are all in the specified compute group

\item {} 
\sphinxAtStartPar
the distance between atoms I,J is between Rjinner and Rjouter

\item {} 
\sphinxAtStartPar
the distance between atoms I,K is between Rkinner and Rkouter

\item {} 
\sphinxAtStartPar
the type of the I atom matches itypeN (one or a range of types)

\item {} 
\sphinxAtStartPar
atoms I,J,K are distinct

\item {} 
\sphinxAtStartPar
the type of the J atom matches jtypeN (one or a range of types)

\item {} 
\sphinxAtStartPar
the type of the K atom matches ktypeN (one or a range of types)

\end{itemize}

\sphinxAtStartPar
Each unique angle satisfying the above criteria is counted only once, regardless
of whether either or both of the neighbor atoms making up the
angle appear in both the J and K lists.
It is OK if a particular angle is included in more than
one individual histogram, due to the way the \sphinxstyleemphasis{itypeN}, \sphinxstyleemphasis{jtypeN}, \sphinxstyleemphasis{ktypeN}
arguments are specified.

\sphinxAtStartPar
The first ADF value for a bin is calculated from the histogram count by
dividing by the total number of triples satisfying the criteria,
so that the integral of the ADF w.r.t. angle is 1, i.e. the ADF
is a probability density function.

\sphinxAtStartPar
The second ADF value is reported as a cumulative sum of
all bins up to the current bins, averaged
over atoms of type \sphinxstyleemphasis{itypeN}. It represents the
number of angles per central atom with angle less
than or equal to the angle of the current bin,
analogous to the coordination
number radial distribution function.

\sphinxAtStartPar
The \sphinxstyleemphasis{ordinate} optional keyword determines
whether the bins are of uniform angular size from zero
to 180 (\sphinxstyleemphasis{degree}), zero to Pi (\sphinxstyleemphasis{radian}), or the
cosine of the angle uniform in the range {[}\sphinxhyphen{}1,1{]} (\sphinxstyleemphasis{cosine}).
\sphinxstyleemphasis{cosine} has the advantage of eliminating the \sphinxstyleemphasis{acos()} function
call, which speeds up the compute by 2\sphinxhyphen{}3x, and it is also preferred
on physical grounds, because the for uniformly distributed particles
in 3D, the angular probability density w.r.t dtheta is
sin(theta)/2, while for d(cos(theta)), it is 1/2,
Regardless of which ordinate is chosen, the first column of ADF
values is normalized w.r.t. the range of that ordinate, so that
the integral is 1.

\sphinxAtStartPar
The simplest way to output the results of the compute adf calculation
to a file is to use the {\hyperref[\detokenize{fix_ave_time::doc}]{\sphinxcrossref{\DUrole{doc}{fix ave/time}}}} command, for
example:

\begin{sphinxVerbatim}[commandchars=\\\{\}]
\PYG{k}{compute }\PYG{n+nv+nvIdentifier}{myADF}\PYG{+w}{ }\PYG{n+nv+nvIdentifier}{all}\PYG{+w}{ }\PYG{n}{adf}\PYG{+w}{ }\PYG{l+m}{32}\PYG{+w}{ }\PYG{l+m}{2}\PYG{+w}{ }\PYG{l+m}{2}\PYG{+w}{ }\PYG{l+m}{2}\PYG{+w}{ }\PYG{l+m}{0.5}\PYG{+w}{ }\PYG{l+m}{3.5}\PYG{+w}{ }\PYG{l+m}{0.5}\PYG{+w}{ }\PYG{l+m}{3.5}
\PYG{k}{fix }\PYG{n+nv+nvIdentifier}{1}\PYG{+w}{ }\PYG{n+nv+nvIdentifier}{all}\PYG{+w}{ }\PYG{n}{ave}\PYG{o}{/}\PYG{n}{time}\PYG{+w}{ }\PYG{l+m}{100}\PYG{+w}{ }\PYG{l+m}{1}\PYG{+w}{ }\PYG{l+m}{100}\PYG{+w}{ }\PYG{n}{c\PYGZus{}myADF[}\PYG{o}{*}\PYG{n}{]}\PYG{+w}{ }\PYG{n}{file}\PYG{+w}{ }\PYG{n}{tmp.adf}\PYG{+w}{ }\PYG{n}{mode}\PYG{+w}{ }\PYG{n}{vector}
\end{sphinxVerbatim}


\subsection{Output info}
\label{\detokenize{compute_adf:output-info}}
\sphinxAtStartPar
This compute calculates a global array with the number of rows =
\sphinxstyleemphasis{Nbins} and the number of columns = \(1 + 2 \times\) \sphinxstyleemphasis{Ntriples}, where \sphinxstyleemphasis{Ntriples}
is the number of I,J,K triples specified.  The first column has the bin
coordinate (angle\sphinxhyphen{}related ordinate at midpoint of bin). Each subsequent column has
the two ADF values for a specific set of (\sphinxstyleemphasis{itypeN},\sphinxstyleemphasis{jtypeN},\sphinxstyleemphasis{ktypeN})
interactions, as described above.  These values can be used
by any command that uses a global values from a compute as input.  See
the {\hyperref[\detokenize{Howto_output::doc}]{\sphinxcrossref{\DUrole{doc}{Howto output}}}} page for an overview of
LAMMPS output options.

\sphinxAtStartPar
The array values calculated by this compute are all “intensive”.

\sphinxAtStartPar
The first column of array values is the angle\sphinxhyphen{}related ordinate, either
the angle in degrees or radians, or the cosine of the angle.  Each
subsequent pair of columns gives the first and second kinds of ADF
for a specific set of (\sphinxstyleemphasis{itypeN},\sphinxstyleemphasis{jtypeN},\sphinxstyleemphasis{ktypeN}). The values
in the first ADF column are normalized numbers \(\ge 0.0\),
whose integral w.r.t. the ordinate is 1,
i.e. the first ADF is a normalized probability distribution.
The values in the second ADF column are also numbers \(\ge 0.0\).
They are the cumulative density distribution of angles per atom.
By definition, this ADF is monotonically increasing from zero to
a maximum value equal to the average total number of
angles per atom satisfying the ADF criteria.


\subsection{Restrictions}
\label{\detokenize{compute_adf:restrictions}}
\sphinxAtStartPar
This compute is part of the EXTRA\sphinxhyphen{}COMPUTE package.  It is only enabled
if LAMMPS was built with that package.  See the {\hyperref[\detokenize{Build_package::doc}]{\sphinxcrossref{\DUrole{doc}{Build package}}}} page for more info.

\sphinxAtStartPar
By default, the ADF is not computed for distances longer than the
largest force cutoff, since the neighbor list creation will only contain
pairs up to that distance (plus neighbor list skin).  If you use outer
cutoffs larger than that, you must use {\hyperref[\detokenize{neighbor::doc}]{\sphinxcrossref{\DUrole{doc}{neighbor style ‘bin’ or
‘nsq’}}}}.

\sphinxAtStartPar
If you want an ADF for a larger outer cutoff, you can also use the
{\hyperref[\detokenize{rerun::doc}]{\sphinxcrossref{\DUrole{doc}{rerun}}}} command to post\sphinxhyphen{}process a dump file, use {\hyperref[\detokenize{pair_zero::doc}]{\sphinxcrossref{\DUrole{doc}{pair
style zero}}}} and set the force cutoff to be larger in the
rerun script.  Note that in the rerun context, the force cutoff is
arbitrary and with pair style zero you are not computing any forces, and
since you are not running dynamics you are not changing the model that
generated the trajectory.

\sphinxAtStartPar
The ADF is not computed for neighbors outside the force cutoff,
since processors (in parallel) don’t know about atom coordinates for
atoms further away than that distance.  If you want an ADF for larger
distances, you can use the {\hyperref[\detokenize{rerun::doc}]{\sphinxcrossref{\DUrole{doc}{rerun}}}} command to post\sphinxhyphen{}process
a dump file and set the cutoff for the potential to be longer in the
rerun script.  Note that in the rerun context, the force cutoff is
arbitrary, since you are not running dynamics and thus are not changing
your model.


\subsection{Related commands}
\label{\detokenize{compute_adf:related-commands}}
\sphinxAtStartPar
{\hyperref[\detokenize{compute_rdf::doc}]{\sphinxcrossref{\DUrole{doc}{compute rdf}}}}, {\hyperref[\detokenize{fix_ave_time::doc}]{\sphinxcrossref{\DUrole{doc}{fix ave/time}}}}, {\hyperref[\detokenize{compute_modify::doc}]{\sphinxcrossref{\DUrole{doc}{compute\_modify}}}}


\subsection{Default}
\label{\detokenize{compute_adf:default}}
\sphinxAtStartPar
The keyword default is ordinate = degree.

\sphinxstepscope

\index{compute angle@\spxentry{compute angle}}\ignorespaces 

\section{compute angle command}
\label{\detokenize{compute_angle:compute-angle-command}}\label{\detokenize{compute_angle:index-0}}\label{\detokenize{compute_angle::doc}}

\subsection{Syntax}
\label{\detokenize{compute_angle:syntax}}
\begin{sphinxVerbatim}[commandchars=\\\{\}]
\PYG{k}{compute }\PYG{n+nv+nvIdentifier}{ID}\PYG{+w}{ }\PYG{n+nv+nvIdentifier}{group\PYGZhy{}ID}\PYG{+w}{ }\PYG{n}{angle}
\end{sphinxVerbatim}
\begin{itemize}
\item {} 
\sphinxAtStartPar
ID, group\sphinxhyphen{}ID are documented in {\hyperref[\detokenize{compute::doc}]{\sphinxcrossref{\DUrole{doc}{compute}}}} command

\item {} 
\sphinxAtStartPar
angle = style name of this compute command

\end{itemize}


\subsection{Examples}
\label{\detokenize{compute_angle:examples}}
\begin{sphinxVerbatim}[commandchars=\\\{\}]
\PYG{k}{compute }\PYG{n+nv+nvIdentifier}{1}\PYG{+w}{ }\PYG{n+nv+nvIdentifier}{all}\PYG{+w}{ }\PYG{n}{angle}
\end{sphinxVerbatim}


\subsection{Description}
\label{\detokenize{compute_angle:description}}
\sphinxAtStartPar
Define a computation that extracts the angle energy calculated by each of the
angle sub\sphinxhyphen{}styles used in the {\hyperref[\detokenize{angle_hybrid::doc}]{\sphinxcrossref{\DUrole{doc}{angle\_style hybrid}}}} command.
These values are made accessible for output or further processing by other
commands.  The group specified for this command is ignored.

\sphinxAtStartPar
This compute is useful when using {\hyperref[\detokenize{angle_hybrid::doc}]{\sphinxcrossref{\DUrole{doc}{angle\_style hybrid}}}} if
you want to know the portion of the total energy contributed by one or more of
the hybrid sub\sphinxhyphen{}styles.


\subsection{Output info}
\label{\detokenize{compute_angle:output-info}}
\sphinxAtStartPar
This compute calculates a global vector of length \sphinxstyleemphasis{N}, where \sphinxstyleemphasis{N} is the number
of sub\_styles defined by the {\hyperref[\detokenize{angle_style::doc}]{\sphinxcrossref{\DUrole{doc}{angle\_style hybrid}}}} command,
which can be accessed by indices 1 through \sphinxstyleemphasis{N}.  These values can be used by
any command that uses global scalar or vector values from a compute as input.
See the {\hyperref[\detokenize{Howto_output::doc}]{\sphinxcrossref{\DUrole{doc}{Howto output}}}} page for an overview of LAMMPS
output options.

\sphinxAtStartPar
The vector values are “extensive” and will be in energy
{\hyperref[\detokenize{units::doc}]{\sphinxcrossref{\DUrole{doc}{units}}}}.


\subsection{Restrictions}
\label{\detokenize{compute_angle:restrictions}}
\sphinxAtStartPar
none


\subsection{Related commands}
\label{\detokenize{compute_angle:related-commands}}
\sphinxAtStartPar
{\hyperref[\detokenize{compute_pe::doc}]{\sphinxcrossref{\DUrole{doc}{compute pe}}}}, {\hyperref[\detokenize{compute_pair::doc}]{\sphinxcrossref{\DUrole{doc}{compute pair}}}}


\subsection{Default}
\label{\detokenize{compute_angle:default}}
\sphinxAtStartPar
none

\sphinxstepscope

\index{compute angle/local@\spxentry{compute angle/local}}\ignorespaces 

\section{compute angle/local command}
\label{\detokenize{compute_angle_local:compute-angle-local-command}}\label{\detokenize{compute_angle_local:index-0}}\label{\detokenize{compute_angle_local::doc}}

\subsection{Syntax}
\label{\detokenize{compute_angle_local:syntax}}
\begin{sphinxVerbatim}[commandchars=\\\{\}]
\PYG{k}{compute }\PYG{n+nv+nvIdentifier}{ID}\PYG{+w}{ }\PYG{n+nv+nvIdentifier}{group\PYGZhy{}ID}\PYG{+w}{ }\PYG{n}{angle}\PYG{o}{/}\PYG{n}{local}\PYG{+w}{ }\PYG{n}{value1}\PYG{+w}{ }\PYG{n}{value2}\PYG{+w}{ }\PYG{n}{...}\PYG{+w}{ }\PYG{n}{keyword}\PYG{+w}{ }\PYG{n}{args}\PYG{+w}{ }\PYG{n}{...}
\end{sphinxVerbatim}
\begin{itemize}
\item {} 
\sphinxAtStartPar
ID, group\sphinxhyphen{}ID are documented in {\hyperref[\detokenize{compute::doc}]{\sphinxcrossref{\DUrole{doc}{compute}}}} command

\item {} 
\sphinxAtStartPar
angle/local = style name of this compute command

\item {} 
\sphinxAtStartPar
one or more values may be appended

\item {} 
\sphinxAtStartPar
value = \sphinxstyleemphasis{theta} or \sphinxstyleemphasis{eng} or \sphinxstyleemphasis{v\_name}
\begin{sphinxalltt}
\sphinxstyleemphasis{theta} = tabulate angles
\sphinxstyleemphasis{eng} = tabulate angle energies
\sphinxstyleemphasis{v\_name} = equal\sphinxhyphen{}style variable with name (see below)
\end{sphinxalltt}

\item {} 
\sphinxAtStartPar
zero or more keyword/args pairs may be appended

\item {} 
\sphinxAtStartPar
keyword = \sphinxstyleemphasis{set}
\begin{sphinxalltt}
\sphinxstyleemphasis{set} args = theta name
  theta = only currently allowed arg
  name = name of variable to set with theta
\end{sphinxalltt}

\end{itemize}


\subsection{Examples}
\label{\detokenize{compute_angle_local:examples}}
\begin{sphinxVerbatim}[commandchars=\\\{\}]
\PYG{k}{compute }\PYG{n+nv+nvIdentifier}{1}\PYG{+w}{ }\PYG{n+nv+nvIdentifier}{all}\PYG{+w}{ }\PYG{n}{angle}\PYG{o}{/}\PYG{n}{local}\PYG{+w}{ }\PYG{n}{theta}
\PYG{k}{compute }\PYG{n+nv+nvIdentifier}{1}\PYG{+w}{ }\PYG{n+nv+nvIdentifier}{all}\PYG{+w}{ }\PYG{n}{angle}\PYG{o}{/}\PYG{n}{local}\PYG{+w}{ }\PYG{n}{eng}\PYG{+w}{ }\PYG{n}{theta}
\PYG{k}{compute }\PYG{n+nv+nvIdentifier}{1}\PYG{+w}{ }\PYG{n+nv+nvIdentifier}{all}\PYG{+w}{ }\PYG{n}{angle}\PYG{o}{/}\PYG{n}{local}\PYG{+w}{ }\PYG{n}{theta}\PYG{+w}{ }\PYG{n}{v\PYGZus{}cos}\PYG{+w}{ }\PYG{n}{set}\PYG{+w}{ }\PYG{n}{theta}\PYG{+w}{ }\PYG{n}{t}
\end{sphinxVerbatim}


\subsection{Description}
\label{\detokenize{compute_angle_local:description}}
\sphinxAtStartPar
Define a computation that calculates properties of individual angle
interactions.  The number of datums generated, aggregated across all
processors, equals the number of angles in the system, modified by the
group parameter as explained below.

\sphinxAtStartPar
The value \sphinxstyleemphasis{theta} is the angle for the three atoms in the interaction.

\sphinxAtStartPar
The value \sphinxstyleemphasis{eng} is the interaction energy for the angle.

\sphinxAtStartPar
The value \sphinxstyleemphasis{v\_name} can be used together with the \sphinxstyleemphasis{set} keyword to
compute a user\sphinxhyphen{}specified function of the angle theta.  The \sphinxstyleemphasis{name}
specified for the \sphinxstyleemphasis{v\_name} value is the name of an {\hyperref[\detokenize{variable::doc}]{\sphinxcrossref{\DUrole{doc}{equal\sphinxhyphen{}style variable}}}} which should evaluate a formula based on a
variable which will store the angle theta.  This other variable must
be an {\hyperref[\detokenize{variable::doc}]{\sphinxcrossref{\DUrole{doc}{internal\sphinxhyphen{}style variable}}}} defined in the input
script; its initial numeric value can be anything.  It must be an
internal\sphinxhyphen{}style variable, because this command resets its value
directly.  The \sphinxstyleemphasis{set} keyword is used to identify the name of this
other variable associated with theta.

\sphinxAtStartPar
Note that the value of theta for each angle which stored in the
internal variable is in radians, not degrees.

\sphinxAtStartPar
As an example, these commands can be added to the bench/in.rhodo
script to compute the cosine and cosine\sphinxhyphen{}squared of every angle in the
system and output the statistics in various ways:

\begin{sphinxVerbatim}[commandchars=\\\{\}]
\PYG{k}{variable }\PYG{n+nv+nvIdentifier}{t}\PYG{+w}{ }\PYG{n}{internal}\PYG{+w}{ }\PYG{l+m}{0.0}
\PYG{k}{variable }\PYG{n+nv+nvIdentifier}{cos}\PYG{+w}{ }\PYG{n}{equal}\PYG{+w}{ }\PYG{n}{cos}\PYG{n+nv}{(}\PYG{n+nv}{v\PYGZus{}t}\PYG{n+nv}{)}
\PYG{k}{variable }\PYG{n+nv+nvIdentifier}{cossq}\PYG{+w}{ }\PYG{n}{equal}\PYG{+w}{ }\PYG{n}{cos}\PYG{n+nv}{(}\PYG{n+nv}{v\PYGZus{}t}\PYG{n+nv}{)}\PYG{o}{*}\PYG{n}{cos}\PYG{n+nv}{(}\PYG{n+nv}{v\PYGZus{}t}\PYG{n+nv}{)}

\PYG{k}{compute }\PYG{n+nv+nvIdentifier}{1}\PYG{+w}{ }\PYG{n+nv+nvIdentifier}{all}\PYG{+w}{ }\PYG{n}{property}\PYG{o}{/}\PYG{n}{local}\PYG{+w}{ }\PYG{n}{aatom1}\PYG{+w}{ }\PYG{n}{aatom2}\PYG{+w}{ }\PYG{n}{aatom3}\PYG{+w}{ }\PYG{n}{atype}
\PYG{k}{compute }\PYG{n+nv+nvIdentifier}{2}\PYG{+w}{ }\PYG{n+nv+nvIdentifier}{all}\PYG{+w}{ }\PYG{n}{angle}\PYG{o}{/}\PYG{n}{local}\PYG{+w}{ }\PYG{n}{eng}\PYG{+w}{ }\PYG{n}{theta}\PYG{+w}{ }\PYG{n}{v\PYGZus{}cos}\PYG{+w}{ }\PYG{n}{v\PYGZus{}cossq}\PYG{+w}{ }\PYG{n}{set}\PYG{+w}{ }\PYG{n}{theta}\PYG{+w}{ }\PYG{n}{t}
\PYG{k}{dump }\PYG{n+nv+nvIdentifier}{1}\PYG{+w}{ }\PYG{n+nv+nvIdentifier}{all}\PYG{+w}{ }\PYG{n}{local}\PYG{+w}{ }\PYG{l+m}{100}\PYG{+w}{ }\PYG{n}{tmp.dump}\PYG{+w}{ }\PYG{n}{c\PYGZus{}1[}\PYG{o}{*}\PYG{n}{]}\PYG{+w}{ }\PYG{n}{c\PYGZus{}2[}\PYG{o}{*}\PYG{n}{]}

\PYG{k}{compute }\PYG{n+nv+nvIdentifier}{3}\PYG{+w}{ }\PYG{n+nv+nvIdentifier}{all}\PYG{+w}{ }\PYG{n}{reduce}\PYG{+w}{ }\PYG{n}{ave}\PYG{+w}{ }\PYG{n}{c\PYGZus{}2[}\PYG{o}{*}\PYG{n}{]}
\PYG{k}{thermo\PYGZus{}style}\PYG{+w}{ }\PYG{n}{custom}\PYG{+w}{ }\PYG{n}{step}\PYG{+w}{ }\PYG{n}{temp}\PYG{+w}{ }\PYG{n}{press}\PYG{+w}{ }\PYG{n}{c\PYGZus{}3[}\PYG{o}{*}\PYG{n}{]}

\PYG{k}{fix }\PYG{n+nv+nvIdentifier}{10}\PYG{+w}{ }\PYG{n+nv+nvIdentifier}{all}\PYG{+w}{ }\PYG{n}{ave}\PYG{o}{/}\PYG{n}{histo}\PYG{+w}{ }\PYG{l+m}{10}\PYG{+w}{ }\PYG{l+m}{10}\PYG{+w}{ }\PYG{l+m}{100}\PYG{+w}{ }\PYG{o}{\PYGZhy{}}\PYG{l+m}{1}\PYG{+w}{ }\PYG{l+m}{1}\PYG{+w}{ }\PYG{l+m}{20}\PYG{+w}{ }\PYG{n}{c\PYGZus{}2[3]}\PYG{+w}{ }\PYG{n}{mode}\PYG{+w}{ }\PYG{n}{vector}\PYG{+w}{ }\PYG{n}{file}\PYG{+w}{ }\PYG{n}{tmp.histo}
\end{sphinxVerbatim}

\sphinxAtStartPar
The {\hyperref[\detokenize{dump::doc}]{\sphinxcrossref{\DUrole{doc}{dump local}}}} command will output the potential energy
(\(\phi\)), the angle (\(\theta\)), \(\cos(\theta\)), and
\(\cos^2(\theta)\) for every angle \(\theta\) in the system.
The {\hyperref[\detokenize{thermo_style::doc}]{\sphinxcrossref{\DUrole{doc}{thermo\_style}}}} command will print the
average of those quantities via the {\hyperref[\detokenize{compute_reduce::doc}]{\sphinxcrossref{\DUrole{doc}{compute reduce}}}}
command with thermo output.  And the {\hyperref[\detokenize{fix_ave_histo::doc}]{\sphinxcrossref{\DUrole{doc}{fix ave/histo}}}}
command will histogram the \(\cos(\theta)\) values and write them to a
file.


\bigskip\hrule\bigskip


\sphinxAtStartPar
The local data stored by this command is generated by looping over all
the atoms owned on a processor and their angles.  An angle will only
be included if all three atoms in the angle are in the specified compute
group.  Any angles that have been broken (see the
{\hyperref[\detokenize{angle_style::doc}]{\sphinxcrossref{\DUrole{doc}{angle\_style}}}} command) by setting their angle type to
0 are not included.  Angles that have been turned off (see the {\hyperref[\detokenize{fix_shake::doc}]{\sphinxcrossref{\DUrole{doc}{fix shake}}}} or {\hyperref[\detokenize{delete_bonds::doc}]{\sphinxcrossref{\DUrole{doc}{delete\_bonds}}}} commands) by
setting their angle type negative are written into the file, but their
energy will be 0.0.

\sphinxAtStartPar
Note that as atoms migrate from processor to processor, there will be
no consistent ordering of the entries within the local vector or array
from one timestep to the next.  The only consistency that is
guaranteed is that the ordering on a particular timestep will be the
same for local vectors or arrays generated by other compute commands.
For example, angle output from the {\hyperref[\detokenize{compute_property_local::doc}]{\sphinxcrossref{\DUrole{doc}{compute property/local}}}} command can be combined
with data from this command and output by the {\hyperref[\detokenize{dump::doc}]{\sphinxcrossref{\DUrole{doc}{dump local}}}}
command in a consistent way.

\sphinxAtStartPar
Here is an example of how to do this:

\begin{sphinxVerbatim}[commandchars=\\\{\}]
\PYG{k}{compute }\PYG{n+nv+nvIdentifier}{1}\PYG{+w}{ }\PYG{n+nv+nvIdentifier}{all}\PYG{+w}{ }\PYG{n}{property}\PYG{o}{/}\PYG{n}{local}\PYG{+w}{ }\PYG{n}{atype}\PYG{+w}{ }\PYG{n}{aatom1}\PYG{+w}{ }\PYG{n}{aatom2}\PYG{+w}{ }\PYG{n}{aatom3}
\PYG{k}{compute }\PYG{n+nv+nvIdentifier}{2}\PYG{+w}{ }\PYG{n+nv+nvIdentifier}{all}\PYG{+w}{ }\PYG{n}{angle}\PYG{o}{/}\PYG{n}{local}\PYG{+w}{ }\PYG{n}{theta}\PYG{+w}{ }\PYG{n}{eng}
\PYG{k}{dump }\PYG{n+nv+nvIdentifier}{1}\PYG{+w}{ }\PYG{n+nv+nvIdentifier}{all}\PYG{+w}{ }\PYG{n}{local}\PYG{+w}{ }\PYG{l+m}{1000}\PYG{+w}{ }\PYG{n}{tmp.dump}\PYG{+w}{ }\PYG{n}{index}\PYG{+w}{ }\PYG{n}{c\PYGZus{}1[1]}\PYG{+w}{ }\PYG{n}{c\PYGZus{}1[2]}\PYG{+w}{ }\PYG{n}{c\PYGZus{}1[3]}\PYG{+w}{ }\PYG{n}{c\PYGZus{}1[4]}\PYG{+w}{ }\PYG{n}{c\PYGZus{}2[1]}\PYG{+w}{ }\PYG{n}{c\PYGZus{}2[2]}
\end{sphinxVerbatim}


\subsection{Output info}
\label{\detokenize{compute_angle_local:output-info}}
\sphinxAtStartPar
This compute calculates a local vector or local array depending on the
number of values.  The length of the vector or number of rows in the
array is the number of angles.  If a single value is specified, a
local vector is produced.  If two or more values are specified, a
local array is produced where the number of columns = the number of
values.  The vector or array can be accessed by any command that uses
local values from a compute as input.  See the {\hyperref[\detokenize{Howto_output::doc}]{\sphinxcrossref{\DUrole{doc}{Howto output}}}} page for an overview of LAMMPS output
options.

\sphinxAtStartPar
The output for \sphinxstyleemphasis{theta} will be in degrees.  The output for \sphinxstyleemphasis{eng} will
be in energy {\hyperref[\detokenize{units::doc}]{\sphinxcrossref{\DUrole{doc}{units}}}}.


\subsection{Restrictions}
\label{\detokenize{compute_angle_local:restrictions}}\begin{quote}

\sphinxAtStartPar
none
\end{quote}


\subsection{Related commands}
\label{\detokenize{compute_angle_local:related-commands}}
\sphinxAtStartPar
{\hyperref[\detokenize{dump::doc}]{\sphinxcrossref{\DUrole{doc}{dump local}}}}, {\hyperref[\detokenize{compute_property_local::doc}]{\sphinxcrossref{\DUrole{doc}{compute property/local}}}}


\subsection{Default}
\label{\detokenize{compute_angle_local:default}}
\sphinxAtStartPar
none

\sphinxstepscope

\index{compute angmom/chunk@\spxentry{compute angmom/chunk}}\ignorespaces 

\section{compute angmom/chunk command}
\label{\detokenize{compute_angmom_chunk:compute-angmom-chunk-command}}\label{\detokenize{compute_angmom_chunk:index-0}}\label{\detokenize{compute_angmom_chunk::doc}}

\subsection{Syntax}
\label{\detokenize{compute_angmom_chunk:syntax}}
\begin{sphinxVerbatim}[commandchars=\\\{\}]
\PYG{k}{compute }\PYG{n+nv+nvIdentifier}{ID}\PYG{+w}{ }\PYG{n+nv+nvIdentifier}{group\PYGZhy{}ID}\PYG{+w}{ }\PYG{n}{angmom}\PYG{o}{/}\PYG{n}{chunk}\PYG{+w}{ }\PYG{n}{chunkID}
\end{sphinxVerbatim}
\begin{itemize}
\item {} 
\sphinxAtStartPar
ID, group\sphinxhyphen{}ID are documented in {\hyperref[\detokenize{compute::doc}]{\sphinxcrossref{\DUrole{doc}{compute}}}} command

\item {} 
\sphinxAtStartPar
angmom/chunk = style name of this compute command

\item {} 
\sphinxAtStartPar
chunkID = ID of {\hyperref[\detokenize{compute_chunk_atom::doc}]{\sphinxcrossref{\DUrole{doc}{compute chunk/atom}}}} command

\end{itemize}


\subsection{Examples}
\label{\detokenize{compute_angmom_chunk:examples}}
\begin{sphinxVerbatim}[commandchars=\\\{\}]
\PYG{k}{compute }\PYG{n+nv+nvIdentifier}{1}\PYG{+w}{ }\PYG{n+nv+nvIdentifier}{fluid}\PYG{+w}{ }\PYG{n}{angmom}\PYG{o}{/}\PYG{n}{chunk}\PYG{+w}{ }\PYG{n}{molchunk}
\end{sphinxVerbatim}


\subsection{Description}
\label{\detokenize{compute_angmom_chunk:description}}
\sphinxAtStartPar
Define a computation that calculates the angular momentum of multiple
chunks of atoms.

\sphinxAtStartPar
In LAMMPS, chunks are collections of atoms defined by a {\hyperref[\detokenize{compute_chunk_atom::doc}]{\sphinxcrossref{\DUrole{doc}{compute chunk/atom}}}} command, which assigns each atom
to a single chunk (or no chunk).  The ID for this command is specified
as chunkID.  For example, a single chunk could be the atoms in a
molecule or atoms in a spatial bin.  See the {\hyperref[\detokenize{compute_chunk_atom::doc}]{\sphinxcrossref{\DUrole{doc}{compute chunk/atom}}}} and {\hyperref[\detokenize{Howto_chunk::doc}]{\sphinxcrossref{\DUrole{doc}{Howto chunk}}}}
doc pages for details of how chunks can be defined and examples of how
they can be used to measure properties of a system.

\sphinxAtStartPar
This compute calculates the 3 components of the angular momentum
vector for each chunk, due to the velocity/momentum of the individual
atoms in the chunk around the center\sphinxhyphen{}of\sphinxhyphen{}mass of the chunk.  The
calculation includes all effects due to atoms passing through periodic
boundaries.

\sphinxAtStartPar
Note that only atoms in the specified group contribute to the
calculation.  The {\hyperref[\detokenize{compute_chunk_atom::doc}]{\sphinxcrossref{\DUrole{doc}{compute chunk/atom}}}} command
defines its own group; atoms will have a chunk ID = 0 if they are not
in that group, signifying they are not assigned to a chunk, and will
thus also not contribute to this calculation.  You can specify the
“all” group for this command if you simply want to include atoms with
non\sphinxhyphen{}zero chunk IDs.

\begin{sphinxadmonition}{note}{Note:}
\sphinxAtStartPar
The coordinates of an atom contribute to the chunk’s angular
momentum in “unwrapped” form, by using the image flags associated with
each atom.  See the {\hyperref[\detokenize{dump::doc}]{\sphinxcrossref{\DUrole{doc}{dump custom}}}} command for a discussion
of “unwrapped” coordinates.  See the Atoms section of the
{\hyperref[\detokenize{read_data::doc}]{\sphinxcrossref{\DUrole{doc}{read\_data}}}} command for a discussion of image flags and
how they are set for each atom.  You can reset the image flags
(e.g. to 0) before invoking this compute by using the {\hyperref[\detokenize{set::doc}]{\sphinxcrossref{\DUrole{doc}{set image}}}} command.
\end{sphinxadmonition}

\sphinxAtStartPar
The simplest way to output the results of the compute angmom/chunk
calculation to a file is to use the {\hyperref[\detokenize{fix_ave_time::doc}]{\sphinxcrossref{\DUrole{doc}{fix ave/time}}}}
command, for example:

\begin{sphinxVerbatim}[commandchars=\\\{\}]
\PYG{k}{compute }\PYG{n+nv+nvIdentifier}{cc1}\PYG{+w}{ }\PYG{n+nv+nvIdentifier}{all}\PYG{+w}{ }\PYG{n}{chunk}\PYG{o}{/}\PYG{n}{atom}\PYG{+w}{ }\PYG{n}{molecule}
\PYG{k}{compute }\PYG{n+nv+nvIdentifier}{myChunk}\PYG{+w}{ }\PYG{n+nv+nvIdentifier}{all}\PYG{+w}{ }\PYG{n}{angmom}\PYG{o}{/}\PYG{n}{chunk}\PYG{+w}{ }\PYG{n}{cc1}
\PYG{k}{fix }\PYG{n+nv+nvIdentifier}{1}\PYG{+w}{ }\PYG{n+nv+nvIdentifier}{all}\PYG{+w}{ }\PYG{n}{ave}\PYG{o}{/}\PYG{n}{time}\PYG{+w}{ }\PYG{l+m}{100}\PYG{+w}{ }\PYG{l+m}{1}\PYG{+w}{ }\PYG{l+m}{100}\PYG{+w}{ }\PYG{n}{c\PYGZus{}myChunk[}\PYG{o}{*}\PYG{n}{]}\PYG{+w}{ }\PYG{n}{file}\PYG{+w}{ }\PYG{n}{tmp.out}\PYG{+w}{ }\PYG{n}{mode}\PYG{+w}{ }\PYG{n}{vector}
\end{sphinxVerbatim}


\subsection{Output info}
\label{\detokenize{compute_angmom_chunk:output-info}}
\sphinxAtStartPar
This compute calculates a global array where the number of rows = the
number of chunks \sphinxstyleemphasis{Nchunk} as calculated by the specified {\hyperref[\detokenize{compute_chunk_atom::doc}]{\sphinxcrossref{\DUrole{doc}{compute chunk/atom}}}} command.  The number of columns = 3 for the three
(\sphinxstyleemphasis{x}, \sphinxstyleemphasis{y}, \sphinxstyleemphasis{z}) components of the angular momentum for each chunk.
These values can be accessed by any command that uses global array
values from a compute as input.  See the {\hyperref[\detokenize{Howto_output::doc}]{\sphinxcrossref{\DUrole{doc}{Howto output}}}} page for an overview of LAMMPS output
options.

\sphinxAtStartPar
The array values are “intensive”.  The array values will be in
mass\sphinxhyphen{}velocity\sphinxhyphen{}distance {\hyperref[\detokenize{units::doc}]{\sphinxcrossref{\DUrole{doc}{units}}}}.


\subsection{Restrictions}
\label{\detokenize{compute_angmom_chunk:restrictions}}\begin{quote}

\sphinxAtStartPar
none
\end{quote}


\subsection{Related commands}
\label{\detokenize{compute_angmom_chunk:related-commands}}
\sphinxAtStartPar
{\hyperref[\detokenize{variable::doc}]{\sphinxcrossref{\DUrole{doc}{variable angmom() function}}}}


\subsection{Default}
\label{\detokenize{compute_angmom_chunk:default}}
\sphinxAtStartPar
none

\sphinxstepscope

\index{compute ave/sphere/atom@\spxentry{compute ave/sphere/atom}}\ignorespaces 
\index{compute ave/sphere/atom/kk@\spxentry{compute ave/sphere/atom/kk}}\ignorespaces 

\section{compute ave/sphere/atom command}
\label{\detokenize{compute_ave_sphere_atom:compute-ave-sphere-atom-command}}\label{\detokenize{compute_ave_sphere_atom:index-1}}\label{\detokenize{compute_ave_sphere_atom:index-0}}\label{\detokenize{compute_ave_sphere_atom::doc}}
\sphinxAtStartPar
Accelerator Variants: \sphinxstyleemphasis{ave/sphere/atom/kk}


\subsection{Syntax}
\label{\detokenize{compute_ave_sphere_atom:syntax}}
\begin{sphinxVerbatim}[commandchars=\\\{\}]
\PYG{k}{compute }\PYG{n+nv+nvIdentifier}{ID}\PYG{+w}{ }\PYG{n+nv+nvIdentifier}{group\PYGZhy{}ID}\PYG{+w}{ }\PYG{n}{ave}\PYG{o}{/}\PYG{n}{sphere}\PYG{o}{/}\PYG{n}{atom}\PYG{+w}{ }\PYG{n}{keyword}\PYG{+w}{ }\PYG{n}{values}\PYG{+w}{ }\PYG{n}{...}
\end{sphinxVerbatim}
\begin{itemize}
\item {} 
\sphinxAtStartPar
ID, group\sphinxhyphen{}ID are documented in {\hyperref[\detokenize{compute::doc}]{\sphinxcrossref{\DUrole{doc}{compute}}}} command

\item {} 
\sphinxAtStartPar
ave/sphere/atom = style name of this compute command

\item {} 
\sphinxAtStartPar
one or more keyword/value pairs may be appended
\begin{sphinxalltt}
keyword = \sphinxstyleemphasis{cutoff}
  \sphinxstyleemphasis{cutoff} value = distance cutoff
\end{sphinxalltt}

\end{itemize}


\subsection{Examples}
\label{\detokenize{compute_ave_sphere_atom:examples}}
\begin{sphinxVerbatim}[commandchars=\\\{\}]
\PYG{k}{compute }\PYG{n+nv+nvIdentifier}{1}\PYG{+w}{ }\PYG{n+nv+nvIdentifier}{all}\PYG{+w}{ }\PYG{n}{ave}\PYG{o}{/}\PYG{n}{sphere}\PYG{o}{/}\PYG{n}{atom}

\PYG{k}{compute }\PYG{n+nv+nvIdentifier}{1}\PYG{+w}{ }\PYG{n+nv+nvIdentifier}{all}\PYG{+w}{ }\PYG{n}{ave}\PYG{o}{/}\PYG{n}{sphere}\PYG{o}{/}\PYG{n}{atom}\PYG{+w}{ }\PYG{n}{cutoff}\PYG{+w}{ }\PYG{l+m}{5.0}
\PYG{k}{comm\PYGZus{}modify}\PYG{+w}{ }\PYG{n}{cutoff}\PYG{+w}{ }\PYG{l+m}{5.0}
\end{sphinxVerbatim}


\subsection{Description}
\label{\detokenize{compute_ave_sphere_atom:description}}
\sphinxAtStartPar
\DUrole{versionmodified,added}{Added in version 7Jan2022.}

\sphinxAtStartPar
Define a computation that calculates the local mass density and
temperature for each atom based on its neighbors inside a spherical
cutoff.  If an atom has \(M\) neighbors, then its local mass density is
calculated as the sum of its mass and its \(M\) neighbor masses, divided
by the volume of the cutoff sphere (or circle in 2d).  The local
temperature of the atom is calculated as the temperature of the
collection of \(M+1\) atoms, after subtracting the center\sphinxhyphen{}of\sphinxhyphen{}mass velocity
of the \(M+1\) atoms from each of the \(M+1\) atom’s velocities.  This
is effectively the thermal velocity of the neighborhood of the central
atom, similar to {\hyperref[\detokenize{compute_temp_com::doc}]{\sphinxcrossref{\DUrole{doc}{compute temp/com}}}}.

\sphinxAtStartPar
The optional keyword \sphinxstyleemphasis{cutoff} defines the distance cutoff used when
searching for neighbors. The default value is the cutoff specified by
the pair style. If no pair style is defined, then a cutoff must be
defined using this keyword. If the specified cutoff is larger than
that of the pair\_style plus neighbor skin (or no pair style is
defined), the \sphinxstyleemphasis{comm\_modify cutoff} option must also be set to match
that of the \sphinxstyleemphasis{cutoff} keyword.

\sphinxAtStartPar
The neighbor list needed to compute this quantity is constructed each
time the calculation is performed (i.e. each time a snapshot of atoms
is dumped).  Thus it can be inefficient to compute/dump this quantity
too frequently.

\begin{sphinxadmonition}{note}{Note:}
\sphinxAtStartPar
If you have a bonded system, then the settings of
{\hyperref[\detokenize{special_bonds::doc}]{\sphinxcrossref{\DUrole{doc}{special\_bonds}}}} command can remove pairwise
interactions between atoms in the same bond, angle, or dihedral.
This is the default setting for the {\hyperref[\detokenize{special_bonds::doc}]{\sphinxcrossref{\DUrole{doc}{special\_bonds}}}} command, and means those pairwise interactions do
not appear in the neighbor list.  Because this compute uses the
neighbor list, it also means those pairs will not be included in
the order parameter.  This difficulty can be circumvented by
writing a dump file, and using the {\hyperref[\detokenize{rerun::doc}]{\sphinxcrossref{\DUrole{doc}{rerun}}}} command to
compute the order parameter for snapshots in the dump file.  The
rerun script can use a {\hyperref[\detokenize{special_bonds::doc}]{\sphinxcrossref{\DUrole{doc}{special\_bonds}}}} command
that includes all pairs in the neighbor list.
\end{sphinxadmonition}


\bigskip\hrule\bigskip


\sphinxAtStartPar
Styles with a \sphinxstyleemphasis{gpu}, \sphinxstyleemphasis{intel}, \sphinxstyleemphasis{kk}, \sphinxstyleemphasis{omp}, or \sphinxstyleemphasis{opt} suffix are
functionally the same as the corresponding style without the suffix.
They have been optimized to run faster, depending on your available
hardware, as discussed on the {\hyperref[\detokenize{Speed_packages::doc}]{\sphinxcrossref{\DUrole{doc}{Accelerator packages}}}}
page.  The accelerated styles take the same arguments and should
produce the same results, except for round\sphinxhyphen{}off and precision issues.

\sphinxAtStartPar
These accelerated styles are part of the GPU, INTEL, KOKKOS,
OPENMP, and OPT packages, respectively.  They are only enabled if
LAMMPS was built with those packages.  See the {\hyperref[\detokenize{Build_package::doc}]{\sphinxcrossref{\DUrole{doc}{Build package}}}} page for more info.

\sphinxAtStartPar
You can specify the accelerated styles explicitly in your input script
by including their suffix, or you can use the {\hyperref[\detokenize{Run_options::doc}]{\sphinxcrossref{\DUrole{doc}{\sphinxhyphen{}suffix command\sphinxhyphen{}line switch}}}} when you invoke LAMMPS, or you can use the
{\hyperref[\detokenize{suffix::doc}]{\sphinxcrossref{\DUrole{doc}{suffix}}}} command in your input script.

\sphinxAtStartPar
See the {\hyperref[\detokenize{Speed_packages::doc}]{\sphinxcrossref{\DUrole{doc}{Accelerator packages}}}} page for more
instructions on how to use the accelerated styles effectively.


\bigskip\hrule\bigskip



\subsection{Output info}
\label{\detokenize{compute_ave_sphere_atom:output-info}}
\sphinxAtStartPar
This compute calculates a per\sphinxhyphen{}atom array with two columns: mass
density in density {\hyperref[\detokenize{units::doc}]{\sphinxcrossref{\DUrole{doc}{units}}}} and temperature in temperature
{\hyperref[\detokenize{units::doc}]{\sphinxcrossref{\DUrole{doc}{units}}}}.

\sphinxAtStartPar
These values can be accessed by any command that uses per\sphinxhyphen{}atom values
from a compute as input.  See the {\hyperref[\detokenize{Howto_output::doc}]{\sphinxcrossref{\DUrole{doc}{Howto output}}}}
doc page for an overview of LAMMPS output options.


\subsection{Restrictions}
\label{\detokenize{compute_ave_sphere_atom:restrictions}}
\sphinxAtStartPar
This compute is part of the EXTRA\sphinxhyphen{}COMPUTE package.  It is only enabled
if LAMMPS was built with that package.  See the {\hyperref[\detokenize{Build_package::doc}]{\sphinxcrossref{\DUrole{doc}{Build package}}}} page for more info.

\sphinxAtStartPar
This compute requires {\hyperref[\detokenize{neighbor::doc}]{\sphinxcrossref{\DUrole{doc}{neighbor styles ‘bin’ or ‘nsq’}}}}.


\subsection{Related commands}
\label{\detokenize{compute_ave_sphere_atom:related-commands}}
\sphinxAtStartPar
{\hyperref[\detokenize{comm_modify::doc}]{\sphinxcrossref{\DUrole{doc}{comm\_modify}}}}


\subsection{Default}
\label{\detokenize{compute_ave_sphere_atom:default}}
\sphinxAtStartPar
The option defaults are \sphinxstyleemphasis{cutoff} = pair style cutoff.

\sphinxstepscope

\index{compute basal/atom@\spxentry{compute basal/atom}}\ignorespaces 

\section{compute basal/atom command}
\label{\detokenize{compute_basal_atom:compute-basal-atom-command}}\label{\detokenize{compute_basal_atom:index-0}}\label{\detokenize{compute_basal_atom::doc}}

\subsection{Syntax}
\label{\detokenize{compute_basal_atom:syntax}}
\begin{sphinxVerbatim}[commandchars=\\\{\}]
\PYG{k}{compute }\PYG{n+nv+nvIdentifier}{ID}\PYG{+w}{ }\PYG{n+nv+nvIdentifier}{group\PYGZhy{}ID}\PYG{+w}{ }\PYG{n}{basal}\PYG{o}{/}\PYG{n}{atom}
\end{sphinxVerbatim}
\begin{itemize}
\item {} 
\sphinxAtStartPar
ID, group\sphinxhyphen{}ID are documented in {\hyperref[\detokenize{compute::doc}]{\sphinxcrossref{\DUrole{doc}{compute}}}} command

\item {} 
\sphinxAtStartPar
basal/atom = style name of this compute command

\end{itemize}


\subsection{Examples}
\label{\detokenize{compute_basal_atom:examples}}
\begin{sphinxVerbatim}[commandchars=\\\{\}]
\PYG{k}{compute }\PYG{n+nv+nvIdentifier}{1}\PYG{+w}{ }\PYG{n+nv+nvIdentifier}{all}\PYG{+w}{ }\PYG{n}{basal}\PYG{o}{/}\PYG{n}{atom}
\end{sphinxVerbatim}


\subsection{Description}
\label{\detokenize{compute_basal_atom:description}}
\sphinxAtStartPar
Defines a computation that calculates the hexagonal close\sphinxhyphen{}packed “c”
lattice vector for each atom in the group.  It does this by
calculating the normal unit vector to the basal plane for each atom.
The results enable efficient identification and characterization of
twins and grains in hexagonal close\sphinxhyphen{}packed structures.

\sphinxAtStartPar
The output of the compute is thus the 3 components of a unit vector
associated with each atom.  The components are set to 0.0 for
atoms not in the group.

\sphinxAtStartPar
Details of the calculation are given in {\hyperref[\detokenize{compute_basal_atom:barrett}]{\sphinxcrossref{\DUrole{std,std-ref}{(Barrett)}}}}.

\sphinxAtStartPar
The neighbor list needed to compute this quantity is constructed each
time the calculation is performed (i.e. each time a snapshot of atoms
is dumped).  Thus it can be inefficient to compute/dump this quantity
too frequently or to have multiple compute/dump commands, each of
which computes this quantity.

\sphinxAtStartPar
An example input script that uses this compute is provided
in examples/PACKAGES/basal.


\subsection{Output info}
\label{\detokenize{compute_basal_atom:output-info}}
\sphinxAtStartPar
This compute calculates a per\sphinxhyphen{}atom array with three columns, which can be
accessed by indices 1\textendash{}3 by any command that uses per\sphinxhyphen{}atom values from
a compute as input.  See the {\hyperref[\detokenize{Howto_output::doc}]{\sphinxcrossref{\DUrole{doc}{Howto output}}}} doc page
for an overview of LAMMPS output options.

\sphinxAtStartPar
The per\sphinxhyphen{}atom vector values are unitless since the three columns represent
components of a unit vector.


\subsection{Restrictions}
\label{\detokenize{compute_basal_atom:restrictions}}
\sphinxAtStartPar
This compute is part of the EXTRA\sphinxhyphen{}COMPUTE package.  It is only enabled if
LAMMPS was built with that package.  See the {\hyperref[\detokenize{Build_package::doc}]{\sphinxcrossref{\DUrole{doc}{Build package}}}} page for more info.

\sphinxAtStartPar
The output of this compute will be meaningless unless the atoms are on
(or near) hcp lattice sites, since the calculation assumes a
well\sphinxhyphen{}defined basal plane.


\subsection{Related commands}
\label{\detokenize{compute_basal_atom:related-commands}}
\sphinxAtStartPar
{\hyperref[\detokenize{compute_centro_atom::doc}]{\sphinxcrossref{\DUrole{doc}{compute centro/atom}}}}, {\hyperref[\detokenize{compute_ackland_atom::doc}]{\sphinxcrossref{\DUrole{doc}{compute ackland/atom}}}}


\subsection{Default}
\label{\detokenize{compute_basal_atom:default}}
\sphinxAtStartPar
none


\bigskip\hrule\bigskip

\phantomsection\label{\detokenize{compute_basal_atom:barrett}}
\sphinxAtStartPar
\sphinxstylestrong{(Barrett)} Barrett, Tschopp, El Kadiri, Scripta Mat. 66, p.666 (2012).

\sphinxstepscope

\index{compute body/local@\spxentry{compute body/local}}\ignorespaces 

\section{compute body/local command}
\label{\detokenize{compute_body_local:compute-body-local-command}}\label{\detokenize{compute_body_local:index-0}}\label{\detokenize{compute_body_local::doc}}

\subsection{Syntax}
\label{\detokenize{compute_body_local:syntax}}
\begin{sphinxVerbatim}[commandchars=\\\{\}]
\PYG{k}{compute }\PYG{n+nv+nvIdentifier}{ID}\PYG{+w}{ }\PYG{n+nv+nvIdentifier}{group\PYGZhy{}ID}\PYG{+w}{ }\PYG{n}{body}\PYG{o}{/}\PYG{n}{local}\PYG{+w}{ }\PYG{n}{input1}\PYG{+w}{ }\PYG{n}{input2}\PYG{+w}{ }\PYG{n}{...}
\end{sphinxVerbatim}
\begin{itemize}
\item {} 
\sphinxAtStartPar
ID, group\sphinxhyphen{}ID are documented in {\hyperref[\detokenize{compute::doc}]{\sphinxcrossref{\DUrole{doc}{compute}}}} command

\item {} 
\sphinxAtStartPar
body/local = style name of this compute command

\item {} 
\sphinxAtStartPar
one or more keywords may be appended

\item {} 
\sphinxAtStartPar
keyword = \sphinxstyleemphasis{id} or \sphinxstyleemphasis{type} or \sphinxstyleemphasis{integer}
\begin{sphinxalltt}
\sphinxstyleemphasis{id} = atom ID of the body particle
\sphinxstyleemphasis{type} = atom type of the body particle
\sphinxstyleemphasis{integer} = 1,2,3,etc = index of fields defined by body style
\end{sphinxalltt}

\end{itemize}


\subsection{Examples}
\label{\detokenize{compute_body_local:examples}}
\begin{sphinxVerbatim}[commandchars=\\\{\}]
\PYG{k}{compute }\PYG{n+nv+nvIdentifier}{1}\PYG{+w}{ }\PYG{n+nv+nvIdentifier}{all}\PYG{+w}{ }\PYG{n}{body}\PYG{o}{/}\PYG{n}{local}\PYG{+w}{ }\PYG{n}{type}\PYG{+w}{ }\PYG{l+m}{1}\PYG{+w}{ }\PYG{l+m}{2}\PYG{+w}{ }\PYG{l+m}{3}
\PYG{k}{compute }\PYG{n+nv+nvIdentifier}{1}\PYG{+w}{ }\PYG{n+nv+nvIdentifier}{all}\PYG{+w}{ }\PYG{n}{body}\PYG{o}{/}\PYG{n}{local}\PYG{+w}{ }\PYG{l+m}{3}\PYG{+w}{ }\PYG{l+m}{6}
\end{sphinxVerbatim}


\subsection{Description}
\label{\detokenize{compute_body_local:description}}
\sphinxAtStartPar
Define a computation that calculates properties of individual body
sub\sphinxhyphen{}particles.  The number of data generated, aggregated across all
processors, equals the number of body sub\sphinxhyphen{}particles plus the number of
non\sphinxhyphen{}body particles in the system, modified by the group parameter as
explained below.  See the {\hyperref[\detokenize{Howto_body::doc}]{\sphinxcrossref{\DUrole{doc}{Howto body}}}} page for
more details on using body particles.

\sphinxAtStartPar
The local data stored by this command is generated by looping over all
the atoms.  An atom will only be included if it is in the group.  If
the atom is a body particle, then its \(N\) sub\sphinxhyphen{}particles will be looped
over, and it will contribute \(N\) data to the count of data.  If it
is not a body particle, it will contribute 1 datum.

\sphinxAtStartPar
For both body particles and non\sphinxhyphen{}body particles, the \sphinxstyleemphasis{id} keyword
will store the ID of the particle.

\sphinxAtStartPar
For both body particles and non\sphinxhyphen{}body particles, the \sphinxstyleemphasis{type} keyword
will store the type of the particle.

\sphinxAtStartPar
The \sphinxstyleemphasis{integer} keywords mean different things for body and non\sphinxhyphen{}body
particles.  If the atom is not a body particle, only its \sphinxstyleemphasis{x}, \sphinxstyleemphasis{y}, \sphinxstyleemphasis{z}
coordinates can be referenced, using the \sphinxstyleemphasis{integer} keywords 1,2,3.
Note that this means that if you want to access more fields than this
for body particles, then you cannot include non\sphinxhyphen{}body particles in the
group.

\sphinxAtStartPar
For a body particle, the \sphinxstyleemphasis{integer} keywords refer to fields calculated
by the body style for each sub\sphinxhyphen{}particle.  The body style, as specified
by the {\hyperref[\detokenize{atom_style::doc}]{\sphinxcrossref{\DUrole{doc}{atom\_style body}}}}, determines how many fields
exist and what they are.  See the {\hyperref[\detokenize{Howto_body::doc}]{\sphinxcrossref{\DUrole{doc}{Howto\_body}}}} doc
page for details of the different styles.

\sphinxAtStartPar
Here is an example of how to output body information using the {\hyperref[\detokenize{dump::doc}]{\sphinxcrossref{\DUrole{doc}{dump local}}}} command with this compute.  If fields 1, 2, and 3 for the
body sub\sphinxhyphen{}particles are (\sphinxstyleemphasis{x}, \sphinxstyleemphasis{y}, \sphinxstyleemphasis{z}) coordinates, then the dump file will be
formatted similar to the output of a {\hyperref[\detokenize{dump::doc}]{\sphinxcrossref{\DUrole{doc}{dump atom or custom}}}}
command.

\begin{sphinxVerbatim}[commandchars=\\\{\}]
\PYG{k}{compute }\PYG{n+nv+nvIdentifier}{1}\PYG{+w}{ }\PYG{n+nv+nvIdentifier}{all}\PYG{+w}{ }\PYG{n}{body}\PYG{o}{/}\PYG{n}{local}\PYG{+w}{ }\PYG{n}{type}\PYG{+w}{ }\PYG{l+m}{1}\PYG{+w}{ }\PYG{l+m}{2}\PYG{+w}{ }\PYG{l+m}{3}
\PYG{k}{dump }\PYG{n+nv+nvIdentifier}{1}\PYG{+w}{ }\PYG{n+nv+nvIdentifier}{all}\PYG{+w}{ }\PYG{n}{local}\PYG{+w}{ }\PYG{l+m}{1000}\PYG{+w}{ }\PYG{n}{tmp.dump}\PYG{+w}{ }\PYG{n}{index}\PYG{+w}{ }\PYG{n}{c\PYGZus{}1[1]}\PYG{+w}{ }\PYG{n}{c\PYGZus{}1[2]}\PYG{+w}{ }\PYG{n}{c\PYGZus{}1[3]}\PYG{+w}{ }\PYG{n}{c\PYGZus{}1[4]}
\end{sphinxVerbatim}


\subsection{Output info}
\label{\detokenize{compute_body_local:output-info}}
\sphinxAtStartPar
This compute calculates a local vector or local array depending on the
number of keywords.  The length of the vector or number of rows in the
array is the number of data as described above.  If a single keyword
is specified, a local vector is produced.  If two or more keywords are
specified, a local array is produced where the number of columns = the
number of keywords.  The vector or array can be accessed by any
command that uses local values from a compute as input.  See the
{\hyperref[\detokenize{Howto_output::doc}]{\sphinxcrossref{\DUrole{doc}{Howto output}}}} page for an overview of LAMMPS
output options.

\sphinxAtStartPar
The {\hyperref[\detokenize{units::doc}]{\sphinxcrossref{\DUrole{doc}{units}}}} for output values depend on the body style.


\subsection{Restrictions}
\label{\detokenize{compute_body_local:restrictions}}\begin{quote}

\sphinxAtStartPar
none
\end{quote}


\subsection{Related commands}
\label{\detokenize{compute_body_local:related-commands}}
\sphinxAtStartPar
{\hyperref[\detokenize{dump::doc}]{\sphinxcrossref{\DUrole{doc}{dump local}}}}


\subsection{Default}
\label{\detokenize{compute_body_local:default}}
\sphinxAtStartPar
none

\sphinxstepscope

\index{compute bond@\spxentry{compute bond}}\ignorespaces 

\section{compute bond command}
\label{\detokenize{compute_bond:compute-bond-command}}\label{\detokenize{compute_bond:index-0}}\label{\detokenize{compute_bond::doc}}

\subsection{Syntax}
\label{\detokenize{compute_bond:syntax}}
\begin{sphinxVerbatim}[commandchars=\\\{\}]
\PYG{k}{compute }\PYG{n+nv+nvIdentifier}{ID}\PYG{+w}{ }\PYG{n+nv+nvIdentifier}{group\PYGZhy{}ID}\PYG{+w}{ }\PYG{n}{bond}
\end{sphinxVerbatim}
\begin{itemize}
\item {} 
\sphinxAtStartPar
ID, group\sphinxhyphen{}ID are documented in {\hyperref[\detokenize{compute::doc}]{\sphinxcrossref{\DUrole{doc}{compute}}}} command

\item {} 
\sphinxAtStartPar
bond = style name of this compute command

\end{itemize}


\subsection{Examples}
\label{\detokenize{compute_bond:examples}}
\begin{sphinxVerbatim}[commandchars=\\\{\}]
\PYG{k}{compute }\PYG{n+nv+nvIdentifier}{1}\PYG{+w}{ }\PYG{n+nv+nvIdentifier}{all}\PYG{+w}{ }\PYG{n}{bond}
\end{sphinxVerbatim}


\subsection{Description}
\label{\detokenize{compute_bond:description}}
\sphinxAtStartPar
Define a computation that extracts the bond energy calculated by each
of the bond sub\sphinxhyphen{}styles used in the {\hyperref[\detokenize{bond_hybrid::doc}]{\sphinxcrossref{\DUrole{doc}{bond\_style hybrid}}}} command.  These values are made accessible
for output or further processing by other commands.  The group
specified for this command is ignored.

\sphinxAtStartPar
This compute is useful when using {\hyperref[\detokenize{bond_hybrid::doc}]{\sphinxcrossref{\DUrole{doc}{bond\_style hybrid}}}}
if you want to know the portion of the total energy contributed by one
or more of the hybrid sub\sphinxhyphen{}styles.


\subsection{Output info}
\label{\detokenize{compute_bond:output-info}}
\sphinxAtStartPar
This compute calculates a global vector of length \(N\), where \(N\)
is the number of sub\_styles defined by the
{\hyperref[\detokenize{bond_style::doc}]{\sphinxcrossref{\DUrole{doc}{bond\_style hybrid}}}} command,
which can be accessed by indices 1 through \(N\).
These values can be used by any command that uses global scalar or
vector values from a compute as input.  See the
{\hyperref[\detokenize{Howto_output::doc}]{\sphinxcrossref{\DUrole{doc}{Howto output}}}} page for an overview of LAMMPS output
options.

\sphinxAtStartPar
The vector values are “extensive” and will be in energy
{\hyperref[\detokenize{units::doc}]{\sphinxcrossref{\DUrole{doc}{units}}}}.


\subsection{Restrictions}
\label{\detokenize{compute_bond:restrictions}}\begin{quote}

\sphinxAtStartPar
none
\end{quote}


\subsection{Related commands}
\label{\detokenize{compute_bond:related-commands}}
\sphinxAtStartPar
{\hyperref[\detokenize{compute_pe::doc}]{\sphinxcrossref{\DUrole{doc}{compute pe}}}}, {\hyperref[\detokenize{compute_pair::doc}]{\sphinxcrossref{\DUrole{doc}{compute pair}}}}


\subsection{Default}
\label{\detokenize{compute_bond:default}}
\sphinxAtStartPar
none

\sphinxstepscope

\index{compute bond/local@\spxentry{compute bond/local}}\ignorespaces 

\section{compute bond/local command}
\label{\detokenize{compute_bond_local:compute-bond-local-command}}\label{\detokenize{compute_bond_local:index-0}}\label{\detokenize{compute_bond_local::doc}}

\subsection{Syntax}
\label{\detokenize{compute_bond_local:syntax}}
\begin{sphinxVerbatim}[commandchars=\\\{\}]
\PYG{k}{compute }\PYG{n+nv+nvIdentifier}{ID}\PYG{+w}{ }\PYG{n+nv+nvIdentifier}{group\PYGZhy{}ID}\PYG{+w}{ }\PYG{n}{bond}\PYG{o}{/}\PYG{n}{local}\PYG{+w}{ }\PYG{n}{value1}\PYG{+w}{ }\PYG{n}{value2}\PYG{+w}{ }\PYG{n}{...}\PYG{+w}{ }\PYG{n}{keyword}\PYG{+w}{ }\PYG{n}{args}\PYG{+w}{ }\PYG{n}{...}
\end{sphinxVerbatim}
\begin{itemize}
\item {} 
\sphinxAtStartPar
ID, group\sphinxhyphen{}ID are documented in {\hyperref[\detokenize{compute::doc}]{\sphinxcrossref{\DUrole{doc}{compute}}}} command

\item {} 
\sphinxAtStartPar
bond/local = style name of this compute command

\item {} 
\sphinxAtStartPar
one or more values may be appended

\item {} 
\sphinxAtStartPar
value = \sphinxstyleemphasis{dist} or \sphinxstyleemphasis{dx} or \sphinxstyleemphasis{dy} or \sphinxstyleemphasis{dz} or \sphinxstyleemphasis{engpot} or \sphinxstyleemphasis{force} or \sphinxstyleemphasis{fx} or \sphinxstyleemphasis{fy} or \sphinxstyleemphasis{fz} or \sphinxstyleemphasis{engvib} or \sphinxstyleemphasis{engrot} or \sphinxstyleemphasis{engtrans} or \sphinxstyleemphasis{omega} or \sphinxstyleemphasis{velvib} or \sphinxstyleemphasis{v\_name} or \sphinxstyleemphasis{bN}

\end{itemize}
\begin{sphinxalltt}
\sphinxstyleemphasis{dist} = bond distance
\sphinxstyleemphasis{engpot} = bond potential energy
\sphinxstyleemphasis{force} = bond force

\sphinxstyleemphasis{dx},\sphinxstyleemphasis{dy},\sphinxstyleemphasis{dz} = components of pairwise distance
\sphinxstyleemphasis{fx},\sphinxstyleemphasis{fy},\sphinxstyleemphasis{fz} = components of bond force
\sphinxstyleemphasis{engvib} = bond kinetic energy of vibration
\sphinxstyleemphasis{engrot} = bond kinetic energy of rotation
\sphinxstyleemphasis{engtrans} = bond kinetic energy of translation
\sphinxstyleemphasis{omega} = magnitude of bond angular velocity
\sphinxstyleemphasis{velvib} = vibrational velocity along the bond length
\sphinxstyleemphasis{v\_name} = equal\sphinxhyphen{}style variable with name (see below)
\sphinxstyleemphasis{bN} = bond style specific quantities for allowed N values
\end{sphinxalltt}
\begin{itemize}
\item {} 
\sphinxAtStartPar
zero or more keyword/args pairs may be appended

\item {} 
\sphinxAtStartPar
keyword = \sphinxstyleemphasis{set}

\end{itemize}
\begin{sphinxalltt}
\sphinxstyleemphasis{set} args = \sphinxstyleemphasis{dist} name
  \sphinxstyleemphasis{dist} = only currently allowed arg
  name = name of variable to set with distance (dist)
\end{sphinxalltt}


\subsection{Examples}
\label{\detokenize{compute_bond_local:examples}}
\begin{sphinxVerbatim}[commandchars=\\\{\}]
\PYG{k}{compute }\PYG{n+nv+nvIdentifier}{1}\PYG{+w}{ }\PYG{n+nv+nvIdentifier}{all}\PYG{+w}{ }\PYG{n}{bond}\PYG{o}{/}\PYG{n}{local}\PYG{+w}{ }\PYG{n}{engpot}
\PYG{k}{compute }\PYG{n+nv+nvIdentifier}{1}\PYG{+w}{ }\PYG{n+nv+nvIdentifier}{all}\PYG{+w}{ }\PYG{n}{bond}\PYG{o}{/}\PYG{n}{local}\PYG{+w}{ }\PYG{n}{dist}\PYG{+w}{ }\PYG{n}{engpot}\PYG{+w}{ }\PYG{n}{force}

\PYG{k}{compute }\PYG{n+nv+nvIdentifier}{1}\PYG{+w}{ }\PYG{n+nv+nvIdentifier}{all}\PYG{+w}{ }\PYG{n}{bond}\PYG{o}{/}\PYG{n}{local}\PYG{+w}{ }\PYG{n}{dist}\PYG{+w}{ }\PYG{n}{fx}\PYG{+w}{ }\PYG{n}{fy}\PYG{+w}{ }\PYG{n}{fz}\PYG{+w}{ }\PYG{n}{b1}\PYG{+w}{ }\PYG{n}{b2}

\PYG{k}{compute }\PYG{n+nv+nvIdentifier}{1}\PYG{+w}{ }\PYG{n+nv+nvIdentifier}{all}\PYG{+w}{ }\PYG{n}{bond}\PYG{o}{/}\PYG{n}{local}\PYG{+w}{ }\PYG{n}{dist}\PYG{+w}{ }\PYG{n}{v\PYGZus{}distsq}\PYG{+w}{ }\PYG{n}{set}\PYG{+w}{ }\PYG{n}{dist}\PYG{+w}{ }\PYG{n}{d}
\end{sphinxVerbatim}


\subsection{Description}
\label{\detokenize{compute_bond_local:description}}
\sphinxAtStartPar
Define a computation that calculates properties of individual bond
interactions.  The number of datums generated, aggregated across all
processors, equals the number of bonds in the system, modified by the
group parameter as explained below.

\sphinxAtStartPar
All these properties are computed for the pair of atoms in a bond,
whether the two atoms represent a simple diatomic molecule, or are part
of some larger molecule.

\sphinxAtStartPar
The value \sphinxstyleemphasis{dist} is the current length of the bond.
The values \sphinxstyleemphasis{dx}, \sphinxstyleemphasis{dy}, and \sphinxstyleemphasis{dz} are the xyz components of the
\sphinxstyleemphasis{distance} between the pair of atoms. This value is always the
distance from the atom of lower to the one with the higher id.

\sphinxAtStartPar
The value \sphinxstyleemphasis{engpot} is the potential energy for the bond,
based on the current separation of the pair of atoms in the bond.

\sphinxAtStartPar
The value \sphinxstyleemphasis{force} is the magnitude of the force acting between the
pair of atoms in the bond.

\sphinxAtStartPar
The values \sphinxstyleemphasis{fx}, \sphinxstyleemphasis{fy}, and \sphinxstyleemphasis{fz} are the xyz components of
\sphinxstyleemphasis{force} between the pair of atoms in the bond. For bond styles that apply
non\sphinxhyphen{}central forces, such as {\hyperref[\detokenize{bond_bpm_rotational::doc}]{\sphinxcrossref{\DUrole{doc}{bond\_style bpm/rotational}}}}, these values only include the \((x,y,z)\)
components of the normal force component.

\sphinxAtStartPar
The remaining properties are all computed for motion of the two atoms
relative to the center of mass (COM) velocity of the two atoms in the
bond.

\sphinxAtStartPar
The value \sphinxstyleemphasis{engvib} is the vibrational kinetic energy of the two atoms
in the bond, which is simply \(\frac12 m_1 v_1^2 + \frac12 m_2 v_2^2,\)
where \(v_1\) and \(v_2\) are the magnitude of the velocity of the two
atoms along the bond direction, after the COM velocity has been subtracted from
each.

\sphinxAtStartPar
The value \sphinxstyleemphasis{engrot} is the rotational kinetic energy of the two atoms
in the bond, which is simply \(\frac12 m_1 v_1^2 + \frac12 m_2 v_2^2,\)
where \(v_1\) and \(v_2\) are the magnitude of the velocity of the two
atoms perpendicular to the bond direction, after the COM velocity has been
subtracted from each.

\sphinxAtStartPar
The value \sphinxstyleemphasis{engtrans} is the translational kinetic energy associated
with the motion of the COM of the system itself, namely \(\frac12(m_1+m_2)
V_{\mathrm{cm}}^2\), where \sphinxtitleref{Vcm} = magnitude of the velocity of the COM.

\sphinxAtStartPar
Note that these three kinetic energy terms are simply a partitioning of
the summed kinetic energy of the two atoms themselves.  That is, the total
kinetic energy is
\(\frac12 m_1 v_1^2 + \frac12 m_2 v_2^2\) = engvib + engrot + engtrans,
where \(v_1\) and \(v_2\) are the magnitude of the velocities of the
two atoms, without any adjustment for the COM velocity.

\sphinxAtStartPar
The value \sphinxstyleemphasis{omega} is the magnitude of the angular velocity of the
two atoms around their COM position.

\sphinxAtStartPar
The value \sphinxstyleemphasis{velvib} is the magnitude of the relative velocity of the
two atoms in the bond towards each other.  A negative value means the
two atoms are moving toward each other; a positive value means they are
moving apart.

\sphinxAtStartPar
The value \sphinxstyleemphasis{v\_name} can be used together with the \sphinxstyleemphasis{set} keyword to
compute a user\sphinxhyphen{}specified function of the bond distance.  The \sphinxstyleemphasis{name}
specified for the \sphinxstyleemphasis{v\_name} value is the name of an {\hyperref[\detokenize{variable::doc}]{\sphinxcrossref{\DUrole{doc}{equal\sphinxhyphen{}style variable}}}} which should evaluate a formula based on a
variable which will store the bond distance.  This other variable must
be an {\hyperref[\detokenize{variable::doc}]{\sphinxcrossref{\DUrole{doc}{internal\sphinxhyphen{}style variable}}}} defined in the input
script; its initial numeric value can be anything.  It must be an
internal\sphinxhyphen{}style variable, because this command resets its value
directly.  The \sphinxstyleemphasis{set} keyword is used to identify the name of this
other variable associated with theta.

\sphinxAtStartPar
As an example, these commands can be added to the bench/in.rhodo
script to compute the length\(^2\) of every bond in the system and
output the statistics in various ways:

\begin{sphinxVerbatim}[commandchars=\\\{\}]
\PYG{k}{variable }\PYG{n+nv+nvIdentifier}{d}\PYG{+w}{ }\PYG{n}{internal}\PYG{+w}{ }\PYG{l+m}{0.0}
\PYG{k}{variable }\PYG{n+nv+nvIdentifier}{dsq}\PYG{+w}{ }\PYG{n}{equal}\PYG{+w}{ }\PYG{n}{v\PYGZus{}d}\PYG{o}{*}\PYG{n}{v\PYGZus{}d}

\PYG{k}{compute }\PYG{n+nv+nvIdentifier}{1}\PYG{+w}{ }\PYG{n+nv+nvIdentifier}{all}\PYG{+w}{ }\PYG{n}{property}\PYG{o}{/}\PYG{n}{local}\PYG{+w}{ }\PYG{n}{batom1}\PYG{+w}{ }\PYG{n}{batom2}\PYG{+w}{ }\PYG{n}{btype}
\PYG{k}{compute }\PYG{n+nv+nvIdentifier}{2}\PYG{+w}{ }\PYG{n+nv+nvIdentifier}{all}\PYG{+w}{ }\PYG{n}{bond}\PYG{o}{/}\PYG{n}{local}\PYG{+w}{ }\PYG{n}{engpot}\PYG{+w}{ }\PYG{n}{dist}\PYG{+w}{ }\PYG{n}{v\PYGZus{}dsq}\PYG{+w}{ }\PYG{n}{set}\PYG{+w}{ }\PYG{n}{dist}\PYG{+w}{ }\PYG{n}{d}
\PYG{k}{dump }\PYG{n+nv+nvIdentifier}{1}\PYG{+w}{ }\PYG{n+nv+nvIdentifier}{all}\PYG{+w}{ }\PYG{n}{local}\PYG{+w}{ }\PYG{l+m}{100}\PYG{+w}{ }\PYG{n}{tmp.dump}\PYG{+w}{ }\PYG{n}{c\PYGZus{}1[}\PYG{o}{*}\PYG{n}{]}\PYG{+w}{ }\PYG{n}{c\PYGZus{}2[}\PYG{o}{*}\PYG{n}{]}

\PYG{k}{compute }\PYG{n+nv+nvIdentifier}{3}\PYG{+w}{ }\PYG{n+nv+nvIdentifier}{all}\PYG{+w}{ }\PYG{n}{reduce}\PYG{+w}{ }\PYG{n}{ave}\PYG{+w}{ }\PYG{n}{c\PYGZus{}2[}\PYG{o}{*}\PYG{n}{]}
\PYG{k}{thermo\PYGZus{}style}\PYG{+w}{ }\PYG{n}{custom}\PYG{+w}{ }\PYG{n}{step}\PYG{+w}{ }\PYG{n}{temp}\PYG{+w}{ }\PYG{n}{press}\PYG{+w}{ }\PYG{n}{c\PYGZus{}3[}\PYG{o}{*}\PYG{n}{]}

\PYG{k}{fix }\PYG{n+nv+nvIdentifier}{10}\PYG{+w}{ }\PYG{n+nv+nvIdentifier}{all}\PYG{+w}{ }\PYG{n}{ave}\PYG{o}{/}\PYG{n}{histo}\PYG{+w}{ }\PYG{l+m}{10}\PYG{+w}{ }\PYG{l+m}{10}\PYG{+w}{ }\PYG{l+m}{100}\PYG{+w}{ }\PYG{l+m}{0}\PYG{+w}{ }\PYG{l+m}{6}\PYG{+w}{ }\PYG{l+m}{20}\PYG{+w}{ }\PYG{n}{c\PYGZus{}2[3]}\PYG{+w}{ }\PYG{n}{mode}\PYG{+w}{ }\PYG{n}{vector}\PYG{+w}{ }\PYG{n}{file}\PYG{+w}{ }\PYG{n}{tmp.histo}
\end{sphinxVerbatim}

\sphinxAtStartPar
The {\hyperref[\detokenize{dump::doc}]{\sphinxcrossref{\DUrole{doc}{dump local}}}} command will output the energy, length,
and length\(^2\) for every bond in the system.  The
{\hyperref[\detokenize{thermo_style::doc}]{\sphinxcrossref{\DUrole{doc}{thermo\_style}}}} command will print the average of
those quantities via the {\hyperref[\detokenize{compute_reduce::doc}]{\sphinxcrossref{\DUrole{doc}{compute reduce}}}} command
with thermo output, and the {\hyperref[\detokenize{fix_ave_histo::doc}]{\sphinxcrossref{\DUrole{doc}{fix ave/histo}}}}
command will histogram the length\(^2\) values and write them to a file.

\sphinxAtStartPar
A bond style may define additional bond quantities which can be
accessed as \sphinxstyleemphasis{b1} to \sphinxstyleemphasis{bN}, where N is defined by the bond style.  Most
bond styles do not define any additional quantities, so N = 0.  An
example of ones that do are the {\hyperref[\detokenize{Howto_bpm::doc}]{\sphinxcrossref{\DUrole{doc}{BPM bond styles}}}}
which store the reference state between two particles. See
individual bond styles for details.

\sphinxAtStartPar
When using \sphinxstyleemphasis{bN} with bond style \sphinxstyleemphasis{hybrid}, the output will be the Nth
quantity from the sub\sphinxhyphen{}style that computes the bonded interaction
(based on bond type).  If that sub\sphinxhyphen{}style does not define a \sphinxstyleemphasis{bN},
the output will be 0.0.  The maximum allowed N is the maximum number
of quantities provided by any sub\sphinxhyphen{}style.


\bigskip\hrule\bigskip


\sphinxAtStartPar
The local data stored by this command is generated by looping over all
the atoms owned on a processor and their bonds.  A bond will only be
included if both atoms in the bond are in the specified compute group.
Any bonds that have been broken (see the {\hyperref[\detokenize{bond_style::doc}]{\sphinxcrossref{\DUrole{doc}{bond\_style}}}}
command) by setting their bond type to 0 are not included.  Bonds that
have been turned off (see the {\hyperref[\detokenize{fix_shake::doc}]{\sphinxcrossref{\DUrole{doc}{fix shake}}}} or
{\hyperref[\detokenize{delete_bonds::doc}]{\sphinxcrossref{\DUrole{doc}{delete\_bonds}}}} commands) by setting their bond type
negative are written into the file, but their energy will be 0.0.

\sphinxAtStartPar
Note that as atoms migrate from processor to processor, there will be
no consistent ordering of the entries within the local vector or array
from one timestep to the next.  The only consistency that is
guaranteed is that the ordering on a particular timestep will be the
same for local vectors or arrays generated by other compute commands.
For example, bond output from the {\hyperref[\detokenize{compute_property_local::doc}]{\sphinxcrossref{\DUrole{doc}{compute property/local}}}} command can be combined
with data from this command and output by the {\hyperref[\detokenize{dump::doc}]{\sphinxcrossref{\DUrole{doc}{dump local}}}}
command in a consistent way.

\sphinxAtStartPar
Here is an example of how to do this:

\begin{sphinxVerbatim}[commandchars=\\\{\}]
\PYG{k}{compute }\PYG{n+nv+nvIdentifier}{1}\PYG{+w}{ }\PYG{n+nv+nvIdentifier}{all}\PYG{+w}{ }\PYG{n}{property}\PYG{o}{/}\PYG{n}{local}\PYG{+w}{ }\PYG{n}{btype}\PYG{+w}{ }\PYG{n}{batom1}\PYG{+w}{ }\PYG{n}{batom2}
\PYG{k}{compute }\PYG{n+nv+nvIdentifier}{2}\PYG{+w}{ }\PYG{n+nv+nvIdentifier}{all}\PYG{+w}{ }\PYG{n}{bond}\PYG{o}{/}\PYG{n}{local}\PYG{+w}{ }\PYG{n}{dist}\PYG{+w}{ }\PYG{n}{engpot}
\PYG{k}{dump }\PYG{n+nv+nvIdentifier}{1}\PYG{+w}{ }\PYG{n+nv+nvIdentifier}{all}\PYG{+w}{ }\PYG{n}{local}\PYG{+w}{ }\PYG{l+m}{1000}\PYG{+w}{ }\PYG{n}{tmp.dump}\PYG{+w}{ }\PYG{n}{index}\PYG{+w}{ }\PYG{n}{c\PYGZus{}1[}\PYG{o}{*}\PYG{n}{]}\PYG{+w}{ }\PYG{n}{c\PYGZus{}2[}\PYG{o}{*}\PYG{n}{]}
\end{sphinxVerbatim}


\subsection{Output info}
\label{\detokenize{compute_bond_local:output-info}}
\sphinxAtStartPar
This compute calculates a local vector or local array depending on the
number of values.  The length of the vector or number of rows in the
array is the number of bonds.  If a single value is specified, a local
vector is produced.  If two or more values are specified, a local
array is produced where the number of columns = the number of values.
The vector or array can be accessed by any command that uses local
values from a compute as input.  See the {\hyperref[\detokenize{Howto_output::doc}]{\sphinxcrossref{\DUrole{doc}{Howto output}}}} page for an overview of LAMMPS output
options.

\sphinxAtStartPar
The output for \sphinxstyleemphasis{dist} will be in distance {\hyperref[\detokenize{units::doc}]{\sphinxcrossref{\DUrole{doc}{units}}}}. The
output for \sphinxstyleemphasis{velvib} will be in velocity {\hyperref[\detokenize{units::doc}]{\sphinxcrossref{\DUrole{doc}{units}}}}. The output
for \sphinxstyleemphasis{omega} will be in velocity/distance {\hyperref[\detokenize{units::doc}]{\sphinxcrossref{\DUrole{doc}{units}}}}. The
output for \sphinxstyleemphasis{engtrans}, \sphinxstyleemphasis{engvib}, \sphinxstyleemphasis{engrot}, and \sphinxstyleemphasis{engpot} will be in
energy {\hyperref[\detokenize{units::doc}]{\sphinxcrossref{\DUrole{doc}{units}}}}. The output for \sphinxstyleemphasis{force} will be in force
{\hyperref[\detokenize{units::doc}]{\sphinxcrossref{\DUrole{doc}{units}}}}.


\subsection{Restrictions}
\label{\detokenize{compute_bond_local:restrictions}}\begin{quote}

\sphinxAtStartPar
none
\end{quote}


\subsection{Related commands}
\label{\detokenize{compute_bond_local:related-commands}}
\sphinxAtStartPar
{\hyperref[\detokenize{dump::doc}]{\sphinxcrossref{\DUrole{doc}{dump local}}}}, {\hyperref[\detokenize{compute_property_local::doc}]{\sphinxcrossref{\DUrole{doc}{compute property/local}}}}


\subsection{Default}
\label{\detokenize{compute_bond_local:default}}
\sphinxAtStartPar
none

\sphinxstepscope

\index{compute born/matrix@\spxentry{compute born/matrix}}\ignorespaces 

\section{compute born/matrix command}
\label{\detokenize{compute_born_matrix:compute-born-matrix-command}}\label{\detokenize{compute_born_matrix:index-0}}\label{\detokenize{compute_born_matrix::doc}}

\subsection{Syntax}
\label{\detokenize{compute_born_matrix:syntax}}
\begin{sphinxVerbatim}[commandchars=\\\{\}]
\PYG{k}{compute }\PYG{n+nv+nvIdentifier}{ID}\PYG{+w}{ }\PYG{n+nv+nvIdentifier}{group\PYGZhy{}ID}\PYG{+w}{ }\PYG{n}{born}\PYG{o}{/}\PYG{n}{matrix}\PYG{+w}{ }\PYG{n}{keyword}\PYG{+w}{ }\PYG{n}{value}\PYG{+w}{ }\PYG{n}{...}
\end{sphinxVerbatim}
\begin{itemize}
\item {} 
\sphinxAtStartPar
ID, group\sphinxhyphen{}ID are documented in {\hyperref[\detokenize{compute::doc}]{\sphinxcrossref{\DUrole{doc}{compute}}}} command

\item {} 
\sphinxAtStartPar
born/matrix = style name of this compute command

\item {} 
\sphinxAtStartPar
zero or more keywords or keyword/value pairs may be appended
\begin{sphinxalltt}
keyword = \sphinxstyleemphasis{numdiff} or \sphinxstyleemphasis{pair} or \sphinxstyleemphasis{bond} or \sphinxstyleemphasis{angle} or \sphinxstyleemphasis{dihedral} or \sphinxstyleemphasis{improper}
  \sphinxstyleemphasis{numdiff} values = delta virial\sphinxhyphen{}ID
    delta = magnitude of strain (dimensionless)
    virial\sphinxhyphen{}ID = ID of pressure compute for virial (string)
    (\sphinxstyleemphasis{numdiff} cannot be used with any other keyword)
  \sphinxstyleemphasis{pair} = compute pair\sphinxhyphen{}wise contributions
  \sphinxstyleemphasis{bond} = compute bonding contributions
  \sphinxstyleemphasis{angle} = compute angle contributions
  \sphinxstyleemphasis{dihedral} = compute dihedral contributions
  \sphinxstyleemphasis{improper} = compute improper contributions
\end{sphinxalltt}

\end{itemize}


\subsection{Examples}
\label{\detokenize{compute_born_matrix:examples}}
\begin{sphinxVerbatim}[commandchars=\\\{\}]
\PYG{k}{compute }\PYG{n+nv+nvIdentifier}{1}\PYG{+w}{ }\PYG{n+nv+nvIdentifier}{all}\PYG{+w}{ }\PYG{n}{born}\PYG{o}{/}\PYG{n}{matrix}
\PYG{k}{compute }\PYG{n+nv+nvIdentifier}{1}\PYG{+w}{ }\PYG{n+nv+nvIdentifier}{all}\PYG{+w}{ }\PYG{n}{born}\PYG{o}{/}\PYG{n}{matrix}\PYG{+w}{ }\PYG{n}{bond}\PYG{+w}{ }\PYG{n}{angle}
\PYG{k}{compute }\PYG{n+nv+nvIdentifier}{1}\PYG{+w}{ }\PYG{n+nv+nvIdentifier}{all}\PYG{+w}{ }\PYG{n}{born}\PYG{o}{/}\PYG{n}{matrix}\PYG{+w}{ }\PYG{n}{numdiff}\PYG{+w}{ }\PYG{l+m}{1.0e\PYGZhy{}4}\PYG{+w}{ }\PYG{n}{myvirial}
\end{sphinxVerbatim}


\subsection{Description}
\label{\detokenize{compute_born_matrix:description}}
\sphinxAtStartPar
\DUrole{versionmodified,added}{Added in version 4May2022.}

\sphinxAtStartPar
Define a compute that calculates
\(\frac{\partial{}^2U}{\partial\varepsilon_{i}\partial\varepsilon_{j}},\) the
second derivatives of the potential energy \(U\) with respect to the strain
tensor \(\varepsilon\) elements. These values are related to:
\begin{equation*}
\begin{split}C^{B}_{i,j}=\frac{1}{V}\frac{\partial{}^2U}{\partial{}\varepsilon_{i}\partial\varepsilon_{j}}\end{split}
\end{equation*}
\sphinxAtStartPar
also called the Born term of elastic constants in the stress\sphinxhyphen{}stress fluctuation
formalism. This quantity can be used to compute the elastic constant tensor.
Using the symmetric Voigt notation, the elastic constant tensor can be written
as a 6x6 symmetric matrix:
\begin{equation*}
\begin{split}C_{i,j} = \langle{}C^{B}_{i,j}\rangle
         + \frac{V}{k_{B}T}\left(\langle\sigma_{i}\sigma_{j}\rangle\right.
         \left.- \langle\sigma_{i}\rangle\langle\sigma_{j}\rangle\right)
         + \frac{Nk_{B}T}{V}
           \left(\delta_{i,j}+(\delta_{1,i}+\delta_{2,i}+\delta_{3,i})\right.
           \left.*(\delta_{1,j}+\delta_{2,j}+\delta_{3,j})\right)\end{split}
\end{equation*}
\sphinxAtStartPar
In the above expression, \(\sigma\) stands for the virial stress
tensor, \(\delta\) is the Kronecker delta and the usual notation apply for
the number of particle, the temperature and volume respectively \(N\),
\(T\) and \(V\). \(k_{B}\) is the Boltzmann constant.

\sphinxAtStartPar
The Born term is a symmetric 6x6 matrix, as is the matrix of second derivatives
of potential energy w.r.t strain,
whose 21 independent elements are output in this order:
\begin{equation*}
\begin{split}\begin{bmatrix}
   C_{1}  & C_{7}   & C_{8}  & C_{9}  & C_{10} & C_{11} \\
   C_{7}  & C_{2}   & C_{12} & C_{13} & C_{14} & C_{15} \\
   \vdots & C_{12}  & C_{3}  & C_{16} & C_{17} & C_{18} \\
   \vdots & C_{13}  & C_{16} & C_{4}  & C_{19} & C_{20} \\
   \vdots & \vdots  & \vdots & C_{19} & C_{5}  & C_{21} \\
   \vdots & \vdots  & \vdots & \vdots & C_{21} & C_{6}
\end{bmatrix}\end{split}
\end{equation*}
\sphinxAtStartPar
in this matrix the indices of \(C_{k}\) value are the corresponding element
\(k\) in the global vector output by this compute. Each term comes from the sum
of the derivatives of every contribution to the potential energy
in the system as explained in {\hyperref[\detokenize{compute_born_matrix:vanworkum}]{\sphinxcrossref{\DUrole{std,std-ref}{(VanWorkum)}}}}.

\sphinxAtStartPar
The output can be accessed using usual Lammps routines:

\begin{sphinxVerbatim}[commandchars=\\\{\}]
\PYG{k}{compute }\PYG{n+nv+nvIdentifier}{1}\PYG{+w}{ }\PYG{n+nv+nvIdentifier}{all}\PYG{+w}{ }\PYG{n}{born}\PYG{o}{/}\PYG{n}{matrix}
\PYG{k}{compute }\PYG{n+nv+nvIdentifier}{2}\PYG{+w}{ }\PYG{n+nv+nvIdentifier}{all}\PYG{+w}{ }\PYG{n}{pressure}\PYG{+w}{ }\PYG{n}{NULL}\PYG{+w}{ }\PYG{n}{virial}
\PYG{k}{variable }\PYG{n+nv+nvIdentifier}{S1}\PYG{+w}{ }\PYG{n}{equal}\PYG{+w}{ }\PYG{o}{\PYGZhy{}}\PYG{n}{c\PYGZus{}2[1]}
\PYG{k}{variable }\PYG{n+nv+nvIdentifier}{S2}\PYG{+w}{ }\PYG{n}{equal}\PYG{+w}{ }\PYG{o}{\PYGZhy{}}\PYG{n}{c\PYGZus{}2[2]}
\PYG{k}{variable }\PYG{n+nv+nvIdentifier}{S3}\PYG{+w}{ }\PYG{n}{equal}\PYG{+w}{ }\PYG{o}{\PYGZhy{}}\PYG{n}{c\PYGZus{}2[3]}
\PYG{k}{variable }\PYG{n+nv+nvIdentifier}{S4}\PYG{+w}{ }\PYG{n}{equal}\PYG{+w}{ }\PYG{o}{\PYGZhy{}}\PYG{n}{c\PYGZus{}2[4]}
\PYG{k}{variable }\PYG{n+nv+nvIdentifier}{S5}\PYG{+w}{ }\PYG{n}{equal}\PYG{+w}{ }\PYG{o}{\PYGZhy{}}\PYG{n}{c\PYGZus{}2[5]}
\PYG{k}{variable }\PYG{n+nv+nvIdentifier}{S6}\PYG{+w}{ }\PYG{n}{equal}\PYG{+w}{ }\PYG{o}{\PYGZhy{}}\PYG{n}{c\PYGZus{}2[6]}
\PYG{k}{fix }\PYG{n+nv+nvIdentifier}{1}\PYG{+w}{ }\PYG{n+nv+nvIdentifier}{all}\PYG{+w}{ }\PYG{n}{ave}\PYG{o}{/}\PYG{n}{time}\PYG{+w}{ }\PYG{l+m}{1}\PYG{+w}{ }\PYG{l+m}{1}\PYG{+w}{ }\PYG{l+m}{1}\PYG{+w}{ }\PYG{n}{v\PYGZus{}S1}\PYG{+w}{ }\PYG{n}{v\PYGZus{}S2}\PYG{+w}{ }\PYG{n}{v\PYGZus{}S3}\PYG{+w}{ }\PYG{n}{v\PYGZus{}S4}\PYG{+w}{ }\PYG{n}{v\PYGZus{}S5}\PYG{+w}{ }\PYG{n}{v\PYGZus{}S6}\PYG{+w}{ }\PYG{n}{c\PYGZus{}1[}\PYG{o}{*}\PYG{n}{]}\PYG{+w}{ }\PYG{n}{file}\PYG{+w}{ }\PYG{n}{born.out}
\end{sphinxVerbatim}

\sphinxAtStartPar
In this example, the file \sphinxstyleemphasis{born.out} will contain the information needed to
compute the first and second terms of the elastic constant matrix in a post
processing procedure. The other required quantities can be accessed using any
other \sphinxstyleemphasis{LAMMPS} usual method. Several examples of this method are
provided in the examples/ELASTIC\_T/BORN\_MATRIX directory
described on the {\hyperref[\detokenize{Examples::doc}]{\sphinxcrossref{\DUrole{doc}{Examples}}}} doc page.

\sphinxAtStartPar
NOTE: In the above \(C_{i,j}\) computation, the fluctuation
term involving the virial stress tensor \(\sigma\) is the
covariance between each elements. In a
solid the stress fluctuations can vary rapidly, while average
fluctuations can be slow to converge.
A detailed analysis of the convergence rate of all the terms in
the elastic tensor
is provided in the paper by Clavier et al. {\hyperref[\detokenize{compute_born_matrix:clavier2}]{\sphinxcrossref{\DUrole{std,std-ref}{(Clavier)}}}}.

\sphinxAtStartPar
Two different computation methods for the Born matrix are implemented in this
compute and are mutually exclusive.

\sphinxAtStartPar
The first one is a direct computation from the analytical formula from the
different terms of the potential used for the simulations {\hyperref[\detokenize{compute_born_matrix:vanworkum}]{\sphinxcrossref{\DUrole{std,std-ref}{(VanWorkum)}}}}. However, the implementation of such derivations must be done
for every potential form. This has not been done yet and can be very
complicated for complex potentials. At the moment a warning message is
displayed for every term that is not supporting the compute at the moment.
This method is the default for now.

\sphinxAtStartPar
The second method uses finite differences of energy to numerically approximate
the second derivatives {\hyperref[\detokenize{compute_born_matrix:zhen}]{\sphinxcrossref{\DUrole{std,std-ref}{(Zhen)}}}}. This is useful when using
interaction styles for which the analytical second derivatives have not been
implemented. In this cases, the compute applies linear strain fields of
magnitude \sphinxstyleemphasis{delta} to all the atoms relative to a point at the center of the
box. The strain fields are in six different directions, corresponding to the
six Cartesian components of the stress tensor defined by LAMMPS. For each
direction it applies the strain field in both the positive and negative senses,
and the new stress virial tensor of the entire system is calculated after each.
The difference in these two virials divided by two times \sphinxstyleemphasis{delta}, approximates
the corresponding components of the second derivative, after applying a
suitable unit conversion.

\begin{sphinxadmonition}{note}{Note:}
\sphinxAtStartPar
It is important to choose a suitable value for delta, the magnitude of
strains that are used to generate finite difference
approximations to the exact virial stress.  For typical systems, a value in
the range of 1 part in 1e5 to 1e6 will be sufficient.
However, the best value will depend on a multitude of factors
including the stiffness of the interatomic potential, the thermodynamic
state of the material being probed, and so on. The only way to be sure
that you have made a good choice is to do a sensitivity study on a
representative atomic configuration, sweeping over a wide range of
values of delta. If delta is too small, the output values will vary
erratically due to truncation effects. If delta is increased beyond a
certain point, the output values will start to vary smoothly with
delta, due to growing contributions from higher order derivatives. In
between these two limits, the numerical virial values should be largely
independent of delta.
\end{sphinxadmonition}

\sphinxAtStartPar
The keyword requires the additional arguments \sphinxstyleemphasis{delta} and \sphinxstyleemphasis{virial\sphinxhyphen{}ID}.
\sphinxstyleemphasis{delta} gives the size of the applied strains. \sphinxstyleemphasis{virial\sphinxhyphen{}ID} gives
the ID string of the pressure compute that provides the virial stress tensor,
requiring that it use the virial keyword e.g.

\begin{sphinxVerbatim}[commandchars=\\\{\}]
\PYG{k}{compute }\PYG{n+nv+nvIdentifier}{myvirial}\PYG{+w}{ }\PYG{n+nv+nvIdentifier}{all}\PYG{+w}{ }\PYG{n}{pressure}\PYG{+w}{ }\PYG{n}{NULL}\PYG{+w}{ }\PYG{n}{virial}
\PYG{k}{compute }\PYG{n+nv+nvIdentifier}{1}\PYG{+w}{ }\PYG{n+nv+nvIdentifier}{all}\PYG{+w}{ }\PYG{n}{born}\PYG{o}{/}\PYG{n}{matrix}\PYG{+w}{ }\PYG{n}{numdiff}\PYG{+w}{ }\PYG{l+m}{1.0e\PYGZhy{}4}\PYG{+w}{ }\PYG{n}{myvirial}
\end{sphinxVerbatim}

\sphinxAtStartPar
\sphinxstylestrong{Output info:}

\sphinxAtStartPar
This compute calculates a global vector with 21 values that are
the second derivatives of the potential energy with respect to strain.
The values are in energy units.
The values are ordered as explained above. These values can be used
by any command that uses global values from a compute as input. See
the {\hyperref[\detokenize{Howto_output::doc}]{\sphinxcrossref{\DUrole{doc}{Howto output}}}} doc page for an overview of
LAMMPS output options.

\sphinxAtStartPar
The array values calculated by this compute are all “extensive”.


\subsection{Restrictions}
\label{\detokenize{compute_born_matrix:restrictions}}
\sphinxAtStartPar
This compute is part of the EXTRA\sphinxhyphen{}COMPUTE package.  It is only enabled if
LAMMPS was built with that package.  See the {\hyperref[\detokenize{Build_package::doc}]{\sphinxcrossref{\DUrole{doc}{Build package}}}} page for more info.  LAMMPS was built with that package.  See
the {\hyperref[\detokenize{Build_package::doc}]{\sphinxcrossref{\DUrole{doc}{Build package}}}} page for more info.

\sphinxAtStartPar
The Born term can be decomposed as a product of two terms. The first one is a
general term which depends on the configuration. The second one is specific to
every interaction composing your force field (non\sphinxhyphen{}bonded, bonds, angle, …).
Currently not all LAMMPS interaction styles implement the \sphinxstyleemphasis{born\_matrix} method
giving first and second order derivatives and LAMMPS will exit with an error if
this compute is used with such interactions unless the \sphinxstyleemphasis{numdiff} option is
also used. The \sphinxstyleemphasis{numdiff} option cannot be used with any other keyword. In this
situation, LAMMPS will also exit with an error.


\subsection{Default}
\label{\detokenize{compute_born_matrix:default}}
\sphinxAtStartPar
none


\bigskip\hrule\bigskip

\phantomsection\label{\detokenize{compute_born_matrix:vanworkum}}
\sphinxAtStartPar
\sphinxstylestrong{(Van Workum)} K. Van Workum et al., J. Chem. Phys. 125 144506 (2006)

\phantomsection\label{\detokenize{compute_born_matrix:clavier2}}
\sphinxAtStartPar
\sphinxstylestrong{(Clavier)} G. Clavier, N. Desbiens, E. Bourasseau, V. Lachet, N. Brusselle\sphinxhyphen{}Dupend and B. Rousseau, Mol Sim, 43, 1413 (2017).

\phantomsection\label{\detokenize{compute_born_matrix:zhen}}
\sphinxAtStartPar
\sphinxstylestrong{(Zhen)} Y. Zhen, C. Chu, Computer Physics Communications 183(2012)261\sphinxhyphen{}265

\sphinxstepscope

\index{compute centro/atom@\spxentry{compute centro/atom}}\ignorespaces 

\section{compute centro/atom command}
\label{\detokenize{compute_centro_atom:compute-centro-atom-command}}\label{\detokenize{compute_centro_atom:index-0}}\label{\detokenize{compute_centro_atom::doc}}

\subsection{Syntax}
\label{\detokenize{compute_centro_atom:syntax}}
\begin{sphinxVerbatim}[commandchars=\\\{\}]
\PYG{k}{compute }\PYG{n+nv+nvIdentifier}{ID}\PYG{+w}{ }\PYG{n+nv+nvIdentifier}{group\PYGZhy{}ID}\PYG{+w}{ }\PYG{n}{centro}\PYG{o}{/}\PYG{n}{atom}\PYG{+w}{ }\PYG{n}{lattice}\PYG{+w}{ }\PYG{n}{keyword}\PYG{+w}{ }\PYG{n}{value}\PYG{+w}{ }\PYG{n}{...}
\end{sphinxVerbatim}
\begin{itemize}
\item {} 
\sphinxAtStartPar
ID, group\sphinxhyphen{}ID are documented in {\hyperref[\detokenize{compute::doc}]{\sphinxcrossref{\DUrole{doc}{compute}}}} command

\item {} 
\sphinxAtStartPar
centro/atom = style name of this compute command

\item {} 
\sphinxAtStartPar
lattice = \sphinxstyleemphasis{fcc} or \sphinxstyleemphasis{bcc} or N = \# of neighbors per atom to include

\item {} 
\sphinxAtStartPar
zero or more keyword/value pairs may be appended

\item {} 
\sphinxAtStartPar
keyword = \sphinxstyleemphasis{axes}
\begin{sphinxalltt}
\sphinxstyleemphasis{axes} value = \sphinxstyleemphasis{no} or \sphinxstyleemphasis{yes}
  \sphinxstyleemphasis{no} = do not calculate 3 symmetry axes
  \sphinxstyleemphasis{yes} = calculate 3 symmetry axes
\end{sphinxalltt}

\end{itemize}


\subsection{Examples}
\label{\detokenize{compute_centro_atom:examples}}
\begin{sphinxVerbatim}[commandchars=\\\{\}]
\PYG{k}{compute }\PYG{n+nv+nvIdentifier}{1}\PYG{+w}{ }\PYG{n+nv+nvIdentifier}{all}\PYG{+w}{ }\PYG{n}{centro}\PYG{o}{/}\PYG{n}{atom}\PYG{+w}{ }\PYG{n}{fcc}

\PYG{k}{compute }\PYG{n+nv+nvIdentifier}{1}\PYG{+w}{ }\PYG{n+nv+nvIdentifier}{all}\PYG{+w}{ }\PYG{n}{centro}\PYG{o}{/}\PYG{n}{atom}\PYG{+w}{ }\PYG{l+m}{8}
\end{sphinxVerbatim}


\subsection{Description}
\label{\detokenize{compute_centro_atom:description}}
\sphinxAtStartPar
Define a computation that calculates the centro\sphinxhyphen{}symmetry parameter for
each atom in the group, for either FCC or BCC lattices, depending on
the choice of the \sphinxstyleemphasis{lattice} argument.  In solid\sphinxhyphen{}state systems the
centro\sphinxhyphen{}symmetry parameter is a useful measure of the local lattice
disorder around an atom and can be used to characterize whether the
atom is part of a perfect lattice, a local defect (e.g. a dislocation
or stacking fault), or at a surface.

\sphinxAtStartPar
The value of the centro\sphinxhyphen{}symmetry parameter will be 0.0 for atoms not
in the specified compute group.

\sphinxAtStartPar
This parameter is computed using the following formula from
{\hyperref[\detokenize{compute_centro_atom:kelchner}]{\sphinxcrossref{\DUrole{std,std-ref}{(Kelchner)}}}}
\begin{equation*}
\begin{split}CS = \sum_{i = 1}^{N/2} | \vec{R}_i + \vec{R}_{i+N/2} |^2\end{split}
\end{equation*}
\sphinxAtStartPar
where the \(N\) nearest neighbors of each atom are identified and
\(\vec{R}_i\) and \(\vec{R}_{i+N/2}\) are vectors from the
central atom to a particular pair of nearest neighbors.  There are
\(N (N-1)/2\) possible neighbor pairs that can contribute to this
formula.  The quantity in the sum is computed for each, and the
\(N/2\) smallest are used.  This will typically be for pairs of
atoms in symmetrically opposite positions with respect to the central
atom; hence the \(i+N/2\) notation.

\sphinxAtStartPar
\(N\) is an input parameter, which should be set to correspond to
the number of nearest neighbors in the underlying lattice of atoms.
If the keyword \sphinxstyleemphasis{fcc} or \sphinxstyleemphasis{bcc} is used, \sphinxstyleemphasis{N} is set to 12 and 8
respectively.  More generally, \sphinxstyleemphasis{N} can be set to a positive, even
integer.

\sphinxAtStartPar
For an atom on a lattice site, surrounded by atoms on a perfect
lattice, the centro\sphinxhyphen{}symmetry parameter will be 0.  It will be near 0
for small thermal perturbations of a perfect lattice.  If a point
defect exists, the symmetry is broken, and the parameter will be a
larger positive value.  An atom at a surface will have a large
positive parameter.  If the atom does not have \(N\) neighbors
(within the potential cutoff), then its centro\sphinxhyphen{}symmetry parameter is
set to 0.0.

\sphinxAtStartPar
If the keyword \sphinxstyleemphasis{axes} has the setting \sphinxstyleemphasis{yes}, then this compute also
estimates three symmetry axes for each atom’s local neighborhood.  The
first two of these are the vectors joining the two pairs of neighbor
atoms with smallest contributions to the centrosymmetry parameter,
i.e. the two most symmetric pairs of atoms.  The third vector is
normal to the first two by the right\sphinxhyphen{}hand rule.  All three vectors are
normalized to unit length.  For FCC crystals, the first two vectors
will lie along a \(\langle110\rangle\) direction, while the third vector
will lie along either a \(\langle100\rangle\) or \(\langle111\rangle\)
direction.  For HCP crystals, the first two vectors will lie along
\(\langle1000\rangle\) directions, while the third vector
will lie along \(\langle0001\rangle\).  This provides a simple way to
measure local orientation in HCP structures.  In general, the \sphinxstyleemphasis{axes} keyword
can be used to estimate the orientation of symmetry axes in the neighborhood
of any atom.

\sphinxAtStartPar
Only atoms within the cutoff of the pairwise neighbor list are
considered as possible neighbors.  Atoms not in the compute group are
included in the \(N\) neighbors used in this calculation.

\sphinxAtStartPar
The neighbor list needed to compute this quantity is constructed each
time the calculation is performed (e.g., each time a snapshot of atoms
is dumped).  Thus it can be inefficient to compute/dump this quantity
too frequently or to have multiple compute/dump commands, each with a
\sphinxstyleemphasis{centro/atom} style.


\subsection{Output info}
\label{\detokenize{compute_centro_atom:output-info}}
\sphinxAtStartPar
By default, this compute calculates the centrosymmetry value for each
atom as a per\sphinxhyphen{}atom vector, which can be accessed by any command that
uses per\sphinxhyphen{}atom values from a compute as input.  See the {\hyperref[\detokenize{Howto_output::doc}]{\sphinxcrossref{\DUrole{doc}{Howto output}}}} page for an overview of LAMMPS output
options.

\sphinxAtStartPar
If the \sphinxstyleemphasis{axes} keyword setting is \sphinxstyleemphasis{yes}, then a per\sphinxhyphen{}atom array is
calculated. The first column is the centrosymmetry parameter.  The
next three columns are the \sphinxstyleemphasis{x}, \sphinxstyleemphasis{y}, and \sphinxstyleemphasis{z} components of the first
symmetry axis, followed by the second, and third symmetry axes in
columns 5\textendash{}7 and 8\textendash{}10.

\sphinxAtStartPar
The centrosymmetry values are unitless values \(\ge 0.0\). Their magnitude
depends on the lattice style due to the number of contributing neighbor
pairs in the summation in the formula above.  And it depends on the
local defects surrounding the central atom, as described above.  For
the \sphinxstyleemphasis{axes yes} case, the vector components are also unitless, since
they represent spatial directions.

\sphinxAtStartPar
Here are typical centro\sphinxhyphen{}symmetry values, from a nanoindentation
simulation into gold (FCC).  These were provided by Jon Zimmerman
(Sandia):

\begin{sphinxVerbatim}[commandchars=\\\{\}]
Bulk lattice = 0
Dislocation core \PYGZti{} 1.0 (0.5 to 1.25)
Stacking faults \PYGZti{} 5.0 (4.0 to 6.0)
Free surface \PYGZti{} 23.0
\end{sphinxVerbatim}

\sphinxAtStartPar
These values are \sphinxstylestrong{not} normalized by the square of the lattice
parameter.  If they were, normalized values would be:

\begin{sphinxVerbatim}[commandchars=\\\{\}]
Bulk lattice = 0
Dislocation core \PYGZti{} 0.06 (0.03 to 0.075)
Stacking faults \PYGZti{} 0.3 (0.24 to 0.36)
Free surface \PYGZti{} 1.38
\end{sphinxVerbatim}

\sphinxAtStartPar
For BCC materials, the values for dislocation cores and free surfaces
would be somewhat different, due to their being only 8 neighbors instead
of 12.


\subsection{Restrictions}
\label{\detokenize{compute_centro_atom:restrictions}}\begin{quote}

\sphinxAtStartPar
none
\end{quote}


\subsection{Related commands}
\label{\detokenize{compute_centro_atom:related-commands}}
\sphinxAtStartPar
{\hyperref[\detokenize{compute_cna_atom::doc}]{\sphinxcrossref{\DUrole{doc}{compute cna/atom}}}}


\subsection{Default}
\label{\detokenize{compute_centro_atom:default}}
\sphinxAtStartPar
The default value for the optional keyword is axes = no.


\bigskip\hrule\bigskip

\phantomsection\label{\detokenize{compute_centro_atom:kelchner}}
\sphinxAtStartPar
\sphinxstylestrong{(Kelchner)} Kelchner, Plimpton, Hamilton, Phys Rev B, 58, 11085 (1998).

\sphinxstepscope

\index{compute chunk/atom@\spxentry{compute chunk/atom}}\ignorespaces 

\section{compute chunk/atom command}
\label{\detokenize{compute_chunk_atom:compute-chunk-atom-command}}\label{\detokenize{compute_chunk_atom:index-0}}\label{\detokenize{compute_chunk_atom::doc}}

\subsection{Syntax}
\label{\detokenize{compute_chunk_atom:syntax}}
\begin{sphinxVerbatim}[commandchars=\\\{\}]
\PYG{k}{compute }\PYG{n+nv+nvIdentifier}{ID}\PYG{+w}{ }\PYG{n+nv+nvIdentifier}{group\PYGZhy{}ID}\PYG{+w}{ }\PYG{n}{chunk}\PYG{o}{/}\PYG{n}{atom}\PYG{+w}{ }\PYG{n}{style}\PYG{+w}{ }\PYG{n}{args}\PYG{+w}{ }\PYG{n}{keyword}\PYG{+w}{ }\PYG{n}{values}\PYG{+w}{ }\PYG{n}{...}
\end{sphinxVerbatim}
\begin{itemize}
\item {} 
\sphinxAtStartPar
ID, group\sphinxhyphen{}ID are documented in {\hyperref[\detokenize{compute::doc}]{\sphinxcrossref{\DUrole{doc}{compute}}}} command

\item {} 
\sphinxAtStartPar
chunk/atom = style name of this compute command
\begin{sphinxalltt}
style = \sphinxstyleemphasis{bin/1d} or \sphinxstyleemphasis{bin/2d} or \sphinxstyleemphasis{bin/3d} or \sphinxstyleemphasis{bin/sphere} or \sphinxstyleemphasis{bin/cylinder} or \sphinxstyleemphasis{type} or \sphinxstyleemphasis{molecule} or c\_ID, c\_ID{[}I{]}, f\_ID, f\_ID{[}I{]}, v\_name
  \sphinxstyleemphasis{bin/1d} args = dim origin delta
    dim = \sphinxstyleemphasis{x} or \sphinxstyleemphasis{y} or \sphinxstyleemphasis{z}
    origin = \sphinxstyleemphasis{lower} or \sphinxstyleemphasis{center} or \sphinxstyleemphasis{upper} or coordinate value (distance units)
    delta = thickness of spatial bins in dim (distance units)
  \sphinxstyleemphasis{bin/2d} args = dim origin delta dim origin delta
    dim = \sphinxstyleemphasis{x} or \sphinxstyleemphasis{y} or \sphinxstyleemphasis{z}
    origin = \sphinxstyleemphasis{lower} or \sphinxstyleemphasis{center} or \sphinxstyleemphasis{upper} or coordinate value (distance units)
    delta = thickness of spatial bins in dim (distance units)
  \sphinxstyleemphasis{bin/3d} args = dim origin delta dim origin delta dim origin delta
    dim = \sphinxstyleemphasis{x} or \sphinxstyleemphasis{y} or \sphinxstyleemphasis{z}
    origin = \sphinxstyleemphasis{lower} or \sphinxstyleemphasis{center} or \sphinxstyleemphasis{upper} or coordinate value (distance units)
    delta = thickness of spatial bins in dim (distance units)
  \sphinxstyleemphasis{bin/sphere} args = xorig yorig zorig rmin rmax nsbin
    xorig,yorig,zorig = center point of sphere
    srmin,srmax = bin from sphere radius rmin to rmax
    nsbin = \# of spherical shell bins between rmin and rmax
  \sphinxstyleemphasis{bin/cylinder} args = dim origin delta c1 c2 rmin rmax ncbin
    dim = \sphinxstyleemphasis{x} or \sphinxstyleemphasis{y} or \sphinxstyleemphasis{z} = axis of cylinder axis
    origin = \sphinxstyleemphasis{lower} or \sphinxstyleemphasis{center} or \sphinxstyleemphasis{upper} or coordinate value (distance units)
    delta = thickness of spatial bins in dim (distance units)
    c1,c2 = coords of cylinder axis in other 2 dimensions (distance units)
    crmin,crmax = bin from cylinder radius rmin to rmax (distance units)
    ncbin = \# of concentric circle bins between rmin and rmax
  \sphinxstyleemphasis{type} args = none
  \sphinxstyleemphasis{molecule} args = none
  c\_ID, c\_ID{[}I{]}, f\_ID, f\_ID{[}I{]}, v\_name args = none
    c\_ID = per\sphinxhyphen{}atom vector calculated by a compute with ID
    c\_ID{[}I{]} = Ith column of per\sphinxhyphen{}atom array calculated by a compute with ID
    f\_ID = per\sphinxhyphen{}atom vector calculated by a fix with ID
    f\_ID{[}I{]} = Ith column of per\sphinxhyphen{}atom array calculated by a fix with ID
    v\_name = per\sphinxhyphen{}atom vector calculated by an atom\sphinxhyphen{}style variable with name
\end{sphinxalltt}

\item {} 
\sphinxAtStartPar
zero or more keyword/values pairs may be appended

\item {} 
\sphinxAtStartPar
keyword = \sphinxstyleemphasis{region} or \sphinxstyleemphasis{nchunk} or \sphinxstyleemphasis{limit} or \sphinxstyleemphasis{ids} or \sphinxstyleemphasis{compress} or \sphinxstyleemphasis{discard} or \sphinxstyleemphasis{bound} or \sphinxstyleemphasis{pbc} or \sphinxstyleemphasis{units}
\begin{sphinxalltt}
\sphinxstyleemphasis{region} value = region\sphinxhyphen{}ID
  region\sphinxhyphen{}ID = ID of region atoms must be in to be part of a chunk
\sphinxstyleemphasis{nchunk} value = \sphinxstyleemphasis{once} or \sphinxstyleemphasis{every}
  once = only compute the number of chunks once
  every = re\sphinxhyphen{}compute the number of chunks whenever invoked
\sphinxstyleemphasis{limit} values = 0 or Nc max or Nc exact
  0 = no limit on the number of chunks
  Nc max = limit number of chunks to be \textless{}= Nc
  Nc exact = set number of chunks to exactly Nc
\sphinxstyleemphasis{ids} value = \sphinxstyleemphasis{once} or \sphinxstyleemphasis{nfreq} or \sphinxstyleemphasis{every}
  once = assign chunk IDs to atoms only once, they persist thereafter
  nfreq = assign chunk IDs to atoms only once every Nfreq steps (if invoked by {\hyperref[\detokenize{fix_ave_chunk::doc}]{\sphinxcrossref{\DUrole{doc}{fix ave/chunk}}}} which sets Nfreq)
  every = assign chunk IDs to atoms whenever invoked
\sphinxstyleemphasis{compress} value = \sphinxstyleemphasis{yes} or \sphinxstyleemphasis{no}
  yes = compress chunk IDs to eliminate IDs with no atoms
  no = do not compress chunk IDs even if some IDs have no atoms
\sphinxstyleemphasis{discard} value = \sphinxstyleemphasis{yes} or \sphinxstyleemphasis{no} or \sphinxstyleemphasis{mixed}
  yes = discard atoms with out\sphinxhyphen{}of\sphinxhyphen{}range chunk IDs by assigning a chunk ID = 0
  no = keep atoms with out\sphinxhyphen{}of\sphinxhyphen{}range chunk IDs by assigning a valid chunk ID
  mixed = keep or discard such atoms according to spatial binning rule
\sphinxstyleemphasis{bound} values = x/y/z lo hi
  x/y/z = \sphinxstyleemphasis{x} or \sphinxstyleemphasis{y} or \sphinxstyleemphasis{z} to bound spatial bins in this dimension
  lo = \sphinxstyleemphasis{lower} or coordinate value (distance units)
  hi = \sphinxstyleemphasis{upper} or coordinate value (distance units)
\sphinxstyleemphasis{pbc} value = \sphinxstyleemphasis{no} or \sphinxstyleemphasis{yes}
  yes = use periodic distance for bin/sphere and bin/cylinder styles
\sphinxstyleemphasis{units} value = \sphinxstyleemphasis{box} or \sphinxstyleemphasis{lattice} or \sphinxstyleemphasis{reduced}
\end{sphinxalltt}

\end{itemize}


\subsection{Examples}
\label{\detokenize{compute_chunk_atom:examples}}
\begin{sphinxVerbatim}[commandchars=\\\{\}]
\PYG{k}{compute }\PYG{n+nv+nvIdentifier}{1}\PYG{+w}{ }\PYG{n+nv+nvIdentifier}{all}\PYG{+w}{ }\PYG{n}{chunk}\PYG{o}{/}\PYG{n}{atom}\PYG{+w}{ }\PYG{n}{type}
\PYG{k}{compute }\PYG{n+nv+nvIdentifier}{1}\PYG{+w}{ }\PYG{n+nv+nvIdentifier}{all}\PYG{+w}{ }\PYG{n}{chunk}\PYG{o}{/}\PYG{n}{atom}\PYG{+w}{ }\PYG{n}{bin}\PYG{o}{/}\PYG{l+m}{1}\PYG{n}{d}\PYG{+w}{ }\PYG{n}{z}\PYG{+w}{ }\PYG{n}{lower}\PYG{+w}{ }\PYG{l+m}{0.02}\PYG{+w}{ }\PYG{n}{units}\PYG{+w}{ }\PYG{n}{reduced}
\PYG{k}{compute }\PYG{n+nv+nvIdentifier}{1}\PYG{+w}{ }\PYG{n+nv+nvIdentifier}{all}\PYG{+w}{ }\PYG{n}{chunk}\PYG{o}{/}\PYG{n}{atom}\PYG{+w}{ }\PYG{n}{bin}\PYG{o}{/}\PYG{l+m}{2}\PYG{n}{d}\PYG{+w}{ }\PYG{n}{z}\PYG{+w}{ }\PYG{n}{lower}\PYG{+w}{ }\PYG{l+m}{1.0}\PYG{+w}{ }\PYG{n}{y}\PYG{+w}{ }\PYG{l+m}{0.0}\PYG{+w}{ }\PYG{l+m}{2.5}
\PYG{k}{compute }\PYG{n+nv+nvIdentifier}{1}\PYG{+w}{ }\PYG{n+nv+nvIdentifier}{all}\PYG{+w}{ }\PYG{n}{chunk}\PYG{o}{/}\PYG{n}{atom}\PYG{+w}{ }\PYG{n}{molecule}\PYG{+w}{ }\PYG{k}{region }\PYG{n+nv+nvIdentifier}{sphere}\PYG{+w}{ }\PYG{n}{nchunk}\PYG{+w}{ }\PYG{n}{once}\PYG{+w}{ }\PYG{n}{ids}\PYG{+w}{ }\PYG{n}{once}\PYG{+w}{ }\PYG{n}{compress}\PYG{+w}{ }\PYG{n}{yes}
\PYG{k}{compute }\PYG{n+nv+nvIdentifier}{1}\PYG{+w}{ }\PYG{n+nv+nvIdentifier}{all}\PYG{+w}{ }\PYG{n}{chunk}\PYG{o}{/}\PYG{n}{atom}\PYG{+w}{ }\PYG{n}{bin}\PYG{o}{/}\PYG{n}{sphere}\PYG{+w}{ }\PYG{l+m}{5}\PYG{+w}{ }\PYG{l+m}{5}\PYG{+w}{ }\PYG{l+m}{5}\PYG{+w}{ }\PYG{l+m}{2.0}\PYG{+w}{ }\PYG{l+m}{5.0}\PYG{+w}{ }\PYG{l+m}{5}\PYG{+w}{ }\PYG{n}{discard}\PYG{+w}{ }\PYG{n}{yes}
\PYG{k}{compute }\PYG{n+nv+nvIdentifier}{1}\PYG{+w}{ }\PYG{n+nv+nvIdentifier}{all}\PYG{+w}{ }\PYG{n}{chunk}\PYG{o}{/}\PYG{n}{atom}\PYG{+w}{ }\PYG{n}{bin}\PYG{o}{/}\PYG{n}{cylinder}\PYG{+w}{ }\PYG{n}{z}\PYG{+w}{ }\PYG{n}{lower}\PYG{+w}{ }\PYG{l+m}{2}\PYG{+w}{ }\PYG{l+m}{10}\PYG{+w}{ }\PYG{l+m}{10}\PYG{+w}{ }\PYG{l+m}{2.0}\PYG{+w}{ }\PYG{l+m}{5.0}\PYG{+w}{ }\PYG{l+m}{3}\PYG{+w}{ }\PYG{n}{discard}\PYG{+w}{ }\PYG{n}{yes}
\PYG{k}{compute }\PYG{n+nv+nvIdentifier}{1}\PYG{+w}{ }\PYG{n+nv+nvIdentifier}{all}\PYG{+w}{ }\PYG{n}{chunk}\PYG{o}{/}\PYG{n}{atom}\PYG{+w}{ }\PYG{n}{c\PYGZus{}cluster}
\end{sphinxVerbatim}


\subsection{Description}
\label{\detokenize{compute_chunk_atom:description}}
\sphinxAtStartPar
Define a computation that calculates an integer chunk ID from 1 to
Nchunk for each atom in the group.  Values of chunk IDs are determined
by the \sphinxstyleemphasis{style} of chunk, which can be based on atom type or molecule
ID or spatial binning or a per\sphinxhyphen{}atom property or value calculated by
another {\hyperref[\detokenize{compute::doc}]{\sphinxcrossref{\DUrole{doc}{compute}}}}, {\hyperref[\detokenize{fix::doc}]{\sphinxcrossref{\DUrole{doc}{fix}}}}, or {\hyperref[\detokenize{variable::doc}]{\sphinxcrossref{\DUrole{doc}{atom\sphinxhyphen{}style variable}}}}.  Per\sphinxhyphen{}atom chunk IDs can be used by other
computes with “chunk” in their style name, such as {\hyperref[\detokenize{compute_com_chunk::doc}]{\sphinxcrossref{\DUrole{doc}{compute com/chunk}}}} or {\hyperref[\detokenize{compute_msd_chunk::doc}]{\sphinxcrossref{\DUrole{doc}{compute msd/chunk}}}}.  Or they can be used by the {\hyperref[\detokenize{fix_ave_chunk::doc}]{\sphinxcrossref{\DUrole{doc}{fix ave/chunk}}}} command to sum and time average a
variety of per\sphinxhyphen{}atom properties over the atoms in each chunk.  Or they
can simply be accessed by any command that uses per\sphinxhyphen{}atom values from a
compute as input, as discussed on the {\hyperref[\detokenize{Howto_output::doc}]{\sphinxcrossref{\DUrole{doc}{Howto output}}}}
doc page.

\sphinxAtStartPar
See the {\hyperref[\detokenize{Howto_chunk::doc}]{\sphinxcrossref{\DUrole{doc}{Howto chunk}}}} page for an overview of how
this compute can be used with a variety of other commands to tabulate
properties of a simulation.  The page gives several examples of input
script commands that can be used to calculate interesting properties.

\sphinxAtStartPar
Conceptually it is important to realize that this compute does two
simple things.  First, it sets the value of \sphinxstyleemphasis{Nchunk} = the number of
chunks, which can be a constant value or change over time.  Second, it
assigns each atom to a chunk via a chunk ID.  Chunk IDs range from 1
to \sphinxstyleemphasis{Nchunk} inclusive; some chunks may have no atoms assigned to them.
Atoms that do not belong to any chunk are assigned a value of 0.  Note
that the two operations are not always performed together.  For
example, spatial bins can be setup once (which sets \sphinxstyleemphasis{Nchunk}), and
atoms assigned to those bins many times thereafter (setting their
chunk IDs).

\sphinxAtStartPar
All other commands in LAMMPS that use chunk IDs assume there are
\sphinxstyleemphasis{Nchunk} number of chunks, and that every atom is assigned to one of
those chunks, or not assigned to any chunk.

\sphinxAtStartPar
There are many options for specifying for how and when \sphinxstyleemphasis{Nchunk} is
calculated, and how and when chunk IDs are assigned to atoms.  The
details depend on the chunk \sphinxstyleemphasis{style} and its \sphinxstyleemphasis{args}, as well as
optional keyword settings.  They can also depend on whether a {\hyperref[\detokenize{fix_ave_chunk::doc}]{\sphinxcrossref{\DUrole{doc}{fix ave/chunk}}}} command is using this compute, since
that command requires \sphinxstyleemphasis{Nchunk} to remain static across windows of
timesteps it specifies, while it accumulates per\sphinxhyphen{}chunk averages.

\sphinxAtStartPar
The details are described below.


\bigskip\hrule\bigskip


\sphinxAtStartPar
The different chunk styles operate as follows.  For each style, how it
calculates \sphinxstyleemphasis{Nchunk} and assigns chunk IDs to atoms is explained.  Note
that using the optional keywords can change both of those actions, as
described further below where the keywords are discussed.


\bigskip\hrule\bigskip


\sphinxAtStartPar
The \sphinxstyleemphasis{binning} styles perform a spatial binning of atoms, and assign an
atom the chunk ID corresponding to the bin number it is in.  \sphinxstyleemphasis{Nchunk}
is set to the number of bins, which can change if the simulation box
size changes.  This also depends on the setting of the \sphinxstyleemphasis{units}
keyword (e.g., for \sphinxstyleemphasis{reduced} units the number of chunks may not change
even if the box size does).

\sphinxAtStartPar
The \sphinxstyleemphasis{bin/1d}, \sphinxstyleemphasis{bin/2d}, and \sphinxstyleemphasis{bin/3d} styles define bins as 1d layers
(slabs), 2d pencils, or 3d boxes.  The \sphinxstyleemphasis{dim}, \sphinxstyleemphasis{origin}, and \sphinxstyleemphasis{delta}
settings are specified 1, 2, or 3 times.  For 2d or 3d bins, there is
no restriction on specifying dim = \sphinxstyleemphasis{x} before dim = \sphinxstyleemphasis{y} or \sphinxstyleemphasis{z}, or dim = \sphinxstyleemphasis{y}
before dim = \sphinxstyleemphasis{z}.  Bins in a particular \sphinxstyleemphasis{dim} have a bin size in that
dimension given by \sphinxstyleemphasis{delta}.  In each dimension, bins are defined
relative to a specified \sphinxstyleemphasis{origin}, which may be the lower/upper edge of
the simulation box (in that dimension), or its center point, or a
specified coordinate value.  Starting at the origin, sufficient bins
are created in both directions to completely span the simulation box
or the bounds specified by the optional \sphinxstyleemphasis{bounds} keyword.

\sphinxAtStartPar
For orthogonal simulation boxes, the bins are layers, pencils, or
boxes aligned with the xyz coordinate axes.  For triclinic
(non\sphinxhyphen{}orthogonal) simulation boxes, the bin faces are parallel to the
tilted faces of the simulation box.  See the {\hyperref[\detokenize{Howto_triclinic::doc}]{\sphinxcrossref{\DUrole{doc}{Howto triclinic}}}} page for a discussion of the
geometry of triclinic boxes in LAMMPS.  As described there, a tilted
simulation box has edge vectors \(\vec a\), \(\vec b\), and
\(\vec c\).  In that nomenclature, bins in
the \sphinxstyleemphasis{x} dimension have faces with normals in the \(\vec b \times \vec c\)
direction, bins in \sphinxstyleemphasis{y} have faces normal to the \(\vec a \times \vec c\)
direction, and bins in \sphinxstyleemphasis{z} have faces normal to the \(\vec a \times \vec b\)
direction.  Note that in order to define the size and position of
these bins in an unambiguous fashion, the \sphinxstyleemphasis{units} option must be set
to \sphinxstyleemphasis{reduced} when using a triclinic simulation box, as noted below.

\sphinxAtStartPar
The meaning of \sphinxstyleemphasis{origin} and \sphinxstyleemphasis{delta} for triclinic boxes is as follows.
Consider a triclinic box with bins that are 1d layers or slabs in the
x dimension.  No matter how the box is tilted, an \sphinxstyleemphasis{origin} of 0.0
means start layers at the lower \(\vec b \times \vec c\) plane of the
simulation box and an \sphinxstyleemphasis{origin} of 1.0 means to start layers at the upper
\(\vec b \times \vec c\) face of the box.  A \sphinxstyleemphasis{delta} value of 0.1 in
\sphinxstyleemphasis{reduced} units means there will be 10 layers from 0.0 to 1.0, regardless of
the current size or shape of the simulation box.

\sphinxAtStartPar
The \sphinxstyleemphasis{bin/sphere} style defines a set of spherical shell bins around
the origin (\sphinxstyleemphasis{xorig},\sphinxstyleemphasis{yorig},\sphinxstyleemphasis{zorig}), using \sphinxstyleemphasis{nsbin} bins with radii
equally spaced between \sphinxstyleemphasis{srmin} and \sphinxstyleemphasis{srmax}.  This is effectively a 1d
vector of bins.  For example, if \sphinxstyleemphasis{srmin} = 1.0 and \sphinxstyleemphasis{srmax} = 10.0 and
\sphinxstyleemphasis{nsbin} = 9, then the first bin spans \(1.0 < r < 2.0\), and the last bin
spans \(9.0 < r < 10.0\).  The geometry of the bins is the same whether the
simulation box is orthogonal or triclinic (i.e., the spherical shells
are not tilted or scaled differently in different dimensions to
transform them into ellipsoidal shells).

\sphinxAtStartPar
The \sphinxstyleemphasis{bin/cylinder} style defines bins for a cylinder oriented along
the axis \sphinxstyleemphasis{dim} with the axis coordinates in the other two radial
dimensions at (\sphinxstyleemphasis{c1},\sphinxstyleemphasis{c2}).  For dim = \sphinxstyleemphasis{x}, \(c_1/c_2 = y/z\);
for dim = \sphinxstyleemphasis{y}, \(c_1/c_2 = x/z\); for dim = \sphinxstyleemphasis{z},
\(c_1/c_2 = x/y\).  This is effectively a 2d array of bins.  The first
dimension is along the cylinder axis, the second dimension is radially outward
from the cylinder axis.  The bin size and positions along the cylinder axis are
specified by the \sphinxstyleemphasis{origin} and \sphinxstyleemphasis{delta} values, the same as for the \sphinxstyleemphasis{bin/1d},
\sphinxstyleemphasis{bin/2d}, and \sphinxstyleemphasis{bin/3d} styles.  There are \sphinxstyleemphasis{ncbin} concentric circle bins in the
radial direction from the cylinder axis with radii equally spaced
between \sphinxstyleemphasis{crmin} and \sphinxstyleemphasis{crmax}.  For example, if \sphinxstyleemphasis{crmin} = 1.0 and
\sphinxstyleemphasis{crmax} = 10.0 and \sphinxstyleemphasis{ncbin} = 9, then the first bin spans \(1.0 < r < 2.0\)
and the last bin spans \(9.0 < r < 10.0\).  The geometry of the bins in
the radial dimensions is the same whether the simulation box is
orthogonal or triclinic (i.e., the concentric circles are not tilted or
scaled differently in the two different dimensions to transform them
into ellipses).

\sphinxAtStartPar
The created bins (and hence the chunk IDs) are numbered consecutively
from 1 to the number of bins = \sphinxstyleemphasis{Nchunk}.  For \sphinxstyleemphasis{bin2d} and \sphinxstyleemphasis{bin3d}, the
numbering varies most rapidly in the first dimension (which could be
\sphinxstyleemphasis{x}, \sphinxstyleemphasis{y}, or \sphinxstyleemphasis{z}), next rapidly in the second dimension, and most slowly in the
third dimension.  For \sphinxstyleemphasis{bin/sphere}, the bin with smallest radii is chunk
1 and the bin with largest radii is chunk Nchunk = \sphinxstyleemphasis{ncbin}.  For
\sphinxstyleemphasis{bin/cylinder}, the numbering varies most rapidly in the dimension
along the cylinder axis and most slowly in the radial direction.

\sphinxAtStartPar
Each time this compute is invoked, each atom is mapped to a bin based
on its current position.  Note that between reneighboring timesteps,
atoms can move outside the current simulation box.  If the box is
periodic (in that dimension) the atom is remapping into the periodic
box for purposes of binning.  If the box in not periodic, the atom may
have moved outside the bounds of all bins.  If an atom is not inside
any bin, the \sphinxstyleemphasis{discard} keyword is used to determine how a chunk ID is
assigned to the atom.


\bigskip\hrule\bigskip


\sphinxAtStartPar
The \sphinxstyleemphasis{type} style uses the atom type as the chunk ID.  \sphinxstyleemphasis{Nchunk} is set
to the number of atom types defined for the simulation (e.g., via the
{\hyperref[\detokenize{create_box::doc}]{\sphinxcrossref{\DUrole{doc}{create\_box}}}} or {\hyperref[\detokenize{read_data::doc}]{\sphinxcrossref{\DUrole{doc}{read\_data}}}} commands).


\bigskip\hrule\bigskip


\sphinxAtStartPar
The \sphinxstyleemphasis{molecule} style uses the molecule ID of each atom as its chunk
ID.  \sphinxstyleemphasis{Nchunk} is set to the largest chunk ID.  Note that this excludes
molecule IDs for atoms which are not in the specified group or
optional region.

\sphinxAtStartPar
There is no requirement that all atoms in a particular molecule are
assigned the same chunk ID (zero or non\sphinxhyphen{}zero), though you probably
want that to be the case, if you wish to compute a per\sphinxhyphen{}molecule
property.  LAMMPS will issue a warning if that is not the case, but
only the first time that \sphinxstyleemphasis{Nchunk} is calculated.

\sphinxAtStartPar
Note that atoms with a molecule ID = 0, which may be non\sphinxhyphen{}molecular
solvent atoms, have an out\sphinxhyphen{}of\sphinxhyphen{}range chunk ID.  These atoms are
discarded (not assigned to any chunk) or assigned to \sphinxstyleemphasis{Nchunk},
depending on the value of the \sphinxstyleemphasis{discard} keyword.


\bigskip\hrule\bigskip


\sphinxAtStartPar
The \sphinxstyleemphasis{compute/fix/variable} styles set the chunk ID of each atom based
on a quantity calculated and stored by a compute, fix, or variable.
In each case, it must be a per\sphinxhyphen{}atom quantity.  In each case the
referenced floating point values are converted to an integer chunk ID
as follows.  The floating point value is truncated (rounded down) to
an integer value.  If the integer value is \(\le 0\), then a chunk ID of 0
is assigned to the atom.  If the integer value is \(> 0\), it becomes the
chunk ID to the atom.  \sphinxstyleemphasis{Nchunk} is set to the largest chunk ID.  Note
that this excludes atoms which are not in the specified group or
optional region.

\sphinxAtStartPar
If the style begins with “c\_”, a compute ID must follow which has been
previously defined in the input script.  If no bracketed integer is
appended, the per\sphinxhyphen{}atom vector calculated by the compute is used.  If a
bracketed integer is appended, the Ith column of the per\sphinxhyphen{}atom array
calculated by the compute is used.  Users can also write code for
their own compute styles and {\hyperref[\detokenize{Modify::doc}]{\sphinxcrossref{\DUrole{doc}{add them to LAMMPS}}}}.

\sphinxAtStartPar
If the style begins with “f\_”, a fix ID must follow which has been
previously defined in the input script.  If no bracketed integer is
appended, the per\sphinxhyphen{}atom vector calculated by the fix is used.  If a
bracketed integer is appended, the Ith column of the per\sphinxhyphen{}atom array
calculated by the fix is used.  Note that some fixes only produce
their values on certain timesteps, which must be compatible with the
timestep on which this compute accesses the fix, else an error
results.  Users can also write code for their own fix styles and {\hyperref[\detokenize{Modify::doc}]{\sphinxcrossref{\DUrole{doc}{add them to LAMMPS}}}}.

\sphinxAtStartPar
If a value begins with “v\_”, a variable name for an \sphinxstyleemphasis{atom} or
\sphinxstyleemphasis{atomfile} style {\hyperref[\detokenize{variable::doc}]{\sphinxcrossref{\DUrole{doc}{variable}}}} must follow which has been
previously defined in the input script.  Variables of style \sphinxstyleemphasis{atom} can
reference thermodynamic keywords and various per\sphinxhyphen{}atom attributes, or
invoke other computes, fixes, or variables when they are evaluated, so
this is a very general means of generating per\sphinxhyphen{}atom quantities to
treat as a chunk ID.


\bigskip\hrule\bigskip


\sphinxAtStartPar
Normally, \sphinxstyleemphasis{Nchunk} = the number of chunks, is re\sphinxhyphen{}calculated every time
this fix is invoked, though the value may or may not change.  As
explained below, the \sphinxstyleemphasis{nchunk} keyword can be set to \sphinxstyleemphasis{once} which means
\sphinxstyleemphasis{Nchunk} will never change.

\sphinxAtStartPar
If a {\hyperref[\detokenize{fix_ave_chunk::doc}]{\sphinxcrossref{\DUrole{doc}{fix ave/chunk}}}} command uses this compute, it
can also turn off the re\sphinxhyphen{}calculation of \sphinxstyleemphasis{Nchunk} for one or more
windows of timesteps.  The extent of the windows, during which Nchunk
is held constant, are determined by the \sphinxstyleemphasis{Nevery}, \sphinxstyleemphasis{Nrepeat}, \sphinxstyleemphasis{Nfreq}
values and the \sphinxstyleemphasis{ave} keyword setting that are used by the {\hyperref[\detokenize{fix_ave_chunk::doc}]{\sphinxcrossref{\DUrole{doc}{fix ave/chunk}}}} command.

\sphinxAtStartPar
Specifically, if \sphinxstyleemphasis{ave} = \sphinxstyleemphasis{one}, then for each span of \sphinxstyleemphasis{Nfreq}
timesteps, \sphinxstyleemphasis{Nchunk} is held constant between the first timestep when
averaging is done (within the Nfreq\sphinxhyphen{}length window), and the last
timestep when averaging is done (multiple of Nfreq).  If \sphinxstyleemphasis{ave} =
\sphinxstyleemphasis{running} or \sphinxstyleemphasis{window}, then \sphinxstyleemphasis{Nchunk} is held constant forever,
starting on the first timestep when the {\hyperref[\detokenize{fix_ave_chunk::doc}]{\sphinxcrossref{\DUrole{doc}{fix ave/chunk}}}} command invokes this compute.

\sphinxAtStartPar
Note that multiple {\hyperref[\detokenize{fix_ave_chunk::doc}]{\sphinxcrossref{\DUrole{doc}{fix ave/chunk}}}} commands can use
the same compute chunk/atom compute.  However, the time windows they
induce for holding \sphinxstyleemphasis{Nchunk} constant must be identical, else an error
will be generated.


\bigskip\hrule\bigskip


\sphinxAtStartPar
The various optional keywords operate as follows.  Note that some of
them function differently or are ignored by different chunk styles.
Some of them also have different default values, depending on
the chunk style, as listed below.

\sphinxAtStartPar
The \sphinxstyleemphasis{region} keyword applies to all chunk styles.  If used, an atom
must be in both the specified group and the specified geometric
{\hyperref[\detokenize{region::doc}]{\sphinxcrossref{\DUrole{doc}{region}}}} to be assigned to a chunk.


\bigskip\hrule\bigskip


\sphinxAtStartPar
The \sphinxstyleemphasis{nchunk} keyword applies to all chunk styles.  It specifies how
often \sphinxstyleemphasis{Nchunk} is recalculated, which in turn can affect the chunk IDs
assigned to individual atoms.

\sphinxAtStartPar
If \sphinxstyleemphasis{nchunk} is set to \sphinxstyleemphasis{once}, then \sphinxstyleemphasis{Nchunk} is only calculated once,
the first time this compute is invoked.  If \sphinxstyleemphasis{nchunk} is set to
\sphinxstyleemphasis{every}, then \sphinxstyleemphasis{Nchunk} is re\sphinxhyphen{}calculated every time the compute is
invoked.  Note that, as described above, the use of this compute
by the {\hyperref[\detokenize{fix_ave_chunk::doc}]{\sphinxcrossref{\DUrole{doc}{fix ave/chunk}}}} command can override
the \sphinxstyleemphasis{every} setting.

\sphinxAtStartPar
The default values for \sphinxstyleemphasis{nchunk} are listed below and depend on the
chunk style and other system and keyword settings.  They attempt to
represent typical use cases for the various chunk styles.  The
\sphinxstyleemphasis{nchunk} value can always be set explicitly if desired.


\bigskip\hrule\bigskip


\sphinxAtStartPar
The \sphinxstyleemphasis{limit} keyword can be used to limit the calculated value of
\sphinxstyleemphasis{Nchunk} = the number of chunks.  The limit is applied each time
\sphinxstyleemphasis{Nchunk} is calculated, which also limits the chunk IDs assigned to
any atom.  The \sphinxstyleemphasis{limit} keyword is used by all chunk styles except the
\sphinxstyleemphasis{binning} styles, which ignore it.  This is because the number of bins
can be tailored using the \sphinxstyleemphasis{bound} keyword (described below) which
effectively limits the size of \sphinxstyleemphasis{Nchunk}.

\sphinxAtStartPar
If \sphinxstyleemphasis{limit} is set to \sphinxstyleemphasis{Nc} = 0, then no limit is imposed on \sphinxstyleemphasis{Nchunk},
though the \sphinxstyleemphasis{compress} keyword can still be used to reduce \sphinxstyleemphasis{Nchunk}, as
described below.

\sphinxAtStartPar
If \sphinxstyleemphasis{Nc} \(>\) 0, then the effect of the \sphinxstyleemphasis{limit} keyword depends on whether
the \sphinxstyleemphasis{compress} keyword is also used with a setting of \sphinxstyleemphasis{yes}, and
whether the \sphinxstyleemphasis{compress} keyword is specified before the \sphinxstyleemphasis{limit} keyword
or after.

\sphinxAtStartPar
In all cases, \sphinxstyleemphasis{Nchunk} is first calculated in the usual way for each
chunk style, as described above.

\sphinxAtStartPar
First, here is what occurs if \sphinxstyleemphasis{compress yes} is not set.  If \sphinxstyleemphasis{limit}
is set to \sphinxstyleemphasis{Nc max}, then \sphinxstyleemphasis{Nchunk} is reset to the smaller of \sphinxstyleemphasis{Nchunk}
and \sphinxstyleemphasis{Nc}.  If \sphinxstyleemphasis{limit} is set to \sphinxstyleemphasis{Nc exact}, then \sphinxstyleemphasis{Nchunk} is reset to
\sphinxstyleemphasis{Nc}, whether the original \sphinxstyleemphasis{Nchunk} was larger or smaller than \sphinxstyleemphasis{Nc}.
If \sphinxstyleemphasis{Nchunk} shrank due to the \sphinxstyleemphasis{limit} setting, then atom chunk IDs \(>\)
\sphinxstyleemphasis{Nchunk} will be reset to 0 or \sphinxstyleemphasis{Nchunk}, depending on the setting of
the \sphinxstyleemphasis{discard} keyword.  If \sphinxstyleemphasis{Nchunk} grew, there will simply be some
chunks with no atoms assigned to them.

\sphinxAtStartPar
If \sphinxstyleemphasis{compress yes} is set, and the \sphinxstyleemphasis{compress} keyword comes before the
\sphinxstyleemphasis{limit} keyword, the compression operation is performed first, as
described below, which resets \sphinxstyleemphasis{Nchunk}.  The \sphinxstyleemphasis{limit} keyword is then
applied to the new \sphinxstyleemphasis{Nchunk} value, exactly as described in the
preceding paragraph.  Note that in this case, all atoms will end up
with chunk IDs \(\le\) \sphinxstyleemphasis{Nc}, but their original values (e.g., molecule ID
or compute/fix/variable) may have been \(>\) \sphinxstyleemphasis{Nc}, because of the
compression operation.

\sphinxAtStartPar
If \sphinxstyleemphasis{compress yes} is set, and the \sphinxstyleemphasis{compress} keyword comes after the
\sphinxstyleemphasis{limit} keyword, then the \sphinxstyleemphasis{limit} value of \sphinxstyleemphasis{Nc} is applied first to
the uncompressed value of \sphinxstyleemphasis{Nchunk}, but only if \sphinxstyleemphasis{Nc} \(<\) \sphinxstyleemphasis{Nchunk}
(whether \sphinxstyleemphasis{Nc max} or \sphinxstyleemphasis{Nc exact} is used).  This effectively means all
atoms with chunk IDs \(>\) \sphinxstyleemphasis{Nc} have their chunk IDs reset to 0 or \sphinxstyleemphasis{Nc},
depending on the setting of the \sphinxstyleemphasis{discard} keyword.  The compression
operation is then performed, which may shrink \sphinxstyleemphasis{Nchunk} further.  If
the new \sphinxstyleemphasis{Nchunk} \(<\) \sphinxstyleemphasis{Nc} and \sphinxstyleemphasis{limit} = \sphinxstyleemphasis{Nc exact} is specified, then
\sphinxstyleemphasis{Nchunk} is reset to \sphinxstyleemphasis{Nc}, which results in extra chunks with no atoms
assigned to them.  Note that in this case, all atoms will end up with
chunk IDs \(\le\) \sphinxstyleemphasis{Nc}, and their original values (e.g., molecule ID or
compute/fix/variable value) will also have been \(\le\) \sphinxstyleemphasis{Nc}.


\bigskip\hrule\bigskip


\sphinxAtStartPar
The \sphinxstyleemphasis{ids} keyword applies to all chunk styles.  If the setting is
\sphinxstyleemphasis{once} then the chunk IDs assigned to atoms the first time this
compute is invoked will be permanent, and never be re\sphinxhyphen{}computed.

\sphinxAtStartPar
If the setting is \sphinxstyleemphasis{nfreq} and if a {\hyperref[\detokenize{fix_ave_chunk::doc}]{\sphinxcrossref{\DUrole{doc}{fix ave/chunk}}}}
command is using this compute, then in each of the \sphinxstyleemphasis{Nchunk} = constant
time windows (discussed above), the chunk ID’s assigned to atoms on
the first step of the time window will persist until the end of the
time window.

\sphinxAtStartPar
If the setting is \sphinxstyleemphasis{every}, which is the default, then chunk IDs are
re\sphinxhyphen{}calculated on any timestep this compute is invoked.

\begin{sphinxadmonition}{note}{Note:}
\sphinxAtStartPar
If you want the persistent chunk\sphinxhyphen{}IDs calculated by this compute
to be continuous when running from a {\hyperref[\detokenize{read_restart::doc}]{\sphinxcrossref{\DUrole{doc}{restart file}}}},
then you should use the same ID for this compute, as in the original
run.  This is so that the fix this compute creates to store per\sphinxhyphen{}atom
quantities will also have the same ID, and thus be initialized
correctly with chunk IDs from the restart file.
\end{sphinxadmonition}


\bigskip\hrule\bigskip


\sphinxAtStartPar
The \sphinxstyleemphasis{compress} keyword applies to all chunk styles and affects how
\sphinxstyleemphasis{Nchunk} is calculated, which in turn affects the chunk IDs assigned
to each atom.  It is useful for converting a “sparse” set of chunk IDs
(with many IDs that have no atoms assigned to them), into a “dense”
set of IDs, where every chunk has one or more atoms assigned to it.

\sphinxAtStartPar
Two possible use cases are as follows.  If a large simulation box is
mostly empty space, then the \sphinxstyleemphasis{binning} style may produce many bins
with no atoms.  If \sphinxstyleemphasis{compress} is set to \sphinxstyleemphasis{yes}, only bins with atoms
will be contribute to \sphinxstyleemphasis{Nchunk}.  Likewise, the \sphinxstyleemphasis{molecule} or
\sphinxstyleemphasis{compute/fix/variable} styles may produce large \sphinxstyleemphasis{Nchunk} values.  For
example, the {\hyperref[\detokenize{compute_cluster_atom::doc}]{\sphinxcrossref{\DUrole{doc}{compute cluster/atom}}}} command
assigns every atom an atom ID for one of the atoms it is clustered
with.  For a million\sphinxhyphen{}atom system with 5 clusters, there would only be
5 unique chunk IDs, but the largest chunk ID might be 1 million,
resulting in \sphinxstyleemphasis{Nchunk} = 1 million.  If \sphinxstyleemphasis{compress} is set to \sphinxstyleemphasis{yes},
\sphinxstyleemphasis{Nchunk} will be reset to 5.

\sphinxAtStartPar
If \sphinxstyleemphasis{compress} is set to \sphinxstyleemphasis{no}, which is the default, no compression is
done.  If it is set to \sphinxstyleemphasis{yes}, all chunk IDs with no atoms are removed
from the list of chunk IDs, and the list is sorted.  The remaining
chunk IDs are renumbered from 1 to \sphinxstyleemphasis{Nchunk} where \sphinxstyleemphasis{Nchunk} is the new
length of the list.  The chunk IDs assigned to each atom reflect
the new renumbering from 1 to \sphinxstyleemphasis{Nchunk}.

\sphinxAtStartPar
The original chunk IDs (before renumbering) can be accessed by the
{\hyperref[\detokenize{compute_property_chunk::doc}]{\sphinxcrossref{\DUrole{doc}{compute property/chunk}}}} command and its
\sphinxstyleemphasis{id} keyword, or by the {\hyperref[\detokenize{fix_ave_chunk::doc}]{\sphinxcrossref{\DUrole{doc}{fix ave/chunk}}}} command
which outputs the original IDs as one of the columns in its global
output array.  For example, using the “compute cluster/atom” command
discussed above, the original 5 unique chunk IDs might be atom IDs
(27,4982,58374,857838,1000000).  After compression, these will be
renumbered to (1,2,3,4,5).  The original values (27,…,1000000) can
be output to a file by the {\hyperref[\detokenize{fix_ave_chunk::doc}]{\sphinxcrossref{\DUrole{doc}{fix ave/chunk}}}} command,
or by using the {\hyperref[\detokenize{fix_ave_time::doc}]{\sphinxcrossref{\DUrole{doc}{fix ave/time}}}} command in
conjunction with the {\hyperref[\detokenize{compute_property_chunk::doc}]{\sphinxcrossref{\DUrole{doc}{compute property/chunk}}}} command.

\begin{sphinxadmonition}{note}{Note:}
\sphinxAtStartPar
The compression operation requires global communication across
all processors to share their chunk ID values.  It can require large
memory on every processor to store them, even after they are
compressed, if there are a large number of unique chunk IDs with
atoms assigned to them.  It uses a STL map to find unique chunk IDs
and store them in sorted order.  Each time an atom is assigned a
compressed chunk ID, it must access the STL map.  All of this means
that compression can be expensive, both in memory and CPU time.  The
use of the \sphinxstyleemphasis{limit} keyword in conjunction with the \sphinxstyleemphasis{compress} keyword
can affect these costs, depending on which keyword is used first.  So
use this option with care.
\end{sphinxadmonition}


\bigskip\hrule\bigskip


\sphinxAtStartPar
The \sphinxstyleemphasis{discard} keyword applies to all chunk styles.  It affects what
chunk IDs are assigned to atoms that do not match one of the valid
chunk IDs from 1 to \sphinxstyleemphasis{Nchunk}.  Note that it does not apply to atoms
that are not in the specified group or optionally specified region.
Those atoms are always assigned a chunk ID = 0.

\sphinxAtStartPar
If the calculated chunk ID for an atom is not within the range 1 to
\sphinxstyleemphasis{Nchunk} then it is a “discard” atom.  Note that \sphinxstyleemphasis{Nchunk} may have
been shrunk by the \sphinxstyleemphasis{limit} keyword.  Or the \sphinxstyleemphasis{compress} keyword may
have eliminated chunk IDs that were valid before the compression took
place, and are now not in the compressed list.  Also note that for the
\sphinxstyleemphasis{molecule} chunk style, if new molecules are added to the system,
their chunk IDs may exceed a previously calculated \sphinxstyleemphasis{Nchunk}.
Likewise, evaluation of a compute/fix/variable on a later timestep may
return chunk IDs that are invalid for the previously calculated
\sphinxstyleemphasis{Nchunk}.

\sphinxAtStartPar
All the chunk styles except the \sphinxstyleemphasis{binning} styles, must use \sphinxstyleemphasis{discard}
set to either \sphinxstyleemphasis{yes} or \sphinxstyleemphasis{no}.  If \sphinxstyleemphasis{discard} is set to \sphinxstyleemphasis{yes}, which is
the default, then every “discard” atom has its chunk ID set to 0.  If
\sphinxstyleemphasis{discard} is set to \sphinxstyleemphasis{no}, every “discard” atom has its chunk ID set to
\sphinxstyleemphasis{Nchunk}.  I.e. it becomes part of the last chunk.

\sphinxAtStartPar
The \sphinxstyleemphasis{binning} styles use the \sphinxstyleemphasis{discard} keyword to decide whether to
discard atoms outside the spatial domain covered by bins, or to assign
them to the bin they are nearest to.

\sphinxAtStartPar
For the \sphinxstyleemphasis{bin/1d}, \sphinxstyleemphasis{bin/2d}, \sphinxstyleemphasis{bin/3d} styles the details are as
follows.  If \sphinxstyleemphasis{discard} is set to \sphinxstyleemphasis{yes}, an out\sphinxhyphen{}of\sphinxhyphen{}domain atom will
have its chunk ID set to 0.  If \sphinxstyleemphasis{discard} is set to \sphinxstyleemphasis{no}, the atom
will have its chunk ID set to the first or last bin in that dimension.
If \sphinxstyleemphasis{discard} is set to \sphinxstyleemphasis{mixed}, which is the default, it will only
have its chunk ID set to the first or last bin if bins extend to the
simulation box boundary in that dimension.  This is the case if the
\sphinxstyleemphasis{bound} keyword settings are \sphinxstyleemphasis{lower} and \sphinxstyleemphasis{upper}, which is the
default.  If the \sphinxstyleemphasis{bound} keyword settings are numeric values, then the
atom will have its chunk ID set to 0 if it is outside the bounds of
any bin.  Note that in this case, it is possible that the first or
last bin extends beyond the numeric \sphinxstyleemphasis{bounds} settings, depending on
the specified \sphinxstyleemphasis{origin}.  If this is the case, the chunk ID of the atom
is only set to 0 if it is outside the first or last bin, not if it is
simply outside the numeric \sphinxstyleemphasis{bounds} setting.

\sphinxAtStartPar
For the \sphinxstyleemphasis{bin/sphere} style the details are as follows.  If \sphinxstyleemphasis{discard}
is set to \sphinxstyleemphasis{yes}, an out\sphinxhyphen{}of\sphinxhyphen{}domain atom will have its chunk ID set to
0.  If \sphinxstyleemphasis{discard} is set to \sphinxstyleemphasis{no} or \sphinxstyleemphasis{mixed}, the atom will have its
chunk ID set to the first or last bin, i.e. the innermost or outermost
spherical shell.  If the distance of the atom from the origin is less
than \sphinxstyleemphasis{rmin}, it will be assigned to the first bin.  If the distance of
the atom from the origin is greater than \sphinxstyleemphasis{rmax}, it will be assigned
to the last bin.

\sphinxAtStartPar
For the \sphinxstyleemphasis{bin/cylinder} style the details are as follows.  If \sphinxstyleemphasis{discard}
is set to \sphinxstyleemphasis{yes}, an out\sphinxhyphen{}of\sphinxhyphen{}domain atom will have its chunk ID set to
0.  If \sphinxstyleemphasis{discard} is set to \sphinxstyleemphasis{no}, the atom will have its chunk ID set
to the first or last bin in both the radial and axis dimensions.  If
\sphinxstyleemphasis{discard} is set to \sphinxstyleemphasis{mixed}, which is the default, the radial
dimension is treated the same as for \sphinxstyleemphasis{discard} = no.  But for the axis
dimension, it will only have its chunk ID set to the first or last
bin if bins extend to the simulation box boundary in the axis
dimension.  This is the case if the \sphinxstyleemphasis{bound} keyword settings are
\sphinxstyleemphasis{lower} and \sphinxstyleemphasis{upper}, which is the default.  If the \sphinxstyleemphasis{bound} keyword
settings are numeric values, then the atom will have its chunk ID set
to 0 if it is outside the bounds of any bin.  Note that in this case,
it is possible that the first or last bin extends beyond the numeric
\sphinxstyleemphasis{bounds} settings, depending on the specified \sphinxstyleemphasis{origin}.  If this is
the case, the chunk ID of the atom is only set to 0 if it is outside
the first or last bin, not if it is simply outside the numeric
\sphinxstyleemphasis{bounds} setting.

\sphinxAtStartPar
If \sphinxstyleemphasis{discard} is set to \sphinxstyleemphasis{no} or \sphinxstyleemphasis{mixed}, the atom will have its
chunk ID set to the first or last bin, i.e. the innermost or outermost
spherical shell.  If the distance of the atom from the origin is less
than \sphinxstyleemphasis{rmin}, it will be assigned to the first bin.  If the distance of
the atom from the origin is greater than \sphinxstyleemphasis{rmax}, it will be assigned
to the last bin.


\bigskip\hrule\bigskip


\sphinxAtStartPar
The \sphinxstyleemphasis{bound} keyword only applies to the \sphinxstyleemphasis{bin/1d}, \sphinxstyleemphasis{bin/2d}, \sphinxstyleemphasis{bin/3d}
styles and to the axis dimension of the \sphinxstyleemphasis{bin/cylinder} style;
otherwise it is ignored.  It can be used one or more times to limit
the extent of bin coverage in a specified dimension, i.e. to only bin
a portion of the box.  If the \sphinxstyleemphasis{lo} setting is \sphinxstyleemphasis{lower} or the \sphinxstyleemphasis{hi}
setting is \sphinxstyleemphasis{upper}, the bin extent in that direction extends to the
box boundary.  If a numeric value is used for \sphinxstyleemphasis{lo} and/or \sphinxstyleemphasis{hi}, then
the bin extent in the \sphinxstyleemphasis{lo} or \sphinxstyleemphasis{hi} direction extends only to that
value, which is assumed to be inside (or at least near) the simulation
box boundaries, though LAMMPS does not check for this.  Note that
using the \sphinxstyleemphasis{bound} keyword typically reduces the total number of bins
and thus the number of chunks \sphinxstyleemphasis{Nchunk}.

\sphinxAtStartPar
The \sphinxstyleemphasis{pbc} keyword only applies to the \sphinxstyleemphasis{bin/sphere} and \sphinxstyleemphasis{bin/cylinder}
styles.  If set to \sphinxstyleemphasis{yes}, the distance an atom is from the sphere
origin or cylinder axis is calculated in a minimum image sense with
respect to periodic dimensions, when determining which bin the atom is
in.  I.e. if x is a periodic dimension and the distance between the
atom and the sphere center in the x dimension is greater than 0.5 *
simulation box length in x, then a box length is subtracted to give a
distance \textless{} 0.5 * simulation box length.  This allosws the sphere or
cylinder center to be near a box edge, and atoms on the other side of
the periodic box will still be close to the center point/axis.  Note
that with a setting of \sphinxstyleemphasis{yes}, the outer sphere or cylinder radius must
also be \textless{}= 0.5 * simulation box length in any periodic dimension
except for the cylinder axis dimension, or an error is generated.

\sphinxAtStartPar
The \sphinxstyleemphasis{units} keyword only applies to the \sphinxstyleemphasis{binning} styles; otherwise it
is ignored.  For the \sphinxstyleemphasis{bin/1d}, \sphinxstyleemphasis{bin/2d}, \sphinxstyleemphasis{bin/3d} styles, it
determines the meaning of the distance units used for the bin sizes
\sphinxstyleemphasis{delta} and for \sphinxstyleemphasis{origin} and \sphinxstyleemphasis{bounds} values if they are coordinate
values.  For the \sphinxstyleemphasis{bin/sphere} style it determines the meaning of the
distance units used for \sphinxstyleemphasis{xorig},\sphinxstyleemphasis{yorig},\sphinxstyleemphasis{zorig} and the radii \sphinxstyleemphasis{srmin}
and \sphinxstyleemphasis{srmax}.  For the \sphinxstyleemphasis{bin/cylinder} style it determines the meaning
of the distance units used for \sphinxstyleemphasis{delta},\sphinxstyleemphasis{c1},\sphinxstyleemphasis{c2} and the radii \sphinxstyleemphasis{crmin}
and \sphinxstyleemphasis{crmax}.

\sphinxAtStartPar
For orthogonal simulation boxes, any of the 3 options may
be used.  For non\sphinxhyphen{}orthogonal (triclinic) simulation boxes, only the
\sphinxstyleemphasis{reduced} option may be used.

\sphinxAtStartPar
A \sphinxstyleemphasis{box} value selects standard distance units as defined by the
{\hyperref[\detokenize{units::doc}]{\sphinxcrossref{\DUrole{doc}{units}}}} command (e.g., \(\AA\) for units = \sphinxstyleemphasis{real} or \sphinxstyleemphasis{metal}).
A \sphinxstyleemphasis{lattice} value means the distance units are in lattice spacings.
The {\hyperref[\detokenize{lattice::doc}]{\sphinxcrossref{\DUrole{doc}{lattice}}}} command must have been previously used to
define the lattice spacing.  A \sphinxstyleemphasis{reduced} value means normalized
unitless values between 0 and 1, which represent the lower and upper
faces of the simulation box respectively.  Thus an \sphinxstyleemphasis{origin} value of
0.5 means the center of the box in any dimension.  A \sphinxstyleemphasis{delta} value of
0.1 means 10 bins span the box in that dimension.

\sphinxAtStartPar
Note that for the \sphinxstyleemphasis{bin/sphere} style, the radii \sphinxstyleemphasis{srmin} and \sphinxstyleemphasis{srmax} are
scaled by the lattice spacing or reduced value of the \sphinxstyleemphasis{x} dimension.

\sphinxAtStartPar
Note that for the \sphinxstyleemphasis{bin/cylinder} style, the radii \sphinxstyleemphasis{crmin} and \sphinxstyleemphasis{crmax}
are scaled by the lattice spacing or reduced value of the first
dimension perpendicular to the cylinder axis (e.g., \sphinxstyleemphasis{y} for an \sphinxstyleemphasis{x}\sphinxhyphen{}axis
cylinder, \sphinxstyleemphasis{x} for a \sphinxstyleemphasis{y}\sphinxhyphen{}axis cylinder, and \sphinxstyleemphasis{x} for a \sphinxstyleemphasis{z}\sphinxhyphen{}axis cylinder).


\bigskip\hrule\bigskip



\subsection{Output info}
\label{\detokenize{compute_chunk_atom:output-info}}
\sphinxAtStartPar
This compute calculates a per\sphinxhyphen{}atom vector (the chunk ID), which can
be accessed by any command that uses per\sphinxhyphen{}atom values from a compute
as input.  It also calculates a global scalar (the number of chunks),
which can be similarly accessed everywhere outside of a per\sphinxhyphen{}atom context.
See the {\hyperref[\detokenize{Howto_output::doc}]{\sphinxcrossref{\DUrole{doc}{Howto output}}}} page for an overview of
LAMMPS output options.

\sphinxAtStartPar
The per\sphinxhyphen{}atom vector values are unitless chunk IDs, ranging from 1 to
\sphinxstyleemphasis{Nchunk} (inclusive) for atoms assigned to chunks, and 0 for atoms not
belonging to a chunk.  The scalar contains the value of \sphinxstyleemphasis{Nchunk}.


\subsection{Restrictions}
\label{\detokenize{compute_chunk_atom:restrictions}}
\sphinxAtStartPar
Even if the \sphinxstyleemphasis{nchunk} keyword is set to \sphinxstyleemphasis{once}, the chunk IDs assigned
to each atom are not stored in a restart files.  This means you cannot
expect those assignments to persist in a restarted simulation.
Instead you must re\sphinxhyphen{}specify this command and assign atoms to chunks when
the restarted simulation begins.


\subsection{Related commands}
\label{\detokenize{compute_chunk_atom:related-commands}}
\sphinxAtStartPar
{\hyperref[\detokenize{fix_ave_chunk::doc}]{\sphinxcrossref{\DUrole{doc}{fix ave/chunk}}}},
{\hyperref[\detokenize{compute_global_atom::doc}]{\sphinxcrossref{\DUrole{doc}{compute global/atom}}}}


\subsection{Default}
\label{\detokenize{compute_chunk_atom:default}}
\sphinxAtStartPar
The option defaults are as follows:
\begin{itemize}
\item {} 
\sphinxAtStartPar
region = none

\item {} 
\sphinxAtStartPar
nchunk = every, if compress is yes, overriding other defaults listed here

\item {} 
\sphinxAtStartPar
nchunk = once, for type style

\item {} 
\sphinxAtStartPar
nchunk = once, for mol style if region is none

\item {} 
\sphinxAtStartPar
nchunk = every, for mol style if region is set

\item {} 
\sphinxAtStartPar
nchunk = once, for binning style if the simulation box size is static or units = reduced

\item {} 
\sphinxAtStartPar
nchunk = every, for binning style if the simulation box size is dynamic and units is lattice or box

\item {} 
\sphinxAtStartPar
nchunk = every, for compute/fix/variable style

\item {} 
\sphinxAtStartPar
limit = 0

\item {} 
\sphinxAtStartPar
ids = every

\item {} 
\sphinxAtStartPar
compress = no

\item {} 
\sphinxAtStartPar
discard = yes, for all styles except binning

\item {} 
\sphinxAtStartPar
discard = mixed, for binning styles

\item {} 
\sphinxAtStartPar
bound = lower and upper in all dimensions

\item {} 
\sphinxAtStartPar
pbc = no

\item {} 
\sphinxAtStartPar
units = lattice

\end{itemize}

\sphinxstepscope

\index{compute chunk/spread/atom@\spxentry{compute chunk/spread/atom}}\ignorespaces 

\section{compute chunk/spread/atom command}
\label{\detokenize{compute_chunk_spread_atom:compute-chunk-spread-atom-command}}\label{\detokenize{compute_chunk_spread_atom:index-0}}\label{\detokenize{compute_chunk_spread_atom::doc}}

\subsection{Syntax}
\label{\detokenize{compute_chunk_spread_atom:syntax}}
\begin{sphinxVerbatim}[commandchars=\\\{\}]
\PYG{k}{compute }\PYG{n+nv+nvIdentifier}{ID}\PYG{+w}{ }\PYG{n+nv+nvIdentifier}{group\PYGZhy{}ID}\PYG{+w}{ }\PYG{n}{chunk}\PYG{o}{/}\PYG{n}{spread}\PYG{o}{/}\PYG{n}{atom}\PYG{+w}{ }\PYG{n}{chunkID}\PYG{+w}{ }\PYG{n}{input1}\PYG{+w}{ }\PYG{n}{input2}\PYG{+w}{ }\PYG{n}{...}
\end{sphinxVerbatim}
\begin{itemize}
\item {} 
\sphinxAtStartPar
ID, group\sphinxhyphen{}ID are documented in {\hyperref[\detokenize{compute::doc}]{\sphinxcrossref{\DUrole{doc}{compute}}}} command

\item {} 
\sphinxAtStartPar
chunk/spread/atom = style name of this compute command

\item {} 
\sphinxAtStartPar
chunkID = ID of {\hyperref[\detokenize{compute_chunk_atom::doc}]{\sphinxcrossref{\DUrole{doc}{compute chunk/atom}}}} command

\item {} 
\sphinxAtStartPar
one or more inputs can be listed

\item {} 
\sphinxAtStartPar
input = c\_ID, c\_ID{[}N{]}, f\_ID, f\_ID{[}N{]}

\begin{sphinxVerbatim}[commandchars=\\\{\}]
c\PYGZus{}ID = global vector calculated by a compute with ID
c\PYGZus{}ID[I] = Ith column of global array calculated by a compute with ID, I can include wildcard (see below)
f\PYGZus{}ID = global vector calculated by a fix with ID
f\PYGZus{}ID[I] = Ith column of global array calculated by a fix with ID, I can include wildcard (see below)
\end{sphinxVerbatim}

\end{itemize}


\subsection{Examples}
\label{\detokenize{compute_chunk_spread_atom:examples}}
\begin{sphinxVerbatim}[commandchars=\\\{\}]
\PYG{k}{compute }\PYG{n+nv+nvIdentifier}{1}\PYG{+w}{ }\PYG{n+nv+nvIdentifier}{all}\PYG{+w}{ }\PYG{n}{chunk}\PYG{o}{/}\PYG{n}{spread}\PYG{o}{/}\PYG{n}{atom}\PYG{+w}{ }\PYG{n}{mychunk}\PYG{+w}{ }\PYG{n}{c\PYGZus{}com[}\PYG{o}{*}\PYG{n}{]}\PYG{+w}{ }\PYG{n}{c\PYGZus{}gyration}
\end{sphinxVerbatim}


\subsection{Description}
\label{\detokenize{compute_chunk_spread_atom:description}}
\sphinxAtStartPar
Define a calculation that “spreads” one or more per\sphinxhyphen{}chunk values to
each atom in the chunk.  This can be useful in several scenarios:
\begin{itemize}
\item {} 
\sphinxAtStartPar
For creating a {\hyperref[\detokenize{dump::doc}]{\sphinxcrossref{\DUrole{doc}{dump file}}}} where each atom lists info about
the chunk it is in, e.g. for post\sphinxhyphen{}processing purposes.

\item {} 
\sphinxAtStartPar
To access chunk value in {\hyperref[\detokenize{variable::doc}]{\sphinxcrossref{\DUrole{doc}{atom\sphinxhyphen{}style variables}}}} that
need info about the chunk each atom is in.

\item {} 
\sphinxAtStartPar
To use the {\hyperref[\detokenize{fix_ave_chunk::doc}]{\sphinxcrossref{\DUrole{doc}{fix ave/chunk}}}} command to spatially
average per\sphinxhyphen{}chunk values calculated by a per\sphinxhyphen{}chunk compute.

\end{itemize}

\sphinxAtStartPar
Examples are given below.

\sphinxAtStartPar
In LAMMPS, chunks are collections of atoms defined by a {\hyperref[\detokenize{compute_chunk_atom::doc}]{\sphinxcrossref{\DUrole{doc}{compute chunk/atom}}}} command, which assigns each atom
to a single chunk (or no chunk).  The ID for this command is specified
as chunkID.  For example, a single chunk could be the atoms in a
molecule or atoms in a spatial bin.  See the {\hyperref[\detokenize{compute_chunk_atom::doc}]{\sphinxcrossref{\DUrole{doc}{compute chunk/atom}}}} and {\hyperref[\detokenize{Howto_chunk::doc}]{\sphinxcrossref{\DUrole{doc}{Howto chunk}}}}
doc pages for details of how chunks can be defined and examples of how
they can be used to measure properties of a system.

\sphinxAtStartPar
For inputs that are computes, they must be a compute that calculates
per\sphinxhyphen{}chunk values.  These are computes whose style names end in
“/chunk”.

\sphinxAtStartPar
For inputs that are fixes, they should be a fix that calculates
per\sphinxhyphen{}chunk values.  For example, {\hyperref[\detokenize{fix_ave_chunk::doc}]{\sphinxcrossref{\DUrole{doc}{fix ave/chunk}}}} or
{\hyperref[\detokenize{fix_ave_time::doc}]{\sphinxcrossref{\DUrole{doc}{fix ave/time}}}} (assuming it is time\sphinxhyphen{}averaging
per\sphinxhyphen{}chunk data).

\sphinxAtStartPar
For each atom, this compute accesses its chunk ID from the specified
\sphinxstyleemphasis{chunkID} compute, then accesses the per\sphinxhyphen{}chunk value in each input.
Those values are copied to this compute to become the output for that
atom.

\sphinxAtStartPar
The values generated by this compute will be 0.0 for atoms not in the
specified compute group \sphinxstyleemphasis{group\sphinxhyphen{}ID}.  They will also be 0.0 if the atom
is not in a chunk, as assigned by the \sphinxstyleemphasis{chunkID} compute.  They will
also be 0.0 if the current chunk ID for the atom is out\sphinxhyphen{}of\sphinxhyphen{}bounds with
respect to the number of chunks stored by a particular input compute
or fix.

\begin{sphinxadmonition}{note}{Note:}
\sphinxAtStartPar
LAMMPS does not check that a compute or fix which calculates
per\sphinxhyphen{}chunk values uses the same definition of chunks as this compute.
It’s up to you to be consistent.  Likewise, for a fix input, LAMMPS
does not check that it is per\sphinxhyphen{}chunk data.  It only checks that the fix
produces a global vector or array.
\end{sphinxadmonition}


\bigskip\hrule\bigskip


\sphinxAtStartPar
Each listed input is operated on independently.

\sphinxAtStartPar
If a bracketed index I is used, it can be specified using a wildcard
asterisk with the index to effectively specify multiple values.  This
takes the form “*” or “*n” or “n*” or “m*n”.  If N = the number of
columns in the array, then an asterisk with no numeric values means
all indices from 1 to N.  A leading asterisk means all indices from 1
to n (inclusive).  A trailing asterisk means all indices from n to N
(inclusive).  A middle asterisk means all indices from m to n
(inclusive).

\sphinxAtStartPar
Using a wildcard is the same as if the individual columns of the array
had been listed one by one.  E.g. these 2 compute chunk/spread/atom
commands are equivalent, since the {\hyperref[\detokenize{compute_com_chunk::doc}]{\sphinxcrossref{\DUrole{doc}{compute com/chunk}}}} command creates a per\sphinxhyphen{}atom array
with 3 columns:

\begin{sphinxVerbatim}[commandchars=\\\{\}]
\PYG{k}{compute }\PYG{n+nv+nvIdentifier}{com}\PYG{+w}{ }\PYG{n+nv+nvIdentifier}{all}\PYG{+w}{ }\PYG{n}{com}\PYG{o}{/}\PYG{n}{chunk}\PYG{+w}{ }\PYG{n}{mychunk}
\PYG{k}{compute }\PYG{n+nv+nvIdentifier}{10}\PYG{+w}{ }\PYG{n+nv+nvIdentifier}{all}\PYG{+w}{ }\PYG{n}{chunk}\PYG{o}{/}\PYG{n}{spread}\PYG{o}{/}\PYG{n}{atom}\PYG{+w}{ }\PYG{n}{mychunk}\PYG{+w}{ }\PYG{n}{c\PYGZus{}com[}\PYG{o}{*}\PYG{n}{]}
\PYG{k}{compute }\PYG{n+nv+nvIdentifier}{10}\PYG{+w}{ }\PYG{n+nv+nvIdentifier}{all}\PYG{+w}{ }\PYG{n}{chunk}\PYG{o}{/}\PYG{n}{spread}\PYG{o}{/}\PYG{n}{atom}\PYG{+w}{ }\PYG{n}{mychunk}\PYG{+w}{ }\PYG{n}{c\PYGZus{}com[1]}\PYG{+w}{ }\PYG{n}{c\PYGZus{}com[2]}\PYG{+w}{ }\PYG{n}{c\PYGZus{}com[3]}
\end{sphinxVerbatim}


\bigskip\hrule\bigskip


\sphinxAtStartPar
Here is an example of writing a dump file the with the center\sphinxhyphen{}of\sphinxhyphen{}mass
(COM) for the chunk each atom is in.  The commands below can be added
to the bench/in.chain script.

\begin{sphinxVerbatim}[commandchars=\\\{\}]
\PYG{k}{compute         }\PYG{n+nv+nvIdentifier}{cmol}\PYG{+w}{ }\PYG{n+nv+nvIdentifier}{all}\PYG{+w}{ }\PYG{n}{chunk}\PYG{o}{/}\PYG{n}{atom}\PYG{+w}{ }\PYG{n}{molecule}
\PYG{k}{compute         }\PYG{n+nv+nvIdentifier}{com}\PYG{+w}{ }\PYG{n+nv+nvIdentifier}{all}\PYG{+w}{ }\PYG{n}{com}\PYG{o}{/}\PYG{n}{chunk}\PYG{+w}{ }\PYG{n}{cmol}
\PYG{k}{compute         }\PYG{n+nv+nvIdentifier}{comchunk}\PYG{+w}{ }\PYG{n+nv+nvIdentifier}{all}\PYG{+w}{ }\PYG{n}{chunk}\PYG{o}{/}\PYG{n}{spread}\PYG{o}{/}\PYG{n}{atom}\PYG{+w}{ }\PYG{n}{cmol}\PYG{+w}{ }\PYG{n}{c\PYGZus{}com[}\PYG{o}{*}\PYG{n}{]}
\PYG{k}{dump            }\PYG{n+nv+nvIdentifier}{1}\PYG{+w}{ }\PYG{n+nv+nvIdentifier}{all}\PYG{+w}{ }\PYG{n}{custom}\PYG{+w}{ }\PYG{l+m}{50}\PYG{+w}{ }\PYG{n}{tmp.dump}\PYG{+w}{ }\PYG{n}{id}\PYG{+w}{ }\PYG{n}{mol}\PYG{+w}{ }\PYG{n}{type}\PYG{+w}{ }\PYG{n}{x}\PYG{+w}{ }\PYG{n}{y}\PYG{+w}{ }\PYG{n}{z}\PYG{+w}{ }\PYG{n}{c\PYGZus{}comchunk[}\PYG{o}{*}\PYG{n}{]}
\PYG{k}{dump\PYGZus{}modify     }\PYG{n+nv+nvIdentifier}{1}\PYG{+w}{ }\PYG{n}{sort}\PYG{+w}{ }\PYG{n}{id}
\end{sphinxVerbatim}

\sphinxAtStartPar
The same per\sphinxhyphen{}chunk data for each atom could be used to define per\sphinxhyphen{}atom
forces for the {\hyperref[\detokenize{fix_addforce::doc}]{\sphinxcrossref{\DUrole{doc}{fix addforce}}}} command.  In this
example the forces act to pull atoms of an extended polymer chain
towards its COM in an attractive manner.

\begin{sphinxVerbatim}[commandchars=\\\{\}]
\PYG{k}{compute         }\PYG{n+nv+nvIdentifier}{prop}\PYG{+w}{ }\PYG{n+nv+nvIdentifier}{all}\PYG{+w}{ }\PYG{n}{property}\PYG{o}{/}\PYG{n}{atom}\PYG{+w}{ }\PYG{n}{xu}\PYG{+w}{ }\PYG{n}{yu}\PYG{+w}{ }\PYG{n}{zu}
\PYG{k}{variable        }\PYG{n+nv+nvIdentifier}{k}\PYG{+w}{ }\PYG{n}{equal}\PYG{+w}{ }\PYG{l+m}{0.1}
\PYG{k}{variable        }\PYG{n+nv+nvIdentifier}{fx}\PYG{+w}{ }\PYG{n}{atom}\PYG{+w}{ }\PYG{n}{v\PYGZus{}k}\PYG{o}{*}\PYG{n+nv}{(}\PYG{n+nv}{c\PYGZus{}comchunk[1]\PYGZhy{}c\PYGZus{}prop[1]}\PYG{n+nv}{)}
\PYG{k}{variable        }\PYG{n+nv+nvIdentifier}{fy}\PYG{+w}{ }\PYG{n}{atom}\PYG{+w}{ }\PYG{n}{v\PYGZus{}k}\PYG{o}{*}\PYG{n+nv}{(}\PYG{n+nv}{c\PYGZus{}comchunk[2]\PYGZhy{}c\PYGZus{}prop[2]}\PYG{n+nv}{)}
\PYG{k}{variable        }\PYG{n+nv+nvIdentifier}{fz}\PYG{+w}{ }\PYG{n}{atom}\PYG{+w}{ }\PYG{n}{v\PYGZus{}k}\PYG{o}{*}\PYG{n+nv}{(}\PYG{n+nv}{c\PYGZus{}comchunk[3]\PYGZhy{}c\PYGZus{}prop[3]}\PYG{n+nv}{)}
\PYG{k}{fix             }\PYG{n+nv+nvIdentifier}{3}\PYG{+w}{ }\PYG{n+nv+nvIdentifier}{all}\PYG{+w}{ }\PYG{n}{addforce}\PYG{+w}{ }\PYG{n}{v\PYGZus{}fx}\PYG{+w}{ }\PYG{n}{v\PYGZus{}fy}\PYG{+w}{ }\PYG{n}{v\PYGZus{}fz}
\end{sphinxVerbatim}

\sphinxAtStartPar
Note that {\hyperref[\detokenize{compute_property_atom::doc}]{\sphinxcrossref{\DUrole{doc}{compute property/atom}}}} is used
to generate unwrapped coordinates for use in the per\sphinxhyphen{}atom force
calculation, so that the effect of periodic boundaries is accounted
for properly.

\sphinxAtStartPar
Over time this applied force could shrink each polymer chain’s radius
of gyration in a polymer mixture simulation.  Here is output from the
bench/in.chain script.  Thermo output is shown for 1000 steps, where
the last column is the average radius of gyration over all 320 chains
in the 32000 atom system:

\begin{sphinxVerbatim}[commandchars=\\\{\}]
compute         gyr all gyration/chunk cmol
variable        ave equal ave(c\PYGZus{}gyr)
thermo\PYGZus{}style    custom step etotal press v\PYGZus{}ave

       0    22.394765    4.6721833     5.128278
     100    22.445002    4.8166709    5.0348372
     200    22.500128    4.8790392    4.9364875
     300    22.534686    4.9183766    4.8590693
     400    22.557196    4.9492211    4.7937849
     500    22.571017    4.9161853    4.7412008
     600    22.573944    5.0229708    4.6931243
     700    22.581804    5.0541301    4.6440647
     800    22.584683    4.9691734    4.6000016
     900     22.59128    5.0247538    4.5611513
    1000    22.586832      4.94697    4.5238362
\end{sphinxVerbatim}


\bigskip\hrule\bigskip


\sphinxAtStartPar
Here is an example for using one set of chunks, defined for molecules,
to compute the dipole moment vector for each chunk.  E.g. for water
molecules. Then spreading those values to each atom in each chunk.
Then defining a second set of chunks based on spatial bins.  And
finally, using the {\hyperref[\detokenize{fix_ave_chunk::doc}]{\sphinxcrossref{\DUrole{doc}{fix ave/chunk}}}} command to
calculate an average dipole moment vector per spatial bin.

\begin{sphinxVerbatim}[commandchars=\\\{\}]
\PYG{k}{compute       }\PYG{n+nv+nvIdentifier}{cmol}\PYG{+w}{ }\PYG{n+nv+nvIdentifier}{all}\PYG{+w}{ }\PYG{n}{chunk}\PYG{o}{/}\PYG{n}{atom}\PYG{+w}{ }\PYG{n}{molecule}
\PYG{k}{compute       }\PYG{n+nv+nvIdentifier}{dipole}\PYG{+w}{ }\PYG{n+nv+nvIdentifier}{all}\PYG{+w}{ }\PYG{n}{dipole}\PYG{o}{/}\PYG{n}{chunk}\PYG{+w}{ }\PYG{n}{cmol}
\PYG{k}{compute       }\PYG{n+nv+nvIdentifier}{spread}\PYG{+w}{ }\PYG{n+nv+nvIdentifier}{all}\PYG{+w}{ }\PYG{n}{chunk}\PYG{o}{/}\PYG{n}{spread}\PYG{o}{/}\PYG{n}{atom}\PYG{+w}{ }\PYG{n}{cmol}\PYG{+w}{ }\PYG{n}{c\PYGZus{}dipole[1]}\PYG{+w}{ }\PYG{n}{c\PYGZus{}dipole[2]}\PYG{+w}{ }\PYG{n}{c\PYGZus{}dipole[3]}
\PYG{k}{compute       }\PYG{n+nv+nvIdentifier}{cspatial}\PYG{+w}{ }\PYG{n+nv+nvIdentifier}{all}\PYG{+w}{ }\PYG{n}{chunk}\PYG{o}{/}\PYG{n}{atom}\PYG{+w}{ }\PYG{n}{bin}\PYG{o}{/}\PYG{l+m}{1}\PYG{n}{d}\PYG{+w}{ }\PYG{n}{z}\PYG{+w}{ }\PYG{n}{lower}\PYG{+w}{ }\PYG{l+m}{0.1}\PYG{+w}{ }\PYG{n}{units}\PYG{+w}{ }\PYG{n}{reduced}
\PYG{k}{fix           }\PYG{n+nv+nvIdentifier}{ave}\PYG{+w}{ }\PYG{n+nv+nvIdentifier}{all}\PYG{+w}{ }\PYG{n}{ave}\PYG{o}{/}\PYG{n}{chunk}\PYG{+w}{ }\PYG{l+m}{100}\PYG{+w}{ }\PYG{l+m}{10}\PYG{+w}{ }\PYG{l+m}{1000}\PYG{+w}{ }\PYG{n}{cspatial}\PYG{+w}{ }\PYG{n}{c\PYGZus{}spread[}\PYG{o}{*}\PYG{n}{]}
\end{sphinxVerbatim}

\sphinxAtStartPar
Note that the {\hyperref[\detokenize{fix_ave_chunk::doc}]{\sphinxcrossref{\DUrole{doc}{fix ave/chunk}}}} command requires
per\sphinxhyphen{}atom values as input.  That is why the compute chunk/spread/atom
command is used to assign per\sphinxhyphen{}chunk values to each atom in the chunk.
If a molecule straddles bin boundaries, each of its atoms contributes
in a weighted manner to the average dipole moment of the spatial bin
it is in.


\bigskip\hrule\bigskip



\subsection{Output info}
\label{\detokenize{compute_chunk_spread_atom:output-info}}
\sphinxAtStartPar
This compute calculates a per\sphinxhyphen{}atom vector or array, which can be
accessed by any command that uses per\sphinxhyphen{}atom values from a compute as
input.  See the {\hyperref[\detokenize{Howto_output::doc}]{\sphinxcrossref{\DUrole{doc}{Howto output}}}} page for an
overview of LAMMPS output options.

\sphinxAtStartPar
The output is a per\sphinxhyphen{}atom vector if a single input value is specified,
otherwise a per\sphinxhyphen{}atom array is output.  The number of columns in the
array is the number of inputs provided.  The per\sphinxhyphen{}atom values for the
vector or each column of the array will be in whatever
{\hyperref[\detokenize{units::doc}]{\sphinxcrossref{\DUrole{doc}{units}}}} the corresponding input value is in.

\sphinxAtStartPar
The vector or array values are “intensive”.


\subsection{Restrictions}
\label{\detokenize{compute_chunk_spread_atom:restrictions}}\begin{quote}

\sphinxAtStartPar
none
\end{quote}


\subsection{Related commands}
\label{\detokenize{compute_chunk_spread_atom:related-commands}}
\sphinxAtStartPar
{\hyperref[\detokenize{compute_chunk_atom::doc}]{\sphinxcrossref{\DUrole{doc}{compute chunk/atom}}}}, {\hyperref[\detokenize{fix_ave_chunk::doc}]{\sphinxcrossref{\DUrole{doc}{fix ave/chunk}}}}, {\hyperref[\detokenize{compute_reduce_chunk::doc}]{\sphinxcrossref{\DUrole{doc}{compute reduce/chunk}}}}


\subsection{Default}
\label{\detokenize{compute_chunk_spread_atom:default}}
\sphinxAtStartPar
none

\sphinxstepscope

\index{compute cluster/atom@\spxentry{compute cluster/atom}}\ignorespaces 
\index{compute fragment/atom@\spxentry{compute fragment/atom}}\ignorespaces 
\index{compute aggregate/atom@\spxentry{compute aggregate/atom}}\ignorespaces 

\section{compute cluster/atom command}
\label{\detokenize{compute_cluster_atom:compute-cluster-atom-command}}\label{\detokenize{compute_cluster_atom:index-2}}\label{\detokenize{compute_cluster_atom:index-1}}\label{\detokenize{compute_cluster_atom:index-0}}\label{\detokenize{compute_cluster_atom::doc}}

\section{compute fragment/atom command}
\label{\detokenize{compute_cluster_atom:compute-fragment-atom-command}}

\section{compute aggregate/atom command}
\label{\detokenize{compute_cluster_atom:compute-aggregate-atom-command}}

\subsection{Syntax}
\label{\detokenize{compute_cluster_atom:syntax}}
\begin{sphinxVerbatim}[commandchars=\\\{\}]
\PYG{k}{compute }\PYG{n+nv+nvIdentifier}{ID}\PYG{+w}{ }\PYG{n+nv+nvIdentifier}{group\PYGZhy{}ID}\PYG{+w}{ }\PYG{n}{cluster}\PYG{o}{/}\PYG{n}{atom}\PYG{+w}{ }\PYG{n}{cutoff}
\PYG{k}{compute }\PYG{n+nv+nvIdentifier}{ID}\PYG{+w}{ }\PYG{n+nv+nvIdentifier}{group\PYGZhy{}ID}\PYG{+w}{ }\PYG{n}{fragment}\PYG{o}{/}\PYG{n}{atom}\PYG{+w}{ }\PYG{n}{keyword}\PYG{+w}{ }\PYG{n}{value}\PYG{+w}{ }\PYG{n}{...}
\PYG{k}{compute }\PYG{n+nv+nvIdentifier}{ID}\PYG{+w}{ }\PYG{n+nv+nvIdentifier}{group\PYGZhy{}ID}\PYG{+w}{ }\PYG{n}{aggregate}\PYG{o}{/}\PYG{n}{atom}\PYG{+w}{ }\PYG{n}{cutoff}
\end{sphinxVerbatim}
\begin{itemize}
\item {} 
\sphinxAtStartPar
ID, group\sphinxhyphen{}ID are documented in {\hyperref[\detokenize{compute::doc}]{\sphinxcrossref{\DUrole{doc}{compute}}}} command

\item {} 
\sphinxAtStartPar
\sphinxstyleemphasis{cluster/atom} or \sphinxstyleemphasis{fragment/atom} or \sphinxstyleemphasis{aggregate/atom} = style name of this compute command

\item {} 
\sphinxAtStartPar
cutoff = distance within which to label atoms as part of same cluster (distance units)

\item {} 
\sphinxAtStartPar
zero or more keyword/value pairs may be appended to \sphinxstyleemphasis{fragment/atom}

\item {} 
\sphinxAtStartPar
keyword = \sphinxstyleemphasis{single}
\begin{quote}
\begin{sphinxalltt}
\sphinxstyleemphasis{single} value = \sphinxstyleemphasis{yes} or \sphinxstyleemphasis{no} to treat single atoms (no bonds) as fragments
\end{sphinxalltt}
\end{quote}

\end{itemize}


\subsection{Examples}
\label{\detokenize{compute_cluster_atom:examples}}
\begin{sphinxVerbatim}[commandchars=\\\{\}]
\PYG{k}{compute }\PYG{n+nv+nvIdentifier}{1}\PYG{+w}{ }\PYG{n+nv+nvIdentifier}{all}\PYG{+w}{ }\PYG{n}{cluster}\PYG{o}{/}\PYG{n}{atom}\PYG{+w}{ }\PYG{l+m}{3.5}
\PYG{k}{compute }\PYG{n+nv+nvIdentifier}{1}\PYG{+w}{ }\PYG{n+nv+nvIdentifier}{all}\PYG{+w}{ }\PYG{n}{fragment}\PYG{o}{/}\PYG{n}{atom}
\PYG{k}{compute }\PYG{n+nv+nvIdentifier}{1}\PYG{+w}{ }\PYG{n+nv+nvIdentifier}{all}\PYG{+w}{ }\PYG{n}{fragment}\PYG{o}{/}\PYG{n}{atom}\PYG{+w}{ }\PYG{n}{single}\PYG{+w}{ }\PYG{n}{no}
\PYG{k}{compute }\PYG{n+nv+nvIdentifier}{1}\PYG{+w}{ }\PYG{n+nv+nvIdentifier}{all}\PYG{+w}{ }\PYG{n}{aggregate}\PYG{o}{/}\PYG{n}{atom}\PYG{+w}{ }\PYG{l+m}{3.5}
\end{sphinxVerbatim}


\subsection{Description}
\label{\detokenize{compute_cluster_atom:description}}
\sphinxAtStartPar
Define a computation that assigns each atom a cluster, fragment, or
aggregate ID.  Only atoms in the compute group are clustered and
assigned cluster IDs. Atoms not in the compute group are assigned an
ID = 0.

\sphinxAtStartPar
A cluster is defined as a set of atoms, each of which is within the
cutoff distance from one or more other atoms in the cluster.  If an
atom has no neighbors within the cutoff distance, then it is a 1\sphinxhyphen{}atom
cluster.

\sphinxAtStartPar
A fragment is similarly defined as a set of atoms, each of which has a
bond to another atom in the fragment.  Bonds can be defined initially
via the {\hyperref[\detokenize{read_data::doc}]{\sphinxcrossref{\DUrole{doc}{data file}}}} or {\hyperref[\detokenize{create_bonds::doc}]{\sphinxcrossref{\DUrole{doc}{create\_bonds}}}} commands, or dynamically by fixes which create or
break bonds like {\hyperref[\detokenize{fix_bond_react::doc}]{\sphinxcrossref{\DUrole{doc}{fix bond/react}}}}, {\hyperref[\detokenize{fix_bond_create::doc}]{\sphinxcrossref{\DUrole{doc}{fix
bond/create}}}}, {\hyperref[\detokenize{fix_bond_swap::doc}]{\sphinxcrossref{\DUrole{doc}{fix bond/swap}}}},
or {\hyperref[\detokenize{fix_bond_break::doc}]{\sphinxcrossref{\DUrole{doc}{fix bond/break}}}}.  The cluster ID or fragment
ID of every atom in the cluster will be set to the smallest atom ID of
any atom in the cluster or fragment, respectively.

\sphinxAtStartPar
For the \sphinxstyleemphasis{fragment/atom} style, the \sphinxstyleemphasis{single} keyword determines whether
single atoms (not bonded to another atom) are treated as one\sphinxhyphen{}atom
fragments or not, based on the \sphinxstyleemphasis{yes} or \sphinxstyleemphasis{no} setting.  If the setting
is \sphinxstyleemphasis{no} (the default), their fragment IDs are set to 0.

\sphinxAtStartPar
An aggregate is defined by combining the rules for clusters and
fragments (i.e., a set of atoms, where each of them is within the cutoff
distance from one or more atoms within a fragment that is part of
the same cluster). This measure can be used to track molecular assemblies
like micelles.

\sphinxAtStartPar
For computes \sphinxstyleemphasis{cluster/atom} and \sphinxstyleemphasis{aggregate/atom} a neighbor list
needed to compute cluster IDs is constructed each time the compute is
invoked.  Thus it can be inefficient to compute/dump this quantity too
frequently or to have multiple \sphinxstyleemphasis{cluster/atom} or \sphinxstyleemphasis{aggregate/atom}
style computes.

\begin{sphinxadmonition}{note}{Note:}
\sphinxAtStartPar
If you have a bonded system, then the settings of
{\hyperref[\detokenize{special_bonds::doc}]{\sphinxcrossref{\DUrole{doc}{special\_bonds}}}} command can remove pairwise
interactions between atoms in the same bond, angle, or dihedral.  This
is the default setting for the {\hyperref[\detokenize{special_bonds::doc}]{\sphinxcrossref{\DUrole{doc}{special\_bonds}}}}
command, and means those pairwise interactions do not appear in the
neighbor list.  Because this fix uses the neighbor list, it also means
those pairs will not be included when computing the clusters. This
does not apply when using long\sphinxhyphen{}range coulomb (\sphinxstyleemphasis{coul/long}, \sphinxstyleemphasis{coul/msm},
\sphinxstyleemphasis{coul/wolf} or similar.  One way to get around this would be to set
special\_bond scaling factors to very tiny numbers that are not exactly
zero (e.g., \(1.0 \times 10^{-50}\)). Another workaround is to write a
dump file and use the {\hyperref[\detokenize{rerun::doc}]{\sphinxcrossref{\DUrole{doc}{rerun}}}} command to compute the clusters
for snapshots in the dump file.  The rerun script can use a
{\hyperref[\detokenize{special_bonds::doc}]{\sphinxcrossref{\DUrole{doc}{special\_bonds}}}} command that includes all pairs in
the neighbor list.
\end{sphinxadmonition}

\begin{sphinxadmonition}{note}{Note:}
\sphinxAtStartPar
For the compute fragment/atom style, each fragment is identified
using the current bond topology.  This will not account for bonds
broken by the {\hyperref[\detokenize{bond_quartic::doc}]{\sphinxcrossref{\DUrole{doc}{bond\_style quartic}}}} command
because it does not perform a full update of the bond topology data
structures within LAMMPS.
\end{sphinxadmonition}


\subsection{Output info}
\label{\detokenize{compute_cluster_atom:output-info}}
\sphinxAtStartPar
This compute calculates a per\sphinxhyphen{}atom vector, which can be accessed by
any command that uses per\sphinxhyphen{}atom values from a compute as input.  See
the {\hyperref[\detokenize{Howto_output::doc}]{\sphinxcrossref{\DUrole{doc}{Howto output}}}} page for an overview of
LAMMPS output options.

\sphinxAtStartPar
The per\sphinxhyphen{}atom vector values will be an ID \(> 0\), as explained above.


\subsection{Restrictions}
\label{\detokenize{compute_cluster_atom:restrictions}}
\sphinxAtStartPar
none


\subsection{Related commands}
\label{\detokenize{compute_cluster_atom:related-commands}}
\sphinxAtStartPar
{\hyperref[\detokenize{compute_coord_atom::doc}]{\sphinxcrossref{\DUrole{doc}{compute coord/atom}}}}


\subsection{Default}
\label{\detokenize{compute_cluster_atom:default}}
\sphinxAtStartPar
The default for fragment/atom is single=no.

\sphinxstepscope

\index{compute cna/atom@\spxentry{compute cna/atom}}\ignorespaces 

\section{compute cna/atom command}
\label{\detokenize{compute_cna_atom:compute-cna-atom-command}}\label{\detokenize{compute_cna_atom:index-0}}\label{\detokenize{compute_cna_atom::doc}}

\subsection{Syntax}
\label{\detokenize{compute_cna_atom:syntax}}
\begin{sphinxVerbatim}[commandchars=\\\{\}]
\PYG{k}{compute }\PYG{n+nv+nvIdentifier}{ID}\PYG{+w}{ }\PYG{n+nv+nvIdentifier}{group\PYGZhy{}ID}\PYG{+w}{ }\PYG{n}{cna}\PYG{o}{/}\PYG{n}{atom}\PYG{+w}{ }\PYG{n}{cutoff}
\end{sphinxVerbatim}
\begin{itemize}
\item {} 
\sphinxAtStartPar
ID, group\sphinxhyphen{}ID are documented in {\hyperref[\detokenize{compute::doc}]{\sphinxcrossref{\DUrole{doc}{compute}}}} command

\item {} 
\sphinxAtStartPar
cna/atom = style name of this compute command

\item {} 
\sphinxAtStartPar
cutoff = cutoff distance for nearest neighbors (distance units)

\end{itemize}


\subsection{Examples}
\label{\detokenize{compute_cna_atom:examples}}
\begin{sphinxVerbatim}[commandchars=\\\{\}]
\PYG{k}{compute }\PYG{n+nv+nvIdentifier}{1}\PYG{+w}{ }\PYG{n+nv+nvIdentifier}{all}\PYG{+w}{ }\PYG{n}{cna}\PYG{o}{/}\PYG{n}{atom}\PYG{+w}{ }\PYG{l+m}{3.08}
\end{sphinxVerbatim}


\subsection{Description}
\label{\detokenize{compute_cna_atom:description}}
\sphinxAtStartPar
Define a computation that calculates the CNA (Common Neighbor
Analysis) pattern for each atom in the group.  In solid\sphinxhyphen{}state systems
the CNA pattern is a useful measure of the local crystal structure
around an atom.  The CNA methodology is described in {\hyperref[\detokenize{compute_cna_atom:faken}]{\sphinxcrossref{\DUrole{std,std-ref}{(Faken)}}}}
and {\hyperref[\detokenize{compute_cna_atom:tsuzuki1}]{\sphinxcrossref{\DUrole{std,std-ref}{(Tsuzuki)}}}}.

\sphinxAtStartPar
Currently, there are five kinds of CNA patterns LAMMPS recognizes:
\begin{itemize}
\item {} 
\sphinxAtStartPar
fcc = 1

\item {} 
\sphinxAtStartPar
hcp = 2

\item {} 
\sphinxAtStartPar
bcc = 3

\item {} 
\sphinxAtStartPar
icosahedral = 4

\item {} 
\sphinxAtStartPar
unknown = 5

\end{itemize}

\sphinxAtStartPar
The value of the CNA pattern will be 0 for atoms not in the specified
compute group.  Note that normally a CNA calculation should only be
performed on mono\sphinxhyphen{}component systems.

\sphinxAtStartPar
The CNA calculation can be sensitive to the specified cutoff value.
You should ensure the appropriate nearest neighbors of an atom are
found within the cutoff distance for the presumed crystal structure
(e.g., 12 nearest neighbor for perfect FCC and HCP crystals, 14 nearest
neighbors for perfect BCC crystals).  These formulas can be used to
obtain a good cutoff distance:
\begin{equation*}
\begin{split}r_{c}^{\mathrm{fcc}} = & \frac{1}{2} \left(\frac{\sqrt{2}}{2} + 1\right) a
  \approx 0.8536 a \\
r_{c}^{\mathrm{bcc}} = & \frac{1}{2}(\sqrt{2} + 1) a
  \approx 1.207 a \\
r_{c}^{\mathrm{hcp}} = & \frac{1}{2}\left(1+\sqrt{\frac{4+2x^{2}}{3}}\right) a\end{split}
\end{equation*}
\sphinxAtStartPar
where \(a\) is the lattice constant for the crystal structure concerned
and in the HCP case, \(x = (c/a) / 1.633\), where 1.633 is the ideal
\(c/a\) for HCP crystals.

\sphinxAtStartPar
Also note that since the CNA calculation in LAMMPS uses the neighbors
of an owned atom to find the nearest neighbors of a ghost atom, the
following relation should also be satisfied:
\begin{equation*}
\begin{split}r_c + r_s > 2*{\rm cutoff}\end{split}
\end{equation*}
\sphinxAtStartPar
where \(r_c\) is the cutoff distance of the potential, \(r_s\)
is the skin
distance as specified by the {\hyperref[\detokenize{neighbor::doc}]{\sphinxcrossref{\DUrole{doc}{neighbor}}}} command, and
cutoff is the argument used with the compute cna/atom command.  LAMMPS
will issue a warning if this is not the case.

\sphinxAtStartPar
The neighbor list needed to compute this quantity is constructed each
time the calculation is performed (e.g. each time a snapshot of atoms
is dumped).  Thus it can be inefficient to compute/dump this quantity
too frequently or to have multiple compute/dump commands, each with a
\sphinxstyleemphasis{cna/atom} style.


\subsection{Output info}
\label{\detokenize{compute_cna_atom:output-info}}
\sphinxAtStartPar
This compute calculates a per\sphinxhyphen{}atom vector, which can be accessed by
any command that uses per\sphinxhyphen{}atom values from a compute as input.  See
the {\hyperref[\detokenize{Howto_output::doc}]{\sphinxcrossref{\DUrole{doc}{Howto output}}}} page for an overview of
LAMMPS output options.

\sphinxAtStartPar
The per\sphinxhyphen{}atom vector values will be a number from 0 to 5, as explained
above.


\subsection{Restrictions}
\label{\detokenize{compute_cna_atom:restrictions}}\begin{quote}

\sphinxAtStartPar
none
\end{quote}


\subsection{Related commands}
\label{\detokenize{compute_cna_atom:related-commands}}
\sphinxAtStartPar
{\hyperref[\detokenize{compute_centro_atom::doc}]{\sphinxcrossref{\DUrole{doc}{compute centro/atom}}}}


\subsection{Default}
\label{\detokenize{compute_cna_atom:default}}
\sphinxAtStartPar
none


\bigskip\hrule\bigskip

\phantomsection\label{\detokenize{compute_cna_atom:faken}}
\sphinxAtStartPar
\sphinxstylestrong{(Faken)} Faken, Jonsson, Comput Mater Sci, 2, 279 (1994).

\phantomsection\label{\detokenize{compute_cna_atom:tsuzuki1}}
\sphinxAtStartPar
\sphinxstylestrong{(Tsuzuki)} Tsuzuki, Branicio, Rino, Comput Phys Comm, 177, 518 (2007).

\sphinxstepscope

\index{compute cnp/atom@\spxentry{compute cnp/atom}}\ignorespaces 

\section{compute cnp/atom command}
\label{\detokenize{compute_cnp_atom:compute-cnp-atom-command}}\label{\detokenize{compute_cnp_atom:index-0}}\label{\detokenize{compute_cnp_atom::doc}}

\subsection{Syntax}
\label{\detokenize{compute_cnp_atom:syntax}}
\begin{sphinxVerbatim}[commandchars=\\\{\}]
\PYG{k}{compute }\PYG{n+nv+nvIdentifier}{ID}\PYG{+w}{ }\PYG{n+nv+nvIdentifier}{group\PYGZhy{}ID}\PYG{+w}{ }\PYG{n}{cnp}\PYG{o}{/}\PYG{n}{atom}\PYG{+w}{ }\PYG{n}{cutoff}
\end{sphinxVerbatim}
\begin{itemize}
\item {} 
\sphinxAtStartPar
ID, group\sphinxhyphen{}ID are documented in {\hyperref[\detokenize{compute::doc}]{\sphinxcrossref{\DUrole{doc}{compute}}}} command

\item {} 
\sphinxAtStartPar
cnp/atom = style name of this compute command

\item {} 
\sphinxAtStartPar
cutoff = cutoff distance for nearest neighbors (distance units)

\end{itemize}


\subsection{Examples}
\label{\detokenize{compute_cnp_atom:examples}}
\begin{sphinxVerbatim}[commandchars=\\\{\}]
\PYG{k}{compute }\PYG{n+nv+nvIdentifier}{1}\PYG{+w}{ }\PYG{n+nv+nvIdentifier}{all}\PYG{+w}{ }\PYG{n}{cnp}\PYG{o}{/}\PYG{n}{atom}\PYG{+w}{ }\PYG{l+m}{3.08}
\end{sphinxVerbatim}


\subsection{Description}
\label{\detokenize{compute_cnp_atom:description}}
\sphinxAtStartPar
Define a computation that calculates the Common Neighborhood
Parameter (CNP) for each atom in the group.  In solid\sphinxhyphen{}state systems
the CNP is a useful measure of the local crystal structure
around an atom and can be used to characterize whether the
atom is part of a perfect lattice, a local defect (e.g., a dislocation
or stacking fault), or at a surface.

\sphinxAtStartPar
The value of the CNP parameter will be 0.0 for atoms not in the
specified compute group.  Note that normally a CNP calculation should
only be performed on single component systems.

\sphinxAtStartPar
This parameter is computed using the following formula from
{\hyperref[\detokenize{compute_cnp_atom:tsuzuki2}]{\sphinxcrossref{\DUrole{std,std-ref}{(Tsuzuki)}}}}
\begin{equation*}
\begin{split}Q_{i} = \frac{1}{n_i}\sum_{j = 1}^{n_i} \left\lVert \sum_{k = 1}^{n_{ij}}  \vec{R}_{ik} + \vec{R}_{jk} \right\rVert^{2}\end{split}
\end{equation*}
\sphinxAtStartPar
where the index \sphinxstyleemphasis{j} goes over the \(n_i\) nearest neighbors of atom
\sphinxstyleemphasis{i}, and the index \sphinxstyleemphasis{k} goes over the \(n_{ij}\) common nearest neighbors
between atom \sphinxstyleemphasis{i} and atom \sphinxstyleemphasis{j}. \(\vec{R}_{ik}\) and
\(\vec{R}_{jk}\) are the vectors connecting atom \sphinxstyleemphasis{k} to atoms \sphinxstyleemphasis{i}
and \sphinxstyleemphasis{j}.  The quantity in the double sum is computed
for each atom.

\sphinxAtStartPar
The CNP calculation is sensitive to the specified cutoff value.
You should ensure that the appropriate nearest neighbors of an atom are
found within the cutoff distance for the presumed crystal structure.
E.g. 12 nearest neighbor for perfect FCC and HCP crystals, 14 nearest
neighbors for perfect BCC crystals.  These formulas can be used to
obtain a good cutoff distance:
\begin{equation*}
\begin{split}r_{c}^{\mathrm{fcc}} = & \frac{1}{2} \left(\frac{\sqrt{2}}{2} + 1\right) a
  \approx 0.8536 a \\
r_{c}^{\mathrm{bcc}} = & \frac{1}{2}(\sqrt{2} + 1) a
  \approx 1.207 a \\
r_{c}^{\mathrm{hcp}} = & \frac{1}{2}\left(1+\sqrt{\frac{4+2x^{2}}{3}}\right) a\end{split}
\end{equation*}
\sphinxAtStartPar
where \(a\) is the lattice constant for the crystal structure concerned
and in the HCP case, \(x = (c/a) / 1.633\), where 1.633 is the ideal
\(c/a\) for HCP crystals.

\sphinxAtStartPar
Also note that since the CNP calculation in LAMMPS uses the neighbors
of an owned atom to find the nearest neighbors of a ghost atom, the
following relation should also be satisfied:
\begin{equation*}
\begin{split}r_c + r_s > 2*{\rm cutoff}\end{split}
\end{equation*}
\sphinxAtStartPar
where \(r_c\) is the cutoff distance of the potential, \(r_s\) is
the skin
distance as specified by the {\hyperref[\detokenize{neighbor::doc}]{\sphinxcrossref{\DUrole{doc}{neighbor}}}} command, and
cutoff is the argument used with the compute cnp/atom command.  LAMMPS
will issue a warning if this is not the case.

\sphinxAtStartPar
The neighbor list needed to compute this quantity is constructed each
time the calculation is performed (e.g., each time a snapshot of atoms
is dumped).  Thus it can be inefficient to compute/dump this quantity
too frequently or to have multiple compute/dump commands, each with a
\sphinxstyleemphasis{cnp/atom} style.


\subsection{Output info}
\label{\detokenize{compute_cnp_atom:output-info}}
\sphinxAtStartPar
This compute calculates a per\sphinxhyphen{}atom vector, which can be accessed by
any command that uses per\sphinxhyphen{}atom values from a compute as input.  See
the {\hyperref[\detokenize{Howto_output::doc}]{\sphinxcrossref{\DUrole{doc}{Howto output}}}} page for an overview of
LAMMPS output options.

\sphinxAtStartPar
The per\sphinxhyphen{}atom vector values will be real positive numbers. Some typical CNP
values:

\begin{sphinxVerbatim}[commandchars=\\\{\}]
FCC lattice = 0.0
BCC lattice = 0.0
HCP lattice = 4.4

FCC (111) surface = 13.0
FCC (100) surface = 26.5
FCC dislocation core = 11
\end{sphinxVerbatim}


\subsection{Restrictions}
\label{\detokenize{compute_cnp_atom:restrictions}}
\sphinxAtStartPar
This compute is part of the EXTRA\sphinxhyphen{}COMPUTE package.  It is only enabled if
LAMMPS was built with that package.  See the {\hyperref[\detokenize{Build_package::doc}]{\sphinxcrossref{\DUrole{doc}{Build package}}}} page for more info.


\subsection{Related commands}
\label{\detokenize{compute_cnp_atom:related-commands}}
\sphinxAtStartPar
{\hyperref[\detokenize{compute_cna_atom::doc}]{\sphinxcrossref{\DUrole{doc}{compute cna/atom}}}}
{\hyperref[\detokenize{compute_centro_atom::doc}]{\sphinxcrossref{\DUrole{doc}{compute centro/atom}}}}


\subsection{Default}
\label{\detokenize{compute_cnp_atom:default}}
\sphinxAtStartPar
none


\bigskip\hrule\bigskip

\phantomsection\label{\detokenize{compute_cnp_atom:tsuzuki2}}
\sphinxAtStartPar
\sphinxstylestrong{(Tsuzuki)} Tsuzuki, Branicio, Rino, Comput Phys Comm, 177, 518 (2007).

\sphinxstepscope

\index{compute com@\spxentry{compute com}}\ignorespaces 

\section{compute com command}
\label{\detokenize{compute_com:compute-com-command}}\label{\detokenize{compute_com:index-0}}\label{\detokenize{compute_com::doc}}

\subsection{Syntax}
\label{\detokenize{compute_com:syntax}}
\begin{sphinxVerbatim}[commandchars=\\\{\}]
\PYG{k}{compute }\PYG{n+nv+nvIdentifier}{ID}\PYG{+w}{ }\PYG{n+nv+nvIdentifier}{group\PYGZhy{}ID}\PYG{+w}{ }\PYG{n}{com}
\end{sphinxVerbatim}
\begin{itemize}
\item {} 
\sphinxAtStartPar
ID, group\sphinxhyphen{}ID are documented in {\hyperref[\detokenize{compute::doc}]{\sphinxcrossref{\DUrole{doc}{compute}}}} command

\item {} 
\sphinxAtStartPar
com = style name of this compute command

\end{itemize}


\subsection{Examples}
\label{\detokenize{compute_com:examples}}
\begin{sphinxVerbatim}[commandchars=\\\{\}]
\PYG{k}{compute }\PYG{n+nv+nvIdentifier}{1}\PYG{+w}{ }\PYG{n+nv+nvIdentifier}{all}\PYG{+w}{ }\PYG{n}{com}
\end{sphinxVerbatim}


\subsection{Description}
\label{\detokenize{compute_com:description}}
\sphinxAtStartPar
Define a computation that calculates the center\sphinxhyphen{}of\sphinxhyphen{}mass of the group
of atoms, including all effects due to atoms passing through periodic
boundaries.

\sphinxAtStartPar
A vector of three quantities is calculated by this compute, which
are the \((x,y,z)\) coordinates of the center of mass.

\begin{sphinxadmonition}{note}{Note:}
\sphinxAtStartPar
The coordinates of an atom contribute to the center\sphinxhyphen{}of\sphinxhyphen{}mass in
“unwrapped” form, by using the image flags associated with each atom.
See the {\hyperref[\detokenize{dump::doc}]{\sphinxcrossref{\DUrole{doc}{dump custom}}}} command for a discussion of
“unwrapped” coordinates.  See the Atoms section of the
{\hyperref[\detokenize{read_data::doc}]{\sphinxcrossref{\DUrole{doc}{read\_data}}}} command for a discussion of image flags and
how they are set for each atom.  You can reset the image flags
(e.g., to 0) before invoking this compute by using the
{\hyperref[\detokenize{set::doc}]{\sphinxcrossref{\DUrole{doc}{set image}}}} command.
\end{sphinxadmonition}


\subsection{Output info}
\label{\detokenize{compute_com:output-info}}
\sphinxAtStartPar
This compute calculates a global vector of length 3, which can be
accessed by indices 1\textendash{}3 by any command that uses global vector values
from a compute as input.  See the {\hyperref[\detokenize{Howto_output::doc}]{\sphinxcrossref{\DUrole{doc}{Howto output}}}} doc
page for an overview of LAMMPS output options.

\sphinxAtStartPar
The vector values are “intensive”.  The vector values will be in
distance {\hyperref[\detokenize{units::doc}]{\sphinxcrossref{\DUrole{doc}{units}}}}.


\subsection{Restrictions}
\label{\detokenize{compute_com:restrictions}}\begin{quote}

\sphinxAtStartPar
none
\end{quote}


\subsection{Related commands}
\label{\detokenize{compute_com:related-commands}}
\sphinxAtStartPar
{\hyperref[\detokenize{compute_com_chunk::doc}]{\sphinxcrossref{\DUrole{doc}{compute com/chunk}}}}


\subsection{Default}
\label{\detokenize{compute_com:default}}
\sphinxAtStartPar
none

\sphinxstepscope

\index{compute com/chunk@\spxentry{compute com/chunk}}\ignorespaces 

\section{compute com/chunk command}
\label{\detokenize{compute_com_chunk:compute-com-chunk-command}}\label{\detokenize{compute_com_chunk:index-0}}\label{\detokenize{compute_com_chunk::doc}}

\subsection{Syntax}
\label{\detokenize{compute_com_chunk:syntax}}
\begin{sphinxVerbatim}[commandchars=\\\{\}]
\PYG{k}{compute }\PYG{n+nv+nvIdentifier}{ID}\PYG{+w}{ }\PYG{n+nv+nvIdentifier}{group\PYGZhy{}ID}\PYG{+w}{ }\PYG{n}{com}\PYG{o}{/}\PYG{n}{chunk}\PYG{+w}{ }\PYG{n}{chunkID}
\end{sphinxVerbatim}
\begin{itemize}
\item {} 
\sphinxAtStartPar
ID, group\sphinxhyphen{}ID are documented in {\hyperref[\detokenize{compute::doc}]{\sphinxcrossref{\DUrole{doc}{compute}}}} command

\item {} 
\sphinxAtStartPar
com/chunk = style name of this compute command

\item {} 
\sphinxAtStartPar
chunkID = ID of {\hyperref[\detokenize{compute_chunk_atom::doc}]{\sphinxcrossref{\DUrole{doc}{compute chunk/atom}}}} command

\end{itemize}


\subsection{Examples}
\label{\detokenize{compute_com_chunk:examples}}
\begin{sphinxVerbatim}[commandchars=\\\{\}]
\PYG{k}{compute }\PYG{n+nv+nvIdentifier}{1}\PYG{+w}{ }\PYG{n+nv+nvIdentifier}{fluid}\PYG{+w}{ }\PYG{n}{com}\PYG{o}{/}\PYG{n}{chunk}\PYG{+w}{ }\PYG{n}{molchunk}
\end{sphinxVerbatim}


\subsection{Description}
\label{\detokenize{compute_com_chunk:description}}
\sphinxAtStartPar
Define a computation that calculates the center\sphinxhyphen{}of\sphinxhyphen{}mass for multiple
chunks of atoms.

\sphinxAtStartPar
In LAMMPS, chunks are collections of atoms defined by a {\hyperref[\detokenize{compute_chunk_atom::doc}]{\sphinxcrossref{\DUrole{doc}{compute chunk/atom}}}} command, which assigns each atom
to a single chunk (or no chunk).  The ID for this command is specified
as chunkID.  For example, a single chunk could be the atoms in a
molecule or atoms in a spatial bin.  See the {\hyperref[\detokenize{compute_chunk_atom::doc}]{\sphinxcrossref{\DUrole{doc}{compute chunk/atom}}}} and {\hyperref[\detokenize{Howto_chunk::doc}]{\sphinxcrossref{\DUrole{doc}{Howto chunk}}}}
doc pages for details of how chunks can be defined and examples of how
they can be used to measure properties of a system.

\sphinxAtStartPar
This compute calculates the \((x,y,z)\) coordinates of the center of mass
for each chunk, which includes all effects due to atoms passing through
periodic boundaries.

\sphinxAtStartPar
Note that only atoms in the specified group contribute to the
calculation.  The {\hyperref[\detokenize{compute_chunk_atom::doc}]{\sphinxcrossref{\DUrole{doc}{compute chunk/atom}}}} command
defines its own group; atoms will have a chunk ID = 0 if they are not
in that group, signifying they are not assigned to a chunk, and will
thus also not contribute to this calculation.  You can specify the
“all” group for this command if you simply want to include atoms with
non\sphinxhyphen{}zero chunk IDs.

\begin{sphinxadmonition}{note}{Note:}
\sphinxAtStartPar
The coordinates of an atom contribute to the chunk’s
center\sphinxhyphen{}of\sphinxhyphen{}mass in “unwrapped” form, by using the image flags
associated with each atom.  See the {\hyperref[\detokenize{dump::doc}]{\sphinxcrossref{\DUrole{doc}{dump custom}}}} command
for a discussion of “unwrapped” coordinates.  See the Atoms section of
the {\hyperref[\detokenize{read_data::doc}]{\sphinxcrossref{\DUrole{doc}{read\_data}}}} command for a discussion of image flags
and how they are set for each atom.  You can reset the image flags
(e.g., to 0) before invoking this compute by using the
{\hyperref[\detokenize{set::doc}]{\sphinxcrossref{\DUrole{doc}{set image}}}} command.
\end{sphinxadmonition}

\sphinxAtStartPar
The simplest way to output the results of the compute com/chunk
calculation to a file is to use the {\hyperref[\detokenize{fix_ave_time::doc}]{\sphinxcrossref{\DUrole{doc}{fix ave/time}}}}
command, for example:

\begin{sphinxVerbatim}[commandchars=\\\{\}]
\PYG{k}{compute }\PYG{n+nv+nvIdentifier}{cc1}\PYG{+w}{ }\PYG{n+nv+nvIdentifier}{all}\PYG{+w}{ }\PYG{n}{chunk}\PYG{o}{/}\PYG{n}{atom}\PYG{+w}{ }\PYG{n}{molecule}
\PYG{k}{compute }\PYG{n+nv+nvIdentifier}{myChunk}\PYG{+w}{ }\PYG{n+nv+nvIdentifier}{all}\PYG{+w}{ }\PYG{n}{com}\PYG{o}{/}\PYG{n}{chunk}\PYG{+w}{ }\PYG{n}{cc1}
\PYG{k}{fix }\PYG{n+nv+nvIdentifier}{1}\PYG{+w}{ }\PYG{n+nv+nvIdentifier}{all}\PYG{+w}{ }\PYG{n}{ave}\PYG{o}{/}\PYG{n}{time}\PYG{+w}{ }\PYG{l+m}{100}\PYG{+w}{ }\PYG{l+m}{1}\PYG{+w}{ }\PYG{l+m}{100}\PYG{+w}{ }\PYG{n}{c\PYGZus{}myChunk[}\PYG{o}{*}\PYG{n}{]}\PYG{+w}{ }\PYG{n}{file}\PYG{+w}{ }\PYG{n}{tmp.out}\PYG{+w}{ }\PYG{n}{mode}\PYG{+w}{ }\PYG{n}{vector}
\end{sphinxVerbatim}


\subsection{Output info}
\label{\detokenize{compute_com_chunk:output-info}}
\sphinxAtStartPar
This compute calculates a global array where the number of rows = the
number of chunks \sphinxstyleemphasis{Nchunk} as calculated by the specified {\hyperref[\detokenize{compute_chunk_atom::doc}]{\sphinxcrossref{\DUrole{doc}{compute chunk/atom}}}} command.  The number of columns is
3 for the \((x,y,z)\) center\sphinxhyphen{}of\sphinxhyphen{}mass coordinates of each chunk.  These
values can be accessed by any command that uses global array values
from a compute as input.  See the {\hyperref[\detokenize{Howto_output::doc}]{\sphinxcrossref{\DUrole{doc}{Howto output}}}} doc
page for an overview of LAMMPS output options.

\sphinxAtStartPar
The array values are “intensive”.  The array values will be in
distance {\hyperref[\detokenize{units::doc}]{\sphinxcrossref{\DUrole{doc}{units}}}}.


\subsection{Restrictions}
\label{\detokenize{compute_com_chunk:restrictions}}\begin{quote}

\sphinxAtStartPar
none
\end{quote}


\subsection{Related commands}
\label{\detokenize{compute_com_chunk:related-commands}}
\sphinxAtStartPar
{\hyperref[\detokenize{compute_com::doc}]{\sphinxcrossref{\DUrole{doc}{compute com}}}}


\subsection{Default}
\label{\detokenize{compute_com_chunk:default}}
\sphinxAtStartPar
none

\sphinxstepscope

\index{compute composition/atom@\spxentry{compute composition/atom}}\ignorespaces 
\index{compute composition/atom/kk@\spxentry{compute composition/atom/kk}}\ignorespaces 

\section{compute composition/atom command}
\label{\detokenize{compute_composition_atom:compute-composition-atom-command}}\label{\detokenize{compute_composition_atom:index-1}}\label{\detokenize{compute_composition_atom:index-0}}\label{\detokenize{compute_composition_atom::doc}}
\sphinxAtStartPar
Accelerator Variants: \sphinxstyleemphasis{composition/atom/kk}


\subsection{Syntax}
\label{\detokenize{compute_composition_atom:syntax}}
\begin{sphinxVerbatim}[commandchars=\\\{\}]
\PYG{k}{compute }\PYG{n+nv+nvIdentifier}{ID}\PYG{+w}{ }\PYG{n+nv+nvIdentifier}{group\PYGZhy{}ID}\PYG{+w}{ }\PYG{n}{composition}\PYG{o}{/}\PYG{n}{atom}\PYG{+w}{ }\PYG{n}{keyword}\PYG{+w}{ }\PYG{n}{values}\PYG{+w}{ }\PYG{n}{...}
\end{sphinxVerbatim}
\begin{itemize}
\item {} 
\sphinxAtStartPar
ID, group\sphinxhyphen{}ID are documented in {\hyperref[\detokenize{compute::doc}]{\sphinxcrossref{\DUrole{doc}{compute}}}} command

\item {} 
\sphinxAtStartPar
composition/atom = style name of this compute command

\item {} 
\sphinxAtStartPar
one or more keyword/value pairs may be appended
\begin{sphinxalltt}
keyword = \sphinxstyleemphasis{cutoff}
  \sphinxstyleemphasis{cutoff} value = distance cutoff
\end{sphinxalltt}

\end{itemize}


\subsection{Examples}
\label{\detokenize{compute_composition_atom:examples}}
\begin{sphinxVerbatim}[commandchars=\\\{\}]
\PYG{k}{compute }\PYG{n+nv+nvIdentifier}{1}\PYG{+w}{ }\PYG{n+nv+nvIdentifier}{all}\PYG{+w}{ }\PYG{n}{composition}\PYG{o}{/}\PYG{n}{atom}

\PYG{k}{compute }\PYG{n+nv+nvIdentifier}{1}\PYG{+w}{ }\PYG{n+nv+nvIdentifier}{all}\PYG{+w}{ }\PYG{n}{composition}\PYG{o}{/}\PYG{n}{atom}\PYG{+w}{ }\PYG{n}{cutoff}\PYG{+w}{ }\PYG{l+m}{9.0}
\PYG{k}{comm\PYGZus{}modify}\PYG{+w}{ }\PYG{n}{cutoff}\PYG{+w}{ }\PYG{l+m}{9.0}
\end{sphinxVerbatim}


\subsection{Description}
\label{\detokenize{compute_composition_atom:description}}
\sphinxAtStartPar
\DUrole{versionmodified,added}{Added in version 21Nov2023.}

\sphinxAtStartPar
Define a computation that calculates a local composition vector for each
atom. For a central atom with \(M\) neighbors within the neighbor cutoff sphere,
composition is defined as the number of atoms of a given type
(including the central atom) divided by (\(M+1\)).  For a given central atom,
the sum of all compositions equals one.

\begin{sphinxadmonition}{note}{Note:}
\sphinxAtStartPar
This compute uses the number of atom types, not chemical species, assigned in
{\hyperref[\detokenize{pair_coeff::doc}]{\sphinxcrossref{\DUrole{doc}{pair\_coeff}}}} command.  If an interatomic potential has two
species (i.e., Cu and Ni) assigned to four different atom types in
{\hyperref[\detokenize{pair_coeff::doc}]{\sphinxcrossref{\DUrole{doc}{pair\_coeff}}}} (i.e., ‘Cu Cu Ni Ni’), the compute will
output four fractional values.  In those cases, the user may desire an extra
calculation step to consolidate per\sphinxhyphen{}type fractions into per\sphinxhyphen{}species fractions.
This calculation can be conducted within LAMMPS using another compute such as
{\hyperref[\detokenize{compute_reduce::doc}]{\sphinxcrossref{\DUrole{doc}{compute reduce}}}}, an atom\sphinxhyphen{}style {\hyperref[\detokenize{variable::doc}]{\sphinxcrossref{\DUrole{doc}{variable command}}}}, or as a
post\sphinxhyphen{}processing step.
\end{sphinxadmonition}


\bigskip\hrule\bigskip


\sphinxAtStartPar
The optional keyword \sphinxstyleemphasis{cutoff} defines the distance cutoff used when
searching for neighbors. The default value is the cutoff specified by
the pair style. If no pair style is defined, then a cutoff must be
defined using this keyword. If the specified cutoff is larger than
that of the pair\_style plus neighbor skin (or no pair style is
defined), the \sphinxstyleemphasis{comm\_modify cutoff} option must also be set to match
that of the \sphinxstyleemphasis{cutoff} keyword.

\sphinxAtStartPar
The neighbor list needed to compute this quantity is constructed each
time the calculation is performed (i.e. each time a snapshot of atoms
is dumped).  Thus it can be inefficient to compute/dump this quantity
too frequently.

\begin{sphinxadmonition}{note}{Note:}
\sphinxAtStartPar
If you have a bonded system, then the settings of
{\hyperref[\detokenize{special_bonds::doc}]{\sphinxcrossref{\DUrole{doc}{special\_bonds}}}} command can remove pairwise
interactions between atoms in the same bond, angle, or dihedral.
This is the default setting for the {\hyperref[\detokenize{special_bonds::doc}]{\sphinxcrossref{\DUrole{doc}{special\_bonds}}}} command, and means those pairwise interactions do
not appear in the neighbor list.  Because this compute uses the
neighbor list, it also means those pairs will not be included in
the order parameter.  This difficulty can be circumvented by
writing a dump file, and using the {\hyperref[\detokenize{rerun::doc}]{\sphinxcrossref{\DUrole{doc}{rerun}}}} command to
compute the order parameter for snapshots in the dump file.  The
rerun script can use a {\hyperref[\detokenize{special_bonds::doc}]{\sphinxcrossref{\DUrole{doc}{special\_bonds}}}} command
that includes all pairs in the neighbor list.
\end{sphinxadmonition}


\bigskip\hrule\bigskip



\subsection{Output info}
\label{\detokenize{compute_composition_atom:output-info}}
\sphinxAtStartPar
This compute calculates a per\sphinxhyphen{}atom array with \(1 + N\) columns, where \(N\)
is the number of atom types. The first column is a count of the number of atoms
used to calculate composition (including the central atom), and each subsequent
column indicates the fraction of that atom type within the cutoff sphere.

\sphinxAtStartPar
These values can be accessed by any command that uses per\sphinxhyphen{}atom values
from a compute as input.  See the {\hyperref[\detokenize{Howto_output::doc}]{\sphinxcrossref{\DUrole{doc}{Howto output}}}}
doc page for an overview of LAMMPS output options.


\subsection{Restrictions}
\label{\detokenize{compute_composition_atom:restrictions}}
\sphinxAtStartPar
This compute is part of the EXTRA\sphinxhyphen{}COMPUTE package.  It is only enabled
if LAMMPS was built with that package.  See the {\hyperref[\detokenize{Build_package::doc}]{\sphinxcrossref{\DUrole{doc}{Build package}}}} page for more info.

\sphinxAtStartPar
This compute requires {\hyperref[\detokenize{neighbor::doc}]{\sphinxcrossref{\DUrole{doc}{neighbor styles ‘bin’ or ‘nsq’}}}}.


\subsection{Related commands}
\label{\detokenize{compute_composition_atom:related-commands}}
\sphinxAtStartPar
{\hyperref[\detokenize{comm_modify::doc}]{\sphinxcrossref{\DUrole{doc}{comm\_modify}}}}


\subsection{Default}
\label{\detokenize{compute_composition_atom:default}}
\sphinxAtStartPar
The option defaults are \sphinxstyleemphasis{cutoff} = pair style cutoff.

\sphinxstepscope

\index{compute contact/atom@\spxentry{compute contact/atom}}\ignorespaces 

\section{compute contact/atom command}
\label{\detokenize{compute_contact_atom:compute-contact-atom-command}}\label{\detokenize{compute_contact_atom:index-0}}\label{\detokenize{compute_contact_atom::doc}}

\subsection{Syntax}
\label{\detokenize{compute_contact_atom:syntax}}
\begin{sphinxVerbatim}[commandchars=\\\{\}]
\PYG{k}{compute }\PYG{n+nv+nvIdentifier}{ID}\PYG{+w}{ }\PYG{n+nv+nvIdentifier}{group\PYGZhy{}ID}\PYG{+w}{ }\PYG{n}{contact}\PYG{o}{/}\PYG{n}{atom}\PYG{+w}{ }\PYG{n}{group2}\PYG{o}{\PYGZhy{}}\PYG{n}{ID}
\end{sphinxVerbatim}
\begin{itemize}
\item {} 
\sphinxAtStartPar
ID, group\sphinxhyphen{}ID are documented in {\hyperref[\detokenize{compute::doc}]{\sphinxcrossref{\DUrole{doc}{compute}}}} command

\item {} 
\sphinxAtStartPar
contact/atom = style name of this compute command

\item {} 
\sphinxAtStartPar
group2\sphinxhyphen{}ID = optional argument to restrict which atoms to consider for contacts (see below)

\end{itemize}


\subsection{Examples}
\label{\detokenize{compute_contact_atom:examples}}
\begin{sphinxVerbatim}[commandchars=\\\{\}]
\PYG{k}{compute }\PYG{n+nv+nvIdentifier}{1}\PYG{+w}{ }\PYG{n+nv+nvIdentifier}{all}\PYG{+w}{ }\PYG{n}{contact}\PYG{o}{/}\PYG{n}{atom}
\PYG{k}{compute }\PYG{n+nv+nvIdentifier}{1}\PYG{+w}{ }\PYG{n+nv+nvIdentifier}{all}\PYG{+w}{ }\PYG{n}{contact}\PYG{o}{/}\PYG{n}{atom}\PYG{+w}{ }\PYG{n}{mygroup}
\end{sphinxVerbatim}


\subsection{Description}
\label{\detokenize{compute_contact_atom:description}}
\sphinxAtStartPar
Define a computation that calculates the number of contacts
for each atom in a group.

\sphinxAtStartPar
The contact number is defined for finite\sphinxhyphen{}size spherical particles as
the number of neighbor atoms which overlap the central particle,
meaning that their distance of separation is less than or equal to the
sum of the radii of the two particles.

\sphinxAtStartPar
The value of the contact number will be 0.0 for atoms not in the
specified compute group.

\sphinxAtStartPar
The optional \sphinxstyleemphasis{group2\sphinxhyphen{}ID} argument allows to specify from which group atoms
contribute to the coordination number. Default setting is group ‘all’.


\subsection{Output info}
\label{\detokenize{compute_contact_atom:output-info}}
\sphinxAtStartPar
This compute calculates a per\sphinxhyphen{}atom vector, whose values can be
accessed by any command that uses per\sphinxhyphen{}atom values from a compute as
input.  See the {\hyperref[\detokenize{Howto_output::doc}]{\sphinxcrossref{\DUrole{doc}{Howto output}}}} page for an
overview of LAMMPS output options.

\sphinxAtStartPar
The per\sphinxhyphen{}atom vector values will be a number \(\ge 0.0\), as explained
above.


\subsection{Restrictions}
\label{\detokenize{compute_contact_atom:restrictions}}
\sphinxAtStartPar
This compute is part of the GRANULAR package.  It is only enabled if
LAMMPS was built with that package.  See the
{\hyperref[\detokenize{Build_package::doc}]{\sphinxcrossref{\DUrole{doc}{Build package}}}} page for more info.

\sphinxAtStartPar
This compute requires that atoms store a radius as defined by the
{\hyperref[\detokenize{atom_style::doc}]{\sphinxcrossref{\DUrole{doc}{atom\_style sphere}}}} command.


\subsection{Related commands}
\label{\detokenize{compute_contact_atom:related-commands}}
\sphinxAtStartPar
{\hyperref[\detokenize{compute_coord_atom::doc}]{\sphinxcrossref{\DUrole{doc}{compute coord/atom}}}}


\subsection{Default}
\label{\detokenize{compute_contact_atom:default}}
\sphinxAtStartPar
\sphinxstyleemphasis{group2\sphinxhyphen{}ID} = all

\sphinxstepscope

\index{compute coord/atom@\spxentry{compute coord/atom}}\ignorespaces 
\index{compute coord/atom/kk@\spxentry{compute coord/atom/kk}}\ignorespaces 

\section{compute coord/atom command}
\label{\detokenize{compute_coord_atom:compute-coord-atom-command}}\label{\detokenize{compute_coord_atom:index-1}}\label{\detokenize{compute_coord_atom:index-0}}\label{\detokenize{compute_coord_atom::doc}}
\sphinxAtStartPar
Accelerator Variants: \sphinxstyleemphasis{coord/atom/kk}


\subsection{Syntax}
\label{\detokenize{compute_coord_atom:syntax}}
\begin{sphinxVerbatim}[commandchars=\\\{\}]
\PYG{k}{compute }\PYG{n+nv+nvIdentifier}{ID}\PYG{+w}{ }\PYG{n+nv+nvIdentifier}{group\PYGZhy{}ID}\PYG{+w}{ }\PYG{n}{coord}\PYG{o}{/}\PYG{n}{atom}\PYG{+w}{ }\PYG{n}{style}\PYG{+w}{ }\PYG{n}{args}\PYG{+w}{ }\PYG{n}{...}
\end{sphinxVerbatim}
\begin{itemize}
\item {} 
\sphinxAtStartPar
ID, group\sphinxhyphen{}ID are documented in {\hyperref[\detokenize{compute::doc}]{\sphinxcrossref{\DUrole{doc}{compute}}}} command

\item {} 
\sphinxAtStartPar
coord/atom = style name of this compute command

\item {} 
\sphinxAtStartPar
style = \sphinxstyleemphasis{cutoff} or \sphinxstyleemphasis{orientorder}
\begin{sphinxalltt}
\sphinxstyleemphasis{cutoff} args = cutoff {[}\sphinxstyleemphasis{group} group2\sphinxhyphen{}ID{]} typeN
  cutoff = distance within which to count coordination neighbors (distance units)
  \sphinxstyleemphasis{group} group2\sphinxhyphen{}ID = select group\sphinxhyphen{}ID to restrict which atoms to consider for coordination number (optional)
  typeN = atom type for Nth coordination count (see asterisk form below)
\sphinxstyleemphasis{orientorder} args = orientorderID threshold
  orientorderID = ID of an orientorder/atom compute
  threshold = minimum value of the product of two "connected" atoms
\end{sphinxalltt}

\end{itemize}


\subsection{Examples}
\label{\detokenize{compute_coord_atom:examples}}
\begin{sphinxVerbatim}[commandchars=\\\{\}]
\PYG{k}{compute }\PYG{n+nv+nvIdentifier}{1}\PYG{+w}{ }\PYG{n+nv+nvIdentifier}{all}\PYG{+w}{ }\PYG{n}{coord}\PYG{o}{/}\PYG{n}{atom}\PYG{+w}{ }\PYG{n}{cutoff}\PYG{+w}{ }\PYG{l+m}{2.0}
\PYG{k}{compute }\PYG{n+nv+nvIdentifier}{1}\PYG{+w}{ }\PYG{n+nv+nvIdentifier}{all}\PYG{+w}{ }\PYG{n}{coord}\PYG{o}{/}\PYG{n}{atom}\PYG{+w}{ }\PYG{n}{cutoff}\PYG{+w}{ }\PYG{l+m}{6.0}\PYG{+w}{ }\PYG{l+m}{1}\PYG{+w}{ }\PYG{l+m}{2}
\PYG{k}{compute }\PYG{n+nv+nvIdentifier}{1}\PYG{+w}{ }\PYG{n+nv+nvIdentifier}{all}\PYG{+w}{ }\PYG{n}{coord}\PYG{o}{/}\PYG{n}{atom}\PYG{+w}{ }\PYG{n}{cutoff}\PYG{+w}{ }\PYG{l+m}{6.0}\PYG{+w}{ }\PYG{l+m}{2}\PYG{o}{*}\PYG{l+m}{4}\PYG{+w}{ }\PYG{l+m}{5}\PYG{o}{*}\PYG{l+m}{8}\PYG{+w}{ }\PYG{o}{*}
\PYG{k}{compute }\PYG{n+nv+nvIdentifier}{1}\PYG{+w}{ }\PYG{n+nv+nvIdentifier}{solute}\PYG{+w}{ }\PYG{n}{coord}\PYG{o}{/}\PYG{n}{atom}\PYG{+w}{ }\PYG{n}{cutoff}\PYG{+w}{ }\PYG{l+m}{2.0}\PYG{+w}{ }\PYG{k}{group }\PYG{n+nv+nvIdentifier}{solvent}
\PYG{k}{compute }\PYG{n+nv+nvIdentifier}{1}\PYG{+w}{ }\PYG{n+nv+nvIdentifier}{all}\PYG{+w}{ }\PYG{n}{coord}\PYG{o}{/}\PYG{n}{atom}\PYG{+w}{ }\PYG{n}{orientorder}\PYG{+w}{ }\PYG{l+m}{2}\PYG{+w}{ }\PYG{l+m}{0.5}
\end{sphinxVerbatim}


\subsection{Description}
\label{\detokenize{compute_coord_atom:description}}
\sphinxAtStartPar
This compute performs calculations between neighboring atoms to
determine a coordination value.  The specific calculation and the
meaning of the resulting value depend on the \sphinxstyleemphasis{cstyle} keyword used.

\sphinxAtStartPar
The \sphinxstyleemphasis{cutoff} cstyle calculates one or more traditional coordination
numbers for each atom.  A coordination number is defined as the number
of neighbor atoms with specified atom type(s), and optionally within
the specified group, that are within the specified cutoff distance from
the central atom. The compute group selects only the central atoms; all
neighboring atoms, unless selected by type, type range, or group option,
are included in the coordination number tally.

\sphinxAtStartPar
The optional \sphinxstyleemphasis{group} keyword allows to specify from which group atoms
contribute to the coordination number. Default setting is group ‘all.’

\sphinxAtStartPar
The \sphinxstyleemphasis{typeN} keywords allow specification of which atom types
contribute to each coordination number.  One coordination number is
computed for each of the \sphinxstyleemphasis{typeN} keywords listed.  If no \sphinxstyleemphasis{typeN}
keywords are listed, a single coordination number is calculated, which
includes atoms of all types (same as the “*” format, see below).

\sphinxAtStartPar
The \sphinxstyleemphasis{typeN} keywords can be specified in one of two ways.  An explicit
numeric value can be used, as in the second example above.  Or a
wild\sphinxhyphen{}card asterisk can be used to specify a range of atom types.  This
takes the form “*” or “*n” or “m*” or “m*n”.  If \(N\) is the number of
atom types, then an asterisk with no numeric values means all types
from 1 to \(N\).  A leading asterisk means all types from 1 to n
(inclusive).  A trailing asterisk means all types from m to \(N\)
(inclusive).  A middle asterisk means all types from m to n
(inclusive).

\sphinxAtStartPar
The \sphinxstyleemphasis{orientorder} cstyle calculates the number of “connected” neighbor
atoms \sphinxstyleemphasis{j} around each central atom \sphinxstyleemphasis{i}.  For this \sphinxstyleemphasis{cstyle}, connected is
defined by the orientational order parameter calculated by the
{\hyperref[\detokenize{compute_orientorder_atom::doc}]{\sphinxcrossref{\DUrole{doc}{compute orientorder/atom}}}} command.
This \sphinxstyleemphasis{cstyle} thus allows one to apply the ten Wolde’s criterion to
identify crystal\sphinxhyphen{}like atoms in a system, as discussed in {\hyperref[\detokenize{compute_coord_atom:tenwolde1}]{\sphinxcrossref{\DUrole{std,std-ref}{ten Wolde}}}}.

\sphinxAtStartPar
The ID of the previously specified {\hyperref[\detokenize{compute_orientorder_atom::doc}]{\sphinxcrossref{\DUrole{doc}{compute orientorder/atom}}}} command is specified as
\sphinxstyleemphasis{orientorderID}.  The compute must invoke its \sphinxstyleemphasis{components} option to
calculate components of the \sphinxstyleemphasis{Ybar\_lm} vector for each atoms, as
described in its documentation.  Note that orientorder/atom compute
defines its own criteria for identifying neighboring atoms.  If the
scalar product (\sphinxstyleemphasis{Ybar\_lm(i)}, \sphinxstyleemphasis{Ybar\_lm(j)}), calculated by the
orientorder/atom compute is larger than the specified \sphinxstyleemphasis{threshold},
then \sphinxstyleemphasis{i} and \sphinxstyleemphasis{j} are connected, and the coordination value of \sphinxstyleemphasis{i} is
incremented by one.

\sphinxAtStartPar
For all \sphinxstyleemphasis{cstyle} settings, all coordination values will be 0.0 for
atoms not in the specified compute group.

\sphinxAtStartPar
The neighbor list needed to compute this quantity is constructed each
time the calculation is performed (i.e., each time a snapshot of atoms
is dumped).  Thus it can be inefficient to compute/dump this quantity
too frequently.

\begin{sphinxadmonition}{note}{Note:}
\sphinxAtStartPar
If you have a bonded system, then the settings of
{\hyperref[\detokenize{special_bonds::doc}]{\sphinxcrossref{\DUrole{doc}{special\_bonds}}}} command can remove pairwise
interactions between atoms in the same bond, angle, or dihedral.  This
is the default setting for the {\hyperref[\detokenize{special_bonds::doc}]{\sphinxcrossref{\DUrole{doc}{special\_bonds}}}}
command, and means those pairwise interactions do not appear in the
neighbor list.  Because this fix uses the neighbor list, it also means
those pairs will not be included in the coordination count.  One way
to get around this, is to write a dump file, and use the
{\hyperref[\detokenize{rerun::doc}]{\sphinxcrossref{\DUrole{doc}{rerun}}}} command to compute the coordination for snapshots
in the dump file.  The rerun script can use a
{\hyperref[\detokenize{special_bonds::doc}]{\sphinxcrossref{\DUrole{doc}{special\_bonds}}}} command that includes all pairs in
the neighbor list.
\end{sphinxadmonition}


\bigskip\hrule\bigskip


\sphinxAtStartPar
Styles with a \sphinxstyleemphasis{gpu}, \sphinxstyleemphasis{intel}, \sphinxstyleemphasis{kk}, \sphinxstyleemphasis{omp}, or \sphinxstyleemphasis{opt} suffix are
functionally the same as the corresponding style without the suffix.
They have been optimized to run faster, depending on your available
hardware, as discussed on the {\hyperref[\detokenize{Speed_packages::doc}]{\sphinxcrossref{\DUrole{doc}{Accelerator packages}}}}
page.  The accelerated styles take the same arguments and should
produce the same results, except for round\sphinxhyphen{}off and precision issues.

\sphinxAtStartPar
These accelerated styles are part of the GPU, INTEL, KOKKOS,
OPENMP, and OPT packages, respectively.  They are only enabled if
LAMMPS was built with those packages.  See the {\hyperref[\detokenize{Build_package::doc}]{\sphinxcrossref{\DUrole{doc}{Build package}}}} page for more info.

\sphinxAtStartPar
You can specify the accelerated styles explicitly in your input script
by including their suffix, or you can use the {\hyperref[\detokenize{Run_options::doc}]{\sphinxcrossref{\DUrole{doc}{\sphinxhyphen{}suffix command\sphinxhyphen{}line switch}}}} when you invoke LAMMPS, or you can use the
{\hyperref[\detokenize{suffix::doc}]{\sphinxcrossref{\DUrole{doc}{suffix}}}} command in your input script.

\sphinxAtStartPar
See the {\hyperref[\detokenize{Speed_packages::doc}]{\sphinxcrossref{\DUrole{doc}{Accelerator packages}}}} page for more
instructions on how to use the accelerated styles effectively.


\bigskip\hrule\bigskip



\subsection{Output info}
\label{\detokenize{compute_coord_atom:output-info}}
\sphinxAtStartPar
For \sphinxstyleemphasis{cstyle} cutoff, this compute can calculate a per\sphinxhyphen{}atom vector or
array.  If single \sphinxstyleemphasis{type1} keyword is specified (or if none are
specified), this compute calculates a per\sphinxhyphen{}atom vector.  If multiple
\sphinxstyleemphasis{typeN} keywords are specified, this compute calculates a per\sphinxhyphen{}atom
array, with \(N\) columns.

\sphinxAtStartPar
For \sphinxstyleemphasis{cstyle} orientorder, this compute calculates a per\sphinxhyphen{}atom vector.

\sphinxAtStartPar
These values can be accessed by any command that uses per\sphinxhyphen{}atom values
from a compute as input.  See the {\hyperref[\detokenize{Howto_output::doc}]{\sphinxcrossref{\DUrole{doc}{Howto output}}}} doc
page for an overview of LAMMPS output options.

\sphinxAtStartPar
The per\sphinxhyphen{}atom vector or array values will be a number \(\ge 0.0\), as
explained above.


\subsection{Restrictions}
\label{\detokenize{compute_coord_atom:restrictions}}\begin{quote}

\sphinxAtStartPar
none
\end{quote}


\subsection{Related commands}
\label{\detokenize{compute_coord_atom:related-commands}}
\sphinxAtStartPar
{\hyperref[\detokenize{compute_cluster_atom::doc}]{\sphinxcrossref{\DUrole{doc}{compute cluster/atom}}}}
{\hyperref[\detokenize{compute_orientorder_atom::doc}]{\sphinxcrossref{\DUrole{doc}{compute orientorder/atom}}}}


\subsection{Default}
\label{\detokenize{compute_coord_atom:default}}
\sphinxAtStartPar
group = all


\bigskip\hrule\bigskip

\phantomsection\label{\detokenize{compute_coord_atom:tenwolde1}}
\sphinxAtStartPar
\sphinxstylestrong{(tenWolde)} P. R. ten Wolde, M. J. Ruiz\sphinxhyphen{}Montero, D. Frenkel,
J. Chem. Phys. 104, 9932 (1996).

\sphinxstepscope

\index{compute count/type@\spxentry{compute count/type}}\ignorespaces 

\section{compute count/type command}
\label{\detokenize{compute_count_type:compute-count-type-command}}\label{\detokenize{compute_count_type:index-0}}\label{\detokenize{compute_count_type::doc}}

\subsection{Syntax}
\label{\detokenize{compute_count_type:syntax}}
\begin{sphinxVerbatim}[commandchars=\\\{\}]
\PYG{k}{compute }\PYG{n+nv+nvIdentifier}{ID}\PYG{+w}{ }\PYG{n+nv+nvIdentifier}{group\PYGZhy{}ID}\PYG{+w}{ }\PYG{n}{count}\PYG{o}{/}\PYG{n}{type}\PYG{+w}{ }\PYG{n}{mode}
\end{sphinxVerbatim}
\begin{itemize}
\item {} 
\sphinxAtStartPar
ID, group\sphinxhyphen{}ID are documented in {\hyperref[\detokenize{compute::doc}]{\sphinxcrossref{\DUrole{doc}{compute}}}} command

\item {} 
\sphinxAtStartPar
count/type = style name of this compute command

\item {} 
\sphinxAtStartPar
mode = \sphinxstyleemphasis{atom} or \sphinxstyleemphasis{bond} or \sphinxstyleemphasis{angle} or \sphinxstyleemphasis{dihedral} or \sphinxstyleemphasis{improper}

\end{itemize}


\subsection{Examples}
\label{\detokenize{compute_count_type:examples}}
\begin{sphinxVerbatim}[commandchars=\\\{\}]
\PYG{k}{compute }\PYG{n+nv+nvIdentifier}{1}\PYG{+w}{ }\PYG{n+nv+nvIdentifier}{all}\PYG{+w}{ }\PYG{n}{count}\PYG{o}{/}\PYG{n}{type}\PYG{+w}{ }\PYG{n}{atom}
\PYG{k}{compute }\PYG{n+nv+nvIdentifier}{1}\PYG{+w}{ }\PYG{n+nv+nvIdentifier}{flowmols}\PYG{+w}{ }\PYG{n}{count}\PYG{o}{/}\PYG{n}{type}\PYG{+w}{ }\PYG{n}{bond}
\end{sphinxVerbatim}


\subsection{Description}
\label{\detokenize{compute_count_type:description}}
\sphinxAtStartPar
\DUrole{versionmodified,added}{Added in version 15Jun2023.}

\sphinxAtStartPar
Define a computation that counts the current number of atoms for each
atom type.  Or the number of bonds (angles, dihedrals, impropers) for
each bond (angle, dihedral, improper) type.

\sphinxAtStartPar
The former can be useful if atoms are added to or deleted from the
system in random ways, e.g. via the {\hyperref[\detokenize{fix_deposit::doc}]{\sphinxcrossref{\DUrole{doc}{fix deposit}}}},
{\hyperref[\detokenize{fix_pour::doc}]{\sphinxcrossref{\DUrole{doc}{fix pour}}}}, or {\hyperref[\detokenize{fix_evaporate::doc}]{\sphinxcrossref{\DUrole{doc}{fix evaporate}}}}
commands.  The latter can be useful in reactive simulations where
molecular bonds are broken or created, as well as angles, dihedrals,
impropers.

\sphinxAtStartPar
Note that for this command, bonds (angles, etc) are the topological
kind enumerated in a data file, initially read by the {\hyperref[\detokenize{read_data::doc}]{\sphinxcrossref{\DUrole{doc}{read\_data}}}} command or defined by the {\hyperref[\detokenize{molecule::doc}]{\sphinxcrossref{\DUrole{doc}{molecule}}}}
command.  They do not refer to implicit bonds defined on\sphinxhyphen{}the\sphinxhyphen{}fly by
bond\sphinxhyphen{}order or reactive pair styles based on the current conformation
of small clusters of atoms.

\sphinxAtStartPar
These commands can turn off topological bonds (angles, etc) by setting
their bond (angle, etc) types to negative values.  This command
includes the turned\sphinxhyphen{}off bonds (angles, etc) in the count for each
type:
\begin{itemize}
\item {} 
\sphinxAtStartPar
{\hyperref[\detokenize{fix_shake::doc}]{\sphinxcrossref{\DUrole{doc}{fix shake}}}}

\item {} 
\sphinxAtStartPar
{\hyperref[\detokenize{delete_bonds::doc}]{\sphinxcrossref{\DUrole{doc}{delete\_bonds}}}}

\end{itemize}

\sphinxAtStartPar
These commands can create and/or break topological bonds (angles,
etc).  In the case of breaking, they remove the bond (angle, etc) from
the system, so that they no longer exist ({\hyperref[\detokenize{bond_quartic::doc}]{\sphinxcrossref{\DUrole{doc}{bond\_style quartic}}}} and {\hyperref[\detokenize{Howto_bpm::doc}]{\sphinxcrossref{\DUrole{doc}{BPM bond styles}}}} are exceptions,
see the discussion below).  Thus they are not included in the counts
for each type:
\begin{itemize}
\item {} 
\sphinxAtStartPar
{\hyperref[\detokenize{delete_bonds::doc}]{\sphinxcrossref{\DUrole{doc}{delete\_bonds remove}}}}

\item {} 
\sphinxAtStartPar
{\hyperref[\detokenize{bond_quartic::doc}]{\sphinxcrossref{\DUrole{doc}{bond\_style quartic}}}}

\item {} 
\sphinxAtStartPar
{\hyperref[\detokenize{fix_bond_react::doc}]{\sphinxcrossref{\DUrole{doc}{fix bond/react}}}}

\item {} 
\sphinxAtStartPar
{\hyperref[\detokenize{fix_bond_create::doc}]{\sphinxcrossref{\DUrole{doc}{fix bond/create}}}}

\item {} 
\sphinxAtStartPar
{\hyperref[\detokenize{fix_bond_break::doc}]{\sphinxcrossref{\DUrole{doc}{fix bond/break}}}}

\item {} 
\sphinxAtStartPar
{\hyperref[\detokenize{Howto_bpm::doc}]{\sphinxcrossref{\DUrole{doc}{BPM package}}}} bond styles

\end{itemize}


\bigskip\hrule\bigskip


\sphinxAtStartPar
If the \sphinxstyleemphasis{mode} setting is \sphinxstyleemphasis{atom} then the count of atoms for each atom
type is tallied.  Only atoms in the specified group are counted.

\sphinxAtStartPar
The atom count for each type can be normalized by the total number of
atoms like so:

\begin{sphinxVerbatim}[commandchars=\\\{\}]
\PYG{k}{compute }\PYG{n+nv+nvIdentifier}{typevec}\PYG{+w}{ }\PYG{n+nv+nvIdentifier}{all}\PYG{+w}{ }\PYG{n}{count}\PYG{o}{/}\PYG{n}{type}\PYG{+w}{ }\PYG{n}{atom}\PYG{+w}{ }\PYG{c}{\PYGZsh{} number of atoms of each type}
\PYG{k}{variable }\PYG{n+nv+nvIdentifier}{normtypes}\PYG{+w}{ }\PYG{n}{vector}\PYG{+w}{ }\PYG{n}{c\PYGZus{}typevec}\PYG{o}{/}\PYG{n}{atoms}\PYG{+w}{ }\PYG{c}{\PYGZsh{} divide by total number of atoms}
\PYG{k}{variable }\PYG{n+nv+nvIdentifier}{ntypes}\PYG{+w}{ }\PYG{n}{equal}\PYG{+w}{ }\PYG{n}{extract\PYGZus{}setting}\PYG{n+nv}{(}\PYG{n+nv}{ntypes}\PYG{n+nv}{)}\PYG{+w}{ }\PYG{c}{\PYGZsh{} number of atom types}
\PYG{k}{thermo\PYGZus{}style}\PYG{+w}{ }\PYG{n}{custom}\PYG{+w}{ }\PYG{n}{step}\PYG{+w}{ }\PYG{n}{v\PYGZus{}normtypes[}\PYG{o}{*}\PYG{n+nv}{\PYGZdl{}\PYGZob{}}\PYG{n+nv}{ntypes}\PYG{n+nv}{\PYGZcb{}}\PYG{n}{]}\PYG{+w}{ }\PYG{c}{\PYGZsh{} vector variable needs upper limit}
\end{sphinxVerbatim}

\sphinxAtStartPar
Similarly, bond counts can be normalized by the total number of bonds.
The same goes for angles, dihedrals, and impropers (see below).

\sphinxAtStartPar
If the \sphinxstyleemphasis{mode} setting is \sphinxstyleemphasis{bond} then the count of bonds for each bond
type is tallied.  Only bonds with both atoms in the specified group
are counted.

\sphinxAtStartPar
For \sphinxstyleemphasis{mode} = \sphinxstyleemphasis{bond}, broken bonds with a bond type of zero are also
counted.  The {\hyperref[\detokenize{bond_quartic::doc}]{\sphinxcrossref{\DUrole{doc}{bond\_style quartic}}}} and {\hyperref[\detokenize{Howto_bpm::doc}]{\sphinxcrossref{\DUrole{doc}{BPM
bond styles}}}} break bonds by doing this.  See the
{\hyperref[\detokenize{Howto_broken_bonds::doc}]{\sphinxcrossref{\DUrole{doc}{Howto broken bonds}}}} doc page for more details.
Note that the group setting is ignored for broken bonds; all broken
bonds in the system are counted.

\sphinxAtStartPar
If the \sphinxstyleemphasis{mode} setting is \sphinxstyleemphasis{angle} then the count of angles for each
angle type is tallied.  Only angles with all 3 atoms in the specified
group are counted.

\sphinxAtStartPar
If the \sphinxstyleemphasis{mode} setting is \sphinxstyleemphasis{dihedral} then the count of dihedrals for
each dihedral type is tallied.  Only dihedrals with all 4 atoms in the
specified group are counted.

\sphinxAtStartPar
If the \sphinxstyleemphasis{mode} setting is \sphinxstyleemphasis{improper} then the count of impropers for
each improper type is tallied.  Only impropers with all 4 atoms in the
specified group are counted.


\bigskip\hrule\bigskip



\subsection{Output info}
\label{\detokenize{compute_count_type:output-info}}
\sphinxAtStartPar
This compute calculates a global vector of counts.  If the mode is
\sphinxstyleemphasis{atom} or \sphinxstyleemphasis{bond} or \sphinxstyleemphasis{angle} or \sphinxstyleemphasis{dihedral} or \sphinxstyleemphasis{improper}, then the
vector length is the number of atom types or bond types or angle types
or dihedral types or improper types, respectively.

\sphinxAtStartPar
If the mode is \sphinxstyleemphasis{bond} this compute also calculates a global scalar
which is the number of broken bonds with type = 0, as explained above.

\sphinxAtStartPar
These values can be used by any command that uses global scalar or
vector values from a compute as input.  See the {\hyperref[\detokenize{Howto_output::doc}]{\sphinxcrossref{\DUrole{doc}{Howto output}}}} page for an overview of LAMMPS output options.

\sphinxAtStartPar
The scalar and vector values calculated by this compute are both
“intensive”.


\subsection{Restrictions}
\label{\detokenize{compute_count_type:restrictions}}
\sphinxAtStartPar
none


\subsection{Related commands}
\label{\detokenize{compute_count_type:related-commands}}
\sphinxAtStartPar
none


\subsection{Default}
\label{\detokenize{compute_count_type:default}}
\sphinxAtStartPar
none

\sphinxstepscope

\index{compute damage/atom@\spxentry{compute damage/atom}}\ignorespaces 

\section{compute damage/atom command}
\label{\detokenize{compute_damage_atom:compute-damage-atom-command}}\label{\detokenize{compute_damage_atom:index-0}}\label{\detokenize{compute_damage_atom::doc}}

\subsection{Syntax}
\label{\detokenize{compute_damage_atom:syntax}}
\begin{sphinxVerbatim}[commandchars=\\\{\}]
\PYG{k}{compute }\PYG{n+nv+nvIdentifier}{ID}\PYG{+w}{ }\PYG{n+nv+nvIdentifier}{group\PYGZhy{}ID}\PYG{+w}{ }\PYG{n}{damage}\PYG{o}{/}\PYG{n}{atom}
\end{sphinxVerbatim}
\begin{itemize}
\item {} 
\sphinxAtStartPar
ID, group\sphinxhyphen{}ID are documented in {\hyperref[\detokenize{compute::doc}]{\sphinxcrossref{\DUrole{doc}{compute}}}} command

\item {} 
\sphinxAtStartPar
damage/atom = style name of this compute command

\end{itemize}


\subsection{Examples}
\label{\detokenize{compute_damage_atom:examples}}
\begin{sphinxVerbatim}[commandchars=\\\{\}]
\PYG{k}{compute }\PYG{n+nv+nvIdentifier}{1}\PYG{+w}{ }\PYG{n+nv+nvIdentifier}{all}\PYG{+w}{ }\PYG{n}{damage}\PYG{o}{/}\PYG{n}{atom}
\end{sphinxVerbatim}


\subsection{Description}
\label{\detokenize{compute_damage_atom:description}}
\sphinxAtStartPar
Define a computation that calculates the per\sphinxhyphen{}atom damage for each atom
in a group.  This is a quantity relevant for {\hyperref[\detokenize{pair_peri::doc}]{\sphinxcrossref{\DUrole{doc}{Peridynamics models}}}}.  See \sphinxhref{PDF/PDLammps\_overview.pdf}{this document} for an
overview of LAMMPS commands for Peridynamics modeling.

\sphinxAtStartPar
The “damage” of a Peridynamics particles is based on the bond breakage
between the particle and its neighbors.  If all the bonds are broken
the particle is considered to be fully damaged.

\sphinxAtStartPar
See the {\hyperref[\detokenize{Howto_peri::doc}]{\sphinxcrossref{\DUrole{doc}{Peridynamics Howto}}}} for a formal definition
of “damage” and more details about Peridynamics as it is implemented in
LAMMPS.

\sphinxAtStartPar
This command can be used with all the Peridynamic pair styles.

\sphinxAtStartPar
The damage value will be 0.0 for atoms not in the specified compute
group.


\subsection{Output info}
\label{\detokenize{compute_damage_atom:output-info}}
\sphinxAtStartPar
This compute calculates a per\sphinxhyphen{}atom vector, which can be accessed by
any command that uses per\sphinxhyphen{}atom values from a compute as input.  See
the {\hyperref[\detokenize{Howto_output::doc}]{\sphinxcrossref{\DUrole{doc}{Howto output}}}} page for an overview of
LAMMPS output options.

\sphinxAtStartPar
The per\sphinxhyphen{}atom vector values are unitless numbers (damage) \(\ge 0.0\).


\subsection{Restrictions}
\label{\detokenize{compute_damage_atom:restrictions}}
\sphinxAtStartPar
This compute is part of the PERI package.  It is only enabled if LAMMPS
was built with that package.  See the {\hyperref[\detokenize{Build_package::doc}]{\sphinxcrossref{\DUrole{doc}{Build package}}}} page for more info.


\subsection{Related commands}
\label{\detokenize{compute_damage_atom:related-commands}}
\sphinxAtStartPar
{\hyperref[\detokenize{compute_dilatation_atom::doc}]{\sphinxcrossref{\DUrole{doc}{compute dilatation/atom}}}},
{\hyperref[\detokenize{compute_plasticity_atom::doc}]{\sphinxcrossref{\DUrole{doc}{compute plasticity/atom}}}}


\subsection{Default}
\label{\detokenize{compute_damage_atom:default}}
\sphinxAtStartPar
none

\sphinxstepscope

\index{compute dihedral@\spxentry{compute dihedral}}\ignorespaces 

\section{compute dihedral command}
\label{\detokenize{compute_dihedral:compute-dihedral-command}}\label{\detokenize{compute_dihedral:index-0}}\label{\detokenize{compute_dihedral::doc}}

\subsection{Syntax}
\label{\detokenize{compute_dihedral:syntax}}
\begin{sphinxVerbatim}[commandchars=\\\{\}]
\PYG{k}{compute }\PYG{n+nv+nvIdentifier}{ID}\PYG{+w}{ }\PYG{n+nv+nvIdentifier}{group\PYGZhy{}ID}\PYG{+w}{ }\PYG{n}{dihedral}
\end{sphinxVerbatim}
\begin{itemize}
\item {} 
\sphinxAtStartPar
ID, group\sphinxhyphen{}ID are documented in {\hyperref[\detokenize{compute::doc}]{\sphinxcrossref{\DUrole{doc}{compute}}}} command

\item {} 
\sphinxAtStartPar
dihedral = style name of this compute command

\end{itemize}


\subsection{Examples}
\label{\detokenize{compute_dihedral:examples}}
\begin{sphinxVerbatim}[commandchars=\\\{\}]
\PYG{k}{compute }\PYG{n+nv+nvIdentifier}{1}\PYG{+w}{ }\PYG{n+nv+nvIdentifier}{all}\PYG{+w}{ }\PYG{n}{dihedral}
\end{sphinxVerbatim}


\subsection{Description}
\label{\detokenize{compute_dihedral:description}}
\sphinxAtStartPar
Define a computation that extracts the dihedral energy calculated by
each of the dihedral sub\sphinxhyphen{}styles used in the {\hyperref[\detokenize{dihedral_hybrid::doc}]{\sphinxcrossref{\DUrole{doc}{dihedral\_style hybrid}}}} command.  These values are made
accessible for output or further processing by other commands.  The
group specified for this command is ignored.

\sphinxAtStartPar
This compute is useful when using {\hyperref[\detokenize{dihedral_hybrid::doc}]{\sphinxcrossref{\DUrole{doc}{dihedral\_style hybrid}}}} if you want to know the portion of the
total energy contributed by one or more of the hybrid sub\sphinxhyphen{}styles.


\subsection{Output info}
\label{\detokenize{compute_dihedral:output-info}}
\sphinxAtStartPar
This compute calculates a global vector of length \(N\), where \(N\)
is the number of sub\_styles defined by the
{\hyperref[\detokenize{dihedral_style::doc}]{\sphinxcrossref{\DUrole{doc}{dihedral\_style hybrid}}}} command, which can be accessed by
the indices 1 through \(N\).  These values can be used by any command that
uses global scalar or vector values from a compute as input.  See the
{\hyperref[\detokenize{Howto_output::doc}]{\sphinxcrossref{\DUrole{doc}{Howto output}}}} page for an overview of LAMMPS output
options.

\sphinxAtStartPar
The vector values are “extensive” and will be in energy
{\hyperref[\detokenize{units::doc}]{\sphinxcrossref{\DUrole{doc}{units}}}}.


\subsection{Restrictions}
\label{\detokenize{compute_dihedral:restrictions}}\begin{quote}

\sphinxAtStartPar
none
\end{quote}


\subsection{Related commands}
\label{\detokenize{compute_dihedral:related-commands}}
\sphinxAtStartPar
{\hyperref[\detokenize{compute_pe::doc}]{\sphinxcrossref{\DUrole{doc}{compute pe}}}}, {\hyperref[\detokenize{compute_pair::doc}]{\sphinxcrossref{\DUrole{doc}{compute pair}}}}


\subsection{Default}
\label{\detokenize{compute_dihedral:default}}
\sphinxAtStartPar
none

\sphinxstepscope

\index{compute dihedral/local@\spxentry{compute dihedral/local}}\ignorespaces 

\section{compute dihedral/local command}
\label{\detokenize{compute_dihedral_local:compute-dihedral-local-command}}\label{\detokenize{compute_dihedral_local:index-0}}\label{\detokenize{compute_dihedral_local::doc}}

\subsection{Syntax}
\label{\detokenize{compute_dihedral_local:syntax}}
\begin{sphinxVerbatim}[commandchars=\\\{\}]
\PYG{k}{compute }\PYG{n+nv+nvIdentifier}{ID}\PYG{+w}{ }\PYG{n+nv+nvIdentifier}{group\PYGZhy{}ID}\PYG{+w}{ }\PYG{n}{dihedral}\PYG{o}{/}\PYG{n}{local}\PYG{+w}{ }\PYG{n}{value1}\PYG{+w}{ }\PYG{n}{value2}\PYG{+w}{ }\PYG{n}{...}\PYG{+w}{ }\PYG{n}{keyword}\PYG{+w}{ }\PYG{n}{args}\PYG{+w}{ }\PYG{n}{...}
\end{sphinxVerbatim}
\begin{itemize}
\item {} 
\sphinxAtStartPar
ID, group\sphinxhyphen{}ID are documented in {\hyperref[\detokenize{compute::doc}]{\sphinxcrossref{\DUrole{doc}{compute}}}} command

\item {} 
\sphinxAtStartPar
dihedral/local = style name of this compute command

\item {} 
\sphinxAtStartPar
one or more values may be appended

\item {} 
\sphinxAtStartPar
value = \sphinxstyleemphasis{phi} or \sphinxstyleemphasis{v\_name}
\begin{sphinxalltt}
\sphinxstyleemphasis{phi} = tabulate dihedral angles
\sphinxstyleemphasis{v\_name} = equal\sphinxhyphen{}style variable with name (see below)
\end{sphinxalltt}

\item {} 
\sphinxAtStartPar
zero or more keyword/args pairs may be appended

\item {} 
\sphinxAtStartPar
keyword = \sphinxstyleemphasis{set}
\begin{sphinxalltt}
\sphinxstyleemphasis{set} args = phi name
  phi = only currently allowed arg
  name = name of variable to set with phi
\end{sphinxalltt}

\end{itemize}


\subsection{Examples}
\label{\detokenize{compute_dihedral_local:examples}}
\begin{sphinxVerbatim}[commandchars=\\\{\}]
\PYG{k}{compute }\PYG{n+nv+nvIdentifier}{1}\PYG{+w}{ }\PYG{n+nv+nvIdentifier}{all}\PYG{+w}{ }\PYG{n}{dihedral}\PYG{o}{/}\PYG{n}{local}\PYG{+w}{ }\PYG{n}{phi}
\PYG{k}{compute }\PYG{n+nv+nvIdentifier}{1}\PYG{+w}{ }\PYG{n+nv+nvIdentifier}{all}\PYG{+w}{ }\PYG{n}{dihedral}\PYG{o}{/}\PYG{n}{local}\PYG{+w}{ }\PYG{n}{phi}\PYG{+w}{ }\PYG{n}{v\PYGZus{}cos}\PYG{+w}{ }\PYG{n}{set}\PYG{+w}{ }\PYG{n}{phi}\PYG{+w}{ }\PYG{n}{p}
\end{sphinxVerbatim}


\subsection{Description}
\label{\detokenize{compute_dihedral_local:description}}
\sphinxAtStartPar
Define a computation that calculates properties of individual dihedral
interactions.  The number of datums generated, aggregated across all
processors, equals the number of dihedral angles in the system, modified
by the group parameter as explained below.

\sphinxAtStartPar
The value \sphinxstyleemphasis{phi} (\(\phi\)) is the dihedral angle, as defined in the diagram
on the {\hyperref[\detokenize{dihedral_style::doc}]{\sphinxcrossref{\DUrole{doc}{dihedral\_style}}}} doc page.

\sphinxAtStartPar
The value \sphinxstyleemphasis{v\_name} can be used together with the \sphinxstyleemphasis{set} keyword to compute a
user\sphinxhyphen{}specified function of the dihedral angle \(\phi\).  The \sphinxstyleemphasis{name}
specified for the \sphinxstyleemphasis{v\_name} value is the name of an
{\hyperref[\detokenize{variable::doc}]{\sphinxcrossref{\DUrole{doc}{equal\sphinxhyphen{}style variable}}}} which should evaluate a formula based on
a variable which will store the angle \(\phi\).  This other variable must
be an {\hyperref[\detokenize{variable::doc}]{\sphinxcrossref{\DUrole{doc}{internal\sphinxhyphen{}style variable}}}} defined in the input
script; its initial numeric value can be anything.  It must be an
internal\sphinxhyphen{}style variable, because this command resets its value
directly.  The \sphinxstyleemphasis{set} keyword is used to identify the name of this
other variable associated with \(\phi\).

\sphinxAtStartPar
Note that the value of \(\phi\) for each angle which stored in the internal
variable is in radians, not degrees.

\sphinxAtStartPar
As an example, these commands can be added to the bench/in.rhodo
script to compute the \(\cos\phi\) and \(\cos^2\phi\) of every dihedral
angle in the system and output the statistics in various ways:

\begin{sphinxVerbatim}[commandchars=\\\{\}]
\PYG{k}{variable }\PYG{n+nv+nvIdentifier}{p}\PYG{+w}{ }\PYG{n}{internal}\PYG{+w}{ }\PYG{l+m}{0.0}
\PYG{k}{variable }\PYG{n+nv+nvIdentifier}{cos}\PYG{+w}{ }\PYG{n}{equal}\PYG{+w}{ }\PYG{n}{cos}\PYG{n+nv}{(}\PYG{n+nv}{v\PYGZus{}p}\PYG{n+nv}{)}
\PYG{k}{variable }\PYG{n+nv+nvIdentifier}{cossq}\PYG{+w}{ }\PYG{n}{equal}\PYG{+w}{ }\PYG{n}{cos}\PYG{n+nv}{(}\PYG{n+nv}{v\PYGZus{}p}\PYG{n+nv}{)}\PYG{o}{*}\PYG{n}{cos}\PYG{n+nv}{(}\PYG{n+nv}{v\PYGZus{}p}\PYG{n+nv}{)}

\PYG{k}{compute }\PYG{n+nv+nvIdentifier}{1}\PYG{+w}{ }\PYG{n+nv+nvIdentifier}{all}\PYG{+w}{ }\PYG{n}{property}\PYG{o}{/}\PYG{n}{local}\PYG{+w}{ }\PYG{n}{datom1}\PYG{+w}{ }\PYG{n}{datom2}\PYG{+w}{ }\PYG{n}{datom3}\PYG{+w}{ }\PYG{n}{datom4}\PYG{+w}{ }\PYG{n}{dtype}
\PYG{k}{compute }\PYG{n+nv+nvIdentifier}{2}\PYG{+w}{ }\PYG{n+nv+nvIdentifier}{all}\PYG{+w}{ }\PYG{n}{dihedral}\PYG{o}{/}\PYG{n}{local}\PYG{+w}{ }\PYG{n}{phi}\PYG{+w}{ }\PYG{n}{v\PYGZus{}cos}\PYG{+w}{ }\PYG{n}{v\PYGZus{}cossq}\PYG{+w}{ }\PYG{n}{set}\PYG{+w}{ }\PYG{n}{phi}\PYG{+w}{ }\PYG{n}{p}
\PYG{k}{dump }\PYG{n+nv+nvIdentifier}{1}\PYG{+w}{ }\PYG{n+nv+nvIdentifier}{all}\PYG{+w}{ }\PYG{n}{local}\PYG{+w}{ }\PYG{l+m}{100}\PYG{+w}{ }\PYG{n}{tmp.dump}\PYG{+w}{ }\PYG{n}{c\PYGZus{}1[}\PYG{o}{*}\PYG{n}{]}\PYG{+w}{ }\PYG{n}{c\PYGZus{}2[}\PYG{o}{*}\PYG{n}{]}

\PYG{k}{compute }\PYG{n+nv+nvIdentifier}{3}\PYG{+w}{ }\PYG{n+nv+nvIdentifier}{all}\PYG{+w}{ }\PYG{n}{reduce}\PYG{+w}{ }\PYG{n}{ave}\PYG{+w}{ }\PYG{n}{c\PYGZus{}2[}\PYG{o}{*}\PYG{n}{]}
\PYG{k}{thermo\PYGZus{}style}\PYG{+w}{ }\PYG{n}{custom}\PYG{+w}{ }\PYG{n}{step}\PYG{+w}{ }\PYG{n}{temp}\PYG{+w}{ }\PYG{n}{press}\PYG{+w}{ }\PYG{n}{c\PYGZus{}3[}\PYG{o}{*}\PYG{n}{]}

\PYG{k}{fix }\PYG{n+nv+nvIdentifier}{10}\PYG{+w}{ }\PYG{n+nv+nvIdentifier}{all}\PYG{+w}{ }\PYG{n}{ave}\PYG{o}{/}\PYG{n}{histo}\PYG{+w}{ }\PYG{l+m}{10}\PYG{+w}{ }\PYG{l+m}{10}\PYG{+w}{ }\PYG{l+m}{100}\PYG{+w}{ }\PYG{o}{\PYGZhy{}}\PYG{l+m}{1}\PYG{+w}{ }\PYG{l+m}{1}\PYG{+w}{ }\PYG{l+m}{20}\PYG{+w}{ }\PYG{n}{c\PYGZus{}2[2]}\PYG{+w}{ }\PYG{n}{mode}\PYG{+w}{ }\PYG{n}{vector}\PYG{+w}{ }\PYG{n}{file}\PYG{+w}{ }\PYG{n}{tmp.histo}
\end{sphinxVerbatim}

\sphinxAtStartPar
The {\hyperref[\detokenize{dump::doc}]{\sphinxcrossref{\DUrole{doc}{dump local}}}} command will output the angle (\(\phi\)),
\(\cos(\phi)\), and \(\cos^2(\phi)\) for every dihedral in the system.
The {\hyperref[\detokenize{thermo_style::doc}]{\sphinxcrossref{\DUrole{doc}{thermo\_style}}}} command will print the average of
those quantities via the {\hyperref[\detokenize{compute_reduce::doc}]{\sphinxcrossref{\DUrole{doc}{compute reduce}}}} command
with thermo output.  And the {\hyperref[\detokenize{fix_ave_histo::doc}]{\sphinxcrossref{\DUrole{doc}{fix ave/histo}}}}
command will histogram the cosine(angle) values and write them to a file.


\bigskip\hrule\bigskip


\sphinxAtStartPar
The local data stored by this command is generated by looping over all
the atoms owned on a processor and their dihedrals.  A dihedral will
only be included if all four atoms in the dihedral are in the specified
compute group.

\sphinxAtStartPar
Note that as atoms migrate from processor to processor, there will be
no consistent ordering of the entries within the local vector or array
from one timestep to the next.  The only consistency that is
guaranteed is that the ordering on a particular timestep will be the
same for local vectors or arrays generated by other compute commands.
For example, dihedral output from the
{\hyperref[\detokenize{compute_property_local::doc}]{\sphinxcrossref{\DUrole{doc}{compute property/local}}}} command can be combined
with data from this command and output by the {\hyperref[\detokenize{dump::doc}]{\sphinxcrossref{\DUrole{doc}{dump local}}}}
command in a consistent way.

\sphinxAtStartPar
Here is an example of how to do this:

\begin{sphinxVerbatim}[commandchars=\\\{\}]
\PYG{k}{compute }\PYG{n+nv+nvIdentifier}{1}\PYG{+w}{ }\PYG{n+nv+nvIdentifier}{all}\PYG{+w}{ }\PYG{n}{property}\PYG{o}{/}\PYG{n}{local}\PYG{+w}{ }\PYG{n}{dtype}\PYG{+w}{ }\PYG{n}{datom1}\PYG{+w}{ }\PYG{n}{datom2}\PYG{+w}{ }\PYG{n}{datom3}\PYG{+w}{ }\PYG{n}{datom4}
\PYG{k}{compute }\PYG{n+nv+nvIdentifier}{2}\PYG{+w}{ }\PYG{n+nv+nvIdentifier}{all}\PYG{+w}{ }\PYG{n}{dihedral}\PYG{o}{/}\PYG{n}{local}\PYG{+w}{ }\PYG{n}{phi}
\PYG{k}{dump }\PYG{n+nv+nvIdentifier}{1}\PYG{+w}{ }\PYG{n+nv+nvIdentifier}{all}\PYG{+w}{ }\PYG{n}{local}\PYG{+w}{ }\PYG{l+m}{1000}\PYG{+w}{ }\PYG{n}{tmp.dump}\PYG{+w}{ }\PYG{n}{index}\PYG{+w}{ }\PYG{n}{c\PYGZus{}1[1]}\PYG{+w}{ }\PYG{n}{c\PYGZus{}1[2]}\PYG{+w}{ }\PYG{n}{c\PYGZus{}1[3]}\PYG{+w}{ }\PYG{n}{c\PYGZus{}1[4]}\PYG{+w}{ }\PYG{n}{c\PYGZus{}1[5]}\PYG{+w}{ }\PYG{n}{c\PYGZus{}2[1]}
\end{sphinxVerbatim}


\subsection{Output info}
\label{\detokenize{compute_dihedral_local:output-info}}
\sphinxAtStartPar
This compute calculates a local vector or local array depending on the
number of values.  The length of the vector or number of rows in the
array is the number of dihedrals.  If a single value is specified, a
local vector is produced.  If two or more values are specified, a
local array is produced where the number of columns is equal to the number of
values.  The vector or array can be accessed by any command that uses
local values from a compute as input.  See the
{\hyperref[\detokenize{Howto_output::doc}]{\sphinxcrossref{\DUrole{doc}{Howto output}}}} page for an overview of LAMMPS output
options.

\sphinxAtStartPar
The output for \sphinxstyleemphasis{phi} will be in degrees.


\subsection{Restrictions}
\label{\detokenize{compute_dihedral_local:restrictions}}\begin{quote}

\sphinxAtStartPar
none
\end{quote}


\subsection{Related commands}
\label{\detokenize{compute_dihedral_local:related-commands}}
\sphinxAtStartPar
{\hyperref[\detokenize{dump::doc}]{\sphinxcrossref{\DUrole{doc}{dump local}}}}, {\hyperref[\detokenize{compute_property_local::doc}]{\sphinxcrossref{\DUrole{doc}{compute property/local}}}}


\subsection{Default}
\label{\detokenize{compute_dihedral_local:default}}
\sphinxAtStartPar
none

\sphinxstepscope

\index{compute dilatation/atom@\spxentry{compute dilatation/atom}}\ignorespaces 

\section{compute dilatation/atom command}
\label{\detokenize{compute_dilatation_atom:compute-dilatation-atom-command}}\label{\detokenize{compute_dilatation_atom:index-0}}\label{\detokenize{compute_dilatation_atom::doc}}

\subsection{Syntax}
\label{\detokenize{compute_dilatation_atom:syntax}}
\begin{sphinxVerbatim}[commandchars=\\\{\}]
\PYG{k}{compute }\PYG{n+nv+nvIdentifier}{ID}\PYG{+w}{ }\PYG{n+nv+nvIdentifier}{group\PYGZhy{}ID}\PYG{+w}{ }\PYG{n}{dilatation}\PYG{o}{/}\PYG{n}{atom}
\end{sphinxVerbatim}
\begin{itemize}
\item {} 
\sphinxAtStartPar
ID, group\sphinxhyphen{}ID are documented in compute command

\item {} 
\sphinxAtStartPar
dilatation/atom = style name of this compute command

\end{itemize}


\subsection{Examples}
\label{\detokenize{compute_dilatation_atom:examples}}
\begin{sphinxVerbatim}[commandchars=\\\{\}]
\PYG{k}{compute }\PYG{n+nv+nvIdentifier}{1}\PYG{+w}{ }\PYG{n+nv+nvIdentifier}{all}\PYG{+w}{ }\PYG{n}{dilatation}\PYG{o}{/}\PYG{n}{atom}
\end{sphinxVerbatim}


\subsection{Description}
\label{\detokenize{compute_dilatation_atom:description}}
\sphinxAtStartPar
Define a computation that calculates the per\sphinxhyphen{}atom dilatation for each
atom in a group.  This is a quantity relevant for {\hyperref[\detokenize{pair_peri::doc}]{\sphinxcrossref{\DUrole{doc}{Peridynamics
models}}}}.  See \sphinxhref{PDF/PDLammps\_overview.pdf}{this document}
for an overview of LAMMPS commands for Peridynamics modeling.

\sphinxAtStartPar
For small deformation, dilatation of is the measure of the volumetric
strain.

\sphinxAtStartPar
The dilatation \(\theta\) for each peridynamic particle \(i\) is
calculated as a sum over its neighbors with unbroken bonds, where the
contribution of the \(ij\) pair is a function of the change in bond
length (versus the initial length in the reference state), the volume
fraction of the particles and an influence function.  See the
{\hyperref[\detokenize{Howto_peri::doc}]{\sphinxcrossref{\DUrole{doc}{Peridynamics Howto}}}} for a formal definition of
dilatation.

\sphinxAtStartPar
This command can only be used with a subset of the Peridynamic
{\hyperref[\detokenize{pair_peri::doc}]{\sphinxcrossref{\DUrole{doc}{pair styles}}}}: \sphinxstyleemphasis{peri/lps}, \sphinxstyleemphasis{peri/ves}, and \sphinxstyleemphasis{peri/eps}.

\sphinxAtStartPar
The dilatation value will be 0.0 for atoms not in the specified
compute group.


\subsection{Output info}
\label{\detokenize{compute_dilatation_atom:output-info}}
\sphinxAtStartPar
This compute calculates a per\sphinxhyphen{}atom vector, which can be accessed by
any command that uses per\sphinxhyphen{}atom values from a compute as input.  See
the {\hyperref[\detokenize{Howto_output::doc}]{\sphinxcrossref{\DUrole{doc}{Howto output}}}} page for an overview of
LAMMPS output options.

\sphinxAtStartPar
The per\sphinxhyphen{}atom vector values are unitless numbers \((\theta \ge 0.0)\).


\subsection{Restrictions}
\label{\detokenize{compute_dilatation_atom:restrictions}}
\sphinxAtStartPar
This compute is part of the PERI package.  It is only enabled if LAMMPS
was built with that package.  See the {\hyperref[\detokenize{Build_package::doc}]{\sphinxcrossref{\DUrole{doc}{Build package}}}} page for more info.


\subsection{Related commands}
\label{\detokenize{compute_dilatation_atom:related-commands}}
\sphinxAtStartPar
{\hyperref[\detokenize{compute_damage_atom::doc}]{\sphinxcrossref{\DUrole{doc}{compute damage/atom}}}},
{\hyperref[\detokenize{compute_plasticity_atom::doc}]{\sphinxcrossref{\DUrole{doc}{compute plasticity/atom}}}}


\subsection{Default}
\label{\detokenize{compute_dilatation_atom:default}}
\sphinxAtStartPar
none

\sphinxstepscope

\index{compute dipole@\spxentry{compute dipole}}\ignorespaces 
\index{compute dipole/tip4p@\spxentry{compute dipole/tip4p}}\ignorespaces 

\section{compute dipole command}
\label{\detokenize{compute_dipole:compute-dipole-command}}\label{\detokenize{compute_dipole:index-1}}\label{\detokenize{compute_dipole:index-0}}\label{\detokenize{compute_dipole::doc}}

\section{compute dipole/tip4p command}
\label{\detokenize{compute_dipole:compute-dipole-tip4p-command}}

\subsection{Syntax}
\label{\detokenize{compute_dipole:syntax}}
\begin{sphinxVerbatim}[commandchars=\\\{\}]
\PYG{k}{compute }\PYG{n+nv+nvIdentifier}{ID}\PYG{+w}{ }\PYG{n+nv+nvIdentifier}{group\PYGZhy{}ID}\PYG{+w}{ }\PYG{n}{style}\PYG{+w}{ }\PYG{n}{arg}
\end{sphinxVerbatim}
\begin{itemize}
\item {} 
\sphinxAtStartPar
ID, group\sphinxhyphen{}ID are documented in {\hyperref[\detokenize{compute::doc}]{\sphinxcrossref{\DUrole{doc}{compute}}}} command

\item {} 
\sphinxAtStartPar
style = \sphinxstyleemphasis{dipole} or \sphinxstyleemphasis{dipole/tip4p}

\item {} 
\sphinxAtStartPar
arg = \sphinxstyleemphasis{mass} or \sphinxstyleemphasis{geometry} = use COM or geometric center for charged chunk correction (optional)

\end{itemize}


\subsection{Examples}
\label{\detokenize{compute_dipole:examples}}
\begin{sphinxVerbatim}[commandchars=\\\{\}]
\PYG{k}{compute }\PYG{n+nv+nvIdentifier}{1}\PYG{+w}{ }\PYG{n+nv+nvIdentifier}{fluid}\PYG{+w}{ }\PYG{n}{dipole}
\PYG{k}{compute }\PYG{n+nv+nvIdentifier}{dw}\PYG{+w}{ }\PYG{n+nv+nvIdentifier}{water}\PYG{+w}{ }\PYG{n}{dipole}\PYG{+w}{ }\PYG{n}{geometry}
\PYG{k}{compute }\PYG{n+nv+nvIdentifier}{dw}\PYG{+w}{ }\PYG{n+nv+nvIdentifier}{water}\PYG{+w}{ }\PYG{n}{dipole}\PYG{o}{/}\PYG{n}{tip4p}
\end{sphinxVerbatim}


\subsection{Description}
\label{\detokenize{compute_dipole:description}}
\sphinxAtStartPar
Define a computation that calculates the dipole vector and total dipole
for a group of atoms.

\sphinxAtStartPar
These computes calculate the x,y,z coordinates of the dipole vector and
the total dipole moment for the atoms in the compute group.  This
includes all effects due to atoms passing through periodic boundaries.
For a group with a net charge the resulting dipole is made position
independent by subtracting the position vector of the center of mass or
geometric center times the net charge from the computed dipole
vector.  Both per\sphinxhyphen{}atom charges and per\sphinxhyphen{}atom dipole moments, if present,
contribute to the computed dipole.

\sphinxAtStartPar
\DUrole{versionmodified,added}{Added in version 28Mar2023.}

\sphinxAtStartPar
Compute \sphinxstyleemphasis{dipole/tip4p} includes adjustments for the charge carrying
point M in molecules with TIP4P water geometry.  The corresponding
parameters are extracted from the pair style.

\begin{sphinxadmonition}{note}{Note:}
\sphinxAtStartPar
The coordinates of an atom contribute to the dipole in “unwrapped”
form, by using the image flags associated with each atom.  See the
{\hyperref[\detokenize{dump::doc}]{\sphinxcrossref{\DUrole{doc}{dump custom}}}} command for a discussion of “unwrapped”
coordinates.  See the Atoms section of the {\hyperref[\detokenize{read_data::doc}]{\sphinxcrossref{\DUrole{doc}{read\_data}}}} command for a discussion of image flags and how they are
set for each atom.  You can reset the image flags (e.g., to 0) before
invoking this compute by using the {\hyperref[\detokenize{set::doc}]{\sphinxcrossref{\DUrole{doc}{set image}}}} command.
\end{sphinxadmonition}


\subsection{Output info}
\label{\detokenize{compute_dipole:output-info}}
\sphinxAtStartPar
These computes calculate a global scalar containing the magnitude of the
computed dipole moment and a global vector of length 3 with the dipole
vector.  See the {\hyperref[\detokenize{Howto_output::doc}]{\sphinxcrossref{\DUrole{doc}{Howto output}}}} page for an overview
of LAMMPS output options.

\sphinxAtStartPar
The computed values are “intensive”.  The array values will be in
dipole units (i.e., charge {\hyperref[\detokenize{units::doc}]{\sphinxcrossref{\DUrole{doc}{units}}}} times distance
{\hyperref[\detokenize{units::doc}]{\sphinxcrossref{\DUrole{doc}{units}}}}).


\subsection{Restrictions}
\label{\detokenize{compute_dipole:restrictions}}
\sphinxAtStartPar
Compute style \sphinxstyleemphasis{dipole/tip4p} is part of the EXTRA\sphinxhyphen{}COMPUTE package. It is
only enabled if LAMMPS was built with that package.  See the {\hyperref[\detokenize{Build_package::doc}]{\sphinxcrossref{\DUrole{doc}{Build
package}}}} page for more info.

\sphinxAtStartPar
Compute style \sphinxstyleemphasis{dipole/tip4p} can only be used with tip4p pair styles.


\subsection{Related commands}
\label{\detokenize{compute_dipole:related-commands}}
\sphinxAtStartPar
{\hyperref[\detokenize{compute_dipole_chunk::doc}]{\sphinxcrossref{\DUrole{doc}{compute dipole/chunk}}}}


\subsection{Default}
\label{\detokenize{compute_dipole:default}}
\sphinxAtStartPar
Using the center of mass is the default setting for the net charge correction.

\sphinxstepscope

\index{compute dipole/chunk@\spxentry{compute dipole/chunk}}\ignorespaces 
\index{compute dipole/tip4p/chunk@\spxentry{compute dipole/tip4p/chunk}}\ignorespaces 

\section{compute dipole/chunk command}
\label{\detokenize{compute_dipole_chunk:compute-dipole-chunk-command}}\label{\detokenize{compute_dipole_chunk:index-1}}\label{\detokenize{compute_dipole_chunk:index-0}}\label{\detokenize{compute_dipole_chunk::doc}}

\section{compute dipole/tip4p/chunk command}
\label{\detokenize{compute_dipole_chunk:compute-dipole-tip4p-chunk-command}}

\subsection{Syntax}
\label{\detokenize{compute_dipole_chunk:syntax}}
\begin{sphinxVerbatim}[commandchars=\\\{\}]
\PYG{k}{compute }\PYG{n+nv+nvIdentifier}{ID}\PYG{+w}{ }\PYG{n+nv+nvIdentifier}{group\PYGZhy{}ID}\PYG{+w}{ }\PYG{n}{style}\PYG{+w}{ }\PYG{n}{chunkID}\PYG{+w}{ }\PYG{n}{arg}
\end{sphinxVerbatim}
\begin{itemize}
\item {} 
\sphinxAtStartPar
ID, group\sphinxhyphen{}ID are documented in {\hyperref[\detokenize{compute::doc}]{\sphinxcrossref{\DUrole{doc}{compute}}}} command

\item {} 
\sphinxAtStartPar
style = \sphinxstyleemphasis{dipole/chunk} or \sphinxstyleemphasis{dipole/tip4p/chunk}

\item {} 
\sphinxAtStartPar
chunkID = ID of {\hyperref[\detokenize{compute_chunk_atom::doc}]{\sphinxcrossref{\DUrole{doc}{compute chunk/atom}}}} command

\item {} 
\sphinxAtStartPar
arg = \sphinxstyleemphasis{mass} or \sphinxstyleemphasis{geometry} = use COM or geometric center for charged chunk correction (optional)

\end{itemize}


\subsection{Examples}
\label{\detokenize{compute_dipole_chunk:examples}}
\begin{sphinxVerbatim}[commandchars=\\\{\}]
\PYG{k}{compute }\PYG{n+nv+nvIdentifier}{1}\PYG{+w}{ }\PYG{n+nv+nvIdentifier}{fluid}\PYG{+w}{ }\PYG{n}{dipole}\PYG{o}{/}\PYG{n}{chunk}\PYG{+w}{ }\PYG{n}{molchunk}
\PYG{k}{compute }\PYG{n+nv+nvIdentifier}{dw}\PYG{+w}{ }\PYG{n+nv+nvIdentifier}{water}\PYG{+w}{ }\PYG{n}{dipole}\PYG{o}{/}\PYG{n}{chunk}\PYG{+w}{ }\PYG{l+m}{1}\PYG{+w}{ }\PYG{n}{geometry}
\end{sphinxVerbatim}


\subsection{Description}
\label{\detokenize{compute_dipole_chunk:description}}
\sphinxAtStartPar
Define a computation that calculates the dipole vector and total dipole
for multiple chunks of atoms.

\sphinxAtStartPar
In LAMMPS, chunks are collections of atoms defined by a {\hyperref[\detokenize{compute_chunk_atom::doc}]{\sphinxcrossref{\DUrole{doc}{compute
chunk/atom}}}} command, which assigns each atom to a
single chunk (or no chunk).  The ID for this command is specified as
chunkID.  For example, a single chunk could be the atoms in a molecule
or atoms in a spatial bin.  See the {\hyperref[\detokenize{compute_chunk_atom::doc}]{\sphinxcrossref{\DUrole{doc}{compute chunk/atom}}}} and {\hyperref[\detokenize{Howto_chunk::doc}]{\sphinxcrossref{\DUrole{doc}{Howto chunk}}}} doc pages for
details of how chunks can be defined and examples of how they can be
used to measure properties of a system.

\sphinxAtStartPar
These computes calculate the \((x,y,z)\) coordinates of the dipole
vector and the total dipole moment for each chunk, which includes all
effects due to atoms passing through periodic boundaries.  For chunks
with a net charge the resulting dipole is made position independent by
subtracting the position vector of the center of mass or geometric
center times the net charge from the computed dipole vector.  Both
per\sphinxhyphen{}atom charges and per\sphinxhyphen{}atom dipole moments, if present, contribute to
the computed dipole.

\sphinxAtStartPar
\DUrole{versionmodified,added}{Added in version 28Mar2023.}

\sphinxAtStartPar
Compute \sphinxstyleemphasis{dipole/tip4p/chunk} includes adjustments for the charge
carrying point M in molecules with TIP4P water geometry.  The
corresponding parameters are extracted from the pair style.

\sphinxAtStartPar
Note that only atoms in the specified group contribute to the
calculation.  The {\hyperref[\detokenize{compute_chunk_atom::doc}]{\sphinxcrossref{\DUrole{doc}{compute chunk/atom}}}} command
defines its own group; atoms will have a chunk ID = 0 if they are not in
that group, signifying they are not assigned to a chunk, and will thus
also not contribute to this calculation.  You can specify the “all”
group for this command if you simply want to include atoms with non\sphinxhyphen{}zero
chunk IDs.

\begin{sphinxadmonition}{note}{Note:}
\sphinxAtStartPar
The coordinates of an atom contribute to the chunk’s dipole in
“unwrapped” form, by using the image flags associated with each atom.
See the {\hyperref[\detokenize{dump::doc}]{\sphinxcrossref{\DUrole{doc}{dump custom}}}} command for a discussion of
“unwrapped” coordinates.  See the Atoms section of the
{\hyperref[\detokenize{read_data::doc}]{\sphinxcrossref{\DUrole{doc}{read\_data}}}} command for a discussion of image flags
and how they are set for each atom.  You can reset the image flags
(e.g., to 0) before invoking this compute by using the {\hyperref[\detokenize{set::doc}]{\sphinxcrossref{\DUrole{doc}{set image}}}} command.
\end{sphinxadmonition}

\sphinxAtStartPar
The simplest way to output the results of the compute com/chunk
calculation to a file is to use the {\hyperref[\detokenize{fix_ave_time::doc}]{\sphinxcrossref{\DUrole{doc}{fix ave/time}}}}
command, for example:

\begin{sphinxVerbatim}[commandchars=\\\{\}]
\PYG{k}{compute }\PYG{n+nv+nvIdentifier}{cc1}\PYG{+w}{ }\PYG{n+nv+nvIdentifier}{all}\PYG{+w}{ }\PYG{n}{chunk}\PYG{o}{/}\PYG{n}{atom}\PYG{+w}{ }\PYG{n}{molecule}
\PYG{k}{compute }\PYG{n+nv+nvIdentifier}{myChunk}\PYG{+w}{ }\PYG{n+nv+nvIdentifier}{all}\PYG{+w}{ }\PYG{n}{dipole}\PYG{o}{/}\PYG{n}{chunk}\PYG{+w}{ }\PYG{n}{cc1}
\PYG{k}{fix }\PYG{n+nv+nvIdentifier}{1}\PYG{+w}{ }\PYG{n+nv+nvIdentifier}{all}\PYG{+w}{ }\PYG{n}{ave}\PYG{o}{/}\PYG{n}{time}\PYG{+w}{ }\PYG{l+m}{100}\PYG{+w}{ }\PYG{l+m}{1}\PYG{+w}{ }\PYG{l+m}{100}\PYG{+w}{ }\PYG{n}{c\PYGZus{}myChunk[}\PYG{o}{*}\PYG{n}{]}\PYG{+w}{ }\PYG{n}{file}\PYG{+w}{ }\PYG{n}{tmp.out}\PYG{+w}{ }\PYG{n}{mode}\PYG{+w}{ }\PYG{n}{vector}
\end{sphinxVerbatim}


\subsection{Output info}
\label{\detokenize{compute_dipole_chunk:output-info}}
\sphinxAtStartPar
These computes calculate a global array where the number of rows = the
number of chunks \sphinxstyleemphasis{Nchunk} as calculated by the specified {\hyperref[\detokenize{compute_chunk_atom::doc}]{\sphinxcrossref{\DUrole{doc}{compute
chunk/atom}}}} command.  The number of columns is 4
for the \((x,y,z)\) dipole vector components and the total dipole of
each chunk.  These values can be accessed by any command that uses
global array values from a compute as input.  See the {\hyperref[\detokenize{Howto_output::doc}]{\sphinxcrossref{\DUrole{doc}{Howto output}}}} page for an overview of LAMMPS output options.

\sphinxAtStartPar
The array values are “intensive”.  The array values will be in
dipole units (i.e., charge {\hyperref[\detokenize{units::doc}]{\sphinxcrossref{\DUrole{doc}{units}}}} times distance
{\hyperref[\detokenize{units::doc}]{\sphinxcrossref{\DUrole{doc}{units}}}}).


\subsection{Restrictions}
\label{\detokenize{compute_dipole_chunk:restrictions}}
\sphinxAtStartPar
Compute style \sphinxstyleemphasis{dipole/tip4p/chunk} is part of the EXTRA\sphinxhyphen{}COMPUTE
package. It is only enabled if LAMMPS was built with that package.  See
the {\hyperref[\detokenize{Build_package::doc}]{\sphinxcrossref{\DUrole{doc}{Build package}}}} page for more info.

\sphinxAtStartPar
Compute style \sphinxstyleemphasis{dipole/tip4p/chunk} can only be used with tip4p pair
styles.


\subsection{Related commands}
\label{\detokenize{compute_dipole_chunk:related-commands}}
\sphinxAtStartPar
{\hyperref[\detokenize{compute_com_chunk::doc}]{\sphinxcrossref{\DUrole{doc}{compute com/chunk}}}},
{\hyperref[\detokenize{compute_dipole::doc}]{\sphinxcrossref{\DUrole{doc}{compute dipole}}}}


\subsection{Default}
\label{\detokenize{compute_dipole_chunk:default}}
\sphinxAtStartPar
Using the center of mass is the default setting for the net charge correction.

\sphinxstepscope

\index{compute displace/atom@\spxentry{compute displace/atom}}\ignorespaces 

\section{compute displace/atom command}
\label{\detokenize{compute_displace_atom:compute-displace-atom-command}}\label{\detokenize{compute_displace_atom:index-0}}\label{\detokenize{compute_displace_atom::doc}}

\subsection{Syntax}
\label{\detokenize{compute_displace_atom:syntax}}
\begin{sphinxVerbatim}[commandchars=\\\{\}]
\PYG{k}{compute }\PYG{n+nv+nvIdentifier}{ID}\PYG{+w}{ }\PYG{n+nv+nvIdentifier}{group\PYGZhy{}ID}\PYG{+w}{ }\PYG{n}{displace}\PYG{o}{/}\PYG{n}{atom}
\end{sphinxVerbatim}
\begin{itemize}
\item {} 
\sphinxAtStartPar
ID, group\sphinxhyphen{}ID are documented in {\hyperref[\detokenize{compute::doc}]{\sphinxcrossref{\DUrole{doc}{compute}}}} command

\item {} 
\sphinxAtStartPar
displace/atom = style name of this compute command

\item {} 
\sphinxAtStartPar
zero or more keyword/arg pairs may be appended

\item {} 
\sphinxAtStartPar
keyword = \sphinxstyleemphasis{refresh}
\begin{sphinxalltt}
\sphinxstyleemphasis{refresh} arg = name of per\sphinxhyphen{}atom variable
\end{sphinxalltt}

\end{itemize}


\subsection{Examples}
\label{\detokenize{compute_displace_atom:examples}}
\begin{sphinxVerbatim}[commandchars=\\\{\}]
\PYG{k}{compute }\PYG{n+nv+nvIdentifier}{1}\PYG{+w}{ }\PYG{n+nv+nvIdentifier}{all}\PYG{+w}{ }\PYG{n}{displace}\PYG{o}{/}\PYG{n}{atom}
\PYG{k}{compute }\PYG{n+nv+nvIdentifier}{1}\PYG{+w}{ }\PYG{n+nv+nvIdentifier}{all}\PYG{+w}{ }\PYG{n}{displace}\PYG{o}{/}\PYG{n}{atom}\PYG{+w}{ }\PYG{n}{refresh}\PYG{+w}{ }\PYG{n}{myVar}
\end{sphinxVerbatim}


\subsection{Description}
\label{\detokenize{compute_displace_atom:description}}
\sphinxAtStartPar
Define a computation that calculates the current displacement of each
atom in the group from its original (reference) coordinates, including
all effects due to atoms passing through periodic boundaries.

\sphinxAtStartPar
A vector of four quantities per atom is calculated by this compute.
The first three elements of the vector are the \((dx,dy,dz)\)
displacements.  The fourth component is the total displacement
(i.e., \(\sqrt{dx^2 + dy^2 + dz^2}\)).

\sphinxAtStartPar
The displacement of an atom is from its original position at the time
the compute command was issued.  The value of the displacement will be
0.0 for atoms not in the specified compute group.

\begin{sphinxadmonition}{note}{Note:}
\sphinxAtStartPar
Initial coordinates are stored in “unwrapped” form, by using the
image flags associated with each atom.  See the {\hyperref[\detokenize{dump::doc}]{\sphinxcrossref{\DUrole{doc}{dump custom}}}} command for a discussion of “unwrapped” coordinates.  See
the Atoms section of the {\hyperref[\detokenize{read_data::doc}]{\sphinxcrossref{\DUrole{doc}{read\_data}}}} command for a
discussion of image flags and how they are set for each atom.  You
can reset the image flags (e.g., to 0) before invoking this compute
by using the {\hyperref[\detokenize{set::doc}]{\sphinxcrossref{\DUrole{doc}{set image}}}} command.
\end{sphinxadmonition}

\begin{sphinxadmonition}{note}{Note:}
\sphinxAtStartPar
If you want the quantities calculated by this compute to be
continuous when running from a {\hyperref[\detokenize{read_restart::doc}]{\sphinxcrossref{\DUrole{doc}{restart file}}}}, then
you should use the same ID for this compute, as in the original run.
This is so that the fix this compute creates to store per\sphinxhyphen{}atom
quantities will also have the same ID, and thus be initialized
correctly with time = 0 atom coordinates from the restart file.
\end{sphinxadmonition}


\bigskip\hrule\bigskip


\sphinxAtStartPar
The \sphinxstyleemphasis{refresh} option can be used in conjunction with the “dump\_modify
refresh” command to generate incremental dump files.

\sphinxAtStartPar
The definition and motivation of an incremental dump file is as
follows.  Instead of outputting all atoms at each snapshot (with some
associated values), you may only wish to output the subset of atoms
with a value that has changed in some way compared to the value the
last time that atom was output.  In some scenarios this can result in
a dramatically smaller dump file.  If desired, by post\sphinxhyphen{}processing the
sequence of snapshots, the values for all atoms at all timesteps can
be inferred.

\sphinxAtStartPar
A concrete example using this compute, is a simulation of atom
diffusion in a solid, represented as atoms on a lattice.  Diffusive
hops are rare.  Imagine that when a hop occurs an atom moves more than
a distance \sphinxstyleemphasis{Dhop}.  For any snapshot we only want to output atoms that
have hopped since the last snapshot.  This can be accomplished with
something like the following commands:

\begin{sphinxVerbatim}[commandchars=\\\{\}]
\PYG{k}{write\PYGZus{}dump      }\PYG{n+nv+nvIdentifier}{all}\PYG{+w}{ }\PYG{n}{custom}\PYG{+w}{ }\PYG{n}{tmp.dump}\PYG{+w}{ }\PYG{n}{id}\PYG{+w}{ }\PYG{n}{type}\PYG{+w}{ }\PYG{n}{x}\PYG{+w}{ }\PYG{n}{y}\PYG{+w}{ }\PYG{n}{z}\PYG{+w}{    }\PYG{c}{\PYGZsh{} see comment below}

\PYG{k}{variable        }\PYG{n+nv+nvIdentifier}{Dhop}\PYG{+w}{ }\PYG{n}{equal}\PYG{+w}{ }\PYG{l+m}{0.6}
\PYG{k}{variable        }\PYG{n+nv+nvIdentifier}{check}\PYG{+w}{ }\PYG{n}{atom}\PYG{+w}{ }\PYG{l+s}{\PYGZdq{}}\PYG{l+s}{c\PYGZus{}dsp[4] \PYGZgt{} v\PYGZus{}Dhop}\PYG{l+s}{\PYGZdq{}}
\PYG{k}{compute         }\PYG{n+nv+nvIdentifier}{dsp}\PYG{+w}{ }\PYG{n+nv+nvIdentifier}{all}\PYG{+w}{ }\PYG{n}{displace}\PYG{o}{/}\PYG{n}{atom}\PYG{+w}{ }\PYG{n}{refresh}\PYG{+w}{ }\PYG{n}{check}
\PYG{k}{dump            }\PYG{n+nv+nvIdentifier}{1}\PYG{+w}{ }\PYG{n+nv+nvIdentifier}{all}\PYG{+w}{ }\PYG{n}{custom}\PYG{+w}{ }\PYG{l+m}{100}\PYG{+w}{ }\PYG{n}{tmp.dump}\PYG{+w}{ }\PYG{n}{id}\PYG{+w}{ }\PYG{n}{type}\PYG{+w}{ }\PYG{n}{x}\PYG{+w}{ }\PYG{n}{y}\PYG{+w}{ }\PYG{n}{z}
\PYG{k}{dump\PYGZus{}modify     }\PYG{n+nv+nvIdentifier}{1}\PYG{+w}{ }\PYG{n}{append}\PYG{+w}{ }\PYG{n}{yes}\PYG{+w}{ }\PYG{n}{thresh}\PYG{+w}{ }\PYG{n}{c\PYGZus{}dsp[4]}\PYG{+w}{ }\PYG{o}{\PYGZgt{}}\PYG{+w}{ }\PYG{n+nv}{\PYGZdl{}\PYGZob{}}\PYG{n+nv}{Dhop}\PYG{n+nv}{\PYGZcb{}}\PYG{+w}{ }\PYG{o}{\PYGZam{}}
\PYG{+w}{                }\PYG{n}{refresh}\PYG{+w}{ }\PYG{n}{c\PYGZus{}dsp}\PYG{+w}{ }\PYG{n}{delay}\PYG{+w}{ }\PYG{l+m}{100}
\end{sphinxVerbatim}

\sphinxAtStartPar
The {\hyperref[\detokenize{dump_modify::doc}]{\sphinxcrossref{\DUrole{doc}{dump\_modify thresh}}}} command will only output
atoms that have displaced more than \(0.6~\AA\) on each
snapshot (assuming metal units).  The dump\_modify \sphinxstyleemphasis{refresh} option triggers a
call to this compute at the end of every dump.

\sphinxAtStartPar
The \sphinxstyleemphasis{refresh} argument for this compute is the ID of an
{\hyperref[\detokenize{variable::doc}]{\sphinxcrossref{\DUrole{doc}{atom\sphinxhyphen{}style variable}}}} which calculates a Boolean value (0 or 1)
based on the same criterion used by dump\_modify thresh.  This compute
evaluates the atom\sphinxhyphen{}style variable.  For each atom that returns 1 (true),
the original (reference) coordinates of the atom (stored by
this compute) are updated.

\sphinxAtStartPar
The effect of these commands is that a particular atom will only be
output in the dump file on the snapshot after it makes a diffusive
hop.  It will not be output again until it makes another hop.

\sphinxAtStartPar
Note that in the first snapshot of a subsequent run, no atoms will be
typically be output.  That is because the initial displacement for all
atoms is 0.0.  If an initial dump snapshot is desired, containing the
initial reference positions of all atoms, one way to do this is
illustrated above.  An initial write\_dump command can be used before
the first run.  It will contain the positions of all the atoms,
Options in the {\hyperref[\detokenize{dump_modify::doc}]{\sphinxcrossref{\DUrole{doc}{dump\_modify}}}} command above will
append new output to that same file and delay the output until a later
timestep.  The \sphinxstyleemphasis{delay} setting avoids a second time = 0 snapshot which
would be empty.


\bigskip\hrule\bigskip



\subsection{Output info}
\label{\detokenize{compute_displace_atom:output-info}}
\sphinxAtStartPar
This compute calculates a per\sphinxhyphen{}atom array with four columns, which can be
accessed by indices 1\textendash{}4 by any command that uses per\sphinxhyphen{}atom values from
a compute as input.  See the {\hyperref[\detokenize{Howto_output::doc}]{\sphinxcrossref{\DUrole{doc}{Howto output}}}} doc page
for an overview of LAMMPS output options.

\sphinxAtStartPar
The per\sphinxhyphen{}atom array values will be in distance {\hyperref[\detokenize{units::doc}]{\sphinxcrossref{\DUrole{doc}{units}}}}.

\sphinxAtStartPar
This compute supports the \sphinxstyleemphasis{refresh} option as explained above, for use
in conjunction with {\hyperref[\detokenize{dump_modify::doc}]{\sphinxcrossref{\DUrole{doc}{dump\_modify refresh}}}} to generate
incremental dump files.


\subsection{Restrictions}
\label{\detokenize{compute_displace_atom:restrictions}}\begin{quote}

\sphinxAtStartPar
none
\end{quote}


\subsection{Related commands}
\label{\detokenize{compute_displace_atom:related-commands}}
\sphinxAtStartPar
{\hyperref[\detokenize{compute_msd::doc}]{\sphinxcrossref{\DUrole{doc}{compute msd}}}}, {\hyperref[\detokenize{dump::doc}]{\sphinxcrossref{\DUrole{doc}{dump custom}}}}, {\hyperref[\detokenize{fix_store_state::doc}]{\sphinxcrossref{\DUrole{doc}{fix store/state}}}}


\subsection{Default}
\label{\detokenize{compute_displace_atom:default}}
\sphinxAtStartPar
none

\sphinxstepscope

\index{compute dpd@\spxentry{compute dpd}}\ignorespaces 

\section{compute dpd command}
\label{\detokenize{compute_dpd:compute-dpd-command}}\label{\detokenize{compute_dpd:index-0}}\label{\detokenize{compute_dpd::doc}}

\subsection{Syntax}
\label{\detokenize{compute_dpd:syntax}}
\begin{sphinxVerbatim}[commandchars=\\\{\}]
\PYG{k}{compute }\PYG{n+nv+nvIdentifier}{ID}\PYG{+w}{ }\PYG{n+nv+nvIdentifier}{group\PYGZhy{}ID}\PYG{+w}{ }\PYG{n}{dpd}
\end{sphinxVerbatim}
\begin{itemize}
\item {} 
\sphinxAtStartPar
ID, group\sphinxhyphen{}ID are documented in {\hyperref[\detokenize{compute::doc}]{\sphinxcrossref{\DUrole{doc}{compute}}}} command

\item {} 
\sphinxAtStartPar
dpd = style name of this compute command

\end{itemize}


\subsection{Examples}
\label{\detokenize{compute_dpd:examples}}
\begin{sphinxVerbatim}[commandchars=\\\{\}]
\PYG{k}{compute }\PYG{n+nv+nvIdentifier}{1}\PYG{+w}{ }\PYG{n+nv+nvIdentifier}{all}\PYG{+w}{ }\PYG{n}{dpd}
\end{sphinxVerbatim}


\subsection{Description}
\label{\detokenize{compute_dpd:description}}
\sphinxAtStartPar
Define a computation that accumulates the total internal conductive
energy (\(U^{\text{cond}}\)), the total internal mechanical energy
(\(U^{\text{mech}}\)), the total chemical energy (\(U^\text{chem}\))
and the \sphinxstyleemphasis{harmonic} average of the internal temperature (\(\theta_\text{avg}\))
for the entire system of particles.  See the
{\hyperref[\detokenize{compute_dpd_atom::doc}]{\sphinxcrossref{\DUrole{doc}{compute dpd/atom}}}} command if you want
per\sphinxhyphen{}particle internal energies and internal temperatures.

\sphinxAtStartPar
The system internal properties are computed according to the following
relations:
\begin{equation*}
\begin{split}U^\text{cond} = & \sum_{i=1}^{N} u_{i}^\text{cond} \\
U^\text{mech} = & \sum_{i=1}^{N} u_{i}^\text{mech} \\
U^\text{chem} = & \sum_{i=1}^{N} u_{i}^\text{chem} \\
            U = & \sum_{i=1}^{N} (u_{i}^\text{cond}
                  + u_{i}^\text{mech} + u_{i}^\text{chem}) \\
\theta_{avg} = & \biggl(\frac{1}{N}\sum_{i=1}^{N}
                       \frac{1}{\theta_{i}}\biggr)^{-1} \\\end{split}
\end{equation*}
\sphinxAtStartPar
where \(N\) is the number of particles in the system.


\bigskip\hrule\bigskip



\subsection{Output info}
\label{\detokenize{compute_dpd:output-info}}
\sphinxAtStartPar
This compute calculates a global vector of length 5 (\(U^\text{cond}\),
\(U^\text{mech}\), \(U^\text{chem}\), \(\theta_\text{avg}\),
\(N\)), which can be accessed by indices 1 through 5.
See the {\hyperref[\detokenize{Howto_output::doc}]{\sphinxcrossref{\DUrole{doc}{Howto output}}}} page for an overview of
LAMMPS output options.

\sphinxAtStartPar
The vector values will be in energy and temperature {\hyperref[\detokenize{units::doc}]{\sphinxcrossref{\DUrole{doc}{units}}}}.


\subsection{Restrictions}
\label{\detokenize{compute_dpd:restrictions}}
\sphinxAtStartPar
This command is part of the DPD\sphinxhyphen{}REACT package.  It is only enabled if
LAMMPS was built with that package.
See the {\hyperref[\detokenize{Build_package::doc}]{\sphinxcrossref{\DUrole{doc}{Build package}}}} page for more info.

\sphinxAtStartPar
This command also requires use of the {\hyperref[\detokenize{atom_style::doc}]{\sphinxcrossref{\DUrole{doc}{atom\_style dpd}}}}
command.


\subsection{Related commands}
\label{\detokenize{compute_dpd:related-commands}}
\sphinxAtStartPar
{\hyperref[\detokenize{compute_dpd_atom::doc}]{\sphinxcrossref{\DUrole{doc}{compute dpd/atom}}}},
{\hyperref[\detokenize{thermo_style::doc}]{\sphinxcrossref{\DUrole{doc}{thermo\_style}}}}


\subsection{Default}
\label{\detokenize{compute_dpd:default}}
\sphinxAtStartPar
none


\bigskip\hrule\bigskip

\phantomsection\label{\detokenize{compute_dpd:larentzos1}}
\sphinxAtStartPar
\sphinxstylestrong{(Larentzos)} J.P. Larentzos, J.K. Brennan, J.D. Moore, and
W.D. Mattson, “LAMMPS Implementation of Constant Energy Dissipative
Particle Dynamics (DPD\sphinxhyphen{}E)”, ARL\sphinxhyphen{}TR\sphinxhyphen{}6863, U.S. Army Research
Laboratory, Aberdeen Proving Ground, MD (2014).

\sphinxstepscope

\index{compute dpd/atom@\spxentry{compute dpd/atom}}\ignorespaces 

\section{compute dpd/atom command}
\label{\detokenize{compute_dpd_atom:compute-dpd-atom-command}}\label{\detokenize{compute_dpd_atom:index-0}}\label{\detokenize{compute_dpd_atom::doc}}

\subsection{Syntax}
\label{\detokenize{compute_dpd_atom:syntax}}
\begin{sphinxVerbatim}[commandchars=\\\{\}]
\PYG{k}{compute }\PYG{n+nv+nvIdentifier}{ID}\PYG{+w}{ }\PYG{n+nv+nvIdentifier}{group\PYGZhy{}ID}\PYG{+w}{ }\PYG{n}{dpd}\PYG{o}{/}\PYG{n}{atom}
\end{sphinxVerbatim}
\begin{itemize}
\item {} 
\sphinxAtStartPar
ID, group\sphinxhyphen{}ID are documented in {\hyperref[\detokenize{compute::doc}]{\sphinxcrossref{\DUrole{doc}{compute}}}} command

\item {} 
\sphinxAtStartPar
dpd/atom = style name of this compute command

\end{itemize}


\subsection{Examples}
\label{\detokenize{compute_dpd_atom:examples}}
\begin{sphinxVerbatim}[commandchars=\\\{\}]
\PYG{k}{compute }\PYG{n+nv+nvIdentifier}{1}\PYG{+w}{ }\PYG{n+nv+nvIdentifier}{all}\PYG{+w}{ }\PYG{n}{dpd}\PYG{o}{/}\PYG{n}{atom}
\end{sphinxVerbatim}


\subsection{Description}
\label{\detokenize{compute_dpd_atom:description}}
\sphinxAtStartPar
Define a computation that accesses the per\sphinxhyphen{}particle internal conductive energy
(\(u^\text{cond}\)), internal mechanical energy (\(u^\text{mech}\)),
internal chemical energy (\(u^\text{chem}\)) and internal temperatures
(\(\theta\)) for each particle in a group.
See the {\hyperref[\detokenize{compute_dpd::doc}]{\sphinxcrossref{\DUrole{doc}{compute dpd}}}} command if you want the total
internal conductive energy, the total internal mechanical energy, the
total chemical energy and average internal temperature of the entire system or
group of dpd particles.


\subsection{Output info}
\label{\detokenize{compute_dpd_atom:output-info}}
\sphinxAtStartPar
This compute calculates a per\sphinxhyphen{}particle array with four columns
(\(u^\text{cond}\), \(u^\text{mech}\), \(u^\text{chem}\),
\(\theta\)), which can be accessed by indices 1\textendash{}4 by any
command that uses per\sphinxhyphen{}particle values from a compute as input.  See
the {\hyperref[\detokenize{Howto_output::doc}]{\sphinxcrossref{\DUrole{doc}{Howto output}}}} page for an overview of
LAMMPS output options.

\sphinxAtStartPar
The per\sphinxhyphen{}particle array values will be in energy (\(u^\text{cond}\),
\(u^\text{mech}\), \(u^\text{chem}\))
and temperature (\(\theta\)) {\hyperref[\detokenize{units::doc}]{\sphinxcrossref{\DUrole{doc}{units}}}}.


\subsection{Restrictions}
\label{\detokenize{compute_dpd_atom:restrictions}}
\sphinxAtStartPar
This command is part of the DPD\sphinxhyphen{}REACT package.  It is only enabled if
LAMMPS was built with that package.  See the {\hyperref[\detokenize{Build_package::doc}]{\sphinxcrossref{\DUrole{doc}{Build package}}}} page for more info.

\sphinxAtStartPar
This command also requires use of the {\hyperref[\detokenize{atom_style::doc}]{\sphinxcrossref{\DUrole{doc}{atom\_style dpd}}}}
command.


\subsection{Related commands}
\label{\detokenize{compute_dpd_atom:related-commands}}
\sphinxAtStartPar
{\hyperref[\detokenize{dump::doc}]{\sphinxcrossref{\DUrole{doc}{dump custom}}}}, {\hyperref[\detokenize{compute_dpd::doc}]{\sphinxcrossref{\DUrole{doc}{compute dpd}}}}


\subsection{Default}
\label{\detokenize{compute_dpd_atom:default}}
\sphinxAtStartPar
none


\bigskip\hrule\bigskip

\phantomsection\label{\detokenize{compute_dpd_atom:larentzos2}}
\sphinxAtStartPar
\sphinxstylestrong{(Larentzos)} J.P. Larentzos, J.K. Brennan, J.D. Moore, and
W.D. Mattson, “LAMMPS Implementation of Constant Energy Dissipative
Particle Dynamics (DPD\sphinxhyphen{}E)”, ARL\sphinxhyphen{}TR\sphinxhyphen{}6863, U.S. Army Research
Laboratory, Aberdeen Proving Ground, MD (2014).

\sphinxstepscope

\index{compute edpd/temp/atom@\spxentry{compute edpd/temp/atom}}\ignorespaces 

\section{compute edpd/temp/atom command}
\label{\detokenize{compute_edpd_temp_atom:compute-edpd-temp-atom-command}}\label{\detokenize{compute_edpd_temp_atom:index-0}}\label{\detokenize{compute_edpd_temp_atom::doc}}

\subsection{Syntax}
\label{\detokenize{compute_edpd_temp_atom:syntax}}
\begin{sphinxVerbatim}[commandchars=\\\{\}]
\PYG{k}{compute }\PYG{n+nv+nvIdentifier}{ID}\PYG{+w}{ }\PYG{n+nv+nvIdentifier}{group\PYGZhy{}ID}\PYG{+w}{ }\PYG{n}{edpd}\PYG{o}{/}\PYG{n}{temp}\PYG{o}{/}\PYG{n}{atom}
\end{sphinxVerbatim}
\begin{itemize}
\item {} 
\sphinxAtStartPar
ID, group\sphinxhyphen{}ID are documented in {\hyperref[\detokenize{compute::doc}]{\sphinxcrossref{\DUrole{doc}{compute}}}} command

\item {} 
\sphinxAtStartPar
edpd/temp/atom = style name of this compute command

\end{itemize}


\subsection{Examples}
\label{\detokenize{compute_edpd_temp_atom:examples}}
\begin{sphinxVerbatim}[commandchars=\\\{\}]
\PYG{k}{compute }\PYG{n+nv+nvIdentifier}{1}\PYG{+w}{ }\PYG{n+nv+nvIdentifier}{all}\PYG{+w}{ }\PYG{n}{edpd}\PYG{o}{/}\PYG{n}{temp}\PYG{o}{/}\PYG{n}{atom}
\end{sphinxVerbatim}


\subsection{Description}
\label{\detokenize{compute_edpd_temp_atom:description}}
\sphinxAtStartPar
Define a computation that calculates the per\sphinxhyphen{}atom temperature
for each eDPD particle in a group.

\sphinxAtStartPar
The temperature is a local temperature derived from the internal energy
of each eDPD particle based on the local equilibrium hypothesis.
For more details please see {\hyperref[\detokenize{compute_edpd_temp_atom:espanol1997}]{\sphinxcrossref{\DUrole{std,std-ref}{(Espanol1997)}}}} and
{\hyperref[\detokenize{compute_edpd_temp_atom:li2014a}]{\sphinxcrossref{\DUrole{std,std-ref}{(Li2014)}}}}.


\subsection{Output info}
\label{\detokenize{compute_edpd_temp_atom:output-info}}
\sphinxAtStartPar
This compute calculates a per\sphinxhyphen{}atom vector, which can be accessed by
any command that uses per\sphinxhyphen{}atom values from a compute as input. See the
{\hyperref[\detokenize{Howto_output::doc}]{\sphinxcrossref{\DUrole{doc}{Howto output}}}} page for an overview of LAMMPS
output options.

\sphinxAtStartPar
The per\sphinxhyphen{}atom vector values will be in temperature {\hyperref[\detokenize{units::doc}]{\sphinxcrossref{\DUrole{doc}{units}}}}.


\subsection{Restrictions}
\label{\detokenize{compute_edpd_temp_atom:restrictions}}
\sphinxAtStartPar
This compute is part of the DPD\sphinxhyphen{}MESO package.  It is only enabled if
LAMMPS was built with that package.  See the {\hyperref[\detokenize{Build_package::doc}]{\sphinxcrossref{\DUrole{doc}{Build package}}}} page for more info.


\subsection{Related commands}
\label{\detokenize{compute_edpd_temp_atom:related-commands}}
\sphinxAtStartPar
{\hyperref[\detokenize{pair_mesodpd::doc}]{\sphinxcrossref{\DUrole{doc}{pair\_style edpd}}}}


\subsection{Default}
\label{\detokenize{compute_edpd_temp_atom:default}}
\sphinxAtStartPar
none


\bigskip\hrule\bigskip

\phantomsection\label{\detokenize{compute_edpd_temp_atom:espanol1997}}
\sphinxAtStartPar
\sphinxstylestrong{(Espanol1997)} Espanol, Europhys Lett, 40(6): 631\sphinxhyphen{}636 (1997).  DOI:
10.1209/epl/i1997\sphinxhyphen{}00515\sphinxhyphen{}8

\phantomsection\label{\detokenize{compute_edpd_temp_atom:li2014a}}
\sphinxAtStartPar
\sphinxstylestrong{(Li2014)} Li, Tang, Lei, Caswell, Karniadakis, J Comput Phys, 265:
113\sphinxhyphen{}127 (2014).  DOI: 10.1016/j.jcp.2014.02.003.

\sphinxstepscope

\index{compute efield/atom@\spxentry{compute efield/atom}}\ignorespaces 

\section{compute efield/atom command}
\label{\detokenize{compute_efield_atom:compute-efield-atom-command}}\label{\detokenize{compute_efield_atom:index-0}}\label{\detokenize{compute_efield_atom::doc}}

\subsection{Syntax}
\label{\detokenize{compute_efield_atom:syntax}}
\begin{sphinxVerbatim}[commandchars=\\\{\}]
\PYG{k}{compute }\PYG{n+nv+nvIdentifier}{ID}\PYG{+w}{ }\PYG{n+nv+nvIdentifier}{group\PYGZhy{}ID}\PYG{+w}{ }\PYG{n}{efield}\PYG{o}{/}\PYG{n}{atom}\PYG{+w}{ }\PYG{n}{keyword}\PYG{+w}{ }\PYG{n}{val}
\end{sphinxVerbatim}
\begin{itemize}
\item {} 
\sphinxAtStartPar
ID, group\sphinxhyphen{}ID are documented in {\hyperref[\detokenize{compute::doc}]{\sphinxcrossref{\DUrole{doc}{compute}}}} command

\item {} 
\sphinxAtStartPar
efield/atom = style name of this compute command

\item {} 
\sphinxAtStartPar
zero or more keyword/value pairs may be appended

\item {} 
\sphinxAtStartPar
keyword = \sphinxstyleemphasis{pair} or \sphinxstyleemphasis{kspace}
\begin{sphinxalltt}
\sphinxstyleemphasis{pair} args = \sphinxstyleemphasis{yes} or \sphinxstyleemphasis{no}
\sphinxstyleemphasis{kspace} args = \sphinxstyleemphasis{yes} or \sphinxstyleemphasis{no}
\end{sphinxalltt}

\end{itemize}


\subsection{Examples}
\label{\detokenize{compute_efield_atom:examples}}
\begin{sphinxVerbatim}[commandchars=\\\{\}]
\PYG{k}{compute }\PYG{n+nv+nvIdentifier}{1}\PYG{+w}{ }\PYG{n+nv+nvIdentifier}{all}\PYG{+w}{ }\PYG{n}{efield}\PYG{o}{/}\PYG{n}{atom}
\PYG{k}{compute }\PYG{n+nv+nvIdentifier}{1}\PYG{+w}{ }\PYG{n+nv+nvIdentifier}{all}\PYG{+w}{ }\PYG{n}{efield}\PYG{o}{/}\PYG{n}{atom}\PYG{+w}{ }\PYG{n}{pair}\PYG{+w}{ }\PYG{n}{yes}\PYG{+w}{ }\PYG{n}{kspace}\PYG{+w}{ }\PYG{n}{no}
\end{sphinxVerbatim}

\sphinxAtStartPar
Used in input scripts:

\begin{sphinxVerbatim}[commandchars=\\\{\}]
examples/PACKAGES/dielectric/in.confined
examples/PACKAGES/dielectric/in.nopbc
\end{sphinxVerbatim}


\subsection{Description}
\label{\detokenize{compute_efield_atom:description}}
\sphinxAtStartPar
Define a computation that calculates the electric field at each atom in a group.
The compute should only enabled with pair and kspace styles that are provided
by the DIELECTRIC package because only these styles compute the per\sphinxhyphen{}atom
electric field at every time step.

\sphinxAtStartPar
The electric field is a 3\sphinxhyphen{}component vector.  The value of the electric field
components will be 0.0 for atoms not in the specified compute group.


\bigskip\hrule\bigskip


\sphinxAtStartPar
The keyword/value option pairs are used in the following ways.

\sphinxAtStartPar
For the \sphinxstyleemphasis{pair} and \sphinxstyleemphasis{kspace} keywords, the real\sphinxhyphen{}space and reciprocal\sphinxhyphen{}space
contributions to the electric field can be turned off and on.


\subsection{Output info}
\label{\detokenize{compute_efield_atom:output-info}}
\sphinxAtStartPar
This compute calculates a per\sphinxhyphen{}atom vector, which can be accessed by
any command that uses per\sphinxhyphen{}atom values from a compute as input.  See
the {\hyperref[\detokenize{Howto_output::doc}]{\sphinxcrossref{\DUrole{doc}{Howto output}}}} page for an overview of
LAMMPS output options.

\sphinxAtStartPar
The per\sphinxhyphen{}atom vector values will be in electric field {\hyperref[\detokenize{units::doc}]{\sphinxcrossref{\DUrole{doc}{units}}}}.


\subsection{Restrictions}
\label{\detokenize{compute_efield_atom:restrictions}}
\sphinxAtStartPar
This compute is part of the DIELECTRIC package. It is only enabled if
LAMMPS was built with that package.


\subsection{Related commands}
\label{\detokenize{compute_efield_atom:related-commands}}
\sphinxAtStartPar
{\hyperref[\detokenize{dump::doc}]{\sphinxcrossref{\DUrole{doc}{dump custom}}}}


\subsection{Default}
\label{\detokenize{compute_efield_atom:default}}
\sphinxAtStartPar
The option defaults are pair = yes and kspace = yes.

\sphinxstepscope

\index{compute efield/wolf/atom@\spxentry{compute efield/wolf/atom}}\ignorespaces 

\section{compute efield/wolf/atom command}
\label{\detokenize{compute_efield_wolf_atom:compute-efield-wolf-atom-command}}\label{\detokenize{compute_efield_wolf_atom:index-0}}\label{\detokenize{compute_efield_wolf_atom::doc}}

\subsection{Syntax}
\label{\detokenize{compute_efield_wolf_atom:syntax}}
\begin{sphinxVerbatim}[commandchars=\\\{\}]
\PYG{k}{compute }\PYG{n+nv+nvIdentifier}{ID}\PYG{+w}{ }\PYG{n+nv+nvIdentifier}{group\PYGZhy{}ID}\PYG{+w}{ }\PYG{n}{efield}\PYG{o}{/}\PYG{n}{wolf}\PYG{o}{/}\PYG{n}{atom}\PYG{+w}{ }\PYG{n}{alpha}\PYG{+w}{ }\PYG{n}{keyword}\PYG{+w}{ }\PYG{n}{val}
\end{sphinxVerbatim}
\begin{itemize}
\item {} 
\sphinxAtStartPar
ID, group\sphinxhyphen{}ID are documented in {\hyperref[\detokenize{compute::doc}]{\sphinxcrossref{\DUrole{doc}{compute}}}} command

\item {} 
\sphinxAtStartPar
efield/atom/wolf = style name of this compute command

\item {} 
\sphinxAtStartPar
alpha = damping parameter (inverse distance units)

\item {} 
\sphinxAtStartPar
zero or more keyword/value pairs may be appended

\item {} 
\sphinxAtStartPar
keyword = \sphinxstyleemphasis{limit} or \sphinxstyleemphasis{cutoff}
\begin{sphinxalltt}
\sphinxstyleemphasis{limit} group2\sphinxhyphen{}ID = limit computing the electric field contributions to a group (default: all)
\sphinxstyleemphasis{cutoff} value = set cutoff for computing contributions to this value (default: maximum cutoff of pair style)
\end{sphinxalltt}

\end{itemize}


\subsection{Examples}
\label{\detokenize{compute_efield_wolf_atom:examples}}
\begin{sphinxVerbatim}[commandchars=\\\{\}]
\PYG{k}{compute }\PYG{n+nv+nvIdentifier}{1}\PYG{+w}{ }\PYG{n+nv+nvIdentifier}{all}\PYG{+w}{ }\PYG{n}{efield}\PYG{o}{/}\PYG{n}{wolf}\PYG{o}{/}\PYG{n}{atom}\PYG{+w}{ }\PYG{l+m}{0.2}
\PYG{k}{compute }\PYG{n+nv+nvIdentifier}{1}\PYG{+w}{ }\PYG{n+nv+nvIdentifier}{mols}\PYG{+w}{ }\PYG{n}{efield}\PYG{o}{/}\PYG{n}{wolf}\PYG{o}{/}\PYG{n}{atom}\PYG{+w}{ }\PYG{l+m}{0.25}\PYG{+w}{ }\PYG{n}{limit}\PYG{+w}{ }\PYG{n}{water}\PYG{+w}{ }\PYG{n}{cutoff}\PYG{+w}{ }\PYG{l+m}{10.0}
\end{sphinxVerbatim}


\subsection{Description}
\label{\detokenize{compute_efield_wolf_atom:description}}
\sphinxAtStartPar
\DUrole{versionmodified,added}{Added in version 8Feb2023.}

\sphinxAtStartPar
Define a computation that approximates the electric field at each atom in a group.
\begin{equation*}
\begin{split}\vec{E}_i = \frac{\vec{F}coul_i}{q_i} = \sum_{j \neq i} \frac{q_j}{r_{ij}^2} \qquad r < r_c\end{split}
\end{equation*}
\sphinxAtStartPar
The electric field at the position of the atom \sphinxstyleemphasis{i} is the coulomb force
on a unit charge at that point, which is equivalent to dividing the
Coulomb force by the charge of the individual atom.

\sphinxAtStartPar
In this compute the electric field is approximated as the derivative of
the potential energy using the Wolf summation method, described in
{\hyperref[\detokenize{compute_efield_wolf_atom:wolf4}]{\sphinxcrossref{\DUrole{std,std-ref}{Wolf}}}}, given by:
\begin{equation*}
\begin{split}E_i = \frac{1}{2} \sum_{j \neq i}
\frac{q_i q_j {\rm erfc}(\alpha r_{ij})}{r_{ij}} +
\frac{1}{2} \sum_{j \neq i}
\frac{q_i q_j {\rm erf}(\alpha r_{ij})}{r_{ij}} \qquad r < r_c\end{split}
\end{equation*}
\sphinxAtStartPar
where \(\alpha\) is the damping parameter, and \sphinxstyleemphasis{erf()} and \sphinxstyleemphasis{erfc()}
are error\sphinxhyphen{}function and complementary error\sphinxhyphen{}function terms.  This
potential is essentially a short\sphinxhyphen{}range, spherically\sphinxhyphen{}truncated,
charge\sphinxhyphen{}neutralized, force\sphinxhyphen{}shifted, pairwise \sphinxstyleemphasis{1/r} summation.  With a
manipulation of adding and subtracting a self term (for i = j) to the
first and second term on the right\sphinxhyphen{}hand\sphinxhyphen{}side, respectively, and a small
enough \(\alpha\) damping parameter, the second term shrinks and the
potential becomes a rapidly\sphinxhyphen{}converging real\sphinxhyphen{}space summation.  With a
long enough cutoff and small enough \(\alpha\) parameter, the
electric field calculated by the Wolf summation method approaches that
computed using the Ewald sum.

\sphinxAtStartPar
The value of the electric field components will be 0.0 for atoms not in
the specified compute group.

\sphinxAtStartPar
When the \sphinxstyleemphasis{limit} keyword is used, only contributions from atoms in the
selected group will be considered, otherwise contributions from all
atoms within the cutoff are included.

\sphinxAtStartPar
When the \sphinxstyleemphasis{cutoff} keyword is used, the cutoff used for the electric
field approximation can be set explicitly.  By default it is the largest
cutoff of any pair style force computation.

\begin{sphinxadmonition}{note}{Computational Efficiency}

\sphinxAtStartPar
This compute will loop over a full neighbor list just like a pair
style does when computing forces, thus it can be quite time\sphinxhyphen{}consuming
and slow down a calculation significantly when its data is used in
every time step.  The {\hyperref[\detokenize{compute_efield_atom::doc}]{\sphinxcrossref{\DUrole{doc}{compute efield/atom}}}} command of the DIELECTRIC package is more
efficient in comparison, since the electric field data is collected
and stored as part of the force computation at next to no extra
computational cost.
\end{sphinxadmonition}


\subsection{Output info}
\label{\detokenize{compute_efield_wolf_atom:output-info}}
\sphinxAtStartPar
This compute calculates a per\sphinxhyphen{}atom vector, which can be accessed by
any command that uses per\sphinxhyphen{}atom values from a compute as input.  See
the {\hyperref[\detokenize{Howto_output::doc}]{\sphinxcrossref{\DUrole{doc}{Howto output}}}} page for an overview of
LAMMPS output options.

\sphinxAtStartPar
The vector contains 3 values per atom which are the x\sphinxhyphen{}, y\sphinxhyphen{}, and
z\sphinxhyphen{}direction electric field components in force units.


\subsection{Restrictions}
\label{\detokenize{compute_efield_wolf_atom:restrictions}}
\sphinxAtStartPar
This compute is part of the EXTRA\sphinxhyphen{}COMPUTE package. It is only enabled if
LAMMPS was built with that package.

\sphinxAtStartPar
This compute requires {\hyperref[\detokenize{neighbor::doc}]{\sphinxcrossref{\DUrole{doc}{neighbor styles ‘bin’ or ‘nsq’}}}}.


\subsection{Related commands}
\label{\detokenize{compute_efield_wolf_atom:related-commands}}
\sphinxAtStartPar
{\hyperref[\detokenize{pair_coul::doc}]{\sphinxcrossref{\DUrole{doc}{pair\_style coul/wolf}}}},
{\hyperref[\detokenize{compute_efield_atom::doc}]{\sphinxcrossref{\DUrole{doc}{compute efield/atom}}}}


\subsection{Default}
\label{\detokenize{compute_efield_wolf_atom:default}}
\sphinxAtStartPar
The option defaults are \sphinxstyleemphasis{limit} = all and \sphinxstyleemphasis{cutoff} = largest cutoff
for pair styles.


\bigskip\hrule\bigskip

\phantomsection\label{\detokenize{compute_efield_wolf_atom:wolf4}}
\sphinxAtStartPar
\sphinxstylestrong{(Wolf)} D. Wolf, P. Keblinski, S. R. Phillpot, J. Eggebrecht, J Chem
Phys, 110, 8254 (1999).

\sphinxstepscope

\index{compute entropy/atom@\spxentry{compute entropy/atom}}\ignorespaces 

\section{compute entropy/atom command}
\label{\detokenize{compute_entropy_atom:compute-entropy-atom-command}}\label{\detokenize{compute_entropy_atom:index-0}}\label{\detokenize{compute_entropy_atom::doc}}

\subsection{Syntax}
\label{\detokenize{compute_entropy_atom:syntax}}
\begin{sphinxVerbatim}[commandchars=\\\{\}]
\PYG{k}{compute }\PYG{n+nv+nvIdentifier}{ID}\PYG{+w}{ }\PYG{n+nv+nvIdentifier}{group\PYGZhy{}ID}\PYG{+w}{ }\PYG{n}{entropy}\PYG{o}{/}\PYG{n}{atom}\PYG{+w}{ }\PYG{n}{sigma}\PYG{+w}{ }\PYG{n}{cutoff}\PYG{+w}{ }\PYG{n}{keyword}\PYG{+w}{ }\PYG{n}{value}\PYG{+w}{ }\PYG{n}{...}
\end{sphinxVerbatim}
\begin{itemize}
\item {} 
\sphinxAtStartPar
ID, group\sphinxhyphen{}ID are documented in {\hyperref[\detokenize{compute::doc}]{\sphinxcrossref{\DUrole{doc}{compute}}}} command

\item {} 
\sphinxAtStartPar
entropy/atom = style name of this compute command

\item {} 
\sphinxAtStartPar
sigma = width of Gaussians used in the \(g(r)\) smoothing

\item {} 
\sphinxAtStartPar
cutoff = cutoff for the \(g(r)\) calculation

\item {} 
\sphinxAtStartPar
one or more keyword/value pairs may be appended

\end{itemize}
\begin{sphinxalltt}
keyword = \sphinxstyleemphasis{avg} or \sphinxstyleemphasis{local}
  \sphinxstyleemphasis{avg} args = neigh cutoff2
    neigh value = \sphinxstyleemphasis{yes} or \sphinxstyleemphasis{no} = whether to average the pair entropy over neighbors
    cutoff2 = cutoff for the averaging over neighbors
  \sphinxstyleemphasis{local} arg = \sphinxstyleemphasis{yes} or \sphinxstyleemphasis{no} = use the local density around each atom to normalize the g(r)
\end{sphinxalltt}


\subsection{Examples}
\label{\detokenize{compute_entropy_atom:examples}}
\begin{sphinxVerbatim}[commandchars=\\\{\}]
\PYG{k}{compute }\PYG{n+nv+nvIdentifier}{1}\PYG{+w}{ }\PYG{n+nv+nvIdentifier}{all}\PYG{+w}{ }\PYG{n}{entropy}\PYG{o}{/}\PYG{n}{atom}\PYG{+w}{ }\PYG{l+m}{0.25}\PYG{+w}{ }\PYG{l+m}{5}\PYG{n}{.}
\PYG{k}{compute }\PYG{n+nv+nvIdentifier}{1}\PYG{+w}{ }\PYG{n+nv+nvIdentifier}{all}\PYG{+w}{ }\PYG{n}{entropy}\PYG{o}{/}\PYG{n}{atom}\PYG{+w}{ }\PYG{l+m}{0.25}\PYG{+w}{ }\PYG{l+m}{5}\PYG{n}{.}\PYG{+w}{ }\PYG{n}{avg}\PYG{+w}{ }\PYG{n}{yes}\PYG{+w}{ }\PYG{l+m}{5}\PYG{n}{.}
\PYG{k}{compute }\PYG{n+nv+nvIdentifier}{1}\PYG{+w}{ }\PYG{n+nv+nvIdentifier}{all}\PYG{+w}{ }\PYG{n}{entropy}\PYG{o}{/}\PYG{n}{atom}\PYG{+w}{ }\PYG{l+m}{0.125}\PYG{+w}{ }\PYG{l+m}{7.3}\PYG{+w}{ }\PYG{n}{avg}\PYG{+w}{ }\PYG{n}{yes}\PYG{+w}{ }\PYG{l+m}{5.1}\PYG{+w}{ }\PYG{n}{local}\PYG{+w}{ }\PYG{n}{yes}
\end{sphinxVerbatim}


\subsection{Description}
\label{\detokenize{compute_entropy_atom:description}}
\sphinxAtStartPar
Define a computation that calculates the pair entropy fingerprint for
each atom in the group. The fingerprint is useful to distinguish between
ordered and disordered environments, for instance liquid and solid\sphinxhyphen{}like
environments, or glassy and crystalline\sphinxhyphen{}like environments. Some
applications could be the identification of grain boundaries, a
melt\sphinxhyphen{}solid interface, or a solid cluster emerging from the melt.
The advantage of this parameter over others is that no a priori
information about the solid structure is required.

\sphinxAtStartPar
This parameter for atom i is computed using the following formula from
{\hyperref[\detokenize{compute_entropy_atom:piaggi}]{\sphinxcrossref{\DUrole{std,std-ref}{(Piaggi)}}}} and {\hyperref[\detokenize{compute_entropy_atom:nettleton}]{\sphinxcrossref{\DUrole{std,std-ref}{(Nettleton)}}}} ,
\begin{equation*}
\begin{split}s_S^i=-2\pi\rho k_B \int\limits_0^{r_m} \left [ g(r) \ln g(r) - g(r) + 1 \right ] r^2 dr\end{split}
\end{equation*}
\sphinxAtStartPar
where \(r\) is a distance, \(g(r)\) is the radial distribution function
of atom \(i\), and \(\rho\) is the density of the system.
The \(g(r)\) computed for each atom \(i\) can be noisy and therefore it
is smoothed using
\begin{equation*}
\begin{split}g_m^i(r) = \frac{1}{4 \pi \rho r^2} \sum\limits_{j} \frac{1}{\sqrt{2 \pi \sigma^2}} e^{-(r-r_{ij})^2/(2\sigma^2)}\end{split}
\end{equation*}
\sphinxAtStartPar
where the sum over \(j\) goes through the neighbors of atom \(i\) and
\(\sigma\) is a parameter to control the smoothing.

\sphinxAtStartPar
The input parameters are \sphinxstyleemphasis{sigma} the smoothing parameter \(\sigma\),
and the \sphinxstyleemphasis{cutoff} for the calculation of \(g(r)\).

\sphinxAtStartPar
If the keyword \sphinxstyleemphasis{avg} has the setting \sphinxstyleemphasis{yes}, then this compute also
averages the parameter over the neighbors  of atom \(i\) according to
\begin{equation*}
\begin{split}\left< s_S^i \right>  = \frac{\sum_j s_S^j + s_S^i}{N + 1},\end{split}
\end{equation*}
\sphinxAtStartPar
where the sum over \(j\) goes over the neighbors of atom \(i\) and
\(N\) is the number of neighbors. This procedure provides a sharper
distinction between order and disorder environments. In this case the input
parameter \sphinxstyleemphasis{cutoff2} is the cutoff for the averaging over the neighbors and
must also be specified.

\sphinxAtStartPar
If the \sphinxstyleemphasis{avg yes} option is used, the effective cutoff of the neighbor
list should be \sphinxstyleemphasis{cutoff}+\sphinxstyleemphasis{cutoff2} and therefore it might be necessary
to increase the skin of the neighbor list with:

\begin{sphinxVerbatim}[commandchars=\\\{\}]
neighbor \PYGZlt{}skin distance\PYGZgt{} bin
\end{sphinxVerbatim}

\sphinxAtStartPar
See {\hyperref[\detokenize{neighbor::doc}]{\sphinxcrossref{\DUrole{doc}{neighbor}}}} for details.

\sphinxAtStartPar
If the \sphinxstyleemphasis{local yes} option is used, the \(g(r)\) is normalized by the
local density around each atom, that is to say the density around each
atom  is the number of neighbors within the neighbor list cutoff divided
by the corresponding volume. This option can be useful when dealing with
inhomogeneous systems such as those that have surfaces.

\sphinxAtStartPar
Here are typical input parameters for fcc aluminum (lattice
constant \(4.05~\AA\)),

\begin{sphinxVerbatim}[commandchars=\\\{\}]
compute 1 all entropy/atom 0.25 5.7 avg yes 3.7
\end{sphinxVerbatim}

\sphinxAtStartPar
and for bcc sodium (lattice constant \(4.23~\AA\)),

\begin{sphinxVerbatim}[commandchars=\\\{\}]
compute 1 all entropy/atom 0.25 7.3 avg yes 5.1
\end{sphinxVerbatim}


\subsection{Output info}
\label{\detokenize{compute_entropy_atom:output-info}}
\sphinxAtStartPar
By default, this compute calculates the pair entropy value for each
atom as a per\sphinxhyphen{}atom vector, which can be accessed by any command that
uses per\sphinxhyphen{}atom values from a compute as input.  See the
{\hyperref[\detokenize{Howto_output::doc}]{\sphinxcrossref{\DUrole{doc}{Howto output}}}} page for an overview of LAMMPS output
options.

\sphinxAtStartPar
The pair entropy values have units of the Boltzmann constant. They are
always negative, and lower values (lower entropy) correspond to more
ordered environments.


\subsection{Restrictions}
\label{\detokenize{compute_entropy_atom:restrictions}}
\sphinxAtStartPar
This compute is part of the EXTRA\sphinxhyphen{}COMPUTE package.  It is only enabled if
LAMMPS was built with that package.  See the {\hyperref[\detokenize{Build_package::doc}]{\sphinxcrossref{\DUrole{doc}{Build package}}}} page for more info.


\subsection{Related commands}
\label{\detokenize{compute_entropy_atom:related-commands}}
\sphinxAtStartPar
{\hyperref[\detokenize{compute_cna_atom::doc}]{\sphinxcrossref{\DUrole{doc}{compute cna/atom}}}}
{\hyperref[\detokenize{compute_centro_atom::doc}]{\sphinxcrossref{\DUrole{doc}{compute centro/atom}}}}


\subsection{Default}
\label{\detokenize{compute_entropy_atom:default}}
\sphinxAtStartPar
The default values for the optional keywords are avg = no and local = no.


\bigskip\hrule\bigskip

\phantomsection\label{\detokenize{compute_entropy_atom:piaggi}}
\sphinxAtStartPar
\sphinxstylestrong{(Piaggi)} Piaggi and Parrinello, J Chem Phys, 147, 114112 (2017).

\phantomsection\label{\detokenize{compute_entropy_atom:nettleton}}
\sphinxAtStartPar
\sphinxstylestrong{(Nettleton)} Nettleton and Green, J Chem Phys, 29, 6 (1958).

\sphinxstepscope

\index{compute erotate/asphere@\spxentry{compute erotate/asphere}}\ignorespaces 

\section{compute erotate/asphere command}
\label{\detokenize{compute_erotate_asphere:compute-erotate-asphere-command}}\label{\detokenize{compute_erotate_asphere:index-0}}\label{\detokenize{compute_erotate_asphere::doc}}

\subsection{Syntax}
\label{\detokenize{compute_erotate_asphere:syntax}}
\begin{sphinxVerbatim}[commandchars=\\\{\}]
\PYG{k}{compute }\PYG{n+nv+nvIdentifier}{ID}\PYG{+w}{ }\PYG{n+nv+nvIdentifier}{group\PYGZhy{}ID}\PYG{+w}{ }\PYG{n}{erotate}\PYG{o}{/}\PYG{n}{asphere}
\end{sphinxVerbatim}
\begin{itemize}
\item {} 
\sphinxAtStartPar
ID, group\sphinxhyphen{}ID are documented in {\hyperref[\detokenize{compute::doc}]{\sphinxcrossref{\DUrole{doc}{compute}}}} command

\item {} 
\sphinxAtStartPar
erotate/asphere = style name of this compute command

\end{itemize}


\subsection{Examples}
\label{\detokenize{compute_erotate_asphere:examples}}
\begin{sphinxVerbatim}[commandchars=\\\{\}]
\PYG{k}{compute }\PYG{n+nv+nvIdentifier}{1}\PYG{+w}{ }\PYG{n+nv+nvIdentifier}{all}\PYG{+w}{ }\PYG{n}{erotate}\PYG{o}{/}\PYG{n}{asphere}
\end{sphinxVerbatim}


\subsection{Description}
\label{\detokenize{compute_erotate_asphere:description}}
\sphinxAtStartPar
Define a computation that calculates the rotational kinetic energy of
a group of aspherical particles.  The aspherical particles can be
ellipsoids, or line segments, or triangles.  See the
{\hyperref[\detokenize{atom_style::doc}]{\sphinxcrossref{\DUrole{doc}{atom\_style}}}} and {\hyperref[\detokenize{read_data::doc}]{\sphinxcrossref{\DUrole{doc}{read\_data}}}} commands
for descriptions of these options.

\sphinxAtStartPar
For all 3 types of particles, the rotational kinetic energy is
computed as \(\frac12 I \omega^2\), where \(I\) is the inertia tensor
for the aspherical particle and \(\omega\) is its angular velocity, which
is computed from its angular momentum if needed.

\begin{sphinxadmonition}{note}{Note:}
\sphinxAtStartPar
For {\hyperref[\detokenize{dimension::doc}]{\sphinxcrossref{\DUrole{doc}{2d models}}}}, ellipsoidal particles are
treated as ellipsoids, not ellipses, meaning their moments of inertia
will be the same as in 3d.
\end{sphinxadmonition}


\subsection{Output info}
\label{\detokenize{compute_erotate_asphere:output-info}}
\sphinxAtStartPar
This compute calculates a global scalar (the KE).  This value can be
used by any command that uses a global scalar value from a compute as
input.  See the {\hyperref[\detokenize{Howto_output::doc}]{\sphinxcrossref{\DUrole{doc}{Howto output}}}} page for an
overview of LAMMPS output options.

\sphinxAtStartPar
The scalar value calculated by this compute is “extensive”.  The
scalar value will be in energy {\hyperref[\detokenize{units::doc}]{\sphinxcrossref{\DUrole{doc}{units}}}}.


\subsection{Restrictions}
\label{\detokenize{compute_erotate_asphere:restrictions}}
\sphinxAtStartPar
This compute requires that ellipsoidal particles atoms store a shape
and quaternion orientation and angular momentum as defined by the
{\hyperref[\detokenize{atom_style::doc}]{\sphinxcrossref{\DUrole{doc}{atom\_style ellipsoid}}}} command.

\sphinxAtStartPar
This compute requires that line segment particles atoms store a length
and orientation and angular velocity as defined by the {\hyperref[\detokenize{atom_style::doc}]{\sphinxcrossref{\DUrole{doc}{atom\_style line}}}} command.

\sphinxAtStartPar
This compute requires that triangular particles atoms store a size and
shape and quaternion orientation and angular momentum as defined by
the {\hyperref[\detokenize{atom_style::doc}]{\sphinxcrossref{\DUrole{doc}{atom\_style tri}}}} command.

\sphinxAtStartPar
All particles in the group must be of finite size.  They cannot be point
particles.


\subsection{Related commands}
\label{\detokenize{compute_erotate_asphere:related-commands}}
\sphinxAtStartPar
none

\sphinxAtStartPar
{\hyperref[\detokenize{compute_erotate_sphere::doc}]{\sphinxcrossref{\DUrole{doc}{compute erotate/sphere}}}}


\subsection{Default}
\label{\detokenize{compute_erotate_asphere:default}}
\sphinxAtStartPar
none

\sphinxstepscope

\index{compute erotate/rigid@\spxentry{compute erotate/rigid}}\ignorespaces 

\section{compute erotate/rigid command}
\label{\detokenize{compute_erotate_rigid:compute-erotate-rigid-command}}\label{\detokenize{compute_erotate_rigid:index-0}}\label{\detokenize{compute_erotate_rigid::doc}}

\subsection{Syntax}
\label{\detokenize{compute_erotate_rigid:syntax}}
\begin{sphinxVerbatim}[commandchars=\\\{\}]
\PYG{k}{compute }\PYG{n+nv+nvIdentifier}{ID}\PYG{+w}{ }\PYG{n+nv+nvIdentifier}{group\PYGZhy{}ID}\PYG{+w}{ }\PYG{n}{erotate}\PYG{o}{/}\PYG{n}{rigid}\PYG{+w}{ }\PYG{n}{fix}\PYG{o}{\PYGZhy{}}\PYG{n}{ID}
\end{sphinxVerbatim}
\begin{itemize}
\item {} 
\sphinxAtStartPar
ID, group\sphinxhyphen{}ID are documented in {\hyperref[\detokenize{compute::doc}]{\sphinxcrossref{\DUrole{doc}{compute}}}} command

\item {} 
\sphinxAtStartPar
erotate/rigid = style name of this compute command

\item {} 
\sphinxAtStartPar
fix\sphinxhyphen{}ID = ID of rigid body fix

\end{itemize}


\subsection{Examples}
\label{\detokenize{compute_erotate_rigid:examples}}
\begin{sphinxVerbatim}[commandchars=\\\{\}]
\PYG{k}{compute }\PYG{n+nv+nvIdentifier}{1}\PYG{+w}{ }\PYG{n+nv+nvIdentifier}{all}\PYG{+w}{ }\PYG{n}{erotate}\PYG{o}{/}\PYG{n}{rigid}\PYG{+w}{ }\PYG{n}{myRigid}
\end{sphinxVerbatim}


\subsection{Description}
\label{\detokenize{compute_erotate_rigid:description}}
\sphinxAtStartPar
Define a computation that calculates the rotational kinetic energy of
a collection of rigid bodies, as defined by one of the
{\hyperref[\detokenize{fix_rigid::doc}]{\sphinxcrossref{\DUrole{doc}{fix rigid}}}} command variants.

\sphinxAtStartPar
The rotational energy of each rigid body is computed as
\(\frac12 I \omega_\text{body}^2\),
where \(I\) is the inertia tensor for the rigid body and
\(\omega_\text{body}\) is its angular velocity vector.
Both \(I\) and \(\omega_\text{body}\) are in the frame of
reference of the rigid body (i.e., \(I\) is diagonal).

\sphinxAtStartPar
The \sphinxstyleemphasis{fix\sphinxhyphen{}ID} should be the ID of one of the {\hyperref[\detokenize{fix_rigid::doc}]{\sphinxcrossref{\DUrole{doc}{fix rigid}}}}
commands which defines the rigid bodies.  The group specified in the
compute command is ignored.  The rotational energy of all the rigid
bodies defined by the fix rigid command in included in the calculation.


\subsection{Output info}
\label{\detokenize{compute_erotate_rigid:output-info}}
\sphinxAtStartPar
This compute calculates a global scalar (the summed rotational energy
of all the rigid bodies).  This value can be used by any command that
uses a global scalar value from a compute as input.  See the {\hyperref[\detokenize{Howto_output::doc}]{\sphinxcrossref{\DUrole{doc}{Howto output}}}} page for an overview of LAMMPS output
options.

\sphinxAtStartPar
The scalar value calculated by this compute is “extensive”.  The
scalar value will be in energy {\hyperref[\detokenize{units::doc}]{\sphinxcrossref{\DUrole{doc}{units}}}}.


\subsection{Restrictions}
\label{\detokenize{compute_erotate_rigid:restrictions}}
\sphinxAtStartPar
This compute is part of the RIGID package.  It is only enabled if
LAMMPS was built with that package.  See the
{\hyperref[\detokenize{Build_package::doc}]{\sphinxcrossref{\DUrole{doc}{Build package}}}} page for more info.


\subsection{Related commands}
\label{\detokenize{compute_erotate_rigid:related-commands}}
\sphinxAtStartPar
{\hyperref[\detokenize{compute_ke_rigid::doc}]{\sphinxcrossref{\DUrole{doc}{compute ke/rigid}}}}


\subsection{Default}
\label{\detokenize{compute_erotate_rigid:default}}
\sphinxAtStartPar
none

\sphinxstepscope

\index{compute erotate/sphere@\spxentry{compute erotate/sphere}}\ignorespaces 
\index{compute erotate/sphere/kk@\spxentry{compute erotate/sphere/kk}}\ignorespaces 

\section{compute erotate/sphere command}
\label{\detokenize{compute_erotate_sphere:compute-erotate-sphere-command}}\label{\detokenize{compute_erotate_sphere:index-1}}\label{\detokenize{compute_erotate_sphere:index-0}}\label{\detokenize{compute_erotate_sphere::doc}}
\sphinxAtStartPar
Accelerator Variants: \sphinxstyleemphasis{erotate/sphere/kk}


\subsection{Syntax}
\label{\detokenize{compute_erotate_sphere:syntax}}
\begin{sphinxVerbatim}[commandchars=\\\{\}]
\PYG{k}{compute }\PYG{n+nv+nvIdentifier}{ID}\PYG{+w}{ }\PYG{n+nv+nvIdentifier}{group\PYGZhy{}ID}\PYG{+w}{ }\PYG{n}{erotate}\PYG{o}{/}\PYG{n}{sphere}
\end{sphinxVerbatim}
\begin{itemize}
\item {} 
\sphinxAtStartPar
ID, group\sphinxhyphen{}ID are documented in {\hyperref[\detokenize{compute::doc}]{\sphinxcrossref{\DUrole{doc}{compute}}}} command

\item {} 
\sphinxAtStartPar
erotate/sphere = style name of this compute command

\end{itemize}


\subsection{Examples}
\label{\detokenize{compute_erotate_sphere:examples}}
\begin{sphinxVerbatim}[commandchars=\\\{\}]
\PYG{k}{compute }\PYG{n+nv+nvIdentifier}{1}\PYG{+w}{ }\PYG{n+nv+nvIdentifier}{all}\PYG{+w}{ }\PYG{n}{erotate}\PYG{o}{/}\PYG{n}{sphere}
\end{sphinxVerbatim}


\subsection{Description}
\label{\detokenize{compute_erotate_sphere:description}}
\sphinxAtStartPar
Define a computation that calculates the rotational kinetic energy of
a group of spherical particles.

\sphinxAtStartPar
The rotational energy is computed as \(\frac12 I \omega^2\),
where \(I\) is the moment of inertia for a sphere and \(\omega\)
is the particle’s angular velocity.

\begin{sphinxadmonition}{note}{Note:}
\sphinxAtStartPar
For {\hyperref[\detokenize{dimension::doc}]{\sphinxcrossref{\DUrole{doc}{2d models}}}}, particles are treated as
spheres, not disks, meaning their moment of inertia will be the same
as in 3d.
\end{sphinxadmonition}


\bigskip\hrule\bigskip


\sphinxAtStartPar
Styles with a \sphinxstyleemphasis{gpu}, \sphinxstyleemphasis{intel}, \sphinxstyleemphasis{kk}, \sphinxstyleemphasis{omp}, or \sphinxstyleemphasis{opt} suffix are
functionally the same as the corresponding style without the suffix.
They have been optimized to run faster, depending on your available
hardware, as discussed on the {\hyperref[\detokenize{Speed_packages::doc}]{\sphinxcrossref{\DUrole{doc}{Accelerator packages}}}}
page.  The accelerated styles take the same arguments and should
produce the same results, except for round\sphinxhyphen{}off and precision issues.

\sphinxAtStartPar
These accelerated styles are part of the GPU, INTEL, KOKKOS,
OPENMP, and OPT packages, respectively.  They are only enabled if
LAMMPS was built with those packages.  See the {\hyperref[\detokenize{Build_package::doc}]{\sphinxcrossref{\DUrole{doc}{Build package}}}} page for more info.

\sphinxAtStartPar
You can specify the accelerated styles explicitly in your input script
by including their suffix, or you can use the {\hyperref[\detokenize{Run_options::doc}]{\sphinxcrossref{\DUrole{doc}{\sphinxhyphen{}suffix command\sphinxhyphen{}line switch}}}} when you invoke LAMMPS, or you can use the
{\hyperref[\detokenize{suffix::doc}]{\sphinxcrossref{\DUrole{doc}{suffix}}}} command in your input script.

\sphinxAtStartPar
See the {\hyperref[\detokenize{Speed_packages::doc}]{\sphinxcrossref{\DUrole{doc}{Accelerator packages}}}} page for more
instructions on how to use the accelerated styles effectively.


\bigskip\hrule\bigskip



\subsection{Output info}
\label{\detokenize{compute_erotate_sphere:output-info}}
\sphinxAtStartPar
This compute calculates a global scalar (the KE).  This value can be
used by any command that uses a global scalar value from a compute as
input.  See the {\hyperref[\detokenize{Howto_output::doc}]{\sphinxcrossref{\DUrole{doc}{Howto output}}}} page for an
overview of LAMMPS output options.

\sphinxAtStartPar
The scalar value calculated by this compute is “extensive”.  The
scalar value will be in energy {\hyperref[\detokenize{units::doc}]{\sphinxcrossref{\DUrole{doc}{units}}}}.


\subsection{Restrictions}
\label{\detokenize{compute_erotate_sphere:restrictions}}
\sphinxAtStartPar
This compute requires that atoms store a radius and angular velocity
(omega) as defined by the {\hyperref[\detokenize{atom_style::doc}]{\sphinxcrossref{\DUrole{doc}{atom\_style sphere}}}} command.

\sphinxAtStartPar
All particles in the group must be finite\sphinxhyphen{}size spheres or point
particles.  They cannot be aspherical.  Point particles will not
contribute to the rotational energy.


\subsection{Related commands}
\label{\detokenize{compute_erotate_sphere:related-commands}}
\sphinxAtStartPar
{\hyperref[\detokenize{compute_erotate_asphere::doc}]{\sphinxcrossref{\DUrole{doc}{compute erotate/asphere}}}}


\subsection{Default}
\label{\detokenize{compute_erotate_sphere:default}}
\sphinxAtStartPar
none

\sphinxstepscope

\index{compute erotate/sphere/atom@\spxentry{compute erotate/sphere/atom}}\ignorespaces 

\section{compute erotate/sphere/atom command}
\label{\detokenize{compute_erotate_sphere_atom:compute-erotate-sphere-atom-command}}\label{\detokenize{compute_erotate_sphere_atom:index-0}}\label{\detokenize{compute_erotate_sphere_atom::doc}}

\subsection{Syntax}
\label{\detokenize{compute_erotate_sphere_atom:syntax}}
\begin{sphinxVerbatim}[commandchars=\\\{\}]
\PYG{k}{compute }\PYG{n+nv+nvIdentifier}{ID}\PYG{+w}{ }\PYG{n+nv+nvIdentifier}{group\PYGZhy{}ID}\PYG{+w}{ }\PYG{n}{erotate}\PYG{o}{/}\PYG{n}{sphere}\PYG{o}{/}\PYG{n}{atom}
\end{sphinxVerbatim}
\begin{itemize}
\item {} 
\sphinxAtStartPar
ID, group\sphinxhyphen{}ID are documented in {\hyperref[\detokenize{compute::doc}]{\sphinxcrossref{\DUrole{doc}{compute}}}} command

\item {} 
\sphinxAtStartPar
erotate/sphere/atom = style name of this compute command

\end{itemize}


\subsection{Examples}
\label{\detokenize{compute_erotate_sphere_atom:examples}}
\begin{sphinxVerbatim}[commandchars=\\\{\}]
\PYG{k}{compute }\PYG{n+nv+nvIdentifier}{1}\PYG{+w}{ }\PYG{n+nv+nvIdentifier}{all}\PYG{+w}{ }\PYG{n}{erotate}\PYG{o}{/}\PYG{n}{sphere}\PYG{o}{/}\PYG{n}{atom}
\end{sphinxVerbatim}


\subsection{Description}
\label{\detokenize{compute_erotate_sphere_atom:description}}
\sphinxAtStartPar
Define a computation that calculates the rotational kinetic energy for
each particle in a group.

\sphinxAtStartPar
The rotational energy is computed as \(\frac12 I \omega^2\), where
\(I\) is the moment of inertia for a sphere and \(\omega\) is the
particle’s angular velocity.

\begin{sphinxadmonition}{note}{Note:}
\sphinxAtStartPar
For {\hyperref[\detokenize{dimension::doc}]{\sphinxcrossref{\DUrole{doc}{2d models}}}}, particles are treated as
spheres, not disks, meaning their moment of inertia will be the same
as in 3d.
\end{sphinxadmonition}

\sphinxAtStartPar
The value of the rotational kinetic energy will be 0.0 for atoms not
in the specified compute group or for point particles with a radius of 0.0.


\subsection{Output info}
\label{\detokenize{compute_erotate_sphere_atom:output-info}}
\sphinxAtStartPar
This compute calculates a per\sphinxhyphen{}atom vector, which can be accessed by
any command that uses per\sphinxhyphen{}atom values from a compute as input.  See
the {\hyperref[\detokenize{Howto_output::doc}]{\sphinxcrossref{\DUrole{doc}{Howto output}}}} page for an overview of
LAMMPS output options.

\sphinxAtStartPar
The per\sphinxhyphen{}atom vector values will be in energy {\hyperref[\detokenize{units::doc}]{\sphinxcrossref{\DUrole{doc}{units}}}}.


\subsection{Restrictions}
\label{\detokenize{compute_erotate_sphere_atom:restrictions}}\begin{quote}

\sphinxAtStartPar
none
\end{quote}


\subsection{Related commands}
\label{\detokenize{compute_erotate_sphere_atom:related-commands}}
\sphinxAtStartPar
{\hyperref[\detokenize{dump::doc}]{\sphinxcrossref{\DUrole{doc}{dump custom}}}}


\subsection{Default}
\label{\detokenize{compute_erotate_sphere_atom:default}}
\sphinxAtStartPar
none

\sphinxstepscope

\index{compute event/displace@\spxentry{compute event/displace}}\ignorespaces 

\section{compute event/displace command}
\label{\detokenize{compute_event_displace:compute-event-displace-command}}\label{\detokenize{compute_event_displace:index-0}}\label{\detokenize{compute_event_displace::doc}}

\subsection{Syntax}
\label{\detokenize{compute_event_displace:syntax}}
\begin{sphinxVerbatim}[commandchars=\\\{\}]
\PYG{k}{compute }\PYG{n+nv+nvIdentifier}{ID}\PYG{+w}{ }\PYG{n+nv+nvIdentifier}{group\PYGZhy{}ID}\PYG{+w}{ }\PYG{n}{event}\PYG{o}{/}\PYG{n}{displace}\PYG{+w}{ }\PYG{n}{threshold}
\end{sphinxVerbatim}
\begin{itemize}
\item {} 
\sphinxAtStartPar
ID, group\sphinxhyphen{}ID are documented in {\hyperref[\detokenize{compute::doc}]{\sphinxcrossref{\DUrole{doc}{compute}}}} command

\item {} 
\sphinxAtStartPar
event/displace = style name of this compute command

\item {} 
\sphinxAtStartPar
threshold = minimum distance any particle must move to trigger an event (distance units)

\end{itemize}


\subsection{Examples}
\label{\detokenize{compute_event_displace:examples}}
\begin{sphinxVerbatim}[commandchars=\\\{\}]
\PYG{k}{compute }\PYG{n+nv+nvIdentifier}{1}\PYG{+w}{ }\PYG{n+nv+nvIdentifier}{all}\PYG{+w}{ }\PYG{n}{event}\PYG{o}{/}\PYG{n}{displace}\PYG{+w}{ }\PYG{l+m}{0.5}
\end{sphinxVerbatim}


\subsection{Description}
\label{\detokenize{compute_event_displace:description}}
\sphinxAtStartPar
Define a computation that flags an “event” if any particle in the
group has moved a distance greater than the specified threshold
distance when compared to a previously stored reference state
(i.e., the previous event).  This compute is typically used in
conjunction with the {\hyperref[\detokenize{prd::doc}]{\sphinxcrossref{\DUrole{doc}{prd}}}} and {\hyperref[\detokenize{tad::doc}]{\sphinxcrossref{\DUrole{doc}{tad}}}} commands,
to detect if a transition to a new minimum energy basin has occurred.

\sphinxAtStartPar
This value calculated by the compute is equal to 0 if no particle has
moved far enough, and equal to 1 if one or more particles have moved
further than the threshold distance.

\begin{sphinxadmonition}{note}{Note:}
\sphinxAtStartPar
If the system is undergoing significant center\sphinxhyphen{}of\sphinxhyphen{}mass motion,
due to thermal motion, an external force, or an initial net momentum,
then this compute will not be able to distinguish that motion from
local atom displacements and may generate “false positives”.
\end{sphinxadmonition}


\subsection{Output info}
\label{\detokenize{compute_event_displace:output-info}}
\sphinxAtStartPar
This compute calculates a global scalar (the flag).  This value can be
used by any command that uses a global scalar value from a compute as
input.  See the {\hyperref[\detokenize{Howto_output::doc}]{\sphinxcrossref{\DUrole{doc}{Howto output}}}} page for an
overview of LAMMPS output options.

\sphinxAtStartPar
The scalar value calculated by this compute is “intensive”.  The
scalar value will be a 0 or 1 as explained above.


\subsection{Restrictions}
\label{\detokenize{compute_event_displace:restrictions}}
\sphinxAtStartPar
This command can only be used if LAMMPS was built with the REPLICA
package.  See the {\hyperref[\detokenize{Build_package::doc}]{\sphinxcrossref{\DUrole{doc}{Build package}}}} doc
page for more info.


\subsection{Related commands}
\label{\detokenize{compute_event_displace:related-commands}}
\sphinxAtStartPar
{\hyperref[\detokenize{prd::doc}]{\sphinxcrossref{\DUrole{doc}{prd}}}}, {\hyperref[\detokenize{tad::doc}]{\sphinxcrossref{\DUrole{doc}{tad}}}}


\subsection{Default}
\label{\detokenize{compute_event_displace:default}}
\sphinxAtStartPar
none

\sphinxstepscope

\index{compute fabric@\spxentry{compute fabric}}\ignorespaces 

\section{compute fabric command}
\label{\detokenize{compute_fabric:compute-fabric-command}}\label{\detokenize{compute_fabric:index-0}}\label{\detokenize{compute_fabric::doc}}

\subsection{Syntax}
\label{\detokenize{compute_fabric:syntax}}
\begin{sphinxVerbatim}[commandchars=\\\{\}]
\PYG{k}{compute }\PYG{n+nv+nvIdentifier}{ID}\PYG{+w}{ }\PYG{n+nv+nvIdentifier}{group\PYGZhy{}ID}\PYG{+w}{ }\PYG{n}{fabric}\PYG{+w}{ }\PYG{n}{cutoff}\PYG{+w}{ }\PYG{n}{attribute}\PYG{+w}{ }\PYG{n}{...}\PYG{+w}{ }\PYG{n}{keyword}\PYG{+w}{ }\PYG{n}{values}\PYG{+w}{ }\PYG{n}{...}
\end{sphinxVerbatim}
\begin{itemize}
\item {} 
\sphinxAtStartPar
ID, group\sphinxhyphen{}ID are documented in {\hyperref[\detokenize{compute::doc}]{\sphinxcrossref{\DUrole{doc}{compute}}}} command

\item {} 
\sphinxAtStartPar
fabric = style name of this compute command

\item {} 
\sphinxAtStartPar
cutoff = \sphinxstyleemphasis{type} or \sphinxstyleemphasis{radius}
\begin{sphinxalltt}
\sphinxstyleemphasis{type} = cutoffs determined based on atom types
\sphinxstyleemphasis{radius} = cutoffs determined based on atom diameters (atom style sphere)
\end{sphinxalltt}

\item {} 
\sphinxAtStartPar
one or more attributes may be appended

\item {} 
\sphinxAtStartPar
attribute = \sphinxstyleemphasis{contact} or \sphinxstyleemphasis{branch} or \sphinxstyleemphasis{force/normal} or \sphinxstyleemphasis{force/tangential}
\begin{sphinxalltt}
\sphinxstyleemphasis{contact} = contact tensor
\sphinxstyleemphasis{branch} = branch tensor
\sphinxstyleemphasis{force/normal} = normal force tensor
\sphinxstyleemphasis{force/tangential} = tangential force tensor
\end{sphinxalltt}

\item {} 
\sphinxAtStartPar
zero or more keyword/value pairs may be appended

\item {} 
\sphinxAtStartPar
keyword = \sphinxstyleemphasis{type/include}
\begin{sphinxalltt}
\sphinxstyleemphasis{type/include} value = arg1 arg2
  arg = separate lists of types (see below)
\end{sphinxalltt}

\end{itemize}


\subsection{Examples}
\label{\detokenize{compute_fabric:examples}}
\begin{sphinxVerbatim}[commandchars=\\\{\}]
\PYG{k}{compute }\PYG{n+nv+nvIdentifier}{1}\PYG{+w}{ }\PYG{n+nv+nvIdentifier}{all}\PYG{+w}{ }\PYG{n}{fabric}\PYG{+w}{ }\PYG{n}{type}\PYG{+w}{ }\PYG{n}{contact}\PYG{+w}{ }\PYG{n}{force}\PYG{o}{/}\PYG{n}{normal}\PYG{+w}{ }\PYG{n}{type}\PYG{o}{/}\PYG{n}{include}\PYG{+w}{ }\PYG{l+m}{1},2\PYG{+w}{ }\PYG{l+m}{3}\PYG{o}{*}\PYG{l+m}{4}
\PYG{k}{compute }\PYG{n+nv+nvIdentifier}{1}\PYG{+w}{ }\PYG{n+nv+nvIdentifier}{all}\PYG{+w}{ }\PYG{n}{fabric}\PYG{+w}{ }\PYG{n}{radius}\PYG{+w}{ }\PYG{n}{force}\PYG{o}{/}\PYG{n}{normal}\PYG{+w}{ }\PYG{n}{force}\PYG{o}{/}\PYG{n}{tangential}
\end{sphinxVerbatim}


\subsection{Description}
\label{\detokenize{compute_fabric:description}}
\sphinxAtStartPar
Define a compute that calculates various fabric tensors for pairwise
interaction {\hyperref[\detokenize{compute_fabric:ouadfel}]{\sphinxcrossref{\DUrole{std,std-ref}{(Ouadfel)}}}}. Fabric tensors are commonly used
to quantify the anisotropy or orientation of granular contacts but can also
be used to characterize the direction of pairwise interactions in general
systems. The \sphinxstyleemphasis{type} and \sphinxstyleemphasis{radius} settings are used to select whether interactions
cutoffs are determined by atom types or by the sum of atomic radii (atom
style sphere), respectively. Calling this compute is roughly the cost of a
pair style invocation as it involves a loop over the neighbor list. If the
normal or tangential force tensors are requested, it will be more expensive
than a pair style invocation as it will also recalculate all pair forces.

\sphinxAtStartPar
Four fabric tensors are available: the contact, branch, normal force, or
tangential force tensor. The contact tensor is calculated as
\begin{equation*}
\begin{split}C_{ab}  =  \frac{15}{2} (\phi_{ab} - \frac{1}{3} \mathrm{Tr}(\phi) \delta_{ab})\end{split}
\end{equation*}
\sphinxAtStartPar
where \(a\) and \(b\) are the \(x\), \(y\), \(z\)
directions, \(\delta_{ab}\) is the Kronecker delta function, and
the tensor \(\phi\) is defined as
\begin{equation*}
\begin{split}\phi_{ab}  =  \sum_{n = 1}^{N_p} \frac{r_{a} r_{b}}{r^2}\end{split}
\end{equation*}
\sphinxAtStartPar
where \(n\) loops over the \(N_p\) pair interactions in the simulation,
\(r_{a}\) is the \(a\) component of the radial vector between the
two pairwise interacting particles, and \(r\) is the magnitude of the
radial vector.

\sphinxAtStartPar
The branch tensor is calculated as
\begin{equation*}
\begin{split}B_{ab}  =  \frac{15}{2\, \mathrm{Tr}(D)} (D_{ab} - \frac{1}{3} \mathrm{Tr}(D) \delta_{ab})\end{split}
\end{equation*}
\sphinxAtStartPar
where the tensor \(D\) is defined as
\begin{equation*}
\begin{split}D_{ab}  =  \sum_{n = 1}^{N_p}
             \frac{1}{N_c (r^2 + C_{cd} r_c r_d)}
             \frac{r_{a} r_{b}}{r}\end{split}
\end{equation*}
\sphinxAtStartPar
where \(N_c\) is the total number of contacts in the system and the
subscripts \(c\) and \(d\) indices are summed according to Einstein
notation.

\sphinxAtStartPar
The normal force fabric tensor is calculated as
\begin{equation*}
\begin{split}F^n_{ab}  =  \frac{15}{2\, \mathrm{Tr}(N)} (N_{ab} - \frac{1}{3} \mathrm{Tr}(N) \delta_{ab})\end{split}
\end{equation*}
\sphinxAtStartPar
where the tensor \(N\) is defined as
\begin{equation*}
\begin{split}N_{ab}  =  \sum_{n = 1}^{N_p}
             \frac{1}{N_c (r^2 + C_{cd} r_c r_d)}
             \frac{r_{a} r_{b}}{r^2} f_n\end{split}
\end{equation*}
\sphinxAtStartPar
and \(f_n\) is the magnitude of the normal, central\sphinxhyphen{}body force between the two atoms.

\sphinxAtStartPar
Finally, the tangential force fabric tensor is only defined for pair styles that
apply tangential forces to particles, namely granular pair styles. It is calculated
as
\begin{equation*}
\begin{split}F^t_{ab}  =  \frac{5}{\mathrm{Tr}(N)} (T_{ab} - \frac{1}{3} \mathrm{Tr}(T) \delta_{ab})\end{split}
\end{equation*}
\sphinxAtStartPar
where the tensor \(T\) is defined as
\begin{equation*}
\begin{split}T_{ab}  =  \sum_{n = 1}^{N_p}
             \frac{1}{N_c (r^2 + C_{cd} r_c r_d)}
             \frac{r_{a} r_{b}}{r^2} f_t\end{split}
\end{equation*}
\sphinxAtStartPar
and \(f_t\) is the magnitude of the tangential force between the two atoms.

\sphinxAtStartPar
The \sphinxstyleemphasis{type/include} keyword filters interactions based on the types of the two atoms.
Interactions between two atoms are only included in calculations if the atom types
are in the two lists. Each list consists of a series of type
ranges separated by commas. The range can be specified as a
single numeric value, or a wildcard asterisk can be used to specify a range
of values.  This takes the form “*” or “*n” or “m*” or “m*n”.  For
example, if \(M\) is the number of atom types, then an asterisk with no
numeric values means all types from 1 to \(M\).  A leading asterisk means
all types from 1 to n (inclusive).  A trailing asterisk means all types from
m to \(M\) (inclusive).  A middle asterisk means all types from m to n
(inclusive).  Multiple \sphinxstyleemphasis{type/include} keywords may be added.


\subsection{Output info}
\label{\detokenize{compute_fabric:output-info}}
\sphinxAtStartPar
This compute calculates a global vector of doubles and a global scalar. The
vector stores the unique components of the first requested tensor in the
order \(xx\), \(yy\), \(zz\), \(xy\), \(xz\),
\(yz\) followed by the same components for all subsequent tensors.
The length of the vector is therefore six times the number of requested
tensors. The scalar output is the number of pairwise interactions included
in the calculation of the fabric tensor.


\subsection{Restrictions}
\label{\detokenize{compute_fabric:restrictions}}
\sphinxAtStartPar
This fix is part of the GRANULAR package.  It is only enabled if LAMMPS
was built with that package.  See the {\hyperref[\detokenize{Build_package::doc}]{\sphinxcrossref{\DUrole{doc}{Build package}}}}
doc page for more info.

\sphinxAtStartPar
Currently, compute \sphinxstyleemphasis{fabric} does not support pair styles
with many\sphinxhyphen{}body interactions.  It also does not
support models with long\sphinxhyphen{}range Coulombic or dispersion forces,
i.e. the kspace\_style command in LAMMPS.  It also does not support the
following fixes which add rigid\sphinxhyphen{}body constraints: {\hyperref[\detokenize{fix_shake::doc}]{\sphinxcrossref{\DUrole{doc}{fix shake}}}}, {\hyperref[\detokenize{fix_shake::doc}]{\sphinxcrossref{\DUrole{doc}{fix rattle}}}}, {\hyperref[\detokenize{fix_rigid::doc}]{\sphinxcrossref{\DUrole{doc}{fix rigid}}}}, {\hyperref[\detokenize{fix_rigid::doc}]{\sphinxcrossref{\DUrole{doc}{fix rigid/small}}}}. It does not support
granular pair styles that extend beyond the contact of atomic radii
(e.g., JKR and DMT).


\subsection{Related commands}
\label{\detokenize{compute_fabric:related-commands}}
\sphinxAtStartPar
none


\subsection{Default}
\label{\detokenize{compute_fabric:default}}
\sphinxAtStartPar
none


\bigskip\hrule\bigskip

\phantomsection\label{\detokenize{compute_fabric:ouadfel}}
\sphinxAtStartPar
\sphinxstylestrong{(Ouadfel)} Ouadfel and Rothenburg
“Stress\sphinxhyphen{}force\sphinxhyphen{}fabric relationship for assemblies of ellipsoids”,
Mechanics of Materials (2001). (\sphinxhref{https://doi.org/10.1016/S0167-6636(00)00057-0}{link to paper})

\sphinxstepscope

\index{compute fep@\spxentry{compute fep}}\ignorespaces 

\section{compute fep command}
\label{\detokenize{compute_fep:compute-fep-command}}\label{\detokenize{compute_fep:index-0}}\label{\detokenize{compute_fep::doc}}

\subsection{Syntax}
\label{\detokenize{compute_fep:syntax}}
\begin{sphinxVerbatim}[commandchars=\\\{\}]
\PYG{k}{compute }\PYG{n+nv+nvIdentifier}{ID}\PYG{+w}{ }\PYG{n+nv+nvIdentifier}{group\PYGZhy{}ID}\PYG{+w}{ }\PYG{n}{fep}\PYG{+w}{ }\PYG{n}{temp}\PYG{+w}{ }\PYG{n}{attribute}\PYG{+w}{ }\PYG{n}{args}\PYG{+w}{ }\PYG{n}{...}\PYG{+w}{ }\PYG{n}{keyword}\PYG{+w}{ }\PYG{n}{value}\PYG{+w}{ }\PYG{n}{...}
\end{sphinxVerbatim}
\begin{itemize}
\item {} 
\sphinxAtStartPar
ID, group\sphinxhyphen{}ID are documented in the {\hyperref[\detokenize{compute::doc}]{\sphinxcrossref{\DUrole{doc}{compute}}}} command

\item {} 
\sphinxAtStartPar
fep = name of this compute command

\item {} 
\sphinxAtStartPar
temp = external temperature (as specified for constant\sphinxhyphen{}temperature run)

\item {} 
\sphinxAtStartPar
one or more attributes with args may be appended

\item {} 
\sphinxAtStartPar
attribute = \sphinxstyleemphasis{pair} or \sphinxstyleemphasis{atom}
\begin{sphinxalltt}
\sphinxstyleemphasis{pair} args = pstyle pparam I J v\_delta
  pstyle = pair style name (e.g., \sphinxstyleemphasis{lj/cut})
  pparam = parameter to perturb
  I,J = type pair(s) to set parameter for
  v\_delta = variable with perturbation to apply (in the units of the parameter)
\sphinxstyleemphasis{atom} args = aparam I v\_delta
  aparam = \sphinxstyleemphasis{charge} = parameter to perturb
  I = type to set parameter for
  v\_delta = variable with perturbation to apply (in the units of the parameter)
\end{sphinxalltt}

\item {} 
\sphinxAtStartPar
zero or more keyword/value pairs may be appended

\item {} 
\sphinxAtStartPar
keyword = \sphinxstyleemphasis{tail} or \sphinxstyleemphasis{volume}
\begin{sphinxalltt}
\sphinxstyleemphasis{tail} value = \sphinxstyleemphasis{no} or \sphinxstyleemphasis{yes}
  \sphinxstyleemphasis{no} = ignore tail correction to pair energies (usually small in fep)
  \sphinxstyleemphasis{yes} = include tail correction to pair energies
\sphinxstyleemphasis{volume} value = \sphinxstyleemphasis{no} or \sphinxstyleemphasis{yes}
  \sphinxstyleemphasis{no} = ignore volume changes (e.g., in \sphinxstyleemphasis{NVE} or \sphinxstyleemphasis{NVT} trajectories)
  \sphinxstyleemphasis{yes} = include volume changes (e.g., in \sphinxstyleemphasis{NPT} trajectories)
\end{sphinxalltt}

\end{itemize}


\subsection{Examples}
\label{\detokenize{compute_fep:examples}}
\begin{sphinxVerbatim}[commandchars=\\\{\}]
\PYG{k}{compute }\PYG{n+nv+nvIdentifier}{1}\PYG{+w}{ }\PYG{n+nv+nvIdentifier}{all}\PYG{+w}{ }\PYG{n}{fep}\PYG{+w}{ }\PYG{l+m}{298}\PYG{+w}{ }\PYG{n}{pair}\PYG{+w}{ }\PYG{n}{lj}\PYG{o}{/}\PYG{n}{cut}\PYG{+w}{ }\PYG{n}{epsilon}\PYG{+w}{ }\PYG{l+m}{1}\PYG{+w}{ }\PYG{o}{*}\PYG{+w}{ }\PYG{n}{v\PYGZus{}delta}\PYG{+w}{ }\PYG{n}{pair}\PYG{+w}{ }\PYG{n}{lj}\PYG{o}{/}\PYG{n}{cut}\PYG{+w}{ }\PYG{n}{sigma}\PYG{+w}{ }\PYG{l+m}{1}\PYG{+w}{ }\PYG{o}{*}\PYG{+w}{ }\PYG{n}{v\PYGZus{}delta}\PYG{+w}{ }\PYG{n}{volume}\PYG{+w}{ }\PYG{n}{yes}
\PYG{k}{compute }\PYG{n+nv+nvIdentifier}{1}\PYG{+w}{ }\PYG{n+nv+nvIdentifier}{all}\PYG{+w}{ }\PYG{n}{fep}\PYG{+w}{ }\PYG{l+m}{300}\PYG{+w}{ }\PYG{n}{atom}\PYG{+w}{ }\PYG{n}{charge}\PYG{+w}{ }\PYG{l+m}{2}\PYG{+w}{ }\PYG{n}{v\PYGZus{}delta}
\end{sphinxVerbatim}

\sphinxAtStartPar
Example input scripts available: examples/PACKAGES/fep


\subsection{Description}
\label{\detokenize{compute_fep:description}}
\sphinxAtStartPar
Apply a perturbation to parameters of the interaction potential and
recalculate the pair potential energy without changing the atomic
coordinates from those of the reference, unperturbed system. This
compute can be used to calculate free energy differences using several
methods, such as free\sphinxhyphen{}energy perturbation (FEP), finite\sphinxhyphen{}difference
thermodynamic integration (FDTI) or Bennet’s acceptance ratio method
(BAR).

\sphinxAtStartPar
The potential energy of the system is decomposed in three terms: a
background term corresponding to interaction sites whose parameters
remain constant, a reference term \(U_0\) corresponding to the
initial interactions of the atoms that will undergo perturbation, and
a term \(U_1\) corresponding to the final interactions of
these atoms:
\begin{equation*}
\begin{split}U(\lambda) = U_{\mathrm{bg}} + U_1(\lambda) + U_0(\lambda)\end{split}
\end{equation*}
\sphinxAtStartPar
A coupling parameter \(\lambda\) varying from 0 to 1 connects the
reference and perturbed systems:
\begin{equation*}
\begin{split}\lambda &= 0 \quad\Rightarrow\quad U = U_{\mathrm{bg}} + U_0 \\
\lambda &= 1 \quad\Rightarrow\quad U = U_{\mathrm{bg}} + U_1\end{split}
\end{equation*}
\sphinxAtStartPar
It is possible but not necessary that the coupling parameter (or a
function thereof) appears as a multiplication factor of the potential
energy. Therefore, this compute can apply perturbations to interaction
parameters that are not directly proportional to the potential energy
(e.g., \(\sigma\) in Lennard\sphinxhyphen{}Jones potentials).

\sphinxAtStartPar
This command can be combined with {\hyperref[\detokenize{fix_adapt::doc}]{\sphinxcrossref{\DUrole{doc}{fix adapt}}}} to
perform multistage free\sphinxhyphen{}energy perturbation calculations along
stepwise alchemical transformations during a simulation run:
\begin{equation*}
\begin{split}\Delta_0^1 A = \sum_{i=0}^{n-1} \Delta_{\lambda_i}^{\lambda_{i+1}} A = - k_B T
\sum_{i=0}^{n-1} \ln \left< \exp \left( - \frac{U(\lambda_{i+1}) -
U(\lambda_i)}{k_B T} \right) \right>_{\lambda_i}\end{split}
\end{equation*}
\sphinxAtStartPar
This compute is suitable for the finite\sphinxhyphen{}difference thermodynamic
integration (FDTI) method {\hyperref[\detokenize{compute_fep:mezei}]{\sphinxcrossref{\DUrole{std,std-ref}{(Mezei)}}}}, which is based on an
evaluation of the numerical derivative of the free energy by a
perturbation method using a very small \(\delta\):
\begin{equation*}
\begin{split}\Delta_0^1 A = \int_{\lambda=0}^{\lambda=1} \left( \frac{\partial
A(\lambda)}{\partial\lambda} \right)_\lambda \mathrm{d}\lambda \approx
\sum_{i=0}^{n-1} w_i \frac{A(\lambda_{i} + \delta) - A(\lambda_i)}{\delta}\end{split}
\end{equation*}
\sphinxAtStartPar
where \(w_i\) are weights of a numerical quadrature. The
{\hyperref[\detokenize{fix_adapt::doc}]{\sphinxcrossref{\DUrole{doc}{fix adapt}}}} command can be used to define the stages of
\(\lambda\) at which the derivative is calculated and averaged.

\sphinxAtStartPar
The compute fep calculates the exponential Boltzmann term and also the
potential energy difference \(U_1 -U_0\). By
choosing a very small perturbation \(\delta\) the thermodynamic
integration method can be implemented using a numerical evaluation of
the derivative of the potential energy with respect to \(\lambda\):
\begin{equation*}
\begin{split}\Delta_0^1 A = \int_{\lambda=0}^{\lambda=1} \left< \frac{\partial
U(\lambda)}{\partial\lambda} \right>_\lambda \mathrm{d}\lambda \approx
\sum_{i=0}^{n-1} w_i \left< \frac{U(\lambda_{i} + \delta) -
U(\lambda_i)}{\delta} \right>_{\lambda_i}\end{split}
\end{equation*}
\sphinxAtStartPar
Another technique to calculate free energy differences is the
acceptance ratio method {\hyperref[\detokenize{compute_fep:bennet}]{\sphinxcrossref{\DUrole{std,std-ref}{(Bennet)}}}}, which can be implemented
by calculating the potential energy differences with \(\delta = 1.0\) on
both the forward and reverse routes:
\begin{equation*}
\begin{split}\left< \frac{1}{1 + \exp\left[\left(U_1 - U_0 - \Delta_0^1A \right) /k_B T
\right]} \right>_0 = \left< \frac{1}{1 + \exp\left[\left(U_0 - U_1 +
\Delta_0^1A \right) /k_B T \right]} \right>_1\end{split}
\end{equation*}
\sphinxAtStartPar
The value of the free energy difference is determined by numerical
root finding to establish the equality.

\sphinxAtStartPar
Concerning the choice of how the atomic parameters are perturbed in
order to setup an alchemical transformation route, several strategies
are available, such as single\sphinxhyphen{}topology or double\sphinxhyphen{}topology strategies
{\hyperref[\detokenize{compute_fep:pearlman}]{\sphinxcrossref{\DUrole{std,std-ref}{(Pearlman)}}}}. The latter does not require modification of
bond lengths, angles or other internal coordinates.

\sphinxAtStartPar
NOTES: This compute command does not take kinetic energy into account,
therefore the masses of the particles should not be modified between
the reference and perturbed states, or along the alchemical
transformation route.  This compute command does not change bond
lengths or other internal coordinates {\hyperref[\detokenize{compute_fep:boreschkarplus}]{\sphinxcrossref{\DUrole{std,std-ref}{(Boresch, Karplus)}}}}.


\bigskip\hrule\bigskip


\sphinxAtStartPar
The \sphinxstyleemphasis{pair} attribute enables various parameters of potentials defined
by the {\hyperref[\detokenize{pair_style::doc}]{\sphinxcrossref{\DUrole{doc}{pair\_style}}}} and {\hyperref[\detokenize{pair_coeff::doc}]{\sphinxcrossref{\DUrole{doc}{pair\_coeff}}}}
commands to be changed, if the pair style supports it.

\sphinxAtStartPar
The \sphinxstyleemphasis{pstyle} argument is the name of the pair style. For example,
\sphinxstyleemphasis{pstyle} could be specified as “lj/cut”.  The \sphinxstyleemphasis{pparam} argument is the
name of the parameter to change.  This is a list of
pair styles and parameters that can be used with this compute.  See
the doc pages for individual pair styles and their energy formulas for
the meaning of these parameters:


\begin{savenotes}\sphinxattablestart
\sphinxthistablewithglobalstyle
\centering
\begin{tabulary}{\linewidth}[t]{TTT}
\sphinxtoprule
\sphinxtableatstartofbodyhook
\sphinxAtStartPar
{\hyperref[\detokenize{pair_born::doc}]{\sphinxcrossref{\DUrole{doc}{born}}}}
&
\sphinxAtStartPar
a,b,c
&
\sphinxAtStartPar
type pairs
\\
\sphinxhline
\sphinxAtStartPar
{\hyperref[\detokenize{pair_buck::doc}]{\sphinxcrossref{\DUrole{doc}{buck, buck/coul/cut, buck/coul/long, buck/coul/msm}}}}
&
\sphinxAtStartPar
a,c
&
\sphinxAtStartPar
type pairs
\\
\sphinxhline
\sphinxAtStartPar
{\hyperref[\detokenize{pair_mdf::doc}]{\sphinxcrossref{\DUrole{doc}{buck/mdf}}}}
&
\sphinxAtStartPar
a,c
&
\sphinxAtStartPar
type pairs
\\
\sphinxhline
\sphinxAtStartPar
{\hyperref[\detokenize{pair_coul::doc}]{\sphinxcrossref{\DUrole{doc}{coul/cut}}}}
&
\sphinxAtStartPar
scale
&
\sphinxAtStartPar
type pairs
\\
\sphinxhline
\sphinxAtStartPar
{\hyperref[\detokenize{pair_fep_soft::doc}]{\sphinxcrossref{\DUrole{doc}{coul/cut/soft}}}}
&
\sphinxAtStartPar
lambda
&
\sphinxAtStartPar
type pairs
\\
\sphinxhline
\sphinxAtStartPar
{\hyperref[\detokenize{pair_coul::doc}]{\sphinxcrossref{\DUrole{doc}{coul/long, coul/msm}}}}
&
\sphinxAtStartPar
scale
&
\sphinxAtStartPar
type pairs
\\
\sphinxhline
\sphinxAtStartPar
{\hyperref[\detokenize{pair_fep_soft::doc}]{\sphinxcrossref{\DUrole{doc}{coul/long/soft}}}}
&
\sphinxAtStartPar
scale, lambda
&
\sphinxAtStartPar
type pairs
\\
\sphinxhline
\sphinxAtStartPar
{\hyperref[\detokenize{pair_eam::doc}]{\sphinxcrossref{\DUrole{doc}{eam}}}}
&
\sphinxAtStartPar
scale
&
\sphinxAtStartPar
type pairs
\\
\sphinxhline
\sphinxAtStartPar
{\hyperref[\detokenize{pair_gauss::doc}]{\sphinxcrossref{\DUrole{doc}{gauss}}}}
&
\sphinxAtStartPar
a
&
\sphinxAtStartPar
type pairs
\\
\sphinxhline
\sphinxAtStartPar
{\hyperref[\detokenize{pair_mdf::doc}]{\sphinxcrossref{\DUrole{doc}{lennard/mdf}}}}
&
\sphinxAtStartPar
a,b
&
\sphinxAtStartPar
type pairs
\\
\sphinxhline
\sphinxAtStartPar
{\hyperref[\detokenize{pair_class2::doc}]{\sphinxcrossref{\DUrole{doc}{lj/class2}}}}
&
\sphinxAtStartPar
epsilon,sigma
&
\sphinxAtStartPar
type pairs
\\
\sphinxhline
\sphinxAtStartPar
{\hyperref[\detokenize{pair_class2::doc}]{\sphinxcrossref{\DUrole{doc}{lj/class2/coul/cut, lj/class2/coul/long}}}}
&
\sphinxAtStartPar
epsilon,sigma
&
\sphinxAtStartPar
type pairs
\\
\sphinxhline
\sphinxAtStartPar
{\hyperref[\detokenize{pair_lj::doc}]{\sphinxcrossref{\DUrole{doc}{lj/cut}}}}
&
\sphinxAtStartPar
epsilon,sigma
&
\sphinxAtStartPar
type pairs
\\
\sphinxhline
\sphinxAtStartPar
{\hyperref[\detokenize{pair_fep_soft::doc}]{\sphinxcrossref{\DUrole{doc}{lj/cut/soft}}}}
&
\sphinxAtStartPar
epsilon,sigma,lambda
&
\sphinxAtStartPar
type pairs
\\
\sphinxhline
\sphinxAtStartPar
{\hyperref[\detokenize{pair_lj_cut_coul::doc}]{\sphinxcrossref{\DUrole{doc}{lj/cut/coul/cut, lj/cut/coul/long, lj/cut/coul/msm}}}}
&
\sphinxAtStartPar
epsilon,sigma
&
\sphinxAtStartPar
type pairs
\\
\sphinxhline
\sphinxAtStartPar
{\hyperref[\detokenize{pair_fep_soft::doc}]{\sphinxcrossref{\DUrole{doc}{lj/cut/coul/cut/soft, lj/cut/coul/long/soft}}}}
&
\sphinxAtStartPar
epsilon,sigma,lambda
&
\sphinxAtStartPar
type pairs
\\
\sphinxhline
\sphinxAtStartPar
{\hyperref[\detokenize{pair_lj_cut_tip4p::doc}]{\sphinxcrossref{\DUrole{doc}{lj/cut/tip4p/cut, lj/cut/tip4p/long}}}}
&
\sphinxAtStartPar
epsilon,sigma
&
\sphinxAtStartPar
type pairs
\\
\sphinxhline
\sphinxAtStartPar
{\hyperref[\detokenize{pair_fep_soft::doc}]{\sphinxcrossref{\DUrole{doc}{lj/cut/tip4p/long/soft}}}}
&
\sphinxAtStartPar
epsilon,sigma,lambda
&
\sphinxAtStartPar
type pairs
\\
\sphinxhline
\sphinxAtStartPar
{\hyperref[\detokenize{pair_lj_expand::doc}]{\sphinxcrossref{\DUrole{doc}{lj/expand}}}}
&
\sphinxAtStartPar
epsilon,sigma,delta
&
\sphinxAtStartPar
type pairs
\\
\sphinxhline
\sphinxAtStartPar
{\hyperref[\detokenize{pair_mdf::doc}]{\sphinxcrossref{\DUrole{doc}{lj/mdf}}}}
&
\sphinxAtStartPar
epsilon,sigma
&
\sphinxAtStartPar
type pairs
\\
\sphinxhline
\sphinxAtStartPar
{\hyperref[\detokenize{pair_dipole::doc}]{\sphinxcrossref{\DUrole{doc}{lj/sf/dipole/sf}}}}
&
\sphinxAtStartPar
epsilon,sigma,scale
&
\sphinxAtStartPar
type pairs
\\
\sphinxhline
\sphinxAtStartPar
{\hyperref[\detokenize{pair_mie::doc}]{\sphinxcrossref{\DUrole{doc}{mie/cut}}}}
&
\sphinxAtStartPar
epsilon,sigma,gamR,gamA
&
\sphinxAtStartPar
type pairs
\\
\sphinxhline
\sphinxAtStartPar
{\hyperref[\detokenize{pair_morse::doc}]{\sphinxcrossref{\DUrole{doc}{morse, morse/smooth/linear}}}}
&
\sphinxAtStartPar
d0,r0,alpha
&
\sphinxAtStartPar
type pairs
\\
\sphinxhline
\sphinxAtStartPar
{\hyperref[\detokenize{pair_morse::doc}]{\sphinxcrossref{\DUrole{doc}{morse/soft}}}}
&
\sphinxAtStartPar
d0,r0,alpha,lambda
&
\sphinxAtStartPar
type pairs
\\
\sphinxhline
\sphinxAtStartPar
{\hyperref[\detokenize{pair_nm::doc}]{\sphinxcrossref{\DUrole{doc}{nm/cut}}}}
&
\sphinxAtStartPar
e0,r0,nn,mm
&
\sphinxAtStartPar
type pairs
\\
\sphinxhline
\sphinxAtStartPar
{\hyperref[\detokenize{pair_nm::doc}]{\sphinxcrossref{\DUrole{doc}{nm/cut/coul/cut, nm/cut/coul/long}}}}
&
\sphinxAtStartPar
e0,r0,nn,mm
&
\sphinxAtStartPar
type pairs
\\
\sphinxhline
\sphinxAtStartPar
{\hyperref[\detokenize{pair_ufm::doc}]{\sphinxcrossref{\DUrole{doc}{ufm}}}}
&
\sphinxAtStartPar
epsilon,sigma,scale
&
\sphinxAtStartPar
type pairs
\\
\sphinxhline
\sphinxAtStartPar
{\hyperref[\detokenize{pair_soft::doc}]{\sphinxcrossref{\DUrole{doc}{soft}}}}
&
\sphinxAtStartPar
a
&
\sphinxAtStartPar
type pairs
\\
\sphinxbottomrule
\end{tabulary}
\sphinxtableafterendhook\par
\sphinxattableend\end{savenotes}

\sphinxAtStartPar
Note that it is easy to add new potentials and their parameters to
this list.  All it typically takes is adding an extract() method to
the pair\_*.cpp file associated with the potential.

\sphinxAtStartPar
Similar to the {\hyperref[\detokenize{pair_coeff::doc}]{\sphinxcrossref{\DUrole{doc}{pair\_coeff}}}} command, I and J can be
specified in one of two ways.  Explicit numeric values can be used for
each, as in the first example above.  I \(\le\) J is required.  LAMMPS sets
the coefficients for the symmetric J,I interaction to the same
values. A wild\sphinxhyphen{}card asterisk can be used in place of or in conjunction
with the I,J arguments to set the coefficients for multiple pairs of
atom types.  This takes the form “*” or “*n” or “m*” or “m*n”.  If
\(N\) is the number of atom types, then an asterisk with no numeric values
means all types from 1 to \(N\).   A leading asterisk means all types from
1 to n (inclusive).  A trailing asterisk means all types from m to N
(inclusive).  A middle asterisk means all types from m to n
(inclusive).  Note that only type pairs with I \(\le\) J are considered; if
asterisks imply type pairs where J \(<\) I, they are ignored.

\sphinxAtStartPar
If {\hyperref[\detokenize{pair_hybrid::doc}]{\sphinxcrossref{\DUrole{doc}{pair\_style hybrid or hybrid/overlay}}}} is being
used, then the \sphinxstyleemphasis{pstyle} will be a sub\sphinxhyphen{}style name.  You must specify
I,J arguments that correspond to type pair values defined (via the
{\hyperref[\detokenize{pair_coeff::doc}]{\sphinxcrossref{\DUrole{doc}{pair\_coeff}}}} command) for that sub\sphinxhyphen{}style.

\sphinxAtStartPar
The \sphinxstyleemphasis{v\_name} argument for keyword \sphinxstyleemphasis{pair} is the name of an
{\hyperref[\detokenize{variable::doc}]{\sphinxcrossref{\DUrole{doc}{equal\sphinxhyphen{}style variable}}}} which will be evaluated each time
this compute is invoked.  It should be specified as v\_name, where name
is the variable name.


\bigskip\hrule\bigskip


\sphinxAtStartPar
The \sphinxstyleemphasis{atom} attribute enables atom properties to be changed.  The
\sphinxstyleemphasis{aparam} argument is the name of the parameter to change.  This is the
current list of atom parameters that can be used with this compute:
\begin{itemize}
\item {} 
\sphinxAtStartPar
charge = charge on particle

\end{itemize}

\sphinxAtStartPar
The \sphinxstyleemphasis{v\_name} argument for keyword \sphinxstyleemphasis{pair} is the name of an
{\hyperref[\detokenize{variable::doc}]{\sphinxcrossref{\DUrole{doc}{equal\sphinxhyphen{}style variable}}}} which will be evaluated each time
this compute is invoked.  It should be specified as v\_name, where name
is the variable name.


\bigskip\hrule\bigskip


\sphinxAtStartPar
The \sphinxstyleemphasis{tail} keyword controls the calculation of the tail correction to
“van der Waals” pair energies beyond the cutoff, if this has been
activated via the {\hyperref[\detokenize{pair_modify::doc}]{\sphinxcrossref{\DUrole{doc}{pair\_modify}}}} command. If the
perturbation is small, the tail contribution to the energy difference
between the reference and perturbed systems should be negligible.

\sphinxAtStartPar
If the keyword \sphinxstyleemphasis{volume} = \sphinxstyleemphasis{yes}, then the Boltzmann term is multiplied
by the volume so that correct ensemble averaging can be performed over
trajectories during which the volume fluctuates or changes {\hyperref[\detokenize{compute_fep:allentildesley}]{\sphinxcrossref{\DUrole{std,std-ref}{(Allen and Tildesley)}}}}:
\begin{equation*}
\begin{split}\Delta_0^1 A = - k_B T \sum_{i=0}^{n-1} \ln \frac{\left< V \exp \left( -
\frac{U(\lambda_{i+1}) - U(\lambda_i)}{k_B T} \right)
\right>_{\lambda_i}}{\left< V \right>_{\lambda_i}}\end{split}
\end{equation*}

\bigskip\hrule\bigskip



\subsection{Output info}
\label{\detokenize{compute_fep:output-info}}
\sphinxAtStartPar
This compute calculates a global vector of length 3 which contains the
energy difference ( \(U_1-U_0\) ) as c\_ID{[}1{]}, the
Boltzmann factor \(\exp(-(U_1-U_0)/k_B T)\), or
\(V \exp(-(U_1-U_0)/k_B T)\), as c\_ID{[}2{]} and the
volume of the simulation box \(V\) as c\_ID{[}3{]}. \(U_1\) is the
pair potential energy obtained with the perturbed parameters and
\(U_0\) is the pair potential energy obtained with the
unperturbed parameters. The energies include kspace terms if these
are used in the simulation.

\sphinxAtStartPar
These output results can be used by any command that uses a global
scalar or vector from a compute as input.  See the {\hyperref[\detokenize{Howto_output::doc}]{\sphinxcrossref{\DUrole{doc}{Howto output}}}} page for an overview of LAMMPS output
options. For example, the computed values can be averaged using {\hyperref[\detokenize{fix_ave_time::doc}]{\sphinxcrossref{\DUrole{doc}{fix ave/time}}}}.

\sphinxAtStartPar
The values calculated by this compute are “extensive”.


\subsection{Restrictions}
\label{\detokenize{compute_fep:restrictions}}
\sphinxAtStartPar
This compute is distributed as the FEP package.  It is only
enabled if LAMMPS was built with that package.  See the {\hyperref[\detokenize{Build_package::doc}]{\sphinxcrossref{\DUrole{doc}{Build package}}}} page for more info.


\subsection{Related commands}
\label{\detokenize{compute_fep:related-commands}}
\sphinxAtStartPar
{\hyperref[\detokenize{fix_adapt_fep::doc}]{\sphinxcrossref{\DUrole{doc}{fix adapt/fep}}}}, {\hyperref[\detokenize{fix_ave_time::doc}]{\sphinxcrossref{\DUrole{doc}{fix ave/time}}}},
{\hyperref[\detokenize{pair_fep_soft::doc}]{\sphinxcrossref{\DUrole{doc}{pair\_style …/soft}}}}


\subsection{Default}
\label{\detokenize{compute_fep:default}}
\sphinxAtStartPar
The option defaults are \sphinxstyleemphasis{tail} = \sphinxstyleemphasis{no}, \sphinxstyleemphasis{volume} = \sphinxstyleemphasis{no}.


\bigskip\hrule\bigskip

\phantomsection\label{\detokenize{compute_fep:pearlman}}
\sphinxAtStartPar
\sphinxstylestrong{(Pearlman)} Pearlman, J Chem Phys, 98, 1487 (1994)

\phantomsection\label{\detokenize{compute_fep:mezei}}
\sphinxAtStartPar
\sphinxstylestrong{(Mezei)} Mezei, J Chem Phys, 86, 7084 (1987)

\phantomsection\label{\detokenize{compute_fep:bennet}}
\sphinxAtStartPar
\sphinxstylestrong{(Bennet)} Bennet, J Comput Phys, 22, 245 (1976)

\phantomsection\label{\detokenize{compute_fep:boreschkarplus}}
\sphinxAtStartPar
\sphinxstylestrong{(BoreschKarplus)} Boresch and Karplus, J Phys Chem A, 103, 103 (1999)

\phantomsection\label{\detokenize{compute_fep:allentildesley}}
\sphinxAtStartPar
\sphinxstylestrong{(AllenTildesley)} Allen and Tildesley, Computer Simulation of
Liquids, Oxford University Press (1987)

\sphinxstepscope

\index{compute fep/ta@\spxentry{compute fep/ta}}\ignorespaces 

\section{compute fep/ta command}
\label{\detokenize{compute_fep_ta:compute-fep-ta-command}}\label{\detokenize{compute_fep_ta:index-0}}\label{\detokenize{compute_fep_ta::doc}}

\subsection{Syntax}
\label{\detokenize{compute_fep_ta:syntax}}
\begin{sphinxVerbatim}[commandchars=\\\{\}]
\PYG{k}{compute }\PYG{n+nv+nvIdentifier}{ID}\PYG{+w}{ }\PYG{n+nv+nvIdentifier}{group\PYGZhy{}ID}\PYG{+w}{ }\PYG{n}{fep}\PYG{o}{/}\PYG{n}{ta}\PYG{+w}{ }\PYG{n}{temp}\PYG{+w}{ }\PYG{n}{plane}\PYG{+w}{ }\PYG{n}{scale\PYGZus{}factor}\PYG{+w}{ }\PYG{n}{keyword}\PYG{+w}{ }\PYG{n}{value}\PYG{+w}{ }\PYG{n}{...}
\end{sphinxVerbatim}
\begin{itemize}
\item {} 
\sphinxAtStartPar
ID, group\sphinxhyphen{}ID are documented in the {\hyperref[\detokenize{compute::doc}]{\sphinxcrossref{\DUrole{doc}{compute}}}} command

\item {} 
\sphinxAtStartPar
fep/ta = name of this compute command

\item {} 
\sphinxAtStartPar
temp = external temperature (as specified for constant\sphinxhyphen{}temperature run)

\item {} 
\sphinxAtStartPar
plane = \sphinxstyleemphasis{xy} or \sphinxstyleemphasis{xz} or \sphinxstyleemphasis{yz}

\item {} 
\sphinxAtStartPar
scale\_factor = multiplicative factor for change in plane area

\item {} 
\sphinxAtStartPar
zero or more keyword/value pairs may be appended

\item {} 
\sphinxAtStartPar
keyword = \sphinxstyleemphasis{tail}
\begin{sphinxalltt}
\sphinxstyleemphasis{tail} value = \sphinxstyleemphasis{no} or \sphinxstyleemphasis{yes}
  \sphinxstyleemphasis{no} = ignore tail correction to pair energies (usually small in fep)
  \sphinxstyleemphasis{yes} = include tail correction to pair energies
\end{sphinxalltt}

\end{itemize}


\subsection{Examples}
\label{\detokenize{compute_fep_ta:examples}}
\begin{sphinxVerbatim}[commandchars=\\\{\}]
\PYG{k}{compute }\PYG{n+nv+nvIdentifier}{1}\PYG{+w}{ }\PYG{n+nv+nvIdentifier}{all}\PYG{+w}{ }\PYG{n}{fep}\PYG{o}{/}\PYG{n}{ta}\PYG{+w}{ }\PYG{l+m}{298}\PYG{+w}{ }\PYG{n}{xy}\PYG{+w}{ }\PYG{l+m}{1.0005}
\end{sphinxVerbatim}


\subsection{Description}
\label{\detokenize{compute_fep_ta:description}}
\sphinxAtStartPar
\DUrole{versionmodified,added}{Added in version 4May2022.}

\sphinxAtStartPar
Define a computation that calculates the change in the free energy due
to a test\sphinxhyphen{}area (TA) perturbation {\hyperref[\detokenize{compute_fep_ta:gloor}]{\sphinxcrossref{\DUrole{std,std-ref}{(Gloor)}}}}. The test\sphinxhyphen{}area
approach can be used to determine the interfacial tension of the system
in a single simulation:
\begin{equation*}
\begin{split}\gamma = \lim_{\Delta \mathcal{A} \to 0} \left( \frac{\Delta A_{0 \to 1 }}{\Delta \mathcal{A}}\right)_{N,V,T}
= - \frac{k_B T}{\Delta \mathcal{A}} \ln \left\langle \exp\left(\frac{-(U_1 - U_0)}{k_B T}\right) \right\rangle_0\end{split}
\end{equation*}
\sphinxAtStartPar
During the perturbation, both axes of \sphinxstyleemphasis{plane} are scaled by multiplying
\(\sqrt{\mathrm{scale\_factor}}\), while the other axis divided by
\(\mathrm{scale\_factor}\) such that the overall volume of the system is
maintained.

\sphinxAtStartPar
The \sphinxstyleemphasis{tail} keyword controls the calculation of the tail correction to
“van der Waals” pair energies beyond the cutoff, if this has been
activated via the {\hyperref[\detokenize{pair_modify::doc}]{\sphinxcrossref{\DUrole{doc}{pair\_modify}}}} command. If the
perturbation is small, the tail contribution to the energy difference
between the reference and perturbed systems should be negligible.


\bigskip\hrule\bigskip



\subsection{Output info}
\label{\detokenize{compute_fep_ta:output-info}}
\sphinxAtStartPar
This compute calculates a global vector of length 3 which contains the
energy difference \((U_1-U_0)\) as c\_ID{[}1{]}, the Boltzmann factor
\(\exp\bigl(-(U_1-U_0)/k_B T\bigr)\), as c\_ID{[}2{]} and the change in the \sphinxstyleemphasis{plane}
area \(\Delta \mathcal{A}\) as c\_ID{[}3{]}. \(U_1\) is the potential
energy of the perturbed state and \(U_0\) is the potential energy of
the reference state.  The energies include kspace terms if these are
used in the simulation.

\sphinxAtStartPar
These output results can be used by any command that uses a global
scalar or vector from a compute as input.  See the {\hyperref[\detokenize{Howto_output::doc}]{\sphinxcrossref{\DUrole{doc}{Howto output}}}} page for an overview of LAMMPS output options. For
example, the computed values can be averaged using {\hyperref[\detokenize{fix_ave_time::doc}]{\sphinxcrossref{\DUrole{doc}{fix ave/time}}}}.


\subsection{Restrictions}
\label{\detokenize{compute_fep_ta:restrictions}}
\sphinxAtStartPar
Constraints, like fix shake, may lead to incorrect values for energy difference.

\sphinxAtStartPar
This compute is distributed as the FEP package.  It is only enabled if
LAMMPS was built with that package.  See the {\hyperref[\detokenize{Build_package::doc}]{\sphinxcrossref{\DUrole{doc}{Build package}}}} page for more info.


\subsection{Related commands}
\label{\detokenize{compute_fep_ta:related-commands}}
\sphinxAtStartPar
{\hyperref[\detokenize{compute_fep::doc}]{\sphinxcrossref{\DUrole{doc}{compute fep}}}}


\subsection{Default}
\label{\detokenize{compute_fep_ta:default}}
\sphinxAtStartPar
The option defaults are \sphinxstyleemphasis{tail} = \sphinxstyleemphasis{no}.


\bigskip\hrule\bigskip

\phantomsection\label{\detokenize{compute_fep_ta:gloor}}
\sphinxAtStartPar
\sphinxstylestrong{(Gloor)} Gloor, J Chem Phys, 123, 134703 (2005)

\sphinxstepscope

\index{compute global/atom@\spxentry{compute global/atom}}\ignorespaces 

\section{compute global/atom command}
\label{\detokenize{compute_global_atom:compute-global-atom-command}}\label{\detokenize{compute_global_atom:index-0}}\label{\detokenize{compute_global_atom::doc}}

\subsection{Syntax}
\label{\detokenize{compute_global_atom:syntax}}
\begin{sphinxVerbatim}[commandchars=\\\{\}]
\PYG{k}{compute }\PYG{n+nv+nvIdentifier}{ID}\PYG{+w}{ }\PYG{n+nv+nvIdentifier}{group\PYGZhy{}ID}\PYG{+w}{ }\PYG{n}{style}\PYG{+w}{ }\PYG{n}{index}\PYG{+w}{ }\PYG{n}{input1}\PYG{+w}{ }\PYG{n}{input2}\PYG{+w}{ }\PYG{n}{...}
\end{sphinxVerbatim}
\begin{itemize}
\item {} 
\sphinxAtStartPar
ID, group\sphinxhyphen{}ID are documented in {\hyperref[\detokenize{compute::doc}]{\sphinxcrossref{\DUrole{doc}{compute}}}} command

\item {} 
\sphinxAtStartPar
global/atom = style name of this compute command

\item {} 
\sphinxAtStartPar
index = c\_ID, c\_ID{[}N{]}, f\_ID, f\_ID{[}N{]}, v\_name

\begin{sphinxVerbatim}[commandchars=\\\{\}]
c\PYGZus{}ID = per\PYGZhy{}atom vector calculated by a compute with ID
c\PYGZus{}ID[I] = Ith column of per\PYGZhy{}atom array calculated by a compute with ID
f\PYGZus{}ID = per\PYGZhy{}atom vector calculated by a fix with ID
f\PYGZus{}ID[I] = Ith column of per\PYGZhy{}atom array calculated by a fix with ID
v\PYGZus{}name = per\PYGZhy{}atom vector calculated by an atom\PYGZhy{}style variable with name
\end{sphinxVerbatim}

\item {} 
\sphinxAtStartPar
one or more inputs can be listed

\item {} 
\sphinxAtStartPar
input = c\_ID, c\_ID{[}N{]}, f\_ID, f\_ID{[}N{]}, v\_name

\begin{sphinxVerbatim}[commandchars=\\\{\}]
c\PYGZus{}ID = global vector calculated by a compute with ID
c\PYGZus{}ID[I] = Ith column of global array calculated by a compute with ID, I can include wildcard (see below)
f\PYGZus{}ID = global vector calculated by a fix with ID
f\PYGZus{}ID[I] = Ith column of global array calculated by a fix with ID, I can include wildcard (see below)
v\PYGZus{}name = global vector calculated by a vector\PYGZhy{}style variable with name
\end{sphinxVerbatim}

\end{itemize}


\subsection{Examples}
\label{\detokenize{compute_global_atom:examples}}
\begin{sphinxVerbatim}[commandchars=\\\{\}]
\PYG{k}{compute }\PYG{n+nv+nvIdentifier}{1}\PYG{+w}{ }\PYG{n+nv+nvIdentifier}{all}\PYG{+w}{ }\PYG{n}{global}\PYG{o}{/}\PYG{n}{atom}\PYG{+w}{ }\PYG{n}{c\PYGZus{}chunk}\PYG{+w}{ }\PYG{n}{c\PYGZus{}com[1]}\PYG{+w}{ }\PYG{n}{c\PYGZus{}com[2]}\PYG{+w}{ }\PYG{n}{c\PYGZus{}com[3]}
\PYG{k}{compute }\PYG{n+nv+nvIdentifier}{1}\PYG{+w}{ }\PYG{n+nv+nvIdentifier}{all}\PYG{+w}{ }\PYG{n}{global}\PYG{o}{/}\PYG{n}{atom}\PYG{+w}{ }\PYG{n}{c\PYGZus{}chunk}\PYG{+w}{ }\PYG{n}{c\PYGZus{}com[}\PYG{o}{*}\PYG{n}{]}
\end{sphinxVerbatim}


\subsection{Description}
\label{\detokenize{compute_global_atom:description}}
\sphinxAtStartPar
Define a calculation that assigns global values to each atom from
vectors or arrays of global values.  The specified \sphinxstyleemphasis{index} parameter
is used to determine which global value is assigned to each atom.

\sphinxAtStartPar
The \sphinxstyleemphasis{index} parameter must reference a per\sphinxhyphen{}atom vector or array from a
{\hyperref[\detokenize{compute::doc}]{\sphinxcrossref{\DUrole{doc}{compute}}}} or {\hyperref[\detokenize{fix::doc}]{\sphinxcrossref{\DUrole{doc}{fix}}}} or the evaluation of an
atom\sphinxhyphen{}style {\hyperref[\detokenize{variable::doc}]{\sphinxcrossref{\DUrole{doc}{variable}}}}.  Each \sphinxstyleemphasis{input} value must
reference a global vector or array from a {\hyperref[\detokenize{compute::doc}]{\sphinxcrossref{\DUrole{doc}{compute}}}} or
{\hyperref[\detokenize{fix::doc}]{\sphinxcrossref{\DUrole{doc}{fix}}}} or the evaluation of an vector\sphinxhyphen{}style
{\hyperref[\detokenize{variable::doc}]{\sphinxcrossref{\DUrole{doc}{variable}}}}.  Details are given below.

\sphinxAtStartPar
The \sphinxstyleemphasis{index} value for an atom is used as an index \(I\) (from 1 to
\(N\), where \(N\) is the number of atoms) into the vector
associated with each of the input values.  The \(I\)th value
from the input vector becomes one output value for that atom.  If the
atom is not in the specified group, or the index \(I < 1\) or
\(I > M\), where \(M\) is the actual length of the input vector,
then an output value of 0.0 is assigned to the atom.

\sphinxAtStartPar
An example of how this command is useful, is in the context of
“chunks” which are static or dynamic subsets of atoms.  The {\hyperref[\detokenize{compute_chunk_atom::doc}]{\sphinxcrossref{\DUrole{doc}{compute chunk/atom}}}} command assigns unique chunk IDs
to each atom.  Its output can be used as the \sphinxstyleemphasis{index} parameter for
this command.  Various other computes with “chunk” in their style
name, such as {\hyperref[\detokenize{compute_com_chunk::doc}]{\sphinxcrossref{\DUrole{doc}{compute com/chunk}}}} or
{\hyperref[\detokenize{compute_msd_chunk::doc}]{\sphinxcrossref{\DUrole{doc}{compute msd/chunk}}}}, calculate properties for each
chunk.  The output of these commands are global vectors or arrays,
with one or more values per chunk, and can be used as input values for
this command.  This command will then assign the global chunk value to
each atom in the chunk, producing a per\sphinxhyphen{}atom vector or per\sphinxhyphen{}atom array
as output.  The per\sphinxhyphen{}atom values can then be output to a dump file or
used by any command that uses per\sphinxhyphen{}atom values from a compute as input,
as discussed on the {\hyperref[\detokenize{Howto_output::doc}]{\sphinxcrossref{\DUrole{doc}{Howto output}}}} doc page.

\sphinxAtStartPar
As a concrete example, these commands will calculate the displacement
of each atom from the center\sphinxhyphen{}of\sphinxhyphen{}mass of the molecule it is in, and
dump those values to a dump file.  In this case, each molecule is a
chunk.

\begin{sphinxVerbatim}[commandchars=\\\{\}]
\PYG{k}{compute }\PYG{n+nv+nvIdentifier}{cc1}\PYG{+w}{ }\PYG{n+nv+nvIdentifier}{all}\PYG{+w}{ }\PYG{n}{chunk}\PYG{o}{/}\PYG{n}{atom}\PYG{+w}{ }\PYG{n}{molecule}
\PYG{k}{compute }\PYG{n+nv+nvIdentifier}{myChunk}\PYG{+w}{ }\PYG{n+nv+nvIdentifier}{all}\PYG{+w}{ }\PYG{n}{com}\PYG{o}{/}\PYG{n}{chunk}\PYG{+w}{ }\PYG{n}{cc1}
\PYG{k}{compute }\PYG{n+nv+nvIdentifier}{prop}\PYG{+w}{ }\PYG{n+nv+nvIdentifier}{all}\PYG{+w}{ }\PYG{n}{property}\PYG{o}{/}\PYG{n}{atom}\PYG{+w}{ }\PYG{n}{xu}\PYG{+w}{ }\PYG{n}{yu}\PYG{+w}{ }\PYG{n}{zu}
\PYG{k}{compute }\PYG{n+nv+nvIdentifier}{glob}\PYG{+w}{ }\PYG{n+nv+nvIdentifier}{all}\PYG{+w}{ }\PYG{n}{global}\PYG{o}{/}\PYG{n}{atom}\PYG{+w}{ }\PYG{n}{c\PYGZus{}cc1}\PYG{+w}{ }\PYG{n}{c\PYGZus{}myChunk[}\PYG{o}{*}\PYG{n}{]}
\PYG{k}{variable }\PYG{n+nv+nvIdentifier}{dx}\PYG{+w}{ }\PYG{n}{atom}\PYG{+w}{ }\PYG{n}{c\PYGZus{}prop[1]}\PYG{o}{\PYGZhy{}}\PYG{n}{c\PYGZus{}glob[1]}
\PYG{k}{variable }\PYG{n+nv+nvIdentifier}{dy}\PYG{+w}{ }\PYG{n}{atom}\PYG{+w}{ }\PYG{n}{c\PYGZus{}prop[2]}\PYG{o}{\PYGZhy{}}\PYG{n}{c\PYGZus{}glob[2]}
\PYG{k}{variable }\PYG{n+nv+nvIdentifier}{dz}\PYG{+w}{ }\PYG{n}{atom}\PYG{+w}{ }\PYG{n}{c\PYGZus{}prop[3]}\PYG{o}{\PYGZhy{}}\PYG{n}{c\PYGZus{}glob[3]}
\PYG{k}{variable }\PYG{n+nv+nvIdentifier}{dist}\PYG{+w}{ }\PYG{n}{atom}\PYG{+w}{ }\PYG{n}{sqrt}\PYG{n+nv}{(}\PYG{n+nv}{v\PYGZus{}dx*v\PYGZus{}dx+v\PYGZus{}dy*v\PYGZus{}dy+v\PYGZus{}dz*v\PYGZus{}dz}\PYG{n+nv}{)}
\PYG{k}{dump }\PYG{n+nv+nvIdentifier}{1}\PYG{+w}{ }\PYG{n+nv+nvIdentifier}{all}\PYG{+w}{ }\PYG{n}{custom}\PYG{+w}{ }\PYG{l+m}{100}\PYG{+w}{ }\PYG{n}{tmp.dump}\PYG{+w}{ }\PYG{n}{id}\PYG{+w}{ }\PYG{n}{xu}\PYG{+w}{ }\PYG{n}{yu}\PYG{+w}{ }\PYG{n}{zu}\PYG{+w}{ }\PYG{n}{c\PYGZus{}glob[1]}\PYG{+w}{ }\PYG{n}{c\PYGZus{}glob[2]}\PYG{+w}{ }\PYG{n}{c\PYGZus{}glob[3]}\PYG{+w}{ }\PYG{o}{\PYGZam{}}
\PYG{+w}{     }\PYG{n}{v\PYGZus{}dx}\PYG{+w}{ }\PYG{n}{v\PYGZus{}dy}\PYG{+w}{ }\PYG{n}{v\PYGZus{}dz}\PYG{+w}{ }\PYG{n}{v\PYGZus{}dist}
\PYG{k}{dump\PYGZus{}modify }\PYG{n+nv+nvIdentifier}{1}\PYG{+w}{ }\PYG{n}{sort}\PYG{+w}{ }\PYG{n}{id}
\end{sphinxVerbatim}

\sphinxAtStartPar
You can add these commands to the bench/in.chain script to see how
they work.


\bigskip\hrule\bigskip


\sphinxAtStartPar
Note that for input values from a compute or fix, the bracketed index
I can be specified using a wildcard asterisk with the index to
effectively specify multiple values.  This takes the form “*” or “*n”
or “m*” or “m*n”.  If \(N\) is the size of the vector
(for \sphinxstyleemphasis{mode} = scalar)
or the number of columns in the array (for \sphinxstyleemphasis{mode} = vector), then an
asterisk with no numeric values means all indices from 1 to \(N\).
A leading asterisk means all indices from 1 to n (inclusive).  A
trailing asterisk means all indices from m to \(N\) (inclusive).
A middle asterisk means all indices from m to n (inclusive).

\sphinxAtStartPar
Using a wildcard is the same as if the individual columns of the array
had been listed one by one. For example, the following two compute global/atom
commands are equivalent, since the {\hyperref[\detokenize{compute_com_chunk::doc}]{\sphinxcrossref{\DUrole{doc}{compute com/chunk}}}}
command creates a global array with three columns:

\begin{sphinxVerbatim}[commandchars=\\\{\}]
\PYG{k}{compute }\PYG{n+nv+nvIdentifier}{cc1}\PYG{+w}{ }\PYG{n+nv+nvIdentifier}{all}\PYG{+w}{ }\PYG{n}{chunk}\PYG{o}{/}\PYG{n}{atom}\PYG{+w}{ }\PYG{n}{molecule}
\PYG{k}{compute }\PYG{n+nv+nvIdentifier}{com}\PYG{+w}{ }\PYG{n+nv+nvIdentifier}{all}\PYG{+w}{ }\PYG{n}{com}\PYG{o}{/}\PYG{n}{chunk}\PYG{+w}{ }\PYG{n}{cc1}
\PYG{k}{compute }\PYG{n+nv+nvIdentifier}{1}\PYG{+w}{ }\PYG{n+nv+nvIdentifier}{all}\PYG{+w}{ }\PYG{n}{global}\PYG{o}{/}\PYG{n}{atom}\PYG{+w}{ }\PYG{n}{c\PYGZus{}cc1}\PYG{+w}{ }\PYG{n}{c\PYGZus{}com[1]}\PYG{+w}{ }\PYG{n}{c\PYGZus{}com[2]}\PYG{+w}{ }\PYG{n}{c\PYGZus{}com[3]}
\PYG{k}{compute }\PYG{n+nv+nvIdentifier}{1}\PYG{+w}{ }\PYG{n+nv+nvIdentifier}{all}\PYG{+w}{ }\PYG{n}{global}\PYG{o}{/}\PYG{n}{atom}\PYG{+w}{ }\PYG{n}{c\PYGZus{}cc1}\PYG{+w}{ }\PYG{n}{c\PYGZus{}com[}\PYG{o}{*}\PYG{n}{]}
\end{sphinxVerbatim}


\bigskip\hrule\bigskip


\sphinxAtStartPar
This section explains the \sphinxstyleemphasis{index} parameter.  Note that it must
reference per\sphinxhyphen{}atom values, as contrasted with the \sphinxstyleemphasis{input} values, which
must reference global values.

\sphinxAtStartPar
Note that all of these options generate floating point values.  When
they are used as an index into the specified input vectors, they
simple rounded down to convert the value to integer indices.  The
final values should range from 1 to \(N\) (inclusive), since they are used
to access values from \(N\)\sphinxhyphen{}length vectors.

\sphinxAtStartPar
If \sphinxstyleemphasis{index} begins with “c\_”, a compute ID must follow which has been
previously defined in the input script.  The compute must generate
per\sphinxhyphen{}atom quantities.  See the individual {\hyperref[\detokenize{compute::doc}]{\sphinxcrossref{\DUrole{doc}{compute}}}} doc
page for details.  If no bracketed integer is appended, the per\sphinxhyphen{}atom
vector calculated by the compute is used.  If a bracketed integer is
appended, the Ith column of the per\sphinxhyphen{}atom array calculated by the
compute is used.  Users can also write code for their own compute
styles and {\hyperref[\detokenize{Modify::doc}]{\sphinxcrossref{\DUrole{doc}{add them to LAMMPS}}}}.  See the
discussion above for how I can be specified with a wildcard asterisk
to effectively specify multiple values.

\sphinxAtStartPar
If \sphinxstyleemphasis{index} begins with “f\_”, a fix ID must follow which has been
previously defined in the input script.  The Fix must generate
per\sphinxhyphen{}atom quantities.  See the individual {\hyperref[\detokenize{fix::doc}]{\sphinxcrossref{\DUrole{doc}{fix}}}} page for
details.  Note that some fixes only produce their values on certain
timesteps, which must be compatible with when compute global/atom
references the values, else an error results.  If no bracketed integer
is appended, the per\sphinxhyphen{}atom vector calculated by the fix is used.  If a
bracketed integer is appended, the Ith column of the per\sphinxhyphen{}atom array
calculated by the fix is used.  Users can also write code for their
own fix style and {\hyperref[\detokenize{Modify::doc}]{\sphinxcrossref{\DUrole{doc}{add them to LAMMPS}}}}.  See the
discussion above for how I can be specified with a wildcard asterisk
to effectively specify multiple values.

\sphinxAtStartPar
If \sphinxstyleemphasis{index} begins with “v\_”, a variable name must follow which has
been previously defined in the input script.  It must be an
{\hyperref[\detokenize{variable::doc}]{\sphinxcrossref{\DUrole{doc}{atom\sphinxhyphen{}style variable}}}}.  Atom\sphinxhyphen{}style variables can
reference thermodynamic keywords and various per\sphinxhyphen{}atom attributes, or
invoke other computes, fixes, or variables when they are evaluated, so
this is a very general means of generating per\sphinxhyphen{}atom quantities to use
as \sphinxstyleemphasis{index}.


\bigskip\hrule\bigskip


\sphinxAtStartPar
This section explains the kinds of \sphinxstyleemphasis{input} values that can be used.
Note that inputs reference global values, as contrasted with the
\sphinxstyleemphasis{index} parameter which must reference per\sphinxhyphen{}atom values.

\sphinxAtStartPar
If a value begins with “c\_”, a compute ID must follow which has been
previously defined in the input script.  The compute must generate a
global vector or array.  See the individual {\hyperref[\detokenize{compute::doc}]{\sphinxcrossref{\DUrole{doc}{compute}}}} doc
page for details.  If no bracketed integer is appended, the vector
calculated by the compute is used.  If a bracketed integer is
appended, the Ith column of the array calculated by the compute is
used.  Users can also write code for their own compute styles and
{\hyperref[\detokenize{Modify::doc}]{\sphinxcrossref{\DUrole{doc}{add them to LAMMPS}}}}.  See the discussion above for how
I can be specified with a wildcard asterisk to effectively specify
multiple values.

\sphinxAtStartPar
If a value begins with “f\_”, a fix ID must follow which has been
previously defined in the input script.  The fix must generate a
global vector or array.  See the individual {\hyperref[\detokenize{fix::doc}]{\sphinxcrossref{\DUrole{doc}{fix}}}} doc page
for details.  Note that some fixes only produce their values on
certain timesteps, which must be compatible with when compute
global/atom references the values, else an error results.  If no
bracketed integer is appended, the vector calculated by the fix is
used.  If a bracketed integer is appended, the Ith column of the array
calculated by the fix is used.  Users can also write code for their
own fix style and {\hyperref[\detokenize{Modify::doc}]{\sphinxcrossref{\DUrole{doc}{add them to LAMMPS}}}}.  See the
discussion above for how I can be specified with a wildcard asterisk
to effectively specify multiple values.

\sphinxAtStartPar
If a value begins with “v\_”, a variable name must follow which has
been previously defined in the input script.  It must be a
{\hyperref[\detokenize{variable::doc}]{\sphinxcrossref{\DUrole{doc}{vector\sphinxhyphen{}style variable}}}}.  Vector\sphinxhyphen{}style variables can
reference thermodynamic keywords and various other attributes of
atoms, or invoke other computes, fixes, or variables when they are
evaluated, so this is a very general means of generating a vector of
global quantities which the \sphinxstyleemphasis{index} parameter will reference for
assignment of global values to atoms.


\bigskip\hrule\bigskip



\subsection{Output info}
\label{\detokenize{compute_global_atom:output-info}}
\sphinxAtStartPar
If a single input is specified this compute produces a per\sphinxhyphen{}atom
vector.  If multiple inputs are specified, this compute produces a
per\sphinxhyphen{}atom array values, where the number of columns is equal to the
number of inputs specified.  These values can be used by any command
that uses per\sphinxhyphen{}atom vector or array values from a compute as input.
See the {\hyperref[\detokenize{Howto_output::doc}]{\sphinxcrossref{\DUrole{doc}{Howto output}}}} page for an overview of
LAMMPS output options.

\sphinxAtStartPar
The per\sphinxhyphen{}atom vector or array values will be in whatever units the
corresponding input values are in.


\subsection{Restrictions}
\label{\detokenize{compute_global_atom:restrictions}}\begin{quote}

\sphinxAtStartPar
none
\end{quote}


\subsection{Related commands}
\label{\detokenize{compute_global_atom:related-commands}}
\sphinxAtStartPar
{\hyperref[\detokenize{compute::doc}]{\sphinxcrossref{\DUrole{doc}{compute}}}}, {\hyperref[\detokenize{fix::doc}]{\sphinxcrossref{\DUrole{doc}{fix}}}}, {\hyperref[\detokenize{variable::doc}]{\sphinxcrossref{\DUrole{doc}{variable}}}},
{\hyperref[\detokenize{compute_chunk_atom::doc}]{\sphinxcrossref{\DUrole{doc}{compute chunk/atom}}}}, {\hyperref[\detokenize{compute_reduce::doc}]{\sphinxcrossref{\DUrole{doc}{compute reduce}}}}


\subsection{Default}
\label{\detokenize{compute_global_atom:default}}
\sphinxAtStartPar
none

\sphinxstepscope

\index{compute group/group@\spxentry{compute group/group}}\ignorespaces 

\section{compute group/group command}
\label{\detokenize{compute_group_group:compute-group-group-command}}\label{\detokenize{compute_group_group:index-0}}\label{\detokenize{compute_group_group::doc}}

\subsection{Syntax}
\label{\detokenize{compute_group_group:syntax}}
\begin{sphinxVerbatim}[commandchars=\\\{\}]
\PYG{k}{compute }\PYG{n+nv+nvIdentifier}{ID}\PYG{+w}{ }\PYG{n+nv+nvIdentifier}{group\PYGZhy{}ID}\PYG{+w}{ }\PYG{n}{group}\PYG{o}{/}\PYG{k}{group }\PYG{n+nv+nvIdentifier}{group2\PYGZhy{}ID}\PYG{+w}{ }\PYG{n}{keyword}\PYG{+w}{ }\PYG{n}{value}\PYG{+w}{ }\PYG{n}{...}
\end{sphinxVerbatim}
\begin{itemize}
\item {} 
\sphinxAtStartPar
ID, group\sphinxhyphen{}ID are documented in {\hyperref[\detokenize{compute::doc}]{\sphinxcrossref{\DUrole{doc}{compute}}}} command

\item {} 
\sphinxAtStartPar
group/group = style name of this compute command

\item {} 
\sphinxAtStartPar
group2\sphinxhyphen{}ID = group ID of second (or same) group

\item {} 
\sphinxAtStartPar
zero or more keyword/value pairs may be appended

\item {} 
\sphinxAtStartPar
keyword = \sphinxstyleemphasis{pair} or \sphinxstyleemphasis{kspace} or \sphinxstyleemphasis{boundary} or \sphinxstyleemphasis{molecule}
\begin{sphinxalltt}
\sphinxstyleemphasis{pair} value = \sphinxstyleemphasis{yes} or \sphinxstyleemphasis{no}
\sphinxstyleemphasis{kspace} value = \sphinxstyleemphasis{yes} or \sphinxstyleemphasis{no}
\sphinxstyleemphasis{boundary} value = \sphinxstyleemphasis{yes} or \sphinxstyleemphasis{no}
\sphinxstyleemphasis{molecule} value = \sphinxstyleemphasis{off} or \sphinxstyleemphasis{inter} or \sphinxstyleemphasis{intra}
\end{sphinxalltt}

\end{itemize}


\subsection{Examples}
\label{\detokenize{compute_group_group:examples}}
\begin{sphinxVerbatim}[commandchars=\\\{\}]
\PYG{k}{compute }\PYG{n+nv+nvIdentifier}{1}\PYG{+w}{ }\PYG{n+nv+nvIdentifier}{lower}\PYG{+w}{ }\PYG{n}{group}\PYG{o}{/}\PYG{k}{group }\PYG{n+nv+nvIdentifier}{upper}
\PYG{k}{compute }\PYG{n+nv+nvIdentifier}{1}\PYG{+w}{ }\PYG{n+nv+nvIdentifier}{lower}\PYG{+w}{ }\PYG{n}{group}\PYG{o}{/}\PYG{k}{group }\PYG{n+nv+nvIdentifier}{upper}\PYG{+w}{ }\PYG{n}{kspace}\PYG{+w}{ }\PYG{n}{yes}
\PYG{k}{compute }\PYG{n+nv+nvIdentifier}{mine}\PYG{+w}{ }\PYG{n+nv+nvIdentifier}{fluid}\PYG{+w}{ }\PYG{n}{group}\PYG{o}{/}\PYG{k}{group }\PYG{n+nv+nvIdentifier}{wall}
\end{sphinxVerbatim}


\subsection{Description}
\label{\detokenize{compute_group_group:description}}
\sphinxAtStartPar
Define a computation that calculates the total energy and force
interaction between two groups of atoms: the compute group and the
specified group2.  The two groups can be the same.

\sphinxAtStartPar
If the \sphinxstyleemphasis{pair} keyword is set to \sphinxstyleemphasis{yes}, which is the default, then the
the interaction energy will include a pair component which is defined
as the pairwise energy between all pairs of atoms where one atom in
the pair is in the first group and the other is in the second group.
Likewise, the interaction force calculated by this compute will
include the force on the compute group atoms due to pairwise
interactions with atoms in the specified group2.

\begin{sphinxadmonition}{note}{Note:}
\sphinxAtStartPar
The energies computed by the \sphinxstyleemphasis{pair} keyword do not include tail
corrections, even if they are enabled via the
{\hyperref[\detokenize{pair_modify::doc}]{\sphinxcrossref{\DUrole{doc}{pair\_modify}}}} command.
\end{sphinxadmonition}

\sphinxAtStartPar
If the \sphinxstyleemphasis{molecule} keyword is set to \sphinxstyleemphasis{inter} or \sphinxstyleemphasis{intra} than an
additional check is made based on the molecule IDs of the two atoms in
each pair before including their pairwise interaction energy and
force.  For the \sphinxstyleemphasis{inter} setting, the two atoms must be in different
molecules.  For the \sphinxstyleemphasis{intra} setting, the two atoms must be in the same
molecule.

\sphinxAtStartPar
If the \sphinxstyleemphasis{kspace} keyword is set to \sphinxstyleemphasis{yes}, which is not the default, and
if a {\hyperref[\detokenize{kspace_style::doc}]{\sphinxcrossref{\DUrole{doc}{kspace\_style}}}} is defined, then the interaction
energy will include a Kspace component which is the long\sphinxhyphen{}range
Coulombic energy between all the atoms in the first group and all the
atoms in the second group.  Likewise, the interaction force calculated by
this compute will include the force on the compute group atoms due to
long\sphinxhyphen{}range Coulombic interactions with atoms in the specified group2.

\sphinxAtStartPar
Normally the long\sphinxhyphen{}range Coulombic energy converges only when the net
charge of the unit cell is zero.  However, one can assume the net
charge of the system is neutralized by a uniform background plasma,
and a correction to the system energy can be applied to reduce
artifacts. For more information see {\hyperref[\detokenize{compute_group_group:bogusz}]{\sphinxcrossref{\DUrole{std,std-ref}{(Bogusz)}}}}.  If the
\sphinxstyleemphasis{boundary} keyword is set to \sphinxstyleemphasis{yes}, which is the default, and \sphinxstyleemphasis{kspace}
contributions are included, then this energy correction term will be
added to the total group\sphinxhyphen{}group energy.  This correction term does not
affect the force calculation and will be zero if one or both of the
groups are charge neutral.  This energy correction term is the same as
that included in the regular Ewald and PPPM routines.

\begin{sphinxadmonition}{note}{Note:}
\sphinxAtStartPar
The \sphinxstyleemphasis{molecule} setting only affects the group/group
contributions calculated by the \sphinxstyleemphasis{pair} keyword.  It does not affect
the group/group contributions calculated by the \sphinxstyleemphasis{kspace} keyword.
\end{sphinxadmonition}

\sphinxAtStartPar
This compute does not calculate any bond or angle or dihedral or
improper interactions between atoms in the two groups.


\bigskip\hrule\bigskip


\sphinxAtStartPar
The pairwise contributions to the group\sphinxhyphen{}group interactions are
calculated by looping over a neighbor list.  The Kspace contribution
to the group\sphinxhyphen{}group interactions require essentially the same amount of
work (FFTs, Ewald summation) as computing long\sphinxhyphen{}range forces for the
entire system.  Thus it can be costly to invoke this compute too
frequently.

\begin{sphinxadmonition}{note}{Note:}
\sphinxAtStartPar
If you have a bonded system, then the settings of
{\hyperref[\detokenize{special_bonds::doc}]{\sphinxcrossref{\DUrole{doc}{special\_bonds}}}} command can remove pairwise
interactions between atoms in the same bond, angle, or dihedral.  This is
the default setting for the {\hyperref[\detokenize{special_bonds::doc}]{\sphinxcrossref{\DUrole{doc}{special\_bonds}}}} command,
and means those pairwise interactions do not appear in the neighbor list.
Because this compute uses a neighbor list, it also means those pairs will
not be included in the group/group interaction.  This does not apply when
using long\sphinxhyphen{}range Coulomb interactions
(\sphinxstyleemphasis{coul/long}, \sphinxstyleemphasis{coul/msm}, \sphinxstyleemphasis{coul/wolf} or similar).  One way to get
around this would be to set \sphinxstyleemphasis{special\_bond} scaling factors to very tiny
numbers that are not exactly zero (e.g., \(1.0 \times 10^{-50}\)).
Another workaround would be to write a dump file and use the
{\hyperref[\detokenize{rerun::doc}]{\sphinxcrossref{\DUrole{doc}{rerun}}}} command to compute the group/group interactions for
snapshots in the dump file. The rerun script can use a
{\hyperref[\detokenize{special_bonds::doc}]{\sphinxcrossref{\DUrole{doc}{special\_bonds}}}} command that includes all pairs in the
neighbor list.
\end{sphinxadmonition}

\sphinxAtStartPar
If you desire a breakdown of the interactions into a pairwise and
Kspace component, simply invoke the compute twice with the appropriate
yes/no settings for the \sphinxstyleemphasis{pair} and \sphinxstyleemphasis{kspace} keywords.  This is no more
costly than using a single compute with both keywords set to \sphinxstyleemphasis{yes}.
The individual contributions can be summed in a
{\hyperref[\detokenize{variable::doc}]{\sphinxcrossref{\DUrole{doc}{variable}}}} if desired.

\sphinxAtStartPar
This \sphinxhref{PDF/kspace.pdf}{document} describes how the long\sphinxhyphen{}range
group\sphinxhyphen{}group calculations are performed.


\bigskip\hrule\bigskip



\subsection{Output info}
\label{\detokenize{compute_group_group:output-info}}
\sphinxAtStartPar
This compute calculates a global scalar (the energy) and a global
vector of length 3 (force), which can be accessed by indices 1\textendash{}3.
These values can be used by any command that uses global scalar or
vector values from a compute as input.  See the
{\hyperref[\detokenize{Howto_output::doc}]{\sphinxcrossref{\DUrole{doc}{Howto output}}}} page for an overview of LAMMPS output
options.

\sphinxAtStartPar
Both the scalar and vector values calculated by this compute are
“extensive”.  The scalar value will be in energy {\hyperref[\detokenize{units::doc}]{\sphinxcrossref{\DUrole{doc}{units}}}}.
The vector values will be in force {\hyperref[\detokenize{units::doc}]{\sphinxcrossref{\DUrole{doc}{units}}}}.


\subsection{Restrictions}
\label{\detokenize{compute_group_group:restrictions}}
\sphinxAtStartPar
Not all pair styles can be evaluated in a pairwise mode as required by
this compute.  For example, three\sphinxhyphen{}body and other many\sphinxhyphen{}body potentials,
such as {\hyperref[\detokenize{pair_tersoff::doc}]{\sphinxcrossref{\DUrole{doc}{Tersoff}}}} and
{\hyperref[\detokenize{pair_sw::doc}]{\sphinxcrossref{\DUrole{doc}{Stillinger\sphinxhyphen{}Weber}}}} cannot be used.  {\hyperref[\detokenize{pair_eam::doc}]{\sphinxcrossref{\DUrole{doc}{EAM}}}}
potentials will re\sphinxhyphen{}use previously computed embedding term contributions,
so the computed pairwise forces and energies are based on the whole
system and not valid if particles have been moved since.

\sphinxAtStartPar
Not all {\hyperref[\detokenize{kspace_style::doc}]{\sphinxcrossref{\DUrole{doc}{Kspace styles}}}} support the calculation of
group/group interactions. The regular \sphinxstyleemphasis{ewald} and \sphinxstyleemphasis{pppm} styles do.


\subsection{Related commands}
\label{\detokenize{compute_group_group:related-commands}}
\sphinxAtStartPar
none


\subsection{Default}
\label{\detokenize{compute_group_group:default}}
\sphinxAtStartPar
The option defaults are pair = yes, kspace = no, boundary = yes,
molecule = off.


\bigskip\hrule\bigskip

\phantomsection\label{\detokenize{compute_group_group:bogusz}}
\sphinxAtStartPar
Bogusz et al, J Chem Phys, 108, 7070 (1998)

\sphinxstepscope

\index{compute gyration@\spxentry{compute gyration}}\ignorespaces 

\section{compute gyration command}
\label{\detokenize{compute_gyration:compute-gyration-command}}\label{\detokenize{compute_gyration:index-0}}\label{\detokenize{compute_gyration::doc}}

\subsection{Syntax}
\label{\detokenize{compute_gyration:syntax}}
\begin{sphinxVerbatim}[commandchars=\\\{\}]
\PYG{k}{compute }\PYG{n+nv+nvIdentifier}{ID}\PYG{+w}{ }\PYG{n+nv+nvIdentifier}{group\PYGZhy{}ID}\PYG{+w}{ }\PYG{n}{gyration}
\end{sphinxVerbatim}
\begin{itemize}
\item {} 
\sphinxAtStartPar
ID, group\sphinxhyphen{}ID are documented in {\hyperref[\detokenize{compute::doc}]{\sphinxcrossref{\DUrole{doc}{compute}}}} command

\item {} 
\sphinxAtStartPar
gyration = style name of this compute command

\end{itemize}


\subsection{Examples}
\label{\detokenize{compute_gyration:examples}}
\begin{sphinxVerbatim}[commandchars=\\\{\}]
\PYG{k}{compute }\PYG{n+nv+nvIdentifier}{1}\PYG{+w}{ }\PYG{n+nv+nvIdentifier}{molecule}\PYG{+w}{ }\PYG{n}{gyration}
\end{sphinxVerbatim}


\subsection{Description}
\label{\detokenize{compute_gyration:description}}
\sphinxAtStartPar
Define a computation that calculates the radius of gyration \(R_g\) of the
group of atoms, including all effects due to atoms passing through
periodic boundaries.

\sphinxAtStartPar
\(R_g\) is a measure of the size of the group of atoms, and is computed as
the square root of the \(R_g^2\) value in this formula
\begin{equation*}
\begin{split}R_g^2 = \frac{1}{M} \sum_i m_i (r_i - r_{\text{cm}})^2\end{split}
\end{equation*}
\sphinxAtStartPar
where \(M\) is the total mass of the group, \(r_{\text{cm}}\) is the
center\sphinxhyphen{}of\sphinxhyphen{}mass position of the group, and the sum is over all atoms in
the group.

\sphinxAtStartPar
A \(R_g^2\) tensor, stored as a 6\sphinxhyphen{}element vector, is also calculated
by this compute.  The formula for the components of the tensor is the
same as the above formula, except that \((r_i - r_{\text{cm}})^2\) is
replaced by
\((r_{i,x} - r_{\text{cm},x}) \cdot (r_{i,y} - r_{\text{cm},y})\) for the
\(xy\) component, and so on.  The six components of the vector are ordered
\(xx\), \(yy\), \(zz\), \(xy\), \(xz\), \(yz\).
Note that unlike the scalar \(R_g\), each of the six values of the tensor
is effectively a “squared” value, since the cross\sphinxhyphen{}terms may be negative
and taking a square root would be invalid.

\begin{sphinxadmonition}{note}{Note:}
\sphinxAtStartPar
The coordinates of an atom contribute to \(R_g\) in “unwrapped” form,
by using the image flags associated with each atom.  See the
{\hyperref[\detokenize{dump::doc}]{\sphinxcrossref{\DUrole{doc}{dump custom}}}} command for a discussion of “unwrapped” coordinates.
See the Atoms section of the {\hyperref[\detokenize{read_data::doc}]{\sphinxcrossref{\DUrole{doc}{read\_data}}}} command for a
discussion of image flags and how they are set for each atom.  You can
reset the image flags (e.g., to 0) before invoking this compute by
using the {\hyperref[\detokenize{set::doc}]{\sphinxcrossref{\DUrole{doc}{set image}}}} command.
\end{sphinxadmonition}


\subsection{Output info}
\label{\detokenize{compute_gyration:output-info}}
\sphinxAtStartPar
This compute calculates a global scalar (\(R_g\)) and a global vector of
length 6 (\(R_g^2\) tensor), which can be accessed by indices 1\textendash{}6.  These
values can be used by any command that uses a global scalar value or
vector values from a compute as input.  See the {\hyperref[\detokenize{Howto_output::doc}]{\sphinxcrossref{\DUrole{doc}{Howto output}}}} page for an overview of LAMMPS output
options.

\sphinxAtStartPar
The scalar and vector values calculated by this compute are
“intensive”.  The scalar and vector values will be in distance and
distance\(^2\) {\hyperref[\detokenize{units::doc}]{\sphinxcrossref{\DUrole{doc}{units}}}}, respectively.


\subsection{Restrictions}
\label{\detokenize{compute_gyration:restrictions}}\begin{quote}

\sphinxAtStartPar
none
\end{quote}


\subsection{Related commands}
\label{\detokenize{compute_gyration:related-commands}}
\sphinxAtStartPar
{\hyperref[\detokenize{compute_gyration_chunk::doc}]{\sphinxcrossref{\DUrole{doc}{compute gyration/chunk}}}},
{\hyperref[\detokenize{compute_gyration_shape::doc}]{\sphinxcrossref{\DUrole{doc}{compute gyration/shape}}}}


\subsection{Default}
\label{\detokenize{compute_gyration:default}}
\sphinxAtStartPar
none

\sphinxstepscope

\index{compute gyration/chunk@\spxentry{compute gyration/chunk}}\ignorespaces 

\section{compute gyration/chunk command}
\label{\detokenize{compute_gyration_chunk:compute-gyration-chunk-command}}\label{\detokenize{compute_gyration_chunk:index-0}}\label{\detokenize{compute_gyration_chunk::doc}}

\subsection{Syntax}
\label{\detokenize{compute_gyration_chunk:syntax}}
\begin{sphinxVerbatim}[commandchars=\\\{\}]
\PYG{k}{compute }\PYG{n+nv+nvIdentifier}{ID}\PYG{+w}{ }\PYG{n+nv+nvIdentifier}{group\PYGZhy{}ID}\PYG{+w}{ }\PYG{n}{gyration}\PYG{o}{/}\PYG{n}{chunk}\PYG{+w}{ }\PYG{n}{chunkID}\PYG{+w}{ }\PYG{n}{keyword}\PYG{+w}{ }\PYG{n}{value}\PYG{+w}{ }\PYG{n}{...}
\end{sphinxVerbatim}
\begin{itemize}
\item {} 
\sphinxAtStartPar
ID, group\sphinxhyphen{}ID are documented in {\hyperref[\detokenize{compute::doc}]{\sphinxcrossref{\DUrole{doc}{compute}}}} command

\item {} 
\sphinxAtStartPar
gyration/chunk = style name of this compute command

\item {} 
\sphinxAtStartPar
chunkID = ID of {\hyperref[\detokenize{compute_chunk_atom::doc}]{\sphinxcrossref{\DUrole{doc}{compute chunk/atom}}}} command

\item {} 
\sphinxAtStartPar
zero or more keyword/value pairs may be appended

\item {} 
\sphinxAtStartPar
keyword = \sphinxstyleemphasis{tensor}
\begin{sphinxalltt}
\sphinxstyleemphasis{tensor} value = none
\end{sphinxalltt}

\end{itemize}


\subsection{Examples}
\label{\detokenize{compute_gyration_chunk:examples}}
\begin{sphinxVerbatim}[commandchars=\\\{\}]
\PYG{k}{compute }\PYG{n+nv+nvIdentifier}{1}\PYG{+w}{ }\PYG{n+nv+nvIdentifier}{molecule}\PYG{+w}{ }\PYG{n}{gyration}\PYG{o}{/}\PYG{n}{chunk}\PYG{+w}{ }\PYG{n}{molchunk}
\PYG{k}{compute }\PYG{n+nv+nvIdentifier}{2}\PYG{+w}{ }\PYG{n+nv+nvIdentifier}{molecule}\PYG{+w}{ }\PYG{n}{gyration}\PYG{o}{/}\PYG{n}{chunk}\PYG{+w}{ }\PYG{n}{molchunk}\PYG{+w}{ }\PYG{n}{tensor}
\end{sphinxVerbatim}


\subsection{Description}
\label{\detokenize{compute_gyration_chunk:description}}
\sphinxAtStartPar
Define a computation that calculates the radius of gyration \(R_g\) for
multiple chunks of atoms.

\sphinxAtStartPar
In LAMMPS, chunks are collections of atoms defined by a
{\hyperref[\detokenize{compute_chunk_atom::doc}]{\sphinxcrossref{\DUrole{doc}{compute chunk/atom}}}} command, which assigns each atom
to a single chunk (or no chunk).  The ID for this command is specified
as chunkID.  For example, a single chunk could be the atoms in a
molecule or atoms in a spatial bin.
See the {\hyperref[\detokenize{compute_chunk_atom::doc}]{\sphinxcrossref{\DUrole{doc}{compute chunk/atom}}}} and
{\hyperref[\detokenize{Howto_chunk::doc}]{\sphinxcrossref{\DUrole{doc}{Howto chunk}}}}
doc pages for details of how chunks can be defined and examples of how
they can be used to measure properties of a system.

\sphinxAtStartPar
This compute calculates the radius of gyration \(R_g\) for each chunk,
which includes all effects due to atoms passing through periodic
boundaries.

\sphinxAtStartPar
\(R_g\) is a measure of the size of a chunk, and is computed by the
formula
\begin{equation*}
\begin{split}R_g^2 = \frac{1}{M} \sum_i m_i (r_i - r_{\text{cm}})^2\end{split}
\end{equation*}
\sphinxAtStartPar
where \(M\) is the total mass of the chunk, \(r_{\text{cm}}\) is
the center\sphinxhyphen{}of\sphinxhyphen{}mass position of the chunk, and the sum is over all atoms in the
chunk.

\sphinxAtStartPar
Note that only atoms in the specified group contribute to the
calculation.  The {\hyperref[\detokenize{compute_chunk_atom::doc}]{\sphinxcrossref{\DUrole{doc}{compute chunk/atom}}}} command
defines its own group; atoms will have a chunk ID = 0 if they are not
in that group, signifying they are not assigned to a chunk, and will
thus also not contribute to this calculation.  You can specify the
“all” group for this command if you simply want to include atoms with
non\sphinxhyphen{}zero chunk IDs.

\sphinxAtStartPar
If the \sphinxstyleemphasis{tensor} keyword is specified, then the scalar \(R_g\) value is not
calculated, but an \(R_g\) tensor is instead calculated for each chunk.
The formula for the components of the tensor is the same as the above
formula, except that \((r_i - r_{\text{cm}})^2\) is replaced by
\((r_{i,x} - r_{\text{cm},x}) \cdot (r_{i,y} - r_{\text{cm},y})\) for the
\(xy\) component, and so on.  The six components of the tensor are
ordered \(xx\), \(yy\), \(zz\), \(xy\), \(xz\), \(yz\).

\begin{sphinxadmonition}{note}{Note:}
\sphinxAtStartPar
The coordinates of an atom contribute to \(R_g\) in “unwrapped” form,
by using the image flags associated with each atom.  See the {\hyperref[\detokenize{dump::doc}]{\sphinxcrossref{\DUrole{doc}{dump custom}}}}
command for a discussion of “unwrapped” coordinates.
See the Atoms section of the {\hyperref[\detokenize{read_data::doc}]{\sphinxcrossref{\DUrole{doc}{read\_data}}}} command for a
discussion of image flags and how they are set for each atom.  You can
reset the image flags (e.g., to 0) before invoking this compute by
using the {\hyperref[\detokenize{set::doc}]{\sphinxcrossref{\DUrole{doc}{set image}}}} command.
\end{sphinxadmonition}

\sphinxAtStartPar
The simplest way to output the results of the compute gyration/chunk
calculation to a file is to use the {\hyperref[\detokenize{fix_ave_time::doc}]{\sphinxcrossref{\DUrole{doc}{fix ave/time}}}}
command, for example:

\begin{sphinxVerbatim}[commandchars=\\\{\}]
\PYG{k}{compute }\PYG{n+nv+nvIdentifier}{cc1}\PYG{+w}{ }\PYG{n+nv+nvIdentifier}{all}\PYG{+w}{ }\PYG{n}{chunk}\PYG{o}{/}\PYG{n}{atom}\PYG{+w}{ }\PYG{n}{molecule}
\PYG{k}{compute }\PYG{n+nv+nvIdentifier}{myChunk}\PYG{+w}{ }\PYG{n+nv+nvIdentifier}{all}\PYG{+w}{ }\PYG{n}{gyration}\PYG{o}{/}\PYG{n}{chunk}\PYG{+w}{ }\PYG{n}{cc1}
\PYG{k}{fix }\PYG{n+nv+nvIdentifier}{1}\PYG{+w}{ }\PYG{n+nv+nvIdentifier}{all}\PYG{+w}{ }\PYG{n}{ave}\PYG{o}{/}\PYG{n}{time}\PYG{+w}{ }\PYG{l+m}{100}\PYG{+w}{ }\PYG{l+m}{1}\PYG{+w}{ }\PYG{l+m}{100}\PYG{+w}{ }\PYG{n}{c\PYGZus{}myChunk}\PYG{+w}{ }\PYG{n}{file}\PYG{+w}{ }\PYG{n}{tmp.out}\PYG{+w}{ }\PYG{n}{mode}\PYG{+w}{ }\PYG{n}{vector}
\end{sphinxVerbatim}


\subsection{Output info}
\label{\detokenize{compute_gyration_chunk:output-info}}
\sphinxAtStartPar
This compute calculates a global vector if the \sphinxstyleemphasis{tensor} keyword is not
specified and a global array if it is.  The length of the vector or
number of rows in the array = the number of chunks \sphinxstyleemphasis{Nchunk} as
calculated by the specified {\hyperref[\detokenize{compute_chunk_atom::doc}]{\sphinxcrossref{\DUrole{doc}{compute chunk/atom}}}}
command.  If the \sphinxstyleemphasis{tensor} keyword is specified, the global array has six
columns.  The vector or array can
be accessed by any command that uses global values from a compute as
input.  See the {\hyperref[\detokenize{Howto_output::doc}]{\sphinxcrossref{\DUrole{doc}{Howto output}}}} page for an
overview of LAMMPS output options.

\sphinxAtStartPar
All the vector or array values calculated by this compute are
“intensive”.  The vector or array values will be in distance
{\hyperref[\detokenize{units::doc}]{\sphinxcrossref{\DUrole{doc}{units}}}}, since they are the square root of values
represented by the formula above.


\subsection{Restrictions}
\label{\detokenize{compute_gyration_chunk:restrictions}}\begin{quote}

\sphinxAtStartPar
none
\end{quote}


\subsection{Related commands}
\label{\detokenize{compute_gyration_chunk:related-commands}}
\sphinxAtStartPar
none

\sphinxAtStartPar
{\hyperref[\detokenize{compute_gyration::doc}]{\sphinxcrossref{\DUrole{doc}{compute gyration}}}}


\subsection{Default}
\label{\detokenize{compute_gyration_chunk:default}}
\sphinxAtStartPar
none

\sphinxstepscope

\index{compute gyration/shape@\spxentry{compute gyration/shape}}\ignorespaces 

\section{compute gyration/shape command}
\label{\detokenize{compute_gyration_shape:compute-gyration-shape-command}}\label{\detokenize{compute_gyration_shape:index-0}}\label{\detokenize{compute_gyration_shape::doc}}

\subsection{Syntax}
\label{\detokenize{compute_gyration_shape:syntax}}
\begin{sphinxVerbatim}[commandchars=\\\{\}]
\PYG{k}{compute }\PYG{n+nv+nvIdentifier}{ID}\PYG{+w}{ }\PYG{n+nv+nvIdentifier}{group\PYGZhy{}ID}\PYG{+w}{ }\PYG{n}{gyration}\PYG{o}{/}\PYG{n}{shape}\PYG{+w}{ }\PYG{n}{compute}\PYG{o}{\PYGZhy{}}\PYG{n}{ID}
\end{sphinxVerbatim}
\begin{itemize}
\item {} 
\sphinxAtStartPar
ID, group\sphinxhyphen{}ID are documented in {\hyperref[\detokenize{compute::doc}]{\sphinxcrossref{\DUrole{doc}{compute}}}} command

\item {} 
\sphinxAtStartPar
gyration/shape = style name of this compute command

\item {} 
\sphinxAtStartPar
compute\sphinxhyphen{}ID = ID of {\hyperref[\detokenize{compute_gyration::doc}]{\sphinxcrossref{\DUrole{doc}{compute gyration}}}} command

\end{itemize}


\subsection{Examples}
\label{\detokenize{compute_gyration_shape:examples}}
\begin{sphinxVerbatim}[commandchars=\\\{\}]
\PYG{k}{compute }\PYG{n+nv+nvIdentifier}{1}\PYG{+w}{ }\PYG{n+nv+nvIdentifier}{molecule}\PYG{+w}{ }\PYG{n}{gyration}\PYG{o}{/}\PYG{n}{shape}\PYG{+w}{ }\PYG{n}{pe}
\end{sphinxVerbatim}


\subsection{Description}
\label{\detokenize{compute_gyration_shape:description}}
\sphinxAtStartPar
Define a computation that calculates the eigenvalues of the gyration tensor of a
group of atoms and three shape parameters. The computation includes all effects
due to atoms passing through periodic boundaries.

\sphinxAtStartPar
The three computed shape parameters are the asphericity, \(b\),
the acylindricity, \(c\), and the relative shape anisotropy, \(k\),
viz.,
\begin{equation*}
\begin{split}b &= l_z - \frac12(l_y+l_x) \\
c &= l_y - l_x \\
k &= \frac{3}{2} \frac{l_x^2+l_y^2+l_z^2}{(l_x+l_y+l_z)^2} - \frac{1}{2}\end{split}
\end{equation*}
\sphinxAtStartPar
where \(l_x \le l_y \le l_z\) are the three eigenvalues of the gyration
tensor. A general description of these parameters is provided in
{\hyperref[\detokenize{compute_gyration_shape:mattice1}]{\sphinxcrossref{\DUrole{std,std-ref}{(Mattice)}}}} while an application to polymer systems
can be found in {\hyperref[\detokenize{compute_gyration_shape:theodorou1}]{\sphinxcrossref{\DUrole{std,std-ref}{(Theodorou)}}}}.
The asphericity is always non\sphinxhyphen{}negative and zero only when the three principal
moments are equal. This zero condition is met when the distribution of
particles is spherically symmetric (hence the name asphericity) but also
whenever the particle distribution is symmetric with respect to the three
coordinate axes (e.g., when the particles are distributed uniformly on a cube,
tetrahedron or other Platonic solid). The acylindricity is always non\sphinxhyphen{}negative
and zero only when the two principal moments are equal. This zero condition is
met when the distribution of particles is cylindrically symmetric (hence the
name, acylindricity), but also whenever the particle distribution is symmetric
with respect to the two coordinate axes (e.g., when the
particles are distributed uniformly on a regular prism).
The relative shape anisotropy is bounded between zero (if all points are
spherically symmetric) and one (if all points lie on a line).

\begin{sphinxadmonition}{note}{Note:}
\sphinxAtStartPar
The coordinates of an atom contribute to the gyration tensor in
“unwrapped” form, by using the image flags associated with each atom.
See the {\hyperref[\detokenize{dump::doc}]{\sphinxcrossref{\DUrole{doc}{dump custom}}}} command for a discussion of “unwrapped”
coordinates. See the Atoms section of the {\hyperref[\detokenize{read_data::doc}]{\sphinxcrossref{\DUrole{doc}{read\_data}}}}
command for a discussion of image flags and how they are set for each
atom.  You can reset the image flags (e.g., to 0) before invoking this
compute by using the {\hyperref[\detokenize{set::doc}]{\sphinxcrossref{\DUrole{doc}{set image}}}} command.
\end{sphinxadmonition}


\subsection{Output info}
\label{\detokenize{compute_gyration_shape:output-info}}
\sphinxAtStartPar
This compute calculates a global vector of length 6, which can be accessed by
indices 1\textendash{}6. The first three values are the eigenvalues of the gyration tensor
followed by the asphericity, the acylindricity and the relative shape
anisotropy.  The computed values can be used by any command that uses global
vector values from a compute as input.  See the
{\hyperref[\detokenize{Howto_output::doc}]{\sphinxcrossref{\DUrole{doc}{Howto output}}}} page for an overview of LAMMPS output
options.

\sphinxAtStartPar
The vector values calculated by this compute are
“intensive”.  The first five vector values will be in
distance\(2\) {\hyperref[\detokenize{units::doc}]{\sphinxcrossref{\DUrole{doc}{units}}}} while the sixth one is dimensionless.


\subsection{Restrictions}
\label{\detokenize{compute_gyration_shape:restrictions}}
\sphinxAtStartPar
This compute is part of the EXTRA\sphinxhyphen{}COMPUTE package.  It is only enabled if
LAMMPS was built with that package.  See the {\hyperref[\detokenize{Build_package::doc}]{\sphinxcrossref{\DUrole{doc}{Build package}}}} page for more info.


\subsection{Related commands}
\label{\detokenize{compute_gyration_shape:related-commands}}
\sphinxAtStartPar
{\hyperref[\detokenize{compute_gyration::doc}]{\sphinxcrossref{\DUrole{doc}{compute gyration}}}}


\subsection{Default}
\label{\detokenize{compute_gyration_shape:default}}
\sphinxAtStartPar
none


\bigskip\hrule\bigskip

\phantomsection\label{\detokenize{compute_gyration_shape:mattice1}}
\sphinxAtStartPar
\sphinxstylestrong{(Mattice)} Mattice, Suter, Conformational Theory of Large Molecules, Wiley, New York, 1994.

\phantomsection\label{\detokenize{compute_gyration_shape:theodorou1}}
\sphinxAtStartPar
\sphinxstylestrong{(Theodorou)} Theodorou, Suter, Macromolecules, 18, 1206 (1985).

\sphinxstepscope

\index{compute gyration/shape/chunk@\spxentry{compute gyration/shape/chunk}}\ignorespaces 

\section{compute gyration/shape/chunk command}
\label{\detokenize{compute_gyration_shape_chunk:compute-gyration-shape-chunk-command}}\label{\detokenize{compute_gyration_shape_chunk:index-0}}\label{\detokenize{compute_gyration_shape_chunk::doc}}

\subsection{Syntax}
\label{\detokenize{compute_gyration_shape_chunk:syntax}}
\begin{sphinxVerbatim}[commandchars=\\\{\}]
\PYG{k}{compute }\PYG{n+nv+nvIdentifier}{ID}\PYG{+w}{ }\PYG{n+nv+nvIdentifier}{group\PYGZhy{}ID}\PYG{+w}{ }\PYG{n}{gyration}\PYG{o}{/}\PYG{n}{shape}\PYG{o}{/}\PYG{n}{chunk}\PYG{+w}{ }\PYG{n}{compute}\PYG{o}{\PYGZhy{}}\PYG{n}{ID}
\end{sphinxVerbatim}
\begin{itemize}
\item {} 
\sphinxAtStartPar
ID, group\sphinxhyphen{}ID are documented in {\hyperref[\detokenize{compute::doc}]{\sphinxcrossref{\DUrole{doc}{compute}}}} command

\item {} 
\sphinxAtStartPar
gyration/shape/chunk = style name of this compute command

\item {} 
\sphinxAtStartPar
compute\sphinxhyphen{}ID = ID of {\hyperref[\detokenize{compute_gyration_chunk::doc}]{\sphinxcrossref{\DUrole{doc}{compute gyration/chunk}}}} command

\end{itemize}


\subsection{Examples}
\label{\detokenize{compute_gyration_shape_chunk:examples}}
\begin{sphinxVerbatim}[commandchars=\\\{\}]
\PYG{k}{compute }\PYG{n+nv+nvIdentifier}{1}\PYG{+w}{ }\PYG{n+nv+nvIdentifier}{molecule}\PYG{+w}{ }\PYG{n}{gyration}\PYG{o}{/}\PYG{n}{shape}\PYG{o}{/}\PYG{n}{chunk}\PYG{+w}{ }\PYG{n}{pe}
\end{sphinxVerbatim}


\subsection{Description}
\label{\detokenize{compute_gyration_shape_chunk:description}}
\sphinxAtStartPar
Define a computation that calculates the eigenvalues of the gyration tensor and
three shape parameters of multiple chunks of atoms. The computation includes
all effects due to atoms passing through periodic boundaries.

\sphinxAtStartPar
The three computed shape parameters are the asphericity, \(b\),
the acylindricity, \(c\), and the relative shape anisotropy, \(k\),
viz.,
\begin{equation*}
\begin{split}b &= l_z - \frac12(l_y+l_x) \\
c &= l_y - l_x \\
k &= \frac{3}{2} \frac{l_x^2+l_y^2+l_z^2}{(l_x+l_y+l_z)^2} - \frac{1}{2}\end{split}
\end{equation*}
\sphinxAtStartPar
where \(l_x \le l_y \le l_z\) are the three eigenvalues of the gyration
tensor. A general description of these parameters is provided in
{\hyperref[\detokenize{compute_gyration_shape_chunk:mattice2}]{\sphinxcrossref{\DUrole{std,std-ref}{(Mattice)}}}} while an application to polymer systems
can be found in {\hyperref[\detokenize{compute_gyration_shape_chunk:theodorou2}]{\sphinxcrossref{\DUrole{std,std-ref}{(Theodorou)}}}}. The asphericity is always
non\sphinxhyphen{}negative and zero only when the three principal moments are equal.
This zero condition is met when the distribution of particles is spherically
symmetric (hence the name asphericity) but also whenever the particle
distribution is symmetric with respect to the three coordinate axes (e.g.,
when the particles are distributed uniformly on a cube, tetrahedron, or other
Platonic solid). The acylindricity is always non\sphinxhyphen{}negative and zero only when
the two principal moments are equal. This zero condition is met when the
distribution of particles is cylindrically symmetric (hence the name,
acylindricity), but also whenever the particle distribution is symmetric with
respect to the two coordinate axes (e.g., when the particles are distributed
uniformly on a regular prism). The relative shape anisotropy
is bounded between 0 (if all points are spherically symmetric) and 1
(if all points lie on a line).

\sphinxAtStartPar
The tensor keyword must be specified in the compute gyration/chunk command.

\begin{sphinxadmonition}{note}{Note:}
\sphinxAtStartPar
The coordinates of an atom contribute to the gyration tensor in
“unwrapped” form, by using the image flags associated with each atom.
See the {\hyperref[\detokenize{dump::doc}]{\sphinxcrossref{\DUrole{doc}{dump custom}}}} command for a discussion of “unwrapped”
coordinates. See the Atoms section of the {\hyperref[\detokenize{read_data::doc}]{\sphinxcrossref{\DUrole{doc}{read\_data}}}}
command for a discussion of image flags and how they are set for each
atom.  You can reset the image flags (e.g., to 0) before invoking this
compute by using the {\hyperref[\detokenize{set::doc}]{\sphinxcrossref{\DUrole{doc}{set image}}}} command.
\end{sphinxadmonition}


\subsection{Output info}
\label{\detokenize{compute_gyration_shape_chunk:output-info}}
\sphinxAtStartPar
This compute calculates a global array with six columns,
which can be accessed by indices 1\textendash{}6. The first three columns are the
eigenvalues of the gyration tensor followed by the asphericity, the
acylindricity and the relative shape anisotropy.  The computed values can be
used by any command that uses global array values from a compute as input.
See the {\hyperref[\detokenize{Howto_output::doc}]{\sphinxcrossref{\DUrole{doc}{Howto output}}}} page for an overview of LAMMPS
output options.

\sphinxAtStartPar
The array calculated by this compute is
“intensive”.  The first five columns will be in
distance\(^2\) {\hyperref[\detokenize{units::doc}]{\sphinxcrossref{\DUrole{doc}{units}}}} while the sixth one is dimensionless.


\subsection{Restrictions}
\label{\detokenize{compute_gyration_shape_chunk:restrictions}}
\sphinxAtStartPar
This compute is part of the EXTRA\sphinxhyphen{}COMPUTE package.  It is only enabled if
LAMMPS was built with that package.  See the {\hyperref[\detokenize{Build_package::doc}]{\sphinxcrossref{\DUrole{doc}{Build package}}}} page for more info.


\subsection{Related commands}
\label{\detokenize{compute_gyration_shape_chunk:related-commands}}
\sphinxAtStartPar
{\hyperref[\detokenize{compute_gyration_chunk::doc}]{\sphinxcrossref{\DUrole{doc}{compute gyration/chunk}}}}
{\hyperref[\detokenize{compute_gyration_shape::doc}]{\sphinxcrossref{\DUrole{doc}{compute gyration/shape}}}}


\subsection{Default}
\label{\detokenize{compute_gyration_shape_chunk:default}}
\sphinxAtStartPar
none


\bigskip\hrule\bigskip

\phantomsection\label{\detokenize{compute_gyration_shape_chunk:mattice2}}
\sphinxAtStartPar
\sphinxstylestrong{(Mattice)} Mattice, Suter, Conformational Theory of Large Molecules, Wiley, New York, 1994.

\phantomsection\label{\detokenize{compute_gyration_shape_chunk:theodorou2}}
\sphinxAtStartPar
\sphinxstylestrong{(Theodorou)} Theodorou, Suter, Macromolecules, 18, 1206 (1985).

\sphinxstepscope

\index{compute heat/flux@\spxentry{compute heat/flux}}\ignorespaces 

\section{compute heat/flux command}
\label{\detokenize{compute_heat_flux:compute-heat-flux-command}}\label{\detokenize{compute_heat_flux:index-0}}\label{\detokenize{compute_heat_flux::doc}}

\subsection{Syntax}
\label{\detokenize{compute_heat_flux:syntax}}
\begin{sphinxVerbatim}[commandchars=\\\{\}]
\PYG{k}{compute }\PYG{n+nv+nvIdentifier}{ID}\PYG{+w}{ }\PYG{n+nv+nvIdentifier}{group\PYGZhy{}ID}\PYG{+w}{ }\PYG{n}{heat}\PYG{o}{/}\PYG{n}{flux}\PYG{+w}{ }\PYG{n}{ke}\PYG{o}{\PYGZhy{}}\PYG{n}{ID}\PYG{+w}{ }\PYG{n}{pe}\PYG{o}{\PYGZhy{}}\PYG{n}{ID}\PYG{+w}{ }\PYG{n}{stress}\PYG{o}{\PYGZhy{}}\PYG{n}{ID}
\end{sphinxVerbatim}
\begin{itemize}
\item {} 
\sphinxAtStartPar
ID, group\sphinxhyphen{}ID are documented in {\hyperref[\detokenize{compute::doc}]{\sphinxcrossref{\DUrole{doc}{compute}}}} command

\item {} 
\sphinxAtStartPar
heat/flux = style name of this compute command

\item {} 
\sphinxAtStartPar
ke\sphinxhyphen{}ID = ID of a compute that calculates per\sphinxhyphen{}atom kinetic energy

\item {} 
\sphinxAtStartPar
pe\sphinxhyphen{}ID = ID of a compute that calculates per\sphinxhyphen{}atom potential energy

\item {} 
\sphinxAtStartPar
stress\sphinxhyphen{}ID = ID of a compute that calculates per\sphinxhyphen{}atom stress

\end{itemize}


\subsection{Examples}
\label{\detokenize{compute_heat_flux:examples}}
\begin{sphinxVerbatim}[commandchars=\\\{\}]
\PYG{k}{compute }\PYG{n+nv+nvIdentifier}{myFlux}\PYG{+w}{ }\PYG{n+nv+nvIdentifier}{all}\PYG{+w}{ }\PYG{n}{heat}\PYG{o}{/}\PYG{n}{flux}\PYG{+w}{ }\PYG{n}{myKE}\PYG{+w}{ }\PYG{n}{myPE}\PYG{+w}{ }\PYG{n}{myStress}
\end{sphinxVerbatim}


\subsection{Description}
\label{\detokenize{compute_heat_flux:description}}
\sphinxAtStartPar
Define a computation that calculates the heat flux vector based on
contributions from atoms in the specified group.  This can be used by
itself to measure the heat flux through a set of atoms (e.g., a region
between two thermostatted reservoirs held at different temperatures),
or to calculate a thermal conductivity using the equilibrium
Green\sphinxhyphen{}Kubo formalism.

\sphinxAtStartPar
For other non\sphinxhyphen{}equilibrium ways to compute a thermal conductivity, see
the {\hyperref[\detokenize{Howto_kappa::doc}]{\sphinxcrossref{\DUrole{doc}{Howto kappa}}}} doc page.  These include use of
the {\hyperref[\detokenize{fix_thermal_conductivity::doc}]{\sphinxcrossref{\DUrole{doc}{fix thermal/conductivity}}}} command
for the Muller\sphinxhyphen{}Plathe method.  Or the {\hyperref[\detokenize{fix_heat::doc}]{\sphinxcrossref{\DUrole{doc}{fix heat}}}} command
which can add or subtract heat from groups of atoms.

\sphinxAtStartPar
The compute takes three arguments which are IDs of other
{\hyperref[\detokenize{compute::doc}]{\sphinxcrossref{\DUrole{doc}{computes}}}}.  One calculates per\sphinxhyphen{}atom kinetic energy
(\sphinxstyleemphasis{ke\sphinxhyphen{}ID}), one calculates per\sphinxhyphen{}atom potential energy (\sphinxstyleemphasis{pe\sphinxhyphen{}ID)}, and the
third calculates per\sphinxhyphen{}atom stress (\sphinxstyleemphasis{stress\sphinxhyphen{}ID}).

\begin{sphinxadmonition}{note}{Note:}
\sphinxAtStartPar
These other computes should provide values for all the atoms in
the group this compute specifies.  That means the other computes could
use the same group as this compute, or they can just use group “all”
(or any group whose atoms are superset of the atoms in this compute’s
group).  LAMMPS does not check for this.
\end{sphinxadmonition}

\sphinxAtStartPar
In case of two\sphinxhyphen{}body interactions, the heat flux \(\mathbf{J}\) is defined as
\begin{equation*}
\begin{split}\mathbf{J} &= \frac{1}{V} \left[ \sum_i e_i \mathbf{v}_i - \sum_{i} \mathbf{S}_{i} \mathbf{v}_i \right] \\
&= \frac{1}{V} \left[ \sum_i e_i \mathbf{v}_i + \sum_{i<j} \left( \mathbf{F}_{ij} \cdot \mathbf{v}_j \right) \mathbf{r}_{ij} \right] \\
&= \frac{1}{V} \left[ \sum_i e_i \mathbf{v}_i + \frac{1}{2} \sum_{i<j} \bigl( \mathbf{F}_{ij} \cdot \left(\mathbf{v}_i + \mathbf{v}_j \right) \bigr) \mathbf{r}_{ij} \right]\end{split}
\end{equation*}
\sphinxAtStartPar
\(e_i\) in the first term of the equation
is the per\sphinxhyphen{}atom energy (potential and kinetic).
This is calculated by the computes \sphinxstyleemphasis{ke\sphinxhyphen{}ID}
and \sphinxstyleemphasis{pe\sphinxhyphen{}ID}. \(\mathbf{S}_i\) in the second term is the
per\sphinxhyphen{}atom stress tensor calculated by the compute \sphinxstyleemphasis{stress\sphinxhyphen{}ID}.
See {\hyperref[\detokenize{compute_stress_atom::doc}]{\sphinxcrossref{\DUrole{doc}{compute stress/atom}}}}
and {\hyperref[\detokenize{compute_stress_atom::doc}]{\sphinxcrossref{\DUrole{doc}{compute centroid/stress/atom}}}}
for possible definitions of atomic stress \(\mathbf{S}_i\)
in the case of bonded and many\sphinxhyphen{}body interactions.
The tensor multiplies \(\mathbf{v}_i\) by a \(3\times3\) matrix
to yield a vector.
Note that as discussed below, the \(1/V\) scaling factor in the
equation for \(\mathbf{J}\) is \sphinxstylestrong{not} included in the calculation
performed by these computes; you need to add it for a volume appropriate to the
atoms included in the calculation.

\begin{sphinxadmonition}{note}{Note:}
\sphinxAtStartPar
The {\hyperref[\detokenize{compute_pe_atom::doc}]{\sphinxcrossref{\DUrole{doc}{compute pe/atom}}}} and
{\hyperref[\detokenize{compute_stress_atom::doc}]{\sphinxcrossref{\DUrole{doc}{compute stress/atom}}}}
commands have options for which
terms to include in their calculation (pair, bond, etc).  The heat
flux calculation will thus include exactly the same terms. Normally
you should use {\hyperref[\detokenize{compute_stress_atom::doc}]{\sphinxcrossref{\DUrole{doc}{compute stress/atom virial}}}}
or {\hyperref[\detokenize{compute_stress_atom::doc}]{\sphinxcrossref{\DUrole{doc}{compute centroid/stress/atom virial}}}}
so as not to include a kinetic energy term in the heat flux.
\end{sphinxadmonition}

\begin{sphinxadmonition}{warning}{Warning:}
\sphinxAtStartPar
The compute \sphinxstyleemphasis{heat/flux} has been reported to produce unphysical
values for angle, dihedral, improper and constraint force contributions
when used with {\hyperref[\detokenize{compute_stress_atom::doc}]{\sphinxcrossref{\DUrole{doc}{compute stress/atom}}}},
as discussed in {\hyperref[\detokenize{compute_heat_flux:surblys3}]{\sphinxcrossref{\DUrole{std,std-ref}{(Surblys2019)}}}}, {\hyperref[\detokenize{compute_heat_flux:boone}]{\sphinxcrossref{\DUrole{std,std-ref}{(Boone)}}}}
and {\hyperref[\detokenize{compute_heat_flux:surblys4}]{\sphinxcrossref{\DUrole{std,std-ref}{(Surblys2021)}}}}. You are strongly advised to
use {\hyperref[\detokenize{compute_stress_atom::doc}]{\sphinxcrossref{\DUrole{doc}{compute centroid/stress/atom}}}},
which has been implemented specifically for such cases.
\end{sphinxadmonition}

\begin{sphinxadmonition}{warning}{Warning:}
\sphinxAtStartPar
Due to an implementation detail, the \(y\) and \(z\)
components of heat flux from {\hyperref[\detokenize{fix_rigid::doc}]{\sphinxcrossref{\DUrole{doc}{fix rigid}}}}
contribution when computed via {\hyperref[\detokenize{compute_stress_atom::doc}]{\sphinxcrossref{\DUrole{doc}{compute stress/atom}}}}
are highly unphysical and should not be used.
\end{sphinxadmonition}

\sphinxAtStartPar
The Green\textendash{}Kubo formulas relate the ensemble average of the
auto\sphinxhyphen{}correlation of the heat flux \(\mathbf{J}\)
to the thermal conductivity \(\kappa\):
\begin{equation*}
\begin{split}\kappa  = \frac{V}{k_B T^2} \int_0^\infty \langle J_x(0)  J_x(t) \rangle \, \mathrm{d} t = \frac{V}{3 k_B T^2} \int_0^\infty \langle \mathbf{J}(0) \cdot  \mathbf{J}(t)  \rangle \, \mathrm{d}t\end{split}
\end{equation*}

\bigskip\hrule\bigskip


\sphinxAtStartPar
The heat flux can be output every so many timesteps (e.g., via the
{\hyperref[\detokenize{thermo_style::doc}]{\sphinxcrossref{\DUrole{doc}{thermo\_style custom}}}} command).  Then as a
post\sphinxhyphen{}processing operation, an auto\sphinxhyphen{}correlation can be performed, its
integral estimated, and the Green\textendash{}Kubo formula above evaluated.

\sphinxAtStartPar
The {\hyperref[\detokenize{fix_ave_correlate::doc}]{\sphinxcrossref{\DUrole{doc}{fix ave/correlate}}}} command can calculate
the auto\sphinxhyphen{}correlation.  The trap() function in the
{\hyperref[\detokenize{variable::doc}]{\sphinxcrossref{\DUrole{doc}{variable}}}} command can calculate the integral.

\sphinxAtStartPar
An example LAMMPS input script for solid argon is appended below.  The
result should be an average conductivity
\(\approx 0.29~\mathrm{W/m \cdot K}\).


\bigskip\hrule\bigskip



\subsection{Output info}
\label{\detokenize{compute_heat_flux:output-info}}
\sphinxAtStartPar
This compute calculates a global vector of length 6.
The first three components are the \(x\), \(y\), and \(z\)
components of the full heat flux vector
(i.e., \(J_x\), \(J_y\), and \(J_z\)).
The next three components are the \(x\), \(y\), and \(z\)
components of just the convective portion of the flux (i.e., the
first term in the equation for \(\mathbf{J}\)).
Each component can be accessed by indices 1\textendash{}6. These values can be used by any
command that uses global vector values from a compute as input.
See the {\hyperref[\detokenize{Howto_output::doc}]{\sphinxcrossref{\DUrole{doc}{Howto output}}}} documentation for an overview of
LAMMPS output options.

\sphinxAtStartPar
The vector values calculated by this compute are “extensive”, meaning
they scale with the number of atoms in the simulation.  They can be
divided by the appropriate volume to get a flux, which would then be an
“intensive” value, meaning independent of the number of atoms in the
simulation.  Note that if the compute group is “all”, then the
appropriate volume to divide by is the simulation box volume.  However,
if a group with a subset of atoms is used, it should be the volume
containing those atoms.

\sphinxAtStartPar
The vector values will be in energy*velocity {\hyperref[\detokenize{units::doc}]{\sphinxcrossref{\DUrole{doc}{units}}}}.  Once
divided by a volume the units will be that of flux, namely
energy/area/time {\hyperref[\detokenize{units::doc}]{\sphinxcrossref{\DUrole{doc}{units}}}}


\subsection{Restrictions}
\label{\detokenize{compute_heat_flux:restrictions}}\begin{quote}

\sphinxAtStartPar
none
\end{quote}


\subsection{Related commands}
\label{\detokenize{compute_heat_flux:related-commands}}
\sphinxAtStartPar
{\hyperref[\detokenize{fix_thermal_conductivity::doc}]{\sphinxcrossref{\DUrole{doc}{fix thermal/conductivity}}}},
{\hyperref[\detokenize{fix_ave_correlate::doc}]{\sphinxcrossref{\DUrole{doc}{fix ave/correlate}}}},
{\hyperref[\detokenize{variable::doc}]{\sphinxcrossref{\DUrole{doc}{variable}}}}


\subsection{Default}
\label{\detokenize{compute_heat_flux:default}}
\sphinxAtStartPar
none


\bigskip\hrule\bigskip



\subsubsection{Example Input File}
\label{\detokenize{compute_heat_flux:example-input-file}}
\begin{sphinxVerbatim}[commandchars=\\\{\}]
\PYG{c}{\PYGZsh{} Sample LAMMPS input script for thermal conductivity of solid Ar}

\PYG{k}{units}\PYG{+w}{       }\PYG{n}{real}
\PYG{k}{variable    }\PYG{n+nv+nvIdentifier}{T}\PYG{+w}{ }\PYG{n}{equal}\PYG{+w}{ }\PYG{l+m}{70}
\PYG{k}{variable    }\PYG{n+nv+nvIdentifier}{V}\PYG{+w}{ }\PYG{n}{equal}\PYG{+w}{ }\PYG{n}{vol}
\PYG{k}{variable    }\PYG{n+nv+nvIdentifier}{dt}\PYG{+w}{ }\PYG{n}{equal}\PYG{+w}{ }\PYG{l+m}{4.0}
\PYG{k}{variable    }\PYG{n+nv+nvIdentifier}{p}\PYG{+w}{ }\PYG{n}{equal}\PYG{+w}{ }\PYG{l+m}{200}\PYG{+w}{     }\PYG{c}{\PYGZsh{} correlation length}
\PYG{k}{variable    }\PYG{n+nv+nvIdentifier}{s}\PYG{+w}{ }\PYG{n}{equal}\PYG{+w}{ }\PYG{l+m}{10}\PYG{+w}{      }\PYG{c}{\PYGZsh{} sample interval}
\PYG{k}{variable    }\PYG{n+nv+nvIdentifier}{d}\PYG{+w}{ }\PYG{n}{equal}\PYG{+w}{ }\PYG{n+nv}{\PYGZdl{}p}\PYG{o}{*}\PYG{n+nv}{\PYGZdl{}s}\PYG{+w}{   }\PYG{c}{\PYGZsh{} dump interval}

\PYG{c}{\PYGZsh{} convert from LAMMPS real units to SI}

\PYG{k}{variable    }\PYG{n+nv+nvIdentifier}{kB}\PYG{+w}{ }\PYG{n}{equal}\PYG{+w}{ }\PYG{l+m}{1.3806504e\PYGZhy{}23}\PYG{+w}{    }\PYG{c}{\PYGZsh{} [J/K] Boltzmann}
\PYG{k}{variable    }\PYG{n+nv+nvIdentifier}{kCal2J}\PYG{+w}{ }\PYG{n}{equal}\PYG{+w}{ }\PYG{l+m}{4186.0}\PYG{o}{/}\PYG{l+m}{6.02214e23}
\PYG{k}{variable    }\PYG{n+nv+nvIdentifier}{A2m}\PYG{+w}{ }\PYG{n}{equal}\PYG{+w}{ }\PYG{l+m}{1.0e\PYGZhy{}10}
\PYG{k}{variable    }\PYG{n+nv+nvIdentifier}{fs2s}\PYG{+w}{ }\PYG{n}{equal}\PYG{+w}{ }\PYG{l+m}{1.0e\PYGZhy{}15}
\PYG{k}{variable    }\PYG{n+nv+nvIdentifier}{convert}\PYG{+w}{ }\PYG{n}{equal}\PYG{+w}{ }\PYG{n+nv}{\PYGZdl{}\PYGZob{}}\PYG{n+nv}{kCal2J}\PYG{n+nv}{\PYGZcb{}}\PYG{o}{*}\PYG{n+nv}{\PYGZdl{}\PYGZob{}}\PYG{n+nv}{kCal2J}\PYG{n+nv}{\PYGZcb{}}\PYG{o}{/}\PYG{n+nv}{\PYGZdl{}\PYGZob{}}\PYG{n+nv}{fs2s}\PYG{n+nv}{\PYGZcb{}}\PYG{o}{/}\PYG{n+nv}{\PYGZdl{}\PYGZob{}}\PYG{n+nv}{A2m}\PYG{n+nv}{\PYGZcb{}}

\PYG{c}{\PYGZsh{} setup problem}

\PYG{k}{dimension}\PYG{+w}{    }\PYG{l+m}{3}
\PYG{k}{boundary}\PYG{+w}{     }\PYG{n}{p}\PYG{+w}{ }\PYG{n}{p}\PYG{+w}{ }\PYG{n}{p}
\PYG{k}{lattice}\PYG{+w}{      }\PYG{n}{fcc}\PYG{+w}{ }\PYG{l+m}{5.376}\PYG{+w}{ }\PYG{n}{orient}\PYG{+w}{ }\PYG{n}{x}\PYG{+w}{ }\PYG{l+m}{1}\PYG{+w}{ }\PYG{l+m}{0}\PYG{+w}{ }\PYG{l+m}{0}\PYG{+w}{ }\PYG{n}{orient}\PYG{+w}{ }\PYG{n}{y}\PYG{+w}{ }\PYG{l+m}{0}\PYG{+w}{ }\PYG{l+m}{1}\PYG{+w}{ }\PYG{l+m}{0}\PYG{+w}{ }\PYG{n}{orient}\PYG{+w}{ }\PYG{n}{z}\PYG{+w}{ }\PYG{l+m}{0}\PYG{+w}{ }\PYG{l+m}{0}\PYG{+w}{ }\PYG{l+m}{1}
\PYG{k}{region       }\PYG{n+nv+nvIdentifier}{box}\PYG{+w}{ }\PYG{n}{block}\PYG{+w}{ }\PYG{l+m}{0}\PYG{+w}{ }\PYG{l+m}{4}\PYG{+w}{ }\PYG{l+m}{0}\PYG{+w}{ }\PYG{l+m}{4}\PYG{+w}{ }\PYG{l+m}{0}\PYG{+w}{ }\PYG{l+m}{4}
\PYG{k}{create\PYGZus{}box   }\PYG{n+nv+nvIdentifier}{1}\PYG{+w}{ }\PYG{n+nv+nvIdentifier}{box}
\PYG{k}{create\PYGZus{}atoms}\PYG{+w}{ }\PYG{l+m}{1}\PYG{+w}{ }\PYG{n}{box}
\PYG{k}{mass}\PYG{+w}{         }\PYG{l+m}{1}\PYG{+w}{ }\PYG{l+m}{39.948}
\PYG{k}{pair\PYGZus{}style}\PYG{+w}{   }\PYG{n}{lj}\PYG{o}{/}\PYG{n}{cut}\PYG{+w}{ }\PYG{l+m}{13.0}
\PYG{k}{pair\PYGZus{}coeff}\PYG{+w}{   }\PYG{o}{*}\PYG{+w}{ }\PYG{o}{*}\PYG{+w}{ }\PYG{l+m}{0.2381}\PYG{+w}{ }\PYG{l+m}{3.405}
\PYG{k}{timestep}\PYG{+w}{     }\PYG{n+nv}{\PYGZdl{}\PYGZob{}}\PYG{n+nv}{dt}\PYG{n+nv}{\PYGZcb{}}
\PYG{k}{thermo}\PYG{+w}{       }\PYG{n+nv}{\PYGZdl{}d}

\PYG{c}{\PYGZsh{} equilibration and thermalization}

\PYG{k}{velocity}\PYG{+w}{     }\PYG{n}{all}\PYG{+w}{ }\PYG{n}{create}\PYG{+w}{ }\PYG{n+nv}{\PYGZdl{}T}\PYG{+w}{ }\PYG{l+m}{102486}\PYG{+w}{ }\PYG{n}{mom}\PYG{+w}{ }\PYG{n}{yes}\PYG{+w}{ }\PYG{n}{rot}\PYG{+w}{ }\PYG{n}{yes}\PYG{+w}{ }\PYG{n}{dist}\PYG{+w}{ }\PYG{n}{gaussian}
\PYG{k}{fix          }\PYG{n+nv+nvIdentifier}{NVT}\PYG{+w}{ }\PYG{n+nv+nvIdentifier}{all}\PYG{+w}{ }\PYG{n}{nvt}\PYG{+w}{ }\PYG{n}{temp}\PYG{+w}{ }\PYG{n+nv}{\PYGZdl{}T}\PYG{+w}{ }\PYG{n+nv}{\PYGZdl{}T}\PYG{+w}{ }\PYG{l+m}{10}\PYG{+w}{ }\PYG{n}{drag}\PYG{+w}{ }\PYG{l+m}{0.2}
\PYG{k}{run}\PYG{+w}{          }\PYG{l+m}{8000}

\PYG{c}{\PYGZsh{} thermal conductivity calculation, switch to NVE if desired}

\PYG{c}{\PYGZsh{}unfix       NVT}
\PYG{c}{\PYGZsh{}fix         NVE all nve}

\PYG{k}{reset\PYGZus{}timestep}\PYG{+w}{ }\PYG{l+m}{0}
\PYG{k}{compute      }\PYG{n+nv+nvIdentifier}{myKE}\PYG{+w}{ }\PYG{n+nv+nvIdentifier}{all}\PYG{+w}{ }\PYG{n}{ke}\PYG{o}{/}\PYG{n}{atom}
\PYG{k}{compute      }\PYG{n+nv+nvIdentifier}{myPE}\PYG{+w}{ }\PYG{n+nv+nvIdentifier}{all}\PYG{+w}{ }\PYG{n}{pe}\PYG{o}{/}\PYG{n}{atom}
\PYG{k}{compute      }\PYG{n+nv+nvIdentifier}{myStress}\PYG{+w}{ }\PYG{n+nv+nvIdentifier}{all}\PYG{+w}{ }\PYG{n}{stress}\PYG{o}{/}\PYG{n}{atom}\PYG{+w}{ }\PYG{n}{NULL}\PYG{+w}{ }\PYG{n}{virial}
\PYG{k}{compute      }\PYG{n+nv+nvIdentifier}{flux}\PYG{+w}{ }\PYG{n+nv+nvIdentifier}{all}\PYG{+w}{ }\PYG{n}{heat}\PYG{o}{/}\PYG{n}{flux}\PYG{+w}{ }\PYG{n}{myKE}\PYG{+w}{ }\PYG{n}{myPE}\PYG{+w}{ }\PYG{n}{myStress}
\PYG{k}{variable     }\PYG{n+nv+nvIdentifier}{Jx}\PYG{+w}{ }\PYG{n}{equal}\PYG{+w}{ }\PYG{n}{c\PYGZus{}flux[1]}\PYG{o}{/}\PYG{n}{vol}
\PYG{k}{variable     }\PYG{n+nv+nvIdentifier}{Jy}\PYG{+w}{ }\PYG{n}{equal}\PYG{+w}{ }\PYG{n}{c\PYGZus{}flux[2]}\PYG{o}{/}\PYG{n}{vol}
\PYG{k}{variable     }\PYG{n+nv+nvIdentifier}{Jz}\PYG{+w}{ }\PYG{n}{equal}\PYG{+w}{ }\PYG{n}{c\PYGZus{}flux[3]}\PYG{o}{/}\PYG{n}{vol}
\PYG{k}{fix          }\PYG{n+nv+nvIdentifier}{JJ}\PYG{+w}{ }\PYG{n+nv+nvIdentifier}{all}\PYG{+w}{ }\PYG{n}{ave}\PYG{o}{/}\PYG{n}{correlate}\PYG{+w}{ }\PYG{n+nv}{\PYGZdl{}s}\PYG{+w}{ }\PYG{n+nv}{\PYGZdl{}p}\PYG{+w}{ }\PYG{n+nv}{\PYGZdl{}d}\PYG{+w}{ }\PYG{o}{\PYGZam{}}
\PYG{+w}{             }\PYG{n}{c\PYGZus{}flux[1]}\PYG{+w}{ }\PYG{n}{c\PYGZus{}flux[2]}\PYG{+w}{ }\PYG{n}{c\PYGZus{}flux[3]}\PYG{+w}{ }\PYG{n}{type}\PYG{+w}{ }\PYG{n}{auto}\PYG{+w}{ }\PYG{n}{file}\PYG{+w}{ }\PYG{n}{J0Jt.dat}\PYG{+w}{ }\PYG{n}{ave}\PYG{+w}{ }\PYG{n}{running}
\PYG{k}{variable     }\PYG{n+nv+nvIdentifier}{scale}\PYG{+w}{ }\PYG{n}{equal}\PYG{+w}{ }\PYG{n+nv}{\PYGZdl{}\PYGZob{}}\PYG{n+nv}{convert}\PYG{n+nv}{\PYGZcb{}}\PYG{o}{/}\PYG{n+nv}{\PYGZdl{}\PYGZob{}}\PYG{n+nv}{kB}\PYG{n+nv}{\PYGZcb{}}\PYG{o}{/}\PYG{n+nv}{\PYGZdl{}T}\PYG{o}{/}\PYG{n+nv}{\PYGZdl{}T}\PYG{o}{/}\PYG{n+nv}{\PYGZdl{}V}\PYG{o}{*}\PYG{n+nv}{\PYGZdl{}s}\PYG{o}{*}\PYG{n+nv}{\PYGZdl{}\PYGZob{}}\PYG{n+nv}{dt}\PYG{n+nv}{\PYGZcb{}}
\PYG{k}{variable     }\PYG{n+nv+nvIdentifier}{k11}\PYG{+w}{ }\PYG{n}{equal}\PYG{+w}{ }\PYG{n}{trap}\PYG{n+nv}{(}\PYG{n+nv}{f\PYGZus{}JJ[3]}\PYG{n+nv}{)}\PYG{o}{*}\PYG{n+nv}{\PYGZdl{}\PYGZob{}}\PYG{n+nv}{scale}\PYG{n+nv}{\PYGZcb{}}
\PYG{k}{variable     }\PYG{n+nv+nvIdentifier}{k22}\PYG{+w}{ }\PYG{n}{equal}\PYG{+w}{ }\PYG{n}{trap}\PYG{n+nv}{(}\PYG{n+nv}{f\PYGZus{}JJ[4]}\PYG{n+nv}{)}\PYG{o}{*}\PYG{n+nv}{\PYGZdl{}\PYGZob{}}\PYG{n+nv}{scale}\PYG{n+nv}{\PYGZcb{}}
\PYG{k}{variable     }\PYG{n+nv+nvIdentifier}{k33}\PYG{+w}{ }\PYG{n}{equal}\PYG{+w}{ }\PYG{n}{trap}\PYG{n+nv}{(}\PYG{n+nv}{f\PYGZus{}JJ[5]}\PYG{n+nv}{)}\PYG{o}{*}\PYG{n+nv}{\PYGZdl{}\PYGZob{}}\PYG{n+nv}{scale}\PYG{n+nv}{\PYGZcb{}}
\PYG{k}{thermo\PYGZus{}style}\PYG{+w}{ }\PYG{n}{custom}\PYG{+w}{ }\PYG{n}{step}\PYG{+w}{ }\PYG{n}{temp}\PYG{+w}{ }\PYG{n}{v\PYGZus{}Jx}\PYG{+w}{ }\PYG{n}{v\PYGZus{}Jy}\PYG{+w}{ }\PYG{n}{v\PYGZus{}Jz}\PYG{+w}{ }\PYG{n}{v\PYGZus{}k11}\PYG{+w}{ }\PYG{n}{v\PYGZus{}k22}\PYG{+w}{ }\PYG{n}{v\PYGZus{}k33}
\PYG{k}{run}\PYG{+w}{          }\PYG{l+m}{100000}
\PYG{k}{variable     }\PYG{n+nv+nvIdentifier}{k}\PYG{+w}{ }\PYG{n}{equal}\PYG{+w}{ }\PYG{n+nv}{(}\PYG{n+nv}{v\PYGZus{}k11+v\PYGZus{}k22+v\PYGZus{}k33}\PYG{n+nv}{)}\PYG{o}{/}\PYG{l+m}{3.0}
\PYG{k}{variable     }\PYG{n+nv+nvIdentifier}{ndens}\PYG{+w}{ }\PYG{n}{equal}\PYG{+w}{ }\PYG{n}{count}\PYG{n+nv}{(}\PYG{n+nv}{all}\PYG{n+nv}{)}\PYG{o}{/}\PYG{n}{vol}
\PYG{k}{print}\PYG{+w}{        }\PYG{l+s}{\PYGZdq{}}\PYG{l+s}{average conductivity: \PYGZdl{}k[W/mK] @ \PYGZdl{}T K, \PYGZdl{}\PYGZob{}ndens\PYGZcb{} /A\PYGZbs{}\PYGZca{}3}\PYG{l+s}{\PYGZdq{}}
\end{sphinxVerbatim}


\bigskip\hrule\bigskip

\phantomsection\label{\detokenize{compute_heat_flux:surblys3}}
\sphinxAtStartPar
\sphinxstylestrong{(Surblys2019)} Surblys, Matsubara, Kikugawa, Ohara, Phys Rev E, 99, 051301(R) (2019).

\phantomsection\label{\detokenize{compute_heat_flux:boone}}
\sphinxAtStartPar
\sphinxstylestrong{(Boone)} Boone, Babaei, Wilmer, J Chem Theory Comput, 15, 5579\textendash{}5587 (2019).

\phantomsection\label{\detokenize{compute_heat_flux:surblys4}}
\sphinxAtStartPar
\sphinxstylestrong{(Surblys2021)} Surblys, Matsubara, Kikugawa, Ohara, J Appl Phys 130, 215104 (2021).

\sphinxstepscope

\index{compute hexorder/atom@\spxentry{compute hexorder/atom}}\ignorespaces 

\section{compute hexorder/atom command}
\label{\detokenize{compute_hexorder_atom:compute-hexorder-atom-command}}\label{\detokenize{compute_hexorder_atom:index-0}}\label{\detokenize{compute_hexorder_atom::doc}}

\subsection{Syntax}
\label{\detokenize{compute_hexorder_atom:syntax}}
\begin{sphinxVerbatim}[commandchars=\\\{\}]
\PYG{k}{compute }\PYG{n+nv+nvIdentifier}{ID}\PYG{+w}{ }\PYG{n+nv+nvIdentifier}{group\PYGZhy{}ID}\PYG{+w}{ }\PYG{n}{hexorder}\PYG{o}{/}\PYG{n}{atom}\PYG{+w}{ }\PYG{n}{keyword}\PYG{+w}{ }\PYG{n}{values}\PYG{+w}{ }\PYG{n}{...}
\end{sphinxVerbatim}
\begin{itemize}
\item {} 
\sphinxAtStartPar
ID, group\sphinxhyphen{}ID are documented in {\hyperref[\detokenize{compute::doc}]{\sphinxcrossref{\DUrole{doc}{compute}}}} command

\item {} 
\sphinxAtStartPar
hexorder/atom = style name of this compute command

\item {} 
\sphinxAtStartPar
one or more keyword/value pairs may be appended
\begin{sphinxalltt}
keyword = \sphinxstyleemphasis{degree} or \sphinxstyleemphasis{nnn} or \sphinxstyleemphasis{cutoff}
  \sphinxstyleemphasis{cutoff} value = distance cutoff
  \sphinxstyleemphasis{nnn} value = number of nearest neighbors
  \sphinxstyleemphasis{degree} value = degree \sphinxstyleemphasis{n} of order parameter
\end{sphinxalltt}

\end{itemize}


\subsection{Examples}
\label{\detokenize{compute_hexorder_atom:examples}}
\begin{sphinxVerbatim}[commandchars=\\\{\}]
\PYG{k}{compute }\PYG{n+nv+nvIdentifier}{1}\PYG{+w}{ }\PYG{n+nv+nvIdentifier}{all}\PYG{+w}{ }\PYG{n}{hexorder}\PYG{o}{/}\PYG{n}{atom}
\PYG{k}{compute }\PYG{n+nv+nvIdentifier}{1}\PYG{+w}{ }\PYG{n+nv+nvIdentifier}{all}\PYG{+w}{ }\PYG{n}{hexorder}\PYG{o}{/}\PYG{n}{atom}\PYG{+w}{ }\PYG{n}{degree}\PYG{+w}{ }\PYG{l+m}{4}\PYG{+w}{ }\PYG{n}{nnn}\PYG{+w}{ }\PYG{l+m}{4}\PYG{+w}{ }\PYG{n}{cutoff}\PYG{+w}{ }\PYG{l+m}{1.2}
\end{sphinxVerbatim}


\subsection{Description}
\label{\detokenize{compute_hexorder_atom:description}}
\sphinxAtStartPar
Define a computation that calculates \(q_n\) the bond\sphinxhyphen{}orientational
order parameter for each atom in a group. The hexatic (\sphinxstyleemphasis{n} = 6) order
parameter was introduced by {\hyperref[\detokenize{compute_hexorder_atom:nelson}]{\sphinxcrossref{\DUrole{std,std-ref}{Nelson and Halperin}}}} as a way to detect
hexagonal symmetry in two\sphinxhyphen{}dimensional systems. For each atom, \(q_n\)
is a complex number (stored as two real numbers) defined as follows:
\begin{equation*}
\begin{split}q_n = \frac{1}{nnn}\sum_{j = 1}^{nnn} e^{n i \theta({\bf r}_{ij})}\end{split}
\end{equation*}
\sphinxAtStartPar
where the sum is over the \sphinxstyleemphasis{nnn} nearest neighbors
of the central atom. The angle \(\theta\)
is formed by the bond vector \(r_{ij}\) and the \sphinxstyleemphasis{x} axis.
\(\theta\) is calculated only using the \sphinxstyleemphasis{x} and \sphinxstyleemphasis{y} components,
whereas the distance from the
central atom is calculated using all three
\sphinxstyleemphasis{x}, \sphinxstyleemphasis{y}, and \sphinxstyleemphasis{z} components of the bond vector.
Neighbor atoms not in the group
are included in the order parameter of atoms in the group.

\sphinxAtStartPar
The optional keyword \sphinxstyleemphasis{cutoff} defines the distance cutoff
used when searching for neighbors. The default value, also
the maximum allowable value, is the cutoff specified
by the pair style.

\sphinxAtStartPar
The optional keyword \sphinxstyleemphasis{nnn} defines the number of nearest
neighbors used to calculate \(q_n\). The default value is 6.
If the value is NULL, then all neighbors up to the
distance cutoff are used.

\sphinxAtStartPar
The optional keyword \sphinxstyleemphasis{degree} sets the degree \sphinxstyleemphasis{n} of the order parameter.
The default value is 6. For a perfect hexagonal lattice with
\sphinxstyleemphasis{nnn} = 6,
\(q_6 = e^{6 i \phi}\) for all atoms, where the constant \(0 < \phi < \frac{\pi}{3}\)
depends only on the orientation of the lattice relative to the \sphinxstyleemphasis{x} axis.
In an isotropic liquid, local neighborhoods may still exhibit
weak hexagonal symmetry, but because the orientational correlation
decays quickly with distance, the value of \(\phi\) will be different for
different atoms, and so when \(q_6\) is averaged over all the atoms
in the system, \(| \left< q_6 \right> | << 1\).

\sphinxAtStartPar
The value of \(q_n\) is set to zero for atoms not in the
specified compute group, as well as for atoms that have less than
\sphinxstyleemphasis{nnn} neighbors within the distance cutoff.

\sphinxAtStartPar
The neighbor list needed to compute this quantity is constructed each
time the calculation is performed (i.e. each time a snapshot of atoms
is dumped).  Thus it can be inefficient to compute/dump this quantity
too frequently.

\begin{sphinxadmonition}{note}{Note:}
\sphinxAtStartPar
If you have a bonded system, then the settings of
{\hyperref[\detokenize{special_bonds::doc}]{\sphinxcrossref{\DUrole{doc}{special\_bonds}}}} command can remove pairwise
interactions between atoms in the same bond, angle, or dihedral.  This
is the default setting for the {\hyperref[\detokenize{special_bonds::doc}]{\sphinxcrossref{\DUrole{doc}{special\_bonds}}}}
command, and means those pairwise interactions do not appear in the
neighbor list.  Because this fix uses the neighbor list, it also means
those pairs will not be included in the order parameter.  This
difficulty can be circumvented by writing a dump file, and using the
{\hyperref[\detokenize{rerun::doc}]{\sphinxcrossref{\DUrole{doc}{rerun}}}} command to compute the order parameter for
snapshots in the dump file.  The rerun script can use a
{\hyperref[\detokenize{special_bonds::doc}]{\sphinxcrossref{\DUrole{doc}{special\_bonds}}}} command that includes all pairs in
the neighbor list.
\end{sphinxadmonition}


\subsection{Output info}
\label{\detokenize{compute_hexorder_atom:output-info}}
\sphinxAtStartPar
This compute calculates a per\sphinxhyphen{}atom array with 2 columns, giving the
real and imaginary parts \(q_n\), a complex number restricted to the
unit disk of the complex plane (i.e., \(\Re(q_n)^2 + \Im(q_n)^2 \le 1\)).

\sphinxAtStartPar
These values can be accessed by any command that uses per\sphinxhyphen{}atom values
from a compute as input.  See the {\hyperref[\detokenize{Howto_output::doc}]{\sphinxcrossref{\DUrole{doc}{Howto output}}}} doc
page for an overview of LAMMPS output options.


\subsection{Restrictions}
\label{\detokenize{compute_hexorder_atom:restrictions}}
\sphinxAtStartPar
This compute is part of the EXTRA\sphinxhyphen{}COMPUTE package.  It is only enabled if
LAMMPS was built with that package.  See the {\hyperref[\detokenize{Build_package::doc}]{\sphinxcrossref{\DUrole{doc}{Build package}}}} page for more info.


\subsection{Related commands}
\label{\detokenize{compute_hexorder_atom:related-commands}}
\sphinxAtStartPar
{\hyperref[\detokenize{compute_orientorder_atom::doc}]{\sphinxcrossref{\DUrole{doc}{compute orientorder/atom}}}}, {\hyperref[\detokenize{compute_coord_atom::doc}]{\sphinxcrossref{\DUrole{doc}{compute coord/atom}}}}, {\hyperref[\detokenize{compute_centro_atom::doc}]{\sphinxcrossref{\DUrole{doc}{compute centro/atom}}}}


\subsection{Default}
\label{\detokenize{compute_hexorder_atom:default}}
\sphinxAtStartPar
The option defaults are \sphinxstyleemphasis{cutoff} = pair style cutoff, \sphinxstyleemphasis{nnn} = 6, \sphinxstyleemphasis{degree} = 6


\bigskip\hrule\bigskip

\phantomsection\label{\detokenize{compute_hexorder_atom:nelson}}
\sphinxAtStartPar
\sphinxstylestrong{(Nelson)} Nelson, Halperin, Phys Rev B, 19, 2457 (1979).

\sphinxstepscope

\index{compute hma@\spxentry{compute hma}}\ignorespaces 

\section{compute hma command}
\label{\detokenize{compute_hma:compute-hma-command}}\label{\detokenize{compute_hma:index-0}}\label{\detokenize{compute_hma::doc}}

\subsection{Syntax}
\label{\detokenize{compute_hma:syntax}}
\begin{sphinxVerbatim}[commandchars=\\\{\}]
\PYG{k}{compute }\PYG{n+nv+nvIdentifier}{ID}\PYG{+w}{ }\PYG{n+nv+nvIdentifier}{group\PYGZhy{}ID}\PYG{+w}{ }\PYG{n}{hma}\PYG{+w}{ }\PYG{n}{temp}\PYG{o}{\PYGZhy{}}\PYG{n}{ID}\PYG{+w}{ }\PYG{n}{keyword}\PYG{+w}{ }\PYG{n}{...}
\end{sphinxVerbatim}
\begin{itemize}
\item {} 
\sphinxAtStartPar
ID, group\sphinxhyphen{}ID are documented in {\hyperref[\detokenize{compute::doc}]{\sphinxcrossref{\DUrole{doc}{compute}}}} command

\item {} 
\sphinxAtStartPar
hma = style name of this compute command

\item {} 
\sphinxAtStartPar
temp\sphinxhyphen{}ID = ID of fix that specifies the set temperature during canonical simulation

\item {} 
\sphinxAtStartPar
one or more keywords or keyword/argument pairs must be appended

\item {} 
\sphinxAtStartPar
keyword = \sphinxstyleemphasis{anharmonic} or \sphinxstyleemphasis{u} or \sphinxstyleemphasis{p} or \sphinxstyleemphasis{cv}

\end{itemize}
\begin{sphinxalltt}
\sphinxstyleemphasis{anharmonic} = compute will return anharmonic property values
\sphinxstyleemphasis{u} = compute will return potential energy
\sphinxstyleemphasis{p} value = Pharm = compute will return pressure
   Pharm = difference between the harmonic pressure and lattice pressure
           as described below
\sphinxstyleemphasis{cv} = compute will return the heat capacity
\end{sphinxalltt}


\subsection{Examples}
\label{\detokenize{compute_hma:examples}}
\begin{sphinxVerbatim}[commandchars=\\\{\}]
\PYG{k}{compute }\PYG{n+nv+nvIdentifier}{2}\PYG{+w}{ }\PYG{n+nv+nvIdentifier}{all}\PYG{+w}{ }\PYG{n}{hma}\PYG{+w}{ }\PYG{l+m}{1}\PYG{+w}{ }\PYG{n}{u}
\PYG{k}{compute }\PYG{n+nv+nvIdentifier}{2}\PYG{+w}{ }\PYG{n+nv+nvIdentifier}{all}\PYG{+w}{ }\PYG{n}{hma}\PYG{+w}{ }\PYG{l+m}{1}\PYG{+w}{ }\PYG{n}{anharmonic}\PYG{+w}{ }\PYG{n}{u}\PYG{+w}{ }\PYG{n}{p}\PYG{+w}{ }\PYG{l+m}{0.9}
\PYG{k}{compute }\PYG{n+nv+nvIdentifier}{2}\PYG{+w}{ }\PYG{n+nv+nvIdentifier}{all}\PYG{+w}{ }\PYG{n}{hma}\PYG{+w}{ }\PYG{l+m}{1}\PYG{+w}{ }\PYG{n}{u}\PYG{+w}{ }\PYG{n}{cv}
\end{sphinxVerbatim}


\subsection{Description}
\label{\detokenize{compute_hma:description}}
\sphinxAtStartPar
Define a computation that calculates the properties of a solid (potential
energy, pressure or heat capacity), using the harmonically\sphinxhyphen{}mapped averaging
(HMA) method.
This command yields much higher precision than the equivalent compute commands
({\hyperref[\detokenize{compute_pe::doc}]{\sphinxcrossref{\DUrole{doc}{compute pe}}}}, {\hyperref[\detokenize{compute_pressure::doc}]{\sphinxcrossref{\DUrole{doc}{compute pressure}}}}, etc.)
commands during a canonical simulation of an atomic crystal. Specifically,
near melting HMA can yield averages of a given precision an order of magnitude
faster than conventional methods, and this only improves as the temperatures is
lowered.  This is particularly important for evaluating the free energy by
thermodynamic integration, where the low\sphinxhyphen{}temperature contributions are the
greatest source of statistical uncertainty.  Moreover, HMA has other
advantages, including smaller potential\sphinxhyphen{}truncation effects, finite\sphinxhyphen{}size
effects, smaller timestep inaccuracy, faster equilibration and shorter
decorrelation time.

\sphinxAtStartPar
HMA should not be used if atoms are expected to diffuse.  It is also
restricted to simulations in the NVT ensemble.  While this compute may
be used with any potential in LAMMPS, it will provide inaccurate results
for potentials that do not go to 0 at the truncation distance;
{\hyperref[\detokenize{pair_lj_smooth_linear::doc}]{\sphinxcrossref{\DUrole{doc}{pair\_style lj/smooth/linear}}}} and Ewald
summation should work fine, while {\hyperref[\detokenize{pair_lj::doc}]{\sphinxcrossref{\DUrole{doc}{pair\_style lj/cut}}}}
will perform poorly unless the potential is shifted (via
{\hyperref[\detokenize{pair_modify::doc}]{\sphinxcrossref{\DUrole{doc}{pair\_modify}}}} shift) or the cutoff is large.
Furthermore, computation of the heat capacity with this compute is
restricted to those that implement the \sphinxstyleemphasis{single\_hessian} method in Pair.
Implementing \sphinxstyleemphasis{single\_hessian} in additional pair styles is simple.
Please contact Andrew Schultz (ajs42 at buffalo.edu) and David Kofke
(kofke at buffalo.edu) if your desired pair style does not have this
method.  This is the list of pair styles that currently implement
\sphinxstyleemphasis{single\_hessian}:
\begin{itemize}
\item {} 
\sphinxAtStartPar
{\hyperref[\detokenize{pair_lj_smooth_linear::doc}]{\sphinxcrossref{\DUrole{doc}{pair\_style lj/smooth/linear}}}}

\end{itemize}

\sphinxAtStartPar
In this method, the analytically known harmonic behavior of a crystal is removed from the traditional ensemble
averages, which leads to an accurate and precise measurement of the anharmonic contributions without contamination
by noise produced by the already\sphinxhyphen{}known harmonic behavior.
A detailed description of this method can be found in ({\hyperref[\detokenize{compute_hma:hma-moustafa}]{\sphinxcrossref{\DUrole{std,std-ref}{Moustafa}}}}). The potential energy is computed by the formula:
\begin{equation*}
\begin{split}\left< U\right>_\text{HMA} = \frac{d}{2} (N-1) k_B T  + \left< U + \frac{1}{2} \vec F\cdot\Delta \vec r \right>\end{split}
\end{equation*}
\sphinxAtStartPar
where \(N\) is the number of atoms in the system, \(k_B\) is Boltzmann’s
constant, \(T\) is the temperature, \(d\) is the dimensionality of the
system (2 or 3 for 2d/3d), \(\vec F\cdot\Delta\vec r\) is the sum of dot
products of the atomic force vectors and displacement (from lattice sites)
vectors, and \(U\) is the sum of pair, bond, angle, dihedral, improper,
kspace (long\sphinxhyphen{}range), and fix energies.

\sphinxAtStartPar
The pressure is computed by the formula:
\begin{equation*}
\begin{split}\left< P\right>_{HMA} = \Delta \hat P + \left< P_\text{vir}
+ \frac{\beta \Delta \hat P - \rho}{d(N-1)} \vec F\cdot\Delta \vec r \right>\end{split}
\end{equation*}
\sphinxAtStartPar
where \(\rho\) is the number density of the system, \(\Delta \hat P\)
is the difference between the harmonic and lattice pressure,
\(P_\text{vir}\) is the virial pressure computed as the sum of pair, bond,
angle, dihedral, improper, kspace (long\sphinxhyphen{}range), and fix contributions to the
force on each atom, and \(k_B=1/k_B T\).  Although the method will work for
any value of \(\Delta \hat P\)
specified (use pressure {\hyperref[\detokenize{units::doc}]{\sphinxcrossref{\DUrole{doc}{units}}}}), the precision of the resultant
pressure is sensitive to \(\Delta \hat P\); the precision tends to be
best when \(\Delta \hat P\) is the actual the difference between the
lattice pressure and harmonic pressure.
\begin{equation*}
\begin{split}\left<C_V \right>_\text{HMA} = \frac{d}{2} (N-1) k_B
 + \frac{1}{k_B T^2} \left( \left<U_\text{HMA}^2 \right>
                      - \left<U_\text{HMA}\right>^2 \right) + \frac{1}{4 T}
\left<\vec F\cdot\Delta\vec r + \Delta r \cdot\Phi\cdot \Delta\vec r\right>\end{split}
\end{equation*}
\sphinxAtStartPar
where \(\Phi\) is the Hessian matrix. The compute hma command
computes the full expression for \(C_V\) except for the
\(\left<U_\text{HMA}\right>^2\) in the variance term, which can be obtained
by passing the \sphinxstyleemphasis{u} keyword; you must add this extra contribution to the
\(C_V\) value reported by this compute.  The variance term can cause
significant round\sphinxhyphen{}off error when computing \(C_V\).  To address this, the
\sphinxstyleemphasis{anharmonic} keyword can be passed and/or the output format can be specified
with more digits.

\begin{sphinxVerbatim}[commandchars=\\\{\}]
\PYG{k}{thermo\PYGZus{}modify}\PYG{+w}{ }\PYG{n}{format}\PYG{+w}{ }\PYG{n}{float}\PYG{+w}{ }\PYG{l+s}{\PYGZsq{}}\PYG{l+s}{\PYGZpc{}22.15e}\PYG{l+s}{\PYGZsq{}}
\end{sphinxVerbatim}

\sphinxAtStartPar
The \sphinxstyleemphasis{anharmonic} keyword will instruct the compute to return anharmonic
properties rather than the full properties, which include lattice, harmonic
and anharmonic contributions.
When using this keyword, the compute must be first active (it must be included
via a {\hyperref[\detokenize{thermo_style::doc}]{\sphinxcrossref{\DUrole{doc}{thermo\_style custom}}}} command) while the atoms are
still at their lattice sites (before equilibration).

\sphinxAtStartPar
The temp\sphinxhyphen{}ID specified with compute hma command should be same as the fix\sphinxhyphen{}ID of
the Nose\textendash{}Hoover ({\hyperref[\detokenize{fix_nh::doc}]{\sphinxcrossref{\DUrole{doc}{fix nvt}}}}) or
Berendsen ({\hyperref[\detokenize{fix_temp_berendsen::doc}]{\sphinxcrossref{\DUrole{doc}{fix temp/berendsen}}}}) thermostat used for
the simulation. While using this command, the Langevin thermostat
({\hyperref[\detokenize{fix_langevin::doc}]{\sphinxcrossref{\DUrole{doc}{fix langevin}}}})
should be avoided as its extra forces interfere with the HMA implementation.

\begin{sphinxadmonition}{note}{Note:}
\sphinxAtStartPar
Compute hma command should be used right after the energy minimization, when the atoms are at their lattice sites.
The simulation should not be started before this command has been used in the input script.
\end{sphinxadmonition}

\sphinxAtStartPar
The following example illustrates the placement of this command in the input script:

\begin{sphinxVerbatim}[commandchars=\\\{\}]
\PYG{k}{min\PYGZus{}style}\PYG{+w}{ }\PYG{n}{cg}
\PYG{k}{minimize}\PYG{+w}{ }\PYG{l+m}{1e\PYGZhy{}35}\PYG{+w}{ }\PYG{l+m}{1e\PYGZhy{}15}\PYG{+w}{ }\PYG{l+m}{50000}\PYG{+w}{ }\PYG{l+m}{500000}
\PYG{k}{compute }\PYG{n+nv+nvIdentifier}{1}\PYG{+w}{ }\PYG{n+nv+nvIdentifier}{all}\PYG{+w}{ }\PYG{n}{hma}\PYG{+w}{ }\PYG{n}{thermostatid}\PYG{+w}{ }\PYG{n}{u}
\PYG{k}{fix }\PYG{n+nv+nvIdentifier}{thermostatid}\PYG{+w}{ }\PYG{n+nv+nvIdentifier}{all}\PYG{+w}{ }\PYG{n}{nvt}\PYG{+w}{ }\PYG{n}{temp}\PYG{+w}{ }\PYG{l+m}{600.0}\PYG{+w}{ }\PYG{l+m}{600.0}\PYG{+w}{ }\PYG{l+m}{100.0}
\end{sphinxVerbatim}

\begin{sphinxadmonition}{note}{Note:}
\sphinxAtStartPar
Compute hma should be used when the atoms of the solid do not diffuse. Diffusion will reduce the precision in the potential energy computation.
\end{sphinxadmonition}

\begin{sphinxadmonition}{note}{Note:}
\sphinxAtStartPar
The {\hyperref[\detokenize{fix_modify::doc}]{\sphinxcrossref{\DUrole{doc}{fix\_modify energy yes}}}} command must also be specified if a fix is to contribute potential energy to this command.
\end{sphinxadmonition}

\sphinxAtStartPar
An example input script that uses this compute is included in
examples/PACKAGES/hma/ along with corresponding LAMMPS output showing that the HMA
properties fluctuate less than the corresponding conventional properties.


\subsection{Output info}
\label{\detokenize{compute_hma:output-info}}
\sphinxAtStartPar
This compute calculates a global vector that includes the n properties
requested as arguments to the command (the potential energy, pressure and/or heat
capacity).  The elements of the vector can be accessed by indices 1\textendash{}n by any
command that uses global vector values as input.  See the {\hyperref[\detokenize{Howto_output::doc}]{\sphinxcrossref{\DUrole{doc}{Howto output}}}} page for an overview of LAMMPS output options.

\sphinxAtStartPar
The vector values calculated by this compute are “extensive”.  The
scalar value will be in energy {\hyperref[\detokenize{units::doc}]{\sphinxcrossref{\DUrole{doc}{units}}}}.


\subsection{Restrictions}
\label{\detokenize{compute_hma:restrictions}}
\sphinxAtStartPar
This compute is part of the EXTRA\sphinxhyphen{}COMPUTE package.  It is enabled only
if LAMMPS was built with that package.  See the {\hyperref[\detokenize{Build_package::doc}]{\sphinxcrossref{\DUrole{doc}{Build package}}}} page for more info.

\sphinxAtStartPar
Usage restricted to canonical (NVT) ensemble simulation only.


\subsection{Related commands}
\label{\detokenize{compute_hma:related-commands}}
\sphinxAtStartPar
{\hyperref[\detokenize{compute_pe::doc}]{\sphinxcrossref{\DUrole{doc}{compute pe}}}}, {\hyperref[\detokenize{compute_pressure::doc}]{\sphinxcrossref{\DUrole{doc}{compute pressure}}}}

\sphinxAtStartPar
{\hyperref[\detokenize{dynamical_matrix::doc}]{\sphinxcrossref{\DUrole{doc}{dynamical matrix}}}} provides a finite difference
formulation of the Hessian provided by Pair’s single\_hessian, which is used by
this compute.


\subsection{Default}
\label{\detokenize{compute_hma:default}}
\sphinxAtStartPar
none


\bigskip\hrule\bigskip

\phantomsection\label{\detokenize{compute_hma:hma-moustafa}}
\sphinxAtStartPar
\sphinxstylestrong{(Moustafa)} Sabry G. Moustafa, Andrew J. Schultz, and David A. Kofke, \sphinxstyleemphasis{Very fast averaging of thermal properties of crystals by molecular simulation},
\sphinxhref{https://link.aps.org/doi/10.1103/PhysRevE.92.043303}{Phys. Rev. E {[}92{]}, 043303 (2015)}

\sphinxstepscope

\index{compute improper@\spxentry{compute improper}}\ignorespaces 

\section{compute improper command}
\label{\detokenize{compute_improper:compute-improper-command}}\label{\detokenize{compute_improper:index-0}}\label{\detokenize{compute_improper::doc}}

\subsection{Syntax}
\label{\detokenize{compute_improper:syntax}}
\begin{sphinxVerbatim}[commandchars=\\\{\}]
\PYG{k}{compute }\PYG{n+nv+nvIdentifier}{ID}\PYG{+w}{ }\PYG{n+nv+nvIdentifier}{group\PYGZhy{}ID}\PYG{+w}{ }\PYG{n}{improper}
\end{sphinxVerbatim}
\begin{itemize}
\item {} 
\sphinxAtStartPar
ID, group\sphinxhyphen{}ID are documented in {\hyperref[\detokenize{compute::doc}]{\sphinxcrossref{\DUrole{doc}{compute}}}} command

\item {} 
\sphinxAtStartPar
improper = style name of this compute command

\end{itemize}


\subsection{Examples}
\label{\detokenize{compute_improper:examples}}
\begin{sphinxVerbatim}[commandchars=\\\{\}]
\PYG{k}{compute }\PYG{n+nv+nvIdentifier}{1}\PYG{+w}{ }\PYG{n+nv+nvIdentifier}{all}\PYG{+w}{ }\PYG{n}{improper}
\end{sphinxVerbatim}


\subsection{Description}
\label{\detokenize{compute_improper:description}}
\sphinxAtStartPar
Define a computation that extracts the improper energy calculated by
each of the improper sub\sphinxhyphen{}styles used in the {\hyperref[\detokenize{improper_hybrid::doc}]{\sphinxcrossref{\DUrole{doc}{improper\_style hybrid}}}} command.  These values are made
accessible for output or further processing by other commands.  The
group specified for this command is ignored.

\sphinxAtStartPar
This compute is useful when using {\hyperref[\detokenize{improper_hybrid::doc}]{\sphinxcrossref{\DUrole{doc}{improper\_style hybrid}}}} if you want to know the portion of the
total energy contributed by one or more of the hybrid sub\sphinxhyphen{}styles.


\subsection{Output info}
\label{\detokenize{compute_improper:output-info}}
\sphinxAtStartPar
This compute calculates a global vector of length \(N\), where \(N\) is
the number of sub\_styles defined by the
{\hyperref[\detokenize{improper_style::doc}]{\sphinxcrossref{\DUrole{doc}{improper\_style hybrid}}}} command.
These styles can be accessed by the indices 1 through \(N\).
These values can be used by any command that uses global scalar or vector
values from a compute as input.  See the {\hyperref[\detokenize{Howto_output::doc}]{\sphinxcrossref{\DUrole{doc}{Howto output}}}}
page for an overview of LAMMPS output options.

\sphinxAtStartPar
The vector values are “extensive” and will be in energy
{\hyperref[\detokenize{units::doc}]{\sphinxcrossref{\DUrole{doc}{units}}}}.


\subsection{Restrictions}
\label{\detokenize{compute_improper:restrictions}}\begin{quote}

\sphinxAtStartPar
none
\end{quote}


\subsection{Related commands}
\label{\detokenize{compute_improper:related-commands}}
\sphinxAtStartPar
{\hyperref[\detokenize{compute_pe::doc}]{\sphinxcrossref{\DUrole{doc}{compute pe}}}}, {\hyperref[\detokenize{compute_pair::doc}]{\sphinxcrossref{\DUrole{doc}{compute pair}}}}


\subsection{Default}
\label{\detokenize{compute_improper:default}}
\sphinxAtStartPar
none

\sphinxstepscope

\index{compute improper/local@\spxentry{compute improper/local}}\ignorespaces 

\section{compute improper/local command}
\label{\detokenize{compute_improper_local:compute-improper-local-command}}\label{\detokenize{compute_improper_local:index-0}}\label{\detokenize{compute_improper_local::doc}}

\subsection{Syntax}
\label{\detokenize{compute_improper_local:syntax}}
\begin{sphinxVerbatim}[commandchars=\\\{\}]
\PYG{k}{compute }\PYG{n+nv+nvIdentifier}{ID}\PYG{+w}{ }\PYG{n+nv+nvIdentifier}{group\PYGZhy{}ID}\PYG{+w}{ }\PYG{n}{improper}\PYG{o}{/}\PYG{n}{local}\PYG{+w}{ }\PYG{n}{value1}\PYG{+w}{ }\PYG{n}{value2}\PYG{+w}{ }\PYG{n}{...}
\end{sphinxVerbatim}
\begin{itemize}
\item {} 
\sphinxAtStartPar
ID, group\sphinxhyphen{}ID are documented in {\hyperref[\detokenize{compute::doc}]{\sphinxcrossref{\DUrole{doc}{compute}}}} command

\item {} 
\sphinxAtStartPar
improper/local = style name of this compute command

\item {} 
\sphinxAtStartPar
one or more values may be appended

\item {} 
\sphinxAtStartPar
value = \sphinxstyleemphasis{chi}
\begin{sphinxalltt}
\sphinxstyleemphasis{chi} = tabulate improper angles
\end{sphinxalltt}

\end{itemize}


\subsection{Examples}
\label{\detokenize{compute_improper_local:examples}}
\begin{sphinxVerbatim}[commandchars=\\\{\}]
\PYG{k}{compute }\PYG{n+nv+nvIdentifier}{1}\PYG{+w}{ }\PYG{n+nv+nvIdentifier}{all}\PYG{+w}{ }\PYG{n}{improper}\PYG{o}{/}\PYG{n}{local}\PYG{+w}{ }\PYG{n}{chi}
\end{sphinxVerbatim}


\subsection{Description}
\label{\detokenize{compute_improper_local:description}}
\sphinxAtStartPar
Define a computation that calculates properties of individual improper
interactions.  The number of datums generated, aggregated across all
processors, equals the number of impropers in the system, modified by
the group parameter as explained below.

\sphinxAtStartPar
The value \sphinxstyleemphasis{chi} is the improper angle, as defined in the doc pages for
the individual improper styles listed on
{\hyperref[\detokenize{improper_style::doc}]{\sphinxcrossref{\DUrole{doc}{improper\_style}}}} doc page.

\sphinxAtStartPar
The local data stored by this command is generated by looping over all
the atoms owned on a processor and their impropers.  An improper will
only be included if all four atoms in the improper are in the specified
compute group.

\sphinxAtStartPar
Note that as atoms migrate from processor to processor, there will be
no consistent ordering of the entries within the local vector or array
from one timestep to the next.  The only consistency that is
guaranteed is that the ordering on a particular timestep will be the
same for local vectors or arrays generated by other compute commands.
For example, improper output from the {\hyperref[\detokenize{compute_property_local::doc}]{\sphinxcrossref{\DUrole{doc}{compute property/local}}}} command can be combined
with data from this command and output by the {\hyperref[\detokenize{dump::doc}]{\sphinxcrossref{\DUrole{doc}{dump local}}}}
command in a consistent way.

\sphinxAtStartPar
Here is an example of how to do this:

\begin{sphinxVerbatim}[commandchars=\\\{\}]
\PYG{k}{compute }\PYG{n+nv+nvIdentifier}{1}\PYG{+w}{ }\PYG{n+nv+nvIdentifier}{all}\PYG{+w}{ }\PYG{n}{property}\PYG{o}{/}\PYG{n}{local}\PYG{+w}{ }\PYG{n}{itype}\PYG{+w}{ }\PYG{n}{iatom1}\PYG{+w}{ }\PYG{n}{iatom2}\PYG{+w}{ }\PYG{n}{iatom3}\PYG{+w}{ }\PYG{n}{iatom4}
\PYG{k}{compute }\PYG{n+nv+nvIdentifier}{2}\PYG{+w}{ }\PYG{n+nv+nvIdentifier}{all}\PYG{+w}{ }\PYG{n}{improper}\PYG{o}{/}\PYG{n}{local}\PYG{+w}{ }\PYG{n}{chi}
\PYG{k}{dump }\PYG{n+nv+nvIdentifier}{1}\PYG{+w}{ }\PYG{n+nv+nvIdentifier}{all}\PYG{+w}{ }\PYG{n}{local}\PYG{+w}{ }\PYG{l+m}{1000}\PYG{+w}{ }\PYG{n}{tmp.dump}\PYG{+w}{ }\PYG{n}{index}\PYG{+w}{ }\PYG{n}{c\PYGZus{}1[1]}\PYG{+w}{ }\PYG{n}{c\PYGZus{}1[2]}\PYG{+w}{ }\PYG{n}{c\PYGZus{}1[3]}\PYG{+w}{ }\PYG{n}{c\PYGZus{}1[4]}\PYG{+w}{ }\PYG{n}{c\PYGZus{}1[5]}\PYG{+w}{ }\PYG{n}{c\PYGZus{}2[1]}
\end{sphinxVerbatim}


\subsection{Output info}
\label{\detokenize{compute_improper_local:output-info}}
\sphinxAtStartPar
This compute calculates a local vector or local array depending on the
number of keywords.  The length of the vector or number of rows in the
array is the number of impropers.  If a single keyword is specified, a
local vector is produced.  If two or more keywords are specified, a
local array is produced where the number of columns = the number of
keywords.  The vector or array can be accessed by any command that
uses local values from a compute as input.  See the
{\hyperref[\detokenize{Howto_output::doc}]{\sphinxcrossref{\DUrole{doc}{Howto output}}}} page for an overview of LAMMPS output
options.

\sphinxAtStartPar
The output for \sphinxstyleemphasis{chi} will be in degrees.


\subsection{Restrictions}
\label{\detokenize{compute_improper_local:restrictions}}\begin{quote}

\sphinxAtStartPar
none
\end{quote}


\subsection{Related commands}
\label{\detokenize{compute_improper_local:related-commands}}
\sphinxAtStartPar
{\hyperref[\detokenize{dump::doc}]{\sphinxcrossref{\DUrole{doc}{dump local}}}}, {\hyperref[\detokenize{compute_property_local::doc}]{\sphinxcrossref{\DUrole{doc}{compute property/local}}}}


\subsection{Default}
\label{\detokenize{compute_improper_local:default}}
\sphinxAtStartPar
none

\sphinxstepscope

\index{compute inertia/chunk@\spxentry{compute inertia/chunk}}\ignorespaces 

\section{compute inertia/chunk command}
\label{\detokenize{compute_inertia_chunk:compute-inertia-chunk-command}}\label{\detokenize{compute_inertia_chunk:index-0}}\label{\detokenize{compute_inertia_chunk::doc}}

\subsection{Syntax}
\label{\detokenize{compute_inertia_chunk:syntax}}
\begin{sphinxVerbatim}[commandchars=\\\{\}]
\PYG{k}{compute }\PYG{n+nv+nvIdentifier}{ID}\PYG{+w}{ }\PYG{n+nv+nvIdentifier}{group\PYGZhy{}ID}\PYG{+w}{ }\PYG{n}{inertia}\PYG{o}{/}\PYG{n}{chunk}\PYG{+w}{ }\PYG{n}{chunkID}
\end{sphinxVerbatim}
\begin{itemize}
\item {} 
\sphinxAtStartPar
ID, group\sphinxhyphen{}ID are documented in {\hyperref[\detokenize{compute::doc}]{\sphinxcrossref{\DUrole{doc}{compute}}}} command

\item {} 
\sphinxAtStartPar
inertia/chunk = style name of this compute command

\item {} 
\sphinxAtStartPar
chunkID = ID of {\hyperref[\detokenize{compute_chunk_atom::doc}]{\sphinxcrossref{\DUrole{doc}{compute chunk/atom}}}} command

\end{itemize}


\subsection{Examples}
\label{\detokenize{compute_inertia_chunk:examples}}
\begin{sphinxVerbatim}[commandchars=\\\{\}]
\PYG{k}{compute }\PYG{n+nv+nvIdentifier}{1}\PYG{+w}{ }\PYG{n+nv+nvIdentifier}{fluid}\PYG{+w}{ }\PYG{n}{inertia}\PYG{o}{/}\PYG{n}{chunk}\PYG{+w}{ }\PYG{n}{molchunk}
\end{sphinxVerbatim}


\subsection{Description}
\label{\detokenize{compute_inertia_chunk:description}}
\sphinxAtStartPar
Define a computation that calculates the inertia tensor for multiple
chunks of atoms.

\sphinxAtStartPar
In LAMMPS, chunks are collections of atoms defined by a
{\hyperref[\detokenize{compute_chunk_atom::doc}]{\sphinxcrossref{\DUrole{doc}{compute chunk/atom}}}} command, which assigns each atom
to a single chunk (or no chunk).  The ID for this command is specified
as chunkID.  For example, a single chunk could be the atoms in a
molecule or atoms in a spatial bin.  See the
{\hyperref[\detokenize{compute_chunk_atom::doc}]{\sphinxcrossref{\DUrole{doc}{compute chunk/atom}}}} and
{\hyperref[\detokenize{Howto_chunk::doc}]{\sphinxcrossref{\DUrole{doc}{Howto chunk}}}}
doc pages for details of how chunks can be defined and examples of how
they can be used to measure properties of a system.

\sphinxAtStartPar
This compute calculates the six components of the symmetric inertia
tensor for each chunk, ordered
\(I_{xx},I_{yy},I_{zz},I_{xy},I_{yz},I_{xz}\).
The calculation includes all effects due to atoms passing through periodic
boundaries.

\sphinxAtStartPar
Note that only atoms in the specified group contribute to the
calculation.  The {\hyperref[\detokenize{compute_chunk_atom::doc}]{\sphinxcrossref{\DUrole{doc}{compute chunk/atom}}}} command
defines its own group; atoms will have a chunk ID = 0 if they are not
in that group, signifying they are not assigned to a chunk, and will
thus also not contribute to this calculation.  You can specify the
“all” group for this command if you simply want to include atoms with
non\sphinxhyphen{}zero chunk IDs.

\begin{sphinxadmonition}{note}{Note:}
\sphinxAtStartPar
The coordinates of an atom contribute to the chunk’s inertia
tensor in “unwrapped” form, by using the image flags associated with
each atom.  See the {\hyperref[\detokenize{dump::doc}]{\sphinxcrossref{\DUrole{doc}{dump custom}}}} command for a discussion
of “unwrapped” coordinates.  See the Atoms section of the
{\hyperref[\detokenize{read_data::doc}]{\sphinxcrossref{\DUrole{doc}{read\_data}}}} command for a discussion of image flags and
how they are set for each atom.  You can reset the image flags
(e.g., to 0) before invoking this compute by using the
{\hyperref[\detokenize{set::doc}]{\sphinxcrossref{\DUrole{doc}{set image}}}} command.
\end{sphinxadmonition}

\sphinxAtStartPar
The simplest way to output the results of the compute inertia/chunk
calculation to a file is to use the {\hyperref[\detokenize{fix_ave_time::doc}]{\sphinxcrossref{\DUrole{doc}{fix ave/time}}}}
command, for example:

\begin{sphinxVerbatim}[commandchars=\\\{\}]
\PYG{k}{compute }\PYG{n+nv+nvIdentifier}{cc1}\PYG{+w}{ }\PYG{n+nv+nvIdentifier}{all}\PYG{+w}{ }\PYG{n}{chunk}\PYG{o}{/}\PYG{n}{atom}\PYG{+w}{ }\PYG{n}{molecule}
\PYG{k}{compute }\PYG{n+nv+nvIdentifier}{myChunk}\PYG{+w}{ }\PYG{n+nv+nvIdentifier}{all}\PYG{+w}{ }\PYG{n}{inertia}\PYG{o}{/}\PYG{n}{chunk}\PYG{+w}{ }\PYG{n}{cc1}
\PYG{k}{fix }\PYG{n+nv+nvIdentifier}{1}\PYG{+w}{ }\PYG{n+nv+nvIdentifier}{all}\PYG{+w}{ }\PYG{n}{ave}\PYG{o}{/}\PYG{n}{time}\PYG{+w}{ }\PYG{l+m}{100}\PYG{+w}{ }\PYG{l+m}{1}\PYG{+w}{ }\PYG{l+m}{100}\PYG{+w}{ }\PYG{n}{c\PYGZus{}myChunk[}\PYG{o}{*}\PYG{n}{]}\PYG{+w}{ }\PYG{n}{file}\PYG{+w}{ }\PYG{n}{tmp.out}\PYG{+w}{ }\PYG{n}{mode}\PYG{+w}{ }\PYG{n}{vector}
\end{sphinxVerbatim}


\subsection{Output info}
\label{\detokenize{compute_inertia_chunk:output-info}}
\sphinxAtStartPar
This compute calculates a global array where the number of rows = the
number of chunks \sphinxstyleemphasis{Nchunk} as calculated by the specified
{\hyperref[\detokenize{compute_chunk_atom::doc}]{\sphinxcrossref{\DUrole{doc}{compute chunk/atom}}}} command.
The number of columns is 6, one for each of the 6 components of the inertia
tensor for each chunk, ordered as listed above.  These values can be accessed
by any command that uses global array values from a compute as input.  See the
{\hyperref[\detokenize{Howto_output::doc}]{\sphinxcrossref{\DUrole{doc}{Howto output}}}} page for an overview of LAMMPS output
options.

\sphinxAtStartPar
The array values are “intensive”.  The array values will be in
mass*distance\(^2\) {\hyperref[\detokenize{units::doc}]{\sphinxcrossref{\DUrole{doc}{units}}}}.


\subsection{Restrictions}
\label{\detokenize{compute_inertia_chunk:restrictions}}\begin{quote}

\sphinxAtStartPar
none
\end{quote}


\subsection{Related commands}
\label{\detokenize{compute_inertia_chunk:related-commands}}
\sphinxAtStartPar
{\hyperref[\detokenize{variable::doc}]{\sphinxcrossref{\DUrole{doc}{variable inertia() function}}}}


\subsection{Default}
\label{\detokenize{compute_inertia_chunk:default}}
\sphinxAtStartPar
none

\sphinxstepscope

\index{compute ke@\spxentry{compute ke}}\ignorespaces 

\section{compute ke command}
\label{\detokenize{compute_ke:compute-ke-command}}\label{\detokenize{compute_ke:index-0}}\label{\detokenize{compute_ke::doc}}

\subsection{Syntax}
\label{\detokenize{compute_ke:syntax}}
\begin{sphinxVerbatim}[commandchars=\\\{\}]
\PYG{k}{compute }\PYG{n+nv+nvIdentifier}{ID}\PYG{+w}{ }\PYG{n+nv+nvIdentifier}{group\PYGZhy{}ID}\PYG{+w}{ }\PYG{n}{ke}
\end{sphinxVerbatim}
\begin{itemize}
\item {} 
\sphinxAtStartPar
ID, group\sphinxhyphen{}ID are documented in {\hyperref[\detokenize{compute::doc}]{\sphinxcrossref{\DUrole{doc}{compute}}}} command

\item {} 
\sphinxAtStartPar
ke = style name of this compute command

\end{itemize}


\subsection{Examples}
\label{\detokenize{compute_ke:examples}}
\begin{sphinxVerbatim}[commandchars=\\\{\}]
\PYG{k}{compute }\PYG{n+nv+nvIdentifier}{1}\PYG{+w}{ }\PYG{n+nv+nvIdentifier}{all}\PYG{+w}{ }\PYG{n}{ke}
\end{sphinxVerbatim}


\subsection{Description}
\label{\detokenize{compute_ke:description}}
\sphinxAtStartPar
Define a computation that calculates the translational kinetic energy
of a group of particles.

\sphinxAtStartPar
The kinetic energy of each particle is computed as \(\frac{1}{2} m
v^2\), where \sphinxstyleemphasis{m} and \sphinxstyleemphasis{v} are the mass and velocity of the particle,
respectively.

\sphinxAtStartPar
There is a subtle difference between the quantity calculated by this
compute and the kinetic energy calculated by the \sphinxstyleemphasis{ke} or \sphinxstyleemphasis{etotal}
keyword used in thermodynamic output, as specified by the
{\hyperref[\detokenize{thermo_style::doc}]{\sphinxcrossref{\DUrole{doc}{thermo\_style}}}} command.  For this compute, kinetic
energy is “translational” kinetic energy, calculated by the simple
formula above.  For thermodynamic output, the \sphinxstyleemphasis{ke} keyword infers
kinetic energy from the temperature of the system with
\(\frac{1}{2} k_B T\) of energy for each degree of freedom.  For the
default temperature computation via the {\hyperref[\detokenize{compute_temp::doc}]{\sphinxcrossref{\DUrole{doc}{compute temp}}}} command, these are the same.
However, different computes that calculate temperature can subtract out
different non\sphinxhyphen{}thermal components of velocity and/or include different degrees
of freedom (translational, rotational, etc.).


\subsection{Output info}
\label{\detokenize{compute_ke:output-info}}
\sphinxAtStartPar
This compute calculates a global scalar (the summed KE).  This value
can be used by any command that uses a global scalar value from a
compute as input.  See the {\hyperref[\detokenize{Howto_output::doc}]{\sphinxcrossref{\DUrole{doc}{Howto output}}}} doc page
for an overview of LAMMPS output options.

\sphinxAtStartPar
The scalar value calculated by this compute is “extensive”.  The
scalar value will be in energy {\hyperref[\detokenize{units::doc}]{\sphinxcrossref{\DUrole{doc}{units}}}}.


\subsection{Restrictions}
\label{\detokenize{compute_ke:restrictions}}\begin{quote}

\sphinxAtStartPar
none
\end{quote}


\subsection{Related commands}
\label{\detokenize{compute_ke:related-commands}}
\sphinxAtStartPar
{\hyperref[\detokenize{compute_erotate_sphere::doc}]{\sphinxcrossref{\DUrole{doc}{compute erotate/sphere}}}}


\subsection{Default}
\label{\detokenize{compute_ke:default}}
\sphinxAtStartPar
none

\sphinxstepscope

\index{compute ke/atom@\spxentry{compute ke/atom}}\ignorespaces 

\section{compute ke/atom command}
\label{\detokenize{compute_ke_atom:compute-ke-atom-command}}\label{\detokenize{compute_ke_atom:index-0}}\label{\detokenize{compute_ke_atom::doc}}

\subsection{Syntax}
\label{\detokenize{compute_ke_atom:syntax}}
\begin{sphinxVerbatim}[commandchars=\\\{\}]
\PYG{k}{compute }\PYG{n+nv+nvIdentifier}{ID}\PYG{+w}{ }\PYG{n+nv+nvIdentifier}{group\PYGZhy{}ID}\PYG{+w}{ }\PYG{n}{ke}\PYG{o}{/}\PYG{n}{atom}
\end{sphinxVerbatim}
\begin{itemize}
\item {} 
\sphinxAtStartPar
ID, group\sphinxhyphen{}ID are documented in {\hyperref[\detokenize{compute::doc}]{\sphinxcrossref{\DUrole{doc}{compute}}}} command

\item {} 
\sphinxAtStartPar
ke/atom = style name of this compute command

\end{itemize}


\subsection{Examples}
\label{\detokenize{compute_ke_atom:examples}}
\begin{sphinxVerbatim}[commandchars=\\\{\}]
\PYG{k}{compute }\PYG{n+nv+nvIdentifier}{1}\PYG{+w}{ }\PYG{n+nv+nvIdentifier}{all}\PYG{+w}{ }\PYG{n}{ke}\PYG{o}{/}\PYG{n}{atom}
\end{sphinxVerbatim}


\subsection{Description}
\label{\detokenize{compute_ke_atom:description}}
\sphinxAtStartPar
Define a computation that calculates the per\sphinxhyphen{}atom translational
kinetic energy for each atom in a group.

\sphinxAtStartPar
The kinetic energy is simply \(\frac12 m v^2\), where \(m\) is the mass
and \(v\) is the velocity of each atom.

\sphinxAtStartPar
The value of the kinetic energy will be 0.0 for atoms not in the
specified compute group.


\subsection{Output info}
\label{\detokenize{compute_ke_atom:output-info}}
\sphinxAtStartPar
This compute calculates a per\sphinxhyphen{}atom vector, which can be accessed by
any command that uses per\sphinxhyphen{}atom values from a compute as input.  See
the {\hyperref[\detokenize{Howto_output::doc}]{\sphinxcrossref{\DUrole{doc}{Howto output}}}} page for an overview of
LAMMPS output options.

\sphinxAtStartPar
The per\sphinxhyphen{}atom vector values will be in energy {\hyperref[\detokenize{units::doc}]{\sphinxcrossref{\DUrole{doc}{units}}}}.


\subsection{Restrictions}
\label{\detokenize{compute_ke_atom:restrictions}}\begin{quote}

\sphinxAtStartPar
none
\end{quote}


\subsection{Related commands}
\label{\detokenize{compute_ke_atom:related-commands}}
\sphinxAtStartPar
{\hyperref[\detokenize{dump::doc}]{\sphinxcrossref{\DUrole{doc}{dump custom}}}}


\subsection{Default}
\label{\detokenize{compute_ke_atom:default}}
\sphinxAtStartPar
none

\sphinxstepscope

\index{compute ke/atom/eff@\spxentry{compute ke/atom/eff}}\ignorespaces 

\section{compute ke/atom/eff command}
\label{\detokenize{compute_ke_atom_eff:compute-ke-atom-eff-command}}\label{\detokenize{compute_ke_atom_eff:index-0}}\label{\detokenize{compute_ke_atom_eff::doc}}

\subsection{Syntax}
\label{\detokenize{compute_ke_atom_eff:syntax}}
\begin{sphinxVerbatim}[commandchars=\\\{\}]
\PYG{k}{compute }\PYG{n+nv+nvIdentifier}{ID}\PYG{+w}{ }\PYG{n+nv+nvIdentifier}{group\PYGZhy{}ID}\PYG{+w}{ }\PYG{n}{ke}\PYG{o}{/}\PYG{n}{atom}\PYG{o}{/}\PYG{n}{eff}
\end{sphinxVerbatim}
\begin{itemize}
\item {} 
\sphinxAtStartPar
ID, group\sphinxhyphen{}ID are documented in {\hyperref[\detokenize{compute::doc}]{\sphinxcrossref{\DUrole{doc}{compute}}}} command

\item {} 
\sphinxAtStartPar
ke/atom/eff = style name of this compute command

\end{itemize}


\subsection{Examples}
\label{\detokenize{compute_ke_atom_eff:examples}}
\begin{sphinxVerbatim}[commandchars=\\\{\}]
\PYG{k}{compute }\PYG{n+nv+nvIdentifier}{1}\PYG{+w}{ }\PYG{n+nv+nvIdentifier}{all}\PYG{+w}{ }\PYG{n}{ke}\PYG{o}{/}\PYG{n}{atom}\PYG{o}{/}\PYG{n}{eff}
\end{sphinxVerbatim}


\subsection{Description}
\label{\detokenize{compute_ke_atom_eff:description}}
\sphinxAtStartPar
Define a computation that calculates the per\sphinxhyphen{}atom translational
(nuclei and electrons) and radial kinetic energy (electron only) in a
group.  The particles are assumed to be nuclei and electrons modeled
with the {\hyperref[\detokenize{pair_eff::doc}]{\sphinxcrossref{\DUrole{doc}{electronic force field}}}}.

\sphinxAtStartPar
The kinetic energy for each nucleus is computed as \(\frac{1}{2} m
v^2\), where \sphinxstyleemphasis{m} corresponds to the corresponding nuclear mass, and the
kinetic energy for each electron is computed as \(\frac{1}{2} (m_e
v^2 + \frac{3}{4} m_e s^2)\), where \(m_e\) and \sphinxstyleemphasis{v} correspond to the mass
and translational velocity of each electron, and \sphinxstyleemphasis{s} to its radial
velocity, respectively.

\sphinxAtStartPar
There is a subtle difference between the quantity calculated by this
compute and the kinetic energy calculated by the \sphinxstyleemphasis{ke} or \sphinxstyleemphasis{etotal}
keyword used in thermodynamic output, as specified by the
{\hyperref[\detokenize{thermo_style::doc}]{\sphinxcrossref{\DUrole{doc}{thermo\_style}}}} command. For this compute, kinetic
energy is “translational” plus electronic “radial” kinetic energy,
calculated by the simple formula above. For thermodynamic output, the
\sphinxstyleemphasis{ke} keyword infers kinetic energy from the temperature of the system
with \(\frac{1}{2} k_B T\) of energy for each (nuclear\sphinxhyphen{}only) degree
of freedom in eFF.

\begin{sphinxadmonition}{note}{Note:}
\sphinxAtStartPar
The temperature in eFF should be monitored via the {\hyperref[\detokenize{compute_temp_eff::doc}]{\sphinxcrossref{\DUrole{doc}{compute temp/eff}}}} command, which can be printed with
thermodynamic output by using the {\hyperref[\detokenize{thermo_modify::doc}]{\sphinxcrossref{\DUrole{doc}{thermo\_modify}}}}
command, as shown in the following example:
\end{sphinxadmonition}

\begin{sphinxVerbatim}[commandchars=\\\{\}]
\PYG{k}{compute         }\PYG{n+nv+nvIdentifier}{effTemp}\PYG{+w}{ }\PYG{n+nv+nvIdentifier}{all}\PYG{+w}{ }\PYG{n}{temp}\PYG{o}{/}\PYG{n}{eff}
\PYG{k}{thermo\PYGZus{}style}\PYG{+w}{    }\PYG{n}{custom}\PYG{+w}{ }\PYG{n}{step}\PYG{+w}{ }\PYG{n}{etotal}\PYG{+w}{ }\PYG{n}{pe}\PYG{+w}{ }\PYG{n}{ke}\PYG{+w}{ }\PYG{n}{temp}\PYG{+w}{ }\PYG{n}{press}
\PYG{k}{thermo\PYGZus{}modify}\PYG{+w}{   }\PYG{n}{temp}\PYG{+w}{ }\PYG{n}{effTemp}
\end{sphinxVerbatim}

\sphinxAtStartPar
The value of the kinetic energy will be 0.0 for atoms (nuclei or
electrons) not in the specified compute group.


\subsection{Output info}
\label{\detokenize{compute_ke_atom_eff:output-info}}
\sphinxAtStartPar
This compute calculates a scalar quantity for each atom, which can be
accessed by any command that uses per\sphinxhyphen{}atom computes as input.  See the
{\hyperref[\detokenize{Howto_output::doc}]{\sphinxcrossref{\DUrole{doc}{Howto output}}}} page for an overview of LAMMPS
output options.

\sphinxAtStartPar
The per\sphinxhyphen{}atom vector values will be in energy {\hyperref[\detokenize{units::doc}]{\sphinxcrossref{\DUrole{doc}{units}}}}.


\subsection{Restrictions}
\label{\detokenize{compute_ke_atom_eff:restrictions}}
\sphinxAtStartPar
This compute is part of the EFF package.  It is only enabled if
LAMMPS was built with that package.  See the {\hyperref[\detokenize{Build_package::doc}]{\sphinxcrossref{\DUrole{doc}{Build package}}}} page for more info.


\subsection{Related commands}
\label{\detokenize{compute_ke_atom_eff:related-commands}}
\sphinxAtStartPar
{\hyperref[\detokenize{dump::doc}]{\sphinxcrossref{\DUrole{doc}{dump custom}}}}


\subsection{Default}
\label{\detokenize{compute_ke_atom_eff:default}}
\sphinxAtStartPar
none

\sphinxstepscope

\index{compute ke/eff@\spxentry{compute ke/eff}}\ignorespaces 

\section{compute ke/eff command}
\label{\detokenize{compute_ke_eff:compute-ke-eff-command}}\label{\detokenize{compute_ke_eff:index-0}}\label{\detokenize{compute_ke_eff::doc}}

\subsection{Syntax}
\label{\detokenize{compute_ke_eff:syntax}}
\begin{sphinxVerbatim}[commandchars=\\\{\}]
\PYG{k}{compute }\PYG{n+nv+nvIdentifier}{ID}\PYG{+w}{ }\PYG{n+nv+nvIdentifier}{group\PYGZhy{}ID}\PYG{+w}{ }\PYG{n}{ke}\PYG{o}{/}\PYG{n}{eff}
\end{sphinxVerbatim}
\begin{itemize}
\item {} 
\sphinxAtStartPar
ID, group\sphinxhyphen{}ID are documented in {\hyperref[\detokenize{compute::doc}]{\sphinxcrossref{\DUrole{doc}{compute}}}} command

\item {} 
\sphinxAtStartPar
ke/eff = style name of this compute command

\end{itemize}


\subsection{Examples}
\label{\detokenize{compute_ke_eff:examples}}
\begin{sphinxVerbatim}[commandchars=\\\{\}]
\PYG{k}{compute }\PYG{n+nv+nvIdentifier}{1}\PYG{+w}{ }\PYG{n+nv+nvIdentifier}{all}\PYG{+w}{ }\PYG{n}{ke}\PYG{o}{/}\PYG{n}{eff}
\end{sphinxVerbatim}


\subsection{Description}
\label{\detokenize{compute_ke_eff:description}}
\sphinxAtStartPar
Define a computation that calculates the kinetic energy of motion of a
group of eFF particles (nuclei and electrons), as modeled with the
{\hyperref[\detokenize{pair_eff::doc}]{\sphinxcrossref{\DUrole{doc}{electronic force field}}}}.

\sphinxAtStartPar
The kinetic energy for each nucleus is computed as \(\frac{1}{2} m
v^2\) and the kinetic energy for each electron is computed as
\(\frac{1}{2}(m_e v^2 + \frac{3}{4} m_e s^2)\), where \(m\)
corresponds to the nuclear mass, \(m_e\) to the electron mass, \(v\)
to the translational velocity of each particle, and \(s\) to the radial
velocity of the electron, respectively.

\sphinxAtStartPar
There is a subtle difference between the quantity calculated by this
compute and the kinetic energy calculated by the \sphinxstyleemphasis{ke} or \sphinxstyleemphasis{etotal}
keyword used in thermodynamic output, as specified by the
{\hyperref[\detokenize{thermo_style::doc}]{\sphinxcrossref{\DUrole{doc}{thermo\_style}}}} command.  For this compute, kinetic
energy is “translational” and “radial” (only for electrons) kinetic
energy, calculated by the simple formula above.  For thermodynamic
output, the \sphinxstyleemphasis{ke} keyword infers kinetic energy from the temperature of
the system with \(\frac{1}{2} k_B T\) of energy for each degree of
freedom.  For the eFF temperature computation via the {\hyperref[\detokenize{compute_temp_eff::doc}]{\sphinxcrossref{\DUrole{doc}{compute
temp\_eff}}}} command, these are the same.  But
different computes that calculate temperature can subtract out different
non\sphinxhyphen{}thermal components of velocity and/or include other degrees of
freedom.

\begin{sphinxadmonition}{warning}{Warning:}
\sphinxAtStartPar
The temperature in eFF models should be monitored via
the {\hyperref[\detokenize{compute_temp_eff::doc}]{\sphinxcrossref{\DUrole{doc}{compute temp/eff}}}} command, which can be
printed with thermodynamic output by using the
{\hyperref[\detokenize{thermo_modify::doc}]{\sphinxcrossref{\DUrole{doc}{thermo\_modify}}}} command, as shown in the following
example:
\end{sphinxadmonition}

\begin{sphinxVerbatim}[commandchars=\\\{\}]
\PYG{k}{compute         }\PYG{n+nv+nvIdentifier}{effTemp}\PYG{+w}{ }\PYG{n+nv+nvIdentifier}{all}\PYG{+w}{ }\PYG{n}{temp}\PYG{o}{/}\PYG{n}{eff}
\PYG{k}{thermo\PYGZus{}style}\PYG{+w}{    }\PYG{n}{custom}\PYG{+w}{ }\PYG{n}{step}\PYG{+w}{ }\PYG{n}{etotal}\PYG{+w}{ }\PYG{n}{pe}\PYG{+w}{ }\PYG{n}{ke}\PYG{+w}{ }\PYG{n}{temp}\PYG{+w}{ }\PYG{n}{press}
\PYG{k}{thermo\PYGZus{}modify}\PYG{+w}{   }\PYG{n}{temp}\PYG{+w}{ }\PYG{n}{effTemp}
\end{sphinxVerbatim}

\sphinxAtStartPar
See {\hyperref[\detokenize{compute_temp_eff::doc}]{\sphinxcrossref{\DUrole{doc}{compute temp/eff}}}}.


\subsection{Output info}
\label{\detokenize{compute_ke_eff:output-info}}
\sphinxAtStartPar
This compute calculates a global scalar (the KE).  This value can be
used by any command that uses a global scalar value from a compute as
input.  See the {\hyperref[\detokenize{Howto_output::doc}]{\sphinxcrossref{\DUrole{doc}{Howto output}}}} page for an
overview of LAMMPS output options.

\sphinxAtStartPar
The scalar value calculated by this compute is “extensive”.  The
scalar value will be in energy {\hyperref[\detokenize{units::doc}]{\sphinxcrossref{\DUrole{doc}{units}}}}.


\subsection{Restrictions}
\label{\detokenize{compute_ke_eff:restrictions}}
\sphinxAtStartPar
This compute is part of the EFF package.  It is only enabled if
LAMMPS was built with that package.  See the {\hyperref[\detokenize{Build_package::doc}]{\sphinxcrossref{\DUrole{doc}{Build package}}}} page for more info.


\subsection{Related commands}
\label{\detokenize{compute_ke_eff:related-commands}}
\sphinxAtStartPar
none


\subsection{Default}
\label{\detokenize{compute_ke_eff:default}}
\sphinxAtStartPar
none

\sphinxstepscope

\index{compute ke/rigid@\spxentry{compute ke/rigid}}\ignorespaces 

\section{compute ke/rigid command}
\label{\detokenize{compute_ke_rigid:compute-ke-rigid-command}}\label{\detokenize{compute_ke_rigid:index-0}}\label{\detokenize{compute_ke_rigid::doc}}

\subsection{Syntax}
\label{\detokenize{compute_ke_rigid:syntax}}
\begin{sphinxVerbatim}[commandchars=\\\{\}]
\PYG{k}{compute }\PYG{n+nv+nvIdentifier}{ID}\PYG{+w}{ }\PYG{n+nv+nvIdentifier}{group\PYGZhy{}ID}\PYG{+w}{ }\PYG{n}{ke}\PYG{o}{/}\PYG{n}{rigid}\PYG{+w}{ }\PYG{n}{fix}\PYG{o}{\PYGZhy{}}\PYG{n}{ID}
\end{sphinxVerbatim}
\begin{itemize}
\item {} 
\sphinxAtStartPar
ID, group\sphinxhyphen{}ID are documented in {\hyperref[\detokenize{compute::doc}]{\sphinxcrossref{\DUrole{doc}{compute}}}} command

\item {} 
\sphinxAtStartPar
ke = style name of this compute command

\item {} 
\sphinxAtStartPar
fix\sphinxhyphen{}ID = ID of rigid body fix

\end{itemize}


\subsection{Examples}
\label{\detokenize{compute_ke_rigid:examples}}
\begin{sphinxVerbatim}[commandchars=\\\{\}]
\PYG{k}{compute }\PYG{n+nv+nvIdentifier}{1}\PYG{+w}{ }\PYG{n+nv+nvIdentifier}{all}\PYG{+w}{ }\PYG{n}{ke}\PYG{o}{/}\PYG{n}{rigid}\PYG{+w}{ }\PYG{n}{myRigid}
\end{sphinxVerbatim}


\subsection{Description}
\label{\detokenize{compute_ke_rigid:description}}
\sphinxAtStartPar
Define a computation that calculates the translational kinetic energy
of a collection of rigid bodies, as defined by one of the
{\hyperref[\detokenize{fix_rigid::doc}]{\sphinxcrossref{\DUrole{doc}{fix rigid}}}} command variants.

\sphinxAtStartPar
The kinetic energy of each rigid body is computed as
\(\frac12 M V_\text{cm}^2\),
where \(M\) is the total mass of the rigid body, and \(V_\text{cm}\)
is its center\sphinxhyphen{}of\sphinxhyphen{}mass velocity.

\sphinxAtStartPar
The \sphinxstyleemphasis{fix\sphinxhyphen{}ID} should be the ID of one of the {\hyperref[\detokenize{fix_rigid::doc}]{\sphinxcrossref{\DUrole{doc}{fix rigid}}}}
commands which defines the rigid bodies.  The group specified in the
compute command is ignored.  The kinetic energy of all the rigid
bodies defined by the fix rigid command in included in the
calculation.


\subsection{Output info}
\label{\detokenize{compute_ke_rigid:output-info}}
\sphinxAtStartPar
This compute calculates a global scalar (the summed KE of all the
rigid bodies).  This value can be used by any command that uses a
global scalar value from a compute as input.  See the
{\hyperref[\detokenize{Howto_output::doc}]{\sphinxcrossref{\DUrole{doc}{Howto output}}}} page for an overview of LAMMPS output
options.

\sphinxAtStartPar
The scalar value calculated by this compute is “extensive”.  The
scalar value will be in energy {\hyperref[\detokenize{units::doc}]{\sphinxcrossref{\DUrole{doc}{units}}}}.


\subsection{Restrictions}
\label{\detokenize{compute_ke_rigid:restrictions}}
\sphinxAtStartPar
This compute is part of the RIGID package.  It is only enabled if
LAMMPS was built with that package.  See the {\hyperref[\detokenize{Build_package::doc}]{\sphinxcrossref{\DUrole{doc}{Build package}}}} page for more info.


\subsection{Related commands}
\label{\detokenize{compute_ke_rigid:related-commands}}
\sphinxAtStartPar
{\hyperref[\detokenize{compute_erotate_rigid::doc}]{\sphinxcrossref{\DUrole{doc}{compute erotate/rigid}}}}


\subsection{Default}
\label{\detokenize{compute_ke_rigid:default}}
\sphinxAtStartPar
none

\sphinxstepscope

\index{compute mliap@\spxentry{compute mliap}}\ignorespaces 

\section{compute mliap command}
\label{\detokenize{compute_mliap:compute-mliap-command}}\label{\detokenize{compute_mliap:index-0}}\label{\detokenize{compute_mliap::doc}}

\subsection{Syntax}
\label{\detokenize{compute_mliap:syntax}}
\begin{sphinxVerbatim}[commandchars=\\\{\}]
\PYG{k}{compute }\PYG{n+nv+nvIdentifier}{ID}\PYG{+w}{ }\PYG{n+nv+nvIdentifier}{group\PYGZhy{}ID}\PYG{+w}{ }\PYG{n}{mliap}\PYG{+w}{ }\PYG{n}{...}\PYG{+w}{ }\PYG{n}{keyword}\PYG{+w}{ }\PYG{n}{values}\PYG{+w}{ }\PYG{n}{...}
\end{sphinxVerbatim}
\begin{itemize}
\item {} 
\sphinxAtStartPar
ID, group\sphinxhyphen{}ID are documented in {\hyperref[\detokenize{compute::doc}]{\sphinxcrossref{\DUrole{doc}{compute}}}} command

\item {} 
\sphinxAtStartPar
mliap = style name of this compute command

\item {} 
\sphinxAtStartPar
two or more keyword/value pairs must be appended

\item {} 
\sphinxAtStartPar
keyword = \sphinxstyleemphasis{model} or \sphinxstyleemphasis{descriptor} or \sphinxstyleemphasis{gradgradflag}
\begin{sphinxalltt}
\sphinxstyleemphasis{model} values = style
  style = \sphinxstyleemphasis{linear} or \sphinxstyleemphasis{quadratic} or \sphinxstyleemphasis{mliappy}
\sphinxstyleemphasis{descriptor} values = style filename
  style = \sphinxstyleemphasis{sna} or \sphinxstyleemphasis{ace}
  filename = name of file containing descriptor definitions
\sphinxstyleemphasis{gradgradflag} value = 0/1
  toggle gradgrad method for force gradient
\end{sphinxalltt}

\end{itemize}


\subsection{Examples}
\label{\detokenize{compute_mliap:examples}}
\begin{sphinxVerbatim}[commandchars=\\\{\}]
\PYG{k}{compute }\PYG{n+nv+nvIdentifier}{mliap}\PYG{+w}{ }\PYG{n+nv+nvIdentifier}{model}\PYG{+w}{ }\PYG{n}{linear}\PYG{+w}{ }\PYG{n}{descriptor}\PYG{+w}{ }\PYG{n}{sna}\PYG{+w}{ }\PYG{n}{Ta06A.mliap.descriptor}
\PYG{k}{compute }\PYG{n+nv+nvIdentifier}{mliap}\PYG{+w}{ }\PYG{n+nv+nvIdentifier}{model}\PYG{+w}{ }\PYG{n}{linear}\PYG{+w}{ }\PYG{n}{descriptor}\PYG{+w}{ }\PYG{n}{ace}\PYG{+w}{ }\PYG{n}{H\PYGZus{}N\PYGZus{}O\PYGZus{}ccs.yace}\PYG{+w}{ }\PYG{n}{gradgradflag}\PYG{+w}{ }\PYG{l+m}{1}
\end{sphinxVerbatim}


\subsection{Description}
\label{\detokenize{compute_mliap:description}}
\sphinxAtStartPar
Compute style \sphinxstyleemphasis{mliap} provides a general interface to the gradient
of machine\sphinxhyphen{}learning interatomic potentials with respect to model parameters.
It is used primarily for calculating the gradient of energy, force, and
stress components with respect to model parameters, which is useful when
training {\hyperref[\detokenize{pair_mliap::doc}]{\sphinxcrossref{\DUrole{doc}{mliap pair\_style}}}} models to match target data.
It provides separate definitions of the interatomic potential functional
form (\sphinxstyleemphasis{model}) and the geometric quantities that characterize the atomic
positions (\sphinxstyleemphasis{descriptor}). By defining \sphinxstyleemphasis{model} and \sphinxstyleemphasis{descriptor} separately,
it is possible to use many different models with a given descriptor,
or many different descriptors with a given model. Currently, the compute
supports \sphinxstyleemphasis{linear} and \sphinxstyleemphasis{quadratic} SNAP descriptor computes used in
{\hyperref[\detokenize{pair_snap::doc}]{\sphinxcrossref{\DUrole{doc}{pair\_style snap}}}}, \sphinxstyleemphasis{linear} SO3 descriptor computes, and
\sphinxstyleemphasis{linear} ACE descriptor computes used in {\hyperref[\detokenize{pair_pace::doc}]{\sphinxcrossref{\DUrole{doc}{pair\_style pace}}}},
and it is straightforward to add new descriptor styles.

\sphinxAtStartPar
The compute \sphinxstyleemphasis{mliap} command must be followed by two keywords
\sphinxstyleemphasis{model} and \sphinxstyleemphasis{descriptor} in either order.

\sphinxAtStartPar
The \sphinxstyleemphasis{model} keyword is followed by the model style (\sphinxstyleemphasis{linear},
\sphinxstyleemphasis{quadratic} or \sphinxstyleemphasis{mliappy}).  The \sphinxstyleemphasis{mliappy} model is only available if
LAMMPS is built with the \sphinxstyleemphasis{mliappy} Python module. There are
{\hyperref[\detokenize{Build_extras:mliap}]{\sphinxcrossref{\DUrole{std,std-ref}{specific installation instructions}}}} for that module.
For the \sphinxstyleemphasis{mliap} compute, specifying a \sphinxstyleemphasis{linear} model will compute the
specified descriptors and gradients with respect to linear model parameters
whereas \sphinxstyleemphasis{quadratic} will do the same, but for the quadratic products of
descriptors.

\sphinxAtStartPar
The \sphinxstyleemphasis{descriptor} keyword is followed by a descriptor style, and
additional arguments.  The compute currently supports three descriptor
styles: \sphinxstyleemphasis{sna}, \sphinxstyleemphasis{so3}, and \sphinxstyleemphasis{ace}, but it is is straightforward to add
additional descriptor styles.  The SNAP descriptor style \sphinxstyleemphasis{sna} is the
same as that used by {\hyperref[\detokenize{pair_snap::doc}]{\sphinxcrossref{\DUrole{doc}{pair\_style snap}}}}, including the
linear, quadratic, and chem variants.  A single additional argument
specifies the descriptor filename containing the parameters and setting used
by the SNAP descriptor.  The descriptor filename usually ends in the
\sphinxstyleemphasis{.mliap.descriptor} extension.  The format of this file is identical to
the descriptor file in the {\hyperref[\detokenize{pair_mliap::doc}]{\sphinxcrossref{\DUrole{doc}{pair\_style mliap}}}}, and is
described in detail there.

\sphinxAtStartPar
The ACE descriptor style \sphinxstyleemphasis{ace} is the same as {\hyperref[\detokenize{pair_pace::doc}]{\sphinxcrossref{\DUrole{doc}{pair\_style pace}}}}.
A single additional argument specifies the \sphinxstyleemphasis{ace} descriptor filename
that contains parameters and settings for the ACE descriptors. This file
format differs from the SNAP or SO3 descriptor files, and has a \sphinxstyleemphasis{.yace} or
\sphinxstyleemphasis{.ace} extension. However, as with other mliap descriptor styles, this file
is identical to the ace descriptor file in {\hyperref[\detokenize{pair_mliap::doc}]{\sphinxcrossref{\DUrole{doc}{pair\_style mliap}}}},
where it is described in further detail.

\begin{sphinxadmonition}{note}{Note:}
\sphinxAtStartPar
The number of LAMMPS atom types (and the value of \sphinxstyleemphasis{nelems} in the model)
must match the value of \sphinxstyleemphasis{nelems} in the descriptor file.
\end{sphinxadmonition}

\sphinxAtStartPar
Compute \sphinxstyleemphasis{mliap} calculates a global array containing gradient information.
The number of columns in the array is \sphinxstyleemphasis{nelems} \(\times\) \sphinxstyleemphasis{nparams} + 1.
The first row of the array contain the derivative of potential energy with
respect to. to each parameter and each element. The last six rows
of the array contain the corresponding derivatives of the
virial stress tensor, listed in Voigt notation: \sphinxstyleemphasis{pxx}, \sphinxstyleemphasis{pyy}, \sphinxstyleemphasis{pzz},
\sphinxstyleemphasis{pyz}, \sphinxstyleemphasis{pxz}, and \sphinxstyleemphasis{pxy}. In between the energy and stress rows are
the \(3N\) rows containing the derivatives of the force components.
See section below on output for a detailed description of how
rows and columns are ordered.

\sphinxAtStartPar
The element in the last column of each row contains
the potential energy, force, or stress, according to the row.
These quantities correspond to the user\sphinxhyphen{}specified reference potential
that must be subtracted from the target data when training a model.
The potential energy calculation uses the built in compute \sphinxstyleemphasis{thermo\_pe}.
The stress calculation uses a compute called \sphinxstyleemphasis{mliap\_press} that is
automatically created behind the scenes, according to the following
command:

\begin{sphinxVerbatim}[commandchars=\\\{\}]
\PYG{k}{compute }\PYG{n+nv+nvIdentifier}{mliap\PYGZus{}press}\PYG{+w}{ }\PYG{n+nv+nvIdentifier}{all}\PYG{+w}{ }\PYG{n}{pressure}\PYG{+w}{ }\PYG{n}{NULL}\PYG{+w}{ }\PYG{n}{virial}
\end{sphinxVerbatim}

\sphinxAtStartPar
See section below on output for a detailed explanation of the data
layout in the global array.

\sphinxAtStartPar
The optional keyword \sphinxstyleemphasis{gradgradflag} controls how the force
gradient is calculated. A value of 1 requires that the model provide
the matrix of double gradients of energy with respect to both parameters
and descriptors. For the linear and quadratic models this matrix is
sparse and so is easily calculated and stored. For other models, this
matrix may be prohibitively expensive to calculate and store.
A value of 0 requires that the descriptor provide the derivative
of the descriptors with respect to the position of every neighbor atom.
This is not optimal for linear and quadratic models, but may be
a better choice for more complex models.

\sphinxAtStartPar
Atoms not in the group do not contribute to this compute.
Neighbor atoms not in the group do not contribute to this compute.
The neighbor list needed to compute this quantity is constructed each
time the calculation is performed (i.e., each time a snapshot of atoms
is dumped).  Thus it can be inefficient to compute/dump this quantity
too frequently.

\begin{sphinxadmonition}{note}{Note:}
\sphinxAtStartPar
If the user\sphinxhyphen{}specified reference potentials includes bonded and
non\sphinxhyphen{}bonded pairwise interactions, then the settings of
{\hyperref[\detokenize{special_bonds::doc}]{\sphinxcrossref{\DUrole{doc}{special\_bonds}}}} command can remove pairwise
interactions between atoms in the same bond, angle, or dihedral.  This
is the default setting for the {\hyperref[\detokenize{special_bonds::doc}]{\sphinxcrossref{\DUrole{doc}{special\_bonds}}}}
command, and means those pairwise interactions do not appear in the
neighbor list.  Because this fix uses the neighbor list, it also means
those pairs will not be included in the calculation. The {\hyperref[\detokenize{rerun::doc}]{\sphinxcrossref{\DUrole{doc}{rerun}}}}
command is not an option here, since the reference potential is required
for the last column of the global array. A work\sphinxhyphen{}around is to prevent
pairwise interactions from being removed by explicitly adding a
\sphinxstyleemphasis{tiny} positive value for every pairwise interaction that would otherwise be
set to zero in the {\hyperref[\detokenize{special_bonds::doc}]{\sphinxcrossref{\DUrole{doc}{special\_bonds}}}} command.
\end{sphinxadmonition}


\bigskip\hrule\bigskip



\subsection{Output info}
\label{\detokenize{compute_mliap:output-info}}
\sphinxAtStartPar
Compute \sphinxstyleemphasis{mliap} evaluates a global array.  The columns are arranged into
\sphinxstyleemphasis{nelems} blocks, listed in order of element \sphinxstyleemphasis{I}. Each block
contains one column for each of the \sphinxstyleemphasis{nparams} model parameters.
A final column contains the corresponding energy, force component
on an atom, or virial stress component. The rows of the array appear
in the following order:
\begin{itemize}
\item {} 
\sphinxAtStartPar
1 row: Derivatives of potential energy with respect to each parameter of each element.

\item {} 
\sphinxAtStartPar
\(3N\) rows: Derivatives of force components; the \sphinxstyleemphasis{x}, \sphinxstyleemphasis{y}, and \sphinxstyleemphasis{z}
components of the force on atom \sphinxstyleemphasis{i} appear in consecutive rows. The atoms are
sorted based on atom ID.

\item {} 
\sphinxAtStartPar
6 rows: Derivatives of the virial stress tensor with respect to each
parameter of each element. The ordering of the rows follows Voigt notation:
\sphinxstyleemphasis{pxx}, \sphinxstyleemphasis{pyy}, \sphinxstyleemphasis{pzz}, \sphinxstyleemphasis{pyz}, \sphinxstyleemphasis{pxz}, \sphinxstyleemphasis{pxy}.

\end{itemize}

\sphinxAtStartPar
These values can be accessed by any command that uses a global array
from a compute as input.  See the {\hyperref[\detokenize{Howto_output::doc}]{\sphinxcrossref{\DUrole{doc}{Howto output}}}} doc
page for an overview of LAMMPS output options. To see how this command
can be used within a Python workflow to train machine\sphinxhyphen{}learning interatomic
potentials, see the examples in \sphinxhref{https://github.com/FitSNAP/FitSNAP}{FitSNAP}.


\subsection{Restrictions}
\label{\detokenize{compute_mliap:restrictions}}
\sphinxAtStartPar
This compute is part of the ML\sphinxhyphen{}IAP package.  It is only enabled if
LAMMPS was built with that package. In addition, building LAMMPS with
the ML\sphinxhyphen{}IAP package requires building LAMMPS with the ML\sphinxhyphen{}SNAP package.
The \sphinxstyleemphasis{mliappy} model also requires building LAMMPS with the PYTHON
package. The \sphinxstyleemphasis{ace} descriptor also requires building LAMMPS with the
ML\sphinxhyphen{}PACE package. See the {\hyperref[\detokenize{Build_package::doc}]{\sphinxcrossref{\DUrole{doc}{Build package}}}} page for
more info. Note that \sphinxtitleref{kk} (KOKKOS) accelerated variants of SNAP and
ACE descriptors are not compatible with \sphinxtitleref{mliap descriptor}.


\subsection{Related commands}
\label{\detokenize{compute_mliap:related-commands}}
\sphinxAtStartPar
{\hyperref[\detokenize{pair_mliap::doc}]{\sphinxcrossref{\DUrole{doc}{pair\_style mliap}}}}


\subsection{Default}
\label{\detokenize{compute_mliap:default}}
\sphinxAtStartPar
The keyword defaults are gradgradflag = 1

\sphinxstepscope

\index{compute momentum@\spxentry{compute momentum}}\ignorespaces 

\section{compute momentum command}
\label{\detokenize{compute_momentum:compute-momentum-command}}\label{\detokenize{compute_momentum:index-0}}\label{\detokenize{compute_momentum::doc}}

\subsection{Syntax}
\label{\detokenize{compute_momentum:syntax}}
\begin{sphinxVerbatim}[commandchars=\\\{\}]
\PYG{k}{compute }\PYG{n+nv+nvIdentifier}{ID}\PYG{+w}{ }\PYG{n+nv+nvIdentifier}{group\PYGZhy{}ID}\PYG{+w}{ }\PYG{n}{momentum}
\end{sphinxVerbatim}
\begin{itemize}
\item {} 
\sphinxAtStartPar
ID, group\sphinxhyphen{}ID are documented in {\hyperref[\detokenize{compute::doc}]{\sphinxcrossref{\DUrole{doc}{compute}}}} command

\item {} 
\sphinxAtStartPar
momentum = style name of this compute command

\end{itemize}


\subsection{Examples}
\label{\detokenize{compute_momentum:examples}}
\begin{sphinxVerbatim}[commandchars=\\\{\}]
\PYG{k}{compute }\PYG{n+nv+nvIdentifier}{1}\PYG{+w}{ }\PYG{n+nv+nvIdentifier}{all}\PYG{+w}{ }\PYG{n}{momentum}
\end{sphinxVerbatim}


\subsection{Description}
\label{\detokenize{compute_momentum:description}}
\sphinxAtStartPar
Define a computation that calculates the translational momentum \sphinxstyleemphasis{p}
of a group of particles.  It is computed as the sum
\(\vec{p} = \sum_i m_i \cdot \vec{v}_i\)
over all particles in the compute group, where \sphinxstyleemphasis{m} and \sphinxstyleemphasis{v} are
the mass and velocity vector of the particle, respectively.


\subsection{Output info}
\label{\detokenize{compute_momentum:output-info}}
\sphinxAtStartPar
This compute calculates a global vector (the summed momentum) of
length 3. This value can be used by any command that uses a global
vector value from a compute as input. See the {\hyperref[\detokenize{Howto_output::doc}]{\sphinxcrossref{\DUrole{doc}{Howto output}}}} page for an overview of LAMMPS output
options.

\sphinxAtStartPar
The vector value calculated by this compute is “extensive”. The vector
value will be in mass*velocity {\hyperref[\detokenize{units::doc}]{\sphinxcrossref{\DUrole{doc}{units}}}}.


\subsection{Restrictions}
\label{\detokenize{compute_momentum:restrictions}}
\sphinxAtStartPar
This compute is part of the EXTRA\sphinxhyphen{}COMPUTE package.  It is only enabled if
LAMMPS was built with that package.  See the {\hyperref[\detokenize{Build_package::doc}]{\sphinxcrossref{\DUrole{doc}{Build package}}}} page for more info.


\subsection{Related commands}
\label{\detokenize{compute_momentum:related-commands}}

\subsection{Default}
\label{\detokenize{compute_momentum:default}}
\sphinxAtStartPar
none

\sphinxstepscope

\index{compute msd@\spxentry{compute msd}}\ignorespaces 

\section{compute msd command}
\label{\detokenize{compute_msd:compute-msd-command}}\label{\detokenize{compute_msd:index-0}}\label{\detokenize{compute_msd::doc}}

\subsection{Syntax}
\label{\detokenize{compute_msd:syntax}}
\begin{sphinxVerbatim}[commandchars=\\\{\}]
\PYG{k}{compute }\PYG{n+nv+nvIdentifier}{ID}\PYG{+w}{ }\PYG{n+nv+nvIdentifier}{group\PYGZhy{}ID}\PYG{+w}{ }\PYG{n}{msd}\PYG{+w}{ }\PYG{n}{keyword}\PYG{+w}{ }\PYG{n}{values}\PYG{+w}{ }\PYG{n}{...}
\end{sphinxVerbatim}
\begin{itemize}
\item {} 
\sphinxAtStartPar
ID, group\sphinxhyphen{}ID are documented in {\hyperref[\detokenize{compute::doc}]{\sphinxcrossref{\DUrole{doc}{compute}}}} command

\item {} 
\sphinxAtStartPar
msd = style name of this compute command

\item {} 
\sphinxAtStartPar
zero or more keyword/value pairs may be appended

\item {} 
\sphinxAtStartPar
keyword = \sphinxstyleemphasis{com} or \sphinxstyleemphasis{average}
\begin{sphinxalltt}
\sphinxstyleemphasis{com} value = \sphinxstyleemphasis{yes} or \sphinxstyleemphasis{no}
\sphinxstyleemphasis{average} value = \sphinxstyleemphasis{yes} or \sphinxstyleemphasis{no}
\end{sphinxalltt}

\end{itemize}


\subsection{Examples}
\label{\detokenize{compute_msd:examples}}
\begin{sphinxVerbatim}[commandchars=\\\{\}]
\PYG{k}{compute }\PYG{n+nv+nvIdentifier}{1}\PYG{+w}{ }\PYG{n+nv+nvIdentifier}{all}\PYG{+w}{ }\PYG{n}{msd}
\PYG{k}{compute }\PYG{n+nv+nvIdentifier}{1}\PYG{+w}{ }\PYG{n+nv+nvIdentifier}{upper}\PYG{+w}{ }\PYG{n}{msd}\PYG{+w}{ }\PYG{n}{com}\PYG{+w}{ }\PYG{n}{yes}\PYG{+w}{ }\PYG{n}{average}\PYG{+w}{ }\PYG{n}{yes}
\end{sphinxVerbatim}


\subsection{Description}
\label{\detokenize{compute_msd:description}}
\sphinxAtStartPar
Define a computation that calculates the mean\sphinxhyphen{}squared displacement
(MSD) of the group of atoms, including all effects due to atoms
passing through periodic boundaries.  For computation of the non\sphinxhyphen{}Gaussian
parameter of mean\sphinxhyphen{}squared displacement, see the
{\hyperref[\detokenize{compute_msd_nongauss::doc}]{\sphinxcrossref{\DUrole{doc}{compute msd/nongauss}}}} command.

\sphinxAtStartPar
A vector of four quantities is calculated by this compute.  The first three
elements of the vector are the squared \sphinxstyleemphasis{dx}, \sphinxstyleemphasis{dy}, and \sphinxstyleemphasis{dz} displacements,
summed and averaged over atoms in the group.  The fourth element is the total
squared displacement (i.e., \(dx^2 + dy^2 + dz^2\)), summed and
averaged over atoms in the group.

\sphinxAtStartPar
The slope of the mean\sphinxhyphen{}squared displacement (MSD) versus time is
proportional to the diffusion coefficient of the diffusing atoms.

\sphinxAtStartPar
The displacement of an atom is from its reference position. This is
normally the original position at the time
the compute command was issued, unless the \sphinxstyleemphasis{average} keyword is set to \sphinxstyleemphasis{yes}.
The value of the displacement will be
0.0 for atoms not in the specified compute group.

\sphinxAtStartPar
If the \sphinxstyleemphasis{com} option is set to \sphinxstyleemphasis{yes} then the effect of any drift in
the center\sphinxhyphen{}of\sphinxhyphen{}mass of the group of atoms is subtracted out before the
displacement of each atom is calculated.

\sphinxAtStartPar
If the \sphinxstyleemphasis{average} option is set to \sphinxstyleemphasis{yes} then the reference position of
an atom is based on the average position of that atom, corrected for
center\sphinxhyphen{}of\sphinxhyphen{}mass motion if requested.  The average position is a running
average over all previous calls to the compute, including the current
call. So on the first call it is current position, on the second call
it is the arithmetic average of the current position and the position
on the first call, and so on.  Note that when using this option, the
precise value of the mean square displacement will depend on the
number of times the compute is called. So, for example, changing the
frequency of thermo output may change the computed displacement. Also,
the precise values will be changed if a single simulation is broken up
into two parts, using either multiple run commands or a restart
file. It only makes sense to use this option if the atoms are not
diffusing, so that their average positions relative to the center of
mass of the system are stationary. The most common case is crystalline
solids undergoing thermal motion.

\begin{sphinxadmonition}{note}{Note:}
\sphinxAtStartPar
Initial coordinates are stored in “unwrapped” form, by using the
image flags associated with each atom.  See the {\hyperref[\detokenize{dump::doc}]{\sphinxcrossref{\DUrole{doc}{dump custom}}}} command for a discussion of “unwrapped” coordinates.  See the
Atoms section of the {\hyperref[\detokenize{read_data::doc}]{\sphinxcrossref{\DUrole{doc}{read\_data}}}} command for a
discussion of image flags and how they are set for each atom.  You
can reset the image flags (e.g. to 0) before invoking this compute by
using the {\hyperref[\detokenize{set::doc}]{\sphinxcrossref{\DUrole{doc}{set image}}}} command.
\end{sphinxadmonition}

\begin{sphinxadmonition}{note}{Note:}
\sphinxAtStartPar
If you want the quantities calculated by this compute to be
continuous when running from a {\hyperref[\detokenize{read_restart::doc}]{\sphinxcrossref{\DUrole{doc}{restart file}}}}, then
you should use the same ID for this compute, as in the original run.
This is so that the fix this compute creates to store per\sphinxhyphen{}atom
quantities will also have the same ID, and thus be initialized
correctly with atom reference positions from the restart file.  When
\sphinxstyleemphasis{average} is set to yes, then the atom reference positions are
restored correctly, but not the number of samples used obtain them. As
a result, the reference positions from the restart file are combined
with subsequent positions as if they were from a single sample,
instead of many, which will change the values of msd somewhat.
\end{sphinxadmonition}


\subsection{Output info}
\label{\detokenize{compute_msd:output-info}}
\sphinxAtStartPar
This compute calculates a global vector of length 4, which can be
accessed by indices 1\textendash{}4 by any command that uses global vector values
from a compute as input.  See the {\hyperref[\detokenize{Howto_output::doc}]{\sphinxcrossref{\DUrole{doc}{Howto output}}}} doc
page for an overview of LAMMPS output options.

\sphinxAtStartPar
The vector values are “intensive”.  The vector values will be in
distance\(^2\) {\hyperref[\detokenize{units::doc}]{\sphinxcrossref{\DUrole{doc}{units}}}}.


\subsection{Restrictions}
\label{\detokenize{compute_msd:restrictions}}
\sphinxAtStartPar
Compute \sphinxstyleemphasis{msd} cannot be used with a dynamic group.


\subsection{Related commands}
\label{\detokenize{compute_msd:related-commands}}
\sphinxAtStartPar
{\hyperref[\detokenize{compute_msd_nongauss::doc}]{\sphinxcrossref{\DUrole{doc}{compute msd/nongauss}}}}, {\hyperref[\detokenize{compute_displace_atom::doc}]{\sphinxcrossref{\DUrole{doc}{compute displace\_atom}}}}, {\hyperref[\detokenize{fix_store_state::doc}]{\sphinxcrossref{\DUrole{doc}{fix store/state}}}}, {\hyperref[\detokenize{compute_msd_chunk::doc}]{\sphinxcrossref{\DUrole{doc}{compute msd/chunk}}}}


\subsection{Default}
\label{\detokenize{compute_msd:default}}
\sphinxAtStartPar
The option default are com = no, average = no.

\sphinxstepscope

\index{compute msd/chunk@\spxentry{compute msd/chunk}}\ignorespaces 

\section{compute msd/chunk command}
\label{\detokenize{compute_msd_chunk:compute-msd-chunk-command}}\label{\detokenize{compute_msd_chunk:index-0}}\label{\detokenize{compute_msd_chunk::doc}}

\subsection{Syntax}
\label{\detokenize{compute_msd_chunk:syntax}}
\begin{sphinxVerbatim}[commandchars=\\\{\}]
\PYG{k}{compute }\PYG{n+nv+nvIdentifier}{ID}\PYG{+w}{ }\PYG{n+nv+nvIdentifier}{group\PYGZhy{}ID}\PYG{+w}{ }\PYG{n}{msd}\PYG{o}{/}\PYG{n}{chunk}\PYG{+w}{ }\PYG{n}{chunkID}
\end{sphinxVerbatim}
\begin{itemize}
\item {} 
\sphinxAtStartPar
ID, group\sphinxhyphen{}ID are documented in {\hyperref[\detokenize{compute::doc}]{\sphinxcrossref{\DUrole{doc}{compute}}}} command

\item {} 
\sphinxAtStartPar
msd/chunk = style name of this compute command

\item {} 
\sphinxAtStartPar
chunkID = ID of {\hyperref[\detokenize{compute_chunk_atom::doc}]{\sphinxcrossref{\DUrole{doc}{compute chunk/atom}}}} command

\end{itemize}


\subsection{Examples}
\label{\detokenize{compute_msd_chunk:examples}}
\begin{sphinxVerbatim}[commandchars=\\\{\}]
\PYG{k}{compute }\PYG{n+nv+nvIdentifier}{1}\PYG{+w}{ }\PYG{n+nv+nvIdentifier}{all}\PYG{+w}{ }\PYG{n}{msd}\PYG{o}{/}\PYG{n}{chunk}\PYG{+w}{ }\PYG{n}{molchunk}
\end{sphinxVerbatim}


\subsection{Description}
\label{\detokenize{compute_msd_chunk:description}}
\sphinxAtStartPar
Define a computation that calculates the mean\sphinxhyphen{}squared displacement
(MSD) for multiple chunks of atoms.

\sphinxAtStartPar
In LAMMPS, chunks are collections of atoms defined by a
{\hyperref[\detokenize{compute_chunk_atom::doc}]{\sphinxcrossref{\DUrole{doc}{compute chunk/atom}}}} command, which assigns each atom
to a single chunk (or no chunk).  The ID for this command is specified
as chunkID.  For example, a single chunk could be the atoms in a
molecule or atoms in a spatial bin.  See the
{\hyperref[\detokenize{compute_chunk_atom::doc}]{\sphinxcrossref{\DUrole{doc}{compute chunk/atom}}}} and
{\hyperref[\detokenize{Howto_chunk::doc}]{\sphinxcrossref{\DUrole{doc}{Howto chunk}}}}
doc pages for details of how chunks can be defined and examples of how
they can be used to measure properties of a system.

\sphinxAtStartPar
Four quantities are calculated by this compute for each chunk.  The
first 3 quantities are the squared \sphinxstyleemphasis{dx}, \sphinxstyleemphasis{dy}, and \sphinxstyleemphasis{dz} displacements of the
center\sphinxhyphen{}of\sphinxhyphen{}mass.  The fourth component is the total squared displacement
(i.e., \(dx^2 + dy^2 + dz^2\)) of the center\sphinxhyphen{}of\sphinxhyphen{}mass.  These calculations
include all effects due to atoms passing through periodic boundaries.

\sphinxAtStartPar
Note that only atoms in the specified group contribute to the
calculation.  The {\hyperref[\detokenize{compute_chunk_atom::doc}]{\sphinxcrossref{\DUrole{doc}{compute chunk/atom}}}} command
defines its own group; atoms will have a chunk ID = 0 if they are not
in that group, signifying they are not assigned to a chunk, and will
thus also not contribute to this calculation.  You can specify the
“all” group for this command if you simply want to include atoms with
non\sphinxhyphen{}zero chunk IDs.

\sphinxAtStartPar
The slope of the mean\sphinxhyphen{}squared displacement (MSD) versus time is
proportional to the diffusion coefficient of the diffusing chunks.

\sphinxAtStartPar
The displacement of the center\sphinxhyphen{}of\sphinxhyphen{}mass of the chunk is from its
original center\sphinxhyphen{}of\sphinxhyphen{}mass position, calculated on the timestep this
compute command was first invoked.

\begin{sphinxadmonition}{note}{Note:}
\sphinxAtStartPar
The number of chunks \sphinxstyleemphasis{Nchunk} calculated by the
{\hyperref[\detokenize{compute_chunk_atom::doc}]{\sphinxcrossref{\DUrole{doc}{compute chunk/atom}}}} command must remain constant
each time this compute is invoked, so that the displacement for each chunk
from its original position can be computed consistently.  If \sphinxstyleemphasis{Nchunk}
does not remain constant, an error will be generated.  If needed, you
can enforce a constant \sphinxstyleemphasis{Nchunk} by using the \sphinxstyleemphasis{nchunk once} or \sphinxstyleemphasis{ids once}
options when specifying the {\hyperref[\detokenize{compute_chunk_atom::doc}]{\sphinxcrossref{\DUrole{doc}{compute chunk/atom}}}}
command.
\end{sphinxadmonition}

\begin{sphinxadmonition}{note}{Note:}
\sphinxAtStartPar
This compute stores the original position (of the
center\sphinxhyphen{}of\sphinxhyphen{}mass) of each chunk.  When a displacement is calculated on a
later timestep, it is assumed that the same atoms are assigned to the
same chunk ID.  However LAMMPS has no simple way to ensure this is the
case, though you can use the \sphinxstyleemphasis{ids once} option when specifying the
{\hyperref[\detokenize{compute_chunk_atom::doc}]{\sphinxcrossref{\DUrole{doc}{compute chunk/atom}}}} command.  Note that if
this is not the case, the MSD calculation does not have a sensible
meaning.
\end{sphinxadmonition}

\begin{sphinxadmonition}{note}{Note:}
\sphinxAtStartPar
The initial coordinates of the atoms in each chunk are stored in
“unwrapped” form, by using the image flags associated with each atom.
See the {\hyperref[\detokenize{dump::doc}]{\sphinxcrossref{\DUrole{doc}{dump custom}}}} command for a discussion of
“unwrapped” coordinates.  See the Atoms section of the
{\hyperref[\detokenize{read_data::doc}]{\sphinxcrossref{\DUrole{doc}{read\_data}}}} command for a discussion of image flags and
how they are set for each atom.  You can reset the image flags
(e.g., to 0) before invoking this compute by using the
{\hyperref[\detokenize{set::doc}]{\sphinxcrossref{\DUrole{doc}{set image}}}} command.
\end{sphinxadmonition}

\begin{sphinxadmonition}{note}{Note:}
\sphinxAtStartPar
If you want the quantities calculated by this compute to be
continuous when running from a {\hyperref[\detokenize{read_restart::doc}]{\sphinxcrossref{\DUrole{doc}{restart file}}}}, then
you should use the same ID for this compute, as in the original run.
This is so that the fix this compute creates to store per\sphinxhyphen{}chunk
quantities will also have the same ID, and thus be initialized
correctly with chunk reference positions from the restart file.
\end{sphinxadmonition}

\sphinxAtStartPar
The simplest way to output the results of the compute msd/chunk
calculation to a file is to use the {\hyperref[\detokenize{fix_ave_time::doc}]{\sphinxcrossref{\DUrole{doc}{fix ave/time}}}}
command, for example:

\begin{sphinxVerbatim}[commandchars=\\\{\}]
\PYG{k}{compute }\PYG{n+nv+nvIdentifier}{cc1}\PYG{+w}{ }\PYG{n+nv+nvIdentifier}{all}\PYG{+w}{ }\PYG{n}{chunk}\PYG{o}{/}\PYG{n}{atom}\PYG{+w}{ }\PYG{n}{molecule}
\PYG{k}{compute }\PYG{n+nv+nvIdentifier}{myChunk}\PYG{+w}{ }\PYG{n+nv+nvIdentifier}{all}\PYG{+w}{ }\PYG{n}{msd}\PYG{o}{/}\PYG{n}{chunk}\PYG{+w}{ }\PYG{n}{cc1}
\PYG{k}{fix }\PYG{n+nv+nvIdentifier}{1}\PYG{+w}{ }\PYG{n+nv+nvIdentifier}{all}\PYG{+w}{ }\PYG{n}{ave}\PYG{o}{/}\PYG{n}{time}\PYG{+w}{ }\PYG{l+m}{100}\PYG{+w}{ }\PYG{l+m}{1}\PYG{+w}{ }\PYG{l+m}{100}\PYG{+w}{ }\PYG{n}{c\PYGZus{}myChunk[}\PYG{o}{*}\PYG{n}{]}\PYG{+w}{ }\PYG{n}{file}\PYG{+w}{ }\PYG{n}{tmp.out}\PYG{+w}{ }\PYG{n}{mode}\PYG{+w}{ }\PYG{n}{vector}
\end{sphinxVerbatim}


\subsection{Output info}
\label{\detokenize{compute_msd_chunk:output-info}}
\sphinxAtStartPar
This compute calculates a global array where the number of rows = the
number of chunks \sphinxstyleemphasis{Nchunk} as calculated by the specified
{\hyperref[\detokenize{compute_chunk_atom::doc}]{\sphinxcrossref{\DUrole{doc}{compute chunk/atom}}}} command.
The number of columns = 4 for \sphinxstyleemphasis{dx}, \sphinxstyleemphasis{dy}, \sphinxstyleemphasis{dz}, and the total displacement.
These values can be accessed by any command that uses global array values from
a compute as input.  See the {\hyperref[\detokenize{Howto_output::doc}]{\sphinxcrossref{\DUrole{doc}{Howto output}}}} page for an
overview of LAMMPS output options.

\sphinxAtStartPar
The array values are “intensive”.  The array values will be in
distance\(^2\) {\hyperref[\detokenize{units::doc}]{\sphinxcrossref{\DUrole{doc}{units}}}}.


\subsection{Restrictions}
\label{\detokenize{compute_msd_chunk:restrictions}}\begin{quote}

\sphinxAtStartPar
none
\end{quote}


\subsection{Related commands}
\label{\detokenize{compute_msd_chunk:related-commands}}
\sphinxAtStartPar
{\hyperref[\detokenize{compute_msd::doc}]{\sphinxcrossref{\DUrole{doc}{compute msd}}}}


\subsection{Default}
\label{\detokenize{compute_msd_chunk:default}}
\sphinxAtStartPar
none

\sphinxstepscope

\index{compute msd/nongauss@\spxentry{compute msd/nongauss}}\ignorespaces 

\section{compute msd/nongauss command}
\label{\detokenize{compute_msd_nongauss:compute-msd-nongauss-command}}\label{\detokenize{compute_msd_nongauss:index-0}}\label{\detokenize{compute_msd_nongauss::doc}}

\subsection{Syntax}
\label{\detokenize{compute_msd_nongauss:syntax}}
\begin{sphinxVerbatim}[commandchars=\\\{\}]
\PYG{k}{compute }\PYG{n+nv+nvIdentifier}{ID}\PYG{+w}{ }\PYG{n+nv+nvIdentifier}{group\PYGZhy{}ID}\PYG{+w}{ }\PYG{n}{msd}\PYG{o}{/}\PYG{n}{nongauss}\PYG{+w}{ }\PYG{n}{keyword}\PYG{+w}{ }\PYG{n}{values}\PYG{+w}{ }\PYG{n}{...}
\end{sphinxVerbatim}
\begin{itemize}
\item {} 
\sphinxAtStartPar
ID, group\sphinxhyphen{}ID are documented in {\hyperref[\detokenize{compute::doc}]{\sphinxcrossref{\DUrole{doc}{compute}}}} command

\item {} 
\sphinxAtStartPar
msd/nongauss = style name of this compute command

\item {} 
\sphinxAtStartPar
zero or more keyword/value pairs may be appended

\item {} 
\sphinxAtStartPar
keyword = \sphinxstyleemphasis{com}
\begin{sphinxalltt}
\sphinxstyleemphasis{com} value = \sphinxstyleemphasis{yes} or \sphinxstyleemphasis{no}
\end{sphinxalltt}

\end{itemize}


\subsection{Examples}
\label{\detokenize{compute_msd_nongauss:examples}}
\begin{sphinxVerbatim}[commandchars=\\\{\}]
\PYG{k}{compute }\PYG{n+nv+nvIdentifier}{1}\PYG{+w}{ }\PYG{n+nv+nvIdentifier}{all}\PYG{+w}{ }\PYG{n}{msd}\PYG{o}{/}\PYG{n}{nongauss}
\PYG{k}{compute }\PYG{n+nv+nvIdentifier}{1}\PYG{+w}{ }\PYG{n+nv+nvIdentifier}{upper}\PYG{+w}{ }\PYG{n}{msd}\PYG{o}{/}\PYG{n}{nongauss}\PYG{+w}{ }\PYG{n}{com}\PYG{+w}{ }\PYG{n}{yes}
\end{sphinxVerbatim}


\subsection{Description}
\label{\detokenize{compute_msd_nongauss:description}}
\sphinxAtStartPar
Define a computation that calculates the mean\sphinxhyphen{}squared displacement
(MSD) and non\sphinxhyphen{}Gaussian parameter (NGP) of the group of atoms,
including all effects due to atoms passing through periodic boundaries.

\sphinxAtStartPar
A vector of three quantities is calculated by this compute.  The first
element of the vector is the total squared displacement,
\(dr^2 = dx^2 + dy^2 + dz^2\), of the atoms, and the second is the
fourth power of these displacements, \(dr^4 = (dx^2 + dy^2 + dz^2)^2\),
summed and averaged over atoms in the group.  The third component is the
non\sphinxhyphen{}Gaussian diffusion parameter NGP,
\begin{equation*}
\begin{split}\text{NGP}(t) = \frac{3\left\langle(r(t)-r(0))^4\right\rangle}
                     {5\left\langle(r(t)-r(0))^2\right\rangle^2} - 1.\end{split}
\end{equation*}
\sphinxAtStartPar
The NGP is a commonly used quantity in studies of dynamical
heterogeneity.  Its minimum theoretical value \((-0.4)\) occurs when all
atoms have the same displacement magnitude.  \(\text{NGP}=0\) for Brownian
diffusion, while \(\text{NGP} > 0\) when some mobile atoms move faster than
others.

\sphinxAtStartPar
If the \sphinxstyleemphasis{com} option is set to \sphinxstyleemphasis{yes} then the effect of any drift in
the center\sphinxhyphen{}of\sphinxhyphen{}mass of the group of atoms is subtracted out before the
displacement of each atom is calculated.

\sphinxAtStartPar
See the {\hyperref[\detokenize{compute_msd::doc}]{\sphinxcrossref{\DUrole{doc}{compute msd}}}} page for further important
NOTEs, which also apply to this compute.


\subsection{Output info}
\label{\detokenize{compute_msd_nongauss:output-info}}
\sphinxAtStartPar
This compute calculates a global vector of length 3, which can be
accessed by indices 1\textendash{}3 by any command that uses global vector values
from a compute as input.  See the {\hyperref[\detokenize{Howto_output::doc}]{\sphinxcrossref{\DUrole{doc}{Howto output}}}} doc
page for an overview of LAMMPS output options.

\sphinxAtStartPar
The vector values are “intensive”.  The first vector value will be in
distance\(^2\) {\hyperref[\detokenize{units::doc}]{\sphinxcrossref{\DUrole{doc}{units}}}}, the second is in
distance\(^4\) units, and the third is dimensionless.


\subsection{Restrictions}
\label{\detokenize{compute_msd_nongauss:restrictions}}
\sphinxAtStartPar
Compute \sphinxstyleemphasis{msd/nongauss} cannot be used with a dynamic group.

\sphinxAtStartPar
This compute is part of the EXTRA\sphinxhyphen{}COMPUTE package.  It is only enabled
if LAMMPS was built with that package.  See the {\hyperref[\detokenize{Build_package::doc}]{\sphinxcrossref{\DUrole{doc}{Build package}}}} page for more info.


\subsection{Related commands}
\label{\detokenize{compute_msd_nongauss:related-commands}}
\sphinxAtStartPar
{\hyperref[\detokenize{compute_msd::doc}]{\sphinxcrossref{\DUrole{doc}{compute msd}}}}


\subsection{Default}
\label{\detokenize{compute_msd_nongauss:default}}
\sphinxAtStartPar
The option default is com = no.

\sphinxstepscope

\index{compute nbond/atom@\spxentry{compute nbond/atom}}\ignorespaces 

\section{compute nbond/atom command}
\label{\detokenize{compute_nbond_atom:compute-nbond-atom-command}}\label{\detokenize{compute_nbond_atom:index-0}}\label{\detokenize{compute_nbond_atom::doc}}

\subsection{Syntax}
\label{\detokenize{compute_nbond_atom:syntax}}
\begin{sphinxVerbatim}[commandchars=\\\{\}]
\PYG{k}{compute }\PYG{n+nv+nvIdentifier}{ID}\PYG{+w}{ }\PYG{n+nv+nvIdentifier}{group\PYGZhy{}ID}\PYG{+w}{ }\PYG{n}{nbond}\PYG{o}{/}\PYG{n}{atom}\PYG{+w}{ }\PYG{n}{keyword}\PYG{+w}{ }\PYG{n}{value}
\end{sphinxVerbatim}
\begin{itemize}
\item {} 
\sphinxAtStartPar
ID, group\sphinxhyphen{}ID are documented in {\hyperref[\detokenize{compute::doc}]{\sphinxcrossref{\DUrole{doc}{compute}}}} command

\item {} 
\sphinxAtStartPar
nbond/atom = style name of this compute command

\item {} 
\sphinxAtStartPar
zero or more keyword/value pairs may be appended

\item {} 
\sphinxAtStartPar
keyword = \sphinxstyleemphasis{bond/type}
\begin{sphinxalltt}
\sphinxstyleemphasis{bond/type} value = \sphinxstyleemphasis{btype}
  \sphinxstyleemphasis{btype} = bond type included in count
\end{sphinxalltt}

\end{itemize}


\subsection{Examples}
\label{\detokenize{compute_nbond_atom:examples}}
\begin{sphinxVerbatim}[commandchars=\\\{\}]
\PYG{k}{compute }\PYG{n+nv+nvIdentifier}{1}\PYG{+w}{ }\PYG{n+nv+nvIdentifier}{all}\PYG{+w}{ }\PYG{n}{nbond}\PYG{o}{/}\PYG{n}{atom}
\PYG{k}{compute }\PYG{n+nv+nvIdentifier}{1}\PYG{+w}{ }\PYG{n+nv+nvIdentifier}{all}\PYG{+w}{ }\PYG{n}{nbond}\PYG{o}{/}\PYG{n}{atom}\PYG{+w}{ }\PYG{n}{bond}\PYG{o}{/}\PYG{n}{type}\PYG{+w}{ }\PYG{l+m}{2}
\end{sphinxVerbatim}


\subsection{Description}
\label{\detokenize{compute_nbond_atom:description}}
\sphinxAtStartPar
\DUrole{versionmodified,added}{Added in version 4May2022.}

\sphinxAtStartPar
Define a computation that computes the number of bonds each atom is
part of.  Bonds which are broken are not counted in the tally.  See
the {\hyperref[\detokenize{Howto_bpm::doc}]{\sphinxcrossref{\DUrole{doc}{Howto broken bonds}}}} page for more information.
The number of bonds will be zero for atoms not in the specified
compute group. This compute does not depend on Newton bond settings.

\sphinxAtStartPar
If the keyword \sphinxstyleemphasis{bond/type} is specified, only bonds of \sphinxstyleemphasis{btype} are
counted.


\subsection{Output info}
\label{\detokenize{compute_nbond_atom:output-info}}
\sphinxAtStartPar
This compute calculates a per\sphinxhyphen{}atom vector, which can be accessed by
any command that uses per\sphinxhyphen{}atom values from a compute as input.  See
the {\hyperref[\detokenize{Howto_output::doc}]{\sphinxcrossref{\DUrole{doc}{Howto output}}}} doc page for an overview of
LAMMPS output options.


\subsection{Restrictions}
\label{\detokenize{compute_nbond_atom:restrictions}}
\sphinxAtStartPar
This compute is part of the BPM package.  It is only enabled if LAMMPS was
built with that package.  See the {\hyperref[\detokenize{Build_package::doc}]{\sphinxcrossref{\DUrole{doc}{Build package}}}}
page for more info.


\subsection{Related commands}
\label{\detokenize{compute_nbond_atom:related-commands}}

\subsection{Default}
\label{\detokenize{compute_nbond_atom:default}}
\sphinxAtStartPar
none

\sphinxstepscope

\index{compute omega/chunk@\spxentry{compute omega/chunk}}\ignorespaces 

\section{compute omega/chunk command}
\label{\detokenize{compute_omega_chunk:compute-omega-chunk-command}}\label{\detokenize{compute_omega_chunk:index-0}}\label{\detokenize{compute_omega_chunk::doc}}

\subsection{Syntax}
\label{\detokenize{compute_omega_chunk:syntax}}
\begin{sphinxVerbatim}[commandchars=\\\{\}]
\PYG{k}{compute }\PYG{n+nv+nvIdentifier}{ID}\PYG{+w}{ }\PYG{n+nv+nvIdentifier}{group\PYGZhy{}ID}\PYG{+w}{ }\PYG{n}{omega}\PYG{o}{/}\PYG{n}{chunk}\PYG{+w}{ }\PYG{n}{chunkID}
\end{sphinxVerbatim}
\begin{itemize}
\item {} 
\sphinxAtStartPar
ID, group\sphinxhyphen{}ID are documented in {\hyperref[\detokenize{compute::doc}]{\sphinxcrossref{\DUrole{doc}{compute}}}} command

\item {} 
\sphinxAtStartPar
omega/chunk = style name of this compute command

\item {} 
\sphinxAtStartPar
chunkID = ID of {\hyperref[\detokenize{compute_chunk_atom::doc}]{\sphinxcrossref{\DUrole{doc}{compute chunk/atom}}}} command

\end{itemize}


\subsection{Examples}
\label{\detokenize{compute_omega_chunk:examples}}
\begin{sphinxVerbatim}[commandchars=\\\{\}]
\PYG{k}{compute }\PYG{n+nv+nvIdentifier}{1}\PYG{+w}{ }\PYG{n+nv+nvIdentifier}{fluid}\PYG{+w}{ }\PYG{n}{omega}\PYG{o}{/}\PYG{n}{chunk}\PYG{+w}{ }\PYG{n}{molchunk}
\end{sphinxVerbatim}


\subsection{Description}
\label{\detokenize{compute_omega_chunk:description}}
\sphinxAtStartPar
Define a computation that calculates the angular velocity (omega) of
multiple chunks of atoms.

\sphinxAtStartPar
In LAMMPS, chunks are collections of atoms defined by a
{\hyperref[\detokenize{compute_chunk_atom::doc}]{\sphinxcrossref{\DUrole{doc}{compute chunk/atom}}}} command, which assigns each atom
to a single chunk (or no chunk).  The ID for this command is specified
as chunkID.  For example, a single chunk could be the atoms in a
molecule or atoms in a spatial bin.  See the
{\hyperref[\detokenize{compute_chunk_atom::doc}]{\sphinxcrossref{\DUrole{doc}{compute chunk/atom}}}} and
{\hyperref[\detokenize{Howto_chunk::doc}]{\sphinxcrossref{\DUrole{doc}{Howto chunk}}}}
doc pages for details of how chunks can be defined and examples of how
they can be used to measure properties of a system.

\sphinxAtStartPar
This compute calculates the three components of the angular velocity
vector for each chunk via the formula
\(\vec L = \mathrm{I}\cdot \vec\omega\), where \(\vec L\) is the
angular momentum vector of the chunk, \(\mathrm{I}\) is its moment of
inertia tensor, and \(\omega\) is the angular velocity of the chunk.
The calculation includes all effects due to atoms passing through periodic
boundaries.

\sphinxAtStartPar
Note that only atoms in the specified group contribute to the
calculation.  The {\hyperref[\detokenize{compute_chunk_atom::doc}]{\sphinxcrossref{\DUrole{doc}{compute chunk/atom}}}} command
defines its own group; atoms will have a chunk ID = 0 if they are not
in that group, signifying they are not assigned to a chunk, and will
thus also not contribute to this calculation.  You can specify the
“all” group for this command if you simply want to include atoms with
non\sphinxhyphen{}zero chunk IDs.

\begin{sphinxadmonition}{note}{Note:}
\sphinxAtStartPar
The coordinates of an atom contribute to the chunk’s angular
velocity in “unwrapped” form, by using the image flags associated with
each atom.  See the {\hyperref[\detokenize{dump::doc}]{\sphinxcrossref{\DUrole{doc}{dump custom}}}} command for a discussion
of “unwrapped” coordinates.  See the Atoms section of the
{\hyperref[\detokenize{read_data::doc}]{\sphinxcrossref{\DUrole{doc}{read\_data}}}} command for a discussion of image flags and
how they are set for each atom.  You can reset the image flags
(e.g., to 0) before invoking this compute by using the
{\hyperref[\detokenize{set::doc}]{\sphinxcrossref{\DUrole{doc}{set image}}}} command.
\end{sphinxadmonition}

\sphinxAtStartPar
The simplest way to output the results of the compute omega/chunk
calculation to a file is to use the {\hyperref[\detokenize{fix_ave_time::doc}]{\sphinxcrossref{\DUrole{doc}{fix ave/time}}}}
command, for example:

\begin{sphinxVerbatim}[commandchars=\\\{\}]
\PYG{k}{compute }\PYG{n+nv+nvIdentifier}{cc1}\PYG{+w}{ }\PYG{n+nv+nvIdentifier}{all}\PYG{+w}{ }\PYG{n}{chunk}\PYG{o}{/}\PYG{n}{atom}\PYG{+w}{ }\PYG{n}{molecule}
\PYG{k}{compute }\PYG{n+nv+nvIdentifier}{myChunk}\PYG{+w}{ }\PYG{n+nv+nvIdentifier}{all}\PYG{+w}{ }\PYG{n}{omega}\PYG{o}{/}\PYG{n}{chunk}\PYG{+w}{ }\PYG{n}{cc1}
\PYG{k}{fix }\PYG{n+nv+nvIdentifier}{1}\PYG{+w}{ }\PYG{n+nv+nvIdentifier}{all}\PYG{+w}{ }\PYG{n}{ave}\PYG{o}{/}\PYG{n}{time}\PYG{+w}{ }\PYG{l+m}{100}\PYG{+w}{ }\PYG{l+m}{1}\PYG{+w}{ }\PYG{l+m}{100}\PYG{+w}{ }\PYG{n}{c\PYGZus{}myChunk[}\PYG{o}{*}\PYG{n}{]}\PYG{+w}{ }\PYG{n}{file}\PYG{+w}{ }\PYG{n}{tmp.out}\PYG{+w}{ }\PYG{n}{mode}\PYG{+w}{ }\PYG{n}{vector}
\end{sphinxVerbatim}


\subsection{Output info}
\label{\detokenize{compute_omega_chunk:output-info}}
\sphinxAtStartPar
This compute calculates a global array where the number of rows is the
number of chunks \sphinxstyleemphasis{Nchunk} as calculated by the specified {\hyperref[\detokenize{compute_chunk_atom::doc}]{\sphinxcrossref{\DUrole{doc}{compute chunk/atom}}}} command.  The number of columns is 3 for the three
(\sphinxstyleemphasis{x}, \sphinxstyleemphasis{y}, \sphinxstyleemphasis{z}) components of the angular velocity for each chunk.
These values can be accessed by any command that uses global array
values from a compute as input.  See the {\hyperref[\detokenize{Howto_output::doc}]{\sphinxcrossref{\DUrole{doc}{Howto output}}}}
page for an overview of LAMMPS output options.

\sphinxAtStartPar
The array values are “intensive”.  The array values will be in
velocity/distance {\hyperref[\detokenize{units::doc}]{\sphinxcrossref{\DUrole{doc}{units}}}}.


\subsection{Restrictions}
\label{\detokenize{compute_omega_chunk:restrictions}}\begin{quote}

\sphinxAtStartPar
none
\end{quote}


\subsection{Related commands}
\label{\detokenize{compute_omega_chunk:related-commands}}
\sphinxAtStartPar
{\hyperref[\detokenize{variable::doc}]{\sphinxcrossref{\DUrole{doc}{variable omega() function}}}}


\subsection{Default}
\label{\detokenize{compute_omega_chunk:default}}
\sphinxAtStartPar
none

\sphinxstepscope

\index{compute orientorder/atom@\spxentry{compute orientorder/atom}}\ignorespaces 
\index{compute orientorder/atom/kk@\spxentry{compute orientorder/atom/kk}}\ignorespaces 

\section{compute orientorder/atom command}
\label{\detokenize{compute_orientorder_atom:compute-orientorder-atom-command}}\label{\detokenize{compute_orientorder_atom:index-1}}\label{\detokenize{compute_orientorder_atom:index-0}}\label{\detokenize{compute_orientorder_atom::doc}}
\sphinxAtStartPar
Accelerator Variants: \sphinxstyleemphasis{orientorder/atom/kk}


\subsection{Syntax}
\label{\detokenize{compute_orientorder_atom:syntax}}
\begin{sphinxVerbatim}[commandchars=\\\{\}]
\PYG{k}{compute }\PYG{n+nv+nvIdentifier}{ID}\PYG{+w}{ }\PYG{n+nv+nvIdentifier}{group\PYGZhy{}ID}\PYG{+w}{ }\PYG{n}{orientorder}\PYG{o}{/}\PYG{n}{atom}\PYG{+w}{ }\PYG{n}{keyword}\PYG{+w}{ }\PYG{n}{values}\PYG{+w}{ }\PYG{n}{...}
\end{sphinxVerbatim}
\begin{itemize}
\item {} 
\sphinxAtStartPar
ID, group\sphinxhyphen{}ID are documented in {\hyperref[\detokenize{compute::doc}]{\sphinxcrossref{\DUrole{doc}{compute}}}} command

\item {} 
\sphinxAtStartPar
orientorder/atom = style name of this compute command

\item {} 
\sphinxAtStartPar
one or more keyword/value pairs may be appended
\begin{sphinxalltt}
keyword = \sphinxstyleemphasis{cutoff} or \sphinxstyleemphasis{nnn} or \sphinxstyleemphasis{degrees} or \sphinxstyleemphasis{wl} or \sphinxstyleemphasis{wl/hat} or \sphinxstyleemphasis{components} or \sphinxstyleemphasis{chunksize}
  \sphinxstyleemphasis{cutoff} value = distance cutoff
  \sphinxstyleemphasis{nnn} value = number of nearest neighbors
  \sphinxstyleemphasis{degrees} values = nlvalues, l1, l2,...
  \sphinxstyleemphasis{wl} value = \sphinxstyleemphasis{yes} or \sphinxstyleemphasis{no}
  \sphinxstyleemphasis{wl/hat} value = \sphinxstyleemphasis{yes} or \sphinxstyleemphasis{no}
  \sphinxstyleemphasis{components} value = ldegree
  \sphinxstyleemphasis{chunksize} value = number of atoms in each pass
\end{sphinxalltt}

\end{itemize}


\subsection{Examples}
\label{\detokenize{compute_orientorder_atom:examples}}
\begin{sphinxVerbatim}[commandchars=\\\{\}]
\PYG{k}{compute }\PYG{n+nv+nvIdentifier}{1}\PYG{+w}{ }\PYG{n+nv+nvIdentifier}{all}\PYG{+w}{ }\PYG{n}{orientorder}\PYG{o}{/}\PYG{n}{atom}
\PYG{k}{compute }\PYG{n+nv+nvIdentifier}{1}\PYG{+w}{ }\PYG{n+nv+nvIdentifier}{all}\PYG{+w}{ }\PYG{n}{orientorder}\PYG{o}{/}\PYG{n}{atom}\PYG{+w}{ }\PYG{n}{degrees}\PYG{+w}{ }\PYG{l+m}{5}\PYG{+w}{ }\PYG{l+m}{4}\PYG{+w}{ }\PYG{l+m}{6}\PYG{+w}{ }\PYG{l+m}{8}\PYG{+w}{ }\PYG{l+m}{10}\PYG{+w}{ }\PYG{l+m}{12}\PYG{+w}{ }\PYG{n}{nnn}\PYG{+w}{ }\PYG{n}{NULL}\PYG{+w}{ }\PYG{n}{cutoff}\PYG{+w}{ }\PYG{l+m}{1.5}
\PYG{k}{compute }\PYG{n+nv+nvIdentifier}{1}\PYG{+w}{ }\PYG{n+nv+nvIdentifier}{all}\PYG{+w}{ }\PYG{n}{orientorder}\PYG{o}{/}\PYG{n}{atom}\PYG{+w}{ }\PYG{n}{wl}\PYG{o}{/}\PYG{n}{hat}\PYG{+w}{ }\PYG{n}{yes}
\PYG{k}{compute }\PYG{n+nv+nvIdentifier}{1}\PYG{+w}{ }\PYG{n+nv+nvIdentifier}{all}\PYG{+w}{ }\PYG{n}{orientorder}\PYG{o}{/}\PYG{n}{atom}\PYG{+w}{ }\PYG{n}{components}\PYG{+w}{ }\PYG{l+m}{6}
\end{sphinxVerbatim}


\subsection{Description}
\label{\detokenize{compute_orientorder_atom:description}}
\sphinxAtStartPar
Define a computation that calculates a set of bond\sphinxhyphen{}orientational
order parameters \(Q_\ell\) for each atom in a group. These order parameters
were introduced by {\hyperref[\detokenize{compute_orientorder_atom:steinhardt}]{\sphinxcrossref{\DUrole{std,std-ref}{Steinhardt et al.}}}} as a way to
characterize the local orientational order in atomic structures.
For each atom, \(Q_\ell\) is a real number defined as follows:
\begin{equation*}
\begin{split}\bar{Y}_{\ell m} = & \frac{1}{nnn}\sum_{j = 1}^{nnn} Y_{\ell m}\bigl( \theta( {\bf r}_{ij} ), \phi( {\bf r}_{ij} ) \bigr) \\
Q_\ell  = & \sqrt{\frac{4 \pi}{2 \ell  + 1} \sum_{m = -\ell }^{m = \ell } \bar{Y}_{\ell m} \bar{Y}^*_{\ell m}}\end{split}
\end{equation*}
\sphinxAtStartPar
The first equation defines the local order parameters as averages
of the spherical harmonics \(Y_{\ell m}\) for each neighbor.
These are complex number components of the 3D analog of the 2D order
parameter \(q_n\), which is implemented as LAMMPS compute
{\hyperref[\detokenize{compute_hexorder_atom::doc}]{\sphinxcrossref{\DUrole{doc}{hexorder/atom}}}}.
The summation is over the \sphinxstyleemphasis{nnn} nearest
neighbors of the central atom.  The angles \(\theta\) and \(\phi\) are
the standard spherical polar angles
defining the direction of the bond vector \(r_{ij}\).
The phase and sign of \(Y_{\ell m}\) follow the standard conventions,
so that \(\mathrm{sign}(Y_{\ell\ell}(0,0)) = (-1)^\ell\).
The second equation defines \(Q_\ell\), which is a
rotationally invariant non\sphinxhyphen{}negative amplitude obtained by summing
over all the components of degree \(\ell\).

\sphinxAtStartPar
The optional keyword \sphinxstyleemphasis{cutoff} defines the distance cutoff
used when searching for neighbors. The default value, also
the maximum allowable value, is the cutoff specified
by the pair style.

\sphinxAtStartPar
The optional keyword \sphinxstyleemphasis{nnn} defines the number of nearest
neighbors used to calculate \(Q_\ell\). The default value is 12.
If the value is NULL, then all neighbors up to the
specified distance cutoff are used.

\sphinxAtStartPar
The optional keyword \sphinxstyleemphasis{degrees} defines the list of order parameters to
be computed.  The first argument \sphinxstyleemphasis{nlvalues} is the number of order
parameters. This is followed by that number of non\sphinxhyphen{}negative integers giving the
degree of each order parameter. Because \(Q_2\) and all odd\sphinxhyphen{}degree order
parameters are zero for atoms in cubic crystals (see
{\hyperref[\detokenize{compute_orientorder_atom:steinhardt}]{\sphinxcrossref{\DUrole{std,std-ref}{Steinhardt}}}}), the default order parameters are \(Q_4\),
\(Q_6\), \(Q_8\), \(Q_{10}\), and \(Q_{12}\). For the FCC
crystal with \sphinxstyleemphasis{nnn} =12,
\begin{equation*}
\begin{split}Q_4 = \sqrt{\frac{7}{192}} \approx 0.19094\end{split}
\end{equation*}
\sphinxAtStartPar
The numerical values of all order
parameters up to \(Q_{12}\) for a range of commonly encountered
high\sphinxhyphen{}symmetry structures are given in Table I of {\hyperref[\detokenize{compute_orientorder_atom:mickel}]{\sphinxcrossref{\DUrole{std,std-ref}{Mickel et al.}}}},
and these can be reproduced with this compute.

\sphinxAtStartPar
The optional keyword \sphinxstyleemphasis{wl} will output the third\sphinxhyphen{}order invariants \(W_\ell\)
(see Eq. 1.4 in {\hyperref[\detokenize{compute_orientorder_atom:steinhardt}]{\sphinxcrossref{\DUrole{std,std-ref}{Steinhardt}}}}) for the same degrees as
for the \(Q_\ell\) parameters. For the FCC crystal with \sphinxstyleemphasis{nnn} = 12,
\begin{equation*}
\begin{split}W_4 = -\sqrt{\frac{14}{143}} \left(\frac{49}{4096}\right) \pi^{-3/2} \approx -0.0006722136\end{split}
\end{equation*}
\sphinxAtStartPar
The optional keyword \sphinxstyleemphasis{wl/hat} will output the normalized third\sphinxhyphen{}order
invariants \(\hat{W}_\ell\) (see Eq. 2.2 in {\hyperref[\detokenize{compute_orientorder_atom:steinhardt}]{\sphinxcrossref{\DUrole{std,std-ref}{Steinhardt}}}})
for the same degrees as for the \(Q_\ell\) parameters. For the FCC crystal
with \sphinxstyleemphasis{nnn} =12,
\begin{equation*}
\begin{split}\hat{W}_4 = -\frac{7}{3} \sqrt{\frac{2}{429}} \approx -0.159317\end{split}
\end{equation*}
\sphinxAtStartPar
The numerical values of \(\hat{W}_\ell\) for a range of commonly
encountered high\sphinxhyphen{}symmetry structures are given in Table I of
{\hyperref[\detokenize{compute_orientorder_atom:steinhardt}]{\sphinxcrossref{\DUrole{std,std-ref}{Steinhardt}}}}, and these can be reproduced with this keyword.

\sphinxAtStartPar
The optional keyword \sphinxstyleemphasis{components} will output the components of the
\sphinxstyleemphasis{normalized} complex vector
\(\hat{Y}_{\ell m} = \bar{Y}_{\ell m}/|\bar{Y}_{\ell m}|\)
of degree \sphinxstyleemphasis{ldegree}, which must be included in the list of order parameters to
be computed. This option can be used in conjunction with
{\hyperref[\detokenize{compute_coord_atom::doc}]{\sphinxcrossref{\DUrole{doc}{compute coord\_atom}}}} to calculate the ten Wolde’s
criterion to identify crystal\sphinxhyphen{}like particles, as discussed in
{\hyperref[\detokenize{compute_orientorder_atom:tenwolde2}]{\sphinxcrossref{\DUrole{std,std-ref}{ten Wolde}}}}.

\sphinxAtStartPar
The optional keyword \sphinxstyleemphasis{chunksize} is only applicable when using the
the KOKKOS package and is ignored otherwise. This keyword controls
the number of atoms in each pass used to compute the bond\sphinxhyphen{}orientational
order parameters and is used to avoid running out of memory. For example
if there are 32768 atoms in the simulation and the \sphinxstyleemphasis{chunksize}
is set to 16384, the parameter calculation will be broken up
into two passes.

\sphinxAtStartPar
The value of \(Q_\ell\) is set to zero for atoms not in the
specified compute group, as well as for atoms that have less than
\sphinxstyleemphasis{nnn} neighbors within the distance cutoff, unless \sphinxstyleemphasis{nnn} is NULL.

\sphinxAtStartPar
The neighbor list needed to compute this quantity is constructed each
time the calculation is performed (i.e., each time a snapshot of atoms
is dumped).  Thus it can be inefficient to compute/dump this quantity
too frequently.

\begin{sphinxadmonition}{note}{Note:}
\sphinxAtStartPar
If you have a bonded system, then the settings of
{\hyperref[\detokenize{special_bonds::doc}]{\sphinxcrossref{\DUrole{doc}{special\_bonds}}}} command can remove pairwise
interactions between atoms in the same bond, angle, or dihedral.  This
is the default setting for the {\hyperref[\detokenize{special_bonds::doc}]{\sphinxcrossref{\DUrole{doc}{special\_bonds}}}}
command, and means those pairwise interactions do not appear in the
neighbor list.  Because this fix uses the neighbor list, it also means
those pairs will not be included in the order parameter.  This
difficulty can be circumvented by writing a dump file, and using the
{\hyperref[\detokenize{rerun::doc}]{\sphinxcrossref{\DUrole{doc}{rerun}}}} command to compute the order parameter for
snapshots in the dump file.  The rerun script can use a
{\hyperref[\detokenize{special_bonds::doc}]{\sphinxcrossref{\DUrole{doc}{special\_bonds}}}} command that includes all pairs in
the neighbor list.
\end{sphinxadmonition}


\bigskip\hrule\bigskip


\sphinxAtStartPar
Styles with a \sphinxstyleemphasis{gpu}, \sphinxstyleemphasis{intel}, \sphinxstyleemphasis{kk}, \sphinxstyleemphasis{omp}, or \sphinxstyleemphasis{opt} suffix are
functionally the same as the corresponding style without the suffix.
They have been optimized to run faster, depending on your available
hardware, as discussed on the {\hyperref[\detokenize{Speed_packages::doc}]{\sphinxcrossref{\DUrole{doc}{Accelerator packages}}}}
page.  The accelerated styles take the same arguments and should
produce the same results, except for round\sphinxhyphen{}off and precision issues.

\sphinxAtStartPar
These accelerated styles are part of the GPU, INTEL, KOKKOS,
OPENMP, and OPT packages, respectively.  They are only enabled if
LAMMPS was built with those packages.  See the {\hyperref[\detokenize{Build_package::doc}]{\sphinxcrossref{\DUrole{doc}{Build package}}}} page for more info.

\sphinxAtStartPar
You can specify the accelerated styles explicitly in your input script
by including their suffix, or you can use the {\hyperref[\detokenize{Run_options::doc}]{\sphinxcrossref{\DUrole{doc}{\sphinxhyphen{}suffix command\sphinxhyphen{}line switch}}}} when you invoke LAMMPS, or you can use the
{\hyperref[\detokenize{suffix::doc}]{\sphinxcrossref{\DUrole{doc}{suffix}}}} command in your input script.

\sphinxAtStartPar
See the {\hyperref[\detokenize{Speed_packages::doc}]{\sphinxcrossref{\DUrole{doc}{Accelerator packages}}}} page for more
instructions on how to use the accelerated styles effectively.


\bigskip\hrule\bigskip



\subsection{Output info}
\label{\detokenize{compute_orientorder_atom:output-info}}
\sphinxAtStartPar
This compute calculates a per\sphinxhyphen{}atom array with \sphinxstyleemphasis{nlvalues} columns,
giving the \(Q_\ell\) values for each atom, which are real numbers in the
range \(0 \le Q_\ell \le 1\).

\sphinxAtStartPar
If the keyword \sphinxstyleemphasis{wl} is set to yes, then the \(W_\ell\) values for each
atom will be added to the output array, which are real numbers.

\sphinxAtStartPar
If the keyword \sphinxstyleemphasis{wl/hat} is set to yes, then the \(\hat{W}_\ell\)
values for each atom will be added to the output array, which are real numbers.

\sphinxAtStartPar
If the keyword \sphinxstyleemphasis{components} is set, then the real and imaginary parts
of each component of \sphinxstyleemphasis{normalized} \(\hat{Y}_{\ell m}\) will be added to the
output array in the following order:
\(\Re(\hat{Y}_{-m}),\) \(\Im(\hat{Y}_{-m}),\)
\(\Re(\hat{Y}_{-m+1}),\) \(\Im(\hat{Y}_{-m+1}), \dotsc,\)
\(\Re(\hat{Y}_m),\) \(\Im(\hat{Y}_m).\)

\sphinxAtStartPar
In summary, the per\sphinxhyphen{}atom array will contain \sphinxstyleemphasis{nlvalues} columns, followed by
an additional \sphinxstyleemphasis{nlvalues} columns if \sphinxstyleemphasis{wl} is set to yes, followed by
an additional \sphinxstyleemphasis{nlvalues} columns if \sphinxstyleemphasis{wl/hat} is set to yes, followed
by an additional 2*(2* \sphinxstyleemphasis{ldegree}+1) columns if the \sphinxstyleemphasis{components}
keyword is set.

\sphinxAtStartPar
These values can be accessed by any command that uses per\sphinxhyphen{}atom values
from a compute as input.  See the {\hyperref[\detokenize{Howto_output::doc}]{\sphinxcrossref{\DUrole{doc}{Howto output}}}} doc
page for an overview of LAMMPS output options.


\subsection{Restrictions}
\label{\detokenize{compute_orientorder_atom:restrictions}}
\sphinxAtStartPar
none


\subsection{Related commands}
\label{\detokenize{compute_orientorder_atom:related-commands}}
\sphinxAtStartPar
{\hyperref[\detokenize{compute_coord_atom::doc}]{\sphinxcrossref{\DUrole{doc}{compute coord/atom}}}}, {\hyperref[\detokenize{compute_centro_atom::doc}]{\sphinxcrossref{\DUrole{doc}{compute centro/atom}}}}, {\hyperref[\detokenize{compute_hexorder_atom::doc}]{\sphinxcrossref{\DUrole{doc}{compute hexorder/atom}}}}


\subsection{Default}
\label{\detokenize{compute_orientorder_atom:default}}
\sphinxAtStartPar
The option defaults are \sphinxstyleemphasis{cutoff} = pair style cutoff, \sphinxstyleemphasis{nnn} = 12,
\sphinxstyleemphasis{degrees} = 5 4 6 8 10 12 (i.e., \(Q_4\), \(Q_6\), \(Q_8\), \(Q_{10}\), and \(Q_{12}\)),
\sphinxstyleemphasis{wl} = no, \sphinxstyleemphasis{wl/hat} = no, \sphinxstyleemphasis{components} off, and \sphinxstyleemphasis{chunksize} = 16384


\bigskip\hrule\bigskip

\phantomsection\label{\detokenize{compute_orientorder_atom:steinhardt}}
\sphinxAtStartPar
\sphinxstylestrong{(Steinhardt)} P. Steinhardt, D. Nelson, and M. Ronchetti,
Phys. Rev. B 28, 784 (1983).

\phantomsection\label{\detokenize{compute_orientorder_atom:mickel}}
\sphinxAtStartPar
\sphinxstylestrong{(Mickel)} W. Mickel, S. C. Kapfer, G. E. Schroeder\sphinxhyphen{}Turkand, K. Mecke,
J. Chem. Phys. 138, 044501 (2013).

\phantomsection\label{\detokenize{compute_orientorder_atom:tenwolde2}}
\sphinxAtStartPar
\sphinxstylestrong{(tenWolde)} P. R. ten Wolde, M. J. Ruiz\sphinxhyphen{}Montero, D. Frenkel,
J. Chem. Phys. 104, 9932 (1996).

\sphinxstepscope

\index{compute pace@\spxentry{compute pace}}\ignorespaces 

\section{compute pace command}
\label{\detokenize{compute_pace:compute-pace-command}}\label{\detokenize{compute_pace:index-0}}\label{\detokenize{compute_pace::doc}}

\subsection{Syntax}
\label{\detokenize{compute_pace:syntax}}
\begin{sphinxVerbatim}[commandchars=\\\{\}]
\PYG{k}{compute }\PYG{n+nv+nvIdentifier}{ID}\PYG{+w}{ }\PYG{n+nv+nvIdentifier}{group\PYGZhy{}ID}\PYG{+w}{ }\PYG{n}{pace}\PYG{+w}{ }\PYG{n}{ace\PYGZus{}potential\PYGZus{}filename}\PYG{+w}{ }\PYG{n}{...}\PYG{+w}{ }\PYG{n}{keyword}\PYG{+w}{ }\PYG{n}{values}\PYG{+w}{ }\PYG{n}{...}
\end{sphinxVerbatim}
\begin{itemize}
\item {} 
\sphinxAtStartPar
ID, group\sphinxhyphen{}ID are documented in {\hyperref[\detokenize{compute::doc}]{\sphinxcrossref{\DUrole{doc}{compute}}}} command

\item {} 
\sphinxAtStartPar
pace = style name of this compute command

\item {} 
\sphinxAtStartPar
ace\_potential\_filename = file name (in the .yace or .ace format from {\hyperref[\detokenize{pair_pace::doc}]{\sphinxcrossref{\DUrole{doc}{pace pair\_style}}}}) including ACE hyper\sphinxhyphen{}parameters, bonds, and generalized coupling coefficients

\item {} 
\sphinxAtStartPar
keyword = \sphinxstyleemphasis{bikflag} or \sphinxstyleemphasis{dgradflag}
\begin{sphinxalltt}
\sphinxstyleemphasis{bikflag} value = \sphinxstyleemphasis{0} or \sphinxstyleemphasis{1}
   \sphinxstyleemphasis{0} = descriptors are summed over atoms of each type
   \sphinxstyleemphasis{1} = descriptors are listed separately for each atom
\sphinxstyleemphasis{dgradflag} value = \sphinxstyleemphasis{0} or \sphinxstyleemphasis{1}
   \sphinxstyleemphasis{0} = descriptor gradients are summed over atoms of each type
   \sphinxstyleemphasis{1} = descriptor gradients are listed separately for each atom pair
\end{sphinxalltt}

\end{itemize}


\subsection{Examples}
\label{\detokenize{compute_pace:examples}}
\begin{sphinxVerbatim}[commandchars=\\\{\}]
\PYG{k}{compute }\PYG{n+nv+nvIdentifier}{pace}\PYG{+w}{ }\PYG{n+nv+nvIdentifier}{all}\PYG{+w}{ }\PYG{n}{pace}\PYG{+w}{ }\PYG{n}{coupling\PYGZus{}coefficients.yace}
\PYG{k}{compute }\PYG{n+nv+nvIdentifier}{pace}\PYG{+w}{ }\PYG{n+nv+nvIdentifier}{all}\PYG{+w}{ }\PYG{n}{pace}\PYG{+w}{ }\PYG{n}{coupling\PYGZus{}coefficients.yace}\PYG{+w}{ }\PYG{l+m}{0}\PYG{+w}{ }\PYG{l+m}{1}
\PYG{k}{compute }\PYG{n+nv+nvIdentifier}{pace}\PYG{+w}{ }\PYG{n+nv+nvIdentifier}{all}\PYG{+w}{ }\PYG{n}{pace}\PYG{+w}{ }\PYG{n}{coupling\PYGZus{}coefficients.yace}\PYG{+w}{ }\PYG{l+m}{1}\PYG{+w}{ }\PYG{l+m}{1}
\end{sphinxVerbatim}


\subsection{Description}
\label{\detokenize{compute_pace:description}}
\sphinxAtStartPar
\DUrole{versionmodified,added}{Added in version 7Feb2024.}

\sphinxAtStartPar
This compute calculates a set of quantities related to the atomic
cluster expansion (ACE) descriptors of the atoms in a group.  ACE
descriptors are highly general atomic descriptors, encoding the radial
and angular distribution of neighbor atoms, up to arbitrary bond order
(rank).  The detailed mathematical definition is given in the paper by
{\hyperref[\detokenize{compute_pace:drautz19}]{\sphinxcrossref{\DUrole{std,std-ref}{(Drautz)}}}}.  These descriptors are used in the
{\hyperref[\detokenize{pair_pace::doc}]{\sphinxcrossref{\DUrole{doc}{pace pair\_style}}}}.  Quantities obtained from \sphinxtitleref{compute
pace} are related to those used in {\hyperref[\detokenize{pair_pace::doc}]{\sphinxcrossref{\DUrole{doc}{pace pair\_style}}}} to
evaluate atomic energies, forces, and stresses for linear ACE models.

\sphinxAtStartPar
For example, the energy for a linear ACE model is calculated as:
\(E=\sum_i^{N\_atoms} \sum_{\boldsymbol{\nu}} c_{\boldsymbol{\nu}}
B_{i,\boldsymbol{\boldsymbol{\nu}}}\).  The ACE descriptors for atom \sphinxtitleref{i}
\(B_{i,\boldsymbol{\nu}}\), and \(c_{\nu}\) are linear model
parameters.  The detailed definition and indexing convention for ACE
descriptors is given in {\hyperref[\detokenize{compute_pace:drautz19}]{\sphinxcrossref{\DUrole{std,std-ref}{(Drautz)}}}}.  In short, body
order \(N\), angular character, radial character, and chemical
elements in the \sphinxstyleemphasis{N\sphinxhyphen{}body} descriptor are encoded by \(\nu\).  In the
{\hyperref[\detokenize{pair_pace::doc}]{\sphinxcrossref{\DUrole{doc}{pace pair\_style}}}}, the linear model parameters and the
ACE descriptors are combined for efficient evaluation of energies and
forces.  The details and benefits of this efficient implementation are
given in {\hyperref[\detokenize{compute_pace:lysogorskiy21}]{\sphinxcrossref{\DUrole{std,std-ref}{(Lysogorskiy)}}}}, but the combined
descriptors and linear model parameters for the purposes of \sphinxtitleref{compute
pace} may be expressed in terms of the ACE descriptors mentioned above.

\sphinxAtStartPar
\(c_{\boldsymbol{\nu}} B_{i,\boldsymbol{\nu}}= \sum_{\boldsymbol{\nu}' \in \boldsymbol{\nu} } \big[ c_{\boldsymbol{\nu}} C(\boldsymbol{\nu}') \big] A_{i,\boldsymbol{\nu}'}\)

\sphinxAtStartPar
where the bracketed terms on the right\sphinxhyphen{}hand side are the combined functions
with linear model parameters typically provided in the \sphinxtitleref{\textless{}name\textgreater{}.yace} potential
file for \sphinxtitleref{pace pair\_style}. When these bracketed terms are multiplied by the
products of the atomic base from {\hyperref[\detokenize{compute_pace:drautz19}]{\sphinxcrossref{\DUrole{std,std-ref}{(Drautz)}}}},
\(A_{i,\boldsymbol{\nu'}}\), the ACE descriptors are recovered but they
are also scaled by linear model parameters. The generalized coupling coefficients,
written in short\sphinxhyphen{}hand here as \(C(\boldsymbol{\nu}')\), are the generalized
Clebsch\sphinxhyphen{}Gordan or generalized Wigner symbols. It may be desirable to reverse the
combination of these descriptors and the linear model parameters so that the
ACE descriptors themselves may be used. The ACE descriptors and their gradients
are often used when training ACE models, performing custom data analysis,
generalizing ACE model forms, and other tasks that involve direct computation of
descriptors. The key utility of \sphinxtitleref{compute pace} is that it can compute the ACE
descriptors and gradients so that these tasks can be performed during a LAMMPS
simulation or so that LAMMPS can be used as a driver for tasks like ACE model
parameterization. To see how this command can be used within a Python workflow
to train ACE potentials, see the examples in
\sphinxhref{https://github.com/FitSNAP/FitSNAP}{FitSNAP}. Examples on using outputs from
this compute to construct general ACE potential forms are demonstrated in
{\hyperref[\detokenize{compute_pace:goff23}]{\sphinxcrossref{\DUrole{std,std-ref}{(Goff)}}}}. The various keywords and inputs to \sphinxtitleref{compute pace}
determine what ACE descriptors and related quantities are returned in a compute
array.

\sphinxAtStartPar
The coefficient file, \sphinxtitleref{\textless{}name\textgreater{}.yace}, ultimately defines the number of ACE
descriptors to be computed, their maximum body\sphinxhyphen{}order, the degree of angular
character they have, the degree of radial character they have, the chemical
character (which element\sphinxhyphen{}element interactions are encoded by descriptors),
and other hyper\sphinxhyphen{}parameters defined in {\hyperref[\detokenize{compute_pace:drautz19}]{\sphinxcrossref{\DUrole{std,std-ref}{(Drautz)}}}}. These may
be modeled after the potential files in {\hyperref[\detokenize{pair_pace::doc}]{\sphinxcrossref{\DUrole{doc}{pace pair\_style}}}},
and have the same format. Details on how to generate the coefficient files
to train ACE models may be found in \sphinxhref{https://github.com/FitSNAP/FitSNAP}{FitSNAP}.

\sphinxAtStartPar
The keyword \sphinxstyleemphasis{bikflag} determines whether or not to list the descriptors of
each atom separately, or sum them together and list in a single row. If
\sphinxstyleemphasis{bikflag} is set to \sphinxstyleemphasis{0} then a single descriptor row is used, which contains
the per\sphinxhyphen{}atom ACE descriptors \(B_{i,\boldsymbol{\nu}}\) summed over all
atoms \sphinxstyleemphasis{i} to produce \(B_{\boldsymbol{\nu}}\). If \sphinxstyleemphasis{bikflag} is set to
\sphinxstyleemphasis{1} this is replaced by a separate per\sphinxhyphen{}atom ACE descriptor row for each atom.
In this case, the entries in the final column for these rows are set to zero.

\sphinxAtStartPar
The keyword \sphinxstyleemphasis{dgradflag} determines whether to sum atom gradients or list
them separately. If \sphinxstyleemphasis{dgradflag} is set to 0, the ACE
descriptor gradients w.r.t. atom \sphinxstyleemphasis{j} are summed over all atoms \sphinxstyleemphasis{i’}
of, which may be useful when training linear ACE models on atomic forces.
If \sphinxstyleemphasis{dgradflag} is set to 1, gradients are listed separately for each pair of atoms.
Each row corresponds
to a single term \(\frac{\partial {B_{i,\boldsymbol{\nu}}}}{\partial {r}^a_j}\)
where \({r}^a_j\) is the \sphinxstyleemphasis{a\sphinxhyphen{}th} position coordinate of the atom with global
index \sphinxstyleemphasis{j}. This also changes the number of columns to be equal to the number of
ACE descriptors, with 3 additional columns representing the indices \(i\),
\(j\), and \(a\), as explained more in the Output info section below.
The option \sphinxstyleemphasis{dgradflag=1} requires that \sphinxstyleemphasis{bikflag=1}.

\begin{sphinxadmonition}{note}{Note:}
\sphinxAtStartPar
It is noted here that in contrast to {\hyperref[\detokenize{pair_pace::doc}]{\sphinxcrossref{\DUrole{doc}{pace pair\_style}}}},
the \sphinxstyleemphasis{.yace} file for \sphinxtitleref{compute pace} typically should not contain linear
parameters for an ACE potential. If \(c_{\nu}\) are included,
the value of the descriptor will not be returned in the \sphinxtitleref{compute} array,
but instead, the energy contribution from that descriptor will be returned.
Do not do this unless it is the desired behavior.
\sphinxstyleemphasis{In short, you should not plug in a ‘.yace’ for a pace potential into this
compute to evaluate descriptors.}
\end{sphinxadmonition}

\begin{sphinxadmonition}{note}{Note:}
\sphinxAtStartPar
\sphinxstyleemphasis{Generalized Clebsch\sphinxhyphen{}Gordan or Generalized Wigner symbols (with appropriate
factors) must be used to evaluate ACE descriptors with this compute.} There
are multiple ways to define the generalized coupling coefficients. Because
of this, this compute will not revert your potential file to a coupling
coefficient file. Instead this compute allows the user to supply coupling
coefficients that follow any convention.
\end{sphinxadmonition}

\begin{sphinxadmonition}{note}{Note:}
\sphinxAtStartPar
Using \sphinxstyleemphasis{dgradflag} = 1 produces a global array with \(N + 3N^2 + 1\) rows
which becomes expensive for systems with more than 1000 atoms.
\end{sphinxadmonition}

\begin{sphinxadmonition}{note}{Note:}
\sphinxAtStartPar
If you have a bonded system, then the settings of {\hyperref[\detokenize{special_bonds::doc}]{\sphinxcrossref{\DUrole{doc}{special\_bonds}}}} command can remove pairwise interactions between
atoms in the same bond, angle, or dihedral.  This is the default
setting for the {\hyperref[\detokenize{special_bonds::doc}]{\sphinxcrossref{\DUrole{doc}{special\_bonds}}}} command, and
means those pairwise interactions do not appear in the neighbor list.
Because this fix uses the neighbor list, it also means those pairs
will not be included in the calculation.  One way to get around this,
is to write a dump file, and use the {\hyperref[\detokenize{rerun::doc}]{\sphinxcrossref{\DUrole{doc}{rerun}}}} command to
compute the ACE descriptors for snapshots in the dump file.
The rerun script can use a {\hyperref[\detokenize{special_bonds::doc}]{\sphinxcrossref{\DUrole{doc}{special\_bonds}}}}
command that includes all pairs in the neighbor list.
\end{sphinxadmonition}


\bigskip\hrule\bigskip



\subsection{Output info}
\label{\detokenize{compute_pace:output-info}}
\sphinxAtStartPar
Compute \sphinxstyleemphasis{pace} evaluates a global array.  The columns are arranged into
\sphinxstyleemphasis{ntypes} blocks, listed in order of atom type \sphinxstyleemphasis{I}. Each block contains
one column for each ACE descriptor, the same as for compute
\sphinxstyleemphasis{sna/atom}in {\hyperref[\detokenize{compute_sna_atom::doc}]{\sphinxcrossref{\DUrole{doc}{compute snap}}}}. A final column contains the corresponding energy, force
component on an atom, or virial stress component. The rows of the array
appear in the following order:
\begin{itemize}
\item {} 
\sphinxAtStartPar
1 row: \sphinxstyleemphasis{pace} average descriptor values for all atoms of type \sphinxstyleemphasis{I}

\item {} 
\sphinxAtStartPar
3*\sphinxstyleemphasis{n} force rows: quantities, with derivatives w.r.t. x, y, and z coordinate of atom \sphinxstyleemphasis{i} appearing in consecutive rows. The atoms are sorted based on atom ID and run up to the total number of atoms, \sphinxstyleemphasis{n}.

\item {} 
\sphinxAtStartPar
6 rows: \sphinxstyleemphasis{virial} quantities summed for all atoms of type \sphinxstyleemphasis{I}

\end{itemize}

\sphinxAtStartPar
For example, if \(\# \; B_{i, \boldsymbol{\nu}}\) =30 and ntypes=1, the number of columns in the
The number of columns in the global array generated by \sphinxstyleemphasis{pace} are 31, and
931, respectively, while the number of rows is 1+3*\sphinxstyleemphasis{n}+6, where \sphinxstyleemphasis{n}
is the total number of atoms.

\sphinxAtStartPar
If the \sphinxstyleemphasis{bik} keyword is set to 1, the structure of the pace array is expanded.
The first \(N\) rows of the pace array
correspond to \(\# \; B_{i,\boldsymbol{\nu}}\) instead of a single row summed over atoms \(i\).
In this case, the entries in the final column for these rows
are set to zero. Also, each row contains only non\sphinxhyphen{}zero entries for the
columns corresponding to the type of that atom. This is not true in the case
of \sphinxstyleemphasis{dgradflag} keyword = 1 (see below).

\sphinxAtStartPar
If the \sphinxstyleemphasis{dgradflag} keyword is set to 1, this changes the structure of the
global array completely.
Here the per\sphinxhyphen{}atom quantities are replaced with rows corresponding to
descriptor gradient components on single atoms:
\begin{equation*}
\begin{split}\frac{\partial {B_{i,\boldsymbol{\nu}}  }}{\partial {r}^a_j}\end{split}
\end{equation*}
\sphinxAtStartPar
where \({r}^a_j\) is the \sphinxstyleemphasis{a\sphinxhyphen{}th} position coordinate of the atom with global
index \sphinxstyleemphasis{j}. The rows are
organized in chunks, where each chunk corresponds to an atom with global index
\(j\). The rows in an atom \(j\) chunk correspond to
atoms with global index \(i\). The total number of rows for
these descriptor gradients is therefore \(3N^2\).
The number of columns is equal to the number of ACE descriptors,
plus 3 additional left\sphinxhyphen{}most columns representing the global atom indices
\(i\), \(j\),
and Cartesian direction \(a\)  (0, 1, 2, for x, y, z).
The first 3 columns of the first \(N\) rows belong to the reference
potential force components. The remaining K columns contain the
\(B_{i,\boldsymbol{\nu}}\) per\sphinxhyphen{}atom descriptors corresponding to the non\sphinxhyphen{}zero entries
obtained when \sphinxstyleemphasis{bikflag} = 1.
The first column of the last row, after the first
\(N + 3N^2\) rows, contains the reference potential
energy. The virial components are not used with this option. The total number of
rows is therefore \(N + 3N^2 + 1\) and the number of columns is \(K + 3\).

\sphinxAtStartPar
These values can be accessed by any command that uses global values
from a compute as input.  See the {\hyperref[\detokenize{Howto_output::doc}]{\sphinxcrossref{\DUrole{doc}{Howto output}}}} doc
page for an overview of LAMMPS output options.


\subsection{Restrictions}
\label{\detokenize{compute_pace:restrictions}}
\sphinxAtStartPar
These computes are part of the ML\sphinxhyphen{}PACE package.  They are only enabled
if LAMMPS was built with that package.  See the {\hyperref[\detokenize{Build_package::doc}]{\sphinxcrossref{\DUrole{doc}{Build package}}}} page for more info.


\subsection{Related commands}
\label{\detokenize{compute_pace:related-commands}}
\sphinxAtStartPar
{\hyperref[\detokenize{pair_pace::doc}]{\sphinxcrossref{\DUrole{doc}{pair\_style pace}}}}
{\hyperref[\detokenize{pair_snap::doc}]{\sphinxcrossref{\DUrole{doc}{pair\_style snap}}}}
{\hyperref[\detokenize{compute_sna_atom::doc}]{\sphinxcrossref{\DUrole{doc}{compute snap}}}}


\subsection{Default}
\label{\detokenize{compute_pace:default}}
\sphinxAtStartPar
The optional keyword defaults are \sphinxstyleemphasis{bikflag} = 0,
\sphinxstyleemphasis{dgradflag} = 0


\bigskip\hrule\bigskip

\phantomsection\label{\detokenize{compute_pace:drautz19}}
\sphinxAtStartPar
\sphinxstylestrong{(Drautz)} Drautz, Phys Rev B, 99, 014104 (2019).

\phantomsection\label{\detokenize{compute_pace:lysogorskiy21}}
\sphinxAtStartPar
\sphinxstylestrong{(Lysogorskiy)} Lysogorskiy, van der Oord, Bochkarev, Menon, Rinaldi, Hammerschmidt, Mrovec, Thompson, Csanyi, Ortner, Drautz, npj Comp Mat, 7, 97 (2021).

\phantomsection\label{\detokenize{compute_pace:goff23}}
\sphinxAtStartPar
\sphinxstylestrong{(Goff)} Goff, Zhang, Negre, Rohskopf, Niklasson, Journal of Chemical Theory and Computation 19, no. 13 (2023).

\sphinxstepscope

\index{compute pair@\spxentry{compute pair}}\ignorespaces 

\section{compute pair command}
\label{\detokenize{compute_pair:compute-pair-command}}\label{\detokenize{compute_pair:index-0}}\label{\detokenize{compute_pair::doc}}

\subsection{Syntax}
\label{\detokenize{compute_pair:syntax}}
\begin{sphinxVerbatim}[commandchars=\\\{\}]
\PYG{k}{compute }\PYG{n+nv+nvIdentifier}{ID}\PYG{+w}{ }\PYG{n+nv+nvIdentifier}{group\PYGZhy{}ID}\PYG{+w}{ }\PYG{n}{pair}\PYG{+w}{ }\PYG{n}{pstyle}\PYG{+w}{ }\PYG{n}{[nstyle]}\PYG{+w}{ }\PYG{n}{[evalue]}
\end{sphinxVerbatim}
\begin{itemize}
\item {} 
\sphinxAtStartPar
ID, group\sphinxhyphen{}ID are documented in {\hyperref[\detokenize{compute::doc}]{\sphinxcrossref{\DUrole{doc}{compute}}}} command

\item {} 
\sphinxAtStartPar
pair = style name of this compute command

\item {} 
\sphinxAtStartPar
pstyle = style name of a pair style that calculates additional values

\item {} 
\sphinxAtStartPar
nsub = \sphinxstyleemphasis{n}\sphinxhyphen{}instance of a sub\sphinxhyphen{}style, if a pair style is used multiple times in a hybrid style

\item {} 
\sphinxAtStartPar
\sphinxstyleemphasis{evalue} = \sphinxstyleemphasis{epair} or \sphinxstyleemphasis{evdwl} or \sphinxstyleemphasis{ecoul} or blank (optional)

\end{itemize}


\subsection{Examples}
\label{\detokenize{compute_pair:examples}}
\begin{sphinxVerbatim}[commandchars=\\\{\}]
\PYG{k}{compute }\PYG{n+nv+nvIdentifier}{1}\PYG{+w}{ }\PYG{n+nv+nvIdentifier}{all}\PYG{+w}{ }\PYG{n}{pair}\PYG{+w}{ }\PYG{n}{gauss}
\PYG{k}{compute }\PYG{n+nv+nvIdentifier}{1}\PYG{+w}{ }\PYG{n+nv+nvIdentifier}{all}\PYG{+w}{ }\PYG{n}{pair}\PYG{+w}{ }\PYG{n}{lj}\PYG{o}{/}\PYG{n}{cut}\PYG{o}{/}\PYG{n}{coul}\PYG{o}{/}\PYG{n}{cut}\PYG{+w}{ }\PYG{n}{ecoul}
\PYG{k}{compute }\PYG{n+nv+nvIdentifier}{1}\PYG{+w}{ }\PYG{n+nv+nvIdentifier}{all}\PYG{+w}{ }\PYG{n}{pair}\PYG{+w}{ }\PYG{n}{tersoff}\PYG{+w}{ }\PYG{l+m}{2}\PYG{+w}{ }\PYG{n}{epair}
\PYG{k}{compute }\PYG{n+nv+nvIdentifier}{1}\PYG{+w}{ }\PYG{n+nv+nvIdentifier}{all}\PYG{+w}{ }\PYG{n}{pair}\PYG{+w}{ }\PYG{n}{reaxff}
\end{sphinxVerbatim}


\subsection{Description}
\label{\detokenize{compute_pair:description}}
\sphinxAtStartPar
Define a computation that extracts additional values calculated by a
pair style, and makes them accessible for output or further processing
by other commands.

\begin{sphinxadmonition}{note}{Note:}
\sphinxAtStartPar
The group specified for this command is \sphinxstylestrong{ignored}.
\end{sphinxadmonition}

\sphinxAtStartPar
The specified \sphinxstyleemphasis{pstyle} must be a pair style used in your simulation
either by itself or as a sub\sphinxhyphen{}style in a {\hyperref[\detokenize{pair_hybrid::doc}]{\sphinxcrossref{\DUrole{doc}{pair\_style hybrid or hybrid/overlay}}}} command. If the sub\sphinxhyphen{}style is
used more than once, an additional number \sphinxstyleemphasis{nsub} has to be specified
in order to choose which instance of the sub\sphinxhyphen{}style will be used by
the compute. Not specifying the number in this case will cause the
compute to fail.

\sphinxAtStartPar
The \sphinxstyleemphasis{evalue} setting is optional.  All
pair styles tally a potential energy \sphinxstyleemphasis{epair} which may be broken into
two parts: \sphinxstyleemphasis{evdwl} and \sphinxstyleemphasis{ecoul} such that \sphinxstyleemphasis{epair} = \sphinxstyleemphasis{evdwl} + \sphinxstyleemphasis{ecoul}.
If the pair style calculates Coulombic interactions, their energy will
be tallied in \sphinxstyleemphasis{ecoul}.  Everything else (whether it is a Lennard\sphinxhyphen{}Jones
style van der Waals interaction or not) is tallied in \sphinxstyleemphasis{evdwl}.  If
\sphinxstyleemphasis{evalue} is blank or specified as \sphinxstyleemphasis{epair}, then \sphinxstyleemphasis{epair} is stored
as a global scalar by this compute.  This is useful when using
{\hyperref[\detokenize{pair_hybrid::doc}]{\sphinxcrossref{\DUrole{doc}{pair\_style hybrid}}}} if you want to know the portion
of the total energy contributed by one sub\sphinxhyphen{}style.  If \sphinxstyleemphasis{evalue} is
specified as \sphinxstyleemphasis{evdwl} or \sphinxstyleemphasis{ecoul}, then just that portion of the energy
is stored as a global scalar.

\begin{sphinxadmonition}{note}{Note:}
\sphinxAtStartPar
The energy returned by the \sphinxstyleemphasis{evdwl} keyword does not include tail
corrections, even if they are enabled via the
{\hyperref[\detokenize{pair_modify::doc}]{\sphinxcrossref{\DUrole{doc}{pair\_modify}}}} command.
\end{sphinxadmonition}

\sphinxAtStartPar
Some pair styles tally additional quantities, e.g. a breakdown of
potential energy into 14 components is tallied by the
{\hyperref[\detokenize{pair_reaxff::doc}]{\sphinxcrossref{\DUrole{doc}{pair\_style reaxff}}}} command.  These values (1 or more)
are stored as a global vector by this compute.  See the page for
{\hyperref[\detokenize{pair_style::doc}]{\sphinxcrossref{\DUrole{doc}{individual pair styles}}}} for info on these values.


\subsection{Output info}
\label{\detokenize{compute_pair:output-info}}
\sphinxAtStartPar
This compute calculates a global scalar which is \sphinxstyleemphasis{epair} or \sphinxstyleemphasis{evdwl} or
\sphinxstyleemphasis{ecoul}.  If the pair style supports it, it also calculates a global
vector of length \(\ge\) 1, as determined by the pair style.  These values
can be used by any command that uses global scalar or vector values
from a compute as input.  See the {\hyperref[\detokenize{Howto_output::doc}]{\sphinxcrossref{\DUrole{doc}{Howto output}}}} doc
page for an overview of LAMMPS output options.

\sphinxAtStartPar
The scalar and vector values calculated by this compute are
“extensive”.

\sphinxAtStartPar
The scalar value will be in energy {\hyperref[\detokenize{units::doc}]{\sphinxcrossref{\DUrole{doc}{units}}}}.  The vector
values will typically also be in energy {\hyperref[\detokenize{units::doc}]{\sphinxcrossref{\DUrole{doc}{units}}}}, but see
the page for the pair style for details.


\subsection{Restrictions}
\label{\detokenize{compute_pair:restrictions}}\begin{quote}

\sphinxAtStartPar
none
\end{quote}


\subsection{Related commands}
\label{\detokenize{compute_pair:related-commands}}
\sphinxAtStartPar
{\hyperref[\detokenize{compute_pe::doc}]{\sphinxcrossref{\DUrole{doc}{compute pe}}}}, {\hyperref[\detokenize{compute_bond::doc}]{\sphinxcrossref{\DUrole{doc}{compute bond}}}},
{\hyperref[\detokenize{fix_pair::doc}]{\sphinxcrossref{\DUrole{doc}{fix pair}}}}


\subsection{Default}
\label{\detokenize{compute_pair:default}}
\sphinxAtStartPar
The keyword defaults are \sphinxstyleemphasis{evalue} = \sphinxstyleemphasis{epair}, nsub = 0.

\sphinxstepscope

\index{compute pair/local@\spxentry{compute pair/local}}\ignorespaces 

\section{compute pair/local command}
\label{\detokenize{compute_pair_local:compute-pair-local-command}}\label{\detokenize{compute_pair_local:index-0}}\label{\detokenize{compute_pair_local::doc}}

\subsection{Syntax}
\label{\detokenize{compute_pair_local:syntax}}
\begin{sphinxVerbatim}[commandchars=\\\{\}]
\PYG{k}{compute }\PYG{n+nv+nvIdentifier}{ID}\PYG{+w}{ }\PYG{n+nv+nvIdentifier}{group\PYGZhy{}ID}\PYG{+w}{ }\PYG{n}{pair}\PYG{o}{/}\PYG{n}{local}\PYG{+w}{ }\PYG{n}{value1}\PYG{+w}{ }\PYG{n}{value2}\PYG{+w}{ }\PYG{n}{...}\PYG{+w}{ }\PYG{n}{keyword}\PYG{+w}{ }\PYG{n}{args}\PYG{+w}{ }\PYG{n}{...}
\end{sphinxVerbatim}
\begin{itemize}
\item {} 
\sphinxAtStartPar
ID, group\sphinxhyphen{}ID are documented in {\hyperref[\detokenize{compute::doc}]{\sphinxcrossref{\DUrole{doc}{compute}}}} command

\item {} 
\sphinxAtStartPar
pair/local = style name of this compute command

\item {} 
\sphinxAtStartPar
one or more values may be appended

\item {} 
\sphinxAtStartPar
value = \sphinxstyleemphasis{dist} or \sphinxstyleemphasis{dx} or \sphinxstyleemphasis{dy} or \sphinxstyleemphasis{dz} or \sphinxstyleemphasis{eng} or \sphinxstyleemphasis{force} or \sphinxstyleemphasis{fx} or \sphinxstyleemphasis{fy} or \sphinxstyleemphasis{fz} or \sphinxstyleemphasis{p1} or \sphinxstyleemphasis{p2} or …
\begin{sphinxalltt}
\sphinxstyleemphasis{dist} = pairwise distance
\sphinxstyleemphasis{dx},\sphinxstyleemphasis{dy},\sphinxstyleemphasis{dz} = components of pairwise distance
\sphinxstyleemphasis{eng} = pairwise energy
\sphinxstyleemphasis{force} = pairwise force
\sphinxstyleemphasis{fx},\sphinxstyleemphasis{fy},\sphinxstyleemphasis{fz} = components of pairwise force
\sphinxstyleemphasis{p1}, \sphinxstyleemphasis{p2}, ... = pair style specific quantities for allowed N values
\end{sphinxalltt}

\item {} 
\sphinxAtStartPar
zero or more keyword/arg pairs may be appended

\item {} 
\sphinxAtStartPar
keyword = \sphinxstyleemphasis{cutoff}
\begin{sphinxalltt}
\sphinxstyleemphasis{cutoff} arg = \sphinxstyleemphasis{type} or \sphinxstyleemphasis{radius}
\end{sphinxalltt}

\end{itemize}


\subsection{Examples}
\label{\detokenize{compute_pair_local:examples}}
\begin{sphinxVerbatim}[commandchars=\\\{\}]
\PYG{k}{compute }\PYG{n+nv+nvIdentifier}{1}\PYG{+w}{ }\PYG{n+nv+nvIdentifier}{all}\PYG{+w}{ }\PYG{n}{pair}\PYG{o}{/}\PYG{n}{local}\PYG{+w}{ }\PYG{n}{eng}
\PYG{k}{compute }\PYG{n+nv+nvIdentifier}{1}\PYG{+w}{ }\PYG{n+nv+nvIdentifier}{all}\PYG{+w}{ }\PYG{n}{pair}\PYG{o}{/}\PYG{n}{local}\PYG{+w}{ }\PYG{n}{dist}\PYG{+w}{ }\PYG{n}{eng}\PYG{+w}{ }\PYG{n}{force}
\PYG{k}{compute }\PYG{n+nv+nvIdentifier}{1}\PYG{+w}{ }\PYG{n+nv+nvIdentifier}{all}\PYG{+w}{ }\PYG{n}{pair}\PYG{o}{/}\PYG{n}{local}\PYG{+w}{ }\PYG{n}{dist}\PYG{+w}{ }\PYG{n}{eng}\PYG{+w}{ }\PYG{n}{fx}\PYG{+w}{ }\PYG{n}{fy}\PYG{+w}{ }\PYG{n}{fz}
\PYG{k}{compute }\PYG{n+nv+nvIdentifier}{1}\PYG{+w}{ }\PYG{n+nv+nvIdentifier}{all}\PYG{+w}{ }\PYG{n}{pair}\PYG{o}{/}\PYG{n}{local}\PYG{+w}{ }\PYG{n}{dist}\PYG{+w}{ }\PYG{n}{fx}\PYG{+w}{ }\PYG{n}{fy}\PYG{+w}{ }\PYG{n}{fz}\PYG{+w}{ }\PYG{n}{p1}\PYG{+w}{ }\PYG{n}{p2}\PYG{+w}{ }\PYG{n}{p3}
\end{sphinxVerbatim}


\subsection{Description}
\label{\detokenize{compute_pair_local:description}}
\sphinxAtStartPar
Define a computation that calculates properties of individual pairwise
interactions.  The number of datums generated, aggregated across all
processors, equals the number of pairwise interactions in the system.

\sphinxAtStartPar
The local data stored by this command is generated by looping over the
pairwise neighbor list.  Info about an individual pairwise interaction
will only be included if both atoms in the pair are in the specified
compute group, and if the current pairwise distance is less than the
force cutoff distance for that interaction, as defined by the
{\hyperref[\detokenize{pair_style::doc}]{\sphinxcrossref{\DUrole{doc}{pair\_style}}}} and {\hyperref[\detokenize{pair_coeff::doc}]{\sphinxcrossref{\DUrole{doc}{pair\_coeff}}}}
commands.

\sphinxAtStartPar
The value \sphinxstyleemphasis{dist} is the distance between the pair of atoms.
The values \sphinxstyleemphasis{dx}, \sphinxstyleemphasis{dy}, and \sphinxstyleemphasis{dz} are the \((x,y,z)\) components of the
\sphinxstyleemphasis{distance} between the pair of atoms. This value is always the
distance from the atom of higher to the one with the lower atom ID.

\sphinxAtStartPar
The value \sphinxstyleemphasis{eng} is the interaction energy for the pair of atoms.

\sphinxAtStartPar
The value \sphinxstyleemphasis{force} is the force acting between the pair of atoms, which
is positive for a repulsive force and negative for an attractive
force.  The values \sphinxstyleemphasis{fx}, \sphinxstyleemphasis{fy}, and \sphinxstyleemphasis{fz} are the \((x,y,z)\) components of
\sphinxstyleemphasis{force} on atom I. For pair styles that apply non\sphinxhyphen{}central forces,
such as {\hyperref[\detokenize{pair_gran::doc}]{\sphinxcrossref{\DUrole{doc}{granular pair styles}}}}, these values only include
the \((x,y,z)\) components of the normal force component.

\sphinxAtStartPar
A pair style may define additional pairwise quantities which can be
accessed as \sphinxstyleemphasis{p1} to \sphinxstyleemphasis{pN}, where \(N\) is defined by the pair style.
Most pair styles do not define any additional quantities, so \(N = 0\).
An example of ones that do are the {\hyperref[\detokenize{pair_gran::doc}]{\sphinxcrossref{\DUrole{doc}{granular pair styles}}}}
which calculate the tangential force between two particles and return
its components and magnitude acting on atom \(I\) for
\(N \in \{1,2,3,4\}\).  See individual pair styles for details.

\sphinxAtStartPar
When using \sphinxstyleemphasis{pN} with pair style \sphinxstyleemphasis{hybrid}, the output will be the Nth
quantity from the sub\sphinxhyphen{}style that computes the pairwise interaction
(based on atom types).  If that sub\sphinxhyphen{}style does not define a \sphinxstyleemphasis{pN},
the output will be 0.0.  The maximum allowed \(N\) is the maximum number
of quantities provided by any sub\sphinxhyphen{}style.

\sphinxAtStartPar
When using \sphinxstyleemphasis{pN} with pair style \sphinxstyleemphasis{hybrid/overlay} the quantities
from all sub\sphinxhyphen{}styles that provide them are concatenated together
into one long list. For example, if there are 3 sub\sphinxhyphen{}styles and
2 of them have additional output (with 3 and 4 quantities,
respectively), then 7 values (\sphinxstyleemphasis{p1} up to \sphinxstyleemphasis{p7}) are defined.
The values \sphinxstyleemphasis{p1} to \sphinxstyleemphasis{p3} refer to quantities defined by the first
of the two sub\sphinxhyphen{}styles.  Values \sphinxstyleemphasis{p4} to \sphinxstyleemphasis{p7} refer to quantities
from the second of the two sub\sphinxhyphen{}styles.  If the referenced \sphinxstyleemphasis{pN}
is not computed for the specific pairwise interaction (based on
atom types), then the output will be 0.0.

\sphinxAtStartPar
The value \sphinxstyleemphasis{dist}, \sphinxstyleemphasis{dx}, \sphinxstyleemphasis{dy} and \sphinxstyleemphasis{dz} will be in distance {\hyperref[\detokenize{units::doc}]{\sphinxcrossref{\DUrole{doc}{units}}}}.
The value \sphinxstyleemphasis{eng} will be in energy {\hyperref[\detokenize{units::doc}]{\sphinxcrossref{\DUrole{doc}{units}}}}.
The values \sphinxstyleemphasis{force}, \sphinxstyleemphasis{fx}, \sphinxstyleemphasis{fy}, and \sphinxstyleemphasis{fz} will be in force {\hyperref[\detokenize{units::doc}]{\sphinxcrossref{\DUrole{doc}{units}}}}.
The values \sphinxstyleemphasis{pN} will be in whatever units the pair style defines.

\sphinxAtStartPar
The optional \sphinxstyleemphasis{cutoff} keyword determines how the force cutoff distance
for an interaction is determined.  For the default setting of \sphinxstyleemphasis{type},
the pairwise cutoff defined by the {\hyperref[\detokenize{pair_style::doc}]{\sphinxcrossref{\DUrole{doc}{pair\_style}}}}
command for the types of the two atoms is used.  For the \sphinxstyleemphasis{radius}
setting, the sum of the radii of the two particles is used as a
cutoff.  For example, this is appropriate for granular particles which
only interact when they are overlapping, as computed by
{\hyperref[\detokenize{pair_gran::doc}]{\sphinxcrossref{\DUrole{doc}{granular pair styles}}}}.
Note that if a granular model defines atom
types such that all particles of a specific type are monodisperse
(same diameter), then the two settings are effectively identical.

\sphinxAtStartPar
Note that as atoms migrate from processor to processor, there will be
no consistent ordering of the entries within the local vector or array
from one timestep to the next.  The only consistency that is
guaranteed is that the ordering on a particular timestep will be the
same for local vectors or arrays generated by other compute commands.
For example, pair output from the
{\hyperref[\detokenize{compute_property_local::doc}]{\sphinxcrossref{\DUrole{doc}{compute property/local}}}} command can be combined
with data from this command and output by the {\hyperref[\detokenize{dump::doc}]{\sphinxcrossref{\DUrole{doc}{dump local}}}}
command in a consistent way.

\sphinxAtStartPar
Here is an example of how to do this:

\begin{sphinxVerbatim}[commandchars=\\\{\}]
\PYG{k}{compute }\PYG{n+nv+nvIdentifier}{1}\PYG{+w}{ }\PYG{n+nv+nvIdentifier}{all}\PYG{+w}{ }\PYG{n}{property}\PYG{o}{/}\PYG{n}{local}\PYG{+w}{ }\PYG{n}{patom1}\PYG{+w}{ }\PYG{n}{patom2}
\PYG{k}{compute }\PYG{n+nv+nvIdentifier}{2}\PYG{+w}{ }\PYG{n+nv+nvIdentifier}{all}\PYG{+w}{ }\PYG{n}{pair}\PYG{o}{/}\PYG{n}{local}\PYG{+w}{ }\PYG{n}{dist}\PYG{+w}{ }\PYG{n}{eng}\PYG{+w}{ }\PYG{n}{force}
\PYG{k}{dump }\PYG{n+nv+nvIdentifier}{1}\PYG{+w}{ }\PYG{n+nv+nvIdentifier}{all}\PYG{+w}{ }\PYG{n}{local}\PYG{+w}{ }\PYG{l+m}{1000}\PYG{+w}{ }\PYG{n}{tmp.dump}\PYG{+w}{ }\PYG{n}{index}\PYG{+w}{ }\PYG{n}{c\PYGZus{}1[1]}\PYG{+w}{ }\PYG{n}{c\PYGZus{}1[2]}\PYG{+w}{ }\PYG{n}{c\PYGZus{}2[1]}\PYG{+w}{ }\PYG{n}{c\PYGZus{}2[2]}\PYG{+w}{ }\PYG{n}{c\PYGZus{}2[3]}
\end{sphinxVerbatim}

\begin{sphinxadmonition}{note}{Note:}
\sphinxAtStartPar
For pairs, if two atoms I,J are involved in 1\textendash{}2, 1\textendash{}3, and 1\textendash{}4
interactions within the molecular topology, their pairwise interaction
may be turned off, and thus they may not appear in the neighbor list,
and will not be part of the local data created by this command.  More
specifically, this will be true of I,J pairs with a weighting factor
of 0.0; pairs with a non\sphinxhyphen{}zero weighting factor are included.  The
weighting factors for 1\textendash{}2, 1\textendash{}3, and 1\textendash{}4 pairwise interactions are set
by the {\hyperref[\detokenize{special_bonds::doc}]{\sphinxcrossref{\DUrole{doc}{special\_bonds}}}} command.  An exception is if
long\sphinxhyphen{}range Coulombics are being computed via the
{\hyperref[\detokenize{kspace_style::doc}]{\sphinxcrossref{\DUrole{doc}{kspace\_style}}}} command, then atom pairs with
weighting factors of zero are still included in the neighbor list, so
that a portion of the long\sphinxhyphen{}range interaction contribution can be
computed in the pair style.  Hence in that case, those atom pairs will
be part of the local data created by this command.
\end{sphinxadmonition}


\subsection{Output info}
\label{\detokenize{compute_pair_local:output-info}}
\sphinxAtStartPar
This compute calculates a local vector or local array depending on the
number of keywords.  The length of the vector or number of rows in the
array is the number of pairs.  If a single keyword is specified, a
local vector is produced.  If two or more keywords are specified, a
local array is produced where the number of columns = the number of
keywords.  The vector or array can be accessed by any command that
uses local values from a compute as input.  See the {\hyperref[\detokenize{Howto_output::doc}]{\sphinxcrossref{\DUrole{doc}{Howto output}}}} page for an overview of LAMMPS output
options.

\sphinxAtStartPar
The output for \sphinxstyleemphasis{dist} will be in distance {\hyperref[\detokenize{units::doc}]{\sphinxcrossref{\DUrole{doc}{units}}}}.  The
output for \sphinxstyleemphasis{eng} will be in energy {\hyperref[\detokenize{units::doc}]{\sphinxcrossref{\DUrole{doc}{units}}}}.  The output for
\sphinxstyleemphasis{force}, \sphinxstyleemphasis{fx}, \sphinxstyleemphasis{fy}, and \sphinxstyleemphasis{fz} will be in force {\hyperref[\detokenize{units::doc}]{\sphinxcrossref{\DUrole{doc}{units}}}}.
The output for \sphinxstyleemphasis{pN} will be in whatever units the pair style defines.


\subsection{Restrictions}
\label{\detokenize{compute_pair_local:restrictions}}\begin{quote}

\sphinxAtStartPar
none
\end{quote}


\subsection{Related commands}
\label{\detokenize{compute_pair_local:related-commands}}
\sphinxAtStartPar
{\hyperref[\detokenize{dump::doc}]{\sphinxcrossref{\DUrole{doc}{dump local}}}}, {\hyperref[\detokenize{compute_property_local::doc}]{\sphinxcrossref{\DUrole{doc}{compute property/local}}}}


\subsection{Default}
\label{\detokenize{compute_pair_local:default}}
\sphinxAtStartPar
The keyword default is cutoff = type.

\sphinxstepscope

\index{compute pe@\spxentry{compute pe}}\ignorespaces 

\section{compute pe command}
\label{\detokenize{compute_pe:compute-pe-command}}\label{\detokenize{compute_pe:index-0}}\label{\detokenize{compute_pe::doc}}

\subsection{Syntax}
\label{\detokenize{compute_pe:syntax}}
\begin{sphinxVerbatim}[commandchars=\\\{\}]
\PYG{k}{compute }\PYG{n+nv+nvIdentifier}{ID}\PYG{+w}{ }\PYG{n+nv+nvIdentifier}{group\PYGZhy{}ID}\PYG{+w}{ }\PYG{n}{pe}\PYG{+w}{ }\PYG{n}{keyword}\PYG{+w}{ }\PYG{n}{...}
\end{sphinxVerbatim}
\begin{itemize}
\item {} 
\sphinxAtStartPar
ID, group\sphinxhyphen{}ID are documented in {\hyperref[\detokenize{compute::doc}]{\sphinxcrossref{\DUrole{doc}{compute}}}} command

\item {} 
\sphinxAtStartPar
pe = style name of this compute command

\item {} 
\sphinxAtStartPar
zero or more keywords may be appended

\item {} 
\sphinxAtStartPar
keyword = \sphinxstyleemphasis{pair} or \sphinxstyleemphasis{bond} or \sphinxstyleemphasis{angle} or \sphinxstyleemphasis{dihedral} or \sphinxstyleemphasis{improper} or \sphinxstyleemphasis{kspace} or \sphinxstyleemphasis{fix}

\end{itemize}


\subsection{Examples}
\label{\detokenize{compute_pe:examples}}
\begin{sphinxVerbatim}[commandchars=\\\{\}]
\PYG{k}{compute }\PYG{n+nv+nvIdentifier}{1}\PYG{+w}{ }\PYG{n+nv+nvIdentifier}{all}\PYG{+w}{ }\PYG{n}{pe}
\PYG{k}{compute }\PYG{n+nv+nvIdentifier}{molPE}\PYG{+w}{ }\PYG{n+nv+nvIdentifier}{all}\PYG{+w}{ }\PYG{n}{pe}\PYG{+w}{ }\PYG{n}{bond}\PYG{+w}{ }\PYG{n}{angle}\PYG{+w}{ }\PYG{n}{dihedral}\PYG{+w}{ }\PYG{n}{improper}
\end{sphinxVerbatim}


\subsection{Description}
\label{\detokenize{compute_pe:description}}
\sphinxAtStartPar
Define a computation that calculates the potential energy of the
entire system of atoms.  The specified group must be “all”.  See the
{\hyperref[\detokenize{compute_pe_atom::doc}]{\sphinxcrossref{\DUrole{doc}{compute pe/atom}}}} command if you want per\sphinxhyphen{}atom
energies.  These per\sphinxhyphen{}atom values could be summed for a group of atoms
via the {\hyperref[\detokenize{compute_reduce::doc}]{\sphinxcrossref{\DUrole{doc}{compute reduce}}}} command.

\sphinxAtStartPar
The energy is calculated by the various pair, bond, etc. potentials
defined for the simulation.  If no extra keywords are listed, then the
potential energy is the sum of pair, bond, angle, dihedral, improper,
\(k\)\sphinxhyphen{}space (long\sphinxhyphen{}range), and fix energy (i.e., it is as though all the
keywords were listed).  If any extra keywords are listed, then only
those components are summed to compute the potential energy.

\sphinxAtStartPar
The \(k\)\sphinxhyphen{}space contribution requires 1 extra FFT each timestep the energy
is calculated, if using the PPPM solver via the {\hyperref[\detokenize{kspace_style::doc}]{\sphinxcrossref{\DUrole{doc}{kspace\_style pppm}}}} command.  Thus it can increase the cost of the
PPPM calculation if it is needed on a large fraction of the simulation
timesteps.

\sphinxAtStartPar
Various fixes can contribute to the total potential energy of the
system if the \sphinxstyleemphasis{fix} contribution is included.  See the doc pages for
{\hyperref[\detokenize{fix::doc}]{\sphinxcrossref{\DUrole{doc}{individual fixes}}}} for details of which ones compute a
potential energy.

\begin{sphinxadmonition}{note}{Note:}
\sphinxAtStartPar
The {\hyperref[\detokenize{fix_modify::doc}]{\sphinxcrossref{\DUrole{doc}{fix\_modify energy yes}}}} command must also be
specified if a fix is to contribute potential energy to this command.
\end{sphinxadmonition}

\sphinxAtStartPar
A compute of this style with the ID of “thermo\_pe” is created when
LAMMPS starts up, as if this command were in the input script:

\begin{sphinxVerbatim}[commandchars=\\\{\}]
\PYG{k}{compute }\PYG{n+nv+nvIdentifier}{thermo\PYGZus{}pe}\PYG{+w}{ }\PYG{n+nv+nvIdentifier}{all}\PYG{+w}{ }\PYG{n}{pe}
\end{sphinxVerbatim}

\sphinxAtStartPar
See the “thermo\_style” command for more details.


\bigskip\hrule\bigskip



\subsection{Output info}
\label{\detokenize{compute_pe:output-info}}
\sphinxAtStartPar
This compute calculates a global scalar (the potential energy).  This
value can be used by any command that uses a global scalar value from
a compute as input.  See the {\hyperref[\detokenize{Howto_output::doc}]{\sphinxcrossref{\DUrole{doc}{Howto output}}}} doc page
for an overview of LAMMPS output options.

\sphinxAtStartPar
The scalar value calculated by this compute is “extensive”.  The
scalar value will be in energy {\hyperref[\detokenize{units::doc}]{\sphinxcrossref{\DUrole{doc}{units}}}}.


\subsection{Restrictions}
\label{\detokenize{compute_pe:restrictions}}\begin{quote}

\sphinxAtStartPar
none
\end{quote}


\subsection{Related commands}
\label{\detokenize{compute_pe:related-commands}}
\sphinxAtStartPar
{\hyperref[\detokenize{compute_pe_atom::doc}]{\sphinxcrossref{\DUrole{doc}{compute pe/atom}}}}


\subsection{Default}
\label{\detokenize{compute_pe:default}}
\sphinxAtStartPar
none

\sphinxstepscope

\index{compute pe/atom@\spxentry{compute pe/atom}}\ignorespaces 

\section{compute pe/atom command}
\label{\detokenize{compute_pe_atom:compute-pe-atom-command}}\label{\detokenize{compute_pe_atom:index-0}}\label{\detokenize{compute_pe_atom::doc}}

\subsection{Syntax}
\label{\detokenize{compute_pe_atom:syntax}}
\begin{sphinxVerbatim}[commandchars=\\\{\}]
\PYG{k}{compute }\PYG{n+nv+nvIdentifier}{ID}\PYG{+w}{ }\PYG{n+nv+nvIdentifier}{group\PYGZhy{}ID}\PYG{+w}{ }\PYG{n}{pe}\PYG{o}{/}\PYG{n}{atom}\PYG{+w}{ }\PYG{n}{keyword}\PYG{+w}{ }\PYG{n}{...}
\end{sphinxVerbatim}
\begin{itemize}
\item {} 
\sphinxAtStartPar
ID, group\sphinxhyphen{}ID are documented in {\hyperref[\detokenize{compute::doc}]{\sphinxcrossref{\DUrole{doc}{compute}}}} command

\item {} 
\sphinxAtStartPar
pe/atom = style name of this compute command

\item {} 
\sphinxAtStartPar
zero or more keywords may be appended

\item {} 
\sphinxAtStartPar
keyword = \sphinxstyleemphasis{pair} or \sphinxstyleemphasis{bond} or \sphinxstyleemphasis{angle} or \sphinxstyleemphasis{dihedral} or \sphinxstyleemphasis{improper} or \sphinxstyleemphasis{kspace} or \sphinxstyleemphasis{fix}

\end{itemize}


\subsection{Examples}
\label{\detokenize{compute_pe_atom:examples}}
\begin{sphinxVerbatim}[commandchars=\\\{\}]
\PYG{k}{compute }\PYG{n+nv+nvIdentifier}{1}\PYG{+w}{ }\PYG{n+nv+nvIdentifier}{all}\PYG{+w}{ }\PYG{n}{pe}\PYG{o}{/}\PYG{n}{atom}
\PYG{k}{compute }\PYG{n+nv+nvIdentifier}{1}\PYG{+w}{ }\PYG{n+nv+nvIdentifier}{all}\PYG{+w}{ }\PYG{n}{pe}\PYG{o}{/}\PYG{n}{atom}\PYG{+w}{ }\PYG{n}{pair}
\PYG{k}{compute }\PYG{n+nv+nvIdentifier}{1}\PYG{+w}{ }\PYG{n+nv+nvIdentifier}{all}\PYG{+w}{ }\PYG{n}{pe}\PYG{o}{/}\PYG{n}{atom}\PYG{+w}{ }\PYG{n}{pair}\PYG{+w}{ }\PYG{n}{bond}
\end{sphinxVerbatim}


\subsection{Description}
\label{\detokenize{compute_pe_atom:description}}
\sphinxAtStartPar
Define a computation that computes the per\sphinxhyphen{}atom potential energy for
each atom in a group.  See the {\hyperref[\detokenize{compute_pe::doc}]{\sphinxcrossref{\DUrole{doc}{compute pe}}}} command if
you want the potential energy of the entire system.

\sphinxAtStartPar
The per\sphinxhyphen{}atom energy is calculated by the various pair, bond, etc
potentials defined for the simulation.  If no extra keywords are
listed, then the potential energy is the sum of pair, bond, angle,
dihedral, improper, \(k\)\sphinxhyphen{}space (long\sphinxhyphen{}range), and fix energy (i.e., it is as
though all the keywords were listed).  If any extra keywords are listed,
then only those components are summed to compute the potential energy.

\sphinxAtStartPar
Note that the energy of each atom is due to its interaction with all
other atoms in the simulation, not just with other atoms in the group.

\sphinxAtStartPar
For an energy contribution produced by a small set of atoms (e.g., 4
atoms in a dihedral or 3 atoms in a Tersoff 3\sphinxhyphen{}body interaction), that
energy is assigned in equal portions to each atom in the set (e.g., 1/4 of the
dihedral energy to each of the four atoms).

\sphinxAtStartPar
The {\hyperref[\detokenize{dihedral_charmm::doc}]{\sphinxcrossref{\DUrole{doc}{dihedral\_style charmm}}}} style calculates
pairwise interactions between 1\textendash{}4 atoms.  The energy contribution of
these terms is included in the pair energy, not the dihedral energy.

\sphinxAtStartPar
The KSpace contribution is calculated using the method in
{\hyperref[\detokenize{compute_pe_atom:heyes1}]{\sphinxcrossref{\DUrole{std,std-ref}{(Heyes)}}}} for the Ewald method and a related method for PPPM,
as specified by the {\hyperref[\detokenize{kspace_style::doc}]{\sphinxcrossref{\DUrole{doc}{kspace\_style pppm}}}} command.
For PPPM, the calculation requires 1 extra FFT each timestep that
per\sphinxhyphen{}atom energy is calculated.  This \sphinxhref{PDF/kspace.pdf}{document}
describes how the long\sphinxhyphen{}range per\sphinxhyphen{}atom energy calculation is performed.

\sphinxAtStartPar
Various fixes can contribute to the per\sphinxhyphen{}atom potential energy of the
system if the \sphinxstyleemphasis{fix} contribution is included.  See the doc pages for
{\hyperref[\detokenize{fix::doc}]{\sphinxcrossref{\DUrole{doc}{individual fixes}}}} for details of which ones compute a
per\sphinxhyphen{}atom potential energy.

\begin{sphinxadmonition}{note}{Note:}
\sphinxAtStartPar
The {\hyperref[\detokenize{fix_modify::doc}]{\sphinxcrossref{\DUrole{doc}{fix\_modify energy yes}}}} command must also be
specified if a fix is to contribute per\sphinxhyphen{}atom potential energy to this
command.
\end{sphinxadmonition}

\sphinxAtStartPar
As an example of per\sphinxhyphen{}atom potential energy compared to total potential
energy, these lines in an input script should yield the same result
in the last 2 columns of thermo output:

\begin{sphinxVerbatim}[commandchars=\\\{\}]
\PYG{k}{compute        }\PYG{n+nv+nvIdentifier}{peratom}\PYG{+w}{ }\PYG{n+nv+nvIdentifier}{all}\PYG{+w}{ }\PYG{n}{pe}\PYG{o}{/}\PYG{n}{atom}
\PYG{k}{compute        }\PYG{n+nv+nvIdentifier}{pe}\PYG{+w}{ }\PYG{n+nv+nvIdentifier}{all}\PYG{+w}{ }\PYG{n}{reduce}\PYG{+w}{ }\PYG{n}{sum}\PYG{+w}{ }\PYG{n}{c\PYGZus{}peratom}
\PYG{k}{thermo\PYGZus{}style}\PYG{+w}{   }\PYG{n}{custom}\PYG{+w}{ }\PYG{n}{step}\PYG{+w}{ }\PYG{n}{temp}\PYG{+w}{ }\PYG{n}{etotal}\PYG{+w}{ }\PYG{n}{press}\PYG{+w}{ }\PYG{n}{pe}\PYG{+w}{ }\PYG{n}{c\PYGZus{}pe}
\end{sphinxVerbatim}

\begin{sphinxadmonition}{note}{Note:}
\sphinxAtStartPar
The per\sphinxhyphen{}atom energy does not include any Lennard\sphinxhyphen{}Jones tail
corrections to the energy added by the
{\hyperref[\detokenize{pair_modify::doc}]{\sphinxcrossref{\DUrole{doc}{pair\_modify tail yes}}}} command, since those are
contributions to the global system energy.
\end{sphinxadmonition}


\subsection{Output info}
\label{\detokenize{compute_pe_atom:output-info}}
\sphinxAtStartPar
This compute calculates a per\sphinxhyphen{}atom vector, which can be accessed by
any command that uses per\sphinxhyphen{}atom values from a compute as input.  See
the {\hyperref[\detokenize{Howto_output::doc}]{\sphinxcrossref{\DUrole{doc}{Howto output}}}} page for an overview of
LAMMPS output options.

\sphinxAtStartPar
The per\sphinxhyphen{}atom vector values will be in energy {\hyperref[\detokenize{units::doc}]{\sphinxcrossref{\DUrole{doc}{units}}}}.


\subsection{Restrictions}
\label{\detokenize{compute_pe_atom:restrictions}}

\subsection{Related commands}
\label{\detokenize{compute_pe_atom:related-commands}}
\sphinxAtStartPar
{\hyperref[\detokenize{compute_pe::doc}]{\sphinxcrossref{\DUrole{doc}{compute pe}}}}, {\hyperref[\detokenize{compute_stress_atom::doc}]{\sphinxcrossref{\DUrole{doc}{compute stress/atom}}}}


\subsection{Default}
\label{\detokenize{compute_pe_atom:default}}
\sphinxAtStartPar
none


\bigskip\hrule\bigskip

\phantomsection\label{\detokenize{compute_pe_atom:heyes1}}
\sphinxAtStartPar
\sphinxstylestrong{(Heyes)} Heyes, Phys Rev B 49, 755 (1994),

\sphinxstepscope

\index{compute plasticity/atom@\spxentry{compute plasticity/atom}}\ignorespaces 

\section{compute plasticity/atom command}
\label{\detokenize{compute_plasticity_atom:compute-plasticity-atom-command}}\label{\detokenize{compute_plasticity_atom:index-0}}\label{\detokenize{compute_plasticity_atom::doc}}

\subsection{Syntax}
\label{\detokenize{compute_plasticity_atom:syntax}}
\begin{sphinxVerbatim}[commandchars=\\\{\}]
\PYG{k}{compute }\PYG{n+nv+nvIdentifier}{ID}\PYG{+w}{ }\PYG{n+nv+nvIdentifier}{group\PYGZhy{}ID}\PYG{+w}{ }\PYG{n}{plasticity}\PYG{o}{/}\PYG{n}{atom}
\end{sphinxVerbatim}
\begin{itemize}
\item {} 
\sphinxAtStartPar
ID, group\sphinxhyphen{}ID are documented in compute command

\item {} 
\sphinxAtStartPar
plasticity/atom = style name of this compute command

\end{itemize}


\subsection{Examples}
\label{\detokenize{compute_plasticity_atom:examples}}
\begin{sphinxVerbatim}[commandchars=\\\{\}]
\PYG{k}{compute }\PYG{n+nv+nvIdentifier}{1}\PYG{+w}{ }\PYG{n+nv+nvIdentifier}{all}\PYG{+w}{ }\PYG{n}{plasticity}\PYG{o}{/}\PYG{n}{atom}
\end{sphinxVerbatim}


\subsection{Description}
\label{\detokenize{compute_plasticity_atom:description}}
\sphinxAtStartPar
Define a computation that calculates the per\sphinxhyphen{}atom plasticity for each
atom in a group.  This is a quantity relevant for
{\hyperref[\detokenize{pair_peri::doc}]{\sphinxcrossref{\DUrole{doc}{Peridynamics models}}}}.
See \sphinxhref{PDF/PDLammps\_overview.pdf}{this document}
for an overview of LAMMPS commands for Peridynamics modeling.

\sphinxAtStartPar
The plasticity for a Peridynamic particle is the so\sphinxhyphen{}called consistency
parameter (\(\lambda\)).  For elastic deformation, \(\lambda = 0\),
otherwise \(\lambda > 0\) for plastic deformation.  For details, see
{\hyperref[\detokenize{compute_plasticity_atom:mitchell}]{\sphinxcrossref{\DUrole{std,std-ref}{(Mitchell)}}}} and the PDF doc included in the LAMMPS
distribution in \sphinxhref{PDF/PDLammps\_EPS.pdf}{doc/PDF/PDLammps\_EPS.pdf}.

\sphinxAtStartPar
This command can be invoked for one of the Peridynamic
{\hyperref[\detokenize{pair_peri::doc}]{\sphinxcrossref{\DUrole{doc}{pair styles}}}}: peri/eps.

\sphinxAtStartPar
The plasticity value will be 0.0 for atoms not in the specified
compute group.


\subsection{Output info}
\label{\detokenize{compute_plasticity_atom:output-info}}
\sphinxAtStartPar
This compute calculates a per\sphinxhyphen{}atom vector, which can be accessed by
any command that uses per\sphinxhyphen{}atom values from a compute as input.  See
the {\hyperref[\detokenize{Howto_output::doc}]{\sphinxcrossref{\DUrole{doc}{Howto output}}}} page for an overview of
LAMMPS output options.

\sphinxAtStartPar
The per\sphinxhyphen{}atom vector values are unitless numbers \(\lambda \ge 0.0\).


\subsection{Restrictions}
\label{\detokenize{compute_plasticity_atom:restrictions}}
\sphinxAtStartPar
This compute is part of the PERI package.  It is only enabled if
LAMMPS was built with that package.  See the {\hyperref[\detokenize{Build_package::doc}]{\sphinxcrossref{\DUrole{doc}{Build package}}}} page for more info.


\subsection{Related commands}
\label{\detokenize{compute_plasticity_atom:related-commands}}
\sphinxAtStartPar
{\hyperref[\detokenize{compute_damage_atom::doc}]{\sphinxcrossref{\DUrole{doc}{compute damage/atom}}}},
{\hyperref[\detokenize{compute_dilatation_atom::doc}]{\sphinxcrossref{\DUrole{doc}{compute dilatation/atom}}}}


\subsection{Default}
\label{\detokenize{compute_plasticity_atom:default}}
\sphinxAtStartPar
none


\bigskip\hrule\bigskip

\phantomsection\label{\detokenize{compute_plasticity_atom:mitchell}}
\sphinxAtStartPar
\sphinxstylestrong{(Mitchell)} Mitchell, “A non\sphinxhyphen{}local, ordinary\sphinxhyphen{}state\sphinxhyphen{}based
viscoelasticity model for peridynamics”, Sandia National Lab Report,
8064:1\sphinxhyphen{}28 (2011).

\sphinxstepscope

\index{compute pod/atom@\spxentry{compute pod/atom}}\ignorespaces 
\index{compute podd/atom@\spxentry{compute podd/atom}}\ignorespaces 
\index{compute pod/local@\spxentry{compute pod/local}}\ignorespaces 
\index{compute pod/global@\spxentry{compute pod/global}}\ignorespaces 

\section{compute pod/atom command}
\label{\detokenize{compute_pod_atom:compute-pod-atom-command}}\label{\detokenize{compute_pod_atom:index-3}}\label{\detokenize{compute_pod_atom:index-2}}\label{\detokenize{compute_pod_atom:index-1}}\label{\detokenize{compute_pod_atom:index-0}}\label{\detokenize{compute_pod_atom::doc}}

\section{compute podd/atom command}
\label{\detokenize{compute_pod_atom:compute-podd-atom-command}}

\section{compute pod/local command}
\label{\detokenize{compute_pod_atom:compute-pod-local-command}}

\section{compute pod/global command}
\label{\detokenize{compute_pod_atom:compute-pod-global-command}}

\subsection{Syntax}
\label{\detokenize{compute_pod_atom:syntax}}
\begin{sphinxVerbatim}[commandchars=\\\{\}]
\PYG{k}{compute }\PYG{n+nv+nvIdentifier}{ID}\PYG{+w}{ }\PYG{n+nv+nvIdentifier}{group\PYGZhy{}ID}\PYG{+w}{ }\PYG{n}{pod}\PYG{o}{/}\PYG{n}{atom}\PYG{+w}{ }\PYG{n}{param.pod}\PYG{+w}{ }\PYG{n}{coefficients.pod}
\PYG{k}{compute }\PYG{n+nv+nvIdentifier}{ID}\PYG{+w}{ }\PYG{n+nv+nvIdentifier}{group\PYGZhy{}ID}\PYG{+w}{ }\PYG{n}{podd}\PYG{o}{/}\PYG{n}{atom}\PYG{+w}{ }\PYG{n}{param.pod}\PYG{+w}{ }\PYG{n}{coefficients.pod}
\PYG{k}{compute }\PYG{n+nv+nvIdentifier}{ID}\PYG{+w}{ }\PYG{n+nv+nvIdentifier}{group\PYGZhy{}ID}\PYG{+w}{ }\PYG{n}{pod}\PYG{o}{/}\PYG{n}{local}\PYG{+w}{ }\PYG{n}{param.pod}\PYG{+w}{ }\PYG{n}{coefficients.pod}
\PYG{k}{compute }\PYG{n+nv+nvIdentifier}{ID}\PYG{+w}{ }\PYG{n+nv+nvIdentifier}{group\PYGZhy{}ID}\PYG{+w}{ }\PYG{n}{pod}\PYG{o}{/}\PYG{n}{global}\PYG{+w}{ }\PYG{n}{param.pod}\PYG{+w}{ }\PYG{n}{coefficients.pod}
\end{sphinxVerbatim}
\begin{itemize}
\item {} 
\sphinxAtStartPar
ID, group\sphinxhyphen{}ID are documented in {\hyperref[\detokenize{compute::doc}]{\sphinxcrossref{\DUrole{doc}{compute}}}} command

\item {} 
\sphinxAtStartPar
pod/atom = style name of this compute command

\item {} 
\sphinxAtStartPar
param.pod = the parameter file specifies parameters of the POD descriptors

\item {} 
\sphinxAtStartPar
coefficients.pod = the coefficient file specifies coefficients of the POD potential

\end{itemize}


\subsection{Examples}
\label{\detokenize{compute_pod_atom:examples}}
\begin{sphinxVerbatim}[commandchars=\\\{\}]
\PYG{k}{compute }\PYG{n+nv+nvIdentifier}{d}\PYG{+w}{ }\PYG{n+nv+nvIdentifier}{all}\PYG{+w}{ }\PYG{n}{pod}\PYG{o}{/}\PYG{n}{atom}\PYG{+w}{ }\PYG{n}{Ta\PYGZus{}param.pod}
\PYG{k}{compute }\PYG{n+nv+nvIdentifier}{dd}\PYG{+w}{ }\PYG{n+nv+nvIdentifier}{all}\PYG{+w}{ }\PYG{n}{podd}\PYG{o}{/}\PYG{n}{atom}\PYG{+w}{ }\PYG{n}{Ta\PYGZus{}param.pod}
\PYG{k}{compute }\PYG{n+nv+nvIdentifier}{ldd}\PYG{+w}{ }\PYG{n+nv+nvIdentifier}{all}\PYG{+w}{ }\PYG{n}{pod}\PYG{o}{/}\PYG{n}{local}\PYG{+w}{ }\PYG{n}{Ta\PYGZus{}param.pod}
\PYG{k}{compute }\PYG{n+nv+nvIdentifier}{gdd}\PYG{+w}{ }\PYG{n+nv+nvIdentifier}{all}\PYG{+w}{ }\PYG{n}{podd}\PYG{o}{/}\PYG{n}{global}\PYG{+w}{ }\PYG{n}{Ta\PYGZus{}param.pod}
\PYG{k}{compute }\PYG{n+nv+nvIdentifier}{d}\PYG{+w}{ }\PYG{n+nv+nvIdentifier}{all}\PYG{+w}{ }\PYG{n}{pod}\PYG{o}{/}\PYG{n}{atom}\PYG{+w}{ }\PYG{n}{Ta\PYGZus{}param.pod}\PYG{+w}{ }\PYG{n}{Ta\PYGZus{}coefficients.pod}
\PYG{k}{compute }\PYG{n+nv+nvIdentifier}{dd}\PYG{+w}{ }\PYG{n+nv+nvIdentifier}{all}\PYG{+w}{ }\PYG{n}{podd}\PYG{o}{/}\PYG{n}{atom}\PYG{+w}{ }\PYG{n}{Ta\PYGZus{}param.pod}\PYG{+w}{ }\PYG{n}{Ta\PYGZus{}coefficients.pod}
\PYG{k}{compute }\PYG{n+nv+nvIdentifier}{ldd}\PYG{+w}{ }\PYG{n+nv+nvIdentifier}{all}\PYG{+w}{ }\PYG{n}{pod}\PYG{o}{/}\PYG{n}{local}\PYG{+w}{ }\PYG{n}{Ta\PYGZus{}param.pod}\PYG{+w}{ }\PYG{n}{Ta\PYGZus{}coefficients.pod}
\PYG{k}{compute }\PYG{n+nv+nvIdentifier}{gdd}\PYG{+w}{ }\PYG{n+nv+nvIdentifier}{all}\PYG{+w}{ }\PYG{n}{podd}\PYG{o}{/}\PYG{n}{global}\PYG{+w}{ }\PYG{n}{Ta\PYGZus{}param.pod}\PYG{+w}{ }\PYG{n}{Ta\PYGZus{}coefficients.pod}
\end{sphinxVerbatim}


\subsection{Description}
\label{\detokenize{compute_pod_atom:description}}
\sphinxAtStartPar
\DUrole{versionmodified,added}{Added in version 27June2024.}

\sphinxAtStartPar
Define a computation that calculates a set of quantities related to the
POD descriptors of the atoms in a group. These computes are used
primarily for calculating the dependence of energy and force components
on the linear coefficients in the {\hyperref[\detokenize{pair_pod::doc}]{\sphinxcrossref{\DUrole{doc}{pod pair\_style}}}},
which is useful when training a POD potential to match target data. POD
descriptors of an atom are characterized by the radial and angular
distribution of neighbor atoms. The detailed mathematical definition is
given in the papers by {\hyperref[\detokenize{compute_pod_atom:nguyen20222c}]{\sphinxcrossref{\DUrole{std,std-ref}{(Nguyen and Rohskopf)}}}},
{\hyperref[\detokenize{compute_pod_atom:nguyen20232c}]{\sphinxcrossref{\DUrole{std,std-ref}{(Nguyen2023)}}}}, {\hyperref[\detokenize{compute_pod_atom:nguyen20242c}]{\sphinxcrossref{\DUrole{std,std-ref}{(Nguyen2024)}}}},
and {\hyperref[\detokenize{compute_pod_atom:nguyen20243c}]{\sphinxcrossref{\DUrole{std,std-ref}{(Nguyen and Sema)}}}}.

\sphinxAtStartPar
Compute \sphinxstyleemphasis{pod/atom} calculates the per\sphinxhyphen{}atom POD descriptors.

\sphinxAtStartPar
Compute \sphinxstyleemphasis{podd/atom} calculates derivatives of the per\sphinxhyphen{}atom POD
descriptors with respect to atom positions.

\sphinxAtStartPar
Compute \sphinxstyleemphasis{pod/local} calculates the per\sphinxhyphen{}atom POD descriptors and their
derivatives with respect to atom positions.

\sphinxAtStartPar
Compute \sphinxstyleemphasis{pod/global} calculates the global POD descriptors and their
derivatives with respect to atom positions.

\sphinxAtStartPar
Examples how to use Compute POD commands are found in the directory
\sphinxcode{\sphinxupquote{examples/PACKAGES/pod}}.

\begin{sphinxadmonition}{warning}{Warning:}
\sphinxAtStartPar
All of these compute styles produce \sphinxstyleemphasis{very} large per\sphinxhyphen{}atom output
arrays that scale with the total number of atoms in the system.
This will result in \sphinxstyleemphasis{very} large memory consumption for systems
with a large number of atoms.
\end{sphinxadmonition}


\bigskip\hrule\bigskip



\subsection{Output info}
\label{\detokenize{compute_pod_atom:output-info}}
\sphinxAtStartPar
Compute \sphinxstyleemphasis{pod/atom} produces an 2D array of size \(N \times M\),
where \(N\) is the number of atoms and \(M\) is the number of
descriptors. Each column corresponds to a particular POD descriptor.

\sphinxAtStartPar
Compute \sphinxstyleemphasis{podd/atom} produces an 2D array of size \(N \times (M * 3
N)\). Each column corresponds to a particular derivative of a POD
descriptor.

\sphinxAtStartPar
Compute \sphinxstyleemphasis{pod/local} produces an 2D array of size \((1 + 3N) \times
(M * N)\).  The first row contains the per\sphinxhyphen{}atom descriptors, and the last
3N rows contain the derivatives of the per\sphinxhyphen{}atom descriptors with respect
to atom positions.

\sphinxAtStartPar
Compute \sphinxstyleemphasis{pod/global} produces an 2D array of size \((1 + 3N) \times
(M)\).  The first row contains the global descriptors, and the last 3N
rows contain the derivatives of the global descriptors with respect to
atom positions.


\subsection{Restrictions}
\label{\detokenize{compute_pod_atom:restrictions}}
\sphinxAtStartPar
These computes are part of the ML\sphinxhyphen{}POD package.  They are only enabled
if LAMMPS was built with that package.  See the {\hyperref[\detokenize{Build_package::doc}]{\sphinxcrossref{\DUrole{doc}{Build package}}}} page for more info.


\subsection{Related commands}
\label{\detokenize{compute_pod_atom:related-commands}}
\sphinxAtStartPar
{\hyperref[\detokenize{fitpod_command::doc}]{\sphinxcrossref{\DUrole{doc}{fitpod}}}},
{\hyperref[\detokenize{pair_pod::doc}]{\sphinxcrossref{\DUrole{doc}{pair\_style pod}}}}


\subsection{Default}
\label{\detokenize{compute_pod_atom:default}}
\sphinxAtStartPar
none


\bigskip\hrule\bigskip

\phantomsection\label{\detokenize{compute_pod_atom:nguyen20222c}}
\sphinxAtStartPar
\sphinxstylestrong{(Nguyen and Rohskopf)} Nguyen and Rohskopf,  Journal of Computational Physics, 480, 112030, (2023).

\phantomsection\label{\detokenize{compute_pod_atom:nguyen20232c}}
\sphinxAtStartPar
\sphinxstylestrong{(Nguyen2023)} Nguyen, Physical Review B, 107(14), 144103, (2023).

\phantomsection\label{\detokenize{compute_pod_atom:nguyen20242c}}
\sphinxAtStartPar
\sphinxstylestrong{(Nguyen2024)} Nguyen, Journal of Computational Physics, 113102, (2024).

\phantomsection\label{\detokenize{compute_pod_atom:nguyen20243c}}
\sphinxAtStartPar
\sphinxstylestrong{(Nguyen and Sema)} Nguyen and Sema, \sphinxurl{https://arxiv.org/abs/2405.00306}, (2024).

\sphinxstepscope

\index{compute pressure@\spxentry{compute pressure}}\ignorespaces 

\section{compute pressure command}
\label{\detokenize{compute_pressure:compute-pressure-command}}\label{\detokenize{compute_pressure:index-0}}\label{\detokenize{compute_pressure::doc}}

\subsection{Syntax}
\label{\detokenize{compute_pressure:syntax}}
\begin{sphinxVerbatim}[commandchars=\\\{\}]
\PYG{k}{compute }\PYG{n+nv+nvIdentifier}{ID}\PYG{+w}{ }\PYG{n+nv+nvIdentifier}{group\PYGZhy{}ID}\PYG{+w}{ }\PYG{n}{pressure}\PYG{+w}{ }\PYG{n}{temp}\PYG{o}{\PYGZhy{}}\PYG{n}{ID}\PYG{+w}{ }\PYG{n}{keyword}\PYG{+w}{ }\PYG{n}{...}
\end{sphinxVerbatim}
\begin{itemize}
\item {} 
\sphinxAtStartPar
ID, group\sphinxhyphen{}ID are documented in {\hyperref[\detokenize{compute::doc}]{\sphinxcrossref{\DUrole{doc}{compute}}}} command

\item {} 
\sphinxAtStartPar
pressure = style name of this compute command

\item {} 
\sphinxAtStartPar
temp\sphinxhyphen{}ID = ID of compute that calculates temperature, can be NULL if not needed

\item {} 
\sphinxAtStartPar
zero or more keywords may be appended

\item {} 
\sphinxAtStartPar
keyword = \sphinxstyleemphasis{ke} or \sphinxstyleemphasis{pair} or \sphinxstyleemphasis{bond} or \sphinxstyleemphasis{angle} or \sphinxstyleemphasis{dihedral} or \sphinxstyleemphasis{improper} or \sphinxstyleemphasis{kspace} or \sphinxstyleemphasis{fix} or \sphinxstyleemphasis{virial} or \sphinxstyleemphasis{pair/hybrid}

\end{itemize}


\subsection{Examples}
\label{\detokenize{compute_pressure:examples}}
\begin{sphinxVerbatim}[commandchars=\\\{\}]
\PYG{k}{compute }\PYG{n+nv+nvIdentifier}{1}\PYG{+w}{ }\PYG{n+nv+nvIdentifier}{all}\PYG{+w}{ }\PYG{n}{pressure}\PYG{+w}{ }\PYG{n}{thermo\PYGZus{}temp}
\PYG{k}{compute }\PYG{n+nv+nvIdentifier}{1}\PYG{+w}{ }\PYG{n+nv+nvIdentifier}{all}\PYG{+w}{ }\PYG{n}{pressure}\PYG{+w}{ }\PYG{n}{NULL}\PYG{+w}{ }\PYG{n}{pair}\PYG{+w}{ }\PYG{n}{bond}
\PYG{k}{compute }\PYG{n+nv+nvIdentifier}{1}\PYG{+w}{ }\PYG{n+nv+nvIdentifier}{all}\PYG{+w}{ }\PYG{n}{pressure}\PYG{+w}{ }\PYG{n}{NULL}\PYG{+w}{ }\PYG{n}{pair}\PYG{o}{/}\PYG{n}{hybrid}\PYG{+w}{ }\PYG{n}{lj}\PYG{o}{/}\PYG{n}{cut}
\end{sphinxVerbatim}


\subsection{Description}
\label{\detokenize{compute_pressure:description}}
\sphinxAtStartPar
Define a computation that calculates the pressure of the entire system
of atoms.  The specified group must be “all”.  See the
{\hyperref[\detokenize{compute_stress_atom::doc}]{\sphinxcrossref{\DUrole{doc}{compute stress/atom}}}} command if you want per\sphinxhyphen{}atom
pressure (stress).  These per\sphinxhyphen{}atom values could be summed for a group
of atoms via the {\hyperref[\detokenize{compute_reduce::doc}]{\sphinxcrossref{\DUrole{doc}{compute reduce}}}} command.

\sphinxAtStartPar
The pressure is computed by the formula
\begin{equation*}
\begin{split}P = \frac{N k_B T}{V} + \frac{1}{V d}\sum_{i=1}^{N'} \vec r_i \cdot \vec f_i\end{split}
\end{equation*}
\sphinxAtStartPar
where \sphinxstyleemphasis{N} is the number of atoms in the system (see discussion of DOF
below), \(k_B\) is the Boltzmann constant, \(T\) is the
temperature, \sphinxstyleemphasis{d} is the dimensionality of the system (2 for 2d, 3 for
3d), and \sphinxstyleemphasis{V} is the system volume (or area in 2d).  The second term is
the virial, equal to \(-dU/dV\), computed for all pairwise as well
as 2\sphinxhyphen{}body, 3\sphinxhyphen{}body, 4\sphinxhyphen{}body, many\sphinxhyphen{}body, and long\sphinxhyphen{}range interactions, where
\(\vec r_i\) and \(\vec f_i\) are the position and force vector
of atom \sphinxstyleemphasis{i}, and the dot indicates the dot product (scalar product).
This is computed in parallel for each subdomain and then summed over
all parallel processes. Thus \(N'\) necessarily includes atoms from
neighboring subdomains (so\sphinxhyphen{}called ghost atoms) and the position and
force vectors of ghost atoms are thus included in the summation.  Only
when running in serial and without periodic boundary conditions is
\(N' = N\) the number of atoms in the system.  {\hyperref[\detokenize{fix::doc}]{\sphinxcrossref{\DUrole{doc}{Fixes}}}}
that impose constraints (e.g., the {\hyperref[\detokenize{fix_shake::doc}]{\sphinxcrossref{\DUrole{doc}{fix shake}}}} command)
may also contribute to the virial term.

\sphinxAtStartPar
A symmetric pressure tensor, stored as a 6\sphinxhyphen{}element vector, is also
calculated by this compute.  The six components of the vector are
ordered \(xx,\) \(yy,\) \(zz,\) \(xy,\) \(xz,\)
\(yz.\) The equation for the \((I,J)\) components (where
\(I\) and \(J\) are \(x\), \(y\), or \(z\)) is
similar to the above formula, except that the first term uses
components related to the kinetic energy tensor and the second term
uses components of the virial tensor:
\begin{equation*}
\begin{split}P_{IJ} = \frac{1}{V}\sum_{k=1}^{N} m_k v_{k_I} v_{k_J} +
\frac{1}{V}\sum_{k=1}^{N'} r_{k_I} f_{k_J}.\end{split}
\end{equation*}
\sphinxAtStartPar
If no extra keywords are listed, the entire equations above are
calculated.  This includes a kinetic energy (temperature) term and the
virial as the sum of pair, bond, angle, dihedral, improper, kspace
(long\sphinxhyphen{}range), and fix contributions to the force on each atom.  If any
extra keywords are listed, then only those components are summed to
compute temperature or ke and/or the virial.  The \sphinxstyleemphasis{virial} keyword means
include all terms except the kinetic energy \sphinxstyleemphasis{ke}.

\sphinxAtStartPar
The \sphinxstyleemphasis{pair/hybrid} keyword means to only include contribution
from a sub\sphinxhyphen{}style in a \sphinxstyleemphasis{hybrid} or \sphinxstyleemphasis{hybrid/overlay} pair style.

\sphinxAtStartPar
Details of how LAMMPS computes the virial efficiently for the entire
system, including for many\sphinxhyphen{}body potentials and accounting for the
effects of periodic boundary conditions are discussed in
{\hyperref[\detokenize{compute_pressure:thompson1}]{\sphinxcrossref{\DUrole{std,std-ref}{(Thompson)}}}}.

\sphinxAtStartPar
The temperature and kinetic energy tensor are not calculated by this
compute, but rather by the temperature compute specified with the
command.  See the doc pages for individual compute temp variants for an
explanation of how they calculate temperature and a symmetric tensor
(6\sphinxhyphen{}element vector) whose components are twice that of the traditional KE
tensor.  That tensor is what appears in the pressure tensor formula
above.

\sphinxAtStartPar
If the kinetic energy is not included in the pressure, than the
temperature compute is not used and can be specified as NULL.  Normally
the temperature compute used by compute pressure should calculate the
temperature of all atoms for consistency with the virial term, but any
compute style that calculates temperature can be used (e.g., one that
excludes frozen atoms or other degrees of freedom).

\sphinxAtStartPar
Note that if desired the specified temperature compute can be one that
subtracts off a bias to calculate a temperature using only the thermal
velocity of the atoms (e.g., by subtracting a background streaming
velocity).  See the doc pages for individual {\hyperref[\detokenize{compute::doc}]{\sphinxcrossref{\DUrole{doc}{compute commands}}}} to determine which ones include a bias.

\sphinxAtStartPar
Also note that the \(N\) in the first formula above is really
degrees\sphinxhyphen{}of\sphinxhyphen{}freedom divided by \(d\) = dimensionality, where the
DOF value is calculated by the temperature compute.  See the various
{\hyperref[\detokenize{compute::doc}]{\sphinxcrossref{\DUrole{doc}{compute temperature}}}} styles for details.

\sphinxAtStartPar
A compute of this style with the ID of thermo\_press is created when
LAMMPS starts up, as if this command were in the input script:

\begin{sphinxVerbatim}[commandchars=\\\{\}]
\PYG{k}{compute }\PYG{n+nv+nvIdentifier}{thermo\PYGZus{}press}\PYG{+w}{ }\PYG{n+nv+nvIdentifier}{all}\PYG{+w}{ }\PYG{n}{pressure}\PYG{+w}{ }\PYG{n}{thermo\PYGZus{}temp}
\end{sphinxVerbatim}

\sphinxAtStartPar
where thermo\_temp is the ID of a similarly defined compute of style
“temp”.  See the {\hyperref[\detokenize{thermo_style::doc}]{\sphinxcrossref{\DUrole{doc}{thermo\_style}}}} command for more details.


\bigskip\hrule\bigskip


\sphinxAtStartPar
Styles with a \sphinxstyleemphasis{gpu}, \sphinxstyleemphasis{intel}, \sphinxstyleemphasis{kk}, \sphinxstyleemphasis{omp}, or \sphinxstyleemphasis{opt} suffix are
functionally the same as the corresponding style without the suffix.
They have been optimized to run faster, depending on your available
hardware, as discussed on the {\hyperref[\detokenize{Speed_packages::doc}]{\sphinxcrossref{\DUrole{doc}{Accelerator packages}}}}
page.  The accelerated styles take the same arguments and should
produce the same results, except for round\sphinxhyphen{}off and precision issues.

\sphinxAtStartPar
These accelerated styles are part of the GPU, INTEL, KOKKOS,
OPENMP, and OPT packages, respectively.  They are only enabled if
LAMMPS was built with those packages.  See the {\hyperref[\detokenize{Build_package::doc}]{\sphinxcrossref{\DUrole{doc}{Build package}}}} page for more info.

\sphinxAtStartPar
You can specify the accelerated styles explicitly in your input script
by including their suffix, or you can use the {\hyperref[\detokenize{Run_options::doc}]{\sphinxcrossref{\DUrole{doc}{\sphinxhyphen{}suffix command\sphinxhyphen{}line switch}}}} when you invoke LAMMPS, or you can use the
{\hyperref[\detokenize{suffix::doc}]{\sphinxcrossref{\DUrole{doc}{suffix}}}} command in your input script.

\sphinxAtStartPar
See the {\hyperref[\detokenize{Speed_packages::doc}]{\sphinxcrossref{\DUrole{doc}{Accelerator packages}}}} page for more
instructions on how to use the accelerated styles effectively.


\bigskip\hrule\bigskip



\subsection{Output info}
\label{\detokenize{compute_pressure:output-info}}
\sphinxAtStartPar
This compute calculates a global scalar (the pressure) and a global
vector of length 6 (pressure tensor), which can be accessed by indices
1\textendash{}6.  These values can be used by any command that uses global scalar
or vector values from a compute as input.  See the {\hyperref[\detokenize{Howto_output::doc}]{\sphinxcrossref{\DUrole{doc}{Howto output}}}} page for an overview of LAMMPS output options.

\sphinxAtStartPar
The ordering of values in the symmetric pressure tensor is as follows:
\(p_{xx},\) \(p_{yy},\) \(p_{zz},\) \(p_{xy},\)
\(p_{xz},\) \(p_{yz}.\)

\sphinxAtStartPar
The scalar and vector values calculated by this compute are “intensive”.
The scalar and vector values will be in pressure {\hyperref[\detokenize{units::doc}]{\sphinxcrossref{\DUrole{doc}{units}}}}.


\subsection{Restrictions}
\label{\detokenize{compute_pressure:restrictions}}\begin{quote}

\sphinxAtStartPar
none
\end{quote}


\subsection{Related commands}
\label{\detokenize{compute_pressure:related-commands}}
\sphinxAtStartPar
{\hyperref[\detokenize{compute_temp::doc}]{\sphinxcrossref{\DUrole{doc}{compute temp}}}}, {\hyperref[\detokenize{compute_stress_atom::doc}]{\sphinxcrossref{\DUrole{doc}{compute stress/atom}}}},
{\hyperref[\detokenize{thermo_style::doc}]{\sphinxcrossref{\DUrole{doc}{thermo\_style}}}}, {\hyperref[\detokenize{fix_numdiff_virial::doc}]{\sphinxcrossref{\DUrole{doc}{fix numdiff/virial}}}},


\subsection{Default}
\label{\detokenize{compute_pressure:default}}
\sphinxAtStartPar
By default the compute includes contributions from the keywords:
\sphinxcode{\sphinxupquote{ke pair bond angle dihedral improper kspace fix}}


\bigskip\hrule\bigskip

\phantomsection\label{\detokenize{compute_pressure:thompson1}}
\sphinxAtStartPar
\sphinxstylestrong{(Thompson)} Thompson, Plimpton, Mattson, J Chem Phys, 131, 154107 (2009).

\sphinxstepscope

\index{compute pressure/alchemy@\spxentry{compute pressure/alchemy}}\ignorespaces 

\section{compute pressure/alchemy command}
\label{\detokenize{compute_pressure_alchemy:compute-pressure-alchemy-command}}\label{\detokenize{compute_pressure_alchemy:index-0}}\label{\detokenize{compute_pressure_alchemy::doc}}

\subsection{Syntax}
\label{\detokenize{compute_pressure_alchemy:syntax}}
\begin{sphinxVerbatim}[commandchars=\\\{\}]
\PYG{k}{compute }\PYG{n+nv+nvIdentifier}{ID}\PYG{+w}{ }\PYG{n+nv+nvIdentifier}{group\PYGZhy{}ID}\PYG{+w}{ }\PYG{n}{pressure}\PYG{o}{/}\PYG{n}{alchemy}\PYG{+w}{ }\PYG{n}{fix}\PYG{o}{\PYGZhy{}}\PYG{n}{ID}
\end{sphinxVerbatim}
\begin{itemize}
\item {} 
\sphinxAtStartPar
ID, group\sphinxhyphen{}ID are documented in {\hyperref[\detokenize{compute::doc}]{\sphinxcrossref{\DUrole{doc}{compute}}}} command

\item {} 
\sphinxAtStartPar
pressure/alchemy = style name of this compute command

\item {} 
\sphinxAtStartPar
fix\sphinxhyphen{}ID = ID of {\hyperref[\detokenize{fix_alchemy::doc}]{\sphinxcrossref{\DUrole{doc}{fix alchemy}}}} command

\end{itemize}


\subsection{Examples}
\label{\detokenize{compute_pressure_alchemy:examples}}
\begin{sphinxVerbatim}[commandchars=\\\{\}]
\PYG{k}{fix }\PYG{n+nv+nvIdentifier}{trans}\PYG{+w}{ }\PYG{n+nv+nvIdentifier}{all}\PYG{+w}{ }\PYG{n}{alchemy}
\PYG{k}{compute }\PYG{n+nv+nvIdentifier}{mixed}\PYG{+w}{ }\PYG{n+nv+nvIdentifier}{all}\PYG{+w}{ }\PYG{n}{pressure}\PYG{o}{/}\PYG{n}{alchemy}\PYG{+w}{ }\PYG{n}{trans}
\PYG{k}{thermo\PYGZus{}modify}\PYG{+w}{ }\PYG{n}{press}\PYG{+w}{ }\PYG{n}{mixed}
\end{sphinxVerbatim}


\subsection{Description}
\label{\detokenize{compute_pressure_alchemy:description}}
\sphinxAtStartPar
\DUrole{versionmodified,added}{Added in version 28Mar2023.}

\sphinxAtStartPar
Define a compute style that makes the “mixed” system pressure available
for a system that uses the {\hyperref[\detokenize{fix_alchemy::doc}]{\sphinxcrossref{\DUrole{doc}{fix alchemy}}}} command to
transform one topology to another.  This can be used in combination with
either {\hyperref[\detokenize{thermo_modify::doc}]{\sphinxcrossref{\DUrole{doc}{thermo\_modify press}}}} or {\hyperref[\detokenize{fix_modify::doc}]{\sphinxcrossref{\DUrole{doc}{fix\_modify
press}}}} to output and access a pressure consistent with the
simulated combined two topology system.

\sphinxAtStartPar
The actual pressure is determined with {\hyperref[\detokenize{compute_pressure::doc}]{\sphinxcrossref{\DUrole{doc}{compute pressure}}}} commands that are internally used by {\hyperref[\detokenize{fix_alchemy::doc}]{\sphinxcrossref{\DUrole{doc}{fix
alchemy}}}} for each topology individually and then combined.
This command just extracts the information from the fix.

\sphinxAtStartPar
The \sphinxcode{\sphinxupquote{examples/PACKAGES/alchemy}} folder contains an example input for this command.


\bigskip\hrule\bigskip



\subsection{Output info}
\label{\detokenize{compute_pressure_alchemy:output-info}}
\sphinxAtStartPar
This compute calculates a global scalar (the pressure) and a global
vector of length 6 (the pressure tensor), which can be accessed by
indices 1\textendash{}6.  These values can be used by any command that uses global
scalar or vector values from a compute as input.  See the {\hyperref[\detokenize{Howto_output::doc}]{\sphinxcrossref{\DUrole{doc}{Howto
output}}}} page for an overview of LAMMPS output options.

\sphinxAtStartPar
The ordering of values in the symmetric pressure tensor is as follows:
\(p_{xx},\) \(p_{yy},\) \(p_{zz},\) \(p_{xy},\)
\(p_{xz},\) \(p_{yz}.\)

\sphinxAtStartPar
The scalar and vector values calculated by this compute are “intensive”.
The scalar and vector values will be in pressure {\hyperref[\detokenize{units::doc}]{\sphinxcrossref{\DUrole{doc}{units}}}}.


\subsection{Restrictions}
\label{\detokenize{compute_pressure_alchemy:restrictions}}
\sphinxAtStartPar
This compute is part of the REPLICA package.  It is only enabled if
LAMMPS was built with that package.  See the {\hyperref[\detokenize{Build_package::doc}]{\sphinxcrossref{\DUrole{doc}{Build package}}}} page for more info.


\subsection{Related commands}
\label{\detokenize{compute_pressure_alchemy:related-commands}}
\sphinxAtStartPar
{\hyperref[\detokenize{fix_alchemy::doc}]{\sphinxcrossref{\DUrole{doc}{fix alchemy}}}}, {\hyperref[\detokenize{compute_pressure::doc}]{\sphinxcrossref{\DUrole{doc}{compute pressure}}}},
{\hyperref[\detokenize{thermo_modify::doc}]{\sphinxcrossref{\DUrole{doc}{thermo\_modify}}}}, {\hyperref[\detokenize{fix_modify::doc}]{\sphinxcrossref{\DUrole{doc}{fix\_modify}}}}


\subsection{Default}
\label{\detokenize{compute_pressure_alchemy:default}}
\sphinxAtStartPar
none

\sphinxstepscope

\index{compute pressure/uef@\spxentry{compute pressure/uef}}\ignorespaces 

\section{compute pressure/uef command}
\label{\detokenize{compute_pressure_uef:compute-pressure-uef-command}}\label{\detokenize{compute_pressure_uef:index-0}}\label{\detokenize{compute_pressure_uef::doc}}

\subsection{Syntax}
\label{\detokenize{compute_pressure_uef:syntax}}
\begin{sphinxVerbatim}[commandchars=\\\{\}]
\PYG{k}{compute }\PYG{n+nv+nvIdentifier}{ID}\PYG{+w}{ }\PYG{n+nv+nvIdentifier}{group\PYGZhy{}ID}\PYG{+w}{ }\PYG{n}{pressure}\PYG{o}{/}\PYG{n}{uef}\PYG{+w}{ }\PYG{n}{temp}\PYG{o}{\PYGZhy{}}\PYG{n}{ID}\PYG{+w}{ }\PYG{n}{keyword}\PYG{+w}{ }\PYG{n}{...}
\end{sphinxVerbatim}
\begin{itemize}
\item {} 
\sphinxAtStartPar
ID, group\sphinxhyphen{}ID are documented in {\hyperref[\detokenize{compute::doc}]{\sphinxcrossref{\DUrole{doc}{compute}}}} command

\item {} 
\sphinxAtStartPar
pressure/uef = style name of this compute command

\item {} 
\sphinxAtStartPar
temp\sphinxhyphen{}ID = ID of compute that calculates temperature, can be NULL if not needed

\item {} 
\sphinxAtStartPar
zero or more keywords may be appended

\item {} 
\sphinxAtStartPar
keyword = \sphinxstyleemphasis{ke} or \sphinxstyleemphasis{pair} or \sphinxstyleemphasis{bond} or \sphinxstyleemphasis{angle} or \sphinxstyleemphasis{dihedral} or \sphinxstyleemphasis{improper} or \sphinxstyleemphasis{kspace} or \sphinxstyleemphasis{fix} or \sphinxstyleemphasis{virial}

\end{itemize}


\subsection{Examples}
\label{\detokenize{compute_pressure_uef:examples}}
\begin{sphinxVerbatim}[commandchars=\\\{\}]
\PYG{k}{compute }\PYG{n+nv+nvIdentifier}{1}\PYG{+w}{ }\PYG{n+nv+nvIdentifier}{all}\PYG{+w}{ }\PYG{n}{pressure}\PYG{o}{/}\PYG{n}{uef}\PYG{+w}{ }\PYG{n}{my\PYGZus{}temp\PYGZus{}uef}
\PYG{k}{compute }\PYG{n+nv+nvIdentifier}{2}\PYG{+w}{ }\PYG{n+nv+nvIdentifier}{all}\PYG{+w}{ }\PYG{n}{pressure}\PYG{o}{/}\PYG{n}{uef}\PYG{+w}{ }\PYG{n}{my\PYGZus{}temp\PYGZus{}uef}\PYG{+w}{ }\PYG{n}{virial}
\end{sphinxVerbatim}


\subsection{Description}
\label{\detokenize{compute_pressure_uef:description}}
\sphinxAtStartPar
This command is used to compute the pressure tensor in
the reference frame of the applied flow field when
{\hyperref[\detokenize{fix_nh_uef::doc}]{\sphinxcrossref{\DUrole{doc}{fix nvt/uef}}}} or
{\hyperref[\detokenize{fix_nh_uef::doc}]{\sphinxcrossref{\DUrole{doc}{fix npt/uef}}}} is used.
It is not necessary to use this command to compute the scalar
value of the pressure. A {\hyperref[\detokenize{compute_pressure::doc}]{\sphinxcrossref{\DUrole{doc}{compute pressure}}}}
may be used for that purpose.

\sphinxAtStartPar
The keywords and output information are documented in
{\hyperref[\detokenize{compute_pressure::doc}]{\sphinxcrossref{\DUrole{doc}{compute\_pressure}}}}.


\subsection{Restrictions}
\label{\detokenize{compute_pressure_uef:restrictions}}
\sphinxAtStartPar
This fix is part of the UEF package. It is only enabled if LAMMPS
was built with that package. See the {\hyperref[\detokenize{Build_package::doc}]{\sphinxcrossref{\DUrole{doc}{Build package}}}} page
for more info.

\sphinxAtStartPar
This command can only be used when {\hyperref[\detokenize{fix_nh_uef::doc}]{\sphinxcrossref{\DUrole{doc}{fix nvt/uef}}}}
or {\hyperref[\detokenize{fix_nh_uef::doc}]{\sphinxcrossref{\DUrole{doc}{fix npt/uef}}}} is active.

\sphinxAtStartPar
The kinetic contribution to the pressure tensor
will be accurate only when the compute specified by \sphinxstyleemphasis{temp\sphinxhyphen{}ID} is a
{\hyperref[\detokenize{compute_temp_uef::doc}]{\sphinxcrossref{\DUrole{doc}{compute temp/uef}}}}.


\subsection{Related commands}
\label{\detokenize{compute_pressure_uef:related-commands}}
\sphinxAtStartPar
{\hyperref[\detokenize{compute_pressure::doc}]{\sphinxcrossref{\DUrole{doc}{compute pressure}}}},
{\hyperref[\detokenize{fix_nh_uef::doc}]{\sphinxcrossref{\DUrole{doc}{fix nvt/uef}}}},
{\hyperref[\detokenize{compute_temp_uef::doc}]{\sphinxcrossref{\DUrole{doc}{compute temp/uef}}}}


\subsection{Default}
\label{\detokenize{compute_pressure_uef:default}}
\sphinxAtStartPar
none

\sphinxstepscope

\index{compute property/atom@\spxentry{compute property/atom}}\ignorespaces 

\section{compute property/atom command}
\label{\detokenize{compute_property_atom:compute-property-atom-command}}\label{\detokenize{compute_property_atom:index-0}}\label{\detokenize{compute_property_atom::doc}}

\subsection{Syntax}
\label{\detokenize{compute_property_atom:syntax}}
\begin{sphinxVerbatim}[commandchars=\\\{\}]
\PYG{k}{compute }\PYG{n+nv+nvIdentifier}{ID}\PYG{+w}{ }\PYG{n+nv+nvIdentifier}{group\PYGZhy{}ID}\PYG{+w}{ }\PYG{n}{property}\PYG{o}{/}\PYG{n}{atom}\PYG{+w}{ }\PYG{n}{input1}\PYG{+w}{ }\PYG{n}{input2}\PYG{+w}{ }\PYG{n}{...}
\end{sphinxVerbatim}
\begin{itemize}
\item {} 
\sphinxAtStartPar
ID, group\sphinxhyphen{}ID are documented in {\hyperref[\detokenize{compute::doc}]{\sphinxcrossref{\DUrole{doc}{compute}}}} command

\item {} 
\sphinxAtStartPar
property/atom = style name of this compute command

\item {} 
\sphinxAtStartPar
input = one or more atom attributes

\begin{sphinxVerbatim}[commandchars=\\\{\}]
possible attributes = id, mol, proc, type, mass,
                      x, y, z, xs, ys, zs, xu, yu, zu, ix, iy, iz,
                      vx, vy, vz, fx, fy, fz,
                      q, mux, muy, muz, mu,
                      spx, spy, spz, sp, fmx, fmy, fmz,
                      nbonds,
                      radius, diameter, omegax, omegay, omegaz,
                      temperature, heatflow,
                      angmomx, angmomy, angmomz,
                      shapex, shapey, shapez,
                      quatw, quati, quatj, quatk, tqx, tqy, tqz,
                      end1x, end1y, end1z, end2x, end2y, end2z,
                      corner1x, corner1y, corner1z,
                      corner2x, corner2y, corner2z,
                      corner3x, corner3y, corner3z,
                      i\PYGZus{}name, d\PYGZus{}name, i2\PYGZus{}name[I], d2\PYGZus{}name[I],
                      vfrac, s0, espin, eradius, ervel, erforce,
                      rho, drho, e, de, cv, buckling,
\end{sphinxVerbatim}
\begin{sphinxalltt}
\sphinxstyleemphasis{id} = atom ID
\sphinxstyleemphasis{mol} = molecule ID
\sphinxstyleemphasis{proc} = ID of processor that owns atom
\sphinxstyleemphasis{type} = atom type
\sphinxstyleemphasis{mass} = atom mass
\sphinxstyleemphasis{x,y,z} = unscaled atom coordinates
\sphinxstyleemphasis{xs,ys,zs} = scaled atom coordinates
\sphinxstyleemphasis{xu,yu,zu} = unwrapped atom coordinates
\sphinxstyleemphasis{ix,iy,iz} = box image that the atom is in
\sphinxstyleemphasis{vx,vy,vz} = atom velocities
\sphinxstyleemphasis{fx,fy,fz} = forces on atoms
\sphinxstyleemphasis{q} = atom charge
\sphinxstyleemphasis{mux,muy,muz} = orientation of dipole moment of atom
\sphinxstyleemphasis{mu} = magnitude of dipole moment of atom
\sphinxstyleemphasis{spx, spy, spz} = direction of the atomic magnetic spin
\sphinxstyleemphasis{sp} = magintude of atomic magnetic spin moment
\sphinxstyleemphasis{fmx, fmy, fmz} = magnetic force
\sphinxstyleemphasis{nbonds} = number of bonds assigned to an atom
\sphinxstyleemphasis{radius,diameter} = radius,diameter of spherical particle
\sphinxstyleemphasis{omegax,omegay,omegaz} = angular velocity of spherical particle
\sphinxstyleemphasis{temperature} = internal temperature of spherical particle
\sphinxstyleemphasis{heatflow} = internal heat flow of spherical particle
\sphinxstyleemphasis{angmomx,angmomy,angmomz} = angular momentum of aspherical particle
\sphinxstyleemphasis{shapex,shapey,shapez} = 3 diameters of aspherical particle
\sphinxstyleemphasis{quatw,quati,quatj,quatk} = quaternion components for aspherical or body particles
\sphinxstyleemphasis{tqx,tqy,tqz} = torque on finite\sphinxhyphen{}size particles
\sphinxstyleemphasis{end12x, end12y, end12z} = end points of line segment
\sphinxstyleemphasis{corner123x, corner123y, corner123z} = corner points of triangle
\sphinxstyleemphasis{i\_name} = custom integer vector with name
\sphinxstyleemphasis{d\_name} = custom floating point vector with name
\sphinxstyleemphasis{i2\_name{[}I{]}} = Ith column of custom integer array with name
\sphinxstyleemphasis{d2\_name{[}I{]}} = Ith column of custom floating\sphinxhyphen{}point array with name
\end{sphinxalltt}

\begin{sphinxVerbatim}[commandchars=\\\{\}]
PERI package per\PYGZhy{}atom properties:
vfrac = volume fraction
s0 = max stretch of any bond a particle is part of
\end{sphinxVerbatim}

\begin{sphinxVerbatim}[commandchars=\\\{\}]
EFF and AWPMD package per\PYGZhy{}atom properties:
espin = electron spin
eradius = electron radius
ervel = electron radial velocity
erforce = electron radial force
\end{sphinxVerbatim}

\begin{sphinxVerbatim}[commandchars=\\\{\}]
SPH package per\PYGZhy{}atom properties:
rho = density of SPH particles
drho = change in density
e = energy
de = change in thermal energy
cv = heat capacity
\end{sphinxVerbatim}

\end{itemize}


\subsection{Examples}
\label{\detokenize{compute_property_atom:examples}}
\begin{sphinxVerbatim}[commandchars=\\\{\}]
\PYG{k}{compute }\PYG{n+nv+nvIdentifier}{1}\PYG{+w}{ }\PYG{n+nv+nvIdentifier}{all}\PYG{+w}{ }\PYG{n}{property}\PYG{o}{/}\PYG{n}{atom}\PYG{+w}{ }\PYG{n}{xs}\PYG{+w}{ }\PYG{n}{vx}\PYG{+w}{ }\PYG{n}{fx}\PYG{+w}{ }\PYG{n}{mux}
\PYG{k}{compute }\PYG{n+nv+nvIdentifier}{2}\PYG{+w}{ }\PYG{n+nv+nvIdentifier}{all}\PYG{+w}{ }\PYG{n}{property}\PYG{o}{/}\PYG{n}{atom}\PYG{+w}{ }\PYG{n}{type}
\PYG{k}{compute }\PYG{n+nv+nvIdentifier}{1}\PYG{+w}{ }\PYG{n+nv+nvIdentifier}{all}\PYG{+w}{ }\PYG{n}{property}\PYG{o}{/}\PYG{n}{atom}\PYG{+w}{ }\PYG{n}{ix}\PYG{+w}{ }\PYG{n}{iy}\PYG{+w}{ }\PYG{n}{iz}
\PYG{k}{compute }\PYG{n+nv+nvIdentifier}{3}\PYG{+w}{ }\PYG{n+nv+nvIdentifier}{all}\PYG{+w}{ }\PYG{n}{property}\PYG{o}{/}\PYG{n}{atom}\PYG{+w}{ }\PYG{n}{sp}\PYG{+w}{ }\PYG{n}{spx}\PYG{+w}{ }\PYG{n}{spy}\PYG{+w}{ }\PYG{n}{spz}
\PYG{k}{compute }\PYG{n+nv+nvIdentifier}{1}\PYG{+w}{ }\PYG{n+nv+nvIdentifier}{all}\PYG{+w}{ }\PYG{n}{property}\PYG{o}{/}\PYG{n}{atom}\PYG{+w}{ }\PYG{n}{i\PYGZus{}myFlag}\PYG{+w}{ }\PYG{n}{d\PYGZus{}Sxyz[1]}\PYG{+w}{ }\PYG{n}{d\PYGZus{}Sxyz[3]}
\end{sphinxVerbatim}


\subsection{Description}
\label{\detokenize{compute_property_atom:description}}
\sphinxAtStartPar
Define a computation that simply stores atom attributes for each atom
in the group.  This is useful so that the values can be used by other
{\hyperref[\detokenize{Howto_output::doc}]{\sphinxcrossref{\DUrole{doc}{output commands}}}} that take computes as inputs.
See for example, the {\hyperref[\detokenize{compute_reduce::doc}]{\sphinxcrossref{\DUrole{doc}{compute reduce}}}}, {\hyperref[\detokenize{fix_ave_atom::doc}]{\sphinxcrossref{\DUrole{doc}{fix
ave/atom}}}}, {\hyperref[\detokenize{fix_ave_histo::doc}]{\sphinxcrossref{\DUrole{doc}{fix ave/histo}}}},
{\hyperref[\detokenize{fix_ave_chunk::doc}]{\sphinxcrossref{\DUrole{doc}{fix ave/chunk}}}}, and {\hyperref[\detokenize{variable::doc}]{\sphinxcrossref{\DUrole{doc}{atom\sphinxhyphen{}style variable}}}} commands.

\sphinxAtStartPar
The list of possible attributes is essentially the same as that used
by the {\hyperref[\detokenize{dump::doc}]{\sphinxcrossref{\DUrole{doc}{dump custom}}}} command, which describes their
meaning, with some additional quantities that are only defined for
certain {\hyperref[\detokenize{atom_style::doc}]{\sphinxcrossref{\DUrole{doc}{atom styles}}}}.  The goal of this augmented
list gives an input script access to any per\sphinxhyphen{}atom quantity stored by
LAMMPS.

\sphinxAtStartPar
The values are stored in a per\sphinxhyphen{}atom vector or array as discussed
below.  Zeroes are stored for atoms not in the specified group or for
quantities that are not defined for a particular particle in the group
(e.g., \sphinxstyleemphasis{shapex} if the particle is not an ellipsoid).

\sphinxAtStartPar
Attributes \sphinxstyleemphasis{i\_name}, \sphinxstyleemphasis{d\_name}, \sphinxstyleemphasis{i2\_name}, \sphinxstyleemphasis{d2\_name} refer to custom
per\sphinxhyphen{}atom integer and floating\sphinxhyphen{}point vectors or arrays that have been
added via the {\hyperref[\detokenize{fix_property_atom::doc}]{\sphinxcrossref{\DUrole{doc}{fix property/atom}}}} command.
When that command is used specific names are given to each attribute
which are the “name” portion of these attributes.  For arrays
\sphinxstyleemphasis{i2\_name} and \sphinxstyleemphasis{d2\_name}, the column of the array must also be included
following the name in brackets (e.g., d2\_xyz{[}2{]} or i2\_mySpin{[}3{]}).

\sphinxAtStartPar
The additional quantities only accessible via this command, and not
directly via the {\hyperref[\detokenize{dump::doc}]{\sphinxcrossref{\DUrole{doc}{dump custom}}}} command, are as follows.

\sphinxAtStartPar
\sphinxstyleemphasis{Nbonds} is available for all molecular atom styles and refers to the
number of explicit bonds assigned to an atom.  Note that if the
{\hyperref[\detokenize{newton::doc}]{\sphinxcrossref{\DUrole{doc}{newton bond}}}} command is set to \sphinxstyleemphasis{on}, which is the
default, then every bond in the system is assigned to only one of the
two atoms in the bond.  Thus a bond between atoms \(I\) and \(J\) may
be tallied for either atom \(I\) or atom \(J\).
If {\hyperref[\detokenize{newton::doc}]{\sphinxcrossref{\DUrole{doc}{newton bond off}}}} is set, it will be tallied with both atom
\(I\) and atom \(J\).

\sphinxAtStartPar
The quantities \sphinxstyleemphasis{shapex}, \sphinxstyleemphasis{shapey}, and \sphinxstyleemphasis{shapez} are defined for ellipsoidal
particles and define the 3d shape of each particle.

\sphinxAtStartPar
The quantities \sphinxstyleemphasis{quatw}, \sphinxstyleemphasis{quati}, \sphinxstyleemphasis{quatj}, and \sphinxstyleemphasis{quatk} are defined for
ellipsoidal particles and body particles and store the 4\sphinxhyphen{}vector quaternion
representing the orientation of each particle.  See the {\hyperref[\detokenize{set::doc}]{\sphinxcrossref{\DUrole{doc}{set}}}}
command for an explanation of the quaternion vector.

\sphinxAtStartPar
\sphinxstyleemphasis{End1x}, \sphinxstyleemphasis{end1y}, \sphinxstyleemphasis{end1z}, \sphinxstyleemphasis{end2x}, \sphinxstyleemphasis{end2y}, \sphinxstyleemphasis{end2z}, are defined for line
segment particles and define the end points of each line segment.

\sphinxAtStartPar
\sphinxstyleemphasis{Corner1x}, \sphinxstyleemphasis{corner1y}, \sphinxstyleemphasis{corner1z}, \sphinxstyleemphasis{corner2x}, \sphinxstyleemphasis{corner2y},
\sphinxstyleemphasis{corner2z}, \sphinxstyleemphasis{corner3x}, \sphinxstyleemphasis{corner3y}, \sphinxstyleemphasis{corner3z}, are defined for
triangular particles and define the corner points of each triangle.

\sphinxAtStartPar
In addition, the various per\sphinxhyphen{}atom quantities listed above for specific
packages are only accessible by this command.

\sphinxAtStartPar
\DUrole{versionmodified,changed}{Changed in version 15Sep2022: }The \sphinxstyleemphasis{espin} property was previously called \sphinxstyleemphasis{spin}.


\subsection{Output info}
\label{\detokenize{compute_property_atom:output-info}}
\sphinxAtStartPar
This compute calculates a per\sphinxhyphen{}atom vector or per\sphinxhyphen{}atom array depending
on the number of input values.  If a single input is specified, a
per\sphinxhyphen{}atom vector is produced.  If two or more inputs are specified, a
per\sphinxhyphen{}atom array is produced where the number of columns = the number of
inputs.  The vector or array can be accessed by any command that uses
per\sphinxhyphen{}atom values from a compute as input.  See the {\hyperref[\detokenize{Howto_output::doc}]{\sphinxcrossref{\DUrole{doc}{Howto output}}}} page for an overview of LAMMPS output options.

\sphinxAtStartPar
The vector or array values will be in whatever {\hyperref[\detokenize{units::doc}]{\sphinxcrossref{\DUrole{doc}{units}}}} the
corresponding attribute is in (e.g., velocity units for \sphinxstyleemphasis{vx}, charge
units for \sphinxstyleemphasis{q}).

\sphinxAtStartPar
For the spin quantities, \sphinxstyleemphasis{sp} is in the units of the Bohr magneton;
\sphinxstyleemphasis{spx}, \sphinxstyleemphasis{spy}, and \sphinxstyleemphasis{spz} are unitless quantities; and \sphinxstyleemphasis{fmx}, \sphinxstyleemphasis{fmy}, and \sphinxstyleemphasis{fmz}
are given in rad/THz.


\subsection{Restrictions}
\label{\detokenize{compute_property_atom:restrictions}}\begin{quote}

\sphinxAtStartPar
none
\end{quote}


\subsection{Related commands}
\label{\detokenize{compute_property_atom:related-commands}}
\sphinxAtStartPar
{\hyperref[\detokenize{dump::doc}]{\sphinxcrossref{\DUrole{doc}{dump custom}}}}, {\hyperref[\detokenize{compute_reduce::doc}]{\sphinxcrossref{\DUrole{doc}{compute reduce}}}},
{\hyperref[\detokenize{fix_ave_atom::doc}]{\sphinxcrossref{\DUrole{doc}{fix ave/atom}}}}, {\hyperref[\detokenize{fix_ave_chunk::doc}]{\sphinxcrossref{\DUrole{doc}{fix ave/chunk}}}},
{\hyperref[\detokenize{fix_property_atom::doc}]{\sphinxcrossref{\DUrole{doc}{fix property/atom}}}}


\subsection{Default}
\label{\detokenize{compute_property_atom:default}}
\sphinxAtStartPar
none

\sphinxstepscope

\index{compute property/chunk@\spxentry{compute property/chunk}}\ignorespaces 

\section{compute property/chunk command}
\label{\detokenize{compute_property_chunk:compute-property-chunk-command}}\label{\detokenize{compute_property_chunk:index-0}}\label{\detokenize{compute_property_chunk::doc}}

\subsection{Syntax}
\label{\detokenize{compute_property_chunk:syntax}}
\begin{sphinxVerbatim}[commandchars=\\\{\}]
\PYG{k}{compute }\PYG{n+nv+nvIdentifier}{ID}\PYG{+w}{ }\PYG{n+nv+nvIdentifier}{group\PYGZhy{}ID}\PYG{+w}{ }\PYG{n}{property}\PYG{o}{/}\PYG{n}{chunk}\PYG{+w}{ }\PYG{n}{chunkID}\PYG{+w}{ }\PYG{n}{input1}\PYG{+w}{ }\PYG{n}{input2}\PYG{+w}{ }\PYG{n}{...}
\end{sphinxVerbatim}
\begin{itemize}
\item {} 
\sphinxAtStartPar
ID, group\sphinxhyphen{}ID are documented in {\hyperref[\detokenize{compute::doc}]{\sphinxcrossref{\DUrole{doc}{compute}}}} command

\item {} 
\sphinxAtStartPar
property/chunk = style name of this compute command

\item {} 
\sphinxAtStartPar
chunkID = ID of {\hyperref[\detokenize{compute_chunk_atom::doc}]{\sphinxcrossref{\DUrole{doc}{compute chunk/atom}}}} command that defines the chunks

\item {} 
\sphinxAtStartPar
input1,etc = one or more attributes
\begin{sphinxalltt}
attributes = count, id, coord1, coord2, coord3
  count = \# of atoms in chunk
  id = original chunk IDs before compression by {\hyperref[\detokenize{compute_chunk_atom::doc}]{\sphinxcrossref{\DUrole{doc}{compute chunk/atom}}}}
  coord123 = coordinates for spatial bins calculated by {\hyperref[\detokenize{compute_chunk_atom::doc}]{\sphinxcrossref{\DUrole{doc}{compute chunk/atom}}}}
\end{sphinxalltt}

\end{itemize}


\subsection{Examples}
\label{\detokenize{compute_property_chunk:examples}}
\begin{sphinxVerbatim}[commandchars=\\\{\}]
\PYG{k}{compute }\PYG{n+nv+nvIdentifier}{1}\PYG{+w}{ }\PYG{n+nv+nvIdentifier}{all}\PYG{+w}{ }\PYG{n}{property}\PYG{o}{/}\PYG{n}{chunk}\PYG{+w}{ }\PYG{n}{bin2d}\PYG{+w}{ }\PYG{n}{id}\PYG{+w}{ }\PYG{n}{count}
\PYG{k}{compute }\PYG{n+nv+nvIdentifier}{1}\PYG{+w}{ }\PYG{n+nv+nvIdentifier}{all}\PYG{+w}{ }\PYG{n}{property}\PYG{o}{/}\PYG{n}{chunk}\PYG{+w}{ }\PYG{n}{myChunks}\PYG{+w}{ }\PYG{n}{id}\PYG{+w}{ }\PYG{n}{coord1}
\end{sphinxVerbatim}


\subsection{Description}
\label{\detokenize{compute_property_chunk:description}}
\sphinxAtStartPar
Define a computation that stores the specified attributes of chunks of
atoms.

\sphinxAtStartPar
In LAMMPS, chunks are collections of atoms defined by a {\hyperref[\detokenize{compute_chunk_atom::doc}]{\sphinxcrossref{\DUrole{doc}{compute
chunk/atom}}}} command, which assigns each atom to a
single chunk (or no chunk).  The ID for this command is specified as
chunkID.  For example, a single chunk could be the atoms in a molecule
or atoms in a spatial bin.  See the {\hyperref[\detokenize{compute_chunk_atom::doc}]{\sphinxcrossref{\DUrole{doc}{compute chunk/atom}}}} and {\hyperref[\detokenize{Howto_chunk::doc}]{\sphinxcrossref{\DUrole{doc}{Howto chunk}}}} doc pages
for details of how chunks can be defined and examples of how they can
be used to measure properties of a system.

\sphinxAtStartPar
This compute calculates and stores the specified attributes of chunks
as global data so they can be accessed by other {\hyperref[\detokenize{Howto_output::doc}]{\sphinxcrossref{\DUrole{doc}{output commands}}}} and used in conjunction with other commands that
generate per\sphinxhyphen{}chunk data, such as {\hyperref[\detokenize{compute_com_chunk::doc}]{\sphinxcrossref{\DUrole{doc}{compute com/chunk}}}} or {\hyperref[\detokenize{compute_msd_chunk::doc}]{\sphinxcrossref{\DUrole{doc}{compute msd/chunk}}}}.

\sphinxAtStartPar
Note that only atoms in the specified group contribute to the
calculation of the \sphinxstyleemphasis{count} attribute.  The {\hyperref[\detokenize{compute_chunk_atom::doc}]{\sphinxcrossref{\DUrole{doc}{compute chunk/atom}}}} command defines its own group; atoms will have a
chunk ID = 0 if they are not in that group, signifying they are not
assigned to a chunk, and will thus also not contribute to this
calculation.  You can specify the “all” group for this command if you
simply want to include atoms with non\sphinxhyphen{}zero chunk IDs.

\sphinxAtStartPar
The \sphinxstyleemphasis{count} attribute is the number of atoms in the chunk.

\sphinxAtStartPar
The \sphinxstyleemphasis{id} attribute stores the original chunk ID for each chunk.  It
can only be used if the \sphinxstyleemphasis{compress} keyword was set to \sphinxstyleemphasis{yes} for the
{\hyperref[\detokenize{compute_chunk_atom::doc}]{\sphinxcrossref{\DUrole{doc}{compute chunk/atom}}}} command referenced by
chunkID.  This means that the original chunk IDs (e.g., molecule IDs)
will have been compressed to remove chunk IDs with no atoms assigned
to them.  Thus a compressed chunk ID of 3 may correspond to an
original chunk ID (molecule ID in this case) of 415.  The \sphinxstyleemphasis{id}
attribute will then be 415 for the third chunk.

\sphinxAtStartPar
The \sphinxstyleemphasis{coordN} attributes can only be used if a \sphinxstyleemphasis{binning} style was used

\sphinxAtStartPar
in the {\hyperref[\detokenize{compute_chunk_atom::doc}]{\sphinxcrossref{\DUrole{doc}{compute chunk/atom}}}} command
referenced by chunkID.  For \sphinxstyleemphasis{bin/1d}, \sphinxstyleemphasis{bin/2d}, and \sphinxstyleemphasis{bin/3d} styles
the attribute is the center point of the bin in the corresponding
dimension.  Style \sphinxstyleemphasis{bin/1d} only defines a \sphinxstyleemphasis{coord1} attribute.  Style
\sphinxstyleemphasis{bin/2d} adds a \sphinxstyleemphasis{coord2} attribute.  Style \sphinxstyleemphasis{bin/3d} adds a \sphinxstyleemphasis{coord3}
attribute.

\sphinxAtStartPar
Note that if the value of the \sphinxstyleemphasis{units} keyword used in the
{\hyperref[\detokenize{compute_chunk_atom::doc}]{\sphinxcrossref{\DUrole{doc}{compute chunk/atom command}}}} is \sphinxstyleemphasis{box} or
\sphinxstyleemphasis{lattice}, the \sphinxstyleemphasis{coordN} attributes will be in distance {\hyperref[\detokenize{units::doc}]{\sphinxcrossref{\DUrole{doc}{units}}}}.  If the value of the \sphinxstyleemphasis{units} keyword is \sphinxstyleemphasis{reduced}, the
\sphinxstyleemphasis{coordN} attributes will be in unitless reduced units (0\sphinxhyphen{}1).

\sphinxAtStartPar
The simplest way to output the results of the compute property/chunk
calculation to a file is to use the {\hyperref[\detokenize{fix_ave_time::doc}]{\sphinxcrossref{\DUrole{doc}{fix ave/time}}}}
command, for example:

\begin{sphinxVerbatim}[commandchars=\\\{\}]
\PYG{k}{compute }\PYG{n+nv+nvIdentifier}{cc1}\PYG{+w}{ }\PYG{n+nv+nvIdentifier}{all}\PYG{+w}{ }\PYG{n}{chunk}\PYG{o}{/}\PYG{n}{atom}\PYG{+w}{ }\PYG{n}{molecule}
\PYG{k}{compute }\PYG{n+nv+nvIdentifier}{myChunk1}\PYG{+w}{ }\PYG{n+nv+nvIdentifier}{all}\PYG{+w}{ }\PYG{n}{property}\PYG{o}{/}\PYG{n}{chunk}\PYG{+w}{ }\PYG{n}{cc1}\PYG{+w}{ }\PYG{n}{count}
\PYG{k}{compute }\PYG{n+nv+nvIdentifier}{myChunk2}\PYG{+w}{ }\PYG{n+nv+nvIdentifier}{all}\PYG{+w}{ }\PYG{n}{com}\PYG{o}{/}\PYG{n}{chunk}\PYG{+w}{ }\PYG{n}{cc1}
\PYG{k}{fix }\PYG{n+nv+nvIdentifier}{1}\PYG{+w}{ }\PYG{n+nv+nvIdentifier}{all}\PYG{+w}{ }\PYG{n}{ave}\PYG{o}{/}\PYG{n}{time}\PYG{+w}{ }\PYG{l+m}{100}\PYG{+w}{ }\PYG{l+m}{1}\PYG{+w}{ }\PYG{l+m}{100}\PYG{+w}{ }\PYG{n}{c\PYGZus{}myChunk1}\PYG{+w}{ }\PYG{n}{c\PYGZus{}myChunk2[}\PYG{o}{*}\PYG{n}{]}\PYG{+w}{ }\PYG{n}{file}\PYG{+w}{ }\PYG{n}{tmp.out}\PYG{+w}{ }\PYG{n}{mode}\PYG{+w}{ }\PYG{n}{vector}
\end{sphinxVerbatim}


\subsection{Output info}
\label{\detokenize{compute_property_chunk:output-info}}
\sphinxAtStartPar
This compute calculates a global vector or global array depending on
the number of input values.  The length of the vector or number of
rows in the array is the number of chunks.

\sphinxAtStartPar
This compute calculates a global vector or global array where the
number of rows = the number of chunks \sphinxstyleemphasis{Nchunk} as calculated by the
specified {\hyperref[\detokenize{compute_chunk_atom::doc}]{\sphinxcrossref{\DUrole{doc}{compute chunk/atom}}}} command.  If a
single input is specified, a global vector is produced.  If two or
more inputs are specified, a global array is produced where the number
of columns = the number of inputs.  The vector or array can be
accessed by any command that uses global values from a compute as
input.  See the {\hyperref[\detokenize{Howto_output::doc}]{\sphinxcrossref{\DUrole{doc}{Howto output}}}} page for an
overview of LAMMPS output options.

\sphinxAtStartPar
The vector or array values are “intensive”.  The values will be
unitless or in the units discussed above.


\subsection{Restrictions}
\label{\detokenize{compute_property_chunk:restrictions}}\begin{quote}

\sphinxAtStartPar
none
\end{quote}


\subsection{Related commands}
\label{\detokenize{compute_property_chunk:related-commands}}
\sphinxAtStartPar
{\hyperref[\detokenize{fix_ave_chunk::doc}]{\sphinxcrossref{\DUrole{doc}{fix ave/chunk}}}}


\subsection{Default}
\label{\detokenize{compute_property_chunk:default}}
\sphinxAtStartPar
none

\sphinxstepscope

\index{compute property/grid@\spxentry{compute property/grid}}\ignorespaces 

\section{compute property/grid command}
\label{\detokenize{compute_property_grid:compute-property-grid-command}}\label{\detokenize{compute_property_grid:index-0}}\label{\detokenize{compute_property_grid::doc}}

\subsection{Syntax}
\label{\detokenize{compute_property_grid:syntax}}
\begin{sphinxVerbatim}[commandchars=\\\{\}]
\PYG{k}{compute }\PYG{n+nv+nvIdentifier}{ID}\PYG{+w}{ }\PYG{n+nv+nvIdentifier}{group\PYGZhy{}ID}\PYG{+w}{ }\PYG{n}{property}\PYG{o}{/}\PYG{n}{grid}\PYG{+w}{ }\PYG{n}{Nx}\PYG{+w}{ }\PYG{n}{Ny}\PYG{+w}{ }\PYG{n}{Nz}\PYG{+w}{ }\PYG{n}{input1}\PYG{+w}{ }\PYG{n}{input2}\PYG{+w}{ }\PYG{n}{...}
\end{sphinxVerbatim}
\begin{itemize}
\item {} 
\sphinxAtStartPar
ID, group\sphinxhyphen{}ID are documented in {\hyperref[\detokenize{compute::doc}]{\sphinxcrossref{\DUrole{doc}{compute}}}} command

\item {} 
\sphinxAtStartPar
property/grid = style name of this compute command

\item {} 
\sphinxAtStartPar
Nx, Ny, Nz = grid size in each dimension

\item {} 
\sphinxAtStartPar
input1,etc = one or more attributes

\begin{sphinxVerbatim}[commandchars=\\\{\}]
attributes = id, ix, iy, iz, x, y, z, xs, ys, zs, xc, yc, zc, xsc, ysc, zsc
  id = ID of grid cell, x fastest, y next, z slowest
  proc = processor ID (0 to Nprocs\PYGZhy{}1) which owns the grid cell
  ix,iy,iz = grid indices in each dimension (1 to N inclusive)
  x,y,z = coords of lower left corner of grid cell
  xs,ys,zs = scaled coords of lower left corner of grid cell (0.0 to 1.0)
  xc,yc,zc = coords of center point of grid cell
  xsc,ysc,zsc = scaled coords of center point of grid cell (0.0 to 1.0)
\end{sphinxVerbatim}

\end{itemize}


\subsection{Examples}
\label{\detokenize{compute_property_grid:examples}}
\begin{sphinxVerbatim}[commandchars=\\\{\}]
\PYG{k}{compute }\PYG{n+nv+nvIdentifier}{1}\PYG{+w}{ }\PYG{n+nv+nvIdentifier}{all}\PYG{+w}{ }\PYG{n}{property}\PYG{o}{/}\PYG{n}{grid}\PYG{+w}{ }\PYG{l+m}{10}\PYG{+w}{ }\PYG{l+m}{10}\PYG{+w}{ }\PYG{l+m}{20}\PYG{+w}{ }\PYG{n}{id}\PYG{+w}{ }\PYG{n}{ix}\PYG{+w}{ }\PYG{n}{iy}\PYG{+w}{ }\PYG{n}{iz}
\PYG{k}{compute }\PYG{n+nv+nvIdentifier}{1}\PYG{+w}{ }\PYG{n+nv+nvIdentifier}{all}\PYG{+w}{ }\PYG{n}{property}\PYG{o}{/}\PYG{n}{grid}\PYG{+w}{ }\PYG{l+m}{100}\PYG{+w}{ }\PYG{l+m}{100}\PYG{+w}{ }\PYG{l+m}{1}\PYG{+w}{ }\PYG{n}{id}\PYG{+w}{ }\PYG{n}{xc}\PYG{+w}{ }\PYG{n}{yc}\PYG{+w}{ }\PYG{n}{zc}
\end{sphinxVerbatim}


\subsection{Description}
\label{\detokenize{compute_property_grid:description}}
\sphinxAtStartPar
Define a computation that stores the specified attributes of a
distributed grid.  In LAMMPS, distributed grids are regular 2d or 3d
grids which overlay a 2d or 3d simulation domain.  Each processor owns
the grid cells whose center points lie within its subdomain.  See the
{\hyperref[\detokenize{Howto_grid::doc}]{\sphinxcrossref{\DUrole{doc}{Howto grid}}}} doc page for details of how distributed
grids can be defined by various commands and referenced.

\sphinxAtStartPar
This compute stores the specified attributes of grids as per\sphinxhyphen{}grid data
so they can be accessed by other {\hyperref[\detokenize{Howto_output::doc}]{\sphinxcrossref{\DUrole{doc}{output commands}}}}
such as {\hyperref[\detokenize{dump::doc}]{\sphinxcrossref{\DUrole{doc}{dump grid}}}}.

\sphinxAtStartPar
\sphinxstyleemphasis{Nx}, \sphinxstyleemphasis{Ny}, and \sphinxstyleemphasis{Nz} define the size of the grid.  For a 2d simulation
\sphinxstyleemphasis{Nz} must be 1.  When this compute is used by {\hyperref[\detokenize{dump::doc}]{\sphinxcrossref{\DUrole{doc}{dump grid}}}},
to output per\sphinxhyphen{}grid values from other computes of fixes, the grid size
specified for this command must be consistent with the grid sizes
used by the other commands.

\sphinxAtStartPar
The \sphinxstyleemphasis{id} attribute is the grid ID for each grid cell.  For a global
grid of size Nx by Ny by Nz (in 3d simulations) the grid IDs range
from 1 to Nx*Ny*Nz.  They are ordered with the X index of the 3d grid
varying fastest, then Y, then Z slowest.  For 2d grids (in 2d
simulations), the grid IDs range from 1 to Nx*Ny, with X varying
fastest and Y slowest.

\sphinxAtStartPar
\DUrole{versionmodified,added}{Added in version 21Nov2023.}

\sphinxAtStartPar
The \sphinxstyleemphasis{proc} attribute is the ID of the processor which owns the grid
cell.  Processor IDs range from 0 to Nprocs \sphinxhyphen{} 1, where Nprocs is the
number of processors the simulation is running on.  Each grid cell is
owned by a single processor.

\sphinxAtStartPar
The \sphinxstyleemphasis{ix}, \sphinxstyleemphasis{iy}, \sphinxstyleemphasis{iz} attributes are the indices of a grid cell in
each dimension.  They range from 1 to Nx inclusive in the X dimension,
and similar for Y and Z.

\sphinxAtStartPar
The \sphinxstyleemphasis{x}, \sphinxstyleemphasis{y}, \sphinxstyleemphasis{z} attributes are the coordinates of the lower left
corner point of each grid cell.

\sphinxAtStartPar
The \sphinxstyleemphasis{xs}, \sphinxstyleemphasis{ys}, \sphinxstyleemphasis{zs} attributes are also coordinates of the lower left
corner point of each grid cell, except in scaled coordinates, where
the lower\sphinxhyphen{}left corner of the entire simulation box is (0,0,0) and the
upper right corner is (1,1,1).

\sphinxAtStartPar
The \sphinxstyleemphasis{xc}, \sphinxstyleemphasis{yc}, \sphinxstyleemphasis{zc} attributes are the coordinates of the center
point of each grid cell.

\sphinxAtStartPar
The \sphinxstyleemphasis{xsc}, \sphinxstyleemphasis{ysc}, \sphinxstyleemphasis{zsc} attributes are also coordinates of the center
point each grid cell, except in scaled coordinates, where the
lower\sphinxhyphen{}left corner of the entire simulation box is (0,0,0) and the upper
right corner is (1,1,1).

\sphinxAtStartPar
For {\hyperref[\detokenize{Howto_triclinic::doc}]{\sphinxcrossref{\DUrole{doc}{triclinic simulation boxes}}}}, the grid
point coordinates for (x,y,z) and (xc,yc,zc) will reflect the
triclinic geometry.  For (xs,yz,zs) and (xsc,ysc,zsc), the coordinates
are the same for orthogonal versus triclinic boxes.


\subsection{Output info}
\label{\detokenize{compute_property_grid:output-info}}
\sphinxAtStartPar
This compute calculates a per\sphinxhyphen{}grid vector or array depending on the
number of input values.  The length of the vector or number of array
rows (distributed across all processors) is Nx * Ny * Nz.  For access
by other commands, the name of the single grid produced by this
command is “grid”.  The name of its per\sphinxhyphen{}grid data is “data”.

\sphinxAtStartPar
The (x,y,z) and (xc,yc,zc) coordinates are in distance {\hyperref[\detokenize{units::doc}]{\sphinxcrossref{\DUrole{doc}{units}}}}.


\subsection{Restrictions}
\label{\detokenize{compute_property_grid:restrictions}}
\sphinxAtStartPar
For 2d simulations, the attributes which refer to
the Z dimension cannot be used.


\subsection{Related commands}
\label{\detokenize{compute_property_grid:related-commands}}
\sphinxAtStartPar
{\hyperref[\detokenize{dump::doc}]{\sphinxcrossref{\DUrole{doc}{dump grid}}}}


\subsection{Default}
\label{\detokenize{compute_property_grid:default}}
\sphinxAtStartPar
none

\sphinxstepscope

\index{compute property/local@\spxentry{compute property/local}}\ignorespaces 

\section{compute property/local command}
\label{\detokenize{compute_property_local:compute-property-local-command}}\label{\detokenize{compute_property_local:index-0}}\label{\detokenize{compute_property_local::doc}}

\subsection{Syntax}
\label{\detokenize{compute_property_local:syntax}}
\begin{sphinxVerbatim}[commandchars=\\\{\}]
\PYG{k}{compute }\PYG{n+nv+nvIdentifier}{ID}\PYG{+w}{ }\PYG{n+nv+nvIdentifier}{group\PYGZhy{}ID}\PYG{+w}{ }\PYG{n}{property}\PYG{o}{/}\PYG{n}{local}\PYG{+w}{ }\PYG{n}{attribute1}\PYG{+w}{ }\PYG{n}{attribute2}\PYG{+w}{ }\PYG{n}{...}\PYG{+w}{ }\PYG{n}{keyword}\PYG{+w}{ }\PYG{n}{args}\PYG{+w}{ }\PYG{n}{...}
\end{sphinxVerbatim}
\begin{itemize}
\item {} 
\sphinxAtStartPar
ID, group\sphinxhyphen{}ID are documented in {\hyperref[\detokenize{compute::doc}]{\sphinxcrossref{\DUrole{doc}{compute}}}} command

\item {} 
\sphinxAtStartPar
property/local = style name of this compute command

\item {} 
\sphinxAtStartPar
one or more attributes of the same type (neighbor, pair, bond, angle,
dihedral, or improper) may be appended

\begin{sphinxVerbatim}[commandchars=\\\{\}]
possible attributes = natom1, natom2, ntype1, ntype2,
                      patom1, patom2, ptype1, ptype2,
                      batom1, batom2, btype,
                      aatom1, aatom2, aatom3, atype,
                      datom1, datom2, datom3, datom4, dtype,
                      iatom1, iatom2, iatom3, iatom4, itype
\end{sphinxVerbatim}
\begin{itemize}
\item {} 
\sphinxAtStartPar
Neighbor attributes

\end{itemize}

\begin{sphinxVerbatim}[commandchars=\\\{\}]
natom1, natom2 = store IDs of two atoms in each pair (within neighbor cutoff)
ntype1, ntype2 = store types of two atoms in each pair (within neighbor cutoff)
\end{sphinxVerbatim}
\begin{itemize}
\item {} 
\sphinxAtStartPar
Pair attributes

\end{itemize}

\begin{sphinxVerbatim}[commandchars=\\\{\}]
patom1, patom2 = store IDs of two atoms in each pair (within force cutoff)
ptype1, ptype2 = store types of two atoms in each pair (within force cutoff)
\end{sphinxVerbatim}
\begin{itemize}
\item {} 
\sphinxAtStartPar
Bond attributes

\end{itemize}

\begin{sphinxVerbatim}[commandchars=\\\{\}]
batom1, batom2 = store IDs of two atoms in each bond
btype = store bond type of each bond
\end{sphinxVerbatim}
\begin{itemize}
\item {} 
\sphinxAtStartPar
Angle attributes

\end{itemize}

\begin{sphinxVerbatim}[commandchars=\\\{\}]
aatom1, aatom2, aatom3 = store IDs of three atoms in each angle
atype = store angle type of each angle
\end{sphinxVerbatim}
\begin{itemize}
\item {} 
\sphinxAtStartPar
Dihedral attributes

\end{itemize}

\begin{sphinxVerbatim}[commandchars=\\\{\}]
datom1, datom2, datom3, datom4 = store IDs of 4 atoms in each dihedral
dtype = store dihedral type of each dihedral
\end{sphinxVerbatim}
\begin{itemize}
\item {} 
\sphinxAtStartPar
Improper attributes

\end{itemize}

\begin{sphinxVerbatim}[commandchars=\\\{\}]
iatom1, iatom2, iatom3, iatom4 = store IDs of 4 atoms in each improper
itype = store improper type of each improper
\end{sphinxVerbatim}

\item {} 
\sphinxAtStartPar
zero or more keyword/arg pairs may be appended

\item {} 
\sphinxAtStartPar
keyword = \sphinxstyleemphasis{cutoff}
\begin{sphinxalltt}
\sphinxstyleemphasis{cutoff} arg = \sphinxstyleemphasis{type} or \sphinxstyleemphasis{radius}
\end{sphinxalltt}

\end{itemize}


\subsection{Examples}
\label{\detokenize{compute_property_local:examples}}
\begin{sphinxVerbatim}[commandchars=\\\{\}]
\PYG{k}{compute }\PYG{n+nv+nvIdentifier}{1}\PYG{+w}{ }\PYG{n+nv+nvIdentifier}{all}\PYG{+w}{ }\PYG{n}{property}\PYG{o}{/}\PYG{n}{local}\PYG{+w}{ }\PYG{n}{btype}\PYG{+w}{ }\PYG{n}{batom1}\PYG{+w}{ }\PYG{n}{batom2}
\PYG{k}{compute }\PYG{n+nv+nvIdentifier}{1}\PYG{+w}{ }\PYG{n+nv+nvIdentifier}{all}\PYG{+w}{ }\PYG{n}{property}\PYG{o}{/}\PYG{n}{local}\PYG{+w}{ }\PYG{n}{atype}\PYG{+w}{ }\PYG{n}{aatom2}
\end{sphinxVerbatim}


\subsection{Description}
\label{\detokenize{compute_property_local:description}}
\sphinxAtStartPar
Define a computation that stores the specified attributes as local
data so it can be accessed by other {\hyperref[\detokenize{Howto_output::doc}]{\sphinxcrossref{\DUrole{doc}{output commands}}}}.  If the input attributes refer to bond information,
then the number of datums generated, aggregated across all processors,
equals the number of bonds in the system.  Ditto for pairs, angles,
etc.

\sphinxAtStartPar
If multiple attributes are specified then they must all generate the
same amount of information, so that the resulting local array has the
same number of rows for each column.  This means that only bond
attributes can be specified together, or angle attributes, etc.  Bond
and angle attributes cannot be mixed in the same compute
property/local command.

\sphinxAtStartPar
If the inputs are pair attributes, the local data is generated by
looping over the pairwise neighbor list.  Info about an individual
pairwise interaction will only be included if both atoms in the pair
are in the specified compute group.  For \sphinxstyleemphasis{natom1} and \sphinxstyleemphasis{natom2}, all
atom pairs in the neighbor list are considered (out to the neighbor
cutoff = force cutoff + {\hyperref[\detokenize{neighbor::doc}]{\sphinxcrossref{\DUrole{doc}{neighbor skin}}}}).  For \sphinxstyleemphasis{patom1}
and \sphinxstyleemphasis{patom2}, the distance between the atoms must be less than the
force cutoff distance for that pair to be included, as defined by the
{\hyperref[\detokenize{pair_style::doc}]{\sphinxcrossref{\DUrole{doc}{pair\_style}}}} and {\hyperref[\detokenize{pair_coeff::doc}]{\sphinxcrossref{\DUrole{doc}{pair\_coeff}}}}
commands.

\sphinxAtStartPar
The optional \sphinxstyleemphasis{cutoff} keyword determines how the force cutoff distance
for an interaction is determined for the \sphinxstyleemphasis{patom1} and \sphinxstyleemphasis{patom2}
attributes.  For the default setting of \sphinxstyleemphasis{type}, the pairwise cutoff
defined by the {\hyperref[\detokenize{pair_style::doc}]{\sphinxcrossref{\DUrole{doc}{pair\_style}}}} command for the types of
the two atoms is used.  For the \sphinxstyleemphasis{radius} setting, the sum of the radii
of the two particles is used as a cutoff.  For example, this is
appropriate for granular particles which only interact when they are
overlapping, as computed by {\hyperref[\detokenize{pair_gran::doc}]{\sphinxcrossref{\DUrole{doc}{granular pair styles}}}}.
Note that if a granular model defines atom types such that all
particles of a specific type are monodisperse (same diameter), then
the two settings are effectively identical.

\sphinxAtStartPar
If the inputs are bond, angle, etc attributes, the local data is
generated by looping over all the atoms owned on a processor and
extracting bond, angle, etc info.  For bonds, info about an individual
bond will only be included if both atoms in the bond are in the
specified compute group.  Likewise for angles, dihedrals, etc.

\sphinxAtStartPar
For bonds and angles, a bonds/angles that have been broken by setting
their bond/angle type to 0 will not be included.  Bonds/angles that
have been turned off (see the {\hyperref[\detokenize{fix_shake::doc}]{\sphinxcrossref{\DUrole{doc}{fix shake}}}} or
{\hyperref[\detokenize{delete_bonds::doc}]{\sphinxcrossref{\DUrole{doc}{delete\_bonds}}}} commands) by setting their
bond/angle type negative are written into the file.  This is
consistent with the {\hyperref[\detokenize{compute_bond_local::doc}]{\sphinxcrossref{\DUrole{doc}{compute bond/local}}}} and
{\hyperref[\detokenize{compute_angle_local::doc}]{\sphinxcrossref{\DUrole{doc}{compute angle/local}}}} commands

\sphinxAtStartPar
Note that as atoms migrate from processor to processor, there will be
no consistent ordering of the entries within the local vector or array
from one timestep to the next.  The only consistency that is
guaranteed is that the ordering on a particular timestep will be the
same for local vectors or arrays generated by other compute commands.
For example, output from the {\hyperref[\detokenize{compute_bond_local::doc}]{\sphinxcrossref{\DUrole{doc}{compute bond/local}}}} command can be combined with bond atom indices
from this command and output by the {\hyperref[\detokenize{dump::doc}]{\sphinxcrossref{\DUrole{doc}{dump local}}}} command
in a consistent way.

\sphinxAtStartPar
The \sphinxstyleemphasis{natom1} and \sphinxstyleemphasis{natom2} or \sphinxstyleemphasis{patom1} and \sphinxstyleemphasis{patom2} attributes refer
to the atom IDs of the two atoms in each pairwise interaction computed
by the {\hyperref[\detokenize{pair_style::doc}]{\sphinxcrossref{\DUrole{doc}{pair\_style}}}} command.  The \sphinxstyleemphasis{ntype1} and
\sphinxstyleemphasis{ntype2} or \sphinxstyleemphasis{ptype1} and \sphinxstyleemphasis{ptype2} attributes refer to the atom types
of the two atoms in each pairwise interaction.

\begin{sphinxadmonition}{note}{Note:}
\sphinxAtStartPar
For pairs, if two atoms \(I,J\) are involved in 1\textendash{}2, 1\textendash{}3, 1\textendash{}4
interactions within the molecular topology, their pairwise interaction
may be turned off, and thus they may not appear in the neighbor list,
and will not be part of the local data created by this command.  More
specifically, this may be true of \(I,J\) pairs with a weighting factor
of 0.0; pairs with a non\sphinxhyphen{}zero weighting factor are included.  The
weighting factors for 1\textendash{}2, 1\textendash{}3, and 1\textendash{}4 pairwise interactions are set
by the {\hyperref[\detokenize{special_bonds::doc}]{\sphinxcrossref{\DUrole{doc}{special\_bonds}}}} command.
\end{sphinxadmonition}

\sphinxAtStartPar
The \sphinxstyleemphasis{batom1} and \sphinxstyleemphasis{batom2} attributes refer to the atom IDs of the 2
atoms in each {\hyperref[\detokenize{bond_style::doc}]{\sphinxcrossref{\DUrole{doc}{bond}}}}.  The \sphinxstyleemphasis{btype} attribute refers to
the type of the bond, from 1 to Nbtypes = \# of bond types.  The number
of bond types is defined in the data file read by the
{\hyperref[\detokenize{read_data::doc}]{\sphinxcrossref{\DUrole{doc}{read\_data}}}} command.

\sphinxAtStartPar
The attributes that start with “a”, “d”, and “i” refer to similar values
for {\hyperref[\detokenize{angle_style::doc}]{\sphinxcrossref{\DUrole{doc}{angles}}}}, {\hyperref[\detokenize{dihedral_style::doc}]{\sphinxcrossref{\DUrole{doc}{dihedrals}}}}, and
{\hyperref[\detokenize{improper_style::doc}]{\sphinxcrossref{\DUrole{doc}{impropers}}}}.


\subsection{Output info}
\label{\detokenize{compute_property_local:output-info}}
\sphinxAtStartPar
This compute calculates a local vector or local array depending on the
number of input values.  The length of the vector or number of rows in
the array is the number of bonds, angles, etc.  If a single input is
specified, a local vector is produced.  If two or more inputs are
specified, a local array is produced where the number of columns = the
number of inputs.  The vector or array can be accessed by any command
that uses local values from a compute as input.  See the {\hyperref[\detokenize{Howto_output::doc}]{\sphinxcrossref{\DUrole{doc}{Howto
output}}}} page for an overview of LAMMPS output options.

\sphinxAtStartPar
The vector or array values will be integers that correspond to the
specified attribute.


\subsection{Restrictions}
\label{\detokenize{compute_property_local:restrictions}}\begin{quote}

\sphinxAtStartPar
none
\end{quote}


\subsection{Related commands}
\label{\detokenize{compute_property_local:related-commands}}
\sphinxAtStartPar
{\hyperref[\detokenize{dump::doc}]{\sphinxcrossref{\DUrole{doc}{dump local}}}}, {\hyperref[\detokenize{compute_reduce::doc}]{\sphinxcrossref{\DUrole{doc}{compute reduce}}}}


\subsection{Default}
\label{\detokenize{compute_property_local:default}}
\sphinxAtStartPar
The keyword default is cutoff = type.

\sphinxstepscope

\index{compute ptm/atom@\spxentry{compute ptm/atom}}\ignorespaces 

\section{compute ptm/atom command}
\label{\detokenize{compute_ptm_atom:compute-ptm-atom-command}}\label{\detokenize{compute_ptm_atom:index-0}}\label{\detokenize{compute_ptm_atom::doc}}

\subsection{Syntax}
\label{\detokenize{compute_ptm_atom:syntax}}
\begin{sphinxVerbatim}[commandchars=\\\{\}]
\PYG{k}{compute }\PYG{n+nv+nvIdentifier}{ID}\PYG{+w}{ }\PYG{n+nv+nvIdentifier}{group\PYGZhy{}ID}\PYG{+w}{ }\PYG{n}{ptm}\PYG{o}{/}\PYG{n}{atom}\PYG{+w}{ }\PYG{n}{structures}\PYG{+w}{ }\PYG{n}{threshold}\PYG{+w}{ }\PYG{n}{group2}\PYG{o}{\PYGZhy{}}\PYG{n}{ID}
\end{sphinxVerbatim}
\begin{itemize}
\item {} 
\sphinxAtStartPar
ID, group\sphinxhyphen{}ID are documented in {\hyperref[\detokenize{compute::doc}]{\sphinxcrossref{\DUrole{doc}{compute}}}} command

\item {} 
\sphinxAtStartPar
ptm/atom = style name of this compute command

\item {} 
\sphinxAtStartPar
structures = \sphinxstyleemphasis{default} or \sphinxstyleemphasis{all} or any hyphen\sphinxhyphen{}separated combination of \sphinxstyleemphasis{fcc}, \sphinxstyleemphasis{hcp}, \sphinxstyleemphasis{bcc}, \sphinxstyleemphasis{ico}, \sphinxstyleemphasis{sc}, \sphinxstyleemphasis{dcub}, \sphinxstyleemphasis{dhex}, or \sphinxstyleemphasis{graphene} = structure types to search for

\item {} 
\sphinxAtStartPar
threshold = lattice distortion threshold (RMSD)

\item {} 
\sphinxAtStartPar
group2\sphinxhyphen{}ID determines which group is used for neighbor selection (optional, default “all”)

\end{itemize}


\subsection{Examples}
\label{\detokenize{compute_ptm_atom:examples}}
\begin{sphinxVerbatim}[commandchars=\\\{\}]
\PYG{k}{compute }\PYG{n+nv+nvIdentifier}{1}\PYG{+w}{ }\PYG{n+nv+nvIdentifier}{all}\PYG{+w}{ }\PYG{n}{ptm}\PYG{o}{/}\PYG{n}{atom}\PYG{+w}{ }\PYG{n}{default}\PYG{+w}{ }\PYG{l+m}{0.1}\PYG{+w}{ }\PYG{n}{all}
\PYG{k}{compute }\PYG{n+nv+nvIdentifier}{1}\PYG{+w}{ }\PYG{n+nv+nvIdentifier}{all}\PYG{+w}{ }\PYG{n}{ptm}\PYG{o}{/}\PYG{n}{atom}\PYG{+w}{ }\PYG{n}{fcc}\PYG{o}{\PYGZhy{}}\PYG{n}{hcp}\PYG{o}{\PYGZhy{}}\PYG{n}{dcub}\PYG{o}{\PYGZhy{}}\PYG{n}{dhex}\PYG{+w}{ }\PYG{l+m}{0.15}\PYG{+w}{ }\PYG{n}{all}
\PYG{k}{compute }\PYG{n+nv+nvIdentifier}{1}\PYG{+w}{ }\PYG{n+nv+nvIdentifier}{all}\PYG{+w}{ }\PYG{n}{ptm}\PYG{o}{/}\PYG{n}{atom}\PYG{+w}{ }\PYG{n}{all}\PYG{+w}{ }\PYG{l+m}{0}
\end{sphinxVerbatim}


\subsection{Description}
\label{\detokenize{compute_ptm_atom:description}}
\sphinxAtStartPar
Define a computation that determines the local lattice structure
around an atom using the PTM (Polyhedral Template Matching) method.
The PTM method is described in {\hyperref[\detokenize{compute_ptm_atom:larsen}]{\sphinxcrossref{\DUrole{std,std-ref}{(Larsen)}}}}.

\sphinxAtStartPar
Currently, there are seven lattice structures PTM recognizes:
\begin{itemize}
\item {} 
\sphinxAtStartPar
fcc = 1

\item {} 
\sphinxAtStartPar
hcp = 2

\item {} 
\sphinxAtStartPar
bcc = 3

\item {} 
\sphinxAtStartPar
ico (icosahedral) = 4

\item {} 
\sphinxAtStartPar
sc (simple cubic) = 5

\item {} 
\sphinxAtStartPar
dcub (diamond cubic) = 6

\item {} 
\sphinxAtStartPar
dhex (diamond hexagonal) = 7

\item {} 
\sphinxAtStartPar
graphene = 8

\end{itemize}

\sphinxAtStartPar
The value of the PTM structure will be 0 for unknown types and \(-1\) for
atoms not in the specified compute group.  The choice of structures to search
for can be specified using the “structures” argument, which is a
hyphen\sphinxhyphen{}separated list of structure keywords.
Two convenient pre\sphinxhyphen{}set options are provided:
\begin{itemize}
\item {} 
\sphinxAtStartPar
default: fcc\sphinxhyphen{}hcp\sphinxhyphen{}bcc\sphinxhyphen{}ico

\item {} 
\sphinxAtStartPar
all: fcc\sphinxhyphen{}hcp\sphinxhyphen{}bcc\sphinxhyphen{}ico\sphinxhyphen{}sc\sphinxhyphen{}dcub\sphinxhyphen{}dhex\sphinxhyphen{}graphene

\end{itemize}

\sphinxAtStartPar
The ‘default’ setting detects the same structures as the Common Neighbor Analysis method.
The ‘all’ setting searches for all structure types.  A performance penalty is
incurred for the diamond and graphene structures, so it is not recommended to use this option if
it is known that the simulation does not contain these structures.

\sphinxAtStartPar
PTM identifies structures using two steps.  First, a graph isomorphism test is used
to identify potential structure matches.  Next, the deviation is computed between the
local structure (in the simulation) and a template of the ideal lattice structure.
The deviation is calculated as:
\begin{equation*}
\begin{split}\text{RMSD}(\mathbf{u}, \mathbf{v})
 = \min_{s, \mathbf{Q}} \sqrt{\frac{1}{N} \sum\limits_{i=1}^{N}
{\left\lVert s[\vec{u}_i - \mathbf{\bar{u}}]
              - \mathbf{Q} \cdot \vec{v}_i \right\rVert}^2}\end{split}
\end{equation*}
\sphinxAtStartPar
Here, \(\vec u\) and \(\vec v\) contain the coordinates of the local
and ideal structures respectively, \(s\) is a scale factor, and
\(\mathbf Q\) is a rotation.  The best match is identified by the lowest
RMSD value, using the optimal scaling, rotation, and correspondence between the
points.

\sphinxAtStartPar
The \sphinxstyleemphasis{threshold} keyword sets an upper limit on the maximum permitted deviation
before a local structure is identified as disordered.  Typical values are in
the range 0.1\textendash{}0.15, but larger values may be desirable at higher temperatures.
A value of 0 is equivalent to infinity and can be used if no threshold is
desired.

\sphinxAtStartPar
The neighbor list needed to compute this quantity is constructed each
time the calculation is performed (e.g., each time a snapshot of atoms
is dumped).  Thus it can be inefficient to compute/dump this quantity
too frequently or to have multiple compute/dump commands, each with a
\sphinxstyleemphasis{ptm/atom} style. By default the compute processes \sphinxstylestrong{all} neighbors
unless the optional \sphinxstyleemphasis{group2\sphinxhyphen{}ID} argument is given, then only members
of that group are considered as neighbors.


\subsection{Output info}
\label{\detokenize{compute_ptm_atom:output-info}}
\sphinxAtStartPar
This compute calculates a per\sphinxhyphen{}atom array, which can be accessed by
any command that uses per\sphinxhyphen{}atom values from a compute as input.  See
the {\hyperref[\detokenize{Howto_output::doc}]{\sphinxcrossref{\DUrole{doc}{Howto output}}}} page for an overview of
LAMMPS output options.

\sphinxAtStartPar
Results are stored in the per\sphinxhyphen{}atom array in the following order:
\begin{itemize}
\item {} 
\sphinxAtStartPar
type

\item {} 
\sphinxAtStartPar
rmsd

\item {} 
\sphinxAtStartPar
interatomic distance

\item {} 
\sphinxAtStartPar
qw

\item {} 
\sphinxAtStartPar
qx

\item {} 
\sphinxAtStartPar
qy

\item {} 
\sphinxAtStartPar
qz

\end{itemize}

\sphinxAtStartPar
The type is a number from \(-1\) to 8. The rmsd is a positive real number.
The interatomic distance is computed from the scale factor in the RMSD equation.
The \((qw,qx,qy,qz)\) parameters represent the orientation of the local
structure in quaternion form.  The reference coordinates for each template
(from which the orientation is determined) can be found in the
\sphinxstyleemphasis{ptm\_constants.h} file in the PTM source directory.
For atoms that are not within the compute group\sphinxhyphen{}ID, all values are set to zero.


\subsection{Restrictions}
\label{\detokenize{compute_ptm_atom:restrictions}}
\sphinxAtStartPar
This fix is part of the PTM package.  It is only enabled if
LAMMPS was built with that package.  See the {\hyperref[\detokenize{Build_package::doc}]{\sphinxcrossref{\DUrole{doc}{Build package}}}} page for more info.


\subsection{Related commands}
\label{\detokenize{compute_ptm_atom:related-commands}}
\sphinxAtStartPar
{\hyperref[\detokenize{compute_centro_atom::doc}]{\sphinxcrossref{\DUrole{doc}{compute centro/atom}}}}
{\hyperref[\detokenize{compute_cna_atom::doc}]{\sphinxcrossref{\DUrole{doc}{compute cna/atom}}}}


\subsection{Default}
\label{\detokenize{compute_ptm_atom:default}}
\sphinxAtStartPar
none


\bigskip\hrule\bigskip

\phantomsection\label{\detokenize{compute_ptm_atom:larsen}}
\sphinxAtStartPar
\sphinxstylestrong{(Larsen)} Larsen, Schmidt, Schiotz, Modelling Simul Mater Sci Eng, 24, 055007 (2016).

\sphinxstepscope

\index{compute rattlers/atom@\spxentry{compute rattlers/atom}}\ignorespaces 

\section{compute rattlers/atom command}
\label{\detokenize{compute_rattlers_atom:compute-rattlers-atom-command}}\label{\detokenize{compute_rattlers_atom:index-0}}\label{\detokenize{compute_rattlers_atom::doc}}

\subsection{Syntax}
\label{\detokenize{compute_rattlers_atom:syntax}}
\begin{sphinxVerbatim}[commandchars=\\\{\}]
compute ID group\PYGZhy{}ID rattlers/atom cutoff zmin ntries
\end{sphinxVerbatim}
\begin{itemize}
\item {} 
\sphinxAtStartPar
ID, group\sphinxhyphen{}ID are documented in {\hyperref[\detokenize{compute::doc}]{\sphinxcrossref{\DUrole{doc}{compute}}}} command

\item {} 
\sphinxAtStartPar
rattlers/atom = style name of this compute command

\item {} 
\sphinxAtStartPar
cutoff = \sphinxstyleemphasis{type} or \sphinxstyleemphasis{radius}
\begin{sphinxalltt}
\sphinxstyleemphasis{type} = cutoffs determined based on atom types
\sphinxstyleemphasis{radius} = cutoffs determined based on atom diameters (atom style sphere)
\end{sphinxalltt}

\item {} 
\sphinxAtStartPar
zmin = minimum coordination for a non\sphinxhyphen{}rattler atom

\item {} 
\sphinxAtStartPar
ntries = maximum number of iterations to remove rattlers

\end{itemize}


\subsection{Examples}
\label{\detokenize{compute_rattlers_atom:examples}}
\begin{sphinxVerbatim}[commandchars=\\\{\}]
\PYG{k}{compute }\PYG{n+nv+nvIdentifier}{1}\PYG{+w}{ }\PYG{n+nv+nvIdentifier}{all}\PYG{+w}{ }\PYG{n}{rattlers}\PYG{o}{/}\PYG{n}{atom}\PYG{+w}{ }\PYG{n}{type}\PYG{+w}{ }\PYG{l+m}{4}\PYG{+w}{ }\PYG{l+m}{10}
\end{sphinxVerbatim}


\subsection{Description}
\label{\detokenize{compute_rattlers_atom:description}}
\sphinxAtStartPar
\DUrole{versionmodified,added}{Added in version 7Feb2024.}

\sphinxAtStartPar
Define a compute that identifies rattlers in a system. Rattlers are often
identified in granular or glassy packings as under\sphinxhyphen{}coordinated atoms that
do not have the required number of contacts to constrain their translational
degrees of freedom. Such atoms are not considered rigid and can often freely
rattle around in the system. This compute identifies rattlers which can be
helpful for excluding them from analysis or providing extra damping forces
to accelerate relaxation processes.

\sphinxAtStartPar
Rattlers are identified using an interactive approach. The coordination
number of all atoms is first calculated.  The \sphinxstyleemphasis{type} and \sphinxstyleemphasis{radius} settings
are used to select whether interaction cutoffs are determined by atom
types or by the sum of atomic radii (atom style sphere), respectively.
Rattlers are then identified as atoms with a coordination number less
than \sphinxstyleemphasis{zmin} and are removed from consideration. Atomic coordination
numbers are then recalculated, excluding previously identified rattlers,
to identify a new set of rattlers. This process is iterated up to a maximum
of \sphinxstyleemphasis{ntries} or until no new rattlers are identified and the remaining
atoms form a stable network of contacts.

\sphinxAtStartPar
In dense homogeneous systems where the average atom coordination number
is expected to be larger than \sphinxstyleemphasis{zmin}, this process usually only takes a few
iterations and a value of \sphinxstyleemphasis{ntries} around ten may be sufficient. In systems
with significant heterogeneity or average coordination numbers less than
\sphinxstyleemphasis{zmin}, an appropriate value of \sphinxstyleemphasis{ntries} depends heavily on the specific
system. For instance, a linear chain of N rattler atoms with a \sphinxstyleemphasis{zmin} of 2
would take N/2 iterations to identify that all the atoms are rattlers.


\subsection{Output info}
\label{\detokenize{compute_rattlers_atom:output-info}}
\sphinxAtStartPar
This compute calculates a per\sphinxhyphen{}atom vector and a global scalar. The vector
designates which atoms are rattlers, indicated by a value 1. Non\sphinxhyphen{}rattlers
have a value of 0. The global scalar returns the total number of rattlers
in the system. See the {\hyperref[\detokenize{Howto_output::doc}]{\sphinxcrossref{\DUrole{doc}{Howto output}}}} page for an
overview of LAMMPS output options.


\subsection{Restrictions}
\label{\detokenize{compute_rattlers_atom:restrictions}}
\sphinxAtStartPar
This compute is part of the EXTRA\sphinxhyphen{}COMPUTE package.  It is only enabled if
LAMMPS was built with that package.  See the
{\hyperref[\detokenize{Build_package::doc}]{\sphinxcrossref{\DUrole{doc}{Build package}}}} page for more info.

\sphinxAtStartPar
The \sphinxstyleemphasis{radius} cutoff option requires that atoms store a radius as defined by the
{\hyperref[\detokenize{atom_style::doc}]{\sphinxcrossref{\DUrole{doc}{atom\_style sphere}}}} or similar commands.


\subsection{Related commands}
\label{\detokenize{compute_rattlers_atom:related-commands}}
\sphinxAtStartPar
{\hyperref[\detokenize{compute_coord_atom::doc}]{\sphinxcrossref{\DUrole{doc}{compute coord/atom}}}}
{\hyperref[\detokenize{compute_contact_atom::doc}]{\sphinxcrossref{\DUrole{doc}{compute contact/atom}}}}


\subsection{Default}
\label{\detokenize{compute_rattlers_atom:default}}
\sphinxAtStartPar
none

\sphinxstepscope

\index{compute rdf@\spxentry{compute rdf}}\ignorespaces 

\section{compute rdf command}
\label{\detokenize{compute_rdf:compute-rdf-command}}\label{\detokenize{compute_rdf:index-0}}\label{\detokenize{compute_rdf::doc}}

\subsection{Syntax}
\label{\detokenize{compute_rdf:syntax}}
\begin{sphinxVerbatim}[commandchars=\\\{\}]
\PYG{k}{compute }\PYG{n+nv+nvIdentifier}{ID}\PYG{+w}{ }\PYG{n+nv+nvIdentifier}{group\PYGZhy{}ID}\PYG{+w}{ }\PYG{n}{rdf}\PYG{+w}{ }\PYG{n}{Nbin}\PYG{+w}{ }\PYG{n}{itype1}\PYG{+w}{ }\PYG{n}{jtype1}\PYG{+w}{ }\PYG{n}{itype2}\PYG{+w}{ }\PYG{n}{jtype2}\PYG{+w}{ }\PYG{n}{...}\PYG{+w}{ }\PYG{n}{keyword}\PYG{o}{/}\PYG{n}{value}\PYG{+w}{ }\PYG{n}{...}
\end{sphinxVerbatim}
\begin{itemize}
\item {} 
\sphinxAtStartPar
ID, group\sphinxhyphen{}ID are documented in {\hyperref[\detokenize{compute::doc}]{\sphinxcrossref{\DUrole{doc}{compute}}}} command

\item {} 
\sphinxAtStartPar
rdf = style name of this compute command

\item {} 
\sphinxAtStartPar
Nbin = number of RDF bins

\item {} 
\sphinxAtStartPar
itypeN = central atom type for Nth RDF histogram (integer, type label, or asterisk form)

\item {} 
\sphinxAtStartPar
jtypeN = distribution atom type for Nth RDF histogram (integer, type label, or asterisk form)

\item {} 
\sphinxAtStartPar
zero or more keyword/value pairs may be appended

\item {} 
\sphinxAtStartPar
keyword = \sphinxstyleemphasis{cutoff}
\begin{sphinxalltt}
\sphinxstyleemphasis{cutoff} value = Rcut
  Rcut = cutoff distance for RDF computation (distance units)
\end{sphinxalltt}

\end{itemize}


\subsection{Examples}
\label{\detokenize{compute_rdf:examples}}
\begin{sphinxVerbatim}[commandchars=\\\{\}]
\PYG{k}{compute }\PYG{n+nv+nvIdentifier}{1}\PYG{+w}{ }\PYG{n+nv+nvIdentifier}{all}\PYG{+w}{ }\PYG{n}{rdf}\PYG{+w}{ }\PYG{l+m}{100}
\PYG{k}{compute }\PYG{n+nv+nvIdentifier}{1}\PYG{+w}{ }\PYG{n+nv+nvIdentifier}{all}\PYG{+w}{ }\PYG{n}{rdf}\PYG{+w}{ }\PYG{l+m}{100}\PYG{+w}{ }\PYG{l+m}{1}\PYG{+w}{ }\PYG{l+m}{1}
\PYG{k}{compute }\PYG{n+nv+nvIdentifier}{1}\PYG{+w}{ }\PYG{n+nv+nvIdentifier}{all}\PYG{+w}{ }\PYG{n}{rdf}\PYG{+w}{ }\PYG{l+m}{100}\PYG{+w}{ }\PYG{o}{*}\PYG{+w}{ }\PYG{l+m}{3}\PYG{+w}{ }\PYG{n}{cutoff}\PYG{+w}{ }\PYG{l+m}{5.0}
\PYG{k}{compute }\PYG{n+nv+nvIdentifier}{1}\PYG{+w}{ }\PYG{n+nv+nvIdentifier}{fluid}\PYG{+w}{ }\PYG{n}{rdf}\PYG{+w}{ }\PYG{l+m}{500}\PYG{+w}{ }\PYG{l+m}{1}\PYG{+w}{ }\PYG{l+m}{1}\PYG{+w}{ }\PYG{l+m}{1}\PYG{+w}{ }\PYG{l+m}{2}\PYG{+w}{ }\PYG{l+m}{2}\PYG{+w}{ }\PYG{l+m}{1}\PYG{+w}{ }\PYG{l+m}{2}\PYG{+w}{ }\PYG{l+m}{2}
\PYG{k}{compute }\PYG{n+nv+nvIdentifier}{1}\PYG{+w}{ }\PYG{n+nv+nvIdentifier}{fluid}\PYG{+w}{ }\PYG{n}{rdf}\PYG{+w}{ }\PYG{l+m}{500}\PYG{+w}{ }\PYG{l+m}{1}\PYG{o}{*}\PYG{l+m}{3}\PYG{+w}{ }\PYG{l+m}{2}\PYG{+w}{ }\PYG{l+m}{5}\PYG{+w}{ }\PYG{o}{*}\PYG{l+m}{10}\PYG{+w}{ }\PYG{n}{cutoff}\PYG{+w}{ }\PYG{l+m}{3.5}
\end{sphinxVerbatim}


\subsection{Description}
\label{\detokenize{compute_rdf:description}}
\sphinxAtStartPar
Define a computation that calculates the radial distribution function
(RDF), also called \(g(r)\), and the coordination number for a group of
particles.  Both are calculated in histogram form by binning pairwise
distances into \sphinxstyleemphasis{Nbin} bins from 0.0 to the maximum force cutoff
defined by the {\hyperref[\detokenize{pair_style::doc}]{\sphinxcrossref{\DUrole{doc}{pair\_style}}}} command or the cutoff
distance \sphinxstyleemphasis{Rcut} specified via the \sphinxstyleemphasis{cutoff} keyword.  The bins are of
uniform size in radial distance.  Thus a single bin encompasses a thin
shell of distances in 3d and a thin ring of distances in 2d.

\begin{sphinxadmonition}{note}{Note:}
\sphinxAtStartPar
If you have a bonded system, then the settings of
{\hyperref[\detokenize{special_bonds::doc}]{\sphinxcrossref{\DUrole{doc}{special\_bonds}}}} command can remove pairwise
interactions between atoms in the same bond, angle, or dihedral.  This
is the default setting for the {\hyperref[\detokenize{special_bonds::doc}]{\sphinxcrossref{\DUrole{doc}{special\_bonds}}}}
command, and means those pairwise interactions do not appear in the
neighbor list.  Because this fix uses a neighbor list, it also means
those pairs will not be included in the RDF. This does not apply when
using long\sphinxhyphen{}range coulomb interactions (\sphinxstyleemphasis{coul/long}, \sphinxstyleemphasis{coul/msm},
\sphinxstyleemphasis{coul/wolf} or similar.  One way to get around this would be to set
special\_bond scaling factors to very tiny numbers that are not exactly
zero (e.g., \(1.0 \times 10^{-50}\)).  Another workaround is to write a
dump file, and use the {\hyperref[\detokenize{rerun::doc}]{\sphinxcrossref{\DUrole{doc}{rerun}}}} command to compute the RDF for
snapshots in the dump file.  The rerun script can use a
{\hyperref[\detokenize{special_bonds::doc}]{\sphinxcrossref{\DUrole{doc}{special\_bonds}}}} command that includes all pairs in
the neighbor list.
\end{sphinxadmonition}

\sphinxAtStartPar
By default the RDF is computed out to the maximum force cutoff defined
by the {\hyperref[\detokenize{pair_style::doc}]{\sphinxcrossref{\DUrole{doc}{pair\_style}}}} command.  If the \sphinxstyleemphasis{cutoff} keyword
is used, then the RDF is computed accurately out to the \sphinxstyleemphasis{Rcut} \(> 0.0\)
distance specified.

\begin{sphinxadmonition}{note}{Note:}
\sphinxAtStartPar
Normally, you should only use the \sphinxstyleemphasis{cutoff} keyword if no pair
style is defined (e.g., the {\hyperref[\detokenize{rerun::doc}]{\sphinxcrossref{\DUrole{doc}{rerun}}}} command is being used to
post\sphinxhyphen{}process a dump file of snapshots) or if you really want the RDF
for distances beyond the pair\_style force cutoff and cannot easily
post\sphinxhyphen{}process a dump file to calculate it.  This is because using the
\sphinxstyleemphasis{cutoff} keyword incurs extra computation and possibly communication,
which may slow down your simulation.  If you specify \sphinxstyleemphasis{Rcut} \(\le\)
force cutoff, you will force an additional neighbor list to be built at
every timestep this command is invoked (or every reneighboring
timestep, whichever is less frequent), which is inefficient.  LAMMPS
will warn you if this is the case.  If you specify a \sphinxstyleemphasis{Rcut} \textgreater{} force
cutoff, you must ensure ghost atom information out to \sphinxstyleemphasis{Rcut} + \sphinxstyleemphasis{skin}
is communicated, via the {\hyperref[\detokenize{comm_modify::doc}]{\sphinxcrossref{\DUrole{doc}{comm\_modify cutoff}}}}
command, else the RDF computation cannot be performed, and LAMMPS will
give an error message.  The \sphinxstyleemphasis{skin} value is what is specified with the
{\hyperref[\detokenize{neighbor::doc}]{\sphinxcrossref{\DUrole{doc}{neighbor}}}} command.  In this case, you are forcing a
large neighbor list to be built just for the RDF computation, and
extra communication to be performed every timestep.
\end{sphinxadmonition}

\sphinxAtStartPar
The \sphinxstyleemphasis{itypeN} and \sphinxstyleemphasis{jtypeN} arguments are optional.  These arguments
must come in pairs.  If no pairs are listed, then a single histogram
is computed for \(g(r)\) between all atom types.  If one or more pairs are
listed, then a separate histogram is generated for each
\sphinxstyleemphasis{itype},\sphinxstyleemphasis{jtype} pair.

\sphinxAtStartPar
The \sphinxstyleemphasis{itypeN} and \sphinxstyleemphasis{jtypeN} settings can be specified in one of three
ways.  One or both of the types in the I,J pair can be a
{\hyperref[\detokenize{Howto_type_labels::doc}]{\sphinxcrossref{\DUrole{doc}{type label}}}}.  Or an explicit numeric value can be
used, as in the fourth example above. Or a wild\sphinxhyphen{}card asterisk can be used
to specify a range of atom types.  This takes the form “*” or “*n” or
“m*” or “m*n”. If \(N\) is the number of atom types, then an asterisk
with no numeric values means all types from 1 to \(N\). A leading
asterisk means all types from 1 to n (inclusive).  A trailing asterisk
means all types from m to \(N\) (inclusive).  A middle asterisk means
all types from m to n (inclusive).

\sphinxAtStartPar
If both \sphinxstyleemphasis{itypeN} and \sphinxstyleemphasis{jtypeN} are single values, as in the fourth example
above, this means that a \(g(r)\) is computed where atoms of type \sphinxstyleemphasis{itypeN}
are the central atom, and atoms of type \sphinxstyleemphasis{jtypeN} are the distribution
atom.  If either \sphinxstyleemphasis{itypeN} and \sphinxstyleemphasis{jtypeN} represent a range of values via
the wild\sphinxhyphen{}card asterisk, as in the fifth example above, this means that a
\(g(r)\) is computed where atoms of any of the range of types represented
by \sphinxstyleemphasis{itypeN} are the central atom, and atoms of any of the range of
types represented by \sphinxstyleemphasis{jtypeN} are the distribution atom.

\sphinxAtStartPar
Pairwise distances are generated by looping over a pairwise neighbor
list, just as they would be in a {\hyperref[\detokenize{pair_style::doc}]{\sphinxcrossref{\DUrole{doc}{pair\_style}}}}
computation.  The distance between two atoms \(I\) and \(J\) is
included in a specific histogram if the following criteria are met:
\begin{itemize}
\item {} 
\sphinxAtStartPar
atoms \(I\) and \(J\) are both in the specified compute group

\item {} 
\sphinxAtStartPar
the distance between atoms \(I\) and \(J\) is less than the maximum
force cutoff

\item {} 
\sphinxAtStartPar
the type of the \(I\) atom matches \sphinxstyleemphasis{itypeN} (one or a range of types)

\item {} 
\sphinxAtStartPar
the type of the \(J\) atom matches \sphinxstyleemphasis{jtypeN} (one or a range of types)

\end{itemize}

\sphinxAtStartPar
It is OK if a particular pairwise distance is included in more than
one individual histogram, due to the way the \sphinxstyleemphasis{itypeN} and \sphinxstyleemphasis{jtypeN}
arguments are specified.

\sphinxAtStartPar
The \(g(r)\) value for a bin is calculated from the histogram count by
scaling it by the idealized number of how many counts there would be
if atoms of type \sphinxstyleemphasis{jtypeN} were uniformly distributed.  Thus it
involves the count of \sphinxstyleemphasis{itypeN} atoms, the count of \sphinxstyleemphasis{jtypeN} atoms, the
volume of the entire simulation box, and the volume of the bin’s thin
shell in 3d (or the area of the bin’s thin ring in 2d).

\sphinxAtStartPar
A coordination number \(\mathrm{coord}(r)\) is also calculated, which is
the number of atoms of type \sphinxstyleemphasis{jtypeN} within the current bin or closer, averaged
over atoms of type \sphinxstyleemphasis{itypeN}.  This is calculated as the area\sphinxhyphen{} or
volume\sphinxhyphen{}weighted sum of \(g(r)\) values over all bins up to and including
the current bin, multiplied by the global average volume density of
atoms of type \sphinxstyleemphasis{jtypeN}.

\sphinxAtStartPar
The simplest way to output the results of the compute rdf calculation
to a file is to use the {\hyperref[\detokenize{fix_ave_time::doc}]{\sphinxcrossref{\DUrole{doc}{fix ave/time}}}} command, for
example:

\begin{sphinxVerbatim}[commandchars=\\\{\}]
\PYG{k}{compute }\PYG{n+nv+nvIdentifier}{myRDF}\PYG{+w}{ }\PYG{n+nv+nvIdentifier}{all}\PYG{+w}{ }\PYG{n}{rdf}\PYG{+w}{ }\PYG{l+m}{50}
\PYG{k}{fix }\PYG{n+nv+nvIdentifier}{1}\PYG{+w}{ }\PYG{n+nv+nvIdentifier}{all}\PYG{+w}{ }\PYG{n}{ave}\PYG{o}{/}\PYG{n}{time}\PYG{+w}{ }\PYG{l+m}{100}\PYG{+w}{ }\PYG{l+m}{1}\PYG{+w}{ }\PYG{l+m}{100}\PYG{+w}{ }\PYG{n}{c\PYGZus{}myRDF[}\PYG{o}{*}\PYG{n}{]}\PYG{+w}{ }\PYG{n}{file}\PYG{+w}{ }\PYG{n}{tmp.rdf}\PYG{+w}{ }\PYG{n}{mode}\PYG{+w}{ }\PYG{n}{vector}
\end{sphinxVerbatim}


\subsection{Output info}
\label{\detokenize{compute_rdf:output-info}}
\sphinxAtStartPar
This compute calculates a global array in which the number of rows is
\sphinxstyleemphasis{Nbins} and the number of columns is \(1 + 2N_\text{pairs}\), where
\(N_\text{pairs}\) is the number of \(I,J\) pairings specified.
The first column has the bin coordinate (center of the bin), and each
successive set of two columns has the \(g(r)\) and \(\text{coord}(r)\)
values for a specific set of \sphinxstyleemphasis{itypeN} versus \sphinxstyleemphasis{jtypeN} interactions,
as described above.  These values can be used
by any command that uses a global values from a compute as input.  See
the {\hyperref[\detokenize{Howto_output::doc}]{\sphinxcrossref{\DUrole{doc}{Howto output}}}} page for an overview of
LAMMPS output options.

\sphinxAtStartPar
The array values calculated by this compute are all “intensive”.

\sphinxAtStartPar
The first column of array values will be in distance
{\hyperref[\detokenize{units::doc}]{\sphinxcrossref{\DUrole{doc}{units}}}}.  The \(g(r)\) columns of array values are normalized
numbers \(\ge 0.0\).  The coordination number columns of array values are
also numbers \(\ge 0.0\).


\subsection{Restrictions}
\label{\detokenize{compute_rdf:restrictions}}
\sphinxAtStartPar
By default, the RDF is not computed for distances longer than the
largest force cutoff, since the neighbor list creation will only contain
pairs up to that distance (plus neighbor list skin).  This distance can
be increased using the \sphinxstyleemphasis{cutoff} keyword but this keyword is only valid
with {\hyperref[\detokenize{neighbor::doc}]{\sphinxcrossref{\DUrole{doc}{neighbor styles ‘bin’ and ‘nsq’}}}}.

\sphinxAtStartPar
If you want an RDF for larger distances, you can also use the
{\hyperref[\detokenize{rerun::doc}]{\sphinxcrossref{\DUrole{doc}{rerun}}}} command to post\sphinxhyphen{}process a dump file, use {\hyperref[\detokenize{pair_zero::doc}]{\sphinxcrossref{\DUrole{doc}{pair
style zero}}}} and set the force cutoff to be longer in the
rerun script.  Note that in the rerun context, the force cutoff is
arbitrary and with pair style zero you are not computing any forces, and
you are not running dynamics you are not changing the model that
generated the trajectory.

\sphinxAtStartPar
The definition of \(g(r)\) used by LAMMPS is only appropriate for
characterizing atoms that are uniformly distributed throughout the
simulation cell. In such cases, the coordination number is still correct
and meaningful.  As an example, if a large simulation cell contains only
one atom of type \sphinxstyleemphasis{itypeN} and one of \sphinxstyleemphasis{jtypeN}, then \(g(r)\) will
register an arbitrarily large spike at whatever distance they happen to
be at, and zero everywhere else.  The function \(\text{coord}(r)\)
will show a step change from zero to one at the location of the spike in
\(g(r)\).

\begin{sphinxadmonition}{note}{Note:}
\sphinxAtStartPar
compute rdf can handle dynamic groups and systems where atoms
are added or removed, but this causes that certain normalization
parameters need to be re\sphinxhyphen{}computed in every step and include collective
communication operations. This will reduce performance and limit
parallel efficiency and scaling. For systems, where only the type
of atoms changes (e.g., when using {\hyperref[\detokenize{fix_atom_swap::doc}]{\sphinxcrossref{\DUrole{doc}{fix atom/swap}}}}),
you need to explicitly request the dynamic normalization updates
via {\hyperref[\detokenize{compute_modify::doc}]{\sphinxcrossref{\DUrole{doc}{compute\_modify dynamic/dof yes}}}}
\end{sphinxadmonition}


\subsection{Related commands}
\label{\detokenize{compute_rdf:related-commands}}
\sphinxAtStartPar
{\hyperref[\detokenize{fix_ave_time::doc}]{\sphinxcrossref{\DUrole{doc}{fix ave/time}}}}, {\hyperref[\detokenize{compute_modify::doc}]{\sphinxcrossref{\DUrole{doc}{compute\_modify}}}},
{\hyperref[\detokenize{compute_adf::doc}]{\sphinxcrossref{\DUrole{doc}{compute adf}}}}


\subsection{Default}
\label{\detokenize{compute_rdf:default}}
\sphinxAtStartPar
The keyword defaults are cutoff = 0.0 (use the pairwise force cutoff).

\sphinxstepscope

\index{compute reaxff/atom@\spxentry{compute reaxff/atom}}\ignorespaces 
\index{compute reaxff/atom/kk@\spxentry{compute reaxff/atom/kk}}\ignorespaces 

\section{compute reaxff/atom command}
\label{\detokenize{compute_reaxff_atom:compute-reaxff-atom-command}}\label{\detokenize{compute_reaxff_atom:index-1}}\label{\detokenize{compute_reaxff_atom:index-0}}\label{\detokenize{compute_reaxff_atom::doc}}
\sphinxAtStartPar
Accelerator Variants: \sphinxstyleemphasis{reaxff/atom/kk}


\subsection{Syntax}
\label{\detokenize{compute_reaxff_atom:syntax}}
\begin{sphinxVerbatim}[commandchars=\\\{\}]
\PYG{k}{compute }\PYG{n+nv+nvIdentifier}{ID}\PYG{+w}{ }\PYG{n+nv+nvIdentifier}{group\PYGZhy{}ID}\PYG{+w}{ }\PYG{n}{reaxff}\PYG{o}{/}\PYG{n}{atom}\PYG{+w}{ }\PYG{n}{attribute}\PYG{+w}{ }\PYG{n}{args}\PYG{+w}{ }\PYG{n}{...}\PYG{+w}{ }\PYG{n}{keyword}\PYG{+w}{ }\PYG{n}{value}\PYG{+w}{ }\PYG{n}{...}
\end{sphinxVerbatim}
\begin{itemize}
\item {} 
\sphinxAtStartPar
ID, group\sphinxhyphen{}ID are documented in {\hyperref[\detokenize{compute::doc}]{\sphinxcrossref{\DUrole{doc}{compute}}}} command

\item {} 
\sphinxAtStartPar
reaxff/atom = name of this compute command

\item {} 
\sphinxAtStartPar
attribute = \sphinxstyleemphasis{pair}
\begin{sphinxalltt}
\sphinxstyleemphasis{pair} args = nsub
  nsub = \sphinxstyleemphasis{n}\sphinxhyphen{}instance of a sub\sphinxhyphen{}style, if a pair style is used multiple times in a hybrid style
\end{sphinxalltt}

\item {} 
\sphinxAtStartPar
keyword = \sphinxstyleemphasis{bonds}
\begin{sphinxalltt}
\sphinxstyleemphasis{bonds} value = \sphinxstyleemphasis{no} or \sphinxstyleemphasis{yes}
  \sphinxstyleemphasis{no} = ignore list of local bonds
  \sphinxstyleemphasis{yes} = include list of local bonds
\end{sphinxalltt}

\end{itemize}


\subsection{Examples}
\label{\detokenize{compute_reaxff_atom:examples}}
\begin{sphinxVerbatim}[commandchars=\\\{\}]
\PYG{k}{compute }\PYG{n+nv+nvIdentifier}{1}\PYG{+w}{ }\PYG{n+nv+nvIdentifier}{all}\PYG{+w}{ }\PYG{n}{reaxff}\PYG{o}{/}\PYG{n}{atom}\PYG{+w}{ }\PYG{n}{bonds}\PYG{+w}{ }\PYG{n}{yes}
\end{sphinxVerbatim}


\subsection{Description}
\label{\detokenize{compute_reaxff_atom:description}}
\sphinxAtStartPar
\DUrole{versionmodified,added}{Added in version 7Feb2024.}

\sphinxAtStartPar
Define a computation that extracts bond information computed by the ReaxFF
potential specified by {\hyperref[\detokenize{pair_reaxff::doc}]{\sphinxcrossref{\DUrole{doc}{pair\_style reaxff}}}}.

\sphinxAtStartPar
By default, it produces per\sphinxhyphen{}atom data that includes the following columns:
\begin{itemize}
\item {} 
\sphinxAtStartPar
abo = atom bond order (sum of all bonds)

\item {} 
\sphinxAtStartPar
nlp = number of lone pairs

\item {} 
\sphinxAtStartPar
nb = number of bonds

\end{itemize}

\sphinxAtStartPar
Bonds will only be included if its atoms are in the group.

\sphinxAtStartPar
In addition, if \sphinxcode{\sphinxupquote{bonds}} is set to \sphinxcode{\sphinxupquote{yes}}, the compute will also produce a
local array of all bonds on the current processor whose atoms are in the group.
The columns of each entry of this local array are:
\begin{itemize}
\item {} 
\sphinxAtStartPar
id\_i = atom i id of bond

\item {} 
\sphinxAtStartPar
id\_j = atom j id of bond

\item {} 
\sphinxAtStartPar
bo = bond order of bond

\end{itemize}


\subsection{Output info}
\label{\detokenize{compute_reaxff_atom:output-info}}
\sphinxAtStartPar
This compute calculates a per\sphinxhyphen{}atom array and local array depending on the
number of keywords. The number of rows in the local array is the number of
bonds as described above. Both per\sphinxhyphen{}atom and local array have 3 columns.

\sphinxAtStartPar
The arrays can be accessed by any command that uses local and per\sphinxhyphen{}atom values
from a compute as input.  See the {\hyperref[\detokenize{Howto_output::doc}]{\sphinxcrossref{\DUrole{doc}{Howto output}}}} page for
an overview of LAMMPS output options.


\bigskip\hrule\bigskip


\sphinxAtStartPar
Styles with a \sphinxstyleemphasis{gpu}, \sphinxstyleemphasis{intel}, \sphinxstyleemphasis{kk}, \sphinxstyleemphasis{omp}, or \sphinxstyleemphasis{opt} suffix are
functionally the same as the corresponding style without the suffix.
They have been optimized to run faster, depending on your available
hardware, as discussed on the {\hyperref[\detokenize{Speed_packages::doc}]{\sphinxcrossref{\DUrole{doc}{Accelerator packages}}}}
page.  The accelerated styles take the same arguments and should
produce the same results, except for round\sphinxhyphen{}off and precision issues.

\sphinxAtStartPar
These accelerated styles are part of the GPU, INTEL, KOKKOS,
OPENMP, and OPT packages, respectively.  They are only enabled if
LAMMPS was built with those packages.  See the {\hyperref[\detokenize{Build_package::doc}]{\sphinxcrossref{\DUrole{doc}{Build package}}}} page for more info.

\sphinxAtStartPar
You can specify the accelerated styles explicitly in your input script
by including their suffix, or you can use the {\hyperref[\detokenize{Run_options::doc}]{\sphinxcrossref{\DUrole{doc}{\sphinxhyphen{}suffix command\sphinxhyphen{}line switch}}}} when you invoke LAMMPS, or you can use the
{\hyperref[\detokenize{suffix::doc}]{\sphinxcrossref{\DUrole{doc}{suffix}}}} command in your input script.

\sphinxAtStartPar
See the {\hyperref[\detokenize{Speed_packages::doc}]{\sphinxcrossref{\DUrole{doc}{Accelerator packages}}}} page for more
instructions on how to use the accelerated styles effectively.


\bigskip\hrule\bigskip



\subsection{Restrictions}
\label{\detokenize{compute_reaxff_atom:restrictions}}
\sphinxAtStartPar
The compute reaxff/atom command requires that the {\hyperref[\detokenize{pair_reaxff::doc}]{\sphinxcrossref{\DUrole{doc}{pair\_style reaxff}}}} is invoked.  This fix is part of the REAXFF package.  It is only
enabled if LAMMPS was built with that package.  See the {\hyperref[\detokenize{Build_package::doc}]{\sphinxcrossref{\DUrole{doc}{Build package}}}} page for more info.


\subsection{Related commands}
\label{\detokenize{compute_reaxff_atom:related-commands}}
\sphinxAtStartPar
{\hyperref[\detokenize{pair_reaxff::doc}]{\sphinxcrossref{\DUrole{doc}{pair\_style reaxff}}}}


\subsection{Default}
\label{\detokenize{compute_reaxff_atom:default}}
\sphinxAtStartPar
The option defaults are \sphinxstyleemphasis{bonds} = \sphinxstyleemphasis{no}.

\sphinxstepscope

\index{compute reduce@\spxentry{compute reduce}}\ignorespaces 
\index{compute reduce/region@\spxentry{compute reduce/region}}\ignorespaces 

\section{compute reduce command}
\label{\detokenize{compute_reduce:compute-reduce-command}}\label{\detokenize{compute_reduce:index-1}}\label{\detokenize{compute_reduce:index-0}}\label{\detokenize{compute_reduce::doc}}

\section{compute reduce/region command}
\label{\detokenize{compute_reduce:compute-reduce-region-command}}

\subsection{Syntax}
\label{\detokenize{compute_reduce:syntax}}
\begin{sphinxVerbatim}[commandchars=\\\{\}]
\PYG{k}{compute }\PYG{n+nv+nvIdentifier}{ID}\PYG{+w}{ }\PYG{n+nv+nvIdentifier}{group\PYGZhy{}ID}\PYG{+w}{ }\PYG{n}{style}\PYG{+w}{ }\PYG{n}{arg}\PYG{+w}{ }\PYG{n}{mode}\PYG{+w}{ }\PYG{n}{input1}\PYG{+w}{ }\PYG{n}{input2}\PYG{+w}{ }\PYG{n}{...}\PYG{+w}{ }\PYG{n}{keyword}\PYG{+w}{ }\PYG{n}{args}\PYG{+w}{ }\PYG{n}{...}
\end{sphinxVerbatim}
\begin{itemize}
\item {} 
\sphinxAtStartPar
ID, group\sphinxhyphen{}ID are documented in {\hyperref[\detokenize{compute::doc}]{\sphinxcrossref{\DUrole{doc}{compute}}}} command

\item {} 
\sphinxAtStartPar
style = \sphinxstyleemphasis{reduce} or \sphinxstyleemphasis{reduce/region}
\begin{sphinxalltt}
\sphinxstyleemphasis{reduce} arg = none
\sphinxstyleemphasis{reduce/region} arg = region\sphinxhyphen{}ID
  region\sphinxhyphen{}ID = ID of region to use for choosing atoms
\end{sphinxalltt}

\item {} 
\sphinxAtStartPar
mode = \sphinxstyleemphasis{sum} or \sphinxstyleemphasis{min} or \sphinxstyleemphasis{minabs} or \sphinxstyleemphasis{max} or \sphinxstyleemphasis{maxabs} or \sphinxstyleemphasis{ave} or \sphinxstyleemphasis{sumsq} or \sphinxstyleemphasis{avesq} or \sphinxstyleemphasis{sumabs} or \sphinxstyleemphasis{aveabs}

\item {} 
\sphinxAtStartPar
one or more inputs can be listed

\item {} 
\sphinxAtStartPar
input = \sphinxstyleemphasis{x} or \sphinxstyleemphasis{y} or \sphinxstyleemphasis{z} or \sphinxstyleemphasis{vx} or \sphinxstyleemphasis{vy} or \sphinxstyleemphasis{vz} or \sphinxstyleemphasis{fx} or \sphinxstyleemphasis{fy} or \sphinxstyleemphasis{fz} or c\_ID or c\_ID{[}N{]} or f\_ID or f\_ID{[}N{]} or v\_name
\begin{sphinxalltt}
\sphinxstyleemphasis{x},\sphinxstyleemphasis{y},\sphinxstyleemphasis{z},\sphinxstyleemphasis{vx},\sphinxstyleemphasis{vy},\sphinxstyleemphasis{vz},\sphinxstyleemphasis{fx},\sphinxstyleemphasis{fy},\sphinxstyleemphasis{fz} = atom attribute (position, velocity, force component)
c\_ID = per\sphinxhyphen{}atom or local vector calculated by a compute with ID
c\_ID{[}I{]} = Ith column of per\sphinxhyphen{}atom or local array calculated by a compute with ID, I can include wildcard (see below)
f\_ID = per\sphinxhyphen{}atom or local vector calculated by a fix with ID
f\_ID{[}I{]} = Ith column of per\sphinxhyphen{}atom or local array calculated by a fix with ID, I can include wildcard (see below)
v\_name = per\sphinxhyphen{}atom vector calculated by an atom\sphinxhyphen{}style variable with name
\end{sphinxalltt}

\item {} 
\sphinxAtStartPar
zero or more keyword/args pairs may be appended

\item {} 
\sphinxAtStartPar
keyword = \sphinxstyleemphasis{replace} or \sphinxstyleemphasis{inputs}
\begin{sphinxalltt}
\sphinxstyleemphasis{replace} args = vec1 vec2
  vec1 = reduced value from this input vector will be replaced
  vec2 = replace it with vec1{[}N{]} where N is index of max/min value from vec2
\sphinxstyleemphasis{inputs} arg = peratom or local
  peratom = all inputs are per\sphinxhyphen{}atom quantities (default)
  local = all input are local quantities
\end{sphinxalltt}

\end{itemize}


\subsection{Examples}
\label{\detokenize{compute_reduce:examples}}
\begin{sphinxVerbatim}[commandchars=\\\{\}]
\PYG{k}{compute }\PYG{n+nv+nvIdentifier}{1}\PYG{+w}{ }\PYG{n+nv+nvIdentifier}{all}\PYG{+w}{ }\PYG{n}{reduce}\PYG{+w}{ }\PYG{n}{sum}\PYG{+w}{ }\PYG{n}{c\PYGZus{}force}
\PYG{k}{compute }\PYG{n+nv+nvIdentifier}{1}\PYG{+w}{ }\PYG{n+nv+nvIdentifier}{all}\PYG{+w}{ }\PYG{n}{reduce}\PYG{o}{/}\PYG{k}{region }\PYG{n+nv+nvIdentifier}{subbox}\PYG{+w}{ }\PYG{n}{sum}\PYG{+w}{ }\PYG{n}{c\PYGZus{}force}
\PYG{k}{compute }\PYG{n+nv+nvIdentifier}{2}\PYG{+w}{ }\PYG{n+nv+nvIdentifier}{all}\PYG{+w}{ }\PYG{n}{reduce}\PYG{+w}{ }\PYG{n}{min}\PYG{+w}{ }\PYG{n}{c\PYGZus{}press[2]}\PYG{+w}{ }\PYG{n}{f\PYGZus{}ave}\PYG{+w}{ }\PYG{n}{v\PYGZus{}myKE}
\PYG{k}{compute }\PYG{n+nv+nvIdentifier}{2}\PYG{+w}{ }\PYG{n+nv+nvIdentifier}{all}\PYG{+w}{ }\PYG{n}{reduce}\PYG{+w}{ }\PYG{n}{min}\PYG{+w}{ }\PYG{n}{c\PYGZus{}press[}\PYG{o}{*}\PYG{n}{]}\PYG{+w}{ }\PYG{n}{f\PYGZus{}ave}\PYG{+w}{ }\PYG{n}{v\PYGZus{}myKE}\PYG{+w}{ }\PYG{n}{inputs}\PYG{+w}{ }\PYG{n}{peratom}
\PYG{k}{compute }\PYG{n+nv+nvIdentifier}{3}\PYG{+w}{ }\PYG{n+nv+nvIdentifier}{fluid}\PYG{+w}{ }\PYG{n}{reduce}\PYG{+w}{ }\PYG{n}{max}\PYG{+w}{ }\PYG{n}{c\PYGZus{}index[1]}\PYG{+w}{ }\PYG{n}{c\PYGZus{}index[2]}\PYG{+w}{ }\PYG{n}{c\PYGZus{}dist}\PYG{+w}{ }\PYG{n}{replace}\PYG{+w}{ }\PYG{l+m}{1}\PYG{+w}{ }\PYG{l+m}{3}\PYG{+w}{ }\PYG{n}{replace}\PYG{+w}{ }\PYG{l+m}{2}\PYG{+w}{ }\PYG{l+m}{3}
\PYG{k}{compute }\PYG{n+nv+nvIdentifier}{4}\PYG{+w}{ }\PYG{n+nv+nvIdentifier}{all}\PYG{+w}{ }\PYG{n}{reduce}\PYG{+w}{ }\PYG{n}{max}\PYG{+w}{ }\PYG{n}{c\PYGZus{}bond}\PYG{+w}{ }\PYG{n}{inputs}\PYG{+w}{ }\PYG{n}{local}
\end{sphinxVerbatim}


\subsection{Description}
\label{\detokenize{compute_reduce:description}}
\sphinxAtStartPar
Define a calculation that “reduces” one or more vector inputs into
scalar values, one per listed input.  For the compute reduce command,
the inputs can be either per\sphinxhyphen{}atom or local quantities and must all be
of the same kind (per\sphinxhyphen{}atom or local); see discussion of the optional
\sphinxstyleemphasis{inputs} keyword below.  The compute reduce/region command can only be
used with per\sphinxhyphen{}atom inputs.

\sphinxAtStartPar
Atom attributes are per\sphinxhyphen{}atom quantities, {\hyperref[\detokenize{compute::doc}]{\sphinxcrossref{\DUrole{doc}{computes}}}} and
{\hyperref[\detokenize{fix::doc}]{\sphinxcrossref{\DUrole{doc}{fixes}}}} can generate either per\sphinxhyphen{}atom or local quantities,
and {\hyperref[\detokenize{variable::doc}]{\sphinxcrossref{\DUrole{doc}{atom\sphinxhyphen{}style variables}}}} generate per\sphinxhyphen{}atom
quantities.  See the {\hyperref[\detokenize{variable::doc}]{\sphinxcrossref{\DUrole{doc}{variable}}}} command and its
special functions which can perform the same reduction operations as
the compute reduce command on global vectors.

\sphinxAtStartPar
The reduction operation is specified by the \sphinxstyleemphasis{mode} setting.  The \sphinxstyleemphasis{sum}
option adds the values in the vector into a global total.  The \sphinxstyleemphasis{min}
or \sphinxstyleemphasis{max} options find the minimum or maximum value across all vector
values.  The \sphinxstyleemphasis{minabs} or \sphinxstyleemphasis{maxabs} options find the minimum or maximum
value across all absolute vector values.  The \sphinxstyleemphasis{ave} setting adds the
vector values into a global total, then divides by the number of
values in the vector.  The \sphinxstyleemphasis{sumsq} option sums the square of the
values in the vector into a global total.  The \sphinxstyleemphasis{avesq} setting does
the same as \sphinxstyleemphasis{sumsq}, then divides the sum of squares by the number of
values.  The last two options can be useful for calculating the
variance of some quantity (e.g., variance = sumsq \(-\) ave\(^2\)).  The \sphinxstyleemphasis{sumabs} option sums the absolute values in the
vector into a global total.  The \sphinxstyleemphasis{aveabs} setting does the same as
\sphinxstyleemphasis{sumabs}, then divides the sum of absolute values by the number of
values.

\sphinxAtStartPar
Each listed input is operated on independently.  For per\sphinxhyphen{}atom inputs,
the group specified with this command means only atoms within the
group contribute to the result.  Likewise for per\sphinxhyphen{}atom inputs, if the
compute reduce/region command is used, the atoms must also currently
be within the region.  Note that an input that produces per\sphinxhyphen{}atom
quantities may define its own group which affects the quantities it
returns.  For example, if a compute is used as an input which
generates a per\sphinxhyphen{}atom vector, it will generate values of 0.0 for atoms
that are not in the group specified for that compute.

\sphinxAtStartPar
Each listed input can be an atom attribute (position, velocity, force
component) or can be the result of a {\hyperref[\detokenize{compute::doc}]{\sphinxcrossref{\DUrole{doc}{compute}}}} or
{\hyperref[\detokenize{fix::doc}]{\sphinxcrossref{\DUrole{doc}{fix}}}} or the evaluation of an atom\sphinxhyphen{}style
{\hyperref[\detokenize{variable::doc}]{\sphinxcrossref{\DUrole{doc}{variable}}}}.

\sphinxAtStartPar
Note that for values from a compute or fix, the bracketed index \(I\) can
be specified using a wildcard asterisk with the index to effectively
specify multiple values.  This takes the form “*” or “*n” or “m*” or
“m*n”.  If \(N\) is the size of the vector (for \sphinxstyleemphasis{mode} = scalar) or the
number of columns in the array (for \sphinxstyleemphasis{mode} = vector), then an asterisk
with no numeric values means all indices from 1 to \(N\).  A leading
asterisk means all indices from 1 to n (inclusive).  A trailing
asterisk means all indices from m to \(N\) (inclusive).  A middle asterisk
means all indices from m to n (inclusive).

\sphinxAtStartPar
Using a wildcard is the same as if the individual columns of the array
had been listed one by one. For example, the following two compute reduce
commands are equivalent, since the
{\hyperref[\detokenize{compute_stress_atom::doc}]{\sphinxcrossref{\DUrole{doc}{compute stress/atom}}}} command creates a per\sphinxhyphen{}atom
array with six columns:

\begin{sphinxVerbatim}[commandchars=\\\{\}]
\PYG{k}{compute }\PYG{n+nv+nvIdentifier}{myPress}\PYG{+w}{ }\PYG{n+nv+nvIdentifier}{all}\PYG{+w}{ }\PYG{n}{stress}\PYG{o}{/}\PYG{n}{atom}\PYG{+w}{ }\PYG{n}{NULL}
\PYG{k}{compute }\PYG{n+nv+nvIdentifier}{2}\PYG{+w}{ }\PYG{n+nv+nvIdentifier}{all}\PYG{+w}{ }\PYG{n}{reduce}\PYG{+w}{ }\PYG{n}{min}\PYG{+w}{ }\PYG{n}{c\PYGZus{}myPress[}\PYG{o}{*}\PYG{n}{]}
\PYG{k}{compute }\PYG{n+nv+nvIdentifier}{2}\PYG{+w}{ }\PYG{n+nv+nvIdentifier}{all}\PYG{+w}{ }\PYG{n}{reduce}\PYG{+w}{ }\PYG{n}{min}\PYG{+w}{ }\PYG{n}{c\PYGZus{}myPress[1]}\PYG{+w}{ }\PYG{n}{c\PYGZus{}myPress[2]}\PYG{+w}{ }\PYG{n}{c\PYGZus{}myPress[3]}\PYG{+w}{ }\PYG{o}{\PYGZam{}}
\PYG{+w}{                         }\PYG{n}{c\PYGZus{}myPress[4]}\PYG{+w}{ }\PYG{n}{c\PYGZus{}myPress[5]}\PYG{+w}{ }\PYG{n}{c\PYGZus{}myPress[6]}
\end{sphinxVerbatim}


\bigskip\hrule\bigskip


\sphinxAtStartPar
The atom attribute values (\sphinxstyleemphasis{x}, \sphinxstyleemphasis{y}, \sphinxstyleemphasis{z}, \sphinxstyleemphasis{vx}, \sphinxstyleemphasis{vy}, \sphinxstyleemphasis{vz}, \sphinxstyleemphasis{fx},
\sphinxstyleemphasis{fy}, and \sphinxstyleemphasis{fz}) are self\sphinxhyphen{}explanatory.  Note that other atom attributes
can be used as inputs to this fix by using the {\hyperref[\detokenize{compute_property_atom::doc}]{\sphinxcrossref{\DUrole{doc}{compute
property/atom}}}} command and then specifying an
input value from that compute.

\sphinxAtStartPar
If a value begins with “c\_”, a compute ID must follow which has been
previously defined in the input script.  Valid computes can generate
per\sphinxhyphen{}atom or local quantities.  See the individual {\hyperref[\detokenize{compute::doc}]{\sphinxcrossref{\DUrole{doc}{compute}}}} page for details.  If no bracketed integer is appended, the
vector calculated by the compute is used.  If a bracketed integer is
appended, the Ith column of the array calculated by the compute is
used.  Users can also write code for their own compute styles and
{\hyperref[\detokenize{Modify::doc}]{\sphinxcrossref{\DUrole{doc}{add them to LAMMPS}}}}.  See the discussion above for how
\(I\) can be specified with a wildcard asterisk to effectively
specify multiple values.

\sphinxAtStartPar
If a value begins with “f\_”, a fix ID must follow which has been
previously defined in the input script.  Valid fixes can generate
per\sphinxhyphen{}atom or local quantities.  See the individual {\hyperref[\detokenize{fix::doc}]{\sphinxcrossref{\DUrole{doc}{fix}}}}
page for details.  Note that some fixes only produce their values on
certain timesteps, which must be compatible with when compute reduce
references the values, else an error results.  If no bracketed integer
is appended, the vector calculated by the fix is used.  If a bracketed
integer is appended, the Ith column of the array calculated by the fix
is used.  Users can also write code for their own fix style and
{\hyperref[\detokenize{Modify::doc}]{\sphinxcrossref{\DUrole{doc}{add them to LAMMPS}}}}.  See the discussion above for how
\(I\) can be specified with a wildcard asterisk to effectively
specify multiple values.

\sphinxAtStartPar
If a value begins with “v\_”, a variable name must follow which has
been previously defined in the input script.  It must be an
{\hyperref[\detokenize{variable::doc}]{\sphinxcrossref{\DUrole{doc}{atom\sphinxhyphen{}style variable}}}}.  Atom\sphinxhyphen{}style variables can
reference thermodynamic keywords and various per\sphinxhyphen{}atom attributes, or
invoke other computes, fixes, or variables when they are evaluated, so
this is a very general means of generating per\sphinxhyphen{}atom quantities to
reduce.


\bigskip\hrule\bigskip


\sphinxAtStartPar
If the \sphinxstyleemphasis{replace} keyword is used, two indices \sphinxstyleemphasis{vec1} and \sphinxstyleemphasis{vec2} are
specified, where each index ranges from 1 to the number of input
values.  The replace keyword can only be used if the \sphinxstyleemphasis{mode} is \sphinxstyleemphasis{min}
or \sphinxstyleemphasis{max}.  It works as follows.  A min/max is computed as usual on
the \sphinxstyleemphasis{vec2} input vector.  The index \(N\) of that value within
\sphinxstyleemphasis{vec2} is also stored.  Then, instead of performing a min/max on the
\sphinxstyleemphasis{vec1} input vector, the stored index is used to select the \(N\)th element of the \sphinxstyleemphasis{vec1} vector.

\sphinxAtStartPar
Thus, for example, if you wish to use this compute to find the bond
with maximum stretch, you can do it as follows:

\begin{sphinxVerbatim}[commandchars=\\\{\}]
\PYG{k}{compute }\PYG{n+nv+nvIdentifier}{1}\PYG{+w}{ }\PYG{n+nv+nvIdentifier}{all}\PYG{+w}{ }\PYG{n}{property}\PYG{o}{/}\PYG{n}{local}\PYG{+w}{ }\PYG{n}{batom1}\PYG{+w}{ }\PYG{n}{batom2}
\PYG{k}{compute }\PYG{n+nv+nvIdentifier}{2}\PYG{+w}{ }\PYG{n+nv+nvIdentifier}{all}\PYG{+w}{ }\PYG{n}{bond}\PYG{o}{/}\PYG{n}{local}\PYG{+w}{ }\PYG{n}{dist}
\PYG{k}{compute }\PYG{n+nv+nvIdentifier}{3}\PYG{+w}{ }\PYG{n+nv+nvIdentifier}{all}\PYG{+w}{ }\PYG{n}{reduce}\PYG{+w}{ }\PYG{n}{max}\PYG{+w}{ }\PYG{n}{c\PYGZus{}1[1]}\PYG{+w}{ }\PYG{n}{c\PYGZus{}1[2]}\PYG{+w}{ }\PYG{n}{c\PYGZus{}2}\PYG{+w}{ }\PYG{n}{replace}\PYG{+w}{ }\PYG{l+m}{1}\PYG{+w}{ }\PYG{l+m}{3}\PYG{+w}{ }\PYG{n}{replace}\PYG{+w}{ }\PYG{l+m}{2}\PYG{+w}{ }\PYG{l+m}{3}
\PYG{k}{thermo\PYGZus{}style}\PYG{+w}{ }\PYG{n}{custom}\PYG{+w}{ }\PYG{n}{step}\PYG{+w}{ }\PYG{n}{temp}\PYG{+w}{ }\PYG{n}{c\PYGZus{}3[1]}\PYG{+w}{ }\PYG{n}{c\PYGZus{}3[2]}\PYG{+w}{ }\PYG{n}{c\PYGZus{}3[3]}
\end{sphinxVerbatim}

\sphinxAtStartPar
The first two input values in the compute reduce command are vectors
with the IDs of the two atoms in each bond, using the
{\hyperref[\detokenize{compute_property_local::doc}]{\sphinxcrossref{\DUrole{doc}{compute property/local}}}} command.  The last input
value is bond distance, using the
{\hyperref[\detokenize{compute_bond_local::doc}]{\sphinxcrossref{\DUrole{doc}{compute bond/local}}}} command.  Instead of taking the
max of the two atom ID vectors, which does not yield useful
information in this context, the \sphinxstyleemphasis{replace} keywords will extract the
atom IDs for the two atoms in the bond of maximum stretch.  These atom
IDs and the bond stretch will be printed with thermodynamic output.

\sphinxAtStartPar
\DUrole{versionmodified,added}{Added in version 21Nov2023.}

\sphinxAtStartPar
The \sphinxstyleemphasis{inputs} keyword allows selection of whether all the inputs are
per\sphinxhyphen{}atom or local quantities.  As noted above, all the inputs must be
the same kind (per\sphinxhyphen{}atom or local).  Per\sphinxhyphen{}atom is the default setting.
If a compute or fix is specified as an input, it must produce per\sphinxhyphen{}atom
or local data to match this setting.  If it produces both, e.g. for
the {\hyperref[\detokenize{compute_voronoi_atom::doc}]{\sphinxcrossref{\DUrole{doc}{compute voronoi/atom}}}} command, then
this keyword selects between them.


\bigskip\hrule\bigskip


\sphinxAtStartPar
If a single input is specified this compute produces a global scalar
value.  If multiple inputs are specified, this compute produces a
global vector of values, the length of which is equal to the number of
inputs specified.

\sphinxAtStartPar
As discussed below, for the \sphinxstyleemphasis{sum}, \sphinxstyleemphasis{sumabs}, and \sphinxstyleemphasis{sumsq} modes, the
value(s) produced by this compute are all “extensive”, meaning their
value scales linearly with the number of atoms involved.  If
normalized values are desired, this compute can be accessed by the
{\hyperref[\detokenize{thermo_style::doc}]{\sphinxcrossref{\DUrole{doc}{thermo\_style custom}}}} command with
{\hyperref[\detokenize{thermo_modify::doc}]{\sphinxcrossref{\DUrole{doc}{thermo\_modify norm yes}}}} set as an option.  Or it
can be accessed by a {\hyperref[\detokenize{variable::doc}]{\sphinxcrossref{\DUrole{doc}{variable}}}} that divides by the
appropriate atom count.


\bigskip\hrule\bigskip



\subsection{Output info}
\label{\detokenize{compute_reduce:output-info}}
\sphinxAtStartPar
This compute calculates a global scalar if a single input value is
specified or a global vector of length \(N\), where \(N\) is
the number of inputs, and which can be accessed by indices 1 to
\(N\).  These values can be used by any command that uses global
scalar or vector values from a compute as input.  See the {\hyperref[\detokenize{Howto_output::doc}]{\sphinxcrossref{\DUrole{doc}{Howto
output}}}} doc page for an overview of LAMMPS output
options.

\sphinxAtStartPar
All the scalar or vector values calculated by this compute are
“intensive”, except when the \sphinxstyleemphasis{sum}, \sphinxstyleemphasis{sumabs}, or \sphinxstyleemphasis{sumsq} modes are used on
per\sphinxhyphen{}atom or local vectors, in which case the calculated values are
“extensive”.

\sphinxAtStartPar
The scalar or vector values will be in whatever {\hyperref[\detokenize{units::doc}]{\sphinxcrossref{\DUrole{doc}{units}}}}
the quantities being reduced are in.


\subsection{Restrictions}
\label{\detokenize{compute_reduce:restrictions}}
\sphinxAtStartPar
As noted above, the compute reduce/region command can only be used
with per\sphinxhyphen{}atom inputs.


\subsection{Related commands}
\label{\detokenize{compute_reduce:related-commands}}
\sphinxAtStartPar
{\hyperref[\detokenize{compute::doc}]{\sphinxcrossref{\DUrole{doc}{compute}}}}, {\hyperref[\detokenize{fix::doc}]{\sphinxcrossref{\DUrole{doc}{fix}}}}, {\hyperref[\detokenize{variable::doc}]{\sphinxcrossref{\DUrole{doc}{variable}}}}


\subsection{Default}
\label{\detokenize{compute_reduce:default}}
\sphinxAtStartPar
The default value for the \sphinxstyleemphasis{inputs} keyword is peratom.

\sphinxstepscope

\index{compute reduce/chunk@\spxentry{compute reduce/chunk}}\ignorespaces 

\section{compute reduce/chunk command}
\label{\detokenize{compute_reduce_chunk:compute-reduce-chunk-command}}\label{\detokenize{compute_reduce_chunk:index-0}}\label{\detokenize{compute_reduce_chunk::doc}}

\subsection{Syntax}
\label{\detokenize{compute_reduce_chunk:syntax}}
\begin{sphinxVerbatim}[commandchars=\\\{\}]
\PYG{k}{compute }\PYG{n+nv+nvIdentifier}{ID}\PYG{+w}{ }\PYG{n+nv+nvIdentifier}{group\PYGZhy{}ID}\PYG{+w}{ }\PYG{n}{reduce}\PYG{o}{/}\PYG{n}{chunk}\PYG{+w}{ }\PYG{n}{chunkID}\PYG{+w}{ }\PYG{n}{mode}\PYG{+w}{ }\PYG{n}{input1}\PYG{+w}{ }\PYG{n}{input2}\PYG{+w}{ }\PYG{n}{...}
\end{sphinxVerbatim}
\begin{itemize}
\item {} 
\sphinxAtStartPar
ID, group\sphinxhyphen{}ID are documented in {\hyperref[\detokenize{compute::doc}]{\sphinxcrossref{\DUrole{doc}{compute}}}} command

\item {} 
\sphinxAtStartPar
reduce/chunk = style name of this compute command

\item {} 
\sphinxAtStartPar
chunkID = ID of {\hyperref[\detokenize{compute_chunk_atom::doc}]{\sphinxcrossref{\DUrole{doc}{compute chunk/atom}}}} command

\item {} 
\sphinxAtStartPar
mode = \sphinxstyleemphasis{sum} or \sphinxstyleemphasis{min} or \sphinxstyleemphasis{max}

\item {} 
\sphinxAtStartPar
one or more inputs can be listed

\item {} 
\sphinxAtStartPar
input = c\_ID, c\_ID{[}N{]}, f\_ID, f\_ID{[}N{]}, v\_ID

\begin{sphinxVerbatim}[commandchars=\\\{\}]
c\PYGZus{}ID = per\PYGZhy{}atom vector calculated by a compute with ID
c\PYGZus{}ID[I] = Ith column of per\PYGZhy{}atom array calculated by a compute with ID, I can include wildcard (see below)
f\PYGZus{}ID = per\PYGZhy{}atom vector calculated by a fix with ID
f\PYGZus{}ID[I] = Ith column of per\PYGZhy{}atom array calculated by a fix with ID, I can include wildcard (see below)
v\PYGZus{}name = per\PYGZhy{}atom vector calculated by an atom\PYGZhy{}style variable with name
\end{sphinxVerbatim}

\end{itemize}


\subsection{Examples}
\label{\detokenize{compute_reduce_chunk:examples}}
\begin{sphinxVerbatim}[commandchars=\\\{\}]
\PYG{k}{compute }\PYG{n+nv+nvIdentifier}{1}\PYG{+w}{ }\PYG{n+nv+nvIdentifier}{all}\PYG{+w}{ }\PYG{n}{reduce}\PYG{o}{/}\PYG{n}{chunk}\PYG{+w}{ }\PYG{n}{mychunk}\PYG{+w}{ }\PYG{n}{min}\PYG{+w}{ }\PYG{n}{c\PYGZus{}cluster}
\end{sphinxVerbatim}


\subsection{Description}
\label{\detokenize{compute_reduce_chunk:description}}
\sphinxAtStartPar
Define a calculation that reduces one or more per\sphinxhyphen{}atom vectors into
per\sphinxhyphen{}chunk values.  This can be useful for diagnostic output.  Or when
used in conjunction with the {\hyperref[\detokenize{compute_chunk_spread_atom::doc}]{\sphinxcrossref{\DUrole{doc}{compute chunk/spread/atom}}}} command it can be
used to create per\sphinxhyphen{}atom values that induce a new set of chunks with a
second {\hyperref[\detokenize{compute_chunk_atom::doc}]{\sphinxcrossref{\DUrole{doc}{compute chunk/atom}}}} command.  An
example is given below.

\sphinxAtStartPar
In LAMMPS, chunks are collections of atoms defined by a {\hyperref[\detokenize{compute_chunk_atom::doc}]{\sphinxcrossref{\DUrole{doc}{compute chunk/atom}}}} command, which assigns each atom
to a single chunk (or no chunk).  The ID for this command is specified
as chunkID.  For example, a single chunk could be the atoms in a
molecule or atoms in a spatial bin.  See the {\hyperref[\detokenize{compute_chunk_atom::doc}]{\sphinxcrossref{\DUrole{doc}{compute chunk/atom}}}} and {\hyperref[\detokenize{Howto_chunk::doc}]{\sphinxcrossref{\DUrole{doc}{Howto chunk}}}}
doc pages for details of how chunks can be defined and examples of how
they can be used to measure properties of a system.

\sphinxAtStartPar
For each atom, this compute accesses its chunk ID from the specified
\sphinxstyleemphasis{chunkID} compute.  The per\sphinxhyphen{}atom value from an input contributes
to a per\sphinxhyphen{}chunk value corresponding the the chunk ID.

\sphinxAtStartPar
The reduction operation is specified by the \sphinxstyleemphasis{mode} setting and is
performed over all the per\sphinxhyphen{}atom values from the atoms in each chunk.
The \sphinxstyleemphasis{sum} option adds the pre\sphinxhyphen{}atom values to a per\sphinxhyphen{}chunk total.  The
\sphinxstyleemphasis{min} or \sphinxstyleemphasis{max} options find the minimum or maximum value of the
per\sphinxhyphen{}atom values for each chunk.

\sphinxAtStartPar
Note that only atoms in the specified group contribute to the
reduction operation.  If the \sphinxstyleemphasis{chunkID} compute returns a 0 for the
chunk ID of an atom (i.e., the atom is not in a chunk defined by the
{\hyperref[\detokenize{compute_chunk_atom::doc}]{\sphinxcrossref{\DUrole{doc}{compute chunk/atom}}}} command), that atom will
also not contribute to the reduction operation.  An input that is a
compute or fix may define its own group which affects the quantities
it returns.  For example, a compute with return a zero value for atoms
that are not in the group specified for that compute.

\sphinxAtStartPar
Each listed input is operated on independently.  Each input can be the
result of a {\hyperref[\detokenize{compute::doc}]{\sphinxcrossref{\DUrole{doc}{compute}}}} or {\hyperref[\detokenize{fix::doc}]{\sphinxcrossref{\DUrole{doc}{fix}}}} or the evaluation
of an atom\sphinxhyphen{}style {\hyperref[\detokenize{variable::doc}]{\sphinxcrossref{\DUrole{doc}{variable}}}}.

\sphinxAtStartPar
Note that for values from a compute or fix, the bracketed index I can
be specified using a wildcard asterisk with the index to effectively
specify multiple values.  This takes the form “*” or “*n” or “m*” or
“m*n”.  If \(N\) is the size of the vector (for \sphinxstyleemphasis{mode} = scalar) or the
number of columns in the array (for \sphinxstyleemphasis{mode} = vector), then an asterisk
with no numeric values means all indices from 1 to \(N\).  A leading
asterisk means all indices from 1 to n (inclusive).  A trailing
asterisk means all indices from n to \(N\) (inclusive).  A middle asterisk
means all indices from m to n (inclusive).

\sphinxAtStartPar
Using a wildcard is the same as if the individual columns of the array
had been listed one by one.  For example, the following two compute reduce/chunk
commands are equivalent, since the
{\hyperref[\detokenize{compute_property_chunk::doc}]{\sphinxcrossref{\DUrole{doc}{compute property/chunk}}}} command creates a per\sphinxhyphen{}atom
array with 3 columns:

\begin{sphinxVerbatim}[commandchars=\\\{\}]
\PYG{k}{compute }\PYG{n+nv+nvIdentifier}{prop}\PYG{+w}{ }\PYG{n+nv+nvIdentifier}{all}\PYG{+w}{ }\PYG{n}{property}\PYG{o}{/}\PYG{n}{atom}\PYG{+w}{ }\PYG{n}{vx}\PYG{+w}{ }\PYG{n}{vy}\PYG{+w}{ }\PYG{n}{vz}
\PYG{k}{compute }\PYG{n+nv+nvIdentifier}{10}\PYG{+w}{ }\PYG{n+nv+nvIdentifier}{all}\PYG{+w}{ }\PYG{n}{reduce}\PYG{o}{/}\PYG{n}{chunk}\PYG{+w}{ }\PYG{n}{mychunk}\PYG{+w}{ }\PYG{n}{max}\PYG{+w}{ }\PYG{n}{c\PYGZus{}prop[}\PYG{o}{*}\PYG{n}{]}
\PYG{k}{compute }\PYG{n+nv+nvIdentifier}{10}\PYG{+w}{ }\PYG{n+nv+nvIdentifier}{all}\PYG{+w}{ }\PYG{n}{reduce}\PYG{o}{/}\PYG{n}{chunk}\PYG{+w}{ }\PYG{n}{mychunk}\PYG{+w}{ }\PYG{n}{max}\PYG{+w}{ }\PYG{n}{c\PYGZus{}prop[1]}\PYG{+w}{ }\PYG{n}{c\PYGZus{}prop[2]}\PYG{+w}{ }\PYG{n}{c\PYGZus{}prop[3]}
\end{sphinxVerbatim}


\bigskip\hrule\bigskip


\sphinxAtStartPar
Here is an example of using this compute, in conjunction with the
compute chunk/spread/atom command to identify self\sphinxhyphen{}assembled micelles.
The commands below can be added to the examples/in.micelle script.

\sphinxAtStartPar
Imagine a collection of polymer chains or small molecules with
hydrophobic end groups.  All the hydrophobic (HP) atoms are assigned
to a group called “phobic”.

\sphinxAtStartPar
These commands will assign a unique cluster ID to all HP atoms within
a specified distance of each other.  A cluster will contain all HP
atoms in a single molecule, but also the HP atoms in nearby molecules
(e.g., molecules that have clumped to form a micelle due to the
attraction induced by the hydrophobicity).  The output of the
chunk/reduce command will be a cluster ID per chunk (molecule).
Molecules with the same cluster ID are in the same micelle.

\begin{sphinxVerbatim}[commandchars=\\\{\}]
\PYG{k}{group }\PYG{n+nv+nvIdentifier}{phobic}\PYG{+w}{ }\PYG{n}{type}\PYG{+w}{ }\PYG{l+m}{4}\PYG{+w}{     }\PYG{c}{\PYGZsh{} specific to in.micelle model}
\PYG{k}{compute }\PYG{n+nv+nvIdentifier}{cluster}\PYG{+w}{ }\PYG{n+nv+nvIdentifier}{phobic}\PYG{+w}{ }\PYG{n}{cluster}\PYG{o}{/}\PYG{n}{atom}\PYG{+w}{ }\PYG{l+m}{2.0}
\PYG{k}{compute }\PYG{n+nv+nvIdentifier}{cmol}\PYG{+w}{ }\PYG{n+nv+nvIdentifier}{all}\PYG{+w}{ }\PYG{n}{chunk}\PYG{o}{/}\PYG{n}{atom}\PYG{+w}{ }\PYG{n}{molecule}
\PYG{k}{compute }\PYG{n+nv+nvIdentifier}{reduce}\PYG{+w}{ }\PYG{n+nv+nvIdentifier}{phobic}\PYG{+w}{ }\PYG{n}{reduce}\PYG{o}{/}\PYG{n}{chunk}\PYG{+w}{ }\PYG{n}{cmol}\PYG{+w}{ }\PYG{n}{min}\PYG{+w}{ }\PYG{n}{c\PYGZus{}cluster}
\end{sphinxVerbatim}

\sphinxAtStartPar
This per\sphinxhyphen{}chunk info could be output in at least two ways:

\begin{sphinxVerbatim}[commandchars=\\\{\}]
\PYG{k}{fix }\PYG{n+nv+nvIdentifier}{10}\PYG{+w}{ }\PYG{n+nv+nvIdentifier}{all}\PYG{+w}{ }\PYG{n}{ave}\PYG{o}{/}\PYG{n}{time}\PYG{+w}{ }\PYG{l+m}{1000}\PYG{+w}{ }\PYG{l+m}{1}\PYG{+w}{ }\PYG{l+m}{1000}\PYG{+w}{ }\PYG{n}{c\PYGZus{}reduce}\PYG{+w}{ }\PYG{n}{file}\PYG{+w}{ }\PYG{n}{tmp.phobic}\PYG{+w}{ }\PYG{n}{mode}\PYG{+w}{ }\PYG{n}{vector}

\PYG{k}{compute }\PYG{n+nv+nvIdentifier}{spread}\PYG{+w}{ }\PYG{n+nv+nvIdentifier}{all}\PYG{+w}{ }\PYG{n}{chunk}\PYG{o}{/}\PYG{n}{spread}\PYG{o}{/}\PYG{n}{atom}\PYG{+w}{ }\PYG{n}{cmol}\PYG{+w}{ }\PYG{n}{c\PYGZus{}reduce}
\PYG{k}{dump }\PYG{n+nv+nvIdentifier}{1}\PYG{+w}{ }\PYG{n+nv+nvIdentifier}{all}\PYG{+w}{ }\PYG{n}{custom}\PYG{+w}{ }\PYG{l+m}{1000}\PYG{+w}{ }\PYG{n}{tmp.dump}\PYG{+w}{ }\PYG{n}{id}\PYG{+w}{ }\PYG{n}{type}\PYG{+w}{ }\PYG{n}{mol}\PYG{+w}{ }\PYG{n}{x}\PYG{+w}{ }\PYG{n}{y}\PYG{+w}{ }\PYG{n}{z}\PYG{+w}{ }\PYG{n}{c\PYGZus{}cluster}\PYG{+w}{ }\PYG{n}{c\PYGZus{}spread}
\PYG{k}{dump\PYGZus{}modify }\PYG{n+nv+nvIdentifier}{1}\PYG{+w}{ }\PYG{n}{sort}\PYG{+w}{ }\PYG{n}{id}
\end{sphinxVerbatim}

\sphinxAtStartPar
In the first case, each snapshot in the tmp.phobic file will contain
one line per molecule.  Molecules with the same value are in the same
micelle.  In the second case each dump snapshot contains all atoms,
each with a final field with the cluster ID of the micelle that the HP
atoms of that atom’s molecule belong to.

\sphinxAtStartPar
The result from compute chunk/spread/atom can be used to define a new
set of chunks, where all the atoms in all the molecules in the same
micelle are assigned to the same chunk (i.e., one chunk per micelle).

\begin{sphinxVerbatim}[commandchars=\\\{\}]
\PYG{k}{compute }\PYG{n+nv+nvIdentifier}{micelle}\PYG{+w}{ }\PYG{n+nv+nvIdentifier}{all}\PYG{+w}{ }\PYG{n}{chunk}\PYG{o}{/}\PYG{n}{atom}\PYG{+w}{ }\PYG{n}{c\PYGZus{}spread}\PYG{+w}{ }\PYG{n}{compress}\PYG{+w}{ }\PYG{n}{yes}
\end{sphinxVerbatim}

\sphinxAtStartPar
Further analysis on a per\sphinxhyphen{}micelle basis can now be performed using any
of the per\sphinxhyphen{}chunk computes listed on the {\hyperref[\detokenize{Howto_chunk::doc}]{\sphinxcrossref{\DUrole{doc}{Howto chunk}}}}
doc page (e.g., count the number of atoms in each micelle, calculate
its center or mass, shape/moments of inertia, and radius of gyration).

\begin{sphinxVerbatim}[commandchars=\\\{\}]
\PYG{k}{compute }\PYG{n+nv+nvIdentifier}{prop}\PYG{+w}{ }\PYG{n+nv+nvIdentifier}{all}\PYG{+w}{ }\PYG{n}{property}\PYG{o}{/}\PYG{n}{chunk}\PYG{+w}{ }\PYG{n}{micelle}\PYG{+w}{ }\PYG{n}{count}
\PYG{k}{fix }\PYG{n+nv+nvIdentifier}{20}\PYG{+w}{ }\PYG{n+nv+nvIdentifier}{all}\PYG{+w}{ }\PYG{n}{ave}\PYG{o}{/}\PYG{n}{time}\PYG{+w}{ }\PYG{l+m}{1000}\PYG{+w}{ }\PYG{l+m}{1}\PYG{+w}{ }\PYG{l+m}{1000}\PYG{+w}{ }\PYG{n}{c\PYGZus{}prop}\PYG{+w}{ }\PYG{n}{file}\PYG{+w}{ }\PYG{n}{tmp.micelle}\PYG{+w}{ }\PYG{n}{mode}\PYG{+w}{ }\PYG{n}{vector}
\end{sphinxVerbatim}

\sphinxAtStartPar
Each snapshot in the tmp.micelle file will have one line per micelle
with its count of atoms, plus a first line for a chunk with all the
solvent atoms.  By the time 50000 steps have elapsed, there are a
handful of large micelles.


\bigskip\hrule\bigskip



\subsection{Output info}
\label{\detokenize{compute_reduce_chunk:output-info}}
\sphinxAtStartPar
This compute calculates a global vector if a single input value is
specified, otherwise a global array is output.  The number of columns
in the array is the number of inputs provided.  The length of the
vector or the number of vector elements or array rows = the number of
chunks \sphinxstyleemphasis{Nchunk} as calculated by the specified {\hyperref[\detokenize{compute_chunk_atom::doc}]{\sphinxcrossref{\DUrole{doc}{compute chunk/atom}}}} command.  The vector or array can
be accessed by any command that uses global values from a compute as
input.  See the {\hyperref[\detokenize{Howto_output::doc}]{\sphinxcrossref{\DUrole{doc}{Howto output}}}} page for an
overview of LAMMPS output options.

\sphinxAtStartPar
The per\sphinxhyphen{}atom values for the vector or each column of the array will be
in whatever {\hyperref[\detokenize{units::doc}]{\sphinxcrossref{\DUrole{doc}{units}}}} the corresponding input value is in.
The vector or array values are “intensive”.


\subsection{Restrictions}
\label{\detokenize{compute_reduce_chunk:restrictions}}\begin{quote}

\sphinxAtStartPar
none
\end{quote}


\subsection{Related commands}
\label{\detokenize{compute_reduce_chunk:related-commands}}
\sphinxAtStartPar
{\hyperref[\detokenize{compute_chunk_atom::doc}]{\sphinxcrossref{\DUrole{doc}{compute chunk/atom}}}}, {\hyperref[\detokenize{compute_reduce::doc}]{\sphinxcrossref{\DUrole{doc}{compute reduce}}}}, {\hyperref[\detokenize{compute_chunk_spread_atom::doc}]{\sphinxcrossref{\DUrole{doc}{compute chunk/spread/atom}}}}


\subsection{Default}
\label{\detokenize{compute_reduce_chunk:default}}
\sphinxAtStartPar
none

\sphinxstepscope

\index{compute rheo/property/atom@\spxentry{compute rheo/property/atom}}\ignorespaces 

\section{compute rheo/property/atom command}
\label{\detokenize{compute_rheo_property_atom:compute-rheo-property-atom-command}}\label{\detokenize{compute_rheo_property_atom:index-0}}\label{\detokenize{compute_rheo_property_atom::doc}}

\subsection{Syntax}
\label{\detokenize{compute_rheo_property_atom:syntax}}
\begin{sphinxVerbatim}[commandchars=\\\{\}]
\PYG{k}{compute }\PYG{n+nv+nvIdentifier}{ID}\PYG{+w}{ }\PYG{n+nv+nvIdentifier}{group\PYGZhy{}ID}\PYG{+w}{ }\PYG{n}{rheo}\PYG{o}{/}\PYG{n}{property}\PYG{o}{/}\PYG{n}{atom}\PYG{+w}{ }\PYG{n}{input1}\PYG{+w}{ }\PYG{n}{input2}\PYG{+w}{ }\PYG{n}{...}
\end{sphinxVerbatim}
\begin{itemize}
\item {} 
\sphinxAtStartPar
ID, group\sphinxhyphen{}ID are documented in {\hyperref[\detokenize{compute::doc}]{\sphinxcrossref{\DUrole{doc}{compute}}}} command

\item {} 
\sphinxAtStartPar
rheo/property/atom = style name of this compute command

\item {} 
\sphinxAtStartPar
input = one or more atom attributes

\begin{sphinxVerbatim}[commandchars=\\\{\}]
possible attributes = phase, surface, surface/r,
                      surface/divr, surface/n/a, coordination,
                      shift/v/a, energy, temperature, heatflow,
                      conductivity, cv, viscosity, pressure, rho,
                      grad/v/ab, stress/v/ab, stress/t/ab, nbond/shell
\end{sphinxVerbatim}
\begin{sphinxalltt}
\sphinxstyleemphasis{phase} = atom phase state
\sphinxstyleemphasis{surface} = atom surface status
\sphinxstyleemphasis{surface/r} = atom distance from the surface
\sphinxstyleemphasis{surface/divr} = divergence of position at atom position
\sphinxstyleemphasis{surface/n/a} = a\sphinxhyphen{}component of surface normal vector
\sphinxstyleemphasis{coordination} = coordination number
\sphinxstyleemphasis{shift/v/a} = a\sphinxhyphen{}component of atom shifting velocity
\sphinxstyleemphasis{energy} = atom energy
\sphinxstyleemphasis{temperature} = atom temperature
\sphinxstyleemphasis{heatflow} = atom heat flow
\sphinxstyleemphasis{conductivity} = atom conductivity
\sphinxstyleemphasis{cv} = atom specific heat
\sphinxstyleemphasis{viscosity} = atom viscosity
\sphinxstyleemphasis{pressure} = atom pressure
\sphinxstyleemphasis{rho} = atom density
\sphinxstyleemphasis{grad/v/ab} = ab\sphinxhyphen{}component of atom velocity gradient tensor
\sphinxstyleemphasis{stress/v/ab} = ab\sphinxhyphen{}component of atom viscous stress tensor
\sphinxstyleemphasis{stress/t/ab} = ab\sphinxhyphen{}component of atom total stress tensor (pressure and viscous)
\sphinxstyleemphasis{nbond/shell} = number of oxide bonds
\end{sphinxalltt}

\end{itemize}


\subsection{Examples}
\label{\detokenize{compute_rheo_property_atom:examples}}
\begin{sphinxVerbatim}[commandchars=\\\{\}]
\PYG{k}{compute }\PYG{n+nv+nvIdentifier}{1}\PYG{+w}{ }\PYG{n+nv+nvIdentifier}{all}\PYG{+w}{ }\PYG{n}{rheo}\PYG{o}{/}\PYG{n}{property}\PYG{o}{/}\PYG{n}{atom}\PYG{+w}{ }\PYG{n}{phase}\PYG{+w}{ }\PYG{n}{surface}\PYG{o}{/}\PYG{n}{r}\PYG{+w}{ }\PYG{n}{surface}\PYG{o}{/}\PYG{n}{n}\PYG{o}{/}\PYG{o}{*}\PYG{+w}{ }\PYG{n}{pressure}
\PYG{k}{compute }\PYG{n+nv+nvIdentifier}{2}\PYG{+w}{ }\PYG{n+nv+nvIdentifier}{all}\PYG{+w}{ }\PYG{n}{rheo}\PYG{o}{/}\PYG{n}{property}\PYG{o}{/}\PYG{n}{atom}\PYG{+w}{ }\PYG{n}{shift}\PYG{o}{/}\PYG{n}{v}\PYG{o}{/}\PYG{n}{x}\PYG{+w}{ }\PYG{n}{grad}\PYG{o}{/}\PYG{n}{v}\PYG{o}{/}\PYG{n}{xx}\PYG{+w}{ }\PYG{n}{stress}\PYG{o}{/}\PYG{n}{v}\PYG{o}{/}\PYG{o}{*}
\end{sphinxVerbatim}


\subsection{Description}
\label{\detokenize{compute_rheo_property_atom:description}}
\sphinxAtStartPar
\DUrole{versionmodified,added}{Added in version 29Aug2024.}

\sphinxAtStartPar
Define a computation that stores atom attributes specific to the RHEO
package for each atom in the group.  This is useful so that the values
can be used by other {\hyperref[\detokenize{Howto_output::doc}]{\sphinxcrossref{\DUrole{doc}{output commands}}}} that take
computes as inputs. See for example, the
{\hyperref[\detokenize{compute_reduce::doc}]{\sphinxcrossref{\DUrole{doc}{compute reduce}}}},
{\hyperref[\detokenize{fix_ave_atom::doc}]{\sphinxcrossref{\DUrole{doc}{fix ave/atom}}}},
{\hyperref[\detokenize{fix_ave_histo::doc}]{\sphinxcrossref{\DUrole{doc}{fix ave/histo}}}},
{\hyperref[\detokenize{fix_ave_chunk::doc}]{\sphinxcrossref{\DUrole{doc}{fix ave/chunk}}}}, and
{\hyperref[\detokenize{variable::doc}]{\sphinxcrossref{\DUrole{doc}{atom\sphinxhyphen{}style variable}}}} commands.

\sphinxAtStartPar
For vector attributes, e.g. \sphinxstyleemphasis{shift/v/}\(\alpha\), one must specify
\(\alpha\) as the \sphinxstyleemphasis{x}, \sphinxstyleemphasis{y}, or \sphinxstyleemphasis{z} component, e.g. \sphinxstyleemphasis{shift/v/x}.
Alternatively, a wild card * will include all components, \sphinxstyleemphasis{x} and \sphinxstyleemphasis{y} in
2D or \sphinxstyleemphasis{x}, \sphinxstyleemphasis{y}, and \sphinxstyleemphasis{z} in 3D.

\sphinxAtStartPar
For tensor attributes, e.g. \sphinxstyleemphasis{grad/v/}\(\alpha \beta\), one must specify
both \(\alpha\) and \(\beta\) as  \sphinxstyleemphasis{x}, \sphinxstyleemphasis{y}, or \sphinxstyleemphasis{z}, e.g. \sphinxstyleemphasis{grad/v/xy}.
Alternatively, a wild card * will include all components. In 2D, this
includes \sphinxstyleemphasis{xx}, \sphinxstyleemphasis{xy}, \sphinxstyleemphasis{yx}, and \sphinxstyleemphasis{yy}. In 3D, this includes \sphinxstyleemphasis{xx}, \sphinxstyleemphasis{xy}, \sphinxstyleemphasis{xz},
\sphinxstyleemphasis{yx}, \sphinxstyleemphasis{yy}, \sphinxstyleemphasis{yz}, \sphinxstyleemphasis{zx}, \sphinxstyleemphasis{zy}, and \sphinxstyleemphasis{zz}.

\sphinxAtStartPar
Many properties require their respective fixes, listed below in related
commands, be defined. For instance, the \sphinxstyleemphasis{viscosity} attribute is the
viscosity of a particle calculated by
{\hyperref[\detokenize{fix_rheo_viscosity::doc}]{\sphinxcrossref{\DUrole{doc}{fix rheo/viscous}}}}. The meaning of less obvious
properties is described below.

\sphinxAtStartPar
The \sphinxstyleemphasis{phase} property indicates whether the particle is in a fluid state,
a value of 0, or a solid state, a value of 1.

\sphinxAtStartPar
The \sphinxstyleemphasis{surface} property indicates the surface designation produced by
the \sphinxstyleemphasis{interface/reconstruct} option of {\hyperref[\detokenize{fix_rheo::doc}]{\sphinxcrossref{\DUrole{doc}{fix rheo}}}}. Bulk
particles have a value of 0, surface particles have a value of 1, and
splash particles have a value of 2. The \sphinxstyleemphasis{surface/r} property is the
distance from the surface, up to the kernel cutoff length. Surface particles
have a value of 0. The \sphinxstyleemphasis{surface/n/}\(\alpha\) properties are the
components of the surface normal vector.

\sphinxAtStartPar
The \sphinxstyleemphasis{shift/v/}\(\alpha\) properties are the components of the shifting
velocity produced by the \sphinxstyleemphasis{shift} option of {\hyperref[\detokenize{fix_rheo::doc}]{\sphinxcrossref{\DUrole{doc}{fix rheo}}}}.

\sphinxAtStartPar
The \sphinxstyleemphasis{nbond/shell} property is the number of shell bonds that have been
activated from {\hyperref[\detokenize{bond_rheo_shell::doc}]{\sphinxcrossref{\DUrole{doc}{bond style rheo/shell}}}}.

\sphinxAtStartPar
The values are stored in a per\sphinxhyphen{}atom vector or array as discussed
below.  Zeroes are stored for atoms not in the specified group or for
quantities that are not defined for a particular particle in the group


\subsection{Output info}
\label{\detokenize{compute_rheo_property_atom:output-info}}
\sphinxAtStartPar
This compute calculates a per\sphinxhyphen{}atom vector or per\sphinxhyphen{}atom array depending
on the number of input values.  Generally, if a single input is specified,
a per\sphinxhyphen{}atom vector is produced.  If two or more inputs are specified, a
per\sphinxhyphen{}atom array is produced where the number of columns = the number of
inputs. However, if a wild card * is used for a vector or tensor, then
the number of inputs is considered to be incremented by the dimension or
the dimension squared, respectively. The vector or array can be accessed
by any command that uses per\sphinxhyphen{}atom values from a compute as input.  See the
{\hyperref[\detokenize{Howto_output::doc}]{\sphinxcrossref{\DUrole{doc}{Howto output}}}} page for an overview of LAMMPS output
options.

\sphinxAtStartPar
The vector or array values will be in whatever {\hyperref[\detokenize{units::doc}]{\sphinxcrossref{\DUrole{doc}{units}}}} the
corresponding attribute is in (e.g., density units for \sphinxstyleemphasis{rho}).


\subsection{Restrictions}
\label{\detokenize{compute_rheo_property_atom:restrictions}}\begin{quote}

\sphinxAtStartPar
none
\end{quote}


\subsection{Related commands}
\label{\detokenize{compute_rheo_property_atom:related-commands}}
\sphinxAtStartPar
{\hyperref[\detokenize{dump::doc}]{\sphinxcrossref{\DUrole{doc}{dump custom}}}}, {\hyperref[\detokenize{compute_reduce::doc}]{\sphinxcrossref{\DUrole{doc}{compute reduce}}}},
{\hyperref[\detokenize{fix_ave_atom::doc}]{\sphinxcrossref{\DUrole{doc}{fix ave/atom}}}}, {\hyperref[\detokenize{fix_ave_chunk::doc}]{\sphinxcrossref{\DUrole{doc}{fix ave/chunk}}}},
{\hyperref[\detokenize{fix_rheo_viscosity::doc}]{\sphinxcrossref{\DUrole{doc}{fix rheo/viscosity}}}},
{\hyperref[\detokenize{fix_rheo_pressure::doc}]{\sphinxcrossref{\DUrole{doc}{fix rheo/pressure}}}},
{\hyperref[\detokenize{fix_rheo_thermal::doc}]{\sphinxcrossref{\DUrole{doc}{fix rheo/thermal}}}},
{\hyperref[\detokenize{fix_rheo_oxidation::doc}]{\sphinxcrossref{\DUrole{doc}{fix rheo/oxdiation}}}},
{\hyperref[\detokenize{fix_rheo::doc}]{\sphinxcrossref{\DUrole{doc}{fix rheo}}}}


\subsection{Default}
\label{\detokenize{compute_rheo_property_atom:default}}
\sphinxAtStartPar
none

\sphinxstepscope

\index{compute rigid/local@\spxentry{compute rigid/local}}\ignorespaces 

\section{compute rigid/local command}
\label{\detokenize{compute_rigid_local:compute-rigid-local-command}}\label{\detokenize{compute_rigid_local:index-0}}\label{\detokenize{compute_rigid_local::doc}}

\subsection{Syntax}
\label{\detokenize{compute_rigid_local:syntax}}
\begin{sphinxVerbatim}[commandchars=\\\{\}]
\PYG{k}{compute }\PYG{n+nv+nvIdentifier}{ID}\PYG{+w}{ }\PYG{n+nv+nvIdentifier}{group\PYGZhy{}ID}\PYG{+w}{ }\PYG{n}{rigid}\PYG{o}{/}\PYG{n}{local}\PYG{+w}{ }\PYG{n}{rigidID}\PYG{+w}{ }\PYG{n}{input1}\PYG{+w}{ }\PYG{n}{input2}\PYG{+w}{ }\PYG{n}{...}
\end{sphinxVerbatim}
\begin{itemize}
\item {} 
\sphinxAtStartPar
ID, group\sphinxhyphen{}ID are documented in {\hyperref[\detokenize{compute::doc}]{\sphinxcrossref{\DUrole{doc}{compute}}}} command

\item {} 
\sphinxAtStartPar
rigid/local = style name of this compute command

\item {} 
\sphinxAtStartPar
rigidID = ID of fix rigid/small command or one of its variants

\item {} 
\sphinxAtStartPar
input = one or more rigid body attributes

\begin{sphinxVerbatim}[commandchars=\\\{\}]
possible attributes = id, mol, mass,
                      x, y, z, xu, yu, zu, ix, iy, iz
                      vx, vy, vz, fx, fy, fz,
                      omegax, omegay, omegaz,
                      angmomx, angmomy, angmomz,
                      quatw, quati, quatj, quatk,
                      tqx, tqy, tqz,
                      inertiax, inertiay, inertiaz
\end{sphinxVerbatim}
\begin{sphinxalltt}
id = atom ID of atom within body which owns body properties
mol = molecule ID used to define body in {\hyperref[\detokenize{fix_rigid::doc}]{\sphinxcrossref{\DUrole{doc}{fix rigid/small}}}} command
mass = total mass of body
x,y,z = center of mass coords of body
xu,yu,zu = unwrapped center of mass coords of body
ix,iy,iz = box image that the center of mass is in
vx,vy,vz = center of mass velocities
fx,fy,fz = force of center of mass
omegax,omegay,omegaz = angular velocity of body
angmomx,angmomy,angmomz = angular momentum of body
quatw,quati,quatj,quatk = quaternion components for body
tqx,tqy,tqz = torque on body
inertiax,inertiay,inertiaz = diagonalized moments of inertia of body
\end{sphinxalltt}

\end{itemize}


\subsection{Examples}
\label{\detokenize{compute_rigid_local:examples}}
\begin{sphinxVerbatim}[commandchars=\\\{\}]
\PYG{k}{compute }\PYG{n+nv+nvIdentifier}{1}\PYG{+w}{ }\PYG{n+nv+nvIdentifier}{all}\PYG{+w}{ }\PYG{n}{rigid}\PYG{o}{/}\PYG{n}{local}\PYG{+w}{ }\PYG{n}{myRigid}\PYG{+w}{ }\PYG{n}{mol}\PYG{+w}{ }\PYG{n}{x}\PYG{+w}{ }\PYG{n}{y}\PYG{+w}{ }\PYG{n}{z}
\end{sphinxVerbatim}


\subsection{Description}
\label{\detokenize{compute_rigid_local:description}}
\sphinxAtStartPar
Define a computation that simply stores rigid body attributes for
rigid bodies defined by the {\hyperref[\detokenize{fix_rigid::doc}]{\sphinxcrossref{\DUrole{doc}{fix rigid/small}}}} command
or one of its NVE, NVT, NPT, NPH variants.  The data is stored as
local data so it can be accessed by other {\hyperref[\detokenize{Howto_output::doc}]{\sphinxcrossref{\DUrole{doc}{output commands}}}} that process local data, such as the
{\hyperref[\detokenize{compute_reduce::doc}]{\sphinxcrossref{\DUrole{doc}{compute reduce}}}} or {\hyperref[\detokenize{dump::doc}]{\sphinxcrossref{\DUrole{doc}{dump local}}}}
commands.

\sphinxAtStartPar
Note that this command only works with the {\hyperref[\detokenize{fix_rigid::doc}]{\sphinxcrossref{\DUrole{doc}{fix rigid/small}}}} command or its variants, not the fix rigid
command and its variants.  The ID of the {\hyperref[\detokenize{fix_rigid::doc}]{\sphinxcrossref{\DUrole{doc}{fix rigid/small}}}} command used to define rigid bodies must
be specified as \sphinxstyleemphasis{rigidID}.  The {\hyperref[\detokenize{fix_rigid::doc}]{\sphinxcrossref{\DUrole{doc}{fix rigid}}}} command is
typically used to define a handful of (potentially very large) rigid
bodies.  It outputs similar per\sphinxhyphen{}body information as this command
directly from the fix as global data; see the {\hyperref[\detokenize{fix_rigid::doc}]{\sphinxcrossref{\DUrole{doc}{fix rigid}}}} page for details

\sphinxAtStartPar
The local data stored by this command is generated by looping over all
the atoms owned on a processor.  If the atom is not in the specified
\sphinxstyleemphasis{group\sphinxhyphen{}ID} or is not part of a rigid body it is skipped.  If it is not
the atom within a body that is assigned to store the body information
it is skipped (only one atom per body is so assigned).  If it is the
assigned atom, then the info for that body is output.  This means that
information for \(N\) bodies is generated.  \(N\) may be less than the
number of bodies defined by the fix rigid command, if the atoms in some bodies
are not in the \sphinxstyleemphasis{group\sphinxhyphen{}ID}.

\begin{sphinxadmonition}{note}{Note:}
\sphinxAtStartPar
Which atom in a body owns the body info is determined internal
to LAMMPS; it’s the one nearest the geometric center of the body.
Typically you should avoid this complication, by defining the group
associated with this fix to include/exclude entire bodies.
\end{sphinxadmonition}

\sphinxAtStartPar
Note that as atoms and bodies migrate from processor to processor,
there will be no consistent ordering of the entries within the local
vector or array from one timestep to the next.

\sphinxAtStartPar
Here is an example of how to use this compute to dump rigid body info
to a file:

\begin{sphinxVerbatim}[commandchars=\\\{\}]
\PYG{k}{compute }\PYG{n+nv+nvIdentifier}{1}\PYG{+w}{ }\PYG{n+nv+nvIdentifier}{all}\PYG{+w}{ }\PYG{n}{rigid}\PYG{o}{/}\PYG{n}{local}\PYG{+w}{ }\PYG{n}{myRigid}\PYG{+w}{ }\PYG{n}{mol}\PYG{+w}{ }\PYG{n}{x}\PYG{+w}{ }\PYG{n}{y}\PYG{+w}{ }\PYG{n}{z}\PYG{+w}{ }\PYG{n}{fx}\PYG{+w}{ }\PYG{n}{fy}\PYG{+w}{ }\PYG{n}{fz}
\PYG{k}{dump }\PYG{n+nv+nvIdentifier}{1}\PYG{+w}{ }\PYG{n+nv+nvIdentifier}{all}\PYG{+w}{ }\PYG{n}{local}\PYG{+w}{ }\PYG{l+m}{1000}\PYG{+w}{ }\PYG{n}{tmp.dump}\PYG{+w}{ }\PYG{n}{index}\PYG{+w}{ }\PYG{n}{c\PYGZus{}1[1]}\PYG{+w}{ }\PYG{n}{c\PYGZus{}1[2]}\PYG{+w}{ }\PYG{n}{c\PYGZus{}1[3]}\PYG{+w}{ }\PYG{n}{c\PYGZus{}1[4]}\PYG{+w}{ }\PYG{n}{c\PYGZus{}1[5]}\PYG{+w}{ }\PYG{n}{c\PYGZus{}1[6]}\PYG{+w}{ }\PYG{n}{c\PYGZus{}1[7]}
\end{sphinxVerbatim}


\bigskip\hrule\bigskip


\sphinxAtStartPar
This section explains the rigid body attributes that can be specified.

\sphinxAtStartPar
The \sphinxstyleemphasis{id} attribute is the atom\sphinxhyphen{}ID of the atom which owns the rigid body, which is
assigned by the {\hyperref[\detokenize{fix_rigid::doc}]{\sphinxcrossref{\DUrole{doc}{fix rigid/small}}}} command.

\sphinxAtStartPar
The \sphinxstyleemphasis{mol} attribute is the molecule ID of the rigid body.  It should
be the molecule ID which all of the atoms in the body belong to, since
that is how the {\hyperref[\detokenize{fix_rigid::doc}]{\sphinxcrossref{\DUrole{doc}{fix rigid/small}}}} command defines its
rigid bodies.

\sphinxAtStartPar
The \sphinxstyleemphasis{mass} attribute is the total mass of the rigid body.

\sphinxAtStartPar
There are two options for outputting the coordinates of the center of
mass (COM) of the body.  The \sphinxstyleemphasis{x}, \sphinxstyleemphasis{y}, \sphinxstyleemphasis{z} attributes write the COM
“unscaled”, in the appropriate distance {\hyperref[\detokenize{units::doc}]{\sphinxcrossref{\DUrole{doc}{units}}}}
(\(\AA\), \(\sigma\), etc).  Use \sphinxstyleemphasis{xu}, \sphinxstyleemphasis{yu}, \sphinxstyleemphasis{zu} if you want the COM “unwrapped” by
the image flags for each body.  Unwrapped means that if the body
COM has passed through a periodic boundary one or more times, the value
is generated what the COM coordinate would be if it had not been
wrapped back into the periodic box.

\sphinxAtStartPar
The image flags for the body can be generated directly using the \sphinxstyleemphasis{ix},
\sphinxstyleemphasis{iy}, \sphinxstyleemphasis{iz} attributes.  For periodic dimensions, they specify which
image of the simulation box the COM is considered to be in.  An image
of 0 means it is inside the box as defined.  A value of 2 means add 2
box lengths to get the true value.  A value of \(-1\) means subtract 1 box
length to get the true value.  LAMMPS updates these flags as the rigid
body COMs cross periodic boundaries during the simulation.

\sphinxAtStartPar
The \sphinxstyleemphasis{vx}, \sphinxstyleemphasis{vy}, \sphinxstyleemphasis{vz}, \sphinxstyleemphasis{fx}, \sphinxstyleemphasis{fy}, \sphinxstyleemphasis{fz} attributes are components of
the COM velocity and force on the COM of the body.

\sphinxAtStartPar
The \sphinxstyleemphasis{omegax}, \sphinxstyleemphasis{omegay}, and \sphinxstyleemphasis{omegaz} attributes are the angular
velocity components of the body in the system frame around its COM.

\sphinxAtStartPar
The \sphinxstyleemphasis{angmomx}, \sphinxstyleemphasis{angmomy}, and \sphinxstyleemphasis{angmomz} attributes are the angular
momentum components of the body in the system frame around its COM.

\sphinxAtStartPar
The \sphinxstyleemphasis{quatw}, \sphinxstyleemphasis{quati}, \sphinxstyleemphasis{quatj}, and \sphinxstyleemphasis{quatk} attributes are the
components of the 4\sphinxhyphen{}vector quaternion representing the orientation of
the rigid body.  See the {\hyperref[\detokenize{set::doc}]{\sphinxcrossref{\DUrole{doc}{set}}}} command for an explanation of
the quaternion vector.

\sphinxAtStartPar
The \sphinxstyleemphasis{tqx}, \sphinxstyleemphasis{tqy}, \sphinxstyleemphasis{tqz} attributes are components of the torque acting
on the body around its COM.

\sphinxAtStartPar
The \sphinxstyleemphasis{inertiax}, \sphinxstyleemphasis{inertiay}, \sphinxstyleemphasis{inertiaz} attributes are components of
diagonalized inertia tensor for the body (i.e., the three moments of inertia
for the body around its principal axes), as computed internally by
LAMMPS.


\bigskip\hrule\bigskip



\subsection{Output info}
\label{\detokenize{compute_rigid_local:output-info}}
\sphinxAtStartPar
This compute calculates a local vector or local array depending on the
number of keywords.  The length of the vector or number of rows in the
array is the number of rigid bodies.  If a single keyword is
specified, a local vector is produced.  If two or more keywords are
specified, a local array is produced where the number of columns = the
number of keywords.  The vector or array can be accessed by any
command that uses local values from a compute as input.  See the
{\hyperref[\detokenize{Howto_output::doc}]{\sphinxcrossref{\DUrole{doc}{Howto output}}}} page for an overview of LAMMPS
output options.

\sphinxAtStartPar
The vector or array values will be in whatever {\hyperref[\detokenize{units::doc}]{\sphinxcrossref{\DUrole{doc}{units}}}} the
corresponding attribute is in:
\begin{itemize}
\item {} 
\sphinxAtStartPar
id,mol = unitless

\item {} 
\sphinxAtStartPar
mass = mass units

\item {} 
\sphinxAtStartPar
x,y,z and xy,yu,zu = distance units

\item {} 
\sphinxAtStartPar
vx,vy,vz = velocity units

\item {} 
\sphinxAtStartPar
fx,fy,fz = force units

\item {} 
\sphinxAtStartPar
omegax,omegay,omegaz = radians/time units

\item {} 
\sphinxAtStartPar
angmomx,angmomy,angmomz = mass*distance\(^2\)/time units

\item {} 
\sphinxAtStartPar
quatw,quati,quatj,quatk = unitless

\item {} 
\sphinxAtStartPar
tqx,tqy,tqz = torque units

\item {} 
\sphinxAtStartPar
inertiax,inertiay,inertiaz = mass*distance\(^2\) units

\end{itemize}


\subsection{Restrictions}
\label{\detokenize{compute_rigid_local:restrictions}}
\sphinxAtStartPar
This compute is part of the RIGID package.  It is only enabled if
LAMMPS was built with that package.  See the {\hyperref[\detokenize{Build_package::doc}]{\sphinxcrossref{\DUrole{doc}{Build package}}}} page for more info.


\subsection{Related commands}
\label{\detokenize{compute_rigid_local:related-commands}}
\sphinxAtStartPar
{\hyperref[\detokenize{dump::doc}]{\sphinxcrossref{\DUrole{doc}{dump local}}}}, {\hyperref[\detokenize{compute_reduce::doc}]{\sphinxcrossref{\DUrole{doc}{compute reduce}}}}


\subsection{Default}
\label{\detokenize{compute_rigid_local:default}}
\sphinxAtStartPar
none

\sphinxstepscope

\index{compute saed@\spxentry{compute saed}}\ignorespaces 

\section{compute saed command}
\label{\detokenize{compute_saed:compute-saed-command}}\label{\detokenize{compute_saed:index-0}}\label{\detokenize{compute_saed::doc}}

\subsection{Syntax}
\label{\detokenize{compute_saed:syntax}}
\begin{sphinxVerbatim}[commandchars=\\\{\}]
\PYG{k}{compute }\PYG{n+nv+nvIdentifier}{ID}\PYG{+w}{ }\PYG{n+nv+nvIdentifier}{group\PYGZhy{}ID}\PYG{+w}{ }\PYG{n}{saed}\PYG{+w}{ }\PYG{n}{lambda}\PYG{+w}{ }\PYG{n}{type1}\PYG{+w}{ }\PYG{n}{type2}\PYG{+w}{ }\PYG{n}{...}\PYG{+w}{ }\PYG{n}{typeN}\PYG{+w}{ }\PYG{n}{keyword}\PYG{+w}{ }\PYG{n}{value}\PYG{+w}{ }\PYG{n}{...}
\end{sphinxVerbatim}
\begin{itemize}
\item {} 
\sphinxAtStartPar
ID, group\sphinxhyphen{}ID are documented in {\hyperref[\detokenize{compute::doc}]{\sphinxcrossref{\DUrole{doc}{compute}}}} command

\item {} 
\sphinxAtStartPar
saed = style name of this compute command

\item {} 
\sphinxAtStartPar
lambda = wavelength of incident radiation (length units)

\item {} 
\sphinxAtStartPar
type1 type2 … typeN = chemical symbol of each atom type (see valid options below)

\item {} 
\sphinxAtStartPar
zero or more keyword/value pairs may be appended

\item {} 
\sphinxAtStartPar
keyword = \sphinxstyleemphasis{Kmax} or \sphinxstyleemphasis{Zone} or \sphinxstyleemphasis{dR\_Ewald} or \sphinxstyleemphasis{c} or \sphinxstyleemphasis{manual} or \sphinxstyleemphasis{echo}
\begin{sphinxalltt}
\sphinxstyleemphasis{Kmax} value = Maximum distance explored from reciprocal space origin
               (inverse length units)
\sphinxstyleemphasis{Zone} values = z1 z2 z3
  z1,z2,z3 = Zone axis of incident radiation. If z1=z2=z3=0 all
             reciprocal space will be meshed up to \sphinxstyleemphasis{Kmax}
\sphinxstyleemphasis{dR\_Ewald} value = Thickness of Ewald sphere slice intercepting
                   reciprocal space (inverse length units)
\sphinxstyleemphasis{c} values = c1 c2 c3
  c1,c2,c3 = parameters to adjust the spacing of the reciprocal
             lattice nodes in the h, k, and l directions respectively
\sphinxstyleemphasis{manual} = flag to use manual spacing of reciprocal lattice points
           based on the values of the \sphinxstyleemphasis{c} parameters
\sphinxstyleemphasis{echo} = flag to provide extra output for debugging purposes
\end{sphinxalltt}

\end{itemize}


\subsection{Examples}
\label{\detokenize{compute_saed:examples}}
\begin{sphinxVerbatim}[commandchars=\\\{\}]
\PYG{k}{compute }\PYG{n+nv+nvIdentifier}{1}\PYG{+w}{ }\PYG{n+nv+nvIdentifier}{all}\PYG{+w}{ }\PYG{n}{saed}\PYG{+w}{ }\PYG{l+m}{0.0251}\PYG{+w}{ }\PYG{n}{Al}\PYG{+w}{ }\PYG{n}{O}\PYG{+w}{ }\PYG{n}{Kmax}\PYG{+w}{ }\PYG{l+m}{1.70}\PYG{+w}{ }\PYG{n}{Zone}\PYG{+w}{ }\PYG{l+m}{0}\PYG{+w}{ }\PYG{l+m}{0}\PYG{+w}{ }\PYG{l+m}{1}\PYG{+w}{ }\PYG{n}{dR\PYGZus{}Ewald}\PYG{+w}{ }\PYG{l+m}{0.01}\PYG{+w}{ }\PYG{n}{c}\PYG{+w}{ }\PYG{l+m}{0.5}\PYG{+w}{ }\PYG{l+m}{0.5}\PYG{+w}{ }\PYG{l+m}{0.5}
\PYG{k}{compute }\PYG{n+nv+nvIdentifier}{2}\PYG{+w}{ }\PYG{n+nv+nvIdentifier}{all}\PYG{+w}{ }\PYG{n}{saed}\PYG{+w}{ }\PYG{l+m}{0.0251}\PYG{+w}{ }\PYG{n}{Ni}\PYG{+w}{ }\PYG{n}{Kmax}\PYG{+w}{ }\PYG{l+m}{1.70}\PYG{+w}{ }\PYG{n}{Zone}\PYG{+w}{ }\PYG{l+m}{0}\PYG{+w}{ }\PYG{l+m}{0}\PYG{+w}{ }\PYG{l+m}{0}\PYG{+w}{ }\PYG{n}{c}\PYG{+w}{ }\PYG{l+m}{0.05}\PYG{+w}{ }\PYG{l+m}{0.05}\PYG{+w}{ }\PYG{l+m}{0.05}\PYG{+w}{ }\PYG{n}{manual}\PYG{+w}{ }\PYG{n}{echo}

\PYG{k}{fix }\PYG{n+nv+nvIdentifier}{1}\PYG{+w}{ }\PYG{n+nv+nvIdentifier}{all}\PYG{+w}{ }\PYG{n}{saed}\PYG{o}{/}\PYG{n}{vtk}\PYG{+w}{ }\PYG{l+m}{1}\PYG{+w}{ }\PYG{l+m}{1}\PYG{+w}{ }\PYG{l+m}{1}\PYG{+w}{ }\PYG{n}{c\PYGZus{}1}\PYG{+w}{ }\PYG{n}{file}\PYG{+w}{ }\PYG{n}{Al2O3\PYGZus{}001.saed}
\PYG{k}{fix }\PYG{n+nv+nvIdentifier}{2}\PYG{+w}{ }\PYG{n+nv+nvIdentifier}{all}\PYG{+w}{ }\PYG{n}{saed}\PYG{o}{/}\PYG{n}{vtk}\PYG{+w}{ }\PYG{l+m}{1}\PYG{+w}{ }\PYG{l+m}{1}\PYG{+w}{ }\PYG{l+m}{1}\PYG{+w}{ }\PYG{n}{c\PYGZus{}2}\PYG{+w}{ }\PYG{n}{file}\PYG{+w}{ }\PYG{n}{Ni\PYGZus{}000.saed}
\end{sphinxVerbatim}


\subsection{Description}
\label{\detokenize{compute_saed:description}}
\sphinxAtStartPar
Define a computation that calculates electron diffraction intensity as
described in {\hyperref[\detokenize{compute_saed:saed-coleman}]{\sphinxcrossref{\DUrole{std,std-ref}{(Coleman)}}}} on a mesh of reciprocal lattice nodes
defined by the entire simulation domain (or manually) using simulated
radiation of wavelength lambda.

\sphinxAtStartPar
The electron diffraction intensity I at each reciprocal lattice point
is computed from the structure factor F using the equations:
\begin{equation*}
\begin{split}I = & \frac{F^{*}F}{N} \\
F(\mathbf{k}) = & \sum_{j=1}^{N}f_j(\theta)exp(2\pi i \mathbf{k} \cdot \mathbf{r}_j)\end{split}
\end{equation*}
\sphinxAtStartPar
Here, K is the location of the reciprocal lattice node, \(r_j\) is the
position of each atom, \(f_j\) are atomic scattering factors.

\sphinxAtStartPar
Diffraction intensities are calculated on a three\sphinxhyphen{}dimensional mesh of
reciprocal lattice nodes. The mesh spacing is defined either (a)  by
the entire simulation domain or (b) manually using selected values as
shown in the 2D diagram below.

\noindent{\hspace*{\fill}\sphinxincludegraphics[scale=0.75]{{saed_mesh}.png}\hspace*{\fill}}

\sphinxAtStartPar
For a mesh defined by the simulation domain, a rectilinear grid is
constructed with spacing \sphinxstyleemphasis{c}*inv(A) along each reciprocal lattice
axis. Where A are the vectors corresponding to the edges of the
simulation cell. If one or two directions has non\sphinxhyphen{}periodic boundary
conditions, then the spacing in these directions is defined from the
average of the (inversed) box lengths with periodic boundary conditions.
Meshes defined by the simulation domain must contain at least one periodic
boundary.

\sphinxAtStartPar
If the \sphinxstyleemphasis{manual} flag is included, the mesh of reciprocal lattice nodes
will defined using the \sphinxstyleemphasis{c} values for the spacing along each reciprocal
lattice axis. Note that manual mapping of the reciprocal space mesh is
good for comparing diffraction results from  multiple simulations; however
it can reduce the likelihood that Bragg reflections will be satisfied
unless small spacing parameters (\(<0.05~\AA^-1\))
are implemented.  Meshes with manual spacing do not require a periodic
boundary.

\sphinxAtStartPar
The limits of the reciprocal lattice mesh are determined by the use of
the \sphinxstyleemphasis{Kmax}, \sphinxstyleemphasis{Zone}, and \sphinxstyleemphasis{dR\_Ewald} parameters.  The rectilinear mesh
created about the origin of reciprocal space is terminated at the
boundary of a sphere of radius \sphinxstyleemphasis{Kmax} centered at the origin.  If
\sphinxstyleemphasis{Zone} parameters \sphinxstyleemphasis{z1} = \sphinxstyleemphasis{z2} = \sphinxstyleemphasis{z3} = 0 are used, diffraction intensities are
computed throughout the entire spherical volume \sphinxhyphen{} note this can
greatly increase the cost of computation.  Otherwise, \sphinxstyleemphasis{Zone}
parameters will denote the \(z1=h\), \(z2=k\), and \(z3=\ell\)
(in a global sense) zone axis of an intersecting Ewald sphere.  Diffraction
intensities will only be computed at the intersection of the reciprocal lattice
mesh and a \sphinxstyleemphasis{dR\_Ewald} thick surface of the Ewald sphere.  See the
example 3D intensity data and the intersection of a {[}010{]} zone axis
in the below image.

\noindent{\hspace*{\fill}\sphinxincludegraphics[scale=0.75]{{saed_ewald_intersect}.jpg}\hspace*{\fill}}

\sphinxAtStartPar
The atomic scattering factors, fj, accounts for the reduction in
diffraction intensity due to Compton scattering.  Compute saed uses
analytical approximations of the atomic scattering factors that vary
for each atom type (type1 type2 … typeN) and angle of diffraction.
The analytic approximation is computed using the formula
{\hyperref[\detokenize{compute_saed:brown}]{\sphinxcrossref{\DUrole{std,std-ref}{(Brown)}}}}:
\begin{equation*}
\begin{split}f_j\left ( \frac{sin(\theta)}{\lambda} \right )=\sum_{i}^{5}
a_i exp\left ( -b_i \frac{sin^{2}(\theta)}{\lambda^{2}} \right )\end{split}
\end{equation*}
\sphinxAtStartPar
Coefficients parameterized by {\hyperref[\detokenize{compute_saed:fox}]{\sphinxcrossref{\DUrole{std,std-ref}{(Fox)}}}} are assigned for each
atom type designating the chemical symbol and charge of each atom
type. Valid chemical symbols for compute saed are:


\begin{savenotes}\sphinxattablestart
\sphinxthistablewithglobalstyle
\centering
\begin{tabulary}{\linewidth}[t]{TTTTTTTTTTTTTTTTTTTT}
\sphinxtoprule
\sphinxtableatstartofbodyhook
\sphinxAtStartPar
H
&
\sphinxAtStartPar
He
&
\sphinxAtStartPar
Li
&
\sphinxAtStartPar
Be
&
\sphinxAtStartPar
B
&
\sphinxAtStartPar
C
&
\sphinxAtStartPar
N
&
\sphinxAtStartPar
O
&
\sphinxAtStartPar
F
&
\sphinxAtStartPar
Ne
&
\sphinxAtStartPar
Na
&
\sphinxAtStartPar
Mg
&
\sphinxAtStartPar
Al
&
\sphinxAtStartPar
Si
&
\sphinxAtStartPar
P
&
\sphinxAtStartPar
S
&
\sphinxAtStartPar
Cl
&
\sphinxAtStartPar
Ar
&
\sphinxAtStartPar
K
&
\sphinxAtStartPar
Ca
\\
\sphinxhline
\sphinxAtStartPar
Sc
&
\sphinxAtStartPar
Ti
&
\sphinxAtStartPar
V
&
\sphinxAtStartPar
Cr
&
\sphinxAtStartPar
Mn
&
\sphinxAtStartPar
Fe
&
\sphinxAtStartPar
Co
&
\sphinxAtStartPar
Ni
&
\sphinxAtStartPar
Cu
&
\sphinxAtStartPar
Zn
&
\sphinxAtStartPar
Ga
&
\sphinxAtStartPar
Ge
&
\sphinxAtStartPar
As
&
\sphinxAtStartPar
Se
&
\sphinxAtStartPar
Br
&
\sphinxAtStartPar
Kr
&
\sphinxAtStartPar
Rb
&
\sphinxAtStartPar
Sr
&
\sphinxAtStartPar
Y
&
\sphinxAtStartPar
Zr
\\
\sphinxhline
\sphinxAtStartPar
Nb
&
\sphinxAtStartPar
Mo
&
\sphinxAtStartPar
Tc
&
\sphinxAtStartPar
Ru
&
\sphinxAtStartPar
Rh
&
\sphinxAtStartPar
Pd
&
\sphinxAtStartPar
Ag
&
\sphinxAtStartPar
Cd
&
\sphinxAtStartPar
In
&
\sphinxAtStartPar
Sn
&
\sphinxAtStartPar
Sb
&
\sphinxAtStartPar
Te
&
\sphinxAtStartPar
I
&
\sphinxAtStartPar
Xe
&
\sphinxAtStartPar
Cs
&
\sphinxAtStartPar
Ba
&
\sphinxAtStartPar
La
&
\sphinxAtStartPar
Ce
&
\sphinxAtStartPar
Pr
&
\sphinxAtStartPar
Nd
\\
\sphinxhline
\sphinxAtStartPar
Pm
&
\sphinxAtStartPar
Sm
&
\sphinxAtStartPar
Eu
&
\sphinxAtStartPar
Gd
&
\sphinxAtStartPar
Tb
&
\sphinxAtStartPar
Dy
&
\sphinxAtStartPar
Ho
&
\sphinxAtStartPar
Er
&
\sphinxAtStartPar
Tm
&
\sphinxAtStartPar
Yb
&
\sphinxAtStartPar
Lu
&
\sphinxAtStartPar
Hf
&
\sphinxAtStartPar
Ta
&
\sphinxAtStartPar
W
&
\sphinxAtStartPar
Re
&
\sphinxAtStartPar
Os
&
\sphinxAtStartPar
Ir
&
\sphinxAtStartPar
Pt
&
\sphinxAtStartPar
Au
&
\sphinxAtStartPar
Hg
\\
\sphinxhline
\sphinxAtStartPar
Tl
&
\sphinxAtStartPar
Pb
&
\sphinxAtStartPar
Bi
&
\sphinxAtStartPar
Po
&
\sphinxAtStartPar
At
&
\sphinxAtStartPar
Rn
&
\sphinxAtStartPar
Fr
&
\sphinxAtStartPar
Ra
&
\sphinxAtStartPar
Ac
&
\sphinxAtStartPar
Th
&
\sphinxAtStartPar
Pa
&
\sphinxAtStartPar
U
&
\sphinxAtStartPar
Np
&
\sphinxAtStartPar
Pu
&
\sphinxAtStartPar
Am
&
\sphinxAtStartPar
Cm
&
\sphinxAtStartPar
Bk
&
\sphinxAtStartPar
Cf
&&\\
\sphinxbottomrule
\end{tabulary}
\sphinxtableafterendhook\par
\sphinxattableend\end{savenotes}

\sphinxAtStartPar
If the \sphinxstyleemphasis{echo} keyword is specified, compute saed will provide extra
reporting information to the screen.


\subsection{Output info}
\label{\detokenize{compute_saed:output-info}}
\sphinxAtStartPar
This compute calculates a global vector.  The length of the vector is
the number of reciprocal lattice nodes that are explored by the mesh.
The entries of the global vector are the computed diffraction
intensities as described above.

\sphinxAtStartPar
The vector can be accessed by any command that uses global values from
a compute as input.  See the {\hyperref[\detokenize{Howto_output::doc}]{\sphinxcrossref{\DUrole{doc}{Howto output}}}} doc page
for an overview of LAMMPS output options.

\sphinxAtStartPar
All array values calculated by this compute are “intensive”.


\subsection{Restrictions}
\label{\detokenize{compute_saed:restrictions}}
\sphinxAtStartPar
This compute is part of the DIFFRACTION package.  It is only enabled if
LAMMPS was built with that package.  See the {\hyperref[\detokenize{Build_package::doc}]{\sphinxcrossref{\DUrole{doc}{Build package}}}} page for more info.

\sphinxAtStartPar
The compute\_saed command does not work for triclinic cells.


\subsection{Related commands}
\label{\detokenize{compute_saed:related-commands}}
\sphinxAtStartPar
{\hyperref[\detokenize{fix_saed_vtk::doc}]{\sphinxcrossref{\DUrole{doc}{fix saed\_vtk}}}}, {\hyperref[\detokenize{compute_xrd::doc}]{\sphinxcrossref{\DUrole{doc}{compute xrd}}}}


\subsection{Default}
\label{\detokenize{compute_saed:default}}
\sphinxAtStartPar
The option defaults are Kmax = 1.70, Zone 1 0 0, c 1 1 1, dR\_Ewald =
0.01.


\bigskip\hrule\bigskip

\phantomsection\label{\detokenize{compute_saed:saed-coleman}}
\sphinxAtStartPar
\sphinxstylestrong{(Coleman)} Coleman, Spearot, Capolungo, MSMSE, 21, 055020
(2013).

\phantomsection\label{\detokenize{compute_saed:brown}}
\sphinxAtStartPar
\sphinxstylestrong{(Brown)} Brown et al. International Tables for Crystallography
Volume C: Mathematical and Chemical Tables, 554\sphinxhyphen{}95 (2004).

\phantomsection\label{\detokenize{compute_saed:fox}}
\sphinxAtStartPar
\sphinxstylestrong{(Fox)} Fox, O’Keefe, Tabbernor, Acta Crystallogr. A, 45, 786\sphinxhyphen{}93
(1989).

\sphinxstepscope

\index{compute sans@\spxentry{compute sans}}\ignorespaces 

\section{compute sans command}
\label{\detokenize{compute_sans:compute-sans-command}}\label{\detokenize{compute_sans:index-0}}\label{\detokenize{compute_sans::doc}}

\subsection{Syntax}
\label{\detokenize{compute_sans:syntax}}
\begin{sphinxVerbatim}[commandchars=\\\{\}]
\PYG{k}{compute }\PYG{n+nv+nvIdentifier}{ID}\PYG{+w}{ }\PYG{n+nv+nvIdentifier}{group\PYGZhy{}ID}\PYG{+w}{ }\PYG{n}{sans}\PYG{+w}{ }\PYG{n}{kmax}\PYG{+w}{ }\PYG{n}{keyword}\PYG{+w}{ }\PYG{n}{value}\PYG{+w}{ }\PYG{n}{...}
\end{sphinxVerbatim}
\begin{itemize}
\item {} 
\sphinxAtStartPar
ID, group\sphinxhyphen{}ID are documented in {\hyperref[\detokenize{compute::doc}]{\sphinxcrossref{\DUrole{doc}{compute}}}} command

\item {} 
\sphinxAtStartPar
sans = style name of this compute command

\item {} 
\sphinxAtStartPar
kmax = maximum wave vector magnitude to explore (inverse length units)

\item {} 
\sphinxAtStartPar
zero or more keyword/value pairs may be appended

\item {} 
\sphinxAtStartPar
keyword = \sphinxstyleemphasis{qmin} or \sphinxstyleemphasis{qmax} or \sphinxstyleemphasis{Nq} or \sphinxstyleemphasis{dR\_Ewald} or \sphinxstyleemphasis{maxdeg} or \sphinxstyleemphasis{logdist} or \sphinxstyleemphasis{lengthpath}
\begin{sphinxalltt}
\sphinxstyleemphasis{qmin} value = minimum wave vector magnitude to calculate (inverse length units)
               (default: 1.0 Å⁻\(\sp{\text{1}}\))
\sphinxstyleemphasis{qmax} value = maximum wave vector magnitude to calculate (inverse length units)
               (default: 30.0 Å⁻\(\sp{\text{1}}\))
\sphinxstyleemphasis{Nq} value = number of wave vectors distributed between qmin and qmax
             (default: 100)
\sphinxstyleemphasis{dR\_Ewald} value = thickness of Ewald sphere slice around target q values
                   (inverse length units)
                   (default: 0.2 Å⁻\(\sp{\text{1}}\))
\sphinxstyleemphasis{maxdeg} value = maximum degeneracy allowed for wave vector selection
                 (default: varies)
\sphinxstyleemphasis{logdist} = flag to use logarithmic distribution of wave vectors instead
            of linear distribution
\sphinxstyleemphasis{lengthpath} file = path to file containing custom neutron scattering lengths
\end{sphinxalltt}

\end{itemize}


\subsection{Examples}
\label{\detokenize{compute_sans:examples}}
\begin{sphinxVerbatim}[commandchars=\\\{\}]
\PYG{k}{compute }\PYG{n+nv+nvIdentifier}{1}\PYG{+w}{ }\PYG{n+nv+nvIdentifier}{all}\PYG{+w}{ }\PYG{n}{sans}\PYG{+w}{ }\PYG{l+m}{10.0}\PYG{+w}{ }\PYG{n}{qmin}\PYG{+w}{ }\PYG{l+m}{1.0}\PYG{+w}{ }\PYG{n}{qmax}\PYG{+w}{ }\PYG{l+m}{30.0}\PYG{+w}{ }\PYG{n}{Nq}\PYG{+w}{ }\PYG{l+m}{50}
\PYG{k}{compute }\PYG{n+nv+nvIdentifier}{2}\PYG{+w}{ }\PYG{n+nv+nvIdentifier}{all}\PYG{+w}{ }\PYG{n}{sans}\PYG{+w}{ }\PYG{l+m}{8.0}\PYG{+w}{ }\PYG{n}{qmax}\PYG{+w}{ }\PYG{l+m}{25.0}\PYG{+w}{ }\PYG{n}{dR\PYGZus{}Ewald}\PYG{+w}{ }\PYG{l+m}{0.15}\PYG{+w}{ }\PYG{n}{logdist}
\PYG{k}{compute }\PYG{n+nv+nvIdentifier}{3}\PYG{+w}{ }\PYG{n+nv+nvIdentifier}{all}\PYG{+w}{ }\PYG{n}{sans}\PYG{+w}{ }\PYG{l+m}{12.0}\PYG{+w}{ }\PYG{n}{Nq}\PYG{+w}{ }\PYG{l+m}{100}\PYG{+w}{ }\PYG{n}{lengthpath}\PYG{+w}{ }\PYG{n}{scattering\PYGZus{}lengths.txt}
\end{sphinxVerbatim}


\subsection{Description}
\label{\detokenize{compute_sans:description}}
\sphinxAtStartPar
Define a computation that calculates small\sphinxhyphen{}angle neutron scattering (SANS)
intensity as a function of wave vector magnitude (q), based on the atomic
structure of the system. This computation is particularly useful for
comparing molecular dynamics simulations with neutron diffraction
experiments.

\sphinxAtStartPar
The SANS intensity I(q) at each wave vector magnitude is computed from
the structure factor F(q) using:
\begin{equation*}
\begin{split}I(q) = \frac{1}{N} \left| \sum_{i}^{N} b_i exp(2\pi i \mathbf{q} \cdot \mathbf{r}_i) \right|^2\end{split}
\end{equation*}
\sphinxAtStartPar
where N is the number of atoms in the group, \(\mathbf{r}_i\) is the
position of atom i, \(b_i\) is the neutron scattering length for atom
type i, and \(\mathbf{q}\) is the scattering wave vector with magnitude q.

\sphinxAtStartPar
\sphinxstylestrong{Wave Vector Generation}

\sphinxAtStartPar
The compute generates a set of wave vectors distributed within reciprocal
space. By default, wave vectors are distributed linearly between \sphinxstyleemphasis{qmin}
and \sphinxstyleemphasis{qmax} in steps determined by \sphinxstyleemphasis{Nq}. If the \sphinxstyleemphasis{logdist} flag is specified,
wave vectors are instead distributed on a logarithmic scale, which is often
more suitable for exploring a wide range of length scales.

\sphinxAtStartPar
The parameter \sphinxstyleemphasis{kmax} determines the maximum distance explored from the
origin of reciprocal space in units of inverse length. Wave vectors are
selected to lie approximately at the specified q magnitudes, within a
tolerance defined by \sphinxstyleemphasis{dR\_Ewald} (the thickness of the Ewald sphere slice).

\sphinxAtStartPar
\sphinxstylestrong{Scattering Lengths}

\sphinxAtStartPar
By default, scattering lengths are set to 1.0 for all atom types, which
gives an unweighted calculation. For comparison with actual neutron
scattering experiments, custom neutron scattering lengths can be provided
via the \sphinxstyleemphasis{lengthpath} keyword, which specifies a file containing scattering
length values for each atom type.

\sphinxAtStartPar
The scattering lengths file should contain one line per atom type with the
format:

\begin{sphinxVerbatim}[commandchars=\\\{\}]
type\PYGZus{}number scattering\PYGZus{}length
\end{sphinxVerbatim}

\sphinxAtStartPar
where type\_number is the LAMMPS atom type (starting from 1) and
scattering\_length is the neutron scattering length in inverse length
units (typically in fm or 10⁻\(\sp{\text{1}}\)\(\sp{\text{5}}\) m for neutrons).

\sphinxAtStartPar
\sphinxstylestrong{Degeneracy and Ewald Sphere}

\sphinxAtStartPar
The \sphinxstyleemphasis{dR\_Ewald} parameter controls how many wave vectors are selected
around each target q value. A smaller value results in fewer selected
wave vectors but more precise targeting of specific q values. A larger
value samples a thicker sphere in reciprocal space around each q value.

\sphinxAtStartPar
The \sphinxstyleemphasis{maxdeg} parameter can be used to limit the degeneracy (number of
equivalent wave vectors) selected for each q value.

\sphinxAtStartPar
\sphinxstylestrong{Parallel Calculation}

\sphinxAtStartPar
This compute uses MPI to efficiently distribute the scattering calculation
across multiple processors. The cos and sin components of the structure
factor are accumulated separately on each processor and then reduced
globally before computing the final intensity.


\subsection{Output info}
\label{\detokenize{compute_sans:output-info}}
\sphinxAtStartPar
This compute calculates a global array with dimensions of Nq rows and
3 columns. Each row contains:
\begin{itemize}
\item {} 
\sphinxAtStartPar
Column 1: Index (i) of the wave vector

\item {} 
\sphinxAtStartPar
Column 2: Wave vector magnitude (q) in inverse length units

\item {} 
\sphinxAtStartPar
Column 3: Normalized SANS intensity I(q)/N

\end{itemize}

\sphinxAtStartPar
The array can be accessed by any command that uses global values from
a compute as input. See the {\hyperref[\detokenize{Howto_output::doc}]{\sphinxcrossref{\DUrole{doc}{Howto output}}}} doc page
for an overview of LAMMPS output options.

\sphinxAtStartPar
All array values calculated by this compute are “intensive”.


\subsection{Restrictions}
\label{\detokenize{compute_sans:restrictions}}
\sphinxAtStartPar
This compute is part of the DIFFRACTION package. It is only enabled if
LAMMPS was built with that package. See the {\hyperref[\detokenize{Build_package::doc}]{\sphinxcrossref{\DUrole{doc}{Build package}}}} page for more info.

\sphinxAtStartPar
The compute\_sans command does not work for triclinic cells.

\sphinxAtStartPar
The compute\_sans command only works with 3\sphinxhyphen{}dimensional systems.


\subsection{Related commands}
\label{\detokenize{compute_sans:related-commands}}
\sphinxAtStartPar
{\hyperref[\detokenize{compute_saed::doc}]{\sphinxcrossref{\DUrole{doc}{compute saed}}}}, {\hyperref[\detokenize{compute_xrd::doc}]{\sphinxcrossref{\DUrole{doc}{compute xrd}}}}


\subsection{Default}
\label{\detokenize{compute_sans:default}}
\sphinxAtStartPar
The option defaults are qmin = 1.0, qmax = 30.0, Nq = 100, dR\_Ewald = 0.2,
logdist = off (linear distribution)

\sphinxstepscope

\index{compute slcsa/atom@\spxentry{compute slcsa/atom}}\ignorespaces 

\section{compute slcsa/atom command}
\label{\detokenize{compute_slcsa_atom:compute-slcsa-atom-command}}\label{\detokenize{compute_slcsa_atom:index-0}}\label{\detokenize{compute_slcsa_atom::doc}}

\subsection{Syntax}
\label{\detokenize{compute_slcsa_atom:syntax}}
\begin{sphinxVerbatim}[commandchars=\\\{\}]
\PYG{k}{compute }\PYG{n+nv+nvIdentifier}{ID}\PYG{+w}{ }\PYG{n+nv+nvIdentifier}{group\PYGZhy{}ID}\PYG{+w}{ }\PYG{n}{slcsa}\PYG{o}{/}\PYG{n}{atom}\PYG{+w}{ }\PYG{n}{twojmax}\PYG{+w}{ }\PYG{n}{nclasses}\PYG{+w}{ }\PYG{n}{db\PYGZus{}mean\PYGZus{}descriptor\PYGZus{}file}\PYG{+w}{ }\PYG{n}{lda\PYGZus{}file}\PYG{+w}{ }\PYG{n}{lr\PYGZus{}decision\PYGZus{}file}\PYG{+w}{ }\PYG{n}{lr\PYGZus{}bias\PYGZus{}file}\PYG{+w}{ }\PYG{n}{maha\PYGZus{}file}\PYG{+w}{ }\PYG{n}{value}
\end{sphinxVerbatim}
\begin{itemize}
\item {} 
\sphinxAtStartPar
ID, group\sphinxhyphen{}ID are documented in {\hyperref[\detokenize{compute::doc}]{\sphinxcrossref{\DUrole{doc}{compute}}}} command

\item {} 
\sphinxAtStartPar
slcsa/atom = style name of this compute command

\item {} 
\sphinxAtStartPar
twojmax = band limit for bispectrum components (non\sphinxhyphen{}negative integer)

\item {} 
\sphinxAtStartPar
nclasses = number of crystal structures used in the database for the classifier SL\sphinxhyphen{}CSA

\item {} 
\sphinxAtStartPar
db\_mean\_descriptor\_file = file name of file containing the database mean descriptor

\item {} 
\sphinxAtStartPar
lda\_file = file name of file containing the linear discriminant analysis matrix for dimension reduction

\item {} 
\sphinxAtStartPar
lr\_decision\_file = file name of file containing the scaling matrix for logistic regression classification

\item {} 
\sphinxAtStartPar
lr\_bias\_file = file name of file containing the bias vector for logistic regression classification

\item {} 
\sphinxAtStartPar
maha\_file = file name of file containing for each crystal structure: the Mahalanobis distance threshold for sanity check purposes, the average reduced descriptor and the inverse of the corresponding covariance matrix

\item {} 
\sphinxAtStartPar
c\_ID{[}*{]} = compute ID of previously required \sphinxstyleemphasis{compute sna/atom} command

\end{itemize}


\subsection{Examples}
\label{\detokenize{compute_slcsa_atom:examples}}
\begin{sphinxVerbatim}[commandchars=\\\{\}]
\PYG{k}{compute }\PYG{n+nv+nvIdentifier}{b1}\PYG{+w}{ }\PYG{n+nv+nvIdentifier}{all}\PYG{+w}{ }\PYG{n}{sna}\PYG{o}{/}\PYG{n}{atom}\PYG{+w}{ }\PYG{l+m}{9.0}\PYG{+w}{ }\PYG{l+m}{0.99363}\PYG{+w}{ }\PYG{l+m}{8}\PYG{+w}{ }\PYG{l+m}{0.5}\PYG{+w}{ }\PYG{l+m}{1.0}\PYG{+w}{ }\PYG{n}{rmin0}\PYG{+w}{ }\PYG{l+m}{0.0}\PYG{+w}{ }\PYG{n}{nnn}\PYG{+w}{ }\PYG{l+m}{24}\PYG{+w}{ }\PYG{n}{wmode}\PYG{+w}{ }\PYG{l+m}{1}\PYG{+w}{ }\PYG{n}{delta}\PYG{+w}{ }\PYG{l+m}{0.3}
\PYG{k}{compute }\PYG{n+nv+nvIdentifier}{b2}\PYG{+w}{ }\PYG{n+nv+nvIdentifier}{all}\PYG{+w}{ }\PYG{n}{slcsa}\PYG{o}{/}\PYG{n}{atom}\PYG{+w}{ }\PYG{l+m}{8}\PYG{+w}{ }\PYG{l+m}{4}\PYG{+w}{ }\PYG{n}{mean\PYGZus{}descriptors.dat}\PYG{+w}{ }\PYG{n}{lda\PYGZus{}scalings.dat}\PYG{+w}{ }\PYG{n}{lr\PYGZus{}decision.dat}\PYG{+w}{ }\PYG{n}{lr\PYGZus{}bias.dat}\PYG{+w}{ }\PYG{n}{maha\PYGZus{}thresholds.dat}\PYG{+w}{ }\PYG{n}{c\PYGZus{}b1[}\PYG{o}{*}\PYG{n}{]}
\end{sphinxVerbatim}


\subsection{Description}
\label{\detokenize{compute_slcsa_atom:description}}
\sphinxAtStartPar
\DUrole{versionmodified,added}{Added in version 7Feb2024.}

\sphinxAtStartPar
Define a computation that performs the Supervised Learning Crystal
Structure Analysis (SL\sphinxhyphen{}CSA) from {\hyperref[\detokenize{compute_slcsa_atom:lafourcade2023-1}]{\sphinxcrossref{\DUrole{std,std-ref}{(Lafourcade)}}}}
for each atom in the group. The SL\sphinxhyphen{}CSA tool takes as an input a per\sphinxhyphen{}atom
descriptor (bispectrum) that is computed through the \sphinxstyleemphasis{compute sna/atom}
command and then proceeds to a dimension reduction step followed by a
logistic regression in order to assign a probable crystal structure to
each atom in the group. The SL\sphinxhyphen{}CSA tool is pre\sphinxhyphen{}trained on a database
containing \(C\) distinct crystal structures from which a crystal
structure classifier is derived and a tutorial to build such a tool is
available at \sphinxhref{https://github.com/lafourcadep/SL-CSA}{SL\sphinxhyphen{}CSA}.

\sphinxAtStartPar
The first step of the SL\sphinxhyphen{}CSA tool consists in performing a dimension
reduction of the per\sphinxhyphen{}atom descriptor \(\mathbf{B}^i \in
\mathbb{R}^{D}\) through the Linear Discriminant Analysis (LDA) method,
leading to a new projected descriptor
\(\mathbf{x}^i=\mathrm{P}_\mathrm{LDA}(\mathbf{B}^i):\mathbb{R}^D
\rightarrow \mathbb{R}^{d=C-1}\):
\begin{equation*}
\begin{split}\mathbf{x}^i = \mathbf{C}^T_\mathrm{LDA} \cdot (\mathbf{B}^i - \mu^\mathbf{B}_\mathrm{db})\end{split}
\end{equation*}
\sphinxAtStartPar
where \(\mathbf{C}^T_\mathrm{LDA} \in \mathbb{R}^{D \times d}\) is
the reduction coefficients matrix of the LDA model read in file
\sphinxstyleemphasis{lda\_file}, \(\mathbf{B}^i \in \mathbb{R}^{D}\) is the bispectrum of
atom \(i\) and \(\mu^\mathbf{B}_\mathrm{db} \in \mathbb{R}^{D}\)
is the average descriptor of the entire database. The latter is computed
from the average descriptors of each crystal structure read from the
file \sphinxstyleemphasis{mean\_descriptors\_file}.

\sphinxAtStartPar
The new projected descriptor with dimension \(d=C-1\) allows for a
good separation of different crystal structures fingerprints in the
latent space.

\sphinxAtStartPar
Once the dimension reduction step is performed by means of LDA, the new
descriptor \(\mathbf{x}^i \in \mathbb{R}^{d=C-1}\) is taken as an
input for performing a multinomial logistic regression (LR) which
provides a score vector
\(\mathbf{s}^i=\mathrm{P}_\mathrm{LR}(\mathbf{x}^i):\mathbb{R}^d
\rightarrow \mathbb{R}^C\) defined as:
\begin{equation*}
\begin{split}\mathbf{s}^i = \mathbf{b}_\mathrm{LR} + \mathbf{D}_\mathrm{LR} \cdot {\mathbf{x}^i}^T\end{split}
\end{equation*}
\sphinxAtStartPar
with \(\mathbf{b}_\mathrm{LR} \in \mathbb{R}^C\) and
\(\mathbf{D}_\mathrm{LR} \in \mathbb{R}^{C \times d}\) the bias
vector and decision matrix of the LR model after training both read in
files \sphinxstyleemphasis{lr\_fil1} and \sphinxstyleemphasis{lr\_file2} respectively.

\sphinxAtStartPar
Finally, a probability vector
\(\mathbf{p}^i=\mathrm{P}_\mathrm{LR}(\mathbf{x}^i):\mathbb{R}^d
\rightarrow \mathbb{R}^C\) is defined as:
\begin{equation*}
\begin{split}\mathbf{p}^i = \frac{\mathrm{exp}(\mathbf{s}^i)}{\sum\limits_{j} \mathrm{exp}(s^i_j) }\end{split}
\end{equation*}
\sphinxAtStartPar
from which the crystal structure assigned to each atom with descriptor
\(\mathbf{B}^i\) and projected descriptor \(\mathbf{x}^i\) is
computed as the \sphinxstyleemphasis{argmax} of the probability vector
\(\mathbf{p}^i\). Since the logistic regression step systematically
attributes a crystal structure to each atom, a sanity check is needed to
avoid misclassification. To this end, a per\sphinxhyphen{}atom Mahalanobis distance to
each crystal structure \sphinxstyleemphasis{CS} present in the database is computed:
\begin{equation*}
\begin{split}d_\mathrm{Mahalanobis}^{i \rightarrow \mathrm{CS}} = \sqrt{(\mathbf{x}^i - \mathbf{\mu}^\mathbf{x}_\mathrm{CS})^\mathrm{T} \cdot \mathbf{\Sigma}^{-1}_\mathrm{CS} \cdot (\mathbf{x}^i - \mathbf{\mu}^\mathbf{x}_\mathrm{CS}) }\end{split}
\end{equation*}
\sphinxAtStartPar
where \(\mathbf{\mu}^\mathbf{x}_\mathrm{CS} \in \mathbb{R}^{d}\) is
the average projected descriptor of crystal structure \sphinxstyleemphasis{CS} in the
database and where \(\mathbf{\Sigma}_\mathrm{CS} \in \mathbb{R}^{d
\times d}\) is the corresponding covariance matrix. Finally, if the
Mahalanobis distance to crystal structure \sphinxstyleemphasis{CS} for atom \sphinxstyleemphasis{i} is greater
than the pre\sphinxhyphen{}determined threshold, no crystal structure is assigned to
atom \sphinxstyleemphasis{i}. The Mahalanobis distance thresholds are read in file
\sphinxstyleemphasis{maha\_file} while the covariance matrices are read in file
\sphinxstyleemphasis{covmat\_file}.

\sphinxAtStartPar
The \sphinxhref{https://github.com/lafourcadep/SL-CSA}{SL\sphinxhyphen{}CSA} framework provides
an automatic computation of the different matrices and thresholds
required for a proper classification and writes down all the required
files for calling the \sphinxstyleemphasis{compute slcsa/atom} command.

\sphinxAtStartPar
The \sphinxstyleemphasis{compute slcsa/atom} command requires that the {\hyperref[\detokenize{compute_sna_atom::doc}]{\sphinxcrossref{\DUrole{doc}{compute
sna/atom}}}} command is called before as it takes the
resulting per\sphinxhyphen{}atom bispectrum as an input. In addition, it is crucial
that the value \sphinxstyleemphasis{twojmax} is set to the same value of the value \sphinxstyleemphasis{twojmax}
used in the \sphinxstyleemphasis{compute sna/atom} command, as well as that the value
\sphinxstyleemphasis{nclasses} is set to the number of crystal structures used in the
database to train the SL\sphinxhyphen{}CSA tool.


\subsection{Output info}
\label{\detokenize{compute_slcsa_atom:output-info}}
\sphinxAtStartPar
By default, this compute computes the Mahalanobis distances to the
different crystal structures present in the database in addition to
assigning a crystal structure for each atom as a per\sphinxhyphen{}atom vector, which
can be accessed by any command that uses per\sphinxhyphen{}atom values from a compute
as input.  See the {\hyperref[\detokenize{Howto_output::doc}]{\sphinxcrossref{\DUrole{doc}{Howto output}}}} page for an
overview of LAMMPS output options.


\subsection{Restrictions}
\label{\detokenize{compute_slcsa_atom:restrictions}}
\sphinxAtStartPar
This compute is part of the EXTRA\sphinxhyphen{}COMPUTE package.  It is only enabled
if LAMMPS was built with that package.  See the {\hyperref[\detokenize{Build_package::doc}]{\sphinxcrossref{\DUrole{doc}{Build package}}}} page for more info.


\subsection{Related commands}
\label{\detokenize{compute_slcsa_atom:related-commands}}
\sphinxAtStartPar
{\hyperref[\detokenize{compute_sna_atom::doc}]{\sphinxcrossref{\DUrole{doc}{compute sna/atom}}}}


\subsection{Default}
\label{\detokenize{compute_slcsa_atom:default}}
\sphinxAtStartPar
none


\bigskip\hrule\bigskip

\phantomsection\label{\detokenize{compute_slcsa_atom:lafourcade2023-1}}
\sphinxAtStartPar
\sphinxstylestrong{(Lafourcade)} Lafourcade, Maillet, Denoual, Duval, Allera, Goryaeva, and Marinica,
\sphinxhref{https://doi.org/10.1016/j.commatsci.2023.112534}{Comp. Mat. Science, 230, 112534 (2023)}

\sphinxstepscope

\index{compute slice@\spxentry{compute slice}}\ignorespaces 

\section{compute slice command}
\label{\detokenize{compute_slice:compute-slice-command}}\label{\detokenize{compute_slice:index-0}}\label{\detokenize{compute_slice::doc}}

\subsection{Syntax}
\label{\detokenize{compute_slice:syntax}}
\begin{sphinxVerbatim}[commandchars=\\\{\}]
\PYG{k}{compute }\PYG{n+nv+nvIdentifier}{ID}\PYG{+w}{ }\PYG{n+nv+nvIdentifier}{group\PYGZhy{}ID}\PYG{+w}{ }\PYG{n}{slice}\PYG{+w}{ }\PYG{n}{Nstart}\PYG{+w}{ }\PYG{n}{Nstop}\PYG{+w}{ }\PYG{n}{Nskip}\PYG{+w}{ }\PYG{n}{input1}\PYG{+w}{ }\PYG{n}{input2}\PYG{+w}{ }\PYG{n}{...}
\end{sphinxVerbatim}
\begin{itemize}
\item {} 
\sphinxAtStartPar
ID, group\sphinxhyphen{}ID are documented in {\hyperref[\detokenize{compute::doc}]{\sphinxcrossref{\DUrole{doc}{compute}}}} command

\item {} 
\sphinxAtStartPar
slice = style name of this compute command

\item {} 
\sphinxAtStartPar
Nstart = starting index within input vector(s)

\item {} 
\sphinxAtStartPar
Nstop = stopping index within input vector(s)

\item {} 
\sphinxAtStartPar
Nskip = extract every Nskip elements from input vector(s)

\item {} 
\sphinxAtStartPar
input = c\_ID, c\_ID{[}N{]}, f\_ID, f\_ID{[}N{]}

\begin{sphinxVerbatim}[commandchars=\\\{\}]
c\PYGZus{}ID = global vector calculated by a compute with ID
c\PYGZus{}ID[I] = Ith column of global array calculated by a compute with ID
f\PYGZus{}ID = global vector calculated by a fix with ID
f\PYGZus{}ID[I] = Ith column of global array calculated by a fix with ID
v\PYGZus{}name = vector calculated by an vector\PYGZhy{}style variable with name
\end{sphinxVerbatim}

\end{itemize}


\subsection{Examples}
\label{\detokenize{compute_slice:examples}}
\begin{sphinxVerbatim}[commandchars=\\\{\}]
\PYG{k}{compute }\PYG{n+nv+nvIdentifier}{1}\PYG{+w}{ }\PYG{n+nv+nvIdentifier}{all}\PYG{+w}{ }\PYG{n}{slice}\PYG{+w}{ }\PYG{l+m}{1}\PYG{+w}{ }\PYG{l+m}{100}\PYG{+w}{ }\PYG{l+m}{10}\PYG{+w}{ }\PYG{n}{c\PYGZus{}msdmol[4]}
\PYG{k}{compute }\PYG{n+nv+nvIdentifier}{1}\PYG{+w}{ }\PYG{n+nv+nvIdentifier}{all}\PYG{+w}{ }\PYG{n}{slice}\PYG{+w}{ }\PYG{l+m}{301}\PYG{+w}{ }\PYG{l+m}{400}\PYG{+w}{ }\PYG{l+m}{1}\PYG{+w}{ }\PYG{n}{c\PYGZus{}msdmol[4]}\PYG{+w}{ }\PYG{n}{v\PYGZus{}myVec}
\end{sphinxVerbatim}


\subsection{Description}
\label{\detokenize{compute_slice:description}}
\sphinxAtStartPar
Define a calculation that “slices” one or more vector inputs into
smaller vectors, one per listed input.  The inputs can be global
quantities; they cannot be per\sphinxhyphen{}atom or local quantities.
{\hyperref[\detokenize{compute::doc}]{\sphinxcrossref{\DUrole{doc}{Computes}}}} and {\hyperref[\detokenize{fix::doc}]{\sphinxcrossref{\DUrole{doc}{fixes}}}} and vector\sphinxhyphen{}style
{\hyperref[\detokenize{variable::doc}]{\sphinxcrossref{\DUrole{doc}{variables}}}} can generate such global quantities.  The
group specified with this command is ignored.

\sphinxAtStartPar
The values extracted from the input vector(s) are determined by the
\sphinxstyleemphasis{Nstart}, \sphinxstyleemphasis{Nstop}, and \sphinxstyleemphasis{Nskip} parameters.  The elements of an input
vector of length N are indexed from 1 to N.  Starting at element
\sphinxstyleemphasis{Nstart}, every Mth element is extracted, where M = \sphinxstyleemphasis{Nskip}, until
element \sphinxstyleemphasis{Nstop} is reached.  The extracted quantities are stored as a
vector, which is typically shorter than the input vector.

\sphinxAtStartPar
Each listed input is operated on independently to produce one output
vector.  Each listed input must be a global vector or column of a
global array calculated by another {\hyperref[\detokenize{compute::doc}]{\sphinxcrossref{\DUrole{doc}{compute}}}} or
{\hyperref[\detokenize{fix::doc}]{\sphinxcrossref{\DUrole{doc}{fix}}}}.

\sphinxAtStartPar
If an input value begins with “c\_”, a compute ID must follow which has
been previously defined in the input script and which generates a
global vector or array.  See the individual {\hyperref[\detokenize{compute::doc}]{\sphinxcrossref{\DUrole{doc}{compute}}}} doc
page for details.  If no bracketed integer is appended, the vector
calculated by the compute is used.  If a bracketed integer is
appended, the Ith column of the array calculated by the compute is
used.  Users can also write code for their own compute styles and {\hyperref[\detokenize{Modify::doc}]{\sphinxcrossref{\DUrole{doc}{add them to LAMMPS}}}}.

\sphinxAtStartPar
If a value begins with “f\_”, a fix ID must follow which has been
previously defined in the input script and which generates a global
vector or array.  See the individual {\hyperref[\detokenize{fix::doc}]{\sphinxcrossref{\DUrole{doc}{fix}}}} page for
details.  Note that some fixes only produce their values on certain
timesteps, which must be compatible with when compute slice references
the values, else an error results.  If no bracketed integer is
appended, the vector calculated by the fix is used.  If a bracketed
integer is appended, the Ith column of the array calculated by the fix
is used.  Users can also write code for their own fix style and {\hyperref[\detokenize{Modify::doc}]{\sphinxcrossref{\DUrole{doc}{add them to LAMMPS}}}}.

\sphinxAtStartPar
If an input value begins with “v\_”, a variable name must follow which
has been previously defined in the input script.  Only vector\sphinxhyphen{}style
variables can be referenced.  See the {\hyperref[\detokenize{variable::doc}]{\sphinxcrossref{\DUrole{doc}{variable}}}} command
for details.  Note that variables of style \sphinxstyleemphasis{vector} define a formula
which can reference individual atom properties or thermodynamic
keywords, or they can invoke other computes, fixes, or variables when
they are evaluated, so this is a very general means of specifying
quantities to slice.

\sphinxAtStartPar
If a single input is specified this compute produces a global vector,
even if the length of the vector is 1.  If multiple inputs are
specified, then a global array of values is produced, with the number
of columns equal to the number of inputs specified.


\bigskip\hrule\bigskip



\subsection{Output info}
\label{\detokenize{compute_slice:output-info}}
\sphinxAtStartPar
This compute calculates a global vector if a single input value is
specified or a global array with N columns where N is the number of
inputs.  The length of the vector or the number of rows in the array
is equal to the number of values extracted from each input vector.
These values can be used by any command that uses global vector or
array values from a compute as input.  See the {\hyperref[\detokenize{Howto_output::doc}]{\sphinxcrossref{\DUrole{doc}{Howto output}}}} page for an overview of LAMMPS output
options.

\sphinxAtStartPar
The vector or array values calculated by this compute are simply
copies of values generated by computes or fixes or variables that are
input vectors to this compute.  If there is a single input vector of
intensive and/or extensive values, then each value in the vector of
values calculated by this compute will be “intensive” or “extensive”,
depending on the corresponding input value.  If there are multiple
input vectors, and all the values in them are intensive, then the
array values calculated by this compute are “intensive”.  If there are
multiple input vectors, and any value in them is extensive, then the
array values calculated by this compute are “extensive”.  Values
produced by a variable are treated as intensive.

\sphinxAtStartPar
The vector or array values will be in whatever {\hyperref[\detokenize{units::doc}]{\sphinxcrossref{\DUrole{doc}{units}}}} the
input quantities are in.


\subsection{Restrictions}
\label{\detokenize{compute_slice:restrictions}}\begin{quote}

\sphinxAtStartPar
none
\end{quote}


\subsection{Related commands}
\label{\detokenize{compute_slice:related-commands}}
\sphinxAtStartPar
{\hyperref[\detokenize{compute::doc}]{\sphinxcrossref{\DUrole{doc}{compute}}}}, {\hyperref[\detokenize{fix::doc}]{\sphinxcrossref{\DUrole{doc}{fix}}}}, {\hyperref[\detokenize{compute_reduce::doc}]{\sphinxcrossref{\DUrole{doc}{compute reduce}}}}


\subsection{Default}
\label{\detokenize{compute_slice:default}}
\sphinxAtStartPar
none

\sphinxstepscope

\index{compute smd/contact/radius@\spxentry{compute smd/contact/radius}}\ignorespaces 

\section{compute smd/contact/radius command}
\label{\detokenize{compute_smd_contact_radius:compute-smd-contact-radius-command}}\label{\detokenize{compute_smd_contact_radius:index-0}}\label{\detokenize{compute_smd_contact_radius::doc}}

\subsection{Syntax}
\label{\detokenize{compute_smd_contact_radius:syntax}}
\begin{sphinxVerbatim}[commandchars=\\\{\}]
\PYG{k}{compute }\PYG{n+nv+nvIdentifier}{ID}\PYG{+w}{ }\PYG{n+nv+nvIdentifier}{group\PYGZhy{}ID}\PYG{+w}{ }\PYG{n}{smd}\PYG{o}{/}\PYG{n}{contact}\PYG{o}{/}\PYG{n}{radius}
\end{sphinxVerbatim}
\begin{itemize}
\item {} 
\sphinxAtStartPar
ID, group\sphinxhyphen{}ID are documented in {\hyperref[\detokenize{compute::doc}]{\sphinxcrossref{\DUrole{doc}{compute}}}} command

\item {} 
\sphinxAtStartPar
smd/contact/radius = style name of this compute command

\end{itemize}


\subsection{Examples}
\label{\detokenize{compute_smd_contact_radius:examples}}
\begin{sphinxVerbatim}[commandchars=\\\{\}]
\PYG{k}{compute }\PYG{n+nv+nvIdentifier}{1}\PYG{+w}{ }\PYG{n+nv+nvIdentifier}{all}\PYG{+w}{ }\PYG{n}{smd}\PYG{o}{/}\PYG{n}{contact}\PYG{o}{/}\PYG{n}{radius}
\end{sphinxVerbatim}


\subsection{Description}
\label{\detokenize{compute_smd_contact_radius:description}}
\sphinxAtStartPar
Define a computation which outputs the contact radius, i.e., the
radius used to prevent particles from penetrating each other.  The
contact radius is used only to prevent particles belonging to
different physical bodies from penetrating each other. It is used by
the contact pair styles, e.g., smd/hertz and smd/tri\_surface.

\sphinxAtStartPar
See \sphinxhref{PDF/MACHDYN\_LAMMPS\_userguide.pdf}{this PDF guide} to using Smooth
Mach Dynamics in LAMMPS.

\sphinxAtStartPar
The value of the contact radius will be 0.0 for particles not in the
specified compute group.


\subsection{Output info}
\label{\detokenize{compute_smd_contact_radius:output-info}}
\sphinxAtStartPar
This compute calculates a per\sphinxhyphen{}particle vector, which can be accessed
by any command that uses per\sphinxhyphen{}particle values from a compute as input.
See the {\hyperref[\detokenize{Howto_output::doc}]{\sphinxcrossref{\DUrole{doc}{Howto output}}}} page for an overview of
LAMMPS output options.

\sphinxAtStartPar
The per\sphinxhyphen{}particle vector values will be in distance {\hyperref[\detokenize{units::doc}]{\sphinxcrossref{\DUrole{doc}{units}}}}.


\subsection{Restrictions}
\label{\detokenize{compute_smd_contact_radius:restrictions}}
\sphinxAtStartPar
This compute is part of the MACHDYN package.  It is only enabled if
LAMMPS was built with that package.  See the {\hyperref[\detokenize{Build_package::doc}]{\sphinxcrossref{\DUrole{doc}{Build package}}}} page for more info.


\subsection{Related commands}
\label{\detokenize{compute_smd_contact_radius:related-commands}}
\sphinxAtStartPar
{\hyperref[\detokenize{dump::doc}]{\sphinxcrossref{\DUrole{doc}{dump custom}}}} smd/hertz smd/tri\_surface


\subsection{Default}
\label{\detokenize{compute_smd_contact_radius:default}}
\sphinxAtStartPar
none

\sphinxstepscope

\index{compute smd/damage@\spxentry{compute smd/damage}}\ignorespaces 

\section{compute smd/damage command}
\label{\detokenize{compute_smd_damage:compute-smd-damage-command}}\label{\detokenize{compute_smd_damage:index-0}}\label{\detokenize{compute_smd_damage::doc}}

\subsection{Syntax}
\label{\detokenize{compute_smd_damage:syntax}}
\begin{sphinxVerbatim}[commandchars=\\\{\}]
\PYG{k}{compute }\PYG{n+nv+nvIdentifier}{ID}\PYG{+w}{ }\PYG{n+nv+nvIdentifier}{group\PYGZhy{}ID}\PYG{+w}{ }\PYG{n}{smd}\PYG{o}{/}\PYG{n}{damage}
\end{sphinxVerbatim}
\begin{itemize}
\item {} 
\sphinxAtStartPar
ID, group\sphinxhyphen{}ID are documented in {\hyperref[\detokenize{compute::doc}]{\sphinxcrossref{\DUrole{doc}{compute}}}} command

\item {} 
\sphinxAtStartPar
smd/damage = style name of this compute command

\end{itemize}


\subsection{Examples}
\label{\detokenize{compute_smd_damage:examples}}
\begin{sphinxVerbatim}[commandchars=\\\{\}]
\PYG{k}{compute }\PYG{n+nv+nvIdentifier}{1}\PYG{+w}{ }\PYG{n+nv+nvIdentifier}{all}\PYG{+w}{ }\PYG{n}{smd}\PYG{o}{/}\PYG{n}{damage}
\end{sphinxVerbatim}


\subsection{Description}
\label{\detokenize{compute_smd_damage:description}}
\sphinxAtStartPar
Define a computation that calculates the damage status of SPH particles
according to the damage model which is defined via the SMD SPH pair styles, e.g., the maximum plastic strain failure criterion.

\sphinxAtStartPar
See \sphinxhref{PDF/MACHDYN\_LAMMPS\_userguide.pdf}{this PDF guide} to use Smooth Mach Dynamics in LAMMPS.

\sphinxAtStartPar
\sphinxstylestrong{Output Info:}

\sphinxAtStartPar
This compute calculates a per\sphinxhyphen{}particle vector, which can be accessed
by any command that uses per\sphinxhyphen{}particle values from a compute as input.
See the {\hyperref[\detokenize{Howto_output::doc}]{\sphinxcrossref{\DUrole{doc}{Howto output}}}} page for an overview of
LAMMPS output options.

\sphinxAtStartPar
The per\sphinxhyphen{}particle values are dimensionless an in the range of zero to one.


\subsection{Restrictions}
\label{\detokenize{compute_smd_damage:restrictions}}
\sphinxAtStartPar
This compute is part of the MACHDYN package.  It is only enabled if
LAMMPS was built with that package.  See the “Build


\subsection{Related commands}
\label{\detokenize{compute_smd_damage:related-commands}}
\sphinxAtStartPar
{\hyperref[\detokenize{compute_smd_plastic_strain::doc}]{\sphinxcrossref{\DUrole{doc}{smd/plastic\_strain}}}}, {\hyperref[\detokenize{compute_smd_tlsph_stress::doc}]{\sphinxcrossref{\DUrole{doc}{smd/tlsph\_stress}}}}


\subsection{Default}
\label{\detokenize{compute_smd_damage:default}}
\sphinxAtStartPar
none

\sphinxstepscope

\index{compute smd/hourglass/error@\spxentry{compute smd/hourglass/error}}\ignorespaces 

\section{compute smd/hourglass/error command}
\label{\detokenize{compute_smd_hourglass_error:compute-smd-hourglass-error-command}}\label{\detokenize{compute_smd_hourglass_error:index-0}}\label{\detokenize{compute_smd_hourglass_error::doc}}

\subsection{Syntax}
\label{\detokenize{compute_smd_hourglass_error:syntax}}
\begin{sphinxVerbatim}[commandchars=\\\{\}]
\PYG{k}{compute }\PYG{n+nv+nvIdentifier}{ID}\PYG{+w}{ }\PYG{n+nv+nvIdentifier}{group\PYGZhy{}ID}\PYG{+w}{ }\PYG{n}{smd}\PYG{o}{/}\PYG{n}{hourglass}\PYG{o}{/}\PYG{n}{error}
\end{sphinxVerbatim}
\begin{itemize}
\item {} 
\sphinxAtStartPar
ID, group\sphinxhyphen{}ID are documented in {\hyperref[\detokenize{compute::doc}]{\sphinxcrossref{\DUrole{doc}{compute}}}} command

\item {} 
\sphinxAtStartPar
smd/hourglass/error = style name of this compute command

\end{itemize}


\subsection{Examples}
\label{\detokenize{compute_smd_hourglass_error:examples}}
\begin{sphinxVerbatim}[commandchars=\\\{\}]
\PYG{k}{compute }\PYG{n+nv+nvIdentifier}{1}\PYG{+w}{ }\PYG{n+nv+nvIdentifier}{all}\PYG{+w}{ }\PYG{n}{smd}\PYG{o}{/}\PYG{n}{hourglass}\PYG{o}{/}\PYG{n}{error}
\end{sphinxVerbatim}


\subsection{Description}
\label{\detokenize{compute_smd_hourglass_error:description}}
\sphinxAtStartPar
Define a computation which outputs the error of the approximated
relative separation with respect to the actual relative separation of
the particles i and j. Ideally, if the deformation gradient is exact,
and there exists a unique mapping between all particles’ positions
within the neighborhood of the central node and the deformation
gradient, the approximated relative separation will coincide with the
actual relative separation of the particles i and j in the deformed
configuration.  This compute is only really useful for debugging the
hourglass control mechanism which is part of the Total\sphinxhyphen{}Lagrangian SPH
pair style.

\sphinxAtStartPar
See \sphinxhref{PDF/MACHDYN\_LAMMPS\_userguide.pdf}{this PDF guide} to use Smooth
Mach Dynamics in LAMMPS.

\sphinxAtStartPar
\sphinxstylestrong{Output Info:}

\sphinxAtStartPar
This compute calculates a per\sphinxhyphen{}particle vector, which can be accessed
by any command that uses per\sphinxhyphen{}particle values from a compute as input.
See the {\hyperref[\detokenize{Howto_output::doc}]{\sphinxcrossref{\DUrole{doc}{Howto output}}}} page for an overview of
LAMMPS output options.

\sphinxAtStartPar
The per\sphinxhyphen{}particle vector values will are dimensionless. See
{\hyperref[\detokenize{units::doc}]{\sphinxcrossref{\DUrole{doc}{units}}}}.


\subsection{Restrictions}
\label{\detokenize{compute_smd_hourglass_error:restrictions}}
\sphinxAtStartPar
This compute is part of the MACHDYN package.  It is only enabled if
LAMMPS was built with that package.  See the {\hyperref[\detokenize{Build_package::doc}]{\sphinxcrossref{\DUrole{doc}{Build package}}}} page for more info.

\sphinxAtStartPar
This quantity will be computed only for particles which interact with
tlsph pair style.


\subsection{Related commands}
\label{\detokenize{compute_smd_hourglass_error:related-commands}}
\sphinxAtStartPar
{\hyperref[\detokenize{compute_smd_tlsph_defgrad::doc}]{\sphinxcrossref{\DUrole{doc}{smd/tlsph\_defgrad}}}}


\subsection{Default}
\label{\detokenize{compute_smd_hourglass_error:default}}
\sphinxstepscope

\index{compute smd/internal/energy@\spxentry{compute smd/internal/energy}}\ignorespaces 

\section{compute smd/internal/energy command}
\label{\detokenize{compute_smd_internal_energy:compute-smd-internal-energy-command}}\label{\detokenize{compute_smd_internal_energy:index-0}}\label{\detokenize{compute_smd_internal_energy::doc}}

\subsection{Syntax}
\label{\detokenize{compute_smd_internal_energy:syntax}}
\begin{sphinxVerbatim}[commandchars=\\\{\}]
\PYG{k}{compute }\PYG{n+nv+nvIdentifier}{ID}\PYG{+w}{ }\PYG{n+nv+nvIdentifier}{group\PYGZhy{}ID}\PYG{+w}{ }\PYG{n}{smd}\PYG{o}{/}\PYG{n}{internal}\PYG{o}{/}\PYG{n}{energy}
\end{sphinxVerbatim}
\begin{itemize}
\item {} 
\sphinxAtStartPar
ID, group\sphinxhyphen{}ID are documented in {\hyperref[\detokenize{compute::doc}]{\sphinxcrossref{\DUrole{doc}{compute}}}} command

\item {} 
\sphinxAtStartPar
smd/smd/internal/energy = style name of this compute command

\end{itemize}


\subsection{Examples}
\label{\detokenize{compute_smd_internal_energy:examples}}
\begin{sphinxVerbatim}[commandchars=\\\{\}]
\PYG{k}{compute }\PYG{n+nv+nvIdentifier}{1}\PYG{+w}{ }\PYG{n+nv+nvIdentifier}{all}\PYG{+w}{ }\PYG{n}{smd}\PYG{o}{/}\PYG{n}{internal}\PYG{o}{/}\PYG{n}{energy}
\end{sphinxVerbatim}


\subsection{Description}
\label{\detokenize{compute_smd_internal_energy:description}}
\sphinxAtStartPar
Define a computation which outputs the per\sphinxhyphen{}particle enthalpy, i.e.,
the sum of potential energy and heat.

\sphinxAtStartPar
See \sphinxhref{PDF/MACHDYN\_LAMMPS\_userguide.pdf}{this PDF guide} to use Smooth
Mach Dynamics in LAMMPS.


\subsection{Output Info}
\label{\detokenize{compute_smd_internal_energy:output-info}}
\sphinxAtStartPar
This compute calculates a per\sphinxhyphen{}particle vector, which can be accessed
by any command that uses per\sphinxhyphen{}particle values from a compute as input.
See the {\hyperref[\detokenize{Howto_output::doc}]{\sphinxcrossref{\DUrole{doc}{Howto output}}}} page for an overview of
LAMMPS output options.

\sphinxAtStartPar
The per\sphinxhyphen{}particle vector values will be given in {\hyperref[\detokenize{units::doc}]{\sphinxcrossref{\DUrole{doc}{units}}}} of energy.


\subsection{Restrictions}
\label{\detokenize{compute_smd_internal_energy:restrictions}}
\sphinxAtStartPar
This compute is part of the MACHDYN package.  It is only enabled if
LAMMPS was built with that package.  See the {\hyperref[\detokenize{Build_package::doc}]{\sphinxcrossref{\DUrole{doc}{Build package}}}} page for more info. This compute can
only be used for particles which interact via the updated Lagrangian
or total Lagrangian SPH pair styles.


\subsection{Related commands}
\label{\detokenize{compute_smd_internal_energy:related-commands}}
\sphinxAtStartPar
none


\subsection{Default}
\label{\detokenize{compute_smd_internal_energy:default}}
\sphinxstepscope

\index{compute smd/plastic/strain@\spxentry{compute smd/plastic/strain}}\ignorespaces 

\section{compute smd/plastic/strain command}
\label{\detokenize{compute_smd_plastic_strain:compute-smd-plastic-strain-command}}\label{\detokenize{compute_smd_plastic_strain:index-0}}\label{\detokenize{compute_smd_plastic_strain::doc}}

\subsection{Syntax}
\label{\detokenize{compute_smd_plastic_strain:syntax}}
\begin{sphinxVerbatim}[commandchars=\\\{\}]
\PYG{k}{compute }\PYG{n+nv+nvIdentifier}{ID}\PYG{+w}{ }\PYG{n+nv+nvIdentifier}{group\PYGZhy{}ID}\PYG{+w}{ }\PYG{n}{smd}\PYG{o}{/}\PYG{n}{plastic}\PYG{o}{/}\PYG{n}{strain}
\end{sphinxVerbatim}
\begin{itemize}
\item {} 
\sphinxAtStartPar
ID, group\sphinxhyphen{}ID are documented in {\hyperref[\detokenize{compute::doc}]{\sphinxcrossref{\DUrole{doc}{compute}}}} command

\item {} 
\sphinxAtStartPar
smd/plastic/strain = style name of this compute command

\end{itemize}


\subsection{Examples}
\label{\detokenize{compute_smd_plastic_strain:examples}}
\begin{sphinxVerbatim}[commandchars=\\\{\}]
\PYG{k}{compute }\PYG{n+nv+nvIdentifier}{1}\PYG{+w}{ }\PYG{n+nv+nvIdentifier}{all}\PYG{+w}{ }\PYG{n}{smd}\PYG{o}{/}\PYG{n}{plastic}\PYG{o}{/}\PYG{n}{strain}
\end{sphinxVerbatim}


\subsection{Description}
\label{\detokenize{compute_smd_plastic_strain:description}}
\sphinxAtStartPar
Define a computation that outputs the equivalent plastic strain per
particle.  This command is only meaningful if a material model with
plasticity is defined.

\sphinxAtStartPar
See \sphinxhref{PDF/MACHDYN\_LAMMPS\_userguide.pdf}{this PDF guide} to use Smooth
Mach Dynamics in LAMMPS.

\sphinxAtStartPar
\sphinxstylestrong{Output Info:}

\sphinxAtStartPar
This compute calculates a per\sphinxhyphen{}particle vector, which can be accessed
by any command that uses per\sphinxhyphen{}particle values from a compute as input.
See the {\hyperref[\detokenize{Howto_output::doc}]{\sphinxcrossref{\DUrole{doc}{Howto output}}}} page for an overview of
LAMMPS output options.

\sphinxAtStartPar
The per\sphinxhyphen{}particle values will be given dimensionless. See {\hyperref[\detokenize{units::doc}]{\sphinxcrossref{\DUrole{doc}{units}}}}.


\subsection{Restrictions}
\label{\detokenize{compute_smd_plastic_strain:restrictions}}
\sphinxAtStartPar
This compute is part of the MACHDYN package.  It is only enabled if
LAMMPS was built with that package.  See the {\hyperref[\detokenize{Build_package::doc}]{\sphinxcrossref{\DUrole{doc}{Build package}}}} page for more info. This compute can
only be used for particles which interact via the updated Lagrangian
or total Lagrangian SPH pair styles.


\subsection{Related commands}
\label{\detokenize{compute_smd_plastic_strain:related-commands}}
\sphinxAtStartPar
{\hyperref[\detokenize{compute_smd_tlsph_strain::doc}]{\sphinxcrossref{\DUrole{doc}{smd/plastic/strain/rate}}}},
{\hyperref[\detokenize{compute_smd_tlsph_strain_rate::doc}]{\sphinxcrossref{\DUrole{doc}{smd/tlsph/strain/rate}}}},
{\hyperref[\detokenize{compute_smd_tlsph_strain::doc}]{\sphinxcrossref{\DUrole{doc}{smd/tlsph/strain}}}}


\subsection{Default}
\label{\detokenize{compute_smd_plastic_strain:default}}
\sphinxAtStartPar
none

\sphinxstepscope

\index{compute smd/plastic/strain/rate@\spxentry{compute smd/plastic/strain/rate}}\ignorespaces 

\section{compute smd/plastic/strain/rate command}
\label{\detokenize{compute_smd_plastic_strain_rate:compute-smd-plastic-strain-rate-command}}\label{\detokenize{compute_smd_plastic_strain_rate:index-0}}\label{\detokenize{compute_smd_plastic_strain_rate::doc}}

\subsection{Syntax}
\label{\detokenize{compute_smd_plastic_strain_rate:syntax}}
\begin{sphinxVerbatim}[commandchars=\\\{\}]
\PYG{k}{compute }\PYG{n+nv+nvIdentifier}{ID}\PYG{+w}{ }\PYG{n+nv+nvIdentifier}{group\PYGZhy{}ID}\PYG{+w}{ }\PYG{n}{smd}\PYG{o}{/}\PYG{n}{plastic}\PYG{o}{/}\PYG{n}{strain}\PYG{o}{/}\PYG{n}{rate}
\end{sphinxVerbatim}
\begin{itemize}
\item {} 
\sphinxAtStartPar
ID, group\sphinxhyphen{}ID are documented in {\hyperref[\detokenize{compute::doc}]{\sphinxcrossref{\DUrole{doc}{compute}}}} command

\item {} 
\sphinxAtStartPar
smd/plastic/strain/rate = style name of this compute command

\end{itemize}


\subsection{Examples}
\label{\detokenize{compute_smd_plastic_strain_rate:examples}}
\begin{sphinxVerbatim}[commandchars=\\\{\}]
\PYG{k}{compute }\PYG{n+nv+nvIdentifier}{1}\PYG{+w}{ }\PYG{n+nv+nvIdentifier}{all}\PYG{+w}{ }\PYG{n}{smd}\PYG{o}{/}\PYG{n}{plastic}\PYG{o}{/}\PYG{n}{strain}\PYG{o}{/}\PYG{n}{rate}
\end{sphinxVerbatim}


\subsection{Description}
\label{\detokenize{compute_smd_plastic_strain_rate:description}}
\sphinxAtStartPar
Define a computation that outputs the time rate of the equivalent
plastic strain.  This command is only meaningful if a material model
with plasticity is defined.

\sphinxAtStartPar
See \sphinxhref{PDF/MACHDYN\_LAMMPS\_userguide.pdf}{this PDF guide} to use Smooth
Mach Dynamics in LAMMPS.

\sphinxAtStartPar
\sphinxstylestrong{Output Info:}

\sphinxAtStartPar
This compute calculates a per\sphinxhyphen{}particle vector, which can be accessed
by any command that uses per\sphinxhyphen{}particle values from a compute as input.
See the {\hyperref[\detokenize{Howto_output::doc}]{\sphinxcrossref{\DUrole{doc}{Howto output}}}} page for an overview of
LAMMPS output options.

\sphinxAtStartPar
The per\sphinxhyphen{}particle values will be given in {\hyperref[\detokenize{units::doc}]{\sphinxcrossref{\DUrole{doc}{units}}}} of one over time.


\subsection{Restrictions}
\label{\detokenize{compute_smd_plastic_strain_rate:restrictions}}
\sphinxAtStartPar
This compute is part of the MACHDYN package.  It is only enabled if
LAMMPS was built with that package.  See the {\hyperref[\detokenize{Build_package::doc}]{\sphinxcrossref{\DUrole{doc}{Build package}}}} page for more info. This compute can
only be used for particles which interact via the updated Lagrangian
or total Lagrangian SPH pair styles.


\subsection{Related commands}
\label{\detokenize{compute_smd_plastic_strain_rate:related-commands}}
\sphinxAtStartPar
{\hyperref[\detokenize{compute_smd_plastic_strain::doc}]{\sphinxcrossref{\DUrole{doc}{smd/plastic/strain}}}},
{\hyperref[\detokenize{compute_smd_tlsph_strain_rate::doc}]{\sphinxcrossref{\DUrole{doc}{smd/tlsph/strain/rate}}}},
{\hyperref[\detokenize{compute_smd_tlsph_strain::doc}]{\sphinxcrossref{\DUrole{doc}{smd/tlsph/strain}}}}


\subsection{Default}
\label{\detokenize{compute_smd_plastic_strain_rate:default}}
\sphinxAtStartPar
none

\sphinxstepscope

\index{compute smd/rho@\spxentry{compute smd/rho}}\ignorespaces 

\section{compute smd/rho command}
\label{\detokenize{compute_smd_rho:compute-smd-rho-command}}\label{\detokenize{compute_smd_rho:index-0}}\label{\detokenize{compute_smd_rho::doc}}

\subsection{Syntax}
\label{\detokenize{compute_smd_rho:syntax}}
\begin{sphinxVerbatim}[commandchars=\\\{\}]
\PYG{k}{compute }\PYG{n+nv+nvIdentifier}{ID}\PYG{+w}{ }\PYG{n+nv+nvIdentifier}{group\PYGZhy{}ID}\PYG{+w}{ }\PYG{n}{smd}\PYG{o}{/}\PYG{n}{rho}
\end{sphinxVerbatim}
\begin{itemize}
\item {} 
\sphinxAtStartPar
ID, group\sphinxhyphen{}ID are documented in {\hyperref[\detokenize{compute::doc}]{\sphinxcrossref{\DUrole{doc}{compute}}}} command

\item {} 
\sphinxAtStartPar
smd/rho = style name of this compute command

\end{itemize}


\subsection{Examples}
\label{\detokenize{compute_smd_rho:examples}}
\begin{sphinxVerbatim}[commandchars=\\\{\}]
\PYG{k}{compute }\PYG{n+nv+nvIdentifier}{1}\PYG{+w}{ }\PYG{n+nv+nvIdentifier}{all}\PYG{+w}{ }\PYG{n}{smd}\PYG{o}{/}\PYG{n}{rho}
\end{sphinxVerbatim}


\subsection{Description}
\label{\detokenize{compute_smd_rho:description}}
\sphinxAtStartPar
Define a computation that calculates the per\sphinxhyphen{}particle mass density.
The mass density is the mass of a particle which is constant during
the course of a simulation, divided by its volume, which can change
due to mechanical deformation.

\sphinxAtStartPar
See \sphinxhref{PDF/MACHDYN\_LAMMPS\_userguide.pdf}{this PDF guide} to use Smooth
Mach Dynamics in LAMMPS.


\subsection{Output info}
\label{\detokenize{compute_smd_rho:output-info}}
\sphinxAtStartPar
This compute calculates a per\sphinxhyphen{}particle vector, which can be accessed
by any command that uses per\sphinxhyphen{}particle values from a compute as input.
See the {\hyperref[\detokenize{Howto_output::doc}]{\sphinxcrossref{\DUrole{doc}{Howto output}}}} page for an overview of
LAMMPS output options.

\sphinxAtStartPar
The per\sphinxhyphen{}particle values will be in {\hyperref[\detokenize{units::doc}]{\sphinxcrossref{\DUrole{doc}{units}}}} of mass over volume.


\subsection{Restrictions}
\label{\detokenize{compute_smd_rho:restrictions}}
\sphinxAtStartPar
This compute is part of the MACHDYN package. It is only enabled if
LAMMPS was built with that package. See the {\hyperref[\detokenize{Build_package::doc}]{\sphinxcrossref{\DUrole{doc}{Build package}}}} page for more info.


\subsection{Related commands}
\label{\detokenize{compute_smd_rho:related-commands}}
\sphinxAtStartPar
{\hyperref[\detokenize{compute_smd_vol::doc}]{\sphinxcrossref{\DUrole{doc}{compute smd/vol}}}}


\subsection{Default}
\label{\detokenize{compute_smd_rho:default}}
\sphinxAtStartPar
none

\sphinxstepscope

\index{compute smd/tlsph/defgrad@\spxentry{compute smd/tlsph/defgrad}}\ignorespaces 

\section{compute smd/tlsph/defgrad command}
\label{\detokenize{compute_smd_tlsph_defgrad:compute-smd-tlsph-defgrad-command}}\label{\detokenize{compute_smd_tlsph_defgrad:index-0}}\label{\detokenize{compute_smd_tlsph_defgrad::doc}}

\subsection{Syntax}
\label{\detokenize{compute_smd_tlsph_defgrad:syntax}}
\begin{sphinxVerbatim}[commandchars=\\\{\}]
\PYG{k}{compute }\PYG{n+nv+nvIdentifier}{ID}\PYG{+w}{ }\PYG{n+nv+nvIdentifier}{group\PYGZhy{}ID}\PYG{+w}{ }\PYG{n}{smd}\PYG{o}{/}\PYG{n}{tlsph}\PYG{o}{/}\PYG{n}{defgrad}
\end{sphinxVerbatim}
\begin{itemize}
\item {} 
\sphinxAtStartPar
ID, group\sphinxhyphen{}ID are documented in {\hyperref[\detokenize{compute::doc}]{\sphinxcrossref{\DUrole{doc}{compute}}}} command

\item {} 
\sphinxAtStartPar
smd/tlsph/defgrad = style name of this compute command

\end{itemize}


\subsection{Examples}
\label{\detokenize{compute_smd_tlsph_defgrad:examples}}
\begin{sphinxVerbatim}[commandchars=\\\{\}]
\PYG{k}{compute }\PYG{n+nv+nvIdentifier}{1}\PYG{+w}{ }\PYG{n+nv+nvIdentifier}{all}\PYG{+w}{ }\PYG{n}{smd}\PYG{o}{/}\PYG{n}{tlsph}\PYG{o}{/}\PYG{n}{defgrad}
\end{sphinxVerbatim}


\subsection{Description}
\label{\detokenize{compute_smd_tlsph_defgrad:description}}
\sphinxAtStartPar
Define a computation that calculates the deformation gradient.  It is
only meaningful for particles which interact according to the
Total\sphinxhyphen{}Lagrangian SPH pair style.

\sphinxAtStartPar
See \sphinxhref{PDF/MACHDYN\_LAMMPS\_userguide.pdf}{this PDF guide} to use Smooth
Mach Dynamics in LAMMPS.


\subsection{Output info}
\label{\detokenize{compute_smd_tlsph_defgrad:output-info}}
\sphinxAtStartPar
This compute outputs a per\sphinxhyphen{}particle vector of vectors (tensors),
which can be accessed by any command that uses per\sphinxhyphen{}particle values
from a compute as input. See the {\hyperref[\detokenize{Howto_output::doc}]{\sphinxcrossref{\DUrole{doc}{Howto output}}}} doc
page for an overview of LAMMPS output options.

\sphinxAtStartPar
The per\sphinxhyphen{}particle vector values will be given dimensionless. See
{\hyperref[\detokenize{units::doc}]{\sphinxcrossref{\DUrole{doc}{units}}}}.  The per\sphinxhyphen{}particle vector has 10 entries. The first
nine entries correspond to the xx, xy, xz, yx, yy, yz, zx, zy, zz
components of the asymmetric deformation gradient tensor. The tenth
entry is the determinant of the deformation gradient.


\subsection{Restrictions}
\label{\detokenize{compute_smd_tlsph_defgrad:restrictions}}
\sphinxAtStartPar
This compute is part of the MACHDYN package.  It is only enabled if
LAMMPS was built with that package. See the {\hyperref[\detokenize{Build_package::doc}]{\sphinxcrossref{\DUrole{doc}{Build package}}}} page for more info. TThis compute can
only be used for particles which interact via the total Lagrangian SPH
pair style.


\subsection{Related commands}
\label{\detokenize{compute_smd_tlsph_defgrad:related-commands}}
\sphinxAtStartPar
{\hyperref[\detokenize{compute_smd_hourglass_error::doc}]{\sphinxcrossref{\DUrole{doc}{smd/hourglass/error}}}}


\subsection{Default}
\label{\detokenize{compute_smd_tlsph_defgrad:default}}
\sphinxAtStartPar
none

\sphinxstepscope

\index{compute smd/tlsph/dt@\spxentry{compute smd/tlsph/dt}}\ignorespaces 

\section{compute smd/tlsph/dt command}
\label{\detokenize{compute_smd_tlsph_dt:compute-smd-tlsph-dt-command}}\label{\detokenize{compute_smd_tlsph_dt:index-0}}\label{\detokenize{compute_smd_tlsph_dt::doc}}

\subsection{Syntax}
\label{\detokenize{compute_smd_tlsph_dt:syntax}}
\begin{sphinxVerbatim}[commandchars=\\\{\}]
\PYG{k}{compute }\PYG{n+nv+nvIdentifier}{ID}\PYG{+w}{ }\PYG{n+nv+nvIdentifier}{group\PYGZhy{}ID}\PYG{+w}{ }\PYG{n}{smd}\PYG{o}{/}\PYG{n}{tlsph}\PYG{o}{/}\PYG{n}{dt}
\end{sphinxVerbatim}
\begin{itemize}
\item {} 
\sphinxAtStartPar
ID, group\sphinxhyphen{}ID are documented in {\hyperref[\detokenize{compute::doc}]{\sphinxcrossref{\DUrole{doc}{compute}}}} command

\item {} 
\sphinxAtStartPar
smd/tlsph/dt = style name of this compute command

\end{itemize}


\subsection{Examples}
\label{\detokenize{compute_smd_tlsph_dt:examples}}
\begin{sphinxVerbatim}[commandchars=\\\{\}]
\PYG{k}{compute }\PYG{n+nv+nvIdentifier}{1}\PYG{+w}{ }\PYG{n+nv+nvIdentifier}{all}\PYG{+w}{ }\PYG{n}{smd}\PYG{o}{/}\PYG{n}{tlsph}\PYG{o}{/}\PYG{n}{dt}
\end{sphinxVerbatim}


\subsection{Description}
\label{\detokenize{compute_smd_tlsph_dt:description}}
\sphinxAtStartPar
Define a computation that outputs the CFL\sphinxhyphen{}stable time increment per
particle.  This time increment is essentially given by the speed of
sound, divided by the SPH smoothing length.  Because both the speed of
sound and the smoothing length typically change during the course of a
simulation, the stable time increment needs to be re\sphinxhyphen{}computed every
time step.  This calculation is performed automatically in the
relevant SPH pair styles and this compute only serves to make the
stable time increment accessible for output purposes.

\sphinxAtStartPar
See \sphinxhref{PDF/MACHDYN\_LAMMPS\_userguide.pdf}{this PDF guide} to using Smooth
Mach Dynamics in LAMMPS.


\subsection{Output info}
\label{\detokenize{compute_smd_tlsph_dt:output-info}}
\sphinxAtStartPar
This compute calculates a per\sphinxhyphen{}particle vector, which can be accessed
by any command that uses per\sphinxhyphen{}particle values from a compute as input.
See the {\hyperref[\detokenize{Howto_output::doc}]{\sphinxcrossref{\DUrole{doc}{Howto output}}}} page for an overview of
LAMMPS output options.

\sphinxAtStartPar
The per\sphinxhyphen{}particle values will be given in {\hyperref[\detokenize{units::doc}]{\sphinxcrossref{\DUrole{doc}{units}}}} of time.


\subsection{Restrictions}
\label{\detokenize{compute_smd_tlsph_dt:restrictions}}
\sphinxAtStartPar
This compute is part of the MACHDYN package.  It is only enabled if
LAMMPS was built with that package.  See the {\hyperref[\detokenize{Build_package::doc}]{\sphinxcrossref{\DUrole{doc}{Build package}}}} page for more info.

\sphinxAtStartPar
This compute can only be used for particles interacting with the
Total\sphinxhyphen{}Lagrangian SPH pair style.


\subsection{Related commands}
\label{\detokenize{compute_smd_tlsph_dt:related-commands}}
\sphinxAtStartPar
{\hyperref[\detokenize{fix_smd_adjust_dt::doc}]{\sphinxcrossref{\DUrole{doc}{smd/adjust/dt}}}}


\subsection{Default}
\label{\detokenize{compute_smd_tlsph_dt:default}}
\sphinxAtStartPar
none

\sphinxstepscope

\index{compute smd/tlsph/num/neighs@\spxentry{compute smd/tlsph/num/neighs}}\ignorespaces 

\section{compute smd/tlsph/num/neighs command}
\label{\detokenize{compute_smd_tlsph_num_neighs:compute-smd-tlsph-num-neighs-command}}\label{\detokenize{compute_smd_tlsph_num_neighs:index-0}}\label{\detokenize{compute_smd_tlsph_num_neighs::doc}}

\subsection{Syntax}
\label{\detokenize{compute_smd_tlsph_num_neighs:syntax}}
\begin{sphinxVerbatim}[commandchars=\\\{\}]
\PYG{k}{compute }\PYG{n+nv+nvIdentifier}{ID}\PYG{+w}{ }\PYG{n+nv+nvIdentifier}{group\PYGZhy{}ID}\PYG{+w}{ }\PYG{n}{smd}\PYG{o}{/}\PYG{n}{tlsph}\PYG{o}{/}\PYG{n}{num}\PYG{o}{/}\PYG{n}{neighs}
\end{sphinxVerbatim}
\begin{itemize}
\item {} 
\sphinxAtStartPar
ID, group\sphinxhyphen{}ID are documented in {\hyperref[\detokenize{compute::doc}]{\sphinxcrossref{\DUrole{doc}{compute}}}} command

\item {} 
\sphinxAtStartPar
smd/tlsph/num/neighs = style name of this compute command

\end{itemize}


\subsection{Examples}
\label{\detokenize{compute_smd_tlsph_num_neighs:examples}}
\begin{sphinxVerbatim}[commandchars=\\\{\}]
\PYG{k}{compute }\PYG{n+nv+nvIdentifier}{1}\PYG{+w}{ }\PYG{n+nv+nvIdentifier}{all}\PYG{+w}{ }\PYG{n}{smd}\PYG{o}{/}\PYG{n}{tlsph}\PYG{o}{/}\PYG{n}{num}\PYG{o}{/}\PYG{n}{neighs}
\end{sphinxVerbatim}


\subsection{Description}
\label{\detokenize{compute_smd_tlsph_num_neighs:description}}
\sphinxAtStartPar
Define a computation that calculates the number of particles inside of
the smoothing kernel radius for particles interacting via the
Total\sphinxhyphen{}Lagrangian SPH pair style.

\sphinxAtStartPar
See \sphinxhref{PDF/MACHDYN\_LAMMPS\_userguide.pdf}{this PDF guide} to using Smooth
Mach Dynamics in LAMMPS.


\subsection{Output info}
\label{\detokenize{compute_smd_tlsph_num_neighs:output-info}}
\sphinxAtStartPar
This compute calculates a per\sphinxhyphen{}particle vector, which can be accessed
by any command that uses per\sphinxhyphen{}particle values from a compute as input.
See the {\hyperref[\detokenize{Howto_output::doc}]{\sphinxcrossref{\DUrole{doc}{Howto output}}}} page for an overview of
LAMMPS output options.

\sphinxAtStartPar
The per\sphinxhyphen{}particle values are dimensionless. See {\hyperref[\detokenize{units::doc}]{\sphinxcrossref{\DUrole{doc}{units}}}}.


\subsection{Restrictions}
\label{\detokenize{compute_smd_tlsph_num_neighs:restrictions}}
\sphinxAtStartPar
This compute is part of the MACHDYN package.  It is only enabled if
LAMMPS was built with that package.  See the {\hyperref[\detokenize{Build_package::doc}]{\sphinxcrossref{\DUrole{doc}{Build package}}}} page for more info.

\sphinxAtStartPar
This quantity will be computed only for particles which interact with
the Total\sphinxhyphen{}Lagrangian pair style.


\subsection{Related commands}
\label{\detokenize{compute_smd_tlsph_num_neighs:related-commands}}
\sphinxAtStartPar
{\hyperref[\detokenize{compute_smd_ulsph_num_neighs::doc}]{\sphinxcrossref{\DUrole{doc}{smd/ulsph/num/neighs}}}}


\subsection{Default}
\label{\detokenize{compute_smd_tlsph_num_neighs:default}}
\sphinxAtStartPar
none

\sphinxstepscope

\index{compute smd/tlsph/shape@\spxentry{compute smd/tlsph/shape}}\ignorespaces 

\section{compute smd/tlsph/shape command}
\label{\detokenize{compute_smd_tlsph_shape:compute-smd-tlsph-shape-command}}\label{\detokenize{compute_smd_tlsph_shape:index-0}}\label{\detokenize{compute_smd_tlsph_shape::doc}}

\subsection{Syntax}
\label{\detokenize{compute_smd_tlsph_shape:syntax}}
\begin{sphinxVerbatim}[commandchars=\\\{\}]
\PYG{k}{compute }\PYG{n+nv+nvIdentifier}{ID}\PYG{+w}{ }\PYG{n+nv+nvIdentifier}{group\PYGZhy{}ID}\PYG{+w}{ }\PYG{n}{smd}\PYG{o}{/}\PYG{n}{tlsph}\PYG{o}{/}\PYG{n}{shape}
\end{sphinxVerbatim}
\begin{itemize}
\item {} 
\sphinxAtStartPar
ID, group\sphinxhyphen{}ID are documented in {\hyperref[\detokenize{compute::doc}]{\sphinxcrossref{\DUrole{doc}{compute}}}} command

\item {} 
\sphinxAtStartPar
smd/tlsph/shape = style name of this compute command

\end{itemize}


\subsection{Examples}
\label{\detokenize{compute_smd_tlsph_shape:examples}}
\begin{sphinxVerbatim}[commandchars=\\\{\}]
\PYG{k}{compute }\PYG{n+nv+nvIdentifier}{1}\PYG{+w}{ }\PYG{n+nv+nvIdentifier}{all}\PYG{+w}{ }\PYG{n}{smd}\PYG{o}{/}\PYG{n}{tlsph}\PYG{o}{/}\PYG{n}{shape}
\end{sphinxVerbatim}


\subsection{Description}
\label{\detokenize{compute_smd_tlsph_shape:description}}
\sphinxAtStartPar
Define a computation that outputs the current shape of the volume
associated with a particle as a rotated ellipsoid.  It is only
meaningful for particles which interact according to the
Total\sphinxhyphen{}Lagrangian SPH pair style.

\sphinxAtStartPar
See \sphinxhref{PDF/MACHDYN\_LAMMPS\_userguide.pdf}{this PDF guide} to use Smooth
Mach Dynamics in LAMMPS.


\subsection{Output info}
\label{\detokenize{compute_smd_tlsph_shape:output-info}}
\sphinxAtStartPar
This compute calculates a per\sphinxhyphen{}particle vector of vectors, which can be
accessed by any command that uses per\sphinxhyphen{}particle values from a compute
as input. See the {\hyperref[\detokenize{Howto_output::doc}]{\sphinxcrossref{\DUrole{doc}{Howto output}}}} page for an
overview of LAMMPS output options.

\sphinxAtStartPar
The per\sphinxhyphen{}particle vector has 7 entries. The first three entries
correspond to the lengths of the ellipsoid’s axes and have units of
length.  These axis values are computed as the contact radius times the
xx, yy, or zz components of the Green\sphinxhyphen{}Lagrange strain tensor
associated with the particle.  The next 4 values are quaternions
(order: q, x, y, z) which describe the spatial rotation of the
particle relative to its initial state.


\subsection{Restrictions}
\label{\detokenize{compute_smd_tlsph_shape:restrictions}}
\sphinxAtStartPar
This compute is part of the MACHDYN package.  It is only enabled if
LAMMPS was built with that package. See the {\hyperref[\detokenize{Build_package::doc}]{\sphinxcrossref{\DUrole{doc}{Build package}}}} page for more info.

\sphinxAtStartPar
This quantity will be computed only for particles which interact with
the Total\sphinxhyphen{}Lagrangian SPH pair style.


\subsection{Related commands}
\label{\detokenize{compute_smd_tlsph_shape:related-commands}}
\sphinxAtStartPar
{\hyperref[\detokenize{compute_smd_contact_radius::doc}]{\sphinxcrossref{\DUrole{doc}{smd/contact/radius}}}}


\subsection{Default}
\label{\detokenize{compute_smd_tlsph_shape:default}}
\sphinxAtStartPar
none

\sphinxstepscope

\index{compute smd/tlsph/strain@\spxentry{compute smd/tlsph/strain}}\ignorespaces 

\section{compute smd/tlsph/strain command}
\label{\detokenize{compute_smd_tlsph_strain:compute-smd-tlsph-strain-command}}\label{\detokenize{compute_smd_tlsph_strain:index-0}}\label{\detokenize{compute_smd_tlsph_strain::doc}}

\subsection{Syntax}
\label{\detokenize{compute_smd_tlsph_strain:syntax}}
\begin{sphinxVerbatim}[commandchars=\\\{\}]
\PYG{k}{compute }\PYG{n+nv+nvIdentifier}{ID}\PYG{+w}{ }\PYG{n+nv+nvIdentifier}{group\PYGZhy{}ID}\PYG{+w}{ }\PYG{n}{smd}\PYG{o}{/}\PYG{n}{tlsph}\PYG{o}{/}\PYG{n}{strain}
\end{sphinxVerbatim}
\begin{itemize}
\item {} 
\sphinxAtStartPar
ID, group\sphinxhyphen{}ID are documented in {\hyperref[\detokenize{compute::doc}]{\sphinxcrossref{\DUrole{doc}{compute}}}} command

\item {} 
\sphinxAtStartPar
smd/tlsph/strain = style name of this compute command

\end{itemize}


\subsection{Examples}
\label{\detokenize{compute_smd_tlsph_strain:examples}}
\begin{sphinxVerbatim}[commandchars=\\\{\}]
\PYG{k}{compute }\PYG{n+nv+nvIdentifier}{1}\PYG{+w}{ }\PYG{n+nv+nvIdentifier}{all}\PYG{+w}{ }\PYG{n}{smd}\PYG{o}{/}\PYG{n}{tlsph}\PYG{o}{/}\PYG{n}{strain}
\end{sphinxVerbatim}


\subsection{Description}
\label{\detokenize{compute_smd_tlsph_strain:description}}
\sphinxAtStartPar
Define a computation that calculates the Green\sphinxhyphen{}Lagrange strain tensor
for particles interacting via the Total\sphinxhyphen{}Lagrangian SPH pair style.

\sphinxAtStartPar
See \sphinxhref{PDF/MACHDYN\_LAMMPS\_userguide.pdf}{this PDF guide} to using Smooth
Mach Dynamics in LAMMPS.


\subsection{Output info}
\label{\detokenize{compute_smd_tlsph_strain:output-info}}
\sphinxAtStartPar
This compute calculates a per\sphinxhyphen{}particle vector of vectors (tensors),
which can be accessed by any command that uses per\sphinxhyphen{}particle values
from a compute as input.  See the {\hyperref[\detokenize{Howto_output::doc}]{\sphinxcrossref{\DUrole{doc}{Howto output}}}} doc
page for an overview of LAMMPS output options.

\sphinxAtStartPar
The per\sphinxhyphen{}particle tensor values will be given dimensionless. See
{\hyperref[\detokenize{units::doc}]{\sphinxcrossref{\DUrole{doc}{units}}}}.

\sphinxAtStartPar
The per\sphinxhyphen{}particle vector has 6 entries, corresponding to the xx, yy,
zz, xy, xz, yz components of the symmetric strain tensor.


\subsection{Restrictions}
\label{\detokenize{compute_smd_tlsph_strain:restrictions}}
\sphinxAtStartPar
This compute is part of the MACHDYN package.  It is only enabled if
LAMMPS was built with that package.  See the {\hyperref[\detokenize{Build_package::doc}]{\sphinxcrossref{\DUrole{doc}{Build package}}}} page for more info.

\sphinxAtStartPar
This quantity will be computed only for particles which interact with
the Total\sphinxhyphen{}Lagrangian SPH pair style.


\subsection{Related commands}
\label{\detokenize{compute_smd_tlsph_strain:related-commands}}
\sphinxAtStartPar
{\hyperref[\detokenize{compute_smd_tlsph_strain_rate::doc}]{\sphinxcrossref{\DUrole{doc}{smd/tlsph/strain/rate}}}},
{\hyperref[\detokenize{compute_smd_tlsph_stress::doc}]{\sphinxcrossref{\DUrole{doc}{smd/tlsph/stress}}}}


\subsection{Default}
\label{\detokenize{compute_smd_tlsph_strain:default}}
\sphinxAtStartPar
none

\sphinxstepscope

\index{compute smd/tlsph/strain/rate@\spxentry{compute smd/tlsph/strain/rate}}\ignorespaces 

\section{compute smd/tlsph/strain/rate command}
\label{\detokenize{compute_smd_tlsph_strain_rate:compute-smd-tlsph-strain-rate-command}}\label{\detokenize{compute_smd_tlsph_strain_rate:index-0}}\label{\detokenize{compute_smd_tlsph_strain_rate::doc}}

\subsection{Syntax}
\label{\detokenize{compute_smd_tlsph_strain_rate:syntax}}
\begin{sphinxVerbatim}[commandchars=\\\{\}]
\PYG{k}{compute }\PYG{n+nv+nvIdentifier}{ID}\PYG{+w}{ }\PYG{n+nv+nvIdentifier}{group\PYGZhy{}ID}\PYG{+w}{ }\PYG{n}{smd}\PYG{o}{/}\PYG{n}{tlsph}\PYG{o}{/}\PYG{n}{strain}\PYG{o}{/}\PYG{n}{rate}
\end{sphinxVerbatim}
\begin{itemize}
\item {} 
\sphinxAtStartPar
ID, group\sphinxhyphen{}ID are documented in {\hyperref[\detokenize{compute::doc}]{\sphinxcrossref{\DUrole{doc}{compute}}}} command

\item {} 
\sphinxAtStartPar
smd/tlsph/strain/rate = style name of this compute command

\end{itemize}


\subsection{Examples}
\label{\detokenize{compute_smd_tlsph_strain_rate:examples}}
\begin{sphinxVerbatim}[commandchars=\\\{\}]
\PYG{k}{compute }\PYG{n+nv+nvIdentifier}{1}\PYG{+w}{ }\PYG{n+nv+nvIdentifier}{all}\PYG{+w}{ }\PYG{n}{smd}\PYG{o}{/}\PYG{n}{tlsph}\PYG{o}{/}\PYG{n}{strain}\PYG{o}{/}\PYG{n}{rate}
\end{sphinxVerbatim}


\subsection{Description}
\label{\detokenize{compute_smd_tlsph_strain_rate:description}}
\sphinxAtStartPar
Define a computation that calculates the rate of the strain tensor for
particles interacting via the Total\sphinxhyphen{}Lagrangian SPH pair style.

\sphinxAtStartPar
See \sphinxhref{PDF/MACHDYN\_LAMMPS\_userguide.pdf}{this PDF guide} to using Smooth
Mach Dynamics in LAMMPS.


\subsection{Output info}
\label{\detokenize{compute_smd_tlsph_strain_rate:output-info}}
\sphinxAtStartPar
This compute calculates a per\sphinxhyphen{}particle vector of vectors (tensors),
which can be accessed by any command that uses per\sphinxhyphen{}particle values
from a compute as input. See the {\hyperref[\detokenize{Howto_output::doc}]{\sphinxcrossref{\DUrole{doc}{Howto output}}}} doc
page for an overview of LAMMPS output options.

\sphinxAtStartPar
The values will be given in {\hyperref[\detokenize{units::doc}]{\sphinxcrossref{\DUrole{doc}{units}}}} of one over time.

\sphinxAtStartPar
The per\sphinxhyphen{}particle vector has 6 entries, corresponding to the xx, yy,
zz, xy, xz, yz components of the symmetric strain rate tensor.


\subsection{Restrictions}
\label{\detokenize{compute_smd_tlsph_strain_rate:restrictions}}
\sphinxAtStartPar
This compute is part of the MACHDYN package.  It is only enabled if
LAMMPS was built with that package.  See the {\hyperref[\detokenize{Build_package::doc}]{\sphinxcrossref{\DUrole{doc}{Build package}}}} page for more info.

\sphinxAtStartPar
This quantity will be computed only for particles which interact with
Total\sphinxhyphen{}Lagrangian SPH pair style.


\subsection{Related commands}
\label{\detokenize{compute_smd_tlsph_strain_rate:related-commands}}
\sphinxAtStartPar
{\hyperref[\detokenize{compute_smd_tlsph_strain::doc}]{\sphinxcrossref{\DUrole{doc}{compute smd/tlsph/strain}}}}, {\hyperref[\detokenize{compute_smd_tlsph_stress::doc}]{\sphinxcrossref{\DUrole{doc}{compute smd/tlsph/stress}}}}


\subsection{Default}
\label{\detokenize{compute_smd_tlsph_strain_rate:default}}
\sphinxAtStartPar
none

\sphinxstepscope

\index{compute smd/tlsph/stress@\spxentry{compute smd/tlsph/stress}}\ignorespaces 

\section{compute smd/tlsph/stress command}
\label{\detokenize{compute_smd_tlsph_stress:compute-smd-tlsph-stress-command}}\label{\detokenize{compute_smd_tlsph_stress:index-0}}\label{\detokenize{compute_smd_tlsph_stress::doc}}

\subsection{Syntax}
\label{\detokenize{compute_smd_tlsph_stress:syntax}}
\begin{sphinxVerbatim}[commandchars=\\\{\}]
\PYG{k}{compute }\PYG{n+nv+nvIdentifier}{ID}\PYG{+w}{ }\PYG{n+nv+nvIdentifier}{group\PYGZhy{}ID}\PYG{+w}{ }\PYG{n}{smd}\PYG{o}{/}\PYG{n}{tlsph}\PYG{o}{/}\PYG{n}{stress}
\end{sphinxVerbatim}
\begin{itemize}
\item {} 
\sphinxAtStartPar
ID, group\sphinxhyphen{}ID are documented in {\hyperref[\detokenize{compute::doc}]{\sphinxcrossref{\DUrole{doc}{compute}}}} command

\item {} 
\sphinxAtStartPar
smd/tlsph/stress = style name of this compute command

\end{itemize}


\subsection{Examples}
\label{\detokenize{compute_smd_tlsph_stress:examples}}
\begin{sphinxVerbatim}[commandchars=\\\{\}]
\PYG{k}{compute }\PYG{n+nv+nvIdentifier}{1}\PYG{+w}{ }\PYG{n+nv+nvIdentifier}{all}\PYG{+w}{ }\PYG{n}{smd}\PYG{o}{/}\PYG{n}{tlsph}\PYG{o}{/}\PYG{n}{stress}
\end{sphinxVerbatim}


\subsection{Description}
\label{\detokenize{compute_smd_tlsph_stress:description}}
\sphinxAtStartPar
Define a computation that outputs the Cauchy stress tensor for
particles interacting via the Total\sphinxhyphen{}Lagrangian SPH pair style.

\sphinxAtStartPar
See \sphinxhref{PDF/MACHDYN\_LAMMPS\_userguide.pdf}{this PDF guide} to using Smooth
Mach Dynamics in LAMMPS.


\subsection{Output info}
\label{\detokenize{compute_smd_tlsph_stress:output-info}}
\sphinxAtStartPar
This compute calculates a per\sphinxhyphen{}particle vector of vectors (tensors),
which can be accessed by any command that uses per\sphinxhyphen{}particle values
from a compute as input. See the {\hyperref[\detokenize{Howto_output::doc}]{\sphinxcrossref{\DUrole{doc}{Howto output}}}} doc
page for an overview of LAMMPS output options.

\sphinxAtStartPar
The values will be given in {\hyperref[\detokenize{units::doc}]{\sphinxcrossref{\DUrole{doc}{units}}}} of pressure.

\sphinxAtStartPar
The per\sphinxhyphen{}particle vector has 7 entries. The first six entries
correspond to the xx, yy, zz, xy, xz and yz components of the
symmetric Cauchy stress tensor. The seventh entry is the second
invariant of the stress tensor, i.e., the von Mises equivalent stress.


\subsection{Restrictions}
\label{\detokenize{compute_smd_tlsph_stress:restrictions}}
\sphinxAtStartPar
This compute is part of the MACHDYN package.  It is only enabled if
LAMMPS was built with that package.  See the {\hyperref[\detokenize{Build_package::doc}]{\sphinxcrossref{\DUrole{doc}{Build package}}}} page for more info.

\sphinxAtStartPar
This quantity will be computed only for particles which interact with
the Total\sphinxhyphen{}Lagrangian SPH pair style.


\subsection{Related commands}
\label{\detokenize{compute_smd_tlsph_stress:related-commands}}
\sphinxAtStartPar
{\hyperref[\detokenize{compute_smd_tlsph_strain::doc}]{\sphinxcrossref{\DUrole{doc}{compute smd/tlsph/strain}}}}, {\hyperref[\detokenize{compute_smd_tlsph_strain_rate::doc}]{\sphinxcrossref{\DUrole{doc}{cmopute smd/tlsph/strain/rate}}}}


\subsection{Default}
\label{\detokenize{compute_smd_tlsph_stress:default}}
\sphinxAtStartPar
none

\sphinxstepscope

\index{compute smd/triangle/vertices@\spxentry{compute smd/triangle/vertices}}\ignorespaces 

\section{compute smd/triangle/vertices command}
\label{\detokenize{compute_smd_triangle_vertices:compute-smd-triangle-vertices-command}}\label{\detokenize{compute_smd_triangle_vertices:index-0}}\label{\detokenize{compute_smd_triangle_vertices::doc}}

\subsection{Syntax}
\label{\detokenize{compute_smd_triangle_vertices:syntax}}
\begin{sphinxVerbatim}[commandchars=\\\{\}]
\PYG{k}{compute }\PYG{n+nv+nvIdentifier}{ID}\PYG{+w}{ }\PYG{n+nv+nvIdentifier}{group\PYGZhy{}ID}\PYG{+w}{ }\PYG{n}{smd}\PYG{o}{/}\PYG{n}{triangle}\PYG{o}{/}\PYG{n}{vertices}
\end{sphinxVerbatim}
\begin{itemize}
\item {} 
\sphinxAtStartPar
ID, group\sphinxhyphen{}ID are documented in {\hyperref[\detokenize{compute::doc}]{\sphinxcrossref{\DUrole{doc}{compute}}}} command

\item {} 
\sphinxAtStartPar
smd/triangle/vertices = style name of this compute command

\end{itemize}


\subsection{Examples}
\label{\detokenize{compute_smd_triangle_vertices:examples}}
\begin{sphinxVerbatim}[commandchars=\\\{\}]
\PYG{k}{compute }\PYG{n+nv+nvIdentifier}{1}\PYG{+w}{ }\PYG{n+nv+nvIdentifier}{all}\PYG{+w}{ }\PYG{n}{smd}\PYG{o}{/}\PYG{n}{triangle}\PYG{o}{/}\PYG{n}{vertices}
\end{sphinxVerbatim}


\subsection{Description}
\label{\detokenize{compute_smd_triangle_vertices:description}}
\sphinxAtStartPar
Define a computation that returns the coordinates of the vertices
corresponding to the triangle\sphinxhyphen{}elements of a mesh created by the {\hyperref[\detokenize{fix_smd_wall_surface::doc}]{\sphinxcrossref{\DUrole{doc}{fix smd/wall\_surface}}}}.

\sphinxAtStartPar
See \sphinxhref{PDF/MACHDYN\_LAMMPS\_userguide.pdf}{this PDF guide} to using Smooth
Mach Dynamics in LAMMPS.


\subsection{Output info}
\label{\detokenize{compute_smd_triangle_vertices:output-info}}
\sphinxAtStartPar
This compute returns a per\sphinxhyphen{}particle vector of vectors, which can be
accessed by any command that uses per\sphinxhyphen{}particle values from a compute
as input. See the {\hyperref[\detokenize{Howto_output::doc}]{\sphinxcrossref{\DUrole{doc}{Howto output}}}} page for an
overview of LAMMPS output options.

\sphinxAtStartPar
The per\sphinxhyphen{}particle vector has nine entries, (x1/y1/z1), (x2/y2/z2), and
(x3/y3/z3) corresponding to the first, second, and third vertex of
each triangle.

\sphinxAtStartPar
It is only meaningful to use this compute for a group of particles
which is created via the {\hyperref[\detokenize{fix_smd_wall_surface::doc}]{\sphinxcrossref{\DUrole{doc}{fix smd/wall\_surface}}}} command.

\sphinxAtStartPar
The output of this compute can be used with the dump2vtk\_tris tool to
generate a VTK representation of the smd/wall\_surface mesh for
visualization purposes.

\sphinxAtStartPar
The values will be given in {\hyperref[\detokenize{units::doc}]{\sphinxcrossref{\DUrole{doc}{units}}}} of distance.


\subsection{Restrictions}
\label{\detokenize{compute_smd_triangle_vertices:restrictions}}
\sphinxAtStartPar
This compute is part of the MACHDYN package.  It is only enabled if
LAMMPS was built with that package.  See the {\hyperref[\detokenize{Build_package::doc}]{\sphinxcrossref{\DUrole{doc}{Build package}}}} page for more info.


\subsection{Related commands}
\label{\detokenize{compute_smd_triangle_vertices:related-commands}}
\sphinxAtStartPar
{\hyperref[\detokenize{fix_smd_move_triangulated_surface::doc}]{\sphinxcrossref{\DUrole{doc}{fix smd/move/tri/surf}}}},
{\hyperref[\detokenize{fix_smd_wall_surface::doc}]{\sphinxcrossref{\DUrole{doc}{fix smd/wall\_surface}}}}


\subsection{Default}
\label{\detokenize{compute_smd_triangle_vertices:default}}
\sphinxAtStartPar
none

\sphinxstepscope

\index{compute smd/ulsph/effm@\spxentry{compute smd/ulsph/effm}}\ignorespaces 

\section{compute smd/ulsph/effm command}
\label{\detokenize{compute_smd_ulsph_effm:compute-smd-ulsph-effm-command}}\label{\detokenize{compute_smd_ulsph_effm:index-0}}\label{\detokenize{compute_smd_ulsph_effm::doc}}

\subsection{Syntax}
\label{\detokenize{compute_smd_ulsph_effm:syntax}}
\begin{sphinxVerbatim}[commandchars=\\\{\}]
\PYG{k}{compute }\PYG{n+nv+nvIdentifier}{ID}\PYG{+w}{ }\PYG{n+nv+nvIdentifier}{group\PYGZhy{}ID}\PYG{+w}{ }\PYG{n}{smd}\PYG{o}{/}\PYG{n}{ulsph}\PYG{o}{/}\PYG{n}{effm}
\end{sphinxVerbatim}
\begin{itemize}
\item {} 
\sphinxAtStartPar
ID, group\sphinxhyphen{}ID are documented in {\hyperref[\detokenize{compute::doc}]{\sphinxcrossref{\DUrole{doc}{compute}}}} command

\item {} 
\sphinxAtStartPar
smd/ulsph/effm = style name of this compute command

\end{itemize}


\subsection{Examples}
\label{\detokenize{compute_smd_ulsph_effm:examples}}
\begin{sphinxVerbatim}[commandchars=\\\{\}]
\PYG{k}{compute }\PYG{n+nv+nvIdentifier}{1}\PYG{+w}{ }\PYG{n+nv+nvIdentifier}{all}\PYG{+w}{ }\PYG{n}{smd}\PYG{o}{/}\PYG{n}{ulsph}\PYG{o}{/}\PYG{n}{effm}
\end{sphinxVerbatim}


\subsection{Description}
\label{\detokenize{compute_smd_ulsph_effm:description}}
\sphinxAtStartPar
Define a computation that outputs the effective shear modulus for
particles interacting via the updated Lagrangian SPH pair style.

\sphinxAtStartPar
See \sphinxhref{PDF/MACHDYN\_LAMMPS\_userguide.pdf}{this PDF guide} to using Smooth
Mach Dynamics in LAMMPS.


\subsection{Output info}
\label{\detokenize{compute_smd_ulsph_effm:output-info}}
\sphinxAtStartPar
This compute calculates a per\sphinxhyphen{}particle vector, which can be accessed
by any command that uses per\sphinxhyphen{}particle values from a compute as input.
See the {\hyperref[\detokenize{Howto_output::doc}]{\sphinxcrossref{\DUrole{doc}{Howto output}}}} page for an overview of
LAMMPS output options.

\sphinxAtStartPar
The per\sphinxhyphen{}particle vector contains the current effective per atom shear
modulus as computed by the {\hyperref[\detokenize{pair_smd_ulsph::doc}]{\sphinxcrossref{\DUrole{doc}{pair smd/ulsph}}}} pair
style.


\subsection{Restrictions}
\label{\detokenize{compute_smd_ulsph_effm:restrictions}}
\sphinxAtStartPar
This compute is part of the MACHDYN package.  It is only enabled if
LAMMPS was built with that package. See the {\hyperref[\detokenize{Build_package::doc}]{\sphinxcrossref{\DUrole{doc}{Build package}}}} page for more info. This compute can
only be used for particles which interact with the updated Lagrangian
SPH pair style.


\subsection{Related commands}
\label{\detokenize{compute_smd_ulsph_effm:related-commands}}
\sphinxAtStartPar
{\hyperref[\detokenize{pair_smd_ulsph::doc}]{\sphinxcrossref{\DUrole{doc}{pair smd/ulsph}}}}


\subsection{Default}
\label{\detokenize{compute_smd_ulsph_effm:default}}
\sphinxAtStartPar
none

\sphinxstepscope

\index{compute smd/ulsph/num/neighs@\spxentry{compute smd/ulsph/num/neighs}}\ignorespaces 

\section{compute smd/ulsph/num/neighs command}
\label{\detokenize{compute_smd_ulsph_num_neighs:compute-smd-ulsph-num-neighs-command}}\label{\detokenize{compute_smd_ulsph_num_neighs:index-0}}\label{\detokenize{compute_smd_ulsph_num_neighs::doc}}

\subsection{Syntax}
\label{\detokenize{compute_smd_ulsph_num_neighs:syntax}}
\begin{sphinxVerbatim}[commandchars=\\\{\}]
\PYG{k}{compute }\PYG{n+nv+nvIdentifier}{ID}\PYG{+w}{ }\PYG{n+nv+nvIdentifier}{group\PYGZhy{}ID}\PYG{+w}{ }\PYG{n}{smd}\PYG{o}{/}\PYG{n}{ulsph}\PYG{o}{/}\PYG{n}{num}\PYG{o}{/}\PYG{n}{neighs}
\end{sphinxVerbatim}
\begin{itemize}
\item {} 
\sphinxAtStartPar
ID, group\sphinxhyphen{}ID are documented in {\hyperref[\detokenize{compute::doc}]{\sphinxcrossref{\DUrole{doc}{compute}}}} command

\item {} 
\sphinxAtStartPar
smd/ulsph/num/neighs = style name of this compute command

\end{itemize}


\subsection{Examples}
\label{\detokenize{compute_smd_ulsph_num_neighs:examples}}
\begin{sphinxVerbatim}[commandchars=\\\{\}]
\PYG{k}{compute }\PYG{n+nv+nvIdentifier}{1}\PYG{+w}{ }\PYG{n+nv+nvIdentifier}{all}\PYG{+w}{ }\PYG{n}{smd}\PYG{o}{/}\PYG{n}{ulsph}\PYG{o}{/}\PYG{n}{num}\PYG{o}{/}\PYG{n}{neighs}
\end{sphinxVerbatim}


\subsection{Description}
\label{\detokenize{compute_smd_ulsph_num_neighs:description}}
\sphinxAtStartPar
Define a computation that returns the number of neighbor particles
inside of the smoothing kernel radius for particles interacting via
the updated Lagrangian SPH pair style.

\sphinxAtStartPar
See \sphinxhref{PDF/MACHDYN\_LAMMPS\_userguide.pdf}{this PDF guide} to using Smooth
Mach Dynamics in LAMMPS.


\subsection{Output info}
\label{\detokenize{compute_smd_ulsph_num_neighs:output-info}}
\sphinxAtStartPar
This compute returns a per\sphinxhyphen{}particle vector, which can be accessed by
any command that uses per\sphinxhyphen{}particle values from a compute as input.
See the {\hyperref[\detokenize{Howto_output::doc}]{\sphinxcrossref{\DUrole{doc}{Howto output}}}} page for an overview of
LAMMPS output options.

\sphinxAtStartPar
The per\sphinxhyphen{}particle values will be given dimensionless, see {\hyperref[\detokenize{units::doc}]{\sphinxcrossref{\DUrole{doc}{units}}}}.


\subsection{Restrictions}
\label{\detokenize{compute_smd_ulsph_num_neighs:restrictions}}
\sphinxAtStartPar
This compute is part of the MACHDYN package.  It is only enabled if
LAMMPS was built with that package.  See the {\hyperref[\detokenize{Build_package::doc}]{\sphinxcrossref{\DUrole{doc}{Build package}}}} page for more info.  This compute can
only be used for particles which interact with the updated Lagrangian
SPH pair style.


\subsection{Related commands}
\label{\detokenize{compute_smd_ulsph_num_neighs:related-commands}}
\sphinxAtStartPar
{\hyperref[\detokenize{compute_smd_tlsph_num_neighs::doc}]{\sphinxcrossref{\DUrole{doc}{compute smd/tlsph/num/neighs}}}}


\subsection{Default}
\label{\detokenize{compute_smd_ulsph_num_neighs:default}}
\sphinxAtStartPar
none

\sphinxstepscope

\index{compute smd/ulsph/strain@\spxentry{compute smd/ulsph/strain}}\ignorespaces 

\section{compute smd/ulsph/strain command}
\label{\detokenize{compute_smd_ulsph_strain:compute-smd-ulsph-strain-command}}\label{\detokenize{compute_smd_ulsph_strain:index-0}}\label{\detokenize{compute_smd_ulsph_strain::doc}}

\subsection{Syntax}
\label{\detokenize{compute_smd_ulsph_strain:syntax}}
\begin{sphinxVerbatim}[commandchars=\\\{\}]
\PYG{k}{compute }\PYG{n+nv+nvIdentifier}{ID}\PYG{+w}{ }\PYG{n+nv+nvIdentifier}{group\PYGZhy{}ID}\PYG{+w}{ }\PYG{n}{smd}\PYG{o}{/}\PYG{n}{ulsph}\PYG{o}{/}\PYG{n}{strain}
\end{sphinxVerbatim}
\begin{itemize}
\item {} 
\sphinxAtStartPar
ID, group\sphinxhyphen{}ID are documented in {\hyperref[\detokenize{compute::doc}]{\sphinxcrossref{\DUrole{doc}{compute}}}} command

\item {} 
\sphinxAtStartPar
smd/ulsph/strain = style name of this compute command

\end{itemize}


\subsection{Examples}
\label{\detokenize{compute_smd_ulsph_strain:examples}}
\begin{sphinxVerbatim}[commandchars=\\\{\}]
\PYG{k}{compute }\PYG{n+nv+nvIdentifier}{1}\PYG{+w}{ }\PYG{n+nv+nvIdentifier}{all}\PYG{+w}{ }\PYG{n}{smd}\PYG{o}{/}\PYG{n}{ulsph}\PYG{o}{/}\PYG{n}{strain}
\end{sphinxVerbatim}


\subsection{Description}
\label{\detokenize{compute_smd_ulsph_strain:description}}
\sphinxAtStartPar
Define a computation that outputs the logarithmic strain tensor.  for
particles interacting via the updated Lagrangian SPH pair style.

\sphinxAtStartPar
See \sphinxhref{PDF/MACHDYN\_LAMMPS\_userguide.pdf}{this PDF guide} to using Smooth
Mach Dynamics in LAMMPS.


\subsection{Output info}
\label{\detokenize{compute_smd_ulsph_strain:output-info}}
\sphinxAtStartPar
This compute calculates a per\sphinxhyphen{}particle tensor, which can be accessed
by any command that uses per\sphinxhyphen{}particle values from a compute as input.
See the {\hyperref[\detokenize{Howto_output::doc}]{\sphinxcrossref{\DUrole{doc}{Howto output}}}} page for an overview of
LAMMPS output options.

\sphinxAtStartPar
The per\sphinxhyphen{}particle vector has 6 entries, corresponding to the xx, yy,
zz, xy, xz, yz components of the symmetric strain rate tensor.

\sphinxAtStartPar
The per\sphinxhyphen{}particle tensor values will be given dimensionless, see
{\hyperref[\detokenize{units::doc}]{\sphinxcrossref{\DUrole{doc}{units}}}}.


\subsection{Restrictions}
\label{\detokenize{compute_smd_ulsph_strain:restrictions}}
\sphinxAtStartPar
This compute is part of the MACHDYN package.  It is only enabled if
LAMMPS was built with that package. See the {\hyperref[\detokenize{Build_package::doc}]{\sphinxcrossref{\DUrole{doc}{Build package}}}} page for more info. This compute can
only be used for particles which interact with the updated Lagrangian
SPH pair style.


\subsection{Related commands}
\label{\detokenize{compute_smd_ulsph_strain:related-commands}}
\sphinxAtStartPar
{\hyperref[\detokenize{compute_smd_tlsph_strain::doc}]{\sphinxcrossref{\DUrole{doc}{compute smd/tlsph/strain}}}}


\subsection{Default}
\label{\detokenize{compute_smd_ulsph_strain:default}}
\sphinxAtStartPar
none

\sphinxstepscope

\index{compute smd/ulsph/strain/rate@\spxentry{compute smd/ulsph/strain/rate}}\ignorespaces 

\section{compute smd/ulsph/strain/rate command}
\label{\detokenize{compute_smd_ulsph_strain_rate:compute-smd-ulsph-strain-rate-command}}\label{\detokenize{compute_smd_ulsph_strain_rate:index-0}}\label{\detokenize{compute_smd_ulsph_strain_rate::doc}}

\subsection{Syntax}
\label{\detokenize{compute_smd_ulsph_strain_rate:syntax}}
\begin{sphinxVerbatim}[commandchars=\\\{\}]
\PYG{k}{compute }\PYG{n+nv+nvIdentifier}{ID}\PYG{+w}{ }\PYG{n+nv+nvIdentifier}{group\PYGZhy{}ID}\PYG{+w}{ }\PYG{n}{smd}\PYG{o}{/}\PYG{n}{ulsph}\PYG{o}{/}\PYG{n}{strain}\PYG{o}{/}\PYG{n}{rate}
\end{sphinxVerbatim}
\begin{itemize}
\item {} 
\sphinxAtStartPar
ID, group\sphinxhyphen{}ID are documented in {\hyperref[\detokenize{compute::doc}]{\sphinxcrossref{\DUrole{doc}{compute}}}} command

\item {} 
\sphinxAtStartPar
smd/ulsph/strain/rate = style name of this compute command

\end{itemize}


\subsection{Examples}
\label{\detokenize{compute_smd_ulsph_strain_rate:examples}}
\begin{sphinxVerbatim}[commandchars=\\\{\}]
\PYG{k}{compute }\PYG{n+nv+nvIdentifier}{1}\PYG{+w}{ }\PYG{n+nv+nvIdentifier}{all}\PYG{+w}{ }\PYG{n}{smd}\PYG{o}{/}\PYG{n}{ulsph}\PYG{o}{/}\PYG{n}{strain}\PYG{o}{/}\PYG{n}{rate}
\end{sphinxVerbatim}


\subsection{Description}
\label{\detokenize{compute_smd_ulsph_strain_rate:description}}
\sphinxAtStartPar
Define a computation that outputs the rate of the logarithmic strain
tensor for particles interacting via the updated Lagrangian SPH pair
style.

\sphinxAtStartPar
See \sphinxhref{PDF/MACHDYN\_LAMMPS\_userguide.pdf}{this PDF guide} to using Smooth
Mach Dynamics in LAMMPS.


\subsection{Output info}
\label{\detokenize{compute_smd_ulsph_strain_rate:output-info}}
\sphinxAtStartPar
This compute calculates a per\sphinxhyphen{}particle vector of vectors (tensors),
which can be accessed by any command that uses per\sphinxhyphen{}particle values
from a compute as input. See the {\hyperref[\detokenize{Howto_output::doc}]{\sphinxcrossref{\DUrole{doc}{Howto output}}}} doc
page for an overview of LAMMPS output options.

\sphinxAtStartPar
The values will be given in {\hyperref[\detokenize{units::doc}]{\sphinxcrossref{\DUrole{doc}{units}}}} of one over time.

\sphinxAtStartPar
The per\sphinxhyphen{}particle vector has 6 entries, corresponding to the xx, yy,
zz, xy, xz, yz components of the symmetric strain rate tensor.


\subsection{Restrictions}
\label{\detokenize{compute_smd_ulsph_strain_rate:restrictions}}
\sphinxAtStartPar
This compute is part of the MACHDYN package.  It is only enabled if
LAMMPS was built with that package.  See the {\hyperref[\detokenize{Build_package::doc}]{\sphinxcrossref{\DUrole{doc}{Build package}}}} page for more info.

\sphinxAtStartPar
This compute can only be used for particles which interact with the
updated Lagrangian SPH pair style.


\subsection{Related commands}
\label{\detokenize{compute_smd_ulsph_strain_rate:related-commands}}
\sphinxAtStartPar
{\hyperref[\detokenize{compute_smd_tlsph_strain_rate::doc}]{\sphinxcrossref{\DUrole{doc}{compute smd/tlsph/strain/rate}}}}


\subsection{Default}
\label{\detokenize{compute_smd_ulsph_strain_rate:default}}
\sphinxAtStartPar
none

\sphinxstepscope

\index{compute smd/ulsph/stress@\spxentry{compute smd/ulsph/stress}}\ignorespaces 

\section{compute smd/ulsph/stress command}
\label{\detokenize{compute_smd_ulsph_stress:compute-smd-ulsph-stress-command}}\label{\detokenize{compute_smd_ulsph_stress:index-0}}\label{\detokenize{compute_smd_ulsph_stress::doc}}

\subsection{Syntax}
\label{\detokenize{compute_smd_ulsph_stress:syntax}}
\begin{sphinxVerbatim}[commandchars=\\\{\}]
\PYG{k}{compute }\PYG{n+nv+nvIdentifier}{ID}\PYG{+w}{ }\PYG{n+nv+nvIdentifier}{group\PYGZhy{}ID}\PYG{+w}{ }\PYG{n}{smd}\PYG{o}{/}\PYG{n}{ulsph}\PYG{o}{/}\PYG{n}{stress}
\end{sphinxVerbatim}
\begin{itemize}
\item {} 
\sphinxAtStartPar
ID, group\sphinxhyphen{}ID are documented in {\hyperref[\detokenize{compute::doc}]{\sphinxcrossref{\DUrole{doc}{compute}}}} command

\item {} 
\sphinxAtStartPar
smd/ulsph/stress = style name of this compute command

\end{itemize}


\subsection{Examples}
\label{\detokenize{compute_smd_ulsph_stress:examples}}
\begin{sphinxVerbatim}[commandchars=\\\{\}]
\PYG{k}{compute }\PYG{n+nv+nvIdentifier}{1}\PYG{+w}{ }\PYG{n+nv+nvIdentifier}{all}\PYG{+w}{ }\PYG{n}{smd}\PYG{o}{/}\PYG{n}{ulsph}\PYG{o}{/}\PYG{n}{stress}
\end{sphinxVerbatim}


\subsection{Description}
\label{\detokenize{compute_smd_ulsph_stress:description}}
\sphinxAtStartPar
Define a computation that outputs the Cauchy stress tensor.

\sphinxAtStartPar
See \sphinxhref{PDF/MACHDYN\_LAMMPS\_userguide.pdf}{this PDF guide} to using Smooth
Mach Dynamics in LAMMPS.


\subsection{Output info}
\label{\detokenize{compute_smd_ulsph_stress:output-info}}
\sphinxAtStartPar
This compute calculates a per\sphinxhyphen{}particle vector of vectors (tensors),
which can be accessed by any command that uses per\sphinxhyphen{}particle values
from a compute as input. See the {\hyperref[\detokenize{Howto_output::doc}]{\sphinxcrossref{\DUrole{doc}{Howto output}}}} doc
page for an overview of LAMMPS output options.

\sphinxAtStartPar
The values will be given in {\hyperref[\detokenize{units::doc}]{\sphinxcrossref{\DUrole{doc}{units}}}} of pressure.

\sphinxAtStartPar
The per\sphinxhyphen{}particle vector has 7 entries. The first six entries
correspond to the xx, yy, zz, xy, xz, yz components of the symmetric
Cauchy stress tensor. The seventh entry is the second invariant of the
stress tensor, i.e., the von Mises equivalent stress.


\subsection{Restrictions}
\label{\detokenize{compute_smd_ulsph_stress:restrictions}}
\sphinxAtStartPar
This compute is part of the MACHDYN package.  It is only enabled if
LAMMPS was built with that package. See the {\hyperref[\detokenize{Build_package::doc}]{\sphinxcrossref{\DUrole{doc}{Build package}}}} page for more info. This compute can
only be used for particles which interact with the updated Lagrangian
SPH pair style.


\subsection{Related commands}
\label{\detokenize{compute_smd_ulsph_stress:related-commands}}
\sphinxAtStartPar
{\hyperref[\detokenize{compute_smd_ulsph_strain::doc}]{\sphinxcrossref{\DUrole{doc}{compute smd/ulsph/strain}}}}, {\hyperref[\detokenize{compute_smd_ulsph_strain_rate::doc}]{\sphinxcrossref{\DUrole{doc}{compute smd/ulsph/strain/rate}}}} {\hyperref[\detokenize{compute_smd_tlsph_stress::doc}]{\sphinxcrossref{\DUrole{doc}{compute smd/tlsph/stress}}}}


\subsection{Default}
\label{\detokenize{compute_smd_ulsph_stress:default}}
\sphinxAtStartPar
none

\sphinxstepscope

\index{compute smd/vol@\spxentry{compute smd/vol}}\ignorespaces 

\section{compute smd/vol command}
\label{\detokenize{compute_smd_vol:compute-smd-vol-command}}\label{\detokenize{compute_smd_vol:index-0}}\label{\detokenize{compute_smd_vol::doc}}

\subsection{Syntax}
\label{\detokenize{compute_smd_vol:syntax}}
\begin{sphinxVerbatim}[commandchars=\\\{\}]
\PYG{k}{compute }\PYG{n+nv+nvIdentifier}{ID}\PYG{+w}{ }\PYG{n+nv+nvIdentifier}{group\PYGZhy{}ID}\PYG{+w}{ }\PYG{n}{smd}\PYG{o}{/}\PYG{n}{vol}
\end{sphinxVerbatim}
\begin{itemize}
\item {} 
\sphinxAtStartPar
ID, group\sphinxhyphen{}ID are documented in {\hyperref[\detokenize{compute::doc}]{\sphinxcrossref{\DUrole{doc}{compute}}}} command

\item {} 
\sphinxAtStartPar
smd/vol = style name of this compute command

\end{itemize}


\subsection{Examples}
\label{\detokenize{compute_smd_vol:examples}}
\begin{sphinxVerbatim}[commandchars=\\\{\}]
\PYG{k}{compute }\PYG{n+nv+nvIdentifier}{1}\PYG{+w}{ }\PYG{n+nv+nvIdentifier}{all}\PYG{+w}{ }\PYG{n}{smd}\PYG{o}{/}\PYG{n}{vol}
\end{sphinxVerbatim}


\subsection{Description}
\label{\detokenize{compute_smd_vol:description}}
\sphinxAtStartPar
Define a computation that provides the per\sphinxhyphen{}particle volume and the sum
of the per\sphinxhyphen{}particle volumes of the group for which the compute is defined.

\sphinxAtStartPar
See \sphinxhref{PDF/MACHDYN\_LAMMPS\_userguide.pdf}{this PDF guide} to using Smooth
Mach Dynamics in LAMMPS.


\subsection{Output info}
\label{\detokenize{compute_smd_vol:output-info}}
\sphinxAtStartPar
This compute calculates a per\sphinxhyphen{}particle vector, which can be accessed
by any command that uses per\sphinxhyphen{}particle values from a compute as input.
See the {\hyperref[\detokenize{Howto_output::doc}]{\sphinxcrossref{\DUrole{doc}{Howto output}}}} page for an overview of
LAMMPS output options.

\sphinxAtStartPar
The per\sphinxhyphen{}particle vector values will be given in {\hyperref[\detokenize{units::doc}]{\sphinxcrossref{\DUrole{doc}{units}}}} of
volume.

\sphinxAtStartPar
Additionally, the compute returns a scalar, which is the sum of the
per\sphinxhyphen{}particle volumes of the group for which the compute is defined.


\subsection{Restrictions}
\label{\detokenize{compute_smd_vol:restrictions}}
\sphinxAtStartPar
This compute is part of the MACHDYN package.  It is only enabled if
LAMMPS was built with that package. See the {\hyperref[\detokenize{Build_package::doc}]{\sphinxcrossref{\DUrole{doc}{Build package}}}} page for more info.


\subsection{Related commands}
\label{\detokenize{compute_smd_vol:related-commands}}
\sphinxAtStartPar
{\hyperref[\detokenize{compute_smd_rho::doc}]{\sphinxcrossref{\DUrole{doc}{compute smd/rho}}}}


\subsection{Default}
\label{\detokenize{compute_smd_vol:default}}
\sphinxAtStartPar
none

\sphinxstepscope

\index{compute sna/atom@\spxentry{compute sna/atom}}\ignorespaces 
\index{compute snad/atom@\spxentry{compute snad/atom}}\ignorespaces 
\index{compute snav/atom@\spxentry{compute snav/atom}}\ignorespaces 
\index{compute snap@\spxentry{compute snap}}\ignorespaces 
\index{compute sna/grid@\spxentry{compute sna/grid}}\ignorespaces 
\index{compute sna/grid/local@\spxentry{compute sna/grid/local}}\ignorespaces 

\section{compute sna/atom command}
\label{\detokenize{compute_sna_atom:compute-sna-atom-command}}\label{\detokenize{compute_sna_atom:index-5}}\label{\detokenize{compute_sna_atom:index-4}}\label{\detokenize{compute_sna_atom:index-3}}\label{\detokenize{compute_sna_atom:index-2}}\label{\detokenize{compute_sna_atom:index-1}}\label{\detokenize{compute_sna_atom:index-0}}\label{\detokenize{compute_sna_atom::doc}}

\section{compute snad/atom command}
\label{\detokenize{compute_sna_atom:compute-snad-atom-command}}

\section{compute snav/atom command}
\label{\detokenize{compute_sna_atom:compute-snav-atom-command}}

\section{compute snap command}
\label{\detokenize{compute_sna_atom:compute-snap-command}}

\section{compute sna/grid command}
\label{\detokenize{compute_sna_atom:compute-sna-grid-command}}

\section{compute sna/grid/local command}
\label{\detokenize{compute_sna_atom:compute-sna-grid-local-command}}

\subsection{Syntax}
\label{\detokenize{compute_sna_atom:syntax}}
\begin{sphinxVerbatim}[commandchars=\\\{\}]
\PYG{k}{compute }\PYG{n+nv+nvIdentifier}{ID}\PYG{+w}{ }\PYG{n+nv+nvIdentifier}{group\PYGZhy{}ID}\PYG{+w}{ }\PYG{n}{sna}\PYG{o}{/}\PYG{n}{atom}\PYG{+w}{ }\PYG{n}{rcutfac}\PYG{+w}{ }\PYG{n}{rfac0}\PYG{+w}{ }\PYG{n}{twojmax}\PYG{+w}{ }\PYG{n}{R\PYGZus{}1}\PYG{+w}{ }\PYG{n}{R\PYGZus{}2}\PYG{+w}{ }\PYG{n}{...}\PYG{+w}{ }\PYG{n}{w\PYGZus{}1}\PYG{+w}{ }\PYG{n}{w\PYGZus{}2}\PYG{+w}{ }\PYG{n}{...}\PYG{+w}{ }\PYG{n}{keyword}\PYG{+w}{ }\PYG{n}{values}\PYG{+w}{ }\PYG{n}{...}
\PYG{k}{compute }\PYG{n+nv+nvIdentifier}{ID}\PYG{+w}{ }\PYG{n+nv+nvIdentifier}{group\PYGZhy{}ID}\PYG{+w}{ }\PYG{n}{snad}\PYG{o}{/}\PYG{n}{atom}\PYG{+w}{ }\PYG{n}{rcutfac}\PYG{+w}{ }\PYG{n}{rfac0}\PYG{+w}{ }\PYG{n}{twojmax}\PYG{+w}{ }\PYG{n}{R\PYGZus{}1}\PYG{+w}{ }\PYG{n}{R\PYGZus{}2}\PYG{+w}{ }\PYG{n}{...}\PYG{+w}{ }\PYG{n}{w\PYGZus{}1}\PYG{+w}{ }\PYG{n}{w\PYGZus{}2}\PYG{+w}{ }\PYG{n}{...}\PYG{+w}{ }\PYG{n}{keyword}\PYG{+w}{ }\PYG{n}{values}\PYG{+w}{ }\PYG{n}{...}
\PYG{k}{compute }\PYG{n+nv+nvIdentifier}{ID}\PYG{+w}{ }\PYG{n+nv+nvIdentifier}{group\PYGZhy{}ID}\PYG{+w}{ }\PYG{n}{snav}\PYG{o}{/}\PYG{n}{atom}\PYG{+w}{ }\PYG{n}{rcutfac}\PYG{+w}{ }\PYG{n}{rfac0}\PYG{+w}{ }\PYG{n}{twojmax}\PYG{+w}{ }\PYG{n}{R\PYGZus{}1}\PYG{+w}{ }\PYG{n}{R\PYGZus{}2}\PYG{+w}{ }\PYG{n}{...}\PYG{+w}{ }\PYG{n}{w\PYGZus{}1}\PYG{+w}{ }\PYG{n}{w\PYGZus{}2}\PYG{+w}{ }\PYG{n}{...}\PYG{+w}{ }\PYG{n}{keyword}\PYG{+w}{ }\PYG{n}{values}\PYG{+w}{ }\PYG{n}{...}
\PYG{k}{compute }\PYG{n+nv+nvIdentifier}{ID}\PYG{+w}{ }\PYG{n+nv+nvIdentifier}{group\PYGZhy{}ID}\PYG{+w}{ }\PYG{n}{snap}\PYG{+w}{ }\PYG{n}{rcutfac}\PYG{+w}{ }\PYG{n}{rfac0}\PYG{+w}{ }\PYG{n}{twojmax}\PYG{+w}{ }\PYG{n}{R\PYGZus{}1}\PYG{+w}{ }\PYG{n}{R\PYGZus{}2}\PYG{+w}{ }\PYG{n}{...}\PYG{+w}{ }\PYG{n}{w\PYGZus{}1}\PYG{+w}{ }\PYG{n}{w\PYGZus{}2}\PYG{+w}{ }\PYG{n}{...}\PYG{+w}{ }\PYG{n}{keyword}\PYG{+w}{ }\PYG{n}{values}\PYG{+w}{ }\PYG{n}{...}
\PYG{k}{compute }\PYG{n+nv+nvIdentifier}{ID}\PYG{+w}{ }\PYG{n+nv+nvIdentifier}{group\PYGZhy{}ID}\PYG{+w}{ }\PYG{n}{snap}\PYG{+w}{ }\PYG{n}{rcutfac}\PYG{+w}{ }\PYG{n}{rfac0}\PYG{+w}{ }\PYG{n}{twojmax}\PYG{+w}{ }\PYG{n}{R\PYGZus{}1}\PYG{+w}{ }\PYG{n}{R\PYGZus{}2}\PYG{+w}{ }\PYG{n}{...}\PYG{+w}{ }\PYG{n}{w\PYGZus{}1}\PYG{+w}{ }\PYG{n}{w\PYGZus{}2}\PYG{+w}{ }\PYG{n}{...}\PYG{+w}{ }\PYG{n}{keyword}\PYG{+w}{ }\PYG{n}{values}\PYG{+w}{ }\PYG{n}{...}
\PYG{k}{compute }\PYG{n+nv+nvIdentifier}{ID}\PYG{+w}{ }\PYG{n+nv+nvIdentifier}{group\PYGZhy{}ID}\PYG{+w}{ }\PYG{n}{sna}\PYG{o}{/}\PYG{n}{grid}\PYG{+w}{ }\PYG{n}{nx}\PYG{+w}{ }\PYG{n}{ny}\PYG{+w}{ }\PYG{n}{nz}\PYG{+w}{ }\PYG{n}{rcutfac}\PYG{+w}{ }\PYG{n}{rfac0}\PYG{+w}{ }\PYG{n}{twojmax}\PYG{+w}{ }\PYG{n}{R\PYGZus{}1}\PYG{+w}{ }\PYG{n}{R\PYGZus{}2}\PYG{+w}{ }\PYG{n}{...}\PYG{+w}{ }\PYG{n}{w\PYGZus{}1}\PYG{+w}{ }\PYG{n}{w\PYGZus{}2}\PYG{+w}{ }\PYG{n}{...}\PYG{+w}{ }\PYG{n}{keyword}\PYG{+w}{ }\PYG{n}{values}\PYG{+w}{ }\PYG{n}{...}
\PYG{k}{compute }\PYG{n+nv+nvIdentifier}{ID}\PYG{+w}{ }\PYG{n+nv+nvIdentifier}{group\PYGZhy{}ID}\PYG{+w}{ }\PYG{n}{sna}\PYG{o}{/}\PYG{n}{grid}\PYG{o}{/}\PYG{n}{local}\PYG{+w}{ }\PYG{n}{nx}\PYG{+w}{ }\PYG{n}{ny}\PYG{+w}{ }\PYG{n}{nz}\PYG{+w}{ }\PYG{n}{rcutfac}\PYG{+w}{ }\PYG{n}{rfac0}\PYG{+w}{ }\PYG{n}{twojmax}\PYG{+w}{ }\PYG{n}{R\PYGZus{}1}\PYG{+w}{ }\PYG{n}{R\PYGZus{}2}\PYG{+w}{ }\PYG{n}{...}\PYG{+w}{ }\PYG{n}{w\PYGZus{}1}\PYG{+w}{ }\PYG{n}{w\PYGZus{}2}\PYG{+w}{ }\PYG{n}{...}\PYG{+w}{ }\PYG{n}{keyword}\PYG{+w}{ }\PYG{n}{values}\PYG{+w}{ }\PYG{n}{...}
\end{sphinxVerbatim}
\begin{itemize}
\item {} 
\sphinxAtStartPar
ID, group\sphinxhyphen{}ID are documented in {\hyperref[\detokenize{compute::doc}]{\sphinxcrossref{\DUrole{doc}{compute}}}} command

\item {} 
\sphinxAtStartPar
sna/atom = style name of this compute command

\item {} 
\sphinxAtStartPar
rcutfac = scale factor applied to all cutoff radii (positive real)

\item {} 
\sphinxAtStartPar
rfac0 = parameter in distance to angle conversion (0 \textless{} rcutfac \textless{} 1)

\item {} 
\sphinxAtStartPar
twojmax = band limit for bispectrum components (non\sphinxhyphen{}negative integer)

\item {} 
\sphinxAtStartPar
R\_1, R\_2,… = list of cutoff radii, one for each type (distance units)

\item {} 
\sphinxAtStartPar
w\_1, w\_2,… = list of neighbor weights, one for each type

\item {} 
\sphinxAtStartPar
nx, ny, nz = number of grid points in x, y, and z directions (positive integer)

\item {} 
\sphinxAtStartPar
zero or more keyword/value pairs may be appended

\item {} 
\sphinxAtStartPar
keyword = \sphinxstyleemphasis{rmin0} or \sphinxstyleemphasis{switchflag} or \sphinxstyleemphasis{bzeroflag} or \sphinxstyleemphasis{quadraticflag} or \sphinxstyleemphasis{chem} or \sphinxstyleemphasis{bnormflag} or \sphinxstyleemphasis{wselfallflag} or \sphinxstyleemphasis{bikflag} or \sphinxstyleemphasis{switchinnerflag} or \sphinxstyleemphasis{sinner} or \sphinxstyleemphasis{dinner} or \sphinxstyleemphasis{dgradflag} or \sphinxstyleemphasis{nnn} or \sphinxstyleemphasis{wmode} or \sphinxstyleemphasis{delta}
\begin{sphinxalltt}
\sphinxstyleemphasis{rmin0} value = parameter in distance to angle conversion (distance units)
\sphinxstyleemphasis{switchflag} value = \sphinxstyleemphasis{0} or \sphinxstyleemphasis{1}
   \sphinxstyleemphasis{0} = do not use switching function
   \sphinxstyleemphasis{1} = use switching function
\sphinxstyleemphasis{bzeroflag} value = \sphinxstyleemphasis{0} or \sphinxstyleemphasis{1}
   \sphinxstyleemphasis{0} = do not subtract B0
   \sphinxstyleemphasis{1} = subtract B0
\sphinxstyleemphasis{quadraticflag} value = \sphinxstyleemphasis{0} or \sphinxstyleemphasis{1}
   \sphinxstyleemphasis{0} = do not generate quadratic terms
   \sphinxstyleemphasis{1} = generate quadratic terms
\sphinxstyleemphasis{chem} values = \sphinxstyleemphasis{nelements} \sphinxstyleemphasis{elementlist}
   \sphinxstyleemphasis{nelements} = number of SNAP elements
   \sphinxstyleemphasis{elementlist} = \sphinxstyleemphasis{ntypes} integers in range {[}0, \sphinxstyleemphasis{nelements})
\sphinxstyleemphasis{bnormflag} value = \sphinxstyleemphasis{0} or \sphinxstyleemphasis{1}
   \sphinxstyleemphasis{0} = do not normalize
   \sphinxstyleemphasis{1} = normalize bispectrum components
\sphinxstyleemphasis{wselfallflag} value = \sphinxstyleemphasis{0} or \sphinxstyleemphasis{1}
   \sphinxstyleemphasis{0} = self\sphinxhyphen{}contribution only for element of central atom
   \sphinxstyleemphasis{1} = self\sphinxhyphen{}contribution for all elements
\sphinxstyleemphasis{switchinnerflag} value = \sphinxstyleemphasis{0} or \sphinxstyleemphasis{1}
   \sphinxstyleemphasis{0} = do not use inner switching function
   \sphinxstyleemphasis{1} = use inner switching function
\sphinxstyleemphasis{sinner} values = \sphinxstyleemphasis{sinnerlist}
   \sphinxstyleemphasis{sinnerlist} = \sphinxstyleemphasis{ntypes} values of \sphinxstyleemphasis{Sinner} (distance units)
\sphinxstyleemphasis{dinner} values = \sphinxstyleemphasis{dinnerlist}
   \sphinxstyleemphasis{dinnerlist} = \sphinxstyleemphasis{ntypes} values of \sphinxstyleemphasis{Dinner} (distance units)
\sphinxstyleemphasis{bikflag} value = \sphinxstyleemphasis{0} or \sphinxstyleemphasis{1} (only implemented for compute snap)
   \sphinxstyleemphasis{0} = descriptors are summed over atoms of each type
   \sphinxstyleemphasis{1} = descriptors are listed separately for each atom
\sphinxstyleemphasis{dgradflag} value = \sphinxstyleemphasis{0} or \sphinxstyleemphasis{1} (only implemented for compute snap)
   \sphinxstyleemphasis{0} = descriptor gradients are summed over atoms of each type
   \sphinxstyleemphasis{1} = descriptor gradients are listed separately for each atom pair
\end{sphinxalltt}

\item {} 
\sphinxAtStartPar
additional keyword = \sphinxstyleemphasis{nnn} or \sphinxstyleemphasis{wmode} or \sphinxstyleemphasis{delta}
\begin{sphinxalltt}
\sphinxstyleemphasis{nnn} value = number of considered nearest neighbors to compute the bispectrum over a target specific number of neighbors (only implemented for compute sna/atom)
\sphinxstyleemphasis{wmode} value = weight function for finding optimal cutoff to match the target number of neighbors (required if nnn used, only implemented for compute sna/atom)
   \sphinxstyleemphasis{0} = heavyside weight function
   \sphinxstyleemphasis{1} = hyperbolic tangent weight function
\sphinxstyleemphasis{delta} value = transition interval centered at cutoff distance for hyperbolic tangent weight function (ignored if wmode=0, required if wmode=1, only implemented for compute sna/atom)
\end{sphinxalltt}

\end{itemize}


\subsection{Examples}
\label{\detokenize{compute_sna_atom:examples}}
\begin{sphinxVerbatim}[commandchars=\\\{\}]
\PYG{k}{compute }\PYG{n+nv+nvIdentifier}{b}\PYG{+w}{ }\PYG{n+nv+nvIdentifier}{all}\PYG{+w}{ }\PYG{n}{sna}\PYG{o}{/}\PYG{n}{atom}\PYG{+w}{ }\PYG{l+m}{1.4}\PYG{+w}{ }\PYG{l+m}{0.99363}\PYG{+w}{ }\PYG{l+m}{6}\PYG{+w}{ }\PYG{l+m}{2.0}\PYG{+w}{ }\PYG{l+m}{2.4}\PYG{+w}{ }\PYG{l+m}{0.75}\PYG{+w}{ }\PYG{l+m}{1.0}\PYG{+w}{ }\PYG{n}{rmin0}\PYG{+w}{ }\PYG{l+m}{0.0}
\PYG{k}{compute }\PYG{n+nv+nvIdentifier}{db}\PYG{+w}{ }\PYG{n+nv+nvIdentifier}{all}\PYG{+w}{ }\PYG{n}{sna}\PYG{o}{/}\PYG{n}{atom}\PYG{+w}{ }\PYG{l+m}{1.4}\PYG{+w}{ }\PYG{l+m}{0.95}\PYG{+w}{ }\PYG{l+m}{6}\PYG{+w}{ }\PYG{l+m}{2.0}\PYG{+w}{ }\PYG{l+m}{1.0}
\PYG{k}{compute }\PYG{n+nv+nvIdentifier}{vb}\PYG{+w}{ }\PYG{n+nv+nvIdentifier}{all}\PYG{+w}{ }\PYG{n}{sna}\PYG{o}{/}\PYG{n}{atom}\PYG{+w}{ }\PYG{l+m}{1.4}\PYG{+w}{ }\PYG{l+m}{0.95}\PYG{+w}{ }\PYG{l+m}{6}\PYG{+w}{ }\PYG{l+m}{2.0}\PYG{+w}{ }\PYG{l+m}{1.0}
\PYG{k}{compute }\PYG{n+nv+nvIdentifier}{snap}\PYG{+w}{ }\PYG{n+nv+nvIdentifier}{all}\PYG{+w}{ }\PYG{n}{snap}\PYG{+w}{ }\PYG{l+m}{1.4}\PYG{+w}{ }\PYG{l+m}{0.95}\PYG{+w}{ }\PYG{l+m}{6}\PYG{+w}{ }\PYG{l+m}{2.0}\PYG{+w}{ }\PYG{l+m}{1.0}
\PYG{k}{compute }\PYG{n+nv+nvIdentifier}{snap}\PYG{+w}{ }\PYG{n+nv+nvIdentifier}{all}\PYG{+w}{ }\PYG{n}{snap}\PYG{+w}{ }\PYG{l+m}{1.0}\PYG{+w}{ }\PYG{l+m}{0.99363}\PYG{+w}{ }\PYG{l+m}{6}\PYG{+w}{ }\PYG{l+m}{3.81}\PYG{+w}{ }\PYG{l+m}{3.83}\PYG{+w}{ }\PYG{l+m}{1.0}\PYG{+w}{ }\PYG{l+m}{0.93}\PYG{+w}{ }\PYG{n}{chem}\PYG{+w}{ }\PYG{l+m}{2}\PYG{+w}{ }\PYG{l+m}{0}\PYG{+w}{ }\PYG{l+m}{1}
\PYG{k}{compute }\PYG{n+nv+nvIdentifier}{snap}\PYG{+w}{ }\PYG{n+nv+nvIdentifier}{all}\PYG{+w}{ }\PYG{n}{snap}\PYG{+w}{ }\PYG{l+m}{1.0}\PYG{+w}{ }\PYG{l+m}{0.99363}\PYG{+w}{ }\PYG{l+m}{6}\PYG{+w}{ }\PYG{l+m}{3.81}\PYG{+w}{ }\PYG{l+m}{3.83}\PYG{+w}{ }\PYG{l+m}{1.0}\PYG{+w}{ }\PYG{l+m}{0.93}\PYG{+w}{ }\PYG{n}{switchinnerflag}\PYG{+w}{ }\PYG{l+m}{1}\PYG{+w}{ }\PYG{n}{sinner}\PYG{+w}{ }\PYG{l+m}{1.35}\PYG{+w}{ }\PYG{l+m}{1.6}\PYG{+w}{ }\PYG{n}{dinner}\PYG{+w}{ }\PYG{l+m}{0.25}\PYG{+w}{ }\PYG{l+m}{0.3}
\PYG{k}{compute }\PYG{n+nv+nvIdentifier}{bgrid}\PYG{+w}{ }\PYG{n+nv+nvIdentifier}{all}\PYG{+w}{ }\PYG{n}{sna}\PYG{o}{/}\PYG{n}{grid}\PYG{o}{/}\PYG{n}{local}\PYG{+w}{ }\PYG{l+m}{200}\PYG{+w}{ }\PYG{l+m}{200}\PYG{+w}{ }\PYG{l+m}{200}\PYG{+w}{ }\PYG{l+m}{1.4}\PYG{+w}{ }\PYG{l+m}{0.95}\PYG{+w}{ }\PYG{l+m}{6}\PYG{+w}{ }\PYG{l+m}{2.0}\PYG{+w}{ }\PYG{l+m}{1.0}
\PYG{k}{compute }\PYG{n+nv+nvIdentifier}{bnnn}\PYG{+w}{ }\PYG{n+nv+nvIdentifier}{all}\PYG{+w}{ }\PYG{n}{sna}\PYG{o}{/}\PYG{n}{atom}\PYG{+w}{ }\PYG{l+m}{9.0}\PYG{+w}{ }\PYG{l+m}{0.99363}\PYG{+w}{ }\PYG{l+m}{8}\PYG{+w}{ }\PYG{l+m}{0.5}\PYG{+w}{ }\PYG{l+m}{1.0}\PYG{+w}{ }\PYG{n}{rmin0}\PYG{+w}{ }\PYG{l+m}{0.0}\PYG{+w}{ }\PYG{n}{nnn}\PYG{+w}{ }\PYG{l+m}{24}\PYG{+w}{ }\PYG{n}{wmode}\PYG{+w}{ }\PYG{l+m}{1}\PYG{+w}{ }\PYG{n}{delta}\PYG{+w}{ }\PYG{l+m}{0.2}
\end{sphinxVerbatim}


\subsection{Description}
\label{\detokenize{compute_sna_atom:description}}
\sphinxAtStartPar
Define a computation that calculates a set of quantities related to the
bispectrum components of the atoms in a group. These computes are used
primarily for calculating the dependence of energy, force, and stress
components on the linear coefficients in the {\hyperref[\detokenize{pair_snap::doc}]{\sphinxcrossref{\DUrole{doc}{snap pair\_style}}}}, which is useful when training a SNAP potential to match
target data.

\sphinxAtStartPar
Bispectrum components of an atom are order parameters characterizing the
radial and angular distribution of neighbor atoms. The detailed
mathematical definition is given in the paper by Thompson et
al. {\hyperref[\detokenize{compute_sna_atom:thompson20141}]{\sphinxcrossref{\DUrole{std,std-ref}{(Thompson)}}}}

\sphinxAtStartPar
The position of a neighbor atom \sphinxstyleemphasis{i’} relative to a central atom \sphinxstyleemphasis{i} is a
point within the 3D ball of radius \(R_{ii'}\) = \sphinxstyleemphasis{rcutfac}
\((R_i + R_i')\)

\sphinxAtStartPar
Bartok et al. {\hyperref[\detokenize{compute_sna_atom:bartok20101}]{\sphinxcrossref{\DUrole{std,std-ref}{(Bartok)}}}}, proposed mapping this 3D
ball onto the 3\sphinxhyphen{}sphere, the surface of the unit ball in a
four\sphinxhyphen{}dimensional space.  The radial distance \sphinxstyleemphasis{r} within \sphinxstyleemphasis{R\_ii’} is
mapped on to a third polar angle \(\theta_0\) defined by,
\begin{equation*}
\begin{split}\theta_0 = {\sf rfac0} \frac{r-r_{min0}}{R_{ii'}-r_{min0}} \pi\end{split}
\end{equation*}
\sphinxAtStartPar
In this way, all possible neighbor positions are mapped on to a subset
of the 3\sphinxhyphen{}sphere.  Points south of the latitude \(\theta_0\) =
\sphinxstyleemphasis{rfac0} \(\pi\) are excluded.

\sphinxAtStartPar
The natural basis for functions on the 3\sphinxhyphen{}sphere is formed by the
representatives of \sphinxstyleemphasis{SU(2)}, the matrices \(U^j_{m,m'}(\theta, \phi,
\theta_0)\).  These functions are better known as \(D^j_{m,m'}\), the
elements of the Wigner \sphinxstyleemphasis{D}\sphinxhyphen{}matrices {\hyperref[\detokenize{compute_sna_atom:meremianin2006}]{\sphinxcrossref{\DUrole{std,std-ref}{(Meremianin}}}}, {\hyperref[\detokenize{compute_sna_atom:varshalovich1987}]{\sphinxcrossref{\DUrole{std,std-ref}{Varshalovich}}}}, {\hyperref[\detokenize{compute_sna_atom:mason2009}]{\sphinxcrossref{\DUrole{std,std-ref}{Mason)}}}} The density of neighbors on the 3\sphinxhyphen{}sphere can be written as
a sum of Dirac\sphinxhyphen{}delta functions, one for each neighbor, weighted by
species and radial distance. Expanding this density function as a
generalized Fourier series in the basis functions, we can write each
Fourier coefficient as
\begin{equation*}
\begin{split}u^j_{m,m'} = U^j_{m,m'}(0,0,0) + \sum_{r_{ii'} < R_{ii'}}{f_c(r_{ii'}) w_{\mu_{i'}} U^j_{m,m'}(\theta_0,\theta,\phi)}\end{split}
\end{equation*}
\sphinxAtStartPar
The \(w_{\mu_{i'}}\) neighbor weights are dimensionless numbers that
depend on \(\mu_{i'}\), the SNAP element of atom \sphinxstyleemphasis{i’}, while the
central atom is arbitrarily assigned a unit weight.  The function
\(f_c(r)\) ensures that the contribution of each neighbor atom goes
smoothly to zero at \(R_{ii'}\):
\begin{equation*}
\begin{split}f_c(r)   = & \frac{1}{2}(\cos(\pi \frac{r-r_{min0}}{R_{ii'}-r_{min0}}) + 1), r \leq R_{ii'} \\
         = & 0,  r > R_{ii'}\end{split}
\end{equation*}
\sphinxAtStartPar
The expansion coefficients \(u^j_{m,m'}\) are complex\sphinxhyphen{}valued and
they are not directly useful as descriptors, because they are not
invariant under rotation of the polar coordinate frame. However, the
following scalar triple products of expansion coefficients can be shown
to be real\sphinxhyphen{}valued and invariant under rotation {\hyperref[\detokenize{compute_sna_atom:bartok20101}]{\sphinxcrossref{\DUrole{std,std-ref}{(Bartok)}}}}.
\begin{equation*}
\begin{split}B_{j_1,j_2,j}  =
\sum_{m_1,m'_1=-j_1}^{j_1}\sum_{m_2,m'_2=-j_2}^{j_2}\sum_{m,m'=-j}^{j} (u^j_{m,m'})^*
H {\scriptscriptstyle \begin{array}{l} {j} {m} {m'} \\
     {j_1} {m_1} {m'_1} \\
     {j_2} {m_2} {m'_2} \end{array}}
     u^{j_1}_{m_1,m'_1} u^{j_2}_{m_2,m'_2}\end{split}
\end{equation*}
\sphinxAtStartPar
The constants \(H^{jmm'}_{j_1 m_1 m_{1'},j_2 m_ 2m_{2'}}\) are
coupling coefficients, analogous to Clebsch\sphinxhyphen{}Gordan coefficients for
rotations on the 2\sphinxhyphen{}sphere. These invariants are the components of the
bispectrum and these are the quantities calculated by the compute
\sphinxstyleemphasis{sna/atom}. They characterize the strength of density correlations at
three points on the 3\sphinxhyphen{}sphere. The j2=0 subset form the power spectrum,
which characterizes the correlations of two points. The lowest\sphinxhyphen{}order
components describe the coarsest features of the density function, while
higher\sphinxhyphen{}order components reflect finer detail. Each bispectrum component
contains terms that depend on the positions of up to 4 atoms (3
neighbors and the central atom).

\sphinxAtStartPar
Compute \sphinxstyleemphasis{snad/atom} calculates the derivative of the bispectrum
components summed separately for each LAMMPS atom type:
\begin{equation*}
\begin{split}-\sum_{i' \in I} \frac{\partial {B^{i'}_{j_1,j_2,j}  }}{\partial {\bf r}_i}\end{split}
\end{equation*}
\sphinxAtStartPar
The sum is over all atoms \sphinxstyleemphasis{i’} of atom type \sphinxstyleemphasis{I}.  For each atom \sphinxstyleemphasis{i},
this compute evaluates the above expression for each direction, each
atom type, and each bispectrum component.  See section below on output
for a detailed explanation.

\sphinxAtStartPar
Compute \sphinxstyleemphasis{snav/atom} calculates the virial contribution due to the
derivatives:
\begin{equation*}
\begin{split}-{\bf r}_i \otimes \sum_{i' \in I} \frac{\partial {B^{i'}_{j_1,j_2,j}}}{\partial {\bf r}_i}\end{split}
\end{equation*}
\sphinxAtStartPar
Again, the sum is over all atoms \sphinxstyleemphasis{i’} of atom type \sphinxstyleemphasis{I}.  For each atom
\sphinxstyleemphasis{i}, this compute evaluates the above expression for each of the six
virial components, each atom type, and each bispectrum component.  See
section below on output for a detailed explanation.

\sphinxAtStartPar
Compute \sphinxstyleemphasis{snap} calculates a global array containing information related
to all three of the above per\sphinxhyphen{}atom computes \sphinxstyleemphasis{sna/atom}, \sphinxstyleemphasis{snad/atom},
and \sphinxstyleemphasis{snav/atom}. The first row of the array contains the summation of
\sphinxstyleemphasis{sna/atom} over all atoms, but broken out by type. The last six rows of
the array contain the summation of \sphinxstyleemphasis{snav/atom} over all atoms, broken
out by type. In between these are 3*\sphinxstyleemphasis{N} rows containing the same
values computed by \sphinxstyleemphasis{snad/atom} (these are already summed over all atoms
and broken out by type). The element in the last column of each row
contains the potential energy, force, or stress, according to the row.
These quantities correspond to the user\sphinxhyphen{}specified reference potential
that must be subtracted from the target data when fitting SNAP.  The
potential energy calculation uses the built in compute \sphinxstyleemphasis{thermo\_pe}.  The
stress calculation uses a compute called \sphinxstyleemphasis{snap\_press} that is
automatically created behind the scenes, according to the following
command:

\begin{sphinxVerbatim}[commandchars=\\\{\}]
\PYG{k}{compute }\PYG{n+nv+nvIdentifier}{snap\PYGZus{}press}\PYG{+w}{ }\PYG{n+nv+nvIdentifier}{all}\PYG{+w}{ }\PYG{n}{pressure}\PYG{+w}{ }\PYG{n}{NULL}\PYG{+w}{ }\PYG{n}{virial}
\end{sphinxVerbatim}

\sphinxAtStartPar
See section below on output for a detailed explanation of the data
layout in the global array.

\sphinxAtStartPar
\DUrole{versionmodified,added}{Added in version 3Aug2022.}

\sphinxAtStartPar
The compute \sphinxstyleemphasis{sna/grid} and \sphinxstyleemphasis{sna/grid/local} commands calculate
bispectrum components for a regular grid of points.  These are
calculated from the local density of nearby atoms \sphinxstyleemphasis{i’} around each grid
point, as if there was a central atom \sphinxstyleemphasis{i} at the grid point. This is
useful for characterizing fine\sphinxhyphen{}scale structure in a configuration of
atoms, and it is used in the \sphinxhref{https://github.com/casus/mala}{MALA package} to build machine\sphinxhyphen{}learning surrogates
for finite\sphinxhyphen{}temperature Kohn\sphinxhyphen{}Sham density functional theory ({\hyperref[\detokenize{compute_sna_atom:ellis2021}]{\sphinxcrossref{\DUrole{std,std-ref}{Ellis
et al.}}}}) Neighbor atoms not in the group do not contribute
to the bispectrum components of the grid points. The distance cutoff
\(R_{ii'}\) assumes that \sphinxstyleemphasis{i} has the same type as the neighbor atom
\sphinxstyleemphasis{i’}.

\sphinxAtStartPar
Compute \sphinxstyleemphasis{sna/grid} calculates a global array containing bispectrum
components for a regular grid of points.
The grid is aligned with the current box dimensions, with the
first point at the box origin, and forming a regular 3d array with
\sphinxstyleemphasis{nx}, \sphinxstyleemphasis{ny}, and \sphinxstyleemphasis{nz} points in the x, y, and z directions. For triclinic
boxes, the array is congruent with the periodic lattice vectors
a, b, and c. The array contains one row for each of the
\(nx \times ny \times nz\) grid points, looping over the index for \sphinxstyleemphasis{ix} fastest,
then \sphinxstyleemphasis{iy}, and \sphinxstyleemphasis{iz} slowest.  Each row of the array contains the \sphinxstyleemphasis{x}, \sphinxstyleemphasis{y},
and \sphinxstyleemphasis{z} coordinates of the grid point, followed by the bispectrum
components. See section below on output for a detailed explanation of the data
layout in the global array.

\sphinxAtStartPar
Compute \sphinxstyleemphasis{sna/grid/local} calculates bispectrum components of a regular
grid of points similarly to compute \sphinxstyleemphasis{sna/grid} described above.
However, because the array is local, it contains only rows for grid points
that are local to the processor subdomain. The global grid
of \(nx \times ny \times nz\) points is still laid out in space the same as for \sphinxstyleemphasis{sna/grid},
but grid points are strictly partitioned, so that every grid point appears in
one and only one local array.  The array contains one row for each of the
local grid points, looping over the global index \sphinxstyleemphasis{ix} fastest,
then \sphinxstyleemphasis{iy}, and \sphinxstyleemphasis{iz} slowest.  Each row of the array contains
the global indexes \sphinxstyleemphasis{ix}, \sphinxstyleemphasis{iy}, and \sphinxstyleemphasis{iz} first, followed by the \sphinxstyleemphasis{x}, \sphinxstyleemphasis{y},
and \sphinxstyleemphasis{z} coordinates of the grid point, followed by the bispectrum
components. See section below on output for a detailed explanation of the data
layout in the global array.

\sphinxAtStartPar
The value of all bispectrum components will be zero for atoms not in
the group. Neighbor atoms not in the group do not contribute to the
bispectrum of atoms in the group.

\sphinxAtStartPar
The neighbor list needed to compute this quantity is constructed each
time the calculation is performed (i.e. each time a snapshot of atoms
is dumped).  Thus it can be inefficient to compute/dump this quantity
too frequently.

\sphinxAtStartPar
The argument \sphinxstyleemphasis{rcutfac} is a scale factor that controls the ratio of
atomic radius to radial cutoff distance.

\sphinxAtStartPar
The argument \sphinxstyleemphasis{rfac0} and the optional keyword \sphinxstyleemphasis{rmin0} define the
linear mapping from radial distance to polar angle \(theta_0\) on the
3\sphinxhyphen{}sphere, given above.

\sphinxAtStartPar
The argument \sphinxstyleemphasis{twojmax} defines which
bispectrum components are generated. See section below on output for a
detailed explanation of the number of bispectrum components and the
ordered in which they are listed.

\sphinxAtStartPar
The keyword \sphinxstyleemphasis{switchflag} can be used to turn off the switching
function \(f_c(r)\).

\sphinxAtStartPar
The keyword \sphinxstyleemphasis{bzeroflag} determines whether or not \sphinxstyleemphasis{B0}, the bispectrum
components of an atom with no neighbors, are subtracted from the
calculated bispectrum components. This optional keyword normally only
affects compute \sphinxstyleemphasis{sna/atom}. However, when \sphinxstyleemphasis{quadraticflag} is on, it
also affects \sphinxstyleemphasis{snad/atom} and \sphinxstyleemphasis{snav/atom}.

\sphinxAtStartPar
The keyword \sphinxstyleemphasis{quadraticflag} determines whether or not the quadratic
combinations of bispectrum quantities are generated.  These are formed
by taking the outer product of the vector of bispectrum components with
itself.  See section below on output for a detailed explanation of the
number of quadratic terms and the ordered in which they are listed.

\sphinxAtStartPar
The keyword \sphinxstyleemphasis{chem} activates the explicit multi\sphinxhyphen{}element variant of the
SNAP bispectrum components. The argument \sphinxstyleemphasis{nelements} specifies the
number of SNAP elements that will be handled.  This is followed by
\sphinxstyleemphasis{elementlist}, a list of integers of length \sphinxstyleemphasis{ntypes}, with values in the
range {[}0, \sphinxstyleemphasis{nelements} ), which maps each LAMMPS type to one of the SNAP
elements.  Note that multiple LAMMPS types can be mapped to the same
element, and some elements may be mapped by no LAMMPS type. However, in
typical use cases (training SNAP potentials) the mapping from LAMMPS
types to elements is one\sphinxhyphen{}to\sphinxhyphen{}one.

\sphinxAtStartPar
The explicit multi\sphinxhyphen{}element variant invoked by the \sphinxstyleemphasis{chem} keyword
partitions the density of neighbors into partial densities for each
chemical element.  This is described in detail in the paper by
{\hyperref[\detokenize{compute_sna_atom:cusentino2020}]{\sphinxcrossref{\DUrole{std,std-ref}{Cusentino et al.}}}} The bispectrum components are
indexed on ordered triplets of elements:
\begin{equation*}
\begin{split}B_{j_1,j_2,j}^{\kappa\lambda\mu} =
\sum_{m_1,m'_1=-j_1}^{j_1}\sum_{m_2,m'_2=-j_2}^{j_2}\sum_{m,m'=-j}^{j} (u^{\mu}_{j,m,m'})^*
H {\scriptscriptstyle \begin{array}{l} {j} {m} {m'} \\
     {j_1} {m_1} {m'_1} \\
     {j_2} {m_2} {m'_2} \end{array}}
     u^{\kappa}_{j_1,m_1,m'_1} u^{\lambda}_{j_2,m_2,m'_2}\end{split}
\end{equation*}
\sphinxAtStartPar
where \(u^{\mu}_{j,m,m'}\) is an expansion coefficient for the partial density of neighbors
of element \(\mu\)
\begin{equation*}
\begin{split}u^{\mu}_{j,m,m'} =  w^{self}_{\mu_{i}\mu} U^{j,m,m'}(0,0,0) + \sum_{r_{ii'} < R_{ii'}}{\delta_{\mu\mu_{i'}}f_c(r_{ii'}) w_{\mu_{i'}} U^{j,m,m'}(\theta_0,\theta,\phi)}\end{split}
\end{equation*}
\sphinxAtStartPar
where \(w^{self}_{\mu_{i}\mu}\) is the self\sphinxhyphen{}contribution, which is
either 1 or 0 (see keyword \sphinxstyleemphasis{wselfallflag} below),
\(\delta_{\mu\mu_{i'}}\) indicates that the sum is only over
neighbor atoms of element \(\mu\), and all other quantities are the
same as those appearing in the original equation for \(u^j_{m,m'}\)
given above.

\sphinxAtStartPar
The keyword \sphinxstyleemphasis{wselfallflag} defines the rule used for the
self\sphinxhyphen{}contribution.  If \sphinxstyleemphasis{wselfallflag} is on, then
\(w^{self}_{\mu_{i}\mu}\) = 1. If it is off then
\(w^{self}_{\mu_{i}\mu}\) = 0, except in the case of
\({\mu_{i}=\mu}\), when \(w^{self}_{\mu_{i}\mu}\) = 1.  When the
\sphinxstyleemphasis{chem} keyword is not used, this keyword has no effect.

\sphinxAtStartPar
The keyword \sphinxstyleemphasis{bnormflag} determines whether or not the bispectrum
component \(B_{j_1,j_2,j}\) is divided by a factor of \(2j+1\).
This normalization simplifies force calculations because of the
following symmetry relation
\begin{equation*}
\begin{split}\frac{B_{j_1,j_2,j}}{2j+1} = \frac{B_{j,j_2,j_1}}{2j_1+1} = \frac{B_{j_1,j,j_2}}{2j_2+1}\end{split}
\end{equation*}
\sphinxAtStartPar
This option is typically used in conjunction with the \sphinxstyleemphasis{chem} keyword,
and LAMMPS will generate a warning if both \sphinxstyleemphasis{chem} and \sphinxstyleemphasis{bnormflag}
are not both set or not both unset.

\sphinxAtStartPar
The keyword \sphinxstyleemphasis{switchinnerflag} with value 1
activates an additional radial switching
function similar to \(f_c(r)\) above, but acting to switch off
smoothly contributions from neighbor atoms at short separation distances.
This is useful when SNAP is used in combination with a simple
repulsive potential. For a neighbor atom at
distance \(r\), its contribution is scaled by a multiplicative
factor \(f_{inner}(r)\) defined as follows:
\begin{equation*}
\begin{split}             = & 0,  r \leq S_{inner} - D_{inner} \\
f_{inner}(r) = & \frac{1}{2}(1 - \cos(\frac{\pi}{2} (1 + \frac{r-S_{inner}}{D_{inner}})), S_{inner} - D_{inner} < r \leq S_{inner} + D_{inner} \\
             = & 1,  r > S_{inner} + D_{inner}\end{split}
\end{equation*}
\sphinxAtStartPar
where the switching region is centered at \(S_{inner}\) and it extends a distance \(D_{inner}\)
to the left and to the right of this.
With this option, additional keywords \sphinxstyleemphasis{sinner} and \sphinxstyleemphasis{dinner} must be used,
each followed by \sphinxstyleemphasis{ntypes}
values for \(S_{inner}\) and \(D_{inner}\), respectively.
When the central atom and the neighbor atom have different types,
the values of \(S_{inner}\) and \(D_{inner}\) are
the arithmetic means of the values for both types.

\sphinxAtStartPar
The keywords \sphinxstyleemphasis{bikflag} and \sphinxstyleemphasis{dgradflag} are only used by compute \sphinxstyleemphasis{snap}.
The keyword \sphinxstyleemphasis{bikflag} determines whether or not to list the descriptors
of each atom separately, or sum them together and list in a single row.
If \sphinxstyleemphasis{bikflag} is set
to \sphinxstyleemphasis{0} then a single bispectrum row is used, which contains the per\sphinxhyphen{}atom bispectrum
descriptors \(B_{i,k}\) summed over all atoms \sphinxstyleemphasis{i} to produce
\(B_k\).  If \sphinxstyleemphasis{bikflag} is set
to \sphinxstyleemphasis{1} this is replaced by a separate per\sphinxhyphen{}atom bispectrum row for each atom.
In this case, the entries in the final column for these rows
are set to zero.

\sphinxAtStartPar
The keyword \sphinxstyleemphasis{dgradflag} determines whether to sum atom gradients or list
them separately. If \sphinxstyleemphasis{dgradflag} is set to 0, the bispectrum
descriptor gradients w.r.t. atom \sphinxstyleemphasis{j} are summed over all atoms \sphinxstyleemphasis{i’}
of type \sphinxstyleemphasis{I} (similar to \sphinxstyleemphasis{snad/atom} above).
If \sphinxstyleemphasis{dgradflag} is set to 1, gradients are listed separately for each pair of atoms.
Each row corresponds
to a single term \(\frac{\partial {B_{i,k}  }}{\partial {r}^a_j}\)
where \({r}^a_j\) is the \sphinxstyleemphasis{a\sphinxhyphen{}th} position coordinate of the atom with global
index \sphinxstyleemphasis{j}. This also changes
the number of columns to be equal to the number of bispectrum components, with 3
additional columns representing the indices \(i\), \(j\), and \(a\),
as explained more in the Output info section below. The option \sphinxstyleemphasis{dgradflag=1}
requires that \sphinxstyleemphasis{bikflag=1}.

\begin{sphinxadmonition}{note}{Note:}
\sphinxAtStartPar
Using \sphinxstyleemphasis{dgradflag} = 1 produces a global array with \(N + 3N^2 + 1\) rows
which becomes expensive for systems with more than 1000 atoms.
\end{sphinxadmonition}

\begin{sphinxadmonition}{note}{Note:}
\sphinxAtStartPar
If you have a bonded system, then the settings of {\hyperref[\detokenize{special_bonds::doc}]{\sphinxcrossref{\DUrole{doc}{special\_bonds}}}} command can remove pairwise interactions between
atoms in the same bond, angle, or dihedral.  This is the default
setting for the {\hyperref[\detokenize{special_bonds::doc}]{\sphinxcrossref{\DUrole{doc}{special\_bonds}}}} command, and
means those pairwise interactions do not appear in the neighbor list.
Because this fix uses the neighbor list, it also means those pairs
will not be included in the calculation.  One way to get around this,
is to write a dump file, and use the {\hyperref[\detokenize{rerun::doc}]{\sphinxcrossref{\DUrole{doc}{rerun}}}} command to
compute the bispectrum components for snapshots in the dump file.
The rerun script can use a {\hyperref[\detokenize{special_bonds::doc}]{\sphinxcrossref{\DUrole{doc}{special\_bonds}}}}
command that includes all pairs in the neighbor list.
\end{sphinxadmonition}

\sphinxAtStartPar
The keyword \sphinxstyleemphasis{nnn} allows for the calculation of the bispectrum over a
specific target number of neighbors. This option is only implemented for
the compute \sphinxstyleemphasis{sna/atom}.  An optimal cutoff radius for defining the
neighborhood of the central atom is calculated by means of a dichotomy
algorithm.  This iterative process allows to assign weights to
neighboring atoms in order to match the total sum of weights with the
target number of neighbors.  Depending on the radial weight function
used in that process, the cutoff radius can fluctuate a lot in the
presence of thermal noise.  Therefore, in addition to the \sphinxstyleemphasis{nnn} keyword,
the keyword \sphinxstyleemphasis{wmode} allows to choose whether a Heaviside (\sphinxstyleemphasis{wmode} = 0)
function or a Hyperbolic tangent function (\sphinxstyleemphasis{wmode} = 1) should be used.
If the Heaviside function is used, the cutoff radius exactly matches the
distance between the central atom an its \sphinxstyleemphasis{nnn}’th neighbor.  However, in
the case of the hyperbolic tangent function, the dichotomy algorithm
allows to span the weights over a distance \sphinxstyleemphasis{delta} in order to reduce
fluctuations in the resulting local atomic environment fingerprint.  The
detailed formalism is given in the paper by Lafourcade et
al. {\hyperref[\detokenize{compute_sna_atom:lafourcade2023-2}]{\sphinxcrossref{\DUrole{std,std-ref}{(Lafourcade)}}}}.


\bigskip\hrule\bigskip



\subsection{Output info}
\label{\detokenize{compute_sna_atom:output-info}}
\sphinxAtStartPar
Compute \sphinxstyleemphasis{sna/atom} calculates a per\sphinxhyphen{}atom array, each column
corresponding to a particular bispectrum component.  The total number of
columns and the identity of the bispectrum component contained in each
column depend of the value of \sphinxstyleemphasis{twojmax}, as described by the following
piece of python code:

\begin{sphinxVerbatim}[commandchars=\\\{\}]
for j1 in range(0,twojmax+1):
    for j2 in range(0,j1+1):
        for j in range(j1\PYGZhy{}j2,min(twojmax,j1+j2)+1,2):
            if (j\PYGZgt{}=j1): print j1/2.,j2/2.,j/2.
\end{sphinxVerbatim}

\sphinxAtStartPar
There are \(m(m+1)/2\) descriptors with last index \sphinxstyleemphasis{j},
where \sphinxstyleemphasis{m} = \(\lfloor j \rfloor + 1\).
Hence, for even \sphinxstyleemphasis{twojmax} = 2(\sphinxstyleemphasis{m}\sphinxhyphen{}1), \(K = m(m+1)(2m+1)/6\), the \sphinxstyleemphasis{m}\sphinxhyphen{}th pyramidal number,
and for odd \sphinxstyleemphasis{twojmax} = 2 \sphinxstyleemphasis{m}\sphinxhyphen{}1, \(K = m(m+1)(m+2)/3\), twice the \sphinxstyleemphasis{m}\sphinxhyphen{}th tetrahedral number.

\begin{sphinxadmonition}{note}{Note:}
\sphinxAtStartPar
the \sphinxstyleemphasis{diagonal} keyword allowing other possible choices
for the number of bispectrum components was removed in 2019,
since all potentials use the value of 3, corresponding to the
above set of bispectrum components.
\end{sphinxadmonition}

\sphinxAtStartPar
Compute \sphinxstyleemphasis{snad/atom} evaluates a per\sphinxhyphen{}atom array. The columns are arranged
into \sphinxstyleemphasis{ntypes} blocks, listed in order of atom type \sphinxstyleemphasis{I}.  Each block
contains three sub\sphinxhyphen{}blocks corresponding to the \sphinxstyleemphasis{x}, \sphinxstyleemphasis{y}, and \sphinxstyleemphasis{z}
components of the atom position.  Each of these sub\sphinxhyphen{}blocks contains \sphinxstyleemphasis{K}
columns for the \sphinxstyleemphasis{K} bispectrum components, the same as for compute
\sphinxstyleemphasis{sna/atom}

\sphinxAtStartPar
Compute \sphinxstyleemphasis{snav/atom} evaluates a per\sphinxhyphen{}atom array. The columns are arranged
into \sphinxstyleemphasis{ntypes} blocks, listed in order of atom type \sphinxstyleemphasis{I}.  Each block
contains six sub\sphinxhyphen{}blocks corresponding to the \sphinxstyleemphasis{xx}, \sphinxstyleemphasis{yy}, \sphinxstyleemphasis{zz},
\sphinxstyleemphasis{yz}, \sphinxstyleemphasis{xz}, and \sphinxstyleemphasis{xy} components of the virial tensor in Voigt
notation.  Each of these sub\sphinxhyphen{}blocks contains \sphinxstyleemphasis{K} columns for the \sphinxstyleemphasis{K}
bispectrum components, the same as for compute \sphinxstyleemphasis{sna/atom}

\sphinxAtStartPar
Compute \sphinxstyleemphasis{snap} evaluates a global array.  The columns are arranged into
\sphinxstyleemphasis{ntypes} blocks, listed in order of atom type \sphinxstyleemphasis{I}. Each block contains
one column for each bispectrum component, the same as for compute
\sphinxstyleemphasis{sna/atom}. A final column contains the corresponding energy, force
component on an atom, or virial stress component. The rows of the array
appear in the following order:
\begin{itemize}
\item {} 
\sphinxAtStartPar
1 row: \sphinxstyleemphasis{sna/atom} quantities summed for all atoms of type \sphinxstyleemphasis{I}

\item {} 
\sphinxAtStartPar
3*\sphinxstyleemphasis{N} rows: \sphinxstyleemphasis{snad/atom} quantities, with derivatives w.r.t. x, y, and z coordinate of atom \sphinxstyleemphasis{i} appearing in consecutive rows. The atoms are sorted based on atom ID.

\item {} 
\sphinxAtStartPar
6 rows: \sphinxstyleemphasis{snav/atom} quantities summed for all atoms of type \sphinxstyleemphasis{I}

\end{itemize}

\sphinxAtStartPar
For example, if \sphinxstyleemphasis{K} =30 and ntypes=1, the number of columns in the
per\sphinxhyphen{}atom arrays generated by \sphinxstyleemphasis{sna/atom}, \sphinxstyleemphasis{snad/atom}, and
\sphinxstyleemphasis{snav/atom} are 30, 90, and 180, respectively. With \sphinxstyleemphasis{quadratic} value=1,
the numbers of columns are 930, 2790, and 5580, respectively.  The
number of columns in the global array generated by \sphinxstyleemphasis{snap} are 31, and
931, respectively, while the number of rows is 1+3*\sphinxstyleemphasis{N}+6, where \sphinxstyleemphasis{N}
is the total number of atoms.

\sphinxAtStartPar
Compute \sphinxstyleemphasis{sna/grid} evaluates a global array.
The array contains one row for each of the
\(nx \times ny \times nz\) grid points, looping over the index for \sphinxstyleemphasis{ix} fastest,
then \sphinxstyleemphasis{iy}, and \sphinxstyleemphasis{iz} slowest.  Each row of the array contains the \sphinxstyleemphasis{x}, \sphinxstyleemphasis{y},
and \sphinxstyleemphasis{z} coordinates of the grid point, followed by the bispectrum
components.

\sphinxAtStartPar
Compute \sphinxstyleemphasis{sna/grid/local} evaluates a local array.
The array contains one row for each of the
local grid points, looping over the global index \sphinxstyleemphasis{ix} fastest,
then \sphinxstyleemphasis{iy}, and \sphinxstyleemphasis{iz} slowest.  Each row of the array contains
the global indexes \sphinxstyleemphasis{ix}, \sphinxstyleemphasis{iy}, and \sphinxstyleemphasis{iz} first, followed by the \sphinxstyleemphasis{x}, \sphinxstyleemphasis{y},
and \sphinxstyleemphasis{z} coordinates of the grid point, followed by the bispectrum
components.

\sphinxAtStartPar
If the \sphinxstyleemphasis{quadratic} keyword value is set to 1, then additional columns
are generated, corresponding to the products of all distinct pairs of
bispectrum components. If the number of bispectrum components is \sphinxstyleemphasis{K},
then the number of distinct pairs is \sphinxstyleemphasis{K}(\sphinxstyleemphasis{K}+1)/2.  For compute
\sphinxstyleemphasis{sna/atom} these columns are appended to existing \sphinxstyleemphasis{K} columns.  The
ordering of quadratic terms is upper\sphinxhyphen{}triangular, (1,1),(1,2)…(1,\sphinxstyleemphasis{K}),(2,1)…(\sphinxstyleemphasis{K}\sphinxhyphen{}1,\sphinxstyleemphasis{K}\sphinxhyphen{}1),(\sphinxstyleemphasis{K}\sphinxhyphen{}1,\sphinxstyleemphasis{K}),(\sphinxstyleemphasis{K},\sphinxstyleemphasis{K}).
For computes \sphinxstyleemphasis{snad/atom} and \sphinxstyleemphasis{snav/atom} each set of \sphinxstyleemphasis{K}(\sphinxstyleemphasis{K}+1)/2
additional columns is inserted directly after each of sub\sphinxhyphen{}block of
linear terms i.e. linear and quadratic terms are contiguous.  So the
nesting order from inside to outside is bispectrum component, linear
then quadratic, vector/tensor component, type.

\sphinxAtStartPar
If the \sphinxstyleemphasis{chem} keyword is used, then the data is arranged into
\(N_{elem}^3\) sub\sphinxhyphen{}blocks, each sub\sphinxhyphen{}block corresponding to a
particular chemical labeling \(\kappa\lambda\mu\) with the last
label changing fastest.  Each sub\sphinxhyphen{}block contains \sphinxstyleemphasis{K} bispectrum
components. For the purposes of handling contributions to force, virial,
and quadratic combinations, these \(N_{elem}^3\) sub\sphinxhyphen{}blocks are
treated as a single block of \(K N_{elem}^3\) columns.

\sphinxAtStartPar
If the \sphinxstyleemphasis{bik} keyword is set to 1, the structure of the snap array is expanded.
The first \(N\) rows of the snap array
correspond to \(B_{i,k}\) instead of a single row summed over atoms \(i\).
In this case, the entries in the final column for these rows
are set to zero. Also, each row contains only non\sphinxhyphen{}zero entries for the
columns corresponding to the type of that atom. This is not true in the case
of \sphinxstyleemphasis{dgradflag} keyword = 1 (see below).

\sphinxAtStartPar
If the \sphinxstyleemphasis{dgradflag} keyword is set to 1, this changes the structure of the
global array completely.
Here the \sphinxstyleemphasis{snad/atom} quantities are replaced with rows corresponding to
descriptor gradient components on single atoms:
\begin{equation*}
\begin{split}\frac{\partial {B_{i,k}  }}{\partial {r}^a_j}\end{split}
\end{equation*}
\sphinxAtStartPar
where \({r}^a_j\) is the \sphinxstyleemphasis{a\sphinxhyphen{}th} position coordinate of the atom with global
index \sphinxstyleemphasis{j}. The rows are
organized in chunks, where each chunk corresponds to an atom with global index
\(j\). The rows in an atom \(j\) chunk correspond to
atoms with global index \(i\). The total number of rows for
these descriptor gradients is therefore \(3N^2\).
The number of columns is equal to the number of bispectrum components,
plus 3 additional left\sphinxhyphen{}most columns representing the global atom indices
\(i\), \(j\),
and Cartesian direction \(a\)  (0, 1, 2, for x, y, z).
The first 3 columns of the first \(N\) rows belong to the reference
potential force components. The remaining K columns contain the
\(B_{i,k}\) per\sphinxhyphen{}atom descriptors corresponding to the non\sphinxhyphen{}zero entries
obtained when \sphinxstyleemphasis{bikflag} = 1.
The first column of the last row, after the first
\(N + 3N^2\) rows, contains the reference potential
energy. The virial components are not used with this option. The total number of
rows is therefore \(N + 3N^2 + 1\) and the number of columns is \(K + 3\).

\sphinxAtStartPar
These values can be accessed by any command that uses per\sphinxhyphen{}atom values
from a compute as input.  See the {\hyperref[\detokenize{Howto_output::doc}]{\sphinxcrossref{\DUrole{doc}{Howto output}}}} doc
page for an overview of LAMMPS output options. To see how this command
can be used within a Python workflow to train SNAP potentials, see the
examples in \sphinxhref{https://github.com/FitSNAP/FitSNAP}{FitSNAP}.


\subsection{Restrictions}
\label{\detokenize{compute_sna_atom:restrictions}}
\sphinxAtStartPar
These computes are part of the ML\sphinxhyphen{}SNAP package.  They are only enabled
if LAMMPS was built with that package.  See the {\hyperref[\detokenize{Build_package::doc}]{\sphinxcrossref{\DUrole{doc}{Build package}}}} page for more info.


\subsection{Related commands}
\label{\detokenize{compute_sna_atom:related-commands}}
\sphinxAtStartPar
{\hyperref[\detokenize{pair_snap::doc}]{\sphinxcrossref{\DUrole{doc}{pair\_style snap}}}}
{\hyperref[\detokenize{compute_slcsa_atom::doc}]{\sphinxcrossref{\DUrole{doc}{compute slcsa/atom}}}}


\subsection{Default}
\label{\detokenize{compute_sna_atom:default}}
\sphinxAtStartPar
The optional keyword defaults are \sphinxstyleemphasis{rmin0} = 0,
\sphinxstyleemphasis{switchflag} = 1, \sphinxstyleemphasis{bzeroflag} = 1, \sphinxstyleemphasis{quadraticflag} = 0,
\sphinxstyleemphasis{bnormflag} = 0, \sphinxstyleemphasis{wselfallflag} = 0, \sphinxstyleemphasis{switchinnerflag} = 0,
\sphinxstyleemphasis{nnn} = \sphinxhyphen{}1, \sphinxstyleemphasis{wmode} = 0, \sphinxstyleemphasis{delta} = 1.e\sphinxhyphen{}3


\bigskip\hrule\bigskip

\phantomsection\label{\detokenize{compute_sna_atom:thompson20141}}
\sphinxAtStartPar
\sphinxstylestrong{(Thompson)} Thompson, Swiler, Trott, Foiles, Tucker, J Comp Phys, 285, 316, (2015).

\phantomsection\label{\detokenize{compute_sna_atom:bartok20101}}
\sphinxAtStartPar
\sphinxstylestrong{(Bartok)} Bartok, Payne, Risi, Csanyi, Phys Rev Lett, 104, 136403 (2010).

\phantomsection\label{\detokenize{compute_sna_atom:meremianin2006}}
\sphinxAtStartPar
\sphinxstylestrong{(Meremianin)} Meremianin, J. Phys. A,  39, 3099 (2006).

\phantomsection\label{\detokenize{compute_sna_atom:varshalovich1987}}
\sphinxAtStartPar
\sphinxstylestrong{(Varshalovich)} Varshalovich, Moskalev, Khersonskii, Quantum Theory
of Angular Momentum, World Scientific, Singapore (1987).

\phantomsection\label{\detokenize{compute_sna_atom:mason2009}}
\sphinxAtStartPar
\sphinxstylestrong{(Mason)} J. K. Mason, Acta Cryst A65, 259 (2009).

\phantomsection\label{\detokenize{compute_sna_atom:cusentino2020}}
\sphinxAtStartPar
\sphinxstylestrong{(Cusentino)} Cusentino, Wood, Thompson, J Phys Chem A, 124, 5456, (2020)

\phantomsection\label{\detokenize{compute_sna_atom:ellis2021}}
\sphinxAtStartPar
\sphinxstylestrong{(Ellis)} Ellis, Fiedler, Popoola, Modine, Stephens, Thompson, Cangi, Rajamanickam,  Phys Rev B, 104, 035120, (2021)

\phantomsection\label{\detokenize{compute_sna_atom:lafourcade2023-2}}
\sphinxAtStartPar
\sphinxstylestrong{(Lafourcade)} Lafourcade, Maillet, Denoual, Duval, Allera, Goryaeva, and Marinica,
\sphinxhref{https://doi.org/10.1016/j.commatsci.2023.112534}{Comp. Mat. Science, 230, 112534 (2023)}

\sphinxstepscope

\index{compute sph/e/atom@\spxentry{compute sph/e/atom}}\ignorespaces 

\section{compute sph/e/atom command}
\label{\detokenize{compute_sph_e_atom:compute-sph-e-atom-command}}\label{\detokenize{compute_sph_e_atom:index-0}}\label{\detokenize{compute_sph_e_atom::doc}}

\subsection{Syntax}
\label{\detokenize{compute_sph_e_atom:syntax}}
\begin{sphinxVerbatim}[commandchars=\\\{\}]
\PYG{k}{compute }\PYG{n+nv+nvIdentifier}{ID}\PYG{+w}{ }\PYG{n+nv+nvIdentifier}{group\PYGZhy{}ID}\PYG{+w}{ }\PYG{n}{sph}\PYG{o}{/}\PYG{n}{e}\PYG{o}{/}\PYG{n}{atom}
\end{sphinxVerbatim}
\begin{itemize}
\item {} 
\sphinxAtStartPar
ID, group\sphinxhyphen{}ID are documented in {\hyperref[\detokenize{compute::doc}]{\sphinxcrossref{\DUrole{doc}{compute}}}} command

\item {} 
\sphinxAtStartPar
sph/e/atom = style name of this compute command

\end{itemize}


\subsection{Examples}
\label{\detokenize{compute_sph_e_atom:examples}}
\begin{sphinxVerbatim}[commandchars=\\\{\}]
\PYG{k}{compute }\PYG{n+nv+nvIdentifier}{1}\PYG{+w}{ }\PYG{n+nv+nvIdentifier}{all}\PYG{+w}{ }\PYG{n}{sph}\PYG{o}{/}\PYG{n}{e}\PYG{o}{/}\PYG{n}{atom}
\end{sphinxVerbatim}


\subsection{Description}
\label{\detokenize{compute_sph_e_atom:description}}
\sphinxAtStartPar
Define a computation that calculates the per\sphinxhyphen{}atom internal energy
for each atom in a group.

\sphinxAtStartPar
The internal energy is the energy associated with the internal degrees
of freedom of an SPH particle, i.e. a Smooth\sphinxhyphen{}Particle Hydrodynamics
particle.

\sphinxAtStartPar
See \sphinxhref{PDF/SPH\_LAMMPS\_userguide.pdf}{this PDF guide} to using SPH in
LAMMPS.

\sphinxAtStartPar
The value of the internal energy will be 0.0 for atoms not in the
specified compute group.


\subsection{Output info}
\label{\detokenize{compute_sph_e_atom:output-info}}
\sphinxAtStartPar
This compute calculates a per\sphinxhyphen{}atom vector, which can be accessed by
any command that uses per\sphinxhyphen{}atom values from a compute as input.  See
the {\hyperref[\detokenize{Howto_output::doc}]{\sphinxcrossref{\DUrole{doc}{Howto output}}}} page for an overview of
LAMMPS output options.

\sphinxAtStartPar
The per\sphinxhyphen{}atom vector values will be in energy {\hyperref[\detokenize{units::doc}]{\sphinxcrossref{\DUrole{doc}{units}}}}.


\subsection{Restrictions}
\label{\detokenize{compute_sph_e_atom:restrictions}}
\sphinxAtStartPar
This compute is part of the SPH package.  It is only enabled if
LAMMPS was built with that package.  See the {\hyperref[\detokenize{Build_package::doc}]{\sphinxcrossref{\DUrole{doc}{Build package}}}} page for more info.


\subsection{Related commands}
\label{\detokenize{compute_sph_e_atom:related-commands}}
\sphinxAtStartPar
{\hyperref[\detokenize{dump::doc}]{\sphinxcrossref{\DUrole{doc}{dump custom}}}}


\subsection{Default}
\label{\detokenize{compute_sph_e_atom:default}}
\sphinxAtStartPar
none

\sphinxstepscope

\index{compute sph/rho/atom@\spxentry{compute sph/rho/atom}}\ignorespaces 

\section{compute sph/rho/atom command}
\label{\detokenize{compute_sph_rho_atom:compute-sph-rho-atom-command}}\label{\detokenize{compute_sph_rho_atom:index-0}}\label{\detokenize{compute_sph_rho_atom::doc}}

\subsection{Syntax}
\label{\detokenize{compute_sph_rho_atom:syntax}}
\begin{sphinxVerbatim}[commandchars=\\\{\}]
\PYG{k}{compute }\PYG{n+nv+nvIdentifier}{ID}\PYG{+w}{ }\PYG{n+nv+nvIdentifier}{group\PYGZhy{}ID}\PYG{+w}{ }\PYG{n}{sph}\PYG{o}{/}\PYG{n}{rho}\PYG{o}{/}\PYG{n}{atom}
\end{sphinxVerbatim}
\begin{itemize}
\item {} 
\sphinxAtStartPar
ID, group\sphinxhyphen{}ID are documented in {\hyperref[\detokenize{compute::doc}]{\sphinxcrossref{\DUrole{doc}{compute}}}} command

\item {} 
\sphinxAtStartPar
sph/rho/atom = style name of this compute command

\end{itemize}


\subsection{Examples}
\label{\detokenize{compute_sph_rho_atom:examples}}
\begin{sphinxVerbatim}[commandchars=\\\{\}]
\PYG{k}{compute }\PYG{n+nv+nvIdentifier}{1}\PYG{+w}{ }\PYG{n+nv+nvIdentifier}{all}\PYG{+w}{ }\PYG{n}{sph}\PYG{o}{/}\PYG{n}{rho}\PYG{o}{/}\PYG{n}{atom}
\end{sphinxVerbatim}


\subsection{Description}
\label{\detokenize{compute_sph_rho_atom:description}}
\sphinxAtStartPar
Define a computation that calculates the per\sphinxhyphen{}atom SPH density for each
atom in a group, i.e. a Smooth\sphinxhyphen{}Particle Hydrodynamics density.

\sphinxAtStartPar
The SPH density is the mass density of an SPH particle, calculated by
kernel function interpolation using “pair style sph/rhosum”.

\sphinxAtStartPar
See \sphinxhref{PDF/SPH\_LAMMPS\_userguide.pdf}{this PDF guide} to using SPH in
LAMMPS.

\sphinxAtStartPar
The value of the SPH density will be 0.0 for atoms not in the
specified compute group.


\subsection{Output info}
\label{\detokenize{compute_sph_rho_atom:output-info}}
\sphinxAtStartPar
This compute calculates a per\sphinxhyphen{}atom vector, which can be accessed by
any command that uses per\sphinxhyphen{}atom values from a compute as input.  See
the {\hyperref[\detokenize{Howto_output::doc}]{\sphinxcrossref{\DUrole{doc}{Howto output}}}} page for an overview of
LAMMPS output options.

\sphinxAtStartPar
The per\sphinxhyphen{}atom vector values will be in mass/volume {\hyperref[\detokenize{units::doc}]{\sphinxcrossref{\DUrole{doc}{units}}}}.


\subsection{Restrictions}
\label{\detokenize{compute_sph_rho_atom:restrictions}}
\sphinxAtStartPar
This compute is part of the SPH package.  It is only enabled if
LAMMPS was built with that package.  See the {\hyperref[\detokenize{Build_package::doc}]{\sphinxcrossref{\DUrole{doc}{Build package}}}} page for more info.


\subsection{Related commands}
\label{\detokenize{compute_sph_rho_atom:related-commands}}
\sphinxAtStartPar
{\hyperref[\detokenize{dump::doc}]{\sphinxcrossref{\DUrole{doc}{dump custom}}}}


\subsection{Default}
\label{\detokenize{compute_sph_rho_atom:default}}
\sphinxAtStartPar
none

\sphinxstepscope

\index{compute sph/t/atom@\spxentry{compute sph/t/atom}}\ignorespaces 

\section{compute sph/t/atom command}
\label{\detokenize{compute_sph_t_atom:compute-sph-t-atom-command}}\label{\detokenize{compute_sph_t_atom:index-0}}\label{\detokenize{compute_sph_t_atom::doc}}

\subsection{Syntax}
\label{\detokenize{compute_sph_t_atom:syntax}}
\begin{sphinxVerbatim}[commandchars=\\\{\}]
\PYG{k}{compute }\PYG{n+nv+nvIdentifier}{ID}\PYG{+w}{ }\PYG{n+nv+nvIdentifier}{group\PYGZhy{}ID}\PYG{+w}{ }\PYG{n}{sph}\PYG{o}{/}\PYG{n}{t}\PYG{o}{/}\PYG{n}{atom}
\end{sphinxVerbatim}
\begin{itemize}
\item {} 
\sphinxAtStartPar
ID, group\sphinxhyphen{}ID are documented in {\hyperref[\detokenize{compute::doc}]{\sphinxcrossref{\DUrole{doc}{compute}}}} command

\item {} 
\sphinxAtStartPar
sph/t/atom = style name of this compute command

\end{itemize}


\subsection{Examples}
\label{\detokenize{compute_sph_t_atom:examples}}
\begin{sphinxVerbatim}[commandchars=\\\{\}]
\PYG{k}{compute }\PYG{n+nv+nvIdentifier}{1}\PYG{+w}{ }\PYG{n+nv+nvIdentifier}{all}\PYG{+w}{ }\PYG{n}{sph}\PYG{o}{/}\PYG{n}{t}\PYG{o}{/}\PYG{n}{atom}
\end{sphinxVerbatim}


\subsection{Description}
\label{\detokenize{compute_sph_t_atom:description}}
\sphinxAtStartPar
Define a computation that calculates the per\sphinxhyphen{}atom internal temperature
for each atom in a group.

\sphinxAtStartPar
The internal temperature is the ratio of internal energy over the heat
capacity associated with the internal degrees of freedom of an SPH
particles, i.e. a Smooth\sphinxhyphen{}Particle Hydrodynamics particle.
\begin{equation*}
\begin{split}T_{int} = E_{int} / C_{V,int}\end{split}
\end{equation*}
\sphinxAtStartPar
See \sphinxhref{PDF/SPH\_LAMMPS\_userguide.pdf}{this PDF guide} to using SPH in
LAMMPS.

\sphinxAtStartPar
The value of the internal energy will be 0.0 for atoms not in the
specified compute group.


\subsection{Output info}
\label{\detokenize{compute_sph_t_atom:output-info}}
\sphinxAtStartPar
This compute calculates a per\sphinxhyphen{}atom vector, which can be accessed by
any command that uses per\sphinxhyphen{}atom values from a compute as input.  See
the {\hyperref[\detokenize{Howto_output::doc}]{\sphinxcrossref{\DUrole{doc}{Howto output}}}} page for an overview of
LAMMPS output options.

\sphinxAtStartPar
The per\sphinxhyphen{}atom vector values will be in temperature {\hyperref[\detokenize{units::doc}]{\sphinxcrossref{\DUrole{doc}{units}}}}.


\subsection{Restrictions}
\label{\detokenize{compute_sph_t_atom:restrictions}}
\sphinxAtStartPar
This compute is part of the SPH package.  It is only enabled if
LAMMPS was built with that package.  See the {\hyperref[\detokenize{Build_package::doc}]{\sphinxcrossref{\DUrole{doc}{Build package}}}} page for more info.


\subsection{Related commands}
\label{\detokenize{compute_sph_t_atom:related-commands}}
\sphinxAtStartPar
{\hyperref[\detokenize{dump::doc}]{\sphinxcrossref{\DUrole{doc}{dump custom}}}}


\subsection{Default}
\label{\detokenize{compute_sph_t_atom:default}}
\sphinxAtStartPar
none

\sphinxstepscope

\index{compute spin@\spxentry{compute spin}}\ignorespaces 

\section{compute spin command}
\label{\detokenize{compute_spin:compute-spin-command}}\label{\detokenize{compute_spin:index-0}}\label{\detokenize{compute_spin::doc}}

\subsection{Syntax}
\label{\detokenize{compute_spin:syntax}}
\begin{sphinxVerbatim}[commandchars=\\\{\}]
\PYG{k}{compute }\PYG{n+nv+nvIdentifier}{ID}\PYG{+w}{ }\PYG{n+nv+nvIdentifier}{group\PYGZhy{}ID}\PYG{+w}{ }\PYG{n}{spin}
\end{sphinxVerbatim}
\begin{itemize}
\item {} 
\sphinxAtStartPar
ID, group\sphinxhyphen{}ID are documented in {\hyperref[\detokenize{compute::doc}]{\sphinxcrossref{\DUrole{doc}{compute}}}} command

\item {} 
\sphinxAtStartPar
spin = style name of this compute command

\end{itemize}


\subsection{Examples}
\label{\detokenize{compute_spin:examples}}
\begin{sphinxVerbatim}[commandchars=\\\{\}]
\PYG{k}{compute }\PYG{n+nv+nvIdentifier}{out\PYGZus{}mag}\PYG{+w}{ }\PYG{n+nv+nvIdentifier}{all}\PYG{+w}{ }\PYG{n}{spin}
\end{sphinxVerbatim}


\subsection{Description}
\label{\detokenize{compute_spin:description}}
\sphinxAtStartPar
Define a computation that calculates magnetic quantities for a system
of atoms having spins.

\sphinxAtStartPar
This compute calculates the following 6 magnetic quantities:
\begin{itemize}
\item {} 
\sphinxAtStartPar
the three first quantities are the x,y and z coordinates of the total
magnetization,

\item {} 
\sphinxAtStartPar
the fourth quantity is the norm of the total magnetization,

\item {} 
\sphinxAtStartPar
The fifth quantity is the magnetic energy (in eV),

\item {} 
\sphinxAtStartPar
The sixth one is referred to as the spin temperature, according
to the work of {\hyperref[\detokenize{compute_spin:nurdin1}]{\sphinxcrossref{\DUrole{std,std-ref}{(Nurdin)}}}}.

\end{itemize}

\sphinxAtStartPar
The simplest way to output the results of the compute spin calculation
is to define some of the quantities as variables, and to use the thermo and
thermo\_style commands, for example:

\begin{sphinxVerbatim}[commandchars=\\\{\}]
\PYG{k}{compute }\PYG{n+nv+nvIdentifier}{out\PYGZus{}mag}\PYG{+w}{         }\PYG{n+nv+nvIdentifier}{all}\PYG{+w}{ }\PYG{n}{spin}

\PYG{k}{variable }\PYG{n+nv+nvIdentifier}{mag\PYGZus{}z}\PYG{+w}{          }\PYG{n}{equal}\PYG{+w}{ }\PYG{n}{c\PYGZus{}out\PYGZus{}mag[3]}
\PYG{k}{variable }\PYG{n+nv+nvIdentifier}{mag\PYGZus{}norm}\PYG{+w}{       }\PYG{n}{equal}\PYG{+w}{ }\PYG{n}{c\PYGZus{}out\PYGZus{}mag[4]}
\PYG{k}{variable }\PYG{n+nv+nvIdentifier}{temp\PYGZus{}mag}\PYG{+w}{       }\PYG{n}{equal}\PYG{+w}{ }\PYG{n}{c\PYGZus{}out\PYGZus{}mag[6]}

\PYG{k}{thermo}\PYG{+w}{                  }\PYG{l+m}{10}
\PYG{k}{thermo\PYGZus{}style}\PYG{+w}{            }\PYG{n}{custom}\PYG{+w}{ }\PYG{n}{step}\PYG{+w}{ }\PYG{n}{v\PYGZus{}mag\PYGZus{}z}\PYG{+w}{ }\PYG{n}{v\PYGZus{}mag\PYGZus{}norm}\PYG{+w}{ }\PYG{n}{v\PYGZus{}temp\PYGZus{}mag}
\end{sphinxVerbatim}

\sphinxAtStartPar
This series of commands evaluates the total magnetization along z, the norm of
the total magnetization, and the magnetic temperature. Three variables are
assigned to those quantities. The thermo and thermo\_style commands print them
every 10 timesteps.


\subsection{Output info}
\label{\detokenize{compute_spin:output-info}}
\sphinxAtStartPar
The array values are “intensive”.  The array values will be in
metal units ({\hyperref[\detokenize{units::doc}]{\sphinxcrossref{\DUrole{doc}{units}}}}).


\subsection{Restrictions}
\label{\detokenize{compute_spin:restrictions}}
\sphinxAtStartPar
The \sphinxstyleemphasis{spin} compute is part of the SPIN package.  This compute is only
enabled if LAMMPS was built with this package.  See the {\hyperref[\detokenize{Build_package::doc}]{\sphinxcrossref{\DUrole{doc}{Build package}}}} page for more info.  The atom\_style
has to be “spin” for this compute to be valid.

\sphinxAtStartPar
\sphinxstylestrong{Related commands:}

\sphinxAtStartPar
none


\subsection{Default}
\label{\detokenize{compute_spin:default}}
\sphinxAtStartPar
none


\bigskip\hrule\bigskip

\phantomsection\label{\detokenize{compute_spin:nurdin1}}
\sphinxAtStartPar
\sphinxstylestrong{(Nurdin)} Nurdin and Schotte Phys Rev E, 61(4), 3579 (2000)

\sphinxstepscope

\index{compute stress/atom@\spxentry{compute stress/atom}}\ignorespaces 
\index{compute centroid/stress/atom@\spxentry{compute centroid/stress/atom}}\ignorespaces 

\section{compute stress/atom command}
\label{\detokenize{compute_stress_atom:compute-stress-atom-command}}\label{\detokenize{compute_stress_atom:index-1}}\label{\detokenize{compute_stress_atom:index-0}}\label{\detokenize{compute_stress_atom::doc}}

\section{compute centroid/stress/atom command}
\label{\detokenize{compute_stress_atom:compute-centroid-stress-atom-command}}

\subsection{Syntax}
\label{\detokenize{compute_stress_atom:syntax}}
\begin{sphinxVerbatim}[commandchars=\\\{\}]
\PYG{k}{compute }\PYG{n+nv+nvIdentifier}{ID}\PYG{+w}{ }\PYG{n+nv+nvIdentifier}{group\PYGZhy{}ID}\PYG{+w}{ }\PYG{n}{style}\PYG{+w}{ }\PYG{n}{temp}\PYG{o}{\PYGZhy{}}\PYG{n}{ID}\PYG{+w}{ }\PYG{n}{keyword}\PYG{+w}{ }\PYG{n}{...}
\end{sphinxVerbatim}
\begin{itemize}
\item {} 
\sphinxAtStartPar
ID, group\sphinxhyphen{}ID are documented in {\hyperref[\detokenize{compute::doc}]{\sphinxcrossref{\DUrole{doc}{compute}}}} command

\item {} 
\sphinxAtStartPar
style = \sphinxstyleemphasis{stress/atom} or \sphinxstyleemphasis{centroid/stress/atom}

\item {} 
\sphinxAtStartPar
temp\sphinxhyphen{}ID = ID of compute that calculates temperature, can be NULL if not needed

\item {} 
\sphinxAtStartPar
zero or more keywords may be appended

\item {} 
\sphinxAtStartPar
keyword = \sphinxstyleemphasis{ke} or \sphinxstyleemphasis{pair} or \sphinxstyleemphasis{bond} or \sphinxstyleemphasis{angle} or \sphinxstyleemphasis{dihedral} or \sphinxstyleemphasis{improper} or \sphinxstyleemphasis{kspace} or \sphinxstyleemphasis{fix} or \sphinxstyleemphasis{virial}

\end{itemize}


\subsection{Examples}
\label{\detokenize{compute_stress_atom:examples}}
\begin{sphinxVerbatim}[commandchars=\\\{\}]
\PYG{k}{compute }\PYG{n+nv+nvIdentifier}{1}\PYG{+w}{ }\PYG{n+nv+nvIdentifier}{mobile}\PYG{+w}{ }\PYG{n}{stress}\PYG{o}{/}\PYG{n}{atom}\PYG{+w}{ }\PYG{n}{NULL}
\PYG{k}{compute }\PYG{n+nv+nvIdentifier}{1}\PYG{+w}{ }\PYG{n+nv+nvIdentifier}{mobile}\PYG{+w}{ }\PYG{n}{stress}\PYG{o}{/}\PYG{n}{atom}\PYG{+w}{ }\PYG{n}{myRamp}
\PYG{k}{compute }\PYG{n+nv+nvIdentifier}{1}\PYG{+w}{ }\PYG{n+nv+nvIdentifier}{all}\PYG{+w}{ }\PYG{n}{stress}\PYG{o}{/}\PYG{n}{atom}\PYG{+w}{ }\PYG{n}{NULL}\PYG{+w}{ }\PYG{n}{pair}\PYG{+w}{ }\PYG{n}{bond}
\PYG{k}{compute }\PYG{n+nv+nvIdentifier}{1}\PYG{+w}{ }\PYG{n+nv+nvIdentifier}{all}\PYG{+w}{ }\PYG{n}{centroid}\PYG{o}{/}\PYG{n}{stress}\PYG{o}{/}\PYG{n}{atom}\PYG{+w}{ }\PYG{n}{NULL}\PYG{+w}{ }\PYG{n}{bond}\PYG{+w}{ }\PYG{n}{dihedral}\PYG{+w}{ }\PYG{n}{improper}
\end{sphinxVerbatim}


\subsection{Description}
\label{\detokenize{compute_stress_atom:description}}
\sphinxAtStartPar
Define a computation that computes per\sphinxhyphen{}atom stress tensor for each
atom in a group.  In case of compute \sphinxstyleemphasis{stress/atom}, the tensor for
each atom is symmetric with 6 components and is stored as a 6\sphinxhyphen{}element
vector in the following order: \(xx\), \(yy\), \(zz\),
\(xy\), \(xz\), \(yz\).  In case of compute
\sphinxstyleemphasis{centroid/stress/atom}, the tensor for each atom is asymmetric with 9
components and is stored as a 9\sphinxhyphen{}element vector in the following order:
\(xx\), \(yy\), \(zz\), \(xy\), \(xz\),
\(yz\), \(yx\), \(zx\), \(zy\).  See the {\hyperref[\detokenize{compute_pressure::doc}]{\sphinxcrossref{\DUrole{doc}{compute
pressure}}}} command if you want the stress tensor
(pressure) of the entire system.

\sphinxAtStartPar
The stress tensor for atom \(I\) is given by the following
formula, where \(a\) and \(b\) take on values \(x\),
\(y\), \(z\) to generate the components of the tensor:
\begin{equation*}
\begin{split}S_{ab}  =  - m v_a v_b - W_{ab}\end{split}
\end{equation*}
\sphinxAtStartPar
The first term is a kinetic energy contribution for atom \(I\).
See details below on how the specified \sphinxstyleemphasis{temp\sphinxhyphen{}ID} can affect the
velocities used in this calculation. The second term is the virial
contribution due to intra and intermolecular interactions, where the
exact computation details are determined by the compute style.

\sphinxAtStartPar
In case of compute \sphinxstyleemphasis{stress/atom}, the virial contribution is:
\begin{equation*}
\begin{split} W_{ab} & = \frac{1}{2} \sum_{n = 1}^{N_p} (r_{1_a} F_{1_b} + r_{2_a} F_{2_b}) + \frac{1}{2} \sum_{n = 1}^{N_b} (r_{1_a} F_{1_b} + r_{2_a} F_{2_b})  \\
& + \frac{1}{3} \sum_{n = 1}^{N_a} (r_{1_a} F_{1_b} + r_{2_a} F_{2_b} + r_{3_a} F_{3_b}) + \frac{1}{4} \sum_{n = 1}^{N_d} (r_{1_a} F_{1_b} + r_{2_a} F_{2_b} + r_{3_a} F_{3_b} + r_{4_a} F_{4_b}) \\
& + \frac{1}{4} \sum_{n = 1}^{N_i} (r_{1_a} F_{1_b} + r_{2_a} F_{2_b} + r_{3_a} F_{3_b} + r_{4_a} F_{4_b}) + {\rm Kspace}(r_{i_a},F_{i_b}) + \sum_{n = 1}^{N_f} r_{i_a} F_{i_b}\end{split}
\end{equation*}
\sphinxAtStartPar
The first term is a pairwise energy contribution where \(n\) loops
over the \(N_p\) neighbors of atom \(I\), \(\mathbf{r}_1\)
and \(\mathbf{r}_2\) are the positions of the two atoms in the
pairwise interaction, and \(\mathbf{F}_1\) and
\(\mathbf{F}_2\) are the forces on the two atoms resulting from the
pairwise interaction.  The second term is a bond contribution of
similar form for the \(N_b\) bonds which atom \(I\) is part
of.  There are similar terms for the \(N_a\) angle, \(N_d\)
dihedral, and \(N_i\) improper interactions atom \(I\) is part
of.  There is also a term for the KSpace contribution from long\sphinxhyphen{}range
Coulombic interactions, if defined.  Finally, there is a term for the
\(N_f\) {\hyperref[\detokenize{fix::doc}]{\sphinxcrossref{\DUrole{doc}{fixes}}}} that apply internal constraint forces
to atom \(I\). Currently, only the {\hyperref[\detokenize{fix_shake::doc}]{\sphinxcrossref{\DUrole{doc}{fix shake}}}}
and {\hyperref[\detokenize{fix_rigid::doc}]{\sphinxcrossref{\DUrole{doc}{fix rigid}}}} commands contribute to this term.  As
the coefficients in the formula imply, a virial contribution produced
by a small set of atoms (e.g. 4 atoms in a dihedral or 3 atoms in a
Tersoff 3\sphinxhyphen{}body interaction) is assigned in equal portions to each atom
in the set.  E.g. 1/4 of the dihedral virial to each of the 4 atoms,
or 1/3 of the fix virial due to SHAKE constraints applied to atoms in
a water molecule via the {\hyperref[\detokenize{fix_shake::doc}]{\sphinxcrossref{\DUrole{doc}{fix shake}}}} command.
As an exception, the virial contribution from
constraint forces in {\hyperref[\detokenize{fix_rigid::doc}]{\sphinxcrossref{\DUrole{doc}{fix rigid}}}} on each atom
is computed from the constraint force acting on the corresponding atom
and its position, i.e. the total virial is not equally distributed.

\sphinxAtStartPar
In case of compute \sphinxstyleemphasis{centroid/stress/atom}, the virial contribution is:
\begin{equation*}
\begin{split} W_{ab} & = \sum_{n = 1}^{N_p} r_{I0_a} F_{I_b} + \sum_{n = 1}^{N_b} r_{I0_a} F_{I_b} + \sum_{n = 1}^{N_a} r_{I0_a}  F_{I_b} + \sum_{n = 1}^{N_d} r_{I0_a} F_{I_b} + \sum_{n = 1}^{N_i} r_{I0_a} F_{I_b} \\
& + {\rm Kspace}(r_{i_a},F_{i_b}) + \sum_{n = 1}^{N_f} r_{i_a} F_{i_b}\end{split}
\end{equation*}
\sphinxAtStartPar
As with compute \sphinxstyleemphasis{stress/atom}, the first, second, third, fourth and
fifth terms are pairwise, bond, angle, dihedral and improper
contributions, but instead of assigning the virial contribution
equally to each atom, only the force \(\mathbf{F}_I\) acting on
atom \(I\) due to the interaction and the relative position
\(\mathbf{r}_{I0}\) of the atom \(I\) to the geometric center
of the interacting atoms, i.e. centroid, is used.  As the geometric
center is different for each interaction, the \(\mathbf{r}_{I0}\)
also differs. The sixth term, Kspace contribution,
is computed identically to compute \sphinxstyleemphasis{stress/atom}.
The seventh term is handed differently depending on
if the constraint forces are due to {\hyperref[\detokenize{fix_shake::doc}]{\sphinxcrossref{\DUrole{doc}{fix shake}}}}
or {\hyperref[\detokenize{fix_rigid::doc}]{\sphinxcrossref{\DUrole{doc}{fix rigid}}}}.
In case of SHAKE constraints, each distance constraint is
handed as a pairwise interaction.
E.g. in case of a water molecule, two OH and one HH distance
constraints are treated as three pairwise interactions.
In case of {\hyperref[\detokenize{fix_rigid::doc}]{\sphinxcrossref{\DUrole{doc}{fix rigid}}}},
all constraint forces in the molecule are treated
as a single many\sphinxhyphen{}body interaction with a single centroid position.
In case of water molecule, the formula expression would become
identical to that of the three\sphinxhyphen{}body angle interaction.
Although the total system virial is the same as
compute \sphinxstyleemphasis{stress/atom}, compute \sphinxstyleemphasis{centroid/stress/atom} is know to
result in more consistent heat flux values for angle, dihedrals,
improper and constraint force contributions
when computed via {\hyperref[\detokenize{compute_heat_flux::doc}]{\sphinxcrossref{\DUrole{doc}{compute heat/flux}}}}.

\sphinxAtStartPar
If no extra keywords are listed, the kinetic contribution \sphinxstyleemphasis{and} all
of the virial contribution terms are included in the per\sphinxhyphen{}atom stress
tensor.  If any extra keywords are listed, only those terms are
summed to compute the tensor.  The \sphinxstyleemphasis{virial} keyword means include all
terms except the kinetic energy \sphinxstyleemphasis{ke}.

\sphinxAtStartPar
Note that the stress for each atom is due to its interaction with all
other atoms in the simulation, not just with other atoms in the group.

\sphinxAtStartPar
Details of how compute \sphinxstyleemphasis{stress/atom} obtains the virial for individual
atoms for either pairwise or many\sphinxhyphen{}body potentials, and including the
effects of periodic boundary conditions is discussed in
{\hyperref[\detokenize{compute_stress_atom:thompson2}]{\sphinxcrossref{\DUrole{std,std-ref}{(Thompson)}}}}.  The basic idea for many\sphinxhyphen{}body
potentials is to treat each component of the force computation between
a small cluster of atoms in the same manner as in the formula above
for bond, angle, dihedral, etc interactions.  Namely the quantity
\(\mathbf{r} \cdot \mathbf{F}\) is summed over the atoms in the
interaction, with the \(r\) vectors unwrapped by periodic
boundaries so that the cluster of atoms is close together.  The total
contribution for the cluster interaction is divided evenly among those
atoms.

\sphinxAtStartPar
Details of how compute \sphinxstyleemphasis{centroid/stress/atom} obtains the virial for
individual atoms are given in {\hyperref[\detokenize{compute_stress_atom:surblys1}]{\sphinxcrossref{\DUrole{std,std-ref}{(Surblys2019)}}}} and
{\hyperref[\detokenize{compute_stress_atom:surblys2}]{\sphinxcrossref{\DUrole{std,std-ref}{(Surblys2021)}}}}, where the
idea is that the virial of the atom \(I\) is the result of only
the force \(\mathbf{F}_I\) on the atom due to the interaction and
its positional vector \(\mathbf{r}_{I0}\), relative to the
geometric center of the interacting atoms, regardless of the number of
participating atoms.  The periodic boundary treatment is identical to
that of compute \sphinxstyleemphasis{stress/atom}, and both of them reduce to identical
expressions for two\sphinxhyphen{}body interactions, i.e. computed values for
contributions from bonds and two\sphinxhyphen{}body pair styles, such as
{\hyperref[\detokenize{pair_lj::doc}]{\sphinxcrossref{\DUrole{doc}{Lennard\sphinxhyphen{}Jones}}}}, will be the same, while contributions
from angles, dihedrals and impropers will be different.

\sphinxAtStartPar
The {\hyperref[\detokenize{dihedral_charmm::doc}]{\sphinxcrossref{\DUrole{doc}{dihedral\_style charmm}}}} style calculates
pairwise interactions between 1\sphinxhyphen{}4 atoms.  The virial contribution of
these terms is included in the pair virial, not the dihedral virial.

\sphinxAtStartPar
The KSpace contribution is calculated using the method in
{\hyperref[\detokenize{compute_stress_atom:heyes2}]{\sphinxcrossref{\DUrole{std,std-ref}{(Heyes)}}}} for the Ewald method and by the methodology
described in {\hyperref[\detokenize{compute_stress_atom:sirk1}]{\sphinxcrossref{\DUrole{std,std-ref}{(Sirk)}}}} for PPPM.  The choice of KSpace
solver is specified by the {\hyperref[\detokenize{kspace_style::doc}]{\sphinxcrossref{\DUrole{doc}{kspace\_style pppm}}}}
command.  Note that for PPPM, the calculation requires 6 extra FFTs
each timestep that per\sphinxhyphen{}atom stress is calculated.  Thus it can
significantly increase the cost of the PPPM calculation if it is
needed on a large fraction of the simulation timesteps.

\sphinxAtStartPar
The \sphinxstyleemphasis{temp\sphinxhyphen{}ID} argument can be used to affect the per\sphinxhyphen{}atom velocities
used in the kinetic energy contribution to the total stress.  If the
kinetic energy is not included in the stress, than the temperature
compute is not used and can be specified as NULL.  If the kinetic
energy is included and you wish to use atom velocities as\sphinxhyphen{}is, then
\sphinxstyleemphasis{temp\sphinxhyphen{}ID} can also be specified as NULL.  If desired, the specified
temperature compute can be one that subtracts off a bias to leave each
atom with only a thermal velocity to use in the formula above, e.g. by
subtracting a background streaming velocity.  See the doc pages for
individual {\hyperref[\detokenize{compute::doc}]{\sphinxcrossref{\DUrole{doc}{compute commands}}}} to determine which ones
include a bias.


\bigskip\hrule\bigskip


\sphinxAtStartPar
Note that as defined in the formula, per\sphinxhyphen{}atom stress is the negative
of the per\sphinxhyphen{}atom pressure tensor.  It is also really a stress*volume
formulation, meaning the computed quantity is in units of
pressure*volume.  It would need to be divided by a per\sphinxhyphen{}atom volume to
have units of stress (pressure), but an individual atom’s volume is
not well defined or easy to compute in a deformed solid or a liquid.
See the {\hyperref[\detokenize{compute_voronoi_atom::doc}]{\sphinxcrossref{\DUrole{doc}{compute voronoi/atom}}}} command for
one possible way to estimate a per\sphinxhyphen{}atom volume.

\sphinxAtStartPar
Thus, if the diagonal components of the per\sphinxhyphen{}atom stress tensor are
summed for all atoms in the system and the sum is divided by
\(dV\), where \(d\) = dimension and \(V\) is the volume of
the system, the result should be \(-P\), where \(P\) is the
total pressure of the system.

\sphinxAtStartPar
These lines in an input script for a 3d system should yield that
result. I.e. the last 2 columns of thermo output will be the same:

\begin{sphinxVerbatim}[commandchars=\\\{\}]
\PYG{k}{compute        }\PYG{n+nv+nvIdentifier}{peratom}\PYG{+w}{ }\PYG{n+nv+nvIdentifier}{all}\PYG{+w}{ }\PYG{n}{stress}\PYG{o}{/}\PYG{n}{atom}\PYG{+w}{ }\PYG{n}{NULL}
\PYG{k}{compute        }\PYG{n+nv+nvIdentifier}{p}\PYG{+w}{ }\PYG{n+nv+nvIdentifier}{all}\PYG{+w}{ }\PYG{n}{reduce}\PYG{+w}{ }\PYG{n}{sum}\PYG{+w}{ }\PYG{n}{c\PYGZus{}peratom[1]}\PYG{+w}{ }\PYG{n}{c\PYGZus{}peratom[2]}\PYG{+w}{ }\PYG{n}{c\PYGZus{}peratom[3]}
\PYG{k}{variable       }\PYG{n+nv+nvIdentifier}{press}\PYG{+w}{ }\PYG{n}{equal}\PYG{+w}{ }\PYG{o}{\PYGZhy{}}\PYG{n+nv}{(}\PYG{n+nv}{c\PYGZus{}p[1]+c\PYGZus{}p[2]+c\PYGZus{}p[3]}\PYG{n+nv}{)}\PYG{o}{/}\PYG{n+nv}{(}\PYG{n+nv}{3*vol}\PYG{n+nv}{)}
\PYG{k}{thermo\PYGZus{}style}\PYG{+w}{   }\PYG{n}{custom}\PYG{+w}{ }\PYG{n}{step}\PYG{+w}{ }\PYG{n}{temp}\PYG{+w}{ }\PYG{n}{etotal}\PYG{+w}{ }\PYG{n}{press}\PYG{+w}{ }\PYG{n}{v\PYGZus{}press}
\end{sphinxVerbatim}

\begin{sphinxadmonition}{note}{Note:}
\sphinxAtStartPar
The per\sphinxhyphen{}atom stress does not include any Lennard\sphinxhyphen{}Jones tail
corrections to the pressure added by the {\hyperref[\detokenize{pair_modify::doc}]{\sphinxcrossref{\DUrole{doc}{pair\_modify tail yes}}}} command, since those are contributions to the global
system pressure.
\end{sphinxadmonition}

\sphinxAtStartPar
The compute stress/atom can be used in a number of ways.  Here is an
example to compute a 1\sphinxhyphen{}d pressure profile in x\sphinxhyphen{}direction across the
complete simulation box. You will need to adjust the number of bins and the
selections for time averaging to your specific simulation.  This assumes
that the dimensions of the simulation cell does not change.

\begin{sphinxVerbatim}[commandchars=\\\{\}]
\PYG{c}{\PYGZsh{} set number of bins}
\PYG{k}{variable }\PYG{n+nv+nvIdentifier}{nbins}\PYG{+w}{ }\PYG{n}{index}\PYG{+w}{ }\PYG{l+m}{20}
\PYG{k}{variable }\PYG{n+nv+nvIdentifier}{fraction}\PYG{+w}{ }\PYG{n}{equal}\PYG{+w}{ }\PYG{l+m}{1.0}\PYG{o}{/}\PYG{n}{v\PYGZus{}nbins}
\PYG{c}{\PYGZsh{} define bins as chunks}
\PYG{k}{compute }\PYG{n+nv+nvIdentifier}{cchunk}\PYG{+w}{ }\PYG{n+nv+nvIdentifier}{all}\PYG{+w}{ }\PYG{n}{chunk}\PYG{o}{/}\PYG{n}{atom}\PYG{+w}{ }\PYG{n}{bin}\PYG{o}{/}\PYG{l+m}{1}\PYG{n}{d}\PYG{+w}{ }\PYG{n}{x}\PYG{+w}{ }\PYG{n}{lower}\PYG{+w}{ }\PYG{n+nv}{\PYGZdl{}\PYGZob{}}\PYG{n+nv}{fraction}\PYG{n+nv}{\PYGZcb{}}\PYG{+w}{ }\PYG{n}{units}\PYG{+w}{ }\PYG{n}{reduced}
\PYG{k}{compute }\PYG{n+nv+nvIdentifier}{stress}\PYG{+w}{ }\PYG{n+nv+nvIdentifier}{all}\PYG{+w}{ }\PYG{n}{stress}\PYG{o}{/}\PYG{n}{atom}\PYG{+w}{ }\PYG{n}{NULL}
\PYG{c}{\PYGZsh{} apply conversion to pressure early since we have no variable style for processing chunks}
\PYG{k}{variable }\PYG{n+nv+nvIdentifier}{press}\PYG{+w}{ }\PYG{n}{atom}\PYG{+w}{ }\PYG{o}{\PYGZhy{}}\PYG{n+nv}{(}\PYG{n+nv}{c\PYGZus{}stress[1]+c\PYGZus{}stress[2]+c\PYGZus{}stress[3]}\PYG{n+nv}{)}\PYG{o}{/}\PYG{n+nv}{(}\PYG{n+nv}{3.0*vol*\PYGZdl{}\PYGZob{}fraction\PYGZcb{}}\PYG{n+nv}{)}
\PYG{k}{compute }\PYG{n+nv+nvIdentifier}{binpress}\PYG{+w}{ }\PYG{n+nv+nvIdentifier}{all}\PYG{+w}{ }\PYG{n}{reduce}\PYG{o}{/}\PYG{n}{chunk}\PYG{+w}{ }\PYG{n}{cchunk}\PYG{+w}{ }\PYG{n}{sum}\PYG{+w}{ }\PYG{n}{v\PYGZus{}press}
\PYG{k}{fix }\PYG{n+nv+nvIdentifier}{avg}\PYG{+w}{ }\PYG{n+nv+nvIdentifier}{all}\PYG{+w}{ }\PYG{n}{ave}\PYG{o}{/}\PYG{n}{time}\PYG{+w}{ }\PYG{l+m}{10}\PYG{+w}{ }\PYG{l+m}{40}\PYG{+w}{ }\PYG{l+m}{400}\PYG{+w}{ }\PYG{n}{c\PYGZus{}binpress}\PYG{+w}{ }\PYG{n}{mode}\PYG{+w}{ }\PYG{n}{vector}\PYG{+w}{ }\PYG{n}{file}\PYG{+w}{ }\PYG{n}{ave\PYGZus{}stress.txt}
\end{sphinxVerbatim}


\subsection{Output info}
\label{\detokenize{compute_stress_atom:output-info}}
\sphinxAtStartPar
Compute \sphinxstyleemphasis{stress/atom} calculates a per\sphinxhyphen{}atom array with 6 columns,
which can be accessed by indices 1\sphinxhyphen{}6 by any command that uses per\sphinxhyphen{}atom
values from a compute as input.  Compute \sphinxstyleemphasis{centroid/stress/atom}
produces a per\sphinxhyphen{}atom array with 9 columns, but otherwise can be used in
an identical manner to compute \sphinxstyleemphasis{stress/atom}.  See the {\hyperref[\detokenize{Howto_output::doc}]{\sphinxcrossref{\DUrole{doc}{Howto
output}}}} page for an overview of LAMMPS output
options.

\sphinxAtStartPar
The ordering of the 6 columns for \sphinxstyleemphasis{stress/atom} is as follows: xx, yy,
zz, xy, xz, yz.  The ordering of the 9 columns for
\sphinxstyleemphasis{centroid/stress/atom} is as follows: xx, yy, zz, xy, xz, yz, yx, zx,
zy.

\sphinxAtStartPar
The per\sphinxhyphen{}atom array values will be in pressure*volume {\hyperref[\detokenize{units::doc}]{\sphinxcrossref{\DUrole{doc}{units}}}} as discussed above.


\subsection{Restrictions}
\label{\detokenize{compute_stress_atom:restrictions}}
\sphinxAtStartPar
Currently, compute \sphinxstyleemphasis{centroid/stress/atom} does not support pair styles
with many\sphinxhyphen{}body interactions ({\hyperref[\detokenize{pair_eam::doc}]{\sphinxcrossref{\DUrole{doc}{EAM}}}} is an exception,
since its computations are performed pairwise), nor granular pair
styles with pairwise forces which are not aligned with the vector
between the pair of particles.  All bond styles are supported.  All
angle, dihedral, improper styles are supported with the exception of
INTEL and KOKKOS variants of specific styles.  It also does not
support models with long\sphinxhyphen{}range Coulombic or dispersion forces,
i.e. the kspace\_style command in LAMMPS.  It also does not implement the
following fixes which add rigid\sphinxhyphen{}body constraints:
{\hyperref[\detokenize{fix_rigid::doc}]{\sphinxcrossref{\DUrole{doc}{fix rigid/*}}}} and the OpenMP accelerated version of {\hyperref[\detokenize{fix_rigid::doc}]{\sphinxcrossref{\DUrole{doc}{fix rigid/small}}}},
while all other {\hyperref[\detokenize{fix_rigid::doc}]{\sphinxcrossref{\DUrole{doc}{fix rigid/*/small}}}} are implemented.

\sphinxAtStartPar
LAMMPS will generate an error if one of these options is included in
your model.  Extension of centroid stress calculations to these force
and fix styles is planned for the future.


\subsection{Related commands}
\label{\detokenize{compute_stress_atom:related-commands}}
\sphinxAtStartPar
{\hyperref[\detokenize{compute_pe::doc}]{\sphinxcrossref{\DUrole{doc}{compute pe}}}}, {\hyperref[\detokenize{compute_pressure::doc}]{\sphinxcrossref{\DUrole{doc}{compute pressure}}}}


\subsection{Default}
\label{\detokenize{compute_stress_atom:default}}
\sphinxAtStartPar
By default the compute includes contributions from the keywords:
\sphinxcode{\sphinxupquote{ke pair bond angle dihedral improper kspace fix}}


\bigskip\hrule\bigskip

\phantomsection\label{\detokenize{compute_stress_atom:heyes2}}
\sphinxAtStartPar
\sphinxstylestrong{(Heyes)} Heyes, Phys Rev B, 49, 755 (1994).

\phantomsection\label{\detokenize{compute_stress_atom:sirk1}}
\sphinxAtStartPar
\sphinxstylestrong{(Sirk)} Sirk, Moore, Brown, J Chem Phys, 138, 064505 (2013).

\phantomsection\label{\detokenize{compute_stress_atom:thompson2}}
\sphinxAtStartPar
\sphinxstylestrong{(Thompson)} Thompson, Plimpton, Mattson, J Chem Phys, 131, 154107 (2009).

\phantomsection\label{\detokenize{compute_stress_atom:surblys1}}
\sphinxAtStartPar
\sphinxstylestrong{(Surblys2019)} Surblys, Matsubara, Kikugawa, Ohara, Phys Rev E, 99, 051301(R) (2019).

\phantomsection\label{\detokenize{compute_stress_atom:surblys2}}
\sphinxAtStartPar
\sphinxstylestrong{(Surblys2021)} Surblys, Matsubara, Kikugawa, Ohara, J Appl Phys 130, 215104 (2021).

\sphinxstepscope

\index{compute stress/cartesian@\spxentry{compute stress/cartesian}}\ignorespaces 

\section{compute stress/cartesian command}
\label{\detokenize{compute_stress_cartesian:compute-stress-cartesian-command}}\label{\detokenize{compute_stress_cartesian:index-0}}\label{\detokenize{compute_stress_cartesian::doc}}

\subsection{Syntax}
\label{\detokenize{compute_stress_cartesian:syntax}}
\begin{sphinxVerbatim}[commandchars=\\\{\}]
\PYG{k}{compute }\PYG{n+nv+nvIdentifier}{ID}\PYG{+w}{ }\PYG{n+nv+nvIdentifier}{group\PYGZhy{}ID}\PYG{+w}{ }\PYG{n}{stress}\PYG{o}{/}\PYG{n}{cartesian}\PYG{+w}{ }\PYG{n}{args}
\end{sphinxVerbatim}
\begin{itemize}
\item {} 
\sphinxAtStartPar
ID, group\sphinxhyphen{}ID are documented in {\hyperref[\detokenize{compute::doc}]{\sphinxcrossref{\DUrole{doc}{compute}}}} command

\item {} 
\sphinxAtStartPar
args = argument specific to the compute style

\end{itemize}
\begin{sphinxalltt}
\sphinxstyleemphasis{stress/cartesian} args = dim1 bin\_width1 dim2 bin\_width2 keyword
  dim1 = \sphinxstyleemphasis{x} or \sphinxstyleemphasis{y} or \sphinxstyleemphasis{z}
  bin\_width1 = width of the bin
  dim2 = \sphinxstyleemphasis{x} or \sphinxstyleemphasis{y} or \sphinxstyleemphasis{z} or \sphinxstyleemphasis{NULL}
  bin\_width2 = width of the bin
  keyword = \sphinxstyleemphasis{ke} or \sphinxstyleemphasis{pair} or \sphinxstyleemphasis{bond}
\end{sphinxalltt}


\subsection{Examples}
\label{\detokenize{compute_stress_cartesian:examples}}
\begin{sphinxVerbatim}[commandchars=\\\{\}]
\PYG{k}{compute }\PYG{n+nv+nvIdentifier}{1}\PYG{+w}{ }\PYG{n+nv+nvIdentifier}{all}\PYG{+w}{ }\PYG{n}{stress}\PYG{o}{/}\PYG{n}{cartesian}\PYG{+w}{ }\PYG{n}{x}\PYG{+w}{ }\PYG{l+m}{0.1}\PYG{+w}{ }\PYG{n}{NULL}\PYG{+w}{ }\PYG{l+m}{0}
\PYG{k}{compute }\PYG{n+nv+nvIdentifier}{1}\PYG{+w}{ }\PYG{n+nv+nvIdentifier}{all}\PYG{+w}{ }\PYG{n}{stress}\PYG{o}{/}\PYG{n}{cartesian}\PYG{+w}{ }\PYG{n}{y}\PYG{+w}{ }\PYG{l+m}{0.1}\PYG{+w}{ }\PYG{n}{z}\PYG{+w}{ }\PYG{l+m}{0.1}
\PYG{k}{compute }\PYG{n+nv+nvIdentifier}{1}\PYG{+w}{ }\PYG{n+nv+nvIdentifier}{all}\PYG{+w}{ }\PYG{n}{stress}\PYG{o}{/}\PYG{n}{cartesian}\PYG{+w}{ }\PYG{n}{x}\PYG{+w}{ }\PYG{l+m}{0.1}\PYG{+w}{ }\PYG{n}{NULL}\PYG{+w}{ }\PYG{l+m}{0}\PYG{+w}{ }\PYG{n}{ke}\PYG{+w}{ }\PYG{n}{pair}
\end{sphinxVerbatim}


\subsection{Description}
\label{\detokenize{compute_stress_cartesian:description}}
\sphinxAtStartPar
Compute \sphinxstyleemphasis{stress/cartesian} defines computations that calculate profiles of the
diagonal components of the local stress tensor over one or two Cartesian
dimensions, as described in {\hyperref[\detokenize{compute_stress_cartesian:ikeshoji2}]{\sphinxcrossref{\DUrole{std,std-ref}{(Ikeshoji)}}}}. The stress tensor is
split into a kinetic contribution \(P^k\) and a virial contribution
\(P^v\). The sum gives the total stress tensor \(P = P^k+P^v\).
This compute obeys momentum balance through fluid interfaces. They use the
Irving\textendash{}Kirkwood contour, which is the straight line between particle pairs.

\sphinxAtStartPar
\DUrole{versionmodified,added}{Added in version 15Jun2023: }Added support for bond styles

\sphinxAtStartPar
This compute only supports pair and bond (no angle, dihedral, improper,
or kspace) forces. By default, if no extra keywords are specified, all
supported contributions to the stress are included (ke, pair, bond). If any
keywords are specified, then only those components are summed.


\subsection{Output info}
\label{\detokenize{compute_stress_cartesian:output-info}}
\sphinxAtStartPar
The output columns for \sphinxstyleemphasis{stress/cartesian} are the position of the
center of the local volume in the first and second dimensions, number
density, \(P^k_{xx}\), \(P^k_{yy}\), \(P^k_{zz}\),
\(P^v_{xx}\), \(P^v_{yy}\), and \(P^v_{zz}\). There are 8
columns when one dimension is specified and 9 columns when two
dimensions are specified. The number of bins (rows) is
\((L_1/b_1)(L_2/b_2)\), where \(L_1\) and \(L_2\) are the lengths
of the simulation box in the specified dimensions and \(b_1\) and
\(b_2\) are the specified bin widths. When only one dimension is
specified, the number of bins (rows) is \(L_1/b_1\).

\sphinxAtStartPar
This array can be output with {\hyperref[\detokenize{fix_ave_time::doc}]{\sphinxcrossref{\DUrole{doc}{fix ave/time}}}},

\begin{sphinxVerbatim}[commandchars=\\\{\}]
\PYG{k}{compute }\PYG{n+nv+nvIdentifier}{p}\PYG{+w}{ }\PYG{n+nv+nvIdentifier}{all}\PYG{+w}{ }\PYG{n}{stress}\PYG{o}{/}\PYG{n}{cartesian}\PYG{+w}{ }\PYG{n}{x}\PYG{+w}{ }\PYG{l+m}{0.1}
\PYG{k}{fix }\PYG{n+nv+nvIdentifier}{2}\PYG{+w}{ }\PYG{n+nv+nvIdentifier}{all}\PYG{+w}{ }\PYG{n}{ave}\PYG{o}{/}\PYG{n}{time}\PYG{+w}{ }\PYG{l+m}{100}\PYG{+w}{ }\PYG{l+m}{1}\PYG{+w}{ }\PYG{l+m}{100}\PYG{+w}{ }\PYG{n}{c\PYGZus{}p[}\PYG{o}{*}\PYG{n}{]}\PYG{+w}{ }\PYG{n}{file}\PYG{+w}{ }\PYG{n}{dump\PYGZus{}p.out}\PYG{+w}{ }\PYG{n}{mode}\PYG{+w}{ }\PYG{n}{vector}
\end{sphinxVerbatim}

\sphinxAtStartPar
The values calculated by this compute are “intensive”.  The stress
values will be in pressure {\hyperref[\detokenize{units::doc}]{\sphinxcrossref{\DUrole{doc}{units}}}}. The number density
values are in inverse volume {\hyperref[\detokenize{units::doc}]{\sphinxcrossref{\DUrole{doc}{units}}}}.

\sphinxAtStartPar
NOTE 1: The local stress does not include any Lennard\sphinxhyphen{}Jones tail
corrections to the stress added by the
{\hyperref[\detokenize{pair_modify::doc}]{\sphinxcrossref{\DUrole{doc}{pair\_modify tail yes}}}}
command, since those are contributions to the global system pressure.

\sphinxAtStartPar
NOTE 2: The local stress profiles generated by these computes are
similar to those obtained by the
{\hyperref[\detokenize{compute_stress_mop::doc}]{\sphinxcrossref{\DUrole{doc}{method\sphinxhyphen{}of\sphinxhyphen{}planes (MOP)}}}}.
A key difference is that compute
{\hyperref[\detokenize{compute_stress_mop::doc}]{\sphinxcrossref{\DUrole{doc}{stress/mop/profile}}}}
considers particles crossing a set of planes, while
\sphinxstyleemphasis{stress/cartesian} computes averages for a set of small volumes.
Moreover, this compute computes the diagonal components of the stress
tensor \(P_{xx}\), \(P_{yy}\), and \(P_{zz}\), while
\sphinxstyleemphasis{stress/mop/profile} computes the components
\(P_{ix}\), \(P_{iy}\), and \(P_{iz}\), where \(i\) is the
direction normal to the plane.

\sphinxAtStartPar
More information on the similarities and differences can be found in
{\hyperref[\detokenize{compute_stress_cartesian:ikeshoji2}]{\sphinxcrossref{\DUrole{std,std-ref}{(Ikeshoji)}}}}.


\subsection{Restrictions}
\label{\detokenize{compute_stress_cartesian:restrictions}}
\sphinxAtStartPar
These computes calculate the stress tensor contributions for pair and bond
forces only (no angle, dihedral, improper, or kspace force).
It requires pairwise force calculations not available for most
many\sphinxhyphen{}body pair styles.

\sphinxAtStartPar
These computes are part of the EXTRA\sphinxhyphen{}COMPUTE package.  They are only
enabled if LAMMPS was built with that package.  See the {\hyperref[\detokenize{Build_package::doc}]{\sphinxcrossref{\DUrole{doc}{Build
package}}}} doc page for more info.


\subsection{Related commands}
\label{\detokenize{compute_stress_cartesian:related-commands}}
\sphinxAtStartPar
{\hyperref[\detokenize{compute_stress_atom::doc}]{\sphinxcrossref{\DUrole{doc}{compute stress/atom}}}}, {\hyperref[\detokenize{compute_pressure::doc}]{\sphinxcrossref{\DUrole{doc}{compute pressure}}}},
{\hyperref[\detokenize{compute_stress_mop::doc}]{\sphinxcrossref{\DUrole{doc}{compute stress/mop/profile}}}}, {\hyperref[\detokenize{compute_stress_curvilinear::doc}]{\sphinxcrossref{\DUrole{doc}{compute stress/spherical}}}},
{\hyperref[\detokenize{compute_stress_curvilinear::doc}]{\sphinxcrossref{\DUrole{doc}{compute stress/cylinder}}}}


\bigskip\hrule\bigskip

\phantomsection\label{\detokenize{compute_stress_cartesian:ikeshoji2}}
\sphinxAtStartPar
\sphinxstylestrong{(Ikeshoji)} Ikeshoji, Hafskjold, Furuholt, Mol Sim, 29, 101\sphinxhyphen{}109, (2003).

\sphinxstepscope

\index{compute stress/cylinder@\spxentry{compute stress/cylinder}}\ignorespaces 
\index{compute stress/spherical@\spxentry{compute stress/spherical}}\ignorespaces 

\section{compute stress/cylinder command}
\label{\detokenize{compute_stress_curvilinear:compute-stress-cylinder-command}}\label{\detokenize{compute_stress_curvilinear:index-1}}\label{\detokenize{compute_stress_curvilinear:index-0}}\label{\detokenize{compute_stress_curvilinear::doc}}

\section{compute stress/spherical command}
\label{\detokenize{compute_stress_curvilinear:compute-stress-spherical-command}}

\subsection{Syntax}
\label{\detokenize{compute_stress_curvilinear:syntax}}
\begin{sphinxVerbatim}[commandchars=\\\{\}]
\PYG{k}{compute }\PYG{n+nv+nvIdentifier}{ID}\PYG{+w}{ }\PYG{n+nv+nvIdentifier}{group\PYGZhy{}ID}\PYG{+w}{ }\PYG{n}{style}\PYG{+w}{ }\PYG{n}{args}
\end{sphinxVerbatim}
\begin{itemize}
\item {} 
\sphinxAtStartPar
ID, group\sphinxhyphen{}ID are documented in {\hyperref[\detokenize{compute::doc}]{\sphinxcrossref{\DUrole{doc}{compute}}}} command

\item {} 
\sphinxAtStartPar
style = stress/spherical or stress/cylinder

\item {} 
\sphinxAtStartPar
args = argument specific to the compute style

\end{itemize}
\begin{sphinxalltt}
\sphinxstyleemphasis{stress/cylinder} args = zlo zh Rmax bin\_width keyword
  zlo = minimum z\sphinxhyphen{}boundary for cylinder
  zhi = maximum z\sphinxhyphen{}boundary for cylinder
  Rmax = maximum radius to perform calculation to
  bin\_width = width of radial bins to use for calculation
  keyword = ke (zero or one can be specified)
    ke = yes or no
\sphinxstyleemphasis{stress/spherical}
  x0, y0, z0 = origin of the spherical coordinate system
  bin\_width = width of spherical shells
  Rmax = maximum radius of spherical shells
\end{sphinxalltt}


\subsection{Examples}
\label{\detokenize{compute_stress_curvilinear:examples}}
\begin{sphinxVerbatim}[commandchars=\\\{\}]
\PYG{k}{compute }\PYG{n+nv+nvIdentifier}{1}\PYG{+w}{ }\PYG{n+nv+nvIdentifier}{all}\PYG{+w}{ }\PYG{n}{stress}\PYG{o}{/}\PYG{n}{cylinder}\PYG{+w}{ }\PYG{o}{\PYGZhy{}}\PYG{l+m}{10.0}\PYG{+w}{ }\PYG{l+m}{10.0}\PYG{+w}{ }\PYG{l+m}{15.0}\PYG{+w}{ }\PYG{l+m}{0.25}
\PYG{k}{compute }\PYG{n+nv+nvIdentifier}{1}\PYG{+w}{ }\PYG{n+nv+nvIdentifier}{all}\PYG{+w}{ }\PYG{n}{stress}\PYG{o}{/}\PYG{n}{cylinder}\PYG{+w}{ }\PYG{o}{\PYGZhy{}}\PYG{l+m}{10.0}\PYG{+w}{ }\PYG{l+m}{10.0}\PYG{+w}{ }\PYG{l+m}{15.0}\PYG{+w}{ }\PYG{l+m}{0.25}\PYG{+w}{ }\PYG{n}{ke}\PYG{+w}{ }\PYG{n}{no}
\PYG{k}{compute }\PYG{n+nv+nvIdentifier}{1}\PYG{+w}{ }\PYG{n+nv+nvIdentifier}{all}\PYG{+w}{ }\PYG{n}{stress}\PYG{o}{/}\PYG{n}{spherical}\PYG{+w}{ }\PYG{l+m}{0}\PYG{+w}{ }\PYG{l+m}{0}\PYG{+w}{ }\PYG{l+m}{0}\PYG{+w}{ }\PYG{l+m}{0.1}\PYG{+w}{ }\PYG{l+m}{10}
\end{sphinxVerbatim}


\subsection{Description}
\label{\detokenize{compute_stress_curvilinear:description}}
\sphinxAtStartPar
Compute \sphinxstyleemphasis{stress/cylinder}, and compute
\sphinxstyleemphasis{stress/spherical} define computations that calculate profiles of the
diagonal components of the local stress tensor in the specified
coordinate system. The stress tensor is split into a kinetic
contribution \(P^k\) and a virial contribution \(P^v\). The sum
gives the total stress tensor \(P = P^k+P^v\). These computes can
for example be used to calculate the diagonal components of the local
stress tensor of surfaces with cylindrical or spherical
symmetry. These computes obeys momentum balance through fluid
interfaces. They use the Irving\textendash{}Kirkwood contour, which is the straight
line between particle pairs.

\sphinxAtStartPar
The compute \sphinxstyleemphasis{stress/cylinder} computes the stress profile along the
radial direction in cylindrical coordinates, as described in
{\hyperref[\detokenize{compute_stress_curvilinear:addington1}]{\sphinxcrossref{\DUrole{std,std-ref}{(Addington)}}}}. The compute \sphinxstyleemphasis{stress/spherical}
computes the stress profile along the radial direction in spherical
coordinates, as described in {\hyperref[\detokenize{compute_stress_curvilinear:ikeshoji4}]{\sphinxcrossref{\DUrole{std,std-ref}{(Ikeshoji)}}}}.


\subsection{Output info}
\label{\detokenize{compute_stress_curvilinear:output-info}}
\sphinxAtStartPar
The default output columns for \sphinxstyleemphasis{stress/cylinder} are the radius to the
center of the cylindrical shell, number density, \(P^k_{rr}\),
\(P^k_{\phi\phi}\), \(P^k_{zz}\), \(P^v_{rr}\),
\(P^v_{\phi\phi}\), and \(P^v_{zz}\). When the keyword \sphinxstyleemphasis{ke} is
set to \sphinxstyleemphasis{no}, the kinetic contributions are not calculated, and
consequently there are only 5 columns: the position of the center of the
cylindrical shell, the number density, \(P^v_{rr}\),
\(P^v_{\phi\phi}\), and \(P^v_{zz}\). The number of bins (rows) is
\(R_\text{max}/b\), where \(b\) is the specified bin width.

\sphinxAtStartPar
The output columns for \sphinxstyleemphasis{stress/spherical} are the position of the center
of the spherical shell, the number density, \(P^k_{rr}\),
\(P^k_{\theta\theta}\), \(P^k_{\phi\phi}\), \(P^v_{rr}\),
\(P^v_{\theta\theta}\), and \(P^v_{\phi\phi}\). There are 8
columns and the number of bins (rows) is \(R_\text{max}/b\), where
\(b\) is the specified bin width.

\sphinxAtStartPar
This array can be output with {\hyperref[\detokenize{fix_ave_time::doc}]{\sphinxcrossref{\DUrole{doc}{fix ave/time}}}},

\begin{sphinxVerbatim}[commandchars=\\\{\}]
\PYG{k}{compute }\PYG{n+nv+nvIdentifier}{1}\PYG{+w}{ }\PYG{n+nv+nvIdentifier}{all}\PYG{+w}{ }\PYG{n}{stress}\PYG{o}{/}\PYG{n}{spherical}\PYG{+w}{ }\PYG{l+m}{0}\PYG{+w}{ }\PYG{l+m}{0}\PYG{+w}{ }\PYG{l+m}{0}\PYG{+w}{ }\PYG{l+m}{0.1}\PYG{+w}{ }\PYG{l+m}{10}
\PYG{k}{fix }\PYG{n+nv+nvIdentifier}{2}\PYG{+w}{ }\PYG{n+nv+nvIdentifier}{all}\PYG{+w}{ }\PYG{n}{ave}\PYG{o}{/}\PYG{n}{time}\PYG{+w}{ }\PYG{l+m}{100}\PYG{+w}{ }\PYG{l+m}{1}\PYG{+w}{ }\PYG{l+m}{100}\PYG{+w}{ }\PYG{n}{c\PYGZus{}p[}\PYG{o}{*}\PYG{n}{]}\PYG{+w}{ }\PYG{n}{file}\PYG{+w}{ }\PYG{n}{dump\PYGZus{}p.out}\PYG{+w}{ }\PYG{n}{mode}\PYG{+w}{ }\PYG{n}{vector}
\end{sphinxVerbatim}

\sphinxAtStartPar
The values calculated by this compute are “intensive”.  The stress
values will be in pressure {\hyperref[\detokenize{units::doc}]{\sphinxcrossref{\DUrole{doc}{units}}}}. The number density
values are in inverse volume {\hyperref[\detokenize{units::doc}]{\sphinxcrossref{\DUrole{doc}{units}}}}.

\sphinxAtStartPar
NOTE 1: The local stress does not include any Lennard\sphinxhyphen{}Jones tail
corrections to the stress added by the
{\hyperref[\detokenize{pair_modify::doc}]{\sphinxcrossref{\DUrole{doc}{pair\_modify tail yes}}}}
command, since those are contributions to the global system pressure.


\subsection{Restrictions}
\label{\detokenize{compute_stress_curvilinear:restrictions}}
\sphinxAtStartPar
These computes calculate the stress tensor contributions for pair styles
only (i.e., no bond, angle, dihedral, etc. contributions, and in the
presence of bonded interactions, the result may be incorrect due to
exclusions for {\hyperref[\detokenize{special_bonds::doc}]{\sphinxcrossref{\DUrole{doc}{special bonds}}}} excluding pairs of atoms
completely). It requires pairwise force calculations not available for most
many\sphinxhyphen{}body pair styles.  Note that \(k\)\sphinxhyphen{}space calculations are also excluded.

\sphinxAtStartPar
These computes are part of the EXTRA\sphinxhyphen{}COMPUTE package.  They are only
enabled if LAMMPS was built with that package.  See the {\hyperref[\detokenize{Build_package::doc}]{\sphinxcrossref{\DUrole{doc}{Build
package}}}} doc page for more info.


\subsection{Related commands}
\label{\detokenize{compute_stress_curvilinear:related-commands}}
\sphinxAtStartPar
{\hyperref[\detokenize{compute_stress_atom::doc}]{\sphinxcrossref{\DUrole{doc}{compute stress/atom}}}}, {\hyperref[\detokenize{compute_pressure::doc}]{\sphinxcrossref{\DUrole{doc}{compute pressure}}}},
{\hyperref[\detokenize{compute_stress_mop::doc}]{\sphinxcrossref{\DUrole{doc}{compute stress/mop/profile}}}}, {\hyperref[\detokenize{compute_stress_cartesian::doc}]{\sphinxcrossref{\DUrole{doc}{compute stress/cartesian}}}}


\subsection{Default}
\label{\detokenize{compute_stress_curvilinear:default}}
\sphinxAtStartPar
The keyword default for ke in style \sphinxstyleemphasis{stress/cylinder} is yes.


\bigskip\hrule\bigskip

\phantomsection\label{\detokenize{compute_stress_curvilinear:ikeshoji4}}
\sphinxAtStartPar
\sphinxstylestrong{(Ikeshoji)} Ikeshoji, Hafskjold, Furuholt, Mol Sim, 29, 101\sphinxhyphen{}109, (2003).

\phantomsection\label{\detokenize{compute_stress_curvilinear:addington1}}
\sphinxAtStartPar
\sphinxstylestrong{(Addington)} Addington, Long, Gubbins, J Chem Phys, 149, 084109 (2018).

\sphinxstepscope

\index{compute stress/mop@\spxentry{compute stress/mop}}\ignorespaces 
\index{compute stress/mop/profile@\spxentry{compute stress/mop/profile}}\ignorespaces 

\section{compute stress/mop command}
\label{\detokenize{compute_stress_mop:compute-stress-mop-command}}\label{\detokenize{compute_stress_mop:index-1}}\label{\detokenize{compute_stress_mop:index-0}}\label{\detokenize{compute_stress_mop::doc}}

\section{compute stress/mop/profile command}
\label{\detokenize{compute_stress_mop:compute-stress-mop-profile-command}}

\subsection{Syntax}
\label{\detokenize{compute_stress_mop:syntax}}
\begin{sphinxVerbatim}[commandchars=\\\{\}]
\PYG{k}{compute }\PYG{n+nv+nvIdentifier}{ID}\PYG{+w}{ }\PYG{n+nv+nvIdentifier}{group\PYGZhy{}ID}\PYG{+w}{ }\PYG{n}{style}\PYG{+w}{ }\PYG{n}{dir}\PYG{+w}{ }\PYG{n}{args}\PYG{+w}{ }\PYG{n}{keywords}\PYG{+w}{ }\PYG{n}{...}
\end{sphinxVerbatim}
\begin{itemize}
\item {} 
\sphinxAtStartPar
ID, group\sphinxhyphen{}ID are documented in {\hyperref[\detokenize{compute::doc}]{\sphinxcrossref{\DUrole{doc}{compute}}}} command

\item {} 
\sphinxAtStartPar
style = \sphinxstyleemphasis{stress/mop} or \sphinxstyleemphasis{stress/mop/profile}

\item {} 
\sphinxAtStartPar
dir = \sphinxstyleemphasis{x} or \sphinxstyleemphasis{y} or \sphinxstyleemphasis{z} is the direction normal to the plane

\item {} 
\sphinxAtStartPar
args = argument specific to the compute style

\item {} 
\sphinxAtStartPar
keywords = \sphinxstyleemphasis{kin} or \sphinxstyleemphasis{conf} or \sphinxstyleemphasis{total} or \sphinxstyleemphasis{pair} or \sphinxstyleemphasis{bond} or \sphinxstyleemphasis{angle} or \sphinxstyleemphasis{dihedral} (one or more can be specified)

\end{itemize}
\begin{sphinxalltt}
\sphinxstyleemphasis{stress/mop} args = pos
  pos = \sphinxstyleemphasis{lower} or \sphinxstyleemphasis{center} or \sphinxstyleemphasis{upper} or coordinate value (distance units) is the position of the plane
\sphinxstyleemphasis{stress/mop/profile} args = origin delta
  origin = \sphinxstyleemphasis{lower} or \sphinxstyleemphasis{center} or \sphinxstyleemphasis{upper} or coordinate value (distance units) is the position of the first plane
  delta = value (distance units) is the distance between planes
\end{sphinxalltt}


\subsection{Examples}
\label{\detokenize{compute_stress_mop:examples}}
\begin{sphinxVerbatim}[commandchars=\\\{\}]
\PYG{k}{compute }\PYG{n+nv+nvIdentifier}{1}\PYG{+w}{ }\PYG{n+nv+nvIdentifier}{all}\PYG{+w}{ }\PYG{n}{stress}\PYG{o}{/}\PYG{n}{mop}\PYG{+w}{ }\PYG{n}{x}\PYG{+w}{ }\PYG{n}{lower}\PYG{+w}{ }\PYG{n}{total}
\PYG{k}{compute }\PYG{n+nv+nvIdentifier}{1}\PYG{+w}{ }\PYG{n+nv+nvIdentifier}{liquid}\PYG{+w}{ }\PYG{n}{stress}\PYG{o}{/}\PYG{n}{mop}\PYG{+w}{ }\PYG{n}{z}\PYG{+w}{ }\PYG{l+m}{0.0}\PYG{+w}{ }\PYG{n}{kin}\PYG{+w}{ }\PYG{n}{conf}
\PYG{k}{fix }\PYG{n+nv+nvIdentifier}{1}\PYG{+w}{ }\PYG{n+nv+nvIdentifier}{all}\PYG{+w}{ }\PYG{n}{ave}\PYG{o}{/}\PYG{n}{time}\PYG{+w}{ }\PYG{l+m}{10}\PYG{+w}{ }\PYG{l+m}{1000}\PYG{+w}{ }\PYG{l+m}{10000}\PYG{+w}{ }\PYG{n}{c\PYGZus{}1[}\PYG{o}{*}\PYG{n}{]}\PYG{+w}{ }\PYG{n}{file}\PYG{+w}{ }\PYG{n}{mop.time}
\PYG{k}{fix }\PYG{n+nv+nvIdentifier}{1}\PYG{+w}{ }\PYG{n+nv+nvIdentifier}{all}\PYG{+w}{ }\PYG{n}{ave}\PYG{o}{/}\PYG{n}{time}\PYG{+w}{ }\PYG{l+m}{10}\PYG{+w}{ }\PYG{l+m}{1000}\PYG{+w}{ }\PYG{l+m}{10000}\PYG{+w}{ }\PYG{n}{c\PYGZus{}1[2]}\PYG{+w}{ }\PYG{n}{file}\PYG{+w}{ }\PYG{n}{mop.time}

\PYG{k}{compute }\PYG{n+nv+nvIdentifier}{1}\PYG{+w}{ }\PYG{n+nv+nvIdentifier}{all}\PYG{+w}{ }\PYG{n}{stress}\PYG{o}{/}\PYG{n}{mop}\PYG{o}{/}\PYG{n}{profile}\PYG{+w}{ }\PYG{n}{x}\PYG{+w}{ }\PYG{n}{lower}\PYG{+w}{ }\PYG{l+m}{0.1}\PYG{+w}{ }\PYG{n}{total}
\PYG{k}{compute }\PYG{n+nv+nvIdentifier}{1}\PYG{+w}{ }\PYG{n+nv+nvIdentifier}{liquid}\PYG{+w}{ }\PYG{n}{stress}\PYG{o}{/}\PYG{n}{mop}\PYG{o}{/}\PYG{n}{profile}\PYG{+w}{ }\PYG{n}{z}\PYG{+w}{ }\PYG{l+m}{0.0}\PYG{+w}{ }\PYG{l+m}{0.25}\PYG{+w}{ }\PYG{n}{kin}\PYG{+w}{ }\PYG{n}{conf}
\PYG{k}{fix }\PYG{n+nv+nvIdentifier}{1}\PYG{+w}{ }\PYG{n+nv+nvIdentifier}{all}\PYG{+w}{ }\PYG{n}{ave}\PYG{o}{/}\PYG{n}{time}\PYG{+w}{ }\PYG{l+m}{500}\PYG{+w}{ }\PYG{l+m}{20}\PYG{+w}{ }\PYG{l+m}{10000}\PYG{+w}{ }\PYG{n}{c\PYGZus{}1[}\PYG{o}{*}\PYG{n}{]}\PYG{+w}{ }\PYG{n}{ave}\PYG{+w}{ }\PYG{n}{running}\PYG{+w}{ }\PYG{n}{overwrite}\PYG{+w}{ }\PYG{n}{file}\PYG{+w}{ }\PYG{n}{mopp.time}\PYG{+w}{ }\PYG{n}{mode}\PYG{+w}{ }\PYG{n}{vector}
\end{sphinxVerbatim}


\subsection{Description}
\label{\detokenize{compute_stress_mop:description}}
\sphinxAtStartPar
Compute \sphinxstyleemphasis{stress/mop} and compute \sphinxstyleemphasis{stress/mop/profile} define
computations that calculate components of the local stress tensor using
the method of planes {\hyperref[\detokenize{compute_stress_mop:mop-todd}]{\sphinxcrossref{\DUrole{std,std-ref}{(Todd)}}}}.  Specifically in compute
\sphinxstyleemphasis{stress/mop} calculates 3 components are computed in directions \sphinxstyleemphasis{dir},\sphinxstyleemphasis{x}; \sphinxstyleemphasis{dir},\sphinxstyleemphasis{y}; and \sphinxstyleemphasis{dir},\sphinxstyleemphasis{z}; where \sphinxstyleemphasis{dir} is the direction
normal to the plane, while in compute \sphinxstyleemphasis{stress/mop/profile} the profile
of the stress is computed.

\sphinxAtStartPar
Contrary to methods based on histograms of atomic stress (i.e., using
{\hyperref[\detokenize{compute_stress_atom::doc}]{\sphinxcrossref{\DUrole{doc}{compute stress/atom}}}}), the method of planes
is compatible with mechanical balance in heterogeneous systems and at
interfaces {\hyperref[\detokenize{compute_stress_mop:mop-todd}]{\sphinxcrossref{\DUrole{std,std-ref}{(Todd)}}}}.

\sphinxAtStartPar
The stress tensor is the sum of a kinetic term and a configurational
term, which are given respectively by Eq. (21) and Eq. (16) in
{\hyperref[\detokenize{compute_stress_mop:mop-todd}]{\sphinxcrossref{\DUrole{std,std-ref}{(Todd)}}}}. For the kinetic part, the algorithm considers
that atoms have crossed the plane if their positions at times
\(t-\Delta t\) and \(t\) are one on either side of the plane,
and uses the velocity at time \(t-\Delta t/2\) given by the velocity
Verlet algorithm.

\sphinxAtStartPar
\DUrole{versionmodified,added}{Added in version 15Jun2023: }contributions from bond, angle and dihedral potentials

\sphinxAtStartPar
Between one and seven keywords can be used to indicate which contributions
to the stress must be computed: total stress (total), kinetic stress
(kin), configurational stress (conf), stress due to bond stretching
(bond), stress due to angle bending (angle), stress due to dihedral terms (dihedral)
and/or due to pairwise non\sphinxhyphen{}bonded interactions (pair).

\sphinxAtStartPar
NOTE 1: The configurational stress is computed considering all pairs of
atoms where at least one atom belongs to group group\sphinxhyphen{}ID.

\sphinxAtStartPar
NOTE 2: The local stress does not include any Lennard\sphinxhyphen{}Jones tail
corrections to the stress added by the {\hyperref[\detokenize{pair_modify::doc}]{\sphinxcrossref{\DUrole{doc}{pair\_modify tail yes}}}} command, since those are contributions to the global
system pressure.

\sphinxAtStartPar
NOTE 3: The local stress profile generated by compute
\sphinxstyleemphasis{stress/mop/profile} is similar to that obtained by compute
{\hyperref[\detokenize{compute_stress_cartesian::doc}]{\sphinxcrossref{\DUrole{doc}{stress/cartesian}}}}.
A key difference is that compute \sphinxstyleemphasis{stress/mop/profile}
considers particles crossing a set of planes, while
\sphinxstyleemphasis{stress/cartesian} computes averages for a set
of small volumes.
Moreover, \sphinxstyleemphasis{stress/cartesian} compute computes the diagonal components of the stress
tensor \(P_{xx}\), \(P_{yy}\), and \(P_{zz}\), while
\sphinxstyleemphasis{stress/mop/profile} computes the components
\(P_{ix}\), \(P_{iy}\), and \(P_{iz}\), where \(i\) is the
direction normal to the plane.


\subsection{Output info}
\label{\detokenize{compute_stress_mop:output-info}}
\sphinxAtStartPar
Compute \sphinxstyleemphasis{stress/mop} calculates a global vector (indices starting at 1),
with 3 values for each declared keyword (in the order the keywords have
been declared). For each keyword, the stress tensor components are
ordered as follows: stress\_dir,x, stress\_dir,y, and stress\_dir,z.

\sphinxAtStartPar
Compute \sphinxstyleemphasis{stress/mop/profile} instead calculates a global array, with 1
column giving the position of the planes where the stress tensor was
computed, and with 3 columns of values for each declared keyword (in the
order the keywords have been declared). For each keyword, the profiles
of stress tensor components are ordered as follows: stress\_dir,x;
stress\_dir,y; and stress\_dir,z.

\sphinxAtStartPar
The values are in pressure {\hyperref[\detokenize{units::doc}]{\sphinxcrossref{\DUrole{doc}{units}}}}.

\sphinxAtStartPar
The values produced by this compute can be accessed by various
{\hyperref[\detokenize{Howto_output::doc}]{\sphinxcrossref{\DUrole{doc}{output commands}}}}.  For instance, the results can be
written to a file using the {\hyperref[\detokenize{fix_ave_time::doc}]{\sphinxcrossref{\DUrole{doc}{fix ave/time}}}}
command. Please see the example in the examples/PACKAGES/mop folder.


\subsection{Restrictions}
\label{\detokenize{compute_stress_mop:restrictions}}
\sphinxAtStartPar
These styles are part of the EXTRA\sphinxhyphen{}COMPUTE package. They are only
enabled if LAMMPS is built with that package. See the {\hyperref[\detokenize{Build_package::doc}]{\sphinxcrossref{\DUrole{doc}{Build
package}}}} doc page on for more info.

\sphinxAtStartPar
The method is implemented for orthogonal simulation boxes whose
size does not change in time, and axis\sphinxhyphen{}aligned planes.

\sphinxAtStartPar
The method only works with two\sphinxhyphen{}body pair interactions, because it
requires the class method \sphinxcode{\sphinxupquote{Pair::single()}} to be implemented, which is
not possible for manybody potentials.  In particular, compute
\sphinxstyleemphasis{stress/mop/profile} and \sphinxstyleemphasis{stress/mop} do not work with more than two\sphinxhyphen{}body
pair interactions, long range (kspace) interactions and
improper intramolecular interactions.

\sphinxAtStartPar
The impact of fixes that affect the stress (e.g. fix langevin) is
also not included in the stress computed here.


\subsection{Related commands}
\label{\detokenize{compute_stress_mop:related-commands}}
\sphinxAtStartPar
{\hyperref[\detokenize{compute_stress_atom::doc}]{\sphinxcrossref{\DUrole{doc}{compute stress/atom}}}},
{\hyperref[\detokenize{compute_pressure::doc}]{\sphinxcrossref{\DUrole{doc}{compute pressure}}}},
{\hyperref[\detokenize{compute_stress_cartesian::doc}]{\sphinxcrossref{\DUrole{doc}{compute stress/cartesian}}}},
{\hyperref[\detokenize{compute_stress_curvilinear::doc}]{\sphinxcrossref{\DUrole{doc}{compute stress/cylinder}}}},
{\hyperref[\detokenize{compute_stress_curvilinear::doc}]{\sphinxcrossref{\DUrole{doc}{compute stress/spherical}}}}


\subsection{Default}
\label{\detokenize{compute_stress_mop:default}}
\sphinxAtStartPar
none


\bigskip\hrule\bigskip

\phantomsection\label{\detokenize{compute_stress_mop:mop-todd}}
\sphinxAtStartPar
\sphinxstylestrong{(Todd)} B. D. Todd, Denis J. Evans, and Peter J. Daivis: “Pressure tensor for inhomogeneous fluids”,
Phys. Rev. E 52, 1627 (1995).

\phantomsection\label{\detokenize{compute_stress_mop:ikeshoji3}}
\sphinxAtStartPar
\sphinxstylestrong{(Ikeshoji)} Ikeshoji, Hafskjold, Furuholt, Mol Sim, 29, 101\sphinxhyphen{}109, (2003).

\sphinxstepscope

\index{compute force/tally@\spxentry{compute force/tally}}\ignorespaces 
\index{compute heat/flux/tally@\spxentry{compute heat/flux/tally}}\ignorespaces 
\index{compute heat/flux/virial/tally@\spxentry{compute heat/flux/virial/tally}}\ignorespaces 
\index{compute pe/tally@\spxentry{compute pe/tally}}\ignorespaces 
\index{compute pe/mol/tally@\spxentry{compute pe/mol/tally}}\ignorespaces 
\index{compute stress/tally@\spxentry{compute stress/tally}}\ignorespaces 

\section{compute force/tally command}
\label{\detokenize{compute_tally:compute-force-tally-command}}\label{\detokenize{compute_tally:index-5}}\label{\detokenize{compute_tally:index-4}}\label{\detokenize{compute_tally:index-3}}\label{\detokenize{compute_tally:index-2}}\label{\detokenize{compute_tally:index-1}}\label{\detokenize{compute_tally:index-0}}\label{\detokenize{compute_tally::doc}}

\section{compute heat/flux/tally command}
\label{\detokenize{compute_tally:compute-heat-flux-tally-command}}

\section{compute heat/flux/virial/tally command}
\label{\detokenize{compute_tally:compute-heat-flux-virial-tally-command}}

\section{compute pe/tally command}
\label{\detokenize{compute_tally:compute-pe-tally-command}}

\section{compute pe/mol/tally command}
\label{\detokenize{compute_tally:compute-pe-mol-tally-command}}

\section{compute stress/tally command}
\label{\detokenize{compute_tally:compute-stress-tally-command}}

\subsection{Syntax}
\label{\detokenize{compute_tally:syntax}}
\begin{sphinxVerbatim}[commandchars=\\\{\}]
\PYG{k}{compute }\PYG{n+nv+nvIdentifier}{ID}\PYG{+w}{ }\PYG{n+nv+nvIdentifier}{group\PYGZhy{}ID}\PYG{+w}{ }\PYG{n}{style}\PYG{+w}{ }\PYG{n}{group2}\PYG{o}{\PYGZhy{}}\PYG{n}{ID}
\end{sphinxVerbatim}
\begin{itemize}
\item {} 
\sphinxAtStartPar
ID, group\sphinxhyphen{}ID are documented in {\hyperref[\detokenize{compute::doc}]{\sphinxcrossref{\DUrole{doc}{compute}}}} command

\item {} 
\sphinxAtStartPar
style = \sphinxstyleemphasis{force/tally} or \sphinxstyleemphasis{heat/flux/tally} or \sphinxstyleemphasis{heat/flux/virial/tally} or \sphinxstyleemphasis{pe/tally} or \sphinxstyleemphasis{pe/mol/tally} or \sphinxstyleemphasis{stress/tally}

\item {} 
\sphinxAtStartPar
group2\sphinxhyphen{}ID = group ID of second (or same) group

\end{itemize}


\subsection{Examples}
\label{\detokenize{compute_tally:examples}}
\begin{sphinxVerbatim}[commandchars=\\\{\}]
\PYG{k}{compute }\PYG{n+nv+nvIdentifier}{1}\PYG{+w}{ }\PYG{n+nv+nvIdentifier}{lower}\PYG{+w}{ }\PYG{n}{force}\PYG{o}{/}\PYG{n}{tally}\PYG{+w}{ }\PYG{n}{upper}
\PYG{k}{compute }\PYG{n+nv+nvIdentifier}{1}\PYG{+w}{ }\PYG{n+nv+nvIdentifier}{left}\PYG{+w}{ }\PYG{n}{pe}\PYG{o}{/}\PYG{n}{tally}\PYG{+w}{ }\PYG{n}{right}
\PYG{k}{compute }\PYG{n+nv+nvIdentifier}{1}\PYG{+w}{ }\PYG{n+nv+nvIdentifier}{lower}\PYG{+w}{ }\PYG{n}{stress}\PYG{o}{/}\PYG{n}{tally}\PYG{+w}{ }\PYG{n}{lower}
\PYG{k}{compute }\PYG{n+nv+nvIdentifier}{1}\PYG{+w}{ }\PYG{n+nv+nvIdentifier}{subregion}\PYG{+w}{ }\PYG{n}{heat}\PYG{o}{/}\PYG{n}{flux}\PYG{o}{/}\PYG{n}{tally}\PYG{+w}{ }\PYG{n}{all}
\PYG{k}{compute }\PYG{n+nv+nvIdentifier}{1}\PYG{+w}{ }\PYG{n+nv+nvIdentifier}{liquid}\PYG{+w}{ }\PYG{n}{heat}\PYG{o}{/}\PYG{n}{flux}\PYG{o}{/}\PYG{n}{virial}\PYG{o}{/}\PYG{n}{tally}\PYG{+w}{ }\PYG{n}{solid}
\end{sphinxVerbatim}


\subsection{Description}
\label{\detokenize{compute_tally:description}}
\sphinxAtStartPar
Define a computation that calculates properties between two groups of
atoms by accumulating them from pairwise non\sphinxhyphen{}bonded computations.
Except for \sphinxstyleemphasis{heat/flux/virial/tally}, the two groups can be the same.
This is similar to {\hyperref[\detokenize{compute_group_group::doc}]{\sphinxcrossref{\DUrole{doc}{compute group/group}}}}
only that the data is
accumulated directly during the non\sphinxhyphen{}bonded force computation. The
computes \sphinxstyleemphasis{force/tally}, \sphinxstyleemphasis{pe/tally}, \sphinxstyleemphasis{stress/tally}, and
\sphinxstyleemphasis{heat/flux/tally} are primarily provided as example how to program
additional, more sophisticated computes using the tally callback
mechanism. Compute \sphinxstyleemphasis{pe/mol/tally} is one such style, that can—through using
this mechanism—separately tally intermolecular
and intramolecular energies. Something that would otherwise be
impossible without integrating this as a core functionality into
the base classes of LAMMPS.


\bigskip\hrule\bigskip


\sphinxAtStartPar
Compute \sphinxstyleemphasis{heat/flux/tally} obtains the heat flux
(strictly speaking, heat flow) inside the first group,
which is the sum of the convective contribution
due to atoms in the first group and the virial contribution
due to interaction between the first and second groups:
\begin{equation*}
\begin{split}\mathbf{Q}=  \sum_{i \in \text{group 1}} e_i \mathbf{v}_i + \frac{1}{2} \sum_{i \in \text{group 1}} \sum_{\substack{j \in \text{group 2} \\ j \neq i } } \left( \mathbf{F}_{ij} \cdot \mathbf{v}_j \right) \mathbf{r}_{ij}\end{split}
\end{equation*}
\sphinxAtStartPar
When the second group in \sphinxstyleemphasis{heat/flux/tally} is set to “all”,
the resulting values will be identical
to that obtained by {\hyperref[\detokenize{compute_heat_flux::doc}]{\sphinxcrossref{\DUrole{doc}{compute heat/flux}}}},
provided only pairwise interactions exist.

\sphinxAtStartPar
Compute \sphinxstyleemphasis{heat/flux/virial/tally} obtains the total virial heat flux
(strictly speaking, heat flow) into the first group due to interaction
with the second group, and is defined as:
\begin{equation*}
\begin{split}Q = \frac{1}{2} \sum_{i \in \text{group 1}} \sum_{j \in \text{group 2}} \mathbf{F}_{ij} \cdot \left(\mathbf{v}_i + \mathbf{v}_j \right)\end{split}
\end{equation*}
\sphinxAtStartPar
Although, the \sphinxstyleemphasis{heat/flux/virial/tally} compute
does not include the convective term,
it can be used to obtain the total heat flux over control surfaces,
when there are no particles crossing over,
such as is often in solid\textendash{}solid and solid\textendash{}liquid interfaces.
This would be identical to the method of planes method.
Note that the \sphinxstyleemphasis{heat/flux/virial/tally} compute is distinctly different
from the \sphinxstyleemphasis{heat/flux} and \sphinxstyleemphasis{heat/flux/tally} computes,
that are essentially volume averaging methods.
The following example demonstrates the difference:

\begin{sphinxVerbatim}[commandchars=\\\{\}]
\PYG{c}{\PYGZsh{} System with only pairwise interactions.}
\PYG{c}{\PYGZsh{} Non\PYGZhy{}periodic boundaries in the x direction.}
\PYG{c}{\PYGZsh{} Has LeftLiquid and RightWall groups along x direction.}

\PYG{c}{\PYGZsh{} Heat flux over the solid\PYGZhy{}liquid interface}
\PYG{k}{compute }\PYG{n+nv+nvIdentifier}{hflow\PYGZus{}hfvt}\PYG{+w}{ }\PYG{n+nv+nvIdentifier}{RightWall}\PYG{+w}{ }\PYG{n}{heat}\PYG{o}{/}\PYG{n}{flux}\PYG{o}{/}\PYG{n}{virial}\PYG{o}{/}\PYG{n}{tally}\PYG{+w}{ }\PYG{n}{LeftLiquid}
\PYG{k}{variable }\PYG{n+nv+nvIdentifier}{hflux\PYGZus{}hfvt}\PYG{+w}{ }\PYG{n}{equal}\PYG{+w}{ }\PYG{n}{c\PYGZus{}hflow\PYGZus{}hfvt}\PYG{o}{/}\PYG{n+nv}{(}\PYG{n+nv}{ly*lz}\PYG{n+nv}{)}

\PYG{c}{\PYGZsh{} x component of approximate heat flux vector inside the liquid region,}
\PYG{c}{\PYGZsh{} two approaches.}
\PYG{c}{\PYGZsh{}}
\PYG{k}{compute }\PYG{n+nv+nvIdentifier}{myKE}\PYG{+w}{ }\PYG{n+nv+nvIdentifier}{all}\PYG{+w}{ }\PYG{n}{ke}\PYG{o}{/}\PYG{n}{atom}
\PYG{k}{compute }\PYG{n+nv+nvIdentifier}{myPE}\PYG{+w}{ }\PYG{n+nv+nvIdentifier}{all}\PYG{+w}{ }\PYG{n}{pe}\PYG{o}{/}\PYG{n}{atom}
\PYG{k}{compute }\PYG{n+nv+nvIdentifier}{myStress}\PYG{+w}{ }\PYG{n+nv+nvIdentifier}{all}\PYG{+w}{ }\PYG{n}{stress}\PYG{o}{/}\PYG{n}{atom}\PYG{+w}{ }\PYG{n}{NULL}\PYG{+w}{ }\PYG{n}{virial}
\PYG{k}{compute }\PYG{n+nv+nvIdentifier}{hflow\PYGZus{}hf}\PYG{+w}{ }\PYG{n+nv+nvIdentifier}{LeftLiquid}\PYG{+w}{ }\PYG{n}{heat}\PYG{o}{/}\PYG{n}{flux}\PYG{+w}{ }\PYG{n}{myKE}\PYG{+w}{ }\PYG{n}{myPE}\PYG{+w}{ }\PYG{n}{myStress}
\PYG{k}{variable }\PYG{n+nv+nvIdentifier}{hflux\PYGZus{}hf}\PYG{+w}{ }\PYG{n}{equal}\PYG{+w}{ }\PYG{n}{c\PYGZus{}hflow\PYGZus{}hf[1]}\PYG{o}{/}\PYG{n+nv}{\PYGZdl{}\PYGZob{}}\PYG{n+nv}{volLiq}\PYG{n+nv}{\PYGZcb{}}
\PYG{c}{\PYGZsh{}}
\PYG{k}{compute }\PYG{n+nv+nvIdentifier}{hflow\PYGZus{}hft}\PYG{+w}{ }\PYG{n+nv+nvIdentifier}{LeftLiquid}\PYG{+w}{ }\PYG{n}{heat}\PYG{o}{/}\PYG{n}{flux}\PYG{o}{/}\PYG{n}{tally}\PYG{+w}{ }\PYG{n}{all}
\PYG{k}{variable }\PYG{n+nv+nvIdentifier}{hflux\PYGZus{}hft}\PYG{+w}{ }\PYG{n}{equal}\PYG{+w}{ }\PYG{n}{c\PYGZus{}hflow\PYGZus{}hft[1]}\PYG{o}{/}\PYG{n+nv}{\PYGZdl{}\PYGZob{}}\PYG{n+nv}{volLiq}\PYG{n+nv}{\PYGZcb{}}

\PYG{c}{\PYGZsh{} Pressure over the solid\PYGZhy{}liquid interface, three approaches.}
\PYG{c}{\PYGZsh{}}
\PYG{k}{compute }\PYG{n+nv+nvIdentifier}{force\PYGZus{}gg}\PYG{+w}{ }\PYG{n+nv+nvIdentifier}{RightWall}\PYG{+w}{ }\PYG{n}{group}\PYG{o}{/}\PYG{k}{group }\PYG{n+nv+nvIdentifier}{LeftLiquid}
\PYG{k}{variable }\PYG{n+nv+nvIdentifier}{press\PYGZus{}gg}\PYG{+w}{ }\PYG{n}{equal}\PYG{+w}{ }\PYG{n}{c\PYGZus{}force\PYGZus{}gg[1]}\PYG{o}{/}\PYG{n+nv}{(}\PYG{n+nv}{ly*lz}\PYG{n+nv}{)}
\PYG{c}{\PYGZsh{}}
\PYG{k}{compute }\PYG{n+nv+nvIdentifier}{force\PYGZus{}ft}\PYG{+w}{ }\PYG{n+nv+nvIdentifier}{RightWall}\PYG{+w}{ }\PYG{n}{force}\PYG{o}{/}\PYG{n}{tally}\PYG{+w}{ }\PYG{n}{LeftLiquid}
\PYG{k}{compute }\PYG{n+nv+nvIdentifier}{rforce\PYGZus{}ft}\PYG{+w}{ }\PYG{n+nv+nvIdentifier}{RightWall}\PYG{+w}{ }\PYG{n}{reduce}\PYG{+w}{ }\PYG{n}{sum}\PYG{+w}{ }\PYG{n}{c\PYGZus{}force\PYGZus{}ft[1]}
\PYG{k}{variable }\PYG{n+nv+nvIdentifier}{press\PYGZus{}ft}\PYG{+w}{ }\PYG{n}{equal}\PYG{+w}{ }\PYG{n}{c\PYGZus{}rforce\PYGZus{}ft}\PYG{o}{/}\PYG{n+nv}{(}\PYG{n+nv}{ly*lz}\PYG{n+nv}{)}
\PYG{c}{\PYGZsh{}}
\PYG{k}{compute }\PYG{n+nv+nvIdentifier}{rforce\PYGZus{}hfvt}\PYG{+w}{ }\PYG{n+nv+nvIdentifier}{all}\PYG{+w}{ }\PYG{n}{reduce}\PYG{+w}{ }\PYG{n}{sum}\PYG{+w}{ }\PYG{n}{c\PYGZus{}hflow\PYGZus{}hfvt[1]}
\PYG{k}{variable }\PYG{n+nv+nvIdentifier}{press\PYGZus{}hfvt}\PYG{+w}{ }\PYG{n}{equal}\PYG{+w}{ }\PYG{n}{c\PYGZus{}rforce\PYGZus{}hfvt}\PYG{o}{/}\PYG{n+nv}{(}\PYG{n+nv}{ly*lz}\PYG{n+nv}{)}
\end{sphinxVerbatim}


\bigskip\hrule\bigskip


\sphinxAtStartPar
The pairwise contributions are computing via a callback that the
compute registers with the non\sphinxhyphen{}bonded pairwise force computation.
This limits the use to systems that have no bonds, no Kspace, and no
many\sphinxhyphen{}body interactions. On the other hand, the computation does not
have to compute forces or energies a second time and thus can be much
more efficient. The callback mechanism allows to write more complex
pairwise property computations.


\bigskip\hrule\bigskip



\subsection{Output info}
\label{\detokenize{compute_tally:output-info}}\begin{itemize}
\item {} 
\sphinxAtStartPar
Compute \sphinxstyleemphasis{pe/tally} calculates a global scalar (the energy) and a per
atom scalar (the contributions of the single atom to the global
scalar).

\item {} 
\sphinxAtStartPar
Compute \sphinxstyleemphasis{pe/mol/tally} calculates a global four\sphinxhyphen{}element vector containing
(in this order): \sphinxstyleemphasis{evdwl} and \sphinxstyleemphasis{ecoul} for intramolecular pairs and
\sphinxstyleemphasis{evdwl} and \sphinxstyleemphasis{ecoul} for intermolecular pairs. Since molecules are
identified by their molecule IDs, the partitioning does not have to be
related to molecules, but the energies are tallied into the respective
slots depending on whether the molecule IDs of a pair are the same or
different.

\item {} 
\sphinxAtStartPar
Compute \sphinxstyleemphasis{force/tally} calculates a global scalar (the force magnitude)
and a per atom 3\sphinxhyphen{}element vector (force contribution from each atom).

\item {} 
\sphinxAtStartPar
Compute \sphinxstyleemphasis{stress/tally} calculates a global scalar
(average of the diagonal elements of the stress tensor) and a per atom
vector (the six elements of stress tensor contributions from the
individual atom).

\item {} 
\sphinxAtStartPar
As in {\hyperref[\detokenize{compute_heat_flux::doc}]{\sphinxcrossref{\DUrole{doc}{compute heat/flux}}}},
compute \sphinxstyleemphasis{heat/flux/tally} calculates a global vector of length 6,
where the first three components are the \(x\), \(y\), \(z\)
components of the full heat flow vector,
and the next three components are the corresponding components
of just the convective portion of the flow (i.e., the
first term in the equation for \(\mathbf{Q}\)).

\item {} 
\sphinxAtStartPar
Compute \sphinxstyleemphasis{heat/flux/virial/tally} calculates a global scalar (heat flow)
and a per atom three\sphinxhyphen{}element vector
(contribution to the force acting over atoms in the first group
from individual atoms in both groups).

\end{itemize}

\sphinxAtStartPar
Both the scalar and vector values calculated by this compute are
“extensive”.


\subsection{Restrictions}
\label{\detokenize{compute_tally:restrictions}}
\sphinxAtStartPar
This compute is part of the TALLY package.  It is only enabled if LAMMPS
was built with that package.  See the {\hyperref[\detokenize{Build_package::doc}]{\sphinxcrossref{\DUrole{doc}{Build package}}}} page for more info.

\sphinxAtStartPar
Not all pair styles can be evaluated in a pairwise mode as required by
this compute.  For example, 3\sphinxhyphen{}body and other many\sphinxhyphen{}body potentials, such
as {\hyperref[\detokenize{pair_tersoff::doc}]{\sphinxcrossref{\DUrole{doc}{Tersoff}}}} and {\hyperref[\detokenize{pair_sw::doc}]{\sphinxcrossref{\DUrole{doc}{Stillinger\sphinxhyphen{}Weber}}}}
cannot be used.  {\hyperref[\detokenize{pair_eam::doc}]{\sphinxcrossref{\DUrole{doc}{EAM}}}} potentials only include the pair
potential portion of the EAM interaction when used by this compute, not
the embedding term.  Also bonded or Kspace interactions do not
contribute to this compute.

\sphinxAtStartPar
These computes are not compatible with accelerated pair styles from the
GPU, INTEL, KOKKOS, or OPENMP packages. They will either create an error
or print a warning when required data was not tallied in the required way
and thus the data acquisition functions from these computes not called.

\sphinxAtStartPar
When used with dynamic groups, a {\hyperref[\detokenize{run::doc}]{\sphinxcrossref{\DUrole{doc}{run 0}}}} command needs to
be inserted in order to initialize the dynamic groups before accessing
the computes.


\subsection{Related commands}
\label{\detokenize{compute_tally:related-commands}}\begin{itemize}
\item {} 
\sphinxAtStartPar
{\hyperref[\detokenize{compute_group_group::doc}]{\sphinxcrossref{\DUrole{doc}{compute group/group}}}}

\item {} 
\sphinxAtStartPar
{\hyperref[\detokenize{compute_heat_flux::doc}]{\sphinxcrossref{\DUrole{doc}{compute heat/flux}}}}

\end{itemize}


\subsection{Default}
\label{\detokenize{compute_tally:default}}
\sphinxAtStartPar
none

\sphinxstepscope

\index{compute tdpd/cc/atom@\spxentry{compute tdpd/cc/atom}}\ignorespaces 

\section{compute tdpd/cc/atom command}
\label{\detokenize{compute_tdpd_cc_atom:compute-tdpd-cc-atom-command}}\label{\detokenize{compute_tdpd_cc_atom:index-0}}\label{\detokenize{compute_tdpd_cc_atom::doc}}

\subsection{Syntax}
\label{\detokenize{compute_tdpd_cc_atom:syntax}}
\begin{sphinxVerbatim}[commandchars=\\\{\}]
\PYG{k}{compute }\PYG{n+nv+nvIdentifier}{ID}\PYG{+w}{ }\PYG{n+nv+nvIdentifier}{group\PYGZhy{}ID}\PYG{+w}{ }\PYG{n}{tdpd}\PYG{o}{/}\PYG{n}{cc}\PYG{o}{/}\PYG{n}{atom}\PYG{+w}{ }\PYG{n}{index}
\end{sphinxVerbatim}
\begin{itemize}
\item {} 
\sphinxAtStartPar
ID, group\sphinxhyphen{}ID are documented in {\hyperref[\detokenize{compute::doc}]{\sphinxcrossref{\DUrole{doc}{compute}}}} command

\item {} 
\sphinxAtStartPar
tdpd/cc/atom = style name of this compute command

\item {} 
\sphinxAtStartPar
index = index of chemical species (1 to Nspecies)

\end{itemize}


\subsection{Examples}
\label{\detokenize{compute_tdpd_cc_atom:examples}}
\begin{sphinxVerbatim}[commandchars=\\\{\}]
\PYG{k}{compute }\PYG{n+nv+nvIdentifier}{1}\PYG{+w}{ }\PYG{n+nv+nvIdentifier}{all}\PYG{+w}{ }\PYG{n}{tdpd}\PYG{o}{/}\PYG{n}{cc}\PYG{o}{/}\PYG{n}{atom}\PYG{+w}{ }\PYG{l+m}{2}
\end{sphinxVerbatim}


\subsection{Description}
\label{\detokenize{compute_tdpd_cc_atom:description}}
\sphinxAtStartPar
Define a computation that calculates the per\sphinxhyphen{}atom chemical
concentration of a specified species for each tDPD particle in a
group.

\sphinxAtStartPar
The chemical concentration of each species is defined as the number of
molecules carried by a tDPD particle for dilute solution.  For more
details see {\hyperref[\detokenize{compute_tdpd_cc_atom:li2015a}]{\sphinxcrossref{\DUrole{std,std-ref}{(Li2015)}}}}.


\subsection{Output info}
\label{\detokenize{compute_tdpd_cc_atom:output-info}}
\sphinxAtStartPar
This compute calculates a per\sphinxhyphen{}atom vector, which can be accessed by
any command that uses per\sphinxhyphen{}atom values from a compute as input. See the
{\hyperref[\detokenize{Howto_output::doc}]{\sphinxcrossref{\DUrole{doc}{Howto output}}}} page for an overview of LAMMPS
output options.

\sphinxAtStartPar
The per\sphinxhyphen{}atom vector values will be in the units of chemical species
per unit mass.


\subsection{Restrictions}
\label{\detokenize{compute_tdpd_cc_atom:restrictions}}
\sphinxAtStartPar
This compute is part of the DPD\sphinxhyphen{}MESO package.  It is only enabled if
LAMMPS was built with that package.  See the {\hyperref[\detokenize{Build_package::doc}]{\sphinxcrossref{\DUrole{doc}{Build package}}}} page for more info.


\subsection{Related commands}
\label{\detokenize{compute_tdpd_cc_atom:related-commands}}
\sphinxAtStartPar
{\hyperref[\detokenize{pair_mesodpd::doc}]{\sphinxcrossref{\DUrole{doc}{pair\_style tdpd}}}}


\subsection{Default}
\label{\detokenize{compute_tdpd_cc_atom:default}}
\sphinxAtStartPar
none


\bigskip\hrule\bigskip

\phantomsection\label{\detokenize{compute_tdpd_cc_atom:li2015a}}
\sphinxAtStartPar
\sphinxstylestrong{(Li2015)} Li, Yazdani, Tartakovsky, Karniadakis, J Chem Phys, 143:
014101 (2015).  DOI: 10.1063/1.4923254

\sphinxstepscope

\index{compute temp@\spxentry{compute temp}}\ignorespaces 
\index{compute temp/kk@\spxentry{compute temp/kk}}\ignorespaces 

\section{compute temp command}
\label{\detokenize{compute_temp:compute-temp-command}}\label{\detokenize{compute_temp:index-1}}\label{\detokenize{compute_temp:index-0}}\label{\detokenize{compute_temp::doc}}
\sphinxAtStartPar
Accelerator Variants: \sphinxstyleemphasis{temp/kk}


\subsection{Syntax}
\label{\detokenize{compute_temp:syntax}}
\begin{sphinxVerbatim}[commandchars=\\\{\}]
\PYG{k}{compute }\PYG{n+nv+nvIdentifier}{ID}\PYG{+w}{ }\PYG{n+nv+nvIdentifier}{group\PYGZhy{}ID}\PYG{+w}{ }\PYG{n}{temp}
\end{sphinxVerbatim}
\begin{itemize}
\item {} 
\sphinxAtStartPar
ID, group\sphinxhyphen{}ID are documented in {\hyperref[\detokenize{compute::doc}]{\sphinxcrossref{\DUrole{doc}{compute}}}} command

\item {} 
\sphinxAtStartPar
temp = style name of this compute command

\end{itemize}


\subsection{Examples}
\label{\detokenize{compute_temp:examples}}
\begin{sphinxVerbatim}[commandchars=\\\{\}]
\PYG{k}{compute }\PYG{n+nv+nvIdentifier}{1}\PYG{+w}{ }\PYG{n+nv+nvIdentifier}{all}\PYG{+w}{ }\PYG{n}{temp}
\PYG{k}{compute }\PYG{n+nv+nvIdentifier}{myTemp}\PYG{+w}{ }\PYG{n+nv+nvIdentifier}{mobile}\PYG{+w}{ }\PYG{n}{temp}
\end{sphinxVerbatim}


\subsection{Description}
\label{\detokenize{compute_temp:description}}
\sphinxAtStartPar
Define a computation that calculates the temperature of a group of
atoms.  A compute of this style can be used by any command that
computes a temperature, e.g. {\hyperref[\detokenize{thermo_modify::doc}]{\sphinxcrossref{\DUrole{doc}{thermo\_modify}}}},
{\hyperref[\detokenize{fix_temp_rescale::doc}]{\sphinxcrossref{\DUrole{doc}{fix temp/rescale}}}}, {\hyperref[\detokenize{fix_nh::doc}]{\sphinxcrossref{\DUrole{doc}{fix npt}}}},
etc.

\sphinxAtStartPar
The temperature is calculated by the formula
\begin{equation*}
\begin{split}T  = \frac{2 E_\mathrm{kin}}{N_\mathrm{DOF} k_B} \quad \mathrm{with} \quad
E_\mathrm{kin} = \sum^{N_\mathrm{atoms}}_{i=1} \frac{1}{2} m_i v^2_i \quad \mathrm{and} \quad
N_\mathrm{DOF} = n_\mathrm{dim} N_\mathrm{atoms} - n_\mathrm{dim} - N_\mathrm{fix DOFs}\end{split}
\end{equation*}
\sphinxAtStartPar
where \(E_\mathrm{kin}\) is the total kinetic energy of the group of
atoms, \(n_\mathrm{dim}\) is the dimensionality of the simulation
(i.e. either 2 or 3), \(N_\mathrm{atoms}\) is the number of atoms in
the group, \(N_\mathrm{fix DOFs}\) is the number of degrees of
freedom removed by fix commands (see below), \(k_B\) is the
Boltzmann constant, and \(T\) is the resulting computed temperature.

\sphinxAtStartPar
A symmetric tensor, stored as a six\sphinxhyphen{}element vector, is also calculated
by this compute for use in the computation of a pressure tensor by the
{\hyperref[\detokenize{compute_pressure::doc}]{\sphinxcrossref{\DUrole{doc}{compute pressue}}}} command.  The formula for
the components of the tensor is the same as the above expression for
\(E_\mathrm{kin}\), except that the 1/2 factor is NOT included and
the \(v_i^2\) is replaced by \(v_{i,x} v_{i,y}\) for the
\(xy\) component, and so on.  Note that because it lacks the 1/2
factor, these tensor components are twice those of the traditional
kinetic energy tensor.  The six components of the vector are ordered
\(xx\), \(yy\), \(zz\), \(xy\), \(xz\),
\(yz\).

\sphinxAtStartPar
The number of atoms contributing to the temperature is assumed to be
constant for the duration of the run; use the \sphinxstyleemphasis{dynamic} option of the
{\hyperref[\detokenize{compute_modify::doc}]{\sphinxcrossref{\DUrole{doc}{compute\_modify}}}} command if this is not the case.

\sphinxAtStartPar
This compute subtracts out degrees\sphinxhyphen{}of\sphinxhyphen{}freedom due to fixes that
constrain molecular motion, such as {\hyperref[\detokenize{fix_shake::doc}]{\sphinxcrossref{\DUrole{doc}{fix shake}}}} and
{\hyperref[\detokenize{fix_rigid::doc}]{\sphinxcrossref{\DUrole{doc}{fix rigid}}}}.  This means the temperature of groups of
atoms that include these constraints will be computed correctly.  If
needed, the subtracted degrees\sphinxhyphen{}of\sphinxhyphen{}freedom can be altered using the
\sphinxstyleemphasis{extra} option of the {\hyperref[\detokenize{compute_modify::doc}]{\sphinxcrossref{\DUrole{doc}{compute\_modify}}}} command.
By default this \sphinxstyleemphasis{extra} component is initialized to
\(n_\mathrm{dim}\) (as shown in the formula above) to represent the
degrees of freedom removed from a system due to its translation
invariance due to periodic boundary conditions.

\sphinxAtStartPar
A compute of this style with the ID of “thermo\_temp” is created when
LAMMPS starts up, as if this command were in the input script:

\begin{sphinxVerbatim}[commandchars=\\\{\}]
\PYG{k}{compute }\PYG{n+nv+nvIdentifier}{thermo\PYGZus{}temp}\PYG{+w}{ }\PYG{n+nv+nvIdentifier}{all}\PYG{+w}{ }\PYG{n}{temp}
\end{sphinxVerbatim}

\sphinxAtStartPar
See the “thermo\_style” command for more details.

\sphinxAtStartPar
See the {\hyperref[\detokenize{Howto_thermostat::doc}]{\sphinxcrossref{\DUrole{doc}{Howto thermostat}}}} page for a
discussion of different ways to compute temperature and perform
thermostatting.


\bigskip\hrule\bigskip


\sphinxAtStartPar
Styles with a \sphinxstyleemphasis{gpu}, \sphinxstyleemphasis{intel}, \sphinxstyleemphasis{kk}, \sphinxstyleemphasis{omp}, or \sphinxstyleemphasis{opt} suffix are
functionally the same as the corresponding style without the suffix.
They have been optimized to run faster, depending on your available
hardware, as discussed on the {\hyperref[\detokenize{Speed_packages::doc}]{\sphinxcrossref{\DUrole{doc}{Accelerator packages}}}}
page.  The accelerated styles take the same arguments and should
produce the same results, except for round\sphinxhyphen{}off and precision issues.

\sphinxAtStartPar
These accelerated styles are part of the GPU, INTEL, KOKKOS,
OPENMP, and OPT packages, respectively.  They are only enabled if
LAMMPS was built with those packages.  See the {\hyperref[\detokenize{Build_package::doc}]{\sphinxcrossref{\DUrole{doc}{Build package}}}} page for more info.

\sphinxAtStartPar
You can specify the accelerated styles explicitly in your input script
by including their suffix, or you can use the {\hyperref[\detokenize{Run_options::doc}]{\sphinxcrossref{\DUrole{doc}{\sphinxhyphen{}suffix command\sphinxhyphen{}line switch}}}} when you invoke LAMMPS, or you can use the
{\hyperref[\detokenize{suffix::doc}]{\sphinxcrossref{\DUrole{doc}{suffix}}}} command in your input script.

\sphinxAtStartPar
See the {\hyperref[\detokenize{Speed_packages::doc}]{\sphinxcrossref{\DUrole{doc}{Accelerator packages}}}} page for more
instructions on how to use the accelerated styles effectively.


\bigskip\hrule\bigskip



\subsection{Output info}
\label{\detokenize{compute_temp:output-info}}
\sphinxAtStartPar
This compute calculates a global scalar (the temperature) and a global
vector of length six (symmetric tensor), which can be accessed by
indices 1\textendash{}6.  These values can be used by any command that uses
global scalar or vector values from a compute as input.  See the
{\hyperref[\detokenize{Howto_output::doc}]{\sphinxcrossref{\DUrole{doc}{Howto output}}}} page for an overview of LAMMPS
output options.

\sphinxAtStartPar
The scalar value calculated by this compute is “intensive”.  The
vector values are “extensive”.

\sphinxAtStartPar
The scalar value is in temperature {\hyperref[\detokenize{units::doc}]{\sphinxcrossref{\DUrole{doc}{units}}}}.  The vector
values are in energy {\hyperref[\detokenize{units::doc}]{\sphinxcrossref{\DUrole{doc}{units}}}}.


\subsection{Restrictions}
\label{\detokenize{compute_temp:restrictions}}\begin{quote}

\sphinxAtStartPar
none
\end{quote}


\subsection{Related commands}
\label{\detokenize{compute_temp:related-commands}}
\sphinxAtStartPar
{\hyperref[\detokenize{compute_temp_partial::doc}]{\sphinxcrossref{\DUrole{doc}{compute temp/partial}}}},
{\hyperref[\detokenize{compute_temp_region::doc}]{\sphinxcrossref{\DUrole{doc}{compute temp/region}}}},
{\hyperref[\detokenize{compute_pressure::doc}]{\sphinxcrossref{\DUrole{doc}{compute pressure}}}}


\subsection{Default}
\label{\detokenize{compute_temp:default}}
\sphinxAtStartPar
none

\sphinxstepscope

\index{compute temp/asphere@\spxentry{compute temp/asphere}}\ignorespaces 

\section{compute temp/asphere command}
\label{\detokenize{compute_temp_asphere:compute-temp-asphere-command}}\label{\detokenize{compute_temp_asphere:index-0}}\label{\detokenize{compute_temp_asphere::doc}}

\subsection{Syntax}
\label{\detokenize{compute_temp_asphere:syntax}}
\begin{sphinxVerbatim}[commandchars=\\\{\}]
\PYG{k}{compute }\PYG{n+nv+nvIdentifier}{ID}\PYG{+w}{ }\PYG{n+nv+nvIdentifier}{group\PYGZhy{}ID}\PYG{+w}{ }\PYG{n}{temp}\PYG{o}{/}\PYG{n}{asphere}\PYG{+w}{ }\PYG{n}{keyword}\PYG{+w}{ }\PYG{n}{value}\PYG{+w}{ }\PYG{n}{...}
\end{sphinxVerbatim}
\begin{itemize}
\item {} 
\sphinxAtStartPar
ID, group\sphinxhyphen{}ID are documented in {\hyperref[\detokenize{compute::doc}]{\sphinxcrossref{\DUrole{doc}{compute}}}} command

\item {} 
\sphinxAtStartPar
temp/asphere = style name of this compute command

\item {} 
\sphinxAtStartPar
zero or more keyword/value pairs may be appended

\item {} 
\sphinxAtStartPar
keyword = \sphinxstyleemphasis{bias} or \sphinxstyleemphasis{dof}
\begin{sphinxalltt}
\sphinxstyleemphasis{bias} value = bias\sphinxhyphen{}ID
  bias\sphinxhyphen{}ID = ID of a temperature compute that removes a velocity bias
\sphinxstyleemphasis{dof} value = \sphinxstyleemphasis{all} or \sphinxstyleemphasis{rotate}
  all = compute temperature of translational and rotational degrees of freedom
  rotate = compute temperature of just rotational degrees of freedom
\end{sphinxalltt}

\end{itemize}


\subsection{Examples}
\label{\detokenize{compute_temp_asphere:examples}}
\begin{sphinxVerbatim}[commandchars=\\\{\}]
\PYG{k}{compute }\PYG{n+nv+nvIdentifier}{1}\PYG{+w}{ }\PYG{n+nv+nvIdentifier}{all}\PYG{+w}{ }\PYG{n}{temp}\PYG{o}{/}\PYG{n}{asphere}
\PYG{k}{compute }\PYG{n+nv+nvIdentifier}{myTemp}\PYG{+w}{ }\PYG{n+nv+nvIdentifier}{mobile}\PYG{+w}{ }\PYG{n}{temp}\PYG{o}{/}\PYG{n}{asphere}\PYG{+w}{ }\PYG{n}{bias}\PYG{+w}{ }\PYG{n}{tempCOM}
\PYG{k}{compute }\PYG{n+nv+nvIdentifier}{myTemp}\PYG{+w}{ }\PYG{n+nv+nvIdentifier}{mobile}\PYG{+w}{ }\PYG{n}{temp}\PYG{o}{/}\PYG{n}{asphere}\PYG{+w}{ }\PYG{n}{dof}\PYG{+w}{ }\PYG{n}{rotate}
\end{sphinxVerbatim}


\subsection{Description}
\label{\detokenize{compute_temp_asphere:description}}
\sphinxAtStartPar
Define a computation that calculates the temperature of a group of
aspherical particles, including a contribution from both their
translational and rotational kinetic energy.  This differs from the
usual {\hyperref[\detokenize{compute_temp::doc}]{\sphinxcrossref{\DUrole{doc}{compute temp}}}} command, which assumes point
particles with only translational kinetic energy.

\sphinxAtStartPar
Only finite\sphinxhyphen{}size particles (aspherical or spherical) can be included in
the group.  For 3d finite\sphinxhyphen{}size particles, each has six degrees of
freedom (three translational, three rotational).  For 2d finite\sphinxhyphen{}size
particles, each has three degrees of freedom (two translational, one
rotational).

\begin{sphinxadmonition}{note}{Note:}
\sphinxAtStartPar
This choice for degrees of freedom (DOF) assumes that all finite\sphinxhyphen{}size
aspherical or spherical particles in your model will freely rotate,
sampling all their rotational DOF.  It is possible to use a
combination of interaction potentials and fixes that induce no torque
or otherwise constrain some of all of your particles so that this is
not the case.  Then there are fewer DOF and you should use the
{\hyperref[\detokenize{compute_modify::doc}]{\sphinxcrossref{\DUrole{doc}{compute\_modify extra/dof}}}} command to adjust
the DOF accordingly.
\end{sphinxadmonition}

\sphinxAtStartPar
For example, an aspherical particle with all three of its shape
parameters the same is a sphere.  If it does not rotate, then it should
have 3 DOF instead of 6 in 3d (or two instead of three in 2d).  A
uniaxial aspherical particle has two of its three shape parameters the
same.  If it does not rotate around the axis perpendicular to its
circular cross section, then it should have 5 DOF instead of 6 in 3d.
The latter is the case for uniaxial ellipsoids in a {\hyperref[\detokenize{pair_gayberne::doc}]{\sphinxcrossref{\DUrole{doc}{GayBerne model}}}} since there is no induced torque around the optical
axis.  It will also be the case for biaxial ellipsoids when exactly two
of the semiaxes have the same length and the corresponding relative well
depths are equal.

\sphinxAtStartPar
The translational kinetic energy is computed the same as is described by
the {\hyperref[\detokenize{compute_temp::doc}]{\sphinxcrossref{\DUrole{doc}{compute temp}}}} command.  The rotational kinetic
energy is computed as \(\frac12 I \omega^2\), where \(I\) is the
inertia tensor for the aspherical particle and \(\omega\) is its
angular velocity, which is computed from its angular momentum.

\begin{sphinxadmonition}{note}{Note:}
\sphinxAtStartPar
For {\hyperref[\detokenize{dimension::doc}]{\sphinxcrossref{\DUrole{doc}{2d models}}}}, particles are treated as
ellipsoids, not ellipses, meaning their moments of inertia will be
the same as in 3d.
\end{sphinxadmonition}

\sphinxAtStartPar
A kinetic energy tensor, stored as a six\sphinxhyphen{}element vector, is also
calculated by this compute.  The formula for the components of the
tensor is the same as the above formula, except that \(v^2\) and
\(\omega^2\) are replaced by \(v_x v_y\) and \(\omega_x
\omega_y\) for the \(xy\) component, and the appropriate elements of
the moment of inertia tensor are used.  The six components of the vector
are ordered \(xx\), \(yy\), \(zz\), \(xy\), \(xz\),
\(yz\).

\sphinxAtStartPar
A symmetric tensor, stored as a six\sphinxhyphen{}element vector, is also calculated
by this compute for use in the computation of a pressure tensor by the
{\hyperref[\detokenize{compute_pressure::doc}]{\sphinxcrossref{\DUrole{doc}{compute pressue}}}} command.  The formula for the
components of the tensor is the same as the above expression for
\(E_\mathrm{kin}\), except that the 1/2 factor is NOT included and
the \(v_i^2\) and \(\omega^2\) are replaced by \(v_x v_y\)
and \(\omega_x \omega_y\) for the \(xy\) component, and so on.
And the appropriate elements of the moment of inertia tensor are used.
Note that because it lacks the 1/2 factor, these tensor components are
twice those of the traditional kinetic energy tensor.  The six
components of the vector are ordered \(xx\), \(yy\), \(zz\),
\(xy\), \(xz\), \(yz\).

\sphinxAtStartPar
The number of atoms contributing to the temperature is assumed to be
constant for the duration of the run; use the \sphinxstyleemphasis{dynamic/dof} option of
the {\hyperref[\detokenize{compute_modify::doc}]{\sphinxcrossref{\DUrole{doc}{compute\_modify}}}} command if this is not the
case.

\sphinxAtStartPar
This compute subtracts out translational degrees\sphinxhyphen{}of\sphinxhyphen{}freedom due to fixes
that constrain molecular motion, such as {\hyperref[\detokenize{fix_shake::doc}]{\sphinxcrossref{\DUrole{doc}{fix shake}}}}
and {\hyperref[\detokenize{fix_rigid::doc}]{\sphinxcrossref{\DUrole{doc}{fix rigid}}}}.  This means the temperature of groups
of atoms that include these constraints will be computed correctly.  If
needed, the subtracted degrees\sphinxhyphen{}of\sphinxhyphen{}freedom can be altered using the
\sphinxstyleemphasis{extra/dof} option of the {\hyperref[\detokenize{compute_modify::doc}]{\sphinxcrossref{\DUrole{doc}{compute\_modify}}}}
command.

\sphinxAtStartPar
See the {\hyperref[\detokenize{Howto_thermostat::doc}]{\sphinxcrossref{\DUrole{doc}{Howto thermostat}}}} page for a
discussion of different ways to compute temperature and perform
thermostatting.


\bigskip\hrule\bigskip


\sphinxAtStartPar
The keyword/value option pairs are used in the following ways.

\sphinxAtStartPar
For the \sphinxstyleemphasis{bias} keyword, \sphinxstyleemphasis{bias\sphinxhyphen{}ID} refers to the ID of a temperature
compute that removes a “bias” velocity from each atom.  This allows
compute temp/sphere to compute its thermal temperature after the
translational kinetic energy components have been altered in a
prescribed way (e.g., to remove a flow velocity profile).  Thermostats
that use this compute will work with this bias term.  See the doc
pages for individual computes that calculate a temperature and the doc
pages for fixes that perform thermostatting for more details.

\sphinxAtStartPar
For the \sphinxstyleemphasis{dof} keyword, a setting of \sphinxstyleemphasis{all} calculates a temperature
that includes both translational and rotational degrees of freedom.
A setting of \sphinxstyleemphasis{rotate} calculates a temperature that includes only
rotational degrees of freedom.


\bigskip\hrule\bigskip



\subsection{Output info}
\label{\detokenize{compute_temp_asphere:output-info}}
\sphinxAtStartPar
This compute calculates a global scalar (the temperature) and a global
vector of length 6 (symmetric tensor), which can be accessed by indices
1\textendash{}6.  These values can be used by any command that uses global scalar
or vector values from a compute as input.  See the {\hyperref[\detokenize{Howto_output::doc}]{\sphinxcrossref{\DUrole{doc}{Howto output}}}} page for an overview of LAMMPS output options.

\sphinxAtStartPar
The scalar value calculated by this compute is “intensive”.  The vector
values are “extensive”.

\sphinxAtStartPar
The scalar value is in temperature {\hyperref[\detokenize{units::doc}]{\sphinxcrossref{\DUrole{doc}{units}}}}.  The vector
values are in energy {\hyperref[\detokenize{units::doc}]{\sphinxcrossref{\DUrole{doc}{units}}}}.


\subsection{Restrictions}
\label{\detokenize{compute_temp_asphere:restrictions}}
\sphinxAtStartPar
This compute is part of the ASPHERE package.  It is only enabled if
LAMMPS was built with that package.  See the {\hyperref[\detokenize{Build_package::doc}]{\sphinxcrossref{\DUrole{doc}{Build package}}}} page for more info.

\sphinxAtStartPar
This compute requires that atoms store angular momentum and a quaternion
as defined by the {\hyperref[\detokenize{atom_style::doc}]{\sphinxcrossref{\DUrole{doc}{atom\_style ellipsoid}}}} command.

\sphinxAtStartPar
All particles in the group must be finite\sphinxhyphen{}size.  They cannot be point
particles, but they can be aspherical or spherical as defined by their
shape attribute.


\subsection{Related commands}
\label{\detokenize{compute_temp_asphere:related-commands}}
\sphinxAtStartPar
{\hyperref[\detokenize{compute_temp::doc}]{\sphinxcrossref{\DUrole{doc}{compute temp}}}}


\subsection{Default}
\label{\detokenize{compute_temp_asphere:default}}
\sphinxAtStartPar
none

\sphinxstepscope

\index{compute temp/body@\spxentry{compute temp/body}}\ignorespaces 

\section{compute temp/body command}
\label{\detokenize{compute_temp_body:compute-temp-body-command}}\label{\detokenize{compute_temp_body:index-0}}\label{\detokenize{compute_temp_body::doc}}

\subsection{Syntax}
\label{\detokenize{compute_temp_body:syntax}}
\begin{sphinxVerbatim}[commandchars=\\\{\}]
\PYG{k}{compute }\PYG{n+nv+nvIdentifier}{ID}\PYG{+w}{ }\PYG{n+nv+nvIdentifier}{group\PYGZhy{}ID}\PYG{+w}{ }\PYG{n}{temp}\PYG{o}{/}\PYG{n}{body}\PYG{+w}{ }\PYG{n}{keyword}\PYG{+w}{ }\PYG{n}{value}\PYG{+w}{ }\PYG{n}{...}
\end{sphinxVerbatim}
\begin{itemize}
\item {} 
\sphinxAtStartPar
ID, group\sphinxhyphen{}ID are documented in {\hyperref[\detokenize{compute::doc}]{\sphinxcrossref{\DUrole{doc}{compute}}}} command

\item {} 
\sphinxAtStartPar
temp/body = style name of this compute command

\item {} 
\sphinxAtStartPar
zero or more keyword/value pairs may be appended

\item {} 
\sphinxAtStartPar
keyword = \sphinxstyleemphasis{bias} or \sphinxstyleemphasis{dof}
\begin{sphinxalltt}
\sphinxstyleemphasis{bias} value = bias\sphinxhyphen{}ID
  bias\sphinxhyphen{}ID = ID of a temperature compute that removes a velocity bias
\sphinxstyleemphasis{dof} value = \sphinxstyleemphasis{all} or \sphinxstyleemphasis{rotate}
  all = compute temperature of translational and rotational degrees of freedom
  rotate = compute temperature of just rotational degrees of freedom
\end{sphinxalltt}

\end{itemize}


\subsection{Examples}
\label{\detokenize{compute_temp_body:examples}}
\begin{sphinxVerbatim}[commandchars=\\\{\}]
\PYG{k}{compute }\PYG{n+nv+nvIdentifier}{1}\PYG{+w}{ }\PYG{n+nv+nvIdentifier}{all}\PYG{+w}{ }\PYG{n}{temp}\PYG{o}{/}\PYG{n}{body}
\PYG{k}{compute }\PYG{n+nv+nvIdentifier}{myTemp}\PYG{+w}{ }\PYG{n+nv+nvIdentifier}{mobile}\PYG{+w}{ }\PYG{n}{temp}\PYG{o}{/}\PYG{n}{body}\PYG{+w}{ }\PYG{n}{bias}\PYG{+w}{ }\PYG{n}{tempCOM}
\PYG{k}{compute }\PYG{n+nv+nvIdentifier}{myTemp}\PYG{+w}{ }\PYG{n+nv+nvIdentifier}{mobile}\PYG{+w}{ }\PYG{n}{temp}\PYG{o}{/}\PYG{n}{body}\PYG{+w}{ }\PYG{n}{dof}\PYG{+w}{ }\PYG{n}{rotate}
\end{sphinxVerbatim}


\subsection{Description}
\label{\detokenize{compute_temp_body:description}}
\sphinxAtStartPar
Define a computation that calculates the temperature of a group of
body particles, including a contribution from both their
translational and rotational kinetic energy.  This differs from the
usual {\hyperref[\detokenize{compute_temp::doc}]{\sphinxcrossref{\DUrole{doc}{compute temp}}}} command, which assumes point
particles with only translational kinetic energy.

\sphinxAtStartPar
Only body particles can be included in the group.  For 3d particles,
each has 6 degrees of freedom (3 translational, 3 rotational).  For 2d
body particles, each has 3 degrees of freedom (2 translational, 1
rotational).

\begin{sphinxadmonition}{note}{Note:}
\sphinxAtStartPar
This choice for degrees of freedom (DOF) assumes that all body
particles in your model will freely rotate, sampling all their
rotational DOF.  It is possible to use a combination of interaction
potentials and fixes that induce no torque or otherwise constrain
some of all of your particles so that this is not the case.  Then
there are less DOF and you should use the {\hyperref[\detokenize{compute_modify::doc}]{\sphinxcrossref{\DUrole{doc}{compute\_modify
extra/dof}}}} command to adjust the DOF accordingly.
\end{sphinxadmonition}

\sphinxAtStartPar
The translational kinetic energy is computed the same as is described
by the {\hyperref[\detokenize{compute_temp::doc}]{\sphinxcrossref{\DUrole{doc}{compute temp}}}} command.  The rotational
kinetic energy is computed as \(\frac12 I \omega^2\), where \(I\)
is the moment of inertia tensor for the aspherical particle and \(\omega\)
is its angular velocity, which is computed from its angular momentum.

\sphinxAtStartPar
A symmetric tensor, stored as a six\sphinxhyphen{}element vector, is also calculated
by this compute for use in the computation of a pressure tensor by the
{\hyperref[\detokenize{compute_pressure::doc}]{\sphinxcrossref{\DUrole{doc}{compute pressue}}}} command.  The formula for
the components of the tensor is the same as the above expression for
\(E_\mathrm{kin}\), except that the 1/2 factor is NOT included and
the \(v_i^2\) and \(\omega^2\) are replaced by \(v_x v_y\)
and \(\omega_x \omega_y\) for the \(xy\) component, and so on.
And the appropriate elements of the moment of inertia tensor are used.
Note that because it lacks the 1/2 factor, these tensor components are
twice those of the traditional kinetic energy tensor.  The six
components of the vector are ordered \(xx\), \(yy\),
\(zz\), \(xy\), \(xz\), \(yz\).

\sphinxAtStartPar
The number of atoms contributing to the temperature is assumed to be
constant for the duration of the run; use the \sphinxstyleemphasis{dynamic/dof} option of
the {\hyperref[\detokenize{compute_modify::doc}]{\sphinxcrossref{\DUrole{doc}{compute\_modify}}}} command if this is not the
case.

\sphinxAtStartPar
This compute subtracts out translational degrees\sphinxhyphen{}of\sphinxhyphen{}freedom due to
fixes that constrain molecular motion, such as {\hyperref[\detokenize{fix_shake::doc}]{\sphinxcrossref{\DUrole{doc}{fix shake}}}}
and {\hyperref[\detokenize{fix_rigid::doc}]{\sphinxcrossref{\DUrole{doc}{fix rigid}}}}.  This means the
temperature of groups of atoms that include these constraints will be
computed correctly.  If needed, the subtracted degrees\sphinxhyphen{}of\sphinxhyphen{}freedom can
be altered using the \sphinxstyleemphasis{extra/dof} option of the
{\hyperref[\detokenize{compute_modify::doc}]{\sphinxcrossref{\DUrole{doc}{compute\_modify}}}} command.

\sphinxAtStartPar
See the {\hyperref[\detokenize{Howto_thermostat::doc}]{\sphinxcrossref{\DUrole{doc}{Howto thermostat}}}} page for a
discussion of different ways to compute temperature and perform
thermostatting.


\bigskip\hrule\bigskip


\sphinxAtStartPar
The keyword/value option pairs are used in the following ways.

\sphinxAtStartPar
For the \sphinxstyleemphasis{bias} keyword, \sphinxstyleemphasis{bias\sphinxhyphen{}ID} refers to the ID of a temperature
compute that removes a “bias” velocity from each atom.  This allows
compute temp/sphere to compute its thermal temperature after the
translational kinetic energy components have been altered in a
prescribed way (e.g., to remove a flow velocity profile).  Thermostats
that use this compute will work with this bias term.  See the doc
pages for individual computes that calculate a temperature and the doc
pages for fixes that perform thermostatting for more details.

\sphinxAtStartPar
For the \sphinxstyleemphasis{dof} keyword, a setting of \sphinxstyleemphasis{all} calculates a temperature
that includes both translational and rotational degrees of freedom.  A
setting of \sphinxstyleemphasis{rotate} calculates a temperature that includes only
rotational degrees of freedom.


\bigskip\hrule\bigskip



\subsection{Output info}
\label{\detokenize{compute_temp_body:output-info}}
\sphinxAtStartPar
This compute calculates a global scalar (the temperature) and a global
vector of length 6 (symmetric tensor), which can be accessed by
indices 1\textendash{}6.  These values can be used by any command that uses
global scalar or vector values from a compute as input.  See the
{\hyperref[\detokenize{Howto_output::doc}]{\sphinxcrossref{\DUrole{doc}{Howto output}}}} page for an overview of LAMMPS
output options.

\sphinxAtStartPar
The scalar value calculated by this compute is “intensive”.  The
vector values are “extensive”.

\sphinxAtStartPar
The scalar value is in temperature {\hyperref[\detokenize{units::doc}]{\sphinxcrossref{\DUrole{doc}{units}}}}.  The vector
values are in energy {\hyperref[\detokenize{units::doc}]{\sphinxcrossref{\DUrole{doc}{units}}}}.


\subsection{Restrictions}
\label{\detokenize{compute_temp_body:restrictions}}
\sphinxAtStartPar
This compute is part of the BODY package.  It is only enabled if
LAMMPS was built with that package.
See the {\hyperref[\detokenize{Build_package::doc}]{\sphinxcrossref{\DUrole{doc}{Build package}}}} page for more info.

\sphinxAtStartPar
This compute requires that atoms store angular momentum and a
quaternion as defined by the {\hyperref[\detokenize{atom_style::doc}]{\sphinxcrossref{\DUrole{doc}{atom\_style body}}}}
command.


\subsection{Related commands}
\label{\detokenize{compute_temp_body:related-commands}}
\sphinxAtStartPar
{\hyperref[\detokenize{compute_temp::doc}]{\sphinxcrossref{\DUrole{doc}{compute temp}}}}


\subsection{Default}
\label{\detokenize{compute_temp_body:default}}
\sphinxAtStartPar
none

\sphinxstepscope

\index{compute temp/chunk@\spxentry{compute temp/chunk}}\ignorespaces 

\section{compute temp/chunk command}
\label{\detokenize{compute_temp_chunk:compute-temp-chunk-command}}\label{\detokenize{compute_temp_chunk:index-0}}\label{\detokenize{compute_temp_chunk::doc}}

\subsection{Syntax}
\label{\detokenize{compute_temp_chunk:syntax}}
\begin{sphinxVerbatim}[commandchars=\\\{\}]
\PYG{k}{compute }\PYG{n+nv+nvIdentifier}{ID}\PYG{+w}{ }\PYG{n+nv+nvIdentifier}{group\PYGZhy{}ID}\PYG{+w}{ }\PYG{n}{temp}\PYG{o}{/}\PYG{n}{chunk}\PYG{+w}{ }\PYG{n}{chunkID}\PYG{+w}{ }\PYG{n}{value1}\PYG{+w}{ }\PYG{n}{value2}\PYG{+w}{ }\PYG{n}{...}\PYG{+w}{ }\PYG{n}{keyword}\PYG{+w}{ }\PYG{n}{value}\PYG{+w}{ }\PYG{n}{...}
\end{sphinxVerbatim}
\begin{itemize}
\item {} 
\sphinxAtStartPar
ID, group\sphinxhyphen{}ID are documented in {\hyperref[\detokenize{compute::doc}]{\sphinxcrossref{\DUrole{doc}{compute}}}} command

\item {} 
\sphinxAtStartPar
temp/chunk = style name of this compute command

\item {} 
\sphinxAtStartPar
chunkID = ID of {\hyperref[\detokenize{compute_chunk_atom::doc}]{\sphinxcrossref{\DUrole{doc}{compute chunk/atom}}}} command

\item {} 
\sphinxAtStartPar
zero or more values can be listed as value1,value2,etc.

\item {} 
\sphinxAtStartPar
value = \sphinxstyleemphasis{temp} or \sphinxstyleemphasis{kecom} or \sphinxstyleemphasis{internal}

\begin{sphinxVerbatim}[commandchars=\\\{\}]
temp = temperature of each chunk
kecom = kinetic energy of each chunk based on velocity of center of mass
internal = internal kinetic energy of each chunk
\end{sphinxVerbatim}

\item {} 
\sphinxAtStartPar
zero or more keyword/value pairs may be appended

\item {} 
\sphinxAtStartPar
keyword = \sphinxstyleemphasis{com} or \sphinxstyleemphasis{bias} or \sphinxstyleemphasis{adof} or \sphinxstyleemphasis{cdof}
\begin{sphinxalltt}
\sphinxstyleemphasis{com} value = \sphinxstyleemphasis{yes} or \sphinxstyleemphasis{no}
  yes = subtract center\sphinxhyphen{}of\sphinxhyphen{}mass velocity from each chunk before calculating temperature
  no = do not subtract center\sphinxhyphen{}of\sphinxhyphen{}mass velocity
\sphinxstyleemphasis{bias} value = bias\sphinxhyphen{}ID
  bias\sphinxhyphen{}ID = ID of a temperature compute that removes a velocity bias
\sphinxstyleemphasis{adof} value = dof\_per\_atom
  dof\_per\_atom = define this many degrees\sphinxhyphen{}of\sphinxhyphen{}freedom per atom
\sphinxstyleemphasis{cdof} value = dof\_per\_chunk
  dof\_per\_chunk = define this many degrees\sphinxhyphen{}of\sphinxhyphen{}freedom per chunk
\end{sphinxalltt}

\end{itemize}


\subsection{Examples}
\label{\detokenize{compute_temp_chunk:examples}}
\begin{sphinxVerbatim}[commandchars=\\\{\}]
\PYG{k}{compute }\PYG{n+nv+nvIdentifier}{1}\PYG{+w}{ }\PYG{n+nv+nvIdentifier}{fluid}\PYG{+w}{ }\PYG{n}{temp}\PYG{o}{/}\PYG{n}{chunk}\PYG{+w}{ }\PYG{n}{molchunk}
\PYG{k}{compute }\PYG{n+nv+nvIdentifier}{1}\PYG{+w}{ }\PYG{n+nv+nvIdentifier}{fluid}\PYG{+w}{ }\PYG{n}{temp}\PYG{o}{/}\PYG{n}{chunk}\PYG{+w}{ }\PYG{n}{molchunk}\PYG{+w}{ }\PYG{n}{temp}\PYG{+w}{ }\PYG{n}{internal}
\PYG{k}{compute }\PYG{n+nv+nvIdentifier}{1}\PYG{+w}{ }\PYG{n+nv+nvIdentifier}{fluid}\PYG{+w}{ }\PYG{n}{temp}\PYG{o}{/}\PYG{n}{chunk}\PYG{+w}{ }\PYG{n}{molchunk}\PYG{+w}{ }\PYG{n}{bias}\PYG{+w}{ }\PYG{n}{tpartial}\PYG{+w}{ }\PYG{n}{adof}\PYG{+w}{ }\PYG{l+m}{2.0}
\end{sphinxVerbatim}


\subsection{Description}
\label{\detokenize{compute_temp_chunk:description}}
\sphinxAtStartPar
Define a computation that calculates the temperature of a group of
atoms that are also in chunks, after optionally subtracting out the
center\sphinxhyphen{}of\sphinxhyphen{}mass velocity of each chunk.  By specifying optional values,
it can also calculate the per\sphinxhyphen{}chunk temperature or energies of the
multiple chunks of atoms.

\sphinxAtStartPar
In LAMMPS, chunks are collections of atoms defined by a
{\hyperref[\detokenize{compute_chunk_atom::doc}]{\sphinxcrossref{\DUrole{doc}{compute chunk/atom}}}} command, which assigns each atom
to a single chunk (or no chunk).  The ID for this command is specified
as chunkID.  For example, a single chunk could be the atoms in a
molecule or atoms in a spatial bin.  See the
{\hyperref[\detokenize{compute_chunk_atom::doc}]{\sphinxcrossref{\DUrole{doc}{compute chunk/atom}}}} and
{\hyperref[\detokenize{Howto_chunk::doc}]{\sphinxcrossref{\DUrole{doc}{Howto chunk}}}}
doc pages for details of how chunks can be defined and examples of how
they can be used to measure properties of a system.

\sphinxAtStartPar
The temperature is calculated by the formula
\begin{equation*}
\begin{split}\text{KE} = \frac{\text{DOF}}{2} k_B T,\end{split}
\end{equation*}
\sphinxAtStartPar
where KE is the total kinetic energy of all atoms assigned to chunks
(sum of \(\frac12 m v^2\)), DOF is the the total number of degrees of
freedom for those atoms, \(k_B\) is Boltzmann constant, and \(T\) is the
absolute temperature.

\sphinxAtStartPar
The DOF is calculated as \(N\times\)\sphinxstyleemphasis{adof}
+ \(N_\text{chunk}\times\)\sphinxstyleemphasis{cdof},
where \(N\) is the number of atoms contributing to the kinetic energy,
\sphinxstyleemphasis{adof} is the number of degrees of freedom per atom, and
\sphinxstyleemphasis{cdof} is the number of degrees of freedom per chunk.
By default, \sphinxstyleemphasis{adof} = 2 or 3 = dimensionality of system, as set via the
{\hyperref[\detokenize{dimension::doc}]{\sphinxcrossref{\DUrole{doc}{dimension}}}} command, and \sphinxstyleemphasis{cdof} = 0.0.
This gives the usual formula for temperature.

\sphinxAtStartPar
A symmetric tensor, stored as a six\sphinxhyphen{}element vector, is also calculated
by this compute.  The formula for the components of the tensor is the
same as the above expression for \(E_\mathrm{kin}\), except that
the 1/2 factor is NOT included and the \(v_i^2\) is replaced by
\(v_{i,x} v_{i,y}\) for the \(xy\) component, and so on.  Note
that because it lacks the 1/2 factor, these tensor components are
twice those of the traditional kinetic energy tensor.  The six
components of the vector are ordered \(xx\), \(yy\),
\(zz\), \(xy\), \(xz\), \(yz\).

\sphinxAtStartPar
Note that the number of atoms contributing to the temperature is
calculated each time the temperature is evaluated since it is assumed
the atoms may be dynamically assigned to chunks.  Thus there is no
need to use the \sphinxstyleemphasis{dynamic} option of the
{\hyperref[\detokenize{compute_modify::doc}]{\sphinxcrossref{\DUrole{doc}{compute\_modify}}}} command for this compute style.

\sphinxAtStartPar
If any optional values are specified, then per\sphinxhyphen{}chunk quantities are
also calculated and stored in a global array, as described below.

\sphinxAtStartPar
The \sphinxstyleemphasis{temp} value calculates the temperature for each chunk by the
formula
\begin{equation*}
\begin{split}\text{KE} = \frac{\text{DOF}}{2} k_B T,\end{split}
\end{equation*}
\sphinxAtStartPar
where KE is the total kinetic energy of the chunk of atoms (sum of
\(\frac12 m v^2\)), DOF is the total number of degrees of freedom for all
atoms in the chunk, \(k_B\) is the Boltzmann constant, and \(T\) is the
absolute temperature.

\sphinxAtStartPar
The number of degrees of freedom (DOF) in this case is calculated as
\(N\times\)\sphinxstyleemphasis{adof} + \sphinxstyleemphasis{cdof}, where \(N\) is the number
of atoms in the chunk, \sphinxstyleemphasis{adof} is the number of degrees of freedom
per atom, and \sphinxstyleemphasis{cdof} is the number of degrees of freedom per
chunk.  By default, \sphinxstyleemphasis{cdof} = 2 or 3 = dimensionality of system, as set
via the {\hyperref[\detokenize{dimension::doc}]{\sphinxcrossref{\DUrole{doc}{dimension}}}} command, and \sphinxstyleemphasis{cdof} = 0.0.
This gives the usual formula for temperature.

\sphinxAtStartPar
The \sphinxstyleemphasis{kecom} value calculates the kinetic energy of each chunk as if
all its atoms were moving with the velocity of the center\sphinxhyphen{}of\sphinxhyphen{}mass of
the chunk.

\sphinxAtStartPar
The \sphinxstyleemphasis{internal} value calculates the internal kinetic energy of each
chunk.  The interal KE is summed over the atoms in the chunk using an
internal “thermal” velocity for each atom, which is its velocity minus
the center\sphinxhyphen{}of\sphinxhyphen{}mass velocity of the chunk.


\bigskip\hrule\bigskip


\sphinxAtStartPar
Note that currently the global and per\sphinxhyphen{}chunk temperatures calculated
by this compute only include translational degrees of freedom for each
atom.  No rotational degrees of freedom are included for finite\sphinxhyphen{}size
particles.  Also no degrees of freedom are subtracted for any velocity
bias or constraints that are applied, such as
{\hyperref[\detokenize{compute_temp_partial::doc}]{\sphinxcrossref{\DUrole{doc}{compute temp/partial}}}}, or
{\hyperref[\detokenize{fix_shake::doc}]{\sphinxcrossref{\DUrole{doc}{fix shake}}}} or {\hyperref[\detokenize{fix_rigid::doc}]{\sphinxcrossref{\DUrole{doc}{fix rigid}}}}.
This is because those degrees of
freedom (e.g., a constrained bond) could apply to sets of atoms that
are both included and excluded from a specific chunk, and hence the
concept is somewhat ill\sphinxhyphen{}defined.  In some cases, you can use the
\sphinxstyleemphasis{adof} and \sphinxstyleemphasis{cdof} keywords to adjust the calculated degrees of freedom
appropriately, as explained below.

\sphinxAtStartPar
Note that the per\sphinxhyphen{}chunk temperature calculated by this compute and the
{\hyperref[\detokenize{fix_ave_chunk::doc}]{\sphinxcrossref{\DUrole{doc}{fix ave/chunk temp}}}} command can be different.
This compute calculates the temperature for each chunk for a single
snapshot.  Fix ave/chunk can do that but can also time average those
values over many snapshots, or it can compute a temperature as if the
atoms in the chunk on different timesteps were collected together as
one set of atoms to calculate their temperature.  This compute allows
the center\sphinxhyphen{}of\sphinxhyphen{}mass velocity of each chunk to be subtracted before
calculating the temperature; fix ave/chunk does not.

\begin{sphinxadmonition}{note}{Note:}
\sphinxAtStartPar
Only atoms in the specified group contribute to the calculations performed
by this compute.  The {\hyperref[\detokenize{compute_chunk_atom::doc}]{\sphinxcrossref{\DUrole{doc}{compute chunk/atom}}}}
command defines its own group; atoms will have a chunk ID = 0 if they are
not in that group, signifying they are not assigned to a chunk, and will
thus also not contribute to this calculation.  You can specify the “all”
group for this command if you simply want to include atoms with non\sphinxhyphen{}zero
chunk IDs.
\end{sphinxadmonition}

\sphinxAtStartPar
The simplest way to output the per\sphinxhyphen{}chunk results of the compute
temp/chunk calculation to a file is to use the
{\hyperref[\detokenize{fix_ave_time::doc}]{\sphinxcrossref{\DUrole{doc}{fix ave/time}}}} command, for example:

\begin{sphinxVerbatim}[commandchars=\\\{\}]
\PYG{k}{compute }\PYG{n+nv+nvIdentifier}{cc1}\PYG{+w}{ }\PYG{n+nv+nvIdentifier}{all}\PYG{+w}{ }\PYG{n}{chunk}\PYG{o}{/}\PYG{n}{atom}\PYG{+w}{ }\PYG{n}{molecule}
\PYG{k}{compute }\PYG{n+nv+nvIdentifier}{myChunk}\PYG{+w}{ }\PYG{n+nv+nvIdentifier}{all}\PYG{+w}{ }\PYG{n}{temp}\PYG{o}{/}\PYG{n}{chunk}\PYG{+w}{ }\PYG{n}{cc1}\PYG{+w}{ }\PYG{n}{temp}
\PYG{k}{fix }\PYG{n+nv+nvIdentifier}{1}\PYG{+w}{ }\PYG{n+nv+nvIdentifier}{all}\PYG{+w}{ }\PYG{n}{ave}\PYG{o}{/}\PYG{n}{time}\PYG{+w}{ }\PYG{l+m}{100}\PYG{+w}{ }\PYG{l+m}{1}\PYG{+w}{ }\PYG{l+m}{100}\PYG{+w}{ }\PYG{n}{c\PYGZus{}myChunk[1]}\PYG{+w}{ }\PYG{n}{file}\PYG{+w}{ }\PYG{n}{tmp.out}\PYG{+w}{ }\PYG{n}{mode}\PYG{+w}{ }\PYG{n}{vector}
\end{sphinxVerbatim}


\bigskip\hrule\bigskip


\sphinxAtStartPar
The keyword/value option pairs are used in the following ways.

\sphinxAtStartPar
The \sphinxstyleemphasis{com} keyword can be used with a value of \sphinxstyleemphasis{yes} to subtract the
velocity of the center\sphinxhyphen{}of\sphinxhyphen{}mass for each chunk from the velocity of the
atoms in that chunk, before calculating either the global or per\sphinxhyphen{}chunk
temperature.  This can be useful if the atoms are streaming or
otherwise moving collectively, and you wish to calculate only the
thermal temperature.

\sphinxAtStartPar
For the \sphinxstyleemphasis{bias} keyword, \sphinxstyleemphasis{bias\sphinxhyphen{}ID} refers to the ID of a temperature
compute that removes a “bias” velocity from each atom.  This also
allows calculation of the global or per\sphinxhyphen{}chunk temperature using only
the thermal temperature of atoms in each chunk after the translational
kinetic energy components have been altered in a prescribed way
(e.g., to remove a velocity profile).  It also applies to the calculation
of the other per\sphinxhyphen{}chunk values, such as \sphinxstyleemphasis{kecom} or \sphinxstyleemphasis{internal}, which
involve the center\sphinxhyphen{}of\sphinxhyphen{}mass velocity of each chunk, which is calculated
after the velocity bias is removed from each atom.  Note that the
temperature compute will apply its bias globally to the entire system,
not on a per\sphinxhyphen{}chunk basis.

\sphinxAtStartPar
The \sphinxstyleemphasis{adof} and \sphinxstyleemphasis{cdof} keywords define the values used in the degree of
freedom (DOF) formulas used for the global or per\sphinxhyphen{}chunk temperature,
as described above.  They can be used to calculate a more appropriate
temperature for some kinds of chunks.  Here are three examples:

\sphinxAtStartPar
If spatially binned chunks contain some number of water molecules and
{\hyperref[\detokenize{fix_shake::doc}]{\sphinxcrossref{\DUrole{doc}{fix shake}}}} is used to make each molecule rigid, then
you could calculate a temperature with six degrees of freedom (DOF) (three
translational, three rotational) per molecule by setting \sphinxstyleemphasis{adof} to 2.0.

\sphinxAtStartPar
If {\hyperref[\detokenize{compute_temp_partial::doc}]{\sphinxcrossref{\DUrole{doc}{compute temp/partial}}}} is used with the
\sphinxstyleemphasis{bias} keyword to only allow the x component of velocity to contribute
to the temperature, then \sphinxstyleemphasis{adof} = 1.0 would be appropriate.

\sphinxAtStartPar
If each chunk consists of a large molecule, with some number of its
bonds constrained by {\hyperref[\detokenize{fix_shake::doc}]{\sphinxcrossref{\DUrole{doc}{fix shake}}}} or the entire molecule
by {\hyperref[\detokenize{fix_rigid::doc}]{\sphinxcrossref{\DUrole{doc}{fix rigid/small}}}}, \sphinxstyleemphasis{adof} = 0.0 and \sphinxstyleemphasis{cdof} could be
set to the remaining degrees of freedom for the entire molecule
(entire chunk in this case; i.e., 6 for 3d, or 3 for 2d, for a rigid
molecule).


\bigskip\hrule\bigskip



\subsection{Output info}
\label{\detokenize{compute_temp_chunk:output-info}}
\sphinxAtStartPar
This compute calculates a global scalar (the temperature) and a global
vector of length 6 (symmetric tensor), which can be accessed by
indices 1\textendash{}6.  These values can be used by any command that uses
global scalar or vector values from a compute as input.  See the
{\hyperref[\detokenize{Howto_output::doc}]{\sphinxcrossref{\DUrole{doc}{Howto output}}}} page for an overview of LAMMPS
output options.

\sphinxAtStartPar
This compute also optionally calculates a global array, if one or more
of the optional values are specified.  The number of rows in the array
is the number of chunks \sphinxstyleemphasis{Nchunk} as calculated by the specified
{\hyperref[\detokenize{compute_chunk_atom::doc}]{\sphinxcrossref{\DUrole{doc}{compute chunk/atom}}}} command.  The number of
columns is the number of specified values (1 or more).  These values
can be accessed by any command that uses global array values from a
compute as input.  Again, see the {\hyperref[\detokenize{Howto_output::doc}]{\sphinxcrossref{\DUrole{doc}{Howto output}}}} doc
page for an overview of LAMMPS output options.

\sphinxAtStartPar
The scalar value calculated by this compute is “intensive”.  The
vector values are “extensive”.  The array values are “intensive”.

\sphinxAtStartPar
The scalar value is in temperature {\hyperref[\detokenize{units::doc}]{\sphinxcrossref{\DUrole{doc}{units}}}}.  The vector
values are in energy {\hyperref[\detokenize{units::doc}]{\sphinxcrossref{\DUrole{doc}{units}}}}.  The array values will be
in temperature {\hyperref[\detokenize{units::doc}]{\sphinxcrossref{\DUrole{doc}{units}}}} for the \sphinxstyleemphasis{temp} value, and in
energy {\hyperref[\detokenize{units::doc}]{\sphinxcrossref{\DUrole{doc}{units}}}} for the \sphinxstyleemphasis{kecom} and \sphinxstyleemphasis{internal} values.


\subsection{Restrictions}
\label{\detokenize{compute_temp_chunk:restrictions}}
\sphinxAtStartPar
The \sphinxstyleemphasis{com} and \sphinxstyleemphasis{bias} keywords cannot be used together.


\subsection{Related commands}
\label{\detokenize{compute_temp_chunk:related-commands}}
\sphinxAtStartPar
{\hyperref[\detokenize{compute_temp::doc}]{\sphinxcrossref{\DUrole{doc}{compute temp}}}}, {\hyperref[\detokenize{fix_ave_chunk::doc}]{\sphinxcrossref{\DUrole{doc}{fix ave/chunk temp}}}}


\subsection{Default}
\label{\detokenize{compute_temp_chunk:default}}
\sphinxAtStartPar
The option defaults are com no, no bias, adof = dimensionality of the
system (2 or 3), and cdof = 0.0.

\sphinxstepscope

\index{compute temp/com@\spxentry{compute temp/com}}\ignorespaces 

\section{compute temp/com command}
\label{\detokenize{compute_temp_com:compute-temp-com-command}}\label{\detokenize{compute_temp_com:index-0}}\label{\detokenize{compute_temp_com::doc}}

\subsection{Syntax}
\label{\detokenize{compute_temp_com:syntax}}
\begin{sphinxVerbatim}[commandchars=\\\{\}]
\PYG{k}{compute }\PYG{n+nv+nvIdentifier}{ID}\PYG{+w}{ }\PYG{n+nv+nvIdentifier}{group\PYGZhy{}ID}\PYG{+w}{ }\PYG{n}{temp}\PYG{o}{/}\PYG{n}{com}
\end{sphinxVerbatim}
\begin{itemize}
\item {} 
\sphinxAtStartPar
ID, group\sphinxhyphen{}ID are documented in {\hyperref[\detokenize{compute::doc}]{\sphinxcrossref{\DUrole{doc}{compute}}}} command

\item {} 
\sphinxAtStartPar
temp/com = style name of this compute command

\end{itemize}


\subsection{Examples}
\label{\detokenize{compute_temp_com:examples}}
\begin{sphinxVerbatim}[commandchars=\\\{\}]
\PYG{k}{compute }\PYG{n+nv+nvIdentifier}{1}\PYG{+w}{ }\PYG{n+nv+nvIdentifier}{all}\PYG{+w}{ }\PYG{n}{temp}\PYG{o}{/}\PYG{n}{com}
\PYG{k}{compute }\PYG{n+nv+nvIdentifier}{myTemp}\PYG{+w}{ }\PYG{n+nv+nvIdentifier}{mobile}\PYG{+w}{ }\PYG{n}{temp}\PYG{o}{/}\PYG{n}{com}
\end{sphinxVerbatim}


\subsection{Description}
\label{\detokenize{compute_temp_com:description}}
\sphinxAtStartPar
Define a computation that calculates the temperature of a group of
atoms, after subtracting out the center\sphinxhyphen{}of\sphinxhyphen{}mass velocity of the group.
This is useful if the group is expected to have a non\sphinxhyphen{}zero net
velocity for some reason.  A compute of this style can be used by any
command that computes a temperature,
(e.g., {\hyperref[\detokenize{thermo_modify::doc}]{\sphinxcrossref{\DUrole{doc}{thermo\_modify}}}},
{\hyperref[\detokenize{fix_temp_rescale::doc}]{\sphinxcrossref{\DUrole{doc}{fix temp/rescale}}}}, {\hyperref[\detokenize{fix_nh::doc}]{\sphinxcrossref{\DUrole{doc}{fix npt}}}}).

\sphinxAtStartPar
After the center\sphinxhyphen{}of\sphinxhyphen{}mass velocity has been subtracted from each atom,
the temperature is calculated by the formula
\begin{equation*}
\begin{split}\text{KE} = \frac{\text{dim}}{2} N k_B T,\end{split}
\end{equation*}
\sphinxAtStartPar
where KE is the total kinetic energy of the group of atoms (sum of
\(\frac12 m v^2\)), dim = 2 or 3 is the dimensionality of the
simulation, \(N\) is number of atoms in the group, \(k_B\) is
the Boltzmann constant, and \(T\) is the absolute temperature.

\sphinxAtStartPar
A symmetric tensor, stored as a six\sphinxhyphen{}element vector, is also calculated
by this compute for use in the computation of a pressure tensor by the
{\hyperref[\detokenize{compute_pressure::doc}]{\sphinxcrossref{\DUrole{doc}{compute pressue}}}} command.  The formula for
the components of the tensor is the same as the above expression for
\(E_\mathrm{kin}\), except that the 1/2 factor is NOT included and
the \(v_i^2\) is replaced by \(v_{i,x} v_{i,y}\) for the
\(xy\) component, and so on.  Note that because it lacks the 1/2
factor, these tensor components are twice those of the traditional
kinetic energy tensor.  The six components of the vector are ordered
\(xx\), \(yy\), \(zz\), \(xy\), \(xz\),
\(yz\).

\sphinxAtStartPar
The number of atoms contributing to the temperature is assumed to be
constant for the duration of the run; use the \sphinxstyleemphasis{dynamic} option of the
{\hyperref[\detokenize{compute_modify::doc}]{\sphinxcrossref{\DUrole{doc}{compute\_modify}}}} command if this is not the case.

\sphinxAtStartPar
The removal of the center\sphinxhyphen{}of\sphinxhyphen{}mass velocity by this fix is essentially
computing the temperature after a “bias” has been removed from the
velocity of the atoms.  If this compute is used with a fix command
that performs thermostatting then this bias will be subtracted from
each atom, thermostatting of the remaining thermal velocity will be
performed, and the bias will be added back in.  Thermostatting fixes
that work in this way include {\hyperref[\detokenize{fix_nh::doc}]{\sphinxcrossref{\DUrole{doc}{fix nvt}}}},
{\hyperref[\detokenize{fix_temp_rescale::doc}]{\sphinxcrossref{\DUrole{doc}{fix temp/rescale}}}},
{\hyperref[\detokenize{fix_temp_berendsen::doc}]{\sphinxcrossref{\DUrole{doc}{fix temp/berendsen}}}}, and
{\hyperref[\detokenize{fix_langevin::doc}]{\sphinxcrossref{\DUrole{doc}{fix langevin}}}}.

\sphinxAtStartPar
This compute subtracts out degrees\sphinxhyphen{}of\sphinxhyphen{}freedom due to fixes that
constrain molecular motion, such as {\hyperref[\detokenize{fix_shake::doc}]{\sphinxcrossref{\DUrole{doc}{fix shake}}}} and
{\hyperref[\detokenize{fix_rigid::doc}]{\sphinxcrossref{\DUrole{doc}{fix rigid}}}}.  This means the temperature of groups of
atoms that include these constraints will be computed correctly.
If needed, the subtracted degrees\sphinxhyphen{}of\sphinxhyphen{}freedom can be altered using the
\sphinxstyleemphasis{extra} option of the {\hyperref[\detokenize{compute_modify::doc}]{\sphinxcrossref{\DUrole{doc}{compute\_modify}}}} command.

\sphinxAtStartPar
See the {\hyperref[\detokenize{Howto_thermostat::doc}]{\sphinxcrossref{\DUrole{doc}{Howto thermostat}}}} page for a
discussion of different ways to compute temperature and perform
thermostatting.


\subsection{Output info}
\label{\detokenize{compute_temp_com:output-info}}
\sphinxAtStartPar
This compute calculates a global scalar (the temperature) and a global
vector of length 6 (symmetric tensor), which can be accessed by
indices 1\textendash{}6.  These values can be used by any command that uses
global scalar or vector values from a compute as input.  See the
{\hyperref[\detokenize{Howto_output::doc}]{\sphinxcrossref{\DUrole{doc}{Howto output}}}} page for an overview of LAMMPS
output options.

\sphinxAtStartPar
The scalar value calculated by this compute is “intensive”.  The
vector values are “extensive”.

\sphinxAtStartPar
The scalar value is in temperature {\hyperref[\detokenize{units::doc}]{\sphinxcrossref{\DUrole{doc}{units}}}}.  The vector
values is in energy {\hyperref[\detokenize{units::doc}]{\sphinxcrossref{\DUrole{doc}{units}}}}.


\subsection{Restrictions}
\label{\detokenize{compute_temp_com:restrictions}}\begin{quote}

\sphinxAtStartPar
none
\end{quote}


\subsection{Related commands}
\label{\detokenize{compute_temp_com:related-commands}}
\sphinxAtStartPar
{\hyperref[\detokenize{compute_temp::doc}]{\sphinxcrossref{\DUrole{doc}{compute temp}}}}


\subsection{Default}
\label{\detokenize{compute_temp_com:default}}
\sphinxAtStartPar
none

\sphinxstepscope

\index{compute temp/cs@\spxentry{compute temp/cs}}\ignorespaces 

\section{compute temp/cs command}
\label{\detokenize{compute_temp_cs:compute-temp-cs-command}}\label{\detokenize{compute_temp_cs:index-0}}\label{\detokenize{compute_temp_cs::doc}}

\subsection{Syntax}
\label{\detokenize{compute_temp_cs:syntax}}
\begin{sphinxVerbatim}[commandchars=\\\{\}]
\PYG{k}{compute }\PYG{n+nv+nvIdentifier}{ID}\PYG{+w}{ }\PYG{n+nv+nvIdentifier}{group\PYGZhy{}ID}\PYG{+w}{ }\PYG{n}{temp}\PYG{o}{/}\PYG{n}{cs}\PYG{+w}{ }\PYG{n}{group1}\PYG{+w}{ }\PYG{n}{group2}
\end{sphinxVerbatim}
\begin{itemize}
\item {} 
\sphinxAtStartPar
ID, group\sphinxhyphen{}ID are documented in {\hyperref[\detokenize{compute::doc}]{\sphinxcrossref{\DUrole{doc}{compute}}}} command

\item {} 
\sphinxAtStartPar
temp/cs = style name of this compute command

\item {} 
\sphinxAtStartPar
group1 = group\sphinxhyphen{}ID of either cores or shells

\item {} 
\sphinxAtStartPar
group2 = group\sphinxhyphen{}ID of either shells or cores

\end{itemize}


\subsection{Examples}
\label{\detokenize{compute_temp_cs:examples}}
\begin{sphinxVerbatim}[commandchars=\\\{\}]
\PYG{k}{compute }\PYG{n+nv+nvIdentifier}{oxygen\PYGZus{}c\PYGZhy{}s}\PYG{+w}{ }\PYG{n+nv+nvIdentifier}{all}\PYG{+w}{ }\PYG{n}{temp}\PYG{o}{/}\PYG{n}{cs}\PYG{+w}{ }\PYG{n}{O\PYGZus{}core}\PYG{+w}{ }\PYG{n}{O\PYGZus{}shell}
\PYG{k}{compute }\PYG{n+nv+nvIdentifier}{core\PYGZus{}shells}\PYG{+w}{ }\PYG{n+nv+nvIdentifier}{all}\PYG{+w}{ }\PYG{n}{temp}\PYG{o}{/}\PYG{n}{cs}\PYG{+w}{ }\PYG{n}{cores}\PYG{+w}{ }\PYG{n}{shells}
\end{sphinxVerbatim}


\subsection{Description}
\label{\detokenize{compute_temp_cs:description}}
\sphinxAtStartPar
Define a computation that calculates the temperature of a system based
on the center\sphinxhyphen{}of\sphinxhyphen{}mass velocity of atom pairs that are bonded to each
other.  This compute is designed to be used with the adiabatic
core/shell model of {\hyperref[\detokenize{compute_temp_cs:mitchellfincham1}]{\sphinxcrossref{\DUrole{std,std-ref}{(Mitchell and Fincham)}}}}.
See the {\hyperref[\detokenize{Howto_coreshell::doc}]{\sphinxcrossref{\DUrole{doc}{Howto coreshell}}}} page for an overview of
the model as implemented in LAMMPS.  Specifically, this compute enables
correct temperature calculation and thermostatting of core/shell pairs
where it is desirable for the internal degrees of freedom of the
core/shell pairs to not be influenced by a thermostat.  A compute of
this style can be used by any command that computes a temperature via
{\hyperref[\detokenize{fix_modify::doc}]{\sphinxcrossref{\DUrole{doc}{fix\_modify}}}} (e.g., {\hyperref[\detokenize{fix_temp_rescale::doc}]{\sphinxcrossref{\DUrole{doc}{fix temp/rescale}}}}, {\hyperref[\detokenize{fix_nh::doc}]{\sphinxcrossref{\DUrole{doc}{fix npt}}}}).

\sphinxAtStartPar
Note that this compute does not require all ions to be polarized, hence
defined as core/shell pairs.  One can mix core/shell pairs and ions
without a satellite particle if desired. The compute will consider the
non\sphinxhyphen{}polarized ions according to the physical system.

\sphinxAtStartPar
For this compute, core and shell particles are specified by two
respective group IDs, which can be defined using the {\hyperref[\detokenize{group::doc}]{\sphinxcrossref{\DUrole{doc}{group}}}} command.  The number of atoms in the two groups must be the
same and there should be one bond defined between a pair of atoms in the
two groups.  Non\sphinxhyphen{}polarized ions which might also be included in the
treated system should not be included into either of these groups, they
are taken into account by the \sphinxstyleemphasis{group\sphinxhyphen{}ID} (second argument) of the
compute.

\sphinxAtStartPar
The temperature is calculated by the formula
\begin{equation*}
\begin{split}\text{KE} = \frac{\text{dim}}{2} N k_B T,\end{split}
\end{equation*}
\sphinxAtStartPar
where KE is the total kinetic energy of the group of atoms (sum of
\(\frac12 m v^2\)), dim = 2 or 3 is the dimensionality of the
simulation, \(N\) is the number of atoms in the group, \(k_B\)
is the Boltzmann constant, and \(T\) is the absolute temperature.
Note that the velocity of each core or shell atom used in the KE
calculation is the velocity of the center\sphinxhyphen{}of\sphinxhyphen{}mass (COM) of the
core/shell pair the atom is part of.

\sphinxAtStartPar
A symmetric tensor, stored as a six\sphinxhyphen{}element vector, is also calculated
by this compute for use in the computation of a pressure tensor by the
{\hyperref[\detokenize{compute_pressure::doc}]{\sphinxcrossref{\DUrole{doc}{compute pressue}}}} command.  The formula for the
components of the tensor is the same as the above expression for
\(E_\mathrm{kin}\), except that the 1/2 factor is NOT included and
the \(v_i^2\) is replaced by \(v_{i,x} v_{i,y}\) for the
\(xy\) component, and so on.  Note that because it lacks the 1/2
factor, these tensor components are twice those of the traditional
kinetic energy tensor.  The six components of the vector are ordered
\(xx\), \(yy\), \(zz\), \(xy\), \(xz\), \(yz\).

\sphinxAtStartPar
The change this fix makes to core/shell atom velocities is essentially
computing the temperature after a “bias” has been removed from the
velocity of the atoms.  This “bias” is the velocity of the atom relative
to the center\sphinxhyphen{}of\sphinxhyphen{}mass velocity of the core/shell pair.  If this compute
is used with a fix command that performs thermostatting then this bias
will be subtracted from each atom, thermostatting of the remaining
center\sphinxhyphen{}of\sphinxhyphen{}mass velocity will be performed, and the bias will be added
back in.  This means the thermostatting will effectively be performed on
the core/shell pairs, instead of on the individual core and shell atoms.
Thermostatting fixes that work in this way include {\hyperref[\detokenize{fix_nh::doc}]{\sphinxcrossref{\DUrole{doc}{fix nvt}}}}, {\hyperref[\detokenize{fix_temp_rescale::doc}]{\sphinxcrossref{\DUrole{doc}{fix temp/rescale}}}}, {\hyperref[\detokenize{fix_temp_berendsen::doc}]{\sphinxcrossref{\DUrole{doc}{fix
temp/berendsen}}}}, and {\hyperref[\detokenize{fix_langevin::doc}]{\sphinxcrossref{\DUrole{doc}{fix langevin}}}}.

\sphinxAtStartPar
The internal energy of core/shell pairs can be calculated by the
{\hyperref[\detokenize{compute_temp_chunk::doc}]{\sphinxcrossref{\DUrole{doc}{compute temp/chunk}}}} command, if chunks are
defined as core/shell pairs.  See the {\hyperref[\detokenize{Howto_coreshell::doc}]{\sphinxcrossref{\DUrole{doc}{Howto coreshell}}}} doc page for more discussion on how to do this.


\subsection{Output info}
\label{\detokenize{compute_temp_cs:output-info}}
\sphinxAtStartPar
This compute calculates a global scalar (the temperature) and a global
vector of length 6 (symmetric tensor), which can be accessed by indices
1\textendash{}6.  These values can be used by any command that uses global scalar
or vector values from a compute as input.

\sphinxAtStartPar
The scalar value calculated by this compute is “intensive”.  The vector
values are “extensive”.

\sphinxAtStartPar
The scalar value is in temperature {\hyperref[\detokenize{units::doc}]{\sphinxcrossref{\DUrole{doc}{units}}}}.  The vector
values are in energy {\hyperref[\detokenize{units::doc}]{\sphinxcrossref{\DUrole{doc}{units}}}}.


\subsection{Restrictions}
\label{\detokenize{compute_temp_cs:restrictions}}
\sphinxAtStartPar
The number of core/shell pairs contributing to the temperature is
assumed to be constant for the duration of the run.  No fixes should
be used which generate new molecules or atoms during a simulation.


\subsection{Related commands}
\label{\detokenize{compute_temp_cs:related-commands}}
\sphinxAtStartPar
{\hyperref[\detokenize{compute_temp::doc}]{\sphinxcrossref{\DUrole{doc}{compute temp}}}},
{\hyperref[\detokenize{compute_temp_chunk::doc}]{\sphinxcrossref{\DUrole{doc}{compute temp/chunk}}}}


\subsection{Default}
\label{\detokenize{compute_temp_cs:default}}
\sphinxAtStartPar
none


\bigskip\hrule\bigskip

\phantomsection\label{\detokenize{compute_temp_cs:mitchellfincham1}}
\sphinxAtStartPar
\sphinxstylestrong{(Mitchell and Fincham)} Mitchell, Fincham, J Phys Condensed Matter,
5, 1031\sphinxhyphen{}1038 (1993).

\sphinxstepscope

\index{compute temp/deform@\spxentry{compute temp/deform}}\ignorespaces 
\index{compute temp/deform/kk@\spxentry{compute temp/deform/kk}}\ignorespaces 

\section{compute temp/deform command}
\label{\detokenize{compute_temp_deform:compute-temp-deform-command}}\label{\detokenize{compute_temp_deform:index-1}}\label{\detokenize{compute_temp_deform:index-0}}\label{\detokenize{compute_temp_deform::doc}}
\sphinxAtStartPar
Accelerator Variants: \sphinxstyleemphasis{temp/deform/kk}


\subsection{Syntax}
\label{\detokenize{compute_temp_deform:syntax}}
\begin{sphinxVerbatim}[commandchars=\\\{\}]
\PYG{k}{compute }\PYG{n+nv+nvIdentifier}{ID}\PYG{+w}{ }\PYG{n+nv+nvIdentifier}{group\PYGZhy{}ID}\PYG{+w}{ }\PYG{n}{temp}\PYG{o}{/}\PYG{n}{deform}
\end{sphinxVerbatim}
\begin{itemize}
\item {} 
\sphinxAtStartPar
ID, group\sphinxhyphen{}ID are documented in {\hyperref[\detokenize{compute::doc}]{\sphinxcrossref{\DUrole{doc}{compute}}}} command

\item {} 
\sphinxAtStartPar
temp/deform = style name of this compute command

\end{itemize}


\subsection{Examples}
\label{\detokenize{compute_temp_deform:examples}}
\begin{sphinxVerbatim}[commandchars=\\\{\}]
\PYG{k}{compute }\PYG{n+nv+nvIdentifier}{myTemp}\PYG{+w}{ }\PYG{n+nv+nvIdentifier}{all}\PYG{+w}{ }\PYG{n}{temp}\PYG{o}{/}\PYG{n}{deform}
\end{sphinxVerbatim}


\subsection{Description}
\label{\detokenize{compute_temp_deform:description}}
\sphinxAtStartPar
Define a computation that calculates the temperature of a group of
atoms, after subtracting out a streaming velocity induced by the
simulation box changing size and/or shape, for example in a
non\sphinxhyphen{}equilibrium MD (NEMD) simulation.  The size/shape change is
induced by use of the {\hyperref[\detokenize{fix_deform::doc}]{\sphinxcrossref{\DUrole{doc}{fix deform}}}} command.  A compute
of this style is created by the {\hyperref[\detokenize{fix_nvt_sllod::doc}]{\sphinxcrossref{\DUrole{doc}{fix nvt/sllod}}}}
command to compute the thermal temperature of atoms for thermostatting
purposes.  A compute of this style can also be used by any command
that computes a temperature (e.g., {\hyperref[\detokenize{thermo_modify::doc}]{\sphinxcrossref{\DUrole{doc}{thermo\_modify}}}},
{\hyperref[\detokenize{fix_temp_rescale::doc}]{\sphinxcrossref{\DUrole{doc}{fix temp/rescale}}}}, {\hyperref[\detokenize{fix_nh::doc}]{\sphinxcrossref{\DUrole{doc}{fix npt}}}}).

\sphinxAtStartPar
The deformation fix changes the box size and/or shape over time, so
each atom in the simulation box can be thought of as having a
“streaming” velocity.  For example, if the box is being sheared in \sphinxstyleemphasis{x},
relative to \sphinxstyleemphasis{y}, then atoms at the bottom of the box (low \sphinxstyleemphasis{y}) have a
small \sphinxstyleemphasis{x} velocity, while atoms at the top of the box (high \sphinxstyleemphasis{y}) have a
large \sphinxstyleemphasis{x} velocity.  This position\sphinxhyphen{}dependent streaming velocity is
subtracted from each atom’s actual velocity to yield a thermal
velocity, which is then used to compute the temperature.

\begin{sphinxadmonition}{note}{Note:}
\sphinxAtStartPar
{\hyperref[\detokenize{fix_deform::doc}]{\sphinxcrossref{\DUrole{doc}{Fix deform}}}} has an option for remapping either
atom coordinates or velocities to the changing simulation box.  When
using this compute in conjunction with a deforming box, fix deform
should NOT remap atom positions, but rather should let atoms respond
to the changing box by adjusting their own velocities (or let
{\hyperref[\detokenize{fix_deform::doc}]{\sphinxcrossref{\DUrole{doc}{fix deform}}}} remap the atom velocities; see its remap
option).  If fix deform does remap atom positions, then they appear to
move with the box but their velocity is not changed, and thus they do
NOT have the streaming velocity assumed by this compute.  LAMMPS will
warn you if fix deform is defined and its remap setting is not
consistent with this compute.
\end{sphinxadmonition}

\sphinxAtStartPar
After the streaming velocity has been subtracted from each atom, the
temperature is calculated by the formula
\begin{equation*}
\begin{split}\text{KE} = \frac{\text{dim}}{2} N k_B T,\end{split}
\end{equation*}
\sphinxAtStartPar
where KE is the total kinetic energy of the group of atoms (sum of
\(\frac12 m v^2\), dim = 2 or 3 is the dimensionality of the
simulation, \(N\) is the number of atoms in the group, \(k_B\)
is the Boltzmann constant, and \(T\) is the temperature.  Note that
\(v\) in the kinetic energy formula is the atom’s velocity.

\sphinxAtStartPar
A symmetric tensor, stored as a six\sphinxhyphen{}element vector, is also calculated
by this compute for use in the computation of a pressure tensor by the
{\hyperref[\detokenize{compute_pressure::doc}]{\sphinxcrossref{\DUrole{doc}{compute pressue}}}} command.  The formula for
the components of the tensor is the same as the above expression for
\(E_\mathrm{kin}\), except that the 1/2 factor is NOT included and
the \(v_i^2\) is replaced by \(v_{i,x} v_{i,y}\) for the
\(xy\) component, and so on.  Note that because it lacks the 1/2
factor, these tensor components are twice those of the traditional
kinetic energy tensor.  The six components of the vector are ordered
\(xx\), \(yy\), \(zz\), \(xy\), \(xz\),
\(yz\).

\sphinxAtStartPar
The number of atoms contributing to the temperature is assumed to be
constant for the duration of the run; use the \sphinxstyleemphasis{dynamic} option of the
{\hyperref[\detokenize{compute_modify::doc}]{\sphinxcrossref{\DUrole{doc}{compute\_modify}}}} command if this is not the case.

\sphinxAtStartPar
The removal of the box deformation velocity component by this fix is
essentially computing the temperature after a “bias” has been removed
from the velocity of the atoms.  If this compute is used with a fix
command that performs thermostatting then this bias will be subtracted
from each atom, thermostatting of the remaining thermal velocity will
be performed, and the bias will be added back in.  Thermostatting
fixes that work in this way include {\hyperref[\detokenize{fix_nh::doc}]{\sphinxcrossref{\DUrole{doc}{fix nvt}}}},
{\hyperref[\detokenize{fix_temp_rescale::doc}]{\sphinxcrossref{\DUrole{doc}{fix temp/rescale}}}},
{\hyperref[\detokenize{fix_temp_berendsen::doc}]{\sphinxcrossref{\DUrole{doc}{fix temp/berendsen}}}}, and
{\hyperref[\detokenize{fix_langevin::doc}]{\sphinxcrossref{\DUrole{doc}{fix langevin}}}}.

\begin{sphinxadmonition}{note}{Note:}
\sphinxAtStartPar
The temperature calculated by this compute is only accurate if
the atoms are indeed moving with a stream velocity profile that
matches the box deformation.  If not, then the compute will subtract
off an incorrect stream velocity, yielding a bogus thermal
temperature.  You should \sphinxstylestrong{not} assume that your atoms are streaming at
the same rate the box is deforming.  Rather, you should monitor their
velocity profiles (e.g., via the {\hyperref[\detokenize{fix_ave_chunk::doc}]{\sphinxcrossref{\DUrole{doc}{fix ave/chunk}}}}
command).  You can also compare the results of this compute to
{\hyperref[\detokenize{compute_temp_profile::doc}]{\sphinxcrossref{\DUrole{doc}{compute temp/profile}}}}, which actually
calculates the stream profile before subtracting it.  If the two computes do
not give roughly the same temperature, then your atoms are not streaming
consistent with the box deformation.  See the {\hyperref[\detokenize{fix_deform::doc}]{\sphinxcrossref{\DUrole{doc}{fix deform}}}}
command for more details on ways to get atoms to stream consistently with
the box deformation.
\end{sphinxadmonition}

\sphinxAtStartPar
This compute subtracts out degrees\sphinxhyphen{}of\sphinxhyphen{}freedom due to fixes that
constrain molecular motion, such as {\hyperref[\detokenize{fix_shake::doc}]{\sphinxcrossref{\DUrole{doc}{fix shake}}}} and
{\hyperref[\detokenize{fix_rigid::doc}]{\sphinxcrossref{\DUrole{doc}{fix rigid}}}}.  This means the temperature of groups of
atoms that include these constraints will be computed correctly.  If
needed, the subtracted degrees\sphinxhyphen{}of\sphinxhyphen{}freedom can be altered using the
\sphinxstyleemphasis{extra} option of the {\hyperref[\detokenize{compute_modify::doc}]{\sphinxcrossref{\DUrole{doc}{compute\_modify}}}} command.

\sphinxAtStartPar
See the {\hyperref[\detokenize{Howto_thermostat::doc}]{\sphinxcrossref{\DUrole{doc}{Howto thermostat}}}} page for a
discussion of different ways to compute temperature and perform
thermostatting.


\subsection{Output info}
\label{\detokenize{compute_temp_deform:output-info}}
\sphinxAtStartPar
This compute calculates a global scalar (the temperature) and a global
vector of length 6 (symmetric tensor), which can be accessed by
indices 1\textendash{}6.  These values can be used by any command that uses
global scalar or vector values from a compute as input.  See the
{\hyperref[\detokenize{Howto_output::doc}]{\sphinxcrossref{\DUrole{doc}{Howto output}}}} page for an overview of LAMMPS
output options.

\sphinxAtStartPar
The scalar value calculated by this compute is “intensive”.  The
vector values are “extensive”.

\sphinxAtStartPar
The scalar value is in temperature {\hyperref[\detokenize{units::doc}]{\sphinxcrossref{\DUrole{doc}{units}}}}.  The vector
values are in energy {\hyperref[\detokenize{units::doc}]{\sphinxcrossref{\DUrole{doc}{units}}}}.


\subsection{Restrictions}
\label{\detokenize{compute_temp_deform:restrictions}}\begin{quote}

\sphinxAtStartPar
none
\end{quote}


\subsection{Related commands}
\label{\detokenize{compute_temp_deform:related-commands}}
\sphinxAtStartPar
{\hyperref[\detokenize{compute_temp_ramp::doc}]{\sphinxcrossref{\DUrole{doc}{compute temp/ramp}}}},
{\hyperref[\detokenize{compute_temp_profile::doc}]{\sphinxcrossref{\DUrole{doc}{compute temp/profile}}}},
{\hyperref[\detokenize{fix_deform::doc}]{\sphinxcrossref{\DUrole{doc}{fix deform}}}},
{\hyperref[\detokenize{fix_nvt_sllod::doc}]{\sphinxcrossref{\DUrole{doc}{fix nvt/sllod}}}}


\subsection{Default}
\label{\detokenize{compute_temp_deform:default}}
\sphinxAtStartPar
none

\sphinxstepscope

\index{compute temp/deform/eff@\spxentry{compute temp/deform/eff}}\ignorespaces 

\section{compute temp/deform/eff command}
\label{\detokenize{compute_temp_deform_eff:compute-temp-deform-eff-command}}\label{\detokenize{compute_temp_deform_eff:index-0}}\label{\detokenize{compute_temp_deform_eff::doc}}

\subsection{Syntax}
\label{\detokenize{compute_temp_deform_eff:syntax}}
\begin{sphinxVerbatim}[commandchars=\\\{\}]
\PYG{k}{compute }\PYG{n+nv+nvIdentifier}{ID}\PYG{+w}{ }\PYG{n+nv+nvIdentifier}{group\PYGZhy{}ID}\PYG{+w}{ }\PYG{n}{temp}\PYG{o}{/}\PYG{n}{deform}\PYG{o}{/}\PYG{n}{eff}
\end{sphinxVerbatim}
\begin{itemize}
\item {} 
\sphinxAtStartPar
ID, group\sphinxhyphen{}ID are documented in {\hyperref[\detokenize{compute::doc}]{\sphinxcrossref{\DUrole{doc}{compute}}}} command

\item {} 
\sphinxAtStartPar
temp/deform/eff = style name of this compute command

\end{itemize}


\subsection{Examples}
\label{\detokenize{compute_temp_deform_eff:examples}}
\begin{sphinxVerbatim}[commandchars=\\\{\}]
\PYG{k}{compute }\PYG{n+nv+nvIdentifier}{myTemp}\PYG{+w}{ }\PYG{n+nv+nvIdentifier}{all}\PYG{+w}{ }\PYG{n}{temp}\PYG{o}{/}\PYG{n}{deform}\PYG{o}{/}\PYG{n}{eff}
\end{sphinxVerbatim}


\subsection{Description}
\label{\detokenize{compute_temp_deform_eff:description}}
\sphinxAtStartPar
Define a computation that calculates the temperature of a group of
nuclei and electrons in the {\hyperref[\detokenize{pair_eff::doc}]{\sphinxcrossref{\DUrole{doc}{electron force field}}}}
model, after subtracting out a streaming velocity induced by the
simulation box changing size and/or shape, for example in a
non\sphinxhyphen{}equilibrium MD (NEMD) simulation.  The size/shape change is
induced by use of the {\hyperref[\detokenize{fix_deform::doc}]{\sphinxcrossref{\DUrole{doc}{fix deform}}}} command.  A
compute of this style is created by the {\hyperref[\detokenize{fix_nvt_sllod_eff::doc}]{\sphinxcrossref{\DUrole{doc}{fix nvt/sllod/eff}}}} command to compute the thermal temperature of
atoms for thermostatting purposes.  A compute of this style can also
be used by any command that computes a temperature (e.g.,
{\hyperref[\detokenize{thermo_modify::doc}]{\sphinxcrossref{\DUrole{doc}{thermo\_modify}}}}, {\hyperref[\detokenize{fix_nh_eff::doc}]{\sphinxcrossref{\DUrole{doc}{fix npt/eff}}}}).

\sphinxAtStartPar
The calculation performed by this compute is exactly like that
described by the {\hyperref[\detokenize{compute_temp_deform::doc}]{\sphinxcrossref{\DUrole{doc}{compute temp/deform}}}}
command, except that the formulas for the temperature (scalar) and
diagonal components of the symmetric tensor (vector) include the
radial electron velocity contributions, as discussed by the
{\hyperref[\detokenize{compute_temp_eff::doc}]{\sphinxcrossref{\DUrole{doc}{compute temp/eff}}}} command.  Note that only
the translational degrees of freedom for each nuclei or electron are
affected by the streaming velocity adjustment.  The radial velocity
component of the electrons is not affected.


\subsection{Output info}
\label{\detokenize{compute_temp_deform_eff:output-info}}
\sphinxAtStartPar
This compute calculates a global scalar (the temperature) and a global
vector of length 6 (symmetric tensor), which can be accessed by
indices 1\textendash{}6.  These values can be used by any command that uses
global scalar or vector values from a compute as input.  See the
{\hyperref[\detokenize{Howto_output::doc}]{\sphinxcrossref{\DUrole{doc}{Howto output}}}} page for an overview of LAMMPS
output options.

\sphinxAtStartPar
The scalar value calculated by this compute is “intensive”.  The
vector values are “extensive”.

\sphinxAtStartPar
The scalar value is in temperature {\hyperref[\detokenize{units::doc}]{\sphinxcrossref{\DUrole{doc}{units}}}}.  The vector
values are in energy {\hyperref[\detokenize{units::doc}]{\sphinxcrossref{\DUrole{doc}{units}}}}.


\subsection{Restrictions}
\label{\detokenize{compute_temp_deform_eff:restrictions}}
\sphinxAtStartPar
This compute is part of the EFF package.  It is only enabled if
LAMMPS was built with that package.  See the
{\hyperref[\detokenize{Build_package::doc}]{\sphinxcrossref{\DUrole{doc}{Build package}}}} page for more info.


\subsection{Related commands}
\label{\detokenize{compute_temp_deform_eff:related-commands}}
\sphinxAtStartPar
{\hyperref[\detokenize{compute_temp_ramp::doc}]{\sphinxcrossref{\DUrole{doc}{compute temp/ramp}}}}, {\hyperref[\detokenize{fix_deform::doc}]{\sphinxcrossref{\DUrole{doc}{fix deform}}}},
{\hyperref[\detokenize{fix_nvt_sllod_eff::doc}]{\sphinxcrossref{\DUrole{doc}{fix nvt/sllod/eff}}}}


\subsection{Default}
\label{\detokenize{compute_temp_deform_eff:default}}
\sphinxAtStartPar
none

\sphinxstepscope

\index{compute temp/drude@\spxentry{compute temp/drude}}\ignorespaces 

\section{compute temp/drude command}
\label{\detokenize{compute_temp_drude:compute-temp-drude-command}}\label{\detokenize{compute_temp_drude:index-0}}\label{\detokenize{compute_temp_drude::doc}}

\subsection{Syntax}
\label{\detokenize{compute_temp_drude:syntax}}
\begin{sphinxVerbatim}[commandchars=\\\{\}]
\PYG{k}{compute }\PYG{n+nv+nvIdentifier}{ID}\PYG{+w}{ }\PYG{n+nv+nvIdentifier}{group\PYGZhy{}ID}\PYG{+w}{ }\PYG{n}{temp}\PYG{o}{/}\PYG{n}{drude}
\end{sphinxVerbatim}
\begin{itemize}
\item {} 
\sphinxAtStartPar
ID, group\sphinxhyphen{}ID are documented in {\hyperref[\detokenize{compute::doc}]{\sphinxcrossref{\DUrole{doc}{compute}}}} command

\item {} 
\sphinxAtStartPar
temp/drude = style name of this compute command

\end{itemize}


\subsection{Examples}
\label{\detokenize{compute_temp_drude:examples}}
\begin{sphinxVerbatim}[commandchars=\\\{\}]
\PYG{k}{compute }\PYG{n+nv+nvIdentifier}{TDRUDE}\PYG{+w}{ }\PYG{n+nv+nvIdentifier}{all}\PYG{+w}{ }\PYG{n}{temp}\PYG{o}{/}\PYG{n}{drude}
\end{sphinxVerbatim}

\sphinxAtStartPar
Example input scripts available: \sphinxcode{\sphinxupquote{examples/PACKAGES/drude}}.


\subsection{Description}
\label{\detokenize{compute_temp_drude:description}}
\sphinxAtStartPar
Define a computation that calculates the temperatures of core\textendash{}Drude
pairs. This compute is designed to be used with the
{\hyperref[\detokenize{Howto_drude::doc}]{\sphinxcrossref{\DUrole{doc}{thermalized Drude oscillator model}}}}.
Polarizable models in LAMMPS
are described on the {\hyperref[\detokenize{Howto_polarizable::doc}]{\sphinxcrossref{\DUrole{doc}{Howto polarizable}}}} doc page.

\sphinxAtStartPar
Drude oscillators consist of a core particle and a Drude particle
connected by a harmonic bond, and the relative motion of these Drude
oscillators is usually maintained cold by a specific thermostat that
acts on the relative motion of the core\textendash{}Drude particle
pairs. Therefore, because LAMMPS considers Drude particles as normal
atoms in its default temperature compute ({\hyperref[\detokenize{compute_temp::doc}]{\sphinxcrossref{\DUrole{doc}{compute temp}}}}
command), the reduced temperature of the core\textendash{}Drude particle pairs is not
calculated correctly.

\sphinxAtStartPar
By contrast, this compute calculates the temperature of the cores
using center\sphinxhyphen{}of\sphinxhyphen{}mass velocities of the core\textendash{}Drude pairs, and the
reduced temperature of the Drude particles using the relative
velocities of the Drude particles with respect to their cores.
Non\sphinxhyphen{}polarizable atoms are considered as cores.  Their velocities
contribute to the temperature of the cores.


\subsection{Output info}
\label{\detokenize{compute_temp_drude:output-info}}
\sphinxAtStartPar
This compute calculates a global scalar (the temperature) and a global
vector of length 6, which can be accessed by indices 1\textendash{}6, whose components
are
\begin{enumerate}
\sphinxsetlistlabels{\arabic}{enumi}{enumii}{}{.}%
\item {} 
\sphinxAtStartPar
temperature of the centers of mass (temperature units)

\item {} 
\sphinxAtStartPar
temperature of the dipoles (temperature units)

\item {} 
\sphinxAtStartPar
number of degrees of freedom of the centers of mass

\item {} 
\sphinxAtStartPar
number of degrees of freedom of the dipoles

\item {} 
\sphinxAtStartPar
kinetic energy of the centers of mass (energy units)

\item {} 
\sphinxAtStartPar
kinetic energy of the dipoles (energy units)

\end{enumerate}

\sphinxAtStartPar
These values can be used by any command that uses global scalar or
vector values from a compute as input.  See the
{\hyperref[\detokenize{Howto_output::doc}]{\sphinxcrossref{\DUrole{doc}{Howto output}}}} page for an overview of LAMMPS output
options.

\sphinxAtStartPar
Both the scalar value and the first two values of the vector
calculated by this compute are “intensive”.  The other four vector values
are “extensive”.


\subsection{Restrictions}
\label{\detokenize{compute_temp_drude:restrictions}}
\sphinxAtStartPar
The number of degrees of freedom contributing to the temperature is
assumed to be constant for the duration of the run unless the
{\hyperref[\detokenize{fix_modify::doc}]{\sphinxcrossref{\DUrole{doc}{fix\_modify command}}}} sets the option \sphinxstyleemphasis{dynamic/dof
yes}.


\subsection{Related commands}
\label{\detokenize{compute_temp_drude:related-commands}}
\sphinxAtStartPar
{\hyperref[\detokenize{fix_drude::doc}]{\sphinxcrossref{\DUrole{doc}{fix drude}}}}, {\hyperref[\detokenize{fix_langevin_drude::doc}]{\sphinxcrossref{\DUrole{doc}{fix langevin/drude}}}},
{\hyperref[\detokenize{fix_drude_transform::doc}]{\sphinxcrossref{\DUrole{doc}{fix drude/transform}}}},
{\hyperref[\detokenize{pair_thole::doc}]{\sphinxcrossref{\DUrole{doc}{pair\_style thole}}}}, {\hyperref[\detokenize{compute_temp::doc}]{\sphinxcrossref{\DUrole{doc}{compute temp}}}}


\subsection{Default}
\label{\detokenize{compute_temp_drude:default}}
\sphinxAtStartPar
none

\sphinxstepscope

\index{compute temp/eff@\spxentry{compute temp/eff}}\ignorespaces 

\section{compute temp/eff command}
\label{\detokenize{compute_temp_eff:compute-temp-eff-command}}\label{\detokenize{compute_temp_eff:index-0}}\label{\detokenize{compute_temp_eff::doc}}

\subsection{Syntax}
\label{\detokenize{compute_temp_eff:syntax}}
\begin{sphinxVerbatim}[commandchars=\\\{\}]
\PYG{k}{compute }\PYG{n+nv+nvIdentifier}{ID}\PYG{+w}{ }\PYG{n+nv+nvIdentifier}{group\PYGZhy{}ID}\PYG{+w}{ }\PYG{n}{temp}\PYG{o}{/}\PYG{n}{eff}
\end{sphinxVerbatim}
\begin{itemize}
\item {} 
\sphinxAtStartPar
ID, group\sphinxhyphen{}ID are documented in {\hyperref[\detokenize{compute::doc}]{\sphinxcrossref{\DUrole{doc}{compute}}}} command

\item {} 
\sphinxAtStartPar
temp/eff = style name of this compute command

\end{itemize}


\subsection{Examples}
\label{\detokenize{compute_temp_eff:examples}}
\begin{sphinxVerbatim}[commandchars=\\\{\}]
\PYG{k}{compute }\PYG{n+nv+nvIdentifier}{1}\PYG{+w}{ }\PYG{n+nv+nvIdentifier}{all}\PYG{+w}{ }\PYG{n}{temp}\PYG{o}{/}\PYG{n}{eff}
\PYG{k}{compute }\PYG{n+nv+nvIdentifier}{myTemp}\PYG{+w}{ }\PYG{n+nv+nvIdentifier}{mobile}\PYG{+w}{ }\PYG{n}{temp}\PYG{o}{/}\PYG{n}{eff}
\end{sphinxVerbatim}


\subsection{Description}
\label{\detokenize{compute_temp_eff:description}}
\sphinxAtStartPar
Define a computation that calculates the temperature of a group of
nuclei and electrons in the {\hyperref[\detokenize{pair_eff::doc}]{\sphinxcrossref{\DUrole{doc}{electron force field}}}}
model.  A compute of this style can be used by commands that compute a
temperature (e.g., {\hyperref[\detokenize{thermo_modify::doc}]{\sphinxcrossref{\DUrole{doc}{thermo\_modify}}}},
{\hyperref[\detokenize{fix_nh_eff::doc}]{\sphinxcrossref{\DUrole{doc}{fix npt/eff}}}}).

\sphinxAtStartPar
The temperature is calculated by the formula
\begin{equation*}
\begin{split}\text{KE} = \frac{\text{dim}}{2} N k_B T,\end{split}
\end{equation*}
\sphinxAtStartPar
where KE is the total kinetic energy of the group of atoms (sum of
\(\frac12 m v^2\) for nuclei and sum of
\(\frac12 (m v^2 + \frac34 m s^2\)) for electrons, where \(s\)
includes the radial electron velocity contributions), dim = 2 or 3 is the
dimensionality of the simulation, \(N\) is the number of atoms (only total
number of nuclei in the eFF (see the {\hyperref[\detokenize{pair_style::doc}]{\sphinxcrossref{\DUrole{doc}{pair\_eff}}}}
command) in the group, \(k_B\) is the Boltzmann constant, and \(T\) is
the absolute temperature.  This expression is summed over all nuclear and
electronic degrees of freedom, essentially by setting the kinetic contribution
to the heat capacity to \(\frac32 k\) (where only nuclei contribute). This
subtlety is valid for temperatures well below the Fermi temperature, which for
densities two to five times the density of liquid hydrogen ranges from
86,000 to 170,000 K.

\begin{sphinxadmonition}{note}{Note:}
\sphinxAtStartPar
For eFF models, in order to override the default temperature
reported by LAMMPS in the thermodynamic quantities reported via the
{\hyperref[\detokenize{thermo::doc}]{\sphinxcrossref{\DUrole{doc}{thermo}}}} command, the user should apply a
{\hyperref[\detokenize{thermo_modify::doc}]{\sphinxcrossref{\DUrole{doc}{thermo\_modify}}}} command, as shown in the following
example:
\end{sphinxadmonition}

\begin{sphinxVerbatim}[commandchars=\\\{\}]
\PYG{k}{compute         }\PYG{n+nv+nvIdentifier}{effTemp}\PYG{+w}{ }\PYG{n+nv+nvIdentifier}{all}\PYG{+w}{ }\PYG{n}{temp}\PYG{o}{/}\PYG{n}{eff}
\PYG{k}{thermo\PYGZus{}style}\PYG{+w}{    }\PYG{n}{custom}\PYG{+w}{ }\PYG{n}{step}\PYG{+w}{ }\PYG{n}{etotal}\PYG{+w}{ }\PYG{n}{pe}\PYG{+w}{ }\PYG{n}{ke}\PYG{+w}{ }\PYG{n}{temp}\PYG{+w}{ }\PYG{n}{press}
\PYG{k}{thermo\PYGZus{}modify}\PYG{+w}{   }\PYG{n}{temp}\PYG{+w}{ }\PYG{n}{effTemp}
\end{sphinxVerbatim}

\sphinxAtStartPar
A six\sphinxhyphen{}component kinetic energy tensor is also calculated by this compute
for use in the computation of a pressure tensor.  The formula for the
components of the tensor is the same as the above formula, except that
\(v^2\) is replaced by \(v_x v_y\) for the \(xy\) component, etc.
For the eFF, again, the radial electronic velocities are also considered.

\sphinxAtStartPar
The number of atoms contributing to the temperature is assumed to be
constant for the duration of the run; use the \sphinxstyleemphasis{dynamic} option of the
{\hyperref[\detokenize{compute_modify::doc}]{\sphinxcrossref{\DUrole{doc}{compute\_modify}}}} command if this is not the case.

\sphinxAtStartPar
This compute subtracts out degrees\sphinxhyphen{}of\sphinxhyphen{}freedom due to fixes that
constrain molecular motion, such as {\hyperref[\detokenize{fix_shake::doc}]{\sphinxcrossref{\DUrole{doc}{fix shake}}}} and
{\hyperref[\detokenize{fix_rigid::doc}]{\sphinxcrossref{\DUrole{doc}{fix rigid}}}}.  This means the temperature of groups of
atoms that include these constraints will be computed correctly.  If
needed, the subtracted degrees\sphinxhyphen{}of\sphinxhyphen{}freedom can be altered using the
\sphinxstyleemphasis{extra} option of the {\hyperref[\detokenize{compute_modify::doc}]{\sphinxcrossref{\DUrole{doc}{compute\_modify}}}} command.

\sphinxAtStartPar
See the {\hyperref[\detokenize{Howto_thermostat::doc}]{\sphinxcrossref{\DUrole{doc}{Howto thermostat}}}} page for a
discussion of different ways to compute temperature and perform
thermostatting.


\subsection{Output info}
\label{\detokenize{compute_temp_eff:output-info}}
\sphinxAtStartPar
The scalar value calculated by this compute is “intensive”, meaning it
is independent of the number of atoms in the simulation.  The vector
values are “extensive”, meaning they scale with the number of atoms in
the simulation.


\subsection{Restrictions}
\label{\detokenize{compute_temp_eff:restrictions}}
\sphinxAtStartPar
This compute is part of the EFF package.  It is only enabled if
LAMMPS was built with that package.  See the {\hyperref[\detokenize{Build_package::doc}]{\sphinxcrossref{\DUrole{doc}{Build package}}}} page for more info.


\subsection{Related commands}
\label{\detokenize{compute_temp_eff:related-commands}}
\sphinxAtStartPar
{\hyperref[\detokenize{compute_temp_partial::doc}]{\sphinxcrossref{\DUrole{doc}{compute temp/partial}}}},
{\hyperref[\detokenize{compute_temp_region::doc}]{\sphinxcrossref{\DUrole{doc}{compute temp/region}}}},
{\hyperref[\detokenize{compute_pressure::doc}]{\sphinxcrossref{\DUrole{doc}{compute pressure}}}}


\subsection{Default}
\label{\detokenize{compute_temp_eff:default}}
\sphinxAtStartPar
none

\sphinxstepscope

\index{compute temp/partial@\spxentry{compute temp/partial}}\ignorespaces 

\section{compute temp/partial command}
\label{\detokenize{compute_temp_partial:compute-temp-partial-command}}\label{\detokenize{compute_temp_partial:index-0}}\label{\detokenize{compute_temp_partial::doc}}

\subsection{Syntax}
\label{\detokenize{compute_temp_partial:syntax}}
\begin{sphinxVerbatim}[commandchars=\\\{\}]
\PYG{k}{compute }\PYG{n+nv+nvIdentifier}{ID}\PYG{+w}{ }\PYG{n+nv+nvIdentifier}{group\PYGZhy{}ID}\PYG{+w}{ }\PYG{n}{temp}\PYG{o}{/}\PYG{n}{partial}\PYG{+w}{ }\PYG{n}{xflag}\PYG{+w}{ }\PYG{n}{yflag}\PYG{+w}{ }\PYG{n}{zflag}
\end{sphinxVerbatim}
\begin{itemize}
\item {} 
\sphinxAtStartPar
ID, group\sphinxhyphen{}ID are documented in {\hyperref[\detokenize{compute::doc}]{\sphinxcrossref{\DUrole{doc}{compute}}}} command

\item {} 
\sphinxAtStartPar
temp/partial = style name of this compute command

\item {} 
\sphinxAtStartPar
xflag,yflag,zflag = 0/1 for whether to exclude/include this dimension

\end{itemize}


\subsection{Examples}
\label{\detokenize{compute_temp_partial:examples}}
\begin{sphinxVerbatim}[commandchars=\\\{\}]
\PYG{k}{compute }\PYG{n+nv+nvIdentifier}{newT}\PYG{+w}{ }\PYG{n+nv+nvIdentifier}{flow}\PYG{+w}{ }\PYG{n}{temp}\PYG{o}{/}\PYG{n}{partial}\PYG{+w}{ }\PYG{l+m}{1}\PYG{+w}{ }\PYG{l+m}{1}\PYG{+w}{ }\PYG{l+m}{0}
\end{sphinxVerbatim}


\subsection{Description}
\label{\detokenize{compute_temp_partial:description}}
\sphinxAtStartPar
Define a computation that calculates the temperature of a group of
atoms, after excluding one or more velocity components.  A compute of
this style can be used by any command that computes a temperature
(e.g. {\hyperref[\detokenize{thermo_modify::doc}]{\sphinxcrossref{\DUrole{doc}{thermo\_modify}}}},
{\hyperref[\detokenize{fix_temp_rescale::doc}]{\sphinxcrossref{\DUrole{doc}{fix temp/rescale}}}}, {\hyperref[\detokenize{fix_nh::doc}]{\sphinxcrossref{\DUrole{doc}{fix npt}}}}).

\sphinxAtStartPar
The temperature is calculated by the formula
\begin{equation*}
\begin{split}\text{KE} = \frac{\text{dim}}{2} N k_B T,\end{split}
\end{equation*}
\sphinxAtStartPar
where KE is the total kinetic energy of the group of atoms (sum of
\(\frac12 m v^2\)), dim = 2 or 3 is the dimensionality of the simulation,
\(N\) is the number of atoms in the group, \(k_B\) is the Boltzmann
constant, and \(T\) = temperature.  The calculation of KE excludes the
\(x\), \(y\), or \(z\) dimensions if \sphinxstyleemphasis{xflag}, \sphinxstyleemphasis{yflag}, or \sphinxstyleemphasis{zflag}
is 0.  The dim parameter is adjusted to give the correct number of
degrees of freedom.

\sphinxAtStartPar
A symmetric tensor, stored as a six\sphinxhyphen{}element vector, is also calculated
by this compute for use in the computation of a pressure tensor by the
{\hyperref[\detokenize{compute_pressure::doc}]{\sphinxcrossref{\DUrole{doc}{compute pressue}}}} command.  The formula for
the components of the tensor is the same as the above expression for
\(E_\mathrm{kin}\), except that the 1/2 factor is NOT included and
the \(v_i^2\) is replaced by \(v_{i,x} v_{i,y}\) for the
\(xy\) component, and so on.  Note that because it lacks the 1/2
factor, these tensor components are twice those of the traditional
kinetic energy tensor.  The six components of the vector are ordered
\(xx\), \(yy\), \(zz\), \(xy\), \(xz\),
\(yz\).

\sphinxAtStartPar
The number of atoms contributing to the temperature is assumed to be
constant for the duration of the run; use the \sphinxstyleemphasis{dynamic} option of the
{\hyperref[\detokenize{compute_modify::doc}]{\sphinxcrossref{\DUrole{doc}{compute\_modify}}}} command if this is not the case.

\sphinxAtStartPar
The removal of velocity components by this fix is essentially
computing the temperature after a “bias” has been removed from the
velocity of the atoms.  If this compute is used with a fix command
that performs thermostatting then this bias will be subtracted from
each atom, thermostatting of the remaining thermal velocity will be
performed, and the bias will be added back in.  Thermostatting fixes
that work in this way include {\hyperref[\detokenize{fix_nh::doc}]{\sphinxcrossref{\DUrole{doc}{fix nvt}}}},
{\hyperref[\detokenize{fix_temp_rescale::doc}]{\sphinxcrossref{\DUrole{doc}{fix temp/rescale}}}},
{\hyperref[\detokenize{fix_temp_berendsen::doc}]{\sphinxcrossref{\DUrole{doc}{fix temp/berendsen}}}}, and
{\hyperref[\detokenize{fix_langevin::doc}]{\sphinxcrossref{\DUrole{doc}{fix langevin}}}}.

\sphinxAtStartPar
This compute subtracts out degrees\sphinxhyphen{}of\sphinxhyphen{}freedom due to fixes that
constrain molecular motion, such as {\hyperref[\detokenize{fix_shake::doc}]{\sphinxcrossref{\DUrole{doc}{fix shake}}}} and
{\hyperref[\detokenize{fix_rigid::doc}]{\sphinxcrossref{\DUrole{doc}{fix rigid}}}}.  This means the temperature of groups of
atoms that include these constraints will be computed correctly.  If
needed, the subtracted degrees\sphinxhyphen{}of\sphinxhyphen{}freedom can be altered using the
\sphinxstyleemphasis{extra} option of the {\hyperref[\detokenize{compute_modify::doc}]{\sphinxcrossref{\DUrole{doc}{compute\_modify}}}} command.

\sphinxAtStartPar
See the {\hyperref[\detokenize{Howto_thermostat::doc}]{\sphinxcrossref{\DUrole{doc}{Howto thermostat}}}} page for a
discussion of different ways to compute temperature and perform
thermostatting.


\bigskip\hrule\bigskip


\sphinxAtStartPar
Styles with a \sphinxstyleemphasis{gpu}, \sphinxstyleemphasis{intel}, \sphinxstyleemphasis{kk}, \sphinxstyleemphasis{omp}, or \sphinxstyleemphasis{opt} suffix are
functionally the same as the corresponding style without the suffix.
They have been optimized to run faster, depending on your available
hardware, as discussed on the {\hyperref[\detokenize{Speed_packages::doc}]{\sphinxcrossref{\DUrole{doc}{Accelerator packages}}}}
page.  The accelerated styles take the same arguments and should
produce the same results, except for round\sphinxhyphen{}off and precision issues.

\sphinxAtStartPar
These accelerated styles are part of the GPU, INTEL, KOKKOS,
OPENMP, and OPT packages, respectively.  They are only enabled if
LAMMPS was built with those packages.  See the {\hyperref[\detokenize{Build_package::doc}]{\sphinxcrossref{\DUrole{doc}{Build package}}}} page for more info.

\sphinxAtStartPar
You can specify the accelerated styles explicitly in your input script
by including their suffix, or you can use the {\hyperref[\detokenize{Run_options::doc}]{\sphinxcrossref{\DUrole{doc}{\sphinxhyphen{}suffix command\sphinxhyphen{}line switch}}}} when you invoke LAMMPS, or you can use the
{\hyperref[\detokenize{suffix::doc}]{\sphinxcrossref{\DUrole{doc}{suffix}}}} command in your input script.

\sphinxAtStartPar
See the {\hyperref[\detokenize{Speed_packages::doc}]{\sphinxcrossref{\DUrole{doc}{Accelerator packages}}}} page for more
instructions on how to use the accelerated styles effectively.


\bigskip\hrule\bigskip



\subsection{Output info}
\label{\detokenize{compute_temp_partial:output-info}}
\sphinxAtStartPar
This compute calculates a global scalar (the temperature) and a global
vector of length 6 (symmetric tensor), which can be accessed by
indices 1\textendash{}6.  These values can be used by any command that uses
global scalar or vector values from a compute as input.  See the
{\hyperref[\detokenize{Howto_output::doc}]{\sphinxcrossref{\DUrole{doc}{Howto output}}}} page for an overview of LAMMPS
output options.

\sphinxAtStartPar
The scalar value calculated by this compute is “intensive”.  The
vector values are “extensive”.

\sphinxAtStartPar
The scalar value is in temperature {\hyperref[\detokenize{units::doc}]{\sphinxcrossref{\DUrole{doc}{units}}}}.  The vector
values are in energy {\hyperref[\detokenize{units::doc}]{\sphinxcrossref{\DUrole{doc}{units}}}}.


\subsection{Restrictions}
\label{\detokenize{compute_temp_partial:restrictions}}\begin{quote}

\sphinxAtStartPar
none
\end{quote}


\subsection{Related commands}
\label{\detokenize{compute_temp_partial:related-commands}}
\sphinxAtStartPar
{\hyperref[\detokenize{compute_temp::doc}]{\sphinxcrossref{\DUrole{doc}{compute temp}}}}, {\hyperref[\detokenize{compute_temp_region::doc}]{\sphinxcrossref{\DUrole{doc}{compute temp/region}}}}, {\hyperref[\detokenize{compute_pressure::doc}]{\sphinxcrossref{\DUrole{doc}{compute pressure}}}}


\subsection{Default}
\label{\detokenize{compute_temp_partial:default}}
\sphinxAtStartPar
none

\sphinxstepscope

\index{compute temp/profile@\spxentry{compute temp/profile}}\ignorespaces 

\section{compute temp/profile command}
\label{\detokenize{compute_temp_profile:compute-temp-profile-command}}\label{\detokenize{compute_temp_profile:index-0}}\label{\detokenize{compute_temp_profile::doc}}

\subsection{Syntax}
\label{\detokenize{compute_temp_profile:syntax}}
\begin{sphinxVerbatim}[commandchars=\\\{\}]
\PYG{k}{compute }\PYG{n+nv+nvIdentifier}{ID}\PYG{+w}{ }\PYG{n+nv+nvIdentifier}{group\PYGZhy{}ID}\PYG{+w}{ }\PYG{n}{temp}\PYG{o}{/}\PYG{n}{profile}\PYG{+w}{ }\PYG{n}{xflag}\PYG{+w}{ }\PYG{n}{yflag}\PYG{+w}{ }\PYG{n}{zflag}\PYG{+w}{ }\PYG{n}{binstyle}\PYG{+w}{ }\PYG{n}{args}
\end{sphinxVerbatim}
\begin{itemize}
\item {} 
\sphinxAtStartPar
ID, group\sphinxhyphen{}ID are documented in {\hyperref[\detokenize{compute::doc}]{\sphinxcrossref{\DUrole{doc}{compute}}}} command

\item {} 
\sphinxAtStartPar
temp/profile = style name of this compute command

\item {} 
\sphinxAtStartPar
xflag,yflag,zflag = 0/1 for whether to exclude/include this dimension

\item {} 
\sphinxAtStartPar
binstyle = \sphinxstyleemphasis{x} or \sphinxstyleemphasis{y} or \sphinxstyleemphasis{z} or \sphinxstyleemphasis{xy} or \sphinxstyleemphasis{yz} or \sphinxstyleemphasis{xz} or \sphinxstyleemphasis{xyz}
\begin{sphinxalltt}
\sphinxstyleemphasis{x} arg = Nx
\sphinxstyleemphasis{y} arg = Ny
\sphinxstyleemphasis{z} arg = Nz
\sphinxstyleemphasis{xy} args = Nx Ny
\sphinxstyleemphasis{yz} args = Ny Nz
\sphinxstyleemphasis{xz} args = Nx Nz
\sphinxstyleemphasis{xyz} args = Nx Ny Nz
  Nx, Ny, Nz = number of velocity bins in \sphinxstyleemphasis{x}, \sphinxstyleemphasis{y}, \sphinxstyleemphasis{z} dimensions
\end{sphinxalltt}

\item {} 
\sphinxAtStartPar
zero or more keyword/value pairs may be appended

\item {} 
\sphinxAtStartPar
keyword = \sphinxstyleemphasis{out}
\begin{sphinxalltt}
\sphinxstyleemphasis{out} value = \sphinxstyleemphasis{tensor} or \sphinxstyleemphasis{bin}
\end{sphinxalltt}

\end{itemize}


\subsection{Examples}
\label{\detokenize{compute_temp_profile:examples}}
\begin{sphinxVerbatim}[commandchars=\\\{\}]
\PYG{k}{compute }\PYG{n+nv+nvIdentifier}{myTemp}\PYG{+w}{ }\PYG{n+nv+nvIdentifier}{flow}\PYG{+w}{ }\PYG{n}{temp}\PYG{o}{/}\PYG{n}{profile}\PYG{+w}{ }\PYG{l+m}{1}\PYG{+w}{ }\PYG{l+m}{1}\PYG{+w}{ }\PYG{l+m}{1}\PYG{+w}{ }\PYG{n}{x}\PYG{+w}{ }\PYG{l+m}{10}
\PYG{k}{compute }\PYG{n+nv+nvIdentifier}{myTemp}\PYG{+w}{ }\PYG{n+nv+nvIdentifier}{flow}\PYG{+w}{ }\PYG{n}{temp}\PYG{o}{/}\PYG{n}{profile}\PYG{+w}{ }\PYG{l+m}{1}\PYG{+w}{ }\PYG{l+m}{1}\PYG{+w}{ }\PYG{l+m}{1}\PYG{+w}{ }\PYG{n}{x}\PYG{+w}{ }\PYG{l+m}{10}\PYG{+w}{ }\PYG{n}{out}\PYG{+w}{ }\PYG{n}{bin}
\PYG{k}{compute }\PYG{n+nv+nvIdentifier}{myTemp}\PYG{+w}{ }\PYG{n+nv+nvIdentifier}{flow}\PYG{+w}{ }\PYG{n}{temp}\PYG{o}{/}\PYG{n}{profile}\PYG{+w}{ }\PYG{l+m}{0}\PYG{+w}{ }\PYG{l+m}{1}\PYG{+w}{ }\PYG{l+m}{1}\PYG{+w}{ }\PYG{n}{xyz}\PYG{+w}{ }\PYG{l+m}{20}\PYG{+w}{ }\PYG{l+m}{20}\PYG{+w}{ }\PYG{l+m}{20}
\end{sphinxVerbatim}


\subsection{Description}
\label{\detokenize{compute_temp_profile:description}}
\sphinxAtStartPar
Define a computation that calculates the temperature of a group of
atoms, after subtracting out a spatially\sphinxhyphen{}averaged center\sphinxhyphen{}of\sphinxhyphen{}mass
velocity field, before computing the kinetic energy.  This can be
useful for thermostatting a collection of atoms undergoing a complex
flow (e.g. via a profile\sphinxhyphen{}unbiased thermostat (PUT) as described in
{\hyperref[\detokenize{compute_temp_profile:evans1}]{\sphinxcrossref{\DUrole{std,std-ref}{(Evans)}}}}).  A compute of this style can be used by any command
that computes a temperature (e.g. {\hyperref[\detokenize{thermo_modify::doc}]{\sphinxcrossref{\DUrole{doc}{thermo\_modify}}}},
{\hyperref[\detokenize{fix_temp_rescale::doc}]{\sphinxcrossref{\DUrole{doc}{fix temp/rescale}}}}, {\hyperref[\detokenize{fix_nh::doc}]{\sphinxcrossref{\DUrole{doc}{fix npt}}}}).

\sphinxAtStartPar
The \sphinxstyleemphasis{xflag}, \sphinxstyleemphasis{yflag}, \sphinxstyleemphasis{zflag} settings determine which components of
average velocity are subtracted out.

\sphinxAtStartPar
The \sphinxstyleemphasis{binstyle} setting and its \sphinxstyleemphasis{Nx}, \sphinxstyleemphasis{Ny}, \sphinxstyleemphasis{Nz} arguments determine how bins
are setup to perform spatial averaging.  “Bins” can be 1d slabs, 2d pencils,
or 3d bricks depending on which \sphinxstyleemphasis{binstyle} is used.  The simulation box is
partitioned conceptually into \sphinxstyleemphasis{Nx} \(\times\) \sphinxstyleemphasis{Ny} \(\times\) \sphinxstyleemphasis{Nz}
bins.  Depending on the \sphinxstyleemphasis{binstyle}, you may only specify one or two of these
values; the others are effectively set to 1 (no binning in that dimension).
For non\sphinxhyphen{}orthogonal (triclinic) simulation boxes, the bins are “tilted” slabs or
pencils or bricks that are parallel to the tilted faces of the box.  See the
{\hyperref[\detokenize{region::doc}]{\sphinxcrossref{\DUrole{doc}{region prism}}}} command for a discussion of the geometry of tilted
boxes in LAMMPS.

\sphinxAtStartPar
When a temperature is computed, the center\sphinxhyphen{}of\sphinxhyphen{}mass velocity for the
set of atoms that are both in the compute group and in the same
spatial bin is calculated.  This bias velocity is then subtracted from
the velocities of individual atoms in the bin to yield a thermal
velocity for each atom.  Note that if there is only one atom in the
bin, its thermal velocity will thus be 0.0.

\sphinxAtStartPar
After the spatially\sphinxhyphen{}averaged velocity field has been subtracted from
each atom, the temperature is calculated by the formula
\begin{equation*}
\begin{split}\text{KE} = \left( \frac{\text{dim}}{N} - N_s N_x N_y N_z
                      - \text{extra} \right) \frac{k_B T}{2},\end{split}
\end{equation*}
\sphinxAtStartPar
where KE is the total kinetic energy of the group of atoms (sum of
\(\frac12 m v^2\); dim = 2 or 3 is the dimensionality of the simulation;
\(N_s =\) 0, 1, 2, or 3 for streaming velocity subtracted in 0, 1, 2, or 3
dimensions, respectively; \sphinxstyleemphasis{extra} is the number of  extra degrees of freedom;
\sphinxstyleemphasis{N} is the number of atoms in the group; \(k_B\) is the Boltzmann constant,
and \(T\) is the absolute temperature.  The \(N_s N_x N_y N_z\) term is
the number of degrees of freedom subtracted to adjust for the removal of the
center\sphinxhyphen{}of\sphinxhyphen{}mass velocity in each direction of the \sphinxstyleemphasis{Nx*Ny*Nz} bins, as
discussed in the {\hyperref[\detokenize{compute_temp_profile:evans1}]{\sphinxcrossref{\DUrole{std,std-ref}{(Evans)}}}} paper.  The extra term defaults to
\(\text{dim} - N_s\) and accounts for overall conservation of
center\sphinxhyphen{}of\sphinxhyphen{}mass velocity across the group in directions where streaming velocity
is \sphinxstyleemphasis{not} subtracted. This can be altered using the \sphinxstyleemphasis{extra} option of the
{\hyperref[\detokenize{compute_modify::doc}]{\sphinxcrossref{\DUrole{doc}{compute\_modify}}}} command.

\sphinxAtStartPar
If the \sphinxstyleemphasis{out} keyword is used with a \sphinxstyleemphasis{tensor} value, which is the
default, then a symmetric tensor, stored as a six\sphinxhyphen{}element vector, is
also calculated by this compute for use in the computation of a
pressure tensor by the {\hyperref[\detokenize{compute_pressure::doc}]{\sphinxcrossref{\DUrole{doc}{compute pressue}}}}
command.  The formula for the components of the tensor is the same as
the above expression for \(E_\mathrm{kin}\), except that the 1/2
factor is NOT included and the \(v_i^2\) is replaced by
\(v_{i,x} v_{i,y}\) for the \(xy\) component, and so on.  Note
that because it lacks the 1/2 factor, these tensor components are
twice those of the traditional kinetic energy tensor.  The six
components of the vector are ordered \(xx\), \(yy\),
\(zz\), \(xy\), \(xz\), \(yz\).

\sphinxAtStartPar
If the \sphinxstyleemphasis{out} keyword is used with a \sphinxstyleemphasis{bin} value, the count of atoms
and computed temperature for each bin are stored for output, as an
array of values, as described below.  The temperature of each bin is
calculated as described above, where the bias velocity is subtracted
and only the remaining thermal velocity of atoms in the bin
contributes to the temperature.  See the note below for how the
temperature is normalized by the degrees\sphinxhyphen{}of\sphinxhyphen{}freedom of atoms in the
bin.

\sphinxAtStartPar
The number of atoms contributing to the temperature is assumed to be
constant for the duration of the run; use the \sphinxstyleemphasis{dynamic} option of the
{\hyperref[\detokenize{compute_modify::doc}]{\sphinxcrossref{\DUrole{doc}{compute\_modify}}}} command if this is not the case.

\sphinxAtStartPar
The removal of the spatially\sphinxhyphen{}averaged velocity field by this fix is
essentially computing the temperature after a “bias” has been removed
from the velocity of the atoms.  If this compute is used with a fix
command that performs thermostatting then this bias will be subtracted
from each atom, thermostatting of the remaining thermal velocity will
be performed, and the bias will be added back in.  Thermostatting
fixes that work in this way include {\hyperref[\detokenize{fix_nh::doc}]{\sphinxcrossref{\DUrole{doc}{fix nvt}}}},
{\hyperref[\detokenize{fix_temp_rescale::doc}]{\sphinxcrossref{\DUrole{doc}{fix temp/rescale}}}},
{\hyperref[\detokenize{fix_temp_berendsen::doc}]{\sphinxcrossref{\DUrole{doc}{fix temp/berendsen}}}},
and {\hyperref[\detokenize{fix_langevin::doc}]{\sphinxcrossref{\DUrole{doc}{fix langevin}}}}.

\sphinxAtStartPar
This compute subtracts out degrees\sphinxhyphen{}of\sphinxhyphen{}freedom due to fixes that constrain
molecular motion, such as {\hyperref[\detokenize{fix_shake::doc}]{\sphinxcrossref{\DUrole{doc}{fix shake}}}} and
{\hyperref[\detokenize{fix_rigid::doc}]{\sphinxcrossref{\DUrole{doc}{fix rigid}}}}.  This means the temperature of groups of atoms
that include these constraints will be computed correctly.  If needed, the
subtracted degrees\sphinxhyphen{}of\sphinxhyphen{}freedom can be altered using the \sphinxstyleemphasis{extra} option of the
{\hyperref[\detokenize{compute_modify::doc}]{\sphinxcrossref{\DUrole{doc}{compute\_modify}}}} command.

\begin{sphinxadmonition}{note}{Note:}
\sphinxAtStartPar
When using the \sphinxstyleemphasis{out} keyword with a value of \sphinxstyleemphasis{bin}, the
calculated temperature for each bin includes the degrees\sphinxhyphen{}of\sphinxhyphen{}freedom
adjustment described in the preceding paragraph for fixes that
constrain molecular motion, as well as the adjustment due to
the \sphinxstyleemphasis{extra} option (which defaults to \sphinxstyleemphasis{dim} \sphinxhyphen{} \sphinxstyleemphasis{Ns} as described above),
by fractionally applying them based on the fraction of atoms in each
bin. As a result, the bin degrees\sphinxhyphen{}of\sphinxhyphen{}freedom summed over all bins exactly
equals the degrees\sphinxhyphen{}of\sphinxhyphen{}freedom used in the scalar temperature calculation,
\(\Sigma N_{\text{DOF}_i} = N_\text{DOF}\) and the corresponding
relation for temperature is also satisfied
(\(\Sigma N_{\text{DOF}_i} T_i = N_\text{DOF} T\)).
These relations will break down in cases for which the adjustment
exceeds the actual number of degrees of freedom in a bin. This could happen
if a bin is empty or in situations in which rigid molecules
are non\sphinxhyphen{}uniformly distributed, in which case the reported
temperature within a bin may not be accurate.
\end{sphinxadmonition}

\sphinxAtStartPar
See the {\hyperref[\detokenize{Howto_thermostat::doc}]{\sphinxcrossref{\DUrole{doc}{Howto thermostat}}}} page for a
discussion of different ways to compute temperature and perform
thermostatting.  Using this compute in conjunction with a
thermostatting fix, as explained there, will effectively implement a
profile\sphinxhyphen{}unbiased thermostat (PUT), as described in {\hyperref[\detokenize{compute_temp_profile:evans1}]{\sphinxcrossref{\DUrole{std,std-ref}{(Evans)}}}}.


\subsection{Output info}
\label{\detokenize{compute_temp_profile:output-info}}
\sphinxAtStartPar
This compute calculates a global scalar (the temperature).  Depending
on the setting of the \sphinxstyleemphasis{out} keyword, it also calculates a global
vector or array.  For \sphinxstyleemphasis{out} = \sphinxstyleemphasis{tensor}, it calculates a vector of
length 6 (symmetric tensor), which can be accessed by indices 1\textendash{}6.
For \sphinxstyleemphasis{out} = \sphinxstyleemphasis{bin} it calculates a global array which has 2 columns and
\(N\) rows, where \(N\) is the number of bins.  The first
column contains the number of atoms in that bin.  The second contains
the temperature of that bin, calculated as described above.  The
ordering of rows in the array is as follows.  Bins in \(x\) vary
fastest, then \(y\), then \(z\).  Thus for a \(10\times
10\times 10\) 3d array of bins, there will be 1000 rows.  The bin with
indices \((i_x,i_y,i_z) = (2,3,4)\) would map to row \(M =
10^2(i_z-1) + 10(i_y-1) + i_x = 322\), where the rows are numbered from
1 to 1000 and the bin indices are numbered from 1 to 10 in each
dimension.

\sphinxAtStartPar
These values can be used by any command that uses global scalar or
vector or array values from a compute as input.  See the
{\hyperref[\detokenize{Howto_output::doc}]{\sphinxcrossref{\DUrole{doc}{Howto output}}}} page for an overview of LAMMPS output
options.

\sphinxAtStartPar
The scalar value calculated by this compute is “intensive”.  The
vector values are “extensive”.  The array values are “intensive”.

\sphinxAtStartPar
The scalar value us in temperature {\hyperref[\detokenize{units::doc}]{\sphinxcrossref{\DUrole{doc}{units}}}}.  The vector
values are in energy {\hyperref[\detokenize{units::doc}]{\sphinxcrossref{\DUrole{doc}{units}}}}.  The first column of array
values are counts; the values in the second column will be in
temperature {\hyperref[\detokenize{units::doc}]{\sphinxcrossref{\DUrole{doc}{units}}}}.


\subsection{Restrictions}
\label{\detokenize{compute_temp_profile:restrictions}}
\sphinxAtStartPar
You should not use too large a velocity\sphinxhyphen{}binning grid, especially in
3d.  In the current implementation, the binned velocity averages are
summed across all processors, so this will be inefficient if the grid
is too large, and the operation is performed every timestep, as it
will be for most thermostats.


\subsection{Related commands}
\label{\detokenize{compute_temp_profile:related-commands}}
\sphinxAtStartPar
{\hyperref[\detokenize{compute_temp::doc}]{\sphinxcrossref{\DUrole{doc}{compute temp}}}},
{\hyperref[\detokenize{compute_temp_ramp::doc}]{\sphinxcrossref{\DUrole{doc}{compute temp/ramp}}}},
{\hyperref[\detokenize{compute_temp_deform::doc}]{\sphinxcrossref{\DUrole{doc}{compute temp/deform}}}},
{\hyperref[\detokenize{compute_pressure::doc}]{\sphinxcrossref{\DUrole{doc}{compute pressure}}}}


\subsection{Default}
\label{\detokenize{compute_temp_profile:default}}
\sphinxAtStartPar
The option default is out = tensor.


\bigskip\hrule\bigskip

\phantomsection\label{\detokenize{compute_temp_profile:evans1}}
\sphinxAtStartPar
\sphinxstylestrong{(Evans)} Evans and Morriss, Phys Rev Lett, 56, 2172\sphinxhyphen{}2175 (1986).

\sphinxstepscope

\index{compute temp/ramp@\spxentry{compute temp/ramp}}\ignorespaces 

\section{compute temp/ramp command}
\label{\detokenize{compute_temp_ramp:compute-temp-ramp-command}}\label{\detokenize{compute_temp_ramp:index-0}}\label{\detokenize{compute_temp_ramp::doc}}

\subsection{Syntax}
\label{\detokenize{compute_temp_ramp:syntax}}
\begin{sphinxVerbatim}[commandchars=\\\{\}]
\PYG{k}{compute }\PYG{n+nv+nvIdentifier}{ID}\PYG{+w}{ }\PYG{n+nv+nvIdentifier}{group\PYGZhy{}ID}\PYG{+w}{ }\PYG{n}{temp}\PYG{o}{/}\PYG{n}{ramp}\PYG{+w}{ }\PYG{n}{vdim}\PYG{+w}{ }\PYG{n}{vlo}\PYG{+w}{ }\PYG{n}{vhi}\PYG{+w}{ }\PYG{n}{dim}\PYG{+w}{ }\PYG{n}{clo}\PYG{+w}{ }\PYG{n}{chi}\PYG{+w}{ }\PYG{n}{keyword}\PYG{+w}{ }\PYG{n}{value}\PYG{+w}{ }\PYG{n}{...}
\end{sphinxVerbatim}
\begin{itemize}
\item {} 
\sphinxAtStartPar
ID, group\sphinxhyphen{}ID are documented in {\hyperref[\detokenize{compute::doc}]{\sphinxcrossref{\DUrole{doc}{compute}}}} command

\item {} 
\sphinxAtStartPar
temp/ramp = style name of this compute command

\item {} 
\sphinxAtStartPar
vdim = \sphinxstyleemphasis{vx} or \sphinxstyleemphasis{vy} or \sphinxstyleemphasis{vz}

\item {} 
\sphinxAtStartPar
vlo,vhi = subtract velocities between vlo and vhi (velocity units)

\item {} 
\sphinxAtStartPar
dim = \sphinxstyleemphasis{x} or \sphinxstyleemphasis{y} or \sphinxstyleemphasis{z}

\item {} 
\sphinxAtStartPar
clo,chi = lower and upper bound of domain to subtract from (distance units)

\item {} 
\sphinxAtStartPar
zero or more keyword/value pairs may be appended

\item {} 
\sphinxAtStartPar
keyword = \sphinxstyleemphasis{units}

\end{itemize}
\begin{sphinxalltt}
\sphinxstyleemphasis{units} value = \sphinxstyleemphasis{lattice} or \sphinxstyleemphasis{box}
\end{sphinxalltt}


\subsection{Examples}
\label{\detokenize{compute_temp_ramp:examples}}
\begin{sphinxVerbatim}[commandchars=\\\{\}]
\PYG{k}{compute }\PYG{n+nv+nvIdentifier}{2nd}\PYG{+w}{ }\PYG{n+nv+nvIdentifier}{middle}\PYG{+w}{ }\PYG{n}{temp}\PYG{o}{/}\PYG{n}{ramp}\PYG{+w}{ }\PYG{n}{vx}\PYG{+w}{ }\PYG{l+m}{0}\PYG{+w}{ }\PYG{l+m}{8}\PYG{+w}{ }\PYG{n}{y}\PYG{+w}{ }\PYG{l+m}{2}\PYG{+w}{ }\PYG{l+m}{12}\PYG{+w}{ }\PYG{n}{units}\PYG{+w}{ }\PYG{n}{lattice}
\end{sphinxVerbatim}


\subsection{Description}
\label{\detokenize{compute_temp_ramp:description}}
\sphinxAtStartPar
Define a computation that calculates the temperature of a group of atoms,
after subtracting out an ramped velocity profile before computing the kinetic
energy.  A compute of this style can be used by any command that computes a
temperature (e.g. {\hyperref[\detokenize{thermo_modify::doc}]{\sphinxcrossref{\DUrole{doc}{thermo\_modify}}}},
{\hyperref[\detokenize{fix_temp_rescale::doc}]{\sphinxcrossref{\DUrole{doc}{fix temp/rescale}}}}, {\hyperref[\detokenize{fix_nh::doc}]{\sphinxcrossref{\DUrole{doc}{fix npt}}}}).

\sphinxAtStartPar
The meaning of the arguments for this command which define the
velocity ramp are the same as for the {\hyperref[\detokenize{velocity::doc}]{\sphinxcrossref{\DUrole{doc}{velocity ramp}}}}
command which was presumably used to impose the velocity.

\sphinxAtStartPar
After the ramp velocity has been subtracted from the specified
dimension for each atom, the temperature is calculated by the formula
\begin{equation*}
\begin{split}\text{KE} = \frac{\text{dim}}{2} N k_B T,\end{split}
\end{equation*}
\sphinxAtStartPar
where KE is the total kinetic energy of the group of atoms (sum of
\(\frac12 m v^2\)), dim = 2 or 3 is the dimensionality of the simulation,
\(N\) is the number of atoms in the group, \(k_B\) is the Boltzmann
constant, and \(T\) is the absolute temperature.

\sphinxAtStartPar
The \sphinxstyleemphasis{units} keyword determines the meaning of the distance units used
for coordinates (\sphinxstyleemphasis{clo}, \sphinxstyleemphasis{chi}) and velocities (\sphinxstyleemphasis{vlo}, \sphinxstyleemphasis{vhi}).  A \sphinxstyleemphasis{box} value
selects standard distance units as defined by the {\hyperref[\detokenize{units::doc}]{\sphinxcrossref{\DUrole{doc}{units}}}}
command (e.g., \(\AA\) for units = real or metal).  A
\sphinxstyleemphasis{lattice} value means the distance units are in lattice spacings (i.e.,
velocity in lattice spacings per unit time).  The {\hyperref[\detokenize{lattice::doc}]{\sphinxcrossref{\DUrole{doc}{lattice}}}}
command must have been previously used to define the lattice spacing.

\sphinxAtStartPar
A symmetric tensor, stored as a six\sphinxhyphen{}element vector, is also calculated
by this compute for use in the computation of a pressure tensor by the
{\hyperref[\detokenize{compute_pressure::doc}]{\sphinxcrossref{\DUrole{doc}{compute pressue}}}} command.  The formula for
the components of the tensor is the same as the above expression for
\(E_\mathrm{kin}\), except that the 1/2 factor is NOT included and
the \(v_i^2\) is replaced by \(v_{i,x} v_{i,y}\) for the
\(xy\) component, and so on.  Note that because it lacks the 1/2
factor, these tensor components are twice those of the traditional
kinetic energy tensor.  The six components of the vector are ordered
\(xx\), \(yy\), \(zz\), \(xy\), \(xz\),
\(yz\).

\sphinxAtStartPar
The number of atoms contributing to the temperature is assumed to be constant
for the duration of the run; use the \sphinxstyleemphasis{dynamic} option of the
{\hyperref[\detokenize{compute_modify::doc}]{\sphinxcrossref{\DUrole{doc}{compute\_modify}}}} command if this is not the case.

\sphinxAtStartPar
The removal of the ramped velocity component by this fix is
essentially computing the temperature after a “bias” has been removed
from the velocity of the atoms.  If this compute is used with a fix
command that performs thermostatting then this bias will be subtracted
from each atom, thermostatting of the remaining thermal velocity will
be performed, and the bias will be added back in.  Thermostatting
fixes that work in this way include {\hyperref[\detokenize{fix_nh::doc}]{\sphinxcrossref{\DUrole{doc}{fix nvt}}}},
{\hyperref[\detokenize{fix_temp_rescale::doc}]{\sphinxcrossref{\DUrole{doc}{fix temp/rescale}}}},
{\hyperref[\detokenize{fix_temp_berendsen::doc}]{\sphinxcrossref{\DUrole{doc}{fix temp/berendsen}}}}, and
{\hyperref[\detokenize{fix_langevin::doc}]{\sphinxcrossref{\DUrole{doc}{fix langevin}}}}.

\sphinxAtStartPar
This compute subtracts out degrees\sphinxhyphen{}of\sphinxhyphen{}freedom due to fixes that
constrain molecular motion, such as {\hyperref[\detokenize{fix_shake::doc}]{\sphinxcrossref{\DUrole{doc}{fix shake}}}} and
{\hyperref[\detokenize{fix_rigid::doc}]{\sphinxcrossref{\DUrole{doc}{fix rigid}}}}.  This means the temperature of groups of
atoms that include these constraints will be computed correctly.  If
needed, the subtracted degrees\sphinxhyphen{}of\sphinxhyphen{}freedom can be altered using the
\sphinxstyleemphasis{extra} option of the {\hyperref[\detokenize{compute_modify::doc}]{\sphinxcrossref{\DUrole{doc}{compute\_modify}}}} command.

\sphinxAtStartPar
See the {\hyperref[\detokenize{Howto_thermostat::doc}]{\sphinxcrossref{\DUrole{doc}{Howto thermostat}}}} page for a
discussion of different ways to compute temperature and perform
thermostatting.


\subsection{Output info}
\label{\detokenize{compute_temp_ramp:output-info}}
\sphinxAtStartPar
This compute calculates a global scalar (the temperature) and a global
vector of length 6 (symmetric tensor), which can be accessed by
indices 1\textendash{}6.  These values can be used by any command that uses
global scalar or vector values from a compute as input.  See the
{\hyperref[\detokenize{Howto_output::doc}]{\sphinxcrossref{\DUrole{doc}{Howto output}}}} page for an overview of LAMMPS
output options.

\sphinxAtStartPar
The scalar value calculated by this compute is “intensive”.  The
vector values are “extensive”.

\sphinxAtStartPar
The scalar value is in temperature {\hyperref[\detokenize{units::doc}]{\sphinxcrossref{\DUrole{doc}{units}}}}.  The vector
values are in energy {\hyperref[\detokenize{units::doc}]{\sphinxcrossref{\DUrole{doc}{units}}}}.


\subsection{Restrictions}
\label{\detokenize{compute_temp_ramp:restrictions}}\begin{quote}

\sphinxAtStartPar
none
\end{quote}


\subsection{Related commands}
\label{\detokenize{compute_temp_ramp:related-commands}}
\sphinxAtStartPar
{\hyperref[\detokenize{compute_temp::doc}]{\sphinxcrossref{\DUrole{doc}{compute temp}}}},
{\hyperref[\detokenize{compute_temp_profile::doc}]{\sphinxcrossref{\DUrole{doc}{compute temp/profile}}}},
{\hyperref[\detokenize{compute_temp_deform::doc}]{\sphinxcrossref{\DUrole{doc}{compute temp/deform}}}},
{\hyperref[\detokenize{compute_pressure::doc}]{\sphinxcrossref{\DUrole{doc}{compute pressure}}}}


\subsection{Default}
\label{\detokenize{compute_temp_ramp:default}}
\sphinxAtStartPar
The option default is units = lattice.

\sphinxstepscope

\index{compute temp/region@\spxentry{compute temp/region}}\ignorespaces 

\section{compute temp/region command}
\label{\detokenize{compute_temp_region:compute-temp-region-command}}\label{\detokenize{compute_temp_region:index-0}}\label{\detokenize{compute_temp_region::doc}}

\subsection{Syntax}
\label{\detokenize{compute_temp_region:syntax}}
\begin{sphinxVerbatim}[commandchars=\\\{\}]
\PYG{k}{compute }\PYG{n+nv+nvIdentifier}{ID}\PYG{+w}{ }\PYG{n+nv+nvIdentifier}{group\PYGZhy{}ID}\PYG{+w}{ }\PYG{n}{temp}\PYG{o}{/}\PYG{k}{region }\PYG{n+nv+nvIdentifier}{region\PYGZhy{}ID}
\end{sphinxVerbatim}
\begin{itemize}
\item {} 
\sphinxAtStartPar
ID, group\sphinxhyphen{}ID are documented in {\hyperref[\detokenize{compute::doc}]{\sphinxcrossref{\DUrole{doc}{compute}}}} command

\item {} 
\sphinxAtStartPar
temp/region = style name of this compute command

\item {} 
\sphinxAtStartPar
region\sphinxhyphen{}ID = ID of region to use for choosing atoms

\end{itemize}


\subsection{Examples}
\label{\detokenize{compute_temp_region:examples}}
\begin{sphinxVerbatim}[commandchars=\\\{\}]
\PYG{k}{compute }\PYG{n+nv+nvIdentifier}{mine}\PYG{+w}{ }\PYG{n+nv+nvIdentifier}{flow}\PYG{+w}{ }\PYG{n}{temp}\PYG{o}{/}\PYG{k}{region }\PYG{n+nv+nvIdentifier}{boundary}
\end{sphinxVerbatim}


\subsection{Description}
\label{\detokenize{compute_temp_region:description}}
\sphinxAtStartPar
Define a computation that calculates the temperature of a group of atoms in a
geometric region.  This can be useful for thermostatting one portion of the
simulation box.  For example, a McDLT simulation where one side is cooled, and
the other side is heated.  A compute of this style can be used by any command
that computes a temperature (e.g., {\hyperref[\detokenize{thermo_modify::doc}]{\sphinxcrossref{\DUrole{doc}{thermo\_modify}}}},
{\hyperref[\detokenize{fix_temp_rescale::doc}]{\sphinxcrossref{\DUrole{doc}{fix temp/rescale}}}}).

\sphinxAtStartPar
Note that a \sphinxstyleemphasis{region}\sphinxhyphen{}style temperature can be used to thermostat with
{\hyperref[\detokenize{fix_temp_rescale::doc}]{\sphinxcrossref{\DUrole{doc}{fix temp/rescale}}}} or
{\hyperref[\detokenize{fix_langevin::doc}]{\sphinxcrossref{\DUrole{doc}{fix langevin}}}}, but should probably not be used with
Nose\textendash{}Hoover style fixes ({\hyperref[\detokenize{fix_nh::doc}]{\sphinxcrossref{\DUrole{doc}{fix nvt}}}}, {\hyperref[\detokenize{fix_nh::doc}]{\sphinxcrossref{\DUrole{doc}{fix npt}}}},
or {\hyperref[\detokenize{fix_nh::doc}]{\sphinxcrossref{\DUrole{doc}{fix nph}}}}) if the degrees of freedom included in the computed
temperature vary with time.

\sphinxAtStartPar
The temperature is calculated by the formula
\begin{equation*}
\begin{split}\text{KE} = \frac{\text{dim}}{2} N k_B T,\end{split}
\end{equation*}
\sphinxAtStartPar
where KE = is the total kinetic energy of the group of atoms (sum of
\(\frac12 m v^2\)), dim = 2 or 3 is the dimensionality of the simulation,
\(N\) is the  number of atoms in both the group and region, \(k_B\) is
the Boltzmann constant, and \(T\) temperature.

\sphinxAtStartPar
A symmetric tensor, stored as a six\sphinxhyphen{}element vector, is also calculated
by this compute for use in the computation of a pressure tensor by the
{\hyperref[\detokenize{compute_pressure::doc}]{\sphinxcrossref{\DUrole{doc}{compute pressue}}}} command.  The formula for
the components of the tensor is the same as the above expression for
\(E_\mathrm{kin}\), except that the 1/2 factor is NOT included and
the \(v_i^2\) is replaced by \(v_{i,x} v_{i,y}\) for the
\(xy\) component, and so on.  Note that because it lacks the 1/2
factor, these tensor components are twice those of the traditional
kinetic energy tensor.  The six components of the vector are ordered
\(xx\), \(yy\), \(zz\), \(xy\), \(xz\),
\(yz\).

\sphinxAtStartPar
The number of atoms contributing to the temperature is calculated each
time the temperature is evaluated since it is assumed atoms can
enter/leave the region.  Thus there is no need to use the \sphinxstyleemphasis{dynamic}
option of the {\hyperref[\detokenize{compute_modify::doc}]{\sphinxcrossref{\DUrole{doc}{compute\_modify}}}} command for this
compute style.

\sphinxAtStartPar
The removal of atoms outside the region by this fix is essentially
computing the temperature after a “bias” has been removed, which in
this case is the velocity of any atoms outside the region.  If this
compute is used with a fix command that performs thermostatting then
this bias will be subtracted from each atom, thermostatting of the
remaining thermal velocity will be performed, and the bias will be
added back in.  Thermostatting fixes that work in this way include
{\hyperref[\detokenize{fix_nh::doc}]{\sphinxcrossref{\DUrole{doc}{fix nvt}}}}, {\hyperref[\detokenize{fix_temp_rescale::doc}]{\sphinxcrossref{\DUrole{doc}{fix temp/rescale}}}},
{\hyperref[\detokenize{fix_temp_berendsen::doc}]{\sphinxcrossref{\DUrole{doc}{fix temp/berendsen}}}}, and
{\hyperref[\detokenize{fix_langevin::doc}]{\sphinxcrossref{\DUrole{doc}{fix langevin}}}}.  This means that when this compute
is used to calculate the temperature for any of the thermostatting
fixes via the {\hyperref[\detokenize{fix_modify::doc}]{\sphinxcrossref{\DUrole{doc}{fix modify temp}}}} command, the thermostat
will operate only on atoms that are currently in the geometric region.

\sphinxAtStartPar
Unlike other compute styles that calculate temperature, this compute
does not subtract out degrees\sphinxhyphen{}of\sphinxhyphen{}freedom due to fixes that constrain
motion, such as {\hyperref[\detokenize{fix_shake::doc}]{\sphinxcrossref{\DUrole{doc}{fix shake}}}} and {\hyperref[\detokenize{fix_rigid::doc}]{\sphinxcrossref{\DUrole{doc}{fix rigid}}}}.  This is because those degrees of freedom (e.g., a
constrained bond) could apply to sets of atoms that straddle the
region boundary, and hence the concept is somewhat ill\sphinxhyphen{}defined.  If
needed the number of subtracted degrees of freedom can be set
explicitly using the \sphinxstyleemphasis{extra} option of the {\hyperref[\detokenize{compute_modify::doc}]{\sphinxcrossref{\DUrole{doc}{compute\_modify}}}} command.

\sphinxAtStartPar
See the {\hyperref[\detokenize{Howto_thermostat::doc}]{\sphinxcrossref{\DUrole{doc}{Howto thermostat}}}} page for a
discussion of different ways to compute temperature and perform
thermostatting.


\subsection{Output info}
\label{\detokenize{compute_temp_region:output-info}}
\sphinxAtStartPar
This compute calculates a global scalar (the temperature) and a global
vector of length 6 (symmetric tensor), which can be accessed by
indices 1\textendash{}6.  These values can be used by any command that uses
global scalar or vector values from a compute as input.  See the
{\hyperref[\detokenize{Howto_output::doc}]{\sphinxcrossref{\DUrole{doc}{Howto output}}}} page for an overview of LAMMPS
output options.

\sphinxAtStartPar
The scalar value calculated by this compute is “intensive”.  The
vector values are “extensive”.

\sphinxAtStartPar
The scalar value is in temperature {\hyperref[\detokenize{units::doc}]{\sphinxcrossref{\DUrole{doc}{units}}}}.  The vector
values are in energy {\hyperref[\detokenize{units::doc}]{\sphinxcrossref{\DUrole{doc}{units}}}}.


\subsection{Restrictions}
\label{\detokenize{compute_temp_region:restrictions}}\begin{quote}

\sphinxAtStartPar
none
\end{quote}


\subsection{Related commands}
\label{\detokenize{compute_temp_region:related-commands}}
\sphinxAtStartPar
{\hyperref[\detokenize{compute_temp::doc}]{\sphinxcrossref{\DUrole{doc}{compute temp}}}}, {\hyperref[\detokenize{compute_pressure::doc}]{\sphinxcrossref{\DUrole{doc}{compute pressure}}}}


\subsection{Default}
\label{\detokenize{compute_temp_region:default}}
\sphinxAtStartPar
none

\sphinxstepscope

\index{compute temp/region/eff@\spxentry{compute temp/region/eff}}\ignorespaces 

\section{compute temp/region/eff command}
\label{\detokenize{compute_temp_region_eff:compute-temp-region-eff-command}}\label{\detokenize{compute_temp_region_eff:index-0}}\label{\detokenize{compute_temp_region_eff::doc}}

\subsection{Syntax}
\label{\detokenize{compute_temp_region_eff:syntax}}
\begin{sphinxVerbatim}[commandchars=\\\{\}]
\PYG{k}{compute }\PYG{n+nv+nvIdentifier}{ID}\PYG{+w}{ }\PYG{n+nv+nvIdentifier}{group\PYGZhy{}ID}\PYG{+w}{ }\PYG{n}{temp}\PYG{o}{/}\PYG{n}{region}\PYG{o}{/}\PYG{n}{eff}\PYG{+w}{ }\PYG{n}{region}\PYG{o}{\PYGZhy{}}\PYG{n}{ID}
\end{sphinxVerbatim}
\begin{itemize}
\item {} 
\sphinxAtStartPar
ID, group\sphinxhyphen{}ID are documented in {\hyperref[\detokenize{compute::doc}]{\sphinxcrossref{\DUrole{doc}{compute}}}} command

\item {} 
\sphinxAtStartPar
temp/region/eff = style name of this compute command

\item {} 
\sphinxAtStartPar
region\sphinxhyphen{}ID = ID of region to use for choosing atoms

\end{itemize}


\subsection{Examples}
\label{\detokenize{compute_temp_region_eff:examples}}
\begin{sphinxVerbatim}[commandchars=\\\{\}]
\PYG{k}{compute }\PYG{n+nv+nvIdentifier}{mine}\PYG{+w}{ }\PYG{n+nv+nvIdentifier}{flow}\PYG{+w}{ }\PYG{n}{temp}\PYG{o}{/}\PYG{n}{region}\PYG{o}{/}\PYG{n}{eff}\PYG{+w}{ }\PYG{n}{boundary}
\end{sphinxVerbatim}


\subsection{Description}
\label{\detokenize{compute_temp_region_eff:description}}
\sphinxAtStartPar
Define a computation that calculates the temperature of a group of
nuclei and electrons in the {\hyperref[\detokenize{pair_eff::doc}]{\sphinxcrossref{\DUrole{doc}{electron force field}}}}
model, within a geometric region using the electron force field.
A compute of this style can be used by commands that compute a
temperature (e.g., {\hyperref[\detokenize{thermo_modify::doc}]{\sphinxcrossref{\DUrole{doc}{thermo\_modify}}}}).

\sphinxAtStartPar
The operation of this compute is exactly like that described by the
{\hyperref[\detokenize{compute_temp_region::doc}]{\sphinxcrossref{\DUrole{doc}{compute temp/region}}}} command, except that
the formulas for the temperature (scalar) and diagonal components of
the symmetric tensor (vector) include the radial electron velocity
contributions, as discussed by the {\hyperref[\detokenize{compute_temp_eff::doc}]{\sphinxcrossref{\DUrole{doc}{compute temp/eff}}}} command.


\subsection{Output info}
\label{\detokenize{compute_temp_region_eff:output-info}}
\sphinxAtStartPar
This compute calculates a global scalar (the temperature) and a global
vector of length 6 (symmetric tensor), which can be accessed by
indices 1\textendash{}6.  These values can be used by any command that uses
global scalar or vector values from a compute as input.  See the
{\hyperref[\detokenize{Howto_output::doc}]{\sphinxcrossref{\DUrole{doc}{Howto output}}}} page for an overview of LAMMPS
output options.

\sphinxAtStartPar
The scalar value calculated by this compute is “intensive”.  The
vector values are “extensive”.

\sphinxAtStartPar
The scalar value is in temperature {\hyperref[\detokenize{units::doc}]{\sphinxcrossref{\DUrole{doc}{units}}}}.  The vector
values are in energy {\hyperref[\detokenize{units::doc}]{\sphinxcrossref{\DUrole{doc}{units}}}}.


\subsection{Restrictions}
\label{\detokenize{compute_temp_region_eff:restrictions}}
\sphinxAtStartPar
This compute is part of the EFF package.  It is only enabled if LAMMPS
was built with that package.  See the {\hyperref[\detokenize{Build_package::doc}]{\sphinxcrossref{\DUrole{doc}{Build package}}}} page for more info.


\subsection{Related commands}
\label{\detokenize{compute_temp_region_eff:related-commands}}
\sphinxAtStartPar
{\hyperref[\detokenize{compute_temp_region::doc}]{\sphinxcrossref{\DUrole{doc}{compute temp/region}}}},
{\hyperref[\detokenize{compute_temp_eff::doc}]{\sphinxcrossref{\DUrole{doc}{compute temp/eff}}}},
{\hyperref[\detokenize{compute_pressure::doc}]{\sphinxcrossref{\DUrole{doc}{compute pressure}}}}


\subsection{Default}
\label{\detokenize{compute_temp_region_eff:default}}
\sphinxAtStartPar
none

\sphinxstepscope

\index{compute temp/rotate@\spxentry{compute temp/rotate}}\ignorespaces 

\section{compute temp/rotate command}
\label{\detokenize{compute_temp_rotate:compute-temp-rotate-command}}\label{\detokenize{compute_temp_rotate:index-0}}\label{\detokenize{compute_temp_rotate::doc}}

\subsection{Syntax}
\label{\detokenize{compute_temp_rotate:syntax}}
\begin{sphinxVerbatim}[commandchars=\\\{\}]
\PYG{k}{compute }\PYG{n+nv+nvIdentifier}{ID}\PYG{+w}{ }\PYG{n+nv+nvIdentifier}{group\PYGZhy{}ID}\PYG{+w}{ }\PYG{n}{temp}\PYG{o}{/}\PYG{n}{rotate}
\end{sphinxVerbatim}
\begin{itemize}
\item {} 
\sphinxAtStartPar
ID, group\sphinxhyphen{}ID are documented in {\hyperref[\detokenize{compute::doc}]{\sphinxcrossref{\DUrole{doc}{compute}}}} command

\item {} 
\sphinxAtStartPar
temp/rotate = style name of this compute command

\end{itemize}


\subsection{Examples}
\label{\detokenize{compute_temp_rotate:examples}}
\begin{sphinxVerbatim}[commandchars=\\\{\}]
\PYG{k}{compute }\PYG{n+nv+nvIdentifier}{Tbead}\PYG{+w}{ }\PYG{n+nv+nvIdentifier}{bead}\PYG{+w}{ }\PYG{n}{temp}\PYG{o}{/}\PYG{n}{rotate}
\end{sphinxVerbatim}


\subsection{Description}
\label{\detokenize{compute_temp_rotate:description}}
\sphinxAtStartPar
Define a computation that calculates the temperature of a group of
atoms, after subtracting out the center\sphinxhyphen{}of\sphinxhyphen{}mass velocity and angular velocity
of the group.  This is useful if the group is expected to have a non\sphinxhyphen{}zero net
velocity and/or global rotation motion for some reason.  A compute of this
style can be used by any command that computes a temperature
(e.g., {\hyperref[\detokenize{thermo_modify::doc}]{\sphinxcrossref{\DUrole{doc}{thermo\_modify}}}},
{\hyperref[\detokenize{fix_temp_rescale::doc}]{\sphinxcrossref{\DUrole{doc}{fix temp/rescale}}}}, {\hyperref[\detokenize{fix_nh::doc}]{\sphinxcrossref{\DUrole{doc}{fix npt}}}}).

\sphinxAtStartPar
After the center\sphinxhyphen{}of\sphinxhyphen{}mass velocity and angular velocity has been subtracted from
each atom, the temperature is calculated by the formula
\begin{equation*}
\begin{split}\text{KE} = \frac{\text{dim}}{2} N k_B T,\end{split}
\end{equation*}
\sphinxAtStartPar
where KE is the total kinetic energy of the group of atoms (sum of
\(\frac12 m v^2\)), dim = 2 or 3 is the dimensionality of the simulation,
\(N\) is the number of atoms in the group, \(k_B\) is the Boltzmann
constant, and \(T\) is the absolute temperature.

\sphinxAtStartPar
A symmetric tensor, stored as a six\sphinxhyphen{}element vector, is also calculated
by this compute for use in the computation of a pressure tensor by the
{\hyperref[\detokenize{compute_pressure::doc}]{\sphinxcrossref{\DUrole{doc}{compute pressue}}}} command.  The formula for
the components of the tensor is the same as the above expression for
\(E_\mathrm{kin}\), except that the 1/2 factor is NOT included and
the \(v_i^2\) is replaced by \(v_{i,x} v_{i,y}\) for the
\(xy\) component, and so on.  Note that because it lacks the 1/2
factor, these tensor components are twice those of the traditional
kinetic energy tensor.  The six components of the vector are ordered
\(xx\), \(yy\), \(zz\), \(xy\), \(xz\),
\(yz\).

\sphinxAtStartPar
The number of atoms contributing to the temperature is assumed to be
constant for the duration of the run; use the \sphinxstyleemphasis{dynamic} option of the
{\hyperref[\detokenize{compute_modify::doc}]{\sphinxcrossref{\DUrole{doc}{compute\_modify}}}} command if this is not the case.

\sphinxAtStartPar
The removal of the center\sphinxhyphen{}of\sphinxhyphen{}mass velocity and angular velocity by this fix is
essentially computing the temperature after a “bias” has been removed from the
velocity of the atoms.  If this compute is used with a fix command
that performs thermostatting then this bias will be subtracted from
each atom, thermostatting of the remaining thermal velocity will be
performed, and the bias will be added back in.  Thermostatting fixes
that work in this way include {\hyperref[\detokenize{fix_nh::doc}]{\sphinxcrossref{\DUrole{doc}{fix nvt}}}},
{\hyperref[\detokenize{fix_temp_rescale::doc}]{\sphinxcrossref{\DUrole{doc}{fix temp/rescale}}}},
{\hyperref[\detokenize{fix_temp_berendsen::doc}]{\sphinxcrossref{\DUrole{doc}{fix temp/berendsen}}}}, and
{\hyperref[\detokenize{fix_langevin::doc}]{\sphinxcrossref{\DUrole{doc}{fix langevin}}}}.

\sphinxAtStartPar
This compute subtracts out degrees\sphinxhyphen{}of\sphinxhyphen{}freedom due to fixes that
constrain molecular motion, such as {\hyperref[\detokenize{fix_shake::doc}]{\sphinxcrossref{\DUrole{doc}{fix shake}}}} and
{\hyperref[\detokenize{fix_rigid::doc}]{\sphinxcrossref{\DUrole{doc}{fix rigid}}}}.  This means the temperature of groups of
atoms that include these constraints will be computed correctly.  If
needed, the subtracted degrees\sphinxhyphen{}of\sphinxhyphen{}freedom can be altered using the
\sphinxstyleemphasis{extra} option of the {\hyperref[\detokenize{compute_modify::doc}]{\sphinxcrossref{\DUrole{doc}{compute\_modify}}}} command.

\sphinxAtStartPar
See the {\hyperref[\detokenize{Howto_thermostat::doc}]{\sphinxcrossref{\DUrole{doc}{Howto thermostat}}}} page for a
discussion of different ways to compute temperature and perform
thermostatting.


\subsection{Output info}
\label{\detokenize{compute_temp_rotate:output-info}}
\sphinxAtStartPar
This compute calculates a global scalar (the temperature) and a global
vector of length 6 (symmetric tensor), which can be accessed by
indices 1\sphinxhyphen{}6.  These values can be used by any command that uses global
scalar or vector values from a compute as input.  See the {\hyperref[\detokenize{Howto_output::doc}]{\sphinxcrossref{\DUrole{doc}{Howto
output}}}} page for an overview of LAMMPS output options.

\sphinxAtStartPar
The scalar value calculated by this compute is “intensive”.  The
vector values are “extensive”.

\sphinxAtStartPar
The scalar value is in temperature {\hyperref[\detokenize{units::doc}]{\sphinxcrossref{\DUrole{doc}{units}}}}.  The vector
values are in energy {\hyperref[\detokenize{units::doc}]{\sphinxcrossref{\DUrole{doc}{units}}}}.


\subsection{Restrictions}
\label{\detokenize{compute_temp_rotate:restrictions}}
\sphinxAtStartPar
This compute is part of the EXTRA\sphinxhyphen{}COMPUTE package.  It is only enabled if
LAMMPS was built with that package.
See the {\hyperref[\detokenize{Build_package::doc}]{\sphinxcrossref{\DUrole{doc}{Build package}}}} page for more info.


\subsection{Related commands}
\label{\detokenize{compute_temp_rotate:related-commands}}
\sphinxAtStartPar
{\hyperref[\detokenize{compute_temp::doc}]{\sphinxcrossref{\DUrole{doc}{compute temp}}}}


\subsection{Default}
\label{\detokenize{compute_temp_rotate:default}}
\sphinxAtStartPar
none

\sphinxstepscope

\index{compute temp/sphere@\spxentry{compute temp/sphere}}\ignorespaces 

\section{compute temp/sphere command}
\label{\detokenize{compute_temp_sphere:compute-temp-sphere-command}}\label{\detokenize{compute_temp_sphere:index-0}}\label{\detokenize{compute_temp_sphere::doc}}

\subsection{Syntax}
\label{\detokenize{compute_temp_sphere:syntax}}
\begin{sphinxVerbatim}[commandchars=\\\{\}]
\PYG{k}{compute }\PYG{n+nv+nvIdentifier}{ID}\PYG{+w}{ }\PYG{n+nv+nvIdentifier}{group\PYGZhy{}ID}\PYG{+w}{ }\PYG{n}{temp}\PYG{o}{/}\PYG{n}{sphere}\PYG{+w}{ }\PYG{n}{keyword}\PYG{+w}{ }\PYG{n}{value}\PYG{+w}{ }\PYG{n}{...}
\end{sphinxVerbatim}
\begin{itemize}
\item {} 
\sphinxAtStartPar
ID, group\sphinxhyphen{}ID are documented in {\hyperref[\detokenize{compute::doc}]{\sphinxcrossref{\DUrole{doc}{compute}}}} command

\item {} 
\sphinxAtStartPar
temp/sphere = style name of this compute command

\item {} 
\sphinxAtStartPar
zero or more keyword/value pairs may be appended

\item {} 
\sphinxAtStartPar
keyword = \sphinxstyleemphasis{bias} or \sphinxstyleemphasis{dof}
\begin{sphinxalltt}
\sphinxstyleemphasis{bias} value = bias\sphinxhyphen{}ID
  bias\sphinxhyphen{}ID = ID of a temperature compute that removes a velocity bias
\sphinxstyleemphasis{dof} value = \sphinxstyleemphasis{all} or \sphinxstyleemphasis{rotate}
  all = compute temperature of translational and rotational degrees of freedom
  rotate = compute temperature of just rotational degrees of freedom
\end{sphinxalltt}

\end{itemize}


\subsection{Examples}
\label{\detokenize{compute_temp_sphere:examples}}
\begin{sphinxVerbatim}[commandchars=\\\{\}]
\PYG{k}{compute }\PYG{n+nv+nvIdentifier}{1}\PYG{+w}{ }\PYG{n+nv+nvIdentifier}{all}\PYG{+w}{ }\PYG{n}{temp}\PYG{o}{/}\PYG{n}{sphere}
\PYG{k}{compute }\PYG{n+nv+nvIdentifier}{myTemp}\PYG{+w}{ }\PYG{n+nv+nvIdentifier}{mobile}\PYG{+w}{ }\PYG{n}{temp}\PYG{o}{/}\PYG{n}{sphere}\PYG{+w}{ }\PYG{n}{bias}\PYG{+w}{ }\PYG{n}{tempCOM}
\PYG{k}{compute }\PYG{n+nv+nvIdentifier}{myTemp}\PYG{+w}{ }\PYG{n+nv+nvIdentifier}{mobile}\PYG{+w}{ }\PYG{n}{temp}\PYG{o}{/}\PYG{n}{sphere}\PYG{+w}{ }\PYG{n}{dof}\PYG{+w}{ }\PYG{n}{rotate}
\end{sphinxVerbatim}


\subsection{Description}
\label{\detokenize{compute_temp_sphere:description}}
\sphinxAtStartPar
Define a computation that calculates the temperature of a group of
spherical particles, including a contribution from both their
translational and rotational kinetic energy.  This differs from the
usual {\hyperref[\detokenize{compute_temp::doc}]{\sphinxcrossref{\DUrole{doc}{compute temp}}}} command, which assumes point
particles with only translational kinetic energy.

\sphinxAtStartPar
Both point and finite\sphinxhyphen{}size particles can be included in the group.
Point particles do not rotate, so they have only three translational
degrees of freedom.  For 3d spherical particles, each has six degrees of
freedom (three translational, three rotational).  For 2d spherical particles,
each has three degrees of freedom (two translational, one rotational).

\begin{sphinxadmonition}{note}{Note:}
\sphinxAtStartPar
This choice for degrees of freedom (DOF) assumes that all finite\sphinxhyphen{}size
spherical particles in your model will freely rotate, sampling all
their rotational DOF.  It is possible to use a combination of
interaction potentials and fixes that induce no torque or otherwise
constrain some of all of your particles so that this is not the case.
Then there are less DOF and you should use the {\hyperref[\detokenize{compute_modify::doc}]{\sphinxcrossref{\DUrole{doc}{compute\_modify
extra/dof}}}} command to adjust the DOF accordingly.
\end{sphinxadmonition}

\sphinxAtStartPar
The translational kinetic energy is computed the same as is described
by the {\hyperref[\detokenize{compute_temp::doc}]{\sphinxcrossref{\DUrole{doc}{compute temp}}}} command.  The rotational
kinetic energy is computed as \(\frac12 I \omega^2\), where \(I\) is
the moment of inertia for a sphere and \(\omega\) is the particle’s angular
velocity.

\begin{sphinxadmonition}{note}{Note:}
\sphinxAtStartPar
For {\hyperref[\detokenize{dimension::doc}]{\sphinxcrossref{\DUrole{doc}{2d models}}}}, particles are treated as
spheres, not disks, meaning their moment of inertia will be the same
as in 3d.
\end{sphinxadmonition}

\sphinxAtStartPar
A kinetic energy tensor, stored as a six\sphinxhyphen{}element vector, is also
calculated by this compute.  The formula for the components of the
tensor is the same as the above formulas, except that \(v^2\) and
\(\omega^2\) are replaced by \(v_x v_y\) and \(\omega_x
\omega_y\) for the \(xy\) component.  The six components of the
vector are ordered \(xx\), \(yy\), \(zz\), \(xy\),
\(xz\), \(yz\).

\sphinxAtStartPar
A symmetric tensor, stored as a six\sphinxhyphen{}element vector, is also calculated
by this compute for use in the computation of a pressure tensor by the
{\hyperref[\detokenize{compute_pressure::doc}]{\sphinxcrossref{\DUrole{doc}{compute pressue}}}} command.  The formula for
the components of the tensor is the same as the above expression for
\(E_\mathrm{kin}\), except that the 1/2 factor is NOT included and
the \(v_i^2\) and \(\omega^2\) are replaced by \(v_x v_y\)
and \(\omega_x \omega_y\) for the \(xy\) component, and so on.
Note that because it lacks the 1/2 factor, these tensor components are
twice those of the traditional kinetic energy tensor.  The six
components of the vector are ordered \(xx\), \(yy\),
\(zz\), \(xy\), \(xz\), \(yz\).

\sphinxAtStartPar
The number of atoms contributing to the temperature is assumed to be
constant for the duration of the run; use the \sphinxstyleemphasis{dynamic} option of the
{\hyperref[\detokenize{compute_modify::doc}]{\sphinxcrossref{\DUrole{doc}{compute\_modify}}}} command if this is not the case.

\sphinxAtStartPar
This compute subtracts out translational degrees\sphinxhyphen{}of\sphinxhyphen{}freedom due to fixes
that constrain molecular motion, such as {\hyperref[\detokenize{fix_shake::doc}]{\sphinxcrossref{\DUrole{doc}{fix shake}}}}
and {\hyperref[\detokenize{fix_rigid::doc}]{\sphinxcrossref{\DUrole{doc}{fix rigid}}}}.  This means the temperature of groups
of atoms that include these constraints will be computed correctly.  If
needed, the subtracted degrees of freedom can be altered using the
\sphinxstyleemphasis{extra/dof} option of the {\hyperref[\detokenize{compute_modify::doc}]{\sphinxcrossref{\DUrole{doc}{compute\_modify}}}}
command.

\sphinxAtStartPar
See the {\hyperref[\detokenize{Howto_thermostat::doc}]{\sphinxcrossref{\DUrole{doc}{Howto thermostat}}}} page for a
discussion of different ways to compute temperature and perform
thermostatting.


\bigskip\hrule\bigskip


\sphinxAtStartPar
The keyword/value option pairs are used in the following ways.

\sphinxAtStartPar
For the \sphinxstyleemphasis{bias} keyword, \sphinxstyleemphasis{bias\sphinxhyphen{}ID} refers to the ID of a temperature
compute that removes a “bias” velocity from each atom.  This allows
compute temp/sphere to compute its thermal temperature after the
translational kinetic energy components have been altered in a
prescribed way (e.g., to remove a flow velocity profile).  Thermostats
that use this compute will work with this bias term.  See the doc
pages for individual computes that calculate a temperature and the doc
pages for fixes that perform thermostatting for more details.

\sphinxAtStartPar
For the \sphinxstyleemphasis{dof} keyword, a setting of \sphinxstyleemphasis{all} calculates a temperature
that includes both translational and rotational degrees of freedom.
A setting of \sphinxstyleemphasis{rotate} calculates a temperature that includes only
rotational degrees of freedom.


\bigskip\hrule\bigskip



\subsection{Output info}
\label{\detokenize{compute_temp_sphere:output-info}}
\sphinxAtStartPar
This compute calculates a global scalar (the temperature) and a global
vector of length 6 (symmetric tensor), which can be accessed by
indices 1\textendash{}6.  These values can be used by any command that uses
global scalar or vector values from a compute as input.  See the
{\hyperref[\detokenize{Howto_output::doc}]{\sphinxcrossref{\DUrole{doc}{Howto output}}}} page for an overview of LAMMPS
output options.

\sphinxAtStartPar
The scalar value calculated by this compute is “intensive”.  The
vector values are “extensive”.

\sphinxAtStartPar
The scalar value is in temperature {\hyperref[\detokenize{units::doc}]{\sphinxcrossref{\DUrole{doc}{units}}}}.  The vector
values are in energy {\hyperref[\detokenize{units::doc}]{\sphinxcrossref{\DUrole{doc}{units}}}}.


\subsection{Restrictions}
\label{\detokenize{compute_temp_sphere:restrictions}}
\sphinxAtStartPar
This fix requires that atoms store torque and angular velocity (omega)
and a radius as defined by the {\hyperref[\detokenize{atom_style::doc}]{\sphinxcrossref{\DUrole{doc}{atom\_style sphere}}}}
command.

\sphinxAtStartPar
All particles in the group must be finite\sphinxhyphen{}size spheres, or point
particles with radius = 0.0.


\subsection{Related commands}
\label{\detokenize{compute_temp_sphere:related-commands}}
\sphinxAtStartPar
{\hyperref[\detokenize{compute_temp::doc}]{\sphinxcrossref{\DUrole{doc}{compute temp}}}}, {\hyperref[\detokenize{compute_temp::doc}]{\sphinxcrossref{\DUrole{doc}{compute temp/asphere}}}}


\subsection{Default}
\label{\detokenize{compute_temp_sphere:default}}
\sphinxAtStartPar
The option defaults are no bias and dof = all.

\sphinxstepscope

\index{compute temp/uef@\spxentry{compute temp/uef}}\ignorespaces 

\section{compute temp/uef command}
\label{\detokenize{compute_temp_uef:compute-temp-uef-command}}\label{\detokenize{compute_temp_uef:index-0}}\label{\detokenize{compute_temp_uef::doc}}

\subsection{Syntax}
\label{\detokenize{compute_temp_uef:syntax}}
\begin{sphinxVerbatim}[commandchars=\\\{\}]
\PYG{k}{compute }\PYG{n+nv+nvIdentifier}{ID}\PYG{+w}{ }\PYG{n+nv+nvIdentifier}{group\PYGZhy{}ID}\PYG{+w}{ }\PYG{n}{temp}\PYG{o}{/}\PYG{n}{uef}
\end{sphinxVerbatim}
\begin{itemize}
\item {} 
\sphinxAtStartPar
ID, group\sphinxhyphen{}ID are documented in {\hyperref[\detokenize{compute::doc}]{\sphinxcrossref{\DUrole{doc}{compute}}}} command

\item {} 
\sphinxAtStartPar
temp/uef = style name of this compute command

\end{itemize}


\subsection{Examples}
\label{\detokenize{compute_temp_uef:examples}}
\begin{sphinxVerbatim}[commandchars=\\\{\}]
\PYG{k}{compute }\PYG{n+nv+nvIdentifier}{1}\PYG{+w}{ }\PYG{n+nv+nvIdentifier}{all}\PYG{+w}{ }\PYG{n}{temp}\PYG{o}{/}\PYG{n}{uef}
\PYG{k}{compute }\PYG{n+nv+nvIdentifier}{2}\PYG{+w}{ }\PYG{n+nv+nvIdentifier}{sel}\PYG{+w}{ }\PYG{n}{temp}\PYG{o}{/}\PYG{n}{uef}
\end{sphinxVerbatim}


\subsection{Description}
\label{\detokenize{compute_temp_uef:description}}
\sphinxAtStartPar
This command is used to compute the kinetic energy tensor in
the reference frame of the applied flow field when
{\hyperref[\detokenize{fix_nh_uef::doc}]{\sphinxcrossref{\DUrole{doc}{fix nvt/uef}}}} or
{\hyperref[\detokenize{fix_nh_uef::doc}]{\sphinxcrossref{\DUrole{doc}{fix npt/uef}}}} is used.
It is not necessary to use this command to compute the scalar
value of the temperature. A {\hyperref[\detokenize{compute_temp::doc}]{\sphinxcrossref{\DUrole{doc}{compute temp}}}}
may be used for that purpose.

\sphinxAtStartPar
Output information for this command can be found in the
documentation for {\hyperref[\detokenize{compute_temp::doc}]{\sphinxcrossref{\DUrole{doc}{compute temp}}}}.


\subsection{Restrictions}
\label{\detokenize{compute_temp_uef:restrictions}}
\sphinxAtStartPar
This fix is part of the UEF package. It is only enabled if LAMMPS was built
with that package. See the {\hyperref[\detokenize{Build_package::doc}]{\sphinxcrossref{\DUrole{doc}{Build package}}}} page for more
info.

\sphinxAtStartPar
This command can only be used when {\hyperref[\detokenize{fix_nh_uef::doc}]{\sphinxcrossref{\DUrole{doc}{fix nvt/uef}}}}
or {\hyperref[\detokenize{fix_nh_uef::doc}]{\sphinxcrossref{\DUrole{doc}{fix npt/uef}}}} is active.


\subsection{Related commands}
\label{\detokenize{compute_temp_uef:related-commands}}
\sphinxAtStartPar
{\hyperref[\detokenize{compute_temp::doc}]{\sphinxcrossref{\DUrole{doc}{compute temp}}}},
{\hyperref[\detokenize{fix_nh_uef::doc}]{\sphinxcrossref{\DUrole{doc}{fix nvt/uef}}}},
{\hyperref[\detokenize{compute_pressure_uef::doc}]{\sphinxcrossref{\DUrole{doc}{compute pressure/uef}}}}


\subsection{Default}
\label{\detokenize{compute_temp_uef:default}}
\sphinxAtStartPar
none

\sphinxstepscope

\index{compute ti@\spxentry{compute ti}}\ignorespaces 

\section{compute ti command}
\label{\detokenize{compute_ti:compute-ti-command}}\label{\detokenize{compute_ti:index-0}}\label{\detokenize{compute_ti::doc}}

\subsection{Syntax}
\label{\detokenize{compute_ti:syntax}}
\begin{sphinxVerbatim}[commandchars=\\\{\}]
\PYG{k}{compute }\PYG{n+nv+nvIdentifier}{ID}\PYG{+w}{ }\PYG{n+nv+nvIdentifier}{group}\PYG{+w}{ }\PYG{n}{ti}\PYG{+w}{ }\PYG{n}{keyword}\PYG{+w}{ }\PYG{n}{args}\PYG{+w}{ }\PYG{n}{...}
\end{sphinxVerbatim}
\begin{itemize}
\item {} 
\sphinxAtStartPar
ID, group\sphinxhyphen{}ID are documented in {\hyperref[\detokenize{compute::doc}]{\sphinxcrossref{\DUrole{doc}{compute}}}} command

\item {} 
\sphinxAtStartPar
ti = style name of this compute command

\item {} 
\sphinxAtStartPar
one or more attribute/arg pairs may be appended

\item {} 
\sphinxAtStartPar
keyword = pair style (lj/cut, gauss, born, etc.) or \sphinxstyleemphasis{tail} or \sphinxstyleemphasis{kspace}
\begin{sphinxalltt}
pair style args = atype v\_name1 v\_name2
  atype = atom type (see asterisk form below)
  v\_name1 = variable with name1 that is energy scale factor and function of lambda
  v\_name2 = variable with name2 that is derivative of v\_name1 with respect to lambda
\sphinxstyleemphasis{tail} args = atype v\_name1 v\_name2
  atype = atom type (see asterisk form below)
  v\_name1 = variable with name1 that is energy tail correction scale factor and function of lambda
  v\_name2 = variable with name2 that is derivative of v\_name1 with respect to lambda
\sphinxstyleemphasis{kspace} args = atype v\_name1 v\_name2
  atype = atom type (see asterisk form below)
  v\_name1 = variable with name1 that is K\sphinxhyphen{}Space scale factor and function of lambda
  v\_name2 = variable with name2 that is derivative of v\_name1 with respect to lambda
\end{sphinxalltt}

\end{itemize}


\subsection{Examples}
\label{\detokenize{compute_ti:examples}}
\begin{sphinxVerbatim}[commandchars=\\\{\}]
\PYG{k}{compute }\PYG{n+nv+nvIdentifier}{1}\PYG{+w}{ }\PYG{n+nv+nvIdentifier}{all}\PYG{+w}{ }\PYG{n}{ti}\PYG{+w}{ }\PYG{n}{lj}\PYG{o}{/}\PYG{n}{cut}\PYG{+w}{ }\PYG{l+m}{1}\PYG{+w}{ }\PYG{n}{v\PYGZus{}lj}\PYG{+w}{ }\PYG{n}{v\PYGZus{}dlj}\PYG{+w}{ }\PYG{n}{coul}\PYG{o}{/}\PYG{n}{long}\PYG{+w}{ }\PYG{l+m}{2}\PYG{+w}{ }\PYG{n}{v\PYGZus{}c}\PYG{+w}{ }\PYG{n}{v\PYGZus{}dc}\PYG{+w}{ }\PYG{n}{kspace}\PYG{+w}{ }\PYG{l+m}{1}\PYG{+w}{ }\PYG{n}{v\PYGZus{}ks}\PYG{+w}{ }\PYG{n}{v\PYGZus{}dks}
\PYG{k}{compute }\PYG{n+nv+nvIdentifier}{1}\PYG{+w}{ }\PYG{n+nv+nvIdentifier}{all}\PYG{+w}{ }\PYG{n}{ti}\PYG{+w}{ }\PYG{n}{lj}\PYG{o}{/}\PYG{n}{cut}\PYG{+w}{ }\PYG{l+m}{1}\PYG{o}{*}\PYG{l+m}{3}\PYG{+w}{ }\PYG{n}{v\PYGZus{}lj}\PYG{+w}{ }\PYG{n}{v\PYGZus{}dlj}\PYG{+w}{ }\PYG{n}{coul}\PYG{o}{/}\PYG{n}{long}\PYG{+w}{ }\PYG{o}{*}\PYG{+w}{ }\PYG{n}{v\PYGZus{}c}\PYG{+w}{ }\PYG{n}{v\PYGZus{}dc}\PYG{+w}{ }\PYG{n}{kspace}\PYG{+w}{ }\PYG{o}{*}\PYG{+w}{ }\PYG{n}{v\PYGZus{}ks}\PYG{+w}{ }\PYG{n}{v\PYGZus{}dks}
\end{sphinxVerbatim}


\subsection{Description}
\label{\detokenize{compute_ti:description}}
\sphinxAtStartPar
Define a computation that calculates the derivative of the interaction
potential with respect to \sphinxstyleemphasis{lambda}, the coupling parameter used in a
thermodynamic integration.  This derivative can be used to infer a
free energy difference resulting from an alchemical simulation, as
described in {\hyperref[\detokenize{compute_ti:eike}]{\sphinxcrossref{\DUrole{std,std-ref}{Eike}}}}.

\sphinxAtStartPar
Typically this compute will be used in conjunction with the
{\hyperref[\detokenize{fix_adapt::doc}]{\sphinxcrossref{\DUrole{doc}{fix adapt}}}} command which can perform alchemical
transformations by adjusting the strength of an interaction potential
as a simulation runs, as defined by one or more
{\hyperref[\detokenize{pair_style::doc}]{\sphinxcrossref{\DUrole{doc}{pair\_style}}}} or {\hyperref[\detokenize{kspace_style::doc}]{\sphinxcrossref{\DUrole{doc}{kspace\_style}}}}
commands.  This scaling is done via a prefactor on the energy, forces,
virial calculated by the pair or \(k\)\sphinxhyphen{}space style.  The prefactor is
often a function of a \sphinxstyleemphasis{lambda} parameter which may be adjusted from 0 to 1
(or vice versa) over the course of a {\hyperref[\detokenize{run::doc}]{\sphinxcrossref{\DUrole{doc}{run}}}}.
The time\sphinxhyphen{}dependent adjustment is what the {\hyperref[\detokenize{fix_adapt::doc}]{\sphinxcrossref{\DUrole{doc}{fix adapt}}}}
command does.

\sphinxAtStartPar
Assume that the unscaled energy of a pair\_style or kspace\_style is
given by \(U\).  Then the scaled energy is
\begin{equation*}
\begin{split}U_s = f(\lambda) U\end{split}
\end{equation*}
\sphinxAtStartPar
where \(f\) is some function of \(\lambda\).  What this compute
calculates is
\begin{equation*}
\begin{split}\frac{dU_s}{d\lambda} = U \frac{df(\lambda)}{d\lambda}
  = \frac{U_s}{f(\lambda)} \frac{df(\lambda)}{d\lambda},\end{split}
\end{equation*}
\sphinxAtStartPar
which is the derivative of the system’s scaled potential energy \(U_s\)
with respect to \(\lambda\).

\sphinxAtStartPar
To perform this calculation, you provide one or more atom types as
\sphinxstyleemphasis{atype}.  The variable \sphinxstyleemphasis{atype} can be specified in one of two ways.
An explicit numeric value can be used, as in the first example above, or a
wildcard asterisk can be used in place of or in conjunction with the
\sphinxstyleemphasis{atype} argument to select multiple atom types.  This takes the form
“*” or “*n” or “m*” or “m*n”.  If \(N\) is the number of atom types,
then an asterisk with no numeric values means all types from 1 to \(N\).
A leading asterisk means all types from 1 to n (inclusive).  A trailing
asterisk means all types from m to N (inclusive).  A middle asterisk
means all types from m to n (inclusive).

\sphinxAtStartPar
You also specify two functions, as {\hyperref[\detokenize{variable::doc}]{\sphinxcrossref{\DUrole{doc}{equal\sphinxhyphen{}style variables}}}}.
The first is specified as \sphinxstyleemphasis{v\_name1}, where \sphinxstyleemphasis{name1} is the name of the
variable, and is \(f(\lambda)\) in the notation above.  The second is
specified as \sphinxstyleemphasis{v\_name2}, where \sphinxstyleemphasis{name2} is the name of the variable, and is
\(df(\lambda)/d\lambda\) in the notation above (i.e., it is the analytic
derivative of \(f\) with respect to \(\lambda\)).
Note that the \sphinxstyleemphasis{name1} variable is also typically given as an
argument to the {\hyperref[\detokenize{fix_adapt::doc}]{\sphinxcrossref{\DUrole{doc}{fix adapt}}}} command.

\sphinxAtStartPar
An alchemical simulation may use several pair potentials together,
invoked via the {\hyperref[\detokenize{pair_hybrid::doc}]{\sphinxcrossref{\DUrole{doc}{pair\_style hybrid or hybrid/overlay}}}}
command.  The total \(dU_s/d\lambda\) for the overall system is calculated
as the sum of each contributing term as listed by the keywords in the
{\hyperref[\detokenize{compute_ti::doc}]{\sphinxcrossref{\DUrole{doc}{compute ti}}}} command.  Individual pair potentials can be
listed, which will be sub\sphinxhyphen{}styles in the hybrid case.  You can also include a
\(k\)\sphinxhyphen{}space term via the \sphinxstyleemphasis{kspace} keyword.  You can also include a pairwise
long\sphinxhyphen{}range tail correction to the energy via the \sphinxstyleemphasis{tail} keyword.

\sphinxAtStartPar
For each term, you can specify a different (or the same) scale factor
by the two variables that you list.  Again, these will typically
correspond toe the scale factors applied to these various potentials
and the \(k\)\sphinxhyphen{}space contribution via the {\hyperref[\detokenize{fix_adapt::doc}]{\sphinxcrossref{\DUrole{doc}{fix adapt}}}}
command.

\sphinxAtStartPar
More details about the exact functional forms for the computation of
\(du/dl\) can be found in the paper by {\hyperref[\detokenize{compute_ti:eike}]{\sphinxcrossref{\DUrole{std,std-ref}{Eike}}}}.


\bigskip\hrule\bigskip



\subsection{Output info}
\label{\detokenize{compute_ti:output-info}}
\sphinxAtStartPar
This compute calculates a global scalar, namely \(dU_s/d\lambda\).  This
value can be used by any command that uses a global scalar value from
a compute as input.  See the {\hyperref[\detokenize{Howto_output::doc}]{\sphinxcrossref{\DUrole{doc}{Howto output}}}} doc page
for an overview of LAMMPS output options.

\sphinxAtStartPar
The scalar value calculated by this compute is “extensive”.

\sphinxAtStartPar
The scalar value will be in energy {\hyperref[\detokenize{units::doc}]{\sphinxcrossref{\DUrole{doc}{units}}}}.


\subsection{Restrictions}
\label{\detokenize{compute_ti:restrictions}}
\sphinxAtStartPar
This compute is part of the EXTRA\sphinxhyphen{}COMPUTE package.  It is only enabled if
LAMMPS was built with that package.  See the
{\hyperref[\detokenize{Build_package::doc}]{\sphinxcrossref{\DUrole{doc}{Build package}}}} page for more info.


\subsection{Related commands}
\label{\detokenize{compute_ti:related-commands}}
\sphinxAtStartPar
{\hyperref[\detokenize{fix_adapt::doc}]{\sphinxcrossref{\DUrole{doc}{fix adapt}}}}


\subsection{Default}
\label{\detokenize{compute_ti:default}}
\sphinxAtStartPar
none


\bigskip\hrule\bigskip

\phantomsection\label{\detokenize{compute_ti:eike}}
\sphinxAtStartPar
\sphinxstylestrong{(Eike)} Eike and Maginn, Journal of Chemical Physics, 124, 164503 (2006).

\sphinxstepscope

\index{compute torque/chunk@\spxentry{compute torque/chunk}}\ignorespaces 

\section{compute torque/chunk command}
\label{\detokenize{compute_torque_chunk:compute-torque-chunk-command}}\label{\detokenize{compute_torque_chunk:index-0}}\label{\detokenize{compute_torque_chunk::doc}}

\subsection{Syntax}
\label{\detokenize{compute_torque_chunk:syntax}}
\begin{sphinxVerbatim}[commandchars=\\\{\}]
\PYG{k}{compute }\PYG{n+nv+nvIdentifier}{ID}\PYG{+w}{ }\PYG{n+nv+nvIdentifier}{group\PYGZhy{}ID}\PYG{+w}{ }\PYG{n}{torque}\PYG{o}{/}\PYG{n}{chunk}\PYG{+w}{ }\PYG{n}{chunkID}
\end{sphinxVerbatim}
\begin{itemize}
\item {} 
\sphinxAtStartPar
ID, group\sphinxhyphen{}ID are documented in {\hyperref[\detokenize{compute::doc}]{\sphinxcrossref{\DUrole{doc}{compute}}}} command

\item {} 
\sphinxAtStartPar
torque/chunk = style name of this compute command

\item {} 
\sphinxAtStartPar
chunkID = ID of {\hyperref[\detokenize{compute_chunk_atom::doc}]{\sphinxcrossref{\DUrole{doc}{compute chunk/atom}}}} command

\end{itemize}


\subsection{Examples}
\label{\detokenize{compute_torque_chunk:examples}}
\begin{sphinxVerbatim}[commandchars=\\\{\}]
\PYG{k}{compute }\PYG{n+nv+nvIdentifier}{1}\PYG{+w}{ }\PYG{n+nv+nvIdentifier}{fluid}\PYG{+w}{ }\PYG{n}{torque}\PYG{o}{/}\PYG{n}{chunk}\PYG{+w}{ }\PYG{n}{molchunk}
\end{sphinxVerbatim}


\subsection{Description}
\label{\detokenize{compute_torque_chunk:description}}
\sphinxAtStartPar
Define a computation that calculates the torque on multiple chunks of
atoms.

\sphinxAtStartPar
In LAMMPS, chunks are collections of atoms defined by a
{\hyperref[\detokenize{compute_chunk_atom::doc}]{\sphinxcrossref{\DUrole{doc}{compute chunk/atom}}}} command, which assigns each atom
to a single chunk (or no chunk).  The ID for this command is specified
as chunkID.  For example, a single chunk could be the atoms in a
molecule or atoms in a spatial bin.  See the
{\hyperref[\detokenize{compute_chunk_atom::doc}]{\sphinxcrossref{\DUrole{doc}{compute chunk/atom}}}} and
{\hyperref[\detokenize{Howto_chunk::doc}]{\sphinxcrossref{\DUrole{doc}{Howto chunk}}}}
doc pages for details of how chunks can be defined and examples of how
they can be used to measure properties of a system.

\sphinxAtStartPar
This compute calculates the three components of the torque vector for eqch
chunk, due to the forces on the individual atoms in the chunk around
the center\sphinxhyphen{}of\sphinxhyphen{}mass of the chunk.  The calculation includes all effects
due to atoms passing through periodic boundaries.

\sphinxAtStartPar
Note that only atoms in the specified group contribute to the
calculation.  The {\hyperref[\detokenize{compute_chunk_atom::doc}]{\sphinxcrossref{\DUrole{doc}{compute chunk/atom}}}} command
defines its own group; atoms will have a chunk ID = 0 if they are not
in that group, signifying they are not assigned to a chunk, and will
thus also not contribute to this calculation.  You can specify the
“all” group for this command if you simply want to include atoms with
non\sphinxhyphen{}zero chunk IDs.

\begin{sphinxadmonition}{note}{Note:}
\sphinxAtStartPar
The coordinates of an atom contribute to the chunk’s torque in
“unwrapped” form, by using the image flags associated with each atom.
See the {\hyperref[\detokenize{dump::doc}]{\sphinxcrossref{\DUrole{doc}{dump custom}}}} command for a discussion of
“unwrapped” coordinates.  See the Atoms section of the
{\hyperref[\detokenize{read_data::doc}]{\sphinxcrossref{\DUrole{doc}{read\_data}}}} command for a discussion of image flags and
how they are set for each atom.  You can reset the image flags
(e.g., to 0) before invoking this compute by using the
{\hyperref[\detokenize{set::doc}]{\sphinxcrossref{\DUrole{doc}{set image}}}} command.
\end{sphinxadmonition}

\sphinxAtStartPar
The simplest way to output the results of the compute torque/chunk
calculation to a file is to use the {\hyperref[\detokenize{fix_ave_time::doc}]{\sphinxcrossref{\DUrole{doc}{fix ave/time}}}}
command, for example:

\begin{sphinxVerbatim}[commandchars=\\\{\}]
\PYG{k}{compute }\PYG{n+nv+nvIdentifier}{cc1}\PYG{+w}{ }\PYG{n+nv+nvIdentifier}{all}\PYG{+w}{ }\PYG{n}{chunk}\PYG{o}{/}\PYG{n}{atom}\PYG{+w}{ }\PYG{n}{molecule}
\PYG{k}{compute }\PYG{n+nv+nvIdentifier}{myChunk}\PYG{+w}{ }\PYG{n+nv+nvIdentifier}{all}\PYG{+w}{ }\PYG{n}{torque}\PYG{o}{/}\PYG{n}{chunk}\PYG{+w}{ }\PYG{n}{cc1}
\PYG{k}{fix }\PYG{n+nv+nvIdentifier}{1}\PYG{+w}{ }\PYG{n+nv+nvIdentifier}{all}\PYG{+w}{ }\PYG{n}{ave}\PYG{o}{/}\PYG{n}{time}\PYG{+w}{ }\PYG{l+m}{100}\PYG{+w}{ }\PYG{l+m}{1}\PYG{+w}{ }\PYG{l+m}{100}\PYG{+w}{ }\PYG{n}{c\PYGZus{}myChunk[}\PYG{o}{*}\PYG{n}{]}\PYG{+w}{ }\PYG{n}{file}\PYG{+w}{ }\PYG{n}{tmp.out}\PYG{+w}{ }\PYG{n}{mode}\PYG{+w}{ }\PYG{n}{vector}
\end{sphinxVerbatim}


\subsection{Output info}
\label{\detokenize{compute_torque_chunk:output-info}}
\sphinxAtStartPar
This compute calculates a global array where the number of rows is equal to the
number of chunks \sphinxstyleemphasis{Nchunk} as calculated by the specified
{\hyperref[\detokenize{compute_chunk_atom::doc}]{\sphinxcrossref{\DUrole{doc}{compute chunk/atom}}}} command.  The number of columns
is three for the \(x\), \(y\), and \(z\) components of the torque
for each chunk.  These values can be accessed by any command that uses global
array values from a compute as input.
See the {\hyperref[\detokenize{Howto_output::doc}]{\sphinxcrossref{\DUrole{doc}{Howto output}}}} doc page
for an overview of LAMMPS output options.

\sphinxAtStartPar
The array values are “intensive”.  The array values will be in
force\sphinxhyphen{}distance {\hyperref[\detokenize{units::doc}]{\sphinxcrossref{\DUrole{doc}{units}}}}.


\subsection{Restrictions}
\label{\detokenize{compute_torque_chunk:restrictions}}\begin{quote}

\sphinxAtStartPar
none
\end{quote}


\subsection{Related commands}
\label{\detokenize{compute_torque_chunk:related-commands}}
\sphinxAtStartPar
{\hyperref[\detokenize{variable::doc}]{\sphinxcrossref{\DUrole{doc}{variable torque() function}}}}


\subsection{Default}
\label{\detokenize{compute_torque_chunk:default}}
\sphinxAtStartPar
none

\sphinxstepscope

\index{compute vacf@\spxentry{compute vacf}}\ignorespaces 

\section{compute vacf command}
\label{\detokenize{compute_vacf:compute-vacf-command}}\label{\detokenize{compute_vacf:index-0}}\label{\detokenize{compute_vacf::doc}}

\subsection{Syntax}
\label{\detokenize{compute_vacf:syntax}}
\begin{sphinxVerbatim}[commandchars=\\\{\}]
\PYG{k}{compute }\PYG{n+nv+nvIdentifier}{ID}\PYG{+w}{ }\PYG{n+nv+nvIdentifier}{group\PYGZhy{}ID}\PYG{+w}{ }\PYG{n}{vacf}
\end{sphinxVerbatim}
\begin{itemize}
\item {} 
\sphinxAtStartPar
ID, group\sphinxhyphen{}ID are documented in {\hyperref[\detokenize{compute::doc}]{\sphinxcrossref{\DUrole{doc}{compute}}}} command

\item {} 
\sphinxAtStartPar
vacf = style name of this compute command

\end{itemize}


\subsection{Examples}
\label{\detokenize{compute_vacf:examples}}
\begin{sphinxVerbatim}[commandchars=\\\{\}]
\PYG{k}{compute }\PYG{n+nv+nvIdentifier}{1}\PYG{+w}{ }\PYG{n+nv+nvIdentifier}{all}\PYG{+w}{ }\PYG{n}{vacf}
\PYG{k}{compute }\PYG{n+nv+nvIdentifier}{1}\PYG{+w}{ }\PYG{n+nv+nvIdentifier}{upper}\PYG{+w}{ }\PYG{n}{vacf}
\end{sphinxVerbatim}


\subsection{Description}
\label{\detokenize{compute_vacf:description}}
\sphinxAtStartPar
Define a computation that calculates the velocity auto\sphinxhyphen{}correlation
function (VACF), averaged over a group of atoms.  Each atom’s
contribution to the VACF is its current velocity vector dotted into
its initial velocity vector at the time the compute was specified.

\sphinxAtStartPar
A vector of four quantities is calculated by this compute.  The first three
elements of the vector are \(v_x v_{x,0}\) (and similar for the
\(y\) and \(z\) components), summed and averaged over atoms in the
group, where \(v_x\) is the current \(x\)\sphinxhyphen{}component of the velocity of
the atom and \(v_{x,0}\) is the initial \(x\)\sphinxhyphen{}component of the velocity
of the atom.  The fourth element of the vector is the total VACF
(i.e., \((v_x v_{x,0} + v_y v_{y,0} + v_z v_{z,0})\)),
summed and averaged over atoms in the group.

\sphinxAtStartPar
The integral of the VACF versus time is proportional to the diffusion
coefficient of the diffusing atoms.  This can be computed in the
following manner, using the {\hyperref[\detokenize{variable::doc}]{\sphinxcrossref{\DUrole{doc}{variable trap()}}}} function:

\begin{sphinxVerbatim}[commandchars=\\\{\}]
\PYG{k}{compute         }\PYG{n+nv+nvIdentifier}{2}\PYG{+w}{ }\PYG{n+nv+nvIdentifier}{all}\PYG{+w}{ }\PYG{n}{vacf}
\PYG{k}{fix             }\PYG{n+nv+nvIdentifier}{5}\PYG{+w}{ }\PYG{n+nv+nvIdentifier}{all}\PYG{+w}{ }\PYG{n}{vector}\PYG{+w}{ }\PYG{l+m}{1}\PYG{+w}{ }\PYG{n}{c\PYGZus{}2[4]}
\PYG{k}{variable        }\PYG{n+nv+nvIdentifier}{diff}\PYG{+w}{ }\PYG{n}{equal}\PYG{+w}{ }\PYG{n}{dt}\PYG{o}{*}\PYG{n}{trap}\PYG{n+nv}{(}\PYG{n+nv}{f\PYGZus{}5}\PYG{n+nv}{)}
\PYG{k}{thermo\PYGZus{}style}\PYG{+w}{    }\PYG{n}{custom}\PYG{+w}{ }\PYG{n}{step}\PYG{+w}{ }\PYG{n}{v\PYGZus{}diff}
\end{sphinxVerbatim}

\begin{sphinxadmonition}{note}{Note:}
\sphinxAtStartPar
If you want the quantities calculated by this compute to be
continuous when running from a {\hyperref[\detokenize{read_restart::doc}]{\sphinxcrossref{\DUrole{doc}{restart file}}}}, then
you should use the same ID for this compute, as in the original run.
This is so that the fix this compute creates to store per\sphinxhyphen{}atom
quantities will also have the same ID, and thus be initialized
correctly with time=0 atom velocities from the restart file.
\end{sphinxadmonition}


\subsection{Output info}
\label{\detokenize{compute_vacf:output-info}}
\sphinxAtStartPar
This compute calculates a global vector of length 4, which can be
accessed by indices 1\textendash{}4 by any command that uses global vector values
from a compute as input.  See the {\hyperref[\detokenize{Howto_output::doc}]{\sphinxcrossref{\DUrole{doc}{Howto output}}}} doc
page for an overview of LAMMPS output options.

\sphinxAtStartPar
The vector values are “intensive”.  The vector values will be in
velocity\(^2\) {\hyperref[\detokenize{units::doc}]{\sphinxcrossref{\DUrole{doc}{units}}}}.


\subsection{Restrictions}
\label{\detokenize{compute_vacf:restrictions}}\begin{quote}

\sphinxAtStartPar
none
\end{quote}


\subsection{Related commands}
\label{\detokenize{compute_vacf:related-commands}}
\sphinxAtStartPar
{\hyperref[\detokenize{compute_msd::doc}]{\sphinxcrossref{\DUrole{doc}{compute msd}}}}


\subsection{Default}
\label{\detokenize{compute_vacf:default}}
\sphinxAtStartPar
none

\sphinxstepscope

\index{compute vcm/chunk@\spxentry{compute vcm/chunk}}\ignorespaces 

\section{compute vcm/chunk command}
\label{\detokenize{compute_vcm_chunk:compute-vcm-chunk-command}}\label{\detokenize{compute_vcm_chunk:index-0}}\label{\detokenize{compute_vcm_chunk::doc}}

\subsection{Syntax}
\label{\detokenize{compute_vcm_chunk:syntax}}
\begin{sphinxVerbatim}[commandchars=\\\{\}]
\PYG{k}{compute }\PYG{n+nv+nvIdentifier}{ID}\PYG{+w}{ }\PYG{n+nv+nvIdentifier}{group\PYGZhy{}ID}\PYG{+w}{ }\PYG{n}{vcm}\PYG{o}{/}\PYG{n}{chunk}\PYG{+w}{ }\PYG{n}{chunkID}
\end{sphinxVerbatim}
\begin{itemize}
\item {} 
\sphinxAtStartPar
ID, group\sphinxhyphen{}ID are documented in {\hyperref[\detokenize{compute::doc}]{\sphinxcrossref{\DUrole{doc}{compute}}}} command

\item {} 
\sphinxAtStartPar
vcm/chunk = style name of this compute command

\item {} 
\sphinxAtStartPar
chunkID = ID of {\hyperref[\detokenize{compute_chunk_atom::doc}]{\sphinxcrossref{\DUrole{doc}{compute chunk/atom}}}} command

\end{itemize}


\subsection{Examples}
\label{\detokenize{compute_vcm_chunk:examples}}
\begin{sphinxVerbatim}[commandchars=\\\{\}]
\PYG{k}{compute }\PYG{n+nv+nvIdentifier}{1}\PYG{+w}{ }\PYG{n+nv+nvIdentifier}{fluid}\PYG{+w}{ }\PYG{n}{vcm}\PYG{o}{/}\PYG{n}{chunk}\PYG{+w}{ }\PYG{n}{molchunk}
\end{sphinxVerbatim}


\subsection{Description}
\label{\detokenize{compute_vcm_chunk:description}}
\sphinxAtStartPar
Define a computation that calculates the center\sphinxhyphen{}of\sphinxhyphen{}mass velocity for
multiple chunks of atoms.

\sphinxAtStartPar
In LAMMPS, chunks are collections of atoms defined by a
{\hyperref[\detokenize{compute_chunk_atom::doc}]{\sphinxcrossref{\DUrole{doc}{compute chunk/atom}}}} command, which assigns each atom
to a single chunk (or no chunk).  The ID for this command is specified as
chunkID.  For example, a single chunk could be the atoms in a molecule or atoms
in a spatial bin.  See the {\hyperref[\detokenize{compute_chunk_atom::doc}]{\sphinxcrossref{\DUrole{doc}{compute chunk/atom}}}} and
{\hyperref[\detokenize{Howto_chunk::doc}]{\sphinxcrossref{\DUrole{doc}{Howto chunk}}}} doc pages for details of how chunks can be
defined and examples of how they can be used to measure properties of a system.

\sphinxAtStartPar
This compute calculates the \((x,y,z)\) components of the center\sphinxhyphen{}of\sphinxhyphen{}mass
velocity for each chunk.  This is done by summing mass*velocity for
each atom in the chunk and dividing the sum by the total mass of the
chunk.

\sphinxAtStartPar
Note that only atoms in the specified group contribute to the
calculation.  The {\hyperref[\detokenize{compute_chunk_atom::doc}]{\sphinxcrossref{\DUrole{doc}{compute chunk/atom}}}} command
defines its own group; atoms will have a chunk ID = 0 if they are not
in that group, signifying they are not assigned to a chunk, and will
thus also not contribute to this calculation.  You can specify the
“all” group for this command if you simply want to include atoms with
non\sphinxhyphen{}zero chunk IDs.

\sphinxAtStartPar
The simplest way to output the results of the compute vcm/chunk
calculation to a file is to use the {\hyperref[\detokenize{fix_ave_time::doc}]{\sphinxcrossref{\DUrole{doc}{fix ave/time}}}}
command, for example:

\begin{sphinxVerbatim}[commandchars=\\\{\}]
\PYG{k}{compute }\PYG{n+nv+nvIdentifier}{cc1}\PYG{+w}{ }\PYG{n+nv+nvIdentifier}{all}\PYG{+w}{ }\PYG{n}{chunk}\PYG{o}{/}\PYG{n}{atom}\PYG{+w}{ }\PYG{n}{molecule}
\PYG{k}{compute }\PYG{n+nv+nvIdentifier}{myChunk}\PYG{+w}{ }\PYG{n+nv+nvIdentifier}{all}\PYG{+w}{ }\PYG{n}{vcm}\PYG{o}{/}\PYG{n}{chunk}\PYG{+w}{ }\PYG{n}{cc1}
\PYG{k}{fix }\PYG{n+nv+nvIdentifier}{1}\PYG{+w}{ }\PYG{n+nv+nvIdentifier}{all}\PYG{+w}{ }\PYG{n}{ave}\PYG{o}{/}\PYG{n}{time}\PYG{+w}{ }\PYG{l+m}{100}\PYG{+w}{ }\PYG{l+m}{1}\PYG{+w}{ }\PYG{l+m}{100}\PYG{+w}{ }\PYG{n}{c\PYGZus{}myChunk[}\PYG{o}{*}\PYG{n}{]}\PYG{+w}{ }\PYG{n}{file}\PYG{+w}{ }\PYG{n}{tmp.out}\PYG{+w}{ }\PYG{n}{mode}\PYG{+w}{ }\PYG{n}{vector}
\end{sphinxVerbatim}


\subsection{Output info}
\label{\detokenize{compute_vcm_chunk:output-info}}
\sphinxAtStartPar
This compute calculates a global array where the number of rows is the
number of chunks \sphinxstyleemphasis{Nchunk} as calculated by the specified
{\hyperref[\detokenize{compute_chunk_atom::doc}]{\sphinxcrossref{\DUrole{doc}{compute chunk/atom}}}} command.  The number of
columns is 3 for the \((x,y,z)\) center\sphinxhyphen{}of\sphinxhyphen{}mass velocity coordinates of
each chunk. These values can be accessed by any command that uses global array
values from a compute as input.  See the {\hyperref[\detokenize{Howto_output::doc}]{\sphinxcrossref{\DUrole{doc}{Howto output}}}}
page for an overview of LAMMPS output options.

\sphinxAtStartPar
The array values are “intensive”.  The array values will be in
velocity {\hyperref[\detokenize{units::doc}]{\sphinxcrossref{\DUrole{doc}{units}}}}.


\subsection{Restrictions}
\label{\detokenize{compute_vcm_chunk:restrictions}}
\sphinxAtStartPar
none


\subsection{Related commands}
\label{\detokenize{compute_vcm_chunk:related-commands}}
\sphinxAtStartPar
none


\subsection{Default}
\label{\detokenize{compute_vcm_chunk:default}}
\sphinxAtStartPar
none

\sphinxstepscope

\index{compute viscosity/cos@\spxentry{compute viscosity/cos}}\ignorespaces 

\section{compute viscosity/cos command}
\label{\detokenize{compute_viscosity_cos:compute-viscosity-cos-command}}\label{\detokenize{compute_viscosity_cos:index-0}}\label{\detokenize{compute_viscosity_cos::doc}}

\subsection{Syntax}
\label{\detokenize{compute_viscosity_cos:syntax}}
\begin{sphinxVerbatim}[commandchars=\\\{\}]
\PYG{k}{compute }\PYG{n+nv+nvIdentifier}{ID}\PYG{+w}{ }\PYG{n+nv+nvIdentifier}{group\PYGZhy{}ID}\PYG{+w}{ }\PYG{n}{viscosity}\PYG{o}{/}\PYG{n}{cos}
\end{sphinxVerbatim}
\begin{itemize}
\item {} 
\sphinxAtStartPar
ID, group\sphinxhyphen{}ID are documented in {\hyperref[\detokenize{compute::doc}]{\sphinxcrossref{\DUrole{doc}{compute}}}} command

\item {} 
\sphinxAtStartPar
viscosity/cos = style name of this compute command

\end{itemize}


\subsection{Examples}
\label{\detokenize{compute_viscosity_cos:examples}}
\begin{sphinxVerbatim}[commandchars=\\\{\}]
\PYG{k}{units}\PYG{+w}{    }\PYG{n}{real}
\PYG{k}{compute  }\PYG{n+nv+nvIdentifier}{cos}\PYG{+w}{ }\PYG{n+nv+nvIdentifier}{all}\PYG{+w}{ }\PYG{n}{viscosity}\PYG{o}{/}\PYG{n}{cos}
\PYG{k}{variable }\PYG{n+nv+nvIdentifier}{V}\PYG{+w}{ }\PYG{n}{equal}\PYG{+w}{ }\PYG{n}{c\PYGZus{}cos[7]}
\PYG{k}{variable }\PYG{n+nv+nvIdentifier}{A}\PYG{+w}{ }\PYG{n}{equal}\PYG{+w}{ }\PYG{l+m}{0.02E\PYGZhy{}5}\PYG{+w}{  }\PYG{c}{\PYGZsh{} A/fs\PYGZca{}2}
\PYG{k}{variable }\PYG{n+nv+nvIdentifier}{density}\PYG{+w}{ }\PYG{n}{equal}\PYG{+w}{ }\PYG{n}{density}
\PYG{k}{variable }\PYG{n+nv+nvIdentifier}{lz}\PYG{+w}{ }\PYG{n}{equal}\PYG{+w}{ }\PYG{n}{lz}
\PYG{k}{variable }\PYG{n+nv+nvIdentifier}{reciprocalViscosity}\PYG{+w}{ }\PYG{n}{equal}\PYG{+w}{ }\PYG{n}{v\PYGZus{}V}\PYG{o}{/}\PYG{n+nv}{\PYGZdl{}\PYGZob{}}\PYG{n+nv}{A}\PYG{n+nv}{\PYGZcb{}}\PYG{o}{/}\PYG{n}{v\PYGZus{}density}\PYG{o}{*}\PYG{l+m}{39.4784}\PYG{o}{/}\PYG{n}{v\PYGZus{}lz}\PYG{o}{/}\PYG{n}{v\PYGZus{}lz}\PYG{o}{*}\PYG{l+m}{100}\PYG{+w}{  }\PYG{c}{\PYGZsh{} 1/(Pa*s)}
\end{sphinxVerbatim}


\subsection{Description}
\label{\detokenize{compute_viscosity_cos:description}}
\sphinxAtStartPar
Define a computation that calculates the velocity amplitude of a group of atoms
with an cosine\sphinxhyphen{}shaped velocity profile and the temperature of them
after subtracting out the velocity profile before computing the kinetic energy.
A compute of this style can be used by any command that computes a temperature
(e.g., {\hyperref[\detokenize{thermo_modify::doc}]{\sphinxcrossref{\DUrole{doc}{thermo\_modify}}}}, {\hyperref[\detokenize{fix_nh::doc}]{\sphinxcrossref{\DUrole{doc}{fix npt}}}}).

\sphinxAtStartPar
This command together with {\hyperref[\detokenize{fix_accelerate_cos::doc}]{\sphinxcrossref{\DUrole{doc}{fix\_accelerate/cos}}}}
enables viscosity calculation with periodic perturbation method,
as described by {\hyperref[\detokenize{compute_viscosity_cos:hess1}]{\sphinxcrossref{\DUrole{std,std-ref}{Hess}}}}.
An acceleration along the \(x\)\sphinxhyphen{}direction is applied to the simulation
system by using {\hyperref[\detokenize{fix_accelerate_cos::doc}]{\sphinxcrossref{\DUrole{doc}{fix\_accelerate/cos}}}} command.
The acceleration is a periodic function along the \(z\)\sphinxhyphen{}direction:
\begin{equation*}
\begin{split}a_{x}(z) = A \cos \left(\frac{2 \pi z}{l_{z}}\right)\end{split}
\end{equation*}
\sphinxAtStartPar
where \(A\) is the acceleration amplitude, \(l_z\) is the
\(z\)\sphinxhyphen{}length of the simulation box. At steady state, the acceleration
generates a velocity profile:
\begin{equation*}
\begin{split}v_{x}(z) = V \cos \left(\frac{2 \pi z}{l_{z}}\right)\end{split}
\end{equation*}
\sphinxAtStartPar
The generated velocity amplitude \(V\) is related to the
shear viscosity \(\eta\) by
\begin{equation*}
\begin{split}V = \frac{A \rho}{\eta}\left(\frac{l_{z}}{2 \pi}\right)^{2},\end{split}
\end{equation*}
\sphinxAtStartPar
and it can be obtained from ensemble average of the velocity profile via
\begin{equation*}
\begin{split}V = \frac{\sum\limits_i 2 m_{i} v_{i, x} \cos \left(\frac{2 \pi z_i}{l_{z}}\right)}{\sum\limits_i m_{i}}\end{split}
\end{equation*}
\sphinxAtStartPar
where \(m_i\), \(v_{i,x}\) and \(z_i\) are the mass,
\(x\)\sphinxhyphen{}component velocity, and \(z\)\sphinxhyphen{}coordinate of a particle,
respectively.

\sphinxAtStartPar
After the cosine\sphinxhyphen{}shaped collective velocity in the \(x\)\sphinxhyphen{}direction has been
subtracted for each atom, the temperature is calculated by the formula
\begin{equation*}
\begin{split}\text{KE} = \frac{\text{dim}}{2} N k_B T,\end{split}
\end{equation*}
\sphinxAtStartPar
where KE is the total kinetic energy of the group of atoms (sum of
\(\frac12 m v^2\)), dim = 2 or 3 is the dimensionality of the simulation,
\(N\) is the number of atoms in the group, \(k_B\) is the Boltzmann
constant, and \(T\) is the absolute temperature.

\sphinxAtStartPar
A symmetric tensor, stored as a six\sphinxhyphen{}element vector, is also calculated
by this compute for use in the computation of a pressure tensor by the
{\hyperref[\detokenize{compute_pressure::doc}]{\sphinxcrossref{\DUrole{doc}{compute pressue}}}} command.  The formula for
the components of the tensor is the same as the above expression for
\(E_\mathrm{kin}\), except that the 1/2 factor is NOT included and
the \(v_i^2\) is replaced by \(v_{i,x} v_{i,y}\) for the
\(xy\) component, and so on.  Note that because it lacks the 1/2
factor, these tensor components are twice those of the traditional
kinetic energy tensor.  The six components of the vector are ordered
\(xx\), \(yy\), \(zz\), \(xy\), \(xz\),
\(yz\).

\sphinxAtStartPar
The number of atoms contributing to the temperature is assumed to be
constant for the duration of the run; use the \sphinxstyleemphasis{dynamic} option of the
{\hyperref[\detokenize{compute_modify::doc}]{\sphinxcrossref{\DUrole{doc}{compute\_modify}}}} command if this is not the case.
However, in order to get meaningful results, the group ID of this compute should
be all.

\sphinxAtStartPar
The removal of the cosine\sphinxhyphen{}shaped velocity component by this command is
essentially computing the temperature after a “bias” has been removed
from the velocity of the atoms.  If this compute is used with a fix
command that performs thermostatting then this bias will be subtracted
from each atom, thermostatting of the remaining thermal velocity will
be performed, and the bias will be added back in.  Thermostatting
fixes that work in this way include {\hyperref[\detokenize{fix_nh::doc}]{\sphinxcrossref{\DUrole{doc}{fix nvt}}}},
{\hyperref[\detokenize{fix_temp_rescale::doc}]{\sphinxcrossref{\DUrole{doc}{fix temp/rescale}}}},
{\hyperref[\detokenize{fix_temp_berendsen::doc}]{\sphinxcrossref{\DUrole{doc}{fix temp/berendsen}}}}, and
{\hyperref[\detokenize{fix_langevin::doc}]{\sphinxcrossref{\DUrole{doc}{fix langevin}}}}.

\sphinxAtStartPar
This compute subtracts out degrees of freedom due to fixes that
constrain molecular motion, such as {\hyperref[\detokenize{fix_shake::doc}]{\sphinxcrossref{\DUrole{doc}{fix shake}}}} and
{\hyperref[\detokenize{fix_rigid::doc}]{\sphinxcrossref{\DUrole{doc}{fix rigid}}}}.  This means that the temperature of groups of
atoms that include these constraints will be computed correctly.  If
needed, the subtracted degrees of freedom can be altered using the
\sphinxstyleemphasis{extra} option of the {\hyperref[\detokenize{compute_modify::doc}]{\sphinxcrossref{\DUrole{doc}{compute\_modify}}}} command.

\sphinxAtStartPar
See the {\hyperref[\detokenize{Howto_thermostat::doc}]{\sphinxcrossref{\DUrole{doc}{Howto thermostat}}}} page for a discussion of
different ways to compute temperature and perform thermostatting.


\bigskip\hrule\bigskip



\subsection{Output info}
\label{\detokenize{compute_viscosity_cos:output-info}}
\sphinxAtStartPar
This compute calculates a global scalar (the temperature) and a global
vector of length 7, which can be accessed by indices 1\textendash{}7.  The first
six elements of the vector are those of the symmetric tensor discussed
above.  The seventh is the cosine\sphinxhyphen{}shaped velocity amplitude \(V\),
which can be used to calculate the reciprocal viscosity, as shown in
the example.  These values can be used by any command that uses global
scalar or vector values from a compute as input.  See the {\hyperref[\detokenize{Howto_output::doc}]{\sphinxcrossref{\DUrole{doc}{Howto
output}}}} page for an overview of LAMMPS output options.

\sphinxAtStartPar
The scalar value calculated by this compute is “intensive”.  The
first six elements of vector values are “extensive”,
and the seventh element of vector values is “intensive”.

\sphinxAtStartPar
The scalar value is in temperature {\hyperref[\detokenize{units::doc}]{\sphinxcrossref{\DUrole{doc}{units}}}}.  The first
six elements of vector values are in energy {\hyperref[\detokenize{units::doc}]{\sphinxcrossref{\DUrole{doc}{units}}}}.  The
seventh element of vector value us in velocity {\hyperref[\detokenize{units::doc}]{\sphinxcrossref{\DUrole{doc}{units}}}}.


\subsection{Restrictions}
\label{\detokenize{compute_viscosity_cos:restrictions}}
\sphinxAtStartPar
This compute is part of the MISC package.  It is only enabled if LAMMPS
was built with that package.  See the {\hyperref[\detokenize{Build_package::doc}]{\sphinxcrossref{\DUrole{doc}{Build package}}}} page for more info.

\sphinxAtStartPar
Since this compute depends on {\hyperref[\detokenize{fix_accelerate_cos::doc}]{\sphinxcrossref{\DUrole{doc}{fix accelerate/cos}}}}
which can only work for 3d systems, it cannot be used for 2d systems.


\subsection{Related commands}
\label{\detokenize{compute_viscosity_cos:related-commands}}
\sphinxAtStartPar
{\hyperref[\detokenize{fix_accelerate_cos::doc}]{\sphinxcrossref{\DUrole{doc}{fix accelerate/cos}}}}


\subsection{Default}
\label{\detokenize{compute_viscosity_cos:default}}\begin{quote}

\sphinxAtStartPar
none
\end{quote}


\bigskip\hrule\bigskip

\phantomsection\label{\detokenize{compute_viscosity_cos:hess1}}
\sphinxAtStartPar
\sphinxstylestrong{(Hess)} Hess, B. The Journal of Chemical Physics 2002, 116 (1), 209\sphinxhyphen{}217.

\sphinxstepscope

\index{compute voronoi/atom@\spxentry{compute voronoi/atom}}\ignorespaces 

\section{compute voronoi/atom command}
\label{\detokenize{compute_voronoi_atom:compute-voronoi-atom-command}}\label{\detokenize{compute_voronoi_atom:index-0}}\label{\detokenize{compute_voronoi_atom::doc}}

\subsection{Syntax}
\label{\detokenize{compute_voronoi_atom:syntax}}
\begin{sphinxVerbatim}[commandchars=\\\{\}]
\PYG{k}{compute }\PYG{n+nv+nvIdentifier}{ID}\PYG{+w}{ }\PYG{n+nv+nvIdentifier}{group\PYGZhy{}ID}\PYG{+w}{ }\PYG{n}{voronoi}\PYG{o}{/}\PYG{n}{atom}\PYG{+w}{ }\PYG{n}{keyword}\PYG{+w}{ }\PYG{n}{arg}\PYG{+w}{ }\PYG{n}{...}
\end{sphinxVerbatim}
\begin{itemize}
\item {} 
\sphinxAtStartPar
ID, group\sphinxhyphen{}ID are documented in {\hyperref[\detokenize{compute::doc}]{\sphinxcrossref{\DUrole{doc}{compute}}}} command

\item {} 
\sphinxAtStartPar
voronoi/atom = style name of this compute command

\item {} 
\sphinxAtStartPar
zero or more keyword/value pairs may be appended

\item {} 
\sphinxAtStartPar
keyword = \sphinxstyleemphasis{only\_group} or \sphinxstyleemphasis{occupation} or \sphinxstyleemphasis{surface} or \sphinxstyleemphasis{radius} or \sphinxstyleemphasis{edge\_histo} or \sphinxstyleemphasis{edge\_threshold} or \sphinxstyleemphasis{face\_threshold} or \sphinxstyleemphasis{neighbors}
\begin{sphinxalltt}
\sphinxstyleemphasis{only\_group} = no arg
\sphinxstyleemphasis{occupation} = no arg
\sphinxstyleemphasis{surface} arg = sgroup\sphinxhyphen{}ID
  sgroup\sphinxhyphen{}ID = compute the dividing surface between group\sphinxhyphen{}ID and sgroup\sphinxhyphen{}ID
    this keyword adds a third column to the compute output
\sphinxstyleemphasis{radius} arg = v\_r
  v\_r = radius atom style variable for a poly\sphinxhyphen{}disperse Voronoi tessellation
\sphinxstyleemphasis{edge\_histo} arg = maxedge
  maxedge = maximum number of Voronoi cell edges to be accounted in the histogram
\sphinxstyleemphasis{edge\_threshold} arg = minlength
  minlength = minimum length for an edge to be counted
\sphinxstyleemphasis{face\_threshold} arg = minarea
  minarea = minimum area for a face to be counted
\sphinxstyleemphasis{neighbors} value = \sphinxstyleemphasis{yes} or \sphinxstyleemphasis{no} = store list of all neighbors or no
\end{sphinxalltt}

\end{itemize}


\subsection{Examples}
\label{\detokenize{compute_voronoi_atom:examples}}
\begin{sphinxVerbatim}[commandchars=\\\{\}]
\PYG{k}{compute }\PYG{n+nv+nvIdentifier}{1}\PYG{+w}{ }\PYG{n+nv+nvIdentifier}{all}\PYG{+w}{ }\PYG{n}{voronoi}\PYG{o}{/}\PYG{n}{atom}
\PYG{k}{compute }\PYG{n+nv+nvIdentifier}{2}\PYG{+w}{ }\PYG{n+nv+nvIdentifier}{precipitate}\PYG{+w}{ }\PYG{n}{voronoi}\PYG{o}{/}\PYG{n}{atom}\PYG{+w}{ }\PYG{n}{surface}\PYG{+w}{ }\PYG{n}{matrix}
\PYG{k}{compute }\PYG{n+nv+nvIdentifier}{3b}\PYG{+w}{ }\PYG{n+nv+nvIdentifier}{precipitate}\PYG{+w}{ }\PYG{n}{voronoi}\PYG{o}{/}\PYG{n}{atom}\PYG{+w}{ }\PYG{n}{radius}\PYG{+w}{ }\PYG{n}{v\PYGZus{}r}
\PYG{k}{compute }\PYG{n+nv+nvIdentifier}{4}\PYG{+w}{ }\PYG{n+nv+nvIdentifier}{solute}\PYG{+w}{ }\PYG{n}{voronoi}\PYG{o}{/}\PYG{n}{atom}\PYG{+w}{ }\PYG{n}{only\PYGZus{}group}
\PYG{k}{compute }\PYG{n+nv+nvIdentifier}{5}\PYG{+w}{ }\PYG{n+nv+nvIdentifier}{defects}\PYG{+w}{ }\PYG{n}{voronoi}\PYG{o}{/}\PYG{n}{atom}\PYG{+w}{ }\PYG{n}{occupation}
\PYG{k}{compute }\PYG{n+nv+nvIdentifier}{6}\PYG{+w}{ }\PYG{n+nv+nvIdentifier}{all}\PYG{+w}{ }\PYG{n}{voronoi}\PYG{o}{/}\PYG{n}{atom}\PYG{+w}{ }\PYG{n}{neighbors}\PYG{+w}{ }\PYG{n}{yes}
\end{sphinxVerbatim}


\subsection{Description}
\label{\detokenize{compute_voronoi_atom:description}}
\sphinxAtStartPar
Define a computation that calculates the Voronoi tessellation of the
atoms in the simulation box.  The tessellation is calculated using all
atoms in the simulation, but non\sphinxhyphen{}zero values are only stored for atoms
in the group.

\sphinxAtStartPar
Two per\sphinxhyphen{}atom quantities are calculated by this compute.  The first is
the volume of the Voronoi cell around each atom.  Any point in an
atom’s Voronoi cell is closer to that atom than any other.  The second
is the number of faces of the Voronoi cell. This is equal to the
number of nearest neighbors of the central atom, plus any exterior
faces (see note below).


\bigskip\hrule\bigskip


\sphinxAtStartPar
If the \sphinxstyleemphasis{only\_group} keyword is specified the tessellation is performed
only with respect to the atoms contained in the compute group. This is
equivalent to deleting all atoms not contained in the group prior to
evaluating the tessellation.

\sphinxAtStartPar
If the \sphinxstyleemphasis{surface} keyword is specified a third quantity per atom is
computed: the Voronoi cell surface of the given atom. \sphinxstyleemphasis{surface} takes
a group ID as an argument. If a group other than \sphinxstyleemphasis{all} is specified,
only the Voronoi cell facets facing a neighbor atom from the specified
group are counted towards the surface area.

\sphinxAtStartPar
In the example above, a precipitate embedded in a matrix, only atoms
at the surface of the precipitate will have non\sphinxhyphen{}zero surface area, and
only the outward facing facets of the Voronoi cells are counted (the
hull of the precipitate). The total surface area of the precipitate
can be obtained by running a “reduce sum” compute on c\_2{[}3{]}.

\sphinxAtStartPar
If the \sphinxstyleemphasis{radius} keyword is specified with an atom style variable as
the argument, a poly\sphinxhyphen{}disperse Voronoi tessellation is
performed. Examples for radius variables are

\begin{sphinxVerbatim}[commandchars=\\\{\}]
\PYG{k}{variable }\PYG{n+nv+nvIdentifier}{r1}\PYG{+w}{ }\PYG{n}{atom}\PYG{+w}{ }\PYG{n+nv}{(}\PYG{n+nv}{type==1}\PYG{n+nv}{)}\PYG{o}{*}\PYG{l+m}{0.1}\PYG{o}{+}\PYG{n+nv}{(}\PYG{n+nv}{type==2}\PYG{n+nv}{)}\PYG{o}{*}\PYG{l+m}{0.4}
\PYG{k}{compute }\PYG{n+nv+nvIdentifier}{radius}\PYG{+w}{ }\PYG{n+nv+nvIdentifier}{all}\PYG{+w}{ }\PYG{n}{property}\PYG{o}{/}\PYG{n}{atom}\PYG{+w}{ }\PYG{n}{radius}
\PYG{k}{variable }\PYG{n+nv+nvIdentifier}{r2}\PYG{+w}{ }\PYG{n}{atom}\PYG{+w}{ }\PYG{n}{c\PYGZus{}radius}
\end{sphinxVerbatim}

\sphinxAtStartPar
Here v\_r1 specifies a per\sphinxhyphen{}type radius of 0.1 units for type 1 atoms
and 0.4 units for type 2 atoms, and v\_r2 accesses the radius property
present in atom\_style sphere for granular models.

\sphinxAtStartPar
The \sphinxstyleemphasis{edge\_histo} keyword activates the compilation of a histogram of
number of edges on the faces of the Voronoi cells in the compute
group. The argument \sphinxstyleemphasis{maxedge} of the this keyword is the largest
number of edges on a single Voronoi cell face expected to occur in the
sample. This keyword generates output of a global vector by this
compute with \sphinxstyleemphasis{maxedge}+1 entries. The last entry in the vector
contains the number of faces with more than \sphinxstyleemphasis{maxedge} edges. Since the
polygon with the smallest amount of edges is a triangle, entries 1 and
2 of the vector will always be zero.

\sphinxAtStartPar
The \sphinxstyleemphasis{edge\_threshold} and \sphinxstyleemphasis{face\_threshold} keywords allow the
suppression of edges below a given minimum length and faces below a
given minimum area. Ultra short edges and ultra small faces can occur
as artifacts of the Voronoi tessellation. These keywords will affect
the neighbor count and edge histogram outputs.

\sphinxAtStartPar
If the \sphinxstyleemphasis{occupation} keyword is specified the tessellation is only
performed for the first invocation of the compute and then stored.
For all following invocations of the compute the number of atoms in
each Voronoi cell in the stored tessellation is counted. In this mode
the compute returns a per\sphinxhyphen{}atom array with 2 columns. The first column
is the number of atoms currently in the Voronoi volume defined by this
atom at the time of the first invocation of the compute (note that the
atom may have moved significantly). The second column contains the
total number of atoms sharing the Voronoi cell of the stored
tessellation at the location of the current atom. Numbers in column
one can be any positive integer including zero, while column two
values will always be greater than zero. Column one data can be used
to locate vacancies (the coordinates are given by the atom coordinates
at the time step when the compute was first invoked), while column two
data can be used to identify interstitial atoms.

\sphinxAtStartPar
If the \sphinxstyleemphasis{neighbors} value is set to yes, then this compute also creates
a local array with 3 columns. There is one row for each face of each
Voronoi cell. The 3 columns are the atom ID of the atom that owns the
cell, the atom ID of the atom in the neighboring cell (or zero if the
face is external), and the area of the face.  The array can be
accessed by any command that uses local values from a compute as
input.  See the {\hyperref[\detokenize{Howto_output::doc}]{\sphinxcrossref{\DUrole{doc}{Howto output}}}} page for an
overview of LAMMPS output options. More specifically, the array can be
accessed by a {\hyperref[\detokenize{dump::doc}]{\sphinxcrossref{\DUrole{doc}{dump local}}}} command to write a file
containing all the Voronoi neighbors in a system:

\begin{sphinxVerbatim}[commandchars=\\\{\}]
\PYG{k}{compute }\PYG{n+nv+nvIdentifier}{6}\PYG{+w}{ }\PYG{n+nv+nvIdentifier}{all}\PYG{+w}{ }\PYG{n}{voronoi}\PYG{o}{/}\PYG{n}{atom}\PYG{+w}{ }\PYG{n}{neighbors}\PYG{+w}{ }\PYG{n}{yes}
\PYG{k}{dump }\PYG{n+nv+nvIdentifier}{d2}\PYG{+w}{ }\PYG{n+nv+nvIdentifier}{all}\PYG{+w}{ }\PYG{n}{local}\PYG{+w}{ }\PYG{l+m}{1}\PYG{+w}{ }\PYG{n}{dump.neighbors}\PYG{+w}{ }\PYG{n}{index}\PYG{+w}{ }\PYG{n}{c\PYGZus{}6[1]}\PYG{+w}{ }\PYG{n}{c\PYGZus{}6[2]}\PYG{+w}{ }\PYG{n}{c\PYGZus{}6[3]}
\end{sphinxVerbatim}

\sphinxAtStartPar
If the \sphinxstyleemphasis{face\_threshold} keyword is used, then only faces with areas
greater than the threshold are stored.


\bigskip\hrule\bigskip


\sphinxAtStartPar
The Voronoi calculation is performed by the freely available
\sphinxhref{https://math.lbl.gov/voro++/}{Voro++ package}, written by Chris Rycroft at UC Berkeley and LBL,
which must be installed on your system when building LAMMPS for use
with this compute.  See instructions on obtaining and installing the
Voro++ software in the src/VORONOI/README file.

\begin{sphinxadmonition}{note}{Note:}
\sphinxAtStartPar
The calculation of Voronoi volumes is performed by each processor
for the atoms it owns, and includes the effect of ghost atoms
stored by the processor.  This assumes that the Voronoi cells of
owned atoms are not affected by atoms beyond the ghost atom cut\sphinxhyphen{}off
distance.  This is usually a good assumption for liquid and solid
systems, but may lead to underestimation of Voronoi volumes in low
density systems.  By default, the set of ghost atoms stored by each
processor is determined by the cutoff used for {\hyperref[\detokenize{pair_style::doc}]{\sphinxcrossref{\DUrole{doc}{pair\_style}}}} interactions.  The cutoff can be set explicitly via
the {\hyperref[\detokenize{comm_modify::doc}]{\sphinxcrossref{\DUrole{doc}{comm\_modify cutoff}}}} command.  The Voronoi
cells for atoms adjacent to empty regions will extend into those
regions up to the communication cutoff in \(x\), \(y\), or
\(z\).  In that situation, an exterior face is created at the
cutoff distance normal to the \(x\), \(y\), or \(z\)
direction.  For triclinic systems, the exterior face is parallel to
the corresponding reciprocal lattice vector.
\end{sphinxadmonition}

\begin{sphinxadmonition}{note}{Note:}
\sphinxAtStartPar
The Voro++ package performs its calculation in 3d.  This will still
work for a 2d LAMMPS simulation, provided all the atoms have the
same \(z\)\sphinxhyphen{}coordinate. The Voronoi cell of each atom will be a
columnar polyhedron with constant cross\sphinxhyphen{}sectional area along the
\(z\)\sphinxhyphen{}direction and two exterior faces at the top and bottom of
the simulation box. If the atoms do not all have the same
\(z\)\sphinxhyphen{}coordinate, then the columnar cells will be accordingly
distorted. The cross\sphinxhyphen{}sectional area of each Voronoi cell can be
obtained by dividing its volume by the \(z\) extent of the
simulation box.  Note that you define the \(z\) extent of the
simulation box for 2d simulations when using the {\hyperref[\detokenize{create_box::doc}]{\sphinxcrossref{\DUrole{doc}{create\_box}}}} or {\hyperref[\detokenize{read_data::doc}]{\sphinxcrossref{\DUrole{doc}{read\_data}}}} commands.
\end{sphinxadmonition}


\subsection{Output info}
\label{\detokenize{compute_voronoi_atom:output-info}}
\sphinxAtStartPar
\DUrole{versionmodified,deprecated}{Deprecated since version 21Nov2023: }The \sphinxstyleemphasis{peratom} keyword was removed as it is no longer required.

\sphinxAtStartPar
This compute calculates a per\sphinxhyphen{}atom array with two columns. In regular
dynamic tessellation mode the first column is the Voronoi volume, the
second is the neighbor count, as described above (read above for the
output data in case the \sphinxstyleemphasis{occupation} keyword is specified).  These
values can be accessed by any command that uses per\sphinxhyphen{}atom values from a
compute as input.  See the {\hyperref[\detokenize{Howto_output::doc}]{\sphinxcrossref{\DUrole{doc}{Howto output}}}} page for
an overview of LAMMPS output options.

\sphinxAtStartPar
If the \sphinxstyleemphasis{edge\_histo} keyword is used, then this compute generates a
global vector of length \sphinxstyleemphasis{maxedge}+1, containing a histogram of the
number of edges per face.

\sphinxAtStartPar
If the \sphinxstyleemphasis{neighbors} value is set to \sphinxstyleemphasis{yes}, then this compute calculates a
local array with three columns. There is one row for each face of each
Voronoi cell.

\sphinxAtStartPar
The Voronoi cell volume will be in distance {\hyperref[\detokenize{units::doc}]{\sphinxcrossref{\DUrole{doc}{units}}}} cubed.
The Voronoi face area will be in distance {\hyperref[\detokenize{units::doc}]{\sphinxcrossref{\DUrole{doc}{units}}}} squared.


\subsection{Restrictions}
\label{\detokenize{compute_voronoi_atom:restrictions}}
\sphinxAtStartPar
This compute is part of the VORONOI package.  It is only enabled if
LAMMPS was built with that package.  See the {\hyperref[\detokenize{Build_package::doc}]{\sphinxcrossref{\DUrole{doc}{Build package}}}} page for more info.

\sphinxAtStartPar
It also requires you have a copy of the Voro++ library built and
installed on your system.  See instructions on obtaining and
installing the Voro++ software in the src/VORONOI/README file.


\subsection{Related commands}
\label{\detokenize{compute_voronoi_atom:related-commands}}
\sphinxAtStartPar
{\hyperref[\detokenize{dump::doc}]{\sphinxcrossref{\DUrole{doc}{dump custom}}}}, {\hyperref[\detokenize{dump::doc}]{\sphinxcrossref{\DUrole{doc}{dump local}}}}


\subsection{Default}
\label{\detokenize{compute_voronoi_atom:default}}
\sphinxAtStartPar
The default for the neighbors keyword is no.

\sphinxstepscope

\index{compute xrd@\spxentry{compute xrd}}\ignorespaces 

\section{compute xrd command}
\label{\detokenize{compute_xrd:compute-xrd-command}}\label{\detokenize{compute_xrd:index-0}}\label{\detokenize{compute_xrd::doc}}

\subsection{Syntax}
\label{\detokenize{compute_xrd:syntax}}
\begin{sphinxVerbatim}[commandchars=\\\{\}]
\PYG{k}{compute }\PYG{n+nv+nvIdentifier}{ID}\PYG{+w}{ }\PYG{n+nv+nvIdentifier}{group\PYGZhy{}ID}\PYG{+w}{ }\PYG{n}{xrd}\PYG{+w}{ }\PYG{n}{lambda}\PYG{+w}{ }\PYG{n}{type1}\PYG{+w}{ }\PYG{n}{type2}\PYG{+w}{ }\PYG{n}{...}\PYG{+w}{ }\PYG{n}{typeN}\PYG{+w}{ }\PYG{n}{keyword}\PYG{+w}{ }\PYG{n}{value}\PYG{+w}{ }\PYG{n}{...}
\end{sphinxVerbatim}
\begin{itemize}
\item {} 
\sphinxAtStartPar
ID, group\sphinxhyphen{}ID are documented in {\hyperref[\detokenize{compute::doc}]{\sphinxcrossref{\DUrole{doc}{compute}}}} command

\item {} 
\sphinxAtStartPar
xrd = style name of this compute command

\item {} 
\sphinxAtStartPar
lambda = wavelength of incident radiation (length units)

\item {} 
\sphinxAtStartPar
type1 type2 … typeN = chemical symbol of each atom type (see valid options below)

\item {} 
\sphinxAtStartPar
zero or more keyword/value pairs may be appended

\item {} 
\sphinxAtStartPar
keyword = \sphinxstyleemphasis{2Theta} or \sphinxstyleemphasis{c} or \sphinxstyleemphasis{LP} or \sphinxstyleemphasis{manual} or \sphinxstyleemphasis{echo}
\begin{sphinxalltt}
\sphinxstyleemphasis{2Theta} values = Min2Theta Max2Theta
  Min2Theta,Max2Theta = minimum and maximum 2 theta range to explore
  (radians or degrees)
\sphinxstyleemphasis{c} values = c1 c2 c3
  c1,c2,c3 = parameters to adjust the spacing of the reciprocal
             lattice nodes in the h, k, and l directions respectively
\sphinxstyleemphasis{LP} value = switch to apply Lorentz\sphinxhyphen{}polarization factor
  0/1 = off/on
\sphinxstyleemphasis{manual} = flag to use manual spacing of reciprocal lattice points
           based on the values of the \sphinxstyleemphasis{c} parameters
\sphinxstyleemphasis{echo} = flag to provide extra output for debugging purposes
\end{sphinxalltt}

\end{itemize}


\subsection{Examples}
\label{\detokenize{compute_xrd:examples}}
\begin{sphinxVerbatim}[commandchars=\\\{\}]
\PYG{k}{compute }\PYG{n+nv+nvIdentifier}{1}\PYG{+w}{ }\PYG{n+nv+nvIdentifier}{all}\PYG{+w}{ }\PYG{n}{xrd}\PYG{+w}{ }\PYG{l+m}{1.541838}\PYG{+w}{ }\PYG{n}{Al}\PYG{+w}{ }\PYG{n}{O}\PYG{+w}{ }\PYG{l+m}{2}\PYG{n}{Theta}\PYG{+w}{ }\PYG{l+m}{0.087}\PYG{+w}{ }\PYG{l+m}{0.87}\PYG{+w}{ }\PYG{n}{c}\PYG{+w}{ }\PYG{l+m}{1}\PYG{+w}{ }\PYG{l+m}{1}\PYG{+w}{ }\PYG{l+m}{1}\PYG{+w}{ }\PYG{n}{LP}\PYG{+w}{ }\PYG{l+m}{1}\PYG{+w}{ }\PYG{n}{echo}
\PYG{k}{compute }\PYG{n+nv+nvIdentifier}{2}\PYG{+w}{ }\PYG{n+nv+nvIdentifier}{all}\PYG{+w}{ }\PYG{n}{xrd}\PYG{+w}{ }\PYG{l+m}{1.541838}\PYG{+w}{ }\PYG{n}{Al}\PYG{+w}{ }\PYG{n}{O}\PYG{+w}{ }\PYG{l+m}{2}\PYG{n}{Theta}\PYG{+w}{ }\PYG{l+m}{10}\PYG{+w}{ }\PYG{l+m}{100}\PYG{+w}{ }\PYG{n}{c}\PYG{+w}{ }\PYG{l+m}{0.05}\PYG{+w}{ }\PYG{l+m}{0.05}\PYG{+w}{ }\PYG{l+m}{0.05}\PYG{+w}{ }\PYG{n}{LP}\PYG{+w}{ }\PYG{l+m}{1}\PYG{+w}{ }\PYG{n}{manual}

\PYG{k}{fix }\PYG{n+nv+nvIdentifier}{1}\PYG{+w}{ }\PYG{n+nv+nvIdentifier}{all}\PYG{+w}{ }\PYG{n}{ave}\PYG{o}{/}\PYG{n}{histo}\PYG{o}{/}\PYG{n}{weight}\PYG{+w}{ }\PYG{l+m}{1}\PYG{+w}{ }\PYG{l+m}{1}\PYG{+w}{ }\PYG{l+m}{1}\PYG{+w}{ }\PYG{l+m}{0.087}\PYG{+w}{ }\PYG{l+m}{0.87}\PYG{+w}{ }\PYG{l+m}{250}\PYG{+w}{ }\PYG{n}{c\PYGZus{}1[1]}\PYG{+w}{ }\PYG{n}{c\PYGZus{}1[2]}\PYG{+w}{ }\PYG{n}{mode}\PYG{+w}{ }\PYG{n}{vector}\PYG{+w}{ }\PYG{n}{file}\PYG{+w}{ }\PYG{n}{Rad2Theta.xrd}
\PYG{k}{fix }\PYG{n+nv+nvIdentifier}{2}\PYG{+w}{ }\PYG{n+nv+nvIdentifier}{all}\PYG{+w}{ }\PYG{n}{ave}\PYG{o}{/}\PYG{n}{histo}\PYG{o}{/}\PYG{n}{weight}\PYG{+w}{ }\PYG{l+m}{1}\PYG{+w}{ }\PYG{l+m}{1}\PYG{+w}{ }\PYG{l+m}{1}\PYG{+w}{ }\PYG{l+m}{10}\PYG{+w}{ }\PYG{l+m}{100}\PYG{+w}{ }\PYG{l+m}{250}\PYG{+w}{ }\PYG{n}{c\PYGZus{}2[1]}\PYG{+w}{ }\PYG{n}{c\PYGZus{}2[2]}\PYG{+w}{ }\PYG{n}{mode}\PYG{+w}{ }\PYG{n}{vector}\PYG{+w}{ }\PYG{n}{file}\PYG{+w}{ }\PYG{n}{Deg2Theta.xrd}
\end{sphinxVerbatim}


\subsection{Description}
\label{\detokenize{compute_xrd:description}}
\sphinxAtStartPar
Define a computation that calculates X\sphinxhyphen{}ray diffraction intensity as described
in {\hyperref[\detokenize{compute_xrd:xrd-coleman}]{\sphinxcrossref{\DUrole{std,std-ref}{(Coleman)}}}} on a mesh of reciprocal lattice nodes defined
by the entire simulation domain (or manually) using a simulated radiation
of wavelength \sphinxstyleemphasis{lambda}.

\sphinxAtStartPar
The X\sphinxhyphen{}ray diffraction intensity, \(I\), at each reciprocal lattice point,
\(k\), is computed from the structure factor, \(F\), using the
equations:
\begin{equation*}
\begin{split}I &= L_p(\theta)\frac{F^{*}F}{N} \\
F(\mathbf{k}) &= \sum_{j=1}^{N}f_j(\theta)exp(2\pi i \mathbf{k}\cdot \mathbf{r}_j) \\
L_p(\theta) &= \frac{1+\cos^2(2\theta)}{\cos(\theta)\sin^2(\theta)} \\
\frac{\sin(\theta)}{\lambda} &= \frac{\left\lVert\mathbf{k}\right\rVert}{2}\end{split}
\end{equation*}
\sphinxAtStartPar
Here, \(\mathbf{k}\) is the location of the reciprocal lattice node,
\(r_j\) is the position of each atom, \(f_j\) are atomic
scattering factors, \sphinxstyleemphasis{Lp} is the Lorentz\sphinxhyphen{}polarization factor, and
\(\theta\) is the scattering angle of diffraction.  The
Lorentz\sphinxhyphen{}polarization factor can be turned off using the optional \sphinxstyleemphasis{LP}
keyword.

\sphinxAtStartPar
Diffraction intensities are calculated on a three\sphinxhyphen{}dimensional mesh of
reciprocal lattice nodes. The mesh spacing is defined either (a) by the
entire simulation domain or (b) manually using selected values as shown
in the 2D diagram below.

\noindent{\hspace*{\fill}\sphinxincludegraphics[scale=0.75]{{xrd_mesh}.png}\hspace*{\fill}}

\sphinxAtStartPar
For a mesh defined by the simulation domain, a rectilinear grid is
constructed with spacing \(c A^{-1}\) along each reciprocal lattice
axis, where \(A\) is a matrix containing the vectors corresponding
to the edges of the simulation cell. If one or two directions has
non\sphinxhyphen{}periodic boundary conditions, then the spacing in these directions
is defined from the average of the (inversed) box lengths with periodic
boundary conditions.  Meshes defined by the simulation domain must
contain at least one periodic boundary.

\sphinxAtStartPar
If the \sphinxstyleemphasis{manual} flag is included, the mesh of reciprocal lattice nodes
will be defined using the \sphinxstyleemphasis{c} values for the spacing along each
reciprocal lattice axis. Note that manual mapping of the reciprocal
space mesh is good for comparing diffraction results from multiple
simulations; however, it can reduce the likelihood that Bragg
reflections will be satisfied unless small spacing parameters
(\(< 0.05~\AA^{-1}\)) are implemented.
Meshes with manual spacing do not require a periodic boundary.

\sphinxAtStartPar
The limits of the reciprocal lattice mesh are determined by range of
scattering angles explored.  The \sphinxstyleemphasis{2Theta} parameter allows the user
to reduce the scattering angle range to only the region of interest
which reduces the cost of the computation.

\sphinxAtStartPar
The atomic scattering factor, \(f_j\), accounts for the reduction in
diffraction intensity due to Compton scattering.  Compute xrd uses
analytical approximations of the atomic scattering factors that vary
for each atom type (type1 type2 … typeN) and angle of diffraction.
The analytic approximation is computed using the formula
{\hyperref[\detokenize{compute_xrd:colliex}]{\sphinxcrossref{\DUrole{std,std-ref}{(Colliex)}}}}:
\begin{equation*}
\begin{split}f_j\left ( \frac{\sin(\theta)}{\lambda} \right )=\sum_{i=1}^{4}
a_i \exp\left ( -b_i \frac{\sin^{2}(\theta)}{\lambda^{2}} \right )+c\end{split}
\end{equation*}
\sphinxAtStartPar
Coefficients parameterized by {\hyperref[\detokenize{compute_xrd:peng}]{\sphinxcrossref{\DUrole{std,std-ref}{(Peng)}}}} are assigned for each
atom type designating the chemical symbol and charge of each atom
type. Valid chemical symbols for compute xrd are:


\begin{savenotes}
\sphinxatlongtablestart
\sphinxthistablewithglobalstyle
\makeatletter
  \LTleft \@totalleftmargin plus1fill
  \LTright\dimexpr\columnwidth-\@totalleftmargin-\linewidth\relax plus1fill
\makeatother
\begin{longtable}{lllll}
\sphinxtoprule
\endfirsthead

\multicolumn{5}{c}{\sphinxnorowcolor
    \makebox[0pt]{\sphinxtablecontinued{\tablename\ \thetable{} \textendash{} continued from previous page}}%
}\\
\sphinxtoprule
\endhead

\sphinxbottomrule
\multicolumn{5}{r}{\sphinxnorowcolor
    \makebox[0pt][r]{\sphinxtablecontinued{continues on next page}}%
}\\
\endfoot

\endlastfoot
\sphinxtableatstartofbodyhook

\sphinxAtStartPar
H
&
\sphinxAtStartPar
He1\sphinxhyphen{}
&
\sphinxAtStartPar
He
&
\sphinxAtStartPar
Li
&
\sphinxAtStartPar
Li1+
\\
\sphinxhline
\sphinxAtStartPar
Be
&
\sphinxAtStartPar
Be2+
&
\sphinxAtStartPar
B
&
\sphinxAtStartPar
C
&
\sphinxAtStartPar
Cval
\\
\sphinxhline
\sphinxAtStartPar
N
&
\sphinxAtStartPar
O
&
\sphinxAtStartPar
O1\sphinxhyphen{}
&
\sphinxAtStartPar
F
&
\sphinxAtStartPar
F1\sphinxhyphen{}
\\
\sphinxhline
\sphinxAtStartPar
Ne
&
\sphinxAtStartPar
Na
&
\sphinxAtStartPar
Na1+
&
\sphinxAtStartPar
Mg
&
\sphinxAtStartPar
Mg2+
\\
\sphinxhline
\sphinxAtStartPar
Al
&
\sphinxAtStartPar
Al3+
&
\sphinxAtStartPar
Si
&
\sphinxAtStartPar
Sival
&
\sphinxAtStartPar
Si4+
\\
\sphinxhline
\sphinxAtStartPar
P
&
\sphinxAtStartPar
S
&
\sphinxAtStartPar
Cl
&
\sphinxAtStartPar
Cl1\sphinxhyphen{}
&
\sphinxAtStartPar
Ar
\\
\sphinxhline
\sphinxAtStartPar
K
&
\sphinxAtStartPar
Ca
&
\sphinxAtStartPar
Ca2+
&
\sphinxAtStartPar
Sc
&
\sphinxAtStartPar
Sc3+
\\
\sphinxhline
\sphinxAtStartPar
Ti
&
\sphinxAtStartPar
Ti2+
&
\sphinxAtStartPar
Ti3+
&
\sphinxAtStartPar
Ti4+
&
\sphinxAtStartPar
V
\\
\sphinxhline
\sphinxAtStartPar
V2+
&
\sphinxAtStartPar
V3+
&
\sphinxAtStartPar
V5+
&
\sphinxAtStartPar
Cr
&
\sphinxAtStartPar
Cr2+
\\
\sphinxhline
\sphinxAtStartPar
Cr3+
&
\sphinxAtStartPar
Mn
&
\sphinxAtStartPar
Mn2+
&
\sphinxAtStartPar
Mn3+
&
\sphinxAtStartPar
Mn4+
\\
\sphinxhline
\sphinxAtStartPar
Fe
&
\sphinxAtStartPar
Fe2+
&
\sphinxAtStartPar
Fe3+
&
\sphinxAtStartPar
Co
&
\sphinxAtStartPar
Co2+
\\
\sphinxhline
\sphinxAtStartPar
Co
&
\sphinxAtStartPar
Ni
&
\sphinxAtStartPar
Ni2+
&
\sphinxAtStartPar
Ni3+
&
\sphinxAtStartPar
Cu
\\
\sphinxhline
\sphinxAtStartPar
Cu1+
&
\sphinxAtStartPar
Cu2+
&
\sphinxAtStartPar
Zn
&
\sphinxAtStartPar
Zn2+
&
\sphinxAtStartPar
Ga
\\
\sphinxhline
\sphinxAtStartPar
Ga3+
&
\sphinxAtStartPar
Ge
&
\sphinxAtStartPar
Ge4+
&
\sphinxAtStartPar
As
&
\sphinxAtStartPar
Se
\\
\sphinxhline
\sphinxAtStartPar
Br
&
\sphinxAtStartPar
Br1\sphinxhyphen{}
&
\sphinxAtStartPar
Kr
&
\sphinxAtStartPar
Rb
&
\sphinxAtStartPar
Rb1+
\\
\sphinxhline
\sphinxAtStartPar
Sr
&
\sphinxAtStartPar
Sr2+
&
\sphinxAtStartPar
Y
&
\sphinxAtStartPar
Y3+
&
\sphinxAtStartPar
Zr
\\
\sphinxhline
\sphinxAtStartPar
Zr4+
&
\sphinxAtStartPar
Nb
&
\sphinxAtStartPar
Nb3+
&
\sphinxAtStartPar
Nb5+
&
\sphinxAtStartPar
Mo
\\
\sphinxhline
\sphinxAtStartPar
Mo3+
&
\sphinxAtStartPar
Mo5+
&
\sphinxAtStartPar
Mo6+
&
\sphinxAtStartPar
Tc
&
\sphinxAtStartPar
Ru
\\
\sphinxhline
\sphinxAtStartPar
Ru3+
&
\sphinxAtStartPar
Ru4+
&
\sphinxAtStartPar
Rh
&
\sphinxAtStartPar
Rh3+
&
\sphinxAtStartPar
Rh4+
\\
\sphinxhline
\sphinxAtStartPar
Pd
&
\sphinxAtStartPar
Pd2+
&
\sphinxAtStartPar
Pd4+
&
\sphinxAtStartPar
Ag
&
\sphinxAtStartPar
Ag1+
\\
\sphinxhline
\sphinxAtStartPar
Ag2+
&
\sphinxAtStartPar
Cd
&
\sphinxAtStartPar
Cd2+
&
\sphinxAtStartPar
In
&
\sphinxAtStartPar
In3+
\\
\sphinxhline
\sphinxAtStartPar
Sn
&
\sphinxAtStartPar
Sn2+
&
\sphinxAtStartPar
Sn4+
&
\sphinxAtStartPar
Sb
&
\sphinxAtStartPar
Sb3+
\\
\sphinxhline
\sphinxAtStartPar
Sb5+
&
\sphinxAtStartPar
Te
&
\sphinxAtStartPar
I
&
\sphinxAtStartPar
I1\sphinxhyphen{}
&
\sphinxAtStartPar
Xe
\\
\sphinxhline
\sphinxAtStartPar
Cs
&
\sphinxAtStartPar
Cs1+
&
\sphinxAtStartPar
Ba
&
\sphinxAtStartPar
Ba2+
&
\sphinxAtStartPar
La
\\
\sphinxhline
\sphinxAtStartPar
La3+
&
\sphinxAtStartPar
Ce
&
\sphinxAtStartPar
Ce3+
&
\sphinxAtStartPar
Ce4+
&
\sphinxAtStartPar
Pr
\\
\sphinxhline
\sphinxAtStartPar
Pr3+
&
\sphinxAtStartPar
Pr4+
&
\sphinxAtStartPar
Nd
&
\sphinxAtStartPar
Nd3+
&
\sphinxAtStartPar
Pm
\\
\sphinxhline
\sphinxAtStartPar
Pm3+
&
\sphinxAtStartPar
Sm
&
\sphinxAtStartPar
Sm3+
&
\sphinxAtStartPar
Eu
&
\sphinxAtStartPar
Eu2+
\\
\sphinxhline
\sphinxAtStartPar
Eu3+
&
\sphinxAtStartPar
Gd
&
\sphinxAtStartPar
Gd3+
&
\sphinxAtStartPar
Tb
&
\sphinxAtStartPar
Tb3+
\\
\sphinxhline
\sphinxAtStartPar
Dy
&
\sphinxAtStartPar
Dy3+
&
\sphinxAtStartPar
Ho
&
\sphinxAtStartPar
Ho3+
&
\sphinxAtStartPar
Er
\\
\sphinxhline
\sphinxAtStartPar
Er3+
&
\sphinxAtStartPar
Tm
&
\sphinxAtStartPar
Tm3+
&
\sphinxAtStartPar
Yb
&
\sphinxAtStartPar
Yb2+
\\
\sphinxhline
\sphinxAtStartPar
Yb3+
&
\sphinxAtStartPar
Lu
&
\sphinxAtStartPar
Lu3+
&
\sphinxAtStartPar
Hf
&
\sphinxAtStartPar
Hf4+
\\
\sphinxhline
\sphinxAtStartPar
Ta
&
\sphinxAtStartPar
Ta5+
&
\sphinxAtStartPar
W
&
\sphinxAtStartPar
W6+
&
\sphinxAtStartPar
Re
\\
\sphinxhline
\sphinxAtStartPar
Os
&
\sphinxAtStartPar
Os4+
&
\sphinxAtStartPar
Ir
&
\sphinxAtStartPar
Ir3+
&
\sphinxAtStartPar
Ir4+
\\
\sphinxhline
\sphinxAtStartPar
Pt
&
\sphinxAtStartPar
Pt2+
&
\sphinxAtStartPar
Pt4+
&
\sphinxAtStartPar
Au
&
\sphinxAtStartPar
Au1+
\\
\sphinxhline
\sphinxAtStartPar
Au3+
&
\sphinxAtStartPar
Hg
&
\sphinxAtStartPar
Hg1+
&
\sphinxAtStartPar
Hg2+
&
\sphinxAtStartPar
Tl
\\
\sphinxhline
\sphinxAtStartPar
Tl1+
&
\sphinxAtStartPar
Tl3+
&
\sphinxAtStartPar
Pb
&
\sphinxAtStartPar
Pb2+
&
\sphinxAtStartPar
Pb4+
\\
\sphinxhline
\sphinxAtStartPar
Bi
&
\sphinxAtStartPar
Bi3+
&
\sphinxAtStartPar
Bi5+
&
\sphinxAtStartPar
Po
&
\sphinxAtStartPar
At
\\
\sphinxhline
\sphinxAtStartPar
Rn
&
\sphinxAtStartPar
Fr
&
\sphinxAtStartPar
Ra
&
\sphinxAtStartPar
Ra2+
&
\sphinxAtStartPar
Ac
\\
\sphinxhline
\sphinxAtStartPar
Ac3+
&
\sphinxAtStartPar
Th
&
\sphinxAtStartPar
Th4+
&
\sphinxAtStartPar
Pa
&
\sphinxAtStartPar
U
\\
\sphinxhline
\sphinxAtStartPar
U3+
&
\sphinxAtStartPar
U4+
&
\sphinxAtStartPar
U6+
&
\sphinxAtStartPar
Np
&
\sphinxAtStartPar
Np3+
\\
\sphinxhline
\sphinxAtStartPar
Np4+
&
\sphinxAtStartPar
Np6+
&
\sphinxAtStartPar
Pu
&
\sphinxAtStartPar
Pu3+
&
\sphinxAtStartPar
Pu4+
\\
\sphinxhline
\sphinxAtStartPar
Pu6+
&
\sphinxAtStartPar
Am
&
\sphinxAtStartPar
Cm
&
\sphinxAtStartPar
Bk
&
\sphinxAtStartPar
Cf
\\
\sphinxbottomrule
\end{longtable}
\sphinxtableafterendhook
\sphinxatlongtableend
\end{savenotes}

\sphinxAtStartPar
If the \sphinxstyleemphasis{echo} keyword is specified, compute xrd will provide extra
reporting information to the screen.


\subsection{Output info}
\label{\detokenize{compute_xrd:output-info}}
\sphinxAtStartPar
This compute calculates a global array.  The number of rows in the
array is the number of reciprocal lattice nodes that are explored
which by the mesh.  The global array has two columns.

\sphinxAtStartPar
The first column contains the diffraction angle in the units (radians
or degrees) provided with the \sphinxstyleemphasis{2Theta} values. The second column contains
the computed diffraction intensities as described above.

\sphinxAtStartPar
The array can be accessed by any command that uses global values from
a compute as input.  See the {\hyperref[\detokenize{Howto_output::doc}]{\sphinxcrossref{\DUrole{doc}{Howto output}}}} doc page
for an overview of LAMMPS output options.

\sphinxAtStartPar
All array values calculated by this compute are “intensive”.


\subsection{Restrictions}
\label{\detokenize{compute_xrd:restrictions}}
\sphinxAtStartPar
This compute is part of the DIFFRACTION package.  It is only
enabled if LAMMPS was built with that package.  See the {\hyperref[\detokenize{Build_package::doc}]{\sphinxcrossref{\DUrole{doc}{Build package}}}} page for more info.

\sphinxAtStartPar
The compute\_xrd command does not work for triclinic cells.


\subsection{Related commands}
\label{\detokenize{compute_xrd:related-commands}}
\sphinxAtStartPar
{\hyperref[\detokenize{fix_ave_histo::doc}]{\sphinxcrossref{\DUrole{doc}{fix ave/histo}}}},
{\hyperref[\detokenize{compute_saed::doc}]{\sphinxcrossref{\DUrole{doc}{compute saed}}}}


\subsection{Default}
\label{\detokenize{compute_xrd:default}}
\sphinxAtStartPar
The option defaults are \sphinxstyleemphasis{2Theta} = 1 179 (degrees), \sphinxstyleemphasis{c} = 1 1 1, \sphinxstyleemphasis{LP} = 1,
no manual flag, no echo flag.


\bigskip\hrule\bigskip

\phantomsection\label{\detokenize{compute_xrd:xrd-coleman}}
\sphinxAtStartPar
\sphinxstylestrong{(Coleman)} Coleman, Spearot, Capolungo, MSMSE, 21, 055020
(2013).

\phantomsection\label{\detokenize{compute_xrd:colliex}}
\sphinxAtStartPar
\sphinxstylestrong{(Colliex)} Colliex et al. International Tables for Crystallography
Volume C: Mathematical and Chemical Tables, 249\sphinxhyphen{}429 (2004).

\phantomsection\label{\detokenize{compute_xrd:peng}}
\sphinxAtStartPar
\sphinxstylestrong{(Peng)} Peng, Ren, Dudarev, Whelan, Acta Crystallogr. A, 52, 257\sphinxhyphen{}76
(1996).

\sphinxstepscope


\chapter{Pair Styles}
\label{\detokenize{pairs:pair-styles}}\label{\detokenize{pairs::doc}}
\sphinxstepscope

\index{pair\_style adp@\spxentry{pair\_style adp}}\ignorespaces 
\index{pair\_style adp/kk@\spxentry{pair\_style adp/kk}}\ignorespaces 
\index{pair\_style adp/omp@\spxentry{pair\_style adp/omp}}\ignorespaces 

\section{pair\_style adp command}
\label{\detokenize{pair_adp:pair-style-adp-command}}\label{\detokenize{pair_adp:index-2}}\label{\detokenize{pair_adp:index-1}}\label{\detokenize{pair_adp:index-0}}\label{\detokenize{pair_adp::doc}}
\sphinxAtStartPar
Accelerator Variants: \sphinxstyleemphasis{adp/kk}, \sphinxstyleemphasis{adp/omp}


\subsection{Syntax}
\label{\detokenize{pair_adp:syntax}}
\begin{sphinxVerbatim}[commandchars=\\\{\}]
\PYG{k}{pair\PYGZus{}style}\PYG{+w}{ }\PYG{n}{adp}
\end{sphinxVerbatim}


\subsection{Examples}
\label{\detokenize{pair_adp:examples}}
\begin{sphinxVerbatim}[commandchars=\\\{\}]
\PYG{k}{pair\PYGZus{}style}\PYG{+w}{ }\PYG{n}{adp}
\PYG{k}{pair\PYGZus{}coeff}\PYG{+w}{ }\PYG{o}{*}\PYG{+w}{ }\PYG{o}{*}\PYG{+w}{ }\PYG{n}{Ta.adp}\PYG{+w}{ }\PYG{n}{Ta}
\PYG{k}{pair\PYGZus{}coeff}\PYG{+w}{ }\PYG{o}{*}\PYG{+w}{ }\PYG{o}{*}\PYG{+w}{ }\PYG{n}{..}\PYG{o}{/}\PYG{n}{potentials}\PYG{o}{/}\PYG{n}{AlCu.adp}\PYG{+w}{ }\PYG{n}{Al}\PYG{+w}{ }\PYG{n}{Al}\PYG{+w}{ }\PYG{n}{Cu}
\end{sphinxVerbatim}


\subsection{Description}
\label{\detokenize{pair_adp:description}}
\sphinxAtStartPar
Style \sphinxstyleemphasis{adp} computes pairwise interactions for metals and metal alloys
using the angular dependent potential (ADP) of {\hyperref[\detokenize{pair_adp:mishin}]{\sphinxcrossref{\DUrole{std,std-ref}{(Mishin)}}}},
which is a generalization of the {\hyperref[\detokenize{pair_eam::doc}]{\sphinxcrossref{\DUrole{doc}{embedded atom method (EAM) potential}}}}.  The LAMMPS implementation is discussed in
{\hyperref[\detokenize{pair_adp:singh}]{\sphinxcrossref{\DUrole{std,std-ref}{(Singh)}}}}.  The total energy Ei of an atom I is given by
\begin{equation*}
\begin{split}E_i            & = F_\alpha \left( \sum_{j\neq i} \rho_\beta (r_{ij}) \right) + \frac{1}{2} \sum_{j\neq i}\phi_{\alpha\beta}(r_{ij})+ \frac{1}{2} \sum_s (\mu_i^s)^2 + \frac{1}{2} \sum_{s,t} (\lambda_i^{st})^2 - \frac{1}{6} \nu_i^2 \\
\mu_i^s        & = \sum_{j\neq i}u_{\alpha\beta}(r_{ij})r_{ij}^s\\
\lambda_i^{st} & = \sum_{j\neq i}w_{\alpha\beta}(r_{ij})r_{ij}^sr_{ij}^t\\
\nu_i          & = \sum_s\lambda_i^{ss}\end{split}
\end{equation*}
\sphinxAtStartPar
where \(F\) is the embedding energy which is a function of the atomic
electron density \(\rho\), \(\phi\) is a pair potential interaction,
\(\alpha\) and \(\beta\) are the element types of atoms \(I\) and
\(J\), and \(s\) and \(t = 1,2,3\) and refer to the cartesian
coordinates.  The \(\mu\) and \(\lambda\) terms represent the dipole
and quadruple distortions of the local atomic environment which extend the
original EAM framework by introducing angular forces.

\sphinxAtStartPar
Note that unlike for other potentials, cutoffs for ADP potentials are
not set in the pair\_style or pair\_coeff command; they are specified in
the ADP potential files themselves.  Likewise, the ADP potential files
list atomic masses; thus you do not need to use the {\hyperref[\detokenize{mass::doc}]{\sphinxcrossref{\DUrole{doc}{mass}}}}
command to specify them.

\sphinxAtStartPar
\sphinxstylestrong{ADP potentials are available from:}
\begin{itemize}
\item {} 
\sphinxAtStartPar
The NIST WWW site at \sphinxurl{https://www.ctcms.nist.gov/potentials}.
Note that ADP potentials obtained from NIST must be converted
into the extended DYNAMO \sphinxstyleemphasis{setfl} format discussed below.

\item {} 
\sphinxAtStartPar
The OpenKIM Project at
\sphinxurl{https://openkim.org/browse/models/by-type}
provides ADP potentials that can be used directly in LAMMPS with the
{\hyperref[\detokenize{kim_commands::doc}]{\sphinxcrossref{\DUrole{doc}{kim command}}}} interface.

\end{itemize}


\bigskip\hrule\bigskip


\sphinxAtStartPar
Only a single pair\_coeff command is used with the \sphinxstyleemphasis{adp} style which
specifies an extended DYNAMO \sphinxstyleemphasis{setfl} file, which contains information
for \(M\) elements.  These are mapped to LAMMPS atom types by specifying \(N\)
additional arguments after the filename in the pair\_coeff command,
where \(N\) is the number of LAMMPS atom types:
\begin{itemize}
\item {} 
\sphinxAtStartPar
filename

\item {} 
\sphinxAtStartPar
\(N\) element names = mapping of extended \sphinxstyleemphasis{setfl} elements to atom types

\end{itemize}

\sphinxAtStartPar
See the {\hyperref[\detokenize{pair_coeff::doc}]{\sphinxcrossref{\DUrole{doc}{pair\_coeff}}}} page for alternate ways to
specify the path for the potential file.

\sphinxAtStartPar
As an example, the potentials/AlCu.adp file, included in the
potentials directory of the LAMMPS distribution, is an extended \sphinxstyleemphasis{setfl}
file which has tabulated ADP values for w elements and their alloy
interactions: Cu and Al.  If your LAMMPS simulation has 4 atoms types
and you want the first 3 to be Al, and the fourth to be Cu, you would use
the following pair\_coeff command:

\begin{sphinxVerbatim}[commandchars=\\\{\}]
\PYG{k}{pair\PYGZus{}coeff}\PYG{+w}{ }\PYG{o}{*}\PYG{+w}{ }\PYG{o}{*}\PYG{+w}{ }\PYG{n}{AlCu.adp}\PYG{+w}{ }\PYG{n}{Al}\PYG{+w}{ }\PYG{n}{Al}\PYG{+w}{ }\PYG{n}{Al}\PYG{+w}{ }\PYG{n}{Cu}
\end{sphinxVerbatim}

\sphinxAtStartPar
The first 2 arguments must be * * so as to span all LAMMPS atom types.
The first three Al arguments map LAMMPS atom types 1,2,3 to the Al
element in the extended \sphinxstyleemphasis{setfl} file.  The final Cu argument maps
LAMMPS atom type 4 to the Al element in the extended \sphinxstyleemphasis{setfl} file.
Note that there is no requirement that your simulation use all the
elements specified by the extended \sphinxstyleemphasis{setfl} file.

\sphinxAtStartPar
If a mapping value is specified as NULL, the mapping is not performed.
This can be used when an \sphinxstyleemphasis{adp} potential is used as part of the
\sphinxstyleemphasis{hybrid} pair style.  The NULL values are placeholders for atom types
that will be used with other potentials.

\sphinxAtStartPar
\sphinxstyleemphasis{Adp} files in the \sphinxstyleemphasis{potentials} directory of the LAMMPS distribution
have an “.adp” suffix.  A DYNAMO \sphinxstyleemphasis{setfl} file extended for ADP is
formatted as follows.  Basically it is the standard \sphinxstyleemphasis{setfl} format
with additional tabulated functions u and w added to the file after
the tabulated pair potentials.  See the {\hyperref[\detokenize{pair_eam::doc}]{\sphinxcrossref{\DUrole{doc}{pair\_eam}}}}
command for further details on the \sphinxstyleemphasis{setfl} format.
\begin{itemize}
\item {} 
\sphinxAtStartPar
lines 1,2,3 = comments (ignored)

\item {} 
\sphinxAtStartPar
line 4: \(N_{\text{elements}}\) Element1 Element2 … ElementN

\item {} 
\sphinxAtStartPar
line 5: \(N_{\rho}\), \(d_{\rho}\), \(N_r\), \(d_r\), cutoff

\end{itemize}

\sphinxAtStartPar
Following the 5 header lines are \(N_{\text{elements}}\) sections, one for each
element, each with the following format:
\begin{itemize}
\item {} 
\sphinxAtStartPar
line 1 = atomic number, mass, lattice constant, lattice type (e.g. FCC)

\item {} 
\sphinxAtStartPar
embedding function \(F(\rho)\) (\(N_{\rho}\) values)

\item {} 
\sphinxAtStartPar
density function \(\rho(r)\) (\(N_r\) values)

\end{itemize}

\sphinxAtStartPar
Following the \(N_{\text{elements}}\) sections, \(N_r\) values for each pair potential
\(\phi(r)\) array are listed for all \(i,j\) element pairs in the same format
as other arrays.  Since these interactions are symmetric (\(i,j = j,i\))
only \(\phi\) arrays with \(i \geq j\) are listed, in the following order:
\begin{equation*}
\begin{split}i,j = (1,1), (2,1), (2,2), (3,1), (3,2), (3,3), (4,1), ..., (N_{\text{elements}},N_{\text{elements}}).\end{split}
\end{equation*}
\sphinxAtStartPar
The tabulated values for each \(\phi\) function are listed as
\(r*\phi\) (in units of eV\sphinxhyphen{}Angstroms), since they are for atom pairs, the
same as for {\hyperref[\detokenize{pair_eam::doc}]{\sphinxcrossref{\DUrole{doc}{other EAM files}}}}.

\sphinxAtStartPar
After the \(\phi(r)\) arrays, each of the \(u(r)\) arrays are listed in the
same order with the same assumptions of symmetry.  Directly following
the \(u(r)\), the \(w(r)\) arrays are listed.  Note that \(\phi(r)\) is the only
array tabulated with a scaling by \(r\).


\bigskip\hrule\bigskip


\sphinxAtStartPar
Styles with a \sphinxstyleemphasis{gpu}, \sphinxstyleemphasis{intel}, \sphinxstyleemphasis{kk}, \sphinxstyleemphasis{omp}, or \sphinxstyleemphasis{opt} suffix are
functionally the same as the corresponding style without the suffix.
They have been optimized to run faster, depending on your available
hardware, as discussed on the {\hyperref[\detokenize{Speed_packages::doc}]{\sphinxcrossref{\DUrole{doc}{Accelerator packages}}}}
page.  The accelerated styles take the same arguments and should
produce the same results, except for round\sphinxhyphen{}off and precision issues.

\sphinxAtStartPar
These accelerated styles are part of the GPU, INTEL, KOKKOS,
OPENMP, and OPT packages, respectively.  They are only enabled if
LAMMPS was built with those packages.  See the {\hyperref[\detokenize{Build_package::doc}]{\sphinxcrossref{\DUrole{doc}{Build package}}}} page for more info.

\sphinxAtStartPar
You can specify the accelerated styles explicitly in your input script
by including their suffix, or you can use the {\hyperref[\detokenize{Run_options::doc}]{\sphinxcrossref{\DUrole{doc}{\sphinxhyphen{}suffix command\sphinxhyphen{}line switch}}}} when you invoke LAMMPS, or you can use the
{\hyperref[\detokenize{suffix::doc}]{\sphinxcrossref{\DUrole{doc}{suffix}}}} command in your input script.

\sphinxAtStartPar
See the {\hyperref[\detokenize{Speed_packages::doc}]{\sphinxcrossref{\DUrole{doc}{Accelerator packages}}}} page for more
instructions on how to use the accelerated styles effectively.


\bigskip\hrule\bigskip



\subsection{Mixing, shift, table, tail correction, restart, rRESPA info}
\label{\detokenize{pair_adp:mixing-shift-table-tail-correction-restart-rrespa-info}}
\sphinxAtStartPar
For atom type pairs I,J and I != J, where types I and J correspond to
two different element types, no special mixing rules are needed, since
the ADP potential files specify alloy interactions explicitly.

\sphinxAtStartPar
This pair style does not support the {\hyperref[\detokenize{pair_modify::doc}]{\sphinxcrossref{\DUrole{doc}{pair\_modify}}}}
shift, table, and tail options.

\sphinxAtStartPar
This pair style does not write its information to {\hyperref[\detokenize{restart::doc}]{\sphinxcrossref{\DUrole{doc}{binary restart files}}}}, since it is stored in tabulated potential files.
Thus, you need to re\sphinxhyphen{}specify the pair\_style and pair\_coeff commands in
an input script that reads a restart file.

\sphinxAtStartPar
This pair style can only be used via the \sphinxstyleemphasis{pair} keyword of the
{\hyperref[\detokenize{run_style::doc}]{\sphinxcrossref{\DUrole{doc}{run\_style respa}}}} command.  It does not support the
\sphinxstyleemphasis{inner}, \sphinxstyleemphasis{middle}, \sphinxstyleemphasis{outer} keywords.


\bigskip\hrule\bigskip



\subsection{Restrictions}
\label{\detokenize{pair_adp:restrictions}}
\sphinxAtStartPar
This pair style is part of the MANYBODY package.  It is only enabled
if LAMMPS was built with that package.


\subsection{Related commands}
\label{\detokenize{pair_adp:related-commands}}
\sphinxAtStartPar
{\hyperref[\detokenize{pair_coeff::doc}]{\sphinxcrossref{\DUrole{doc}{pair\_coeff}}}}, {\hyperref[\detokenize{pair_eam::doc}]{\sphinxcrossref{\DUrole{doc}{pair\_eam}}}}


\subsection{Default}
\label{\detokenize{pair_adp:default}}
\sphinxAtStartPar
none


\bigskip\hrule\bigskip

\phantomsection\label{\detokenize{pair_adp:mishin}}
\sphinxAtStartPar
\sphinxstylestrong{(Mishin)} Mishin, Mehl, and Papaconstantopoulos, Acta Mater, 53, 4029
(2005).

\phantomsection\label{\detokenize{pair_adp:singh}}
\sphinxAtStartPar
\sphinxstylestrong{(Singh)} Singh and Warner, Acta Mater, 58, 5797\sphinxhyphen{}5805 (2010),

\sphinxstepscope

\index{pair\_style agni@\spxentry{pair\_style agni}}\ignorespaces 
\index{pair\_style agni/omp@\spxentry{pair\_style agni/omp}}\ignorespaces 

\section{pair\_style agni command}
\label{\detokenize{pair_agni:pair-style-agni-command}}\label{\detokenize{pair_agni:index-1}}\label{\detokenize{pair_agni:index-0}}\label{\detokenize{pair_agni::doc}}
\sphinxAtStartPar
Accelerator Variants: \sphinxstyleemphasis{agni/omp}


\subsection{Syntax}
\label{\detokenize{pair_agni:syntax}}
\begin{sphinxVerbatim}[commandchars=\\\{\}]
\PYG{k}{pair\PYGZus{}style}\PYG{+w}{ }\PYG{n}{agni}
\end{sphinxVerbatim}


\subsection{Examples}
\label{\detokenize{pair_agni:examples}}
\begin{sphinxVerbatim}[commandchars=\\\{\}]
\PYG{k}{pair\PYGZus{}style}\PYG{+w}{      }\PYG{n}{agni}
\PYG{k}{pair\PYGZus{}coeff}\PYG{+w}{      }\PYG{o}{*}\PYG{+w}{ }\PYG{o}{*}\PYG{+w}{ }\PYG{n}{Al.agni}\PYG{+w}{ }\PYG{n}{Al}
\end{sphinxVerbatim}


\subsection{Description}
\label{\detokenize{pair_agni:description}}
\sphinxAtStartPar
Style \sphinxstyleemphasis{agni} style computes the many\sphinxhyphen{}body vectorial force components for
an atom as
\begin{equation*}
\begin{split}F_i^u                  & = \sum_t^{N_t}\alpha_t \cdot \exp\left[-\frac{\left(d_{i,t}^u\right)^2}{2l^2}\right] \\
d_{i,t}^u              & = \left|\left| V_i^u(\eta) - V_t^u(\eta) \right|\right| \\
V_i^u(\eta)            & = \sum_{j \neq i}\frac{r^u_{ij}}{r_{ij}} \cdot e^{-\left(\frac{r_{ij}}{\eta} \right)^2} \cdot f_d\left(r_{ij}\right) \\
f_d\left(r_{ij}\right) & = \frac{1}{2} \left[\cos\left(\frac{\pi r_{ij}}{R_c}\right) + 1 \right]\end{split}
\end{equation*}
\sphinxAtStartPar
\(u\) labels the individual components, i.e. \(x\), \(y\) or \(z\), and \(V\) is the
corresponding atomic fingerprint. \(d\) is the Euclidean distance between
any two atomic fingerprints. A total of \(N_t\) reference atomic
environments are considered to construct the force field file. \(\alpha_t\)
and \(l\) are the weight coefficients and length scale parameter of the
non\sphinxhyphen{}linear regression model.

\sphinxAtStartPar
The method implements the recently proposed machine learning access to
atomic forces as discussed extensively in the following publications \sphinxhyphen{}
{\hyperref[\detokenize{pair_agni:botu2015adaptive}]{\sphinxcrossref{\DUrole{std,std-ref}{(Botu1)}}}} and {\hyperref[\detokenize{pair_agni:botu2015learning}]{\sphinxcrossref{\DUrole{std,std-ref}{(Botu2)}}}}. The premise
of the method is to map the atomic environment numerically into a
fingerprint, and use machine learning methods to create a mapping to the
vectorial atomic forces.

\sphinxAtStartPar
Only a single pair\_coeff command is used with the \sphinxstyleemphasis{agni} style which
specifies an AGNI potential file containing the parameters of the
force field for the needed elements. These are mapped to LAMMPS atom
types by specifying \(N\) additional arguments after the filename in the
pair\_coeff command, where \(N\) is the number of LAMMPS atom types:
\begin{itemize}
\item {} 
\sphinxAtStartPar
filename

\item {} 
\sphinxAtStartPar
\(N\) element names = mapping of AGNI elements to atom types

\end{itemize}

\sphinxAtStartPar
See the {\hyperref[\detokenize{pair_coeff::doc}]{\sphinxcrossref{\DUrole{doc}{pair\_coeff}}}} page for alternate ways
to specify the path for the force field file.

\sphinxAtStartPar
An AGNI force field is fully specified by the filename which contains the
parameters of the force field, i.e., the reference training environments
used to construct the machine learning force field. Example force field
and input files are provided in the examples/PACKAGES/agni directory.


\bigskip\hrule\bigskip


\sphinxAtStartPar
Styles with a \sphinxstyleemphasis{gpu}, \sphinxstyleemphasis{intel}, \sphinxstyleemphasis{kk}, \sphinxstyleemphasis{omp}, or \sphinxstyleemphasis{opt} suffix are
functionally the same as the corresponding style without the suffix.
They have been optimized to run faster, depending on your available
hardware, as discussed on the {\hyperref[\detokenize{Speed_packages::doc}]{\sphinxcrossref{\DUrole{doc}{Accelerator packages}}}}
page.  The accelerated styles take the same arguments and should
produce the same results, except for round\sphinxhyphen{}off and precision issues.

\sphinxAtStartPar
These accelerated styles are part of the GPU, INTEL, KOKKOS,
OPENMP, and OPT packages, respectively.  They are only enabled if
LAMMPS was built with those packages.  See the {\hyperref[\detokenize{Build_package::doc}]{\sphinxcrossref{\DUrole{doc}{Build package}}}} page for more info.

\sphinxAtStartPar
You can specify the accelerated styles explicitly in your input script
by including their suffix, or you can use the {\hyperref[\detokenize{Run_options::doc}]{\sphinxcrossref{\DUrole{doc}{\sphinxhyphen{}suffix command\sphinxhyphen{}line switch}}}} when you invoke LAMMPS, or you can use the
{\hyperref[\detokenize{suffix::doc}]{\sphinxcrossref{\DUrole{doc}{suffix}}}} command in your input script.

\sphinxAtStartPar
See the {\hyperref[\detokenize{Speed_packages::doc}]{\sphinxcrossref{\DUrole{doc}{Accelerator packages}}}} page for more
instructions on how to use the accelerated styles effectively.


\bigskip\hrule\bigskip



\subsection{Mixing, shift, table, tail correction, restart, rRESPA info}
\label{\detokenize{pair_agni:mixing-shift-table-tail-correction-restart-rrespa-info}}
\sphinxAtStartPar
This pair style does not support the {\hyperref[\detokenize{pair_modify::doc}]{\sphinxcrossref{\DUrole{doc}{pair\_modify}}}}
shift, table, and tail options.

\sphinxAtStartPar
This pair style does not write its information to {\hyperref[\detokenize{restart::doc}]{\sphinxcrossref{\DUrole{doc}{binary restart files}}}}, since it is stored in potential files.  Thus, you
need to re\sphinxhyphen{}specify the pair\_style and pair\_coeff commands in an input
script that reads a restart file.

\sphinxAtStartPar
This pair style can only be used via the \sphinxstyleemphasis{pair} keyword of the
{\hyperref[\detokenize{run_style::doc}]{\sphinxcrossref{\DUrole{doc}{run\_style respa}}}} command.  It does not support the
\sphinxstyleemphasis{inner}, \sphinxstyleemphasis{middle}, \sphinxstyleemphasis{outer} keywords.


\bigskip\hrule\bigskip



\subsection{Restrictions}
\label{\detokenize{pair_agni:restrictions}}
\sphinxAtStartPar
Currently, only elemental systems are implemented. Also, the method
only provides access to the forces and not energies or stresses.
The lack of potential energy data makes this pair style incompatible with
several of the {\hyperref[\detokenize{min_style::doc}]{\sphinxcrossref{\DUrole{doc}{minimizer algorthms}}}} like \sphinxstyleemphasis{cg} or \sphinxstyleemphasis{sd}.
It should work with damped dynamics based minimizers like \sphinxstyleemphasis{fire} or
\sphinxstyleemphasis{quickmin}.  However, one can access the energy via thermodynamic
integration of the forces as discussed in
{\hyperref[\detokenize{pair_agni:botu2016construct}]{\sphinxcrossref{\DUrole{std,std-ref}{(Botu3)}}}}.  This pair style is part of the
MISC package. It is only enabled if LAMMPS was built with that
package. See the {\hyperref[\detokenize{Build_package::doc}]{\sphinxcrossref{\DUrole{doc}{Build package}}}} page for more
info.

\sphinxAtStartPar
The AGNI force field files provided with LAMMPS (see the
potentials directory) are parameterized for metal {\hyperref[\detokenize{units::doc}]{\sphinxcrossref{\DUrole{doc}{units}}}}.
You can use the AGNI potential with any LAMMPS units, but you would need
to create your own AGNI potential file with coefficients listed in the
appropriate units if your simulation does not use “metal” units.


\subsection{Related commands}
\label{\detokenize{pair_agni:related-commands}}
\sphinxAtStartPar
{\hyperref[\detokenize{pair_coeff::doc}]{\sphinxcrossref{\DUrole{doc}{pair\_coeff}}}}


\subsection{Default}
\label{\detokenize{pair_agni:default}}
\sphinxAtStartPar
none


\bigskip\hrule\bigskip

\phantomsection\label{\detokenize{pair_agni:botu2015adaptive}}
\sphinxAtStartPar
\sphinxstylestrong{(Botu1)} V. Botu and R. Ramprasad, Int. J. Quant. Chem., 115(16), 1074 (2015).

\phantomsection\label{\detokenize{pair_agni:botu2015learning}}
\sphinxAtStartPar
\sphinxstylestrong{(Botu2)} V. Botu and R. Ramprasad, Phys. Rev. B, 92(9), 094306 (2015).

\phantomsection\label{\detokenize{pair_agni:botu2016construct}}
\sphinxAtStartPar
\sphinxstylestrong{(Botu3)} V. Botu, R. Batra, J. Chapman and R. Ramprasad, \sphinxurl{https://arxiv.org/abs/1610.02098} (2016).

\sphinxstepscope

\index{pair\_style aip/water/2dm@\spxentry{pair\_style aip/water/2dm}}\ignorespaces 
\index{pair\_style aip/water/2dm/opt@\spxentry{pair\_style aip/water/2dm/opt}}\ignorespaces 

\section{pair\_style aip/water/2dm command}
\label{\detokenize{pair_aip_water_2dm:pair-style-aip-water-2dm-command}}\label{\detokenize{pair_aip_water_2dm:index-1}}\label{\detokenize{pair_aip_water_2dm:index-0}}\label{\detokenize{pair_aip_water_2dm::doc}}
\sphinxAtStartPar
Accelerator Variant: \sphinxstyleemphasis{aip/water/2dm/opt}


\subsection{Syntax}
\label{\detokenize{pair_aip_water_2dm:syntax}}
\begin{sphinxVerbatim}[commandchars=\\\{\}]
\PYG{k}{pair\PYGZus{}style}\PYG{+w}{ }\PYG{n}{[hybrid}\PYG{o}{/}\PYG{n}{overlay}\PYG{+w}{ }\PYG{n}{...]}\PYG{+w}{ }\PYG{n}{aip}\PYG{o}{/}\PYG{n}{water}\PYG{o}{/}\PYG{l+m}{2}\PYG{n}{dm}\PYG{+w}{ }\PYG{n}{cutoff}\PYG{+w}{ }\PYG{n}{tap\PYGZus{}flag}
\end{sphinxVerbatim}
\begin{itemize}
\item {} 
\sphinxAtStartPar
cutoff = global cutoff (distance units)

\item {} 
\sphinxAtStartPar
tap\_flag = 0/1 to turn off/on the taper function

\end{itemize}


\subsection{Examples}
\label{\detokenize{pair_aip_water_2dm:examples}}
\begin{sphinxVerbatim}[commandchars=\\\{\}]
\PYG{k}{pair\PYGZus{}style}\PYG{+w}{  }\PYG{n}{hybrid}\PYG{o}{/}\PYG{n}{overlay}\PYG{+w}{ }\PYG{n}{aip}\PYG{o}{/}\PYG{n}{water}\PYG{o}{/}\PYG{l+m}{2}\PYG{n}{dm}\PYG{+w}{ }\PYG{l+m}{16.0}\PYG{+w}{ }\PYG{l+m}{1}
\PYG{k}{pair\PYGZus{}coeff}\PYG{+w}{  }\PYG{o}{*}\PYG{+w}{ }\PYG{o}{*}\PYG{+w}{ }\PYG{n}{aip}\PYG{o}{/}\PYG{n}{water}\PYG{o}{/}\PYG{l+m}{2}\PYG{n}{dm}\PYG{+w}{  }\PYG{n}{CBNOH.aip.water.2dm}\PYG{+w}{ }\PYG{n}{C}\PYG{+w}{ }\PYG{n}{Ow}\PYG{+w}{ }\PYG{n}{Hw}

\PYG{k}{pair\PYGZus{}style}\PYG{+w}{  }\PYG{n}{hybrid}\PYG{o}{/}\PYG{n}{overlay}\PYG{+w}{ }\PYG{n}{aip}\PYG{o}{/}\PYG{n}{water}\PYG{o}{/}\PYG{l+m}{2}\PYG{n}{dm}\PYG{+w}{ }\PYG{l+m}{16.0}\PYG{+w}{ }\PYG{n}{lj}\PYG{o}{/}\PYG{n}{cut}\PYG{o}{/}\PYG{n}{tip4p}\PYG{o}{/}\PYG{n}{long}\PYG{+w}{ }\PYG{l+m}{2}\PYG{+w}{ }\PYG{l+m}{3}\PYG{+w}{ }\PYG{l+m}{1}\PYG{+w}{ }\PYG{l+m}{1}\PYG{+w}{ }\PYG{l+m}{0.1546}\PYG{+w}{ }\PYG{l+m}{10}\PYG{+w}{ }\PYG{l+m}{8.5}
\PYG{k}{pair\PYGZus{}coeff}\PYG{+w}{  }\PYG{l+m}{2}\PYG{+w}{ }\PYG{l+m}{2}\PYG{+w}{   }\PYG{n}{lj}\PYG{o}{/}\PYG{n}{cut}\PYG{o}{/}\PYG{n}{tip4p}\PYG{o}{/}\PYG{n}{long}\PYG{+w}{    }\PYG{l+m}{8.0313e\PYGZhy{}3}\PYG{+w}{  }\PYG{l+m}{3.1589}\PYG{+w}{      }\PYG{c}{\PYGZsh{} O\PYGZhy{}O}
\PYG{k}{pair\PYGZus{}coeff}\PYG{+w}{  }\PYG{l+m}{2}\PYG{+w}{ }\PYG{l+m}{3}\PYG{+w}{   }\PYG{n}{lj}\PYG{o}{/}\PYG{n}{cut}\PYG{o}{/}\PYG{n}{tip4p}\PYG{o}{/}\PYG{n}{long}\PYG{+w}{    }\PYG{l+m}{0.0}\PYG{+w}{        }\PYG{l+m}{0.0}\PYG{+w}{         }\PYG{c}{\PYGZsh{} O\PYGZhy{}H}
\PYG{k}{pair\PYGZus{}coeff}\PYG{+w}{  }\PYG{l+m}{3}\PYG{+w}{ }\PYG{l+m}{3}\PYG{+w}{   }\PYG{n}{lj}\PYG{o}{/}\PYG{n}{cut}\PYG{o}{/}\PYG{n}{tip4p}\PYG{o}{/}\PYG{n}{long}\PYG{+w}{    }\PYG{l+m}{0.0}\PYG{+w}{        }\PYG{l+m}{0.0}\PYG{+w}{         }\PYG{c}{\PYGZsh{} H\PYGZhy{}H}
\PYG{k}{pair\PYGZus{}coeff}\PYG{+w}{  }\PYG{o}{*}\PYG{+w}{ }\PYG{o}{*}\PYG{+w}{   }\PYG{n}{aip}\PYG{o}{/}\PYG{n}{water}\PYG{o}{/}\PYG{l+m}{2}\PYG{n}{dm}\PYG{+w}{        }\PYG{n}{CBNOH.aip.water.2dm}\PYG{+w}{    }\PYG{n}{C}\PYG{+w}{ }\PYG{n}{Ow}\PYG{+w}{ }\PYG{n}{Hw}

\PYG{k}{pair\PYGZus{}style}\PYG{+w}{  }\PYG{n}{hybrid}\PYG{o}{/}\PYG{n}{overlay}\PYG{+w}{ }\PYG{n}{aip}\PYG{o}{/}\PYG{n}{water}\PYG{o}{/}\PYG{l+m}{2}\PYG{n}{dm}\PYG{+w}{ }\PYG{l+m}{16.0}\PYG{+w}{ }\PYG{n}{lj}\PYG{o}{/}\PYG{n}{cut}\PYG{o}{/}\PYG{n}{tip4p}\PYG{o}{/}\PYG{n}{long}\PYG{+w}{ }\PYG{l+m}{3}\PYG{+w}{ }\PYG{l+m}{4}\PYG{+w}{ }\PYG{l+m}{1}\PYG{+w}{ }\PYG{l+m}{1}\PYG{+w}{ }\PYG{l+m}{0.1546}\PYG{+w}{ }\PYG{l+m}{10}\PYG{+w}{ }\PYG{l+m}{8.5}\PYG{+w}{ }\PYG{n}{coul}\PYG{o}{/}\PYG{n}{shield}\PYG{+w}{ }\PYG{l+m}{16.0}\PYG{+w}{ }\PYG{l+m}{1}
\PYG{k}{pair\PYGZus{}coeff}\PYG{+w}{  }\PYG{l+m}{1}\PYG{o}{*}\PYG{l+m}{2}\PYG{+w}{ }\PYG{l+m}{1}\PYG{o}{*}\PYG{l+m}{2}\PYG{+w}{   }\PYG{n}{none}
\PYG{k}{pair\PYGZus{}coeff}\PYG{+w}{  }\PYG{l+m}{3}\PYG{+w}{ }\PYG{l+m}{3}\PYG{+w}{   }\PYG{n}{lj}\PYG{o}{/}\PYG{n}{cut}\PYG{o}{/}\PYG{n}{tip4p}\PYG{o}{/}\PYG{n}{long}\PYG{+w}{    }\PYG{l+m}{8.0313e\PYGZhy{}3}\PYG{+w}{  }\PYG{l+m}{3.1589}\PYG{+w}{      }\PYG{c}{\PYGZsh{} O\PYGZhy{}O}
\PYG{k}{pair\PYGZus{}coeff}\PYG{+w}{  }\PYG{l+m}{3}\PYG{+w}{ }\PYG{l+m}{4}\PYG{+w}{   }\PYG{n}{lj}\PYG{o}{/}\PYG{n}{cut}\PYG{o}{/}\PYG{n}{tip4p}\PYG{o}{/}\PYG{n}{long}\PYG{+w}{    }\PYG{l+m}{0.0}\PYG{+w}{        }\PYG{l+m}{0.0}\PYG{+w}{         }\PYG{c}{\PYGZsh{} O\PYGZhy{}H}
\PYG{k}{pair\PYGZus{}coeff}\PYG{+w}{  }\PYG{l+m}{4}\PYG{+w}{ }\PYG{l+m}{4}\PYG{+w}{   }\PYG{n}{lj}\PYG{o}{/}\PYG{n}{cut}\PYG{o}{/}\PYG{n}{tip4p}\PYG{o}{/}\PYG{n}{long}\PYG{+w}{    }\PYG{l+m}{0.0}\PYG{+w}{        }\PYG{l+m}{0.0}\PYG{+w}{         }\PYG{c}{\PYGZsh{} H\PYGZhy{}H}
\PYG{k}{pair\PYGZus{}coeff}\PYG{+w}{  }\PYG{o}{*}\PYG{+w}{ }\PYG{o}{*}\PYG{+w}{   }\PYG{n}{aip}\PYG{o}{/}\PYG{n}{water}\PYG{o}{/}\PYG{l+m}{2}\PYG{n}{dm}\PYG{+w}{        }\PYG{n}{CBNOH.aip.water.2dm}\PYG{+w}{  }\PYG{n}{B}\PYG{+w}{ }\PYG{n}{N}\PYG{+w}{ }\PYG{n}{Ow}\PYG{+w}{ }\PYG{n}{Hw}
\PYG{k}{pair\PYGZus{}coeff}\PYG{+w}{  }\PYG{l+m}{1}\PYG{+w}{ }\PYG{l+m}{3}\PYG{+w}{   }\PYG{n}{coul}\PYG{o}{/}\PYG{n}{shield}\PYG{+w}{          }\PYG{l+m}{1.333}
\PYG{k}{pair\PYGZus{}coeff}\PYG{+w}{  }\PYG{l+m}{1}\PYG{+w}{ }\PYG{l+m}{4}\PYG{+w}{   }\PYG{n}{coul}\PYG{o}{/}\PYG{n}{shield}\PYG{+w}{          }\PYG{l+m}{1.333}
\PYG{k}{pair\PYGZus{}coeff}\PYG{+w}{  }\PYG{l+m}{2}\PYG{+w}{ }\PYG{l+m}{3}\PYG{+w}{   }\PYG{n}{coul}\PYG{o}{/}\PYG{n}{shield}\PYG{+w}{          }\PYG{l+m}{1.333}
\PYG{k}{pair\PYGZus{}coeff}\PYG{+w}{  }\PYG{l+m}{2}\PYG{+w}{ }\PYG{l+m}{4}\PYG{+w}{   }\PYG{n}{coul}\PYG{o}{/}\PYG{n}{shield}\PYG{+w}{          }\PYG{l+m}{1.333}
\end{sphinxVerbatim}


\subsection{Description}
\label{\detokenize{pair_aip_water_2dm:description}}
\sphinxAtStartPar
\DUrole{versionmodified,added}{Added in version 15Jun2023.}

\sphinxAtStartPar
The \sphinxstyleemphasis{aip/water/2dm} style computes the anisotropic interfacial potential
(AIP) potential for interfaces of water with two\sphinxhyphen{}dimensional (2D)
materials as described in {\hyperref[\detokenize{pair_aip_water_2dm:feng1}]{\sphinxcrossref{\DUrole{std,std-ref}{(Feng1)}}}} and {\hyperref[\detokenize{pair_aip_water_2dm:feng2}]{\sphinxcrossref{\DUrole{std,std-ref}{(Feng2)}}}}.
\begin{equation*}
\begin{split}E  = & \frac{1}{2} \sum_i \sum_{j \neq i} V_{ij} \\
V_{ij}  = & {\rm Tap}(r_{ij})\left \{ e^{-\alpha (r_{ij}/\beta -1)}
             \left [ \epsilon + f(\rho_{ij}) + f(\rho_{ji})\right ] -
              \frac{1}{1+e^{-d\left [ \left ( r_{ij}/\left (s_R \cdot r^{eff} \right ) \right )-1 \right ]}}
              \cdot \frac{C_6}{r^6_{ij}} \right \}\\
\rho_{ij}^2 = & r_{ij}^2 - ({\bf r}_{ij} \cdot {\bf n}_i)^2 \\
\rho_{ji}^2  = & r_{ij}^2 - ({\bf r}_{ij} \cdot {\bf n}_j)^2 \\
f(\rho)  = &  C e^{ -( \rho / \delta )^2 } \\
{\rm Tap}(r_{ij})  = & 20\left ( \frac{r_{ij}}{R_{cut}} \right )^7 -
                        70\left ( \frac{r_{ij}}{R_{cut}} \right )^6 +
                        84\left ( \frac{r_{ij}}{R_{cut}} \right )^5 -
                        35\left ( \frac{r_{ij}}{R_{cut}} \right )^4 + 1\end{split}
\end{equation*}
\sphinxAtStartPar
Where \(\mathrm{Tap}(r_{ij})\) is the taper function which provides
a continuous cutoff (up to third derivative) for interatomic separations
larger than \(r_c\) {\hyperref[\detokenize{pair_ilp_graphene_hbn::doc}]{\sphinxcrossref{\DUrole{doc}{pair\_style ilp\_graphene\_hbn}}}}.

\begin{sphinxadmonition}{note}{Note:}
\sphinxAtStartPar
This pair style uses the atomic normal vector definition from
{\hyperref[\detokenize{pair_aip_water_2dm:feng1}]{\sphinxcrossref{\DUrole{std,std-ref}{(Feng1)}}}}), where the atomic normal vectors of the
hydrogen atoms are assumed to lie along the corresponding
oxygen\sphinxhyphen{}hydrogen bonds and the normal vector of the central oxygen
atom is defined as their average.
\end{sphinxadmonition}

\sphinxAtStartPar
The provided parameter file, \sphinxcode{\sphinxupquote{CBNOH.aip.water.2dm}}, is intended for use
with \sphinxstyleemphasis{metal} {\hyperref[\detokenize{units::doc}]{\sphinxcrossref{\DUrole{doc}{units}}}}, with energies in meV.  Two additional
parameters, \sphinxstyleemphasis{S}, and \sphinxstyleemphasis{rcut} are included in the parameter file. \sphinxstyleemphasis{S} is
designed to facilitate scaling of energies; \sphinxstyleemphasis{rcut} is the cutoff for an
internal, short distance neighbor list that is generated for speeding up
the calculation of the normals for all atom pairs.

\begin{sphinxadmonition}{note}{Note:}
\sphinxAtStartPar
The parameters presented in the provided parameter file,
\sphinxcode{\sphinxupquote{CBNOH.aip.water.2dm}}, are fitted with the taper function enabled by
setting the cutoff equal to 16.0 Angstrom.  Using a different cutoff
or taper function setting should be carefully checked as they can
lead to significant errors.  These parameters provide a good
description in both short\sphinxhyphen{} and long\sphinxhyphen{}range interaction regimes. This
is essential for simulations in high pressure regime (i.e., the
interlayer distance is smaller than the equilibrium distance).
\end{sphinxadmonition}

\sphinxAtStartPar
This potential must be used in combination with hybrid/overlay.  Other
interactions can be set to zero using {\hyperref[\detokenize{pair_coeff::doc}]{\sphinxcrossref{\DUrole{doc}{pair\_coeff settings}}}} with the pair style set to \sphinxstyleemphasis{none}.

\sphinxAtStartPar
This pair style tallies a breakdown of the total interlayer potential
energy into sub\sphinxhyphen{}categories, which can be accessed via the {\hyperref[\detokenize{compute_pair::doc}]{\sphinxcrossref{\DUrole{doc}{compute
pair}}}} command as a vector of values of length 2.  The 2
values correspond to the following sub\sphinxhyphen{}categories:
\begin{enumerate}
\sphinxsetlistlabels{\arabic}{enumi}{enumii}{}{.}%
\item {} 
\sphinxAtStartPar
\sphinxstyleemphasis{E\_vdW} = vdW (attractive) energy

\item {} 
\sphinxAtStartPar
\sphinxstyleemphasis{E\_Rep} = Repulsive energy

\end{enumerate}

\sphinxAtStartPar
To print these quantities to the log file (with descriptive column
headings) the following commands could be included in an input script:

\begin{sphinxVerbatim}[commandchars=\\\{\}]
\PYG{k}{compute }\PYG{n+nv+nvIdentifier}{0}\PYG{+w}{ }\PYG{n+nv+nvIdentifier}{all}\PYG{+w}{ }\PYG{n}{pair}\PYG{+w}{ }\PYG{n}{aip}\PYG{o}{/}\PYG{n}{water}\PYG{o}{/}\PYG{l+m}{2}\PYG{n}{dm}
\PYG{k}{variable }\PYG{n+nv+nvIdentifier}{Evdw}\PYG{+w}{  }\PYG{n}{equal}\PYG{+w}{ }\PYG{n}{c\PYGZus{}0[1]}
\PYG{k}{variable }\PYG{n+nv+nvIdentifier}{Erep}\PYG{+w}{  }\PYG{n}{equal}\PYG{+w}{ }\PYG{n}{c\PYGZus{}0[2]}
\PYG{k}{thermo\PYGZus{}style}\PYG{+w}{ }\PYG{n}{custom}\PYG{+w}{ }\PYG{n}{step}\PYG{+w}{ }\PYG{n}{temp}\PYG{+w}{ }\PYG{n}{epair}\PYG{+w}{ }\PYG{n}{v\PYGZus{}Erep}\PYG{+w}{ }\PYG{n}{v\PYGZus{}Evdw}
\end{sphinxVerbatim}


\bigskip\hrule\bigskip


\sphinxAtStartPar
Styles with a \sphinxstyleemphasis{gpu}, \sphinxstyleemphasis{intel}, \sphinxstyleemphasis{kk}, \sphinxstyleemphasis{omp}, or \sphinxstyleemphasis{opt} suffix are
functionally the same as the corresponding style without the suffix.
They have been optimized to run faster, depending on your available
hardware, as discussed on the {\hyperref[\detokenize{Speed_packages::doc}]{\sphinxcrossref{\DUrole{doc}{Accelerator packages}}}}
page.  The accelerated styles take the same arguments and should
produce the same results, except for round\sphinxhyphen{}off and precision issues.

\sphinxAtStartPar
These accelerated styles are part of the GPU, INTEL, KOKKOS,
OPENMP, and OPT packages, respectively.  They are only enabled if
LAMMPS was built with those packages.  See the {\hyperref[\detokenize{Build_package::doc}]{\sphinxcrossref{\DUrole{doc}{Build package}}}} page for more info.

\sphinxAtStartPar
You can specify the accelerated styles explicitly in your input script
by including their suffix, or you can use the {\hyperref[\detokenize{Run_options::doc}]{\sphinxcrossref{\DUrole{doc}{\sphinxhyphen{}suffix command\sphinxhyphen{}line switch}}}} when you invoke LAMMPS, or you can use the
{\hyperref[\detokenize{suffix::doc}]{\sphinxcrossref{\DUrole{doc}{suffix}}}} command in your input script.

\sphinxAtStartPar
See the {\hyperref[\detokenize{Speed_packages::doc}]{\sphinxcrossref{\DUrole{doc}{Accelerator packages}}}} page for more
instructions on how to use the accelerated styles effectively.


\bigskip\hrule\bigskip



\subsection{Mixing, shift, table, tail correction, restart, rRESPA info}
\label{\detokenize{pair_aip_water_2dm:mixing-shift-table-tail-correction-restart-rrespa-info}}
\sphinxAtStartPar
This pair style does not support the pair\_modify mix, shift, table, and
tail options.

\sphinxAtStartPar
This pair style does not write their information to binary restart
files, since it is stored in potential files. Thus, you need to
re\sphinxhyphen{}specify the pair\_style and pair\_coeff commands in an input script
that reads a restart file.


\subsection{Restrictions}
\label{\detokenize{pair_aip_water_2dm:restrictions}}
\sphinxAtStartPar
This pair style is part of the INTERLAYER package.  It is only enabled
if LAMMPS was built with that package.  See the {\hyperref[\detokenize{Build_package::doc}]{\sphinxcrossref{\DUrole{doc}{Build package}}}} page for more info.

\sphinxAtStartPar
This pair style requires the newton setting to be \sphinxstyleemphasis{on} for pair
interactions.

\sphinxAtStartPar
The \sphinxcode{\sphinxupquote{CBNOH.aip.water.2dm}} potential file provided with LAMMPS is
parameterized for \sphinxstyleemphasis{metal} units.  You can use this pair style with any
LAMMPS units, but you would need to create your own potential file with
parameters in the appropriate units, if your simulation does not use
\sphinxstyleemphasis{metal} units.


\subsection{Related commands}
\label{\detokenize{pair_aip_water_2dm:related-commands}}
\sphinxAtStartPar
{\hyperref[\detokenize{pair_coeff::doc}]{\sphinxcrossref{\DUrole{doc}{pair\_coeff}}}},
{\hyperref[\detokenize{pair_none::doc}]{\sphinxcrossref{\DUrole{doc}{pair\_none}}}},
{\hyperref[\detokenize{pair_hybrid::doc}]{\sphinxcrossref{\DUrole{doc}{pair\_style hybrid/overlay}}}},
{\hyperref[\detokenize{pair_drip::doc}]{\sphinxcrossref{\DUrole{doc}{pair\_style drip}}}},
{\hyperref[\detokenize{pair_ilp_tmd::doc}]{\sphinxcrossref{\DUrole{doc}{pair\_style ilp\_tmd}}}},
{\hyperref[\detokenize{pair_saip_metal::doc}]{\sphinxcrossref{\DUrole{doc}{pair\_style saip\_metal}}}},
{\hyperref[\detokenize{pair_ilp_graphene_hbn::doc}]{\sphinxcrossref{\DUrole{doc}{pair\_style ilp\_graphene\_hbn}}}},
{\hyperref[\detokenize{pair_kolmogorov_crespi_z::doc}]{\sphinxcrossref{\DUrole{doc}{pair\_style pair\_kolmogorov\_crespi\_z}}}},
{\hyperref[\detokenize{pair_kolmogorov_crespi_full::doc}]{\sphinxcrossref{\DUrole{doc}{pair\_style pair\_kolmogorov\_crespi\_full}}}},
{\hyperref[\detokenize{pair_lebedeva_z::doc}]{\sphinxcrossref{\DUrole{doc}{pair\_style pair\_lebedeva\_z}}}},
{\hyperref[\detokenize{pair_coul_shield::doc}]{\sphinxcrossref{\DUrole{doc}{pair\_style pair\_coul\_shield}}}}.


\subsection{Default}
\label{\detokenize{pair_aip_water_2dm:default}}
\sphinxAtStartPar
tap\_flag = 1


\bigskip\hrule\bigskip

\phantomsection\label{\detokenize{pair_aip_water_2dm:feng1}}
\sphinxAtStartPar
\sphinxstylestrong{(Feng1)} Z. Feng, …, and W. Ouyang, J. Phys. Chem. C. 127(18), 8704\sphinxhyphen{}8713 (2023).

\phantomsection\label{\detokenize{pair_aip_water_2dm:feng2}}
\sphinxAtStartPar
\sphinxstylestrong{(Feng2)} Z. Feng, …, and W. Ouyang, Langmuir 39(50), 18198\sphinxhyphen{}18207 (2023).

\sphinxstepscope

\index{pair\_style airebo@\spxentry{pair\_style airebo}}\ignorespaces 
\index{pair\_style airebo/intel@\spxentry{pair\_style airebo/intel}}\ignorespaces 
\index{pair\_style airebo/omp@\spxentry{pair\_style airebo/omp}}\ignorespaces 
\index{pair\_style airebo/morse@\spxentry{pair\_style airebo/morse}}\ignorespaces 
\index{pair\_style airebo/morse/intel@\spxentry{pair\_style airebo/morse/intel}}\ignorespaces 
\index{pair\_style airebo/morse/omp@\spxentry{pair\_style airebo/morse/omp}}\ignorespaces 
\index{pair\_style rebo@\spxentry{pair\_style rebo}}\ignorespaces 
\index{pair\_style rebo/intel@\spxentry{pair\_style rebo/intel}}\ignorespaces 
\index{pair\_style rebo/omp@\spxentry{pair\_style rebo/omp}}\ignorespaces 

\section{pair\_style airebo command}
\label{\detokenize{pair_airebo:pair-style-airebo-command}}\label{\detokenize{pair_airebo:index-8}}\label{\detokenize{pair_airebo:index-7}}\label{\detokenize{pair_airebo:index-6}}\label{\detokenize{pair_airebo:index-5}}\label{\detokenize{pair_airebo:index-4}}\label{\detokenize{pair_airebo:index-3}}\label{\detokenize{pair_airebo:index-2}}\label{\detokenize{pair_airebo:index-1}}\label{\detokenize{pair_airebo:index-0}}\label{\detokenize{pair_airebo::doc}}
\sphinxAtStartPar
Accelerator Variants: \sphinxstyleemphasis{airebo/intel}, \sphinxstyleemphasis{airebo/omp}


\section{pair\_style airebo/morse command}
\label{\detokenize{pair_airebo:pair-style-airebo-morse-command}}
\sphinxAtStartPar
Accelerator Variants: \sphinxstyleemphasis{airebo/morse/intel}, \sphinxstyleemphasis{airebo/morse/omp}


\section{pair\_style rebo command}
\label{\detokenize{pair_airebo:pair-style-rebo-command}}
\sphinxAtStartPar
Accelerator Variants: \sphinxstyleemphasis{rebo/intel}, \sphinxstyleemphasis{rebo/omp}


\subsection{Syntax}
\label{\detokenize{pair_airebo:syntax}}
\begin{sphinxVerbatim}[commandchars=\\\{\}]
\PYG{k}{pair\PYGZus{}style}\PYG{+w}{ }\PYG{n}{style}\PYG{+w}{ }\PYG{n}{cutoff}\PYG{+w}{ }\PYG{n}{LJ\PYGZus{}flag}\PYG{+w}{ }\PYG{n}{TORSION\PYGZus{}flag}\PYG{+w}{ }\PYG{n}{cutoff\PYGZus{}min}
\end{sphinxVerbatim}
\begin{itemize}
\item {} 
\sphinxAtStartPar
style = \sphinxstyleemphasis{airebo} or \sphinxstyleemphasis{airebo/morse} or \sphinxstyleemphasis{rebo}

\item {} 
\sphinxAtStartPar
cutoff = LJ or Morse cutoff (\(\sigma\) scale factor) (AIREBO and AIREBO\sphinxhyphen{}M only)

\item {} 
\sphinxAtStartPar
LJ\_flag = 0/1 to turn off/on the LJ or Morse term (AIREBO and AIREBO\sphinxhyphen{}M only, optional)

\item {} 
\sphinxAtStartPar
TORSION\_flag = 0/1 to turn off/on the torsion term (AIREBO and AIREBO\sphinxhyphen{}M only, optional)

\item {} 
\sphinxAtStartPar
cutoff\_min = Start of the transition region of cutoff (\(\sigma\) scale factor) (AIREBO and AIREBO\sphinxhyphen{}M only, optional)

\end{itemize}


\subsection{Examples}
\label{\detokenize{pair_airebo:examples}}
\begin{sphinxVerbatim}[commandchars=\\\{\}]
\PYG{k}{pair\PYGZus{}style}\PYG{+w}{ }\PYG{n}{airebo}\PYG{+w}{ }\PYG{l+m}{3.0}
\PYG{k}{pair\PYGZus{}style}\PYG{+w}{ }\PYG{n}{airebo}\PYG{+w}{ }\PYG{l+m}{2.5}\PYG{+w}{ }\PYG{l+m}{1}\PYG{+w}{ }\PYG{l+m}{0}
\PYG{k}{pair\PYGZus{}coeff}\PYG{+w}{ }\PYG{o}{*}\PYG{+w}{ }\PYG{o}{*}\PYG{+w}{ }\PYG{n}{..}\PYG{o}{/}\PYG{n}{potentials}\PYG{o}{/}\PYG{n}{CH.airebo}\PYG{+w}{ }\PYG{n}{H}\PYG{+w}{ }\PYG{n}{C}

\PYG{k}{pair\PYGZus{}style}\PYG{+w}{ }\PYG{n}{airebo}\PYG{o}{/}\PYG{n}{morse}\PYG{+w}{ }\PYG{l+m}{3.0}
\PYG{k}{pair\PYGZus{}coeff}\PYG{+w}{ }\PYG{o}{*}\PYG{+w}{ }\PYG{o}{*}\PYG{+w}{ }\PYG{n}{..}\PYG{o}{/}\PYG{n}{potentials}\PYG{o}{/}\PYG{n}{CH.airebo}\PYG{o}{\PYGZhy{}}\PYG{n}{m}\PYG{+w}{ }\PYG{n}{H}\PYG{+w}{ }\PYG{n}{C}

\PYG{k}{pair\PYGZus{}style}\PYG{+w}{ }\PYG{n}{rebo}
\PYG{k}{pair\PYGZus{}coeff}\PYG{+w}{ }\PYG{o}{*}\PYG{+w}{ }\PYG{o}{*}\PYG{+w}{ }\PYG{n}{..}\PYG{o}{/}\PYG{n}{potentials}\PYG{o}{/}\PYG{n}{CH.rebo}\PYG{+w}{ }\PYG{n}{H}\PYG{+w}{ }\PYG{n}{C}
\end{sphinxVerbatim}


\subsection{Description}
\label{\detokenize{pair_airebo:description}}
\sphinxAtStartPar
The \sphinxstyleemphasis{airebo} pair style computes the Adaptive Intermolecular Reactive
Empirical Bond Order (AIREBO) Potential of {\hyperref[\detokenize{pair_airebo:stuart}]{\sphinxcrossref{\DUrole{std,std-ref}{(Stuart)}}}} for a
system of carbon and/or hydrogen atoms.  Note that this is the initial
formulation of AIREBO from 2000, not the later formulation.

\sphinxAtStartPar
The \sphinxstyleemphasis{airebo/morse} pair style computes the AIREBO\sphinxhyphen{}M potential, which
is equivalent to AIREBO, but replaces the LJ term with a Morse potential.
The Morse potentials are parameterized by high\sphinxhyphen{}quality quantum chemistry
(MP2) calculations and do not diverge as quickly as particle density
increases. This allows AIREBO\sphinxhyphen{}M to retain accuracy to much higher pressures
than AIREBO (up to 40 GPa for Polyethylene). Details for this potential
and its parameterization are given in {\hyperref[\detokenize{pair_airebo:oconnor}]{\sphinxcrossref{\DUrole{std,std-ref}{(O’Conner)}}}}.

\sphinxAtStartPar
The \sphinxstyleemphasis{rebo} pair style computes the Reactive Empirical Bond Order (REBO)
Potential of {\hyperref[\detokenize{pair_airebo:brenner}]{\sphinxcrossref{\DUrole{std,std-ref}{(Brenner)}}}}. Note that this is the so\sphinxhyphen{}called
second generation REBO from 2002, not the original REBO from 1990.
As discussed below, second generation REBO is closely related to the
initial AIREBO; it is just a subset of the potential energy terms
with a few slightly different parameters

\sphinxAtStartPar
The AIREBO potential consists of three terms:
\begin{equation*}
\begin{split}E & = \frac{1}{2} \sum_i \sum_{j \neq i}
\left[ E^{\text{REBO}}_{ij} + E^{\text{LJ}}_{ij} +
 \sum_{k \neq i,j} \sum_{l \neq i,j,k} E^{\text{TORSION}}_{kijl} \right] \\\end{split}
\end{equation*}
\sphinxAtStartPar
By default, all three terms are included.  For the \sphinxstyleemphasis{airebo} style, if
the first two optional flag arguments to the pair\_style command are
included, the LJ and torsional terms can be turned off.  Note that
both or neither of the flags must be included.  If both of the LJ an
torsional terms are turned off, it becomes the second\sphinxhyphen{}generation REBO
potential, with a small caveat on the spline fitting procedure
mentioned below.  This can be specified directly as pair\_style \sphinxstyleemphasis{rebo}
with no additional arguments.

\sphinxAtStartPar
The detailed formulas for this potential are given in
{\hyperref[\detokenize{pair_airebo:stuart}]{\sphinxcrossref{\DUrole{std,std-ref}{(Stuart)}}}}; here we provide only a brief description.

\sphinxAtStartPar
The \(E^{\text{REBO}}\) term has the same functional form as the hydrocarbon REBO
potential developed in {\hyperref[\detokenize{pair_airebo:brenner}]{\sphinxcrossref{\DUrole{std,std-ref}{(Brenner)}}}}.  The coefficients for
\(E^{\text{REBO}}\) in AIREBO are essentially the same as Brenner’s potential, but
a few fitted spline values are slightly different.  For most cases the
\(E^{\text{REBO}}\) term in AIREBO will produce the same energies, forces and
statistical averages as the original REBO potential from which it was
derived.  The \(E^{\text{REBO}}\) term in the AIREBO potential gives the model its
reactive capabilities and only describes short\sphinxhyphen{}ranged C\sphinxhyphen{}C, C\sphinxhyphen{}H and H\sphinxhyphen{}H
interactions (\(r < 2 \AA\)). These interactions have strong
coordination\sphinxhyphen{}dependence through a bond order parameter, which adjusts
the attraction between the I,J atoms based on the position of other
nearby atoms and thus has 3\sphinxhyphen{} and 4\sphinxhyphen{}body dependence.

\sphinxAtStartPar
The \(E^{\text{LJ}}\) term adds longer\sphinxhyphen{}ranged interactions (\(2 < r < \text{cutoff}\)) using a
form similar to the standard {\hyperref[\detokenize{pair_lj::doc}]{\sphinxcrossref{\DUrole{doc}{Lennard Jones potential}}}}.
The \(E^{\text{LJ}}\) term in AIREBO contains a series of switching functions so
that the short\sphinxhyphen{}ranged LJ repulsion (\(1/r^{12}\)) does not interfere with
the energetics captured by the \(E^{\text{REBO}}\) term.  The extent of the \(E^{\text{LJ}}\)
interactions is determined by the \sphinxstyleemphasis{cutoff} argument to the pair\_style
command which is a scale factor.  For each type pair (C\sphinxhyphen{}C, C\sphinxhyphen{}H, H\sphinxhyphen{}H)
the cutoff is obtained by multiplying the scale factor by the sigma
value defined in the potential file for that type pair.  In the
standard AIREBO potential, \(\sigma_{CC} = 3.4 \AA\), so with a scale
factor of 3.0 (the argument in pair\_style), the resulting \(E^{\text{LJ}}\) cutoff
would be \(10.2 \AA\).

\sphinxAtStartPar
By default, the longer\sphinxhyphen{}ranged interaction is smoothly switched off
between 2.16 and 3.0 \(\sigma\). By specifying \sphinxstyleemphasis{cutoff\_min} in addition
to \sphinxstyleemphasis{cutoff}, the switching can be configured to take place between
\sphinxstyleemphasis{cutoff\_min} and \sphinxstyleemphasis{cutoff}. \sphinxstyleemphasis{cutoff\_min} can only be specified if all
optional arguments are given.

\sphinxAtStartPar
The \(E^{\text{TORSION}}\) term is an explicit 4\sphinxhyphen{}body potential that describes
various dihedral angle preferences in hydrocarbon configurations.


\bigskip\hrule\bigskip


\sphinxAtStartPar
Only a single pair\_coeff command is used with the \sphinxstyleemphasis{airebo}, \sphinxstyleemphasis{airebo}
or \sphinxstyleemphasis{rebo} style which specifies an AIREBO, REBO, or AIREBO\sphinxhyphen{}M potential
file with parameters for C and H.  Note that as of LAMMPS version
15 May 2019 the \sphinxstyleemphasis{rebo} style in LAMMPS uses its own potential
file (CH.rebo).  These are mapped to LAMMPS atom types by specifying
N additional arguments after the filename in the pair\_coeff command,
where N is the number of LAMMPS atom types:
\begin{itemize}
\item {} 
\sphinxAtStartPar
filename

\item {} 
\sphinxAtStartPar
\(N\) element names = mapping of AIREBO elements to atom types

\end{itemize}

\sphinxAtStartPar
See the {\hyperref[\detokenize{pair_coeff::doc}]{\sphinxcrossref{\DUrole{doc}{pair\_coeff}}}} page for alternate ways
to specify the path for the potential file.

\sphinxAtStartPar
As an example, if your LAMMPS simulation has 4 atom types and you want
the first 3 to be C, and the fourth to be H, you would use the following
pair\_coeff command:

\begin{sphinxVerbatim}[commandchars=\\\{\}]
\PYG{k}{pair\PYGZus{}coeff}\PYG{+w}{ }\PYG{o}{*}\PYG{+w}{ }\PYG{o}{*}\PYG{+w}{ }\PYG{n}{CH.airebo}\PYG{+w}{ }\PYG{n}{C}\PYG{+w}{ }\PYG{n}{C}\PYG{+w}{ }\PYG{n}{C}\PYG{+w}{ }\PYG{n}{H}
\end{sphinxVerbatim}

\sphinxAtStartPar
The first 2 arguments must be * * so as to span all LAMMPS atom types.
The first three C arguments map LAMMPS atom types 1,2,3 to the C
element in the AIREBO file.  The final H argument maps LAMMPS atom
type 4 to the H element in the AIREBO file.  If a mapping value is
specified as NULL, the mapping is not performed.  This can be used
when a \sphinxstyleemphasis{airebo} potential is used as part of the \sphinxstyleemphasis{hybrid} pair style.
The NULL values are placeholders for atom types that will be used with
other potentials.

\sphinxAtStartPar
The parameters/coefficients for the AIREBO potentials are listed in
the CH.airebo file to agree with the original {\hyperref[\detokenize{pair_airebo:stuart}]{\sphinxcrossref{\DUrole{std,std-ref}{(Stuart)}}}}
paper.  Thus the parameters are specific to this potential and the way
it was fit, so modifying the file should be done cautiously.

\sphinxAtStartPar
Similarly the parameters/coefficients for the AIREBO\sphinxhyphen{}M potentials are
listed in the CH.airebo\sphinxhyphen{}m file to agree with the {\hyperref[\detokenize{pair_airebo:oconnor}]{\sphinxcrossref{\DUrole{std,std-ref}{(O’Connor)}}}}
paper. Thus the parameters are specific to this potential and the way
it was fit, so modifying the file should be done cautiously. The
AIREBO\sphinxhyphen{}M Morse potentials were parameterized using a cutoff of
3.0 (\(\sigma\)). Modifying this cutoff may impact simulation accuracy.

\sphinxAtStartPar
This pair style tallies a breakdown of the total AIREBO potential
energy into sub\sphinxhyphen{}categories, which can be accessed via the {\hyperref[\detokenize{compute_pair::doc}]{\sphinxcrossref{\DUrole{doc}{compute pair}}}} command as a vector of values of length 3.
The 3 values correspond to the following sub\sphinxhyphen{}categories:
\begin{enumerate}
\sphinxsetlistlabels{\arabic}{enumi}{enumii}{}{.}%
\item {} 
\sphinxAtStartPar
\(E_{\text{REBO}}\) = REBO energy

\item {} 
\sphinxAtStartPar
\(E_{\text{LJ}}\) = Lennard\sphinxhyphen{}Jones energy

\item {} 
\sphinxAtStartPar
\(E_{\text{TORSION}}\) = Torsion energy

\end{enumerate}

\sphinxAtStartPar
To print these quantities to the log file (with descriptive column
headings) the following commands could be included in an input script:

\begin{sphinxVerbatim}[commandchars=\\\{\}]
\PYG{k}{compute }\PYG{n+nv+nvIdentifier}{0}\PYG{+w}{ }\PYG{n+nv+nvIdentifier}{all}\PYG{+w}{ }\PYG{n}{pair}\PYG{+w}{ }\PYG{n}{airebo}
\PYG{k}{variable }\PYG{n+nv+nvIdentifier}{REBO}\PYG{+w}{     }\PYG{n}{equal}\PYG{+w}{ }\PYG{n}{c\PYGZus{}0[1]}
\PYG{k}{variable }\PYG{n+nv+nvIdentifier}{LJ}\PYG{+w}{       }\PYG{n}{equal}\PYG{+w}{ }\PYG{n}{c\PYGZus{}0[2]}
\PYG{k}{variable }\PYG{n+nv+nvIdentifier}{TORSION}\PYG{+w}{  }\PYG{n}{equal}\PYG{+w}{ }\PYG{n}{c\PYGZus{}0[3]}
\PYG{k}{thermo\PYGZus{}style}\PYG{+w}{ }\PYG{n}{custom}\PYG{+w}{ }\PYG{n}{step}\PYG{+w}{ }\PYG{n}{temp}\PYG{+w}{ }\PYG{n}{epair}\PYG{+w}{ }\PYG{n}{v\PYGZus{}REBO}\PYG{+w}{ }\PYG{n}{v\PYGZus{}LJ}\PYG{+w}{ }\PYG{n}{v\PYGZus{}TORSION}
\end{sphinxVerbatim}


\bigskip\hrule\bigskip


\sphinxAtStartPar
Styles with a \sphinxstyleemphasis{gpu}, \sphinxstyleemphasis{intel}, \sphinxstyleemphasis{kk}, \sphinxstyleemphasis{omp}, or \sphinxstyleemphasis{opt} suffix are
functionally the same as the corresponding style without the suffix.
They have been optimized to run faster, depending on your available
hardware, as discussed on the {\hyperref[\detokenize{Speed_packages::doc}]{\sphinxcrossref{\DUrole{doc}{Accelerator packages}}}}
page.  The accelerated styles take the same arguments and should
produce the same results, except for round\sphinxhyphen{}off and precision issues.

\sphinxAtStartPar
These accelerated styles are part of the GPU, INTEL, KOKKOS,
OPENMP, and OPT packages, respectively.  They are only enabled if
LAMMPS was built with those packages.  See the {\hyperref[\detokenize{Build_package::doc}]{\sphinxcrossref{\DUrole{doc}{Build package}}}} page for more info.

\sphinxAtStartPar
You can specify the accelerated styles explicitly in your input script
by including their suffix, or you can use the {\hyperref[\detokenize{Run_options::doc}]{\sphinxcrossref{\DUrole{doc}{\sphinxhyphen{}suffix command\sphinxhyphen{}line switch}}}} when you invoke LAMMPS, or you can use the
{\hyperref[\detokenize{suffix::doc}]{\sphinxcrossref{\DUrole{doc}{suffix}}}} command in your input script.

\sphinxAtStartPar
See the {\hyperref[\detokenize{Speed_packages::doc}]{\sphinxcrossref{\DUrole{doc}{Accelerator packages}}}} page for more
instructions on how to use the accelerated styles effectively.


\bigskip\hrule\bigskip



\subsection{Mixing, shift, table, tail correction, restart, rRESPA info}
\label{\detokenize{pair_airebo:mixing-shift-table-tail-correction-restart-rrespa-info}}
\sphinxAtStartPar
These pair styles do not support the {\hyperref[\detokenize{pair_modify::doc}]{\sphinxcrossref{\DUrole{doc}{pair\_modify}}}}
mix, shift, table, and tail options.

\sphinxAtStartPar
These pair styles do not write their information to {\hyperref[\detokenize{restart::doc}]{\sphinxcrossref{\DUrole{doc}{binary restart files}}}}, since it is stored in potential files.  Thus, you
need to re\sphinxhyphen{}specify the pair\_style and pair\_coeff commands in an input
script that reads a restart file.

\sphinxAtStartPar
These pair styles can only be used via the \sphinxstyleemphasis{pair} keyword of the
{\hyperref[\detokenize{run_style::doc}]{\sphinxcrossref{\DUrole{doc}{run\_style respa}}}} command.  They do not support the
\sphinxstyleemphasis{inner}, \sphinxstyleemphasis{middle}, \sphinxstyleemphasis{outer} keywords.


\subsection{Restrictions}
\label{\detokenize{pair_airebo:restrictions}}
\sphinxAtStartPar
These pair styles are part of the MANYBODY package.  They are only
enabled if LAMMPS was built with that package.  See the {\hyperref[\detokenize{Build_package::doc}]{\sphinxcrossref{\DUrole{doc}{Build package}}}} page for more info.

\sphinxAtStartPar
These pair potentials require the {\hyperref[\detokenize{newton::doc}]{\sphinxcrossref{\DUrole{doc}{newton}}}} setting to be
“on” for pair interactions.

\sphinxAtStartPar
The CH.airebo and CH.airebo\sphinxhyphen{}m potential files provided with LAMMPS (see
the potentials directory) are parameterized for metal {\hyperref[\detokenize{units::doc}]{\sphinxcrossref{\DUrole{doc}{units}}}}.  You can use the pair styles with \sphinxstyleemphasis{any} LAMMPS units, but you
would need to create your own AIREBO or AIREBO\sphinxhyphen{}M potential file with
coefficients listed in the appropriate units, if your simulation does
not use “metal” units.

\sphinxAtStartPar
The pair styles provided here \sphinxstylestrong{only} support potential files parameterized
for the elements carbon and hydrogen (designated with “C” and “H” in the
\sphinxstyleemphasis{pair\_coeff} command.  Using potential files for other elements will trigger
an error.


\subsection{Related commands}
\label{\detokenize{pair_airebo:related-commands}}
\sphinxAtStartPar
{\hyperref[\detokenize{pair_coeff::doc}]{\sphinxcrossref{\DUrole{doc}{pair\_coeff}}}}


\subsection{Default}
\label{\detokenize{pair_airebo:default}}
\sphinxAtStartPar
none


\bigskip\hrule\bigskip

\phantomsection\label{\detokenize{pair_airebo:stuart}}
\sphinxAtStartPar
\sphinxstylestrong{(Stuart)} Stuart, Tutein, Harrison, J Chem Phys, 112, 6472\sphinxhyphen{}6486
(2000).

\phantomsection\label{\detokenize{pair_airebo:brenner}}
\sphinxAtStartPar
\sphinxstylestrong{(Brenner)} Brenner, Shenderova, Harrison, Stuart, Ni, Sinnott, J
Physics: Condensed Matter, 14, 783\sphinxhyphen{}802 (2002).

\phantomsection\label{\detokenize{pair_airebo:oconnor}}
\sphinxAtStartPar
\sphinxstylestrong{(O’Connor)} O’Connor et al., J. Chem. Phys. 142, 024903 (2015).

\sphinxstepscope

\index{pair\_style amoeba@\spxentry{pair\_style amoeba}}\ignorespaces 
\index{pair\_style amoeba/gpu@\spxentry{pair\_style amoeba/gpu}}\ignorespaces 
\index{pair\_style hippo@\spxentry{pair\_style hippo}}\ignorespaces 
\index{pair\_style hippo/gpu@\spxentry{pair\_style hippo/gpu}}\ignorespaces 

\section{pair\_style amoeba command}
\label{\detokenize{pair_amoeba:pair-style-amoeba-command}}\label{\detokenize{pair_amoeba:index-3}}\label{\detokenize{pair_amoeba:index-2}}\label{\detokenize{pair_amoeba:index-1}}\label{\detokenize{pair_amoeba:index-0}}\label{\detokenize{pair_amoeba::doc}}
\sphinxAtStartPar
Accelerator Variants: \sphinxstyleemphasis{amoeba/gpu}


\section{pair\_style hippo command}
\label{\detokenize{pair_amoeba:pair-style-hippo-command}}
\sphinxAtStartPar
Accelerator Variants: \sphinxstyleemphasis{hippo/gpu}


\subsection{Syntax}
\label{\detokenize{pair_amoeba:syntax}}
\begin{sphinxVerbatim}[commandchars=\\\{\}]
\PYG{k}{pair\PYGZus{}style}\PYG{+w}{ }\PYG{n}{style}
\end{sphinxVerbatim}
\begin{itemize}
\item {} 
\sphinxAtStartPar
style = \sphinxstyleemphasis{amoeba} or \sphinxstyleemphasis{hippo}

\end{itemize}


\subsection{Examples}
\label{\detokenize{pair_amoeba:examples}}
\begin{sphinxVerbatim}[commandchars=\\\{\}]
\PYG{k}{pair\PYGZus{}style}\PYG{+w}{ }\PYG{n}{amoeba}
\PYG{k}{pair\PYGZus{}coeff}\PYG{+w}{ }\PYG{o}{*}\PYG{+w}{ }\PYG{o}{*}\PYG{+w}{ }\PYG{n}{protein.prm.amoeba}\PYG{+w}{ }\PYG{n}{protein.key.amoeba}
\end{sphinxVerbatim}

\begin{sphinxVerbatim}[commandchars=\\\{\}]
\PYG{k}{pair\PYGZus{}style}\PYG{+w}{ }\PYG{n}{hippo}
\PYG{k}{pair\PYGZus{}coeff}\PYG{+w}{ }\PYG{o}{*}\PYG{+w}{ }\PYG{o}{*}\PYG{+w}{ }\PYG{n}{water.prm.hippo}\PYG{+w}{ }\PYG{n}{water.key.hippo}
\end{sphinxVerbatim}


\subsection{Additional info}
\label{\detokenize{pair_amoeba:additional-info}}\begin{itemize}
\item {} 
\sphinxAtStartPar
{\hyperref[\detokenize{Howto_amoeba::doc}]{\sphinxcrossref{\DUrole{doc}{Howto amoeba}}}}

\item {} 
\sphinxAtStartPar
examples/amoeba

\item {} 
\sphinxAtStartPar
tools/amoeba

\item {} 
\sphinxAtStartPar
potentials/*.amoeba

\item {} 
\sphinxAtStartPar
potentials/*.hippo

\end{itemize}


\subsection{Description}
\label{\detokenize{pair_amoeba:description}}
\sphinxAtStartPar
The \sphinxstyleemphasis{amoeba} style computes the AMOEBA polarizable field formulated
by Jay Ponder’s group at the U Washington at St Louis {\hyperref[\detokenize{pair_amoeba:amoeba-ren}]{\sphinxcrossref{\DUrole{std,std-ref}{(Ren)}}}}, {\hyperref[\detokenize{pair_amoeba:amoeba-shi}]{\sphinxcrossref{\DUrole{std,std-ref}{(Shi)}}}}.  The \sphinxstyleemphasis{hippo} style computes
the HIPPO polarizable force field, an extension to AMOEBA, formulated
by Josh Rackers and collaborators in the Ponder group {\hyperref[\detokenize{pair_amoeba:amoeba-rackers}]{\sphinxcrossref{\DUrole{std,std-ref}{(Rackers)}}}}.

\sphinxAtStartPar
These force fields can be used when polarization effects are desired
in simulations of water, organic molecules, and biomolecules including
proteins, provided that parameterizations (Tinker PRM force field
files) are available for the systems you are interested in.  Files in
the LAMMPS potentials directory with a “amoeba” or “hippo” suffix can
be used.  The Tinker distribution and website have additional force
field files as well.

\sphinxAtStartPar
As discussed on the {\hyperref[\detokenize{Howto_amoeba::doc}]{\sphinxcrossref{\DUrole{doc}{Howto amoeba}}}} doc page, the
intermolecular (non\sphinxhyphen{}bonded) portion of the AMOEBA force field contains
these terms:
\begin{equation*}
\begin{split}U_{amoeba} = U_{multipole} + U_{polar} + U_{hal}\end{split}
\end{equation*}
\sphinxAtStartPar
while the HIPPO force field contains these terms:
\begin{equation*}
\begin{split}U_{hippo} = U_{multipole} + U_{polar} + U_{qxfer} + U_{repulsion} + U_{dispersion}\end{split}
\end{equation*}
\sphinxAtStartPar
Conceptually, these terms compute the following interactions:
\begin{itemize}
\item {} 
\sphinxAtStartPar
\(U_{hal}\) = buffered 14\sphinxhyphen{}7 van der Waals with offsets applied to hydrogen atoms

\item {} 
\sphinxAtStartPar
\(U_{repulsion}\) = Pauli repulsion due to rearrangement of electron density

\item {} 
\sphinxAtStartPar
\(U_{dispersion}\) = dispersion between correlated, instantaneous induced dipole moments

\item {} 
\sphinxAtStartPar
\(U_{multipole}\) = electrostatics between permanent point charges, dipoles, and quadrupoles

\item {} 
\sphinxAtStartPar
\(U_{polar}\) = electronic polarization between induced point dipoles

\item {} 
\sphinxAtStartPar
\(U_{qxfer}\) = charge transfer effects

\end{itemize}

\sphinxAtStartPar
Note that the AMOEBA versus HIPPO force fields typically compute the
same term differently using their own formulas.  The references on
this doc page give full details for both force fields.

\sphinxAtStartPar
The formulas for the AMOEBA energy terms are:
\begin{equation*}
\begin{split}U_{hal} = & \epsilon_{ij} \left( \frac{1.07}{\rho_{ij} + 0.07} \right)^7 \left( \frac{1.12}{\rho_{ij}^7 + 0.12} - 2 \right) \\
U_{multipole} = & \vec{M}_i\boldsymbol{T_{ij}}\vec{M}_j, \quad \mbox{with} \quad
   \vec{M} = \left(q, \vec{\mu}_{perm}, \boldsymbol{\Theta} \right) \\
U_{polar} = & \frac{1}{2}\vec{\mu}_i^{ind} \vec{E}_i^{perm}\end{split}
\end{equation*}
\sphinxAtStartPar
The formulas for the HIPPO energy terms are:
\begin{equation*}
\begin{split}U_{multipole} = & Z_i \frac{1}{r_{ij}} Z_j + Z_i T_{ij}^{damp} \vec{M}_j + Z_j T_{ji}^{damp} \vec{M}_i + \vec{M}_i T_{ij}^{damp} \vec{M}_j, \quad \mbox{with} \quad
   \vec{M} = \left(q, \vec{\mu}_{perm}, \boldsymbol{\Theta} \right) \\
U_{polar} = & \frac{1}{2}\vec{\mu}_i^{ind} \vec{E}_i^{perm} \\
U_{qxfer} = & \epsilon_i e^{-\eta_j r_{ij}} + \epsilon_j e^{-\eta_i r_{ij}} \\
U_{repulsion} = & \frac{K_i K_j}{r_{ij}} S^2
   S^2 = \left( \int{\phi_i \phi_j} dv \right)^2 = \vec{M}_i\boldsymbol{T_{ij}^{repulsion}}\vec{M}_j \\
U_{dispersion} = & -\frac{C_6^iC_6^j}{r_{ij}^6} \left( f_{damp}^{dispersion} \right)_{ij}^2\end{split}
\end{equation*}
\begin{sphinxadmonition}{note}{Note:}
\sphinxAtStartPar
The AMOEBA and HIPPO force fields compute long\sphinxhyphen{}range charge, dipole,
and quadrupole interactions as well as long\sphinxhyphen{}range dispersion
effects.  However, unlike other models with long\sphinxhyphen{}range interactions
in LAMMPS, this does not require use of a KSpace style via the
{\hyperref[\detokenize{kspace_style::doc}]{\sphinxcrossref{\DUrole{doc}{kspace\_style}}}} command.  That is because for
AMOEBA and HIPPO the long\sphinxhyphen{}range computations are intertwined with
the pairwise computations.  So these pair style include both short\sphinxhyphen{}
and long\sphinxhyphen{}range computations.  This means the energy and virial
computed by the pair style as well as the “Pair” timing reported by
LAMMPS will include the long\sphinxhyphen{}range calculations.
\end{sphinxadmonition}

\sphinxAtStartPar
The implementation of the AMOEBA and HIPPO force fields in LAMMPS was
done using F90 code provided by the Ponder group from their \sphinxhref{https://dasher.wustl.edu/tinker/}{Tinker MD
code}.

\sphinxAtStartPar
The current implementation (July 2022) of AMOEBA in LAMMPS matches the
version discussed in {\hyperref[\detokenize{pair_amoeba:amoeba-ponder}]{\sphinxcrossref{\DUrole{std,std-ref}{(Ponder)}}}}, {\hyperref[\detokenize{pair_amoeba:amoeba-ren}]{\sphinxcrossref{\DUrole{std,std-ref}{(Ren)}}}}, and {\hyperref[\detokenize{pair_amoeba:amoeba-shi}]{\sphinxcrossref{\DUrole{std,std-ref}{(Shi)}}}}.  Likewise the current
implementation of HIPPO in LAMMPS matches the version discussed in
{\hyperref[\detokenize{pair_amoeba:amoeba-rackers}]{\sphinxcrossref{\DUrole{std,std-ref}{(Rackers)}}}}.

\sphinxAtStartPar
\DUrole{versionmodified,added}{Added in version 8Feb2023.}

\sphinxAtStartPar
Accelerator support via the GPU package is available.


\bigskip\hrule\bigskip


\sphinxAtStartPar
Only a single pair\_coeff command is used with either the \sphinxstyleemphasis{amoeba} and
\sphinxstyleemphasis{hippo} styles which specifies two Tinker files, a PRM and KEY file.

\begin{sphinxVerbatim}[commandchars=\\\{\}]
\PYG{k}{pair\PYGZus{}coeff}\PYG{+w}{ }\PYG{o}{*}\PYG{+w}{ }\PYG{o}{*}\PYG{+w}{ }\PYG{n}{..}\PYG{o}{/}\PYG{n}{potentials}\PYG{o}{/}\PYG{n}{protein.prm.amoeba}\PYG{+w}{ }\PYG{n}{..}\PYG{o}{/}\PYG{n}{potentials}\PYG{o}{/}\PYG{n}{protein.key.amoeba}
\PYG{k}{pair\PYGZus{}coeff}\PYG{+w}{ }\PYG{o}{*}\PYG{+w}{ }\PYG{o}{*}\PYG{+w}{ }\PYG{n}{..}\PYG{o}{/}\PYG{n}{potentials}\PYG{o}{/}\PYG{n}{water.prm.hippo}\PYG{+w}{ }\PYG{n}{..}\PYG{o}{/}\PYG{n}{potentials}\PYG{o}{/}\PYG{n}{water.key.hippo}
\end{sphinxVerbatim}

\sphinxAtStartPar
Examples of the PRM files are in the potentials directory with an
*.amoeba or *.hippo suffix.  The examples/amoeba directory has
examples of both PRM and KEY files.

\sphinxAtStartPar
A Tinker PRM file is composed of sections, each of which has multiple
lines.  A Tinker KEY file is composed of lines, each of which has a
keyword followed by zero or more parameters.

\sphinxAtStartPar
The list of PRM sections and KEY keywords which LAMMPS recognizes are
listed on the {\hyperref[\detokenize{Howto_amoeba::doc}]{\sphinxcrossref{\DUrole{doc}{Howto amoeba}}}} doc page.  If not
recognized, the section or keyword is skipped.

\sphinxAtStartPar
Note that if the KEY file is specified as NULL, then no file is
required; default values for various AMOEBA/HIPPO settings are used.
The {\hyperref[\detokenize{Howto_amoeba::doc}]{\sphinxcrossref{\DUrole{doc}{Howto amoeba}}}} doc page also gives the default
settings.


\bigskip\hrule\bigskip


\sphinxAtStartPar
\DUrole{versionmodified,added}{Added in version 3Nov2022.}

\sphinxAtStartPar
The \sphinxstyleemphasis{amoeba} and \sphinxstyleemphasis{hippo} pair styles support extraction of two per\sphinxhyphen{}atom
quantities by the {\hyperref[\detokenize{fix_pair::doc}]{\sphinxcrossref{\DUrole{doc}{fix pair}}}} command.  This allows the
quantities to be output to files by the {\hyperref[\detokenize{dump::doc}]{\sphinxcrossref{\DUrole{doc}{dump}}}} or otherwise
processed by other LAMMPS commands.

\sphinxAtStartPar
The names of the two quantities are “uind” and “uinp” for the induced
dipole moments for each atom.  Neither quantity needs to be triggered by
the {\hyperref[\detokenize{fix_pair::doc}]{\sphinxcrossref{\DUrole{doc}{fix pair}}}} command in order for these pair styles to
calculate it.


\bigskip\hrule\bigskip



\subsection{Mixing, shift, table, tail correction, restart, rRESPA info}
\label{\detokenize{pair_amoeba:mixing-shift-table-tail-correction-restart-rrespa-info}}
\sphinxAtStartPar
These pair styles do not support the {\hyperref[\detokenize{pair_modify::doc}]{\sphinxcrossref{\DUrole{doc}{pair\_modify}}}}
mix, shift, table, and tail options.

\sphinxAtStartPar
These pair styles do not write their information to {\hyperref[\detokenize{restart::doc}]{\sphinxcrossref{\DUrole{doc}{binary
restart files}}}}, since it is stored in potential files.
Thus, you need to re\sphinxhyphen{}specify the pair\_style and pair\_coeff commands in
an input script that reads a restart file.

\sphinxAtStartPar
These pair styles can only be used via the \sphinxstyleemphasis{pair} keyword of the
{\hyperref[\detokenize{run_style::doc}]{\sphinxcrossref{\DUrole{doc}{run\_style respa}}}} command.  They do not support the
\sphinxstyleemphasis{inner}, \sphinxstyleemphasis{middle}, \sphinxstyleemphasis{outer} keywords.


\bigskip\hrule\bigskip


\sphinxAtStartPar
Styles with a \sphinxstyleemphasis{gpu}, \sphinxstyleemphasis{intel}, \sphinxstyleemphasis{kk}, \sphinxstyleemphasis{omp}, or \sphinxstyleemphasis{opt} suffix are
functionally the same as the corresponding style without the suffix.
They have been optimized to run faster, depending on your available
hardware, as discussed on the {\hyperref[\detokenize{Speed_packages::doc}]{\sphinxcrossref{\DUrole{doc}{Accelerator packages}}}}
page.  The accelerated styles take the same arguments and should
produce the same results, except for round\sphinxhyphen{}off and precision issues.

\sphinxAtStartPar
These accelerated styles are part of the GPU, INTEL, KOKKOS,
OPENMP, and OPT packages, respectively.  They are only enabled if
LAMMPS was built with those packages.  See the {\hyperref[\detokenize{Build_package::doc}]{\sphinxcrossref{\DUrole{doc}{Build package}}}} page for more info.

\sphinxAtStartPar
You can specify the accelerated styles explicitly in your input script
by including their suffix, or you can use the {\hyperref[\detokenize{Run_options::doc}]{\sphinxcrossref{\DUrole{doc}{\sphinxhyphen{}suffix command\sphinxhyphen{}line switch}}}} when you invoke LAMMPS, or you can use the
{\hyperref[\detokenize{suffix::doc}]{\sphinxcrossref{\DUrole{doc}{suffix}}}} command in your input script.

\sphinxAtStartPar
See the {\hyperref[\detokenize{Speed_packages::doc}]{\sphinxcrossref{\DUrole{doc}{Accelerator packages}}}} page for more
instructions on how to use the accelerated styles effectively.

\begin{sphinxadmonition}{note}{Note:}
\sphinxAtStartPar
Using the GPU accelerated pair styles ‘amoeba/gpu’ or ‘hippo/gpu’
when compiling the GPU package for OpenCL has a few known issues
when running on integrated GPUs and the calculation may crash.

\sphinxAtStartPar
The GPU accelerated pair styles are also not (yet) compatible
with single precision FFTs.
\end{sphinxadmonition}


\bigskip\hrule\bigskip



\subsection{Restrictions}
\label{\detokenize{pair_amoeba:restrictions}}
\sphinxAtStartPar
These pair styles are part of the AMOEBA package.  They are only
enabled if LAMMPS was built with that package.  See the {\hyperref[\detokenize{Build_package::doc}]{\sphinxcrossref{\DUrole{doc}{Build
package}}}} doc page for more info.

\sphinxAtStartPar
The AMOEBA and HIPPO potential (PRM) and KEY files provided with
LAMMPS in the potentials and examples/amoeba directories are Tinker
files parameterized for Tinker units.  Their numeric parameters are
converted by LAMMPS to its real units {\hyperref[\detokenize{units::doc}]{\sphinxcrossref{\DUrole{doc}{units}}}}.  Thus you
can only use these pair styles with real units.

\sphinxAtStartPar
These potentials do not yet calculate per\sphinxhyphen{}atom energy or virial
contributions.

\sphinxAtStartPar
As explained on the {\hyperref[\detokenize{Howto_amoeba::doc}]{\sphinxcrossref{\DUrole{doc}{AMOEBA and HIPPO howto}}}} page,
use of these pair styles to run a simulation with the AMOEBA or HIPPO
force fields requires several things.

\sphinxAtStartPar
The first is a data file generated by the tools/tinker/tinker2lmp.py
conversion script which uses Tinker file force field file input to
create a data file compatible with LAMMPS.

\sphinxAtStartPar
The second is use of these commands:
\begin{itemize}
\item {} 
\sphinxAtStartPar
{\hyperref[\detokenize{atom_style::doc}]{\sphinxcrossref{\DUrole{doc}{atom\_style amoeba}}}}

\item {} 
\sphinxAtStartPar
{\hyperref[\detokenize{fix_property_atom::doc}]{\sphinxcrossref{\DUrole{doc}{fix property/atom}}}}

\item {} 
\sphinxAtStartPar
{\hyperref[\detokenize{special_bonds::doc}]{\sphinxcrossref{\DUrole{doc}{special\_bonds one/five}}}}

\end{itemize}

\sphinxAtStartPar
And third, depending on the model being simulated, these
commands for intramolecular interactions may also be required:
\begin{itemize}
\item {} 
\sphinxAtStartPar
{\hyperref[\detokenize{bond_class2::doc}]{\sphinxcrossref{\DUrole{doc}{bond\_style class2}}}}

\item {} 
\sphinxAtStartPar
{\hyperref[\detokenize{angle_amoeba::doc}]{\sphinxcrossref{\DUrole{doc}{angle\_style amoeba}}}}

\item {} 
\sphinxAtStartPar
{\hyperref[\detokenize{dihedral_fourier::doc}]{\sphinxcrossref{\DUrole{doc}{dihedral\_style fourier}}}}

\item {} 
\sphinxAtStartPar
{\hyperref[\detokenize{improper_amoeba::doc}]{\sphinxcrossref{\DUrole{doc}{improper\_style amoeba}}}}

\item {} 
\sphinxAtStartPar
{\hyperref[\detokenize{fix_amoeba_pitorsion::doc}]{\sphinxcrossref{\DUrole{doc}{fix amoeba/pitorsion}}}}

\item {} 
\sphinxAtStartPar
{\hyperref[\detokenize{fix_amoeba_bitorsion::doc}]{\sphinxcrossref{\DUrole{doc}{fix amoeba/bitorsion}}}}

\end{itemize}


\bigskip\hrule\bigskip



\subsection{Related commands}
\label{\detokenize{pair_amoeba:related-commands}}
\sphinxAtStartPar
{\hyperref[\detokenize{atom_style::doc}]{\sphinxcrossref{\DUrole{doc}{atom\_style amoeba}}}},
{\hyperref[\detokenize{bond_class2::doc}]{\sphinxcrossref{\DUrole{doc}{bond\_style class2}}}},
{\hyperref[\detokenize{angle_amoeba::doc}]{\sphinxcrossref{\DUrole{doc}{angle\_style amoeba}}}},
{\hyperref[\detokenize{dihedral_fourier::doc}]{\sphinxcrossref{\DUrole{doc}{dihedral\_style fourier}}}},
{\hyperref[\detokenize{improper_amoeba::doc}]{\sphinxcrossref{\DUrole{doc}{improper\_style amoeba}}}},
{\hyperref[\detokenize{fix_amoeba_pitorsion::doc}]{\sphinxcrossref{\DUrole{doc}{fix amoeba/pitorsion}}}},
{\hyperref[\detokenize{fix_amoeba_bitorsion::doc}]{\sphinxcrossref{\DUrole{doc}{fix amoeba/bitorsion}}}},
{\hyperref[\detokenize{special_bonds::doc}]{\sphinxcrossref{\DUrole{doc}{special\_bonds one/five}}}},
{\hyperref[\detokenize{fix_property_atom::doc}]{\sphinxcrossref{\DUrole{doc}{fix property/atom}}}}


\subsection{Default}
\label{\detokenize{pair_amoeba:default}}
\sphinxAtStartPar
none


\bigskip\hrule\bigskip

\phantomsection\label{\detokenize{pair_amoeba:amoeba-ponder}}
\sphinxAtStartPar
\sphinxstylestrong{(Ponder)} Ponder, Wu, Ren, Pande, Chodera, Schnieders, Haque, Mobley, Lambrecht, DiStasio Jr, M. Head\sphinxhyphen{}Gordon, Clark, Johnson, T. Head\sphinxhyphen{}Gordon, J Phys Chem B, 114, 2549\sphinxhyphen{}2564 (2010).

\phantomsection\label{\detokenize{pair_amoeba:amoeba-rackers}}
\sphinxAtStartPar
\sphinxstylestrong{(Rackers)} Rackers, Silva, Wang, Ponder, J Chem Theory Comput, 17, 7056\sphinxhyphen{}7084 (2021).

\phantomsection\label{\detokenize{pair_amoeba:amoeba-ren}}
\sphinxAtStartPar
\sphinxstylestrong{(Ren)} Ren and Ponder, J Phys Chem B, 107, 5933 (2003).

\phantomsection\label{\detokenize{pair_amoeba:amoeba-shi}}
\sphinxAtStartPar
\sphinxstylestrong{(Shi)} Shi, Xia, Zhang, Best, Wu, Ponder, Ren, J Chem Theory Comp, 9, 4046, 2013.

\sphinxstepscope

\index{pair\_style atm@\spxentry{pair\_style atm}}\ignorespaces 

\section{pair\_style atm command}
\label{\detokenize{pair_atm:pair-style-atm-command}}\label{\detokenize{pair_atm:index-0}}\label{\detokenize{pair_atm::doc}}

\subsection{Syntax}
\label{\detokenize{pair_atm:syntax}}
\begin{sphinxVerbatim}[commandchars=\\\{\}]
\PYG{k}{pair\PYGZus{}style}\PYG{+w}{ }\PYG{n}{atm}\PYG{+w}{ }\PYG{n}{cutoff}\PYG{+w}{ }\PYG{n}{cutoff\PYGZus{}triple}
\end{sphinxVerbatim}
\begin{itemize}
\item {} 
\sphinxAtStartPar
cutoff = cutoff for each pair in 3\sphinxhyphen{}body interaction (distance units)

\item {} 
\sphinxAtStartPar
cutoff\_triple = additional cutoff applied to product of 3 pairwise distances (distance units)

\end{itemize}


\subsection{Examples}
\label{\detokenize{pair_atm:examples}}
\begin{sphinxVerbatim}[commandchars=\\\{\}]
\PYG{k}{pair\PYGZus{}style}\PYG{+w}{ }\PYG{n}{atm}\PYG{+w}{ }\PYG{l+m}{4.5}\PYG{+w}{ }\PYG{l+m}{2.5}
\PYG{k}{pair\PYGZus{}coeff}\PYG{+w}{ }\PYG{o}{*}\PYG{+w}{ }\PYG{o}{*}\PYG{+w}{ }\PYG{o}{*}\PYG{+w}{ }\PYG{l+m}{0.072}

\PYG{k}{pair\PYGZus{}style}\PYG{+w}{ }\PYG{n}{hybrid}\PYG{o}{/}\PYG{n}{overlay}\PYG{+w}{ }\PYG{n}{lj}\PYG{o}{/}\PYG{n}{cut}\PYG{+w}{ }\PYG{l+m}{6.5}\PYG{+w}{ }\PYG{n}{atm}\PYG{+w}{ }\PYG{l+m}{4.5}\PYG{+w}{ }\PYG{l+m}{2.5}
\PYG{k}{pair\PYGZus{}coeff}\PYG{+w}{ }\PYG{o}{*}\PYG{+w}{ }\PYG{o}{*}\PYG{+w}{ }\PYG{n}{lj}\PYG{o}{/}\PYG{n}{cut}\PYG{+w}{ }\PYG{l+m}{1.0}\PYG{+w}{ }\PYG{l+m}{1.0}
\PYG{k}{pair\PYGZus{}coeff}\PYG{+w}{ }\PYG{l+m}{1}\PYG{+w}{ }\PYG{l+m}{1}\PYG{+w}{ }\PYG{n}{atm}\PYG{+w}{ }\PYG{l+m}{1}\PYG{+w}{ }\PYG{l+m}{0.064}
\PYG{k}{pair\PYGZus{}coeff}\PYG{+w}{ }\PYG{l+m}{1}\PYG{+w}{ }\PYG{l+m}{1}\PYG{+w}{ }\PYG{n}{atm}\PYG{+w}{ }\PYG{l+m}{2}\PYG{+w}{ }\PYG{l+m}{0.080}
\PYG{k}{pair\PYGZus{}coeff}\PYG{+w}{ }\PYG{l+m}{1}\PYG{+w}{ }\PYG{l+m}{2}\PYG{+w}{ }\PYG{n}{atm}\PYG{+w}{ }\PYG{l+m}{2}\PYG{+w}{ }\PYG{l+m}{0.100}
\PYG{k}{pair\PYGZus{}coeff}\PYG{+w}{ }\PYG{l+m}{2}\PYG{+w}{ }\PYG{l+m}{2}\PYG{+w}{ }\PYG{n}{atm}\PYG{+w}{ }\PYG{l+m}{2}\PYG{+w}{ }\PYG{l+m}{0.125}
\end{sphinxVerbatim}


\subsection{Description}
\label{\detokenize{pair_atm:description}}
\sphinxAtStartPar
The \sphinxstyleemphasis{atm} style computes a 3\sphinxhyphen{}body {\hyperref[\detokenize{pair_atm:axilrod}]{\sphinxcrossref{\DUrole{std,std-ref}{Axilrod\sphinxhyphen{}Teller\sphinxhyphen{}Muto}}}}
potential for the energy E of a system of atoms as
\begin{equation*}
\begin{split}E & = \nu\frac{1+3\cos\gamma_1\cos\gamma_2\cos\gamma_3}{r_{12}^3r_{23}^3r_{31}^3} \\\end{split}
\end{equation*}
\sphinxAtStartPar
where \(\nu\) is the three\sphinxhyphen{}body interaction strength.  The distances
between pairs of atoms \(r_{12}\), \(r_{23}\), \(r_{31}\) and the angles \(\gamma_1\), \(\gamma_2\),
\(\gamma_3\) are as shown in this diagram:

\noindent{\hspace*{\fill}\sphinxincludegraphics{{pair_atm_dia}.jpg}\hspace*{\fill}}

\sphinxAtStartPar
Note that for the interaction between a triplet of atoms \(I,J,K\), there
is no “central” atom.  The interaction is symmetric with respect to
permutation of the three atoms. Thus the \(\nu\) value is
the same for all those permutations of the atom types of \(I,J,K\)
and needs to be specified only once, as discussed below.

\sphinxAtStartPar
The \sphinxstyleemphasis{atm} potential is typically used in combination with a two\sphinxhyphen{}body
potential using the {\hyperref[\detokenize{pair_hybrid::doc}]{\sphinxcrossref{\DUrole{doc}{pair\_style hybrid/overlay}}}}
command as in the example above.

\sphinxAtStartPar
The potential for a triplet of atom is calculated only if all 3 distances
\(r_{12}\), \(r_{23}\), \(r_{31}\) between the three atoms satisfy
\(r_{IJ} < \text{cutoff}\).  In addition, the product of the 3 distances
\(r_{12} r_{23} r_{31}\) \textless{} cutoff\_triple \(^3\) is required, which
excludes from calculation the triplets with small contribution to the
interaction.

\sphinxAtStartPar
The following coefficients must be defined for each pair of atoms
types via the {\hyperref[\detokenize{pair_coeff::doc}]{\sphinxcrossref{\DUrole{doc}{pair\_coeff}}}} command as in the examples
above, or in the restart files read by the
{\hyperref[\detokenize{read_restart::doc}]{\sphinxcrossref{\DUrole{doc}{read\_restart}}}} commands:
\begin{itemize}
\item {} 
\sphinxAtStartPar
\(K\) = atom type of the third atom (1 to \(N_{\text{types}}\))

\item {} 
\sphinxAtStartPar
\(\nu\) = prefactor (energy/distance\textasciicircum{}9 units)

\end{itemize}

\sphinxAtStartPar
\(K\) can be specified in one of two ways.  An explicit numeric value
or type label can be used, as in the second example above.   LAMMPS
sets the coefficients for the other 5 symmetric interactions to the same
values.  E.g. if \(I = 1\), \(J = 2\), \(K = 3\), then these 6
values are set to the specified \(\nu\): \(\nu_{123}\),
\(\nu_{132}\), \(\nu_{213}\), \(\nu_{231}\), \(\nu_{312}\),
\(\nu_{321}\).  This enforces the symmetry discussed above.

\sphinxAtStartPar
A wildcard asterisk can be used for K to set the coefficients for
multiple triplets of atom types.  This takes the form “*” or “*n” or
“n*” or “m*n”.  If \(N\) equals the number of atom types, then an asterisk with
no numeric values means all types from 1 to \(N\).  A leading asterisk
means all types from 1 to \(n\) (inclusive).  A trailing asterisk means
all types from \(n\) to \(N\) (inclusive).  A middle asterisk means all types
from \(m\) to \(n\) (inclusive).  Note that only type triplets with \(J \leq K\) are
considered; if asterisks imply type triplets where \(K < J\), they are
ignored.

\sphinxAtStartPar
Note that a pair\_coeff command can override a previous setting for the
same \(I,J,K\) triplet.  For example, these commands set \(\nu\) for all \(I,J.K\)
triplets, then overwrite nu for just the \(I,J,K = 2,3,4\) triplet:

\begin{sphinxVerbatim}[commandchars=\\\{\}]
\PYG{k}{pair\PYGZus{}coeff}\PYG{+w}{ }\PYG{o}{*}\PYG{+w}{ }\PYG{o}{*}\PYG{+w}{ }\PYG{o}{*}\PYG{+w}{ }\PYG{l+m}{0.25}
\PYG{k}{pair\PYGZus{}coeff}\PYG{+w}{ }\PYG{l+m}{2}\PYG{+w}{ }\PYG{l+m}{3}\PYG{+w}{ }\PYG{l+m}{4}\PYG{+w}{ }\PYG{l+m}{0.1}
\end{sphinxVerbatim}

\sphinxAtStartPar
Note that for a simulation with a single atom type, only a single
entry is required, e.g.

\begin{sphinxVerbatim}[commandchars=\\\{\}]
\PYG{k}{pair\PYGZus{}coeff}\PYG{+w}{ }\PYG{l+m}{1}\PYG{+w}{ }\PYG{l+m}{1}\PYG{+w}{ }\PYG{l+m}{1}\PYG{+w}{ }\PYG{l+m}{0.25}
\end{sphinxVerbatim}

\sphinxAtStartPar
For a simulation with two atom types, four pair\_coeff commands will
specify all possible nu values:

\begin{sphinxVerbatim}[commandchars=\\\{\}]
\PYG{k}{pair\PYGZus{}coeff}\PYG{+w}{ }\PYG{l+m}{1}\PYG{+w}{ }\PYG{l+m}{1}\PYG{+w}{ }\PYG{l+m}{1}\PYG{+w}{ }\PYG{n}{nu1}
\PYG{k}{pair\PYGZus{}coeff}\PYG{+w}{ }\PYG{l+m}{1}\PYG{+w}{ }\PYG{l+m}{1}\PYG{+w}{ }\PYG{l+m}{2}\PYG{+w}{ }\PYG{n}{nu2}
\PYG{k}{pair\PYGZus{}coeff}\PYG{+w}{ }\PYG{l+m}{1}\PYG{+w}{ }\PYG{l+m}{2}\PYG{+w}{ }\PYG{l+m}{2}\PYG{+w}{ }\PYG{n}{nu3}
\PYG{k}{pair\PYGZus{}coeff}\PYG{+w}{ }\PYG{l+m}{2}\PYG{+w}{ }\PYG{l+m}{2}\PYG{+w}{ }\PYG{l+m}{2}\PYG{+w}{ }\PYG{n}{nu4}
\end{sphinxVerbatim}

\sphinxAtStartPar
For a simulation with three atom types, ten pair\_coeff commands will
specify all possible nu values:

\begin{sphinxVerbatim}[commandchars=\\\{\}]
\PYG{k}{pair\PYGZus{}coeff}\PYG{+w}{ }\PYG{l+m}{1}\PYG{+w}{ }\PYG{l+m}{1}\PYG{+w}{ }\PYG{l+m}{1}\PYG{+w}{ }\PYG{n}{nu1}
\PYG{k}{pair\PYGZus{}coeff}\PYG{+w}{ }\PYG{l+m}{1}\PYG{+w}{ }\PYG{l+m}{1}\PYG{+w}{ }\PYG{l+m}{2}\PYG{+w}{ }\PYG{n}{nu2}
\PYG{k}{pair\PYGZus{}coeff}\PYG{+w}{ }\PYG{l+m}{1}\PYG{+w}{ }\PYG{l+m}{1}\PYG{+w}{ }\PYG{l+m}{3}\PYG{+w}{ }\PYG{n}{nu3}
\PYG{k}{pair\PYGZus{}coeff}\PYG{+w}{ }\PYG{l+m}{1}\PYG{+w}{ }\PYG{l+m}{2}\PYG{+w}{ }\PYG{l+m}{2}\PYG{+w}{ }\PYG{n}{nu4}
\PYG{k}{pair\PYGZus{}coeff}\PYG{+w}{ }\PYG{l+m}{1}\PYG{+w}{ }\PYG{l+m}{2}\PYG{+w}{ }\PYG{l+m}{3}\PYG{+w}{ }\PYG{n}{nu5}
\PYG{k}{pair\PYGZus{}coeff}\PYG{+w}{ }\PYG{l+m}{1}\PYG{+w}{ }\PYG{l+m}{3}\PYG{+w}{ }\PYG{l+m}{3}\PYG{+w}{ }\PYG{n}{nu6}
\PYG{k}{pair\PYGZus{}coeff}\PYG{+w}{ }\PYG{l+m}{2}\PYG{+w}{ }\PYG{l+m}{2}\PYG{+w}{ }\PYG{l+m}{2}\PYG{+w}{ }\PYG{n}{nu7}
\PYG{k}{pair\PYGZus{}coeff}\PYG{+w}{ }\PYG{l+m}{2}\PYG{+w}{ }\PYG{l+m}{2}\PYG{+w}{ }\PYG{l+m}{3}\PYG{+w}{ }\PYG{n}{nu8}
\PYG{k}{pair\PYGZus{}coeff}\PYG{+w}{ }\PYG{l+m}{2}\PYG{+w}{ }\PYG{l+m}{3}\PYG{+w}{ }\PYG{l+m}{3}\PYG{+w}{ }\PYG{n}{nu9}
\PYG{k}{pair\PYGZus{}coeff}\PYG{+w}{ }\PYG{l+m}{3}\PYG{+w}{ }\PYG{l+m}{3}\PYG{+w}{ }\PYG{l+m}{3}\PYG{+w}{ }\PYG{n}{nu10}
\end{sphinxVerbatim}

\sphinxAtStartPar
By default the \(\nu\) value for all triplets is set to 0.0.  Thus it is
not required to provide pair\_coeff commands that enumerate triplet
interactions for all \(K\) types.  If some \(I,J,K\) combination is not
specified, then there will be no 3\sphinxhyphen{}body ATM interactions for that
combination and all its permutations.  However, as with all pair
styles, it is required to specify a pair\_coeff command for all \(I,J\)
combinations, else an error will result.


\bigskip\hrule\bigskip



\subsection{Mixing, shift, table, tail correction, restart, rRESPA info}
\label{\detokenize{pair_atm:mixing-shift-table-tail-correction-restart-rrespa-info}}
\sphinxAtStartPar
This pair style do not support the {\hyperref[\detokenize{pair_modify::doc}]{\sphinxcrossref{\DUrole{doc}{pair\_modify}}}}
mix, shift, table, and tail options.

\sphinxAtStartPar
This pair style writes its information to {\hyperref[\detokenize{restart::doc}]{\sphinxcrossref{\DUrole{doc}{binary restart files}}}}, so pair\_style and pair\_coeff commands do not need to be specified
in an input script that reads a restart file.  However, if the \sphinxstyleemphasis{atm} potential
is used in combination with other potentials using the {\hyperref[\detokenize{pair_hybrid::doc}]{\sphinxcrossref{\DUrole{doc}{pair\_style
hybrid/overlay}}}} command then pair\_coeff commands need to be
re\sphinxhyphen{}specified in the restart input script.

\sphinxAtStartPar
This pair style can only be used via the \sphinxstyleemphasis{pair} keyword of the
{\hyperref[\detokenize{run_style::doc}]{\sphinxcrossref{\DUrole{doc}{run\_style respa}}}} command.  It does not support the
\sphinxstyleemphasis{inner}, \sphinxstyleemphasis{middle}, and \sphinxstyleemphasis{outer} keywords.


\bigskip\hrule\bigskip



\subsection{Restrictions}
\label{\detokenize{pair_atm:restrictions}}
\sphinxAtStartPar
This pair style is part of the MANYBODY package.  It is only enabled
if LAMMPS was built with that package.  See the {\hyperref[\detokenize{Build_package::doc}]{\sphinxcrossref{\DUrole{doc}{Build package}}}} page for more info.


\subsection{Related commands}
\label{\detokenize{pair_atm:related-commands}}
\sphinxAtStartPar
{\hyperref[\detokenize{pair_coeff::doc}]{\sphinxcrossref{\DUrole{doc}{pair\_coeff}}}}


\subsection{Default}
\label{\detokenize{pair_atm:default}}
\sphinxAtStartPar
none


\bigskip\hrule\bigskip

\phantomsection\label{\detokenize{pair_atm:axilrod}}
\sphinxAtStartPar
\sphinxstylestrong{(Axilrod)}
Axilrod and Teller, J Chem Phys, 11, 299 (1943);
Muto, Nippon Sugaku\sphinxhyphen{}Buturigakkwaishi 17, 629 (1943).

\sphinxstepscope

\index{pair\_style awpmd/cut@\spxentry{pair\_style awpmd/cut}}\ignorespaces 

\section{pair\_style awpmd/cut command}
\label{\detokenize{pair_awpmd:pair-style-awpmd-cut-command}}\label{\detokenize{pair_awpmd:index-0}}\label{\detokenize{pair_awpmd::doc}}

\subsection{Syntax}
\label{\detokenize{pair_awpmd:syntax}}
\begin{sphinxVerbatim}[commandchars=\\\{\}]
\PYG{k}{pair\PYGZus{}style}\PYG{+w}{ }\PYG{n}{awpmd}\PYG{o}{/}\PYG{n}{cut}\PYG{+w}{ }\PYG{n}{Rc}\PYG{+w}{ }\PYG{n}{keyword}\PYG{+w}{ }\PYG{n}{value}\PYG{+w}{ }\PYG{n}{...}
\end{sphinxVerbatim}
\begin{itemize}
\item {} 
\sphinxAtStartPar
Rc = global cutoff, \sphinxhyphen{}1 means cutoff of half the shortest box length

\item {} 
\sphinxAtStartPar
zero or more keyword/value pairs may be appended

\item {} 
\sphinxAtStartPar
keyword = \sphinxstyleemphasis{hartree} or \sphinxstyleemphasis{dproduct} or \sphinxstyleemphasis{uhf} or \sphinxstyleemphasis{free} or \sphinxstyleemphasis{pbc} or \sphinxstyleemphasis{fix} or \sphinxstyleemphasis{harm} or \sphinxstyleemphasis{ermscale} or \sphinxstyleemphasis{flex\_press}
\begin{sphinxalltt}
\sphinxstyleemphasis{hartree} value = none
\sphinxstyleemphasis{dproduct} value = none
\sphinxstyleemphasis{uhf} value = none
\sphinxstyleemphasis{free} value = none
\sphinxstyleemphasis{pbc} value = Plen
  Plen = periodic width of electron = \sphinxhyphen{}1 or positive value (distance units)
\sphinxstyleemphasis{fix} value = Flen
  Flen = fixed width of electron = \sphinxhyphen{}1 or positive value (distance units)
\sphinxstyleemphasis{harm} value = width
  width = harmonic width constraint
\sphinxstyleemphasis{ermscale} value = factor
  factor = scaling between electron mass and width variable mass
\sphinxstyleemphasis{flex\_press} value = none
\end{sphinxalltt}

\end{itemize}


\subsection{Examples}
\label{\detokenize{pair_awpmd:examples}}
\begin{sphinxVerbatim}[commandchars=\\\{\}]
\PYG{k}{pair\PYGZus{}style}\PYG{+w}{ }\PYG{n}{awpmd}\PYG{o}{/}\PYG{n}{cut}\PYG{+w}{ }\PYG{o}{\PYGZhy{}}\PYG{l+m}{1}
\PYG{k}{pair\PYGZus{}style}\PYG{+w}{ }\PYG{n}{awpmd}\PYG{o}{/}\PYG{n}{cut}\PYG{+w}{ }\PYG{l+m}{40.0}\PYG{+w}{ }\PYG{n}{uhf}\PYG{+w}{ }\PYG{n}{free}
\PYG{k}{pair\PYGZus{}coeff}\PYG{+w}{ }\PYG{o}{*}\PYG{+w}{ }\PYG{o}{*}
\PYG{k}{pair\PYGZus{}coeff}\PYG{+w}{ }\PYG{l+m}{2}\PYG{+w}{ }\PYG{l+m}{2}\PYG{+w}{ }\PYG{l+m}{20.0}
\end{sphinxVerbatim}


\subsection{Description}
\label{\detokenize{pair_awpmd:description}}
\sphinxAtStartPar
This pair style contains an implementation of the Antisymmetrized Wave
Packet Molecular Dynamics (AWPMD) method.  Need citation here.  Need
basic formulas here.  Could be links to other documents.

\sphinxAtStartPar
Rc is the cutoff.

\sphinxAtStartPar
The pair\_style command allows for several optional keywords
to be specified.

\sphinxAtStartPar
The \sphinxstyleemphasis{hartree}, \sphinxstyleemphasis{dproduct}, and \sphinxstyleemphasis{uhf} keywords specify the form of the
initial trial wave function for the system.  If the \sphinxstyleemphasis{hartree} keyword
is used, then a Hartree multielectron trial wave function is used.  If
the \sphinxstyleemphasis{dproduct} keyword is used, then a trial function which is a
product of two determinants for each spin type is used.  If the \sphinxstyleemphasis{uhf}
keyword is used, then an unrestricted Hartree\sphinxhyphen{}Fock trial wave function
is used.

\sphinxAtStartPar
The \sphinxstyleemphasis{free}, \sphinxstyleemphasis{pbc}, and \sphinxstyleemphasis{fix} keywords specify a width constraint on
the electron wave packets.  If the \sphinxstyleemphasis{free} keyword is specified, then there is no
constraint.  If the \sphinxstyleemphasis{pbc} keyword is used and \sphinxstyleemphasis{Plen} is specified as
\sphinxhyphen{}1, then the maximum width is half the shortest box length.  If \sphinxstyleemphasis{Plen}
is a positive value, then the value is the maximum width.  If the
\sphinxstyleemphasis{fix} keyword is used and \sphinxstyleemphasis{Flen} is specified as \sphinxhyphen{}1, then electrons
have a constant width that is read from the data file.  If \sphinxstyleemphasis{Flen} is a
positive value, then the constant width for all electrons is set to
\sphinxstyleemphasis{Flen}.

\sphinxAtStartPar
The \sphinxstyleemphasis{harm} keyword allow oscillations in the width of the
electron wave packets.  More details are needed.

\sphinxAtStartPar
The \sphinxstyleemphasis{ermscale} keyword specifies a unitless scaling factor
between the electron masses and the width variable mass.  More
details needed.

\sphinxAtStartPar
If the \sphinxstyleemphasis{flex\_press} keyword is used, then a contribution from the
electrons is added to the total virial and pressure of the system.

\sphinxAtStartPar
This potential is designed to be used with {\hyperref[\detokenize{atom_style::doc}]{\sphinxcrossref{\DUrole{doc}{atom\_style wavepacket}}}} definitions, in order to handle the
description of systems with interacting nuclei and explicit electrons.

\sphinxAtStartPar
The following coefficients must be defined for each pair of atoms
types via the {\hyperref[\detokenize{pair_coeff::doc}]{\sphinxcrossref{\DUrole{doc}{pair\_coeff}}}} command as in the examples
above, or in the data file or restart files read by the
{\hyperref[\detokenize{read_data::doc}]{\sphinxcrossref{\DUrole{doc}{read\_data}}}} or {\hyperref[\detokenize{read_restart::doc}]{\sphinxcrossref{\DUrole{doc}{read\_restart}}}}
commands, or by mixing as described below:
\begin{itemize}
\item {} 
\sphinxAtStartPar
cutoff (distance units)

\end{itemize}

\sphinxAtStartPar
For \sphinxstyleemphasis{awpmd/cut}, the cutoff coefficient is optional.  If it is not
used (as in some of the examples above), the default global value
specified in the pair\_style command is used.


\bigskip\hrule\bigskip



\subsection{Mixing, shift, table, tail correction, restart, rRESPA info}
\label{\detokenize{pair_awpmd:mixing-shift-table-tail-correction-restart-rrespa-info}}
\sphinxAtStartPar
The {\hyperref[\detokenize{pair_modify::doc}]{\sphinxcrossref{\DUrole{doc}{pair\_modify}}}} mix, shift, table, and tail options
are not relevant for this pair style.

\sphinxAtStartPar
This pair style writes its information to {\hyperref[\detokenize{restart::doc}]{\sphinxcrossref{\DUrole{doc}{binary restart files}}}}, so pair\_style and pair\_coeff commands do not need
to be specified in an input script that reads a restart file.

\sphinxAtStartPar
This pair style can only be used via the \sphinxstyleemphasis{pair} keyword of the
{\hyperref[\detokenize{run_style::doc}]{\sphinxcrossref{\DUrole{doc}{run\_style respa}}}} command.  It does not support the
\sphinxstyleemphasis{inner}, \sphinxstyleemphasis{middle}, \sphinxstyleemphasis{outer} keywords.


\bigskip\hrule\bigskip



\subsection{Restrictions}
\label{\detokenize{pair_awpmd:restrictions}}\begin{quote}

\sphinxAtStartPar
none
\end{quote}


\subsection{Related commands}
\label{\detokenize{pair_awpmd:related-commands}}
\sphinxAtStartPar
{\hyperref[\detokenize{pair_coeff::doc}]{\sphinxcrossref{\DUrole{doc}{pair\_coeff}}}}


\subsection{Default}
\label{\detokenize{pair_awpmd:default}}
\sphinxAtStartPar
These are the defaults for the pair\_style keywords: \sphinxstyleemphasis{hartree} for the
initial wave function, \sphinxstyleemphasis{free} for the wave packet width.

\sphinxstepscope

\index{pair\_style beck@\spxentry{pair\_style beck}}\ignorespaces 
\index{pair\_style beck/gpu@\spxentry{pair\_style beck/gpu}}\ignorespaces 
\index{pair\_style beck/omp@\spxentry{pair\_style beck/omp}}\ignorespaces 

\section{pair\_style beck command}
\label{\detokenize{pair_beck:pair-style-beck-command}}\label{\detokenize{pair_beck:index-2}}\label{\detokenize{pair_beck:index-1}}\label{\detokenize{pair_beck:index-0}}\label{\detokenize{pair_beck::doc}}
\sphinxAtStartPar
Accelerator Variants: \sphinxstyleemphasis{beck/gpu}, \sphinxstyleemphasis{beck/omp}


\subsection{Syntax}
\label{\detokenize{pair_beck:syntax}}
\begin{sphinxVerbatim}[commandchars=\\\{\}]
\PYG{k}{pair\PYGZus{}style}\PYG{+w}{ }\PYG{n}{beck}\PYG{+w}{ }\PYG{n}{Rc}
\end{sphinxVerbatim}
\begin{itemize}
\item {} 
\sphinxAtStartPar
Rc = cutoff for interactions (distance units)

\end{itemize}


\subsection{Examples}
\label{\detokenize{pair_beck:examples}}
\begin{sphinxVerbatim}[commandchars=\\\{\}]
\PYG{k}{pair\PYGZus{}style}\PYG{+w}{ }\PYG{n}{beck}\PYG{+w}{ }\PYG{l+m}{8.0}
\PYG{k}{pair\PYGZus{}coeff}\PYG{+w}{ }\PYG{o}{*}\PYG{+w}{ }\PYG{o}{*}\PYG{+w}{ }\PYG{l+m}{399.671876712}\PYG{+w}{ }\PYG{l+m}{0.0000867636112694}\PYG{+w}{ }\PYG{l+m}{0.675}\PYG{+w}{ }\PYG{l+m}{4.390}\PYG{+w}{ }\PYG{l+m}{0.0003746}
\PYG{k}{pair\PYGZus{}coeff}\PYG{+w}{ }\PYG{l+m}{1}\PYG{+w}{ }\PYG{l+m}{1}\PYG{+w}{ }\PYG{l+m}{399.671876712}\PYG{+w}{ }\PYG{l+m}{0.0000867636112694}\PYG{+w}{ }\PYG{l+m}{0.675}\PYG{+w}{ }\PYG{l+m}{4.390}\PYG{+w}{ }\PYG{l+m}{0.0003746}\PYG{+w}{ }\PYG{l+m}{6.0}
\end{sphinxVerbatim}


\subsection{Description}
\label{\detokenize{pair_beck:description}}
\sphinxAtStartPar
Style \sphinxstyleemphasis{beck} computes interactions based on the potential by
{\hyperref[\detokenize{pair_beck:beck}]{\sphinxcrossref{\DUrole{std,std-ref}{(Beck)}}}}, originally designed for simulation of Helium.  It
includes truncation at a cutoff distance \(r_c\).
\begin{equation*}
\begin{split}E(r) &= A \exp\left[-\alpha r - \beta r^6\right] - \frac{B}{\left(r^2+a^2\right)^3} \left(1+\frac{2.709+3a^2}{r^2+a^2}\right) \qquad r < r_c \\\end{split}
\end{equation*}
\sphinxAtStartPar
The following coefficients must be defined for each pair of atoms
types via the {\hyperref[\detokenize{pair_coeff::doc}]{\sphinxcrossref{\DUrole{doc}{pair\_coeff}}}} command as in the examples
above, or in the data file or restart files read by the
{\hyperref[\detokenize{read_data::doc}]{\sphinxcrossref{\DUrole{doc}{read\_data}}}} or {\hyperref[\detokenize{read_restart::doc}]{\sphinxcrossref{\DUrole{doc}{read\_restart}}}}
commands.
\begin{itemize}
\item {} 
\sphinxAtStartPar
\(A\) (energy units)

\item {} 
\sphinxAtStartPar
\(B\) (energy\sphinxhyphen{}distance\textasciicircum{}6 units)

\item {} 
\sphinxAtStartPar
\(a\) (distance units)

\item {} 
\sphinxAtStartPar
\(\alpha\) (1/distance units)

\item {} 
\sphinxAtStartPar
\(\beta\)  (1/distance\textasciicircum{}6 units)

\item {} 
\sphinxAtStartPar
cutoff (distance units)

\end{itemize}

\sphinxAtStartPar
The last coefficient is optional.  If not specified, the global cutoff
\(r_c\) is used.


\bigskip\hrule\bigskip


\sphinxAtStartPar
Styles with a \sphinxstyleemphasis{gpu}, \sphinxstyleemphasis{intel}, \sphinxstyleemphasis{kk}, \sphinxstyleemphasis{omp}, or \sphinxstyleemphasis{opt} suffix are
functionally the same as the corresponding style without the suffix.
They have been optimized to run faster, depending on your available
hardware, as discussed on the {\hyperref[\detokenize{Speed_packages::doc}]{\sphinxcrossref{\DUrole{doc}{Accelerator packages}}}}
page.  The accelerated styles take the same arguments and should
produce the same results, except for round\sphinxhyphen{}off and precision issues.

\sphinxAtStartPar
These accelerated styles are part of the GPU, INTEL, KOKKOS,
OPENMP, and OPT packages, respectively.  They are only enabled if
LAMMPS was built with those packages.  See the {\hyperref[\detokenize{Build_package::doc}]{\sphinxcrossref{\DUrole{doc}{Build package}}}} page for more info.

\sphinxAtStartPar
You can specify the accelerated styles explicitly in your input script
by including their suffix, or you can use the {\hyperref[\detokenize{Run_options::doc}]{\sphinxcrossref{\DUrole{doc}{\sphinxhyphen{}suffix command\sphinxhyphen{}line switch}}}} when you invoke LAMMPS, or you can use the
{\hyperref[\detokenize{suffix::doc}]{\sphinxcrossref{\DUrole{doc}{suffix}}}} command in your input script.

\sphinxAtStartPar
See the {\hyperref[\detokenize{Speed_packages::doc}]{\sphinxcrossref{\DUrole{doc}{Accelerator packages}}}} page for more
instructions on how to use the accelerated styles effectively.


\bigskip\hrule\bigskip



\subsection{Mixing, shift, table, tail correction, restart, rRESPA info}
\label{\detokenize{pair_beck:mixing-shift-table-tail-correction-restart-rrespa-info}}
\sphinxAtStartPar
For atom type pairs I,J and I != J, coefficients must be specified.
No default mixing rules are used.

\sphinxAtStartPar
This pair style does not support the {\hyperref[\detokenize{pair_modify::doc}]{\sphinxcrossref{\DUrole{doc}{pair\_modify}}}} shift
option for the energy of the pair interaction.

\sphinxAtStartPar
The {\hyperref[\detokenize{pair_modify::doc}]{\sphinxcrossref{\DUrole{doc}{pair\_modify}}}} table option is not relevant
for this pair style.

\sphinxAtStartPar
This pair style does not support the {\hyperref[\detokenize{pair_modify::doc}]{\sphinxcrossref{\DUrole{doc}{pair\_modify}}}}
tail option for adding long\sphinxhyphen{}range tail corrections.

\sphinxAtStartPar
This pair style writes its information to {\hyperref[\detokenize{restart::doc}]{\sphinxcrossref{\DUrole{doc}{binary restart files}}}}, so pair\_style and pair\_coeff commands do not need
to be specified in an input script that reads a restart file.

\sphinxAtStartPar
This pair style can only be used via the \sphinxstyleemphasis{pair} keyword of the
{\hyperref[\detokenize{run_style::doc}]{\sphinxcrossref{\DUrole{doc}{run\_style respa}}}} command.  It does not support the
\sphinxstyleemphasis{inner}, \sphinxstyleemphasis{middle}, \sphinxstyleemphasis{outer} keywords.


\bigskip\hrule\bigskip



\subsection{Restrictions}
\label{\detokenize{pair_beck:restrictions}}
\sphinxAtStartPar
This pair style is part of the EXTRA\sphinxhyphen{}PAIR package.  It is only enabled if
LAMMPS was built with that package.  See the
{\hyperref[\detokenize{Build_package::doc}]{\sphinxcrossref{\DUrole{doc}{Build package}}}} page for more info.


\subsection{Related commands}
\label{\detokenize{pair_beck:related-commands}}
\sphinxAtStartPar
{\hyperref[\detokenize{pair_coeff::doc}]{\sphinxcrossref{\DUrole{doc}{pair\_coeff}}}}


\subsection{Default}
\label{\detokenize{pair_beck:default}}
\sphinxAtStartPar
none


\bigskip\hrule\bigskip

\phantomsection\label{\detokenize{pair_beck:beck}}
\sphinxAtStartPar
\sphinxstylestrong{(Beck)} Beck, Molecular Physics, 14, 311 (1968).

\sphinxstepscope

\index{pair\_style body/nparticle@\spxentry{pair\_style body/nparticle}}\ignorespaces 

\section{pair\_style body/nparticle command}
\label{\detokenize{pair_body_nparticle:pair-style-body-nparticle-command}}\label{\detokenize{pair_body_nparticle:index-0}}\label{\detokenize{pair_body_nparticle::doc}}

\subsection{Syntax}
\label{\detokenize{pair_body_nparticle:syntax}}
\begin{sphinxVerbatim}[commandchars=\\\{\}]
\PYG{k}{pair\PYGZus{}style}\PYG{+w}{ }\PYG{n}{body}\PYG{o}{/}\PYG{n}{nparticle}\PYG{+w}{ }\PYG{n}{cutoff}
\end{sphinxVerbatim}

\sphinxAtStartPar
cutoff = global cutoff for interactions (distance units)


\subsection{Examples}
\label{\detokenize{pair_body_nparticle:examples}}
\begin{sphinxVerbatim}[commandchars=\\\{\}]
\PYG{k}{pair\PYGZus{}style}\PYG{+w}{ }\PYG{n}{body}\PYG{o}{/}\PYG{n}{nparticle}\PYG{+w}{ }\PYG{l+m}{3.0}
\PYG{k}{pair\PYGZus{}coeff}\PYG{+w}{ }\PYG{o}{*}\PYG{+w}{ }\PYG{o}{*}\PYG{+w}{ }\PYG{l+m}{1.0}\PYG{+w}{ }\PYG{l+m}{1.0}
\PYG{k}{pair\PYGZus{}coeff}\PYG{+w}{ }\PYG{l+m}{1}\PYG{+w}{ }\PYG{l+m}{1}\PYG{+w}{ }\PYG{l+m}{1.0}\PYG{+w}{ }\PYG{l+m}{1.5}\PYG{+w}{ }\PYG{l+m}{2.5}
\end{sphinxVerbatim}


\subsection{Description}
\label{\detokenize{pair_body_nparticle:description}}
\sphinxAtStartPar
Style \sphinxstyleemphasis{body/nparticle} is for use with body particles and calculates
pairwise body/body interactions as well as interactions between body
and point\sphinxhyphen{}particles.  See the {\hyperref[\detokenize{Howto_body::doc}]{\sphinxcrossref{\DUrole{doc}{Howto body}}}} doc page
for more details on using body particles.

\sphinxAtStartPar
This pair style is designed for use with the “nparticle” body style,
which is specified as an argument to the “atom\sphinxhyphen{}style body” command.
See the {\hyperref[\detokenize{Howto_body::doc}]{\sphinxcrossref{\DUrole{doc}{Howto body}}}} page for more details about
the body styles LAMMPS supports.  The “nparticle” style treats a body
particle as a rigid body composed of N sub\sphinxhyphen{}particles.

\sphinxAtStartPar
The coordinates of a body particle are its center\sphinxhyphen{}of\sphinxhyphen{}mass (COM).  If
the COMs of a pair of body particles are within the cutoff (global or
type\sphinxhyphen{}specific, as specified above), then all interactions between
pairs of sub\sphinxhyphen{}particles in the two body particles are computed.
E.g. if the first body particle has 3 sub\sphinxhyphen{}particles, and the second
has 10, then 30 interactions are computed and summed to yield the
total force and torque on each body particle.

\begin{sphinxadmonition}{note}{Note:}
\sphinxAtStartPar
In the example just described, all 30 interactions are computed
even if the distance between a particular pair of sub\sphinxhyphen{}particles is
greater than the cutoff.  Likewise, no interaction between two body
particles is computed if the two COMs are further apart than the
cutoff, even if the distance between some pairs of their sub\sphinxhyphen{}particles
is within the cutoff.  Thus care should be used in defining the cutoff
distances for body particles, depending on their shape and size.
\end{sphinxadmonition}

\sphinxAtStartPar
Similar rules apply for a body particle interacting with a point
particle.  The distance between the two particles is calculated using
the COM of the body particle and the position of the point particle.
If the distance is within the cutoff and the body particle has N
sub\sphinxhyphen{}particles, then N interactions with the point particle are
computed and summed.  If the distance is not within the cutoff, no
interactions between the body and point particle are computed.

\sphinxAtStartPar
The interaction between two sub\sphinxhyphen{}particles, or a sub\sphinxhyphen{}particle and point
particle, or between two point particles is computed as a Lennard\sphinxhyphen{}Jones
interaction, using the standard formula
\begin{equation*}
\begin{split}E & = 4 \epsilon \left[ \left(\frac{\sigma}{r}\right)^{12} -
                    \left(\frac{\sigma}{r}\right)^6 \right]
                    \qquad r < R_c \\\end{split}
\end{equation*}
\sphinxAtStartPar
where \(R_c\) is the cutoff.  As explained above, an interaction involving
one or two body sub\sphinxhyphen{}particles may be computed even for \(r > R_c\).

\sphinxAtStartPar
For style \sphinxstyleemphasis{body}, the following coefficients must be defined for each
pair of atoms types via the {\hyperref[\detokenize{pair_coeff::doc}]{\sphinxcrossref{\DUrole{doc}{pair\_coeff}}}} command as in
the examples above, or in the data file or restart files read by the
{\hyperref[\detokenize{read_data::doc}]{\sphinxcrossref{\DUrole{doc}{read\_data}}}} or {\hyperref[\detokenize{read_restart::doc}]{\sphinxcrossref{\DUrole{doc}{read\_restart}}}}
commands:
\begin{itemize}
\item {} 
\sphinxAtStartPar
\(\epsilon\) (energy units)

\item {} 
\sphinxAtStartPar
\(\sigma\) (distance units)

\item {} 
\sphinxAtStartPar
cutoff (distance units)

\end{itemize}

\sphinxAtStartPar
The last coefficient is optional.  If not specified, the global cutoff
is used.


\bigskip\hrule\bigskip



\subsection{Mixing, shift, table, tail correction, restart, rRESPA info}
\label{\detokenize{pair_body_nparticle:mixing-shift-table-tail-correction-restart-rrespa-info}}
\sphinxAtStartPar
For atom type pairs I,J and I != J, the epsilon and sigma coefficients
and cutoff distance for all of this pair style can be mixed.  The
default mix value is \sphinxstyleemphasis{geometric}.  See the {\hyperref[\detokenize{pair_modify::doc}]{\sphinxcrossref{\DUrole{doc}{pair\_modify}}}} command for
details.

\sphinxAtStartPar
This pair style does not support the {\hyperref[\detokenize{pair_modify::doc}]{\sphinxcrossref{\DUrole{doc}{pair\_modify}}}}
shift, table, and tail options.

\sphinxAtStartPar
This pair style does not write its information to {\hyperref[\detokenize{restart::doc}]{\sphinxcrossref{\DUrole{doc}{binary restart files}}}}.

\sphinxAtStartPar
This pair style can only be used via the \sphinxstyleemphasis{pair} keyword of the
{\hyperref[\detokenize{run_style::doc}]{\sphinxcrossref{\DUrole{doc}{run\_style respa}}}} command.  It does not support the
\sphinxstyleemphasis{inner}, \sphinxstyleemphasis{middle}, \sphinxstyleemphasis{outer} keywords.


\bigskip\hrule\bigskip



\subsection{Restrictions}
\label{\detokenize{pair_body_nparticle:restrictions}}
\sphinxAtStartPar
This style is part of the BODY package.  It is only enabled if LAMMPS
was built with that package.  See the {\hyperref[\detokenize{Build_package::doc}]{\sphinxcrossref{\DUrole{doc}{Build package}}}} page for more info.

\sphinxAtStartPar
Defining particles to be bodies so they participate in body/body or
body/particle interactions requires the use of the {\hyperref[\detokenize{atom_style::doc}]{\sphinxcrossref{\DUrole{doc}{atom\_style body}}}} command.


\subsection{Related commands}
\label{\detokenize{pair_body_nparticle:related-commands}}
\sphinxAtStartPar
{\hyperref[\detokenize{pair_coeff::doc}]{\sphinxcrossref{\DUrole{doc}{pair\_coeff}}}}, {\hyperref[\detokenize{fix_rigid::doc}]{\sphinxcrossref{\DUrole{doc}{fix rigid}}}}


\subsection{Default}
\label{\detokenize{pair_body_nparticle:default}}
\sphinxAtStartPar
none

\sphinxstepscope

\index{pair\_style body/rounded/polygon@\spxentry{pair\_style body/rounded/polygon}}\ignorespaces 

\section{pair\_style body/rounded/polygon command}
\label{\detokenize{pair_body_rounded_polygon:pair-style-body-rounded-polygon-command}}\label{\detokenize{pair_body_rounded_polygon:index-0}}\label{\detokenize{pair_body_rounded_polygon::doc}}

\subsection{Syntax}
\label{\detokenize{pair_body_rounded_polygon:syntax}}
\begin{sphinxVerbatim}[commandchars=\\\{\}]
\PYG{k}{pair\PYGZus{}style}\PYG{+w}{ }\PYG{n}{body}\PYG{o}{/}\PYG{n}{rounded}\PYG{o}{/}\PYG{n}{polygon}\PYG{+w}{ }\PYG{n}{c\PYGZus{}n}\PYG{+w}{ }\PYG{n}{c\PYGZus{}t}\PYG{+w}{ }\PYG{n}{mu}\PYG{+w}{ }\PYG{n}{delta\PYGZus{}ua}\PYG{+w}{ }\PYG{n}{cutoff}
\end{sphinxVerbatim}

\begin{sphinxVerbatim}[commandchars=\\\{\}]
c\PYGZus{}n = normal damping coefficient
c\PYGZus{}t = tangential damping coefficient
mu = normal friction coefficient during gross sliding
delta\PYGZus{}ua = multiple contact scaling factor
cutoff = global separation cutoff for interactions (distance units), see below for definition
\end{sphinxVerbatim}


\subsection{Examples}
\label{\detokenize{pair_body_rounded_polygon:examples}}
\begin{sphinxVerbatim}[commandchars=\\\{\}]
\PYG{k}{pair\PYGZus{}style}\PYG{+w}{ }\PYG{n}{body}\PYG{o}{/}\PYG{n}{rounded}\PYG{o}{/}\PYG{n}{polygon}\PYG{+w}{ }\PYG{l+m}{20.0}\PYG{+w}{ }\PYG{l+m}{5.0}\PYG{+w}{ }\PYG{l+m}{0.0}\PYG{+w}{ }\PYG{l+m}{1.0}\PYG{+w}{ }\PYG{l+m}{0.5}
\PYG{k}{pair\PYGZus{}coeff}\PYG{+w}{ }\PYG{o}{*}\PYG{+w}{ }\PYG{o}{*}\PYG{+w}{ }\PYG{l+m}{100.0}\PYG{+w}{ }\PYG{l+m}{1.0}
\PYG{k}{pair\PYGZus{}coeff}\PYG{+w}{ }\PYG{l+m}{1}\PYG{+w}{ }\PYG{l+m}{1}\PYG{+w}{ }\PYG{l+m}{100.0}\PYG{+w}{ }\PYG{l+m}{1.0}
\end{sphinxVerbatim}


\subsection{Description}
\label{\detokenize{pair_body_rounded_polygon:description}}
\sphinxAtStartPar
Style \sphinxstyleemphasis{body/rounded/polygon} is for use with 2d models of body
particles of style \sphinxstyleemphasis{rounded/polygon}.  It calculates pairwise
body/body interactions which can include body particles modeled as
1\sphinxhyphen{}vertex circular disks with a specified diameter.  See the
{\hyperref[\detokenize{Howto_body::doc}]{\sphinxcrossref{\DUrole{doc}{Howto body}}}} page for more details on using body
rounded/polygon particles.

\sphinxAtStartPar
This pairwise interaction between rounded polygons is described in
{\hyperref[\detokenize{pair_body_rounded_polygon:pair-fraige}]{\sphinxcrossref{\DUrole{std,std-ref}{Fraige}}}}, where a polygon does not have sharp corners,
but is rounded at its vertices by circles centered on each vertex with
a specified diameter.  The edges of the polygon are defined between
pairs of adjacent vertices.  The circle diameter for each polygon is
specified in the data file read by the {\hyperref[\detokenize{read_data::doc}]{\sphinxcrossref{\DUrole{doc}{read data}}}}
command.  This is a 2d discrete element model (DEM) which allows for
multiple contact points.

\sphinxAtStartPar
Note that when two particles interact, the effective surface of each
polygon particle is displaced outward from each of its vertices and
edges by half its circle diameter (as in the diagram below of a gray
and yellow square particle).  The interaction forces and energies
between two particles are defined with respect to the separation of
their respective rounded surfaces, not by the separation of the
vertices and edges themselves.

\sphinxAtStartPar
This means that the specified cutoff in the pair\_style command is the
cutoff distance, \(r_c\), for the surface separation, \(\delta_n\) (see figure
below).  This is the distance at which two particles no longer
interact.  If \(r_c\) is specified as 0.0, then it is a contact\sphinxhyphen{}only
interaction.  I.e. the two particles must overlap in order to exert a
repulsive force on each other.  If \(r_c > 0.0\), then the force between
two particles will be attractive for surface separations from 0 to
\(r_c\), and repulsive once the particles overlap.

\sphinxAtStartPar
Note that unlike for other pair styles, the specified cutoff is not
the distance between the centers of two particles at which they stop
interacting.  This center\sphinxhyphen{}to\sphinxhyphen{}center distance depends on the shape and
size of the two particles and their relative orientation.  LAMMPS
takes that into account when computing the surface separation distance
and applying the \(r_c\) cutoff.

\sphinxAtStartPar
The forces between vertex\sphinxhyphen{}vertex, vertex\sphinxhyphen{}edge, and edge\sphinxhyphen{}edge overlaps
are given by:
\begin{equation*}
\begin{split}F_n &= \begin{cases}
        k_n \delta_n - c_n v_n     &  \delta_n \le 0 \\
       -k_{na} \delta_n - c_n v_n  &  0 < \delta_n \le r_c \\
       0                           & \delta_n > r_c \\
       \end{cases} \\
F_t &= \begin{cases}
       \mu k_n \delta_n - c_t v_t & \delta_n \le 0 \\
       0                          & \delta_n > 0
       \end{cases}\end{split}
\end{equation*}
\noindent{\hspace*{\fill}\sphinxincludegraphics{{pair_body_rounded}.jpg}\hspace*{\fill}}

\sphinxAtStartPar
Note that \(F_n\) and \(F_t\) are functions of the surface separation
\(\delta_n = d - (R_i + R_j)\).  In this model, when
\((R_i + R_j) < d < (R_i + R_j) + r_c\), that is, \(0 < \delta_n < r_c\),
the cohesive region of the two surfaces overlap and the two surfaces are
attractive to each other.

\sphinxAtStartPar
In {\hyperref[\detokenize{pair_body_rounded_polygon:pair-fraige}]{\sphinxcrossref{\DUrole{std,std-ref}{Fraige}}}}, the tangential friction force between two
particles that are in contact is modeled differently prior to gross
sliding (i.e. static friction) and during gross\sphinxhyphen{}sliding (kinetic
friction).  The latter takes place when the tangential deformation
exceeds the Coulomb frictional limit.  In the current implementation,
however, we do not take into account frictional history, i.e. we do
not keep track of how many time steps the two particles have been in
contact nor calculate the tangential deformation.  Instead, we assume
that gross sliding takes place as soon as two particles are in
contact.

\sphinxAtStartPar
The following coefficients must be defined for each pair of atom types
via the {\hyperref[\detokenize{pair_coeff::doc}]{\sphinxcrossref{\DUrole{doc}{pair\_coeff}}}} command as in the examples above,
or in the data file read by the {\hyperref[\detokenize{read_data::doc}]{\sphinxcrossref{\DUrole{doc}{read\_data}}}} command:
\begin{itemize}
\item {} 
\sphinxAtStartPar
\(k_n\) (energy/distance\textasciicircum{}2 units)

\item {} 
\sphinxAtStartPar
\(k_{na}\) (energy/distance\textasciicircum{}2 units)

\end{itemize}

\sphinxAtStartPar
Effectively, \(k_n\) and \(k_{na}\) are the slopes of the red lines in the plot
above for force versus surface separation, for \(\delta_n < 0\) and
\(0 < \delta_n < r_c\) respectively.


\subsection{Mixing, shift, table, tail correction, restart, rRESPA info}
\label{\detokenize{pair_body_rounded_polygon:mixing-shift-table-tail-correction-restart-rrespa-info}}
\sphinxAtStartPar
This pair style does not support the {\hyperref[\detokenize{pair_modify::doc}]{\sphinxcrossref{\DUrole{doc}{pair\_modify}}}}
mix, shift, table, and tail options.

\sphinxAtStartPar
This pair style does not write its information to {\hyperref[\detokenize{restart::doc}]{\sphinxcrossref{\DUrole{doc}{binary restart files}}}}.  Thus, you need to re\sphinxhyphen{}specify the pair\_style and
pair\_coeff commands in an input script that reads a restart file.

\sphinxAtStartPar
This pair style can only be used via the \sphinxstyleemphasis{pair} keyword of the
{\hyperref[\detokenize{run_style::doc}]{\sphinxcrossref{\DUrole{doc}{run\_style respa}}}} command.  It does not support the
\sphinxstyleemphasis{inner}, \sphinxstyleemphasis{middle}, \sphinxstyleemphasis{outer} keywords.


\subsection{Restrictions}
\label{\detokenize{pair_body_rounded_polygon:restrictions}}
\sphinxAtStartPar
These pair styles are part of the BODY package.  They are only enabled
if LAMMPS was built with that package.  See the {\hyperref[\detokenize{Build_package::doc}]{\sphinxcrossref{\DUrole{doc}{Build package}}}} page for more info.

\sphinxAtStartPar
This pair style requires the {\hyperref[\detokenize{newton::doc}]{\sphinxcrossref{\DUrole{doc}{newton}}}} setting to be “on”
for pair interactions.


\subsection{Related commands}
\label{\detokenize{pair_body_rounded_polygon:related-commands}}
\sphinxAtStartPar
{\hyperref[\detokenize{pair_coeff::doc}]{\sphinxcrossref{\DUrole{doc}{pair\_coeff}}}}


\subsection{Default}
\label{\detokenize{pair_body_rounded_polygon:default}}
\sphinxAtStartPar
none

\phantomsection\label{\detokenize{pair_body_rounded_polygon:pair-fraige}}
\sphinxAtStartPar
\sphinxstylestrong{(Fraige)} F. Y. Fraige, P. A. Langston, A. J. Matchett, J. Dodds,
Particuology, 6, 455 (2008).

\sphinxstepscope

\index{pair\_style body/rounded/polyhedron@\spxentry{pair\_style body/rounded/polyhedron}}\ignorespaces 

\section{pair\_style body/rounded/polyhedron command}
\label{\detokenize{pair_body_rounded_polyhedron:pair-style-body-rounded-polyhedron-command}}\label{\detokenize{pair_body_rounded_polyhedron:index-0}}\label{\detokenize{pair_body_rounded_polyhedron::doc}}

\subsection{Syntax}
\label{\detokenize{pair_body_rounded_polyhedron:syntax}}
\begin{sphinxVerbatim}[commandchars=\\\{\}]
\PYG{k}{pair\PYGZus{}style}\PYG{+w}{ }\PYG{n}{body}\PYG{o}{/}\PYG{n}{rounded}\PYG{o}{/}\PYG{n}{polyhedron}\PYG{+w}{ }\PYG{n}{c\PYGZus{}n}\PYG{+w}{ }\PYG{n}{c\PYGZus{}t}\PYG{+w}{ }\PYG{n}{mu}\PYG{+w}{ }\PYG{n}{delta\PYGZus{}ua}\PYG{+w}{ }\PYG{n}{cutoff}
\end{sphinxVerbatim}

\begin{sphinxVerbatim}[commandchars=\\\{\}]
c\PYGZus{}n = normal damping coefficient
c\PYGZus{}t = tangential damping coefficient
mu = normal friction coefficient during gross sliding
delta\PYGZus{}ua = multiple contact scaling factor
cutoff = global separation cutoff for interactions (distance units), see below for definition
\end{sphinxVerbatim}


\subsection{Examples}
\label{\detokenize{pair_body_rounded_polyhedron:examples}}
\begin{sphinxVerbatim}[commandchars=\\\{\}]
\PYG{k}{pair\PYGZus{}style}\PYG{+w}{ }\PYG{n}{body}\PYG{o}{/}\PYG{n}{rounded}\PYG{o}{/}\PYG{n}{polyhedron}\PYG{+w}{ }\PYG{l+m}{20.0}\PYG{+w}{ }\PYG{l+m}{5.0}\PYG{+w}{ }\PYG{l+m}{0.0}\PYG{+w}{ }\PYG{l+m}{1.0}\PYG{+w}{ }\PYG{l+m}{0.5}
\PYG{k}{pair\PYGZus{}coeff}\PYG{+w}{ }\PYG{o}{*}\PYG{+w}{ }\PYG{o}{*}\PYG{+w}{ }\PYG{l+m}{100.0}\PYG{+w}{ }\PYG{l+m}{1.0}
\PYG{k}{pair\PYGZus{}coeff}\PYG{+w}{ }\PYG{l+m}{1}\PYG{+w}{ }\PYG{l+m}{1}\PYG{+w}{ }\PYG{l+m}{100.0}\PYG{+w}{ }\PYG{l+m}{1.0}
\end{sphinxVerbatim}


\subsection{Description}
\label{\detokenize{pair_body_rounded_polyhedron:description}}
\sphinxAtStartPar
Style \sphinxstyleemphasis{body/rounded/polygon} is for use with 3d models of body
particles of style \sphinxstyleemphasis{rounded/polyhedron}.  It calculates pairwise
body/body interactions which can include body particles modeled as
1\sphinxhyphen{}vertex spheres with a specified diameter.  See the
{\hyperref[\detokenize{Howto_body::doc}]{\sphinxcrossref{\DUrole{doc}{Howto body}}}} page for more details on using body
rounded/polyhedron particles.

\sphinxAtStartPar
This pairwise interaction between the rounded polyhedra is described
in {\hyperref[\detokenize{pair_body_rounded_polyhedron:pair-wang}]{\sphinxcrossref{\DUrole{std,std-ref}{Wang}}}}, where a polyhedron does not have sharp corners
and edges, but is rounded at its vertices and edges by spheres
centered on each vertex with a specified diameter.  The edges of the
polyhedron are defined between pairs of adjacent vertices.  Its faces
are defined by a loop of edges.  The sphere diameter for each polygon
is specified in the data file read by the {\hyperref[\detokenize{read_data::doc}]{\sphinxcrossref{\DUrole{doc}{read data}}}}
command.  This is a discrete element model (DEM) which allows for
multiple contact points.

\sphinxAtStartPar
Note that when two particles interact, the effective surface of each
polyhedron particle is displaced outward from each of its vertices,
edges, and faces by half its sphere diameter.  The interaction forces
and energies between two particles are defined with respect to the
separation of their respective rounded surfaces, not by the separation
of the vertices, edges, and faces themselves.

\sphinxAtStartPar
This means that the specified cutoff in the pair\_style command is the
cutoff distance, \(r_c\), for the surface separation, \(\delta_n\) (see figure
below).  This is the distance at which two particles no longer
interact.  If \(r_c\) is specified as 0.0, then it is a contact\sphinxhyphen{}only
interaction.  I.e. the two particles must overlap in order to exert a
repulsive force on each other.  If \(r_c > 0.0\), then the force between
two particles will be attractive for surface separations from 0 to
\(r_c\), and repulsive once the particles overlap.

\sphinxAtStartPar
Note that unlike for other pair styles, the specified cutoff is not
the distance between the centers of two particles at which they stop
interacting.  This center\sphinxhyphen{}to\sphinxhyphen{}center distance depends on the shape and
size of the two particles and their relative orientation.  LAMMPS
takes that into account when computing the surface separation distance
and applying the \(r_c\) cutoff.

\sphinxAtStartPar
The forces between vertex\sphinxhyphen{}vertex, vertex\sphinxhyphen{}edge, vertex\sphinxhyphen{}face, edge\sphinxhyphen{}edge,
and edge\sphinxhyphen{}face overlaps are given by:
\begin{equation*}
\begin{split}F_n &= \begin{cases}
       k_n \delta_n - c_n v_n,    & \delta_n \le 0 \\
      -k_{na} \delta_n - c_n v_n  & 0 < \delta_n \le r_c \\
       0                          & \delta_n > r_c \\
       \end{cases} \\
F_t &= \begin{cases}
       \mu k_n \delta_n - c_t v_t & \delta_n \le 0 \\
       0                          & \delta_n > 0
       \end{cases}\end{split}
\end{equation*}
\noindent{\hspace*{\fill}\sphinxincludegraphics{{pair_body_rounded}.jpg}\hspace*{\fill}}

\sphinxAtStartPar
In {\hyperref[\detokenize{pair_body_rounded_polyhedron:pair-wang}]{\sphinxcrossref{\DUrole{std,std-ref}{Wang}}}}, the tangential friction force between two
particles that are in contact is modeled differently prior to gross
sliding (i.e. static friction) and during gross\sphinxhyphen{}sliding (kinetic
friction).  The latter takes place when the tangential deformation
exceeds the Coulomb frictional limit.  In the current implementation,
however, we do not take into account frictional history, i.e. we do
not keep track of how many time steps the two particles have been in
contact nor calculate the tangential deformation.  Instead, we assume
that gross sliding takes place as soon as two particles are in
contact.

\sphinxAtStartPar
The following coefficients must be defined for each pair of atom types
via the {\hyperref[\detokenize{pair_coeff::doc}]{\sphinxcrossref{\DUrole{doc}{pair\_coeff}}}} command as in the examples above,
or in the data file read by the {\hyperref[\detokenize{read_data::doc}]{\sphinxcrossref{\DUrole{doc}{read\_data}}}} command:
\begin{itemize}
\item {} 
\sphinxAtStartPar
\(k_n\) (energy/distance\textasciicircum{}2 units)

\item {} 
\sphinxAtStartPar
\(k_{na}\) (energy/distance\textasciicircum{}2 units)

\end{itemize}

\sphinxAtStartPar
Effectively, \(k_n\) and \(k_{na}\) are the slopes of the red lines in the plot
above for force versus surface separation, for \(\delta_n\) \textless{} 0 and
\(0 < \delta_n < r_c\) respectively.


\subsection{Mixing, shift, table, tail correction, restart, rRESPA info}
\label{\detokenize{pair_body_rounded_polyhedron:mixing-shift-table-tail-correction-restart-rrespa-info}}
\sphinxAtStartPar
This pair style does not support the {\hyperref[\detokenize{pair_modify::doc}]{\sphinxcrossref{\DUrole{doc}{pair\_modify}}}}
mix, shift, table, and tail options.

\sphinxAtStartPar
This pair style does not write its information to {\hyperref[\detokenize{restart::doc}]{\sphinxcrossref{\DUrole{doc}{binary restart files}}}}.
Thus, you need to re\sphinxhyphen{}specify the pair\_style and pair\_coeff
commands in an input script that reads a restart file.

\sphinxAtStartPar
This pair style can only be used via the \sphinxstyleemphasis{pair} keyword of the
{\hyperref[\detokenize{run_style::doc}]{\sphinxcrossref{\DUrole{doc}{run\_style respa}}}} command.  It does not support the
\sphinxstyleemphasis{inner}, \sphinxstyleemphasis{middle}, \sphinxstyleemphasis{outer} keywords.


\subsection{Restrictions}
\label{\detokenize{pair_body_rounded_polyhedron:restrictions}}
\sphinxAtStartPar
These pair styles are part of the BODY package.  They are only enabled
if LAMMPS was built with that package.  See the {\hyperref[\detokenize{Build_package::doc}]{\sphinxcrossref{\DUrole{doc}{Build package}}}} page for more info.

\sphinxAtStartPar
This pair style requires the {\hyperref[\detokenize{newton::doc}]{\sphinxcrossref{\DUrole{doc}{newton}}}} setting to be “on”
for pair interactions.


\subsection{Related commands}
\label{\detokenize{pair_body_rounded_polyhedron:related-commands}}
\sphinxAtStartPar
{\hyperref[\detokenize{pair_coeff::doc}]{\sphinxcrossref{\DUrole{doc}{pair\_coeff}}}}


\subsection{Default}
\label{\detokenize{pair_body_rounded_polyhedron:default}}
\sphinxAtStartPar
none

\phantomsection\label{\detokenize{pair_body_rounded_polyhedron:pair-wang}}
\sphinxAtStartPar
\sphinxstylestrong{(Wang)} J. Wang, H. S. Yu, P. A. Langston, F. Y. Fraige, Granular
Matter, 13, 1 (2011).

\sphinxstepscope

\index{pair\_style bop@\spxentry{pair\_style bop}}\ignorespaces 

\section{pair\_style bop command}
\label{\detokenize{pair_bop:pair-style-bop-command}}\label{\detokenize{pair_bop:index-0}}\label{\detokenize{pair_bop::doc}}

\subsection{Syntax}
\label{\detokenize{pair_bop:syntax}}
\begin{sphinxVerbatim}[commandchars=\\\{\}]
\PYG{k}{pair\PYGZus{}style}\PYG{+w}{ }\PYG{n}{bop}\PYG{+w}{ }\PYG{n}{keyword}\PYG{+w}{ }\PYG{n}{...}
\end{sphinxVerbatim}
\begin{itemize}
\item {} 
\sphinxAtStartPar
zero or more keywords may be appended

\item {} 
\sphinxAtStartPar
keyword = \sphinxstyleemphasis{save}

\end{itemize}

\begin{sphinxVerbatim}[commandchars=\\\{\}]
save = pre\PYGZhy{}compute and save some values
\end{sphinxVerbatim}


\subsection{Examples}
\label{\detokenize{pair_bop:examples}}
\begin{sphinxVerbatim}[commandchars=\\\{\}]
\PYG{k}{pair\PYGZus{}style}\PYG{+w}{ }\PYG{n}{bop}
\PYG{k}{pair\PYGZus{}coeff}\PYG{+w}{ }\PYG{o}{*}\PYG{+w}{ }\PYG{o}{*}\PYG{+w}{ }\PYG{n}{..}\PYG{o}{/}\PYG{n}{potentials}\PYG{o}{/}\PYG{n}{CdTe\PYGZus{}bop}\PYG{+w}{ }\PYG{n}{Cd}\PYG{+w}{ }\PYG{n}{Te}
\PYG{k}{pair\PYGZus{}style}\PYG{+w}{ }\PYG{n}{bop}\PYG{+w}{ }\PYG{n}{save}
\PYG{k}{pair\PYGZus{}coeff}\PYG{+w}{ }\PYG{o}{*}\PYG{+w}{ }\PYG{o}{*}\PYG{+w}{ }\PYG{n}{..}\PYG{o}{/}\PYG{n}{potentials}\PYG{o}{/}\PYG{n}{CdTe.bop.table}\PYG{+w}{ }\PYG{n}{Cd}\PYG{+w}{ }\PYG{n}{Te}\PYG{+w}{ }\PYG{n}{Te}
\PYG{k}{comm\PYGZus{}modify}\PYG{+w}{ }\PYG{n}{cutoff}\PYG{+w}{ }\PYG{l+m}{14.70}
\end{sphinxVerbatim}


\subsection{Description}
\label{\detokenize{pair_bop:description}}
\sphinxAtStartPar
The \sphinxstyleemphasis{bop} pair style computes Bond\sphinxhyphen{}Order Potentials (BOP) based on
quantum mechanical theory incorporating both \(\sigma\) and \(\pi\) bonding.
By analytically deriving the BOP from quantum mechanical theory its
transferability to different phases can approach that of quantum
mechanical methods.  This potential is similar to the original BOP
developed by Pettifor ({\hyperref[\detokenize{pair_bop:pettifor-1}]{\sphinxcrossref{\DUrole{std,std-ref}{Pettifor\_1}}}},
{\hyperref[\detokenize{pair_bop:pettifor-2}]{\sphinxcrossref{\DUrole{std,std-ref}{Pettifor\_2}}}}, {\hyperref[\detokenize{pair_bop:pettifor-3}]{\sphinxcrossref{\DUrole{std,std-ref}{Pettifor\_3}}}}) and later updated
by Murdick, Zhou, and Ward ({\hyperref[\detokenize{pair_bop:murdick}]{\sphinxcrossref{\DUrole{std,std-ref}{Murdick}}}}, {\hyperref[\detokenize{pair_bop:ward}]{\sphinxcrossref{\DUrole{std,std-ref}{Ward}}}}).
Currently, BOP potential files for these systems are provided with
LAMMPS: AlCu, CCu, CdTe, CdTeSe, CdZnTe, CuH, GaAs.  A system with
only a subset of these elements, including a single element (e.g. C or
Cu or Al or Ga or Zn or CdZn), can also be modeled by using the
appropriate alloy file and assigning all atom types to the
single element or subset of elements via the pair\_coeff command, as
discussed below.

\sphinxAtStartPar
The BOP potential consists of three terms:
\begin{equation*}
\begin{split}E = \frac{1}{2} \sum_{i=1}^{N} \sum_{j=i_1}^{i_N} \phi_{ij} \left( r_{ij} \right) - \sum_{i=1}^{N} \sum_{j=i_1}^{i_N} \beta_{\sigma,ij} \left( r_{ij} \right) \cdot \Theta_{\sigma,ij} - \sum_{i=1}^{N} \sum_{j=i_1}^{i_N} \beta_{\pi,ij} \left( r_{ij} \right) \cdot \Theta_{\pi,ij} + U_{prom}\end{split}
\end{equation*}
\sphinxAtStartPar
where \(\phi_{ij}(r_{ij})\) is a short\sphinxhyphen{}range two\sphinxhyphen{}body function
representing the repulsion between a pair of ion cores,
\(\beta_{\sigma,ij}(r_{ij})\) and \(\beta_{\sigma,ij}(r_{ij})\)
are respectively sigma and \(\pi\) bond integrals, \(\Theta_{\sigma,ij}\)
and \(\Theta_{\pi,ij}\) are \(\sigma\) and \(\pi\)
bond\sphinxhyphen{}orders, and U\_prom is the promotion energy for sp\sphinxhyphen{}valent systems.

\sphinxAtStartPar
The detailed formulas for this potential are given in Ward
({\hyperref[\detokenize{pair_bop:ward}]{\sphinxcrossref{\DUrole{std,std-ref}{Ward}}}}); here we provide only a brief description.

\sphinxAtStartPar
The repulsive energy \(\phi_{ij}(r_{ij})\) and the bond integrals
\(\beta_{\sigma,ij}(r_{ij})\) and \(\beta_{\phi,ij}(r_{ij})\) are functions of the
interatomic distance \(r_{ij}\) between atom \sphinxstyleemphasis{i} and \sphinxstyleemphasis{j}.  Each of these
potentials has a smooth cutoff at a radius of \(r_{cut,ij}\).  These
smooth cutoffs ensure stable behavior at situations with high sampling
near the cutoff such as melts and surfaces.

\sphinxAtStartPar
The bond\sphinxhyphen{}orders can be viewed as environment\sphinxhyphen{}dependent local variables
that are ij bond specific.  The maximum value of the \(\sigma\)
bond\sphinxhyphen{}order (\(\Theta_{\sigma}\) is 1, while that of the \(\pi\)
bond\sphinxhyphen{}order (\(\Theta_{\pi}\)) is 2, attributing to a maximum value
of the total bond\sphinxhyphen{}order (\(\Theta_{\sigma}+\Theta_{\pi}\)) of 3.
The \(\sigma\) and \(\pi\) bond\sphinxhyphen{}orders reflect the ubiquitous
single\sphinxhyphen{}, double\sphinxhyphen{}, and triple\sphinxhyphen{} bond behavior of chemistry. Their
analytical expressions can be derived from tight\sphinxhyphen{} binding theory by
recursively expanding an inter\sphinxhyphen{}site Green’s function as a continued
fraction. To accurately represent the bonding with a computationally
efficient potential formulation suitable for MD simulations, the derived
BOP only takes (and retains) the first two levels of the recursive
representations for both the \(\sigma\) and the \(\pi\) bond\sphinxhyphen{}orders. Bond\sphinxhyphen{}order
terms can be understood in terms of molecular orbital hopping paths
based upon the Cyrot\sphinxhyphen{}Lackmann theorem ({\hyperref[\detokenize{pair_bop:pettifor-1}]{\sphinxcrossref{\DUrole{std,std-ref}{Pettifor\_1}}}}).
The \(\sigma\) bond\sphinxhyphen{}order with a half\sphinxhyphen{}full valence shell is used to
interpolate the bond\sphinxhyphen{}order expression that incorporated explicit valance
band filling.  This \(\pi\) bond\sphinxhyphen{}order expression also contains also contains
a three\sphinxhyphen{}member ring term that allows implementation of an asymmetric
density of states, which helps to either stabilize or destabilize
close\sphinxhyphen{}packed structures.  The \(\pi\) bond\sphinxhyphen{}order includes hopping paths of
length 4.  This enables the incorporation of dihedral angles effects.

\begin{sphinxadmonition}{note}{Note:}
\sphinxAtStartPar
Note that unlike for other potentials, cutoffs for BOP
potentials are not set in the pair\_style or pair\_coeff command; they
are specified in the BOP potential files themselves.  Likewise, the
BOP potential files list atomic masses; thus you do not need to use
the {\hyperref[\detokenize{mass::doc}]{\sphinxcrossref{\DUrole{doc}{mass}}}} command to specify them.  Note that for BOP
potentials with hydrogen, you will likely want to set the mass of H
atoms to be 10x or 20x larger to avoid having to use a tiny timestep.
You can do this by using the {\hyperref[\detokenize{mass::doc}]{\sphinxcrossref{\DUrole{doc}{mass}}}} command after using the
{\hyperref[\detokenize{pair_coeff::doc}]{\sphinxcrossref{\DUrole{doc}{pair\_coeff}}}} command to read the BOP potential
file.
\end{sphinxadmonition}

\sphinxAtStartPar
One option can be specified as a keyword with the pair\_style command.

\sphinxAtStartPar
The \sphinxstyleemphasis{save} keyword gives you the option to calculate in advance and
store a set of distances, angles, and derivatives of angles.  The
default is to not do this, but to calculate them on\sphinxhyphen{}the\sphinxhyphen{}fly each time
they are needed.  The former may be faster, but takes more memory.
The latter requires less memory, but may be slower.  It is best to
test this option to optimize the speed of BOP for your particular
system configuration.


\bigskip\hrule\bigskip


\sphinxAtStartPar
Only a single pair\_coeff command is used with the \sphinxstyleemphasis{bop} style which
specifies a BOP potential file, with parameters for all needed
elements.  These are mapped to LAMMPS atom types by specifying
N additional arguments after the filename in the pair\_coeff command,
where N is the number of LAMMPS atom types:
\begin{itemize}
\item {} 
\sphinxAtStartPar
filename

\item {} 
\sphinxAtStartPar
N element names = mapping of BOP elements to atom types

\end{itemize}

\sphinxAtStartPar
As an example, imagine the CdTe.bop file has BOP values for Cd
and Te.  If your LAMMPS simulation has 4 atoms types and you want the
first 3 to be Cd, and the fourth to be Te, you would use the following
pair\_coeff command:

\begin{sphinxVerbatim}[commandchars=\\\{\}]
\PYG{k}{pair\PYGZus{}coeff}\PYG{+w}{ }\PYG{o}{*}\PYG{+w}{ }\PYG{o}{*}\PYG{+w}{ }\PYG{n}{CdTe}\PYG{+w}{ }\PYG{n}{Cd}\PYG{+w}{ }\PYG{n}{Cd}\PYG{+w}{ }\PYG{n}{Cd}\PYG{+w}{ }\PYG{n}{Te}
\end{sphinxVerbatim}

\sphinxAtStartPar
The first 2 arguments must be * * so as to span all LAMMPS atom types.
The first three Cd arguments map LAMMPS atom types 1,2,3 to the Cd
element in the BOP file.  The final Te argument maps LAMMPS atom type
4 to the Te element in the BOP file.

\sphinxAtStartPar
BOP files in the \sphinxstyleemphasis{potentials} directory of the LAMMPS distribution
have a “.bop” suffix.  The potentials are in tabulated form containing
pre\sphinxhyphen{}tabulated pair functions for phi\_ij(r\_ij), beta\_(sigma,ij)(r\_ij),
and beta\_pi,ij)(r\_ij).

\sphinxAtStartPar
The parameters/coefficients format for the different kinds of BOP
files are given below with variables matching the formulation of Ward
({\hyperref[\detokenize{pair_bop:ward}]{\sphinxcrossref{\DUrole{std,std-ref}{Ward}}}}) and Zhou ({\hyperref[\detokenize{pair_bop:zhou1}]{\sphinxcrossref{\DUrole{std,std-ref}{Zhou}}}}). Each header line containing a
“:” is preceded by a blank line.


\bigskip\hrule\bigskip


\sphinxAtStartPar
\sphinxstylestrong{No angular table file format}:

\sphinxAtStartPar
The parameters/coefficients format for the BOP potentials input file
containing pre\sphinxhyphen{}tabulated functions of g is given below with variables
matching the formulation of Ward ({\hyperref[\detokenize{pair_bop:ward}]{\sphinxcrossref{\DUrole{std,std-ref}{Ward}}}}).  This format also
assumes the angular functions have the formulation of ({\hyperref[\detokenize{pair_bop:ward}]{\sphinxcrossref{\DUrole{std,std-ref}{Ward}}}}).
\begin{itemize}
\item {} 
\sphinxAtStartPar
Line 1: \# elements N

\end{itemize}

\sphinxAtStartPar
The first line is followed by N lines containing the atomic
number, mass, and element symbol of each element.

\sphinxAtStartPar
Following the definition of the elements several global variables for
the tabulated functions are given.
\begin{itemize}
\item {} 
\sphinxAtStartPar
Line 1: nr, nBOt (nr is the number of divisions the radius is broken
into for function tables and MUST be a factor of 5; nBOt is the number
of divisions for the tabulated values of THETA\_(S,ij)

\item {} 
\sphinxAtStartPar
Line 2: delta\_1\sphinxhyphen{}delta\_7 (if all are not used in the particular

\item {} 
\sphinxAtStartPar
formulation, set unused values to 0.0)

\end{itemize}

\sphinxAtStartPar
Following this N lines for e\_1\sphinxhyphen{}e\_N containing p\_pi.
\begin{itemize}
\item {} 
\sphinxAtStartPar
Line 3: p\_pi (for e\_1)

\item {} 
\sphinxAtStartPar
Line 4: p\_pi (for e\_2 and continues to e\_N)

\end{itemize}

\sphinxAtStartPar
The next section contains several pair constants for the number of
interaction types e\_i\sphinxhyphen{}e\_j, with i=1\sphinxhyphen{}\textgreater{}N, j=i\sphinxhyphen{}\textgreater{}N
\begin{itemize}
\item {} 
\sphinxAtStartPar
Line 1: r\_cut (for e\_1\sphinxhyphen{}e\_1 interactions)

\item {} 
\sphinxAtStartPar
Line 2: c\_sigma, a\_sigma, c\_pi, a\_pi

\item {} 
\sphinxAtStartPar
Line 3: delta\_sigma, delta\_pi

\item {} 
\sphinxAtStartPar
Line 4: f\_sigma, k\_sigma, delta\_3 (This delta\_3 is similar to that of
the previous section but is interaction type dependent)

\end{itemize}

\sphinxAtStartPar
The next section contains a line for each three body interaction type
e\_j\sphinxhyphen{}e\_i\sphinxhyphen{}e\_k with i=0\sphinxhyphen{}\textgreater{}N, j=0\sphinxhyphen{}\textgreater{}N, k=j\sphinxhyphen{}\textgreater{}N
\begin{itemize}
\item {} 
\sphinxAtStartPar
Line 1: g\_(sigma0), g\_(sigma1), g\_(sigma2) (These are coefficients for
g\_(sigma,jik)(THETA\_ijk) for e\_1\sphinxhyphen{}e\_1\sphinxhyphen{}e\_1 interaction. {\hyperref[\detokenize{pair_bop:ward}]{\sphinxcrossref{\DUrole{std,std-ref}{Ward}}}}
contains the full expressions for the constants as functions of
b\_(sigma,ijk), p\_(sigma,ijk), u\_(sigma,ijk))

\item {} 
\sphinxAtStartPar
Line 2: g\_(sigma0), g\_(sigma1), g\_(sigma2) (for e\_1\sphinxhyphen{}e\_1\sphinxhyphen{}e\_2)

\end{itemize}

\sphinxAtStartPar
The next section contains a block for each interaction type for the
phi\_ij(r\_ij).  Each block has nr entries with 5 entries per line.
\begin{itemize}
\item {} 
\sphinxAtStartPar
Line 1: phi(r1), phi(r2), phi(r3), phi(r4), phi(r5) (for the e\_1\sphinxhyphen{}e\_1
interaction type)

\item {} 
\sphinxAtStartPar
Line 2: phi(r6), phi(r7), phi(r8), phi(r9), phi(r10) (this continues
until nr)

\item {} 
\sphinxAtStartPar
…

\item {} 
\sphinxAtStartPar
Line nr/5\_1: phi(r1), phi(r2), phi(r3), phi(r4), phi(r5), (for the
e\_1\sphinxhyphen{}e\_1 interaction type)

\end{itemize}

\sphinxAtStartPar
The next section contains a block for each interaction type for the
beta\_(sigma,ij)(r\_ij).  Each block has nr entries with 5 entries per
line.
\begin{itemize}
\item {} 
\sphinxAtStartPar
Line 1: beta\_sigma(r1), beta\_sigma(r2), beta\_sigma(r3), beta\_sigma(r4),
beta\_sigma(r5) (for the e\_1\sphinxhyphen{}e\_1 interaction type)

\item {} 
\sphinxAtStartPar
Line 2: beta\_sigma(r6), beta\_sigma(r7), beta\_sigma(r8), beta\_sigma(r9),
beta\_sigma(r10) (this continues until nr)

\item {} 
\sphinxAtStartPar
…

\item {} 
\sphinxAtStartPar
Line nr/5+1: beta\_sigma(r1), beta\_sigma(r2), beta\_sigma(r3),
beta\_sigma(r4), beta\_sigma(r5) (for the e\_1\sphinxhyphen{}e\_2 interaction type)

\end{itemize}

\sphinxAtStartPar
The next section contains a block for each interaction type for
beta\_(pi,ij)(r\_ij).  Each block has nr entries with 5 entries per line.
\begin{itemize}
\item {} 
\sphinxAtStartPar
Line 1: beta\_pi(r1), beta\_pi(r2), beta\_pi(r3), beta\_pi(r4), beta\_pi(r5)
(for the e\_1\sphinxhyphen{}e\_1 interaction type)

\item {} 
\sphinxAtStartPar
Line 2: beta\_pi(r6), beta\_pi(r7), beta\_pi(r8), beta\_pi(r9),
beta\_pi(r10) (this continues until nr)

\item {} 
\sphinxAtStartPar
…

\item {} 
\sphinxAtStartPar
Line nr/5+1: beta\_pi(r1), beta\_pi(r2), beta\_pi(r3), beta\_pi(r4),
beta\_pi(r5) (for the e\_1\sphinxhyphen{}e\_2 interaction type)

\end{itemize}

\sphinxAtStartPar
The next section contains a block for each interaction type for the
THETA\_(S,ij)((THETA\_(sigma,ij))\textasciicircum{}(1/2), f\_(sigma,ij)).  Each block has
nBOt entries with 5 entries per line.
\begin{itemize}
\item {} 
\sphinxAtStartPar
Line 1: THETA\_(S,ij)(r1), THETA\_(S,ij)(r2), THETA\_(S,ij)(r3),
THETA\_(S,ij)(r4), THETA\_(S,ij)(r5) (for the e\_1\sphinxhyphen{}e\_2 interaction type)

\item {} 
\sphinxAtStartPar
Line 2: THETA\_(S,ij)(r6), THETA\_(S,ij)(r7), THETA\_(S,ij)(r8),
THETA\_(S,ij)(r9), THETA\_(S,ij)(r10) (this continues until nBOt)

\item {} 
\sphinxAtStartPar
…

\item {} 
\sphinxAtStartPar
Line nBOt/5+1: THETA\_(S,ij)(r1), THETA\_(S,ij)(r2), THETA\_(S,ij)(r3),
THETA\_(S,ij)(r4), THETA\_(S,ij)(r5) (for the e\_1\sphinxhyphen{}e\_2 interaction type)

\end{itemize}

\sphinxAtStartPar
The next section contains a block of N lines for e\_1\sphinxhyphen{}e\_N
\begin{itemize}
\item {} 
\sphinxAtStartPar
Line 1: delta\textasciicircum{}mu (for e\_1)

\item {} 
\sphinxAtStartPar
Line 2: delta\textasciicircum{}mu (for e\_2 and repeats to e\_N)

\end{itemize}

\sphinxAtStartPar
The last section contains more constants for e\_i\sphinxhyphen{}e\_j interactions with
i=0\sphinxhyphen{}\textgreater{}N, j=i\sphinxhyphen{}\textgreater{}N
\begin{itemize}
\item {} 
\sphinxAtStartPar
Line 1: (A\_ij)\textasciicircum{}(mu*nu) (for e1\sphinxhyphen{}e1)

\item {} 
\sphinxAtStartPar
Line 2: (A\_ij)\textasciicircum{}(mu*nu) (for e1\sphinxhyphen{}e2 and repeats as above)

\end{itemize}


\bigskip\hrule\bigskip


\sphinxAtStartPar
\sphinxstylestrong{Angular spline table file format}:

\sphinxAtStartPar
The parameters/coefficients format for the BOP potentials input file
containing pre\sphinxhyphen{}tabulated functions of g is given below with variables
matching the formulation of Ward ({\hyperref[\detokenize{pair_bop:ward}]{\sphinxcrossref{\DUrole{std,std-ref}{Ward}}}}).  This format also
assumes the angular functions have the formulation of ({\hyperref[\detokenize{pair_bop:zhou1}]{\sphinxcrossref{\DUrole{std,std-ref}{Zhou}}}}).
\begin{itemize}
\item {} 
\sphinxAtStartPar
Line 1: \# elements N

\end{itemize}

\sphinxAtStartPar
The first line is followed by N lines containing the atomic
number, mass, and element symbol of each element.

\sphinxAtStartPar
Following the definition of the elements several global variables for
the tabulated functions are given.
\begin{itemize}
\item {} 
\sphinxAtStartPar
Line 1: nr, ntheta, nBOt (nr is the number of divisions the radius is broken
into for function tables and MUST be a factor of 5; ntheta is the power of the
power of the spline used to fit the angular function; nBOt is the number
of divisions for the tabulated values of THETA\_(S,ij)

\item {} 
\sphinxAtStartPar
Line 2: delta\_1\sphinxhyphen{}delta\_7 (if all are not used in the particular

\item {} 
\sphinxAtStartPar
formulation, set unused values to 0.0)

\end{itemize}

\sphinxAtStartPar
Following this N lines for e\_1\sphinxhyphen{}e\_N containing p\_pi.
\begin{itemize}
\item {} 
\sphinxAtStartPar
Line 3: p\_pi (for e\_1)

\item {} 
\sphinxAtStartPar
Line 4: p\_pi (for e\_2 and continues to e\_N)

\end{itemize}

\sphinxAtStartPar
The next section contains several pair constants for the number of
interaction types e\_i\sphinxhyphen{}e\_j, with i=1\sphinxhyphen{}\textgreater{}N, j=i\sphinxhyphen{}\textgreater{}N
\begin{itemize}
\item {} 
\sphinxAtStartPar
Line 1: r\_cut (for e\_1\sphinxhyphen{}e\_1 interactions)

\item {} 
\sphinxAtStartPar
Line 2: c\_sigma, a\_sigma, c\_pi, a\_pi

\item {} 
\sphinxAtStartPar
Line 3: delta\_sigma, delta\_pi

\item {} 
\sphinxAtStartPar
Line 4: f\_sigma, k\_sigma, delta\_3 (This delta\_3 is similar to that of
the previous section but is interaction type dependent)

\end{itemize}

\sphinxAtStartPar
The next section contains a line for each three body interaction type
e\_j\sphinxhyphen{}e\_i\sphinxhyphen{}e\_k with i=0\sphinxhyphen{}\textgreater{}N, j=0\sphinxhyphen{}\textgreater{}N, k=j\sphinxhyphen{}\textgreater{}N
\begin{itemize}
\item {} 
\sphinxAtStartPar
Line 1: g0, g1, g2… (These are coefficients for the angular spline
of the g\_(sigma,jik)(THETA\_ijk) for e\_1\sphinxhyphen{}e\_1\sphinxhyphen{}e\_1 interaction.  The
function can contain up to 10 term thus 10 constants.  The first line
can contain up to five constants.  If the spline has more than five
terms the second line will contain the remaining constants The
following lines will then contain the constants for the remaining g0,
g1, g2… (for e\_1\sphinxhyphen{}e\_1\sphinxhyphen{}e\_2) and the other three body
interactions

\end{itemize}

\sphinxAtStartPar
The rest of the table has the same structure as the previous section
(see above).


\bigskip\hrule\bigskip


\sphinxAtStartPar
\sphinxstylestrong{Angular no\sphinxhyphen{}spline table file format}:

\sphinxAtStartPar
The parameters/coefficients format for the BOP potentials input file
containing pre\sphinxhyphen{}tabulated functions of g is given below with variables
matching the formulation of Ward ({\hyperref[\detokenize{pair_bop:ward}]{\sphinxcrossref{\DUrole{std,std-ref}{Ward}}}}).  This format also
assumes the angular functions have the formulation of ({\hyperref[\detokenize{pair_bop:zhou1}]{\sphinxcrossref{\DUrole{std,std-ref}{Zhou}}}}).
\begin{itemize}
\item {} 
\sphinxAtStartPar
Line 1: \# elements N

\end{itemize}

\sphinxAtStartPar
The first two lines are followed by N lines containing the atomic
number, mass, and element symbol of each element.

\sphinxAtStartPar
Following the definition of the elements several global variables for
the tabulated functions are given.
\begin{itemize}
\item {} 
\sphinxAtStartPar
Line 1: nr, ntheta, nBOt (nr is the number of divisions the radius is broken
into for function tables and MUST be a factor of 5; ntheta is the number of
divisions for the tabulated values of the g angular function; nBOt is the number
of divisions for the tabulated values of THETA\_(S,ij)

\item {} 
\sphinxAtStartPar
Line 2: delta\_1\sphinxhyphen{}delta\_7 (if all are not used in the particular

\item {} 
\sphinxAtStartPar
formulation, set unused values to 0.0)

\end{itemize}

\sphinxAtStartPar
Following this N lines for e\_1\sphinxhyphen{}e\_N containing p\_pi.
\begin{itemize}
\item {} 
\sphinxAtStartPar
Line 3: p\_pi (for e\_1)

\item {} 
\sphinxAtStartPar
Line 4: p\_pi (for e\_2 and continues to e\_N)

\end{itemize}

\sphinxAtStartPar
The next section contains several pair constants for the number of
interaction types e\_i\sphinxhyphen{}e\_j, with i=1\sphinxhyphen{}\textgreater{}N, j=i\sphinxhyphen{}\textgreater{}N
\begin{itemize}
\item {} 
\sphinxAtStartPar
Line 1: r\_cut (for e\_1\sphinxhyphen{}e\_1 interactions)

\item {} 
\sphinxAtStartPar
Line 2: c\_sigma, a\_sigma, c\_pi, a\_pi

\item {} 
\sphinxAtStartPar
Line 3: delta\_sigma, delta\_pi

\item {} 
\sphinxAtStartPar
Line 4: f\_sigma, k\_sigma, delta\_3 (This delta\_3 is similar to that of
the previous section but is interaction type dependent)

\end{itemize}

\sphinxAtStartPar
The next section contains a line for each three body interaction type
e\_j\sphinxhyphen{}e\_i\sphinxhyphen{}e\_k with i=0\sphinxhyphen{}\textgreater{}N, j=0\sphinxhyphen{}\textgreater{}N, k=j\sphinxhyphen{}\textgreater{}N
\begin{itemize}
\item {} 
\sphinxAtStartPar
Line 1: g(theta1), g(theta2), g(theta3), g(theta4), g(theta5) (for the e\_1\sphinxhyphen{}e\_1\sphinxhyphen{}e\_1
interaction type)

\item {} 
\sphinxAtStartPar
Line 2: g(theta6), g(theta7), g(theta8), g(theta9), g(theta10) (this continues
until ntheta)

\item {} 
\sphinxAtStartPar
…

\item {} 
\sphinxAtStartPar
Line ntheta/5+1: g(theta1), g(theta2), g(theta3), g(theta4), g(theta5), (for the
e\_1\sphinxhyphen{}e\_1\sphinxhyphen{}e\_2 interaction type)

\end{itemize}

\sphinxAtStartPar
The rest of the table has the same structure as the previous section (see above).


\bigskip\hrule\bigskip



\subsection{Mixing, shift, table, tail correction, restart, rRESPA info}
\label{\detokenize{pair_bop:mixing-shift-table-tail-correction-restart-rrespa-info}}
\sphinxAtStartPar
This pair style does not support the {\hyperref[\detokenize{pair_modify::doc}]{\sphinxcrossref{\DUrole{doc}{pair\_modify}}}}
mix, shift, table, and tail options.

\sphinxAtStartPar
This pair style does not write its information to {\hyperref[\detokenize{restart::doc}]{\sphinxcrossref{\DUrole{doc}{binary restart files}}}}, since it is stored in potential files.  Thus, you
need to re\sphinxhyphen{}specify the pair\_style and pair\_coeff commands in an input
script that reads a restart file.

\sphinxAtStartPar
This pair style can only be used via the \sphinxstyleemphasis{pair} keyword of the
{\hyperref[\detokenize{run_style::doc}]{\sphinxcrossref{\DUrole{doc}{run\_style respa}}}} command.  It does not support the
\sphinxstyleemphasis{inner}, \sphinxstyleemphasis{middle}, \sphinxstyleemphasis{outer} keywords.


\bigskip\hrule\bigskip



\subsection{Restrictions}
\label{\detokenize{pair_bop:restrictions}}
\sphinxAtStartPar
These pair styles are part of the MANYBODY package.  They are only
enabled if LAMMPS was built with that package.  See the {\hyperref[\detokenize{Build_package::doc}]{\sphinxcrossref{\DUrole{doc}{Build package}}}} page for more info.

\sphinxAtStartPar
These pair potentials require the {\hyperref[\detokenize{newton::doc}]{\sphinxcrossref{\DUrole{doc}{newtion}}}} setting to be
“on” for pair interactions.

\sphinxAtStartPar
Pair style bop is not compatible with being used as a sub\sphinxhyphen{}style with
doc:\sphinxtitleref{hybrid pair styles \textless{}pair\_hybrid\textgreater{}}. Pair style bop is also not
compatible with {\hyperref[\detokenize{neighbor::doc}]{\sphinxcrossref{\DUrole{doc}{multi\sphinxhyphen{}cutoff neighbor lists}}}} or
{\hyperref[\detokenize{comm_modify::doc}]{\sphinxcrossref{\DUrole{doc}{multi\sphinxhyphen{}cutoff communitcation}}}}.

\sphinxAtStartPar
The .bop.table potential files provided with LAMMPS (see the
potentials directory) are parameterized for metal {\hyperref[\detokenize{units::doc}]{\sphinxcrossref{\DUrole{doc}{units}}}}.
You can use the BOP potential with any LAMMPS units, but you would need
to create your own BOP potential file with coefficients listed in the
appropriate units if your simulation does not use “metal” units.


\subsection{Related commands}
\label{\detokenize{pair_bop:related-commands}}
\sphinxAtStartPar
{\hyperref[\detokenize{pair_coeff::doc}]{\sphinxcrossref{\DUrole{doc}{pair\_coeff}}}}


\subsection{Default}
\label{\detokenize{pair_bop:default}}
\sphinxAtStartPar
non\sphinxhyphen{}tabulated potential file, a\_0 is non\sphinxhyphen{}zero.


\bigskip\hrule\bigskip

\phantomsection\label{\detokenize{pair_bop:pettifor-1}}
\sphinxAtStartPar
\sphinxstylestrong{(Pettifor\_1)} D.G. Pettifor and I.I. Oleinik, Phys. Rev. B, 59, 8487
(1999).

\phantomsection\label{\detokenize{pair_bop:pettifor-2}}
\sphinxAtStartPar
\sphinxstylestrong{(Pettifor\_2)} D.G. Pettifor and I.I. Oleinik, Phys. Rev. Lett., 84,
4124 (2000).

\phantomsection\label{\detokenize{pair_bop:pettifor-3}}
\sphinxAtStartPar
\sphinxstylestrong{(Pettifor\_3)} D.G. Pettifor and I.I. Oleinik, Phys. Rev. B, 65, 172103
(2002).

\phantomsection\label{\detokenize{pair_bop:murdick}}
\sphinxAtStartPar
\sphinxstylestrong{(Murdick)} D.A. Murdick, X.W. Zhou, H.N.G. Wadley, D. Nguyen\sphinxhyphen{}Manh, R.
Drautz, and D.G. Pettifor, Phys. Rev. B, 73, 45206 (2006).

\phantomsection\label{\detokenize{pair_bop:ward}}
\sphinxAtStartPar
\sphinxstylestrong{(Ward)} D.K. Ward, X.W. Zhou, B.M. Wong, F.P. Doty, and J.A.
Zimmerman, Phys. Rev. B, 85,115206 (2012).

\phantomsection\label{\detokenize{pair_bop:zhou1}}
\sphinxAtStartPar
\sphinxstylestrong{(Zhou)} X.W. Zhou, D.K. Ward, M. Foster (TBP).

\sphinxstepscope

\index{pair\_style born@\spxentry{pair\_style born}}\ignorespaces 
\index{pair\_style born/omp@\spxentry{pair\_style born/omp}}\ignorespaces 
\index{pair\_style born/gpu@\spxentry{pair\_style born/gpu}}\ignorespaces 
\index{pair\_style born/coul/long@\spxentry{pair\_style born/coul/long}}\ignorespaces 
\index{pair\_style born/coul/long/gpu@\spxentry{pair\_style born/coul/long/gpu}}\ignorespaces 
\index{pair\_style born/coul/long/omp@\spxentry{pair\_style born/coul/long/omp}}\ignorespaces 
\index{pair\_style born/coul/msm@\spxentry{pair\_style born/coul/msm}}\ignorespaces 
\index{pair\_style born/coul/msm/omp@\spxentry{pair\_style born/coul/msm/omp}}\ignorespaces 
\index{pair\_style born/coul/wolf@\spxentry{pair\_style born/coul/wolf}}\ignorespaces 
\index{pair\_style born/coul/wolf/gpu@\spxentry{pair\_style born/coul/wolf/gpu}}\ignorespaces 
\index{pair\_style born/coul/wolf/omp@\spxentry{pair\_style born/coul/wolf/omp}}\ignorespaces 
\index{pair\_style born/coul/dsf@\spxentry{pair\_style born/coul/dsf}}\ignorespaces 

\section{pair\_style born command}
\label{\detokenize{pair_born:pair-style-born-command}}\label{\detokenize{pair_born:index-11}}\label{\detokenize{pair_born:index-10}}\label{\detokenize{pair_born:index-9}}\label{\detokenize{pair_born:index-8}}\label{\detokenize{pair_born:index-7}}\label{\detokenize{pair_born:index-6}}\label{\detokenize{pair_born:index-5}}\label{\detokenize{pair_born:index-4}}\label{\detokenize{pair_born:index-3}}\label{\detokenize{pair_born:index-2}}\label{\detokenize{pair_born:index-1}}\label{\detokenize{pair_born:index-0}}\label{\detokenize{pair_born::doc}}
\sphinxAtStartPar
Accelerator Variants: \sphinxstyleemphasis{born/omp}, \sphinxstyleemphasis{born/gpu}


\section{pair\_style born/coul/long command}
\label{\detokenize{pair_born:pair-style-born-coul-long-command}}
\sphinxAtStartPar
Accelerator Variants: \sphinxstyleemphasis{born/coul/long/gpu}, \sphinxstyleemphasis{born/coul/long/omp}


\section{pair\_style born/coul/msm command}
\label{\detokenize{pair_born:pair-style-born-coul-msm-command}}
\sphinxAtStartPar
Accelerator Variants: \sphinxstyleemphasis{born/coul/msm/omp}


\section{pair\_style born/coul/wolf command}
\label{\detokenize{pair_born:pair-style-born-coul-wolf-command}}
\sphinxAtStartPar
Accelerator Variants: \sphinxstyleemphasis{born/coul/wolf/gpu}, \sphinxstyleemphasis{born/coul/wolf/omp}


\section{pair\_style born/coul/dsf command}
\label{\detokenize{pair_born:pair-style-born-coul-dsf-command}}

\subsection{Syntax}
\label{\detokenize{pair_born:syntax}}
\begin{sphinxVerbatim}[commandchars=\\\{\}]
\PYG{k}{pair\PYGZus{}style}\PYG{+w}{ }\PYG{n}{style}\PYG{+w}{ }\PYG{n}{args}
\end{sphinxVerbatim}
\begin{itemize}
\item {} 
\sphinxAtStartPar
style = \sphinxstyleemphasis{born} or \sphinxstyleemphasis{born/coul/long} or \sphinxstyleemphasis{born/coul/msm} or \sphinxstyleemphasis{born/coul/wolf}

\item {} 
\sphinxAtStartPar
args = list of arguments for a particular style

\end{itemize}
\begin{sphinxalltt}
\sphinxstyleemphasis{born} args = cutoff
  cutoff = global cutoff for non\sphinxhyphen{}Coulombic interactions (distance units)
\sphinxstyleemphasis{born/coul/long} args = cutoff (cutoff2)
  cutoff = global cutoff for non\sphinxhyphen{}Coulombic (and Coulombic if only 1 arg) (distance units)
  cutoff2 = global cutoff for Coulombic (optional) (distance units)
\sphinxstyleemphasis{born/coul/msm} args = cutoff (cutoff2)
  cutoff = global cutoff for non\sphinxhyphen{}Coulombic (and Coulombic if only 1 arg) (distance units)
  cutoff2 = global cutoff for Coulombic (optional) (distance units)
\sphinxstyleemphasis{born/coul/wolf} args = alpha cutoff (cutoff2)
  alpha = damping parameter (inverse distance units)
  cutoff = global cutoff for non\sphinxhyphen{}Coulombic (and Coulombic if only 1 arg) (distance units)
  cutoff2 = global cutoff for Coulombic (optional) (distance units)
\sphinxstyleemphasis{born/coul/dsf} args = alpha cutoff (cutoff2)
  alpha = damping parameter (inverse distance units)
  cutoff = global cutoff for non\sphinxhyphen{}Coulombic (and Coulombic if only 1 arg) (distance units)
  cutoff2 = global cutoff for Coulombic (distance units)
\end{sphinxalltt}


\subsection{Examples}
\label{\detokenize{pair_born:examples}}
\begin{sphinxVerbatim}[commandchars=\\\{\}]
\PYG{k}{pair\PYGZus{}style}\PYG{+w}{ }\PYG{n}{born}\PYG{+w}{ }\PYG{l+m}{10.0}
\PYG{k}{pair\PYGZus{}coeff}\PYG{+w}{ }\PYG{o}{*}\PYG{+w}{ }\PYG{o}{*}\PYG{+w}{ }\PYG{l+m}{6.08}\PYG{+w}{ }\PYG{l+m}{0.317}\PYG{+w}{ }\PYG{l+m}{2.340}\PYG{+w}{ }\PYG{l+m}{24.18}\PYG{+w}{ }\PYG{l+m}{11.51}
\PYG{k}{pair\PYGZus{}coeff}\PYG{+w}{ }\PYG{l+m}{1}\PYG{+w}{ }\PYG{l+m}{1}\PYG{+w}{ }\PYG{l+m}{6.08}\PYG{+w}{ }\PYG{l+m}{0.317}\PYG{+w}{ }\PYG{l+m}{2.340}\PYG{+w}{ }\PYG{l+m}{24.18}\PYG{+w}{ }\PYG{l+m}{11.51}

\PYG{k}{pair\PYGZus{}style}\PYG{+w}{ }\PYG{n}{born}\PYG{o}{/}\PYG{n}{coul}\PYG{o}{/}\PYG{n}{long}\PYG{+w}{ }\PYG{l+m}{10.0}
\PYG{k}{pair\PYGZus{}style}\PYG{+w}{ }\PYG{n}{born}\PYG{o}{/}\PYG{n}{coul}\PYG{o}{/}\PYG{n}{long}\PYG{+w}{ }\PYG{l+m}{10.0}\PYG{+w}{ }\PYG{l+m}{8}\PYG{n}{.}
\PYG{k}{pair\PYGZus{}coeff}\PYG{+w}{ }\PYG{o}{*}\PYG{+w}{ }\PYG{o}{*}\PYG{+w}{ }\PYG{l+m}{6.08}\PYG{+w}{ }\PYG{l+m}{0.317}\PYG{+w}{ }\PYG{l+m}{2.340}\PYG{+w}{ }\PYG{l+m}{24.18}\PYG{+w}{ }\PYG{l+m}{11.51}
\PYG{k}{pair\PYGZus{}coeff}\PYG{+w}{ }\PYG{l+m}{1}\PYG{+w}{ }\PYG{l+m}{1}\PYG{+w}{ }\PYG{l+m}{6.08}\PYG{+w}{ }\PYG{l+m}{0.317}\PYG{+w}{ }\PYG{l+m}{2.340}\PYG{+w}{ }\PYG{l+m}{24.18}\PYG{+w}{ }\PYG{l+m}{11.51}

\PYG{k}{pair\PYGZus{}style}\PYG{+w}{ }\PYG{n}{born}\PYG{o}{/}\PYG{n}{coul}\PYG{o}{/}\PYG{n}{msm}\PYG{+w}{ }\PYG{l+m}{10.0}
\PYG{k}{pair\PYGZus{}style}\PYG{+w}{ }\PYG{n}{born}\PYG{o}{/}\PYG{n}{coul}\PYG{o}{/}\PYG{n}{msm}\PYG{+w}{ }\PYG{l+m}{10.0}\PYG{+w}{ }\PYG{l+m}{8.0}
\PYG{k}{pair\PYGZus{}coeff}\PYG{+w}{ }\PYG{o}{*}\PYG{+w}{ }\PYG{o}{*}\PYG{+w}{ }\PYG{l+m}{6.08}\PYG{+w}{ }\PYG{l+m}{0.317}\PYG{+w}{ }\PYG{l+m}{2.340}\PYG{+w}{ }\PYG{l+m}{24.18}\PYG{+w}{ }\PYG{l+m}{11.51}
\PYG{k}{pair\PYGZus{}coeff}\PYG{+w}{ }\PYG{l+m}{1}\PYG{+w}{ }\PYG{l+m}{1}\PYG{+w}{ }\PYG{l+m}{6.08}\PYG{+w}{ }\PYG{l+m}{0.317}\PYG{+w}{ }\PYG{l+m}{2.340}\PYG{+w}{ }\PYG{l+m}{24.18}\PYG{+w}{ }\PYG{l+m}{11.51}

\PYG{k}{pair\PYGZus{}style}\PYG{+w}{ }\PYG{n}{born}\PYG{o}{/}\PYG{n}{coul}\PYG{o}{/}\PYG{n}{wolf}\PYG{+w}{ }\PYG{l+m}{0.25}\PYG{+w}{ }\PYG{l+m}{10.0}
\PYG{k}{pair\PYGZus{}style}\PYG{+w}{ }\PYG{n}{born}\PYG{o}{/}\PYG{n}{coul}\PYG{o}{/}\PYG{n}{wolf}\PYG{+w}{ }\PYG{l+m}{0.25}\PYG{+w}{ }\PYG{l+m}{10.0}\PYG{+w}{ }\PYG{l+m}{9.0}
\PYG{k}{pair\PYGZus{}coeff}\PYG{+w}{ }\PYG{o}{*}\PYG{+w}{ }\PYG{o}{*}\PYG{+w}{ }\PYG{l+m}{6.08}\PYG{+w}{ }\PYG{l+m}{0.317}\PYG{+w}{ }\PYG{l+m}{2.340}\PYG{+w}{ }\PYG{l+m}{24.18}\PYG{+w}{ }\PYG{l+m}{11.51}
\PYG{k}{pair\PYGZus{}coeff}\PYG{+w}{ }\PYG{l+m}{1}\PYG{+w}{ }\PYG{l+m}{1}\PYG{+w}{ }\PYG{l+m}{6.08}\PYG{+w}{ }\PYG{l+m}{0.317}\PYG{+w}{ }\PYG{l+m}{2.340}\PYG{+w}{ }\PYG{l+m}{24.18}\PYG{+w}{ }\PYG{l+m}{11.51}

\PYG{k}{pair\PYGZus{}style}\PYG{+w}{ }\PYG{n}{born}\PYG{o}{/}\PYG{n}{coul}\PYG{o}{/}\PYG{n}{dsf}\PYG{+w}{ }\PYG{l+m}{0.1}\PYG{+w}{ }\PYG{l+m}{10.0}\PYG{+w}{ }\PYG{l+m}{12.0}
\PYG{k}{pair\PYGZus{}coeff}\PYG{+w}{ }\PYG{o}{*}\PYG{+w}{ }\PYG{o}{*}\PYG{+w}{   }\PYG{l+m}{0.0}\PYG{+w}{ }\PYG{l+m}{1.00}\PYG{+w}{ }\PYG{l+m}{0.00}\PYG{+w}{ }\PYG{l+m}{0.00}\PYG{+w}{ }\PYG{l+m}{0.00}
\PYG{k}{pair\PYGZus{}coeff}\PYG{+w}{ }\PYG{l+m}{1}\PYG{+w}{ }\PYG{l+m}{1}\PYG{+w}{ }\PYG{l+m}{480.0}\PYG{+w}{ }\PYG{l+m}{0.25}\PYG{+w}{ }\PYG{l+m}{0.00}\PYG{+w}{ }\PYG{l+m}{1.05}\PYG{+w}{ }\PYG{l+m}{0.50}
\end{sphinxVerbatim}


\subsection{Description}
\label{\detokenize{pair_born:description}}
\sphinxAtStartPar
The \sphinxstyleemphasis{born} style computes the Born\sphinxhyphen{}Mayer\sphinxhyphen{}Huggins or Tosi/Fumi
potential described in {\hyperref[\detokenize{pair_born:fumitosi}]{\sphinxcrossref{\DUrole{std,std-ref}{(Fumi and Tosi)}}}}, given by
\begin{equation*}
\begin{split}E = A \exp \left(\frac{\sigma - r}{\rho} \right) -
\frac{C}{r^6} + \frac{D}{r^8} \qquad r < r_c\end{split}
\end{equation*}
\sphinxAtStartPar
where \(\sigma\) is an interaction\sphinxhyphen{}dependent length parameter,
\(\rho\) is an ionic\sphinxhyphen{}pair dependent length parameter, and
\(r_c\) is the cutoff.

\sphinxAtStartPar
The styles with \sphinxstyleemphasis{coul/long} or \sphinxstyleemphasis{coul/msm} add a Coulombic term as
described for the {\hyperref[\detokenize{pair_lj::doc}]{\sphinxcrossref{\DUrole{doc}{lj/cut}}}} pair styles.  An additional
damping factor is applied to the Coulombic term so it can be used in
conjunction with the {\hyperref[\detokenize{kspace_style::doc}]{\sphinxcrossref{\DUrole{doc}{kspace\_style}}}} command and its
\sphinxstyleemphasis{ewald} or \sphinxstyleemphasis{pppm} of \sphinxstyleemphasis{msm} option.  The Coulombic cutoff specified for
this style means that pairwise interactions within this distance are
computed directly; interactions outside that distance are computed in
reciprocal space.

\sphinxAtStartPar
If one cutoff is specified for the \sphinxstyleemphasis{born/coul/long} and
\sphinxstyleemphasis{born/coul/msm} style, it is used for both the A,C,D and Coulombic
terms.  If two cutoffs are specified, the first is used as the cutoff
for the A,C,D terms, and the second is the cutoff for the Coulombic
term.

\sphinxAtStartPar
The \sphinxstyleemphasis{born/coul/wolf} style adds a Coulombic term as described for the
Wolf potential in the {\hyperref[\detokenize{pair_coul::doc}]{\sphinxcrossref{\DUrole{doc}{coul/wolf}}}} pair style.

\sphinxAtStartPar
The \sphinxstyleemphasis{born/coul/dsf} style computes the Coulomb contribution with the
damped shifted force model as in the {\hyperref[\detokenize{pair_coul::doc}]{\sphinxcrossref{\DUrole{doc}{coul/dsf}}}} style.

\sphinxAtStartPar
Note that these potentials are related to the {\hyperref[\detokenize{pair_buck::doc}]{\sphinxcrossref{\DUrole{doc}{Buckingham potential}}}}.

\sphinxAtStartPar
The following coefficients must be defined for each pair of atoms
types via the {\hyperref[\detokenize{pair_coeff::doc}]{\sphinxcrossref{\DUrole{doc}{pair\_coeff}}}} command as in the examples
above, or in the data file or restart files read by the
{\hyperref[\detokenize{read_data::doc}]{\sphinxcrossref{\DUrole{doc}{read\_data}}}} or {\hyperref[\detokenize{read_restart::doc}]{\sphinxcrossref{\DUrole{doc}{read\_restart}}}}
commands, or by mixing as described below:
\begin{itemize}
\item {} 
\sphinxAtStartPar
A (energy units)

\item {} 
\sphinxAtStartPar
\(\rho\) (distance units)

\item {} 
\sphinxAtStartPar
\(\sigma\) (distance units)

\item {} 
\sphinxAtStartPar
C (energy units * distance units\textasciicircum{}6)

\item {} 
\sphinxAtStartPar
D (energy units * distance units\textasciicircum{}8)

\item {} 
\sphinxAtStartPar
cutoff (distance units)

\end{itemize}

\sphinxAtStartPar
The second coefficient, rho, must be greater than zero.

\sphinxAtStartPar
The last coefficient is optional.  If not specified, the global A,C,D
cutoff specified in the pair\_style command is used.

\sphinxAtStartPar
For \sphinxstyleemphasis{born/coul/long}, \sphinxstyleemphasis{born/coul/wolf} and \sphinxstyleemphasis{born/coul/dsf} no
Coulombic cutoff can be specified for an individual I,J type pair.
All type pairs use the same global Coulombic cutoff specified in the
pair\_style command.


\bigskip\hrule\bigskip


\sphinxAtStartPar
Styles with a \sphinxstyleemphasis{gpu}, \sphinxstyleemphasis{intel}, \sphinxstyleemphasis{kk}, \sphinxstyleemphasis{omp}, or \sphinxstyleemphasis{opt} suffix are
functionally the same as the corresponding style without the suffix.
They have been optimized to run faster, depending on your available
hardware, as discussed on the {\hyperref[\detokenize{Speed_packages::doc}]{\sphinxcrossref{\DUrole{doc}{Accelerator packages}}}}
page.  The accelerated styles take the same arguments and should
produce the same results, except for round\sphinxhyphen{}off and precision issues.

\sphinxAtStartPar
These accelerated styles are part of the GPU, INTEL, KOKKOS,
OPENMP, and OPT packages, respectively.  They are only enabled if
LAMMPS was built with those packages.  See the {\hyperref[\detokenize{Build_package::doc}]{\sphinxcrossref{\DUrole{doc}{Build package}}}} page for more info.

\sphinxAtStartPar
You can specify the accelerated styles explicitly in your input script
by including their suffix, or you can use the {\hyperref[\detokenize{Run_options::doc}]{\sphinxcrossref{\DUrole{doc}{\sphinxhyphen{}suffix command\sphinxhyphen{}line switch}}}} when you invoke LAMMPS, or you can use the
{\hyperref[\detokenize{suffix::doc}]{\sphinxcrossref{\DUrole{doc}{suffix}}}} command in your input script.

\sphinxAtStartPar
See the {\hyperref[\detokenize{Speed_packages::doc}]{\sphinxcrossref{\DUrole{doc}{Accelerator packages}}}} page for more
instructions on how to use the accelerated styles effectively.


\bigskip\hrule\bigskip



\subsection{Mixing, shift, table, tail correction, restart, rRESPA info}
\label{\detokenize{pair_born:mixing-shift-table-tail-correction-restart-rrespa-info}}
\sphinxAtStartPar
These pair styles do not support mixing.  Thus, coefficients for all
I,J pairs must be specified explicitly.

\sphinxAtStartPar
These styles support the {\hyperref[\detokenize{pair_modify::doc}]{\sphinxcrossref{\DUrole{doc}{pair\_modify}}}} shift option
for the energy of the exp(), 1/r\textasciicircum{}6, and 1/r\textasciicircum{}8 portion of the pair
interaction.

\sphinxAtStartPar
The \sphinxstyleemphasis{born/coul/long} pair style supports the
{\hyperref[\detokenize{pair_modify::doc}]{\sphinxcrossref{\DUrole{doc}{pair\_modify}}}} table option to tabulate the
short\sphinxhyphen{}range portion of the long\sphinxhyphen{}range Coulombic interaction.

\sphinxAtStartPar
These styles support the pair\_modify tail option for adding long\sphinxhyphen{}range
tail corrections to energy and pressure.

\sphinxAtStartPar
Thess styles writes thei information to binary {\hyperref[\detokenize{restart::doc}]{\sphinxcrossref{\DUrole{doc}{restart}}}}
files, so pair\_style and pair\_coeff commands do not need to be
specified in an input script that reads a restart file.

\sphinxAtStartPar
These styles can only be used via the \sphinxstyleemphasis{pair} keyword of the {\hyperref[\detokenize{run_style::doc}]{\sphinxcrossref{\DUrole{doc}{run\_style respa}}}} command.  They do not support the \sphinxstyleemphasis{inner},
\sphinxstyleemphasis{middle}, \sphinxstyleemphasis{outer} keywords.


\bigskip\hrule\bigskip



\subsection{Restrictions}
\label{\detokenize{pair_born:restrictions}}
\sphinxAtStartPar
The \sphinxstyleemphasis{born/coul/long} style is part of the KSPACE package.  It is only
enabled if LAMMPS was built with that package.  See the
{\hyperref[\detokenize{Build_package::doc}]{\sphinxcrossref{\DUrole{doc}{Build package}}}} page for more info.

\sphinxAtStartPar
The \sphinxstyleemphasis{born/coul/dsf} and \sphinxstyleemphasis{born/coul/wolf} pair styles are part of the
EXTRA\sphinxhyphen{}PAIR package.  They are only enabled if LAMMPS was built with
that package.  See the {\hyperref[\detokenize{Build_package::doc}]{\sphinxcrossref{\DUrole{doc}{Build package}}}} page
for more info.


\subsection{Related commands}
\label{\detokenize{pair_born:related-commands}}
\sphinxAtStartPar
{\hyperref[\detokenize{pair_coeff::doc}]{\sphinxcrossref{\DUrole{doc}{pair\_coeff}}}}, {\hyperref[\detokenize{pair_buck::doc}]{\sphinxcrossref{\DUrole{doc}{pair\_style buck}}}}


\subsection{Default}
\label{\detokenize{pair_born:default}}
\sphinxAtStartPar
none


\bigskip\hrule\bigskip

\phantomsection\label{\detokenize{pair_born:fumitosi}}
\sphinxAtStartPar
Fumi and Tosi, J Phys Chem Solids, 25, 31 (1964),
Fumi and Tosi, J Phys Chem Solids, 25, 45 (1964).

\sphinxstepscope

\index{pair\_style born/gauss@\spxentry{pair\_style born/gauss}}\ignorespaces 

\section{pair\_style born/gauss command}
\label{\detokenize{pair_born_gauss:pair-style-born-gauss-command}}\label{\detokenize{pair_born_gauss:index-0}}\label{\detokenize{pair_born_gauss::doc}}

\subsection{Syntax}
\label{\detokenize{pair_born_gauss:syntax}}
\begin{sphinxVerbatim}[commandchars=\\\{\}]
\PYG{k}{pair\PYGZus{}style}\PYG{+w}{ }\PYG{n}{born}\PYG{o}{/}\PYG{n}{gauss}\PYG{+w}{ }\PYG{n}{cutoff}
\end{sphinxVerbatim}
\begin{itemize}
\item {} 
\sphinxAtStartPar
born/gauss = name of the pair style

\item {} 
\sphinxAtStartPar
cutoff = global cutoff (distance units)

\end{itemize}


\subsection{Examples}
\label{\detokenize{pair_born_gauss:examples}}
\begin{sphinxVerbatim}[commandchars=\\\{\}]
\PYG{k}{pair\PYGZus{}style}\PYG{+w}{ }\PYG{n}{born}\PYG{o}{/}\PYG{n}{gauss}\PYG{+w}{ }\PYG{l+m}{10.0}
\PYG{k}{pair\PYGZus{}coeff}\PYG{+w}{ }\PYG{l+m}{1}\PYG{+w}{ }\PYG{l+m}{1}\PYG{+w}{ }\PYG{l+m}{8.2464e13}\PYG{+w}{ }\PYG{l+m}{12.48}\PYG{+w}{ }\PYG{l+m}{0.042644277}\PYG{+w}{ }\PYG{l+m}{0.44}\PYG{+w}{ }\PYG{l+m}{3.56}
\end{sphinxVerbatim}


\subsection{Description}
\label{\detokenize{pair_born_gauss:description}}
\sphinxAtStartPar
\DUrole{versionmodified,added}{Added in version 28Mar2023.}

\sphinxAtStartPar
Pair style \sphinxstyleemphasis{born/gauss} computes pairwise interactions from a combination of a Born\sphinxhyphen{}Mayer
repulsive term and a Gaussian attractive term according to {\hyperref[\detokenize{pair_born_gauss:bomont}]{\sphinxcrossref{\DUrole{std,std-ref}{(Bomont)}}}}:
\begin{equation*}
\begin{split}E = A_0 \exp \left( -\alpha r \right) - A_1 \exp\left[ -\beta \left(r - r_0 \right)^2 \right]
    \qquad r < r_c\end{split}
\end{equation*}
\sphinxAtStartPar
\(r_c\) is the cutoff.

\sphinxAtStartPar
The following coefficients must be defined for each pair of atoms
types via the {\hyperref[\detokenize{pair_coeff::doc}]{\sphinxcrossref{\DUrole{doc}{pair\_coeff}}}} command as in the examples
above, or in the data file or restart files read by the
{\hyperref[\detokenize{read_data::doc}]{\sphinxcrossref{\DUrole{doc}{read\_data}}}} or {\hyperref[\detokenize{read_restart::doc}]{\sphinxcrossref{\DUrole{doc}{read\_restart}}}}
commands:
\begin{itemize}
\item {} 
\sphinxAtStartPar
\(A_0\) (energy units)

\item {} 
\sphinxAtStartPar
\(\alpha\) (1/distance units)

\item {} 
\sphinxAtStartPar
\(A_1\) (energy units)

\item {} 
\sphinxAtStartPar
\(\beta\) (1/(distance units)\textasciicircum{}2)

\item {} 
\sphinxAtStartPar
\(r_0\) (distance units)

\item {} 
\sphinxAtStartPar
cutoff (distance units)

\end{itemize}

\sphinxAtStartPar
The last coefficient is optional.  If not specified, the global cutoff is used.


\bigskip\hrule\bigskip



\subsection{Mixing, shift, table, tail correction, restart, rRESPA info}
\label{\detokenize{pair_born_gauss:mixing-shift-table-tail-correction-restart-rrespa-info}}
\sphinxAtStartPar
This pair style does not support mixing.  Thus, coefficients for all I,J
pairs must be specified explicitly.

\sphinxAtStartPar
This pair style supports the {\hyperref[\detokenize{pair_modify::doc}]{\sphinxcrossref{\DUrole{doc}{pair\_modify}}}} shift
option for the energy of the pair interaction.

\sphinxAtStartPar
The {\hyperref[\detokenize{pair_modify::doc}]{\sphinxcrossref{\DUrole{doc}{pair\_modify}}}} table options are not relevant for
this pair style.

\sphinxAtStartPar
This pair style does not support the {\hyperref[\detokenize{pair_modify::doc}]{\sphinxcrossref{\DUrole{doc}{pair\_modify}}}}
tail option for adding long\sphinxhyphen{}range tail corrections to energy and
pressure.

\sphinxAtStartPar
This pair style writes its information to {\hyperref[\detokenize{restart::doc}]{\sphinxcrossref{\DUrole{doc}{binary restart files}}}}, so pair\_style and pair\_coeff commands do not need to be
specified in an input script that reads a restart file.

\sphinxAtStartPar
This pair style can only be used via the \sphinxstyleemphasis{pair} keyword of the
{\hyperref[\detokenize{run_style::doc}]{\sphinxcrossref{\DUrole{doc}{run\_style respa}}}} command.  It does not support the
\sphinxstyleemphasis{inner}, \sphinxstyleemphasis{middle}, \sphinxstyleemphasis{outer} keywords.


\bigskip\hrule\bigskip



\subsection{Restrictions}
\label{\detokenize{pair_born_gauss:restrictions}}
\sphinxAtStartPar
This pair style is only enabled if LAMMPS was built with the EXTRA\sphinxhyphen{}PAIR
package.  See the {\hyperref[\detokenize{Build_package::doc}]{\sphinxcrossref{\DUrole{doc}{Build package}}}} page for more
info.


\subsection{Related commands}
\label{\detokenize{pair_born_gauss:related-commands}}
\sphinxAtStartPar
{\hyperref[\detokenize{pair_coeff::doc}]{\sphinxcrossref{\DUrole{doc}{pair\_coeff}}}}, {\hyperref[\detokenize{pair_born::doc}]{\sphinxcrossref{\DUrole{doc}{pair\_style born}}}}


\subsection{Default}
\label{\detokenize{pair_born_gauss:default}}
\sphinxAtStartPar
none


\bigskip\hrule\bigskip

\phantomsection\label{\detokenize{pair_born_gauss:bomont}}
\sphinxAtStartPar
\sphinxstylestrong{(Bomont)} Bomont, Bretonnet, J. Chem. Phys. 124, 054504 (2006)

\sphinxstepscope

\index{pair\_style bpm/spring@\spxentry{pair\_style bpm/spring}}\ignorespaces 

\section{pair\_style bpm/spring command}
\label{\detokenize{pair_bpm_spring:pair-style-bpm-spring-command}}\label{\detokenize{pair_bpm_spring:index-0}}\label{\detokenize{pair_bpm_spring::doc}}

\subsection{Syntax}
\label{\detokenize{pair_bpm_spring:syntax}}
\begin{sphinxVerbatim}[commandchars=\\\{\}]
\PYG{k}{pair\PYGZus{}style}\PYG{+w}{ }\PYG{n}{bpm}\PYG{o}{/}\PYG{n}{spring}
\end{sphinxVerbatim}


\subsection{Examples}
\label{\detokenize{pair_bpm_spring:examples}}
\begin{sphinxVerbatim}[commandchars=\\\{\}]
\PYG{k}{pair\PYGZus{}style}\PYG{+w}{ }\PYG{n}{bpm}\PYG{o}{/}\PYG{n}{spring}
\PYG{k}{pair\PYGZus{}coeff}\PYG{+w}{ }\PYG{o}{*}\PYG{+w}{ }\PYG{o}{*}\PYG{+w}{ }\PYG{l+m}{1.0}\PYG{+w}{ }\PYG{l+m}{1.0}\PYG{+w}{ }\PYG{l+m}{1.0}
\PYG{k}{pair\PYGZus{}coeff}\PYG{+w}{ }\PYG{l+m}{1}\PYG{+w}{ }\PYG{l+m}{1}\PYG{+w}{ }\PYG{l+m}{1.0}\PYG{+w}{ }\PYG{l+m}{1.0}\PYG{+w}{ }\PYG{l+m}{1.0}
\end{sphinxVerbatim}


\subsection{Description}
\label{\detokenize{pair_bpm_spring:description}}
\sphinxAtStartPar
\DUrole{versionmodified,added}{Added in version 4May2022.}

\sphinxAtStartPar
Style \sphinxstyleemphasis{bpm/spring} computes pairwise forces with the formula
\begin{equation*}
\begin{split}F = k (r - r_c)\end{split}
\end{equation*}
\sphinxAtStartPar
where \(k\) is a stiffness and \(r_c\) is the cutoff length.
An additional damping force is also applied to interacting
particles. The force is proportional to the difference in the
normal velocity of particles
\begin{equation*}
\begin{split}F_D = - \gamma w (\hat{r} \bullet \vec{v})\end{split}
\end{equation*}
\sphinxAtStartPar
where \(\gamma\) is the damping strength, \(\hat{r}\) is the
radial normal vector, \(\vec{v}\) is the velocity difference
between the two particles, and \(w\) is a smoothing factor.
This smoothing factor is constructed such that damping forces go to zero
as particles come out of contact to avoid discontinuities. It is
given by
\begin{equation*}
\begin{split}w = 1.0 - \left( \frac{r}{r_c} \right)^8 .\end{split}
\end{equation*}
\sphinxAtStartPar
This pair style is designed for use in a spring\sphinxhyphen{}based bonded particle
model.  It mirrors the construction of the {\hyperref[\detokenize{bond_bpm_spring::doc}]{\sphinxcrossref{\DUrole{doc}{bpm/spring}}}} bond style.

\sphinxAtStartPar
This pair interaction is always applied to pairs of non\sphinxhyphen{}bonded particles
that are within the interaction distance. For pairs of bonded particles
that are within the interaction distance, there is the option to either
include this pair interaction and overlay the pair force over the bond
force or to exclude this pair interaction such that the two particles
only interact via the bond force. See discussion of the \sphinxstyleemphasis{overlay/pair}
option for BPM bond styles and the {\hyperref[\detokenize{special_bonds::doc}]{\sphinxcrossref{\DUrole{doc}{special\_bonds}}}}
command in the {\hyperref[\detokenize{Howto_bpm::doc}]{\sphinxcrossref{\DUrole{doc}{how to}}}} page on BPMs for more details.

\sphinxAtStartPar
The following coefficients must be defined for each pair of atom types
via the {\hyperref[\detokenize{pair_coeff::doc}]{\sphinxcrossref{\DUrole{doc}{pair\_coeff}}}} command as in the examples
above, or in the data file or restart files read by the
{\hyperref[\detokenize{read_data::doc}]{\sphinxcrossref{\DUrole{doc}{read\_data}}}} or {\hyperref[\detokenize{read_restart::doc}]{\sphinxcrossref{\DUrole{doc}{read\_restart}}}}
commands, or by mixing as described below:
\begin{itemize}
\item {} 
\sphinxAtStartPar
\(k\)             (force/distance units)

\item {} 
\sphinxAtStartPar
\(r_c\)           (distance units)

\item {} 
\sphinxAtStartPar
\(\gamma\)        (force/velocity units)

\end{itemize}


\bigskip\hrule\bigskip



\subsection{Mixing, shift, table, tail correction, restart, rRESPA info}
\label{\detokenize{pair_bpm_spring:mixing-shift-table-tail-correction-restart-rrespa-info}}
\sphinxAtStartPar
For atom type pairs I,J and I != J, the A coefficient and cutoff
distance for this pair style can be mixed.  A is always mixed via a
\sphinxstyleemphasis{geometric} rule.  The cutoff is mixed according to the pair\_modify
mix value.  The default mix value is \sphinxstyleemphasis{geometric}.  See the
“pair\_modify” command for details.

\sphinxAtStartPar
This pair style does not support the {\hyperref[\detokenize{pair_modify::doc}]{\sphinxcrossref{\DUrole{doc}{pair\_modify}}}}
shift option, since the pair interaction goes to 0.0 at the cutoff.

\sphinxAtStartPar
The {\hyperref[\detokenize{pair_modify::doc}]{\sphinxcrossref{\DUrole{doc}{pair\_modify}}}} table and tail options are not
relevant for this pair style.

\sphinxAtStartPar
This pair style writes its information to {\hyperref[\detokenize{restart::doc}]{\sphinxcrossref{\DUrole{doc}{binary restart files}}}}, so pair\_style and pair\_coeff commands do not need to be
specified in an input script that reads a restart file.

\sphinxAtStartPar
This pair style can only be used via the \sphinxstyleemphasis{pair} keyword of the
{\hyperref[\detokenize{run_style::doc}]{\sphinxcrossref{\DUrole{doc}{run\_style respa}}}} command.  It does not support the
\sphinxstyleemphasis{inner}, \sphinxstyleemphasis{middle}, \sphinxstyleemphasis{outer} keywords.


\bigskip\hrule\bigskip



\subsection{Restrictions}
\label{\detokenize{pair_bpm_spring:restrictions}}
\sphinxAtStartPar
This pair style is part of the BPM package.  It is only enabled if
LAMMPS was built with that package.  See the {\hyperref[\detokenize{Build_package::doc}]{\sphinxcrossref{\DUrole{doc}{Build package}}}} page for more info.


\subsection{Related commands}
\label{\detokenize{pair_bpm_spring:related-commands}}
\sphinxAtStartPar
{\hyperref[\detokenize{pair_coeff::doc}]{\sphinxcrossref{\DUrole{doc}{pair\_coeff}}}}, {\hyperref[\detokenize{bond_bpm_spring::doc}]{\sphinxcrossref{\DUrole{doc}{bond bpm/spring}}}}


\subsection{Default}
\label{\detokenize{pair_bpm_spring:default}}
\sphinxAtStartPar
none

\sphinxstepscope

\index{pair\_style brownian@\spxentry{pair\_style brownian}}\ignorespaces 
\index{pair\_style brownian/omp@\spxentry{pair\_style brownian/omp}}\ignorespaces 
\index{pair\_style brownian/poly@\spxentry{pair\_style brownian/poly}}\ignorespaces 
\index{pair\_style brownian/poly/omp@\spxentry{pair\_style brownian/poly/omp}}\ignorespaces 

\section{pair\_style brownian command}
\label{\detokenize{pair_brownian:pair-style-brownian-command}}\label{\detokenize{pair_brownian:index-3}}\label{\detokenize{pair_brownian:index-2}}\label{\detokenize{pair_brownian:index-1}}\label{\detokenize{pair_brownian:index-0}}\label{\detokenize{pair_brownian::doc}}
\sphinxAtStartPar
Accelerator Variants: \sphinxstyleemphasis{brownian/omp}


\section{pair\_style brownian/poly command}
\label{\detokenize{pair_brownian:pair-style-brownian-poly-command}}
\sphinxAtStartPar
Accelerator Variants: \sphinxstyleemphasis{brownian/poly/omp}


\subsection{Syntax}
\label{\detokenize{pair_brownian:syntax}}
\begin{sphinxVerbatim}[commandchars=\\\{\}]
\PYG{k}{pair\PYGZus{}style}\PYG{+w}{ }\PYG{n}{style}\PYG{+w}{ }\PYG{n}{mu}\PYG{+w}{ }\PYG{n}{flaglog}\PYG{+w}{ }\PYG{n}{flagfld}\PYG{+w}{ }\PYG{n}{cutinner}\PYG{+w}{ }\PYG{n}{cutoff}\PYG{+w}{ }\PYG{n}{t\PYGZus{}target}\PYG{+w}{ }\PYG{n}{seed}\PYG{+w}{ }\PYG{n}{flagHI}\PYG{+w}{ }\PYG{n}{flagVF}
\end{sphinxVerbatim}
\begin{itemize}
\item {} 
\sphinxAtStartPar
style = \sphinxstyleemphasis{brownian} or \sphinxstyleemphasis{brownian/poly}

\item {} 
\sphinxAtStartPar
mu = dynamic viscosity (dynamic viscosity units)

\item {} 
\sphinxAtStartPar
flaglog = 0/1 log terms in the lubrication approximation on/off

\item {} 
\sphinxAtStartPar
flagfld = 0/1 to include/exclude Fast Lubrication Dynamics effects

\item {} 
\sphinxAtStartPar
cutinner = inner cutoff distance (distance units)

\item {} 
\sphinxAtStartPar
cutoff = outer cutoff for interactions (distance units)

\item {} 
\sphinxAtStartPar
t\_target = target temp of the system (temperature units)

\item {} 
\sphinxAtStartPar
seed = seed for the random number generator (positive integer)

\item {} 
\sphinxAtStartPar
flagHI (optional) = 0/1 to include/exclude 1/r hydrodynamic interactions

\item {} 
\sphinxAtStartPar
flagVF (optional) = 0/1 to include/exclude volume fraction corrections in the long\sphinxhyphen{}range isotropic terms

\end{itemize}


\subsection{Examples}
\label{\detokenize{pair_brownian:examples}}
\begin{sphinxVerbatim}[commandchars=\\\{\}]
\PYG{k}{pair\PYGZus{}style}\PYG{+w}{ }\PYG{n}{brownian}\PYG{+w}{ }\PYG{l+m}{1.5}\PYG{+w}{ }\PYG{l+m}{1}\PYG{+w}{ }\PYG{l+m}{1}\PYG{+w}{ }\PYG{l+m}{2.01}\PYG{+w}{ }\PYG{l+m}{2.5}\PYG{+w}{ }\PYG{l+m}{2.0}\PYG{+w}{ }\PYG{l+m}{5878567}\PYG{+w}{ }\PYG{c}{\PYGZsh{} (assuming radius = 1)}
\PYG{k}{pair\PYGZus{}coeff}\PYG{+w}{ }\PYG{l+m}{1}\PYG{+w}{ }\PYG{l+m}{1}\PYG{+w}{ }\PYG{l+m}{2.05}\PYG{+w}{ }\PYG{l+m}{2.8}
\PYG{k}{pair\PYGZus{}coeff}\PYG{+w}{ }\PYG{o}{*}\PYG{+w}{ }\PYG{o}{*}
\end{sphinxVerbatim}


\subsection{Description}
\label{\detokenize{pair_brownian:description}}
\sphinxAtStartPar
Styles \sphinxstyleemphasis{brownian} and \sphinxstyleemphasis{brownian/poly} compute Brownian forces and
torques on finite\sphinxhyphen{}size spherical particles.  The former requires
monodisperse spherical particles; the latter allows for polydisperse
spherical particles.

\sphinxAtStartPar
These pair styles are designed to be used with either the
{\hyperref[\detokenize{pair_lubricate::doc}]{\sphinxcrossref{\DUrole{doc}{pair\_style lubricate}}}} or {\hyperref[\detokenize{pair_lubricateU::doc}]{\sphinxcrossref{\DUrole{doc}{pair\_style
lubricateU}}}} commands to provide thermostatting when
dissipative lubrication forces are acting.  Thus the parameters \sphinxstyleemphasis{mu},
\sphinxstyleemphasis{flaglog}, \sphinxstyleemphasis{flagfld}, \sphinxstyleemphasis{cutinner}, and \sphinxstyleemphasis{cutoff} should be specified
consistent with the settings in the lubrication pair styles.  For
details, refer to either of the lubrication pair styles.

\sphinxAtStartPar
The \sphinxstyleemphasis{t\_target} setting is used to specify the target temperature of
the system.  The random number \sphinxstyleemphasis{seed} is used to generate random
numbers for the thermostatting procedure.

\sphinxAtStartPar
The \sphinxstyleemphasis{flagHI} and \sphinxstyleemphasis{flagVF} settings are optional.  Neither should be
used, or both must be defined.


\bigskip\hrule\bigskip


\sphinxAtStartPar
The following coefficients must be defined for each pair of atoms
types via the {\hyperref[\detokenize{pair_coeff::doc}]{\sphinxcrossref{\DUrole{doc}{pair\_coeff}}}} command as in the examples
above, or in the data file or restart files read by the
{\hyperref[\detokenize{read_data::doc}]{\sphinxcrossref{\DUrole{doc}{read\_data}}}} or {\hyperref[\detokenize{read_restart::doc}]{\sphinxcrossref{\DUrole{doc}{read\_restart}}}}
commands, or by mixing as described below:
\begin{itemize}
\item {} 
\sphinxAtStartPar
cutinner (distance units)

\item {} 
\sphinxAtStartPar
cutoff (distance units)

\end{itemize}

\sphinxAtStartPar
The two coefficients are optional.  If neither is specified, the two
cutoffs specified in the pair\_style command are used.  Otherwise both
must be specified.


\bigskip\hrule\bigskip


\sphinxAtStartPar
Styles with a \sphinxstyleemphasis{gpu}, \sphinxstyleemphasis{intel}, \sphinxstyleemphasis{kk}, \sphinxstyleemphasis{omp}, or \sphinxstyleemphasis{opt} suffix are
functionally the same as the corresponding style without the suffix.
They have been optimized to run faster, depending on your available
hardware, as discussed on the {\hyperref[\detokenize{Speed_packages::doc}]{\sphinxcrossref{\DUrole{doc}{Accelerator packages}}}}
page.  The accelerated styles take the same arguments and should
produce the same results, except for round\sphinxhyphen{}off and precision issues.

\sphinxAtStartPar
These accelerated styles are part of the GPU, INTEL, KOKKOS,
OPENMP, and OPT packages, respectively.  They are only enabled if
LAMMPS was built with those packages.  See the {\hyperref[\detokenize{Build_package::doc}]{\sphinxcrossref{\DUrole{doc}{Build package}}}} page for more info.

\sphinxAtStartPar
You can specify the accelerated styles explicitly in your input script
by including their suffix, or you can use the {\hyperref[\detokenize{Run_options::doc}]{\sphinxcrossref{\DUrole{doc}{\sphinxhyphen{}suffix command\sphinxhyphen{}line switch}}}} when you invoke LAMMPS, or you can use the
{\hyperref[\detokenize{suffix::doc}]{\sphinxcrossref{\DUrole{doc}{suffix}}}} command in your input script.

\sphinxAtStartPar
See the {\hyperref[\detokenize{Speed_packages::doc}]{\sphinxcrossref{\DUrole{doc}{Accelerator packages}}}} page for more
instructions on how to use the accelerated styles effectively.


\bigskip\hrule\bigskip



\subsection{Mixing, shift, table, tail correction, restart, rRESPA info}
\label{\detokenize{pair_brownian:mixing-shift-table-tail-correction-restart-rrespa-info}}
\sphinxAtStartPar
For atom type pairs I,J and I != J, the two cutoff distances for these
pair styles can be mixed.  The default mix value is \sphinxstyleemphasis{geometric}.  See
the “pair\_modify” command for details.

\sphinxAtStartPar
These pair styles do not support the {\hyperref[\detokenize{pair_modify::doc}]{\sphinxcrossref{\DUrole{doc}{pair\_modify}}}}
shift option for the energy of the pair interaction.

\sphinxAtStartPar
The {\hyperref[\detokenize{pair_modify::doc}]{\sphinxcrossref{\DUrole{doc}{pair\_modify}}}} table option is not relevant
for these pair styles.

\sphinxAtStartPar
These pair styles do not support the {\hyperref[\detokenize{pair_modify::doc}]{\sphinxcrossref{\DUrole{doc}{pair\_modify}}}}
tail option for adding long\sphinxhyphen{}range tail corrections to energy and
pressure.

\sphinxAtStartPar
These pair styles write their information to {\hyperref[\detokenize{restart::doc}]{\sphinxcrossref{\DUrole{doc}{binary restart files}}}}, so pair\_style and pair\_coeff commands do not need to be
specified in an input script that reads a restart file.

\sphinxAtStartPar
These pair styles can only be used via the \sphinxstyleemphasis{pair} keyword of the
{\hyperref[\detokenize{run_style::doc}]{\sphinxcrossref{\DUrole{doc}{run\_style respa}}}} command.  They do not support the
\sphinxstyleemphasis{inner}, \sphinxstyleemphasis{middle}, \sphinxstyleemphasis{outer} keywords.


\bigskip\hrule\bigskip



\subsection{Restrictions}
\label{\detokenize{pair_brownian:restrictions}}
\sphinxAtStartPar
These styles are part of the COLLOID package.  They are only enabled if
LAMMPS was built with that package.  See the {\hyperref[\detokenize{Build_package::doc}]{\sphinxcrossref{\DUrole{doc}{Build package}}}} page for more info.

\sphinxAtStartPar
Only spherical monodisperse particles are allowed for pair\_style
brownian.

\sphinxAtStartPar
Only spherical particles are allowed for pair\_style brownian/poly.

\sphinxAtStartPar
These pair styles are only compatible with the following wall fixes:
doc:\sphinxtitleref{fix wall/lj93, fix wall/lj126, fix wall/lj1043, fix wall/colloid,
fix wall/harmonic, fix wall/lepton, fix wall/morse, fix wall/table
\textless{}fix\_wall\textgreater{}}.


\subsection{Related commands}
\label{\detokenize{pair_brownian:related-commands}}
\sphinxAtStartPar
{\hyperref[\detokenize{pair_coeff::doc}]{\sphinxcrossref{\DUrole{doc}{pair\_coeff}}}}, {\hyperref[\detokenize{pair_lubricate::doc}]{\sphinxcrossref{\DUrole{doc}{pair\_style lubricate}}}}, {\hyperref[\detokenize{pair_lubricateU::doc}]{\sphinxcrossref{\DUrole{doc}{pair\_style lubricateU}}}}


\subsection{Default}
\label{\detokenize{pair_brownian:default}}
\sphinxAtStartPar
The default settings for the optional args are flagHI = 1 and flagVF =
1.

\sphinxstepscope

\index{pair\_style buck@\spxentry{pair\_style buck}}\ignorespaces 
\index{pair\_style buck/gpu@\spxentry{pair\_style buck/gpu}}\ignorespaces 
\index{pair\_style buck/intel@\spxentry{pair\_style buck/intel}}\ignorespaces 
\index{pair\_style buck/kk@\spxentry{pair\_style buck/kk}}\ignorespaces 
\index{pair\_style buck/omp@\spxentry{pair\_style buck/omp}}\ignorespaces 
\index{pair\_style buck/coul/cut@\spxentry{pair\_style buck/coul/cut}}\ignorespaces 
\index{pair\_style buck/coul/cut/gpu@\spxentry{pair\_style buck/coul/cut/gpu}}\ignorespaces 
\index{pair\_style buck/coul/cut/intel@\spxentry{pair\_style buck/coul/cut/intel}}\ignorespaces 
\index{pair\_style buck/coul/cut/kk@\spxentry{pair\_style buck/coul/cut/kk}}\ignorespaces 
\index{pair\_style buck/coul/cut/omp@\spxentry{pair\_style buck/coul/cut/omp}}\ignorespaces 
\index{pair\_style buck/coul/long@\spxentry{pair\_style buck/coul/long}}\ignorespaces 
\index{pair\_style buck/coul/long/gpu@\spxentry{pair\_style buck/coul/long/gpu}}\ignorespaces 
\index{pair\_style buck/coul/long/intel@\spxentry{pair\_style buck/coul/long/intel}}\ignorespaces 
\index{pair\_style buck/coul/long/kk@\spxentry{pair\_style buck/coul/long/kk}}\ignorespaces 
\index{pair\_style buck/coul/long/omp@\spxentry{pair\_style buck/coul/long/omp}}\ignorespaces 
\index{pair\_style buck/coul/msm@\spxentry{pair\_style buck/coul/msm}}\ignorespaces 
\index{pair\_style buck/coul/msm/omp@\spxentry{pair\_style buck/coul/msm/omp}}\ignorespaces 

\section{pair\_style buck command}
\label{\detokenize{pair_buck:pair-style-buck-command}}\label{\detokenize{pair_buck:index-16}}\label{\detokenize{pair_buck:index-15}}\label{\detokenize{pair_buck:index-14}}\label{\detokenize{pair_buck:index-13}}\label{\detokenize{pair_buck:index-12}}\label{\detokenize{pair_buck:index-11}}\label{\detokenize{pair_buck:index-10}}\label{\detokenize{pair_buck:index-9}}\label{\detokenize{pair_buck:index-8}}\label{\detokenize{pair_buck:index-7}}\label{\detokenize{pair_buck:index-6}}\label{\detokenize{pair_buck:index-5}}\label{\detokenize{pair_buck:index-4}}\label{\detokenize{pair_buck:index-3}}\label{\detokenize{pair_buck:index-2}}\label{\detokenize{pair_buck:index-1}}\label{\detokenize{pair_buck:index-0}}\label{\detokenize{pair_buck::doc}}
\sphinxAtStartPar
Accelerator Variants: \sphinxstyleemphasis{buck/gpu}, \sphinxstyleemphasis{buck/intel}, \sphinxstyleemphasis{buck/kk}, \sphinxstyleemphasis{buck/omp}


\section{pair\_style buck/coul/cut command}
\label{\detokenize{pair_buck:pair-style-buck-coul-cut-command}}
\sphinxAtStartPar
Accelerator Variants: \sphinxstyleemphasis{buck/coul/cut/gpu}, \sphinxstyleemphasis{buck/coul/cut/intel}, \sphinxstyleemphasis{buck/coul/cut/kk}, \sphinxstyleemphasis{buck/coul/cut/omp}


\section{pair\_style buck/coul/long command}
\label{\detokenize{pair_buck:pair-style-buck-coul-long-command}}
\sphinxAtStartPar
Accelerator Variants: \sphinxstyleemphasis{buck/coul/long/gpu}, \sphinxstyleemphasis{buck/coul/long/intel}, \sphinxstyleemphasis{buck/coul/long/kk}, \sphinxstyleemphasis{buck/coul/long/omp}


\section{pair\_style buck/coul/msm command}
\label{\detokenize{pair_buck:pair-style-buck-coul-msm-command}}
\sphinxAtStartPar
Accelerator Variants: \sphinxstyleemphasis{buck/coul/msm/omp}


\subsection{Syntax}
\label{\detokenize{pair_buck:syntax}}
\begin{sphinxVerbatim}[commandchars=\\\{\}]
\PYG{k}{pair\PYGZus{}style}\PYG{+w}{ }\PYG{n}{style}\PYG{+w}{ }\PYG{n}{args}
\end{sphinxVerbatim}
\begin{itemize}
\item {} 
\sphinxAtStartPar
style = \sphinxstyleemphasis{buck} or \sphinxstyleemphasis{buck/coul/cut} or \sphinxstyleemphasis{buck/coul/long} or \sphinxstyleemphasis{buck/coul/msm}

\item {} 
\sphinxAtStartPar
args = list of arguments for a particular style

\end{itemize}
\begin{sphinxalltt}
\sphinxstyleemphasis{buck} args = cutoff
  cutoff = global cutoff for Buckingham interactions (distance units)
\sphinxstyleemphasis{buck/coul/cut} args = cutoff (cutoff2)
  cutoff = global cutoff for Buckingham (and Coulombic if only 1 arg) (distance units)
  cutoff2 = global cutoff for Coulombic (optional) (distance units)
\sphinxstyleemphasis{buck/coul/long} args = cutoff (cutoff2)
  cutoff = global cutoff for Buckingham (and Coulombic if only 1 arg) (distance units)
  cutoff2 = global cutoff for Coulombic (optional) (distance units)
\sphinxstyleemphasis{buck/coul/msm} args = cutoff (cutoff2)
  cutoff = global cutoff for Buckingham (and Coulombic if only 1 arg) (distance units)
  cutoff2 = global cutoff for Coulombic (optional) (distance units)
\end{sphinxalltt}


\subsection{Examples}
\label{\detokenize{pair_buck:examples}}
\begin{sphinxVerbatim}[commandchars=\\\{\}]
\PYG{k}{pair\PYGZus{}style}\PYG{+w}{ }\PYG{n}{buck}\PYG{+w}{ }\PYG{l+m}{2.5}
\PYG{k}{pair\PYGZus{}coeff}\PYG{+w}{ }\PYG{o}{*}\PYG{+w}{ }\PYG{o}{*}\PYG{+w}{ }\PYG{l+m}{100.0}\PYG{+w}{ }\PYG{l+m}{1.5}\PYG{+w}{ }\PYG{l+m}{200.0}
\PYG{k}{pair\PYGZus{}coeff}\PYG{+w}{ }\PYG{o}{*}\PYG{+w}{ }\PYG{o}{*}\PYG{+w}{ }\PYG{l+m}{100.0}\PYG{+w}{ }\PYG{l+m}{1.5}\PYG{+w}{ }\PYG{l+m}{200.0}\PYG{+w}{ }\PYG{l+m}{3.0}

\PYG{k}{pair\PYGZus{}style}\PYG{+w}{ }\PYG{n}{buck}\PYG{o}{/}\PYG{n}{coul}\PYG{o}{/}\PYG{n}{cut}\PYG{+w}{ }\PYG{l+m}{10.0}
\PYG{k}{pair\PYGZus{}style}\PYG{+w}{ }\PYG{n}{buck}\PYG{o}{/}\PYG{n}{coul}\PYG{o}{/}\PYG{n}{cut}\PYG{+w}{ }\PYG{l+m}{10.0}\PYG{+w}{ }\PYG{l+m}{8.0}
\PYG{k}{pair\PYGZus{}coeff}\PYG{+w}{ }\PYG{o}{*}\PYG{+w}{ }\PYG{o}{*}\PYG{+w}{ }\PYG{l+m}{100.0}\PYG{+w}{ }\PYG{l+m}{1.5}\PYG{+w}{ }\PYG{l+m}{200.0}
\PYG{k}{pair\PYGZus{}coeff}\PYG{+w}{ }\PYG{l+m}{1}\PYG{+w}{ }\PYG{l+m}{1}\PYG{+w}{ }\PYG{l+m}{100.0}\PYG{+w}{ }\PYG{l+m}{1.5}\PYG{+w}{ }\PYG{l+m}{200.0}\PYG{+w}{ }\PYG{l+m}{9.0}
\PYG{k}{pair\PYGZus{}coeff}\PYG{+w}{ }\PYG{l+m}{1}\PYG{+w}{ }\PYG{l+m}{1}\PYG{+w}{ }\PYG{l+m}{100.0}\PYG{+w}{ }\PYG{l+m}{1.5}\PYG{+w}{ }\PYG{l+m}{200.0}\PYG{+w}{ }\PYG{l+m}{9.0}\PYG{+w}{ }\PYG{l+m}{8.0}

\PYG{k}{pair\PYGZus{}style}\PYG{+w}{ }\PYG{n}{buck}\PYG{o}{/}\PYG{n}{coul}\PYG{o}{/}\PYG{n}{long}\PYG{+w}{ }\PYG{l+m}{10.0}
\PYG{k}{pair\PYGZus{}style}\PYG{+w}{ }\PYG{n}{buck}\PYG{o}{/}\PYG{n}{coul}\PYG{o}{/}\PYG{n}{long}\PYG{+w}{ }\PYG{l+m}{10.0}\PYG{+w}{ }\PYG{l+m}{8.0}
\PYG{k}{pair\PYGZus{}coeff}\PYG{+w}{ }\PYG{o}{*}\PYG{+w}{ }\PYG{o}{*}\PYG{+w}{ }\PYG{l+m}{100.0}\PYG{+w}{ }\PYG{l+m}{1.5}\PYG{+w}{ }\PYG{l+m}{200.0}
\PYG{k}{pair\PYGZus{}coeff}\PYG{+w}{ }\PYG{l+m}{1}\PYG{+w}{ }\PYG{l+m}{1}\PYG{+w}{ }\PYG{l+m}{100.0}\PYG{+w}{ }\PYG{l+m}{1.5}\PYG{+w}{ }\PYG{l+m}{200.0}\PYG{+w}{ }\PYG{l+m}{9.0}

\PYG{k}{pair\PYGZus{}style}\PYG{+w}{ }\PYG{n}{buck}\PYG{o}{/}\PYG{n}{coul}\PYG{o}{/}\PYG{n}{msm}\PYG{+w}{ }\PYG{l+m}{10.0}
\PYG{k}{pair\PYGZus{}style}\PYG{+w}{ }\PYG{n}{buck}\PYG{o}{/}\PYG{n}{coul}\PYG{o}{/}\PYG{n}{msm}\PYG{+w}{ }\PYG{l+m}{10.0}\PYG{+w}{ }\PYG{l+m}{8.0}
\PYG{k}{pair\PYGZus{}coeff}\PYG{+w}{ }\PYG{o}{*}\PYG{+w}{ }\PYG{o}{*}\PYG{+w}{ }\PYG{l+m}{100.0}\PYG{+w}{ }\PYG{l+m}{1.5}\PYG{+w}{ }\PYG{l+m}{200.0}
\PYG{k}{pair\PYGZus{}coeff}\PYG{+w}{ }\PYG{l+m}{1}\PYG{+w}{ }\PYG{l+m}{1}\PYG{+w}{ }\PYG{l+m}{100.0}\PYG{+w}{ }\PYG{l+m}{1.5}\PYG{+w}{ }\PYG{l+m}{200.0}\PYG{+w}{ }\PYG{l+m}{9.0}
\end{sphinxVerbatim}


\subsection{Description}
\label{\detokenize{pair_buck:description}}
\sphinxAtStartPar
The \sphinxstyleemphasis{buck} style computes a Buckingham potential (exp/6 instead of
Lennard\sphinxhyphen{}Jones 12/6) given by
\begin{equation*}
\begin{split}E = A e^{-r / \rho} - \frac{C}{r^6} \qquad r < r_c\end{split}
\end{equation*}
\sphinxAtStartPar
where \(\rho\) is an ionic\sphinxhyphen{}pair dependent length parameter, and
\(r_c\) is the cutoff on both terms.

\sphinxAtStartPar
The styles with \sphinxstyleemphasis{coul/cut} or \sphinxstyleemphasis{coul/long} or \sphinxstyleemphasis{coul/msm} add a
Coulombic term as described for the {\hyperref[\detokenize{pair_lj::doc}]{\sphinxcrossref{\DUrole{doc}{lj/cut}}}} pair styles.
For \sphinxstyleemphasis{buck/coul/long} and \sphinxstyleemphasis{buc/coul/msm}, an additional damping factor
is applied to the Coulombic term so it can be used in conjunction with
the {\hyperref[\detokenize{kspace_style::doc}]{\sphinxcrossref{\DUrole{doc}{kspace\_style}}}} command and its \sphinxstyleemphasis{ewald} or \sphinxstyleemphasis{pppm}
or \sphinxstyleemphasis{msm} option.  The Coulombic cutoff specified for this style means
that pairwise interactions within this distance are computed directly;
interactions outside that distance are computed in reciprocal space.

\sphinxAtStartPar
If one cutoff is specified for the \sphinxstyleemphasis{born/coul/cut} and
\sphinxstyleemphasis{born/coul/long} and \sphinxstyleemphasis{born/coul/msm} styles, it is used for both the
A,C and Coulombic terms.  If two cutoffs are specified, the first is
used as the cutoff for the A,C terms, and the second is the cutoff for
the Coulombic term.

\sphinxAtStartPar
Note that these potentials are related to the {\hyperref[\detokenize{pair_born::doc}]{\sphinxcrossref{\DUrole{doc}{Born\sphinxhyphen{}Mayer\sphinxhyphen{}Huggins potential}}}}.

\begin{sphinxadmonition}{note}{Note:}
\sphinxAtStartPar
For all these pair styles, the terms with A and C are always
cutoff.  The additional Coulombic term can be cutoff or long\sphinxhyphen{}range (no
cutoff) depending on whether the style name includes coul/cut or
coul/long or coul/msm.  If you wish the C/r\textasciicircum{}6 term to be long\sphinxhyphen{}range
(no cutoff), then see the {\hyperref[\detokenize{pair_buck_long::doc}]{\sphinxcrossref{\DUrole{doc}{pair\_style buck/long/coul/long}}}} command.
\end{sphinxadmonition}

\sphinxAtStartPar
The following coefficients must be defined for each pair of atoms
types via the {\hyperref[\detokenize{pair_coeff::doc}]{\sphinxcrossref{\DUrole{doc}{pair\_coeff}}}} command as in the examples
above, or in the data file or restart files read by the
{\hyperref[\detokenize{read_data::doc}]{\sphinxcrossref{\DUrole{doc}{read\_data}}}} or {\hyperref[\detokenize{read_restart::doc}]{\sphinxcrossref{\DUrole{doc}{read\_restart}}}}
commands:
\begin{itemize}
\item {} 
\sphinxAtStartPar
A (energy units)

\item {} 
\sphinxAtStartPar
\(\rho\) (distance units)

\item {} 
\sphinxAtStartPar
C (energy\sphinxhyphen{}distance\textasciicircum{}6 units)

\item {} 
\sphinxAtStartPar
cutoff (distance units)

\item {} 
\sphinxAtStartPar
cutoff2 (distance units)

\end{itemize}

\sphinxAtStartPar
The second coefficient, \(\rho\), must be greater than zero.
The coefficients A, \(\rho\), and C can be written as analytical expressions
of \(\epsilon\) and \(\sigma\), in analogy to the Lennard\sphinxhyphen{}Jones potential
{\hyperref[\detokenize{pair_buck:khrapak}]{\sphinxcrossref{\DUrole{std,std-ref}{(Khrapak)}}}}.

\sphinxAtStartPar
The latter 2 coefficients are optional.  If not specified, the global
A,C and Coulombic cutoffs are used.  If only one cutoff is specified,
it is used as the cutoff for both A,C and Coulombic interactions for
this type pair.  If both coefficients are specified, they are used as
the A,C and Coulombic cutoffs for this type pair.  You cannot specify
2 cutoffs for style \sphinxstyleemphasis{buck}, since it has no Coulombic terms.
For \sphinxstyleemphasis{buck/coul/long} only the LJ cutoff can be specified since a
Coulombic cutoff cannot be specified for an individual I,J type pair.
All type pairs use the same global Coulombic cutoff specified in the
pair\_style command.


\bigskip\hrule\bigskip


\sphinxAtStartPar
Styles with a \sphinxstyleemphasis{gpu}, \sphinxstyleemphasis{intel}, \sphinxstyleemphasis{kk}, \sphinxstyleemphasis{omp}, or \sphinxstyleemphasis{opt} suffix are
functionally the same as the corresponding style without the suffix.
They have been optimized to run faster, depending on your available
hardware, as discussed on the {\hyperref[\detokenize{Speed_packages::doc}]{\sphinxcrossref{\DUrole{doc}{Accelerator packages}}}}
page.  The accelerated styles take the same arguments and should
produce the same results, except for round\sphinxhyphen{}off and precision issues.

\sphinxAtStartPar
These accelerated styles are part of the GPU, INTEL, KOKKOS,
OPENMP, and OPT packages, respectively.  They are only enabled if
LAMMPS was built with those packages.  See the {\hyperref[\detokenize{Build_package::doc}]{\sphinxcrossref{\DUrole{doc}{Build package}}}} page for more info.

\sphinxAtStartPar
You can specify the accelerated styles explicitly in your input script
by including their suffix, or you can use the {\hyperref[\detokenize{Run_options::doc}]{\sphinxcrossref{\DUrole{doc}{\sphinxhyphen{}suffix command\sphinxhyphen{}line switch}}}} when you invoke LAMMPS, or you can use the
{\hyperref[\detokenize{suffix::doc}]{\sphinxcrossref{\DUrole{doc}{suffix}}}} command in your input script.

\sphinxAtStartPar
See the {\hyperref[\detokenize{Speed_packages::doc}]{\sphinxcrossref{\DUrole{doc}{Accelerator packages}}}} page for more
instructions on how to use the accelerated styles effectively.


\bigskip\hrule\bigskip



\subsection{Mixing, shift, table, tail correction, restart, rRESPA info}
\label{\detokenize{pair_buck:mixing-shift-table-tail-correction-restart-rrespa-info}}
\sphinxAtStartPar
These pair styles do not support mixing.  Thus, coefficients for all
I,J pairs must be specified explicitly.

\sphinxAtStartPar
These styles support the {\hyperref[\detokenize{pair_modify::doc}]{\sphinxcrossref{\DUrole{doc}{pair\_modify}}}} shift option
for the energy of the exp() and 1/r\textasciicircum{}6 portion of the pair interaction.

\sphinxAtStartPar
The \sphinxstyleemphasis{buck/coul/long} pair style supports the
{\hyperref[\detokenize{pair_modify::doc}]{\sphinxcrossref{\DUrole{doc}{pair\_modify}}}} table option to tabulate the
short\sphinxhyphen{}range portion of the long\sphinxhyphen{}range Coulombic interaction.

\sphinxAtStartPar
These styles support the pair\_modify tail option for adding long\sphinxhyphen{}range
tail corrections to energy and pressure for the A,C terms in the
pair interaction.

\sphinxAtStartPar
These styles write their information to {\hyperref[\detokenize{restart::doc}]{\sphinxcrossref{\DUrole{doc}{binary restart files}}}}, so pair\_style and pair\_coeff commands do not need
to be specified in an input script that reads a restart file.

\sphinxAtStartPar
These styles can only be used via the \sphinxstyleemphasis{pair} keyword of the {\hyperref[\detokenize{run_style::doc}]{\sphinxcrossref{\DUrole{doc}{run\_style respa}}}} command.  They do not support the \sphinxstyleemphasis{inner},
\sphinxstyleemphasis{middle}, \sphinxstyleemphasis{outer} keywords.


\subsection{Restrictions}
\label{\detokenize{pair_buck:restrictions}}
\sphinxAtStartPar
The \sphinxstyleemphasis{buck/coul/long} style is part of the KSPACE package.  They are
only enabled if LAMMPS was built with that package.  See the {\hyperref[\detokenize{Build_package::doc}]{\sphinxcrossref{\DUrole{doc}{Build package}}}} page for more info.


\subsection{Related commands}
\label{\detokenize{pair_buck:related-commands}}
\sphinxAtStartPar
{\hyperref[\detokenize{pair_coeff::doc}]{\sphinxcrossref{\DUrole{doc}{pair\_coeff}}}}, {\hyperref[\detokenize{pair_born::doc}]{\sphinxcrossref{\DUrole{doc}{pair\_style born}}}}


\subsection{Default}
\label{\detokenize{pair_buck:default}}
\sphinxAtStartPar
none

\phantomsection\label{\detokenize{pair_buck:khrapak}}
\sphinxAtStartPar
\sphinxstylestrong{(Khrapak)} Khrapak, Chaudhuri, and Morfill, J Chem Phys, 134, 054120 (2011).

\sphinxstepscope

\index{pair\_style buck6d/coul/gauss/dsf@\spxentry{pair\_style buck6d/coul/gauss/dsf}}\ignorespaces 
\index{pair\_style buck6d/coul/gauss/long@\spxentry{pair\_style buck6d/coul/gauss/long}}\ignorespaces 

\section{pair\_style buck6d/coul/gauss/dsf command}
\label{\detokenize{pair_buck6d_coul_gauss:pair-style-buck6d-coul-gauss-dsf-command}}\label{\detokenize{pair_buck6d_coul_gauss:index-1}}\label{\detokenize{pair_buck6d_coul_gauss:index-0}}\label{\detokenize{pair_buck6d_coul_gauss::doc}}

\section{pair\_style buck6d/coul/gauss/long command}
\label{\detokenize{pair_buck6d_coul_gauss:pair-style-buck6d-coul-gauss-long-command}}

\subsection{Syntax}
\label{\detokenize{pair_buck6d_coul_gauss:syntax}}
\begin{sphinxVerbatim}[commandchars=\\\{\}]
\PYG{k}{pair\PYGZus{}style}\PYG{+w}{ }\PYG{n}{style}\PYG{+w}{ }\PYG{n}{args}
\end{sphinxVerbatim}
\begin{itemize}
\item {} 
\sphinxAtStartPar
style = \sphinxstyleemphasis{buck6d/coul/gauss/dsf} or \sphinxstyleemphasis{buck6d/coul/gauss/long}

\item {} 
\sphinxAtStartPar
args = list of arguments for a particular style

\end{itemize}
\begin{sphinxalltt}
\sphinxstyleemphasis{buck6d/coul/gauss/dsf} args = smooth cutoff (cutoff2)
  smooth  = smoothing onset within Buckingham cutoff (ratio)
  cutoff  = global cutoff for Buckingham (and Coulombic if only 1 arg) (distance units)
  cutoff2 = global cutoff for Coulombic (optional) (distance units)
\sphinxstyleemphasis{buck6d/coul/gauss/long} args = smooth smooth2 cutoff (cutoff2)
  smooth   = smoothing onset within Buckingham cutoff (ratio)
  smooth2  = smoothing onset within Coulombic cutoff (ratio)
  cutoff   = global cutoff for Buckingham (and Coulombic if only 1 arg) (distance units)
  cutoff2  = global cutoff for Coulombic (optional) (distance units)
\end{sphinxalltt}


\subsection{Examples}
\label{\detokenize{pair_buck6d_coul_gauss:examples}}
\begin{sphinxVerbatim}[commandchars=\\\{\}]
\PYG{k}{pair\PYGZus{}style}\PYG{+w}{ }\PYG{n}{buck6d}\PYG{o}{/}\PYG{n}{coul}\PYG{o}{/}\PYG{n}{gauss}\PYG{o}{/}\PYG{n}{dsf}\PYG{+w}{    }\PYG{l+m}{0.9000}\PYG{+w}{    }\PYG{l+m}{12.0000}
\PYG{k}{pair\PYGZus{}coeff}\PYG{+w}{ }\PYG{l+m}{1}\PYG{+w}{  }\PYG{l+m}{1}\PYG{+w}{  }\PYG{l+m}{1030}\PYG{n}{.}\PYG{+w}{  }\PYG{l+m}{3.061}\PYG{+w}{  }\PYG{l+m}{457.179}\PYG{+w}{  }\PYG{l+m}{4.521}\PYG{+w}{  }\PYG{l+m}{0.608}

\PYG{k}{pair\PYGZus{}style}\PYG{+w}{ }\PYG{n}{buck6d}\PYG{o}{/}\PYG{n}{coul}\PYG{o}{/}\PYG{n}{gauss}\PYG{o}{/}\PYG{n}{long}\PYG{+w}{   }\PYG{l+m}{0.9000}\PYG{+w}{  }\PYG{l+m}{1.0000}\PYG{+w}{  }\PYG{l+m}{12.0000}
\PYG{k}{pair\PYGZus{}coeff}\PYG{+w}{ }\PYG{l+m}{1}\PYG{+w}{  }\PYG{l+m}{1}\PYG{+w}{  }\PYG{l+m}{1030}\PYG{n}{.}\PYG{+w}{  }\PYG{l+m}{3.061}\PYG{+w}{  }\PYG{l+m}{457.179}\PYG{+w}{  }\PYG{l+m}{4.521}\PYG{+w}{  }\PYG{l+m}{0.608}
\end{sphinxVerbatim}


\subsection{Description}
\label{\detokenize{pair_buck6d_coul_gauss:description}}
\sphinxAtStartPar
The \sphinxstyleemphasis{buck6d/coul/gauss} styles evaluate vdW and Coulomb
interactions following the MOF\sphinxhyphen{}FF force field after
{\hyperref[\detokenize{pair_buck6d_coul_gauss:schmid}]{\sphinxcrossref{\DUrole{std,std-ref}{(Schmid)}}}}. The vdW term of the \sphinxstyleemphasis{buck6d} styles
computes a dispersion damped Buckingham potential:
\begin{equation*}
\begin{split}E = A e^{-\kappa r} - \frac{C}{r^6} \cdot \frac{1}{1 + D r^{14}} \qquad r < r_c \\\end{split}
\end{equation*}
\sphinxAtStartPar
where A and C are a force constant, \(\kappa\) is an ionic\sphinxhyphen{}pair dependent
reciprocal length parameter, D is a dispersion correction parameter,
and the cutoff \(r_c\) truncates the interaction distance.
The first term in the potential corresponds to the Buckingham
repulsion term and the second term to the dispersion attraction with
a damping correction analog to the Grimme correction used in DFT.
The latter corrects for artifacts occurring at short distances which
become an issue for soft vdW potentials.

\sphinxAtStartPar
The \sphinxstyleemphasis{buck6d} styles include a smoothing function which is invoked
according to the global smoothing parameter within the specified
cutoff.  Hereby a parameter of i.e. 0.9 invokes the smoothing
within 90\% of the cutoff.  No smoothing is applied at a value
of 1.0. For the \sphinxstyleemphasis{gauss/dsf} style this smoothing is only applicable
for the dispersion damped Buckingham potential. For the \sphinxstyleemphasis{gauss/long}
styles the smoothing function can also be invoked for the real
space coulomb interactions which enforce continuous energies and
forces at the cutoff.

\sphinxAtStartPar
Both styles \sphinxstyleemphasis{buck6d/coul/gauss/dsf} and \sphinxstyleemphasis{buck6d/coul/gauss/long}
evaluate a Coulomb potential using spherical Gaussian type charge
distributions which effectively dampen electrostatic interactions
for high charges at close distances.  The electrostatic potential
is thus evaluated as:
\begin{equation*}
\begin{split}E = \frac{C_{q_i q_j}}{\epsilon r_{ij}}\,\, \textrm{erf}\left(\alpha_{ij} r_{ij}\right)\quad\quad\quad r < r_c\end{split}
\end{equation*}
\sphinxAtStartPar
where C is an energy\sphinxhyphen{}conversion constant, \(q_i\) and \(q_j\)
are the charges on the two atoms, epsilon is the dielectric constant which
can be set by the {\hyperref[\detokenize{dielectric::doc}]{\sphinxcrossref{\DUrole{doc}{dielectric}}}} command, \(\alpha\)
is the ion pair dependent damping parameter and erf() is the
error\sphinxhyphen{}function.  The cutoff \(r_c\) truncates the interaction distance.

\sphinxAtStartPar
The style \sphinxstyleemphasis{buck6d/coul/gauss/dsf} computes the Coulomb interaction
via the damped shifted force model described in {\hyperref[\detokenize{pair_buck6d_coul_gauss:fennell}]{\sphinxcrossref{\DUrole{std,std-ref}{(Fennell)}}}}
approximating an Ewald sum similar to the {\hyperref[\detokenize{pair_coul::doc}]{\sphinxcrossref{\DUrole{doc}{pair coul/dsf}}}}
styles. In \sphinxstyleemphasis{buck6d/coul/gauss/long} an additional damping factor is
applied to the Coulombic term so it can be used in conjunction with the
{\hyperref[\detokenize{kspace_style::doc}]{\sphinxcrossref{\DUrole{doc}{kspace\_style}}}} command and its \sphinxstyleemphasis{ewald} or \sphinxstyleemphasis{pppm}
options. The Coulombic cutoff in this case separates the real and
reciprocal space evaluation of the Ewald sum.

\sphinxAtStartPar
If one cutoff is specified it is used for both the vdW and Coulomb
terms.  If two cutoffs are specified, the first is used as the cutoff
for the vdW terms, and the second is the cutoff for the Coulombic term.

\sphinxAtStartPar
The following coefficients must be defined for each pair of atoms
types via the {\hyperref[\detokenize{pair_coeff::doc}]{\sphinxcrossref{\DUrole{doc}{pair\_coeff}}}} command as in the examples
above, or in the data file or restart files read by the
{\hyperref[\detokenize{read_data::doc}]{\sphinxcrossref{\DUrole{doc}{read\_data}}}} or {\hyperref[\detokenize{read_restart::doc}]{\sphinxcrossref{\DUrole{doc}{read\_restart}}}}
commands:
\begin{itemize}
\item {} 
\sphinxAtStartPar
A (energy units)

\item {} 
\sphinxAtStartPar
\(\rho\) (distance\textasciicircum{}\sphinxhyphen{}1 units)

\item {} 
\sphinxAtStartPar
C (energy\sphinxhyphen{}distance\textasciicircum{}6 units)

\item {} 
\sphinxAtStartPar
D (distance\textasciicircum{}14 units)

\item {} 
\sphinxAtStartPar
\(\alpha\) (distance\textasciicircum{}\sphinxhyphen{}1 units)

\item {} 
\sphinxAtStartPar
cutoff (distance units)

\end{itemize}

\sphinxAtStartPar
The second coefficient, \(\rho\), must be greater than zero. The
latter coefficient is optional.  If not specified, the global vdW cutoff
is used.


\bigskip\hrule\bigskip



\subsection{Mixing, shift, table, tail correction, restart, rRESPA info}
\label{\detokenize{pair_buck6d_coul_gauss:mixing-shift-table-tail-correction-restart-rrespa-info}}
\sphinxAtStartPar
These pair styles do not support mixing.  Thus, coefficients for all
I,J pairs must be specified explicitly.

\sphinxAtStartPar
These styles do not support the {\hyperref[\detokenize{pair_modify::doc}]{\sphinxcrossref{\DUrole{doc}{pair\_modify}}}} shift
option for the energy. Instead the smoothing function should be applied
by setting the global smoothing parameter to a value \textless{} 1.0.

\sphinxAtStartPar
These styles write their information to {\hyperref[\detokenize{restart::doc}]{\sphinxcrossref{\DUrole{doc}{binary restart files}}}}, so pair\_style and pair\_coeff commands do not need
to be specified in an input script that reads a restart file.


\subsection{Restrictions}
\label{\detokenize{pair_buck6d_coul_gauss:restrictions}}
\sphinxAtStartPar
These styles are part of the MOFFF package.  They are only
enabled if LAMMPS was built with that package.  See the {\hyperref[\detokenize{Build_package::doc}]{\sphinxcrossref{\DUrole{doc}{Build package}}}} page for more info.


\subsection{Related commands}
\label{\detokenize{pair_buck6d_coul_gauss:related-commands}}
\sphinxAtStartPar
{\hyperref[\detokenize{pair_coeff::doc}]{\sphinxcrossref{\DUrole{doc}{pair\_coeff}}}}


\subsection{Default}
\label{\detokenize{pair_buck6d_coul_gauss:default}}
\sphinxAtStartPar
none

\phantomsection\label{\detokenize{pair_buck6d_coul_gauss:schmid}}\phantomsection\label{\detokenize{pair_buck6d_coul_gauss:fennell}}
\sphinxAtStartPar
\sphinxstylestrong{(Schmid)} S. Bureekaew, S. Amirjalayer, M. Tafipolsky, C. Spickermann, T.K. Roy and R. Schmid, Phys. Status Solidi B, 6, 1128 (2013).

\sphinxAtStartPar
\sphinxstylestrong{(Fennell)} C. J. Fennell, J. D. Gezelter, J Chem Phys, 124, 234104 (2006).

\sphinxstepscope

\index{pair\_style buck/long/coul/long@\spxentry{pair\_style buck/long/coul/long}}\ignorespaces 
\index{pair\_style buck/long/coul/long/omp@\spxentry{pair\_style buck/long/coul/long/omp}}\ignorespaces 

\section{pair\_style buck/long/coul/long command}
\label{\detokenize{pair_buck_long:pair-style-buck-long-coul-long-command}}\label{\detokenize{pair_buck_long:index-1}}\label{\detokenize{pair_buck_long:index-0}}\label{\detokenize{pair_buck_long::doc}}
\sphinxAtStartPar
Accelerator Variants: \sphinxstyleemphasis{buck/long/coul/long/omp}


\subsection{Syntax}
\label{\detokenize{pair_buck_long:syntax}}
\begin{sphinxVerbatim}[commandchars=\\\{\}]
\PYG{k}{pair\PYGZus{}style}\PYG{+w}{ }\PYG{n}{buck}\PYG{o}{/}\PYG{n}{long}\PYG{o}{/}\PYG{n}{coul}\PYG{o}{/}\PYG{n}{long}\PYG{+w}{ }\PYG{n}{flag\PYGZus{}buck}\PYG{+w}{ }\PYG{n}{flag\PYGZus{}coul}\PYG{+w}{ }\PYG{n}{cutoff}\PYG{+w}{ }\PYG{n+nv}{(}\PYG{n+nv}{cutoff2}\PYG{n+nv}{)}
\end{sphinxVerbatim}
\begin{itemize}
\item {} 
\sphinxAtStartPar
flag\_buck = \sphinxstyleemphasis{long} or \sphinxstyleemphasis{cut}
\begin{sphinxalltt}
\sphinxstyleemphasis{long} = use Kspace long\sphinxhyphen{}range summation for the dispersion term 1/r\textasciicircum{}6
\sphinxstyleemphasis{cut} = use a cutoff
\end{sphinxalltt}

\item {} 
\sphinxAtStartPar
flag\_coul = \sphinxstyleemphasis{long} or \sphinxstyleemphasis{off}
\begin{sphinxalltt}
\sphinxstyleemphasis{long} = use Kspace long\sphinxhyphen{}range summation for the Coulombic term 1/r
\sphinxstyleemphasis{off} = omit the Coulombic term
\end{sphinxalltt}

\item {} 
\sphinxAtStartPar
cutoff = global cutoff for Buckingham (and Coulombic if only 1 cutoff) (distance units)

\item {} 
\sphinxAtStartPar
cutoff2 = global cutoff for Coulombic (optional) (distance units)

\end{itemize}


\subsection{Examples}
\label{\detokenize{pair_buck_long:examples}}
\begin{sphinxVerbatim}[commandchars=\\\{\}]
\PYG{k}{pair\PYGZus{}style}\PYG{+w}{ }\PYG{n}{buck}\PYG{o}{/}\PYG{n}{long}\PYG{o}{/}\PYG{n}{coul}\PYG{o}{/}\PYG{n}{long}\PYG{+w}{ }\PYG{n}{cut}\PYG{+w}{ }\PYG{n}{off}\PYG{+w}{ }\PYG{l+m}{2.5}
\PYG{k}{pair\PYGZus{}style}\PYG{+w}{ }\PYG{n}{buck}\PYG{o}{/}\PYG{n}{long}\PYG{o}{/}\PYG{n}{coul}\PYG{o}{/}\PYG{n}{long}\PYG{+w}{ }\PYG{n}{cut}\PYG{+w}{ }\PYG{n}{long}\PYG{+w}{ }\PYG{l+m}{2.5}\PYG{+w}{ }\PYG{l+m}{4.0}
\PYG{k}{pair\PYGZus{}style}\PYG{+w}{ }\PYG{n}{buck}\PYG{o}{/}\PYG{n}{long}\PYG{o}{/}\PYG{n}{coul}\PYG{o}{/}\PYG{n}{long}\PYG{+w}{ }\PYG{n}{long}\PYG{+w}{ }\PYG{n}{long}\PYG{+w}{ }\PYG{l+m}{4.0}
\PYG{k}{pair\PYGZus{}coeff}\PYG{+w}{ }\PYG{o}{*}\PYG{+w}{ }\PYG{o}{*}\PYG{+w}{ }\PYG{l+m}{1}\PYG{+w}{ }\PYG{l+m}{1}
\PYG{k}{pair\PYGZus{}coeff}\PYG{+w}{ }\PYG{l+m}{1}\PYG{+w}{ }\PYG{l+m}{1}\PYG{+w}{ }\PYG{l+m}{1}\PYG{+w}{ }\PYG{l+m}{3}\PYG{+w}{ }\PYG{l+m}{4}
\end{sphinxVerbatim}


\subsection{Description}
\label{\detokenize{pair_buck_long:description}}
\sphinxAtStartPar
The \sphinxstyleemphasis{buck/long/coul/long} style computes a Buckingham potential (exp/6
instead of Lennard\sphinxhyphen{}Jones 12/6) and Coulombic potential, given by
\begin{equation*}
\begin{split}E = & A e^{-r / \rho} - \frac{C}{r^6} \qquad r < r_c \\
E = & \frac{C q_i q_j}{\epsilon  r} \qquad r < r_c\end{split}
\end{equation*}
\sphinxAtStartPar
\(r_c\) is the cutoff.  If one cutoff is specified in the pair\_style
command, it is used for both the Buckingham and Coulombic terms.  If
two cutoffs are specified, they are used as cutoffs for the Buckingham
and Coulombic terms respectively.

\sphinxAtStartPar
The purpose of this pair style is to capture long\sphinxhyphen{}range interactions
resulting from both attractive 1/r\textasciicircum{}6 Buckingham and Coulombic 1/r
interactions.  This is done by use of the \sphinxstyleemphasis{flag\_buck} and \sphinxstyleemphasis{flag\_coul}
settings.  The {\hyperref[\detokenize{pair_buck_long:ismail}]{\sphinxcrossref{\DUrole{std,std-ref}{Ismail}}}} paper has more details on when it is
appropriate to include long\sphinxhyphen{}range 1/r\textasciicircum{}6 interactions, using this
potential.

\sphinxAtStartPar
If \sphinxstyleemphasis{flag\_buck} is set to \sphinxstyleemphasis{long}, no cutoff is used on the Buckingham
1/r\textasciicircum{}6 dispersion term.  The long\sphinxhyphen{}range portion can be calculated by
using the {\hyperref[\detokenize{kspace_style::doc}]{\sphinxcrossref{\DUrole{doc}{kspace\_style ewald/disp or pppm/disp}}}}
commands.  The specified Buckingham cutoff then determines which
portion of the Buckingham interactions are computed directly by the
pair potential versus which part is computed in reciprocal space via
the Kspace style.  If \sphinxstyleemphasis{flag\_buck} is set to \sphinxstyleemphasis{cut}, the Buckingham
interactions are simply cutoff, as with {\hyperref[\detokenize{pair_buck::doc}]{\sphinxcrossref{\DUrole{doc}{pair\_style buck}}}}.

\sphinxAtStartPar
If \sphinxstyleemphasis{flag\_coul} is set to \sphinxstyleemphasis{long}, no cutoff is used on the Coulombic
interactions.  The long\sphinxhyphen{}range portion can calculated by using any of
several {\hyperref[\detokenize{kspace_style::doc}]{\sphinxcrossref{\DUrole{doc}{kspace\_style}}}} command options such as
\sphinxstyleemphasis{pppm} or \sphinxstyleemphasis{ewald}.  Note that if \sphinxstyleemphasis{flag\_buck} is also set to long, then
the \sphinxstyleemphasis{ewald/disp} or \sphinxstyleemphasis{pppm/disp} Kspace style needs to be used to
perform the long\sphinxhyphen{}range calculations for both the Buckingham and
Coulombic interactions.  If \sphinxstyleemphasis{flag\_coul} is set to \sphinxstyleemphasis{off}, Coulombic
interactions are not computed.

\sphinxAtStartPar
The following coefficients must be defined for each pair of atoms
types via the {\hyperref[\detokenize{pair_coeff::doc}]{\sphinxcrossref{\DUrole{doc}{pair\_coeff}}}} command as in the examples
above, or in the data file or restart files read by the
{\hyperref[\detokenize{read_data::doc}]{\sphinxcrossref{\DUrole{doc}{read\_data}}}} or {\hyperref[\detokenize{read_restart::doc}]{\sphinxcrossref{\DUrole{doc}{read\_restart}}}}
commands:
\begin{itemize}
\item {} 
\sphinxAtStartPar
A (energy units)

\item {} 
\sphinxAtStartPar
rho (distance units)

\item {} 
\sphinxAtStartPar
C (energy\sphinxhyphen{}distance\textasciicircum{}6 units)

\item {} 
\sphinxAtStartPar
cutoff (distance units)

\item {} 
\sphinxAtStartPar
cutoff2 (distance units)

\end{itemize}

\sphinxAtStartPar
The second coefficient, rho, must be greater than zero.

\sphinxAtStartPar
The latter 2 coefficients are optional.  If not specified, the global
Buckingham and Coulombic cutoffs specified in the pair\_style command
are used.  If only one cutoff is specified, it is used as the cutoff
for both Buckingham and Coulombic interactions for this type pair.  If
both coefficients are specified, they are used as the Buckingham and
Coulombic cutoffs for this type pair.  Note that if you are using
\sphinxstyleemphasis{flag\_buck} set to \sphinxstyleemphasis{long}, you cannot specify a Buckingham cutoff for
an atom type pair, since only one global Buckingham cutoff is allowed.
Similarly, if you are using \sphinxstyleemphasis{flag\_coul} set to \sphinxstyleemphasis{long}, you cannot
specify a Coulombic cutoff for an atom type pair, since only one
global Coulombic cutoff is allowed.


\bigskip\hrule\bigskip


\sphinxAtStartPar
Styles with a \sphinxstyleemphasis{gpu}, \sphinxstyleemphasis{intel}, \sphinxstyleemphasis{kk}, \sphinxstyleemphasis{omp}, or \sphinxstyleemphasis{opt} suffix are
functionally the same as the corresponding style without the suffix.
They have been optimized to run faster, depending on your available
hardware, as discussed on the {\hyperref[\detokenize{Speed_packages::doc}]{\sphinxcrossref{\DUrole{doc}{Accelerator packages}}}}
page.  The accelerated styles take the same arguments and should
produce the same results, except for round\sphinxhyphen{}off and precision issues.

\sphinxAtStartPar
These accelerated styles are part of the GPU, INTEL, KOKKOS,
OPENMP, and OPT packages, respectively.  They are only enabled if
LAMMPS was built with those packages.  See the {\hyperref[\detokenize{Build_package::doc}]{\sphinxcrossref{\DUrole{doc}{Build package}}}} page for more info.

\sphinxAtStartPar
You can specify the accelerated styles explicitly in your input script
by including their suffix, or you can use the {\hyperref[\detokenize{Run_options::doc}]{\sphinxcrossref{\DUrole{doc}{\sphinxhyphen{}suffix command\sphinxhyphen{}line switch}}}} when you invoke LAMMPS, or you can use the
{\hyperref[\detokenize{suffix::doc}]{\sphinxcrossref{\DUrole{doc}{suffix}}}} command in your input script.

\sphinxAtStartPar
See the {\hyperref[\detokenize{Speed_packages::doc}]{\sphinxcrossref{\DUrole{doc}{Accelerator packages}}}} page for more
instructions on how to use the accelerated styles effectively.


\bigskip\hrule\bigskip



\subsection{Mixing, shift, table, tail correction, restart, rRESPA info}
\label{\detokenize{pair_buck_long:mixing-shift-table-tail-correction-restart-rrespa-info}}
\sphinxAtStartPar
This pair style does not support mixing.  Thus, coefficients for all
I,J pairs must be specified explicitly.

\sphinxAtStartPar
This pair style supports the {\hyperref[\detokenize{pair_modify::doc}]{\sphinxcrossref{\DUrole{doc}{pair\_modify}}}} shift
option for the energy of the exp() and 1/r\textasciicircum{}6 portion of the pair
interaction, assuming \sphinxstyleemphasis{flag\_buck} is \sphinxstyleemphasis{cut}.

\sphinxAtStartPar
This pair style does not support the {\hyperref[\detokenize{pair_modify::doc}]{\sphinxcrossref{\DUrole{doc}{pair\_modify}}}}
shift option for the energy of the Buckingham portion of the pair
interaction.

\sphinxAtStartPar
This pair style supports the {\hyperref[\detokenize{pair_modify::doc}]{\sphinxcrossref{\DUrole{doc}{pair\_modify}}}} table and
table/disp options since they can tabulate the short\sphinxhyphen{}range portion of
the long\sphinxhyphen{}range Coulombic and dispersion interactions.

\sphinxAtStartPar
This pair style write its information to {\hyperref[\detokenize{restart::doc}]{\sphinxcrossref{\DUrole{doc}{binary restart files}}}}, so pair\_style and pair\_coeff commands do not need
to be specified in an input script that reads a restart file.

\sphinxAtStartPar
This pair style supports the use of the \sphinxstyleemphasis{inner}, \sphinxstyleemphasis{middle}, and \sphinxstyleemphasis{outer}
keywords of the {\hyperref[\detokenize{run_style::doc}]{\sphinxcrossref{\DUrole{doc}{run\_style respa}}}} command, meaning the
pairwise forces can be partitioned by distance at different levels of
the rRESPA hierarchy.  See the {\hyperref[\detokenize{run_style::doc}]{\sphinxcrossref{\DUrole{doc}{run\_style}}}} command for
details.


\bigskip\hrule\bigskip



\subsection{Restrictions}
\label{\detokenize{pair_buck_long:restrictions}}
\sphinxAtStartPar
This style is part of the KSPACE package.  It is only enabled if
LAMMPS was built with that package.  See the {\hyperref[\detokenize{Build_package::doc}]{\sphinxcrossref{\DUrole{doc}{Build package}}}} page for more info.


\subsection{Related commands}
\label{\detokenize{pair_buck_long:related-commands}}
\sphinxAtStartPar
{\hyperref[\detokenize{pair_coeff::doc}]{\sphinxcrossref{\DUrole{doc}{pair\_coeff}}}}


\subsection{Default}
\label{\detokenize{pair_buck_long:default}}
\sphinxAtStartPar
none


\bigskip\hrule\bigskip

\phantomsection\label{\detokenize{pair_buck_long:ismail}}
\sphinxAtStartPar
\sphinxstylestrong{(Ismail)} Ismail, Tsige, In ‘t Veld, Grest, Molecular Physics
(accepted) (2007).

\sphinxstepscope

\index{pair\_style lj/charmm/coul/charmm@\spxentry{pair\_style lj/charmm/coul/charmm}}\ignorespaces 
\index{pair\_style lj/charmm/coul/charmm/gpu@\spxentry{pair\_style lj/charmm/coul/charmm/gpu}}\ignorespaces 
\index{pair\_style lj/charmm/coul/charmm/intel@\spxentry{pair\_style lj/charmm/coul/charmm/intel}}\ignorespaces 
\index{pair\_style lj/charmm/coul/charmm/kk@\spxentry{pair\_style lj/charmm/coul/charmm/kk}}\ignorespaces 
\index{pair\_style lj/charmm/coul/charmm/omp@\spxentry{pair\_style lj/charmm/coul/charmm/omp}}\ignorespaces 
\index{pair\_style lj/charmm/coul/charmm/implicit@\spxentry{pair\_style lj/charmm/coul/charmm/implicit}}\ignorespaces 
\index{pair\_style lj/charmm/coul/charmm/implicit/kk@\spxentry{pair\_style lj/charmm/coul/charmm/implicit/kk}}\ignorespaces 
\index{pair\_style lj/charmm/coul/charmm/implicit/omp@\spxentry{pair\_style lj/charmm/coul/charmm/implicit/omp}}\ignorespaces 
\index{pair\_style lj/charmm/coul/long@\spxentry{pair\_style lj/charmm/coul/long}}\ignorespaces 
\index{pair\_style lj/charmm/coul/long/gpu@\spxentry{pair\_style lj/charmm/coul/long/gpu}}\ignorespaces 
\index{pair\_style lj/charmm/coul/long/intel@\spxentry{pair\_style lj/charmm/coul/long/intel}}\ignorespaces 
\index{pair\_style lj/charmm/coul/long/kk@\spxentry{pair\_style lj/charmm/coul/long/kk}}\ignorespaces 
\index{pair\_style lj/charmm/coul/long/opt@\spxentry{pair\_style lj/charmm/coul/long/opt}}\ignorespaces 
\index{pair\_style lj/charmm/coul/long/omp@\spxentry{pair\_style lj/charmm/coul/long/omp}}\ignorespaces 
\index{pair\_style lj/charmm/coul/msm@\spxentry{pair\_style lj/charmm/coul/msm}}\ignorespaces 
\index{pair\_style lj/charmm/coul/msm/omp@\spxentry{pair\_style lj/charmm/coul/msm/omp}}\ignorespaces 
\index{pair\_style lj/charmmfsw/coul/charmmfsh@\spxentry{pair\_style lj/charmmfsw/coul/charmmfsh}}\ignorespaces 
\index{pair\_style lj/charmmfsw/coul/long@\spxentry{pair\_style lj/charmmfsw/coul/long}}\ignorespaces 
\index{pair\_style lj/charmmfsw/coul/long/kk@\spxentry{pair\_style lj/charmmfsw/coul/long/kk}}\ignorespaces 

\section{pair\_style lj/charmm/coul/charmm command}
\label{\detokenize{pair_charmm:pair-style-lj-charmm-coul-charmm-command}}\label{\detokenize{pair_charmm:index-18}}\label{\detokenize{pair_charmm:index-17}}\label{\detokenize{pair_charmm:index-16}}\label{\detokenize{pair_charmm:index-15}}\label{\detokenize{pair_charmm:index-14}}\label{\detokenize{pair_charmm:index-13}}\label{\detokenize{pair_charmm:index-12}}\label{\detokenize{pair_charmm:index-11}}\label{\detokenize{pair_charmm:index-10}}\label{\detokenize{pair_charmm:index-9}}\label{\detokenize{pair_charmm:index-8}}\label{\detokenize{pair_charmm:index-7}}\label{\detokenize{pair_charmm:index-6}}\label{\detokenize{pair_charmm:index-5}}\label{\detokenize{pair_charmm:index-4}}\label{\detokenize{pair_charmm:index-3}}\label{\detokenize{pair_charmm:index-2}}\label{\detokenize{pair_charmm:index-1}}\label{\detokenize{pair_charmm:index-0}}\label{\detokenize{pair_charmm::doc}}
\sphinxAtStartPar
Accelerator Variants: \sphinxstyleemphasis{lj/charmm/coul/charmm/gpu}, \sphinxstyleemphasis{lj/charmm/coul/charmm/intel}, \sphinxstyleemphasis{lj/charmm/coul/charmm/kk}, \sphinxstyleemphasis{lj/charmm/coul/charmm/omp}


\section{pair\_style lj/charmm/coul/charmm/implicit command}
\label{\detokenize{pair_charmm:pair-style-lj-charmm-coul-charmm-implicit-command}}
\sphinxAtStartPar
Accelerator Variants: \sphinxstyleemphasis{lj/charmm/coul/charmm/implicit/kk}, \sphinxstyleemphasis{lj/charmm/coul/charmm/implicit/omp}


\section{pair\_style lj/charmm/coul/long command}
\label{\detokenize{pair_charmm:pair-style-lj-charmm-coul-long-command}}
\sphinxAtStartPar
Accelerator Variants: \sphinxstyleemphasis{lj/charmm/coul/long/gpu}, \sphinxstyleemphasis{lj/charmm/coul/long/intel}, \sphinxstyleemphasis{lj/charmm/coul/long/kk}, \sphinxstyleemphasis{lj/charmm/coul/long/opt}, \sphinxstyleemphasis{lj/charmm/coul/long/omp}


\section{pair\_style lj/charmm/coul/msm command}
\label{\detokenize{pair_charmm:pair-style-lj-charmm-coul-msm-command}}
\sphinxAtStartPar
Accelerator Variants: \sphinxstyleemphasis{lj/charmm/coul/msm/omp}


\section{pair\_style lj/charmmfsw/coul/charmmfsh command}
\label{\detokenize{pair_charmm:pair-style-lj-charmmfsw-coul-charmmfsh-command}}

\section{pair\_style lj/charmmfsw/coul/long command}
\label{\detokenize{pair_charmm:pair-style-lj-charmmfsw-coul-long-command}}
\sphinxAtStartPar
Accelerator Variants: \sphinxstyleemphasis{lj/charmmfsw/coul/long/kk}


\subsection{Syntax}
\label{\detokenize{pair_charmm:syntax}}
\begin{sphinxVerbatim}[commandchars=\\\{\}]
\PYG{k}{pair\PYGZus{}style}\PYG{+w}{ }\PYG{n}{style}\PYG{+w}{ }\PYG{n}{args}
\end{sphinxVerbatim}
\begin{itemize}
\item {} 
\sphinxAtStartPar
style = \sphinxstyleemphasis{lj/charmm/coul/charmm} or \sphinxstyleemphasis{lj/charmm/coul/charmm/implicit} or \sphinxstyleemphasis{lj/charmm/coul/long} or \sphinxstyleemphasis{lj/charmm/coul/msm} or \sphinxstyleemphasis{lj/charmmfsw/coul/charmmfsh} or \sphinxstyleemphasis{lj/charmmfsw/coul/long}

\item {} 
\sphinxAtStartPar
args = list of arguments for a particular style

\end{itemize}
\begin{sphinxalltt}
\sphinxstyleemphasis{lj/charmm/coul/charmm} args = inner outer (inner2) (outer2)
  inner, outer = global switching cutoffs for Lennard Jones (and Coulombic if only 2 args)
  inner2, outer2 = global switching cutoffs for Coulombic (optional)
\sphinxstyleemphasis{lj/charmm/coul/charmm/implicit} args = inner outer (inner2) (outer2)
  inner, outer = global switching cutoffs for LJ (and Coulombic if only 2 args)
  inner2, outer2 = global switching cutoffs for Coulombic (optional)
\sphinxstyleemphasis{lj/charmm/coul/long} args = inner outer (cutoff)
  inner, outer = global switching cutoffs for LJ (and Coulombic if only 2 args)
  cutoff = global cutoff for Coulombic (optional, outer is Coulombic cutoff if only 2 args)
\sphinxstyleemphasis{lj/charmm/coul/msm} args = inner outer (cutoff)
  inner, outer = global switching cutoffs for LJ (and Coulombic if only 2 args)
  cutoff = global cutoff for Coulombic (optional, outer is Coulombic cutoff if only 2 args)
\sphinxstyleemphasis{lj/charmmfsw/coul/charmmfsh} args = inner outer (cutoff)
  inner, outer = global cutoffs for LJ (and Coulombic if only 2 args)
  cutoff = global cutoff for Coulombic (optional, outer is Coulombic cutoff if only 2 args)
\sphinxstyleemphasis{lj/charmmfsw/coul/long} args = inner outer (cutoff)
  inner, outer = global cutoffs for LJ (and Coulombic if only 2 args)
  cutoff = global cutoff for Coulombic (optional, outer is Coulombic cutoff if only 2 args)
\end{sphinxalltt}


\subsection{Examples}
\label{\detokenize{pair_charmm:examples}}
\begin{sphinxVerbatim}[commandchars=\\\{\}]
\PYG{k}{pair\PYGZus{}style}\PYG{+w}{ }\PYG{n}{lj}\PYG{o}{/}\PYG{n}{charmm}\PYG{o}{/}\PYG{n}{coul}\PYG{o}{/}\PYG{n}{charmm}\PYG{+w}{ }\PYG{l+m}{8.0}\PYG{+w}{ }\PYG{l+m}{10.0}
\PYG{k}{pair\PYGZus{}style}\PYG{+w}{ }\PYG{n}{lj}\PYG{o}{/}\PYG{n}{charmm}\PYG{o}{/}\PYG{n}{coul}\PYG{o}{/}\PYG{n}{charmm}\PYG{+w}{ }\PYG{l+m}{8.0}\PYG{+w}{ }\PYG{l+m}{10.0}\PYG{+w}{ }\PYG{l+m}{7.0}\PYG{+w}{ }\PYG{l+m}{9.0}
\PYG{k}{pair\PYGZus{}style}\PYG{+w}{ }\PYG{n}{lj}\PYG{o}{/}\PYG{n}{charmmfsw}\PYG{o}{/}\PYG{n}{coul}\PYG{o}{/}\PYG{n}{charmmfsh}\PYG{+w}{ }\PYG{l+m}{10.0}\PYG{+w}{ }\PYG{l+m}{12.0}
\PYG{k}{pair\PYGZus{}style}\PYG{+w}{ }\PYG{n}{lj}\PYG{o}{/}\PYG{n}{charmmfsw}\PYG{o}{/}\PYG{n}{coul}\PYG{o}{/}\PYG{n}{charmmfsh}\PYG{+w}{ }\PYG{l+m}{10.0}\PYG{+w}{ }\PYG{l+m}{12.0}\PYG{+w}{ }\PYG{l+m}{9.0}
\PYG{k}{pair\PYGZus{}coeff}\PYG{+w}{ }\PYG{o}{*}\PYG{+w}{ }\PYG{o}{*}\PYG{+w}{ }\PYG{l+m}{100.0}\PYG{+w}{ }\PYG{l+m}{2.0}
\PYG{k}{pair\PYGZus{}coeff}\PYG{+w}{ }\PYG{l+m}{1}\PYG{+w}{ }\PYG{l+m}{1}\PYG{+w}{ }\PYG{l+m}{100.0}\PYG{+w}{ }\PYG{l+m}{2.0}\PYG{+w}{ }\PYG{l+m}{150.0}\PYG{+w}{ }\PYG{l+m}{3.5}

\PYG{k}{pair\PYGZus{}style}\PYG{+w}{ }\PYG{n}{lj}\PYG{o}{/}\PYG{n}{charmm}\PYG{o}{/}\PYG{n}{coul}\PYG{o}{/}\PYG{n}{charmm}\PYG{o}{/}\PYG{n}{implicit}\PYG{+w}{ }\PYG{l+m}{8.0}\PYG{+w}{ }\PYG{l+m}{10.0}
\PYG{k}{pair\PYGZus{}style}\PYG{+w}{ }\PYG{n}{lj}\PYG{o}{/}\PYG{n}{charmm}\PYG{o}{/}\PYG{n}{coul}\PYG{o}{/}\PYG{n}{charmm}\PYG{o}{/}\PYG{n}{implicit}\PYG{+w}{ }\PYG{l+m}{8.0}\PYG{+w}{ }\PYG{l+m}{10.0}\PYG{+w}{ }\PYG{l+m}{7.0}\PYG{+w}{ }\PYG{l+m}{9.0}
\PYG{k}{pair\PYGZus{}coeff}\PYG{+w}{ }\PYG{o}{*}\PYG{+w}{ }\PYG{o}{*}\PYG{+w}{ }\PYG{l+m}{100.0}\PYG{+w}{ }\PYG{l+m}{2.0}
\PYG{k}{pair\PYGZus{}coeff}\PYG{+w}{ }\PYG{l+m}{1}\PYG{+w}{ }\PYG{l+m}{1}\PYG{+w}{ }\PYG{l+m}{100.0}\PYG{+w}{ }\PYG{l+m}{2.0}\PYG{+w}{ }\PYG{l+m}{150.0}\PYG{+w}{ }\PYG{l+m}{3.5}

\PYG{k}{pair\PYGZus{}style}\PYG{+w}{ }\PYG{n}{lj}\PYG{o}{/}\PYG{n}{charmm}\PYG{o}{/}\PYG{n}{coul}\PYG{o}{/}\PYG{n}{long}\PYG{+w}{ }\PYG{l+m}{8.0}\PYG{+w}{ }\PYG{l+m}{10.0}
\PYG{k}{pair\PYGZus{}style}\PYG{+w}{ }\PYG{n}{lj}\PYG{o}{/}\PYG{n}{charmm}\PYG{o}{/}\PYG{n}{coul}\PYG{o}{/}\PYG{n}{long}\PYG{+w}{ }\PYG{l+m}{8.0}\PYG{+w}{ }\PYG{l+m}{10.0}\PYG{+w}{ }\PYG{l+m}{9.0}
\PYG{k}{pair\PYGZus{}style}\PYG{+w}{ }\PYG{n}{lj}\PYG{o}{/}\PYG{n}{charmmfsw}\PYG{o}{/}\PYG{n}{coul}\PYG{o}{/}\PYG{n}{long}\PYG{+w}{ }\PYG{l+m}{8.0}\PYG{+w}{ }\PYG{l+m}{10.0}
\PYG{k}{pair\PYGZus{}style}\PYG{+w}{ }\PYG{n}{lj}\PYG{o}{/}\PYG{n}{charmmfsw}\PYG{o}{/}\PYG{n}{coul}\PYG{o}{/}\PYG{n}{long}\PYG{+w}{ }\PYG{l+m}{8.0}\PYG{+w}{ }\PYG{l+m}{10.0}\PYG{+w}{ }\PYG{l+m}{9.0}
\PYG{k}{pair\PYGZus{}coeff}\PYG{+w}{ }\PYG{o}{*}\PYG{+w}{ }\PYG{o}{*}\PYG{+w}{ }\PYG{l+m}{100.0}\PYG{+w}{ }\PYG{l+m}{2.0}
\PYG{k}{pair\PYGZus{}coeff}\PYG{+w}{ }\PYG{l+m}{1}\PYG{+w}{ }\PYG{l+m}{1}\PYG{+w}{ }\PYG{l+m}{100.0}\PYG{+w}{ }\PYG{l+m}{2.0}\PYG{+w}{ }\PYG{l+m}{150.0}\PYG{+w}{ }\PYG{l+m}{3.5}

\PYG{k}{pair\PYGZus{}style}\PYG{+w}{ }\PYG{n}{lj}\PYG{o}{/}\PYG{n}{charmm}\PYG{o}{/}\PYG{n}{coul}\PYG{o}{/}\PYG{n}{msm}\PYG{+w}{ }\PYG{l+m}{8.0}\PYG{+w}{ }\PYG{l+m}{10.0}
\PYG{k}{pair\PYGZus{}style}\PYG{+w}{ }\PYG{n}{lj}\PYG{o}{/}\PYG{n}{charmm}\PYG{o}{/}\PYG{n}{coul}\PYG{o}{/}\PYG{n}{msm}\PYG{+w}{ }\PYG{l+m}{8.0}\PYG{+w}{ }\PYG{l+m}{10.0}\PYG{+w}{ }\PYG{l+m}{9.0}
\PYG{k}{pair\PYGZus{}coeff}\PYG{+w}{ }\PYG{o}{*}\PYG{+w}{ }\PYG{o}{*}\PYG{+w}{ }\PYG{l+m}{100.0}\PYG{+w}{ }\PYG{l+m}{2.0}
\PYG{k}{pair\PYGZus{}coeff}\PYG{+w}{ }\PYG{l+m}{1}\PYG{+w}{ }\PYG{l+m}{1}\PYG{+w}{ }\PYG{l+m}{100.0}\PYG{+w}{ }\PYG{l+m}{2.0}\PYG{+w}{ }\PYG{l+m}{150.0}\PYG{+w}{ }\PYG{l+m}{3.5}
\end{sphinxVerbatim}


\subsection{Description}
\label{\detokenize{pair_charmm:description}}
\sphinxAtStartPar
These pair styles compute Lennard Jones (LJ) and Coulombic
interactions with additional switching or shifting functions that ramp
the energy and/or force smoothly to zero between an inner and outer
cutoff.  They implement the widely used CHARMM force field, see
{\hyperref[\detokenize{Howto_bioFF::doc}]{\sphinxcrossref{\DUrole{doc}{Howto discussion on biomolecular force fields}}}} for
details.

\sphinxAtStartPar
The styles with \sphinxstyleemphasis{charmm} (not \sphinxstyleemphasis{charmmfsw} or \sphinxstyleemphasis{charmmfsh}) in their
name are the older, original LAMMPS implementations.  They compute the
LJ and Coulombic interactions with an energy switching function which
ramps the energy smoothly to zero between the inner and outer cutoff.
This can cause irregularities in pairwise forces (due to the discontinuous
second derivative of energy at the boundaries of the switching region),
which in some cases can result in detectable artifacts in an MD simulation.

\sphinxAtStartPar
The newer styles with \sphinxstyleemphasis{charmmfsw} or \sphinxstyleemphasis{charmmfsh} in their name replace
the energy switching with force switching (fsw) and force shifting
(fsh) functions, for LJ and Coulombic interactions respectively.

\begin{sphinxadmonition}{note}{Note:}
\sphinxAtStartPar
The newer \sphinxstyleemphasis{charmmfsw} or \sphinxstyleemphasis{charmmfsh} styles were released in
March 2017.  We recommend they be used instead of the older \sphinxstyleemphasis{charmm}
styles.  This includes the newer {\hyperref[\detokenize{dihedral_charmm::doc}]{\sphinxcrossref{\DUrole{doc}{dihedral\_style charmmfsw}}}} command.  Eventually code from the new
styles will propagate into the related pair styles (e.g. implicit,
accelerator, free energy variants).
\end{sphinxadmonition}

\begin{sphinxadmonition}{note}{Note:}
\sphinxAtStartPar
The newest CHARMM pair styles reset the Coulombic energy
conversion factor used internally in the code, from the LAMMPS value
to the CHARMM value, as if it were effectively a parameter of the
force field.  This is because the CHARMM code uses a slightly
different value for the this conversion factor in {\hyperref[\detokenize{units::doc}]{\sphinxcrossref{\DUrole{doc}{real units}}}} (kcal/mol), namely CHARMM = 332.0716, LAMMPS =
332.06371.  This is to enable more precise agreement by LAMMPS with
the CHARMM force field energies and forces, when using one of these
two CHARMM pair styles.
\end{sphinxadmonition}

\sphinxAtStartPar
When using the \sphinxstyleemphasis{lj/charmm/coul/charmm styles}, both the LJ and
Coulombic terms require an inner and outer cutoff. They can be the
same for both formulas or different depending on whether 2 or 4
arguments are used in the pair\_style command.  For the
\sphinxstyleemphasis{lj/charmmfsw/coul/charmmfsh} style, the LJ term requires both an
inner and outer cutoff, while the Coulombic term requires only one
cutoff.  If the Coulombic cutoff is not specified (2 instead of 3
arguments), the LJ outer cutoff is used for the Coulombic cutoff.  In
all cases where an inner and outer cutoff are specified, the inner
cutoff distance must be less than the outer cutoff.  It is typical to
make the difference between the inner and outer cutoffs about 2.0
Angstroms.

\sphinxAtStartPar
Style \sphinxstyleemphasis{lj/charmm/coul/charmm/implicit} computes the same formulas as
style \sphinxstyleemphasis{lj/charmm/coul/charmm} except that an additional 1/r term is
included in the Coulombic formula.  The Coulombic energy thus varies
as 1/r\textasciicircum{}2.  This is effectively a distance\sphinxhyphen{}dependent dielectric term
which is a simple model for an implicit solvent with additional
screening.  It is designed for use in a simulation of an unsolvated
biomolecule (no explicit water molecules).

\sphinxAtStartPar
Styles \sphinxstyleemphasis{lj/charmm/coul/long} and \sphinxstyleemphasis{lj/charmm/coul/msm} compute the same
formulas as style \sphinxstyleemphasis{lj/charmm/coul/charmm} and style
\sphinxstyleemphasis{lj/charmmfsw/coul/long} computes the same formulas as style
\sphinxstyleemphasis{lj/charmmfsw/coul/charmmfsh}, except that an additional damping
factor is applied to the Coulombic term, so it can be used in
conjunction with the {\hyperref[\detokenize{kspace_style::doc}]{\sphinxcrossref{\DUrole{doc}{kspace\_style}}}} command and its
\sphinxstyleemphasis{ewald} or \sphinxstyleemphasis{pppm} or \sphinxstyleemphasis{msm} option.  Only one Coulombic cutoff is
specified for these styles; if only 2 arguments are used in the
pair\_style command, then the outer LJ cutoff is used as the single
Coulombic cutoff.  The Coulombic cutoff specified for these styles
means that pairwise interactions within this distance are computed
directly; interactions outside that distance are computed in
reciprocal space.

\sphinxAtStartPar
The following coefficients must be defined for each pair of atoms
types via the {\hyperref[\detokenize{pair_coeff::doc}]{\sphinxcrossref{\DUrole{doc}{pair\_coeff}}}} command as in the examples
above, or in the data file or restart files read by the
{\hyperref[\detokenize{read_data::doc}]{\sphinxcrossref{\DUrole{doc}{read\_data}}}} or {\hyperref[\detokenize{read_restart::doc}]{\sphinxcrossref{\DUrole{doc}{read\_restart}}}}
commands, or by mixing as described below:
\begin{itemize}
\item {} 
\sphinxAtStartPar
\(\epsilon\) (energy units)

\item {} 
\sphinxAtStartPar
\(\sigma\) (distance units)

\item {} 
\sphinxAtStartPar
\(\epsilon_{14}\) (energy units)

\item {} 
\sphinxAtStartPar
\(\sigma_{14}\) (distance units)

\end{itemize}

\sphinxAtStartPar
Note that \(\sigma\) is defined in the LJ formula as the zero\sphinxhyphen{}crossing
distance for the potential, not as the energy minimum at \(2^{1/6} \sigma\).

\sphinxAtStartPar
The latter 2 coefficients are optional.  If they are specified, they
are used in the LJ formula between two atoms of these types which are
also first and fourth atoms in any dihedral.  No cutoffs are specified
because the CHARMM force field does not allow varying cutoffs for
individual atom pairs; all pairs use the global cutoff(s) specified in
the pair\_style command.


\bigskip\hrule\bigskip


\sphinxAtStartPar
Styles with a \sphinxstyleemphasis{gpu}, \sphinxstyleemphasis{intel}, \sphinxstyleemphasis{kk}, \sphinxstyleemphasis{omp}, or \sphinxstyleemphasis{opt} suffix are
functionally the same as the corresponding style without the suffix.
They have been optimized to run faster, depending on your available
hardware, as discussed on the {\hyperref[\detokenize{Speed_packages::doc}]{\sphinxcrossref{\DUrole{doc}{Accelerator packages}}}}
page.  The accelerated styles take the same arguments and should
produce the same results, except for round\sphinxhyphen{}off and precision issues.

\sphinxAtStartPar
These accelerated styles are part of the GPU, INTEL, KOKKOS,
OPENMP, and OPT packages, respectively.  They are only enabled if
LAMMPS was built with those packages.  See the {\hyperref[\detokenize{Build_package::doc}]{\sphinxcrossref{\DUrole{doc}{Build package}}}} page for more info.

\sphinxAtStartPar
You can specify the accelerated styles explicitly in your input script
by including their suffix, or you can use the {\hyperref[\detokenize{Run_options::doc}]{\sphinxcrossref{\DUrole{doc}{\sphinxhyphen{}suffix command\sphinxhyphen{}line switch}}}} when you invoke LAMMPS, or you can use the
{\hyperref[\detokenize{suffix::doc}]{\sphinxcrossref{\DUrole{doc}{suffix}}}} command in your input script.

\sphinxAtStartPar
See the {\hyperref[\detokenize{Speed_packages::doc}]{\sphinxcrossref{\DUrole{doc}{Accelerator packages}}}} page for more
instructions on how to use the accelerated styles effectively.


\bigskip\hrule\bigskip



\subsection{Mixing, shift, table, tail correction, restart, rRESPA info}
\label{\detokenize{pair_charmm:mixing-shift-table-tail-correction-restart-rrespa-info}}
\sphinxAtStartPar
For atom type pairs I,J and I != J, the epsilon, sigma, epsilon\_14,
and sigma\_14 coefficients for all of the lj/charmm pair styles can be
mixed.  The default mix value is \sphinxstyleemphasis{arithmetic} to coincide with the
usual settings for the CHARMM force field.  See the “pair\_modify”
command for details.

\sphinxAtStartPar
None of the \sphinxstyleemphasis{lj/charmm} or \sphinxstyleemphasis{lj/charmmfsw} pair styles support the
{\hyperref[\detokenize{pair_modify::doc}]{\sphinxcrossref{\DUrole{doc}{pair\_modify}}}} shift option, since the Lennard\sphinxhyphen{}Jones
portion of the pair interaction is smoothed to 0.0 at the cutoff.

\sphinxAtStartPar
The \sphinxstyleemphasis{lj/charmm/coul/long} and \sphinxstyleemphasis{lj/charmmfsw/coul/long} styles support
the {\hyperref[\detokenize{pair_modify::doc}]{\sphinxcrossref{\DUrole{doc}{pair\_modify}}}} table option since they can
tabulate the short\sphinxhyphen{}range portion of the long\sphinxhyphen{}range Coulombic
interaction.

\sphinxAtStartPar
None of the \sphinxstyleemphasis{lj/charmm} or \sphinxstyleemphasis{lj/charmmfsw} pair styles support the
{\hyperref[\detokenize{pair_modify::doc}]{\sphinxcrossref{\DUrole{doc}{pair\_modify}}}} tail option for adding long\sphinxhyphen{}range tail
corrections to energy and pressure, since the Lennard\sphinxhyphen{}Jones portion of
the pair interaction is smoothed to 0.0 at the cutoff.

\sphinxAtStartPar
All of the \sphinxstyleemphasis{lj/charmm} and \sphinxstyleemphasis{lj/charmmfsw} pair styles write their
information to {\hyperref[\detokenize{restart::doc}]{\sphinxcrossref{\DUrole{doc}{binary restart files}}}}, so pair\_style and
pair\_coeff commands do not need to be specified in an input script
that reads a restart file.

\sphinxAtStartPar
The \sphinxstyleemphasis{lj/charmm/coul/long} and \sphinxstyleemphasis{lj/charmmfsw/coul/long} pair styles
support the use of the \sphinxstyleemphasis{inner}, \sphinxstyleemphasis{middle}, and \sphinxstyleemphasis{outer} keywords of the
{\hyperref[\detokenize{run_style::doc}]{\sphinxcrossref{\DUrole{doc}{run\_style respa}}}} command, meaning the pairwise forces
can be partitioned by distance at different levels of the rRESPA
hierarchy.  The other styles only support the \sphinxstyleemphasis{pair} keyword of
run\_style respa.  See the {\hyperref[\detokenize{run_style::doc}]{\sphinxcrossref{\DUrole{doc}{run\_style}}}} command for
details.


\bigskip\hrule\bigskip



\subsection{Restrictions}
\label{\detokenize{pair_charmm:restrictions}}
\sphinxAtStartPar
All the styles with \sphinxstyleemphasis{coul/charmm} or \sphinxstyleemphasis{coul/charmmfsh} styles are part
of the MOLECULE package.  All the styles with \sphinxstyleemphasis{coul/long} style are
part of the KSPACE package.  They are only enabled if LAMMPS was built
with those packages.  See the {\hyperref[\detokenize{Build_package::doc}]{\sphinxcrossref{\DUrole{doc}{Build package}}}} doc
page for more info.


\subsection{Related commands}
\label{\detokenize{pair_charmm:related-commands}}
\sphinxAtStartPar
{\hyperref[\detokenize{pair_coeff::doc}]{\sphinxcrossref{\DUrole{doc}{pair\_coeff}}}}, {\hyperref[\detokenize{angle_charmm::doc}]{\sphinxcrossref{\DUrole{doc}{angle\_style charmm}}}},
{\hyperref[\detokenize{dihedral_charmm::doc}]{\sphinxcrossref{\DUrole{doc}{dihedral\_style charmm}}}},
{\hyperref[\detokenize{dihedral_charmm::doc}]{\sphinxcrossref{\DUrole{doc}{dihedral\_style charmmfsw}}}}, {\hyperref[\detokenize{fix_cmap::doc}]{\sphinxcrossref{\DUrole{doc}{fix cmap}}}}


\subsection{Default}
\label{\detokenize{pair_charmm:default}}
\sphinxAtStartPar
none


\bigskip\hrule\bigskip

\phantomsection\label{\detokenize{pair_charmm:brooks1}}
\sphinxAtStartPar
\sphinxstylestrong{(Brooks)} Brooks, et al, J Comput Chem, 30, 1545 (2009).

\phantomsection\label{\detokenize{pair_charmm:pair-mackerell}}
\sphinxAtStartPar
\sphinxstylestrong{(MacKerell)} MacKerell, Bashford, Bellott, Dunbrack, Evanseck, Field,
Fischer, Gao, Guo, Ha, et al, J Phys Chem, 102, 3586 (1998).

\phantomsection\label{\detokenize{pair_charmm:steinbach}}
\sphinxAtStartPar
\sphinxstylestrong{(Steinbach)} Steinbach, Brooks, J Comput Chem, 15, 667 (1994).

\sphinxstepscope

\index{pair\_style lj/class2@\spxentry{pair\_style lj/class2}}\ignorespaces 
\index{pair\_style lj/class2/gpu@\spxentry{pair\_style lj/class2/gpu}}\ignorespaces 
\index{pair\_style lj/class2/kk@\spxentry{pair\_style lj/class2/kk}}\ignorespaces 
\index{pair\_style lj/class2/omp@\spxentry{pair\_style lj/class2/omp}}\ignorespaces 
\index{pair\_style lj/class2/coul/cut@\spxentry{pair\_style lj/class2/coul/cut}}\ignorespaces 
\index{pair\_style lj/class2/coul/cut/kk@\spxentry{pair\_style lj/class2/coul/cut/kk}}\ignorespaces 
\index{pair\_style lj/class2/coul/cut/omp@\spxentry{pair\_style lj/class2/coul/cut/omp}}\ignorespaces 
\index{pair\_style lj/class2/coul/long@\spxentry{pair\_style lj/class2/coul/long}}\ignorespaces 
\index{pair\_style lj/class2/coul/long/gpu@\spxentry{pair\_style lj/class2/coul/long/gpu}}\ignorespaces 
\index{pair\_style lj/class2/coul/long/kk@\spxentry{pair\_style lj/class2/coul/long/kk}}\ignorespaces 
\index{pair\_style lj/class2/coul/long/omp@\spxentry{pair\_style lj/class2/coul/long/omp}}\ignorespaces 

\section{pair\_style lj/class2 command}
\label{\detokenize{pair_class2:pair-style-lj-class2-command}}\label{\detokenize{pair_class2:index-10}}\label{\detokenize{pair_class2:index-9}}\label{\detokenize{pair_class2:index-8}}\label{\detokenize{pair_class2:index-7}}\label{\detokenize{pair_class2:index-6}}\label{\detokenize{pair_class2:index-5}}\label{\detokenize{pair_class2:index-4}}\label{\detokenize{pair_class2:index-3}}\label{\detokenize{pair_class2:index-2}}\label{\detokenize{pair_class2:index-1}}\label{\detokenize{pair_class2:index-0}}\label{\detokenize{pair_class2::doc}}
\sphinxAtStartPar
Accelerator Variants: \sphinxstyleemphasis{lj/class2/gpu}, \sphinxstyleemphasis{lj/class2/kk}, \sphinxstyleemphasis{lj/class2/omp}


\section{pair\_style lj/class2/coul/cut command}
\label{\detokenize{pair_class2:pair-style-lj-class2-coul-cut-command}}
\sphinxAtStartPar
Accelerator Variants: \sphinxstyleemphasis{lj/class2/coul/cut/kk}, \sphinxstyleemphasis{lj/class2/coul/cut/omp}


\section{pair\_style lj/class2/coul/long command}
\label{\detokenize{pair_class2:pair-style-lj-class2-coul-long-command}}
\sphinxAtStartPar
Accelerator Variants: \sphinxstyleemphasis{lj/class2/coul/long/gpu}, \sphinxstyleemphasis{lj/class2/coul/long/kk}, \sphinxstyleemphasis{lj/class2/coul/long/omp}


\subsection{Syntax}
\label{\detokenize{pair_class2:syntax}}
\begin{sphinxVerbatim}[commandchars=\\\{\}]
\PYG{k}{pair\PYGZus{}style}\PYG{+w}{ }\PYG{n}{style}\PYG{+w}{ }\PYG{n}{args}
\end{sphinxVerbatim}
\begin{itemize}
\item {} 
\sphinxAtStartPar
style = \sphinxstyleemphasis{lj/class2} or \sphinxstyleemphasis{lj/class2/coul/cut} or \sphinxstyleemphasis{lj/class2/coul/long}

\item {} 
\sphinxAtStartPar
args = list of arguments for a particular style

\end{itemize}
\begin{sphinxalltt}
\sphinxstyleemphasis{lj/class2} args = cutoff
  cutoff = global cutoff for class 2 interactions (distance units)
\sphinxstyleemphasis{lj/class2/coul/cut} args = cutoff (cutoff2)
  cutoff = global cutoff for class 2 (and Coulombic if only 1 arg) (distance units)
  cutoff2 = global cutoff for Coulombic (optional) (distance units)
\sphinxstyleemphasis{lj/class2/coul/long} args = cutoff (cutoff2)
  cutoff = global cutoff for class 2 (and Coulombic if only 1 arg) (distance units)
  cutoff2 = global cutoff for Coulombic (optional) (distance units)
\end{sphinxalltt}


\subsection{Examples}
\label{\detokenize{pair_class2:examples}}
\begin{sphinxVerbatim}[commandchars=\\\{\}]
\PYG{k}{pair\PYGZus{}style}\PYG{+w}{ }\PYG{n}{lj}\PYG{o}{/}\PYG{n}{class2}\PYG{+w}{ }\PYG{l+m}{10.0}
\PYG{k}{pair\PYGZus{}coeff}\PYG{+w}{ }\PYG{o}{*}\PYG{+w}{ }\PYG{o}{*}\PYG{+w}{ }\PYG{l+m}{100.0}\PYG{+w}{ }\PYG{l+m}{2.5}
\PYG{k}{pair\PYGZus{}coeff}\PYG{+w}{ }\PYG{l+m}{1}\PYG{+w}{ }\PYG{l+m}{2}\PYG{o}{*}\PYG{+w}{ }\PYG{l+m}{100.0}\PYG{+w}{ }\PYG{l+m}{2.5}\PYG{+w}{ }\PYG{l+m}{9.0}

\PYG{k}{pair\PYGZus{}style}\PYG{+w}{ }\PYG{n}{lj}\PYG{o}{/}\PYG{n}{class2}\PYG{o}{/}\PYG{n}{coul}\PYG{o}{/}\PYG{n}{cut}\PYG{+w}{ }\PYG{l+m}{10.0}
\PYG{k}{pair\PYGZus{}style}\PYG{+w}{ }\PYG{n}{lj}\PYG{o}{/}\PYG{n}{class2}\PYG{o}{/}\PYG{n}{coul}\PYG{o}{/}\PYG{n}{cut}\PYG{+w}{ }\PYG{l+m}{10.0}\PYG{+w}{ }\PYG{l+m}{8.0}
\PYG{k}{pair\PYGZus{}coeff}\PYG{+w}{ }\PYG{o}{*}\PYG{+w}{ }\PYG{o}{*}\PYG{+w}{ }\PYG{l+m}{100.0}\PYG{+w}{ }\PYG{l+m}{3.0}
\PYG{k}{pair\PYGZus{}coeff}\PYG{+w}{ }\PYG{l+m}{1}\PYG{+w}{ }\PYG{l+m}{1}\PYG{+w}{ }\PYG{l+m}{100.0}\PYG{+w}{ }\PYG{l+m}{3.5}\PYG{+w}{ }\PYG{l+m}{9.0}
\PYG{k}{pair\PYGZus{}coeff}\PYG{+w}{ }\PYG{l+m}{1}\PYG{+w}{ }\PYG{l+m}{1}\PYG{+w}{ }\PYG{l+m}{100.0}\PYG{+w}{ }\PYG{l+m}{3.5}\PYG{+w}{ }\PYG{l+m}{9.0}\PYG{+w}{ }\PYG{l+m}{9.0}

\PYG{k}{pair\PYGZus{}style}\PYG{+w}{ }\PYG{n}{lj}\PYG{o}{/}\PYG{n}{class2}\PYG{o}{/}\PYG{n}{coul}\PYG{o}{/}\PYG{n}{long}\PYG{+w}{ }\PYG{l+m}{10.0}
\PYG{k}{pair\PYGZus{}style}\PYG{+w}{ }\PYG{n}{lj}\PYG{o}{/}\PYG{n}{class2}\PYG{o}{/}\PYG{n}{coul}\PYG{o}{/}\PYG{n}{long}\PYG{+w}{ }\PYG{l+m}{10.0}\PYG{+w}{ }\PYG{l+m}{8.0}
\PYG{k}{pair\PYGZus{}coeff}\PYG{+w}{ }\PYG{o}{*}\PYG{+w}{ }\PYG{o}{*}\PYG{+w}{ }\PYG{l+m}{100.0}\PYG{+w}{ }\PYG{l+m}{3.0}
\PYG{k}{pair\PYGZus{}coeff}\PYG{+w}{ }\PYG{l+m}{1}\PYG{+w}{ }\PYG{l+m}{1}\PYG{+w}{ }\PYG{l+m}{100.0}\PYG{+w}{ }\PYG{l+m}{3.5}\PYG{+w}{ }\PYG{l+m}{9.0}
\end{sphinxVerbatim}


\subsection{Description}
\label{\detokenize{pair_class2:description}}
\sphinxAtStartPar
The \sphinxstyleemphasis{lj/class2} styles compute a 6/9 Lennard\sphinxhyphen{}Jones potential given by
\begin{equation*}
\begin{split}E = \epsilon \left[ 2 \left(\frac{\sigma}{r}\right)^9 -
  3 \left(\frac{\sigma}{r}\right)^6 \right]
\qquad r < r_c\end{split}
\end{equation*}
\sphinxAtStartPar
\(r_c\) is the cutoff.

\sphinxAtStartPar
The \sphinxstyleemphasis{lj/class2/coul/cut} and \sphinxstyleemphasis{lj/class2/coul/long} styles add a
Coulombic term as described for the {\hyperref[\detokenize{pair_lj::doc}]{\sphinxcrossref{\DUrole{doc}{lj/cut}}}} pair styles.

\sphinxAtStartPar
See {\hyperref[\detokenize{pair_class2:pair-sun}]{\sphinxcrossref{\DUrole{std,std-ref}{(Sun)}}}} for a description of the COMPASS class2 force field.

\sphinxAtStartPar
The following coefficients must be defined for each pair of atoms
types via the {\hyperref[\detokenize{pair_coeff::doc}]{\sphinxcrossref{\DUrole{doc}{pair\_coeff}}}} command as in the examples
above, or in the data file or restart files read by the
{\hyperref[\detokenize{read_data::doc}]{\sphinxcrossref{\DUrole{doc}{read\_data}}}} or {\hyperref[\detokenize{read_restart::doc}]{\sphinxcrossref{\DUrole{doc}{read\_restart}}}}
commands, or by mixing as described below:
\begin{itemize}
\item {} 
\sphinxAtStartPar
\(\epsilon\) (energy units)

\item {} 
\sphinxAtStartPar
\(\sigma\) (distance units)

\item {} 
\sphinxAtStartPar
cutoff1 (distance units)

\item {} 
\sphinxAtStartPar
cutoff2 (distance units)

\end{itemize}

\sphinxAtStartPar
The latter 2 coefficients are optional.  If not specified, the global
class 2 and Coulombic cutoffs are used.  If only one cutoff is
specified, it is used as the cutoff for both class 2 and Coulombic
interactions for this type pair.  If both coefficients are specified,
they are used as the class 2 and Coulombic cutoffs for this type pair.
You cannot specify 2 cutoffs for style \sphinxstyleemphasis{lj/class2}, since it has no
Coulombic terms.

\sphinxAtStartPar
For \sphinxstyleemphasis{lj/class2/coul/long} only the class 2 cutoff can be specified
since a Coulombic cutoff cannot be specified for an individual I,J
type pair.  All type pairs use the same global Coulombic cutoff
specified in the pair\_style command.


\bigskip\hrule\bigskip


\sphinxAtStartPar
If the pair\_coeff command is not used to define coefficients for a
particular I != J type pair, the mixing rule for \(\epsilon\) and
\(\sigma\) for all class2 potentials is to use the \sphinxstyleemphasis{sixthpower}
formulas documented by the {\hyperref[\detokenize{pair_modify::doc}]{\sphinxcrossref{\DUrole{doc}{pair\_modify}}}} command.
The {\hyperref[\detokenize{pair_modify::doc}]{\sphinxcrossref{\DUrole{doc}{pair\_modify mix}}}} setting is thus ignored for
class2 potentials for epsilon and sigma.  However it is still followed
for mixing the cutoff distance.


\bigskip\hrule\bigskip


\sphinxAtStartPar
A version of these styles with a soft core, \sphinxstyleemphasis{lj/cut/soft}, suitable for use in
free energy calculations, is part of the FEP package and is documented with
the {\hyperref[\detokenize{pair_fep_soft::doc}]{\sphinxcrossref{\DUrole{doc}{pair\_style */soft}}}} styles. The version with soft core is
only available if LAMMPS was built with that package. See the {\hyperref[\detokenize{Build_package::doc}]{\sphinxcrossref{\DUrole{doc}{Build package}}}} page for more info.


\bigskip\hrule\bigskip


\sphinxAtStartPar
Styles with a \sphinxstyleemphasis{gpu}, \sphinxstyleemphasis{intel}, \sphinxstyleemphasis{kk}, \sphinxstyleemphasis{omp}, or \sphinxstyleemphasis{opt} suffix are
functionally the same as the corresponding style without the suffix.
They have been optimized to run faster, depending on your available
hardware, as discussed on the {\hyperref[\detokenize{Speed_packages::doc}]{\sphinxcrossref{\DUrole{doc}{Accelerator packages}}}}
page.  The accelerated styles take the same arguments and should
produce the same results, except for round\sphinxhyphen{}off and precision issues.

\sphinxAtStartPar
These accelerated styles are part of the GPU, INTEL, KOKKOS,
OPENMP, and OPT packages, respectively.  They are only enabled if
LAMMPS was built with those packages.  See the {\hyperref[\detokenize{Build_package::doc}]{\sphinxcrossref{\DUrole{doc}{Build package}}}} page for more info.

\sphinxAtStartPar
You can specify the accelerated styles explicitly in your input script
by including their suffix, or you can use the {\hyperref[\detokenize{Run_options::doc}]{\sphinxcrossref{\DUrole{doc}{\sphinxhyphen{}suffix command\sphinxhyphen{}line switch}}}} when you invoke LAMMPS, or you can use the
{\hyperref[\detokenize{suffix::doc}]{\sphinxcrossref{\DUrole{doc}{suffix}}}} command in your input script.

\sphinxAtStartPar
See the {\hyperref[\detokenize{Speed_packages::doc}]{\sphinxcrossref{\DUrole{doc}{Accelerator packages}}}} page for more
instructions on how to use the accelerated styles effectively.


\bigskip\hrule\bigskip



\subsection{Mixing, shift, table, tail correction, restart, rRESPA info}
\label{\detokenize{pair_class2:mixing-shift-table-tail-correction-restart-rrespa-info}}
\sphinxAtStartPar
For atom type pairs I,J and I != J, the epsilon and sigma coefficients
and cutoff distance for all of the lj/class2 pair styles can be mixed.
Epsilon and sigma are always mixed with the value \sphinxstyleemphasis{sixthpower}.  The
cutoff distance is mixed by whatever option is set by the pair\_modify
command (default = geometric).  See the “pair\_modify” command for
details.

\sphinxAtStartPar
All of the lj/class2 pair styles support the
{\hyperref[\detokenize{pair_modify::doc}]{\sphinxcrossref{\DUrole{doc}{pair\_modify}}}} shift option for the energy of the
Lennard\sphinxhyphen{}Jones portion of the pair interaction.

\sphinxAtStartPar
The \sphinxstyleemphasis{lj/class2/coul/long} pair style does not support the
{\hyperref[\detokenize{pair_modify::doc}]{\sphinxcrossref{\DUrole{doc}{pair\_modify}}}} table option since a tabulation
capability has not yet been added to this potential.

\sphinxAtStartPar
All of the lj/class2 pair styles support the
{\hyperref[\detokenize{pair_modify::doc}]{\sphinxcrossref{\DUrole{doc}{pair\_modify}}}} tail option for adding a long\sphinxhyphen{}range
tail correction to the energy and pressure of the Lennard\sphinxhyphen{}Jones
portion of the pair interaction.

\sphinxAtStartPar
All of the lj/class2 pair styles write their information to {\hyperref[\detokenize{restart::doc}]{\sphinxcrossref{\DUrole{doc}{binary restart files}}}}, so pair\_style and pair\_coeff commands do
not need to be specified in an input script that reads a restart file.

\sphinxAtStartPar
Only the \sphinxstyleemphasis{lj/class2} and \sphinxstyleemphasis{lj/class2/coul/long} pair styles support the use of the
\sphinxstyleemphasis{inner}, \sphinxstyleemphasis{middle}, and \sphinxstyleemphasis{outer} keywords of the {\hyperref[\detokenize{run_style::doc}]{\sphinxcrossref{\DUrole{doc}{run\_style respa}}}} command, meaning the pairwise forces can be
partitioned by distance at different levels of the rRESPA hierarchy.
The other styles only support the \sphinxstyleemphasis{pair} keyword of run\_style respa.
See the {\hyperref[\detokenize{run_style::doc}]{\sphinxcrossref{\DUrole{doc}{run\_style}}}} command for details.


\subsection{Restrictions}
\label{\detokenize{pair_class2:restrictions}}
\sphinxAtStartPar
These styles are part of the CLASS2 package.  They are only enabled if
LAMMPS was built with that package.  See the {\hyperref[\detokenize{Build_package::doc}]{\sphinxcrossref{\DUrole{doc}{Build package}}}} page for more info.


\subsection{Related commands}
\label{\detokenize{pair_class2:related-commands}}
\sphinxAtStartPar
{\hyperref[\detokenize{pair_coeff::doc}]{\sphinxcrossref{\DUrole{doc}{pair\_coeff}}}}, {\hyperref[\detokenize{pair_fep_soft::doc}]{\sphinxcrossref{\DUrole{doc}{pair\_style */soft}}}}


\subsection{Default}
\label{\detokenize{pair_class2:default}}
\sphinxAtStartPar
none


\bigskip\hrule\bigskip

\phantomsection\label{\detokenize{pair_class2:pair-sun}}
\sphinxAtStartPar
\sphinxstylestrong{(Sun)} Sun, J Phys Chem B 102, 7338\sphinxhyphen{}7364 (1998).

\sphinxstepscope

\index{pair\_style colloid@\spxentry{pair\_style colloid}}\ignorespaces 
\index{pair\_style colloid/gpu@\spxentry{pair\_style colloid/gpu}}\ignorespaces 
\index{pair\_style colloid/omp@\spxentry{pair\_style colloid/omp}}\ignorespaces 

\section{pair\_style colloid command}
\label{\detokenize{pair_colloid:pair-style-colloid-command}}\label{\detokenize{pair_colloid:index-2}}\label{\detokenize{pair_colloid:index-1}}\label{\detokenize{pair_colloid:index-0}}\label{\detokenize{pair_colloid::doc}}
\sphinxAtStartPar
Accelerator Variants: \sphinxstyleemphasis{colloid/gpu}, \sphinxstyleemphasis{colloid/omp}


\subsection{Syntax}
\label{\detokenize{pair_colloid:syntax}}
\begin{sphinxVerbatim}[commandchars=\\\{\}]
\PYG{k}{pair\PYGZus{}style}\PYG{+w}{ }\PYG{n}{colloid}\PYG{+w}{ }\PYG{n}{cutoff}
\end{sphinxVerbatim}
\begin{itemize}
\item {} 
\sphinxAtStartPar
cutoff = global cutoff for colloidal interactions (distance units)

\end{itemize}


\subsection{Examples}
\label{\detokenize{pair_colloid:examples}}
\begin{sphinxVerbatim}[commandchars=\\\{\}]
\PYG{k}{pair\PYGZus{}style}\PYG{+w}{ }\PYG{n}{colloid}\PYG{+w}{ }\PYG{l+m}{10.0}
\PYG{k}{pair\PYGZus{}coeff}\PYG{+w}{ }\PYG{o}{*}\PYG{+w}{ }\PYG{o}{*}\PYG{+w}{  }\PYG{l+m}{25}\PYG{+w}{ }\PYG{l+m}{1.0}\PYG{+w}{ }\PYG{l+m}{10.0}\PYG{+w}{ }\PYG{l+m}{10.0}
\PYG{k}{pair\PYGZus{}coeff}\PYG{+w}{ }\PYG{l+m}{1}\PYG{+w}{ }\PYG{l+m}{1}\PYG{+w}{ }\PYG{l+m}{144}\PYG{+w}{ }\PYG{l+m}{1.0}\PYG{+w}{ }\PYG{l+m}{0.0}\PYG{+w}{ }\PYG{l+m}{0.0}\PYG{+w}{ }\PYG{l+m}{3.0}
\PYG{k}{pair\PYGZus{}coeff}\PYG{+w}{ }\PYG{l+m}{1}\PYG{+w}{ }\PYG{l+m}{2}\PYG{+w}{  }\PYG{l+m}{75.398}\PYG{+w}{ }\PYG{l+m}{1.0}\PYG{+w}{ }\PYG{l+m}{0.0}\PYG{+w}{ }\PYG{l+m}{10.0}\PYG{+w}{ }\PYG{l+m}{9.0}
\PYG{k}{pair\PYGZus{}coeff}\PYG{+w}{ }\PYG{l+m}{2}\PYG{+w}{ }\PYG{l+m}{2}\PYG{+w}{  }\PYG{l+m}{39.478}\PYG{+w}{ }\PYG{l+m}{1.0}\PYG{+w}{ }\PYG{l+m}{10.0}\PYG{+w}{ }\PYG{l+m}{10.0}\PYG{+w}{ }\PYG{l+m}{25.0}
\end{sphinxVerbatim}


\subsection{Description}
\label{\detokenize{pair_colloid:description}}
\sphinxAtStartPar
Style \sphinxstyleemphasis{colloid} computes pairwise interactions between large colloidal
particles and small solvent particles using 3 formulas.  A colloidal
particle has a size \textgreater{} sigma; a solvent particle is the usual
Lennard\sphinxhyphen{}Jones particle of size sigma.

\sphinxAtStartPar
The colloid\sphinxhyphen{}colloid interaction energy is given by
\begin{equation*}
\begin{split} U_A = & - \frac{A_{cc}}{6} \left[
 \frac{2 a_1 a_2}{r^2-\left(a_1+a_2\right)^2}
 + \frac{2 a_1 a_2}{r^2 - \left(a_1 - a_2\right)^2}
   + \mathrm{ln}
     \left(
\frac{r^2-\left(a_1+a_2\right)^2}{r^2-\left(a_1-a_2\right)^2}
 \right)
\right] \\
  & \\
  U_R = & \frac{A_{cc}}{37800}  \frac{\sigma^6}{r}
  \biggl[ \frac{r^2-7r\left(a_1+a_2\right)+6\left(a_1^2+7a_1a_2+a_2^2\right)}
{\left(r-a_1-a_2\right)^7} \\
 &\qquad              +\frac{r^2+7r\left(a_1+a_2\right)+6\left(a_1^2+7a_1a_2+a_2^2\right)}
{\left(r+a_1+a_2\right)^7}  \\
&\qquad               -\frac{r^2+7r\left(a_1-a_2\right)+6\left(a_1^2-7a_1a_2+a_2^2\right)}
{\left(r+a_1-a_2\right)^7} \\
&\qquad       \left.  -\frac{r^2-7r\left(a_1-a_2\right)+6\left(a_1^2-7a_1a_2+a_2^2\right)}
{\left(r-a_1+a_2\right)^7}
\right]  \\
& \\
U = & U_A + U_R, \qquad r < r_c\end{split}
\end{equation*}
\sphinxAtStartPar
where \(A_{cc}\) is the Hamaker constant, \(a_1\) and \(a_2\) are the
radii of the two colloidal particles, and \(r_c\) is the cutoff.  This
equation results from describing each colloidal particle as an
integrated collection of Lennard\sphinxhyphen{}Jones particles of size sigma and is
derived in {\hyperref[\detokenize{pair_colloid:everaers1}]{\sphinxcrossref{\DUrole{std,std-ref}{(Everaers)}}}}.

\sphinxAtStartPar
The colloid\sphinxhyphen{}solvent interaction energy is given by
\begin{equation*}
\begin{split}U = \frac{2 ~ a^3 ~ \sigma^3 ~ A_{cs}}{9 \left( a^2 - r^2 \right)^3}
\left[ 1 - \frac{\left(5 ~ a^6+45~a^4~r^2+63~a^2~r^4+15~r^6\right) \sigma^6}
{15 \left(a-r\right)^6 \left( a+r \right)^6} \right], \quad r < r_c\end{split}
\end{equation*}
\sphinxAtStartPar
where \(A_{cs}\) is the Hamaker constant, \sphinxstyleemphasis{a} is the radius of the colloidal
particle, and \(r_c\) is the cutoff.  This formula is derived from the
colloid\sphinxhyphen{}colloid interaction, letting one of the particle sizes go to
zero.

\sphinxAtStartPar
The solvent\sphinxhyphen{}solvent interaction energy is given by the usual
Lennard\sphinxhyphen{}Jones formula
\begin{equation*}
\begin{split}U = \frac{A_{ss}}{36} \left[ \left( \frac{\sigma}{r}
     \right)^{12} - \left( \frac{ \sigma}{r} \right)^6 \right], \quad
     r < r_c\end{split}
\end{equation*}
\sphinxAtStartPar
with \(A_{ss}\) set appropriately, which results from letting both
particle sizes go to zero.

\sphinxAtStartPar
When used in combination with {\hyperref[\detokenize{pair_yukawa_colloid::doc}]{\sphinxcrossref{\DUrole{doc}{pair\_style yukawa/colloid}}}}, the two terms become the so\sphinxhyphen{}called DLVO potential,
which combines electrostatic repulsion and van der Waals attraction.

\sphinxAtStartPar
The following coefficients must be defined for each pair of atoms
types via the {\hyperref[\detokenize{pair_coeff::doc}]{\sphinxcrossref{\DUrole{doc}{pair\_coeff}}}} command as in the examples
above, or in the data file or restart files read by the
{\hyperref[\detokenize{read_data::doc}]{\sphinxcrossref{\DUrole{doc}{read\_data}}}} or {\hyperref[\detokenize{read_restart::doc}]{\sphinxcrossref{\DUrole{doc}{read\_restart}}}}
commands, or by mixing as described below:
\begin{itemize}
\item {} 
\sphinxAtStartPar
A (energy units)

\item {} 
\sphinxAtStartPar
\(\sigma\) (distance units)

\item {} 
\sphinxAtStartPar
d1 (distance units)

\item {} 
\sphinxAtStartPar
d2 (distance units)

\item {} 
\sphinxAtStartPar
cutoff (distance units)

\end{itemize}

\sphinxAtStartPar
A is the Hamaker energy prefactor and should typically be set as
follows:
\begin{itemize}
\item {} 
\sphinxAtStartPar
\(A_{cc}\) = colloid/colloid = \(4 \pi^2 = 39.5\)

\item {} 
\sphinxAtStartPar
\(A_{cs}\) = colloid/solvent = \(\sqrt{A_{cc} A_{ss}}\)

\item {} 
\sphinxAtStartPar
\(A_{ss}\) = solvent/solvent = 144 (assuming epsilon = 1, so that 144/36 = 4)

\end{itemize}

\sphinxAtStartPar
\(\sigma\) is the size of the solvent particle or the constituent
particles integrated over in the colloidal particle and should typically
be set as follows:
\begin{itemize}
\item {} 
\sphinxAtStartPar
\(\sigma_{cc}\) = colloid/colloid = 1.0

\item {} 
\sphinxAtStartPar
\(\sigma_{cs}\) = colloid/solvent = arithmetic mixing between colloid \(\sigma\) and solvent \(\sigma\)

\item {} 
\sphinxAtStartPar
\(\sigma_{ss}\) = solvent/solvent = 1.0 or whatever size the solvent particle is

\end{itemize}

\sphinxAtStartPar
Thus typically \(\sigma_{cs} = 1.0\), unless the solvent particle’s size !=
1.0.

\sphinxAtStartPar
D1 and d2 are particle diameters, so that d1 = 2*a1 and d2 = 2*a2 in
the formulas above.  Both d1 and d2 must be values \textgreater{}= 0.  If d1 \textgreater{} 0
and d2 \textgreater{} 0, then the pair interacts via the colloid\sphinxhyphen{}colloid formula
above.  If d1 = 0 and d2 = 0, then the pair interacts via the
solvent\sphinxhyphen{}solvent formula.  I.e. a d value of 0 is a Lennard\sphinxhyphen{}Jones
particle of size \(\sigma\).  If either d1 = 0 or d2 = 0 and the other is
larger, then the pair interacts via the colloid\sphinxhyphen{}solvent formula.

\sphinxAtStartPar
Note that the diameter of a particular particle type may appear in
multiple pair\_coeff commands, as it interacts with other particle
types.  You should ensure the particle diameter is specified
consistently each time it appears.

\sphinxAtStartPar
The last coefficient is optional.  If not specified, the global cutoff
specified in the pair\_style command is used.  However, you typically
want different cutoffs for interactions between different particle
sizes.  E.g. if colloidal particles of diameter 10 are used with
solvent particles of diameter 1, then a solvent\sphinxhyphen{}solvent cutoff of 2.5
would correspond to a colloid\sphinxhyphen{}colloid cutoff of 25.  A good
rule\sphinxhyphen{}of\sphinxhyphen{}thumb is to use a colloid\sphinxhyphen{}solvent cutoff that is half the big
diameter + 4 times the small diameter.  I.e. 9 = 5 + 4 for the
colloid\sphinxhyphen{}solvent cutoff in this case.

\begin{sphinxadmonition}{note}{Note:}
\sphinxAtStartPar
When using pair\_style colloid for a mixture with 2 (or more)
widely different particles sizes (e.g. sigma=10 colloids in a
background sigma=1 LJ fluid), you will likely want to use these
commands for efficiency: {\hyperref[\detokenize{neighbor::doc}]{\sphinxcrossref{\DUrole{doc}{neighbor multi}}}} and
{\hyperref[\detokenize{comm_modify::doc}]{\sphinxcrossref{\DUrole{doc}{comm\_modify multi}}}}.
\end{sphinxadmonition}


\bigskip\hrule\bigskip


\sphinxAtStartPar
Styles with a \sphinxstyleemphasis{gpu}, \sphinxstyleemphasis{intel}, \sphinxstyleemphasis{kk}, \sphinxstyleemphasis{omp}, or \sphinxstyleemphasis{opt} suffix are
functionally the same as the corresponding style without the suffix.
They have been optimized to run faster, depending on your available
hardware, as discussed on the {\hyperref[\detokenize{Speed_packages::doc}]{\sphinxcrossref{\DUrole{doc}{Accelerator packages}}}}
page.  The accelerated styles take the same arguments and should
produce the same results, except for round\sphinxhyphen{}off and precision issues.

\sphinxAtStartPar
These accelerated styles are part of the GPU, INTEL, KOKKOS,
OPENMP, and OPT packages, respectively.  They are only enabled if
LAMMPS was built with those packages.  See the {\hyperref[\detokenize{Build_package::doc}]{\sphinxcrossref{\DUrole{doc}{Build package}}}} page for more info.

\sphinxAtStartPar
You can specify the accelerated styles explicitly in your input script
by including their suffix, or you can use the {\hyperref[\detokenize{Run_options::doc}]{\sphinxcrossref{\DUrole{doc}{\sphinxhyphen{}suffix command\sphinxhyphen{}line switch}}}} when you invoke LAMMPS, or you can use the
{\hyperref[\detokenize{suffix::doc}]{\sphinxcrossref{\DUrole{doc}{suffix}}}} command in your input script.

\sphinxAtStartPar
See the {\hyperref[\detokenize{Speed_packages::doc}]{\sphinxcrossref{\DUrole{doc}{Accelerator packages}}}} page for more
instructions on how to use the accelerated styles effectively.


\bigskip\hrule\bigskip



\subsection{Mixing, shift, table, tail correction, restart, rRESPA info}
\label{\detokenize{pair_colloid:mixing-shift-table-tail-correction-restart-rrespa-info}}
\sphinxAtStartPar
For atom type pairs I,J and I != J, the A, sigma, d1, and d2
coefficients and cutoff distance for this pair style can be mixed.  A
is an energy value mixed like a LJ epsilon.  D1 and d2 are distance
values and are mixed like sigma.  The default mix value is
\sphinxstyleemphasis{geometric}.  See the “pair\_modify” command for details.

\sphinxAtStartPar
This pair style supports the {\hyperref[\detokenize{pair_modify::doc}]{\sphinxcrossref{\DUrole{doc}{pair\_modify}}}} shift
option for the energy of the pair interaction.

\sphinxAtStartPar
The {\hyperref[\detokenize{pair_modify::doc}]{\sphinxcrossref{\DUrole{doc}{pair\_modify}}}} table option is not relevant
for this pair style.

\sphinxAtStartPar
This pair style does not support the {\hyperref[\detokenize{pair_modify::doc}]{\sphinxcrossref{\DUrole{doc}{pair\_modify}}}}
tail option for adding long\sphinxhyphen{}range tail corrections to energy and
pressure.

\sphinxAtStartPar
This pair style writes its information to {\hyperref[\detokenize{restart::doc}]{\sphinxcrossref{\DUrole{doc}{binary restart files}}}}, so pair\_style and pair\_coeff commands do not need
to be specified in an input script that reads a restart file.

\sphinxAtStartPar
This pair style can only be used via the \sphinxstyleemphasis{pair} keyword of the
{\hyperref[\detokenize{run_style::doc}]{\sphinxcrossref{\DUrole{doc}{run\_style respa}}}} command.  It does not support the
\sphinxstyleemphasis{inner}, \sphinxstyleemphasis{middle}, \sphinxstyleemphasis{outer} keywords.


\bigskip\hrule\bigskip



\subsection{Restrictions}
\label{\detokenize{pair_colloid:restrictions}}
\sphinxAtStartPar
This style is part of the COLLOID package.  It is only enabled if
LAMMPS was built with that package.  See the {\hyperref[\detokenize{Build_package::doc}]{\sphinxcrossref{\DUrole{doc}{Build package}}}} page for more info.

\sphinxAtStartPar
Normally, this pair style should be used with finite\sphinxhyphen{}size particles
which have a diameter, e.g. see the {\hyperref[\detokenize{atom_style::doc}]{\sphinxcrossref{\DUrole{doc}{atom\_style sphere}}}} command.  However, this is not a requirement,
since the only definition of particle size is via the pair\_coeff
parameters for each type.  In other words, the physical radius of the
particle is ignored.  Thus you should ensure that the d1,d2 parameters
you specify are consistent with the physical size of the particles of
that type.

\sphinxAtStartPar
Per\sphinxhyphen{}particle polydispersity is not yet supported by this pair style;
only per\sphinxhyphen{}type polydispersity is enabled via the pair\_coeff parameters.


\subsection{Related commands}
\label{\detokenize{pair_colloid:related-commands}}
\sphinxAtStartPar
{\hyperref[\detokenize{pair_coeff::doc}]{\sphinxcrossref{\DUrole{doc}{pair\_coeff}}}}


\subsection{Default}
\label{\detokenize{pair_colloid:default}}
\sphinxAtStartPar
none


\bigskip\hrule\bigskip

\phantomsection\label{\detokenize{pair_colloid:everaers1}}
\sphinxAtStartPar
\sphinxstylestrong{(Everaers)} Everaers, Ejtehadi, Phys Rev E, 67, 041710 (2003).

\sphinxstepscope

\index{pair\_style comb@\spxentry{pair\_style comb}}\ignorespaces 
\index{pair\_style comb/omp@\spxentry{pair\_style comb/omp}}\ignorespaces 
\index{pair\_style comb3@\spxentry{pair\_style comb3}}\ignorespaces 

\section{pair\_style comb command}
\label{\detokenize{pair_comb:pair-style-comb-command}}\label{\detokenize{pair_comb:index-2}}\label{\detokenize{pair_comb:index-1}}\label{\detokenize{pair_comb:index-0}}\label{\detokenize{pair_comb::doc}}
\sphinxAtStartPar
Accelerator Variants: \sphinxstyleemphasis{comb/omp}


\section{pair\_style comb3 command}
\label{\detokenize{pair_comb:pair-style-comb3-command}}

\subsection{Syntax}
\label{\detokenize{pair_comb:syntax}}
\begin{sphinxVerbatim}[commandchars=\\\{\}]
\PYG{k}{pair\PYGZus{}style}\PYG{+w}{ }\PYG{n}{comb}
\PYG{k}{pair\PYGZus{}style}\PYG{+w}{ }\PYG{n}{comb3}\PYG{+w}{ }\PYG{n}{keyword}
\end{sphinxVerbatim}
\begin{sphinxalltt}
keyword = \sphinxstyleemphasis{polar}
  \sphinxstyleemphasis{polar} value = \sphinxstyleemphasis{polar\_on} or \sphinxstyleemphasis{polar\_off} = whether or not to include atomic polarization
\end{sphinxalltt}


\subsection{Examples}
\label{\detokenize{pair_comb:examples}}
\begin{sphinxVerbatim}[commandchars=\\\{\}]
\PYG{k}{pair\PYGZus{}style}\PYG{+w}{ }\PYG{n}{comb}
\PYG{k}{pair\PYGZus{}coeff}\PYG{+w}{ }\PYG{o}{*}\PYG{+w}{ }\PYG{o}{*}\PYG{+w}{ }\PYG{n}{..}\PYG{o}{/}\PYG{n}{potentials}\PYG{o}{/}\PYG{n}{ffield.comb}\PYG{+w}{ }\PYG{n}{Si}
\PYG{k}{pair\PYGZus{}coeff}\PYG{+w}{ }\PYG{o}{*}\PYG{+w}{ }\PYG{o}{*}\PYG{+w}{ }\PYG{n}{..}\PYG{o}{/}\PYG{n}{potentials}\PYG{o}{/}\PYG{n}{ffield.comb}\PYG{+w}{ }\PYG{n}{Hf}\PYG{+w}{ }\PYG{n}{Si}\PYG{+w}{ }\PYG{n}{O}

\PYG{k}{pair\PYGZus{}style}\PYG{+w}{ }\PYG{n}{comb3}\PYG{+w}{ }\PYG{n}{polar\PYGZus{}off}
\PYG{k}{pair\PYGZus{}coeff}\PYG{+w}{ }\PYG{o}{*}\PYG{+w}{ }\PYG{o}{*}\PYG{+w}{ }\PYG{n}{..}\PYG{o}{/}\PYG{n}{potentials}\PYG{o}{/}\PYG{n}{ffield.comb3}\PYG{+w}{ }\PYG{n}{O}\PYG{+w}{ }\PYG{n}{Cu}\PYG{+w}{ }\PYG{n}{N}\PYG{+w}{ }\PYG{n}{C}\PYG{+w}{ }\PYG{n}{O}
\end{sphinxVerbatim}


\subsection{Description}
\label{\detokenize{pair_comb:description}}
\sphinxAtStartPar
Style \sphinxstyleemphasis{comb} computes the second generation variable charge COMB
(Charge\sphinxhyphen{}Optimized Many\sphinxhyphen{}Body) potential.  Style \sphinxstyleemphasis{comb3} computes the
third\sphinxhyphen{}generation COMB potential.  These COMB potentials are described
in {\hyperref[\detokenize{pair_comb:comb}]{\sphinxcrossref{\DUrole{std,std-ref}{(COMB)}}}} and {\hyperref[\detokenize{pair_comb:comb3}]{\sphinxcrossref{\DUrole{std,std-ref}{(COMB3)}}}}.  Briefly, the
total energy \(E_T\) of a system of atoms is given by
\begin{equation*}
\begin{split}E_T  = & \sum_i [ E_i^{self} (q_i) + \sum_{j>i} [E_{ij}^{short} (r_{ij}, q_i, q_j) + E_{ij}^{Coul} (r_{ij}, q_i, q_j)] + \\
       & E^{polar} (q_i, r_{ij}) + E^{vdW} (r_{ij}) + E^{barr} (q_i) + E^{corr} (r_{ij}, \theta_{jik})]\end{split}
\end{equation*}
\sphinxAtStartPar
where \(E_i^{self}\) is the self\sphinxhyphen{}energy of atom \sphinxstyleemphasis{i} (including
atomic ionization energies and electron affinities),
\(E_{ij}^{short}\) is the bond\sphinxhyphen{}order potential between atoms \sphinxstyleemphasis{i} and
\sphinxstyleemphasis{j}, \(E_{ij}^{Coul}\) is the Coulomb interactions,
\(E^{polar}\) is the polarization term for organic systems (style
\sphinxstyleemphasis{comb3} only), \(E^{vdW}\) is the van der Waals energy (style
\sphinxstyleemphasis{comb3} only), \(E^{barr}\) is a charge barrier function, and
\(E^{corr}\) are angular correction terms.

\sphinxAtStartPar
The COMB potentials (styles \sphinxstyleemphasis{comb} and \sphinxstyleemphasis{comb3}) are variable charge
potentials.  The equilibrium charge on each atom is calculated by the
electronegativity equalization (QEq) method.  See {\hyperref[\detokenize{pair_comb:rick2}]{\sphinxcrossref{\DUrole{std,std-ref}{Rick}}}} for
further details.  This is implemented by the {\hyperref[\detokenize{fix_qeq_comb::doc}]{\sphinxcrossref{\DUrole{doc}{fix qeq/comb}}}} command, which should normally be
specified in the input script when running a model with the COMB
potential.  The {\hyperref[\detokenize{fix_qeq_comb::doc}]{\sphinxcrossref{\DUrole{doc}{fix qeq/comb}}}} command has options
that determine how often charge equilibration is performed, its
convergence criterion, and which atoms are included in the
calculation.

\sphinxAtStartPar
Only a single pair\_coeff command is used with the \sphinxstyleemphasis{comb} and \sphinxstyleemphasis{comb3}
styles which specifies the COMB potential file with parameters for all
needed elements.  These are mapped to LAMMPS atom types by specifying
N additional arguments after the potential file in the pair\_coeff
command, where N is the number of LAMMPS atom types.

\sphinxAtStartPar
For example, if your LAMMPS simulation of a Si/SiO2/
HfO2 interface has 4 atom types, and you want the first and
last to be Si, the second to be Hf, and the third to be O, and you would
use the following pair\_coeff command:

\begin{sphinxVerbatim}[commandchars=\\\{\}]
\PYG{k}{pair\PYGZus{}coeff}\PYG{+w}{ }\PYG{o}{*}\PYG{+w}{ }\PYG{o}{*}\PYG{+w}{ }\PYG{n}{..}\PYG{o}{/}\PYG{n}{potentials}\PYG{o}{/}\PYG{n}{ffield.comb}\PYG{+w}{ }\PYG{n}{Si}\PYG{+w}{ }\PYG{n}{Hf}\PYG{+w}{ }\PYG{n}{O}\PYG{+w}{ }\PYG{n}{Si}
\end{sphinxVerbatim}

\sphinxAtStartPar
The first two arguments must be * * so as to span all LAMMPS atom
types.  The first and last Si arguments map LAMMPS atom types 1 and 4
to the Si element in the \sphinxstyleemphasis{ffield.comb} file.  The second Hf argument
maps LAMMPS atom type 2 to the Hf element, and the third O argument
maps LAMMPS atom type 3 to the O element in the potential file.  If a
mapping value is specified as NULL, the mapping is not performed.
This can be used when a \sphinxstyleemphasis{comb} potential is used as part of the
\sphinxstyleemphasis{hybrid} pair style.  The NULL values are placeholders for atom types
that will be used with other potentials.

\sphinxAtStartPar
For style \sphinxstyleemphasis{comb}, the provided potential file \sphinxstyleemphasis{ffield.comb} contains
all currently\sphinxhyphen{}available second generation COMB parameterizations: for Si,
Cu, Hf, Ti, O, their oxides and Zr, Zn and U metals.  For style
\sphinxstyleemphasis{comb3}, the potential file \sphinxstyleemphasis{ffield.comb3} contains all
currently\sphinxhyphen{}available third generation COMB parameterizations: O, Cu, N, C,
H, Ti, Zn and Zr.  The status of the optimization of the compounds, for
example Cu2O, TiN and hydrocarbons, are given in the
following table:


\begin{savenotes}\sphinxattablestart
\sphinxthistablewithglobalstyle
\centering
\begin{tabulary}{\linewidth}[t]{TTTTTTTTT}
\sphinxtoprule
\sphinxtableatstartofbodyhook&
\sphinxAtStartPar
O
&
\sphinxAtStartPar
Cu
&
\sphinxAtStartPar
N
&
\sphinxAtStartPar
C
&
\sphinxAtStartPar
H
&
\sphinxAtStartPar
Ti
&
\sphinxAtStartPar
Zn
&
\sphinxAtStartPar
Zr
\\
\sphinxhline
\sphinxAtStartPar
O
&
\sphinxAtStartPar
F
&
\sphinxAtStartPar
F
&
\sphinxAtStartPar
F
&
\sphinxAtStartPar
F
&
\sphinxAtStartPar
F
&
\sphinxAtStartPar
F
&
\sphinxAtStartPar
F
&
\sphinxAtStartPar
F
\\
\sphinxhline
\sphinxAtStartPar
Cu
&
\sphinxAtStartPar
F
&
\sphinxAtStartPar
F
&
\sphinxAtStartPar
P
&
\sphinxAtStartPar
F
&
\sphinxAtStartPar
F
&
\sphinxAtStartPar
P
&
\sphinxAtStartPar
F
&
\sphinxAtStartPar
P
\\
\sphinxhline
\sphinxAtStartPar
N
&
\sphinxAtStartPar
F
&
\sphinxAtStartPar
P
&
\sphinxAtStartPar
F
&
\sphinxAtStartPar
M
&
\sphinxAtStartPar
F
&
\sphinxAtStartPar
P
&
\sphinxAtStartPar
P
&
\sphinxAtStartPar
P
\\
\sphinxhline
\sphinxAtStartPar
C
&
\sphinxAtStartPar
F
&
\sphinxAtStartPar
F
&
\sphinxAtStartPar
M
&
\sphinxAtStartPar
F
&
\sphinxAtStartPar
F
&
\sphinxAtStartPar
M
&
\sphinxAtStartPar
M
&
\sphinxAtStartPar
M
\\
\sphinxhline
\sphinxAtStartPar
H
&
\sphinxAtStartPar
F
&
\sphinxAtStartPar
F
&
\sphinxAtStartPar
F
&
\sphinxAtStartPar
F
&
\sphinxAtStartPar
F
&
\sphinxAtStartPar
M
&
\sphinxAtStartPar
M
&
\sphinxAtStartPar
F
\\
\sphinxhline
\sphinxAtStartPar
Ti
&
\sphinxAtStartPar
F
&
\sphinxAtStartPar
P
&
\sphinxAtStartPar
P
&
\sphinxAtStartPar
M
&
\sphinxAtStartPar
M
&
\sphinxAtStartPar
F
&
\sphinxAtStartPar
P
&
\sphinxAtStartPar
P
\\
\sphinxhline
\sphinxAtStartPar
Zn
&
\sphinxAtStartPar
F
&
\sphinxAtStartPar
F
&
\sphinxAtStartPar
P
&
\sphinxAtStartPar
M
&
\sphinxAtStartPar
M
&
\sphinxAtStartPar
P
&
\sphinxAtStartPar
F
&
\sphinxAtStartPar
P
\\
\sphinxhline
\sphinxAtStartPar
Zr
&
\sphinxAtStartPar
F
&
\sphinxAtStartPar
P
&
\sphinxAtStartPar
P
&
\sphinxAtStartPar
M
&
\sphinxAtStartPar
F
&
\sphinxAtStartPar
P
&
\sphinxAtStartPar
P
&
\sphinxAtStartPar
F
\\
\sphinxbottomrule
\end{tabulary}
\sphinxtableafterendhook\par
\sphinxattableend\end{savenotes}
\begin{itemize}
\item {} 
\sphinxAtStartPar
F = Fully optimized

\item {} 
\sphinxAtStartPar
M = Only optimized for dimer molecule

\item {} 
\sphinxAtStartPar
P = in progress, but have it from mixing rule

\end{itemize}

\sphinxAtStartPar
For style \sphinxstyleemphasis{comb3}, in addition to ffield.comb3, a special parameter
file, \sphinxstyleemphasis{lib.comb3}, that is exclusively used for C/O/H systems, will be
automatically loaded if carbon atom is detected in LAMMPS input
structure.  This file must be in your working directory or in the
directories listed in the environment variable \sphinxcode{\sphinxupquote{LAMMPS\_POTENTIALS}}, as
described on the {\hyperref[\detokenize{pair_coeff::doc}]{\sphinxcrossref{\DUrole{doc}{pair\_coeff}}}} command doc page.

\sphinxAtStartPar
The keyword \sphinxstyleemphasis{polar} indicates whether the force field includes
the atomic polarization.  Since the equilibration of the polarization
has not yet been implemented, it can only set polar\_off at present.

\begin{sphinxadmonition}{note}{Note:}
\sphinxAtStartPar
You can not use potential file \sphinxstyleemphasis{ffield.comb} with style \sphinxstyleemphasis{comb3},
nor file \sphinxstyleemphasis{ffield.comb3} with style \sphinxstyleemphasis{comb}.
\end{sphinxadmonition}


\bigskip\hrule\bigskip


\sphinxAtStartPar
Styles with a \sphinxstyleemphasis{gpu}, \sphinxstyleemphasis{intel}, \sphinxstyleemphasis{kk}, \sphinxstyleemphasis{omp}, or \sphinxstyleemphasis{opt} suffix are
functionally the same as the corresponding style without the suffix.
They have been optimized to run faster, depending on your available
hardware, as discussed on the {\hyperref[\detokenize{Speed_packages::doc}]{\sphinxcrossref{\DUrole{doc}{Accelerator packages}}}}
page.  The accelerated styles take the same arguments and should
produce the same results, except for round\sphinxhyphen{}off and precision issues.

\sphinxAtStartPar
These accelerated styles are part of the GPU, INTEL, KOKKOS,
OPENMP, and OPT packages, respectively.  They are only enabled if
LAMMPS was built with those packages.  See the {\hyperref[\detokenize{Build_package::doc}]{\sphinxcrossref{\DUrole{doc}{Build package}}}} page for more info.

\sphinxAtStartPar
You can specify the accelerated styles explicitly in your input script
by including their suffix, or you can use the {\hyperref[\detokenize{Run_options::doc}]{\sphinxcrossref{\DUrole{doc}{\sphinxhyphen{}suffix command\sphinxhyphen{}line switch}}}} when you invoke LAMMPS, or you can use the
{\hyperref[\detokenize{suffix::doc}]{\sphinxcrossref{\DUrole{doc}{suffix}}}} command in your input script.

\sphinxAtStartPar
See the {\hyperref[\detokenize{Speed_packages::doc}]{\sphinxcrossref{\DUrole{doc}{Accelerator packages}}}} page for more
instructions on how to use the accelerated styles effectively.


\bigskip\hrule\bigskip



\subsection{Mixing, shift, table, tail correction, restart, rRESPA info}
\label{\detokenize{pair_comb:mixing-shift-table-tail-correction-restart-rrespa-info}}
\sphinxAtStartPar
For atom type pairs I,J and I != J, where types I and J correspond to
two different element types, mixing is performed by LAMMPS as
described above from values in the potential file.

\sphinxAtStartPar
These pair styles does not support the {\hyperref[\detokenize{pair_modify::doc}]{\sphinxcrossref{\DUrole{doc}{pair\_modify}}}}
shift, table, and tail options.

\sphinxAtStartPar
These pair styles do not write its information to {\hyperref[\detokenize{restart::doc}]{\sphinxcrossref{\DUrole{doc}{binary restart files}}}}, since it is stored in potential files.  Thus, you
need to re\sphinxhyphen{}specify the pair\_style, pair\_coeff, and {\hyperref[\detokenize{fix_qeq_comb::doc}]{\sphinxcrossref{\DUrole{doc}{fix qeq/comb}}}} commands in an input script that reads a
restart file.

\sphinxAtStartPar
These pair styles can only be used via the \sphinxstyleemphasis{pair} keyword of the
{\hyperref[\detokenize{run_style::doc}]{\sphinxcrossref{\DUrole{doc}{run\_style respa}}}} command.  It does not support the
\sphinxstyleemphasis{inner}, \sphinxstyleemphasis{middle}, \sphinxstyleemphasis{outer} keywords.


\bigskip\hrule\bigskip



\subsection{Restrictions}
\label{\detokenize{pair_comb:restrictions}}
\sphinxAtStartPar
These pair styles are part of the MANYBODY package.  It is only
enabled if LAMMPS was built with that package.  See the {\hyperref[\detokenize{Build_package::doc}]{\sphinxcrossref{\DUrole{doc}{Build package}}}} page for more info.

\sphinxAtStartPar
These pair styles requires the {\hyperref[\detokenize{newton::doc}]{\sphinxcrossref{\DUrole{doc}{newton}}}} setting to be “on”
for pair interactions.

\sphinxAtStartPar
The COMB potentials in the \sphinxstyleemphasis{ffield.comb} and \sphinxstyleemphasis{ffield.comb3} files provided
with LAMMPS (see the potentials directory) are parameterized for metal
{\hyperref[\detokenize{units::doc}]{\sphinxcrossref{\DUrole{doc}{units}}}}.  You can use the COMB potential with any LAMMPS
units, but you would need to create your own COMB potential file with
coefficients listed in the appropriate units if your simulation
does not use “metal” units.


\subsection{Related commands}
\label{\detokenize{pair_comb:related-commands}}
\sphinxAtStartPar
{\hyperref[\detokenize{pair_style::doc}]{\sphinxcrossref{\DUrole{doc}{pair\_style}}}}, {\hyperref[\detokenize{pair_coeff::doc}]{\sphinxcrossref{\DUrole{doc}{pair\_coeff}}}},
{\hyperref[\detokenize{fix_qeq_comb::doc}]{\sphinxcrossref{\DUrole{doc}{fix qeq/comb}}}}


\subsection{Default}
\label{\detokenize{pair_comb:default}}
\sphinxAtStartPar
none


\bigskip\hrule\bigskip

\phantomsection\label{\detokenize{pair_comb:comb}}
\sphinxAtStartPar
\sphinxstylestrong{(COMB)}  T.\sphinxhyphen{}R. Shan, B. D. Devine, T. W. Kemper, S. B. Sinnott, and
S. R. Phillpot, Phys. Rev. B 81, 125328 (2010)

\phantomsection\label{\detokenize{pair_comb:comb3}}
\sphinxAtStartPar
\sphinxstylestrong{(COMB3)} T. Liang, T.\sphinxhyphen{}R. Shan, Y.\sphinxhyphen{}T. Cheng, B. D. Devine, M. Noordhoek,
Y. Li, Z. Lu, S. R. Phillpot, and S. B. Sinnott, Mat. Sci. \& Eng: R 74,
255\sphinxhyphen{}279 (2013).

\phantomsection\label{\detokenize{pair_comb:rick2}}
\sphinxAtStartPar
\sphinxstylestrong{(Rick)} S. W. Rick, S. J. Stuart, B. J. Berne, J Chem Phys 101, 6141
(1994).

\sphinxstepscope

\index{pair\_style cosine/squared@\spxentry{pair\_style cosine/squared}}\ignorespaces 

\section{pair\_style cosine/squared command}
\label{\detokenize{pair_cosine_squared:pair-style-cosine-squared-command}}\label{\detokenize{pair_cosine_squared:index-0}}\label{\detokenize{pair_cosine_squared::doc}}

\subsection{Syntax}
\label{\detokenize{pair_cosine_squared:syntax}}
\begin{sphinxVerbatim}[commandchars=\\\{\}]
\PYG{k}{pair\PYGZus{}style}\PYG{+w}{ }\PYG{n}{cosine}\PYG{o}{/}\PYG{n}{squared}\PYG{+w}{ }\PYG{n}{cutoff}
\end{sphinxVerbatim}
\begin{itemize}
\item {} 
\sphinxAtStartPar
cutoff = global cutoff for cosine\sphinxhyphen{}squared interactions (distance units)

\end{itemize}

\begin{sphinxVerbatim}[commandchars=\\\{\}]
\PYG{k}{pair\PYGZus{}coeff}\PYG{+w}{ }\PYG{n}{I}\PYG{+w}{ }\PYG{n}{J}\PYG{+w}{ }\PYG{n}{eps}\PYG{+w}{ }\PYG{n}{sigma}
\PYG{k}{pair\PYGZus{}coeff}\PYG{+w}{ }\PYG{n}{I}\PYG{+w}{ }\PYG{n}{J}\PYG{+w}{ }\PYG{n}{eps}\PYG{+w}{ }\PYG{n}{sigma}\PYG{+w}{ }\PYG{n}{cutoff}
\PYG{k}{pair\PYGZus{}coeff}\PYG{+w}{ }\PYG{n}{I}\PYG{+w}{ }\PYG{n}{J}\PYG{+w}{ }\PYG{n}{eps}\PYG{+w}{ }\PYG{n}{sigma}\PYG{+w}{ }\PYG{n}{wca}
\PYG{k}{pair\PYGZus{}coeff}\PYG{+w}{ }\PYG{n}{I}\PYG{+w}{ }\PYG{n}{J}\PYG{+w}{ }\PYG{n}{eps}\PYG{+w}{ }\PYG{n}{sigma}\PYG{+w}{ }\PYG{n}{cutoff}\PYG{+w}{ }\PYG{n}{wca}
\end{sphinxVerbatim}
\begin{itemize}
\item {} 
\sphinxAtStartPar
I, J = a particle type

\item {} 
\sphinxAtStartPar
eps = interaction strength, i.e. the depth of the potential minimum (energy units)

\item {} 
\sphinxAtStartPar
sigma = distance of the potential minimum from 0

\item {} 
\sphinxAtStartPar
cutoff = the cutoff distance for this pair type, if different from global (distance units)

\item {} 
\sphinxAtStartPar
wca = if specified a Weeks\sphinxhyphen{}Chandler\sphinxhyphen{}Andersen potential (with eps strength and minimum at sigma) is added, otherwise not

\end{itemize}


\subsection{Examples}
\label{\detokenize{pair_cosine_squared:examples}}
\begin{sphinxVerbatim}[commandchars=\\\{\}]
\PYG{k}{pair\PYGZus{}style}\PYG{+w}{ }\PYG{n}{cosine}\PYG{o}{/}\PYG{n}{squared}\PYG{+w}{ }\PYG{l+m}{3.0}
\PYG{k}{pair\PYGZus{}coeff}\PYG{+w}{ }\PYG{o}{*}\PYG{+w}{ }\PYG{o}{*}\PYG{+w}{ }\PYG{l+m}{1.0}\PYG{+w}{ }\PYG{l+m}{1.3}
\PYG{k}{pair\PYGZus{}coeff}\PYG{+w}{ }\PYG{l+m}{1}\PYG{+w}{ }\PYG{l+m}{3}\PYG{+w}{ }\PYG{l+m}{1.0}\PYG{+w}{ }\PYG{l+m}{1.3}\PYG{+w}{ }\PYG{l+m}{2.0}
\PYG{k}{pair\PYGZus{}coeff}\PYG{+w}{ }\PYG{l+m}{1}\PYG{+w}{ }\PYG{l+m}{3}\PYG{+w}{ }\PYG{l+m}{1.0}\PYG{+w}{ }\PYG{l+m}{1.3}\PYG{+w}{ }\PYG{n}{wca}
\PYG{k}{pair\PYGZus{}coeff}\PYG{+w}{ }\PYG{l+m}{1}\PYG{+w}{ }\PYG{l+m}{3}\PYG{+w}{ }\PYG{l+m}{1.0}\PYG{+w}{ }\PYG{l+m}{1.3}\PYG{+w}{ }\PYG{l+m}{2.0}\PYG{+w}{ }\PYG{n}{wca}
\end{sphinxVerbatim}


\subsection{Description}
\label{\detokenize{pair_cosine_squared:description}}
\sphinxAtStartPar
Style \sphinxstyleemphasis{cosine/squared} computes a potential of the form
\begin{equation*}
\begin{split}E =
\begin{cases}
-\epsilon& \quad r < \sigma \\
-\epsilon\cos\left(\frac{\pi\left(r - \sigma\right)}{2\left(r_c - \sigma\right)}\right)^2&\quad \sigma \leq r < r_c \\
0& \quad r \geq r_c
\end{cases}\end{split}
\end{equation*}
\sphinxAtStartPar
between two point particles, where (\(\sigma, -\epsilon\)) is the
location of the (rightmost) minimum of the potential, as explained in
the syntax section above.

\sphinxAtStartPar
This potential was first used in {\hyperref[\detokenize{pair_cosine_squared:ckd}]{\sphinxcrossref{\DUrole{std,std-ref}{(Cooke)}}}} for a coarse\sphinxhyphen{}grained lipid
membrane model.  It is generally very useful as a non\sphinxhyphen{}specific
interaction potential because it is fully adjustable in depth and width
while joining the minimum at (sigma, \sphinxhyphen{}epsilon) and zero at (cutoff, 0)
smoothly, requiring no shifting and causing no related artifacts, tail
energy calculations etc. This evidently requires \sphinxstyleemphasis{cutoff} to be larger
than \sphinxstyleemphasis{sigma}.

\sphinxAtStartPar
If the \sphinxstyleemphasis{wca} option is used then a Weeks\sphinxhyphen{}Chandler\sphinxhyphen{}Andersen potential
{\hyperref[\detokenize{pair_cosine_squared:wca}]{\sphinxcrossref{\DUrole{std,std-ref}{(Weeks)}}}} is added to the above specified cosine\sphinxhyphen{}squared potential,
specifically the following:
\begin{equation*}
\begin{split}E = \epsilon \left[ \left(\frac{\sigma}{r}\right)^{12} -
                      2\left(\frac{\sigma}{r}\right)^6 + 1\right]
                      , \quad r < \sigma\end{split}
\end{equation*}
\sphinxAtStartPar
In this case, and this case only, the \(\sigma\) parameter can be equal to
\sphinxstyleemphasis{cutoff} (\(\sigma =\) cutoff) which will result in ONLY the WCA potential
being used (and print a warning), so the minimum will be attained at
(sigma, 0). This is a convenience feature that enables a purely
repulsive potential to be used without a need to define an additional
pair style and use the hybrid styles.

\sphinxAtStartPar
The energy and force of this pair style for parameters epsilon = 1.0,
sigma = 1.0, cutoff = 2.5, with and without the WCA potential, are shown
in the graphs below:

\noindent{\hspace*{\fill}\sphinxincludegraphics{{pair_cosine_squared_graphs}.jpg}\hspace*{\fill}}


\bigskip\hrule\bigskip



\subsection{Mixing, shift, table, tail correction, restart, rRESPA info}
\label{\detokenize{pair_cosine_squared:mixing-shift-table-tail-correction-restart-rrespa-info}}
\sphinxAtStartPar
Mixing is not supported for this style.

\sphinxAtStartPar
The \sphinxstyleemphasis{shift}, \sphinxstyleemphasis{table} and \sphinxstyleemphasis{tail} options are not relevant for this style.

\sphinxAtStartPar
This pair style writes its information to {\hyperref[\detokenize{restart::doc}]{\sphinxcrossref{\DUrole{doc}{binary restart files}}}}, so pair\_style and pair\_coeff commands do not need
to be specified in an input script that reads a restart file.

\sphinxAtStartPar
These pair styles can only be used via the \sphinxstyleemphasis{pair} keyword of the
{\hyperref[\detokenize{run_style::doc}]{\sphinxcrossref{\DUrole{doc}{run\_style respa}}}} command.  They do not support the
\sphinxstyleemphasis{inner}, \sphinxstyleemphasis{middle}, \sphinxstyleemphasis{outer} keywords.


\bigskip\hrule\bigskip



\subsection{Restrictions}
\label{\detokenize{pair_cosine_squared:restrictions}}
\sphinxAtStartPar
The \sphinxstyleemphasis{cosine/squared} style is part of the EXTRA\sphinxhyphen{}PAIR package. It is only
enabled if LAMMPS is build with that package.  See the {\hyperref[\detokenize{Build_package::doc}]{\sphinxcrossref{\DUrole{doc}{Build package}}}} page for more info.


\subsection{Related commands}
\label{\detokenize{pair_cosine_squared:related-commands}}
\sphinxAtStartPar
{\hyperref[\detokenize{pair_coeff::doc}]{\sphinxcrossref{\DUrole{doc}{pair\_coeff}}}},
{\hyperref[\detokenize{pair_lj::doc}]{\sphinxcrossref{\DUrole{doc}{pair\_style lj/cut}}}}


\subsection{Default}
\label{\detokenize{pair_cosine_squared:default}}
\sphinxAtStartPar
none

\phantomsection\label{\detokenize{pair_cosine_squared:ckd}}
\sphinxAtStartPar
\sphinxstylestrong{(Cooke)} “Cooke, Kremer and Deserno, Phys. Rev. E, 72, 011506 (2005)”

\phantomsection\label{\detokenize{pair_cosine_squared:wca}}
\sphinxAtStartPar
\sphinxstylestrong{(Weeks)} “Weeks, Chandler and Andersen, J. Chem. Phys., 54, 5237 (1971)”

\sphinxstepscope

\index{pair\_style coul/cut@\spxentry{pair\_style coul/cut}}\ignorespaces 
\index{pair\_style coul/cut/gpu@\spxentry{pair\_style coul/cut/gpu}}\ignorespaces 
\index{pair\_style coul/cut/kk@\spxentry{pair\_style coul/cut/kk}}\ignorespaces 
\index{pair\_style coul/cut/omp@\spxentry{pair\_style coul/cut/omp}}\ignorespaces 
\index{pair\_style coul/cut/global@\spxentry{pair\_style coul/cut/global}}\ignorespaces 
\index{pair\_style coul/cut/global/omp@\spxentry{pair\_style coul/cut/global/omp}}\ignorespaces 
\index{pair\_style coul/debye@\spxentry{pair\_style coul/debye}}\ignorespaces 
\index{pair\_style coul/debye/gpu@\spxentry{pair\_style coul/debye/gpu}}\ignorespaces 
\index{pair\_style coul/debye/kk@\spxentry{pair\_style coul/debye/kk}}\ignorespaces 
\index{pair\_style coul/debye/omp@\spxentry{pair\_style coul/debye/omp}}\ignorespaces 
\index{pair\_style coul/dsf@\spxentry{pair\_style coul/dsf}}\ignorespaces 
\index{pair\_style coul/dsf/gpu@\spxentry{pair\_style coul/dsf/gpu}}\ignorespaces 
\index{pair\_style coul/dsf/kk@\spxentry{pair\_style coul/dsf/kk}}\ignorespaces 
\index{pair\_style coul/dsf/omp@\spxentry{pair\_style coul/dsf/omp}}\ignorespaces 
\index{pair\_style coul/exclude@\spxentry{pair\_style coul/exclude}}\ignorespaces 
\index{pair\_style coul/long@\spxentry{pair\_style coul/long}}\ignorespaces 
\index{pair\_style coul/long/omp@\spxentry{pair\_style coul/long/omp}}\ignorespaces 
\index{pair\_style coul/long/kk@\spxentry{pair\_style coul/long/kk}}\ignorespaces 
\index{pair\_style coul/long/gpu@\spxentry{pair\_style coul/long/gpu}}\ignorespaces 
\index{pair\_style coul/msm@\spxentry{pair\_style coul/msm}}\ignorespaces 
\index{pair\_style coul/msm/omp@\spxentry{pair\_style coul/msm/omp}}\ignorespaces 
\index{pair\_style coul/streitz@\spxentry{pair\_style coul/streitz}}\ignorespaces 
\index{pair\_style coul/wolf@\spxentry{pair\_style coul/wolf}}\ignorespaces 
\index{pair\_style coul/wolf/kk@\spxentry{pair\_style coul/wolf/kk}}\ignorespaces 
\index{pair\_style coul/wolf/omp@\spxentry{pair\_style coul/wolf/omp}}\ignorespaces 
\index{pair\_style tip4p/cut@\spxentry{pair\_style tip4p/cut}}\ignorespaces 
\index{pair\_style tip4p/cut/omp@\spxentry{pair\_style tip4p/cut/omp}}\ignorespaces 
\index{pair\_style tip4p/long@\spxentry{pair\_style tip4p/long}}\ignorespaces 
\index{pair\_style tip4p/long/omp@\spxentry{pair\_style tip4p/long/omp}}\ignorespaces 

\section{pair\_style coul/cut command}
\label{\detokenize{pair_coul:pair-style-coul-cut-command}}\label{\detokenize{pair_coul:index-28}}\label{\detokenize{pair_coul:index-27}}\label{\detokenize{pair_coul:index-26}}\label{\detokenize{pair_coul:index-25}}\label{\detokenize{pair_coul:index-24}}\label{\detokenize{pair_coul:index-23}}\label{\detokenize{pair_coul:index-22}}\label{\detokenize{pair_coul:index-21}}\label{\detokenize{pair_coul:index-20}}\label{\detokenize{pair_coul:index-19}}\label{\detokenize{pair_coul:index-18}}\label{\detokenize{pair_coul:index-17}}\label{\detokenize{pair_coul:index-16}}\label{\detokenize{pair_coul:index-15}}\label{\detokenize{pair_coul:index-14}}\label{\detokenize{pair_coul:index-13}}\label{\detokenize{pair_coul:index-12}}\label{\detokenize{pair_coul:index-11}}\label{\detokenize{pair_coul:index-10}}\label{\detokenize{pair_coul:index-9}}\label{\detokenize{pair_coul:index-8}}\label{\detokenize{pair_coul:index-7}}\label{\detokenize{pair_coul:index-6}}\label{\detokenize{pair_coul:index-5}}\label{\detokenize{pair_coul:index-4}}\label{\detokenize{pair_coul:index-3}}\label{\detokenize{pair_coul:index-2}}\label{\detokenize{pair_coul:index-1}}\label{\detokenize{pair_coul:index-0}}\label{\detokenize{pair_coul::doc}}
\sphinxAtStartPar
Accelerator Variants: \sphinxstyleemphasis{coul/cut/gpu}, \sphinxstyleemphasis{coul/cut/kk}, \sphinxstyleemphasis{coul/cut/omp}


\section{pair\_style coul/cut/global command}
\label{\detokenize{pair_coul:pair-style-coul-cut-global-command}}
\sphinxAtStartPar
Accelerator Variants: \sphinxstyleemphasis{coul/cut/omp}


\section{pair\_style coul/debye command}
\label{\detokenize{pair_coul:pair-style-coul-debye-command}}
\sphinxAtStartPar
Accelerator Variants: \sphinxstyleemphasis{coul/debye/gpu}, \sphinxstyleemphasis{coul/debye/kk}, \sphinxstyleemphasis{coul/debye/omp}


\section{pair\_style coul/dsf command}
\label{\detokenize{pair_coul:pair-style-coul-dsf-command}}
\sphinxAtStartPar
Accelerator Variants: \sphinxstyleemphasis{coul/dsf/gpu}, \sphinxstyleemphasis{coul/dsf/kk}, \sphinxstyleemphasis{coul/dsf/omp}


\section{pair\_style coul/exclude command}
\label{\detokenize{pair_coul:pair-style-coul-exclude-command}}

\section{pair\_style coul/long command}
\label{\detokenize{pair_coul:pair-style-coul-long-command}}
\sphinxAtStartPar
Accelerator Variants: \sphinxstyleemphasis{coul/long/omp}, \sphinxstyleemphasis{coul/long/kk}, \sphinxstyleemphasis{coul/long/gpu}


\section{pair\_style coul/msm command}
\label{\detokenize{pair_coul:pair-style-coul-msm-command}}
\sphinxAtStartPar
Accelerator Variants: \sphinxstyleemphasis{coul/msm/omp}


\section{pair\_style coul/streitz command}
\label{\detokenize{pair_coul:pair-style-coul-streitz-command}}

\section{pair\_style coul/wolf command}
\label{\detokenize{pair_coul:pair-style-coul-wolf-command}}
\sphinxAtStartPar
Accelerator Variants: \sphinxstyleemphasis{coul/wolf/kk}, \sphinxstyleemphasis{coul/wolf/omp}


\section{pair\_style tip4p/cut command}
\label{\detokenize{pair_coul:pair-style-tip4p-cut-command}}
\sphinxAtStartPar
Accelerator Variants: \sphinxstyleemphasis{tip4p/cut/omp}


\section{pair\_style tip4p/long command}
\label{\detokenize{pair_coul:pair-style-tip4p-long-command}}
\sphinxAtStartPar
Accelerator Variants: \sphinxstyleemphasis{tip4p/long/omp}


\subsection{Syntax}
\label{\detokenize{pair_coul:syntax}}
\begin{sphinxVerbatim}[commandchars=\\\{\}]
\PYG{k}{pair\PYGZus{}style}\PYG{+w}{ }\PYG{n}{coul}\PYG{o}{/}\PYG{n}{cut}\PYG{+w}{ }\PYG{n}{cutoff}
\PYG{k}{pair\PYGZus{}style}\PYG{+w}{ }\PYG{n}{coul}\PYG{o}{/}\PYG{n}{cut}\PYG{o}{/}\PYG{n}{global}\PYG{+w}{ }\PYG{n}{cutoff}
\PYG{k}{pair\PYGZus{}style}\PYG{+w}{ }\PYG{n}{coul}\PYG{o}{/}\PYG{n}{debye}\PYG{+w}{ }\PYG{n}{kappa}\PYG{+w}{ }\PYG{n}{cutoff}
\PYG{k}{pair\PYGZus{}style}\PYG{+w}{ }\PYG{n}{coul}\PYG{o}{/}\PYG{n}{dsf}\PYG{+w}{ }\PYG{n}{alpha}\PYG{+w}{ }\PYG{n}{cutoff}
\PYG{k}{pair\PYGZus{}style}\PYG{+w}{ }\PYG{n}{coul}\PYG{o}{/}\PYG{n}{exclude}\PYG{+w}{ }\PYG{n}{cutoff}
\PYG{k}{pair\PYGZus{}style}\PYG{+w}{ }\PYG{n}{coul}\PYG{o}{/}\PYG{n}{long}\PYG{+w}{ }\PYG{n}{cutoff}
\PYG{k}{pair\PYGZus{}style}\PYG{+w}{ }\PYG{n}{coul}\PYG{o}{/}\PYG{n}{wolf}\PYG{+w}{ }\PYG{n}{alpha}\PYG{+w}{ }\PYG{n}{cutoff}
\PYG{k}{pair\PYGZus{}style}\PYG{+w}{ }\PYG{n}{coul}\PYG{o}{/}\PYG{n}{streitz}\PYG{+w}{ }\PYG{n}{cutoff}\PYG{+w}{ }\PYG{n}{keyword}\PYG{+w}{ }\PYG{n}{alpha}

\PYG{o}{*}\PYG{+w}{ }\PYG{n}{cutoff}\PYG{+w}{ }\PYG{o}{=}\PYG{+w}{ }\PYG{n}{global}\PYG{+w}{ }\PYG{n}{cutoff}\PYG{+w}{ }\PYG{n}{for}\PYG{+w}{ }\PYG{n}{Coulombic}\PYG{+w}{ }\PYG{n}{interactions}
\PYG{o}{*}\PYG{+w}{ }\PYG{n}{kappa}\PYG{+w}{ }\PYG{o}{=}\PYG{+w}{ }\PYG{n}{Debye}\PYG{+w}{ }\PYG{n}{length}\PYG{+w}{ }\PYG{n+nv}{(}\PYG{n+nv}{inverse distance units}\PYG{n+nv}{)}
\PYG{o}{*}\PYG{+w}{ }\PYG{n}{alpha}\PYG{+w}{ }\PYG{o}{=}\PYG{+w}{ }\PYG{n}{damping}\PYG{+w}{ }\PYG{n}{parameter}\PYG{+w}{ }\PYG{n+nv}{(}\PYG{n+nv}{inverse distance units}\PYG{n+nv}{)}
\end{sphinxVerbatim}

\begin{sphinxVerbatim}[commandchars=\\\{\}]
\PYG{k}{pair\PYGZus{}style}\PYG{+w}{ }\PYG{n}{tip4p}\PYG{o}{/}\PYG{n}{cut}\PYG{+w}{ }\PYG{n}{otype}\PYG{+w}{ }\PYG{n}{htype}\PYG{+w}{ }\PYG{n}{btype}\PYG{+w}{ }\PYG{n}{atype}\PYG{+w}{ }\PYG{n}{qdist}\PYG{+w}{ }\PYG{n}{cutoff}
\PYG{k}{pair\PYGZus{}style}\PYG{+w}{ }\PYG{n}{tip4p}\PYG{o}{/}\PYG{n}{long}\PYG{+w}{ }\PYG{n}{otype}\PYG{+w}{ }\PYG{n}{htype}\PYG{+w}{ }\PYG{n}{btype}\PYG{+w}{ }\PYG{n}{atype}\PYG{+w}{ }\PYG{n}{qdist}\PYG{+w}{ }\PYG{n}{cutoff}

\PYG{o}{*}\PYG{+w}{ }\PYG{n}{otype},htype\PYG{+w}{ }\PYG{o}{=}\PYG{+w}{ }\PYG{n}{atom}\PYG{+w}{ }\PYG{n}{types}\PYG{+w}{ }\PYG{n+nv}{(}\PYG{n+nv}{numeric or type label}\PYG{n+nv}{)}\PYG{+w}{ }\PYG{n}{for}\PYG{+w}{ }\PYG{n}{TIP4P}\PYG{+w}{ }\PYG{n}{O}\PYG{+w}{ }\PYG{n}{and}\PYG{+w}{ }\PYG{n}{H}
\PYG{o}{*}\PYG{+w}{ }\PYG{n}{btype},atype\PYG{+w}{ }\PYG{o}{=}\PYG{+w}{ }\PYG{n}{bond}\PYG{+w}{ }\PYG{n}{and}\PYG{+w}{ }\PYG{n}{angle}\PYG{+w}{ }\PYG{n}{types}\PYG{+w}{ }\PYG{n+nv}{(}\PYG{n+nv}{numeric or type label}\PYG{n+nv}{)}\PYG{+w}{ }\PYG{n}{for}\PYG{+w}{ }\PYG{n}{TIP4P}\PYG{+w}{ }\PYG{n}{waters}
\PYG{o}{*}\PYG{+w}{ }\PYG{n}{qdist}\PYG{+w}{ }\PYG{o}{=}\PYG{+w}{ }\PYG{n}{distance}\PYG{+w}{ }\PYG{n}{from}\PYG{+w}{ }\PYG{n}{O}\PYG{+w}{ }\PYG{n}{atom}\PYG{+w}{ }\PYG{n}{to}\PYG{+w}{ }\PYG{n}{massless}\PYG{+w}{ }\PYG{n}{charge}\PYG{+w}{ }\PYG{n+nv}{(}\PYG{n+nv}{distance units}\PYG{n+nv}{)}
\end{sphinxVerbatim}


\subsection{Examples}
\label{\detokenize{pair_coul:examples}}
\begin{sphinxVerbatim}[commandchars=\\\{\}]
\PYG{k}{pair\PYGZus{}style}\PYG{+w}{ }\PYG{n}{coul}\PYG{o}{/}\PYG{n}{cut}\PYG{+w}{ }\PYG{l+m}{2.5}
\PYG{k}{pair\PYGZus{}coeff}\PYG{+w}{ }\PYG{o}{*}\PYG{+w}{ }\PYG{o}{*}
\PYG{k}{pair\PYGZus{}coeff}\PYG{+w}{ }\PYG{l+m}{2}\PYG{+w}{ }\PYG{l+m}{2}\PYG{+w}{ }\PYG{l+m}{3.5}

\PYG{k}{pair\PYGZus{}style}\PYG{+w}{ }\PYG{n}{coul}\PYG{o}{/}\PYG{n}{debye}\PYG{+w}{ }\PYG{l+m}{1.4}\PYG{+w}{ }\PYG{l+m}{3.0}
\PYG{k}{pair\PYGZus{}coeff}\PYG{+w}{ }\PYG{o}{*}\PYG{+w}{ }\PYG{o}{*}
\PYG{k}{pair\PYGZus{}coeff}\PYG{+w}{ }\PYG{l+m}{2}\PYG{+w}{ }\PYG{l+m}{2}\PYG{+w}{ }\PYG{l+m}{3.5}

\PYG{k}{pair\PYGZus{}style}\PYG{+w}{ }\PYG{n}{coul}\PYG{o}{/}\PYG{n}{dsf}\PYG{+w}{ }\PYG{l+m}{0.05}\PYG{+w}{ }\PYG{l+m}{10.0}
\PYG{k}{pair\PYGZus{}coeff}\PYG{+w}{ }\PYG{o}{*}\PYG{+w}{ }\PYG{o}{*}

\PYG{k}{pair\PYGZus{}style}\PYG{+w}{ }\PYG{n}{hybrid}\PYG{o}{/}\PYG{n}{overlay}\PYG{+w}{ }\PYG{n}{coul}\PYG{o}{/}\PYG{n}{exclude}\PYG{+w}{ }\PYG{l+m}{10.0}\PYG{+w}{ }\PYG{n}{...}
\PYG{k}{pair\PYGZus{}coeff}\PYG{+w}{ }\PYG{o}{*}\PYG{+w}{ }\PYG{o}{*}\PYG{+w}{ }\PYG{n}{coul}\PYG{o}{/}\PYG{n}{exclude}

\PYG{k}{pair\PYGZus{}style}\PYG{+w}{ }\PYG{n}{coul}\PYG{o}{/}\PYG{n}{long}\PYG{+w}{ }\PYG{l+m}{10.0}
\PYG{k}{pair\PYGZus{}coeff}\PYG{+w}{ }\PYG{o}{*}\PYG{+w}{ }\PYG{o}{*}

\PYG{k}{pair\PYGZus{}style}\PYG{+w}{ }\PYG{n}{coul}\PYG{o}{/}\PYG{n}{msm}\PYG{+w}{ }\PYG{l+m}{10.0}
\PYG{k}{pair\PYGZus{}coeff}\PYG{+w}{ }\PYG{o}{*}\PYG{+w}{ }\PYG{o}{*}

\PYG{k}{pair\PYGZus{}style}\PYG{+w}{ }\PYG{n}{coul}\PYG{o}{/}\PYG{n}{wolf}\PYG{+w}{ }\PYG{l+m}{0.2}\PYG{+w}{ }\PYG{l+m}{9.0}
\PYG{k}{pair\PYGZus{}coeff}\PYG{+w}{ }\PYG{o}{*}\PYG{+w}{ }\PYG{o}{*}

\PYG{k}{pair\PYGZus{}style}\PYG{+w}{ }\PYG{n}{coul}\PYG{o}{/}\PYG{n}{streitz}\PYG{+w}{ }\PYG{l+m}{12.0}\PYG{+w}{ }\PYG{n}{ewald}
\PYG{k}{pair\PYGZus{}style}\PYG{+w}{ }\PYG{n}{coul}\PYG{o}{/}\PYG{n}{streitz}\PYG{+w}{ }\PYG{l+m}{12.0}\PYG{+w}{ }\PYG{n}{wolf}\PYG{+w}{ }\PYG{l+m}{0.30}
\PYG{k}{pair\PYGZus{}coeff}\PYG{+w}{ }\PYG{o}{*}\PYG{+w}{ }\PYG{o}{*}\PYG{+w}{ }\PYG{n}{AlO.streitz}\PYG{+w}{ }\PYG{n}{Al}\PYG{+w}{ }\PYG{n}{O}

\PYG{k}{pair\PYGZus{}style}\PYG{+w}{ }\PYG{n}{tip4p}\PYG{o}{/}\PYG{n}{cut}\PYG{+w}{ }\PYG{l+m}{1}\PYG{+w}{ }\PYG{l+m}{2}\PYG{+w}{ }\PYG{l+m}{7}\PYG{+w}{ }\PYG{l+m}{8}\PYG{+w}{ }\PYG{l+m}{0.15}\PYG{+w}{ }\PYG{l+m}{12.0}
\PYG{k}{pair\PYGZus{}coeff}\PYG{+w}{ }\PYG{o}{*}\PYG{+w}{ }\PYG{o}{*}

\PYG{k}{pair\PYGZus{}style}\PYG{+w}{ }\PYG{n}{tip4p}\PYG{o}{/}\PYG{n}{long}\PYG{+w}{ }\PYG{l+m}{1}\PYG{+w}{ }\PYG{l+m}{2}\PYG{+w}{ }\PYG{l+m}{7}\PYG{+w}{ }\PYG{l+m}{8}\PYG{+w}{ }\PYG{l+m}{0.15}\PYG{+w}{ }\PYG{l+m}{10.0}
\PYG{k}{pair\PYGZus{}coeff}\PYG{+w}{ }\PYG{o}{*}\PYG{+w}{ }\PYG{o}{*}

\PYG{k}{pair\PYGZus{}style}\PYG{+w}{ }\PYG{n}{tip4p}\PYG{o}{/}\PYG{n}{cut}\PYG{+w}{ }\PYG{n}{OW}\PYG{+w}{ }\PYG{n}{HW}\PYG{+w}{ }\PYG{n}{HW}\PYG{o}{\PYGZhy{}}\PYG{n}{OW}\PYG{+w}{ }\PYG{n}{HW}\PYG{o}{\PYGZhy{}}\PYG{n}{OW}\PYG{o}{\PYGZhy{}}\PYG{n}{HW}\PYG{+w}{ }\PYG{l+m}{0.15}\PYG{+w}{ }\PYG{l+m}{12.0}
\PYG{k}{labelmap}\PYG{+w}{ }\PYG{n}{atom}\PYG{+w}{ }\PYG{l+m}{1}\PYG{+w}{ }\PYG{n}{OW}\PYG{+w}{ }\PYG{l+m}{2}\PYG{+w}{ }\PYG{n}{HW}
\PYG{k}{labelmap}\PYG{+w}{ }\PYG{n}{bond}\PYG{+w}{ }\PYG{l+m}{1}\PYG{+w}{ }\PYG{n}{HW}\PYG{o}{\PYGZhy{}}\PYG{n}{OW}
\PYG{k}{labelmap}\PYG{+w}{ }\PYG{n}{angle}\PYG{+w}{ }\PYG{l+m}{1}\PYG{+w}{ }\PYG{n}{HW}\PYG{o}{\PYGZhy{}}\PYG{n}{OW}\PYG{o}{\PYGZhy{}}\PYG{n}{HW}
\PYG{k}{pair\PYGZus{}coeff}\PYG{+w}{ }\PYG{o}{*}\PYG{+w}{ }\PYG{o}{*}
\end{sphinxVerbatim}


\subsection{Description}
\label{\detokenize{pair_coul:description}}
\sphinxAtStartPar
The \sphinxstyleemphasis{coul/cut} style computes the standard Coulombic interaction
potential given by
\begin{equation*}
\begin{split}E = \frac{C q_i q_j}{\epsilon  r} \qquad r < r_c\end{split}
\end{equation*}
\sphinxAtStartPar
where C is an energy\sphinxhyphen{}conversion constant, Qi and Qj are the charges on
the two atoms, and \(\epsilon\) is the dielectric constant which can be set
by the {\hyperref[\detokenize{dielectric::doc}]{\sphinxcrossref{\DUrole{doc}{dielectric}}}} command.  The cutoff \(r_c\) truncates
the interaction distance.

\sphinxAtStartPar
Pair style \sphinxstyleemphasis{coul/cut/global} computes the same Coulombic interactions
as style \sphinxstyleemphasis{coul/cut} except that it allows only a single global cutoff
and thus makes it compatible for use in combination with long\sphinxhyphen{}range
coulomb styles in {\hyperref[\detokenize{pair_hybrid::doc}]{\sphinxcrossref{\DUrole{doc}{hybrid pair styles}}}}.


\bigskip\hrule\bigskip


\sphinxAtStartPar
Style \sphinxstyleemphasis{coul/debye} adds an additional exp() damping factor to the
Coulombic term, given by
\begin{equation*}
\begin{split}E = \frac{C q_i q_j}{\epsilon  r} \exp(- \kappa r) \qquad r < r_c\end{split}
\end{equation*}
\sphinxAtStartPar
where \(\kappa\) is the Debye length.  This potential is another way to
mimic the screening effect of a polar solvent.


\bigskip\hrule\bigskip


\sphinxAtStartPar
Style \sphinxstyleemphasis{coul/dsf} computes Coulombic interactions via the damped
shifted force model described in {\hyperref[\detokenize{pair_coul:fennell1}]{\sphinxcrossref{\DUrole{std,std-ref}{Fennell}}}}, given by:
\begin{equation*}
\begin{split}E = q_iq_j \left[ \frac{\mbox{erfc} (\alpha r)}{r} -  \frac{\mbox{erfc} (\alpha r_c)}{r_c} +
\left( \frac{\mbox{erfc} (\alpha r_c)}{r_c^2} +  \frac{2\alpha}{\sqrt{\pi}}\frac{\exp (-\alpha^2    r^2_c)}{r_c} \right)(r-r_c) \right] \qquad r < r_c\end{split}
\end{equation*}
\sphinxAtStartPar
where \(\alpha\) is the damping parameter and \sphinxstyleemphasis{erfc()} is the
complementary error\sphinxhyphen{}function. The potential corrects issues in the Wolf
model (described below) to provide consistent forces and energies (the
Wolf potential is not differentiable at the cutoff) and smooth decay to
zero.


\bigskip\hrule\bigskip


\sphinxAtStartPar
Style \sphinxstyleemphasis{coul/wolf} computes Coulombic interactions via the Wolf
summation method, described in {\hyperref[\detokenize{pair_coul:wolf1}]{\sphinxcrossref{\DUrole{std,std-ref}{Wolf}}}}, given by:
\begin{equation*}
\begin{split}E_i = \frac{1}{2} \sum_{j \neq i}
\frac{q_i q_j {\rm erfc}(\alpha r_{ij})}{r_{ij}} +
\frac{1}{2} \sum_{j \neq i}
\frac{q_i q_j {\rm erf}(\alpha r_{ij})}{r_{ij}} \qquad r < r_c\end{split}
\end{equation*}
\sphinxAtStartPar
where \(\alpha\) is the damping parameter, and \sphinxstyleemphasis{erf()} and \sphinxstyleemphasis{erfc()}
are error\sphinxhyphen{}function and complementary error\sphinxhyphen{}function terms.  This
potential is essentially a short\sphinxhyphen{}range, spherically\sphinxhyphen{}truncated,
charge\sphinxhyphen{}neutralized, shifted, pairwise \sphinxstyleemphasis{1/r} summation.  With a
manipulation of adding and subtracting a self term (for i = j) to the
first and second term on the right\sphinxhyphen{}hand\sphinxhyphen{}side, respectively, and a small
enough \(\alpha\) damping parameter, the second term shrinks and the
potential becomes a rapidly\sphinxhyphen{}converging real\sphinxhyphen{}space summation.  With a
long enough cutoff and small enough \(\alpha\) parameter, the energy
and forces calculated by the Wolf summation method approach those of the
Ewald sum.  So it is a means of getting effective long\sphinxhyphen{}range
interactions with a short\sphinxhyphen{}range potential.


\bigskip\hrule\bigskip


\sphinxAtStartPar
Style \sphinxstyleemphasis{coul/streitz} is the Coulomb pair interaction defined as part of
the Streitz\sphinxhyphen{}Mintmire potential, as described in {\hyperref[\detokenize{pair_coul:streitz2}]{\sphinxcrossref{\DUrole{std,std-ref}{this paper}}}}, in which charge distribution about an atom is modeled as a
Slater 1\sphinxstyleemphasis{s} orbital.  More details can be found in the referenced
paper.  To fully reproduce the published Streitz\sphinxhyphen{}Mintmire potential,
which is a variable charge potential, style \sphinxstyleemphasis{coul/streitz} must be used
with {\hyperref[\detokenize{pair_eam::doc}]{\sphinxcrossref{\DUrole{doc}{pair\_style eam/alloy}}}} (or some other short\sphinxhyphen{}range
potential that has been parameterized appropriately) via the
{\hyperref[\detokenize{pair_hybrid::doc}]{\sphinxcrossref{\DUrole{doc}{pair\_style hybrid/overlay}}}} command.  Likewise,
charge equilibration must be performed via the {\hyperref[\detokenize{fix_qeq::doc}]{\sphinxcrossref{\DUrole{doc}{fix qeq/slater}}}} command. For example:

\begin{sphinxVerbatim}[commandchars=\\\{\}]
\PYG{k}{pair\PYGZus{}style}\PYG{+w}{ }\PYG{n}{hybrid}\PYG{o}{/}\PYG{n}{overlay}\PYG{+w}{ }\PYG{n}{coul}\PYG{o}{/}\PYG{n}{streitz}\PYG{+w}{ }\PYG{l+m}{12.0}\PYG{+w}{ }\PYG{n}{wolf}\PYG{+w}{ }\PYG{l+m}{0.31}\PYG{+w}{ }\PYG{n}{eam}\PYG{o}{/}\PYG{n}{alloy}
\PYG{k}{pair\PYGZus{}coeff}\PYG{+w}{ }\PYG{o}{*}\PYG{+w}{ }\PYG{o}{*}\PYG{+w}{ }\PYG{n}{coul}\PYG{o}{/}\PYG{n}{streitz}\PYG{+w}{ }\PYG{n}{AlO.streitz}\PYG{+w}{ }\PYG{n}{Al}\PYG{+w}{ }\PYG{n}{O}
\PYG{k}{pair\PYGZus{}coeff}\PYG{+w}{ }\PYG{o}{*}\PYG{+w}{ }\PYG{o}{*}\PYG{+w}{ }\PYG{n}{eam}\PYG{o}{/}\PYG{n}{alloy}\PYG{+w}{ }\PYG{n}{AlO.eam.alloy}\PYG{+w}{ }\PYG{n}{Al}\PYG{+w}{ }\PYG{n}{O}
\PYG{k}{fix }\PYG{n+nv+nvIdentifier}{1}\PYG{+w}{ }\PYG{n+nv+nvIdentifier}{all}\PYG{+w}{ }\PYG{n}{qeq}\PYG{o}{/}\PYG{n}{slater}\PYG{+w}{ }\PYG{l+m}{1}\PYG{+w}{ }\PYG{l+m}{12.0}\PYG{+w}{ }\PYG{l+m}{1.0e\PYGZhy{}6}\PYG{+w}{ }\PYG{l+m}{100}\PYG{+w}{ }\PYG{n}{coul}\PYG{o}{/}\PYG{n}{streitz}
\end{sphinxVerbatim}

\sphinxAtStartPar
The keyword \sphinxstyleemphasis{wolf} in the coul/streitz command denotes computing
Coulombic interactions via Wolf summation.  An additional damping
parameter is required for the Wolf summation, as described for the
coul/wolf potential above.  Alternatively, Coulombic interactions can
be computed via an Ewald summation.  For example:

\begin{sphinxVerbatim}[commandchars=\\\{\}]
\PYG{k}{pair\PYGZus{}style}\PYG{+w}{ }\PYG{n}{hybrid}\PYG{o}{/}\PYG{n}{overlay}\PYG{+w}{ }\PYG{n}{coul}\PYG{o}{/}\PYG{n}{streitz}\PYG{+w}{ }\PYG{l+m}{12.0}\PYG{+w}{ }\PYG{n}{ewald}\PYG{+w}{ }\PYG{n}{eam}\PYG{o}{/}\PYG{n}{alloy}
\PYG{k}{kspace\PYGZus{}style}\PYG{+w}{ }\PYG{n}{ewald}\PYG{+w}{ }\PYG{l+m}{1e\PYGZhy{}6}
\end{sphinxVerbatim}

\sphinxAtStartPar
Keyword \sphinxstyleemphasis{ewald} does not need a damping parameter, but a
{\hyperref[\detokenize{kspace_style::doc}]{\sphinxcrossref{\DUrole{doc}{kspace\_style}}}} must be defined, which can be style
\sphinxstyleemphasis{ewald} or \sphinxstyleemphasis{pppm}.  The Ewald method was used in Streitz and
Mintmire’s original paper, but a Wolf summation offers a speed\sphinxhyphen{}up in
some cases.

\sphinxAtStartPar
For the fix qeq/slater command, the \sphinxstyleemphasis{qfile} can be a filename that
contains QEq parameters as discussed on the {\hyperref[\detokenize{fix_qeq::doc}]{\sphinxcrossref{\DUrole{doc}{fix qeq}}}}
command doc page.  Alternatively \sphinxstyleemphasis{qfile} can be replaced by
“coul/streitz”, in which case the fix will extract QEq parameters from
the coul/streitz pair style itself.

\sphinxAtStartPar
See the examples/strietz directory for an example input script that
uses the Streitz\sphinxhyphen{}Mintmire potential.  The potentials directory has the
AlO.eam.alloy and AlO.streitz potential files used by the example.

\sphinxAtStartPar
Note that the Streiz\sphinxhyphen{}Mintmire potential is generally used for oxides,
but there is no conceptual problem with extending it to nitrides and
carbides (such as SiC, TiN).  Pair coul/strietz used by itself or with
any other pair style such as EAM, MEAM, Tersoff, or LJ in
hybrid/overlay mode.  To do this, you would need to provide a
Streitz\sphinxhyphen{}Mintmire parameterization for the material being modeled.


\bigskip\hrule\bigskip


\sphinxAtStartPar
Pair style \sphinxstyleemphasis{coul/exclude} computes Coulombic interactions like \sphinxstyleemphasis{coul/cut}
but \sphinxstylestrong{only} applies them to excluded pairs using a scaling factor
of \(\gamma - 1.0\) with \(\gamma\) being the factor assigned
to that excluded pair via the {\hyperref[\detokenize{special_bonds::doc}]{\sphinxcrossref{\DUrole{doc}{special\_bonds coul}}}}
setting. With this it is possible to treat Coulomb interactions for
molecular systems with {\hyperref[\detokenize{kspace_style::doc}]{\sphinxcrossref{\DUrole{doc}{kspace style scafacos}}}},
which always computes the \sphinxstyleemphasis{full} Coulomb interactions without exclusions.
Pair style \sphinxstyleemphasis{coul/exclude} will then \sphinxstyleemphasis{subtract} the excluded interactions
accordingly. So to achieve the same forces as with \sphinxcode{\sphinxupquote{pair\_style lj/cut/coul/long 12.0}}
with \sphinxcode{\sphinxupquote{kspace\_style pppm 1.0e\sphinxhyphen{}6}}, one would use
\sphinxcode{\sphinxupquote{pair\_style hybrid/overlay lj/cut 12.0 coul/exclude 12.0}} with
\sphinxcode{\sphinxupquote{kspace\_style scafacos p3m 1.0e\sphinxhyphen{}6}}.

\sphinxAtStartPar
Styles \sphinxstyleemphasis{coul/long} and \sphinxstyleemphasis{coul/msm} compute the same Coulombic
interactions as style \sphinxstyleemphasis{coul/cut} except that an additional damping
factor is applied so it can be used in conjunction with the
{\hyperref[\detokenize{kspace_style::doc}]{\sphinxcrossref{\DUrole{doc}{kspace\_style}}}} command and its \sphinxstyleemphasis{ewald} or \sphinxstyleemphasis{pppm}
option.  The Coulombic cutoff specified for this style means that
pairwise interactions within this distance are computed directly;
interactions outside that distance are computed in reciprocal space.

\sphinxAtStartPar
Styles \sphinxstyleemphasis{tip4p/cut} and \sphinxstyleemphasis{tip4p/long} implement the Coulomb part of
the TIP4P water model of {\hyperref[\detokenize{pair_coul:jorgensen3}]{\sphinxcrossref{\DUrole{std,std-ref}{(Jorgensen)}}}}, which introduces
a massless site located a short distance away from the oxygen atom
along the bisector of the HOH angle.  The atomic types of the oxygen and
hydrogen atoms, the bond and angle types for OH and HOH interactions,
and the distance to the massless charge site are specified as
pair\_style arguments.  Style \sphinxstyleemphasis{tip4p/cut} uses a global cutoff for
Coulomb interactions; style \sphinxstyleemphasis{tip4p/long} is for use with a long\sphinxhyphen{}range
Coulombic solver (Ewald or PPPM).

\begin{sphinxadmonition}{note}{Note:}
\sphinxAtStartPar
For each TIP4P water molecule in your system, the atom IDs for
the O and 2 H atoms must be consecutive, with the O atom first.  This
is to enable LAMMPS to “find” the 2 H atoms associated with each O
atom.  For example, if the atom ID of an O atom in a TIP4P water
molecule is 500, then its 2 H atoms must have IDs 501 and 502.
\end{sphinxadmonition}

\begin{sphinxadmonition}{note}{Note:}
\sphinxAtStartPar
If using type labels, the type labels must be defined before calling
the {\hyperref[\detokenize{pair_coeff::doc}]{\sphinxcrossref{\DUrole{doc}{pair\_coeff}}}} command.
\end{sphinxadmonition}

\sphinxAtStartPar
See the {\hyperref[\detokenize{Howto_tip4p::doc}]{\sphinxcrossref{\DUrole{doc}{Howto tip4p}}}} page for more information
on how to use the TIP4P pair styles and lists of parameters to set.
Note that the neighbor list cutoff for Coulomb interactions is
effectively extended by a distance 2*qdist when using the TIP4P pair
style, to account for the offset distance of the fictitious charges on
O atoms in water molecules.  Thus it is typically best in an
efficiency sense to use a LJ cutoff \textgreater{}= Coulombic cutoff + 2*qdist, to
shrink the size of the neighbor list.  This leads to slightly larger
cost for the long\sphinxhyphen{}range calculation, so you can test the trade\sphinxhyphen{}off for
your model.


\bigskip\hrule\bigskip


\sphinxAtStartPar
Note that these potentials are designed to be combined with other pair
potentials via the {\hyperref[\detokenize{pair_hybrid::doc}]{\sphinxcrossref{\DUrole{doc}{pair\_style hybrid/overlay}}}}
command.  This is because they have no repulsive core.  Hence if they
are used by themselves, there will be no repulsion to keep two
oppositely charged particles from moving arbitrarily close to each
other.

\sphinxAtStartPar
The following coefficients must be defined for each pair of atoms
types via the {\hyperref[\detokenize{pair_coeff::doc}]{\sphinxcrossref{\DUrole{doc}{pair\_coeff}}}} command as in the examples
above, or in the data or restart files read by the
{\hyperref[\detokenize{read_data::doc}]{\sphinxcrossref{\DUrole{doc}{read\_data}}}} or {\hyperref[\detokenize{read_restart::doc}]{\sphinxcrossref{\DUrole{doc}{read\_restart}}}}
commands, or by mixing as described below:
\begin{itemize}
\item {} 
\sphinxAtStartPar
cutoff (distance units)

\end{itemize}

\sphinxAtStartPar
For \sphinxstyleemphasis{coul/cut} and \sphinxstyleemphasis{coul/debye} the cutoff coefficient is optional.
If it is not used (as in some of the examples above), the default
global value specified in the pair\_style command is used.

\sphinxAtStartPar
For \sphinxstyleemphasis{coul/cut/global}, \sphinxstyleemphasis{coul/long} and \sphinxstyleemphasis{coul/msm} no cutoff can be
specified for an individual I,J type pair via the pair\_coeff command.
All type pairs use the same global Coulomb cutoff specified in the
pair\_style command.


\bigskip\hrule\bigskip


\sphinxAtStartPar
Styles with a \sphinxstyleemphasis{gpu}, \sphinxstyleemphasis{intel}, \sphinxstyleemphasis{kk}, \sphinxstyleemphasis{omp}, or \sphinxstyleemphasis{opt} suffix are
functionally the same as the corresponding style without the suffix.
They have been optimized to run faster, depending on your available
hardware, as discussed on the {\hyperref[\detokenize{Speed_packages::doc}]{\sphinxcrossref{\DUrole{doc}{Accelerator packages}}}}
page.  The accelerated styles take the same arguments and should
produce the same results, except for round\sphinxhyphen{}off and precision issues.

\sphinxAtStartPar
These accelerated styles are part of the GPU, INTEL, KOKKOS,
OPENMP, and OPT packages, respectively.  They are only enabled if
LAMMPS was built with those packages.  See the {\hyperref[\detokenize{Build_package::doc}]{\sphinxcrossref{\DUrole{doc}{Build package}}}} page for more info.

\sphinxAtStartPar
You can specify the accelerated styles explicitly in your input script
by including their suffix, or you can use the {\hyperref[\detokenize{Run_options::doc}]{\sphinxcrossref{\DUrole{doc}{\sphinxhyphen{}suffix command\sphinxhyphen{}line switch}}}} when you invoke LAMMPS, or you can use the
{\hyperref[\detokenize{suffix::doc}]{\sphinxcrossref{\DUrole{doc}{suffix}}}} command in your input script.

\sphinxAtStartPar
See the {\hyperref[\detokenize{Speed_packages::doc}]{\sphinxcrossref{\DUrole{doc}{Accelerator packages}}}} page for more
instructions on how to use the accelerated styles effectively.


\bigskip\hrule\bigskip



\subsection{Mixing, shift, table, tail correction, restart, rRESPA info}
\label{\detokenize{pair_coul:mixing-shift-table-tail-correction-restart-rrespa-info}}
\sphinxAtStartPar
For atom type pairs I,J and I != J, the cutoff distance for the
\sphinxstyleemphasis{coul/cut} style can be mixed.  The default mix value is \sphinxstyleemphasis{geometric}.
See the “pair\_modify” command for details.

\sphinxAtStartPar
The {\hyperref[\detokenize{pair_modify::doc}]{\sphinxcrossref{\DUrole{doc}{pair\_modify}}}} shift option is not relevant
for these pair styles.

\sphinxAtStartPar
The \sphinxstyleemphasis{coul/long} style supports the {\hyperref[\detokenize{pair_modify::doc}]{\sphinxcrossref{\DUrole{doc}{pair\_modify}}}}
table option for tabulation of the short\sphinxhyphen{}range portion of the
long\sphinxhyphen{}range Coulombic interaction.

\sphinxAtStartPar
These pair styles do not support the {\hyperref[\detokenize{pair_modify::doc}]{\sphinxcrossref{\DUrole{doc}{pair\_modify}}}}
tail option for adding long\sphinxhyphen{}range tail corrections to energy and
pressure.

\sphinxAtStartPar
These pair styles write their information to {\hyperref[\detokenize{restart::doc}]{\sphinxcrossref{\DUrole{doc}{binary restart files}}}}, so pair\_style and pair\_coeff commands do not need
to be specified in an input script that reads a restart file.

\sphinxAtStartPar
These pair styles can only be used via the \sphinxstyleemphasis{pair} keyword of the
{\hyperref[\detokenize{run_style::doc}]{\sphinxcrossref{\DUrole{doc}{run\_style respa}}}} command.  They do not support the
\sphinxstyleemphasis{inner}, \sphinxstyleemphasis{middle}, \sphinxstyleemphasis{outer} keywords.


\bigskip\hrule\bigskip



\subsection{Restrictions}
\label{\detokenize{pair_coul:restrictions}}
\sphinxAtStartPar
The \sphinxstyleemphasis{coul/long}, \sphinxstyleemphasis{coul/msm}, \sphinxstyleemphasis{coul/streitz}, and \sphinxstyleemphasis{tip4p/long} styles are
part of the KSPACE package.  The \sphinxstyleemphasis{coul/cut/global}, \sphinxstyleemphasis{coul/exclude} styles are
part of the EXTRA\sphinxhyphen{}PAIR package.  The \sphinxstyleemphasis{tip4p/cut} style is part of the MOLECULE
package.  A pair style is only enabled if LAMMPS was built with its
corresponding package.  See the {\hyperref[\detokenize{Build_package::doc}]{\sphinxcrossref{\DUrole{doc}{Build package}}}}
doc page for more info.


\subsection{Related commands}
\label{\detokenize{pair_coul:related-commands}}
\sphinxAtStartPar
{\hyperref[\detokenize{pair_coeff::doc}]{\sphinxcrossref{\DUrole{doc}{pair\_coeff}}}}, {\hyperref[\detokenize{pair_hybrid::doc}]{\sphinxcrossref{\DUrole{doc}{pair\_style, hybrid/overlay}}}}, {\hyperref[\detokenize{kspace_style::doc}]{\sphinxcrossref{\DUrole{doc}{kspace\_style}}}}


\subsection{Default}
\label{\detokenize{pair_coul:default}}
\sphinxAtStartPar
none


\bigskip\hrule\bigskip

\phantomsection\label{\detokenize{pair_coul:wolf1}}
\sphinxAtStartPar
\sphinxstylestrong{(Wolf)} D. Wolf, P. Keblinski, S. R. Phillpot, J. Eggebrecht, J Chem
Phys, 110, 8254 (1999).

\phantomsection\label{\detokenize{pair_coul:fennell1}}
\sphinxAtStartPar
\sphinxstylestrong{(Fennell)} C. J. Fennell, J. D. Gezelter, J Chem Phys, 124,
234104 (2006).

\phantomsection\label{\detokenize{pair_coul:streitz2}}
\sphinxAtStartPar
\sphinxstylestrong{(Streitz)} F. H. Streitz, J. W. Mintmire, Phys Rev B, 50, 11996\sphinxhyphen{}12003
(1994).

\phantomsection\label{\detokenize{pair_coul:jorgensen3}}
\sphinxAtStartPar
\sphinxstylestrong{(Jorgensen)} Jorgensen, Chandrasekhar, Madura, Impey, Klein, J Chem
Phys, 79, 926 (1983).

\sphinxstepscope

\index{pair\_style coul/diel@\spxentry{pair\_style coul/diel}}\ignorespaces 
\index{pair\_style coul/diel/omp@\spxentry{pair\_style coul/diel/omp}}\ignorespaces 

\section{pair\_style coul/diel command}
\label{\detokenize{pair_coul_diel:pair-style-coul-diel-command}}\label{\detokenize{pair_coul_diel:index-1}}\label{\detokenize{pair_coul_diel:index-0}}\label{\detokenize{pair_coul_diel::doc}}
\sphinxAtStartPar
Accelerator Variants: \sphinxstyleemphasis{coul/diel/omp}


\subsection{Syntax}
\label{\detokenize{pair_coul_diel:syntax}}
\begin{sphinxVerbatim}[commandchars=\\\{\}]
\PYG{k}{pair\PYGZus{}style}\PYG{+w}{ }\PYG{n}{coul}\PYG{o}{/}\PYG{n}{diel}\PYG{+w}{ }\PYG{n}{cutoff}
\end{sphinxVerbatim}

\sphinxAtStartPar
cutoff = global cutoff (distance units)


\subsection{Examples}
\label{\detokenize{pair_coul_diel:examples}}
\begin{sphinxVerbatim}[commandchars=\\\{\}]
\PYG{k}{pair\PYGZus{}style}\PYG{+w}{ }\PYG{n}{coul}\PYG{o}{/}\PYG{n}{diel}\PYG{+w}{ }\PYG{l+m}{3.5}
\PYG{k}{pair\PYGZus{}coeff}\PYG{+w}{ }\PYG{l+m}{1}\PYG{+w}{ }\PYG{l+m}{4}\PYG{+w}{ }\PYG{l+m}{78}\PYG{n}{.}\PYG{+w}{ }\PYG{l+m}{1.375}\PYG{+w}{ }\PYG{l+m}{0.112}
\end{sphinxVerbatim}


\subsection{Description}
\label{\detokenize{pair_coul_diel:description}}
\sphinxAtStartPar
Style \sphinxstyleemphasis{coul/diel} computes a Coulomb correction for implicit solvent
ion interactions in which the dielectric permittivity is distance dependent.
The dielectric permittivity \(\epsilon_D(r)\) connects to limiting regimes:
One limit is defined by a small dielectric permittivity (close to vacuum)
at or close to contact separation between the ions. At larger separations
the dielectric permittivity reaches a bulk value used in the regular Coulomb
interaction coul/long or coul/cut.
The transition is modeled by a hyperbolic function which is incorporated
in the Coulomb correction term for small ion separations as follows
\begin{equation*}
\begin{split}E  = & \frac{Cq_iq_j}{\epsilon r} \left( \frac{\epsilon}{\epsilon_D(r)}-1\right)                       \qquad r < r_c \\
\epsilon_D(r)  = & \frac{5.2+\epsilon}{2} +  \frac{\epsilon-5.2}{2}\tanh\left(\frac{r-r_{me}}{\sigma_e}\right)\end{split}
\end{equation*}
\sphinxAtStartPar
where \(r_{me}\) is the inflection point of \(\epsilon_D(r)\) and \(\sigma_e\) is a slope
defining length scale. C is the same Coulomb conversion factor as in the
pair\_styles coul/cut, coul/long, and coul/debye. In this way the Coulomb
interaction between ions is corrected at small distances r. The lower
limit of \(\epsilon_D(r \to 0) = 5.2\) due to dielectric saturation {\hyperref[\detokenize{pair_coul_diel:stiles}]{\sphinxcrossref{\DUrole{std,std-ref}{(Stiles)}}}}
while the Coulomb interaction reaches its bulk limit by setting
\(\epsilon_D(r \to \infty) = \epsilon\), the bulk value of the solvent which is 78
for water at 298K.

\sphinxAtStartPar
Examples of the use of this type of Coulomb interaction include implicit
solvent simulations of salt ions
{\hyperref[\detokenize{pair_coul_diel:lenart1}]{\sphinxcrossref{\DUrole{std,std-ref}{(Lenart)}}}} and of ionic surfactants {\hyperref[\detokenize{pair_coul_diel:jusufi1}]{\sphinxcrossref{\DUrole{std,std-ref}{(Jusufi)}}}}.
Note that this potential is only reasonable for implicit solvent simulations
and in combination with coul/cut or coul/long. It is also usually combined
with gauss/cut, see {\hyperref[\detokenize{pair_coul_diel:lenart1}]{\sphinxcrossref{\DUrole{std,std-ref}{(Lenart)}}}} or {\hyperref[\detokenize{pair_coul_diel:jusufi1}]{\sphinxcrossref{\DUrole{std,std-ref}{(Jusufi)}}}}.

\sphinxAtStartPar
The following coefficients must be defined for each pair of atom
types via the {\hyperref[\detokenize{pair_coeff::doc}]{\sphinxcrossref{\DUrole{doc}{pair\_coeff}}}} command as in the example
above, or in the data file or restart files read by the
{\hyperref[\detokenize{read_data::doc}]{\sphinxcrossref{\DUrole{doc}{read\_data}}}} or {\hyperref[\detokenize{read_restart::doc}]{\sphinxcrossref{\DUrole{doc}{read\_restart}}}}
commands:
\begin{itemize}
\item {} 
\sphinxAtStartPar
\(\epsilon\) (no units)

\item {} 
\sphinxAtStartPar
\(r_{me}\) (distance units)

\item {} 
\sphinxAtStartPar
\(\sigma_e\) (distance units)

\end{itemize}

\sphinxAtStartPar
The global cutoff (\(r_c\)) specified in the pair\_style command is used.


\bigskip\hrule\bigskip



\subsection{Mixing, shift, table, tail correction, restart, rRESPA info}
\label{\detokenize{pair_coul_diel:mixing-shift-table-tail-correction-restart-rrespa-info}}
\sphinxAtStartPar
This pair style does not support parameter mixing. Coefficients must
be given explicitly for each type of particle pairs.

\sphinxAtStartPar
This pair style supports the {\hyperref[\detokenize{pair_modify::doc}]{\sphinxcrossref{\DUrole{doc}{pair\_modify}}}} shift
option for the energy of the Gauss\sphinxhyphen{}potential portion of the pair
interaction.

\sphinxAtStartPar
The {\hyperref[\detokenize{pair_modify::doc}]{\sphinxcrossref{\DUrole{doc}{pair\_modify}}}} table option is not relevant
for this pair style.

\sphinxAtStartPar
This pair style does not support the {\hyperref[\detokenize{pair_modify::doc}]{\sphinxcrossref{\DUrole{doc}{pair\_modify}}}}
tail option for adding long\sphinxhyphen{}range tail corrections to energy and
pressure.

\sphinxAtStartPar
This pair style can only be used via the \sphinxstyleemphasis{pair} keyword of the
{\hyperref[\detokenize{run_style::doc}]{\sphinxcrossref{\DUrole{doc}{run\_style respa}}}} command.  It does not support the
\sphinxstyleemphasis{inner}, \sphinxstyleemphasis{middle}, \sphinxstyleemphasis{outer} keywords.


\subsection{Restrictions}
\label{\detokenize{pair_coul_diel:restrictions}}
\sphinxAtStartPar
This style is part of the EXTRA\sphinxhyphen{}PAIR package.  It is only enabled if
LAMMPS was built with that package.  See the {\hyperref[\detokenize{Build_package::doc}]{\sphinxcrossref{\DUrole{doc}{Build package}}}} page for more info.


\subsection{Related commands}
\label{\detokenize{pair_coul_diel:related-commands}}
\sphinxAtStartPar
{\hyperref[\detokenize{pair_coeff::doc}]{\sphinxcrossref{\DUrole{doc}{pair\_coeff}}}}
{\hyperref[\detokenize{pair_gauss::doc}]{\sphinxcrossref{\DUrole{doc}{pair\_style gauss/cut}}}}


\subsection{Default}
\label{\detokenize{pair_coul_diel:default}}
\sphinxAtStartPar
none


\bigskip\hrule\bigskip

\phantomsection\label{\detokenize{pair_coul_diel:stiles}}
\sphinxAtStartPar
\sphinxstylestrong{(Stiles)} Stiles, Hubbard, and Kayser, J Chem Phys, 77,
6189 (1982).

\phantomsection\label{\detokenize{pair_coul_diel:lenart1}}
\sphinxAtStartPar
\sphinxstylestrong{(Lenart)} Lenart, Jusufi, and Panagiotopoulos, J Chem Phys, 126,
044509 (2007).

\phantomsection\label{\detokenize{pair_coul_diel:jusufi1}}
\sphinxAtStartPar
\sphinxstylestrong{(Jusufi)} Jusufi, Hynninen, and Panagiotopoulos, J Phys Chem B, 112,
13783 (2008).

\sphinxstepscope

\index{pair\_style coul/shield@\spxentry{pair\_style coul/shield}}\ignorespaces 

\section{pair\_style coul/shield command}
\label{\detokenize{pair_coul_shield:pair-style-coul-shield-command}}\label{\detokenize{pair_coul_shield:index-0}}\label{\detokenize{pair_coul_shield::doc}}

\subsection{Syntax}
\label{\detokenize{pair_coul_shield:syntax}}
\begin{sphinxVerbatim}[commandchars=\\\{\}]
\PYG{k}{pair\PYGZus{}style}\PYG{+w}{ }\PYG{n}{coul}\PYG{o}{/}\PYG{n}{shield}\PYG{+w}{ }\PYG{n}{cutoff}\PYG{+w}{ }\PYG{n}{tap\PYGZus{}flag}
\end{sphinxVerbatim}
\begin{itemize}
\item {} 
\sphinxAtStartPar
cutoff = global cutoff (distance units)

\item {} 
\sphinxAtStartPar
tap\_flag = 0/1 to turn off/on the taper function

\end{itemize}


\subsection{Examples}
\label{\detokenize{pair_coul_shield:examples}}
\begin{sphinxVerbatim}[commandchars=\\\{\}]
\PYG{k}{pair\PYGZus{}style}\PYG{+w}{ }\PYG{n}{coul}\PYG{o}{/}\PYG{n}{shield}\PYG{+w}{ }\PYG{l+m}{16.0}\PYG{+w}{ }\PYG{l+m}{1}
\PYG{k}{pair\PYGZus{}coeff}\PYG{+w}{ }\PYG{l+m}{1}\PYG{+w}{ }\PYG{l+m}{2}\PYG{+w}{ }\PYG{l+m}{0.70}
\end{sphinxVerbatim}


\subsection{Description}
\label{\detokenize{pair_coul_shield:description}}
\sphinxAtStartPar
Style \sphinxstyleemphasis{coul/shield} computes a Coulomb interaction for boron and
nitrogen atoms located in different layers of hexagonal boron
nitride. This potential is designed be used in combination with
the pair style {\hyperref[\detokenize{pair_ilp_graphene_hbn::doc}]{\sphinxcrossref{\DUrole{doc}{ilp/graphene/hbn}}}}

\begin{sphinxadmonition}{note}{Note:}
\sphinxAtStartPar
This potential is intended for electrostatic interactions between
two different layers of hexagonal boron nitride. Therefore, to avoid
interaction within the same layers, each layer should have a separate
molecule id and is recommended to use the “full” atom style, so that
charge and molecule ID information is included.
\end{sphinxadmonition}
\begin{equation*}
\begin{split}E      = & \frac{1}{2} \sum_i \sum_{j \neq i} V_{ij} \\
V_{ij} = & {\rm Tap}(r_{ij})\frac{\kappa q_i q_j}{\sqrt[3]{r_{ij}^3+(1/\lambda_{ij})^3}}\\
{\rm Tap}(r_{ij}) = & 20\left ( \frac{r_{ij}}{R_{cut}} \right )^7 -
                       70\left ( \frac{r_{ij}}{R_{cut}} \right )^6 +
                       84\left ( \frac{r_{ij}}{R_{cut}} \right )^5 -
                       35\left ( \frac{r_{ij}}{R_{cut}} \right )^4 + 1\end{split}
\end{equation*}
\sphinxAtStartPar
Where Tap(\(r_{ij}\)) is the taper function which provides a continuous cutoff
(up to third derivative) for inter\sphinxhyphen{}atomic separations larger than \(r_c\)
{\hyperref[\detokenize{pair_coul_shield:leven3}]{\sphinxcrossref{\DUrole{std,std-ref}{(Leven1)}}}}, {\hyperref[\detokenize{pair_coul_shield:leven4}]{\sphinxcrossref{\DUrole{std,std-ref}{(Leven2)}}}} and {\hyperref[\detokenize{pair_coul_shield:maaravi1}]{\sphinxcrossref{\DUrole{std,std-ref}{(Maaravi)}}}}.
Here \(\lambda\) is the shielding parameter that
eliminates the short\sphinxhyphen{}range singularity of the classical mono\sphinxhyphen{}polar
electrostatic interaction expression {\hyperref[\detokenize{pair_coul_shield:maaravi1}]{\sphinxcrossref{\DUrole{std,std-ref}{(Maaravi)}}}}.

\sphinxAtStartPar
The shielding parameter \(\lambda\) (1/distance units) must be defined for
each pair of atom types via the {\hyperref[\detokenize{pair_coeff::doc}]{\sphinxcrossref{\DUrole{doc}{pair\_coeff}}}} command as
in the example above, or in the data file or restart files read by the
{\hyperref[\detokenize{read_data::doc}]{\sphinxcrossref{\DUrole{doc}{read\_data}}}} or {\hyperref[\detokenize{read_restart::doc}]{\sphinxcrossref{\DUrole{doc}{read\_restart}}}} commands:

\sphinxAtStartPar
The global cutoff (\(r_c\)) specified in the pair\_style command is used.


\bigskip\hrule\bigskip



\subsection{Mixing, shift, table, tail correction, restart, rRESPA info}
\label{\detokenize{pair_coul_shield:mixing-shift-table-tail-correction-restart-rrespa-info}}
\sphinxAtStartPar
This pair style does not support parameter mixing. Coefficients must
be given explicitly for each type of particle pairs.

\sphinxAtStartPar
The {\hyperref[\detokenize{pair_modify::doc}]{\sphinxcrossref{\DUrole{doc}{pair\_modify}}}} \sphinxstyleemphasis{table} option is not relevant
for this pair style.

\sphinxAtStartPar
This pair style does not support the {\hyperref[\detokenize{pair_modify::doc}]{\sphinxcrossref{\DUrole{doc}{pair\_modify}}}}
\sphinxstyleemphasis{tail} option for adding long\sphinxhyphen{}range tail corrections to energy and
pressure.

\sphinxAtStartPar
This pair style can only be used via the \sphinxstyleemphasis{pair} keyword of the
{\hyperref[\detokenize{run_style::doc}]{\sphinxcrossref{\DUrole{doc}{run\_style respa}}}} command.  It does not support the
\sphinxstyleemphasis{inner}, \sphinxstyleemphasis{middle}, \sphinxstyleemphasis{outer} keywords.


\subsection{Restrictions}
\label{\detokenize{pair_coul_shield:restrictions}}
\sphinxAtStartPar
This pair style is part of the INTERLAYER package.  It is only enabled
if LAMMPS was built with that package.  See the {\hyperref[\detokenize{Build_package::doc}]{\sphinxcrossref{\DUrole{doc}{Build package}}}} page for more info.


\subsection{Related commands}
\label{\detokenize{pair_coul_shield:related-commands}}
\sphinxAtStartPar
{\hyperref[\detokenize{pair_coeff::doc}]{\sphinxcrossref{\DUrole{doc}{pair\_coeff}}}}
{\hyperref[\detokenize{pair_ilp_graphene_hbn::doc}]{\sphinxcrossref{\DUrole{doc}{pair\_style ilp/graphene/hbn}}}}


\subsection{Default}
\label{\detokenize{pair_coul_shield:default}}
\sphinxAtStartPar
tap\_flag = 1


\bigskip\hrule\bigskip

\phantomsection\label{\detokenize{pair_coul_shield:leven3}}
\sphinxAtStartPar
\sphinxstylestrong{(Leven1)} I. Leven, I. Azuri, L. Kronik and O. Hod, J. Chem. Phys. 140, 104106 (2014).

\phantomsection\label{\detokenize{pair_coul_shield:leven4}}
\sphinxAtStartPar
\sphinxstylestrong{(Leven2)} I. Leven et al, J. Chem.Theory Comput. 12, 2896\sphinxhyphen{}905 (2016).

\phantomsection\label{\detokenize{pair_coul_shield:maaravi1}}
\sphinxAtStartPar
\sphinxstylestrong{(Maaravi)} T. Maaravi et al, J. Phys. Chem. C 121, 22826\sphinxhyphen{}22835 (2017).

\sphinxstepscope

\index{pair\_style coul/slater@\spxentry{pair\_style coul/slater}}\ignorespaces 
\index{pair\_style coul/slater/cut@\spxentry{pair\_style coul/slater/cut}}\ignorespaces 
\index{pair\_style coul/slater/long@\spxentry{pair\_style coul/slater/long}}\ignorespaces 
\index{pair\_style coul/slater/long/gpu@\spxentry{pair\_style coul/slater/long/gpu}}\ignorespaces 

\section{pair\_style coul/slater command}
\label{\detokenize{pair_coul_slater:pair-style-coul-slater-command}}\label{\detokenize{pair_coul_slater:index-3}}\label{\detokenize{pair_coul_slater:index-2}}\label{\detokenize{pair_coul_slater:index-1}}\label{\detokenize{pair_coul_slater:index-0}}\label{\detokenize{pair_coul_slater::doc}}

\section{pair\_style coul/slater/cut command}
\label{\detokenize{pair_coul_slater:pair-style-coul-slater-cut-command}}

\section{pair\_style coul/slater/long command}
\label{\detokenize{pair_coul_slater:pair-style-coul-slater-long-command}}
\sphinxAtStartPar
Accelerator Variants: \sphinxstyleemphasis{coul/slater/long/gpu}


\subsection{Syntax}
\label{\detokenize{pair_coul_slater:syntax}}
\begin{sphinxVerbatim}[commandchars=\\\{\}]
\PYG{k}{pair\PYGZus{}style}\PYG{+w}{ }\PYG{n}{coul}\PYG{o}{/}\PYG{n}{slater}\PYG{o}{/}\PYG{n}{cut}\PYG{+w}{ }\PYG{n}{lambda}\PYG{+w}{ }\PYG{n}{cutoff}
\PYG{k}{pair\PYGZus{}style}\PYG{+w}{ }\PYG{n}{coul}\PYG{o}{/}\PYG{n}{slater}\PYG{o}{/}\PYG{n}{long}\PYG{+w}{ }\PYG{n}{lambda}\PYG{+w}{ }\PYG{n}{cutoff}
\end{sphinxVerbatim}

\sphinxAtStartPar
lambda = decay length of the charge (distance units)
cutoff = cutoff (distance units)


\subsection{Examples}
\label{\detokenize{pair_coul_slater:examples}}
\begin{sphinxVerbatim}[commandchars=\\\{\}]
\PYG{k}{pair\PYGZus{}style}\PYG{+w}{ }\PYG{n}{coul}\PYG{o}{/}\PYG{n}{slater}\PYG{o}{/}\PYG{n}{cut}\PYG{+w}{ }\PYG{l+m}{1.0}\PYG{+w}{ }\PYG{l+m}{3.5}
\PYG{k}{pair\PYGZus{}coeff}\PYG{+w}{ }\PYG{o}{*}\PYG{+w}{ }\PYG{o}{*}
\PYG{k}{pair\PYGZus{}coeff}\PYG{+w}{ }\PYG{l+m}{2}\PYG{+w}{ }\PYG{l+m}{2}\PYG{+w}{ }\PYG{l+m}{2.5}

\PYG{k}{pair\PYGZus{}style}\PYG{+w}{ }\PYG{n}{coul}\PYG{o}{/}\PYG{n}{slater}\PYG{o}{/}\PYG{n}{long}\PYG{+w}{ }\PYG{l+m}{1.0}\PYG{+w}{ }\PYG{l+m}{12.0}
\PYG{k}{pair\PYGZus{}coeff}\PYG{+w}{ }\PYG{o}{*}\PYG{+w}{ }\PYG{o}{*}
\PYG{k}{pair\PYGZus{}coeff}\PYG{+w}{ }\PYG{l+m}{1}\PYG{+w}{ }\PYG{l+m}{1}\PYG{+w}{ }\PYG{l+m}{5.0}
\end{sphinxVerbatim}


\subsection{Description}
\label{\detokenize{pair_coul_slater:description}}
\sphinxAtStartPar
Styles \sphinxstyleemphasis{coul/slater/*} compute electrostatic interactions in mesoscopic models
which employ potentials without explicit excluded\sphinxhyphen{}volume interactions.
The goal is to prevent artificial ionic pair formation by including a charge
distribution in the Coulomb potential, following the formulation of
{\hyperref[\detokenize{pair_coul_slater:melchor}]{\sphinxcrossref{\DUrole{std,std-ref}{(Melchor)}}}}:
\begin{equation*}
\begin{split}E  =  \frac{Cq_iq_j}{\epsilon r} \left( 1- \left( 1 + \frac{r_{ij}}{\lambda} exp\left( -2r_{ij}/\lambda \right) \right) \right)                       \qquad r < r_c\end{split}
\end{equation*}
\sphinxAtStartPar
where \(r_c\) is the cutoff distance and \(\lambda\) is the decay length of the charge.
C is the same Coulomb conversion factor as in the pair\_styles coul/cut and coul/long. In this way the Coulomb
interaction between ions is corrected at small distances r.
For the \sphinxstyleemphasis{coul/slater/cut} style, the potential energy for distances larger than the cutoff is zero,
while for the \sphinxstyleemphasis{coul/slater/long}, the long\sphinxhyphen{}range interactions are computed either by the Ewald or the PPPM technique.

\sphinxAtStartPar
Phenomena that can be captured at a mesoscopic level using this type of electrostatic
interactions include the formation of polyelectrolyte\sphinxhyphen{}surfactant aggregates,
charge stabilization of colloidal suspensions, and the formation of
complexes driven by charged species in biological systems. {\hyperref[\detokenize{pair_coul_slater:vaiwala}]{\sphinxcrossref{\DUrole{std,std-ref}{(Vaiwala)}}}}.

\sphinxAtStartPar
The cutoff distance is optional. If it is not used,
the default global value specified in the pair\_style command is used.
For each pair of atom types, a specific cutoff distance can be defined via the {\hyperref[\detokenize{pair_coeff::doc}]{\sphinxcrossref{\DUrole{doc}{pair\_coeff}}}} command as in the example
above, or in the data file or restart files read by the
{\hyperref[\detokenize{read_data::doc}]{\sphinxcrossref{\DUrole{doc}{read\_data}}}} or {\hyperref[\detokenize{read_restart::doc}]{\sphinxcrossref{\DUrole{doc}{read\_restart}}}}
commands:
\begin{itemize}
\item {} 
\sphinxAtStartPar
\(r_c\) (distance units)

\end{itemize}

\sphinxAtStartPar
The global decay length of the charge (\(\lambda\)) specified in the pair\_style command is used for all pairs.


\bigskip\hrule\bigskip



\subsection{Mixing, shift, table, tail correction, restart, rRESPA info}
\label{\detokenize{pair_coul_slater:mixing-shift-table-tail-correction-restart-rrespa-info}}
\sphinxAtStartPar
For atom type pairs I,J and I != J, the cutoff distance for the
\sphinxstyleemphasis{coul/slater} styles can be mixed.  The default mix value is \sphinxstyleemphasis{geometric}.
See the “pair\_modify” command for details.

\sphinxAtStartPar
The {\hyperref[\detokenize{pair_modify::doc}]{\sphinxcrossref{\DUrole{doc}{pair\_modify}}}} shift and table options are not relevant
for these pair styles.

\sphinxAtStartPar
These pair styles do not support the {\hyperref[\detokenize{pair_modify::doc}]{\sphinxcrossref{\DUrole{doc}{pair\_modify}}}}
tail option for adding long\sphinxhyphen{}range tail corrections to energy and
pressure.

\sphinxAtStartPar
These pair styles write their information to {\hyperref[\detokenize{restart::doc}]{\sphinxcrossref{\DUrole{doc}{binary restart files}}}}, so pair\_style and pair\_coeff commands do not need
to be specified in an input script that reads a restart file.

\sphinxAtStartPar
This pair style can only be used via the \sphinxstyleemphasis{pair} keyword of the
{\hyperref[\detokenize{run_style::doc}]{\sphinxcrossref{\DUrole{doc}{run\_style respa}}}} command.  It does not support the
\sphinxstyleemphasis{inner}, \sphinxstyleemphasis{middle}, \sphinxstyleemphasis{outer} keywords.


\subsection{Restrictions}
\label{\detokenize{pair_coul_slater:restrictions}}
\sphinxAtStartPar
The  \sphinxstyleemphasis{coul/slater/long} style requires the long\sphinxhyphen{}range solvers included in the KSPACE package.

\sphinxAtStartPar
These styles are part of the EXTRA\sphinxhyphen{}PAIR package.  They are only enabled if
LAMMPS was built with that package.  See the {\hyperref[\detokenize{Build_package::doc}]{\sphinxcrossref{\DUrole{doc}{Build package}}}} page for more info.


\subsection{Related commands}
\label{\detokenize{pair_coul_slater:related-commands}}
\sphinxAtStartPar
{\hyperref[\detokenize{pair_coeff::doc}]{\sphinxcrossref{\DUrole{doc}{pair\_coeff}}}}, {\hyperref[\detokenize{pair_hybrid::doc}]{\sphinxcrossref{\DUrole{doc}{pair\_style, hybrid/overlay}}}}, {\hyperref[\detokenize{kspace_style::doc}]{\sphinxcrossref{\DUrole{doc}{kspace\_style}}}}


\subsection{Default}
\label{\detokenize{pair_coul_slater:default}}
\sphinxAtStartPar
none


\bigskip\hrule\bigskip

\phantomsection\label{\detokenize{pair_coul_slater:melchor}}
\sphinxAtStartPar
\sphinxstylestrong{(Melchor)} Gonzalez\sphinxhyphen{}Melchor, Mayoral, Velazquez, and Alejandre, J Chem Phys, 125, 224107 (2006).

\phantomsection\label{\detokenize{pair_coul_slater:vaiwala}}
\sphinxAtStartPar
\sphinxstylestrong{(Vaiwala)} Vaiwala, Jadhav, and Thaokar, J Chem Phys, 146, 124904 (2017).

\sphinxstepscope

\index{pair\_style coul/tt@\spxentry{pair\_style coul/tt}}\ignorespaces 

\section{pair\_style coul/tt command}
\label{\detokenize{pair_coul_tt:pair-style-coul-tt-command}}\label{\detokenize{pair_coul_tt:index-0}}\label{\detokenize{pair_coul_tt::doc}}

\subsection{Syntax}
\label{\detokenize{pair_coul_tt:syntax}}
\begin{sphinxVerbatim}[commandchars=\\\{\}]
\PYG{k}{pair\PYGZus{}style}\PYG{+w}{ }\PYG{n}{style}\PYG{+w}{ }\PYG{n}{args}
\end{sphinxVerbatim}
\begin{itemize}
\item {} 
\sphinxAtStartPar
style = \sphinxstyleemphasis{coul/tt}

\item {} 
\sphinxAtStartPar
args = list of arguments for a particular style

\end{itemize}
\begin{sphinxalltt}
\sphinxstyleemphasis{coul/tt} args = n cutoff
  n = degree of polynomial
  cutoff = global cutoff (distance units)
\end{sphinxalltt}


\subsection{Examples}
\label{\detokenize{pair_coul_tt:examples}}
\begin{sphinxVerbatim}[commandchars=\\\{\}]
\PYG{k}{pair\PYGZus{}style}\PYG{+w}{ }\PYG{n}{hybrid}\PYG{o}{/}\PYG{n}{overlay}\PYG{+w}{ }\PYG{n}{...}\PYG{+w}{ }\PYG{n}{coul}\PYG{o}{/}\PYG{n}{tt}\PYG{+w}{ }\PYG{l+m}{4}\PYG{+w}{ }\PYG{l+m}{12.0}
\PYG{k}{pair\PYGZus{}coeff}\PYG{+w}{ }\PYG{l+m}{1}\PYG{+w}{ }\PYG{l+m}{2}\PYG{+w}{  }\PYG{n}{coul}\PYG{o}{/}\PYG{n}{tt}\PYG{+w}{ }\PYG{l+m}{4.5}\PYG{+w}{ }\PYG{l+m}{1.0}
\PYG{k}{pair\PYGZus{}coeff}\PYG{+w}{ }\PYG{l+m}{1}\PYG{+w}{ }\PYG{l+m}{2}\PYG{+w}{  }\PYG{n}{coul}\PYG{o}{/}\PYG{n}{tt}\PYG{+w}{ }\PYG{l+m}{4.0}\PYG{+w}{ }\PYG{l+m}{1.0}\PYG{+w}{ }\PYG{l+m}{4}\PYG{+w}{ }\PYG{l+m}{12.0}
\PYG{k}{pair\PYGZus{}coeff}\PYG{+w}{ }\PYG{l+m}{1}\PYG{+w}{ }\PYG{l+m}{3}\PYG{o}{*}\PYG{+w}{ }\PYG{n}{coul}\PYG{o}{/}\PYG{n}{tt}\PYG{+w}{ }\PYG{l+m}{4.5}\PYG{+w}{ }\PYG{l+m}{1.0}\PYG{+w}{ }\PYG{l+m}{4}
\end{sphinxVerbatim}

\sphinxAtStartPar
Example input scripts available: examples/PACKAGES/drude


\subsection{Description}
\label{\detokenize{pair_coul_tt:description}}
\sphinxAtStartPar
The \sphinxstyleemphasis{coul/tt} pair style is meant to be used with force fields that
include explicit polarization through Drude dipoles.

\sphinxAtStartPar
The \sphinxstyleemphasis{coul/tt} pair style should be used as a sub\sphinxhyphen{}style within in the
{\hyperref[\detokenize{pair_hybrid::doc}]{\sphinxcrossref{\DUrole{doc}{pair\_style hybrid/overlay}}}} command, in conjunction with a
main pair style including Coulomb interactions and \sphinxstyleemphasis{thole} pair style,
or with \sphinxstyleemphasis{lj/cut/thole/long} pair style that is equivalent to the combination
of preceding two.

\sphinxAtStartPar
The \sphinxstyleemphasis{coul/tt} pair styles compute the charge\sphinxhyphen{}dipole Coulomb interaction damped
at short distances by a function
\begin{equation*}
\begin{split}f_{n,ij}(r) = 1 - c_{ij} \cdot e^{-b_{ij} r} \sum_{k=0}^n \frac{(b_{ij} r)^k}{k!}\end{split}
\end{equation*}
\sphinxAtStartPar
This function results from an adaptation to the Coulomb interaction {\hyperref[\detokenize{pair_coul_tt:salanne1}]{\sphinxcrossref{\DUrole{std,std-ref}{(Salanne)}}}} of the damping function originally proposed
by {\hyperref[\detokenize{pair_coul_tt:tangtoennies1}]{\sphinxcrossref{\DUrole{std,std-ref}{Tang Toennies}}}} for van der Waals interactions.

\sphinxAtStartPar
The polynomial takes the degree 4 for damping the Coulomb interaction.
The parameters \(b_{ij}\) and \(c_{ij}\) could be determined from
first\sphinxhyphen{}principle calculations for small, mainly mono\sphinxhyphen{}atomic, ions {\hyperref[\detokenize{pair_coul_tt:salanne1}]{\sphinxcrossref{\DUrole{std,std-ref}{(Salanne)}}}}, or else treated as empirical for large molecules.

\sphinxAtStartPar
In pair styles with Drude induced dipoles, this damping function is typically
applied to the interactions between a Drude charge (either \(q_{D,i}\) on
a Drude particle or \(-q_{D,i}\) on the respective
Drude core)) and a charge on a non\sphinxhyphen{}polarizable atom, \(q_{j}\).

\sphinxAtStartPar
The Tang\sphinxhyphen{}Toennies function could also be used to damp electrostatic
interactions between the (non\sphinxhyphen{}polarizable part of the) charge of a core,
\(q_{i}-q_{D,i}\), and the Drude charge of another, \(-q_{D,j}\).
The \(b_{ij}\) and \(c_{ij}\) are equal  to \(b_{ji}\) and
\(c_{ji}\) in the case of core\sphinxhyphen{}core interactions.

\sphinxAtStartPar
For pair\_style \sphinxstyleemphasis{coul/tt}, the following coefficients must be defined for
each pair of atoms types via the {\hyperref[\detokenize{pair_coeff::doc}]{\sphinxcrossref{\DUrole{doc}{pair\_coeff}}}} command
as in the example above.
\begin{itemize}
\item {} 
\sphinxAtStartPar
\(b_{ij}\)

\item {} 
\sphinxAtStartPar
\(c_{ij}\)

\item {} 
\sphinxAtStartPar
degree of polynomial (positive integer)

\item {} 
\sphinxAtStartPar
cutoff (distance units)

\end{itemize}

\sphinxAtStartPar
The last two coefficients are optional.  If not specified the global
degree of the polynomial or the global cutoff specified in the pair\_style
command are used. In order to specify a cutoff (forth argument), the degree of
the polynomial (third argument) must also be specified.


\bigskip\hrule\bigskip



\subsection{Mixing, shift, table, tail correction, restart, rRESPA info}
\label{\detokenize{pair_coul_tt:mixing-shift-table-tail-correction-restart-rrespa-info}}
\sphinxAtStartPar
The \sphinxstyleemphasis{coul/tt} pair style does not support mixing.  Thus, coefficients
for all I,J pairs must be specified explicitly.


\subsection{Restrictions}
\label{\detokenize{pair_coul_tt:restrictions}}
\sphinxAtStartPar
These pair styles are part of the DRUDE package. They are only
enabled if LAMMPS was built with that package. See the {\hyperref[\detokenize{Build_package::doc}]{\sphinxcrossref{\DUrole{doc}{Build package}}}} page for more info.

\sphinxAtStartPar
This pair\_style should currently not be used with the {\hyperref[\detokenize{dihedral_charmm::doc}]{\sphinxcrossref{\DUrole{doc}{charmm dihedral
style}}}} if the latter has non\sphinxhyphen{}zero 1\sphinxhyphen{}4 weighting
factors. This is because the \sphinxstyleemphasis{coul/tt} pair style does not know which
pairs are 1\sphinxhyphen{}4 partners of which dihedrals.


\subsection{Related commands}
\label{\detokenize{pair_coul_tt:related-commands}}
\sphinxAtStartPar
{\hyperref[\detokenize{fix_drude::doc}]{\sphinxcrossref{\DUrole{doc}{fix drude}}}}, {\hyperref[\detokenize{fix_langevin_drude::doc}]{\sphinxcrossref{\DUrole{doc}{fix langevin/drude}}}},
{\hyperref[\detokenize{fix_drude_transform::doc}]{\sphinxcrossref{\DUrole{doc}{fix drude/transform}}}},
{\hyperref[\detokenize{compute_temp_drude::doc}]{\sphinxcrossref{\DUrole{doc}{compute temp/drude}}}},
{\hyperref[\detokenize{pair_thole::doc}]{\sphinxcrossref{\DUrole{doc}{pair\_style thole}}}}


\subsection{Default}
\label{\detokenize{pair_coul_tt:default}}
\sphinxAtStartPar
none


\bigskip\hrule\bigskip

\phantomsection\label{\detokenize{pair_coul_tt:thole3}}
\sphinxAtStartPar
\sphinxstylestrong{(Thole)} Chem Phys, 59, 341 (1981).

\phantomsection\label{\detokenize{pair_coul_tt:salanne1}}
\sphinxAtStartPar
\sphinxstylestrong{(Salanne)} Salanne, Rotenberg, Jahn, Vuilleumier, Simon, Christian and Madden, Theor Chem Acc, 131, 1143 (2012).

\phantomsection\label{\detokenize{pair_coul_tt:tangtoennies1}}
\sphinxAtStartPar
\sphinxstylestrong{(Tang and Toennies)} J Chem Phys, 80, 3726 (1984).

\sphinxstepscope

\index{pair\_style born/coul/dsf/cs@\spxentry{pair\_style born/coul/dsf/cs}}\ignorespaces 
\index{pair\_style born/coul/long/cs@\spxentry{pair\_style born/coul/long/cs}}\ignorespaces 
\index{pair\_style born/coul/long/cs/gpu@\spxentry{pair\_style born/coul/long/cs/gpu}}\ignorespaces 
\index{pair\_style born/coul/wolf/cs@\spxentry{pair\_style born/coul/wolf/cs}}\ignorespaces 
\index{pair\_style born/coul/wolf/cs/gpu@\spxentry{pair\_style born/coul/wolf/cs/gpu}}\ignorespaces 
\index{pair\_style buck/coul/long/cs@\spxentry{pair\_style buck/coul/long/cs}}\ignorespaces 
\index{pair\_style coul/long/cs@\spxentry{pair\_style coul/long/cs}}\ignorespaces 
\index{pair\_style coul/long/cs/gpu@\spxentry{pair\_style coul/long/cs/gpu}}\ignorespaces 
\index{pair\_style coul/wolf/cs@\spxentry{pair\_style coul/wolf/cs}}\ignorespaces 
\index{pair\_style lj/cut/coul/long/cs@\spxentry{pair\_style lj/cut/coul/long/cs}}\ignorespaces 
\index{pair\_style lj/class2/coul/long/cs@\spxentry{pair\_style lj/class2/coul/long/cs}}\ignorespaces 

\section{pair\_style born/coul/dsf/cs command}
\label{\detokenize{pair_cs:pair-style-born-coul-dsf-cs-command}}\label{\detokenize{pair_cs:index-10}}\label{\detokenize{pair_cs:index-9}}\label{\detokenize{pair_cs:index-8}}\label{\detokenize{pair_cs:index-7}}\label{\detokenize{pair_cs:index-6}}\label{\detokenize{pair_cs:index-5}}\label{\detokenize{pair_cs:index-4}}\label{\detokenize{pair_cs:index-3}}\label{\detokenize{pair_cs:index-2}}\label{\detokenize{pair_cs:index-1}}\label{\detokenize{pair_cs:index-0}}\label{\detokenize{pair_cs::doc}}

\section{pair\_style born/coul/long/cs command}
\label{\detokenize{pair_cs:pair-style-born-coul-long-cs-command}}
\sphinxAtStartPar
Accelerator Variants: \sphinxstyleemphasis{born/coul/long/cs/gpu}


\section{pair\_style born/coul/wolf/cs command}
\label{\detokenize{pair_cs:pair-style-born-coul-wolf-cs-command}}
\sphinxAtStartPar
Accelerator Variants: \sphinxstyleemphasis{born/coul/wolf/cs/gpu}


\section{pair\_style buck/coul/long/cs command}
\label{\detokenize{pair_cs:pair-style-buck-coul-long-cs-command}}

\section{pair\_style coul/long/cs command}
\label{\detokenize{pair_cs:pair-style-coul-long-cs-command}}
\sphinxAtStartPar
Accelerator Variants: \sphinxstyleemphasis{coul/long/cs/gpu}


\section{pair\_style coul/wolf/cs command}
\label{\detokenize{pair_cs:pair-style-coul-wolf-cs-command}}

\section{pair\_style lj/cut/coul/long/cs command}
\label{\detokenize{pair_cs:pair-style-lj-cut-coul-long-cs-command}}

\section{pair\_style lj/class2/coul/long/cs command}
\label{\detokenize{pair_cs:pair-style-lj-class2-coul-long-cs-command}}

\subsection{Syntax}
\label{\detokenize{pair_cs:syntax}}
\begin{sphinxVerbatim}[commandchars=\\\{\}]
\PYG{k}{pair\PYGZus{}style}\PYG{+w}{ }\PYG{n}{style}\PYG{+w}{ }\PYG{n}{args}
\end{sphinxVerbatim}
\begin{itemize}
\item {} 
\sphinxAtStartPar
style = \sphinxstyleemphasis{born/coul/dsf/cs} or \sphinxstyleemphasis{born/coul/long/cs} or \sphinxstyleemphasis{born/coul/wolf/cs} or \sphinxstyleemphasis{buck/coul/long/cs} or \sphinxstyleemphasis{coul/long/cs} or \sphinxstyleemphasis{coul/wolf/cs} or \sphinxstyleemphasis{lj/cut/coul/long/cs} or \sphinxstyleemphasis{lj/class2/coul/long/cs}

\item {} 
\sphinxAtStartPar
args = list of arguments for a particular style

\end{itemize}
\begin{sphinxalltt}
\sphinxstyleemphasis{born/coul/dsf/cs} args = alpha cutoff (cutoff2)
  alpha = damping parameter (inverse distance units)
  cutoff = global cutoff for non\sphinxhyphen{}Coulombic (and Coulombic if only 1 arg) (distance units)
  cutoff2 = global cutoff for Coulombic (distance units)
\sphinxstyleemphasis{born/coul/long/cs} args = cutoff (cutoff2)
  cutoff = global cutoff for non\sphinxhyphen{}Coulombic (and Coulombic if only 1 arg) (distance units)
  cutoff2 = global cutoff for Coulombic (optional) (distance units)
\sphinxstyleemphasis{born/coul/wolf/cs} args = alpha cutoff (cutoff2)
  alpha = damping parameter (inverse distance units)
  cutoff = global cutoff for Buckingham (and Coulombic if only 1 arg) (distance units)
  cutoff2 = global cutoff for Coulombic (optional) (distance units)
\sphinxstyleemphasis{buck/coul/long/cs} args = cutoff (cutoff2)
  cutoff = global cutoff for Buckingham (and Coulombic if only 1 arg) (distance units)
  cutoff2 = global cutoff for Coulombic (optional) (distance units)
\sphinxstyleemphasis{coul/long} args = cutoff
  cutoff = global cutoff for Coulombic (distance units)
\sphinxstyleemphasis{coul/wolf} args = alpha cutoff
  alpha = damping parameter (inverse distance units)
  cutoff = global cutoff for Coulombic (distance units)
\sphinxstyleemphasis{lj/cut/coul/long/cs} args = cutoff (cutoff2)
  cutoff = global cutoff for LJ (and Coulombic if only 1 arg) (distance units)
  cutoff2 = global cutoff for Coulombic (optional) (distance units)
\sphinxstyleemphasis{lj/class2/coul/long/cs} args = cutoff (cutoff2)
  cutoff = global cutoff for LJ (and Coulombic if only 1 arg) (distance units)
  cutoff2 = global cutoff for Coulombic (optional) (distance units)
\end{sphinxalltt}


\subsection{Examples}
\label{\detokenize{pair_cs:examples}}
\begin{sphinxVerbatim}[commandchars=\\\{\}]
\PYG{k}{pair\PYGZus{}style}\PYG{+w}{ }\PYG{n}{born}\PYG{o}{/}\PYG{n}{coul}\PYG{o}{/}\PYG{n}{dsf}\PYG{o}{/}\PYG{n}{cs}\PYG{+w}{ }\PYG{l+m}{0.1}\PYG{+w}{ }\PYG{l+m}{10.0}\PYG{+w}{ }\PYG{l+m}{12.0}
\PYG{k}{pair\PYGZus{}coeff}\PYG{+w}{ }\PYG{o}{*}\PYG{+w}{ }\PYG{o}{*}\PYG{+w}{   }\PYG{l+m}{0.0}\PYG{+w}{ }\PYG{l+m}{1.00}\PYG{+w}{ }\PYG{l+m}{0.00}\PYG{+w}{ }\PYG{l+m}{0.00}\PYG{+w}{ }\PYG{l+m}{0.00}
\PYG{k}{pair\PYGZus{}coeff}\PYG{+w}{ }\PYG{l+m}{1}\PYG{+w}{ }\PYG{l+m}{1}\PYG{+w}{ }\PYG{l+m}{480.0}\PYG{+w}{ }\PYG{l+m}{0.25}\PYG{+w}{ }\PYG{l+m}{0.00}\PYG{+w}{ }\PYG{l+m}{1.05}\PYG{+w}{ }\PYG{l+m}{0.50}

\PYG{k}{pair\PYGZus{}style}\PYG{+w}{ }\PYG{n}{born}\PYG{o}{/}\PYG{n}{coul}\PYG{o}{/}\PYG{n}{long}\PYG{o}{/}\PYG{n}{cs}\PYG{+w}{ }\PYG{l+m}{10.0}\PYG{+w}{ }\PYG{l+m}{8.0}
\PYG{k}{pair\PYGZus{}coeff}\PYG{+w}{ }\PYG{l+m}{1}\PYG{+w}{ }\PYG{l+m}{1}\PYG{+w}{ }\PYG{l+m}{6.08}\PYG{+w}{ }\PYG{l+m}{0.317}\PYG{+w}{ }\PYG{l+m}{2.340}\PYG{+w}{ }\PYG{l+m}{24.18}\PYG{+w}{ }\PYG{l+m}{11.51}

\PYG{k}{pair\PYGZus{}style}\PYG{+w}{ }\PYG{n}{born}\PYG{o}{/}\PYG{n}{coul}\PYG{o}{/}\PYG{n}{wolf}\PYG{o}{/}\PYG{n}{cs}\PYG{+w}{ }\PYG{l+m}{0.25}\PYG{+w}{ }\PYG{l+m}{10.0}\PYG{+w}{ }\PYG{l+m}{12.0}
\PYG{k}{pair\PYGZus{}coeff}\PYG{+w}{ }\PYG{o}{*}\PYG{+w}{ }\PYG{o}{*}\PYG{+w}{   }\PYG{l+m}{0.0}\PYG{+w}{ }\PYG{l+m}{1.00}\PYG{+w}{ }\PYG{l+m}{0.00}\PYG{+w}{ }\PYG{l+m}{0.00}\PYG{+w}{ }\PYG{l+m}{0.00}
\PYG{k}{pair\PYGZus{}coeff}\PYG{+w}{ }\PYG{l+m}{1}\PYG{+w}{ }\PYG{l+m}{1}\PYG{+w}{ }\PYG{l+m}{480.0}\PYG{+w}{ }\PYG{l+m}{0.25}\PYG{+w}{ }\PYG{l+m}{0.00}\PYG{+w}{ }\PYG{l+m}{1.05}\PYG{+w}{ }\PYG{l+m}{0.50}

\PYG{k}{pair\PYGZus{}style}\PYG{+w}{ }\PYG{n}{buck}\PYG{o}{/}\PYG{n}{coul}\PYG{o}{/}\PYG{n}{long}\PYG{o}{/}\PYG{n}{cs}\PYG{+w}{ }\PYG{l+m}{10.0}
\PYG{k}{pair\PYGZus{}style}\PYG{+w}{ }\PYG{n}{buck}\PYG{o}{/}\PYG{n}{coul}\PYG{o}{/}\PYG{n}{long}\PYG{o}{/}\PYG{n}{cs}\PYG{+w}{ }\PYG{l+m}{10.0}\PYG{+w}{ }\PYG{l+m}{8.0}
\PYG{k}{pair\PYGZus{}coeff}\PYG{+w}{ }\PYG{o}{*}\PYG{+w}{ }\PYG{o}{*}\PYG{+w}{ }\PYG{l+m}{100.0}\PYG{+w}{ }\PYG{l+m}{1.5}\PYG{+w}{ }\PYG{l+m}{200.0}
\PYG{k}{pair\PYGZus{}coeff}\PYG{+w}{ }\PYG{l+m}{1}\PYG{+w}{ }\PYG{l+m}{1}\PYG{+w}{ }\PYG{l+m}{100.0}\PYG{+w}{ }\PYG{l+m}{1.5}\PYG{+w}{ }\PYG{l+m}{200.0}\PYG{+w}{ }\PYG{l+m}{9.0}

\PYG{k}{pair\PYGZus{}style}\PYG{+w}{ }\PYG{n}{coul}\PYG{o}{/}\PYG{n}{long}\PYG{o}{/}\PYG{n}{cs}\PYG{+w}{ }\PYG{l+m}{10.0}
\PYG{k}{pair\PYGZus{}coeff}\PYG{+w}{ }\PYG{o}{*}\PYG{+w}{ }\PYG{o}{*}

\PYG{k}{pair\PYGZus{}style}\PYG{+w}{ }\PYG{n}{coul}\PYG{o}{/}\PYG{n}{wolf}\PYG{o}{/}\PYG{n}{cs}\PYG{+w}{ }\PYG{l+m}{0.2}\PYG{+w}{ }\PYG{l+m}{9.0}
\PYG{k}{pair\PYGZus{}coeff}\PYG{+w}{ }\PYG{o}{*}\PYG{+w}{ }\PYG{o}{*}

\PYG{k}{pair\PYGZus{}style}\PYG{+w}{ }\PYG{n}{lj}\PYG{o}{/}\PYG{n}{cut}\PYG{o}{/}\PYG{n}{coul}\PYG{o}{/}\PYG{n}{long}\PYG{o}{/}\PYG{n}{cs}\PYG{+w}{ }\PYG{l+m}{10.0}
\PYG{k}{pair\PYGZus{}style}\PYG{+w}{ }\PYG{n}{lj}\PYG{o}{/}\PYG{n}{cut}\PYG{o}{/}\PYG{n}{coul}\PYG{o}{/}\PYG{n}{long}\PYG{o}{/}\PYG{n}{cs}\PYG{+w}{ }\PYG{l+m}{10.0}\PYG{+w}{ }\PYG{l+m}{8.0}
\PYG{k}{pair\PYGZus{}coeff}\PYG{+w}{ }\PYG{o}{*}\PYG{+w}{ }\PYG{o}{*}\PYG{+w}{ }\PYG{l+m}{100.0}\PYG{+w}{ }\PYG{l+m}{3.0}
\PYG{k}{pair\PYGZus{}coeff}\PYG{+w}{ }\PYG{l+m}{1}\PYG{+w}{ }\PYG{l+m}{1}\PYG{+w}{ }\PYG{l+m}{100.0}\PYG{+w}{ }\PYG{l+m}{3.5}\PYG{+w}{ }\PYG{l+m}{9.0}
\end{sphinxVerbatim}


\subsection{Description}
\label{\detokenize{pair_cs:description}}
\sphinxAtStartPar
These pair styles are designed to be used with the adiabatic
core/shell model of {\hyperref[\detokenize{pair_cs:mitchellfincham3}]{\sphinxcrossref{\DUrole{std,std-ref}{(Mitchell and Fincham)}}}}.  See
the {\hyperref[\detokenize{Howto_coreshell::doc}]{\sphinxcrossref{\DUrole{doc}{Howto coreshell}}}} page for an overview of
the model as implemented in LAMMPS.

\sphinxAtStartPar
All the styles are identical to the corresponding pair style without
the “/cs” in the name:
\begin{itemize}
\item {} 
\sphinxAtStartPar
{\hyperref[\detokenize{pair_born::doc}]{\sphinxcrossref{\DUrole{doc}{pair\_style born/coul/dsf}}}}

\item {} 
\sphinxAtStartPar
{\hyperref[\detokenize{pair_born::doc}]{\sphinxcrossref{\DUrole{doc}{pair\_style born/coul/long}}}}

\item {} 
\sphinxAtStartPar
{\hyperref[\detokenize{pair_born::doc}]{\sphinxcrossref{\DUrole{doc}{pair\_style born/coul/wolf}}}}

\item {} 
\sphinxAtStartPar
{\hyperref[\detokenize{pair_buck::doc}]{\sphinxcrossref{\DUrole{doc}{pair\_style buck/coul/long}}}}

\item {} 
\sphinxAtStartPar
{\hyperref[\detokenize{pair_coul::doc}]{\sphinxcrossref{\DUrole{doc}{pair\_style coul/long}}}}

\item {} 
\sphinxAtStartPar
{\hyperref[\detokenize{pair_coul::doc}]{\sphinxcrossref{\DUrole{doc}{pair\_style coul/wolf}}}}

\item {} 
\sphinxAtStartPar
{\hyperref[\detokenize{pair_lj_cut_coul::doc}]{\sphinxcrossref{\DUrole{doc}{pair\_style lj/cut/coul/long}}}}

\item {} 
\sphinxAtStartPar
{\hyperref[\detokenize{pair_class2::doc}]{\sphinxcrossref{\DUrole{doc}{pair\_style lj/class2/coul/long}}}}

\end{itemize}

\sphinxAtStartPar
except that they correctly treat the special case where the distance
between two charged core and shell atoms in the same core/shell pair
approach r = 0.0.

\sphinxAtStartPar
Styles with a “/long” in the name are used with a long\sphinxhyphen{}range solver
for Coulombic interactions via the {\hyperref[\detokenize{kspace_style::doc}]{\sphinxcrossref{\DUrole{doc}{kspace\_style}}}}
command.  They require special treatment of the short\sphinxhyphen{}range Coulombic
interactions within the cor/shell model.

\sphinxAtStartPar
Specifically, the short\sphinxhyphen{}range Coulomb interaction between a core and
its shell should be turned off using the
{\hyperref[\detokenize{special_bonds::doc}]{\sphinxcrossref{\DUrole{doc}{special\_bonds}}}} command by setting the 1\sphinxhyphen{}2 weight
to 0.0, which works because the core and shell atoms are bonded to
each other.  This induces a long\sphinxhyphen{}range correction approximation which
fails at small distances (\textasciitilde{}\textless{} 10e\sphinxhyphen{}8). Therefore, the Coulomb term which
is used to calculate the correction factor is extended by a minimal
distance (r\_min = 1.0\sphinxhyphen{}6) when the interaction between a core/shell
pair is treated, as follows
\begin{equation*}
\begin{split}E = \frac{C q_i q_j}{\epsilon (r + r_{min})} \qquad r \rightarrow 0\end{split}
\end{equation*}
\sphinxAtStartPar
where C is an energy\sphinxhyphen{}conversion constant, \(q_i\) and \(q_j\)
are the charges on the core and shell, epsilon is the dielectric
constant and \(r_{min}\) is the minimal distance.

\sphinxAtStartPar
For styles that are not used with a long\sphinxhyphen{}range solver, i.e. those with
“/dsf” or “/wolf” in the name, the only correction is the addition of
a minimal distance to avoid the possible r = 0.0 case for a core/shell
pair.


\bigskip\hrule\bigskip


\sphinxAtStartPar
Styles with a \sphinxstyleemphasis{gpu}, \sphinxstyleemphasis{intel}, \sphinxstyleemphasis{kk}, \sphinxstyleemphasis{omp}, or \sphinxstyleemphasis{opt} suffix are
functionally the same as the corresponding style without the suffix.
They have been optimized to run faster, depending on your available
hardware, as discussed on the {\hyperref[\detokenize{Speed_packages::doc}]{\sphinxcrossref{\DUrole{doc}{Accelerator packages}}}}
page.  The accelerated styles take the same arguments and should
produce the same results, except for round\sphinxhyphen{}off and precision issues.

\sphinxAtStartPar
These accelerated styles are part of the GPU, INTEL, KOKKOS,
OPENMP, and OPT packages, respectively.  They are only enabled if
LAMMPS was built with those packages.  See the {\hyperref[\detokenize{Build_package::doc}]{\sphinxcrossref{\DUrole{doc}{Build package}}}} page for more info.

\sphinxAtStartPar
You can specify the accelerated styles explicitly in your input script
by including their suffix, or you can use the {\hyperref[\detokenize{Run_options::doc}]{\sphinxcrossref{\DUrole{doc}{\sphinxhyphen{}suffix command\sphinxhyphen{}line switch}}}} when you invoke LAMMPS, or you can use the
{\hyperref[\detokenize{suffix::doc}]{\sphinxcrossref{\DUrole{doc}{suffix}}}} command in your input script.

\sphinxAtStartPar
See the {\hyperref[\detokenize{Speed_packages::doc}]{\sphinxcrossref{\DUrole{doc}{Accelerator packages}}}} page for more
instructions on how to use the accelerated styles effectively.


\bigskip\hrule\bigskip



\subsection{Mixing, shift, table, tail correction, restart, rRESPA info}
\label{\detokenize{pair_cs:mixing-shift-table-tail-correction-restart-rrespa-info}}
\sphinxAtStartPar
See the corresponding doc pages for pair styles without the “cs”
suffix to see how mixing, shifting, tabulation, tail correction,
restarting, and rRESPA are handled by theses pair styles.


\bigskip\hrule\bigskip



\subsection{Restrictions}
\label{\detokenize{pair_cs:restrictions}}
\sphinxAtStartPar
These pair styles are part of the CORESHELL package.  They are only
enabled if LAMMPS was built with that package.  See the {\hyperref[\detokenize{Build_package::doc}]{\sphinxcrossref{\DUrole{doc}{Build package}}}} page for more info.


\subsection{Related commands}
\label{\detokenize{pair_cs:related-commands}}
\sphinxAtStartPar
{\hyperref[\detokenize{pair_coeff::doc}]{\sphinxcrossref{\DUrole{doc}{pair\_coeff}}}}, {\hyperref[\detokenize{pair_born::doc}]{\sphinxcrossref{\DUrole{doc}{pair\_style born}}}},
{\hyperref[\detokenize{pair_buck::doc}]{\sphinxcrossref{\DUrole{doc}{pair\_style buck}}}}


\subsection{Default}
\label{\detokenize{pair_cs:default}}
\sphinxAtStartPar
none


\bigskip\hrule\bigskip

\phantomsection\label{\detokenize{pair_cs:mitchellfincham3}}
\sphinxAtStartPar
\sphinxstylestrong{(Mitchell and Fincham)} Mitchell, Fincham, J Phys Condensed Matter,
5, 1031\sphinxhyphen{}1038 (1993).

\sphinxstepscope

\index{pair\_style coul/cut/dielectric@\spxentry{pair\_style coul/cut/dielectric}}\ignorespaces 
\index{pair\_style coul/long/dielectric@\spxentry{pair\_style coul/long/dielectric}}\ignorespaces 
\index{pair\_style lj/cut/coul/cut/dielectric@\spxentry{pair\_style lj/cut/coul/cut/dielectric}}\ignorespaces 
\index{pair\_style lj/cut/coul/cut/dielectric/omp@\spxentry{pair\_style lj/cut/coul/cut/dielectric/omp}}\ignorespaces 
\index{pair\_style lj/cut/coul/debye/dielectric@\spxentry{pair\_style lj/cut/coul/debye/dielectric}}\ignorespaces 
\index{pair\_style lj/cut/coul/debye/dielectric/omp@\spxentry{pair\_style lj/cut/coul/debye/dielectric/omp}}\ignorespaces 
\index{pair\_style lj/cut/coul/long/dielectric@\spxentry{pair\_style lj/cut/coul/long/dielectric}}\ignorespaces 
\index{pair\_style lj/cut/coul/long/dielectric/omp@\spxentry{pair\_style lj/cut/coul/long/dielectric/omp}}\ignorespaces 
\index{pair\_style lj/cut/coul/msm/dielectric@\spxentry{pair\_style lj/cut/coul/msm/dielectric}}\ignorespaces 
\index{pair\_style lj/long/coul/long/dielectric@\spxentry{pair\_style lj/long/coul/long/dielectric}}\ignorespaces 

\section{pair\_style coul/cut/dielectric command}
\label{\detokenize{pair_dielectric:pair-style-coul-cut-dielectric-command}}\label{\detokenize{pair_dielectric:index-9}}\label{\detokenize{pair_dielectric:index-8}}\label{\detokenize{pair_dielectric:index-7}}\label{\detokenize{pair_dielectric:index-6}}\label{\detokenize{pair_dielectric:index-5}}\label{\detokenize{pair_dielectric:index-4}}\label{\detokenize{pair_dielectric:index-3}}\label{\detokenize{pair_dielectric:index-2}}\label{\detokenize{pair_dielectric:index-1}}\label{\detokenize{pair_dielectric:index-0}}\label{\detokenize{pair_dielectric::doc}}

\section{pair\_style coul/long/dielectric command}
\label{\detokenize{pair_dielectric:pair-style-coul-long-dielectric-command}}

\section{pair\_style lj/cut/coul/cut/dielectric command}
\label{\detokenize{pair_dielectric:pair-style-lj-cut-coul-cut-dielectric-command}}
\sphinxAtStartPar
Accelerator Variants: \sphinxstyleemphasis{lj/cut/coul/cut/dielectric/omp}


\section{pair\_style lj/cut/coul/debye/dielectric command}
\label{\detokenize{pair_dielectric:pair-style-lj-cut-coul-debye-dielectric-command}}
\sphinxAtStartPar
Accelerator Variants: \sphinxstyleemphasis{lj/cut/coul/debye/dielectric/omp}


\section{pair\_style lj/cut/coul/long/dielectric command}
\label{\detokenize{pair_dielectric:pair-style-lj-cut-coul-long-dielectric-command}}
\sphinxAtStartPar
Accelerator Variants: \sphinxstyleemphasis{lj/cut/coul/long/dielectric/omp}


\section{pair\_style lj/cut/coul/msm/dielectric command}
\label{\detokenize{pair_dielectric:pair-style-lj-cut-coul-msm-dielectric-command}}

\section{pair\_style lj/long/coul/long/dielectric command}
\label{\detokenize{pair_dielectric:pair-style-lj-long-coul-long-dielectric-command}}

\subsection{Syntax}
\label{\detokenize{pair_dielectric:syntax}}
\begin{sphinxVerbatim}[commandchars=\\\{\}]
\PYG{k}{pair\PYGZus{}style}\PYG{+w}{ }\PYG{n}{style}\PYG{+w}{ }\PYG{n}{args}
\end{sphinxVerbatim}
\begin{itemize}
\item {} 
\sphinxAtStartPar
style = \sphinxstyleemphasis{lj/cut/coul/cut/dielectric} or \sphinxstyleemphasis{lj/cut/coul/long/dielectric} or \sphinxstyleemphasis{lj/cut/coul/msm/dielectric} or \sphinxstyleemphasis{lj/long/coul/msm/dielectric}

\item {} 
\sphinxAtStartPar
args = list of arguments for a particular style

\end{itemize}


\subsection{Examples}
\label{\detokenize{pair_dielectric:examples}}
\begin{sphinxVerbatim}[commandchars=\\\{\}]
\PYG{k}{pair\PYGZus{}style}\PYG{+w}{ }\PYG{n}{coul}\PYG{o}{/}\PYG{n}{cut}\PYG{o}{/}\PYG{n}{dielectric}\PYG{+w}{ }\PYG{l+m}{10.0}
\PYG{k}{pair\PYGZus{}coeff}\PYG{+w}{ }\PYG{o}{*}\PYG{+w}{ }\PYG{o}{*}
\PYG{k}{pair\PYGZus{}coeff}\PYG{+w}{ }\PYG{l+m}{1}\PYG{+w}{ }\PYG{l+m}{1}\PYG{+w}{ }\PYG{l+m}{9.0}

\PYG{k}{pair\PYGZus{}style}\PYG{+w}{ }\PYG{n}{lj}\PYG{o}{/}\PYG{n}{cut}\PYG{o}{/}\PYG{n}{coul}\PYG{o}{/}\PYG{n}{cut}\PYG{o}{/}\PYG{n}{dielectric}\PYG{+w}{ }\PYG{l+m}{10.0}
\PYG{k}{pair\PYGZus{}style}\PYG{+w}{ }\PYG{n}{lj}\PYG{o}{/}\PYG{n}{cut}\PYG{o}{/}\PYG{n}{coul}\PYG{o}{/}\PYG{n}{cut}\PYG{o}{/}\PYG{n}{dielectric}\PYG{+w}{ }\PYG{l+m}{10.0}\PYG{+w}{ }\PYG{l+m}{8.0}
\PYG{k}{pair\PYGZus{}coeff}\PYG{+w}{ }\PYG{o}{*}\PYG{+w}{ }\PYG{o}{*}\PYG{+w}{ }\PYG{l+m}{100.0}\PYG{+w}{ }\PYG{l+m}{3.0}
\PYG{k}{pair\PYGZus{}coeff}\PYG{+w}{ }\PYG{l+m}{1}\PYG{+w}{ }\PYG{l+m}{1}\PYG{+w}{ }\PYG{l+m}{100.0}\PYG{+w}{ }\PYG{l+m}{3.5}\PYG{+w}{ }\PYG{l+m}{9.0}

\PYG{k}{pair\PYGZus{}style}\PYG{+w}{ }\PYG{n}{lj}\PYG{o}{/}\PYG{n}{cut}\PYG{o}{/}\PYG{n}{coul}\PYG{o}{/}\PYG{n}{long}\PYG{o}{/}\PYG{n}{dielectric}\PYG{+w}{ }\PYG{l+m}{10.0}
\PYG{k}{pair\PYGZus{}style}\PYG{+w}{ }\PYG{n}{lj}\PYG{o}{/}\PYG{n}{cut}\PYG{o}{/}\PYG{n}{coul}\PYG{o}{/}\PYG{n}{long}\PYG{o}{/}\PYG{n}{dielectric}\PYG{+w}{ }\PYG{l+m}{10.0}\PYG{+w}{ }\PYG{l+m}{8.0}
\PYG{k}{pair\PYGZus{}coeff}\PYG{+w}{ }\PYG{o}{*}\PYG{+w}{ }\PYG{o}{*}\PYG{+w}{ }\PYG{l+m}{100.0}\PYG{+w}{ }\PYG{l+m}{3.0}
\PYG{k}{pair\PYGZus{}coeff}\PYG{+w}{ }\PYG{l+m}{1}\PYG{+w}{ }\PYG{l+m}{1}\PYG{+w}{ }\PYG{l+m}{100.0}\PYG{+w}{ }\PYG{l+m}{3.5}\PYG{+w}{ }\PYG{l+m}{9.0}
\end{sphinxVerbatim}

\sphinxAtStartPar
Used in input scripts:
\begin{quote}

\begin{sphinxVerbatim}[commandchars=\\\{\}]
examples/PACKAGES/dielectric/in.confined
examples/PACKAGES/dielectric/in.nopbc
\end{sphinxVerbatim}
\end{quote}


\subsection{Description}
\label{\detokenize{pair_dielectric:description}}
\sphinxAtStartPar
All these pair styles are derived from the corresponding pair styles
without the \sphinxstyleemphasis{dielectric} suffix. In addition to computing atom forces
and energies, these pair styles compute the electric field vector at
each atom, which are intended to be used by the {\hyperref[\detokenize{fix_polarize::doc}]{\sphinxcrossref{\DUrole{doc}{fix polarize}}}} commands to compute induced charges at interfaces
between two regions of different dielectric constant.

\sphinxAtStartPar
These pair styles should be used with {\hyperref[\detokenize{atom_style::doc}]{\sphinxcrossref{\DUrole{doc}{atom\_style dielectric}}}}.

\sphinxAtStartPar
The styles lj/cut/coul/long/dielectric, lj/cut/coul/msm/dielectric, and
lj/long/coul/long/dielectric should be used with their kspace style
counterparts, namely, pppm/dielectric, pppm/disp/dielectric, and
msm/dielectric, respectively.


\bigskip\hrule\bigskip


\sphinxAtStartPar
Styles with a \sphinxstyleemphasis{gpu}, \sphinxstyleemphasis{intel}, \sphinxstyleemphasis{kk}, \sphinxstyleemphasis{omp}, or \sphinxstyleemphasis{opt} suffix are
functionally the same as the corresponding style without the suffix.
They have been optimized to run faster, depending on your available
hardware, as discussed on the {\hyperref[\detokenize{Speed_packages::doc}]{\sphinxcrossref{\DUrole{doc}{Accelerator packages}}}}
page.  The accelerated styles take the same arguments and should
produce the same results, except for round\sphinxhyphen{}off and precision issues.

\sphinxAtStartPar
These accelerated styles are part of the GPU, INTEL, KOKKOS,
OPENMP, and OPT packages, respectively.  They are only enabled if
LAMMPS was built with those packages.  See the {\hyperref[\detokenize{Build_package::doc}]{\sphinxcrossref{\DUrole{doc}{Build package}}}} page for more info.

\sphinxAtStartPar
You can specify the accelerated styles explicitly in your input script
by including their suffix, or you can use the {\hyperref[\detokenize{Run_options::doc}]{\sphinxcrossref{\DUrole{doc}{\sphinxhyphen{}suffix command\sphinxhyphen{}line switch}}}} when you invoke LAMMPS, or you can use the
{\hyperref[\detokenize{suffix::doc}]{\sphinxcrossref{\DUrole{doc}{suffix}}}} command in your input script.

\sphinxAtStartPar
See the {\hyperref[\detokenize{Speed_packages::doc}]{\sphinxcrossref{\DUrole{doc}{Accelerator packages}}}} page for more
instructions on how to use the accelerated styles effectively.


\bigskip\hrule\bigskip



\subsection{Mixing, shift, table, tail correction, restart, rRESPA info}
\label{\detokenize{pair_dielectric:mixing-shift-table-tail-correction-restart-rrespa-info}}
\sphinxAtStartPar
For atom type pairs I,J and I != J, the epsilon and sigma coefficients
and cutoff distances for this pair style can be mixed.  The default mix
algorithm is \sphinxstyleemphasis{geometric}.  See the {\hyperref[\detokenize{pair_modify::doc}]{\sphinxcrossref{\DUrole{doc}{pair\_modify}}}}”
command for details.

\sphinxAtStartPar
The {\hyperref[\detokenize{pair_modify::doc}]{\sphinxcrossref{\DUrole{doc}{pair\_modify}}}} table option is not relevant
for this pair style.

\sphinxAtStartPar
These pair styles write its information to {\hyperref[\detokenize{restart::doc}]{\sphinxcrossref{\DUrole{doc}{binary restart files}}}}, so pair\_style and pair\_coeff commands do not need to be
specified in an input script that reads a restart file.

\sphinxAtStartPar
These pair styles can only be used via the \sphinxstyleemphasis{pair} keyword of the
{\hyperref[\detokenize{run_style::doc}]{\sphinxcrossref{\DUrole{doc}{run\_style respa}}}} command.  It does not support the
\sphinxstyleemphasis{inner}, \sphinxstyleemphasis{middle}, \sphinxstyleemphasis{outer} keywords.


\subsection{Restrictions}
\label{\detokenize{pair_dielectric:restrictions}}
\sphinxAtStartPar
These styles are part of the DIELECTRIC package.  They are only enabled
if LAMMPS was built with that package.  See the {\hyperref[\detokenize{Build_package::doc}]{\sphinxcrossref{\DUrole{doc}{Build package}}}} page for more info.


\subsection{Related commands}
\label{\detokenize{pair_dielectric:related-commands}}
\sphinxAtStartPar
{\hyperref[\detokenize{pair_coeff::doc}]{\sphinxcrossref{\DUrole{doc}{pair\_coeff}}}}, {\hyperref[\detokenize{fix_polarize::doc}]{\sphinxcrossref{\DUrole{doc}{fix polarize}}}}, {\hyperref[\detokenize{read_data::doc}]{\sphinxcrossref{\DUrole{doc}{read\_data}}}}


\subsection{Default}
\label{\detokenize{pair_dielectric:default}}
\sphinxAtStartPar
none

\sphinxstepscope

\index{pair\_style lj/cut/dipole/cut@\spxentry{pair\_style lj/cut/dipole/cut}}\ignorespaces 
\index{pair\_style lj/cut/dipole/cut/gpu@\spxentry{pair\_style lj/cut/dipole/cut/gpu}}\ignorespaces 
\index{pair\_style lj/cut/dipole/cut/kk@\spxentry{pair\_style lj/cut/dipole/cut/kk}}\ignorespaces 
\index{pair\_style lj/cut/dipole/cut/omp@\spxentry{pair\_style lj/cut/dipole/cut/omp}}\ignorespaces 
\index{pair\_style lj/sf/dipole/sf@\spxentry{pair\_style lj/sf/dipole/sf}}\ignorespaces 
\index{pair\_style lj/sf/dipole/sf/gpu@\spxentry{pair\_style lj/sf/dipole/sf/gpu}}\ignorespaces 
\index{pair\_style lj/sf/dipole/sf/omp@\spxentry{pair\_style lj/sf/dipole/sf/omp}}\ignorespaces 
\index{pair\_style lj/cut/dipole/long@\spxentry{pair\_style lj/cut/dipole/long}}\ignorespaces 
\index{pair\_style lj/cut/dipole/long/gpu@\spxentry{pair\_style lj/cut/dipole/long/gpu}}\ignorespaces 
\index{pair\_style lj/long/dipole/long@\spxentry{pair\_style lj/long/dipole/long}}\ignorespaces 

\section{pair\_style lj/cut/dipole/cut command}
\label{\detokenize{pair_dipole:pair-style-lj-cut-dipole-cut-command}}\label{\detokenize{pair_dipole:index-9}}\label{\detokenize{pair_dipole:index-8}}\label{\detokenize{pair_dipole:index-7}}\label{\detokenize{pair_dipole:index-6}}\label{\detokenize{pair_dipole:index-5}}\label{\detokenize{pair_dipole:index-4}}\label{\detokenize{pair_dipole:index-3}}\label{\detokenize{pair_dipole:index-2}}\label{\detokenize{pair_dipole:index-1}}\label{\detokenize{pair_dipole:index-0}}\label{\detokenize{pair_dipole::doc}}
\sphinxAtStartPar
Accelerator Variants: \sphinxstyleemphasis{lj/cut/dipole/cut/gpu}, \sphinxstyleemphasis{lj/cut/dipole/cut/kk}, \sphinxstyleemphasis{lj/cut/dipole/cut/omp}


\section{pair\_style lj/sf/dipole/sf command}
\label{\detokenize{pair_dipole:pair-style-lj-sf-dipole-sf-command}}
\sphinxAtStartPar
Accelerator Variants: \sphinxstyleemphasis{lj/sf/dipole/sf/gpu}, \sphinxstyleemphasis{lj/sf/dipole/sf/omp}


\section{pair\_style lj/cut/dipole/long command}
\label{\detokenize{pair_dipole:pair-style-lj-cut-dipole-long-command}}
\sphinxAtStartPar
Accelerator Variants: \sphinxstyleemphasis{lj/cut/dipole/long/gpu}


\section{pair\_style lj/long/dipole/long command}
\label{\detokenize{pair_dipole:pair-style-lj-long-dipole-long-command}}

\subsection{Syntax}
\label{\detokenize{pair_dipole:syntax}}
\begin{sphinxVerbatim}[commandchars=\\\{\}]
\PYG{k}{pair\PYGZus{}style}\PYG{+w}{ }\PYG{n}{lj}\PYG{o}{/}\PYG{n}{cut}\PYG{o}{/}\PYG{n}{dipole}\PYG{o}{/}\PYG{n}{cut}\PYG{+w}{ }\PYG{n}{cutoff}\PYG{+w}{ }\PYG{n+nv}{(}\PYG{n+nv}{cutoff2}\PYG{n+nv}{)}
\PYG{k}{pair\PYGZus{}style}\PYG{+w}{ }\PYG{n}{lj}\PYG{o}{/}\PYG{n}{sf}\PYG{o}{/}\PYG{n}{dipole}\PYG{o}{/}\PYG{n}{sf}\PYG{+w}{ }\PYG{n}{cutoff}\PYG{+w}{ }\PYG{n+nv}{(}\PYG{n+nv}{cutoff2}\PYG{n+nv}{)}
\PYG{k}{pair\PYGZus{}style}\PYG{+w}{ }\PYG{n}{lj}\PYG{o}{/}\PYG{n}{cut}\PYG{o}{/}\PYG{n}{dipole}\PYG{o}{/}\PYG{n}{long}\PYG{+w}{ }\PYG{n}{cutoff}\PYG{+w}{ }\PYG{n+nv}{(}\PYG{n+nv}{cutoff2}\PYG{n+nv}{)}
\PYG{k}{pair\PYGZus{}style}\PYG{+w}{ }\PYG{n}{lj}\PYG{o}{/}\PYG{n}{long}\PYG{o}{/}\PYG{n}{dipole}\PYG{o}{/}\PYG{n}{long}\PYG{+w}{ }\PYG{n}{flag\PYGZus{}lj}\PYG{+w}{ }\PYG{n}{flag\PYGZus{}coul}\PYG{+w}{ }\PYG{n}{cutoff}\PYG{+w}{ }\PYG{n+nv}{(}\PYG{n+nv}{cutoff2}\PYG{n+nv}{)}
\end{sphinxVerbatim}
\begin{itemize}
\item {} 
\sphinxAtStartPar
cutoff = global cutoff LJ (and Coulombic if only 1 arg) (distance units)

\item {} 
\sphinxAtStartPar
cutoff2 = global cutoff for Coulombic and dipole (optional) (distance units)

\item {} 
\sphinxAtStartPar
flag\_lj = \sphinxstyleemphasis{long} or \sphinxstyleemphasis{cut} or \sphinxstyleemphasis{off}
\begin{sphinxalltt}
\sphinxstyleemphasis{long} = use long\sphinxhyphen{}range damping on dispersion 1/r\textasciicircum{}6 term
\sphinxstyleemphasis{cut} = use a cutoff on dispersion 1/r\textasciicircum{}6 term
\sphinxstyleemphasis{off} = omit disperion 1/r\textasciicircum{}6 term entirely
\end{sphinxalltt}

\item {} 
\sphinxAtStartPar
flag\_coul = \sphinxstyleemphasis{long} or \sphinxstyleemphasis{off}
\begin{sphinxalltt}
\sphinxstyleemphasis{long} = use long\sphinxhyphen{}range damping on Coulombic 1/r and point\sphinxhyphen{}dipole terms
\sphinxstyleemphasis{off} = omit Coulombic and point\sphinxhyphen{}dipole terms entirely
\end{sphinxalltt}

\end{itemize}


\subsection{Examples}
\label{\detokenize{pair_dipole:examples}}
\begin{sphinxVerbatim}[commandchars=\\\{\}]
\PYG{k}{pair\PYGZus{}style}\PYG{+w}{ }\PYG{n}{lj}\PYG{o}{/}\PYG{n}{cut}\PYG{o}{/}\PYG{n}{dipole}\PYG{o}{/}\PYG{n}{cut}\PYG{+w}{ }\PYG{l+m}{2.5}\PYG{+w}{ }\PYG{l+m}{5.0}
\PYG{k}{pair\PYGZus{}coeff}\PYG{+w}{ }\PYG{o}{*}\PYG{+w}{ }\PYG{o}{*}\PYG{+w}{ }\PYG{l+m}{1.0}\PYG{+w}{ }\PYG{l+m}{1.0}
\PYG{k}{pair\PYGZus{}coeff}\PYG{+w}{ }\PYG{l+m}{2}\PYG{+w}{ }\PYG{l+m}{3}\PYG{+w}{ }\PYG{l+m}{0.8}\PYG{+w}{ }\PYG{l+m}{1.0}\PYG{+w}{ }\PYG{l+m}{2.5}\PYG{+w}{ }\PYG{l+m}{4.0}

\PYG{k}{pair\PYGZus{}style}\PYG{+w}{ }\PYG{n}{lj}\PYG{o}{/}\PYG{n}{sf}\PYG{o}{/}\PYG{n}{dipole}\PYG{o}{/}\PYG{n}{sf}\PYG{+w}{ }\PYG{l+m}{9.0}
\PYG{k}{pair\PYGZus{}coeff}\PYG{+w}{ }\PYG{o}{*}\PYG{+w}{ }\PYG{o}{*}\PYG{+w}{ }\PYG{l+m}{1.0}\PYG{+w}{ }\PYG{l+m}{1.0}
\PYG{k}{pair\PYGZus{}coeff}\PYG{+w}{ }\PYG{l+m}{2}\PYG{+w}{ }\PYG{l+m}{3}\PYG{+w}{ }\PYG{l+m}{1.0}\PYG{+w}{ }\PYG{l+m}{1.0}\PYG{+w}{ }\PYG{l+m}{2.5}\PYG{+w}{ }\PYG{l+m}{4.0}\PYG{+w}{ }\PYG{n}{scale}\PYG{+w}{ }\PYG{l+m}{0.5}
\PYG{k}{pair\PYGZus{}coeff}\PYG{+w}{ }\PYG{l+m}{2}\PYG{+w}{ }\PYG{l+m}{3}\PYG{+w}{ }\PYG{l+m}{0.8}\PYG{+w}{ }\PYG{l+m}{1.0}\PYG{+w}{ }\PYG{l+m}{2.5}\PYG{+w}{ }\PYG{l+m}{4.0}

\PYG{k}{pair\PYGZus{}style}\PYG{+w}{ }\PYG{n}{lj}\PYG{o}{/}\PYG{n}{cut}\PYG{o}{/}\PYG{n}{dipole}\PYG{o}{/}\PYG{n}{long}\PYG{+w}{ }\PYG{l+m}{2.5}\PYG{+w}{ }\PYG{l+m}{3.5}
\PYG{k}{pair\PYGZus{}coeff}\PYG{+w}{ }\PYG{o}{*}\PYG{+w}{ }\PYG{o}{*}\PYG{+w}{ }\PYG{l+m}{1.0}\PYG{+w}{ }\PYG{l+m}{1.0}
\PYG{k}{pair\PYGZus{}coeff}\PYG{+w}{ }\PYG{l+m}{2}\PYG{+w}{ }\PYG{l+m}{3}\PYG{+w}{ }\PYG{l+m}{0.8}\PYG{+w}{ }\PYG{l+m}{1.0}\PYG{+w}{ }\PYG{l+m}{3.0}

\PYG{k}{pair\PYGZus{}style}\PYG{+w}{ }\PYG{n}{lj}\PYG{o}{/}\PYG{n}{long}\PYG{o}{/}\PYG{n}{dipole}\PYG{o}{/}\PYG{n}{long}\PYG{+w}{ }\PYG{n}{long}\PYG{+w}{ }\PYG{n}{long}\PYG{+w}{ }\PYG{l+m}{3.5}
\PYG{k}{pair\PYGZus{}coeff}\PYG{+w}{ }\PYG{o}{*}\PYG{+w}{ }\PYG{o}{*}\PYG{+w}{ }\PYG{l+m}{1.0}\PYG{+w}{ }\PYG{l+m}{1.0}
\PYG{k}{pair\PYGZus{}coeff}\PYG{+w}{ }\PYG{l+m}{2}\PYG{+w}{ }\PYG{l+m}{3}\PYG{+w}{ }\PYG{l+m}{0.8}\PYG{+w}{ }\PYG{l+m}{1.0}

\PYG{k}{pair\PYGZus{}style}\PYG{+w}{ }\PYG{n}{lj}\PYG{o}{/}\PYG{n}{long}\PYG{o}{/}\PYG{n}{dipole}\PYG{o}{/}\PYG{n}{long}\PYG{+w}{ }\PYG{n}{cut}\PYG{+w}{ }\PYG{n}{long}\PYG{+w}{ }\PYG{l+m}{2.5}\PYG{+w}{ }\PYG{l+m}{3.5}
\PYG{k}{pair\PYGZus{}coeff}\PYG{+w}{ }\PYG{o}{*}\PYG{+w}{ }\PYG{o}{*}\PYG{+w}{ }\PYG{l+m}{1.0}\PYG{+w}{ }\PYG{l+m}{1.0}
\PYG{k}{pair\PYGZus{}coeff}\PYG{+w}{ }\PYG{l+m}{2}\PYG{+w}{ }\PYG{l+m}{3}\PYG{+w}{ }\PYG{l+m}{0.8}\PYG{+w}{ }\PYG{l+m}{1.0}\PYG{+w}{ }\PYG{l+m}{3.0}
\end{sphinxVerbatim}


\subsection{Description}
\label{\detokenize{pair_dipole:description}}
\sphinxAtStartPar
Style \sphinxstyleemphasis{lj/cut/dipole/cut} computes interactions between pairs of
particles that each have a charge and/or a point dipole moment.  In
addition to the usual Lennard\sphinxhyphen{}Jones interaction between the particles
(Elj) the charge\sphinxhyphen{}charge (Eqq), charge\sphinxhyphen{}dipole (Eqp), and dipole\sphinxhyphen{}dipole
(Epp) interactions are computed by these formulas for the energy (E),
force (F), and torque (T) between particles I and J.
\begin{equation*}
\begin{split}E_{LJ}  = & 4 \epsilon \left[ \left(\frac{\sigma}{r}\right)^{12} -
                     \left(\frac{\sigma}{r}\right)^6 \right] \\
E_{qq}  = & \frac{q_i q_j}{r} \\
E_{qp}  = & \frac{q}{r^3} (p \bullet \vec{r}) \\
E_{pp}  = & \frac{1}{r^3} (\vec{p}_i \bullet \vec{p}_j) -
          \frac{3}{r^5} (\vec{p}_i \bullet \vec{r}) (\vec{p}_j \bullet \vec{r}) \\
          & \\
F_{qq}  = & \frac{q_i q_j}{r^3} \vec{r} \\
F_{qp}  = & -\frac{q}{r^3} \vec{p} + \frac{3q}{r^5}
          (\vec{p} \bullet \vec{r}) \vec{r} \\
F_{pp}  = & \frac{3}{r^5} (\vec{p}_i \bullet \vec{p}_j) \vec{r} -
          \frac{15}{r^7} (\vec{p}_i \bullet \vec{r})
          (\vec{p}_j \bullet \vec{r}) \vec{r} +
          \frac{3}{r^5} \left[ (\vec{p}_j \bullet \vec{r}) \vec{p}_i +
          (\vec{p}_i \bullet \vec{r}) \vec{p}_j \right] \\
          & \\
T_{pq} = T_{ij}  = & \frac{q_j}{r^3} (\vec{p}_i \times \vec{r}) \\
T_{qp} = T_{ji}  = & - \frac{q_i}{r^3} (\vec{p}_j \times \vec{r}) \\
T_{pp} = T_{ij}  = & -\frac{1}{r^3} (\vec{p}_i \times \vec{p}_j) +
                   \frac{3}{r^5} (\vec{p}_j \bullet \vec{r})
                   (\vec{p}_i \times \vec{r}) \\
T_{pp} = T_{ji}  = & -\frac{1}{r^3} (\vec{p}_j \times \vec{p}_i) +
                   \frac{3}{r^5} (\vec{p}_i \bullet \vec{r})
                   (\vec{p}_j \times \vec{r})\end{split}
\end{equation*}
\sphinxAtStartPar
where \(q_i\) and \(q_j\) are the charges on the two
particles, \(\vec{p}_i\) and \(\vec{p}_j\) are the dipole
moment vectors of the two particles, r is their separation distance,
and the vector r = Ri \sphinxhyphen{} Rj is the separation vector between the two
particles.  Note that Eqq and Fqq are simply Coulombic energy and
force, Fij = \sphinxhyphen{}Fji as symmetric forces, and Tij != \sphinxhyphen{}Tji since the
torques do not act symmetrically.  These formulas are discussed in
{\hyperref[\detokenize{pair_dipole:allen2}]{\sphinxcrossref{\DUrole{std,std-ref}{(AllenTildesley)}}}} and in {\hyperref[\detokenize{pair_dipole:toukmaji2}]{\sphinxcrossref{\DUrole{std,std-ref}{(Toukmaji)}}}}.

\sphinxAtStartPar
Also note, that in the code, all of these terms (except Elj) have a
\(C/\epsilon\) prefactor, the same as the Coulombic term in the
LJ + Coulombic pair styles discussed {\hyperref[\detokenize{pair_lj::doc}]{\sphinxcrossref{\DUrole{doc}{here}}}}.  C is an
energy\sphinxhyphen{}conversion constant and \(\epsilon\) is the dielectric constant
which can be set by the {\hyperref[\detokenize{dielectric::doc}]{\sphinxcrossref{\DUrole{doc}{dielectric}}}} command.  The
same is true of the equations that follow for other dipole pair
styles.

\sphinxAtStartPar
Style \sphinxstyleemphasis{lj/sf/dipole/sf} computes “shifted\sphinxhyphen{}force” interactions between
pairs of particles that each have a charge and/or a point dipole
moment. In general, a shifted\sphinxhyphen{}force potential is a (slightly) modified
potential containing extra terms that make both the energy and its
derivative go to zero at the cutoff distance; this removes
(cutoff\sphinxhyphen{}related) problems in energy conservation and any numerical
instability in the equations of motion {\hyperref[\detokenize{pair_dipole:allen2}]{\sphinxcrossref{\DUrole{std,std-ref}{(AllenTildesley)}}}}. Shifted\sphinxhyphen{}force interactions for the Lennard\sphinxhyphen{}Jones (E\_LJ),
charge\sphinxhyphen{}charge (Eqq), charge\sphinxhyphen{}dipole (Eqp), dipole\sphinxhyphen{}charge (Epq) and
dipole\sphinxhyphen{}dipole (Epp) potentials are computed by these formulas for the
energy (E), force (F), and torque (T) between particles I and J:
\begin{align*}\!\begin{aligned}
E_{LJ}  = &  4\epsilon \left\{ \left[ \left( \frac{\sigma}{r} \right)^{\!12} -
\left( \frac{\sigma}{r} \right)^{\!6}  \right] +
\left[ 6\left( \frac{\sigma}{r_c} \right)^{\!12} -
3\left(\frac{\sigma}{r_c}\right)^{\!6}\right]\left(\frac{r}{r_c}\right)^{\!2}
- 7\left( \frac{\sigma}{r_c} \right)^{\!12} +
4\left( \frac{\sigma}{r_c} \right)^{\!6}\right\} \\
E_{qq}  = & \frac{q_i q_j}{r}\left(1-\frac{r}{r_c}\right)^{\!2} \\
E_{pq}  = & E_{ji} = -\frac{q}{r^3} \left[ 1 -
3\left(\frac{r}{r_c}\right)^{\!2} +
2\left(\frac{r}{r_c}\right)^{\!3}\right] (\vec{p}\bullet\vec{r}) \\
E_{qp}  = & E_{ij} = \frac{q}{r^3} \left[ 1 -
3\left(\frac{r}{r_c}\right)^{\!2} +
2\left(\frac{r}{r_c}\right)^{\!3}\right] (\vec{p}\bullet\vec{r}) \\
E_{pp} = & \left[1-4\left(\frac{r}{r_c}\right)^{\!3} +
3\left(\frac{r}{r_c}\right)^{\!4}\right]\left[\frac{1}{r^3}
(\vec{p}_i \bullet \vec{p}_j) - \frac{3}{r^5}
(\vec{p}_i \bullet \vec{r}) (\vec{p}_j \bullet \vec{r})\right] \\
         & \\\\
F_{LJ}  = & \left\{\left[48\epsilon \left(\frac{\sigma}{r}\right)^{\!12} -
24\epsilon \left(\frac{\sigma}{r}\right)^{\!6} \right]\frac{1}{r^2} -
\left[48\epsilon \left(\frac{\sigma}{r_c}\right)^{\!12} - 24\epsilon
\left(\frac{\sigma}{r_c}\right)^{\!6} \right]\frac{1}{r_c^2}\right\}\vec{r}\\
F_{qq}  = & \frac{q_i q_j}{r}\left(\frac{1}{r^2} -
\frac{1}{r_c^2}\right)\vec{r} \\
F_{pq} = & F_{ij } =  -\frac{3q}{r^5} \left[ 1 -
\left(\frac{r}{r_c}\right)^{\!2}\right](\vec{p}\bullet\vec{r})\vec{r} +
\frac{q}{r^3}\left[1-3\left(\frac{r}{r_c}\right)^{\!2} +
2\left(\frac{r}{r_c}\right)^{\!3}\right] \vec{p} \\
F_{qp} = & F_{ij}  = \frac{3q}{r^5} \left[ 1 -
\left(\frac{r}{r_c}\right)^{\!2}\right] (\vec{p}\bullet\vec{r})\vec{r} -
\frac{q}{r^3}\left[1-3\left(\frac{r}{r_c}\right)^{\!2} +
2\left(\frac{r}{r_c}\right)^{\!3}\right] \vec{p} \\
F_{pp}  = &\frac{3}{r^5}\Bigg\{\left[1-\left(\frac{r}{r_c}\right)^{\!4}\right]
\left[(\vec{p}_i\bullet\vec{p}_j) - \frac{3}{r^2} (\vec{p}_i\bullet\vec{r})
(\vec{p}_j \bullet \vec{r})\right] \vec{r} + \\
  & \left[1 -
4\left(\frac{r}{r_c}\right)^{\!3}+3\left(\frac{r}{r_c}\right)^{\!4}\right]
\left[ (\vec{p}_j \bullet \vec{r}) \vec{p}_i + (\vec{p}_i \bullet \vec{r})
\vec{p}_j -\frac{2}{r^2} (\vec{p}_i \bullet \vec{r})
(\vec{p}_j \bullet \vec{r})\vec{r}\right] \Bigg\}\\
\end{aligned}\end{align*}\begin{equation*}
\begin{split} T_{pq} = T_{ij}  = & \frac{q_j}{r^3} \left[ 1 -
3\left(\frac{r}{r_c}\right)^{\!2} +
2\left(\frac{r}{r_c}\right)^{\!3}\right] (\vec{p}_i\times\vec{r}) \\
T_{qp} = T_{ji}  = & - \frac{q_i}{r^3} \left[ 1 -
3\left(\frac{r}{r_c}\right)^{\!2} +
2\left(\frac{r}{r_c}\right)^{\!3} \right] (\vec{p}_j\times\vec{r}) \\
T_{pp} = T_{ij}  = & -\frac{1}{r^3}\left[1-4\left(\frac{r}{r_c}\right)^{\!3} +
e3\left(\frac{r}{r_c}\right)^{\!4}\right] (\vec{p}_i \times \vec{p}_j) + \\
                   & \frac{3}{r^5}\left[1-4\left(\frac{r}{r_c}\right)^{\!3} +
3\left(\frac{r}{r_c}\right)^{\!4}\right] (\vec{p}_j\bullet\vec{r})
(\vec{p}_i \times \vec{r}) \\
T_{pp} = T_{ji} = & -\frac{1}{r^3}\left[1-4\left(\frac{r}{r_c}\right)^{\!3} +
3\left(\frac{r}{r_c}\right)^{\!4}\right](\vec{p}_j \times \vec{p}_i) + \\
                   & \frac{3}{r^5}\left[1-4\left(\frac{r}{r_c}\right)^{\!3} +
3\left(\frac{r}{r_c}\right)^{\!4}\right] (\vec{p}_i \bullet \vec{r})
(\vec{p}_j \times \vec{r})\end{split}
\end{equation*}
\sphinxAtStartPar
where \(\epsilon\) and \(\sigma\) are the standard LJ
parameters, \(r_c\) is the cutoff, \(q_i\) and \(q_j\) are
the charges on the two particles, \(\vec{p}_i\) and
\(\vec{p}_j\) are the dipole moment vectors of the two particles,
r is their separation distance, and the vector r = Ri \sphinxhyphen{} Rj is the
separation vector between the two particles.  Note that Eqq and Fqq
are simply Coulombic energy and force, Fij = \sphinxhyphen{}Fji as symmetric forces,
and Tij != \sphinxhyphen{}Tji since the torques do not act symmetrically.  The
shifted\sphinxhyphen{}force formula for the Lennard\sphinxhyphen{}Jones potential is reported in
{\hyperref[\detokenize{pair_dipole:stoddard}]{\sphinxcrossref{\DUrole{std,std-ref}{(Stoddard)}}}}.  The original (non\sphinxhyphen{}shifted) formulas for
the electrostatic potentials, forces and torques can be found in
{\hyperref[\detokenize{pair_dipole:price2}]{\sphinxcrossref{\DUrole{std,std-ref}{(Price)}}}}. The shifted\sphinxhyphen{}force electrostatic potentials
have been obtained by applying equation 5.13 of {\hyperref[\detokenize{pair_dipole:allen2}]{\sphinxcrossref{\DUrole{std,std-ref}{(AllenTildesley)}}}}. The formulas for the corresponding forces and torques have
been obtained by applying the ‘chain rule’ as in appendix C.3 of
{\hyperref[\detokenize{pair_dipole:allen2}]{\sphinxcrossref{\DUrole{std,std-ref}{(AllenTildesley)}}}}.

\sphinxAtStartPar
If one cutoff is specified in the pair\_style command, it is used for
both the LJ and Coulombic (q,p) terms.  If two cutoffs are specified,
they are used as cutoffs for the LJ and Coulombic (q,p) terms
respectively. This pair style also supports an optional \sphinxstyleemphasis{scale}
keyword as part of a pair\_coeff statement, where the interactions can
be scaled according to this factor. This scale factor is also made
available for use with fix adapt.

\sphinxAtStartPar
Style \sphinxstyleemphasis{lj/cut/dipole/long} computes the short\sphinxhyphen{}range portion of
point\sphinxhyphen{}dipole interactions as discussed in {\hyperref[\detokenize{pair_dipole:toukmaji2}]{\sphinxcrossref{\DUrole{std,std-ref}{(Toukmaji)}}}}. Dipole\sphinxhyphen{}dipole, dipole\sphinxhyphen{}charge, and charge\sphinxhyphen{}charge
interactions are all supported, along with the standard 12/6
Lennard\sphinxhyphen{}Jones interactions, which are computed with a cutoff.  A
{\hyperref[\detokenize{kspace_style::doc}]{\sphinxcrossref{\DUrole{doc}{kspace\_style}}}} must be defined to use this pair
style.  If only dipoles (not point charges) are included in the model,
the kspace style can be one of these 3 options, all of which compute
the long\sphinxhyphen{}range portion of dipole\sphinxhyphen{}dipole interactions.  If the model
includes point charges (in addition to dipoles), then only the first
of these kspace styles can be used:
\begin{itemize}
\item {} 
\sphinxAtStartPar
{\hyperref[\detokenize{kspace_style::doc}]{\sphinxcrossref{\DUrole{doc}{kspace\_style ewald/disp}}}}

\item {} 
\sphinxAtStartPar
{\hyperref[\detokenize{kspace_style::doc}]{\sphinxcrossref{\DUrole{doc}{kspace\_style ewald/dipole}}}}

\item {} 
\sphinxAtStartPar
{\hyperref[\detokenize{kspace_style::doc}]{\sphinxcrossref{\DUrole{doc}{kspace\_style pppm/dipole}}}}

\end{itemize}

\sphinxAtStartPar
Style \sphinxstyleemphasis{lj/long/dipole/long} has the same functionality as style
\sphinxstyleemphasis{lj/cut/dipole/long}, except it also has an option to compute 12/6
Lennard\sphinxhyphen{}Jones interactions for use with a long\sphinxhyphen{}range dispersion kspace
style.  This is done by setting its \sphinxstyleemphasis{flag\_lj} argument to \sphinxstyleemphasis{long}.  For
long\sphinxhyphen{}range LJ interactions, the {\hyperref[\detokenize{kspace_style::doc}]{\sphinxcrossref{\DUrole{doc}{kspace\_style ewald/disp}}}} command must be used.


\bigskip\hrule\bigskip


\sphinxAtStartPar
The following coefficients must be defined for each pair of atoms types
via the {\hyperref[\detokenize{pair_coeff::doc}]{\sphinxcrossref{\DUrole{doc}{pair\_coeff}}}} command as in the examples above,
or in the data file or restart files read by the {\hyperref[\detokenize{read_data::doc}]{\sphinxcrossref{\DUrole{doc}{read\_data}}}} or {\hyperref[\detokenize{read_restart::doc}]{\sphinxcrossref{\DUrole{doc}{read\_restart}}}} commands, or by
mixing as described below:
\begin{itemize}
\item {} 
\sphinxAtStartPar
\(\epsilon\) (energy units)

\item {} 
\sphinxAtStartPar
\(\sigma\) (distance units)

\item {} 
\sphinxAtStartPar
cutoff1 (distance units)

\item {} 
\sphinxAtStartPar
cutoff2 (distance units)

\end{itemize}

\sphinxAtStartPar
The latter 2 coefficients are optional.  If not specified, the global LJ
and Coulombic cutoffs specified in the pair\_style command are used.  If
only one cutoff is specified, it is used as the cutoff for both LJ and
Coulombic interactions for this type pair.  If both coefficients are
specified, they are used as the LJ and Coulombic cutoffs for this type
pair.  When using a long\sphinxhyphen{}rang Coulomb solver, only a global Coulomb
cutoff may be used and only the LJ cutoff may be changed with the
{\hyperref[\detokenize{pair_coeff::doc}]{\sphinxcrossref{\DUrole{doc}{pair\_coeff}}}} command.  When using the
\sphinxstyleemphasis{lj/long/dipole/long} pair style with \sphinxstyleemphasis{long} \sphinxstyleemphasis{long} setting, only a
single global cutoff may be provided and no cutoff for the
{\hyperref[\detokenize{pair_coeff::doc}]{\sphinxcrossref{\DUrole{doc}{pair\_coeff}}}} command.


\bigskip\hrule\bigskip


\sphinxAtStartPar
Note that for systems using these pair styles, typically particles
should be able to exert torque on each other via their dipole moments
so that the particle and its dipole moment can rotate.  This requires
they not be point particles, but finite\sphinxhyphen{}size spheres.  Thus you should
use a command like {\hyperref[\detokenize{atom_style::doc}]{\sphinxcrossref{\DUrole{doc}{atom\_style hybrid sphere dipole}}}}
to use particles with both attributes.

\sphinxAtStartPar
The magnitude and orientation of the dipole moment for each particle
can be defined by the {\hyperref[\detokenize{set::doc}]{\sphinxcrossref{\DUrole{doc}{set}}}} command or in the “Atoms”
section of the data file read in by the {\hyperref[\detokenize{read_data::doc}]{\sphinxcrossref{\DUrole{doc}{read\_data}}}}
command.

\sphinxAtStartPar
Rotating finite\sphinxhyphen{}size particles have 6 degrees of freedom (DOFs),
translation and rotational.  You can use the {\hyperref[\detokenize{compute_temp_sphere::doc}]{\sphinxcrossref{\DUrole{doc}{compute temp/sphere}}}} command to monitor a temperature which includes
all these DOFs.

\sphinxAtStartPar
Finite\sphinxhyphen{}size particles with dipole moments should be integrated using
one of these options:
\begin{itemize}
\item {} 
\sphinxAtStartPar
{\hyperref[\detokenize{fix_nve_sphere::doc}]{\sphinxcrossref{\DUrole{doc}{fix nve/sphere update dipole}}}}

\item {} 
\sphinxAtStartPar
{\hyperref[\detokenize{fix_nve_sphere::doc}]{\sphinxcrossref{\DUrole{doc}{fix nve/sphere update dipole}}}} plus {\hyperref[\detokenize{fix_langevin::doc}]{\sphinxcrossref{\DUrole{doc}{fix langevin omega yes}}}}

\item {} 
\sphinxAtStartPar
{\hyperref[\detokenize{fix_nvt_sphere::doc}]{\sphinxcrossref{\DUrole{doc}{fix nvt/sphere update dipole}}}}

\item {} 
\sphinxAtStartPar
{\hyperref[\detokenize{fix_npt_sphere::doc}]{\sphinxcrossref{\DUrole{doc}{fix npt/sphere update dipole}}}}

\end{itemize}

\sphinxAtStartPar
In all cases the “update dipole” setting ensures the dipole moments
are also rotated when the finite\sphinxhyphen{}size spheres rotate.  The 2nd and 3rd
bullets perform thermostatting; in the case of a Langevin thermostat
the “omega yes” option also thermostats the rotational degrees of
freedom (if desired).  The 4th bullet performs thermostatting and
barostatting.


\bigskip\hrule\bigskip


\sphinxAtStartPar
Styles with a \sphinxstyleemphasis{gpu}, \sphinxstyleemphasis{intel}, \sphinxstyleemphasis{kk}, \sphinxstyleemphasis{omp}, or \sphinxstyleemphasis{opt} suffix are
functionally the same as the corresponding style without the suffix.
They have been optimized to run faster, depending on your available
hardware, as discussed on the {\hyperref[\detokenize{Speed_packages::doc}]{\sphinxcrossref{\DUrole{doc}{Accelerator packages}}}}
page.  The accelerated styles take the same arguments and should
produce the same results, except for round\sphinxhyphen{}off and precision issues.

\sphinxAtStartPar
These accelerated styles are part of the GPU, INTEL, KOKKOS,
OPENMP, and OPT packages, respectively.  They are only enabled if
LAMMPS was built with those packages.  See the {\hyperref[\detokenize{Build_package::doc}]{\sphinxcrossref{\DUrole{doc}{Build package}}}} page for more info.

\sphinxAtStartPar
You can specify the accelerated styles explicitly in your input script
by including their suffix, or you can use the {\hyperref[\detokenize{Run_options::doc}]{\sphinxcrossref{\DUrole{doc}{\sphinxhyphen{}suffix command\sphinxhyphen{}line switch}}}} when you invoke LAMMPS, or you can use the
{\hyperref[\detokenize{suffix::doc}]{\sphinxcrossref{\DUrole{doc}{suffix}}}} command in your input script.

\sphinxAtStartPar
See the {\hyperref[\detokenize{Speed_packages::doc}]{\sphinxcrossref{\DUrole{doc}{Accelerator packages}}}} page for more
instructions on how to use the accelerated styles effectively.


\bigskip\hrule\bigskip



\subsection{Mixing, shift, table, tail correction, restart, rRESPA info}
\label{\detokenize{pair_dipole:mixing-shift-table-tail-correction-restart-rrespa-info}}
\sphinxAtStartPar
For atom type pairs I,J and I != J, the epsilon and sigma coefficients
and cutoff distances for this pair style can be mixed.  The default
mix value is \sphinxstyleemphasis{geometric}.  See the “pair\_modify” command for details.

\sphinxAtStartPar
For atom type pairs I,J and I != J, the A, sigma, d1, and d2
coefficients and cutoff distance for this pair style can be mixed.  A
is an energy value mixed like a LJ epsilon.  D1 and d2 are distance
values and are mixed like sigma.  The default mix value is
\sphinxstyleemphasis{geometric}.  See the “pair\_modify” command for details.

\sphinxAtStartPar
This pair style does not support the {\hyperref[\detokenize{pair_modify::doc}]{\sphinxcrossref{\DUrole{doc}{pair\_modify}}}}
shift option for the energy of the Lennard\sphinxhyphen{}Jones portion of the pair
interaction; such energy goes to zero at the cutoff by construction.

\sphinxAtStartPar
The {\hyperref[\detokenize{pair_modify::doc}]{\sphinxcrossref{\DUrole{doc}{pair\_modify}}}} table option is not relevant
for this pair style.

\sphinxAtStartPar
This pair style does not support the {\hyperref[\detokenize{pair_modify::doc}]{\sphinxcrossref{\DUrole{doc}{pair\_modify}}}}
tail option for adding long\sphinxhyphen{}range tail corrections to energy and
pressure.

\sphinxAtStartPar
This pair style writes its information to {\hyperref[\detokenize{restart::doc}]{\sphinxcrossref{\DUrole{doc}{binary restart files}}}}, so pair\_style and pair\_coeff commands do not need
to be specified in an input script that reads a restart file.

\sphinxAtStartPar
This pair style can only be used via the \sphinxstyleemphasis{pair} keyword of the
{\hyperref[\detokenize{run_style::doc}]{\sphinxcrossref{\DUrole{doc}{run\_style respa}}}} command.  It does not support the
\sphinxstyleemphasis{inner}, \sphinxstyleemphasis{middle}, \sphinxstyleemphasis{outer} keywords.


\subsection{Restrictions}
\label{\detokenize{pair_dipole:restrictions}}
\sphinxAtStartPar
The \sphinxstyleemphasis{lj/cut/dipole/cut}, \sphinxstyleemphasis{lj/cut/dipole/long}, \sphinxstyleemphasis{lj/long/dipole/long},
and lj/sf/dipole/sf* styles are part of the DIPOLE package.  They are only
enabled if LAMMPS was built with that package.  See the
{\hyperref[\detokenize{Build_package::doc}]{\sphinxcrossref{\DUrole{doc}{Build package}}}} page for more info.

\sphinxAtStartPar
Using dipole pair styles with \sphinxstyleemphasis{electron} {\hyperref[\detokenize{units::doc}]{\sphinxcrossref{\DUrole{doc}{units}}}} is not currently supported.


\subsection{Related commands}
\label{\detokenize{pair_dipole:related-commands}}
\sphinxAtStartPar
{\hyperref[\detokenize{pair_coeff::doc}]{\sphinxcrossref{\DUrole{doc}{pair\_coeff}}}}, {\hyperref[\detokenize{set::doc}]{\sphinxcrossref{\DUrole{doc}{set}}}}, {\hyperref[\detokenize{read_data::doc}]{\sphinxcrossref{\DUrole{doc}{read\_data}}}},
{\hyperref[\detokenize{fix_nve_sphere::doc}]{\sphinxcrossref{\DUrole{doc}{fix nve/sphere}}}}, {\hyperref[\detokenize{fix_nvt_sphere::doc}]{\sphinxcrossref{\DUrole{doc}{fix nvt/sphere}}}}


\subsection{Default}
\label{\detokenize{pair_dipole:default}}
\sphinxAtStartPar
none


\bigskip\hrule\bigskip

\phantomsection\label{\detokenize{pair_dipole:allen2}}
\sphinxAtStartPar
\sphinxstylestrong{(AllenTildesley)} Allen and Tildesley, Computer Simulation of Liquids,
Clarendon Press, Oxford, 1987.

\phantomsection\label{\detokenize{pair_dipole:toukmaji2}}
\sphinxAtStartPar
\sphinxstylestrong{(Toukmaji)} Toukmaji, Sagui, Board, and Darden, J Chem Phys, 113,
10913 (2000).

\phantomsection\label{\detokenize{pair_dipole:stoddard}}
\sphinxAtStartPar
\sphinxstylestrong{(Stoddard)} Stoddard and Ford, Phys Rev A, 8, 1504 (1973).

\phantomsection\label{\detokenize{pair_dipole:price2}}
\sphinxAtStartPar
\sphinxstylestrong{(Price)} Price, Stone and Alderton, Mol Phys, 52, 987 (1984).

\sphinxstepscope

\index{pair\_style dpd@\spxentry{pair\_style dpd}}\ignorespaces 
\index{pair\_style dpd/gpu@\spxentry{pair\_style dpd/gpu}}\ignorespaces 
\index{pair\_style dpd/intel@\spxentry{pair\_style dpd/intel}}\ignorespaces 
\index{pair\_style dpd/kk@\spxentry{pair\_style dpd/kk}}\ignorespaces 
\index{pair\_style dpd/omp@\spxentry{pair\_style dpd/omp}}\ignorespaces 
\index{pair\_style dpd/tstat@\spxentry{pair\_style dpd/tstat}}\ignorespaces 
\index{pair\_style dpd/tstat/gpu@\spxentry{pair\_style dpd/tstat/gpu}}\ignorespaces 
\index{pair\_style dpd/tstat/kk@\spxentry{pair\_style dpd/tstat/kk}}\ignorespaces 
\index{pair\_style dpd/tstat/omp@\spxentry{pair\_style dpd/tstat/omp}}\ignorespaces 

\section{pair\_style dpd command}
\label{\detokenize{pair_dpd:pair-style-dpd-command}}\label{\detokenize{pair_dpd:index-8}}\label{\detokenize{pair_dpd:index-7}}\label{\detokenize{pair_dpd:index-6}}\label{\detokenize{pair_dpd:index-5}}\label{\detokenize{pair_dpd:index-4}}\label{\detokenize{pair_dpd:index-3}}\label{\detokenize{pair_dpd:index-2}}\label{\detokenize{pair_dpd:index-1}}\label{\detokenize{pair_dpd:index-0}}\label{\detokenize{pair_dpd::doc}}
\sphinxAtStartPar
Accelerator Variants: \sphinxstyleemphasis{dpd/gpu}, \sphinxstyleemphasis{dpd/intel}, \sphinxstyleemphasis{dpd/kk}, \sphinxstyleemphasis{dpd/omp}


\section{pair\_style dpd/tstat command}
\label{\detokenize{pair_dpd:pair-style-dpd-tstat-command}}
\sphinxAtStartPar
Accelerator Variants: \sphinxstyleemphasis{dpd/tstat/gpu}, \sphinxstyleemphasis{dpd/tstat/kk}, \sphinxstyleemphasis{dpd/tstat/omp}


\subsection{Syntax}
\label{\detokenize{pair_dpd:syntax}}
\begin{sphinxVerbatim}[commandchars=\\\{\}]
\PYG{k}{pair\PYGZus{}style}\PYG{+w}{ }\PYG{n}{dpd}\PYG{+w}{ }\PYG{n}{T}\PYG{+w}{ }\PYG{n}{cutoff}\PYG{+w}{ }\PYG{n}{seed}
\PYG{k}{pair\PYGZus{}style}\PYG{+w}{ }\PYG{n}{dpd}\PYG{o}{/}\PYG{n}{tstat}\PYG{+w}{ }\PYG{n}{Tstart}\PYG{+w}{ }\PYG{n}{Tstop}\PYG{+w}{ }\PYG{n}{cutoff}\PYG{+w}{ }\PYG{n}{seed}
\end{sphinxVerbatim}
\begin{itemize}
\item {} 
\sphinxAtStartPar
T = temperature (temperature units) (dpd only)

\item {} 
\sphinxAtStartPar
Tstart,Tstop = desired temperature at start/end of run (temperature units) (dpd/tstat only)

\item {} 
\sphinxAtStartPar
cutoff = global cutoff for DPD interactions (distance units)

\item {} 
\sphinxAtStartPar
seed = random \# seed (positive integer)

\end{itemize}


\subsection{Examples}
\label{\detokenize{pair_dpd:examples}}
\begin{sphinxVerbatim}[commandchars=\\\{\}]
\PYG{k}{pair\PYGZus{}style}\PYG{+w}{ }\PYG{n}{dpd}\PYG{+w}{ }\PYG{l+m}{1.0}\PYG{+w}{ }\PYG{l+m}{2.5}\PYG{+w}{ }\PYG{l+m}{34387}
\PYG{k}{pair\PYGZus{}coeff}\PYG{+w}{ }\PYG{o}{*}\PYG{+w}{ }\PYG{o}{*}\PYG{+w}{ }\PYG{l+m}{3.0}\PYG{+w}{ }\PYG{l+m}{1.0}
\PYG{k}{pair\PYGZus{}coeff}\PYG{+w}{ }\PYG{l+m}{1}\PYG{+w}{ }\PYG{l+m}{1}\PYG{+w}{ }\PYG{l+m}{3.0}\PYG{+w}{ }\PYG{l+m}{1.0}\PYG{+w}{ }\PYG{l+m}{1.0}

\PYG{k}{pair\PYGZus{}style}\PYG{+w}{ }\PYG{n}{hybrid}\PYG{o}{/}\PYG{n}{overlay}\PYG{+w}{ }\PYG{n}{lj}\PYG{o}{/}\PYG{n}{cut}\PYG{+w}{ }\PYG{l+m}{2.5}\PYG{+w}{ }\PYG{n}{dpd}\PYG{o}{/}\PYG{n}{tstat}\PYG{+w}{ }\PYG{l+m}{1.0}\PYG{+w}{ }\PYG{l+m}{1.0}\PYG{+w}{ }\PYG{l+m}{2.5}\PYG{+w}{ }\PYG{l+m}{34387}
\PYG{k}{pair\PYGZus{}coeff}\PYG{+w}{ }\PYG{o}{*}\PYG{+w}{ }\PYG{o}{*}\PYG{+w}{ }\PYG{n}{lj}\PYG{o}{/}\PYG{n}{cut}\PYG{+w}{ }\PYG{l+m}{1.0}\PYG{+w}{ }\PYG{l+m}{1.0}
\PYG{k}{pair\PYGZus{}coeff}\PYG{+w}{ }\PYG{o}{*}\PYG{+w}{ }\PYG{o}{*}\PYG{+w}{ }\PYG{n}{dpd}\PYG{o}{/}\PYG{n}{tstat}\PYG{+w}{ }\PYG{l+m}{1.0}
\end{sphinxVerbatim}


\subsection{Description}
\label{\detokenize{pair_dpd:description}}
\sphinxAtStartPar
Style \sphinxstyleemphasis{dpd} computes a force field for dissipative particle dynamics
(DPD) following the exposition in {\hyperref[\detokenize{pair_dpd:groot1}]{\sphinxcrossref{\DUrole{std,std-ref}{(Groot)}}}}.

\sphinxAtStartPar
Style \sphinxstyleemphasis{dpd/tstat} invokes a DPD thermostat on pairwise interactions,
which is equivalent to the non\sphinxhyphen{}conservative portion of the DPD force
field.  This pairwise thermostat can be used in conjunction with any
{\hyperref[\detokenize{pair_style::doc}]{\sphinxcrossref{\DUrole{doc}{pair style}}}}, and instead of per\sphinxhyphen{}particle thermostats
like {\hyperref[\detokenize{fix_langevin::doc}]{\sphinxcrossref{\DUrole{doc}{fix langevin}}}} or ensemble thermostats like
Nose Hoover as implemented by {\hyperref[\detokenize{fix_nh::doc}]{\sphinxcrossref{\DUrole{doc}{fix nvt}}}}.  To use
\sphinxstyleemphasis{dpd/tstat} as a thermostat for another pair style, use the
{\hyperref[\detokenize{pair_hybrid::doc}]{\sphinxcrossref{\DUrole{doc}{pair\_style hybrid/overlay}}}} command to compute both
the desired pair interaction and the thermostat for each pair of
particles.

\sphinxAtStartPar
For style \sphinxstyleemphasis{dpd}, the force on atom I due to atom J is given as a sum
of 3 terms
\begin{equation*}
\begin{split}\vec{f}  = & (F^C + F^D + F^R) \hat{r_{ij}} \qquad \qquad r < r_c \\
F^C      = & A w(r) \\
F^D      = & - \gamma w^2(r) (\hat{r_{ij}} \bullet \vec{v}_{ij}) \\
F^R      = & \sigma w(r) \alpha (\Delta t)^{-1/2} \\
w(r)     = & 1 - \frac{r}{r_c}\end{split}
\end{equation*}
\sphinxAtStartPar
where \(F^C\) is a conservative force, \(F^D\) is a dissipative
force, and \(F^R\) is a random force.  \(\hat{r_{ij}}\) is a
unit vector in the direction \(r_i - r_j\), \(\vec{v}_{ij}\) is
the vector difference in velocities of the two atoms \(\vec{v}_i -
\vec{v}_j\), \(\alpha\) is a Gaussian random number with zero mean
and unit variance, \sphinxstyleemphasis{dt} is the timestep size, and \(w(r)\) is a
weighting factor that varies between 0 and 1.  \(r_c\) is the
pairwise cutoff.  \(\sigma\) is set equal to \(\sqrt{2 k_B T
\gamma}\), where \(k_B\) is the Boltzmann constant and \sphinxstyleemphasis{T} is the
temperature parameter in the pair\_style command.

\sphinxAtStartPar
For style \sphinxstyleemphasis{dpd/tstat}, the force on atom I due to atom J is the same as
the above equation, except that the conservative \(F^C\) term is
dropped.  Also, during the run, \sphinxstyleemphasis{T} is set each timestep to a ramped
value from \sphinxstyleemphasis{Tstart} to \sphinxstyleemphasis{Tstop}.

\sphinxAtStartPar
For style \sphinxstyleemphasis{dpd}, the pairwise energy associated with style \sphinxstyleemphasis{dpd} is only
due to the conservative force term \(F^C\), and is shifted to be
zero at the cutoff distance \(r_c\).  The pairwise virial is
calculated using all 3 terms.  For style \sphinxstyleemphasis{dpd/tstat} there is no
pairwise energy, but the last two terms of the formula make a
contribution to the virial.

\sphinxAtStartPar
For style \sphinxstyleemphasis{dpd}, the following coefficients must be defined for each
pair of atoms types via the {\hyperref[\detokenize{pair_coeff::doc}]{\sphinxcrossref{\DUrole{doc}{pair\_coeff}}}} command as in
the examples above, or in the data file or restart files read by the
{\hyperref[\detokenize{read_data::doc}]{\sphinxcrossref{\DUrole{doc}{read\_data}}}} or {\hyperref[\detokenize{read_restart::doc}]{\sphinxcrossref{\DUrole{doc}{read\_restart}}}}
commands:
\begin{itemize}
\item {} 
\sphinxAtStartPar
A (force units)

\item {} 
\sphinxAtStartPar
\(\gamma\) (force/velocity units)

\item {} 
\sphinxAtStartPar
cutoff (distance units)

\end{itemize}

\sphinxAtStartPar
The cutoff coefficient is optional.  If not specified, the global DPD
cutoff is used.  Note that sigma is set equal to sqrt(2 T gamma),
where T is the temperature set by the {\hyperref[\detokenize{pair_style::doc}]{\sphinxcrossref{\DUrole{doc}{pair\_style}}}}
command so it does not need to be specified.

\sphinxAtStartPar
For style \sphinxstyleemphasis{dpd/tstat}, the coefficients defined for each pair of
atoms types via the {\hyperref[\detokenize{pair_coeff::doc}]{\sphinxcrossref{\DUrole{doc}{pair\_coeff}}}} command are:
\begin{itemize}
\item {} 
\sphinxAtStartPar
\(\gamma\) (force/velocity units)

\item {} 
\sphinxAtStartPar
cutoff (distance units)

\end{itemize}

\sphinxAtStartPar
The cutoff coefficient is optional.

\sphinxAtStartPar
Styles with a \sphinxstyleemphasis{gpu} suffix are implemented based on
the work of {\hyperref[\detokenize{pair_dpd:afshar}]{\sphinxcrossref{\DUrole{std,std-ref}{(Afshar)}}}} and {\hyperref[\detokenize{pair_dpd:phillips}]{\sphinxcrossref{\DUrole{std,std-ref}{(Phillips)}}}}.

\begin{sphinxadmonition}{note}{Note:}
\sphinxAtStartPar
If you are modeling DPD polymer chains, you may want to use the
{\hyperref[\detokenize{pair_srp::doc}]{\sphinxcrossref{\DUrole{doc}{pair\_style srp}}}} command in conjunction with these pair
styles.  It is a soft segmental repulsive potential (SRP) that can
prevent DPD polymer chains from crossing each other.
\end{sphinxadmonition}

\begin{sphinxadmonition}{note}{Note:}
\sphinxAtStartPar
The virial calculation for pressure when using these pair styles
includes all the components of force listed above, including the
random force.  Since the random force depends on random numbers,
everything that changes the order of atoms in the neighbor list
(e.g. different number of MPI ranks or a different neighbor list
skin distance) will also change the sequence in which the random
numbers are applied and thus the individual forces and therefore
also the virial/pressure.
\end{sphinxadmonition}

\begin{sphinxadmonition}{note}{Note:}
\sphinxAtStartPar
For more consistent time integration and force computation you may
consider using {\hyperref[\detokenize{fix_mvv_dpd::doc}]{\sphinxcrossref{\DUrole{doc}{fix mvv/dpd}}}} instead of {\hyperref[\detokenize{fix_nve::doc}]{\sphinxcrossref{\DUrole{doc}{fix
nve}}}}.
\end{sphinxadmonition}


\bigskip\hrule\bigskip


\sphinxAtStartPar
Styles with a \sphinxstyleemphasis{gpu}, \sphinxstyleemphasis{intel}, \sphinxstyleemphasis{kk}, \sphinxstyleemphasis{omp}, or \sphinxstyleemphasis{opt} suffix are
functionally the same as the corresponding style without the suffix.
They have been optimized to run faster, depending on your available
hardware, as discussed on the {\hyperref[\detokenize{Speed_packages::doc}]{\sphinxcrossref{\DUrole{doc}{Accelerator packages}}}}
page.  The accelerated styles take the same arguments and should
produce the same results, except for round\sphinxhyphen{}off and precision issues.

\sphinxAtStartPar
These accelerated styles are part of the GPU, INTEL, KOKKOS,
OPENMP, and OPT packages, respectively.  They are only enabled if
LAMMPS was built with those packages.  See the {\hyperref[\detokenize{Build_package::doc}]{\sphinxcrossref{\DUrole{doc}{Build package}}}} page for more info.

\sphinxAtStartPar
You can specify the accelerated styles explicitly in your input script
by including their suffix, or you can use the {\hyperref[\detokenize{Run_options::doc}]{\sphinxcrossref{\DUrole{doc}{\sphinxhyphen{}suffix command\sphinxhyphen{}line switch}}}} when you invoke LAMMPS, or you can use the
{\hyperref[\detokenize{suffix::doc}]{\sphinxcrossref{\DUrole{doc}{suffix}}}} command in your input script.

\sphinxAtStartPar
See the {\hyperref[\detokenize{Speed_packages::doc}]{\sphinxcrossref{\DUrole{doc}{Accelerator packages}}}} page for more
instructions on how to use the accelerated styles effectively.


\bigskip\hrule\bigskip



\subsection{Mixing, shift, table, tail correction, restart, rRESPA info}
\label{\detokenize{pair_dpd:mixing-shift-table-tail-correction-restart-rrespa-info}}
\sphinxAtStartPar
These pair styles do not support mixing.  Thus, coefficients for all
I,J pairs must be specified explicitly.

\sphinxAtStartPar
These pair styles do not support the {\hyperref[\detokenize{pair_modify::doc}]{\sphinxcrossref{\DUrole{doc}{pair\_modify}}}}
shift option for the energy of the pair interaction.  Note that as
discussed above, the energy due to the conservative \(F^C\) term is already
shifted to be 0.0 at the cutoff distance \(r_c\).

\sphinxAtStartPar
The {\hyperref[\detokenize{pair_modify::doc}]{\sphinxcrossref{\DUrole{doc}{pair\_modify}}}} table option is not relevant
for these pair styles.

\sphinxAtStartPar
These pair styles do not support the {\hyperref[\detokenize{pair_modify::doc}]{\sphinxcrossref{\DUrole{doc}{pair\_modify}}}}
tail option for adding long\sphinxhyphen{}range tail corrections to energy and
pressure.

\sphinxAtStartPar
These pair styles write their information to {\hyperref[\detokenize{restart::doc}]{\sphinxcrossref{\DUrole{doc}{binary restart files}}}}, so pair\_style and pair\_coeff commands do not need to be
specified in an input script that reads a restart file.  Note that the
user\sphinxhyphen{}specified random number seed is stored in the restart file, so when
a simulation is restarted, each processor will re\sphinxhyphen{}initialize its random
number generator the same way it did initially.  This means the random
forces will be random, but will not be the same as they would have been
if the original simulation had continued past the restart time.

\sphinxAtStartPar
These pair styles can only be used via the \sphinxstyleemphasis{pair} keyword of the
{\hyperref[\detokenize{run_style::doc}]{\sphinxcrossref{\DUrole{doc}{run\_style respa}}}} command.  They do not support the
\sphinxstyleemphasis{inner}, \sphinxstyleemphasis{middle}, \sphinxstyleemphasis{outer} keywords.

\sphinxAtStartPar
The \sphinxstyleemphasis{dpd/tstat} style can ramp its target temperature over multiple
runs, using the \sphinxstyleemphasis{start} and \sphinxstyleemphasis{stop} keywords of the {\hyperref[\detokenize{run::doc}]{\sphinxcrossref{\DUrole{doc}{run}}}}
command.  See the {\hyperref[\detokenize{run::doc}]{\sphinxcrossref{\DUrole{doc}{run}}}} command for details of how to do
this.


\bigskip\hrule\bigskip



\subsection{Restrictions}
\label{\detokenize{pair_dpd:restrictions}}
\sphinxAtStartPar
These styles are part of the DPD\sphinxhyphen{}BASIC package.  They are only enabled if
LAMMPS was built with that package.  See the {\hyperref[\detokenize{Build_package::doc}]{\sphinxcrossref{\DUrole{doc}{Build package}}}} page for more info.

\sphinxAtStartPar
The default frequency for rebuilding neighbor lists is every 10 steps
(see the {\hyperref[\detokenize{neigh_modify::doc}]{\sphinxcrossref{\DUrole{doc}{neigh\_modify}}}} command). This may be too
infrequent for style \sphinxstyleemphasis{dpd} simulations since particles move rapidly and
can overlap by large amounts.  If this setting yields a non\sphinxhyphen{}zero number
of “dangerous” reneighborings (printed at the end of a simulation), you
should experiment with forcing reneighboring more often and see if
system energies/trajectories change.

\sphinxAtStartPar
These pair styles requires you to use the {\hyperref[\detokenize{comm_modify::doc}]{\sphinxcrossref{\DUrole{doc}{comm\_modify vel yes}}}} command so that velocities are stored by ghost atoms.

\sphinxAtStartPar
These pair styles will not restart exactly when using the
{\hyperref[\detokenize{read_restart::doc}]{\sphinxcrossref{\DUrole{doc}{read\_restart}}}} command, though they should provide
statistically similar results.  This is because the forces they compute
depend on atom velocities.  See the {\hyperref[\detokenize{read_restart::doc}]{\sphinxcrossref{\DUrole{doc}{read\_restart}}}}
command for more details.


\subsection{Related commands}
\label{\detokenize{pair_dpd:related-commands}}
\sphinxAtStartPar
{\hyperref[\detokenize{pair_dpd_ext::doc}]{\sphinxcrossref{\DUrole{doc}{pair\_style dpd/ext}}}}, {\hyperref[\detokenize{pair_coeff::doc}]{\sphinxcrossref{\DUrole{doc}{pair\_coeff}}}},
{\hyperref[\detokenize{fix_nh::doc}]{\sphinxcrossref{\DUrole{doc}{fix nvt}}}}, {\hyperref[\detokenize{fix_langevin::doc}]{\sphinxcrossref{\DUrole{doc}{fix langevin}}}},
{\hyperref[\detokenize{pair_srp::doc}]{\sphinxcrossref{\DUrole{doc}{pair\_style srp}}}}, {\hyperref[\detokenize{fix_mvv_dpd::doc}]{\sphinxcrossref{\DUrole{doc}{fix mvv/dpd}}}}.


\subsection{Default}
\label{\detokenize{pair_dpd:default}}
\sphinxAtStartPar
none


\bigskip\hrule\bigskip

\phantomsection\label{\detokenize{pair_dpd:groot1}}
\sphinxAtStartPar
\sphinxstylestrong{(Groot)} Groot and Warren, J Chem Phys, 107, 4423\sphinxhyphen{}35 (1997).

\phantomsection\label{\detokenize{pair_dpd:afshar}}
\sphinxAtStartPar
\sphinxstylestrong{(Afshar)} Afshar, F. Schmid, A. Pishevar, S. Worley, Comput Phys
Comm, 184, 1119\sphinxhyphen{}1128 (2013).

\phantomsection\label{\detokenize{pair_dpd:phillips}}
\sphinxAtStartPar
\sphinxstylestrong{(Phillips)} C. L. Phillips, J. A. Anderson, S. C. Glotzer, Comput
Phys Comm, 230, 7191\sphinxhyphen{}7201 (2011).

\sphinxstepscope

\index{pair\_style dpd/coul/slater/long@\spxentry{pair\_style dpd/coul/slater/long}}\ignorespaces 
\index{pair\_style dpd/coul/slater/long/gpu@\spxentry{pair\_style dpd/coul/slater/long/gpu}}\ignorespaces 

\section{pair\_style dpd/coul/slater/long command}
\label{\detokenize{pair_dpd_coul_slater_long:pair-style-dpd-coul-slater-long-command}}\label{\detokenize{pair_dpd_coul_slater_long:index-1}}\label{\detokenize{pair_dpd_coul_slater_long:index-0}}\label{\detokenize{pair_dpd_coul_slater_long::doc}}
\sphinxAtStartPar
Accelerator Variants: \sphinxstyleemphasis{dpd/coul/slater/long/gpu}


\subsection{Syntax}
\label{\detokenize{pair_dpd_coul_slater_long:syntax}}
\begin{sphinxVerbatim}[commandchars=\\\{\}]
\PYG{k}{pair\PYGZus{}style}\PYG{+w}{ }\PYG{n}{dpd}\PYG{o}{/}\PYG{n}{coul}\PYG{o}{/}\PYG{n}{slater}\PYG{o}{/}\PYG{n}{long}\PYG{+w}{ }\PYG{n}{T}\PYG{+w}{ }\PYG{n}{cutoff\PYGZus{}DPD}\PYG{+w}{ }\PYG{n}{seed}\PYG{+w}{ }\PYG{n}{lambda}\PYG{+w}{ }\PYG{n}{cutoff\PYGZus{}coul}
\end{sphinxVerbatim}
\begin{itemize}
\item {} 
\sphinxAtStartPar
T = temperature (temperature units)

\item {} 
\sphinxAtStartPar
cutoff\_DPD = global cutoff for DPD interactions (distance units)

\item {} 
\sphinxAtStartPar
seed = random \# seed (positive integer)

\item {} 
\sphinxAtStartPar
lambda = decay length of the charge (distance units)

\item {} 
\sphinxAtStartPar
cutoff\_coul = global cutoff for Coulombic interactions (distance units)

\end{itemize}


\subsection{Examples}
\label{\detokenize{pair_dpd_coul_slater_long:examples}}
\begin{sphinxVerbatim}[commandchars=\\\{\}]
\PYG{k}{pair\PYGZus{}style}\PYG{+w}{ }\PYG{n}{dpd}\PYG{o}{/}\PYG{n}{coul}\PYG{o}{/}\PYG{n}{slater}\PYG{o}{/}\PYG{n}{long}\PYG{+w}{ }\PYG{l+m}{1.0}\PYG{+w}{ }\PYG{l+m}{2.5}\PYG{+w}{ }\PYG{l+m}{34387}\PYG{+w}{ }\PYG{l+m}{0.25}\PYG{+w}{ }\PYG{l+m}{3.0}

\PYG{k}{pair\PYGZus{}coeff}\PYG{+w}{ }\PYG{l+m}{1}\PYG{+w}{ }\PYG{l+m}{1}\PYG{+w}{ }\PYG{l+m}{78.0}\PYG{+w}{ }\PYG{l+m}{4.5}\PYG{+w}{          }\PYG{c}{\PYGZsh{} not charged by default}
\PYG{k}{pair\PYGZus{}coeff}\PYG{+w}{ }\PYG{l+m}{2}\PYG{+w}{ }\PYG{l+m}{2}\PYG{+w}{ }\PYG{l+m}{78.0}\PYG{+w}{ }\PYG{l+m}{4.5}\PYG{+w}{ }\PYG{n}{yes}
\end{sphinxVerbatim}


\subsection{Description}
\label{\detokenize{pair_dpd_coul_slater_long:description}}
\sphinxAtStartPar
\DUrole{versionmodified,added}{Added in version 27June2024.}

\sphinxAtStartPar
Style \sphinxstyleemphasis{dpd/coul/slater/long} computes a force field for dissipative
particle dynamics (DPD) following the exposition in {\hyperref[\detokenize{pair_dpd_coul_slater_long:groot5}]{\sphinxcrossref{\DUrole{std,std-ref}{(Groot)}}}}.  It also allows for the use of charged particles in the
model by adding a long\sphinxhyphen{}range Coulombic term to the DPD interactions.
The short\sphinxhyphen{}range portion of the Coulombics is calculated by this pair
style.  The long\sphinxhyphen{}range Coulombics are computed by use of the
{\hyperref[\detokenize{kspace_style::doc}]{\sphinxcrossref{\DUrole{doc}{kspace\_style}}}} command, e.g. using the Ewald or
PPPM styles.

\sphinxAtStartPar
Coulombic forces in mesoscopic models such as DPD employ potentials
without explicit excluded\sphinxhyphen{}volume interactions.  The goal is to prevent
artificial ionic pair formation by including a charge distribution in
the Coulomb potential, following the formulation in {\hyperref[\detokenize{pair_dpd_coul_slater_long:melchor1}]{\sphinxcrossref{\DUrole{std,std-ref}{(Melchor1)}}}}.

\begin{sphinxadmonition}{note}{Note:}
\sphinxAtStartPar
This pair style is effectively the combination of the
{\hyperref[\detokenize{pair_dpd::doc}]{\sphinxcrossref{\DUrole{doc}{pair\_style dpd}}}} and {\hyperref[\detokenize{pair_coul_slater::doc}]{\sphinxcrossref{\DUrole{doc}{pair\_style
coul/slater/long}}}} commands, but should be more
efficient (especially on GPUs) than using {\hyperref[\detokenize{pair_hybrid::doc}]{\sphinxcrossref{\DUrole{doc}{pair\_style
hybrid/overlay dpd coul/slater/long}}}}. That is
particularly true for the GPU package version of the pair style since
this version is compatible with computing neighbor lists on the GPU
instead of the CPU as is required for hybrid styles.
\end{sphinxadmonition}

\sphinxAtStartPar
In the charged DPD model, the force on bead I due to bead J is given
as a sum of 4 terms:
\begin{equation*}
\begin{split}\vec{f}  = & (F^C + F^D + F^R + F^E) \hat{r_{ij}} \\
F^C      = & A w(r) \qquad \qquad \qquad \qquad \qquad r < r_{DPD} \\
F^D      = & - \gamma w^2(r) (\hat{r_{ij}} \bullet \vec{v}_{ij}) \qquad \qquad r < r_{DPD} \\
F^R      = & \sigma w(r) \alpha (\Delta t)^{-1/2} \qquad \qquad \qquad r < r_{DPD} \\
w(r)     = & 1 - \frac{r}{r_{DPD}} \\
F^E      = & \frac{C q_iq_j}{\epsilon r^2} \left( 1- exp\left( \frac{2r_{ij}}{\lambda} \right) \left( 1 + \frac{2r_{ij}}{\lambda} \left( 1 + \frac{r_{ij}}{\lambda} \right)\right) \right)\end{split}
\end{equation*}
\sphinxAtStartPar
where \(F^C\) is a conservative force, \(F^D\) is a
dissipative force, \(F^R\) is a random force, and \(F^E\) is
an electrostatic force.  \(\hat{r_{ij}}\) is a unit vector in the
direction \(r_i - r_j\), \(\vec{v}_{ij}\) is the vector
difference in velocities of the two atoms \(\vec{v}_i -
\vec{v}_j\), \(\alpha\) is a Gaussian random number with zero mean
and unit variance, \sphinxstyleemphasis{dt} is the timestep size, and \(w(r)\) is a
weighting factor that varies between 0 and 1.

\sphinxAtStartPar
\(\sigma\) is set equal to \(\sqrt{2 k_B T \gamma}\), where
\(k_B\) is the Boltzmann constant and \sphinxstyleemphasis{T} is the temperature
parameter in the pair\_style command.

\sphinxAtStartPar
\(r_{DPD}\) is the pairwise cutoff for the first 3 DPD terms in
the formula as specified by \sphinxstyleemphasis{cutoff\_DPD}.  For the \(F^E\) term,
pairwise interactions within the specified \sphinxstyleemphasis{cutoff\_coul} distance are
computed directly; interactions beyond that distance are computed in
reciprocal space.  \sphinxstyleemphasis{C} is the same Coulomb conversion factor used in
the Coulombic formulas described on the {\hyperref[\detokenize{pair_coul::doc}]{\sphinxcrossref{\DUrole{doc}{pair\_coul}}}}
doc page.

\sphinxAtStartPar
The following parameters must be defined for each pair of atoms types
via the {\hyperref[\detokenize{pair_coeff::doc}]{\sphinxcrossref{\DUrole{doc}{pair\_coeff}}}} command as in the examples
above, or in the data file or restart files read by the
{\hyperref[\detokenize{read_data::doc}]{\sphinxcrossref{\DUrole{doc}{read\_data}}}} or {\hyperref[\detokenize{read_restart::doc}]{\sphinxcrossref{\DUrole{doc}{read\_restart}}}}
commands:
\begin{itemize}
\item {} 
\sphinxAtStartPar
A (force units)

\item {} 
\sphinxAtStartPar
\(\gamma\) (force/velocity units)

\item {} 
\sphinxAtStartPar
is\_charged (optional boolean, default = no)

\end{itemize}

\sphinxAtStartPar
The \sphinxstyleemphasis{is\_charged} parameter is optional and can be specified as \sphinxstyleemphasis{yes} or
\sphinxstyleemphasis{no}.  \sphinxstyleemphasis{Yes} should be used for interactions between two types of
charged particles.  \sphinxstyleemphasis{No} is the default and should be used for
interactions between two types of particles when one or both are
uncharged.


\bigskip\hrule\bigskip


\sphinxAtStartPar
Styles with a \sphinxstyleemphasis{gpu}, \sphinxstyleemphasis{intel}, \sphinxstyleemphasis{kk}, \sphinxstyleemphasis{omp}, or \sphinxstyleemphasis{opt} suffix are
functionally the same as the corresponding style without the suffix.
They have been optimized to run faster, depending on your available
hardware, as discussed on the {\hyperref[\detokenize{Speed_packages::doc}]{\sphinxcrossref{\DUrole{doc}{Accelerator packages}}}}
page.  The accelerated styles take the same arguments and should
produce the same results, except for round\sphinxhyphen{}off and precision issues.

\sphinxAtStartPar
These accelerated styles are part of the GPU, INTEL, KOKKOS,
OPENMP, and OPT packages, respectively.  They are only enabled if
LAMMPS was built with those packages.  See the {\hyperref[\detokenize{Build_package::doc}]{\sphinxcrossref{\DUrole{doc}{Build package}}}} page for more info.

\sphinxAtStartPar
You can specify the accelerated styles explicitly in your input script
by including their suffix, or you can use the {\hyperref[\detokenize{Run_options::doc}]{\sphinxcrossref{\DUrole{doc}{\sphinxhyphen{}suffix command\sphinxhyphen{}line switch}}}} when you invoke LAMMPS, or you can use the
{\hyperref[\detokenize{suffix::doc}]{\sphinxcrossref{\DUrole{doc}{suffix}}}} command in your input script.

\sphinxAtStartPar
See the {\hyperref[\detokenize{Speed_packages::doc}]{\sphinxcrossref{\DUrole{doc}{Accelerator packages}}}} page for more
instructions on how to use the accelerated styles effectively.


\bigskip\hrule\bigskip



\subsection{Mixing, shift, table, tail correction, restart, rRESPA info}
\label{\detokenize{pair_dpd_coul_slater_long:mixing-shift-table-tail-correction-restart-rrespa-info}}
\sphinxAtStartPar
This pair style does not support mixing.  Thus, coefficients for all
I,J pairs must be specified explicitly.

\sphinxAtStartPar
This pair style does not support the {\hyperref[\detokenize{pair_modify::doc}]{\sphinxcrossref{\DUrole{doc}{pair\_modify}}}}
shift option for the energy of the pair interaction.

\sphinxAtStartPar
The {\hyperref[\detokenize{pair_modify::doc}]{\sphinxcrossref{\DUrole{doc}{pair\_modify}}}} table option is not relevant
for this pair style.

\sphinxAtStartPar
This pair style does not support the {\hyperref[\detokenize{pair_modify::doc}]{\sphinxcrossref{\DUrole{doc}{pair\_modify}}}}
tail option for adding long\sphinxhyphen{}range tail corrections to energy and
pressure.

\sphinxAtStartPar
This pair style writes its information to {\hyperref[\detokenize{restart::doc}]{\sphinxcrossref{\DUrole{doc}{binary restart files}}}}, so pair\_style and pair\_coeff commands do not need to be
specified in an input script that reads a restart file.  Note that the
user\sphinxhyphen{}specified random number seed is stored in the restart file, so
when a simulation is restarted, each processor will re\sphinxhyphen{}initialize its
random number generator the same way it did initially.  This means the
random forces will be random, but will not be the same as they would
have been if the original simulation had continued past the restart
time.

\sphinxAtStartPar
This pair style can only be used via the \sphinxstyleemphasis{pair} keyword of the
{\hyperref[\detokenize{run_style::doc}]{\sphinxcrossref{\DUrole{doc}{run\_style respa}}}} command.  It does not support the
\sphinxstyleemphasis{inner}, \sphinxstyleemphasis{middle}, \sphinxstyleemphasis{outer} keywords.


\bigskip\hrule\bigskip



\subsection{Restrictions}
\label{\detokenize{pair_dpd_coul_slater_long:restrictions}}
\sphinxAtStartPar
This style is part of the DPD\sphinxhyphen{}BASIC package.  It is only enabled if
LAMMPS was built with that package.  See the {\hyperref[\detokenize{Build_package::doc}]{\sphinxcrossref{\DUrole{doc}{Build package}}}} page for more info.

\sphinxAtStartPar
The default frequency for rebuilding neighbor lists is every 10 steps
(see the {\hyperref[\detokenize{neigh_modify::doc}]{\sphinxcrossref{\DUrole{doc}{neigh\_modify}}}} command). This may be too
infrequent since particles move rapidly and can overlap by large
amounts.  If this setting yields a non\sphinxhyphen{}zero number of “dangerous”
reneighborings (printed at the end of a simulation), you should
experiment with forcing reneighboring more often and see if system
energies/trajectories change.

\sphinxAtStartPar
This pair style requires use of the {\hyperref[\detokenize{comm_modify::doc}]{\sphinxcrossref{\DUrole{doc}{comm\_modify vel yes}}}} command so that velocities are stored by ghost atoms.

\sphinxAtStartPar
This pair style also requires use of a long\sphinxhyphen{}range solvers from the
KSPACE package.

\sphinxAtStartPar
This pair style will not restart exactly when using the
{\hyperref[\detokenize{read_restart::doc}]{\sphinxcrossref{\DUrole{doc}{read\_restart}}}} command, though they should provide
statistically similar results.  This is because the forces they compute
depend on atom velocities.  See the {\hyperref[\detokenize{read_restart::doc}]{\sphinxcrossref{\DUrole{doc}{read\_restart}}}}
command for more details.


\subsection{Related commands}
\label{\detokenize{pair_dpd_coul_slater_long:related-commands}}
\sphinxAtStartPar
{\hyperref[\detokenize{pair_dpd::doc}]{\sphinxcrossref{\DUrole{doc}{pair\_style dpd}}}}, {\hyperref[\detokenize{pair_coul_slater::doc}]{\sphinxcrossref{\DUrole{doc}{pair\_style coul/slater/long}}}},


\subsection{Default}
\label{\detokenize{pair_dpd_coul_slater_long:default}}
\sphinxAtStartPar
For the pair\_coeff command, the default is is\_charged = no.


\bigskip\hrule\bigskip

\phantomsection\label{\detokenize{pair_dpd_coul_slater_long:groot5}}
\sphinxAtStartPar
\sphinxstylestrong{(Groot)} Groot and Warren, J Chem Phys, 107, 4423\sphinxhyphen{}35 (1997).

\phantomsection\label{\detokenize{pair_dpd_coul_slater_long:melchor1}}
\sphinxAtStartPar
\sphinxstylestrong{(Melchor)} Gonzalez\sphinxhyphen{}Melchor, Mayoral, Velazquez, and Alejandre, J Chem Phys, 125, 224107 (2006).

\sphinxstepscope

\index{pair\_style dpd/ext@\spxentry{pair\_style dpd/ext}}\ignorespaces 
\index{pair\_style dpd/ext/kk@\spxentry{pair\_style dpd/ext/kk}}\ignorespaces 
\index{pair\_style dpd/ext/omp@\spxentry{pair\_style dpd/ext/omp}}\ignorespaces 
\index{pair\_style dpd/ext/tstat@\spxentry{pair\_style dpd/ext/tstat}}\ignorespaces 
\index{pair\_style dpd/ext/tstat/kk@\spxentry{pair\_style dpd/ext/tstat/kk}}\ignorespaces 
\index{pair\_style dpd/ext/tstat/omp@\spxentry{pair\_style dpd/ext/tstat/omp}}\ignorespaces 

\section{pair\_style dpd/ext command}
\label{\detokenize{pair_dpd_ext:pair-style-dpd-ext-command}}\label{\detokenize{pair_dpd_ext:index-5}}\label{\detokenize{pair_dpd_ext:index-4}}\label{\detokenize{pair_dpd_ext:index-3}}\label{\detokenize{pair_dpd_ext:index-2}}\label{\detokenize{pair_dpd_ext:index-1}}\label{\detokenize{pair_dpd_ext:index-0}}\label{\detokenize{pair_dpd_ext::doc}}
\sphinxAtStartPar
Accelerator Variants: dpd/ext/kk dpd/ext/omp


\section{pair\_style dpd/ext/tstat command}
\label{\detokenize{pair_dpd_ext:pair-style-dpd-ext-tstat-command}}
\sphinxAtStartPar
Accelerator Variants: dpd/ext/tstat/kk dpd/ext/tstat/omp


\subsection{Syntax}
\label{\detokenize{pair_dpd_ext:syntax}}
\begin{sphinxVerbatim}[commandchars=\\\{\}]
\PYG{k}{pair\PYGZus{}style}\PYG{+w}{ }\PYG{n}{dpd}\PYG{o}{/}\PYG{n}{ext}\PYG{+w}{ }\PYG{n}{T}\PYG{+w}{ }\PYG{n}{cutoff}\PYG{+w}{ }\PYG{n}{seed}
\PYG{k}{pair\PYGZus{}style}\PYG{+w}{ }\PYG{n}{dpd}\PYG{o}{/}\PYG{n}{ext}\PYG{o}{/}\PYG{n}{tstat}\PYG{+w}{ }\PYG{n}{Tstart}\PYG{+w}{ }\PYG{n}{Tstop}\PYG{+w}{ }\PYG{n}{cutoff}\PYG{+w}{ }\PYG{n}{seed}
\end{sphinxVerbatim}
\begin{itemize}
\item {} 
\sphinxAtStartPar
T = temperature (temperature units)

\item {} 
\sphinxAtStartPar
Tstart,Tstop = desired temperature at start/end of run (temperature units)

\item {} 
\sphinxAtStartPar
cutoff = global cutoff for DPD interactions (distance units)

\item {} 
\sphinxAtStartPar
seed = random \# seed (positive integer)

\end{itemize}


\subsection{Examples}
\label{\detokenize{pair_dpd_ext:examples}}
\begin{sphinxVerbatim}[commandchars=\\\{\}]
\PYG{k}{pair\PYGZus{}style}\PYG{+w}{ }\PYG{n}{dpd}\PYG{o}{/}\PYG{n}{ext}\PYG{+w}{ }\PYG{l+m}{1.0}\PYG{+w}{ }\PYG{l+m}{2.5}\PYG{+w}{ }\PYG{l+m}{34387}
\PYG{k}{pair\PYGZus{}coeff}\PYG{+w}{ }\PYG{l+m}{1}\PYG{+w}{ }\PYG{l+m}{1}\PYG{+w}{ }\PYG{l+m}{25.0}\PYG{+w}{ }\PYG{l+m}{4.5}\PYG{+w}{ }\PYG{l+m}{4.5}\PYG{+w}{ }\PYG{l+m}{0.5}\PYG{+w}{ }\PYG{l+m}{0.5}\PYG{+w}{ }\PYG{l+m}{1.2}
\PYG{k}{pair\PYGZus{}coeff}\PYG{+w}{ }\PYG{l+m}{1}\PYG{+w}{ }\PYG{l+m}{2}\PYG{+w}{ }\PYG{l+m}{40.0}\PYG{+w}{ }\PYG{l+m}{4.5}\PYG{+w}{ }\PYG{l+m}{4.5}\PYG{+w}{ }\PYG{l+m}{0.5}\PYG{+w}{ }\PYG{l+m}{0.5}\PYG{+w}{ }\PYG{l+m}{1.2}

\PYG{k}{pair\PYGZus{}style}\PYG{+w}{ }\PYG{n}{hybrid}\PYG{o}{/}\PYG{n}{overlay}\PYG{+w}{ }\PYG{n}{lj}\PYG{o}{/}\PYG{n}{cut}\PYG{+w}{ }\PYG{l+m}{2.5}\PYG{+w}{ }\PYG{n}{dpd}\PYG{o}{/}\PYG{n}{ext}\PYG{o}{/}\PYG{n}{tstat}\PYG{+w}{ }\PYG{l+m}{1.0}\PYG{+w}{ }\PYG{l+m}{1.0}\PYG{+w}{ }\PYG{l+m}{2.5}\PYG{+w}{ }\PYG{l+m}{34387}
\PYG{k}{pair\PYGZus{}coeff}\PYG{+w}{ }\PYG{o}{*}\PYG{+w}{ }\PYG{o}{*}\PYG{+w}{ }\PYG{n}{lj}\PYG{o}{/}\PYG{n}{cut}\PYG{+w}{ }\PYG{l+m}{1.0}\PYG{+w}{ }\PYG{l+m}{1.0}
\PYG{k}{pair\PYGZus{}coeff}\PYG{+w}{ }\PYG{o}{*}\PYG{+w}{ }\PYG{o}{*}\PYG{+w}{ }\PYG{l+m}{4.5}\PYG{+w}{ }\PYG{l+m}{4.5}\PYG{+w}{ }\PYG{l+m}{0.5}\PYG{+w}{ }\PYG{l+m}{0.5}\PYG{+w}{ }\PYG{l+m}{1.2}
\end{sphinxVerbatim}


\subsection{Description}
\label{\detokenize{pair_dpd_ext:description}}
\sphinxAtStartPar
The style \sphinxstyleemphasis{dpd/ext} computes an extended force field for dissipative
particle dynamics (DPD) following the exposition in {\hyperref[\detokenize{pair_dpd_ext:groot}]{\sphinxcrossref{\DUrole{std,std-ref}{(Groot)}}}}, {\hyperref[\detokenize{pair_dpd_ext:junghans}]{\sphinxcrossref{\DUrole{std,std-ref}{(Junghans)}}}}.

\sphinxAtStartPar
Style \sphinxstyleemphasis{dpd/ext/tstat} invokes an extended DPD thermostat on pairwise
interactions, equivalent to the non\sphinxhyphen{}conservative portion of the extended
DPD force field. To use \sphinxstyleemphasis{dpd/ext/tstat} as a thermostat for another pair
style, use the {\hyperref[\detokenize{pair_hybrid::doc}]{\sphinxcrossref{\DUrole{doc}{pair\_style hybrid/overlay}}}} command to
compute both the desired pair interaction and the thermostat for each
pair of particles.

\sphinxAtStartPar
For the style \sphinxstyleemphasis{dpd/ext}, the force on atom I due to atom J is given as
a sum of 3 terms
\begin{equation*}
\begin{split}\mathbf{f}  = & f^C + f^D + f^R \qquad \qquad r < r_c \\
f^C      = & A_{ij} w(r) \hat{\mathbf{r}}_{ij} \\
f^D      = & - \gamma_{\parallel} w_{\parallel}^2(r) (\hat{\mathbf{r}}_{ij} \cdot \mathbf{v}_{ij}) \hat{\mathbf{r}}_{ij}  - \gamma_{\perp} w_{\perp}^2 (r) ( \mathbf{I} - \hat{\mathbf{r}}_{ij} \hat{\mathbf{r}}_{ij}^{\rm T} ) \mathbf{v}_{ij} \\
f^R      = & \sigma_{\parallel} w_{\parallel}(r) \frac{\alpha}{\sqrt{\Delta t}} \hat{\mathbf{r}}_{ij}  + \sigma_{\perp} w_{\perp} (r) ( \mathbf{I} - \hat{\mathbf{r}}_{ij} \hat{\mathbf{r}}_{ij}^{\rm T} ) \frac{\mathbf{\xi}_{ij}}{\sqrt{\Delta t}}\\
w(r)     = & 1 - r/r_c \\\end{split}
\end{equation*}
\sphinxAtStartPar
where \(\mathbf{f}^C\) is a conservative force, \(\mathbf{f}^D\)
is a dissipative force, and \(\mathbf{f}^R\) is a random
force. \(A_{ij}\) is the maximum repulsion between the two atoms,
\(\hat{\mathbf{r}}_{ij}\) is a unit vector in the direction
\(\mathbf{r}_i - \mathbf{r}_j\), \(\mathbf{v}_{ij} =
\mathbf{v}_i - \mathbf{v}_j\) is the vector difference in velocities of
the two atoms, \(\alpha\) and \(\mathbf{\xi}_{ij}\) are Gaussian
random numbers with zero mean and unit variance, \(\Delta t\) is the
timestep, \(w (r) = 1 - r / r_c\) is a weight function for the
conservative interactions that varies between 0 and 1, \(r_c\) is
the corresponding cutoff, \(w_{\alpha} ( r ) = ( 1 - r / \bar{r}_c
)^{s_{\alpha}}\), \(\alpha \equiv ( \parallel, \perp )\), are weight
functions with coefficients \(s_\alpha\) that vary between 0 and 1,
\(\bar{r}_c\) is the corresponding cutoff, \(\mathbf{I}\) is the
unit matrix, \(\sigma_{\alpha} = \sqrt{2 k_B T \gamma_{\alpha}}\),
where \(k_B\) is the Boltzmann constant and \(T\) is the
temperature in the pair\_style command.

\sphinxAtStartPar
For the style \sphinxstyleemphasis{dpd/ext/tstat}, the force on atom I due to atom J is
the same as the above equation, except that the conservative
\(\mathbf{f}^C\) term is dropped. Also, during the run, T is set
each timestep to a ramped value from Tstart to Tstop.

\sphinxAtStartPar
For the style \sphinxstyleemphasis{dpd/ext}, the pairwise energy associated with style
\sphinxstyleemphasis{dpd/ext} is only due to the conservative force term
\(\mathbf{f}^C\), and is shifted to be zero at the cutoff distance
\(r_c\). The pairwise virial is calculated using all three
terms. There is no pairwise energy for style \sphinxstyleemphasis{dpd/ext/tstat}, but the
last two terms of the formula contribute the virial.

\sphinxAtStartPar
For the style \sphinxstyleemphasis{dpd/ext/tstat}, the force on atom I due to atom J is the
same as the above equation, except that the conservative
\(\mathbf{f}^C\) term is dropped.  Also, during the run, T is set
each timestep to a ramped value from Tstart to Tstop.

\sphinxAtStartPar
For the style \sphinxstyleemphasis{dpd/ext}, the pairwise energy associated with style
\sphinxstyleemphasis{dpd/ext} is only due to the conservative force term
\(\mathbf{f}^C\), and is shifted to be zero at the cutoff distance
\(r_c\). The pairwise virial is calculated using all three
terms. There is no pairwise energy for style \sphinxstyleemphasis{dpd/ext/tstat}, but the
last two terms of the formula contribute the virial.

\sphinxAtStartPar
For the style \sphinxstyleemphasis{dpd/ext}, the following coefficients must be defined for
each pair of atoms types via the {\hyperref[\detokenize{pair_coeff::doc}]{\sphinxcrossref{\DUrole{doc}{pair\_coeff}}}} command
as in the examples above:
\begin{itemize}
\item {} 
\sphinxAtStartPar
A (force units)

\item {} 
\sphinxAtStartPar
\(\gamma_{\parallel}\) (force/velocity units)

\item {} 
\sphinxAtStartPar
\(\gamma_{\perp}\) (force/velocity units)

\item {} 
\sphinxAtStartPar
\(s_{\parallel}\) (unitless)

\item {} 
\sphinxAtStartPar
\(s_{\perp}\) (unitless)

\item {} 
\sphinxAtStartPar
\(r_c\) (distance units)

\end{itemize}

\sphinxAtStartPar
The last coefficient is optional. If not specified, the global DPD
cutoff is used. Note that \(\sigma\)’s are set equal to
\(\sqrt{2 k_B T \gamma}\), where \(T\) is the temperature set by
the {\hyperref[\detokenize{pair_style::doc}]{\sphinxcrossref{\DUrole{doc}{pair\_style}}}} command so it does not need to be
specified.

\sphinxAtStartPar
For the style \sphinxstyleemphasis{dpd/ext/tstat}, the coefficients defined for each pair of
atoms types via the {\hyperref[\detokenize{pair_coeff::doc}]{\sphinxcrossref{\DUrole{doc}{pair\_coeff}}}} command are:
\begin{itemize}
\item {} 
\sphinxAtStartPar
\(\gamma_{\parallel}\) (force/velocity units)

\item {} 
\sphinxAtStartPar
\(\gamma_{\perp}\) (force/velocity units)

\item {} 
\sphinxAtStartPar
\(s_{\parallel}\) (unitless)

\item {} 
\sphinxAtStartPar
\(s_{\perp}\) (unitless)

\item {} 
\sphinxAtStartPar
\(r_c\) (distance units)

\end{itemize}

\sphinxAtStartPar
The last coefficient is optional.

\begin{sphinxadmonition}{note}{Note:}
\sphinxAtStartPar
If you are modeling DPD polymer chains, you may want to use the
{\hyperref[\detokenize{pair_srp::doc}]{\sphinxcrossref{\DUrole{doc}{pair\_style srp}}}} command in conjunction with these pair
styles.  It is a soft segmental repulsive potential (SRP) that can
prevent DPD polymer chains from crossing each other.
\end{sphinxadmonition}

\begin{sphinxadmonition}{note}{Note:}
\sphinxAtStartPar
The virial calculation for pressure when using these pair styles
includes all the components of force listed above, including the
random force.  Since the random force depends on random numbers,
everything that changes the order of atoms in the neighbor list
(e.g. different number of MPI ranks or a different neighbor list
skin distance) will also change the sequence in which the random
numbers are applied and thus the individual forces and therefore
also the virial/pressure.
\end{sphinxadmonition}

\begin{sphinxadmonition}{note}{Note:}
\sphinxAtStartPar
For more consistent time integration and force computation you may
consider using {\hyperref[\detokenize{fix_mvv_dpd::doc}]{\sphinxcrossref{\DUrole{doc}{fix mvv/dpd}}}} instead of {\hyperref[\detokenize{fix_nve::doc}]{\sphinxcrossref{\DUrole{doc}{fix
nve}}}}.
\end{sphinxadmonition}


\bigskip\hrule\bigskip


\sphinxAtStartPar
Styles with a \sphinxstyleemphasis{gpu}, \sphinxstyleemphasis{intel}, \sphinxstyleemphasis{kk}, \sphinxstyleemphasis{omp}, or \sphinxstyleemphasis{opt} suffix are
functionally the same as the corresponding style without the suffix.
They have been optimized to run faster, depending on your available
hardware, as discussed on the {\hyperref[\detokenize{Speed_packages::doc}]{\sphinxcrossref{\DUrole{doc}{Accelerator packages}}}}
page.  The accelerated styles take the same arguments and should
produce the same results, except for round\sphinxhyphen{}off and precision issues.

\sphinxAtStartPar
These accelerated styles are part of the GPU, INTEL, KOKKOS,
OPENMP, and OPT packages, respectively.  They are only enabled if
LAMMPS was built with those packages.  See the {\hyperref[\detokenize{Build_package::doc}]{\sphinxcrossref{\DUrole{doc}{Build package}}}} page for more info.

\sphinxAtStartPar
You can specify the accelerated styles explicitly in your input script
by including their suffix, or you can use the {\hyperref[\detokenize{Run_options::doc}]{\sphinxcrossref{\DUrole{doc}{\sphinxhyphen{}suffix command\sphinxhyphen{}line switch}}}} when you invoke LAMMPS, or you can use the
{\hyperref[\detokenize{suffix::doc}]{\sphinxcrossref{\DUrole{doc}{suffix}}}} command in your input script.

\sphinxAtStartPar
See the {\hyperref[\detokenize{Speed_packages::doc}]{\sphinxcrossref{\DUrole{doc}{Accelerator packages}}}} page for more
instructions on how to use the accelerated styles effectively.


\bigskip\hrule\bigskip


\sphinxAtStartPar
\sphinxstylestrong{Mixing, shift, table, tail correction, restart, rRESPA info}:

\sphinxAtStartPar
The style \sphinxstyleemphasis{dpd/ext} does not support mixing. Thus, coefficients for all
I,J pairs must be specified explicitly.

\sphinxAtStartPar
The pair styles do not support the {\hyperref[\detokenize{pair_modify::doc}]{\sphinxcrossref{\DUrole{doc}{pair\_modify}}}}
shift option for the energy of the pair interaction. Note that as
discussed above, the energy due to the conservative \(\mathbf{f}^C\)
term is already shifted to be zero at the cutoff distance \(r_c\).

\sphinxAtStartPar
The {\hyperref[\detokenize{pair_modify::doc}]{\sphinxcrossref{\DUrole{doc}{pair\_modify}}}} table option is not relevant for
the style \sphinxstyleemphasis{dpd/ext}.

\sphinxAtStartPar
The style \sphinxstyleemphasis{dpd/ext} does not support the {\hyperref[\detokenize{pair_modify::doc}]{\sphinxcrossref{\DUrole{doc}{pair\_modify}}}} tail option for adding long\sphinxhyphen{}range tail corrections to
energy and pressure.

\sphinxAtStartPar
The pair styles can only be used via the pair keyword of the
{\hyperref[\detokenize{run_style::doc}]{\sphinxcrossref{\DUrole{doc}{run\_style respa}}}} command. They do not support the
\sphinxstyleemphasis{inner}, \sphinxstyleemphasis{middle}, and \sphinxstyleemphasis{outer}keywords.

\sphinxAtStartPar
The style \sphinxstyleemphasis{dpd/ext/tstat} can ramp its target temperature over multiple
runs, using the start and stop keywords of the {\hyperref[\detokenize{run::doc}]{\sphinxcrossref{\DUrole{doc}{run}}}}
command. See the {\hyperref[\detokenize{run::doc}]{\sphinxcrossref{\DUrole{doc}{run}}}} command for details of how to do this.


\bigskip\hrule\bigskip



\subsection{Restrictions}
\label{\detokenize{pair_dpd_ext:restrictions}}
\sphinxAtStartPar
These styles are part of the DPD\sphinxhyphen{}BASIC package.  They are only enabled if
LAMMPS was built with that package.  See the {\hyperref[\detokenize{Build_package::doc}]{\sphinxcrossref{\DUrole{doc}{Build package}}}} page for more info.

\sphinxAtStartPar
The default frequency for rebuilding neighbor lists is every 10 steps
(see the {\hyperref[\detokenize{neigh_modify::doc}]{\sphinxcrossref{\DUrole{doc}{neigh\_modify}}}} command). This may be too
infrequent for style \sphinxstyleemphasis{dpd/ext} simulations since particles move rapidly
and can overlap by large amounts. If this setting yields a non\sphinxhyphen{}zero
number of say\{dangerous\} reneighborings (printed at the end of a
simulation), you should experiment with forcing reneighboring more often
and see if system energies/trajectories change.

\sphinxAtStartPar
The pair styles require to use the {\hyperref[\detokenize{comm_modify::doc}]{\sphinxcrossref{\DUrole{doc}{comm\_modify vel yes}}}} command so that velocities are stored by ghost atoms.

\sphinxAtStartPar
The pair styles will not restart exactly when using the
{\hyperref[\detokenize{read_restart::doc}]{\sphinxcrossref{\DUrole{doc}{read\_restart}}}} command, though they should provide
statistically similar results. This is because the forces they compute
depend on atom velocities. See the {\hyperref[\detokenize{read_restart::doc}]{\sphinxcrossref{\DUrole{doc}{read\_restart}}}}
command for more details.


\subsection{Related commands}
\label{\detokenize{pair_dpd_ext:related-commands}}
\sphinxAtStartPar
{\hyperref[\detokenize{pair_dpd::doc}]{\sphinxcrossref{\DUrole{doc}{pair\_style dpd}}}}, {\hyperref[\detokenize{pair_coeff::doc}]{\sphinxcrossref{\DUrole{doc}{pair\_coeff}}}},
{\hyperref[\detokenize{fix_nh::doc}]{\sphinxcrossref{\DUrole{doc}{fix nvt}}}}, {\hyperref[\detokenize{fix_langevin::doc}]{\sphinxcrossref{\DUrole{doc}{fix langevin}}}},
{\hyperref[\detokenize{pair_srp::doc}]{\sphinxcrossref{\DUrole{doc}{pair\_style srp}}}}, {\hyperref[\detokenize{fix_mvv_dpd::doc}]{\sphinxcrossref{\DUrole{doc}{fix mvv/dpd}}}}.

\sphinxAtStartPar
\sphinxstylestrong{Default:} none


\bigskip\hrule\bigskip

\phantomsection\label{\detokenize{pair_dpd_ext:groot}}
\sphinxAtStartPar
\sphinxstylestrong{(Groot)} Groot and Warren, J Chem Phys, 107, 4423\sphinxhyphen{}35 (1997).

\phantomsection\label{\detokenize{pair_dpd_ext:junghans}}
\sphinxAtStartPar
\sphinxstylestrong{(Junghans)} Junghans, Praprotnik and Kremer, Soft Matter 4, 156, 1119\sphinxhyphen{}1128 (2008).

\sphinxstepscope

\index{pair\_style dpd/fdt@\spxentry{pair\_style dpd/fdt}}\ignorespaces 
\index{pair\_style dpd/fdt/energy@\spxentry{pair\_style dpd/fdt/energy}}\ignorespaces 
\index{pair\_style dpd/fdt/energy/kk@\spxentry{pair\_style dpd/fdt/energy/kk}}\ignorespaces 

\section{pair\_style dpd/fdt command}
\label{\detokenize{pair_dpd_fdt:pair-style-dpd-fdt-command}}\label{\detokenize{pair_dpd_fdt:index-2}}\label{\detokenize{pair_dpd_fdt:index-1}}\label{\detokenize{pair_dpd_fdt:index-0}}\label{\detokenize{pair_dpd_fdt::doc}}

\section{pair\_style dpd/fdt/energy command}
\label{\detokenize{pair_dpd_fdt:pair-style-dpd-fdt-energy-command}}
\sphinxAtStartPar
Accelerator Variants: \sphinxstyleemphasis{dpd/fdt/energy/kk}


\subsection{Syntax}
\label{\detokenize{pair_dpd_fdt:syntax}}
\begin{sphinxVerbatim}[commandchars=\\\{\}]
\PYG{k}{pair\PYGZus{}style}\PYG{+w}{ }\PYG{n}{style}\PYG{+w}{ }\PYG{n}{args}
\end{sphinxVerbatim}
\begin{itemize}
\item {} 
\sphinxAtStartPar
style = \sphinxstyleemphasis{dpd/fdt} or \sphinxstyleemphasis{dpd/fdt/energy}

\item {} 
\sphinxAtStartPar
args = list of arguments for a particular style

\end{itemize}
\begin{sphinxalltt}
\sphinxstyleemphasis{dpd/fdt} args = T cutoff seed
  T = temperature (temperature units)
  cutoff = global cutoff for DPD interactions (distance units)
  seed = random \# seed (positive integer)
\sphinxstyleemphasis{dpd/fdt/energy} args = cutoff seed
  cutoff = global cutoff for DPD interactions (distance units)
  seed = random \# seed (positive integer)
\end{sphinxalltt}


\subsection{Examples}
\label{\detokenize{pair_dpd_fdt:examples}}
\begin{sphinxVerbatim}[commandchars=\\\{\}]
\PYG{k}{pair\PYGZus{}style}\PYG{+w}{ }\PYG{n}{dpd}\PYG{o}{/}\PYG{n}{fdt}\PYG{+w}{ }\PYG{l+m}{300.0}\PYG{+w}{ }\PYG{l+m}{2.5}\PYG{+w}{ }\PYG{l+m}{34387}
\PYG{k}{pair\PYGZus{}coeff}\PYG{+w}{ }\PYG{o}{*}\PYG{+w}{ }\PYG{o}{*}\PYG{+w}{ }\PYG{l+m}{3.0}\PYG{+w}{ }\PYG{l+m}{1.0}\PYG{+w}{ }\PYG{l+m}{2.5}

\PYG{k}{pair\PYGZus{}style}\PYG{+w}{ }\PYG{n}{dpd}\PYG{o}{/}\PYG{n}{fdt}\PYG{o}{/}\PYG{n}{energy}\PYG{+w}{ }\PYG{l+m}{2.5}\PYG{+w}{ }\PYG{l+m}{34387}
\PYG{k}{pair\PYGZus{}coeff}\PYG{+w}{ }\PYG{o}{*}\PYG{+w}{ }\PYG{o}{*}\PYG{+w}{ }\PYG{l+m}{3.0}\PYG{+w}{ }\PYG{l+m}{1.0}\PYG{+w}{ }\PYG{l+m}{0.1}\PYG{+w}{ }\PYG{l+m}{2.5}
\end{sphinxVerbatim}


\subsection{Description}
\label{\detokenize{pair_dpd_fdt:description}}
\sphinxAtStartPar
Styles \sphinxstyleemphasis{dpd/fdt} and \sphinxstyleemphasis{dpd/fdt/energy} compute the force for dissipative
particle dynamics (DPD) simulations.  The \sphinxstyleemphasis{dpd/fdt} style is used to
perform DPD simulations under isothermal and isobaric conditions,
while the \sphinxstyleemphasis{dpd/fdt/energy} style is used to perform DPD simulations
under isoenergetic and isoenthalpic conditions (see {\hyperref[\detokenize{pair_dpd_fdt:lisal3}]{\sphinxcrossref{\DUrole{std,std-ref}{(Lisal)}}}}).
For DPD simulations in general, the force on atom I due to atom J is
given as a sum of 3 terms
\begin{equation*}
\begin{split}\vec{f}  = & (F^C + F^D + F^R) \hat{r_{ij}} \qquad \qquad r < r_c \\
F^C      = & A w(r) \\
F^D      = & - \gamma w^2(r) (\hat{r_{ij}} \bullet \vec{v}_{ij}) \\
F^R      = & \sigma w(r) \alpha (\Delta t)^{-1/2} \\
w(r)     = & 1 - \frac{r}{r_c}\end{split}
\end{equation*}
\sphinxAtStartPar
where \(F^C\) is a conservative force, \(F^D\) is a dissipative
force, and \(F^R\) is a random force.  \(\hat{r_{ij}}\) is a
unit vector in the direction \(r_i - r_j\), \(\vec{v}_{ij}\) is
the vector difference in velocities of the two atoms, \(\vec{v}_i -
\vec{v}_j\), \(\alpha\) is a Gaussian random number with zero mean
and unit variance, \sphinxstyleemphasis{dt} is the timestep size, and \(w(r)\) is a
weighting factor that varies between 0 and 1, \(r_c\) is the
pairwise cutoff.  Note that alternative definitions of the weighting
function exist, but would have to be implemented as a separate pair
style command.

\sphinxAtStartPar
For style \sphinxstyleemphasis{dpd/fdt}, the fluctuation\sphinxhyphen{}dissipation theorem defines
\(\gamma\) to be set equal to \(\sigma^2/(2 T)\), where \sphinxstyleemphasis{T} is the
set point temperature specified as a pair style parameter in the above
examples.  The following coefficients must be defined for each pair of
atoms types via the {\hyperref[\detokenize{pair_coeff::doc}]{\sphinxcrossref{\DUrole{doc}{pair\_coeff}}}} command as in the
examples above, or in the data file or restart files read by the
{\hyperref[\detokenize{read_data::doc}]{\sphinxcrossref{\DUrole{doc}{read\_data}}}} or {\hyperref[\detokenize{read_restart::doc}]{\sphinxcrossref{\DUrole{doc}{read\_restart}}}}
commands:
\begin{itemize}
\item {} 
\sphinxAtStartPar
A (force units)

\item {} 
\sphinxAtStartPar
\(\sigma\) (force*time\textasciicircum{}(1/2) units)

\item {} 
\sphinxAtStartPar
cutoff (distance units)

\end{itemize}

\sphinxAtStartPar
The last coefficient is optional.  If not specified, the global DPD
cutoff is used.

\sphinxAtStartPar
Style \sphinxstyleemphasis{dpd/fdt/energy} is used to perform DPD simulations under
isoenergetic and isoenthalpic conditions.  The fluctuation\sphinxhyphen{}dissipation
theorem defines \(\gamma\) to be set equal to \(\sigma^2/(2
\theta)\), where \(\theta\) is the average internal temperature for
the pair. The particle internal temperature is related to the particle
internal energy through a mesoparticle equation of state (see {\hyperref[\detokenize{fix::doc}]{\sphinxcrossref{\DUrole{doc}{fix
eos}}}}). The differential internal conductive and mechanical
energies are computed within style \sphinxstyleemphasis{dpd/fdt/energy} as:
\begin{equation*}
\begin{split}du_{i}^{cond}  = & \kappa_{ij}(\frac{1}{\theta_{i}}-\frac{1}{\theta_{j}})\omega_{ij}^{2} + \alpha_{ij}\omega_{ij}\zeta_{ij}^{q}(\Delta{t})^{-1/2} \\
du_{i}^{mech}  = & -\frac{1}{2}\gamma_{ij}\omega_{ij}^{2}(\frac{\vec{r}_{ij}}{r_{ij}}\bullet\vec{v}_{ij})^{2} -
\frac{\sigma^{2}_{ij}}{4}(\frac{1}{m_{i}}+\frac{1}{m_{j}})\omega_{ij}^{2} -
\frac{1}{2}\sigma_{ij}\omega_{ij}(\frac{\vec{r}_{ij}}{r_{ij}}\bullet\vec{v}_{ij})\zeta_{ij}(\Delta{t})^{-1/2}\end{split}
\end{equation*}
\sphinxAtStartPar
where
\begin{equation*}
\begin{split}\alpha_{ij}^{2}  = & 2k_{B}\kappa_{ij} \\
\sigma^{2}_{ij}  = & 2\gamma_{ij}k_{B}\Theta_{ij} \\
\Theta_{ij}^{-1}  = & \frac{1}{2}(\frac{1}{\theta_{i}}+\frac{1}{\theta_{j}})\end{split}
\end{equation*}
\sphinxAtStartPar
\(\zeta_ij^q\) is a second Gaussian random number with zero mean and
unit variance that is used to compute the internal conductive
energy. The fluctuation\sphinxhyphen{}dissipation theorem defines \(alpha^2\) to
be set equal to \(2k_B\kappa\), where \(\kappa\) is the
mesoparticle thermal conductivity parameter.  The following coefficients
must be defined for each pair of atoms types via the {\hyperref[\detokenize{pair_coeff::doc}]{\sphinxcrossref{\DUrole{doc}{pair\_coeff}}}} command as in the examples above, or in the data file or
restart files read by the {\hyperref[\detokenize{read_data::doc}]{\sphinxcrossref{\DUrole{doc}{read\_data}}}} or
{\hyperref[\detokenize{read_restart::doc}]{\sphinxcrossref{\DUrole{doc}{read\_restart}}}} commands:
\begin{itemize}
\item {} 
\sphinxAtStartPar
A (force units)

\item {} 
\sphinxAtStartPar
\(\sigma\) (force*time\textasciicircum{}(1/2) units)

\item {} 
\sphinxAtStartPar
\(\kappa\) (energy*temperature/time units)

\item {} 
\sphinxAtStartPar
cutoff (distance units)

\end{itemize}

\sphinxAtStartPar
The last coefficient is optional.  If not specified, the global DPD
cutoff is used.

\sphinxAtStartPar
The pairwise energy associated with styles \sphinxstyleemphasis{dpd/fdt} and
\sphinxstyleemphasis{dpd/fdt/energy} is only due to the conservative force term \(F^C\),
and is shifted to be zero at the cutoff distance \(r_c\).  The
pairwise virial is calculated using only the conservative term.

\sphinxAtStartPar
The forces computed through the \sphinxstyleemphasis{dpd/fdt} and \sphinxstyleemphasis{dpd/fdt/energy} styles
can be integrated with the velocity\sphinxhyphen{}Verlet integration scheme or the
Shardlow splitting integration scheme described by {\hyperref[\detokenize{pair_dpd_fdt:lisal3}]{\sphinxcrossref{\DUrole{std,std-ref}{(Lisal)}}}}.  In the cases when these pair styles are combined with the
{\hyperref[\detokenize{fix_shardlow::doc}]{\sphinxcrossref{\DUrole{doc}{fix shardlow}}}}, these pair styles differ from the
other dpd styles in that the dissipative and random forces are split
from the force calculation and are not computed within the pair style.
Thus, only the conservative force is computed by the pair style, while
the stochastic integration of the dissipative and random forces are
handled through the Shardlow splitting algorithm approach.  The Shardlow
splitting algorithm is advantageous, especially when performing DPD
under isoenergetic conditions, as it allows significantly larger
timesteps to be taken.


\bigskip\hrule\bigskip


\sphinxAtStartPar
Styles with a \sphinxstyleemphasis{gpu}, \sphinxstyleemphasis{intel}, \sphinxstyleemphasis{kk}, \sphinxstyleemphasis{omp}, or \sphinxstyleemphasis{opt} suffix are
functionally the same as the corresponding style without the suffix.
They have been optimized to run faster, depending on your available
hardware, as discussed on the {\hyperref[\detokenize{Speed_packages::doc}]{\sphinxcrossref{\DUrole{doc}{Accelerator packages}}}}
page.  The accelerated styles take the same arguments and should
produce the same results, except for round\sphinxhyphen{}off and precision issues.

\sphinxAtStartPar
These accelerated styles are part of the GPU, INTEL, KOKKOS,
OPENMP, and OPT packages, respectively.  They are only enabled if
LAMMPS was built with those packages.  See the {\hyperref[\detokenize{Build_package::doc}]{\sphinxcrossref{\DUrole{doc}{Build package}}}} page for more info.

\sphinxAtStartPar
You can specify the accelerated styles explicitly in your input script
by including their suffix, or you can use the {\hyperref[\detokenize{Run_options::doc}]{\sphinxcrossref{\DUrole{doc}{\sphinxhyphen{}suffix command\sphinxhyphen{}line switch}}}} when you invoke LAMMPS, or you can use the
{\hyperref[\detokenize{suffix::doc}]{\sphinxcrossref{\DUrole{doc}{suffix}}}} command in your input script.

\sphinxAtStartPar
See the {\hyperref[\detokenize{Speed_packages::doc}]{\sphinxcrossref{\DUrole{doc}{Accelerator packages}}}} page for more
instructions on how to use the accelerated styles effectively.


\bigskip\hrule\bigskip



\subsection{Restrictions}
\label{\detokenize{pair_dpd_fdt:restrictions}}
\sphinxAtStartPar
These commands are part of the DPD\sphinxhyphen{}REACT package.  They are only enabled
if LAMMPS was built with that package.  See the {\hyperref[\detokenize{Build_package::doc}]{\sphinxcrossref{\DUrole{doc}{Build package}}}} page for more info.

\sphinxAtStartPar
Pair styles \sphinxstyleemphasis{dpd/fdt} and \sphinxstyleemphasis{dpd/fdt/energy} require use of the
{\hyperref[\detokenize{comm_modify::doc}]{\sphinxcrossref{\DUrole{doc}{comm\_modify vel yes}}}} option so that velocities are
stored by ghost atoms.

\sphinxAtStartPar
Pair style \sphinxstyleemphasis{dpd/fdt/energy} requires {\hyperref[\detokenize{atom_style::doc}]{\sphinxcrossref{\DUrole{doc}{atom\_style dpd}}}}
to be used in order to properly account for the particle internal
energies and temperatures.


\subsection{Related commands}
\label{\detokenize{pair_dpd_fdt:related-commands}}
\sphinxAtStartPar
{\hyperref[\detokenize{pair_coeff::doc}]{\sphinxcrossref{\DUrole{doc}{pair\_coeff}}}}, {\hyperref[\detokenize{fix_shardlow::doc}]{\sphinxcrossref{\DUrole{doc}{fix shardlow}}}}


\subsection{Default}
\label{\detokenize{pair_dpd_fdt:default}}
\sphinxAtStartPar
none


\bigskip\hrule\bigskip

\phantomsection\label{\detokenize{pair_dpd_fdt:lisal3}}
\sphinxAtStartPar
\sphinxstylestrong{(Lisal)} M. Lisal, J.K. Brennan, J. Bonet Avalos, J. Chem. Phys., 135, 204105 (2011).

\sphinxstepscope

\index{pair\_style drip@\spxentry{pair\_style drip}}\ignorespaces 

\section{pair\_style drip command}
\label{\detokenize{pair_drip:pair-style-drip-command}}\label{\detokenize{pair_drip:index-0}}\label{\detokenize{pair_drip::doc}}

\subsection{Syntax}
\label{\detokenize{pair_drip:syntax}}
\begin{sphinxVerbatim}[commandchars=\\\{\}]
\PYG{k}{pair\PYGZus{}style}\PYG{+w}{ }\PYG{n}{hybrid}\PYG{o}{/}\PYG{n}{overlay}\PYG{+w}{ }\PYG{n}{drip}\PYG{+w}{ }\PYG{n}{[styles}\PYG{+w}{ }\PYG{n}{...]}
\end{sphinxVerbatim}
\begin{itemize}
\item {} 
\sphinxAtStartPar
styles = other styles to be overlayed with drip (optional)

\end{itemize}


\subsection{Examples}
\label{\detokenize{pair_drip:examples}}
\begin{sphinxVerbatim}[commandchars=\\\{\}]
\PYG{k}{pair\PYGZus{}style}\PYG{+w}{ }\PYG{n}{hybrid}\PYG{o}{/}\PYG{n}{overlay}\PYG{+w}{ }\PYG{n}{drip}
\PYG{k}{pair\PYGZus{}coeff}\PYG{+w}{ }\PYG{o}{*}\PYG{+w}{ }\PYG{o}{*}\PYG{+w}{ }\PYG{n}{none}
\PYG{k}{pair\PYGZus{}coeff}\PYG{+w}{ }\PYG{o}{*}\PYG{+w}{ }\PYG{o}{*}\PYG{+w}{ }\PYG{n}{drip}\PYG{+w}{  }\PYG{n}{C.drip}\PYG{+w}{  }\PYG{n}{C}

\PYG{k}{pair\PYGZus{}style}\PYG{+w}{ }\PYG{n}{hybrid}\PYG{o}{/}\PYG{n}{overlay}\PYG{+w}{ }\PYG{n}{drip}\PYG{+w}{ }\PYG{n}{rebo}
\PYG{k}{pair\PYGZus{}coeff}\PYG{+w}{ }\PYG{o}{*}\PYG{+w}{ }\PYG{o}{*}\PYG{+w}{ }\PYG{n}{drip}\PYG{+w}{  }\PYG{n}{C.drip}\PYG{+w}{     }\PYG{n}{C}
\PYG{k}{pair\PYGZus{}coeff}\PYG{+w}{ }\PYG{o}{*}\PYG{+w}{ }\PYG{o}{*}\PYG{+w}{ }\PYG{n}{rebo}\PYG{+w}{  }\PYG{n}{CH.airebo}\PYG{+w}{  }\PYG{n}{C}

\PYG{k}{pair\PYGZus{}style}\PYG{+w}{ }\PYG{n}{hybrid}\PYG{o}{/}\PYG{n}{overlay}\PYG{+w}{ }\PYG{n}{drip}\PYG{+w}{ }\PYG{n}{rebo}
\PYG{k}{pair\PYGZus{}coeff}\PYG{+w}{ }\PYG{o}{*}\PYG{+w}{ }\PYG{o}{*}\PYG{+w}{ }\PYG{n}{drip}\PYG{+w}{  }\PYG{n}{C.drip}\PYG{+w}{     }\PYG{n}{C}\PYG{+w}{ }\PYG{n}{NULL}
\PYG{k}{pair\PYGZus{}coeff}\PYG{+w}{ }\PYG{o}{*}\PYG{+w}{ }\PYG{o}{*}\PYG{+w}{ }\PYG{n}{rebo}\PYG{+w}{  }\PYG{n}{CH.airebo}\PYG{+w}{  }\PYG{n}{C}\PYG{+w}{ }\PYG{n}{H}
\end{sphinxVerbatim}


\subsection{Description}
\label{\detokenize{pair_drip:description}}
\sphinxAtStartPar
Style \sphinxstyleemphasis{drip} computes the interlayer interactions of layered materials using
the dihedral\sphinxhyphen{}angle\sphinxhyphen{}corrected registry\sphinxhyphen{}dependent (DRIP) potential as described
in {\hyperref[\detokenize{pair_drip:wen2018}]{\sphinxcrossref{\DUrole{std,std-ref}{(Wen)}}}}, which is based on the {\hyperref[\detokenize{pair_drip:kolmogorov2005}]{\sphinxcrossref{\DUrole{std,std-ref}{(Kolmogorov)}}}}
potential and provides an improved prediction for forces.
The total potential energy of a system is
\begin{equation*}
\begin{split}E = & \frac{1}{2} \sum_{i} \sum_{j\notin\text{layer}\,i} \phi_{ij} \\
\phi_{ij} = &f_\text{c}(x_r) \left[ e^{-\lambda(r_{ij} - z_0 )} \left[C+f(\rho_{ij})+  g(\rho_{ij}, \{\alpha_{ij}^{(m)}\}) \right]- A\left (\frac{z_0}{r_{ij}} \right)^6 \right]\end{split}
\end{equation*}
\sphinxAtStartPar
where the \(r^{-6}\) term models the attractive London dispersion,
the exponential term is designed to capture the registry effect due to
overlapping \sphinxstyleemphasis{pi} bonds, and \sphinxstyleemphasis{fc} is a cutoff function.

\sphinxAtStartPar
This potential (DRIP) only provides the interlayer interactions between
graphene layers. So, to perform a realistic simulation, it should be used in
combination with an intralayer potential such as {\hyperref[\detokenize{pair_airebo::doc}]{\sphinxcrossref{\DUrole{doc}{REBO}}}} and
{\hyperref[\detokenize{pair_tersoff::doc}]{\sphinxcrossref{\DUrole{doc}{Tersoff}}}}.
To keep the intralayer interactions unaffected, we should avoid applying DRIP
to contribute energy to intralayer interactions. This can be achieved by
assigning different molecular IDs to atoms in different layers, and DRIP is
implemented such that only atoms with different molecular ID can interact with
each other. For this purpose, {\hyperref[\detokenize{atom_style::doc}]{\sphinxcrossref{\DUrole{doc}{atom style}}}} “molecular” or
“full” has to be used.

\sphinxAtStartPar
On the other way around, {\hyperref[\detokenize{pair_airebo::doc}]{\sphinxcrossref{\DUrole{doc}{REBO}}}} ({\hyperref[\detokenize{pair_tersoff::doc}]{\sphinxcrossref{\DUrole{doc}{Tersoff}}}}
or any other potential used to provide the intralayer interactions) should not
interfere with the interlayer interactions described by DRIP. This is typically
automatically achieved using the commands provided in the \sphinxstyleemphasis{Examples} section
above, since the cutoff distance for carbon\sphinxhyphen{}carbon interaction in the intralayer
potentials (e.g. 2 Angstrom for {\hyperref[\detokenize{pair_airebo::doc}]{\sphinxcrossref{\DUrole{doc}{REBO}}}}) is much smaller than
the equilibrium layer distance of graphene layers (about 3.4 Angstrom).
If you want, you can enforce this by assigning different atom types to atoms in
different layers, and apply an intralayer potential to one atom type.
See {\hyperref[\detokenize{pair_hybrid::doc}]{\sphinxcrossref{\DUrole{doc}{pair\_hybrid}}}} for details.


\bigskip\hrule\bigskip


\sphinxAtStartPar
The {\hyperref[\detokenize{pair_coeff::doc}]{\sphinxcrossref{\DUrole{doc}{pair\_coeff}}}} command for DRIP takes \sphinxstyleemphasis{4+N} arguments, where
\sphinxstyleemphasis{N} is the number of LAMMPS atom types. The fist three arguments must be fixed
to be \sphinxstyleemphasis{* * drip}, the fourth argument is the path to the DRIP parameter file,
and the remaining N arguments specifying the mapping between element in the
parameter file and atom types. For example, if your LAMMPS simulation has 3 atom
types and you want all of them to be C, you would use the following pair\_coeff
command:

\begin{sphinxVerbatim}[commandchars=\\\{\}]
\PYG{k}{pair\PYGZus{}coeff}\PYG{+w}{ }\PYG{o}{*}\PYG{+w}{ }\PYG{o}{*}\PYG{+w}{ }\PYG{n}{drip}\PYG{+w}{  }\PYG{n}{C.drip}\PYG{+w}{  }\PYG{n}{C}\PYG{+w}{ }\PYG{n}{C}\PYG{+w}{ }\PYG{n}{C}
\end{sphinxVerbatim}

\sphinxAtStartPar
If a mapping value is specified as NULL, the mapping is not performed. This
could be useful when DRIP is used to model part of the system where other
element exists. Suppose you have a hydrocarbon system, with C of atom type 1
and H of atom type 2, you can use the following command to inform DRIP not to
model H atoms:

\begin{sphinxVerbatim}[commandchars=\\\{\}]
\PYG{k}{pair\PYGZus{}style}\PYG{+w}{ }\PYG{n}{hybrid}\PYG{o}{/}\PYG{n}{overlay}\PYG{+w}{ }\PYG{n}{drip}\PYG{+w}{ }\PYG{n}{rebo}
\PYG{k}{pair\PYGZus{}coeff}\PYG{+w}{ }\PYG{o}{*}\PYG{+w}{ }\PYG{o}{*}\PYG{+w}{ }\PYG{n}{drip}\PYG{+w}{  }\PYG{n}{C.drip}\PYG{+w}{     }\PYG{n}{C}\PYG{+w}{ }\PYG{n}{NULL}
\PYG{k}{pair\PYGZus{}coeff}\PYG{+w}{ }\PYG{o}{*}\PYG{+w}{ }\PYG{o}{*}\PYG{+w}{ }\PYG{n}{rebo}\PYG{+w}{  }\PYG{n}{CH.airebo}\PYG{+w}{  }\PYG{n}{C}\PYG{+w}{ }\PYG{n}{H}
\end{sphinxVerbatim}

\begin{sphinxadmonition}{note}{Note:}
\sphinxAtStartPar
The potential parameters developed in {\hyperref[\detokenize{pair_drip:wen2018}]{\sphinxcrossref{\DUrole{std,std-ref}{(Wen)}}}} are provided with
LAMMPS (see the “potentials” directory). Besides those in {\hyperref[\detokenize{pair_drip:wen2018}]{\sphinxcrossref{\DUrole{std,std-ref}{Wen}}}}, an
additional parameter “normal\_cutoff”, specific to the LAMMPS implementation, is
used to find the three nearest neighbors of an atom to construct the normal.
\end{sphinxadmonition}


\bigskip\hrule\bigskip



\subsection{Mixing, shift, table, tail correction, restart, rRESPA info}
\label{\detokenize{pair_drip:mixing-shift-table-tail-correction-restart-rrespa-info}}
\sphinxAtStartPar
This pair style does not support the pair\_modify mix, shift, table,
and tail options.

\sphinxAtStartPar
This pair style does not write their information to binary restart files, since
it is stored in potential files. Thus, you need to re\sphinxhyphen{}specify the pair\_style and
pair\_coeff commands in an input script that reads a restart file.


\subsection{Restrictions}
\label{\detokenize{pair_drip:restrictions}}
\sphinxAtStartPar
This pair style is part of the INTERLAYER package. It is only enabled if
LAMMPS was built with that package.  See the {\hyperref[\detokenize{Build_package::doc}]{\sphinxcrossref{\DUrole{doc}{Build package}}}} page for more info.

\sphinxAtStartPar
This pair style requires the {\hyperref[\detokenize{newton::doc}]{\sphinxcrossref{\DUrole{doc}{newton}}}} setting to be “on” for
pair interactions.

\sphinxAtStartPar
The \sphinxstyleemphasis{C.drip} parameter file provided with LAMMPS (see the “potentials”
directory) is parameterized for metal {\hyperref[\detokenize{units::doc}]{\sphinxcrossref{\DUrole{doc}{units}}}}. You can use
the DRIP potential with any LAMMPS units, but you would need to create
your own custom parameter file with coefficients listed in the
appropriate units, if your simulation does not use “metal” units.


\subsection{Related commands}
\label{\detokenize{pair_drip:related-commands}}
\sphinxAtStartPar
{\hyperref[\detokenize{pair_lebedeva_z::doc}]{\sphinxcrossref{\DUrole{doc}{pair\_style lebedeva\_z}}}},
{\hyperref[\detokenize{pair_kolmogorov_crespi_z::doc}]{\sphinxcrossref{\DUrole{doc}{pair\_style kolmogorov/crespi/z}}}},
{\hyperref[\detokenize{pair_kolmogorov_crespi_full::doc}]{\sphinxcrossref{\DUrole{doc}{pair\_style kolmogorov/crespi/full}}}},
{\hyperref[\detokenize{pair_ilp_graphene_hbn::doc}]{\sphinxcrossref{\DUrole{doc}{pair\_style ilp/graphene/hbn}}}}.


\bigskip\hrule\bigskip

\phantomsection\label{\detokenize{pair_drip:wen2018}}
\sphinxAtStartPar
\sphinxstylestrong{(Wen)} M. Wen, S. Carr, S. Fang, E. Kaxiras, and E. B. Tadmor, Phys. Rev. B,
98, 235404 (2018)

\phantomsection\label{\detokenize{pair_drip:kolmogorov2005}}
\sphinxAtStartPar
\sphinxstylestrong{(Kolmogorov)} A. N. Kolmogorov, V. H. Crespi, Phys. Rev. B 71, 235415 (2005)

\sphinxstepscope

\index{pair\_style dsmc@\spxentry{pair\_style dsmc}}\ignorespaces 

\section{pair\_style dsmc command}
\label{\detokenize{pair_dsmc:pair-style-dsmc-command}}\label{\detokenize{pair_dsmc:index-0}}\label{\detokenize{pair_dsmc::doc}}

\subsection{Syntax}
\label{\detokenize{pair_dsmc:syntax}}
\begin{sphinxVerbatim}[commandchars=\\\{\}]
\PYG{k}{pair\PYGZus{}style}\PYG{+w}{ }\PYG{n}{dsmc}\PYG{+w}{ }\PYG{n}{max\PYGZus{}cell\PYGZus{}size}\PYG{+w}{ }\PYG{n}{seed}\PYG{+w}{ }\PYG{n}{weighting}\PYG{+w}{ }\PYG{n}{Tref}\PYG{+w}{ }\PYG{n}{Nrecompute}\PYG{+w}{ }\PYG{n}{Nsample}
\end{sphinxVerbatim}
\begin{itemize}
\item {} 
\sphinxAtStartPar
max\_cell\_size = global maximum cell size for DSMC interactions (distance units)

\item {} 
\sphinxAtStartPar
seed = random \# seed (positive integer)

\item {} 
\sphinxAtStartPar
weighting = macroparticle weighting

\item {} 
\sphinxAtStartPar
Tref = reference temperature (temperature units)

\item {} 
\sphinxAtStartPar
Nrecompute = re\sphinxhyphen{}compute v*sigma\_max every this many timesteps (timesteps)

\item {} 
\sphinxAtStartPar
Nsample = sample this many times in recomputing v*sigma\_max

\end{itemize}


\subsection{Examples}
\label{\detokenize{pair_dsmc:examples}}
\begin{sphinxVerbatim}[commandchars=\\\{\}]
\PYG{k}{pair\PYGZus{}style}\PYG{+w}{ }\PYG{n}{dsmc}\PYG{+w}{ }\PYG{l+m}{2.5}\PYG{+w}{ }\PYG{l+m}{34387}\PYG{+w}{ }\PYG{l+m}{10}\PYG{+w}{ }\PYG{l+m}{1.0}\PYG{+w}{ }\PYG{l+m}{100}\PYG{+w}{ }\PYG{l+m}{20}
\PYG{k}{pair\PYGZus{}coeff}\PYG{+w}{ }\PYG{o}{*}\PYG{+w}{ }\PYG{o}{*}\PYG{+w}{ }\PYG{l+m}{1.0}
\PYG{k}{pair\PYGZus{}coeff}\PYG{+w}{ }\PYG{l+m}{1}\PYG{+w}{ }\PYG{l+m}{1}\PYG{+w}{ }\PYG{l+m}{1.0}
\end{sphinxVerbatim}


\subsection{Description}
\label{\detokenize{pair_dsmc:description}}
\sphinxAtStartPar
Style \sphinxstyleemphasis{dsmc} computes collisions between pairs of particles for a
direct simulation Monte Carlo (DSMC) model following the exposition in
{\hyperref[\detokenize{pair_dsmc:bird}]{\sphinxcrossref{\DUrole{std,std-ref}{(Bird)}}}}.  Each collision resets the velocities of the two
particles involved.  The number of pairwise collisions for each pair
or particle types and the length scale within which they occur are
determined by the parameters of the pair\_style and pair\_coeff
commands.

\sphinxAtStartPar
Stochastic collisions are performed using the variable hard sphere
(VHS) approach, with the user\sphinxhyphen{}defined \sphinxstyleemphasis{max\_cell\_size} value used as
the maximum DSMC cell size, and reference cross\sphinxhyphen{}sections for
collisions given using the pair\_coeff command.

\sphinxAtStartPar
There is no pairwise energy or virial contributions associated with
this pair style.

\sphinxAtStartPar
The following coefficient must be defined for each pair of atoms types
via the {\hyperref[\detokenize{pair_coeff::doc}]{\sphinxcrossref{\DUrole{doc}{pair\_coeff}}}} command as in the examples above,
or in the data file or restart files read by the
{\hyperref[\detokenize{read_data::doc}]{\sphinxcrossref{\DUrole{doc}{read\_data}}}} or {\hyperref[\detokenize{read_restart::doc}]{\sphinxcrossref{\DUrole{doc}{read\_restart}}}}
commands:
\begin{itemize}
\item {} 
\sphinxAtStartPar
sigma (area units, i.e. distance\sphinxhyphen{}squared)

\end{itemize}

\sphinxAtStartPar
The global DSMC \sphinxstyleemphasis{max\_cell\_size} determines the maximum cell length
used in the DSMC calculation.  A structured mesh is overlayed on the
simulation box such that an integer number of cells are created in
each direction for each processor’s subdomain.  Cell lengths are
adjusted up to the user\sphinxhyphen{}specified maximum cell size.


\bigskip\hrule\bigskip


\sphinxAtStartPar
To perform a DSMC simulation with LAMMPS, several additional options
should be set in your input script, though LAMMPS does not check for
these settings.

\sphinxAtStartPar
Since this pair style does not compute particle forces, you should use
the “fix nve/noforce” time integration fix for the DSMC particles,
e.g.

\begin{sphinxVerbatim}[commandchars=\\\{\}]
\PYG{k}{fix }\PYG{n+nv+nvIdentifier}{1}\PYG{+w}{ }\PYG{n+nv+nvIdentifier}{all}\PYG{+w}{ }\PYG{n}{nve}\PYG{o}{/}\PYG{n}{noforce}
\end{sphinxVerbatim}

\sphinxAtStartPar
This pair style assumes that all particles will communicated to
neighboring processors every timestep as they move.  This makes it
possible to perform all collisions between pairs of particles that are
on the same processor.  To ensure this occurs, you should use
these commands:

\begin{sphinxVerbatim}[commandchars=\\\{\}]
\PYG{k}{neighbor}\PYG{+w}{ }\PYG{l+m}{0.0}\PYG{+w}{ }\PYG{n}{bin}
\PYG{k}{neigh\PYGZus{}modify}\PYG{+w}{ }\PYG{n}{every}\PYG{+w}{ }\PYG{l+m}{1}\PYG{+w}{ }\PYG{n}{delay}\PYG{+w}{ }\PYG{l+m}{0}\PYG{+w}{ }\PYG{n}{check}\PYG{+w}{ }\PYG{n}{no}
\PYG{k}{atom\PYGZus{}modify}\PYG{+w}{ }\PYG{n}{sort}\PYG{+w}{ }\PYG{l+m}{0}\PYG{+w}{ }\PYG{l+m}{0.0}
\PYG{n}{communicate}\PYG{+w}{ }\PYG{n}{single}\PYG{+w}{ }\PYG{n}{cutoff}\PYG{+w}{ }\PYG{l+m}{0.0}
\end{sphinxVerbatim}

\sphinxAtStartPar
These commands ensure that LAMMPS communicates particles to
neighboring processors every timestep and that no ghost atoms are
created.  The output statistics for a simulation run should indicate
there are no ghost particles or neighbors.

\sphinxAtStartPar
In order to get correct DSMC collision statistics, users should
specify a Gaussian velocity distribution when populating the
simulation domain. Note that the default velocity distribution is
uniform, which will not give good DSMC collision rates. Specify
“dist gaussian” when using the {\hyperref[\detokenize{velocity::doc}]{\sphinxcrossref{\DUrole{doc}{velocity}}}} command
as in the following:

\begin{sphinxVerbatim}[commandchars=\\\{\}]
\PYG{k}{velocity }\PYG{n+nv+nvIdentifier}{all}\PYG{+w}{ }\PYG{n}{create}\PYG{+w}{ }\PYG{l+m}{594.6}\PYG{+w}{ }\PYG{l+m}{87287}\PYG{+w}{ }\PYG{n}{loop}\PYG{+w}{ }\PYG{n}{geom}\PYG{+w}{ }\PYG{n}{dist}\PYG{+w}{ }\PYG{n}{gaussian}
\end{sphinxVerbatim}


\bigskip\hrule\bigskip



\subsection{Mixing, shift, table, tail correction, restart, rRESPA info}
\label{\detokenize{pair_dsmc:mixing-shift-table-tail-correction-restart-rrespa-info}}
\sphinxAtStartPar
This pair style does not support mixing.  Thus, coefficients for all
I,J pairs must be specified explicitly.

\sphinxAtStartPar
This pair style does not support the {\hyperref[\detokenize{pair_modify::doc}]{\sphinxcrossref{\DUrole{doc}{pair\_modify}}}}
shift option for the energy of the pair interaction.

\sphinxAtStartPar
The {\hyperref[\detokenize{pair_modify::doc}]{\sphinxcrossref{\DUrole{doc}{pair\_modify}}}} table option is not relevant
for this pair style.

\sphinxAtStartPar
This pair style does not support the {\hyperref[\detokenize{pair_modify::doc}]{\sphinxcrossref{\DUrole{doc}{pair\_modify}}}}
tail option for adding long\sphinxhyphen{}range tail corrections to energy and
pressure.

\sphinxAtStartPar
This pair style writes its information to {\hyperref[\detokenize{restart::doc}]{\sphinxcrossref{\DUrole{doc}{binary restart files}}}}, so pair\_style and pair\_coeff commands do not need
to be specified in an input script that reads a restart file.  Note
that the user\sphinxhyphen{}specified random number seed is stored in the restart
file, so when a simulation is restarted, each processor will
re\sphinxhyphen{}initialize its random number generator the same way it did
initially.  This means the random forces will be random, but will not
be the same as they would have been if the original simulation had
continued past the restart time.

\sphinxAtStartPar
This pair style can only be used via the \sphinxstyleemphasis{pair} keyword of the
{\hyperref[\detokenize{run_style::doc}]{\sphinxcrossref{\DUrole{doc}{run\_style respa}}}} command.  It does not support the
\sphinxstyleemphasis{inner}, \sphinxstyleemphasis{middle}, \sphinxstyleemphasis{outer} keywords.


\bigskip\hrule\bigskip



\subsection{Restrictions}
\label{\detokenize{pair_dsmc:restrictions}}
\sphinxAtStartPar
This pair style is part of the MC package.  It is only enabled if LAMMPS
was built with that package.  See the {\hyperref[\detokenize{Build_package::doc}]{\sphinxcrossref{\DUrole{doc}{Build package}}}} page for more info.

\sphinxAtStartPar
This pair style requires an {\hyperref[\detokenize{atom_style::doc}]{\sphinxcrossref{\DUrole{doc}{atom style}}}} with per
atom type masses.


\subsection{Related commands}
\label{\detokenize{pair_dsmc:related-commands}}
\sphinxAtStartPar
{\hyperref[\detokenize{pair_coeff::doc}]{\sphinxcrossref{\DUrole{doc}{pair\_coeff}}}}, {\hyperref[\detokenize{fix_nve_noforce::doc}]{\sphinxcrossref{\DUrole{doc}{fix nve/noforce}}}},
{\hyperref[\detokenize{neigh_modify::doc}]{\sphinxcrossref{\DUrole{doc}{neigh\_modify}}}}, {\hyperref[\detokenize{neighbor::doc}]{\sphinxcrossref{\DUrole{doc}{neighbor}}}},
{\hyperref[\detokenize{comm_modify::doc}]{\sphinxcrossref{\DUrole{doc}{comm\_modify}}}}


\subsection{Default}
\label{\detokenize{pair_dsmc:default}}
\sphinxAtStartPar
none


\bigskip\hrule\bigskip

\phantomsection\label{\detokenize{pair_dsmc:bird}}
\sphinxAtStartPar
\sphinxstylestrong{(Bird)} G. A. Bird, “Molecular Gas Dynamics and the Direct Simulation
of Gas Flows” (1994).

\sphinxstepscope

\index{pair\_style e3b@\spxentry{pair\_style e3b}}\ignorespaces 

\section{pair\_style e3b command}
\label{\detokenize{pair_e3b:pair-style-e3b-command}}\label{\detokenize{pair_e3b:index-0}}\label{\detokenize{pair_e3b::doc}}

\subsection{Syntax}
\label{\detokenize{pair_e3b:syntax}}
\begin{sphinxVerbatim}[commandchars=\\\{\}]
\PYG{k}{pair\PYGZus{}style}\PYG{+w}{ }\PYG{n}{e3b}\PYG{+w}{ }\PYG{n}{Otype}
\end{sphinxVerbatim}
\begin{itemize}
\item {} 
\sphinxAtStartPar
Otype = atom type (numeric or type label) for oxygen

\end{itemize}

\begin{sphinxVerbatim}[commandchars=\\\{\}]
\PYG{k}{pair\PYGZus{}coeff}\PYG{+w}{ }\PYG{o}{*}\PYG{+w}{ }\PYG{o}{*}\PYG{+w}{ }\PYG{n}{keyword}
\end{sphinxVerbatim}
\begin{itemize}
\item {} 
\sphinxAtStartPar
one or more keyword/value pairs must be appended.

\item {} 
\sphinxAtStartPar
keyword = \sphinxstyleemphasis{preset} or \sphinxstyleemphasis{Ea} or \sphinxstyleemphasis{Eb} or \sphinxstyleemphasis{Ec} or \sphinxstyleemphasis{E2} or \sphinxstyleemphasis{K3} or \sphinxstyleemphasis{K2} or \sphinxstyleemphasis{Rs} or \sphinxstyleemphasis{Rc3} or \sphinxstyleemphasis{Rc2} or \sphinxstyleemphasis{bondL} or \sphinxstyleemphasis{neigh}

\item {} 
\sphinxAtStartPar
If the \sphinxstyleemphasis{preset} keyword is given, no others are needed.
Otherwise, all are mandatory except for \sphinxstyleemphasis{neigh}.
The \sphinxstyleemphasis{neigh} keyword is always optional.

\end{itemize}
\begin{sphinxalltt}
\sphinxstyleemphasis{preset} arg = \sphinxstyleemphasis{2011} or \sphinxstyleemphasis{2015} = which set of predefined parameters to use
         2011 = use the potential parameters from {\hyperref[\detokenize{pair_e3b:tainter2011}]{\sphinxcrossref{\DUrole{std,std-ref}{(Tainter 2011)}}}}
         2015 = use the potential parameters from {\hyperref[\detokenize{pair_e3b:tainter2015}]{\sphinxcrossref{\DUrole{std,std-ref}{(Tainter 2015)}}}}
\sphinxstyleemphasis{Ea} arg = three\sphinxhyphen{}body energy for type A hydrogen bonding interactions (energy units)
\sphinxstyleemphasis{Eb} arg = three\sphinxhyphen{}body energy for type B hydrogen bonding interactions (energy units)
\sphinxstyleemphasis{Ec} arg = three\sphinxhyphen{}body energy for type C hydrogen bonding interactions (energy units)
\sphinxstyleemphasis{E2} arg = two\sphinxhyphen{}body energy correction (energy units)
\sphinxstyleemphasis{K3} arg = three\sphinxhyphen{}body exponential constant (inverse distance units)
\sphinxstyleemphasis{K2} arg = two\sphinxhyphen{}body exponential constant (inverse distance units)
\sphinxstyleemphasis{Rc3} arg = three\sphinxhyphen{}body cutoff (distance units)
\sphinxstyleemphasis{Rc2} arg = two\sphinxhyphen{}body cutoff (distance units)
\sphinxstyleemphasis{Rs} arg = three\sphinxhyphen{}body switching function cutoff (distance units)
\sphinxstyleemphasis{bondL} arg = intramolecular OH bond length (distance units)
\sphinxstyleemphasis{neigh} arg = approximate integer number of molecules within Rc3 of an oxygen atom
\end{sphinxalltt}


\subsection{Examples}
\label{\detokenize{pair_e3b:examples}}
\begin{sphinxVerbatim}[commandchars=\\\{\}]
\PYG{k}{pair\PYGZus{}style}\PYG{+w}{ }\PYG{n}{e3b}\PYG{+w}{ }\PYG{l+m}{1}
\PYG{k}{pair\PYGZus{}coeff}\PYG{+w}{ }\PYG{o}{*}\PYG{+w}{ }\PYG{o}{*}\PYG{+w}{ }\PYG{n}{Ea}\PYG{+w}{ }\PYG{l+m}{35.85}\PYG{+w}{ }\PYG{n}{Eb}\PYG{+w}{ }\PYG{o}{\PYGZhy{}}\PYG{l+m}{240.2}\PYG{+w}{ }\PYG{n}{Ec}\PYG{+w}{ }\PYG{l+m}{449.3}\PYG{+w}{ }\PYG{n}{E2}\PYG{+w}{ }\PYG{l+m}{108269.9}\PYG{+w}{ }\PYG{n}{K3}\PYG{+w}{ }\PYG{l+m}{1.907}\PYG{+w}{ }\PYG{n}{K2}\PYG{+w}{ }\PYG{l+m}{4.872}\PYG{+w}{ }\PYG{n}{Rc3}\PYG{+w}{ }\PYG{l+m}{5.2}\PYG{+w}{ }\PYG{n}{Rc2}\PYG{+w}{ }\PYG{l+m}{5.2}\PYG{+w}{ }\PYG{n}{Rs}\PYG{+w}{ }\PYG{l+m}{5.0}\PYG{+w}{ }\PYG{n}{bondL}\PYG{+w}{ }\PYG{l+m}{0.9572}

\PYG{k}{pair\PYGZus{}style}\PYG{+w}{ }\PYG{n}{hybrid}\PYG{o}{/}\PYG{n}{overlay}\PYG{+w}{ }\PYG{n}{e3b}\PYG{+w}{ }\PYG{l+m}{1}\PYG{+w}{ }\PYG{n}{lj}\PYG{o}{/}\PYG{n}{cut}\PYG{o}{/}\PYG{n}{tip4p}\PYG{o}{/}\PYG{n}{long}\PYG{+w}{ }\PYG{l+m}{1}\PYG{+w}{ }\PYG{l+m}{2}\PYG{+w}{ }\PYG{l+m}{1}\PYG{+w}{ }\PYG{l+m}{1}\PYG{+w}{ }\PYG{l+m}{0.15}\PYG{+w}{ }\PYG{l+m}{8.5}
\PYG{k}{pair\PYGZus{}coeff}\PYG{+w}{ }\PYG{o}{*}\PYG{+w}{ }\PYG{o}{*}\PYG{+w}{ }\PYG{n}{e3b}\PYG{+w}{ }\PYG{n}{preset}\PYG{+w}{ }\PYG{l+m}{2011}

\PYG{k}{pair\PYGZus{}style}\PYG{+w}{ }\PYG{n}{e3b}\PYG{+w}{ }\PYG{n}{OW}
\PYG{k}{labelmap}\PYG{+w}{ }\PYG{n}{atom}\PYG{+w}{ }\PYG{l+m}{1}\PYG{+w}{ }\PYG{n}{C}\PYG{+w}{ }\PYG{l+m}{2}\PYG{+w}{ }\PYG{n}{H}\PYG{+w}{ }\PYG{l+m}{3}\PYG{+w}{ }\PYG{n}{O}\PYG{+w}{ }\PYG{l+m}{4}\PYG{+w}{ }\PYG{n}{N}\PYG{+w}{ }\PYG{l+m}{5}\PYG{+w}{ }\PYG{n}{OW}\PYG{+w}{ }\PYG{l+m}{6}\PYG{+w}{ }\PYG{n}{HW}
\PYG{k}{pair\PYGZus{}coeff}\PYG{+w}{ }\PYG{o}{*}\PYG{+w}{ }\PYG{o}{*}\PYG{+w}{ }\PYG{n}{Ea}\PYG{+w}{ }\PYG{l+m}{35.85}\PYG{+w}{ }\PYG{n}{Eb}\PYG{+w}{ }\PYG{o}{\PYGZhy{}}\PYG{l+m}{240.2}\PYG{+w}{ }\PYG{n}{Ec}\PYG{+w}{ }\PYG{l+m}{449.3}\PYG{+w}{ }\PYG{n}{E2}\PYG{+w}{ }\PYG{l+m}{108269.9}\PYG{+w}{ }\PYG{n}{K3}\PYG{+w}{ }\PYG{l+m}{1.907}\PYG{+w}{ }\PYG{n}{K2}\PYG{+w}{ }\PYG{l+m}{4.872}\PYG{+w}{ }\PYG{n}{Rc3}\PYG{+w}{ }\PYG{l+m}{5.2}\PYG{+w}{ }\PYG{n}{Rc2}\PYG{+w}{ }\PYG{l+m}{5.2}\PYG{+w}{ }\PYG{n}{Rs}\PYG{+w}{ }\PYG{l+m}{5.0}\PYG{+w}{ }\PYG{n}{bondL}\PYG{+w}{ }\PYG{l+m}{0.9572}
\end{sphinxVerbatim}

\sphinxAtStartPar
Used in example input script:

\begin{sphinxVerbatim}[commandchars=\\\{\}]
examples/PACKAGES/e3b/in.e3b\PYGZhy{}tip4p2005
\end{sphinxVerbatim}


\subsection{Description}
\label{\detokenize{pair_e3b:description}}
\sphinxAtStartPar
The \sphinxstyleemphasis{e3b} style computes an "explicit three\sphinxhyphen{}body" (E3B) potential for water {\hyperref[\detokenize{pair_e3b:kumar}]{\sphinxcrossref{\DUrole{std,std-ref}{(Kumar 2008)}}}}.
\begin{equation*}
\begin{split}E =& E_2 \sum_{i,j}e^{-k_2 r_{ij}} + E_A \sum_{\substack{i,j,k,\ell \\
\in \textrm{type A}}} f(r_{ij})f(r_{k\ell}) + E_B \sum_{\substack{i,j,k,\ell \\
\in \textrm{type B}}} f(r_{ij})f(r_{k\ell}) + E_C \sum_{\substack{i,j,k,\ell \\
\in \textrm{type C}}} f(r_{ij})f(r_{k\ell}) \\
f(r) =& e^{-k_3 r}s(r) \\
s(r) =& \begin{cases}
1 & r<R_s \\
\displaystyle\frac{(R_f-r)^2(R_f-3R_s+2r)}{(R_f-R_s)^3} & R_s\leq r\leq R_f \\
0 & r>R_f\\
\end{cases}\end{split}
\end{equation*}
\sphinxAtStartPar
This potential was developed as a water model that includes the
three\sphinxhyphen{}body cooperativity of hydrogen bonding explicitly.  To use it in
this way, it must be applied in conjunction with a conventional two\sphinxhyphen{}body
water model, through pair style {\hyperref[\detokenize{pair_hybrid::doc}]{\sphinxcrossref{\DUrole{doc}{hybrid/overlay}}}}.  The
three body interactions are split into three types: A, B, and C.  Type A
corresponds to anti\sphinxhyphen{}cooperative double hydrogen bond donor interactions.
Type B corresponds to the cooperative interaction of molecules that both
donate and accept a hydrogen bond.  Type C corresponds to
anti\sphinxhyphen{}cooperative double hydrogen bond acceptor interactions.  The
three\sphinxhyphen{}body interactions are smoothly cutoff by the switching function
s(r) between Rs and Rc3.  The two\sphinxhyphen{}body interactions are designed to
correct for the effective many\sphinxhyphen{}body interactions implicitly included in
the conventional two\sphinxhyphen{}body potential.  The two\sphinxhyphen{}body interactions are cut
off sharply at Rc2, because K3 is typically significantly smaller than
K2.  See {\hyperref[\detokenize{pair_e3b:kumar}]{\sphinxcrossref{\DUrole{std,std-ref}{(Kumar 2008)}}}} for more details.

\sphinxAtStartPar
Only a single {\hyperref[\detokenize{pair_coeff::doc}]{\sphinxcrossref{\DUrole{doc}{pair\_coeff}}}} command is used with the
\sphinxstyleemphasis{e3b} style and the first two arguments must be * *.  The oxygen atom
type for the pair style is passed as the only argument to the
\sphinxstyleemphasis{pair\_style} command, not in the \sphinxstyleemphasis{pair\_coeff} command.  The hydrogen
atom type is inferred from the ordering of the atoms.

\begin{sphinxadmonition}{note}{Note:}
\sphinxAtStartPar
Every atom of type Otype must be part of a water molecule.
Each water molecule must have consecutive IDs with the oxygen first.
This pair style does not test that this criteria is met.
\end{sphinxadmonition}

\begin{sphinxadmonition}{note}{Note:}
\sphinxAtStartPar
If using type labels, the type labels must be defined before calling
the {\hyperref[\detokenize{pair_coeff::doc}]{\sphinxcrossref{\DUrole{doc}{pair\_coeff}}}} command.
\end{sphinxadmonition}

\sphinxAtStartPar
The \sphinxstyleemphasis{pair\_coeff} command must have at least one keyword/value pair, as
described above.  The \sphinxstyleemphasis{preset} keyword sets the potential parameters to
the values used in {\hyperref[\detokenize{pair_e3b:tainter2011}]{\sphinxcrossref{\DUrole{std,std-ref}{(Tainter 2011)}}}} or
{\hyperref[\detokenize{pair_e3b:tainter2015}]{\sphinxcrossref{\DUrole{std,std-ref}{(Tainter 2015)}}}}.  To use the water models defined in
those references, the \sphinxstyleemphasis{e3b} style should always be used in conjunction
with an \sphinxstyleemphasis{lj/cut/tip4p/long} style through \sphinxstyleemphasis{pair\_style hybrid/overlay},
as demonstrated in the second example above.  The \sphinxstyleemphasis{preset 2011} option
should be used with the {\hyperref[\detokenize{Howto_tip4p::doc}]{\sphinxcrossref{\DUrole{doc}{TIP4P water model}}}}.  The
\sphinxstyleemphasis{preset 2015} option should be used with the {\hyperref[\detokenize{Howto_tip4p::doc}]{\sphinxcrossref{\DUrole{doc}{TIP4P/2005 water
model}}}}.  If the \sphinxstyleemphasis{preset} keyword is used, no other keyword
is needed.  Changes to the preset parameters can be made by specifying
the \sphinxstyleemphasis{preset} keyword followed by the specific parameter to change, like
\sphinxstyleemphasis{Ea}.  Note that the other keywords must come after \sphinxstyleemphasis{preset} in the
pair\_style command.  The \sphinxstyleemphasis{e3b} style can also be used to implement any
three\sphinxhyphen{}body potential of the same form by specifying all the keywords
except \sphinxstyleemphasis{neigh}: \sphinxstyleemphasis{Ea}, \sphinxstyleemphasis{Eb}, \sphinxstyleemphasis{Ec}, \sphinxstyleemphasis{E2}, \sphinxstyleemphasis{K3}, \sphinxstyleemphasis{K2}, \sphinxstyleemphasis{Rc3}, \sphinxstyleemphasis{Rc2},
\sphinxstyleemphasis{Rs}, and \sphinxstyleemphasis{bondL}.  The keyword \sphinxstyleemphasis{bondL} specifies the intramolecular
OH bond length of the water model being used.  This is needed to include
H atoms that are within the cutoff even when the attached oxygen atom is
not.

\sphinxAtStartPar
This pair style allocates arrays sized according to the number of
pairwise interactions within Rc3.  To do this it needs an estimate for
the number of water molecules within Rc3 of an oxygen atom.  This
estimate defaults to 10 and can be changed using the \sphinxstyleemphasis{neigh} keyword,
which takes an integer as an argument.  If the neigh setting is too
small, the simulation will fail with the error “neigh is too small”.  If
the neigh setting is too large, the pair style will use more memory than
necessary.

\sphinxAtStartPar
This pair style tallies a breakdown of the total E3B potential energy
into sub\sphinxhyphen{}categories, which can be accessed via the {\hyperref[\detokenize{compute_pair::doc}]{\sphinxcrossref{\DUrole{doc}{compute pair}}}} command as a vector of values of length 4.  The 4 values
correspond to the terms in the first equation above: the E2 term, the Ea
term, the Eb term, and the Ec term.

\sphinxAtStartPar
See the examples/PACKAGES/e3b directory for a complete example script.


\bigskip\hrule\bigskip



\subsection{Mixing, shift, table, tail correction, restart, rRESPA info}
\label{\detokenize{pair_e3b:mixing-shift-table-tail-correction-restart-rrespa-info}}
\sphinxAtStartPar
This pair style does not support the {\hyperref[\detokenize{pair_modify::doc}]{\sphinxcrossref{\DUrole{doc}{pair\_modify}}}}
shift, table, and tail options.

\sphinxAtStartPar
This pair style does not write its information to {\hyperref[\detokenize{restart::doc}]{\sphinxcrossref{\DUrole{doc}{binary restart files}}}}.  Thus, you
need to re\sphinxhyphen{}specify the pair\_style and pair\_coeff commands in an input
script that reads a restart file.

\sphinxAtStartPar
This pair style is incompatible with {\hyperref[\detokenize{run_style::doc}]{\sphinxcrossref{\DUrole{doc}{respa}}}}.


\bigskip\hrule\bigskip



\subsection{Restrictions}
\label{\detokenize{pair_e3b:restrictions}}
\sphinxAtStartPar
This pair style is part of the EXTRA\sphinxhyphen{}PAIR package.  It is only enabled
if LAMMPS was built with that package.  See the {\hyperref[\detokenize{Build_package::doc}]{\sphinxcrossref{\DUrole{doc}{Build package}}}} page for more info.

\sphinxAtStartPar
This pair style requires the {\hyperref[\detokenize{newton::doc}]{\sphinxcrossref{\DUrole{doc}{newton}}}} setting to be “on”
for pair interactions.

\sphinxAtStartPar
This pair style requires a fixed number of atoms in the simulation, so it is incompatible with fixes like {\hyperref[\detokenize{fix_deposit::doc}]{\sphinxcrossref{\DUrole{doc}{fix deposit}}}}.
If the number of atoms changes between runs, this pair style must be re\sphinxhyphen{}initialized by calling the \sphinxstyleemphasis{pair\_style} and \sphinxstyleemphasis{pair\_coeffs} commands.
This is not a fundamental limitation of the pair style, but the code currently does not support a variable number of atoms.

\sphinxAtStartPar
The \sphinxstyleemphasis{preset} keyword currently only works with real, metal, si, and cgs {\hyperref[\detokenize{units::doc}]{\sphinxcrossref{\DUrole{doc}{units}}}}.


\subsection{Related commands}
\label{\detokenize{pair_e3b:related-commands}}
\sphinxAtStartPar
{\hyperref[\detokenize{pair_coeff::doc}]{\sphinxcrossref{\DUrole{doc}{pair\_coeff}}}}, {\hyperref[\detokenize{compute_pair::doc}]{\sphinxcrossref{\DUrole{doc}{compute pair}}}}


\subsection{Default}
\label{\detokenize{pair_e3b:default}}
\sphinxAtStartPar
The option default for the \sphinxstyleemphasis{neigh} keyword is 10.


\bigskip\hrule\bigskip

\phantomsection\label{\detokenize{pair_e3b:kumar}}\phantomsection\label{\detokenize{pair_e3b:tainter2011}}
\sphinxAtStartPar
\sphinxstylestrong{(Kumar)} Kumar and Skinner, J. Phys. Chem. B, 112, 8311 (2008)

\phantomsection\label{\detokenize{pair_e3b:tainter2015}}
\sphinxAtStartPar
\sphinxstylestrong{(Tainter 2011)} Tainter, Pieniazek, Lin, and Skinner, J. Chem. Phys., 134, 184501 (2011)

\sphinxAtStartPar
\sphinxstylestrong{(Tainter 2015)} Tainter, Shi, and Skinner, 11, 2268 (2015)

\sphinxstepscope

\index{pair\_style eam@\spxentry{pair\_style eam}}\ignorespaces 
\index{pair\_style eam/gpu@\spxentry{pair\_style eam/gpu}}\ignorespaces 
\index{pair\_style eam/intel@\spxentry{pair\_style eam/intel}}\ignorespaces 
\index{pair\_style eam/kk@\spxentry{pair\_style eam/kk}}\ignorespaces 
\index{pair\_style eam/omp@\spxentry{pair\_style eam/omp}}\ignorespaces 
\index{pair\_style eam/opt@\spxentry{pair\_style eam/opt}}\ignorespaces 
\index{pair\_style eam/alloy@\spxentry{pair\_style eam/alloy}}\ignorespaces 
\index{pair\_style eam/alloy/gpu@\spxentry{pair\_style eam/alloy/gpu}}\ignorespaces 
\index{pair\_style eam/alloy/intel@\spxentry{pair\_style eam/alloy/intel}}\ignorespaces 
\index{pair\_style eam/alloy/kk@\spxentry{pair\_style eam/alloy/kk}}\ignorespaces 
\index{pair\_style eam/alloy/omp@\spxentry{pair\_style eam/alloy/omp}}\ignorespaces 
\index{pair\_style eam/alloy/opt@\spxentry{pair\_style eam/alloy/opt}}\ignorespaces 
\index{pair\_style eam/cd@\spxentry{pair\_style eam/cd}}\ignorespaces 
\index{pair\_style eam/cd/old@\spxentry{pair\_style eam/cd/old}}\ignorespaces 
\index{pair\_style eam/fs@\spxentry{pair\_style eam/fs}}\ignorespaces 
\index{pair\_style eam/fs/gpu@\spxentry{pair\_style eam/fs/gpu}}\ignorespaces 
\index{pair\_style eam/fs/intel@\spxentry{pair\_style eam/fs/intel}}\ignorespaces 
\index{pair\_style eam/fs/kk@\spxentry{pair\_style eam/fs/kk}}\ignorespaces 
\index{pair\_style eam/fs/omp@\spxentry{pair\_style eam/fs/omp}}\ignorespaces 
\index{pair\_style eam/fs/opt@\spxentry{pair\_style eam/fs/opt}}\ignorespaces 
\index{pair\_style eam/he@\spxentry{pair\_style eam/he}}\ignorespaces 

\section{pair\_style eam command}
\label{\detokenize{pair_eam:pair-style-eam-command}}\label{\detokenize{pair_eam:index-20}}\label{\detokenize{pair_eam:index-19}}\label{\detokenize{pair_eam:index-18}}\label{\detokenize{pair_eam:index-17}}\label{\detokenize{pair_eam:index-16}}\label{\detokenize{pair_eam:index-15}}\label{\detokenize{pair_eam:index-14}}\label{\detokenize{pair_eam:index-13}}\label{\detokenize{pair_eam:index-12}}\label{\detokenize{pair_eam:index-11}}\label{\detokenize{pair_eam:index-10}}\label{\detokenize{pair_eam:index-9}}\label{\detokenize{pair_eam:index-8}}\label{\detokenize{pair_eam:index-7}}\label{\detokenize{pair_eam:index-6}}\label{\detokenize{pair_eam:index-5}}\label{\detokenize{pair_eam:index-4}}\label{\detokenize{pair_eam:index-3}}\label{\detokenize{pair_eam:index-2}}\label{\detokenize{pair_eam:index-1}}\label{\detokenize{pair_eam:index-0}}\label{\detokenize{pair_eam::doc}}
\sphinxAtStartPar
Accelerator Variants: \sphinxstyleemphasis{eam/gpu}, \sphinxstyleemphasis{eam/intel}, \sphinxstyleemphasis{eam/kk}, \sphinxstyleemphasis{eam/omp}, \sphinxstyleemphasis{eam/opt}


\section{pair\_style eam/alloy command}
\label{\detokenize{pair_eam:pair-style-eam-alloy-command}}
\sphinxAtStartPar
Accelerator Variants: \sphinxstyleemphasis{eam/alloy/gpu}, \sphinxstyleemphasis{eam/alloy/intel}, \sphinxstyleemphasis{eam/alloy/kk}, \sphinxstyleemphasis{eam/alloy/omp}, \sphinxstyleemphasis{eam/alloy/opt}


\section{pair\_style eam/cd command}
\label{\detokenize{pair_eam:pair-style-eam-cd-command}}

\section{pair\_style eam/cd/old command}
\label{\detokenize{pair_eam:pair-style-eam-cd-old-command}}

\section{pair\_style eam/fs command}
\label{\detokenize{pair_eam:pair-style-eam-fs-command}}

\section{pair\_style eam/he command}
\label{\detokenize{pair_eam:pair-style-eam-he-command}}
\sphinxAtStartPar
Accelerator Variants: \sphinxstyleemphasis{eam/fs/gpu}, \sphinxstyleemphasis{eam/fs/intel}, \sphinxstyleemphasis{eam/fs/kk}, \sphinxstyleemphasis{eam/fs/omp}, \sphinxstyleemphasis{eam/fs/opt}


\subsection{Syntax}
\label{\detokenize{pair_eam:syntax}}
\begin{sphinxVerbatim}[commandchars=\\\{\}]
\PYG{k}{pair\PYGZus{}style}\PYG{+w}{ }\PYG{n}{style}
\end{sphinxVerbatim}
\begin{itemize}
\item {} 
\sphinxAtStartPar
style = \sphinxstyleemphasis{eam} or \sphinxstyleemphasis{eam/alloy} or \sphinxstyleemphasis{eam/cd} or \sphinxstyleemphasis{eam/cd/old} or \sphinxstyleemphasis{eam/fs} or \sphinxstyleemphasis{eam/he}

\end{itemize}


\subsection{Examples}
\label{\detokenize{pair_eam:examples}}
\begin{sphinxVerbatim}[commandchars=\\\{\}]
\PYG{k}{pair\PYGZus{}style}\PYG{+w}{ }\PYG{n}{eam}
\PYG{k}{pair\PYGZus{}coeff}\PYG{+w}{ }\PYG{o}{*}\PYG{+w}{ }\PYG{o}{*}\PYG{+w}{ }\PYG{n}{cuu3}
\PYG{k}{pair\PYGZus{}coeff}\PYG{+w}{ }\PYG{l+m}{1}\PYG{o}{*}\PYG{l+m}{3}\PYG{+w}{ }\PYG{l+m}{1}\PYG{o}{*}\PYG{l+m}{3}\PYG{+w}{ }\PYG{n}{niu3.eam}

\PYG{k}{pair\PYGZus{}style}\PYG{+w}{ }\PYG{n}{eam}\PYG{o}{/}\PYG{n}{alloy}
\PYG{k}{pair\PYGZus{}coeff}\PYG{+w}{ }\PYG{o}{*}\PYG{+w}{ }\PYG{o}{*}\PYG{+w}{ }\PYG{n}{..}\PYG{o}{/}\PYG{n}{potentials}\PYG{o}{/}\PYG{n}{NiAlH\PYGZus{}jea.eam.alloy}\PYG{+w}{ }\PYG{n}{Ni}\PYG{+w}{ }\PYG{n}{Al}\PYG{+w}{ }\PYG{n}{Ni}\PYG{+w}{ }\PYG{n}{Ni}

\PYG{k}{pair\PYGZus{}style}\PYG{+w}{ }\PYG{n}{eam}\PYG{o}{/}\PYG{n}{cd}
\PYG{k}{pair\PYGZus{}coeff}\PYG{+w}{ }\PYG{o}{*}\PYG{+w}{ }\PYG{o}{*}\PYG{+w}{ }\PYG{n}{..}\PYG{o}{/}\PYG{n}{potentials}\PYG{o}{/}\PYG{n}{FeCr.cdeam}\PYG{+w}{ }\PYG{n}{Fe}\PYG{+w}{ }\PYG{n}{Cr}

\PYG{k}{pair\PYGZus{}style}\PYG{+w}{ }\PYG{n}{eam}\PYG{o}{/}\PYG{n}{fs}
\PYG{k}{pair\PYGZus{}coeff}\PYG{+w}{ }\PYG{o}{*}\PYG{+w}{ }\PYG{o}{*}\PYG{+w}{ }\PYG{n}{NiAlH\PYGZus{}jea.eam.fs}\PYG{+w}{ }\PYG{n}{Ni}\PYG{+w}{ }\PYG{n}{Al}\PYG{+w}{ }\PYG{n}{Ni}\PYG{+w}{ }\PYG{n}{Ni}

\PYG{k}{pair\PYGZus{}style}\PYG{+w}{ }\PYG{n}{eam}\PYG{o}{/}\PYG{n}{he}
\PYG{k}{pair\PYGZus{}coeff}\PYG{+w}{ }\PYG{o}{*}\PYG{+w}{ }\PYG{o}{*}\PYG{+w}{ }\PYG{n}{PdHHe.eam.he}\PYG{+w}{ }\PYG{n}{Pd}\PYG{+w}{ }\PYG{n}{H}\PYG{+w}{ }\PYG{n}{He}
\end{sphinxVerbatim}


\subsection{Description}
\label{\detokenize{pair_eam:description}}
\sphinxAtStartPar
Style \sphinxstyleemphasis{eam} computes pairwise interactions for metals and metal alloys
using embedded\sphinxhyphen{}atom method (EAM) potentials {\hyperref[\detokenize{pair_eam:daw}]{\sphinxcrossref{\DUrole{std,std-ref}{(Daw)}}}}.  The total
energy Ei of an atom I is given by
\begin{equation*}
\begin{split}E_i = F_\alpha \left(\sum_{j \neq i}\ \rho_\beta (r_{ij})\right) +
      \frac{1}{2} \sum_{j \neq i} \phi_{\alpha\beta} (r_{ij})\end{split}
\end{equation*}
\sphinxAtStartPar
where F is the embedding energy which is a function of the atomic
electron density rho, phi is a pair potential interaction, and alpha
and beta are the element types of atoms I and J.  The multi\sphinxhyphen{}body
nature of the EAM potential is a result of the embedding energy term.
Both summations in the formula are over all neighbors J of atom I
within the cutoff distance.

\sphinxAtStartPar
The cutoff distance and the tabulated values of the functionals F,
rho, and phi are listed in one or more files which are specified by
the {\hyperref[\detokenize{pair_coeff::doc}]{\sphinxcrossref{\DUrole{doc}{pair\_coeff}}}} command.  These are ASCII text files
in a DYNAMO\sphinxhyphen{}style format which is described below.  DYNAMO was the
original serial EAM MD code, written by the EAM originators.  Several
DYNAMO potential files for different metals are included in the
“potentials” directory of the LAMMPS distribution.  All of these files
are parameterized in terms of LAMMPS {\hyperref[\detokenize{units::doc}]{\sphinxcrossref{\DUrole{doc}{metal units}}}}.

\begin{sphinxadmonition}{note}{Note:}
\sphinxAtStartPar
The \sphinxstyleemphasis{eam} style reads single\sphinxhyphen{}element EAM potentials in the
DYNAMO \sphinxstyleemphasis{funcfl} format.  Either single element or alloy systems can be
modeled using multiple \sphinxstyleemphasis{funcfl} files and style \sphinxstyleemphasis{eam}.  For the alloy
case LAMMPS mixes the single\sphinxhyphen{}element potentials to produce alloy
potentials, the same way that DYNAMO does.  Alternatively, a single
DYNAMO \sphinxstyleemphasis{setfl} file or Finnis/Sinclair EAM file can be used by LAMMPS
to model alloy systems by invoking the \sphinxstyleemphasis{eam/alloy} or \sphinxstyleemphasis{eam/cd} or
\sphinxstyleemphasis{eam/fs} or \sphinxstyleemphasis{eam/he} styles as described below.  These files require no
mixing since they specify alloy interactions explicitly.
\end{sphinxadmonition}

\begin{sphinxadmonition}{note}{Note:}
\sphinxAtStartPar
Note that unlike for other potentials, cutoffs for EAM potentials are not
set in the pair\_style or pair\_coeff command; they are specified in the EAM
potential files themselves.  Likewise, valid EAM potential files usually
contain atomic masses; thus you may not need to use the {\hyperref[\detokenize{mass::doc}]{\sphinxcrossref{\DUrole{doc}{mass}}}}
command to specify them, unless the potential file uses a dummy value (e.g.
0.0). LAMMPS will print a warning, if this is the case.
\end{sphinxadmonition}

\sphinxAtStartPar
There are web sites that distribute and document EAM potentials stored
in DYNAMO or other formats:
\begin{itemize}
\item {} 
\sphinxAtStartPar
\sphinxurl{https://www.ctcms.nist.gov/potentials}

\item {} 
\sphinxAtStartPar
\sphinxurl{https://openkim.org}

\end{itemize}

\sphinxAtStartPar
These potentials should be usable with LAMMPS, though the alternate
formats would need to be converted to the DYNAMO format used by LAMMPS
and described on this page.  The NIST site is maintained by Chandler
Becker (cbecker at nist.gov) who is good resource for info on
interatomic potentials and file formats.

\sphinxAtStartPar
The OpenKIM Project at
\sphinxurl{https://openkim.org/browse/models/by-type}
provides EAM potentials that can be used directly in LAMMPS with the
{\hyperref[\detokenize{kim_commands::doc}]{\sphinxcrossref{\DUrole{doc}{kim command}}}} interface.

\begin{sphinxadmonition}{warning}{Warning:}
\sphinxAtStartPar
The EAM potential files tabulate the embedding energy as a function
of the local electron density \(\rho\).  When atoms get too
close, this electron density may exceed the range for which the
embedding energy was tabulated for.  To avoid crashes, LAMMPS will
assume a linearly increasing embedding energy for electron densities
beyond the maximum tabulated value.  LAMMPS will print a warning when
this happens.  It may be acceptable at the beginning of an
equilibration (e.g. when using randomized coordinates) but would be a
big concern for accuracy if it happens during production runs.  The
EAM potential file triggering the warning during production is thus
not a good choice, and the EAM model in general not likely a good
model for the kind of system under investigation.
\end{sphinxadmonition}


\bigskip\hrule\bigskip


\sphinxAtStartPar
For style \sphinxstyleemphasis{eam}, potential values are read from a file that is in the
DYNAMO single\sphinxhyphen{}element \sphinxstyleemphasis{funcfl} format.  If the DYNAMO file was created
by a Fortran program, it cannot have “D” values in it for exponents.
C only recognizes “e” or “E” for scientific notation.

\sphinxAtStartPar
For style \sphinxstyleemphasis{eam} a potential file must be assigned to each I,I pair of
atom types by using one or more pair\_coeff commands, each with a
single argument:
\begin{itemize}
\item {} 
\sphinxAtStartPar
filename

\end{itemize}

\sphinxAtStartPar
Thus the following command

\begin{sphinxVerbatim}[commandchars=\\\{\}]
\PYG{k}{pair\PYGZus{}coeff}\PYG{+w}{ }\PYG{o}{*}\PYG{l+m}{2}\PYG{+w}{ }\PYG{l+m}{1}\PYG{o}{*}\PYG{l+m}{2}\PYG{+w}{ }\PYG{n}{cuu3.eam}
\end{sphinxVerbatim}

\sphinxAtStartPar
will read the cuu3 potential file and use the tabulated Cu values for
F, phi, rho that it contains for type pairs 1,1 and 2,2 (type pairs
1,2 and 2,1 are ignored).  See the {\hyperref[\detokenize{pair_coeff::doc}]{\sphinxcrossref{\DUrole{doc}{pair\_coeff}}}} doc
page for alternate ways to specify the path for the potential file.
In effect, this makes atom types 1 and 2 in LAMMPS be Cu atoms.
Different single\sphinxhyphen{}element files can be assigned to different atom types
to model an alloy system.  The mixing to create alloy potentials for
type pairs with I != J is done automatically the same way that the
serial DYNAMO code originally did it; you do not need to specify
coefficients for these type pairs.

\sphinxAtStartPar
\sphinxstyleemphasis{Funcfl} files in the \sphinxstyleemphasis{potentials} directory of the LAMMPS
distribution have an “.eam” suffix.  A DYNAMO single\sphinxhyphen{}element \sphinxstyleemphasis{funcfl}
file is formatted as follows:
\begin{itemize}
\item {} 
\sphinxAtStartPar
line 1: comment (ignored)

\item {} 
\sphinxAtStartPar
line 2: atomic number, mass, lattice constant, lattice type (e.g. FCC)

\item {} 
\sphinxAtStartPar
line 3: Nrho, drho, Nr, dr, cutoff

\end{itemize}

\sphinxAtStartPar
On line 2, all values but the mass are ignored by LAMMPS.  The mass is
in mass {\hyperref[\detokenize{units::doc}]{\sphinxcrossref{\DUrole{doc}{units}}}}, e.g. mass number or grams/mole for metal
units.  The cubic lattice constant is in Angstroms.  On line 3, Nrho
and Nr are the number of tabulated values in the subsequent arrays,
drho and dr are the spacing in density and distance space for the
values in those arrays, and the specified cutoff becomes the pairwise
cutoff used by LAMMPS for the potential.  The units of dr are
Angstroms; I’m not sure of the units for drho \sphinxhyphen{} some measure of
electron density.

\sphinxAtStartPar
Following the three header lines are three arrays of tabulated values:
\begin{itemize}
\item {} 
\sphinxAtStartPar
embedding function F(rho) (Nrho values)

\item {} 
\sphinxAtStartPar
effective charge function Z(r) (Nr values)

\item {} 
\sphinxAtStartPar
density function rho(r) (Nr values)

\end{itemize}

\sphinxAtStartPar
The values for each array can be listed as multiple values per line,
so long as each array starts on a new line.  For example, the
individual Z(r) values are for r = 0,dr,2*dr, … (Nr\sphinxhyphen{}1)*dr.

\sphinxAtStartPar
The units for the embedding function F are eV.  The units for the
density function rho are the same as for drho (see above, electron
density).  The units for the effective charge Z are “atomic charge” or
sqrt(Hartree * Bohr\sphinxhyphen{}radii).  For two interacting atoms i,j this is used
by LAMMPS to compute the pair potential term in the EAM energy
expression as r*phi, in units of eV\sphinxhyphen{}Angstroms, via the formula
\begin{equation*}
\begin{split}r \cdot \phi = 27.2 \cdot 0.529 \cdot Z_i \cdot Z_j\end{split}
\end{equation*}
\sphinxAtStartPar
where 1 Hartree = 27.2 eV and 1 Bohr = 0.529 Angstroms.


\bigskip\hrule\bigskip


\sphinxAtStartPar
Style \sphinxstyleemphasis{eam/alloy} computes pairwise interactions using the same
formula as style \sphinxstyleemphasis{eam}.  However the associated
{\hyperref[\detokenize{pair_coeff::doc}]{\sphinxcrossref{\DUrole{doc}{pair\_coeff}}}} command reads a DYNAMO \sphinxstyleemphasis{setfl} file
instead of a \sphinxstyleemphasis{funcfl} file.  \sphinxstyleemphasis{Setfl} files can be used to model a
single\sphinxhyphen{}element or alloy system.  In the alloy case, as explained
above, \sphinxstyleemphasis{setfl} files contain explicit tabulated values for alloy
interactions.  Thus they allow more generality than \sphinxstyleemphasis{funcfl} files for
modeling alloys.

\sphinxAtStartPar
For style \sphinxstyleemphasis{eam/alloy}, potential values are read from a file that is
in the DYNAMO multi\sphinxhyphen{}element \sphinxstyleemphasis{setfl} format, except that element names
(Ni, Cu, etc) are added to one of the lines in the file.  If the
DYNAMO file was created by a Fortran program, it cannot have “D”
values in it for exponents.  C only recognizes “e” or “E” for
scientific notation.

\sphinxAtStartPar
Only a single pair\_coeff command is used with the \sphinxstyleemphasis{eam/alloy} style
which specifies a DYNAMO \sphinxstyleemphasis{setfl} file, which contains information for
M elements.  These are mapped to LAMMPS atom types by specifying N
additional arguments after the filename in the pair\_coeff command,
where N is the number of LAMMPS atom types:
\begin{itemize}
\item {} 
\sphinxAtStartPar
filename

\item {} 
\sphinxAtStartPar
N element names = mapping of \sphinxstyleemphasis{setfl} elements to atom types

\end{itemize}

\sphinxAtStartPar
As an example, the potentials/NiAlH\_jea.eam.alloy file is a \sphinxstyleemphasis{setfl}
file which has tabulated EAM values for 3 elements and their alloy
interactions: Ni, Al, and H.  See the {\hyperref[\detokenize{pair_coeff::doc}]{\sphinxcrossref{\DUrole{doc}{pair\_coeff}}}} doc
page for alternate ways to specify the path for the potential file.
If your LAMMPS simulation has 4 atoms types and you want the first 3 to
be Ni, and the fourth to be Al, you would use the following pair\_coeff
command:

\begin{sphinxVerbatim}[commandchars=\\\{\}]
\PYG{k}{pair\PYGZus{}coeff}\PYG{+w}{ }\PYG{o}{*}\PYG{+w}{ }\PYG{o}{*}\PYG{+w}{ }\PYG{n}{NiAlH\PYGZus{}jea.eam.alloy}\PYG{+w}{ }\PYG{n}{Ni}\PYG{+w}{ }\PYG{n}{Ni}\PYG{+w}{ }\PYG{n}{Ni}\PYG{+w}{ }\PYG{n}{Al}
\end{sphinxVerbatim}

\sphinxAtStartPar
The first 2 arguments must be * * so as to span all LAMMPS atom types.
The first three Ni arguments map LAMMPS atom types 1,2,3 to the Ni
element in the \sphinxstyleemphasis{setfl} file.  The final Al argument maps LAMMPS atom
type 4 to the Al element in the \sphinxstyleemphasis{setfl} file.  Note that there is no
requirement that your simulation use all the elements specified by the
\sphinxstyleemphasis{setfl} file.

\sphinxAtStartPar
If a mapping value is specified as NULL, the mapping is not performed.
This can be used when an \sphinxstyleemphasis{eam/alloy} potential is used as part of the
\sphinxstyleemphasis{hybrid} pair style.  The NULL values are placeholders for atom types
that will be used with other potentials.

\sphinxAtStartPar
\sphinxstyleemphasis{Setfl} files in the \sphinxstyleemphasis{potentials} directory of the LAMMPS distribution
have an “.eam.alloy” suffix.  A DYNAMO multi\sphinxhyphen{}element \sphinxstyleemphasis{setfl} file is
formatted as follows:
\begin{itemize}
\item {} 
\sphinxAtStartPar
lines 1,2,3 = comments (ignored)

\item {} 
\sphinxAtStartPar
line 4: Nelements Element1 Element2 … ElementN

\item {} 
\sphinxAtStartPar
line 5: Nrho, drho, Nr, dr, cutoff

\end{itemize}

\sphinxAtStartPar
In a DYNAMO \sphinxstyleemphasis{setfl} file, line 4 only lists Nelements = the \# of
elements in the \sphinxstyleemphasis{setfl} file.  For LAMMPS, the element name (Ni, Cu,
etc) of each element must be added to the line, in the order the
elements appear in the file.

\sphinxAtStartPar
The meaning and units of the values in line 5 is the same as for the
\sphinxstyleemphasis{funcfl} file described above.  Note that the cutoff (in Angstroms) is
a global value, valid for all pairwise interactions for all element
pairings.

\sphinxAtStartPar
Following the 5 header lines are Nelements sections, one for each
element, each with the following format:
\begin{itemize}
\item {} 
\sphinxAtStartPar
line 1 = atomic number, mass, lattice constant, lattice type (e.g. FCC)

\item {} 
\sphinxAtStartPar
embedding function F(rho) (Nrho values)

\item {} 
\sphinxAtStartPar
density function rho(r) (Nr values)

\end{itemize}

\sphinxAtStartPar
As with the \sphinxstyleemphasis{funcfl} files, only the mass (in mass {\hyperref[\detokenize{units::doc}]{\sphinxcrossref{\DUrole{doc}{units}}}},
e.g. mass number or grams/mole for metal units) is used by LAMMPS from
the first line.  The cubic lattice constant is in Angstroms.  The F and
rho arrays are unique to a single element and have the same format and
units as in a \sphinxstyleemphasis{funcfl} file.

\sphinxAtStartPar
Following the Nelements sections, Nr values for each pair potential
phi(r) array are listed for all i,j element pairs in the same format
as other arrays.  Since these interactions are symmetric (i,j = j,i)
only phi arrays with i \textgreater{}= j are listed, in the following order: i,j =
(1,1), (2,1), (2,2), (3,1), (3,2), (3,3), (4,1), …, (Nelements,
Nelements).  Unlike the effective charge array Z(r) in \sphinxstyleemphasis{funcfl} files,
the tabulated values for each phi function are listed in \sphinxstyleemphasis{setfl} files
directly as r*phi (in units of eV\sphinxhyphen{}Angstroms), since they are for atom
pairs.


\bigskip\hrule\bigskip


\sphinxAtStartPar
Style \sphinxstyleemphasis{eam/cd} is similar to the \sphinxstyleemphasis{eam/alloy} style, except that it
computes alloy pairwise interactions using the concentration\sphinxhyphen{}dependent
embedded\sphinxhyphen{}atom method (CD\sphinxhyphen{}EAM).  This model can reproduce the enthalpy
of mixing of alloys over the full composition range, as described in
{\hyperref[\detokenize{pair_eam:stukowski}]{\sphinxcrossref{\DUrole{std,std-ref}{(Stukowski)}}}}. Style \sphinxstyleemphasis{eam/cd/old} is an older, slightly
different and slower two\sphinxhyphen{}site formulation of the model {\hyperref[\detokenize{pair_eam:caro}]{\sphinxcrossref{\DUrole{std,std-ref}{(Caro)}}}}.

\sphinxAtStartPar
The pair\_coeff command is specified the same as for the \sphinxstyleemphasis{eam/alloy}
style.  However the DYNAMO \sphinxstyleemphasis{setfl} file must has two
lines added to it, at the end of the file:
\begin{itemize}
\item {} 
\sphinxAtStartPar
line 1: Comment line (ignored)

\item {} 
\sphinxAtStartPar
line 2: N Coefficient0 Coefficient1 … CoefficientN

\end{itemize}

\sphinxAtStartPar
The last line begins with the degree \sphinxstyleemphasis{N} of the polynomial function
\sphinxstyleemphasis{h(x)} that modifies the cross interaction between A and B elements.
Then \sphinxstyleemphasis{N+1} coefficients for the terms of the polynomial are then
listed.

\sphinxAtStartPar
Modified EAM \sphinxstyleemphasis{setfl} files used with the \sphinxstyleemphasis{eam/cd} style must contain
exactly two elements, i.e. in the current implementation the \sphinxstyleemphasis{eam/cd}
style only supports binary alloys.  The first and second elements in
the input EAM file are always taken as the \sphinxstyleemphasis{A} and \sphinxstyleemphasis{B} species.

\sphinxAtStartPar
\sphinxstyleemphasis{CD\sphinxhyphen{}EAM} files in the \sphinxstyleemphasis{potentials} directory of the LAMMPS
distribution have a “.cdeam” suffix.


\bigskip\hrule\bigskip


\sphinxAtStartPar
Style \sphinxstyleemphasis{eam/fs} computes pairwise interactions for metals and metal
alloys using a generalized form of EAM potentials due to Finnis and
Sinclair {\hyperref[\detokenize{pair_eam:finnis1}]{\sphinxcrossref{\DUrole{std,std-ref}{(Finnis)}}}}.  Style \sphinxstyleemphasis{eam/he} is similar to
\sphinxstyleemphasis{eam/fs} except that it allows for negative electron density in
order to capture the behavior of helium in metals {\hyperref[\detokenize{pair_eam:zhou6}]{\sphinxcrossref{\DUrole{std,std-ref}{(Zhou6)}}}}.

\sphinxAtStartPar
The total energy Ei of an atom I is given by
\begin{equation*}
\begin{split}E_i = F_\alpha \left(\sum_{j \neq i}\
\rho_{\alpha\beta} (r_{ij})\right) +
\frac{1}{2} \sum_{j \neq i} \phi_{\alpha\beta} (r_{ij})\end{split}
\end{equation*}
\sphinxAtStartPar
where \(\rho_{\alpha\beta}\) refers to the density contributed
by a neighbor atom J of element \(\beta\) at the site of atom I
of element \(\alpha\).
This has the same form as the EAM formula above, except that rho is
now a functional specific to the elements of both atoms I and J,
so that different elements can contribute differently to the total
electron density at an atomic site depending on the identity of the
element at that atomic site.

\sphinxAtStartPar
The associated {\hyperref[\detokenize{pair_coeff::doc}]{\sphinxcrossref{\DUrole{doc}{pair\_coeff}}}} command for style \sphinxstyleemphasis{eam/fs}
or \sphinxstyleemphasis{eam/he} reads a DYNAMO \sphinxstyleemphasis{setfl} file that has been extended to include
additional \(\rho_{\alpha\beta}\) arrays of tabulated values.  A
discussion of how FS EAM differs from conventional EAM alloy potentials is
given in {\hyperref[\detokenize{pair_eam:ackland1}]{\sphinxcrossref{\DUrole{std,std-ref}{(Ackland1)}}}}.  An example of such a potential is the
same author’s Fe\sphinxhyphen{}P FS potential {\hyperref[\detokenize{pair_eam:ackland2}]{\sphinxcrossref{\DUrole{std,std-ref}{(Ackland2)}}}}.  Note that while
FS potentials always specify the embedding energy with a square root
dependence on the total density, the implementation in LAMMPS does not
require that; the user can tabulate any functional form desired in the
FS potential files.

\sphinxAtStartPar
For style \sphinxstyleemphasis{eam/fs} and \sphinxstyleemphasis{eam/he} the form of the pair\_coeff command is exactly
the same as for style \sphinxstyleemphasis{eam/alloy}, e.g.

\begin{sphinxVerbatim}[commandchars=\\\{\}]
\PYG{k}{pair\PYGZus{}coeff}\PYG{+w}{ }\PYG{o}{*}\PYG{+w}{ }\PYG{o}{*}\PYG{+w}{ }\PYG{n}{NiAlH\PYGZus{}jea.eam.fs}\PYG{+w}{ }\PYG{n}{Ni}\PYG{+w}{ }\PYG{n}{Ni}\PYG{+w}{ }\PYG{n}{Ni}\PYG{+w}{ }\PYG{n}{Al}
\end{sphinxVerbatim}

\sphinxAtStartPar
with N additional arguments after the filename, where N is
the number of LAMMPS atom types.  See the {\hyperref[\detokenize{pair_coeff::doc}]{\sphinxcrossref{\DUrole{doc}{pair\_coeff}}}}
doc page for alternate ways to specify the path for the potential
file.  The N values determine the mapping of LAMMPS atom types to EAM
elements in the file, as described above for style \sphinxstyleemphasis{eam/alloy}.  As
with \sphinxstyleemphasis{eam/alloy}, if a mapping value is NULL, the mapping is not
performed.  This can be used when an \sphinxstyleemphasis{eam/fs} or \sphinxstyleemphasis{eam/he} potential is
used as part of a \sphinxstyleemphasis{hybrid} pair style.  The NULL values are used as
placeholders for atom types that will be used with other potentials.

\sphinxAtStartPar
FS EAM and HE EAM files include more information than the DYNAMO \sphinxstyleemphasis{setfl}
format files read by \sphinxstyleemphasis{eam/alloy}, in that i,j density functionals for
all pairs of elements are included as needed by the Finnis/Sinclair
formulation of the EAM.

\sphinxAtStartPar
FS EAM files in the \sphinxstyleemphasis{potentials} directory of the LAMMPS distribution
have an “.eam.fs” suffix.  They are formatted as follows:
\begin{itemize}
\item {} 
\sphinxAtStartPar
lines 1,2,3 = comments (ignored)

\item {} 
\sphinxAtStartPar
line 4: Nelements Element1 Element2 … ElementN

\item {} 
\sphinxAtStartPar
line 5: Nrho, drho, Nr, dr, cutoff

\end{itemize}

\sphinxAtStartPar
The 5\sphinxhyphen{}line header section is identical to an EAM \sphinxstyleemphasis{setfl} file.

\sphinxAtStartPar
Following the header are Nelements sections, one for each element \(\beta\),
each with the following format:
\begin{itemize}
\item {} 
\sphinxAtStartPar
line 1 = atomic number, mass, lattice constant, lattice type (e.g. FCC)

\item {} 
\sphinxAtStartPar
embedding function F(rho) (Nrho values)

\item {} 
\sphinxAtStartPar
density function \(\rho_{1\beta} (r)\) for element \(\beta\) at element 1 (Nr values)

\item {} 
\sphinxAtStartPar
density function \(\rho_{2\beta} (r)\) for element \(\beta\) at element 2

\item {} 
\sphinxAtStartPar
…

\item {} 
\sphinxAtStartPar
density function \(\rho_{N_{elem}\beta} (r)\) for element \(\beta\) at element \(N_{elem}\)

\end{itemize}

\sphinxAtStartPar
The units of these quantities in line 1 are the same as for \sphinxstyleemphasis{setfl}
files.  Note that the rho(r) arrays in Finnis/Sinclair can be
asymmetric (\(\rho_{\alpha\beta} (r) \neq \rho_{\beta\alpha} (r)\) )
so there are Nelements\textasciicircum{}2 of them listed in the file.

\sphinxAtStartPar
Following the Nelements sections, Nr values for each pair potential
phi(r) array are listed in the same manner (r*phi, units of
eV\sphinxhyphen{}Angstroms) as in EAM \sphinxstyleemphasis{setfl} files.  Note that in Finnis/Sinclair,
the phi(r) arrays are still symmetric, so only phi arrays for i \textgreater{}= j
are listed.

\sphinxAtStartPar
HE EAM files in the \sphinxstyleemphasis{potentials} directory of the LAMMPS distribution
have an “.eam.he” suffix.  They are formatted as follows:
\begin{itemize}
\item {} 
\sphinxAtStartPar
lines 1,2,3 = comments (ignored)

\item {} 
\sphinxAtStartPar
line 4: Nelements Element1 Element2 … ElementN

\item {} 
\sphinxAtStartPar
line 5: Nrho, drho, Nr, dr, cutoff, rhomax

\end{itemize}

\sphinxAtStartPar
The 5\sphinxhyphen{}line header section is identical to an FS EAM file
except that line 5 lists an additional value, rhomax. Unlike in FS EAM
files where embedding energies F(rho) are always defined between rho = 0
and rho = (Nrho \sphinxhyphen{}1)drho, F(rho) in HE EAM files are defined between
rho = rhomin and rho = rhomax.  Since drho = (rhomax \sphinxhyphen{} rhomin)/(Nrho \sphinxhyphen{} 1),
rhomin = rhomax \sphinxhyphen{} (Nrho \sphinxhyphen{} 1)drho.  The embedding energies F(rho) are
listed for rho = rhomin, rhomin + drho, rhomin + 2drho, …, rhomax.
This gives users additional flexibility to define a negative rhomin and
therefore an embedding energy function that works for both positive and
negative electron densities.  The format and units of these sections are
identical to the FS EAM files (see above).


\bigskip\hrule\bigskip


\sphinxAtStartPar
\DUrole{versionmodified,added}{Added in version 3Nov2022.}

\sphinxAtStartPar
The \sphinxstyleemphasis{eam}, \sphinxstyleemphasis{eam/alloy}, \sphinxstyleemphasis{eam/fs}, and \sphinxstyleemphasis{eam/he} pair styles support
extraction of two per\sphinxhyphen{}atom quantities by the {\hyperref[\detokenize{fix_pair::doc}]{\sphinxcrossref{\DUrole{doc}{fix pair}}}}
command.  This allows the quantities to be output to files by the
{\hyperref[\detokenize{dump::doc}]{\sphinxcrossref{\DUrole{doc}{dump}}}} or otherwise processed by other LAMMPS commands.

\sphinxAtStartPar
The names of the two quantities are “rho” and “fp” for the density and
derivative of the embedding energy for each atom.  Neither quantity
needs to be triggered by the {\hyperref[\detokenize{fix_pair::doc}]{\sphinxcrossref{\DUrole{doc}{fix pair}}}} command in order
for these pair styles to calculate it.


\bigskip\hrule\bigskip


\sphinxAtStartPar
Styles with a \sphinxstyleemphasis{gpu}, \sphinxstyleemphasis{intel}, \sphinxstyleemphasis{kk}, \sphinxstyleemphasis{omp}, or \sphinxstyleemphasis{opt} suffix are
functionally the same as the corresponding style without the suffix.
They have been optimized to run faster, depending on your available
hardware, as discussed on the {\hyperref[\detokenize{Speed_packages::doc}]{\sphinxcrossref{\DUrole{doc}{Accelerator packages}}}}
page.  The accelerated styles take the same arguments and should
produce the same results, except for round\sphinxhyphen{}off and precision issues.

\sphinxAtStartPar
These accelerated styles are part of the GPU, INTEL, KOKKOS,
OPENMP, and OPT packages, respectively.  They are only enabled if
LAMMPS was built with those packages.  See the {\hyperref[\detokenize{Build_package::doc}]{\sphinxcrossref{\DUrole{doc}{Build package}}}} page for more info.

\sphinxAtStartPar
You can specify the accelerated styles explicitly in your input script
by including their suffix, or you can use the {\hyperref[\detokenize{Run_options::doc}]{\sphinxcrossref{\DUrole{doc}{\sphinxhyphen{}suffix command\sphinxhyphen{}line switch}}}} when you invoke LAMMPS, or you can use the
{\hyperref[\detokenize{suffix::doc}]{\sphinxcrossref{\DUrole{doc}{suffix}}}} command in your input script.

\sphinxAtStartPar
See the {\hyperref[\detokenize{Speed_packages::doc}]{\sphinxcrossref{\DUrole{doc}{Accelerator packages}}}} page for more
instructions on how to use the accelerated styles effectively.


\bigskip\hrule\bigskip



\subsection{Mixing, shift, table, tail correction, restart, rRESPA info}
\label{\detokenize{pair_eam:mixing-shift-table-tail-correction-restart-rrespa-info}}
\sphinxAtStartPar
For atom type pairs I,J and I != J, where types I and J correspond to
two different element types, mixing is performed by LAMMPS as
described above with the individual styles.  You never need to specify
a pair\_coeff command with I != J arguments for the eam styles.

\sphinxAtStartPar
This pair style does not support the {\hyperref[\detokenize{pair_modify::doc}]{\sphinxcrossref{\DUrole{doc}{pair\_modify}}}}
shift, table, and tail options.

\sphinxAtStartPar
The eam pair styles do not write their information to {\hyperref[\detokenize{restart::doc}]{\sphinxcrossref{\DUrole{doc}{binary
restart files}}}}, since it is stored in tabulated potential
files.  Thus, you need to re\sphinxhyphen{}specify the pair\_style and pair\_coeff
commands in an input script that reads a restart file.

\sphinxAtStartPar
The eam pair styles can only be used via the \sphinxstyleemphasis{pair} keyword of the
{\hyperref[\detokenize{run_style::doc}]{\sphinxcrossref{\DUrole{doc}{run\_style respa}}}} command.  They do not support the
\sphinxstyleemphasis{inner}, \sphinxstyleemphasis{middle}, \sphinxstyleemphasis{outer} keywords.


\bigskip\hrule\bigskip



\subsection{Restrictions}
\label{\detokenize{pair_eam:restrictions}}
\sphinxAtStartPar
All of these styles are part of the MANYBODY package.  They are only
enabled if LAMMPS was built with that package.  See the {\hyperref[\detokenize{Build_package::doc}]{\sphinxcrossref{\DUrole{doc}{Build
package}}}} page for more info.


\subsection{Related commands}
\label{\detokenize{pair_eam:related-commands}}
\sphinxAtStartPar
{\hyperref[\detokenize{pair_coeff::doc}]{\sphinxcrossref{\DUrole{doc}{pair\_coeff}}}}


\subsection{Default}
\label{\detokenize{pair_eam:default}}
\sphinxAtStartPar
none


\bigskip\hrule\bigskip

\phantomsection\label{\detokenize{pair_eam:ackland1}}
\sphinxAtStartPar
\sphinxstylestrong{(Ackland1)} Ackland, Condensed Matter (2005).

\phantomsection\label{\detokenize{pair_eam:ackland2}}
\sphinxAtStartPar
\sphinxstylestrong{(Ackland2)} Ackland, Mendelev, Srolovitz, Han and Barashev, Journal
of Physics: Condensed Matter, 16, S2629 (2004).

\phantomsection\label{\detokenize{pair_eam:daw}}
\sphinxAtStartPar
\sphinxstylestrong{(Daw)} Daw, Baskes, Phys Rev Lett, 50, 1285 (1983).
Daw, Baskes, Phys Rev B, 29, 6443 (1984).

\phantomsection\label{\detokenize{pair_eam:finnis1}}
\sphinxAtStartPar
\sphinxstylestrong{(Finnis)} Finnis, Sinclair, Philosophical Magazine A, 50, 45 (1984).

\phantomsection\label{\detokenize{pair_eam:zhou6}}
\sphinxAtStartPar
\sphinxstylestrong{(Zhou6)} Zhou, Bartelt, Sills, Physical Review B, 103, 014108 (2021).

\phantomsection\label{\detokenize{pair_eam:stukowski}}
\sphinxAtStartPar
\sphinxstylestrong{(Stukowski)} Stukowski, Sadigh, Erhart, Caro; Modeling Simulation
Materials Science \& Engineering, 7, 075005 (2009).

\phantomsection\label{\detokenize{pair_eam:caro}}
\sphinxAtStartPar
\sphinxstylestrong{(Caro)} A Caro, DA Crowson, M Caro; Phys Rev Lett, 95, 075702 (2005)

\sphinxstepscope

\index{pair\_style edip@\spxentry{pair\_style edip}}\ignorespaces 
\index{pair\_style edip/omp@\spxentry{pair\_style edip/omp}}\ignorespaces 
\index{pair\_style edip/multi@\spxentry{pair\_style edip/multi}}\ignorespaces 

\section{pair\_style edip command}
\label{\detokenize{pair_edip:pair-style-edip-command}}\label{\detokenize{pair_edip:index-2}}\label{\detokenize{pair_edip:index-1}}\label{\detokenize{pair_edip:index-0}}\label{\detokenize{pair_edip::doc}}
\sphinxAtStartPar
Accelerator Variants: \sphinxstyleemphasis{edip/omp}


\section{pair\_style edip/multi command}
\label{\detokenize{pair_edip:pair-style-edip-multi-command}}

\subsection{Syntax}
\label{\detokenize{pair_edip:syntax}}
\begin{sphinxVerbatim}[commandchars=\\\{\}]
\PYG{k}{pair\PYGZus{}style}\PYG{+w}{ }\PYG{n}{style}
\end{sphinxVerbatim}
\begin{itemize}
\item {} 
\sphinxAtStartPar
style = \sphinxstyleemphasis{edip} or \sphinxstyleemphasis{edip/multi}

\end{itemize}


\subsection{Examples}
\label{\detokenize{pair_edip:examples}}
\begin{sphinxVerbatim}[commandchars=\\\{\}]
\PYG{k}{pair\PYGZus{}style}\PYG{+w}{ }\PYG{n}{edip}
\PYG{k}{pair\PYGZus{}coeff}\PYG{+w}{ }\PYG{o}{*}\PYG{+w}{ }\PYG{o}{*}\PYG{+w}{ }\PYG{n}{Si.edip}\PYG{+w}{ }\PYG{n}{Si}
\end{sphinxVerbatim}


\subsection{Description}
\label{\detokenize{pair_edip:description}}
\sphinxAtStartPar
The \sphinxstyleemphasis{edip} and \sphinxstyleemphasis{edip/multi} styles compute a 3\sphinxhyphen{}body {\hyperref[\detokenize{pair_edip:edip}]{\sphinxcrossref{\DUrole{std,std-ref}{EDIP}}}}
potential which is popular for modeling silicon materials where
it can have advantages over other models such as the
{\hyperref[\detokenize{pair_sw::doc}]{\sphinxcrossref{\DUrole{doc}{Stillinger\sphinxhyphen{}Weber}}}} or {\hyperref[\detokenize{pair_tersoff::doc}]{\sphinxcrossref{\DUrole{doc}{Tersoff}}}}
potentials. The \sphinxstyleemphasis{edip} style has been programmed for single element
potentials, while \sphinxstyleemphasis{edip/multi} supports multi\sphinxhyphen{}element EDIP runs.

\sphinxAtStartPar
In EDIP, the energy E of a system of atoms is
\begin{equation*}
\begin{split}E  = & \sum_{j \ne i} \phi_{2}(R_{ij}, Z_{i}) + \sum_{j \ne i} \sum_{k \ne i,k > j} \phi_{3}(R_{ij}, R_{ik}, Z_{i}) \\
\phi_{2}(r, Z)  = & A\left[\left(\frac{B}{r}\right)^{\rho} - e^{-\beta Z^2}\right]exp{\left(\frac{\sigma}{r-a}\right)} \\
\phi_{3}(R_{ij}, R_{ik}, Z_i)  = & exp{\left(\frac{\gamma}{R_{ij}-a}\right)}exp{\left(\frac{\gamma}{R_{ik}-a}\right)}h(cos\theta_{ijk},Z_i) \\
Z_i  = & \sum_{m \ne i} f(R_{im}) \qquad
f(r) = \begin{cases}
       1 & \quad r<c \\
       \exp\left(\frac{\alpha}{1-x^{-3}}\right) & \quad c<r<a \\
       0 & \quad r>a
       \end{cases} \\
h(l,Z)  = & \lambda [(1-e^{-Q(Z)(l+\tau(Z))^2}) + \eta Q(Z)(l+\tau(Z))^2 ] \\
Q(Z)  = & Q_0 e^{-\mu Z} \qquad \tau(Z) = u_1 + u_2 (u_3 e^{-u_4 Z} - e^{-2u_4 Z})\end{split}
\end{equation*}
\sphinxAtStartPar
where \(\phi_2\) is a two\sphinxhyphen{}body term and \(\phi_3\) is a
three\sphinxhyphen{}body term.  The summations in the formula are over all neighbors J
and K of atom I within a cutoff distance = a.  Both terms depend on the
local environment of atom I through its effective coordination number
defined by Z, which is unity for a cutoff distance \textless{} c and gently goes
to 0 at distance = a.

\sphinxAtStartPar
Only a single pair\_coeff command is used with the \sphinxstyleemphasis{edip} style which
specifies a EDIP potential file with parameters for all
needed elements.  These are mapped to LAMMPS atom types by specifying
N additional arguments after the filename in the pair\_coeff command,
where N is the number of LAMMPS atom types:
\begin{itemize}
\item {} 
\sphinxAtStartPar
filename

\item {} 
\sphinxAtStartPar
N element names = mapping of EDIP elements to atom types

\end{itemize}

\sphinxAtStartPar
See the {\hyperref[\detokenize{pair_coeff::doc}]{\sphinxcrossref{\DUrole{doc}{pair\_coeff}}}} page for alternate ways
to specify the path for the potential file.

\sphinxAtStartPar
As an example, imagine a file Si.edip has EDIP values for Si.

\sphinxAtStartPar
EDIP files in the \sphinxstyleemphasis{potentials} directory of the LAMMPS
distribution have a “.edip” suffix.  Lines that are not blank or
comments (starting with \#) define parameters for a triplet of
elements.  The parameters in a single entry correspond to the two\sphinxhyphen{}body
and three\sphinxhyphen{}body coefficients in the formula above:
\begin{itemize}
\item {} 
\sphinxAtStartPar
element 1 (the center atom in a 3\sphinxhyphen{}body interaction)

\item {} 
\sphinxAtStartPar
element 2

\item {} 
\sphinxAtStartPar
element 3

\item {} 
\sphinxAtStartPar
A (energy units)

\item {} 
\sphinxAtStartPar
B (distance units)

\item {} 
\sphinxAtStartPar
cutoffA (distance units)

\item {} 
\sphinxAtStartPar
cutoffC (distance units)

\item {} 
\sphinxAtStartPar
\(\alpha\)

\item {} 
\sphinxAtStartPar
\(\beta\)

\item {} 
\sphinxAtStartPar
\(\eta\)

\item {} 
\sphinxAtStartPar
\(\gamma\) (distance units)

\item {} 
\sphinxAtStartPar
\(lambda\) (energy units)

\item {} 
\sphinxAtStartPar
\(\mu\)

\item {} 
\sphinxAtStartPar
\(\tau\)

\item {} 
\sphinxAtStartPar
\(\sigma\) (distance units)

\item {} 
\sphinxAtStartPar
Q0

\item {} 
\sphinxAtStartPar
u1

\item {} 
\sphinxAtStartPar
u2

\item {} 
\sphinxAtStartPar
u3

\item {} 
\sphinxAtStartPar
u4

\end{itemize}

\sphinxAtStartPar
The A, B, beta, sigma parameters are used only for two\sphinxhyphen{}body interactions.
The eta, gamma, lambda, mu, Q0 and all u1 to u4 parameters are used only
for three\sphinxhyphen{}body interactions. The alpha and cutoffC parameters are used
for the coordination environment function only.

\sphinxAtStartPar
The EDIP potential file must contain entries for all the
elements listed in the pair\_coeff command.  It can also contain
entries for additional elements not being used in a particular
simulation; LAMMPS ignores those entries.

\sphinxAtStartPar
For a single\sphinxhyphen{}element simulation, only a single entry is required
(e.g. SiSiSi).  For a two\sphinxhyphen{}element simulation, the file must contain 8
entries (for SiSiSi, SiSiC, SiCSi, SiCC, CSiSi, CSiC, CCSi, CCC), that
specify EDIP parameters for all permutations of the two elements
interacting in three\sphinxhyphen{}body configurations.  Thus for 3 elements, 27
entries would be required, etc.

\sphinxAtStartPar
At the moment, only a single element parameterization is
implemented. However, the author is not aware of other
multi\sphinxhyphen{}element EDIP parameterization. If you know any and
you are interest in that, please contact the author of
the EDIP package.


\bigskip\hrule\bigskip


\sphinxAtStartPar
Styles with a \sphinxstyleemphasis{gpu}, \sphinxstyleemphasis{intel}, \sphinxstyleemphasis{kk}, \sphinxstyleemphasis{omp}, or \sphinxstyleemphasis{opt} suffix are
functionally the same as the corresponding style without the suffix.
They have been optimized to run faster, depending on your available
hardware, as discussed on the {\hyperref[\detokenize{Speed_packages::doc}]{\sphinxcrossref{\DUrole{doc}{Accelerator packages}}}}
page.  The accelerated styles take the same arguments and should
produce the same results, except for round\sphinxhyphen{}off and precision issues.

\sphinxAtStartPar
These accelerated styles are part of the GPU, INTEL, KOKKOS,
OPENMP, and OPT packages, respectively.  They are only enabled if
LAMMPS was built with those packages.  See the {\hyperref[\detokenize{Build_package::doc}]{\sphinxcrossref{\DUrole{doc}{Build package}}}} page for more info.

\sphinxAtStartPar
You can specify the accelerated styles explicitly in your input script
by including their suffix, or you can use the {\hyperref[\detokenize{Run_options::doc}]{\sphinxcrossref{\DUrole{doc}{\sphinxhyphen{}suffix command\sphinxhyphen{}line switch}}}} when you invoke LAMMPS, or you can use the
{\hyperref[\detokenize{suffix::doc}]{\sphinxcrossref{\DUrole{doc}{suffix}}}} command in your input script.

\sphinxAtStartPar
See the {\hyperref[\detokenize{Speed_packages::doc}]{\sphinxcrossref{\DUrole{doc}{Accelerator packages}}}} page for more
instructions on how to use the accelerated styles effectively.


\bigskip\hrule\bigskip



\subsection{Mixing, shift, table, tail correction, restart, rRESPA info}
\label{\detokenize{pair_edip:mixing-shift-table-tail-correction-restart-rrespa-info}}
\sphinxAtStartPar
This pair style does not support the {\hyperref[\detokenize{pair_modify::doc}]{\sphinxcrossref{\DUrole{doc}{pair\_modify}}}}
shift, table, and tail options.

\sphinxAtStartPar
This pair style does not write its information to {\hyperref[\detokenize{restart::doc}]{\sphinxcrossref{\DUrole{doc}{binary restart files}}}}, since it is stored in potential files.  Thus, you
need to re\sphinxhyphen{}specify the pair\_style and pair\_coeff commands in an input
script that reads a restart file.

\sphinxAtStartPar
This pair style can only be used via the \sphinxstyleemphasis{pair} keyword of the
{\hyperref[\detokenize{run_style::doc}]{\sphinxcrossref{\DUrole{doc}{run\_style respa}}}} command.  It does not support the
\sphinxstyleemphasis{inner}, \sphinxstyleemphasis{middle}, \sphinxstyleemphasis{outer} keywords.


\bigskip\hrule\bigskip



\subsection{Restrictions}
\label{\detokenize{pair_edip:restrictions}}
\sphinxAtStartPar
This pair style can only be used if LAMMPS was built with the
MANYBODY package.  See the {\hyperref[\detokenize{Build_package::doc}]{\sphinxcrossref{\DUrole{doc}{Build package}}}} doc
page for more info.

\sphinxAtStartPar
This pair style requires the {\hyperref[\detokenize{newton::doc}]{\sphinxcrossref{\DUrole{doc}{newton}}}} setting to be “on”
for pair interactions.

\sphinxAtStartPar
The EDIP potential files provided with LAMMPS (see the potentials directory)
are parameterized for metal {\hyperref[\detokenize{units::doc}]{\sphinxcrossref{\DUrole{doc}{units}}}}.
You can use the EDIP potential with any LAMMPS units, but you would need
to create your own EDIP potential file with coefficients listed in the
appropriate units if your simulation does not use “metal” units.


\subsection{Related commands}
\label{\detokenize{pair_edip:related-commands}}
\sphinxAtStartPar
{\hyperref[\detokenize{pair_coeff::doc}]{\sphinxcrossref{\DUrole{doc}{pair\_coeff}}}}


\subsection{Default}
\label{\detokenize{pair_edip:default}}
\sphinxAtStartPar
none


\bigskip\hrule\bigskip

\phantomsection\label{\detokenize{pair_edip:edip}}
\sphinxAtStartPar
\sphinxstylestrong{(EDIP)} J F Justo et al, Phys Rev B 58, 2539 (1998).

\sphinxstepscope

\index{pair\_style eff/cut@\spxentry{pair\_style eff/cut}}\ignorespaces 

\section{pair\_style eff/cut command}
\label{\detokenize{pair_eff:pair-style-eff-cut-command}}\label{\detokenize{pair_eff:index-0}}\label{\detokenize{pair_eff::doc}}

\subsection{Syntax}
\label{\detokenize{pair_eff:syntax}}
\begin{sphinxVerbatim}[commandchars=\\\{\}]
\PYG{k}{pair\PYGZus{}style}\PYG{+w}{ }\PYG{n}{eff}\PYG{o}{/}\PYG{n}{cut}\PYG{+w}{ }\PYG{n}{cutoff}\PYG{+w}{ }\PYG{n}{keyword}\PYG{+w}{ }\PYG{n}{args}\PYG{+w}{ }\PYG{n}{...}
\end{sphinxVerbatim}
\begin{itemize}
\item {} 
\sphinxAtStartPar
cutoff = global cutoff for Coulombic interactions

\item {} 
\sphinxAtStartPar
zero or more keyword/value pairs may be appended
\begin{sphinxalltt}
keyword = \sphinxstyleemphasis{limit/eradius} or \sphinxstyleemphasis{pressure/evirials} or \sphinxstyleemphasis{ecp}
  \sphinxstyleemphasis{limit/eradius} args = none
  \sphinxstyleemphasis{pressure/evirials} args = none
  \sphinxstyleemphasis{ecp} args = type element type element ...
    type = LAMMPS atom type (1 to Ntypes)
    element = element symbol (e.g. H, Si)
\end{sphinxalltt}

\end{itemize}


\subsection{Examples}
\label{\detokenize{pair_eff:examples}}
\begin{sphinxVerbatim}[commandchars=\\\{\}]
\PYG{k}{pair\PYGZus{}style}\PYG{+w}{ }\PYG{n}{eff}\PYG{o}{/}\PYG{n}{cut}\PYG{+w}{ }\PYG{l+m}{39.7}
\PYG{k}{pair\PYGZus{}style}\PYG{+w}{ }\PYG{n}{eff}\PYG{o}{/}\PYG{n}{cut}\PYG{+w}{ }\PYG{l+m}{40.0}\PYG{+w}{ }\PYG{n}{limit}\PYG{o}{/}\PYG{n}{eradius}
\PYG{k}{pair\PYGZus{}style}\PYG{+w}{ }\PYG{n}{eff}\PYG{o}{/}\PYG{n}{cut}\PYG{+w}{ }\PYG{l+m}{40.0}\PYG{+w}{ }\PYG{n}{limit}\PYG{o}{/}\PYG{n}{eradius}\PYG{+w}{ }\PYG{n}{pressure}\PYG{o}{/}\PYG{n}{evirials}
\PYG{k}{pair\PYGZus{}style}\PYG{+w}{ }\PYG{n}{eff}\PYG{o}{/}\PYG{n}{cut}\PYG{+w}{ }\PYG{l+m}{40.0}\PYG{+w}{ }\PYG{n}{ecp}\PYG{+w}{ }\PYG{l+m}{1}\PYG{+w}{ }\PYG{n}{Si}\PYG{+w}{ }\PYG{l+m}{3}\PYG{+w}{ }\PYG{n}{C}
\PYG{k}{pair\PYGZus{}coeff}\PYG{+w}{ }\PYG{o}{*}\PYG{+w}{ }\PYG{o}{*}
\PYG{k}{pair\PYGZus{}coeff}\PYG{+w}{ }\PYG{l+m}{2}\PYG{+w}{ }\PYG{l+m}{2}\PYG{+w}{ }\PYG{l+m}{20.0}
\PYG{k}{pair\PYGZus{}coeff}\PYG{+w}{ }\PYG{l+m}{1}\PYG{+w}{ }\PYG{n}{s}\PYG{+w}{ }\PYG{l+m}{0.320852}\PYG{+w}{ }\PYG{l+m}{2.283269}\PYG{+w}{ }\PYG{l+m}{0.814857}
\PYG{k}{pair\PYGZus{}coeff}\PYG{+w}{ }\PYG{l+m}{3}\PYG{+w}{ }\PYG{n}{p}\PYG{+w}{ }\PYG{l+m}{22.721015}\PYG{+w}{ }\PYG{l+m}{0.728733}\PYG{+w}{ }\PYG{l+m}{1.103199}\PYG{+w}{ }\PYG{l+m}{17.695345}\PYG{+w}{ }\PYG{l+m}{6.693621}
\end{sphinxVerbatim}


\subsection{Description}
\label{\detokenize{pair_eff:description}}
\sphinxAtStartPar
This pair style contains a LAMMPS implementation of the electron Force
Field (eFF) potential currently under development at Caltech, as
described in {\hyperref[\detokenize{pair_eff:jaramillo-botero}]{\sphinxcrossref{\DUrole{std,std-ref}{(Jaramillo\sphinxhyphen{}Botero)}}}}.  The eFF for Z\textless{}6
was first introduced by {\hyperref[\detokenize{pair_eff:su}]{\sphinxcrossref{\DUrole{std,std-ref}{(Su)}}}} in 2007. It has been extended to
higher Zs by using effective core potentials (ECPs) that now cover up
to second and third row p\sphinxhyphen{}block elements of the periodic table.

\sphinxAtStartPar
eFF can be viewed as an approximation to QM wave packet dynamics and
Fermionic molecular dynamics, combining the ability of electronic
structure methods to describe atomic structure, bonding, and chemistry
in materials, and of plasma methods to describe nonequilibrium
dynamics of large systems with a large number of highly excited
electrons.  Yet, eFF relies on a simplification of the electronic
wave function in which electrons are described as floating Gaussian
wave packets whose position and size respond to the various dynamic
forces between interacting classical nuclear particles and spherical
Gaussian electron wave packets.  The wave function is taken to be a
Hartree product of the wave packets.  To compensate for the lack of
explicit antisymmetry in the resulting wave function, a spin\sphinxhyphen{}dependent
Pauli potential is included in the Hamiltonian.  Substituting this
wave function into the time\sphinxhyphen{}dependent Schrodinger equation produces
equations of motion that correspond \sphinxhyphen{} to second order \sphinxhyphen{} to classical
Hamiltonian relations between electron position and size, and their
conjugate momenta.  The N\sphinxhyphen{}electron wave function is described as a
product of one\sphinxhyphen{}electron Gaussian functions, whose size is a dynamical
variable and whose position is not constrained to a nuclear
center. This form allows for straightforward propagation of the
wave function, with time, using a simple formulation from which the
equations of motion are then integrated with conventional MD
algorithms. In addition to this spin\sphinxhyphen{}dependent Pauli repulsion
potential term between Gaussians, eFF includes the electron kinetic
energy from the Gaussians.  These two terms are based on
first\sphinxhyphen{}principles quantum mechanics.  On the other hand, nuclei are
described as point charges, which interact with other nuclei and
electrons through standard electrostatic potential forms.

\sphinxAtStartPar
The full Hamiltonian (shown below), contains then a standard
description for electrostatic interactions between a set of
delocalized point and Gaussian charges which include, nuclei\sphinxhyphen{}nuclei
(NN), electron\sphinxhyphen{}electron (ee), and nuclei\sphinxhyphen{}electron (Ne). Thus, eFF is a
mixed QM\sphinxhyphen{}classical mechanics method rather than a conventional force
field method (in which electron motions are averaged out into ground
state nuclear motions, i.e a single electronic state, and particle
interactions are described via empirically parameterized interatomic
potential functions). This makes eFF uniquely suited to simulate
materials over a wide range of temperatures and pressures where
electronically excited and ionized states of matter can occur and
coexist.  Furthermore, the interactions between particles \sphinxhyphen{}nuclei and
electrons\sphinxhyphen{} reduce to the sum of a set of effective pairwise potentials
in the eFF formulation.  The \sphinxstyleemphasis{eff/cut} style computes the pairwise
Coulomb interactions between nuclei and electrons (E\_NN,E\_Ne,E\_ee),
and the quantum\sphinxhyphen{}derived Pauli (E\_PR) and Kinetic energy interactions
potentials between electrons (E\_KE) for a total energy expression
given as,
\begin{equation*}
\begin{split}U\left(R,r,s\right) =  E_{NN} \left( R \right) + E_{Ne} \left( {R,r,s} \right) + E_{ee} \left( {r,s} \right) + E_{KE} \left( {r,s} \right) + E_{PR} \left( { \uparrow  \downarrow ,S} \right)\end{split}
\end{equation*}
\sphinxAtStartPar
The individual terms are defined as follows:
\begin{equation*}
\begin{split}E_{KE}  = & \frac{\hbar^2 }{{m_{e} }}\sum\limits_i {\frac{3}{{2s_i^2 }}} \\
E_{NN}  = & \frac{1}{{4\pi \varepsilon _0 }}\sum\limits_{i < j} {\frac{{Z_i Z_j }}{{R_{ij} }}} \\
E_{Ne}  = & - \frac{1}{{4\pi \varepsilon _0 }}\sum\limits_{i,j} {\frac{{Z_i }}{{R_{ij} }}Erf\left( {\frac{{\sqrt 2 R_{ij} }}{{s_j }}} \right)} \\
E_{ee}  = & \frac{1}{{4\pi \varepsilon _0 }}\sum\limits_{i < j} {\frac{1}{{r_{ij} }}Erf\left( {\frac{{\sqrt 2 r_{ij} }}{{\sqrt {s_i^2  + s_j^2 } }}} \right)} \\
E_{Pauli}  = & \sum\limits_{\sigma _i  = \sigma _j } {E\left( { \uparrow  \uparrow } \right)_{ij}}  + \sum\limits_{\sigma _i  \ne \sigma _j } {E\left( { \uparrow  \downarrow } \right)_{ij}} \\\end{split}
\end{equation*}
\sphinxAtStartPar
where, s\_i correspond to the electron sizes, the sigmas i’s to the
fixed spins of the electrons, Z\_i to the charges on the nuclei, R\_ij
to the distances between the nuclei or the nuclei and electrons, and
r\_ij to the distances between electrons.  For additional details see
{\hyperref[\detokenize{pair_eff:jaramillo-botero}]{\sphinxcrossref{\DUrole{std,std-ref}{(Jaramillo\sphinxhyphen{}Botero)}}}}.

\sphinxAtStartPar
The overall electrostatics energy is given in Hartree units of energy
by default and can be modified by an energy\sphinxhyphen{}conversion constant,
according to the units chosen (see {\hyperref[\detokenize{units::doc}]{\sphinxcrossref{\DUrole{doc}{electron\_units}}}}).  The
cutoff Rc, given in Bohrs (by default), truncates the interaction
distance.  The recommended cutoff for this pair style should follow
the minimum image criterion, i.e. half of the minimum unit cell
length.

\sphinxAtStartPar
This potential is designed to be used with {\hyperref[\detokenize{atom_style::doc}]{\sphinxcrossref{\DUrole{doc}{atom\_style electron}}}} definitions, in order to handle the
description of systems with interacting nuclei and explicit electrons.

\sphinxAtStartPar
The following coefficients must be defined for each pair of atoms
types via the {\hyperref[\detokenize{pair_coeff::doc}]{\sphinxcrossref{\DUrole{doc}{pair\_coeff}}}} command as in the examples
above, or in the data file or restart files read by the
{\hyperref[\detokenize{read_data::doc}]{\sphinxcrossref{\DUrole{doc}{read\_data}}}} or {\hyperref[\detokenize{read_restart::doc}]{\sphinxcrossref{\DUrole{doc}{read\_restart}}}}
commands, or by mixing as described below:
\begin{itemize}
\item {} 
\sphinxAtStartPar
cutoff (distance units)

\end{itemize}

\sphinxAtStartPar
For \sphinxstyleemphasis{eff/cut}, the cutoff coefficient is optional.  If it is not used
(as in some of the examples above), the default global value specified
in the pair\_style command is used.


\bigskip\hrule\bigskip


\sphinxAtStartPar
The \sphinxstyleemphasis{limit/eradius} and \sphinxstyleemphasis{pressure/evirials} keywords are optional.
Neither or both must be specified.  If not specified they are unset.

\sphinxAtStartPar
The \sphinxstyleemphasis{limit/eradius} keyword is used to restrain electron size from
becoming excessively diffuse at very high temperatures were the
Gaussian wave packet representation breaks down, and from expanding as
free particles to infinite size.  If unset, electron radius is free to
increase without bounds.  If set, a restraining harmonic potential of
the form E = 1/2k\_ss\textasciicircum{}2 for s \textgreater{} L\_box/2, where k\_s = 1 Hartrees/Bohr\textasciicircum{}2,
is applied on the electron radius.

\sphinxAtStartPar
The \sphinxstyleemphasis{pressure/evirials} keyword is used to control between two types
of pressure computation: if unset, the computed pressure does not
include the electronic radial virials contributions to the total
pressure (scalar or tensor).  If set, the computed pressure will
include the electronic radial virial contributions to the total
pressure (scalar and tensor).

\sphinxAtStartPar
The \sphinxstyleemphasis{ecp} keyword is used to associate an ECP representation for a
particular atom type.  The ECP captures the orbital overlap between a
core pseudo particle and valence electrons within the Pauli repulsion.
A list of type:element\sphinxhyphen{}symbol pairs may be provided for all ECP
representations, after the “ecp” keyword.

\begin{sphinxadmonition}{note}{Note:}
\sphinxAtStartPar
Default ECP parameters are provided for C, N, O, Al, and Si.
Users can modify these using the pair\_coeff command as exemplified
above.  For this, the User must distinguish between two different
functional forms supported, one that captures the orbital overlap
assuming the s\sphinxhyphen{}type core interacts with an s\sphinxhyphen{}like valence electron
(s\sphinxhyphen{}s) and another that assumes the interaction is s\sphinxhyphen{}p.  For systems
that exhibit significant p\sphinxhyphen{}character (e.g. C, N, O) the s\sphinxhyphen{}p form is
recommended. The “s” ECP form requires 3 parameters and the “p” 5
parameters.
\end{sphinxadmonition}

\begin{sphinxadmonition}{note}{Note:}
\sphinxAtStartPar
There are two different pressures that can be reported for eFF
when defining this pair\_style, one (default) that considers electrons
do not contribute radial virial components (i.e. electrons treated as
incompressible ‘rigid’ spheres) and one that does.  The radial
electronic contributions to the virials are only tallied if the
flexible pressure option is set, and this will affect both global and
per\sphinxhyphen{}atom quantities.  In principle, the true pressure of a system is
somewhere in between the rigid and the flexible eFF pressures, but,
for most cases, the difference between these two pressures will not be
significant over long\sphinxhyphen{}term averaged runs (i.e. even though the energy
partitioning changes, the total energy remains similar).
\end{sphinxadmonition}


\bigskip\hrule\bigskip


\begin{sphinxadmonition}{note}{Note:}
\sphinxAtStartPar
This implementation of eFF gives a reasonably accurate description
for systems containing nuclei from Z = 1\sphinxhyphen{}6 in “all electron”
representations.  For systems with increasingly non\sphinxhyphen{}spherical
electrons, Users should use the ECP representations.  ECPs are now
supported and validated for most of the second and third row elements of
the p\sphinxhyphen{}block.  Predefined parameters are provided for C, N, O, Al, and
Si.  The ECP captures the orbital overlap between the core and valence
electrons (i.e. Pauli repulsion) with one of the functional forms:
\end{sphinxadmonition}
\begin{equation*}
\begin{split}E_{Pauli(ECP_s)} = & p_1\exp\left(-\frac{p_2r^2}{p_3+s^2} \right) \\
E_{Pauli(ECP_p)} = & p_1\left( \frac{2}{p_2/s+s/p_2} \right)\left( r-p_3s\right)^2\exp \left[ -\frac{p_4\left( r-p_3s \right)^2}{p_5+s^2} \right]\end{split}
\end{equation*}
\sphinxAtStartPar
Where the first form correspond to core interactions with s\sphinxhyphen{}type valence
electrons and the second to core interactions with p\sphinxhyphen{}type valence
electrons.

\sphinxAtStartPar
The current version adds full support for models with fixed\sphinxhyphen{}core and
ECP definitions.  to enable larger timesteps (i.e. by avoiding the
high frequency vibrational modes \sphinxhyphen{}translational and radial\sphinxhyphen{} of the 2 s
electrons), and in the ECP case to reduce the increased orbital
complexity in higher Z elements (up to Z\textless{}18).  A fixed\sphinxhyphen{}core should be
defined with a mass that includes the corresponding nuclear mass plus
the 2 s electrons in atomic mass units (2x5.4857990943e\sphinxhyphen{}4), and a
radius equivalent to that of minimized 1s electrons (see examples
under /examples/PACKAGES/eff/fixed\sphinxhyphen{}core).  An pseudo\sphinxhyphen{}core should be
described with a mass that includes the corresponding nuclear mass,
plus all the core electrons (i.e no outer shell electrons), and a
radius equivalent to that of a corresponding minimized full\sphinxhyphen{}electron
system.  The charge for a pseudo\sphinxhyphen{}core atom should be given by the
number of outer shell electrons.

\sphinxAtStartPar
In general, eFF excels at computing the properties of materials in
extreme conditions and tracing the system dynamics over multi\sphinxhyphen{}picosecond
timescales; this is particularly relevant where electron excitations
can change significantly the nature of bonding in the system. It can
capture with surprising accuracy the behavior of such systems because
it describes consistently and in an unbiased manner many different
kinds of bonds, including covalent, ionic, multicenter, ionic, and
plasma, and how they interconvert and/or change when they become
excited.  eFF also excels in computing the relative thermochemistry of
isodemic reactions and conformational changes, where the bonds of the
reactants are of the same type as the bonds of the products.  eFF
assumes that kinetic energy differences dominate the overall exchange
energy, which is true when the electrons present are nearly spherical
and nodeless and valid for covalent compounds such as dense hydrogen,
hydrocarbons, and diamond; alkali metals (e.g. lithium), alkali earth
metals (e.g. beryllium) and semimetals such as boron; and various
compounds containing ionic and/or multicenter bonds, such as boron
dihydride.


\bigskip\hrule\bigskip



\subsection{Mixing, shift, table, tail correction, restart, rRESPA info}
\label{\detokenize{pair_eff:mixing-shift-table-tail-correction-restart-rrespa-info}}
\sphinxAtStartPar
For atom type pairs I,J and I != J, the cutoff distance for the
\sphinxstyleemphasis{eff/cut} style can be mixed.  The default mix value is \sphinxstyleemphasis{geometric}.
See the “pair\_modify” command for details.

\sphinxAtStartPar
The {\hyperref[\detokenize{pair_modify::doc}]{\sphinxcrossref{\DUrole{doc}{pair\_modify}}}} shift option is not relevant for
these pair styles.

\sphinxAtStartPar
The \sphinxstyleemphasis{eff/long} (not yet available) style supports the
{\hyperref[\detokenize{pair_modify::doc}]{\sphinxcrossref{\DUrole{doc}{pair\_modify}}}} table option for tabulation of the
short\sphinxhyphen{}range portion of the long\sphinxhyphen{}range Coulombic interaction.

\sphinxAtStartPar
These pair styles do not support the {\hyperref[\detokenize{pair_modify::doc}]{\sphinxcrossref{\DUrole{doc}{pair\_modify}}}}
tail option for adding long\sphinxhyphen{}range tail corrections to energy and
pressure.

\sphinxAtStartPar
These pair styles write their information to {\hyperref[\detokenize{restart::doc}]{\sphinxcrossref{\DUrole{doc}{binary restart files}}}}, so pair\_style and pair\_coeff commands do not need
to be specified in an input script that reads a restart file.

\sphinxAtStartPar
These pair styles can only be used via the \sphinxstyleemphasis{pair} keyword of the
{\hyperref[\detokenize{run_style::doc}]{\sphinxcrossref{\DUrole{doc}{run\_style respa}}}} command.  They do not support the
\sphinxstyleemphasis{inner}, \sphinxstyleemphasis{middle}, \sphinxstyleemphasis{outer} keywords.


\bigskip\hrule\bigskip



\subsection{Restrictions}
\label{\detokenize{pair_eff:restrictions}}
\sphinxAtStartPar
These pair styles will only be enabled if LAMMPS is built with the
EFF package.  It will only be enabled if LAMMPS was built with
that package.  See the {\hyperref[\detokenize{Build_package::doc}]{\sphinxcrossref{\DUrole{doc}{Build package}}}} page for
more info.

\sphinxAtStartPar
These pair styles require that particles store electron attributes
such as radius, radial velocity, and radial force, as defined by the
{\hyperref[\detokenize{atom_style::doc}]{\sphinxcrossref{\DUrole{doc}{atom\_style}}}}.  The \sphinxstyleemphasis{electron} atom style does all of
this.

\sphinxAtStartPar
Thes pair styles require you to use the {\hyperref[\detokenize{comm_modify::doc}]{\sphinxcrossref{\DUrole{doc}{comm\_modify vel yes}}}} command so that velocities are stored by ghost
atoms.


\subsection{Related commands}
\label{\detokenize{pair_eff:related-commands}}
\sphinxAtStartPar
{\hyperref[\detokenize{pair_coeff::doc}]{\sphinxcrossref{\DUrole{doc}{pair\_coeff}}}}


\subsection{Default}
\label{\detokenize{pair_eff:default}}
\sphinxAtStartPar
If not specified, limit\_eradius = 0 and pressure\_with\_evirials = 0.


\bigskip\hrule\bigskip

\phantomsection\label{\detokenize{pair_eff:su}}
\sphinxAtStartPar
\sphinxstylestrong{(Su)} Su and Goddard, Excited Electron Dynamics Modeling of Warm
Dense Matter, Phys Rev Lett, 99:185003 (2007).

\phantomsection\label{\detokenize{pair_eff:jaramillo-botero}}
\sphinxAtStartPar
\sphinxstylestrong{(Jaramillo\sphinxhyphen{}Botero)} Jaramillo\sphinxhyphen{}Botero, Su, Qi, Goddard, Large\sphinxhyphen{}scale,
Long\sphinxhyphen{}term Non\sphinxhyphen{}adiabatic Electron Molecular Dynamics for Describing
Material Properties and Phenomena in Extreme Environments, J Comp
Chem, 32, 497\sphinxhyphen{}512 (2011).

\sphinxstepscope

\index{pair\_style eim@\spxentry{pair\_style eim}}\ignorespaces 
\index{pair\_style eim/omp@\spxentry{pair\_style eim/omp}}\ignorespaces 

\section{pair\_style eim command}
\label{\detokenize{pair_eim:pair-style-eim-command}}\label{\detokenize{pair_eim:index-1}}\label{\detokenize{pair_eim:index-0}}\label{\detokenize{pair_eim::doc}}
\sphinxAtStartPar
Accelerator Variants: \sphinxstyleemphasis{eim/omp}


\subsection{Syntax}
\label{\detokenize{pair_eim:syntax}}
\begin{sphinxVerbatim}[commandchars=\\\{\}]
\PYG{k}{pair\PYGZus{}style}\PYG{+w}{ }\PYG{n}{style}
\end{sphinxVerbatim}
\begin{itemize}
\item {} 
\sphinxAtStartPar
style = \sphinxstyleemphasis{eim}

\end{itemize}


\subsection{Examples}
\label{\detokenize{pair_eim:examples}}
\begin{sphinxVerbatim}[commandchars=\\\{\}]
\PYG{k}{pair\PYGZus{}style}\PYG{+w}{ }\PYG{n}{eim}
\PYG{k}{pair\PYGZus{}coeff}\PYG{+w}{ }\PYG{o}{*}\PYG{+w}{ }\PYG{o}{*}\PYG{+w}{ }\PYG{n}{Na}\PYG{+w}{ }\PYG{n}{Cl}\PYG{+w}{ }\PYG{n}{..}\PYG{o}{/}\PYG{n}{potentials}\PYG{o}{/}\PYG{n}{ffield.eim}\PYG{+w}{ }\PYG{n}{Na}\PYG{+w}{ }\PYG{n}{Cl}
\PYG{k}{pair\PYGZus{}coeff}\PYG{+w}{ }\PYG{o}{*}\PYG{+w}{ }\PYG{o}{*}\PYG{+w}{ }\PYG{n}{Na}\PYG{+w}{ }\PYG{n}{Cl}\PYG{+w}{ }\PYG{n}{ffield.eim}\PYG{+w}{  }\PYG{n}{Na}\PYG{+w}{ }\PYG{n}{Na}\PYG{+w}{ }\PYG{n}{Na}\PYG{+w}{ }\PYG{n}{Cl}
\PYG{k}{pair\PYGZus{}coeff}\PYG{+w}{ }\PYG{o}{*}\PYG{+w}{ }\PYG{o}{*}\PYG{+w}{ }\PYG{n}{Na}\PYG{+w}{ }\PYG{n}{Cl}\PYG{+w}{ }\PYG{n}{..}\PYG{o}{/}\PYG{n}{potentials}\PYG{o}{/}\PYG{n}{ffield.eim}\PYG{+w}{ }\PYG{n}{Cl}\PYG{+w}{ }\PYG{n}{NULL}\PYG{+w}{ }\PYG{n}{Na}
\end{sphinxVerbatim}


\subsection{Description}
\label{\detokenize{pair_eim:description}}
\sphinxAtStartPar
Style \sphinxstyleemphasis{eim} computes pairwise interactions for ionic compounds
using embedded\sphinxhyphen{}ion method (EIM) potentials {\hyperref[\detokenize{pair_eim:zhou2}]{\sphinxcrossref{\DUrole{std,std-ref}{(Zhou)}}}}.  The
energy of the system E is given by
\begin{equation*}
\begin{split}E = \frac{1}{2} \sum_{i=1}^{N} \sum_{j=i_1}^{i_N} \phi_{ij} \left(r_{ij}\right) + \sum_{i=1}^{N}E_i\left(q_i,\sigma_i\right)\end{split}
\end{equation*}
\sphinxAtStartPar
The first term is a double pairwise sum over the J neighbors of all I
atoms, where \(\phi_{ij}\) is a pair potential.  The second term sums over
the embedding energy E\_i of atom I, which is a function of its charge
q\_i and the electrical potential \(\sigma_i\) at its location.  E\_i, q\_i,
and \(sigma_i\) are calculated as
\begin{equation*}
\begin{split}q_i  = & \sum_{j=i_1}^{i_N} \eta_{ji}\left(r_{ij}\right) \\
\sigma_i  = & \sum_{j=i_1}^{i_N} q_j \cdot \psi_{ij} \left(r_{ij}\right) \\
E_i\left(q_i,\sigma_i\right)  = & \frac{1}{2} \cdot q_i \cdot \sigma_i\end{split}
\end{equation*}
\sphinxAtStartPar
where \(\eta_{ji}\) is a pairwise function describing electron flow from atom
I to atom J, and \(\psi_{ij}\) is another pairwise function.  The multi\sphinxhyphen{}body
nature of the EIM potential is a result of the embedding energy term.
A complete list of all the pair functions used in EIM is summarized
below
\begin{equation*}
\begin{split}\phi_{ij}\left(r\right) = & \left\{ \begin{array}{lr}
\left[\frac{E_{b,ij}\beta_{ij}}{\beta_{ij}-\alpha_{ij}}\exp\left(-\alpha_{ij} \frac{r-r_{e,ij}}{r_{e,ij}}\right)-\frac{E_{b,ij}\alpha_{ij}}{\beta_{ij}-\alpha_{ij}}\exp\left(-\beta_{ij} \frac{r-r_{e,ij}}{r_{e,ij}}\right)\right]f_c\left(r,r_{e,ij},r_{c,\phi,ij}\right),& p_{ij}=1 \\
\left[\frac{E_{b,ij}\beta_{ij}}{\beta_{ij}-\alpha_{ij}} \left(\frac{r_{e,ij}}{r}\right)^{\alpha_{ij}}  -\frac{E_{b,ij}\alpha_{ij}}{\beta_{ij}-\alpha_{ij}} \left(\frac{r_{e,ij}}{r}\right)^{\beta_{ij}}\right]f_c\left(r,r_{e,ij},r_{c,\phi,ij}\right),& p_{ij}=2
\end{array}
\right.\\
\eta_{ji} = & A_{\eta,ij}\left(\chi_j-\chi_i\right)f_c\left(r,r_{s,\eta,ij},r_{c,\eta,ij}\right) \\
\psi_{ij}\left(r\right) = & A_{\psi,ij}\exp\left(-\zeta_{ij}r\right)f_c\left(r,r_{s,\psi,ij},r_{c,\psi,ij}\right) \\
f_{c}\left(r,r_p,r_c\right) = & 0.510204 \cdot \mathrm{erfc}\left[\frac{1.64498\left(2r-r_p-r_c\right)}{r_c-r_p}\right] - 0.010204\end{split}
\end{equation*}
\sphinxAtStartPar
Here \(E_b, r_e, r_(c,\phi), \alpha, \beta, A_(\psi), \zeta, r_(s,\psi),
r_(c,\psi), A_(\eta), r_(s,\eta), r_(c,\eta), \chi,\) and pair function type
\sphinxstyleemphasis{p} are parameters, with subscripts \sphinxstyleemphasis{ij} indicating the two species of
atoms in the atomic pair.

\begin{sphinxadmonition}{note}{Note:}
\sphinxAtStartPar
Even though the EIM potential is treating atoms as charged ions,
you should not use a LAMMPS {\hyperref[\detokenize{atom_style::doc}]{\sphinxcrossref{\DUrole{doc}{atom\_style}}}} that stores a
charge on each atom and thus requires you to assign a charge to each
atom, e.g. the \sphinxstyleemphasis{charge} or \sphinxstyleemphasis{full} atom styles.  This is because the
EIM potential infers the charge on an atom from the equation above for
q\_i; you do not assign charges explicitly.
\end{sphinxadmonition}


\bigskip\hrule\bigskip


\sphinxAtStartPar
All the EIM parameters are listed in a potential file which is
specified by the {\hyperref[\detokenize{pair_coeff::doc}]{\sphinxcrossref{\DUrole{doc}{pair\_coeff}}}} command.  This is an
ASCII text file in a format described below.  The “ffield.eim” file
included in the “potentials” directory of the LAMMPS distribution
currently includes nine elements Li, Na, K, Rb, Cs, F, Cl, Br, and I.
A system with any combination of these elements can be modeled.  This
file is parameterized in terms of LAMMPS {\hyperref[\detokenize{units::doc}]{\sphinxcrossref{\DUrole{doc}{metal units}}}}.

\sphinxAtStartPar
Note that unlike other potentials, cutoffs for EIM potentials are not
set in the pair\_style or pair\_coeff command; they are specified in the
EIM potential file itself.  Likewise, the EIM potential file lists
atomic masses; thus you do not need to use the {\hyperref[\detokenize{mass::doc}]{\sphinxcrossref{\DUrole{doc}{mass}}}}
command to specify them.

\sphinxAtStartPar
Only a single pair\_coeff command is used with the \sphinxstyleemphasis{eim} style which
specifies an EIM potential file and the element(s) to extract
information for.  The EIM elements are mapped to LAMMPS atom types by
specifying N additional arguments after the filename in the pair\_coeff
command, where N is the number of LAMMPS atom types:
\begin{itemize}
\item {} 
\sphinxAtStartPar
Elem1, Elem2, …

\item {} 
\sphinxAtStartPar
EIM potential file

\item {} 
\sphinxAtStartPar
N element names = mapping of EIM elements to atom types

\end{itemize}

\sphinxAtStartPar
See the {\hyperref[\detokenize{pair_coeff::doc}]{\sphinxcrossref{\DUrole{doc}{pair\_coeff}}}} page for alternate ways
to specify the path for the potential file.

\sphinxAtStartPar
As an example like one of those above, suppose you want to model a
system with Na and Cl atoms.  If your LAMMPS simulation has 4 atoms
types and you want the first 3 to be Na, and the fourth to be Cl, you would
use the following pair\_coeff command:

\begin{sphinxVerbatim}[commandchars=\\\{\}]
\PYG{k}{pair\PYGZus{}coeff}\PYG{+w}{ }\PYG{o}{*}\PYG{+w}{ }\PYG{o}{*}\PYG{+w}{ }\PYG{n}{Na}\PYG{+w}{ }\PYG{n}{Cl}\PYG{+w}{ }\PYG{n}{ffield.eim}\PYG{+w}{ }\PYG{n}{Na}\PYG{+w}{ }\PYG{n}{Na}\PYG{+w}{ }\PYG{n}{Na}\PYG{+w}{ }\PYG{n}{Cl}
\end{sphinxVerbatim}

\sphinxAtStartPar
The first 2 arguments must be * * so as to span all LAMMPS atom types.
The filename is the EIM potential file.  The Na and Cl arguments
(before the file name) are the two elements for which info will be
extracted from the potential file.  The first three trailing Na
arguments map LAMMPS atom types 1,2,3 to the EIM Na element.  The
final Cl argument maps LAMMPS atom type 4 to the EIM Cl element.

\sphinxAtStartPar
If a mapping value is specified as NULL, the mapping is not performed.
This can be used when an \sphinxstyleemphasis{eim} potential is used as part of the
\sphinxstyleemphasis{hybrid} pair style.  The NULL values are placeholders for atom types
that will be used with other potentials.

\sphinxAtStartPar
The ffield.eim file in the \sphinxstyleemphasis{potentials} directory of the LAMMPS
distribution is formatted as follows:

\sphinxAtStartPar
Lines starting with \# are comments and are ignored by LAMMPS.  Lines
starting with “global:” include three global values. The first value
divides the cations from anions, i.e., any elements with
electronegativity above this value are viewed as anions, and any
elements with electronegativity below this value are viewed as
cations. The second and third values are related to the cutoff
function \sphinxhyphen{} i.e. the 0.510204, 1.64498, and 0.010204 shown in the above
equation can be derived from these values.

\sphinxAtStartPar
Lines starting with “element:” are formatted as follows: name of
element, atomic number, atomic mass, electronic negativity, atomic
radius (LAMMPS ignores it), ionic radius (LAMMPS ignores it), cohesive
energy (LAMMPS ignores it), and q0 (must be 0).

\sphinxAtStartPar
Lines starting with “pair:” are entered as: element 1, element 2,
r\_(c,phi), r\_(c,phi) (redundant for historical reasons), E\_b, r\_e,
alpha, beta, r\_(c,eta), A\_(eta), r\_(s,eta), r\_(c,psi), A\_(psi), zeta,
r\_(s,psi), and p.

\sphinxAtStartPar
The lines in the file can be in any order; LAMMPS extracts the info it
needs.


\bigskip\hrule\bigskip


\sphinxAtStartPar
Styles with a \sphinxstyleemphasis{gpu}, \sphinxstyleemphasis{intel}, \sphinxstyleemphasis{kk}, \sphinxstyleemphasis{omp}, or \sphinxstyleemphasis{opt} suffix are
functionally the same as the corresponding style without the suffix.
They have been optimized to run faster, depending on your available
hardware, as discussed on the {\hyperref[\detokenize{Speed_packages::doc}]{\sphinxcrossref{\DUrole{doc}{Accelerator packages}}}}
page.  The accelerated styles take the same arguments and should
produce the same results, except for round\sphinxhyphen{}off and precision issues.

\sphinxAtStartPar
These accelerated styles are part of the GPU, INTEL, KOKKOS,
OPENMP, and OPT packages, respectively.  They are only enabled if
LAMMPS was built with those packages.  See the {\hyperref[\detokenize{Build_package::doc}]{\sphinxcrossref{\DUrole{doc}{Build package}}}} page for more info.

\sphinxAtStartPar
You can specify the accelerated styles explicitly in your input script
by including their suffix, or you can use the {\hyperref[\detokenize{Run_options::doc}]{\sphinxcrossref{\DUrole{doc}{\sphinxhyphen{}suffix command\sphinxhyphen{}line switch}}}} when you invoke LAMMPS, or you can use the
{\hyperref[\detokenize{suffix::doc}]{\sphinxcrossref{\DUrole{doc}{suffix}}}} command in your input script.

\sphinxAtStartPar
See the {\hyperref[\detokenize{Speed_packages::doc}]{\sphinxcrossref{\DUrole{doc}{Accelerator packages}}}} page for more
instructions on how to use the accelerated styles effectively.


\bigskip\hrule\bigskip



\subsection{Restrictions}
\label{\detokenize{pair_eim:restrictions}}
\sphinxAtStartPar
This style is part of the MANYBODY package.  It is only enabled if
LAMMPS was built with that package.


\subsection{Related commands}
\label{\detokenize{pair_eim:related-commands}}
\sphinxAtStartPar
{\hyperref[\detokenize{pair_coeff::doc}]{\sphinxcrossref{\DUrole{doc}{pair\_coeff}}}}


\subsection{Default}
\label{\detokenize{pair_eim:default}}
\sphinxAtStartPar
none


\bigskip\hrule\bigskip

\phantomsection\label{\detokenize{pair_eim:zhou2}}
\sphinxAtStartPar
\sphinxstylestrong{(Zhou)} Zhou, submitted for publication (2010).  Please contact
Xiaowang Zhou (Sandia) for details via email at xzhou at sandia.gov.

\sphinxstepscope

\index{pair\_style exp6/rx@\spxentry{pair\_style exp6/rx}}\ignorespaces 
\index{pair\_style exp6/rx/kk@\spxentry{pair\_style exp6/rx/kk}}\ignorespaces 

\section{pair\_style exp6/rx command}
\label{\detokenize{pair_exp6_rx:pair-style-exp6-rx-command}}\label{\detokenize{pair_exp6_rx:index-1}}\label{\detokenize{pair_exp6_rx:index-0}}\label{\detokenize{pair_exp6_rx::doc}}
\sphinxAtStartPar
Accelerator Variants: \sphinxstyleemphasis{exp6/rx/kk}


\subsection{Syntax}
\label{\detokenize{pair_exp6_rx:syntax}}
\begin{sphinxVerbatim}[commandchars=\\\{\}]
\PYG{k}{pair\PYGZus{}style}\PYG{+w}{ }\PYG{n}{exp6}\PYG{o}{/}\PYG{n}{rx}\PYG{+w}{ }\PYG{n}{cutoff}\PYG{+w}{ }\PYG{n}{...}
\end{sphinxVerbatim}
\begin{itemize}
\item {} 
\sphinxAtStartPar
cutoff = global cutoff for DPD interactions (distance units)

\item {} 
\sphinxAtStartPar
weighting = fractional or molecular (optional)

\end{itemize}


\subsection{Examples}
\label{\detokenize{pair_exp6_rx:examples}}
\begin{sphinxVerbatim}[commandchars=\\\{\}]
\PYG{k}{pair\PYGZus{}style}\PYG{+w}{ }\PYG{n}{exp6}\PYG{o}{/}\PYG{n}{rx}\PYG{+w}{ }\PYG{l+m}{10.0}
\PYG{k}{pair\PYGZus{}style}\PYG{+w}{ }\PYG{n}{exp6}\PYG{o}{/}\PYG{n}{rx}\PYG{+w}{ }\PYG{l+m}{10.0}\PYG{+w}{ }\PYG{n}{fractional}
\PYG{k}{pair\PYGZus{}style}\PYG{+w}{ }\PYG{n}{exp6}\PYG{o}{/}\PYG{n}{rx}\PYG{+w}{ }\PYG{l+m}{10.0}\PYG{+w}{ }\PYG{n}{molecular}
\PYG{k}{pair\PYGZus{}coeff}\PYG{+w}{ }\PYG{o}{*}\PYG{+w}{ }\PYG{o}{*}\PYG{+w}{ }\PYG{n}{exp6.params}\PYG{+w}{ }\PYG{n}{h2o}\PYG{+w}{ }\PYG{n}{h2o}\PYG{+w}{ }\PYG{n}{exponent}\PYG{+w}{ }\PYG{l+m}{1.0}\PYG{+w}{ }\PYG{l+m}{1.0}\PYG{+w}{ }\PYG{l+m}{10.0}
\PYG{k}{pair\PYGZus{}coeff}\PYG{+w}{ }\PYG{o}{*}\PYG{+w}{ }\PYG{o}{*}\PYG{+w}{ }\PYG{n}{exp6.params}\PYG{+w}{ }\PYG{n}{h2o}\PYG{+w}{ }\PYG{l+m}{1}\PYG{n}{fluid}\PYG{+w}{ }\PYG{n}{exponent}\PYG{+w}{ }\PYG{l+m}{1.0}\PYG{+w}{ }\PYG{l+m}{1.0}\PYG{+w}{ }\PYG{l+m}{10.0}
\PYG{k}{pair\PYGZus{}coeff}\PYG{+w}{ }\PYG{o}{*}\PYG{+w}{ }\PYG{o}{*}\PYG{+w}{ }\PYG{n}{exp6.params}\PYG{+w}{ }\PYG{l+m}{1}\PYG{n}{fluid}\PYG{+w}{ }\PYG{l+m}{1}\PYG{n}{fluid}\PYG{+w}{ }\PYG{n}{exponent}\PYG{+w}{ }\PYG{l+m}{1.0}\PYG{+w}{ }\PYG{l+m}{1.0}\PYG{+w}{ }\PYG{l+m}{10.0}
\PYG{k}{pair\PYGZus{}coeff}\PYG{+w}{ }\PYG{o}{*}\PYG{+w}{ }\PYG{o}{*}\PYG{+w}{ }\PYG{n}{exp6.params}\PYG{+w}{ }\PYG{l+m}{1}\PYG{n}{fluid}\PYG{+w}{ }\PYG{l+m}{1}\PYG{n}{fluid}\PYG{+w}{ }\PYG{n}{none}\PYG{+w}{ }\PYG{l+m}{10.0}
\PYG{k}{pair\PYGZus{}coeff}\PYG{+w}{ }\PYG{o}{*}\PYG{+w}{ }\PYG{o}{*}\PYG{+w}{ }\PYG{n}{exp6.params}\PYG{+w}{ }\PYG{l+m}{1}\PYG{n}{fluid}\PYG{+w}{ }\PYG{l+m}{1}\PYG{n}{fluid}\PYG{+w}{ }\PYG{n}{polynomial}\PYG{+w}{ }\PYG{n}{filename}\PYG{+w}{ }\PYG{l+m}{10.0}
\end{sphinxVerbatim}


\subsection{Description}
\label{\detokenize{pair_exp6_rx:description}}
\sphinxAtStartPar
Style \sphinxstyleemphasis{exp6/rx} is used in reaction DPD simulations, where the
coarse\sphinxhyphen{}grained (CG) particles are composed of \sphinxstyleemphasis{m} species whose
reaction rate kinetics are determined from a set of \sphinxstyleemphasis{n} reaction rate
equations through the {\hyperref[\detokenize{fix_rx::doc}]{\sphinxcrossref{\DUrole{doc}{fix rx}}}} command.  The species of
one CG particle can interact with a species in a neighboring CG
particle through a site\sphinxhyphen{}site interaction potential model.  The
\sphinxstyleemphasis{exp6/rx} style computes an exponential\sphinxhyphen{}6 potential given by
\begin{equation*}
\begin{split}U_{ij}(r) = \frac{\epsilon}{\alpha-6}\{6\exp[\alpha(1-\frac{r_{ij}}{R_{m}})]-\alpha(\frac{R_{m}}{r_{ij}})^6\}\end{split}
\end{equation*}
\sphinxAtStartPar
where the \(\epsilon\) parameter determines the depth of the
potential minimum located at \(R_m\), and \(\alpha\) determines
the softness of the repulsion.

\sphinxAtStartPar
The coefficients must be defined for each species in a given particle
type via the {\hyperref[\detokenize{pair_coeff::doc}]{\sphinxcrossref{\DUrole{doc}{pair\_coeff}}}} command as in the examples
above, where the first argument is the filename that includes the
exponential\sphinxhyphen{}6 parameters for each species.  The file includes the
species tag followed by the \(\alpha, \epsilon\) and \(R_m\)
parameters. The format of the file is described below.

\sphinxAtStartPar
The second and third arguments specify the site\sphinxhyphen{}site interaction
potential between two species contained within two different
particles.  The species tags must either correspond to the species
defined in the reaction kinetics files specified with the {\hyperref[\detokenize{fix_rx::doc}]{\sphinxcrossref{\DUrole{doc}{fix rx}}}} command or they must correspond to the tag “1fluid”,
signifying interaction with a product species mixture determined
through a one\sphinxhyphen{}fluid approximation.  The interaction potential is
weighted by the geometric average of either the mole fraction concentrations
or the number of molecules associated with the interacting coarse\sphinxhyphen{}grained
particles (see the \sphinxstyleemphasis{fractional} or \sphinxstyleemphasis{molecular} weighting pair style options).
The coarse\sphinxhyphen{}grained potential is stored before and after the
reaction kinetics solver is applied, where the difference is defined
to be the internal chemical energy (uChem).

\sphinxAtStartPar
The fourth argument specifies the type of scaling that will be used
to scale the EXP\sphinxhyphen{}6 parameters as reactions occur.  Currently, there
are three scaling options:  \sphinxstyleemphasis{exponent}, \sphinxstyleemphasis{polynomial} and \sphinxstyleemphasis{none}.

\sphinxAtStartPar
Exponent scaling requires two additional arguments for scaling
the \(R_m\) and \(\epsilon\) parameters, respectively.  The scaling factor
is computed by phi\textasciicircum{}exponent, where phi is the number of molecules
represented by the coarse\sphinxhyphen{}grain particle and exponent is specified
as a pair coefficient argument for \(R_m\) and \(\epsilon\), respectively.
The \(R_m\) and \(\epsilon\) parameters are multiplied by the scaling
factor to give the scaled interaction parameters for the CG particle.

\sphinxAtStartPar
Polynomial scaling requires a filename to be specified as a pair
coeff argument.  The file contains the coefficients to a fifth order
polynomial for the \(\alpha\), \(\epsilon\) and \(R_m\) parameters that depend
upon phi (the number of molecules represented by the CG particle).
The format of a polynomial file is provided below.

\sphinxAtStartPar
The \sphinxstyleemphasis{none} option to the scaling does not have any additional pair coeff
arguments.  This is equivalent to specifying the \sphinxstyleemphasis{exponent} option with
\(R_m\) and \(\epsilon\) exponents of 0.0 and 0.0, respectively.

\sphinxAtStartPar
The final argument specifies the interaction cutoff (optional).


\bigskip\hrule\bigskip


\sphinxAtStartPar
The format of a tabulated file is as follows (without the
parenthesized comments):

\begin{sphinxVerbatim}[commandchars=\\\{\}]
\PYGZsh{} exponential\PYGZhy{}6 parameters for various species      (one or more comment or blank lines)

h2o  exp6  11.00 0.02 3.50                          (species, exp6, alpha, Rm, epsilon)
no2  exp6  13.60 0.01 3.70
...
co2  exp6  13.00 0.03 3.20
\end{sphinxVerbatim}

\sphinxAtStartPar
The format of the polynomial scaling file as follows (without the
parenthesized comments):
\begin{sphinxalltt}
\# POLYNOMIAL FILE          (one or more comment or blank lines)

\#  General Functional Form:
\#  A*phi\textasciicircum{}5 + B*phi\textasciicircum{}4 + C*phi\textasciicircum{}3 + D*phi\textasciicircum{}2 + E*phi + F
\#
\#  Parameter  A        B         C        D         E        F
                           (blank)
alpha        0.0000   0.00000   0.00008  0.04955  \sphinxhyphen{}0.73804  13.63201
epsilon      0.0000   0.00478  \sphinxhyphen{}0.06283  0.24486  \sphinxhyphen{}0.33737   2.60097
rm           0.0001  \sphinxhyphen{}0.00118  \sphinxhyphen{}0.00253  0.05812  \sphinxhyphen{}0.00509   1.50106
\end{sphinxalltt}

\sphinxAtStartPar
A section begins with a non\sphinxhyphen{}blank line whose first character is not a
“\#”; blank lines or lines starting with “\#” can be used as comments
between sections.

\sphinxAtStartPar
Following a blank line, the next N lines list the species and their
corresponding parameters.  The first argument is the species tag, the
second argument is the exp6 tag, the third argument is the \(\alpha\)
parameter (energy units), the fourth argument is the \(\epsilon\) parameter
(energy\sphinxhyphen{}distance\textasciicircum{}6 units), and the fifth argument is the \(R_m\) parameter
(distance units).  If a species tag of “1fluid” is listed as a pair
coefficient, a one\sphinxhyphen{}fluid approximation is specified where a
concentration\sphinxhyphen{}dependent combination of the parameters is computed
through the following equations:
\begin{equation*}
\begin{split}R_{m}^{3} = & \sum_{a}\sum_{b} x_{a}x_{b}R_{m,ab}^{3} \\
\epsilon  = & \frac{1}{R_{m}^{3}}\sum_{a}\sum_{b} x_{a}x_{b}\epsilon_{ab}R_{m,ab}^{3} \\
\alpha    = & \frac{1}{\epsilon R_{m}^{3}}\sum_{a}\sum_{b} x_{a}x_{b}\alpha_{ab}\epsilon_{ab}R_{m,ab}^{3}\end{split}
\end{equation*}
\sphinxAtStartPar
where
\begin{equation*}
\begin{split}\epsilon_{ab} = & \sqrt{\epsilon_{a}\epsilon_{b}} \\
R_{m,ab}      = & \frac{R_{m,a}+R_{m,b}}{2} \\
\alpha_{ab}   = & \sqrt{\alpha_{a}\alpha_{b}}\end{split}
\end{equation*}
\sphinxAtStartPar
and \(x_a\) and \(x_b\) are the mole fractions of a and b, respectively, which
comprise the gas mixture.


\bigskip\hrule\bigskip



\subsection{Mixing, shift, table, tail correction, restart, rRESPA info}
\label{\detokenize{pair_exp6_rx:mixing-shift-table-tail-correction-restart-rrespa-info}}
\sphinxAtStartPar
This pair style does not support mixing.  Thus, coefficients for all
I,J pairs must be specified explicitly.

\sphinxAtStartPar
This style does not support the {\hyperref[\detokenize{pair_modify::doc}]{\sphinxcrossref{\DUrole{doc}{pair\_modify}}}} shift option
for the energy of the exp() and 1/r\textasciicircum{}6 portion of the pair interaction.

\sphinxAtStartPar
This style does not support the pair\_modify tail option for adding long\sphinxhyphen{}range
tail corrections to energy and pressure for the A,C terms in the
pair interaction.


\bigskip\hrule\bigskip


\sphinxAtStartPar
Styles with a \sphinxstyleemphasis{gpu}, \sphinxstyleemphasis{intel}, \sphinxstyleemphasis{kk}, \sphinxstyleemphasis{omp}, or \sphinxstyleemphasis{opt} suffix are
functionally the same as the corresponding style without the suffix.
They have been optimized to run faster, depending on your available
hardware, as discussed on the {\hyperref[\detokenize{Speed_packages::doc}]{\sphinxcrossref{\DUrole{doc}{Accelerator packages}}}}
page.  The accelerated styles take the same arguments and should
produce the same results, except for round\sphinxhyphen{}off and precision issues.

\sphinxAtStartPar
These accelerated styles are part of the GPU, INTEL, KOKKOS,
OPENMP, and OPT packages, respectively.  They are only enabled if
LAMMPS was built with those packages.  See the {\hyperref[\detokenize{Build_package::doc}]{\sphinxcrossref{\DUrole{doc}{Build package}}}} page for more info.

\sphinxAtStartPar
You can specify the accelerated styles explicitly in your input script
by including their suffix, or you can use the {\hyperref[\detokenize{Run_options::doc}]{\sphinxcrossref{\DUrole{doc}{\sphinxhyphen{}suffix command\sphinxhyphen{}line switch}}}} when you invoke LAMMPS, or you can use the
{\hyperref[\detokenize{suffix::doc}]{\sphinxcrossref{\DUrole{doc}{suffix}}}} command in your input script.

\sphinxAtStartPar
See the {\hyperref[\detokenize{Speed_packages::doc}]{\sphinxcrossref{\DUrole{doc}{Accelerator packages}}}} page for more
instructions on how to use the accelerated styles effectively.


\bigskip\hrule\bigskip



\subsection{Restrictions}
\label{\detokenize{pair_exp6_rx:restrictions}}
\sphinxAtStartPar
This command is part of the DPD\sphinxhyphen{}REACT package.  It is only enabled if
LAMMPS was built with that package.  See the {\hyperref[\detokenize{Build_package::doc}]{\sphinxcrossref{\DUrole{doc}{Build package}}}} page for more info.


\subsection{Related commands}
\label{\detokenize{pair_exp6_rx:related-commands}}
\sphinxAtStartPar
{\hyperref[\detokenize{pair_coeff::doc}]{\sphinxcrossref{\DUrole{doc}{pair\_coeff}}}}


\subsection{Default}
\label{\detokenize{pair_exp6_rx:default}}
\sphinxAtStartPar
fractional weighting

\sphinxstepscope

\index{pair\_style extep@\spxentry{pair\_style extep}}\ignorespaces 

\section{pair\_style extep command}
\label{\detokenize{pair_extep:pair-style-extep-command}}\label{\detokenize{pair_extep:index-0}}\label{\detokenize{pair_extep::doc}}

\subsection{Syntax}
\label{\detokenize{pair_extep:syntax}}
\begin{sphinxVerbatim}[commandchars=\\\{\}]
\PYG{k}{pair\PYGZus{}style}\PYG{+w}{ }\PYG{n}{extep}
\end{sphinxVerbatim}


\subsection{Examples}
\label{\detokenize{pair_extep:examples}}
\begin{sphinxVerbatim}[commandchars=\\\{\}]
\PYG{k}{pair\PYGZus{}style}\PYG{+w}{ }\PYG{n}{extep}
\PYG{k}{pair\PYGZus{}coeff}\PYG{+w}{ }\PYG{o}{*}\PYG{+w}{ }\PYG{o}{*}\PYG{+w}{ }\PYG{n}{BN.extep}\PYG{+w}{ }\PYG{n}{B}\PYG{+w}{ }\PYG{n}{N}
\end{sphinxVerbatim}


\subsection{Description}
\label{\detokenize{pair_extep:description}}
\sphinxAtStartPar
Style \sphinxstyleemphasis{extep} computes the Extended Tersoff Potential (ExTeP)
interactions as described in {\hyperref[\detokenize{pair_extep:los2017}]{\sphinxcrossref{\DUrole{std,std-ref}{(Los2017)}}}}.


\bigskip\hrule\bigskip



\subsection{Restrictions}
\label{\detokenize{pair_extep:restrictions}}
\sphinxAtStartPar
none


\subsection{Related commands}
\label{\detokenize{pair_extep:related-commands}}
\sphinxAtStartPar
{\hyperref[\detokenize{pair_tersoff::doc}]{\sphinxcrossref{\DUrole{doc}{pair\_tersoff}}}}


\subsection{Default}
\label{\detokenize{pair_extep:default}}
\sphinxAtStartPar
none


\bigskip\hrule\bigskip

\phantomsection\label{\detokenize{pair_extep:los2017}}
\sphinxAtStartPar
\sphinxstylestrong{(Los2017)} J. H. Los et al. “Extended Tersoff potential for boron nitride:
Energetics and elastic properties of pristine and defective h\sphinxhyphen{}BN”,
Phys. Rev. B 96 (184108), 2017.

\sphinxstepscope

\index{pair\_style lj/cut/soft@\spxentry{pair\_style lj/cut/soft}}\ignorespaces 
\index{pair\_style lj/cut/soft/omp@\spxentry{pair\_style lj/cut/soft/omp}}\ignorespaces 
\index{pair\_style lj/cut/coul/cut/soft@\spxentry{pair\_style lj/cut/coul/cut/soft}}\ignorespaces 
\index{pair\_style lj/cut/coul/cut/soft/gpu@\spxentry{pair\_style lj/cut/coul/cut/soft/gpu}}\ignorespaces 
\index{pair\_style lj/cut/coul/cut/soft/omp@\spxentry{pair\_style lj/cut/coul/cut/soft/omp}}\ignorespaces 
\index{pair\_style lj/cut/coul/long/soft@\spxentry{pair\_style lj/cut/coul/long/soft}}\ignorespaces 
\index{pair\_style lj/cut/coul/long/soft/gpu@\spxentry{pair\_style lj/cut/coul/long/soft/gpu}}\ignorespaces 
\index{pair\_style lj/cut/coul/long/soft/omp@\spxentry{pair\_style lj/cut/coul/long/soft/omp}}\ignorespaces 
\index{pair\_style lj/cut/tip4p/long/soft@\spxentry{pair\_style lj/cut/tip4p/long/soft}}\ignorespaces 
\index{pair\_style lj/cut/tip4p/long/soft/omp@\spxentry{pair\_style lj/cut/tip4p/long/soft/omp}}\ignorespaces 
\index{pair\_style lj/charmm/coul/long/soft@\spxentry{pair\_style lj/charmm/coul/long/soft}}\ignorespaces 
\index{pair\_style lj/charmm/coul/long/soft/omp@\spxentry{pair\_style lj/charmm/coul/long/soft/omp}}\ignorespaces 
\index{pair\_style lj/class2/soft@\spxentry{pair\_style lj/class2/soft}}\ignorespaces 
\index{pair\_style lj/class2/coul/cut/soft@\spxentry{pair\_style lj/class2/coul/cut/soft}}\ignorespaces 
\index{pair\_style lj/class2/coul/long/soft@\spxentry{pair\_style lj/class2/coul/long/soft}}\ignorespaces 
\index{pair\_style coul/cut/soft@\spxentry{pair\_style coul/cut/soft}}\ignorespaces 
\index{pair\_style coul/cut/soft/omp@\spxentry{pair\_style coul/cut/soft/omp}}\ignorespaces 
\index{pair\_style coul/long/soft@\spxentry{pair\_style coul/long/soft}}\ignorespaces 
\index{pair\_style coul/long/soft/omp@\spxentry{pair\_style coul/long/soft/omp}}\ignorespaces 
\index{pair\_style tip4p/long/soft@\spxentry{pair\_style tip4p/long/soft}}\ignorespaces 
\index{pair\_style tip4p/long/soft/omp@\spxentry{pair\_style tip4p/long/soft/omp}}\ignorespaces 
\index{pair\_style morse/soft@\spxentry{pair\_style morse/soft}}\ignorespaces 

\section{pair\_style lj/cut/soft command}
\label{\detokenize{pair_fep_soft:pair-style-lj-cut-soft-command}}\label{\detokenize{pair_fep_soft:index-21}}\label{\detokenize{pair_fep_soft:index-20}}\label{\detokenize{pair_fep_soft:index-19}}\label{\detokenize{pair_fep_soft:index-18}}\label{\detokenize{pair_fep_soft:index-17}}\label{\detokenize{pair_fep_soft:index-16}}\label{\detokenize{pair_fep_soft:index-15}}\label{\detokenize{pair_fep_soft:index-14}}\label{\detokenize{pair_fep_soft:index-13}}\label{\detokenize{pair_fep_soft:index-12}}\label{\detokenize{pair_fep_soft:index-11}}\label{\detokenize{pair_fep_soft:index-10}}\label{\detokenize{pair_fep_soft:index-9}}\label{\detokenize{pair_fep_soft:index-8}}\label{\detokenize{pair_fep_soft:index-7}}\label{\detokenize{pair_fep_soft:index-6}}\label{\detokenize{pair_fep_soft:index-5}}\label{\detokenize{pair_fep_soft:index-4}}\label{\detokenize{pair_fep_soft:index-3}}\label{\detokenize{pair_fep_soft:index-2}}\label{\detokenize{pair_fep_soft:index-1}}\label{\detokenize{pair_fep_soft:index-0}}\label{\detokenize{pair_fep_soft::doc}}
\sphinxAtStartPar
Accelerator Variants: \sphinxstyleemphasis{lj/cut/soft/omp}


\section{pair\_style lj/cut/coul/cut/soft command}
\label{\detokenize{pair_fep_soft:pair-style-lj-cut-coul-cut-soft-command}}
\sphinxAtStartPar
Accelerator Variants: \sphinxstyleemphasis{lj/cut/coul/cut/soft/gpu}, \sphinxstyleemphasis{lj/cut/coul/cut/soft/omp}


\section{pair\_style lj/cut/coul/long/soft command}
\label{\detokenize{pair_fep_soft:pair-style-lj-cut-coul-long-soft-command}}
\sphinxAtStartPar
Accelerator Variants: \sphinxstyleemphasis{lj/cut/coul/long/soft/gpu}, \sphinxstyleemphasis{lj/cut/coul/long/soft/omp}


\section{pair\_style lj/cut/tip4p/long/soft command}
\label{\detokenize{pair_fep_soft:pair-style-lj-cut-tip4p-long-soft-command}}
\sphinxAtStartPar
Accelerator Variants: \sphinxstyleemphasis{lj/cut/tip4p/long/soft/omp}


\section{pair\_style lj/charmm/coul/long/soft command}
\label{\detokenize{pair_fep_soft:pair-style-lj-charmm-coul-long-soft-command}}
\sphinxAtStartPar
Accelerator Variants: \sphinxstyleemphasis{lj/charmm/coul/long/soft/omp}


\section{pair\_style lj/class2/soft command}
\label{\detokenize{pair_fep_soft:pair-style-lj-class2-soft-command}}

\section{pair\_style lj/class2/coul/cut/soft command}
\label{\detokenize{pair_fep_soft:pair-style-lj-class2-coul-cut-soft-command}}

\section{pair\_style lj/class2/coul/long/soft command}
\label{\detokenize{pair_fep_soft:pair-style-lj-class2-coul-long-soft-command}}

\section{pair\_style coul/cut/soft command}
\label{\detokenize{pair_fep_soft:pair-style-coul-cut-soft-command}}
\sphinxAtStartPar
Accelerator Variants: \sphinxstyleemphasis{coul/cut/soft/omp}


\section{pair\_style coul/long/soft command}
\label{\detokenize{pair_fep_soft:pair-style-coul-long-soft-command}}
\sphinxAtStartPar
Accelerator Variants: \sphinxstyleemphasis{coul/long/soft/omp}


\section{pair\_style tip4p/long/soft command}
\label{\detokenize{pair_fep_soft:pair-style-tip4p-long-soft-command}}
\sphinxAtStartPar
Accelerator Variants: \sphinxstyleemphasis{tip4p/long/soft/omp}


\section{pair\_style morse/soft command}
\label{\detokenize{pair_fep_soft:pair-style-morse-soft-command}}

\subsection{Syntax}
\label{\detokenize{pair_fep_soft:syntax}}
\begin{sphinxVerbatim}[commandchars=\\\{\}]
\PYG{k}{pair\PYGZus{}style}\PYG{+w}{ }\PYG{n}{style}\PYG{+w}{ }\PYG{n}{args}
\end{sphinxVerbatim}
\begin{itemize}
\item {} 
\sphinxAtStartPar
style = \sphinxstyleemphasis{lj/cut/soft} or \sphinxstyleemphasis{lj/cut/coul/cut/soft} or \sphinxstyleemphasis{lj/cut/coul/long/soft} or \sphinxstyleemphasis{lj/cut/tip4p/long/soft} or \sphinxstyleemphasis{lj/charmm/coul/long/soft} or \sphinxstyleemphasis{lj/class2/soft} or \sphinxstyleemphasis{lj/class2/coul/cut/soft} or \sphinxstyleemphasis{lj/class2/coul/long/soft} or \sphinxstyleemphasis{coul/cut/soft} or \sphinxstyleemphasis{coul/long/soft} or \sphinxstyleemphasis{tip4p/long/soft} or \sphinxstyleemphasis{morse/soft}

\item {} 
\sphinxAtStartPar
args = list of arguments for a particular style

\end{itemize}
\begin{sphinxalltt}
\sphinxstyleemphasis{lj/cut/soft} args = n alpha\_lj cutoff
  n, alpha\_LJ = parameters of soft\sphinxhyphen{}core potential
  cutoff = global cutoff for Lennard\sphinxhyphen{}Jones interactions (distance units)
\sphinxstyleemphasis{lj/cut/coul/cut/soft} args = n alpha\_LJ alpha\_C cutoff (cutoff2)
  n, alpha\_LJ, alpha\_C = parameters of soft\sphinxhyphen{}core potential
  cutoff = global cutoff for LJ (and Coulombic if only 1 arg) (distance units)
  cutoff2 = global cutoff for Coulombic (optional) (distance units)
\sphinxstyleemphasis{lj/cut/coul/long/soft} args = n alpha\_LJ alpha\_C cutoff
  n, alpha\_LJ, alpha\_C = parameters of the soft\sphinxhyphen{}core potential
  cutoff = global cutoff for LJ (and Coulombic if only 1 arg) (distance units)
  cutoff2 = global cutoff for Coulombic (optional) (distance units)
\sphinxstyleemphasis{lj/cut/tip4p/long/soft} args = otype htype btype atype qdist n alpha\_LJ alpha\_C cutoff (cutoff2)
  otype,htype = atom types (numeric or type label) for TIP4P O and H
  btype,atype = bond and angle types (numeric or type label) for TIP4P waters
  qdist = distance from O atom to massless charge (distance units)
  n, alpha\_LJ, alpha\_C = parameters of the soft\sphinxhyphen{}core potential
  cutoff = global cutoff for LJ (and Coulombic if only 1 arg) (distance units)
  cutoff2 = global cutoff for Coulombic (optional) (distance units)
\sphinxstyleemphasis{lj/charmm/coul/long/soft} args = n alpha\_LJ alpha\_C inner outer (cutoff)
  n, alpha\_LJ, alpha\_C = parameters of the soft\sphinxhyphen{}core potential
  inner, outer = global switching cutoffs for LJ (and Coulombic if only 5 args)
  cutoff = global cutoff for Coulombic (optional, outer is Coulombic cutoff if only 5 args)
\sphinxstyleemphasis{lj/class2/soft} args = n alpha\_lj cutoff
  n, alpha\_LJ = parameters of soft\sphinxhyphen{}core potential
  cutoff = global cutoff for Lennard\sphinxhyphen{}Jones interactions (distance units)
\sphinxstyleemphasis{lj/class2/coul/cut/soft} args = n alpha\_LJ alpha\_C cutoff (cutoff2)
  n, alpha\_LJ, alpha\_C = parameters of soft\sphinxhyphen{}core potential
  cutoff = global cutoff for LJ (and Coulombic if only 1 arg) (distance units)
  cutoff2 = global cutoff for Coulombic (optional) (distance units)
\sphinxstyleemphasis{lj/class2/coul/long/soft} args = n alpha\_LJ alpha\_C cutoff (cutoff2)
  n, alpha\_LJ, alpha\_C = parameters of soft\sphinxhyphen{}core potential
  cutoff = global cutoff for LJ (and Coulombic if only 1 arg) (distance units)
  cutoff2 = global cutoff for Coulombic (optional) (distance units)
\sphinxstyleemphasis{coul/cut/soft} args = n alpha\_C cutoff
  n, alpha\_C = parameters of the soft\sphinxhyphen{}core potential
  cutoff = global cutoff for Coulomb interactions (distance units)
\sphinxstyleemphasis{coul/long/soft} args = n alpha\_C cutoff
  n, alpha\_C = parameters of the soft\sphinxhyphen{}core potential
  cutoff = global cutoff for Coulomb interactions (distance units)
\sphinxstyleemphasis{tip4p/long/soft} args = otype htype btype atype qdist n alpha\_C cutoff
  otype,htype = atom types (numeric or type label) for TIP4P O and H
  btype,atype = bond and angle types (numeric or type label) for TIP4P waters
  qdist = distance from O atom to massless charge (distance units)
  n, alpha\_C = parameters of the soft\sphinxhyphen{}core potential
  cutoff = global cutoff for Coulomb interactions (distance units)
\sphinxstyleemphasis{morse/soft} args = n lf cutoff
  n = soft\sphinxhyphen{}core parameter
  lf = transformation range is lf \textless{} lambda \textless{} 1
  cutoff = global cutoff for Morse interactions (distance units)
\end{sphinxalltt}


\subsection{Examples}
\label{\detokenize{pair_fep_soft:examples}}
\begin{sphinxVerbatim}[commandchars=\\\{\}]
\PYG{k}{pair\PYGZus{}style}\PYG{+w}{ }\PYG{n}{lj}\PYG{o}{/}\PYG{n}{cut}\PYG{o}{/}\PYG{n}{soft}\PYG{+w}{ }\PYG{l+m}{2.0}\PYG{+w}{ }\PYG{l+m}{0.5}\PYG{+w}{ }\PYG{l+m}{9.5}
\PYG{k}{pair\PYGZus{}coeff}\PYG{+w}{ }\PYG{o}{*}\PYG{+w}{ }\PYG{o}{*}\PYG{+w}{ }\PYG{l+m}{0.28}\PYG{+w}{ }\PYG{l+m}{3.1}\PYG{+w}{ }\PYG{l+m}{1.0}
\PYG{k}{pair\PYGZus{}coeff}\PYG{+w}{ }\PYG{l+m}{1}\PYG{+w}{ }\PYG{l+m}{1}\PYG{+w}{ }\PYG{l+m}{0.28}\PYG{+w}{ }\PYG{l+m}{3.1}\PYG{+w}{ }\PYG{l+m}{1.0}\PYG{+w}{ }\PYG{l+m}{9.5}

\PYG{k}{pair\PYGZus{}style}\PYG{+w}{ }\PYG{n}{lj}\PYG{o}{/}\PYG{n}{cut}\PYG{o}{/}\PYG{n}{coul}\PYG{o}{/}\PYG{n}{cut}\PYG{o}{/}\PYG{n}{soft}\PYG{+w}{ }\PYG{l+m}{2.0}\PYG{+w}{ }\PYG{l+m}{0.5}\PYG{+w}{ }\PYG{l+m}{10.0}\PYG{+w}{ }\PYG{l+m}{9.5}
\PYG{k}{pair\PYGZus{}style}\PYG{+w}{ }\PYG{n}{lj}\PYG{o}{/}\PYG{n}{cut}\PYG{o}{/}\PYG{n}{coul}\PYG{o}{/}\PYG{n}{cut}\PYG{o}{/}\PYG{n}{soft}\PYG{+w}{ }\PYG{l+m}{2.0}\PYG{+w}{ }\PYG{l+m}{0.5}\PYG{+w}{ }\PYG{l+m}{10.0}\PYG{+w}{ }\PYG{l+m}{9.5}\PYG{+w}{ }\PYG{l+m}{9.5}
\PYG{k}{pair\PYGZus{}coeff}\PYG{+w}{ }\PYG{o}{*}\PYG{+w}{ }\PYG{o}{*}\PYG{+w}{ }\PYG{l+m}{0.28}\PYG{+w}{ }\PYG{l+m}{3.1}\PYG{+w}{ }\PYG{l+m}{1.0}
\PYG{k}{pair\PYGZus{}coeff}\PYG{+w}{ }\PYG{l+m}{1}\PYG{+w}{ }\PYG{l+m}{1}\PYG{+w}{ }\PYG{l+m}{0.28}\PYG{+w}{ }\PYG{l+m}{3.1}\PYG{+w}{ }\PYG{l+m}{0.5}\PYG{+w}{ }\PYG{l+m}{10.0}
\PYG{k}{pair\PYGZus{}coeff}\PYG{+w}{ }\PYG{l+m}{1}\PYG{+w}{ }\PYG{l+m}{1}\PYG{+w}{ }\PYG{l+m}{0.28}\PYG{+w}{ }\PYG{l+m}{3.1}\PYG{+w}{ }\PYG{l+m}{0.5}\PYG{+w}{ }\PYG{l+m}{10.0}\PYG{+w}{ }\PYG{l+m}{9.5}

\PYG{k}{pair\PYGZus{}style}\PYG{+w}{ }\PYG{n}{lj}\PYG{o}{/}\PYG{n}{cut}\PYG{o}{/}\PYG{n}{coul}\PYG{o}{/}\PYG{n}{long}\PYG{o}{/}\PYG{n}{soft}\PYG{+w}{ }\PYG{l+m}{2.0}\PYG{+w}{ }\PYG{l+m}{0.5}\PYG{+w}{ }\PYG{l+m}{10.0}\PYG{+w}{ }\PYG{l+m}{9.5}
\PYG{k}{pair\PYGZus{}style}\PYG{+w}{ }\PYG{n}{lj}\PYG{o}{/}\PYG{n}{cut}\PYG{o}{/}\PYG{n}{coul}\PYG{o}{/}\PYG{n}{long}\PYG{o}{/}\PYG{n}{soft}\PYG{+w}{ }\PYG{l+m}{2.0}\PYG{+w}{ }\PYG{l+m}{0.5}\PYG{+w}{ }\PYG{l+m}{10.0}\PYG{+w}{ }\PYG{l+m}{9.5}\PYG{+w}{ }\PYG{l+m}{9.5}
\PYG{k}{pair\PYGZus{}coeff}\PYG{+w}{ }\PYG{o}{*}\PYG{+w}{ }\PYG{o}{*}\PYG{+w}{ }\PYG{l+m}{0.28}\PYG{+w}{ }\PYG{l+m}{3.1}\PYG{+w}{ }\PYG{l+m}{1.0}
\PYG{k}{pair\PYGZus{}coeff}\PYG{+w}{ }\PYG{l+m}{1}\PYG{+w}{ }\PYG{l+m}{1}\PYG{+w}{ }\PYG{l+m}{0.28}\PYG{+w}{ }\PYG{l+m}{3.1}\PYG{+w}{ }\PYG{l+m}{0.0}\PYG{+w}{ }\PYG{l+m}{10.0}
\PYG{k}{pair\PYGZus{}coeff}\PYG{+w}{ }\PYG{l+m}{1}\PYG{+w}{ }\PYG{l+m}{1}\PYG{+w}{ }\PYG{l+m}{0.28}\PYG{+w}{ }\PYG{l+m}{3.1}\PYG{+w}{ }\PYG{l+m}{0.0}\PYG{+w}{ }\PYG{l+m}{10.0}\PYG{+w}{ }\PYG{l+m}{9.5}

\PYG{k}{pair\PYGZus{}style}\PYG{+w}{ }\PYG{n}{lj}\PYG{o}{/}\PYG{n}{cut}\PYG{o}{/}\PYG{n}{tip4p}\PYG{o}{/}\PYG{n}{long}\PYG{o}{/}\PYG{n}{soft}\PYG{+w}{ }\PYG{l+m}{1}\PYG{+w}{ }\PYG{l+m}{2}\PYG{+w}{ }\PYG{l+m}{7}\PYG{+w}{ }\PYG{l+m}{8}\PYG{+w}{ }\PYG{l+m}{0.15}\PYG{+w}{ }\PYG{l+m}{2.0}\PYG{+w}{ }\PYG{l+m}{0.5}\PYG{+w}{ }\PYG{l+m}{10.0}\PYG{+w}{ }\PYG{l+m}{9.8}
\PYG{k}{pair\PYGZus{}style}\PYG{+w}{ }\PYG{n}{lj}\PYG{o}{/}\PYG{n}{cut}\PYG{o}{/}\PYG{n}{tip4p}\PYG{o}{/}\PYG{n}{long}\PYG{o}{/}\PYG{n}{soft}\PYG{+w}{ }\PYG{l+m}{1}\PYG{+w}{ }\PYG{l+m}{2}\PYG{+w}{ }\PYG{l+m}{7}\PYG{+w}{ }\PYG{l+m}{8}\PYG{+w}{ }\PYG{l+m}{0.15}\PYG{+w}{ }\PYG{l+m}{2.0}\PYG{+w}{ }\PYG{l+m}{0.5}\PYG{+w}{ }\PYG{l+m}{10.0}\PYG{+w}{ }\PYG{l+m}{9.8}\PYG{+w}{ }\PYG{l+m}{9.5}
\PYG{k}{pair\PYGZus{}coeff}\PYG{+w}{ }\PYG{o}{*}\PYG{+w}{ }\PYG{o}{*}\PYG{+w}{ }\PYG{l+m}{0.155}\PYG{+w}{ }\PYG{l+m}{3.1536}\PYG{+w}{ }\PYG{l+m}{1.0}
\PYG{k}{pair\PYGZus{}coeff}\PYG{+w}{ }\PYG{l+m}{1}\PYG{+w}{ }\PYG{l+m}{1}\PYG{+w}{ }\PYG{l+m}{0.155}\PYG{+w}{ }\PYG{l+m}{3.1536}\PYG{+w}{ }\PYG{l+m}{1.0}\PYG{+w}{ }\PYG{l+m}{9.5}

\PYG{k}{pair\PYGZus{}style}\PYG{+w}{ }\PYG{n}{lj}\PYG{o}{/}\PYG{n}{cut}\PYG{o}{/}\PYG{n}{tip4p}\PYG{o}{/}\PYG{n}{long}\PYG{o}{/}\PYG{n}{soft}\PYG{+w}{ }\PYG{n}{OW}\PYG{+w}{ }\PYG{n}{HW}\PYG{+w}{ }\PYG{n}{HW}\PYG{o}{\PYGZhy{}}\PYG{n}{OW}\PYG{+w}{ }\PYG{n}{HW}\PYG{o}{\PYGZhy{}}\PYG{n}{OW}\PYG{o}{\PYGZhy{}}\PYG{n}{HW}\PYG{+w}{ }\PYG{l+m}{0.15}\PYG{+w}{ }\PYG{l+m}{2.0}\PYG{+w}{ }\PYG{l+m}{0.5}\PYG{+w}{ }\PYG{l+m}{10.0}\PYG{+w}{ }\PYG{l+m}{9.8}
\PYG{k}{labelmap}\PYG{+w}{ }\PYG{n}{atom}\PYG{+w}{ }\PYG{l+m}{1}\PYG{+w}{ }\PYG{n}{OW}\PYG{+w}{ }\PYG{l+m}{2}\PYG{+w}{ }\PYG{n}{HW}
\PYG{k}{labelmap}\PYG{+w}{ }\PYG{n}{bond}\PYG{+w}{ }\PYG{l+m}{1}\PYG{+w}{ }\PYG{n}{HW}\PYG{o}{\PYGZhy{}}\PYG{n}{OW}
\PYG{k}{labelmap}\PYG{+w}{ }\PYG{n}{angle}\PYG{+w}{ }\PYG{l+m}{1}\PYG{+w}{ }\PYG{n}{HW}\PYG{o}{\PYGZhy{}}\PYG{n}{OW}\PYG{o}{\PYGZhy{}}\PYG{n}{HW}
\PYG{k}{pair\PYGZus{}coeff}\PYG{+w}{ }\PYG{o}{*}\PYG{+w}{ }\PYG{o}{*}\PYG{+w}{ }\PYG{l+m}{0.155}\PYG{+w}{ }\PYG{l+m}{3.1536}\PYG{+w}{ }\PYG{l+m}{1.0}
\PYG{k}{pair\PYGZus{}coeff}\PYG{+w}{ }\PYG{n}{OW}\PYG{+w}{ }\PYG{n}{OW}\PYG{+w}{ }\PYG{l+m}{0.155}\PYG{+w}{ }\PYG{l+m}{3.1536}\PYG{+w}{ }\PYG{l+m}{1.0}\PYG{+w}{ }\PYG{l+m}{9.5}

\PYG{k}{pair\PYGZus{}style}\PYG{+w}{ }\PYG{n}{lj}\PYG{o}{/}\PYG{n}{charmm}\PYG{o}{/}\PYG{n}{coul}\PYG{o}{/}\PYG{n}{long}\PYG{+w}{ }\PYG{l+m}{2.0}\PYG{+w}{ }\PYG{l+m}{0.5}\PYG{+w}{ }\PYG{l+m}{10.0}\PYG{+w}{ }\PYG{l+m}{8.0}\PYG{+w}{ }\PYG{l+m}{10.0}
\PYG{k}{pair\PYGZus{}style}\PYG{+w}{ }\PYG{n}{lj}\PYG{o}{/}\PYG{n}{charmm}\PYG{o}{/}\PYG{n}{coul}\PYG{o}{/}\PYG{n}{long}\PYG{+w}{ }\PYG{l+m}{2.0}\PYG{+w}{ }\PYG{l+m}{0.5}\PYG{+w}{ }\PYG{l+m}{10.0}\PYG{+w}{ }\PYG{l+m}{8.0}\PYG{+w}{ }\PYG{l+m}{10.0}\PYG{+w}{ }\PYG{l+m}{9.0}
\PYG{k}{pair\PYGZus{}coeff}\PYG{+w}{ }\PYG{o}{*}\PYG{+w}{ }\PYG{o}{*}\PYG{+w}{ }\PYG{l+m}{0.28}\PYG{+w}{ }\PYG{l+m}{3.1}\PYG{+w}{ }\PYG{l+m}{1.0}
\PYG{k}{pair\PYGZus{}coeff}\PYG{+w}{ }\PYG{l+m}{1}\PYG{+w}{ }\PYG{l+m}{1}\PYG{+w}{ }\PYG{l+m}{0.28}\PYG{+w}{ }\PYG{l+m}{3.1}\PYG{+w}{ }\PYG{l+m}{1.0}\PYG{+w}{ }\PYG{l+m}{0.14}\PYG{+w}{ }\PYG{l+m}{3.1}

\PYG{k}{pair\PYGZus{}style}\PYG{+w}{ }\PYG{n}{lj}\PYG{o}{/}\PYG{n}{class2}\PYG{o}{/}\PYG{n}{coul}\PYG{o}{/}\PYG{n}{long}\PYG{o}{/}\PYG{n}{soft}\PYG{+w}{ }\PYG{l+m}{2.0}\PYG{+w}{ }\PYG{l+m}{0.5}\PYG{+w}{ }\PYG{l+m}{10.0}\PYG{+w}{ }\PYG{l+m}{9.5}
\PYG{k}{pair\PYGZus{}style}\PYG{+w}{ }\PYG{n}{lj}\PYG{o}{/}\PYG{n}{class2}\PYG{o}{/}\PYG{n}{coul}\PYG{o}{/}\PYG{n}{long}\PYG{o}{/}\PYG{n}{soft}\PYG{+w}{ }\PYG{l+m}{2.0}\PYG{+w}{ }\PYG{l+m}{0.5}\PYG{+w}{ }\PYG{l+m}{10.0}\PYG{+w}{ }\PYG{l+m}{9.5}\PYG{+w}{ }\PYG{l+m}{9.5}
\PYG{k}{pair\PYGZus{}coeff}\PYG{+w}{ }\PYG{o}{*}\PYG{+w}{ }\PYG{o}{*}\PYG{+w}{ }\PYG{l+m}{0.28}\PYG{+w}{ }\PYG{l+m}{3.1}\PYG{+w}{ }\PYG{l+m}{1.0}
\PYG{k}{pair\PYGZus{}coeff}\PYG{+w}{ }\PYG{l+m}{1}\PYG{+w}{ }\PYG{l+m}{1}\PYG{+w}{ }\PYG{l+m}{0.28}\PYG{+w}{ }\PYG{l+m}{3.1}\PYG{+w}{ }\PYG{l+m}{0.0}\PYG{+w}{ }\PYG{l+m}{10.0}
\PYG{k}{pair\PYGZus{}coeff}\PYG{+w}{ }\PYG{l+m}{1}\PYG{+w}{ }\PYG{l+m}{1}\PYG{+w}{ }\PYG{l+m}{0.28}\PYG{+w}{ }\PYG{l+m}{3.1}\PYG{+w}{ }\PYG{l+m}{0.0}\PYG{+w}{ }\PYG{l+m}{10.0}\PYG{+w}{ }\PYG{l+m}{9.5}

\PYG{k}{pair\PYGZus{}style}\PYG{+w}{ }\PYG{n}{coul}\PYG{o}{/}\PYG{n}{long}\PYG{o}{/}\PYG{n}{soft}\PYG{+w}{ }\PYG{l+m}{1.0}\PYG{+w}{ }\PYG{l+m}{10.0}\PYG{+w}{ }\PYG{l+m}{9.5}
\PYG{k}{pair\PYGZus{}coeff}\PYG{+w}{ }\PYG{o}{*}\PYG{+w}{ }\PYG{o}{*}\PYG{+w}{ }\PYG{l+m}{1.0}
\PYG{k}{pair\PYGZus{}coeff}\PYG{+w}{ }\PYG{l+m}{1}\PYG{+w}{ }\PYG{l+m}{1}\PYG{+w}{ }\PYG{l+m}{1.0}

\PYG{k}{pair\PYGZus{}style}\PYG{+w}{ }\PYG{n}{tip4p}\PYG{o}{/}\PYG{n}{long}\PYG{o}{/}\PYG{n}{soft}\PYG{+w}{ }\PYG{l+m}{1}\PYG{+w}{ }\PYG{l+m}{2}\PYG{+w}{ }\PYG{l+m}{7}\PYG{+w}{ }\PYG{l+m}{8}\PYG{+w}{ }\PYG{l+m}{0.15}\PYG{+w}{ }\PYG{l+m}{2.0}\PYG{+w}{ }\PYG{l+m}{0.5}\PYG{+w}{ }\PYG{l+m}{10.0}\PYG{+w}{ }\PYG{l+m}{9.8}
\PYG{k}{pair\PYGZus{}coeff}\PYG{+w}{ }\PYG{o}{*}\PYG{+w}{ }\PYG{o}{*}\PYG{+w}{ }\PYG{l+m}{1.0}
\PYG{k}{pair\PYGZus{}coeff}\PYG{+w}{ }\PYG{l+m}{1}\PYG{+w}{ }\PYG{l+m}{1}\PYG{+w}{ }\PYG{l+m}{1.0}

\PYG{k}{pair\PYGZus{}style}\PYG{+w}{ }\PYG{n}{morse}\PYG{o}{/}\PYG{n}{soft}\PYG{+w}{ }\PYG{l+m}{4}\PYG{+w}{ }\PYG{l+m}{0.9}\PYG{+w}{ }\PYG{l+m}{10.0}
\PYG{k}{pair\PYGZus{}coeff}\PYG{+w}{ }\PYG{o}{*}\PYG{+w}{ }\PYG{o}{*}\PYG{+w}{ }\PYG{l+m}{100.0}\PYG{+w}{ }\PYG{l+m}{2.0}\PYG{+w}{ }\PYG{l+m}{1.5}\PYG{+w}{ }\PYG{l+m}{1.0}
\PYG{k}{pair\PYGZus{}coeff}\PYG{+w}{ }\PYG{l+m}{1}\PYG{+w}{ }\PYG{l+m}{1}\PYG{+w}{ }\PYG{l+m}{100.0}\PYG{+w}{ }\PYG{l+m}{2.0}\PYG{+w}{ }\PYG{l+m}{1.5}\PYG{+w}{ }\PYG{l+m}{1.0}\PYG{+w}{ }\PYG{l+m}{3.0}
\end{sphinxVerbatim}

\sphinxAtStartPar
Example input scripts available: examples/PACKAGES/fep


\subsection{Description}
\label{\detokenize{pair_fep_soft:description}}
\sphinxAtStartPar
These pair styles have a soft repulsive core, tunable by a parameter lambda,
in order to avoid singularities during free energy calculations when sites are
created or annihilated {\hyperref[\detokenize{pair_fep_soft:beutler}]{\sphinxcrossref{\DUrole{std,std-ref}{(Beutler)}}}}.  When lambda tends to 0 the pair
interaction vanishes with a soft repulsive core.  When lambda tends to 1, the pair
interaction approaches the normal, non\sphinxhyphen{}soft potential. These pair styles
are suited for “alchemical” free energy calculations using the {\hyperref[\detokenize{fix_adapt_fep::doc}]{\sphinxcrossref{\DUrole{doc}{fix adapt/fep}}}} and {\hyperref[\detokenize{compute_fep::doc}]{\sphinxcrossref{\DUrole{doc}{compute fep}}}} commands.

\sphinxAtStartPar
The \sphinxstyleemphasis{lj/cut/soft} style and related sub\sphinxhyphen{}styles compute the 12\sphinxhyphen{}6 Lennard\sphinxhyphen{}Jones
and Coulomb potentials modified by a soft core, with the functional form
\begin{equation*}
\begin{split}E = \lambda^n 4 \epsilon \left\{
\frac{1}{ \left[ \alpha_{\mathrm{LJ}} (1-\lambda)^2 +
\left( \displaystyle\frac{r}{\sigma} \right)^6 \right]^2 } -
\frac{1}{ \alpha_{\mathrm{LJ}} (1-\lambda)^2 +
\left( \displaystyle\frac{r}{\sigma} \right)^6 }
\right\} \qquad r < r_c\end{split}
\end{equation*}
\sphinxAtStartPar
The \sphinxstyleemphasis{lj/class2/soft} style is a 9\sphinxhyphen{}6 potential with the exponent of the
denominator of the first term in brackets taking the value 1.5 instead of 2
(other details differ, see the form of the potential in
{\hyperref[\detokenize{pair_class2::doc}]{\sphinxcrossref{\DUrole{doc}{pair\_style lj/class2}}}}).

\sphinxAtStartPar
Coulomb interactions can also be damped with a soft core at short distance,
\begin{equation*}
\begin{split}E = \lambda^n \frac{ C q_i q_j}{\epsilon \left[ \alpha_{\mathrm{C}}
(1-\lambda)^2 + r^2 \right]^{1/2}} \qquad r < r_c\end{split}
\end{equation*}
\sphinxAtStartPar
In the Coulomb part \(C\) is an energy\sphinxhyphen{}conversion constant, \(q_i\) and
\(q_j\) are the charges on the two atoms, and epsilon is the dielectric
constant which can be set by the {\hyperref[\detokenize{dielectric::doc}]{\sphinxcrossref{\DUrole{doc}{dielectric}}}} command.

\sphinxAtStartPar
The coefficient lambda is an activation parameter. When \(\lambda = 1\) the
pair potential is identical to a Lennard\sphinxhyphen{}Jones term or a Coulomb term or a
combination of both. When \(\lambda = 0\) the interactions are
deactivated. The transition between these two extrema is smoothed by a soft
repulsive core in order to avoid singularities in potential energy and forces
when sites are created or annihilated and can overlap {\hyperref[\detokenize{pair_fep_soft:beutler}]{\sphinxcrossref{\DUrole{std,std-ref}{(Beutler)}}}}.

\sphinxAtStartPar
The parameters \(n\), \(\alpha_\mathrm{LJ}\) and
\(\alpha_\mathrm{C}\) are set in the {\hyperref[\detokenize{pair_style::doc}]{\sphinxcrossref{\DUrole{doc}{pair\_style}}}} command,
before the cutoffs.  Usual choices for the exponent are \(n = 2\) or
\(n = 1\). For the remaining coefficients \(\alpha_\mathrm{LJ} = 0.5\)
and \(\alpha_\mathrm{C} = 10~\text{A}^2\) are appropriate choices. Plots of
the 12\sphinxhyphen{}6 LJ and Coulomb terms are shown below, for lambda ranging from 1 to 0
every 0.1.

\noindent\sphinxincludegraphics{{lj_soft}.jpg}

\noindent\sphinxincludegraphics{{coul_soft}.jpg}

\sphinxAtStartPar
For the \sphinxstyleemphasis{lj/cut/coul/cut/soft} or \sphinxstyleemphasis{lj/cut/coul/long/soft} pair styles, as well
as for the equivalent \sphinxstyleemphasis{class2} versions, the following coefficients must be
defined for each pair of atoms types via the {\hyperref[\detokenize{pair_coeff::doc}]{\sphinxcrossref{\DUrole{doc}{pair\_coeff}}}}
command as in the examples above, or in the data file or restart files read by
the {\hyperref[\detokenize{read_data::doc}]{\sphinxcrossref{\DUrole{doc}{read\_data}}}} or {\hyperref[\detokenize{read_restart::doc}]{\sphinxcrossref{\DUrole{doc}{read\_restart}}}} commands,
or by mixing as described below:
\begin{itemize}
\item {} 
\sphinxAtStartPar
\(\epsilon\) (energy units)

\item {} 
\sphinxAtStartPar
\(\sigma\) (distance units)

\item {} 
\sphinxAtStartPar
\(\lambda\) (activation parameter, between 0 and 1)

\item {} 
\sphinxAtStartPar
cutoff1 (distance units)

\item {} 
\sphinxAtStartPar
cutoff2 (distance units)

\end{itemize}

\sphinxAtStartPar
The latter two coefficients are optional.  If not specified, the global
LJ and Coulombic cutoffs specified in the pair\_style command are used.
If only one cutoff is specified, it is used as the cutoff for both LJ
and Coulombic interactions for this type pair.  If both coefficients
are specified, they are used as the LJ and Coulombic cutoffs for this
type pair.  You cannot specify 2 cutoffs for style \sphinxstyleemphasis{lj/cut/soft},
since it has no Coulombic terms. For the \sphinxstyleemphasis{coul/cut/soft} and
\sphinxstyleemphasis{coul/long/soft} only lambda and the optional cutoff2 are to be
specified.

\sphinxAtStartPar
Style \sphinxstyleemphasis{lj/cut/tip4p/long/soft} implements a soft\sphinxhyphen{}core version of the TIP4P water
model. The usage of the TIP4P pair style is documented in the {\hyperref[\detokenize{pair_lj::doc}]{\sphinxcrossref{\DUrole{doc}{pair\_lj}}}} styles. In the soft version the parameters \(n\),
\(\alpha_\mathrm{LJ}\) and \(\alpha_\mathrm {C}\) are set in the
{\hyperref[\detokenize{pair_style::doc}]{\sphinxcrossref{\DUrole{doc}{pair\_style}}}} command, after the specific parameters of the
TIP4P water model and before the cutoffs. The activation parameter lambda is
supplied as an argument of the {\hyperref[\detokenize{pair_coeff::doc}]{\sphinxcrossref{\DUrole{doc}{pair\_coeff}}}} command, after
epsilon and sigma and before the optional cutoffs.

\begin{sphinxadmonition}{note}{Note:}
\sphinxAtStartPar
If using type labels, the type labels must be defined before calling
the {\hyperref[\detokenize{pair_coeff::doc}]{\sphinxcrossref{\DUrole{doc}{pair\_coeff}}}} command.
\end{sphinxadmonition}

\sphinxAtStartPar
Style \sphinxstyleemphasis{lj/charmm/coul/long/soft} implements a soft\sphinxhyphen{}core version of the modified
12\sphinxhyphen{}6 LJ potential used in CHARMM and documented in the {\hyperref[\detokenize{pair_charmm::doc}]{\sphinxcrossref{\DUrole{doc}{pair\_style
lj/charmm/coul/long}}}} style. In the soft version the parameters
\(n\), \(\alpha_\mathrm{LJ}\) and \(\alpha_\mathrm{C}\) are set in
the {\hyperref[\detokenize{pair_style::doc}]{\sphinxcrossref{\DUrole{doc}{pair\_style}}}} command, before the global cutoffs. The
activation parameter lambda is introduced as an argument of the {\hyperref[\detokenize{pair_coeff::doc}]{\sphinxcrossref{\DUrole{doc}{pair\_coeff}}}} command, after \(\epsilon\) and \(\sigma\) and before the
optional eps14 and sigma14.

\sphinxAtStartPar
Style \sphinxstyleemphasis{lj/class2/soft} implements a soft\sphinxhyphen{}core version of the 9\sphinxhyphen{}6 potential in
{\hyperref[\detokenize{pair_class2::doc}]{\sphinxcrossref{\DUrole{doc}{pair\_style lj/class2}}}}. In the soft version the parameters
\(n\), \(\alpha_\mathrm{LJ}\) and \(\alpha_\mathrm{C}\) are set in the
{\hyperref[\detokenize{pair_style::doc}]{\sphinxcrossref{\DUrole{doc}{pair\_style}}}} command, before the global cutoffs. The
activation parameter lambda is introduced as an argument of the the
{\hyperref[\detokenize{pair_coeff::doc}]{\sphinxcrossref{\DUrole{doc}{pair\_coeff}}}} command, after \(\epsilon\) and
\(\sigma\) and before the optional cutoffs.

\sphinxAtStartPar
The \sphinxstyleemphasis{coul/cut/soft}, \sphinxstyleemphasis{coul/long/soft} and \sphinxstyleemphasis{tip4p/long/soft} sub\sphinxhyphen{}styles are
designed to be combined with other pair potentials via the {\hyperref[\detokenize{pair_hybrid::doc}]{\sphinxcrossref{\DUrole{doc}{pair\_style
hybrid/overlay}}}} command.  This is because they have no repulsive
core.  Hence, if used by themselves, there will be no repulsion to keep two
oppositely charged particles from overlapping each other. In this case, if
\(\lambda = 1\), a singularity may occur.  These sub\sphinxhyphen{}styles are suitable to
represent charges embedded in the Lennard\sphinxhyphen{}Jones radius of another site (for
example hydrogen atoms in several water models). The \(\lambda\) must
be defined for each pair, and \sphinxstyleemphasis{coul/cut/soft} can accept an optional cutoff as
the second coefficient.

\begin{sphinxadmonition}{note}{Note:}
\sphinxAtStartPar
When using the soft\sphinxhyphen{}core Coulomb potentials with long\sphinxhyphen{}range solvers (\sphinxstyleemphasis{coul/long/soft}, \sphinxstyleemphasis{lj/cut/coul/long/soft}, etc.)  in a free energy
calculation in which sites holding electrostatic charges are being created or
annihilated (using {\hyperref[\detokenize{fix_adapt_fep::doc}]{\sphinxcrossref{\DUrole{doc}{fix adapt/fep}}}} and {\hyperref[\detokenize{compute_fep::doc}]{\sphinxcrossref{\DUrole{doc}{compute fep}}}}) it is important to adapt both the \(\lambda\) activation
parameter (from 0 to 1, or the reverse) and the value of the charge (from 0
to its final value, or the reverse). This ensures that long\sphinxhyphen{}range
electrostatic terms (kspace) are correct. It is not necessary to use
soft\sphinxhyphen{}core Coulomb potentials if the van der Waals site is present during the
free\sphinxhyphen{}energy route, thus avoiding overlap of the charges. Examples are
provided in the LAMMPS source directory tree, under examples/PACKAGES/fep.
\end{sphinxadmonition}

\begin{sphinxadmonition}{note}{Note:}
\sphinxAtStartPar
To avoid division by zero do not set \(\sigma = 0\) in the \sphinxstyleemphasis{lj/cut/soft}
and related styles; use the lambda parameter instead to activate/deactivate
interactions, or use \(\epsilon = 0\) and \(\sigma = 1\).
Alternatively, when sites do not interact though the Lennard\sphinxhyphen{}Jones term
the \sphinxstyleemphasis{coul/long/soft} or similar sub\sphinxhyphen{}style can be used via the
{\hyperref[\detokenize{pair_hybrid::doc}]{\sphinxcrossref{\DUrole{doc}{pair\_style hybrid/overlay}}}} command.
\end{sphinxadmonition}


\bigskip\hrule\bigskip


\sphinxAtStartPar
The \sphinxstyleemphasis{morse/soft} variant modifies the {\hyperref[\detokenize{pair_morse::doc}]{\sphinxcrossref{\DUrole{doc}{pair\_morse}}}} style at
short range to have a soft core. The functional form differs from that of the
\sphinxstyleemphasis{lj/soft} styles, and is instead given by:
\begin{equation*}
\begin{split}\begin{split}
s(\lambda) =& (1 - \lambda) / (1 - \lambda_f), \qquad B = -2D e^{-2 \alpha
r_0} (e^{\alpha r_0} - 1) / 3 \\
E =& D_0 \left[ e^{- 2 \alpha (r - r_0)} - 2 e^{- \alpha (r - r_0)} \right] +
s(\lambda) B e^{-3\alpha(r-r_0)}, \qquad \hspace{2.85em}\lambda \geq
\lambda_f,\quad r < r_c \\
E =& \left( D_0 \left[ e^{- 2 \alpha (r - r_0)} - 2 e^{- \alpha (r - r_0)}
\right] + B e^{-3\alpha(r-r_0)} \right)(\lambda/\lambda_f)^n, \qquad \lambda
< \lambda_f,\quad r < r_c
\end{split}\end{split}
\end{equation*}
\sphinxAtStartPar
The \sphinxstyleemphasis{morse/soft} style requires the following pair coefficients:
\begin{itemize}
\item {} 
\sphinxAtStartPar
\(D_0\) (energy units)

\item {} 
\sphinxAtStartPar
\(\alpha\) (1/distance units)

\item {} 
\sphinxAtStartPar
\(r_0\) (distance units)

\item {} 
\sphinxAtStartPar
\(\lambda\) (unitless, between 0.0 and 1.0)

\item {} 
\sphinxAtStartPar
cutoff (distance units)

\end{itemize}

\sphinxAtStartPar
The last coefficient is optional. If not specified, the global morse cutoff is
used.


\bigskip\hrule\bigskip


\sphinxAtStartPar
Styles with a \sphinxstyleemphasis{gpu}, \sphinxstyleemphasis{intel}, \sphinxstyleemphasis{kk}, \sphinxstyleemphasis{omp}, or \sphinxstyleemphasis{opt} suffix are
functionally the same as the corresponding style without the suffix.
They have been optimized to run faster, depending on your available
hardware, as discussed on the {\hyperref[\detokenize{Speed_packages::doc}]{\sphinxcrossref{\DUrole{doc}{Accelerator packages}}}}
page.  The accelerated styles take the same arguments and should
produce the same results, except for round\sphinxhyphen{}off and precision issues.

\sphinxAtStartPar
These accelerated styles are part of the GPU, INTEL, KOKKOS,
OPENMP, and OPT packages, respectively.  They are only enabled if
LAMMPS was built with those packages.  See the {\hyperref[\detokenize{Build_package::doc}]{\sphinxcrossref{\DUrole{doc}{Build package}}}} page for more info.

\sphinxAtStartPar
You can specify the accelerated styles explicitly in your input script
by including their suffix, or you can use the {\hyperref[\detokenize{Run_options::doc}]{\sphinxcrossref{\DUrole{doc}{\sphinxhyphen{}suffix command\sphinxhyphen{}line switch}}}} when you invoke LAMMPS, or you can use the
{\hyperref[\detokenize{suffix::doc}]{\sphinxcrossref{\DUrole{doc}{suffix}}}} command in your input script.

\sphinxAtStartPar
See the {\hyperref[\detokenize{Speed_packages::doc}]{\sphinxcrossref{\DUrole{doc}{Accelerator packages}}}} page for more
instructions on how to use the accelerated styles effectively.


\bigskip\hrule\bigskip



\subsection{Mixing, shift, table, tail correction, restart, rRESPA info}
\label{\detokenize{pair_fep_soft:mixing-shift-table-tail-correction-restart-rrespa-info}}
\sphinxAtStartPar
The different versions of the \sphinxstyleemphasis{lj/cut/soft} pair styles support mixing.  For
atom type pairs I,J and I != J, the \(\epsilon\) and \(\sigma\)
coefficients and cutoff distance for these pair styles can be mixed.  The default
mix value is \sphinxstyleemphasis{geometric} for 12\sphinxhyphen{}6 styles.

\sphinxAtStartPar
The mixing rule for epsilon and sigma for \sphinxstyleemphasis{lj/class2/soft} 9\sphinxhyphen{}6 potentials is to
use the \sphinxstyleemphasis{sixthpower} formulas. The {\hyperref[\detokenize{pair_modify::doc}]{\sphinxcrossref{\DUrole{doc}{pair\_modify mix}}}} setting
is thus ignored for class2 potentials for \(\epsilon\) and
\(\sigma\). However it is still followed for mixing the cutoff distance. See
the {\hyperref[\detokenize{pair_modify::doc}]{\sphinxcrossref{\DUrole{doc}{pair\_modify}}}} command for details.

\sphinxAtStartPar
The \sphinxstyleemphasis{morse/soft} pair style does not support mixing. Thus, coefficients for all
LJ pairs must be specified explicitly.

\sphinxAtStartPar
All of the pair styles with soft core support the {\hyperref[\detokenize{pair_modify::doc}]{\sphinxcrossref{\DUrole{doc}{pair\_modify}}}}
shift option for the energy of the Lennard\sphinxhyphen{}Jones portion of the pair
interaction.

\sphinxAtStartPar
The different versions of the \sphinxstyleemphasis{lj/cut/soft} pair styles support the
{\hyperref[\detokenize{pair_modify::doc}]{\sphinxcrossref{\DUrole{doc}{pair\_modify}}}} tail option for adding a long\sphinxhyphen{}range tail
correction to the energy and pressure for the Lennard\sphinxhyphen{}Jones portion of the pair
interaction.

\begin{sphinxadmonition}{note}{Note:}
\sphinxAtStartPar
The analytical form of the tail corrections for energy and pressure used in
the \sphinxstyleemphasis{lj/cut/soft} potentials are approximate, being identical to that of the
corresponding non\sphinxhyphen{}soft potentials scaled by a factor \(\lambda^n\). The
errors due to this approximation should be negligible. For example, for a
cutoff of \(2.5\sigma\) this approximation leads to maximum relative
errors in tail corrections of the order of 1e\sphinxhyphen{}4 for energy and virial
(\(\alpha_\mathrm{LJ} = 0.5, n = 2\)). The error vanishes when lambda
approaches 0 or 1. Note that these are the errors affecting the long\sphinxhyphen{}range
tail (itself a correction to the interaction energy) which includes other
approximations, namely that the system is homogeneous (local density equal
the average density) beyond the cutoff.
\end{sphinxadmonition}

\sphinxAtStartPar
The \sphinxstyleemphasis{morse/soft} pair style does not support the {\hyperref[\detokenize{pair_modify::doc}]{\sphinxcrossref{\DUrole{doc}{pair\_modify}}}} tail option for adding long\sphinxhyphen{}range tail corrections to energy and
pressure.

\sphinxAtStartPar
All of these pair styles write information to {\hyperref[\detokenize{restart::doc}]{\sphinxcrossref{\DUrole{doc}{binary restart files}}}}, so pair\_style and pair\_coeff commands do not need to be specified
in an input script that reads a restart file.


\bigskip\hrule\bigskip



\subsection{Restrictions}
\label{\detokenize{pair_fep_soft:restrictions}}
\sphinxAtStartPar
The pair styles with soft core are only enabled if LAMMPS was built with the
FEP package. The \sphinxstyleemphasis{long} versions also require the KSPACE package to be
installed. The soft \sphinxstyleemphasis{tip4p} versions also require the MOLECULE package to be
installed. These styles are only enabled if LAMMPS was built with those
packages.  See the {\hyperref[\detokenize{Build_package::doc}]{\sphinxcrossref{\DUrole{doc}{Build package}}}} page for more
info.


\subsection{Related commands}
\label{\detokenize{pair_fep_soft:related-commands}}
\sphinxAtStartPar
{\hyperref[\detokenize{pair_coeff::doc}]{\sphinxcrossref{\DUrole{doc}{pair\_coeff}}}}, {\hyperref[\detokenize{fix_adapt::doc}]{\sphinxcrossref{\DUrole{doc}{fix adapt}}}},
{\hyperref[\detokenize{fix_adapt_fep::doc}]{\sphinxcrossref{\DUrole{doc}{fix adapt/fep}}}}, {\hyperref[\detokenize{compute_fep::doc}]{\sphinxcrossref{\DUrole{doc}{compute fep}}}}


\subsection{Default}
\label{\detokenize{pair_fep_soft:default}}
\sphinxAtStartPar
none


\bigskip\hrule\bigskip

\phantomsection\label{\detokenize{pair_fep_soft:beutler}}
\sphinxAtStartPar
\sphinxstylestrong{(Beutler)} Beutler, Mark, van Schaik, Gerber, van Gunsteren, Chem
Phys Lett, 222, 529 (1994).

\sphinxstepscope

\index{pair\_style gauss@\spxentry{pair\_style gauss}}\ignorespaces 
\index{pair\_style gauss/gpu@\spxentry{pair\_style gauss/gpu}}\ignorespaces 
\index{pair\_style gauss/omp@\spxentry{pair\_style gauss/omp}}\ignorespaces 
\index{pair\_style gauss/cut@\spxentry{pair\_style gauss/cut}}\ignorespaces 
\index{pair\_style gauss/cut/omp@\spxentry{pair\_style gauss/cut/omp}}\ignorespaces 

\section{pair\_style gauss command}
\label{\detokenize{pair_gauss:pair-style-gauss-command}}\label{\detokenize{pair_gauss:index-4}}\label{\detokenize{pair_gauss:index-3}}\label{\detokenize{pair_gauss:index-2}}\label{\detokenize{pair_gauss:index-1}}\label{\detokenize{pair_gauss:index-0}}\label{\detokenize{pair_gauss::doc}}
\sphinxAtStartPar
Accelerator Variants: \sphinxstyleemphasis{gauss/gpu}, \sphinxstyleemphasis{gauss/omp}


\section{pair\_style gauss/cut command}
\label{\detokenize{pair_gauss:pair-style-gauss-cut-command}}
\sphinxAtStartPar
Accelerator Variants: \sphinxstyleemphasis{gauss/cut/omp}


\subsection{Syntax}
\label{\detokenize{pair_gauss:syntax}}
\begin{sphinxVerbatim}[commandchars=\\\{\}]
\PYG{k}{pair\PYGZus{}style}\PYG{+w}{ }\PYG{n}{gauss}\PYG{+w}{ }\PYG{n}{cutoff}
\PYG{k}{pair\PYGZus{}style}\PYG{+w}{ }\PYG{n}{gauss}\PYG{o}{/}\PYG{n}{cut}\PYG{+w}{ }\PYG{n}{cutoff}
\end{sphinxVerbatim}
\begin{itemize}
\item {} 
\sphinxAtStartPar
cutoff = global cutoff for Gauss interactions (distance units)

\end{itemize}


\subsection{Examples}
\label{\detokenize{pair_gauss:examples}}
\begin{sphinxVerbatim}[commandchars=\\\{\}]
\PYG{k}{pair\PYGZus{}style}\PYG{+w}{ }\PYG{n}{gauss}\PYG{+w}{ }\PYG{l+m}{12.0}
\PYG{k}{pair\PYGZus{}coeff}\PYG{+w}{ }\PYG{o}{*}\PYG{+w}{ }\PYG{o}{*}\PYG{+w}{ }\PYG{l+m}{1.0}\PYG{+w}{ }\PYG{l+m}{0.9}
\PYG{k}{pair\PYGZus{}coeff}\PYG{+w}{ }\PYG{l+m}{1}\PYG{+w}{ }\PYG{l+m}{4}\PYG{+w}{ }\PYG{l+m}{1.0}\PYG{+w}{ }\PYG{l+m}{0.9}\PYG{+w}{ }\PYG{l+m}{10.0}

\PYG{k}{pair\PYGZus{}style}\PYG{+w}{ }\PYG{n}{gauss}\PYG{o}{/}\PYG{n}{cut}\PYG{+w}{ }\PYG{l+m}{3.5}
\PYG{k}{pair\PYGZus{}coeff}\PYG{+w}{ }\PYG{l+m}{1}\PYG{+w}{ }\PYG{l+m}{4}\PYG{+w}{ }\PYG{l+m}{0.2805}\PYG{+w}{ }\PYG{l+m}{1.45}\PYG{+w}{ }\PYG{l+m}{0.112}
\end{sphinxVerbatim}


\subsection{Description}
\label{\detokenize{pair_gauss:description}}
\sphinxAtStartPar
Style \sphinxstyleemphasis{gauss} computes a tethering potential of the form
\begin{equation*}
\begin{split}E = - A \exp(-B r^2) \qquad r < r_c\end{split}
\end{equation*}
\sphinxAtStartPar
between an atom and its corresponding tether site which will typically
be a frozen atom in the simulation.  \(r_c\) is the cutoff.

\sphinxAtStartPar
The following coefficients must be defined for each pair of atom types
via the {\hyperref[\detokenize{pair_coeff::doc}]{\sphinxcrossref{\DUrole{doc}{pair\_coeff}}}} command as in the examples above,
or in the data file or restart files read by the
{\hyperref[\detokenize{read_data::doc}]{\sphinxcrossref{\DUrole{doc}{read\_data}}}} or {\hyperref[\detokenize{read_restart::doc}]{\sphinxcrossref{\DUrole{doc}{read\_restart}}}}
commands:
\begin{itemize}
\item {} 
\sphinxAtStartPar
A (energy units)

\item {} 
\sphinxAtStartPar
B (1/distance\textasciicircum{}2 units)

\item {} 
\sphinxAtStartPar
cutoff (distance units)

\end{itemize}

\sphinxAtStartPar
The last coefficient is optional. If not specified, the global cutoff
is used.

\sphinxAtStartPar
Style \sphinxstyleemphasis{gauss/cut} computes a generalized Gaussian interaction potential
between pairs of particles:
\begin{equation*}
\begin{split}E = \frac{H}{\sigma_h\sqrt{2\pi}} \exp\left[-\frac{(r-r_{mh})^2}{2\sigma_h^2}\right]\end{split}
\end{equation*}
\sphinxAtStartPar
where H determines together with the standard deviation \(\sigma_h\)
the peak height of the Gaussian function, and \(r_{mh}\) the peak
position.  Examples of the use of the Gaussian potentials include
implicit solvent simulations of salt ions {\hyperref[\detokenize{pair_gauss:lenart2}]{\sphinxcrossref{\DUrole{std,std-ref}{(Lenart)}}}} and
of surfactants {\hyperref[\detokenize{pair_gauss:jusufi2}]{\sphinxcrossref{\DUrole{std,std-ref}{(Jusufi)}}}}.  In these instances the
Gaussian potential mimics the hydration barrier between a pair of
particles. The hydration barrier is located at \(r_{mh}\) and has a
width of \(\sigma_h\). The prefactor determines the height of the
potential barrier.

\sphinxAtStartPar
The following coefficients must be defined for each pair of atom types
via the {\hyperref[\detokenize{pair_coeff::doc}]{\sphinxcrossref{\DUrole{doc}{pair\_coeff}}}} command as in the example above,
or in the data file or restart files read by the
{\hyperref[\detokenize{read_data::doc}]{\sphinxcrossref{\DUrole{doc}{read\_data}}}} or {\hyperref[\detokenize{read_restart::doc}]{\sphinxcrossref{\DUrole{doc}{read\_restart}}}}
commands:
\begin{itemize}
\item {} 
\sphinxAtStartPar
H (energy * distance units)

\item {} 
\sphinxAtStartPar
\(r_{mh}\) (distance units)

\item {} 
\sphinxAtStartPar
\(\sigma_h\) (distance units)

\item {} 
\sphinxAtStartPar
cutoff (distance units)

\end{itemize}

\sphinxAtStartPar
The last coefficient is optional. If not specified, the global cutoff
is used.


\bigskip\hrule\bigskip


\sphinxAtStartPar
Styles with a \sphinxstyleemphasis{gpu}, \sphinxstyleemphasis{intel}, \sphinxstyleemphasis{kk}, \sphinxstyleemphasis{omp}, or \sphinxstyleemphasis{opt} suffix are
functionally the same as the corresponding style without the suffix.
They have been optimized to run faster, depending on your available
hardware, as discussed on the {\hyperref[\detokenize{Speed_packages::doc}]{\sphinxcrossref{\DUrole{doc}{Accelerator packages}}}}
page.  The accelerated styles take the same arguments and should
produce the same results, except for round\sphinxhyphen{}off and precision issues.

\sphinxAtStartPar
These accelerated styles are part of the GPU, INTEL, KOKKOS,
OPENMP, and OPT packages, respectively.  They are only enabled if
LAMMPS was built with those packages.  See the {\hyperref[\detokenize{Build_package::doc}]{\sphinxcrossref{\DUrole{doc}{Build package}}}} page for more info.

\sphinxAtStartPar
You can specify the accelerated styles explicitly in your input script
by including their suffix, or you can use the {\hyperref[\detokenize{Run_options::doc}]{\sphinxcrossref{\DUrole{doc}{\sphinxhyphen{}suffix command\sphinxhyphen{}line switch}}}} when you invoke LAMMPS, or you can use the
{\hyperref[\detokenize{suffix::doc}]{\sphinxcrossref{\DUrole{doc}{suffix}}}} command in your input script.

\sphinxAtStartPar
See the {\hyperref[\detokenize{Speed_packages::doc}]{\sphinxcrossref{\DUrole{doc}{Accelerator packages}}}} page for more
instructions on how to use the accelerated styles effectively.


\bigskip\hrule\bigskip



\subsection{Mixing, shift, table, tail correction, restart, rRESPA info}
\label{\detokenize{pair_gauss:mixing-shift-table-tail-correction-restart-rrespa-info}}
\sphinxAtStartPar
For atom type pairs I,J and I != J, the A, B, H, sigma\_h, r\_mh
parameters, and the cutoff distance for these pair styles can be mixed:
\begin{itemize}
\item {} 
\sphinxAtStartPar
A (energy units)

\item {} 
\sphinxAtStartPar
\(\sqrt{\frac{1}{B}}\) (distance units, see below)

\item {} 
\sphinxAtStartPar
H (energy units)

\item {} 
\sphinxAtStartPar
\(r_{mh}\) (distance units)

\item {} 
\sphinxAtStartPar
\(\sigma_h\) (distance units)

\item {} 
\sphinxAtStartPar
cutoff (distance units)

\end{itemize}

\sphinxAtStartPar
The default mix value is \sphinxstyleemphasis{geometric}.
Only \sphinxstyleemphasis{arithmetic} and \sphinxstyleemphasis{geometric} mix values are supported.
See the “pair\_modify” command for details.

\sphinxAtStartPar
The A and H parameters are mixed using the same rules normally
used to mix the “epsilon” parameter in a Lennard Jones interaction.
The sigma\_h, r\_mh, and the cutoff distance are mixed using the same
rules used to mix the “sigma” parameter in a Lennard Jones interaction.
The B parameter is converted to a distance (sigma), before mixing
(using sigma=B\textasciicircum{}\sphinxhyphen{}0.5), and converted back to a coefficient
afterwards (using B=sigma\textasciicircum{}2).
Negative A values are converted to positive A values (using abs(A))
before mixing, and converted back after mixing
(by multiplying by min(sign(Ai),sign(Aj))).
This way, if either particle is repulsive (if Ai\textless{}0 or Aj\textless{}0),
then the default interaction between both particles will be repulsive.

\sphinxAtStartPar
For the \sphinxstyleemphasis{gauss} style there is no effect due to the Gaussian well
beyond the cutoff; hence reasonable cutoffs need to be specified.

\sphinxAtStartPar
The \sphinxstyleemphasis{gauss/cut} style supports the {\hyperref[\detokenize{pair_modify::doc}]{\sphinxcrossref{\DUrole{doc}{pair\_modify}}}} shift
option for the energy of the Gauss\sphinxhyphen{}potential portion of the pair
interaction.

\sphinxAtStartPar
The {\hyperref[\detokenize{pair_modify::doc}]{\sphinxcrossref{\DUrole{doc}{pair\_modify}}}} table and tail options are not
relevant for these pair styles.

\sphinxAtStartPar
These pair styles write their information to {\hyperref[\detokenize{restart::doc}]{\sphinxcrossref{\DUrole{doc}{binary restart files}}}}, so pair\_style and pair\_coeff commands do not need to be
specified in an input script that reads a restart file.

\sphinxAtStartPar
These pair styles can only be used via the \sphinxstyleemphasis{pair} keyword of the
{\hyperref[\detokenize{run_style::doc}]{\sphinxcrossref{\DUrole{doc}{run\_style respa}}}} command.  They do not support the
\sphinxstyleemphasis{inner}, \sphinxstyleemphasis{middle}, \sphinxstyleemphasis{outer} keywords.

\sphinxAtStartPar
The \sphinxstyleemphasis{gauss} pair style tallies an “occupancy” count of how many
Gaussian\sphinxhyphen{}well sites have an atom within the distance at which the force
is a maximum = sqrt(0.5/b).  This quantity can be accessed via the
{\hyperref[\detokenize{compute_pair::doc}]{\sphinxcrossref{\DUrole{doc}{compute pair}}}} command as a vector of values of
length 1.

\sphinxAtStartPar
To print this quantity to the log file (with a descriptive column
heading) the following commands could be included in an input script:

\begin{sphinxVerbatim}[commandchars=\\\{\}]
\PYG{k}{compute }\PYG{n+nv+nvIdentifier}{gauss}\PYG{+w}{ }\PYG{n+nv+nvIdentifier}{all}\PYG{+w}{ }\PYG{n}{pair}\PYG{+w}{ }\PYG{n}{gauss}
\PYG{k}{variable }\PYG{n+nv+nvIdentifier}{occ}\PYG{+w}{ }\PYG{n}{equal}\PYG{+w}{ }\PYG{n}{c\PYGZus{}gauss[1]}
\PYG{k}{thermo\PYGZus{}style}\PYG{+w}{ }\PYG{n}{custom}\PYG{+w}{ }\PYG{n}{step}\PYG{+w}{ }\PYG{n}{temp}\PYG{+w}{ }\PYG{n}{epair}\PYG{+w}{ }\PYG{n}{v\PYGZus{}occ}
\end{sphinxVerbatim}


\bigskip\hrule\bigskip



\subsection{Restrictions}
\label{\detokenize{pair_gauss:restrictions}}
\sphinxAtStartPar
The \sphinxstyleemphasis{gauss} and \sphinxstyleemphasis{gauss/cut} styles are part of the EXTRA\sphinxhyphen{}PAIR package.
They are only enabled if LAMMPS is build with that package.  See the
{\hyperref[\detokenize{Build_package::doc}]{\sphinxcrossref{\DUrole{doc}{Build package}}}} page for more info.

\sphinxAtStartPar
\DUrole{versionmodified,changed}{Changed in version 28Mar2023.}

\sphinxAtStartPar
Prior to this version, the \sphinxstyleemphasis{gauss} pair style did not apply
{\hyperref[\detokenize{special_bonds::doc}]{\sphinxcrossref{\DUrole{doc}{special\_bonds}}}} factors.


\subsection{Related commands}
\label{\detokenize{pair_gauss:related-commands}}
\sphinxAtStartPar
{\hyperref[\detokenize{pair_coeff::doc}]{\sphinxcrossref{\DUrole{doc}{pair\_coeff}}}},
{\hyperref[\detokenize{pair_coul_diel::doc}]{\sphinxcrossref{\DUrole{doc}{pair\_style coul/diel}}}}


\subsection{Default}
\label{\detokenize{pair_gauss:default}}
\sphinxAtStartPar
none

\phantomsection\label{\detokenize{pair_gauss:lenart2}}
\sphinxAtStartPar
\sphinxstylestrong{(Lenart)} Lenart , Jusufi, and Panagiotopoulos, J Chem Phys, 126,
044509 (2007).

\phantomsection\label{\detokenize{pair_gauss:jusufi2}}
\sphinxAtStartPar
\sphinxstylestrong{(Jusufi)} Jusufi, Hynninen, and Panagiotopoulos, J Phys Chem B, 112,
13783 (2008).

\sphinxstepscope

\index{pair\_style gayberne@\spxentry{pair\_style gayberne}}\ignorespaces 
\index{pair\_style gayberne/gpu@\spxentry{pair\_style gayberne/gpu}}\ignorespaces 
\index{pair\_style gayberne/intel@\spxentry{pair\_style gayberne/intel}}\ignorespaces 
\index{pair\_style gayberne/omp@\spxentry{pair\_style gayberne/omp}}\ignorespaces 

\section{pair\_style gayberne command}
\label{\detokenize{pair_gayberne:pair-style-gayberne-command}}\label{\detokenize{pair_gayberne:index-3}}\label{\detokenize{pair_gayberne:index-2}}\label{\detokenize{pair_gayberne:index-1}}\label{\detokenize{pair_gayberne:index-0}}\label{\detokenize{pair_gayberne::doc}}
\sphinxAtStartPar
Accelerator Variants: \sphinxstyleemphasis{gayberne/gpu}, \sphinxstyleemphasis{gayberne/intel}, \sphinxstyleemphasis{gayberne/omp}


\subsection{Syntax}
\label{\detokenize{pair_gayberne:syntax}}
\begin{sphinxVerbatim}[commandchars=\\\{\}]
\PYG{k}{pair\PYGZus{}style}\PYG{+w}{ }\PYG{n}{gayberne}\PYG{+w}{ }\PYG{n}{gamma}\PYG{+w}{ }\PYG{n}{upsilon}\PYG{+w}{ }\PYG{n}{mu}\PYG{+w}{ }\PYG{n}{cutoff}
\end{sphinxVerbatim}
\begin{itemize}
\item {} 
\sphinxAtStartPar
gamma = shift for potential minimum (typically 1)

\item {} 
\sphinxAtStartPar
upsilon = exponent for eta orientation\sphinxhyphen{}dependent energy function

\item {} 
\sphinxAtStartPar
mu = exponent for chi orientation\sphinxhyphen{}dependent energy function

\item {} 
\sphinxAtStartPar
cutoff = global cutoff for interactions (distance units)

\end{itemize}


\subsection{Examples}
\label{\detokenize{pair_gayberne:examples}}
\begin{sphinxVerbatim}[commandchars=\\\{\}]
\PYG{k}{pair\PYGZus{}style}\PYG{+w}{ }\PYG{n}{gayberne}\PYG{+w}{ }\PYG{l+m}{1.0}\PYG{+w}{ }\PYG{l+m}{1.0}\PYG{+w}{ }\PYG{l+m}{1.0}\PYG{+w}{ }\PYG{l+m}{10.0}
\PYG{k}{pair\PYGZus{}coeff}\PYG{+w}{ }\PYG{o}{*}\PYG{+w}{ }\PYG{o}{*}\PYG{+w}{ }\PYG{l+m}{1.0}\PYG{+w}{ }\PYG{l+m}{1.7}\PYG{+w}{ }\PYG{l+m}{1.7}\PYG{+w}{ }\PYG{l+m}{3.4}\PYG{+w}{ }\PYG{l+m}{3.4}\PYG{+w}{ }\PYG{l+m}{1.0}\PYG{+w}{ }\PYG{l+m}{1.0}\PYG{+w}{ }\PYG{l+m}{1.0}
\end{sphinxVerbatim}


\subsection{Description}
\label{\detokenize{pair_gayberne:description}}
\sphinxAtStartPar
The \sphinxstyleemphasis{gayberne} styles compute a Gay\sphinxhyphen{}Berne anisotropic LJ interaction
{\hyperref[\detokenize{pair_gayberne:berardi}]{\sphinxcrossref{\DUrole{std,std-ref}{(Berardi)}}}} between pairs of ellipsoidal particles or an
ellipsoidal and spherical particle via the formulas
\begin{equation*}
\begin{split}U ( \mathbf{A}_1, \mathbf{A}_2, \mathbf{r}_{12} ) = & U_r (
\mathbf{A}_1, \mathbf{A}_2, \mathbf{r}_{12}, \gamma ) \cdot \eta_{12} (
\mathbf{A}_1, \mathbf{A}_2, \upsilon ) \cdot \chi_{12} ( \mathbf{A}_1,
\mathbf{A}_2, \mathbf{r}_{12}, \mu ) \\
U_r = & 4 \epsilon ( \varrho^{12} - \varrho^6) \\
\varrho = & \frac{\sigma}{ h_{12} + \gamma \sigma}\end{split}
\end{equation*}
\sphinxAtStartPar
where \(\mathbf{A}_1\) and \(\mathbf{A}_2\) are the
transformation matrices from the simulation box frame to the body frame
and \(r_{12}\) is the center to center vector between the particles.
\(U_r\) controls the shifted distance dependent interaction based on
the distance of closest approach of the two particles (\(h_{12}\))
and the user\sphinxhyphen{}specified shift parameter \(\gamma\).  When both
particles are spherical, the formula reduces to the usual Lennard\sphinxhyphen{}Jones
interaction (see details below for when Gay\sphinxhyphen{}Berne treats a particle as
“spherical”).

\sphinxAtStartPar
For large uniform molecules it has been shown that the energy
parameters are approximately representable in terms of local contact
curvatures {\hyperref[\detokenize{pair_gayberne:everaers2}]{\sphinxcrossref{\DUrole{std,std-ref}{(Everaers)}}}}:
\begin{equation*}
\begin{split}\epsilon_a = \sigma \cdot { \frac{a}{ b \cdot c } }; \epsilon_b =
\sigma \cdot { \frac{b}{ a \cdot c } }; \epsilon_c = \sigma \cdot {
\frac{c}{ a \cdot b } }\end{split}
\end{equation*}
\sphinxAtStartPar
The variable names utilized as potential parameters are for the most
part taken from {\hyperref[\detokenize{pair_gayberne:everaers2}]{\sphinxcrossref{\DUrole{std,std-ref}{(Everaers)}}}} in order to be consistent with
the {\hyperref[\detokenize{pair_resquared::doc}]{\sphinxcrossref{\DUrole{doc}{RE\sphinxhyphen{}squared pair potential}}}}.  Details on the
upsilon and mu parameters are given
\sphinxhref{PDF/pair\_resquared\_extra.pdf}{here}.

\sphinxAtStartPar
More details of the Gay\sphinxhyphen{}Berne formulation are given in the references
listed below and in \sphinxhref{PDF/pair\_gayberne\_extra.pdf}{this supplementary document}.

\sphinxAtStartPar
Use of this pair style requires the NVE, NVT, or NPT fixes with the
\sphinxstyleemphasis{asphere} extension (e.g. {\hyperref[\detokenize{fix_nve_asphere::doc}]{\sphinxcrossref{\DUrole{doc}{fix nve/asphere}}}}) in
order to integrate particle rotation.  Additionally, {\hyperref[\detokenize{atom_style::doc}]{\sphinxcrossref{\DUrole{doc}{atom\_style
ellipsoid}}}} should be used since it defines the rotational
state and the size and shape of each ellipsoidal particle.

\sphinxAtStartPar
The following coefficients must be defined for each pair of atoms
types via the {\hyperref[\detokenize{pair_coeff::doc}]{\sphinxcrossref{\DUrole{doc}{pair\_coeff}}}} command as in the examples
above, or in the data file or restart files read by the
{\hyperref[\detokenize{read_data::doc}]{\sphinxcrossref{\DUrole{doc}{read\_data}}}} or {\hyperref[\detokenize{read_restart::doc}]{\sphinxcrossref{\DUrole{doc}{read\_restart}}}}
commands, or by mixing as described below:
\begin{itemize}
\item {} 
\sphinxAtStartPar
\(\epsilon\) = well depth (energy units)

\item {} 
\sphinxAtStartPar
\(\sigma\) = minimum effective particle radii (distance units)

\item {} 
\sphinxAtStartPar
\(\epsilon_{i,a}\) = relative well depth of type I for side\sphinxhyphen{}to\sphinxhyphen{}side interactions

\item {} 
\sphinxAtStartPar
\(\epsilon_{i,b}\) = relative well depth of type I for face\sphinxhyphen{}to\sphinxhyphen{}face interactions

\item {} 
\sphinxAtStartPar
\(\epsilon_{i,c}\) = relative well depth of type I for end\sphinxhyphen{}to\sphinxhyphen{}end interactions

\item {} 
\sphinxAtStartPar
\(\epsilon_{j,a}\) = relative well depth of type J for side\sphinxhyphen{}to\sphinxhyphen{}side interactions

\item {} 
\sphinxAtStartPar
\(\epsilon_{j,b}\) = relative well depth of type J for face\sphinxhyphen{}to\sphinxhyphen{}face interactions

\item {} 
\sphinxAtStartPar
\(\epsilon_{j,c}\) = relative well depth of type J for end\sphinxhyphen{}to\sphinxhyphen{}end interactions

\item {} 
\sphinxAtStartPar
cutoff (distance units)

\end{itemize}

\sphinxAtStartPar
The last coefficient is optional.  If not specified, the global
cutoff specified in the pair\_style command is used.

\sphinxAtStartPar
It is typical with the Gay\sphinxhyphen{}Berne potential to define \(\sigma\) as
the minimum of the 3 shape diameters of the particles involved in an I,I
interaction, though this is not required.  Note that this is a different
meaning for \(\sigma\) than the {\hyperref[\detokenize{pair_resquared::doc}]{\sphinxcrossref{\DUrole{doc}{pair\_style resquared}}}} potential uses.

\sphinxAtStartPar
The \(\epsilon_i\) and \(\epsilon_j\) coefficients are actually
defined for atom types, not for pairs of atom types.  Thus, in a series
of pair\_coeff commands, they only need to be specified once for each
atom type.

\sphinxAtStartPar
Specifically, if any of \(\epsilon_{i,a}\), \(\epsilon_{i,b}\),
\(\epsilon_{i,c}\) are non\sphinxhyphen{}zero, the three values are assigned to
atom type I.  If all the \(\epsilon_i\) values are zero, they are
ignored.  If any of \(\epsilon_{j,a}\), \(\epsilon_{j,b}\),
\(\epsilon_{j,c}\) are non\sphinxhyphen{}zero, the three values are assigned to
atom type J.  If all three epsilon\_j values are zero, they are ignored.
Thus the typical way to define the \(\epsilon_i\) and
\(\epsilon_j\) coefficients is to list their values in “pair\_coeff
I J” commands when I = J, but set them to 0.0 when I != J.  If you do
list them when I != J, you should ensure they are consistent with their
values in other pair\_coeff commands, since only the last setting will
be in effect.

\sphinxAtStartPar
Note that if this potential is being used as a sub\sphinxhyphen{}style of
{\hyperref[\detokenize{pair_hybrid::doc}]{\sphinxcrossref{\DUrole{doc}{pair\_style hybrid}}}}, and there is no “pair\_coeff I I”
setting made for Gay\sphinxhyphen{}Berne for a particular type I (because I\sphinxhyphen{}I
interactions are computed by another hybrid pair potential), then you
still need to ensure the \(\epsilon\) a,b,c coefficients are assigned to
that type. e.g. in a “pair\_coeff I J” command.

\begin{sphinxadmonition}{note}{Note:}
\sphinxAtStartPar
If the \(\epsilon_{a}\) = \(\epsilon_{b}\) = \(\epsilon_{c}\) for an atom type, and if the shape
of the particle itself is spherical, meaning its 3 shape parameters
are all the same, then the particle is treated as an LJ sphere by the
Gay\sphinxhyphen{}Berne potential.  This is significant because if two LJ spheres
interact, then the simple Lennard\sphinxhyphen{}Jones formula is used to compute
their interaction energy/force using the specified epsilon and sigma
as the standard LJ parameters.  This is much cheaper to compute than
the full Gay\sphinxhyphen{}Berne formula.  To treat the particle as a LJ sphere
with sigma = D, you should normally set \(\epsilon_{a}\) = \(\epsilon_{b}\) = \(\epsilon_{c}\) =
1.0, set the pair\_coeff \(\sigma = D\), and also set the 3 shape
parameters for the particle to D.  The one exception is that if the 3
shape parameters are set to 0.0, which is a valid way in LAMMPS to
specify a point particle, then the Gay\sphinxhyphen{}Berne potential will treat
that as shape parameters of 1.0 (i.e. a LJ particle with
\(\sigma = 1\)), since it requires finite\sphinxhyphen{}size particles.  In
this case you should still set the pair\_coeff \(\sigma\) to 1.0
as well.
\end{sphinxadmonition}


\bigskip\hrule\bigskip


\sphinxAtStartPar
Styles with a \sphinxstyleemphasis{gpu}, \sphinxstyleemphasis{intel}, \sphinxstyleemphasis{kk}, \sphinxstyleemphasis{omp}, or \sphinxstyleemphasis{opt} suffix are
functionally the same as the corresponding style without the suffix.
They have been optimized to run faster, depending on your available
hardware, as discussed on the {\hyperref[\detokenize{Speed_packages::doc}]{\sphinxcrossref{\DUrole{doc}{Accelerator packages}}}}
page.  The accelerated styles take the same arguments and should
produce the same results, except for round\sphinxhyphen{}off and precision issues.

\sphinxAtStartPar
These accelerated styles are part of the GPU, INTEL, KOKKOS,
OPENMP, and OPT packages, respectively.  They are only enabled if
LAMMPS was built with those packages.  See the {\hyperref[\detokenize{Build_package::doc}]{\sphinxcrossref{\DUrole{doc}{Build package}}}} page for more info.

\sphinxAtStartPar
You can specify the accelerated styles explicitly in your input script
by including their suffix, or you can use the {\hyperref[\detokenize{Run_options::doc}]{\sphinxcrossref{\DUrole{doc}{\sphinxhyphen{}suffix command\sphinxhyphen{}line switch}}}} when you invoke LAMMPS, or you can use the
{\hyperref[\detokenize{suffix::doc}]{\sphinxcrossref{\DUrole{doc}{suffix}}}} command in your input script.

\sphinxAtStartPar
See the {\hyperref[\detokenize{Speed_packages::doc}]{\sphinxcrossref{\DUrole{doc}{Accelerator packages}}}} page for more
instructions on how to use the accelerated styles effectively.


\bigskip\hrule\bigskip



\subsection{Mixing, shift, table, tail correction, restart, rRESPA info}
\label{\detokenize{pair_gayberne:mixing-shift-table-tail-correction-restart-rrespa-info}}
\sphinxAtStartPar
For atom type pairs I,J and I != J, the epsilon and sigma coefficients
and cutoff distance for this pair style can be mixed.  The default mix
value is \sphinxstyleemphasis{geometric}.  See the “pair\_modify” command for details.

\sphinxAtStartPar
This pair style supports the {\hyperref[\detokenize{pair_modify::doc}]{\sphinxcrossref{\DUrole{doc}{pair\_modify}}}} shift
option for the energy of the Lennard\sphinxhyphen{}Jones portion of the pair
interaction, but only for sphere\sphinxhyphen{}sphere interactions.  There is no
shifting performed for ellipsoidal interactions due to the anisotropic
dependence of the interaction.

\sphinxAtStartPar
The {\hyperref[\detokenize{pair_modify::doc}]{\sphinxcrossref{\DUrole{doc}{pair\_modify}}}} table option is not relevant
for this pair style.

\sphinxAtStartPar
This pair style does not support the {\hyperref[\detokenize{pair_modify::doc}]{\sphinxcrossref{\DUrole{doc}{pair\_modify}}}}
tail option for adding long\sphinxhyphen{}range tail corrections to energy and
pressure.

\sphinxAtStartPar
This pair style writes its information to {\hyperref[\detokenize{restart::doc}]{\sphinxcrossref{\DUrole{doc}{binary restart files}}}}, so pair\_style and pair\_coeff commands do not need
to be specified in an input script that reads a restart file.

\sphinxAtStartPar
This pair style can only be used via the \sphinxstyleemphasis{pair} keyword of the
{\hyperref[\detokenize{run_style::doc}]{\sphinxcrossref{\DUrole{doc}{run\_style respa}}}} command.  It does not support the
\sphinxstyleemphasis{inner}, \sphinxstyleemphasis{middle}, \sphinxstyleemphasis{outer} keywords.


\bigskip\hrule\bigskip



\subsection{Restrictions}
\label{\detokenize{pair_gayberne:restrictions}}
\sphinxAtStartPar
The \sphinxstyleemphasis{gayberne} style is part of the ASPHERE package.  It is only
enabled if LAMMPS was built with that package.  See the {\hyperref[\detokenize{Build_package::doc}]{\sphinxcrossref{\DUrole{doc}{Build package}}}} page for more info.

\sphinxAtStartPar
These pair styles require that atoms store torque and a quaternion to
represent their orientation, as defined by the
{\hyperref[\detokenize{atom_style::doc}]{\sphinxcrossref{\DUrole{doc}{atom\_style}}}}.  It also require they store a per\sphinxhyphen{}type
{\hyperref[\detokenize{set::doc}]{\sphinxcrossref{\DUrole{doc}{shape}}}}.  The particles cannot store a per\sphinxhyphen{}particle
diameter.

\sphinxAtStartPar
This pair style requires that atoms be ellipsoids as defined by the
{\hyperref[\detokenize{atom_style::doc}]{\sphinxcrossref{\DUrole{doc}{atom\_style ellipsoid}}}} command.

\sphinxAtStartPar
Particles acted on by the potential can be finite\sphinxhyphen{}size aspherical or
spherical particles, or point particles.  Spherical particles have all
3 of their shape parameters equal to each other.  Point particles have
all 3 of their shape parameters equal to 0.0.

\sphinxAtStartPar
The Gay\sphinxhyphen{}Berne potential does not become isotropic as r increases
{\hyperref[\detokenize{pair_gayberne:everaers2}]{\sphinxcrossref{\DUrole{std,std-ref}{(Everaers)}}}}.  The distance\sphinxhyphen{}of\sphinxhyphen{}closest\sphinxhyphen{}approach
approximation used by LAMMPS becomes less accurate when high\sphinxhyphen{}aspect
ratio ellipsoids are used.


\subsection{Related commands}
\label{\detokenize{pair_gayberne:related-commands}}
\sphinxAtStartPar
{\hyperref[\detokenize{pair_coeff::doc}]{\sphinxcrossref{\DUrole{doc}{pair\_coeff}}}}, {\hyperref[\detokenize{fix_nve_asphere::doc}]{\sphinxcrossref{\DUrole{doc}{fix nve/asphere}}}},
{\hyperref[\detokenize{compute_temp_asphere::doc}]{\sphinxcrossref{\DUrole{doc}{compute temp/asphere}}}}, {\hyperref[\detokenize{pair_resquared::doc}]{\sphinxcrossref{\DUrole{doc}{pair\_style resquared}}}}


\subsection{Default}
\label{\detokenize{pair_gayberne:default}}
\sphinxAtStartPar
none


\bigskip\hrule\bigskip

\phantomsection\label{\detokenize{pair_gayberne:everaers2}}
\sphinxAtStartPar
\sphinxstylestrong{(Everaers)} Everaers and Ejtehadi, Phys Rev E, 67, 041710 (2003).

\phantomsection\label{\detokenize{pair_gayberne:berardi}}
\sphinxAtStartPar
\sphinxstylestrong{(Berardi)} Berardi, Fava, Zannoni, Chem Phys Lett, 297, 8\sphinxhyphen{}14 (1998).
Berardi, Muccioli, Zannoni, J Chem Phys, 128, 024905 (2008).

\phantomsection\label{\detokenize{pair_gayberne:perram}}
\sphinxAtStartPar
\sphinxstylestrong{(Perram)} Perram and Rasmussen, Phys Rev E, 54, 6565\sphinxhyphen{}6572 (1996).

\phantomsection\label{\detokenize{pair_gayberne:allen3}}
\sphinxAtStartPar
\sphinxstylestrong{(Allen)} Allen and Germano, Mol Phys 104, 3225\sphinxhyphen{}3235 (2006).

\sphinxstepscope

\index{pair\_style gran/hooke@\spxentry{pair\_style gran/hooke}}\ignorespaces 
\index{pair\_style gran/hooke/omp@\spxentry{pair\_style gran/hooke/omp}}\ignorespaces 
\index{pair\_style gran/hooke/history@\spxentry{pair\_style gran/hooke/history}}\ignorespaces 
\index{pair\_style gran/hooke/history/omp@\spxentry{pair\_style gran/hooke/history/omp}}\ignorespaces 
\index{pair\_style gran/hooke/history/kk@\spxentry{pair\_style gran/hooke/history/kk}}\ignorespaces 
\index{pair\_style gran/hertz/history@\spxentry{pair\_style gran/hertz/history}}\ignorespaces 
\index{pair\_style gran/hertz/history/omp@\spxentry{pair\_style gran/hertz/history/omp}}\ignorespaces 

\section{pair\_style gran/hooke command}
\label{\detokenize{pair_gran:pair-style-gran-hooke-command}}\label{\detokenize{pair_gran:index-6}}\label{\detokenize{pair_gran:index-5}}\label{\detokenize{pair_gran:index-4}}\label{\detokenize{pair_gran:index-3}}\label{\detokenize{pair_gran:index-2}}\label{\detokenize{pair_gran:index-1}}\label{\detokenize{pair_gran:index-0}}\label{\detokenize{pair_gran::doc}}
\sphinxAtStartPar
Accelerator Variants: \sphinxstyleemphasis{gran/hooke/omp}


\section{pair\_style gran/hooke/history command}
\label{\detokenize{pair_gran:pair-style-gran-hooke-history-command}}
\sphinxAtStartPar
Accelerator Variants: \sphinxstyleemphasis{gran/hooke/history/omp}, \sphinxstyleemphasis{gran/hooke/history/kk}


\section{pair\_style gran/hertz/history command}
\label{\detokenize{pair_gran:pair-style-gran-hertz-history-command}}
\sphinxAtStartPar
Accelerator Variants: \sphinxstyleemphasis{gran/hertz/history/omp}


\subsection{Syntax}
\label{\detokenize{pair_gran:syntax}}
\begin{sphinxVerbatim}[commandchars=\\\{\}]
\PYG{k}{pair\PYGZus{}style}\PYG{+w}{ }\PYG{n}{style}\PYG{+w}{ }\PYG{n}{Kn}\PYG{+w}{ }\PYG{n}{Kt}\PYG{+w}{ }\PYG{n}{gamma\PYGZus{}n}\PYG{+w}{ }\PYG{n}{gamma\PYGZus{}t}\PYG{+w}{ }\PYG{n}{xmu}\PYG{+w}{ }\PYG{n}{dampflag}\PYG{+w}{ }\PYG{n}{keyword}
\end{sphinxVerbatim}
\begin{itemize}
\item {} 
\sphinxAtStartPar
style = \sphinxstyleemphasis{gran/hooke} or \sphinxstyleemphasis{gran/hooke/history} or \sphinxstyleemphasis{gran/hertz/history}

\item {} 
\sphinxAtStartPar
Kn = elastic constant for normal particle repulsion (force/distance units or pressure units \sphinxhyphen{} see discussion below)

\item {} 
\sphinxAtStartPar
Kt = elastic constant for tangential contact (force/distance units or pressure units \sphinxhyphen{} see discussion below)

\item {} 
\sphinxAtStartPar
gamma\_n = damping coefficient for collisions in normal direction (1/time units or 1/time\sphinxhyphen{}distance units \sphinxhyphen{} see discussion below)

\item {} 
\sphinxAtStartPar
gamma\_t = damping coefficient for collisions in tangential direction (1/time units or 1/time\sphinxhyphen{}distance units \sphinxhyphen{} see discussion below)

\item {} 
\sphinxAtStartPar
xmu = static yield criterion (unitless value between 0.0 and 1.0e4)

\item {} 
\sphinxAtStartPar
dampflag = 0 or 1 if tangential damping force is excluded or included

\item {} 
\sphinxAtStartPar
keyword = \sphinxstyleemphasis{limit\_damping}
\begin{sphinxalltt}
\sphinxstyleemphasis{limit\_damping} value = none
   limit damping to prevent attractive interaction
\end{sphinxalltt}

\end{itemize}

\begin{sphinxadmonition}{note}{Note:}
\sphinxAtStartPar
Versions of LAMMPS before 9Jan09 had different style names for
granular force fields.  This is to emphasize the fact that the
Hertzian equation has changed to model polydispersity more accurately.
A side effect of the change is that the Kn, Kt, gamma\_n, and gamma\_t
coefficients in the pair\_style command must be specified with
different values in order to reproduce calculations made with earlier
versions of LAMMPS, even for monodisperse systems.  See the NOTE below
for details.
\end{sphinxadmonition}


\subsection{Examples}
\label{\detokenize{pair_gran:examples}}
\begin{sphinxVerbatim}[commandchars=\\\{\}]
\PYG{k}{pair\PYGZus{}style}\PYG{+w}{ }\PYG{n}{gran}\PYG{o}{/}\PYG{n}{hooke}\PYG{o}{/}\PYG{n}{history}\PYG{+w}{ }\PYG{l+m}{200000.0}\PYG{+w}{ }\PYG{n}{NULL}\PYG{+w}{ }\PYG{l+m}{50.0}\PYG{+w}{ }\PYG{n}{NULL}\PYG{+w}{ }\PYG{l+m}{0.5}\PYG{+w}{ }\PYG{l+m}{1}
\PYG{k}{pair\PYGZus{}style}\PYG{+w}{ }\PYG{n}{gran}\PYG{o}{/}\PYG{n}{hooke}\PYG{+w}{ }\PYG{l+m}{200000.0}\PYG{+w}{ }\PYG{l+m}{70000.0}\PYG{+w}{ }\PYG{l+m}{50.0}\PYG{+w}{ }\PYG{l+m}{30.0}\PYG{+w}{ }\PYG{l+m}{0.5}\PYG{+w}{ }\PYG{l+m}{0}
\PYG{k}{pair\PYGZus{}style}\PYG{+w}{ }\PYG{n}{gran}\PYG{o}{/}\PYG{n}{hooke}\PYG{+w}{ }\PYG{l+m}{200000.0}\PYG{+w}{ }\PYG{l+m}{70000.0}\PYG{+w}{ }\PYG{l+m}{50.0}\PYG{+w}{ }\PYG{l+m}{30.0}\PYG{+w}{ }\PYG{l+m}{0.5}\PYG{+w}{ }\PYG{l+m}{0}\PYG{+w}{ }\PYG{n}{limit\PYGZus{}damping}
\end{sphinxVerbatim}


\subsection{Description}
\label{\detokenize{pair_gran:description}}
\sphinxAtStartPar
The \sphinxstyleemphasis{gran} styles use the following formulas for the frictional force
between two granular particles, as described in
{\hyperref[\detokenize{pair_gran:brilliantov}]{\sphinxcrossref{\DUrole{std,std-ref}{(Brilliantov)}}}}, {\hyperref[\detokenize{pair_gran:silbert}]{\sphinxcrossref{\DUrole{std,std-ref}{(Silbert)}}}}, and
{\hyperref[\detokenize{pair_gran:zhang3}]{\sphinxcrossref{\DUrole{std,std-ref}{(Zhang)}}}}, when the distance r between two particles of radii
Ri and Rj is less than their contact distance d = Ri + Rj.  There is
no force between the particles when r \textgreater{} d.

\sphinxAtStartPar
The two Hookean styles use this formula:
\begin{equation*}
\begin{split}F_{hk} = (k_n \delta \mathbf{n}_{ij} -
m_{eff} \gamma_n\mathbf{ v}_n) -
(k_t \boldsymbol{\Delta} \mathbf{s}_t +
m_{eff} \gamma_t \mathbf{v}_t)\end{split}
\end{equation*}
\sphinxAtStartPar
The Hertzian style uses this formula:
\begin{equation*}
\begin{split}F_{hz} = \sqrt{\delta} \sqrt{\frac{R_i R_j}{R_i + R_j}} F_{hk} =
  \sqrt{\delta} \sqrt{\frac{R_i R_j}{R_i + R_j}}
  \Big[ (k_n \delta \mathbf{n}_{ij} -
    m_{eff} \: \gamma_n \mathbf{ v}_n) -
    (k_t \boldsymbol{\Delta} \mathbf{s}_t +
    m_{eff} \: \gamma_t \mathbf{v}_t) \Big]\end{split}
\end{equation*}
\sphinxAtStartPar
In both equations the first parenthesized term is the normal force
between the two particles and the second parenthesized term is the
tangential force.  The normal force has 2 terms, a contact force and a
damping force.  The tangential force also has 2 terms: a shear force
and a damping force.  The shear force is a “history” effect that
accounts for the tangential displacement between the particles for the
duration of the time they are in contact.  This term is included in
pair styles \sphinxstyleemphasis{hooke/history} and \sphinxstyleemphasis{hertz/history}, but is not included
in pair style \sphinxstyleemphasis{hooke}.  The tangential damping force term is included
in all three pair styles if \sphinxstyleemphasis{dampflag} is set to 1; it is not included
if \sphinxstyleemphasis{dampflag} is set to 0.

\sphinxAtStartPar
The other quantities in the equations are as follows:
\begin{itemize}
\item {} 
\sphinxAtStartPar
\(\delta\) = d \sphinxhyphen{} r = overlap distance of 2 particles

\item {} 
\sphinxAtStartPar
\(K_n\) = elastic constant for normal contact

\item {} 
\sphinxAtStartPar
\(K_t\) = elastic constant for tangential contact

\item {} 
\sphinxAtStartPar
\(\gamma_n\) = viscoelastic damping constant for normal contact

\item {} 
\sphinxAtStartPar
\(\gamma_t\) = viscoelastic damping constant for tangential contact

\item {} 
\sphinxAtStartPar
\(m_{eff} = M_i M_j / (M_i + M_j) =\) effective mass of 2 particles of mass M\_i and M\_j

\item {} 
\sphinxAtStartPar
\(\boldsymbol{\Delta} \mathbf{s}_t =\) tangential displacement vector between 2 particles which is truncated to satisfy a frictional yield criterion

\item {} 
\sphinxAtStartPar
\(n_{ij} =\) unit vector along the line connecting the centers of the 2 particles

\item {} 
\sphinxAtStartPar
\(V_n =\) normal component of the relative velocity of the 2 particles

\item {} 
\sphinxAtStartPar
\(V_t =\) tangential component of the relative velocity of the 2 particles

\end{itemize}

\sphinxAtStartPar
The \(K_n\), \(K_t\), \(\gamma_n\), and \(\gamma_t\)
coefficients are specified as parameters to the pair\_style command.  If
a NULL is used for \(K_t\), then a default value is used where
\(K_t = 2/7 K_n\).  If a NULL is used for \(\gamma_t\), then a
default value is used where \(\gamma_t = 1/2 \gamma_n\).

\sphinxAtStartPar
The interpretation and units for these 4 coefficients are different in
the Hookean versus Hertzian equations.

\sphinxAtStartPar
The Hookean model is one where the normal push\sphinxhyphen{}back force for two
overlapping particles is a linear function of the overlap distance.
Thus the specified \(K_n\) is in units of (force/distance).  Note
that this push\sphinxhyphen{}back force is independent of absolute particle size (in
the monodisperse case) and of the relative sizes of the two particles
(in the polydisperse case).  This model also applies to the other terms
in the force equation so that the specified \(\gamma_n\) is in units
of (1/time), \(K_t\) is in units of (force/distance), and
\(\gamma_t\) is in units of (1/time).

\sphinxAtStartPar
The Hertzian model is one where the normal push\sphinxhyphen{}back force for two
overlapping particles is proportional to the area of overlap of the
two particles, and is thus a non\sphinxhyphen{}linear function of overlap distance.
Thus Kn has units of force per area and is thus specified in units of
(pressure).  The effects of absolute particle size (monodispersity)
and relative size (polydispersity) are captured in the radii\sphinxhyphen{}dependent
prefactors.  When these prefactors are carried through to the other
terms in the force equation it means that the specified \(\gamma_n\) is in
units of (1/(time*distance)), \(K_t\) is in units of (pressure), and
\(\gamma_t\) is in units of (1/(time*distance)).

\sphinxAtStartPar
Note that in the Hookean case, \(K_n\) can be thought of as a linear
spring constant with units of force/distance.  In the Hertzian case,
\(K_n\) is like a non\sphinxhyphen{}linear spring constant with units of
force/area or pressure, and as shown in the {\hyperref[\detokenize{pair_gran:zhang3}]{\sphinxcrossref{\DUrole{std,std-ref}{(Zhang)}}}}
paper, \(K_n = 4G / (3(1-\nu))\) where \(\nu =\) the Poisson ratio,
G = shear modulus = \(E / (2(1+\nu))\), and E = Young’s modulus.  Similarly,
\(K_t = 4G / (2-\nu)\).  (NOTE: in an earlier version of the manual, we incorrectly
stated that \(K_t = 8G / (2-\nu)\).)

\sphinxAtStartPar
Thus in the Hertzian case \(K_n\) and \(K_t\) can be set to
values that corresponds to properties of the material being modeled.
This is also true in the Hookean case, except that a spring constant
must be chosen that is appropriate for the absolute size of particles in
the model.  Since relative particle sizes are not accounted for, the
Hookean styles may not be a suitable model for polydisperse systems.

\begin{sphinxadmonition}{note}{Note:}
\sphinxAtStartPar
In versions of LAMMPS before 9Jan09, the equation for Hertzian
interactions did not include the \(\sqrt{r_i r_j / (r_i + r_j)}\)
term and thus was not as accurate for polydisperse systems.  For
monodisperse systems, \(\sqrt{ r_i r_j /(r_i+r_j)}\) is a
constant factor that effectively scales all 4 coefficients:
\(K_n, K_t, \gamma_n, \gamma_t\).  Thus you can set the values of
these 4 coefficients appropriately in the current code to reproduce
the results of a previous Hertzian monodisperse calculation.  For
example, for the common case of a monodisperse system with particles
of diameter 1, all 4 of these coefficients should now be set 2x
larger than they were previously.
\end{sphinxadmonition}

\sphinxAtStartPar
Xmu is also specified in the pair\_style command and is the upper limit
of the tangential force through the Coulomb criterion Ft = xmu*Fn,
where Ft and Fn are the total tangential and normal force components
in the formulas above.  Thus in the Hookean case, the tangential force
between 2 particles grows according to a tangential spring and
dash\sphinxhyphen{}pot model until Ft/Fn = xmu and is then held at Ft = Fn*xmu until
the particles lose contact.  In the Hertzian case, a similar analogy
holds, though the spring is no longer linear.

\begin{sphinxadmonition}{note}{Note:}
\sphinxAtStartPar
Normally, xmu should be specified as a fractional value between
0.0 and 1.0, however LAMMPS allows large values (up to 1.0e4) to allow
for modeling of systems which can sustain very large tangential
forces.
\end{sphinxadmonition}

\sphinxAtStartPar
The effective mass \sphinxstyleemphasis{m\_eff} is given by the formula above for two
isolated particles.  If either particle is part of a rigid body, its
mass is replaced by the mass of the rigid body in the formula above.
This is determined by searching for a {\hyperref[\detokenize{fix_rigid::doc}]{\sphinxcrossref{\DUrole{doc}{fix rigid}}}}
command (or its variants).

\sphinxAtStartPar
For granular styles there are no additional coefficients to set for
each pair of atom types via the {\hyperref[\detokenize{pair_coeff::doc}]{\sphinxcrossref{\DUrole{doc}{pair\_coeff}}}} command.
All settings are global and are made via the pair\_style command.
However you must still use the {\hyperref[\detokenize{pair_coeff::doc}]{\sphinxcrossref{\DUrole{doc}{pair\_coeff}}}} for all
pairs of granular atom types.  For example the command

\begin{sphinxVerbatim}[commandchars=\\\{\}]
\PYG{k}{pair\PYGZus{}coeff}\PYG{+w}{ }\PYG{o}{*}\PYG{+w}{ }\PYG{o}{*}
\end{sphinxVerbatim}

\sphinxAtStartPar
should be used if all atoms in the simulation interact via a granular
potential (i.e. one of the pair styles above is used).  If a granular
potential is used as a sub\sphinxhyphen{}style of {\hyperref[\detokenize{pair_hybrid::doc}]{\sphinxcrossref{\DUrole{doc}{pair\_style hybrid}}}}, then specific atom types can be used in the
pair\_coeff command to determine which atoms interact via a granular
potential.

\sphinxAtStartPar
If two particles are moving away from each other while in contact, there
is a possibility that the particles could experience an effective attractive
force due to damping. If the \sphinxstyleemphasis{limit\_damping} keyword is used, this option
will zero out the normal component of the force if there is an effective
attractive force.


\bigskip\hrule\bigskip


\sphinxAtStartPar
Styles with a \sphinxstyleemphasis{gpu}, \sphinxstyleemphasis{intel}, \sphinxstyleemphasis{kk}, \sphinxstyleemphasis{omp}, or \sphinxstyleemphasis{opt} suffix are
functionally the same as the corresponding style without the suffix.
They have been optimized to run faster, depending on your available
hardware, as discussed on the {\hyperref[\detokenize{Speed_packages::doc}]{\sphinxcrossref{\DUrole{doc}{Accelerator packages}}}}
page.  The accelerated styles take the same arguments and should
produce the same results, except for round\sphinxhyphen{}off and precision issues.

\sphinxAtStartPar
These accelerated styles are part of the GPU, INTEL, KOKKOS,
OPENMP, and OPT packages, respectively.  They are only enabled if
LAMMPS was built with those packages.  See the {\hyperref[\detokenize{Build_package::doc}]{\sphinxcrossref{\DUrole{doc}{Build package}}}} page for more info.

\sphinxAtStartPar
You can specify the accelerated styles explicitly in your input script
by including their suffix, or you can use the {\hyperref[\detokenize{Run_options::doc}]{\sphinxcrossref{\DUrole{doc}{\sphinxhyphen{}suffix command\sphinxhyphen{}line switch}}}} when you invoke LAMMPS, or you can use the
{\hyperref[\detokenize{suffix::doc}]{\sphinxcrossref{\DUrole{doc}{suffix}}}} command in your input script.

\sphinxAtStartPar
See the {\hyperref[\detokenize{Speed_packages::doc}]{\sphinxcrossref{\DUrole{doc}{Accelerator packages}}}} page for more
instructions on how to use the accelerated styles effectively.


\bigskip\hrule\bigskip



\subsection{Mixing, shift, table, tail correction, restart, rRESPA info}
\label{\detokenize{pair_gran:mixing-shift-table-tail-correction-restart-rrespa-info}}
\sphinxAtStartPar
The {\hyperref[\detokenize{pair_modify::doc}]{\sphinxcrossref{\DUrole{doc}{pair\_modify}}}} mix, shift, table, and tail options
are not relevant for granular pair styles.

\sphinxAtStartPar
These pair styles write their information to {\hyperref[\detokenize{restart::doc}]{\sphinxcrossref{\DUrole{doc}{binary restart files}}}}, so a pair\_style command does not need to be
specified in an input script that reads a restart file.

\sphinxAtStartPar
These pair styles can only be used via the \sphinxstyleemphasis{pair} keyword of the
{\hyperref[\detokenize{run_style::doc}]{\sphinxcrossref{\DUrole{doc}{run\_style respa}}}} command.  They do not support the
\sphinxstyleemphasis{inner}, \sphinxstyleemphasis{middle}, \sphinxstyleemphasis{outer} keywords.

\sphinxAtStartPar
The single() function of these pair styles returns 0.0 for the energy
of a pairwise interaction, since energy is not conserved in these
dissipative potentials.  It also returns only the normal component of
the pairwise interaction force.  However, the single() function also
calculates 10 extra pairwise quantities.  The first 3 are the
components of the tangential force between particles I and J, acting
on particle I.  The fourth is the magnitude of this tangential force.
The next 3 (5\sphinxhyphen{}7) are the components of the relative velocity in the
normal direction (along the line joining the 2 sphere centers).  The
last 3 (8\sphinxhyphen{}10) the components of the relative velocity in the
tangential direction.

\sphinxAtStartPar
These extra quantities can be accessed by the {\hyperref[\detokenize{compute_pair_local::doc}]{\sphinxcrossref{\DUrole{doc}{compute pair/local}}}} command, as \sphinxstyleemphasis{p1}, \sphinxstyleemphasis{p2}, …,
\sphinxstyleemphasis{p10}.


\bigskip\hrule\bigskip



\subsection{Restrictions}
\label{\detokenize{pair_gran:restrictions}}
\sphinxAtStartPar
All the granular pair styles are part of the GRANULAR package.  It is
only enabled if LAMMPS was built with that package.  See the {\hyperref[\detokenize{Build_package::doc}]{\sphinxcrossref{\DUrole{doc}{Build package}}}} page for more info.

\sphinxAtStartPar
These pair styles require that atoms store torque and angular velocity
(omega) as defined by the {\hyperref[\detokenize{atom_style::doc}]{\sphinxcrossref{\DUrole{doc}{atom\_style}}}}.  They also
require a per\sphinxhyphen{}particle radius is stored.  The \sphinxstyleemphasis{sphere} atom style does
all of this.

\sphinxAtStartPar
This pair style requires you to use the {\hyperref[\detokenize{comm_modify::doc}]{\sphinxcrossref{\DUrole{doc}{comm\_modify vel yes}}}} command so that velocities are stored by ghost
atoms.

\sphinxAtStartPar
These pair styles will not restart exactly when using the
{\hyperref[\detokenize{read_restart::doc}]{\sphinxcrossref{\DUrole{doc}{read\_restart}}}} command, though they should provide
statistically similar results.  This is because the forces they
compute depend on atom velocities.  See the
{\hyperref[\detokenize{read_restart::doc}]{\sphinxcrossref{\DUrole{doc}{read\_restart}}}} command for more details.

\sphinxAtStartPar
Accumulated values for individual contacts are saved to to restart
files but are not saved to data files. Therefore, forces may
differ significantly when a system is reloaded using A
{\hyperref[\detokenize{read_data::doc}]{\sphinxcrossref{\DUrole{doc}{read\_data}}}} command.


\subsection{Related commands}
\label{\detokenize{pair_gran:related-commands}}
\sphinxAtStartPar
{\hyperref[\detokenize{pair_coeff::doc}]{\sphinxcrossref{\DUrole{doc}{pair\_coeff}}}}


\subsection{Default}
\label{\detokenize{pair_gran:default}}
\sphinxAtStartPar
none


\bigskip\hrule\bigskip

\phantomsection\label{\detokenize{pair_gran:brilliantov}}
\sphinxAtStartPar
\sphinxstylestrong{(Brilliantov)} Brilliantov, Spahn, Hertzsch, Poschel, Phys Rev E, 53,
p 5382\sphinxhyphen{}5392 (1996).

\phantomsection\label{\detokenize{pair_gran:silbert}}
\sphinxAtStartPar
\sphinxstylestrong{(Silbert)} Silbert, Ertas, Grest, Halsey, Levine, Plimpton, Phys Rev
E, 64, p 051302 (2001).

\phantomsection\label{\detokenize{pair_gran:zhang3}}
\sphinxAtStartPar
\sphinxstylestrong{(Zhang)} Zhang and Makse, Phys Rev E, 72, p 011301 (2005).

\sphinxstepscope

\index{pair\_style granular@\spxentry{pair\_style granular}}\ignorespaces 

\section{pair\_style granular command}
\label{\detokenize{pair_granular:pair-style-granular-command}}\label{\detokenize{pair_granular:index-0}}\label{\detokenize{pair_granular::doc}}

\subsection{Syntax}
\label{\detokenize{pair_granular:syntax}}
\begin{sphinxVerbatim}[commandchars=\\\{\}]
\PYG{k}{pair\PYGZus{}style}\PYG{+w}{ }\PYG{n}{granular}\PYG{+w}{ }\PYG{n}{cutoff}
\end{sphinxVerbatim}
\begin{itemize}
\item {} 
\sphinxAtStartPar
cutoff = global cutoff (optional).  See discussion below.

\end{itemize}


\subsection{Examples}
\label{\detokenize{pair_granular:examples}}
\begin{sphinxVerbatim}[commandchars=\\\{\}]
\PYG{k}{pair\PYGZus{}style}\PYG{+w}{ }\PYG{n}{granular}
\PYG{k}{pair\PYGZus{}coeff}\PYG{+w}{ }\PYG{o}{*}\PYG{+w}{ }\PYG{o}{*}\PYG{+w}{ }\PYG{n}{hooke}\PYG{+w}{ }\PYG{l+m}{1000.0}\PYG{+w}{ }\PYG{l+m}{50.0}\PYG{+w}{ }\PYG{n}{tangential}\PYG{+w}{ }\PYG{n}{linear\PYGZus{}nohistory}\PYG{+w}{ }\PYG{l+m}{1.0}\PYG{+w}{ }\PYG{l+m}{0.4}\PYG{+w}{ }\PYG{n}{damping}\PYG{+w}{ }\PYG{n}{mass\PYGZus{}velocity}

\PYG{k}{pair\PYGZus{}style}\PYG{+w}{ }\PYG{n}{granular}
\PYG{k}{pair\PYGZus{}coeff}\PYG{+w}{ }\PYG{o}{*}\PYG{+w}{ }\PYG{o}{*}\PYG{+w}{ }\PYG{n}{hooke}\PYG{+w}{ }\PYG{l+m}{1000.0}\PYG{+w}{ }\PYG{l+m}{50.0}\PYG{+w}{ }\PYG{n}{tangential}\PYG{+w}{ }\PYG{n}{linear\PYGZus{}history}\PYG{+w}{ }\PYG{l+m}{500.0}\PYG{+w}{ }\PYG{l+m}{1.0}\PYG{+w}{ }\PYG{l+m}{0.4}\PYG{+w}{ }\PYG{n}{damping}\PYG{+w}{ }\PYG{n}{mass\PYGZus{}velocity}

\PYG{k}{pair\PYGZus{}style}\PYG{+w}{ }\PYG{n}{granular}
\PYG{k}{pair\PYGZus{}coeff}\PYG{+w}{ }\PYG{o}{*}\PYG{+w}{ }\PYG{o}{*}\PYG{+w}{ }\PYG{n}{hertz}\PYG{+w}{ }\PYG{l+m}{1000.0}\PYG{+w}{ }\PYG{l+m}{50.0}\PYG{+w}{ }\PYG{n}{tangential}\PYG{+w}{ }\PYG{n}{mindlin}\PYG{+w}{ }\PYG{l+m}{1000.0}\PYG{+w}{ }\PYG{l+m}{1.0}\PYG{+w}{ }\PYG{l+m}{0.4}\PYG{+w}{ }\PYG{n}{limit\PYGZus{}damping}

\PYG{k}{pair\PYGZus{}style}\PYG{+w}{ }\PYG{n}{granular}
\PYG{k}{pair\PYGZus{}coeff}\PYG{+w}{ }\PYG{o}{*}\PYG{+w}{ }\PYG{o}{*}\PYG{+w}{ }\PYG{n}{hertz}\PYG{o}{/}\PYG{n}{material}\PYG{+w}{ }\PYG{l+m}{1e8}\PYG{+w}{ }\PYG{l+m}{0.3}\PYG{+w}{ }\PYG{l+m}{0.3}\PYG{+w}{ }\PYG{n}{tangential}\PYG{+w}{ }\PYG{n}{mindlin\PYGZus{}rescale}\PYG{+w}{ }\PYG{n}{NULL}\PYG{+w}{ }\PYG{l+m}{1.0}\PYG{+w}{ }\PYG{l+m}{0.4}\PYG{+w}{ }\PYG{n}{damping}\PYG{+w}{ }\PYG{n}{tsuji}

\PYG{k}{pair\PYGZus{}style}\PYG{+w}{ }\PYG{n}{granular}
\PYG{k}{pair\PYGZus{}coeff}\PYG{+w}{ }\PYG{l+m}{1}\PYG{+w}{ }\PYG{o}{*}\PYG{+w}{ }\PYG{n}{jkr}\PYG{+w}{ }\PYG{l+m}{1000.0}\PYG{+w}{ }\PYG{l+m}{500.0}\PYG{+w}{ }\PYG{l+m}{0.3}\PYG{+w}{ }\PYG{l+m}{10}\PYG{+w}{ }\PYG{n}{tangential}\PYG{+w}{ }\PYG{n}{mindlin}\PYG{+w}{ }\PYG{l+m}{800.0}\PYG{+w}{ }\PYG{l+m}{1.0}\PYG{+w}{ }\PYG{l+m}{0.5}\PYG{+w}{ }\PYG{n}{rolling}\PYG{+w}{ }\PYG{n}{sds}\PYG{+w}{ }\PYG{l+m}{500.0}\PYG{+w}{ }\PYG{l+m}{200.0}\PYG{+w}{ }\PYG{l+m}{0.5}\PYG{+w}{ }\PYG{n}{twisting}\PYG{+w}{ }\PYG{n}{marshall}
\PYG{k}{pair\PYGZus{}coeff}\PYG{+w}{ }\PYG{l+m}{2}\PYG{+w}{ }\PYG{l+m}{2}\PYG{+w}{ }\PYG{n}{hertz}\PYG{+w}{ }\PYG{l+m}{200.0}\PYG{+w}{ }\PYG{l+m}{100.0}\PYG{+w}{ }\PYG{n}{tangential}\PYG{+w}{ }\PYG{n}{linear\PYGZus{}history}\PYG{+w}{ }\PYG{l+m}{300.0}\PYG{+w}{ }\PYG{l+m}{1.0}\PYG{+w}{ }\PYG{l+m}{0.1}\PYG{+w}{ }\PYG{n}{rolling}\PYG{+w}{ }\PYG{n}{sds}\PYG{+w}{ }\PYG{l+m}{200.0}\PYG{+w}{ }\PYG{l+m}{100.0}\PYG{+w}{ }\PYG{l+m}{0.1}\PYG{+w}{ }\PYG{n}{twisting}\PYG{+w}{ }\PYG{n}{marshall}

\PYG{k}{pair\PYGZus{}style}\PYG{+w}{ }\PYG{n}{granular}
\PYG{k}{pair\PYGZus{}coeff}\PYG{+w}{ }\PYG{l+m}{1}\PYG{+w}{ }\PYG{l+m}{1}\PYG{+w}{ }\PYG{n}{dmt}\PYG{+w}{ }\PYG{l+m}{1000.0}\PYG{+w}{ }\PYG{l+m}{50.0}\PYG{+w}{ }\PYG{l+m}{0.3}\PYG{+w}{ }\PYG{l+m}{0.0}\PYG{+w}{ }\PYG{n}{tangential}\PYG{+w}{ }\PYG{n}{mindlin}\PYG{+w}{ }\PYG{n}{NULL}\PYG{+w}{ }\PYG{l+m}{0.5}\PYG{+w}{ }\PYG{l+m}{0.5}\PYG{+w}{ }\PYG{n}{rolling}\PYG{+w}{ }\PYG{n}{sds}\PYG{+w}{ }\PYG{l+m}{500.0}\PYG{+w}{ }\PYG{l+m}{200.0}\PYG{+w}{ }\PYG{l+m}{0.5}\PYG{+w}{ }\PYG{n}{twisting}\PYG{+w}{ }\PYG{n}{marshall}
\PYG{k}{pair\PYGZus{}coeff}\PYG{+w}{ }\PYG{l+m}{2}\PYG{+w}{ }\PYG{l+m}{2}\PYG{+w}{ }\PYG{n}{dmt}\PYG{+w}{ }\PYG{l+m}{1000.0}\PYG{+w}{ }\PYG{l+m}{50.0}\PYG{+w}{ }\PYG{l+m}{0.3}\PYG{+w}{ }\PYG{l+m}{10.0}\PYG{+w}{ }\PYG{n}{tangential}\PYG{+w}{ }\PYG{n}{mindlin}\PYG{+w}{ }\PYG{n}{NULL}\PYG{+w}{ }\PYG{l+m}{0.5}\PYG{+w}{ }\PYG{l+m}{0.1}\PYG{+w}{ }\PYG{n}{rolling}\PYG{+w}{ }\PYG{n}{sds}\PYG{+w}{ }\PYG{l+m}{500.0}\PYG{+w}{ }\PYG{l+m}{200.0}\PYG{+w}{ }\PYG{l+m}{0.1}\PYG{+w}{ }\PYG{n}{twisting}\PYG{+w}{ }\PYG{n}{marshall}

\PYG{k}{pair\PYGZus{}style}\PYG{+w}{ }\PYG{n}{granular}
\PYG{k}{pair\PYGZus{}coeff}\PYG{+w}{ }\PYG{o}{*}\PYG{+w}{ }\PYG{o}{*}\PYG{+w}{ }\PYG{n}{hertz}\PYG{+w}{ }\PYG{l+m}{1000.0}\PYG{+w}{ }\PYG{l+m}{50.0}\PYG{+w}{ }\PYG{n}{tangential}\PYG{+w}{ }\PYG{n}{mindlin}\PYG{+w}{ }\PYG{l+m}{1000.0}\PYG{+w}{ }\PYG{l+m}{1.0}\PYG{+w}{ }\PYG{l+m}{0.4}\PYG{+w}{ }\PYG{n}{heat}\PYG{+w}{ }\PYG{n}{area}\PYG{+w}{ }\PYG{l+m}{0.1}
\end{sphinxVerbatim}


\subsection{Description}
\label{\detokenize{pair_granular:description}}
\sphinxAtStartPar
The \sphinxstyleemphasis{granular} styles support a variety of options for the normal,
tangential, rolling and twisting forces resulting from contact between
two granular particles. This expands on the options offered by the
{\hyperref[\detokenize{pair_gran::doc}]{\sphinxcrossref{\DUrole{doc}{pair gran/*}}}} pair styles. The total computed forces
and torques are the sum of various models selected for the normal,
tangential, rolling and twisting modes of motion.

\sphinxAtStartPar
All model choices and parameters are entered in the
{\hyperref[\detokenize{pair_coeff::doc}]{\sphinxcrossref{\DUrole{doc}{pair\_coeff}}}} command, as described below.  Unlike
e.g. {\hyperref[\detokenize{pair_gran::doc}]{\sphinxcrossref{\DUrole{doc}{pair gran/hooke}}}}, coefficient values are not
global, but can be set to different values for different combinations
of particle types, as determined by the {\hyperref[\detokenize{pair_coeff::doc}]{\sphinxcrossref{\DUrole{doc}{pair\_coeff}}}}
command.  If the contact model choice is the same for two particle
types, the mixing for the cross\sphinxhyphen{}coefficients can be carried out
automatically. This is shown in the last example, where model
choices are the same for type 1 \sphinxhyphen{} type 1 as for type 2 \sphinxhyphen{} type2
interactions, but coefficients are different. In this case, the
mixed coefficients for type 1 \sphinxhyphen{} type 2 interactions can be determined from
mixing rules discussed below.  For additional flexibility,
coefficients as well as model forms can vary between particle types,
as shown in the fourth example: type 1 \sphinxhyphen{} type 1 interactions are based
on a Johnson\sphinxhyphen{}Kendall\sphinxhyphen{}Roberts normal contact model and 2\sphinxhyphen{}2 interactions
are based on a DMT cohesive model (see below).  In that example, 1\sphinxhyphen{}1
and 2\sphinxhyphen{}2 interactions have different model forms, in which case mixing of
coefficients cannot be determined, so 1\sphinxhyphen{}2 interactions must be
explicitly defined via the \sphinxstyleemphasis{pair\_coeff 1 *} command, otherwise an
error would result.


\bigskip\hrule\bigskip


\sphinxAtStartPar
The first required keyword for the \sphinxstyleemphasis{pair\_coeff} command is the normal
contact model. Currently supported options for normal contact models
and their required arguments are:
\begin{enumerate}
\sphinxsetlistlabels{\arabic}{enumi}{enumii}{}{.}%
\item {} 
\sphinxAtStartPar
\sphinxstyleemphasis{hooke} : \(k_n\), \(\eta_{n0}\) (or \(e\))

\item {} 
\sphinxAtStartPar
\sphinxstyleemphasis{hertz} : \(k_n\), \(\eta_{n0}\) (or \(e\))

\item {} 
\sphinxAtStartPar
\sphinxstyleemphasis{hertz/material} : E, \(\eta_{n0}\) (or \(e\)), \(\nu\)

\item {} 
\sphinxAtStartPar
\sphinxstyleemphasis{dmt} : E, \(\eta_{n0}\) (or \(e\)), \(\nu\), \(\gamma\)

\item {} 
\sphinxAtStartPar
\sphinxstyleemphasis{jkr} : E, \(\eta_{n0}\) (or \(e\)), \(\nu\), \(\gamma\)

\end{enumerate}

\sphinxAtStartPar
Here, \(k_n\) is spring stiffness (with units that depend on model
choice, see below); \(\eta_{n0}\) is a damping prefactor (or, in its
place a coefficient of restitution \(e\), depending on the choice of
damping mode, see below); E is Young’s modulus in units of
\sphinxstyleemphasis{force}/\sphinxstyleemphasis{length}\textasciicircum{}2, i.e. \sphinxstyleemphasis{pressure}; \(\nu\) is Poisson’s ratio and
\(\gamma\) is a surface energy density, in units of
\sphinxstyleemphasis{energy}/\sphinxstyleemphasis{length}\textasciicircum{}2.

\sphinxAtStartPar
For the \sphinxstyleemphasis{hooke} model, the normal, elastic component of force acting
on particle \sphinxstyleemphasis{i} due to contact with particle \sphinxstyleemphasis{j} is given by:
\begin{equation*}
\begin{split}\mathbf{F}_{ne, Hooke} = k_n \delta_{ij} \mathbf{n}\end{split}
\end{equation*}
\sphinxAtStartPar
Where \(\delta_{ij} = R_i + R_j - \|\mathbf{r}_{ij}\|\) is the particle
overlap, \(R_i, R_j\) are the particle radii, \(\mathbf{r}_{ij} = \mathbf{r}_i - \mathbf{r}_j\) is the vector separating the two
particle centers (note the i\sphinxhyphen{}j ordering so that \(\mathbf{F}_{ne}\) is
positive for repulsion), and \(\mathbf{n} = \frac{\mathbf{r}_{ij}}{\|\mathbf{r}_{ij}\|}\).  Therefore,
for \sphinxstyleemphasis{hooke}, the units of the spring constant \(k_n\) are
\sphinxstyleemphasis{force}/\sphinxstyleemphasis{distance}, or equivalently \sphinxstyleemphasis{mass}/\sphinxstyleemphasis{time\textasciicircum{}2}.

\sphinxAtStartPar
For the \sphinxstyleemphasis{hertz} model, the normal component of force is given by:
\begin{equation*}
\begin{split}\mathbf{F}_{ne, Hertz} = k_n R_{eff}^{1/2}\delta_{ij}^{3/2} \mathbf{n}\end{split}
\end{equation*}
\sphinxAtStartPar
Here, \(R_{eff} = R = \frac{R_i R_j}{R_i + R_j}\) is the effective
radius, denoted for simplicity as \sphinxstyleemphasis{R} from here on.  For \sphinxstyleemphasis{hertz}, the
units of the spring constant \(k_n\) are \sphinxstyleemphasis{force}/\sphinxstyleemphasis{length}\textasciicircum{}2, or
equivalently \sphinxstyleemphasis{pressure}.

\sphinxAtStartPar
For the \sphinxstyleemphasis{hertz/material} model, the force is given by:
\begin{equation*}
\begin{split}\mathbf{F}_{ne, Hertz/material} = \frac{4}{3} E_{eff} R^{1/2}\delta_{ij}^{3/2} \mathbf{n}\end{split}
\end{equation*}
\sphinxAtStartPar
Here, \(E_{eff} = E = \left(\frac{1-\nu_i^2}{E_i} + \frac{1-\nu_j^2}{E_j}\right)^{-1}\)
is the effective Young’s modulus, with \(\nu_i, \nu_j\) the Poisson ratios
of the particles of types \sphinxstyleemphasis{i} and \sphinxstyleemphasis{j}. \(E_{eff}\) is denoted as \sphinxstyleemphasis{E} from here on.
Note that if the elastic modulus and the shear modulus of the two particles are the
same, the \sphinxstyleemphasis{hertz/material} model is equivalent to the \sphinxstyleemphasis{hertz} model with
\(k_n = 4/3 E\)

\sphinxAtStartPar
The \sphinxstyleemphasis{dmt} model corresponds to the
{\hyperref[\detokenize{pair_granular:dmt1975}]{\sphinxcrossref{\DUrole{std,std-ref}{(Derjaguin\sphinxhyphen{}Muller\sphinxhyphen{}Toporov)}}}} cohesive model, where the force
is simply Hertz with an additional attractive cohesion term:
\begin{equation*}
\begin{split}\mathbf{F}_{ne, dmt} = \left(\frac{4}{3} E R^{1/2}\delta_{ij}^{3/2} - 4\pi\gamma R\right)\mathbf{n}\end{split}
\end{equation*}
\sphinxAtStartPar
The \sphinxstyleemphasis{jkr} model is the {\hyperref[\detokenize{pair_granular:jkr1971}]{\sphinxcrossref{\DUrole{std,std-ref}{(Johnson\sphinxhyphen{}Kendall\sphinxhyphen{}Roberts)}}}} model,
where the force is computed as:
\begin{equation*}
\begin{split}\mathbf{F}_{ne, jkr} = \left(\frac{4Ea^3}{3R} - 2\pi a^2\sqrt{\frac{4\gamma E}{\pi a}}\right)\mathbf{n}\end{split}
\end{equation*}
\sphinxAtStartPar
Here, \(a\) is the radius of the contact zone, related to the overlap
\(\delta\) according to:
\begin{equation*}
\begin{split}\delta = a^2/R - 2\sqrt{\pi \gamma a/E}\end{split}
\end{equation*}
\sphinxAtStartPar
LAMMPS internally inverts the equation above to solve for \sphinxstyleemphasis{a} in terms
of \(\delta\), then solves for the force in the previous
equation. Additionally, note that the JKR model allows for a tensile
force beyond contact (i.e. for \(\delta < 0\)), up to a maximum of
\(3\pi\gamma R\) (also known as the ‘pull\sphinxhyphen{}off’ force).  Note that this
is a hysteretic effect, where particles that are not contacting
initially will not experience force until they come into contact
\(\delta \geq 0\); as they move apart and (\(\delta < 0\)), they
experience a tensile force up to \(3\pi\gamma R\), at which point they
lose contact.


\bigskip\hrule\bigskip


\sphinxAtStartPar
In addition, the normal force is augmented by a damping term of the
following general form:
\begin{equation*}
\begin{split}\mathbf{F}_{n,damp} = -\eta_n \mathbf{v}_{n,rel}\end{split}
\end{equation*}
\sphinxAtStartPar
Here, \(\mathbf{v}_{n,rel} = (\mathbf{v}_j - \mathbf{v}_i) \cdot \mathbf{n}\ \mathbf{n}\) is the component of relative velocity along
\(\mathbf{n}\).

\sphinxAtStartPar
The optional \sphinxstyleemphasis{damping} keyword to the \sphinxstyleemphasis{pair\_coeff} command followed by
a keyword determines the model form of the damping factor \(\eta_n\),
and the interpretation of the \(\eta_{n0}\) or \(e\) coefficients
specified as part of the normal contact model settings. The \sphinxstyleemphasis{damping}
keyword and corresponding model form selection may be appended
anywhere in the \sphinxstyleemphasis{pair coeff} command.  Note that the choice of damping
model affects both the normal and tangential damping (and depending on
other settings, potentially also the twisting damping).  The options
for the damping model currently supported are:
\begin{enumerate}
\sphinxsetlistlabels{\arabic}{enumi}{enumii}{}{.}%
\item {} 
\sphinxAtStartPar
\sphinxstyleemphasis{velocity}

\item {} 
\sphinxAtStartPar
\sphinxstyleemphasis{mass\_velocity}

\item {} 
\sphinxAtStartPar
\sphinxstyleemphasis{viscoelastic}

\item {} 
\sphinxAtStartPar
\sphinxstyleemphasis{tsuji}

\item {} 
\sphinxAtStartPar
\sphinxstyleemphasis{coeff\_restitution}

\end{enumerate}

\sphinxAtStartPar
If the \sphinxstyleemphasis{damping} keyword is not specified, the \sphinxstyleemphasis{viscoelastic} model is
used by default.

\sphinxAtStartPar
For \sphinxstyleemphasis{damping velocity}, the normal damping is simply equal to the
user\sphinxhyphen{}specified damping coefficient in the \sphinxstyleemphasis{normal} model:
\begin{equation*}
\begin{split}\eta_n = \eta_{n0}\end{split}
\end{equation*}
\sphinxAtStartPar
Here, \(\eta_{n0}\) is the damping coefficient specified for the normal
contact model, in units of \sphinxstyleemphasis{mass}/\sphinxstyleemphasis{time}.

\sphinxAtStartPar
For \sphinxstyleemphasis{damping mass\_velocity}, the normal damping is given by:
\begin{equation*}
\begin{split}\eta_n = \eta_{n0} m_{eff}\end{split}
\end{equation*}
\sphinxAtStartPar
Here, \(\eta_{n0}\) is the damping coefficient specified for the normal
contact model, in units of 1/\sphinxstyleemphasis{time} and
\(m_{eff} = m_i m_j/(m_i + m_j)\) is the effective mass.
Use \sphinxstyleemphasis{damping mass\_velocity} to reproduce the damping behavior of
\sphinxstyleemphasis{pair gran/hooke/*}.

\sphinxAtStartPar
The \sphinxstyleemphasis{damping viscoelastic} model is based on the viscoelastic
treatment of {\hyperref[\detokenize{pair_granular:brill1996}]{\sphinxcrossref{\DUrole{std,std-ref}{(Brilliantov et al)}}}}, where the normal
damping is given by:
\begin{equation*}
\begin{split}\eta_n = \eta_{n0}\ a m_{eff}\end{split}
\end{equation*}
\sphinxAtStartPar
Here, \sphinxstyleemphasis{a} is the contact radius, given by \(a =\sqrt{R\delta}\)
for all models except \sphinxstyleemphasis{jkr}, for which it is given implicitly according
to \(\delta = a^2/R - 2\sqrt{\pi \gamma a/E}\).  For \sphinxstyleemphasis{damping viscoelastic},
\(\eta_{n0}\) is in units of 1/(\sphinxstyleemphasis{time}*\sphinxstyleemphasis{distance}).

\sphinxAtStartPar
The \sphinxstyleemphasis{tsuji} model is based on the work of {\hyperref[\detokenize{pair_granular:tsuji1992}]{\sphinxcrossref{\DUrole{std,std-ref}{(Tsuji et al)}}}}.
Here, the damping coefficient specified as part of the normal model is interpreted
as a restitution coefficient \(e\). The damping constant \(\eta_n\) is
given by:
\begin{equation*}
\begin{split}\eta_n = \alpha (m_{eff}k_n)^{1/2}\end{split}
\end{equation*}
\sphinxAtStartPar
For normal contact models based on material parameters, \(k_n = 4/3Ea\). This
damping model is not compatible with cohesive normal models such as \sphinxstyleemphasis{JKR} or \sphinxstyleemphasis{DMT}.
The parameter \(\alpha\) is related to the restitution coefficient \sphinxstyleemphasis{e}
according to:
\begin{equation*}
\begin{split}\alpha = 1.2728-4.2783e+11.087e^2-22.348e^3+27.467e^4-18.022e^5+4.8218e^6\end{split}
\end{equation*}
\sphinxAtStartPar
The dimensionless coefficient of restitution \(e\) specified as part
of the normal contact model parameters should be between 0 and 1, but
no error check is performed on this.

\sphinxAtStartPar
The \sphinxstyleemphasis{coeff\_restitution} model is useful when a specific normal coefficient of
restitution \(e\) is required. In these models, the normal coefficient of
restitution \(e\) is specified as an input. Following the approach of
{\hyperref[\detokenize{pair_granular:brill1996}]{\sphinxcrossref{\DUrole{std,std-ref}{(Brilliantov et al)}}}}, when using the \sphinxstyleemphasis{hooke} normal model,
\sphinxstyleemphasis{coeff\_restitution} calculates the damping coefficient as:
\begin{equation*}
\begin{split}\eta_n = \sqrt{\frac{4m_{eff}k_n}{1+\left( \frac{\pi}{\log(e)}\right)^2}} ,\end{split}
\end{equation*}
\sphinxAtStartPar
For any other normal model, e.g. the \sphinxstyleemphasis{hertz} and \sphinxstyleemphasis{hertz/material} models, the damping
coefficient is:
\begin{equation*}
\begin{split}\eta_n = -2\sqrt{\frac{5}{6}}\frac{\log(e)}{\sqrt{\pi^2+(\log(e))^2}}(R_{eff} \delta_{ij})^{\frac{1}{4}}\sqrt{\frac{3}{2}k_n m_{eff}} ,\end{split}
\end{equation*}
\sphinxAtStartPar
where \(k_n = \frac{4}{3} E_{eff}\) for the \sphinxstyleemphasis{hertz/material} model. Since
\sphinxstyleemphasis{coeff\_restitution} accounts for the effective mass, effective radius, and
pairwise overlaps (except when used with the \sphinxstyleemphasis{hooke} normal model) when calculating
the damping coefficient, it accurately reproduces the specified coefficient of
restitution for both monodisperse and polydisperse particle pairs.  This damping
model is not compatible with cohesive normal models such as \sphinxstyleemphasis{JKR} or \sphinxstyleemphasis{DMT}.

\sphinxAtStartPar
The total normal force is computed as the sum of the elastic and
damping components:
\begin{equation*}
\begin{split}\mathbf{F}_n = \mathbf{F}_{ne} + \mathbf{F}_{n,damp}\end{split}
\end{equation*}

\bigskip\hrule\bigskip


\sphinxAtStartPar
The \sphinxstyleemphasis{pair\_coeff} command also requires specification of the tangential
contact model. The required keyword \sphinxstyleemphasis{tangential} is expected, followed
by the model choice and associated parameters. Currently supported
tangential model choices and their expected parameters are as follows:
\begin{enumerate}
\sphinxsetlistlabels{\arabic}{enumi}{enumii}{}{.}%
\item {} 
\sphinxAtStartPar
\sphinxstyleemphasis{linear\_nohistory} : \(x_{\gamma,t}\), \(\mu_s\)

\item {} 
\sphinxAtStartPar
\sphinxstyleemphasis{linear\_history} : \(k_t\), \(x_{\gamma,t}\), \(\mu_s\)

\item {} 
\sphinxAtStartPar
\sphinxstyleemphasis{mindlin} : \(k_t\) or NULL, \(x_{\gamma,t}\), \(\mu_s\)

\item {} 
\sphinxAtStartPar
\sphinxstyleemphasis{mindlin/force} : \(k_t\) or NULL, \(x_{\gamma,t}\), \(\mu_s\)

\item {} 
\sphinxAtStartPar
\sphinxstyleemphasis{mindlin\_rescale} : \(k_t\) or NULL, \(x_{\gamma,t}\), \(\mu_s\)

\item {} 
\sphinxAtStartPar
\sphinxstyleemphasis{mindlin\_rescale/force} : \(k_t\) or NULL, \(x_{\gamma,t}\), \(\mu_s\)

\end{enumerate}

\sphinxAtStartPar
Here, \(x_{\gamma,t}\) is a dimensionless multiplier for the normal
damping \(\eta_n\) that determines the magnitude of the tangential
damping, \(\mu_t\) is the tangential (or sliding) friction
coefficient, and \(k_t\) is the tangential stiffness coefficient.

\sphinxAtStartPar
For \sphinxstyleemphasis{tangential linear\_nohistory}, a simple velocity\sphinxhyphen{}dependent Coulomb
friction criterion is used, which mimics the behavior of the \sphinxstyleemphasis{pair
gran/hooke} style. The tangential force \(\mathbf{F}_t\) is given by:
\begin{equation*}
\begin{split}\mathbf{F}_t =  -\min(\mu_t F_{n0}, \|\mathbf{F}_\mathrm{t,damp}\|) \mathbf{t}\end{split}
\end{equation*}
\sphinxAtStartPar
The tangential damping force \(\mathbf{F}_\mathrm{t,damp}\) is given by:
\begin{equation*}
\begin{split}\mathbf{F}_\mathrm{t,damp} = -\eta_t \mathbf{v}_{t,rel}\end{split}
\end{equation*}
\sphinxAtStartPar
The tangential damping prefactor \(\eta_t\) is calculated by scaling
the normal damping \(\eta_n\) (see above):
\begin{equation*}
\begin{split}\eta_t = -x_{\gamma,t} \eta_n\end{split}
\end{equation*}
\sphinxAtStartPar
The normal damping prefactor \(\eta_n\) is determined by the choice
of the \sphinxstyleemphasis{damping} keyword, as discussed above.  Thus, the \sphinxstyleemphasis{damping}
keyword also affects the tangential damping.  The parameter
\(x_{\gamma,t}\) is a scaling coefficient. Several works in the
literature use \(x_{\gamma,t} = 1\) ({\hyperref[\detokenize{pair_granular:marshall2009}]{\sphinxcrossref{\DUrole{std,std-ref}{Marshall}}}},
{\hyperref[\detokenize{pair_granular:tsuji1992}]{\sphinxcrossref{\DUrole{std,std-ref}{Tsuji et al}}}}, {\hyperref[\detokenize{pair_granular:silbert2001}]{\sphinxcrossref{\DUrole{std,std-ref}{Silbert et al}}}}).  The relative
tangential velocity at the point of contact is given by
\(\mathbf{v}_{t, rel} = \mathbf{v}_{t} - (R_i\boldsymbol{\Omega}_i + R_j\boldsymbol{\Omega}_j) \times \mathbf{n}\), where \(\mathbf{v}_{t} = \mathbf{v}_r - \mathbf{v}_r\cdot\mathbf{n}\ \mathbf{n}\),
\(\mathbf{v}_r = \mathbf{v}_j - \mathbf{v}_i\) .
The direction of the applied force is \(\mathbf{t} = \mathbf{v_{t,rel}}/\|\mathbf{v_{t,rel}}\|\) .

\sphinxAtStartPar
The normal force value \(F_{n0}\) used to compute the critical force
depends on the form of the contact model. For non\sphinxhyphen{}cohesive models
(\sphinxstyleemphasis{hertz}, \sphinxstyleemphasis{hertz/material}, \sphinxstyleemphasis{hooke}), it is given by the magnitude of
the normal force:
\begin{equation*}
\begin{split}F_{n0} = \|\mathbf{F}_n\|\end{split}
\end{equation*}
\sphinxAtStartPar
For cohesive models such as \sphinxstyleemphasis{jkr} and \sphinxstyleemphasis{dmt}, the critical force is
adjusted so that the critical tangential force approaches \(\mu_t F_{pulloff}\), see {\hyperref[\detokenize{pair_granular:marshall2009}]{\sphinxcrossref{\DUrole{std,std-ref}{Marshall}}}}, equation 43, and
{\hyperref[\detokenize{pair_granular:thornton1991}]{\sphinxcrossref{\DUrole{std,std-ref}{Thornton}}}}.  For both models, \(F_{n0}\) takes the
form:
\begin{equation*}
\begin{split}F_{n0} = \|\mathbf{F}_{ne} + 2 F_{pulloff}\|\end{split}
\end{equation*}
\sphinxAtStartPar
Where \(F_{pulloff} = 3\pi \gamma R\) for \sphinxstyleemphasis{jkr}, and
\(F_{pulloff} = 4\pi \gamma R\) for \sphinxstyleemphasis{dmt}.

\sphinxAtStartPar
The remaining tangential options all use accumulated tangential
displacement (i.e. contact history), except for the options
\sphinxstyleemphasis{mindlin/force} and \sphinxstyleemphasis{mindlin\_rescale/force}, that use accumulated
tangential force instead, and are discussed further below.
The accumulated tangential displacement is discussed in details below
in the context of the \sphinxstyleemphasis{linear\_history} option. The same treatment of
the accumulated displacement applies to the other options as well.

\sphinxAtStartPar
For \sphinxstyleemphasis{tangential linear\_history}, the tangential force is given by:
\begin{equation*}
\begin{split}\mathbf{F}_t =  -\min(\mu_t F_{n0}, \|-k_t\mathbf{\xi} + \mathbf{F}_\mathrm{t,damp}\|) \mathbf{t}\end{split}
\end{equation*}
\sphinxAtStartPar
Here, \(\mathbf{\xi}\) is the tangential displacement accumulated
during the entire duration of the contact:
\begin{equation*}
\begin{split}\mathbf{\xi} = \int_{t0}^t \mathbf{v}_{t,rel}(\tau) \mathrm{d}\tau\end{split}
\end{equation*}
\sphinxAtStartPar
This accumulated tangential displacement must be adjusted to account
for changes in the frame of reference of the contacting pair of
particles during contact. This occurs due to the overall motion of the
contacting particles in a rigid\sphinxhyphen{}body\sphinxhyphen{}like fashion during the duration
of the contact. There are two modes of motion that are relevant: the
‘tumbling’ rotation of the contacting pair, which changes the
orientation of the plane in which tangential displacement occurs; and
‘spinning’ rotation of the contacting pair about the vector connecting
their centers of mass (\(\mathbf{n}\)).  Corrections due to the
former mode of motion are made by rotating the accumulated
displacement into the plane that is tangential to the contact vector
at each step, or equivalently removing any component of the tangential
displacement that lies along \(\mathbf{n}\), and rescaling to
preserve the magnitude.  This follows the discussion in
{\hyperref[\detokenize{pair_granular:luding2008}]{\sphinxcrossref{\DUrole{std,std-ref}{Luding}}}}, see equation 17 and relevant discussion in that
work:
\begin{equation*}
\begin{split}\mathbf{\xi} = \left(\mathbf{\xi'} - (\mathbf{n} \cdot \mathbf{\xi'})\mathbf{n}\right) \frac{\|\mathbf{\xi'}\|}{\|\mathbf{\xi'} - (\mathbf{n}\cdot\mathbf{\xi'})\mathbf{n}\|}\end{split}
\end{equation*}
\sphinxAtStartPar
Here, \(\mathbf{\xi'}\) is the accumulated displacement prior to the
current time step and \(\mathbf{\xi}\) is the corrected
displacement. Corrections to the displacement due to the second mode
of motion described above (rotations about \(\mathbf{n}\)) are not
currently implemented, but are expected to be minor for most
simulations.

\sphinxAtStartPar
Furthermore, when the tangential force exceeds the critical force, the
tangential displacement is re\sphinxhyphen{}scaled to match the value for the
critical force (see {\hyperref[\detokenize{pair_granular:luding2008}]{\sphinxcrossref{\DUrole{std,std-ref}{Luding}}}}, equation 20 and related
discussion):
\begin{equation*}
\begin{split}\mathbf{\xi} = -\frac{1}{k_t}\left(\mu_t F_{n0}\mathbf{t} - \mathbf{F}_{t,damp}\right)\end{split}
\end{equation*}
\sphinxAtStartPar
The tangential force is added to the total normal force (elastic plus
damping) to produce the total force on the particle. The tangential
force also acts at the contact point (defined as the center of the
overlap region) to induce a torque on each particle according to:
\begin{equation*}
\begin{split}\mathbf{\tau}_i = -(R_i - 0.5 \delta) \mathbf{n} \times \mathbf{F}_t\end{split}
\end{equation*}\begin{equation*}
\begin{split}\mathbf{\tau}_j = -(R_j - 0.5 \delta) \mathbf{n} \times \mathbf{F}_t\end{split}
\end{equation*}
\sphinxAtStartPar
For \sphinxstyleemphasis{tangential mindlin}, the {\hyperref[\detokenize{pair_granular:mindlin1949}]{\sphinxcrossref{\DUrole{std,std-ref}{Mindlin}}}} no\sphinxhyphen{}slip solution
is used which differs from the \sphinxstyleemphasis{linear\_history} option by an additional factor
of \(a\), the radius of the contact region. The tangential force is given by:
\begin{equation*}
\begin{split}\mathbf{F}_t =  -\min(\mu_t F_{n0}, \|-k_t a \mathbf{\xi} + \mathbf{F}_\mathrm{t,damp}\|) \mathbf{t}\end{split}
\end{equation*}
\sphinxAtStartPar
Here, \(a\) is the radius of the contact region, given by \(a =\sqrt{R\delta}\)
for all normal contact models, except for \sphinxstyleemphasis{jkr}, where it is given
implicitly by \(\delta = a^2/R - 2\sqrt{\pi \gamma a/E}\), see
discussion above. To match the Mindlin solution, one should set
\(k_t = 8G_{eff}\), where \(G_{eff}\) is the effective shear modulus given by:
\begin{equation*}
\begin{split}G_{eff} = \left(\frac{2-\nu_i}{G_i} + \frac{2-\nu_j}{G_j}\right)^{-1}\end{split}
\end{equation*}
\sphinxAtStartPar
where \(G_i\) is the shear modulus of a particle of type \(i\), related to Young’s
modulus \(E_i\) and Poisson’s ratio \(\nu_i\) by \(G_i = E_i/(2(1+\nu_i))\).
This can also be achieved by specifying \sphinxstyleemphasis{NULL} for \(k_t\), in which case a
normal contact model that specifies material parameters \(E_i\) and
\(\nu_i\) is required (e.g. \sphinxstyleemphasis{hertz/material}, \sphinxstyleemphasis{dmt} or \sphinxstyleemphasis{jkr}). In this
case, mixing of the shear modulus for different particle types \sphinxstyleemphasis{i} and
\sphinxstyleemphasis{j} is done according to the formula above.

\begin{sphinxadmonition}{note}{Note:}
\sphinxAtStartPar
The radius of the contact region \(a\) depends on the normal overlap.
As a result, the tangential force for \sphinxstyleemphasis{mindlin} can change due to
a variation in normal overlap, even with no change in tangential displacement.
\end{sphinxadmonition}

\sphinxAtStartPar
For \sphinxstyleemphasis{tangential mindlin/force}, the accumulated elastic tangential force
characterizes the contact history, instead of the accumulated tangential
displacement. This prevents the dependence of the tangential force on the
normal overlap as noted above. The tangential force is given by:
\begin{equation*}
\begin{split}\mathbf{F}_t =  -\min(\mu_t F_{n0}, \|\mathbf{F}_{te} + \mathbf{F}_\mathrm{t,damp}\|) \mathbf{t}\end{split}
\end{equation*}
\sphinxAtStartPar
The increment of the elastic component of the tangential force
\(\mathbf{F}_{te}\) is given by:
\begin{equation*}
\begin{split}\mathrm{d}\mathbf{F}_{te} = -k_t a \mathbf{v}_{t,rel} \mathrm{d}\tau\end{split}
\end{equation*}
\sphinxAtStartPar
The changes in frame of reference of the contacting pair of particles during
contact are accounted for by the same formula as above, replacing the
accumulated tangential displacement \(\xi\), by the accumulated tangential
elastic force \(F_{te}\). When the tangential force exceeds the critical
force, the tangential force is directly re\sphinxhyphen{}scaled to match the value for
the critical force:
\begin{equation*}
\begin{split}\mathbf{F}_{te} = - \mu_t F_{n0}\mathbf{t} + \mathbf{F}_{t,damp}\end{split}
\end{equation*}
\sphinxAtStartPar
The same rules as those described for \sphinxstyleemphasis{mindlin} apply regarding the tangential
stiffness and mixing of the shear modulus for different particle types.

\sphinxAtStartPar
The \sphinxstyleemphasis{mindlin\_rescale} option uses the same form as \sphinxstyleemphasis{mindlin}, but the
magnitude of the tangential displacement is re\sphinxhyphen{}scaled as the contact
unloads, i.e. if \(a < a_{t_{n-1}}\):
\begin{equation*}
\begin{split}\mathbf{\xi} = \mathbf{\xi_{t_{n-1}}} \frac{a}{a_{t_{n-1}}}\end{split}
\end{equation*}
\sphinxAtStartPar
Here, \(t_{n-1}\) indicates the value at the previous time
step. This rescaling accounts for the fact that a decrease in the
contact area upon unloading leads to the contact being unable to
support the previous tangential loading, and spurious energy is
created without the rescaling above ({\hyperref[\detokenize{pair_granular:waltonpc}]{\sphinxcrossref{\DUrole{std,std-ref}{Walton}}}} ).

\begin{sphinxadmonition}{note}{Note:}
\sphinxAtStartPar
For \sphinxstyleemphasis{mindlin}, a decrease in the tangential force already occurs as the
contact unloads, due to the dependence of the tangential force on the normal
force described above. By re\sphinxhyphen{}scaling \(\xi\), \sphinxstyleemphasis{mindlin\_rescale}
effectively re\sphinxhyphen{}scales the tangential force twice, i.e., proportionally to
\(a^2\). This peculiar behavior results from use of the accumulated
tangential displacement to characterize the contact history. Although
\sphinxstyleemphasis{mindlin\_rescale} remains available for historic reasons and backward
compatibility purposes, it should be avoided in favor of \sphinxstyleemphasis{mindlin\_rescale/force}.
\end{sphinxadmonition}

\sphinxAtStartPar
The \sphinxstyleemphasis{mindlin\_rescale/force} option uses the same form as \sphinxstyleemphasis{mindlin/force},
but the magnitude of the tangential elastic force is re\sphinxhyphen{}scaled as the contact
unloads, i.e. if \(a < a_{t_{n-1}}\):
\begin{equation*}
\begin{split}\mathbf{F}_{te} = \mathbf{F}_{te, t_{n-1}} \frac{a}{a_{t_{n-1}}}\end{split}
\end{equation*}
\sphinxAtStartPar
This approach provides a better approximation of the {\hyperref[\detokenize{pair_granular:mindlin1953}]{\sphinxcrossref{\DUrole{std,std-ref}{Mindlin\sphinxhyphen{}Deresiewicz}}}}
laws and is more consistent than \sphinxstyleemphasis{mindlin\_rescale}. See discussions in
{\hyperref[\detokenize{pair_granular:thornton2013}]{\sphinxcrossref{\DUrole{std,std-ref}{Thornton et al, 2013}}}}, particularly equation 18(b) of that
work and associated discussion, and {\hyperref[\detokenize{pair_granular:agnolinroux2007}]{\sphinxcrossref{\DUrole{std,std-ref}{Agnolin and Roux, 2007}}}},
particularly Appendix A.


\bigskip\hrule\bigskip


\sphinxAtStartPar
The optional \sphinxstyleemphasis{rolling} keyword enables rolling friction, which resists
pure rolling motion of particles. The options currently supported are:
\begin{enumerate}
\sphinxsetlistlabels{\arabic}{enumi}{enumii}{}{.}%
\item {} 
\sphinxAtStartPar
\sphinxstyleemphasis{none}

\item {} 
\sphinxAtStartPar
\sphinxstyleemphasis{sds} : \(k_{roll}\), \(\gamma_{roll}\), \(\mu_{roll}\)

\end{enumerate}

\sphinxAtStartPar
If the \sphinxstyleemphasis{rolling} keyword is not specified, the model defaults to \sphinxstyleemphasis{none}.

\sphinxAtStartPar
For \sphinxstyleemphasis{rolling sds}, rolling friction is computed via a
spring\sphinxhyphen{}dashpot\sphinxhyphen{}slider, using a ‘pseudo\sphinxhyphen{}force’ formulation, as detailed
by {\hyperref[\detokenize{pair_granular:luding2008}]{\sphinxcrossref{\DUrole{std,std-ref}{Luding}}}}. Unlike the formulation in
{\hyperref[\detokenize{pair_granular:marshall2009}]{\sphinxcrossref{\DUrole{std,std-ref}{Marshall}}}}, this allows for the required adjustment of
rolling displacement due to changes in the frame of reference of the
contacting pair.  The rolling pseudo\sphinxhyphen{}force is computed analogously to
the tangential force:
\begin{equation*}
\begin{split}\mathbf{F}_{roll,0} =  k_{roll} \mathbf{\xi}_{roll}  - \gamma_{roll} \mathbf{v}_{roll}\end{split}
\end{equation*}
\sphinxAtStartPar
Here, \(\mathbf{v}_{roll} = -R(\boldsymbol{\Omega}_i - \boldsymbol{\Omega}_j) \times \mathbf{n}\) is the relative rolling
velocity, as given in {\hyperref[\detokenize{pair_granular:wang2015}]{\sphinxcrossref{\DUrole{std,std-ref}{Wang et al}}}} and
{\hyperref[\detokenize{pair_granular:luding2008}]{\sphinxcrossref{\DUrole{std,std-ref}{Luding}}}}. This differs from the expressions given by {\hyperref[\detokenize{pair_granular:kuhn2004}]{\sphinxcrossref{\DUrole{std,std-ref}{Kuhn and Bagi}}}} and used in {\hyperref[\detokenize{pair_granular:marshall2009}]{\sphinxcrossref{\DUrole{std,std-ref}{Marshall}}}}; see {\hyperref[\detokenize{pair_granular:wang2015}]{\sphinxcrossref{\DUrole{std,std-ref}{Wang et al}}}} for details. The rolling displacement is given by:
\begin{equation*}
\begin{split}\mathbf{\xi}_{roll} = \int_{t_0}^t \mathbf{v}_{roll} (\tau) \mathrm{d} \tau\end{split}
\end{equation*}
\sphinxAtStartPar
A Coulomb friction criterion truncates the rolling pseudo\sphinxhyphen{}force if it
exceeds a critical value:
\begin{equation*}
\begin{split}\mathbf{F}_{roll} =  \min(\mu_{roll} F_{n,0}, \|\mathbf{F}_{roll,0}\|)\mathbf{k}\end{split}
\end{equation*}
\sphinxAtStartPar
Here, \(\mathbf{k} = \mathbf{v}_{roll}/\|\mathbf{v}_{roll}\|\) is the direction of
the pseudo\sphinxhyphen{}force.  As with tangential displacement, the rolling
displacement is rescaled when the critical force is exceeded, so that
the spring length corresponds the critical force. Additionally, the
displacement is adjusted to account for rotations of the frame of
reference of the two contacting particles in a manner analogous to the
tangential displacement.

\sphinxAtStartPar
The rolling pseudo\sphinxhyphen{}force does not contribute to the total force on
either particle (hence ‘pseudo’), but acts only to induce an equal and
opposite torque on each particle, according to:
\begin{equation*}
\begin{split}\tau_{roll,i} =  R \mathbf{n} \times \mathbf{F}_{roll}\end{split}
\end{equation*}\begin{equation*}
\begin{split}\tau_{roll,j} =  -\tau_{roll,i}\end{split}
\end{equation*}

\bigskip\hrule\bigskip


\sphinxAtStartPar
The optional \sphinxstyleemphasis{twisting} keyword enables twisting friction, which
resists rotation of two contacting particles about the vector
\(\mathbf{n}\) that connects their centers. The options currently
supported are:
\begin{enumerate}
\sphinxsetlistlabels{\arabic}{enumi}{enumii}{}{.}%
\item {} 
\sphinxAtStartPar
\sphinxstyleemphasis{none}

\item {} 
\sphinxAtStartPar
\sphinxstyleemphasis{sds} : \(k_{twist}\), \(\gamma_{twist}\), \(\mu_{twist}\)

\item {} 
\sphinxAtStartPar
\sphinxstyleemphasis{marshall}

\end{enumerate}

\sphinxAtStartPar
If the \sphinxstyleemphasis{twisting} keyword is not specified, the model defaults to \sphinxstyleemphasis{none}.

\sphinxAtStartPar
For both \sphinxstyleemphasis{twisting sds} and \sphinxstyleemphasis{twisting marshall}, a history\sphinxhyphen{}dependent
spring\sphinxhyphen{}dashpot\sphinxhyphen{}slider is used to compute the twisting torque. Because
twisting displacement is a scalar, there is no need to adjust for
changes in the frame of reference due to rotations of the particle
pair. The formulation in {\hyperref[\detokenize{pair_granular:marshall2009}]{\sphinxcrossref{\DUrole{std,std-ref}{Marshall}}}} therefore provides
the most straightforward treatment:
\begin{equation*}
\begin{split}\tau_{twist,0} = -k_{twist}\xi_{twist} - \gamma_{twist}\Omega_{twist}\end{split}
\end{equation*}
\sphinxAtStartPar
Here \(\xi_{twist} = \int_{t_0}^t \Omega_{twist} (\tau) \mathrm{d}\tau\) is the twisting angular displacement, and
\(\Omega_{twist} = (\mathbf{\Omega}_i - \mathbf{\Omega}_j) \cdot \mathbf{n}\) is the relative twisting angular velocity. The torque
is then truncated according to:
\begin{equation*}
\begin{split}\tau_{twist} = \min(\mu_{twist} F_{n,0}, \tau_{twist,0})\end{split}
\end{equation*}
\sphinxAtStartPar
Similar to the sliding and rolling displacement, the angular
displacement is rescaled so that it corresponds to the critical value
if the twisting torque exceeds this critical value:
\begin{equation*}
\begin{split}\xi_{twist} = \frac{1}{k_{twist}} (\mu_{twist} F_{n,0}sgn(\Omega_{twist}) - \gamma_{twist}\Omega_{twist})\end{split}
\end{equation*}
\sphinxAtStartPar
For \sphinxstyleemphasis{twisting sds}, the coefficients \(k_{twist}, \gamma_{twist}\)
and \(\mu_{twist}\) are simply the user input parameters that follow
the \sphinxstyleemphasis{twisting sds} keywords in the \sphinxstyleemphasis{pair\_coeff} command.

\sphinxAtStartPar
For \sphinxstyleemphasis{twisting\_marshall}, the coefficients are expressed in terms of
sliding friction coefficients, as discussed in
{\hyperref[\detokenize{pair_granular:marshall2009}]{\sphinxcrossref{\DUrole{std,std-ref}{Marshall}}}} (see equations 32 and 33 of that work):
\begin{equation*}
\begin{split}k_{twist} = 0.5k_ta^2\end{split}
\end{equation*}\begin{equation*}
\begin{split}\eta_{twist} = 0.5\eta_ta^2\end{split}
\end{equation*}\begin{equation*}
\begin{split}\mu_{twist} = \frac{2}{3}a\mu_t\end{split}
\end{equation*}
\sphinxAtStartPar
Finally, the twisting torque on each particle is given by:
\begin{equation*}
\begin{split}\mathbf{\tau}_{twist,i} = \tau_{twist}\mathbf{n}\end{split}
\end{equation*}\begin{equation*}
\begin{split}\mathbf{\tau}_{twist,j} = -\mathbf{\tau}_{twist,i}\end{split}
\end{equation*}

\bigskip\hrule\bigskip


\sphinxAtStartPar
If two particles are moving away from each other while in contact, there
is a possibility that the particles could experience an effective attractive
force due to damping. If the optional \sphinxstyleemphasis{limit\_damping} keyword is used, this option
will zero out the normal component of the force if there is an effective
attractive force. This keyword cannot be used with the JKR or DMT models.


\bigskip\hrule\bigskip


\sphinxAtStartPar
The optional \sphinxstyleemphasis{heat} keyword enables heat conduction. The options currently
supported are:
\begin{enumerate}
\sphinxsetlistlabels{\arabic}{enumi}{enumii}{}{.}%
\item {} 
\sphinxAtStartPar
\sphinxstyleemphasis{none}

\item {} 
\sphinxAtStartPar
\sphinxstyleemphasis{radius} : \(k_{s}\)

\item {} 
\sphinxAtStartPar
\sphinxstyleemphasis{area} : \(h_{s}\)

\end{enumerate}

\sphinxAtStartPar
If the \sphinxstyleemphasis{heat} keyword is not specified, the model defaults to \sphinxstyleemphasis{none}.

\sphinxAtStartPar
For \sphinxstyleemphasis{heat} \sphinxstyleemphasis{radius}, the heat
\(Q\) conducted between two particles is given by
\begin{equation*}
\begin{split}Q = 2 k_{s} a \Delta T\end{split}
\end{equation*}
\sphinxAtStartPar
where \(\Delta T\) is the difference in the two particles’ temperature,
\(k_{s}\) is a non\sphinxhyphen{}negative numeric value for the conductivity (in units
of power/(length*temperature)), and \(a\) is the radius of the contact and
depends on the normal force model. This is the model proposed by
{\hyperref[\detokenize{pair_granular:vargasmccarthy2001}]{\sphinxcrossref{\DUrole{std,std-ref}{Vargas and McCarthy}}}}.

\sphinxAtStartPar
For \sphinxstyleemphasis{heat} \sphinxstyleemphasis{area}, the heat
\(Q\) conducted between two particles is given by
\begin{equation*}
\begin{split}Q = h_{s} A \Delta T\end{split}
\end{equation*}
\sphinxAtStartPar
where \(\Delta T\) is the difference in the two particles’ temperature,
\(h_{s}\) is a non\sphinxhyphen{}negative numeric value for the heat transfer
coefficient (in units of power/(area*temperature)), and \(A=\pi a^2\) is
the area of the contact and depends on the normal force model.

\sphinxAtStartPar
Note that the option \sphinxstyleemphasis{none} must either be used in all or none of of the
\sphinxstyleemphasis{pair\_coeff} calls. See {\hyperref[\detokenize{fix_heat_flow::doc}]{\sphinxcrossref{\DUrole{doc}{fix heat/flow}}}} and
{\hyperref[\detokenize{fix_property_atom::doc}]{\sphinxcrossref{\DUrole{doc}{fix property/atom}}}} for more information on this
option.


\bigskip\hrule\bigskip


\sphinxAtStartPar
The \sphinxstyleemphasis{granular} pair style can reproduce the behavior of the
\sphinxstyleemphasis{pair gran/*} styles with the appropriate settings (some very
minor differences can be expected due to corrections in
displacement history frame\sphinxhyphen{}of\sphinxhyphen{}reference, and the application
of the torque at the center of the contact rather than
at each particle). The first example above
is equivalent to \sphinxstyleemphasis{pair gran/hooke 1000.0 NULL 50.0 50.0 0.4 1}.
The second example is equivalent to
\sphinxstyleemphasis{pair gran/hooke/history 1000.0 500.0 50.0 50.0 0.4 1}.
The third example is equivalent to
\sphinxstyleemphasis{pair gran/hertz/history 1000.0 500.0 50.0 50.0 0.4 1 limit\_damping}.


\bigskip\hrule\bigskip


\sphinxAtStartPar
LAMMPS automatically sets pairwise cutoff values for \sphinxstyleemphasis{pair\_style
granular} based on particle radii (and in the case of \sphinxstyleemphasis{jkr} pull\sphinxhyphen{}off
distances). In the vast majority of situations, this is adequate.
However, a cutoff value can optionally be appended to the \sphinxstyleemphasis{pair\_style
granular} command to specify a global cutoff (i.e. a cutoff for all
atom types). Additionally, the optional \sphinxstyleemphasis{cutoff} keyword can be passed
to the \sphinxstyleemphasis{pair\_coeff} command, followed by a cutoff value.  This will
set a pairwise cutoff for the atom types in the \sphinxstyleemphasis{pair\_coeff} command.
These options may be useful in some rare cases where the automatic
cutoff determination is not sufficient, e.g.  if particle diameters
are being modified via the \sphinxstyleemphasis{fix adapt} command. In that case, the
global cutoff specified as part of the \sphinxstyleemphasis{pair\_style granular} command
is applied to all atom types, unless it is overridden for a given atom
type combination by the \sphinxstyleemphasis{cutoff} value specified in the \sphinxstyleemphasis{pair coeff}
command.  If \sphinxstyleemphasis{cutoff} is only specified in the \sphinxstyleemphasis{pair coeff} command
and no global cutoff is appended to the \sphinxstyleemphasis{pair\_style granular} command,
then LAMMPS will use that cutoff for the specified atom type
combination, and automatically set pairwise cutoffs for the remaining
atom types.


\bigskip\hrule\bigskip


\sphinxAtStartPar
Styles with a \sphinxstyleemphasis{gpu}, \sphinxstyleemphasis{intel}, \sphinxstyleemphasis{kk}, \sphinxstyleemphasis{omp}, or \sphinxstyleemphasis{opt} suffix are
functionally the same as the corresponding style without the suffix.
They have been optimized to run faster, depending on your available
hardware, as discussed on the {\hyperref[\detokenize{Speed_packages::doc}]{\sphinxcrossref{\DUrole{doc}{Accelerator packages}}}}
page.  The accelerated styles take the same arguments and should
produce the same results, except for round\sphinxhyphen{}off and precision issues.

\sphinxAtStartPar
These accelerated styles are part of the GPU, INTEL, KOKKOS,
OPENMP, and OPT packages, respectively.  They are only enabled if
LAMMPS was built with those packages.  See the {\hyperref[\detokenize{Build_package::doc}]{\sphinxcrossref{\DUrole{doc}{Build package}}}} page for more info.

\sphinxAtStartPar
You can specify the accelerated styles explicitly in your input script
by including their suffix, or you can use the {\hyperref[\detokenize{Run_options::doc}]{\sphinxcrossref{\DUrole{doc}{\sphinxhyphen{}suffix command\sphinxhyphen{}line switch}}}} when you invoke LAMMPS, or you can use the
{\hyperref[\detokenize{suffix::doc}]{\sphinxcrossref{\DUrole{doc}{suffix}}}} command in your input script.

\sphinxAtStartPar
See the {\hyperref[\detokenize{Speed_packages::doc}]{\sphinxcrossref{\DUrole{doc}{Accelerator packages}}}} page for more
instructions on how to use the accelerated styles effectively.


\bigskip\hrule\bigskip



\subsection{Mixing, shift, table, tail correction, restart, rRESPA info}
\label{\detokenize{pair_granular:mixing-shift-table-tail-correction-restart-rrespa-info}}
\sphinxAtStartPar
The {\hyperref[\detokenize{pair_modify::doc}]{\sphinxcrossref{\DUrole{doc}{pair\_modify}}}} mix, shift, table, and tail options
are not relevant for granular pair styles.

\sphinxAtStartPar
Mixing of coefficients is carried out using geometric averaging for
most quantities, e.g. if friction coefficient for type 1\sphinxhyphen{}type 1
interactions is set to \(\mu_1\), and friction coefficient for type
2\sphinxhyphen{}type 2 interactions is set to \(\mu_2\), the friction coefficient
for type1\sphinxhyphen{}type2 interactions is computed as \(\sqrt{\mu_1\mu_2}\)
(unless explicitly specified to a different value by a \sphinxstyleemphasis{pair\_coeff 1 2
…} command). The exception to this is elastic modulus, only
applicable to \sphinxstyleemphasis{hertz/material}, \sphinxstyleemphasis{dmt} and \sphinxstyleemphasis{jkr} normal contact
models. In that case, the effective elastic modulus is computed as:
\begin{equation*}
\begin{split}E_{eff,ij} = \left(\frac{1-\nu_i^2}{E_i} + \frac{1-\nu_j^2}{E_j}\right)^{-1}\end{split}
\end{equation*}
\sphinxAtStartPar
If the \sphinxstyleemphasis{i\sphinxhyphen{}j} coefficients \(E_{ij}\) and \(\nu_{ij}\) are
explicitly specified, the effective modulus is computed as:
\begin{equation*}
\begin{split}E_{eff,ij} = \left(\frac{1-\nu_{ij}^2}{E_{ij}} + \frac{1-\nu_{ij}^2}{E_{ij}}\right)^{-1}\end{split}
\end{equation*}
\sphinxAtStartPar
or
\begin{equation*}
\begin{split}E_{eff,ij} = \frac{E_{ij}}{2(1-\nu_{ij}^2)}\end{split}
\end{equation*}
\sphinxAtStartPar
These pair styles write their information to {\hyperref[\detokenize{restart::doc}]{\sphinxcrossref{\DUrole{doc}{binary restart files}}}}, so a pair\_style command does not need to be
specified in an input script that reads a restart file.

\sphinxAtStartPar
These pair styles can only be used via the \sphinxstyleemphasis{pair} keyword of the
{\hyperref[\detokenize{run_style::doc}]{\sphinxcrossref{\DUrole{doc}{run\_style respa}}}} command.  They do not support the
\sphinxstyleemphasis{inner}, \sphinxstyleemphasis{middle}, \sphinxstyleemphasis{outer} keywords.

\sphinxAtStartPar
The single() function of these pair styles returns 0.0 for the energy
of a pairwise interaction, since energy is not conserved in these
dissipative potentials.  It also returns only the normal component of
the pairwise interaction force.  However, the single() function also
calculates 13 extra pairwise quantities.  The first 3 are the
components of the tangential force between particles I and J, acting
on particle I.  The fourth is the magnitude of this tangential force.
The next 3 (5\sphinxhyphen{}7) are the components of the rolling torque acting on
particle I. The next entry (8) is the magnitude of the rolling torque.
The next entry (9) is the magnitude of the twisting torque acting
about the vector connecting the two particle centers.
The next 3 (10\sphinxhyphen{}12) are the components of the vector connecting
the centers of the two particles (x\_I \sphinxhyphen{} x\_J). The last quantity (13)
is the heat flow between the two particles, set to 0 if no heat model
is active.

\sphinxAtStartPar
These extra quantities can be accessed by the {\hyperref[\detokenize{compute_pair_local::doc}]{\sphinxcrossref{\DUrole{doc}{compute pair/local}}}} command, as \sphinxstyleemphasis{p1}, \sphinxstyleemphasis{p2}, …,
\sphinxstyleemphasis{p12}.


\bigskip\hrule\bigskip



\subsection{Restrictions}
\label{\detokenize{pair_granular:restrictions}}
\sphinxAtStartPar
This pair style is part of the GRANULAR package.  It is
only enabled if LAMMPS was built with that package.
See the {\hyperref[\detokenize{Build_package::doc}]{\sphinxcrossref{\DUrole{doc}{Build package}}}} page for more info.

\sphinxAtStartPar
This pair style requires that atoms store per\sphinxhyphen{}particle radius,
torque, and angular velocity (omega) as defined by the
{\hyperref[\detokenize{atom_style::doc}]{\sphinxcrossref{\DUrole{doc}{atom\_style sphere}}}}.

\sphinxAtStartPar
This pair style requires you to use the {\hyperref[\detokenize{comm_modify::doc}]{\sphinxcrossref{\DUrole{doc}{comm\_modify vel yes}}}}
command so that velocities are stored by ghost atoms.

\sphinxAtStartPar
This pair style will not restart exactly when using the
{\hyperref[\detokenize{read_restart::doc}]{\sphinxcrossref{\DUrole{doc}{read\_restart}}}} command, though it should provide
statistically similar results.  This is because the forces it
computes depend on atom velocities and the atom velocities have
been propagated half a timestep between the force computation and
when the restart is written, due to using Velocity Verlet time
integration. See the {\hyperref[\detokenize{read_restart::doc}]{\sphinxcrossref{\DUrole{doc}{read\_restart}}}} command
for more details.

\sphinxAtStartPar
Accumulated values for individual contacts are saved to restart
files but are not saved to data files. Therefore, forces may
differ significantly when a system is reloaded using the
{\hyperref[\detokenize{read_data::doc}]{\sphinxcrossref{\DUrole{doc}{read\_data}}}} command.


\subsection{Related commands}
\label{\detokenize{pair_granular:related-commands}}
\sphinxAtStartPar
{\hyperref[\detokenize{pair_coeff::doc}]{\sphinxcrossref{\DUrole{doc}{pair\_coeff}}}}
{\hyperref[\detokenize{pair_gran::doc}]{\sphinxcrossref{\DUrole{doc}{pair gran/*}}}}


\subsection{Default}
\label{\detokenize{pair_granular:default}}
\sphinxAtStartPar
For the \sphinxstyleemphasis{pair\_coeff} settings: \sphinxstyleemphasis{damping viscoelastic}, \sphinxstyleemphasis{rolling none},
\sphinxstyleemphasis{twisting none}.


\subsection{References}
\label{\detokenize{pair_granular:references}}\phantomsection\label{\detokenize{pair_granular:brill1996}}
\sphinxAtStartPar
\sphinxstylestrong{(Brilliantov et al, 1996)} Brilliantov, N. V., Spahn, F., Hertzsch,
J. M., \& Poschel, T. (1996).  Model for collisions in granular
gases. Physical review E, 53(5), 5382.

\phantomsection\label{\detokenize{pair_granular:tsuji1992}}
\sphinxAtStartPar
\sphinxstylestrong{(Tsuji et al, 1992)} Tsuji, Y., Tanaka, T., \& Ishida,
T. (1992). Lagrangian numerical simulation of plug flow of
cohesionless particles in a horizontal pipe. Powder technology, 71(3),
239\sphinxhyphen{}250.

\phantomsection\label{\detokenize{pair_granular:jkr1971}}
\sphinxAtStartPar
\sphinxstylestrong{(Johnson et al, 1971)} Johnson, K. L., Kendall, K., \& Roberts,
A. D. (1971).  Surface energy and the contact of elastic
solids. Proc. R. Soc. Lond. A, 324(1558), 301\sphinxhyphen{}313.

\phantomsection\label{\detokenize{pair_granular:dmt1975}}
\sphinxAtStartPar
\sphinxstylestrong{Derjaguin et al, 1975)} Derjaguin, B. V., Muller, V. M., \& Toporov,
Y. P. (1975). Effect of contact deformations on the adhesion of
particles. Journal of Colloid and interface science, 53(2), 314\sphinxhyphen{}326.

\phantomsection\label{\detokenize{pair_granular:luding2008}}
\sphinxAtStartPar
\sphinxstylestrong{(Luding, 2008)} Luding, S. (2008). Cohesive, frictional powders:
contact models for tension. Granular matter, 10(4), 235.

\phantomsection\label{\detokenize{pair_granular:marshall2009}}
\sphinxAtStartPar
\sphinxstylestrong{(Marshall, 2009)} Marshall, J. S. (2009). Discrete\sphinxhyphen{}element modeling
of particulate aerosol flows.  Journal of Computational Physics,
228(5), 1541\sphinxhyphen{}1561.

\phantomsection\label{\detokenize{pair_granular:silbert2001}}
\sphinxAtStartPar
\sphinxstylestrong{(Silbert, 2001)} Silbert, L. E., Ertas, D., Grest, G. S., Halsey,
T. C., Levine, D., \& Plimpton, S. J. (2001).  Granular flow down an
inclined plane: Bagnold scaling and rheology. Physical Review E,
64(5), 051302.

\phantomsection\label{\detokenize{pair_granular:kuhn2004}}
\sphinxAtStartPar
\sphinxstylestrong{(Kuhn and Bagi, 2005)} Kuhn, M. R., \& Bagi, K. (2004). Contact
rolling and deformation in granular media.  International journal of
solids and structures, 41(21), 5793\sphinxhyphen{}5820.

\phantomsection\label{\detokenize{pair_granular:wang2015}}
\sphinxAtStartPar
\sphinxstylestrong{(Wang et al, 2015)} Wang, Y., Alonso\sphinxhyphen{}Marroquin, F., \& Guo,
W. W. (2015).  Rolling and sliding in 3\sphinxhyphen{}D discrete element
models. Particuology, 23, 49\sphinxhyphen{}55.

\phantomsection\label{\detokenize{pair_granular:thornton1991}}
\sphinxAtStartPar
\sphinxstylestrong{(Thornton, 1991)} Thornton, C. (1991). Interparticle sliding in the
presence of adhesion.  J. Phys. D: Appl. Phys. 24 1942

\phantomsection\label{\detokenize{pair_granular:mindlin1949}}
\sphinxAtStartPar
\sphinxstylestrong{(Mindlin, 1949)} Mindlin, R. D. (1949). Compliance of elastic bodies
in contact.  J. Appl. Mech., ASME 16, 259\sphinxhyphen{}268.

\phantomsection\label{\detokenize{pair_granular:thornton2013}}
\sphinxAtStartPar
\sphinxstylestrong{(Thornton et al, 2013)} Thornton, C., Cummins, S. J., \& Cleary,
P. W. (2013).  An investigation of the comparative behavior of
alternative contact force models during inelastic collisions. Powder
Technology, 233, 30\sphinxhyphen{}46.

\phantomsection\label{\detokenize{pair_granular:waltonpc}}
\sphinxAtStartPar
\sphinxstylestrong{(Otis R. Walton)} Walton, O.R., Personal Communication

\phantomsection\label{\detokenize{pair_granular:mindlin1953}}
\sphinxAtStartPar
\sphinxstylestrong{(Mindlin and Deresiewicz, 1953)} Mindlin, R.D., \& Deresiewicz, H (1953).
Elastic Spheres in Contact under Varying Oblique Force.
J. Appl. Mech., ASME 20, 327\sphinxhyphen{}344.

\phantomsection\label{\detokenize{pair_granular:agnolinroux2007}}
\sphinxAtStartPar
\sphinxstylestrong{(Agnolin and Roux 2007)} Agnolin, I. \& Roux, J\sphinxhyphen{}N. (2007).
Internal states of model isotropic granular packings.
I. Assembling process, geometry, and contact networks. Phys. Rev. E, 76, 061302.

\phantomsection\label{\detokenize{pair_granular:vargasmccarthy2001}}
\sphinxAtStartPar
\sphinxstylestrong{(Vargas and McCarthy 2001)} Vargas, W.L. and McCarthy, J.J. (2001).
Heat conduction in granular materials.
AIChE Journal, 47(5), 1052\sphinxhyphen{}1059.

\sphinxstepscope

\index{pair\_style lj/gromacs@\spxentry{pair\_style lj/gromacs}}\ignorespaces 
\index{pair\_style lj/gromacs/gpu@\spxentry{pair\_style lj/gromacs/gpu}}\ignorespaces 
\index{pair\_style lj/gromacs/kk@\spxentry{pair\_style lj/gromacs/kk}}\ignorespaces 
\index{pair\_style lj/gromacs/omp@\spxentry{pair\_style lj/gromacs/omp}}\ignorespaces 
\index{pair\_style lj/gromacs/coul/gromacs@\spxentry{pair\_style lj/gromacs/coul/gromacs}}\ignorespaces 
\index{pair\_style lj/gromacs/coul/gromacs/kk@\spxentry{pair\_style lj/gromacs/coul/gromacs/kk}}\ignorespaces 
\index{pair\_style lj/gromacs/coul/gromacs/omp@\spxentry{pair\_style lj/gromacs/coul/gromacs/omp}}\ignorespaces 

\section{pair\_style lj/gromacs command}
\label{\detokenize{pair_gromacs:pair-style-lj-gromacs-command}}\label{\detokenize{pair_gromacs:index-6}}\label{\detokenize{pair_gromacs:index-5}}\label{\detokenize{pair_gromacs:index-4}}\label{\detokenize{pair_gromacs:index-3}}\label{\detokenize{pair_gromacs:index-2}}\label{\detokenize{pair_gromacs:index-1}}\label{\detokenize{pair_gromacs:index-0}}\label{\detokenize{pair_gromacs::doc}}
\sphinxAtStartPar
Accelerator Variants: \sphinxstyleemphasis{lj/gromacs/gpu}, \sphinxstyleemphasis{lj/gromacs/kk}, \sphinxstyleemphasis{lj/gromacs/omp}


\section{pair\_style lj/gromacs/coul/gromacs command}
\label{\detokenize{pair_gromacs:pair-style-lj-gromacs-coul-gromacs-command}}
\sphinxAtStartPar
Accelerator Variants: \sphinxstyleemphasis{lj/gromacs/coul/gromacs/kk}, \sphinxstyleemphasis{lj/gromacs/coul/gromacs/omp}


\subsection{Syntax}
\label{\detokenize{pair_gromacs:syntax}}
\begin{sphinxVerbatim}[commandchars=\\\{\}]
\PYG{k}{pair\PYGZus{}style}\PYG{+w}{ }\PYG{n}{style}\PYG{+w}{ }\PYG{n}{args}
\end{sphinxVerbatim}
\begin{itemize}
\item {} 
\sphinxAtStartPar
style = \sphinxstyleemphasis{lj/gromacs} or \sphinxstyleemphasis{lj/gromacs/coul/gromacs}

\item {} 
\sphinxAtStartPar
args = list of arguments for a particular style

\end{itemize}
\begin{sphinxalltt}
\sphinxstyleemphasis{lj/gromacs} args = inner outer
  inner, outer = global switching cutoffs for Lennard Jones
\sphinxstyleemphasis{lj/gromacs/coul/gromacs} args = inner outer (inner2) (outer2)
  inner, outer = global switching cutoffs for Lennard Jones (and Coulombic if only 2 args)
  inner2, outer2 = global switching cutoffs for Coulombic (optional)
\end{sphinxalltt}


\subsection{Examples}
\label{\detokenize{pair_gromacs:examples}}
\begin{sphinxVerbatim}[commandchars=\\\{\}]
\PYG{k}{pair\PYGZus{}style}\PYG{+w}{ }\PYG{n}{lj}\PYG{o}{/}\PYG{n}{gromacs}\PYG{+w}{ }\PYG{l+m}{9.0}\PYG{+w}{ }\PYG{l+m}{12.0}
\PYG{k}{pair\PYGZus{}coeff}\PYG{+w}{ }\PYG{o}{*}\PYG{+w}{ }\PYG{o}{*}\PYG{+w}{ }\PYG{l+m}{100.0}\PYG{+w}{ }\PYG{l+m}{2.0}
\PYG{k}{pair\PYGZus{}coeff}\PYG{+w}{ }\PYG{l+m}{2}\PYG{+w}{ }\PYG{l+m}{2}\PYG{+w}{ }\PYG{l+m}{100.0}\PYG{+w}{ }\PYG{l+m}{2.0}\PYG{+w}{ }\PYG{l+m}{8.0}\PYG{+w}{ }\PYG{l+m}{10.0}

\PYG{k}{pair\PYGZus{}style}\PYG{+w}{ }\PYG{n}{lj}\PYG{o}{/}\PYG{n}{gromacs}\PYG{o}{/}\PYG{n}{coul}\PYG{o}{/}\PYG{n}{gromacs}\PYG{+w}{ }\PYG{l+m}{9.0}\PYG{+w}{ }\PYG{l+m}{12.0}
\PYG{k}{pair\PYGZus{}style}\PYG{+w}{ }\PYG{n}{lj}\PYG{o}{/}\PYG{n}{gromacs}\PYG{o}{/}\PYG{n}{coul}\PYG{o}{/}\PYG{n}{gromacs}\PYG{+w}{ }\PYG{l+m}{8.0}\PYG{+w}{ }\PYG{l+m}{10.0}\PYG{+w}{ }\PYG{l+m}{7.0}\PYG{+w}{ }\PYG{l+m}{9.0}
\PYG{k}{pair\PYGZus{}coeff}\PYG{+w}{ }\PYG{o}{*}\PYG{+w}{ }\PYG{o}{*}\PYG{+w}{ }\PYG{l+m}{100.0}\PYG{+w}{ }\PYG{l+m}{2.0}
\end{sphinxVerbatim}


\subsection{Description}
\label{\detokenize{pair_gromacs:description}}
\sphinxAtStartPar
The \sphinxstyleemphasis{lj/gromacs} styles compute shifted LJ and Coulombic interactions
with an additional switching function S(r) that ramps the energy and force
smoothly to zero between an inner and outer cutoff.  It is a commonly
used potential in the \sphinxhref{https://www.gromacs.org}{GROMACS} MD code and for
the coarse\sphinxhyphen{}grained models of {\hyperref[\detokenize{pair_gromacs:marrink}]{\sphinxcrossref{\DUrole{std,std-ref}{(Marrink)}}}}.
\begin{equation*}
\begin{split}E_{LJ} = & 4 \epsilon \left[ \left(\frac{\sigma}{r}\right)^{12} -
         \left(\frac{\sigma}{r}\right)^6 \right] + S_{LJ}(r)
                    \qquad r < r_c \\
E_C  = & \frac{C q_i q_j}{\epsilon  r} + S_C(r) \qquad r < r_c \\
S(r) = & C \qquad r < r_1 \\
S(r) = & \frac{A}{3} (r - r_1)^3 + \frac{B}{4} (r - r_1)^4 + C \qquad  r_1 < r < r_c \\
A = & (-3 E'(r_c) + (r_c - r_1) E''(r_c))/(r_c - r_1)^2 \\
B = & (2 E'(r_c) - (r_c - r_1) E''(r_c))/(r_c - r_1)^3 \\
C = & -E(r_c) + \frac{1}{2} (r_c - r_1) E'(r_c) - \frac{1}{12} (r_c - r_1)^2 E''(r_c)\end{split}
\end{equation*}
\sphinxAtStartPar
\(r_1\) is the inner cutoff; \(r_c\) is the outer cutoff.  The
coefficients A, B, and C are computed by LAMMPS to perform the shifting
and smoothing.  The function S(r) is actually applied once to each term
of the LJ formula and once to the Coulombic formula, so there are 2 or 3
sets of A,B,C coefficients depending on which pair\_style is used.  The
boundary conditions applied to the smoothing function are as follows:
\(S'(r_1) = S''(r_1) = 0, S(r_c) = -E(r_c), S'(r_c) = -E'(r_c)\),
and \(S''(r_c) = -E''(r_c)\), where E(r) is the corresponding term
in the LJ or Coulombic potential energy function.  Single and double
primes denote first and second derivatives with respect to r,
respectively.

\sphinxAtStartPar
The inner and outer cutoff for the LJ and Coulombic terms can be the
same or different depending on whether 2 or 4 arguments are used in
the pair\_style command.  The inner LJ cutoff must be \textgreater{} 0, but the
inner Coulombic cutoff can be \textgreater{}= 0.

\sphinxAtStartPar
The following coefficients must be defined for each pair of atoms
types via the {\hyperref[\detokenize{pair_coeff::doc}]{\sphinxcrossref{\DUrole{doc}{pair\_coeff}}}} command as in the examples
above, or in the data file or restart files read by the
{\hyperref[\detokenize{read_data::doc}]{\sphinxcrossref{\DUrole{doc}{read\_data}}}} or {\hyperref[\detokenize{read_restart::doc}]{\sphinxcrossref{\DUrole{doc}{read\_restart}}}}
commands, or by mixing as described below:
\begin{itemize}
\item {} 
\sphinxAtStartPar
\(\epsilon\) (energy units)

\item {} 
\sphinxAtStartPar
\(\sigma\) (distance units)

\item {} 
\sphinxAtStartPar
inner (distance units)

\item {} 
\sphinxAtStartPar
outer (distance units)

\end{itemize}

\sphinxAtStartPar
Note that sigma is defined in the LJ formula as the zero\sphinxhyphen{}crossing
distance for the potential, not as the energy minimum at \(2^{1/6} \sigma\).

\sphinxAtStartPar
The last 2 coefficients are optional inner and outer cutoffs for style
\sphinxstyleemphasis{lj/gromacs}.  If not specified, the global \sphinxstyleemphasis{inner} and \sphinxstyleemphasis{outer} values
are used.

\sphinxAtStartPar
The last 2 coefficients cannot be used with style
\sphinxstyleemphasis{lj/gromacs/coul/gromacs} because this force field does not allow
varying cutoffs for individual atom pairs; all pairs use the global
cutoff(s) specified in the pair\_style command.


\bigskip\hrule\bigskip


\sphinxAtStartPar
Styles with a \sphinxstyleemphasis{gpu}, \sphinxstyleemphasis{intel}, \sphinxstyleemphasis{kk}, \sphinxstyleemphasis{omp}, or \sphinxstyleemphasis{opt} suffix are
functionally the same as the corresponding style without the suffix.
They have been optimized to run faster, depending on your available
hardware, as discussed on the {\hyperref[\detokenize{Speed_packages::doc}]{\sphinxcrossref{\DUrole{doc}{Accelerator packages}}}}
page.  The accelerated styles take the same arguments and should
produce the same results, except for round\sphinxhyphen{}off and precision issues.

\sphinxAtStartPar
These accelerated styles are part of the GPU, INTEL, KOKKOS,
OPENMP, and OPT packages, respectively.  They are only enabled if
LAMMPS was built with those packages.  See the {\hyperref[\detokenize{Build_package::doc}]{\sphinxcrossref{\DUrole{doc}{Build package}}}} page for more info.

\sphinxAtStartPar
You can specify the accelerated styles explicitly in your input script
by including their suffix, or you can use the {\hyperref[\detokenize{Run_options::doc}]{\sphinxcrossref{\DUrole{doc}{\sphinxhyphen{}suffix command\sphinxhyphen{}line switch}}}} when you invoke LAMMPS, or you can use the
{\hyperref[\detokenize{suffix::doc}]{\sphinxcrossref{\DUrole{doc}{suffix}}}} command in your input script.

\sphinxAtStartPar
See the {\hyperref[\detokenize{Speed_packages::doc}]{\sphinxcrossref{\DUrole{doc}{Accelerator packages}}}} page for more
instructions on how to use the accelerated styles effectively.


\bigskip\hrule\bigskip



\subsection{Mixing, shift, table, tail correction, restart, rRESPA info}
\label{\detokenize{pair_gromacs:mixing-shift-table-tail-correction-restart-rrespa-info}}
\sphinxAtStartPar
For atom type pairs I,J and I != J, the epsilon and sigma coefficients
and cutoff distance for all of the lj/cut pair styles can be mixed.
The default mix value is \sphinxstyleemphasis{geometric}.  See the “pair\_modify” command
for details.

\sphinxAtStartPar
None of the GROMACS pair styles support the
{\hyperref[\detokenize{pair_modify::doc}]{\sphinxcrossref{\DUrole{doc}{pair\_modify}}}} shift option, since the Lennard\sphinxhyphen{}Jones
portion of the pair interaction is already smoothed to 0.0 at the
cutoff.

\sphinxAtStartPar
The {\hyperref[\detokenize{pair_modify::doc}]{\sphinxcrossref{\DUrole{doc}{pair\_modify}}}} table option is not relevant
for this pair style.

\sphinxAtStartPar
None of the GROMACS pair styles support the
{\hyperref[\detokenize{pair_modify::doc}]{\sphinxcrossref{\DUrole{doc}{pair\_modify}}}} tail option for adding long\sphinxhyphen{}range tail
corrections to energy and pressure, since there are no corrections for
a potential that goes to 0.0 at the cutoff.

\sphinxAtStartPar
All of the GROMACS pair styles write their information to {\hyperref[\detokenize{restart::doc}]{\sphinxcrossref{\DUrole{doc}{binary restart files}}}}, so pair\_style and pair\_coeff commands do
not need to be specified in an input script that reads a restart file.

\sphinxAtStartPar
All of the GROMACS pair styles can only be used via the \sphinxstyleemphasis{pair}
keyword of the {\hyperref[\detokenize{run_style::doc}]{\sphinxcrossref{\DUrole{doc}{run\_style respa}}}} command.  They do not
support the \sphinxstyleemphasis{inner}, \sphinxstyleemphasis{middle}, \sphinxstyleemphasis{outer} keywords.


\bigskip\hrule\bigskip



\subsection{Restrictions}
\label{\detokenize{pair_gromacs:restrictions}}
\sphinxAtStartPar
This pair style is part of the EXTRA\sphinxhyphen{}PAIR package.  It is only enabled if
LAMMPS was built with that package.  See the
{\hyperref[\detokenize{Build_package::doc}]{\sphinxcrossref{\DUrole{doc}{Build package}}}} page for more info.


\subsection{Related commands}
\label{\detokenize{pair_gromacs:related-commands}}
\sphinxAtStartPar
{\hyperref[\detokenize{pair_coeff::doc}]{\sphinxcrossref{\DUrole{doc}{pair\_coeff}}}}


\subsection{Default}
\label{\detokenize{pair_gromacs:default}}
\sphinxAtStartPar
none


\bigskip\hrule\bigskip

\phantomsection\label{\detokenize{pair_gromacs:marrink}}
\sphinxAtStartPar
\sphinxstylestrong{(Marrink)} Marrink, de Vries, Mark, J Phys Chem B, 108, 750\sphinxhyphen{}760 (2004).

\sphinxstepscope

\index{pair\_style gw@\spxentry{pair\_style gw}}\ignorespaces 
\index{pair\_style gw/zbl@\spxentry{pair\_style gw/zbl}}\ignorespaces 

\section{pair\_style gw command}
\label{\detokenize{pair_gw:pair-style-gw-command}}\label{\detokenize{pair_gw:index-1}}\label{\detokenize{pair_gw:index-0}}\label{\detokenize{pair_gw::doc}}

\section{pair\_style gw/zbl command}
\label{\detokenize{pair_gw:pair-style-gw-zbl-command}}

\subsection{Syntax}
\label{\detokenize{pair_gw:syntax}}
\begin{sphinxVerbatim}[commandchars=\\\{\}]
\PYG{k}{pair\PYGZus{}style}\PYG{+w}{ }\PYG{n}{style}
\end{sphinxVerbatim}
\begin{itemize}
\item {} 
\sphinxAtStartPar
style = \sphinxstyleemphasis{gw} or \sphinxstyleemphasis{gw/zbl}

\end{itemize}


\subsection{Examples}
\label{\detokenize{pair_gw:examples}}
\begin{sphinxVerbatim}[commandchars=\\\{\}]
\PYG{k}{pair\PYGZus{}style}\PYG{+w}{ }\PYG{n}{gw}
\PYG{k}{pair\PYGZus{}coeff}\PYG{+w}{ }\PYG{o}{*}\PYG{+w}{ }\PYG{o}{*}\PYG{+w}{ }\PYG{n}{SiC.gw}\PYG{+w}{ }\PYG{n}{Si}\PYG{+w}{ }\PYG{n}{C}\PYG{+w}{ }\PYG{n}{C}

\PYG{k}{pair\PYGZus{}style}\PYG{+w}{ }\PYG{n}{gw}\PYG{o}{/}\PYG{n}{zbl}
\PYG{k}{pair\PYGZus{}coeff}\PYG{+w}{ }\PYG{o}{*}\PYG{+w}{ }\PYG{o}{*}\PYG{+w}{ }\PYG{n}{SiC.gw.zbl}\PYG{+w}{ }\PYG{n}{C}\PYG{+w}{ }\PYG{n}{Si}
\end{sphinxVerbatim}


\subsection{Description}
\label{\detokenize{pair_gw:description}}
\sphinxAtStartPar
The \sphinxstyleemphasis{gw} style computes a 3\sphinxhyphen{}body {\hyperref[\detokenize{pair_gw:gao}]{\sphinxcrossref{\DUrole{std,std-ref}{Gao\sphinxhyphen{}Weber}}}} potential;
similarly \sphinxstyleemphasis{gw/zbl} combines this potential with a modified
repulsive ZBL core function in a similar fashion as implemented
in the {\hyperref[\detokenize{pair_tersoff_zbl::doc}]{\sphinxcrossref{\DUrole{doc}{tersoff/zbl}}}} pair style.

\sphinxAtStartPar
Unfortunately the author of this contributed code has not been
able to submit a suitable documentation explaining the details
of the potentials. The LAMMPS developers thus have finally decided
to release the code anyway with only the technical explanations.
For details of the model and the parameters, please refer to the
linked publication.

\sphinxAtStartPar
Only a single pair\_coeff command is used with the \sphinxstyleemphasis{gw} and \sphinxstyleemphasis{gw/zbl}
styles which specifies a Gao\sphinxhyphen{}Weber potential file with parameters
for all needed elements.  These are mapped to LAMMPS atom types by
specifying N additional arguments after the filename in the pair\_coeff
command, where N is the number of LAMMPS atom types:
\begin{itemize}
\item {} 
\sphinxAtStartPar
filename

\item {} 
\sphinxAtStartPar
N element names = mapping of GW elements to atom types

\end{itemize}

\sphinxAtStartPar
See the {\hyperref[\detokenize{pair_coeff::doc}]{\sphinxcrossref{\DUrole{doc}{pair\_coeff}}}} page for alternate ways
to specify the path for the potential file.

\sphinxAtStartPar
As an example, imagine a file SiC.gw has Gao\sphinxhyphen{}Weber values for Si and C.
If your LAMMPS simulation has 4 atoms types and you want the first 3 to
be Si, and the fourth to be C, you would use the following pair\_coeff command:

\begin{sphinxVerbatim}[commandchars=\\\{\}]
\PYG{k}{pair\PYGZus{}coeff}\PYG{+w}{ }\PYG{o}{*}\PYG{+w}{ }\PYG{o}{*}\PYG{+w}{ }\PYG{n}{SiC.gw}\PYG{+w}{ }\PYG{n}{Si}\PYG{+w}{ }\PYG{n}{Si}\PYG{+w}{ }\PYG{n}{Si}\PYG{+w}{ }\PYG{n}{C}
\end{sphinxVerbatim}

\sphinxAtStartPar
The first 2 arguments must be * * so as to span all LAMMPS atom types.
The first three Si arguments map LAMMPS atom types 1,2,3 to the Si
element in the GW file.  The final C argument maps LAMMPS atom type 4
to the C element in the GW file.  If a mapping value is specified as
NULL, the mapping is not performed.  This can be used when a \sphinxstyleemphasis{gw}
potential is used as part of the \sphinxstyleemphasis{hybrid} pair style.  The NULL values
are placeholders for atom types that will be used with other
potentials.

\sphinxAtStartPar
Gao\sphinxhyphen{}Weber files in the \sphinxstyleemphasis{potentials} directory of the LAMMPS
distribution have a “.gw” suffix.  Gao\sphinxhyphen{}Weber with ZBL files
have a “.gz.zbl” suffix. The structure of the potential files
is similar to other many\sphinxhyphen{}body potentials supported by LAMMPS.
You have to refer to the comments in the files and the literature
to learn more details.


\bigskip\hrule\bigskip



\subsection{Mixing, shift, table, tail correction, restart, rRESPA info}
\label{\detokenize{pair_gw:mixing-shift-table-tail-correction-restart-rrespa-info}}
\sphinxAtStartPar
For atom type pairs I,J and I != J, where types I and J correspond to
two different element types, mixing is performed by LAMMPS as
described above from values in the potential file.

\sphinxAtStartPar
This pair style does not support the {\hyperref[\detokenize{pair_modify::doc}]{\sphinxcrossref{\DUrole{doc}{pair\_modify}}}}
shift, table, and tail options.

\sphinxAtStartPar
This pair style does not write its information to {\hyperref[\detokenize{restart::doc}]{\sphinxcrossref{\DUrole{doc}{binary restart files}}}}, since it is stored in potential files.  Thus, you
need to re\sphinxhyphen{}specify the pair\_style and pair\_coeff commands in an input
script that reads a restart file.

\sphinxAtStartPar
This pair style can only be used via the \sphinxstyleemphasis{pair} keyword of the
{\hyperref[\detokenize{run_style::doc}]{\sphinxcrossref{\DUrole{doc}{run\_style respa}}}} command.  It does not support the
\sphinxstyleemphasis{inner}, \sphinxstyleemphasis{middle}, \sphinxstyleemphasis{outer} keywords.


\bigskip\hrule\bigskip



\subsection{Restrictions}
\label{\detokenize{pair_gw:restrictions}}
\sphinxAtStartPar
This pair style is part of the MANYBODY package. It is only enabled if
LAMMPS was built with that package.  See the {\hyperref[\detokenize{Build_package::doc}]{\sphinxcrossref{\DUrole{doc}{Build package}}}} page for more info.

\sphinxAtStartPar
This pair style requires the {\hyperref[\detokenize{newton::doc}]{\sphinxcrossref{\DUrole{doc}{newton}}}} setting to be “on”
for pair interactions.

\sphinxAtStartPar
The Gao\sphinxhyphen{}Weber potential files provided with LAMMPS (see the
potentials directory) are parameterized for metal {\hyperref[\detokenize{units::doc}]{\sphinxcrossref{\DUrole{doc}{units}}}}.
You can use the GW potential with any LAMMPS units, but you would need
to create your own GW potential file with coefficients listed in the
appropriate units if your simulation does not use “metal” units.


\subsection{Related commands}
\label{\detokenize{pair_gw:related-commands}}
\sphinxAtStartPar
{\hyperref[\detokenize{pair_coeff::doc}]{\sphinxcrossref{\DUrole{doc}{pair\_coeff}}}}


\subsection{Default}
\label{\detokenize{pair_gw:default}}
\sphinxAtStartPar
none


\bigskip\hrule\bigskip

\phantomsection\label{\detokenize{pair_gw:gao}}
\sphinxAtStartPar
\sphinxstylestrong{(Gao)} Gao and Weber, Nuclear Instruments and Methods in Physics
Research B 191 (2012) 504.

\sphinxstepscope

\index{pair\_style harmonic/cut@\spxentry{pair\_style harmonic/cut}}\ignorespaces 
\index{pair\_style harmonic/cut/omp@\spxentry{pair\_style harmonic/cut/omp}}\ignorespaces 

\section{pair\_style harmonic/cut command}
\label{\detokenize{pair_harmonic_cut:pair-style-harmonic-cut-command}}\label{\detokenize{pair_harmonic_cut:index-1}}\label{\detokenize{pair_harmonic_cut:index-0}}\label{\detokenize{pair_harmonic_cut::doc}}
\sphinxAtStartPar
Accelerator Variants: \sphinxstyleemphasis{harmonic/cut/omp}


\subsection{Syntax}
\label{\detokenize{pair_harmonic_cut:syntax}}
\begin{sphinxVerbatim}[commandchars=\\\{\}]
\PYG{k}{pair\PYGZus{}style}\PYG{+w}{ }\PYG{n}{style}
\end{sphinxVerbatim}
\begin{itemize}
\item {} 
\sphinxAtStartPar
style = \sphinxstyleemphasis{harmonic/cut}

\end{itemize}


\subsection{Examples}
\label{\detokenize{pair_harmonic_cut:examples}}
\begin{sphinxVerbatim}[commandchars=\\\{\}]
\PYG{k}{pair\PYGZus{}style}\PYG{+w}{ }\PYG{n}{harmonic}\PYG{o}{/}\PYG{n}{cut}
\PYG{k}{pair\PYGZus{}coeff}\PYG{+w}{ }\PYG{o}{*}\PYG{+w}{ }\PYG{o}{*}\PYG{+w}{ }\PYG{l+m}{0.2}\PYG{+w}{ }\PYG{l+m}{2.0}
\PYG{k}{pair\PYGZus{}coeff}\PYG{+w}{ }\PYG{l+m}{1}\PYG{+w}{ }\PYG{l+m}{1}\PYG{+w}{ }\PYG{l+m}{0.5}\PYG{+w}{ }\PYG{l+m}{2.5}
\end{sphinxVerbatim}


\subsection{Description}
\label{\detokenize{pair_harmonic_cut:description}}
\sphinxAtStartPar
\DUrole{versionmodified,added}{Added in version 17Feb2022.}

\sphinxAtStartPar
Style \sphinxstyleemphasis{harmonic/cut} computes pairwise repulsive\sphinxhyphen{}only harmonic interactions with the formula
\begin{equation*}
\begin{split}E = k (r_c - r)^2  \qquad r < r_c\end{split}
\end{equation*}
\sphinxAtStartPar
where \(r_c\) is the cutoff.  Note that the usual 1/2 factor is
included in \(k\).

\sphinxAtStartPar
The following coefficients must be defined for each pair of atoms
types via the {\hyperref[\detokenize{pair_coeff::doc}]{\sphinxcrossref{\DUrole{doc}{pair\_coeff}}}} command as in the examples
above, or in the data file or restart files read by the
{\hyperref[\detokenize{read_data::doc}]{\sphinxcrossref{\DUrole{doc}{read\_data}}}} or {\hyperref[\detokenize{read_restart::doc}]{\sphinxcrossref{\DUrole{doc}{read\_restart}}}}
commands:
\begin{itemize}
\item {} 
\sphinxAtStartPar
\(k\) (energy/distance\textasciicircum{}2 units)

\item {} 
\sphinxAtStartPar
\(r_c\) (distance units)

\end{itemize}


\bigskip\hrule\bigskip


\sphinxAtStartPar
Styles with a \sphinxstyleemphasis{gpu}, \sphinxstyleemphasis{intel}, \sphinxstyleemphasis{kk}, \sphinxstyleemphasis{omp}, or \sphinxstyleemphasis{opt} suffix are
functionally the same as the corresponding style without the suffix.
They have been optimized to run faster, depending on your available
hardware, as discussed on the {\hyperref[\detokenize{Speed_packages::doc}]{\sphinxcrossref{\DUrole{doc}{Accelerator packages}}}}
page.  The accelerated styles take the same arguments and should
produce the same results, except for round\sphinxhyphen{}off and precision issues.

\sphinxAtStartPar
These accelerated styles are part of the GPU, INTEL, KOKKOS,
OPENMP, and OPT packages, respectively.  They are only enabled if
LAMMPS was built with those packages.  See the {\hyperref[\detokenize{Build_package::doc}]{\sphinxcrossref{\DUrole{doc}{Build package}}}} page for more info.

\sphinxAtStartPar
You can specify the accelerated styles explicitly in your input script
by including their suffix, or you can use the {\hyperref[\detokenize{Run_options::doc}]{\sphinxcrossref{\DUrole{doc}{\sphinxhyphen{}suffix command\sphinxhyphen{}line switch}}}} when you invoke LAMMPS, or you can use the
{\hyperref[\detokenize{suffix::doc}]{\sphinxcrossref{\DUrole{doc}{suffix}}}} command in your input script.

\sphinxAtStartPar
See the {\hyperref[\detokenize{Speed_packages::doc}]{\sphinxcrossref{\DUrole{doc}{Accelerator packages}}}} page for more
instructions on how to use the accelerated styles effectively.


\bigskip\hrule\bigskip



\subsection{Mixing, shift, table, tail correction, restart, rRESPA info}
\label{\detokenize{pair_harmonic_cut:mixing-shift-table-tail-correction-restart-rrespa-info}}
\sphinxAtStartPar
For atom type pairs I,J and I != J, the \(k\) and \(r_c\)
coefficients can be mixed. The default mix value is \sphinxstyleemphasis{geometric}.
See the “pair\_modify” command for details.

\sphinxAtStartPar
Since the potential is zero at and beyond the cutoff parameter by
construction, there is no need to support the {\hyperref[\detokenize{pair_modify::doc}]{\sphinxcrossref{\DUrole{doc}{pair\_modify}}}} shift or tail options for the energy and pressure of the
pair interaction.

\sphinxAtStartPar
These pair styles write their information to {\hyperref[\detokenize{restart::doc}]{\sphinxcrossref{\DUrole{doc}{binary restart files}}}},
so pair\_style and pair\_coeff commands do not need to be specified in an input script
that reads a restart file.

\sphinxAtStartPar
These pair styles can only be used via the \sphinxstyleemphasis{pair} keyword of the
{\hyperref[\detokenize{run_style::doc}]{\sphinxcrossref{\DUrole{doc}{run\_style respa}}}} command.  They do not support the
\sphinxstyleemphasis{inner}, \sphinxstyleemphasis{middle}, \sphinxstyleemphasis{outer} keywords.


\bigskip\hrule\bigskip



\subsection{Restrictions}
\label{\detokenize{pair_harmonic_cut:restrictions}}
\sphinxAtStartPar
The \sphinxstyleemphasis{harmonic/cut} pair style is only enabled if LAMMPS was
built with the EXTRA\sphinxhyphen{}PAIR package.
See the {\hyperref[\detokenize{Build_package::doc}]{\sphinxcrossref{\DUrole{doc}{Build package}}}} page for more info.


\subsection{Related commands}
\label{\detokenize{pair_harmonic_cut:related-commands}}
\sphinxAtStartPar
{\hyperref[\detokenize{pair_coeff::doc}]{\sphinxcrossref{\DUrole{doc}{pair\_coeff}}}}


\subsection{Default}
\label{\detokenize{pair_harmonic_cut:default}}
\sphinxAtStartPar
none

\sphinxstepscope

\index{pair\_style hbond/dreiding/lj@\spxentry{pair\_style hbond/dreiding/lj}}\ignorespaces 
\index{pair\_style hbond/dreiding/lj/omp@\spxentry{pair\_style hbond/dreiding/lj/omp}}\ignorespaces 
\index{pair\_style hbond/dreiding/morse@\spxentry{pair\_style hbond/dreiding/morse}}\ignorespaces 
\index{pair\_style hbond/dreiding/morse/omp@\spxentry{pair\_style hbond/dreiding/morse/omp}}\ignorespaces 

\section{pair\_style hbond/dreiding/lj command}
\label{\detokenize{pair_hbond_dreiding:pair-style-hbond-dreiding-lj-command}}\label{\detokenize{pair_hbond_dreiding:index-3}}\label{\detokenize{pair_hbond_dreiding:index-2}}\label{\detokenize{pair_hbond_dreiding:index-1}}\label{\detokenize{pair_hbond_dreiding:index-0}}\label{\detokenize{pair_hbond_dreiding::doc}}
\sphinxAtStartPar
Accelerator Variants: \sphinxstyleemphasis{hbond/dreiding/lj/omp}


\section{pair\_style hbond/dreiding/morse command}
\label{\detokenize{pair_hbond_dreiding:pair-style-hbond-dreiding-morse-command}}
\sphinxAtStartPar
Accelerator Variants: \sphinxstyleemphasis{hbond/dreiding/morse/omp}


\subsection{Syntax}
\label{\detokenize{pair_hbond_dreiding:syntax}}
\begin{sphinxVerbatim}[commandchars=\\\{\}]
\PYG{k}{pair\PYGZus{}style}\PYG{+w}{ }\PYG{n}{style}\PYG{+w}{ }\PYG{n}{N}\PYG{+w}{ }\PYG{n}{inner\PYGZus{}distance\PYGZus{}cutoff}\PYG{+w}{ }\PYG{n}{outer\PYGZus{}distance\PYGZus{}cutoff}\PYG{+w}{ }\PYG{n}{angle\PYGZus{}cutof}
\end{sphinxVerbatim}
\begin{itemize}
\item {} 
\sphinxAtStartPar
style = \sphinxstyleemphasis{hbond/dreiding/lj} or \sphinxstyleemphasis{hbond/dreiding/morse}

\item {} 
\sphinxAtStartPar
n = cosine angle periodicity

\item {} 
\sphinxAtStartPar
inner\_distance\_cutoff = global inner cutoff for Donor\sphinxhyphen{}Acceptor interactions (distance units)

\item {} 
\sphinxAtStartPar
outer\_distance\_cutoff = global cutoff for Donor\sphinxhyphen{}Acceptor interactions (distance units)

\item {} 
\sphinxAtStartPar
angle\_cutoff = global angle cutoff for Acceptor\sphinxhyphen{}Hydrogen\sphinxhyphen{}Donor

\item {} 
\sphinxAtStartPar
interactions (degrees)

\end{itemize}


\subsection{Examples}
\label{\detokenize{pair_hbond_dreiding:examples}}
\begin{sphinxVerbatim}[commandchars=\\\{\}]
\PYG{k}{pair\PYGZus{}style}\PYG{+w}{ }\PYG{n}{hybrid}\PYG{o}{/}\PYG{n}{overlay}\PYG{+w}{ }\PYG{n}{lj}\PYG{o}{/}\PYG{n}{cut}\PYG{+w}{ }\PYG{l+m}{10.0}\PYG{+w}{ }\PYG{n}{hbond}\PYG{o}{/}\PYG{n}{dreiding}\PYG{o}{/}\PYG{n}{lj}\PYG{+w}{ }\PYG{l+m}{4}\PYG{+w}{ }\PYG{l+m}{9.0}\PYG{+w}{ }\PYG{l+m}{11.0}\PYG{+w}{ }\PYG{l+m}{90}
\PYG{k}{pair\PYGZus{}coeff}\PYG{+w}{ }\PYG{l+m}{1}\PYG{+w}{ }\PYG{l+m}{2}\PYG{+w}{ }\PYG{n}{hbond}\PYG{o}{/}\PYG{n}{dreiding}\PYG{o}{/}\PYG{n}{lj}\PYG{+w}{ }\PYG{l+m}{3}\PYG{+w}{ }\PYG{n}{i}\PYG{+w}{ }\PYG{l+m}{9.5}\PYG{+w}{ }\PYG{l+m}{2.75}\PYG{+w}{ }\PYG{l+m}{4}\PYG{+w}{ }\PYG{l+m}{9.0}\PYG{+w}{ }\PYG{l+m}{11.0}\PYG{+w}{ }\PYG{l+m}{90.0}

\PYG{k}{pair\PYGZus{}style}\PYG{+w}{ }\PYG{n}{hybrid}\PYG{o}{/}\PYG{n}{overlay}\PYG{+w}{ }\PYG{n}{lj}\PYG{o}{/}\PYG{n}{cut}\PYG{+w}{ }\PYG{l+m}{10.0}\PYG{+w}{ }\PYG{n}{hbond}\PYG{o}{/}\PYG{n}{dreiding}\PYG{o}{/}\PYG{n}{morse}\PYG{+w}{ }\PYG{l+m}{2}\PYG{+w}{ }\PYG{l+m}{9.0}\PYG{+w}{ }\PYG{l+m}{11.0}\PYG{+w}{ }\PYG{l+m}{90}
\PYG{k}{pair\PYGZus{}coeff}\PYG{+w}{ }\PYG{l+m}{1}\PYG{+w}{ }\PYG{l+m}{2}\PYG{+w}{ }\PYG{n}{hbond}\PYG{o}{/}\PYG{n}{dreiding}\PYG{o}{/}\PYG{n}{morse}\PYG{+w}{ }\PYG{l+m}{3}\PYG{+w}{ }\PYG{n}{i}\PYG{+w}{ }\PYG{l+m}{3.88}\PYG{+w}{ }\PYG{l+m}{1.7241379}\PYG{+w}{ }\PYG{l+m}{2.9}\PYG{+w}{ }\PYG{l+m}{2}\PYG{+w}{ }\PYG{l+m}{9}\PYG{+w}{ }\PYG{l+m}{11}\PYG{+w}{ }\PYG{l+m}{90}

\PYG{k}{labelmap}\PYG{+w}{ }\PYG{n}{atom}\PYG{+w}{ }\PYG{l+m}{1}\PYG{+w}{ }\PYG{n}{C}\PYG{+w}{ }\PYG{l+m}{2}\PYG{+w}{ }\PYG{n}{O}\PYG{+w}{ }\PYG{l+m}{3}\PYG{+w}{ }\PYG{n}{H}
\PYG{k}{pair\PYGZus{}coeff}\PYG{+w}{ }\PYG{n}{C}\PYG{+w}{ }\PYG{n}{O}\PYG{+w}{ }\PYG{n}{hbond}\PYG{o}{/}\PYG{n}{dreiding}\PYG{o}{/}\PYG{n}{morse}\PYG{+w}{ }\PYG{n}{H}\PYG{+w}{ }\PYG{n}{i}\PYG{+w}{ }\PYG{l+m}{3.88}\PYG{+w}{ }\PYG{l+m}{1.7241379}\PYG{+w}{ }\PYG{l+m}{2.9}\PYG{+w}{ }\PYG{l+m}{2}\PYG{+w}{ }\PYG{l+m}{9}\PYG{+w}{ }\PYG{l+m}{11}\PYG{+w}{ }\PYG{l+m}{90}
\end{sphinxVerbatim}


\subsection{Description}
\label{\detokenize{pair_hbond_dreiding:description}}
\sphinxAtStartPar
The \sphinxstyleemphasis{hbond/dreiding} styles compute the Acceptor\sphinxhyphen{}Hydrogen\sphinxhyphen{}Donor (AHD)
3\sphinxhyphen{}body hydrogen bond interaction for the {\hyperref[\detokenize{Howto_bioFF::doc}]{\sphinxcrossref{\DUrole{doc}{DREIDING}}}}
force field, given by:
\begin{equation*}
\begin{split}E  = & \left[LJ(r) | Morse(r) \right] \qquad \qquad \qquad r < r_{\rm in} \\
   = & S(r) * \left[LJ(r) | Morse(r) \right] \qquad \qquad r_{\rm in} < r < r_{\rm out} \\
   = & 0 \qquad \qquad \qquad \qquad \qquad \qquad \qquad r > r_{\rm out} \\
LJ(r)  = & AR^{-12}-BR^{-10}cos^n\theta=
      \epsilon\left\lbrace 5\left[ \frac{\sigma}{r}\right]^{12}-
      6\left[ \frac{\sigma}{r}\right]^{10}  \right\rbrace cos^n\theta\\
Morse(r)  = & D_0\left\lbrace \chi^2 - 2\chi\right\rbrace cos^n\theta=
      D_{0}\left\lbrace e^{- 2 \alpha (r - r_0)} - 2 e^{- \alpha (r - r_0)}
      \right\rbrace cos^n\theta \\
S(r)  = & \frac{ \left[r_{\rm out}^2 - r^2\right]^2
\left[r_{\rm out}^2 + 2r^2 - 3{r_{\rm in}^2}\right]}
{ \left[r_{\rm out}^2 - {r_{\rm in}}^2\right]^3 }\end{split}
\end{equation*}
\sphinxAtStartPar
where \(r_{\rm in}\) is the inner spline distance cutoff,
\(r_{\rm out}\) is the outer distance cutoff, \(\theta_c\) is
the angle cutoff, and \sphinxstyleemphasis{n} is the cosine periodicity.

\sphinxAtStartPar
Here, \sphinxstyleemphasis{r} is the radial distance between the donor (D) and acceptor
(A) atoms and \(\theta\) is the bond angle between the acceptor, the
hydrogen (H) and the donor atoms:

\noindent{\hspace*{\fill}\sphinxincludegraphics{{dreiding_hbond}.jpg}\hspace*{\fill}}

\sphinxAtStartPar
These 3\sphinxhyphen{}body interactions can be defined for pairs of acceptor and
donor atoms, based on atom types.  For each donor/acceptor atom pair,
the third atom in the interaction is a hydrogen permanently bonded to
the donor atom, e.g. in a bond list read in from a data file via the
{\hyperref[\detokenize{read_data::doc}]{\sphinxcrossref{\DUrole{doc}{read\_data}}}} command.  The atom types of possible
hydrogen atoms for each donor/acceptor type pair are specified by the
{\hyperref[\detokenize{pair_coeff::doc}]{\sphinxcrossref{\DUrole{doc}{pair\_coeff}}}} command (see below).

\sphinxAtStartPar
Style \sphinxstyleemphasis{hbond/dreiding/lj} is the original DREIDING potential of
{\hyperref[\detokenize{pair_hbond_dreiding:pair-mayo}]{\sphinxcrossref{\DUrole{std,std-ref}{(Mayo)}}}}.  It uses a LJ 12/10 functional for the Donor\sphinxhyphen{}Acceptor
interactions. To match the results in the original paper, use n = 4.

\sphinxAtStartPar
Style \sphinxstyleemphasis{hbond/dreiding/morse} is an improved version using a Morse
potential for the Donor\sphinxhyphen{}Acceptor interactions. {\hyperref[\detokenize{pair_hbond_dreiding:liu}]{\sphinxcrossref{\DUrole{std,std-ref}{(Liu)}}}} showed
that the Morse form gives improved results for Dendrimer simulations,
when n = 2.

\sphinxAtStartPar
See the {\hyperref[\detokenize{Howto_bioFF::doc}]{\sphinxcrossref{\DUrole{doc}{Howto bioFF}}}} page for more information
on the DREIDING force field.

\begin{sphinxadmonition}{note}{Note:}
\sphinxAtStartPar
Because the Dreiding hydrogen bond potential is only one portion
of an overall force field which typically includes other pairwise
interactions, it is common to use it as a sub\sphinxhyphen{}style in a {\hyperref[\detokenize{pair_hybrid::doc}]{\sphinxcrossref{\DUrole{doc}{pair\_style hybrid/overlay}}}} command, where another pair style
provides the repulsive core interaction between pairs of atoms, e.g. a
1/r\textasciicircum{}12 Lennard\sphinxhyphen{}Jones repulsion.
\end{sphinxadmonition}

\begin{sphinxadmonition}{note}{Note:}
\sphinxAtStartPar
When using the hbond/dreiding pair styles with {\hyperref[\detokenize{pair_hybrid::doc}]{\sphinxcrossref{\DUrole{doc}{pair\_style hybrid/overlay}}}}, you should explicitly define pair
interactions between the donor atom and acceptor atoms, (as well as
between these atoms and ALL other atoms in your system).  Whenever
{\hyperref[\detokenize{pair_hybrid::doc}]{\sphinxcrossref{\DUrole{doc}{pair\_style hybrid/overlay}}}} is used, ordinary mixing
rules are not applied to atoms like the donor and acceptor atoms
because they are typically referenced in multiple pair styles.
Neglecting to do this can cause difficult\sphinxhyphen{}to\sphinxhyphen{}detect physics problems.
\end{sphinxadmonition}

\begin{sphinxadmonition}{note}{Note:}
\sphinxAtStartPar
In the original Dreiding force field paper 1\sphinxhyphen{}4 non\sphinxhyphen{}bonded
interactions ARE allowed.  If this is desired for your model, use the
special\_bonds command (e.g. “special\_bonds lj 0.0 0.0 1.0”) to turn
these interactions on.
\end{sphinxadmonition}


\bigskip\hrule\bigskip


\sphinxAtStartPar
The following coefficients must be defined for pairs of eligible
donor/acceptor types via the {\hyperref[\detokenize{pair_coeff::doc}]{\sphinxcrossref{\DUrole{doc}{pair\_coeff}}}} command as
in the examples above.

\begin{sphinxadmonition}{note}{Note:}
\sphinxAtStartPar
Unlike other pair styles and their associated
{\hyperref[\detokenize{pair_coeff::doc}]{\sphinxcrossref{\DUrole{doc}{pair\_coeff}}}} commands, you do not need to specify
pair\_coeff settings for all possible I,J type pairs.  Only I,J type
pairs for atoms which act as joint donors/acceptors need to be
specified; all other type pairs are assumed to be inactive.
\end{sphinxadmonition}

\begin{sphinxadmonition}{note}{Note:}
\sphinxAtStartPar
A {\hyperref[\detokenize{pair_coeff::doc}]{\sphinxcrossref{\DUrole{doc}{pair\_coeff}}}} command can be specified multiple
times for the same donor/acceptor type pair.  This enables multiple
hydrogen types to be assigned to the same donor/acceptor type pair.
For other pair\_styles, if the pair\_coeff command is re\sphinxhyphen{}used for the
same I.J type pair, the settings for that type pair are overwritten.
For the hydrogen bond potentials this is not the case; the settings
are cumulative.  This means the only way to turn off a previous
setting, is to re\sphinxhyphen{}use the pair\_style command and start over.
\end{sphinxadmonition}

\sphinxAtStartPar
For the \sphinxstyleemphasis{hbond/dreiding/lj} style the list of coefficients is as
follows:
\begin{itemize}
\item {} 
\sphinxAtStartPar
K = hydrogen atom type = 1 to Ntypes, or type label

\item {} 
\sphinxAtStartPar
donor flag = \sphinxstyleemphasis{i} or \sphinxstyleemphasis{j}

\item {} 
\sphinxAtStartPar
\(\epsilon\) (energy units)

\item {} 
\sphinxAtStartPar
\(\sigma\) (distance units)

\item {} 
\sphinxAtStartPar
\sphinxstyleemphasis{n} = exponent in formula above

\item {} 
\sphinxAtStartPar
distance cutoff \(r_{\rm in}\) (distance units)

\item {} 
\sphinxAtStartPar
distance cutoff \(r_{\rm out}\) (distance units)

\item {} 
\sphinxAtStartPar
angle cutoff (degrees)

\end{itemize}

\sphinxAtStartPar
For the \sphinxstyleemphasis{hbond/dreiding/morse} style the list of coefficients is as
follows:
\begin{itemize}
\item {} 
\sphinxAtStartPar
K = hydrogen atom type = 1 to Ntypes, or type label

\item {} 
\sphinxAtStartPar
donor flag = \sphinxstyleemphasis{i} or \sphinxstyleemphasis{j}

\item {} 
\sphinxAtStartPar
\(D_0\) (energy units)

\item {} 
\sphinxAtStartPar
\(\alpha\) (1/distance units)

\item {} 
\sphinxAtStartPar
\(r_0\) (distance units)

\item {} 
\sphinxAtStartPar
\sphinxstyleemphasis{n} = exponent in formula above

\item {} 
\sphinxAtStartPar
distance cutoff \(r_{\rm in}\) (distance units)

\item {} 
\sphinxAtStartPar
distance cutoff \(r_{out}\) (distance units)

\item {} 
\sphinxAtStartPar
angle cutoff (degrees)

\end{itemize}

\sphinxAtStartPar
A single hydrogen atom type K can be specified, or a wild\sphinxhyphen{}card asterisk
can be used in place of or in conjunction with the K arguments to
select multiple types as hydrogen atoms.  This takes the form
“*” or “*n” or “n*” or “m*n”.  See the {\hyperref[\detokenize{pair_coeff::doc}]{\sphinxcrossref{\DUrole{doc}{pair\_coeff}}}}
command page for details.

\sphinxAtStartPar
If the donor flag is \sphinxstyleemphasis{i}, then the atom of type I in the pair\_coeff
command is treated as the donor, and J is the acceptor.  If the donor
flag is \sphinxstyleemphasis{j}, then the atom of type J in the pair\_coeff command is
treated as the donor and I is the donor.  This option is required
because the {\hyperref[\detokenize{pair_coeff::doc}]{\sphinxcrossref{\DUrole{doc}{pair\_coeff}}}} command requires that I \textless{}= J.

\sphinxAtStartPar
\(\epsilon\) and \(\sigma\) are settings for the hydrogen bond
potential based on a Lennard\sphinxhyphen{}Jones functional form.  Note that sigma is
defined as the zero\sphinxhyphen{}crossing distance for the potential, not as the
energy minimum at \(2^{1/6} \sigma\).

\sphinxAtStartPar
\(D_0\) and \(\alpha\) and \(r_0\) are settings for the
hydrogen bond potential based on a Morse functional form.

\sphinxAtStartPar
The last 3 coefficients for both styles are optional.  If not
specified, the global n, distance cutoff, and angle cutoff specified
in the pair\_style command are used.  If you wish to only override the
second or third optional parameter, you must also specify the preceding
optional parameters.


\bigskip\hrule\bigskip


\sphinxAtStartPar
Styles with a \sphinxstyleemphasis{gpu}, \sphinxstyleemphasis{intel}, \sphinxstyleemphasis{kk}, \sphinxstyleemphasis{omp}, or \sphinxstyleemphasis{opt} suffix are
functionally the same as the corresponding style without the suffix.
They have been optimized to run faster, depending on your available
hardware, as discussed on the {\hyperref[\detokenize{Speed_packages::doc}]{\sphinxcrossref{\DUrole{doc}{Accelerator packages}}}}
page.  The accelerated styles take the same arguments and should
produce the same results, except for round\sphinxhyphen{}off and precision issues.

\sphinxAtStartPar
These accelerated styles are part of the GPU, INTEL, KOKKOS,
OPENMP, and OPT packages, respectively.  They are only enabled if
LAMMPS was built with those packages.  See the {\hyperref[\detokenize{Build_package::doc}]{\sphinxcrossref{\DUrole{doc}{Build package}}}} page for more info.

\sphinxAtStartPar
You can specify the accelerated styles explicitly in your input script
by including their suffix, or you can use the {\hyperref[\detokenize{Run_options::doc}]{\sphinxcrossref{\DUrole{doc}{\sphinxhyphen{}suffix command\sphinxhyphen{}line switch}}}} when you invoke LAMMPS, or you can use the
{\hyperref[\detokenize{suffix::doc}]{\sphinxcrossref{\DUrole{doc}{suffix}}}} command in your input script.

\sphinxAtStartPar
See the {\hyperref[\detokenize{Speed_packages::doc}]{\sphinxcrossref{\DUrole{doc}{Accelerator packages}}}} page for more
instructions on how to use the accelerated styles effectively.


\bigskip\hrule\bigskip



\subsection{Mixing, shift, table, tail correction, restart, rRESPA info}
\label{\detokenize{pair_hbond_dreiding:mixing-shift-table-tail-correction-restart-rrespa-info}}
\sphinxAtStartPar
These pair styles do not support mixing. You must explicitly identify
each donor/acceptor type pair.

\sphinxAtStartPar
These styles do not support the {\hyperref[\detokenize{pair_modify::doc}]{\sphinxcrossref{\DUrole{doc}{pair\_modify}}}} shift
option for the energy of the interactions.

\sphinxAtStartPar
The {\hyperref[\detokenize{pair_modify::doc}]{\sphinxcrossref{\DUrole{doc}{pair\_modify}}}} table option is not relevant for
these pair styles.

\sphinxAtStartPar
These pair styles do not support the {\hyperref[\detokenize{pair_modify::doc}]{\sphinxcrossref{\DUrole{doc}{pair\_modify}}}}
tail option for adding long\sphinxhyphen{}range tail corrections to energy and
pressure.

\sphinxAtStartPar
These pair styles do not write their information to {\hyperref[\detokenize{restart::doc}]{\sphinxcrossref{\DUrole{doc}{binary restart files}}}}, so pair\_style and pair\_coeff commands need to be
re\sphinxhyphen{}specified in an input script that reads a restart file.

\sphinxAtStartPar
These pair styles can only be used via the \sphinxstyleemphasis{pair} keyword of the
{\hyperref[\detokenize{run_style::doc}]{\sphinxcrossref{\DUrole{doc}{run\_style respa}}}} command.  They do not support the
\sphinxstyleemphasis{inner}, \sphinxstyleemphasis{middle}, \sphinxstyleemphasis{outer} keywords.

\sphinxAtStartPar
These pair styles tally a count of how many hydrogen bonding
interactions they calculate each timestep and the hbond energy.  These
quantities can be accessed via the {\hyperref[\detokenize{compute_pair::doc}]{\sphinxcrossref{\DUrole{doc}{compute pair}}}}
command as a vector of values of length 2.

\sphinxAtStartPar
To print these quantities to the log file (with a descriptive column
heading) the following commands could be included in an input script:

\begin{sphinxVerbatim}[commandchars=\\\{\}]
\PYG{k}{compute }\PYG{n+nv+nvIdentifier}{hb}\PYG{+w}{ }\PYG{n+nv+nvIdentifier}{all}\PYG{+w}{ }\PYG{n}{pair}\PYG{+w}{ }\PYG{n}{hbond}\PYG{o}{/}\PYG{n}{dreiding}\PYG{o}{/}\PYG{n}{lj}
\PYG{k}{variable }\PYG{n+nv+nvIdentifier}{n\PYGZus{}hbond}\PYG{+w}{ }\PYG{n}{equal}\PYG{+w}{ }\PYG{n}{c\PYGZus{}hb[1]}\PYG{+w}{ }\PYG{c}{\PYGZsh{}number hbonds}
\PYG{k}{variable }\PYG{n+nv+nvIdentifier}{E\PYGZus{}hbond}\PYG{+w}{ }\PYG{n}{equal}\PYG{+w}{ }\PYG{n}{c\PYGZus{}hb[2]}\PYG{+w}{ }\PYG{c}{\PYGZsh{}hbond energy}
\PYG{k}{thermo\PYGZus{}style}\PYG{+w}{ }\PYG{n}{custom}\PYG{+w}{ }\PYG{n}{step}\PYG{+w}{ }\PYG{n}{temp}\PYG{+w}{ }\PYG{n}{epair}\PYG{+w}{ }\PYG{n}{v\PYGZus{}E\PYGZus{}hbond}
\end{sphinxVerbatim}


\bigskip\hrule\bigskip



\subsection{Restrictions}
\label{\detokenize{pair_hbond_dreiding:restrictions}}
\sphinxAtStartPar
This pair style can only be used if LAMMPS was built with the
MOLECULE package.  See the {\hyperref[\detokenize{Build_package::doc}]{\sphinxcrossref{\DUrole{doc}{Build package}}}} doc page
for more info.


\subsection{Related commands}
\label{\detokenize{pair_hbond_dreiding:related-commands}}
\sphinxAtStartPar
{\hyperref[\detokenize{pair_coeff::doc}]{\sphinxcrossref{\DUrole{doc}{pair\_coeff}}}}


\subsection{Default}
\label{\detokenize{pair_hbond_dreiding:default}}
\sphinxAtStartPar
none


\bigskip\hrule\bigskip

\phantomsection\label{\detokenize{pair_hbond_dreiding:pair-mayo}}
\sphinxAtStartPar
\sphinxstylestrong{(Mayo)} Mayo, Olfason, Goddard III, J Phys Chem, 94, 8897\sphinxhyphen{}8909
(1990).

\phantomsection\label{\detokenize{pair_hbond_dreiding:liu}}
\sphinxAtStartPar
\sphinxstylestrong{(Liu)} Liu, Bryantsev, Diallo, Goddard III, J. Am. Chem. Soc 131 (8)
2798 (2009)

\sphinxstepscope

\index{pair\_style hdnnp@\spxentry{pair\_style hdnnp}}\ignorespaces 

\section{pair\_style hdnnp command}
\label{\detokenize{pair_hdnnp:pair-style-hdnnp-command}}\label{\detokenize{pair_hdnnp:index-0}}\label{\detokenize{pair_hdnnp::doc}}

\subsection{Syntax}
\label{\detokenize{pair_hdnnp:syntax}}
\begin{sphinxVerbatim}[commandchars=\\\{\}]
\PYG{k}{pair\PYGZus{}style}\PYG{+w}{ }\PYG{n}{hdnnp}\PYG{+w}{ }\PYG{n}{cutoff}\PYG{+w}{ }\PYG{n}{keyword}\PYG{+w}{ }\PYG{n}{value}\PYG{+w}{ }\PYG{n}{...}
\end{sphinxVerbatim}
\begin{itemize}
\item {} 
\sphinxAtStartPar
cutoff = short\sphinxhyphen{}range cutoff of HDNNP (maximum symmetry function cutoff radius)

\item {} 
\sphinxAtStartPar
zero or more keyword/value pairs may be appended

\item {} 
\sphinxAtStartPar
keyword = \sphinxstyleemphasis{dir} or \sphinxstyleemphasis{showew} or \sphinxstyleemphasis{showewsum} or \sphinxstyleemphasis{maxew} or \sphinxstyleemphasis{resetew} or \sphinxstyleemphasis{cflength} or \sphinxstyleemphasis{cfenergy}

\item {} 
\sphinxAtStartPar
value depends on the preceding keyword:

\end{itemize}
\begin{sphinxalltt}
\sphinxstyleemphasis{dir} value = directory
    directory = Path to HDNNP configuration files
\sphinxstyleemphasis{showew} value = \sphinxstyleemphasis{yes} or \sphinxstyleemphasis{no}
\sphinxstyleemphasis{showewsum} value = summary
          summary = Write EW summary every this many timesteps (\sphinxstyleemphasis{0} turns summary off)
\sphinxstyleemphasis{maxew} value = threshold
      threshold = Maximum number of EWs allowed
\sphinxstyleemphasis{resetew} value = \sphinxstyleemphasis{yes} or \sphinxstyleemphasis{no}
\sphinxstyleemphasis{cflength} value = length
         length = Length unit conversion factor
\sphinxstyleemphasis{cfenergy} value = energy
         energy = Energy unit conversion factor
\end{sphinxalltt}


\subsection{Examples}
\label{\detokenize{pair_hdnnp:examples}}
\begin{sphinxVerbatim}[commandchars=\\\{\}]
\PYG{k}{pair\PYGZus{}style}\PYG{+w}{ }\PYG{n}{hdnnp}\PYG{+w}{ }\PYG{l+m}{6.35}\PYG{+w}{ }\PYG{n}{showew}\PYG{+w}{ }\PYG{n}{yes}\PYG{+w}{ }\PYG{n}{showewsum}\PYG{+w}{ }\PYG{l+m}{100}\PYG{+w}{ }\PYG{n}{maxew}\PYG{+w}{ }\PYG{l+m}{1000}\PYG{+w}{ }\PYG{n}{resetew}\PYG{+w}{ }\PYG{n}{yes}\PYG{+w}{ }\PYG{n}{cflength}\PYG{+w}{ }\PYG{l+m}{1.8897261328}\PYG{+w}{ }\PYG{n}{cfenergy}\PYG{+w}{ }\PYG{l+m}{0.0367493254}
\PYG{k}{pair\PYGZus{}coeff}\PYG{+w}{ }\PYG{o}{*}\PYG{+w}{ }\PYG{o}{*}\PYG{+w}{ }\PYG{n}{H}\PYG{+w}{ }\PYG{n}{O}

\PYG{k}{pair\PYGZus{}style}\PYG{+w}{ }\PYG{n}{hdnnp}\PYG{+w}{ }\PYG{l+m}{6.01}\PYG{+w}{ }\PYG{n}{dir}\PYG{+w}{ }\PYG{l+s}{\PYGZdq{}}\PYG{l+s}{./}\PYG{l+s}{\PYGZdq{}}\PYG{+w}{ }\PYG{n}{showewsum}\PYG{+w}{ }\PYG{l+m}{10000}
\PYG{k}{pair\PYGZus{}coeff}\PYG{+w}{ }\PYG{o}{*}\PYG{+w}{ }\PYG{o}{*}\PYG{+w}{ }\PYG{n}{S}\PYG{+w}{ }\PYG{n}{Cu}\PYG{+w}{ }\PYG{n}{NULL}\PYG{+w}{ }\PYG{n}{Cu}
\end{sphinxVerbatim}


\subsection{Description}
\label{\detokenize{pair_hdnnp:description}}
\sphinxAtStartPar
This pair style adds an interaction based on the high\sphinxhyphen{}dimensional neural network
potential (HDNNP) method as presented in {\hyperref[\detokenize{pair_hdnnp:behler-parrinello-2007}]{\sphinxcrossref{\DUrole{std,std-ref}{(Behler and Parrinello 2007)}}}}. HDNNPs are machine learning potentials which require
careful training of neural networks prior to application in MD simulations. The
pair style uses an interface to the \sphinxstyleemphasis{n2p2} library {\hyperref[\detokenize{pair_hdnnp:singraber-behler-dellago-2019}]{\sphinxcrossref{\DUrole{std,std-ref}{(Singraber, Behler and
Dellago 2019)}}}} which is available on Github
\sphinxhref{https://github.com/CompPhysVienna/n2p2}{here}. Please see the \sphinxstyleemphasis{n2p2}
\sphinxhref{https://compphysvienna.github.io/n2p2/}{documentation} for further details.
\sphinxstyleemphasis{n2p2} (and hence this pair style) is compatible with neural network potentials
trained with its own tools {\hyperref[\detokenize{pair_hdnnp:singraber-et-al-2019}]{\sphinxcrossref{\DUrole{std,std-ref}{(Singraber et al 2019)}}}}
and with \sphinxhref{https://www.uni-goettingen.de/de/560580.html}{RuNNer}.

\sphinxAtStartPar
Only a single \sphinxstyleemphasis{pair\_coeff} command with two asterisk wild\sphinxhyphen{}cards is used with this
pair style. Its additional arguments define the mapping of LAMMPS atom types to
n2p2 elements.

\begin{sphinxVerbatim}[commandchars=\\\{\}]
\PYG{k}{pair\PYGZus{}coeff}\PYG{+w}{ }\PYG{o}{*}\PYG{+w}{ }\PYG{o}{*}\PYG{+w}{ }\PYG{n}{H}\PYG{+w}{ }\PYG{n}{O}
\end{sphinxVerbatim}

\sphinxAtStartPar
In the above example LAMMPS types 1 and 2 are mapped to the elements “H” and “O”
in n2p2, respectively. Multiple types may map to the same element, or some types
may not be mapped at all. For example, if the LAMMPS simulation has four atom
types, the command

\begin{sphinxVerbatim}[commandchars=\\\{\}]
\PYG{k}{pair\PYGZus{}coeff}\PYG{+w}{ }\PYG{o}{*}\PYG{+w}{ }\PYG{o}{*}\PYG{+w}{ }\PYG{n}{H}\PYG{+w}{ }\PYG{n}{H}\PYG{+w}{ }\PYG{n}{O}\PYG{+w}{ }\PYG{n}{NULL}
\end{sphinxVerbatim}

\sphinxAtStartPar
maps atom types 1 and 2 to the element “H”, type 3 to “O” and type 4 is not mapped
(indicated by NULL).  Atoms mapped to NULL are ignored by the HDNNP calculation,
i.e. they do not contribute in any way to the evaluation of HDNNP energies and forces.
This may be useful in a setup with {\hyperref[\detokenize{pair_hybrid::doc}]{\sphinxcrossref{\DUrole{doc}{hybrid pair styles}}}}.


\bigskip\hrule\bigskip


\sphinxAtStartPar
The mandatory pair style argument \sphinxstyleemphasis{cutoff} must match the short\sphinxhyphen{}range cutoff radius
of the HDNNP.  This corresponds to the maximum cutoff radius of all symmetry
functions (the atomic environment descriptors of HDNNPs) used.

\begin{sphinxadmonition}{note}{Note:}
\sphinxAtStartPar
The cutoff must be given in LAMMPS length units, even if the neural network
potential has been trained using a different unit system (see remarks about the
\sphinxstyleemphasis{cflength} and \sphinxstyleemphasis{cfenergy} keywords below for details).
\end{sphinxadmonition}

\sphinxAtStartPar
The numeric value may be slightly larger than the actual maximum symmetry
function cutoff radius (to account for rounding errors when converting units),
but must not be smaller.

\sphinxAtStartPar
Use the \sphinxstyleemphasis{dir} keyword to specify the directory containing the HDNNP configuration
files. The directory must contain \sphinxcode{\sphinxupquote{input.nn}} with neural network and symmetry
function setup, \sphinxcode{\sphinxupquote{scaling.data}} with symmetry function scaling data and
\sphinxcode{\sphinxupquote{weights.???.data}} with weight parameters for each element.

\sphinxAtStartPar
The keyword \sphinxstyleemphasis{showew} can be used to turn on/off the display of extrapolation
warnings (EWs) which are issued whenever a symmetry function value is out of
bounds defined by minimum/maximum values in \sphinxcode{\sphinxupquote{scaling.data}}. An extrapolation
warning may look like this:

\begin{sphinxVerbatim}[commandchars=\\\{\}]
\PYG{c}{\PYGZsh{}\PYGZsh{}\PYGZsh{} NNP EXTRAPOLATION WARNING \PYGZsh{}\PYGZsh{}\PYGZsh{} STRUCTURE:      0 ATOM:       119 ELEMENT: Cu SYMFUNC:   32 TYPE:  3 VALUE:  2.166E\PYGZhy{}02 MIN:  2.003E\PYGZhy{}05 MAX:  1.756E\PYGZhy{}02}
\end{sphinxVerbatim}

\sphinxAtStartPar
stating that the value 2.166E\sphinxhyphen{}02 of symmetry function 32 of type 3 (Narrow Angular symmetry function), element Cu (see the log file for a symmetry function listing) was out
of bounds (maximum in \sphinxcode{\sphinxupquote{scaling.data}} is 1.756E\sphinxhyphen{}02) for atom 119. Here, the
atom index refers to the LAMMPS tag (global index) and the structure index is
used to print out the MPI rank the atom belongs to.

\begin{sphinxadmonition}{note}{Note:}
\sphinxAtStartPar
The \sphinxstyleemphasis{showew} keyword should only be set to \sphinxstyleemphasis{yes} for debugging purposes.
Extrapolation warnings may add lots of overhead as they are communicated each
timestep. Also, if the simulation is run in a region where the HDNNP was not
correctly trained, lots of extrapolation warnings may clog log files and the
console. In a production run use \sphinxstyleemphasis{showewsum} instead.
\end{sphinxadmonition}

\sphinxAtStartPar
The keyword \sphinxstyleemphasis{showewsum} can be used to get an overview of extrapolation warnings
occurring during an MD simulation. The argument specifies the interval at which
extrapolation warning summaries are displayed and logged. An EW summary may look
like this:

\begin{sphinxVerbatim}[commandchars=\\\{\}]
\PYG{c}{\PYGZsh{}\PYGZsh{}\PYGZsh{} NNP EW SUMMARY \PYGZsh{}\PYGZsh{}\PYGZsh{} TS:        100 EW         11 EWPERSTEP  1.100E\PYGZhy{}01}
\end{sphinxVerbatim}

\sphinxAtStartPar
Here, at timestep 100 the occurrence of 11 extrapolation warnings since the last
summary is reported, which corresponds to an EW rate of 0.11 per timestep.
Setting \sphinxstyleemphasis{showewsum} to 0 deactivates the EW summaries.

\sphinxAtStartPar
A maximum number of allowed extrapolation warnings can be specified with the
\sphinxstyleemphasis{maxew} keyword. If the number of EWs exceeds the \sphinxstyleemphasis{maxew} argument the
simulation is stopped. Note however that this is merely an approximate threshold
since the check is only performed at the end of each timestep and each MPI
process counts individually to minimize communication overhead.

\sphinxAtStartPar
The keyword \sphinxstyleemphasis{resetew} alters the behavior of the above mentioned \sphinxstyleemphasis{maxew}
threshold. If \sphinxstyleemphasis{resetew} is set to \sphinxstyleemphasis{yes} the threshold is applied on a
per\sphinxhyphen{}timestep basis and the internal EW counters are reset at the beginning of
each timestep. With \sphinxstyleemphasis{resetew} set to \sphinxstyleemphasis{no} the counters accumulate EWs along the
whole trajectory.

\sphinxAtStartPar
If the training of a neural network potential has been performed with different
physical units for length and energy than those set in LAMMPS, it is still
possible to use the potential when the unit conversion factors are provided via
the \sphinxstyleemphasis{cflength} and \sphinxstyleemphasis{cfenergy} keywords. If for example, the HDNNP was
parameterized with Bohr and Hartree training data and symmetry function
parameters (i.e. distances and energies in “input.nn” are given in Bohr and
Hartree) but LAMMPS is set to use \sphinxstyleemphasis{metal} units (Angstrom and eV) the correct
conversion factors are:

\begin{sphinxVerbatim}[commandchars=\\\{\}]
\PYG{n}{cflength}\PYG{+w}{ }\PYG{l+m}{1.8897261328}

\PYG{n}{cfenergy}\PYG{+w}{ }\PYG{l+m}{0.0367493254}
\end{sphinxVerbatim}

\sphinxAtStartPar
Thus, arguments of \sphinxstyleemphasis{cflength} and \sphinxstyleemphasis{cfenergy} are the multiplicative factors
required to convert lengths and energies given in LAMMPS units to respective
quantities in native HDNNP units (1 Angstrom = 1.8897261328 Bohr, 1 eV =
0.0367493254 Hartree).


\bigskip\hrule\bigskip



\subsection{Mixing, shift, table, tail correction, restart, rRESPA info}
\label{\detokenize{pair_hdnnp:mixing-shift-table-tail-correction-restart-rrespa-info}}
\sphinxAtStartPar
This style does not support mixing. The {\hyperref[\detokenize{pair_coeff::doc}]{\sphinxcrossref{\DUrole{doc}{pair\_coeff}}}} command
should only be invoked with asterisk wild cards (see above).

\sphinxAtStartPar
This style does not support the {\hyperref[\detokenize{pair_modify::doc}]{\sphinxcrossref{\DUrole{doc}{pair\_modify}}}}
shift, table, and tail options.

\sphinxAtStartPar
This style does not write information to {\hyperref[\detokenize{restart::doc}]{\sphinxcrossref{\DUrole{doc}{binary restart files}}}}.
Thus, you need to re\sphinxhyphen{}specify the pair\_style and pair\_coeff commands in an input
script that reads a restart file.

\sphinxAtStartPar
This style can only be used via the \sphinxstyleemphasis{pair} keyword of the {\hyperref[\detokenize{run_style::doc}]{\sphinxcrossref{\DUrole{doc}{run\_style respa}}}} command.  It does not support the \sphinxstyleemphasis{inner}, \sphinxstyleemphasis{middle}, \sphinxstyleemphasis{outer}
keywords.


\subsection{Restrictions}
\label{\detokenize{pair_hdnnp:restrictions}}
\sphinxAtStartPar
This pair style is part of the ML\sphinxhyphen{}HDNNP package.  It is only enabled if LAMMPS
was built with that package.  See the {\hyperref[\detokenize{Build_package::doc}]{\sphinxcrossref{\DUrole{doc}{Build package}}}} doc
page for more info.

\sphinxAtStartPar
Please report bugs and feature requests to the \sphinxhref{https://github.com/CompPhysVienna/n2p2/issues}{n2p2 GitHub issue page}.


\subsubsection{Related commands}
\label{\detokenize{pair_hdnnp:related-commands}}
\sphinxAtStartPar
{\hyperref[\detokenize{pair_coeff::doc}]{\sphinxcrossref{\DUrole{doc}{pair\_coeff}}}}, {\hyperref[\detokenize{pair_hybrid::doc}]{\sphinxcrossref{\DUrole{doc}{pair\_hybrid}}}}, {\hyperref[\detokenize{units::doc}]{\sphinxcrossref{\DUrole{doc}{units}}}}


\subsubsection{Default}
\label{\detokenize{pair_hdnnp:default}}
\sphinxAtStartPar
The default options are \sphinxstyleemphasis{dir} = “hdnnp/”, \sphinxstyleemphasis{showew} = yes, \sphinxstyleemphasis{showewsum} = 0, \sphinxstyleemphasis{maxew}
= 0, \sphinxstyleemphasis{resetew} = no, \sphinxstyleemphasis{cflength} = 1.0, \sphinxstyleemphasis{cfenergy} = 1.0.


\bigskip\hrule\bigskip

\phantomsection\label{\detokenize{pair_hdnnp:behler-parrinello-2007}}
\sphinxAtStartPar
\sphinxstylestrong{(Behler and Parrinello 2007)} Behler, J.; Parrinello, M. Phys. Rev. Lett.  2007, 98 (14), 146401.

\phantomsection\label{\detokenize{pair_hdnnp:singraber-behler-dellago-2019}}
\sphinxAtStartPar
\sphinxstylestrong{(Singraber, Behler and Dellago 2019)} Singraber, A.; Behler, J.; Dellago, C. J., Chem.  Theory Comput. 2019, 15 (3), 1827\sphinxhyphen{}1840

\phantomsection\label{\detokenize{pair_hdnnp:singraber-et-al-2019}}
\sphinxAtStartPar
\sphinxstylestrong{(Singraber et al 2019)} Singraber, A.; Morawietz, T.; Behler, J.; Dellago, C., J. Chem. Theory Comput.  2019, 15 (5), 3075\sphinxhyphen{}3092.

\sphinxstepscope

\index{pair\_style hybrid@\spxentry{pair\_style hybrid}}\ignorespaces 
\index{pair\_style hybrid/kk@\spxentry{pair\_style hybrid/kk}}\ignorespaces 
\index{pair\_style hybrid/omp@\spxentry{pair\_style hybrid/omp}}\ignorespaces 
\index{pair\_style hybrid/molecular@\spxentry{pair\_style hybrid/molecular}}\ignorespaces 
\index{pair\_style hybrid/molecular/omp@\spxentry{pair\_style hybrid/molecular/omp}}\ignorespaces 
\index{pair\_style hybrid/overlay@\spxentry{pair\_style hybrid/overlay}}\ignorespaces 
\index{pair\_style hybrid/overlay/omp@\spxentry{pair\_style hybrid/overlay/omp}}\ignorespaces 
\index{pair\_style hybrid/overlay/kk@\spxentry{pair\_style hybrid/overlay/kk}}\ignorespaces 
\index{pair\_style hybrid/scaled@\spxentry{pair\_style hybrid/scaled}}\ignorespaces 
\index{pair\_style hybrid/scaled/omp@\spxentry{pair\_style hybrid/scaled/omp}}\ignorespaces 

\section{pair\_style hybrid command}
\label{\detokenize{pair_hybrid:pair-style-hybrid-command}}\label{\detokenize{pair_hybrid:index-9}}\label{\detokenize{pair_hybrid:index-8}}\label{\detokenize{pair_hybrid:index-7}}\label{\detokenize{pair_hybrid:index-6}}\label{\detokenize{pair_hybrid:index-5}}\label{\detokenize{pair_hybrid:index-4}}\label{\detokenize{pair_hybrid:index-3}}\label{\detokenize{pair_hybrid:index-2}}\label{\detokenize{pair_hybrid:index-1}}\label{\detokenize{pair_hybrid:index-0}}\label{\detokenize{pair_hybrid::doc}}
\sphinxAtStartPar
Accelerator Variants: \sphinxstyleemphasis{hybrid/kk}, \sphinxstyleemphasis{hybrid/omp}


\section{pair\_style hybrid/molecular command}
\label{\detokenize{pair_hybrid:pair-style-hybrid-molecular-command}}
\sphinxAtStartPar
Accelerator Variant: \sphinxstyleemphasis{hybrid/molecular/omp}


\section{pair\_style hybrid/overlay command}
\label{\detokenize{pair_hybrid:pair-style-hybrid-overlay-command}}
\sphinxAtStartPar
Accelerator Variants: \sphinxstyleemphasis{hybrid/overlay/kk}, \sphinxstyleemphasis{hybrid/overlay/omp}


\section{pair\_style hybrid/scaled command}
\label{\detokenize{pair_hybrid:pair-style-hybrid-scaled-command}}
\sphinxAtStartPar
Accelerator Variant: \sphinxstyleemphasis{hybrid/scaled/omp}


\subsection{Syntax}
\label{\detokenize{pair_hybrid:syntax}}
\begin{sphinxVerbatim}[commandchars=\\\{\}]
\PYG{k}{pair\PYGZus{}style}\PYG{+w}{ }\PYG{n}{hybrid}\PYG{+w}{ }\PYG{n}{style1}\PYG{+w}{ }\PYG{n}{args}\PYG{+w}{ }\PYG{n}{style2}\PYG{+w}{ }\PYG{n}{args}\PYG{+w}{ }\PYG{n}{...}
\PYG{k}{pair\PYGZus{}style}\PYG{+w}{ }\PYG{n}{hybrid}\PYG{o}{/}\PYG{n}{molecular}\PYG{+w}{ }\PYG{n}{factor1}\PYG{+w}{ }\PYG{n}{style1}\PYG{+w}{ }\PYG{n}{args}\PYG{+w}{ }\PYG{n}{factor2}\PYG{+w}{ }\PYG{n}{style}\PYG{+w}{ }\PYG{l+m}{2}\PYG{+w}{ }\PYG{n}{args}
\PYG{k}{pair\PYGZus{}style}\PYG{+w}{ }\PYG{n}{hybrid}\PYG{o}{/}\PYG{n}{overlay}\PYG{+w}{ }\PYG{n}{style1}\PYG{+w}{ }\PYG{n}{args}\PYG{+w}{ }\PYG{n}{style2}\PYG{+w}{ }\PYG{n}{args}\PYG{+w}{ }\PYG{n}{...}
\PYG{k}{pair\PYGZus{}style}\PYG{+w}{ }\PYG{n}{hybrid}\PYG{o}{/}\PYG{n}{scaled}\PYG{+w}{ }\PYG{n}{factor1}\PYG{+w}{ }\PYG{n}{style1}\PYG{+w}{ }\PYG{n}{args}\PYG{+w}{ }\PYG{n}{factor2}\PYG{+w}{ }\PYG{n}{style}\PYG{+w}{ }\PYG{l+m}{2}\PYG{+w}{ }\PYG{n}{args}\PYG{+w}{ }\PYG{n}{...}
\end{sphinxVerbatim}
\begin{itemize}
\item {} 
\sphinxAtStartPar
style1,style2 = list of one or more pair styles and their arguments

\item {} 
\sphinxAtStartPar
factor1,factor2 = scale factors for pair styles, may be a variable

\end{itemize}


\subsection{Examples}
\label{\detokenize{pair_hybrid:examples}}
\begin{sphinxVerbatim}[commandchars=\\\{\}]
\PYG{k}{pair\PYGZus{}style}\PYG{+w}{ }\PYG{n}{hybrid}\PYG{+w}{ }\PYG{n}{lj}\PYG{o}{/}\PYG{n}{cut}\PYG{o}{/}\PYG{n}{coul}\PYG{o}{/}\PYG{n}{cut}\PYG{+w}{ }\PYG{l+m}{10.0}\PYG{+w}{ }\PYG{n}{eam}\PYG{+w}{ }\PYG{n}{lj}\PYG{o}{/}\PYG{n}{cut}\PYG{+w}{ }\PYG{l+m}{5.0}
\PYG{k}{pair\PYGZus{}coeff}\PYG{+w}{ }\PYG{l+m}{1}\PYG{o}{*}\PYG{l+m}{2}\PYG{+w}{ }\PYG{l+m}{1}\PYG{o}{*}\PYG{l+m}{2}\PYG{+w}{ }\PYG{n}{eam}\PYG{+w}{ }\PYG{n}{niu3}
\PYG{k}{pair\PYGZus{}coeff}\PYG{+w}{ }\PYG{l+m}{3}\PYG{+w}{ }\PYG{l+m}{3}\PYG{+w}{ }\PYG{n}{lj}\PYG{o}{/}\PYG{n}{cut}\PYG{o}{/}\PYG{n}{coul}\PYG{o}{/}\PYG{n}{cut}\PYG{+w}{ }\PYG{l+m}{1.0}\PYG{+w}{ }\PYG{l+m}{1.0}
\PYG{k}{pair\PYGZus{}coeff}\PYG{+w}{ }\PYG{l+m}{1}\PYG{o}{*}\PYG{l+m}{2}\PYG{+w}{ }\PYG{l+m}{3}\PYG{+w}{ }\PYG{n}{lj}\PYG{o}{/}\PYG{n}{cut}\PYG{+w}{ }\PYG{l+m}{0.5}\PYG{+w}{ }\PYG{l+m}{1.2}

\PYG{k}{pair\PYGZus{}style}\PYG{+w}{ }\PYG{n}{hybrid}\PYG{o}{/}\PYG{n}{overlay}\PYG{+w}{ }\PYG{n}{lj}\PYG{o}{/}\PYG{n}{cut}\PYG{+w}{ }\PYG{l+m}{2.5}\PYG{+w}{ }\PYG{n}{coul}\PYG{o}{/}\PYG{n}{long}\PYG{+w}{ }\PYG{l+m}{2.0}
\PYG{k}{pair\PYGZus{}coeff}\PYG{+w}{ }\PYG{o}{*}\PYG{+w}{ }\PYG{o}{*}\PYG{+w}{ }\PYG{n}{lj}\PYG{o}{/}\PYG{n}{cut}\PYG{+w}{ }\PYG{l+m}{1.0}\PYG{+w}{ }\PYG{l+m}{1.0}
\PYG{k}{pair\PYGZus{}coeff}\PYG{+w}{ }\PYG{o}{*}\PYG{+w}{ }\PYG{o}{*}\PYG{+w}{ }\PYG{n}{coul}\PYG{o}{/}\PYG{n}{long}

\PYG{k}{pair\PYGZus{}style}\PYG{+w}{ }\PYG{n}{hybrid}\PYG{o}{/}\PYG{n}{scaled}\PYG{+w}{ }\PYG{l+m}{0.5}\PYG{+w}{ }\PYG{n}{tersoff}\PYG{+w}{ }\PYG{l+m}{0.5}\PYG{+w}{ }\PYG{n}{sw}
\PYG{k}{pair\PYGZus{}coeff}\PYG{+w}{ }\PYG{o}{*}\PYG{+w}{ }\PYG{o}{*}\PYG{+w}{ }\PYG{n}{tersoff}\PYG{+w}{ }\PYG{n}{Si.tersoff}\PYG{+w}{ }\PYG{n}{Si}
\PYG{k}{pair\PYGZus{}coeff}\PYG{+w}{ }\PYG{o}{*}\PYG{+w}{ }\PYG{o}{*}\PYG{+w}{ }\PYG{n}{sw}\PYG{+w}{ }\PYG{n}{Si.sw}\PYG{+w}{ }\PYG{n}{Si}

\PYG{k}{pair\PYGZus{}style}\PYG{+w}{ }\PYG{n}{hybrid}\PYG{o}{/}\PYG{n}{molecular}\PYG{+w}{ }\PYG{n}{lj}\PYG{o}{/}\PYG{n}{cut}\PYG{+w}{ }\PYG{l+m}{2.5}\PYG{+w}{ }\PYG{n}{lj}\PYG{o}{/}\PYG{n}{cut}\PYG{+w}{ }\PYG{l+m}{2.5}
\PYG{k}{pair\PYGZus{}coeff}\PYG{+w}{ }\PYG{o}{*}\PYG{+w}{ }\PYG{o}{*}\PYG{+w}{ }\PYG{n}{lj}\PYG{o}{/}\PYG{n}{cut}\PYG{+w}{ }\PYG{l+m}{1}\PYG{+w}{ }\PYG{l+m}{1.0}\PYG{+w}{ }\PYG{l+m}{1.0}
\PYG{k}{pair\PYGZus{}coeff}\PYG{+w}{ }\PYG{o}{*}\PYG{+w}{ }\PYG{o}{*}\PYG{+w}{ }\PYG{n}{lj}\PYG{o}{/}\PYG{n}{cut}\PYG{+w}{ }\PYG{l+m}{2}\PYG{+w}{ }\PYG{l+m}{1.5}\PYG{+w}{ }\PYG{l+m}{1.0}

\PYG{k}{variable }\PYG{n+nv+nvIdentifier}{one}\PYG{+w}{ }\PYG{n}{equal}\PYG{+w}{ }\PYG{n}{ramp}\PYG{n+nv}{(}\PYG{n+nv}{1.0,0.0}\PYG{n+nv}{)}
\PYG{k}{variable }\PYG{n+nv+nvIdentifier}{two}\PYG{+w}{ }\PYG{n}{equal}\PYG{+w}{ }\PYG{l+m}{1.0}\PYG{o}{\PYGZhy{}}\PYG{n}{v\PYGZus{}one}
\PYG{k}{pair\PYGZus{}style}\PYG{+w}{ }\PYG{n}{hybrid}\PYG{o}{/}\PYG{n}{scaled}\PYG{+w}{ }\PYG{n}{v\PYGZus{}one}\PYG{+w}{ }\PYG{n}{lj}\PYG{o}{/}\PYG{n}{cut}\PYG{+w}{ }\PYG{l+m}{2.5}\PYG{+w}{ }\PYG{n}{v\PYGZus{}two}\PYG{+w}{ }\PYG{n}{morse}\PYG{+w}{ }\PYG{l+m}{2.5}
\PYG{k}{pair\PYGZus{}coeff}\PYG{+w}{ }\PYG{l+m}{1}\PYG{+w}{ }\PYG{l+m}{1}\PYG{+w}{ }\PYG{n}{lj}\PYG{o}{/}\PYG{n}{cut}\PYG{+w}{ }\PYG{l+m}{1.0}\PYG{+w}{ }\PYG{l+m}{1.0}\PYG{+w}{ }\PYG{l+m}{2.5}
\PYG{k}{pair\PYGZus{}coeff}\PYG{+w}{ }\PYG{l+m}{1}\PYG{+w}{ }\PYG{l+m}{1}\PYG{+w}{ }\PYG{n}{morse}\PYG{+w}{ }\PYG{l+m}{1.0}\PYG{+w}{ }\PYG{l+m}{1.0}\PYG{+w}{ }\PYG{l+m}{1.0}\PYG{+w}{ }\PYG{l+m}{2.5}
\end{sphinxVerbatim}


\subsection{Description}
\label{\detokenize{pair_hybrid:description}}
\sphinxAtStartPar
The \sphinxstyleemphasis{hybrid}, \sphinxstyleemphasis{hybrid/overlay}, \sphinxstyleemphasis{hybrid/molecular}, and \sphinxstyleemphasis{hybrid/scaled}
styles enable the use of multiple pair styles in one simulation.  With
the \sphinxstyleemphasis{hybrid} style, exactly one pair style is assigned to each pair of
atom types.  With the \sphinxstyleemphasis{hybrid/overlay} and \sphinxstyleemphasis{hybrid/scaled} styles, one
or more pair styles can be assigned to each pair of atom types.  With
the hybrid/molecular style, pair styles are assigned to either intra\sphinxhyphen{}
or inter\sphinxhyphen{}molecular interactions.

\sphinxAtStartPar
The assignment of pair styles to type pairs is made via the
{\hyperref[\detokenize{pair_coeff::doc}]{\sphinxcrossref{\DUrole{doc}{pair\_coeff}}}} command.  The major difference between
the \sphinxstyleemphasis{hybrid/overlay} and \sphinxstyleemphasis{hybrid/scaled} styles is that the
\sphinxstyleemphasis{hybrid/scaled} adds a scale factor for each sub\sphinxhyphen{}style contribution to
forces, energies and stresses.  Because of the added complexity, the
\sphinxstyleemphasis{hybrid/scaled} style has more overhead and thus may be slower than
\sphinxstyleemphasis{hybrid/overlay}.

\sphinxAtStartPar
The \sphinxstyleemphasis{hybrid/molecular} pair style accepts \sphinxstyleemphasis{only} two sub\sphinxhyphen{}styles: the
first is assigned to intra\sphinxhyphen{}molecular interactions (i.e. both atoms
have the same molecule ID), the second to inter\sphinxhyphen{}molecular interactions
(i.e. interacting atoms have different molecule IDs).

\sphinxAtStartPar
Here are two examples of hybrid simulations.  The \sphinxstyleemphasis{hybrid} style could
be used for a simulation of a metal droplet on a LJ surface.  The metal
atoms interact with each other via an \sphinxstyleemphasis{eam} potential, the surface atoms
interact with each other via a \sphinxstyleemphasis{lj/cut} potential, and the metal/surface
interaction is also computed via a \sphinxstyleemphasis{lj/cut} potential.  The
\sphinxstyleemphasis{hybrid/overlay} style could be used as in the second example above,
where multiple potentials are superimposed in an additive fashion to
compute the interaction between atoms.  In this example, using \sphinxstyleemphasis{lj/cut}
and \sphinxstyleemphasis{coul/long} together gives the same result as if the
\sphinxstyleemphasis{lj/cut/coul/long} potential were used by itself.  In this case, it
would be more efficient to use the single combined potential, but in
general any combination of pair potentials can be used together in to
produce an interaction that is not encoded in any single pair\_style
file, e.g. adding Coulombic forces between granular particles.  Another
limitation of using the \sphinxstyleemphasis{hybrid/overlay} variant, that it does not generate
\sphinxstyleemphasis{lj/cut} parameters for mixed atom types from a mixing rule due to
restrictions discussed below.

\sphinxAtStartPar
If the \sphinxstyleemphasis{hybrid/scaled} style is used instead of \sphinxstyleemphasis{hybrid/overlay},
contributions from sub\sphinxhyphen{}styles are weighted by their scale factors, which
may be fractional or even negative.  Furthermore the scale factors may
be variables that may change during a simulation.  This enables
switching smoothly between two different pair styles or two different
parameter sets during a run in a similar fashion as could be done
with {\hyperref[\detokenize{fix_adapt::doc}]{\sphinxcrossref{\DUrole{doc}{fix adapt}}}} or {\hyperref[\detokenize{fix_alchemy::doc}]{\sphinxcrossref{\DUrole{doc}{fix alchemy}}}}.

\sphinxAtStartPar
All pair styles that will be used are listed as “sub\sphinxhyphen{}styles” following
the \sphinxstyleemphasis{hybrid} or \sphinxstyleemphasis{hybrid/overlay} keyword, in any order.  In case of the
\sphinxstyleemphasis{hybrid/scaled} pair style, each sub\sphinxhyphen{}style is prefixed with a scale
factor.  The scale factor is either a floating point number or an equal
style (or equivalent) variable.  Each sub\sphinxhyphen{}style’s name is followed by
its usual arguments, as illustrated in the examples above.  See the doc
pages of the individual pair styles for a listing and explanation of the
appropriate arguments for them.

\sphinxAtStartPar
Note that an individual pair style can be used multiple times as a
sub\sphinxhyphen{}style.  For efficiency reasons this should only be done if your
model requires it.  E.g. if you have different regions of Si and C atoms
and wish to use a Tersoff potential for pure Si for one set of atoms,
and a Tersoff potential for pure C for the other set (presumably with
some third potential for Si\sphinxhyphen{}C interactions), then the sub\sphinxhyphen{}style
\sphinxstyleemphasis{tersoff} could be listed twice.  But if you just want to use a
Lennard\sphinxhyphen{}Jones or other pairwise potential for several different atom
type pairs in your model, then you should just list the sub\sphinxhyphen{}style once
and use the pair\_coeff command to assign parameters for the different
type pairs.

\begin{sphinxadmonition}{note}{Note:}
\sphinxAtStartPar
There is one exception to this option to list an individual
pair style multiple times: GPU\sphinxhyphen{}enabled pair styles in the GPU package.
This is because the GPU package currently assumes that only one
instance of a pair style is being used.
\end{sphinxadmonition}

\sphinxAtStartPar
In the pair\_coeff commands, the name of a pair style must be added
after the I,J type specification, with the remaining coefficients
being those appropriate to that style.  If the pair style is used
multiple times in the pair\_style command, then an additional numeric
argument must also be specified which is a number from 1 to M where M
is the number of times the sub\sphinxhyphen{}style was listed in the pair style
command.  The extra number indicates which instance of the sub\sphinxhyphen{}style
these coefficients apply to.

\sphinxAtStartPar
For example, consider a simulation with 3 atom types: types 1 and 2
are Ni atoms, type 3 are LJ atoms with charges.  The following
commands would set up a hybrid simulation:

\begin{sphinxVerbatim}[commandchars=\\\{\}]
\PYG{k}{pair\PYGZus{}style}\PYG{+w}{ }\PYG{n}{hybrid}\PYG{+w}{ }\PYG{n}{eam}\PYG{o}{/}\PYG{n}{alloy}\PYG{+w}{ }\PYG{n}{lj}\PYG{o}{/}\PYG{n}{cut}\PYG{o}{/}\PYG{n}{coul}\PYG{o}{/}\PYG{n}{cut}\PYG{+w}{ }\PYG{l+m}{10.0}\PYG{+w}{ }\PYG{n}{lj}\PYG{o}{/}\PYG{n}{cut}\PYG{+w}{ }\PYG{l+m}{8.0}
\PYG{k}{pair\PYGZus{}coeff}\PYG{+w}{ }\PYG{o}{*}\PYG{+w}{ }\PYG{o}{*}\PYG{+w}{ }\PYG{n}{eam}\PYG{o}{/}\PYG{n}{alloy}\PYG{+w}{ }\PYG{n}{nialhjea}\PYG{+w}{ }\PYG{n}{Ni}\PYG{+w}{ }\PYG{n}{Ni}\PYG{+w}{ }\PYG{n}{NULL}
\PYG{k}{pair\PYGZus{}coeff}\PYG{+w}{ }\PYG{l+m}{3}\PYG{+w}{ }\PYG{l+m}{3}\PYG{+w}{ }\PYG{n}{lj}\PYG{o}{/}\PYG{n}{cut}\PYG{o}{/}\PYG{n}{coul}\PYG{o}{/}\PYG{n}{cut}\PYG{+w}{ }\PYG{l+m}{1.0}\PYG{+w}{ }\PYG{l+m}{1.0}
\PYG{k}{pair\PYGZus{}coeff}\PYG{+w}{ }\PYG{l+m}{1}\PYG{o}{*}\PYG{l+m}{2}\PYG{+w}{ }\PYG{l+m}{3}\PYG{+w}{ }\PYG{n}{lj}\PYG{o}{/}\PYG{n}{cut}\PYG{+w}{ }\PYG{l+m}{0.8}\PYG{+w}{ }\PYG{l+m}{1.3}
\end{sphinxVerbatim}

\sphinxAtStartPar
As an example of using the same pair style multiple times, consider a
simulation with 2 atom types.  Type 1 is Si, type 2 is C.  The
following commands would model the Si atoms with Tersoff, the C atoms
with Tersoff, and the cross\sphinxhyphen{}interactions with Lennard\sphinxhyphen{}Jones:

\begin{sphinxVerbatim}[commandchars=\\\{\}]
\PYG{k}{pair\PYGZus{}style}\PYG{+w}{ }\PYG{n}{hybrid}\PYG{+w}{ }\PYG{n}{lj}\PYG{o}{/}\PYG{n}{cut}\PYG{+w}{ }\PYG{l+m}{2.5}\PYG{+w}{ }\PYG{n}{tersoff}\PYG{+w}{ }\PYG{n}{tersoff}
\PYG{k}{pair\PYGZus{}coeff}\PYG{+w}{ }\PYG{o}{*}\PYG{+w}{ }\PYG{o}{*}\PYG{+w}{ }\PYG{n}{tersoff}\PYG{+w}{ }\PYG{l+m}{1}\PYG{+w}{ }\PYG{n}{Si.tersoff}\PYG{+w}{ }\PYG{n}{Si}\PYG{+w}{ }\PYG{n}{NULL}
\PYG{k}{pair\PYGZus{}coeff}\PYG{+w}{ }\PYG{o}{*}\PYG{+w}{ }\PYG{o}{*}\PYG{+w}{ }\PYG{n}{tersoff}\PYG{+w}{ }\PYG{l+m}{2}\PYG{+w}{ }\PYG{n}{C.tersoff}\PYG{+w}{ }\PYG{n}{NULL}\PYG{+w}{ }\PYG{n}{C}
\PYG{k}{pair\PYGZus{}coeff}\PYG{+w}{ }\PYG{l+m}{1}\PYG{+w}{ }\PYG{l+m}{2}\PYG{+w}{ }\PYG{n}{lj}\PYG{o}{/}\PYG{n}{cut}\PYG{+w}{ }\PYG{l+m}{1.0}\PYG{+w}{ }\PYG{l+m}{1.5}
\end{sphinxVerbatim}

\sphinxAtStartPar
It is not recommended to read pair coefficients for a hybrid style from a “Pair Coeffs”
or “PairIJ Coeffs” section of a data file via the {\hyperref[\detokenize{read_data::doc}]{\sphinxcrossref{\DUrole{doc}{read\_data}}}} command,
since those sections expect a fixed number of lines, either one line per atom type or
one line pair pair of atom types, respectively.  When reading from a data file, the
lines of the “Pair Coeffs” and “PairIJ Coeffs” are changed in the same way as the \sphinxstyleemphasis{pair\_coeff}
command, i.e. the name of the pair style to which the parameters apply must follow the
atom type (or atom types), e.g.

\begin{sphinxVerbatim}[commandchars=\\\{\}]
Pair Coeffs

1 lj/cut/coul/cut 1.0 1.0
...

PairIJ Coeffs

1 1 lj/cut/coul/cut 1.0 1.0
...
\end{sphinxVerbatim}

\sphinxAtStartPar
Note that the pair\_coeff command for some potentials such as
{\hyperref[\detokenize{pair_eam::doc}]{\sphinxcrossref{\DUrole{doc}{pair\_style eam/alloy}}}} includes a mapping specification
of elements to all atom types, which in the hybrid case, can include
atom types not assigned to the \sphinxstyleemphasis{eam/alloy} potential.  The NULL
keyword is used by many such potentials (eam/alloy, Tersoff, AIREBO,
etc), to denote an atom type that will be assigned to a different
sub\sphinxhyphen{}style.

\sphinxAtStartPar
For the \sphinxstyleemphasis{hybrid} style, each atom type pair I,J is assigned to exactly
one sub\sphinxhyphen{}style.  Just as with a simulation using a single pair style,
if you specify the same atom type pair in a second pair\_coeff command,
the previous assignment will be overwritten.

\sphinxAtStartPar
For the \sphinxstyleemphasis{hybrid/overlay} and \sphinxstyleemphasis{hybrid/scaled} styles, each atom type pair
I,J can be assigned to one or more sub\sphinxhyphen{}styles.  If you specify the same
atom type pair in a second pair\_coeff command with a new sub\sphinxhyphen{}style, then
the second sub\sphinxhyphen{}style is added to the list of potentials that will be
calculated for two interacting atoms of those types.  If you specify the
same atom type pair in a second pair\_coeff command with a sub\sphinxhyphen{}style that
has already been defined for that pair of atoms, then the new pair
coefficients simply override the previous ones, as in the normal usage
of the pair\_coeff command.  E.g. these two sets of commands are the
same:

\begin{sphinxVerbatim}[commandchars=\\\{\}]
\PYG{k}{pair\PYGZus{}style}\PYG{+w}{ }\PYG{n}{lj}\PYG{o}{/}\PYG{n}{cut}\PYG{+w}{ }\PYG{l+m}{2.5}
\PYG{k}{pair\PYGZus{}coeff}\PYG{+w}{ }\PYG{o}{*}\PYG{+w}{ }\PYG{o}{*}\PYG{+w}{ }\PYG{l+m}{1.0}\PYG{+w}{ }\PYG{l+m}{1.0}
\PYG{k}{pair\PYGZus{}coeff}\PYG{+w}{ }\PYG{l+m}{2}\PYG{+w}{ }\PYG{l+m}{2}\PYG{+w}{ }\PYG{l+m}{1.5}\PYG{+w}{ }\PYG{l+m}{0.8}

\PYG{k}{pair\PYGZus{}style}\PYG{+w}{ }\PYG{n}{hybrid}\PYG{o}{/}\PYG{n}{overlay}\PYG{+w}{ }\PYG{n}{lj}\PYG{o}{/}\PYG{n}{cut}\PYG{+w}{ }\PYG{l+m}{2.5}
\PYG{k}{pair\PYGZus{}coeff}\PYG{+w}{ }\PYG{o}{*}\PYG{+w}{ }\PYG{o}{*}\PYG{+w}{ }\PYG{n}{lj}\PYG{o}{/}\PYG{n}{cut}\PYG{+w}{ }\PYG{l+m}{1.0}\PYG{+w}{ }\PYG{l+m}{1.0}
\PYG{k}{pair\PYGZus{}coeff}\PYG{+w}{ }\PYG{l+m}{2}\PYG{+w}{ }\PYG{l+m}{2}\PYG{+w}{ }\PYG{n}{lj}\PYG{o}{/}\PYG{n}{cut}\PYG{+w}{ }\PYG{l+m}{1.5}\PYG{+w}{ }\PYG{l+m}{0.8}
\end{sphinxVerbatim}

\sphinxAtStartPar
Coefficients must be defined for each pair of atoms types via the
{\hyperref[\detokenize{pair_coeff::doc}]{\sphinxcrossref{\DUrole{doc}{pair\_coeff}}}} command as described above, or in the
“Pair Coeffs” or “PairIJ Coeffs” section of the data file read by the
{\hyperref[\detokenize{read_data::doc}]{\sphinxcrossref{\DUrole{doc}{read\_data}}}} command, or by mixing as described below.

\sphinxAtStartPar
For all of the \sphinxstyleemphasis{hybrid}, \sphinxstyleemphasis{hybrid/overlay}, and \sphinxstyleemphasis{hybrid/scaled} styles,
every atom type pair I,J (where I \textless{}= J) must be assigned to at least one
sub\sphinxhyphen{}style via the {\hyperref[\detokenize{pair_coeff::doc}]{\sphinxcrossref{\DUrole{doc}{pair\_coeff}}}} command as in the
examples above, or in the data file read by the {\hyperref[\detokenize{read_data::doc}]{\sphinxcrossref{\DUrole{doc}{read\_data}}}}, or by mixing as described below.  Also all sub\sphinxhyphen{}styles
must be used at least once in a {\hyperref[\detokenize{pair_coeff::doc}]{\sphinxcrossref{\DUrole{doc}{pair\_coeff}}}} command.

\begin{sphinxadmonition}{warning}{Warning:}
\sphinxAtStartPar
With hybrid pair styles the use of mixing to generate pair
coefficients is significantly limited compared to the individual pair
styles.  LAMMPS \sphinxstylestrong{never} performs mixing of parameters from
different sub\sphinxhyphen{}styles, \sphinxstylestrong{even} if they use the same type of
coefficients, e.g. contain a Lennard\sphinxhyphen{}Jones potential variant.  Those
parameters must be provided explicitly.  Also for \sphinxstyleemphasis{hybrid/overlay}
and \sphinxstyleemphasis{hybrid/scaled} mixing is \sphinxstylestrong{only} performed for pairs of atom
types for which only a single pair style is assigned.

\sphinxAtStartPar
Thus it is strongly recommended to provide all mixed terms
explicitly.  For non\sphinxhyphen{}hybrid styles those could be generated and
written out using the {\hyperref[\detokenize{write_coeff::doc}]{\sphinxcrossref{\DUrole{doc}{write\_coeff command}}}} and
then edited as needed to comply with the requirements for hybrid
styles as explained above.
\end{sphinxadmonition}

\sphinxAtStartPar
If you want there to be no interactions between a particular pair of
atom types, you have 3 choices.  You can assign the pair of atom types
to some sub\sphinxhyphen{}style and use the {\hyperref[\detokenize{neigh_modify::doc}]{\sphinxcrossref{\DUrole{doc}{neigh\_modify exclude type}}}}
command.  You can assign it to some sub\sphinxhyphen{}style and set the coefficients
so that there is effectively no interaction (e.g. epsilon = 0.0 in a LJ
potential).  Or, for \sphinxstyleemphasis{hybrid}, \sphinxstyleemphasis{hybrid/overlay}, or \sphinxstyleemphasis{hybrid/scaled}
simulations, you can use this form of the pair\_coeff command in your
input script or the “PairIJ Coeffs” section of your data file:

\begin{sphinxVerbatim}[commandchars=\\\{\}]
\PYG{k}{pair\PYGZus{}coeff}\PYG{+w}{  }\PYG{l+m}{2}\PYG{+w}{ }\PYG{l+m}{3}\PYG{+w}{  }\PYG{n}{none}
\end{sphinxVerbatim}

\sphinxAtStartPar
or this form in the “Pair Coeffs” section of the data file:

\begin{sphinxVerbatim}[commandchars=\\\{\}]
3  none
\end{sphinxVerbatim}

\sphinxAtStartPar
If an assignment to \sphinxstyleemphasis{none} is made in a simulation with the
\sphinxstyleemphasis{hybrid/overlay} or \sphinxstyleemphasis{hybrid/scaled} pair style, it wipes out all
previous assignments of that pair of atom types to sub\sphinxhyphen{}styles.

\sphinxAtStartPar
Note that you may need to use an {\hyperref[\detokenize{atom_style::doc}]{\sphinxcrossref{\DUrole{doc}{atom\_style}}}} hybrid
command in your input script, if atoms in the simulation will need
attributes from several atom styles, due to using multiple pair
styles with different requirements.


\bigskip\hrule\bigskip


\sphinxAtStartPar
Different force fields (e.g. CHARMM vs. AMBER) may have different rules
for applying exclusions or weights that change the strength of pairwise
non\sphinxhyphen{}bonded interactions between pairs of atoms that are also 1\sphinxhyphen{}2, 1\sphinxhyphen{}3,
and 1\sphinxhyphen{}4 neighbors in the molecular bond topology. This is normally a
global setting defined the {\hyperref[\detokenize{special_bonds::doc}]{\sphinxcrossref{\DUrole{doc}{special\_bonds}}}} command.
However, different weights can be assigned to different hybrid
sub\sphinxhyphen{}styles via the {\hyperref[\detokenize{pair_modify::doc}]{\sphinxcrossref{\DUrole{doc}{pair\_modify special}}}} command.
This allows multiple force fields to be used in a model of a hybrid
system, however, there is no consistent approach to determine parameters
automatically for the interactions \sphinxstylestrong{between} atoms of the two force
fields, thus this approach this is only recommended when particles
described by the different force fields do not mix.

\sphinxAtStartPar
Here is an example for combining CHARMM and AMBER: The global \sphinxstyleemphasis{amber}
setting sets the 1\sphinxhyphen{}4 interactions to non\sphinxhyphen{}zero scaling factors and
then overrides them with 0.0 only for CHARMM:

\begin{sphinxVerbatim}[commandchars=\\\{\}]
\PYG{k}{special\PYGZus{}bonds}\PYG{+w}{ }\PYG{n}{amber}
\PYG{k}{pair\PYGZus{}style}\PYG{+w}{ }\PYG{n}{hybrid}\PYG{+w}{ }\PYG{n}{lj}\PYG{o}{/}\PYG{n}{charmm}\PYG{o}{/}\PYG{n}{coul}\PYG{o}{/}\PYG{n}{long}\PYG{+w}{ }\PYG{l+m}{8.0}\PYG{+w}{ }\PYG{l+m}{10.0}\PYG{+w}{ }\PYG{n}{lj}\PYG{o}{/}\PYG{n}{cut}\PYG{o}{/}\PYG{n}{coul}\PYG{o}{/}\PYG{n}{long}\PYG{+w}{ }\PYG{l+m}{10.0}
\PYG{k}{pair\PYGZus{}modify}\PYG{+w}{ }\PYG{n}{pair}\PYG{+w}{ }\PYG{n}{lj}\PYG{o}{/}\PYG{n}{charmm}\PYG{o}{/}\PYG{n}{coul}\PYG{o}{/}\PYG{n}{long}\PYG{+w}{ }\PYG{n}{special}\PYG{+w}{ }\PYG{n}{lj}\PYG{o}{/}\PYG{n}{coul}\PYG{+w}{ }\PYG{l+m}{0.0}\PYG{+w}{ }\PYG{l+m}{0.0}\PYG{+w}{ }\PYG{l+m}{0.0}
\end{sphinxVerbatim}

\sphinxAtStartPar
This input achieves the same effect:

\begin{sphinxVerbatim}[commandchars=\\\{\}]
\PYG{k}{special\PYGZus{}bonds}\PYG{+w}{ }\PYG{l+m}{0.0}\PYG{+w}{ }\PYG{l+m}{0.0}\PYG{+w}{ }\PYG{l+m}{0.1}
\PYG{k}{pair\PYGZus{}style}\PYG{+w}{ }\PYG{n}{hybrid}\PYG{+w}{ }\PYG{n}{lj}\PYG{o}{/}\PYG{n}{charmm}\PYG{o}{/}\PYG{n}{coul}\PYG{o}{/}\PYG{n}{long}\PYG{+w}{ }\PYG{l+m}{8.0}\PYG{+w}{ }\PYG{l+m}{10.0}\PYG{+w}{ }\PYG{n}{lj}\PYG{o}{/}\PYG{n}{cut}\PYG{o}{/}\PYG{n}{coul}\PYG{o}{/}\PYG{n}{long}\PYG{+w}{ }\PYG{l+m}{10.0}
\PYG{k}{pair\PYGZus{}modify}\PYG{+w}{ }\PYG{n}{pair}\PYG{+w}{ }\PYG{n}{lj}\PYG{o}{/}\PYG{n}{cut}\PYG{o}{/}\PYG{n}{coul}\PYG{o}{/}\PYG{n}{long}\PYG{+w}{ }\PYG{n}{special}\PYG{+w}{ }\PYG{n}{lj}\PYG{+w}{ }\PYG{l+m}{0.0}\PYG{+w}{ }\PYG{l+m}{0.0}\PYG{+w}{ }\PYG{l+m}{0.5}
\PYG{k}{pair\PYGZus{}modify}\PYG{+w}{ }\PYG{n}{pair}\PYG{+w}{ }\PYG{n}{lj}\PYG{o}{/}\PYG{n}{cut}\PYG{o}{/}\PYG{n}{coul}\PYG{o}{/}\PYG{n}{long}\PYG{+w}{ }\PYG{n}{special}\PYG{+w}{ }\PYG{n}{coul}\PYG{+w}{ }\PYG{l+m}{0.0}\PYG{+w}{ }\PYG{l+m}{0.0}\PYG{+w}{ }\PYG{l+m}{0.83333333}
\PYG{k}{pair\PYGZus{}modify}\PYG{+w}{ }\PYG{n}{pair}\PYG{+w}{ }\PYG{n}{lj}\PYG{o}{/}\PYG{n}{charmm}\PYG{o}{/}\PYG{n}{coul}\PYG{o}{/}\PYG{n}{long}\PYG{+w}{ }\PYG{n}{special}\PYG{+w}{ }\PYG{n}{lj}\PYG{o}{/}\PYG{n}{coul}\PYG{+w}{ }\PYG{l+m}{0.0}\PYG{+w}{ }\PYG{l+m}{0.0}\PYG{+w}{ }\PYG{l+m}{0.0}
\end{sphinxVerbatim}

\sphinxAtStartPar
Here is an example for combining Tersoff with OPLS/AA based on
a data file that defines bonds for all atoms where \sphinxhyphen{} for the
Tersoff part of the system \sphinxhyphen{} the force constants for the bonded
interactions have been set to 0. Note the global settings are
effectively \sphinxstyleemphasis{lj/coul 0.0 0.0 0.5} as required for OPLS/AA:

\begin{sphinxVerbatim}[commandchars=\\\{\}]
\PYG{k}{special\PYGZus{}bonds}\PYG{+w}{ }\PYG{n}{lj}\PYG{o}{/}\PYG{n}{coul}\PYG{+w}{ }\PYG{l+m}{1e\PYGZhy{}20}\PYG{+w}{ }\PYG{l+m}{1e\PYGZhy{}20}\PYG{+w}{ }\PYG{l+m}{0.5}
\PYG{k}{pair\PYGZus{}style}\PYG{+w}{ }\PYG{n}{hybrid}\PYG{+w}{ }\PYG{n}{tersoff}\PYG{+w}{ }\PYG{n}{lj}\PYG{o}{/}\PYG{n}{cut}\PYG{o}{/}\PYG{n}{coul}\PYG{o}{/}\PYG{n}{long}\PYG{+w}{ }\PYG{l+m}{12.0}
\PYG{k}{pair\PYGZus{}modify}\PYG{+w}{ }\PYG{n}{pair}\PYG{+w}{ }\PYG{n}{tersoff}\PYG{+w}{ }\PYG{n}{special}\PYG{+w}{ }\PYG{n}{lj}\PYG{o}{/}\PYG{n}{coul}\PYG{+w}{ }\PYG{l+m}{1.0}\PYG{+w}{ }\PYG{l+m}{1.0}\PYG{+w}{ }\PYG{l+m}{1.0}
\end{sphinxVerbatim}

\sphinxAtStartPar
For use with the various {\hyperref[\detokenize{compute_tally::doc}]{\sphinxcrossref{\DUrole{doc}{compute */tally}}}}
computes, the {\hyperref[\detokenize{pair_modify::doc}]{\sphinxcrossref{\DUrole{doc}{pair\_modify compute/tally}}}}
command can be used to selectively turn off processing of
the compute tally styles, for example, if those pair styles
(e.g. many\sphinxhyphen{}body styles) do not support this feature.

\sphinxAtStartPar
See the {\hyperref[\detokenize{pair_modify::doc}]{\sphinxcrossref{\DUrole{doc}{pair\_modify}}}} page for details on
the specific syntax, requirements and restrictions.


\bigskip\hrule\bigskip


\sphinxAtStartPar
The potential energy contribution to the overall system due to an
individual sub\sphinxhyphen{}style can be accessed and output via the {\hyperref[\detokenize{compute_pair::doc}]{\sphinxcrossref{\DUrole{doc}{compute
pair}}}} command.  Note that in the case of pair style
\sphinxstyleemphasis{hybrid/scaled} this is the \sphinxstylestrong{unscaled} potential energy of the
selected sub\sphinxhyphen{}style.


\bigskip\hrule\bigskip


\begin{sphinxadmonition}{note}{Note:}
\sphinxAtStartPar
Several of the potentials defined via the pair\_style command in
LAMMPS are really many\sphinxhyphen{}body potentials, such as Tersoff, AIREBO, MEAM,
ReaxFF, etc.  The way to think about using these potentials in a
hybrid setting is as follows.
\end{sphinxadmonition}

\sphinxAtStartPar
A subset of atom types is assigned to the many\sphinxhyphen{}body potential with a
single {\hyperref[\detokenize{pair_coeff::doc}]{\sphinxcrossref{\DUrole{doc}{pair\_coeff}}}} command, using “* *” to include
all types and the NULL keywords described above to exclude specific
types not assigned to that potential.  If types 1,3,4 were assigned in
that way (but not type 2), this means that all many\sphinxhyphen{}body interactions
between all atoms of types 1,3,4 will be computed by that potential.
Pair\_style hybrid allows interactions between type pairs 2\sphinxhyphen{}2, 1\sphinxhyphen{}2,
2\sphinxhyphen{}3, 2\sphinxhyphen{}4 to be specified for computation by other pair styles.  You
could even add a second interaction for 1\sphinxhyphen{}1 to be computed by another
pair style, assuming pair\_style hybrid/overlay is used.

\sphinxAtStartPar
But you should not, as a general rule, attempt to exclude the many\sphinxhyphen{}body
interactions for some subset of the type pairs within the set of 1,3,4
interactions, e.g. exclude 1\sphinxhyphen{}1 or 1\sphinxhyphen{}3 interactions.  That is not
conceptually well\sphinxhyphen{}defined for many\sphinxhyphen{}body interactions, since the
potential will typically calculate energies and foces for small groups
of atoms, e.g. 3 or 4 atoms, using the neighbor lists of the atoms to
find the additional atoms in the group.

\sphinxAtStartPar
However, you can still use the pair\_coeff none setting or the
{\hyperref[\detokenize{neigh_modify::doc}]{\sphinxcrossref{\DUrole{doc}{neigh\_modify exclude}}}} command to exclude certain
type pairs from the neighbor list that will be passed to a many\sphinxhyphen{}body
sub\sphinxhyphen{}style.  This will alter the calculations made by a many\sphinxhyphen{}body
potential beyond the specific pairs, since it builds its list of 3\sphinxhyphen{}body,
4\sphinxhyphen{}body, etc interactions from the pair lists.  You will need to think
\sphinxstylestrong{carefully} as to whether excluding such pairs produces a physically
meaningful result for your model.

\sphinxAtStartPar
For example, imagine you have two atom types in your model, type 1 for
atoms in one surface, and type 2 for atoms in the other, and you wish
to use a Tersoff potential to compute interactions within each
surface, but not between the surfaces.  Then either of these two command
sequences would implement that model:

\begin{sphinxVerbatim}[commandchars=\\\{\}]
\PYG{k}{pair\PYGZus{}style}\PYG{+w}{ }\PYG{n}{hybrid}\PYG{+w}{ }\PYG{n}{tersoff}
\PYG{k}{pair\PYGZus{}coeff}\PYG{+w}{ }\PYG{o}{*}\PYG{+w}{ }\PYG{o}{*}\PYG{+w}{ }\PYG{n}{tersoff}\PYG{+w}{ }\PYG{n}{SiC.tersoff}\PYG{+w}{ }\PYG{n}{C}\PYG{+w}{ }\PYG{n}{C}
\PYG{k}{pair\PYGZus{}coeff}\PYG{+w}{ }\PYG{l+m}{1}\PYG{+w}{ }\PYG{l+m}{2}\PYG{+w}{ }\PYG{n}{none}

\PYG{k}{pair\PYGZus{}style}\PYG{+w}{ }\PYG{n}{tersoff}
\PYG{k}{pair\PYGZus{}coeff}\PYG{+w}{ }\PYG{o}{*}\PYG{+w}{ }\PYG{o}{*}\PYG{+w}{ }\PYG{n}{SiC.tersoff}\PYG{+w}{ }\PYG{n}{C}\PYG{+w}{ }\PYG{n}{C}
\PYG{k}{neigh\PYGZus{}modify}\PYG{+w}{ }\PYG{n}{exclude}\PYG{+w}{ }\PYG{n}{type}\PYG{+w}{ }\PYG{l+m}{1}\PYG{+w}{ }\PYG{l+m}{2}
\end{sphinxVerbatim}

\sphinxAtStartPar
Either way, only neighbor lists with 1\sphinxhyphen{}1 or 2\sphinxhyphen{}2 interactions would be
passed to the Tersoff potential, which means it would compute no
3\sphinxhyphen{}body interactions containing both type 1 and 2 atoms.

\sphinxAtStartPar
Here is another example to use 2 many\sphinxhyphen{}body potentials together in an
overlapping manner using hybrid/overlay.  Imagine you have CNT (C atoms)
on a Si surface.  You want to use Tersoff for Si/Si and Si/C
interactions, and AIREBO for C/C interactions.  Si atoms are type 1; C
atoms are type 2.  Something like this will work:

\begin{sphinxVerbatim}[commandchars=\\\{\}]
\PYG{k}{pair\PYGZus{}style}\PYG{+w}{ }\PYG{n}{hybrid}\PYG{o}{/}\PYG{n}{overlay}\PYG{+w}{ }\PYG{n}{tersoff}\PYG{+w}{ }\PYG{n}{airebo}\PYG{+w}{ }\PYG{l+m}{3.0}
\PYG{k}{pair\PYGZus{}coeff}\PYG{+w}{ }\PYG{o}{*}\PYG{+w}{ }\PYG{o}{*}\PYG{+w}{ }\PYG{n}{tersoff}\PYG{+w}{ }\PYG{n}{SiC.tersoff.custom}\PYG{+w}{ }\PYG{n}{Si}\PYG{+w}{ }\PYG{n}{C}
\PYG{k}{pair\PYGZus{}coeff}\PYG{+w}{ }\PYG{o}{*}\PYG{+w}{ }\PYG{o}{*}\PYG{+w}{ }\PYG{n}{airebo}\PYG{+w}{ }\PYG{n}{CH.airebo}\PYG{+w}{ }\PYG{n}{NULL}\PYG{+w}{ }\PYG{n}{C}
\end{sphinxVerbatim}

\sphinxAtStartPar
Note that to prevent the Tersoff potential from computing C/C
interactions, you would need to \sphinxstylestrong{modify} the SiC.tersoff potential
file to turn off C/C interaction, i.e. by setting the appropriate
coefficients to 0.0.


\bigskip\hrule\bigskip


\sphinxAtStartPar
Styles with a \sphinxstyleemphasis{gpu}, \sphinxstyleemphasis{intel}, \sphinxstyleemphasis{kk}, \sphinxstyleemphasis{omp}, or \sphinxstyleemphasis{opt} suffix are
functionally the same as the corresponding style without the suffix.
They have been optimized to run faster, depending on your available
hardware, as discussed on the {\hyperref[\detokenize{Speed_packages::doc}]{\sphinxcrossref{\DUrole{doc}{Accelerator packages}}}}
page.  The accelerated styles take the same arguments and should
produce the same results, except for round\sphinxhyphen{}off and precision issues.

\sphinxAtStartPar
These accelerated styles are part of the GPU, INTEL, KOKKOS,
OPENMP, and OPT packages, respectively.  They are only enabled if
LAMMPS was built with those packages.  See the {\hyperref[\detokenize{Build_package::doc}]{\sphinxcrossref{\DUrole{doc}{Build package}}}} page for more info.

\sphinxAtStartPar
You can specify the accelerated styles explicitly in your input script
by including their suffix, or you can use the {\hyperref[\detokenize{Run_options::doc}]{\sphinxcrossref{\DUrole{doc}{\sphinxhyphen{}suffix command\sphinxhyphen{}line switch}}}} when you invoke LAMMPS, or you can use the
{\hyperref[\detokenize{suffix::doc}]{\sphinxcrossref{\DUrole{doc}{suffix}}}} command in your input script.

\sphinxAtStartPar
See the {\hyperref[\detokenize{Speed_packages::doc}]{\sphinxcrossref{\DUrole{doc}{Accelerator packages}}}} page for more
instructions on how to use the accelerated styles effectively.

\begin{sphinxadmonition}{note}{Note:}
\sphinxAtStartPar
Since the \sphinxstyleemphasis{hybrid}, \sphinxstyleemphasis{hybrid/overlay}, \sphinxstyleemphasis{hybrid/scaled} styles
delegate computation to the individual sub\sphinxhyphen{}styles, the suffix
versions of the \sphinxstyleemphasis{hybrid} and \sphinxstyleemphasis{hybrid/overlay} styles are used to
propagate the corresponding suffix to all sub\sphinxhyphen{}styles, if those
versions exist. Otherwise the non\sphinxhyphen{}accelerated version will be used.
The individual accelerated sub\sphinxhyphen{}styles are part of the GPU, KOKKOS,
INTEL, OPENMP, and OPT packages, respectively.  They are only
enabled if LAMMPS was built with those packages.  See the
{\hyperref[\detokenize{Build_package::doc}]{\sphinxcrossref{\DUrole{doc}{Build package}}}} page for more info.
\end{sphinxadmonition}


\bigskip\hrule\bigskip



\subsection{Mixing, shift, table, tail correction, restart, rRESPA info}
\label{\detokenize{pair_hybrid:mixing-shift-table-tail-correction-restart-rrespa-info}}
\sphinxAtStartPar
Any pair potential settings made via the
{\hyperref[\detokenize{pair_modify::doc}]{\sphinxcrossref{\DUrole{doc}{pair\_modify}}}} command are passed along to all
sub\sphinxhyphen{}styles of the hybrid potential.

\sphinxAtStartPar
For atom type pairs I,J and I != J, if the sub\sphinxhyphen{}style assigned to I,I and
J,J is the same, and if the sub\sphinxhyphen{}style allows for mixing, then the
coefficients for I,J can be mixed.  This means you do not have to
specify a pair\_coeff command for I,J since the I,J type pair will be
assigned automatically to the sub\sphinxhyphen{}style defined for both I,I and J,J and
its coefficients generated by the mixing rule used by that sub\sphinxhyphen{}style.
For the \sphinxstyleemphasis{hybrid/overlay} and \sphinxstyleemphasis{hybrid/scaled} style, there is an
additional requirement that both the I,I and J,J pairs are assigned to a
single sub\sphinxhyphen{}style.  If this requirement is not met, no I,J coeffs will be
generated, even if the sub\sphinxhyphen{}styles support mixing, and I,J pair
coefficients must be explicitly defined.

\sphinxAtStartPar
See the {\hyperref[\detokenize{pair_modify::doc}]{\sphinxcrossref{\DUrole{doc}{pair\_modify}}}} command for
details of mixing rules.  See the See the page for the sub\sphinxhyphen{}style to
see if allows for mixing.

\sphinxAtStartPar
The hybrid pair styles supports the {\hyperref[\detokenize{pair_modify::doc}]{\sphinxcrossref{\DUrole{doc}{pair\_modify}}}}
shift, table, and tail options for an I,J pair interaction, if the
associated sub\sphinxhyphen{}style supports it.

\sphinxAtStartPar
For the hybrid pair styles, the list of sub\sphinxhyphen{}styles and their respective
settings are written to {\hyperref[\detokenize{restart::doc}]{\sphinxcrossref{\DUrole{doc}{binary restart files}}}}, so a
{\hyperref[\detokenize{pair_style::doc}]{\sphinxcrossref{\DUrole{doc}{pair\_style}}}} command does not need to specified in an
input script that reads a restart file.  However, the coefficient
information is not stored in the restart file.  The same is true for
{\hyperref[\detokenize{write_data::doc}]{\sphinxcrossref{\DUrole{doc}{data files}}}}.  Thus, pair\_coeff commands need to be
re\sphinxhyphen{}specified in the restart input script.  For pair style
\sphinxstyleemphasis{hybrid/scaled} also the names of any variables used as scale factors
are restored, but not the variables themselves, so those may need to be
redefined when continuing from a restart.

\sphinxAtStartPar
These pair styles support the use of the \sphinxstyleemphasis{inner}, \sphinxstyleemphasis{middle}, and
\sphinxstyleemphasis{outer} keywords of the {\hyperref[\detokenize{run_style::doc}]{\sphinxcrossref{\DUrole{doc}{run\_style respa}}}} command, if
their sub\sphinxhyphen{}styles do.


\subsection{Restrictions}
\label{\detokenize{pair_hybrid:restrictions}}
\sphinxAtStartPar
When using a long\sphinxhyphen{}range Coulombic solver (via the
{\hyperref[\detokenize{kspace_style::doc}]{\sphinxcrossref{\DUrole{doc}{kspace\_style}}}} command) with a hybrid pair\_style,
one or more sub\sphinxhyphen{}styles will be of the “long” variety,
e.g. \sphinxstyleemphasis{lj/cut/coul/long} or \sphinxstyleemphasis{buck/coul/long}.  You must ensure that the
short\sphinxhyphen{}range Coulombic cutoff used by each of these long pair styles is
the same or else LAMMPS will generate an error.

\sphinxAtStartPar
Pair style \sphinxstyleemphasis{hybrid/scaled} currently only works for non\sphinxhyphen{}accelerated
pair styles and pair styles from the OPT package.

\sphinxAtStartPar
Pair style \sphinxstyleemphasis{hybrid/molecular} is not compatible with manybody potentials.

\sphinxAtStartPar
When using pair styles from the GPU package they must not be listed
multiple times.  LAMMPS will detect this and abort.


\subsection{Related commands}
\label{\detokenize{pair_hybrid:related-commands}}
\sphinxAtStartPar
{\hyperref[\detokenize{pair_coeff::doc}]{\sphinxcrossref{\DUrole{doc}{pair\_coeff}}}}


\subsection{Default}
\label{\detokenize{pair_hybrid:default}}
\sphinxAtStartPar
none

\sphinxstepscope

\index{pair\_style ilp/graphene/hbn@\spxentry{pair\_style ilp/graphene/hbn}}\ignorespaces 
\index{pair\_style ilp/graphene/hbn/opt@\spxentry{pair\_style ilp/graphene/hbn/opt}}\ignorespaces 

\section{pair\_style ilp/graphene/hbn command}
\label{\detokenize{pair_ilp_graphene_hbn:pair-style-ilp-graphene-hbn-command}}\label{\detokenize{pair_ilp_graphene_hbn:index-1}}\label{\detokenize{pair_ilp_graphene_hbn:index-0}}\label{\detokenize{pair_ilp_graphene_hbn::doc}}
\sphinxAtStartPar
Accelerator Variant: \sphinxstyleemphasis{ilp/graphene/hbn/opt}


\subsection{Syntax}
\label{\detokenize{pair_ilp_graphene_hbn:syntax}}
\begin{sphinxVerbatim}[commandchars=\\\{\}]
\PYG{k}{pair\PYGZus{}style}\PYG{+w}{ }\PYG{n}{[hybrid}\PYG{o}{/}\PYG{n}{overlay}\PYG{+w}{ }\PYG{n}{...]}\PYG{+w}{ }\PYG{n}{ilp}\PYG{o}{/}\PYG{n}{graphene}\PYG{o}{/}\PYG{n}{hbn}\PYG{+w}{ }\PYG{n}{cutoff}\PYG{+w}{ }\PYG{n}{tap\PYGZus{}flag}
\end{sphinxVerbatim}
\begin{itemize}
\item {} 
\sphinxAtStartPar
cutoff = global cutoff (distance units)

\item {} 
\sphinxAtStartPar
tap\_flag = 0/1 to turn off/on the taper function

\end{itemize}


\subsection{Examples}
\label{\detokenize{pair_ilp_graphene_hbn:examples}}
\begin{sphinxVerbatim}[commandchars=\\\{\}]
\PYG{k}{pair\PYGZus{}style}\PYG{+w}{  }\PYG{n}{hybrid}\PYG{o}{/}\PYG{n}{overlay}\PYG{+w}{ }\PYG{n}{ilp}\PYG{o}{/}\PYG{n}{graphene}\PYG{o}{/}\PYG{n}{hbn}\PYG{+w}{ }\PYG{l+m}{16.0}\PYG{+w}{ }\PYG{l+m}{1}
\PYG{k}{pair\PYGZus{}coeff}\PYG{+w}{  }\PYG{o}{*}\PYG{+w}{ }\PYG{o}{*}\PYG{+w}{ }\PYG{n}{ilp}\PYG{o}{/}\PYG{n}{graphene}\PYG{o}{/}\PYG{n}{hbn}\PYG{+w}{  }\PYG{n}{BNCH.ILP}\PYG{+w}{ }\PYG{n}{B}\PYG{+w}{ }\PYG{n}{N}\PYG{+w}{ }\PYG{n}{C}

\PYG{k}{pair\PYGZus{}style}\PYG{+w}{  }\PYG{n}{hybrid}\PYG{o}{/}\PYG{n}{overlay}\PYG{+w}{ }\PYG{n}{rebo}\PYG{+w}{ }\PYG{n}{tersoff}\PYG{+w}{ }\PYG{n}{ilp}\PYG{o}{/}\PYG{n}{graphene}\PYG{o}{/}\PYG{n}{hbn}\PYG{+w}{ }\PYG{l+m}{16.0}\PYG{+w}{ }\PYG{n}{coul}\PYG{o}{/}\PYG{n}{shield}\PYG{+w}{ }\PYG{l+m}{16.0}
\PYG{k}{pair\PYGZus{}coeff}\PYG{+w}{  }\PYG{o}{*}\PYG{+w}{ }\PYG{o}{*}\PYG{+w}{ }\PYG{n}{rebo}\PYG{+w}{              }\PYG{n}{CH.rebo}\PYG{+w}{     }\PYG{n}{NULL}\PYG{+w}{ }\PYG{n}{NULL}\PYG{+w}{ }\PYG{n}{C}
\PYG{k}{pair\PYGZus{}coeff}\PYG{+w}{  }\PYG{o}{*}\PYG{+w}{ }\PYG{o}{*}\PYG{+w}{ }\PYG{n}{tersoff}\PYG{+w}{           }\PYG{n}{BNC.tersoff}\PYG{+w}{ }\PYG{n}{B}\PYG{+w}{    }\PYG{n}{N}\PYG{+w}{    }\PYG{n}{NULL}
\PYG{k}{pair\PYGZus{}coeff}\PYG{+w}{  }\PYG{o}{*}\PYG{+w}{ }\PYG{o}{*}\PYG{+w}{ }\PYG{n}{ilp}\PYG{o}{/}\PYG{n}{graphene}\PYG{o}{/}\PYG{n}{hbn}\PYG{+w}{  }\PYG{n}{BNCH.ILP}\PYG{+w}{    }\PYG{n}{B}\PYG{+w}{    }\PYG{n}{N}\PYG{+w}{    }\PYG{n}{C}
\PYG{k}{pair\PYGZus{}coeff}\PYG{+w}{  }\PYG{l+m}{1}\PYG{+w}{ }\PYG{l+m}{1}\PYG{+w}{ }\PYG{n}{coul}\PYG{o}{/}\PYG{n}{shield}\PYG{+w}{ }\PYG{l+m}{0.70}
\PYG{k}{pair\PYGZus{}coeff}\PYG{+w}{  }\PYG{l+m}{1}\PYG{+w}{ }\PYG{l+m}{2}\PYG{+w}{ }\PYG{n}{coul}\PYG{o}{/}\PYG{n}{shield}\PYG{+w}{ }\PYG{l+m}{0.695}
\PYG{k}{pair\PYGZus{}coeff}\PYG{+w}{  }\PYG{l+m}{2}\PYG{+w}{ }\PYG{l+m}{2}\PYG{+w}{ }\PYG{n}{coul}\PYG{o}{/}\PYG{n}{shield}\PYG{+w}{ }\PYG{l+m}{0.69}
\end{sphinxVerbatim}


\subsection{Description}
\label{\detokenize{pair_ilp_graphene_hbn:description}}
\sphinxAtStartPar
The \sphinxstyleemphasis{ilp/graphene/hbn} style computes the registry\sphinxhyphen{}dependent interlayer
potential (ILP) potential as described in {\hyperref[\detokenize{pair_ilp_graphene_hbn:leven1}]{\sphinxcrossref{\DUrole{std,std-ref}{(Leven1)}}}},
{\hyperref[\detokenize{pair_ilp_graphene_hbn:leven2}]{\sphinxcrossref{\DUrole{std,std-ref}{(Leven2)}}}} and {\hyperref[\detokenize{pair_ilp_graphene_hbn:maaravi2}]{\sphinxcrossref{\DUrole{std,std-ref}{(Maaravi)}}}}.
The normals are calculated in the way as described
in {\hyperref[\detokenize{pair_ilp_graphene_hbn:kolmogorov2}]{\sphinxcrossref{\DUrole{std,std-ref}{(Kolmogorov)}}}}.
\begin{equation*}
\begin{split}E  = & \frac{1}{2} \sum_i \sum_{j \neq i} V_{ij} \\
V_{ij}  = & {\rm Tap}(r_{ij})\left \{ e^{-\alpha (r_{ij}/\beta -1)}
             \left [ \epsilon + f(\rho_{ij}) + f(\rho_{ji})\right ] -
              \frac{1}{1+e^{-d\left [ \left ( r_{ij}/\left (s_R \cdot r^{eff} \right ) \right )-1 \right ]}}
              \cdot \frac{C_6}{r^6_{ij}} \right \}\\
\rho_{ij}^2 = & r_{ij}^2 - ({\bf r}_{ij} \cdot {\bf n}_i)^2 \\
\rho_{ji}^2  = & r_{ij}^2 - ({\bf r}_{ij} \cdot {\bf n}_j)^2 \\
f(\rho)  = &  C e^{ -( \rho / \delta )^2 } \\
{\rm Tap}(r_{ij})  = & 20\left ( \frac{r_{ij}}{R_{cut}} \right )^7 -
                        70\left ( \frac{r_{ij}}{R_{cut}} \right )^6 +
                        84\left ( \frac{r_{ij}}{R_{cut}} \right )^5 -
                        35\left ( \frac{r_{ij}}{R_{cut}} \right )^4 + 1\end{split}
\end{equation*}
\sphinxAtStartPar
Where \(\mathrm{Tap}(r_{ij})\) is the taper function which provides
a continuous cutoff (up to third derivative) for interatomic separations
larger than \(r_c\) {\hyperref[\detokenize{pair_ilp_graphene_hbn:maaravi2}]{\sphinxcrossref{\DUrole{std,std-ref}{(Maaravi)}}}}. The definitions of
each parameter in the above equation can be found in {\hyperref[\detokenize{pair_ilp_graphene_hbn:leven1}]{\sphinxcrossref{\DUrole{std,std-ref}{(Leven1)}}}} and {\hyperref[\detokenize{pair_ilp_graphene_hbn:maaravi2}]{\sphinxcrossref{\DUrole{std,std-ref}{(Maaravi)}}}}.

\sphinxAtStartPar
It is important to include all the pairs to build the neighbor list for
calculating the normals.

\begin{sphinxadmonition}{note}{Note:}
\sphinxAtStartPar
This potential (ILP) is intended for interlayer interactions between two
different layers of graphene, hexagonal boron nitride (h\sphinxhyphen{}BN) and their hetero\sphinxhyphen{}junction.
To perform a realistic simulation, this potential must be used in combination with
intralayer potential, such as {\hyperref[\detokenize{pair_airebo::doc}]{\sphinxcrossref{\DUrole{doc}{AIREBO}}}} or {\hyperref[\detokenize{pair_tersoff::doc}]{\sphinxcrossref{\DUrole{doc}{Tersoff}}}} potential.
To keep the intralayer properties unaffected, the interlayer interaction
within the same layers should be avoided. Hence, each atom has to have a layer
identifier such that atoms residing on the same layer interact via the
appropriate intralayer potential and atoms residing on different layers
interact via the ILP. Here, the molecule id is chosen as the layer identifier,
thus a data file with the “full” atom style is required to use this potential.
\end{sphinxadmonition}

\sphinxAtStartPar
The parameter file (e.g. BNCH.ILP), is intended for use with \sphinxstyleemphasis{metal}
{\hyperref[\detokenize{units::doc}]{\sphinxcrossref{\DUrole{doc}{units}}}}, with energies in meV. Two additional parameters,
\sphinxstyleemphasis{S}, and \sphinxstyleemphasis{rcut} are included in the parameter file. \sphinxstyleemphasis{S} is designed to
facilitate scaling of energies. \sphinxstyleemphasis{rcut} is designed to build the neighbor
list for calculating the normals for each atom pair.

\begin{sphinxadmonition}{note}{Note:}
\sphinxAtStartPar
The parameters presented in the parameter file (e.g. BNCH.ILP),
are fitted with taper function by setting the cutoff equal to 16.0
Angstrom.  Using different cutoff or taper function should be careful.
The parameters for atoms pairs between Boron and Nitrogen are fitted with
a screened Coulomb interaction {\hyperref[\detokenize{pair_coul_shield::doc}]{\sphinxcrossref{\DUrole{doc}{coul/shield}}}}. Therefore,
to simulated the properties of h\sphinxhyphen{}BN correctly, this potential must be used in
combination with the pair style {\hyperref[\detokenize{pair_coul_shield::doc}]{\sphinxcrossref{\DUrole{doc}{coul/shield}}}}.
\end{sphinxadmonition}

\begin{sphinxadmonition}{note}{Note:}
\sphinxAtStartPar
Four new sets of parameters of ILP for 2D layered Materials with bilayer and
bulk configurations are presented in {\hyperref[\detokenize{pair_ilp_graphene_hbn:ouyang1}]{\sphinxcrossref{\DUrole{std,std-ref}{(Ouyang1)}}}} and {\hyperref[\detokenize{pair_ilp_graphene_hbn:ouyang2}]{\sphinxcrossref{\DUrole{std,std-ref}{(Ouyang2)}}}}, respectively.
These parameters provide a good description in both short\sphinxhyphen{} and long\sphinxhyphen{}range interaction regimes.
While the old ILP parameters published in {\hyperref[\detokenize{pair_ilp_graphene_hbn:leven2}]{\sphinxcrossref{\DUrole{std,std-ref}{(Leven2)}}}} and
{\hyperref[\detokenize{pair_ilp_graphene_hbn:maaravi2}]{\sphinxcrossref{\DUrole{std,std-ref}{(Maaravi)}}}} are only suitable for long\sphinxhyphen{}range interaction
regime. This feature is essential for simulations in high pressure
regime (i.e., the interlayer distance is smaller than the equilibrium
distance). The benchmark tests and comparison of these parameters can
be found in {\hyperref[\detokenize{pair_ilp_graphene_hbn:ouyang1}]{\sphinxcrossref{\DUrole{std,std-ref}{(Ouyang1)}}}} and {\hyperref[\detokenize{pair_ilp_graphene_hbn:ouyang2}]{\sphinxcrossref{\DUrole{std,std-ref}{(Ouyang2)}}}}.
\end{sphinxadmonition}

\sphinxAtStartPar
This potential must be used in combination with hybrid/overlay.
Other interactions can be set to zero using pair\_style \sphinxstyleemphasis{none}.

\sphinxAtStartPar
This pair style tallies a breakdown of the total interlayer potential
energy into sub\sphinxhyphen{}categories, which can be accessed via the {\hyperref[\detokenize{compute_pair::doc}]{\sphinxcrossref{\DUrole{doc}{compute pair}}}} command as a vector of values of length 2.
The 2 values correspond to the following sub\sphinxhyphen{}categories:
\begin{enumerate}
\sphinxsetlistlabels{\arabic}{enumi}{enumii}{}{.}%
\item {} 
\sphinxAtStartPar
\sphinxstyleemphasis{E\_vdW} = vdW (attractive) energy

\item {} 
\sphinxAtStartPar
\sphinxstyleemphasis{E\_Rep} = Repulsive energy

\end{enumerate}

\sphinxAtStartPar
To print these quantities to the log file (with descriptive column
headings) the following commands could be included in an input script:

\begin{sphinxVerbatim}[commandchars=\\\{\}]
\PYG{k}{compute }\PYG{n+nv+nvIdentifier}{0}\PYG{+w}{ }\PYG{n+nv+nvIdentifier}{all}\PYG{+w}{ }\PYG{n}{pair}\PYG{+w}{ }\PYG{n}{ilp}\PYG{o}{/}\PYG{n}{graphene}\PYG{o}{/}\PYG{n}{hbn}
\PYG{k}{variable }\PYG{n+nv+nvIdentifier}{Evdw}\PYG{+w}{  }\PYG{n}{equal}\PYG{+w}{ }\PYG{n}{c\PYGZus{}0[1]}
\PYG{k}{variable }\PYG{n+nv+nvIdentifier}{Erep}\PYG{+w}{  }\PYG{n}{equal}\PYG{+w}{ }\PYG{n}{c\PYGZus{}0[2]}
\PYG{k}{thermo\PYGZus{}style}\PYG{+w}{ }\PYG{n}{custom}\PYG{+w}{ }\PYG{n}{step}\PYG{+w}{ }\PYG{n}{temp}\PYG{+w}{ }\PYG{n}{epair}\PYG{+w}{ }\PYG{n}{v\PYGZus{}Erep}\PYG{+w}{ }\PYG{n}{v\PYGZus{}Evdw}
\end{sphinxVerbatim}


\bigskip\hrule\bigskip


\sphinxAtStartPar
Styles with a \sphinxstyleemphasis{gpu}, \sphinxstyleemphasis{intel}, \sphinxstyleemphasis{kk}, \sphinxstyleemphasis{omp}, or \sphinxstyleemphasis{opt} suffix are
functionally the same as the corresponding style without the suffix.
They have been optimized to run faster, depending on your available
hardware, as discussed on the {\hyperref[\detokenize{Speed_packages::doc}]{\sphinxcrossref{\DUrole{doc}{Accelerator packages}}}}
page.  The accelerated styles take the same arguments and should
produce the same results, except for round\sphinxhyphen{}off and precision issues.

\sphinxAtStartPar
These accelerated styles are part of the GPU, INTEL, KOKKOS,
OPENMP, and OPT packages, respectively.  They are only enabled if
LAMMPS was built with those packages.  See the {\hyperref[\detokenize{Build_package::doc}]{\sphinxcrossref{\DUrole{doc}{Build package}}}} page for more info.

\sphinxAtStartPar
You can specify the accelerated styles explicitly in your input script
by including their suffix, or you can use the {\hyperref[\detokenize{Run_options::doc}]{\sphinxcrossref{\DUrole{doc}{\sphinxhyphen{}suffix command\sphinxhyphen{}line switch}}}} when you invoke LAMMPS, or you can use the
{\hyperref[\detokenize{suffix::doc}]{\sphinxcrossref{\DUrole{doc}{suffix}}}} command in your input script.

\sphinxAtStartPar
See the {\hyperref[\detokenize{Speed_packages::doc}]{\sphinxcrossref{\DUrole{doc}{Accelerator packages}}}} page for more
instructions on how to use the accelerated styles effectively.


\bigskip\hrule\bigskip



\subsection{Mixing, shift, table, tail correction, restart, rRESPA info}
\label{\detokenize{pair_ilp_graphene_hbn:mixing-shift-table-tail-correction-restart-rrespa-info}}
\sphinxAtStartPar
This pair style does not support the pair\_modify mix, shift, table, and
tail options.

\sphinxAtStartPar
This pair style does not write their information to binary restart
files, since it is stored in potential files. Thus, you need to
re\sphinxhyphen{}specify the pair\_style and pair\_coeff commands in an input script
that reads a restart file.


\subsection{Restrictions}
\label{\detokenize{pair_ilp_graphene_hbn:restrictions}}
\sphinxAtStartPar
This pair style is part of the INTERLAYER package.  It is only enabled
if LAMMPS was built with that package.  See the {\hyperref[\detokenize{Build_package::doc}]{\sphinxcrossref{\DUrole{doc}{Build package}}}} page for more info.

\sphinxAtStartPar
This pair style requires the newton setting to be \sphinxstyleemphasis{on} for pair
interactions.

\sphinxAtStartPar
The BNCH.ILP potential file provided with LAMMPS (see the potentials
directory) are parameterized for \sphinxstyleemphasis{metal} units.  You can use this
potential with any LAMMPS units, but you would need to create your own
custom BNCH.ILP potential file with coefficients listed in the appropriate
units, if your simulation does not use \sphinxstyleemphasis{metal} units.


\subsection{Related commands}
\label{\detokenize{pair_ilp_graphene_hbn:related-commands}}
\sphinxAtStartPar
{\hyperref[\detokenize{pair_coeff::doc}]{\sphinxcrossref{\DUrole{doc}{pair\_coeff}}}},
{\hyperref[\detokenize{pair_none::doc}]{\sphinxcrossref{\DUrole{doc}{pair\_none}}}},
{\hyperref[\detokenize{pair_hybrid::doc}]{\sphinxcrossref{\DUrole{doc}{pair\_style hybrid/overlay}}}},
{\hyperref[\detokenize{pair_drip::doc}]{\sphinxcrossref{\DUrole{doc}{pair\_style drip}}}},
{\hyperref[\detokenize{pair_ilp_tmd::doc}]{\sphinxcrossref{\DUrole{doc}{pair\_style ilp\_tmd}}}},
{\hyperref[\detokenize{pair_saip_metal::doc}]{\sphinxcrossref{\DUrole{doc}{pair\_style saip\_metal}}}},
{\hyperref[\detokenize{pair_kolmogorov_crespi_z::doc}]{\sphinxcrossref{\DUrole{doc}{pair\_style pair\_kolmogorov\_crespi\_z}}}},
{\hyperref[\detokenize{pair_kolmogorov_crespi_full::doc}]{\sphinxcrossref{\DUrole{doc}{pair\_style pair\_kolmogorov\_crespi\_full}}}},
{\hyperref[\detokenize{pair_lebedeva_z::doc}]{\sphinxcrossref{\DUrole{doc}{pair\_style pair\_lebedeva\_z}}}},
{\hyperref[\detokenize{pair_coul_shield::doc}]{\sphinxcrossref{\DUrole{doc}{pair\_style pair\_coul\_shield}}}}.


\subsection{Default}
\label{\detokenize{pair_ilp_graphene_hbn:default}}
\sphinxAtStartPar
tap\_flag = 1


\bigskip\hrule\bigskip

\phantomsection\label{\detokenize{pair_ilp_graphene_hbn:ouyang1}}
\sphinxAtStartPar
\sphinxstylestrong{(Ouyang1)} W. Ouyang, D. Mandelli, M. Urbakh and O. Hod, Nano Lett. 18, 6009\sphinxhyphen{}6016 (2018).

\phantomsection\label{\detokenize{pair_ilp_graphene_hbn:ouyang2}}
\sphinxAtStartPar
\sphinxstylestrong{(Ouyang2)} W. Ouyang et al., J. Chem. Theory Comput. 16(1), 666\sphinxhyphen{}676 (2020).

\phantomsection\label{\detokenize{pair_ilp_graphene_hbn:leven1}}
\sphinxAtStartPar
\sphinxstylestrong{(Leven1)} I. Leven, I. Azuri, L. Kronik and O. Hod, J. Chem. Phys. 140, 104106 (2014).

\phantomsection\label{\detokenize{pair_ilp_graphene_hbn:leven2}}
\sphinxAtStartPar
\sphinxstylestrong{(Leven2)} I. Leven et al, J. Chem.Theory Comput. 12, 2896\sphinxhyphen{}905 (2016).

\phantomsection\label{\detokenize{pair_ilp_graphene_hbn:maaravi2}}
\sphinxAtStartPar
\sphinxstylestrong{(Maaravi)} T. Maaravi et al, J. Phys. Chem. C 121, 22826\sphinxhyphen{}22835 (2017).

\phantomsection\label{\detokenize{pair_ilp_graphene_hbn:kolmogorov2}}
\sphinxAtStartPar
\sphinxstylestrong{(Kolmogorov)} A. N. Kolmogorov, V. H. Crespi, Phys. Rev. B 71, 235415 (2005).

\sphinxstepscope

\index{pair\_style ilp/tmd@\spxentry{pair\_style ilp/tmd}}\ignorespaces 
\index{pair\_style ilp/tmd/opt@\spxentry{pair\_style ilp/tmd/opt}}\ignorespaces 

\section{pair\_style ilp/tmd command}
\label{\detokenize{pair_ilp_tmd:pair-style-ilp-tmd-command}}\label{\detokenize{pair_ilp_tmd:index-1}}\label{\detokenize{pair_ilp_tmd:index-0}}\label{\detokenize{pair_ilp_tmd::doc}}
\sphinxAtStartPar
Accelerator Variant: \sphinxstyleemphasis{ilp/tmd/opt}


\subsection{Syntax}
\label{\detokenize{pair_ilp_tmd:syntax}}
\begin{sphinxVerbatim}[commandchars=\\\{\}]
\PYG{k}{pair\PYGZus{}style}\PYG{+w}{ }\PYG{n}{[hybrid}\PYG{o}{/}\PYG{n}{overlay}\PYG{+w}{ }\PYG{n}{...]}\PYG{+w}{ }\PYG{n}{ilp}\PYG{o}{/}\PYG{n}{tmd}\PYG{+w}{ }\PYG{n}{cutoff}\PYG{+w}{ }\PYG{n}{tap\PYGZus{}flag}
\end{sphinxVerbatim}
\begin{itemize}
\item {} 
\sphinxAtStartPar
cutoff = global cutoff (distance units)

\item {} 
\sphinxAtStartPar
tap\_flag = 0/1 to turn off/on the taper function

\end{itemize}


\subsection{Examples}
\label{\detokenize{pair_ilp_tmd:examples}}
\begin{sphinxVerbatim}[commandchars=\\\{\}]
\PYG{k}{pair\PYGZus{}style}\PYG{+w}{  }\PYG{n}{hybrid}\PYG{o}{/}\PYG{n}{overlay}\PYG{+w}{ }\PYG{n}{ilp}\PYG{o}{/}\PYG{n}{tmd}\PYG{+w}{ }\PYG{l+m}{16.0}\PYG{+w}{ }\PYG{l+m}{1}
\PYG{k}{pair\PYGZus{}coeff}\PYG{+w}{  }\PYG{o}{*}\PYG{+w}{ }\PYG{o}{*}\PYG{+w}{ }\PYG{n}{ilp}\PYG{o}{/}\PYG{n}{tmd}\PYG{+w}{  }\PYG{n}{TMD.ILP}\PYG{+w}{ }\PYG{n}{Mo}\PYG{+w}{ }\PYG{n}{S}\PYG{+w}{ }\PYG{n}{S}

\PYG{k}{pair\PYGZus{}style}\PYG{+w}{  }\PYG{n}{hybrid}\PYG{o}{/}\PYG{n}{overlay}\PYG{+w}{ }\PYG{n}{sw}\PYG{o}{/}\PYG{n}{mod}\PYG{+w}{ }\PYG{n}{sw}\PYG{o}{/}\PYG{n}{mod}\PYG{+w}{ }\PYG{n}{ilp}\PYG{o}{/}\PYG{n}{tmd}\PYG{+w}{ }\PYG{l+m}{16.0}
\PYG{k}{pair\PYGZus{}coeff}\PYG{+w}{  }\PYG{o}{*}\PYG{+w}{ }\PYG{o}{*}\PYG{+w}{ }\PYG{n}{sw}\PYG{o}{/}\PYG{n}{mod}\PYG{+w}{ }\PYG{l+m}{1}\PYG{+w}{  }\PYG{n}{tmd.sw.mod}\PYG{+w}{ }\PYG{n}{Mo}\PYG{+w}{ }\PYG{n}{S}\PYG{+w}{ }\PYG{n}{S}\PYG{+w}{ }\PYG{n}{NULL}\PYG{+w}{ }\PYG{n}{NULL}\PYG{+w}{ }\PYG{n}{NULL}
\PYG{k}{pair\PYGZus{}coeff}\PYG{+w}{  }\PYG{o}{*}\PYG{+w}{ }\PYG{o}{*}\PYG{+w}{ }\PYG{n}{sw}\PYG{o}{/}\PYG{n}{mod}\PYG{+w}{ }\PYG{l+m}{2}\PYG{+w}{  }\PYG{n}{tmd.sw.mod}\PYG{+w}{ }\PYG{n}{NULL}\PYG{+w}{ }\PYG{n}{NULL}\PYG{+w}{ }\PYG{n}{NULL}\PYG{+w}{ }\PYG{n}{W}\PYG{+w}{ }\PYG{n}{Se}\PYG{+w}{ }\PYG{n}{Se}
\PYG{k}{pair\PYGZus{}coeff}\PYG{+w}{  }\PYG{o}{*}\PYG{+w}{ }\PYG{o}{*}\PYG{+w}{ }\PYG{n}{ilp}\PYG{o}{/}\PYG{n}{tmd}\PYG{+w}{   }\PYG{n}{TMD.ILP}\PYG{+w}{   }\PYG{n}{Mo}\PYG{+w}{ }\PYG{n}{S}\PYG{+w}{ }\PYG{n}{S}\PYG{+w}{ }\PYG{n}{W}\PYG{+w}{ }\PYG{n}{Se}\PYG{+w}{ }\PYG{n}{Se}
\end{sphinxVerbatim}


\subsection{Description}
\label{\detokenize{pair_ilp_tmd:description}}
\sphinxAtStartPar
\DUrole{versionmodified,added}{Added in version 17Feb2022.}

\sphinxAtStartPar
The \sphinxstyleemphasis{ilp/tmd} style computes the registry\sphinxhyphen{}dependent interlayer
potential (ILP) potential for transition metal dichalcogenides (TMD)
as described in {\hyperref[\detokenize{pair_ilp_tmd:ouyang7}]{\sphinxcrossref{\DUrole{std,std-ref}{(Ouyang7)}}}} and {\hyperref[\detokenize{pair_ilp_tmd:jiang}]{\sphinxcrossref{\DUrole{std,std-ref}{(Jiang)}}}}.
\begin{equation*}
\begin{split}E  = & \frac{1}{2} \sum_i \sum_{j \neq i} V_{ij} \\
V_{ij}  = & {\rm Tap}(r_{ij})\left \{ e^{-\alpha (r_{ij}/\beta -1)}
             \left [ \epsilon + f(\rho_{ij}) + f(\rho_{ji})\right ] -
              \frac{1}{1+e^{-d\left [ \left ( r_{ij}/\left (s_R \cdot r^{eff} \right ) \right )-1 \right ]}}
              \cdot \frac{C_6}{r^6_{ij}} \right \}\\
\rho_{ij}^2 = & r_{ij}^2 - ({\bf r}_{ij} \cdot {\bf n}_i)^2 \\
\rho_{ji}^2  = & r_{ij}^2 - ({\bf r}_{ij} \cdot {\bf n}_j)^2 \\
f(\rho)  = &  C e^{ -( \rho / \delta )^2 } \\
{\rm Tap}(r_{ij})  = & 20\left ( \frac{r_{ij}}{R_{cut}} \right )^7 -
                        70\left ( \frac{r_{ij}}{R_{cut}} \right )^6 +
                        84\left ( \frac{r_{ij}}{R_{cut}} \right )^5 -
                        35\left ( \frac{r_{ij}}{R_{cut}} \right )^4 + 1\end{split}
\end{equation*}
\sphinxAtStartPar
Where \(\mathrm{Tap}(r_{ij})\) is the taper function which provides
a continuous cutoff (up to third derivative) for interatomic separations
larger than \(r_c\) {\hyperref[\detokenize{pair_ilp_graphene_hbn::doc}]{\sphinxcrossref{\DUrole{doc}{pair\_style ilp\_graphene\_hbn}}}}.

\sphinxAtStartPar
It is important to include all the pairs to build the neighbor list for
calculating the normals.

\begin{sphinxadmonition}{note}{Note:}
\sphinxAtStartPar
Since each MX2 (M = Mo, W and X = S, Se Te) layer contains two
sub\sphinxhyphen{}layers of X atoms and one sub\sphinxhyphen{}layer of M atoms, the definition of the
normal vectors used for graphene and h\sphinxhyphen{}BN is no longer valid for TMDs.
In {\hyperref[\detokenize{pair_ilp_tmd:ouyang7}]{\sphinxcrossref{\DUrole{std,std-ref}{(Ouyang7)}}}}, a new definition is proposed, where for
each atom \sphinxtitleref{i}, its six nearest neighboring atoms belonging to the same
sub\sphinxhyphen{}layer are chosen to define the normal vector \sphinxtitleref{\{bf n\}\_i}.
\end{sphinxadmonition}

\sphinxAtStartPar
The parameter file (e.g. TMD.ILP), is intended for use with \sphinxstyleemphasis{metal}
{\hyperref[\detokenize{units::doc}]{\sphinxcrossref{\DUrole{doc}{units}}}}, with energies in meV. Two additional parameters,
\sphinxstyleemphasis{S}, and \sphinxstyleemphasis{rcut} are included in the parameter file. \sphinxstyleemphasis{S} is designed to
facilitate scaling of energies. \sphinxstyleemphasis{rcut} is designed to build the neighbor
list for calculating the normals for each atom pair.

\begin{sphinxadmonition}{note}{Note:}
\sphinxAtStartPar
The parameters presented in the parameter file (e.g. TMD.ILP),
are fitted with taper function by setting the cutoff equal to 16.0
Angstrom.  Using different cutoff or taper function should be careful.
These parameters provide a good description in both short\sphinxhyphen{} and long\sphinxhyphen{}range
interaction regimes. This feature is essential for simulations in high pressure
regime (i.e., the interlayer distance is smaller than the equilibrium
distance). The benchmark tests and comparison of these parameters can
be found in {\hyperref[\detokenize{pair_ilp_tmd:ouyang7}]{\sphinxcrossref{\DUrole{std,std-ref}{(Ouyang7)}}}}.
\end{sphinxadmonition}

\sphinxAtStartPar
This potential must be used in combination with hybrid/overlay.
Other interactions can be set to zero using pair\_style \sphinxstyleemphasis{none}.

\sphinxAtStartPar
This pair style tallies a breakdown of the total interlayer potential
energy into sub\sphinxhyphen{}categories, which can be accessed via the {\hyperref[\detokenize{compute_pair::doc}]{\sphinxcrossref{\DUrole{doc}{compute pair}}}} command as a vector of values of length 2.
The 2 values correspond to the following sub\sphinxhyphen{}categories:
\begin{enumerate}
\sphinxsetlistlabels{\arabic}{enumi}{enumii}{}{.}%
\item {} 
\sphinxAtStartPar
\sphinxstyleemphasis{E\_vdW} = vdW (attractive) energy

\item {} 
\sphinxAtStartPar
\sphinxstyleemphasis{E\_Rep} = Repulsive energy

\end{enumerate}

\sphinxAtStartPar
To print these quantities to the log file (with descriptive column
headings) the following commands could be included in an input script:

\begin{sphinxVerbatim}[commandchars=\\\{\}]
\PYG{k}{compute }\PYG{n+nv+nvIdentifier}{0}\PYG{+w}{ }\PYG{n+nv+nvIdentifier}{all}\PYG{+w}{ }\PYG{n}{pair}\PYG{+w}{ }\PYG{n}{ilp}\PYG{o}{/}\PYG{n}{tmd}
\PYG{k}{variable }\PYG{n+nv+nvIdentifier}{Evdw}\PYG{+w}{  }\PYG{n}{equal}\PYG{+w}{ }\PYG{n}{c\PYGZus{}0[1]}
\PYG{k}{variable }\PYG{n+nv+nvIdentifier}{Erep}\PYG{+w}{  }\PYG{n}{equal}\PYG{+w}{ }\PYG{n}{c\PYGZus{}0[2]}
\PYG{k}{thermo\PYGZus{}style}\PYG{+w}{ }\PYG{n}{custom}\PYG{+w}{ }\PYG{n}{step}\PYG{+w}{ }\PYG{n}{temp}\PYG{+w}{ }\PYG{n}{epair}\PYG{+w}{ }\PYG{n}{v\PYGZus{}Erep}\PYG{+w}{ }\PYG{n}{v\PYGZus{}Evdw}
\end{sphinxVerbatim}


\bigskip\hrule\bigskip


\sphinxAtStartPar
Styles with a \sphinxstyleemphasis{gpu}, \sphinxstyleemphasis{intel}, \sphinxstyleemphasis{kk}, \sphinxstyleemphasis{omp}, or \sphinxstyleemphasis{opt} suffix are
functionally the same as the corresponding style without the suffix.
They have been optimized to run faster, depending on your available
hardware, as discussed on the {\hyperref[\detokenize{Speed_packages::doc}]{\sphinxcrossref{\DUrole{doc}{Accelerator packages}}}}
page.  The accelerated styles take the same arguments and should
produce the same results, except for round\sphinxhyphen{}off and precision issues.

\sphinxAtStartPar
These accelerated styles are part of the GPU, INTEL, KOKKOS,
OPENMP, and OPT packages, respectively.  They are only enabled if
LAMMPS was built with those packages.  See the {\hyperref[\detokenize{Build_package::doc}]{\sphinxcrossref{\DUrole{doc}{Build package}}}} page for more info.

\sphinxAtStartPar
You can specify the accelerated styles explicitly in your input script
by including their suffix, or you can use the {\hyperref[\detokenize{Run_options::doc}]{\sphinxcrossref{\DUrole{doc}{\sphinxhyphen{}suffix command\sphinxhyphen{}line switch}}}} when you invoke LAMMPS, or you can use the
{\hyperref[\detokenize{suffix::doc}]{\sphinxcrossref{\DUrole{doc}{suffix}}}} command in your input script.

\sphinxAtStartPar
See the {\hyperref[\detokenize{Speed_packages::doc}]{\sphinxcrossref{\DUrole{doc}{Accelerator packages}}}} page for more
instructions on how to use the accelerated styles effectively.


\bigskip\hrule\bigskip



\subsection{Mixing, shift, table, tail correction, restart, rRESPA info}
\label{\detokenize{pair_ilp_tmd:mixing-shift-table-tail-correction-restart-rrespa-info}}
\sphinxAtStartPar
This pair style does not support the pair\_modify mix, shift, table, and
tail options.

\sphinxAtStartPar
This pair style does not write their information to binary restart
files, since it is stored in potential files. Thus, you need to
re\sphinxhyphen{}specify the pair\_style and pair\_coeff commands in an input script
that reads a restart file.


\subsection{Restrictions}
\label{\detokenize{pair_ilp_tmd:restrictions}}
\sphinxAtStartPar
This pair style is part of the INTERLAYER package.  It is only enabled
if LAMMPS was built with that package.  See the {\hyperref[\detokenize{Build_package::doc}]{\sphinxcrossref{\DUrole{doc}{Build package}}}} page for more info.

\sphinxAtStartPar
This pair style requires the newton setting to be \sphinxstyleemphasis{on} for pair
interactions.

\sphinxAtStartPar
The TMD.ILP potential file provided with LAMMPS (see the potentials
directory) are parameterized for \sphinxstyleemphasis{metal} units.  You can use this
potential with any LAMMPS units, but you would need to create your own
custom TMD.ILP potential file with coefficients listed in the appropriate
units, if your simulation does not use \sphinxstyleemphasis{metal} units.


\subsection{Related commands}
\label{\detokenize{pair_ilp_tmd:related-commands}}
\sphinxAtStartPar
{\hyperref[\detokenize{pair_coeff::doc}]{\sphinxcrossref{\DUrole{doc}{pair\_coeff}}}},
{\hyperref[\detokenize{pair_none::doc}]{\sphinxcrossref{\DUrole{doc}{pair\_none}}}},
{\hyperref[\detokenize{pair_hybrid::doc}]{\sphinxcrossref{\DUrole{doc}{pair\_style hybrid/overlay}}}},
{\hyperref[\detokenize{pair_drip::doc}]{\sphinxcrossref{\DUrole{doc}{pair\_style drip}}}},
{\hyperref[\detokenize{pair_saip_metal::doc}]{\sphinxcrossref{\DUrole{doc}{pair\_style saip\_metal}}}},
{\hyperref[\detokenize{pair_ilp_graphene_hbn::doc}]{\sphinxcrossref{\DUrole{doc}{pair\_style ilp\_graphene\_hbn}}}},
{\hyperref[\detokenize{pair_kolmogorov_crespi_z::doc}]{\sphinxcrossref{\DUrole{doc}{pair\_style pair\_kolmogorov\_crespi\_z}}}},
{\hyperref[\detokenize{pair_kolmogorov_crespi_full::doc}]{\sphinxcrossref{\DUrole{doc}{pair\_style pair\_kolmogorov\_crespi\_full}}}},
{\hyperref[\detokenize{pair_lebedeva_z::doc}]{\sphinxcrossref{\DUrole{doc}{pair\_style pair\_lebedeva\_z}}}},
{\hyperref[\detokenize{pair_coul_shield::doc}]{\sphinxcrossref{\DUrole{doc}{pair\_style pair\_coul\_shield}}}}.


\subsection{Default}
\label{\detokenize{pair_ilp_tmd:default}}
\sphinxAtStartPar
tap\_flag = 1


\bigskip\hrule\bigskip

\phantomsection\label{\detokenize{pair_ilp_tmd:ouyang7}}
\sphinxAtStartPar
\sphinxstylestrong{(Ouyang7)} W. Ouyang, et al., J. Chem. Theory Comput. 17, 7237 (2021).

\phantomsection\label{\detokenize{pair_ilp_tmd:jiang}}
\sphinxAtStartPar
\sphinxstylestrong{(Jiang)} W. Jiang, et al., J. Phys. Chem. A, 127, 46, 9820\sphinxhyphen{}9830 (2023).

\sphinxstepscope

\index{pair\_style kim@\spxentry{pair\_style kim}}\ignorespaces 

\section{pair\_style kim command}
\label{\detokenize{pair_kim:pair-style-kim-command}}\label{\detokenize{pair_kim:index-0}}\label{\detokenize{pair_kim::doc}}

\subsection{Syntax}
\label{\detokenize{pair_kim:syntax}}
\begin{sphinxVerbatim}[commandchars=\\\{\}]
\PYG{k}{pair\PYGZus{}style}\PYG{+w}{ }\PYG{k}{kim }\PYG{l+s+s1}{model}
\end{sphinxVerbatim}

\sphinxAtStartPar
model = name of a KIM model (the KIM ID for models archived in OpenKIM)


\subsection{Examples}
\label{\detokenize{pair_kim:examples}}
\begin{sphinxVerbatim}[commandchars=\\\{\}]
\PYG{k}{pair\PYGZus{}style}\PYG{+w}{ }\PYG{k}{kim }\PYG{l+s+s1}{SW\PYGZus{}StillingerWeber\PYGZus{}1985\PYGZus{}Si\PYGZus{}\PYGZus{}MO\PYGZus{}405512056662\PYGZus{}005}
\PYG{k}{pair\PYGZus{}coeff}\PYG{+w}{ }\PYG{o}{*}\PYG{+w}{ }\PYG{o}{*}\PYG{+w}{ }\PYG{n}{Si}
\end{sphinxVerbatim}


\subsection{Description}
\label{\detokenize{pair_kim:description}}
\sphinxAtStartPar
This pair style is a wrapper on the
\sphinxhref{https://openkim.org}{Open Knowledgebase of Interatomic Models (OpenKIM)}
repository of interatomic potentials to enable their use in LAMMPS scripts.

\sphinxAtStartPar
The preferred interface for using interatomic models archived in
OpenKIM is the {\hyperref[\detokenize{kim_commands::doc}]{\sphinxcrossref{\DUrole{doc}{kim command}}}} interface. That
interface supports both “KIM Portable Models” (PMs) that conform to the
KIM API Portable Model Interface (PMI) and can be used by any
simulation code that conforms to the KIM API/PMI, and
“KIM Simulator Models” (SMs) that are natively implemented within a single
simulation code (like LAMMPS) and can only be used with it.
The \sphinxstyleemphasis{pair\_style kim} command is limited to KIM PMs. It is
used by the {\hyperref[\detokenize{kim_commands::doc}]{\sphinxcrossref{\DUrole{doc}{kim command}}}} interface as needed.

\begin{sphinxadmonition}{note}{Note:}
\sphinxAtStartPar
Since \sphinxstyleemphasis{pair\_style kim} is called by \sphinxstyleemphasis{kim interactions} as needed,
it is not recommended to be directly used in input scripts.
\end{sphinxadmonition}


\bigskip\hrule\bigskip


\sphinxAtStartPar
The argument \sphinxstyleemphasis{model} is the name of the KIM PM.
For potentials archived in OpenKIM
this is the extended KIM ID (see {\hyperref[\detokenize{kim_commands::doc}]{\sphinxcrossref{\DUrole{doc}{kim command}}}}
for details). LAMMPS can invoke any KIM PM, however there can
be incompatibilities (for example due to unit matching issues).
In the event of an incompatibility, the code will terminate with
an error message. Check both the LAMMPS and KIM log files for details.

\sphinxAtStartPar
Only a single \sphinxstyleemphasis{pair\_coeff} command is used with the \sphinxstyleemphasis{kim} style, which
specifies the mapping of LAMMPS atom types to the species supported by
the KIM PM.  This is done by specifying \sphinxstyleemphasis{N} additional arguments
after the * * in the \sphinxstyleemphasis{pair\_coeff} command, where \sphinxstyleemphasis{N} is the number of
LAMMPS atom types:
\begin{itemize}
\item {} 
\sphinxAtStartPar
N element names = mapping of KIM elements to atom types

\end{itemize}

\sphinxAtStartPar
For example, consider a KIM PM that supports Si and C species.
If the LAMMPS simulation has four atom types, where the first three are Si,
and the fourth is C, the following \sphinxstyleemphasis{pair\_coeff} command would be used:

\begin{sphinxVerbatim}[commandchars=\\\{\}]
\PYG{k}{pair\PYGZus{}coeff}\PYG{+w}{ }\PYG{o}{*}\PYG{+w}{ }\PYG{o}{*}\PYG{+w}{ }\PYG{n}{Si}\PYG{+w}{ }\PYG{n}{Si}\PYG{+w}{ }\PYG{n}{Si}\PYG{+w}{ }\PYG{n}{C}
\end{sphinxVerbatim}

\sphinxAtStartPar
The first two arguments must be * * so as to span all LAMMPS atom types.
The first three Si arguments map LAMMPS atom types 1, 2, and 3 to Si as
defined within KIM PM.  The final C argument maps LAMMPS atom type 4 to C.


\bigskip\hrule\bigskip


\sphinxAtStartPar
In addition to the usual LAMMPS error messages, the KIM library itself
may generate errors, which should be printed to the screen.  In this
case it is also useful to check the \sphinxstyleemphasis{kim.log} file for additional error
information.  The file \sphinxstyleemphasis{kim.log} should be generated in the same
directory where LAMMPS is running.

\sphinxAtStartPar
To download, build, and install the KIM library on your system, see
the \sphinxstyleemphasis{lib/kim/README} file.  Once you have done this and built LAMMPS
with the KIM package installed you can run the example input scripts
in \sphinxstyleemphasis{examples/kim}.


\bigskip\hrule\bigskip



\subsection{Mixing, shift, table, tail correction, restart, rRESPA info}
\label{\detokenize{pair_kim:mixing-shift-table-tail-correction-restart-rrespa-info}}
\sphinxAtStartPar
This pair style does not support the {\hyperref[\detokenize{pair_modify::doc}]{\sphinxcrossref{\DUrole{doc}{pair\_modify}}}}
mix, shift, table, and tail options.

\sphinxAtStartPar
This pair style does not write its information to {\hyperref[\detokenize{restart::doc}]{\sphinxcrossref{\DUrole{doc}{binary restart files}}}}, since KIM stores the potential parameters.
Thus, you need to re\sphinxhyphen{}specify the pair\_style and pair\_coeff commands in
an input script that reads a restart file.

\sphinxAtStartPar
This pair style can only be used via the \sphinxstyleemphasis{pair} keyword of the
{\hyperref[\detokenize{run_style::doc}]{\sphinxcrossref{\DUrole{doc}{run\_style respa}}}} command.  It does not support the
\sphinxstyleemphasis{inner}, \sphinxstyleemphasis{middle}, \sphinxstyleemphasis{outer} keywords.


\bigskip\hrule\bigskip



\subsection{Restrictions}
\label{\detokenize{pair_kim:restrictions}}
\sphinxAtStartPar
This pair style is part of the KIM package. See details on
restrictions in {\hyperref[\detokenize{kim_commands::doc}]{\sphinxcrossref{\DUrole{doc}{kim command}}}}.

\sphinxAtStartPar
This current version of pair\_style kim is compatible with the
kim\sphinxhyphen{}api package version 2.0.0 and higher.


\subsection{Related commands}
\label{\detokenize{pair_kim:related-commands}}
\sphinxAtStartPar
{\hyperref[\detokenize{pair_coeff::doc}]{\sphinxcrossref{\DUrole{doc}{pair\_coeff}}}}, {\hyperref[\detokenize{kim_commands::doc}]{\sphinxcrossref{\DUrole{doc}{kim command}}}}


\subsection{Default}
\label{\detokenize{pair_kim:default}}
\sphinxAtStartPar
none

\sphinxstepscope

\index{pair\_style kolmogorov/crespi/full@\spxentry{pair\_style kolmogorov/crespi/full}}\ignorespaces 

\section{pair\_style kolmogorov/crespi/full command}
\label{\detokenize{pair_kolmogorov_crespi_full:pair-style-kolmogorov-crespi-full-command}}\label{\detokenize{pair_kolmogorov_crespi_full:index-0}}\label{\detokenize{pair_kolmogorov_crespi_full::doc}}

\subsection{Syntax}
\label{\detokenize{pair_kolmogorov_crespi_full:syntax}}
\begin{sphinxVerbatim}[commandchars=\\\{\}]
\PYG{k}{pair\PYGZus{}style}\PYG{+w}{ }\PYG{n}{hybrid}\PYG{o}{/}\PYG{n}{overlay}\PYG{+w}{ }\PYG{n}{kolmogorov}\PYG{o}{/}\PYG{n}{crespi}\PYG{o}{/}\PYG{n}{full}\PYG{+w}{ }\PYG{n}{cutoff}\PYG{+w}{ }\PYG{n}{tap\PYGZus{}flag}
\end{sphinxVerbatim}
\begin{itemize}
\item {} 
\sphinxAtStartPar
cutoff = global cutoff (distance units)

\item {} 
\sphinxAtStartPar
tap\_flag = 0/1 to turn off/on the taper function

\end{itemize}


\subsection{Examples}
\label{\detokenize{pair_kolmogorov_crespi_full:examples}}
\begin{sphinxVerbatim}[commandchars=\\\{\}]
\PYG{k}{pair\PYGZus{}style}\PYG{+w}{ }\PYG{n}{hybrid}\PYG{o}{/}\PYG{n}{overlay}\PYG{+w}{ }\PYG{n}{kolmogorov}\PYG{o}{/}\PYG{n}{crespi}\PYG{o}{/}\PYG{n}{full}\PYG{+w}{ }\PYG{l+m}{20.0}\PYG{+w}{ }\PYG{l+m}{0}
\PYG{k}{pair\PYGZus{}coeff}\PYG{+w}{ }\PYG{o}{*}\PYG{+w}{ }\PYG{o}{*}\PYG{+w}{ }\PYG{n}{none}
\PYG{k}{pair\PYGZus{}coeff}\PYG{+w}{ }\PYG{o}{*}\PYG{+w}{ }\PYG{o}{*}\PYG{+w}{ }\PYG{n}{kolmogorov}\PYG{o}{/}\PYG{n}{crespi}\PYG{o}{/}\PYG{n}{full}\PYG{+w}{  }\PYG{n}{CH.KC}\PYG{+w}{   }\PYG{n}{C}\PYG{+w}{ }\PYG{n}{C}

\PYG{k}{pair\PYGZus{}style}\PYG{+w}{ }\PYG{n}{hybrid}\PYG{o}{/}\PYG{n}{overlay}\PYG{+w}{ }\PYG{n}{rebo}\PYG{+w}{ }\PYG{n}{kolmogorov}\PYG{o}{/}\PYG{n}{crespi}\PYG{o}{/}\PYG{n}{full}\PYG{+w}{ }\PYG{l+m}{16.0}\PYG{+w}{ }\PYG{l+m}{1}
\PYG{k}{pair\PYGZus{}coeff}\PYG{+w}{ }\PYG{o}{*}\PYG{+w}{ }\PYG{o}{*}\PYG{+w}{ }\PYG{n}{rebo}\PYG{+w}{                    }\PYG{n}{CH.rebo}\PYG{+w}{      }\PYG{n}{C}\PYG{+w}{ }\PYG{n}{H}
\PYG{k}{pair\PYGZus{}coeff}\PYG{+w}{ }\PYG{o}{*}\PYG{+w}{ }\PYG{o}{*}\PYG{+w}{ }\PYG{n}{kolmogorov}\PYG{o}{/}\PYG{n}{crespi}\PYG{o}{/}\PYG{n}{full}\PYG{+w}{  }\PYG{n}{CH\PYGZus{}taper.KC}\PYG{+w}{  }\PYG{n}{C}\PYG{+w}{ }\PYG{n}{H}
\end{sphinxVerbatim}


\subsection{Description}
\label{\detokenize{pair_kolmogorov_crespi_full:description}}
\sphinxAtStartPar
The \sphinxstyleemphasis{kolmogorov/crespi/full} style computes the Kolmogorov\sphinxhyphen{}Crespi (KC)
interaction potential as described in {\hyperref[\detokenize{pair_kolmogorov_crespi_full:kolmogorov1}]{\sphinxcrossref{\DUrole{std,std-ref}{(Kolmogorov)}}}}.
No simplification is made,
\begin{equation*}
\begin{split} E  = & \frac{1}{2} \sum_i \sum_{j \neq i} V_{ij} \\
 V_{ij}  = & e^{-\lambda (r_{ij} -z_0)} \left [ C + f(\rho_{ij}) + f(\rho_{ji}) \right ] - A \left ( \frac{r_{ij}}{z_0}\right )^{-6} \\
\rho_{ij}^2 = & r_{ij}^2 - ({\bf r}_{ij}\cdot {\bf n}_{i})^2 \\
\rho_{ji}^2 = & r_{ij}^2 - ({\bf r}_{ij}\cdot  {\bf n}_{j})^2 \\
f(\rho) & =  e^{-(\rho/\delta)^2} \sum_{n=0}^2 C_{2n} { (\rho/\delta) }^{2n}\end{split}
\end{equation*}
\sphinxAtStartPar
It is important to have a sufficiently large cutoff to ensure smooth
forces and to include all the pairs to build the neighbor list for
calculating the normals.  Energies are shifted so that they go
continuously to zero at the cutoff assuming that the exponential part of
\(V_{ij}\) (first term) decays sufficiently fast.  This shift is achieved by
the last term in the equation for \(V_{ij}\) above. This is essential only
when the tapper function is turned off. The formula of taper function
can be found in pair style {\hyperref[\detokenize{pair_ilp_graphene_hbn::doc}]{\sphinxcrossref{\DUrole{doc}{ilp/graphene/hbn}}}}.

\begin{sphinxadmonition}{note}{Note:}
\sphinxAtStartPar
This potential (ILP) is intended for interlayer interactions between two
different layers of graphene. To perform a realistic simulation, this potential
must be used in combination with intralayer potential, such as
{\hyperref[\detokenize{pair_airebo::doc}]{\sphinxcrossref{\DUrole{doc}{AIREBO}}}} or {\hyperref[\detokenize{pair_tersoff::doc}]{\sphinxcrossref{\DUrole{doc}{Tersoff}}}} potential.
To keep the intralayer properties unaffected, the interlayer interaction
within the same layers should be avoided. Hence, each atom has to have a layer
identifier such that atoms residing on the same layer interact via the
appropriate intralayer potential and atoms residing on different layers
interact via the ILP. Here, the molecule id is chosen as the layer identifier,
thus a data file with the “full” atom style is required to use this potential.
\end{sphinxadmonition}

\sphinxAtStartPar
The parameter file (e.g. CH.KC), is intended for use with \sphinxstyleemphasis{metal}
{\hyperref[\detokenize{units::doc}]{\sphinxcrossref{\DUrole{doc}{units}}}}, with energies in meV. Two additional parameters, \sphinxstyleemphasis{S},
and \sphinxstyleemphasis{rcut} are included in the parameter file. \sphinxstyleemphasis{S} is designed to
facilitate scaling of energies. \sphinxstyleemphasis{rcut} is designed to build the neighbor
list for calculating the normals for each atom pair.

\begin{sphinxadmonition}{note}{Note:}
\sphinxAtStartPar
Two new sets of parameters of KC potential for hydrocarbons, CH.KC
(without the taper function) and CH\_taper.KC (with the taper function)
are presented in {\hyperref[\detokenize{pair_kolmogorov_crespi_full:ouyang3}]{\sphinxcrossref{\DUrole{std,std-ref}{(Ouyang1)}}}}.  The energy for the KC potential
with the taper function goes continuously to zero at the cutoff.  The
parameters in both CH.KC and CH\_taper.KC provide a good description in
both short\sphinxhyphen{} and long\sphinxhyphen{}range interaction regimes. While the original
parameters (CC.KC) published in {\hyperref[\detokenize{pair_kolmogorov_crespi_full:kolmogorov1}]{\sphinxcrossref{\DUrole{std,std-ref}{(Kolmogorov)}}}} are only
suitable for long\sphinxhyphen{}range interaction regime.  This feature is essential
for simulations in high pressure regime (i.e., the interlayer distance
is smaller than the equilibrium distance).  The benchmark tests and
comparison of these parameters can be found in {\hyperref[\detokenize{pair_kolmogorov_crespi_full:ouyang3}]{\sphinxcrossref{\DUrole{std,std-ref}{(Ouyang1)}}}} and {\hyperref[\detokenize{pair_kolmogorov_crespi_full:ouyang4}]{\sphinxcrossref{\DUrole{std,std-ref}{(Ouyang2)}}}}.
\end{sphinxadmonition}

\sphinxAtStartPar
This potential must be used in combination with hybrid/overlay.
Other interactions can be set to zero using pair\_style \sphinxstyleemphasis{none}.

\sphinxAtStartPar
This pair style tallies a breakdown of the total interlayer potential
energy into sub\sphinxhyphen{}categories, which can be accessed via the {\hyperref[\detokenize{compute_pair::doc}]{\sphinxcrossref{\DUrole{doc}{compute pair}}}} command as a vector of values of length 2.
The 2 values correspond to the following sub\sphinxhyphen{}categories:
\begin{enumerate}
\sphinxsetlistlabels{\arabic}{enumi}{enumii}{}{.}%
\item {} 
\sphinxAtStartPar
\sphinxstyleemphasis{E\_vdW} = vdW (attractive) energy

\item {} 
\sphinxAtStartPar
\sphinxstyleemphasis{E\_Rep} = Repulsive energy

\end{enumerate}

\sphinxAtStartPar
To print these quantities to the log file (with descriptive column
headings) the following commands could be included in an input script:

\begin{sphinxVerbatim}[commandchars=\\\{\}]
\PYG{k}{compute }\PYG{n+nv+nvIdentifier}{0}\PYG{+w}{ }\PYG{n+nv+nvIdentifier}{all}\PYG{+w}{ }\PYG{n}{pair}\PYG{+w}{ }\PYG{n}{kolmogorov}\PYG{o}{/}\PYG{n}{crespi}\PYG{o}{/}\PYG{n}{full}
\PYG{k}{variable }\PYG{n+nv+nvIdentifier}{Evdw}\PYG{+w}{  }\PYG{n}{equal}\PYG{+w}{ }\PYG{n}{c\PYGZus{}0[1]}
\PYG{k}{variable }\PYG{n+nv+nvIdentifier}{Erep}\PYG{+w}{  }\PYG{n}{equal}\PYG{+w}{ }\PYG{n}{c\PYGZus{}0[2]}
\PYG{k}{thermo\PYGZus{}style}\PYG{+w}{ }\PYG{n}{custom}\PYG{+w}{ }\PYG{n}{step}\PYG{+w}{ }\PYG{n}{temp}\PYG{+w}{ }\PYG{n}{epair}\PYG{+w}{ }\PYG{n}{v\PYGZus{}Erep}\PYG{+w}{ }\PYG{n}{v\PYGZus{}Evdw}
\end{sphinxVerbatim}


\bigskip\hrule\bigskip



\subsection{Mixing, shift, table, tail correction, restart, rRESPA info}
\label{\detokenize{pair_kolmogorov_crespi_full:mixing-shift-table-tail-correction-restart-rrespa-info}}
\sphinxAtStartPar
This pair style does not support the pair\_modify mix, shift, table,
and tail options.

\sphinxAtStartPar
This pair style does not write their information to binary restart
files, since it is stored in potential files. Thus, you need to
re\sphinxhyphen{}specify the pair\_style and pair\_coeff commands in an input script
that reads a restart file.


\subsection{Restrictions}
\label{\detokenize{pair_kolmogorov_crespi_full:restrictions}}
\sphinxAtStartPar
This pair style is part of the INTERLAYER package.  It is only enabled
if LAMMPS was built with that package.  See the {\hyperref[\detokenize{Build_package::doc}]{\sphinxcrossref{\DUrole{doc}{Build package}}}} page for more info.

\sphinxAtStartPar
This pair style requires the newton setting to be \sphinxstyleemphasis{on} for pair
interactions.

\sphinxAtStartPar
The CH.KC potential file provided with LAMMPS (see the potentials
folder) is parameterized for metal units.  You can use this pair style
with any LAMMPS units, but you would need to create your own custom
CH.KC potential file with all coefficients converted to the appropriate
units.


\subsection{Related commands}
\label{\detokenize{pair_kolmogorov_crespi_full:related-commands}}
\sphinxAtStartPar
{\hyperref[\detokenize{pair_coeff::doc}]{\sphinxcrossref{\DUrole{doc}{pair\_coeff}}}},
{\hyperref[\detokenize{pair_none::doc}]{\sphinxcrossref{\DUrole{doc}{pair\_none}}}},
{\hyperref[\detokenize{pair_hybrid::doc}]{\sphinxcrossref{\DUrole{doc}{pair\_style hybrid/overlay}}}},
{\hyperref[\detokenize{pair_drip::doc}]{\sphinxcrossref{\DUrole{doc}{pair\_style drip}}}},
{\hyperref[\detokenize{pair_lebedeva_z::doc}]{\sphinxcrossref{\DUrole{doc}{pair\_style pair\_lebedeva\_z}}}},
{\hyperref[\detokenize{pair_kolmogorov_crespi_z::doc}]{\sphinxcrossref{\DUrole{doc}{pair\_style kolmogorov/crespi/z}}}},
{\hyperref[\detokenize{pair_ilp_graphene_hbn::doc}]{\sphinxcrossref{\DUrole{doc}{pair\_style ilp/graphene/hbn}}}}.


\subsection{Default}
\label{\detokenize{pair_kolmogorov_crespi_full:default}}
\sphinxAtStartPar
tap\_flag = 0


\bigskip\hrule\bigskip

\phantomsection\label{\detokenize{pair_kolmogorov_crespi_full:kolmogorov1}}
\sphinxAtStartPar
\sphinxstylestrong{(Kolmogorov)} A. N. Kolmogorov, V. H. Crespi, Phys. Rev. B 71, 235415 (2005)

\phantomsection\label{\detokenize{pair_kolmogorov_crespi_full:ouyang3}}
\sphinxAtStartPar
\sphinxstylestrong{(Ouyang1)} W. Ouyang, D. Mandelli, M. Urbakh and O. Hod, Nano Lett. 18, 6009\sphinxhyphen{}6016 (2018).

\phantomsection\label{\detokenize{pair_kolmogorov_crespi_full:ouyang4}}
\sphinxAtStartPar
\sphinxstylestrong{(Ouyang2)} W. Ouyang et al., J. Chem. Theory Comput. 16(1), 666\sphinxhyphen{}676 (2020).

\sphinxstepscope

\index{pair\_style kolmogorov/crespi/z@\spxentry{pair\_style kolmogorov/crespi/z}}\ignorespaces 

\section{pair\_style kolmogorov/crespi/z command}
\label{\detokenize{pair_kolmogorov_crespi_z:pair-style-kolmogorov-crespi-z-command}}\label{\detokenize{pair_kolmogorov_crespi_z:index-0}}\label{\detokenize{pair_kolmogorov_crespi_z::doc}}

\subsection{Syntax}
\label{\detokenize{pair_kolmogorov_crespi_z:syntax}}
\begin{sphinxVerbatim}[commandchars=\\\{\}]
\PYG{k}{pair\PYGZus{}style}\PYG{+w}{ }\PYG{n}{[hybrid}\PYG{o}{/}\PYG{n}{overlay}\PYG{+w}{ }\PYG{n}{...]}\PYG{+w}{ }\PYG{n}{kolmogorov}\PYG{o}{/}\PYG{n}{crespi}\PYG{o}{/}\PYG{n}{z}\PYG{+w}{ }\PYG{n}{cutoff}
\end{sphinxVerbatim}


\subsection{Examples}
\label{\detokenize{pair_kolmogorov_crespi_z:examples}}
\begin{sphinxVerbatim}[commandchars=\\\{\}]
\PYG{k}{pair\PYGZus{}style}\PYG{+w}{ }\PYG{n}{hybrid}\PYG{o}{/}\PYG{n}{overlay}\PYG{+w}{ }\PYG{n}{kolmogorov}\PYG{o}{/}\PYG{n}{crespi}\PYG{o}{/}\PYG{n}{z}\PYG{+w}{ }\PYG{l+m}{20.0}
\PYG{k}{pair\PYGZus{}coeff}\PYG{+w}{ }\PYG{o}{*}\PYG{+w}{ }\PYG{o}{*}\PYG{+w}{ }\PYG{n}{none}
\PYG{k}{pair\PYGZus{}coeff}\PYG{+w}{ }\PYG{l+m}{1}\PYG{+w}{ }\PYG{l+m}{2}\PYG{+w}{ }\PYG{n}{kolmogorov}\PYG{o}{/}\PYG{n}{crespi}\PYG{o}{/}\PYG{n}{z}\PYG{+w}{  }\PYG{n}{CC.KC}\PYG{+w}{   }\PYG{n}{C}\PYG{+w}{ }\PYG{n}{C}

\PYG{k}{pair\PYGZus{}style}\PYG{+w}{ }\PYG{n}{hybrid}\PYG{o}{/}\PYG{n}{overlay}\PYG{+w}{ }\PYG{n}{rebo}\PYG{+w}{ }\PYG{n}{kolmogorov}\PYG{o}{/}\PYG{n}{crespi}\PYG{o}{/}\PYG{n}{z}\PYG{+w}{ }\PYG{l+m}{14.0}
\PYG{k}{pair\PYGZus{}coeff}\PYG{+w}{ }\PYG{o}{*}\PYG{+w}{ }\PYG{o}{*}\PYG{+w}{ }\PYG{n}{rebo}\PYG{+w}{                 }\PYG{n}{CH.rebo}\PYG{+w}{    }\PYG{n}{C}\PYG{+w}{ }\PYG{n}{C}
\PYG{k}{pair\PYGZus{}coeff}\PYG{+w}{ }\PYG{l+m}{1}\PYG{+w}{ }\PYG{l+m}{2}\PYG{+w}{ }\PYG{n}{kolmogorov}\PYG{o}{/}\PYG{n}{crespi}\PYG{o}{/}\PYG{n}{z}\PYG{+w}{  }\PYG{n}{CC.KC}\PYG{+w}{      }\PYG{n}{C}\PYG{+w}{ }\PYG{n}{C}
\end{sphinxVerbatim}


\subsection{Description}
\label{\detokenize{pair_kolmogorov_crespi_z:description}}
\sphinxAtStartPar
The \sphinxstyleemphasis{kolmogorov/crespi/z} style computes the Kolmogorov\sphinxhyphen{}Crespi interaction
potential as described in {\hyperref[\detokenize{pair_kolmogorov_crespi_z:kc05}]{\sphinxcrossref{\DUrole{std,std-ref}{(Kolmogorov)}}}}. An important simplification is made,
which is to take all normals along the z\sphinxhyphen{}axis.
\begin{equation*}
\begin{split}E  = & \frac{1}{2} \sum_i \sum_{j \neq i} V_{ij} \\
V_{ij}  = & e^{-\lambda(r_{ij} -z_0)} \left[ C + f(\rho_{ij}) + f(\rho_{ji}) \right] - A \left( \frac{r_{ij}}{z_0}\right)^{-6} + A \left( \frac{\textrm{cutoff}}{z_0}\right)^{-6} \\
\rho_{ij}^2 = & \rho_{ji}^2  =  x_{ij}^2 + y_{ij}^2 \qquad \qquad (\mathbf{n}_i \equiv \mathbf{\hat{z}}) \\
f(\rho)  = &  e^{-(\rho/\delta)^2} \sum_{n=0}^2 C_{2n} \left( \rho/\delta \right)^{2n}\end{split}
\end{equation*}
\sphinxAtStartPar
It is important to have a sufficiently large cutoff to ensure smooth forces.
Energies are shifted so that they go continuously to zero at the cutoff assuming
that the exponential part of \(V_{ij}\) (first term) decays sufficiently fast.
This shift is achieved by the last term in the equation for \(V_{ij}\) above.

\sphinxAtStartPar
This potential is intended for interactions between two layers of graphene.
Therefore, to avoid interaction between layers in multi\sphinxhyphen{}layered materials,
each layer should have a separate atom type and interactions should only
be computed between atom types of neighboring layers.

\sphinxAtStartPar
The parameter file (e.g. CC.KC), is intended for use with metal
{\hyperref[\detokenize{units::doc}]{\sphinxcrossref{\DUrole{doc}{units}}}}, with energies in meV. An additional parameter, \sphinxstyleemphasis{S},
is available to facilitate scaling of energies in accordance with
{\hyperref[\detokenize{pair_kolmogorov_crespi_z:vanwijk}]{\sphinxcrossref{\DUrole{std,std-ref}{(vanWijk)}}}}.

\sphinxAtStartPar
This potential must be used in combination with hybrid/overlay.
Other interactions can be set to zero using pair\_style \sphinxstyleemphasis{none}.


\subsection{Restrictions}
\label{\detokenize{pair_kolmogorov_crespi_z:restrictions}}
\sphinxAtStartPar
This fix is part of the INTERLAYER package.  It is only enabled if
LAMMPS was built with that package.  See the {\hyperref[\detokenize{Build_package::doc}]{\sphinxcrossref{\DUrole{doc}{Build package}}}} page for more info.


\subsection{Related commands}
\label{\detokenize{pair_kolmogorov_crespi_z:related-commands}}
\sphinxAtStartPar
{\hyperref[\detokenize{pair_coeff::doc}]{\sphinxcrossref{\DUrole{doc}{pair\_coeff}}}},
{\hyperref[\detokenize{pair_none::doc}]{\sphinxcrossref{\DUrole{doc}{pair\_none}}}},
{\hyperref[\detokenize{pair_hybrid::doc}]{\sphinxcrossref{\DUrole{doc}{pair\_style hybrid/overlay}}}},
{\hyperref[\detokenize{pair_drip::doc}]{\sphinxcrossref{\DUrole{doc}{pair\_style drip}}}},
{\hyperref[\detokenize{pair_ilp_graphene_hbn::doc}]{\sphinxcrossref{\DUrole{doc}{pair\_style ilp/graphene/hbn}}}}.
{\hyperref[\detokenize{pair_kolmogorov_crespi_full::doc}]{\sphinxcrossref{\DUrole{doc}{pair\_style kolmogorov/crespi/full}}}},
{\hyperref[\detokenize{pair_lebedeva_z::doc}]{\sphinxcrossref{\DUrole{doc}{pair\_style lebedeva/z}}}}


\subsection{Default}
\label{\detokenize{pair_kolmogorov_crespi_z:default}}
\sphinxAtStartPar
none


\bigskip\hrule\bigskip

\phantomsection\label{\detokenize{pair_kolmogorov_crespi_z:kc05}}
\sphinxAtStartPar
\sphinxstylestrong{(Kolmogorov)} A. N. Kolmogorov, V. H. Crespi, Phys. Rev. B 71, 235415 (2005)

\phantomsection\label{\detokenize{pair_kolmogorov_crespi_z:vanwijk}}
\sphinxAtStartPar
\sphinxstylestrong{(vanWijk)} M. M. van Wijk, A. Schuring, M. I. Katsnelson, and A. Fasolino, Physical Review Letters, 113, 135504 (2014)

\sphinxstepscope

\index{pair\_style lcbop@\spxentry{pair\_style lcbop}}\ignorespaces 

\section{pair\_style lcbop command}
\label{\detokenize{pair_lcbop:pair-style-lcbop-command}}\label{\detokenize{pair_lcbop:index-0}}\label{\detokenize{pair_lcbop::doc}}

\subsection{Syntax}
\label{\detokenize{pair_lcbop:syntax}}
\begin{sphinxVerbatim}[commandchars=\\\{\}]
\PYG{k}{pair\PYGZus{}style}\PYG{+w}{ }\PYG{n}{lcbop}
\end{sphinxVerbatim}


\subsection{Examples}
\label{\detokenize{pair_lcbop:examples}}
\begin{sphinxVerbatim}[commandchars=\\\{\}]
\PYG{k}{pair\PYGZus{}style}\PYG{+w}{ }\PYG{n}{lcbop}
\PYG{k}{pair\PYGZus{}coeff}\PYG{+w}{ }\PYG{o}{*}\PYG{+w}{ }\PYG{o}{*}\PYG{+w}{ }\PYG{n}{..}\PYG{o}{/}\PYG{n}{potentials}\PYG{o}{/}\PYG{n}{C.lcbop}\PYG{+w}{ }\PYG{n}{C}
\end{sphinxVerbatim}


\subsection{Description}
\label{\detokenize{pair_lcbop:description}}
\sphinxAtStartPar
The \sphinxstyleemphasis{lcbop} pair style computes the long\sphinxhyphen{}range bond\sphinxhyphen{}order potential
for carbon (LCBOP) of {\hyperref[\detokenize{pair_lcbop:los}]{\sphinxcrossref{\DUrole{std,std-ref}{(Los and Fasolino)}}}}.  See section II in
that paper for the analytic equations associated with the potential.

\sphinxAtStartPar
Only a single pair\_coeff command is used with the \sphinxstyleemphasis{lcbop} style which
specifies an LCBOP potential file with parameters for specific
elements.  These are mapped to LAMMPS atom types by specifying N
additional arguments after the filename in the pair\_coeff command,
where N is the number of LAMMPS atom types:
\begin{itemize}
\item {} 
\sphinxAtStartPar
filename

\item {} 
\sphinxAtStartPar
N element names = mapping of LCBOP elements to atom types

\end{itemize}

\sphinxAtStartPar
See the {\hyperref[\detokenize{pair_coeff::doc}]{\sphinxcrossref{\DUrole{doc}{pair\_coeff}}}} page for alternate ways
to specify the path for the potential file.

\sphinxAtStartPar
As an example, if your LAMMPS simulation has 4 atom types and you want
the first 3 to be C you would use the following pair\_coeff command:

\begin{sphinxVerbatim}[commandchars=\\\{\}]
\PYG{k}{pair\PYGZus{}coeff}\PYG{+w}{ }\PYG{o}{*}\PYG{+w}{ }\PYG{o}{*}\PYG{+w}{ }\PYG{n}{C.lcbop}\PYG{+w}{ }\PYG{n}{C}\PYG{+w}{ }\PYG{n}{C}\PYG{+w}{ }\PYG{n}{C}\PYG{+w}{ }\PYG{n}{NULL}
\end{sphinxVerbatim}

\sphinxAtStartPar
The first 2 arguments must be * * so as to span all LAMMPS atom types.
The first C argument maps LAMMPS atom type 1 to the C element in the
LCBOP file. If a mapping value is specified as NULL, the mapping is
not performed.  This can be used when a \sphinxstyleemphasis{lcbop} potential is used as
part of the \sphinxstyleemphasis{hybrid} pair style.  The NULL values are placeholders for
atom types that will be used with other potentials.

\sphinxAtStartPar
The parameters/coefficients for the LCBOP potential as applied to C
are listed in the C.lcbop file to agree with the original {\hyperref[\detokenize{pair_lcbop:los}]{\sphinxcrossref{\DUrole{std,std-ref}{(Los and Fasolino)}}}} paper.  Thus the parameters are specific to this
potential and the way it was fit, so modifying the file should be done
carefully.


\bigskip\hrule\bigskip



\subsection{Mixing, shift, table, tail correction, restart, rRESPA info}
\label{\detokenize{pair_lcbop:mixing-shift-table-tail-correction-restart-rrespa-info}}
\sphinxAtStartPar
This pair style does not support the {\hyperref[\detokenize{pair_modify::doc}]{\sphinxcrossref{\DUrole{doc}{pair\_modify}}}}
mix, shift, table, and tail options.

\sphinxAtStartPar
This pair style does not write its information to {\hyperref[\detokenize{restart::doc}]{\sphinxcrossref{\DUrole{doc}{binary restart files}}}}, since it is stored in potential files.  Thus, you
need to re\sphinxhyphen{}specify the pair\_style and pair\_coeff commands in an input
script that reads a restart file.

\sphinxAtStartPar
This pair style can only be used via the \sphinxstyleemphasis{pair} keyword of the
{\hyperref[\detokenize{run_style::doc}]{\sphinxcrossref{\DUrole{doc}{run\_style respa}}}} command.  It does not support the
\sphinxstyleemphasis{inner}, \sphinxstyleemphasis{middle}, \sphinxstyleemphasis{outer} keywords.


\subsection{Restrictions}
\label{\detokenize{pair_lcbop:restrictions}}
\sphinxAtStartPar
This pair style is part of the MANYBODY package.  It is only enabled
if LAMMPS was built with that package.
See the {\hyperref[\detokenize{Build_package::doc}]{\sphinxcrossref{\DUrole{doc}{Build package}}}} page for more info.

\sphinxAtStartPar
This pair potential requires the {\hyperref[\detokenize{newton::doc}]{\sphinxcrossref{\DUrole{doc}{newton}}}} setting to be
“on” for pair interactions.

\sphinxAtStartPar
The \sphinxcode{\sphinxupquote{C.lcbop}} potential file provided with LAMMPS (see the potentials
directory) is parameterized for {\hyperref[\detokenize{units::doc}]{\sphinxcrossref{\DUrole{doc}{metal units}}}}.  You can use
the LCBOP potential with any LAMMPS units, but you would need to
create your own LCBOP potential file with coefficients listed in the
appropriate units if your simulation does not use “metal” units.


\subsection{Related commands}
\label{\detokenize{pair_lcbop:related-commands}}
\sphinxAtStartPar
{\hyperref[\detokenize{pair_airebo::doc}]{\sphinxcrossref{\DUrole{doc}{pair\_airebo}}}}, {\hyperref[\detokenize{pair_coeff::doc}]{\sphinxcrossref{\DUrole{doc}{pair\_coeff}}}}


\subsection{Default}
\label{\detokenize{pair_lcbop:default}}
\sphinxAtStartPar
none


\bigskip\hrule\bigskip

\phantomsection\label{\detokenize{pair_lcbop:los}}
\sphinxAtStartPar
\sphinxstylestrong{(Los and Fasolino)} J. H. Los and A. Fasolino, Phys. Rev. B 68, 024107
(2003).

\sphinxstepscope

\index{pair\_style lebedeva/z@\spxentry{pair\_style lebedeva/z}}\ignorespaces 

\section{pair\_style lebedeva/z command}
\label{\detokenize{pair_lebedeva_z:pair-style-lebedeva-z-command}}\label{\detokenize{pair_lebedeva_z:index-0}}\label{\detokenize{pair_lebedeva_z::doc}}

\subsection{Syntax}
\label{\detokenize{pair_lebedeva_z:syntax}}
\begin{sphinxVerbatim}[commandchars=\\\{\}]
\PYG{k}{pair\PYGZus{}style}\PYG{+w}{ }\PYG{n}{[hybrid}\PYG{o}{/}\PYG{n}{overlay}\PYG{+w}{ }\PYG{n}{...]}\PYG{+w}{ }\PYG{n}{lebedeva}\PYG{o}{/}\PYG{n}{z}\PYG{+w}{ }\PYG{n}{cutoff}
\end{sphinxVerbatim}


\subsection{Examples}
\label{\detokenize{pair_lebedeva_z:examples}}
\begin{sphinxVerbatim}[commandchars=\\\{\}]
\PYG{k}{pair\PYGZus{}style}\PYG{+w}{ }\PYG{n}{hybrid}\PYG{o}{/}\PYG{n}{overlay}\PYG{+w}{ }\PYG{n}{lebedeva}\PYG{o}{/}\PYG{n}{z}\PYG{+w}{ }\PYG{l+m}{20.0}
\PYG{k}{pair\PYGZus{}coeff}\PYG{+w}{ }\PYG{o}{*}\PYG{+w}{ }\PYG{o}{*}\PYG{+w}{ }\PYG{n}{none}
\PYG{k}{pair\PYGZus{}coeff}\PYG{+w}{ }\PYG{l+m}{1}\PYG{+w}{ }\PYG{l+m}{2}\PYG{+w}{ }\PYG{n}{lebedeva}\PYG{o}{/}\PYG{n}{z}\PYG{+w}{  }\PYG{n}{CC.Lebedeva}\PYG{+w}{   }\PYG{n}{C}\PYG{+w}{ }\PYG{n}{C}

\PYG{k}{pair\PYGZus{}style}\PYG{+w}{ }\PYG{n}{hybrid}\PYG{o}{/}\PYG{n}{overlay}\PYG{+w}{ }\PYG{n}{rebo}\PYG{+w}{ }\PYG{n}{lebedeva}\PYG{o}{/}\PYG{n}{z}\PYG{+w}{ }\PYG{l+m}{14.0}
\PYG{k}{pair\PYGZus{}coeff}\PYG{+w}{ }\PYG{o}{*}\PYG{+w}{ }\PYG{o}{*}\PYG{+w}{ }\PYG{n}{rebo}\PYG{+w}{        }\PYG{n}{CH.rebo}\PYG{+w}{       }\PYG{n}{C}\PYG{+w}{ }\PYG{n}{C}
\PYG{k}{pair\PYGZus{}coeff}\PYG{+w}{ }\PYG{l+m}{1}\PYG{+w}{ }\PYG{l+m}{2}\PYG{+w}{ }\PYG{n}{lebedeva}\PYG{o}{/}\PYG{n}{z}\PYG{+w}{  }\PYG{n}{CC.Lebedeva}\PYG{+w}{   }\PYG{n}{C}\PYG{+w}{ }\PYG{n}{C}
\end{sphinxVerbatim}


\subsection{Description}
\label{\detokenize{pair_lebedeva_z:description}}
\sphinxAtStartPar
The \sphinxstyleemphasis{lebedeva/z} pair style computes the Lebedeva interaction potential
as described in {\hyperref[\detokenize{pair_lebedeva_z:leb01}]{\sphinxcrossref{\DUrole{std,std-ref}{(Lebedeva1)}}}} and {\hyperref[\detokenize{pair_lebedeva_z:leb02}]{\sphinxcrossref{\DUrole{std,std-ref}{(Lebedeva2)}}}}.  An important simplification is made, which is to take all
normals along the z\sphinxhyphen{}axis.

\sphinxAtStartPar
The Lebedeva potential is intended for the description of the interlayer
interaction between graphene layers.  To perform a realistic simulation,
this potential must be used in combination with an intralayer potential
such as {\hyperref[\detokenize{pair_airebo::doc}]{\sphinxcrossref{\DUrole{doc}{AIREBO}}}} or {\hyperref[\detokenize{pair_tersoff::doc}]{\sphinxcrossref{\DUrole{doc}{Tersoff}}}}
facilitated by using pair style {\hyperref[\detokenize{pair_hybrid::doc}]{\sphinxcrossref{\DUrole{doc}{hybrid/overlay}}}}.  To
keep the intralayer properties unaffected, the interlayer interaction
within the same layers should be avoided.  This can be achieved by
assigning different atom types to atoms of different layers (e.g. 1 and
2 in the examples above).

\sphinxAtStartPar
Other interactions can be set to zero using pair\_style \sphinxstyleemphasis{none}.
\begin{equation*}
\begin{split}E       = & \frac{1}{2} \sum_i \sum_{j \neq i} V_{ij}\\
V_{ij}  = & B e^{-\alpha(r_{ij} - z_0)} \\
          & + C(1 + D_1\rho^2_{ij} + D_2\rho^4_{ij}) e^{-\lambda_1\rho^2_{ij}} e^{-\lambda_2 (z^2_{ij} - z^2_0)} \\
          & - A \left(\frac{z_0}{r_ij}\right)^6 + A \left( \frac{z_0}{r_c} \right)^6 \\
\rho^2_{ij} = & x^2_{ij} + y^2_{ij} \qquad (\mathbf{n_i} \equiv \mathbf{\hat{z}})\end{split}
\end{equation*}
\sphinxAtStartPar
It is important to have a sufficiently large cutoff to ensure smooth forces.
Energies are shifted so that they go continuously to zero at the cutoff assuming
that the exponential part of \(V_{ij}\) (first term) decays sufficiently fast.
This shift is achieved by the last term in the equation for \(V_{ij}\) above.

\sphinxAtStartPar
The provided parameter file (CC.Lebedeva) contains two sets of parameters.
\begin{itemize}
\item {} 
\sphinxAtStartPar
The first set (element name “C”) is suitable for normal conditions and
is taken from {\hyperref[\detokenize{pair_lebedeva_z:popov}]{\sphinxcrossref{\DUrole{std,std-ref}{(Popov1)}}}}

\item {} 
\sphinxAtStartPar
The second set (element name “C1”) is suitable for high\sphinxhyphen{}pressure
conditions and is taken from {\hyperref[\detokenize{pair_lebedeva_z:koziol}]{\sphinxcrossref{\DUrole{std,std-ref}{(Koziol1)}}}}

\end{itemize}

\sphinxAtStartPar
Both sets contain an additional parameter, \sphinxstyleemphasis{S}, that can be used to
facilitate scaling of energies and is set to 1.0 by default.


\subsection{Restrictions}
\label{\detokenize{pair_lebedeva_z:restrictions}}
\sphinxAtStartPar
This pair style is part of the INTERLAYER package.  It is only enabled
if LAMMPS was built with that package.  See the {\hyperref[\detokenize{Build_package::doc}]{\sphinxcrossref{\DUrole{doc}{Build package}}}} page for more info.


\subsection{Related commands}
\label{\detokenize{pair_lebedeva_z:related-commands}}
\sphinxAtStartPar
{\hyperref[\detokenize{pair_coeff::doc}]{\sphinxcrossref{\DUrole{doc}{pair\_coeff}}}},
{\hyperref[\detokenize{pair_none::doc}]{\sphinxcrossref{\DUrole{doc}{pair\_style none}}}},
{\hyperref[\detokenize{pair_hybrid::doc}]{\sphinxcrossref{\DUrole{doc}{pair\_style hybrid/overlay}}}},
{\hyperref[\detokenize{pair_drip::doc}]{\sphinxcrossref{\DUrole{doc}{pair\_style drip}}}},
{\hyperref[\detokenize{pair_ilp_graphene_hbn::doc}]{\sphinxcrossref{\DUrole{doc}{pair\_style ilp/graphene/hbd}}}},
{\hyperref[\detokenize{pair_kolmogorov_crespi_z::doc}]{\sphinxcrossref{\DUrole{doc}{pair\_style kolmogorov/crespi/z}}}},
{\hyperref[\detokenize{pair_kolmogorov_crespi_full::doc}]{\sphinxcrossref{\DUrole{doc}{pair\_style kolmogorov/crespi/full}}}}.


\subsection{Default}
\label{\detokenize{pair_lebedeva_z:default}}
\sphinxAtStartPar
none


\bigskip\hrule\bigskip

\phantomsection\label{\detokenize{pair_lebedeva_z:leb01}}
\sphinxAtStartPar
\sphinxstylestrong{(Lebedeva1)} I. V. Lebedeva, A. A. Knizhnik, A. M. Popov, Y. E. Lozovik, B. V. Potapkin, Phys. Rev. B, 84, 245437 (2011)

\phantomsection\label{\detokenize{pair_lebedeva_z:leb02}}
\sphinxAtStartPar
\sphinxstylestrong{(Lebedeva2)} I. V. Lebedeva, A. A. Knizhnik, A. M. Popov, Y. E. Lozovik, B. V. Potapkin, Physica E: 44, 949\sphinxhyphen{}954 (2012)

\phantomsection\label{\detokenize{pair_lebedeva_z:popov}}
\sphinxAtStartPar
\sphinxstylestrong{(Popov1)} A.M. Popov, I. V. Lebedeva, A. A. Knizhnik, Y. E. Lozovik and B. V. Potapkin, Chem. Phys. Lett. 536, 82\sphinxhyphen{}86 (2012).

\phantomsection\label{\detokenize{pair_lebedeva_z:koziol}}
\sphinxAtStartPar
\sphinxstylestrong{(Koziol1)} Z. Koziol, G. Gawlik and J. Jagielski, Chinese Phys. B 28, 096101 (2019).

\sphinxstepscope

\index{pair\_style lepton@\spxentry{pair\_style lepton}}\ignorespaces 
\index{pair\_style lepton/omp@\spxentry{pair\_style lepton/omp}}\ignorespaces 
\index{pair\_style lepton/coul@\spxentry{pair\_style lepton/coul}}\ignorespaces 
\index{pair\_style lepton/coul/omp@\spxentry{pair\_style lepton/coul/omp}}\ignorespaces 
\index{pair\_style lepton/sphere@\spxentry{pair\_style lepton/sphere}}\ignorespaces 
\index{pair\_style lepton/sphere/omp@\spxentry{pair\_style lepton/sphere/omp}}\ignorespaces 

\section{pair\_style lepton command}
\label{\detokenize{pair_lepton:pair-style-lepton-command}}\label{\detokenize{pair_lepton:index-5}}\label{\detokenize{pair_lepton:index-4}}\label{\detokenize{pair_lepton:index-3}}\label{\detokenize{pair_lepton:index-2}}\label{\detokenize{pair_lepton:index-1}}\label{\detokenize{pair_lepton:index-0}}\label{\detokenize{pair_lepton::doc}}
\sphinxAtStartPar
Accelerator Variants: \sphinxstyleemphasis{lepton/omp}, \sphinxstyleemphasis{lepton/coul/comp}, \sphinxstyleemphasis{lepton/sphere/comp}


\subsection{Syntax}
\label{\detokenize{pair_lepton:syntax}}
\begin{sphinxVerbatim}[commandchars=\\\{\}]
\PYG{k}{pair\PYGZus{}style}\PYG{+w}{ }\PYG{n}{style}\PYG{+w}{ }\PYG{n}{args}
\end{sphinxVerbatim}
\begin{itemize}
\item {} 
\sphinxAtStartPar
style = \sphinxstyleemphasis{lepton} or \sphinxstyleemphasis{lepton/coul} or \sphinxstyleemphasis{lepton/sphere}

\item {} 
\sphinxAtStartPar
args = list of arguments for a particular style

\end{itemize}
\begin{sphinxalltt}
\sphinxstyleemphasis{lepton} args = cutoff
  cutoff = global cutoff for the interactions (distance units)
\sphinxstyleemphasis{lepton/coul} args = cutoff keyword
  cutoff = global cutoff for the interactions (distance units)
  zero or more keywords may be appended
  keyword = \sphinxstyleemphasis{ewald} or \sphinxstyleemphasis{pppm} or \sphinxstyleemphasis{msm} or \sphinxstyleemphasis{dispersion} or \sphinxstyleemphasis{tip4p}
\sphinxstyleemphasis{lepton/sphere} args = cutoff
  cutoff = global cutoff for the interactions (distance units)
\end{sphinxalltt}


\subsection{Examples}
\label{\detokenize{pair_lepton:examples}}
\begin{sphinxVerbatim}[commandchars=\\\{\}]
\PYG{k}{pair\PYGZus{}style}\PYG{+w}{ }\PYG{n}{lepton}\PYG{+w}{ }\PYG{l+m}{2.5}

\PYG{k}{pair\PYGZus{}coeff}\PYG{+w}{  }\PYG{o}{*}\PYG{+w}{ }\PYG{o}{*}\PYG{+w}{  }\PYG{l+s}{\PYGZdq{}}\PYG{l+s}{k*((r\PYGZhy{}r0)\PYGZca{}2*step(r0\PYGZhy{}r)); k=200; r0=1.5}\PYG{l+s}{\PYGZdq{}}\PYG{+w}{ }\PYG{l+m}{2.0}
\PYG{k}{pair\PYGZus{}coeff}\PYG{+w}{  }\PYG{l+m}{1}\PYG{+w}{ }\PYG{l+m}{2}\PYG{+w}{  }\PYG{l+s}{\PYGZdq{}}\PYG{l+s}{4.0*eps*((sig/r)\PYGZca{}12 \PYGZhy{} (sig/r)\PYGZca{}6);eps=1.0;sig=1.0}\PYG{l+s}{\PYGZdq{}}\PYG{+w}{ }\PYG{l+m}{1.12246204830937}
\PYG{k}{pair\PYGZus{}coeff}\PYG{+w}{  }\PYG{l+m}{2}\PYG{+w}{ }\PYG{l+m}{2}\PYG{+w}{  }\PYG{l+s}{\PYGZdq{}}\PYG{l+s}{eps*(2.0*(sig/r)\PYGZca{}9 \PYGZhy{} 3.0*(sig/r)\PYGZca{}6);eps=1.0;sig=1.0}\PYG{l+s}{\PYGZdq{}}
\PYG{k}{pair\PYGZus{}coeff}\PYG{+w}{  }\PYG{l+m}{1}\PYG{+w}{ }\PYG{l+m}{3}\PYG{+w}{  }\PYG{l+s}{\PYGZdq{}}\PYG{l+s}{zbl(13,6,r)}\PYG{l+s}{\PYGZdq{}}
\PYG{k}{pair\PYGZus{}coeff}\PYG{+w}{  }\PYG{l+m}{3}\PYG{+w}{ }\PYG{l+m}{3}\PYG{+w}{  }\PYG{l+s}{\PYGZdq{}}\PYG{l+s}{(1.0\PYGZhy{}switch)*zbl(6,6,r)\PYGZhy{}switch*4.0*eps*((sig/r)\PYGZca{}6);switch=0.5*(tanh(10.0*(r\PYGZhy{}sig))+1.0);eps=0.05;sig=3.20723}\PYG{l+s}{\PYGZdq{}}

\PYG{k}{pair\PYGZus{}style}\PYG{+w}{ }\PYG{n}{lepton}\PYG{o}{/}\PYG{n}{coul}\PYG{+w}{ }\PYG{l+m}{2.5}
\PYG{k}{pair\PYGZus{}coeff}\PYG{+w}{ }\PYG{l+m}{1}\PYG{+w}{ }\PYG{l+m}{1}\PYG{+w}{ }\PYG{l+s}{\PYGZdq{}}\PYG{l+s}{qi*qj/r}\PYG{l+s}{\PYGZdq{}}\PYG{+w}{ }\PYG{l+m}{4.0}
\PYG{k}{pair\PYGZus{}coeff}\PYG{+w}{ }\PYG{l+m}{1}\PYG{+w}{ }\PYG{l+m}{2}\PYG{+w}{ }\PYG{l+s}{\PYGZdq{}}\PYG{l+s}{lj+coul; lj=4.0*eps*((sig/r)\PYGZca{}12 \PYGZhy{} (sig/r)\PYGZca{}6); eps=1.0; sig=1.0; coul=qi*qj/r}\PYG{l+s}{\PYGZdq{}}

\PYG{k}{pair\PYGZus{}style}\PYG{+w}{ }\PYG{n}{lepton}\PYG{o}{/}\PYG{n}{coul}\PYG{+w}{ }\PYG{l+m}{2.5}\PYG{+w}{ }\PYG{n}{pppm}
\PYG{k}{kspace\PYGZus{}style}\PYG{+w}{ }\PYG{n}{pppm}\PYG{+w}{ }\PYG{l+m}{1.0e\PYGZhy{}4}
\PYG{k}{pair\PYGZus{}coeff}\PYG{+w}{ }\PYG{l+m}{1}\PYG{+w}{ }\PYG{l+m}{1}\PYG{+w}{ }\PYG{l+s}{\PYGZdq{}}\PYG{l+s}{qi*qj/r*erfc(alpha*r); alpha=1.067}\PYG{l+s}{\PYGZdq{}}

\PYG{k}{pair\PYGZus{}style}\PYG{+w}{ }\PYG{n}{lepton}\PYG{o}{/}\PYG{n}{sphere}\PYG{+w}{ }\PYG{l+m}{2.5}
\PYG{k}{pair\PYGZus{}coeff}\PYG{+w}{ }\PYG{l+m}{1}\PYG{+w}{ }\PYG{o}{*}\PYG{+w}{ }\PYG{l+s}{\PYGZdq{}}\PYG{l+s}{k*((r\PYGZhy{}r0)\PYGZca{}2*step(r0\PYGZhy{}r)); k=200; r0=radi+radj}\PYG{l+s}{\PYGZdq{}}
\PYG{k}{pair\PYGZus{}coeff}\PYG{+w}{ }\PYG{l+m}{2}\PYG{+w}{ }\PYG{l+m}{2}\PYG{+w}{ }\PYG{l+s}{\PYGZdq{}}\PYG{l+s}{4.0*eps*((sig/r)\PYGZca{}12 \PYGZhy{} (sig/r)\PYGZca{}6); eps=1.0; sig=2.0*sqrt(radi*radj)}\PYG{l+s}{\PYGZdq{}}
\end{sphinxVerbatim}


\subsection{Description}
\label{\detokenize{pair_lepton:description}}
\sphinxAtStartPar
\DUrole{versionmodified,added}{Added in version 8Feb2023: }added pair styles \sphinxstyleemphasis{lepton} and \sphinxstyleemphasis{lepton/coul}

\sphinxAtStartPar
\DUrole{versionmodified,changed}{Changed in version 15Jun2023: }added pair style \sphinxstyleemphasis{lepton/sphere}

\sphinxAtStartPar
Pair styles \sphinxstyleemphasis{lepton}, \sphinxstyleemphasis{lepton/coul}, \sphinxstyleemphasis{lepton/sphere} compute pairwise
interactions between particles which depend on the distance and have a
cutoff.  The potential function must be provided as an expression string
using “r” as the distance variable.  With pair style \sphinxstyleemphasis{lepton/coul} one
may additionally reference the charges of the two atoms of the pair with
“qi” and “qj”, respectively.  With pair style \sphinxstyleemphasis{lepton/sphere} one may
instead reference the radii of the two atoms of the pair with “radi” and
“radj”, respectively; this is half of the diameter that can be set in
{\hyperref[\detokenize{read_data::doc}]{\sphinxcrossref{\DUrole{doc}{data files}}}} or the {\hyperref[\detokenize{set::doc}]{\sphinxcrossref{\DUrole{doc}{set command}}}}.

\sphinxAtStartPar
Note that further constants in the expressions can be defined in the
same string as additional expressions separated by semicolons as shown
in the examples above.

\sphinxAtStartPar
The expression \sphinxtitleref{“200.0*(r\sphinxhyphen{}1.5)\textasciicircum{}2”} represents a harmonic potential
around the pairwise distance \(r_0\) of 1.5 distance units and a
force constant \sphinxstyleemphasis{K} of 200.0 energy units:
\begin{equation*}
\begin{split}U_{ij} = K (r-r_0)^2\end{split}
\end{equation*}
\sphinxAtStartPar
The expression \sphinxtitleref{“qi*qj/r”} represents a regular Coulombic potential with cutoff:
\begin{equation*}
\begin{split}U_{ij} = \frac{C q_i q_j}{\epsilon  r} \qquad r < r_c\end{split}
\end{equation*}
\sphinxAtStartPar
The expression \sphinxtitleref{“200.0*(r\sphinxhyphen{}(radi+radj)\textasciicircum{}2”} represents a harmonic potential
that has the equilibrium distance chosen so that the radii of the two
atoms touch:
\begin{equation*}
\begin{split}U_{ij} = K (r-(r_i+r_j))^2\end{split}
\end{equation*}
\sphinxAtStartPar
The \sphinxhref{https://simtk.org/projects/lepton}{Lepton library}, that the
\sphinxstyleemphasis{lepton} pair style interfaces with, evaluates this expression string at
run time to compute the pairwise energy.  It also creates an analytical
representation of the first derivative of this expression with respect
to “r” and then uses that to compute the force between the pairs of
particles within the given cutoff.

\sphinxAtStartPar
The following coefficients must be defined for each pair of atoms types
via the {\hyperref[\detokenize{pair_coeff::doc}]{\sphinxcrossref{\DUrole{doc}{pair\_coeff}}}} command as in the examples above,
or in the data file or restart files read by the {\hyperref[\detokenize{read_data::doc}]{\sphinxcrossref{\DUrole{doc}{read\_data}}}} or {\hyperref[\detokenize{read_restart::doc}]{\sphinxcrossref{\DUrole{doc}{read\_restart}}}} commands:
\begin{itemize}
\item {} 
\sphinxAtStartPar
Lepton expression (energy units)

\item {} 
\sphinxAtStartPar
cutoff (distance units)

\end{itemize}

\sphinxAtStartPar
The Lepton expression must be either enclosed in quotes or must not
contain any whitespace so that LAMMPS recognizes it as a single keyword.
More on valid Lepton expressions below.  The last coefficient is
optional; it allows to set the cutoff for a pair of atom types to a
different value than the global cutoff.

\sphinxAtStartPar
For pair style \sphinxstyleemphasis{lepton} only the “lj” values of the {\hyperref[\detokenize{special_bonds::doc}]{\sphinxcrossref{\DUrole{doc}{special\_bonds}}}} settings apply in case the interacting pair is also
connected with a bond.  The potential energy will \sphinxstyleemphasis{only} be added to the
“evdwl” property.

\sphinxAtStartPar
For pair style \sphinxstyleemphasis{lepton/coul} only the “coul” values of the
{\hyperref[\detokenize{special_bonds::doc}]{\sphinxcrossref{\DUrole{doc}{special\_bonds}}}} settings apply in case the
interacting pair is also connected with a bond.  The potential energy
will \sphinxstyleemphasis{only} be added to the “ecoul” property.

\sphinxAtStartPar
For pair style \sphinxstyleemphasis{lepton/sphere} only the “lj” values of the
{\hyperref[\detokenize{special_bonds::doc}]{\sphinxcrossref{\DUrole{doc}{special\_bonds}}}} settings apply in case the
interacting pair is also connected with a bond.  The potential energy
will \sphinxstyleemphasis{only} be added to the “evdwl” property.

\sphinxAtStartPar
In addition to the functions listed below, both pair styles support in
addition a custom “zbl(zi,zj,r)” function which computes the
Ziegler\sphinxhyphen{}Biersack\sphinxhyphen{}Littmark (ZBL) screened nuclear repulsion for
describing high\sphinxhyphen{}energy collisions between atoms.  For details of the
function please see the documentation for {\hyperref[\detokenize{pair_zbl::doc}]{\sphinxcrossref{\DUrole{doc}{pair style zbl}}}}. The arguments of the function are the atomic numbers of
atom i (zi), atom j (zj) and the distance r.  Please see the examples
above.


\bigskip\hrule\bigskip



\subsection{Lepton expression syntax and features}
\label{\detokenize{pair_lepton:lepton-expression-syntax-and-features}}
\sphinxAtStartPar
Lepton supports the following operators in expressions:


\begin{savenotes}\sphinxattablestart
\sphinxthistablewithglobalstyle
\centering
\begin{tabulary}{\linewidth}[t]{TTTTTTTTTTTTTT}
\sphinxtoprule
\sphinxtableatstartofbodyhook
\sphinxAtStartPar
+
&
\sphinxAtStartPar
Add
&&
\sphinxAtStartPar
\sphinxhyphen{}
&
\sphinxAtStartPar
Subtract
&&
\sphinxAtStartPar
*
&
\sphinxAtStartPar
Multiply
&&
\sphinxAtStartPar
/
&
\sphinxAtStartPar
Divide
&&
\sphinxAtStartPar
\textasciicircum{}
&
\sphinxAtStartPar
Power
\\
\sphinxbottomrule
\end{tabulary}
\sphinxtableafterendhook\par
\sphinxattableend\end{savenotes}

\sphinxAtStartPar
The following mathematical functions are available:


\begin{savenotes}\sphinxattablestart
\sphinxthistablewithglobalstyle
\centering
\begin{tabulary}{\linewidth}[t]{TTTT}
\sphinxtoprule
\sphinxtableatstartofbodyhook
\sphinxAtStartPar
sqrt(x)
&
\sphinxAtStartPar
Square root
&
\sphinxAtStartPar
exp(x)
&
\sphinxAtStartPar
Exponential
\\
\sphinxhline
\sphinxAtStartPar
log(x)
&
\sphinxAtStartPar
Natural logarithm
&
\sphinxAtStartPar
sin(x)
&
\sphinxAtStartPar
Sine (angle in radians)
\\
\sphinxhline
\sphinxAtStartPar
cos(x)
&
\sphinxAtStartPar
Cosine (angle in radians)
&
\sphinxAtStartPar
sec(x)
&
\sphinxAtStartPar
Secant (angle in radians)
\\
\sphinxhline
\sphinxAtStartPar
csc(x)
&
\sphinxAtStartPar
Cosecant (angle in radians)
&
\sphinxAtStartPar
tan(x)
&
\sphinxAtStartPar
Tangent (angle in radians)
\\
\sphinxhline
\sphinxAtStartPar
cot(x)
&
\sphinxAtStartPar
Cotangent (angle in radians)
&
\sphinxAtStartPar
asin(x)
&
\sphinxAtStartPar
Inverse sine (in radians)
\\
\sphinxhline
\sphinxAtStartPar
acos(x)
&
\sphinxAtStartPar
Inverse cosine (in radians)
&
\sphinxAtStartPar
atan(x)
&
\sphinxAtStartPar
Inverse tangent (in radians)
\\
\sphinxhline
\sphinxAtStartPar
sinh(x)
&
\sphinxAtStartPar
Hyperbolic sine
&
\sphinxAtStartPar
cosh(x)
&
\sphinxAtStartPar
Hyperbolic cosine
\\
\sphinxhline
\sphinxAtStartPar
tanh(x)
&
\sphinxAtStartPar
Hyperbolic tangent
&
\sphinxAtStartPar
erf(x)
&
\sphinxAtStartPar
Error function
\\
\sphinxhline
\sphinxAtStartPar
erfc(x)
&
\sphinxAtStartPar
Complementary Error function
&
\sphinxAtStartPar
abs(x)
&
\sphinxAtStartPar
Absolute value
\\
\sphinxhline
\sphinxAtStartPar
min(x,y)
&
\sphinxAtStartPar
Minimum of two values
&
\sphinxAtStartPar
max(x,y)
&
\sphinxAtStartPar
Maximum of two values
\\
\sphinxhline
\sphinxAtStartPar
delta(x)
&
\sphinxAtStartPar
delta(x) is 1 for \sphinxtitleref{x = 0}, otherwise 0
&
\sphinxAtStartPar
step(x)
&
\sphinxAtStartPar
step(x) is 0 for \sphinxtitleref{x \textless{} 0}, otherwise 1
\\
\sphinxbottomrule
\end{tabulary}
\sphinxtableafterendhook\par
\sphinxattableend\end{savenotes}

\sphinxAtStartPar
Numbers may be given in either decimal or exponential form.  All of the
following are valid numbers: \sphinxtitleref{5}, \sphinxtitleref{\sphinxhyphen{}3.1}, \sphinxtitleref{1e6}, and \sphinxtitleref{3.12e\sphinxhyphen{}2}.

\sphinxAtStartPar
As an extension to the standard Lepton syntax, it is also possible to
use LAMMPS {\hyperref[\detokenize{variable::doc}]{\sphinxcrossref{\DUrole{doc}{variables}}}} in the format “v\_name”.  Before
evaluating the expression, “v\_name” will be replaced with the value of
the variable “name”.  This is compatible with all kinds of scalar
variables, but not with vectors, arrays, local, or per\sphinxhyphen{}atom
variables.  If necessary, a custom scalar variable needs to be defined
that can access the desired (single) item from a non\sphinxhyphen{}scalar variable.
As an example, the following lines will instruct LAMMPS to ramp
the force constant for a harmonic bond from 100.0 to 200.0 during the
next run:

\begin{sphinxVerbatim}[commandchars=\\\{\}]
\PYG{k}{variable }\PYG{n+nv+nvIdentifier}{fconst}\PYG{+w}{ }\PYG{n}{equal}\PYG{+w}{ }\PYG{n}{ramp}\PYG{n+nv}{(}\PYG{n+nv}{100.0, 200}\PYG{n+nv}{)}
\PYG{k}{bond\PYGZus{}style}\PYG{+w}{ }\PYG{n}{lepton}
\PYG{k}{bond\PYGZus{}coeff}\PYG{+w}{ }\PYG{l+m}{1}\PYG{+w}{ }\PYG{l+m}{1.5}\PYG{+w}{ }\PYG{l+s}{\PYGZdq{}}\PYG{l+s}{v\PYGZus{}fconst * (r\PYGZca{}2)}\PYG{l+s}{\PYGZdq{}}
\end{sphinxVerbatim}

\sphinxAtStartPar
An expression may be followed by definitions for intermediate values that appear in the
expression. A semicolon “;” is used as a delimiter between value definitions. For example,
the expression:

\begin{sphinxVerbatim}[commandchars=\\\{\}]
\PYG{n}{a}\PYG{o}{\PYGZca{}}\PYG{l+m+mi}{2}\PYG{o}{+}\PYG{n}{a}\PYG{o}{*}\PYG{n}{b}\PYG{o}{+}\PYG{n}{b}\PYG{o}{\PYGZca{}}\PYG{l+m+mi}{2}\PYG{p}{;}\PYG{+w}{ }\PYG{n}{a}\PYG{o}{=}\PYG{n}{a1}\PYG{o}{+}\PYG{n}{a2}\PYG{p}{;}\PYG{+w}{ }\PYG{n}{b}\PYG{o}{=}\PYG{n}{b1}\PYG{o}{+}\PYG{n}{b2}
\end{sphinxVerbatim}

\sphinxAtStartPar
is exactly equivalent to

\begin{sphinxVerbatim}[commandchars=\\\{\}]
\PYG{p}{(}\PYG{n}{a1}\PYG{o}{+}\PYG{n}{a2}\PYG{p}{)}\PYG{o}{\PYGZca{}}\PYG{l+m+mi}{2}\PYG{o}{+}\PYG{p}{(}\PYG{n}{a1}\PYG{o}{+}\PYG{n}{a2}\PYG{p}{)}\PYG{o}{*}\PYG{p}{(}\PYG{n}{b1}\PYG{o}{+}\PYG{n}{b2}\PYG{p}{)}\PYG{o}{+}\PYG{p}{(}\PYG{n}{b1}\PYG{o}{+}\PYG{n}{b2}\PYG{p}{)}\PYG{o}{\PYGZca{}}\PYG{l+m+mi}{2}
\end{sphinxVerbatim}

\sphinxAtStartPar
The definition of an intermediate value may itself involve other
intermediate values. Whitespace and quotation characters (’\textquotesingle{}’ and ‘”’)
are ignored.  All uses of a value must appear \sphinxstyleemphasis{before} that value’s
definition.  For efficiency reasons, the expression string is parsed,
optimized, and then stored in an internal, pre\sphinxhyphen{}parsed representation for
evaluation.

\sphinxAtStartPar
Evaluating a Lepton expression is typically between 2.5 and 5 times
slower than the corresponding compiled and optimized C++ code.  If
additional speed or GPU acceleration (via GPU or KOKKOS) is required,
the interaction can be represented as a table.  Suitable table files
can be created either internally using the {\hyperref[\detokenize{pair_write::doc}]{\sphinxcrossref{\DUrole{doc}{pair\_write}}}}
or {\hyperref[\detokenize{bond_write::doc}]{\sphinxcrossref{\DUrole{doc}{bond\_write}}}} command or through the Python scripts
in the {\hyperref[\detokenize{Tools:tabulate}]{\sphinxcrossref{\DUrole{std,std-ref}{tools/tabulate}}}} folder.


\bigskip\hrule\bigskip


\sphinxAtStartPar
Styles with a \sphinxstyleemphasis{gpu}, \sphinxstyleemphasis{intel}, \sphinxstyleemphasis{kk}, \sphinxstyleemphasis{omp}, or \sphinxstyleemphasis{opt} suffix are
functionally the same as the corresponding style without the suffix.
They have been optimized to run faster, depending on your available
hardware, as discussed on the {\hyperref[\detokenize{Speed_packages::doc}]{\sphinxcrossref{\DUrole{doc}{Accelerator packages}}}}
page.  The accelerated styles take the same arguments and should
produce the same results, except for round\sphinxhyphen{}off and precision issues.

\sphinxAtStartPar
These accelerated styles are part of the GPU, INTEL, KOKKOS,
OPENMP, and OPT packages, respectively.  They are only enabled if
LAMMPS was built with those packages.  See the {\hyperref[\detokenize{Build_package::doc}]{\sphinxcrossref{\DUrole{doc}{Build package}}}} page for more info.

\sphinxAtStartPar
You can specify the accelerated styles explicitly in your input script
by including their suffix, or you can use the {\hyperref[\detokenize{Run_options::doc}]{\sphinxcrossref{\DUrole{doc}{\sphinxhyphen{}suffix command\sphinxhyphen{}line switch}}}} when you invoke LAMMPS, or you can use the
{\hyperref[\detokenize{suffix::doc}]{\sphinxcrossref{\DUrole{doc}{suffix}}}} command in your input script.

\sphinxAtStartPar
See the {\hyperref[\detokenize{Speed_packages::doc}]{\sphinxcrossref{\DUrole{doc}{Accelerator packages}}}} page for more
instructions on how to use the accelerated styles effectively.


\bigskip\hrule\bigskip



\subsection{Mixing, shift, table, tail correction, restart, rRESPA info}
\label{\detokenize{pair_lepton:mixing-shift-table-tail-correction-restart-rrespa-info}}
\sphinxAtStartPar
Pair styles \sphinxstyleemphasis{lepton}, \sphinxstyleemphasis{lepton/coul}, and \sphinxstyleemphasis{lepton/sphere} do not support
mixing.  Thus, expressions for \sphinxstyleemphasis{all} I,J pairs must be specified
explicitly.

\sphinxAtStartPar
Only pair style \sphinxstyleemphasis{lepton} supports the {\hyperref[\detokenize{pair_modify::doc}]{\sphinxcrossref{\DUrole{doc}{pair\_modify shift}}}}
option for shifting the potential energy of the pair interaction so that
it is 0 at the cutoff, pair styles \sphinxstyleemphasis{lepton/coul} and \sphinxstyleemphasis{lepton/sphere} do \sphinxstyleemphasis{not}.

\sphinxAtStartPar
The {\hyperref[\detokenize{pair_modify::doc}]{\sphinxcrossref{\DUrole{doc}{pair\_modify table}}}} options are not relevant for
the these pair styles.

\sphinxAtStartPar
These pair styles do not support the {\hyperref[\detokenize{pair_modify::doc}]{\sphinxcrossref{\DUrole{doc}{pair\_modify tail}}}} option for adding long\sphinxhyphen{}range tail corrections to energy
and pressure.

\sphinxAtStartPar
These pair styles write their information to {\hyperref[\detokenize{restart::doc}]{\sphinxcrossref{\DUrole{doc}{binary restart files}}}}, so pair\_style and pair\_coeff commands do not need to be
specified in an input script that reads a restart file.

\sphinxAtStartPar
These pair styles can only be used via the \sphinxstyleemphasis{pair} keyword of the
{\hyperref[\detokenize{run_style::doc}]{\sphinxcrossref{\DUrole{doc}{run\_style respa}}}} command.  They do not support the
\sphinxstyleemphasis{inner}, \sphinxstyleemphasis{middle}, \sphinxstyleemphasis{outer} keywords.


\bigskip\hrule\bigskip



\subsection{Restrictions}
\label{\detokenize{pair_lepton:restrictions}}
\sphinxAtStartPar
These pair styles are part of the LEPTON package and only enabled if
LAMMPS was built with this package.  See the {\hyperref[\detokenize{Build_package::doc}]{\sphinxcrossref{\DUrole{doc}{Build package}}}} page for more info.

\sphinxAtStartPar
Pair style \sphinxstyleemphasis{lepton/coul} requires that atom atoms have a charge
property, e.g. via {\hyperref[\detokenize{atom_style::doc}]{\sphinxcrossref{\DUrole{doc}{atom\_style charge}}}}.

\sphinxAtStartPar
Pair style \sphinxstyleemphasis{lepton/sphere} requires that atom atoms have a radius
property, e.g. via {\hyperref[\detokenize{atom_style::doc}]{\sphinxcrossref{\DUrole{doc}{atom\_style sphere}}}}.


\subsection{Related commands}
\label{\detokenize{pair_lepton:related-commands}}
\sphinxAtStartPar
{\hyperref[\detokenize{pair_coeff::doc}]{\sphinxcrossref{\DUrole{doc}{pair\_coeff}}}}, {\hyperref[\detokenize{pair_python::doc}]{\sphinxcrossref{\DUrole{doc}{pair\_style python}}}},
{\hyperref[\detokenize{pair_table::doc}]{\sphinxcrossref{\DUrole{doc}{pair\_style table}}}}, {\hyperref[\detokenize{pair_write::doc}]{\sphinxcrossref{\DUrole{doc}{pair\_write}}}}


\subsection{Default}
\label{\detokenize{pair_lepton:default}}
\sphinxAtStartPar
none

\sphinxstepscope

\index{pair\_style line/lj@\spxentry{pair\_style line/lj}}\ignorespaces 

\section{pair\_style line/lj command}
\label{\detokenize{pair_line_lj:pair-style-line-lj-command}}\label{\detokenize{pair_line_lj:index-0}}\label{\detokenize{pair_line_lj::doc}}

\subsection{Syntax}
\label{\detokenize{pair_line_lj:syntax}}
\begin{sphinxVerbatim}[commandchars=\\\{\}]
\PYG{k}{pair\PYGZus{}style}\PYG{+w}{ }\PYG{n}{line}\PYG{o}{/}\PYG{n}{lj}\PYG{+w}{ }\PYG{n}{cutoff}
\end{sphinxVerbatim}

\sphinxAtStartPar
cutoff = global cutoff for interactions (distance units)


\subsection{Examples}
\label{\detokenize{pair_line_lj:examples}}
\begin{sphinxVerbatim}[commandchars=\\\{\}]
\PYG{k}{pair\PYGZus{}style}\PYG{+w}{ }\PYG{n}{line}\PYG{o}{/}\PYG{n}{lj}\PYG{+w}{ }\PYG{l+m}{3.0}
\PYG{k}{pair\PYGZus{}coeff}\PYG{+w}{ }\PYG{o}{*}\PYG{+w}{ }\PYG{o}{*}\PYG{+w}{ }\PYG{l+m}{1.0}\PYG{+w}{ }\PYG{l+m}{1.0}\PYG{+w}{ }\PYG{l+m}{1.0}\PYG{+w}{ }\PYG{l+m}{0.8}\PYG{+w}{ }\PYG{l+m}{1.12}
\PYG{k}{pair\PYGZus{}coeff}\PYG{+w}{ }\PYG{l+m}{1}\PYG{+w}{ }\PYG{l+m}{2}\PYG{+w}{ }\PYG{l+m}{1.0}\PYG{+w}{ }\PYG{l+m}{2.0}\PYG{+w}{ }\PYG{l+m}{1.0}\PYG{+w}{ }\PYG{l+m}{1.5}\PYG{+w}{ }\PYG{l+m}{1.12}\PYG{+w}{ }\PYG{l+m}{5.0}
\PYG{k}{pair\PYGZus{}coeff}\PYG{+w}{ }\PYG{l+m}{1}\PYG{+w}{ }\PYG{l+m}{2}\PYG{+w}{ }\PYG{l+m}{1.0}\PYG{+w}{ }\PYG{l+m}{0.0}\PYG{+w}{ }\PYG{l+m}{1.0}\PYG{+w}{ }\PYG{l+m}{1.0}\PYG{+w}{ }\PYG{l+m}{2.5}
\end{sphinxVerbatim}


\subsection{Description}
\label{\detokenize{pair_line_lj:description}}
\sphinxAtStartPar
Style \sphinxstyleemphasis{line/lj} treats particles which are line segments as a set of
small spherical particles that tile the line segment length as explained
below.  Interactions between two line segments, each with N1 and N2
spherical particles, are calculated as the pairwise sum of N1*N2
Lennard\sphinxhyphen{}Jones interactions.  Interactions between a line segment with N
spherical particles and a point particle are treated as the pairwise sum
of N Lennard\sphinxhyphen{}Jones interactions.  See the {\hyperref[\detokenize{pair_lj::doc}]{\sphinxcrossref{\DUrole{doc}{pair\_style lj/cut}}}} page for the definition of Lennard\sphinxhyphen{}Jones interactions.

\sphinxAtStartPar
The set of non\sphinxhyphen{}overlapping spherical sub\sphinxhyphen{}particles that represent a
line segment are generated in the following manner.  Their size is a
function of the line segment length and the specified sub\sphinxhyphen{}particle
size for that particle type.  If a line segment has a length L and is
of type I, then the number of spheres N that represent the segment is
calculated as N = L/sizeI, rounded up to an integer value.  Thus if L
is not evenly divisible by sizeI, N is incremented to include one
extra sphere.  The centers of the spheres are spaced equally along the
line segment.  Imagine N+1 equally\sphinxhyphen{}space points, which include the 2
end points of the segment.  The sphere centers are halfway between
each pair of points.

\sphinxAtStartPar
The LJ interaction between 2 spheres on different line segments (or a
sphere on a line segment and a point particles) is computed with
sub\sphinxhyphen{}particle \(\epsilon\), \(\sigma\), and \sphinxstyleemphasis{cutoff} values that
are set by the pair\_coeff command, as described below.  If the distance
between the 2 spheres is greater than the sub\sphinxhyphen{}particle cutoff, there is
no interaction.  This means that some pairs of sub\sphinxhyphen{}particles on 2 line
segments may interact, but others may not.

\sphinxAtStartPar
For purposes of creating the neighbor list for pairs of interacting
line segments or lines/point particles, a regular particle\sphinxhyphen{}particle
cutoff is used, as defined by the \sphinxstyleemphasis{cutoff} setting above in the
pair\_style command or overridden with an optional argument in the
pair\_coeff command for a type pair as discussed below.  The distance
between the centers of 2 line segments, or the center of a line
segment and a point particle, must be less than this distance (plus
the neighbor skin; see the {\hyperref[\detokenize{neighbor::doc}]{\sphinxcrossref{\DUrole{doc}{neighbor}}}} command), for
the pair of particles to be included in the neighbor list.

\begin{sphinxadmonition}{note}{Note:}
\sphinxAtStartPar
This means that a too\sphinxhyphen{}short value for the \sphinxstyleemphasis{cutoff} setting can
exclude a pair of particles from the neighbor list even if pairs of
their sub\sphinxhyphen{}particle spheres would interact, based on the sub\sphinxhyphen{}particle
cutoff specified in the pair\_coeff command.  E.g. sub\sphinxhyphen{}particles at the
ends of the line segments that are close to each other.  Which may not
be what you want, since it means the ends of 2 line segments could
pass through each other.  It is up to you to specify a \sphinxstyleemphasis{cutoff}
setting that is consistent with the length of the line segments you
are using and the sub\sphinxhyphen{}particle cutoff settings.
\end{sphinxadmonition}

\sphinxAtStartPar
For style \sphinxstyleemphasis{line/lj}, the following coefficients must be defined for
each pair of atom types via the {\hyperref[\detokenize{pair_coeff::doc}]{\sphinxcrossref{\DUrole{doc}{pair\_coeff}}}} command
as in the examples above, or in the data file or restart files read by
the {\hyperref[\detokenize{read_data::doc}]{\sphinxcrossref{\DUrole{doc}{read\_data}}}} or {\hyperref[\detokenize{read_restart::doc}]{\sphinxcrossref{\DUrole{doc}{read\_restart}}}}
commands:
\begin{itemize}
\item {} 
\sphinxAtStartPar
sizeI (distance units)

\item {} 
\sphinxAtStartPar
sizeJ (distance units)

\item {} 
\sphinxAtStartPar
\(\epsilon\) (energy units)

\item {} 
\sphinxAtStartPar
\(\sigma\) (distance units)

\item {} 
\sphinxAtStartPar
subcutoff (distance units)

\item {} 
\sphinxAtStartPar
cutoff (distance units)

\end{itemize}

\sphinxAtStartPar
The \sphinxstyleemphasis{sizeI} and \sphinxstyleemphasis{sizeJ} coefficients are the sub\sphinxhyphen{}particle sizes for
line particles of type I and type J.  They are used to define the N
sub\sphinxhyphen{}particles per segment as described above.  These coefficients are
actually stored on a per\sphinxhyphen{}type basis.  Thus if there are multiple
pair\_coeff commands that involve type I, as either the first or
second atom type, you should use consistent values for sizeI or sizeJ
in all of them.  If you do not do this, the last value specified for
sizeI will apply to all segments of type I.  If typeI or typeJ refers
to point particles, the corresponding sizeI or sizeJ is ignored; it
can be set to 0.0.

\sphinxAtStartPar
The \(\epsilon\), \(\sigma\), and \sphinxstyleemphasis{subcutoff} coefficients are
used to compute an LJ interactions between a pair of sub\sphinxhyphen{}particles on 2
line segments (of type I and J), or between a sub particle/point
particle pair.  As discussed above, the \sphinxstyleemphasis{subcutoff} and \sphinxstyleemphasis{cutoff} params
are different.  The latter is only used for building the neighbor list
when the distance between centers of two line segments or one segment
and a point particle is calculated.

\sphinxAtStartPar
The \sphinxstyleemphasis{cutoff} coefficient is optional.  If not specified, the global
cutoff is used.


\bigskip\hrule\bigskip



\subsection{Mixing, shift, table, tail correction, restart, rRESPA info}
\label{\detokenize{pair_line_lj:mixing-shift-table-tail-correction-restart-rrespa-info}}
\sphinxAtStartPar
For atom type pairs I,J and I != J, coefficients must be specified.
No default mixing rules are used.

\sphinxAtStartPar
This pair style does not support the {\hyperref[\detokenize{pair_modify::doc}]{\sphinxcrossref{\DUrole{doc}{pair\_modify}}}}
shift, table, and tail options.

\sphinxAtStartPar
This pair style does not write its information to {\hyperref[\detokenize{restart::doc}]{\sphinxcrossref{\DUrole{doc}{binary restart files}}}}.

\sphinxAtStartPar
This pair style can only be used via the \sphinxstyleemphasis{pair} keyword of the
{\hyperref[\detokenize{run_style::doc}]{\sphinxcrossref{\DUrole{doc}{run\_style respa}}}} command.  It does not support the
\sphinxstyleemphasis{inner}, \sphinxstyleemphasis{middle}, \sphinxstyleemphasis{outer} keywords.


\bigskip\hrule\bigskip



\subsection{Restrictions}
\label{\detokenize{pair_line_lj:restrictions}}
\sphinxAtStartPar
This style is part of the ASPHERE package.  It is only enabled if
LAMMPS was built with that package.  See the {\hyperref[\detokenize{Build_package::doc}]{\sphinxcrossref{\DUrole{doc}{Build package}}}} page for more info.

\sphinxAtStartPar
Defining particles to be line segments so they participate in
line/line or line/particle interactions requires the use the
{\hyperref[\detokenize{atom_style::doc}]{\sphinxcrossref{\DUrole{doc}{atom\_style line}}}} command.


\subsection{Related commands}
\label{\detokenize{pair_line_lj:related-commands}}
\sphinxAtStartPar
{\hyperref[\detokenize{pair_coeff::doc}]{\sphinxcrossref{\DUrole{doc}{pair\_coeff}}}}, {\hyperref[\detokenize{pair_tri_lj::doc}]{\sphinxcrossref{\DUrole{doc}{pair\_style tri/lj}}}}


\subsection{Default}
\label{\detokenize{pair_line_lj:default}}
\sphinxAtStartPar
none

\sphinxstepscope

\index{pair\_style list@\spxentry{pair\_style list}}\ignorespaces 

\section{pair\_style list command}
\label{\detokenize{pair_list:pair-style-list-command}}\label{\detokenize{pair_list:index-0}}\label{\detokenize{pair_list::doc}}

\subsection{Syntax}
\label{\detokenize{pair_list:syntax}}
\begin{sphinxVerbatim}[commandchars=\\\{\}]
\PYG{k}{pair\PYGZus{}style}\PYG{+w}{ }\PYG{n}{list}\PYG{+w}{ }\PYG{n}{listfile}\PYG{+w}{ }\PYG{n}{cutoff}\PYG{+w}{ }\PYG{n}{keyword}
\end{sphinxVerbatim}
\begin{itemize}
\item {} 
\sphinxAtStartPar
listfile = name of file with list of pairwise interactions

\item {} 
\sphinxAtStartPar
cutoff = global cutoff (distance units)

\item {} 
\sphinxAtStartPar
keyword = optional flag \sphinxstyleemphasis{nocheck} or \sphinxstyleemphasis{check} (default is \sphinxstyleemphasis{check})

\end{itemize}


\subsection{Examples}
\label{\detokenize{pair_list:examples}}
\begin{sphinxVerbatim}[commandchars=\\\{\}]
\PYG{k}{pair\PYGZus{}style}\PYG{+w}{ }\PYG{n}{list}\PYG{+w}{ }\PYG{n}{restraints.txt}\PYG{+w}{ }\PYG{l+m}{200.0}
\PYG{k}{pair\PYGZus{}coeff}\PYG{+w}{ }\PYG{o}{*}\PYG{+w}{ }\PYG{o}{*}

\PYG{k}{pair\PYGZus{}style}\PYG{+w}{ }\PYG{n}{hybrid}\PYG{o}{/}\PYG{n}{overlay}\PYG{+w}{ }\PYG{n}{lj}\PYG{o}{/}\PYG{n}{cut}\PYG{+w}{ }\PYG{l+m}{1.1225}\PYG{+w}{ }\PYG{n}{list}\PYG{+w}{ }\PYG{n}{pair\PYGZus{}list.txt}\PYG{+w}{ }\PYG{l+m}{300.0}
\PYG{k}{pair\PYGZus{}coeff}\PYG{+w}{ }\PYG{o}{*}\PYG{+w}{ }\PYG{o}{*}\PYG{+w}{ }\PYG{n}{lj}\PYG{o}{/}\PYG{n}{cut}\PYG{+w}{ }\PYG{l+m}{1.0}\PYG{+w}{ }\PYG{l+m}{1.0}
\PYG{k}{pair\PYGZus{}coeff}\PYG{+w}{ }\PYG{l+m}{3}\PYG{o}{*}\PYG{+w}{ }\PYG{l+m}{3}\PYG{o}{*}\PYG{+w}{ }\PYG{n}{list}
\end{sphinxVerbatim}


\subsection{Description}
\label{\detokenize{pair_list:description}}
\sphinxAtStartPar
Style \sphinxstyleemphasis{list} computes interactions between explicitly listed pairs of
atoms with the option to select functional form and parameters for each
individual pair.  Because the parameters are set in the list file, the
pair\_coeff command has no parameters (but still needs to be provided).
The \sphinxstyleemphasis{check} and \sphinxstyleemphasis{nocheck} keywords enable/disable tests that checks
whether all listed pairs of atom IDs were present and the interactions
computed.  If \sphinxstyleemphasis{nocheck} is set and either atom ID is not present, the
interaction is skipped.

\sphinxAtStartPar
This pair style can be thought of as a hybrid between bonded,
non\sphinxhyphen{}bonded, and restraint interactions.  It will typically be used as an
additional interaction within the \sphinxstyleemphasis{hybrid/overlay} pair style.  It
currently supports three interaction styles: a 12\sphinxhyphen{}6 Lennard\sphinxhyphen{}Jones, a
Morse and a harmonic potential.

\sphinxAtStartPar
The format of the list file is as follows:
\begin{itemize}
\item {} 
\sphinxAtStartPar
one line per pair of atoms

\item {} 
\sphinxAtStartPar
empty lines will be ignored

\item {} 
\sphinxAtStartPar
comment text starts with a ‘\#’ character

\item {} 
\sphinxAtStartPar
line syntax: \sphinxstyleemphasis{ID1 ID2 style coeffs cutoff}

\begin{sphinxVerbatim}[commandchars=\\\{\}]
ID1 = atom ID of first atom
ID2 = atom ID of second atom
style = style of interaction
coeffs = list of coeffs
cutoff = cutoff for interaction (optional)
\end{sphinxVerbatim}

\end{itemize}

\sphinxAtStartPar
The cutoff parameter is optional for all but the \sphinxstyleemphasis{quartic} interactions.
If it is not specified, the global cutoff is used.

\sphinxAtStartPar
Here is an example file:

\begin{sphinxVerbatim}[commandchars=\\\{\}]
\PYGZsh{} this is a comment

15 259 lj126     1.0 1.0      50.0
15 603 morse    10.0 1.2 2.0  10.0 \PYGZsh{} and another comment
18 470 harmonic 50.0 1.2       5.0
19 332 quartic  10.0 5.0 \PYGZhy{}1.2 1.2
\end{sphinxVerbatim}

\sphinxAtStartPar
The style \sphinxstyleemphasis{lj126} computes pairwise interactions with the formula
\begin{equation*}
\begin{split}E = 4 \epsilon \left[ \left(\frac{\sigma}{r}\right)^{12} - \left(\frac{\sigma}{r}\right)^6 \right] \qquad r < r_c\end{split}
\end{equation*}
\sphinxAtStartPar
and the coefficients:
\begin{itemize}
\item {} 
\sphinxAtStartPar
\(\epsilon\) (energy units)

\item {} 
\sphinxAtStartPar
\(\sigma\) (distance units)

\end{itemize}

\sphinxAtStartPar
The style \sphinxstyleemphasis{morse} computes pairwise interactions with the formula
\begin{equation*}
\begin{split}E = D_0 \left[ 1 - e^{-\alpha (r - r_0)} \right]^2 \qquad r < r_c\end{split}
\end{equation*}
\sphinxAtStartPar
and the coefficients:
\begin{itemize}
\item {} 
\sphinxAtStartPar
\(D_0\) (energy units)

\item {} 
\sphinxAtStartPar
\(\alpha\) (1/distance units)

\item {} 
\sphinxAtStartPar
\(r_0\) (distance units)

\end{itemize}

\sphinxAtStartPar
The style \sphinxstyleemphasis{harmonic} computes pairwise interactions with the formula
\begin{equation*}
\begin{split}E = K (r - r_0)^2 \qquad r < r_c\end{split}
\end{equation*}
\sphinxAtStartPar
and the coefficients:
\begin{itemize}
\item {} 
\sphinxAtStartPar
\(K\) (energy units)

\item {} 
\sphinxAtStartPar
\(r_0\) (distance units)

\end{itemize}

\sphinxAtStartPar
Note that the usual 1/2 factor is included in \(K\).

\sphinxAtStartPar
The style \sphinxstyleemphasis{quartic} computes pairwise interactions with the formula
\begin{equation*}
\begin{split}E = K (r - r_0)^2 (r - r_0 -b_1) (r - r_0 - b_2) \qquad r < r_c\end{split}
\end{equation*}
\sphinxAtStartPar
and the coefficients:
\begin{itemize}
\item {} 
\sphinxAtStartPar
\(K\) (energy units)

\item {} 
\sphinxAtStartPar
\(r_0\) (distance units)

\item {} 
\sphinxAtStartPar
\(b_1\) (distance units)

\item {} 
\sphinxAtStartPar
\(b_2\) (distance units)

\end{itemize}


\bigskip\hrule\bigskip



\subsection{Mixing, shift, table, tail correction, restart, rRESPA info}
\label{\detokenize{pair_list:mixing-shift-table-tail-correction-restart-rrespa-info}}
\sphinxAtStartPar
This pair style does not support mixing since all parameters are
explicit for each pair.

\sphinxAtStartPar
The {\hyperref[\detokenize{pair_modify::doc}]{\sphinxcrossref{\DUrole{doc}{pair\_modify}}}} shift option is supported by this
pair style.

\sphinxAtStartPar
The {\hyperref[\detokenize{pair_modify::doc}]{\sphinxcrossref{\DUrole{doc}{pair\_modify}}}} table and tail options are not
relevant for this pair style.

\sphinxAtStartPar
This pair style does not write its information to {\hyperref[\detokenize{restart::doc}]{\sphinxcrossref{\DUrole{doc}{binary restart
files}}}}, so pair\_style and pair\_coeff commands need to be
specified in an input script that reads a restart file.

\sphinxAtStartPar
This pair style can only be used via the \sphinxstyleemphasis{pair} keyword of the
{\hyperref[\detokenize{run_style::doc}]{\sphinxcrossref{\DUrole{doc}{run\_style respa}}}} command.  It does not support the
\sphinxstyleemphasis{inner}, \sphinxstyleemphasis{middle}, \sphinxstyleemphasis{outer} keywords.


\bigskip\hrule\bigskip



\subsection{Restrictions}
\label{\detokenize{pair_list:restrictions}}
\sphinxAtStartPar
This pair style does not use a neighbor list and instead identifies
atoms by their IDs.  This has two consequences: 1) The cutoff has to be
chosen sufficiently large, so that the second atom of a pair has to be a
ghost atom on the same node on which the first atom is local; otherwise
the interaction will be skipped.  You can use the \sphinxstyleemphasis{check} option to
detect, if interactions are missing.  2) Unlike other pair styles in
LAMMPS, an atom I will not interact with multiple images of atom J
(assuming the images are within the cutoff distance), but only with the
closest image.

\sphinxAtStartPar
This style is part of the MISC package. It is only enabled if LAMMPS is
build with that package. See the {\hyperref[\detokenize{Build_package::doc}]{\sphinxcrossref{\DUrole{doc}{Build package}}}}
page on for more info.


\subsection{Related commands}
\label{\detokenize{pair_list:related-commands}}
\sphinxAtStartPar
{\hyperref[\detokenize{pair_coeff::doc}]{\sphinxcrossref{\DUrole{doc}{pair\_coeff}}}},
{\hyperref[\detokenize{pair_hybrid::doc}]{\sphinxcrossref{\DUrole{doc}{pair\_style hybrid/overlay}}}},
{\hyperref[\detokenize{pair_lj::doc}]{\sphinxcrossref{\DUrole{doc}{pair\_style lj/cut}}}},
{\hyperref[\detokenize{bond_morse::doc}]{\sphinxcrossref{\DUrole{doc}{bond\_style morse}}}},
{\hyperref[\detokenize{bond_harmonic::doc}]{\sphinxcrossref{\DUrole{doc}{bond\_style harmonic}}}}
{\hyperref[\detokenize{bond_quartic::doc}]{\sphinxcrossref{\DUrole{doc}{bond\_style quartic}}}}


\subsection{Default}
\label{\detokenize{pair_list:default}}
\sphinxAtStartPar
none

\sphinxstepscope

\index{pair\_style lj/cut@\spxentry{pair\_style lj/cut}}\ignorespaces 
\index{pair\_style lj/cut/gpu@\spxentry{pair\_style lj/cut/gpu}}\ignorespaces 
\index{pair\_style lj/cut/intel@\spxentry{pair\_style lj/cut/intel}}\ignorespaces 
\index{pair\_style lj/cut/kk@\spxentry{pair\_style lj/cut/kk}}\ignorespaces 
\index{pair\_style lj/cut/opt@\spxentry{pair\_style lj/cut/opt}}\ignorespaces 
\index{pair\_style lj/cut/omp@\spxentry{pair\_style lj/cut/omp}}\ignorespaces 

\section{pair\_style lj/cut command}
\label{\detokenize{pair_lj:pair-style-lj-cut-command}}\label{\detokenize{pair_lj:index-5}}\label{\detokenize{pair_lj:index-4}}\label{\detokenize{pair_lj:index-3}}\label{\detokenize{pair_lj:index-2}}\label{\detokenize{pair_lj:index-1}}\label{\detokenize{pair_lj:index-0}}\label{\detokenize{pair_lj::doc}}
\sphinxAtStartPar
Accelerator Variants: \sphinxstyleemphasis{lj/cut/gpu}, \sphinxstyleemphasis{lj/cut/intel}, \sphinxstyleemphasis{lj/cut/kk}, \sphinxstyleemphasis{lj/cut/opt}, \sphinxstyleemphasis{lj/cut/omp}


\subsection{Syntax}
\label{\detokenize{pair_lj:syntax}}
\begin{sphinxVerbatim}[commandchars=\\\{\}]
\PYG{k}{pair\PYGZus{}style}\PYG{+w}{ }\PYG{n}{style}\PYG{+w}{ }\PYG{n}{args}
\end{sphinxVerbatim}
\begin{itemize}
\item {} 
\sphinxAtStartPar
style = \sphinxstyleemphasis{lj/cut}

\item {} 
\sphinxAtStartPar
args = list of arguments for a particular style

\end{itemize}
\begin{sphinxalltt}
\sphinxstyleemphasis{lj/cut} args = cutoff
  cutoff = global cutoff for Lennard Jones interactions (distance units)
\end{sphinxalltt}


\subsection{Examples}
\label{\detokenize{pair_lj:examples}}
\begin{sphinxVerbatim}[commandchars=\\\{\}]
\PYG{k}{pair\PYGZus{}style}\PYG{+w}{ }\PYG{n}{lj}\PYG{o}{/}\PYG{n}{cut}\PYG{+w}{ }\PYG{l+m}{2.5}
\PYG{k}{pair\PYGZus{}coeff}\PYG{+w}{ }\PYG{o}{*}\PYG{+w}{ }\PYG{o}{*}\PYG{+w}{ }\PYG{l+m}{1}\PYG{+w}{ }\PYG{l+m}{1}
\PYG{k}{pair\PYGZus{}coeff}\PYG{+w}{ }\PYG{l+m}{1}\PYG{+w}{ }\PYG{l+m}{1}\PYG{+w}{ }\PYG{l+m}{1}\PYG{+w}{ }\PYG{l+m}{1.1}\PYG{+w}{ }\PYG{l+m}{2.8}
\end{sphinxVerbatim}


\subsection{Description}
\label{\detokenize{pair_lj:description}}
\sphinxAtStartPar
The \sphinxstyleemphasis{lj/cut} styles compute the standard 12/6 Lennard\sphinxhyphen{}Jones potential,
given by
\begin{equation*}
\begin{split}E = 4 \epsilon \left[ \left(\frac{\sigma}{r}\right)^{12} -
    \left(\frac{\sigma}{r}\right)^6 \right]   \qquad r < r_c\end{split}
\end{equation*}
\sphinxAtStartPar
\(r_c\) is the cutoff.

\sphinxAtStartPar
See the {\hyperref[\detokenize{pair_lj_cut_coul::doc}]{\sphinxcrossref{\DUrole{doc}{lj/cut/coul}}}} styles to add a Coulombic
pairwise interaction and the {\hyperref[\detokenize{pair_lj_cut_tip4p::doc}]{\sphinxcrossref{\DUrole{doc}{lj/cut/tip4p}}}} styles to
add the TIP4P water model.


\subsection{Coefficients}
\label{\detokenize{pair_lj:coefficients}}
\sphinxAtStartPar
The following coefficients must be defined for each pair of atoms types via the
{\hyperref[\detokenize{pair_coeff::doc}]{\sphinxcrossref{\DUrole{doc}{pair\_coeff}}}} command as in the examples above, or in the data
file or restart files read by the {\hyperref[\detokenize{read_data::doc}]{\sphinxcrossref{\DUrole{doc}{read\_data}}}} or
{\hyperref[\detokenize{read_restart::doc}]{\sphinxcrossref{\DUrole{doc}{read\_restart}}}} commands, or by mixing as described below:
\begin{itemize}
\item {} 
\sphinxAtStartPar
\(\epsilon\) (energy units)

\item {} 
\sphinxAtStartPar
\(\sigma\) (distance units)

\item {} 
\sphinxAtStartPar
LJ cutoff (distance units)

\end{itemize}

\sphinxAtStartPar
The last coefficient is optional.  If not specified, the global
LJ cutoff specified in the pair\_style command is used.

\sphinxAtStartPar
Note that \(\sigma\) is defined in the LJ formula as the
zero\sphinxhyphen{}crossing distance for the potential, \sphinxstyleemphasis{not} as the energy minimum at
\(r_0 = 2^{\frac{1}{6}} \sigma\).  The \_same\_ potential function becomes:
\begin{equation*}
\begin{split}E = \epsilon \left[ \left(\frac{r_0}{r}\right)^{12} -
     2 \left(\frac{r_0}{r}\right)^6 \right]  \qquad r < r_c\end{split}
\end{equation*}
\sphinxAtStartPar
When using the minimum as reference width. In the literature both
formulations are used, but the describe the same potential, only the
\(\sigma\) value must be computed by \(\sigma = r_0 /
2^{\frac{1}{6}}\) for use with LAMMPS, if this latter formulation is
used.


\bigskip\hrule\bigskip


\sphinxAtStartPar
A version of these styles with a soft core, \sphinxstyleemphasis{lj/cut/soft}, suitable
for use in free energy calculations, is part of the FEP package and
is documented with the {\hyperref[\detokenize{pair_fep_soft::doc}]{\sphinxcrossref{\DUrole{doc}{pair\_style */soft}}}}
styles.


\bigskip\hrule\bigskip


\sphinxAtStartPar
Styles with a \sphinxstyleemphasis{gpu}, \sphinxstyleemphasis{intel}, \sphinxstyleemphasis{kk}, \sphinxstyleemphasis{omp}, or \sphinxstyleemphasis{opt} suffix are
functionally the same as the corresponding style without the suffix.
They have been optimized to run faster, depending on your available
hardware, as discussed on the {\hyperref[\detokenize{Speed_packages::doc}]{\sphinxcrossref{\DUrole{doc}{Accelerator packages}}}}
page.  The accelerated styles take the same arguments and should
produce the same results, except for round\sphinxhyphen{}off and precision issues.

\sphinxAtStartPar
These accelerated styles are part of the GPU, INTEL, KOKKOS,
OPENMP, and OPT packages, respectively.  They are only enabled if
LAMMPS was built with those packages.  See the {\hyperref[\detokenize{Build_package::doc}]{\sphinxcrossref{\DUrole{doc}{Build package}}}} page for more info.

\sphinxAtStartPar
You can specify the accelerated styles explicitly in your input script
by including their suffix, or you can use the {\hyperref[\detokenize{Run_options::doc}]{\sphinxcrossref{\DUrole{doc}{\sphinxhyphen{}suffix command\sphinxhyphen{}line switch}}}} when you invoke LAMMPS, or you can use the
{\hyperref[\detokenize{suffix::doc}]{\sphinxcrossref{\DUrole{doc}{suffix}}}} command in your input script.

\sphinxAtStartPar
See the {\hyperref[\detokenize{Speed_packages::doc}]{\sphinxcrossref{\DUrole{doc}{Accelerator packages}}}} page for more
instructions on how to use the accelerated styles effectively.


\bigskip\hrule\bigskip



\subsection{Mixing, shift, table, tail correction, restart, rRESPA info}
\label{\detokenize{pair_lj:mixing-shift-table-tail-correction-restart-rrespa-info}}
\sphinxAtStartPar
For atom type pairs I,J and I != J, the epsilon and sigma coefficients
and cutoff distance for all of the lj/cut pair styles can be mixed.
The default mix value is \sphinxstyleemphasis{geometric}.  See the “pair\_modify” command
for details.

\sphinxAtStartPar
All of the \sphinxstyleemphasis{lj/cut} pair styles support the
{\hyperref[\detokenize{pair_modify::doc}]{\sphinxcrossref{\DUrole{doc}{pair\_modify}}}} shift option for the energy of the
Lennard\sphinxhyphen{}Jones portion of the pair interaction.

\sphinxAtStartPar
All of the \sphinxstyleemphasis{lj/cut} pair styles support the
{\hyperref[\detokenize{pair_modify::doc}]{\sphinxcrossref{\DUrole{doc}{pair\_modify}}}} tail option for adding a long\sphinxhyphen{}range
tail correction to the energy and pressure for the Lennard\sphinxhyphen{}Jones
portion of the pair interaction.

\sphinxAtStartPar
All of the \sphinxstyleemphasis{lj/cut} pair styles write their information to {\hyperref[\detokenize{restart::doc}]{\sphinxcrossref{\DUrole{doc}{binary restart files}}}}, so pair\_style and pair\_coeff commands do
not need to be specified in an input script that reads a restart file.

\sphinxAtStartPar
The \sphinxstyleemphasis{lj/cut} pair styles support the use of the \sphinxstyleemphasis{inner}, \sphinxstyleemphasis{middle}, and
\sphinxstyleemphasis{outer} keywords of the {\hyperref[\detokenize{run_style::doc}]{\sphinxcrossref{\DUrole{doc}{run\_style respa}}}} command,
meaning the pairwise forces can be partitioned by distance at different
levels of the rRESPA hierarchy.  The other styles only support the
\sphinxstyleemphasis{pair} keyword of run\_style respa.  See the {\hyperref[\detokenize{run_style::doc}]{\sphinxcrossref{\DUrole{doc}{run\_style}}}}
command for details.


\bigskip\hrule\bigskip



\subsection{Related commands}
\label{\detokenize{pair_lj:related-commands}}\begin{itemize}
\item {} 
\sphinxAtStartPar
{\hyperref[\detokenize{pair_coeff::doc}]{\sphinxcrossref{\DUrole{doc}{pair\_coeff}}}}

\item {} 
\sphinxAtStartPar
{\hyperref[\detokenize{pair_lj_cut_coul::doc}]{\sphinxcrossref{\DUrole{doc}{pair\_style lj/cut/coul/cut}}}}

\item {} 
\sphinxAtStartPar
{\hyperref[\detokenize{pair_lj_cut_coul::doc}]{\sphinxcrossref{\DUrole{doc}{pair\_style lj/cut/coul/debye}}}}

\item {} 
\sphinxAtStartPar
{\hyperref[\detokenize{pair_lj_cut_coul::doc}]{\sphinxcrossref{\DUrole{doc}{pair\_style lj/cut/coul/dsf}}}}

\item {} 
\sphinxAtStartPar
{\hyperref[\detokenize{pair_lj_cut_coul::doc}]{\sphinxcrossref{\DUrole{doc}{pair\_style lj/cut/coul/long}}}}

\item {} 
\sphinxAtStartPar
{\hyperref[\detokenize{pair_lj_cut_coul::doc}]{\sphinxcrossref{\DUrole{doc}{pair\_style lj/cut/coul/msm}}}}

\item {} 
\sphinxAtStartPar
{\hyperref[\detokenize{pair_lj_cut_coul::doc}]{\sphinxcrossref{\DUrole{doc}{pair\_style lj/cut/coul/wolf}}}}

\item {} 
\sphinxAtStartPar
{\hyperref[\detokenize{pair_lj_cut_tip4p::doc}]{\sphinxcrossref{\DUrole{doc}{pair\_style lj/cut/tip4p/cut}}}}

\item {} 
\sphinxAtStartPar
{\hyperref[\detokenize{pair_lj_cut_tip4p::doc}]{\sphinxcrossref{\DUrole{doc}{pair\_style lj/cut/tip4p/long}}}}

\end{itemize}


\subsection{Default}
\label{\detokenize{pair_lj:default}}
\sphinxAtStartPar
none

\sphinxstepscope

\index{pair\_style lj96/cut@\spxentry{pair\_style lj96/cut}}\ignorespaces 
\index{pair\_style lj96/cut/gpu@\spxentry{pair\_style lj96/cut/gpu}}\ignorespaces 
\index{pair\_style lj96/cut/omp@\spxentry{pair\_style lj96/cut/omp}}\ignorespaces 

\section{pair\_style lj96/cut command}
\label{\detokenize{pair_lj96:pair-style-lj96-cut-command}}\label{\detokenize{pair_lj96:index-2}}\label{\detokenize{pair_lj96:index-1}}\label{\detokenize{pair_lj96:index-0}}\label{\detokenize{pair_lj96::doc}}
\sphinxAtStartPar
Accelerator Variants: \sphinxstyleemphasis{lj96/cut/gpu}, \sphinxstyleemphasis{lj96/cut/omp}


\subsection{Syntax}
\label{\detokenize{pair_lj96:syntax}}
\begin{sphinxVerbatim}[commandchars=\\\{\}]
\PYG{k}{pair\PYGZus{}style}\PYG{+w}{ }\PYG{n}{lj96}\PYG{o}{/}\PYG{n}{cut}\PYG{+w}{ }\PYG{n}{cutoff}
\end{sphinxVerbatim}
\begin{itemize}
\item {} 
\sphinxAtStartPar
cutoff = global cutoff for lj96/cut interactions (distance units)

\end{itemize}


\subsection{Examples}
\label{\detokenize{pair_lj96:examples}}
\begin{sphinxVerbatim}[commandchars=\\\{\}]
\PYG{k}{pair\PYGZus{}style}\PYG{+w}{ }\PYG{n}{lj96}\PYG{o}{/}\PYG{n}{cut}\PYG{+w}{ }\PYG{l+m}{2.5}
\PYG{k}{pair\PYGZus{}coeff}\PYG{+w}{ }\PYG{o}{*}\PYG{+w}{ }\PYG{o}{*}\PYG{+w}{ }\PYG{l+m}{1.0}\PYG{+w}{ }\PYG{l+m}{1.0}\PYG{+w}{ }\PYG{l+m}{4.0}
\PYG{k}{pair\PYGZus{}coeff}\PYG{+w}{ }\PYG{l+m}{1}\PYG{+w}{ }\PYG{l+m}{1}\PYG{+w}{ }\PYG{l+m}{1.0}\PYG{+w}{ }\PYG{l+m}{1.0}
\end{sphinxVerbatim}


\subsection{Description}
\label{\detokenize{pair_lj96:description}}
\sphinxAtStartPar
The \sphinxstyleemphasis{lj96/cut} style compute a 9/6 Lennard\sphinxhyphen{}Jones potential, instead
of the standard 12/6 potential, given by
\begin{equation*}
\begin{split}E = 4 \epsilon \left[ \left(\frac{\sigma}{r}\right)^{9} -
\left(\frac{\sigma}{r}\right)^6 \right]
                    \qquad r < r_c\end{split}
\end{equation*}
\sphinxAtStartPar
\(r_c\) is the cutoff.

\sphinxAtStartPar
The following coefficients must be defined for each pair of atoms
types via the {\hyperref[\detokenize{pair_coeff::doc}]{\sphinxcrossref{\DUrole{doc}{pair\_coeff}}}} command as in the examples
above, or in the data file or restart files read by the
{\hyperref[\detokenize{read_data::doc}]{\sphinxcrossref{\DUrole{doc}{read\_data}}}} or {\hyperref[\detokenize{read_restart::doc}]{\sphinxcrossref{\DUrole{doc}{read\_restart}}}}
commands, or by mixing as described below:
\begin{itemize}
\item {} 
\sphinxAtStartPar
\(\epsilon\) (energy units)

\item {} 
\sphinxAtStartPar
\(\sigma\) (distance units)

\item {} 
\sphinxAtStartPar
cutoff (distance units)

\end{itemize}

\sphinxAtStartPar
The last coefficient is optional.  If not specified, the global LJ
cutoff specified in the pair\_style command is used.


\bigskip\hrule\bigskip


\sphinxAtStartPar
Styles with a \sphinxstyleemphasis{gpu}, \sphinxstyleemphasis{intel}, \sphinxstyleemphasis{kk}, \sphinxstyleemphasis{omp}, or \sphinxstyleemphasis{opt} suffix are
functionally the same as the corresponding style without the suffix.
They have been optimized to run faster, depending on your available
hardware, as discussed on the {\hyperref[\detokenize{Speed_packages::doc}]{\sphinxcrossref{\DUrole{doc}{Accelerator packages}}}}
page.  The accelerated styles take the same arguments and should
produce the same results, except for round\sphinxhyphen{}off and precision issues.

\sphinxAtStartPar
These accelerated styles are part of the GPU, INTEL, KOKKOS,
OPENMP, and OPT packages, respectively.  They are only enabled if
LAMMPS was built with those packages.  See the {\hyperref[\detokenize{Build_package::doc}]{\sphinxcrossref{\DUrole{doc}{Build package}}}} page for more info.

\sphinxAtStartPar
You can specify the accelerated styles explicitly in your input script
by including their suffix, or you can use the {\hyperref[\detokenize{Run_options::doc}]{\sphinxcrossref{\DUrole{doc}{\sphinxhyphen{}suffix command\sphinxhyphen{}line switch}}}} when you invoke LAMMPS, or you can use the
{\hyperref[\detokenize{suffix::doc}]{\sphinxcrossref{\DUrole{doc}{suffix}}}} command in your input script.

\sphinxAtStartPar
See the {\hyperref[\detokenize{Speed_packages::doc}]{\sphinxcrossref{\DUrole{doc}{Accelerator packages}}}} page for more
instructions on how to use the accelerated styles effectively.


\bigskip\hrule\bigskip



\subsection{Mixing, shift, table, tail correction, restart, rRESPA info}
\label{\detokenize{pair_lj96:mixing-shift-table-tail-correction-restart-rrespa-info}}
\sphinxAtStartPar
For atom type pairs I,J and I != J, the epsilon and sigma coefficients
and cutoff distance for all of the lj/cut pair styles can be mixed.
The default mix value is \sphinxstyleemphasis{geometric}.  See the “pair\_modify” command
for details.

\sphinxAtStartPar
This pair style supports the {\hyperref[\detokenize{pair_modify::doc}]{\sphinxcrossref{\DUrole{doc}{pair\_modify}}}} shift
option for the energy of the pair interaction.

\sphinxAtStartPar
The {\hyperref[\detokenize{pair_modify::doc}]{\sphinxcrossref{\DUrole{doc}{pair\_modify}}}} table option is not relevant
for this pair style.

\sphinxAtStartPar
This pair style supports the {\hyperref[\detokenize{pair_modify::doc}]{\sphinxcrossref{\DUrole{doc}{pair\_modify}}}} tail
option for adding a long\sphinxhyphen{}range tail correction to the energy and
pressure of the pair interaction.

\sphinxAtStartPar
This pair style writes its information to {\hyperref[\detokenize{restart::doc}]{\sphinxcrossref{\DUrole{doc}{binary restart files}}}}, so pair\_style and pair\_coeff commands do not need
to be specified in an input script that reads a restart file.

\sphinxAtStartPar
This pair style supports the use of the \sphinxstyleemphasis{inner}, \sphinxstyleemphasis{middle}, and \sphinxstyleemphasis{outer}
keywords of the {\hyperref[\detokenize{run_style::doc}]{\sphinxcrossref{\DUrole{doc}{run\_style respa}}}} command, meaning the
pairwise forces can be partitioned by distance at different levels of
the rRESPA hierarchy.  See the {\hyperref[\detokenize{run_style::doc}]{\sphinxcrossref{\DUrole{doc}{run\_style}}}} command for
details.


\bigskip\hrule\bigskip



\subsection{Restrictions}
\label{\detokenize{pair_lj96:restrictions}}
\sphinxAtStartPar
This pair style is part of the EXTRA\sphinxhyphen{}PAIR package.  It is only enabled if
LAMMPS was built with that package.  See the
{\hyperref[\detokenize{Build_package::doc}]{\sphinxcrossref{\DUrole{doc}{Build package}}}} page for more info.


\subsection{Related commands}
\label{\detokenize{pair_lj96:related-commands}}
\sphinxAtStartPar
{\hyperref[\detokenize{pair_coeff::doc}]{\sphinxcrossref{\DUrole{doc}{pair\_coeff}}}}


\subsection{Default}
\label{\detokenize{pair_lj96:default}}
\sphinxAtStartPar
none

\sphinxstepscope

\index{pair\_style lj/cubic@\spxentry{pair\_style lj/cubic}}\ignorespaces 
\index{pair\_style lj/cubic/gpu@\spxentry{pair\_style lj/cubic/gpu}}\ignorespaces 
\index{pair\_style lj/cubic/omp@\spxentry{pair\_style lj/cubic/omp}}\ignorespaces 

\section{pair\_style lj/cubic command}
\label{\detokenize{pair_lj_cubic:pair-style-lj-cubic-command}}\label{\detokenize{pair_lj_cubic:index-2}}\label{\detokenize{pair_lj_cubic:index-1}}\label{\detokenize{pair_lj_cubic:index-0}}\label{\detokenize{pair_lj_cubic::doc}}
\sphinxAtStartPar
Accelerator Variants: \sphinxstyleemphasis{lj/cubic/gpu}, \sphinxstyleemphasis{lj/cubic/omp}


\subsection{Syntax}
\label{\detokenize{pair_lj_cubic:syntax}}
\begin{sphinxVerbatim}[commandchars=\\\{\}]
\PYG{k}{pair\PYGZus{}style}\PYG{+w}{ }\PYG{n}{lj}\PYG{o}{/}\PYG{n}{cubic}
\end{sphinxVerbatim}


\subsection{Examples}
\label{\detokenize{pair_lj_cubic:examples}}
\begin{sphinxVerbatim}[commandchars=\\\{\}]
\PYG{k}{pair\PYGZus{}style}\PYG{+w}{ }\PYG{n}{lj}\PYG{o}{/}\PYG{n}{cubic}
\PYG{k}{pair\PYGZus{}coeff}\PYG{+w}{ }\PYG{o}{*}\PYG{+w}{ }\PYG{o}{*}\PYG{+w}{ }\PYG{l+m}{1.0}\PYG{+w}{ }\PYG{l+m}{0.8908987}
\end{sphinxVerbatim}


\subsection{Description}
\label{\detokenize{pair_lj_cubic:description}}
\sphinxAtStartPar
The \sphinxstyleemphasis{lj/cubic} style computes a truncated LJ interaction potential
whose energy and force are continuous everywhere.  Inside the
inflection point the interaction is identical to the standard 12/6
{\hyperref[\detokenize{pair_lj::doc}]{\sphinxcrossref{\DUrole{doc}{Lennard\sphinxhyphen{}Jones}}}} potential.  The LJ function outside the
inflection point is replaced with a cubic function of distance. The
energy, force, and second derivative are continuous at the inflection
point.  The cubic coefficient A3 is chosen so that both energy and
force go to zero at the cutoff distance.  Outside the cutoff distance
the energy and force are zero.
\begin{equation*}
\begin{split}E & = u_{LJ}(r) \qquad r \leq r_s \\
  & = u_{LJ}(r_s) + (r-r_s) u'_{LJ}(r_s) - \frac{1}{6} A_3 (r-r_s)^3 \qquad r_s < r \leq r_c \\
  & = 0 \qquad r > r_c\end{split}
\end{equation*}
\sphinxAtStartPar
The location of the inflection point \(r_s\) is defined
by the LJ diameter, \(r_s/\sigma = (26/7)^{1/6}\). The cutoff distance
is defined by \(r_c/r_s = 67/48\) or \(r_c/\sigma = 1.737...\)
The analytic expression for the
the cubic coefficient
\(A_3 r_{min}^3/\epsilon = 27.93...\) is given in the paper by
Holian and Ravelo {\hyperref[\detokenize{pair_lj_cubic:holian}]{\sphinxcrossref{\DUrole{std,std-ref}{(Holian)}}}}.

\sphinxAtStartPar
This potential is commonly used to study the shock mechanics of FCC
solids, as in Ravelo et al. {\hyperref[\detokenize{pair_lj_cubic:ravelo2}]{\sphinxcrossref{\DUrole{std,std-ref}{(Ravelo)}}}}.

\sphinxAtStartPar
The following coefficients must be defined for each pair of atom types
via the {\hyperref[\detokenize{pair_coeff::doc}]{\sphinxcrossref{\DUrole{doc}{pair\_coeff}}}} command as in the example above,
or in the data file or restart files read by the
{\hyperref[\detokenize{read_data::doc}]{\sphinxcrossref{\DUrole{doc}{read\_data}}}} or {\hyperref[\detokenize{read_restart::doc}]{\sphinxcrossref{\DUrole{doc}{read\_restart}}}}
commands, or by mixing as described below:
\begin{itemize}
\item {} 
\sphinxAtStartPar
\(\epsilon\) (energy units)

\item {} 
\sphinxAtStartPar
\(\sigma\) (distance units)

\end{itemize}

\sphinxAtStartPar
Note that \(\sigma\) is defined in the LJ formula as the
zero\sphinxhyphen{}crossing distance for the potential, not as the energy minimum,
which is located at \(r_{min} = 2^{\frac{1}{6}} \sigma\). In the
above example, \(\sigma = 0.8908987\), so \(r_{min} = 1.0\).


\bigskip\hrule\bigskip


\sphinxAtStartPar
Styles with a \sphinxstyleemphasis{gpu}, \sphinxstyleemphasis{intel}, \sphinxstyleemphasis{kk}, \sphinxstyleemphasis{omp}, or \sphinxstyleemphasis{opt} suffix are
functionally the same as the corresponding style without the suffix.
They have been optimized to run faster, depending on your available
hardware, as discussed on the {\hyperref[\detokenize{Speed_packages::doc}]{\sphinxcrossref{\DUrole{doc}{Accelerator packages}}}}
page.  The accelerated styles take the same arguments and should
produce the same results, except for round\sphinxhyphen{}off and precision issues.

\sphinxAtStartPar
These accelerated styles are part of the GPU, INTEL, KOKKOS,
OPENMP, and OPT packages, respectively.  They are only enabled if
LAMMPS was built with those packages.  See the {\hyperref[\detokenize{Build_package::doc}]{\sphinxcrossref{\DUrole{doc}{Build package}}}} page for more info.

\sphinxAtStartPar
You can specify the accelerated styles explicitly in your input script
by including their suffix, or you can use the {\hyperref[\detokenize{Run_options::doc}]{\sphinxcrossref{\DUrole{doc}{\sphinxhyphen{}suffix command\sphinxhyphen{}line switch}}}} when you invoke LAMMPS, or you can use the
{\hyperref[\detokenize{suffix::doc}]{\sphinxcrossref{\DUrole{doc}{suffix}}}} command in your input script.

\sphinxAtStartPar
See the {\hyperref[\detokenize{Speed_packages::doc}]{\sphinxcrossref{\DUrole{doc}{Accelerator packages}}}} page for more
instructions on how to use the accelerated styles effectively.


\bigskip\hrule\bigskip



\subsection{Mixing, shift, table, tail correction, restart, rRESPA info}
\label{\detokenize{pair_lj_cubic:mixing-shift-table-tail-correction-restart-rrespa-info}}
\sphinxAtStartPar
For atom type pairs I,J and I != J, the epsilon and sigma coefficients
and cutoff distance for all of the lj/cut pair styles can be mixed.
The default mix value is \sphinxstyleemphasis{geometric}.  See the “pair\_modify” command
for details.

\sphinxAtStartPar
The lj/cubic pair style does not support the
{\hyperref[\detokenize{pair_modify::doc}]{\sphinxcrossref{\DUrole{doc}{pair\_modify}}}} shift option,
since pair interaction is already smoothed to 0.0 at the
cutoff.

\sphinxAtStartPar
The {\hyperref[\detokenize{pair_modify::doc}]{\sphinxcrossref{\DUrole{doc}{pair\_modify}}}} table option is not relevant
for this pair style.

\sphinxAtStartPar
The lj/cubic pair style does not support the
{\hyperref[\detokenize{pair_modify::doc}]{\sphinxcrossref{\DUrole{doc}{pair\_modify}}}} tail option for adding long\sphinxhyphen{}range tail
corrections to energy and pressure, since there are no corrections for
a potential that goes to 0.0 at the cutoff.

\sphinxAtStartPar
The lj/cubic pair style writes its information to {\hyperref[\detokenize{restart::doc}]{\sphinxcrossref{\DUrole{doc}{binary restart files}}}}, so pair\_style and pair\_coeff commands do
not need to be specified in an input script that reads a restart file.

\sphinxAtStartPar
The lj/cubic pair style can only be used via the \sphinxstyleemphasis{pair}
keyword of the {\hyperref[\detokenize{run_style::doc}]{\sphinxcrossref{\DUrole{doc}{run\_style respa}}}} command.  It does not
support the \sphinxstyleemphasis{inner}, \sphinxstyleemphasis{middle}, \sphinxstyleemphasis{outer} keywords.


\bigskip\hrule\bigskip



\subsection{Restrictions}
\label{\detokenize{pair_lj_cubic:restrictions}}
\sphinxAtStartPar
This pair style is part of the EXTRA\sphinxhyphen{}PAIR package.  It is only enabled if
LAMMPS was built with that package.  See the
{\hyperref[\detokenize{Build_package::doc}]{\sphinxcrossref{\DUrole{doc}{Build package}}}} page for more info.


\subsection{Related commands}
\label{\detokenize{pair_lj_cubic:related-commands}}
\sphinxAtStartPar
{\hyperref[\detokenize{pair_coeff::doc}]{\sphinxcrossref{\DUrole{doc}{pair\_coeff}}}}


\subsection{Default}
\label{\detokenize{pair_lj_cubic:default}}
\sphinxAtStartPar
none


\bigskip\hrule\bigskip

\phantomsection\label{\detokenize{pair_lj_cubic:holian}}\phantomsection\label{\detokenize{pair_lj_cubic:ravelo2}}
\sphinxAtStartPar
\sphinxstylestrong{(Holian)} Holian and Ravelo, Phys Rev B, 51, 11275 (1995).

\sphinxAtStartPar
\sphinxstylestrong{(Ravelo)} Ravelo, Holian, Germann and Lomdahl, Phys Rev B, 70, 014103 (2004).

\sphinxstepscope

\index{pair\_style lj/cut/coul/cut@\spxentry{pair\_style lj/cut/coul/cut}}\ignorespaces 
\index{pair\_style lj/cut/coul/cut/gpu@\spxentry{pair\_style lj/cut/coul/cut/gpu}}\ignorespaces 
\index{pair\_style lj/cut/coul/cut/kk@\spxentry{pair\_style lj/cut/coul/cut/kk}}\ignorespaces 
\index{pair\_style lj/cut/coul/cut/omp@\spxentry{pair\_style lj/cut/coul/cut/omp}}\ignorespaces 
\index{pair\_style lj/cut/coul/debye@\spxentry{pair\_style lj/cut/coul/debye}}\ignorespaces 
\index{pair\_style lj/cut/coul/debye/gpu@\spxentry{pair\_style lj/cut/coul/debye/gpu}}\ignorespaces 
\index{pair\_style lj/cut/coul/debye/kk@\spxentry{pair\_style lj/cut/coul/debye/kk}}\ignorespaces 
\index{pair\_style lj/cut/coul/debye/omp@\spxentry{pair\_style lj/cut/coul/debye/omp}}\ignorespaces 
\index{pair\_style lj/cut/coul/dsf@\spxentry{pair\_style lj/cut/coul/dsf}}\ignorespaces 
\index{pair\_style lj/cut/coul/dsf/gpu@\spxentry{pair\_style lj/cut/coul/dsf/gpu}}\ignorespaces 
\index{pair\_style lj/cut/coul/dsf/kk@\spxentry{pair\_style lj/cut/coul/dsf/kk}}\ignorespaces 
\index{pair\_style lj/cut/coul/dsf/omp@\spxentry{pair\_style lj/cut/coul/dsf/omp}}\ignorespaces 
\index{pair\_style lj/cut/coul/long@\spxentry{pair\_style lj/cut/coul/long}}\ignorespaces 
\index{pair\_style lj/cut/coul/long/gpu@\spxentry{pair\_style lj/cut/coul/long/gpu}}\ignorespaces 
\index{pair\_style lj/cut/coul/long/kk@\spxentry{pair\_style lj/cut/coul/long/kk}}\ignorespaces 
\index{pair\_style lj/cut/coul/long/intel@\spxentry{pair\_style lj/cut/coul/long/intel}}\ignorespaces 
\index{pair\_style lj/cut/coul/long/opt@\spxentry{pair\_style lj/cut/coul/long/opt}}\ignorespaces 
\index{pair\_style lj/cut/coul/long/omp@\spxentry{pair\_style lj/cut/coul/long/omp}}\ignorespaces 
\index{pair\_style lj/cut/coul/msm@\spxentry{pair\_style lj/cut/coul/msm}}\ignorespaces 
\index{pair\_style lj/cut/coul/msm/gpu@\spxentry{pair\_style lj/cut/coul/msm/gpu}}\ignorespaces 
\index{pair\_style lj/cut/coul/msm/omp@\spxentry{pair\_style lj/cut/coul/msm/omp}}\ignorespaces 
\index{pair\_style lj/cut/coul/wolf@\spxentry{pair\_style lj/cut/coul/wolf}}\ignorespaces 
\index{pair\_style lj/cut/coul/wolf/omp@\spxentry{pair\_style lj/cut/coul/wolf/omp}}\ignorespaces 

\section{pair\_style lj/cut/coul/cut command}
\label{\detokenize{pair_lj_cut_coul:pair-style-lj-cut-coul-cut-command}}\label{\detokenize{pair_lj_cut_coul:index-22}}\label{\detokenize{pair_lj_cut_coul:index-21}}\label{\detokenize{pair_lj_cut_coul:index-20}}\label{\detokenize{pair_lj_cut_coul:index-19}}\label{\detokenize{pair_lj_cut_coul:index-18}}\label{\detokenize{pair_lj_cut_coul:index-17}}\label{\detokenize{pair_lj_cut_coul:index-16}}\label{\detokenize{pair_lj_cut_coul:index-15}}\label{\detokenize{pair_lj_cut_coul:index-14}}\label{\detokenize{pair_lj_cut_coul:index-13}}\label{\detokenize{pair_lj_cut_coul:index-12}}\label{\detokenize{pair_lj_cut_coul:index-11}}\label{\detokenize{pair_lj_cut_coul:index-10}}\label{\detokenize{pair_lj_cut_coul:index-9}}\label{\detokenize{pair_lj_cut_coul:index-8}}\label{\detokenize{pair_lj_cut_coul:index-7}}\label{\detokenize{pair_lj_cut_coul:index-6}}\label{\detokenize{pair_lj_cut_coul:index-5}}\label{\detokenize{pair_lj_cut_coul:index-4}}\label{\detokenize{pair_lj_cut_coul:index-3}}\label{\detokenize{pair_lj_cut_coul:index-2}}\label{\detokenize{pair_lj_cut_coul:index-1}}\label{\detokenize{pair_lj_cut_coul:index-0}}\label{\detokenize{pair_lj_cut_coul::doc}}
\sphinxAtStartPar
Accelerator Variants: \sphinxstyleemphasis{lj/cut/coul/cut/gpu}, \sphinxstyleemphasis{lj/cut/coul/cut/kk}, \sphinxstyleemphasis{lj/cut/coul/cut/omp}


\section{pair\_style lj/cut/coul/debye command}
\label{\detokenize{pair_lj_cut_coul:pair-style-lj-cut-coul-debye-command}}
\sphinxAtStartPar
Accelerator Variants: \sphinxstyleemphasis{lj/cut/coul/debye/gpu}, \sphinxstyleemphasis{lj/cut/coul/debye/kk}, \sphinxstyleemphasis{lj/cut/coul/debye/omp}


\section{pair\_style lj/cut/coul/dsf command}
\label{\detokenize{pair_lj_cut_coul:pair-style-lj-cut-coul-dsf-command}}
\sphinxAtStartPar
Accelerator Variants: \sphinxstyleemphasis{lj/cut/coul/dsf/gpu}, \sphinxstyleemphasis{lj/cut/coul/dsf/kk}, \sphinxstyleemphasis{lj/cut/coul/dsf/omp}


\section{pair\_style lj/cut/coul/long command}
\label{\detokenize{pair_lj_cut_coul:pair-style-lj-cut-coul-long-command}}
\sphinxAtStartPar
Accelerator Variants: \sphinxstyleemphasis{lj/cut/coul/long/gpu}, \sphinxstyleemphasis{lj/cut/coul/long/kk}, \sphinxstyleemphasis{lj/cut/coul/long/intel}, \sphinxstyleemphasis{lj/cut/coul/long/opt}, \sphinxstyleemphasis{lj/cut/coul/long/omp}


\section{pair\_style lj/cut/coul/msm command}
\label{\detokenize{pair_lj_cut_coul:pair-style-lj-cut-coul-msm-command}}
\sphinxAtStartPar
Accelerator Variants: \sphinxstyleemphasis{lj/cut/coul/msm/gpu}, \sphinxstyleemphasis{lj/cut/coul/msm/omp}


\section{pair\_style lj/cut/coul/wolf command}
\label{\detokenize{pair_lj_cut_coul:pair-style-lj-cut-coul-wolf-command}}
\sphinxAtStartPar
Accelerator Variants: \sphinxstyleemphasis{lj/cut/coul/wolf/omp}


\subsection{Syntax}
\label{\detokenize{pair_lj_cut_coul:syntax}}
\begin{sphinxVerbatim}[commandchars=\\\{\}]
\PYG{k}{pair\PYGZus{}style}\PYG{+w}{ }\PYG{n}{style}\PYG{+w}{ }\PYG{n}{args}
\end{sphinxVerbatim}
\begin{itemize}
\item {} 
\sphinxAtStartPar
style = \sphinxstyleemphasis{lj/cut/coul/cut} or \sphinxstyleemphasis{lj/cut/coul/debye} or \sphinxstyleemphasis{lj/cut/coul/dsf} or \sphinxstyleemphasis{lj/cut/coul/long} \sphinxstyleemphasis{lj/cut/coul/msm} or \sphinxstyleemphasis{lj/cut/coul/wolf}

\item {} 
\sphinxAtStartPar
args = list of arguments for a particular style

\end{itemize}
\begin{sphinxalltt}
\sphinxstyleemphasis{lj/cut/coul/cut} args = cutoff (cutoff2)
  cutoff = global cutoff for LJ (and Coulombic if only 1 arg) (distance units)
  cutoff2 = global cutoff for Coulombic (optional) (distance units)
\sphinxstyleemphasis{lj/cut/coul/debye} args = kappa cutoff (cutoff2)
  kappa = inverse of the Debye length (inverse distance units)
  cutoff = global cutoff for LJ (and Coulombic if only 1 arg) (distance units)
  cutoff2 = global cutoff for Coulombic (optional) (distance units)
\sphinxstyleemphasis{lj/cut/coul/dsf} args = alpha cutoff (cutoff2)
  alpha = damping parameter (inverse distance units)
  cutoff = global cutoff for LJ (and Coulombic if only 1 arg) (distance units)
  cutoff2 = global cutoff for Coulombic (distance units)
\sphinxstyleemphasis{lj/cut/coul/long} args = cutoff (cutoff2)
  cutoff = global cutoff for LJ (and Coulombic if only 1 arg) (distance units)
  cutoff2 = global cutoff for Coulombic (optional) (distance units)
\sphinxstyleemphasis{lj/cut/coul/msm} args = cutoff (cutoff2)
  cutoff = global cutoff for LJ (and Coulombic if only 1 arg) (distance units)
  cutoff2 = global cutoff for Coulombic (optional) (distance units)
\sphinxstyleemphasis{lj/cut/coul/wolf} args = alpha cutoff (cutoff2)
  alpha = damping parameter (inverse distance units)
  cutoff = global cutoff for LJ (and Coulombic if only 2 arg) (distance units)
  cutoff2 = global cutoff for Coulombic (optional) (distance units)
\end{sphinxalltt}


\subsection{Examples}
\label{\detokenize{pair_lj_cut_coul:examples}}
\begin{sphinxVerbatim}[commandchars=\\\{\}]
\PYG{k}{pair\PYGZus{}style}\PYG{+w}{ }\PYG{n}{lj}\PYG{o}{/}\PYG{n}{cut}\PYG{o}{/}\PYG{n}{coul}\PYG{o}{/}\PYG{n}{cut}\PYG{+w}{ }\PYG{l+m}{10.0}
\PYG{k}{pair\PYGZus{}style}\PYG{+w}{ }\PYG{n}{lj}\PYG{o}{/}\PYG{n}{cut}\PYG{o}{/}\PYG{n}{coul}\PYG{o}{/}\PYG{n}{cut}\PYG{+w}{ }\PYG{l+m}{10.0}\PYG{+w}{ }\PYG{l+m}{8.0}
\PYG{k}{pair\PYGZus{}coeff}\PYG{+w}{ }\PYG{o}{*}\PYG{+w}{ }\PYG{o}{*}\PYG{+w}{ }\PYG{l+m}{100.0}\PYG{+w}{ }\PYG{l+m}{3.0}
\PYG{k}{pair\PYGZus{}coeff}\PYG{+w}{ }\PYG{l+m}{1}\PYG{+w}{ }\PYG{l+m}{1}\PYG{+w}{ }\PYG{l+m}{100.0}\PYG{+w}{ }\PYG{l+m}{3.5}\PYG{+w}{ }\PYG{l+m}{9.0}
\PYG{k}{pair\PYGZus{}coeff}\PYG{+w}{ }\PYG{l+m}{1}\PYG{+w}{ }\PYG{l+m}{1}\PYG{+w}{ }\PYG{l+m}{100.0}\PYG{+w}{ }\PYG{l+m}{3.5}\PYG{+w}{ }\PYG{l+m}{9.0}\PYG{+w}{ }\PYG{l+m}{9.0}

\PYG{k}{pair\PYGZus{}style}\PYG{+w}{ }\PYG{n}{lj}\PYG{o}{/}\PYG{n}{cut}\PYG{o}{/}\PYG{n}{coul}\PYG{o}{/}\PYG{n}{debye}\PYG{+w}{ }\PYG{l+m}{1.5}\PYG{+w}{ }\PYG{l+m}{3.0}
\PYG{k}{pair\PYGZus{}style}\PYG{+w}{ }\PYG{n}{lj}\PYG{o}{/}\PYG{n}{cut}\PYG{o}{/}\PYG{n}{coul}\PYG{o}{/}\PYG{n}{debye}\PYG{+w}{ }\PYG{l+m}{1.5}\PYG{+w}{ }\PYG{l+m}{2.5}\PYG{+w}{ }\PYG{l+m}{5.0}
\PYG{k}{pair\PYGZus{}coeff}\PYG{+w}{ }\PYG{o}{*}\PYG{+w}{ }\PYG{o}{*}\PYG{+w}{ }\PYG{l+m}{1.0}\PYG{+w}{ }\PYG{l+m}{1.0}
\PYG{k}{pair\PYGZus{}coeff}\PYG{+w}{ }\PYG{l+m}{1}\PYG{+w}{ }\PYG{l+m}{1}\PYG{+w}{ }\PYG{l+m}{1.0}\PYG{+w}{ }\PYG{l+m}{1.5}\PYG{+w}{ }\PYG{l+m}{2.5}
\PYG{k}{pair\PYGZus{}coeff}\PYG{+w}{ }\PYG{l+m}{1}\PYG{+w}{ }\PYG{l+m}{1}\PYG{+w}{ }\PYG{l+m}{1.0}\PYG{+w}{ }\PYG{l+m}{1.5}\PYG{+w}{ }\PYG{l+m}{2.5}\PYG{+w}{ }\PYG{l+m}{5.0}

\PYG{k}{pair\PYGZus{}style}\PYG{+w}{ }\PYG{n}{lj}\PYG{o}{/}\PYG{n}{cut}\PYG{o}{/}\PYG{n}{coul}\PYG{o}{/}\PYG{n}{dsf}\PYG{+w}{ }\PYG{l+m}{0.05}\PYG{+w}{ }\PYG{l+m}{2.5}\PYG{+w}{ }\PYG{l+m}{10.0}
\PYG{k}{pair\PYGZus{}coeff}\PYG{+w}{ }\PYG{o}{*}\PYG{+w}{ }\PYG{o}{*}\PYG{+w}{ }\PYG{l+m}{1.0}\PYG{+w}{ }\PYG{l+m}{1.0}
\PYG{k}{pair\PYGZus{}coeff}\PYG{+w}{ }\PYG{l+m}{1}\PYG{+w}{ }\PYG{l+m}{1}\PYG{+w}{ }\PYG{l+m}{1.0}\PYG{+w}{ }\PYG{l+m}{1.0}\PYG{+w}{ }\PYG{l+m}{2.5}

\PYG{k}{pair\PYGZus{}style}\PYG{+w}{ }\PYG{n}{lj}\PYG{o}{/}\PYG{n}{cut}\PYG{o}{/}\PYG{n}{coul}\PYG{o}{/}\PYG{n}{long}\PYG{+w}{ }\PYG{l+m}{10.0}
\PYG{k}{pair\PYGZus{}style}\PYG{+w}{ }\PYG{n}{lj}\PYG{o}{/}\PYG{n}{cut}\PYG{o}{/}\PYG{n}{coul}\PYG{o}{/}\PYG{n}{long}\PYG{+w}{ }\PYG{l+m}{10.0}\PYG{+w}{ }\PYG{l+m}{8.0}
\PYG{k}{pair\PYGZus{}coeff}\PYG{+w}{ }\PYG{o}{*}\PYG{+w}{ }\PYG{o}{*}\PYG{+w}{ }\PYG{l+m}{100.0}\PYG{+w}{ }\PYG{l+m}{3.0}
\PYG{k}{pair\PYGZus{}coeff}\PYG{+w}{ }\PYG{l+m}{1}\PYG{+w}{ }\PYG{l+m}{1}\PYG{+w}{ }\PYG{l+m}{100.0}\PYG{+w}{ }\PYG{l+m}{3.5}\PYG{+w}{ }\PYG{l+m}{9.0}

\PYG{k}{pair\PYGZus{}style}\PYG{+w}{ }\PYG{n}{lj}\PYG{o}{/}\PYG{n}{cut}\PYG{o}{/}\PYG{n}{coul}\PYG{o}{/}\PYG{n}{msm}\PYG{+w}{ }\PYG{l+m}{10.0}
\PYG{k}{pair\PYGZus{}style}\PYG{+w}{ }\PYG{n}{lj}\PYG{o}{/}\PYG{n}{cut}\PYG{o}{/}\PYG{n}{coul}\PYG{o}{/}\PYG{n}{msm}\PYG{+w}{ }\PYG{l+m}{10.0}\PYG{+w}{ }\PYG{l+m}{8.0}
\PYG{k}{pair\PYGZus{}coeff}\PYG{+w}{ }\PYG{o}{*}\PYG{+w}{ }\PYG{o}{*}\PYG{+w}{ }\PYG{l+m}{100.0}\PYG{+w}{ }\PYG{l+m}{3.0}
\PYG{k}{pair\PYGZus{}coeff}\PYG{+w}{ }\PYG{l+m}{1}\PYG{+w}{ }\PYG{l+m}{1}\PYG{+w}{ }\PYG{l+m}{100.0}\PYG{+w}{ }\PYG{l+m}{3.5}\PYG{+w}{ }\PYG{l+m}{9.0}

\PYG{k}{pair\PYGZus{}style}\PYG{+w}{ }\PYG{n}{lj}\PYG{o}{/}\PYG{n}{cut}\PYG{o}{/}\PYG{n}{coul}\PYG{o}{/}\PYG{n}{wolf}\PYG{+w}{ }\PYG{l+m}{0.2}\PYG{+w}{ }\PYG{l+m}{5}\PYG{n}{.}\PYG{+w}{ }\PYG{l+m}{10.0}
\PYG{k}{pair\PYGZus{}coeff}\PYG{+w}{ }\PYG{o}{*}\PYG{+w}{ }\PYG{o}{*}\PYG{+w}{ }\PYG{l+m}{1.0}\PYG{+w}{ }\PYG{l+m}{1.0}
\PYG{k}{pair\PYGZus{}coeff}\PYG{+w}{ }\PYG{l+m}{1}\PYG{+w}{ }\PYG{l+m}{1}\PYG{+w}{ }\PYG{l+m}{1.0}\PYG{+w}{ }\PYG{l+m}{1.0}\PYG{+w}{ }\PYG{l+m}{2.5}
\end{sphinxVerbatim}


\subsection{Description}
\label{\detokenize{pair_lj_cut_coul:description}}
\sphinxAtStartPar
The \sphinxstyleemphasis{lj/cut/coul} styles compute the standard 12/6 Lennard\sphinxhyphen{}Jones potential,
given by
\begin{equation*}
\begin{split}E = 4 \epsilon \left[ \left(\frac{\sigma}{r}\right)^{12} -
    \left(\frac{\sigma}{r}\right)^6 \right]
                    \qquad r < r_c\end{split}
\end{equation*}
\sphinxAtStartPar
\(r_c\) is the cutoff.

\sphinxAtStartPar
Style \sphinxstyleemphasis{lj/cut/coul/cut} adds a Coulombic pairwise interaction given by
\begin{equation*}
\begin{split}E = \frac{C q_i q_j}{\epsilon  r} \qquad r < r_c\end{split}
\end{equation*}
\sphinxAtStartPar
where \(C\) is an energy\sphinxhyphen{}conversion constant, \(q_i\) and \(q_j\)
are the charges on the two atoms, and \(\epsilon\) is the dielectric
constant which can be set by the {\hyperref[\detokenize{dielectric::doc}]{\sphinxcrossref{\DUrole{doc}{dielectric}}}} command.
If one cutoff is specified in the pair\_style command, it is used for
both the LJ and Coulombic terms.  If two cutoffs are specified, they are
used as cutoffs for the LJ and Coulombic terms respectively.

\sphinxAtStartPar
Style \sphinxstyleemphasis{lj/cut/coul/debye} adds an additional exp() damping factor
to the Coulombic term, given by
\begin{equation*}
\begin{split}E = \frac{C q_i q_j}{\epsilon  r} \exp(- \kappa r) \qquad r < r_c\end{split}
\end{equation*}
\sphinxAtStartPar
where \(\kappa\) is the inverse of the Debye length.  This potential
is another way to mimic the screening effect of a polar solvent.

\sphinxAtStartPar
Style \sphinxstyleemphasis{lj/cut/coul/dsf} computes the Coulombic term via the damped
shifted force model described in {\hyperref[\detokenize{pair_lj_cut_coul:fennell2}]{\sphinxcrossref{\DUrole{std,std-ref}{Fennell}}}}, given by:
\begin{equation*}
\begin{split}E =
 q_iq_j \left[ \frac{\mbox{erfc} (\alpha r)}{r} -  \frac{\mbox{erfc} (\alpha r_c)}{r_c} +
\left( \frac{\mbox{erfc} (\alpha r_c)}{r_c^2} +  \frac{2\alpha}{\sqrt{\pi}}\frac{\exp (-\alpha^2    r^2_c)}{r_c} \right)(r-r_c) \right] \qquad r < r_c\end{split}
\end{equation*}
\sphinxAtStartPar
where \(\alpha\) is the damping parameter and erfc() is the complementary
error\sphinxhyphen{}function. This potential is essentially a short\sphinxhyphen{}range,
spherically\sphinxhyphen{}truncated, charge\sphinxhyphen{}neutralized, shifted, pairwise \sphinxstyleemphasis{1/r}
summation.  The potential is based on Wolf summation, proposed as an
alternative to Ewald summation for condensed phase systems where
charge screening causes electrostatic interactions to become
effectively short\sphinxhyphen{}ranged. In order for the electrostatic sum to be
absolutely convergent, charge neutralization within the cutoff radius
is enforced by shifting the potential through placement of image
charges on the cutoff sphere. Convergence can often be improved by
setting \(\alpha\) to a small non\sphinxhyphen{}zero value.

\sphinxAtStartPar
Styles \sphinxstyleemphasis{lj/cut/coul/long} and \sphinxstyleemphasis{lj/cut/coul/msm} compute the same
Coulombic interactions as style \sphinxstyleemphasis{lj/cut/coul/cut} except that an
additional damping factor is applied to the Coulombic term so it can
be used in conjunction with the {\hyperref[\detokenize{kspace_style::doc}]{\sphinxcrossref{\DUrole{doc}{kspace\_style}}}}
command and its \sphinxstyleemphasis{ewald} or \sphinxstyleemphasis{pppm} option.  The Coulombic cutoff
specified for this style means that pairwise interactions within this
distance are computed directly; interactions outside that distance are
computed in reciprocal space.

\sphinxAtStartPar
Style \sphinxstyleemphasis{lj/cut/coul/wolf} adds a Coulombic pairwise interaction via the Wolf
summation method, described in {\hyperref[\detokenize{pair_lj_cut_coul:wolf3}]{\sphinxcrossref{\DUrole{std,std-ref}{Wolf}}}}, given by:
\begin{equation*}
\begin{split}E_i = \frac{1}{2} \sum_{j \neq i}
\frac{q_i q_j {\rm erfc}(\alpha r_{ij})}{r_{ij}} +
\frac{1}{2} \sum_{j \neq i}
\frac{q_i q_j {\rm erf}(\alpha r_{ij})}{r_{ij}} \qquad r < r_c\end{split}
\end{equation*}
\sphinxAtStartPar
where \(\alpha\) is the damping parameter, and erfc() is the
complementary error\sphinxhyphen{}function terms.  This potential is essentially a
short\sphinxhyphen{}range, spherically\sphinxhyphen{}truncated, charge\sphinxhyphen{}neutralized, shifted,
pairwise \sphinxstyleemphasis{1/r} summation.  With a manipulation of adding and subtracting
a self term (for i = j) to the first and second term on the
right\sphinxhyphen{}hand\sphinxhyphen{}side, respectively, and a small enough \(\alpha\) damping
parameter, the second term shrinks and the potential becomes a
rapidly\sphinxhyphen{}converging real\sphinxhyphen{}space summation.  With a long enough cutoff and
small enough \(\alpha\) parameter, the energy and forces calculated by the
Wolf summation method approach those of the Ewald sum.  So it is a means
of getting effective long\sphinxhyphen{}range interactions with a short\sphinxhyphen{}range
potential.


\subsection{Coefficients}
\label{\detokenize{pair_lj_cut_coul:coefficients}}
\sphinxAtStartPar
For all of the \sphinxstyleemphasis{lj/cut/coul} pair styles, the following coefficients must
be defined for each pair of atoms types via the
{\hyperref[\detokenize{pair_coeff::doc}]{\sphinxcrossref{\DUrole{doc}{pair\_coeff}}}} command as in the examples above, or in
the data file or restart files read by the {\hyperref[\detokenize{read_data::doc}]{\sphinxcrossref{\DUrole{doc}{read\_data}}}}
or {\hyperref[\detokenize{read_restart::doc}]{\sphinxcrossref{\DUrole{doc}{read\_restart}}}} commands, or by mixing as
described below:
\begin{itemize}
\item {} 
\sphinxAtStartPar
\(\epsilon\) (energy units)

\item {} 
\sphinxAtStartPar
\(\sigma\) (distance units)

\item {} 
\sphinxAtStartPar
cutoff1 (distance units)

\item {} 
\sphinxAtStartPar
cutoff2 (distance units)

\end{itemize}

\sphinxAtStartPar
Note that \(\sigma\) is defined in the LJ formula as the zero\sphinxhyphen{}crossing
distance for the potential, not as the energy minimum at \(2^{\frac{1}{6}} \sigma\).

\sphinxAtStartPar
The latter 2 coefficients are optional.  If not specified, the global
LJ and Coulombic cutoffs specified in the pair\_style command are used.
If only one cutoff is specified, it is used as the cutoff for both LJ
and Coulombic interactions for this type pair.  If both coefficients
are specified, they are used as the LJ and Coulombic cutoffs for this
type pair.

\sphinxAtStartPar
For \sphinxstyleemphasis{lj/cut/coul/long} and \sphinxstyleemphasis{lj/cut/coul/msm} only the LJ cutoff can be
specified since a Coulombic cutoff cannot be specified for an individual I,J
type pair.  All type pairs use the same global Coulombic cutoff specified in
the pair\_style command.


\bigskip\hrule\bigskip


\sphinxAtStartPar
A version of these styles with a soft core, \sphinxstyleemphasis{lj/cut/coul/soft}and
\sphinxstyleemphasis{lj/cut/coul/long/soft}, suitable for use in free energy calculations, is
part of the FEP package and is documented with the {\hyperref[\detokenize{pair_fep_soft::doc}]{\sphinxcrossref{\DUrole{doc}{pair\_style */soft}}}} styles.


\bigskip\hrule\bigskip


\sphinxAtStartPar
Styles with a \sphinxstyleemphasis{gpu}, \sphinxstyleemphasis{intel}, \sphinxstyleemphasis{kk}, \sphinxstyleemphasis{omp}, or \sphinxstyleemphasis{opt} suffix are
functionally the same as the corresponding style without the suffix.
They have been optimized to run faster, depending on your available
hardware, as discussed on the {\hyperref[\detokenize{Speed_packages::doc}]{\sphinxcrossref{\DUrole{doc}{Accelerator packages}}}}
page.  The accelerated styles take the same arguments and should
produce the same results, except for round\sphinxhyphen{}off and precision issues.

\sphinxAtStartPar
These accelerated styles are part of the GPU, INTEL, KOKKOS,
OPENMP, and OPT packages, respectively.  They are only enabled if
LAMMPS was built with those packages.  See the {\hyperref[\detokenize{Build_package::doc}]{\sphinxcrossref{\DUrole{doc}{Build package}}}} page for more info.

\sphinxAtStartPar
You can specify the accelerated styles explicitly in your input script
by including their suffix, or you can use the {\hyperref[\detokenize{Run_options::doc}]{\sphinxcrossref{\DUrole{doc}{\sphinxhyphen{}suffix command\sphinxhyphen{}line switch}}}} when you invoke LAMMPS, or you can use the
{\hyperref[\detokenize{suffix::doc}]{\sphinxcrossref{\DUrole{doc}{suffix}}}} command in your input script.

\sphinxAtStartPar
See the {\hyperref[\detokenize{Speed_packages::doc}]{\sphinxcrossref{\DUrole{doc}{Accelerator packages}}}} page for more
instructions on how to use the accelerated styles effectively.


\bigskip\hrule\bigskip



\subsection{Mixing, shift, table, tail correction, restart, rRESPA info}
\label{\detokenize{pair_lj_cut_coul:mixing-shift-table-tail-correction-restart-rrespa-info}}
\sphinxAtStartPar
For atom type pairs I,J and I != J, the epsilon and sigma coefficients
and cutoff distance for all of the lj/cut pair styles can be mixed.
The default mix value is \sphinxstyleemphasis{geometric}.  See the “pair\_modify” command
for details.

\sphinxAtStartPar
All of the \sphinxstyleemphasis{lj/cut} pair styles support the
{\hyperref[\detokenize{pair_modify::doc}]{\sphinxcrossref{\DUrole{doc}{pair\_modify}}}} shift option for the energy of the
Lennard\sphinxhyphen{}Jones portion of the pair interaction.

\sphinxAtStartPar
The \sphinxstyleemphasis{lj/cut/coul/long} pair styles support the
{\hyperref[\detokenize{pair_modify::doc}]{\sphinxcrossref{\DUrole{doc}{pair\_modify}}}} table option since they can tabulate
the short\sphinxhyphen{}range portion of the long\sphinxhyphen{}range Coulombic interaction.

\sphinxAtStartPar
All of the \sphinxstyleemphasis{lj/cut} pair styles support the
{\hyperref[\detokenize{pair_modify::doc}]{\sphinxcrossref{\DUrole{doc}{pair\_modify}}}} tail option for adding a long\sphinxhyphen{}range
tail correction to the energy and pressure for the Lennard\sphinxhyphen{}Jones
portion of the pair interaction.

\sphinxAtStartPar
All of the \sphinxstyleemphasis{lj/cut} pair styles write their information to {\hyperref[\detokenize{restart::doc}]{\sphinxcrossref{\DUrole{doc}{binary restart files}}}}, so pair\_style and pair\_coeff commands do
not need to be specified in an input script that reads a restart file.

\sphinxAtStartPar
The \sphinxstyleemphasis{lj/cut/coul/long} pair styles support the use of the
\sphinxstyleemphasis{inner}, \sphinxstyleemphasis{middle}, and \sphinxstyleemphasis{outer} keywords of the {\hyperref[\detokenize{run_style::doc}]{\sphinxcrossref{\DUrole{doc}{run\_style respa}}}} command, meaning the pairwise forces can be
partitioned by distance at different levels of the rRESPA hierarchy.
The other styles only support the \sphinxstyleemphasis{pair} keyword of run\_style respa.
See the {\hyperref[\detokenize{run_style::doc}]{\sphinxcrossref{\DUrole{doc}{run\_style}}}} command for details.


\bigskip\hrule\bigskip



\subsection{Restrictions}
\label{\detokenize{pair_lj_cut_coul:restrictions}}
\sphinxAtStartPar
The \sphinxstyleemphasis{lj/cut/coul/long} and \sphinxstyleemphasis{lj/cut/coul/msm} styles are part of the KSPACE package.

\sphinxAtStartPar
The \sphinxstyleemphasis{lj/cut/coul/debye}, \sphinxstyleemphasis{lj/cut/coul/dsf}, and \sphinxstyleemphasis{lj/cut/coul/wolf} styles are part
of the EXTRA\sphinxhyphen{}PAIR package.

\sphinxAtStartPar
These styles are only enabled if LAMMPS was built with those respective
packages.  See the {\hyperref[\detokenize{Build_package::doc}]{\sphinxcrossref{\DUrole{doc}{Build package}}}} page for
more info.


\subsection{Related commands}
\label{\detokenize{pair_lj_cut_coul:related-commands}}
\sphinxAtStartPar
{\hyperref[\detokenize{pair_coeff::doc}]{\sphinxcrossref{\DUrole{doc}{pair\_coeff}}}}


\subsection{Default}
\label{\detokenize{pair_lj_cut_coul:default}}
\sphinxAtStartPar
none


\bigskip\hrule\bigskip

\phantomsection\label{\detokenize{pair_lj_cut_coul:wolf3}}
\sphinxAtStartPar
\sphinxstylestrong{(Wolf)} D. Wolf, P. Keblinski, S. R. Phillpot, J. Eggebrecht, J Chem
Phys, 110, 8254 (1999).

\phantomsection\label{\detokenize{pair_lj_cut_coul:fennell2}}
\sphinxAtStartPar
\sphinxstylestrong{(Fennell)} C. J. Fennell, J. D. Gezelter, J Chem Phys, 124,
234104 (2006).

\sphinxstepscope

\index{pair\_style lj/cut/sphere@\spxentry{pair\_style lj/cut/sphere}}\ignorespaces 
\index{pair\_style lj/cut/sphere/omp@\spxentry{pair\_style lj/cut/sphere/omp}}\ignorespaces 

\section{pair\_style lj/cut/sphere command}
\label{\detokenize{pair_lj_cut_sphere:pair-style-lj-cut-sphere-command}}\label{\detokenize{pair_lj_cut_sphere:index-1}}\label{\detokenize{pair_lj_cut_sphere:index-0}}\label{\detokenize{pair_lj_cut_sphere::doc}}
\sphinxAtStartPar
Accelerator Variant: \sphinxstyleemphasis{lj/cut/sphere/omp}


\subsection{Syntax}
\label{\detokenize{pair_lj_cut_sphere:syntax}}
\begin{sphinxVerbatim}[commandchars=\\\{\}]
\PYG{k}{pair\PYGZus{}style}\PYG{+w}{ }\PYG{n}{style}\PYG{+w}{ }\PYG{n}{args}
\end{sphinxVerbatim}
\begin{itemize}
\item {} 
\sphinxAtStartPar
style = \sphinxstyleemphasis{lj/cut/sphere}

\item {} 
\sphinxAtStartPar
args = list of arguments for a particular style

\end{itemize}
\begin{sphinxalltt}
\sphinxstyleemphasis{lj/cut/sphere} args = cutoff ratio
  cutoff = global cutoff ratio for Lennard Jones interactions (unitless)
\end{sphinxalltt}


\subsection{Examples}
\label{\detokenize{pair_lj_cut_sphere:examples}}
\begin{sphinxVerbatim}[commandchars=\\\{\}]
\PYG{k}{pair\PYGZus{}style}\PYG{+w}{ }\PYG{n}{lj}\PYG{o}{/}\PYG{n}{cut}\PYG{o}{/}\PYG{n}{sphere}\PYG{+w}{ }\PYG{l+m}{2.5}
\PYG{k}{pair\PYGZus{}coeff}\PYG{+w}{ }\PYG{o}{*}\PYG{+w}{ }\PYG{o}{*}\PYG{+w}{ }\PYG{l+m}{1.0}
\PYG{k}{pair\PYGZus{}coeff}\PYG{+w}{ }\PYG{l+m}{1}\PYG{+w}{ }\PYG{l+m}{1}\PYG{+w}{ }\PYG{l+m}{1.1}\PYG{+w}{ }\PYG{l+m}{2.8}
\end{sphinxVerbatim}


\subsection{Description}
\label{\detokenize{pair_lj_cut_sphere:description}}
\sphinxAtStartPar
\DUrole{versionmodified,added}{Added in version 15Jun2023.}

\sphinxAtStartPar
The \sphinxstyleemphasis{lj/cut/sphere} style compute the standard 12/6 Lennard\sphinxhyphen{}Jones potential,
given by
\begin{equation*}
\begin{split}E = 4 \epsilon \left[ \left(\frac{\sigma_{ij}}{r}\right)^{12} -
    \left(\frac{\sigma_{ij}}{r}\right)^6 \right]  \qquad  r < r_c * \sigma_{ij}\end{split}
\end{equation*}
\sphinxAtStartPar
\(r_c\) is the cutoff ratio.

\sphinxAtStartPar
This is the same potential function used by the {\hyperref[\detokenize{pair_lj::doc}]{\sphinxcrossref{\DUrole{doc}{lj/cut}}}} pair style, but the \(\sigma_{ij}\) parameter is not
set as a per\sphinxhyphen{}type parameter via the {\hyperref[\detokenize{pair_coeff::doc}]{\sphinxcrossref{\DUrole{doc}{pair\_coeff command}}}}.  Instead it is calculated individually for each pair
using the per\sphinxhyphen{}atom diameter attribute of {\hyperref[\detokenize{atom_style::doc}]{\sphinxcrossref{\DUrole{doc}{atom\_style sphere}}}} for the two atoms as \(\sigma_{i}\) and
\(\sigma_{j}\); \(\sigma_{ij}\) is then computed by the mixing
rule for pair coefficients as set by the {\hyperref[\detokenize{pair_modify::doc}]{\sphinxcrossref{\DUrole{doc}{pair\_modify mix}}}} command (defaults to geometric mixing).  The cutoff is
not specified as a distance, but as ratio that is internally
multiplied by \(\sigma_{ij}\) to obtain the actual cutoff for each
pair of atoms.

\sphinxAtStartPar
Note that \(\sigma_{ij}\) is defined in the LJ formula above as the
zero\sphinxhyphen{}crossing distance for the potential, \sphinxstyleemphasis{not} as the energy minimum which
is at \(2^{\frac{1}{6}} \sigma_{ij}\).

\begin{sphinxadmonition}{note}{Notes on cutoffs, neighbor lists, and efficiency}

\sphinxAtStartPar
If your system is mildly polydisperse, meaning the ratio of the
diameter of the largest particle to the smallest is less than 2.0,
then the neighbor lists built by the code should be reasonably
efficient.  Which means they will not contain too many particle
pairs that do not interact.  However, if your system is highly
polydisperse (ratio \textgreater{} 2.0), the neighbor list build and force
computations may be inefficient.  There are two ways to try and
speed up the simulations.

\sphinxAtStartPar
The first is to assign atoms to different atom types so that atoms of
each type are similar in size.  E.g. if particle diameters range from
1 to 5 use 4 atom types, ensuring atoms of type 1 have diameters from
1.0\sphinxhyphen{}2.0, type 2 from 2.0\sphinxhyphen{}3.0, etc.  This will reduce the number of
non\sphinxhyphen{}interacting pairs in the neighbor lists and thus reduce the time
spent on computing pairwise interactions.

\sphinxAtStartPar
The second is to use the {\hyperref[\detokenize{neighbor::doc}]{\sphinxcrossref{\DUrole{doc}{neighbor multi}}}} command
which enabled a different algorithm for building neighbor lists.  This
will also require that you assign multiple atom types according to
diameters, but will in addition use a more efficient size\sphinxhyphen{}dependent
strategy to construct the neighbor lists and thus reduce the time
spent on building neighbor lists.

\sphinxAtStartPar
Here are example input script commands using both ideas for a
highly polydisperse system:

\begin{sphinxVerbatim}[commandchars=\\\{\}]
\PYG{n}{units}\PYG{+w}{           }\PYG{n}{lj}
\PYG{n}{atom\PYGZus{}style}\PYG{+w}{      }\PYG{n}{sphere}
\PYG{n}{lattice}\PYG{+w}{         }\PYG{n}{fcc}\PYG{+w}{ }\PYG{l+m+mf}{0.8442}
\PYG{n}{region}\PYG{+w}{          }\PYG{n}{box}\PYG{+w}{ }\PYG{n}{block}\PYG{+w}{ }\PYG{l+m+mi}{0}\PYG{+w}{ }\PYG{l+m+mi}{10}\PYG{+w}{ }\PYG{l+m+mi}{0}\PYG{+w}{ }\PYG{l+m+mi}{10}\PYG{+w}{ }\PYG{l+m+mi}{0}\PYG{+w}{ }\PYG{l+m+mi}{10}
\PYG{n}{create\PYGZus{}box}\PYG{+w}{      }\PYG{l+m+mi}{2}\PYG{+w}{ }\PYG{n}{box}
\PYG{n}{create\PYGZus{}atoms}\PYG{+w}{    }\PYG{l+m+mi}{1}\PYG{+w}{ }\PYG{n}{box}

\PYG{c+cp}{\PYGZsh{}}\PYG{c+cp}{ create atoms with random diameters from bimodal distribution}
\PYG{n}{variable}\PYG{+w}{ }\PYG{k}{switch}\PYG{+w}{ }\PYG{n}{atom}\PYG{+w}{ }\PYG{n}{random}\PYG{p}{(}\PYG{l+m+mf}{0.0}\PYG{p}{,}\PYG{l+m+mf}{1.0}\PYG{p}{,}\PYG{l+m+mi}{345634}\PYG{p}{)}
\PYG{n}{variable}\PYG{+w}{ }\PYG{n}{diam}\PYG{+w}{ }\PYG{n}{atom}\PYG{+w}{ }\PYG{p}{(}\PYG{n}{v\PYGZus{}switch}\PYG{o}{\PYGZlt{}}\PYG{l+m+mf}{0.75}\PYG{p}{)}\PYG{o}{*}\PYG{n}{normal}\PYG{p}{(}\PYG{l+m+mf}{0.4}\PYG{p}{,}\PYG{l+m+mf}{0.075}\PYG{p}{,}\PYG{l+m+mi}{325}\PYG{p}{)}\PYG{o}{+}\PYG{p}{(}\PYG{n}{v\PYGZus{}switch}\PYG{o}{\PYGZgt{}}\PYG{o}{=}\PYG{l+m+mf}{0.7}\PYG{p}{)}\PYG{o}{*}\PYG{n}{normal}\PYG{p}{(}\PYG{l+m+mf}{1.2}\PYG{p}{,}\PYG{l+m+mf}{0.2}\PYG{p}{,}\PYG{l+m+mi}{453}\PYG{p}{)}
\PYG{n}{set}\PYG{+w}{ }\PYG{n}{group}\PYG{+w}{ }\PYG{n}{all}\PYG{+w}{ }\PYG{n}{diameter}\PYG{+w}{ }\PYG{n}{v\PYGZus{}diam}

\PYG{c+cp}{\PYGZsh{}}\PYG{c+cp}{ assign type 2 to atoms with diameter \PYGZgt{} 0.6}
\PYG{n}{variable}\PYG{+w}{ }\PYG{n}{large}\PYG{+w}{ }\PYG{n}{atom}\PYG{+w}{ }\PYG{p}{(}\PYG{l+m+mf}{2.0}\PYG{o}{*}\PYG{n}{radius}\PYG{p}{)}\PYG{o}{\PYGZgt{}}\PYG{l+m+mf}{0.6}
\PYG{n}{group}\PYG{+w}{ }\PYG{n}{large}\PYG{+w}{ }\PYG{n}{variable}\PYG{+w}{ }\PYG{n}{large}
\PYG{n}{set}\PYG{+w}{ }\PYG{n}{group}\PYG{+w}{ }\PYG{n}{large}\PYG{+w}{ }\PYG{n}{type}\PYG{+w}{ }\PYG{l+m+mi}{2}

\PYG{n}{pair\PYGZus{}style}\PYG{+w}{      }\PYG{n}{lj}\PYG{o}{/}\PYG{n}{cut}\PYG{o}{/}\PYG{n}{sphere}\PYG{+w}{ }\PYG{l+m+mf}{2.5}
\PYG{n}{pair\PYGZus{}coeff}\PYG{+w}{      }\PYG{o}{*}\PYG{+w}{ }\PYG{o}{*}\PYG{+w}{ }\PYG{l+m+mf}{1.0}

\PYG{n}{neighbor}\PYG{+w}{ }\PYG{l+m+mf}{0.3}\PYG{+w}{ }\PYG{n}{multi}
\end{sphinxVerbatim}

\sphinxAtStartPar
Using multiple atom types speeds up the calculation for this example
by more than a factor of 2, and using the multi\sphinxhyphen{}style neighbor list
build causes an additional speedup of about 20 percent.
\end{sphinxadmonition}


\subsection{Coefficients}
\label{\detokenize{pair_lj_cut_sphere:coefficients}}
\sphinxAtStartPar
The following coefficients must be defined for each pair of atoms types via the
{\hyperref[\detokenize{pair_coeff::doc}]{\sphinxcrossref{\DUrole{doc}{pair\_coeff}}}} command as in the examples above, or in the data
file or restart files read by the {\hyperref[\detokenize{read_data::doc}]{\sphinxcrossref{\DUrole{doc}{read\_data}}}} or
{\hyperref[\detokenize{read_restart::doc}]{\sphinxcrossref{\DUrole{doc}{read\_restart}}}} commands, or by mixing as described below:
\begin{itemize}
\item {} 
\sphinxAtStartPar
\(\epsilon\) (energy units)

\item {} 
\sphinxAtStartPar
LJ cutoff ratio (unitless) (optional)

\end{itemize}

\sphinxAtStartPar
The last coefficient is optional.  If not specified, the global LJ
cutoff ratio specified in the {\hyperref[\detokenize{pair_style::doc}]{\sphinxcrossref{\DUrole{doc}{pair\_style command}}}} is
used.

\sphinxAtStartPar
If a repulsive only LJ interaction is desired, the coefficient for the cutoff
ratio should be set to the minimum of the LJ potential using \sphinxcode{\sphinxupquote{\$(2.0\textasciicircum{}(1.0/6.0))}}


\bigskip\hrule\bigskip


\sphinxAtStartPar
Styles with a \sphinxstyleemphasis{gpu}, \sphinxstyleemphasis{intel}, \sphinxstyleemphasis{kk}, \sphinxstyleemphasis{omp}, or \sphinxstyleemphasis{opt} suffix are
functionally the same as the corresponding style without the suffix.
They have been optimized to run faster, depending on your available
hardware, as discussed on the {\hyperref[\detokenize{Speed_packages::doc}]{\sphinxcrossref{\DUrole{doc}{Accelerator packages}}}}
page.  The accelerated styles take the same arguments and should
produce the same results, except for round\sphinxhyphen{}off and precision issues.

\sphinxAtStartPar
These accelerated styles are part of the GPU, INTEL, KOKKOS,
OPENMP, and OPT packages, respectively.  They are only enabled if
LAMMPS was built with those packages.  See the {\hyperref[\detokenize{Build_package::doc}]{\sphinxcrossref{\DUrole{doc}{Build package}}}} page for more info.

\sphinxAtStartPar
You can specify the accelerated styles explicitly in your input script
by including their suffix, or you can use the {\hyperref[\detokenize{Run_options::doc}]{\sphinxcrossref{\DUrole{doc}{\sphinxhyphen{}suffix command\sphinxhyphen{}line switch}}}} when you invoke LAMMPS, or you can use the
{\hyperref[\detokenize{suffix::doc}]{\sphinxcrossref{\DUrole{doc}{suffix}}}} command in your input script.

\sphinxAtStartPar
See the {\hyperref[\detokenize{Speed_packages::doc}]{\sphinxcrossref{\DUrole{doc}{Accelerator packages}}}} page for more
instructions on how to use the accelerated styles effectively.


\bigskip\hrule\bigskip



\subsection{Mixing, shift, table, tail correction, restart, rRESPA info}
\label{\detokenize{pair_lj_cut_sphere:mixing-shift-table-tail-correction-restart-rrespa-info}}
\sphinxAtStartPar
For atom type pairs I,J and I != J, the epsilon coefficients and cutoff
ratio for the \sphinxstyleemphasis{lj/cut/sphere} pair style can be mixed.  The default mixing
style is \sphinxstyleemphasis{geometric}.  See the {\hyperref[\detokenize{pair_modify::doc}]{\sphinxcrossref{\DUrole{doc}{pair\_modify command}}}}
for details.

\sphinxAtStartPar
The \sphinxstyleemphasis{lj/cut/sphere} pair style supports the {\hyperref[\detokenize{pair_modify::doc}]{\sphinxcrossref{\DUrole{doc}{pair\_modify shift}}}}
option for the energy of the Lennard\sphinxhyphen{}Jones portion of the pair interaction.

\sphinxAtStartPar
The \sphinxstyleemphasis{lj/cut/sphere} pair style does \sphinxstyleemphasis{not} support the {\hyperref[\detokenize{pair_modify::doc}]{\sphinxcrossref{\DUrole{doc}{pair\_modify}}}} tail option for adding a long\sphinxhyphen{}range tail corrections to
the energy and pressure.

\sphinxAtStartPar
The \sphinxstyleemphasis{lj/cut/sphere} pair style writes its information to {\hyperref[\detokenize{restart::doc}]{\sphinxcrossref{\DUrole{doc}{binary
restart files}}}}, so pair\_style and pair\_coeff commands do not
need to be specified in an input script that reads a restart file.

\sphinxAtStartPar
This pair style can only be used via the \sphinxstyleemphasis{pair} keyword of the
{\hyperref[\detokenize{run_style::doc}]{\sphinxcrossref{\DUrole{doc}{run\_style respa}}}} command.  It does \sphinxstyleemphasis{not} support the
\sphinxstyleemphasis{inner}, \sphinxstyleemphasis{middle}, \sphinxstyleemphasis{outer} keywords.


\bigskip\hrule\bigskip



\subsection{Restrictions}
\label{\detokenize{pair_lj_cut_sphere:restrictions}}
\sphinxAtStartPar
The \sphinxstyleemphasis{lj/cut/sphere} pair style is only enabled if LAMMPS was built with the
EXTRA\sphinxhyphen{}PAIR package.  See the {\hyperref[\detokenize{Build_package::doc}]{\sphinxcrossref{\DUrole{doc}{Build package}}}} page
for more info.

\sphinxAtStartPar
The \sphinxstyleemphasis{lj/cut/sphere} pair style does not support the \sphinxstyleemphasis{sixthpower} mixing rule.


\bigskip\hrule\bigskip



\subsection{Related commands}
\label{\detokenize{pair_lj_cut_sphere:related-commands}}\begin{itemize}
\item {} 
\sphinxAtStartPar
{\hyperref[\detokenize{pair_coeff::doc}]{\sphinxcrossref{\DUrole{doc}{pair\_coeff}}}}

\item {} 
\sphinxAtStartPar
{\hyperref[\detokenize{pair_lj::doc}]{\sphinxcrossref{\DUrole{doc}{pair\_style lj/cut}}}}

\item {} 
\sphinxAtStartPar
{\hyperref[\detokenize{pair_lj_expand_sphere::doc}]{\sphinxcrossref{\DUrole{doc}{pair\_style lj/expnd/sphere}}}}

\end{itemize}


\subsection{Default}
\label{\detokenize{pair_lj_cut_sphere:default}}
\sphinxAtStartPar
none

\sphinxstepscope

\index{pair\_style lj/cut/tip4p/cut@\spxentry{pair\_style lj/cut/tip4p/cut}}\ignorespaces 
\index{pair\_style lj/cut/tip4p/cut/omp@\spxentry{pair\_style lj/cut/tip4p/cut/omp}}\ignorespaces 
\index{pair\_style lj/cut/tip4p/long@\spxentry{pair\_style lj/cut/tip4p/long}}\ignorespaces 
\index{pair\_style lj/cut/tip4p/long/gpu@\spxentry{pair\_style lj/cut/tip4p/long/gpu}}\ignorespaces 
\index{pair\_style lj/cut/tip4p/long/omp@\spxentry{pair\_style lj/cut/tip4p/long/omp}}\ignorespaces 
\index{pair\_style lj/cut/tip4p/long/opt@\spxentry{pair\_style lj/cut/tip4p/long/opt}}\ignorespaces 

\section{pair\_style lj/cut/tip4p/cut command}
\label{\detokenize{pair_lj_cut_tip4p:pair-style-lj-cut-tip4p-cut-command}}\label{\detokenize{pair_lj_cut_tip4p:index-5}}\label{\detokenize{pair_lj_cut_tip4p:index-4}}\label{\detokenize{pair_lj_cut_tip4p:index-3}}\label{\detokenize{pair_lj_cut_tip4p:index-2}}\label{\detokenize{pair_lj_cut_tip4p:index-1}}\label{\detokenize{pair_lj_cut_tip4p:index-0}}\label{\detokenize{pair_lj_cut_tip4p::doc}}
\sphinxAtStartPar
Accelerator Variants: \sphinxstyleemphasis{lj/cut/tip4p/cut/omp}


\section{pair\_style lj/cut/tip4p/long command}
\label{\detokenize{pair_lj_cut_tip4p:pair-style-lj-cut-tip4p-long-command}}
\sphinxAtStartPar
Accelerator Variants: \sphinxstyleemphasis{lj/cut/tip4p/long/gpu}, \sphinxstyleemphasis{lj/cut/tip4p/long/omp}, \sphinxstyleemphasis{lj/cut/tip4p/long/opt}


\subsection{Syntax}
\label{\detokenize{pair_lj_cut_tip4p:syntax}}
\begin{sphinxVerbatim}[commandchars=\\\{\}]
\PYG{k}{pair\PYGZus{}style}\PYG{+w}{ }\PYG{n}{style}\PYG{+w}{ }\PYG{n}{args}
\end{sphinxVerbatim}
\begin{itemize}
\item {} 
\sphinxAtStartPar
style = \sphinxstyleemphasis{lj/cut/tip4p/cut} or \sphinxstyleemphasis{lj/cut/tip4p/long}

\item {} 
\sphinxAtStartPar
args = list of arguments for a particular style

\end{itemize}
\begin{sphinxalltt}
\sphinxstyleemphasis{lj/cut/tip4p/cut} args = otype htype btype atype qdist cutoff (cutoff2)
  otype,htype = atom types (numeric or type label) for TIP4P O and H
  btype,atype = bond and angle types (numeric or type label) for TIP4P waters
  qdist = distance from O atom to massless charge (distance units)
  cutoff = global cutoff for LJ (and Coulombic if only 1 arg) (distance units)
  cutoff2 = global cutoff for Coulombic (optional) (distance units)
\sphinxstyleemphasis{lj/cut/tip4p/long} args = otype htype btype atype qdist cutoff (cutoff2)
  otype,htype = atom types (numeric or type label) for TIP4P O and H
  btype,atype = bond and angle types (numeric or type label) for TIP4P waters
  qdist = distance from O atom to massless charge (distance units)
  cutoff = global cutoff for LJ (and Coulombic if only 1 arg) (distance units)
  cutoff2 = global cutoff for Coulombic (optional) (distance units)
\end{sphinxalltt}


\subsection{Examples}
\label{\detokenize{pair_lj_cut_tip4p:examples}}
\begin{sphinxVerbatim}[commandchars=\\\{\}]
\PYG{k}{pair\PYGZus{}style}\PYG{+w}{ }\PYG{n}{lj}\PYG{o}{/}\PYG{n}{cut}\PYG{o}{/}\PYG{n}{tip4p}\PYG{o}{/}\PYG{n}{cut}\PYG{+w}{ }\PYG{l+m}{1}\PYG{+w}{ }\PYG{l+m}{2}\PYG{+w}{ }\PYG{l+m}{7}\PYG{+w}{ }\PYG{l+m}{8}\PYG{+w}{ }\PYG{l+m}{0.15}\PYG{+w}{ }\PYG{l+m}{12.0}
\PYG{k}{pair\PYGZus{}style}\PYG{+w}{ }\PYG{n}{lj}\PYG{o}{/}\PYG{n}{cut}\PYG{o}{/}\PYG{n}{tip4p}\PYG{o}{/}\PYG{n}{cut}\PYG{+w}{ }\PYG{l+m}{1}\PYG{+w}{ }\PYG{l+m}{2}\PYG{+w}{ }\PYG{l+m}{7}\PYG{+w}{ }\PYG{l+m}{8}\PYG{+w}{ }\PYG{l+m}{0.15}\PYG{+w}{ }\PYG{l+m}{12.0}\PYG{+w}{ }\PYG{l+m}{10.0}
\PYG{k}{pair\PYGZus{}coeff}\PYG{+w}{ }\PYG{o}{*}\PYG{+w}{ }\PYG{o}{*}\PYG{+w}{ }\PYG{l+m}{100.0}\PYG{+w}{ }\PYG{l+m}{3.0}
\PYG{k}{pair\PYGZus{}coeff}\PYG{+w}{ }\PYG{l+m}{1}\PYG{+w}{ }\PYG{l+m}{1}\PYG{+w}{ }\PYG{l+m}{100.0}\PYG{+w}{ }\PYG{l+m}{3.5}\PYG{+w}{ }\PYG{l+m}{9.0}

\PYG{k}{pair\PYGZus{}style}\PYG{+w}{ }\PYG{n}{lj}\PYG{o}{/}\PYG{n}{cut}\PYG{o}{/}\PYG{n}{tip4p}\PYG{o}{/}\PYG{n}{long}\PYG{+w}{ }\PYG{l+m}{1}\PYG{+w}{ }\PYG{l+m}{2}\PYG{+w}{ }\PYG{l+m}{7}\PYG{+w}{ }\PYG{l+m}{8}\PYG{+w}{ }\PYG{l+m}{0.15}\PYG{+w}{ }\PYG{l+m}{12.0}
\PYG{k}{pair\PYGZus{}style}\PYG{+w}{ }\PYG{n}{lj}\PYG{o}{/}\PYG{n}{cut}\PYG{o}{/}\PYG{n}{tip4p}\PYG{o}{/}\PYG{n}{long}\PYG{+w}{ }\PYG{l+m}{1}\PYG{+w}{ }\PYG{l+m}{2}\PYG{+w}{ }\PYG{l+m}{7}\PYG{+w}{ }\PYG{l+m}{8}\PYG{+w}{ }\PYG{l+m}{0.15}\PYG{+w}{ }\PYG{l+m}{12.0}\PYG{+w}{ }\PYG{l+m}{10.0}
\PYG{k}{pair\PYGZus{}coeff}\PYG{+w}{ }\PYG{o}{*}\PYG{+w}{ }\PYG{o}{*}\PYG{+w}{ }\PYG{l+m}{100.0}\PYG{+w}{ }\PYG{l+m}{3.0}
\PYG{k}{pair\PYGZus{}coeff}\PYG{+w}{ }\PYG{l+m}{1}\PYG{+w}{ }\PYG{l+m}{1}\PYG{+w}{ }\PYG{l+m}{100.0}\PYG{+w}{ }\PYG{l+m}{3.5}\PYG{+w}{ }\PYG{l+m}{9.0}

\PYG{k}{pair\PYGZus{}style}\PYG{+w}{ }\PYG{n}{lj}\PYG{o}{/}\PYG{n}{cut}\PYG{o}{/}\PYG{n}{tip4p}\PYG{o}{/}\PYG{n}{long}\PYG{+w}{ }\PYG{n}{OW}\PYG{+w}{ }\PYG{n}{HW}\PYG{+w}{ }\PYG{n}{HW}\PYG{o}{\PYGZhy{}}\PYG{n}{OW}\PYG{+w}{ }\PYG{n}{HW}\PYG{o}{\PYGZhy{}}\PYG{n}{OW}\PYG{o}{\PYGZhy{}}\PYG{n}{HW}\PYG{+w}{ }\PYG{l+m}{0.15}\PYG{+w}{ }\PYG{l+m}{12.0}
\PYG{k}{labelmap}\PYG{+w}{ }\PYG{n}{atom}\PYG{+w}{ }\PYG{l+m}{1}\PYG{+w}{ }\PYG{n}{OW}\PYG{+w}{ }\PYG{l+m}{2}\PYG{+w}{ }\PYG{n}{HW}
\PYG{k}{labelmap}\PYG{+w}{ }\PYG{n}{bond}\PYG{+w}{ }\PYG{l+m}{1}\PYG{+w}{ }\PYG{n}{HW}\PYG{o}{\PYGZhy{}}\PYG{n}{OW}
\PYG{k}{labelmap}\PYG{+w}{ }\PYG{n}{angle}\PYG{+w}{ }\PYG{l+m}{1}\PYG{+w}{ }\PYG{n}{HW}\PYG{o}{\PYGZhy{}}\PYG{n}{OW}\PYG{o}{\PYGZhy{}}\PYG{n}{HW}
\PYG{k}{pair\PYGZus{}coeff}\PYG{+w}{ }\PYG{o}{*}\PYG{+w}{ }\PYG{o}{*}\PYG{+w}{ }\PYG{l+m}{100.0}\PYG{+w}{ }\PYG{l+m}{3.0}
\PYG{k}{pair\PYGZus{}coeff}\PYG{+w}{ }\PYG{n}{OW}\PYG{+w}{ }\PYG{n}{OW}\PYG{+w}{ }\PYG{l+m}{100.0}\PYG{+w}{ }\PYG{l+m}{3.5}\PYG{+w}{ }\PYG{l+m}{9.0}
\end{sphinxVerbatim}


\subsection{Description}
\label{\detokenize{pair_lj_cut_tip4p:description}}
\sphinxAtStartPar
The \sphinxstyleemphasis{lj/cut/tip4p} styles implement the TIP4P water model of
{\hyperref[\detokenize{pair_lj_cut_tip4p:jorgensen2}]{\sphinxcrossref{\DUrole{std,std-ref}{(Jorgensen)}}}} and similar models, which introduce a
massless site M located a short distance away from the oxygen atom along
the bisector of the HOH angle.  The atomic types of the oxygen and
hydrogen atoms, the bond and angle types for OH and HOH interactions,
and the distance to the massless charge site are specified as pair\_style
arguments and are used to identify the TIP4P\sphinxhyphen{}like molecules and
determine the position of the M site from the positions of the hydrogen
and oxygen atoms of the water molecules.  The M site location is used
for all Coulomb interactions instead of the oxygen atom location, also
with all other atom types, while the location of the oxygen atom is used
for the Lennard\sphinxhyphen{}Jones interactions.  Style \sphinxstyleemphasis{lj/cut/tip4p/cut} uses a
cutoff for Coulomb interactions; style \sphinxstyleemphasis{lj/cut/tip4p/long} is for use
with a long\sphinxhyphen{}range Coulombic solver (Ewald or PPPM).

\begin{sphinxadmonition}{note}{Note:}
\sphinxAtStartPar
For each TIP4P water molecule in your system, the atom IDs for
the O and 2 H atoms must be consecutive, with the O atom first.  This
is to enable LAMMPS to “find” the 2 H atoms associated with each O
atom.  For example, if the atom ID of an O atom in a TIP4P water
molecule is 500, then its 2 H atoms must have IDs 501 and 502.
\end{sphinxadmonition}

\begin{sphinxadmonition}{note}{Note:}
\sphinxAtStartPar
If using type labels, the type labels must be defined before calling
the {\hyperref[\detokenize{pair_coeff::doc}]{\sphinxcrossref{\DUrole{doc}{pair\_coeff}}}} command.
\end{sphinxadmonition}

\sphinxAtStartPar
See the {\hyperref[\detokenize{Howto_tip4p::doc}]{\sphinxcrossref{\DUrole{doc}{Howto tip4p}}}} page for more information
on how to use the TIP4P pair styles and lists of parameters to set.
Note that the neighbor list cutoff for Coulomb interactions is
effectively extended by a distance 2*qdist when using the TIP4P pair
style, to account for the offset distance of the fictitious charges on
O atoms in water molecules.  Thus it is typically best in an
efficiency sense to use a LJ cutoff \textgreater{}= Coulombic cutoff + 2*qdist, to
shrink the size of the neighbor list.  This leads to slightly larger
cost for the long\sphinxhyphen{}range calculation, so you can test the trade\sphinxhyphen{}off for
your model.

\sphinxAtStartPar
The \sphinxstyleemphasis{lj/cut/tip4p} styles compute the standard 12/6 Lennard\sphinxhyphen{}Jones potential,
given by
\begin{equation*}
\begin{split}E = 4 \epsilon \left[ \left(\frac{\sigma}{r}\right)^{12} -
    \left(\frac{\sigma}{r}\right)^6 \right]
                    \qquad r < r_c\end{split}
\end{equation*}
\sphinxAtStartPar
\(r_c\) is the cutoff.

\sphinxAtStartPar
They add Coulombic pairwise interactions given by
\begin{equation*}
\begin{split}E = \frac{C q_i q_j}{\epsilon  r} \qquad r < r_c\end{split}
\end{equation*}
\sphinxAtStartPar
where \(C\) is an energy\sphinxhyphen{}conversion constant, \(q_i\) and \(q_j\)
are the charges on the two atoms, and \(\epsilon\) is the dielectric
constant which can be set by the {\hyperref[\detokenize{dielectric::doc}]{\sphinxcrossref{\DUrole{doc}{dielectric}}}} command.
If one cutoff is specified in the pair\_style command, it is used for
both the LJ and Coulombic terms.  If two cutoffs are specified, they are
used as cutoffs for the LJ and Coulombic terms respectively.

\sphinxAtStartPar
Style \sphinxstyleemphasis{lj/cut/tip4p/long} compute the same
Coulombic interactions as style \sphinxstyleemphasis{lj/cut/tip4p/cut} except that an
additional damping factor is applied to the Coulombic term so it can
be used in conjunction with the {\hyperref[\detokenize{kspace_style::doc}]{\sphinxcrossref{\DUrole{doc}{kspace\_style}}}}
command and its \sphinxstyleemphasis{ewald} or \sphinxstyleemphasis{pppm} option.  The Coulombic cutoff
specified for this style means that pairwise interactions within this
distance are computed directly; interactions outside that distance are
computed in reciprocal space.


\subsection{Coefficients}
\label{\detokenize{pair_lj_cut_tip4p:coefficients}}
\sphinxAtStartPar
For all of the \sphinxstyleemphasis{lj/cut} pair styles, the following coefficients must
be defined for each pair of atoms types via the
{\hyperref[\detokenize{pair_coeff::doc}]{\sphinxcrossref{\DUrole{doc}{pair\_coeff}}}} command as in the examples above, or in
the data file or restart files read by the {\hyperref[\detokenize{read_data::doc}]{\sphinxcrossref{\DUrole{doc}{read\_data}}}}
or {\hyperref[\detokenize{read_restart::doc}]{\sphinxcrossref{\DUrole{doc}{read\_restart}}}} commands, or by mixing as
described below:
\begin{itemize}
\item {} 
\sphinxAtStartPar
\(\epsilon\) (energy units)

\item {} 
\sphinxAtStartPar
\(\sigma\) (distance units)

\item {} 
\sphinxAtStartPar
LJ cutoff (distance units)

\end{itemize}

\sphinxAtStartPar
Note that \(\sigma\) is defined in the LJ formula as the zero\sphinxhyphen{}crossing
distance for the potential, not as the energy minimum at \(2^{\frac{1}{6}} \sigma\).

\sphinxAtStartPar
The last coefficient is optional.  If not specified, the global
LJ cutoff specified in the pair\_style command is used.

\sphinxAtStartPar
For \sphinxstyleemphasis{lj/cut/tip4p/cut} and \sphinxstyleemphasis{lj/cut/tip4p/long} only the LJ cutoff can be
specified since a Coulombic cutoff cannot be specified for an individual I,J
type pair. All type pairs use the same global Coulombic cutoff specified in
the pair\_style command.

\begin{sphinxadmonition}{warning}{Warning:}
\sphinxAtStartPar
Because of how these pair styles implement the coulomb interactions
by implicitly defining a fourth site for the negative charge
of the TIP4P and similar water models, special care must be taken
when using these pair styles with other computations that also use
charges.  Unless they are specially set up to also handle the implicit
definition of the 4th site, results are likely incorrect.  Example:
{\hyperref[\detokenize{compute_dipole_chunk::doc}]{\sphinxcrossref{\DUrole{doc}{compute dipole/chunk}}}}.  For the same
reason, when using one of these pair styles with
{\hyperref[\detokenize{pair_hybrid::doc}]{\sphinxcrossref{\DUrole{doc}{pair\_style  hybrid}}}}, \sphinxstylestrong{all} coulomb interactions
should be handled by a single sub\sphinxhyphen{}style with TIP4P support. All other
instances and styles will “see” the M point charges at the position
of the Oxygen atom and thus compute incorrect forces and energies.
LAMMPS will print a warning when it detects one of these issues.
\end{sphinxadmonition}


\bigskip\hrule\bigskip


\sphinxAtStartPar
A version of these styles with a soft core, \sphinxstyleemphasis{lj/cut/tip4p/long/soft}, suitable
for use in free energy calculations, is part of the FEP package and
is documented with the {\hyperref[\detokenize{pair_fep_soft::doc}]{\sphinxcrossref{\DUrole{doc}{pair\_style */soft}}}}
styles.


\bigskip\hrule\bigskip


\sphinxAtStartPar
Styles with a \sphinxstyleemphasis{gpu}, \sphinxstyleemphasis{intel}, \sphinxstyleemphasis{kk}, \sphinxstyleemphasis{omp}, or \sphinxstyleemphasis{opt} suffix are
functionally the same as the corresponding style without the suffix.
They have been optimized to run faster, depending on your available
hardware, as discussed on the {\hyperref[\detokenize{Speed_packages::doc}]{\sphinxcrossref{\DUrole{doc}{Accelerator packages}}}}
page.  The accelerated styles take the same arguments and should
produce the same results, except for round\sphinxhyphen{}off and precision issues.

\sphinxAtStartPar
These accelerated styles are part of the GPU, INTEL, KOKKOS,
OPENMP, and OPT packages, respectively.  They are only enabled if
LAMMPS was built with those packages.  See the {\hyperref[\detokenize{Build_package::doc}]{\sphinxcrossref{\DUrole{doc}{Build package}}}} page for more info.

\sphinxAtStartPar
You can specify the accelerated styles explicitly in your input script
by including their suffix, or you can use the {\hyperref[\detokenize{Run_options::doc}]{\sphinxcrossref{\DUrole{doc}{\sphinxhyphen{}suffix command\sphinxhyphen{}line switch}}}} when you invoke LAMMPS, or you can use the
{\hyperref[\detokenize{suffix::doc}]{\sphinxcrossref{\DUrole{doc}{suffix}}}} command in your input script.

\sphinxAtStartPar
See the {\hyperref[\detokenize{Speed_packages::doc}]{\sphinxcrossref{\DUrole{doc}{Accelerator packages}}}} page for more
instructions on how to use the accelerated styles effectively.


\bigskip\hrule\bigskip



\subsection{Mixing, shift, table, tail correction, restart, rRESPA info}
\label{\detokenize{pair_lj_cut_tip4p:mixing-shift-table-tail-correction-restart-rrespa-info}}
\sphinxAtStartPar
For atom type pairs I,J and I != J, the epsilon and sigma coefficients
and cutoff distance for all of the lj/cut pair styles can be mixed.
The default mix value is \sphinxstyleemphasis{geometric}.  See the “pair\_modify” command
for details.

\sphinxAtStartPar
All of the \sphinxstyleemphasis{lj/cut} pair styles support the
{\hyperref[\detokenize{pair_modify::doc}]{\sphinxcrossref{\DUrole{doc}{pair\_modify}}}} shift option for the energy of the
Lennard\sphinxhyphen{}Jones portion of the pair interaction.

\sphinxAtStartPar
The \sphinxstyleemphasis{lj/cut/coul/long} and \sphinxstyleemphasis{lj/cut/tip4p/long} pair styles support the
{\hyperref[\detokenize{pair_modify::doc}]{\sphinxcrossref{\DUrole{doc}{pair\_modify}}}} table option since they can tabulate
the short\sphinxhyphen{}range portion of the long\sphinxhyphen{}range Coulombic interaction.

\sphinxAtStartPar
All of the \sphinxstyleemphasis{lj/cut} pair styles support the
{\hyperref[\detokenize{pair_modify::doc}]{\sphinxcrossref{\DUrole{doc}{pair\_modify}}}} tail option for adding a long\sphinxhyphen{}range
tail correction to the energy and pressure for the Lennard\sphinxhyphen{}Jones
portion of the pair interaction.

\sphinxAtStartPar
All of the \sphinxstyleemphasis{lj/cut} pair styles write their information to {\hyperref[\detokenize{restart::doc}]{\sphinxcrossref{\DUrole{doc}{binary restart files}}}}, so pair\_style and pair\_coeff commands do
not need to be specified in an input script that reads a restart file.

\sphinxAtStartPar
The \sphinxstyleemphasis{lj/cut} and \sphinxstyleemphasis{lj/cut/coul/long} pair styles support the use of the
\sphinxstyleemphasis{inner}, \sphinxstyleemphasis{middle}, and \sphinxstyleemphasis{outer} keywords of the {\hyperref[\detokenize{run_style::doc}]{\sphinxcrossref{\DUrole{doc}{run\_style respa}}}} command, meaning the pairwise forces can be
partitioned by distance at different levels of the rRESPA hierarchy.
The other styles only support the \sphinxstyleemphasis{pair} keyword of run\_style respa.
See the {\hyperref[\detokenize{run_style::doc}]{\sphinxcrossref{\DUrole{doc}{run\_style}}}} command for details.


\bigskip\hrule\bigskip



\subsection{Restrictions}
\label{\detokenize{pair_lj_cut_tip4p:restrictions}}
\sphinxAtStartPar
The \sphinxstyleemphasis{lj/cut/tip4p/long} styles are part of the
KSPACE package. The \sphinxstyleemphasis{lj/cut/tip4p/cut} style is part of the MOLECULE
package. These styles are only enabled if LAMMPS was built with those
packages.  See the {\hyperref[\detokenize{Build_package::doc}]{\sphinxcrossref{\DUrole{doc}{Build package}}}} page for
more info.


\subsection{Related commands}
\label{\detokenize{pair_lj_cut_tip4p:related-commands}}
\sphinxAtStartPar
{\hyperref[\detokenize{pair_coeff::doc}]{\sphinxcrossref{\DUrole{doc}{pair\_coeff}}}}


\subsection{Default}
\label{\detokenize{pair_lj_cut_tip4p:default}}
\sphinxAtStartPar
none


\bigskip\hrule\bigskip

\phantomsection\label{\detokenize{pair_lj_cut_tip4p:jorgensen2}}
\sphinxAtStartPar
\sphinxstylestrong{(Jorgensen)} Jorgensen, Chandrasekhar, Madura, Impey, Klein, J Chem
Phys, 79, 926 (1983).

\sphinxstepscope

\index{pair\_style lj/expand@\spxentry{pair\_style lj/expand}}\ignorespaces 
\index{pair\_style lj/expand/gpu@\spxentry{pair\_style lj/expand/gpu}}\ignorespaces 
\index{pair\_style lj/expand/kk@\spxentry{pair\_style lj/expand/kk}}\ignorespaces 
\index{pair\_style lj/expand/omp@\spxentry{pair\_style lj/expand/omp}}\ignorespaces 
\index{pair\_style lj/expand/coul/long@\spxentry{pair\_style lj/expand/coul/long}}\ignorespaces 
\index{pair\_style lj/expand/coul/long/kk@\spxentry{pair\_style lj/expand/coul/long/kk}}\ignorespaces 
\index{pair\_style lj/expand/coul/long/gpu@\spxentry{pair\_style lj/expand/coul/long/gpu}}\ignorespaces 

\section{pair\_style lj/expand command}
\label{\detokenize{pair_lj_expand:pair-style-lj-expand-command}}\label{\detokenize{pair_lj_expand:index-6}}\label{\detokenize{pair_lj_expand:index-5}}\label{\detokenize{pair_lj_expand:index-4}}\label{\detokenize{pair_lj_expand:index-3}}\label{\detokenize{pair_lj_expand:index-2}}\label{\detokenize{pair_lj_expand:index-1}}\label{\detokenize{pair_lj_expand:index-0}}\label{\detokenize{pair_lj_expand::doc}}
\sphinxAtStartPar
Accelerator Variants: \sphinxstyleemphasis{lj/expand/gpu}, \sphinxstyleemphasis{lj/expand/kk}, \sphinxstyleemphasis{lj/expand/omp}


\section{pair\_style lj/expand/coul/long command}
\label{\detokenize{pair_lj_expand:pair-style-lj-expand-coul-long-command}}
\sphinxAtStartPar
Accelerator Variants: \sphinxstyleemphasis{lj/expand/coul/long/gpu}, \sphinxstyleemphasis{lj/expand/coul/long/kk}


\subsection{Syntax}
\label{\detokenize{pair_lj_expand:syntax}}
\begin{sphinxVerbatim}[commandchars=\\\{\}]
\PYG{k}{pair\PYGZus{}style}\PYG{+w}{ }\PYG{n}{lj}\PYG{o}{/}\PYG{n}{expand}\PYG{+w}{ }\PYG{n}{cutoff}
\end{sphinxVerbatim}
\begin{itemize}
\item {} 
\sphinxAtStartPar
cutoff = global cutoff for lj/expand interactions (distance units)

\end{itemize}


\subsection{Examples}
\label{\detokenize{pair_lj_expand:examples}}
\begin{sphinxVerbatim}[commandchars=\\\{\}]
\PYG{k}{pair\PYGZus{}style}\PYG{+w}{ }\PYG{n}{lj}\PYG{o}{/}\PYG{n}{expand}\PYG{+w}{ }\PYG{l+m}{2.5}
\PYG{k}{pair\PYGZus{}coeff}\PYG{+w}{ }\PYG{o}{*}\PYG{+w}{ }\PYG{o}{*}\PYG{+w}{ }\PYG{l+m}{1.0}\PYG{+w}{ }\PYG{l+m}{1.0}\PYG{+w}{ }\PYG{l+m}{0.5}
\PYG{k}{pair\PYGZus{}coeff}\PYG{+w}{ }\PYG{l+m}{1}\PYG{+w}{ }\PYG{l+m}{1}\PYG{+w}{ }\PYG{l+m}{1.0}\PYG{+w}{ }\PYG{l+m}{1.0}\PYG{+w}{ }\PYG{o}{\PYGZhy{}}\PYG{l+m}{0.2}\PYG{+w}{ }\PYG{l+m}{2.0}

\PYG{k}{pair\PYGZus{}style}\PYG{+w}{ }\PYG{n}{lj}\PYG{o}{/}\PYG{n}{expand}\PYG{o}{/}\PYG{n}{coul}\PYG{o}{/}\PYG{n}{long}\PYG{+w}{ }\PYG{l+m}{2.5}
\PYG{k}{pair\PYGZus{}style}\PYG{+w}{ }\PYG{n}{lj}\PYG{o}{/}\PYG{n}{expand}\PYG{o}{/}\PYG{n}{coul}\PYG{o}{/}\PYG{n}{long}\PYG{+w}{ }\PYG{l+m}{2.5}\PYG{+w}{ }\PYG{l+m}{4.0}
\PYG{k}{pair\PYGZus{}coeff}\PYG{+w}{ }\PYG{o}{*}\PYG{+w}{ }\PYG{o}{*}\PYG{+w}{ }\PYG{l+m}{1.0}\PYG{+w}{ }\PYG{l+m}{1.0}\PYG{+w}{ }\PYG{l+m}{0.5}
\PYG{k}{pair\PYGZus{}coeff}\PYG{+w}{ }\PYG{l+m}{1}\PYG{+w}{ }\PYG{l+m}{1}\PYG{+w}{ }\PYG{l+m}{1.0}\PYG{+w}{ }\PYG{l+m}{1.0}\PYG{+w}{ }\PYG{o}{\PYGZhy{}}\PYG{l+m}{0.2}\PYG{+w}{ }\PYG{l+m}{3.0}
\end{sphinxVerbatim}


\subsection{Description}
\label{\detokenize{pair_lj_expand:description}}
\sphinxAtStartPar
Style \sphinxstyleemphasis{lj/expand} computes a LJ interaction with a distance shifted by
delta which can be useful when particles are of different sizes, since
it is different that using different sigma values in a standard LJ
formula:
\begin{equation*}
\begin{split}E = 4 \epsilon \left[ \left(\frac{\sigma}{r - \Delta}\right)^{12} -
  \left(\frac{\sigma}{r - \Delta}\right)^6 \right]
  \qquad r < r_c + \Delta\end{split}
\end{equation*}
\sphinxAtStartPar
\(r_c\) is the cutoff which does not include the \(\Delta\)
distance.  I.e. the actual force cutoff is the sum of \(r_c +
\Delta\).

\sphinxAtStartPar
For all of the \sphinxstyleemphasis{lj/expand} pair styles, the following coefficients must
be defined for each pair of atoms types via the {\hyperref[\detokenize{pair_coeff::doc}]{\sphinxcrossref{\DUrole{doc}{pair\_coeff}}}} command as in the examples above, or in the data file or
restart files read by the {\hyperref[\detokenize{read_data::doc}]{\sphinxcrossref{\DUrole{doc}{read\_data}}}} or
{\hyperref[\detokenize{read_restart::doc}]{\sphinxcrossref{\DUrole{doc}{read\_restart}}}} commands, or by mixing as described
below:
\begin{itemize}
\item {} 
\sphinxAtStartPar
\(\epsilon\) (energy units)

\item {} 
\sphinxAtStartPar
\(\sigma\) (distance units)

\item {} 
\sphinxAtStartPar
\(\Delta\) (distance units)

\item {} 
\sphinxAtStartPar
cutoff (distance units)

\end{itemize}

\sphinxAtStartPar
The \(\Delta\) values can be positive or negative.  The last
coefficient is optional.  If not specified, the global LJ cutoff is
used.

\sphinxAtStartPar
For \sphinxstyleemphasis{lj/expand/coul/long} only the LJ cutoff can be specified since a
Coulombic cutoff cannot be specified for an individual I,J type pair.
All type pairs use the same global Coulombic cutoff specified in the
pair\_style command.


\bigskip\hrule\bigskip


\sphinxAtStartPar
Styles with a \sphinxstyleemphasis{gpu}, \sphinxstyleemphasis{intel}, \sphinxstyleemphasis{kk}, \sphinxstyleemphasis{omp}, or \sphinxstyleemphasis{opt} suffix are
functionally the same as the corresponding style without the suffix.
They have been optimized to run faster, depending on your available
hardware, as discussed on the {\hyperref[\detokenize{Speed_packages::doc}]{\sphinxcrossref{\DUrole{doc}{Accelerator packages}}}}
page.  The accelerated styles take the same arguments and should
produce the same results, except for round\sphinxhyphen{}off and precision issues.

\sphinxAtStartPar
These accelerated styles are part of the GPU, INTEL, KOKKOS,
OPENMP, and OPT packages, respectively.  They are only enabled if
LAMMPS was built with those packages.  See the {\hyperref[\detokenize{Build_package::doc}]{\sphinxcrossref{\DUrole{doc}{Build package}}}} page for more info.

\sphinxAtStartPar
You can specify the accelerated styles explicitly in your input script
by including their suffix, or you can use the {\hyperref[\detokenize{Run_options::doc}]{\sphinxcrossref{\DUrole{doc}{\sphinxhyphen{}suffix command\sphinxhyphen{}line switch}}}} when you invoke LAMMPS, or you can use the
{\hyperref[\detokenize{suffix::doc}]{\sphinxcrossref{\DUrole{doc}{suffix}}}} command in your input script.

\sphinxAtStartPar
See the {\hyperref[\detokenize{Speed_packages::doc}]{\sphinxcrossref{\DUrole{doc}{Accelerator packages}}}} page for more
instructions on how to use the accelerated styles effectively.


\bigskip\hrule\bigskip



\subsection{Mixing, shift, table, tail correction, restart, rRESPA info}
\label{\detokenize{pair_lj_expand:mixing-shift-table-tail-correction-restart-rrespa-info}}
\sphinxAtStartPar
For atom type pairs I,J and I != J, the epsilon, sigma, and shift
coefficients and cutoff distance for this pair style can be mixed.
Shift is always mixed via an \sphinxstyleemphasis{arithmetic} rule.  The other
coefficients are mixed according to the pair\_modify mix value.  The
default mix value is \sphinxstyleemphasis{geometric}.  See the “pair\_modify” command for
details.

\sphinxAtStartPar
This pair style supports the {\hyperref[\detokenize{pair_modify::doc}]{\sphinxcrossref{\DUrole{doc}{pair\_modify}}}} shift
option for the energy of the pair interaction.

\sphinxAtStartPar
The {\hyperref[\detokenize{pair_modify::doc}]{\sphinxcrossref{\DUrole{doc}{pair\_modify}}}} table option is not relevant
for this pair style.

\sphinxAtStartPar
This pair style supports the {\hyperref[\detokenize{pair_modify::doc}]{\sphinxcrossref{\DUrole{doc}{pair\_modify}}}} tail
option for adding a long\sphinxhyphen{}range tail correction to the energy and
pressure of the pair interaction.

\sphinxAtStartPar
This pair style writes its information to {\hyperref[\detokenize{restart::doc}]{\sphinxcrossref{\DUrole{doc}{binary restart files}}}}, so pair\_style and pair\_coeff commands do not need
to be specified in an input script that reads a restart file.

\sphinxAtStartPar
This pair style can only be used via the \sphinxstyleemphasis{pair} keyword of the
{\hyperref[\detokenize{run_style::doc}]{\sphinxcrossref{\DUrole{doc}{run\_style respa}}}} command.  It does not support the
\sphinxstyleemphasis{inner}, \sphinxstyleemphasis{middle}, \sphinxstyleemphasis{outer} keywords.


\bigskip\hrule\bigskip



\subsection{Restrictions}
\label{\detokenize{pair_lj_expand:restrictions}}
\sphinxAtStartPar
none


\subsection{Related commands}
\label{\detokenize{pair_lj_expand:related-commands}}
\sphinxAtStartPar
{\hyperref[\detokenize{pair_coeff::doc}]{\sphinxcrossref{\DUrole{doc}{pair\_coeff}}}}


\subsection{Default}
\label{\detokenize{pair_lj_expand:default}}
\sphinxAtStartPar
none

\sphinxstepscope

\index{pair\_style lj/expand/sphere@\spxentry{pair\_style lj/expand/sphere}}\ignorespaces 
\index{pair\_style lj/expand/sphere/omp@\spxentry{pair\_style lj/expand/sphere/omp}}\ignorespaces 

\section{pair\_style lj/expand/sphere command}
\label{\detokenize{pair_lj_expand_sphere:pair-style-lj-expand-sphere-command}}\label{\detokenize{pair_lj_expand_sphere:index-1}}\label{\detokenize{pair_lj_expand_sphere:index-0}}\label{\detokenize{pair_lj_expand_sphere::doc}}
\sphinxAtStartPar
Accelerator Variant: \sphinxstyleemphasis{lj/expand/sphere/omp}


\subsection{Syntax}
\label{\detokenize{pair_lj_expand_sphere:syntax}}
\begin{sphinxVerbatim}[commandchars=\\\{\}]
\PYG{k}{pair\PYGZus{}style}\PYG{+w}{ }\PYG{n}{style}\PYG{+w}{ }\PYG{n}{args}
\end{sphinxVerbatim}
\begin{itemize}
\item {} 
\sphinxAtStartPar
style = \sphinxstyleemphasis{lj/expand/sphere}

\item {} 
\sphinxAtStartPar
args = list of arguments for a particular style

\end{itemize}
\begin{sphinxalltt}
\sphinxstyleemphasis{lj/expand/sphere} args = cutoff
  cutoff = global cutoff for Lennard Jones interactions (distance units)
\end{sphinxalltt}


\subsection{Examples}
\label{\detokenize{pair_lj_expand_sphere:examples}}
\begin{sphinxVerbatim}[commandchars=\\\{\}]
\PYG{k}{pair\PYGZus{}style}\PYG{+w}{ }\PYG{n}{lj}\PYG{o}{/}\PYG{n}{expand}\PYG{o}{/}\PYG{n}{sphere}\PYG{+w}{ }\PYG{l+m}{2.5}
\PYG{k}{pair\PYGZus{}coeff}\PYG{+w}{ }\PYG{o}{*}\PYG{+w}{ }\PYG{o}{*}\PYG{+w}{ }\PYG{l+m}{1.0}\PYG{+w}{ }\PYG{l+m}{1.0}
\PYG{k}{pair\PYGZus{}coeff}\PYG{+w}{ }\PYG{l+m}{1}\PYG{+w}{ }\PYG{l+m}{1}\PYG{+w}{ }\PYG{l+m}{1.1}\PYG{+w}{ }\PYG{l+m}{0.4}\PYG{+w}{ }\PYG{l+m}{2.8}
\end{sphinxVerbatim}


\subsection{Description}
\label{\detokenize{pair_lj_expand_sphere:description}}
\sphinxAtStartPar
\DUrole{versionmodified,added}{Added in version 15Jun2023.}

\sphinxAtStartPar
The \sphinxstyleemphasis{lj/expand/sphere} style compute a 12/6 Lennard\sphinxhyphen{}Jones potential with
a distance shifted by \(\Delta = \frac{1}{2} (d_i + d_j)\), the
average diameter of both atoms.  This can be used to model particles of
different sizes but same interactions, which is different from using
different sigma values as in {\hyperref[\detokenize{pair_lj_cut_sphere::doc}]{\sphinxcrossref{\DUrole{doc}{pair style lj/cut/sphere}}}}.
\begin{equation*}
\begin{split}E = 4 \epsilon \left[ \left(\frac{\sigma}{r - \Delta}\right)^{12} -
  \left(\frac{\sigma}{r - \Delta}\right)^6 \right]
  \qquad r < r_c + \Delta\end{split}
\end{equation*}
\sphinxAtStartPar
\(r_c\) is the cutoff which does not include the distance \(\Delta\).
I.e. the actual force cutoff is the sum \(r_c + \Delta\).

\sphinxAtStartPar
This is the same potential function used by the {\hyperref[\detokenize{pair_lj::doc}]{\sphinxcrossref{\DUrole{doc}{lj/expand}}}} pair style, but the \(\Delta\) parameter is not
set as a per\sphinxhyphen{}type parameter via the {\hyperref[\detokenize{pair_coeff::doc}]{\sphinxcrossref{\DUrole{doc}{pair\_coeff command}}}}.  Instead it is calculated individually for each pair
using the per\sphinxhyphen{}atom diameter attribute of {\hyperref[\detokenize{atom_style::doc}]{\sphinxcrossref{\DUrole{doc}{atom\_style sphere}}}} for the two atoms as the average diameter, \(\Delta = \frac{1}{2} (d_i + d_j)\)

\sphinxAtStartPar
Note that \(\sigma\) is defined in the LJ formula above as the
zero\sphinxhyphen{}crossing distance for the potential, \sphinxstyleemphasis{not} as the energy minimum which
is at \(2^{\frac{1}{6}} \sigma\).

\begin{sphinxadmonition}{note}{Notes on cutoffs, neighbor lists, and efficiency}

\sphinxAtStartPar
If your system is mildly polydisperse, meaning the ratio of the
diameter of the largest particle to the smallest is less than 2.0,
then the neighbor lists built by the code should be reasonably
efficient.  Which means they will not contain too many particle
pairs that do not interact.  However, if your system is highly
polydisperse (ratio \textgreater{} 2.0), the neighbor list build and force
computations may be inefficient.  There are two ways to try and
speed up the simulations.

\sphinxAtStartPar
The first is to assign atoms to different atom types so that atoms of
each type are similar in size.  E.g. if particle diameters range from
1 to 5 use 4 atom types, ensuring atoms of type 1 have diameters from
1.0\sphinxhyphen{}2.0, type 2 from 2.0\sphinxhyphen{}3.0, etc.  This will reduce the number of
non\sphinxhyphen{}interacting pairs in the neighbor lists and thus reduce the time
spent on computing pairwise interactions.

\sphinxAtStartPar
The second is to use the {\hyperref[\detokenize{neighbor::doc}]{\sphinxcrossref{\DUrole{doc}{neighbor multi}}}} command
which enabled a different algorithm for building neighbor lists.  This
will also require that you assign multiple atom types according to
diameters, but will in addition use a more efficient size\sphinxhyphen{}dependent
strategy to construct the neighbor lists and thus reduce the time
spent on building neighbor lists.

\sphinxAtStartPar
Here are example input script commands using the first option for a
highly polydisperse system:

\begin{sphinxVerbatim}[commandchars=\\\{\}]
\PYG{n}{units}\PYG{+w}{           }\PYG{n}{lj}
\PYG{n}{atom\PYGZus{}style}\PYG{+w}{      }\PYG{n}{sphere}
\PYG{n}{lattice}\PYG{+w}{         }\PYG{n}{fcc}\PYG{+w}{ }\PYG{l+m+mf}{0.8442}
\PYG{n}{region}\PYG{+w}{          }\PYG{n}{box}\PYG{+w}{ }\PYG{n}{block}\PYG{+w}{ }\PYG{l+m+mi}{0}\PYG{+w}{ }\PYG{l+m+mi}{10}\PYG{+w}{ }\PYG{l+m+mi}{0}\PYG{+w}{ }\PYG{l+m+mi}{10}\PYG{+w}{ }\PYG{l+m+mi}{0}\PYG{+w}{ }\PYG{l+m+mi}{10}
\PYG{n}{create\PYGZus{}box}\PYG{+w}{      }\PYG{l+m+mi}{2}\PYG{+w}{ }\PYG{n}{box}
\PYG{n}{create\PYGZus{}atoms}\PYG{+w}{    }\PYG{l+m+mi}{1}\PYG{+w}{ }\PYG{n}{box}

\PYG{c+cp}{\PYGZsh{}}\PYG{c+cp}{ create atoms with random diameters from bimodal distribution}
\PYG{n}{variable}\PYG{+w}{ }\PYG{k}{switch}\PYG{+w}{ }\PYG{n}{atom}\PYG{+w}{ }\PYG{n}{random}\PYG{p}{(}\PYG{l+m+mf}{0.0}\PYG{p}{,}\PYG{l+m+mf}{1.0}\PYG{p}{,}\PYG{l+m+mi}{345634}\PYG{p}{)}
\PYG{n}{variable}\PYG{+w}{ }\PYG{n}{diam}\PYG{+w}{ }\PYG{n}{atom}\PYG{+w}{ }\PYG{p}{(}\PYG{n}{v\PYGZus{}switch}\PYG{o}{\PYGZlt{}}\PYG{l+m+mf}{0.75}\PYG{p}{)}\PYG{o}{*}\PYG{n}{normal}\PYG{p}{(}\PYG{l+m+mf}{0.2}\PYG{p}{,}\PYG{l+m+mf}{0.04}\PYG{p}{,}\PYG{l+m+mi}{325}\PYG{p}{)}\PYG{o}{+}\PYG{p}{(}\PYG{n}{v\PYGZus{}switch}\PYG{o}{\PYGZgt{}}\PYG{o}{=}\PYG{l+m+mf}{0.7}\PYG{p}{)}\PYG{o}{*}\PYG{n}{normal}\PYG{p}{(}\PYG{l+m+mf}{0.6}\PYG{p}{,}\PYG{l+m+mf}{0.2}\PYG{p}{,}\PYG{l+m+mi}{453}\PYG{p}{)}
\PYG{n}{set}\PYG{+w}{ }\PYG{n}{group}\PYG{+w}{ }\PYG{n}{all}\PYG{+w}{ }\PYG{n}{diameter}\PYG{+w}{ }\PYG{n}{v\PYGZus{}diam}

\PYG{c+cp}{\PYGZsh{}}\PYG{c+cp}{ assign type 2 to atoms with diameter \PYGZgt{} 0.35}
\PYG{n}{variable}\PYG{+w}{ }\PYG{n}{large}\PYG{+w}{ }\PYG{n}{atom}\PYG{+w}{ }\PYG{p}{(}\PYG{l+m+mf}{2.0}\PYG{o}{*}\PYG{n}{radius}\PYG{p}{)}\PYG{o}{\PYGZgt{}}\PYG{l+m+mf}{0.35}
\PYG{n}{group}\PYG{+w}{ }\PYG{n}{large}\PYG{+w}{ }\PYG{n}{variable}\PYG{+w}{ }\PYG{n}{large}
\PYG{n}{set}\PYG{+w}{ }\PYG{n}{group}\PYG{+w}{ }\PYG{n}{large}\PYG{+w}{ }\PYG{n}{type}\PYG{+w}{ }\PYG{l+m+mi}{2}

\PYG{n}{pair\PYGZus{}style}\PYG{+w}{      }\PYG{n}{lj}\PYG{o}{/}\PYG{n}{expand}\PYG{o}{/}\PYG{n}{sphere}\PYG{+w}{ }\PYG{l+m+mf}{2.0}
\PYG{n}{pair\PYGZus{}coeff}\PYG{+w}{      }\PYG{o}{*}\PYG{+w}{ }\PYG{o}{*}\PYG{+w}{ }\PYG{l+m+mf}{1.0}\PYG{+w}{ }\PYG{l+m+mf}{0.5}

\PYG{n}{neighbor}\PYG{+w}{ }\PYG{l+m+mf}{0.3}\PYG{+w}{ }\PYG{n}{bin}
\end{sphinxVerbatim}

\sphinxAtStartPar
Using multiple atom types speeds up the calculation for this example
by more than 30 percent, but using the multi\sphinxhyphen{}style neighbor list does
not provide a speedup.
\end{sphinxadmonition}


\subsection{Coefficients}
\label{\detokenize{pair_lj_expand_sphere:coefficients}}
\sphinxAtStartPar
The following coefficients must be defined for each pair of atoms types via the
{\hyperref[\detokenize{pair_coeff::doc}]{\sphinxcrossref{\DUrole{doc}{pair\_coeff}}}} command as in the examples above, or in the data
file or restart files read by the {\hyperref[\detokenize{read_data::doc}]{\sphinxcrossref{\DUrole{doc}{read\_data}}}} or
{\hyperref[\detokenize{read_restart::doc}]{\sphinxcrossref{\DUrole{doc}{read\_restart}}}} commands, or by mixing as described below:
\begin{itemize}
\item {} 
\sphinxAtStartPar
\(\epsilon\) (energy units)

\item {} 
\sphinxAtStartPar
\(\sigma\) (distance units)

\item {} 
\sphinxAtStartPar
LJ cutoff (distance units) (optional)

\end{itemize}

\sphinxAtStartPar
The last coefficient is optional.  If not specified, the global LJ
cutoff specified in the {\hyperref[\detokenize{pair_style::doc}]{\sphinxcrossref{\DUrole{doc}{pair\_style command}}}} is used.


\bigskip\hrule\bigskip


\sphinxAtStartPar
Styles with a \sphinxstyleemphasis{gpu}, \sphinxstyleemphasis{intel}, \sphinxstyleemphasis{kk}, \sphinxstyleemphasis{omp}, or \sphinxstyleemphasis{opt} suffix are
functionally the same as the corresponding style without the suffix.
They have been optimized to run faster, depending on your available
hardware, as discussed on the {\hyperref[\detokenize{Speed_packages::doc}]{\sphinxcrossref{\DUrole{doc}{Accelerator packages}}}}
page.  The accelerated styles take the same arguments and should
produce the same results, except for round\sphinxhyphen{}off and precision issues.

\sphinxAtStartPar
These accelerated styles are part of the GPU, INTEL, KOKKOS,
OPENMP, and OPT packages, respectively.  They are only enabled if
LAMMPS was built with those packages.  See the {\hyperref[\detokenize{Build_package::doc}]{\sphinxcrossref{\DUrole{doc}{Build package}}}} page for more info.

\sphinxAtStartPar
You can specify the accelerated styles explicitly in your input script
by including their suffix, or you can use the {\hyperref[\detokenize{Run_options::doc}]{\sphinxcrossref{\DUrole{doc}{\sphinxhyphen{}suffix command\sphinxhyphen{}line switch}}}} when you invoke LAMMPS, or you can use the
{\hyperref[\detokenize{suffix::doc}]{\sphinxcrossref{\DUrole{doc}{suffix}}}} command in your input script.

\sphinxAtStartPar
See the {\hyperref[\detokenize{Speed_packages::doc}]{\sphinxcrossref{\DUrole{doc}{Accelerator packages}}}} page for more
instructions on how to use the accelerated styles effectively.


\bigskip\hrule\bigskip



\subsection{Mixing, shift, table, tail correction, restart, rRESPA info}
\label{\detokenize{pair_lj_expand_sphere:mixing-shift-table-tail-correction-restart-rrespa-info}}
\sphinxAtStartPar
For atom type pairs I,J and I != J, the epsilon, sigma, and cutoff
coefficients for the \sphinxstyleemphasis{lj/expand/sphere} pair style can be mixed.  The
default mixing style is \sphinxstyleemphasis{geometric}.  See the {\hyperref[\detokenize{pair_modify::doc}]{\sphinxcrossref{\DUrole{doc}{pair\_modify command}}}} for details.

\sphinxAtStartPar
The \sphinxstyleemphasis{lj/expand/sphere} pair style supports the {\hyperref[\detokenize{pair_modify::doc}]{\sphinxcrossref{\DUrole{doc}{pair\_modify shift}}}}
option for the energy of the Lennard\sphinxhyphen{}Jones portion of the pair interaction.

\sphinxAtStartPar
The \sphinxstyleemphasis{lj/expand/sphere} pair style does \sphinxstyleemphasis{not} support the {\hyperref[\detokenize{pair_modify::doc}]{\sphinxcrossref{\DUrole{doc}{pair\_modify}}}} tail option for adding a long\sphinxhyphen{}range tail corrections to
the energy and pressure.

\sphinxAtStartPar
The \sphinxstyleemphasis{lj/expand/sphere} pair style writes its information to {\hyperref[\detokenize{restart::doc}]{\sphinxcrossref{\DUrole{doc}{binary
restart files}}}}, so pair\_style and pair\_coeff commands do not
need to be specified in an input script that reads a restart file.

\sphinxAtStartPar
This pair style can only be used via the \sphinxstyleemphasis{pair} keyword of the
{\hyperref[\detokenize{run_style::doc}]{\sphinxcrossref{\DUrole{doc}{run\_style respa}}}} command.  It does \sphinxstyleemphasis{not} support the
\sphinxstyleemphasis{inner}, \sphinxstyleemphasis{middle}, \sphinxstyleemphasis{outer} keywords.


\bigskip\hrule\bigskip



\subsection{Restrictions}
\label{\detokenize{pair_lj_expand_sphere:restrictions}}
\sphinxAtStartPar
The \sphinxstyleemphasis{lj/expand/sphere} pair style is only enabled if LAMMPS was built with the
EXTRA\sphinxhyphen{}PAIR package.  See the {\hyperref[\detokenize{Build_package::doc}]{\sphinxcrossref{\DUrole{doc}{Build package}}}} page
for more info.


\bigskip\hrule\bigskip



\subsection{Related commands}
\label{\detokenize{pair_lj_expand_sphere:related-commands}}\begin{itemize}
\item {} 
\sphinxAtStartPar
{\hyperref[\detokenize{pair_coeff::doc}]{\sphinxcrossref{\DUrole{doc}{pair\_coeff}}}}

\item {} 
\sphinxAtStartPar
{\hyperref[\detokenize{pair_lj::doc}]{\sphinxcrossref{\DUrole{doc}{pair\_style lj/cut}}}}

\item {} 
\sphinxAtStartPar
{\hyperref[\detokenize{pair_lj_cut_sphere::doc}]{\sphinxcrossref{\DUrole{doc}{pair\_style lj/cut/sphere}}}}

\end{itemize}


\subsection{Default}
\label{\detokenize{pair_lj_expand_sphere:default}}
\sphinxAtStartPar
none

\sphinxstepscope

\index{pair\_style lj/long/coul/long@\spxentry{pair\_style lj/long/coul/long}}\ignorespaces 
\index{pair\_style lj/long/coul/long/intel@\spxentry{pair\_style lj/long/coul/long/intel}}\ignorespaces 
\index{pair\_style lj/long/coul/long/omp@\spxentry{pair\_style lj/long/coul/long/omp}}\ignorespaces 
\index{pair\_style lj/long/coul/long/opt@\spxentry{pair\_style lj/long/coul/long/opt}}\ignorespaces 
\index{pair\_style lj/long/tip4p/long@\spxentry{pair\_style lj/long/tip4p/long}}\ignorespaces 
\index{pair\_style lj/long/tip4p/long/omp@\spxentry{pair\_style lj/long/tip4p/long/omp}}\ignorespaces 

\section{pair\_style lj/long/coul/long command}
\label{\detokenize{pair_lj_long:pair-style-lj-long-coul-long-command}}\label{\detokenize{pair_lj_long:index-5}}\label{\detokenize{pair_lj_long:index-4}}\label{\detokenize{pair_lj_long:index-3}}\label{\detokenize{pair_lj_long:index-2}}\label{\detokenize{pair_lj_long:index-1}}\label{\detokenize{pair_lj_long:index-0}}\label{\detokenize{pair_lj_long::doc}}
\sphinxAtStartPar
Accelerator Variants: \sphinxstyleemphasis{lj/long/coul/long/intel}, \sphinxstyleemphasis{lj/long/coul/long/omp}, \sphinxstyleemphasis{lj/long/coul/long/opt}


\section{pair\_style lj/long/tip4p/long command}
\label{\detokenize{pair_lj_long:pair-style-lj-long-tip4p-long-command}}
\sphinxAtStartPar
Accelerator Variants: \sphinxstyleemphasis{lj/long/tip4p/long/omp}


\subsection{Syntax}
\label{\detokenize{pair_lj_long:syntax}}
\begin{sphinxVerbatim}[commandchars=\\\{\}]
\PYG{k}{pair\PYGZus{}style}\PYG{+w}{ }\PYG{n}{style}\PYG{+w}{ }\PYG{n}{args}
\end{sphinxVerbatim}
\begin{itemize}
\item {} 
\sphinxAtStartPar
style = \sphinxstyleemphasis{lj/long/coul/long} or \sphinxstyleemphasis{lj/long/tip4p/long}

\item {} 
\sphinxAtStartPar
args = list of arguments for a particular style

\end{itemize}
\begin{sphinxalltt}
\sphinxstyleemphasis{lj/long/coul/long} args = flag\_lj flag\_coul cutoff (cutoff2)
  flag\_lj = \sphinxstyleemphasis{long} or \sphinxstyleemphasis{cut} or \sphinxstyleemphasis{off}
    \sphinxstyleemphasis{long} = use Kspace long\sphinxhyphen{}range summation for dispersion 1/r\textasciicircum{}6 term
    \sphinxstyleemphasis{cut} = use a cutoff on dispersion 1/r\textasciicircum{}6 term
    \sphinxstyleemphasis{off} = omit disperion 1/r\textasciicircum{}6 term entirely
  flag\_coul = \sphinxstyleemphasis{long} or \sphinxstyleemphasis{off}
    \sphinxstyleemphasis{long} = use Kspace long\sphinxhyphen{}range summation for Coulombic 1/r term
    \sphinxstyleemphasis{off} = omit Coulombic term
  cutoff = global cutoff for LJ (and Coulombic if only 1 arg) (distance units)
  cutoff2 = global cutoff for Coulombic (optional) (distance units)
\sphinxstyleemphasis{lj/long/tip4p/long} args = flag\_lj flag\_coul otype htype btype atype qdist cutoff (cutoff2)
  flag\_lj = \sphinxstyleemphasis{long} or \sphinxstyleemphasis{cut}
    \sphinxstyleemphasis{long} = use Kspace long\sphinxhyphen{}range summation for dispersion 1/r\textasciicircum{}6 term
    \sphinxstyleemphasis{cut} = use a cutoff
  flag\_coul = \sphinxstyleemphasis{long} or \sphinxstyleemphasis{off}
    \sphinxstyleemphasis{long} = use Kspace long\sphinxhyphen{}range summation for Coulombic 1/r term
    \sphinxstyleemphasis{off} = omit Coulombic term
  otype,htype = atom types (numeric or type label) for TIP4P O and H
  btype,atype = bond and angle types (numeric or type label) for TIP4P waters
  qdist = distance from O atom to massless charge (distance units)
  cutoff = global cutoff for LJ (and Coulombic if only 1 arg) (distance units)
  cutoff2 = global cutoff for Coulombic (optional) (distance units)
\end{sphinxalltt}


\subsection{Examples}
\label{\detokenize{pair_lj_long:examples}}
\begin{sphinxVerbatim}[commandchars=\\\{\}]
\PYG{k}{pair\PYGZus{}style}\PYG{+w}{ }\PYG{n}{lj}\PYG{o}{/}\PYG{n}{long}\PYG{o}{/}\PYG{n}{coul}\PYG{o}{/}\PYG{n}{long}\PYG{+w}{ }\PYG{n}{cut}\PYG{+w}{ }\PYG{n}{off}\PYG{+w}{ }\PYG{l+m}{2.5}
\PYG{k}{pair\PYGZus{}style}\PYG{+w}{ }\PYG{n}{lj}\PYG{o}{/}\PYG{n}{long}\PYG{o}{/}\PYG{n}{coul}\PYG{o}{/}\PYG{n}{long}\PYG{+w}{ }\PYG{n}{cut}\PYG{+w}{ }\PYG{n}{long}\PYG{+w}{ }\PYG{l+m}{2.5}\PYG{+w}{ }\PYG{l+m}{4.0}
\PYG{k}{pair\PYGZus{}style}\PYG{+w}{ }\PYG{n}{lj}\PYG{o}{/}\PYG{n}{long}\PYG{o}{/}\PYG{n}{coul}\PYG{o}{/}\PYG{n}{long}\PYG{+w}{ }\PYG{n}{long}\PYG{+w}{ }\PYG{n}{long}\PYG{+w}{ }\PYG{l+m}{2.5}\PYG{+w}{ }\PYG{l+m}{4.0}
\PYG{k}{pair\PYGZus{}coeff}\PYG{+w}{ }\PYG{o}{*}\PYG{+w}{ }\PYG{o}{*}\PYG{+w}{ }\PYG{l+m}{1}\PYG{+w}{ }\PYG{l+m}{1}
\PYG{k}{pair\PYGZus{}coeff}\PYG{+w}{ }\PYG{l+m}{1}\PYG{+w}{ }\PYG{l+m}{1}\PYG{+w}{ }\PYG{l+m}{1}\PYG{+w}{ }\PYG{l+m}{3}\PYG{+w}{ }\PYG{l+m}{4}

\PYG{k}{pair\PYGZus{}style}\PYG{+w}{ }\PYG{n}{lj}\PYG{o}{/}\PYG{n}{long}\PYG{o}{/}\PYG{n}{tip4p}\PYG{o}{/}\PYG{n}{long}\PYG{+w}{ }\PYG{n}{long}\PYG{+w}{ }\PYG{n}{long}\PYG{+w}{ }\PYG{l+m}{1}\PYG{+w}{ }\PYG{l+m}{2}\PYG{+w}{ }\PYG{l+m}{7}\PYG{+w}{ }\PYG{l+m}{8}\PYG{+w}{ }\PYG{l+m}{0.15}\PYG{+w}{ }\PYG{l+m}{12.0}
\PYG{k}{pair\PYGZus{}style}\PYG{+w}{ }\PYG{n}{lj}\PYG{o}{/}\PYG{n}{long}\PYG{o}{/}\PYG{n}{tip4p}\PYG{o}{/}\PYG{n}{long}\PYG{+w}{ }\PYG{n}{long}\PYG{+w}{ }\PYG{n}{long}\PYG{+w}{ }\PYG{l+m}{1}\PYG{+w}{ }\PYG{l+m}{2}\PYG{+w}{ }\PYG{l+m}{7}\PYG{+w}{ }\PYG{l+m}{8}\PYG{+w}{ }\PYG{l+m}{0.15}\PYG{+w}{ }\PYG{l+m}{12.0}\PYG{+w}{ }\PYG{l+m}{10.0}
\PYG{k}{pair\PYGZus{}coeff}\PYG{+w}{ }\PYG{o}{*}\PYG{+w}{ }\PYG{o}{*}\PYG{+w}{ }\PYG{l+m}{100.0}\PYG{+w}{ }\PYG{l+m}{3.0}
\PYG{k}{pair\PYGZus{}coeff}\PYG{+w}{ }\PYG{l+m}{1}\PYG{+w}{ }\PYG{l+m}{1}\PYG{+w}{ }\PYG{l+m}{100.0}\PYG{+w}{ }\PYG{l+m}{3.5}\PYG{+w}{ }\PYG{l+m}{9.0}

\PYG{k}{pair\PYGZus{}style}\PYG{+w}{ }\PYG{n}{lj}\PYG{o}{/}\PYG{n}{long}\PYG{o}{/}\PYG{n}{tip4p}\PYG{o}{/}\PYG{n}{long}\PYG{+w}{ }\PYG{n}{long}\PYG{+w}{ }\PYG{n}{long}\PYG{+w}{ }\PYG{n}{OW}\PYG{+w}{ }\PYG{n}{HW}\PYG{+w}{ }\PYG{n}{HW}\PYG{o}{\PYGZhy{}}\PYG{n}{OW}\PYG{+w}{ }\PYG{n}{HW}\PYG{o}{\PYGZhy{}}\PYG{n}{OW}\PYG{o}{\PYGZhy{}}\PYG{n}{HW}\PYG{+w}{ }\PYG{l+m}{0.15}\PYG{+w}{ }\PYG{l+m}{12.0}
\PYG{k}{labelmap}\PYG{+w}{ }\PYG{n}{atom}\PYG{+w}{ }\PYG{l+m}{1}\PYG{+w}{ }\PYG{n}{OW}\PYG{+w}{ }\PYG{l+m}{2}\PYG{+w}{ }\PYG{n}{HW}
\PYG{k}{labelmap}\PYG{+w}{ }\PYG{n}{bond}\PYG{+w}{ }\PYG{l+m}{1}\PYG{+w}{ }\PYG{n}{HW}\PYG{o}{\PYGZhy{}}\PYG{n}{OW}
\PYG{k}{labelmap}\PYG{+w}{ }\PYG{n}{angle}\PYG{+w}{ }\PYG{l+m}{1}\PYG{+w}{ }\PYG{n}{HW}\PYG{o}{\PYGZhy{}}\PYG{n}{OW}\PYG{o}{\PYGZhy{}}\PYG{n}{HW}
\PYG{k}{pair\PYGZus{}coeff}\PYG{+w}{ }\PYG{o}{*}\PYG{+w}{ }\PYG{o}{*}\PYG{+w}{ }\PYG{l+m}{100.0}\PYG{+w}{ }\PYG{l+m}{3.0}
\PYG{k}{pair\PYGZus{}coeff}\PYG{+w}{ }\PYG{n}{OW}\PYG{+w}{ }\PYG{n}{OW}\PYG{+w}{ }\PYG{l+m}{100.0}\PYG{+w}{ }\PYG{l+m}{3.5}\PYG{+w}{ }\PYG{l+m}{9.0}
\end{sphinxVerbatim}


\subsection{Description}
\label{\detokenize{pair_lj_long:description}}
\sphinxAtStartPar
Style \sphinxstyleemphasis{lj/long/coul/long} computes the standard 12/6 Lennard\sphinxhyphen{}Jones potential:
\begin{equation*}
\begin{split}E = 4 \epsilon \left[ \left(\frac{\sigma}{r}\right)^{12} -
                    \left(\frac{\sigma}{r}\right)^6 \right]
                    \qquad r < r_c \\\end{split}
\end{equation*}
\sphinxAtStartPar
with \(\epsilon\) and \(\sigma\) being the usual Lennard\sphinxhyphen{}Jones
potential parameters, plus the Coulomb potential, given by:
\begin{equation*}
\begin{split}E = \frac{C q_i q_j}{\epsilon  r} \qquad r < r_c\end{split}
\end{equation*}
\sphinxAtStartPar
where C is an energy\sphinxhyphen{}conversion constant, \(q_i\) and \(q_j\) are the charges on
the two atoms, \(\epsilon\) is the dielectric constant which can be set by
the {\hyperref[\detokenize{dielectric::doc}]{\sphinxcrossref{\DUrole{doc}{dielectric}}}} command, and \(r_c\) is the cutoff.  If
one cutoff is specified in the pair\_style command, it is used for both
the LJ and Coulombic terms.  If two cutoffs are specified, they are
used as cutoffs for the LJ and Coulombic terms respectively.

\sphinxAtStartPar
The purpose of this pair style is to capture long\sphinxhyphen{}range interactions
resulting from both attractive 1/r\textasciicircum{}6 Lennard\sphinxhyphen{}Jones and Coulombic 1/r
interactions.  This is done by use of the \sphinxstyleemphasis{flag\_lj} and \sphinxstyleemphasis{flag\_coul}
settings.  The {\hyperref[\detokenize{pair_lj_long:veld2}]{\sphinxcrossref{\DUrole{std,std-ref}{In ‘t Veld}}}} paper has more details on when it is
appropriate to include long\sphinxhyphen{}range 1/r\textasciicircum{}6 interactions, using this
potential.

\sphinxAtStartPar
Style \sphinxstyleemphasis{lj/long/tip4p/long} implements the TIP4P water model of
{\hyperref[\detokenize{pair_lj_long:jorgensen4}]{\sphinxcrossref{\DUrole{std,std-ref}{(Jorgensen)}}}}, which introduces a massless site located a
short distance away from the oxygen atom along the bisector of the HOH
angle.  The atomic types of the oxygen and hydrogen atoms, the bond
and angle types for OH and HOH interactions, and the distance to the
massless charge site are specified as pair\_style arguments.

\begin{sphinxadmonition}{note}{Note:}
\sphinxAtStartPar
For each TIP4P water molecule in your system, the atom IDs for
the O and 2 H atoms must be consecutive, with the O atom first.  This
is to enable LAMMPS to “find” the 2 H atoms associated with each O
atom.  For example, if the atom ID of an O atom in a TIP4P water
molecule is 500, then its 2 H atoms must have IDs 501 and 502.
\end{sphinxadmonition}

\begin{sphinxadmonition}{note}{Note:}
\sphinxAtStartPar
If using type labels, the type labels must be defined before calling
the {\hyperref[\detokenize{pair_coeff::doc}]{\sphinxcrossref{\DUrole{doc}{pair\_coeff}}}} command.
\end{sphinxadmonition}

\sphinxAtStartPar
See the {\hyperref[\detokenize{Howto_tip4p::doc}]{\sphinxcrossref{\DUrole{doc}{Howto tip4p}}}} page for more
information on how to use the TIP4P pair style.  Note that the
neighbor list cutoff for Coulomb interactions is effectively extended
by a distance 2*qdist when using the TIP4P pair style, to account for
the offset distance of the fictitious charges on O atoms in water
molecules.  Thus it is typically best in an efficiency sense to use a
LJ cutoff \textgreater{}= Coulombic cutoff + 2*qdist, to shrink the size of the
neighbor list.  This leads to slightly larger cost for the long\sphinxhyphen{}range
calculation, so you can test the trade\sphinxhyphen{}off for your model.

\sphinxAtStartPar
If \sphinxstyleemphasis{flag\_lj} is set to \sphinxstyleemphasis{long}, no cutoff is used on the LJ 1/r\textasciicircum{}6
dispersion term.  The long\sphinxhyphen{}range portion can be calculated by using
the {\hyperref[\detokenize{kspace_style::doc}]{\sphinxcrossref{\DUrole{doc}{kspace\_style ewald/disp or pppm/disp}}}} commands.
The specified LJ cutoff then determines which portion of the LJ
interactions are computed directly by the pair potential versus which
part is computed in reciprocal space via the Kspace style.  If
\sphinxstyleemphasis{flag\_lj} is set to \sphinxstyleemphasis{cut}, the LJ interactions are simply cutoff, as
with {\hyperref[\detokenize{pair_lj::doc}]{\sphinxcrossref{\DUrole{doc}{pair\_style lj/cut}}}}.

\sphinxAtStartPar
If \sphinxstyleemphasis{flag\_coul} is set to \sphinxstyleemphasis{long}, no cutoff is used on the Coulombic
interactions.  The long\sphinxhyphen{}range portion can calculated by using any of
several {\hyperref[\detokenize{kspace_style::doc}]{\sphinxcrossref{\DUrole{doc}{kspace\_style}}}} command options such as
\sphinxstyleemphasis{pppm} or \sphinxstyleemphasis{ewald}.  Note that if \sphinxstyleemphasis{flag\_lj} is also set to long, then
the \sphinxstyleemphasis{ewald/disp} or \sphinxstyleemphasis{pppm/disp} Kspace style needs to be used to
perform the long\sphinxhyphen{}range calculations for both the LJ and Coulombic
interactions.  If \sphinxstyleemphasis{flag\_coul} is set to \sphinxstyleemphasis{off}, Coulombic interactions
are not computed.

\sphinxAtStartPar
The following coefficients must be defined for each pair of atoms
types via the {\hyperref[\detokenize{pair_coeff::doc}]{\sphinxcrossref{\DUrole{doc}{pair\_coeff}}}} command as in the examples
above, or in the data file or restart files read by the
{\hyperref[\detokenize{read_data::doc}]{\sphinxcrossref{\DUrole{doc}{read\_data}}}} or {\hyperref[\detokenize{read_restart::doc}]{\sphinxcrossref{\DUrole{doc}{read\_restart}}}}
commands, or by mixing as described below:
\begin{itemize}
\item {} 
\sphinxAtStartPar
\(\epsilon\) (energy units)

\item {} 
\sphinxAtStartPar
\(\sigma\) (distance units)

\item {} 
\sphinxAtStartPar
cutoff1 (distance units)

\item {} 
\sphinxAtStartPar
cutoff2 (distance units)

\end{itemize}

\sphinxAtStartPar
Note that sigma is defined in the LJ formula as the zero\sphinxhyphen{}crossing
distance for the potential, not as the energy minimum at 2\textasciicircum{}(1/6)
sigma.

\sphinxAtStartPar
The latter 2 coefficients are optional.  If not specified, the global
LJ and Coulombic cutoffs specified in the pair\_style command are used.
If only one cutoff is specified, it is used as the cutoff for both LJ
and Coulombic interactions for this type pair.  If both coefficients
are specified, they are used as the LJ and Coulombic cutoffs for this
type pair.

\sphinxAtStartPar
Note that if you are using \sphinxstyleemphasis{flag\_lj} set to \sphinxstyleemphasis{long}, you
cannot specify a LJ cutoff for an atom type pair, since only one
global LJ cutoff is allowed.  Similarly, if you are using \sphinxstyleemphasis{flag\_coul}
set to \sphinxstyleemphasis{long}, you cannot specify a Coulombic cutoff for an atom type
pair, since only one global Coulombic cutoff is allowed.

\sphinxAtStartPar
For \sphinxstyleemphasis{lj/long/tip4p/long} only the LJ cutoff can be specified
since a Coulombic cutoff cannot be specified for an individual I,J
type pair.  All type pairs use the same global Coulombic cutoff
specified in the pair\_style command.


\bigskip\hrule\bigskip


\sphinxAtStartPar
A version of these styles with a soft core, \sphinxstyleemphasis{lj/cut/soft}, suitable for use in
free energy calculations, is part of the FEP package and is documented with
the {\hyperref[\detokenize{pair_fep_soft::doc}]{\sphinxcrossref{\DUrole{doc}{pair\_style */soft}}}} styles. The version with soft core is
only available if LAMMPS was built with that package. See the {\hyperref[\detokenize{Build_package::doc}]{\sphinxcrossref{\DUrole{doc}{Build package}}}} page for more info.


\bigskip\hrule\bigskip


\sphinxAtStartPar
Styles with a \sphinxstyleemphasis{gpu}, \sphinxstyleemphasis{intel}, \sphinxstyleemphasis{kk}, \sphinxstyleemphasis{omp}, or \sphinxstyleemphasis{opt} suffix are
functionally the same as the corresponding style without the suffix.
They have been optimized to run faster, depending on your available
hardware, as discussed on the {\hyperref[\detokenize{Speed_packages::doc}]{\sphinxcrossref{\DUrole{doc}{Accelerator packages}}}}
page.  The accelerated styles take the same arguments and should
produce the same results, except for round\sphinxhyphen{}off and precision issues.

\sphinxAtStartPar
These accelerated styles are part of the GPU, INTEL, KOKKOS,
OPENMP, and OPT packages, respectively.  They are only enabled if
LAMMPS was built with those packages.  See the {\hyperref[\detokenize{Build_package::doc}]{\sphinxcrossref{\DUrole{doc}{Build package}}}} page for more info.

\sphinxAtStartPar
You can specify the accelerated styles explicitly in your input script
by including their suffix, or you can use the {\hyperref[\detokenize{Run_options::doc}]{\sphinxcrossref{\DUrole{doc}{\sphinxhyphen{}suffix command\sphinxhyphen{}line switch}}}} when you invoke LAMMPS, or you can use the
{\hyperref[\detokenize{suffix::doc}]{\sphinxcrossref{\DUrole{doc}{suffix}}}} command in your input script.

\sphinxAtStartPar
See the {\hyperref[\detokenize{Speed_packages::doc}]{\sphinxcrossref{\DUrole{doc}{Accelerator packages}}}} page for more
instructions on how to use the accelerated styles effectively.


\bigskip\hrule\bigskip



\subsection{Mixing, shift, table, tail correction, restart, rRESPA info}
\label{\detokenize{pair_lj_long:mixing-shift-table-tail-correction-restart-rrespa-info}}
\sphinxAtStartPar
For atom type pairs I,J and I != J, the epsilon and sigma coefficients
and cutoff distance for all of the lj/long pair styles can be mixed.
The default mix value is \sphinxstyleemphasis{geometric}.  See the “pair\_modify” command
for details.

\sphinxAtStartPar
These pair styles support the {\hyperref[\detokenize{pair_modify::doc}]{\sphinxcrossref{\DUrole{doc}{pair\_modify}}}} shift
option for the energy of the Lennard\sphinxhyphen{}Jones portion of the pair
interaction, assuming \sphinxstyleemphasis{flag\_lj} is \sphinxstyleemphasis{cut}.

\sphinxAtStartPar
These pair styles support the {\hyperref[\detokenize{pair_modify::doc}]{\sphinxcrossref{\DUrole{doc}{pair\_modify}}}} table and
table/disp options since they can tabulate the short\sphinxhyphen{}range portion of
the long\sphinxhyphen{}range Coulombic and dispersion interactions.

\sphinxAtStartPar
Thes pair styles do not support the {\hyperref[\detokenize{pair_modify::doc}]{\sphinxcrossref{\DUrole{doc}{pair\_modify}}}}
tail option for adding a long\sphinxhyphen{}range tail correction to the
Lennard\sphinxhyphen{}Jones portion of the energy and pressure.

\sphinxAtStartPar
These pair styles write their information to {\hyperref[\detokenize{restart::doc}]{\sphinxcrossref{\DUrole{doc}{binary restart files}}}}, so pair\_style and pair\_coeff commands do not need
to be specified in an input script that reads a restart file.

\sphinxAtStartPar
The pair lj/long/coul/long styles support the use of the \sphinxstyleemphasis{inner},
\sphinxstyleemphasis{middle}, and \sphinxstyleemphasis{outer} keywords of the {\hyperref[\detokenize{run_style::doc}]{\sphinxcrossref{\DUrole{doc}{run\_style respa}}}}
command, meaning the pairwise forces can be partitioned by distance at
different levels of the rRESPA hierarchy.  See the
{\hyperref[\detokenize{run_style::doc}]{\sphinxcrossref{\DUrole{doc}{run\_style}}}} command for details.


\bigskip\hrule\bigskip



\subsection{Restrictions}
\label{\detokenize{pair_lj_long:restrictions}}
\sphinxAtStartPar
These styles are part of the KSPACE package.  They are only enabled if
LAMMPS was built with that package.  See the {\hyperref[\detokenize{Build_package::doc}]{\sphinxcrossref{\DUrole{doc}{Build package}}}} page for more info.


\subsection{Related commands}
\label{\detokenize{pair_lj_long:related-commands}}
\sphinxAtStartPar
{\hyperref[\detokenize{pair_coeff::doc}]{\sphinxcrossref{\DUrole{doc}{pair\_coeff}}}}


\subsection{Default}
\label{\detokenize{pair_lj_long:default}}
\sphinxAtStartPar
none


\bigskip\hrule\bigskip

\phantomsection\label{\detokenize{pair_lj_long:veld2}}
\sphinxAtStartPar
\sphinxstylestrong{(In ‘t Veld)} In ‘t Veld, Ismail, Grest, J Chem Phys, 127, 144711 (2007).

\phantomsection\label{\detokenize{pair_lj_long:jorgensen4}}
\sphinxAtStartPar
\sphinxstylestrong{(Jorgensen)} Jorgensen, Chandrasekhar, Madura, Impey, Klein, J Chem
Phys, 79, 926 (1983).

\sphinxstepscope

\index{pair\_style lj/relres@\spxentry{pair\_style lj/relres}}\ignorespaces 
\index{pair\_style lj/relres/omp@\spxentry{pair\_style lj/relres/omp}}\ignorespaces 

\section{pair\_style lj/relres command}
\label{\detokenize{pair_lj_relres:pair-style-lj-relres-command}}\label{\detokenize{pair_lj_relres:index-1}}\label{\detokenize{pair_lj_relres:index-0}}\label{\detokenize{pair_lj_relres::doc}}
\sphinxAtStartPar
Accelerator Variants: \sphinxstyleemphasis{lj/relres/omp}


\subsection{Syntax}
\label{\detokenize{pair_lj_relres:syntax}}
\begin{sphinxVerbatim}[commandchars=\\\{\}]
\PYG{k}{pair\PYGZus{}style}\PYG{+w}{ }\PYG{n}{lj}\PYG{o}{/}\PYG{n}{relres}\PYG{+w}{ }\PYG{n}{Rsi}\PYG{+w}{ }\PYG{n}{Rso}\PYG{+w}{ }\PYG{n}{Rci}\PYG{+w}{ }\PYG{n}{Rco}
\end{sphinxVerbatim}
\begin{itemize}
\item {} 
\sphinxAtStartPar
Rsi = inner switching cutoff between the fine\sphinxhyphen{}grained and coarse\sphinxhyphen{}grained potentials (distance units)

\item {} 
\sphinxAtStartPar
Rso = outer switching cutoff between the fine\sphinxhyphen{}grained and coarse\sphinxhyphen{}grained potentials (distance units)

\item {} 
\sphinxAtStartPar
Rci = inner cutoff beyond which the force smoothing for all interactions is applied (distance units)

\item {} 
\sphinxAtStartPar
Rco = outer cutoff for all interactions (distance units)

\end{itemize}


\subsection{Examples}
\label{\detokenize{pair_lj_relres:examples}}
\begin{sphinxVerbatim}[commandchars=\\\{\}]
\PYG{k}{pair\PYGZus{}style}\PYG{+w}{ }\PYG{n}{lj}\PYG{o}{/}\PYG{n}{relres}\PYG{+w}{ }\PYG{l+m}{4.0}\PYG{+w}{ }\PYG{l+m}{5.0}\PYG{+w}{ }\PYG{l+m}{8.0}\PYG{+w}{ }\PYG{l+m}{10.0}
\PYG{k}{pair\PYGZus{}coeff}\PYG{+w}{ }\PYG{l+m}{1}\PYG{+w}{ }\PYG{l+m}{1}\PYG{+w}{ }\PYG{l+m}{0.5}\PYG{+w}{ }\PYG{l+m}{1.0}\PYG{+w}{ }\PYG{l+m}{1.5}\PYG{+w}{ }\PYG{l+m}{1.1}
\PYG{k}{pair\PYGZus{}coeff}\PYG{+w}{ }\PYG{l+m}{2}\PYG{+w}{ }\PYG{l+m}{2}\PYG{+w}{ }\PYG{l+m}{0.5}\PYG{+w}{ }\PYG{l+m}{1.0}\PYG{+w}{ }\PYG{l+m}{0.0}\PYG{+w}{ }\PYG{l+m}{0.0}\PYG{+w}{ }\PYG{l+m}{3.0}\PYG{+w}{ }\PYG{l+m}{3.5}\PYG{+w}{ }\PYG{l+m}{6.0}\PYG{+w}{ }\PYG{l+m}{7.0}
\end{sphinxVerbatim}


\subsection{Description}
\label{\detokenize{pair_lj_relres:description}}
\sphinxAtStartPar
Pair style \sphinxstyleemphasis{lj/relres} computes a LJ interaction using the Relative
Resolution (RelRes) framework which applies a fine\sphinxhyphen{}grained (FG)
potential between near neighbors and a coarse\sphinxhyphen{}grained (CG) potential
between far neighbors {\hyperref[\detokenize{pair_lj_relres:chaimovich1}]{\sphinxcrossref{\DUrole{std,std-ref}{(Chaimovich1)}}}}. This approach
can improve the computational efficiency by almost an order of
magnitude, while maintaining the correct static and dynamic behavior of
a reference system {\hyperref[\detokenize{pair_lj_relres:chaimovich2}]{\sphinxcrossref{\DUrole{std,std-ref}{(Chaimovich2)}}}}.
\begin{equation*}
\begin{split}E = \left\{\begin{array}{lr}
     4 \epsilon^{\scriptscriptstyle FG} \left[ \left(\frac{\sigma^{FG}}{r}\right)^{12} - \left(\frac{\sigma^{FG}}{r}\right)^6 \right]-\Gamma_{si}, & \quad\mathrm{if}\quad  r< r_{si}, \\
     \sum_{m=0}^{4} \gamma_{sm}\left(r-r_{si}\right)^m-\Gamma_{so} ,   & \quad\mathrm{if}\quad  r_{si}\leq r< r_{so}, \\
     4 \epsilon^{\scriptscriptstyle CG} \left[ \left(\frac{\sigma^{CG}}{r}\right)^{12} -     \left(\frac{\sigma^{CG}}{r}\right)^6 \right]-\Gamma_c, &  \quad\mathrm{if}\quad  r_{so}\leq r<r_{ci}, \\
     \sum_{m=0}^{4} \gamma_{cm}\left(r-r_{ci}\right)^m -\Gamma_c, & \quad\mathrm{if}\quad  r_{ci}\leq r< r_{co}, \\
     0, & \quad\mathrm{if}\quad  r\geq r_{co}.\end{array}\right.\end{split}
\end{equation*}
\sphinxAtStartPar
The FG parameters of the LJ potential (\(\epsilon^{FG}\) and
\(\sigma^{FG}\)) are applied up to the inner switching cutoff,
\(r_{si}\), while the CG parameters of the LJ potential
(\(\epsilon^{CG}\) and \(\sigma^{CG}\)) are applied beyond the
outer switching cutoff, \(r_{so}\). Between \(r_{si}\) and
\(r_{so}\) a polynomial smoothing function is applied so that the
force and its derivative are continuous between the FG and CG
potentials.  An analogous smoothing function is applied between the
inner and outer cutoffs (\(r_{ci}\) and \(r_{co}\)).
The offsets \(\Gamma_{si}\), \(\Gamma_{so}\) and
\(\Gamma_{c}\) ensure the continuity of the energy over the entire
domain.  The corresponding polynomial coefficients \(\gamma_{sm}\)
and \(\gamma_{cm}\), as well as the offsets are automatically
computed by LAMMPS.

\begin{sphinxadmonition}{note}{Note:}
\sphinxAtStartPar
Energy and force resulting from this methodology can be plotted via the
{\hyperref[\detokenize{pair_write::doc}]{\sphinxcrossref{\DUrole{doc}{pair\_write}}}} command.
\end{sphinxadmonition}

\sphinxAtStartPar
The following coefficients must be defined for each pair of atom types
via the {\hyperref[\detokenize{pair_coeff::doc}]{\sphinxcrossref{\DUrole{doc}{pair\_coeff}}}} command as in the examples above,
or in the data file or restart files read by the {\hyperref[\detokenize{read_data::doc}]{\sphinxcrossref{\DUrole{doc}{read\_data}}}} or {\hyperref[\detokenize{read_restart::doc}]{\sphinxcrossref{\DUrole{doc}{read\_restart}}}} commands, or by
mixing as will be described below:
\begin{itemize}
\item {} 
\sphinxAtStartPar
\(\epsilon^{FG}\) (energy units)

\item {} 
\sphinxAtStartPar
\(\sigma^{FG}\) (distance units)

\item {} 
\sphinxAtStartPar
\(\epsilon^{CG}\) (energy units)

\item {} 
\sphinxAtStartPar
\(\sigma^{CG}\) (distance units)

\end{itemize}

\sphinxAtStartPar
Additional parameters can be defined to specify different
\(r_{si}\), \(r_{so}\), \(r_{ci}\), \(r_{co}\) for
a particular set of atom types:
\begin{itemize}
\item {} 
\sphinxAtStartPar
\(r_{si}\) (distance units)

\item {} 
\sphinxAtStartPar
\(r_{so}\) (distance units)

\item {} 
\sphinxAtStartPar
\(r_{ci}\) (distance units)

\item {} 
\sphinxAtStartPar
\(r_{co}\) (distance units)

\end{itemize}

\sphinxAtStartPar
These parameters are optional, and they are used to override the global
cutoffs as defined in the pair\_style command.  If not specified, the
global values for \(r_{si}\), \(r_{so}\), \(r_{ci}\), and
\(r_{co}\) are used.  If this override option is employed, all four
arguments must be specified.


\bigskip\hrule\bigskip


\sphinxAtStartPar
Here are some guidelines for using the pair\_style \sphinxstyleemphasis{lj/relres} command.

\sphinxAtStartPar
In general, RelRes focuses on the speedup of pairwise interactions between
all LJ sites. Importantly, it works with any settings and flags (e.g.,
{\hyperref[\detokenize{special_bonds::doc}]{\sphinxcrossref{\DUrole{doc}{special\_bonds}}}} settings and {\hyperref[\detokenize{newton::doc}]{\sphinxcrossref{\DUrole{doc}{newton}}}}
flags) that can be used in a molecular simulation with the
conventional LJ potential. In particular, all intramolecular topology
with its energetics (i.e., bonds, angles, etc.) remains unaltered.

\sphinxAtStartPar
At the most basic level in the RelRes framework, all sites are mapped into
clusters. Each cluster is just a collection of sites bonded together (the
bonds themselves are not part of the cluster). In general, a molecule may
be comprised of several clusters, and preferably, no two sites in a cluster
are separated by more than two bonds. There are two categories of sites in
RelRes: “hybrid” sites embody both FG and CG models, while “ordinary” sites
embody just FG characteristics with no CG features. A given cluster has
a single hybrid site (typically its central site) and several ordinary sites
(typically its peripheral sites). Notice that while clusters are necessary
for the RelRes parameterization (discussed below), they are not actually
defined in LAMMPS. Besides, the total number of sites in the cluster are
called the “mapping ratio”, and this substantially impacts the computational
efficiency of RelRes: For a mapping ratio of 3, the efficiency factor is
around 4, and for a mapping ratio of 5, the efficiency factor is around 5
{\hyperref[\detokenize{pair_lj_relres:chaimovich2}]{\sphinxcrossref{\DUrole{std,std-ref}{(Chaimovich2)}}}}.

\sphinxAtStartPar
The flexibility of LAMMPS allows placing any values for the LJ
parameters in the input script. However, here are the optimal
recommendations for the RelRes parameters, which yield the correct
structural and thermal behavior in a system of interest
{\hyperref[\detokenize{pair_lj_relres:chaimovich1}]{\sphinxcrossref{\DUrole{std,std-ref}{(Chaimovich1)}}}}.  One must first assign a complete set of
parameters for the FG interactions that are applicable to all atom types.
Regarding the parameters for the CG interactions, the rules rely on the
site category (if it is a hybrid or an ordinary site). For atom types of
ordinary sites, \(\epsilon^{CG}\) must be set to 0 (zero) while the
specific value of \(\sigma^{CG}\) is irrelevant. For atom types of
hybrid sites, the CG parameters should be generally calculated using the
following equations:
\begin{equation*}
\begin{split}\sigma_I^{CG}=\frac{\left((\sum_{\alpha\in A}\sqrt{\epsilon_\alpha^{FG}\left(\sigma_\alpha^{FG}\right)^{12}}\right)^{1/2}}{\left((\sum_{\alpha\in A}\sqrt{\epsilon_\alpha^{FG}\left(\sigma_\alpha^{FG}\right)^6}\right)^{1/3}}
\quad\mathrm{and}\quad
\epsilon_I^{CG}=\frac{\left((\sum_{\alpha\in A}\sqrt{\epsilon_\alpha^{FG}\left(\sigma_\alpha^{FG}\right)^6}\right)^4}{\left((\sum_{\alpha\in A}\sqrt{\epsilon_\alpha^{FG}\left(\sigma_\alpha^{FG}\right)^{12}}\right)^2}\end{split}
\end{equation*}
\sphinxAtStartPar
where \(I\) is an atom type of a hybrid site of a particular cluster
\(A\), and corresponding with this cluster, the summation proceeds over
all of its sites \(\alpha\).  These equations are derived from the
monopole term in the underlying Taylor series, and they are indeed relevant
only if geometric mixing is applicable for the FG model; if this is not the
case, Ref. {\hyperref[\detokenize{pair_lj_relres:chaimovich2}]{\sphinxcrossref{\DUrole{std,std-ref}{(Chaimovich2)}}}} discusses the alternative
formula, and in such a situation, the pair\_coeff command should be explicitly
used for all combinations of atom types \(I\;!=J\).

\sphinxAtStartPar
The switching distance (the midpoint between inner and outer switching
cutoffs) is another crucial factor in RelRes: decreasing it improves the
computational efficiency, yet if it is too small, the molecular simulations
may not capture the system behavior correctly.  As a rule of thumb,
the switching distance should be approximately \(\,\sim\! 1.5\sigma\)
{\hyperref[\detokenize{pair_lj_relres:chaimovich1}]{\sphinxcrossref{\DUrole{std,std-ref}{(Chaimovich1)}}}}; recommendations can be found in
Ref. {\hyperref[\detokenize{pair_lj_relres:chaimovich2}]{\sphinxcrossref{\DUrole{std,std-ref}{(Chaimovich2)}}}}.
Regarding the switching smoothing zone, \(\,\sim\!0.1\sigma\) is
recommended; if desired, smoothing can be eliminated by setting
the inner switching cutoff, \(r_{si}\), equal to the outer
switching cutoff, \(r_{so}\) (the same is true for the other cutoffs
\(r_{ci}\) and \(r_{co}\)).


\bigskip\hrule\bigskip


\sphinxAtStartPar
As an example, imagine that in your system, a molecule is comprised just
of one cluster such that one atom type (\#1) is associated with
its hybrid site, and another atom type (\#2) is associated with its ordinary
sites (in total, there are 2 atom types). If geometric mixing is applicable,
the following commands should be used:

\begin{sphinxVerbatim}[commandchars=\\\{\}]
\PYG{k}{pair\PYGZus{}style}\PYG{+w}{ }\PYG{n}{lj}\PYG{o}{/}\PYG{n}{relres}\PYG{+w}{ }\PYG{n}{Rsi}\PYG{+w}{ }\PYG{n}{Rso}\PYG{+w}{ }\PYG{n}{Rci}\PYG{+w}{ }\PYG{n}{Rco}
\PYG{k}{pair\PYGZus{}coeff}\PYG{+w}{ }\PYG{l+m}{1}\PYG{+w}{ }\PYG{l+m}{1}\PYG{+w}{ }\PYG{n}{epsilon\PYGZus{}FG1}\PYG{+w}{ }\PYG{n}{sigma\PYGZus{}FG1}\PYG{+w}{ }\PYG{n}{epsilon\PYGZus{}CG1}\PYG{+w}{ }\PYG{n}{sigma\PYGZus{}CG1}
\PYG{k}{pair\PYGZus{}coeff}\PYG{+w}{ }\PYG{l+m}{2}\PYG{+w}{ }\PYG{l+m}{2}\PYG{+w}{ }\PYG{n}{epsilon\PYGZus{}FG2}\PYG{+w}{ }\PYG{n}{sigma\PYGZus{}FG2}\PYG{+w}{ }\PYG{l+m}{0.0}\PYG{+w}{         }\PYG{l+m}{0.0}
\PYG{k}{pair\PYGZus{}modify}\PYG{+w}{ }\PYG{n}{shift}\PYG{+w}{ }\PYG{n}{yes}
\end{sphinxVerbatim}

\sphinxAtStartPar
In a more complex situation, there may be two distinct clusters in a system
(these two clusters may be on same molecule or on different molecules),
each with its own switching cutoffs.  If there are still two atom types
in each cluster as in the earlier example, the commands should be:

\begin{sphinxVerbatim}[commandchars=\\\{\}]
\PYG{k}{pair\PYGZus{}style}\PYG{+w}{ }\PYG{n}{lj}\PYG{o}{/}\PYG{n}{relres}\PYG{+w}{ }\PYG{n}{Rsi}\PYG{+w}{ }\PYG{n}{Rso}\PYG{+w}{ }\PYG{n}{Rci}\PYG{+w}{ }\PYG{n}{Rco}
\PYG{k}{pair\PYGZus{}coeff}\PYG{+w}{ }\PYG{l+m}{1}\PYG{+w}{ }\PYG{l+m}{1}\PYG{+w}{ }\PYG{n}{epsilon\PYGZus{}FG1}\PYG{+w}{ }\PYG{n}{sigma\PYGZus{}FG1}\PYG{+w}{ }\PYG{n}{epsilon\PYGZus{}CG1}\PYG{+w}{ }\PYG{n}{sigma\PYGZus{}CG1}\PYG{+w}{ }\PYG{n}{Rsi1}\PYG{+w}{ }\PYG{n}{Rso1}\PYG{+w}{ }\PYG{n}{Rci}\PYG{+w}{ }\PYG{n}{Rco}
\PYG{k}{pair\PYGZus{}coeff}\PYG{+w}{ }\PYG{l+m}{2}\PYG{+w}{ }\PYG{l+m}{2}\PYG{+w}{ }\PYG{n}{epsilon\PYGZus{}FG2}\PYG{+w}{ }\PYG{n}{sigma\PYGZus{}FG2}\PYG{+w}{ }\PYG{l+m}{0.0}\PYG{+w}{         }\PYG{l+m}{0.0}\PYG{+w}{       }\PYG{n}{Rsi1}\PYG{+w}{ }\PYG{n}{Rso1}\PYG{+w}{ }\PYG{n}{Rci}\PYG{+w}{ }\PYG{n}{Rco}
\PYG{k}{pair\PYGZus{}coeff}\PYG{+w}{ }\PYG{l+m}{3}\PYG{+w}{ }\PYG{l+m}{3}\PYG{+w}{ }\PYG{n}{epsilon\PYGZus{}FG3}\PYG{+w}{ }\PYG{n}{sigma\PYGZus{}FG3}\PYG{+w}{ }\PYG{n}{epsilon\PYGZus{}CG3}\PYG{+w}{ }\PYG{n}{sigma\PYGZus{}CG3}
\PYG{k}{pair\PYGZus{}coeff}\PYG{+w}{ }\PYG{l+m}{4}\PYG{+w}{ }\PYG{l+m}{4}\PYG{+w}{ }\PYG{n}{epsilon\PYGZus{}FG4}\PYG{+w}{ }\PYG{n}{sigma\PYGZus{}FG4}\PYG{+w}{ }\PYG{l+m}{0.0}\PYG{+w}{         }\PYG{l+m}{0.0}
\PYG{k}{pair\PYGZus{}modify}\PYG{+w}{ }\PYG{n}{shift}\PYG{+w}{ }\PYG{n}{yes}
\end{sphinxVerbatim}

\sphinxAtStartPar
In this example, the switching cutoffs for the first cluster (atom types 1
and 2) is defined explicitly in the pair\_coeff command which overrides the
global values, while the second cluster (atom types 3 and 4) uses the global
definition from the pair\_style command. The emphasis here is that the atom
types that belong to a specific cluster should have the same switching/cutoff
arguments.

\sphinxAtStartPar
In the case that geometric mixing is not applicable, for simulating the
system from the previous example, we recommend using the following commands:

\begin{sphinxVerbatim}[commandchars=\\\{\}]
\PYG{k}{pair\PYGZus{}style}\PYG{+w}{ }\PYG{n}{lj}\PYG{o}{/}\PYG{n}{relres}\PYG{+w}{ }\PYG{n}{Rsi}\PYG{+w}{ }\PYG{n}{Rso}\PYG{+w}{ }\PYG{n}{Rci}\PYG{+w}{ }\PYG{n}{Rco}
\PYG{k}{pair\PYGZus{}coeff}\PYG{+w}{ }\PYG{l+m}{1}\PYG{+w}{ }\PYG{l+m}{1}\PYG{+w}{ }\PYG{n}{epsilon\PYGZus{}FG1}\PYG{+w}{  }\PYG{n}{sigma\PYGZus{}FG1}\PYG{+w}{  }\PYG{n}{epsilon\PYGZus{}CG1}\PYG{+w}{  }\PYG{n}{sigma\PYGZus{}CG1}\PYG{+w}{  }\PYG{n}{Rsi1}\PYG{+w}{  }\PYG{n}{Rso1}\PYG{+w}{  }\PYG{n}{Rci}\PYG{+w}{ }\PYG{n}{Rco}
\PYG{k}{pair\PYGZus{}coeff}\PYG{+w}{ }\PYG{l+m}{1}\PYG{+w}{ }\PYG{l+m}{2}\PYG{+w}{ }\PYG{n}{epsilon\PYGZus{}FG12}\PYG{+w}{ }\PYG{n}{sigma\PYGZus{}FG12}\PYG{+w}{ }\PYG{l+m}{0.0}\PYG{+w}{          }\PYG{l+m}{0.0}\PYG{+w}{        }\PYG{n}{Rsi1}\PYG{+w}{  }\PYG{n}{Rso1}\PYG{+w}{  }\PYG{n}{Rci}\PYG{+w}{ }\PYG{n}{Rco}
\PYG{k}{pair\PYGZus{}coeff}\PYG{+w}{ }\PYG{l+m}{1}\PYG{+w}{ }\PYG{l+m}{3}\PYG{+w}{ }\PYG{n}{epsilon\PYGZus{}FG13}\PYG{+w}{ }\PYG{n}{sigma\PYGZus{}FG13}\PYG{+w}{ }\PYG{n}{epsilon\PYGZus{}CG13}\PYG{+w}{ }\PYG{n}{sigma\PYGZus{}CG13}\PYG{+w}{ }\PYG{n}{Rsi13}\PYG{+w}{ }\PYG{n}{Rso13}\PYG{+w}{ }\PYG{n}{Rci}\PYG{+w}{ }\PYG{n}{Rco}
\PYG{k}{pair\PYGZus{}coeff}\PYG{+w}{ }\PYG{l+m}{1}\PYG{+w}{ }\PYG{l+m}{4}\PYG{+w}{ }\PYG{n}{epsilon\PYGZus{}FG14}\PYG{+w}{ }\PYG{n}{sigma\PYGZus{}FG14}\PYG{+w}{ }\PYG{l+m}{0.0}\PYG{+w}{          }\PYG{l+m}{0.0}\PYG{+w}{        }\PYG{n}{Rsi13}\PYG{+w}{ }\PYG{n}{Rso13}\PYG{+w}{ }\PYG{n}{Rci}\PYG{+w}{ }\PYG{n}{Rco}
\PYG{k}{pair\PYGZus{}coeff}\PYG{+w}{ }\PYG{l+m}{2}\PYG{+w}{ }\PYG{l+m}{2}\PYG{+w}{ }\PYG{n}{epsilon\PYGZus{}FG2}\PYG{+w}{  }\PYG{n}{sigma\PYGZus{}FG2}\PYG{+w}{  }\PYG{l+m}{0.0}\PYG{+w}{          }\PYG{l+m}{0.0}\PYG{+w}{        }\PYG{n}{Rsi1}\PYG{+w}{  }\PYG{n}{Rso1}\PYG{+w}{  }\PYG{n}{Rci}\PYG{+w}{ }\PYG{n}{Rco}
\PYG{k}{pair\PYGZus{}coeff}\PYG{+w}{ }\PYG{l+m}{2}\PYG{+w}{ }\PYG{l+m}{3}\PYG{+w}{ }\PYG{n}{epsilon\PYGZus{}FG23}\PYG{+w}{ }\PYG{n}{sigma\PYGZus{}FG23}\PYG{+w}{ }\PYG{l+m}{0.0}\PYG{+w}{          }\PYG{l+m}{0.0}\PYG{+w}{        }\PYG{n}{Rsi13}\PYG{+w}{ }\PYG{n}{Rso13}\PYG{+w}{ }\PYG{n}{Rci}\PYG{+w}{ }\PYG{n}{Rco}
\PYG{k}{pair\PYGZus{}coeff}\PYG{+w}{ }\PYG{l+m}{2}\PYG{+w}{ }\PYG{l+m}{4}\PYG{+w}{ }\PYG{n}{epsilon\PYGZus{}FG24}\PYG{+w}{ }\PYG{n}{sigma\PYGZus{}FG24}\PYG{+w}{ }\PYG{l+m}{0.0}\PYG{+w}{          }\PYG{l+m}{0.0}\PYG{+w}{        }\PYG{n}{Rsi13}\PYG{+w}{ }\PYG{n}{Rso13}\PYG{+w}{ }\PYG{n}{Rci}\PYG{+w}{ }\PYG{n}{Rco}
\PYG{k}{pair\PYGZus{}coeff}\PYG{+w}{ }\PYG{l+m}{3}\PYG{+w}{ }\PYG{l+m}{3}\PYG{+w}{ }\PYG{n}{epsilon\PYGZus{}FG3}\PYG{+w}{  }\PYG{n}{sigma\PYGZus{}FG3}\PYG{+w}{  }\PYG{n}{epsilon\PYGZus{}CG3}\PYG{+w}{  }\PYG{n}{sigma\PYGZus{}CG3}
\PYG{k}{pair\PYGZus{}coeff}\PYG{+w}{ }\PYG{l+m}{3}\PYG{+w}{ }\PYG{l+m}{4}\PYG{+w}{ }\PYG{n}{epsilon\PYGZus{}FG34}\PYG{+w}{ }\PYG{n}{sigma\PYGZus{}FG34}\PYG{+w}{ }\PYG{l+m}{0.0}\PYG{+w}{          }\PYG{l+m}{0.0}
\PYG{k}{pair\PYGZus{}coeff}\PYG{+w}{ }\PYG{l+m}{4}\PYG{+w}{ }\PYG{l+m}{4}\PYG{+w}{ }\PYG{n}{epsilon\PYGZus{}FG4}\PYG{+w}{  }\PYG{n}{sigma\PYGZus{}FG4}\PYG{+w}{  }\PYG{l+m}{0.0}\PYG{+w}{          }\PYG{l+m}{0.0}
\PYG{k}{pair\PYGZus{}modify}\PYG{+w}{ }\PYG{n}{shift}\PYG{+w}{ }\PYG{n}{yes}
\end{sphinxVerbatim}

\sphinxAtStartPar
Notice that the CG parameters are mixed only for interactions between atom
types associated with hybrid sites, and that the cutoffs are
mixed on the cluster basis.

\sphinxAtStartPar
More examples can be found in the \sphinxstyleemphasis{examples/relres} folder.


\bigskip\hrule\bigskip


\sphinxAtStartPar
Styles with a \sphinxstyleemphasis{gpu}, \sphinxstyleemphasis{intel}, \sphinxstyleemphasis{kk}, \sphinxstyleemphasis{omp}, or \sphinxstyleemphasis{opt} suffix are
functionally the same as the corresponding style without the suffix.
They have been optimized to run faster, depending on your available
hardware, as discussed on the {\hyperref[\detokenize{Speed_packages::doc}]{\sphinxcrossref{\DUrole{doc}{Accelerator packages}}}}
page.  The accelerated styles take the same arguments and should
produce the same results, except for round\sphinxhyphen{}off and precision issues.

\sphinxAtStartPar
These accelerated styles are part of the GPU, INTEL, KOKKOS,
OPENMP, and OPT packages, respectively.  They are only enabled if
LAMMPS was built with those packages.  See the {\hyperref[\detokenize{Build_package::doc}]{\sphinxcrossref{\DUrole{doc}{Build package}}}} page for more info.

\sphinxAtStartPar
You can specify the accelerated styles explicitly in your input script
by including their suffix, or you can use the {\hyperref[\detokenize{Run_options::doc}]{\sphinxcrossref{\DUrole{doc}{\sphinxhyphen{}suffix command\sphinxhyphen{}line switch}}}} when you invoke LAMMPS, or you can use the
{\hyperref[\detokenize{suffix::doc}]{\sphinxcrossref{\DUrole{doc}{suffix}}}} command in your input script.

\sphinxAtStartPar
See the {\hyperref[\detokenize{Speed_packages::doc}]{\sphinxcrossref{\DUrole{doc}{Accelerator packages}}}} page for more
instructions on how to use the accelerated styles effectively.


\bigskip\hrule\bigskip



\subsection{Mixing, shift, table, tail correction, restart, rRESPA info}
\label{\detokenize{pair_lj_relres:mixing-shift-table-tail-correction-restart-rrespa-info}}
\sphinxAtStartPar
For atom type pairs \(I,\:J\) with \(I\;!=J\), the
\(\epsilon^{FG}\), \(\sigma^{FG}\), \(\epsilon^{CG}\),
\(\sigma^{CG}\), \(r_{si}\), \(r_{so}\), \(r_{ci}\),
and \(r_{co}\) parameters for this pair style can be mixed, if
not defined explicitly. All parameters are mixed according to the
pair\_modify mix option.  The default mix value is \sphinxstyleemphasis{geometric},
and it is recommended to use with this \sphinxstyleemphasis{lj/relres} style.  See the
“pair\_modify” command for details.

\sphinxAtStartPar
This pair style supports the {\hyperref[\detokenize{pair_modify::doc}]{\sphinxcrossref{\DUrole{doc}{pair\_modify}}}} shift
option for the energy of the pair interaction. It is recommended to set
this option to \sphinxstyleemphasis{yes}.  Otherwise, the offset \(\Gamma_{c}\)
is set to zero. Constants \(\Gamma_{si}\) and \(\Gamma_{so}\) are
not impacted by this option.

\sphinxAtStartPar
The {\hyperref[\detokenize{pair_modify::doc}]{\sphinxcrossref{\DUrole{doc}{pair\_modify}}}} table option is not relevant
for this pair style.

\sphinxAtStartPar
This pair style does not support the {\hyperref[\detokenize{pair_modify::doc}]{\sphinxcrossref{\DUrole{doc}{pair\_modify}}}}
tail option for adding long\sphinxhyphen{}range tail corrections to energy and
pressure, since the energy of the pair interaction is smoothed to 0.0
at the cutoff.

\sphinxAtStartPar
This pair style writes its information to {\hyperref[\detokenize{restart::doc}]{\sphinxcrossref{\DUrole{doc}{binary restart files}}}}, so pair\_style and pair\_coeff commands do not need to be
specified in an input script that reads a restart file.

\sphinxAtStartPar
This pair style can only be used via the \sphinxstyleemphasis{pair} keyword of the
{\hyperref[\detokenize{run_style::doc}]{\sphinxcrossref{\DUrole{doc}{run\_style respa}}}} command.  It does not support the
\sphinxstyleemphasis{inner}, \sphinxstyleemphasis{middle}, \sphinxstyleemphasis{outer} keywords.


\bigskip\hrule\bigskip



\subsection{Restrictions}
\label{\detokenize{pair_lj_relres:restrictions}}
\sphinxAtStartPar
This pair style is part of the EXTRA\sphinxhyphen{}PAIR package.  It is only enabled if
LAMMPS was built with that package.  See the
{\hyperref[\detokenize{Build_package::doc}]{\sphinxcrossref{\DUrole{doc}{Build package}}}} page for more info.


\subsection{Related commands}
\label{\detokenize{pair_lj_relres:related-commands}}
\sphinxAtStartPar
{\hyperref[\detokenize{pair_coeff::doc}]{\sphinxcrossref{\DUrole{doc}{pair\_coeff}}}}


\subsection{Default}
\label{\detokenize{pair_lj_relres:default}}
\sphinxAtStartPar
none


\bigskip\hrule\bigskip

\phantomsection\label{\detokenize{pair_lj_relres:chaimovich1}}
\sphinxAtStartPar
\sphinxstylestrong{(Chaimovich1)} A. Chaimovich, C. Peter and K. Kremer, J. Chem. Phys. 143,
243107 (2015).

\phantomsection\label{\detokenize{pair_lj_relres:chaimovich2}}
\sphinxAtStartPar
\sphinxstylestrong{(Chaimovich2)} M. Chaimovich and A. Chaimovich, J. Chem. Theory Comput. 17,
1045\sphinxhyphen{}1059 (2021).

\sphinxstepscope

\index{pair\_style lj/smooth@\spxentry{pair\_style lj/smooth}}\ignorespaces 
\index{pair\_style lj/smooth/gpu@\spxentry{pair\_style lj/smooth/gpu}}\ignorespaces 
\index{pair\_style lj/smooth/omp@\spxentry{pair\_style lj/smooth/omp}}\ignorespaces 

\section{pair\_style lj/smooth command}
\label{\detokenize{pair_lj_smooth:pair-style-lj-smooth-command}}\label{\detokenize{pair_lj_smooth:index-2}}\label{\detokenize{pair_lj_smooth:index-1}}\label{\detokenize{pair_lj_smooth:index-0}}\label{\detokenize{pair_lj_smooth::doc}}
\sphinxAtStartPar
Accelerator Variants: \sphinxstyleemphasis{lj/smooth/gpu}, \sphinxstyleemphasis{lj/smooth/omp}


\subsection{Syntax}
\label{\detokenize{pair_lj_smooth:syntax}}
\begin{sphinxVerbatim}[commandchars=\\\{\}]
\PYG{k}{pair\PYGZus{}style}\PYG{+w}{ }\PYG{n}{lj}\PYG{o}{/}\PYG{n}{smooth}\PYG{+w}{ }\PYG{n}{Rin}\PYG{+w}{ }\PYG{n}{Rc}
\end{sphinxVerbatim}
\begin{itemize}
\item {} 
\sphinxAtStartPar
Rin = inner cutoff beyond which force smoothing will be applied (distance units)

\item {} 
\sphinxAtStartPar
Rc = outer cutoff for lj/smooth interactions (distance units)

\end{itemize}


\subsection{Examples}
\label{\detokenize{pair_lj_smooth:examples}}
\begin{sphinxVerbatim}[commandchars=\\\{\}]
\PYG{k}{pair\PYGZus{}style}\PYG{+w}{ }\PYG{n}{lj}\PYG{o}{/}\PYG{n}{smooth}\PYG{+w}{ }\PYG{l+m}{8.0}\PYG{+w}{ }\PYG{l+m}{10.0}
\PYG{k}{pair\PYGZus{}coeff}\PYG{+w}{ }\PYG{o}{*}\PYG{+w}{ }\PYG{o}{*}\PYG{+w}{ }\PYG{l+m}{10.0}\PYG{+w}{ }\PYG{l+m}{1.5}
\PYG{k}{pair\PYGZus{}coeff}\PYG{+w}{ }\PYG{l+m}{1}\PYG{+w}{ }\PYG{l+m}{1}\PYG{+w}{ }\PYG{l+m}{20.0}\PYG{+w}{ }\PYG{l+m}{1.3}\PYG{+w}{ }\PYG{l+m}{7.0}\PYG{+w}{ }\PYG{l+m}{9.0}
\end{sphinxVerbatim}


\subsection{Description}
\label{\detokenize{pair_lj_smooth:description}}
\sphinxAtStartPar
Style \sphinxstyleemphasis{lj/smooth} computes a LJ interaction with a force smoothing
applied between the inner and outer cutoff.
\begin{equation*}
\begin{split}E & =  4 \epsilon \left[ \left(\frac{\sigma}{r}\right)^{12} -
                      \left(\frac{\sigma}{r}\right)^6 \right]
                      \qquad r < r_{in} \\
F & =  C_1 + C_2 (r - r_{in}) + C_3 (r - r_{in})^2 + C_4 (r - r_{in})^3
                    \qquad r_{in} < r < r_c\end{split}
\end{equation*}
\sphinxAtStartPar
The polynomial coefficients C1, C2, C3, C4 are computed by LAMMPS to
cause the force to vary smoothly from the inner cutoff \(r_{in}\) to the
outer cutoff \(r_c\).

\sphinxAtStartPar
At the inner cutoff the force and its first derivative
will match the non\sphinxhyphen{}smoothed LJ formula.  At the outer cutoff the force
and its first derivative will be 0.0.  The inner cutoff cannot be 0.0.

\begin{sphinxadmonition}{note}{Note:}
\sphinxAtStartPar
this force smoothing causes the energy to be discontinuous both
in its values and first derivative.  This can lead to poor energy
conservation and may require the use of a thermostat.  Plot the energy
and force resulting from this formula via the
{\hyperref[\detokenize{pair_write::doc}]{\sphinxcrossref{\DUrole{doc}{pair\_write}}}} command to see the effect.
\end{sphinxadmonition}

\sphinxAtStartPar
The following coefficients must be defined for each pair of atoms
types via the {\hyperref[\detokenize{pair_coeff::doc}]{\sphinxcrossref{\DUrole{doc}{pair\_coeff}}}} command as in the examples
above, or in the data file or restart files read by the
{\hyperref[\detokenize{read_data::doc}]{\sphinxcrossref{\DUrole{doc}{read\_data}}}} or {\hyperref[\detokenize{read_restart::doc}]{\sphinxcrossref{\DUrole{doc}{read\_restart}}}}
commands, or by mixing as described below:
\begin{itemize}
\item {} 
\sphinxAtStartPar
\(\epsilon\) (energy units)

\item {} 
\sphinxAtStartPar
\(\sigma\) (distance units)

\item {} 
\sphinxAtStartPar
\(r_{in}\) (distance units)

\item {} 
\sphinxAtStartPar
\(r_c\) (distance units)

\end{itemize}

\sphinxAtStartPar
The last 2 coefficients are optional inner and outer cutoffs.  If not
specified, the global values for \(r_{in}\) and \(r_c\) are used.


\bigskip\hrule\bigskip


\sphinxAtStartPar
Styles with a \sphinxstyleemphasis{gpu}, \sphinxstyleemphasis{intel}, \sphinxstyleemphasis{kk}, \sphinxstyleemphasis{omp}, or \sphinxstyleemphasis{opt} suffix are
functionally the same as the corresponding style without the suffix.
They have been optimized to run faster, depending on your available
hardware, as discussed on the {\hyperref[\detokenize{Speed_packages::doc}]{\sphinxcrossref{\DUrole{doc}{Accelerator packages}}}}
page.  The accelerated styles take the same arguments and should
produce the same results, except for round\sphinxhyphen{}off and precision issues.

\sphinxAtStartPar
These accelerated styles are part of the GPU, INTEL, KOKKOS,
OPENMP, and OPT packages, respectively.  They are only enabled if
LAMMPS was built with those packages.  See the {\hyperref[\detokenize{Build_package::doc}]{\sphinxcrossref{\DUrole{doc}{Build package}}}} page for more info.

\sphinxAtStartPar
You can specify the accelerated styles explicitly in your input script
by including their suffix, or you can use the {\hyperref[\detokenize{Run_options::doc}]{\sphinxcrossref{\DUrole{doc}{\sphinxhyphen{}suffix command\sphinxhyphen{}line switch}}}} when you invoke LAMMPS, or you can use the
{\hyperref[\detokenize{suffix::doc}]{\sphinxcrossref{\DUrole{doc}{suffix}}}} command in your input script.

\sphinxAtStartPar
See the {\hyperref[\detokenize{Speed_packages::doc}]{\sphinxcrossref{\DUrole{doc}{Accelerator packages}}}} page for more
instructions on how to use the accelerated styles effectively.


\bigskip\hrule\bigskip



\subsection{Mixing, shift, table, tail correction, restart, rRESPA info}
\label{\detokenize{pair_lj_smooth:mixing-shift-table-tail-correction-restart-rrespa-info}}
\sphinxAtStartPar
For atom type pairs I,J and I != J, the epsilon, sigma, Rin
coefficients and the cutoff distance for this pair style can be mixed.
Rin is a cutoff value and is mixed like the cutoff.  The other
coefficients are mixed according to the pair\_modify mix option.  The
default mix value is \sphinxstyleemphasis{geometric}.  See the “pair\_modify” command for
details.

\sphinxAtStartPar
This pair style supports the {\hyperref[\detokenize{pair_modify::doc}]{\sphinxcrossref{\DUrole{doc}{pair\_modify}}}} shift
option for the energy of the pair interaction.

\sphinxAtStartPar
The {\hyperref[\detokenize{pair_modify::doc}]{\sphinxcrossref{\DUrole{doc}{pair\_modify}}}} table option is not relevant
for this pair style.

\sphinxAtStartPar
This pair style does not support the {\hyperref[\detokenize{pair_modify::doc}]{\sphinxcrossref{\DUrole{doc}{pair\_modify}}}}
tail option for adding long\sphinxhyphen{}range tail corrections to energy and
pressure, since the energy of the pair interaction is smoothed to 0.0
at the cutoff.

\sphinxAtStartPar
This pair style writes its information to {\hyperref[\detokenize{restart::doc}]{\sphinxcrossref{\DUrole{doc}{binary restart files}}}}, so pair\_style and pair\_coeff commands do not need
to be specified in an input script that reads a restart file.

\sphinxAtStartPar
This pair style can only be used via the \sphinxstyleemphasis{pair} keyword of the
{\hyperref[\detokenize{run_style::doc}]{\sphinxcrossref{\DUrole{doc}{run\_style respa}}}} command.  It does not support the
\sphinxstyleemphasis{inner}, \sphinxstyleemphasis{middle}, \sphinxstyleemphasis{outer} keywords.


\bigskip\hrule\bigskip



\subsection{Restrictions}
\label{\detokenize{pair_lj_smooth:restrictions}}
\sphinxAtStartPar
This pair style is part of the EXTRA\sphinxhyphen{}PAIR package.  It is only enabled if
LAMMPS was built with that package.  See the
{\hyperref[\detokenize{Build_package::doc}]{\sphinxcrossref{\DUrole{doc}{Build package}}}} page for more info.


\subsection{Related commands}
\label{\detokenize{pair_lj_smooth:related-commands}}
\sphinxAtStartPar
{\hyperref[\detokenize{pair_coeff::doc}]{\sphinxcrossref{\DUrole{doc}{pair\_coeff}}}}, {\hyperref[\detokenize{pair_lj_smooth_linear::doc}]{\sphinxcrossref{\DUrole{doc}{pair lj/smooth/linear}}}}


\subsection{Default}
\label{\detokenize{pair_lj_smooth:default}}
\sphinxAtStartPar
none

\sphinxstepscope

\index{pair\_style lj/smooth/linear@\spxentry{pair\_style lj/smooth/linear}}\ignorespaces 
\index{pair\_style lj/smooth/linear/omp@\spxentry{pair\_style lj/smooth/linear/omp}}\ignorespaces 

\section{pair\_style lj/smooth/linear command}
\label{\detokenize{pair_lj_smooth_linear:pair-style-lj-smooth-linear-command}}\label{\detokenize{pair_lj_smooth_linear:index-1}}\label{\detokenize{pair_lj_smooth_linear:index-0}}\label{\detokenize{pair_lj_smooth_linear::doc}}
\sphinxAtStartPar
Accelerator Variants: \sphinxstyleemphasis{lj/smooth/linear/omp}


\subsection{Syntax}
\label{\detokenize{pair_lj_smooth_linear:syntax}}
\begin{sphinxVerbatim}[commandchars=\\\{\}]
\PYG{k}{pair\PYGZus{}style}\PYG{+w}{ }\PYG{n}{lj}\PYG{o}{/}\PYG{n}{smooth}\PYG{o}{/}\PYG{n}{linear}\PYG{+w}{ }\PYG{n}{cutoff}
\end{sphinxVerbatim}
\begin{itemize}
\item {} 
\sphinxAtStartPar
cutoff = global cutoff for Lennard\sphinxhyphen{}Jones interactions (distance units)

\end{itemize}


\subsection{Examples}
\label{\detokenize{pair_lj_smooth_linear:examples}}
\begin{sphinxVerbatim}[commandchars=\\\{\}]
\PYG{k}{pair\PYGZus{}style}\PYG{+w}{ }\PYG{n}{lj}\PYG{o}{/}\PYG{n}{smooth}\PYG{o}{/}\PYG{n}{linear}\PYG{+w}{ }\PYG{l+m}{2.5}
\PYG{k}{pair\PYGZus{}coeff}\PYG{+w}{ }\PYG{o}{*}\PYG{+w}{ }\PYG{o}{*}\PYG{+w}{ }\PYG{l+m}{1.0}\PYG{+w}{ }\PYG{l+m}{1.0}
\PYG{k}{pair\PYGZus{}coeff}\PYG{+w}{ }\PYG{l+m}{1}\PYG{+w}{ }\PYG{l+m}{1}\PYG{+w}{ }\PYG{l+m}{0.3}\PYG{+w}{ }\PYG{l+m}{3.0}\PYG{+w}{ }\PYG{l+m}{9.0}
\end{sphinxVerbatim}


\subsection{Description}
\label{\detokenize{pair_lj_smooth_linear:description}}
\sphinxAtStartPar
Style \sphinxstyleemphasis{lj/smooth/linear} computes a truncated and force\sphinxhyphen{}shifted LJ
interaction (aka Shifted Force Lennard\sphinxhyphen{}Jones) that combines the
standard 12/6 Lennard\sphinxhyphen{}Jones function and subtracts a linear term based
on the cutoff distance, so that both, the potential and the force, go
continuously to zero at the cutoff \(r_c\) {\hyperref[\detokenize{pair_lj_smooth_linear:toxvaerd}]{\sphinxcrossref{\DUrole{std,std-ref}{(Toxvaerd)}}}}:
\begin{equation*}
\begin{split}\phi\left(r\right) & =  4 \epsilon \left[ \left(\frac{\sigma}{r}\right)^{12} -
                    \left(\frac{\sigma}{r}\right)^6 \right] \\
E\left(r\right) & =  \phi\left(r\right)  - \phi\left(r_c\right) - \left(r - r_c\right) \left.\frac{d\phi}{d r} \right|_{r=r_c}       \qquad r < r_c\end{split}
\end{equation*}
\sphinxAtStartPar
The following coefficients must be defined for each pair of atoms
types via the {\hyperref[\detokenize{pair_coeff::doc}]{\sphinxcrossref{\DUrole{doc}{pair\_coeff}}}} command as in the examples
above, or in the data file or restart files read by the
{\hyperref[\detokenize{read_data::doc}]{\sphinxcrossref{\DUrole{doc}{read\_data}}}} or {\hyperref[\detokenize{read_restart::doc}]{\sphinxcrossref{\DUrole{doc}{read\_restart}}}}
commands, or by mixing as described below:
\begin{itemize}
\item {} 
\sphinxAtStartPar
\(\epsilon\) (energy units)

\item {} 
\sphinxAtStartPar
\(\sigma\) (distance units)

\item {} 
\sphinxAtStartPar
cutoff (distance units)

\end{itemize}

\sphinxAtStartPar
The last coefficient is optional. If not specified, the global
LJ cutoff specified in the pair\_style command is used.


\bigskip\hrule\bigskip


\sphinxAtStartPar
Styles with a \sphinxstyleemphasis{gpu}, \sphinxstyleemphasis{intel}, \sphinxstyleemphasis{kk}, \sphinxstyleemphasis{omp}, or \sphinxstyleemphasis{opt} suffix are
functionally the same as the corresponding style without the suffix.
They have been optimized to run faster, depending on your available
hardware, as discussed on the {\hyperref[\detokenize{Speed_packages::doc}]{\sphinxcrossref{\DUrole{doc}{Accelerator packages}}}}
page.  The accelerated styles take the same arguments and should
produce the same results, except for round\sphinxhyphen{}off and precision issues.

\sphinxAtStartPar
These accelerated styles are part of the GPU, INTEL, KOKKOS,
OPENMP, and OPT packages, respectively.  They are only enabled if
LAMMPS was built with those packages.  See the {\hyperref[\detokenize{Build_package::doc}]{\sphinxcrossref{\DUrole{doc}{Build package}}}} page for more info.

\sphinxAtStartPar
You can specify the accelerated styles explicitly in your input script
by including their suffix, or you can use the {\hyperref[\detokenize{Run_options::doc}]{\sphinxcrossref{\DUrole{doc}{\sphinxhyphen{}suffix command\sphinxhyphen{}line switch}}}} when you invoke LAMMPS, or you can use the
{\hyperref[\detokenize{suffix::doc}]{\sphinxcrossref{\DUrole{doc}{suffix}}}} command in your input script.

\sphinxAtStartPar
See the {\hyperref[\detokenize{Speed_packages::doc}]{\sphinxcrossref{\DUrole{doc}{Accelerator packages}}}} page for more
instructions on how to use the accelerated styles effectively.


\bigskip\hrule\bigskip



\subsection{Mixing, shift, table, tail correction, restart, rRESPA info}
\label{\detokenize{pair_lj_smooth_linear:mixing-shift-table-tail-correction-restart-rrespa-info}}
\sphinxAtStartPar
For atom type pairs I,J and I != J, the epsilon and sigma coefficients
and cutoff distance can be mixed. The default mix value is geometric.
See the “pair\_modify” command for details.

\sphinxAtStartPar
This pair style does not support the {\hyperref[\detokenize{pair_modify::doc}]{\sphinxcrossref{\DUrole{doc}{pair\_modify}}}}
shift option for the energy of the pair interaction, since it goes
to 0.0 at the cutoff by construction.

\sphinxAtStartPar
The {\hyperref[\detokenize{pair_modify::doc}]{\sphinxcrossref{\DUrole{doc}{pair\_modify}}}} table option is not relevant
for this pair style.

\sphinxAtStartPar
This pair style does not support the {\hyperref[\detokenize{pair_modify::doc}]{\sphinxcrossref{\DUrole{doc}{pair\_modify}}}}
tail option for adding long\sphinxhyphen{}range tail corrections to energy and
pressure, since the energy of the pair interaction is smoothed to 0.0
at the cutoff.

\sphinxAtStartPar
This pair style writes its information to {\hyperref[\detokenize{restart::doc}]{\sphinxcrossref{\DUrole{doc}{binary restart files}}}},
so pair\_style and pair\_coeff commands do not need to be specified
in an input script that reads a restart file.

\sphinxAtStartPar
This pair style can only be used via the \sphinxstyleemphasis{pair} keyword of the
{\hyperref[\detokenize{run_style::doc}]{\sphinxcrossref{\DUrole{doc}{run\_style respa}}}} command.  It does not support the
\sphinxstyleemphasis{inner}, \sphinxstyleemphasis{middle}, \sphinxstyleemphasis{outer} keywords.


\bigskip\hrule\bigskip



\subsection{Restrictions}
\label{\detokenize{pair_lj_smooth_linear:restrictions}}
\sphinxAtStartPar
This pair style is part of the EXTRA\sphinxhyphen{}PAIR package.  It is only enabled if
LAMMPS was built with that package.  See the
{\hyperref[\detokenize{Build_package::doc}]{\sphinxcrossref{\DUrole{doc}{Build package}}}} page for more info.


\subsection{Related commands}
\label{\detokenize{pair_lj_smooth_linear:related-commands}}
\sphinxAtStartPar
{\hyperref[\detokenize{pair_coeff::doc}]{\sphinxcrossref{\DUrole{doc}{pair\_coeff}}}}, {\hyperref[\detokenize{pair_lj_smooth::doc}]{\sphinxcrossref{\DUrole{doc}{pair lj/smooth}}}}


\subsection{Default}
\label{\detokenize{pair_lj_smooth_linear:default}}
\sphinxAtStartPar
none


\bigskip\hrule\bigskip

\phantomsection\label{\detokenize{pair_lj_smooth_linear:toxvaerd}}
\sphinxAtStartPar
\sphinxstylestrong{(Toxvaerd)} Toxvaerd, Dyre, J Chem Phys, 134, 081102 (2011).

\sphinxstepscope

\index{pair\_style lj/switch3/coulgauss/long@\spxentry{pair\_style lj/switch3/coulgauss/long}}\ignorespaces 
\index{pair\_style mm3/switch3/coulgauss/long@\spxentry{pair\_style mm3/switch3/coulgauss/long}}\ignorespaces 

\section{pair\_style lj/switch3/coulgauss/long command}
\label{\detokenize{pair_lj_switch3_coulgauss_long:pair-style-lj-switch3-coulgauss-long-command}}\label{\detokenize{pair_lj_switch3_coulgauss_long:index-1}}\label{\detokenize{pair_lj_switch3_coulgauss_long:index-0}}\label{\detokenize{pair_lj_switch3_coulgauss_long::doc}}

\section{pair\_style mm3/switch3/coulgauss/long command}
\label{\detokenize{pair_lj_switch3_coulgauss_long:pair-style-mm3-switch3-coulgauss-long-command}}

\subsection{Syntax}
\label{\detokenize{pair_lj_switch3_coulgauss_long:syntax}}
\begin{sphinxVerbatim}[commandchars=\\\{\}]
\PYG{k}{pair\PYGZus{}style}\PYG{+w}{ }\PYG{n}{style}\PYG{+w}{ }\PYG{n}{args}
\end{sphinxVerbatim}
\begin{itemize}
\item {} 
\sphinxAtStartPar
style = \sphinxstyleemphasis{lj/switch3/coulgauss/long} or \sphinxstyleemphasis{mm3/switch3/coulgauss/long}

\item {} 
\sphinxAtStartPar
args = list of arguments for a particular style

\end{itemize}
\begin{sphinxalltt}
\sphinxstyleemphasis{lj/switch3/coulgauss/long} args = cutoff (cutoff2) width
  cutoff  = global cutoff for LJ (and Coulombic if only 1 arg) (distance units)
  cutoff2 = global cutoff for Coulombic (optional) (distance units)
  width  = width parameter of the smoothing function (distance units)

\sphinxstyleemphasis{mm3/switch3/coulgauss/long} args = cutoff (cutoff2) width
  cutoff  = global cutoff for MM3 (and Coulombic if only 1 arg) (distance units)
  cutoff2 = global cutoff for Coulombic (optional) (distance units)
  width  = width parameter of the smoothing function (distance units)
\end{sphinxalltt}


\subsection{Examples}
\label{\detokenize{pair_lj_switch3_coulgauss_long:examples}}
\begin{sphinxVerbatim}[commandchars=\\\{\}]
\PYG{k}{pair\PYGZus{}style}\PYG{+w}{ }\PYG{n}{lj}\PYG{o}{/}\PYG{n}{switch3}\PYG{o}{/}\PYG{n}{coulgauss}\PYG{o}{/}\PYG{n}{long}\PYG{+w}{    }\PYG{l+m}{12.0}\PYG{+w}{ }\PYG{l+m}{3.0}
\PYG{k}{pair\PYGZus{}coeff}\PYG{+w}{ }\PYG{l+m}{1}\PYG{+w}{  }\PYG{l+m}{0.2}\PYG{+w}{ }\PYG{l+m}{2.5}\PYG{+w}{ }\PYG{l+m}{1.2}

\PYG{k}{pair\PYGZus{}style}\PYG{+w}{ }\PYG{n}{lj}\PYG{o}{/}\PYG{n}{switch3}\PYG{o}{/}\PYG{n}{coulgauss}\PYG{o}{/}\PYG{n}{long}\PYG{+w}{   }\PYG{l+m}{12.0}\PYG{+w}{ }\PYG{l+m}{10.0}\PYG{+w}{ }\PYG{l+m}{3.0}
\PYG{k}{pair\PYGZus{}coeff}\PYG{+w}{ }\PYG{l+m}{1}\PYG{+w}{  }\PYG{l+m}{0.2}\PYG{+w}{ }\PYG{l+m}{2.5}\PYG{+w}{ }\PYG{l+m}{1.2}

\PYG{k}{pair\PYGZus{}style}\PYG{+w}{ }\PYG{n}{mm3}\PYG{o}{/}\PYG{n}{switch3}\PYG{o}{/}\PYG{n}{coulgauss}\PYG{o}{/}\PYG{n}{long}\PYG{+w}{    }\PYG{l+m}{12.0}\PYG{+w}{ }\PYG{l+m}{3.0}
\PYG{k}{pair\PYGZus{}coeff}\PYG{+w}{ }\PYG{l+m}{1}\PYG{+w}{  }\PYG{l+m}{0.2}\PYG{+w}{ }\PYG{l+m}{2.5}\PYG{+w}{ }\PYG{l+m}{1.2}

\PYG{k}{pair\PYGZus{}style}\PYG{+w}{ }\PYG{n}{mm3}\PYG{o}{/}\PYG{n}{switch3}\PYG{o}{/}\PYG{n}{coulgauss}\PYG{o}{/}\PYG{n}{long}\PYG{+w}{   }\PYG{l+m}{12.0}\PYG{+w}{ }\PYG{l+m}{10.0}\PYG{+w}{ }\PYG{l+m}{3.0}
\PYG{k}{pair\PYGZus{}coeff}\PYG{+w}{ }\PYG{l+m}{1}\PYG{+w}{  }\PYG{l+m}{0.2}\PYG{+w}{ }\PYG{l+m}{2.5}\PYG{+w}{ }\PYG{l+m}{1.2}
\end{sphinxVerbatim}


\subsection{Description}
\label{\detokenize{pair_lj_switch3_coulgauss_long:description}}
\sphinxAtStartPar
The \sphinxstyleemphasis{lj/switch3/coulgauss} style evaluates the LJ
vdW potential
\begin{equation*}
\begin{split}E = 4\epsilon \left[ \left(\frac{\sigma}{r}\right)^{12}-\left(\frac{\sigma}{r}\right)^{6} \right]\end{split}
\end{equation*}
\sphinxAtStartPar
The \sphinxstyleemphasis{mm3/switch3/coulgauss/long} style evaluates the MM3
vdW potential {\hyperref[\detokenize{bond_mm3:mm3-allinger1989}]{\sphinxcrossref{\DUrole{std,std-ref}{(Allinger)}}}}
\begin{equation*}
\begin{split}E & = \epsilon_{ij} \left[ -2.25 \left(\frac{r_{v,ij}}{r_{ij}}\right)^6 + 1.84(10)^5 \exp\left[-12.0 r_{ij}/r_{v,ij}\right] \right] S_3(r_{ij}) \\
r_{v,ij} & =  r_{v,i} + r_{v,j} \\
\epsilon_{ij} & = \sqrt{\epsilon_i \epsilon_j}\end{split}
\end{equation*}
\sphinxAtStartPar
Both potentials go smoothly to zero at the cutoff r\_c as defined by the
switching function
\begin{equation*}
\begin{split}S_3(r) = \left\lbrace \begin{array}{ll}
                    1 & \quad\mathrm{if}\quad r < r_\mathrm{c} - w \\
                    3x^2 - 2x^3 & \quad\mathrm{if}\quad r < r_\mathrm{c} \quad\mathrm{with\quad} x=\frac{r_\mathrm{c} - r}{w} \\
                    0 & \quad\mathrm{if}\quad r >= r_\mathrm{c}
                \end{array} \right.\end{split}
\end{equation*}
\sphinxAtStartPar
where w is the width defined in the arguments. This potential
is combined with Coulomb interaction between Gaussian charge densities:
\begin{equation*}
\begin{split}E = \frac{q_i q_j \mathrm{erf}\left( r/\sqrt{\gamma_1^2+\gamma_2^2} \right) }{\epsilon r_{ij}}\end{split}
\end{equation*}
\sphinxAtStartPar
where \(q_i\) and \(q_j\) are the charges on the two atoms,
\(\epsilon\) is the dielectric constant which can be set by the
{\hyperref[\detokenize{dielectric::doc}]{\sphinxcrossref{\DUrole{doc}{dielectric}}}} command, \(\gamma_i\) and
\(\gamma_j\) are the widths of the Gaussian charge distribution and
erf() is the error\sphinxhyphen{}function.  This style has to be used in conjunction
with the {\hyperref[\detokenize{kspace_style::doc}]{\sphinxcrossref{\DUrole{doc}{kspace\_style}}}} command

\sphinxAtStartPar
If one cutoff is specified it is used for both the vdW and Coulomb
terms.  If two cutoffs are specified, the first is used as the cutoff
for the vdW terms, and the second is the cutoff for the Coulombic term.

\sphinxAtStartPar
The following coefficients must be defined for each pair of atoms
types via the {\hyperref[\detokenize{pair_coeff::doc}]{\sphinxcrossref{\DUrole{doc}{pair\_coeff}}}} command as in the examples
above, or in the data file or restart files read by the
{\hyperref[\detokenize{read_data::doc}]{\sphinxcrossref{\DUrole{doc}{read\_data}}}} or {\hyperref[\detokenize{read_restart::doc}]{\sphinxcrossref{\DUrole{doc}{read\_restart}}}}
commands:
\begin{itemize}
\item {} 
\sphinxAtStartPar
\(\epsilon\) (energy)

\item {} 
\sphinxAtStartPar
\(\sigma\) (distance)

\item {} 
\sphinxAtStartPar
\(\gamma\) (distance)

\end{itemize}


\bigskip\hrule\bigskip



\subsection{Mixing, shift, table, tail correction, restart, rRESPA info}
\label{\detokenize{pair_lj_switch3_coulgauss_long:mixing-shift-table-tail-correction-restart-rrespa-info}}
\sphinxAtStartPar
For atom type pairs I,J and I != J, the epsilon and sigma coefficients
and cutoff distance for all of the lj/long pair styles can be mixed.
The default mix value is \sphinxstyleemphasis{geometric}.  See the “pair\_modify” command
for details.

\sphinxAtStartPar
Shifting the potential energy is not necessary because the switching
function ensures that the potential is zero at the cut\sphinxhyphen{}off.

\sphinxAtStartPar
These pair styles support the {\hyperref[\detokenize{pair_modify::doc}]{\sphinxcrossref{\DUrole{doc}{pair\_modify}}}} table and
options since they can tabulate the short\sphinxhyphen{}range portion of the
long\sphinxhyphen{}range Coulombic interactions.

\sphinxAtStartPar
Thes pair styles do not support the {\hyperref[\detokenize{pair_modify::doc}]{\sphinxcrossref{\DUrole{doc}{pair\_modify}}}}
tail option for adding a long\sphinxhyphen{}range tail correction to the
Lennard\sphinxhyphen{}Jones portion of the energy and pressure.

\sphinxAtStartPar
These pair styles write their information to {\hyperref[\detokenize{restart::doc}]{\sphinxcrossref{\DUrole{doc}{binary restart files}}}}, so pair\_style and pair\_coeff commands do not need
to be specified in an input script that reads a restart file.

\sphinxAtStartPar
These pair styles can only be used via the \sphinxstyleemphasis{pair} keyword of the
{\hyperref[\detokenize{run_style::doc}]{\sphinxcrossref{\DUrole{doc}{run\_style respa}}}} command.  They do not support the
\sphinxstyleemphasis{inner}, \sphinxstyleemphasis{middle}, \sphinxstyleemphasis{outer} keywords.


\subsection{Restrictions}
\label{\detokenize{pair_lj_switch3_coulgauss_long:restrictions}}
\sphinxAtStartPar
These styles are part of the YAFF package.  They are only enabled
if LAMMPS was built with that package.  See the {\hyperref[\detokenize{Build_package::doc}]{\sphinxcrossref{\DUrole{doc}{Build package}}}} page for more info.


\subsection{Related commands}
\label{\detokenize{pair_lj_switch3_coulgauss_long:related-commands}}
\sphinxAtStartPar
{\hyperref[\detokenize{pair_coeff::doc}]{\sphinxcrossref{\DUrole{doc}{pair\_coeff}}}}


\subsection{Default}
\label{\detokenize{pair_lj_switch3_coulgauss_long:default}}
\sphinxAtStartPar
none

\sphinxstepscope

\index{pair\_style local/density@\spxentry{pair\_style local/density}}\ignorespaces 

\section{pair\_style local/density command}
\label{\detokenize{pair_local_density:pair-style-local-density-command}}\label{\detokenize{pair_local_density:index-0}}\label{\detokenize{pair_local_density::doc}}

\subsection{Syntax}
\label{\detokenize{pair_local_density:syntax}}
\begin{sphinxVerbatim}[commandchars=\\\{\}]
\PYG{k}{pair\PYGZus{}style}\PYG{+w}{ }\PYG{n}{style}\PYG{+w}{ }\PYG{n}{arg}
\end{sphinxVerbatim}
\begin{itemize}
\item {} 
\sphinxAtStartPar
style = \sphinxstyleemphasis{local/density}

\item {} 
\sphinxAtStartPar
arg = name of file containing tabulated values of local density and the potential

\end{itemize}


\subsection{Examples}
\label{\detokenize{pair_local_density:examples}}
\begin{sphinxVerbatim}[commandchars=\\\{\}]
\PYG{k}{pair\PYGZus{}style}\PYG{+w}{ }\PYG{n}{local}\PYG{o}{/}\PYG{n}{density}\PYG{+w}{ }\PYG{n}{benzene\PYGZus{}water.localdensity.table}

\PYG{k}{pair\PYGZus{}style}\PYG{+w}{ }\PYG{n}{hybrid}\PYG{o}{/}\PYG{n}{overlay}\PYG{+w}{ }\PYG{n}{table}\PYG{+w}{ }\PYG{n}{spline}\PYG{+w}{ }\PYG{l+m}{500}\PYG{+w}{ }\PYG{n}{local}\PYG{o}{/}\PYG{n}{density}
\PYG{k}{pair\PYGZus{}coeff}\PYG{+w}{ }\PYG{o}{*}\PYG{+w}{ }\PYG{o}{*}\PYG{+w}{ }\PYG{n}{local}\PYG{o}{/}\PYG{n}{density}\PYG{+w}{  }\PYG{n}{benzene\PYGZus{}water.localdensity.table}
\end{sphinxVerbatim}


\subsection{Description}
\label{\detokenize{pair_local_density:description}}
\sphinxAtStartPar
The local density (LD) potential is a mean\sphinxhyphen{}field manybody potential,
and, in some way, a generalization of embedded atom models (EAM).  The
name “local density potential” arises from the fact that it assigns an
energy to an atom depending on the number of neighboring atoms of a
given type around it within a predefined spherical volume (i.e., within
the cutoff).  The bottom\sphinxhyphen{}up coarse\sphinxhyphen{}graining (CG) literature suggests
that such potentials can be widely useful in capturing effective
multibody forces in a computationally efficient manner and thus improve
the quality of CG models of implicit solvation {\hyperref[\detokenize{pair_local_density:sanyal1}]{\sphinxcrossref{\DUrole{std,std-ref}{(Sanyal1)}}}} and phase\sphinxhyphen{}segregation in liquid mixtures {\hyperref[\detokenize{pair_local_density:sanyal2}]{\sphinxcrossref{\DUrole{std,std-ref}{(Sanyal2)}}}}, and provide guidelines to determine the extent of manybody
correlations present in a CG model {\hyperref[\detokenize{pair_local_density:rosenberger}]{\sphinxcrossref{\DUrole{std,std-ref}{(Rosenberger)}}}}.
The LD potential in LAMMPS is primarily intended to be used as a
corrective potential over traditional pair potentials in bottom\sphinxhyphen{}up CG
models via {\hyperref[\detokenize{pair_hybrid::doc}]{\sphinxcrossref{\DUrole{doc}{hybrid/overlay pair style}}}} with other
explicit pair interaction terms (e.g., tabulated, Lennard\sphinxhyphen{}Jones, Morse
etc.).  Because the LD potential is not a pair potential per se, it is
implemented simply as a single auxiliary file with all specifications
that will be read upon initialization.

\begin{sphinxadmonition}{note}{Note:}
\sphinxAtStartPar
Thus when used as the only interaction in the system, there is no
corresponding pair\_coeff command and when used with other pair styles using the
hybrid/overlay option, the corresponding pair\_coeff command must be supplied
*  * as placeholders for the atom types.
\end{sphinxadmonition}


\bigskip\hrule\bigskip


\sphinxAtStartPar
\sphinxstylestrong{System with a single CG atom type:}

\sphinxAtStartPar
A system of a single atom type (e.g., LJ argon) with a single local density (LD)
potential would have an energy given by:
\begin{equation*}
\begin{split}U_{LD} = \sum_i F(\rho_i)\end{split}
\end{equation*}
\sphinxAtStartPar
where \(\rho_i\) is the LD at atom \sphinxstyleemphasis{i} and \(F(\rho)\) is
similar in spirit to the embedding function used in EAM potentials. The
LD at atom \sphinxstyleemphasis{i} is given by the sum
\begin{equation*}
\begin{split}\rho_i = \sum_{j \neq i} \varphi(r_{ij})\end{split}
\end{equation*}
\sphinxAtStartPar
where \(\varphi\) is an indicator function that is one at r=0 and
zero beyond a cutoff distance R2. The choice of the functional form of
\(\varphi\) is somewhat arbitrary, but the following piecewise cubic
function has proven sufficiently general: {\hyperref[\detokenize{pair_local_density:sanyal1}]{\sphinxcrossref{\DUrole{std,std-ref}{(Sanyal1)}}}},
{\hyperref[\detokenize{pair_local_density:sanyal2}]{\sphinxcrossref{\DUrole{std,std-ref}{(Sanyal2)}}}} {\hyperref[\detokenize{pair_local_density:rosenberger}]{\sphinxcrossref{\DUrole{std,std-ref}{(Rosenberger)}}}}
\begin{equation*}
\begin{split}\varphi(r) =
\begin{cases}
1 & r \le R_1 \\
c_0 + c_2r^2 + c_4r^4 + c_6r^6  & r \in (R_1, R_2) \\
0 & r \ge R_2
\end{cases}\end{split}
\end{equation*}
\sphinxAtStartPar
The constants \sphinxstyleemphasis{c} are chosen so that the indicator function smoothly
interpolates between 1 and 0 between the distances R1 and R2, which are
called the inner and outer cutoffs, respectively. Thus phi satisfies
phi(R1) = 1, phi(R2) = dphi/dr @ (r=R1) =  dphi/dr @ (r=R2) = 0. The embedding
function F(rho) may or may not have a closed\sphinxhyphen{}form expression. To maintain
generality, it is practically represented with a spline\sphinxhyphen{}interpolated table
over a predetermined range of rho. Outside of that range it simply adopts zero
values at the endpoints.

\sphinxAtStartPar
It can be shown that the total force between two atoms due to the LD potential
takes the form of a pair force, which motivates its designation as a LAMMPS
pair style. Please see {\hyperref[\detokenize{pair_local_density:sanyal1}]{\sphinxcrossref{\DUrole{std,std-ref}{(Sanyal1)}}}} for details of the derivation.


\bigskip\hrule\bigskip


\sphinxAtStartPar
\sphinxstylestrong{Systems with arbitrary numbers of atom types:}

\sphinxAtStartPar
The potential is easily generalized to systems involving multiple atom types:
\begin{equation*}
\begin{split}U_{LD} = \sum_i a_\alpha F(\rho_i)\end{split}
\end{equation*}
\sphinxAtStartPar
with the LD expressed as
\begin{equation*}
\begin{split}\rho_i = \sum_{j \neq i} b_\beta \varphi(r_{ij})\end{split}
\end{equation*}
\sphinxAtStartPar
where \(\alpha\) gives the type of atom \sphinxstyleemphasis{i}, \(\beta\) the
type of atom \sphinxstyleemphasis{j}, and the coefficients \sphinxstyleemphasis{a} and \sphinxstyleemphasis{b} filter for atom
types as specified by the user. \sphinxstyleemphasis{a} is called the central atom filter as
it determines to which atoms the potential applies; \(a_{\alpha} =
1\) if the LD potential applies to atom type \(\alpha\) else zero. On the
other hand, \sphinxstyleemphasis{b} is called the neighbor atom filter because it specifies
which atom types to use in the calculation of the LD; \(b_{\beta} =
1\) if atom type \(\beta\) contributes to the LD and zero otherwise.

\begin{sphinxadmonition}{note}{Note:}
\sphinxAtStartPar
Note that the potentials need not be symmetric with respect to atom
types, which is the reason for two distinct sets of coefficients \sphinxstyleemphasis{a}
and \sphinxstyleemphasis{b}. An atom type may contribute to the LD but not the
potential, or to the potential but not the LD. Such decisions are
made by the user and should (ideally) be motivated on physical
grounds for the problem at hand.
\end{sphinxadmonition}


\bigskip\hrule\bigskip


\sphinxAtStartPar
\sphinxstylestrong{General form for implementation in LAMMPS:}

\sphinxAtStartPar
Of course, a system with many atom types may have many different possible LD
potentials, each with their own atom type filters, cutoffs, and embedding
functions. The most general form of this potential as implemented in the
pair\_style local/density is:
\begin{equation*}
\begin{split}U_{LD} = \sum_k U_{LD}^{(k)} = \sum_i \left[ \sum_k a_\alpha^{(k)} F^{(k)} \left(\rho_i^{(k)}\right) \right]\end{split}
\end{equation*}
\sphinxAtStartPar
where, \sphinxstyleemphasis{k} is an index that spans the (arbitrary) number of applied LD
potentials N\_LD. Each LD is calculated as before with:
\begin{equation*}
\begin{split}\rho_i^{(k)} = \sum_j b_\beta^{(k)} \varphi^{(k)} (r_{ij})\end{split}
\end{equation*}
\sphinxAtStartPar
The superscript on the indicator function phi simply indicates that it is
associated with specific values of the cutoff distances R1(k) and R2(k). In
summary, there may be N\_LD distinct LD potentials. With each potential type (k),
one must specify:
\begin{itemize}
\item {} 
\sphinxAtStartPar
the inner and outer cutoffs as R1 and R2

\item {} 
\sphinxAtStartPar
the central type filter a(k), where k = 1,2,…N\_LD

\item {} 
\sphinxAtStartPar
the neighbor type filter b(k), where k = 1,2,…N\_LD

\item {} 
\sphinxAtStartPar
the LD potential function F(k)(rho), typically as a table that is later spline\sphinxhyphen{}interpolated

\end{itemize}


\bigskip\hrule\bigskip


\sphinxAtStartPar
\sphinxstylestrong{Tabulated input file format:}

\begin{sphinxVerbatim}[commandchars=\\\{\}]
Line 1:             comment or blank (ignored)
Line 2:             comment or blank (ignored)
Line 3:             N\PYGZus{}LD N\PYGZus{}rho (\PYGZsh{} of LD potentials and \PYGZsh{} of tabulated values, single space separated)
Line 4:             blank (ignored)
Line 5:             R1(k) R2(k) (lower and upper cutoffs, single space separated)
Line 6:             central\PYGZhy{}types (central atom types, single space separated)
Line 7:             neighbor\PYGZhy{}types (neighbor atom types single space separated)
Line 8:             rho\PYGZus{}min rho\PYGZus{}max drho (min, max and diff. in tabulated rho values, single space separated)
Line 9:             F(k)(rho\PYGZus{}min + 0.drho)
Line 10:            F(k)(rho\PYGZus{}min + 1.drho)
Line 11:            F(k)(rho\PYGZus{}min + 2.drho)
...
Line 9+N\PYGZus{}rho:       F(k)(rho\PYGZus{}min + N\PYGZus{}rho . drho)
Line 10+N\PYGZus{}rho:      blank (ignored)

Block 2

Block 3

Block N\PYGZus{}LD
\end{sphinxVerbatim}

\sphinxAtStartPar
Lines 5 to 9+N\_rho constitute the first block. Thus the input file is separated
(by blank lines) into N\_LD blocks each representing a separate LD potential and
each specifying its own upper and lower cutoffs, central and neighbor atoms,
and potential.  In general, blank lines anywhere are ignored.


\bigskip\hrule\bigskip



\subsection{Mixing, shift, table, tail correction, restart, rRESPA info}
\label{\detokenize{pair_local_density:mixing-shift-table-tail-correction-restart-rrespa-info}}
\sphinxAtStartPar
This pair style does not support automatic mixing. For atom type pairs
\(\alpha\), \(\beta\) and \(\alpha\) != \(\beta\), even
if LD potentials of type (\(\alpha\), \(\alpha\)) and
(\(\beta\), \(\beta\)) are provided, you will need to explicitly
provide LD potential types (\(\alpha\), \(\beta\)) and
(\(\beta\), \(\alpha\)) if need be (Here, the notation
(\(\alpha\), \(\beta\)) means that \(\alpha\) is the central
atom to which the LD potential is applied and \(\beta\) is the
neighbor atom which contributes to the LD potential on \(\alpha\)).

\sphinxAtStartPar
This pair style does not support the {\hyperref[\detokenize{pair_modify::doc}]{\sphinxcrossref{\DUrole{doc}{pair\_modify}}}}
shift, table, and tail options.

\sphinxAtStartPar
The local/density pair style does not write its information to {\hyperref[\detokenize{restart::doc}]{\sphinxcrossref{\DUrole{doc}{binary restart files}}}}, since it is stored in tabulated potential files.
Thus, you need to re\sphinxhyphen{}specify the pair\_style and pair\_coeff commands in
an input script that reads a restart file.


\bigskip\hrule\bigskip



\subsection{Restrictions}
\label{\detokenize{pair_local_density:restrictions}}
\sphinxAtStartPar
The local/density pair style is a part of the MANYBODY package. It is
only enabled if LAMMPS was built with that package.  See the
{\hyperref[\detokenize{Build_package::doc}]{\sphinxcrossref{\DUrole{doc}{Build package}}}} page for more info.


\subsection{Related commands}
\label{\detokenize{pair_local_density:related-commands}}
\sphinxAtStartPar
{\hyperref[\detokenize{pair_coeff::doc}]{\sphinxcrossref{\DUrole{doc}{pair\_coeff}}}}


\subsection{Default}
\label{\detokenize{pair_local_density:default}}
\sphinxAtStartPar
none


\bigskip\hrule\bigskip

\phantomsection\label{\detokenize{pair_local_density:sanyal1}}\phantomsection\label{\detokenize{pair_local_density:sanyal2}}
\sphinxAtStartPar
\sphinxstylestrong{(Sanyal1)} Sanyal and Shell, Journal of Chemical Physics, 2016, 145 (3), 034109.

\sphinxAtStartPar
\sphinxstylestrong{(Sanyal2)} Sanyal and Shell, Journal of Physical Chemistry B, 122 (21), 5678\sphinxhyphen{}5693.

\phantomsection\label{\detokenize{pair_local_density:rosenberger}}
\sphinxAtStartPar
\sphinxstylestrong{(Rosenberger)} Rosenberger, Sanyal, Shell and van der Vegt,  Journal of Chemical Physics, 2019, 151 (4), 044111.

\sphinxstepscope

\index{pair\_style lubricate@\spxentry{pair\_style lubricate}}\ignorespaces 
\index{pair\_style lubricate/omp@\spxentry{pair\_style lubricate/omp}}\ignorespaces 
\index{pair\_style lubricate/poly@\spxentry{pair\_style lubricate/poly}}\ignorespaces 
\index{pair\_style lubricate/poly/omp@\spxentry{pair\_style lubricate/poly/omp}}\ignorespaces 

\section{pair\_style lubricate command}
\label{\detokenize{pair_lubricate:pair-style-lubricate-command}}\label{\detokenize{pair_lubricate:index-3}}\label{\detokenize{pair_lubricate:index-2}}\label{\detokenize{pair_lubricate:index-1}}\label{\detokenize{pair_lubricate:index-0}}\label{\detokenize{pair_lubricate::doc}}
\sphinxAtStartPar
Accelerator Variants: \sphinxstyleemphasis{lubricate/omp}


\section{pair\_style lubricate/poly command}
\label{\detokenize{pair_lubricate:pair-style-lubricate-poly-command}}
\sphinxAtStartPar
Accelerator Variants: \sphinxstyleemphasis{lubricate/poly/omp}


\subsection{Syntax}
\label{\detokenize{pair_lubricate:syntax}}
\begin{sphinxVerbatim}[commandchars=\\\{\}]
\PYG{k}{pair\PYGZus{}style}\PYG{+w}{ }\PYG{n}{style}\PYG{+w}{ }\PYG{n}{mu}\PYG{+w}{ }\PYG{n}{flaglog}\PYG{+w}{ }\PYG{n}{flagfld}\PYG{+w}{ }\PYG{n}{cutinner}\PYG{+w}{ }\PYG{n}{cutoff}\PYG{+w}{ }\PYG{n}{flagHI}\PYG{+w}{ }\PYG{n}{flagVF}
\end{sphinxVerbatim}
\begin{itemize}
\item {} 
\sphinxAtStartPar
style = \sphinxstyleemphasis{lubricate} or \sphinxstyleemphasis{lubricate/poly}

\item {} 
\sphinxAtStartPar
mu = dynamic viscosity (dynamic viscosity units)

\item {} 
\sphinxAtStartPar
flaglog = 0/1 to exclude/include log terms in the lubrication approximation

\item {} 
\sphinxAtStartPar
flagfld = 0/1 to exclude/include Fast Lubrication Dynamics (FLD) effects

\item {} 
\sphinxAtStartPar
cutinner = inner cutoff distance (distance units)

\item {} 
\sphinxAtStartPar
cutoff = outer cutoff for interactions (distance units)

\item {} 
\sphinxAtStartPar
flagHI (optional) = 0/1 to exclude/include 1/r hydrodynamic interactions

\item {} 
\sphinxAtStartPar
flagVF (optional) = 0/1 to exclude/include volume fraction corrections in the long\sphinxhyphen{}range isotropic terms

\end{itemize}


\subsection{Examples}
\label{\detokenize{pair_lubricate:examples}}
\sphinxAtStartPar
(all assume radius = 1)

\begin{sphinxVerbatim}[commandchars=\\\{\}]
\PYG{k}{pair\PYGZus{}style}\PYG{+w}{ }\PYG{n}{lubricate}\PYG{+w}{ }\PYG{l+m}{1.5}\PYG{+w}{ }\PYG{l+m}{1}\PYG{+w}{ }\PYG{l+m}{1}\PYG{+w}{ }\PYG{l+m}{2.01}\PYG{+w}{ }\PYG{l+m}{2.5}
\PYG{k}{pair\PYGZus{}coeff}\PYG{+w}{ }\PYG{l+m}{1}\PYG{+w}{ }\PYG{l+m}{1}\PYG{+w}{ }\PYG{l+m}{2.05}\PYG{+w}{ }\PYG{l+m}{2.8}
\PYG{k}{pair\PYGZus{}coeff}\PYG{+w}{ }\PYG{o}{*}\PYG{+w}{ }\PYG{o}{*}

\PYG{k}{pair\PYGZus{}style}\PYG{+w}{ }\PYG{n}{lubricate}\PYG{+w}{ }\PYG{l+m}{1.5}\PYG{+w}{ }\PYG{l+m}{1}\PYG{+w}{ }\PYG{l+m}{1}\PYG{+w}{ }\PYG{l+m}{2.01}\PYG{+w}{ }\PYG{l+m}{2.5}
\PYG{k}{pair\PYGZus{}coeff}\PYG{+w}{ }\PYG{o}{*}\PYG{+w}{ }\PYG{o}{*}
\PYG{k}{variable }\PYG{n+nv+nvIdentifier}{mu}\PYG{+w}{ }\PYG{n}{equal}\PYG{+w}{ }\PYG{n}{ramp}\PYG{n+nv}{(}\PYG{n+nv}{1,2}\PYG{n+nv}{)}
\PYG{k}{fix }\PYG{n+nv+nvIdentifier}{1}\PYG{+w}{ }\PYG{n+nv+nvIdentifier}{all}\PYG{+w}{ }\PYG{n}{adapt}\PYG{+w}{ }\PYG{l+m}{1}\PYG{+w}{ }\PYG{n}{pair}\PYG{+w}{ }\PYG{n}{lubricate}\PYG{+w}{ }\PYG{n}{mu}\PYG{+w}{ }\PYG{o}{*}\PYG{+w}{ }\PYG{o}{*}\PYG{+w}{ }\PYG{n}{v\PYGZus{}mu}
\end{sphinxVerbatim}


\subsection{Description}
\label{\detokenize{pair_lubricate:description}}
\sphinxAtStartPar
Styles \sphinxstyleemphasis{lubricate} and \sphinxstyleemphasis{lubricate/poly} compute hydrodynamic
interactions between mono\sphinxhyphen{}disperse finite\sphinxhyphen{}size spherical particles in
a pairwise fashion.  The interactions have 2 components.  The first is
Ball\sphinxhyphen{}Melrose lubrication terms via the formulas in {\hyperref[\detokenize{pair_lubricate:ball1}]{\sphinxcrossref{\DUrole{std,std-ref}{(Ball and Melrose)}}}}
\begin{equation*}
\begin{split}W & =  - a_{sq} | (v_1 - v_2) \bullet \mathbf{nn} |^2 -
a_{sh} | (\omega_1 + \omega_2) \bullet
(\mathbf{I} - \mathbf{nn}) - 2 \Omega_N |^2 - \\
&  a_{pu} | (\omega_1 - \omega_2) \bullet (\mathbf{I} - \mathbf{nn}) |^2 -
a_{tw} | (\omega_1 - \omega_2) \bullet \mathbf{nn} |^2  \qquad r < r_c \\
& \\
\Omega_N & = \mathbf{n} \times (v_1 - v_2) / r\end{split}
\end{equation*}
\sphinxAtStartPar
which represents the dissipation W between two nearby particles due to
their relative velocities in the presence of a background solvent with
viscosity \sphinxstyleemphasis{mu}.  Note that this is dynamic viscosity which has units of
mass/distance/time, not kinematic viscosity.

\sphinxAtStartPar
The Asq (squeeze) term is the strongest and is included if \sphinxstyleemphasis{flagHI} is
set to 1 (default). It scales as 1/gap where gap is the separation
between the surfaces of the 2 particles. The Ash (shear) and Apu
(pump) terms are only included if \sphinxstyleemphasis{flaglog} is set to 1. They are the
next strongest interactions, and the only other singular interaction,
and scale as log(gap). Note that \sphinxstyleemphasis{flaglog} = 1 and \sphinxstyleemphasis{flagHI} = 0 is
invalid, and will result in a warning message, after which \sphinxstyleemphasis{flagHI} will
be set to 1. The Atw (twist) term is currently not included. It is
typically a very small contribution to the lubrication forces.

\sphinxAtStartPar
The \sphinxstyleemphasis{flagHI} and \sphinxstyleemphasis{flagVF} settings are optional.  Neither should be
used, or both must be defined.

\sphinxAtStartPar
\sphinxstyleemphasis{Cutinner} sets the minimum center\sphinxhyphen{}to\sphinxhyphen{}center separation that will be
used in calculations irrespective of the actual separation.  \sphinxstyleemphasis{Cutoff}
is the maximum center\sphinxhyphen{}to\sphinxhyphen{}center separation at which an interaction is
computed.  Using a \sphinxstyleemphasis{cutoff} less than 3 radii is recommended if
\sphinxstyleemphasis{flaglog} is set to 1.

\sphinxAtStartPar
The other component is due to the Fast Lubrication Dynamics (FLD)
approximation, described in {\hyperref[\detokenize{pair_lubricate:kumar1}]{\sphinxcrossref{\DUrole{std,std-ref}{(Kumar)}}}}, which can be
represented by the following equation
\begin{equation*}
\begin{split}F^{H} = -R_{FU}(U-U^{\infty}) + R_{FE}E^{\infty}\end{split}
\end{equation*}
\sphinxAtStartPar
where U represents the velocities and angular velocities of the
particles, \(U^{\infty}\) represents the velocity and the angular velocity
of the undisturbed fluid, and \(E^{\infty}\) represents the rate of strain
tensor of the undisturbed fluid with viscosity \sphinxstyleemphasis{mu}. Again, note that
this is dynamic viscosity which has units of mass/distance/time, not
kinematic viscosity. Volume fraction corrections to R\_FU are included
as long as \sphinxstyleemphasis{flagVF} is set to 1 (default).

\begin{sphinxadmonition}{note}{Note:}
\sphinxAtStartPar
When using the FLD terms, these pair styles are designed to be used
with explicit time integration and a correspondingly small timestep.
Thus either {\hyperref[\detokenize{fix_nve_sphere::doc}]{\sphinxcrossref{\DUrole{doc}{fix nve/sphere}}}} or {\hyperref[\detokenize{fix_nve_asphere::doc}]{\sphinxcrossref{\DUrole{doc}{fix
nve/asphere}}}} should be used for time integration.
To perform implicit FLD, see the {\hyperref[\detokenize{pair_lubricateU::doc}]{\sphinxcrossref{\DUrole{doc}{pair\_style lubricateU}}}} command.
\end{sphinxadmonition}

\sphinxAtStartPar
Style \sphinxstyleemphasis{lubricate} requires monodisperse spherical particles; style
\sphinxstyleemphasis{lubricate/poly} allows for polydisperse spherical particles.

\sphinxAtStartPar
The viscosity \sphinxstyleemphasis{mu} can be varied in a time\sphinxhyphen{}dependent manner over the
course of a simulation, in which case in which case the pair\_style
setting for \sphinxstyleemphasis{mu} will be overridden.  See the {\hyperref[\detokenize{fix_adapt::doc}]{\sphinxcrossref{\DUrole{doc}{fix adapt}}}}
command for details.

\sphinxAtStartPar
If the suspension is sheared via the {\hyperref[\detokenize{fix_deform::doc}]{\sphinxcrossref{\DUrole{doc}{fix deform}}}}
command then the pair style uses the shear rate to adjust the
hydrodynamic interactions accordingly. Volume changes due to fix
deform are accounted for when computing the volume fraction
corrections to R\_FU.

\sphinxAtStartPar
When computing the volume fraction corrections to R\_FU, the presence of
walls (whether moving or stationary) will affect the volume fraction
available to colloidal particles. This is currently accounted for with
the following types of walls: {\hyperref[\detokenize{fix_wall::doc}]{\sphinxcrossref{\DUrole{doc}{wall/lj93}}}},
{\hyperref[\detokenize{fix_wall::doc}]{\sphinxcrossref{\DUrole{doc}{wall/lj126}}}}, {\hyperref[\detokenize{fix_wall::doc}]{\sphinxcrossref{\DUrole{doc}{wall/colloid}}}}, and
{\hyperref[\detokenize{fix_wall::doc}]{\sphinxcrossref{\DUrole{doc}{wall/harmonic}}}}.  For these wall styles, the correct
volume fraction will be used when walls do not coincide with the box
boundary, as well as when walls move and thereby cause a change in the
volume fraction.  Other wall styles may still work, but they will result
in the volume fraction being computed based on the box boundaries.
Several wall styles are not compatible with these pair styles and using
them will result in an error.

\sphinxAtStartPar
Since lubrication forces are dissipative, it is usually desirable to
thermostat the system at a constant temperature. If Brownian motion
(at a constant temperature) is desired, it can be set using the
{\hyperref[\detokenize{pair_brownian::doc}]{\sphinxcrossref{\DUrole{doc}{pair\_style brownian}}}} command. These pair styles
and the brownian style should use consistent parameters for \sphinxstyleemphasis{mu},
\sphinxstyleemphasis{flaglog}, \sphinxstyleemphasis{flagfld}, \sphinxstyleemphasis{cutinner}, \sphinxstyleemphasis{cutoff}, \sphinxstyleemphasis{flagHI} and \sphinxstyleemphasis{flagVF}.


\bigskip\hrule\bigskip


\sphinxAtStartPar
The following coefficients must be defined for each pair of atoms
types via the {\hyperref[\detokenize{pair_coeff::doc}]{\sphinxcrossref{\DUrole{doc}{pair\_coeff}}}} command as in the examples
above, or in the data file or restart files read by the
{\hyperref[\detokenize{read_data::doc}]{\sphinxcrossref{\DUrole{doc}{read\_data}}}} or {\hyperref[\detokenize{read_restart::doc}]{\sphinxcrossref{\DUrole{doc}{read\_restart}}}}
commands, or by mixing as described below:
\begin{itemize}
\item {} 
\sphinxAtStartPar
cutinner (distance units)

\item {} 
\sphinxAtStartPar
cutoff (distance units)

\end{itemize}

\sphinxAtStartPar
The two coefficients are optional.  If neither is specified, the two
cutoffs specified in the pair\_style command are used.  Otherwise both
must be specified.


\bigskip\hrule\bigskip


\sphinxAtStartPar
Styles with a \sphinxstyleemphasis{gpu}, \sphinxstyleemphasis{intel}, \sphinxstyleemphasis{kk}, \sphinxstyleemphasis{omp}, or \sphinxstyleemphasis{opt} suffix are
functionally the same as the corresponding style without the suffix.
They have been optimized to run faster, depending on your available
hardware, as discussed on the {\hyperref[\detokenize{Speed_packages::doc}]{\sphinxcrossref{\DUrole{doc}{Accelerator packages}}}}
page.  The accelerated styles take the same arguments and should
produce the same results, except for round\sphinxhyphen{}off and precision issues.

\sphinxAtStartPar
These accelerated styles are part of the GPU, INTEL, KOKKOS,
OPENMP, and OPT packages, respectively.  They are only enabled if
LAMMPS was built with those packages.  See the {\hyperref[\detokenize{Build_package::doc}]{\sphinxcrossref{\DUrole{doc}{Build package}}}} page for more info.

\sphinxAtStartPar
You can specify the accelerated styles explicitly in your input script
by including their suffix, or you can use the {\hyperref[\detokenize{Run_options::doc}]{\sphinxcrossref{\DUrole{doc}{\sphinxhyphen{}suffix command\sphinxhyphen{}line switch}}}} when you invoke LAMMPS, or you can use the
{\hyperref[\detokenize{suffix::doc}]{\sphinxcrossref{\DUrole{doc}{suffix}}}} command in your input script.

\sphinxAtStartPar
See the {\hyperref[\detokenize{Speed_packages::doc}]{\sphinxcrossref{\DUrole{doc}{Accelerator packages}}}} page for more
instructions on how to use the accelerated styles effectively.


\bigskip\hrule\bigskip



\subsection{Mixing, shift, table, tail correction, restart, rRESPA info}
\label{\detokenize{pair_lubricate:mixing-shift-table-tail-correction-restart-rrespa-info}}
\sphinxAtStartPar
For atom type pairs I,J and I != J, the two cutoff distances for this
pair style can be mixed.  The default mix value is \sphinxstyleemphasis{geometric}.  See
the “pair\_modify” command for details.

\sphinxAtStartPar
This pair style does not support the {\hyperref[\detokenize{pair_modify::doc}]{\sphinxcrossref{\DUrole{doc}{pair\_modify}}}}
shift option for the energy of the pair interaction.

\sphinxAtStartPar
The {\hyperref[\detokenize{pair_modify::doc}]{\sphinxcrossref{\DUrole{doc}{pair\_modify}}}} table option is not relevant
for this pair style.

\sphinxAtStartPar
This pair style does not support the {\hyperref[\detokenize{pair_modify::doc}]{\sphinxcrossref{\DUrole{doc}{pair\_modify}}}}
tail option for adding long\sphinxhyphen{}range tail corrections to energy and
pressure.

\sphinxAtStartPar
This pair style writes its information to {\hyperref[\detokenize{restart::doc}]{\sphinxcrossref{\DUrole{doc}{binary restart files}}}}, so pair\_style and pair\_coeff commands do not need
to be specified in an input script that reads a restart file.

\sphinxAtStartPar
This pair style can only be used via the \sphinxstyleemphasis{pair} keyword of the
{\hyperref[\detokenize{run_style::doc}]{\sphinxcrossref{\DUrole{doc}{run\_style respa}}}} command.  It does not support the
\sphinxstyleemphasis{inner}, \sphinxstyleemphasis{middle}, \sphinxstyleemphasis{outer} keywords.


\bigskip\hrule\bigskip



\subsection{Restrictions}
\label{\detokenize{pair_lubricate:restrictions}}
\sphinxAtStartPar
These styles are part of the COLLOID package.  They are only enabled
if LAMMPS was built with that package.  See the {\hyperref[\detokenize{Build_package::doc}]{\sphinxcrossref{\DUrole{doc}{Build package}}}} page for more info.

\sphinxAtStartPar
Only spherical monodisperse particles are allowed for pair\_style
lubricate.

\sphinxAtStartPar
Only spherical particles are allowed for pair\_style lubricate/poly.

\sphinxAtStartPar
These pair styles will not restart exactly when using the
{\hyperref[\detokenize{read_restart::doc}]{\sphinxcrossref{\DUrole{doc}{read\_restart}}}} command, though they should provide
statistically similar results.  This is because the forces they
compute depend on atom velocities.  See the
{\hyperref[\detokenize{read_restart::doc}]{\sphinxcrossref{\DUrole{doc}{read\_restart}}}} command for more details.


\subsection{Related commands}
\label{\detokenize{pair_lubricate:related-commands}}
\sphinxAtStartPar
{\hyperref[\detokenize{pair_coeff::doc}]{\sphinxcrossref{\DUrole{doc}{pair\_coeff}}}}, {\hyperref[\detokenize{pair_lubricateU::doc}]{\sphinxcrossref{\DUrole{doc}{pair\_style lubricateU}}}}


\subsection{Default}
\label{\detokenize{pair_lubricate:default}}
\sphinxAtStartPar
The default settings for the optional args are flagHI = 1 and flagVF =
1.


\bigskip\hrule\bigskip

\phantomsection\label{\detokenize{pair_lubricate:ball1}}
\sphinxAtStartPar
\sphinxstylestrong{(Ball)} Ball and Melrose, Physica A, 247, 444\sphinxhyphen{}472 (1997).

\phantomsection\label{\detokenize{pair_lubricate:kumar1}}
\sphinxAtStartPar
\sphinxstylestrong{(Kumar)} Kumar and Higdon, Phys Rev E, 82, 051401 (2010).  See also
his thesis for more details: A. Kumar, “Microscale Dynamics in
Suspensions of Non\sphinxhyphen{}spherical Particles”, Thesis, University of
Illinois Urbana\sphinxhyphen{}Champaign,
(2010). (\sphinxurl{https://www.ideals.illinois.edu/handle/2142/16032})

\sphinxstepscope

\index{pair\_style lubricateU@\spxentry{pair\_style lubricateU}}\ignorespaces 
\index{pair\_style lubricateU/poly@\spxentry{pair\_style lubricateU/poly}}\ignorespaces 

\section{pair\_style lubricateU command}
\label{\detokenize{pair_lubricateU:pair-style-lubricateu-command}}\label{\detokenize{pair_lubricateU:index-1}}\label{\detokenize{pair_lubricateU:index-0}}\label{\detokenize{pair_lubricateU::doc}}

\section{pair\_style lubricateU/poly command}
\label{\detokenize{pair_lubricateU:pair-style-lubricateu-poly-command}}

\subsection{Syntax}
\label{\detokenize{pair_lubricateU:syntax}}
\begin{sphinxVerbatim}[commandchars=\\\{\}]
\PYG{k}{pair\PYGZus{}style}\PYG{+w}{ }\PYG{n}{style}\PYG{+w}{ }\PYG{n}{mu}\PYG{+w}{ }\PYG{n}{flaglog}\PYG{+w}{ }\PYG{n}{cutinner}\PYG{+w}{ }\PYG{n}{cutoff}\PYG{+w}{ }\PYG{n}{gdot}\PYG{+w}{ }\PYG{n}{flagHI}\PYG{+w}{ }\PYG{n}{flagVF}
\end{sphinxVerbatim}
\begin{itemize}
\item {} 
\sphinxAtStartPar
style = \sphinxstyleemphasis{lubricateU} or \sphinxstyleemphasis{lubricateU/poly}

\item {} 
\sphinxAtStartPar
mu = dynamic viscosity (dynamic viscosity units)

\item {} 
\sphinxAtStartPar
flaglog = 0/1 to exclude/include log terms in the lubrication approximation

\item {} 
\sphinxAtStartPar
cutinner = inner cut off distance (distance units)

\item {} 
\sphinxAtStartPar
cutoff = outer cutoff for interactions (distance units)

\item {} 
\sphinxAtStartPar
gdot = shear rate (1/time units)

\item {} 
\sphinxAtStartPar
flagHI (optional) = 0/1 to exclude/include 1/r hydrodynamic interactions

\item {} 
\sphinxAtStartPar
flagVF (optional) = 0/1 to exclude/include volume fraction corrections in the long\sphinxhyphen{}range isotropic terms

\end{itemize}


\subsection{Examples}
\label{\detokenize{pair_lubricateU:examples}}
\sphinxAtStartPar
(all assume radius = 1)

\begin{sphinxVerbatim}[commandchars=\\\{\}]
\PYG{k}{pair\PYGZus{}style}\PYG{+w}{ }\PYG{n}{lubricateU}\PYG{+w}{ }\PYG{l+m}{1.5}\PYG{+w}{ }\PYG{l+m}{1}\PYG{+w}{ }\PYG{l+m}{2.01}\PYG{+w}{ }\PYG{l+m}{2.5}\PYG{+w}{ }\PYG{l+m}{0.01}\PYG{+w}{ }\PYG{l+m}{1}\PYG{+w}{ }\PYG{l+m}{1}
\PYG{k}{pair\PYGZus{}coeff}\PYG{+w}{ }\PYG{l+m}{1}\PYG{+w}{ }\PYG{l+m}{1}\PYG{+w}{ }\PYG{l+m}{2.05}\PYG{+w}{ }\PYG{l+m}{2.8}
\PYG{k}{pair\PYGZus{}coeff}\PYG{+w}{ }\PYG{o}{*}\PYG{+w}{ }\PYG{o}{*}
\end{sphinxVerbatim}


\subsection{Description}
\label{\detokenize{pair_lubricateU:description}}
\sphinxAtStartPar
Styles \sphinxstyleemphasis{lubricateU} and \sphinxstyleemphasis{lubricateU/poly} compute velocities and
angular velocities for finite\sphinxhyphen{}size spherical particles such that the
hydrodynamic interaction balances the force and torque due to all
other types of interactions.

\sphinxAtStartPar
The interactions have 2 components.  The first is
Ball\sphinxhyphen{}Melrose lubrication terms via the formulas in {\hyperref[\detokenize{pair_lubricateU:ball2}]{\sphinxcrossref{\DUrole{std,std-ref}{(Ball and Melrose)}}}}
\begin{equation*}
\begin{split}W & =  - a_{sq} | (v_1 - v_2) \bullet \mathbf{nn} |^2 -
a_{sh} | (\omega_1 + \omega_2) \bullet
(\mathbf{I} - \mathbf{nn}) - 2 \Omega_N |^2 - \\
&  a_{pu} | (\omega_1 - \omega_2) \bullet (\mathbf{I} - \mathbf{nn}) |^2 -
a_{tw} | (\omega_1 - \omega_2) \bullet \mathbf{nn} |^2  \qquad r < r_c \\
& \\
\Omega_N & = \mathbf{n} \times (v_1 - v_2) / r\end{split}
\end{equation*}
\sphinxAtStartPar
which represents the dissipation W between two nearby particles due to
their relative velocities in the presence of a background solvent with
viscosity \sphinxstyleemphasis{mu}.  Note that this is dynamic viscosity which has units of
mass/distance/time, not kinematic viscosity.

\sphinxAtStartPar
The Asq (squeeze) term is the strongest and is included as long as
\sphinxstyleemphasis{flagHI} is set to 1 (default). It scales as 1/gap where gap is the
separation between the surfaces of the 2 particles. The Ash (shear)
and Apu (pump) terms are only included if \sphinxstyleemphasis{flaglog} is set to 1. They
are the next strongest interactions, and the only other singular
interaction, and scale as log(gap). Note that \sphinxstyleemphasis{flaglog} = 1 and
\sphinxstyleemphasis{flagHI} = 0 is invalid, and will result in a warning message, after
which \sphinxstyleemphasis{flagHI} will be set to 1. The Atw (twist) term is currently not
included. It is typically a very small contribution to the lubrication
forces.

\sphinxAtStartPar
The \sphinxstyleemphasis{flagHI} and \sphinxstyleemphasis{flagVF} settings are optional.  Neither should be
used, or both must be defined.

\sphinxAtStartPar
\sphinxstyleemphasis{Cutinner} sets the minimum center\sphinxhyphen{}to\sphinxhyphen{}center separation that will be
used in calculations irrespective of the actual separation.  \sphinxstyleemphasis{Cutoff}
is the maximum center\sphinxhyphen{}to\sphinxhyphen{}center separation at which an interaction is
computed.  Using a \sphinxstyleemphasis{cutoff} less than 3 radii is recommended if
\sphinxstyleemphasis{flaglog} is set to 1.

\sphinxAtStartPar
The other component is due to the Fast Lubrication Dynamics (FLD)
approximation, described in {\hyperref[\detokenize{pair_lubricateU:kumar2}]{\sphinxcrossref{\DUrole{std,std-ref}{(Kumar)}}}}.  The equation being
solved to balance the forces and torques is
\begin{equation*}
\begin{split}-R_{FU}(U-U^{\infty}) = -R_{FE}E^{\infty} - F^{rest}\end{split}
\end{equation*}
\sphinxAtStartPar
where U represents the velocities and angular velocities of the
particles, \(U^{\infty}\) represents the velocities and the angular
velocities of the undisturbed fluid, and \(E^{\infty}\) represents
the rate of strain tensor of the undisturbed fluid flow with viscosity
\sphinxstyleemphasis{mu}. Again, note that this is dynamic viscosity which has units of
mass/distance/time, not kinematic viscosity.  Volume fraction
corrections to R\_FU are included if \sphinxstyleemphasis{flagVF} is set to 1 (default).

\sphinxAtStartPar
F\sphinxstyleemphasis{rest} represents the forces and torques due to all other types of
interactions, e.g. Brownian, electrostatic etc.  Note that this
algorithm neglects the inertial terms, thereby removing the restriction
of resolving the small interial time scale, which may not be of interest
for colloidal particles.  These pair styles solve for the velocity such
that the hydrodynamic force balances all other types of forces, thereby
resulting in a net zero force (zero inertia limit).  When defining these
pair styles, they must be defined last so that when these styles are
invoked all other types of forces have already been computed.  For the
same reason, they won’t work if additional non\sphinxhyphen{}pair styles are defined
(such as bond or Kspace forces) as they are calculated in LAMMPS after
the pairwise interactions have been computed.

\begin{sphinxadmonition}{note}{Note:}
\sphinxAtStartPar
When using these styles, the pair styles are designed to be used with
implicit time integration and a correspondingly larger timestep.
Thus either {\hyperref[\detokenize{fix_nve_noforce::doc}]{\sphinxcrossref{\DUrole{doc}{fix nve/noforce}}}} should be used
for spherical particles defined via {\hyperref[\detokenize{atom_style::doc}]{\sphinxcrossref{\DUrole{doc}{atom\_style sphere}}}} or {\hyperref[\detokenize{fix_nve_asphere_noforce::doc}]{\sphinxcrossref{\DUrole{doc}{fix nve/asphere/noforce}}}} should be used for spherical particles
defined via {\hyperref[\detokenize{atom_style::doc}]{\sphinxcrossref{\DUrole{doc}{atom\_style ellipsoid}}}}.  This is
because the velocity and angular momentum of each particle is set by
the pair style, and should not be reset by the time integration fix.
\end{sphinxadmonition}

\sphinxAtStartPar
Style \sphinxstyleemphasis{lubricateU} requires monodisperse spherical particles; style
\sphinxstyleemphasis{lubricateU/poly} allows for polydisperse spherical particles.

\sphinxAtStartPar
If the suspension is sheared via the {\hyperref[\detokenize{fix_deform::doc}]{\sphinxcrossref{\DUrole{doc}{fix deform}}}}
command then the pair style uses the shear rate to adjust the
hydrodynamic interactions accordingly. Volume changes due to fix
deform are accounted for when computing the volume fraction
corrections to R\_FU.

\sphinxAtStartPar
When computing the volume fraction corrections to R\_FU, the presence of
walls (whether moving or stationary) will affect the volume fraction
available to colloidal particles. This is currently accounted for with
the following types of walls: {\hyperref[\detokenize{fix_wall::doc}]{\sphinxcrossref{\DUrole{doc}{wall/lj93}}}},
{\hyperref[\detokenize{fix_wall::doc}]{\sphinxcrossref{\DUrole{doc}{wall/lj126}}}}, {\hyperref[\detokenize{fix_wall::doc}]{\sphinxcrossref{\DUrole{doc}{wall/colloid}}}}, and
{\hyperref[\detokenize{fix_wall::doc}]{\sphinxcrossref{\DUrole{doc}{wall/harmonic}}}}.  For these wall styles, the correct
volume fraction will be used when walls do not coincide with the box
boundary, as well as when walls move and thereby cause a change in the
volume fraction. To use these wall styles with pair\_style \sphinxstyleemphasis{lubricateU}
or \sphinxstyleemphasis{lubricateU/poly}, the \sphinxstyleemphasis{fld yes} option must be specified in the fix
wall command.  Other wall styles may still work, but they will result
in the volume fraction being computed based on the box boundaries.
Several wall styles are not compatible with these pair styles and using
them will result in an error.

\sphinxAtStartPar
Since lubrication forces are dissipative, it is usually desirable to
thermostat the system at a constant temperature. If Brownian motion
(at a constant temperature) is desired, it can be set using the
{\hyperref[\detokenize{pair_brownian::doc}]{\sphinxcrossref{\DUrole{doc}{pair\_style brownian}}}} command. These pair styles
and the brownian style should use consistent parameters for \sphinxstyleemphasis{mu},
\sphinxstyleemphasis{flaglog}, \sphinxstyleemphasis{flagfld}, \sphinxstyleemphasis{cutinner}, \sphinxstyleemphasis{cutoff}, \sphinxstyleemphasis{flagHI} and \sphinxstyleemphasis{flagVF}.


\bigskip\hrule\bigskip


\sphinxAtStartPar
The following coefficients must be defined for each pair of atoms
types via the {\hyperref[\detokenize{pair_coeff::doc}]{\sphinxcrossref{\DUrole{doc}{pair\_coeff}}}} command as in the examples
above, or in the data file or restart files read by the
{\hyperref[\detokenize{read_data::doc}]{\sphinxcrossref{\DUrole{doc}{read\_data}}}} or {\hyperref[\detokenize{read_restart::doc}]{\sphinxcrossref{\DUrole{doc}{read\_restart}}}}
commands, or by mixing as described below:
\begin{itemize}
\item {} 
\sphinxAtStartPar
cutinner (distance units)

\item {} 
\sphinxAtStartPar
cutoff (distance units)

\end{itemize}

\sphinxAtStartPar
The two coefficients are optional.  If neither is specified, the two
cutoffs specified in the pair\_style command are used.  Otherwise both
must be specified.


\bigskip\hrule\bigskip



\subsection{Mixing, shift, table, tail correction, restart, rRESPA info}
\label{\detokenize{pair_lubricateU:mixing-shift-table-tail-correction-restart-rrespa-info}}
\sphinxAtStartPar
For atom type pairs I,J and I != J, the two cutoff distances for this
pair style can be mixed.  The default mix value is \sphinxstyleemphasis{geometric}.  See
the “pair\_modify” command for details.

\sphinxAtStartPar
These pair styles do not support the {\hyperref[\detokenize{pair_modify::doc}]{\sphinxcrossref{\DUrole{doc}{pair\_modify}}}}
shift option for the energy of the pair interaction.

\sphinxAtStartPar
The {\hyperref[\detokenize{pair_modify::doc}]{\sphinxcrossref{\DUrole{doc}{pair\_modify}}}} table option is not relevant
for these pair styles.

\sphinxAtStartPar
These pair styles do not support the {\hyperref[\detokenize{pair_modify::doc}]{\sphinxcrossref{\DUrole{doc}{pair\_modify}}}}
tail option for adding long\sphinxhyphen{}range tail corrections to energy and
pressure.

\sphinxAtStartPar
These pair styles write their information to {\hyperref[\detokenize{restart::doc}]{\sphinxcrossref{\DUrole{doc}{binary restart files}}}}, so pair\_style and pair\_coeff commands do not need to be
specified in an input script that reads a restart file.

\sphinxAtStartPar
These pair styles can only be used via the \sphinxstyleemphasis{pair} keyword of the
{\hyperref[\detokenize{run_style::doc}]{\sphinxcrossref{\DUrole{doc}{run\_style respa}}}} command.  They does not support the
\sphinxstyleemphasis{inner}, \sphinxstyleemphasis{middle}, \sphinxstyleemphasis{outer} keywords.


\bigskip\hrule\bigskip



\subsection{Restrictions}
\label{\detokenize{pair_lubricateU:restrictions}}
\sphinxAtStartPar
These styles are part of the COLLOID package.  They are only enabled if
LAMMPS was built with that package.  See the {\hyperref[\detokenize{Build_package::doc}]{\sphinxcrossref{\DUrole{doc}{Build package}}}} page for more info.

\sphinxAtStartPar
Currently, these pair styles assume that all other types of
forces/torques on the particles have been already been computed when it
is invoked.  This requires this style to be defined as the last of the
pair styles, and that no fixes apply additional constraint forces.  One
exception is the {\hyperref[\detokenize{fix_wall::doc}]{\sphinxcrossref{\DUrole{doc}{fix wall/colloid}}}} commands, which has
an “fld” option to apply their wall forces correctly.

\sphinxAtStartPar
Only spherical monodisperse particles are allowed for pair\_style
lubricateU.

\sphinxAtStartPar
Only spherical particles are allowed for pair\_style lubricateU/poly.

\sphinxAtStartPar
For sheared suspensions, it is assumed that the shearing is done in the
xy plane, with x being the velocity direction and y being the
velocity\sphinxhyphen{}gradient direction. In this case, one must use {\hyperref[\detokenize{fix_deform::doc}]{\sphinxcrossref{\DUrole{doc}{fix deform}}}} with the same rate of shear (erate).

\sphinxAtStartPar
These pair styles are only compatible with the following wall fixes:
doc:\sphinxtitleref{fix wall/lj93, fix wall/lj126, fix wall/lj1043, fix wall/colloid,
fix wall/harmonic, fix wall/lepton, fix wall/morse, fix wall/table
\textless{}fix\_wall\textgreater{}}.


\subsection{Related commands}
\label{\detokenize{pair_lubricateU:related-commands}}
\sphinxAtStartPar
{\hyperref[\detokenize{pair_coeff::doc}]{\sphinxcrossref{\DUrole{doc}{pair\_coeff}}}}, {\hyperref[\detokenize{pair_lubricate::doc}]{\sphinxcrossref{\DUrole{doc}{pair\_style lubricate}}}}


\subsection{Default}
\label{\detokenize{pair_lubricateU:default}}
\sphinxAtStartPar
The default settings for the optional args are flagHI = 1 and flagVF =
1.


\bigskip\hrule\bigskip

\phantomsection\label{\detokenize{pair_lubricateU:ball2}}
\sphinxAtStartPar
\sphinxstylestrong{(Ball)} Ball and Melrose, Physica A, 247, 444\sphinxhyphen{}472 (1997).

\phantomsection\label{\detokenize{pair_lubricateU:kumar2}}
\sphinxAtStartPar
\sphinxstylestrong{(Kumar)} Kumar and Higdon, Phys Rev E, 82, 051401 (2010).

\sphinxstepscope

\index{pair\_style lj/mdf@\spxentry{pair\_style lj/mdf}}\ignorespaces 
\index{pair\_style buck/mdf@\spxentry{pair\_style buck/mdf}}\ignorespaces 
\index{pair\_style lennard/mdf@\spxentry{pair\_style lennard/mdf}}\ignorespaces 

\section{pair\_style lj/mdf command}
\label{\detokenize{pair_mdf:pair-style-lj-mdf-command}}\label{\detokenize{pair_mdf:index-2}}\label{\detokenize{pair_mdf:index-1}}\label{\detokenize{pair_mdf:index-0}}\label{\detokenize{pair_mdf::doc}}

\section{pair\_style buck/mdf command}
\label{\detokenize{pair_mdf:pair-style-buck-mdf-command}}

\section{pair\_style lennard/mdf command}
\label{\detokenize{pair_mdf:pair-style-lennard-mdf-command}}

\subsection{Syntax}
\label{\detokenize{pair_mdf:syntax}}
\begin{sphinxVerbatim}[commandchars=\\\{\}]
\PYG{k}{pair\PYGZus{}style}\PYG{+w}{ }\PYG{n}{style}\PYG{+w}{ }\PYG{n}{args}
\end{sphinxVerbatim}
\begin{itemize}
\item {} 
\sphinxAtStartPar
style = \sphinxstyleemphasis{lj/mdf} or \sphinxstyleemphasis{buck/mdf} or \sphinxstyleemphasis{lennard/mdf}

\item {} 
\sphinxAtStartPar
args = list of arguments for a particular style
\begin{sphinxalltt}
\sphinxstyleemphasis{lj/mdf} args = cutoff1 cutoff2
  cutoff1 = inner cutoff for the start of the tapering function
  cutoff1 = out cutoff for the end of the tapering function
\sphinxstyleemphasis{buck/mdf} args = cutoff1 cutoff2
  cutoff1 = inner cutoff for the start of the tapering function
  cutoff1 = out cutoff for the end of the tapering function
\sphinxstyleemphasis{lennard/mdf} args = cutoff1 cutoff2
  cutoff1 = inner cutoff for the start of the tapering function
  cutoff1 = out cutoff for the end of the tapering function
\end{sphinxalltt}

\end{itemize}


\subsection{Examples}
\label{\detokenize{pair_mdf:examples}}
\begin{sphinxVerbatim}[commandchars=\\\{\}]
\PYG{k}{pair\PYGZus{}style}\PYG{+w}{ }\PYG{n}{lj}\PYG{o}{/}\PYG{n}{mdf}\PYG{+w}{ }\PYG{l+m}{2.5}\PYG{+w}{ }\PYG{l+m}{3.0}
\PYG{k}{pair\PYGZus{}coeff}\PYG{+w}{ }\PYG{o}{*}\PYG{+w}{ }\PYG{o}{*}\PYG{+w}{ }\PYG{l+m}{1.0}\PYG{+w}{ }\PYG{l+m}{1.0}
\PYG{k}{pair\PYGZus{}coeff}\PYG{+w}{ }\PYG{l+m}{1}\PYG{+w}{ }\PYG{l+m}{1}\PYG{+w}{ }\PYG{l+m}{1.1}\PYG{+w}{ }\PYG{l+m}{2.8}\PYG{+w}{ }\PYG{l+m}{3.0}\PYG{+w}{ }\PYG{l+m}{3.2}

\PYG{k}{pair\PYGZus{}style}\PYG{+w}{ }\PYG{n}{buck}\PYG{o}{/}\PYG{n}{mdf}\PYG{+w}{ }\PYG{l+m}{2.5}\PYG{+w}{ }\PYG{l+m}{3.0}
\PYG{k}{pair\PYGZus{}coeff}\PYG{+w}{ }\PYG{o}{*}\PYG{+w}{ }\PYG{o}{*}\PYG{+w}{ }\PYG{l+m}{100.0}\PYG{+w}{ }\PYG{l+m}{1.5}\PYG{+w}{ }\PYG{l+m}{200.0}
\PYG{k}{pair\PYGZus{}coeff}\PYG{+w}{ }\PYG{o}{*}\PYG{+w}{ }\PYG{o}{*}\PYG{+w}{ }\PYG{l+m}{100.0}\PYG{+w}{ }\PYG{l+m}{1.5}\PYG{+w}{ }\PYG{l+m}{200.0}\PYG{+w}{ }\PYG{l+m}{3.0}\PYG{+w}{ }\PYG{l+m}{3.5}

\PYG{k}{pair\PYGZus{}style}\PYG{+w}{ }\PYG{n}{lennard}\PYG{o}{/}\PYG{n}{mdf}\PYG{+w}{ }\PYG{l+m}{2.5}\PYG{+w}{ }\PYG{l+m}{3.0}
\PYG{k}{pair\PYGZus{}coeff}\PYG{+w}{ }\PYG{o}{*}\PYG{+w}{ }\PYG{o}{*}\PYG{+w}{ }\PYG{l+m}{1.0}\PYG{+w}{ }\PYG{l+m}{1.0}
\PYG{k}{pair\PYGZus{}coeff}\PYG{+w}{ }\PYG{l+m}{1}\PYG{+w}{ }\PYG{l+m}{1}\PYG{+w}{ }\PYG{l+m}{1021760.3664}\PYG{+w}{ }\PYG{l+m}{2120.317338}\PYG{+w}{ }\PYG{l+m}{3.0}\PYG{+w}{ }\PYG{l+m}{3.2}
\end{sphinxVerbatim}


\subsection{Description}
\label{\detokenize{pair_mdf:description}}
\sphinxAtStartPar
The \sphinxstyleemphasis{lj/mdf}, \sphinxstyleemphasis{buck/mdf} and \sphinxstyleemphasis{lennard/mdf} compute the standard 12\sphinxhyphen{}6
Lennard\sphinxhyphen{}Jones and Buckingham potential with the addition of a taper
function that ramps the energy and force smoothly to zero between an
inner and outer cutoff.
\begin{equation*}
\begin{split}E_{smooth}(r) = E(r)*f(r)\end{split}
\end{equation*}
\sphinxAtStartPar
The tapering, \sphinxstyleemphasis{f(r)}, is done by using the Mei, Davenport, Fernando
function {\hyperref[\detokenize{pair_mdf:mei}]{\sphinxcrossref{\DUrole{std,std-ref}{(Mei)}}}}.
\begin{equation*}
\begin{split}f(r) & = 1.0  \qquad \qquad \mathrm{for} \qquad r < r_m \\
f(r) & = (1 - x)^3*(1+3x+6x^2) \quad \mathrm{for} \qquad r_m < r < r_{cut} \\
f(r) & = 0.0  \qquad \qquad \mathrm{for} \qquad  r >= r_{cut} \\\end{split}
\end{equation*}
\sphinxAtStartPar
where
\begin{equation*}
\begin{split}x = \frac{(r-r_m)}{(r_{cut}-r_m)}\end{split}
\end{equation*}
\sphinxAtStartPar
Here \(r_m\) is the inner cutoff radius and \(r_{cut}\) is the
outer cutoff radius.


\bigskip\hrule\bigskip


\sphinxAtStartPar
For the \sphinxstyleemphasis{lj/mdf} pair\_style, the potential energy, \sphinxstyleemphasis{E(r)}, is the
standard 12\sphinxhyphen{}6 Lennard\sphinxhyphen{}Jones written in the epsilon/sigma form:
\begin{equation*}
\begin{split}E(r) = 4 \epsilon \left[ \left(\frac{\sigma}{r}\right)^{12} -
                         \left(\frac{\sigma}{r}\right)^6 \right]\end{split}
\end{equation*}
\sphinxAtStartPar
Either the first two or all of the following coefficients must be
defined for each pair of atoms types via the pair\_coeff command as in
the examples above, or in the data file read by the {\hyperref[\detokenize{read_data::doc}]{\sphinxcrossref{\DUrole{doc}{read\_data}}}}. The two cutoffs default to the global values and
\(\epsilon\) and \(\sigma\) can also be determined by mixing as
described below:
\begin{itemize}
\item {} 
\sphinxAtStartPar
\(\epsilon\) (energy units)

\item {} 
\sphinxAtStartPar
\(\sigma\) (distance units)

\item {} 
\sphinxAtStartPar
\(r_m\) (distance units)

\item {} 
\sphinxAtStartPar
\(r_{cut}\) (distance units)

\end{itemize}


\bigskip\hrule\bigskip


\sphinxAtStartPar
For the \sphinxstyleemphasis{buck/mdf} pair\_style, the potential energy, \sphinxstyleemphasis{E(r)}, is the
standard Buckingham potential with three required coefficients.
The two cutoffs can be omitted and default to the corresponding
global values:
\begin{equation*}
\begin{split}E(r) = A e^{(-r/\rho)} -\frac{C}{r^6}\end{split}
\end{equation*}\begin{itemize}
\item {} 
\sphinxAtStartPar
\sphinxstyleemphasis{A} (energy units)

\item {} 
\sphinxAtStartPar
\(\rho\) (distance units)

\item {} 
\sphinxAtStartPar
\sphinxstyleemphasis{C} (energy\sphinxhyphen{}distance\textasciicircum{}6 units)

\item {} 
\sphinxAtStartPar
\(r_m\) (distance units)

\item {} 
\sphinxAtStartPar
\(r_{cut}\) (distance units)

\end{itemize}


\bigskip\hrule\bigskip


\sphinxAtStartPar
For the \sphinxstyleemphasis{lennard/mdf} pair\_style, the potential energy, \sphinxstyleemphasis{E(r)}, is the
standard 12\sphinxhyphen{}6 Lennard\sphinxhyphen{}Jones written in the A/B form:
\begin{equation*}
\begin{split}E(r) = \frac{A}{r^{12}} - \frac{B}{r^{6}}\end{split}
\end{equation*}
\sphinxAtStartPar
The following coefficients must be defined for each pair of atoms
types via the pair\_coeff command as in the examples above, or in the
data file read by the read\_data commands, or by mixing as described below.
The two cutoffs default to their global values and must be either both
given or both left out:
\begin{itemize}
\item {} 
\sphinxAtStartPar
\sphinxstyleemphasis{A} (energy\sphinxhyphen{}distance\textasciicircum{}12 units)

\item {} 
\sphinxAtStartPar
\sphinxstyleemphasis{B} (energy\sphinxhyphen{}distance\textasciicircum{}6 units)

\item {} 
\sphinxAtStartPar
\(r_m\) (distance units)

\item {} 
\sphinxAtStartPar
\(r_{cut}\) (distance units)

\end{itemize}


\bigskip\hrule\bigskip



\subsection{Mixing, shift, table, tail correction, restart, rRESPA info}
\label{\detokenize{pair_mdf:mixing-shift-table-tail-correction-restart-rrespa-info}}
\sphinxAtStartPar
For atom type pairs I,J and I != J, the \(\epsilon\) and
\(\sigma\) coefficients and cutoff distances for the lj/mdf pair
style can be mixed.  The default mix value is \sphinxstyleemphasis{geometric}.  See the
“pair\_modify” command for details. The other two pair styles buck/mdf
and lennard/mdf do not support mixing, so all I,J pairs of coefficients
must be specified explicitly.

\sphinxAtStartPar
None of the lj/mdf, buck/mdf, or lennard/mdf pair styles supports
the {\hyperref[\detokenize{pair_modify::doc}]{\sphinxcrossref{\DUrole{doc}{pair\_modify}}}} shift option or long\sphinxhyphen{}range
tail corrections to pressure and energy.

\sphinxAtStartPar
These styles write their information to {\hyperref[\detokenize{restart::doc}]{\sphinxcrossref{\DUrole{doc}{binary restart files}}}}, so pair\_style and pair\_coeff commands do not need
to be specified in an input script that reads a restart file.

\sphinxAtStartPar
These styles can only be used via the \sphinxstyleemphasis{pair} keyword of the {\hyperref[\detokenize{run_style::doc}]{\sphinxcrossref{\DUrole{doc}{run\_style respa}}}} command.  They do not support the \sphinxstyleemphasis{inner},
\sphinxstyleemphasis{middle}, \sphinxstyleemphasis{outer} keywords.


\bigskip\hrule\bigskip



\subsection{Restrictions}
\label{\detokenize{pair_mdf:restrictions}}
\sphinxAtStartPar
These pair styles can only be used if LAMMPS was built with the
EXTRA\sphinxhyphen{}PAIR package.  See the {\hyperref[\detokenize{Build_package::doc}]{\sphinxcrossref{\DUrole{doc}{Build package}}}} doc
page for more info.


\subsection{Related commands}
\label{\detokenize{pair_mdf:related-commands}}
\sphinxAtStartPar
{\hyperref[\detokenize{pair_coeff::doc}]{\sphinxcrossref{\DUrole{doc}{pair\_coeff}}}}


\subsection{Default}
\label{\detokenize{pair_mdf:default}}
\sphinxAtStartPar
none


\bigskip\hrule\bigskip

\phantomsection\label{\detokenize{pair_mdf:mei}}
\sphinxAtStartPar
\sphinxstylestrong{(Mei)} Mei, Davenport, Fernando, Phys Rev B, 43 4653 (1991)

\sphinxstepscope

\index{pair\_style meam@\spxentry{pair\_style meam}}\ignorespaces 
\index{pair\_style meam/kk@\spxentry{pair\_style meam/kk}}\ignorespaces 
\index{pair\_style meam/ms@\spxentry{pair\_style meam/ms}}\ignorespaces 
\index{pair\_style meam/ms/kk@\spxentry{pair\_style meam/ms/kk}}\ignorespaces 

\section{pair\_style meam command}
\label{\detokenize{pair_meam:pair-style-meam-command}}\label{\detokenize{pair_meam:index-3}}\label{\detokenize{pair_meam:index-2}}\label{\detokenize{pair_meam:index-1}}\label{\detokenize{pair_meam:index-0}}\label{\detokenize{pair_meam::doc}}
\sphinxAtStartPar
Accelerator Variants: \sphinxstyleemphasis{meam/kk}


\section{pair\_style meam/ms command}
\label{\detokenize{pair_meam:pair-style-meam-ms-command}}
\sphinxAtStartPar
Accelerator Variants: \sphinxstyleemphasis{meam/ms/kk}


\subsection{Syntax}
\label{\detokenize{pair_meam:syntax}}
\begin{sphinxVerbatim}[commandchars=\\\{\}]
\PYG{k}{pair\PYGZus{}style}\PYG{+w}{ }\PYG{n}{style}
\end{sphinxVerbatim}
\begin{itemize}
\item {} 
\sphinxAtStartPar
style = \sphinxstyleemphasis{meam} or \sphinxstyleemphasis{meam/ms}

\end{itemize}


\subsection{Examples}
\label{\detokenize{pair_meam:examples}}
\begin{sphinxVerbatim}[commandchars=\\\{\}]
\PYG{k}{pair\PYGZus{}style}\PYG{+w}{ }\PYG{n}{meam}
\PYG{k}{pair\PYGZus{}coeff}\PYG{+w}{ }\PYG{o}{*}\PYG{+w}{ }\PYG{o}{*}\PYG{+w}{ }\PYG{n}{..}\PYG{o}{/}\PYG{n}{potentials}\PYG{o}{/}\PYG{n}{library.meam}\PYG{+w}{ }\PYG{n}{Si}\PYG{+w}{ }\PYG{n}{..}\PYG{o}{/}\PYG{n}{potentials}\PYG{o}{/}\PYG{n}{si.meam}\PYG{+w}{ }\PYG{n}{Si}
\PYG{k}{pair\PYGZus{}coeff}\PYG{+w}{ }\PYG{o}{*}\PYG{+w}{ }\PYG{o}{*}\PYG{+w}{ }\PYG{n}{..}\PYG{o}{/}\PYG{n}{potentials}\PYG{o}{/}\PYG{n}{library.meam}\PYG{+w}{ }\PYG{n}{Ni}\PYG{+w}{ }\PYG{n}{Al}\PYG{+w}{ }\PYG{n}{NULL}\PYG{+w}{ }\PYG{n}{Ni}\PYG{+w}{ }\PYG{n}{Al}\PYG{+w}{ }\PYG{n}{Ni}\PYG{+w}{ }\PYG{n}{Ni}

\PYG{k}{pair\PYGZus{}style}\PYG{+w}{ }\PYG{n}{meam}\PYG{o}{/}\PYG{n}{ms}
\PYG{k}{pair\PYGZus{}coeff}\PYG{+w}{ }\PYG{o}{*}\PYG{+w}{ }\PYG{o}{*}\PYG{+w}{ }\PYG{n}{..}\PYG{o}{/}\PYG{n}{potentials}\PYG{o}{/}\PYG{n}{library.msmeam}\PYG{+w}{ }\PYG{n}{H}\PYG{+w}{ }\PYG{n}{Ga}\PYG{+w}{ }\PYG{n}{..}\PYG{o}{/}\PYG{n}{potentials}\PYG{o}{/}\PYG{n}{HGa.meam}\PYG{+w}{ }\PYG{n}{H}\PYG{+w}{ }\PYG{n}{Ga}
\end{sphinxVerbatim}


\subsection{Description}
\label{\detokenize{pair_meam:description}}
\begin{sphinxadmonition}{note}{Note:}
\sphinxAtStartPar
The behavior of the MEAM potential for alloy systems has changed
as of November 2010; see description below of the mixture\_ref\_t
parameter
\end{sphinxadmonition}

\sphinxAtStartPar
Pair style \sphinxstyleemphasis{meam} computes non\sphinxhyphen{}bonded interactions for a variety of
materials using the modified embedded\sphinxhyphen{}atom method (MEAM) {\hyperref[\detokenize{pair_meam:baskes}]{\sphinxcrossref{\DUrole{std,std-ref}{(Baskes)}}}}.  Conceptually, it is an extension to the original {\hyperref[\detokenize{pair_eam::doc}]{\sphinxcrossref{\DUrole{doc}{EAM
method}}}} which adds angular forces.  It is thus suitable for
modeling metals and alloys with fcc, bcc, hcp and diamond cubic
structures, as well as materials with covalent interactions like silicon
and carbon.

\sphinxAtStartPar
The \sphinxstyleemphasis{meam} pair style is a translation of the original Fortran version
to C++. It is functionally equivalent but more efficient and has
additional features. The Fortran version of the \sphinxstyleemphasis{meam} pair style has
been removed from LAMMPS after the 12 December 2018 release.

\sphinxAtStartPar
Pair style \sphinxstyleemphasis{meam/ms} uses the multi\sphinxhyphen{}state MEAM (MS\sphinxhyphen{}MEAM) method
according to {\hyperref[\detokenize{pair_meam:baskes2}]{\sphinxcrossref{\DUrole{std,std-ref}{(Baskes2)}}}}, which is an extension to MEAM.
This pair style is mostly equivalent to \sphinxstyleemphasis{meam} and differs only
where noted in the documentation below.

\sphinxAtStartPar
In the MEAM formulation, the total energy E of a system of atoms is
given by:
\begin{equation*}
\begin{split}E = \sum_i \left\{ F_i(\bar{\rho}_i)
    + \frac{1}{2} \sum_{i \neq j} \phi_{ij} (r_{ij}) \right\}\end{split}
\end{equation*}
\sphinxAtStartPar
where \sphinxstyleemphasis{F} is the embedding energy which is a function of the atomic
electron density \(\rho\), and \(\phi\) is a pair potential
interaction.  The pair interaction is summed over all neighbors J of
atom I within the cutoff distance.  As with EAM, the multi\sphinxhyphen{}body nature
of the MEAM potential is a result of the embedding energy term.  Details
of the computation of the embedding and pair energies, as implemented in
LAMMPS, are given in {\hyperref[\detokenize{pair_meam:gullet}]{\sphinxcrossref{\DUrole{std,std-ref}{(Gullet)}}}} and references therein.

\sphinxAtStartPar
The various parameters in the MEAM formulas are listed in two files
which are specified by the {\hyperref[\detokenize{pair_coeff::doc}]{\sphinxcrossref{\DUrole{doc}{pair\_coeff}}}} command.
These are ASCII text files in a format consistent with other MD codes
that implement MEAM potentials, such as the serial DYNAMO code and
Warp.  Several MEAM potential files with parameters for different
materials are included in the “potentials” directory of the LAMMPS
distribution with a “.meam” suffix.  All of these are parameterized in
terms of LAMMPS {\hyperref[\detokenize{units::doc}]{\sphinxcrossref{\DUrole{doc}{metal units}}}}.

\sphinxAtStartPar
Note that unlike for other potentials, cutoffs for MEAM potentials are
not set in the pair\_style or pair\_coeff command; they are specified in
the MEAM potential files themselves.

\sphinxAtStartPar
Only a single pair\_coeff command is used with the \sphinxstyleemphasis{meam} style which
specifies two MEAM files and the element(s) to extract information
for.  The MEAM elements are mapped to LAMMPS atom types by specifying
N additional arguments after the second filename in the pair\_coeff
command, where N is the number of LAMMPS atom types:
\begin{itemize}
\item {} 
\sphinxAtStartPar
MEAM library file

\item {} 
\sphinxAtStartPar
Element1, Element2, …

\item {} 
\sphinxAtStartPar
MEAM parameter file

\item {} 
\sphinxAtStartPar
N element names = mapping of MEAM elements to atom types

\end{itemize}

\sphinxAtStartPar
See the {\hyperref[\detokenize{pair_coeff::doc}]{\sphinxcrossref{\DUrole{doc}{pair\_coeff}}}} page for alternate ways
to specify the path for the potential files.

\sphinxAtStartPar
As an example, the \sphinxcode{\sphinxupquote{potentials/library.meam}} file has generic MEAM
settings for a variety of elements.  The \sphinxcode{\sphinxupquote{potentials/SiC.meam}} file
has specific parameter settings for a Si and C alloy system.  If your
LAMMPS simulation has 4 atoms types and you want the first 3 to be Si,
and the fourth to be C, you would use the following pair\_coeff command:

\begin{sphinxVerbatim}[commandchars=\\\{\}]
\PYG{k}{pair\PYGZus{}coeff}\PYG{+w}{ }\PYG{o}{*}\PYG{+w}{ }\PYG{o}{*}\PYG{+w}{ }\PYG{n}{library.meam}\PYG{+w}{ }\PYG{n}{Si}\PYG{+w}{ }\PYG{n}{C}\PYG{+w}{ }\PYG{n}{sic.meam}\PYG{+w}{ }\PYG{n}{Si}\PYG{+w}{ }\PYG{n}{Si}\PYG{+w}{ }\PYG{n}{Si}\PYG{+w}{ }\PYG{n}{C}
\end{sphinxVerbatim}

\sphinxAtStartPar
The first 2 arguments must be * * so as to span all LAMMPS atom types.
The first filename is the element library file. The list of elements following
it extracts lines from the library file and assigns numeric indices to these
elements. The second filename is the alloy parameter file, which refers to
elements using the numeric indices assigned before.
The arguments after the parameter file map LAMMPS atom types to elements, i.e.
LAMMPS atom types 1,2,3 to the MEAM Si element.  The final C argument maps
LAMMPS atom type 4 to the MEAM C element.

\sphinxAtStartPar
If the second filename is specified as NULL, no parameter file is read,
which simply means the generic parameters in the library file are
used.  Use of the NULL specification for the parameter file is
discouraged for systems with more than a single element type
(e.g. alloys), since the parameter file is expected to set element
interaction terms that are not captured by the information in the
library file.

\sphinxAtStartPar
If a mapping value is specified as NULL, the mapping is not performed.
This can be used when a \sphinxstyleemphasis{meam} potential is used as part of the
\sphinxstyleemphasis{hybrid} pair style.  The NULL values are placeholders for atom types
that will be used with other potentials.

\begin{sphinxadmonition}{note}{Note:}
\sphinxAtStartPar
If the second filename is NULL, the element names between the two
filenames can appear in any order, e.g. “Si C” or “C Si” in the
example above.  However, if the second filename is \sphinxstylestrong{not} NULL (as in the
example above), it contains settings that are indexed \sphinxstylestrong{by numbers}
for the elements that precede it.  Thus you need to ensure that you list
the elements between the filenames in an order consistent with how the
values in the second filename are indexed.  See details below on the
syntax for settings in the second file.
\end{sphinxadmonition}

\sphinxAtStartPar
The MEAM library file provided with LAMMPS has the name
\sphinxcode{\sphinxupquote{potentials/library.meam}}.  It is the “meamf” file used by other MD
codes.  Aside from blank and comment lines (starting with \# which can
appear anywhere), it is formatted as a series of entries, each of which
has 19 parameters and can span multiple lines:

\sphinxAtStartPar
elt, lat, z, ielement, atwt, alpha, b0, b1, b2, b3, alat, esub, asub,
t0, t1, t2, t3, rozero, ibar

\sphinxAtStartPar
The \sphinxstyleemphasis{elt} and \sphinxstyleemphasis{lat} parameters are text strings, such as \sphinxstyleemphasis{elt} = Si or
Cu and \sphinxstyleemphasis{lat} = dia or fcc.  Because the library file is used by Fortran
MD codes, these strings may be enclosed in single quotes, but this is
not required.  The other numeric parameters match values in the
formulas above.  The value of the \sphinxstyleemphasis{elt} string is what is used in the
pair\_coeff command to identify which settings from the library file
you wish to read in.  There can be multiple entries in the library
file with the same \sphinxstyleemphasis{elt} value; LAMMPS reads the first matching entry it
finds and ignores the rest.

\sphinxAtStartPar
Other parameters in the MEAM library file correspond to single\sphinxhyphen{}element
potential parameters:
\begin{sphinxalltt}
lat      = lattice structure of reference configuration
z        = number of nearest neighbors in the reference structure
ielement = atomic number
atwt     = atomic weight
alat     = lattice constant of reference structure
esub     = energy per atom (eV) in the reference structure at equilibrium
asub     = "A" parameter for MEAM (see e.g. {\hyperref[\detokenize{pair_meam:baskes}]{\sphinxcrossref{\DUrole{std,std-ref}{(Baskes)}}}})
\end{sphinxalltt}

\sphinxAtStartPar
The \sphinxstyleemphasis{alpha}, \sphinxstyleemphasis{b0}, \sphinxstyleemphasis{b1}, \sphinxstyleemphasis{b2}, \sphinxstyleemphasis{b3}, \sphinxstyleemphasis{t0}, \sphinxstyleemphasis{t1}, \sphinxstyleemphasis{t2}, \sphinxstyleemphasis{t3} parameters correspond to the
standard MEAM parameters in the literature {\hyperref[\detokenize{pair_meam:baskes}]{\sphinxcrossref{\DUrole{std,std-ref}{(Baskes)}}}} (the b
parameters are the standard beta parameters). Note that only parameters
normalized to \sphinxstyleemphasis{t0 = 1.0} are supported.  The \sphinxstyleemphasis{rozero} parameter is
an element\sphinxhyphen{}dependent density scaling that weights the reference
background density (see e.g. equation 4.5 in {\hyperref[\detokenize{pair_meam:gullet}]{\sphinxcrossref{\DUrole{std,std-ref}{(Gullet)}}}}) and
is typically 1.0 for single\sphinxhyphen{}element systems.  The \sphinxstyleemphasis{ibar} parameter
selects the form of the function G(Gamma) used to compute the electron
density; options are

\begin{sphinxVerbatim}[commandchars=\\\{\}]
 0 =\PYGZgt{} G = sqrt(1+Gamma)
 1 =\PYGZgt{} G = exp(Gamma/2)
 2 =\PYGZgt{} not implemented
 3 =\PYGZgt{} G = 2/(1+exp(\PYGZhy{}Gamma))
 4 =\PYGZgt{} G = sqrt(1+Gamma)
\PYGZhy{}5 =\PYGZgt{} G = +\PYGZhy{}sqrt(abs(1+Gamma))
\end{sphinxVerbatim}

\sphinxAtStartPar
If used, the MEAM parameter file contains settings that override or
complement the library file settings.  Examples of such parameter
files are in the potentials directory with a “.meam” suffix.  Their
format is the same as is read by other Fortran MD codes.  Aside from
blank and comment lines (start with \# which can appear anywhere), each
line has one of the following forms.  Each line can also have a
trailing comment (starting with \#) which is ignored.

\begin{sphinxVerbatim}[commandchars=\\\{\}]
keyword = value
keyword(I) = value
keyword(I,J) = value
keyword(I,J,K) = value
\end{sphinxVerbatim}

\sphinxAtStartPar
The indices I, J, K correspond to the elements selected from the
MEAM library file numbered in the order of how those elements were
selected starting from 1. Thus for the example given before

\begin{sphinxVerbatim}[commandchars=\\\{\}]
\PYG{k}{pair\PYGZus{}coeff}\PYG{+w}{ }\PYG{o}{*}\PYG{+w}{ }\PYG{o}{*}\PYG{+w}{ }\PYG{n}{library.meam}\PYG{+w}{ }\PYG{n}{Si}\PYG{+w}{ }\PYG{n}{C}\PYG{+w}{ }\PYG{n}{sic.meam}\PYG{+w}{ }\PYG{n}{Si}\PYG{+w}{ }\PYG{n}{Si}\PYG{+w}{ }\PYG{n}{Si}\PYG{+w}{ }\PYG{n}{C}
\end{sphinxVerbatim}

\sphinxAtStartPar
an index of 1 would refer to Si and an index of 2 to C.

\sphinxAtStartPar
The recognized keywords for the parameter file are as follows:
\begin{sphinxalltt}
rc          = cutoff radius for cutoff function; default = 4.0
delr        = length of smoothing distance for cutoff function; default = 0.1
rho0(I)     = relative density for element I (overwrites value
              read from meamf file)
Ec(I,J)     = cohesive energy of reference structure for I\sphinxhyphen{}J mixture
delta(I,J)  = heat of formation for I\sphinxhyphen{}J alloy; if Ec\_IJ is input as
              zero, then LAMMPS sets Ec\_IJ = (Ec\_II + Ec\_JJ)/2 \sphinxhyphen{} delta\_IJ
alpha(I,J)  = alpha parameter for pair potential between I and J (can
              be computed from bulk modulus of reference structure)
re(I,J)     = equilibrium distance between I and J in the reference
              structure
Cmax(I,J,K) = Cmax screening parameter when I\sphinxhyphen{}J pair is screened
              by K (I\textless{}=J); default = 2.8
Cmin(I,J,K) = Cmin screening parameter when I\sphinxhyphen{}J pair is screened
              by K (I\textless{}=J); default = 2.0
lattce(I,J) = lattice structure of I\sphinxhyphen{}J reference structure:
                fcc = face centered cubic
                bcc = body centered cubic
                hcp = hexagonal close\sphinxhyphen{}packed
                dim = dimer
                dia = diamond (interlaced fcc for alloy)
                dia3= diamond structure with primary 1NN and secondary 3NN interaction
                b1  = rock salt (NaCl structure)
                c11 = MoSi2 structure
                l12 = Cu3Au structure (lower case L, followed by 12)
                b2  = CsCl structure (interpenetrating simple cubic)
                ch4 = methane\sphinxhyphen{}like structure, only for binary system
                lin = linear structure (180 degree angle)
                zig = zigzag structure with a uniform angle
                tri = H2O\sphinxhyphen{}like structure that has an angle
                sc  = simple cubic
nn2(I,J)    = turn on second\sphinxhyphen{}nearest neighbor MEAM formulation for
              I\sphinxhyphen{}J pair (see for example {\hyperref[\detokenize{pair_meam:lee}]{\sphinxcrossref{\DUrole{std,std-ref}{(Lee)}}}}).
                0 = second\sphinxhyphen{}nearest neighbor formulation off
                1 = second\sphinxhyphen{}nearest neighbor formulation on
                default = 0
attrac(I,J) = additional cubic attraction term in Rose energy I\sphinxhyphen{}J pair potential
                default = 0
repuls(I,J) = additional cubic repulsive term in Rose energy I\sphinxhyphen{}J pair potential
                default = 0
zbl(I,J)    = blend the MEAM I\sphinxhyphen{}J pair potential with the ZBL potential for small
              atom separations {\hyperref[\detokenize{pair_meam:zbl}]{\sphinxcrossref{\DUrole{std,std-ref}{(ZBL)}}}}
                default = 1
theta(I,J)  = angle between three atoms in line, zigzag, and trimer reference structures in degrees
                default = 180
gsmooth\_factor  = factor determining the length of the G\sphinxhyphen{}function smoothing
                  region; only significant for ibar=0 or ibar=4.
                      99.0 = short smoothing region, sharp step
                      0.5  = long smoothing region, smooth step
                      default = 99.0
augt1           = integer flag for whether to augment t1 parameter by
                  3/5*t3 to account for old vs. new meam formulations;
                    0 = don\textquotesingle{}t augment t1
                    1 = augment t1
                    default = 1
ialloy          = integer flag to use alternative averaging rule for t parameters,
                  for comparison with the DYNAMO MEAM code
                    0 = standard averaging (matches ialloy=0 in DYNAMO)
                    1 = alternative averaging (matches ialloy=1 in DYNAMO)
                    2 = no averaging of t (use single\sphinxhyphen{}element values)
                    default = 0
mixture\_ref\_t   = integer flag to use mixture average of t to compute the background
                  reference density for alloys, instead of the single\sphinxhyphen{}element values
                  (see description and warning elsewhere in this doc page)
                    0 = do not use mixture averaging for t in the reference density
                    1 = use mixture averaging for t in the reference density
                    default = 0
erose\_form      = integer value to select the form of the Rose energy function
                  (see description below).
                    default = 0
emb\_lin\_neg     = integer value to select embedding function for negative densities
                    0 = F(rho)=0
                    1 = F(rho) = \sphinxhyphen{}asub*esub*rho (linear in rho, matches DYNAMO)
                    default = 0
bkgd\_dyn        = integer value to select background density formula
                    0 = rho\_bkgd = rho\_ref\_meam(a) (as in the reference structure)
                    1 = rho\_bkgd = rho0\_meam(a)*Z\_meam(a) (matches DYNAMO)
                    default = 0
\end{sphinxalltt}

\sphinxAtStartPar
\sphinxstyleemphasis{Rc}, \sphinxstyleemphasis{delr}, \sphinxstyleemphasis{re} are in distance units (Angstroms in the case of metal
units).  \sphinxstyleemphasis{Ec} and \sphinxstyleemphasis{delta} are in energy units (eV in the case of metal
units).

\sphinxAtStartPar
Each keyword represents a quantity which is either a scalar, vector,
2d array, or 3d array and must be specified with the correct
corresponding array syntax.  The indices I,J,K each run from 1 to N
where N is the number of MEAM elements being used.

\sphinxAtStartPar
Thus these lines

\begin{sphinxVerbatim}[commandchars=\\\{\}]
rho0(2) = 2.25
alpha(1,2) = 4.37
\end{sphinxVerbatim}

\sphinxAtStartPar
set \sphinxstyleemphasis{rho0} for the second element to the value 2.25 and set \sphinxstyleemphasis{alpha} for the
alloy interaction between elements 1 and 2 to 4.37.

\sphinxAtStartPar
The \sphinxstyleemphasis{augt1} parameter is related to modifications in the MEAM
formulation of the partial electron density function.  In recent
literature, an extra term is included in the expression for the
third\sphinxhyphen{}order density in order to make the densities orthogonal (see for
example {\hyperref[\detokenize{pair_meam:wang2}]{\sphinxcrossref{\DUrole{std,std-ref}{(Wang)}}}}, equation 3d); this term is included in the
MEAM implementation in LAMMPS.  However, in earlier published work
this term was not included when deriving parameters, including most of
those provided in the \sphinxcode{\sphinxupquote{library.meam}} file included with LAMMPS, and to
account for this difference the parameter \sphinxstyleemphasis{t1} must be augmented by
3/5**t3*.  If \sphinxstyleemphasis{augt1} = 1, the default, this augmentation is done
automatically.  When parameter values are fit using the modified
density function, as in more recent literature, augt1 should be set to
0.

\sphinxAtStartPar
The \sphinxstyleemphasis{mixture\_ref\_t} parameter is available to match results with those
of previous versions of LAMMPS (before January 2011).  Newer versions
of LAMMPS, by default, use the single\sphinxhyphen{}element values of the \sphinxstyleemphasis{t}
parameters to compute the background reference density.  This is the
proper way to compute these parameters.  Earlier versions of LAMMPS
used an alloy mixture averaged value of \sphinxstyleemphasis{t} to compute the background
reference density.  Setting \sphinxstyleemphasis{mixture\_ref\_t} = 1 gives the old behavior.
WARNING: using \sphinxstyleemphasis{mixture\_ref\_t} = 1 will give results that are demonstrably
incorrect for second\sphinxhyphen{}neighbor MEAM, and non\sphinxhyphen{}standard for
first\sphinxhyphen{}neighbor MEAM; this option is included only for matching with
previous versions of LAMMPS and should be avoided if possible.

\sphinxAtStartPar
The parameters \sphinxstyleemphasis{attrac} and \sphinxstyleemphasis{repuls}, along with the integer selection
parameter \sphinxstyleemphasis{erose\_form}, can be used to modify the Rose energy function
used to compute the pair potential.  This function gives the energy of
the reference state as a function of interatomic spacing.  The form of
this function is:
\begin{sphinxalltt}
astar = alpha * (r/re \sphinxhyphen{} 1.d0)
if erose\_form = 0: erose = \sphinxhyphen{}Ec*(1+astar+a3*(astar**3)/(r/re))*exp(\sphinxhyphen{}astar)
if erose\_form = 1: erose = \sphinxhyphen{}Ec*(1+astar+(\sphinxhyphen{}attrac+repuls/r)*(astar**3))*exp(\sphinxhyphen{}astar)
if erose\_form = 2: erose = \sphinxhyphen{}Ec*(1 +astar + a3*(astar**3))*exp(\sphinxhyphen{}astar)
a3 = repuls, astar \textless{} 0
a3 = attrac, astar \textgreater{}= 0
\end{sphinxalltt}

\sphinxAtStartPar
Most published MEAM parameter sets use the default values \sphinxstyleemphasis{attrac} = \sphinxstyleemphasis{repulse} = 0.
Setting \sphinxstyleemphasis{repuls} = \sphinxstyleemphasis{attrac} = \sphinxstyleemphasis{delta} corresponds to the form used in several
recent published MEAM parameter sets, such as {\hyperref[\detokenize{pair_meam:valone}]{\sphinxcrossref{\DUrole{std,std-ref}{(Valone)}}}}

\sphinxAtStartPar
Then using \sphinxstyleemphasis{meam/ms} pair style the multi\sphinxhyphen{}state MEAM (MS\sphinxhyphen{}MEAM) method is
activated.  This requires 6 extra parameters in the MEAM library file,
resulting in 25 parameters ordered that are ordered like this:

\sphinxAtStartPar
elt, lat, z, ielement, atwt, alpha, b0, b1, b2, b3, b1m, b2m, b3m, alat, esub, asub,
t0, t1, t2, t3, t1m, t2m, t3m, rozero, ibar

\sphinxAtStartPar
The 6 extra MS\sphinxhyphen{}MEAM parameters are \sphinxstyleemphasis{b1m, b2m, b3m, t1m, t2m, t3m}.
In the LAMMPS \sphinxcode{\sphinxupquote{potentials}} folder, compatible files have an “.msmeam” extension.


\bigskip\hrule\bigskip


\sphinxAtStartPar
Styles with a \sphinxstyleemphasis{gpu}, \sphinxstyleemphasis{intel}, \sphinxstyleemphasis{kk}, \sphinxstyleemphasis{omp}, or \sphinxstyleemphasis{opt} suffix are
functionally the same as the corresponding style without the suffix.
They have been optimized to run faster, depending on your available
hardware, as discussed on the {\hyperref[\detokenize{Speed_packages::doc}]{\sphinxcrossref{\DUrole{doc}{Accelerator packages}}}}
page.  The accelerated styles take the same arguments and should
produce the same results, except for round\sphinxhyphen{}off and precision issues.

\sphinxAtStartPar
These accelerated styles are part of the GPU, INTEL, KOKKOS,
OPENMP, and OPT packages, respectively.  They are only enabled if
LAMMPS was built with those packages.  See the {\hyperref[\detokenize{Build_package::doc}]{\sphinxcrossref{\DUrole{doc}{Build package}}}} page for more info.

\sphinxAtStartPar
You can specify the accelerated styles explicitly in your input script
by including their suffix, or you can use the {\hyperref[\detokenize{Run_options::doc}]{\sphinxcrossref{\DUrole{doc}{\sphinxhyphen{}suffix command\sphinxhyphen{}line switch}}}} when you invoke LAMMPS, or you can use the
{\hyperref[\detokenize{suffix::doc}]{\sphinxcrossref{\DUrole{doc}{suffix}}}} command in your input script.

\sphinxAtStartPar
See the {\hyperref[\detokenize{Speed_packages::doc}]{\sphinxcrossref{\DUrole{doc}{Accelerator packages}}}} page for more
instructions on how to use the accelerated styles effectively.


\bigskip\hrule\bigskip


\begin{sphinxadmonition}{note}{Note:}
\sphinxAtStartPar
The default form of the \sphinxstyleemphasis{erose} expression in LAMMPS was corrected
in March 2009.  The current version is correct, but may show different
behavior compared with earlier versions of LAMMPS with the attrac
and/or repuls parameters are non\sphinxhyphen{}zero.  To obtain the previous default
form, use \sphinxstyleemphasis{erose\_form} = 1 (this form does not seem to appear in the
literature).  An alternative form (see e.g. {\hyperref[\detokenize{pair_meam:lee2}]{\sphinxcrossref{\DUrole{std,std-ref}{(Lee2)}}}}) is
available using \sphinxstyleemphasis{erose\_form} = 2.
\end{sphinxadmonition}


\bigskip\hrule\bigskip



\subsection{Mixing, shift, table, tail correction, restart, rRESPA info}
\label{\detokenize{pair_meam:mixing-shift-table-tail-correction-restart-rrespa-info}}
\sphinxAtStartPar
For atom type pairs I,J and I != J, where types I and J correspond to
two different element types, mixing is performed by LAMMPS with
user\sphinxhyphen{}specifiable parameters as described above.

\sphinxAtStartPar
This pair style does not support the {\hyperref[\detokenize{pair_modify::doc}]{\sphinxcrossref{\DUrole{doc}{pair\_modify}}}}
\sphinxstyleemphasis{shift}, \sphinxstyleemphasis{table}, and \sphinxstyleemphasis{tail} options.

\sphinxAtStartPar
This pair style does not write its information to {\hyperref[\detokenize{restart::doc}]{\sphinxcrossref{\DUrole{doc}{binary restart files}}}},
since it is stored in potential files.  Thus, you
need to re\sphinxhyphen{}specify the pair\_style and pair\_coeff commands in an input
script that reads a restart file.

\sphinxAtStartPar
This pair style can only be used via the \sphinxstyleemphasis{pair} keyword of the
{\hyperref[\detokenize{run_style::doc}]{\sphinxcrossref{\DUrole{doc}{run\_style respa}}}} command.  It does not support the
\sphinxstyleemphasis{inner}, \sphinxstyleemphasis{middle}, \sphinxstyleemphasis{outer} keywords.


\bigskip\hrule\bigskip



\subsection{Restrictions}
\label{\detokenize{pair_meam:restrictions}}
\sphinxAtStartPar
The \sphinxstyleemphasis{meam} and \sphinxstyleemphasis{meam/ms} pair styles are provided in the MEAM
package. They are only enabled if LAMMPS was built with that package.
See the {\hyperref[\detokenize{Build_package::doc}]{\sphinxcrossref{\DUrole{doc}{Build package}}}} page for more info.

\sphinxAtStartPar
The maximum number of elements that can be read from the MEAM library
file is determined at compile time. The default is 8.  If you need
support for more elements, you have to change the the constant ‘MAXELT’
at the beginning of the file \sphinxcode{\sphinxupquote{src/MEAM/meam.h}} and update/recompile
LAMMPS.  There is no limit on the number of atoms types.


\subsection{Related commands}
\label{\detokenize{pair_meam:related-commands}}
\sphinxAtStartPar
{\hyperref[\detokenize{pair_coeff::doc}]{\sphinxcrossref{\DUrole{doc}{pair\_coeff}}}}, {\hyperref[\detokenize{pair_eam::doc}]{\sphinxcrossref{\DUrole{doc}{pair\_style eam}}}},
{\hyperref[\detokenize{pair_meam_spline::doc}]{\sphinxcrossref{\DUrole{doc}{pair\_style meam/spline}}}}


\subsection{Default}
\label{\detokenize{pair_meam:default}}
\sphinxAtStartPar
none


\bigskip\hrule\bigskip

\phantomsection\label{\detokenize{pair_meam:baskes}}
\sphinxAtStartPar
\sphinxstylestrong{(Baskes)} Baskes, Phys Rev B, 46, 2727\sphinxhyphen{}2742 (1992).

\phantomsection\label{\detokenize{pair_meam:baskes2}}
\sphinxAtStartPar
\sphinxstylestrong{(Baskes2)} Baskes, Phys Rev B, 75, 094113 (2007).

\phantomsection\label{\detokenize{pair_meam:gullet}}
\sphinxAtStartPar
\sphinxstylestrong{(Gullet)} Gullet, Wagner, Slepoy, SANDIA Report 2003\sphinxhyphen{}8782 (2003). DOI:10.2172/918395
This report may be accessed on\sphinxhyphen{}line via \sphinxhref{https://download.lammps.org/pdfs/MEAM\_report\_2003.pdf}{this link}.

\phantomsection\label{\detokenize{pair_meam:lee}}
\sphinxAtStartPar
\sphinxstylestrong{(Lee)} Lee, Baskes, Phys. Rev. B, 62, 8564\sphinxhyphen{}8567 (2000).

\phantomsection\label{\detokenize{pair_meam:lee2}}
\sphinxAtStartPar
\sphinxstylestrong{(Lee2)} Lee, Baskes, Kim, Cho.  Phys. Rev. B, 64, 184102 (2001).

\phantomsection\label{\detokenize{pair_meam:valone}}
\sphinxAtStartPar
\sphinxstylestrong{(Valone)} Valone, Baskes, Martin, Phys. Rev. B, 73, 214209 (2006).

\phantomsection\label{\detokenize{pair_meam:wang2}}
\sphinxAtStartPar
\sphinxstylestrong{(Wang)} Wang, Van Hove, Ross, Baskes, J. Chem. Phys., 121, 5410 (2004).

\phantomsection\label{\detokenize{pair_meam:zbl}}
\sphinxAtStartPar
\sphinxstylestrong{(ZBL)} J.F. Ziegler, J.P. Biersack, U. Littmark, “Stopping and Ranges
of Ions in Matter”, Vol 1, 1985, Pergamon Press.

\sphinxstepscope

\index{pair\_style meam/spline@\spxentry{pair\_style meam/spline}}\ignorespaces 
\index{pair\_style meam/spline/omp@\spxentry{pair\_style meam/spline/omp}}\ignorespaces 

\section{pair\_style meam/spline command}
\label{\detokenize{pair_meam_spline:pair-style-meam-spline-command}}\label{\detokenize{pair_meam_spline:index-1}}\label{\detokenize{pair_meam_spline:index-0}}\label{\detokenize{pair_meam_spline::doc}}
\sphinxAtStartPar
Accelerator Variants: \sphinxstyleemphasis{meam/spline/omp}


\subsection{Syntax}
\label{\detokenize{pair_meam_spline:syntax}}
\begin{sphinxVerbatim}[commandchars=\\\{\}]
\PYG{k}{pair\PYGZus{}style}\PYG{+w}{ }\PYG{n}{meam}\PYG{o}{/}\PYG{n}{spline}
\end{sphinxVerbatim}


\subsection{Examples}
\label{\detokenize{pair_meam_spline:examples}}
\begin{sphinxVerbatim}[commandchars=\\\{\}]
\PYG{k}{pair\PYGZus{}style}\PYG{+w}{ }\PYG{n}{meam}\PYG{o}{/}\PYG{n}{spline}
\PYG{k}{pair\PYGZus{}coeff}\PYG{+w}{ }\PYG{o}{*}\PYG{+w}{ }\PYG{o}{*}\PYG{+w}{ }\PYG{n}{Ti.meam.spline}\PYG{+w}{ }\PYG{n}{Ti}
\PYG{k}{pair\PYGZus{}coeff}\PYG{+w}{ }\PYG{o}{*}\PYG{+w}{ }\PYG{o}{*}\PYG{+w}{ }\PYG{n}{Ti.meam.spline}\PYG{+w}{ }\PYG{n}{Ti}\PYG{+w}{ }\PYG{n}{O}
\end{sphinxVerbatim}


\subsection{Description}
\label{\detokenize{pair_meam_spline:description}}
\sphinxAtStartPar
The \sphinxstyleemphasis{meam/spline} style computes pairwise interactions for metals
using a variant of modified embedded\sphinxhyphen{}atom method (MEAM) potentials
{\hyperref[\detokenize{pair_meam_spline:lenosky1}]{\sphinxcrossref{\DUrole{std,std-ref}{(Lenosky)}}}}.  For a single species (“old\sphinxhyphen{}style”) MEAM,
the total energy E is given by
\begin{equation*}
\begin{split}E & =\sum_{i<j}\phi(r_{ij})+\sum_{i}U(n_{i}) \\
n_{i} & =\sum_{j}\rho(r_{ij})+\sum_{\substack{j<k,\\j,k\neq i}}f(r_{ij})f(r_{ik})g[\cos(\theta_{jik})]\end{split}
\end{equation*}
\sphinxAtStartPar
where \(\rho_i\) is the density at atom I, \(\theta_{jik}\) is
the angle between atoms J, I, and K centered on atom I. The five
functions \(\phi, U, \rho, f,\) and \sphinxstyleemphasis{g} are represented by cubic splines.

\sphinxAtStartPar
The \sphinxstyleemphasis{meam/spline} style also supports a new style multicomponent
modified embedded\sphinxhyphen{}atom method (MEAM) potential {\hyperref[\detokenize{pair_meam_spline:zhang4}]{\sphinxcrossref{\DUrole{std,std-ref}{(Zhang)}}}}, where
the total energy E is given by
\begin{equation*}
\begin{split}E &= \sum_{i<j}\phi_{ij}(r_{ij})+\sum_{i}U_i(n_{i}) \\
n_{i} & = \sum_{j\ne i}\rho_j(r_{ij})+\sum_{\substack{j<k,\\j,k\neq i}}f_{j}(r_{ij})f_{k}(r_{ik})g_{jk}[\cos(\theta_{jik})]\end{split}
\end{equation*}
\sphinxAtStartPar
where the five functions \(\phi, U, \rho, f,\) and \sphinxstyleemphasis{g} depend on the
chemistry of the atoms in the interaction.  In particular, if there are
N different chemistries, there are N different \sphinxstyleemphasis{U}, \(\rho\), and
\sphinxstyleemphasis{f} functions, while there are N(N+1)/2 different \(\phi\) and \sphinxstyleemphasis{g}
functions.  The new style multicomponent MEAM potential files are
indicated by the second line in the file starts with “meam/spline”
followed by the number of elements and the name of each element.

\sphinxAtStartPar
The cutoffs and the coefficients for these spline functions are listed
in a parameter file which is specified by the
{\hyperref[\detokenize{pair_coeff::doc}]{\sphinxcrossref{\DUrole{doc}{pair\_coeff}}}} command.  Parameter files for different
elements are included in the “potentials” directory of the LAMMPS
distribution and have a “.meam.spline” file suffix.  All of these
files are parameterized in terms of LAMMPS {\hyperref[\detokenize{units::doc}]{\sphinxcrossref{\DUrole{doc}{metal units}}}}.

\sphinxAtStartPar
Note that unlike for other potentials, cutoffs for spline\sphinxhyphen{}based MEAM
potentials are not set in the pair\_style or pair\_coeff command; they
are specified in the potential files themselves.

\sphinxAtStartPar
Unlike the EAM pair style, which retrieves the atomic mass from the
potential file, the spline\sphinxhyphen{}based MEAM potentials do not include mass
information; thus you need to use the {\hyperref[\detokenize{mass::doc}]{\sphinxcrossref{\DUrole{doc}{mass}}}} command to
specify it.

\sphinxAtStartPar
Only a single pair\_coeff command is used with the \sphinxstyleemphasis{meam/spline} style
which specifies a potential file with parameters for all needed
elements.  These are mapped to LAMMPS atom types by specifying N
additional arguments after the filename in the pair\_coeff command,
where N is the number of LAMMPS atom types:
\begin{itemize}
\item {} 
\sphinxAtStartPar
filename

\item {} 
\sphinxAtStartPar
N element names = mapping of spline\sphinxhyphen{}based MEAM elements to atom types

\end{itemize}

\sphinxAtStartPar
See the {\hyperref[\detokenize{pair_coeff::doc}]{\sphinxcrossref{\DUrole{doc}{pair\_coeff}}}} page for alternate ways
to specify the path for the potential file.

\sphinxAtStartPar
As an example, imagine the Ti.meam.spline file has values for Ti (old style).
In that case your LAMMPS simulation may only have one atom type which has
to be mapped to the Ti element as follows:

\begin{sphinxVerbatim}[commandchars=\\\{\}]
\PYG{k}{pair\PYGZus{}coeff}\PYG{+w}{ }\PYG{o}{*}\PYG{+w}{ }\PYG{o}{*}\PYG{+w}{ }\PYG{n}{Ti.meam.spline}\PYG{+w}{ }\PYG{n}{Ti}
\end{sphinxVerbatim}

\sphinxAtStartPar
The first 2 arguments must be * * and there may be only one element
following or NULL. Systems where there would be multiple atom types
assigned to the same element are \sphinxstylestrong{not} supported by this pair style
due to limitations in its implementation.  If a mapping value is
specified as NULL, the mapping is not performed.  This can be used when
a \sphinxstyleemphasis{meam/spline} potential is used as part of the \sphinxstyleemphasis{hybrid} pair style.
The NULL values are placeholders for atom types that will be used with
other potentials.

\sphinxAtStartPar
An example with a two component spline (new style) is TiO.meam.spline, where
the command

\begin{sphinxVerbatim}[commandchars=\\\{\}]
\PYG{k}{pair\PYGZus{}coeff}\PYG{+w}{ }\PYG{o}{*}\PYG{+w}{ }\PYG{o}{*}\PYG{+w}{ }\PYG{n}{TiO.meam.spline}\PYG{+w}{ }\PYG{n}{Ti}\PYG{+w}{ }\PYG{n}{O}
\end{sphinxVerbatim}

\sphinxAtStartPar
will map the first atom type to Ti and the second atom type to O. Note
in this case that the species names need to match exactly with the
names of the elements in the TiO.meam.spline file; otherwise an
error will be raised. This behavior is different than the old style
MEAM files.


\bigskip\hrule\bigskip


\sphinxAtStartPar
Styles with a \sphinxstyleemphasis{gpu}, \sphinxstyleemphasis{intel}, \sphinxstyleemphasis{kk}, \sphinxstyleemphasis{omp}, or \sphinxstyleemphasis{opt} suffix are
functionally the same as the corresponding style without the suffix.
They have been optimized to run faster, depending on your available
hardware, as discussed on the {\hyperref[\detokenize{Speed_packages::doc}]{\sphinxcrossref{\DUrole{doc}{Accelerator packages}}}}
page.  The accelerated styles take the same arguments and should
produce the same results, except for round\sphinxhyphen{}off and precision issues.

\sphinxAtStartPar
These accelerated styles are part of the GPU, INTEL, KOKKOS,
OPENMP, and OPT packages, respectively.  They are only enabled if
LAMMPS was built with those packages.  See the {\hyperref[\detokenize{Build_package::doc}]{\sphinxcrossref{\DUrole{doc}{Build package}}}} page for more info.

\sphinxAtStartPar
You can specify the accelerated styles explicitly in your input script
by including their suffix, or you can use the {\hyperref[\detokenize{Run_options::doc}]{\sphinxcrossref{\DUrole{doc}{\sphinxhyphen{}suffix command\sphinxhyphen{}line switch}}}} when you invoke LAMMPS, or you can use the
{\hyperref[\detokenize{suffix::doc}]{\sphinxcrossref{\DUrole{doc}{suffix}}}} command in your input script.

\sphinxAtStartPar
See the {\hyperref[\detokenize{Speed_packages::doc}]{\sphinxcrossref{\DUrole{doc}{Accelerator packages}}}} page for more
instructions on how to use the accelerated styles effectively.


\bigskip\hrule\bigskip



\subsection{Mixing, shift, table, tail correction, restart, rRESPA info}
\label{\detokenize{pair_meam_spline:mixing-shift-table-tail-correction-restart-rrespa-info}}
\sphinxAtStartPar
This pair style does not support the {\hyperref[\detokenize{pair_modify::doc}]{\sphinxcrossref{\DUrole{doc}{pair\_modify}}}}
shift, table, and tail options.

\sphinxAtStartPar
The \sphinxstyleemphasis{meam/spline} pair style does not write its information to {\hyperref[\detokenize{restart::doc}]{\sphinxcrossref{\DUrole{doc}{binary restart files}}}}, since it is stored in an external
potential parameter file.  Thus, you need to re\sphinxhyphen{}specify the pair\_style
and pair\_coeff commands in an input script that reads a restart file.

\sphinxAtStartPar
The \sphinxstyleemphasis{meam/spline} pair style can only be used via the \sphinxstyleemphasis{pair} keyword of the
{\hyperref[\detokenize{run_style::doc}]{\sphinxcrossref{\DUrole{doc}{run\_style respa}}}} command.  They do not support the
\sphinxstyleemphasis{inner}, \sphinxstyleemphasis{middle}, \sphinxstyleemphasis{outer} keywords.


\bigskip\hrule\bigskip



\subsection{Restrictions}
\label{\detokenize{pair_meam_spline:restrictions}}
\sphinxAtStartPar
This pair style requires the {\hyperref[\detokenize{newton::doc}]{\sphinxcrossref{\DUrole{doc}{newton}}}} setting to be “on”
for pair interactions.

\sphinxAtStartPar
This pair style does not support mapping multiple atom types to the same element.

\sphinxAtStartPar
This pair style is only enabled if LAMMPS was built with the MANYBODY
package.  See the {\hyperref[\detokenize{Build_package::doc}]{\sphinxcrossref{\DUrole{doc}{Build package}}}} page for more
info.


\subsection{Related commands}
\label{\detokenize{pair_meam_spline:related-commands}}
\sphinxAtStartPar
{\hyperref[\detokenize{pair_coeff::doc}]{\sphinxcrossref{\DUrole{doc}{pair\_coeff}}}}, {\hyperref[\detokenize{pair_meam::doc}]{\sphinxcrossref{\DUrole{doc}{pair\_style meam}}}}


\subsection{Default}
\label{\detokenize{pair_meam_spline:default}}
\sphinxAtStartPar
none


\bigskip\hrule\bigskip

\phantomsection\label{\detokenize{pair_meam_spline:lenosky1}}
\sphinxAtStartPar
\sphinxstylestrong{(Lenosky)} Lenosky, Sadigh, Alonso, Bulatov, de la Rubia, Kim, Voter,
Kress, Modelling Simulation Materials Science Engineering, 8, 825
(2000).

\phantomsection\label{\detokenize{pair_meam_spline:zhang4}}
\sphinxAtStartPar
\sphinxstylestrong{(Zhang)} Zhang and Trinkle, Computational Materials Science, 124, 204\sphinxhyphen{}210 (2016).

\sphinxstepscope

\index{pair\_style meam/sw/spline@\spxentry{pair\_style meam/sw/spline}}\ignorespaces 

\section{pair\_style meam/sw/spline command}
\label{\detokenize{pair_meam_sw_spline:pair-style-meam-sw-spline-command}}\label{\detokenize{pair_meam_sw_spline:index-0}}\label{\detokenize{pair_meam_sw_spline::doc}}

\subsection{Syntax}
\label{\detokenize{pair_meam_sw_spline:syntax}}
\begin{sphinxVerbatim}[commandchars=\\\{\}]
\PYG{k}{pair\PYGZus{}style}\PYG{+w}{ }\PYG{n}{meam}\PYG{o}{/}\PYG{n}{sw}\PYG{o}{/}\PYG{n}{spline}
\end{sphinxVerbatim}


\subsection{Examples}
\label{\detokenize{pair_meam_sw_spline:examples}}
\begin{sphinxVerbatim}[commandchars=\\\{\}]
\PYG{k}{pair\PYGZus{}style}\PYG{+w}{ }\PYG{n}{meam}\PYG{o}{/}\PYG{n}{sw}\PYG{o}{/}\PYG{n}{spline}
\PYG{k}{pair\PYGZus{}coeff}\PYG{+w}{ }\PYG{o}{*}\PYG{+w}{ }\PYG{o}{*}\PYG{+w}{ }\PYG{n}{Ti.meam.sw.spline}\PYG{+w}{ }\PYG{n}{Ti}
\PYG{k}{pair\PYGZus{}coeff}\PYG{+w}{ }\PYG{o}{*}\PYG{+w}{ }\PYG{o}{*}\PYG{+w}{ }\PYG{n}{Ti.meam.sw.spline}\PYG{+w}{ }\PYG{n}{Ti}\PYG{+w}{ }\PYG{n}{Ti}\PYG{+w}{ }\PYG{n}{Ti}
\end{sphinxVerbatim}


\subsection{Description}
\label{\detokenize{pair_meam_sw_spline:description}}
\sphinxAtStartPar
The \sphinxstyleemphasis{meam/sw/spline} style computes pairwise interactions for metals
using a variant of modified embedded\sphinxhyphen{}atom method (MEAM) potentials
{\hyperref[\detokenize{pair_meam_sw_spline:lenosky2}]{\sphinxcrossref{\DUrole{std,std-ref}{(Lenosky)}}}} with an additional Stillinger\sphinxhyphen{}Weber (SW) term
{\hyperref[\detokenize{pair_meam_sw_spline:stillinger1}]{\sphinxcrossref{\DUrole{std,std-ref}{(Stillinger)}}}} in the energy.  This form of the potential
was first proposed by Nicklas, Fellinger, and Park
{\hyperref[\detokenize{pair_meam_sw_spline:nicklas}]{\sphinxcrossref{\DUrole{std,std-ref}{(Nicklas)}}}}.  We refer to it as MEAM+SW.  The total energy E
is given by
\begin{equation*}
\begin{split}E & = E_{MEAM} + E_{SW} \\
E_{MEAM} & =  \sum _{IJ} \phi (r_{IJ}) + \sum _{I} U(\rho _I) \\
E_{SW} & =  \sum _{I} \sum _{JK} F(r_{IJ}) \, F(r_{IK}) \, G(\cos(\theta _{JIK})) \\
\rho _I & = \sum _J \rho(r_{IJ}) + \sum _{JK} f(r_{IJ}) \, f(r_{IK}) \, g(\cos(\theta _{JIK}))\end{split}
\end{equation*}
\sphinxAtStartPar
where \(\rho_I\) is the density at atom I, \(\theta_{JIK}\) is
the angle between atoms J, I, and K centered on atom I. The seven
functions \(\phi, F, G, U, \rho, f,\) and \sphinxstyleemphasis{g} are represented by
cubic splines.

\sphinxAtStartPar
The cutoffs and the coefficients for these spline functions are listed
in a parameter file which is specified by the
{\hyperref[\detokenize{pair_coeff::doc}]{\sphinxcrossref{\DUrole{doc}{pair\_coeff}}}} command.  Parameter files for different
elements are included in the “potentials” directory of the LAMMPS
distribution and have a “.meam.sw.spline” file suffix.  All of these
files are parameterized in terms of LAMMPS {\hyperref[\detokenize{units::doc}]{\sphinxcrossref{\DUrole{doc}{metal units}}}}.

\sphinxAtStartPar
Note that unlike for other potentials, cutoffs for spline\sphinxhyphen{}based
MEAM+SW potentials are not set in the pair\_style or pair\_coeff
command; they are specified in the potential files themselves.

\sphinxAtStartPar
Unlike the EAM pair style, which retrieves the atomic mass from the
potential file, the spline\sphinxhyphen{}based MEAM+SW potentials do not include
mass information; thus you need to use the {\hyperref[\detokenize{mass::doc}]{\sphinxcrossref{\DUrole{doc}{mass}}}} command to
specify it.

\sphinxAtStartPar
Only a single pair\_coeff command is used with the meam/sw/spline style
which specifies a potential file with parameters for all needed
elements.  These are mapped to LAMMPS atom types by specifying N
additional arguments after the filename in the pair\_coeff command,
where N is the number of LAMMPS atom types:
\begin{itemize}
\item {} 
\sphinxAtStartPar
filename

\item {} 
\sphinxAtStartPar
N element names = mapping of spline\sphinxhyphen{}based MEAM+SW elements to atom types

\end{itemize}

\sphinxAtStartPar
See the {\hyperref[\detokenize{pair_coeff::doc}]{\sphinxcrossref{\DUrole{doc}{pair\_coeff}}}} page for alternate ways
to specify the path for the potential file.

\sphinxAtStartPar
As an example, imagine the Ti.meam.sw.spline file has values for Ti.
If your LAMMPS simulation has 3 atoms types and they are all to be
treated with this potential, you would use the following pair\_coeff
command:

\begin{sphinxVerbatim}[commandchars=\\\{\}]
\PYG{k}{pair\PYGZus{}coeff}\PYG{+w}{ }\PYG{o}{*}\PYG{+w}{ }\PYG{o}{*}\PYG{+w}{ }\PYG{n}{Ti.meam.sw.spline}\PYG{+w}{ }\PYG{n}{Ti}\PYG{+w}{ }\PYG{n}{Ti}\PYG{+w}{ }\PYG{n}{Ti}
\end{sphinxVerbatim}

\sphinxAtStartPar
The first 2 arguments must be * * so as to span all LAMMPS atom types.
The three Ti arguments map LAMMPS atom types 1,2,3 to the Ti element
in the potential file. If a mapping value is specified as NULL, the
mapping is not performed. This can be used when a \sphinxstyleemphasis{meam/sw/spline}
potential is used as part of the hybrid pair style. The NULL values
are placeholders for atom types that will be used with other
potentials.

\begin{sphinxadmonition}{note}{Note:}
\sphinxAtStartPar
The \sphinxstyleemphasis{meam/sw/spline} style currently supports only
single\sphinxhyphen{}element MEAM+SW potentials.  It may be extended for alloy
systems in the future.
\end{sphinxadmonition}

\sphinxAtStartPar
Example input scripts that use this pair style are provided
in the examples/PACKAGES/meam\_sw\_spline directory.


\bigskip\hrule\bigskip



\subsection{Mixing, shift, table, tail correction, restart, rRESPA info}
\label{\detokenize{pair_meam_sw_spline:mixing-shift-table-tail-correction-restart-rrespa-info}}
\sphinxAtStartPar
The pair style does not support multiple element types or mixing.
It has been designed for pure elements only.

\sphinxAtStartPar
This pair style does not support the {\hyperref[\detokenize{pair_modify::doc}]{\sphinxcrossref{\DUrole{doc}{pair\_modify}}}}
shift, table, and tail options.

\sphinxAtStartPar
The \sphinxstyleemphasis{meam/sw/spline} pair style does not write its information to
{\hyperref[\detokenize{restart::doc}]{\sphinxcrossref{\DUrole{doc}{binary restart files}}}}, since it is stored in an external
potential parameter file.  Thus, you need to re\sphinxhyphen{}specify the pair\_style
and pair\_coeff commands in an input script that reads a restart file.

\sphinxAtStartPar
The \sphinxstyleemphasis{meam/sw/spline} pair style can only be used via the \sphinxstyleemphasis{pair}
keyword of the {\hyperref[\detokenize{run_style::doc}]{\sphinxcrossref{\DUrole{doc}{run\_style respa}}}} command.  They do not
support the \sphinxstyleemphasis{inner}, \sphinxstyleemphasis{middle}, \sphinxstyleemphasis{outer} keywords.


\bigskip\hrule\bigskip



\subsection{Restrictions}
\label{\detokenize{pair_meam_sw_spline:restrictions}}
\sphinxAtStartPar
This pair style requires the {\hyperref[\detokenize{newton::doc}]{\sphinxcrossref{\DUrole{doc}{newton}}}} setting to be “on”
for pair interactions.

\sphinxAtStartPar
This pair style is only enabled if LAMMPS was built with the MANYBODY
package.  See the {\hyperref[\detokenize{Build_package::doc}]{\sphinxcrossref{\DUrole{doc}{Build package}}}} page for more
info.


\subsection{Related commands}
\label{\detokenize{pair_meam_sw_spline:related-commands}}
\sphinxAtStartPar
{\hyperref[\detokenize{pair_coeff::doc}]{\sphinxcrossref{\DUrole{doc}{pair\_coeff}}}}, {\hyperref[\detokenize{pair_meam::doc}]{\sphinxcrossref{\DUrole{doc}{pair\_style meam}}}},
{\hyperref[\detokenize{pair_meam_spline::doc}]{\sphinxcrossref{\DUrole{doc}{pair\_style meam/spline}}}}


\subsection{Default}
\label{\detokenize{pair_meam_sw_spline:default}}
\sphinxAtStartPar
none


\bigskip\hrule\bigskip

\phantomsection\label{\detokenize{pair_meam_sw_spline:lenosky2}}
\sphinxAtStartPar
\sphinxstylestrong{(Lenosky)} Lenosky, Sadigh, Alonso, Bulatov, de la Rubia, Kim, Voter,
Kress, Modell. Simul. Mater. Sci. Eng. 8, 825 (2000).

\phantomsection\label{\detokenize{pair_meam_sw_spline:stillinger1}}
\sphinxAtStartPar
\sphinxstylestrong{(Stillinger)} Stillinger, Weber, Phys. Rev. B 31, 5262 (1985).

\phantomsection\label{\detokenize{pair_meam_sw_spline:nicklas}}
\sphinxAtStartPar
\sphinxstylestrong{(Nicklas)}
The spline\sphinxhyphen{}based MEAM+SW format was first devised and used to develop
potentials for bcc transition metals by Jeremy Nicklas, Michael Fellinger,
and Hyoungki Park at The Ohio State University.

\sphinxstepscope

\index{pair\_style mesocnt@\spxentry{pair\_style mesocnt}}\ignorespaces 
\index{pair\_style mesocnt/viscous@\spxentry{pair\_style mesocnt/viscous}}\ignorespaces 

\section{pair\_style mesocnt command}
\label{\detokenize{pair_mesocnt:pair-style-mesocnt-command}}\label{\detokenize{pair_mesocnt:index-1}}\label{\detokenize{pair_mesocnt:index-0}}\label{\detokenize{pair_mesocnt::doc}}

\section{pair\_style mesocnt/viscous command}
\label{\detokenize{pair_mesocnt:pair-style-mesocnt-viscous-command}}

\subsection{Syntax}
\label{\detokenize{pair_mesocnt:syntax}}
\begin{sphinxVerbatim}[commandchars=\\\{\}]
\PYG{k}{pair\PYGZus{}style}\PYG{+w}{ }\PYG{n}{style}\PYG{+w}{ }\PYG{n}{neigh\PYGZus{}cutoff}\PYG{+w}{ }\PYG{n}{mode}\PYG{+w}{ }\PYG{n}{neigh\PYGZus{}mode}
\end{sphinxVerbatim}
\begin{itemize}
\item {} 
\sphinxAtStartPar
style = \sphinxstyleemphasis{mesocnt} or \sphinxstyleemphasis{mesocnt/viscous}

\item {} 
\sphinxAtStartPar
neigh\_cutoff = neighbor list cutoff (distance units)

\item {} 
\sphinxAtStartPar
mode = \sphinxstyleemphasis{chain} or \sphinxstyleemphasis{segment} (optional)

\item {} 
\sphinxAtStartPar
neigh\_mode = \sphinxstyleemphasis{id} or \sphinxstyleemphasis{topology} (optional)

\end{itemize}


\subsection{Examples}
\label{\detokenize{pair_mesocnt:examples}}
\begin{sphinxVerbatim}[commandchars=\\\{\}]
\PYG{k}{pair\PYGZus{}style}\PYG{+w}{ }\PYG{n}{mesocnt}\PYG{+w}{ }\PYG{l+m}{30.0}
\PYG{k}{pair\PYGZus{}coeff}\PYG{+w}{ }\PYG{o}{*}\PYG{+w}{ }\PYG{o}{*}\PYG{+w}{ }\PYG{n}{C\PYGZus{}10\PYGZus{}10.mesocnt}\PYG{+w}{ }\PYG{l+m}{2}

\PYG{k}{pair\PYGZus{}style}\PYG{+w}{ }\PYG{n}{mesocnt}\PYG{o}{/}\PYG{n}{viscous}\PYG{+w}{ }\PYG{l+m}{60.0}\PYG{+w}{ }\PYG{n}{chain}\PYG{+w}{ }\PYG{n}{topology}
\PYG{k}{pair\PYGZus{}coeff}\PYG{+w}{ }\PYG{o}{*}\PYG{+w}{ }\PYG{o}{*}\PYG{+w}{ }\PYG{n}{C\PYGZus{}10\PYGZus{}10.mesocnt}\PYG{+w}{ }\PYG{l+m}{0.001}\PYG{+w}{ }\PYG{l+m}{20.0}\PYG{+w}{ }\PYG{l+m}{0.2}\PYG{+w}{ }\PYG{l+m}{2}\PYG{+w}{ }\PYG{l+m}{4}
\end{sphinxVerbatim}


\subsection{Description}
\label{\detokenize{pair_mesocnt:description}}
\sphinxAtStartPar
Style \sphinxstyleemphasis{mesocnt} implements a mesoscopic potential for the interaction
of carbon nanotubes (CNTs), or other quasi\sphinxhyphen{}1D objects such as other
kinds of nanotubes or nanowires. In this potential, CNTs are modelled
as chains of cylindrical segments in which each infinitesimal surface
element interacts with all other CNT surface elements with the
Lennard\sphinxhyphen{}Jones (LJ) term adopted from the {\hyperref[\detokenize{pair_airebo::doc}]{\sphinxcrossref{\DUrole{doc}{airebo}}}}
style. The interaction energy is then computed by integrating over the
surfaces of all interacting CNTs.

\sphinxAtStartPar
In LAMMPS, cylindrical segments are represented by bonds. Each segment
is defined by its two end points (“nodes”) which correspond to atoms
in LAMMPS. For the exact functional form of the potential and
implementation details, the reader is referred to the original papers
{\hyperref[\detokenize{pair_mesocnt:volkov1}]{\sphinxcrossref{\DUrole{std,std-ref}{(Volkov1)}}}} and {\hyperref[\detokenize{pair_mesocnt:volkov2}]{\sphinxcrossref{\DUrole{std,std-ref}{(Volkov2)}}}}.

\sphinxAtStartPar
\DUrole{versionmodified,changed}{Changed in version 15Sep2022.}

\sphinxAtStartPar
The potential supports two modes, \sphinxstyleemphasis{segment} and \sphinxstyleemphasis{chain}. By default,
\sphinxstyleemphasis{chain} mode is enabled.  In \sphinxstyleemphasis{segment} mode, interactions are
pair\sphinxhyphen{}wise between all neighboring segments based on a segment\sphinxhyphen{}segment
approach (keyword \sphinxstyleemphasis{segment} in pair\_style command).  In \sphinxstyleemphasis{chain} mode,
interactions are calculated between each segment and infinitely or
semi\sphinxhyphen{}infinitely long CNTs as described in {\hyperref[\detokenize{pair_mesocnt:volkov1}]{\sphinxcrossref{\DUrole{std,std-ref}{(Volkov1)}}}}.
Chains of segments are converted to these (semi\sphinxhyphen{})infinite CNTs bases
on an approximate chain approach outlined in {\hyperref[\detokenize{pair_mesocnt:volkov2}]{\sphinxcrossref{\DUrole{std,std-ref}{(Volkov2)}}}}. Hence, interactions are calculated on a segment\sphinxhyphen{}chain
basis (keyword \sphinxstyleemphasis{chain} in the pair\_style command).  Using \sphinxstyleemphasis{chain} mode
allows to simplify the computation of the interactions significantly
and reduces the computational times to the same order of magnitude as
for regular bead spring models where beads interact with the standard
{\hyperref[\detokenize{pair_lj::doc}]{\sphinxcrossref{\DUrole{doc}{pair\_lj/cut}}}} potential. However, this method is only
valid when the curvature of the CNTs in the system is small.  When
CNTs are buckled (see {\hyperref[\detokenize{angle_mesocnt::doc}]{\sphinxcrossref{\DUrole{doc}{angle\_mesocnt}}}}), local
curvature can be very high and the pair\_style automatically switches
to \sphinxstyleemphasis{segment} mode for interactions involving buckled CNTs.

\sphinxAtStartPar
The potential further implements two different neighbor list
construction modes. Mode \sphinxstyleemphasis{id} uses atom and mol IDs to construct
neighbor lists while \sphinxstyleemphasis{topology} modes uses only the bond topology of
the system. While \sphinxstyleemphasis{id} mode requires bonded atoms to have consecutive
LAMMPS atom IDs and atoms in different CNTs to have different LAMMPS
molecule IDs, \sphinxstyleemphasis{topology} mode has no such requirement. Using \sphinxstyleemphasis{id} mode
is faster and is enabled by default.

\begin{sphinxadmonition}{note}{Note:}
\sphinxAtStartPar
Neighbor \sphinxstyleemphasis{id} mode requires all CNTs in the system to have distinct
LAMMPS molecule IDs and bonded atoms to have consecutive LAMMPS atom
IDs. If this is not possible (e.g. in simulations of CNT rings),
\sphinxstyleemphasis{topology} mode needs to be enabled in the pair\_style command.
\end{sphinxadmonition}

\sphinxAtStartPar
\DUrole{versionmodified,added}{Added in version 15Sep2022.}

\sphinxAtStartPar
In addition to the LJ interactions described above, style
\sphinxstyleemphasis{mesocnt/viscous} explicitly models friction between neighboring
segments. Friction forces are a function of the relative velocity
between a segment and its neighboring approximate chain (even in
\sphinxstyleemphasis{segment} mode) and only act along the axes of the interacting segment
and chain. In this potential, friction forces acting per unit length
of a nanotube segment are modelled as a shifted logistic function:
\begin{equation*}
\begin{split}F^{\text{FRICTION}}(v) / L = \frac{F^{\text{max}}}{1 +
\exp(-k(v-v_0))} - \frac{F^{\text{max}}}{1 + \exp(k v_0)}\end{split}
\end{equation*}

\bigskip\hrule\bigskip


\sphinxAtStartPar
In the pair\_style command, the modes described above can be toggled
using the \sphinxstyleemphasis{segment} or \sphinxstyleemphasis{chain} keywords.  The neighbor list cutoff
defines the cutoff within which atoms are included in the neighbor
list for constructing neighboring CNT chains.  This is different from
the potential cutoff, which is directly calculated from parameters
specified in the potential file. We recommend using a neighbor list
cutoff of at least 3 times the maximum segment length used in the
simulation to ensure proper neighbor chain construction.

\begin{sphinxadmonition}{note}{Note:}
\sphinxAtStartPar
CNT ends are treated differently by all \sphinxstyleemphasis{mesocnt} styles. Atoms on
CNT ends need to be assigned different LAMMPS atom types than atoms
not on CNT ends.
\end{sphinxadmonition}

\sphinxAtStartPar
Style \sphinxstyleemphasis{mesocnt} requires tabulated data provided in a single ASCII
text file, as well as a list of integers corresponding to all LAMMPS
atom types representing CNT ends:
\begin{itemize}
\item {} 
\sphinxAtStartPar
filename

\item {} 
\sphinxAtStartPar
\(N\) CNT end atom types

\end{itemize}

\sphinxAtStartPar
For example, if your LAMMPS simulation of (10, 10) nanotubes has 4
atom types where atom types 1 and 3 are assigned to ‘inner’ nodes and
atom types 2 and 4 are assigned to CNT end nodes, the pair\_coeff
command would be:

\begin{sphinxVerbatim}[commandchars=\\\{\}]
\PYG{k}{pair\PYGZus{}coeff}\PYG{+w}{ }\PYG{o}{*}\PYG{+w}{ }\PYG{o}{*}\PYG{+w}{ }\PYG{n}{C\PYGZus{}10\PYGZus{}10.mesocnt}\PYG{+w}{ }\PYG{l+m}{2}\PYG{+w}{ }\PYG{l+m}{4}
\end{sphinxVerbatim}

\sphinxAtStartPar
Likewise, style \sphinxstyleemphasis{mesocnt/viscous} also requires the same information
as style \sphinxstyleemphasis{mesocnt}, with the addition of 3 parameters for the viscous
friction forces as listed above:
\begin{itemize}
\item {} 
\sphinxAtStartPar
filename

\item {} 
\sphinxAtStartPar
\(F^{\text{max}}\)

\item {} 
\sphinxAtStartPar
\(k\)

\item {} 
\sphinxAtStartPar
\(v_0\)

\item {} 
\sphinxAtStartPar
\(N\) CNT end atom types

\end{itemize}

\sphinxAtStartPar
Using the same example system as with style \sphinxstyleemphasis{mesocnt} with the
addition of friction, the pair\_coeff command is:

\begin{sphinxVerbatim}[commandchars=\\\{\}]
\PYG{k}{pair\PYGZus{}coeff}\PYG{+w}{ }\PYG{o}{*}\PYG{+w}{ }\PYG{o}{*}\PYG{+w}{ }\PYG{n}{C\PYGZus{}10\PYGZus{}10.mesocnt}\PYG{+w}{ }\PYG{l+m}{0.03}\PYG{+w}{ }\PYG{l+m}{20.0}\PYG{+w}{ }\PYG{l+m}{0.20}\PYG{+w}{ }\PYG{l+m}{2}\PYG{+w}{ }\PYG{l+m}{4}
\end{sphinxVerbatim}

\sphinxAtStartPar
Potential files for CNTs can be readily generated using the freely
available code provided on

\begin{sphinxVerbatim}[commandchars=\\\{\}]
https://github.com/phankl/cntpot
\end{sphinxVerbatim}

\sphinxAtStartPar
Using the same approach, it should also be possible to generate
potential files for other 1D systems mentioned above.

\begin{sphinxadmonition}{note}{Note:}
\sphinxAtStartPar
Because of their size, \sphinxstyleemphasis{mesocnt} style potential files are not
bundled with LAMMPS.  When compiling LAMMPS from source code, the
file \sphinxcode{\sphinxupquote{C\_10\_10.mesocnt}} should be downloaded separately from
\sphinxurl{https://download.lammps.org/potentials/C\_10\_10.mesocnt}

\sphinxAtStartPar
The first line of the potential file provides a time stamp and
general information. The second line lists four integers giving the
number of data points provided in the subsequent four data
tables. The third line lists four floating point numbers: the CNT
radius R, the LJ parameter sigma and two numerical parameters
delta1 and delta2. These four parameters are given in
Angstroms. This is followed by four data tables each separated by a
single empty line. The first two tables have two columns and list
the parameters uInfParallel and Gamma respectively.  The last two
tables have three columns giving data on a quadratic array and list
the parameters Phi and uSemiParallel respectively.  uInfParallel
and uSemiParallel are given in eV/Angstrom, Phi is given in eV and
Gamma is unitless.

\sphinxAtStartPar
If a simulation produces many warnings about segment\sphinxhyphen{}chain
interactions falling outside the interpolation range, we recommend
generating a potential file with lower values of delta1 and delta2.
\end{sphinxadmonition}


\bigskip\hrule\bigskip



\subsection{Mixing, shift, table, tail correction, restart, rRESPA info}
\label{\detokenize{pair_mesocnt:mixing-shift-table-tail-correction-restart-rrespa-info}}
\sphinxAtStartPar
These pair styles does not support mixing.

\sphinxAtStartPar
These pair styles does not support the {\hyperref[\detokenize{pair_modify::doc}]{\sphinxcrossref{\DUrole{doc}{pair\_modify}}}} shift, table, and tail options.

\sphinxAtStartPar
These pair styles do not write their information to {\hyperref[\detokenize{restart::doc}]{\sphinxcrossref{\DUrole{doc}{binary
restart files}}}}, since it is stored in tabulated potential
files.  Thus, you need to re\sphinxhyphen{}specify the pair\_style and pair\_coeff
commands in an input script that reads a restart file.

\sphinxAtStartPar
These pair styles can only be used via the \sphinxstyleemphasis{pair} keyword of the
{\hyperref[\detokenize{run_style::doc}]{\sphinxcrossref{\DUrole{doc}{run\_style respa}}}} command.  They do not support the
\sphinxstyleemphasis{inner}, \sphinxstyleemphasis{middle}, \sphinxstyleemphasis{outer} keywords.


\bigskip\hrule\bigskip



\subsection{Restrictions}
\label{\detokenize{pair_mesocnt:restrictions}}
\sphinxAtStartPar
These styles are part of the MESONT package.  They are only enabled if
LAMMPS was built with that package.  See the {\hyperref[\detokenize{Build_package::doc}]{\sphinxcrossref{\DUrole{doc}{Build package}}}} page for more info.

\sphinxAtStartPar
These pair styles require the {\hyperref[\detokenize{newton::doc}]{\sphinxcrossref{\DUrole{doc}{newton}}}} setting to be
“on” for pair interactions.

\sphinxAtStartPar
These pair styles require all 3 {\hyperref[\detokenize{special_bonds::doc}]{\sphinxcrossref{\DUrole{doc}{special\_bonds lj}}}}
settings to be non\sphinxhyphen{}zero for proper neighbor list construction.

\sphinxAtStartPar
Pair style \sphinxstyleemphasis{mesocnt/viscous} requires you to use the {\hyperref[\detokenize{comm_modify::doc}]{\sphinxcrossref{\DUrole{doc}{comm\_modify
vel yes}}}} command so that velocities are stored by ghost
atoms.


\subsection{Related commands}
\label{\detokenize{pair_mesocnt:related-commands}}
\sphinxAtStartPar
{\hyperref[\detokenize{pair_coeff::doc}]{\sphinxcrossref{\DUrole{doc}{pair\_coeff}}}},
{\hyperref[\detokenize{bond_mesocnt::doc}]{\sphinxcrossref{\DUrole{doc}{bond\_style mesocnt}}}},
{\hyperref[\detokenize{angle_mesocnt::doc}]{\sphinxcrossref{\DUrole{doc}{angle\_style mesocnt}}}}


\subsection{Default}
\label{\detokenize{pair_mesocnt:default}}
\sphinxAtStartPar
mode = chain, neigh\_mode = id


\bigskip\hrule\bigskip

\phantomsection\label{\detokenize{pair_mesocnt:volkov1}}
\sphinxAtStartPar
\sphinxstylestrong{(Volkov1)} Volkov and Zhigilei, J Phys Chem C, 114, 5513 (2010).

\phantomsection\label{\detokenize{pair_mesocnt:volkov2}}
\sphinxAtStartPar
\sphinxstylestrong{(Volkov2)} Volkov, Simov and Zhigilei, APS Meeting Abstracts,
Q31.013 (2008).

\sphinxstepscope

\index{pair\_style edpd@\spxentry{pair\_style edpd}}\ignorespaces 
\index{pair\_style edpd/gpu@\spxentry{pair\_style edpd/gpu}}\ignorespaces 
\index{pair\_style mdpd@\spxentry{pair\_style mdpd}}\ignorespaces 
\index{pair\_style mdpd/gpu@\spxentry{pair\_style mdpd/gpu}}\ignorespaces 
\index{pair\_style mdpd/rhosum@\spxentry{pair\_style mdpd/rhosum}}\ignorespaces 
\index{pair\_style tdpd@\spxentry{pair\_style tdpd}}\ignorespaces 

\section{pair\_style edpd command}
\label{\detokenize{pair_mesodpd:pair-style-edpd-command}}\label{\detokenize{pair_mesodpd:index-5}}\label{\detokenize{pair_mesodpd:index-4}}\label{\detokenize{pair_mesodpd:index-3}}\label{\detokenize{pair_mesodpd:index-2}}\label{\detokenize{pair_mesodpd:index-1}}\label{\detokenize{pair_mesodpd:index-0}}\label{\detokenize{pair_mesodpd::doc}}
\sphinxAtStartPar
Accelerator Variants: \sphinxstyleemphasis{edpd/gpu}


\section{pair\_style mdpd command}
\label{\detokenize{pair_mesodpd:pair-style-mdpd-command}}
\sphinxAtStartPar
Accelerator Variants: \sphinxstyleemphasis{mdpd/gpu}


\section{pair\_style mdpd/rhosum command}
\label{\detokenize{pair_mesodpd:pair-style-mdpd-rhosum-command}}

\section{pair\_style tdpd command}
\label{\detokenize{pair_mesodpd:pair-style-tdpd-command}}

\subsection{Syntax}
\label{\detokenize{pair_mesodpd:syntax}}
\begin{sphinxVerbatim}[commandchars=\\\{\}]
\PYG{k}{pair\PYGZus{}style}\PYG{+w}{ }\PYG{n}{style}\PYG{+w}{ }\PYG{n}{args}
\end{sphinxVerbatim}
\begin{itemize}
\item {} 
\sphinxAtStartPar
style = \sphinxstyleemphasis{edpd} or \sphinxstyleemphasis{mdpd} or \sphinxstyleemphasis{mdpd/rhosum} or \sphinxstyleemphasis{tdpd}

\item {} 
\sphinxAtStartPar
args = list of arguments for a particular style
\begin{sphinxalltt}
\sphinxstyleemphasis{edpd} args = cutoff seed
  cutoff = global cutoff for eDPD interactions (distance units)
  seed = random \# seed (integer) (if \textless{}= 0, eDPD will use current time as the seed)
\sphinxstyleemphasis{mdpd} args = T cutoff seed
  T = temperature (temperature units)
  cutoff = global cutoff for mDPD interactions (distance units)
  seed = random \# seed (integer) (if \textless{}= 0, mDPD will use current time as the seed)
\sphinxstyleemphasis{mdpd/rhosum} args =
\sphinxstyleemphasis{tdpd} args = T cutoff seed
  T = temperature (temperature units)
  cutoff = global cutoff for tDPD interactions (distance units)
  seed = random \# seed (integer) (if \textless{}= 0, tDPD will use current time as the seed)
\end{sphinxalltt}

\end{itemize}


\subsection{Examples}
\label{\detokenize{pair_mesodpd:examples}}
\begin{sphinxVerbatim}[commandchars=\\\{\}]
\PYG{k}{pair\PYGZus{}style}\PYG{+w}{ }\PYG{n}{edpd}\PYG{+w}{ }\PYG{l+m}{1.58}\PYG{+w}{ }\PYG{l+m}{9872598}
\PYG{k}{pair\PYGZus{}coeff}\PYG{+w}{ }\PYG{o}{*}\PYG{+w}{ }\PYG{o}{*}\PYG{+w}{ }\PYG{l+m}{18.75}\PYG{+w}{ }\PYG{l+m}{4.5}\PYG{+w}{ }\PYG{l+m}{0.41}\PYG{+w}{ }\PYG{l+m}{1.58}\PYG{+w}{ }\PYG{l+m}{1.42E\PYGZhy{}5}\PYG{+w}{ }\PYG{l+m}{2.0}\PYG{+w}{ }\PYG{l+m}{1.58}
\PYG{k}{pair\PYGZus{}coeff}\PYG{+w}{ }\PYG{l+m}{1}\PYG{+w}{ }\PYG{l+m}{1}\PYG{+w}{ }\PYG{l+m}{18.75}\PYG{+w}{ }\PYG{l+m}{4.5}\PYG{+w}{ }\PYG{l+m}{0.41}\PYG{+w}{ }\PYG{l+m}{1.58}\PYG{+w}{ }\PYG{l+m}{1.42E\PYGZhy{}5}\PYG{+w}{ }\PYG{l+m}{2.0}\PYG{+w}{ }\PYG{l+m}{1.58}\PYG{+w}{ }\PYG{n}{power}\PYG{+w}{ }\PYG{l+m}{10.54}\PYG{+w}{ }\PYG{o}{\PYGZhy{}}\PYG{l+m}{3.66}\PYG{+w}{ }\PYG{l+m}{3.44}\PYG{+w}{ }\PYG{o}{\PYGZhy{}}\PYG{l+m}{4.10}
\PYG{k}{pair\PYGZus{}coeff}\PYG{+w}{ }\PYG{l+m}{1}\PYG{+w}{ }\PYG{l+m}{1}\PYG{+w}{ }\PYG{l+m}{18.75}\PYG{+w}{ }\PYG{l+m}{4.5}\PYG{+w}{ }\PYG{l+m}{0.41}\PYG{+w}{ }\PYG{l+m}{1.58}\PYG{+w}{ }\PYG{l+m}{1.42E\PYGZhy{}5}\PYG{+w}{ }\PYG{l+m}{2.0}\PYG{+w}{ }\PYG{l+m}{1.58}\PYG{+w}{ }\PYG{n}{power}\PYG{+w}{ }\PYG{l+m}{10.54}\PYG{+w}{ }\PYG{o}{\PYGZhy{}}\PYG{l+m}{3.66}\PYG{+w}{ }\PYG{l+m}{3.44}\PYG{+w}{ }\PYG{o}{\PYGZhy{}}\PYG{l+m}{4.10}\PYG{+w}{ }\PYG{n}{kappa}\PYG{+w}{ }\PYG{o}{\PYGZhy{}}\PYG{l+m}{0.44}\PYG{+w}{ }\PYG{o}{\PYGZhy{}}\PYG{l+m}{3.21}\PYG{+w}{ }\PYG{l+m}{5.04}\PYG{+w}{ }\PYG{l+m}{0.00}

\PYG{k}{pair\PYGZus{}style}\PYG{+w}{ }\PYG{n}{hybrid}\PYG{o}{/}\PYG{n}{overlay}\PYG{+w}{ }\PYG{n}{mdpd}\PYG{o}{/}\PYG{n}{rhosum}\PYG{+w}{ }\PYG{n}{mdpd}\PYG{+w}{ }\PYG{l+m}{1.0}\PYG{+w}{ }\PYG{l+m}{1.0}\PYG{+w}{ }\PYG{l+m}{65689}
\PYG{k}{pair\PYGZus{}coeff}\PYG{+w}{ }\PYG{l+m}{1}\PYG{+w}{ }\PYG{l+m}{1}\PYG{+w}{ }\PYG{n}{mdpd}\PYG{o}{/}\PYG{n}{rhosum}\PYG{+w}{  }\PYG{l+m}{0.75}
\PYG{k}{pair\PYGZus{}coeff}\PYG{+w}{ }\PYG{l+m}{1}\PYG{+w}{ }\PYG{l+m}{1}\PYG{+w}{ }\PYG{n}{mdpd}\PYG{+w}{ }\PYG{o}{\PYGZhy{}}\PYG{l+m}{40.0}\PYG{+w}{ }\PYG{l+m}{25.0}\PYG{+w}{ }\PYG{l+m}{18.0}\PYG{+w}{ }\PYG{l+m}{1.0}\PYG{+w}{ }\PYG{l+m}{0.75}

\PYG{k}{pair\PYGZus{}style}\PYG{+w}{ }\PYG{n}{tdpd}\PYG{+w}{ }\PYG{l+m}{1.0}\PYG{+w}{ }\PYG{l+m}{1.58}\PYG{+w}{ }\PYG{l+m}{935662}
\PYG{k}{pair\PYGZus{}coeff}\PYG{+w}{ }\PYG{o}{*}\PYG{+w}{ }\PYG{o}{*}\PYG{+w}{ }\PYG{l+m}{18.75}\PYG{+w}{ }\PYG{l+m}{4.5}\PYG{+w}{ }\PYG{l+m}{0.41}\PYG{+w}{ }\PYG{l+m}{1.58}\PYG{+w}{ }\PYG{l+m}{1.58}\PYG{+w}{ }\PYG{l+m}{1.0}\PYG{+w}{ }\PYG{l+m}{1.0E\PYGZhy{}5}\PYG{+w}{ }\PYG{l+m}{2.0}
\PYG{k}{pair\PYGZus{}coeff}\PYG{+w}{ }\PYG{l+m}{1}\PYG{+w}{ }\PYG{l+m}{1}\PYG{+w}{ }\PYG{l+m}{18.75}\PYG{+w}{ }\PYG{l+m}{4.5}\PYG{+w}{ }\PYG{l+m}{0.41}\PYG{+w}{ }\PYG{l+m}{1.58}\PYG{+w}{ }\PYG{l+m}{1.58}\PYG{+w}{ }\PYG{l+m}{1.0}\PYG{+w}{ }\PYG{l+m}{1.0E\PYGZhy{}5}\PYG{+w}{ }\PYG{l+m}{2.0}\PYG{+w}{ }\PYG{l+m}{3.0}\PYG{+w}{ }\PYG{l+m}{1.0E\PYGZhy{}5}\PYG{+w}{ }\PYG{l+m}{2.0}
\end{sphinxVerbatim}


\subsection{Description}
\label{\detokenize{pair_mesodpd:description}}
\sphinxAtStartPar
The \sphinxstyleemphasis{edpd} style computes the pairwise interactions and heat fluxes
for eDPD particles following the formulations in
{\hyperref[\detokenize{pair_mesodpd:li2014-jcp}]{\sphinxcrossref{\DUrole{std,std-ref}{(Li2014\_JCP)}}}} and {\hyperref[\detokenize{pair_mesodpd:li2015-cc}]{\sphinxcrossref{\DUrole{std,std-ref}{Li2015\_CC}}}}. The time
evolution of an eDPD particle is governed by the conservation of
momentum and energy given by
\begin{equation*}
\begin{split}\frac{\mathrm{d}^2 \mathbf{r}_i}{\mathrm{d} t^2}=
\frac{\mathrm{d} \mathbf{v}_i}{\mathrm{d} t}
=\mathbf{F}_{i}=\sum_{i\neq j}(\mathbf{F}_{ij}^{C}+\mathbf{F}_{ij}^{D}+\mathbf{F}_{ij}^{R}) \\
C_v\frac{\mathrm{d} T_i}{\mathrm{d} t}= q_{i} = \sum_{i\neq j}(q_{ij}^{C}+q_{ij}^{V}+q_{ij}^{R}),\end{split}
\end{equation*}
\sphinxAtStartPar
where the three components of \(F_{i}\) including the conservative
force \(F_{ij}^C\), dissipative force \(F_{ij}^D\) and random
force \(F_{ij}^R\) are expressed as
\begin{equation*}
\begin{split}\mathbf{F}_{ij}^{C} & = \alpha_{ij}{\omega_{C}}(r_{ij})\mathbf{e}_{ij} \\
\mathbf{F}_{ij}^{D} & = -\gamma {\omega_{D}}(r_{ij})(\mathbf{e}_{ij} \cdot \mathbf{v}_{ij})\mathbf{e}_{ij} \\
\mathbf{F}_{ij}^{R} & = \sigma {\omega_{R}}(r_{ij}){\xi_{ij}}\Delta t^{-1/2} \mathbf{e}_{ij} \\
\omega_{C}(r) & = 1 - r/r_c \\
\alpha_{ij} & = A\cdot k_B(T_i + T_j)/2 \\
\omega_{D}(r) & = \omega^2_{R}(r) = (1-r/r_c)^s \\
\sigma_{ij}^2 & = 4\gamma k_B T_i T_j/(T_i + T_j)\end{split}
\end{equation*}
\sphinxAtStartPar
in which the exponent of the weighting function \sphinxstyleemphasis{s} can be defined as a
temperature\sphinxhyphen{}dependent variable. The heat flux between particles
accounting for the collisional heat flux \(q^C\), viscous heat flux
\(q^V\), and random heat flux \(q^R\) are given by
\begin{equation*}
\begin{split}q_i^C & = \sum_{j \ne i} k_{ij} \omega_{CT}(r_{ij}) \left( \frac{1}{T_i} - \frac{1}{T_j} \right) \\
q_i^V & = \frac{1}{2 C_v}\sum_{j \ne i}{ \left\{ \omega_D(r_{ij})\left[\gamma_{ij} \left( \mathbf{e}_{ij} \cdot \mathbf{v}_{ij} \right)^2 - \frac{\left( \sigma _{ij} \right)^2}{m}\right] - \sigma _{ij} \omega_R(r_{ij})\left( \mathbf{e}_{ij} \cdot \mathbf{v}_{ij} \right){\xi_{ij}} \right\} } \\
q_i^R & = \sum_{j \ne i} \beta _{ij} \omega_{RT}(r_{ij}) d {t^{ - 1/2}} \xi_{ij}^e \\
\omega_{CT}(r) & =\omega_{RT}^2(r)=\left(1-r/r_{ct}\right)^{s_T} \\
k_{ij} & =C_v^2\kappa(T_i + T_j)^2/4k_B \\
\beta_{ij}^2 & = 2k_Bk_{ij}\end{split}
\end{equation*}
\sphinxAtStartPar
where the mesoscopic heat friction \(\kappa\) is given by
\begin{equation*}
\begin{split}\kappa  = \frac{315k_B\upsilon }{2\pi \rho C_v r_{ct}^5}\frac{1}{Pr},\end{split}
\end{equation*}
\sphinxAtStartPar
with \(\upsilon\) being the kinematic viscosity. For more details,
see Eq.(15) in {\hyperref[\detokenize{pair_mesodpd:li2014-jcp}]{\sphinxcrossref{\DUrole{std,std-ref}{(Li2014\_JCP)}}}}.

\sphinxAtStartPar
The following coefficients must be defined in eDPD system for each
pair of atom types via the {\hyperref[\detokenize{pair_coeff::doc}]{\sphinxcrossref{\DUrole{doc}{pair\_coeff}}}} command as in
the examples above.
\begin{itemize}
\item {} 
\sphinxAtStartPar
A (force units)

\item {} 
\sphinxAtStartPar
\(\gamma\) (force/velocity units)

\item {} 
\sphinxAtStartPar
power\_f (positive real)

\item {} 
\sphinxAtStartPar
cutoff (distance units)

\item {} 
\sphinxAtStartPar
kappa (thermal conductivity units)

\item {} 
\sphinxAtStartPar
power\_T (positive real)

\item {} 
\sphinxAtStartPar
cutoff\_T (distance units)

\item {} 
\sphinxAtStartPar
optional keyword = power or kappa

\end{itemize}

\sphinxAtStartPar
The keyword \sphinxstyleemphasis{power} or \sphinxstyleemphasis{kappa} is optional. Both “power” and “kappa”
require 4 parameters \(c_1, c_2, c_3, c_4\) showing the temperature
dependence of the exponent \(s(T) = \mathrm{power}_f ( 1+c_1
(T-1) + c_2 (T-1)^2 + c_3 (T-1)^3 + c_4 (T-1)^4 )\) and of the mesoscopic
heat friction \(s_T(T) = \kappa (1 + c_1 (T-1) + c_2 (T-1)^2 + c_3
(T-1)^3 + c_4 (T-1)^4)\).  If the keyword \sphinxstyleemphasis{power} or \sphinxstyleemphasis{kappa} is not
specified, the eDPD system will use constant power\_f and
\(\kappa\), which is independent to temperature changes.


\bigskip\hrule\bigskip


\sphinxAtStartPar
The \sphinxstyleemphasis{mdpd/rhosum} style computes the local particle mass density
\(\rho\) for mDPD particles by kernel function interpolation.

\sphinxAtStartPar
The following coefficients must be defined for each pair of atom types
via the {\hyperref[\detokenize{pair_coeff::doc}]{\sphinxcrossref{\DUrole{doc}{pair\_coeff}}}} command as in the examples above.
\begin{itemize}
\item {} 
\sphinxAtStartPar
cutoff (distance units)

\end{itemize}


\bigskip\hrule\bigskip


\sphinxAtStartPar
The \sphinxstyleemphasis{mdpd} style computes the many\sphinxhyphen{}body interactions between mDPD
particles following the formulations in
{\hyperref[\detokenize{pair_mesodpd:li2013-pof}]{\sphinxcrossref{\DUrole{std,std-ref}{(Li2013\_POF)}}}}. The dissipative and random forces are in
the form same as the classical DPD, but the conservative force is
local density dependent, which are given by
\begin{equation*}
\begin{split}\mathbf{F}_{ij}^C & = Aw_c(r_{ij})\mathbf{e}_{ij} + B(\rho_i+\rho_j)w_d(r_{ij})\mathbf{e}_{ij} \\
\mathbf{F}_{ij}^{D} & = -\gamma {\omega_{D}}(r_{ij})(\mathbf{e}_{ij} \cdot \mathbf{v}_{ij})\mathbf{e}_{ij} \\
\mathbf{F}_{ij}^{R} & = \sigma {\omega_{R}}(r_{ij}){\xi_{ij}}\Delta t^{-1/2} \mathbf{e}_{ij}\end{split}
\end{equation*}
\sphinxAtStartPar
where the first term in \(F_C\) with a negative coefficient \(A
< 0\) stands for an attractive force within an interaction range
\(r_c\), and the second term with \(B > 0\) is the
density\sphinxhyphen{}dependent repulsive force within an interaction range
\(r_d\).

\sphinxAtStartPar
The following coefficients must be defined for each pair of atom types via the
{\hyperref[\detokenize{pair_coeff::doc}]{\sphinxcrossref{\DUrole{doc}{pair\_coeff}}}} command as in the examples above.
\begin{itemize}
\item {} 
\sphinxAtStartPar
A (force units)

\item {} 
\sphinxAtStartPar
B (force units)

\item {} 
\sphinxAtStartPar
\(\gamma\) (force/velocity units)

\item {} 
\sphinxAtStartPar
cutoff\_c (distance units)

\item {} 
\sphinxAtStartPar
cutoff\_d (distance units)

\end{itemize}


\bigskip\hrule\bigskip


\sphinxAtStartPar
The \sphinxstyleemphasis{tdpd} style computes the pairwise interactions and chemical
concentration fluxes for tDPD particles following the formulations in
{\hyperref[\detokenize{pair_mesodpd:li2015-jcp}]{\sphinxcrossref{\DUrole{std,std-ref}{(Li2015\_JCP)}}}}.  The time evolution of a tDPD particle is
governed by the conservation of momentum and concentration given by
\begin{equation*}
\begin{split}\frac{\mathrm{d}^2 \mathbf{r}_i}{\mathrm{d} t^2} & = \frac{\mathrm{d} \mathbf{v}_i}{\mathrm{d} t}=\mathbf{F}_{i}=\sum_{i\neq j}(\mathbf{F}_{ij}^{C}+\mathbf{F}_{ij}^{D}+\mathbf{F}_{ij}^{R}) \\
\frac{\mathrm{d} C_{i}}{\mathrm{d} t} & = Q_{i} = \sum_{i\neq j}(Q_{ij}^{D}+Q_{ij}^{R}) + Q_{i}^{S}\end{split}
\end{equation*}
\sphinxAtStartPar
where the three components of \(F_{i}\) including the conservative
force \(F_{ij}^C\), dissipative force \(F_{ij}^C\) and random
force \(F_{ij}^C\) are expressed as
\begin{equation*}
\begin{split}\mathbf{F}_{ij}^{C} & = A{\omega_{C}}(r_{ij})\mathbf{e}_{ij} \\
\mathbf{F}_{ij}^{D} & = -\gamma {\omega_{D}}(r_{ij})(\mathbf{e}_{ij} \cdot \mathbf{v}_{ij})\mathbf{e}_{ij}  \\
\mathbf{F}_{ij}^{R} & = \sigma {\omega_{R}}(r_{ij}){\xi_{ij}}\Delta t^{-1/2} \mathbf{e}_{ij} \\
\omega_{C}(r) & = 1 - r/r_c \\
\omega_{D}(r) & = \omega^2_{R}(r) = (1-r/r_c)^{\rm power_f} \\
\sigma^2 = 2\gamma k_B T\end{split}
\end{equation*}
\sphinxAtStartPar
The concentration flux between two tDPD particles includes the Fickian
flux \(Q_{ij}^D\) and random flux \(Q_{ij}^R\), which are given
by
\begin{equation*}
\begin{split}Q_{ij}^D & = -\kappa_{ij} w_{DC}(r_{ij}) \left( C_i - C_j \right) \\
Q_{ij}^R & = \epsilon_{ij}\left( C_i + C_j \right) w_{RC}(r_{ij}) \xi_{ij} \\
w_{DC}(r_{ij}) & =w^2_{RC}(r_{ij}) = (1 - r/r_{cc})^{\rm power_{cc}} \\
\epsilon_{ij}^2 & = m_s^2\kappa_{ij}\rho\end{split}
\end{equation*}
\sphinxAtStartPar
where the parameters kappa and epsilon determine the strength of the
Fickian and random fluxes. \(m_s\) is the mass of a single solute
molecule.  In general, \(m_s\) is much smaller than the mass of a
tDPD particle \sphinxstyleemphasis{m}. For more details, see {\hyperref[\detokenize{pair_mesodpd:li2015-jcp}]{\sphinxcrossref{\DUrole{std,std-ref}{(Li2015\_JCP)}}}}.

\sphinxAtStartPar
The following coefficients must be defined for each pair of atom types via the
{\hyperref[\detokenize{pair_coeff::doc}]{\sphinxcrossref{\DUrole{doc}{pair\_coeff}}}} command as in the examples above.
\begin{itemize}
\item {} 
\sphinxAtStartPar
A (force units)

\item {} 
\sphinxAtStartPar
\(\gamma\) (force/velocity units)

\item {} 
\sphinxAtStartPar
power\_f (positive real)

\item {} 
\sphinxAtStartPar
cutoff (distance units)

\item {} 
\sphinxAtStartPar
cutoff\_CC (distance units)

\item {} 
\sphinxAtStartPar
\(\kappa_i\) (diffusivity units)

\item {} 
\sphinxAtStartPar
\(\epsilon_i\) (diffusivity units)

\item {} 
\sphinxAtStartPar
power\_cc\_i (positive real)

\end{itemize}

\sphinxAtStartPar
The last 3 values must be repeated Nspecies times, so that values for
each of the Nspecies chemical species are specified, as indicated by
the “I” suffix.  In the first pair\_coeff example above for pair\_style
tdpd, Nspecies = 1.  In the second example, Nspecies = 2, so 3
additional coeffs are specified (for species 2).


\bigskip\hrule\bigskip



\subsection{Example scripts}
\label{\detokenize{pair_mesodpd:example-scripts}}
\sphinxAtStartPar
There are example scripts for using all these pair styles in
examples/PACKAGES/mesodpd.  The example for an eDPD simulation models heat
conduction with source terms analog of periodic Poiseuille flow
problem. The setup follows Fig.12 in {\hyperref[\detokenize{pair_mesodpd:li2014-jcp}]{\sphinxcrossref{\DUrole{std,std-ref}{(Li2014\_JCP)}}}}. The
output of the short eDPD simulation (about 2 minutes on a single core)
gives a temperature and density profiles as

\noindent{\hspace*{\fill}\sphinxincludegraphics{{examples_edpd}.jpg}\hspace*{\fill}}

\sphinxAtStartPar
The example for a mDPD simulation models the oscillations of a liquid
droplet started from a liquid film. The mDPD parameters are adopted
from {\hyperref[\detokenize{pair_mesodpd:li2013-pof}]{\sphinxcrossref{\DUrole{std,std-ref}{(Li2013\_POF)}}}}.  The short mDPD run (about 2 minutes
on a single core) generates a particle trajectory which can
be visualized as follows.

\noindent{\hspace*{\fill}\sphinxincludegraphics{{examples_mdpd_first}.jpg}\hspace*{\fill}}

\noindent{\hspace*{\fill}\sphinxincludegraphics{{examples_mdpd_last}.jpg}\hspace*{\fill}}

\sphinxAtStartPar
The first image is the initial state of the simulation.  If you
click it a GIF movie should play in your browser.  The second image
is the final state of the simulation.

\sphinxAtStartPar
The example for a tDPD simulation computes the effective diffusion
coefficient of a tDPD system using a method analogous to the periodic
Poiseuille flow.  The tDPD system is specified with two chemical
species, and the setup follows Fig.1 in
{\hyperref[\detokenize{pair_mesodpd:li2015-jcp}]{\sphinxcrossref{\DUrole{std,std-ref}{(Li2015\_JCP)}}}}. The output of the short tDPD simulation
(about one and a half minutes on a single core) gives the
concentration profiles of the two chemical species as

\noindent{\hspace*{\fill}\sphinxincludegraphics{{examples_tdpd}.jpg}\hspace*{\fill}}


\bigskip\hrule\bigskip



\subsection{Mixing, shift, table, tail correction, restart, rRESPA info}
\label{\detokenize{pair_mesodpd:mixing-shift-table-tail-correction-restart-rrespa-info}}
\sphinxAtStartPar
The styles \sphinxstyleemphasis{edpd}, \sphinxstyleemphasis{mdpd}, \sphinxstyleemphasis{mdpd/rhosum} and \sphinxstyleemphasis{tdpd} do not support
mixing. Thus, coefficients for all I,J pairs must be specified explicitly.

\sphinxAtStartPar
The styles \sphinxstyleemphasis{edpd}, \sphinxstyleemphasis{mdpd}, \sphinxstyleemphasis{mdpd/rhosum} and \sphinxstyleemphasis{tdpd} do not support
the {\hyperref[\detokenize{pair_modify::doc}]{\sphinxcrossref{\DUrole{doc}{pair\_modify}}}} shift, table, and tail options.

\sphinxAtStartPar
The styles \sphinxstyleemphasis{edpd}, \sphinxstyleemphasis{mdpd}, \sphinxstyleemphasis{mdpd/rhosum} and \sphinxstyleemphasis{tdpd} do not write
information to {\hyperref[\detokenize{restart::doc}]{\sphinxcrossref{\DUrole{doc}{binary restart files}}}}. Thus, you need
to re\sphinxhyphen{}specify the pair\_style and pair\_coeff commands in an input script
that reads a restart file.


\subsection{Restrictions}
\label{\detokenize{pair_mesodpd:restrictions}}
\sphinxAtStartPar
The pair styles \sphinxstyleemphasis{edpd}, \sphinxstyleemphasis{mdpd}, \sphinxstyleemphasis{mdpd/rhosum} and \sphinxstyleemphasis{tdpd} are part of
the DPD\sphinxhyphen{}MESO package. They are only enabled if LAMMPS was built with
that package.  See the {\hyperref[\detokenize{Build_package::doc}]{\sphinxcrossref{\DUrole{doc}{Build package}}}} page for
more info.


\subsection{Related commands}
\label{\detokenize{pair_mesodpd:related-commands}}
\sphinxAtStartPar
{\hyperref[\detokenize{pair_coeff::doc}]{\sphinxcrossref{\DUrole{doc}{pair\_coeff}}}}, {\hyperref[\detokenize{fix_mvv_dpd::doc}]{\sphinxcrossref{\DUrole{doc}{fix mvv/dpd}}}},
{\hyperref[\detokenize{fix_mvv_dpd::doc}]{\sphinxcrossref{\DUrole{doc}{fix mvv/edpd}}}}, {\hyperref[\detokenize{fix_mvv_dpd::doc}]{\sphinxcrossref{\DUrole{doc}{fix mvv/tdpd}}}},
{\hyperref[\detokenize{fix_dpd_source::doc}]{\sphinxcrossref{\DUrole{doc}{fix edpd/source}}}}, {\hyperref[\detokenize{fix_dpd_source::doc}]{\sphinxcrossref{\DUrole{doc}{fix tdpd/source}}}},
{\hyperref[\detokenize{compute_edpd_temp_atom::doc}]{\sphinxcrossref{\DUrole{doc}{compute edpd/temp/atom}}}},
{\hyperref[\detokenize{compute_tdpd_cc_atom::doc}]{\sphinxcrossref{\DUrole{doc}{compute tdpd/cc/atom}}}}


\subsection{Default}
\label{\detokenize{pair_mesodpd:default}}
\sphinxAtStartPar
none


\bigskip\hrule\bigskip

\phantomsection\label{\detokenize{pair_mesodpd:li2014-jcp}}
\sphinxAtStartPar
\sphinxstylestrong{(Li2014\_JCP)} Li, Tang, Lei, Caswell, Karniadakis, J Comput Phys,
265: 113\sphinxhyphen{}127 (2014).  DOI: 10.1016/j.jcp.2014.02.003.

\phantomsection\label{\detokenize{pair_mesodpd:li2015-cc}}
\sphinxAtStartPar
\sphinxstylestrong{(Li2015\_CC)} Li, Tang, Li, Karniadakis, Chem Commun, 51: 11038\sphinxhyphen{}11040
(2015).  DOI: 10.1039/C5CC01684C.

\phantomsection\label{\detokenize{pair_mesodpd:li2013-pof}}
\sphinxAtStartPar
\sphinxstylestrong{(Li2013\_POF)} Li, Hu, Wang, Ma, Zhou, Phys Fluids, 25: 072103 (2013).
DOI: 10.1063/1.4812366.

\phantomsection\label{\detokenize{pair_mesodpd:li2015-jcp}}
\sphinxAtStartPar
\sphinxstylestrong{(Li2015\_JCP)} Li, Yazdani, Tartakovsky, Karniadakis, J Chem Phys,
143: 014101 (2015).  DOI: 10.1063/1.4923254.

\sphinxstepscope

\index{pair\_style mgpt@\spxentry{pair\_style mgpt}}\ignorespaces 

\section{pair\_style mgpt command}
\label{\detokenize{pair_mgpt:pair-style-mgpt-command}}\label{\detokenize{pair_mgpt:index-0}}\label{\detokenize{pair_mgpt::doc}}

\subsection{Syntax}
\label{\detokenize{pair_mgpt:syntax}}
\begin{sphinxVerbatim}[commandchars=\\\{\}]
\PYG{k}{pair\PYGZus{}style}\PYG{+w}{ }\PYG{n}{mgpt}
\end{sphinxVerbatim}


\subsection{Examples}
\label{\detokenize{pair_mgpt:examples}}
\begin{sphinxVerbatim}[commandchars=\\\{\}]
\PYG{k}{pair\PYGZus{}style}\PYG{+w}{ }\PYG{n}{mgpt}
\PYG{k}{pair\PYGZus{}coeff}\PYG{+w}{ }\PYG{o}{*}\PYG{+w}{ }\PYG{o}{*}\PYG{+w}{ }\PYG{n}{Ta6.8x.mgpt.parmin}\PYG{+w}{ }\PYG{n}{Ta6.8x.mgpt.potin}\PYG{+w}{ }\PYG{n}{Omega}
\PYG{n}{cp}\PYG{+w}{ }\PYGZti{}/lammps/potentials/Ta6.8x.mgpt.parmin\PYG{+w}{ }\PYG{n}{parmin}
\PYG{n}{cp}\PYG{+w}{ }\PYGZti{}/lammps/potentials/Ta6.8x.mgpt.potin\PYG{+w}{ }\PYG{n}{potin}
\PYG{k}{pair\PYGZus{}coeff}\PYG{+w}{ }\PYG{o}{*}\PYG{+w}{ }\PYG{o}{*}\PYG{+w}{ }\PYG{n}{parmin}\PYG{+w}{ }\PYG{n}{potin}\PYG{+w}{ }\PYG{n}{Omega}\PYG{+w}{ }\PYG{n}{volpress}\PYG{+w}{ }\PYG{n}{yes}\PYG{+w}{ }\PYG{n}{nbody}\PYG{+w}{ }\PYG{l+m}{1234}\PYG{+w}{ }\PYG{n}{precision}\PYG{+w}{ }\PYG{n}{double}
\PYG{k}{pair\PYGZus{}coeff}\PYG{+w}{ }\PYG{o}{*}\PYG{+w}{ }\PYG{o}{*}\PYG{+w}{ }\PYG{n}{parmin}\PYG{+w}{ }\PYG{n}{potin}\PYG{+w}{ }\PYG{n}{Omega}\PYG{+w}{ }\PYG{n}{volpress}\PYG{+w}{ }\PYG{n}{yes}\PYG{+w}{ }\PYG{n}{nbody}\PYG{+w}{ }\PYG{l+m}{12}
\end{sphinxVerbatim}


\subsection{Description}
\label{\detokenize{pair_mgpt:description}}
\sphinxAtStartPar
Within DFT quantum mechanics, generalized pseudopotential theory (GPT)
({\hyperref[\detokenize{pair_mgpt:moriarty1}]{\sphinxcrossref{\DUrole{std,std-ref}{Moriarty1}}}}) provides a first\sphinxhyphen{}principles approach to
multi\sphinxhyphen{}ion interatomic potentials in d\sphinxhyphen{}band transition metals, with a
volume\sphinxhyphen{}dependent, real\sphinxhyphen{}space total\sphinxhyphen{}energy functional for the N\sphinxhyphen{}ion
elemental bulk material in the form
\begin{equation*}
\begin{split}E_{\rm tot}({\bf R}_1 \ldots {\bf R}_N) = NE_{\rm vol}(\Omega )
+ \frac{1}{2} \sum _{i,j} \mbox{}^\prime \ v_2(ij;\Omega )
+ \frac{1}{6} \sum _{i,j,k} \mbox{}^\prime \ v_3(ijk;\Omega )
+ \frac{1}{24} \sum _{i,j,k,l} \mbox{}^\prime \ v_4(ijkl;\Omega )\end{split}
\end{equation*}
\sphinxAtStartPar
where the prime on each summation sign indicates the exclusion of all
self\sphinxhyphen{}interaction terms from the summation.  The leading volume term
E\_vol as well as the two\sphinxhyphen{}ion central\sphinxhyphen{}force pair potential v\_2 and the
three\sphinxhyphen{} and four\sphinxhyphen{}ion angular\sphinxhyphen{}force potentials, v\_3 and v\_4, depend
explicitly on the atomic volume Omega, but are structure independent
and transferable to all bulk ion configurations, either ordered or
disordered, and with of without the presence of point and line
defects.  The simplified model GPT or MGPT ({\hyperref[\detokenize{pair_mgpt:moriarty2}]{\sphinxcrossref{\DUrole{std,std-ref}{Moriarty2}}}},
{\hyperref[\detokenize{pair_mgpt:moriarty3}]{\sphinxcrossref{\DUrole{std,std-ref}{Moriarty3}}}}), which retains the form of E\_tot and permits
more efficient large\sphinxhyphen{}scale atomistic simulations, derives from the GPT
through a series of systematic approximations applied to E\_vol and the
potentials v\_n that are valid for mid\sphinxhyphen{}period transition metals with
nearly half\sphinxhyphen{}filled d bands.

\sphinxAtStartPar
Both analytic ({\hyperref[\detokenize{pair_mgpt:moriarty2}]{\sphinxcrossref{\DUrole{std,std-ref}{Moriarty2}}}}) and matrix
({\hyperref[\detokenize{pair_mgpt:moriarty3}]{\sphinxcrossref{\DUrole{std,std-ref}{Moriarty3}}}}) representations of MGPT have been developed.
In the more general matrix representation, which can also be applied
to f\sphinxhyphen{}band actinide metals and permits both canonical and non\sphinxhyphen{}canonical
d/f bands, the multi\sphinxhyphen{}ion potentials are evaluated on the fly during a
simulation through d\sphinxhyphen{} or f\sphinxhyphen{}state matrix multiplication, and the forces
that move the ions are determined analytically.  Fast matrix\sphinxhyphen{}MGPT
algorithms have been developed independently by Glosli
({\hyperref[\detokenize{pair_mgpt:glosli}]{\sphinxcrossref{\DUrole{std,std-ref}{Glosli}}}}, {\hyperref[\detokenize{pair_mgpt:moriarty3}]{\sphinxcrossref{\DUrole{std,std-ref}{Moriarty3}}}}) and by Oppelstrup
({\hyperref[\detokenize{pair_mgpt:oppelstrup}]{\sphinxcrossref{\DUrole{std,std-ref}{Oppelstrup}}}})

\sphinxAtStartPar
The \sphinxstyleemphasis{mgpt} pair style calculates forces, energies, and the total
energy per atom, E\_tot/N, using the Oppelstrup matrix\sphinxhyphen{}MGPT algorithm.
Input potential and control data are entered through the
{\hyperref[\detokenize{pair_coeff::doc}]{\sphinxcrossref{\DUrole{doc}{pair\_coeff}}}} command.  Each material treated requires
input parmin and potin potential files, as shown in the above
examples, as well as specification by the user of the initial atomic
volume Omega through pair\_coeff.  At the beginning of a time step in
any simulation, the total volume of the simulation cell V should
always be equal to Omega*N, where N is the number of metal ions
present, taking into account the presence of any vacancies and/or
interstitials in the case of a solid.  In a constant\sphinxhyphen{}volume
simulation, which is the normal mode of operation for the \sphinxstyleemphasis{mgpt} pair
style, Omega, V and N all remain constant throughout the simulation
and thus are equal to their initial values.  In a constant\sphinxhyphen{}stress
simulation, the cell volume V will change (slowly) as the simulation
proceeds.  After each time step, the atomic volume should be updated
by the code as Omega = V/N.  In addition, the volume term E\_vol and
the potentials v\_2, v\_3 and v\_4 have to be removed at the end of the
time step, and then respecified at the new value of Omega.  In all
simulations, Omega must remain within the defined volume range for
E\_vol and the potentials for the given material.

\sphinxAtStartPar
The default option volpress yes in the {\hyperref[\detokenize{pair_coeff::doc}]{\sphinxcrossref{\DUrole{doc}{pair\_coeff}}}}
command includes all volume derivatives of E\_tot required to calculate
the stress tensor and pressure correctly.  The option volpress no
disregards the pressure contribution resulting from the volume term
E\_vol, and can be used for testing and analysis purposes.  The
additional optional variable nbody controls the specific terms in
E\_tot that are calculated.  The default option and the normal option
for mid\sphinxhyphen{}period transition and actinide metals is nbody 1234 for which
all four terms in E\_tot are retained.  The option nbody 12, for
example, retains only the volume term and the two\sphinxhyphen{}ion pair potential
term and can be used for GPT series\sphinxhyphen{}end transition metals that can be
well described without v\_3 and v\_4.  The nbody option can also be used
to test or analyze the contribution of any of the four terms in E\_tot
to a given calculated property.

\sphinxAtStartPar
The \sphinxstyleemphasis{mgpt} pair style makes extensive use of matrix algebra and
includes optimized kernels for the BlueGene/Q architecture and the
Intel/AMD (x86) architectures.  When compiled with the appropriate
compiler and compiler switches (\sphinxhyphen{}msse3 on x86, and using the IBM XL
compiler on BG/Q), these optimized routines are used automatically.
For BG/Q machines, building with the default Makefile for that
architecture (e.g., “make bgq”) should enable the optimized algebra
routines.  For x\sphinxhyphen{}86 machines, there is a provided Makefile.mgptfast
which enables the fast algebra routines, i.e. build LAMMPS with “make
mgptfast”.  The user will be informed in the output files of the
matrix kernels in use. To further improve speed, on x86 the option
precision single can be added to the {\hyperref[\detokenize{pair_coeff::doc}]{\sphinxcrossref{\DUrole{doc}{pair\_coeff}}}}
command line, which improves speed (up to a factor of two) at the cost
of doing matrix calculations with 7 digit precision instead of the
default 16. For consistency the default option can be specified
explicitly by the option precision double.

\sphinxAtStartPar
All remaining potential and control data are contained with the parmin
and potin files, including cutoffs, atomic mass, and other basic MGPT
variables.  Specific MGPT potential data for the transition metals
tantalum (Ta4 and Ta6.8x potentials), molybdenum (Mo5.2 potentials),
and vanadium (V6.1 potentials) are contained in the LAMMPS potentials
directory.  The stored files are, respectively, Ta4.mgpt.parmin,
Ta4.mgpt.potin, Ta6.8x.mgpt.parmin, Ta6.8x.mgpt.potin,
Mo5.2.mgpt.parmin, Mo5.2.mgpt.potin, V6.1.mgpt.parmin, and
V6.1.mgpt.potin .  Useful corresponding informational “README” files
on the Ta4, Ta6.8x, Mo5.2 and V6.1 potentials are also included in the
potentials directory.  These latter files indicate the volume mesh and
range for each potential and give appropriate references for the
potentials.  It is expected that MGPT potentials for additional
materials will be added over time.

\sphinxAtStartPar
Useful example MGPT scripts are given in the examples/PACKAGES/mgpt
directory.  These scripts show the necessary steps to perform
constant\sphinxhyphen{}volume calculations and simulations.  It is strongly
recommended that the user work through and understand these examples
before proceeding to more complex simulations.

\begin{sphinxadmonition}{note}{Note:}
\sphinxAtStartPar
For good performance, LAMMPS should be built with the compiler
flags “\sphinxhyphen{}O3 \sphinxhyphen{}msse3 \sphinxhyphen{}funroll\sphinxhyphen{}loops” when including this pair style.  The
src/MAKE/OPTIONS/Makefile.mgptfast is an example machine Makefile with
these options included as part of a standard MPI build.  Note that it
as provided, it will build with whatever low\sphinxhyphen{}level compiler (g++, icc,
etc) is the default for your MPI installation.
\end{sphinxadmonition}


\bigskip\hrule\bigskip



\subsection{Mixing, shift, table, tail correction, restart, rRESPA info}
\label{\detokenize{pair_mgpt:mixing-shift-table-tail-correction-restart-rrespa-info}}
\sphinxAtStartPar
This pair style does not support the {\hyperref[\detokenize{pair_modify::doc}]{\sphinxcrossref{\DUrole{doc}{pair\_modify}}}}
mix, shift, table, and tail options.

\sphinxAtStartPar
This pair style does not write its information to {\hyperref[\detokenize{restart::doc}]{\sphinxcrossref{\DUrole{doc}{binary restart files}}}}, since it is stored in potential files.  Thus, you
needs to re\sphinxhyphen{}specify the pair\_style and pair\_coeff commands in an input
script that reads a restart file.

\sphinxAtStartPar
This pair style can only be used via the \sphinxstyleemphasis{pair} keyword of the
{\hyperref[\detokenize{run_style::doc}]{\sphinxcrossref{\DUrole{doc}{run\_style respa}}}} command.  It does not support the
\sphinxstyleemphasis{inner}, \sphinxstyleemphasis{middle}, \sphinxstyleemphasis{outer} keywords.


\bigskip\hrule\bigskip



\subsection{Restrictions}
\label{\detokenize{pair_mgpt:restrictions}}
\sphinxAtStartPar
This pair style is part of the MGPT package and is only enabled
if LAMMPS is built with that package.  See the {\hyperref[\detokenize{Build_package::doc}]{\sphinxcrossref{\DUrole{doc}{Build package}}}} page for more info.

\sphinxAtStartPar
The MGPT potentials require the {\hyperref[\detokenize{newton::doc}]{\sphinxcrossref{\DUrole{doc}{newtion}}}} setting to be
“on” for pair style interactions.

\sphinxAtStartPar
The stored parmin and potin potential files provided with LAMMPS in
the “potentials” directory are written in Rydberg atomic units, with
energies in Rydbergs and distances in Bohr radii. The \sphinxstyleemphasis{mgpt} pair
style converts Rydbergs to Hartrees to make the potential files
compatible with LAMMPS electron {\hyperref[\detokenize{units::doc}]{\sphinxcrossref{\DUrole{doc}{units}}}}.

\sphinxAtStartPar
The form of E\_tot used in the \sphinxstyleemphasis{mgpt} pair style is only appropriate
for elemental bulk solids and liquids.  This includes solids with
point and extended defects such as vacancies, interstitials, grain
boundaries and dislocations.  Alloys and free surfaces, however,
require significant modifications, which are not included in the
\sphinxstyleemphasis{mgpt} pair style.  Likewise, the \sphinxstyleemphasis{hybrid} pair style is not allowed,
where MGPT would be used for some atoms but not for others.

\sphinxAtStartPar
Electron\sphinxhyphen{}thermal effects are not included in the standard MGPT
potentials provided in the “potentials” directory, where the
potentials have been constructed at zero electron temperature.
Physically, electron\sphinxhyphen{}thermal effects may be important in 3d (e.g., V)
and 4d (e.g., Mo) transition metals at high temperatures near melt and
above.  It is expected that temperature\sphinxhyphen{}dependent MGPT potentials for
such cases will be added over time.


\subsection{Related commands}
\label{\detokenize{pair_mgpt:related-commands}}
\sphinxAtStartPar
{\hyperref[\detokenize{pair_coeff::doc}]{\sphinxcrossref{\DUrole{doc}{pair\_coeff}}}}


\subsection{Default}
\label{\detokenize{pair_mgpt:default}}
\sphinxAtStartPar
The options defaults for the {\hyperref[\detokenize{pair_coeff::doc}]{\sphinxcrossref{\DUrole{doc}{pair\_coeff}}}} command are
volpress yes, nbody 1234, and precision double.


\bigskip\hrule\bigskip

\phantomsection\label{\detokenize{pair_mgpt:moriarty1}}
\sphinxAtStartPar
\sphinxstylestrong{(Moriarty1)} Moriarty, Physical Review B, 38, 3199 (1988).

\phantomsection\label{\detokenize{pair_mgpt:moriarty2}}
\sphinxAtStartPar
\sphinxstylestrong{(Moriarty2)} Moriarty, Physical Review B, 42, 1609 (1990).
Moriarty, Physical Review B 49, 12431 (1994).

\phantomsection\label{\detokenize{pair_mgpt:moriarty3}}
\sphinxAtStartPar
\sphinxstylestrong{(Moriarty3)} Moriarty, Benedict, Glosli, Hood, Orlikowski, Patel, Soderlind, Streitz, Tang, and Yang,
Journal of Materials Research, 21, 563 (2006).

\phantomsection\label{\detokenize{pair_mgpt:glosli}}
\sphinxAtStartPar
\sphinxstylestrong{(Glosli)} Glosli, unpublished, 2005.
Streitz, Glosli, Patel, Chan, Yates, de Supinski, Sexton and Gunnels, Journal of Physics: Conference
Series, 46, 254 (2006).

\phantomsection\label{\detokenize{pair_mgpt:oppelstrup}}
\sphinxAtStartPar
\sphinxstylestrong{(Oppelstrup)} Oppelstrup, unpublished, 2015.
Oppelstrup and Moriarty, to be published.

\sphinxstepscope

\index{pair\_style mie/cut@\spxentry{pair\_style mie/cut}}\ignorespaces 
\index{pair\_style mie/cut/gpu@\spxentry{pair\_style mie/cut/gpu}}\ignorespaces 

\section{pair\_style mie/cut command}
\label{\detokenize{pair_mie:pair-style-mie-cut-command}}\label{\detokenize{pair_mie:index-1}}\label{\detokenize{pair_mie:index-0}}\label{\detokenize{pair_mie::doc}}
\sphinxAtStartPar
Accelerator Variants: \sphinxstyleemphasis{mie/cut/gpu}


\subsection{Syntax}
\label{\detokenize{pair_mie:syntax}}
\begin{sphinxVerbatim}[commandchars=\\\{\}]
\PYG{k}{pair\PYGZus{}style}\PYG{+w}{ }\PYG{n}{mie}\PYG{o}{/}\PYG{n}{cut}\PYG{+w}{ }\PYG{n}{cutoff}
\end{sphinxVerbatim}
\begin{itemize}
\item {} 
\sphinxAtStartPar
cutoff = global cutoff for mie/cut interactions (distance units)

\end{itemize}


\subsection{Examples}
\label{\detokenize{pair_mie:examples}}
\begin{sphinxVerbatim}[commandchars=\\\{\}]
\PYG{k}{pair\PYGZus{}style}\PYG{+w}{ }\PYG{n}{mie}\PYG{o}{/}\PYG{n}{cut}\PYG{+w}{ }\PYG{l+m}{10.0}
\PYG{k}{pair\PYGZus{}coeff}\PYG{+w}{ }\PYG{l+m}{1}\PYG{+w}{ }\PYG{l+m}{1}\PYG{+w}{ }\PYG{l+m}{0.72}\PYG{+w}{ }\PYG{l+m}{3.40}\PYG{+w}{ }\PYG{l+m}{23.00}\PYG{+w}{ }\PYG{l+m}{6.66}
\PYG{k}{pair\PYGZus{}coeff}\PYG{+w}{ }\PYG{l+m}{2}\PYG{+w}{ }\PYG{l+m}{2}\PYG{+w}{ }\PYG{l+m}{0.30}\PYG{+w}{ }\PYG{l+m}{3.55}\PYG{+w}{ }\PYG{l+m}{12.65}\PYG{+w}{ }\PYG{l+m}{6.00}
\PYG{k}{pair\PYGZus{}coeff}\PYG{+w}{ }\PYG{l+m}{1}\PYG{+w}{ }\PYG{l+m}{2}\PYG{+w}{ }\PYG{l+m}{0.46}\PYG{+w}{ }\PYG{l+m}{3.32}\PYG{+w}{ }\PYG{l+m}{16.90}\PYG{+w}{ }\PYG{l+m}{6.31}
\end{sphinxVerbatim}


\subsection{Description}
\label{\detokenize{pair_mie:description}}
\sphinxAtStartPar
The \sphinxstyleemphasis{mie/cut} style computes the Mie potential, given by
\begin{equation*}
\begin{split}E =  C \epsilon \left[ \left(\frac{\sigma}{r}\right)^{\gamma_{rep}} - \left(\frac{\sigma}{r}\right)^{\gamma_{att}} \right]
                      \qquad r < r_c\end{split}
\end{equation*}
\sphinxAtStartPar
\(r_c\) is the cutoff and C is a function that depends on the repulsive and
attractive exponents, given by:
\begin{equation*}
\begin{split}C = \left(\frac{\gamma_{rep}}{\gamma_{rep}-\gamma_{att}}\right) \left(\frac{\gamma_{rep}}{\gamma_{att}}\right)^{\left(\frac{\gamma_{att}}{\gamma_{rep}-\gamma_{att}}\right)}\end{split}
\end{equation*}
\sphinxAtStartPar
Note that for 12/6 exponents, C is equal to 4 and the formula is the
same as the standard Lennard\sphinxhyphen{}Jones potential.

\sphinxAtStartPar
The following coefficients must be defined for each pair of atoms
types via the {\hyperref[\detokenize{pair_coeff::doc}]{\sphinxcrossref{\DUrole{doc}{pair\_coeff}}}} command as in the examples
above, or in the data file or restart files read by the
{\hyperref[\detokenize{read_data::doc}]{\sphinxcrossref{\DUrole{doc}{read\_data}}}} or {\hyperref[\detokenize{read_restart::doc}]{\sphinxcrossref{\DUrole{doc}{read\_restart}}}}
commands, or by mixing as described below:
\begin{itemize}
\item {} 
\sphinxAtStartPar
epsilon (energy units)

\item {} 
\sphinxAtStartPar
sigma (distance units)

\item {} 
\sphinxAtStartPar
gammaR

\item {} 
\sphinxAtStartPar
gammaA

\item {} 
\sphinxAtStartPar
cutoff (distance units)

\end{itemize}

\sphinxAtStartPar
The last coefficient is optional.  If not specified, the global
cutoff specified in the pair\_style command is used.


\bigskip\hrule\bigskip



\subsection{Mixing, shift, table, tail correction, restart, rRESPA info}
\label{\detokenize{pair_mie:mixing-shift-table-tail-correction-restart-rrespa-info}}
\sphinxAtStartPar
For atom type pairs I,J and I != J, the epsilon and sigma coefficients
and cutoff distance for all of the mie/cut pair styles can be mixed.
If not explicitly defined, both the repulsive and attractive gamma
exponents for different atoms will be calculated following the same
mixing rule defined for distances.  The default mix value is
\sphinxstyleemphasis{geometric}. See the “pair\_modify” command for details.

\sphinxAtStartPar
This pair style supports the {\hyperref[\detokenize{pair_modify::doc}]{\sphinxcrossref{\DUrole{doc}{pair\_modify}}}} shift
option for the energy of the pair interaction.

\sphinxAtStartPar
This pair style supports the {\hyperref[\detokenize{pair_modify::doc}]{\sphinxcrossref{\DUrole{doc}{pair\_modify}}}} tail
option for adding a long\sphinxhyphen{}range tail correction to the energy and
pressure of the pair interaction.

\sphinxAtStartPar
This pair style writes its information to {\hyperref[\detokenize{restart::doc}]{\sphinxcrossref{\DUrole{doc}{binary restart files}}}}, so pair\_style and pair\_coeff commands do not need
to be specified in an input script that reads a restart file.

\sphinxAtStartPar
This pair style supports the use of the \sphinxstyleemphasis{inner}, \sphinxstyleemphasis{middle}, and \sphinxstyleemphasis{outer}
keywords of the {\hyperref[\detokenize{run_style::doc}]{\sphinxcrossref{\DUrole{doc}{run\_style respa}}}} command, meaning the
pairwise forces can be partitioned by distance at different levels of
the rRESPA hierarchy.  See the {\hyperref[\detokenize{run_style::doc}]{\sphinxcrossref{\DUrole{doc}{run\_style}}}} command for
details.


\bigskip\hrule\bigskip



\subsection{Restrictions}
\label{\detokenize{pair_mie:restrictions}}
\sphinxAtStartPar
This pair style is part of the EXTRA\sphinxhyphen{}PAIR package.  It is only enabled if
LAMMPS was built with that package.  See the
{\hyperref[\detokenize{Build_package::doc}]{\sphinxcrossref{\DUrole{doc}{Build package}}}} page for more info.


\subsection{Related commands}
\label{\detokenize{pair_mie:related-commands}}
\sphinxAtStartPar
{\hyperref[\detokenize{pair_coeff::doc}]{\sphinxcrossref{\DUrole{doc}{pair\_coeff}}}}


\subsection{Default}
\label{\detokenize{pair_mie:default}}
\sphinxAtStartPar
none


\bigskip\hrule\bigskip

\phantomsection\label{\detokenize{pair_mie:mie}}
\sphinxAtStartPar
\sphinxstylestrong{(Mie)} G. Mie, Ann Phys, 316, 657 (1903).

\phantomsection\label{\detokenize{pair_mie:avendano}}
\sphinxAtStartPar
\sphinxstylestrong{(Avendano)} C. Avendano, T. Lafitte, A. Galindo, C. S. Adjiman,
G. Jackson, E. Muller, J Phys Chem B, 115, 11154 (2011).

\sphinxstepscope

\index{pair\_style mliap@\spxentry{pair\_style mliap}}\ignorespaces 
\index{pair\_style mliap/kk@\spxentry{pair\_style mliap/kk}}\ignorespaces 

\section{pair\_style mliap command}
\label{\detokenize{pair_mliap:pair-style-mliap-command}}\label{\detokenize{pair_mliap:index-1}}\label{\detokenize{pair_mliap:index-0}}\label{\detokenize{pair_mliap::doc}}
\sphinxAtStartPar
Accelerator Variants: \sphinxstyleemphasis{mliap/kk}


\subsection{Syntax}
\label{\detokenize{pair_mliap:syntax}}
\begin{sphinxVerbatim}[commandchars=\\\{\}]
\PYG{k}{pair\PYGZus{}style}\PYG{+w}{ }\PYG{n}{mliap}\PYG{+w}{ }\PYG{n}{...}\PYG{+w}{ }\PYG{n}{keyword}\PYG{+w}{ }\PYG{n}{values}\PYG{+w}{ }\PYG{n}{...}
\end{sphinxVerbatim}
\begin{itemize}
\item {} 
\sphinxAtStartPar
one or two keyword/value pairs must be appended

\item {} 
\sphinxAtStartPar
keyword = \sphinxstyleemphasis{model} or \sphinxstyleemphasis{descriptor} or \sphinxstyleemphasis{unified}
\begin{sphinxalltt}
\sphinxstyleemphasis{model} values = style filename
  style = \sphinxstyleemphasis{linear} or \sphinxstyleemphasis{quadratic} or \sphinxstyleemphasis{nn} or \sphinxstyleemphasis{mliappy}
  filename = name of file containing model definitions
\sphinxstyleemphasis{descriptor} values = style filename
  style = \sphinxstyleemphasis{sna} or \sphinxstyleemphasis{so3} or \sphinxstyleemphasis{ace}
  filename = name of file containing descriptor definitions
\sphinxstyleemphasis{unified} values = filename ghostneigh\_flag
  filename = name of file containing serialized unified Python object
  ghostneigh\_flag = 0/1 to turn off/on inclusion of ghost neighbors in neighbors list
\end{sphinxalltt}

\end{itemize}


\subsection{Examples}
\label{\detokenize{pair_mliap:examples}}
\begin{sphinxVerbatim}[commandchars=\\\{\}]
\PYG{k}{pair\PYGZus{}style}\PYG{+w}{ }\PYG{n}{mliap}\PYG{+w}{ }\PYG{n}{model}\PYG{+w}{ }\PYG{n}{linear}\PYG{+w}{ }\PYG{n}{InP.mliap.model}\PYG{+w}{ }\PYG{n}{descriptor}\PYG{+w}{ }\PYG{n}{sna}\PYG{+w}{ }\PYG{n}{InP.mliap.descriptor}
\PYG{k}{pair\PYGZus{}style}\PYG{+w}{ }\PYG{n}{mliap}\PYG{+w}{ }\PYG{n}{model}\PYG{+w}{ }\PYG{n}{quadratic}\PYG{+w}{ }\PYG{n}{W.mliap.model}\PYG{+w}{ }\PYG{n}{descriptor}\PYG{+w}{ }\PYG{n}{sna}\PYG{+w}{ }\PYG{n}{W.mliap.descriptor}
\PYG{k}{pair\PYGZus{}style}\PYG{+w}{ }\PYG{n}{mliap}\PYG{+w}{ }\PYG{n}{model}\PYG{+w}{ }\PYG{n}{nn}\PYG{+w}{ }\PYG{n}{Si.nn.mliap.model}\PYG{+w}{ }\PYG{n}{descriptor}\PYG{+w}{ }\PYG{n}{so3}\PYG{+w}{ }\PYG{n}{Si.nn.mliap.descriptor}
\PYG{k}{pair\PYGZus{}style}\PYG{+w}{ }\PYG{n}{mliap}\PYG{+w}{ }\PYG{n}{model}\PYG{+w}{ }\PYG{n}{mliappy}\PYG{+w}{ }\PYG{n}{ACE\PYGZus{}NN\PYGZus{}Pytorch.pt}\PYG{+w}{ }\PYG{n}{descriptor}\PYG{+w}{ }\PYG{n}{ace}\PYG{+w}{ }\PYG{n}{ccs\PYGZus{}single\PYGZus{}element.yace}
\PYG{k}{pair\PYGZus{}style}\PYG{+w}{ }\PYG{n}{mliap}\PYG{+w}{ }\PYG{n}{unified}\PYG{+w}{ }\PYG{n}{mliap\PYGZus{}unified\PYGZus{}lj\PYGZus{}Ar.pkl}\PYG{+w}{ }\PYG{l+m}{0}
\PYG{k}{pair\PYGZus{}coeff}\PYG{+w}{ }\PYG{o}{*}\PYG{+w}{ }\PYG{o}{*}\PYG{+w}{ }\PYG{n}{In}\PYG{+w}{ }\PYG{n}{P}
\end{sphinxVerbatim}


\subsection{Description}
\label{\detokenize{pair_mliap:description}}
\sphinxAtStartPar
Pair style \sphinxstyleemphasis{mliap} provides a general interface to families of
machine\sphinxhyphen{}learning interatomic potentials.  It allows separate definitions
of the interatomic potential functional form (\sphinxstyleemphasis{model}) and the geometric
quantities that characterize the atomic positions (\sphinxstyleemphasis{descriptor}).

\sphinxAtStartPar
By defining \sphinxstyleemphasis{model} and \sphinxstyleemphasis{descriptor} separately, it is possible to use
many different models with a given descriptor, or many different
descriptors with a given model.  The pair style currently supports \sphinxstyleemphasis{sna},
\sphinxstyleemphasis{so3} and \sphinxstyleemphasis{ace} descriptor styles, but it is straightforward to add new
descriptor styles.  By using the \sphinxstyleemphasis{unified} keyword, it is possible to
define a Python model that combines functionalities of both \sphinxstyleemphasis{model} and
\sphinxstyleemphasis{descriptor}.

\sphinxAtStartPar
The SNAP descriptor style \sphinxstyleemphasis{sna} is the same as that used by
{\hyperref[\detokenize{pair_snap::doc}]{\sphinxcrossref{\DUrole{doc}{pair\_style snap}}}}, including the linear, quadratic, and
chem variants.  The available models are \sphinxstyleemphasis{linear}, \sphinxstyleemphasis{quadratic}, \sphinxstyleemphasis{nn},
and \sphinxstyleemphasis{mliappy}.  The \sphinxstyleemphasis{mliappy} style can be used to couple python models,
e.g. PyTorch neural network energy models, and requires building LAMMPS
with the PYTHON package (see below).  In order to train a model, it is
useful to know the gradient or derivative of energy, force, and stress
w.r.t. model parameters. This information can be accessed using the
related {\hyperref[\detokenize{compute_mliap::doc}]{\sphinxcrossref{\DUrole{doc}{compute mliap}}}} command.

\sphinxAtStartPar
\DUrole{versionmodified,added}{Added in version 2Jun2022.}

\sphinxAtStartPar
The descriptor style \sphinxstyleemphasis{so3} is a descriptor that is derived from the
the smooth SO(3) power spectrum with the explicit inclusion of a radial
basis {\hyperref[\detokenize{pair_mliap:bartok2013}]{\sphinxcrossref{\DUrole{std,std-ref}{(Bartok)}}}} and {\hyperref[\detokenize{pair_mliap:zagaceta2020}]{\sphinxcrossref{\DUrole{std,std-ref}{(Zagaceta)}}}}.
The available models are \sphinxstyleemphasis{linear} and \sphinxstyleemphasis{nn}.

\sphinxAtStartPar
\DUrole{versionmodified,added}{Added in version 17Apr2024.}

\sphinxAtStartPar
The descriptor style \sphinxstyleemphasis{ace} is a class of highly general atomic
descriptors, atomic cluster expansion descriptors (ACE) from
{\hyperref[\detokenize{compute_pace:drautz19}]{\sphinxcrossref{\DUrole{std,std-ref}{(Drautz)}}}}, that include a radial basis, an angular
basis, and bases for other variables (such as chemical species) if
relevant. In descriptor style \sphinxstyleemphasis{ace}, the \sphinxstyleemphasis{ace} descriptors may be
defined up to an arbitrary body order. This descriptor style is the same
as that used in {\hyperref[\detokenize{pair_pace::doc}]{\sphinxcrossref{\DUrole{doc}{pair\_style pace}}}} and {\hyperref[\detokenize{compute_pace::doc}]{\sphinxcrossref{\DUrole{doc}{compute
pace}}}}.  The available models with \sphinxstyleemphasis{ace} in ML\sphinxhyphen{}IAP are
\sphinxstyleemphasis{linear} and \sphinxstyleemphasis{mliappy}.  The \sphinxstyleemphasis{ace} descriptors and models require
building LAMMPS with the ML\sphinxhyphen{}PACE package (see below).  The \sphinxstyleemphasis{mliappy}
model style may be used with \sphinxstyleemphasis{ace} descriptors, but it requires that
LAMMPS is also built with the PYTHON package.  As with other model
styles, the \sphinxstyleemphasis{mliappy} model style can be used to couple arbitrary python
models that use the \sphinxstyleemphasis{ace} descriptors such as Pytorch NNs.  Note that
\sphinxstyleemphasis{ALL} mliap model styles with \sphinxstyleemphasis{ace} descriptors require that descriptors
and hyperparameters are supplied in a \sphinxtitleref{.yace} or \sphinxtitleref{.ace} file, similar to
{\hyperref[\detokenize{compute_pace::doc}]{\sphinxcrossref{\DUrole{doc}{compute pace}}}}.

\sphinxAtStartPar
The pair\_style \sphinxstyleemphasis{mliap} command must be followed by two keywords \sphinxstyleemphasis{model}
and \sphinxstyleemphasis{descriptor} in either order, or the one keyword \sphinxstyleemphasis{unified}.  A
single \sphinxstyleemphasis{pair\_coeff} command is also required.  The first 2 arguments
must be * * so as to span all LAMMPS atom types.  This is followed by
a list of N arguments that specify the mapping of MLIAP element names to
LAMMPS atom types, where N is the number of LAMMPS atom types.

\sphinxAtStartPar
The \sphinxstyleemphasis{model} keyword is followed by the model style. This is followed by
a single argument specifying the model filename containing the
parameters for a set of elements.  The model filename usually ends in
the \sphinxstyleemphasis{.mliap.model} extension.  It may contain parameters for many
elements. The only requirement is that it contain at least those element
names appearing in the \sphinxstyleemphasis{pair\_coeff} command.

\sphinxAtStartPar
The top of the model file can contain any number of blank and comment
lines (start with \#), but follows a strict format after that. The first
non\sphinxhyphen{}blank non\sphinxhyphen{}comment line must contain two integers:
\begin{itemize}
\item {} 
\sphinxAtStartPar
nelems  = Number of elements

\item {} 
\sphinxAtStartPar
nparams = Number of parameters

\end{itemize}

\sphinxAtStartPar
When the \sphinxstyleemphasis{model} keyword is \sphinxstyleemphasis{linear} or \sphinxstyleemphasis{quadratic}, this is followed by
one block for each of the \sphinxstyleemphasis{nelem} elements.  Each block consists of
\sphinxstyleemphasis{nparams} parameters, one per line.  Note that this format is similar,
but not identical to that used for the {\hyperref[\detokenize{pair_snap::doc}]{\sphinxcrossref{\DUrole{doc}{pair\_style snap}}}} coefficient file.  Specifically, the line containing the
element weight and radius is omitted, since these are handled by the
\sphinxstyleemphasis{descriptor}.

\sphinxAtStartPar
When the \sphinxstyleemphasis{model} keyword is \sphinxstyleemphasis{nn} (neural networks), the model file can
contain blank and comment lines (start with \#) anywhere. The second
non\sphinxhyphen{}blank non\sphinxhyphen{}comment line must contain the string NET, followed by two
integers:
\begin{itemize}
\item {} 
\sphinxAtStartPar
ndescriptors = Number of descriptors

\item {} 
\sphinxAtStartPar
nlayers      = Number of layers (including the hidden layers and the output layer)

\end{itemize}

\sphinxAtStartPar
and followed by a sequence of a string and an integer for each layer:
\begin{itemize}
\item {} 
\sphinxAtStartPar
Activation function (linear, sigmoid, tanh or relu)

\item {} 
\sphinxAtStartPar
nnodes = Number of nodes

\end{itemize}

\sphinxAtStartPar
This is followed by one block for each of the \sphinxstyleemphasis{nelem} elements. Each
block consists of \sphinxstyleemphasis{scale0} minimum value, \sphinxstyleemphasis{scale1} (maximum \sphinxhyphen{} minimum)
value, in order to normalize the descriptors, followed by \sphinxstyleemphasis{nparams}
parameters, including \sphinxstyleemphasis{bias} and \sphinxstyleemphasis{weights} of the model, starting with
the first node of the first layer and so on, with a maximum of 30 values
per line.

\sphinxAtStartPar
The detail of \sphinxstyleemphasis{nn} module implementation can be found at {\hyperref[\detokenize{pair_mliap:yanxon2020}]{\sphinxcrossref{\DUrole{std,std-ref}{(Yanxon)}}}}.

\begin{sphinxadmonition}{note}{Notes on mliappy models}

\sphinxAtStartPar
When the \sphinxstyleemphasis{model} keyword is \sphinxstyleemphasis{mliappy}, if the filename ends in ‘.pt’,
or ‘.pth’, it will be loaded using pytorch; otherwise, it will be
loaded as a pickle file.  To load a model from memory (i.e. an
existing python object), specify the filename as “LATER”, and then
call \sphinxtitleref{lammps.mliap.load\_model(model)} from python before using the
pair style.  When using LAMMPS via the library mode, you will need to
call \sphinxtitleref{lammps.mliappy.activate\_mliappy(lmp)} on the active LAMMPS
object before the pair style is defined.  This call locates and loads
the mliap\sphinxhyphen{}specific python module that is built into LAMMPS.
\end{sphinxadmonition}

\sphinxAtStartPar
The \sphinxstyleemphasis{descriptor} keyword is followed by a descriptor style, and additional arguments.
Currently three descriptor styles are available: \sphinxstyleemphasis{sna}, \sphinxstyleemphasis{so3}, and \sphinxstyleemphasis{ace}.
\begin{itemize}
\item {} 
\sphinxAtStartPar
\sphinxstyleemphasis{sna} indicates the bispectrum component descriptors used by the Spectral
Neighbor Analysis Potential (SNAP) potentials of {\hyperref[\detokenize{pair_snap::doc}]{\sphinxcrossref{\DUrole{doc}{pair\_style snap}}}}.  A single additional argument specifies the descriptor
filename containing the parameters and setting used by the SNAP
descriptor.  The descriptor filename usually ends in the
\sphinxstyleemphasis{.mliap.descriptor} extension.

\item {} 
\sphinxAtStartPar
\sphinxstyleemphasis{so3} indicated the power spectrum component descriptors. A single additional
argument specifies the descriptor filename containing the parameters and setting.

\item {} 
\sphinxAtStartPar
\sphinxstyleemphasis{ace} indicates the atomic cluster expansion (ACE) descriptors. A single
additional argument specifies the filename containing parameters, settings, and
definitions of the ace descriptors (through tabulated basis function indices and
corresponding generalized Clebsch\sphinxhyphen{}Gordan coefficients) in the ctilde file format,
e.g. in the potential file format with \sphinxtitleref{*.ace} or \sphinxtitleref{*.yace} extensions from
{\hyperref[\detokenize{pair_pace::doc}]{\sphinxcrossref{\DUrole{doc}{pair\_style pace}}}}. Note that unlike the potential file, the
Clebsch\sphinxhyphen{}Gordan coefficients in the descriptor file supplied should \sphinxstyleemphasis{NOT} be
multiplied by linear or square root embedding terms.

\end{itemize}

\sphinxAtStartPar
The SNAP descriptor file closely follows the format of the
{\hyperref[\detokenize{pair_snap::doc}]{\sphinxcrossref{\DUrole{doc}{pair\_style snap}}}} parameter file.  The file can contain
blank and comment lines (start with \#) anywhere. Each non\sphinxhyphen{}blank
non\sphinxhyphen{}comment line must contain one keyword/value pair. The required
keywords are \sphinxstyleemphasis{rcutfac} and \sphinxstyleemphasis{twojmax}. There are many optional keywords
that are described on the {\hyperref[\detokenize{pair_snap::doc}]{\sphinxcrossref{\DUrole{doc}{pair\_style snap}}}} doc page.
In addition, the SNAP descriptor file must contain the \sphinxstyleemphasis{nelems},
\sphinxstyleemphasis{elems}, \sphinxstyleemphasis{radelems}, and \sphinxstyleemphasis{welems} keywords.  The \sphinxstyleemphasis{nelems} keyword
specifies the number of elements provided in the other three keywords.
The \sphinxstyleemphasis{elems} keyword is followed by a list of \sphinxstyleemphasis{nelems} element names that
must include the element names appearing in the \sphinxstyleemphasis{pair\_coeff} command,
but can contain other names too.  Similarly, the \sphinxstyleemphasis{radelems} and \sphinxstyleemphasis{welems}
keywords are followed by lists of \sphinxstyleemphasis{nelems} numbers giving the element
radius and element weight of each element. Obviously, the order in which
the elements are listed must be consistent for all three keywords.

\sphinxAtStartPar
The SO3 descriptor file is similar to the SNAP descriptor except that it
contains a few more arguments (e.g., \sphinxstyleemphasis{nmax} and \sphinxstyleemphasis{alpha}). The preparation
of SO3 descriptor and model files can be done with the
\sphinxhref{https://github.com/qzhu2017/PyXtal\_FF}{Pyxtal\_FF} package.

\sphinxAtStartPar
The ACE descriptor file differs from the SNAP and SO3 files. It more
closely resembles the potential file format for linear or square\sphinxhyphen{}root
embedding ACE potentials used in the {\hyperref[\detokenize{pair_pace::doc}]{\sphinxcrossref{\DUrole{doc}{pair\_style pace}}}}.
As noted above, the key difference is that the Clebsch\sphinxhyphen{}Gordan
coefficients in the descriptor file with \sphinxstyleemphasis{mliap descriptor ace} are
\sphinxstyleemphasis{NOT} multiplied by linear or square root embedding terms.  In other
words,the model is separated from the descriptor definitions and
hyperparameters.  In {\hyperref[\detokenize{pair_pace::doc}]{\sphinxcrossref{\DUrole{doc}{pair\_style pace}}}}, they are
combined.  The ACE descriptor files required by \sphinxstyleemphasis{mliap} are generated
automatically in \sphinxhref{https://github.com/FitSNAP/FitSNAP}{FitSNAP} during
linear, pytorch, etc. ACE model fitting. Additional tools are provided
there to prepare \sphinxstyleemphasis{ace} descriptor files and hyperparameters before model
fitting.  The \sphinxstyleemphasis{ace} descriptor files can also be extracted from ACE
model fits in \sphinxhref{https://github.com/ICAMS/python-ace}{python\sphinxhyphen{}ace.}.  It
is important to note that order of the types listed in {\hyperref[\detokenize{pair_coeff::doc}]{\sphinxcrossref{\DUrole{doc}{pair\_coeff}}}} must match the order of the elements/types listed in the
ACE descriptor file for all \sphinxstyleemphasis{mliap} styles when using \sphinxstyleemphasis{ace} descriptors.

\sphinxAtStartPar
See the {\hyperref[\detokenize{pair_coeff::doc}]{\sphinxcrossref{\DUrole{doc}{pair\_coeff}}}} page for alternate ways
to specify the path for these \sphinxstyleemphasis{model} and \sphinxstyleemphasis{descriptor} files.

\begin{sphinxadmonition}{note}{Note:}
\sphinxAtStartPar
To significantly reduce SO3 descriptor/force calculation time,
some properties are pre\sphinxhyphen{}computed and reused during the calculation.
These can consume a significant amount of RAM for simulations of
larger systems since their size depends on the total number of
neighbors per MPI process.
\end{sphinxadmonition}

\sphinxAtStartPar
\DUrole{versionmodified,added}{Added in version 3Nov2022.}

\sphinxAtStartPar
The \sphinxstyleemphasis{unified} keyword is followed by an argument specifying the
filename containing the serialized unified Python object and the “ghostneigh” toggle
(0/1) to disable/enable the construction of neighbors lists including
neighbors of ghost atoms. If the filename ends in ‘.pt’, or ‘.pth’, it will be loaded
using pytorch; otherwise, it will be loaded as a pickle file.
If ghostneigh is enabled, it is recommended to set {\hyperref[\detokenize{comm_modify::doc}]{\sphinxcrossref{\DUrole{doc}{comm\_modify}}}}
cutoff manually, such as in the following example.

\begin{sphinxVerbatim}[commandchars=\\\{\}]
\PYG{k}{variable }\PYG{n+nv+nvIdentifier}{ninteractions}\PYG{+w}{ }\PYG{n}{equal}\PYG{+w}{ }\PYG{l+m}{2}
\PYG{k}{variable }\PYG{n+nv+nvIdentifier}{cutdist}\PYG{+w}{ }\PYG{n}{equal}\PYG{+w}{ }\PYG{l+m}{7.5}
\PYG{k}{variable }\PYG{n+nv+nvIdentifier}{skin}\PYG{+w}{ }\PYG{n}{equal}\PYG{+w}{ }\PYG{l+m}{1.0}
\PYG{k}{variable }\PYG{n+nv+nvIdentifier}{commcut}\PYG{+w}{ }\PYG{n}{equal}\PYG{+w}{ }\PYG{n+nv}{(}\PYG{n+nv}{\PYGZdl{}\PYGZob{}ninteractions\PYGZcb{}*\PYGZdl{}\PYGZob{}cutdist\PYGZcb{}}\PYG{n+nv}{)}\PYG{o}{+}\PYG{n+nv}{\PYGZdl{}\PYGZob{}}\PYG{n+nv}{skin}\PYG{n+nv}{\PYGZcb{}}
\PYG{k}{neighbor}\PYG{+w}{ }\PYG{n+nv}{\PYGZdl{}\PYGZob{}}\PYG{n+nv}{skin}\PYG{n+nv}{\PYGZcb{}}\PYG{+w}{ }\PYG{n}{bin}
\PYG{k}{comm\PYGZus{}modify}\PYG{+w}{ }\PYG{n}{cutoff}\PYG{+w}{ }\PYG{n+nv}{\PYGZdl{}\PYGZob{}}\PYG{n+nv}{commcut}\PYG{n+nv}{\PYGZcb{}}
\end{sphinxVerbatim}

\begin{sphinxadmonition}{note}{Note:}
\sphinxAtStartPar
To load a model from memory
(i.e. an existing python object), call \sphinxtitleref{lammps.mliap.load\_unified(unified)}
from python, and then specify the filename as “EXISTS”. When using LAMMPS via
the library mode, you will need to call \sphinxtitleref{lammps.mliappy.activate\_mliappy(lmp)}
on the active LAMMPS object before the pair style is defined. This call locates
and loads the mliap\sphinxhyphen{}specific python module that is built into LAMMPS.
\end{sphinxadmonition}


\bigskip\hrule\bigskip


\sphinxAtStartPar
Styles with a \sphinxstyleemphasis{gpu}, \sphinxstyleemphasis{intel}, \sphinxstyleemphasis{kk}, \sphinxstyleemphasis{omp}, or \sphinxstyleemphasis{opt} suffix are
functionally the same as the corresponding style without the suffix.
They have been optimized to run faster, depending on your available
hardware, as discussed on the {\hyperref[\detokenize{Speed_packages::doc}]{\sphinxcrossref{\DUrole{doc}{Accelerator packages}}}}
page.  The accelerated styles take the same arguments and should
produce the same results, except for round\sphinxhyphen{}off and precision issues.

\sphinxAtStartPar
These accelerated styles are part of the GPU, INTEL, KOKKOS,
OPENMP, and OPT packages, respectively.  They are only enabled if
LAMMPS was built with those packages.  See the {\hyperref[\detokenize{Build_package::doc}]{\sphinxcrossref{\DUrole{doc}{Build package}}}} page for more info.

\sphinxAtStartPar
You can specify the accelerated styles explicitly in your input script
by including their suffix, or you can use the {\hyperref[\detokenize{Run_options::doc}]{\sphinxcrossref{\DUrole{doc}{\sphinxhyphen{}suffix command\sphinxhyphen{}line switch}}}} when you invoke LAMMPS, or you can use the
{\hyperref[\detokenize{suffix::doc}]{\sphinxcrossref{\DUrole{doc}{suffix}}}} command in your input script.

\sphinxAtStartPar
See the {\hyperref[\detokenize{Speed_packages::doc}]{\sphinxcrossref{\DUrole{doc}{Accelerator packages}}}} page for more
instructions on how to use the accelerated styles effectively.


\bigskip\hrule\bigskip



\subsection{Mixing, shift, table, tail correction, restart, rRESPA info}
\label{\detokenize{pair_mliap:mixing-shift-table-tail-correction-restart-rrespa-info}}
\sphinxAtStartPar
For atom type pairs I,J and I != J, where types I and J correspond to
two different element types, mixing is performed by LAMMPS with
user\sphinxhyphen{}specifiable parameters as described above.  You never need to
specify a pair\_coeff command with I != J arguments for this style.

\sphinxAtStartPar
This pair style does not support the {\hyperref[\detokenize{pair_modify::doc}]{\sphinxcrossref{\DUrole{doc}{pair\_modify}}}}
shift, table, and tail options.

\sphinxAtStartPar
This pair style does not write its information to {\hyperref[\detokenize{restart::doc}]{\sphinxcrossref{\DUrole{doc}{binary restart
files}}}}, since it is stored in potential files.  Thus, you need
to re\sphinxhyphen{}specify the pair\_style and pair\_coeff commands in an input script
that reads a restart file.

\sphinxAtStartPar
This pair style can only be used via the \sphinxstyleemphasis{pair} keyword of the
{\hyperref[\detokenize{run_style::doc}]{\sphinxcrossref{\DUrole{doc}{run\_style respa}}}} command.  It does not support the
\sphinxstyleemphasis{inner}, \sphinxstyleemphasis{middle}, \sphinxstyleemphasis{outer} keywords.


\bigskip\hrule\bigskip



\subsection{Restrictions}
\label{\detokenize{pair_mliap:restrictions}}
\sphinxAtStartPar
This pair style is part of the ML\sphinxhyphen{}IAP package.  It is only enabled if
LAMMPS was built with that package. In addition, building LAMMPS with
the ML\sphinxhyphen{}IAP package requires building LAMMPS with the ML\sphinxhyphen{}SNAP package.
The \sphinxstyleemphasis{mliappy} model requires building LAMMPS with the PYTHON package.
The \sphinxstyleemphasis{ace} descriptor requires building LAMMPS with the ML\sphinxhyphen{}PACE package.
See the {\hyperref[\detokenize{Build_package::doc}]{\sphinxcrossref{\DUrole{doc}{Build package}}}} page for more info. Note
that \sphinxtitleref{pair\_mliap/kk} acceleration will \sphinxstyleemphasis{not} invoke the \sphinxtitleref{kk}
accelerated variants of SNAP or ACE descriptors.


\subsection{Related commands}
\label{\detokenize{pair_mliap:related-commands}}
\sphinxAtStartPar
{\hyperref[\detokenize{pair_snap::doc}]{\sphinxcrossref{\DUrole{doc}{pair\_style snap}}}}, {\hyperref[\detokenize{compute_mliap::doc}]{\sphinxcrossref{\DUrole{doc}{compute mliap}}}}


\subsection{Default}
\label{\detokenize{pair_mliap:default}}
\sphinxAtStartPar
none


\bigskip\hrule\bigskip

\phantomsection\label{\detokenize{pair_mliap:bartok2013}}
\sphinxAtStartPar
\sphinxstylestrong{(Bartok2013)} Bartok, Kondor, Csanyi, Phys Rev B, 87, 184115 (2013).

\phantomsection\label{\detokenize{pair_mliap:zagaceta2020}}
\sphinxAtStartPar
\sphinxstylestrong{(Zagaceta2020)} Zagaceta, Yanxon, Zhu, J Appl Phys, 128, 045113 (2020).

\phantomsection\label{\detokenize{pair_mliap:yanxon2020}}
\sphinxAtStartPar
\sphinxstylestrong{(Yanxon2020)} Yanxon, Zagaceta, Tang, Matteson, Zhu, Mach. Learn.: Sci. Technol. 2, 027001 (2020).

\sphinxstepscope

\index{pair\_style momb@\spxentry{pair\_style momb}}\ignorespaces 

\section{pair\_style momb command}
\label{\detokenize{pair_momb:pair-style-momb-command}}\label{\detokenize{pair_momb:index-0}}\label{\detokenize{pair_momb::doc}}

\subsection{Syntax}
\label{\detokenize{pair_momb:syntax}}
\begin{sphinxVerbatim}[commandchars=\\\{\}]
\PYG{k}{pair\PYGZus{}style}\PYG{+w}{ }\PYG{n}{momb}\PYG{+w}{ }\PYG{n}{cutoff}\PYG{+w}{ }\PYG{n}{s6}\PYG{+w}{ }\PYG{n}{d}
\end{sphinxVerbatim}
\begin{itemize}
\item {} 
\sphinxAtStartPar
cutoff = global cutoff (distance units)

\item {} 
\sphinxAtStartPar
s6 = global scaling factor of the exchange\sphinxhyphen{}correlation functional used (unitless)

\item {} 
\sphinxAtStartPar
d = damping scaling factor of Grimme’s method (unitless)

\end{itemize}


\subsection{Examples}
\label{\detokenize{pair_momb:examples}}
\begin{sphinxVerbatim}[commandchars=\\\{\}]
\PYG{k}{pair\PYGZus{}style}\PYG{+w}{ }\PYG{n}{momb}\PYG{+w}{ }\PYG{l+m}{12.0}\PYG{+w}{ }\PYG{l+m}{0.75}\PYG{+w}{ }\PYG{l+m}{20.0}
\PYG{k}{pair\PYGZus{}style}\PYG{+w}{ }\PYG{n}{hybrid}\PYG{o}{/}\PYG{n}{overlay}\PYG{+w}{ }\PYG{n}{eam}\PYG{o}{/}\PYG{n}{fs}\PYG{+w}{ }\PYG{n}{lj}\PYG{o}{/}\PYG{n}{charmm}\PYG{o}{/}\PYG{n}{coul}\PYG{o}{/}\PYG{n}{long}\PYG{+w}{ }\PYG{l+m}{10.0}\PYG{+w}{ }\PYG{l+m}{12.0}\PYG{+w}{ }\PYG{n}{momb}\PYG{+w}{ }\PYG{l+m}{12.0}\PYG{+w}{ }\PYG{l+m}{0.75}\PYG{+w}{ }\PYG{l+m}{20.0}\PYG{+w}{ }\PYG{n}{morse}\PYG{+w}{ }\PYG{l+m}{5.5}

\PYG{k}{pair\PYGZus{}coeff}\PYG{+w}{ }\PYG{l+m}{1}\PYG{+w}{ }\PYG{l+m}{2}\PYG{+w}{ }\PYG{n}{momb}\PYG{+w}{ }\PYG{l+m}{0.0}\PYG{+w}{ }\PYG{l+m}{1.0}\PYG{+w}{ }\PYG{l+m}{1.0}\PYG{+w}{ }\PYG{l+m}{10.2847}\PYG{+w}{ }\PYG{l+m}{2.361}
\end{sphinxVerbatim}


\subsection{Description}
\label{\detokenize{pair_momb:description}}
\sphinxAtStartPar
Style \sphinxstyleemphasis{momb} computes pairwise van der Waals (vdW) and short\sphinxhyphen{}range
interactions using the Morse potential and {\hyperref[\detokenize{pair_momb:grimme}]{\sphinxcrossref{\DUrole{std,std-ref}{(Grimme)}}}} method
implemented in the Many\sphinxhyphen{}Body Metal\sphinxhyphen{}Organic (MOMB) force field
described comprehensively in {\hyperref[\detokenize{pair_momb:fichthorn}]{\sphinxcrossref{\DUrole{std,std-ref}{(Fichthorn)}}}} and
{\hyperref[\detokenize{pair_momb:zhou5}]{\sphinxcrossref{\DUrole{std,std-ref}{(Zhou)}}}}. Grimme’s method is widely used to correct for
dispersion in density functional theory calculations.
\begin{equation*}
\begin{split} E & = D_0 [\exp^{-2 \alpha (r-r_0)} - 2\exp^{-\alpha (r-r_0)}] - s_6 \frac{C_6}{r^6} f_{damp}(r,R_r) \\
f_{damp}(r,R_r) & = \frac{1}{1 + \exp^{-d(r/R_r - 1)}}\end{split}
\end{equation*}
\sphinxAtStartPar
For the \sphinxstyleemphasis{momb} pair style, the following coefficients must be defined
for each pair of atoms types via the {\hyperref[\detokenize{pair_coeff::doc}]{\sphinxcrossref{\DUrole{doc}{pair\_coeff}}}}
command as in the examples above, or in the data file or restart files
read by the {\hyperref[\detokenize{read_data::doc}]{\sphinxcrossref{\DUrole{doc}{read\_data}}}} as described below:
\begin{itemize}
\item {} 
\sphinxAtStartPar
\(D_0\) (energy units)

\item {} 
\sphinxAtStartPar
\(\alpha\) (1/distance units)

\item {} 
\sphinxAtStartPar
\(r_0\) (distance units)

\item {} 
\sphinxAtStartPar
\(C_6\) (energy*distance\textasciicircum{}6 units)

\item {} 
\sphinxAtStartPar
\(R_r\) (distance units, typically sum of atomic vdW radii)

\end{itemize}


\bigskip\hrule\bigskip



\subsection{Restrictions}
\label{\detokenize{pair_momb:restrictions}}
\sphinxAtStartPar
This style is part of the EXTRA\sphinxhyphen{}PAIR package. It is only enabled if
LAMMPS is built with that package. See the {\hyperref[\detokenize{Build_package::doc}]{\sphinxcrossref{\DUrole{doc}{Build package}}}} page on for more info.


\subsection{Related commands}
\label{\detokenize{pair_momb:related-commands}}
\sphinxAtStartPar
{\hyperref[\detokenize{pair_coeff::doc}]{\sphinxcrossref{\DUrole{doc}{pair\_coeff}}}}, {\hyperref[\detokenize{pair_morse::doc}]{\sphinxcrossref{\DUrole{doc}{pair\_style morse}}}}


\subsection{Default}
\label{\detokenize{pair_momb:default}}
\sphinxAtStartPar
none


\bigskip\hrule\bigskip

\phantomsection\label{\detokenize{pair_momb:grimme}}
\sphinxAtStartPar
\sphinxstylestrong{(Grimme)} Grimme, J Comput Chem, 27(15), 1787\sphinxhyphen{}1799 (2006).

\phantomsection\label{\detokenize{pair_momb:fichthorn}}
\sphinxAtStartPar
\sphinxstylestrong{(Fichthorn)} Fichthorn, Balankura, Qi, CrystEngComm, 18(29), 5410\sphinxhyphen{}5417 (2016).

\phantomsection\label{\detokenize{pair_momb:zhou5}}
\sphinxAtStartPar
\sphinxstylestrong{(Zhou)} Zhou, Saidi, Fichthorn, J Phys Chem C, 118(6), 3366\sphinxhyphen{}3374 (2014).

\sphinxstepscope

\index{pair\_style morse@\spxentry{pair\_style morse}}\ignorespaces 
\index{pair\_style morse/gpu@\spxentry{pair\_style morse/gpu}}\ignorespaces 
\index{pair\_style morse/omp@\spxentry{pair\_style morse/omp}}\ignorespaces 
\index{pair\_style morse/opt@\spxentry{pair\_style morse/opt}}\ignorespaces 
\index{pair\_style morse/smooth/linear@\spxentry{pair\_style morse/smooth/linear}}\ignorespaces 
\index{pair\_style morse/smooth/linear/omp@\spxentry{pair\_style morse/smooth/linear/omp}}\ignorespaces 
\index{pair\_style morse/kk@\spxentry{pair\_style morse/kk}}\ignorespaces 

\section{pair\_style morse command}
\label{\detokenize{pair_morse:pair-style-morse-command}}\label{\detokenize{pair_morse:index-6}}\label{\detokenize{pair_morse:index-5}}\label{\detokenize{pair_morse:index-4}}\label{\detokenize{pair_morse:index-3}}\label{\detokenize{pair_morse:index-2}}\label{\detokenize{pair_morse:index-1}}\label{\detokenize{pair_morse:index-0}}\label{\detokenize{pair_morse::doc}}
\sphinxAtStartPar
Accelerator Variants: \sphinxstyleemphasis{morse/gpu}, \sphinxstyleemphasis{morse/omp}, \sphinxstyleemphasis{morse/opt}, \sphinxstyleemphasis{morse/kk}


\section{pair\_style morse/smooth/linear command}
\label{\detokenize{pair_morse:pair-style-morse-smooth-linear-command}}
\sphinxAtStartPar
Accelerator Variants: \sphinxstyleemphasis{morse/smooth/linear/omp}


\subsection{Syntax}
\label{\detokenize{pair_morse:syntax}}
\begin{sphinxVerbatim}[commandchars=\\\{\}]
\PYG{k}{pair\PYGZus{}style}\PYG{+w}{ }\PYG{n}{style}\PYG{+w}{ }\PYG{n}{args}
\end{sphinxVerbatim}
\begin{itemize}
\item {} 
\sphinxAtStartPar
style = \sphinxstyleemphasis{morse} or \sphinxstyleemphasis{morse/smooth/linear} or \sphinxstyleemphasis{morse/soft}

\item {} 
\sphinxAtStartPar
args = list of arguments for a particular style

\end{itemize}
\begin{sphinxalltt}
\sphinxstyleemphasis{morse} args = cutoff
  cutoff = global cutoff for Morse interactions (distance units)
\sphinxstyleemphasis{morse/smooth/linear} args = cutoff
  cutoff = global cutoff for Morse interactions (distance units)
\end{sphinxalltt}


\subsection{Examples}
\label{\detokenize{pair_morse:examples}}
\begin{sphinxVerbatim}[commandchars=\\\{\}]
\PYG{k}{pair\PYGZus{}style}\PYG{+w}{ }\PYG{n}{morse}\PYG{+w}{ }\PYG{l+m}{2.5}
\PYG{k}{pair\PYGZus{}style}\PYG{+w}{ }\PYG{n}{morse}\PYG{o}{/}\PYG{n}{smooth}\PYG{o}{/}\PYG{n}{linear}\PYG{+w}{ }\PYG{l+m}{2.5}
\PYG{k}{pair\PYGZus{}coeff}\PYG{+w}{ }\PYG{o}{*}\PYG{+w}{ }\PYG{o}{*}\PYG{+w}{ }\PYG{l+m}{100.0}\PYG{+w}{ }\PYG{l+m}{2.0}\PYG{+w}{ }\PYG{l+m}{1.5}
\PYG{k}{pair\PYGZus{}coeff}\PYG{+w}{ }\PYG{l+m}{1}\PYG{+w}{ }\PYG{l+m}{1}\PYG{+w}{ }\PYG{l+m}{100.0}\PYG{+w}{ }\PYG{l+m}{2.0}\PYG{+w}{ }\PYG{l+m}{1.5}\PYG{+w}{ }\PYG{l+m}{3.0}
\end{sphinxVerbatim}


\subsection{Description}
\label{\detokenize{pair_morse:description}}
\sphinxAtStartPar
Style \sphinxstyleemphasis{morse} computes pairwise interactions with the formula
\begin{equation*}
\begin{split}E = D_0 \left[ e^{- 2 \alpha (r - r_0)} - 2 e^{- \alpha (r - r_0)} \right]
    \qquad r < r_c\end{split}
\end{equation*}
\sphinxAtStartPar
\(r_c\) is the cutoff.

\sphinxAtStartPar
The following coefficients must be defined for each pair of atoms
types via the {\hyperref[\detokenize{pair_coeff::doc}]{\sphinxcrossref{\DUrole{doc}{pair\_coeff}}}} command as in the examples
above, or in the data file or restart files read by the
{\hyperref[\detokenize{read_data::doc}]{\sphinxcrossref{\DUrole{doc}{read\_data}}}} or {\hyperref[\detokenize{read_restart::doc}]{\sphinxcrossref{\DUrole{doc}{read\_restart}}}}
commands:
\begin{itemize}
\item {} 
\sphinxAtStartPar
\(D_0\) (energy units)

\item {} 
\sphinxAtStartPar
\(\alpha\) (1/distance units)

\item {} 
\sphinxAtStartPar
\(r_0\) (distance units)

\item {} 
\sphinxAtStartPar
cutoff (distance units)

\end{itemize}

\sphinxAtStartPar
The last coefficient is optional.  If not specified, the global morse
cutoff is used.


\bigskip\hrule\bigskip


\sphinxAtStartPar
The \sphinxstyleemphasis{morse/smooth/linear} variant is similar to the lj/smooth/linear
variant in that it adds to the potential a shift and a linear term
so that both, potential energy and force, go to zero at the cut\sphinxhyphen{}off:
\begin{equation*}
\begin{split}\phi\left(r\right) & =  D_0 \left[ e^{- 2 \alpha (r - r_0)} - 2 e^{- \alpha (r - r_0)} \right] \qquad r < r_c \\
E\left(r\right) & =  \phi\left(r\right)  - \phi\left(r_c\right) - \left(r - r_c\right) \left.\frac{d\phi}{d r} \right|_{r=r_c}       \qquad r < r_c\end{split}
\end{equation*}
\sphinxAtStartPar
The syntax of the pair\_style and pair\_coeff commands are the same for
the \sphinxstyleemphasis{morse} and \sphinxstyleemphasis{morse/smooth/linear} styles.


\bigskip\hrule\bigskip


\sphinxAtStartPar
A version of the \sphinxstyleemphasis{morse} style with a soft core, \sphinxstyleemphasis{morse/soft},
suitable for use in free energy calculations, is part of the FEP
package and is documented with the {\hyperref[\detokenize{pair_fep_soft::doc}]{\sphinxcrossref{\DUrole{doc}{pair\_style */soft}}}} styles. The version with soft core is only available if
LAMMPS was built with that package. See the {\hyperref[\detokenize{Build_package::doc}]{\sphinxcrossref{\DUrole{doc}{Build package}}}} page for more info.


\bigskip\hrule\bigskip


\sphinxAtStartPar
Styles with a \sphinxstyleemphasis{gpu}, \sphinxstyleemphasis{intel}, \sphinxstyleemphasis{kk}, \sphinxstyleemphasis{omp}, or \sphinxstyleemphasis{opt} suffix are
functionally the same as the corresponding style without the suffix.
They have been optimized to run faster, depending on your available
hardware, as discussed on the {\hyperref[\detokenize{Speed_packages::doc}]{\sphinxcrossref{\DUrole{doc}{Accelerator packages}}}}
page.  The accelerated styles take the same arguments and should
produce the same results, except for round\sphinxhyphen{}off and precision issues.

\sphinxAtStartPar
These accelerated styles are part of the GPU, INTEL, KOKKOS,
OPENMP, and OPT packages, respectively.  They are only enabled if
LAMMPS was built with those packages.  See the {\hyperref[\detokenize{Build_package::doc}]{\sphinxcrossref{\DUrole{doc}{Build package}}}} page for more info.

\sphinxAtStartPar
You can specify the accelerated styles explicitly in your input script
by including their suffix, or you can use the {\hyperref[\detokenize{Run_options::doc}]{\sphinxcrossref{\DUrole{doc}{\sphinxhyphen{}suffix command\sphinxhyphen{}line switch}}}} when you invoke LAMMPS, or you can use the
{\hyperref[\detokenize{suffix::doc}]{\sphinxcrossref{\DUrole{doc}{suffix}}}} command in your input script.

\sphinxAtStartPar
See the {\hyperref[\detokenize{Speed_packages::doc}]{\sphinxcrossref{\DUrole{doc}{Accelerator packages}}}} page for more
instructions on how to use the accelerated styles effectively.


\bigskip\hrule\bigskip



\subsection{Mixing, shift, table, tail correction, restart, rRESPA info}
\label{\detokenize{pair_morse:mixing-shift-table-tail-correction-restart-rrespa-info}}
\sphinxAtStartPar
None of these pair styles support mixing.  Thus, coefficients for all
I,J pairs must be specified explicitly.

\sphinxAtStartPar
All of these pair styles support the {\hyperref[\detokenize{pair_modify::doc}]{\sphinxcrossref{\DUrole{doc}{pair\_modify}}}}
shift option for the energy of the pair interaction.

\sphinxAtStartPar
The {\hyperref[\detokenize{pair_modify::doc}]{\sphinxcrossref{\DUrole{doc}{pair\_modify}}}} table options are not relevant for
the Morse pair styles.

\sphinxAtStartPar
None of these pair styles support the {\hyperref[\detokenize{pair_modify::doc}]{\sphinxcrossref{\DUrole{doc}{pair\_modify}}}}
tail option for adding long\sphinxhyphen{}range tail corrections to energy and
pressure.

\sphinxAtStartPar
All of these pair styles write their information to {\hyperref[\detokenize{restart::doc}]{\sphinxcrossref{\DUrole{doc}{binary restart files}}}}, so pair\_style and pair\_coeff commands do not need
to be specified in an input script that reads a restart file.

\sphinxAtStartPar
These pair styles can only be used via the \sphinxstyleemphasis{pair} keyword of the
{\hyperref[\detokenize{run_style::doc}]{\sphinxcrossref{\DUrole{doc}{run\_style respa}}}} command.  They do not support the
\sphinxstyleemphasis{inner}, \sphinxstyleemphasis{middle}, \sphinxstyleemphasis{outer} keywords.


\bigskip\hrule\bigskip



\subsection{Restrictions}
\label{\detokenize{pair_morse:restrictions}}
\sphinxAtStartPar
The \sphinxstyleemphasis{morse/smooth/linear} pair style is only enabled if LAMMPS was
built with the EXTRA\sphinxhyphen{}PAIR package.
See the {\hyperref[\detokenize{Build_package::doc}]{\sphinxcrossref{\DUrole{doc}{Build package}}}} page for more info.


\subsection{Related commands}
\label{\detokenize{pair_morse:related-commands}}
\sphinxAtStartPar
{\hyperref[\detokenize{pair_coeff::doc}]{\sphinxcrossref{\DUrole{doc}{pair\_coeff}}}}, {\hyperref[\detokenize{pair_fep_soft::doc}]{\sphinxcrossref{\DUrole{doc}{pair\_style */soft}}}}


\subsection{Default}
\label{\detokenize{pair_morse:default}}
\sphinxAtStartPar
none

\sphinxstepscope

\index{pair\_style multi/lucy@\spxentry{pair\_style multi/lucy}}\ignorespaces 

\section{pair\_style multi/lucy command}
\label{\detokenize{pair_multi_lucy:pair-style-multi-lucy-command}}\label{\detokenize{pair_multi_lucy:index-0}}\label{\detokenize{pair_multi_lucy::doc}}

\subsection{Syntax}
\label{\detokenize{pair_multi_lucy:syntax}}
\begin{sphinxVerbatim}[commandchars=\\\{\}]
\PYG{k}{pair\PYGZus{}style}\PYG{+w}{ }\PYG{n}{multi}\PYG{o}{/}\PYG{n}{lucy}\PYG{+w}{ }\PYG{n}{style}\PYG{+w}{ }\PYG{n}{N}\PYG{+w}{ }\PYG{n}{keyword}\PYG{+w}{ }\PYG{n}{...}
\end{sphinxVerbatim}
\begin{itemize}
\item {} 
\sphinxAtStartPar
style = \sphinxstyleemphasis{lookup} or \sphinxstyleemphasis{linear} = method of interpolation

\item {} 
\sphinxAtStartPar
N = use N values in \sphinxstyleemphasis{lookup}, \sphinxstyleemphasis{linear} tables

\end{itemize}


\subsection{Examples}
\label{\detokenize{pair_multi_lucy:examples}}
\begin{sphinxVerbatim}[commandchars=\\\{\}]
\PYG{k}{pair\PYGZus{}style}\PYG{+w}{ }\PYG{n}{multi}\PYG{o}{/}\PYG{n}{lucy}\PYG{+w}{ }\PYG{n}{linear}\PYG{+w}{ }\PYG{l+m}{1000}
\PYG{k}{pair\PYGZus{}coeff}\PYG{+w}{ }\PYG{o}{*}\PYG{+w}{ }\PYG{o}{*}\PYG{+w}{ }\PYG{n}{multibody.table}\PYG{+w}{ }\PYG{n}{ENTRY1}\PYG{+w}{ }\PYG{l+m}{7.0}
\end{sphinxVerbatim}


\subsection{Description}
\label{\detokenize{pair_multi_lucy:description}}
\sphinxAtStartPar
Style \sphinxstyleemphasis{multi/lucy} computes a density\sphinxhyphen{}dependent force following from
the many\sphinxhyphen{}body form described in {\hyperref[\detokenize{pair_multi_lucy:moore1}]{\sphinxcrossref{\DUrole{std,std-ref}{(Moore)}}}} and
{\hyperref[\detokenize{pair_multi_lucy:warren1}]{\sphinxcrossref{\DUrole{std,std-ref}{(Warren)}}}} as
\begin{equation*}
\begin{split}F_{i}^{DD}(\rho_i,\rho_j,r_{ij}) = \frac{1}{2} \omega_{DD}\left(r_{ij}\right)
\left[A\left(\rho_i\right) + A\left(\rho_j\right)\right]e_{ij}\end{split}
\end{equation*}
\sphinxAtStartPar
which consists of a density\sphinxhyphen{}dependent function, \(A(\rho)\), and a
radial\sphinxhyphen{}dependent weight function, \(\omega_{DD}(r_{ij})\).  The
radial\sphinxhyphen{}dependent weight function, \(\omega_{DD}(r_{ij})\), is taken
as the Lucy function:
\begin{equation*}
\begin{split}\omega_{DD}\left(r_{ij}\right) = \left(1+\frac{3r_{ij}}{r_{cut}}\right)\left(1+\frac{r_{ij}}{r_{cut}}\right)^3\end{split}
\end{equation*}
\sphinxAtStartPar
The density\sphinxhyphen{}dependent energy for a given particle is given by:
\begin{equation*}
\begin{split}u_{i}^{DD}\left(\rho_{i}\right) = \frac{\pi r_{cut}^4}{84} \int_{\rho_0}^{\rho_i} A\left(\rho'\right) d\rho'\end{split}
\end{equation*}
\sphinxAtStartPar
See the supporting information of {\hyperref[\detokenize{pair_multi_lucy:brennan1}]{\sphinxcrossref{\DUrole{std,std-ref}{(Brennan)}}}} or the
publication by {\hyperref[\detokenize{pair_multi_lucy:moore1}]{\sphinxcrossref{\DUrole{std,std-ref}{(Moore)}}}} for more details on the functional
form.

\sphinxAtStartPar
An interpolation table is used to evaluate the density\sphinxhyphen{}dependent energy
(\(\int A(\rho') d\rho'\)) and force (\(A(\rho')\)).  Note that
the prefactor to the energy is computed after the interpolation, thus
the \(\int A(\rho') d \rho'\) will have units of energy / length\textasciicircum{}4.

\sphinxAtStartPar
The interpolation table is created as a pre\sphinxhyphen{}computation by fitting
cubic splines to the file values and interpolating the
density\sphinxhyphen{}dependent energy and force at each of \sphinxstyleemphasis{N} densities.  During a
simulation, the tables are used to interpolate the density\sphinxhyphen{}dependent
energy and force as needed for each pair of particles separated by a
distance \sphinxstyleemphasis{R}.  The interpolation is done in one of 2 styles: \sphinxstyleemphasis{lookup}
and \sphinxstyleemphasis{linear}.

\sphinxAtStartPar
For the \sphinxstyleemphasis{lookup} style, the density is used to find the nearest table
entry, which is the density\sphinxhyphen{}dependent energy and force.

\sphinxAtStartPar
For the \sphinxstyleemphasis{linear} style, the density is used to find the 2 surrounding
table values from which the density\sphinxhyphen{}dependent energy and force are
computed by linear interpolation.

\sphinxAtStartPar
The following coefficients must be defined for each pair of atoms
types via the {\hyperref[\detokenize{pair_coeff::doc}]{\sphinxcrossref{\DUrole{doc}{pair\_coeff}}}} command as in the examples
above.
\begin{itemize}
\item {} 
\sphinxAtStartPar
filename

\item {} 
\sphinxAtStartPar
keyword

\item {} 
\sphinxAtStartPar
cutoff (distance units)

\end{itemize}

\sphinxAtStartPar
The filename specifies a file containing the tabulated
density\sphinxhyphen{}dependent energy and force.  The keyword specifies a section
of the file.  The cutoff is an optional coefficient.  If not
specified, the outer cutoff in the table itself (see below) will be
used to build an interpolation table that extend to the largest
tabulated distance.  If specified, only file values up to the cutoff
are used to create the interpolation table.  The format of this file
is described below.


\bigskip\hrule\bigskip


\sphinxAtStartPar
The format of a tabulated file is a series of one or more sections,
defined as follows (without the parenthesized comments):
\begin{sphinxalltt}
\# Density\sphinxhyphen{}dependent function (one or more comment or blank lines)

DD\sphinxhyphen{}FUNCTION                (keyword is first text on line)
N 500 R 1.0 10.0           (N, R, RSQ parameters)
                           (blank)
1 1.0 25.5 102.34          (index, density, energy/r\textasciicircum{}4, force)
2 1.02 23.4 98.5
...
500 10.0 0.001 0.003
\end{sphinxalltt}

\sphinxAtStartPar
A section begins with a non\sphinxhyphen{}blank line whose first character is not a
“\#”; blank lines or lines starting with “\#” can be used as comments
between sections.  The first line begins with a keyword which
identifies the section.  The line can contain additional text, but the
initial text must match the argument specified in the pair\_coeff
command.  The next line lists (in any order) one or more parameters
for the table.  Each parameter is a keyword followed by one or more
numeric values.

\sphinxAtStartPar
The parameter “N” is required and its value is the number of table
entries that follow.  Note that this may be different than the \sphinxstyleemphasis{N}
specified in the {\hyperref[\detokenize{pair_multi_lucy::doc}]{\sphinxcrossref{\DUrole{doc}{pair\_style multi/lucy}}}} command.
Let Ntable = \sphinxstyleemphasis{N} in the pair\_style command, and Nfile = “N” in the
tabulated file.  What LAMMPS does is a preliminary interpolation by
creating splines using the Nfile tabulated values as nodal points.  It
uses these to interpolate the density\sphinxhyphen{}dependent energy and force at
Ntable different points.  The resulting tables of length Ntable are
then used as described above, when computing the density\sphinxhyphen{}dependent
energy and force.  This means that if you want the interpolation
tables of length Ntable to match exactly what is in the tabulated file
(with effectively no preliminary interpolation), you should set Ntable
= Nfile, and use the “RSQ” parameter.  This is because the internal
table abscissa is always RSQ (separation distance squared), for
efficient lookup.

\sphinxAtStartPar
All other parameters are optional.  If “R” or “RSQ” does
not appear, then the distances in each line of the table are used
as\sphinxhyphen{}is to perform spline interpolation.  In this case, the table values
can be spaced in \sphinxstyleemphasis{density} uniformly or however you wish to position table
values in regions of large gradients.

\sphinxAtStartPar
If used, the parameters “R” or “RSQ” are followed by 2 values \sphinxstyleemphasis{rlo} and
\sphinxstyleemphasis{rhi}.  If specified, the density associated with each
density\sphinxhyphen{}dependent energy and force value is computed from these 2 values
(at high accuracy), rather than using the (low\sphinxhyphen{}accuracy) value listed in
each line of the table.  The density values in the table file are
ignored in this case.  For “R”, distances uniformly spaced between \sphinxstyleemphasis{rlo}
and \sphinxstyleemphasis{rhi} are computed; for “RSQ”, squared distances uniformly spaced
between \sphinxstyleemphasis{rlo*rlo} and \sphinxstyleemphasis{rhi*rhi} are computed.

\begin{sphinxadmonition}{note}{Note:}
\sphinxAtStartPar
If you use “R” or “RSQ”, the tabulated distance values in the file
are effectively ignored, and replaced by new values as described in
the previous paragraph.  If the density value in the table is not
very close to the new value (i.e. round\sphinxhyphen{}off difference), then you
will be assigning density\sphinxhyphen{}dependent energy and force values to a
different density, which is probably not what you want.  LAMMPS will
warn if this is occurring.
\end{sphinxadmonition}

\sphinxAtStartPar
Following a blank line, the next N lines list the tabulated values.
On each line, the first value is the index from 1 to N, the second value is
r (in density units), the third value is the density\sphinxhyphen{}dependent function value
(in energy units / length\textasciicircum{}4), and the fourth is the force (in force units).  The
density values must increase from one line to the next.

\sphinxAtStartPar
Note that one file can contain many sections, each with a tabulated
potential.  LAMMPS reads the file section by section until it finds
one that matches the specified keyword.


\bigskip\hrule\bigskip



\subsection{Mixing, shift, table, tail correction, restart, rRESPA info}
\label{\detokenize{pair_multi_lucy:mixing-shift-table-tail-correction-restart-rrespa-info}}
\sphinxAtStartPar
This pair style does not support mixing.  Thus, coefficients for all
I,J pairs must be specified explicitly.

\sphinxAtStartPar
The {\hyperref[\detokenize{pair_modify::doc}]{\sphinxcrossref{\DUrole{doc}{pair\_modify}}}} shift, table, and tail options are
not relevant for this pair style.

\sphinxAtStartPar
This pair style writes the settings for the “pair\_style multi/lucy” command
to {\hyperref[\detokenize{restart::doc}]{\sphinxcrossref{\DUrole{doc}{binary restart files}}}}, so a pair\_style command does
not need to specified in an input script that reads a restart file.
However, the coefficient information is not stored in the restart
file, since it is tabulated in the potential files.  Thus, pair\_coeff
commands do need to be specified in the restart input script.

\sphinxAtStartPar
This pair style can only be used via the \sphinxstyleemphasis{pair} keyword of the
{\hyperref[\detokenize{run_style::doc}]{\sphinxcrossref{\DUrole{doc}{run\_style respa}}}} command.  It does not support the
\sphinxstyleemphasis{inner}, \sphinxstyleemphasis{middle}, \sphinxstyleemphasis{outer} keywords.


\bigskip\hrule\bigskip



\subsection{Restrictions}
\label{\detokenize{pair_multi_lucy:restrictions}}
\sphinxAtStartPar
This command is part of the DPD\sphinxhyphen{}REACT package.  It is only enabled if
LAMMPS was built with that package.  See the {\hyperref[\detokenize{Build_package::doc}]{\sphinxcrossref{\DUrole{doc}{Build package}}}} page for more info.


\subsection{Related commands}
\label{\detokenize{pair_multi_lucy:related-commands}}
\sphinxAtStartPar
{\hyperref[\detokenize{pair_coeff::doc}]{\sphinxcrossref{\DUrole{doc}{pair\_coeff}}}}


\subsection{Default}
\label{\detokenize{pair_multi_lucy:default}}
\sphinxAtStartPar
none


\bigskip\hrule\bigskip

\phantomsection\label{\detokenize{pair_multi_lucy:warren1}}
\sphinxAtStartPar
\sphinxstylestrong{(Warren)} Warren, Phys Rev E, 68, 066702 (2003).

\phantomsection\label{\detokenize{pair_multi_lucy:brennan1}}
\sphinxAtStartPar
\sphinxstylestrong{(Brennan)} Brennan, J Chem Phys Lett, 5, 2144\sphinxhyphen{}2149 (2014).

\phantomsection\label{\detokenize{pair_multi_lucy:moore1}}
\sphinxAtStartPar
\sphinxstylestrong{(Moore)} Moore, J Chem Phys, 144, 104501 (2016).

\sphinxstepscope

\index{pair\_style multi/lucy/rx@\spxentry{pair\_style multi/lucy/rx}}\ignorespaces 
\index{pair\_style multi/lucy/rx/kk@\spxentry{pair\_style multi/lucy/rx/kk}}\ignorespaces 

\section{pair\_style multi/lucy/rx command}
\label{\detokenize{pair_multi_lucy_rx:pair-style-multi-lucy-rx-command}}\label{\detokenize{pair_multi_lucy_rx:index-1}}\label{\detokenize{pair_multi_lucy_rx:index-0}}\label{\detokenize{pair_multi_lucy_rx::doc}}
\sphinxAtStartPar
Accelerator Variants: \sphinxstyleemphasis{multi/lucy/rx/kk}


\subsection{Syntax}
\label{\detokenize{pair_multi_lucy_rx:syntax}}
\begin{sphinxVerbatim}[commandchars=\\\{\}]
\PYG{k}{pair\PYGZus{}style}\PYG{+w}{ }\PYG{n}{multi}\PYG{o}{/}\PYG{n}{lucy}\PYG{o}{/}\PYG{n}{rx}\PYG{+w}{ }\PYG{n}{style}\PYG{+w}{ }\PYG{n}{N}\PYG{+w}{ }\PYG{n}{keyword}\PYG{+w}{ }\PYG{n}{...}
\end{sphinxVerbatim}
\begin{itemize}
\item {} 
\sphinxAtStartPar
style = \sphinxstyleemphasis{lookup} or \sphinxstyleemphasis{linear} = method of interpolation

\item {} 
\sphinxAtStartPar
N = use N values in \sphinxstyleemphasis{lookup}, \sphinxstyleemphasis{linear} tables

\item {} 
\sphinxAtStartPar
weighting = fractional or molecular (optional)

\end{itemize}


\subsection{Examples}
\label{\detokenize{pair_multi_lucy_rx:examples}}
\begin{sphinxVerbatim}[commandchars=\\\{\}]
\PYG{k}{pair\PYGZus{}style}\PYG{+w}{ }\PYG{n}{multi}\PYG{o}{/}\PYG{n}{lucy}\PYG{o}{/}\PYG{n}{rx}\PYG{+w}{ }\PYG{n}{linear}\PYG{+w}{ }\PYG{l+m}{1000}
\PYG{k}{pair\PYGZus{}style}\PYG{+w}{ }\PYG{n}{multi}\PYG{o}{/}\PYG{n}{lucy}\PYG{o}{/}\PYG{n}{rx}\PYG{+w}{ }\PYG{n}{linear}\PYG{+w}{ }\PYG{l+m}{1000}\PYG{+w}{ }\PYG{n}{fractional}
\PYG{k}{pair\PYGZus{}style}\PYG{+w}{ }\PYG{n}{multi}\PYG{o}{/}\PYG{n}{lucy}\PYG{o}{/}\PYG{n}{rx}\PYG{+w}{ }\PYG{n}{linear}\PYG{+w}{ }\PYG{l+m}{1000}\PYG{+w}{ }\PYG{n}{molecular}
\PYG{k}{pair\PYGZus{}coeff}\PYG{+w}{ }\PYG{o}{*}\PYG{+w}{ }\PYG{o}{*}\PYG{+w}{ }\PYG{n}{multibody.table}\PYG{+w}{ }\PYG{n}{ENTRY1}\PYG{+w}{ }\PYG{n}{h2o}\PYG{+w}{ }\PYG{n}{h2o}\PYG{+w}{ }\PYG{l+m}{7.0}
\PYG{k}{pair\PYGZus{}coeff}\PYG{+w}{ }\PYG{o}{*}\PYG{+w}{ }\PYG{o}{*}\PYG{+w}{ }\PYG{n}{multibody.table}\PYG{+w}{ }\PYG{n}{ENTRY1}\PYG{+w}{ }\PYG{n}{h2o}\PYG{+w}{ }\PYG{l+m}{1}\PYG{n}{fluid}\PYG{+w}{ }\PYG{l+m}{7.0}
\end{sphinxVerbatim}


\subsection{Description}
\label{\detokenize{pair_multi_lucy_rx:description}}
\sphinxAtStartPar
Style \sphinxstyleemphasis{multi/lucy/rx} is used in reaction DPD simulations, where the
coarse\sphinxhyphen{}grained (CG) particles are composed of \sphinxstyleemphasis{m} species whose
reaction rate kinetics are determined from a set of \sphinxstyleemphasis{n} reaction rate
equations through the {\hyperref[\detokenize{fix_rx::doc}]{\sphinxcrossref{\DUrole{doc}{fix rx}}}} command.  The species of
one CG particle can interact with a species in a neighboring CG
particle through a site\sphinxhyphen{}site interaction potential model.  Style
\sphinxstyleemphasis{multi/lucy/rx} computes the site\sphinxhyphen{}site density\sphinxhyphen{}dependent force
following from the many\sphinxhyphen{}body form described in {\hyperref[\detokenize{pair_multi_lucy_rx:moore2}]{\sphinxcrossref{\DUrole{std,std-ref}{(Moore)}}}} and
{\hyperref[\detokenize{pair_multi_lucy_rx:warren2}]{\sphinxcrossref{\DUrole{std,std-ref}{(Warren)}}}} as
\begin{equation*}
\begin{split}F_{i}^{DD}(\rho_i,\rho_j,r_{ij}) = \frac{1}{2} \omega_{DD}\left(r_{ij}\right)
\left[A\left(\rho_i\right) + A\left(\rho_j\right)\right]e_{ij}\end{split}
\end{equation*}
\sphinxAtStartPar
which consists of a density\sphinxhyphen{}dependent function, \(A(\rho)\), and a
radial\sphinxhyphen{}dependent weight function, \(\omega_{DD}(r_{ij})\).  The
radial\sphinxhyphen{}dependent weight function, \(\omega_{DD}(r_{ij})\), is taken
as the Lucy function:
\begin{equation*}
\begin{split}\omega_{DD}\left(r_{ij}\right) = \left(1+\frac{3r_{ij}}{r_{cut}}\right)\left(1+\frac{r_{ij}}{r_{cut}}\right)^3\end{split}
\end{equation*}
\sphinxAtStartPar
The density\sphinxhyphen{}dependent energy for a given particle is given by:
\begin{equation*}
\begin{split}u_{i}^{DD}\left(\rho_{i}\right) = \frac{\pi r_{cut}^4}{84} \int_{\rho_0}^{\rho_i} A\left(\rho'\right) d\rho'\end{split}
\end{equation*}
\sphinxAtStartPar
See the supporting information of {\hyperref[\detokenize{pair_multi_lucy_rx:brennan2}]{\sphinxcrossref{\DUrole{std,std-ref}{(Brennan)}}}} or the
publication by {\hyperref[\detokenize{pair_multi_lucy_rx:moore2}]{\sphinxcrossref{\DUrole{std,std-ref}{(Moore)}}}} for more details on the functional
form.

\sphinxAtStartPar
An interpolation table is used to evaluate the density\sphinxhyphen{}dependent energy
(\(\int A(\rho') d \rho'\)) and force (\(A(\rho')\)).  Note that
the prefactor to the energy is computed after the interpolation, thus
the \(\int A(\rho') d\rho'\) will have units of energy / length\textasciicircum{}4.

\sphinxAtStartPar
The interpolation table is created as a pre\sphinxhyphen{}computation by fitting
cubic splines to the file values and interpolating the
density\sphinxhyphen{}dependent energy and force at each of \sphinxstyleemphasis{N} densities.  During a
simulation, the tables are used to interpolate the density\sphinxhyphen{}dependent
energy and force as needed for each pair of particles separated by a
distance \sphinxstyleemphasis{R}.  The interpolation is done in one of 2 styles: \sphinxstyleemphasis{lookup}
and \sphinxstyleemphasis{linear}.

\sphinxAtStartPar
For the \sphinxstyleemphasis{lookup} style, the density is used to find the nearest table
entry, which is the density\sphinxhyphen{}dependent energy and force.

\sphinxAtStartPar
For the \sphinxstyleemphasis{linear} style, the density is used to find the 2 surrounding
table values from which the density\sphinxhyphen{}dependent energy and force are
computed by linear interpolation.

\sphinxAtStartPar
The following coefficients must be defined for each pair of atoms
types via the {\hyperref[\detokenize{pair_coeff::doc}]{\sphinxcrossref{\DUrole{doc}{pair\_coeff}}}} command as in the examples
above.
\begin{itemize}
\item {} 
\sphinxAtStartPar
filename

\item {} 
\sphinxAtStartPar
keyword

\item {} 
\sphinxAtStartPar
species1

\item {} 
\sphinxAtStartPar
species2

\item {} 
\sphinxAtStartPar
cutoff (distance units)

\end{itemize}

\sphinxAtStartPar
The filename specifies a file containing the tabulated
density\sphinxhyphen{}dependent energy and force.  The keyword specifies a section
of the file.  The cutoff is an optional coefficient.  If not
specified, the outer cutoff in the table itself (see below) will be
used to build an interpolation table that extend to the largest
tabulated distance.  If specified, only file values up to the cutoff
are used to create the interpolation table.  The format of this file
is described below.

\sphinxAtStartPar
The species tags define the site\sphinxhyphen{}site interaction potential between
two species contained within two different particles.  The species
tags must either correspond to the species defined in the reaction
kinetics files specified with the {\hyperref[\detokenize{fix_rx::doc}]{\sphinxcrossref{\DUrole{doc}{fix rx}}}} command or they
must correspond to the tag “1fluid”, signifying interaction with a
product species mixture determined through a one\sphinxhyphen{}fluid approximation.
The interaction potential is weighted by the geometric average of
either the mole fraction concentrations or the number of molecules
associated with the interacting coarse\sphinxhyphen{}grained particles (see the
\sphinxstyleemphasis{fractional} or \sphinxstyleemphasis{molecular} weighting pair style options). The coarse\sphinxhyphen{}grained potential is
stored before and after the reaction kinetics solver is applied, where
the difference is defined to be the internal chemical energy (uChem).


\bigskip\hrule\bigskip


\sphinxAtStartPar
The format of a tabulated file is a series of one or more sections,
defined as follows (without the parenthesized comments):
\begin{sphinxalltt}
\# Density\sphinxhyphen{}dependent function (one or more comment or blank lines)

DD\sphinxhyphen{}FUNCTION                (keyword is first text on line)
N 500 R 1.0 10.0           (N, R, RSQ parameters)
                           (blank)
1 1.0 25.5 102.34          (index, density, energy/r\textasciicircum{}4, force)
2 1.02 23.4 98.5
...
500 10.0 0.001 0.003
\end{sphinxalltt}

\sphinxAtStartPar
A section begins with a non\sphinxhyphen{}blank line whose first character is not a
“\#”; blank lines or lines starting with “\#” can be used as comments
between sections.  The first line begins with a keyword which
identifies the section.  The line can contain additional text, but the
initial text must match the argument specified in the pair\_coeff
command.  The next line lists (in any order) one or more parameters
for the table.  Each parameter is a keyword followed by one or more
numeric values.

\sphinxAtStartPar
The parameter “N” is required and its value is the number of table
entries that follow.  Note that this may be different than the \sphinxstyleemphasis{N}
specified in the {\hyperref[\detokenize{pair_multi_lucy_rx::doc}]{\sphinxcrossref{\DUrole{doc}{pair\_style multi/lucy/rx}}}}
command.  Let Ntable = \sphinxstyleemphasis{N} in the pair\_style command, and Nfile = “N”
in the tabulated file.  What LAMMPS does is a preliminary
interpolation by creating splines using the Nfile tabulated values as
nodal points.  It uses these to interpolate the density\sphinxhyphen{}dependent
energy and force at Ntable different points.  The resulting tables of
length Ntable are then used as described above, when computing the
density\sphinxhyphen{}dependent energy and force.  This means that if you want the
interpolation tables of length Ntable to match exactly what is in the
tabulated file (with effectively no preliminary interpolation), you
should set Ntable = Nfile, and use the “RSQ” parameter.  This is
because the internal table abscissa is always RSQ (separation distance
squared), for efficient lookup.

\sphinxAtStartPar
All other parameters are optional.  If “R” or “RSQ” does not appear,
then the distances in each line of the table are used as\sphinxhyphen{}is to perform
spline interpolation.  In this case, the table values can be spaced in
\sphinxstyleemphasis{density} uniformly or however you wish to position table values in
regions of large gradients.

\sphinxAtStartPar
If used, the parameters “R” or “RSQ” are followed by 2 values \sphinxstyleemphasis{rlo}
and \sphinxstyleemphasis{rhi}.  If specified, the density associated with each
density\sphinxhyphen{}dependent energy and force value is computed from these 2
values (at high accuracy), rather than using the (low\sphinxhyphen{}accuracy) value
listed in each line of the table.  The density values in the table
file are ignored in this case.  For “R”, distances uniformly spaced
between \sphinxstyleemphasis{rlo} and \sphinxstyleemphasis{rhi} are computed; for “RSQ”, squared distances
uniformly spaced between \sphinxstyleemphasis{rlo*rlo} and \sphinxstyleemphasis{rhi*rhi} are computed.

\begin{sphinxadmonition}{note}{Note:}
\sphinxAtStartPar
If you use “R” or “RSQ”, the tabulated distance values in the
file are effectively ignored, and replaced by new values as described
in the previous paragraph.  If the density value in the table is not
very close to the new value (i.e. round\sphinxhyphen{}off difference), then you will
be assigning density\sphinxhyphen{}dependent energy and force values to a different
density, which is probably not what you want.  LAMMPS will warn if
this is occurring.
\end{sphinxadmonition}

\sphinxAtStartPar
Following a blank line, the next N lines list the tabulated values.
On each line, the first value is the index from 1 to N, the second value is
r (in density units), the third value is the density\sphinxhyphen{}dependent function
value (in energy units / length\textasciicircum{}4), and the fourth is the force (in force
units).  The density values must increase from one line to the next.

\sphinxAtStartPar
Note that one file can contain many sections, each with a tabulated
potential.  LAMMPS reads the file section by section until it finds
one that matches the specified keyword.


\bigskip\hrule\bigskip



\subsection{Mixing, shift, table, tail correction, restart, rRESPA info}
\label{\detokenize{pair_multi_lucy_rx:mixing-shift-table-tail-correction-restart-rrespa-info}}
\sphinxAtStartPar
This pair style does not support mixing.  Thus, coefficients for all
I,J pairs must be specified explicitly.

\sphinxAtStartPar
The {\hyperref[\detokenize{pair_modify::doc}]{\sphinxcrossref{\DUrole{doc}{pair\_modify}}}} shift, table, and tail options are
not relevant for this pair style.

\sphinxAtStartPar
This pair style writes the settings for the “pair\_style multi/lucy/rx” command
to {\hyperref[\detokenize{restart::doc}]{\sphinxcrossref{\DUrole{doc}{binary restart files}}}}, so a pair\_style command does
not need to specified in an input script that reads a restart file.
However, the coefficient information is not stored in the restart
file, since it is tabulated in the potential files.  Thus, pair\_coeff
commands do need to be specified in the restart input script.

\sphinxAtStartPar
This pair style can only be used via the \sphinxstyleemphasis{pair} keyword of the
{\hyperref[\detokenize{run_style::doc}]{\sphinxcrossref{\DUrole{doc}{run\_style respa}}}} command.  It does not support the
\sphinxstyleemphasis{inner}, \sphinxstyleemphasis{middle}, \sphinxstyleemphasis{outer} keywords.


\bigskip\hrule\bigskip


\sphinxAtStartPar
Styles with a \sphinxstyleemphasis{gpu}, \sphinxstyleemphasis{intel}, \sphinxstyleemphasis{kk}, \sphinxstyleemphasis{omp}, or \sphinxstyleemphasis{opt} suffix are
functionally the same as the corresponding style without the suffix.
They have been optimized to run faster, depending on your available
hardware, as discussed on the {\hyperref[\detokenize{Speed_packages::doc}]{\sphinxcrossref{\DUrole{doc}{Accelerator packages}}}}
page.  The accelerated styles take the same arguments and should
produce the same results, except for round\sphinxhyphen{}off and precision issues.

\sphinxAtStartPar
These accelerated styles are part of the GPU, INTEL, KOKKOS,
OPENMP, and OPT packages, respectively.  They are only enabled if
LAMMPS was built with those packages.  See the {\hyperref[\detokenize{Build_package::doc}]{\sphinxcrossref{\DUrole{doc}{Build package}}}} page for more info.

\sphinxAtStartPar
You can specify the accelerated styles explicitly in your input script
by including their suffix, or you can use the {\hyperref[\detokenize{Run_options::doc}]{\sphinxcrossref{\DUrole{doc}{\sphinxhyphen{}suffix command\sphinxhyphen{}line switch}}}} when you invoke LAMMPS, or you can use the
{\hyperref[\detokenize{suffix::doc}]{\sphinxcrossref{\DUrole{doc}{suffix}}}} command in your input script.

\sphinxAtStartPar
See the {\hyperref[\detokenize{Speed_packages::doc}]{\sphinxcrossref{\DUrole{doc}{Accelerator packages}}}} page for more
instructions on how to use the accelerated styles effectively.


\bigskip\hrule\bigskip



\subsection{Restrictions}
\label{\detokenize{pair_multi_lucy_rx:restrictions}}
\sphinxAtStartPar
This command is part of the DPD\sphinxhyphen{}REACT package.  It is only enabled if
LAMMPS was built with that package.  See the {\hyperref[\detokenize{Build_package::doc}]{\sphinxcrossref{\DUrole{doc}{Build package}}}} page for more info.


\subsection{Related commands}
\label{\detokenize{pair_multi_lucy_rx:related-commands}}
\sphinxAtStartPar
{\hyperref[\detokenize{pair_coeff::doc}]{\sphinxcrossref{\DUrole{doc}{pair\_coeff}}}}


\subsection{Default}
\label{\detokenize{pair_multi_lucy_rx:default}}
\sphinxAtStartPar
fractional weighting


\bigskip\hrule\bigskip

\phantomsection\label{\detokenize{pair_multi_lucy_rx:warren2}}
\sphinxAtStartPar
\sphinxstylestrong{(Warren)} Warren, Phys Rev E, 68, 066702 (2003).

\phantomsection\label{\detokenize{pair_multi_lucy_rx:brennan2}}
\sphinxAtStartPar
\sphinxstylestrong{(Brennan)} Brennan, J Chem Phys Lett, 5, 2144\sphinxhyphen{}2149 (2014).

\phantomsection\label{\detokenize{pair_multi_lucy_rx:moore2}}
\sphinxAtStartPar
\sphinxstylestrong{(Moore)} Moore, J Chem Phys, 144, 104501 (2016).

\sphinxstepscope

\index{pair\_style nb3b/harmonic@\spxentry{pair\_style nb3b/harmonic}}\ignorespaces 
\index{pair\_style nb3b/screened@\spxentry{pair\_style nb3b/screened}}\ignorespaces 

\section{pair\_style nb3b/harmonic command}
\label{\detokenize{pair_nb3b:pair-style-nb3b-harmonic-command}}\label{\detokenize{pair_nb3b:index-1}}\label{\detokenize{pair_nb3b:index-0}}\label{\detokenize{pair_nb3b::doc}}

\section{pair\_style nb3b/screened command}
\label{\detokenize{pair_nb3b:pair-style-nb3b-screened-command}}

\subsection{Syntax}
\label{\detokenize{pair_nb3b:syntax}}
\begin{sphinxVerbatim}[commandchars=\\\{\}]
\PYG{k}{pair\PYGZus{}style}\PYG{+w}{ }\PYG{n}{style}
\end{sphinxVerbatim}
\begin{itemize}
\item {} 
\sphinxAtStartPar
style = \sphinxstyleemphasis{nb3b/harmonic} or \sphinxstyleemphasis{nb3b/screened}

\end{itemize}


\subsection{Examples}
\label{\detokenize{pair_nb3b:examples}}
\begin{sphinxVerbatim}[commandchars=\\\{\}]
\PYG{k}{pair\PYGZus{}style}\PYG{+w}{ }\PYG{n}{nb3b}\PYG{o}{/}\PYG{n}{harmonic}
\PYG{k}{pair\PYGZus{}coeff}\PYG{+w}{ }\PYG{o}{*}\PYG{+w}{ }\PYG{o}{*}\PYG{+w}{ }\PYG{n}{MgOH.nb3bharmonic}\PYG{+w}{ }\PYG{n}{Mg}\PYG{+w}{ }\PYG{n}{O}\PYG{+w}{ }\PYG{n}{H}

\PYG{k}{pair\PYGZus{}style}\PYG{+w}{ }\PYG{n}{nb3b}\PYG{o}{/}\PYG{n}{screened}
\PYG{k}{pair\PYGZus{}coeff}\PYG{+w}{ }\PYG{o}{*}\PYG{+w}{ }\PYG{o}{*}\PYG{+w}{ }\PYG{n}{PO.nb3b.screened}\PYG{+w}{ }\PYG{n}{P}\PYG{+w}{ }\PYG{n}{NULL}\PYG{+w}{ }\PYG{n}{O}
\PYG{k}{pair\PYGZus{}coeff}\PYG{+w}{ }\PYG{o}{*}\PYG{+w}{ }\PYG{o}{*}\PYG{+w}{ }\PYG{n}{SiOH.nb3b.screened}\PYG{+w}{ }\PYG{n}{Si}\PYG{+w}{ }\PYG{n}{O}\PYG{+w}{ }\PYG{n}{H}
\end{sphinxVerbatim}


\subsection{Description}
\label{\detokenize{pair_nb3b:description}}
\sphinxAtStartPar
The pair style \sphinxstyleemphasis{nb3b/harmonic} computes a non\sphinxhyphen{}bonded 3\sphinxhyphen{}body harmonic potential for the
energy E of a system of atoms as
\begin{equation*}
\begin{split}E = K (\theta - \theta_0)^2\end{split}
\end{equation*}
\sphinxAtStartPar
where \(\theta_0\) is the equilibrium value of the angle and \sphinxstyleemphasis{K} is a
prefactor. Note that the usual 1/2 factor is included in \sphinxstyleemphasis{K}. The form
of the potential is identical to that used in angle\_style \sphinxstyleemphasis{harmonic},
but in this case, the atoms do not need to be explicitly bonded.

\sphinxAtStartPar
Style \sphinxstyleemphasis{nb3b/screened} adds an additional exponentially decaying factor to
the harmonic term, given by
\begin{equation*}
\begin{split}E = K (\theta - \theta_0)^2 \exp \left(- \frac{r_{ij}}{\rho_{ij}} - \frac{r_{ik}}{\rho_{ik}} \right)\end{split}
\end{equation*}
\sphinxAtStartPar
where \(\rho_ij\) and \(\rho_ik\) are the screening factors along
the two bonds. Note that the usual 1/2 factor is included in \sphinxstyleemphasis{K}.

\sphinxAtStartPar
Only a single pair\_coeff command is used with these styles which
specifies a potential file with parameters for specified elements.
These are mapped to LAMMPS atom types by specifying N additional
arguments after the filename in the pair\_coeff command, where N is the
number of LAMMPS atom types:
\begin{itemize}
\item {} 
\sphinxAtStartPar
filename

\item {} 
\sphinxAtStartPar
N element names = mapping of elements to atom types

\end{itemize}

\sphinxAtStartPar
See the {\hyperref[\detokenize{pair_coeff::doc}]{\sphinxcrossref{\DUrole{doc}{pair\_coeff}}}} page for alternate ways
to specify the path for the potential file.

\sphinxAtStartPar
As an example, imagine a file SiC.nb3b.harmonic has potential values
for Si and C.  If your LAMMPS simulation has 4 atoms types and you
want the first 3 to be Si, and the fourth to be C, you would use the
following pair\_coeff command:

\begin{sphinxVerbatim}[commandchars=\\\{\}]
\PYG{k}{pair\PYGZus{}coeff}\PYG{+w}{ }\PYG{o}{*}\PYG{+w}{ }\PYG{o}{*}\PYG{+w}{ }\PYG{n}{SiC.nb3b.harmonic}\PYG{+w}{ }\PYG{n}{Si}\PYG{+w}{ }\PYG{n}{Si}\PYG{+w}{ }\PYG{n}{Si}\PYG{+w}{ }\PYG{n}{C}
\end{sphinxVerbatim}

\sphinxAtStartPar
The first 2 arguments must be * * so as to span all LAMMPS atom types.
The first three Si arguments map LAMMPS atom types 1,2,3 to the Si
element in the potential file.  The final C argument maps LAMMPS atom
type 4 to the C element in the potential file.  If a mapping value is
specified as NULL, the mapping is not performed.  This can be used
when the potential is used as part of the \sphinxstyleemphasis{hybrid} pair style.  The
NULL values are placeholders for atom types that will be used with
other potentials. Two examples of pair\_coeff command for use with the
\sphinxstyleemphasis{hybrid} pair style are:

\begin{sphinxVerbatim}[commandchars=\\\{\}]
\PYG{k}{pair\PYGZus{}coeff}\PYG{+w}{ }\PYG{o}{*}\PYG{+w}{ }\PYG{o}{*}\PYG{+w}{ }\PYG{n}{nb3b}\PYG{o}{/}\PYG{n}{harmonic}\PYG{+w}{ }\PYG{n}{MgOH.nb3b.harmonic}\PYG{+w}{ }\PYG{n}{Mg}\PYG{+w}{ }\PYG{n}{O}\PYG{+w}{ }\PYG{n}{H}
\end{sphinxVerbatim}

\sphinxAtStartPar
Three\sphinxhyphen{}body non\sphinxhyphen{}bonded harmonic files in the \sphinxstyleemphasis{potentials} directory of
the LAMMPS distribution have a “.nb3b.harmonic” suffix.  Lines that
are not blank or comments (starting with \#) define parameters for a
triplet of elements.

\sphinxAtStartPar
Each entry has six arguments. The first three are atom types as
referenced in the LAMMPS input file. The first argument specifies the
central atom. The fourth argument indicates the \sphinxstyleemphasis{K} parameter. The
fifth argument indicates \(\theta_0\). The sixth argument indicates a
separation cutoff in Angstroms.

\sphinxAtStartPar
For a given entry, if the second and third arguments are identical,
then the entry is for a cutoff for the distance between types 1 and 2
(values for \sphinxstyleemphasis{K} and \(\theta_0\) are irrelevant in this case).

\sphinxAtStartPar
For a given entry, if the first three arguments are all different, then
the entry is for the \sphinxstyleemphasis{K} and \(\theta_0\) parameters (the cutoff in
this case is irrelevant).

\sphinxAtStartPar
It is required that the potential file contains entries for \sphinxstyleemphasis{all}
permutations of the elements listed in the pair\_coeff command.
If certain combinations are not parameterized the corresponding
parameters should be set to zero. The potential file can also
contain entries for additional elements which are not used in
a particular simulation; LAMMPS ignores those entries.


\bigskip\hrule\bigskip



\subsection{Restrictions}
\label{\detokenize{pair_nb3b:restrictions}}
\sphinxAtStartPar
This pair style can only be used if LAMMPS was built with the MANYBODY
package.  See the {\hyperref[\detokenize{Build_package::doc}]{\sphinxcrossref{\DUrole{doc}{Build package}}}} page for more
info.


\subsection{Related commands}
\label{\detokenize{pair_nb3b:related-commands}}
\sphinxAtStartPar
{\hyperref[\detokenize{pair_coeff::doc}]{\sphinxcrossref{\DUrole{doc}{pair\_coeff}}}}


\subsection{Default}
\label{\detokenize{pair_nb3b:default}}
\sphinxAtStartPar
none

\sphinxstepscope

\index{pair\_style nm/cut@\spxentry{pair\_style nm/cut}}\ignorespaces 
\index{pair\_style nm/cut/split@\spxentry{pair\_style nm/cut/split}}\ignorespaces 
\index{pair\_style nm/cut/coul/cut@\spxentry{pair\_style nm/cut/coul/cut}}\ignorespaces 
\index{pair\_style nm/cut/coul/long@\spxentry{pair\_style nm/cut/coul/long}}\ignorespaces 
\index{pair\_style nm/cut/omp@\spxentry{pair\_style nm/cut/omp}}\ignorespaces 
\index{pair\_style nm/cut/coul/cut/omp@\spxentry{pair\_style nm/cut/coul/cut/omp}}\ignorespaces 
\index{pair\_style nm/cut/coul/long/omp@\spxentry{pair\_style nm/cut/coul/long/omp}}\ignorespaces 

\section{pair\_style nm/cut command}
\label{\detokenize{pair_nm:pair-style-nm-cut-command}}\label{\detokenize{pair_nm:index-6}}\label{\detokenize{pair_nm:index-5}}\label{\detokenize{pair_nm:index-4}}\label{\detokenize{pair_nm:index-3}}\label{\detokenize{pair_nm:index-2}}\label{\detokenize{pair_nm:index-1}}\label{\detokenize{pair_nm:index-0}}\label{\detokenize{pair_nm::doc}}
\sphinxAtStartPar
Accelerator Variants: \sphinxstyleemphasis{nm/cut/omp}


\section{pair\_style nm/cut/split command}
\label{\detokenize{pair_nm:pair-style-nm-cut-split-command}}

\section{pair\_style nm/cut/coul/cut command}
\label{\detokenize{pair_nm:pair-style-nm-cut-coul-cut-command}}
\sphinxAtStartPar
Accelerator Variants: \sphinxstyleemphasis{nm/cut/coul/cut/omp}


\section{pair\_style nm/cut/coul/long command}
\label{\detokenize{pair_nm:pair-style-nm-cut-coul-long-command}}
\sphinxAtStartPar
Accelerator Variants: \sphinxstyleemphasis{nm/cut/coul/long/omp}


\subsection{Syntax}
\label{\detokenize{pair_nm:syntax}}
\begin{sphinxVerbatim}[commandchars=\\\{\}]
\PYG{k}{pair\PYGZus{}style}\PYG{+w}{ }\PYG{n}{style}\PYG{+w}{ }\PYG{n}{args}
\end{sphinxVerbatim}
\begin{itemize}
\item {} 
\sphinxAtStartPar
style = \sphinxstyleemphasis{nm/cut} or \sphinxstyleemphasis{nm/cut/split} or \sphinxstyleemphasis{nm/cut/coul/cut} or \sphinxstyleemphasis{nm/cut/coul/long}

\item {} 
\sphinxAtStartPar
args = list of arguments for a particular style
\begin{sphinxalltt}
\sphinxstyleemphasis{nm/cut} args = cutoff
  cutoff = global cutoff for Pair interactions (distance units)
\sphinxstyleemphasis{nm/cut/split} args = cutoff
  cutoff = global cutoff for Pair interactions (distance units)
\sphinxstyleemphasis{nm/cut/coul/cut} args = cutoff (cutoff2)
  cutoff = global cutoff for Pair (and Coulombic if only 1 arg) (distance units)
  cutoff2 = global cutoff for Coulombic (optional) (distance units)
\sphinxstyleemphasis{nm/cut/coul/long} args = cutoff (cutoff2)
  cutoff = global cutoff for Pair (and Coulombic if only 1 arg) (distance units)
  cutoff2 = global cutoff for Coulombic (optional) (distance units)
\end{sphinxalltt}

\end{itemize}


\subsection{Examples}
\label{\detokenize{pair_nm:examples}}
\begin{sphinxVerbatim}[commandchars=\\\{\}]
\PYG{k}{pair\PYGZus{}style}\PYG{+w}{ }\PYG{n}{nm}\PYG{o}{/}\PYG{n}{cut}\PYG{+w}{ }\PYG{l+m}{12.0}
\PYG{k}{pair\PYGZus{}coeff}\PYG{+w}{ }\PYG{o}{*}\PYG{+w}{ }\PYG{o}{*}\PYG{+w}{ }\PYG{l+m}{0.01}\PYG{+w}{ }\PYG{l+m}{5.4}\PYG{+w}{ }\PYG{l+m}{8.0}\PYG{+w}{ }\PYG{l+m}{7.0}
\PYG{k}{pair\PYGZus{}coeff}\PYG{+w}{ }\PYG{l+m}{1}\PYG{+w}{ }\PYG{l+m}{1}\PYG{+w}{ }\PYG{l+m}{0.01}\PYG{+w}{ }\PYG{l+m}{4.4}\PYG{+w}{ }\PYG{l+m}{7.0}\PYG{+w}{ }\PYG{l+m}{6.0}

\PYG{k}{pair\PYGZus{}style}\PYG{+w}{ }\PYG{n}{nm}\PYG{o}{/}\PYG{n}{cut}\PYG{o}{/}\PYG{n}{split}\PYG{+w}{ }\PYG{l+m}{1.12246}
\PYG{k}{pair\PYGZus{}coeff}\PYG{+w}{ }\PYG{l+m}{1}\PYG{+w}{ }\PYG{l+m}{1}\PYG{+w}{ }\PYG{l+m}{1.0}\PYG{+w}{ }\PYG{l+m}{1.1246}\PYG{+w}{ }\PYG{l+m}{12}\PYG{+w}{ }\PYG{l+m}{6}
\PYG{k}{pair\PYGZus{}coeff}\PYG{+w}{ }\PYG{o}{*}\PYG{+w}{ }\PYG{o}{*}\PYG{+w}{ }\PYG{l+m}{1.0}\PYG{+w}{ }\PYG{l+m}{1.1246}\PYG{+w}{ }\PYG{l+m}{11}\PYG{+w}{ }\PYG{l+m}{6}

\PYG{k}{pair\PYGZus{}style}\PYG{+w}{ }\PYG{n}{nm}\PYG{o}{/}\PYG{n}{cut}\PYG{o}{/}\PYG{n}{coul}\PYG{o}{/}\PYG{n}{cut}\PYG{+w}{ }\PYG{l+m}{12.0}\PYG{+w}{ }\PYG{l+m}{15.0}
\PYG{k}{pair\PYGZus{}coeff}\PYG{+w}{ }\PYG{o}{*}\PYG{+w}{ }\PYG{o}{*}\PYG{+w}{ }\PYG{l+m}{0.01}\PYG{+w}{ }\PYG{l+m}{5.4}\PYG{+w}{ }\PYG{l+m}{8.0}\PYG{+w}{ }\PYG{l+m}{7.0}
\PYG{k}{pair\PYGZus{}coeff}\PYG{+w}{ }\PYG{l+m}{1}\PYG{+w}{ }\PYG{l+m}{1}\PYG{+w}{ }\PYG{l+m}{0.01}\PYG{+w}{ }\PYG{l+m}{4.4}\PYG{+w}{ }\PYG{l+m}{7.0}\PYG{+w}{ }\PYG{l+m}{6.0}

\PYG{k}{pair\PYGZus{}style}\PYG{+w}{ }\PYG{n}{nm}\PYG{o}{/}\PYG{n}{cut}\PYG{o}{/}\PYG{n}{coul}\PYG{o}{/}\PYG{n}{long}\PYG{+w}{ }\PYG{l+m}{12.0}\PYG{+w}{ }\PYG{l+m}{15.0}
\PYG{k}{pair\PYGZus{}coeff}\PYG{+w}{ }\PYG{o}{*}\PYG{+w}{ }\PYG{o}{*}\PYG{+w}{ }\PYG{l+m}{0.01}\PYG{+w}{ }\PYG{l+m}{5.4}\PYG{+w}{ }\PYG{l+m}{8.0}\PYG{+w}{ }\PYG{l+m}{7.0}
\PYG{k}{pair\PYGZus{}coeff}\PYG{+w}{ }\PYG{l+m}{1}\PYG{+w}{ }\PYG{l+m}{1}\PYG{+w}{ }\PYG{l+m}{0.01}\PYG{+w}{ }\PYG{l+m}{4.4}\PYG{+w}{ }\PYG{l+m}{7.0}\PYG{+w}{ }\PYG{l+m}{6.0}
\end{sphinxVerbatim}


\subsection{Description}
\label{\detokenize{pair_nm:description}}
\sphinxAtStartPar
Style \sphinxstyleemphasis{nm} computes site\sphinxhyphen{}site interactions based on the N\sphinxhyphen{}M potential
by {\hyperref[\detokenize{pair_nm:clarke}]{\sphinxcrossref{\DUrole{std,std-ref}{Clarke}}}}, mainly used for ionic liquids.  A site can
represent a single atom or a united\sphinxhyphen{}atom site.  The energy of an
interaction has the following form:
\begin{equation*}
\begin{split}E = \frac{E_0}{(n-m)} \left[ m \left(\frac{r_0}{r}\right)^n - n
\left(\frac{r_0}{r}\right)^m \right] \qquad r < r_c\end{split}
\end{equation*}
\sphinxAtStartPar
where \(r_c\) is the cutoff and \(r_0\) is the minimum of the
potential.  Please note that this differs from the convention used for
other Lennard\sphinxhyphen{}Jones potentials in LAMMPS where \(\sigma\) represents
the location where the energy is zero.

\sphinxAtStartPar
Style \sphinxstyleemphasis{nm/cut/split} applies the standard LJ (12\sphinxhyphen{}6) potential above
\(r_0 = 2^\frac{1}{6}\sigma\).  Style \sphinxstyleemphasis{nm/cut/split} is employed in
polymer equilibration protocols that combine core\sphinxhyphen{}softening approaches
with topology\sphinxhyphen{}changing moves {\hyperref[\detokenize{pair_nm:dietz}]{\sphinxcrossref{\DUrole{std,std-ref}{Dietz}}}}.

\sphinxAtStartPar
Style \sphinxstyleemphasis{nm/cut/coul/cut} adds a Coulombic pairwise interaction given by
\begin{equation*}
\begin{split}E = \frac{C q_i q_j}{\epsilon  r} \qquad r < r_c\end{split}
\end{equation*}
\sphinxAtStartPar
where \(C\) is an energy\sphinxhyphen{}conversion constant, \(q_i\) and \(q_j\)
are the charges on the two atoms, and epsilon is the dielectric constant which can
be set by the {\hyperref[\detokenize{dielectric::doc}]{\sphinxcrossref{\DUrole{doc}{dielectric}}}} command.  If one cutoff is
specified in the pair\_style command, it is used for both the N\sphinxhyphen{}M and Coulombic
terms.  If two cutoffs are specified, they are used as cutoffs for the N\sphinxhyphen{}M and
Coulombic terms respectively.

\sphinxAtStartPar
Styles \sphinxstyleemphasis{nm/cut/coul/long} compute the same
Coulombic interactions as style \sphinxstyleemphasis{nm/cut/coul/cut} except that an
additional damping factor is applied to the Coulombic term so it can
be used in conjunction with the {\hyperref[\detokenize{kspace_style::doc}]{\sphinxcrossref{\DUrole{doc}{kspace\_style}}}}
command and its \sphinxstyleemphasis{ewald} or \sphinxstyleemphasis{pppm} option.  The Coulombic cutoff
specified for this style means that pairwise interactions within this
distance are computed directly; interactions outside that distance are
computed in reciprocal space.

\sphinxAtStartPar
For all of the \sphinxstyleemphasis{nm} pair styles, the following coefficients must
be defined for each pair of atoms types
via the {\hyperref[\detokenize{pair_coeff::doc}]{\sphinxcrossref{\DUrole{doc}{pair\_coeff}}}} command as in the
examples above, or in the data file or restart files read by the
{\hyperref[\detokenize{read_data::doc}]{\sphinxcrossref{\DUrole{doc}{read\_data}}}} or {\hyperref[\detokenize{read_restart::doc}]{\sphinxcrossref{\DUrole{doc}{read\_restart}}}}
commands.
\begin{itemize}
\item {} 
\sphinxAtStartPar
\(E_0\) (energy units)

\item {} 
\sphinxAtStartPar
\(r_0\) (distance units)

\item {} 
\sphinxAtStartPar
\(n\) (unitless)

\item {} 
\sphinxAtStartPar
\(m\) (unitless)

\item {} 
\sphinxAtStartPar
cutoff1 (distance units)

\item {} 
\sphinxAtStartPar
cutoff2 (distance units)

\end{itemize}

\sphinxAtStartPar
The latter 2 coefficients are optional.  If not specified, the global
N\sphinxhyphen{}M and Coulombic cutoffs specified in the pair\_style command are used.
If only one cutoff is specified, it is used as the cutoff for both N\sphinxhyphen{}M
and Coulombic interactions for this type pair.  If both coefficients
are specified, they are used as the N\sphinxhyphen{}M and Coulombic cutoffs for this
type pair.  You cannot specify 2 cutoffs for style \sphinxstyleemphasis{nm}, since it
has no Coulombic terms.

\sphinxAtStartPar
For \sphinxstyleemphasis{nm/cut/coul/long} only the N\sphinxhyphen{}M cutoff can be specified since a
Coulombic cutoff cannot be specified for an individual I,J type pair.
All type pairs use the same global Coulombic cutoff specified in the
pair\_style command.


\bigskip\hrule\bigskip



\subsection{Mixing, shift, table, tail correction, restart, rRESPA info}
\label{\detokenize{pair_nm:mixing-shift-table-tail-correction-restart-rrespa-info}}
\sphinxAtStartPar
These pair styles do not support mixing. Thus, coefficients for all
I,J pairs must be specified explicitly.

\sphinxAtStartPar
All of the \sphinxstyleemphasis{nm} pair styles supports the
{\hyperref[\detokenize{pair_modify::doc}]{\sphinxcrossref{\DUrole{doc}{pair\_modify}}}} shift option for the energy of the pair
interaction.

\sphinxAtStartPar
The \sphinxstyleemphasis{nm/cut/coul/long} pair styles support the
{\hyperref[\detokenize{pair_modify::doc}]{\sphinxcrossref{\DUrole{doc}{pair\_modify}}}} table option since they can tabulate
the short\sphinxhyphen{}range portion of the long\sphinxhyphen{}range Coulombic interaction.

\sphinxAtStartPar
All of the \sphinxstyleemphasis{nm} pair styles support the {\hyperref[\detokenize{pair_modify::doc}]{\sphinxcrossref{\DUrole{doc}{pair\_modify}}}}
tail option for adding a long\sphinxhyphen{}range tail correction to the energy and
pressure for the N\sphinxhyphen{}M portion of the pair interaction.

\sphinxAtStartPar
All of the \sphinxstyleemphasis{nm} pair styles write their information to {\hyperref[\detokenize{restart::doc}]{\sphinxcrossref{\DUrole{doc}{binary restart files}}}}, so pair\_style and pair\_coeff commands do not need
to be specified in an input script that reads a restart file.

\sphinxAtStartPar
All of the \sphinxstyleemphasis{nm} pair styles can only be used via the \sphinxstyleemphasis{pair} keyword of
the {\hyperref[\detokenize{run_style::doc}]{\sphinxcrossref{\DUrole{doc}{run\_style respa}}}} command.  They do not support the
\sphinxstyleemphasis{inner}, \sphinxstyleemphasis{middle}, \sphinxstyleemphasis{outer} keywords.


\bigskip\hrule\bigskip


\sphinxAtStartPar
Styles with a \sphinxstyleemphasis{gpu}, \sphinxstyleemphasis{intel}, \sphinxstyleemphasis{kk}, \sphinxstyleemphasis{omp}, or \sphinxstyleemphasis{opt} suffix are
functionally the same as the corresponding style without the suffix.
They have been optimized to run faster, depending on your available
hardware, as discussed on the {\hyperref[\detokenize{Speed_packages::doc}]{\sphinxcrossref{\DUrole{doc}{Accelerator packages}}}}
page.  The accelerated styles take the same arguments and should
produce the same results, except for round\sphinxhyphen{}off and precision issues.

\sphinxAtStartPar
These accelerated styles are part of the GPU, INTEL, KOKKOS,
OPENMP, and OPT packages, respectively.  They are only enabled if
LAMMPS was built with those packages.  See the {\hyperref[\detokenize{Build_package::doc}]{\sphinxcrossref{\DUrole{doc}{Build package}}}} page for more info.

\sphinxAtStartPar
You can specify the accelerated styles explicitly in your input script
by including their suffix, or you can use the {\hyperref[\detokenize{Run_options::doc}]{\sphinxcrossref{\DUrole{doc}{\sphinxhyphen{}suffix command\sphinxhyphen{}line switch}}}} when you invoke LAMMPS, or you can use the
{\hyperref[\detokenize{suffix::doc}]{\sphinxcrossref{\DUrole{doc}{suffix}}}} command in your input script.

\sphinxAtStartPar
See the {\hyperref[\detokenize{Speed_packages::doc}]{\sphinxcrossref{\DUrole{doc}{Accelerator packages}}}} page for more
instructions on how to use the accelerated styles effectively.


\subsection{Restrictions}
\label{\detokenize{pair_nm:restrictions}}
\sphinxAtStartPar
These pair styles are part of the EXTRA\sphinxhyphen{}PAIR package.  They are only enabled if
LAMMPS was built with that package.  See the
{\hyperref[\detokenize{Build_package::doc}]{\sphinxcrossref{\DUrole{doc}{Build package}}}} page for more info.


\subsection{Related commands}
\label{\detokenize{pair_nm:related-commands}}
\sphinxAtStartPar
{\hyperref[\detokenize{pair_coeff::doc}]{\sphinxcrossref{\DUrole{doc}{pair\_coeff}}}}, {\hyperref[\detokenize{pair_lj::doc}]{\sphinxcrossref{\DUrole{doc}{pair style lj/cut}}}}, {\hyperref[\detokenize{bond_fene::doc}]{\sphinxcrossref{\DUrole{doc}{bond style fene/nm}}}}


\subsection{Default}
\label{\detokenize{pair_nm:default}}
\sphinxAtStartPar
none


\bigskip\hrule\bigskip

\phantomsection\label{\detokenize{pair_nm:clarke}}
\sphinxAtStartPar
\sphinxstylestrong{(Clarke)} Clarke and Smith, J Chem Phys, 84, 2290 (1986).

\phantomsection\label{\detokenize{pair_nm:dietz}}
\sphinxAtStartPar
\sphinxstylestrong{(Dietz)} Dietz and Hoy, J. Chem Phys, 156, 014103 (2022).

\sphinxstepscope

\index{pair\_style none@\spxentry{pair\_style none}}\ignorespaces 

\section{pair\_style none command}
\label{\detokenize{pair_none:pair-style-none-command}}\label{\detokenize{pair_none:index-0}}\label{\detokenize{pair_none::doc}}

\subsection{Syntax}
\label{\detokenize{pair_none:syntax}}
\begin{sphinxVerbatim}[commandchars=\\\{\}]
\PYG{k}{pair\PYGZus{}style}\PYG{+w}{ }\PYG{n}{none}
\end{sphinxVerbatim}


\subsection{Examples}
\label{\detokenize{pair_none:examples}}
\begin{sphinxVerbatim}[commandchars=\\\{\}]
\PYG{k}{pair\PYGZus{}style}\PYG{+w}{ }\PYG{n}{none}
\end{sphinxVerbatim}


\subsection{Description}
\label{\detokenize{pair_none:description}}
\sphinxAtStartPar
Using a pair style of \sphinxstyleemphasis{none} means that any previous pair style setting
will be deleted and pairwise forces and energies are not computed.

\sphinxAtStartPar
As a consequence there will be a pairwise force cutoff of 0.0, which has
implications for the default setting of the neighbor list and the
communication cutoff.  Those are the sum of the largest pairwise cutoff
and the neighbor skin distance (see the documentation of the
{\hyperref[\detokenize{neighbor::doc}]{\sphinxcrossref{\DUrole{doc}{neighbor}}}} command and the {\hyperref[\detokenize{comm_modify::doc}]{\sphinxcrossref{\DUrole{doc}{comm\_modify}}}} command).  When you have bonds, angles, dihedrals, or
impropers defined at the same time, you must set the communication
cutoff so that communication cutoff distance is large enough to acquire
and communicate sufficient ghost atoms from neighboring subdomains as
needed for computing bonds, angles, etc.

\sphinxAtStartPar
A pair style of \sphinxstyleemphasis{none} will also not request a pairwise neighbor list.
However if the {\hyperref[\detokenize{neighbor::doc}]{\sphinxcrossref{\DUrole{doc}{neighbor}}}} style is \sphinxstyleemphasis{bin}, data
structures for binning are still allocated.  If the neighbor list cutoff
is small, then these data structures can consume a large amount of
memory.  So you should either set the neighbor style to \sphinxstyleemphasis{nsq} or set the
skin distance to a larger value.

\sphinxAtStartPar
See the {\hyperref[\detokenize{pair_zero::doc}]{\sphinxcrossref{\DUrole{doc}{pair\_style zero}}}} for a way to set a pairwise
cutoff and thus trigger the building of a neighbor lists and setting
a corresponding communication cutoff, but compute no pairwise interactions.


\subsection{Restrictions}
\label{\detokenize{pair_none:restrictions}}
\sphinxAtStartPar
You must not use a {\hyperref[\detokenize{pair_coeff::doc}]{\sphinxcrossref{\DUrole{doc}{pair\_coeff}}}} command with this pair
style.  Since there is no interaction computed, you cannot set any
coefficients for it.


\subsection{Related commands}
\label{\detokenize{pair_none:related-commands}}
\sphinxAtStartPar
{\hyperref[\detokenize{pair_zero::doc}]{\sphinxcrossref{\DUrole{doc}{pair\_style zero}}}}


\subsection{Default}
\label{\detokenize{pair_none:default}}
\sphinxAtStartPar
none

\sphinxstepscope

\index{pair\_style oxdna/excv@\spxentry{pair\_style oxdna/excv}}\ignorespaces 
\index{pair\_style oxdna/stk@\spxentry{pair\_style oxdna/stk}}\ignorespaces 
\index{pair\_style oxdna/hbond@\spxentry{pair\_style oxdna/hbond}}\ignorespaces 
\index{pair\_style oxdna/xstk@\spxentry{pair\_style oxdna/xstk}}\ignorespaces 
\index{pair\_style oxdna/coaxstk@\spxentry{pair\_style oxdna/coaxstk}}\ignorespaces 

\section{pair\_style oxdna/excv command}
\label{\detokenize{pair_oxdna:pair-style-oxdna-excv-command}}\label{\detokenize{pair_oxdna:index-4}}\label{\detokenize{pair_oxdna:index-3}}\label{\detokenize{pair_oxdna:index-2}}\label{\detokenize{pair_oxdna:index-1}}\label{\detokenize{pair_oxdna:index-0}}\label{\detokenize{pair_oxdna::doc}}

\section{pair\_style oxdna/stk command}
\label{\detokenize{pair_oxdna:pair-style-oxdna-stk-command}}

\section{pair\_style oxdna/hbond command}
\label{\detokenize{pair_oxdna:pair-style-oxdna-hbond-command}}

\section{pair\_style oxdna/xstk command}
\label{\detokenize{pair_oxdna:pair-style-oxdna-xstk-command}}

\section{pair\_style oxdna/coaxstk command}
\label{\detokenize{pair_oxdna:pair-style-oxdna-coaxstk-command}}

\subsection{Syntax}
\label{\detokenize{pair_oxdna:syntax}}
\begin{sphinxVerbatim}[commandchars=\\\{\}]
\PYG{k}{pair\PYGZus{}style}\PYG{+w}{ }\PYG{n}{style1}

\PYG{k}{pair\PYGZus{}coeff}\PYG{+w}{ }\PYG{o}{*}\PYG{+w}{ }\PYG{o}{*}\PYG{+w}{ }\PYG{n}{style2}\PYG{+w}{ }\PYG{n}{args}
\end{sphinxVerbatim}
\begin{itemize}
\item {} 
\sphinxAtStartPar
style1 = \sphinxstyleemphasis{hybrid/overlay oxdna/excv oxdna/stk oxdna/hbond oxdna/xstk oxdna/coaxstk}

\item {} 
\sphinxAtStartPar
style2 = \sphinxstyleemphasis{oxdna/excv} or \sphinxstyleemphasis{oxdna/stk} or \sphinxstyleemphasis{oxdna/hbond} or \sphinxstyleemphasis{oxdna/xstk} or \sphinxstyleemphasis{oxdna/coaxstk}

\item {} 
\sphinxAtStartPar
args = list of arguments for these particular styles

\end{itemize}
\begin{sphinxalltt}
\sphinxstyleemphasis{oxdna/stk} args = seq T xi kappa 6.0 0.4 0.9 0.32 0.75 1.3 0 0.8 0.9 0 0.95 0.9 0 0.95 2.0 0.65 2.0 0.65
  seq = seqav (for average sequence stacking strength) or seqdep (for sequence\sphinxhyphen{}dependent stacking strength)
  T = temperature (LJ units: 0.1 = 300 K, real units: 300 = 300 K)
  xi = 1.3448 (LJ units) or 8.01727944817084 (real units), temperature\sphinxhyphen{}independent coefficient in stacking strength
  kappa = 2.6568 (LJ units) or 0.005279604 (real units), coefficient of linear temperature dependence in stacking strength
\sphinxstyleemphasis{oxdna/hbond} args = seq eps 8.0 0.4 0.75 0.34 0.7 1.5 0 0.7 1.5 0 0.7 1.5 0 0.7 0.46 3.141592653589793 0.7 4.0 1.5707963267948966 0.45 4.0 1.5707963267948966 0.45
  seq = seqav (for average sequence base\sphinxhyphen{}pairing strength) or seqdep (for sequence\sphinxhyphen{}dependent base\sphinxhyphen{}pairing strength)
  eps = 1.077 (LJ units) or 6.42073911784652 (real units), average hydrogen bonding strength between A\sphinxhyphen{}T and C\sphinxhyphen{}G Watson\sphinxhyphen{}Crick base pairs, 0 between all other pairs
\end{sphinxalltt}


\subsection{Examples}
\label{\detokenize{pair_oxdna:examples}}
\begin{sphinxVerbatim}[commandchars=\\\{\}]
\PYG{c}{\PYGZsh{} LJ units}
\PYG{k}{pair\PYGZus{}style}\PYG{+w}{ }\PYG{n}{hybrid}\PYG{o}{/}\PYG{n}{overlay}\PYG{+w}{ }\PYG{n}{oxdna}\PYG{o}{/}\PYG{n}{excv}\PYG{+w}{ }\PYG{n}{oxdna}\PYG{o}{/}\PYG{n}{stk}\PYG{+w}{ }\PYG{n}{oxdna}\PYG{o}{/}\PYG{n}{hbond}\PYG{+w}{ }\PYG{n}{oxdna}\PYG{o}{/}\PYG{n}{xstk}\PYG{+w}{ }\PYG{n}{oxdna}\PYG{o}{/}\PYG{n}{coaxstk}
\PYG{k}{pair\PYGZus{}coeff}\PYG{+w}{ }\PYG{o}{*}\PYG{+w}{ }\PYG{o}{*}\PYG{+w}{ }\PYG{n}{oxdna}\PYG{o}{/}\PYG{n}{excv}\PYG{+w}{    }\PYG{l+m}{2.0}\PYG{+w}{ }\PYG{l+m}{0.7}\PYG{+w}{ }\PYG{l+m}{0.675}\PYG{+w}{ }\PYG{l+m}{2.0}\PYG{+w}{ }\PYG{l+m}{0.515}\PYG{+w}{ }\PYG{l+m}{0.5}\PYG{+w}{ }\PYG{l+m}{2.0}\PYG{+w}{ }\PYG{l+m}{0.33}\PYG{+w}{ }\PYG{l+m}{0.32}
\PYG{k}{pair\PYGZus{}coeff}\PYG{+w}{ }\PYG{o}{*}\PYG{+w}{ }\PYG{o}{*}\PYG{+w}{ }\PYG{n}{oxdna}\PYG{o}{/}\PYG{n}{stk}\PYG{+w}{     }\PYG{n}{seqdep}\PYG{+w}{ }\PYG{l+m}{0.1}\PYG{+w}{ }\PYG{l+m}{1.3448}\PYG{+w}{ }\PYG{l+m}{2.6568}\PYG{+w}{ }\PYG{l+m}{6.0}\PYG{+w}{ }\PYG{l+m}{0.4}\PYG{+w}{ }\PYG{l+m}{0.9}\PYG{+w}{ }\PYG{l+m}{0.32}\PYG{+w}{ }\PYG{l+m}{0.75}\PYG{+w}{ }\PYG{l+m}{1.3}\PYG{+w}{ }\PYG{l+m}{0}\PYG{+w}{ }\PYG{l+m}{0.8}\PYG{+w}{ }\PYG{l+m}{0.9}\PYG{+w}{ }\PYG{l+m}{0}\PYG{+w}{ }\PYG{l+m}{0.95}\PYG{+w}{ }\PYG{l+m}{0.9}\PYG{+w}{ }\PYG{l+m}{0}\PYG{+w}{ }\PYG{l+m}{0.95}\PYG{+w}{ }\PYG{l+m}{2.0}\PYG{+w}{ }\PYG{l+m}{0.65}\PYG{+w}{ }\PYG{l+m}{2.0}\PYG{+w}{ }\PYG{l+m}{0.65}
\PYG{k}{pair\PYGZus{}coeff}\PYG{+w}{ }\PYG{o}{*}\PYG{+w}{ }\PYG{o}{*}\PYG{+w}{ }\PYG{n}{oxdna}\PYG{o}{/}\PYG{n}{hbond}\PYG{+w}{   }\PYG{n}{seqdep}\PYG{+w}{ }\PYG{l+m}{0.0}\PYG{+w}{ }\PYG{l+m}{8.0}\PYG{+w}{ }\PYG{l+m}{0.4}\PYG{+w}{ }\PYG{l+m}{0.75}\PYG{+w}{ }\PYG{l+m}{0.34}\PYG{+w}{ }\PYG{l+m}{0.7}\PYG{+w}{ }\PYG{l+m}{1.5}\PYG{+w}{ }\PYG{l+m}{0}\PYG{+w}{ }\PYG{l+m}{0.7}\PYG{+w}{ }\PYG{l+m}{1.5}\PYG{+w}{ }\PYG{l+m}{0}\PYG{+w}{ }\PYG{l+m}{0.7}\PYG{+w}{ }\PYG{l+m}{1.5}\PYG{+w}{ }\PYG{l+m}{0}\PYG{+w}{ }\PYG{l+m}{0.7}\PYG{+w}{ }\PYG{l+m}{0.46}\PYG{+w}{ }\PYG{l+m}{3.141592653589793}\PYG{+w}{ }\PYG{l+m}{0.7}\PYG{+w}{ }\PYG{l+m}{4.0}\PYG{+w}{ }\PYG{l+m}{1.5707963267948966}\PYG{+w}{ }\PYG{l+m}{0.45}\PYG{+w}{ }\PYG{l+m}{4.0}\PYG{+w}{ }\PYG{l+m}{1.5707963267948966}\PYG{+w}{ }\PYG{l+m}{0.45}
\PYG{k}{pair\PYGZus{}coeff}\PYG{+w}{ }\PYG{l+m}{1}\PYG{+w}{ }\PYG{l+m}{4}\PYG{+w}{ }\PYG{n}{oxdna}\PYG{o}{/}\PYG{n}{hbond}\PYG{+w}{   }\PYG{n}{seqdep}\PYG{+w}{ }\PYG{l+m}{1.077}\PYG{+w}{ }\PYG{l+m}{8.0}\PYG{+w}{ }\PYG{l+m}{0.4}\PYG{+w}{ }\PYG{l+m}{0.75}\PYG{+w}{ }\PYG{l+m}{0.34}\PYG{+w}{ }\PYG{l+m}{0.7}\PYG{+w}{ }\PYG{l+m}{1.5}\PYG{+w}{ }\PYG{l+m}{0}\PYG{+w}{ }\PYG{l+m}{0.7}\PYG{+w}{ }\PYG{l+m}{1.5}\PYG{+w}{ }\PYG{l+m}{0}\PYG{+w}{ }\PYG{l+m}{0.7}\PYG{+w}{ }\PYG{l+m}{1.5}\PYG{+w}{ }\PYG{l+m}{0}\PYG{+w}{ }\PYG{l+m}{0.7}\PYG{+w}{ }\PYG{l+m}{0.46}\PYG{+w}{ }\PYG{l+m}{3.141592653589793}\PYG{+w}{ }\PYG{l+m}{0.7}\PYG{+w}{ }\PYG{l+m}{4.0}\PYG{+w}{ }\PYG{l+m}{1.5707963267948966}\PYG{+w}{ }\PYG{l+m}{0.45}\PYG{+w}{ }\PYG{l+m}{4.0}\PYG{+w}{ }\PYG{l+m}{1.5707963267948966}\PYG{+w}{ }\PYG{l+m}{0.45}
\PYG{k}{pair\PYGZus{}coeff}\PYG{+w}{ }\PYG{l+m}{2}\PYG{+w}{ }\PYG{l+m}{3}\PYG{+w}{ }\PYG{n}{oxdna}\PYG{o}{/}\PYG{n}{hbond}\PYG{+w}{   }\PYG{n}{seqdep}\PYG{+w}{ }\PYG{l+m}{1.077}\PYG{+w}{ }\PYG{l+m}{8.0}\PYG{+w}{ }\PYG{l+m}{0.4}\PYG{+w}{ }\PYG{l+m}{0.75}\PYG{+w}{ }\PYG{l+m}{0.34}\PYG{+w}{ }\PYG{l+m}{0.7}\PYG{+w}{ }\PYG{l+m}{1.5}\PYG{+w}{ }\PYG{l+m}{0}\PYG{+w}{ }\PYG{l+m}{0.7}\PYG{+w}{ }\PYG{l+m}{1.5}\PYG{+w}{ }\PYG{l+m}{0}\PYG{+w}{ }\PYG{l+m}{0.7}\PYG{+w}{ }\PYG{l+m}{1.5}\PYG{+w}{ }\PYG{l+m}{0}\PYG{+w}{ }\PYG{l+m}{0.7}\PYG{+w}{ }\PYG{l+m}{0.46}\PYG{+w}{ }\PYG{l+m}{3.141592653589793}\PYG{+w}{ }\PYG{l+m}{0.7}\PYG{+w}{ }\PYG{l+m}{4.0}\PYG{+w}{ }\PYG{l+m}{1.5707963267948966}\PYG{+w}{ }\PYG{l+m}{0.45}\PYG{+w}{ }\PYG{l+m}{4.0}\PYG{+w}{ }\PYG{l+m}{1.5707963267948966}\PYG{+w}{ }\PYG{l+m}{0.45}
\PYG{k}{pair\PYGZus{}coeff}\PYG{+w}{ }\PYG{o}{*}\PYG{+w}{ }\PYG{o}{*}\PYG{+w}{ }\PYG{n}{oxdna}\PYG{o}{/}\PYG{n}{xstk}\PYG{+w}{    }\PYG{l+m}{47.5}\PYG{+w}{ }\PYG{l+m}{0.575}\PYG{+w}{ }\PYG{l+m}{0.675}\PYG{+w}{ }\PYG{l+m}{0.495}\PYG{+w}{ }\PYG{l+m}{0.655}\PYG{+w}{ }\PYG{l+m}{2.25}\PYG{+w}{ }\PYG{l+m}{0.791592653589793}\PYG{+w}{ }\PYG{l+m}{0.58}\PYG{+w}{ }\PYG{l+m}{1.7}\PYG{+w}{ }\PYG{l+m}{1.0}\PYG{+w}{ }\PYG{l+m}{0.68}\PYG{+w}{ }\PYG{l+m}{1.7}\PYG{+w}{ }\PYG{l+m}{1.0}\PYG{+w}{ }\PYG{l+m}{0.68}\PYG{+w}{ }\PYG{l+m}{1.5}\PYG{+w}{ }\PYG{l+m}{0}\PYG{+w}{ }\PYG{l+m}{0.65}\PYG{+w}{ }\PYG{l+m}{1.7}\PYG{+w}{ }\PYG{l+m}{0.875}\PYG{+w}{ }\PYG{l+m}{0.68}\PYG{+w}{ }\PYG{l+m}{1.7}\PYG{+w}{ }\PYG{l+m}{0.875}\PYG{+w}{ }\PYG{l+m}{0.68}
\PYG{k}{pair\PYGZus{}coeff}\PYG{+w}{ }\PYG{o}{*}\PYG{+w}{ }\PYG{o}{*}\PYG{+w}{ }\PYG{n}{oxdna}\PYG{o}{/}\PYG{n}{coaxstk}\PYG{+w}{ }\PYG{l+m}{46.0}\PYG{+w}{ }\PYG{l+m}{0.4}\PYG{+w}{ }\PYG{l+m}{0.6}\PYG{+w}{ }\PYG{l+m}{0.22}\PYG{+w}{ }\PYG{l+m}{0.58}\PYG{+w}{ }\PYG{l+m}{2.0}\PYG{+w}{ }\PYG{l+m}{2.541592653589793}\PYG{+w}{ }\PYG{l+m}{0.65}\PYG{+w}{ }\PYG{l+m}{1.3}\PYG{+w}{ }\PYG{l+m}{0}\PYG{+w}{ }\PYG{l+m}{0.8}\PYG{+w}{ }\PYG{l+m}{0.9}\PYG{+w}{ }\PYG{l+m}{0}\PYG{+w}{ }\PYG{l+m}{0.95}\PYG{+w}{ }\PYG{l+m}{0.9}\PYG{+w}{ }\PYG{l+m}{0}\PYG{+w}{ }\PYG{l+m}{0.95}\PYG{+w}{ }\PYG{l+m}{2.0}\PYG{+w}{ }\PYG{o}{\PYGZhy{}}\PYG{l+m}{0.65}\PYG{+w}{ }\PYG{l+m}{2.0}\PYG{+w}{ }\PYG{o}{\PYGZhy{}}\PYG{l+m}{0.65}

\PYG{k}{pair\PYGZus{}style}\PYG{+w}{ }\PYG{n}{hybrid}\PYG{o}{/}\PYG{n}{overlay}\PYG{+w}{ }\PYG{n}{oxdna}\PYG{o}{/}\PYG{n}{excv}\PYG{+w}{ }\PYG{n}{oxdna}\PYG{o}{/}\PYG{n}{stk}\PYG{+w}{ }\PYG{n}{oxdna}\PYG{o}{/}\PYG{n}{hbond}\PYG{+w}{ }\PYG{n}{oxdna}\PYG{o}{/}\PYG{n}{xstk}\PYG{+w}{ }\PYG{n}{oxdna}\PYG{o}{/}\PYG{n}{coaxstk}
\PYG{k}{pair\PYGZus{}coeff}\PYG{+w}{ }\PYG{o}{*}\PYG{+w}{ }\PYG{o}{*}\PYG{+w}{ }\PYG{n}{oxdna}\PYG{o}{/}\PYG{n}{excv}\PYG{+w}{    }\PYG{n}{oxdna\PYGZus{}lj.cgdna}
\PYG{k}{pair\PYGZus{}coeff}\PYG{+w}{ }\PYG{o}{*}\PYG{+w}{ }\PYG{o}{*}\PYG{+w}{ }\PYG{n}{oxdna}\PYG{o}{/}\PYG{n}{stk}\PYG{+w}{     }\PYG{n}{seqav}\PYG{+w}{ }\PYG{l+m}{0.1}\PYG{+w}{ }\PYG{l+m}{1.3448}\PYG{+w}{ }\PYG{l+m}{2.6568}\PYG{+w}{ }\PYG{n}{oxdna\PYGZus{}lj.cgdna}
\PYG{k}{pair\PYGZus{}coeff}\PYG{+w}{ }\PYG{o}{*}\PYG{+w}{ }\PYG{o}{*}\PYG{+w}{ }\PYG{n}{oxdna}\PYG{o}{/}\PYG{n}{hbond}\PYG{+w}{   }\PYG{n}{seqav}\PYG{+w}{ }\PYG{n}{oxdna\PYGZus{}lj.cgdna}
\PYG{k}{pair\PYGZus{}coeff}\PYG{+w}{ }\PYG{l+m}{1}\PYG{+w}{ }\PYG{l+m}{4}\PYG{+w}{ }\PYG{n}{oxdna}\PYG{o}{/}\PYG{n}{hbond}\PYG{+w}{   }\PYG{n}{seqav}\PYG{+w}{ }\PYG{n}{oxdna\PYGZus{}lj.cgdna}
\PYG{k}{pair\PYGZus{}coeff}\PYG{+w}{ }\PYG{l+m}{2}\PYG{+w}{ }\PYG{l+m}{3}\PYG{+w}{ }\PYG{n}{oxdna}\PYG{o}{/}\PYG{n}{hbond}\PYG{+w}{   }\PYG{n}{seqav}\PYG{+w}{ }\PYG{n}{oxdna\PYGZus{}lj.cgdna}
\PYG{k}{pair\PYGZus{}coeff}\PYG{+w}{ }\PYG{o}{*}\PYG{+w}{ }\PYG{o}{*}\PYG{+w}{ }\PYG{n}{oxdna}\PYG{o}{/}\PYG{n}{xstk}\PYG{+w}{    }\PYG{n}{oxdna\PYGZus{}lj.cgdna}
\PYG{k}{pair\PYGZus{}coeff}\PYG{+w}{ }\PYG{o}{*}\PYG{+w}{ }\PYG{o}{*}\PYG{+w}{ }\PYG{n}{oxdna}\PYG{o}{/}\PYG{n}{coaxstk}\PYG{+w}{ }\PYG{n}{oxdna\PYGZus{}lj.cgdna}

\PYG{c}{\PYGZsh{} Real units}
\PYG{k}{pair\PYGZus{}style}\PYG{+w}{ }\PYG{n}{hybrid}\PYG{o}{/}\PYG{n}{overlay}\PYG{+w}{ }\PYG{n}{oxdna}\PYG{o}{/}\PYG{n}{excv}\PYG{+w}{ }\PYG{n}{oxdna}\PYG{o}{/}\PYG{n}{stk}\PYG{+w}{ }\PYG{n}{oxdna}\PYG{o}{/}\PYG{n}{hbond}\PYG{+w}{ }\PYG{n}{oxdna}\PYG{o}{/}\PYG{n}{xstk}\PYG{+w}{ }\PYG{n}{oxdna}\PYG{o}{/}\PYG{n}{coaxstk}
\PYG{k}{pair\PYGZus{}coeff}\PYG{+w}{ }\PYG{o}{*}\PYG{+w}{ }\PYG{o}{*}\PYG{+w}{ }\PYG{n}{oxdna}\PYG{o}{/}\PYG{n}{excv}\PYG{+w}{    }\PYG{l+m}{11.92337812042065}\PYG{+w}{ }\PYG{l+m}{5.9626}\PYG{+w}{ }\PYG{l+m}{5.74965}\PYG{+w}{ }\PYG{l+m}{11.92337812042065}\PYG{+w}{ }\PYG{l+m}{4.38677}\PYG{+w}{ }\PYG{l+m}{4.259}\PYG{+w}{ }\PYG{l+m}{11.92337812042065}\PYG{+w}{ }\PYG{l+m}{2.81094}\PYG{+w}{ }\PYG{l+m}{2.72576}
\PYG{k}{pair\PYGZus{}coeff}\PYG{+w}{ }\PYG{o}{*}\PYG{+w}{ }\PYG{o}{*}\PYG{+w}{ }\PYG{n}{oxdna}\PYG{o}{/}\PYG{n}{stk}\PYG{+w}{     }\PYG{n}{seqdep}\PYG{+w}{ }\PYG{l+m}{300.0}\PYG{+w}{ }\PYG{l+m}{8.01727944817084}\PYG{+w}{ }\PYG{l+m}{0.005279604}\PYG{+w}{ }\PYG{l+m}{0.70439070204273}\PYG{+w}{ }\PYG{l+m}{3.4072}\PYG{+w}{ }\PYG{l+m}{7.6662}\PYG{+w}{ }\PYG{l+m}{2.72576}\PYG{+w}{ }\PYG{l+m}{6.3885}\PYG{+w}{ }\PYG{l+m}{1.3}\PYG{+w}{ }\PYG{l+m}{0.0}\PYG{+w}{ }\PYG{l+m}{0.8}\PYG{+w}{ }\PYG{l+m}{0.9}\PYG{+w}{ }\PYG{l+m}{0.0}\PYG{+w}{ }\PYG{l+m}{0.95}\PYG{+w}{ }\PYG{l+m}{0.9}\PYG{+w}{ }\PYG{l+m}{0.0}\PYG{+w}{ }\PYG{l+m}{0.95}\PYG{+w}{ }\PYG{l+m}{2.0}\PYG{+w}{ }\PYG{l+m}{0.65}\PYG{+w}{ }\PYG{l+m}{2.0}\PYG{+w}{ }\PYG{l+m}{0.65}
\PYG{k}{pair\PYGZus{}coeff}\PYG{+w}{ }\PYG{o}{*}\PYG{+w}{ }\PYG{o}{*}\PYG{+w}{ }\PYG{n}{oxdna}\PYG{o}{/}\PYG{n}{hbond}\PYG{+w}{   }\PYG{n}{seqdep}\PYG{+w}{ }\PYG{l+m}{0.0}\PYG{+w}{ }\PYG{l+m}{0.93918760272364}\PYG{+w}{ }\PYG{l+m}{3.4072}\PYG{+w}{ }\PYG{l+m}{6.3885}\PYG{+w}{ }\PYG{l+m}{2.89612}\PYG{+w}{ }\PYG{l+m}{5.9626}\PYG{+w}{ }\PYG{l+m}{1.5}\PYG{+w}{ }\PYG{l+m}{0.0}\PYG{+w}{ }\PYG{l+m}{0.7}\PYG{+w}{ }\PYG{l+m}{1.5}\PYG{+w}{ }\PYG{l+m}{0.0}\PYG{+w}{ }\PYG{l+m}{0.7}\PYG{+w}{ }\PYG{l+m}{1.5}\PYG{+w}{ }\PYG{l+m}{0.0}\PYG{+w}{ }\PYG{l+m}{0.7}\PYG{+w}{ }\PYG{l+m}{0.46}\PYG{+w}{ }\PYG{l+m}{3.141592654}\PYG{+w}{ }\PYG{l+m}{0.7}\PYG{+w}{ }\PYG{l+m}{4.0}\PYG{+w}{ }\PYG{l+m}{1.570796327}\PYG{+w}{ }\PYG{l+m}{0.45}\PYG{+w}{ }\PYG{l+m}{4.0}\PYG{+w}{ }\PYG{l+m}{1.570796327}\PYG{+w}{ }\PYG{l+m}{0.45}
\PYG{k}{pair\PYGZus{}coeff}\PYG{+w}{ }\PYG{l+m}{1}\PYG{+w}{ }\PYG{l+m}{4}\PYG{+w}{ }\PYG{n}{oxdna}\PYG{o}{/}\PYG{n}{hbond}\PYG{+w}{   }\PYG{n}{seqdep}\PYG{+w}{ }\PYG{l+m}{6.42073911784652}\PYG{+w}{ }\PYG{l+m}{0.93918760272364}\PYG{+w}{ }\PYG{l+m}{3.4072}\PYG{+w}{ }\PYG{l+m}{6.3885}\PYG{+w}{ }\PYG{l+m}{2.89612}\PYG{+w}{ }\PYG{l+m}{5.9626}\PYG{+w}{ }\PYG{l+m}{1.5}\PYG{+w}{ }\PYG{l+m}{0.0}\PYG{+w}{ }\PYG{l+m}{0.7}\PYG{+w}{ }\PYG{l+m}{1.5}\PYG{+w}{ }\PYG{l+m}{0.0}\PYG{+w}{ }\PYG{l+m}{0.7}\PYG{+w}{ }\PYG{l+m}{1.5}\PYG{+w}{ }\PYG{l+m}{0.0}\PYG{+w}{ }\PYG{l+m}{0.7}\PYG{+w}{ }\PYG{l+m}{0.46}\PYG{+w}{ }\PYG{l+m}{3.141592654}\PYG{+w}{ }\PYG{l+m}{0.7}\PYG{+w}{ }\PYG{l+m}{4.0}\PYG{+w}{ }\PYG{l+m}{1.570796327}\PYG{+w}{ }\PYG{l+m}{0.45}\PYG{+w}{ }\PYG{l+m}{4.0}\PYG{+w}{ }\PYG{l+m}{1.570796327}\PYG{+w}{ }\PYG{l+m}{0.45}
\PYG{k}{pair\PYGZus{}coeff}\PYG{+w}{ }\PYG{l+m}{2}\PYG{+w}{ }\PYG{l+m}{3}\PYG{+w}{ }\PYG{n}{oxdna}\PYG{o}{/}\PYG{n}{hbond}\PYG{+w}{   }\PYG{n}{seqdep}\PYG{+w}{ }\PYG{l+m}{6.42073911784652}\PYG{+w}{ }\PYG{l+m}{0.93918760272364}\PYG{+w}{ }\PYG{l+m}{3.4072}\PYG{+w}{ }\PYG{l+m}{6.3885}\PYG{+w}{ }\PYG{l+m}{2.89612}\PYG{+w}{ }\PYG{l+m}{5.9626}\PYG{+w}{ }\PYG{l+m}{1.5}\PYG{+w}{ }\PYG{l+m}{0.0}\PYG{+w}{ }\PYG{l+m}{0.7}\PYG{+w}{ }\PYG{l+m}{1.5}\PYG{+w}{ }\PYG{l+m}{0.0}\PYG{+w}{ }\PYG{l+m}{0.7}\PYG{+w}{ }\PYG{l+m}{1.5}\PYG{+w}{ }\PYG{l+m}{0.0}\PYG{+w}{ }\PYG{l+m}{0.7}\PYG{+w}{ }\PYG{l+m}{0.46}\PYG{+w}{ }\PYG{l+m}{3.141592654}\PYG{+w}{ }\PYG{l+m}{0.7}\PYG{+w}{ }\PYG{l+m}{4.0}\PYG{+w}{ }\PYG{l+m}{1.570796327}\PYG{+w}{ }\PYG{l+m}{0.45}\PYG{+w}{ }\PYG{l+m}{4.0}\PYG{+w}{ }\PYG{l+m}{1.570796327}\PYG{+w}{ }\PYG{l+m}{0.45}
\PYG{k}{pair\PYGZus{}coeff}\PYG{+w}{ }\PYG{o}{*}\PYG{+w}{ }\PYG{o}{*}\PYG{+w}{ }\PYG{n}{oxdna}\PYG{o}{/}\PYG{n}{xstk}\PYG{+w}{    }\PYG{l+m}{3.9029021145006}\PYG{+w}{ }\PYG{l+m}{4.89785}\PYG{+w}{ }\PYG{l+m}{5.74965}\PYG{+w}{ }\PYG{l+m}{4.21641}\PYG{+w}{ }\PYG{l+m}{5.57929}\PYG{+w}{ }\PYG{l+m}{2.25}\PYG{+w}{ }\PYG{l+m}{0.791592654}\PYG{+w}{ }\PYG{l+m}{0.58}\PYG{+w}{ }\PYG{l+m}{1.7}\PYG{+w}{ }\PYG{l+m}{1.0}\PYG{+w}{ }\PYG{l+m}{0.68}\PYG{+w}{ }\PYG{l+m}{1.7}\PYG{+w}{ }\PYG{l+m}{1.0}\PYG{+w}{ }\PYG{l+m}{0.68}\PYG{+w}{ }\PYG{l+m}{1.5}\PYG{+w}{ }\PYG{l+m}{0.0}\PYG{+w}{ }\PYG{l+m}{0.65}\PYG{+w}{ }\PYG{l+m}{1.7}\PYG{+w}{ }\PYG{l+m}{0.875}\PYG{+w}{ }\PYG{l+m}{0.68}\PYG{+w}{ }\PYG{l+m}{1.7}\PYG{+w}{ }\PYG{l+m}{0.875}\PYG{+w}{ }\PYG{l+m}{0.68}
\PYG{k}{pair\PYGZus{}coeff}\PYG{+w}{ }\PYG{o}{*}\PYG{+w}{ }\PYG{o}{*}\PYG{+w}{ }\PYG{n}{oxdna}\PYG{o}{/}\PYG{n}{coaxstk}\PYG{+w}{ }\PYG{l+m}{3.77965257404268}\PYG{+w}{ }\PYG{l+m}{3.4072}\PYG{+w}{ }\PYG{l+m}{5.1108}\PYG{+w}{ }\PYG{l+m}{1.87396}\PYG{+w}{ }\PYG{l+m}{4.94044}\PYG{+w}{ }\PYG{l+m}{2.0}\PYG{+w}{ }\PYG{l+m}{2.541592654}\PYG{+w}{ }\PYG{l+m}{0.65}\PYG{+w}{ }\PYG{l+m}{1.3}\PYG{+w}{ }\PYG{l+m}{0.0}\PYG{+w}{ }\PYG{l+m}{0.8}\PYG{+w}{ }\PYG{l+m}{0.9}\PYG{+w}{ }\PYG{l+m}{0.0}\PYG{+w}{ }\PYG{l+m}{0.95}\PYG{+w}{ }\PYG{l+m}{0.9}\PYG{+w}{ }\PYG{l+m}{0.0}\PYG{+w}{ }\PYG{l+m}{0.95}\PYG{+w}{ }\PYG{l+m}{2.0}\PYG{+w}{ }\PYG{o}{\PYGZhy{}}\PYG{l+m}{0.65}\PYG{+w}{ }\PYG{l+m}{2.0}\PYG{+w}{ }\PYG{o}{\PYGZhy{}}\PYG{l+m}{0.65}

\PYG{k}{pair\PYGZus{}style}\PYG{+w}{ }\PYG{n}{hybrid}\PYG{o}{/}\PYG{n}{overlay}\PYG{+w}{ }\PYG{n}{oxdna}\PYG{o}{/}\PYG{n}{excv}\PYG{+w}{ }\PYG{n}{oxdna}\PYG{o}{/}\PYG{n}{stk}\PYG{+w}{ }\PYG{n}{oxdna}\PYG{o}{/}\PYG{n}{hbond}\PYG{+w}{ }\PYG{n}{oxdna}\PYG{o}{/}\PYG{n}{xstk}\PYG{+w}{ }\PYG{n}{oxdna}\PYG{o}{/}\PYG{n}{coaxstk}
\PYG{k}{pair\PYGZus{}coeff}\PYG{+w}{ }\PYG{o}{*}\PYG{+w}{ }\PYG{o}{*}\PYG{+w}{ }\PYG{n}{oxdna}\PYG{o}{/}\PYG{n}{excv}\PYG{+w}{    }\PYG{n}{oxdna\PYGZus{}real.cgdna}
\PYG{k}{pair\PYGZus{}coeff}\PYG{+w}{ }\PYG{o}{*}\PYG{+w}{ }\PYG{o}{*}\PYG{+w}{ }\PYG{n}{oxdna}\PYG{o}{/}\PYG{n}{stk}\PYG{+w}{     }\PYG{n}{seqav}\PYG{+w}{ }\PYG{l+m}{300.0}\PYG{+w}{ }\PYG{l+m}{8.01727944817084}\PYG{+w}{ }\PYG{l+m}{0.005279604}\PYG{+w}{ }\PYG{n}{oxdna\PYGZus{}real.cgdna}
\PYG{k}{pair\PYGZus{}coeff}\PYG{+w}{ }\PYG{o}{*}\PYG{+w}{ }\PYG{o}{*}\PYG{+w}{ }\PYG{n}{oxdna}\PYG{o}{/}\PYG{n}{hbond}\PYG{+w}{   }\PYG{n}{seqav}\PYG{+w}{ }\PYG{n}{oxdna\PYGZus{}real.cgdna}
\PYG{k}{pair\PYGZus{}coeff}\PYG{+w}{ }\PYG{l+m}{1}\PYG{+w}{ }\PYG{l+m}{4}\PYG{+w}{ }\PYG{n}{oxdna}\PYG{o}{/}\PYG{n}{hbond}\PYG{+w}{   }\PYG{n}{seqav}\PYG{+w}{ }\PYG{n}{oxdna\PYGZus{}real.cgdna}
\PYG{k}{pair\PYGZus{}coeff}\PYG{+w}{ }\PYG{l+m}{2}\PYG{+w}{ }\PYG{l+m}{3}\PYG{+w}{ }\PYG{n}{oxdna}\PYG{o}{/}\PYG{n}{hbond}\PYG{+w}{   }\PYG{n}{seqav}\PYG{+w}{ }\PYG{n}{oxdna\PYGZus{}real.cgdna}
\PYG{k}{pair\PYGZus{}coeff}\PYG{+w}{ }\PYG{o}{*}\PYG{+w}{ }\PYG{o}{*}\PYG{+w}{ }\PYG{n}{oxdna}\PYG{o}{/}\PYG{n}{xstk}\PYG{+w}{    }\PYG{n}{oxdna\PYGZus{}real.cgdna}
\PYG{k}{pair\PYGZus{}coeff}\PYG{+w}{ }\PYG{o}{*}\PYG{+w}{ }\PYG{o}{*}\PYG{+w}{ }\PYG{n}{oxdna}\PYG{o}{/}\PYG{n}{coaxstk}\PYG{+w}{ }\PYG{n}{oxdna\PYGZus{}real.cgdna}
\end{sphinxVerbatim}

\begin{sphinxadmonition}{note}{Note:}
\sphinxAtStartPar
The coefficients in the above examples are provided in forms
compatible with both \sphinxstyleemphasis{units lj} and \sphinxstyleemphasis{units real} (see documentation
of {\hyperref[\detokenize{units::doc}]{\sphinxcrossref{\DUrole{doc}{units}}}}).  These can also be read from a potential
file with correct unit style by specifying the name of the
file. Several potential files for each unit style are included in the
\sphinxcode{\sphinxupquote{potentials}} directory of the LAMMPS distribution.
\end{sphinxadmonition}


\subsection{Description}
\label{\detokenize{pair_oxdna:description}}
\sphinxAtStartPar
The \sphinxstyleemphasis{oxdna} pair styles compute the pairwise\sphinxhyphen{}additive parts of the oxDNA
force field for coarse\sphinxhyphen{}grained modelling of DNA. The effective
interaction between the nucleotides consists of potentials for the
excluded volume interaction \sphinxstyleemphasis{oxdna/excv}, the stacking \sphinxstyleemphasis{oxdna/stk},
cross\sphinxhyphen{}stacking \sphinxstyleemphasis{oxdna/xstk} and coaxial stacking interaction
\sphinxstyleemphasis{oxdna/coaxstk} as well as the hydrogen\sphinxhyphen{}bonding interaction
\sphinxstyleemphasis{oxdna/hbond} between complementary pairs of nucleotides on opposite
strands. Average sequence or sequence\sphinxhyphen{}dependent stacking and
base\sphinxhyphen{}pairing strengths are supported {\hyperref[\detokenize{pair_oxdna:sulc1}]{\sphinxcrossref{\DUrole{std,std-ref}{(Sulc)}}}}. Quasi\sphinxhyphen{}unique
base\sphinxhyphen{}pairing between nucleotides can be achieved by using more
complementary pairs of atom types like 5\sphinxhyphen{}8 and 6\sphinxhyphen{}7, 9\sphinxhyphen{}12 and 10\sphinxhyphen{}11,
13\sphinxhyphen{}16 and 14\sphinxhyphen{}15, etc.  This prevents the hybridization of in principle
complementary bases within Ntypes/4 bases up and down along the
backbone.

\sphinxAtStartPar
The exact functional form of the pair styles is rather complex.  The
individual potentials consist of products of modulation factors, which
themselves are constructed from a number of more basic potentials
(Morse, Lennard\sphinxhyphen{}Jones, harmonic angle and distance) as well as quadratic
smoothing and modulation terms.  We refer to {\hyperref[\detokenize{pair_oxdna:ouldridge-dphil1}]{\sphinxcrossref{\DUrole{std,std-ref}{(Ouldridge\sphinxhyphen{}DPhil)}}}} and {\hyperref[\detokenize{pair_oxdna:ouldridge1}]{\sphinxcrossref{\DUrole{std,std-ref}{(Ouldridge)}}}} for a detailed
description of the oxDNA force field.

\begin{sphinxadmonition}{note}{Note:}
\sphinxAtStartPar
These pair styles have to be used together with the related oxDNA
bond style \sphinxstyleemphasis{oxdna/fene} for the connectivity of the phosphate
backbone (see also documentation of {\hyperref[\detokenize{bond_oxdna::doc}]{\sphinxcrossref{\DUrole{doc}{bond\_style oxdna/fene}}}}). Most of the coefficients in the above example have to
be kept fixed and cannot be changed without reparameterizing the
entire model.  Exceptions are the first four coefficients after
\sphinxstyleemphasis{oxdna/stk} (seq=seqdep, T=0.1, xi=1.3448 and kappa=2.6568 and
corresponding \sphinxstyleemphasis{real unit} equivalents in the above examples) and the
first coefficient after \sphinxstyleemphasis{oxdna/hbond} (seq=seqdep in the above
example).  When using a Langevin thermostat, e.g. through {\hyperref[\detokenize{fix_langevin::doc}]{\sphinxcrossref{\DUrole{doc}{fix
langevin}}}} or {\hyperref[\detokenize{fix_nve_dotc_langevin::doc}]{\sphinxcrossref{\DUrole{doc}{fix nve/dotc/langevin}}}} the temperature coefficients have to be
matched to the one used in the fix.
\end{sphinxadmonition}

\begin{sphinxadmonition}{note}{Note:}
\sphinxAtStartPar
These pair styles have to be used with the \sphinxstyleemphasis{atom\_style hybrid bond
ellipsoid oxdna} (see documentation of {\hyperref[\detokenize{atom_style::doc}]{\sphinxcrossref{\DUrole{doc}{atom\_style}}}}). The \sphinxstyleemphasis{atom\_style oxdna} stores the 3’\sphinxhyphen{}to\sphinxhyphen{}5’ polarity
of the nucleotide strand, which is set through the bond topology in
the data file. The first (second) atom in a bond definition is
understood to point towards the 3’\sphinxhyphen{}end (5’\sphinxhyphen{}end) of the strand.
\end{sphinxadmonition}

\sphinxAtStartPar
Example input and data files for DNA duplexes can be found in
\sphinxcode{\sphinxupquote{examples/PACKAGES/cgdna/examples/oxDNA/}} and \sphinxcode{\sphinxupquote{.../oxDNA2/}}.  A
simple python setup tool which creates single straight or helical DNA
strands, DNA duplexes or arrays of DNA duplexes can be found in
\sphinxcode{\sphinxupquote{examples/PACKAGES/cgdna/util/}}.

\sphinxAtStartPar
Please cite {\hyperref[\detokenize{pair_oxdna:henrich1}]{\sphinxcrossref{\DUrole{std,std-ref}{(Henrich)}}}} in any publication that uses
this implementation. An updated documentation that contains general
information on the model, its implementation and performance as well as
the structure of the data and input file can be found \sphinxhref{PDF/CG-DNA.pdf}{here}.

\sphinxAtStartPar
Please cite also the relevant oxDNA publications
{\hyperref[\detokenize{pair_oxdna:ouldridge1}]{\sphinxcrossref{\DUrole{std,std-ref}{(Ouldridge)}}}},
{\hyperref[\detokenize{pair_oxdna:ouldridge-dphil1}]{\sphinxcrossref{\DUrole{std,std-ref}{(Ouldridge\sphinxhyphen{}DPhil)}}}}
and {\hyperref[\detokenize{pair_oxdna:sulc1}]{\sphinxcrossref{\DUrole{std,std-ref}{(Sulc)}}}}.


\bigskip\hrule\bigskip



\subsection{Potential file reading}
\label{\detokenize{pair_oxdna:potential-file-reading}}
\sphinxAtStartPar
For each pair style above the first non\sphinxhyphen{}modifiable argument can be a
filename, and if it is, no further arguments should be
supplied. Therefore the following command:

\begin{sphinxVerbatim}[commandchars=\\\{\}]
\PYG{k}{pair\PYGZus{}coeff}\PYG{+w}{ }\PYG{l+m}{1}\PYG{+w}{ }\PYG{l+m}{4}\PYG{+w}{ }\PYG{n}{oxdna}\PYG{o}{/}\PYG{n}{hbond}\PYG{+w}{   }\PYG{n}{seqav}\PYG{+w}{ }\PYG{n}{oxdna\PYGZus{}lj.cgdna}
\end{sphinxVerbatim}

\sphinxAtStartPar
will be interpreted as a request to read the corresponding hydrogen
bonding potential parameters from the file with the given name. The file
can define multiple potential parameters for both bonded and pair
interactions, but for the example pair interaction above there must
exist in the file a line of the form:

\begin{sphinxVerbatim}[commandchars=\\\{\}]
\PYG{l+m}{1}\PYG{+w}{ }\PYG{l+m}{4}\PYG{+w}{ }\PYG{n}{hbond}\PYG{+w}{     }\PYG{o}{\PYGZlt{}}\PYG{n}{coefficients}\PYG{o}{\PYGZgt{}}
\end{sphinxVerbatim}

\sphinxAtStartPar
If potential customization is required, the potential file reading can
be mixed with the manual specification of the potential parameters. For
example, the following command:

\begin{sphinxVerbatim}[commandchars=\\\{\}]
\PYG{k}{pair\PYGZus{}style}\PYG{+w}{ }\PYG{n}{hybrid}\PYG{o}{/}\PYG{n}{overlay}\PYG{+w}{ }\PYG{n}{oxdna}\PYG{o}{/}\PYG{n}{excv}\PYG{+w}{ }\PYG{n}{oxdna}\PYG{o}{/}\PYG{n}{stk}\PYG{+w}{ }\PYG{n}{oxdna}\PYG{o}{/}\PYG{n}{hbond}\PYG{+w}{ }\PYG{n}{oxdna}\PYG{o}{/}\PYG{n}{xstk}\PYG{+w}{ }\PYG{n}{oxdna}\PYG{o}{/}\PYG{n}{coaxstk}
\PYG{k}{pair\PYGZus{}coeff}\PYG{+w}{ }\PYG{o}{*}\PYG{+w}{ }\PYG{o}{*}\PYG{+w}{ }\PYG{n}{oxdna}\PYG{o}{/}\PYG{n}{excv}\PYG{+w}{    }\PYG{n}{oxdna\PYGZus{}lj.cgdna}
\PYG{k}{pair\PYGZus{}coeff}\PYG{+w}{ }\PYG{o}{*}\PYG{+w}{ }\PYG{o}{*}\PYG{+w}{ }\PYG{n}{oxdna}\PYG{o}{/}\PYG{n}{stk}\PYG{+w}{     }\PYG{n}{seqav}\PYG{+w}{ }\PYG{l+m}{0.1}\PYG{+w}{ }\PYG{l+m}{1.3448}\PYG{+w}{ }\PYG{l+m}{2.6568}\PYG{+w}{ }\PYG{l+m}{6.0}\PYG{+w}{ }\PYG{l+m}{0.4}\PYG{+w}{ }\PYG{l+m}{0.9}\PYG{+w}{ }\PYG{l+m}{0.32}\PYG{+w}{ }\PYG{l+m}{0.75}\PYG{+w}{ }\PYG{l+m}{1.3}\PYG{+w}{ }\PYG{l+m}{0}\PYG{+w}{ }\PYG{l+m}{0.8}\PYG{+w}{ }\PYG{l+m}{0.9}\PYG{+w}{ }\PYG{l+m}{0}\PYG{+w}{ }\PYG{l+m}{0.95}\PYG{+w}{ }\PYG{l+m}{0.9}\PYG{+w}{ }\PYG{l+m}{0}\PYG{+w}{ }\PYG{l+m}{0.95}\PYG{+w}{ }\PYG{l+m}{2.0}\PYG{+w}{ }\PYG{l+m}{0.65}\PYG{+w}{ }\PYG{l+m}{2.0}\PYG{+w}{ }\PYG{l+m}{0.65}
\PYG{k}{pair\PYGZus{}coeff}\PYG{+w}{ }\PYG{o}{*}\PYG{+w}{ }\PYG{o}{*}\PYG{+w}{ }\PYG{n}{oxdna}\PYG{o}{/}\PYG{n}{hbond}\PYG{+w}{   }\PYG{n}{seqav}\PYG{+w}{ }\PYG{n}{oxdna\PYGZus{}lj.cgdna}
\PYG{k}{pair\PYGZus{}coeff}\PYG{+w}{ }\PYG{l+m}{1}\PYG{+w}{ }\PYG{l+m}{4}\PYG{+w}{ }\PYG{n}{oxdna}\PYG{o}{/}\PYG{n}{hbond}\PYG{+w}{   }\PYG{n}{seqav}\PYG{+w}{ }\PYG{n}{oxdna\PYGZus{}lj.cgdna}
\PYG{k}{pair\PYGZus{}coeff}\PYG{+w}{ }\PYG{l+m}{2}\PYG{+w}{ }\PYG{l+m}{3}\PYG{+w}{ }\PYG{n}{oxdna}\PYG{o}{/}\PYG{n}{hbond}\PYG{+w}{   }\PYG{n}{seqav}\PYG{+w}{ }\PYG{n}{oxdna\PYGZus{}lj.cgdna}
\PYG{k}{pair\PYGZus{}coeff}\PYG{+w}{ }\PYG{o}{*}\PYG{+w}{ }\PYG{o}{*}\PYG{+w}{ }\PYG{n}{oxdna}\PYG{o}{/}\PYG{n}{xstk}\PYG{+w}{    }\PYG{n}{oxdna\PYGZus{}lj.cgdna}
\PYG{k}{pair\PYGZus{}coeff}\PYG{+w}{ }\PYG{o}{*}\PYG{+w}{ }\PYG{o}{*}\PYG{+w}{ }\PYG{n}{oxdna}\PYG{o}{/}\PYG{n}{coaxstk}\PYG{+w}{ }\PYG{l+m}{46.0}\PYG{+w}{ }\PYG{l+m}{0.4}\PYG{+w}{ }\PYG{l+m}{0.6}\PYG{+w}{ }\PYG{l+m}{0.22}\PYG{+w}{ }\PYG{l+m}{0.58}\PYG{+w}{ }\PYG{l+m}{2.0}\PYG{+w}{ }\PYG{l+m}{2.541592653589793}\PYG{+w}{ }\PYG{l+m}{0.65}\PYG{+w}{ }\PYG{l+m}{1.3}\PYG{+w}{ }\PYG{l+m}{0}\PYG{+w}{ }\PYG{l+m}{0.8}\PYG{+w}{ }\PYG{l+m}{0.9}\PYG{+w}{ }\PYG{l+m}{0}\PYG{+w}{ }\PYG{l+m}{0.95}\PYG{+w}{ }\PYG{l+m}{0.9}\PYG{+w}{ }\PYG{l+m}{0}\PYG{+w}{ }\PYG{l+m}{0.95}\PYG{+w}{ }\PYG{l+m}{2.0}\PYG{+w}{ }\PYG{o}{\PYGZhy{}}\PYG{l+m}{0.65}\PYG{+w}{ }\PYG{l+m}{2.0}\PYG{+w}{ }\PYG{o}{\PYGZhy{}}\PYG{l+m}{0.65}
\end{sphinxVerbatim}

\sphinxAtStartPar
will read the stacking and coaxial stacking potential parameters from
the manual specification and all others from the potential file
\sphinxstyleemphasis{oxdna\_lj.cgdna}.

\sphinxAtStartPar
There are sample potential files for each unit style in the
\sphinxcode{\sphinxupquote{potentials}} directory of the LAMMPS distribution. The potential file
unit system must align with the units defined via the {\hyperref[\detokenize{units::doc}]{\sphinxcrossref{\DUrole{doc}{units}}}} command. For conversion between different \sphinxstyleemphasis{LJ} and \sphinxstyleemphasis{real} unit
systems for oxDNA, the python tool \sphinxstyleemphasis{lj2real.py} located in the
\sphinxcode{\sphinxupquote{examples/PACKAGES/cgdna/util/}} directory can be used. This tool
assumes similar file structure to the examples found in
\sphinxcode{\sphinxupquote{examples/PACKAGES/cgdna/examples/}}.


\bigskip\hrule\bigskip



\subsection{Restrictions}
\label{\detokenize{pair_oxdna:restrictions}}
\sphinxAtStartPar
These pair styles can only be used if LAMMPS was built with the
CG\sphinxhyphen{}DNA package and the MOLECULE and ASPHERE package.  See the
{\hyperref[\detokenize{Build_package::doc}]{\sphinxcrossref{\DUrole{doc}{Build package}}}} page for more info.


\subsection{Related commands}
\label{\detokenize{pair_oxdna:related-commands}}
\sphinxAtStartPar
{\hyperref[\detokenize{bond_oxdna::doc}]{\sphinxcrossref{\DUrole{doc}{bond\_style oxdna/fene}}}}, {\hyperref[\detokenize{pair_coeff::doc}]{\sphinxcrossref{\DUrole{doc}{pair\_coeff}}}},
{\hyperref[\detokenize{bond_oxdna::doc}]{\sphinxcrossref{\DUrole{doc}{bond\_style oxdna2/fene}}}}, {\hyperref[\detokenize{pair_oxdna2::doc}]{\sphinxcrossref{\DUrole{doc}{pair\_style oxdna2/excv}}}},
{\hyperref[\detokenize{bond_oxdna::doc}]{\sphinxcrossref{\DUrole{doc}{bond\_style oxrna2/fene}}}}, {\hyperref[\detokenize{pair_oxrna2::doc}]{\sphinxcrossref{\DUrole{doc}{pair\_style oxrna2/excv}}}},
{\hyperref[\detokenize{atom_style::doc}]{\sphinxcrossref{\DUrole{doc}{atom\_style oxdna}}}}, {\hyperref[\detokenize{fix_nve_dotc_langevin::doc}]{\sphinxcrossref{\DUrole{doc}{fix nve/dotc/langevin}}}}


\subsection{Default}
\label{\detokenize{pair_oxdna:default}}
\sphinxAtStartPar
none


\bigskip\hrule\bigskip

\phantomsection\label{\detokenize{pair_oxdna:henrich1}}
\sphinxAtStartPar
\sphinxstylestrong{(Henrich)} O. Henrich, Y. A. Gutierrez\sphinxhyphen{}Fosado, T. Curk, T. E. Ouldridge, Eur. Phys. J. E 41, 57 (2018).

\phantomsection\label{\detokenize{pair_oxdna:ouldridge-dphil1}}
\sphinxAtStartPar
\sphinxstylestrong{(Ouldridge\sphinxhyphen{}DPhil)} T.E. Ouldridge, Coarse\sphinxhyphen{}grained modelling of DNA and DNA self\sphinxhyphen{}assembly, DPhil. University of Oxford (2011).

\phantomsection\label{\detokenize{pair_oxdna:ouldridge1}}
\sphinxAtStartPar
\sphinxstylestrong{(Ouldridge)} T.E. Ouldridge, A.A. Louis, J.P.K. Doye, J. Chem. Phys. 134, 085101 (2011).

\phantomsection\label{\detokenize{pair_oxdna:sulc1}}
\sphinxAtStartPar
\sphinxstylestrong{(Sulc)} P. Sulc, F. Romano, T.E. Ouldridge, L. Rovigatti, J.P.K. Doye, A.A. Louis, J. Chem. Phys. 137, 135101 (2012).

\sphinxstepscope

\index{pair\_style oxdna2/excv@\spxentry{pair\_style oxdna2/excv}}\ignorespaces 
\index{pair\_style oxdna2/stk@\spxentry{pair\_style oxdna2/stk}}\ignorespaces 
\index{pair\_style oxdna2/hbond@\spxentry{pair\_style oxdna2/hbond}}\ignorespaces 
\index{pair\_style oxdna2/xstk@\spxentry{pair\_style oxdna2/xstk}}\ignorespaces 
\index{pair\_style oxdna2/coaxstk@\spxentry{pair\_style oxdna2/coaxstk}}\ignorespaces 
\index{pair\_style oxdna2/dh@\spxentry{pair\_style oxdna2/dh}}\ignorespaces 

\section{pair\_style oxdna2/excv command}
\label{\detokenize{pair_oxdna2:pair-style-oxdna2-excv-command}}\label{\detokenize{pair_oxdna2:index-5}}\label{\detokenize{pair_oxdna2:index-4}}\label{\detokenize{pair_oxdna2:index-3}}\label{\detokenize{pair_oxdna2:index-2}}\label{\detokenize{pair_oxdna2:index-1}}\label{\detokenize{pair_oxdna2:index-0}}\label{\detokenize{pair_oxdna2::doc}}

\section{pair\_style oxdna2/stk command}
\label{\detokenize{pair_oxdna2:pair-style-oxdna2-stk-command}}

\section{pair\_style oxdna2/hbond command}
\label{\detokenize{pair_oxdna2:pair-style-oxdna2-hbond-command}}

\section{pair\_style oxdna2/xstk command}
\label{\detokenize{pair_oxdna2:pair-style-oxdna2-xstk-command}}

\section{pair\_style oxdna2/coaxstk command}
\label{\detokenize{pair_oxdna2:pair-style-oxdna2-coaxstk-command}}

\section{pair\_style oxdna2/dh command}
\label{\detokenize{pair_oxdna2:pair-style-oxdna2-dh-command}}

\subsection{Syntax}
\label{\detokenize{pair_oxdna2:syntax}}
\begin{sphinxVerbatim}[commandchars=\\\{\}]
\PYG{k}{pair\PYGZus{}style}\PYG{+w}{ }\PYG{n}{style1}

\PYG{k}{pair\PYGZus{}coeff}\PYG{+w}{ }\PYG{o}{*}\PYG{+w}{ }\PYG{o}{*}\PYG{+w}{ }\PYG{n}{style2}\PYG{+w}{ }\PYG{n}{args}
\end{sphinxVerbatim}
\begin{itemize}
\item {} 
\sphinxAtStartPar
style1 = \sphinxstyleemphasis{hybrid/overlay oxdna2/excv oxdna2/stk oxdna2/hbond oxdna2/xstk oxdna2/coaxstk oxdna2/dh}

\item {} 
\sphinxAtStartPar
style2 = \sphinxstyleemphasis{oxdna2/excv} or \sphinxstyleemphasis{oxdna2/stk} or \sphinxstyleemphasis{oxdna2/hbond} or \sphinxstyleemphasis{oxdna2/xstk} or \sphinxstyleemphasis{oxdna2/coaxstk} or \sphinxstyleemphasis{oxdna2/dh}

\item {} 
\sphinxAtStartPar
args = list of arguments for these particular styles

\end{itemize}
\begin{sphinxalltt}
\sphinxstyleemphasis{oxdna2/stk} args = seq T xi kappa 6.0 0.4 0.9 0.32 0.75 1.3 0 0.8 0.9 0 0.95 0.9 0 0.95 2.0 0.65 2.0 0.65
  seq = seqav (for average sequence stacking strength) or seqdep (for sequence\sphinxhyphen{}dependent stacking strength)
  T = temperature (LJ units: 0.1 = 300 K, real units: 300 = 300 K)
  xi = 1.3523 (LJ units) or 8.06199211612242 (real units), temperature\sphinxhyphen{}independent coefficient in stacking strength
  kappa = 2.6717 (LJ units) or 0.005309213 (real units), coefficient of linear temperature dependence in stacking strength
\sphinxstyleemphasis{oxdna2/hbond} args = seq eps 8.0 0.4 0.75 0.34 0.7 1.5 0 0.7 1.5 0 0.7 1.5 0 0.7 0.46 3.141592653589793 0.7 4.0 1.5707963267948966 0.45 4.0 1.5707963267948966 0.45
  seq = seqav (for average sequence base\sphinxhyphen{}pairing strength) or seqdep (for sequence\sphinxhyphen{}dependent base\sphinxhyphen{}pairing strength)
  eps = 1.0678 (LJ units) or 6.36589157849259 (real units), average hydrogen bonding strength between A\sphinxhyphen{}T and C\sphinxhyphen{}G Watson\sphinxhyphen{}Crick base pairs, 0 between all other pairs
\sphinxstyleemphasis{oxdna2/dh} args = T rhos qeff
  T = temperature (LJ units: 0.1 = 300 K, real units: 300 = 300 K)
  rhos = salt concentration (mole per litre)
  qeff = 0.815 (effective charge in elementary charges)
\end{sphinxalltt}


\subsection{Examples}
\label{\detokenize{pair_oxdna2:examples}}
\begin{sphinxVerbatim}[commandchars=\\\{\}]
\PYG{c}{\PYGZsh{} LJ units}
\PYG{k}{pair\PYGZus{}style}\PYG{+w}{ }\PYG{n}{hybrid}\PYG{o}{/}\PYG{n}{overlay}\PYG{+w}{ }\PYG{n}{oxdna2}\PYG{o}{/}\PYG{n}{excv}\PYG{+w}{ }\PYG{n}{oxdna2}\PYG{o}{/}\PYG{n}{stk}\PYG{+w}{ }\PYG{n}{oxdna2}\PYG{o}{/}\PYG{n}{hbond}\PYG{+w}{ }\PYG{n}{oxdna2}\PYG{o}{/}\PYG{n}{xstk}\PYG{+w}{ }\PYG{n}{oxdna2}\PYG{o}{/}\PYG{n}{coaxstk}\PYG{+w}{ }\PYG{n}{oxdna2}\PYG{o}{/}\PYG{n}{dh}
\PYG{k}{pair\PYGZus{}coeff}\PYG{+w}{ }\PYG{o}{*}\PYG{+w}{ }\PYG{o}{*}\PYG{+w}{ }\PYG{n}{oxdna2}\PYG{o}{/}\PYG{n}{excv}\PYG{+w}{    }\PYG{l+m}{2.0}\PYG{+w}{ }\PYG{l+m}{0.7}\PYG{+w}{ }\PYG{l+m}{0.675}\PYG{+w}{ }\PYG{l+m}{2.0}\PYG{+w}{ }\PYG{l+m}{0.515}\PYG{+w}{ }\PYG{l+m}{0.5}\PYG{+w}{ }\PYG{l+m}{2.0}\PYG{+w}{ }\PYG{l+m}{0.33}\PYG{+w}{ }\PYG{l+m}{0.32}
\PYG{k}{pair\PYGZus{}coeff}\PYG{+w}{ }\PYG{o}{*}\PYG{+w}{ }\PYG{o}{*}\PYG{+w}{ }\PYG{n}{oxdna2}\PYG{o}{/}\PYG{n}{stk}\PYG{+w}{     }\PYG{n}{seqdep}\PYG{+w}{ }\PYG{l+m}{0.1}\PYG{+w}{ }\PYG{l+m}{1.3523}\PYG{+w}{ }\PYG{l+m}{2.6717}\PYG{+w}{ }\PYG{l+m}{6.0}\PYG{+w}{ }\PYG{l+m}{0.4}\PYG{+w}{ }\PYG{l+m}{0.9}\PYG{+w}{ }\PYG{l+m}{0.32}\PYG{+w}{ }\PYG{l+m}{0.75}\PYG{+w}{ }\PYG{l+m}{1.3}\PYG{+w}{ }\PYG{l+m}{0}\PYG{+w}{ }\PYG{l+m}{0.8}\PYG{+w}{ }\PYG{l+m}{0.9}\PYG{+w}{ }\PYG{l+m}{0}\PYG{+w}{ }\PYG{l+m}{0.95}\PYG{+w}{ }\PYG{l+m}{0.9}\PYG{+w}{ }\PYG{l+m}{0}\PYG{+w}{ }\PYG{l+m}{0.95}\PYG{+w}{ }\PYG{l+m}{2.0}\PYG{+w}{ }\PYG{l+m}{0.65}\PYG{+w}{ }\PYG{l+m}{2.0}\PYG{+w}{ }\PYG{l+m}{0.65}
\PYG{k}{pair\PYGZus{}coeff}\PYG{+w}{ }\PYG{o}{*}\PYG{+w}{ }\PYG{o}{*}\PYG{+w}{ }\PYG{n}{oxdna2}\PYG{o}{/}\PYG{n}{hbond}\PYG{+w}{   }\PYG{n}{seqdep}\PYG{+w}{ }\PYG{l+m}{0.0}\PYG{+w}{ }\PYG{l+m}{8.0}\PYG{+w}{ }\PYG{l+m}{0.4}\PYG{+w}{ }\PYG{l+m}{0.75}\PYG{+w}{ }\PYG{l+m}{0.34}\PYG{+w}{ }\PYG{l+m}{0.7}\PYG{+w}{ }\PYG{l+m}{1.5}\PYG{+w}{ }\PYG{l+m}{0}\PYG{+w}{ }\PYG{l+m}{0.7}\PYG{+w}{ }\PYG{l+m}{1.5}\PYG{+w}{ }\PYG{l+m}{0}\PYG{+w}{ }\PYG{l+m}{0.7}\PYG{+w}{ }\PYG{l+m}{1.5}\PYG{+w}{ }\PYG{l+m}{0}\PYG{+w}{ }\PYG{l+m}{0.7}\PYG{+w}{ }\PYG{l+m}{0.46}\PYG{+w}{ }\PYG{l+m}{3.141592653589793}\PYG{+w}{ }\PYG{l+m}{0.7}\PYG{+w}{ }\PYG{l+m}{4.0}\PYG{+w}{ }\PYG{l+m}{1.5707963267948966}\PYG{+w}{ }\PYG{l+m}{0.45}\PYG{+w}{ }\PYG{l+m}{4.0}\PYG{+w}{ }\PYG{l+m}{1.5707963267948966}\PYG{+w}{ }\PYG{l+m}{0.45}
\PYG{k}{pair\PYGZus{}coeff}\PYG{+w}{ }\PYG{l+m}{1}\PYG{+w}{ }\PYG{l+m}{4}\PYG{+w}{ }\PYG{n}{oxdna2}\PYG{o}{/}\PYG{n}{hbond}\PYG{+w}{   }\PYG{n}{seqdep}\PYG{+w}{ }\PYG{l+m}{1.0678}\PYG{+w}{ }\PYG{l+m}{8.0}\PYG{+w}{ }\PYG{l+m}{0.4}\PYG{+w}{ }\PYG{l+m}{0.75}\PYG{+w}{ }\PYG{l+m}{0.34}\PYG{+w}{ }\PYG{l+m}{0.7}\PYG{+w}{ }\PYG{l+m}{1.5}\PYG{+w}{ }\PYG{l+m}{0}\PYG{+w}{ }\PYG{l+m}{0.7}\PYG{+w}{ }\PYG{l+m}{1.5}\PYG{+w}{ }\PYG{l+m}{0}\PYG{+w}{ }\PYG{l+m}{0.7}\PYG{+w}{ }\PYG{l+m}{1.5}\PYG{+w}{ }\PYG{l+m}{0}\PYG{+w}{ }\PYG{l+m}{0.7}\PYG{+w}{ }\PYG{l+m}{0.46}\PYG{+w}{ }\PYG{l+m}{3.141592653589793}\PYG{+w}{ }\PYG{l+m}{0.7}\PYG{+w}{ }\PYG{l+m}{4.0}\PYG{+w}{ }\PYG{l+m}{1.5707963267948966}\PYG{+w}{ }\PYG{l+m}{0.45}\PYG{+w}{ }\PYG{l+m}{4.0}\PYG{+w}{ }\PYG{l+m}{1.5707963267948966}\PYG{+w}{ }\PYG{l+m}{0.45}
\PYG{k}{pair\PYGZus{}coeff}\PYG{+w}{ }\PYG{l+m}{2}\PYG{+w}{ }\PYG{l+m}{3}\PYG{+w}{ }\PYG{n}{oxdna2}\PYG{o}{/}\PYG{n}{hbond}\PYG{+w}{   }\PYG{n}{seqdep}\PYG{+w}{ }\PYG{l+m}{1.0678}\PYG{+w}{ }\PYG{l+m}{8.0}\PYG{+w}{ }\PYG{l+m}{0.4}\PYG{+w}{ }\PYG{l+m}{0.75}\PYG{+w}{ }\PYG{l+m}{0.34}\PYG{+w}{ }\PYG{l+m}{0.7}\PYG{+w}{ }\PYG{l+m}{1.5}\PYG{+w}{ }\PYG{l+m}{0}\PYG{+w}{ }\PYG{l+m}{0.7}\PYG{+w}{ }\PYG{l+m}{1.5}\PYG{+w}{ }\PYG{l+m}{0}\PYG{+w}{ }\PYG{l+m}{0.7}\PYG{+w}{ }\PYG{l+m}{1.5}\PYG{+w}{ }\PYG{l+m}{0}\PYG{+w}{ }\PYG{l+m}{0.7}\PYG{+w}{ }\PYG{l+m}{0.46}\PYG{+w}{ }\PYG{l+m}{3.141592653589793}\PYG{+w}{ }\PYG{l+m}{0.7}\PYG{+w}{ }\PYG{l+m}{4.0}\PYG{+w}{ }\PYG{l+m}{1.5707963267948966}\PYG{+w}{ }\PYG{l+m}{0.45}\PYG{+w}{ }\PYG{l+m}{4.0}\PYG{+w}{ }\PYG{l+m}{1.5707963267948966}\PYG{+w}{ }\PYG{l+m}{0.45}
\PYG{k}{pair\PYGZus{}coeff}\PYG{+w}{ }\PYG{o}{*}\PYG{+w}{ }\PYG{o}{*}\PYG{+w}{ }\PYG{n}{oxdna2}\PYG{o}{/}\PYG{n}{xstk}\PYG{+w}{    }\PYG{l+m}{47.5}\PYG{+w}{ }\PYG{l+m}{0.575}\PYG{+w}{ }\PYG{l+m}{0.675}\PYG{+w}{ }\PYG{l+m}{0.495}\PYG{+w}{ }\PYG{l+m}{0.655}\PYG{+w}{ }\PYG{l+m}{2.25}\PYG{+w}{ }\PYG{l+m}{0.791592653589793}\PYG{+w}{ }\PYG{l+m}{0.58}\PYG{+w}{ }\PYG{l+m}{1.7}\PYG{+w}{ }\PYG{l+m}{1.0}\PYG{+w}{ }\PYG{l+m}{0.68}\PYG{+w}{ }\PYG{l+m}{1.7}\PYG{+w}{ }\PYG{l+m}{1.0}\PYG{+w}{ }\PYG{l+m}{0.68}\PYG{+w}{ }\PYG{l+m}{1.5}\PYG{+w}{ }\PYG{l+m}{0}\PYG{+w}{ }\PYG{l+m}{0.65}\PYG{+w}{ }\PYG{l+m}{1.7}\PYG{+w}{ }\PYG{l+m}{0.875}\PYG{+w}{ }\PYG{l+m}{0.68}\PYG{+w}{ }\PYG{l+m}{1.7}\PYG{+w}{ }\PYG{l+m}{0.875}\PYG{+w}{ }\PYG{l+m}{0.68}
\PYG{k}{pair\PYGZus{}coeff}\PYG{+w}{ }\PYG{o}{*}\PYG{+w}{ }\PYG{o}{*}\PYG{+w}{ }\PYG{n}{oxdna2}\PYG{o}{/}\PYG{n}{coaxstk}\PYG{+w}{ }\PYG{l+m}{58.5}\PYG{+w}{ }\PYG{l+m}{0.4}\PYG{+w}{ }\PYG{l+m}{0.6}\PYG{+w}{ }\PYG{l+m}{0.22}\PYG{+w}{ }\PYG{l+m}{0.58}\PYG{+w}{ }\PYG{l+m}{2.0}\PYG{+w}{ }\PYG{l+m}{2.891592653589793}\PYG{+w}{ }\PYG{l+m}{0.65}\PYG{+w}{ }\PYG{l+m}{1.3}\PYG{+w}{ }\PYG{l+m}{0}\PYG{+w}{ }\PYG{l+m}{0.8}\PYG{+w}{ }\PYG{l+m}{0.9}\PYG{+w}{ }\PYG{l+m}{0}\PYG{+w}{ }\PYG{l+m}{0.95}\PYG{+w}{ }\PYG{l+m}{0.9}\PYG{+w}{ }\PYG{l+m}{0}\PYG{+w}{ }\PYG{l+m}{0.95}\PYG{+w}{ }\PYG{l+m}{40.0}\PYG{+w}{ }\PYG{l+m}{3.116592653589793}
\PYG{k}{pair\PYGZus{}coeff}\PYG{+w}{ }\PYG{o}{*}\PYG{+w}{ }\PYG{o}{*}\PYG{+w}{ }\PYG{n}{oxdna2}\PYG{o}{/}\PYG{n}{dh}\PYG{+w}{      }\PYG{l+m}{0.1}\PYG{+w}{ }\PYG{l+m}{0.5}\PYG{+w}{ }\PYG{l+m}{0.815}

\PYG{k}{pair\PYGZus{}style}\PYG{+w}{ }\PYG{n}{hybrid}\PYG{o}{/}\PYG{n}{overlay}\PYG{+w}{ }\PYG{n}{oxdna2}\PYG{o}{/}\PYG{n}{excv}\PYG{+w}{ }\PYG{n}{oxdna2}\PYG{o}{/}\PYG{n}{stk}\PYG{+w}{ }\PYG{n}{oxdna2}\PYG{o}{/}\PYG{n}{hbond}\PYG{+w}{ }\PYG{n}{oxdna2}\PYG{o}{/}\PYG{n}{xstk}\PYG{+w}{ }\PYG{n}{oxdna2}\PYG{o}{/}\PYG{n}{coaxstk}\PYG{+w}{ }\PYG{n}{oxdna2}\PYG{o}{/}\PYG{n}{dh}
\PYG{k}{pair\PYGZus{}coeff}\PYG{+w}{ }\PYG{o}{*}\PYG{+w}{ }\PYG{o}{*}\PYG{+w}{ }\PYG{n}{oxdna2}\PYG{o}{/}\PYG{n}{excv}\PYG{+w}{    }\PYG{n}{oxdna2\PYGZus{}lj.cgdna}
\PYG{k}{pair\PYGZus{}coeff}\PYG{+w}{ }\PYG{o}{*}\PYG{+w}{ }\PYG{o}{*}\PYG{+w}{ }\PYG{n}{oxdna2}\PYG{o}{/}\PYG{n}{stk}\PYG{+w}{     }\PYG{n}{seqdep}\PYG{+w}{ }\PYG{l+m}{0.1}\PYG{+w}{ }\PYG{l+m}{1.3523}\PYG{+w}{ }\PYG{l+m}{2.6717}\PYG{+w}{ }\PYG{n}{oxdna2\PYGZus{}lj.cgdna}
\PYG{k}{pair\PYGZus{}coeff}\PYG{+w}{ }\PYG{o}{*}\PYG{+w}{ }\PYG{o}{*}\PYG{+w}{ }\PYG{n}{oxdna2}\PYG{o}{/}\PYG{n}{hbond}\PYG{+w}{   }\PYG{n}{seqdep}\PYG{+w}{ }\PYG{n}{oxdna2\PYGZus{}lj.cgdna}
\PYG{k}{pair\PYGZus{}coeff}\PYG{+w}{ }\PYG{l+m}{1}\PYG{+w}{ }\PYG{l+m}{4}\PYG{+w}{ }\PYG{n}{oxdna2}\PYG{o}{/}\PYG{n}{hbond}\PYG{+w}{   }\PYG{n}{seqdep}\PYG{+w}{ }\PYG{n}{oxdna2\PYGZus{}lj.cgdna}
\PYG{k}{pair\PYGZus{}coeff}\PYG{+w}{ }\PYG{l+m}{2}\PYG{+w}{ }\PYG{l+m}{3}\PYG{+w}{ }\PYG{n}{oxdna2}\PYG{o}{/}\PYG{n}{hbond}\PYG{+w}{   }\PYG{n}{seqdep}\PYG{+w}{ }\PYG{n}{oxdna2\PYGZus{}lj.cgdna}
\PYG{k}{pair\PYGZus{}coeff}\PYG{+w}{ }\PYG{o}{*}\PYG{+w}{ }\PYG{o}{*}\PYG{+w}{ }\PYG{n}{oxdna2}\PYG{o}{/}\PYG{n}{xstk}\PYG{+w}{    }\PYG{n}{oxdna2\PYGZus{}lj.cgdna}
\PYG{k}{pair\PYGZus{}coeff}\PYG{+w}{ }\PYG{o}{*}\PYG{+w}{ }\PYG{o}{*}\PYG{+w}{ }\PYG{n}{oxdna2}\PYG{o}{/}\PYG{n}{coaxstk}\PYG{+w}{ }\PYG{n}{oxdna2\PYGZus{}lj.cgdna}
\PYG{k}{pair\PYGZus{}coeff}\PYG{+w}{ }\PYG{o}{*}\PYG{+w}{ }\PYG{o}{*}\PYG{+w}{ }\PYG{n}{oxdna2}\PYG{o}{/}\PYG{n}{dh}\PYG{+w}{      }\PYG{l+m}{0.1}\PYG{+w}{ }\PYG{l+m}{0.5}\PYG{+w}{ }\PYG{n}{oxdna2\PYGZus{}lj.cgdna}

\PYG{c}{\PYGZsh{} Real units}
\PYG{k}{pair\PYGZus{}style}\PYG{+w}{ }\PYG{n}{hybrid}\PYG{o}{/}\PYG{n}{overlay}\PYG{+w}{ }\PYG{n}{oxdna2}\PYG{o}{/}\PYG{n}{excv}\PYG{+w}{ }\PYG{n}{oxdna2}\PYG{o}{/}\PYG{n}{stk}\PYG{+w}{ }\PYG{n}{oxdna2}\PYG{o}{/}\PYG{n}{hbond}\PYG{+w}{ }\PYG{n}{oxdna2}\PYG{o}{/}\PYG{n}{xstk}\PYG{+w}{ }\PYG{n}{oxdna2}\PYG{o}{/}\PYG{n}{coaxstk}\PYG{+w}{ }\PYG{n}{oxdna2}\PYG{o}{/}\PYG{n}{dh}
\PYG{k}{pair\PYGZus{}coeff}\PYG{+w}{ }\PYG{o}{*}\PYG{+w}{ }\PYG{o}{*}\PYG{+w}{ }\PYG{n}{oxdna2}\PYG{o}{/}\PYG{n}{excv}\PYG{+w}{    }\PYG{l+m}{11.92337812042065}\PYG{+w}{ }\PYG{l+m}{5.9626}\PYG{+w}{ }\PYG{l+m}{5.74965}\PYG{+w}{ }\PYG{l+m}{11.92337812042065}\PYG{+w}{ }\PYG{l+m}{4.38677}\PYG{+w}{ }\PYG{l+m}{4.259}\PYG{+w}{ }\PYG{l+m}{11.92337812042065}\PYG{+w}{ }\PYG{l+m}{2.81094}\PYG{+w}{ }\PYG{l+m}{2.72576}
\PYG{k}{pair\PYGZus{}coeff}\PYG{+w}{ }\PYG{o}{*}\PYG{+w}{ }\PYG{o}{*}\PYG{+w}{ }\PYG{n}{oxdna2}\PYG{o}{/}\PYG{n}{stk}\PYG{+w}{     }\PYG{n}{seqdep}\PYG{+w}{ }\PYG{l+m}{300.0}\PYG{+w}{ }\PYG{l+m}{8.06199211612242}\PYG{+w}{ }\PYG{l+m}{0.005309213}\PYG{+w}{ }\PYG{l+m}{0.70439070204273}\PYG{+w}{ }\PYG{l+m}{3.4072}\PYG{+w}{ }\PYG{l+m}{7.6662}\PYG{+w}{ }\PYG{l+m}{2.72576}\PYG{+w}{ }\PYG{l+m}{6.3885}\PYG{+w}{ }\PYG{l+m}{1.3}\PYG{+w}{ }\PYG{l+m}{0.0}\PYG{+w}{ }\PYG{l+m}{0.8}\PYG{+w}{ }\PYG{l+m}{0.9}\PYG{+w}{ }\PYG{l+m}{0.0}\PYG{+w}{ }\PYG{l+m}{0.95}\PYG{+w}{ }\PYG{l+m}{0.9}\PYG{+w}{ }\PYG{l+m}{0.0}\PYG{+w}{ }\PYG{l+m}{0.95}\PYG{+w}{ }\PYG{l+m}{2.0}\PYG{+w}{ }\PYG{l+m}{0.65}\PYG{+w}{ }\PYG{l+m}{2.0}\PYG{+w}{ }\PYG{l+m}{0.65}
\PYG{k}{pair\PYGZus{}coeff}\PYG{+w}{ }\PYG{o}{*}\PYG{+w}{ }\PYG{o}{*}\PYG{+w}{ }\PYG{n}{oxdna2}\PYG{o}{/}\PYG{n}{hbond}\PYG{+w}{   }\PYG{n}{seqdep}\PYG{+w}{ }\PYG{l+m}{0.0}\PYG{+w}{ }\PYG{l+m}{0.93918760272364}\PYG{+w}{ }\PYG{l+m}{3.4072}\PYG{+w}{ }\PYG{l+m}{6.3885}\PYG{+w}{ }\PYG{l+m}{2.89612}\PYG{+w}{ }\PYG{l+m}{5.9626}\PYG{+w}{ }\PYG{l+m}{1.5}\PYG{+w}{ }\PYG{l+m}{0.0}\PYG{+w}{ }\PYG{l+m}{0.7}\PYG{+w}{ }\PYG{l+m}{1.5}\PYG{+w}{ }\PYG{l+m}{0.0}\PYG{+w}{ }\PYG{l+m}{0.7}\PYG{+w}{ }\PYG{l+m}{1.5}\PYG{+w}{ }\PYG{l+m}{0.0}\PYG{+w}{ }\PYG{l+m}{0.7}\PYG{+w}{ }\PYG{l+m}{0.46}\PYG{+w}{ }\PYG{l+m}{3.141592654}\PYG{+w}{ }\PYG{l+m}{0.7}\PYG{+w}{ }\PYG{l+m}{4.0}\PYG{+w}{ }\PYG{l+m}{1.570796327}\PYG{+w}{ }\PYG{l+m}{0.45}\PYG{+w}{ }\PYG{l+m}{4.0}\PYG{+w}{ }\PYG{l+m}{1.570796327}\PYG{+w}{ }\PYG{l+m}{0.45}
\PYG{k}{pair\PYGZus{}coeff}\PYG{+w}{ }\PYG{l+m}{1}\PYG{+w}{ }\PYG{l+m}{4}\PYG{+w}{ }\PYG{n}{oxdna2}\PYG{o}{/}\PYG{n}{hbond}\PYG{+w}{   }\PYG{n}{seqdep}\PYG{+w}{ }\PYG{l+m}{6.36589157849259}\PYG{+w}{ }\PYG{l+m}{0.93918760272364}\PYG{+w}{ }\PYG{l+m}{3.4072}\PYG{+w}{ }\PYG{l+m}{6.3885}\PYG{+w}{ }\PYG{l+m}{2.89612}\PYG{+w}{ }\PYG{l+m}{5.9626}\PYG{+w}{ }\PYG{l+m}{1.5}\PYG{+w}{ }\PYG{l+m}{0.0}\PYG{+w}{ }\PYG{l+m}{0.7}\PYG{+w}{ }\PYG{l+m}{1.5}\PYG{+w}{ }\PYG{l+m}{0.0}\PYG{+w}{ }\PYG{l+m}{0.7}\PYG{+w}{ }\PYG{l+m}{1.5}\PYG{+w}{ }\PYG{l+m}{0.0}\PYG{+w}{ }\PYG{l+m}{0.7}\PYG{+w}{ }\PYG{l+m}{0.46}\PYG{+w}{ }\PYG{l+m}{3.141592654}\PYG{+w}{ }\PYG{l+m}{0.7}\PYG{+w}{ }\PYG{l+m}{4.0}\PYG{+w}{ }\PYG{l+m}{1.570796327}\PYG{+w}{ }\PYG{l+m}{0.45}\PYG{+w}{ }\PYG{l+m}{4.0}\PYG{+w}{ }\PYG{l+m}{1.570796327}\PYG{+w}{ }\PYG{l+m}{0.45}
\PYG{k}{pair\PYGZus{}coeff}\PYG{+w}{ }\PYG{l+m}{2}\PYG{+w}{ }\PYG{l+m}{3}\PYG{+w}{ }\PYG{n}{oxdna2}\PYG{o}{/}\PYG{n}{hbond}\PYG{+w}{   }\PYG{n}{seqdep}\PYG{+w}{ }\PYG{l+m}{6.36589157849259}\PYG{+w}{ }\PYG{l+m}{0.93918760272364}\PYG{+w}{ }\PYG{l+m}{3.4072}\PYG{+w}{ }\PYG{l+m}{6.3885}\PYG{+w}{ }\PYG{l+m}{2.89612}\PYG{+w}{ }\PYG{l+m}{5.9626}\PYG{+w}{ }\PYG{l+m}{1.5}\PYG{+w}{ }\PYG{l+m}{0.0}\PYG{+w}{ }\PYG{l+m}{0.7}\PYG{+w}{ }\PYG{l+m}{1.5}\PYG{+w}{ }\PYG{l+m}{0.0}\PYG{+w}{ }\PYG{l+m}{0.7}\PYG{+w}{ }\PYG{l+m}{1.5}\PYG{+w}{ }\PYG{l+m}{0.0}\PYG{+w}{ }\PYG{l+m}{0.7}\PYG{+w}{ }\PYG{l+m}{0.46}\PYG{+w}{ }\PYG{l+m}{3.141592654}\PYG{+w}{ }\PYG{l+m}{0.7}\PYG{+w}{ }\PYG{l+m}{4.0}\PYG{+w}{ }\PYG{l+m}{1.570796327}\PYG{+w}{ }\PYG{l+m}{0.45}\PYG{+w}{ }\PYG{l+m}{4.0}\PYG{+w}{ }\PYG{l+m}{1.570796327}\PYG{+w}{ }\PYG{l+m}{0.45}
\PYG{k}{pair\PYGZus{}coeff}\PYG{+w}{ }\PYG{o}{*}\PYG{+w}{ }\PYG{o}{*}\PYG{+w}{ }\PYG{n}{oxdna2}\PYG{o}{/}\PYG{n}{xstk}\PYG{+w}{    }\PYG{l+m}{3.9029021145006}\PYG{+w}{ }\PYG{l+m}{4.89785}\PYG{+w}{ }\PYG{l+m}{5.74965}\PYG{+w}{ }\PYG{l+m}{4.21641}\PYG{+w}{ }\PYG{l+m}{5.57929}\PYG{+w}{ }\PYG{l+m}{2.25}\PYG{+w}{ }\PYG{l+m}{0.791592654}\PYG{+w}{ }\PYG{l+m}{0.58}\PYG{+w}{ }\PYG{l+m}{1.7}\PYG{+w}{ }\PYG{l+m}{1.0}\PYG{+w}{ }\PYG{l+m}{0.68}\PYG{+w}{ }\PYG{l+m}{1.7}\PYG{+w}{ }\PYG{l+m}{1.0}\PYG{+w}{ }\PYG{l+m}{0.68}\PYG{+w}{ }\PYG{l+m}{1.5}\PYG{+w}{ }\PYG{l+m}{0.0}\PYG{+w}{ }\PYG{l+m}{0.65}\PYG{+w}{ }\PYG{l+m}{1.7}\PYG{+w}{ }\PYG{l+m}{0.875}\PYG{+w}{ }\PYG{l+m}{0.68}\PYG{+w}{ }\PYG{l+m}{1.7}\PYG{+w}{ }\PYG{l+m}{0.875}\PYG{+w}{ }\PYG{l+m}{0.68}
\PYG{k}{pair\PYGZus{}coeff}\PYG{+w}{ }\PYG{o}{*}\PYG{+w}{ }\PYG{o}{*}\PYG{+w}{ }\PYG{n}{oxdna2}\PYG{o}{/}\PYG{n}{coaxstk}\PYG{+w}{ }\PYG{l+m}{4.80673207785863}\PYG{+w}{ }\PYG{l+m}{3.4072}\PYG{+w}{ }\PYG{l+m}{5.1108}\PYG{+w}{ }\PYG{l+m}{1.87396}\PYG{+w}{ }\PYG{l+m}{4.94044}\PYG{+w}{ }\PYG{l+m}{2.0}\PYG{+w}{ }\PYG{l+m}{2.891592653589793}\PYG{+w}{ }\PYG{l+m}{0.65}\PYG{+w}{ }\PYG{l+m}{1.3}\PYG{+w}{ }\PYG{l+m}{0.0}\PYG{+w}{ }\PYG{l+m}{0.8}\PYG{+w}{ }\PYG{l+m}{0.9}\PYG{+w}{ }\PYG{l+m}{0.0}\PYG{+w}{ }\PYG{l+m}{0.95}\PYG{+w}{ }\PYG{l+m}{0.9}\PYG{+w}{ }\PYG{l+m}{0.0}\PYG{+w}{ }\PYG{l+m}{0.95}\PYG{+w}{ }\PYG{l+m}{40.0}\PYG{+w}{ }\PYG{l+m}{3.116592653589793}
\PYG{k}{pair\PYGZus{}coeff}\PYG{+w}{ }\PYG{o}{*}\PYG{+w}{ }\PYG{o}{*}\PYG{+w}{ }\PYG{n}{oxdna2}\PYG{o}{/}\PYG{n}{dh}\PYG{+w}{      }\PYG{l+m}{300.0}\PYG{+w}{ }\PYG{l+m}{0.5}\PYG{+w}{ }\PYG{l+m}{0.815}

\PYG{k}{pair\PYGZus{}style}\PYG{+w}{ }\PYG{n}{hybrid}\PYG{o}{/}\PYG{n}{overlay}\PYG{+w}{ }\PYG{n}{oxdna2}\PYG{o}{/}\PYG{n}{excv}\PYG{+w}{ }\PYG{n}{oxdna2}\PYG{o}{/}\PYG{n}{stk}\PYG{+w}{ }\PYG{n}{oxdna2}\PYG{o}{/}\PYG{n}{hbond}\PYG{+w}{ }\PYG{n}{oxdna2}\PYG{o}{/}\PYG{n}{xstk}\PYG{+w}{ }\PYG{n}{oxdna2}\PYG{o}{/}\PYG{n}{coaxstk}\PYG{+w}{ }\PYG{n}{oxdna2}\PYG{o}{/}\PYG{n}{dh}
\PYG{k}{pair\PYGZus{}coeff}\PYG{+w}{ }\PYG{o}{*}\PYG{+w}{ }\PYG{o}{*}\PYG{+w}{ }\PYG{n}{oxdna2}\PYG{o}{/}\PYG{n}{excv}\PYG{+w}{    }\PYG{n}{oxdna2\PYGZus{}real.cgdna}
\PYG{k}{pair\PYGZus{}coeff}\PYG{+w}{ }\PYG{o}{*}\PYG{+w}{ }\PYG{o}{*}\PYG{+w}{ }\PYG{n}{oxdna2}\PYG{o}{/}\PYG{n}{stk}\PYG{+w}{     }\PYG{n}{seqdep}\PYG{+w}{ }\PYG{l+m}{300.0}\PYG{+w}{ }\PYG{l+m}{8.06199211612242}\PYG{+w}{ }\PYG{l+m}{0.005309213}\PYG{+w}{ }\PYG{n}{oxdna2\PYGZus{}real.cgdna}
\PYG{k}{pair\PYGZus{}coeff}\PYG{+w}{ }\PYG{o}{*}\PYG{+w}{ }\PYG{o}{*}\PYG{+w}{ }\PYG{n}{oxdna2}\PYG{o}{/}\PYG{n}{hbond}\PYG{+w}{   }\PYG{n}{seqdep}\PYG{+w}{ }\PYG{n}{oxdna2\PYGZus{}real.cgdna}
\PYG{k}{pair\PYGZus{}coeff}\PYG{+w}{ }\PYG{l+m}{1}\PYG{+w}{ }\PYG{l+m}{4}\PYG{+w}{ }\PYG{n}{oxdna2}\PYG{o}{/}\PYG{n}{hbond}\PYG{+w}{   }\PYG{n}{seqdep}\PYG{+w}{ }\PYG{n}{oxdna2\PYGZus{}real.cgdna}
\PYG{k}{pair\PYGZus{}coeff}\PYG{+w}{ }\PYG{l+m}{2}\PYG{+w}{ }\PYG{l+m}{3}\PYG{+w}{ }\PYG{n}{oxdna2}\PYG{o}{/}\PYG{n}{hbond}\PYG{+w}{   }\PYG{n}{seqdep}\PYG{+w}{ }\PYG{n}{oxdna2\PYGZus{}real.cgdna}
\PYG{k}{pair\PYGZus{}coeff}\PYG{+w}{ }\PYG{o}{*}\PYG{+w}{ }\PYG{o}{*}\PYG{+w}{ }\PYG{n}{oxdna2}\PYG{o}{/}\PYG{n}{xstk}\PYG{+w}{    }\PYG{n}{oxdna2\PYGZus{}real.cgdna}
\PYG{k}{pair\PYGZus{}coeff}\PYG{+w}{ }\PYG{o}{*}\PYG{+w}{ }\PYG{o}{*}\PYG{+w}{ }\PYG{n}{oxdna2}\PYG{o}{/}\PYG{n}{coaxstk}\PYG{+w}{ }\PYG{n}{oxdna2\PYGZus{}real.cgdna}
\PYG{k}{pair\PYGZus{}coeff}\PYG{+w}{ }\PYG{o}{*}\PYG{+w}{ }\PYG{o}{*}\PYG{+w}{ }\PYG{n}{oxdna2}\PYG{o}{/}\PYG{n}{dh}\PYG{+w}{      }\PYG{l+m}{300.0}\PYG{+w}{ }\PYG{l+m}{0.5}\PYG{+w}{ }\PYG{n}{oxdna2\PYGZus{}real.cgdna}
\end{sphinxVerbatim}

\begin{sphinxadmonition}{note}{Note:}
\sphinxAtStartPar
The coefficients in the above examples are provided in forms
compatible with both \sphinxstyleemphasis{units lj} and \sphinxstyleemphasis{units real} (see documentation
of {\hyperref[\detokenize{units::doc}]{\sphinxcrossref{\DUrole{doc}{units}}}}).  These can also be read from a potential
file with correct unit style by specifying the name of the
file. Several potential files for each unit style are included in the
\sphinxcode{\sphinxupquote{potentials}} directory of the LAMMPS distribution.
\end{sphinxadmonition}


\subsection{Description}
\label{\detokenize{pair_oxdna2:description}}
\sphinxAtStartPar
The \sphinxstyleemphasis{oxdna2} pair styles compute the pairwise\sphinxhyphen{}additive parts of the
oxDNA force field for coarse\sphinxhyphen{}grained modelling of DNA. The effective
interaction between the nucleotides consists of potentials for the
excluded volume interaction \sphinxstyleemphasis{oxdna2/excv}, the stacking \sphinxstyleemphasis{oxdna2/stk},
cross\sphinxhyphen{}stacking \sphinxstyleemphasis{oxdna2/xstk} and coaxial stacking interaction
\sphinxstyleemphasis{oxdna2/coaxstk}, electrostatic Debye\sphinxhyphen{}Hueckel interaction \sphinxstyleemphasis{oxdna2/dh} as
well as the hydrogen\sphinxhyphen{}bonding interaction \sphinxstyleemphasis{oxdna2/hbond} between
complementary pairs of nucleotides on opposite strands. Average sequence
or sequence\sphinxhyphen{}dependent stacking and base\sphinxhyphen{}pairing strengths are supported
{\hyperref[\detokenize{pair_oxdna2:sulc2}]{\sphinxcrossref{\DUrole{std,std-ref}{(Sulc)}}}}. Quasi\sphinxhyphen{}unique base\sphinxhyphen{}pairing between nucleotides can
be achieved by using more complementary pairs of atom types like 5\sphinxhyphen{}8 and
6\sphinxhyphen{}7, 9\sphinxhyphen{}12 and 10\sphinxhyphen{}11, 13\sphinxhyphen{}16 and 14\sphinxhyphen{}15, etc.  This prevents the
hybridization of in principle complementary bases within Ntypes/4 bases
up and down along the backbone.

\sphinxAtStartPar
The exact functional form of the pair styles is rather complex.  The
individual potentials consist of products of modulation factors, which
themselves are constructed from a number of more basic potentials
(Morse, Lennard\sphinxhyphen{}Jones, harmonic angle and distance) as well as quadratic
smoothing and modulation terms.  We refer to {\hyperref[\detokenize{pair_oxdna2:snodin2}]{\sphinxcrossref{\DUrole{std,std-ref}{(Snodin)}}}}
and the original oxDNA publications {\hyperref[\detokenize{pair_oxdna2:ouldridge-dphil2}]{\sphinxcrossref{\DUrole{std,std-ref}{(Ouldridge\sphinxhyphen{}DPhil)}}}} and {\hyperref[\detokenize{pair_oxdna2:ouldridge2}]{\sphinxcrossref{\DUrole{std,std-ref}{(Ouldridge)}}}} for a detailed
description of the oxDNA2 force field.

\begin{sphinxadmonition}{note}{Note:}
\sphinxAtStartPar
These pair styles have to be used together with the related oxDNA2
bond style \sphinxstyleemphasis{oxdna2/fene} for the connectivity of the phosphate
backbone (see also documentation of {\hyperref[\detokenize{bond_oxdna::doc}]{\sphinxcrossref{\DUrole{doc}{bond\_style oxdna2/fene}}}}). Most of the coefficients in the above example have to
be kept fixed and cannot be changed without reparameterizing the
entire model.  Exceptions are the first four coefficients after
\sphinxstyleemphasis{oxdna2/stk} (seq=seqdep, T=0.1, xi=1.3523 and kappa=2.6717 and
corresponding \sphinxstyleemphasis{real unit} equivalents in the above examples).  the
first coefficient after \sphinxstyleemphasis{oxdna2/hbond} (seq=seqdep in the above
example) and the three coefficients after \sphinxstyleemphasis{oxdna2/dh} (T=0.1,
rhos=0.5, qeff=0.815 in the above example). When using a Langevin
thermostat e.g. through {\hyperref[\detokenize{fix_langevin::doc}]{\sphinxcrossref{\DUrole{doc}{fix langevin}}}} or
{\hyperref[\detokenize{fix_nve_dotc_langevin::doc}]{\sphinxcrossref{\DUrole{doc}{fix nve/dotc/langevin}}}} the temperature
coefficients have to be matched to the one used in the fix.
\end{sphinxadmonition}

\begin{sphinxadmonition}{note}{Note:}
\sphinxAtStartPar
These pair styles have to be used with the \sphinxstyleemphasis{atom\_style hybrid bond
ellipsoid oxdna} (see documentation of {\hyperref[\detokenize{atom_style::doc}]{\sphinxcrossref{\DUrole{doc}{atom\_style}}}}). The \sphinxstyleemphasis{atom\_style oxdna} stores the 3’\sphinxhyphen{}to\sphinxhyphen{}5’ polarity
of the nucleotide strand, which is set through the bond topology in
the data file. The first (second) atom in a bond definition is
understood to point towards the 3’\sphinxhyphen{}end (5’\sphinxhyphen{}end) of the strand.
\end{sphinxadmonition}

\sphinxAtStartPar
Example input and data files for DNA duplexes can be found in
\sphinxcode{\sphinxupquote{examples/PACKAGES/cgdna/examples/oxDNA/}} and \sphinxcode{\sphinxupquote{.../oxDNA2/}}.  A
simple python setup tool which creates single straight or helical DNA
strands, DNA duplexes or arrays of DNA duplexes can be found in
\sphinxcode{\sphinxupquote{examples/PACKAGES/cgdna/util/}}.

\sphinxAtStartPar
Please cite {\hyperref[\detokenize{pair_oxdna2:henrich2}]{\sphinxcrossref{\DUrole{std,std-ref}{(Henrich)}}}} in any publication that uses
this implementation. An updated documentation that contains general
information on the model, its implementation and performance as well as
the structure of the data and input file can be found \sphinxhref{PDF/CG-DNA.pdf}{here}.

\sphinxAtStartPar
Please cite also the relevant oxDNA2 publications
{\hyperref[\detokenize{pair_oxdna2:snodin2}]{\sphinxcrossref{\DUrole{std,std-ref}{(Snodin)}}}} and {\hyperref[\detokenize{pair_oxdna2:sulc2}]{\sphinxcrossref{\DUrole{std,std-ref}{(Sulc)}}}}.


\bigskip\hrule\bigskip



\subsection{Potential file reading}
\label{\detokenize{pair_oxdna2:potential-file-reading}}
\sphinxAtStartPar
For each pair style above the first non\sphinxhyphen{}modifiable argument can be a
filename (with exception of Debye\sphinxhyphen{}Hueckel, for which the effective
charge argument can be a filename), and if it is, no further arguments
should be supplied.  Therefore the following command:

\begin{sphinxVerbatim}[commandchars=\\\{\}]
\PYG{k}{pair\PYGZus{}coeff}\PYG{+w}{ }\PYG{l+m}{1}\PYG{+w}{ }\PYG{l+m}{4}\PYG{+w}{ }\PYG{n}{oxdna2}\PYG{o}{/}\PYG{n}{hbond}\PYG{+w}{   }\PYG{n}{seqdep}\PYG{+w}{ }\PYG{n}{oxdna\PYGZus{}real.cgdna}
\end{sphinxVerbatim}

\sphinxAtStartPar
will be interpreted as a request to read the corresponding hydrogen
bonding potential parameters from the file with the given name.  The
file can define multiple potential parameters for both bonded and pair
interactions, but for the example pair interaction above there must
exist in the file a line of the form:

\begin{sphinxVerbatim}[commandchars=\\\{\}]
\PYG{l+m}{1}\PYG{+w}{ }\PYG{l+m}{4}\PYG{+w}{ }\PYG{n}{hbond}\PYG{+w}{     }\PYG{o}{\PYGZlt{}}\PYG{n}{coefficients}\PYG{o}{\PYGZgt{}}
\end{sphinxVerbatim}

\sphinxAtStartPar
If potential customization is required, the potential file reading can
be mixed with the manual specification of the potential parameters. For
example, the following command:

\begin{sphinxVerbatim}[commandchars=\\\{\}]
\PYG{k}{pair\PYGZus{}style}\PYG{+w}{ }\PYG{n}{hybrid}\PYG{o}{/}\PYG{n}{overlay}\PYG{+w}{ }\PYG{n}{oxdna2}\PYG{o}{/}\PYG{n}{excv}\PYG{+w}{ }\PYG{n}{oxdna2}\PYG{o}{/}\PYG{n}{stk}\PYG{+w}{ }\PYG{n}{oxdna2}\PYG{o}{/}\PYG{n}{hbond}\PYG{+w}{ }\PYG{n}{oxdna2}\PYG{o}{/}\PYG{n}{xstk}\PYG{+w}{ }\PYG{n}{oxdna2}\PYG{o}{/}\PYG{n}{coaxstk}\PYG{+w}{ }\PYG{n}{oxdna2}\PYG{o}{/}\PYG{n}{dh}
\PYG{k}{pair\PYGZus{}coeff}\PYG{+w}{ }\PYG{o}{*}\PYG{+w}{ }\PYG{o}{*}\PYG{+w}{ }\PYG{n}{oxdna2}\PYG{o}{/}\PYG{n}{excv}\PYG{+w}{    }\PYG{l+m}{2.0}\PYG{+w}{ }\PYG{l+m}{0.7}\PYG{+w}{ }\PYG{l+m}{0.675}\PYG{+w}{ }\PYG{l+m}{2.0}\PYG{+w}{ }\PYG{l+m}{0.515}\PYG{+w}{ }\PYG{l+m}{0.5}\PYG{+w}{ }\PYG{l+m}{2.0}\PYG{+w}{ }\PYG{l+m}{0.33}\PYG{+w}{ }\PYG{l+m}{0.32}
\PYG{k}{pair\PYGZus{}coeff}\PYG{+w}{ }\PYG{o}{*}\PYG{+w}{ }\PYG{o}{*}\PYG{+w}{ }\PYG{n}{oxdna2}\PYG{o}{/}\PYG{n}{stk}\PYG{+w}{     }\PYG{n}{seqdep}\PYG{+w}{ }\PYG{l+m}{0.1}\PYG{+w}{ }\PYG{l+m}{1.3523}\PYG{+w}{ }\PYG{l+m}{2.6717}\PYG{+w}{ }\PYG{n}{oxdna2\PYGZus{}lj.cgdna}
\PYG{k}{pair\PYGZus{}coeff}\PYG{+w}{ }\PYG{o}{*}\PYG{+w}{ }\PYG{o}{*}\PYG{+w}{ }\PYG{n}{oxdna2}\PYG{o}{/}\PYG{n}{hbond}\PYG{+w}{   }\PYG{n}{seqdep}\PYG{+w}{ }\PYG{n}{oxdna2\PYGZus{}lj.cgdna}
\PYG{k}{pair\PYGZus{}coeff}\PYG{+w}{ }\PYG{l+m}{1}\PYG{+w}{ }\PYG{l+m}{4}\PYG{+w}{ }\PYG{n}{oxdna2}\PYG{o}{/}\PYG{n}{hbond}\PYG{+w}{   }\PYG{n}{seqdep}\PYG{+w}{ }\PYG{n}{oxdna2\PYGZus{}lj.cgdna}
\PYG{k}{pair\PYGZus{}coeff}\PYG{+w}{ }\PYG{l+m}{2}\PYG{+w}{ }\PYG{l+m}{3}\PYG{+w}{ }\PYG{n}{oxdna2}\PYG{o}{/}\PYG{n}{hbond}\PYG{+w}{   }\PYG{n}{seqdep}\PYG{+w}{ }\PYG{n}{oxdna2\PYGZus{}lj.cgdna}
\PYG{k}{pair\PYGZus{}coeff}\PYG{+w}{ }\PYG{o}{*}\PYG{+w}{ }\PYG{o}{*}\PYG{+w}{ }\PYG{n}{oxdna2}\PYG{o}{/}\PYG{n}{xstk}\PYG{+w}{    }\PYG{n}{oxdna2\PYGZus{}lj.cgdna}
\PYG{k}{pair\PYGZus{}coeff}\PYG{+w}{ }\PYG{o}{*}\PYG{+w}{ }\PYG{o}{*}\PYG{+w}{ }\PYG{n}{oxdna2}\PYG{o}{/}\PYG{n}{coaxstk}\PYG{+w}{ }\PYG{n}{oxdna2\PYGZus{}lj.cgdna}
\PYG{k}{pair\PYGZus{}coeff}\PYG{+w}{ }\PYG{o}{*}\PYG{+w}{ }\PYG{o}{*}\PYG{+w}{ }\PYG{n}{oxdna2}\PYG{o}{/}\PYG{n}{dh}\PYG{+w}{      }\PYG{l+m}{0.1}\PYG{+w}{ }\PYG{l+m}{0.5}\PYG{+w}{ }\PYG{l+m}{0.815}
\end{sphinxVerbatim}

\sphinxAtStartPar
will read the excluded volume and Debye\sphinxhyphen{}Hueckel effective charge \sphinxstyleemphasis{qeff}
parameters from the manual specification and all others from the
potential file \sphinxstyleemphasis{oxdna2\_lj.cgdna}.

\sphinxAtStartPar
There are sample potential files for each unit style in the \sphinxcode{\sphinxupquote{potentials}}
directory of the LAMMPS distribution. The potential file unit system
must align with the units defined via the {\hyperref[\detokenize{units::doc}]{\sphinxcrossref{\DUrole{doc}{units}}}}
command. For conversion between different \sphinxstyleemphasis{LJ} and \sphinxstyleemphasis{real} unit systems
for oxDNA, the python tool \sphinxstyleemphasis{lj2real.py} located in the
\sphinxcode{\sphinxupquote{examples/PACKAGES/cgdna/util/}} directory can be used. This tool assumes
similar file structure to the examples found in
\sphinxcode{\sphinxupquote{examples/PACKAGES/cgdna/examples/}}.


\bigskip\hrule\bigskip



\subsection{Restrictions}
\label{\detokenize{pair_oxdna2:restrictions}}
\sphinxAtStartPar
These pair styles can only be used if LAMMPS was built with the
CG\sphinxhyphen{}DNA package and the MOLECULE and ASPHERE package.  See the
{\hyperref[\detokenize{Build_package::doc}]{\sphinxcrossref{\DUrole{doc}{Build package}}}} page for more info.


\subsection{Related commands}
\label{\detokenize{pair_oxdna2:related-commands}}
\sphinxAtStartPar
{\hyperref[\detokenize{bond_oxdna::doc}]{\sphinxcrossref{\DUrole{doc}{bond\_style oxdna2/fene}}}}, {\hyperref[\detokenize{pair_coeff::doc}]{\sphinxcrossref{\DUrole{doc}{pair\_coeff}}}},
{\hyperref[\detokenize{bond_oxdna::doc}]{\sphinxcrossref{\DUrole{doc}{bond\_style oxdna/fene}}}}, {\hyperref[\detokenize{pair_oxdna::doc}]{\sphinxcrossref{\DUrole{doc}{pair\_style oxdna/excv}}}},
{\hyperref[\detokenize{bond_oxdna::doc}]{\sphinxcrossref{\DUrole{doc}{bond\_style oxrna2/fene}}}}, {\hyperref[\detokenize{pair_oxrna2::doc}]{\sphinxcrossref{\DUrole{doc}{pair\_style oxrna2/excv}}}},
{\hyperref[\detokenize{atom_style::doc}]{\sphinxcrossref{\DUrole{doc}{atom\_style oxdna}}}}, {\hyperref[\detokenize{fix_nve_dotc_langevin::doc}]{\sphinxcrossref{\DUrole{doc}{fix nve/dotc/langevin}}}}


\subsection{Default}
\label{\detokenize{pair_oxdna2:default}}
\sphinxAtStartPar
none


\bigskip\hrule\bigskip

\phantomsection\label{\detokenize{pair_oxdna2:henrich2}}
\sphinxAtStartPar
\sphinxstylestrong{(Henrich)} O. Henrich, Y. A. Gutierrez\sphinxhyphen{}Fosado, T. Curk, T. E. Ouldridge, Eur. Phys. J. E 41, 57 (2018).

\phantomsection\label{\detokenize{pair_oxdna2:snodin2}}
\sphinxAtStartPar
\sphinxstylestrong{(Snodin)} B.E. Snodin, F. Randisi, M. Mosayebi, et al., J. Chem. Phys. 142, 234901 (2015).

\phantomsection\label{\detokenize{pair_oxdna2:sulc2}}
\sphinxAtStartPar
\sphinxstylestrong{(Sulc)} P. Sulc, F. Romano, T.E. Ouldridge, L. Rovigatti, J.P.K. Doye, A.A. Louis, J. Chem. Phys. 137, 135101 (2012).

\phantomsection\label{\detokenize{pair_oxdna2:ouldridge-dphil2}}
\sphinxAtStartPar
\sphinxstylestrong{(Ouldridge\sphinxhyphen{}DPhil)} T.E. Ouldridge, Coarse\sphinxhyphen{}grained modelling of DNA and DNA self\sphinxhyphen{}assembly, DPhil. University of Oxford (2011).

\phantomsection\label{\detokenize{pair_oxdna2:ouldridge2}}
\sphinxAtStartPar
\sphinxstylestrong{(Ouldridge)} T.E. Ouldridge, A.A. Louis, J.P.K. Doye, J. Chem. Phys. 134, 085101 (2011).

\sphinxstepscope

\index{pair\_style oxrna2/excv@\spxentry{pair\_style oxrna2/excv}}\ignorespaces 
\index{pair\_style oxrna2/stk@\spxentry{pair\_style oxrna2/stk}}\ignorespaces 
\index{pair\_style oxrna2/hbond@\spxentry{pair\_style oxrna2/hbond}}\ignorespaces 
\index{pair\_style oxrna2/xstk@\spxentry{pair\_style oxrna2/xstk}}\ignorespaces 
\index{pair\_style oxrna2/coaxstk@\spxentry{pair\_style oxrna2/coaxstk}}\ignorespaces 
\index{pair\_style oxrna2/dh@\spxentry{pair\_style oxrna2/dh}}\ignorespaces 

\section{pair\_style oxrna2/excv command}
\label{\detokenize{pair_oxrna2:pair-style-oxrna2-excv-command}}\label{\detokenize{pair_oxrna2:index-5}}\label{\detokenize{pair_oxrna2:index-4}}\label{\detokenize{pair_oxrna2:index-3}}\label{\detokenize{pair_oxrna2:index-2}}\label{\detokenize{pair_oxrna2:index-1}}\label{\detokenize{pair_oxrna2:index-0}}\label{\detokenize{pair_oxrna2::doc}}

\section{pair\_style oxrna2/stk command}
\label{\detokenize{pair_oxrna2:pair-style-oxrna2-stk-command}}

\section{pair\_style oxrna2/hbond command}
\label{\detokenize{pair_oxrna2:pair-style-oxrna2-hbond-command}}

\section{pair\_style oxrna2/xstk command}
\label{\detokenize{pair_oxrna2:pair-style-oxrna2-xstk-command}}

\section{pair\_style oxrna2/coaxstk command}
\label{\detokenize{pair_oxrna2:pair-style-oxrna2-coaxstk-command}}

\section{pair\_style oxrna2/dh command}
\label{\detokenize{pair_oxrna2:pair-style-oxrna2-dh-command}}

\subsection{Syntax}
\label{\detokenize{pair_oxrna2:syntax}}
\begin{sphinxVerbatim}[commandchars=\\\{\}]
\PYG{k}{pair\PYGZus{}style}\PYG{+w}{ }\PYG{n}{style1}

\PYG{k}{pair\PYGZus{}coeff}\PYG{+w}{ }\PYG{o}{*}\PYG{+w}{ }\PYG{o}{*}\PYG{+w}{ }\PYG{n}{style2}\PYG{+w}{ }\PYG{n}{args}
\end{sphinxVerbatim}
\begin{itemize}
\item {} 
\sphinxAtStartPar
style1 = \sphinxstyleemphasis{hybrid/overlay oxrna2/excv oxrna2/stk oxrna2/hbond oxrna2/xstk oxrna2/coaxstk oxrna2/dh}

\item {} 
\sphinxAtStartPar
style2 = \sphinxstyleemphasis{oxrna2/excv} or \sphinxstyleemphasis{oxrna2/stk} or \sphinxstyleemphasis{oxrna2/hbond} or \sphinxstyleemphasis{oxrna2/xstk} or \sphinxstyleemphasis{oxrna2/coaxstk} or \sphinxstyleemphasis{oxrna2/dh}

\item {} 
\sphinxAtStartPar
args = list of arguments for these particular styles

\end{itemize}
\begin{sphinxalltt}
\sphinxstyleemphasis{oxrna2/stk} args = seq T xi kappa 6.0 0.43 0.93 0.35 0.78 0.9 0 0.95 0.9 0 0.95 1.3 0 0.8 1.3 0 0.8 2.0 0.65 2.0 0.65
  seq = seqav (for average sequence stacking strength) or seqdep (for sequence\sphinxhyphen{}dependent stacking strength)
  T = temperature (LJ units: 0.1 = 300 K, real units: 300 = 300 K)
  xi = 1.40206 (LJ units) or 8.35864576375849 (real units), temperature\sphinxhyphen{}independent coefficient in stacking strength
  kappa = 2.77 (LJ units) or 0.005504556 (real units), coefficient of linear temperature dependence in stacking strength
\sphinxstyleemphasis{oxrna2/hbond} args = seq eps 8.0 0.4 0.75 0.34 0.7 1.5 0 0.7 1.5 0 0.7 1.5 0 0.7 0.46 3.141592653589793 0.7 4.0 1.5707963267948966 0.45 4.0 1.5707963267948966 0.45
  seq = seqav (for average sequence base\sphinxhyphen{}pairing strength) or seqdep (for sequence\sphinxhyphen{}dependent base\sphinxhyphen{}pairing strength)
  eps = 0.870439 (LJ units) or 5.18928666388042 (real units), average hydrogen bonding strength between A\sphinxhyphen{}U and C\sphinxhyphen{}G Watson\sphinxhyphen{}Crick and G\sphinxhyphen{}U wobble base pairs, 0 between all other pairs
\sphinxstyleemphasis{oxrna2/dh} args = T rhos qeff
  T = temperature (LJ units: 0.1 = 300 K, real units: 300 = 300 K)
  rhos = salt concentration (mole per litre)
  qeff = 1.02455 (effective charge in elementary charges)
\end{sphinxalltt}


\subsection{Examples}
\label{\detokenize{pair_oxrna2:examples}}
\begin{sphinxVerbatim}[commandchars=\\\{\}]
\PYG{c}{\PYGZsh{} LJ units}
\PYG{k}{pair\PYGZus{}style}\PYG{+w}{ }\PYG{n}{hybrid}\PYG{o}{/}\PYG{n}{overlay}\PYG{+w}{ }\PYG{n}{oxrna2}\PYG{o}{/}\PYG{n}{excv}\PYG{+w}{ }\PYG{n}{oxrna2}\PYG{o}{/}\PYG{n}{stk}\PYG{+w}{ }\PYG{n}{oxrna2}\PYG{o}{/}\PYG{n}{hbond}\PYG{+w}{ }\PYG{n}{oxrna2}\PYG{o}{/}\PYG{n}{xstk}\PYG{+w}{ }\PYG{n}{oxrna2}\PYG{o}{/}\PYG{n}{coaxstk}\PYG{+w}{ }\PYG{n}{oxrna2}\PYG{o}{/}\PYG{n}{dh}
\PYG{k}{pair\PYGZus{}coeff}\PYG{+w}{ }\PYG{o}{*}\PYG{+w}{ }\PYG{o}{*}\PYG{+w}{ }\PYG{n}{oxrna2}\PYG{o}{/}\PYG{n}{excv}\PYG{+w}{    }\PYG{l+m}{2.0}\PYG{+w}{ }\PYG{l+m}{0.7}\PYG{+w}{ }\PYG{l+m}{0.675}\PYG{+w}{ }\PYG{l+m}{2.0}\PYG{+w}{ }\PYG{l+m}{0.515}\PYG{+w}{ }\PYG{l+m}{0.5}\PYG{+w}{ }\PYG{l+m}{2.0}\PYG{+w}{ }\PYG{l+m}{0.33}\PYG{+w}{ }\PYG{l+m}{0.32}
\PYG{k}{pair\PYGZus{}coeff}\PYG{+w}{ }\PYG{o}{*}\PYG{+w}{ }\PYG{o}{*}\PYG{+w}{ }\PYG{n}{oxrna2}\PYG{o}{/}\PYG{n}{stk}\PYG{+w}{     }\PYG{n}{seqdep}\PYG{+w}{ }\PYG{l+m}{0.1}\PYG{+w}{ }\PYG{l+m}{1.40206}\PYG{+w}{ }\PYG{l+m}{2.77}\PYG{+w}{ }\PYG{l+m}{6.0}\PYG{+w}{ }\PYG{l+m}{0.43}\PYG{+w}{ }\PYG{l+m}{0.93}\PYG{+w}{ }\PYG{l+m}{0.35}\PYG{+w}{ }\PYG{l+m}{0.78}\PYG{+w}{ }\PYG{l+m}{0.9}\PYG{+w}{ }\PYG{l+m}{0}\PYG{+w}{ }\PYG{l+m}{0.95}\PYG{+w}{ }\PYG{l+m}{0.9}\PYG{+w}{ }\PYG{l+m}{0}\PYG{+w}{ }\PYG{l+m}{0.95}\PYG{+w}{ }\PYG{l+m}{1.3}\PYG{+w}{ }\PYG{l+m}{0}\PYG{+w}{ }\PYG{l+m}{0.8}\PYG{+w}{ }\PYG{l+m}{1.3}\PYG{+w}{ }\PYG{l+m}{0}\PYG{+w}{ }\PYG{l+m}{0.8}\PYG{+w}{ }\PYG{l+m}{2.0}\PYG{+w}{ }\PYG{l+m}{0.65}\PYG{+w}{ }\PYG{l+m}{2.0}\PYG{+w}{ }\PYG{l+m}{0.65}
\PYG{k}{pair\PYGZus{}coeff}\PYG{+w}{ }\PYG{o}{*}\PYG{+w}{ }\PYG{o}{*}\PYG{+w}{ }\PYG{n}{oxrna2}\PYG{o}{/}\PYG{n}{hbond}\PYG{+w}{   }\PYG{n}{seqdep}\PYG{+w}{ }\PYG{l+m}{0.0}\PYG{+w}{ }\PYG{l+m}{8.0}\PYG{+w}{ }\PYG{l+m}{0.4}\PYG{+w}{ }\PYG{l+m}{0.75}\PYG{+w}{ }\PYG{l+m}{0.34}\PYG{+w}{ }\PYG{l+m}{0.7}\PYG{+w}{ }\PYG{l+m}{1.5}\PYG{+w}{ }\PYG{l+m}{0}\PYG{+w}{ }\PYG{l+m}{0.7}\PYG{+w}{ }\PYG{l+m}{1.5}\PYG{+w}{ }\PYG{l+m}{0}\PYG{+w}{ }\PYG{l+m}{0.7}\PYG{+w}{ }\PYG{l+m}{1.5}\PYG{+w}{ }\PYG{l+m}{0}\PYG{+w}{ }\PYG{l+m}{0.7}\PYG{+w}{ }\PYG{l+m}{0.46}\PYG{+w}{ }\PYG{l+m}{3.141592653589793}\PYG{+w}{ }\PYG{l+m}{0.7}\PYG{+w}{ }\PYG{l+m}{4.0}\PYG{+w}{ }\PYG{l+m}{1.5707963267948966}\PYG{+w}{ }\PYG{l+m}{0.45}\PYG{+w}{ }\PYG{l+m}{4.0}\PYG{+w}{ }\PYG{l+m}{1.5707963267948966}\PYG{+w}{ }\PYG{l+m}{0.45}
\PYG{k}{pair\PYGZus{}coeff}\PYG{+w}{ }\PYG{l+m}{1}\PYG{+w}{ }\PYG{l+m}{4}\PYG{+w}{ }\PYG{n}{oxrna2}\PYG{o}{/}\PYG{n}{hbond}\PYG{+w}{   }\PYG{n}{seqdep}\PYG{+w}{ }\PYG{l+m}{0.870439}\PYG{+w}{ }\PYG{l+m}{8.0}\PYG{+w}{ }\PYG{l+m}{0.4}\PYG{+w}{ }\PYG{l+m}{0.75}\PYG{+w}{ }\PYG{l+m}{0.34}\PYG{+w}{ }\PYG{l+m}{0.7}\PYG{+w}{ }\PYG{l+m}{1.5}\PYG{+w}{ }\PYG{l+m}{0}\PYG{+w}{ }\PYG{l+m}{0.7}\PYG{+w}{ }\PYG{l+m}{1.5}\PYG{+w}{ }\PYG{l+m}{0}\PYG{+w}{ }\PYG{l+m}{0.7}\PYG{+w}{ }\PYG{l+m}{1.5}\PYG{+w}{ }\PYG{l+m}{0}\PYG{+w}{ }\PYG{l+m}{0.7}\PYG{+w}{ }\PYG{l+m}{0.46}\PYG{+w}{ }\PYG{l+m}{3.141592653589793}\PYG{+w}{ }\PYG{l+m}{0.7}\PYG{+w}{ }\PYG{l+m}{4.0}\PYG{+w}{ }\PYG{l+m}{1.5707963267948966}\PYG{+w}{ }\PYG{l+m}{0.45}\PYG{+w}{ }\PYG{l+m}{4.0}\PYG{+w}{ }\PYG{l+m}{1.5707963267948966}\PYG{+w}{ }\PYG{l+m}{0.45}
\PYG{k}{pair\PYGZus{}coeff}\PYG{+w}{ }\PYG{l+m}{2}\PYG{+w}{ }\PYG{l+m}{3}\PYG{+w}{ }\PYG{n}{oxrna2}\PYG{o}{/}\PYG{n}{hbond}\PYG{+w}{   }\PYG{n}{seqdep}\PYG{+w}{ }\PYG{l+m}{0.870439}\PYG{+w}{ }\PYG{l+m}{8.0}\PYG{+w}{ }\PYG{l+m}{0.4}\PYG{+w}{ }\PYG{l+m}{0.75}\PYG{+w}{ }\PYG{l+m}{0.34}\PYG{+w}{ }\PYG{l+m}{0.7}\PYG{+w}{ }\PYG{l+m}{1.5}\PYG{+w}{ }\PYG{l+m}{0}\PYG{+w}{ }\PYG{l+m}{0.7}\PYG{+w}{ }\PYG{l+m}{1.5}\PYG{+w}{ }\PYG{l+m}{0}\PYG{+w}{ }\PYG{l+m}{0.7}\PYG{+w}{ }\PYG{l+m}{1.5}\PYG{+w}{ }\PYG{l+m}{0}\PYG{+w}{ }\PYG{l+m}{0.7}\PYG{+w}{ }\PYG{l+m}{0.46}\PYG{+w}{ }\PYG{l+m}{3.141592653589793}\PYG{+w}{ }\PYG{l+m}{0.7}\PYG{+w}{ }\PYG{l+m}{4.0}\PYG{+w}{ }\PYG{l+m}{1.5707963267948966}\PYG{+w}{ }\PYG{l+m}{0.45}\PYG{+w}{ }\PYG{l+m}{4.0}\PYG{+w}{ }\PYG{l+m}{1.5707963267948966}\PYG{+w}{ }\PYG{l+m}{0.45}
\PYG{k}{pair\PYGZus{}coeff}\PYG{+w}{ }\PYG{l+m}{3}\PYG{+w}{ }\PYG{l+m}{4}\PYG{+w}{ }\PYG{n}{oxrna2}\PYG{o}{/}\PYG{n}{hbond}\PYG{+w}{   }\PYG{n}{seqdep}\PYG{+w}{ }\PYG{l+m}{0.870439}\PYG{+w}{ }\PYG{l+m}{8.0}\PYG{+w}{ }\PYG{l+m}{0.4}\PYG{+w}{ }\PYG{l+m}{0.75}\PYG{+w}{ }\PYG{l+m}{0.34}\PYG{+w}{ }\PYG{l+m}{0.7}\PYG{+w}{ }\PYG{l+m}{1.5}\PYG{+w}{ }\PYG{l+m}{0}\PYG{+w}{ }\PYG{l+m}{0.7}\PYG{+w}{ }\PYG{l+m}{1.5}\PYG{+w}{ }\PYG{l+m}{0}\PYG{+w}{ }\PYG{l+m}{0.7}\PYG{+w}{ }\PYG{l+m}{1.5}\PYG{+w}{ }\PYG{l+m}{0}\PYG{+w}{ }\PYG{l+m}{0.7}\PYG{+w}{ }\PYG{l+m}{0.46}\PYG{+w}{ }\PYG{l+m}{3.141592653589793}\PYG{+w}{ }\PYG{l+m}{0.7}\PYG{+w}{ }\PYG{l+m}{4.0}\PYG{+w}{ }\PYG{l+m}{1.5707963267948966}\PYG{+w}{ }\PYG{l+m}{0.45}\PYG{+w}{ }\PYG{l+m}{4.0}\PYG{+w}{ }\PYG{l+m}{1.5707963267948966}\PYG{+w}{ }\PYG{l+m}{0.45}
\PYG{k}{pair\PYGZus{}coeff}\PYG{+w}{ }\PYG{o}{*}\PYG{+w}{ }\PYG{o}{*}\PYG{+w}{ }\PYG{n}{oxrna2}\PYG{o}{/}\PYG{n}{xstk}\PYG{+w}{    }\PYG{l+m}{59.9626}\PYG{+w}{ }\PYG{l+m}{0.5}\PYG{+w}{ }\PYG{l+m}{0.6}\PYG{+w}{ }\PYG{l+m}{0.42}\PYG{+w}{ }\PYG{l+m}{0.58}\PYG{+w}{ }\PYG{l+m}{2.25}\PYG{+w}{ }\PYG{l+m}{0.505}\PYG{+w}{ }\PYG{l+m}{0.58}\PYG{+w}{ }\PYG{l+m}{1.7}\PYG{+w}{ }\PYG{l+m}{1.266}\PYG{+w}{ }\PYG{l+m}{0.68}\PYG{+w}{ }\PYG{l+m}{1.7}\PYG{+w}{ }\PYG{l+m}{1.266}\PYG{+w}{ }\PYG{l+m}{0.68}\PYG{+w}{ }\PYG{l+m}{1.7}\PYG{+w}{ }\PYG{l+m}{0.309}\PYG{+w}{ }\PYG{l+m}{0.68}\PYG{+w}{ }\PYG{l+m}{1.7}\PYG{+w}{ }\PYG{l+m}{0.309}\PYG{+w}{ }\PYG{l+m}{0.68}
\PYG{k}{pair\PYGZus{}coeff}\PYG{+w}{ }\PYG{o}{*}\PYG{+w}{ }\PYG{o}{*}\PYG{+w}{ }\PYG{n}{oxrna2}\PYG{o}{/}\PYG{n}{coaxstk}\PYG{+w}{ }\PYG{l+m}{80}\PYG{+w}{ }\PYG{l+m}{0.5}\PYG{+w}{ }\PYG{l+m}{0.6}\PYG{+w}{ }\PYG{l+m}{0.42}\PYG{+w}{ }\PYG{l+m}{0.58}\PYG{+w}{ }\PYG{l+m}{2.0}\PYG{+w}{ }\PYG{l+m}{2.592}\PYG{+w}{ }\PYG{l+m}{0.65}\PYG{+w}{ }\PYG{l+m}{1.3}\PYG{+w}{ }\PYG{l+m}{0.151}\PYG{+w}{ }\PYG{l+m}{0.8}\PYG{+w}{ }\PYG{l+m}{0.9}\PYG{+w}{ }\PYG{l+m}{0.685}\PYG{+w}{ }\PYG{l+m}{0.95}\PYG{+w}{ }\PYG{l+m}{0.9}\PYG{+w}{ }\PYG{l+m}{0.685}\PYG{+w}{ }\PYG{l+m}{0.95}\PYG{+w}{ }\PYG{l+m}{2.0}\PYG{+w}{ }\PYG{o}{\PYGZhy{}}\PYG{l+m}{0.65}\PYG{+w}{ }\PYG{l+m}{2.0}\PYG{+w}{ }\PYG{o}{\PYGZhy{}}\PYG{l+m}{0.65}
\PYG{k}{pair\PYGZus{}coeff}\PYG{+w}{ }\PYG{o}{*}\PYG{+w}{ }\PYG{o}{*}\PYG{+w}{ }\PYG{n}{oxrna2}\PYG{o}{/}\PYG{n}{dh}\PYG{+w}{      }\PYG{l+m}{0.1}\PYG{+w}{ }\PYG{l+m}{0.5}\PYG{+w}{ }\PYG{l+m}{1.02455}

\PYG{k}{pair\PYGZus{}style}\PYG{+w}{ }\PYG{n}{hybrid}\PYG{o}{/}\PYG{n}{overlay}\PYG{+w}{ }\PYG{n}{oxrna2}\PYG{o}{/}\PYG{n}{excv}\PYG{+w}{ }\PYG{n}{oxrna2}\PYG{o}{/}\PYG{n}{stk}\PYG{+w}{ }\PYG{n}{oxrna2}\PYG{o}{/}\PYG{n}{hbond}\PYG{+w}{ }\PYG{n}{oxrna2}\PYG{o}{/}\PYG{n}{xstk}\PYG{+w}{ }\PYG{n}{oxrna2}\PYG{o}{/}\PYG{n}{coaxstk}\PYG{+w}{ }\PYG{n}{oxrna2}\PYG{o}{/}\PYG{n}{dh}
\PYG{k}{pair\PYGZus{}coeff}\PYG{+w}{ }\PYG{o}{*}\PYG{+w}{ }\PYG{o}{*}\PYG{+w}{ }\PYG{n}{oxrna2}\PYG{o}{/}\PYG{n}{excv}\PYG{+w}{    }\PYG{n}{oxrna2\PYGZus{}lj.cgdna}
\PYG{k}{pair\PYGZus{}coeff}\PYG{+w}{ }\PYG{o}{*}\PYG{+w}{ }\PYG{o}{*}\PYG{+w}{ }\PYG{n}{oxrna2}\PYG{o}{/}\PYG{n}{stk}\PYG{+w}{     }\PYG{n}{seqdep}\PYG{+w}{ }\PYG{l+m}{0.1}\PYG{+w}{ }\PYG{l+m}{1.40206}\PYG{+w}{ }\PYG{l+m}{2.77}\PYG{+w}{ }\PYG{n}{oxrna2\PYGZus{}lj.cgdna}
\PYG{k}{pair\PYGZus{}coeff}\PYG{+w}{ }\PYG{o}{*}\PYG{+w}{ }\PYG{o}{*}\PYG{+w}{ }\PYG{n}{oxrna2}\PYG{o}{/}\PYG{n}{hbond}\PYG{+w}{   }\PYG{n}{seqdep}\PYG{+w}{ }\PYG{n}{oxrna2\PYGZus{}lj.cgdna}
\PYG{k}{pair\PYGZus{}coeff}\PYG{+w}{ }\PYG{l+m}{1}\PYG{+w}{ }\PYG{l+m}{4}\PYG{+w}{ }\PYG{n}{oxrna2}\PYG{o}{/}\PYG{n}{hbond}\PYG{+w}{   }\PYG{n}{seqdep}\PYG{+w}{ }\PYG{n}{oxrna2\PYGZus{}lj.cgdna}
\PYG{k}{pair\PYGZus{}coeff}\PYG{+w}{ }\PYG{l+m}{2}\PYG{+w}{ }\PYG{l+m}{3}\PYG{+w}{ }\PYG{n}{oxrna2}\PYG{o}{/}\PYG{n}{hbond}\PYG{+w}{   }\PYG{n}{seqdep}\PYG{+w}{ }\PYG{n}{oxrna2\PYGZus{}lj.cgdna}
\PYG{k}{pair\PYGZus{}coeff}\PYG{+w}{ }\PYG{l+m}{3}\PYG{+w}{ }\PYG{l+m}{4}\PYG{+w}{ }\PYG{n}{oxrna2}\PYG{o}{/}\PYG{n}{hbond}\PYG{+w}{   }\PYG{n}{seqdep}\PYG{+w}{ }\PYG{n}{oxrna2\PYGZus{}lj.cgdna}
\PYG{k}{pair\PYGZus{}coeff}\PYG{+w}{ }\PYG{o}{*}\PYG{+w}{ }\PYG{o}{*}\PYG{+w}{ }\PYG{n}{oxrna2}\PYG{o}{/}\PYG{n}{xstk}\PYG{+w}{    }\PYG{n}{oxrna2\PYGZus{}lj.cgdna}
\PYG{k}{pair\PYGZus{}coeff}\PYG{+w}{ }\PYG{o}{*}\PYG{+w}{ }\PYG{o}{*}\PYG{+w}{ }\PYG{n}{oxrna2}\PYG{o}{/}\PYG{n}{coaxstk}\PYG{+w}{ }\PYG{n}{oxrna2\PYGZus{}lj.cgdna}
\PYG{k}{pair\PYGZus{}coeff}\PYG{+w}{ }\PYG{o}{*}\PYG{+w}{ }\PYG{o}{*}\PYG{+w}{ }\PYG{n}{oxrna2}\PYG{o}{/}\PYG{n}{dh}\PYG{+w}{      }\PYG{l+m}{0.1}\PYG{+w}{ }\PYG{l+m}{0.5}\PYG{+w}{ }\PYG{n}{oxrna2\PYGZus{}lj.cgdna}

\PYG{c}{\PYGZsh{} Real units}
\PYG{k}{pair\PYGZus{}style}\PYG{+w}{ }\PYG{n}{hybrid}\PYG{o}{/}\PYG{n}{overlay}\PYG{+w}{ }\PYG{n}{oxrna2}\PYG{o}{/}\PYG{n}{excv}\PYG{+w}{ }\PYG{n}{oxrna2}\PYG{o}{/}\PYG{n}{stk}\PYG{+w}{ }\PYG{n}{oxrna2}\PYG{o}{/}\PYG{n}{hbond}\PYG{+w}{ }\PYG{n}{oxrna2}\PYG{o}{/}\PYG{n}{xstk}\PYG{+w}{ }\PYG{n}{oxrna2}\PYG{o}{/}\PYG{n}{coaxstk}\PYG{+w}{ }\PYG{n}{oxrna2}\PYG{o}{/}\PYG{n}{dh}
\PYG{k}{pair\PYGZus{}coeff}\PYG{+w}{ }\PYG{o}{*}\PYG{+w}{ }\PYG{o}{*}\PYG{+w}{ }\PYG{n}{oxrna2}\PYG{o}{/}\PYG{n}{excv}\PYG{+w}{    }\PYG{l+m}{11.92337812042065}\PYG{+w}{ }\PYG{l+m}{5.9626}\PYG{+w}{ }\PYG{l+m}{5.74965}\PYG{+w}{ }\PYG{l+m}{11.92337812042065}\PYG{+w}{ }\PYG{l+m}{4.38677}\PYG{+w}{ }\PYG{l+m}{4.259}\PYG{+w}{ }\PYG{l+m}{11.92337812042065}\PYG{+w}{ }\PYG{l+m}{2.81094}\PYG{+w}{ }\PYG{l+m}{2.72576}
\PYG{k}{pair\PYGZus{}coeff}\PYG{+w}{ }\PYG{o}{*}\PYG{+w}{ }\PYG{o}{*}\PYG{+w}{ }\PYG{n}{oxrna2}\PYG{o}{/}\PYG{n}{stk}\PYG{+w}{     }\PYG{n}{seqdep}\PYG{+w}{ }\PYG{l+m}{300.0}\PYG{+w}{ }\PYG{l+m}{8.35864576375849}\PYG{+w}{ }\PYG{l+m}{0.005504556}\PYG{+w}{ }\PYG{l+m}{0.70439070204273}\PYG{+w}{ }\PYG{l+m}{3.66274}\PYG{+w}{ }\PYG{l+m}{7.92174}\PYG{+w}{ }\PYG{l+m}{2.9813}\PYG{+w}{ }\PYG{l+m}{6.64404}\PYG{+w}{ }\PYG{l+m}{0.9}\PYG{+w}{ }\PYG{l+m}{0.0}\PYG{+w}{ }\PYG{l+m}{0.95}\PYG{+w}{ }\PYG{l+m}{0.9}\PYG{+w}{ }\PYG{l+m}{0.0}\PYG{+w}{ }\PYG{l+m}{0.95}\PYG{+w}{ }\PYG{l+m}{1.3}\PYG{+w}{ }\PYG{l+m}{0.0}\PYG{+w}{ }\PYG{l+m}{0.8}\PYG{+w}{ }\PYG{l+m}{1.3}\PYG{+w}{ }\PYG{l+m}{0.0}\PYG{+w}{ }\PYG{l+m}{0.8}\PYG{+w}{ }\PYG{l+m}{2.0}\PYG{+w}{ }\PYG{l+m}{0.65}\PYG{+w}{ }\PYG{l+m}{2.0}\PYG{+w}{ }\PYG{l+m}{0.65}
\PYG{k}{pair\PYGZus{}coeff}\PYG{+w}{ }\PYG{o}{*}\PYG{+w}{ }\PYG{o}{*}\PYG{+w}{ }\PYG{n}{oxrna2}\PYG{o}{/}\PYG{n}{hbond}\PYG{+w}{   }\PYG{n}{seqdep}\PYG{+w}{ }\PYG{l+m}{0.0}\PYG{+w}{ }\PYG{l+m}{0.93918760272364}\PYG{+w}{ }\PYG{l+m}{3.4072}\PYG{+w}{ }\PYG{l+m}{6.3885}\PYG{+w}{ }\PYG{l+m}{2.89612}\PYG{+w}{ }\PYG{l+m}{5.9626}\PYG{+w}{ }\PYG{l+m}{1.5}\PYG{+w}{ }\PYG{l+m}{0.0}\PYG{+w}{ }\PYG{l+m}{0.7}\PYG{+w}{ }\PYG{l+m}{1.5}\PYG{+w}{ }\PYG{l+m}{0.0}\PYG{+w}{ }\PYG{l+m}{0.7}\PYG{+w}{ }\PYG{l+m}{1.5}\PYG{+w}{ }\PYG{l+m}{0.0}\PYG{+w}{ }\PYG{l+m}{0.7}\PYG{+w}{ }\PYG{l+m}{0.46}\PYG{+w}{ }\PYG{l+m}{3.141592653589793}\PYG{+w}{ }\PYG{l+m}{0.7}\PYG{+w}{ }\PYG{l+m}{4.0}\PYG{+w}{ }\PYG{l+m}{1.5707963267948966}\PYG{+w}{ }\PYG{l+m}{0.45}\PYG{+w}{ }\PYG{l+m}{4.0}\PYG{+w}{ }\PYG{l+m}{1.5707963267948966}\PYG{+w}{ }\PYG{l+m}{0.45}
\PYG{k}{pair\PYGZus{}coeff}\PYG{+w}{ }\PYG{l+m}{1}\PYG{+w}{ }\PYG{l+m}{4}\PYG{+w}{ }\PYG{n}{oxrna2}\PYG{o}{/}\PYG{n}{hbond}\PYG{+w}{   }\PYG{n}{seqdep}\PYG{+w}{ }\PYG{l+m}{5.18928666388042}\PYG{+w}{ }\PYG{l+m}{0.93918760272364}\PYG{+w}{ }\PYG{l+m}{3.4072}\PYG{+w}{ }\PYG{l+m}{6.3885}\PYG{+w}{ }\PYG{l+m}{2.89612}\PYG{+w}{ }\PYG{l+m}{5.9626}\PYG{+w}{ }\PYG{l+m}{1.5}\PYG{+w}{ }\PYG{l+m}{0.0}\PYG{+w}{ }\PYG{l+m}{0.7}\PYG{+w}{ }\PYG{l+m}{1.5}\PYG{+w}{ }\PYG{l+m}{0.0}\PYG{+w}{ }\PYG{l+m}{0.7}\PYG{+w}{ }\PYG{l+m}{1.5}\PYG{+w}{ }\PYG{l+m}{0.0}\PYG{+w}{ }\PYG{l+m}{0.7}\PYG{+w}{ }\PYG{l+m}{0.46}\PYG{+w}{ }\PYG{l+m}{3.141592653589793}\PYG{+w}{ }\PYG{l+m}{0.7}\PYG{+w}{ }\PYG{l+m}{4.0}\PYG{+w}{ }\PYG{l+m}{1.5707963267948966}\PYG{+w}{ }\PYG{l+m}{0.45}\PYG{+w}{ }\PYG{l+m}{4.0}\PYG{+w}{ }\PYG{l+m}{1.5707963267948966}\PYG{+w}{ }\PYG{l+m}{0.45}
\PYG{k}{pair\PYGZus{}coeff}\PYG{+w}{ }\PYG{l+m}{2}\PYG{+w}{ }\PYG{l+m}{3}\PYG{+w}{ }\PYG{n}{oxrna2}\PYG{o}{/}\PYG{n}{hbond}\PYG{+w}{   }\PYG{n}{seqdep}\PYG{+w}{ }\PYG{l+m}{5.18928666388042}\PYG{+w}{ }\PYG{l+m}{0.93918760272364}\PYG{+w}{ }\PYG{l+m}{3.4072}\PYG{+w}{ }\PYG{l+m}{6.3885}\PYG{+w}{ }\PYG{l+m}{2.89612}\PYG{+w}{ }\PYG{l+m}{5.9626}\PYG{+w}{ }\PYG{l+m}{1.5}\PYG{+w}{ }\PYG{l+m}{0.0}\PYG{+w}{ }\PYG{l+m}{0.7}\PYG{+w}{ }\PYG{l+m}{1.5}\PYG{+w}{ }\PYG{l+m}{0.0}\PYG{+w}{ }\PYG{l+m}{0.7}\PYG{+w}{ }\PYG{l+m}{1.5}\PYG{+w}{ }\PYG{l+m}{0.0}\PYG{+w}{ }\PYG{l+m}{0.7}\PYG{+w}{ }\PYG{l+m}{0.46}\PYG{+w}{ }\PYG{l+m}{3.141592653589793}\PYG{+w}{ }\PYG{l+m}{0.7}\PYG{+w}{ }\PYG{l+m}{4.0}\PYG{+w}{ }\PYG{l+m}{1.5707963267948966}\PYG{+w}{ }\PYG{l+m}{0.45}\PYG{+w}{ }\PYG{l+m}{4.0}\PYG{+w}{ }\PYG{l+m}{1.5707963267948966}\PYG{+w}{ }\PYG{l+m}{0.45}
\PYG{k}{pair\PYGZus{}coeff}\PYG{+w}{ }\PYG{l+m}{3}\PYG{+w}{ }\PYG{l+m}{4}\PYG{+w}{ }\PYG{n}{oxrna2}\PYG{o}{/}\PYG{n}{hbond}\PYG{+w}{   }\PYG{n}{seqdep}\PYG{+w}{ }\PYG{l+m}{5.18928666388042}\PYG{+w}{ }\PYG{l+m}{0.93918760272364}\PYG{+w}{ }\PYG{l+m}{3.4072}\PYG{+w}{ }\PYG{l+m}{6.3885}\PYG{+w}{ }\PYG{l+m}{2.89612}\PYG{+w}{ }\PYG{l+m}{5.9626}\PYG{+w}{ }\PYG{l+m}{1.5}\PYG{+w}{ }\PYG{l+m}{0.0}\PYG{+w}{ }\PYG{l+m}{0.7}\PYG{+w}{ }\PYG{l+m}{1.5}\PYG{+w}{ }\PYG{l+m}{0.0}\PYG{+w}{ }\PYG{l+m}{0.7}\PYG{+w}{ }\PYG{l+m}{1.5}\PYG{+w}{ }\PYG{l+m}{0.0}\PYG{+w}{ }\PYG{l+m}{0.7}\PYG{+w}{ }\PYG{l+m}{0.46}\PYG{+w}{ }\PYG{l+m}{3.141592653589793}\PYG{+w}{ }\PYG{l+m}{0.7}\PYG{+w}{ }\PYG{l+m}{4.0}\PYG{+w}{ }\PYG{l+m}{1.5707963267948966}\PYG{+w}{ }\PYG{l+m}{0.45}\PYG{+w}{ }\PYG{l+m}{4.0}\PYG{+w}{ }\PYG{l+m}{1.5707963267948966}\PYG{+w}{ }\PYG{l+m}{0.45}
\PYG{k}{pair\PYGZus{}coeff}\PYG{+w}{ }\PYG{o}{*}\PYG{+w}{ }\PYG{o}{*}\PYG{+w}{ }\PYG{n}{oxrna2}\PYG{o}{/}\PYG{n}{xstk}\PYG{+w}{    }\PYG{l+m}{4.92690859644113}\PYG{+w}{ }\PYG{l+m}{4.259}\PYG{+w}{ }\PYG{l+m}{5.1108}\PYG{+w}{ }\PYG{l+m}{3.57756}\PYG{+w}{ }\PYG{l+m}{4.94044}\PYG{+w}{ }\PYG{l+m}{2.25}\PYG{+w}{ }\PYG{l+m}{0.505}\PYG{+w}{ }\PYG{l+m}{0.58}\PYG{+w}{ }\PYG{l+m}{1.7}\PYG{+w}{ }\PYG{l+m}{1.266}\PYG{+w}{ }\PYG{l+m}{0.68}\PYG{+w}{ }\PYG{l+m}{1.7}\PYG{+w}{ }\PYG{l+m}{1.266}\PYG{+w}{ }\PYG{l+m}{0.68}\PYG{+w}{ }\PYG{l+m}{1.7}\PYG{+w}{ }\PYG{l+m}{0.309}\PYG{+w}{ }\PYG{l+m}{0.68}\PYG{+w}{ }\PYG{l+m}{1.7}\PYG{+w}{ }\PYG{l+m}{0.309}\PYG{+w}{ }\PYG{l+m}{0.68}
\PYG{k}{pair\PYGZus{}coeff}\PYG{+w}{ }\PYG{o}{*}\PYG{+w}{ }\PYG{o}{*}\PYG{+w}{ }\PYG{n}{oxrna2}\PYG{o}{/}\PYG{n}{coaxstk}\PYG{+w}{ }\PYG{l+m}{6.57330882442206}\PYG{+w}{ }\PYG{l+m}{4.259}\PYG{+w}{ }\PYG{l+m}{5.1108}\PYG{+w}{ }\PYG{l+m}{3.57756}\PYG{+w}{ }\PYG{l+m}{4.94044}\PYG{+w}{ }\PYG{l+m}{2.0}\PYG{+w}{ }\PYG{l+m}{2.592}\PYG{+w}{ }\PYG{l+m}{0.65}\PYG{+w}{ }\PYG{l+m}{1.3}\PYG{+w}{ }\PYG{l+m}{0.151}\PYG{+w}{ }\PYG{l+m}{0.8}\PYG{+w}{ }\PYG{l+m}{0.9}\PYG{+w}{ }\PYG{l+m}{0.685}\PYG{+w}{ }\PYG{l+m}{0.95}\PYG{+w}{ }\PYG{l+m}{0.9}\PYG{+w}{ }\PYG{l+m}{0.685}\PYG{+w}{ }\PYG{l+m}{0.95}\PYG{+w}{ }\PYG{l+m}{2.0}\PYG{+w}{ }\PYG{o}{\PYGZhy{}}\PYG{l+m}{0.65}\PYG{+w}{ }\PYG{l+m}{2.0}\PYG{+w}{ }\PYG{o}{\PYGZhy{}}\PYG{l+m}{0.65}
\PYG{k}{pair\PYGZus{}coeff}\PYG{+w}{ }\PYG{o}{*}\PYG{+w}{ }\PYG{o}{*}\PYG{+w}{ }\PYG{n}{oxrna2}\PYG{o}{/}\PYG{n}{dh}\PYG{+w}{      }\PYG{l+m}{300.0}\PYG{+w}{ }\PYG{l+m}{0.5}\PYG{+w}{ }\PYG{l+m}{1.02455}

\PYG{k}{pair\PYGZus{}style}\PYG{+w}{ }\PYG{n}{hybrid}\PYG{o}{/}\PYG{n}{overlay}\PYG{+w}{ }\PYG{n}{oxrna2}\PYG{o}{/}\PYG{n}{excv}\PYG{+w}{ }\PYG{n}{oxrna2}\PYG{o}{/}\PYG{n}{stk}\PYG{+w}{ }\PYG{n}{oxrna2}\PYG{o}{/}\PYG{n}{hbond}\PYG{+w}{ }\PYG{n}{oxrna2}\PYG{o}{/}\PYG{n}{xstk}\PYG{+w}{ }\PYG{n}{oxrna2}\PYG{o}{/}\PYG{n}{coaxstk}\PYG{+w}{ }\PYG{n}{oxrna2}\PYG{o}{/}\PYG{n}{dh}
\PYG{k}{pair\PYGZus{}coeff}\PYG{+w}{ }\PYG{o}{*}\PYG{+w}{ }\PYG{o}{*}\PYG{+w}{ }\PYG{n}{oxrna2}\PYG{o}{/}\PYG{n}{excv}\PYG{+w}{    }\PYG{n}{oxrna2\PYGZus{}real.cgdna}
\PYG{k}{pair\PYGZus{}coeff}\PYG{+w}{ }\PYG{o}{*}\PYG{+w}{ }\PYG{o}{*}\PYG{+w}{ }\PYG{n}{oxrna2}\PYG{o}{/}\PYG{n}{stk}\PYG{+w}{     }\PYG{n}{seqdep}\PYG{+w}{ }\PYG{l+m}{300.0}\PYG{+w}{ }\PYG{l+m}{8.35864576375849}\PYG{+w}{ }\PYG{l+m}{0.005504556}\PYG{+w}{ }\PYG{n}{oxrna2\PYGZus{}real.cgdna}
\PYG{k}{pair\PYGZus{}coeff}\PYG{+w}{ }\PYG{o}{*}\PYG{+w}{ }\PYG{o}{*}\PYG{+w}{ }\PYG{n}{oxrna2}\PYG{o}{/}\PYG{n}{hbond}\PYG{+w}{   }\PYG{n}{seqdep}\PYG{+w}{ }\PYG{n}{oxrna2\PYGZus{}real.cgdna}
\PYG{k}{pair\PYGZus{}coeff}\PYG{+w}{ }\PYG{l+m}{1}\PYG{+w}{ }\PYG{l+m}{4}\PYG{+w}{ }\PYG{n}{oxrna2}\PYG{o}{/}\PYG{n}{hbond}\PYG{+w}{   }\PYG{n}{seqdep}\PYG{+w}{ }\PYG{n}{oxrna2\PYGZus{}real.cgdna}
\PYG{k}{pair\PYGZus{}coeff}\PYG{+w}{ }\PYG{l+m}{2}\PYG{+w}{ }\PYG{l+m}{3}\PYG{+w}{ }\PYG{n}{oxrna2}\PYG{o}{/}\PYG{n}{hbond}\PYG{+w}{   }\PYG{n}{seqdep}\PYG{+w}{ }\PYG{n}{oxrna2\PYGZus{}real.cgdna}
\PYG{k}{pair\PYGZus{}coeff}\PYG{+w}{ }\PYG{l+m}{3}\PYG{+w}{ }\PYG{l+m}{4}\PYG{+w}{ }\PYG{n}{oxrna2}\PYG{o}{/}\PYG{n}{hbond}\PYG{+w}{   }\PYG{n}{seqdep}\PYG{+w}{ }\PYG{n}{oxrna2\PYGZus{}real.cgdna}
\PYG{k}{pair\PYGZus{}coeff}\PYG{+w}{ }\PYG{o}{*}\PYG{+w}{ }\PYG{o}{*}\PYG{+w}{ }\PYG{n}{oxrna2}\PYG{o}{/}\PYG{n}{xstk}\PYG{+w}{    }\PYG{n}{oxrna2\PYGZus{}real.cgdna}
\PYG{k}{pair\PYGZus{}coeff}\PYG{+w}{ }\PYG{o}{*}\PYG{+w}{ }\PYG{o}{*}\PYG{+w}{ }\PYG{n}{oxrna2}\PYG{o}{/}\PYG{n}{coaxstk}\PYG{+w}{ }\PYG{n}{oxrna2\PYGZus{}real.cgdna}
\PYG{k}{pair\PYGZus{}coeff}\PYG{+w}{ }\PYG{o}{*}\PYG{+w}{ }\PYG{o}{*}\PYG{+w}{ }\PYG{n}{oxrna2}\PYG{o}{/}\PYG{n}{dh}\PYG{+w}{      }\PYG{l+m}{300.0}\PYG{+w}{ }\PYG{l+m}{0.5}\PYG{+w}{ }\PYG{n}{oxrna2\PYGZus{}real.cgdna}
\end{sphinxVerbatim}

\begin{sphinxadmonition}{note}{Note:}
\sphinxAtStartPar
The coefficients in the above examples are provided in forms
compatible with both \sphinxstyleemphasis{units lj} and \sphinxstyleemphasis{units real} (see documentation
of {\hyperref[\detokenize{units::doc}]{\sphinxcrossref{\DUrole{doc}{units}}}}).  These can also be read from a potential
file with correct unit style by specifying the name of the
file. Several potential files for each unit style are included in the
\sphinxcode{\sphinxupquote{potentials}} directory of the LAMMPS distribution.
\end{sphinxadmonition}


\subsection{Description}
\label{\detokenize{pair_oxrna2:description}}
\sphinxAtStartPar
The \sphinxstyleemphasis{oxrna2} pair styles compute the pairwise\sphinxhyphen{}additive parts of the
oxDNA force field for coarse\sphinxhyphen{}grained modelling of RNA. The effective
interaction between the nucleotides consists of potentials for the
excluded volume interaction \sphinxstyleemphasis{oxrna2/excv}, the stacking \sphinxstyleemphasis{oxrna2/stk},
cross\sphinxhyphen{}stacking \sphinxstyleemphasis{oxrna2/xstk} and coaxial stacking interaction
\sphinxstyleemphasis{oxrna2/coaxstk}, electrostatic Debye\sphinxhyphen{}Hueckel interaction \sphinxstyleemphasis{oxrna2/dh} as
well as the hydrogen\sphinxhyphen{}bonding interaction \sphinxstyleemphasis{oxrna2/hbond} between
complementary pairs of nucleotides on opposite strands. Average sequence
or sequence\sphinxhyphen{}dependent stacking and base\sphinxhyphen{}pairing strengths are supported
{\hyperref[\detokenize{pair_oxrna2:sulc32}]{\sphinxcrossref{\DUrole{std,std-ref}{(Sulc2)}}}}. Quasi\sphinxhyphen{}unique base\sphinxhyphen{}pairing between nucleotides
can be achieved by using more complementary pairs of atom types like 5\sphinxhyphen{}8
and 6\sphinxhyphen{}7, 9\sphinxhyphen{}12 and 10\sphinxhyphen{}11, 13\sphinxhyphen{}16 and 14\sphinxhyphen{}15, etc.  This prevents the
hybridization of in principle complementary bases within Ntypes/4 bases
up and down along the backbone.

\sphinxAtStartPar
The exact functional form of the pair styles is rather complex.  The
individual potentials consist of products of modulation factors, which
themselves are constructed from a number of more basic potentials
(Morse, Lennard\sphinxhyphen{}Jones, harmonic angle and distance) as well as quadratic
smoothing and modulation terms.  We refer to {\hyperref[\detokenize{pair_oxrna2:sulc31}]{\sphinxcrossref{\DUrole{std,std-ref}{(Sulc1)}}}} and
the original oxDNA publications {\hyperref[\detokenize{pair_oxrna2:ouldridge-dphil3}]{\sphinxcrossref{\DUrole{std,std-ref}{(Ouldridge\sphinxhyphen{}DPhil)}}}} and {\hyperref[\detokenize{pair_oxrna2:ouldridge3}]{\sphinxcrossref{\DUrole{std,std-ref}{(Ouldridge)}}}} for a detailed
description of the oxRNA2 force field.

\begin{sphinxadmonition}{note}{Note:}
\sphinxAtStartPar
These pair styles have to be used together with the related oxDNA2
bond style \sphinxstyleemphasis{oxrna2/fene} for the connectivity of the phosphate
backbone (see also documentation of {\hyperref[\detokenize{bond_oxdna::doc}]{\sphinxcrossref{\DUrole{doc}{bond\_style oxrna2/fene}}}}). Most of the coefficients in the above example have to
be kept fixed and cannot be changed without reparameterizing the
entire model.  Exceptions are the first four coefficients after
\sphinxstyleemphasis{oxrna2/stk} (seq=seqdep, T=0.1, xi=1.40206 and kappa=2.77 and
corresponding \sphinxstyleemphasis{real unit} equivalents in the above examples), the
first coefficient after \sphinxstyleemphasis{oxrna2/hbond} (seq=seqdep in the above
example) and the three coefficients after \sphinxstyleemphasis{oxrna2/dh} (T=0.1,
rhos=0.5, qeff=1.02455 in the above example). When using a Langevin
thermostat e.g. through {\hyperref[\detokenize{fix_langevin::doc}]{\sphinxcrossref{\DUrole{doc}{fix langevin}}}} or
{\hyperref[\detokenize{fix_nve_dotc_langevin::doc}]{\sphinxcrossref{\DUrole{doc}{fix nve/dotc/langevin}}}} the temperature
coefficients have to be matched to the one used in the fix.
\end{sphinxadmonition}

\begin{sphinxadmonition}{note}{Note:}
\sphinxAtStartPar
These pair styles have to be used with the \sphinxstyleemphasis{atom\_style hybrid bond
ellipsoid oxdna} (see documentation of {\hyperref[\detokenize{atom_style::doc}]{\sphinxcrossref{\DUrole{doc}{atom\_style}}}}). The \sphinxstyleemphasis{atom\_style oxdna} stores the 3’\sphinxhyphen{}to\sphinxhyphen{}5’ polarity
of the nucleotide strand, which is set through the bond topology in
the data file. The first (second) atom in a bond definition is
understood to point towards the 3’\sphinxhyphen{}end (5’\sphinxhyphen{}end) of the strand.
\end{sphinxadmonition}

\sphinxAtStartPar
Example input and data files for DNA duplexes can be found in
\sphinxcode{\sphinxupquote{examples/PACKAGES/cgdna/examples/oxDNA/}} and \sphinxcode{\sphinxupquote{.../oxDNA2/}}.  A simple python
setup tool which creates single straight or helical DNA strands, DNA
duplexes or arrays of DNA duplexes can be found in
\sphinxcode{\sphinxupquote{examples/PACKAGES/cgdna/util/}}.

\sphinxAtStartPar
Please cite {\hyperref[\detokenize{pair_oxrna2:henrich3}]{\sphinxcrossref{\DUrole{std,std-ref}{(Henrich)}}}} in any publication that uses
this implementation.  The article contains general information on the
model, its implementation and performance as well as the structure of
the data and input file. The preprint version of the article can be
found \sphinxhref{PDF/CG-DNA.pdf}{here}.  Please cite also the relevant oxRNA2
publications {\hyperref[\detokenize{pair_oxrna2:sulc31}]{\sphinxcrossref{\DUrole{std,std-ref}{(Sulc1)}}}} and {\hyperref[\detokenize{pair_oxrna2:sulc32}]{\sphinxcrossref{\DUrole{std,std-ref}{(Sulc2)}}}}.


\bigskip\hrule\bigskip



\subsection{Potential file reading}
\label{\detokenize{pair_oxrna2:potential-file-reading}}
\sphinxAtStartPar
For each pair style above the first non\sphinxhyphen{}modifiable argument can be a
filename (with exception of Debye\sphinxhyphen{}Hueckel, for which the effective
charge argument can be a filename), and if it is, no further arguments
should be supplied.  Therefore the following command:

\begin{sphinxVerbatim}[commandchars=\\\{\}]
\PYG{k}{pair\PYGZus{}coeff}\PYG{+w}{ }\PYG{l+m}{3}\PYG{+w}{ }\PYG{l+m}{4}\PYG{+w}{ }\PYG{n}{oxrna2}\PYG{o}{/}\PYG{n}{hbond}\PYG{+w}{   }\PYG{n}{seqdep}\PYG{+w}{ }\PYG{n}{oxrna2\PYGZus{}lj.cgdna}
\end{sphinxVerbatim}

\sphinxAtStartPar
will be interpreted as a request to read the corresponding hydrogen
bonding potential parameters from the file with the given name.  The
file can define multiple potential parameters for both bonded and pair
interactions, but for the example pair interaction above there must
exist in the file a line of the form:

\begin{sphinxVerbatim}[commandchars=\\\{\}]
\PYG{l+m}{3}\PYG{+w}{ }\PYG{l+m}{4}\PYG{+w}{ }\PYG{n}{hbond}\PYG{+w}{     }\PYG{o}{\PYGZlt{}}\PYG{n}{coefficients}\PYG{o}{\PYGZgt{}}
\end{sphinxVerbatim}

\sphinxAtStartPar
If potential customization is required, the potential file reading can
be mixed with the manual specification of the potential parameters. For
example, the following command:

\begin{sphinxVerbatim}[commandchars=\\\{\}]
\PYG{k}{pair\PYGZus{}style}\PYG{+w}{ }\PYG{n}{hybrid}\PYG{o}{/}\PYG{n}{overlay}\PYG{+w}{ }\PYG{n}{oxrna2}\PYG{o}{/}\PYG{n}{excv}\PYG{+w}{ }\PYG{n}{oxrna2}\PYG{o}{/}\PYG{n}{stk}\PYG{+w}{ }\PYG{n}{oxrna2}\PYG{o}{/}\PYG{n}{hbond}\PYG{+w}{ }\PYG{n}{oxrna2}\PYG{o}{/}\PYG{n}{xstk}\PYG{+w}{ }\PYG{n}{oxrna2}\PYG{o}{/}\PYG{n}{coaxstk}\PYG{+w}{ }\PYG{n}{oxrna2}\PYG{o}{/}\PYG{n}{dh}
\PYG{k}{pair\PYGZus{}coeff}\PYG{+w}{ }\PYG{o}{*}\PYG{+w}{ }\PYG{o}{*}\PYG{+w}{ }\PYG{n}{oxrna2}\PYG{o}{/}\PYG{n}{excv}\PYG{+w}{    }\PYG{l+m}{2.0}\PYG{+w}{ }\PYG{l+m}{0.7}\PYG{+w}{ }\PYG{l+m}{0.675}\PYG{+w}{ }\PYG{l+m}{2.0}\PYG{+w}{ }\PYG{l+m}{0.515}\PYG{+w}{ }\PYG{l+m}{0.5}\PYG{+w}{ }\PYG{l+m}{2.0}\PYG{+w}{ }\PYG{l+m}{0.33}\PYG{+w}{ }\PYG{l+m}{0.32}
\PYG{k}{pair\PYGZus{}coeff}\PYG{+w}{ }\PYG{o}{*}\PYG{+w}{ }\PYG{o}{*}\PYG{+w}{ }\PYG{n}{oxrna2}\PYG{o}{/}\PYG{n}{stk}\PYG{+w}{     }\PYG{n}{seqdep}\PYG{+w}{ }\PYG{l+m}{0.1}\PYG{+w}{ }\PYG{l+m}{1.40206}\PYG{+w}{ }\PYG{l+m}{2.77}\PYG{+w}{ }\PYG{n}{oxrna2\PYGZus{}lj.cgdna}
\PYG{k}{pair\PYGZus{}coeff}\PYG{+w}{ }\PYG{o}{*}\PYG{+w}{ }\PYG{o}{*}\PYG{+w}{ }\PYG{n}{oxrna2}\PYG{o}{/}\PYG{n}{hbond}\PYG{+w}{   }\PYG{n}{seqdep}\PYG{+w}{ }\PYG{n}{oxrna2\PYGZus{}lj.cgdna}
\PYG{k}{pair\PYGZus{}coeff}\PYG{+w}{ }\PYG{l+m}{1}\PYG{+w}{ }\PYG{l+m}{4}\PYG{+w}{ }\PYG{n}{oxrna2}\PYG{o}{/}\PYG{n}{hbond}\PYG{+w}{   }\PYG{n}{seqdep}\PYG{+w}{ }\PYG{n}{oxrna2\PYGZus{}lj.cgdna}
\PYG{k}{pair\PYGZus{}coeff}\PYG{+w}{ }\PYG{l+m}{2}\PYG{+w}{ }\PYG{l+m}{3}\PYG{+w}{ }\PYG{n}{oxrna2}\PYG{o}{/}\PYG{n}{hbond}\PYG{+w}{   }\PYG{n}{seqdep}\PYG{+w}{ }\PYG{n}{oxrna2\PYGZus{}lj.cgdna}
\PYG{k}{pair\PYGZus{}coeff}\PYG{+w}{ }\PYG{l+m}{3}\PYG{+w}{ }\PYG{l+m}{4}\PYG{+w}{ }\PYG{n}{oxrna2}\PYG{o}{/}\PYG{n}{hbond}\PYG{+w}{   }\PYG{n}{seqdep}\PYG{+w}{ }\PYG{n}{oxrna2\PYGZus{}lj.cgdna}
\PYG{k}{pair\PYGZus{}coeff}\PYG{+w}{ }\PYG{o}{*}\PYG{+w}{ }\PYG{o}{*}\PYG{+w}{ }\PYG{n}{oxrna2}\PYG{o}{/}\PYG{n}{xstk}\PYG{+w}{    }\PYG{n}{oxrna2\PYGZus{}lj.cgdna}
\PYG{k}{pair\PYGZus{}coeff}\PYG{+w}{ }\PYG{o}{*}\PYG{+w}{ }\PYG{o}{*}\PYG{+w}{ }\PYG{n}{oxrna2}\PYG{o}{/}\PYG{n}{coaxstk}\PYG{+w}{ }\PYG{n}{oxrna2\PYGZus{}lj.cgdna}
\PYG{k}{pair\PYGZus{}coeff}\PYG{+w}{ }\PYG{o}{*}\PYG{+w}{ }\PYG{o}{*}\PYG{+w}{ }\PYG{n}{oxrna2}\PYG{o}{/}\PYG{n}{dh}\PYG{+w}{      }\PYG{l+m}{0.1}\PYG{+w}{ }\PYG{l+m}{0.5}\PYG{+w}{ }\PYG{l+m}{1.02455}
\end{sphinxVerbatim}

\sphinxAtStartPar
will read the excluded volume and Debye\sphinxhyphen{}Hueckel effective charge \sphinxstyleemphasis{qeff}
parameters from the manual specification and all others from the
potential file \sphinxstyleemphasis{oxrna2\_lj.cgdna}.

\sphinxAtStartPar
There are sample potential files for each unit style in the
\sphinxcode{\sphinxupquote{potentials}} directory of the LAMMPS distribution. The potential file
unit system must align with the units defined via the {\hyperref[\detokenize{units::doc}]{\sphinxcrossref{\DUrole{doc}{units}}}} command. For conversion between different \sphinxstyleemphasis{LJ} and \sphinxstyleemphasis{real} unit
systems for oxDNA, the python tool \sphinxstyleemphasis{lj2real.py} located in the
\sphinxcode{\sphinxupquote{examples/PACKAGES/cgdna/util/}} directory can be used. This tool
assumes similar file structure to the examples found in
\sphinxcode{\sphinxupquote{examples/PACKAGES/cgdna/examples/}}.


\bigskip\hrule\bigskip



\subsection{Restrictions}
\label{\detokenize{pair_oxrna2:restrictions}}
\sphinxAtStartPar
These pair styles can only be used if LAMMPS was built with the
CG\sphinxhyphen{}DNA package and the MOLECULE and ASPHERE package.  See the
{\hyperref[\detokenize{Build_package::doc}]{\sphinxcrossref{\DUrole{doc}{Build package}}}} page for more info.


\subsection{Related commands}
\label{\detokenize{pair_oxrna2:related-commands}}
\sphinxAtStartPar
{\hyperref[\detokenize{bond_oxdna::doc}]{\sphinxcrossref{\DUrole{doc}{bond\_style oxrna2/fene}}}}, {\hyperref[\detokenize{pair_coeff::doc}]{\sphinxcrossref{\DUrole{doc}{pair\_coeff}}}},
{\hyperref[\detokenize{bond_oxdna::doc}]{\sphinxcrossref{\DUrole{doc}{bond\_style oxdna/fene}}}}, {\hyperref[\detokenize{pair_oxdna::doc}]{\sphinxcrossref{\DUrole{doc}{pair\_style oxdna/excv}}}},
{\hyperref[\detokenize{bond_oxdna::doc}]{\sphinxcrossref{\DUrole{doc}{bond\_style oxdna2/fene}}}}, {\hyperref[\detokenize{pair_oxdna2::doc}]{\sphinxcrossref{\DUrole{doc}{pair\_style oxdna2/excv}}}},
{\hyperref[\detokenize{atom_style::doc}]{\sphinxcrossref{\DUrole{doc}{atom\_style oxdna}}}}, {\hyperref[\detokenize{fix_nve_dotc_langevin::doc}]{\sphinxcrossref{\DUrole{doc}{fix nve/dotc/langevin}}}}


\subsection{Default}
\label{\detokenize{pair_oxrna2:default}}
\sphinxAtStartPar
none


\bigskip\hrule\bigskip

\phantomsection\label{\detokenize{pair_oxrna2:henrich3}}
\sphinxAtStartPar
\sphinxstylestrong{(Henrich)} O. Henrich, Y. A. Gutierrez\sphinxhyphen{}Fosado, T. Curk, T. E. Ouldridge, Eur. Phys. J. E 41, 57 (2018).

\phantomsection\label{\detokenize{pair_oxrna2:sulc31}}
\sphinxAtStartPar
\sphinxstylestrong{(Sulc1)} P. Sulc, F. Romano, T. E. Ouldridge, et al., J. Chem. Phys. 140, 235102 (2014).

\phantomsection\label{\detokenize{pair_oxrna2:sulc32}}
\sphinxAtStartPar
\sphinxstylestrong{(Sulc2)} P. Sulc, F. Romano, T.E. Ouldridge, L. Rovigatti, J.P.K. Doye, A.A. Louis, J. Chem. Phys. 137, 135101 (2012).

\phantomsection\label{\detokenize{pair_oxrna2:ouldridge-dphil3}}
\sphinxAtStartPar
\sphinxstylestrong{(Ouldridge\sphinxhyphen{}DPhil)} T.E. Ouldridge, Coarse\sphinxhyphen{}grained modelling of DNA and DNA self\sphinxhyphen{}assembly, DPhil. University of Oxford (2011).

\phantomsection\label{\detokenize{pair_oxrna2:ouldridge3}}
\sphinxAtStartPar
\sphinxstylestrong{(Ouldridge)} T.E. Ouldridge, A.A. Louis, J.P.K. Doye, J. Chem. Phys. 134, 085101 (2011).

\sphinxstepscope

\index{pair\_style pace@\spxentry{pair\_style pace}}\ignorespaces 
\index{pair\_style pace/kk@\spxentry{pair\_style pace/kk}}\ignorespaces 
\index{pair\_style pace/extrapolation@\spxentry{pair\_style pace/extrapolation}}\ignorespaces 
\index{pair\_style pace/extrapolation/kk@\spxentry{pair\_style pace/extrapolation/kk}}\ignorespaces 

\section{pair\_style pace command}
\label{\detokenize{pair_pace:pair-style-pace-command}}\label{\detokenize{pair_pace:index-3}}\label{\detokenize{pair_pace:index-2}}\label{\detokenize{pair_pace:index-1}}\label{\detokenize{pair_pace:index-0}}\label{\detokenize{pair_pace::doc}}
\sphinxAtStartPar
Accelerator Variants: \sphinxstyleemphasis{pace/kk}, \sphinxstyleemphasis{pace/extrapolation/kk}


\section{pair\_style pace/extrapolation command}
\label{\detokenize{pair_pace:pair-style-pace-extrapolation-command}}

\subsection{Syntax}
\label{\detokenize{pair_pace:syntax}}
\begin{sphinxVerbatim}[commandchars=\\\{\}]
\PYG{k}{pair\PYGZus{}style}\PYG{+w}{ }\PYG{n}{pace}\PYG{+w}{ }\PYG{n}{...}\PYG{+w}{ }\PYG{n}{keyword}\PYG{+w}{ }\PYG{n}{values}\PYG{+w}{ }\PYG{n}{...}
\end{sphinxVerbatim}
\begin{itemize}
\item {} 
\sphinxAtStartPar
one or more keyword/value pairs may be appended
\begin{sphinxalltt}
keyword = \sphinxstyleemphasis{product} or \sphinxstyleemphasis{recursive} or \sphinxstyleemphasis{chunksize}
  \sphinxstyleemphasis{product} = use product algorithm for basis functions
  \sphinxstyleemphasis{recursive} = use recursive algorithm for basis functions
  \sphinxstyleemphasis{chunksize} value = number of atoms in each pass
\end{sphinxalltt}

\end{itemize}

\begin{sphinxVerbatim}[commandchars=\\\{\}]
\PYG{k}{pair\PYGZus{}style}\PYG{+w}{ }\PYG{n}{pace}\PYG{o}{/}\PYG{n}{extrapolation}
\end{sphinxVerbatim}


\subsection{Examples}
\label{\detokenize{pair_pace:examples}}
\begin{sphinxVerbatim}[commandchars=\\\{\}]
\PYG{k}{pair\PYGZus{}style}\PYG{+w}{ }\PYG{n}{pace}
\PYG{k}{pair\PYGZus{}style}\PYG{+w}{ }\PYG{n}{pace}\PYG{+w}{ }\PYG{n}{product}\PYG{+w}{ }\PYG{n}{chunksize}\PYG{+w}{ }\PYG{l+m}{2048}
\PYG{k}{pair\PYGZus{}coeff}\PYG{+w}{ }\PYG{o}{*}\PYG{+w}{ }\PYG{o}{*}\PYG{+w}{ }\PYG{n}{Cu}\PYG{o}{\PYGZhy{}}\PYG{n}{PBE}\PYG{o}{\PYGZhy{}}\PYG{n}{core}\PYG{o}{\PYGZhy{}}\PYG{n}{rep.ace}\PYG{+w}{ }\PYG{n}{Cu}

\PYG{k}{pair\PYGZus{}style}\PYG{+w}{ }\PYG{n}{pace}
\PYG{k}{pair\PYGZus{}coeff}\PYG{+w}{ }\PYG{o}{*}\PYG{+w}{ }\PYG{o}{*}\PYG{+w}{ }\PYG{n}{Cu.yaml}\PYG{+w}{ }\PYG{n}{Cu}

\PYG{k}{pair\PYGZus{}style}\PYG{+w}{ }\PYG{n}{pace}\PYG{o}{/}\PYG{n}{extrapolation}
\PYG{k}{pair\PYGZus{}coeff}\PYG{+w}{ }\PYG{o}{*}\PYG{+w}{ }\PYG{o}{*}\PYG{+w}{ }\PYG{n}{Cu.yaml}\PYG{+w}{ }\PYG{n}{Cu.asi}\PYG{+w}{ }\PYG{n}{Cu}
\end{sphinxVerbatim}


\subsection{Description}
\label{\detokenize{pair_pace:description}}
\sphinxAtStartPar
Pair style \sphinxstyleemphasis{pace} computes interactions using the Atomic Cluster
Expansion (ACE), which is a general expansion of the atomic energy in
multi\sphinxhyphen{}body basis functions. {\hyperref[\detokenize{pair_pace:drautz20191}]{\sphinxcrossref{\DUrole{std,std-ref}{(Drautz19)}}}}.  The \sphinxstyleemphasis{pace}
pair style provides an efficient implementation that is described in
this paper {\hyperref[\detokenize{pair_pace:lysogorskiy20211}]{\sphinxcrossref{\DUrole{std,std-ref}{(Lysogorskiy21)}}}}.

\sphinxAtStartPar
In ACE, the total energy is decomposed into a sum over atomic
energies. The energy of atom \sphinxstyleemphasis{i} is expressed as a linear or non\sphinxhyphen{}linear
function of one or more density functions.  By projecting the density
onto a local atomic base, the lowest order contributions to the energy
can be expressed as a set of scalar polynomials in basis function
contributions summed over neighbor atoms.

\sphinxAtStartPar
Only a single pair\_coeff command is used with the \sphinxstyleemphasis{pace} style which
specifies an ACE coefficient file followed by N additional arguments
specifying the mapping of ACE elements to LAMMPS atom types, where N is
the number of LAMMPS atom types:
\begin{itemize}
\item {} 
\sphinxAtStartPar
ACE coefficient file (.yaml or .yace/.ace format)

\item {} 
\sphinxAtStartPar
N element names = mapping of ACE elements to atom types

\end{itemize}

\sphinxAtStartPar
Only a single pair\_coeff command is used with the \sphinxstyleemphasis{pace} style which
specifies an ACE file that fully defines the potential.  Note that
unlike for other potentials, cutoffs are not set in the pair\_style or
pair\_coeff command; they are specified in the ACE file.

\sphinxAtStartPar
The pair\_style \sphinxstyleemphasis{pace} command may be followed by the optional keyword
\sphinxstyleemphasis{product} or \sphinxstyleemphasis{recursive}, which determines which of two algorithms is
used for the calculation of basis functions and derivatives.  The
default is \sphinxstyleemphasis{recursive}.

\sphinxAtStartPar
The keyword \sphinxstyleemphasis{chunksize} is only applicable when using the pair style
\sphinxstyleemphasis{pace} with the KOKKOS package on GPUs and is ignored otherwise.  This
keyword controls the number of atoms in each pass used to compute the
atomic cluster expansion and is used to avoid running out of memory.
For example if there are 8192 atoms in the simulation and the
\sphinxstyleemphasis{chunksize} is set to 4096, the ACE calculation will be broken up into
two passes (running on a single GPU).


\subsection{Extrapolation grade}
\label{\detokenize{pair_pace:extrapolation-grade}}
\sphinxAtStartPar
Calculation of extrapolation grade in PACE is implemented in \sphinxtitleref{pair\_style
pace/extrapolation}.  It is based on the MaxVol algorithm similar to
Moment Tensor Potential (MTP) by Shapeev et al.  and is described in
{\hyperref[\detokenize{pair_pace:lysogorskiy2023}]{\sphinxcrossref{\DUrole{std,std-ref}{(Lysogorskiy23)}}}}.  In order to compute
extrapolation grade one needs to provide:
\begin{enumerate}
\sphinxsetlistlabels{\arabic}{enumi}{enumii}{}{.}%
\item {} 
\sphinxAtStartPar
ACE potential in B\sphinxhyphen{}basis form (\sphinxtitleref{.yaml} format) and

\item {} 
\sphinxAtStartPar
Active Set Inverted (ASI) file for corresponding potential (\sphinxtitleref{.asi} format)

\end{enumerate}

\sphinxAtStartPar
Calculation of extrapolation grades requires matrix\sphinxhyphen{}vector
multiplication for each atom and is slower than the usual \sphinxtitleref{pair\_style
pace recursive}, therefore it is \sphinxstyleemphasis{not} computed by default.
Extrapolation grade calculation is involved by \sphinxtitleref{fix pair}, which
requests to compute \sphinxtitleref{gamma}, as shown in example below:

\begin{sphinxVerbatim}[commandchars=\\\{\}]
\PYG{k}{pair\PYGZus{}style}\PYG{+w}{  }\PYG{n}{pace}\PYG{o}{/}\PYG{n}{extrapolation}
\PYG{k}{pair\PYGZus{}coeff}\PYG{+w}{  }\PYG{o}{*}\PYG{+w}{ }\PYG{o}{*}\PYG{+w}{ }\PYG{n}{Cu.yaml}\PYG{+w}{ }\PYG{n}{Cu.asi}\PYG{+w}{ }\PYG{n}{Cu}

\PYG{k}{fix }\PYG{n+nv+nvIdentifier}{pace\PYGZus{}gamma}\PYG{+w}{ }\PYG{n+nv+nvIdentifier}{all}\PYG{+w}{ }\PYG{n}{pair}\PYG{+w}{ }\PYG{l+m}{10}\PYG{+w}{ }\PYG{n}{pace}\PYG{o}{/}\PYG{n}{extrapolation}\PYG{+w}{ }\PYG{n}{gamma}\PYG{+w}{ }\PYG{l+m}{1}

\PYG{k}{compute }\PYG{n+nv+nvIdentifier}{max\PYGZus{}pace\PYGZus{}gamma}\PYG{+w}{ }\PYG{n+nv+nvIdentifier}{all}\PYG{+w}{ }\PYG{n}{reduce}\PYG{+w}{ }\PYG{n}{max}\PYG{+w}{ }\PYG{n}{f\PYGZus{}pace\PYGZus{}gamma}
\PYG{k}{variable }\PYG{n+nv+nvIdentifier}{dump\PYGZus{}skip}\PYG{+w}{ }\PYG{n}{equal}\PYG{+w}{ }\PYG{l+s}{\PYGZdq{}}\PYG{l+s}{c\PYGZus{}max\PYGZus{}pace\PYGZus{}gamma \PYGZlt{} 5}\PYG{l+s}{\PYGZdq{}}

\PYG{k}{dump }\PYG{n+nv+nvIdentifier}{pace\PYGZus{}dump}\PYG{+w}{ }\PYG{n+nv+nvIdentifier}{all}\PYG{+w}{ }\PYG{n}{custom}\PYG{+w}{ }\PYG{l+m}{20}\PYG{+w}{ }\PYG{n}{extrapolative\PYGZus{}structures.dump}\PYG{+w}{ }\PYG{n}{id}\PYG{+w}{ }\PYG{n}{type}\PYG{+w}{ }\PYG{n}{x}\PYG{+w}{ }\PYG{n}{y}\PYG{+w}{ }\PYG{n}{z}\PYG{+w}{ }\PYG{n}{f\PYGZus{}pace\PYGZus{}gamma}
\PYG{k}{dump\PYGZus{}modify }\PYG{n+nv+nvIdentifier}{pace\PYGZus{}dump}\PYG{+w}{ }\PYG{n}{skip}\PYG{+w}{ }\PYG{n}{v\PYGZus{}dump\PYGZus{}skip}

\PYG{k}{variable }\PYG{n+nv+nvIdentifier}{max\PYGZus{}pace\PYGZus{}gamma}\PYG{+w}{ }\PYG{n}{equal}\PYG{+w}{ }\PYG{n}{c\PYGZus{}max\PYGZus{}pace\PYGZus{}gamma}
\PYG{k}{fix }\PYG{n+nv+nvIdentifier}{extreme\PYGZus{}extrapolation}\PYG{+w}{ }\PYG{n+nv+nvIdentifier}{all}\PYG{+w}{ }\PYG{n}{halt}\PYG{+w}{ }\PYG{l+m}{10}\PYG{+w}{ }\PYG{n}{v\PYGZus{}max\PYGZus{}pace\PYGZus{}gamma}\PYG{+w}{ }\PYG{o}{\PYGZgt{}}\PYG{+w}{ }\PYG{l+m}{25}
\end{sphinxVerbatim}

\sphinxAtStartPar
Here extrapolation grade gamma is computed every 10 steps and is stored
in \sphinxtitleref{f\_pace\_gamma} per\sphinxhyphen{}atom variable.  The largest value of extrapolation
grade among all atoms in a structure is reduced to \sphinxtitleref{c\_max\_pace\_gamma}
variable.  Only if this value exceeds extrapolation threshold 5, then
the structure will be dumped into \sphinxtitleref{extrapolative\_structures.dump} file,
but not more often than every 20 steps.

\sphinxAtStartPar
On all other steps \sphinxtitleref{pair\_style pace recursive} will be used.

\sphinxAtStartPar
When using the pair style \sphinxstyleemphasis{pace/extrapolation} with the KOKKOS package on GPUs
product B\sphinxhyphen{}basis evaluator is always used and only \sphinxstyleemphasis{linear} ASI is supported.


\bigskip\hrule\bigskip


\sphinxAtStartPar
See the {\hyperref[\detokenize{pair_coeff::doc}]{\sphinxcrossref{\DUrole{doc}{pair\_coeff}}}} page for alternate ways
to specify the path for the ACE coefficient file.


\subsection{Core repulsion}
\label{\detokenize{pair_pace:core-repulsion}}
\sphinxAtStartPar
The ACE potential can be configured to initiate core\sphinxhyphen{}repulsion from an inner cutoff,
seamlessly transitioning from ACE to ZBL. The core repulsion factor can be accessed
as a per\sphinxhyphen{}atom quantity, as demonstrated in the example below:

\begin{sphinxVerbatim}[commandchars=\\\{\}]
\PYG{k}{pair\PYGZus{}style}\PYG{+w}{  }\PYG{n}{pace}
\PYG{k}{pair\PYGZus{}coeff}\PYG{+w}{  }\PYG{o}{*}\PYG{+w}{ }\PYG{o}{*}\PYG{+w}{ }\PYG{n}{CuNi.yaml}\PYG{+w}{ }\PYG{n}{Cu}\PYG{+w}{ }\PYG{n}{Ni}

\PYG{k}{fix }\PYG{n+nv+nvIdentifier}{pace\PYGZus{}corerep}\PYG{+w}{ }\PYG{n+nv+nvIdentifier}{all}\PYG{+w}{ }\PYG{n}{pair}\PYG{+w}{ }\PYG{l+m}{1}\PYG{+w}{ }\PYG{n}{pace}\PYG{+w}{ }\PYG{n}{corerep}\PYG{+w}{ }\PYG{l+m}{1}
\end{sphinxVerbatim}

\sphinxAtStartPar
In this case, per\sphinxhyphen{}atom \sphinxtitleref{f\_pace\_corerep} quantities represent the fraction of ZBL
core\sphinxhyphen{}repulsion for each atom.


\subsection{Mixing, shift, table, tail correction, restart, rRESPA info}
\label{\detokenize{pair_pace:mixing-shift-table-tail-correction-restart-rrespa-info}}
\sphinxAtStartPar
For atom type pairs I,J and I != J, where types I and J correspond to
two different element types, mixing is performed by LAMMPS with
user\sphinxhyphen{}specifiable parameters as described above.  You never need to
specify a pair\_coeff command with I != J arguments for this style.

\sphinxAtStartPar
This pair style does not support the {\hyperref[\detokenize{pair_modify::doc}]{\sphinxcrossref{\DUrole{doc}{pair\_modify}}}}
shift, table, and tail options.

\sphinxAtStartPar
This pair style does not write its information to {\hyperref[\detokenize{restart::doc}]{\sphinxcrossref{\DUrole{doc}{binary restart
files}}}}, since it is stored in potential files.  Thus, you need
to re\sphinxhyphen{}specify the pair\_style and pair\_coeff commands in an input script
that reads a restart file.

\sphinxAtStartPar
This pair style can only be used via the \sphinxstyleemphasis{pair} keyword of the
{\hyperref[\detokenize{run_style::doc}]{\sphinxcrossref{\DUrole{doc}{run\_style respa}}}} command.  It does not support the
\sphinxstyleemphasis{inner}, \sphinxstyleemphasis{middle}, \sphinxstyleemphasis{outer} keywords.


\bigskip\hrule\bigskip


\sphinxAtStartPar
Styles with a \sphinxstyleemphasis{gpu}, \sphinxstyleemphasis{intel}, \sphinxstyleemphasis{kk}, \sphinxstyleemphasis{omp}, or \sphinxstyleemphasis{opt} suffix are
functionally the same as the corresponding style without the suffix.
They have been optimized to run faster, depending on your available
hardware, as discussed on the {\hyperref[\detokenize{Speed_packages::doc}]{\sphinxcrossref{\DUrole{doc}{Accelerator packages}}}}
page.  The accelerated styles take the same arguments and should
produce the same results, except for round\sphinxhyphen{}off and precision issues.

\sphinxAtStartPar
These accelerated styles are part of the GPU, INTEL, KOKKOS,
OPENMP, and OPT packages, respectively.  They are only enabled if
LAMMPS was built with those packages.  See the {\hyperref[\detokenize{Build_package::doc}]{\sphinxcrossref{\DUrole{doc}{Build package}}}} page for more info.

\sphinxAtStartPar
You can specify the accelerated styles explicitly in your input script
by including their suffix, or you can use the {\hyperref[\detokenize{Run_options::doc}]{\sphinxcrossref{\DUrole{doc}{\sphinxhyphen{}suffix command\sphinxhyphen{}line switch}}}} when you invoke LAMMPS, or you can use the
{\hyperref[\detokenize{suffix::doc}]{\sphinxcrossref{\DUrole{doc}{suffix}}}} command in your input script.

\sphinxAtStartPar
See the {\hyperref[\detokenize{Speed_packages::doc}]{\sphinxcrossref{\DUrole{doc}{Accelerator packages}}}} page for more
instructions on how to use the accelerated styles effectively.


\bigskip\hrule\bigskip



\subsection{Restrictions}
\label{\detokenize{pair_pace:restrictions}}
\sphinxAtStartPar
This pair style is part of the ML\sphinxhyphen{}PACE package.  It is only enabled if
LAMMPS was built with that package.  See the {\hyperref[\detokenize{Build_package::doc}]{\sphinxcrossref{\DUrole{doc}{Build package}}}} page for more info.


\subsection{Related commands}
\label{\detokenize{pair_pace:related-commands}}
\sphinxAtStartPar
{\hyperref[\detokenize{pair_snap::doc}]{\sphinxcrossref{\DUrole{doc}{pair\_style snap}}}},
{\hyperref[\detokenize{fix_pair::doc}]{\sphinxcrossref{\DUrole{doc}{fix pair}}}}


\subsection{Default}
\label{\detokenize{pair_pace:default}}
\sphinxAtStartPar
recursive, chunksize = 4096,

\phantomsection\label{\detokenize{pair_pace:drautz20191}}
\sphinxAtStartPar
\sphinxstylestrong{(Drautz19)} Drautz, Phys Rev B, 99, 014104 (2019).

\phantomsection\label{\detokenize{pair_pace:lysogorskiy20211}}
\sphinxAtStartPar
\sphinxstylestrong{(Lysogorskiy21)} Lysogorskiy, van der Oord, Bochkarev, Menon, Rinaldi, Hammerschmidt, Mrovec, Thompson, Csanyi, Ortner, Drautz, npj Comp Mat, 7, 97 (2021).

\phantomsection\label{\detokenize{pair_pace:lysogorskiy2023}}
\sphinxAtStartPar
\sphinxstylestrong{(Lysogorskiy23)} Lysogorskiy, Bochkarev, Mrovec, Drautz, Phys Rev Mater, 7, 043801 (2023) / arXiv:2212.08716 (2022).

\sphinxstepscope

\index{pair\_style pedone@\spxentry{pair\_style pedone}}\ignorespaces 
\index{pair\_style pedone/omp@\spxentry{pair\_style pedone/omp}}\ignorespaces 

\section{pair\_style pedone command}
\label{\detokenize{pair_pedone:pair-style-pedone-command}}\label{\detokenize{pair_pedone:index-1}}\label{\detokenize{pair_pedone:index-0}}\label{\detokenize{pair_pedone::doc}}
\sphinxAtStartPar
Accelerator Variants: \sphinxstyleemphasis{pedone/omp}


\subsection{Syntax}
\label{\detokenize{pair_pedone:syntax}}
\begin{sphinxVerbatim}[commandchars=\\\{\}]
\PYG{k}{pair\PYGZus{}style}\PYG{+w}{ }\PYG{n}{style}\PYG{+w}{ }\PYG{n}{args}
\end{sphinxVerbatim}
\begin{itemize}
\item {} 
\sphinxAtStartPar
style = pedone*

\item {} 
\sphinxAtStartPar
args = list of arguments for a particular style

\end{itemize}
\begin{sphinxalltt}
\sphinxstyleemphasis{pedone} args = cutoff
  cutoff = global cutoff for Pedone interactions (distance units)
\end{sphinxalltt}


\subsection{Examples}
\label{\detokenize{pair_pedone:examples}}
\begin{sphinxVerbatim}[commandchars=\\\{\}]
\PYG{k}{pair\PYGZus{}style}\PYG{+w}{ }\PYG{n}{hybrid}\PYG{o}{/}\PYG{n}{overlay}\PYG{+w}{ }\PYG{n}{pedone}\PYG{+w}{ }\PYG{l+m}{15.0}\PYG{+w}{ }\PYG{n}{coul}\PYG{o}{/}\PYG{n}{long}\PYG{+w}{ }\PYG{l+m}{15.0}
\PYG{k}{kspace\PYGZus{}style}\PYG{+w}{ }\PYG{n}{pppm}\PYG{+w}{ }\PYG{l+m}{1.0e\PYGZhy{}5}

\PYG{k}{pair\PYGZus{}coeff}\PYG{+w}{ }\PYG{o}{*}\PYG{+w}{ }\PYG{o}{*}\PYG{+w}{ }\PYG{n}{coul}\PYG{o}{/}\PYG{n}{long}
\PYG{k}{pair\PYGZus{}coeff}\PYG{+w}{ }\PYG{l+m}{1}\PYG{+w}{ }\PYG{l+m}{2}\PYG{+w}{ }\PYG{n}{pedone}\PYG{+w}{ }\PYG{l+m}{0.030211}\PYG{+w}{ }\PYG{l+m}{2.241334}\PYG{+w}{ }\PYG{l+m}{2.923245}\PYG{+w}{ }\PYG{l+m}{5.0}
\PYG{k}{pair\PYGZus{}coeff}\PYG{+w}{ }\PYG{l+m}{2}\PYG{+w}{ }\PYG{l+m}{2}\PYG{+w}{ }\PYG{n}{pedone}\PYG{+w}{ }\PYG{l+m}{0.042395}\PYG{+w}{ }\PYG{l+m}{1.379316}\PYG{+w}{ }\PYG{l+m}{3.618701}\PYG{+w}{ }\PYG{l+m}{22.0}
\end{sphinxVerbatim}

\sphinxAtStartPar
Used in input scripts:
\begin{quote}

\begin{sphinxVerbatim}[commandchars=\\\{\}]
examples/PACKAGES/pedone/in.pedone.relax
examples/PACKAGES/pedone/in.pedone.melt
\end{sphinxVerbatim}
\end{quote}


\subsection{Description}
\label{\detokenize{pair_pedone:description}}
\sphinxAtStartPar
\DUrole{versionmodified,added}{Added in version 17Apr2024.}

\sphinxAtStartPar
Pair style \sphinxstyleemphasis{pedone} computes the \sphinxstylestrong{non\sphinxhyphen{}Coulomb} interactions of the Pedone
(or PMMCS) potential {\hyperref[\detokenize{pair_pedone:pedone}]{\sphinxcrossref{\DUrole{std,std-ref}{(Pedone)}}}} which combines Coulomb
interactions, Morse potential, and repulsive \(r^{-12}\)
Lennard\sphinxhyphen{}Jones terms (see below).  The \sphinxstyleemphasis{pedone} pair style is meant
to be used in addition to a {\hyperref[\detokenize{pair_coul::doc}]{\sphinxcrossref{\DUrole{doc}{Coulomb pair style}}}} via
pair style {\hyperref[\detokenize{pair_hybrid::doc}]{\sphinxcrossref{\DUrole{doc}{hybrid/overlay}}}} (see example above).
Using \sphinxstyleemphasis{coul/long} or \sphinxstyleemphasis{could/dsf} (for solids) is recommended.

\sphinxAtStartPar
The full Pedone potential function from {\hyperref[\detokenize{pair_pedone:pedone}]{\sphinxcrossref{\DUrole{std,std-ref}{(Pedone)}}}} for each
pair of atoms is:
\begin{equation*}
\begin{split}E =  \frac{C q_i q_j}{\epsilon  r}
    + D_0 \left[ e^{- 2 \alpha (r - r_0)} - 2 e^{- \alpha (r - r_0)} \right]
    + \frac{B_0}{r^{12}} \qquad r < r_c\end{split}
\end{equation*}
\sphinxAtStartPar
\(r_c\) is the cutoff and \(C\) is a conversion factor that is
specific to the choice of {\hyperref[\detokenize{units::doc}]{\sphinxcrossref{\DUrole{doc}{units}}}} so that the entire
Coulomb term is in energy units with \(q_i\) and \(q_j\) as the
assigned charges in multiples of the elementary charge.

\sphinxAtStartPar
The following coefficients must be defined for the selected pairs of
atom types via the {\hyperref[\detokenize{pair_coeff::doc}]{\sphinxcrossref{\DUrole{doc}{pair\_coeff}}}} command as in the
example above:
\begin{itemize}
\item {} 
\sphinxAtStartPar
\(D_0\) (energy units)

\item {} 
\sphinxAtStartPar
\(\alpha\) (1/distance units)

\item {} 
\sphinxAtStartPar
\(r_0\) (distance units)

\item {} 
\sphinxAtStartPar
\(C_0\) (energy units)

\item {} 
\sphinxAtStartPar
cutoff (distance units)

\end{itemize}

\sphinxAtStartPar
The last coefficient is optional.  If not specified, the global \sphinxstyleemphasis{pedone}
cutoff is used.


\bigskip\hrule\bigskip


\sphinxAtStartPar
Styles with a \sphinxstyleemphasis{gpu}, \sphinxstyleemphasis{intel}, \sphinxstyleemphasis{kk}, \sphinxstyleemphasis{omp}, or \sphinxstyleemphasis{opt} suffix are
functionally the same as the corresponding style without the suffix.
They have been optimized to run faster, depending on your available
hardware, as discussed on the {\hyperref[\detokenize{Speed_packages::doc}]{\sphinxcrossref{\DUrole{doc}{Accelerator packages}}}}
page.  The accelerated styles take the same arguments and should
produce the same results, except for round\sphinxhyphen{}off and precision issues.

\sphinxAtStartPar
These accelerated styles are part of the GPU, INTEL, KOKKOS,
OPENMP, and OPT packages, respectively.  They are only enabled if
LAMMPS was built with those packages.  See the {\hyperref[\detokenize{Build_package::doc}]{\sphinxcrossref{\DUrole{doc}{Build package}}}} page for more info.

\sphinxAtStartPar
You can specify the accelerated styles explicitly in your input script
by including their suffix, or you can use the {\hyperref[\detokenize{Run_options::doc}]{\sphinxcrossref{\DUrole{doc}{\sphinxhyphen{}suffix command\sphinxhyphen{}line switch}}}} when you invoke LAMMPS, or you can use the
{\hyperref[\detokenize{suffix::doc}]{\sphinxcrossref{\DUrole{doc}{suffix}}}} command in your input script.

\sphinxAtStartPar
See the {\hyperref[\detokenize{Speed_packages::doc}]{\sphinxcrossref{\DUrole{doc}{Accelerator packages}}}} page for more
instructions on how to use the accelerated styles effectively.


\bigskip\hrule\bigskip



\subsection{Mixing, shift, table, tail correction, restart, rRESPA info}
\label{\detokenize{pair_pedone:mixing-shift-table-tail-correction-restart-rrespa-info}}
\sphinxAtStartPar
This pair style does not support mixing.

\sphinxAtStartPar
This pair style support the {\hyperref[\detokenize{pair_modify::doc}]{\sphinxcrossref{\DUrole{doc}{pair\_modify}}}} shift
option for the energy of the pair interaction.

\sphinxAtStartPar
This pair style does not support the {\hyperref[\detokenize{pair_modify::doc}]{\sphinxcrossref{\DUrole{doc}{pair\_modify}}}}
tail option for adding long\sphinxhyphen{}range tail corrections to energy and
pressure.

\sphinxAtStartPar
This pair style writes its information to {\hyperref[\detokenize{restart::doc}]{\sphinxcrossref{\DUrole{doc}{binary restart files}}}},
so pair\_style and pair\_coeff commands does not need to be specified in an input
script that reads a restart file.

\sphinxAtStartPar
This pair style can only be used via the \sphinxstyleemphasis{pair} keyword of the
{\hyperref[\detokenize{run_style::doc}]{\sphinxcrossref{\DUrole{doc}{run\_style respa}}}} command.  It does not support the
\sphinxstyleemphasis{inner}, \sphinxstyleemphasis{middle}, or \sphinxstyleemphasis{outer} keywords.


\bigskip\hrule\bigskip



\subsection{Restrictions}
\label{\detokenize{pair_pedone:restrictions}}
\sphinxAtStartPar
The \sphinxstyleemphasis{pedone} pair style is only enabled if LAMMPS was built with the
EXTRA\sphinxhyphen{}PAIR package.  See the {\hyperref[\detokenize{Build_package::doc}]{\sphinxcrossref{\DUrole{doc}{Build package}}}} page
for more info.


\subsection{Related commands}
\label{\detokenize{pair_pedone:related-commands}}
\sphinxAtStartPar
{\hyperref[\detokenize{pair_coeff::doc}]{\sphinxcrossref{\DUrole{doc}{pair\_coeff}}}}, {\hyperref[\detokenize{pair_style::doc}]{\sphinxcrossref{\DUrole{doc}{pair\_style}}}},
{\hyperref[\detokenize{pair_coul::doc}]{\sphinxcrossref{\DUrole{doc}{pair style coul/long and coul/dsf}}}},
{\hyperref[\detokenize{pair_morse::doc}]{\sphinxcrossref{\DUrole{doc}{pair style morse}}}}


\subsection{Default}
\label{\detokenize{pair_pedone:default}}
\sphinxAtStartPar
none


\bigskip\hrule\bigskip

\phantomsection\label{\detokenize{pair_pedone:pedone}}
\sphinxAtStartPar
\sphinxstylestrong{(Pedone)} A. Pedone, G. Malavasi, M. C. Menziani, A. N. Cormack, and U. Segre, J. Phys. Chem. B, 110, 11780 (2006)

\sphinxstepscope

\index{pair\_style peri/pmb@\spxentry{pair\_style peri/pmb}}\ignorespaces 
\index{pair\_style peri/pmb/omp@\spxentry{pair\_style peri/pmb/omp}}\ignorespaces 
\index{pair\_style peri/lps@\spxentry{pair\_style peri/lps}}\ignorespaces 
\index{pair\_style peri/lps/omp@\spxentry{pair\_style peri/lps/omp}}\ignorespaces 
\index{pair\_style peri/ves@\spxentry{pair\_style peri/ves}}\ignorespaces 
\index{pair\_style peri/eps@\spxentry{pair\_style peri/eps}}\ignorespaces 

\section{pair\_style peri/pmb command}
\label{\detokenize{pair_peri:pair-style-peri-pmb-command}}\label{\detokenize{pair_peri:index-5}}\label{\detokenize{pair_peri:index-4}}\label{\detokenize{pair_peri:index-3}}\label{\detokenize{pair_peri:index-2}}\label{\detokenize{pair_peri:index-1}}\label{\detokenize{pair_peri:index-0}}\label{\detokenize{pair_peri::doc}}
\sphinxAtStartPar
Accelerator Variants: \sphinxstyleemphasis{peri/pmb/omp}


\section{pair\_style peri/lps command}
\label{\detokenize{pair_peri:pair-style-peri-lps-command}}
\sphinxAtStartPar
Accelerator Variants: \sphinxstyleemphasis{peri/lps/omp}


\section{pair\_style peri/ves command}
\label{\detokenize{pair_peri:pair-style-peri-ves-command}}

\section{pair\_style peri/eps command}
\label{\detokenize{pair_peri:pair-style-peri-eps-command}}

\subsection{Syntax}
\label{\detokenize{pair_peri:syntax}}
\begin{sphinxVerbatim}[commandchars=\\\{\}]
\PYG{k}{pair\PYGZus{}style}\PYG{+w}{ }\PYG{n}{style}
\end{sphinxVerbatim}
\begin{itemize}
\item {} 
\sphinxAtStartPar
style = \sphinxstyleemphasis{peri/pmb} or \sphinxstyleemphasis{peri/lps} or \sphinxstyleemphasis{peri/ves} or \sphinxstyleemphasis{peri/eps}

\end{itemize}


\subsection{Examples}
\label{\detokenize{pair_peri:examples}}
\begin{sphinxVerbatim}[commandchars=\\\{\}]
\PYG{k}{pair\PYGZus{}style}\PYG{+w}{ }\PYG{n}{peri}\PYG{o}{/}\PYG{n}{pmb}
\PYG{k}{pair\PYGZus{}coeff}\PYG{+w}{ }\PYG{o}{*}\PYG{+w}{ }\PYG{o}{*}\PYG{+w}{ }\PYG{l+m}{1.6863e22}\PYG{+w}{ }\PYG{l+m}{0.0015001}\PYG{+w}{ }\PYG{l+m}{0.0005}\PYG{+w}{ }\PYG{l+m}{0.25}

\PYG{k}{pair\PYGZus{}style}\PYG{+w}{ }\PYG{n}{peri}\PYG{o}{/}\PYG{n}{lps}
\PYG{k}{pair\PYGZus{}coeff}\PYG{+w}{ }\PYG{o}{*}\PYG{+w}{ }\PYG{o}{*}\PYG{+w}{ }\PYG{l+m}{14.9e9}\PYG{+w}{ }\PYG{l+m}{14.9e9}\PYG{+w}{ }\PYG{l+m}{0.0015001}\PYG{+w}{ }\PYG{l+m}{0.0005}\PYG{+w}{ }\PYG{l+m}{0.25}

\PYG{k}{pair\PYGZus{}style}\PYG{+w}{ }\PYG{n}{peri}\PYG{o}{/}\PYG{n}{ves}
\PYG{k}{pair\PYGZus{}coeff}\PYG{+w}{ }\PYG{o}{*}\PYG{+w}{ }\PYG{o}{*}\PYG{+w}{ }\PYG{l+m}{14.9e9}\PYG{+w}{ }\PYG{l+m}{14.9e9}\PYG{+w}{ }\PYG{l+m}{0.0015001}\PYG{+w}{ }\PYG{l+m}{0.0005}\PYG{+w}{ }\PYG{l+m}{0.25}\PYG{+w}{ }\PYG{l+m}{0.5}\PYG{+w}{ }\PYG{l+m}{0.001}

\PYG{k}{pair\PYGZus{}style}\PYG{+w}{ }\PYG{n}{peri}\PYG{o}{/}\PYG{n}{eps}
\PYG{k}{pair\PYGZus{}coeff}\PYG{+w}{ }\PYG{o}{*}\PYG{+w}{ }\PYG{o}{*}\PYG{+w}{ }\PYG{l+m}{14.9e9}\PYG{+w}{ }\PYG{l+m}{14.9e9}\PYG{+w}{ }\PYG{l+m}{0.0015001}\PYG{+w}{ }\PYG{l+m}{0.0005}\PYG{+w}{ }\PYG{l+m}{0.25}\PYG{+w}{ }\PYG{l+m}{118.43}
\end{sphinxVerbatim}


\subsection{Description}
\label{\detokenize{pair_peri:description}}
\sphinxAtStartPar
The peridynamic pair styles implement material models that can be used
at the mesoscopic and macroscopic scales.  See \sphinxhref{PDF/PDLammps\_overview.pdf}{this document} for an overview of LAMMPS commands for
Peridynamics modeling.

\sphinxAtStartPar
Style \sphinxstyleemphasis{peri/pmb} implements the Peridynamic bond\sphinxhyphen{}based prototype
microelastic brittle (PMB) model.

\sphinxAtStartPar
Style \sphinxstyleemphasis{peri/lps} implements the Peridynamic state\sphinxhyphen{}based linear
peridynamic solid (LPS) model.

\sphinxAtStartPar
Style \sphinxstyleemphasis{peri/ves} implements the Peridynamic state\sphinxhyphen{}based linear
peridynamic viscoelastic solid (VES) model.

\sphinxAtStartPar
Style \sphinxstyleemphasis{peri/eps} implements the Peridynamic state\sphinxhyphen{}based elastic\sphinxhyphen{}plastic
solid (EPS) model.

\sphinxAtStartPar
The canonical papers on Peridynamics are {\hyperref[\detokenize{pair_peri:silling2000}]{\sphinxcrossref{\DUrole{std,std-ref}{(Silling 2000)}}}} and {\hyperref[\detokenize{pair_peri:silling2007}]{\sphinxcrossref{\DUrole{std,std-ref}{(Silling 2007)}}}}.  The
implementation of Peridynamics in LAMMPS is described in {\hyperref[\detokenize{pair_peri:parks}]{\sphinxcrossref{\DUrole{std,std-ref}{(Parks)}}}}.  Also see the {\hyperref[\detokenize{Howto_peri::doc}]{\sphinxcrossref{\DUrole{doc}{Peridynamics Howto}}}} for more
details about its implementation.

\sphinxAtStartPar
The peridynamic VES and EPS models in PDLAMMPS were implemented by
R. Rahman and J. T. Foster at University of Texas at San Antonio.  The
original VES formulation is described in “(Mitchell2011)” and the
original EPS formulation is in “(Mitchell2011a)”.  Additional PDF docs
that describe the VES and EPS implementations are include in the LAMMPS
distribution in \sphinxhref{PDF/PDLammps\_VES.pdf}{doc/PDF/PDLammps\_VES.pdf} and
\sphinxhref{PDF/PDLammps\_EPS.pdf}{doc/PDF/PDLammps\_EPS.pdf}.  For questions
regarding the VES and EPS models in LAMMPS you can contact R. Rahman
(rezwanur.rahman at utsa.edu).

\sphinxAtStartPar
The following coefficients must be defined for each pair of atom types
via the {\hyperref[\detokenize{pair_coeff::doc}]{\sphinxcrossref{\DUrole{doc}{pair\_coeff}}}} command as in the examples above,
or in the data file or restart files read by the {\hyperref[\detokenize{read_data::doc}]{\sphinxcrossref{\DUrole{doc}{read\_data}}}} or {\hyperref[\detokenize{read_restart::doc}]{\sphinxcrossref{\DUrole{doc}{read\_restart}}}} commands, or by
mixing as described below.

\sphinxAtStartPar
For the \sphinxstyleemphasis{peri/pmb} style:
\begin{itemize}
\item {} 
\sphinxAtStartPar
c (energy/distance/volume\textasciicircum{}2 units)

\item {} 
\sphinxAtStartPar
horizon (distance units)

\item {} 
\sphinxAtStartPar
s00 (unitless)

\item {} 
\sphinxAtStartPar
\(\alpha\) (unitless)

\end{itemize}

\sphinxAtStartPar
C is the effectively a spring constant for Peridynamic bonds, the
horizon is a cutoff distance for truncating interactions, and s00 and
\(\alpha\) are used as a bond breaking criteria.  The units of c are
such that c/distance = stiffness/volume\textasciicircum{}2, where stiffness is
energy/distance\textasciicircum{}2 and volume is distance\textasciicircum{}3.  See the users guide for
more details.

\sphinxAtStartPar
For the \sphinxstyleemphasis{peri/lps} style:
\begin{itemize}
\item {} 
\sphinxAtStartPar
K (force/area units)

\item {} 
\sphinxAtStartPar
G (force/area units)

\item {} 
\sphinxAtStartPar
horizon (distance units)

\item {} 
\sphinxAtStartPar
s00 (unitless)

\item {} 
\sphinxAtStartPar
\(\alpha\) (unitless)

\end{itemize}

\sphinxAtStartPar
K is the bulk modulus and G is the shear modulus.  The horizon is a
cutoff distance for truncating interactions, and s00 and \(\alpha\)
are used as a bond breaking criteria. See the users guide for more
details.

\sphinxAtStartPar
For the \sphinxstyleemphasis{peri/ves} style:
\begin{itemize}
\item {} 
\sphinxAtStartPar
K (force/area units)

\item {} 
\sphinxAtStartPar
G (force/area units)

\item {} 
\sphinxAtStartPar
horizon (distance units)

\item {} 
\sphinxAtStartPar
s00 (unitless)

\item {} 
\sphinxAtStartPar
\(\alpha\) (unitless)

\item {} 
\sphinxAtStartPar
m\_lambdai (unitless)

\item {} 
\sphinxAtStartPar
m\_taubi (unitless)

\end{itemize}

\sphinxAtStartPar
K is the bulk modulus and G is the shear modulus. The horizon is a
cutoff distance for truncating interactions, and s00 and \(\alpha\)
are used as a bond breaking criteria. m\_lambdai and m\_taubi are the
viscoelastic relaxation parameter and time constant,
respectively. m\_lambdai varies within zero to one. For very small values
of m\_lambdai the viscoelastic model responds very similar to a linear
elastic model. For details please see the description in
“(Mitchell2011)”.

\sphinxAtStartPar
For the \sphinxstyleemphasis{peri/eps} style:
\begin{itemize}
\item {} 
\sphinxAtStartPar
K (force/area units)

\item {} 
\sphinxAtStartPar
G (force/area units)

\item {} 
\sphinxAtStartPar
horizon (distance units)

\item {} 
\sphinxAtStartPar
s00 (unitless)

\item {} 
\sphinxAtStartPar
\(\alpha\) (unitless)

\item {} 
\sphinxAtStartPar
m\_yield\_stress (force/area units)

\end{itemize}

\sphinxAtStartPar
K is the bulk modulus and G is the shear modulus. The horizon is a
cutoff distance and s00 and \(\alpha\) are used as a bond breaking
criteria.  m\_yield\_stress is the yield stress of the material. For
details please see the description in “(Mitchell2011a)”.


\bigskip\hrule\bigskip


\sphinxAtStartPar
Styles with a \sphinxstyleemphasis{gpu}, \sphinxstyleemphasis{intel}, \sphinxstyleemphasis{kk}, \sphinxstyleemphasis{omp}, or \sphinxstyleemphasis{opt} suffix are
functionally the same as the corresponding style without the suffix.
They have been optimized to run faster, depending on your available
hardware, as discussed on the {\hyperref[\detokenize{Speed_packages::doc}]{\sphinxcrossref{\DUrole{doc}{Accelerator packages}}}}
page.  The accelerated styles take the same arguments and should
produce the same results, except for round\sphinxhyphen{}off and precision issues.

\sphinxAtStartPar
These accelerated styles are part of the GPU, INTEL, KOKKOS,
OPENMP, and OPT packages, respectively.  They are only enabled if
LAMMPS was built with those packages.  See the {\hyperref[\detokenize{Build_package::doc}]{\sphinxcrossref{\DUrole{doc}{Build package}}}} page for more info.

\sphinxAtStartPar
You can specify the accelerated styles explicitly in your input script
by including their suffix, or you can use the {\hyperref[\detokenize{Run_options::doc}]{\sphinxcrossref{\DUrole{doc}{\sphinxhyphen{}suffix command\sphinxhyphen{}line switch}}}} when you invoke LAMMPS, or you can use the
{\hyperref[\detokenize{suffix::doc}]{\sphinxcrossref{\DUrole{doc}{suffix}}}} command in your input script.

\sphinxAtStartPar
See the {\hyperref[\detokenize{Speed_packages::doc}]{\sphinxcrossref{\DUrole{doc}{Accelerator packages}}}} page for more
instructions on how to use the accelerated styles effectively.


\bigskip\hrule\bigskip



\subsection{Mixing, shift, table, tail correction, restart, rRESPA info}
\label{\detokenize{pair_peri:mixing-shift-table-tail-correction-restart-rrespa-info}}
\sphinxAtStartPar
These pair styles do not support mixing.  Thus, coefficients for all
I,J pairs must be specified explicitly.

\sphinxAtStartPar
These pair styles do not support the {\hyperref[\detokenize{pair_modify::doc}]{\sphinxcrossref{\DUrole{doc}{pair\_modify}}}}
shift option.

\sphinxAtStartPar
The {\hyperref[\detokenize{pair_modify::doc}]{\sphinxcrossref{\DUrole{doc}{pair\_modify}}}} table and tail options are not
relevant for these pair styles.

\sphinxAtStartPar
These pair styles write their information to {\hyperref[\detokenize{restart::doc}]{\sphinxcrossref{\DUrole{doc}{binary restart files}}}}, so pair\_style and pair\_coeff commands do not need to be
specified in an input script that reads a restart file.

\sphinxAtStartPar
These pair styles can only be used via the \sphinxstyleemphasis{pair} keyword of the
{\hyperref[\detokenize{run_style::doc}]{\sphinxcrossref{\DUrole{doc}{run\_style respa}}}} command.  They do not support the
\sphinxstyleemphasis{inner}, \sphinxstyleemphasis{middle}, \sphinxstyleemphasis{outer} keywords.


\bigskip\hrule\bigskip



\subsection{Restrictions}
\label{\detokenize{pair_peri:restrictions}}
\sphinxAtStartPar
All of these styles are part of the PERI package. They are only enabled
if LAMMPS was built with that package.  See the {\hyperref[\detokenize{Build_package::doc}]{\sphinxcrossref{\DUrole{doc}{Build package}}}} page for more info.


\subsection{Related commands}
\label{\detokenize{pair_peri:related-commands}}
\sphinxAtStartPar
{\hyperref[\detokenize{pair_coeff::doc}]{\sphinxcrossref{\DUrole{doc}{pair\_coeff}}}}


\subsection{Default}
\label{\detokenize{pair_peri:default}}
\sphinxAtStartPar
none


\bigskip\hrule\bigskip

\phantomsection\label{\detokenize{pair_peri:parks}}
\sphinxAtStartPar
\sphinxstylestrong{(Parks)} Parks, Lehoucq, Plimpton, Silling, Comp Phys Comm, 179(11),
777\sphinxhyphen{}783 (2008).

\phantomsection\label{\detokenize{pair_peri:silling2000}}
\sphinxAtStartPar
\sphinxstylestrong{(Silling 2000)} Silling, J Mech Phys Solids, 48, 175\sphinxhyphen{}209 (2000).

\phantomsection\label{\detokenize{pair_peri:silling2007}}
\sphinxAtStartPar
\sphinxstylestrong{(Silling 2007)} Silling, Epton, Weckner, Xu, Askari, J Elasticity,
88, 151\sphinxhyphen{}184 (2007).

\phantomsection\label{\detokenize{pair_peri:mitchell2011}}
\sphinxAtStartPar
\sphinxstylestrong{(Mitchell2011)} Mitchell. A non\sphinxhyphen{}local, ordinary\sphinxhyphen{}state\sphinxhyphen{}based
viscoelasticity model for peridynamics. Sandia National Lab Report,
8064:1\sphinxhyphen{}28 (2011).

\phantomsection\label{\detokenize{pair_peri:mitchell2011a}}
\sphinxAtStartPar
\sphinxstylestrong{(Mitchell2011a)} Mitchell. A Nonlocal, Ordinary, State\sphinxhyphen{}Based
Plasticity Model for Peridynamics. Sandia National Lab Report,
3166:1\sphinxhyphen{}34 (2011).

\sphinxstepscope

\index{pair\_style pod@\spxentry{pair\_style pod}}\ignorespaces 
\index{pair\_style pod/kk@\spxentry{pair\_style pod/kk}}\ignorespaces 

\section{pair\_style pod command}
\label{\detokenize{pair_pod:pair-style-pod-command}}\label{\detokenize{pair_pod:index-1}}\label{\detokenize{pair_pod:index-0}}\label{\detokenize{pair_pod::doc}}
\sphinxAtStartPar
Accelerator Variants: \sphinxstyleemphasis{pod/kk}


\subsection{Syntax}
\label{\detokenize{pair_pod:syntax}}
\begin{sphinxVerbatim}[commandchars=\\\{\}]
\PYG{k}{pair\PYGZus{}style}\PYG{+w}{ }\PYG{n}{pod}
\end{sphinxVerbatim}


\subsection{Examples}
\label{\detokenize{pair_pod:examples}}
\begin{sphinxVerbatim}[commandchars=\\\{\}]
\PYG{k}{pair\PYGZus{}style}\PYG{+w}{ }\PYG{n}{pod}
\PYG{k}{pair\PYGZus{}coeff}\PYG{+w}{ }\PYG{o}{*}\PYG{+w}{ }\PYG{o}{*}\PYG{+w}{ }\PYG{n}{Ta\PYGZus{}param.pod}\PYG{+w}{ }\PYG{n}{Ta\PYGZus{}coefficients.pod}\PYG{+w}{ }\PYG{n}{Ta}
\end{sphinxVerbatim}


\subsection{Description}
\label{\detokenize{pair_pod:description}}
\sphinxAtStartPar
\DUrole{versionmodified,added}{Added in version 22Dec2022.}

\sphinxAtStartPar
Pair style \sphinxstyleemphasis{pod} defines the proper orthogonal descriptor (POD)
potential {\hyperref[\detokenize{pair_pod:nguyen20222b}]{\sphinxcrossref{\DUrole{std,std-ref}{(Nguyen and Rohskopf)}}}},
{\hyperref[\detokenize{pair_pod:nguyen20232b}]{\sphinxcrossref{\DUrole{std,std-ref}{(Nguyen2023)}}}}, {\hyperref[\detokenize{pair_pod:nguyen20242b}]{\sphinxcrossref{\DUrole{std,std-ref}{(Nguyen2024)}}}},
and {\hyperref[\detokenize{pair_pod:nguyen20243b}]{\sphinxcrossref{\DUrole{std,std-ref}{(Nguyen and Sema)}}}}.  The {\hyperref[\detokenize{fitpod_command::doc}]{\sphinxcrossref{\DUrole{doc}{fitpod}}}} is used to fit the POD potential.

\sphinxAtStartPar
Only a single pair\_coeff command is used with the \sphinxstyleemphasis{pod} style which
specifies a POD parameter file followed by a coefficient file, a
projection matrix file, and a centroid file.

\sphinxAtStartPar
The POD parameter file (\sphinxcode{\sphinxupquote{Ta\_param.pod}}) can contain blank and comment
lines (start with \#) anywhere. Each non\sphinxhyphen{}blank non\sphinxhyphen{}comment line must
contain one keyword/value pair. See {\hyperref[\detokenize{fitpod_command::doc}]{\sphinxcrossref{\DUrole{doc}{fitpod}}}} for
the description of all the keywords that can be assigned in the
parameter file.

\sphinxAtStartPar
The coefficient file (\sphinxcode{\sphinxupquote{Ta\_coefficients.pod}}) contains coefficients for
the POD potential. The top of the coefficient file can contain any
number of blank and comment lines (start with \#), but follows a strict
format after that. The first non\sphinxhyphen{}blank non\sphinxhyphen{}comment line must contain:
\begin{itemize}
\item {} 
\sphinxAtStartPar
model\_coefficients: \sphinxstyleemphasis{ncoeff} \sphinxstyleemphasis{nproj} \sphinxstyleemphasis{ncentroid}

\end{itemize}

\sphinxAtStartPar
This is followed by \sphinxstyleemphasis{ncoeff} coefficients, \sphinxstyleemphasis{nproj} projection matrix entries,
and \sphinxstyleemphasis{ncentroid} centroid coordinates, one per line. The coefficient
file is generated after training the POD potential using {\hyperref[\detokenize{fitpod_command::doc}]{\sphinxcrossref{\DUrole{doc}{fitpod}}}}.

\sphinxAtStartPar
As an example, if a LAMMPS indium phosphide simulation has 4 atoms
types, with the first two being indium and the third and fourth being
phophorous, the pair\_coeff command would look like this:

\begin{sphinxVerbatim}[commandchars=\\\{\}]
\PYG{k}{pair\PYGZus{}coeff}\PYG{+w}{ }\PYG{o}{*}\PYG{+w}{ }\PYG{o}{*}\PYG{+w}{ }\PYG{n}{pod}\PYG{+w}{ }\PYG{n}{InP\PYGZus{}param.pod}\PYG{+w}{ }\PYG{n}{InP\PYGZus{}coefficients.pod}\PYG{+w}{ }\PYG{n}{In}\PYG{+w}{ }\PYG{n}{In}\PYG{+w}{ }\PYG{n}{P}\PYG{+w}{ }\PYG{n}{P}
\end{sphinxVerbatim}

\sphinxAtStartPar
The first 2 arguments must be * * so as to span all LAMMPS atom types.
The two filenames are for the parameter and coefficient files, respectively.
The two trailing ‘In’ arguments map LAMMPS atom types 1 and 2 to the
POD ‘In’ element. The two trailing ‘P’ arguments map LAMMPS atom types
3 and 4 to the POD ‘P’ element.

\sphinxAtStartPar
If a POD mapping value is specified as NULL, the mapping is not
performed.  This can be used when a \sphinxstyleemphasis{pod} potential is used as part of
the \sphinxstyleemphasis{hybrid} pair style.  The NULL values are placeholders for atom
types that will be used with other potentials.

\sphinxAtStartPar
Examples about training and using POD potentials are found in the
directory lammps/examples/PACKAGES/pod and the Github repo \sphinxurl{https://github.com/cesmix-mit/pod-examples}.


\bigskip\hrule\bigskip



\subsection{Mixing, shift, table, tail correction, restart, rRESPA info}
\label{\detokenize{pair_pod:mixing-shift-table-tail-correction-restart-rrespa-info}}
\sphinxAtStartPar
For atom type pairs I,J and I != J, where types I and J correspond to
two different element types, mixing is performed by LAMMPS with
user\sphinxhyphen{}specifiable parameters as described above.  You never need to
specify a pair\_coeff command with I != J arguments for this style.

\sphinxAtStartPar
This pair style does not support the {\hyperref[\detokenize{pair_modify::doc}]{\sphinxcrossref{\DUrole{doc}{pair\_modify}}}}
shift, table, and tail options.

\sphinxAtStartPar
This pair style does not write its information to {\hyperref[\detokenize{restart::doc}]{\sphinxcrossref{\DUrole{doc}{binary restart
files}}}}, since it is stored in potential files.  Thus, you need
to re\sphinxhyphen{}specify the pair\_style and pair\_coeff commands in an input script
that reads a restart file.

\sphinxAtStartPar
This pair style can only be used via the \sphinxstyleemphasis{pair} keyword of the
{\hyperref[\detokenize{run_style::doc}]{\sphinxcrossref{\DUrole{doc}{run\_style respa}}}} command.  It does not support the
\sphinxstyleemphasis{inner}, \sphinxstyleemphasis{middle}, \sphinxstyleemphasis{outer} keywords.


\bigskip\hrule\bigskip


\sphinxAtStartPar
Styles with a \sphinxstyleemphasis{gpu}, \sphinxstyleemphasis{intel}, \sphinxstyleemphasis{kk}, \sphinxstyleemphasis{omp}, or \sphinxstyleemphasis{opt} suffix are
functionally the same as the corresponding style without the suffix.
They have been optimized to run faster, depending on your available
hardware, as discussed on the {\hyperref[\detokenize{Speed_packages::doc}]{\sphinxcrossref{\DUrole{doc}{Accelerator packages}}}}
page.  The accelerated styles take the same arguments and should
produce the same results, except for round\sphinxhyphen{}off and precision issues.

\sphinxAtStartPar
These accelerated styles are part of the GPU, INTEL, KOKKOS,
OPENMP, and OPT packages, respectively.  They are only enabled if
LAMMPS was built with those packages.  See the {\hyperref[\detokenize{Build_package::doc}]{\sphinxcrossref{\DUrole{doc}{Build package}}}} page for more info.

\sphinxAtStartPar
You can specify the accelerated styles explicitly in your input script
by including their suffix, or you can use the {\hyperref[\detokenize{Run_options::doc}]{\sphinxcrossref{\DUrole{doc}{\sphinxhyphen{}suffix command\sphinxhyphen{}line switch}}}} when you invoke LAMMPS, or you can use the
{\hyperref[\detokenize{suffix::doc}]{\sphinxcrossref{\DUrole{doc}{suffix}}}} command in your input script.

\sphinxAtStartPar
See the {\hyperref[\detokenize{Speed_packages::doc}]{\sphinxcrossref{\DUrole{doc}{Accelerator packages}}}} page for more
instructions on how to use the accelerated styles effectively.


\bigskip\hrule\bigskip



\subsection{Restrictions}
\label{\detokenize{pair_pod:restrictions}}
\sphinxAtStartPar
This style is part of the ML\sphinxhyphen{}POD package.  It is only enabled if LAMMPS
was built with that package. See the {\hyperref[\detokenize{Build_package::doc}]{\sphinxcrossref{\DUrole{doc}{Build package}}}} page for more info.


\subsection{Related commands}
\label{\detokenize{pair_pod:related-commands}}
\sphinxAtStartPar
{\hyperref[\detokenize{fitpod_command::doc}]{\sphinxcrossref{\DUrole{doc}{fitpod}}}},
{\hyperref[\detokenize{compute_pod_atom::doc}]{\sphinxcrossref{\DUrole{doc}{compute pod/atom}}}},
{\hyperref[\detokenize{compute_pod_atom::doc}]{\sphinxcrossref{\DUrole{doc}{compute podd/atom}}}},
{\hyperref[\detokenize{compute_pod_atom::doc}]{\sphinxcrossref{\DUrole{doc}{compute pod/local}}}},
{\hyperref[\detokenize{compute_pod_atom::doc}]{\sphinxcrossref{\DUrole{doc}{compute pod/global}}}}


\subsection{Default}
\label{\detokenize{pair_pod:default}}
\sphinxAtStartPar
none


\bigskip\hrule\bigskip

\phantomsection\label{\detokenize{pair_pod:nguyen20222b}}
\sphinxAtStartPar
\sphinxstylestrong{(Nguyen and Rohskopf)} Nguyen and Rohskopf,  Journal of Computational Physics, 480, 112030, (2023).

\phantomsection\label{\detokenize{pair_pod:nguyen20232b}}
\sphinxAtStartPar
\sphinxstylestrong{(Nguyen2023)} Nguyen, Physical Review B, 107(14), 144103, (2023).

\phantomsection\label{\detokenize{pair_pod:nguyen20242b}}
\sphinxAtStartPar
\sphinxstylestrong{(Nguyen2024)} Nguyen, Journal of Computational Physics, 113102, (2024).

\phantomsection\label{\detokenize{pair_pod:nguyen20243b}}
\sphinxAtStartPar
\sphinxstylestrong{(Nguyen and Sema)} Nguyen and Sema, \sphinxurl{https://arxiv.org/abs/2405.00306}, (2024).

\sphinxstepscope

\index{pair\_style polymorphic@\spxentry{pair\_style polymorphic}}\ignorespaces 

\section{pair\_style polymorphic command}
\label{\detokenize{pair_polymorphic:pair-style-polymorphic-command}}\label{\detokenize{pair_polymorphic:index-0}}\label{\detokenize{pair_polymorphic::doc}}

\subsection{Syntax}
\label{\detokenize{pair_polymorphic:syntax}}
\begin{sphinxVerbatim}[commandchars=\\\{\}]
\PYG{k}{pair\PYGZus{}style}\PYG{+w}{ }\PYG{n}{polymorphic}
\end{sphinxVerbatim}

\sphinxAtStartPar
style = \sphinxstyleemphasis{polymorphic}


\subsection{Examples}
\label{\detokenize{pair_polymorphic:examples}}
\begin{sphinxVerbatim}[commandchars=\\\{\}]
\PYG{k}{pair\PYGZus{}style}\PYG{+w}{ }\PYG{n}{polymorphic}
\PYG{k}{pair\PYGZus{}coeff}\PYG{+w}{ }\PYG{o}{*}\PYG{+w}{ }\PYG{o}{*}\PYG{+w}{ }\PYG{n}{FeCH\PYGZus{}BOP\PYGZus{}I.poly}\PYG{+w}{ }\PYG{n}{Fe}\PYG{+w}{ }\PYG{n}{C}\PYG{+w}{ }\PYG{n}{H}
\PYG{k}{pair\PYGZus{}coeff}\PYG{+w}{ }\PYG{o}{*}\PYG{+w}{ }\PYG{o}{*}\PYG{+w}{ }\PYG{n}{TlBr\PYGZus{}msw.poly}\PYG{+w}{ }\PYG{n}{Tl}\PYG{+w}{ }\PYG{n}{Br}
\PYG{k}{pair\PYGZus{}coeff}\PYG{+w}{ }\PYG{o}{*}\PYG{+w}{ }\PYG{o}{*}\PYG{+w}{ }\PYG{n}{CuTa\PYGZus{}eam.poly}\PYG{+w}{ }\PYG{n}{Cu}\PYG{+w}{ }\PYG{n}{Ta}
\PYG{k}{pair\PYGZus{}coeff}\PYG{+w}{ }\PYG{o}{*}\PYG{+w}{ }\PYG{o}{*}\PYG{+w}{ }\PYG{n}{GaN\PYGZus{}tersoff.poly}\PYG{+w}{ }\PYG{n}{Ga}\PYG{+w}{ }\PYG{n}{N}
\PYG{k}{pair\PYGZus{}coeff}\PYG{+w}{ }\PYG{o}{*}\PYG{+w}{ }\PYG{o}{*}\PYG{+w}{ }\PYG{n}{GaN\PYGZus{}sw.poly}\PYG{+w}{ }\PYG{n}{Ga}\PYG{+w}{ }\PYG{n}{N}
\end{sphinxVerbatim}


\subsection{Description}
\label{\detokenize{pair_polymorphic:description}}
\sphinxAtStartPar
The \sphinxstyleemphasis{polymorphic} pair style computes a 3\sphinxhyphen{}body free\sphinxhyphen{}form potential
({\hyperref[\detokenize{pair_polymorphic:zhou3}]{\sphinxcrossref{\DUrole{std,std-ref}{Zhou3}}}}) for the energy E of a system of atoms as
\begin{equation*}
\begin{split}E & = \frac{1}{2}\sum_{i=1}^{i=N}\sum_{j=1}^{j=N}\left[\left(1-\delta_{ij}\right)\cdot U_{IJ}\left(r_{ij}\right)-\left(1-\eta_{ij}\right)\cdot F_{IJ}\left(X_{ij}\right)\cdot V_{IJ}\left(r_{ij}\right)\right] \\
X_{ij} & = \sum_{k=i_1,k\neq j}^{i_N}W_{IK}\left(r_{ik}\right)\cdot G_{JIK}\left(\cos\theta_{jik}\right)\cdot P_{JIK}\left(\Delta r_{jik}\right) \\
\Delta r_{jik} & = r_{ij}-\xi_{IJ}\cdot r_{ik}\end{split}
\end{equation*}
\sphinxAtStartPar
where I, J, K represent species of atoms i, j, and k, \(i_1, ...,
i_N\) represents a list of \sphinxstyleemphasis{i}‘s neighbors, \(\delta_{ij}\) is a
Dirac constant (i.e., \(\delta_{ij} = 1\) when \(i = j\), and
\(\delta_{ij} = 0\) otherwise), \(\eta_{ij}\) is similar
constant that can be set either to \(\eta_{ij} = \delta_{ij}\) or
\(\eta_{ij} = 1 - \delta_{ij}\) depending on the potential type,
\(U_{IJ}(r_{ij})\), \(V_{IJ}(r_{ij})\), \(W_{IK}(r_{ik})\)
are pair functions, \(G_{JIK}(\cos\theta_{jik})\) is an angular
function, \(P_{JIK}(\Delta r_{jik})\) is a function of atomic
spacing differential \(\Delta r_{jik} = r_{ij} - \xi_{IJ} \cdot
r_{ik}\) with \(\xi_{IJ}\) being a pair\sphinxhyphen{}dependent parameter, and
\(F_{IJ}(X_{ij})\) is a function of the local environment variable
\(X_{ij}\). This generic potential is fully defined once the
constants \(\eta_{ij}\) and \(\xi_{IJ}\), and the six functions
\(U_{IJ}(r_{ij})\), \(V_{IJ}(r_{ij})\), \(W_{IK}(r_{ik})\),
\(G_{JIK}(\cos\theta_{jik})\), \(P_{JIK}(\Delta r_{jik})\), and
\(F_{IJ}(X_{ij})\) are given. Here LAMMPS uses a global parameter
\(\eta\) to represent \(\eta_{ij}\). When \(\eta = 1\),
\(\eta_{ij} = 1 - \delta_{ij}\), otherwise \(\eta_{ij} =
\delta_{ij}\). Additionally, \(\eta = 3\) indicates that the function
\(P_{JIK}(\Delta r)\) depends on species I, J and K, otherwise
\(P_{JIK}(\Delta r) = P_{IK}(\Delta r)\) only depends on species I
and K. Note that these six functions are all one dimensional, and hence
can be provided in a tabular form. This allows users to design different
potentials solely based on a manipulation of these functions. For
instance, the potential reduces to a Stillinger\sphinxhyphen{}Weber potential
({\hyperref[\detokenize{pair_polymorphic:sw}]{\sphinxcrossref{\DUrole{std,std-ref}{SW}}}}) if we set
\begin{equation*}
\begin{split}\eta_{ij} & = \delta_{ij} (\eta = 2~or~\eta = 0),\xi_{IJ}=0 \\
U_{IJ}\left(r\right) & = A_{IJ}\cdot\epsilon_{IJ}\cdot \left(\frac{\sigma_{IJ}}{r}\right)^q\cdot \left[B_{IJ}\cdot \left(\frac{\sigma_{IJ}}{r}\right)^{p-q}-1\right]\cdot exp\left(\frac{\sigma_{IJ}}{r-a_{IJ}\cdot \sigma_{IJ}}\right) \\
V_{IJ}\left(r\right) & = \sqrt{\lambda_{IJ}\cdot \epsilon_{IJ}}\cdot exp\left(\frac{\gamma_{IJ}\cdot \sigma_{IJ}}{r-a_{IJ}\cdot \sigma_{IJ}}\right) \\
F_{IJ}\left(X\right) & = -X \\
P_{JIK}\left(\Delta r\right) & = P_{IK}\left(\Delta r\right) = 1 \\
W_{IJ}\left(r\right) & = \sqrt{\lambda_{IJ}\cdot \epsilon_{IJ}}\cdot exp\left(\frac{\gamma_{IJ}\cdot \sigma_{IJ}}{r-a_{IJ}\cdot \sigma_{IJ}}\right) \\
G_{JIK}\left(\cos\theta\right) & = \left(\cos\theta+\frac{1}{3}\right)^2\end{split}
\end{equation*}
\sphinxAtStartPar
The potential reduces to a Tersoff potential ({\hyperref[\detokenize{pair_polymorphic:tersoff}]{\sphinxcrossref{\DUrole{std,std-ref}{Tersoff}}}}
or {\hyperref[\detokenize{pair_polymorphic:poly-albe}]{\sphinxcrossref{\DUrole{std,std-ref}{Albe1}}}}) if we set
\begin{equation*}
\begin{split}\eta_{ij} & = \delta_{ij} (\eta = 2~or~\eta = 0),\xi_{IJ}=1 \\
U_{IJ}\left(r\right) & = \frac{D_{e,IJ}}{S_{IJ}-1}\cdot exp\left[-\beta_{IJ}\sqrt{2S_{IJ}}\left(r-r_{e,IJ}\right)\right]\cdot f_{c,IJ}\left(r\right) \\
V_{IJ}\left(r\right) & = \frac{S_{IJ}\cdot D_{e,IJ}}{S_{IJ}-1}\cdot exp\left[-\beta_{IJ}\sqrt{\frac{2}{S_{IJ}}}\left(r-r_{e,IJ}\right)\right]\cdot f_{c,IJ}\left(r\right) \\
F_{IJ}\left(X\right) & = \left(1+X\right)^{-\frac{1}{2}} \\
P_{JIK}\left(\Delta r\right) & = P_{IK}\left(\Delta r\right) = exp\left(2\mu_{IK}\cdot \Delta r\right) \\
W_{IJ}\left(r\right) & = f_{c,IJ}\left(r\right) \\
G_{JIK}\left(\cos\theta\right) & = \gamma_{IK}\left[1+\frac{c_{IK}^2}{d_{IK}^2}-\frac{c_{IK}^2}{d_{IK}^2+\left(h_{IK}+\cos\theta\right)^2}\right]\end{split}
\end{equation*}
\sphinxAtStartPar
where
\begin{equation*}
\begin{split}f_{c,IJ}\left(r\right)=\left\{\begin{array}{l}
1, r\leq R_{IJ}-D_{IJ} \\
\frac{1}{2}+\frac{1}{2}cos\left[\frac{\pi\left(r+D_{IJ}-R_{IJ}\right)}{2D_{IJ}}\right], R_{IJ}-D_{IJ} < r < R_{IJ}+D_{IJ} \\
0, r \geq R_{IJ}+D_{IJ}
\end{array}\right.\end{split}
\end{equation*}
\sphinxAtStartPar
The potential reduces to a modified Stillinger\sphinxhyphen{}Weber potential
({\hyperref[\detokenize{pair_polymorphic:zhou3}]{\sphinxcrossref{\DUrole{std,std-ref}{Zhou3}}}}) if we set
\begin{equation*}
\begin{split}\eta_{ij} & = \delta_{ij} (\eta = 2~or~\eta = 0),\xi_{IJ}=0 \\
U_{IJ}\left(r\right) & = \varphi_{R,IJ}\left(r\right)-\varphi_{A,IJ}\left(r\right) \\
V_{IJ}\left(r\right) & = u_{IJ}\left(r\right) \\
F_{IJ}\left(X\right) & = -X \\
P_{JIK}\left(\Delta r\right) & = P_{IK}\left(\Delta r\right) = 1 \\
W_{IJ}\left(r\right) & = u_{IJ}\left(r\right) \\
G_{JIK}\left(\cos\theta\right) & = g_{JIK}\left(\cos\theta\right)\end{split}
\end{equation*}
\sphinxAtStartPar
The potential reduces to a Rockett\sphinxhyphen{}Tersoff potential ({\hyperref[\detokenize{pair_polymorphic:wang3}]{\sphinxcrossref{\DUrole{std,std-ref}{Wang3}}}}) if we set
\begin{equation*}
\begin{split}\eta_{ij} & = \delta_{ij} (\eta = 2~or~\eta = 0),\xi_{IJ}=1 \\
U_{IJ}\left(r\right) & = A_{IJ}exp\left(-\lambda_{1,IJ}\cdot r\right)f_{c,IJ}\left(r\right)f_{ca,IJ}\left(r\right) \\
V_{IJ}\left(r\right) & = \left\{\begin{array}{l}B_{IJ}exp\left(-\lambda_{2,IJ}\cdot r\right)f_{c,IJ}\left(r\right)+ \\ A_{IJ}exp\left(-\lambda_{1,IJ}\cdot r\right)f_{c,IJ}\left(r\right) \left[1-f_{ca,IJ}\left(r\right)\right]\end{array} \right\} \\
F_{IJ}\left(X\right) & = \left[1+\left(\beta_{IJ}X\right)^{n_{IJ}}\right]^{-\frac{1}{2n_{IJ}}} \\
P_{JIK}\left(\Delta r\right) & = P_{IK}\left(\Delta r\right) = exp\left(\lambda_{3,IK}\cdot \Delta r^3\right) \\
W_{IJ}\left(r\right) & = f_{c,IJ}\left(r\right) \\
G_{JIK}\left(\cos\theta\right) & = 1+\frac{c_{IK}^2}{d_{IK}^2}-\frac{c_{IK}^2}{d_{IK}^2+\left(h_{IK}+\cos\theta\right)^2}\end{split}
\end{equation*}
\sphinxAtStartPar
where \(f_{ca,IJ}(r)\) is similar to the \(f_{c,IJ}(r)\) defined
above:
\begin{equation*}
\begin{split}f_{ca,IJ}\left(r\right)=\left\{\begin{array}{l}
1, r\leq R_{a,IJ}-D_{a,IJ} \\
\frac{1}{2}+\frac{1}{2}cos\left[\frac{\pi\left(r+D_{a,IJ}-R_{a,IJ}\right)}{2D_{a,IJ}}\right], R_{a,IJ}-D_{a,IJ} < r < R_{a,IJ}+D_{a,IJ} \\
0, r \geq R_{a,IJ}+D_{a,IJ}
\end{array}\right.\end{split}
\end{equation*}
\sphinxAtStartPar
The potential becomes the embedded atom method ({\hyperref[\detokenize{pair_polymorphic:poly-daw}]{\sphinxcrossref{\DUrole{std,std-ref}{Daw}}}})
if we set
\begin{equation*}
\begin{split}\eta_{ij} & = 1-\delta_{ij} (\eta = 1),\xi_{IJ}=0 \\
U_{IJ}\left(r\right) & = \phi_{IJ}\left(r\right) \\
V_{IJ}\left(r\right) & = 1 \\
F_{II}\left(X\right) & = -2F_I\left(X\right) \\
P_{JIK}\left(\Delta r\right) & = P_{IK}\left(\Delta r\right) = 1 \\
W_{IJ}\left(r\right) & = f_{J}\left(r\right) \\
G_{JIK}\left(\cos\theta\right) & = 1\end{split}
\end{equation*}
\sphinxAtStartPar
In the embedded atom method case, \(\phi_{IJ}(r)\) is the pair
energy, \(F_I(X)\) is the embedding energy, \sphinxstyleemphasis{X} is the local
electron density, and \(f_J(r)\) is the atomic electron density
function.

\sphinxAtStartPar
The potential reduces to another type of Tersoff potential ({\hyperref[\detokenize{pair_polymorphic:zhou4}]{\sphinxcrossref{\DUrole{std,std-ref}{Zhou4}}}}) if we set
\begin{equation*}
\begin{split}\eta_{ij} & = \delta_{ij} (\eta = 3),\xi_{IJ}=1 \\
U_{IJ}\left(r\right) & = \frac{D_{e,IJ}}{S_{IJ}-1}\cdot exp\left[-\beta_{IJ}\sqrt{2S_{IJ}}\left(r-r_{e,IJ}\right)\right]\cdot f_{c,IJ}\left(r\right) \cdot T_{IJ}\left(r\right)+V_{ZBL,IJ}\left(r\right)\left[1-T_{IJ}\left(r\right)\right] \\
V_{IJ}\left(r\right) & = \frac{S_{IJ}\cdot D_{e,IJ}}{S_{IJ}-1}\cdot exp\left[-\beta_{IJ}\sqrt{\frac{2}{S_{IJ}}}\left(r-r_{e,IJ}\right)\right]\cdot f_{c,IJ}\left(r\right) \cdot T_{IJ}\left(r\right) \\
F_{IJ}\left(X\right) & = \left(1+X\right)^{-\frac{1}{2}} \\
P_{JIK}\left(\Delta r\right) & = \omega_{JIK} \cdot exp\left(\alpha_{JIK}\cdot \Delta r\right) \\
W_{IJ}\left(r\right) & = f_{c,IJ}\left(r\right) \\
G_{JIK}\left(\cos\theta\right) & = \gamma_{JIK}\left[1+\frac{c_{JIK}^2}{d_{JIK}^2}-\frac{c_{JIK}^2}{d_{JIK}^2+\left(h_{JIK}+\cos\theta\right)^2}\right] \\
T_{IJ}\left(r\right) & = \frac{1}{1+exp\left[-b_{f,IJ}\left(r-r_{f,IJ}\right)\right]} \\
V_{ZBL,IJ}\left(r\right) & = 14.4 \cdot \frac{Z_I \cdot Z_J}{r}\sum_{k=1}^{4}\mu_k \cdot exp\left[-\nu_k \left(Z_I^{0.23}+Z_J^{0.23}\right) r\right]\end{split}
\end{equation*}
\sphinxAtStartPar
where \(f_{c,IJ}(r)\) is the same as defined above. This Tersoff
potential differs from the one above because the \(P_{JIK}(\Delta
r)\) function is now dependent on all three species I, J, and K.

\sphinxAtStartPar
If the tabulated functions are created using the parameters of
Stillinger\sphinxhyphen{}Weber, Tersoff, and EAM potentials, the polymorphic pair
style will produce the same global properties (energies and stresses)
and the same forces as the {\hyperref[\detokenize{pair_sw::doc}]{\sphinxcrossref{\DUrole{doc}{sw}}}}, {\hyperref[\detokenize{pair_tersoff::doc}]{\sphinxcrossref{\DUrole{doc}{tersoff}}}}, and {\hyperref[\detokenize{pair_eam::doc}]{\sphinxcrossref{\DUrole{doc}{eam}}}} pair styles. The polymorphic
pair style also produces the same per\sphinxhyphen{}atom properties (energies and
stresses) as the corresponding {\hyperref[\detokenize{pair_tersoff::doc}]{\sphinxcrossref{\DUrole{doc}{tersoff}}}} and
{\hyperref[\detokenize{pair_eam::doc}]{\sphinxcrossref{\DUrole{doc}{eam}}}} pair styles. However, due to a different
partitioning of global properties to per\sphinxhyphen{}atom properties, the
polymorphic pair style will produce different per\sphinxhyphen{}atom properties
(energies and stresses) as the {\hyperref[\detokenize{pair_sw::doc}]{\sphinxcrossref{\DUrole{doc}{sw}}}} pair style. This does
not mean that polymorphic pair style is different from the sw pair
style. It just means that the definitions of the atom energies and atom
stresses are different.

\sphinxAtStartPar
Only a single pair\_coeff command is used with the polymorphic pair style
which specifies a potential file for all needed elements.  These are
mapped to LAMMPS atom types by specifying N additional arguments after
the filename in the pair\_coeff command, where N is the number of LAMMPS
atom types:
\begin{itemize}
\item {} 
\sphinxAtStartPar
filename

\item {} 
\sphinxAtStartPar
N element names = mapping of polymorphic potential elements to atom types

\end{itemize}

\sphinxAtStartPar
See the pair\_coeff page for alternate ways to specify the path for
the potential file. Several files for polymorphic potentials are
included in the potentials directory of the LAMMPS distribution. They
have a “poly” suffix.

\sphinxAtStartPar
As an example, imagine the GaN\_tersoff.poly file has tabulated functions
for Ga\sphinxhyphen{}N tersoff potential. If your LAMMPS simulation has 4 atom types and
you want the first 3 to be Ga, and the fourth to be N, you would use the
following pair\_coeff command:

\begin{sphinxVerbatim}[commandchars=\\\{\}]
\PYG{k}{pair\PYGZus{}coeff}\PYG{+w}{ }\PYG{o}{*}\PYG{+w}{ }\PYG{o}{*}\PYG{+w}{ }\PYG{n}{GaN\PYGZus{}tersoff.poly}\PYG{+w}{ }\PYG{n}{Ga}\PYG{+w}{ }\PYG{n}{Ga}\PYG{+w}{ }\PYG{n}{Ga}\PYG{+w}{ }\PYG{n}{N}
\end{sphinxVerbatim}

\sphinxAtStartPar
The first two arguments must be * * to span all pairs of LAMMPS atom
types. The first three Ga arguments map LAMMPS atom types 1,2,3 to the
Ga element in the polymorphic file. The final N argument maps LAMMPS
atom type 4 to the N element in the polymorphic file. If a mapping value
is specified as NULL, the mapping is not performed. This can be used
when an polymorphic potential is used as part of the hybrid pair
style. The NULL values are placeholders for atom types that will be used
with other potentials.

\sphinxAtStartPar
Potential files in the potentials directory of the LAMMPS distribution
have a “.poly” suffix. At the beginning of the files, an unlimited
number of lines starting with ‘\#’ are used to describe the potential and
are ignored by LAMMPS. The next line lists two numbers:

\begin{sphinxVerbatim}[commandchars=\\\{\}]
ntypes eta
\end{sphinxVerbatim}

\sphinxAtStartPar
Here \sphinxstyleemphasis{ntypes} represent total number of species defined in the potential
file, \(\eta = 1\) reduces to embedded atom method, \(\eta = 3\)
assumes a three species dependent \(P_{JIK}(\Delta r)\) function,
and all other values of \(\eta\) assume a two species dependent
\(P_{JK}(\Delta r)\) function. The value of \sphinxstyleemphasis{ntypes} must equal the
total number of different species defined in the pair\_coeff command. The
next \sphinxstyleemphasis{ntypes} lines each lists two numbers and a character string
representing atomic number, atomic mass, and name of the species of the
ntypes elements:

\begin{sphinxVerbatim}[commandchars=\\\{\}]
atomic\PYGZhy{}number atomic\PYGZhy{}mass element\PYGZhy{}name(1)
atomic\PYGZhy{}number atomic\PYGZhy{}mass element\PYGZhy{}name(2)
...
atomic\PYGZhy{}number atomic\PYGZhy{}mass element\PYGZhy{}name(ntypes)
\end{sphinxVerbatim}

\sphinxAtStartPar
The next line contains four numbers:

\begin{sphinxVerbatim}[commandchars=\\\{\}]
nr ntheta nx xmax
\end{sphinxVerbatim}

\sphinxAtStartPar
Here nr is total number of tabular points for radial functions U, V, W,
P, ntheta is total number of tabular points for the angular function G,
nx is total number of tabular points for the function F, xmax is a
maximum value of the argument of function F. Note that the pair
functions \(U_{IJ}(r)\), \(V_{IJ}(r)\), \(W_{IJ}(r)\) are
uniformly tabulated between 0 and cutoff distance of the IJ pair,
\(G_{JIK}(\cos\theta)\) is uniformly tabulated between \sphinxhyphen{}1 and 1,
\(P_{JIK}(\Delta r)\) is uniformly tabulated between \sphinxhyphen{}rcmax and
rcmax where rcmax is the maximum cutoff distance of all pairs, and
\(F_{IJ}(X)\) is uniformly tabulated between 0 and xmax. Linear
extrapolation is assumed if actual simulations exceed these ranges.

\sphinxAtStartPar
The next ntypes*(ntypes+1)/2 lines contain two numbers:
\begin{sphinxalltt}
cut xi(1)
cut xi(2)
...
cut xi(ntypes*(ntypes+1)/2)
\end{sphinxalltt}

\sphinxAtStartPar
Here cut means the cutoff distance of the pair functions, “xi” is
\(\xi\) as defined in the potential functions above. The
ntypes*(ntypes+1)/2 lines are related to the pairs according to the
sequence of first ii (self) pairs, i = 1, 2, …, ntypes, and then ij
(cross) pairs, i = 1, 2, …, ntypes\sphinxhyphen{}1, and j = i+1, i+2, …, ntypes
(i.e., the sequence of the ij pairs follows 11, 22, …, 12, 13, 14,
…, 23, 24, …).

\sphinxAtStartPar
In the final blocks of the potential file, U, V, W, P, G, and F
functions are listed sequentially. First, U functions are given for each
of the ntypes*(ntypes+1)/2 pairs according to the sequence described
above. For each of the pairs, nr values are listed. Next, similar arrays
are given for V and W functions. If P functions depend only on pair
species, i.e., \(\eta \neq 3\), then P functions are also listed the
same way the next. If P functions depend on three species, i.e.,
\(\eta = 3\), then P functions are listed for all the
ntypes*ntypes*ntypes IJK triplets in a natural sequence I from 1 to
ntypes, J from 1 to ntypes, and K from 1 to ntypes (i.e., IJK = 111,
112, 113, …, 121, 122, 123 …, 211, 212, …). Next, G functions are
listed for all the ntypes*ntypes*ntypes IJK triplets similarly. For each
of the G functions, ntheta values are listed. Finally, F functions are
listed for all the ntypes*(ntypes+1)/2 pairs in the same sequence as
described above.  For each of the F functions, nx values are listed.


\subsection{Mixing, shift, table, tail correction, restart, rRESPA info}
\label{\detokenize{pair_polymorphic:mixing-shift-table-tail-correction-restart-rrespa-info}}
\sphinxAtStartPar
This pair style does not support the {\hyperref[\detokenize{pair_modify::doc}]{\sphinxcrossref{\DUrole{doc}{pair\_modify}}}}
shift, table, and tail options.

\sphinxAtStartPar
This pair style does not write their information to {\hyperref[\detokenize{restart::doc}]{\sphinxcrossref{\DUrole{doc}{binary restart
files}}}}, since it is stored in potential files. Thus, you need
to re\sphinxhyphen{}specify the pair\_style and pair\_coeff commands in an input script
that reads a restart file.


\bigskip\hrule\bigskip



\subsection{Restrictions}
\label{\detokenize{pair_polymorphic:restrictions}}
\sphinxAtStartPar
If using create\_atoms command, atomic masses must be defined in the
input script. If using read\_data, atomic masses must be defined in the
atomic structure data file.

\sphinxAtStartPar
This pair style is part of the MANYBODY package. It is only enabled if
LAMMPS was built with that package. See the {\hyperref[\detokenize{Build_package::doc}]{\sphinxcrossref{\DUrole{doc}{Build package}}}} page for more info.

\sphinxAtStartPar
This pair potential requires the {\hyperref[\detokenize{newton::doc}]{\sphinxcrossref{\DUrole{doc}{newton}}}} setting to be
“on” for pair interactions.

\sphinxAtStartPar
The potential files provided with LAMMPS (see the potentials directory)
are parameterized for metal {\hyperref[\detokenize{units::doc}]{\sphinxcrossref{\DUrole{doc}{units}}}}. You can use any LAMMPS
units, but you would need to create your own potential files.


\subsection{Related commands}
\label{\detokenize{pair_polymorphic:related-commands}}
\sphinxAtStartPar
{\hyperref[\detokenize{pair_coeff::doc}]{\sphinxcrossref{\DUrole{doc}{pair\_coeff}}}}


\bigskip\hrule\bigskip

\phantomsection\label{\detokenize{pair_polymorphic:zhou3}}
\sphinxAtStartPar
\sphinxstylestrong{(Zhou3)} X. W. Zhou, M. E. Foster, R. E. Jones, P. Yang, H. Fan, and F. P. Doty, J. Mater. Sci. Res., 4, 15 (2015).

\phantomsection\label{\detokenize{pair_polymorphic:zhou4}}
\sphinxAtStartPar
\sphinxstylestrong{(Zhou4)} X. W. Zhou, M. E. Foster, J. A. Ronevich, and C. W. San Marchi, J. Comp. Chem., 41, 1299 (2020).

\phantomsection\label{\detokenize{pair_polymorphic:sw}}
\sphinxAtStartPar
\sphinxstylestrong{(SW)} F. H. Stillinger, and T. A. Weber, Phys. Rev. B, 31, 5262 (1985).

\phantomsection\label{\detokenize{pair_polymorphic:tersoff}}
\sphinxAtStartPar
\sphinxstylestrong{(Tersoff)} J. Tersoff, Phys. Rev. B, 39, 5566 (1989).

\phantomsection\label{\detokenize{pair_polymorphic:poly-albe}}
\sphinxAtStartPar
\sphinxstylestrong{(Albe1)} K. Albe, K. Nordlund, J. Nord, and A. Kuronen, Phys. Rev. B, 66, 035205 (2002).

\phantomsection\label{\detokenize{pair_polymorphic:wang3}}
\sphinxAtStartPar
\sphinxstylestrong{(Wang)} J. Wang, and A. Rockett, Phys. Rev. B, 43, 12571 (1991).

\phantomsection\label{\detokenize{pair_polymorphic:poly-daw}}
\sphinxAtStartPar
\sphinxstylestrong{(Daw)} M. S. Daw, and M. I. Baskes, Phys. Rev. B, 29, 6443 (1984).

\sphinxstepscope

\index{pair\_style python@\spxentry{pair\_style python}}\ignorespaces 

\section{pair\_style python command}
\label{\detokenize{pair_python:pair-style-python-command}}\label{\detokenize{pair_python:index-0}}\label{\detokenize{pair_python::doc}}

\subsection{Syntax}
\label{\detokenize{pair_python:syntax}}
\begin{sphinxVerbatim}[commandchars=\\\{\}]
\PYG{k}{pair\PYGZus{}style}\PYG{+w}{ }\PYG{n}{python}\PYG{+w}{ }\PYG{n}{cutoff}
\end{sphinxVerbatim}

\sphinxAtStartPar
cutoff = global cutoff for interactions in python potential classes


\subsection{Examples}
\label{\detokenize{pair_python:examples}}
\begin{sphinxVerbatim}[commandchars=\\\{\}]
\PYG{k}{pair\PYGZus{}style}\PYG{+w}{ }\PYG{n}{python}\PYG{+w}{ }\PYG{l+m}{2.5}
\PYG{k}{pair\PYGZus{}coeff}\PYG{+w}{ }\PYG{o}{*}\PYG{+w}{ }\PYG{o}{*}\PYG{+w}{ }\PYG{n}{py\PYGZus{}pot.LJCutMelt}\PYG{+w}{ }\PYG{n}{lj}

\PYG{k}{pair\PYGZus{}style}\PYG{+w}{ }\PYG{n}{python}\PYG{+w}{ }\PYG{l+m}{10.0}
\PYG{k}{pair\PYGZus{}coeff}\PYG{+w}{ }\PYG{o}{*}\PYG{+w}{ }\PYG{o}{*}\PYG{+w}{ }\PYG{n}{py\PYGZus{}pot.HarmonicCut}\PYG{+w}{ }\PYG{n}{A}\PYG{+w}{ }\PYG{n}{B}

\PYG{k}{pair\PYGZus{}style}\PYG{+w}{ }\PYG{n}{hybrid}\PYG{o}{/}\PYG{n}{overlay}\PYG{+w}{ }\PYG{n}{coul}\PYG{o}{/}\PYG{n}{long}\PYG{+w}{ }\PYG{l+m}{12.0}\PYG{+w}{ }\PYG{n}{python}\PYG{+w}{ }\PYG{l+m}{12.0}
\PYG{k}{pair\PYGZus{}coeff}\PYG{+w}{ }\PYG{o}{*}\PYG{+w}{ }\PYG{o}{*}\PYG{+w}{ }\PYG{n}{coul}\PYG{o}{/}\PYG{n}{long}
\PYG{k}{pair\PYGZus{}coeff}\PYG{+w}{ }\PYG{o}{*}\PYG{+w}{ }\PYG{o}{*}\PYG{+w}{ }\PYG{n}{python}\PYG{+w}{ }\PYG{n}{py\PYGZus{}pot.LJCutSPCE}\PYG{+w}{ }\PYG{n}{OW}\PYG{+w}{ }\PYG{n}{NULL}
\end{sphinxVerbatim}


\subsection{Description}
\label{\detokenize{pair_python:description}}
\sphinxAtStartPar
The \sphinxstyleemphasis{python} pair style provides a way to define pairwise additive
potential functions as python script code that is loaded into LAMMPS
from a python file which must contain specific python class definitions.
This allows to rapidly evaluate different potential functions without
having to modify and re\sphinxhyphen{}compile LAMMPS. Due to python being an
interpreted language, however, the performance of this pair style is
going to be significantly slower (often between 20x and 100x) than
corresponding compiled code. This penalty can be significantly reduced
through generating tabulations from the python code through the
{\hyperref[\detokenize{pair_write::doc}]{\sphinxcrossref{\DUrole{doc}{pair\_write}}}} command, which is supported by this style.

\sphinxAtStartPar
Only a single {\hyperref[\detokenize{pair_coeff::doc}]{\sphinxcrossref{\DUrole{doc}{pair\_coeff}}}} command is used with the
\sphinxstyleemphasis{python} pair style which specifies a python class inside a python module
or a file that LAMMPS will look up in the current directory, a folder
pointed to by the \sphinxcode{\sphinxupquote{LAMMPS\_POTENTIALS}} environment variable or somewhere
in your python path.  A single python module can hold multiple python pair
class definitions.  The class definitions itself have to follow specific
rules that are explained below.

\sphinxAtStartPar
Atom types in the python class are specified through symbolic
constants, typically strings. These are mapped to LAMMPS atom types by
specifying N additional arguments after the class name in the
pair\_coeff command, where N must be the number of currently defined
atom types:

\sphinxAtStartPar
As an example, imagine a file \sphinxstyleemphasis{py\_pot.py} has a python potential class
names \sphinxstyleemphasis{LJCutMelt} with parameters and potential functions for a two
Lennard\sphinxhyphen{}Jones atom types labeled as ‘LJ1’ and ‘LJ2’. In your LAMMPS
input and you would have defined 3 atom types, out of which the first
two are supposed to be using the ‘LJ1’ parameters and the third the
‘LJ2’ parameters, then you would use the following pair\_coeff command:

\begin{sphinxVerbatim}[commandchars=\\\{\}]
\PYG{k}{pair\PYGZus{}coeff}\PYG{+w}{ }\PYG{o}{*}\PYG{+w}{ }\PYG{o}{*}\PYG{+w}{ }\PYG{n}{py\PYGZus{}pot.LJCutMelt}\PYG{+w}{ }\PYG{n}{LJ1}\PYG{+w}{ }\PYG{n}{LJ1}\PYG{+w}{ }\PYG{n}{LJ2}
\end{sphinxVerbatim}

\sphinxAtStartPar
The first two arguments \sphinxstylestrong{must} be * * so as to span all LAMMPS atom
types.  The first two LJ1 arguments map LAMMPS atom types 1 and 2 to
the LJ1 atom type in the LJCutMelt class of the py\_pot.py file.  The
final LJ2 argument maps LAMMPS atom type 3 to the LJ2 atom type the
python file.  If a mapping value is specified as NULL, the mapping is
not performed, any pair interaction with this atom type will be
skipped. This can be used when a \sphinxstyleemphasis{python} potential is used as part of
the \sphinxstyleemphasis{hybrid} or \sphinxstyleemphasis{hybrid/overlay} pair style. The NULL values are then
placeholders for atom types that will be used with other potentials.


\bigskip\hrule\bigskip


\sphinxAtStartPar
The python potential file has to start with the following code:

\begin{sphinxVerbatim}[commandchars=\\\{\}]
\PYG{k+kn}{from}\PYG{+w}{ }\PYG{n+nn}{\PYGZus{}\PYGZus{}future\PYGZus{}\PYGZus{}}\PYG{+w}{ }\PYG{k+kn}{import} \PYG{n}{print\PYGZus{}function}

\PYG{k}{class}\PYG{+w}{ }\PYG{n+nc}{LAMMPSPairPotential}\PYG{p}{(}\PYG{n+nb}{object}\PYG{p}{)}\PYG{p}{:}
    \PYG{k}{def}\PYG{+w}{ }\PYG{n+nf+fm}{\PYGZus{}\PYGZus{}init\PYGZus{}\PYGZus{}}\PYG{p}{(}\PYG{n+nb+bp}{self}\PYG{p}{)}\PYG{p}{:}
        \PYG{n+nb+bp}{self}\PYG{o}{.}\PYG{n}{pmap}\PYG{o}{=}\PYG{n+nb}{dict}\PYG{p}{(}\PYG{p}{)}
        \PYG{n+nb+bp}{self}\PYG{o}{.}\PYG{n}{units}\PYG{o}{=}\PYG{l+s+s1}{\PYGZsq{}}\PYG{l+s+s1}{lj}\PYG{l+s+s1}{\PYGZsq{}}
    \PYG{k}{def}\PYG{+w}{ }\PYG{n+nf}{map\PYGZus{}coeff}\PYG{p}{(}\PYG{n+nb+bp}{self}\PYG{p}{,}\PYG{n}{name}\PYG{p}{,}\PYG{n}{ltype}\PYG{p}{)}\PYG{p}{:}
        \PYG{n+nb+bp}{self}\PYG{o}{.}\PYG{n}{pmap}\PYG{p}{[}\PYG{n}{ltype}\PYG{p}{]}\PYG{o}{=}\PYG{n}{name}
    \PYG{k}{def}\PYG{+w}{ }\PYG{n+nf}{check\PYGZus{}units}\PYG{p}{(}\PYG{n+nb+bp}{self}\PYG{p}{,}\PYG{n}{units}\PYG{p}{)}\PYG{p}{:}
        \PYG{k}{if} \PYG{p}{(}\PYG{n}{units} \PYG{o}{!=} \PYG{n+nb+bp}{self}\PYG{o}{.}\PYG{n}{units}\PYG{p}{)}\PYG{p}{:}
           \PYG{k}{raise} \PYG{n+ne}{Exception}\PYG{p}{(}\PYG{l+s+s2}{\PYGZdq{}}\PYG{l+s+s2}{Conflicting units: }\PYG{l+s+si}{\PYGZpc{}s}\PYG{l+s+s2}{ vs. }\PYG{l+s+si}{\PYGZpc{}s}\PYG{l+s+s2}{\PYGZdq{}} \PYG{o}{\PYGZpc{}} \PYG{p}{(}\PYG{n+nb+bp}{self}\PYG{o}{.}\PYG{n}{units}\PYG{p}{,}\PYG{n}{units}\PYG{p}{)}\PYG{p}{)}
\end{sphinxVerbatim}

\sphinxAtStartPar
Any classes with definitions of specific potentials have to be derived
from this class and should be initialize in a similar fashion to the
example given below.

\begin{sphinxadmonition}{note}{Note:}
\sphinxAtStartPar
The class constructor has to set up a data structure containing
the potential parameters supported by this class.  It should also
define a variable \sphinxstyleemphasis{self.units} containing a string matching one of the
options of LAMMPS’ {\hyperref[\detokenize{units::doc}]{\sphinxcrossref{\DUrole{doc}{units}}}} command, which is used to
verify, that the potential definition in the python class and in the
LAMMPS input match.
\end{sphinxadmonition}

\sphinxAtStartPar
Here is an example for a single type Lennard\sphinxhyphen{}Jones potential class
\sphinxstyleemphasis{LJCutMelt} in reduced units, which defines an atom type \sphinxstyleemphasis{lj} for
which the parameters epsilon and sigma are both 1.0:

\begin{sphinxVerbatim}[commandchars=\\\{\}]
\PYG{k}{class}\PYG{+w}{ }\PYG{n+nc}{LJCutMelt}\PYG{p}{(}\PYG{n}{LAMMPSPairPotential}\PYG{p}{)}\PYG{p}{:}
    \PYG{k}{def}\PYG{+w}{ }\PYG{n+nf+fm}{\PYGZus{}\PYGZus{}init\PYGZus{}\PYGZus{}}\PYG{p}{(}\PYG{n+nb+bp}{self}\PYG{p}{)}\PYG{p}{:}
        \PYG{n+nb}{super}\PYG{p}{(}\PYG{n}{LJCutMelt}\PYG{p}{,}\PYG{n+nb+bp}{self}\PYG{p}{)}\PYG{o}{.}\PYG{n+nf+fm}{\PYGZus{}\PYGZus{}init\PYGZus{}\PYGZus{}}\PYG{p}{(}\PYG{p}{)}
        \PYG{c+c1}{\PYGZsh{} set coeffs: 48*eps*sig**12, 24*eps*sig**6,}
        \PYG{c+c1}{\PYGZsh{}              4*eps*sig**12,  4*eps*sig**6}
        \PYG{n+nb+bp}{self}\PYG{o}{.}\PYG{n}{units} \PYG{o}{=} \PYG{l+s+s1}{\PYGZsq{}}\PYG{l+s+s1}{lj}\PYG{l+s+s1}{\PYGZsq{}}
        \PYG{n+nb+bp}{self}\PYG{o}{.}\PYG{n}{coeff} \PYG{o}{=} \PYG{p}{\PYGZob{}}\PYG{l+s+s1}{\PYGZsq{}}\PYG{l+s+s1}{lj}\PYG{l+s+s1}{\PYGZsq{}}  \PYG{p}{:} \PYG{p}{\PYGZob{}}\PYG{l+s+s1}{\PYGZsq{}}\PYG{l+s+s1}{lj}\PYG{l+s+s1}{\PYGZsq{}}  \PYG{p}{:} \PYG{p}{(}\PYG{l+m+mf}{48.0}\PYG{p}{,}\PYG{l+m+mf}{24.0}\PYG{p}{,}\PYG{l+m+mf}{4.0}\PYG{p}{,}\PYG{l+m+mf}{4.0}\PYG{p}{)}\PYG{p}{\PYGZcb{}}\PYG{p}{\PYGZcb{}}
\end{sphinxVerbatim}

\sphinxAtStartPar
The class also has to provide two methods for the computation of the
potential energy and forces, which have be named \sphinxstyleemphasis{compute\_force},
and \sphinxstyleemphasis{compute\_energy}, which both take 3 numerical arguments:
\begin{itemize}
\item {} 
\sphinxAtStartPar
rsq   = the square of the distance between a pair of atoms (float)

\item {} 
\sphinxAtStartPar
itype = the (numerical) type of the first atom

\item {} 
\sphinxAtStartPar
jtype = the (numerical) type of the second atom

\end{itemize}

\sphinxAtStartPar
This functions need to compute the (scaled) force and the energy,
respectively, and use the result as return value. The functions need
to use the \sphinxstyleemphasis{pmap} dictionary to convert the LAMMPS atom type number
to the symbolic value of the internal potential parameter data structure.
Following the \sphinxstyleemphasis{LJCutMelt} example, here are the two functions:

\begin{sphinxVerbatim}[commandchars=\\\{\}]
\PYG{k}{def}\PYG{+w}{ }\PYG{n+nf}{compute\PYGZus{}force}\PYG{p}{(}\PYG{n+nb+bp}{self}\PYG{p}{,}\PYG{n}{rsq}\PYG{p}{,}\PYG{n}{itype}\PYG{p}{,}\PYG{n}{jtype}\PYG{p}{)}\PYG{p}{:}
     \PYG{n}{coeff} \PYG{o}{=} \PYG{n+nb+bp}{self}\PYG{o}{.}\PYG{n}{coeff}\PYG{p}{[}\PYG{n+nb+bp}{self}\PYG{o}{.}\PYG{n}{pmap}\PYG{p}{[}\PYG{n}{itype}\PYG{p}{]}\PYG{p}{]}\PYG{p}{[}\PYG{n+nb+bp}{self}\PYG{o}{.}\PYG{n}{pmap}\PYG{p}{[}\PYG{n}{jtype}\PYG{p}{]}\PYG{p}{]}
     \PYG{n}{r2inv}  \PYG{o}{=} \PYG{l+m+mf}{1.0}\PYG{o}{/}\PYG{n}{rsq}
     \PYG{n}{r6inv}  \PYG{o}{=} \PYG{n}{r2inv}\PYG{o}{*}\PYG{n}{r2inv}\PYG{o}{*}\PYG{n}{r2inv}
     \PYG{n}{lj1} \PYG{o}{=} \PYG{n}{coeff}\PYG{p}{[}\PYG{l+m+mi}{0}\PYG{p}{]}
     \PYG{n}{lj2} \PYG{o}{=} \PYG{n}{coeff}\PYG{p}{[}\PYG{l+m+mi}{1}\PYG{p}{]}
     \PYG{k}{return} \PYG{p}{(}\PYG{n}{r6inv} \PYG{o}{*} \PYG{p}{(}\PYG{n}{lj1}\PYG{o}{*}\PYG{n}{r6inv} \PYG{o}{\PYGZhy{}} \PYG{n}{lj2}\PYG{p}{)}\PYG{p}{)}\PYG{o}{*}\PYG{n}{r2inv}

 \PYG{k}{def}\PYG{+w}{ }\PYG{n+nf}{compute\PYGZus{}energy}\PYG{p}{(}\PYG{n+nb+bp}{self}\PYG{p}{,}\PYG{n}{rsq}\PYG{p}{,}\PYG{n}{itype}\PYG{p}{,}\PYG{n}{jtype}\PYG{p}{)}\PYG{p}{:}
     \PYG{n}{coeff} \PYG{o}{=} \PYG{n+nb+bp}{self}\PYG{o}{.}\PYG{n}{coeff}\PYG{p}{[}\PYG{n+nb+bp}{self}\PYG{o}{.}\PYG{n}{pmap}\PYG{p}{[}\PYG{n}{itype}\PYG{p}{]}\PYG{p}{]}\PYG{p}{[}\PYG{n+nb+bp}{self}\PYG{o}{.}\PYG{n}{pmap}\PYG{p}{[}\PYG{n}{jtype}\PYG{p}{]}\PYG{p}{]}
     \PYG{n}{r2inv}  \PYG{o}{=} \PYG{l+m+mf}{1.0}\PYG{o}{/}\PYG{n}{rsq}
     \PYG{n}{r6inv}  \PYG{o}{=} \PYG{n}{r2inv}\PYG{o}{*}\PYG{n}{r2inv}\PYG{o}{*}\PYG{n}{r2inv}
     \PYG{n}{lj3} \PYG{o}{=} \PYG{n}{coeff}\PYG{p}{[}\PYG{l+m+mi}{2}\PYG{p}{]}
     \PYG{n}{lj4} \PYG{o}{=} \PYG{n}{coeff}\PYG{p}{[}\PYG{l+m+mi}{3}\PYG{p}{]}
     \PYG{k}{return} \PYG{p}{(}\PYG{n}{r6inv} \PYG{o}{*} \PYG{p}{(}\PYG{n}{lj3}\PYG{o}{*}\PYG{n}{r6inv} \PYG{o}{\PYGZhy{}} \PYG{n}{lj4}\PYG{p}{)}\PYG{p}{)}
\end{sphinxVerbatim}

\begin{sphinxadmonition}{note}{Note:}
\sphinxAtStartPar
for consistency with the C++ pair styles in LAMMPS, the
\sphinxstyleemphasis{compute\_force} function follows the conventions of the Pair::single()
methods and does not return the pairwise force directly, but the force
divided by the distance between the two atoms, so this value only needs
to be  multiplied by delta x, delta y, and delta z to conveniently obtain
the three components of the force vector between these two atoms.
\end{sphinxadmonition}

\sphinxAtStartPar
Below is a more complex example using \sphinxstyleemphasis{real} units and defines an interaction
equivalent to:

\begin{sphinxVerbatim}[commandchars=\\\{\}]
\PYG{k}{units}\PYG{+w}{ }\PYG{n}{real}
\PYG{k}{pair\PYGZus{}style}\PYG{+w}{ }\PYG{n}{harmonic}\PYG{o}{/}\PYG{n}{cut}
\PYG{k}{pair\PYGZus{}coeff}\PYG{+w}{ }\PYG{l+m}{1}\PYG{+w}{ }\PYG{l+m}{1}\PYG{+w}{ }\PYG{l+m}{0.2}\PYG{+w}{ }\PYG{l+m}{9.0}
\PYG{k}{pair\PYGZus{}coeff}\PYG{+w}{ }\PYG{l+m}{2}\PYG{+w}{ }\PYG{l+m}{2}\PYG{+w}{ }\PYG{l+m}{0.4}\PYG{+w}{ }\PYG{l+m}{9.0}
\end{sphinxVerbatim}

\sphinxAtStartPar
This uses the default geometric mixing.  The equivalent input with pair
style \sphinxstyleemphasis{python} is:

\begin{sphinxVerbatim}[commandchars=\\\{\}]
\PYG{k}{units}\PYG{+w}{ }\PYG{n}{real}
\PYG{k}{pair\PYGZus{}style}\PYG{+w}{ }\PYG{n}{python}\PYG{+w}{ }\PYG{l+m}{10.0}
\PYG{k}{pair\PYGZus{}coeff}\PYG{+w}{ }\PYG{o}{*}\PYG{+w}{ }\PYG{o}{*}\PYG{+w}{ }\PYG{n}{py\PYGZus{}pot.Harmonic}\PYG{+w}{ }\PYG{n}{A}\PYG{+w}{ }\PYG{n}{B}
\end{sphinxVerbatim}

\sphinxAtStartPar
Note that while for pair style \sphinxstyleemphasis{harmonic/cut} the cutoff is implicitly
set to the minimum of the harmonic potential, for pair style python a
global cutoff must be set and it must be equal or larger to the implicit
cutoff of the potential in python, which has to explicitly return zero
force and energy beyond the cutoff.  Also, the mixed parameters have to
be explicitly provided.  The corresponding python code is:

\begin{sphinxVerbatim}[commandchars=\\\{\}]
\PYG{k}{class}\PYG{+w}{ }\PYG{n+nc}{Harmonic}\PYG{p}{(}\PYG{n}{LAMMPSPairPotential}\PYG{p}{)}\PYG{p}{:}
    \PYG{k}{def}\PYG{+w}{ }\PYG{n+nf+fm}{\PYGZus{}\PYGZus{}init\PYGZus{}\PYGZus{}}\PYG{p}{(}\PYG{n+nb+bp}{self}\PYG{p}{)}\PYG{p}{:}
        \PYG{n+nb}{super}\PYG{p}{(}\PYG{n}{Harmonic}\PYG{p}{,}\PYG{n+nb+bp}{self}\PYG{p}{)}\PYG{o}{.}\PYG{n+nf+fm}{\PYGZus{}\PYGZus{}init\PYGZus{}\PYGZus{}}\PYG{p}{(}\PYG{p}{)}
        \PYG{n+nb+bp}{self}\PYG{o}{.}\PYG{n}{units} \PYG{o}{=} \PYG{l+s+s1}{\PYGZsq{}}\PYG{l+s+s1}{real}\PYG{l+s+s1}{\PYGZsq{}}
        \PYG{c+c1}{\PYGZsh{} set coeffs: K, r0}
        \PYG{n+nb+bp}{self}\PYG{o}{.}\PYG{n}{coeff} \PYG{o}{=} \PYG{p}{\PYGZob{}}\PYG{l+s+s1}{\PYGZsq{}}\PYG{l+s+s1}{A}\PYG{l+s+s1}{\PYGZsq{}}  \PYG{p}{:} \PYG{p}{\PYGZob{}}\PYG{l+s+s1}{\PYGZsq{}}\PYG{l+s+s1}{A}\PYG{l+s+s1}{\PYGZsq{}}  \PYG{p}{:} \PYG{p}{(}\PYG{l+m+mf}{0.2}\PYG{p}{,}\PYG{l+m+mf}{9.0}\PYG{p}{)}\PYG{p}{,}
                              \PYG{l+s+s1}{\PYGZsq{}}\PYG{l+s+s1}{B}\PYG{l+s+s1}{\PYGZsq{}}  \PYG{p}{:} \PYG{p}{(}\PYG{n}{math}\PYG{o}{.}\PYG{n}{sqrt}\PYG{p}{(}\PYG{l+m+mf}{0.2}\PYG{o}{*}\PYG{l+m+mf}{0.4}\PYG{p}{)}\PYG{p}{,}\PYG{l+m+mf}{9.0}\PYG{p}{)}\PYG{p}{\PYGZcb{}}\PYG{p}{,}
                      \PYG{l+s+s1}{\PYGZsq{}}\PYG{l+s+s1}{B}\PYG{l+s+s1}{\PYGZsq{}}  \PYG{p}{:} \PYG{p}{\PYGZob{}}\PYG{l+s+s1}{\PYGZsq{}}\PYG{l+s+s1}{A}\PYG{l+s+s1}{\PYGZsq{}}  \PYG{p}{:} \PYG{p}{(}\PYG{n}{math}\PYG{o}{.}\PYG{n}{sqrt}\PYG{p}{(}\PYG{l+m+mf}{0.2}\PYG{o}{*}\PYG{l+m+mf}{0.4}\PYG{p}{)}\PYG{p}{,}\PYG{l+m+mf}{9.0}\PYG{p}{)}\PYG{p}{,}
                              \PYG{l+s+s1}{\PYGZsq{}}\PYG{l+s+s1}{B}\PYG{l+s+s1}{\PYGZsq{}}  \PYG{p}{:} \PYG{p}{(}\PYG{l+m+mf}{0.4}\PYG{p}{,}\PYG{l+m+mf}{9.0}\PYG{p}{)}\PYG{p}{\PYGZcb{}}\PYG{p}{\PYGZcb{}}

    \PYG{k}{def}\PYG{+w}{ }\PYG{n+nf}{compute\PYGZus{}force}\PYG{p}{(}\PYG{n+nb+bp}{self}\PYG{p}{,}\PYG{n}{rsq}\PYG{p}{,}\PYG{n}{itype}\PYG{p}{,}\PYG{n}{jtype}\PYG{p}{)}\PYG{p}{:}
        \PYG{n}{coeff} \PYG{o}{=} \PYG{n+nb+bp}{self}\PYG{o}{.}\PYG{n}{coeff}\PYG{p}{[}\PYG{n+nb+bp}{self}\PYG{o}{.}\PYG{n}{pmap}\PYG{p}{[}\PYG{n}{itype}\PYG{p}{]}\PYG{p}{]}\PYG{p}{[}\PYG{n+nb+bp}{self}\PYG{o}{.}\PYG{n}{pmap}\PYG{p}{[}\PYG{n}{jtype}\PYG{p}{]}\PYG{p}{]}
        \PYG{n}{r} \PYG{o}{=} \PYG{n}{math}\PYG{o}{.}\PYG{n}{sqrt}\PYG{p}{(}\PYG{n}{rsq}\PYG{p}{)}
        \PYG{n}{delta} \PYG{o}{=} \PYG{n}{coeff}\PYG{p}{[}\PYG{l+m+mi}{1}\PYG{p}{]}\PYG{o}{\PYGZhy{}}\PYG{n}{r}
        \PYG{k}{if} \PYG{p}{(}\PYG{n}{r} \PYG{o}{\PYGZlt{}} \PYG{n}{coeff}\PYG{p}{[}\PYG{l+m+mi}{1}\PYG{p}{]}\PYG{p}{)}\PYG{p}{:}
            \PYG{k}{return} \PYG{l+m+mf}{2.0}\PYG{o}{*}\PYG{n}{delta}\PYG{o}{*}\PYG{n}{coeff}\PYG{p}{[}\PYG{l+m+mi}{0}\PYG{p}{]}\PYG{o}{/}\PYG{n}{r}
        \PYG{k}{else}\PYG{p}{:}
            \PYG{k}{return} \PYG{l+m+mf}{0.0}

    \PYG{k}{def}\PYG{+w}{ }\PYG{n+nf}{compute\PYGZus{}energy}\PYG{p}{(}\PYG{n+nb+bp}{self}\PYG{p}{,}\PYG{n}{rsq}\PYG{p}{,}\PYG{n}{itype}\PYG{p}{,}\PYG{n}{jtype}\PYG{p}{)}\PYG{p}{:}
        \PYG{n}{coeff} \PYG{o}{=} \PYG{n+nb+bp}{self}\PYG{o}{.}\PYG{n}{coeff}\PYG{p}{[}\PYG{n+nb+bp}{self}\PYG{o}{.}\PYG{n}{pmap}\PYG{p}{[}\PYG{n}{itype}\PYG{p}{]}\PYG{p}{]}\PYG{p}{[}\PYG{n+nb+bp}{self}\PYG{o}{.}\PYG{n}{pmap}\PYG{p}{[}\PYG{n}{jtype}\PYG{p}{]}\PYG{p}{]}
        \PYG{n}{r} \PYG{o}{=} \PYG{n}{math}\PYG{o}{.}\PYG{n}{sqrt}\PYG{p}{(}\PYG{n}{rsq}\PYG{p}{)}
        \PYG{n}{delta} \PYG{o}{=} \PYG{n}{coeff}\PYG{p}{[}\PYG{l+m+mi}{1}\PYG{p}{]}\PYG{o}{\PYGZhy{}}\PYG{n}{r}
        \PYG{k}{if} \PYG{p}{(}\PYG{n}{r} \PYG{o}{\PYGZlt{}} \PYG{n}{coeff}\PYG{p}{[}\PYG{l+m+mi}{1}\PYG{p}{]}\PYG{p}{)}\PYG{p}{:}
            \PYG{k}{return} \PYG{n}{delta}\PYG{o}{*}\PYG{n}{delta}\PYG{o}{*}\PYG{n}{coeff}\PYG{p}{[}\PYG{l+m+mi}{0}\PYG{p}{]}
        \PYG{k}{else}\PYG{p}{:}
            \PYG{k}{return} \PYG{l+m+mf}{0.0}
\end{sphinxVerbatim}


\bigskip\hrule\bigskip


\begin{sphinxadmonition}{note}{Performance Impact}

\sphinxAtStartPar
The evaluation of scripted python code will slow down the computation
of pairwise interactions quite significantly. However, this performance
penalty can be worked around through using the python pair style not
for the actual simulation, but to generate tabulated potentials using
the {\hyperref[\detokenize{pair_write::doc}]{\sphinxcrossref{\DUrole{doc}{pair\_write}}}} command.  This will also enable
GPU or multi\sphinxhyphen{}thread acceleration through the GPU, KOKKOS, or OPENMP
package versions of the \sphinxstyleemphasis{table} pair style.  Please see below for a
LAMMPS input example demonstrating how to build a table file:
\end{sphinxadmonition}

\begin{sphinxVerbatim}[commandchars=\\\{\}]
\PYG{k}{pair\PYGZus{}style}\PYG{+w}{ }\PYG{n}{python}\PYG{+w}{ }\PYG{l+m}{2.5}
\PYG{k}{pair\PYGZus{}coeff}\PYG{+w}{ }\PYG{o}{*}\PYG{+w}{ }\PYG{o}{*}\PYG{+w}{ }\PYG{n}{py\PYGZus{}pot.LJCutMelt}\PYG{+w}{ }\PYG{n}{lj}
\PYG{k}{shell}\PYG{+w}{ }\PYG{n}{rm}\PYG{+w}{ }\PYG{o}{\PYGZhy{}}\PYG{n}{f}\PYG{+w}{ }\PYG{n}{lj.table}
\PYG{k}{pair\PYGZus{}write}\PYG{+w}{  }\PYG{l+m}{1}\PYG{+w}{ }\PYG{l+m}{1}\PYG{+w}{ }\PYG{l+m}{2000}\PYG{+w}{ }\PYG{n}{rsq}\PYG{+w}{ }\PYG{l+m}{0.01}\PYG{+w}{ }\PYG{l+m}{2.5}\PYG{+w}{ }\PYG{n}{lj.table}\PYG{+w}{ }\PYG{n}{lj}
\end{sphinxVerbatim}

\sphinxAtStartPar
Note that it is strongly recommended to try to \sphinxstylestrong{delete} the potential
table file before generating it. Since the \sphinxstyleemphasis{pair\_write} command will
always \sphinxstylestrong{append} to a table file, while pair style table will use the
\sphinxstylestrong{first match}. Thus when changing the potential function in the python
class, the table pair style will still read the old variant unless the
table file is first deleted.

\sphinxAtStartPar
After switching the pair style to \sphinxstyleemphasis{table}, the potential tables need
to be assigned to the LAMMPS atom types like this:

\begin{sphinxVerbatim}[commandchars=\\\{\}]
\PYG{k}{pair\PYGZus{}style}\PYG{+w}{      }\PYG{n}{table}\PYG{+w}{ }\PYG{n}{linear}\PYG{+w}{ }\PYG{l+m}{2000}
\PYG{k}{pair\PYGZus{}coeff}\PYG{+w}{      }\PYG{l+m}{1}\PYG{+w}{  }\PYG{l+m}{1}\PYG{+w}{  }\PYG{n}{lj.table}\PYG{+w}{ }\PYG{n}{lj}
\end{sphinxVerbatim}

\sphinxAtStartPar
This can also be done for more complex systems.  Please see the
\sphinxstyleemphasis{examples/python} folders for a few more examples.


\bigskip\hrule\bigskip



\subsection{Mixing, shift, table, tail correction, restart, rRESPA info}
\label{\detokenize{pair_python:mixing-shift-table-tail-correction-restart-rrespa-info}}
\sphinxAtStartPar
Mixing of potential parameters has to be handled inside the provided
python module. The python pair style simply assumes that force and
energy computation can be correctly performed for all pairs of atom
types as they are mapped to the atom type labels inside the python
potential class.

\sphinxAtStartPar
This pair style does not support the {\hyperref[\detokenize{pair_modify::doc}]{\sphinxcrossref{\DUrole{doc}{pair\_modify}}}}
shift, table, and tail options.

\sphinxAtStartPar
This pair style does not write its information to {\hyperref[\detokenize{restart::doc}]{\sphinxcrossref{\DUrole{doc}{binary restart files}}}}, since it is stored in potential files.  Thus, you
need to re\sphinxhyphen{}specify the pair\_style and pair\_coeff commands in an input
script that reads a restart file.

\sphinxAtStartPar
This pair style can only be used via the \sphinxstyleemphasis{pair} keyword of the
{\hyperref[\detokenize{run_style::doc}]{\sphinxcrossref{\DUrole{doc}{run\_style respa}}}} command.  It does not support the
\sphinxstyleemphasis{inner}, \sphinxstyleemphasis{middle}, \sphinxstyleemphasis{outer} keywords.


\bigskip\hrule\bigskip



\subsection{Restrictions}
\label{\detokenize{pair_python:restrictions}}
\sphinxAtStartPar
This pair style is part of the PYTHON package.  It is only enabled if
LAMMPS was built with that package.  See the {\hyperref[\detokenize{Build_package::doc}]{\sphinxcrossref{\DUrole{doc}{Build package}}}} page for more info.


\subsection{Related commands}
\label{\detokenize{pair_python:related-commands}}
\sphinxAtStartPar
{\hyperref[\detokenize{pair_coeff::doc}]{\sphinxcrossref{\DUrole{doc}{pair\_coeff}}}}, {\hyperref[\detokenize{pair_write::doc}]{\sphinxcrossref{\DUrole{doc}{pair\_write}}}},
{\hyperref[\detokenize{pair_table::doc}]{\sphinxcrossref{\DUrole{doc}{pair style table}}}}


\subsection{Default}
\label{\detokenize{pair_python:default}}
\sphinxAtStartPar
none

\sphinxstepscope

\index{pair\_style quip@\spxentry{pair\_style quip}}\ignorespaces 

\section{pair\_style quip command}
\label{\detokenize{pair_quip:pair-style-quip-command}}\label{\detokenize{pair_quip:index-0}}\label{\detokenize{pair_quip::doc}}

\subsection{Syntax}
\label{\detokenize{pair_quip:syntax}}
\begin{sphinxVerbatim}[commandchars=\\\{\}]
\PYG{k}{pair\PYGZus{}style}\PYG{+w}{ }\PYG{n}{quip}
\end{sphinxVerbatim}


\subsection{Examples}
\label{\detokenize{pair_quip:examples}}
\begin{sphinxVerbatim}[commandchars=\\\{\}]
\PYG{k}{pair\PYGZus{}style}\PYG{+w}{      }\PYG{n}{quip}
\PYG{k}{pair\PYGZus{}coeff}\PYG{+w}{      }\PYG{o}{*}\PYG{+w}{ }\PYG{o}{*}\PYG{+w}{ }\PYG{n}{gap\PYGZus{}example.xml}\PYG{+w}{ }\PYG{l+s}{\PYGZdq{}}\PYG{l+s}{Potential xml\PYGZus{}label=GAP\PYGZus{}2014\PYGZus{}5\PYGZus{}8\PYGZus{}60\PYGZus{}17\PYGZus{}10\PYGZus{}38\PYGZus{}466}\PYG{l+s}{\PYGZdq{}}\PYG{+w}{ }\PYG{l+m}{14}
\PYG{k}{pair\PYGZus{}coeff}\PYG{+w}{      }\PYG{o}{*}\PYG{+w}{ }\PYG{o}{*}\PYG{+w}{ }\PYG{n}{sw\PYGZus{}example.xml}\PYG{+w}{ }\PYG{l+s}{\PYGZdq{}}\PYG{l+s}{IP SW}\PYG{l+s}{\PYGZdq{}}\PYG{+w}{ }\PYG{l+m}{14}
\end{sphinxVerbatim}


\subsection{Description}
\label{\detokenize{pair_quip:description}}
\sphinxAtStartPar
Style \sphinxstyleemphasis{quip} provides an interface for calling potential routines from
the QUIP package. QUIP is built separately, and then linked to
LAMMPS. The most recent version of the QUIP package can be downloaded
from GitHub:
\sphinxurl{https://github.com/libAtoms/QUIP}. The
interface is chiefly intended to be used to run Gaussian Approximation
Potentials (GAP), which are described in the following publications:
{\hyperref[\detokenize{pair_quip:bartok-2010}]{\sphinxcrossref{\DUrole{std,std-ref}{(Bartok et al)}}}} and {\hyperref[\detokenize{pair_quip:bartok-phd}]{\sphinxcrossref{\DUrole{std,std-ref}{(PhD thesis of Bartok)}}}}.

\sphinxAtStartPar
Only a single pair\_coeff command is used with the \sphinxstyleemphasis{quip} style that
specifies a QUIP potential file containing the parameters of the
potential for all needed elements in XML format. This is followed by a
QUIP initialization string. Finally, the QUIP elements are mapped to
LAMMPS atom types by specifying N atomic numbers, where N is the
number of LAMMPS atom types:
\begin{itemize}
\item {} 
\sphinxAtStartPar
QUIP filename

\item {} 
\sphinxAtStartPar
QUIP initialization string

\item {} 
\sphinxAtStartPar
N atomic numbers = mapping of QUIP elements to atom types

\end{itemize}

\sphinxAtStartPar
See the {\hyperref[\detokenize{pair_coeff::doc}]{\sphinxcrossref{\DUrole{doc}{pair\_coeff}}}} page for alternate ways
to specify the path for the potential file.

\sphinxAtStartPar
A QUIP potential is fully specified by the filename which contains the
parameters of the potential in XML format, the initialization string,
and the map of atomic numbers.

\sphinxAtStartPar
GAP potentials can be obtained from the \sphinxhref{https://libatoms.github.io/GAP/data.html}{GAP models and databases page
on the libAtoms homepage \textasciigrave{}https://libatoms.github.io}, where the appropriate
initialization strings are also advised. The list of atomic numbers must
be matched to the LAMMPS atom types specified in the LAMMPS data file or
elsewhere.

\sphinxAtStartPar
Two examples input scripts are provided in the examples/PACKAGES/quip
directory.


\subsection{Mixing, shift, table, tail correction, restart, rRESPA info}
\label{\detokenize{pair_quip:mixing-shift-table-tail-correction-restart-rrespa-info}}
\sphinxAtStartPar
This pair style does not support the {\hyperref[\detokenize{pair_modify::doc}]{\sphinxcrossref{\DUrole{doc}{pair\_modify}}}}
mix, shift, table, and tail options.

\sphinxAtStartPar
This pair style does not write its information to {\hyperref[\detokenize{restart::doc}]{\sphinxcrossref{\DUrole{doc}{binary restart files}}}}, since it is stored in potential files.  Thus, you
need to re\sphinxhyphen{}specify the pair\_style and pair\_coeff commands in an input
script that reads a restart file.

\sphinxAtStartPar
This pair style can only be used via the \sphinxstyleemphasis{pair} keyword of the
{\hyperref[\detokenize{run_style::doc}]{\sphinxcrossref{\DUrole{doc}{run\_style respa}}}} command.  It does not support the
\sphinxstyleemphasis{inner}, \sphinxstyleemphasis{middle}, \sphinxstyleemphasis{outer} keywords.


\subsection{Restrictions}
\label{\detokenize{pair_quip:restrictions}}
\sphinxAtStartPar
This pair style is part of the ML\sphinxhyphen{}QUIP package.  It is only enabled
if LAMMPS was built with that package.  See the {\hyperref[\detokenize{Build_package::doc}]{\sphinxcrossref{\DUrole{doc}{Build package}}}} page for more info.

\sphinxAtStartPar
QUIP potentials are parameterized in electron\sphinxhyphen{}volts and Angstroms and
therefore should be used with LAMMPS metal {\hyperref[\detokenize{units::doc}]{\sphinxcrossref{\DUrole{doc}{units}}}}.

\sphinxAtStartPar
QUIP potentials are generally not designed to work with the scaling
factors set by the {\hyperref[\detokenize{special_bonds::doc}]{\sphinxcrossref{\DUrole{doc}{special\_bonds}}}} command.  The
recommended setting in molecular systems is to include all
interactions, i.e. to use \sphinxstyleemphasis{special\_bonds lj/coul 1.0 1.0 1.0}. Scaling
factors \textgreater{} 0.0 will be ignored and treated as 1.0. The only exception
to this rule is if you know that your QUIP potential needs to exclude
bonded, 1\sphinxhyphen{}3, or 1\sphinxhyphen{}4 interactions and does not already do this exclusion
within QUIP. Then a factor 0.0 needs to be used which will remove such
pairs from the neighbor list. This needs to be very carefully tested,
because it may remove pairs from the neighbor list that are still
required.


\subsection{Related commands}
\label{\detokenize{pair_quip:related-commands}}
\sphinxAtStartPar
{\hyperref[\detokenize{pair_coeff::doc}]{\sphinxcrossref{\DUrole{doc}{pair\_coeff}}}}


\bigskip\hrule\bigskip

\phantomsection\label{\detokenize{pair_quip:bartok-2010}}
\sphinxAtStartPar
\sphinxstylestrong{(Bartok2010)} AP Bartok, MC Payne, R Kondor, and G Csanyi, Physical
Review Letters 104, 136403 (2010).

\phantomsection\label{\detokenize{pair_quip:bartok-phd}}
\sphinxAtStartPar
\sphinxstylestrong{(Bartok\_PhD)} A Bartok\sphinxhyphen{}Partay, PhD Thesis, University of Cambridge,
(2010).

\sphinxstepscope

\index{pair\_style rann@\spxentry{pair\_style rann}}\ignorespaces 

\section{pair\_style rann command}
\label{\detokenize{pair_rann:pair-style-rann-command}}\label{\detokenize{pair_rann:index-0}}\label{\detokenize{pair_rann::doc}}

\subsection{Syntax}
\label{\detokenize{pair_rann:syntax}}
\begin{sphinxVerbatim}[commandchars=\\\{\}]
\PYG{k}{pair\PYGZus{}style}\PYG{+w}{ }\PYG{n}{rann}
\PYG{k}{pair\PYGZus{}coeff}\PYG{+w}{ }\PYG{n}{file}\PYG{+w}{ }\PYG{n}{Type1\PYGZus{}element}\PYG{+w}{ }\PYG{n}{Type2\PYGZus{}element}\PYG{+w}{ }\PYG{n}{Type3\PYGZus{}element...}
\end{sphinxVerbatim}


\subsection{Examples}
\label{\detokenize{pair_rann:examples}}
\begin{sphinxVerbatim}[commandchars=\\\{\}]
\PYG{k}{pair\PYGZus{}style}\PYG{+w}{ }\PYG{n}{rann}
\PYG{k}{pair\PYGZus{}coeff}\PYG{+w}{ }\PYG{o}{*}\PYG{+w}{ }\PYG{o}{*}\PYG{+w}{ }\PYG{n}{Mg.rann}\PYG{+w}{ }\PYG{n}{Mg}
\PYG{k}{pair\PYGZus{}coeff}\PYG{+w}{ }\PYG{o}{*}\PYG{+w}{ }\PYG{o}{*}\PYG{+w}{ }\PYG{n}{MgAlalloy.rann}\PYG{+w}{ }\PYG{n}{Mg}\PYG{+w}{ }\PYG{n}{Mg}\PYG{+w}{ }\PYG{n}{Al}\PYG{+w}{ }\PYG{n}{Mg}
\end{sphinxVerbatim}


\subsection{Description}
\label{\detokenize{pair_rann:description}}
\sphinxAtStartPar
Pair style \sphinxstyleemphasis{rann} computes pairwise interactions for a variety of
materials using rapid atomistic neural network (RANN) potentials
({\hyperref[\detokenize{pair_rann:dickel}]{\sphinxcrossref{\DUrole{std,std-ref}{Dickel}}}} , {\hyperref[\detokenize{pair_rann:nitol}]{\sphinxcrossref{\DUrole{std,std-ref}{Nitol}}}}).  Neural network
potentials work by first generating a series of symmetry functions
i.e. structural fingerprints from the neighbor list and then using these
values as the input layer of a neural network.  There is a single output
neuron in the final layer which is the energy.  Atomic forces are found
by analytical derivatives computed via back\sphinxhyphen{}propagation.  For alloy
systems, each element has a unique network.


\subsection{Potential file syntax}
\label{\detokenize{pair_rann:potential-file-syntax}}
\sphinxAtStartPar
The RANN potential is defined by a single text file which contains all
the fitting parameters for the alloy system.  The potential file also
defines the active fingerprints, network architecture, activation
functions, etc.  The potential file is divided into several sections
which are identified by one of the following keywords:
\begin{itemize}
\item {} 
\sphinxAtStartPar
atomtypes

\item {} 
\sphinxAtStartPar
mass

\item {} 
\sphinxAtStartPar
fingerprintsperelement

\item {} 
\sphinxAtStartPar
fingerprints

\item {} 
\sphinxAtStartPar
fingerprintconstants

\item {} 
\sphinxAtStartPar
screening (optional)

\item {} 
\sphinxAtStartPar
networklayers

\item {} 
\sphinxAtStartPar
layersize

\item {} 
\sphinxAtStartPar
weight

\item {} 
\sphinxAtStartPar
bias

\item {} 
\sphinxAtStartPar
activationfunctions

\item {} 
\sphinxAtStartPar
calibrationparameters (ignored)

\end{itemize}

\sphinxAtStartPar
The ‘\#’ character is treated as a comment marker, similar to LAMMPS
input scripts.  Sections are not required to follow a rigid ordering,
but do require previous definition of prerequisite information.  E.g.,
fingerprintconstants for a particular fingerprint must follow the
fingerprints definition; layersize for a particular layer must follow
the declaration of network layers.

\sphinxAtStartPar
\sphinxstyleemphasis{atomtypes} are defined as follows using element keywords separated by spaces.

\begin{sphinxVerbatim}[commandchars=\\\{\}]
atomtypes:
Fe Mg Al etc.
\end{sphinxVerbatim}

\sphinxAtStartPar
\sphinxstyleemphasis{mass} must be specified for each element keyword as follows:

\begin{sphinxVerbatim}[commandchars=\\\{\}]
mass:Mg:
24.305
mass:Fe:
55.847
mass:Al:
26.982
\end{sphinxVerbatim}

\sphinxAtStartPar
\sphinxstyleemphasis{fingerprintsperelement} specifies how many fingerprints are active for
computing the energy of a given atom.  This number must be specified for
each element keyword.  Active elements for each fingerprint depend upon
the type of the central atom and the neighboring atoms.  Pairwise
fingerprints may be defined for a Mg atom based exclusively on its Al
neighbors, for example.  Bond fingerprints may use two neighbor lists of
different element types.  In computing fingerprintsperelement from all
defined fingerprints, only the fingerprints defined for atoms of a
particular element should be considered, regardless of the elements used
in its neighbor list.  In the following code, for example, some
fingerprints may compute pairwise fingerprints summing contributions
about Fe atoms based on a neighbor list of exclusively Al atoms, but if
there are no fingerprints summing contributions of all neighbors about a
central Al atom, then fingerprintsperelement of Al is zero:

\begin{sphinxVerbatim}[commandchars=\\\{\}]
fingerprintsperelement:Mg:
5
fingerprintsperelement:Fe:
2
fingerprintsperelement:Al:
0
\end{sphinxVerbatim}

\sphinxAtStartPar
\sphinxstyleemphasis{fingerprints} specifies the active fingerprints for a certain element
combination.  Pair fingerprints are specified for two elements, while
bond fingerprints are specified for three elements.  Only one
fingerprints header should be used for an individual combination of
elements.  The ordering of the fingerprints in the network input layer
is determined by the order of element combinations specified by
subsequent \sphinxstyleemphasis{fingerprints} lines, and the order of the fingerprints
defined for each element combination.  Multiple fingerprints of the same
style or different ones may be specified.  If the same style and element
combination is used for multiple fingerprints, they should have
different id numbers.  The first element specifies the atoms for which
this fingerprint is computed while the other(s) specify which atoms to
use in the neighbor lists for the computation.  Switching the second and
third element type in bond fingerprints has no effect on the
computation:

\begin{sphinxVerbatim}[commandchars=\\\{\}]
fingerprints:Mg\PYGZus{}Mg:
radial\PYGZus{}0 radialscreened\PYGZus{}0 radial\PYGZus{}1
fingerprints:Mg\PYGZus{}Al\PYGZus{}Fe:
bond\PYGZus{}0 bondspin\PYGZus{}0
fingerprints:Mg\PYGZus{}Al:
radial\PYGZus{}0 radialscreened\PYGZus{}0
\end{sphinxVerbatim}

\sphinxAtStartPar
The following fingerprint styles are currently defined.  See the
{\hyperref[\detokenize{pair_rann:fingerprints}]{\sphinxcrossref{\DUrole{std,std-ref}{formulation section}}}} below for their definitions:
\begin{itemize}
\item {} 
\sphinxAtStartPar
radial

\item {} 
\sphinxAtStartPar
radialscreened

\item {} 
\sphinxAtStartPar
radialspin

\item {} 
\sphinxAtStartPar
radialscreenedspin

\item {} 
\sphinxAtStartPar
bond

\item {} 
\sphinxAtStartPar
bondscreened

\item {} 
\sphinxAtStartPar
bondspin

\item {} 
\sphinxAtStartPar
bondscreenedspin

\end{itemize}

\sphinxAtStartPar
\sphinxstyleemphasis{fingerprintconstants}  specifies the meta\sphinxhyphen{}parameters for a defined fingerprint.  For all radial styles, re, rc,
alpha, dr, o, and n must be specified.  re should usually be the stable interatomic distance, rc is the cutoff
radius, dr is the cutoff smoothing distance, o is the lowest radial power term (which may be negative), and n
is the highest power term.  The total length of the fingerprint vector is (n\sphinxhyphen{}o+1).  alpha is a list of decay parameters
used for exponential decay of radial contributions.  It may be set proportionally to the bulk modulus similarly
to MEAM potentials, but other values may provided better fitting in special cases.  Bond style fingerprints require
specification of re, rc, alphak, dr, k, and m.  Here m is the power of the bond cosines and k is the number of
decay parameters.  Cosine powers go from 0 to m\sphinxhyphen{}1 and are each computed for all values of alphak.  Thus the total
length of the fingerprint vector is m*k.

\begin{sphinxVerbatim}[commandchars=\\\{\}]
fingerprintconstants:Mg\PYGZus{}Mg:radialscreened\PYGZus{}0:re:
3.193592
fingerprintconstants:Mg\PYGZus{}Mg:radialscreened\PYGZus{}0:rc:
6.000000
fingerprintconstants:Mg\PYGZus{}Mg:radialscreened\PYGZus{}0:alpha:
5.520000 5.520000 5.520000 5.520000 5.520000
fingerprintconstants:Mg\PYGZus{}Mg:radialscreened\PYGZus{}0:dr:
2.806408
fingerprintconstants:Mg\PYGZus{}Mg:radialscreened\PYGZus{}0:o:
\PYGZhy{}1
fingerprintconstants:Mg\PYGZus{}Mg:radialscreened\PYGZus{}0:n:
3
\end{sphinxVerbatim}

\sphinxAtStartPar
\sphinxstyleemphasis{screening} specifies the Cmax and Cmin values used in the screening
fingerprints. Contributions form neighbors to the fingerprint are
omitted if they are blocked by a closer neighbor, and reduced if they
are partially blocked.  Larger values of Cmin correspond to neighbors
being blocked more easily.  Cmax cannot be greater than 3, and Cmin
cannot be greater than Cmax or less than zero.  Screening may be omitted
in which case the default values Cmax = 2.8, Cmin = 0.8 are used.  Since
screening is a bond computation, it is specified separately for each
combination of three elements in which the latter two may be
interchanged with no effect.

\begin{sphinxVerbatim}[commandchars=\\\{\}]
screening:Mg\PYGZus{}Mg\PYGZus{}Mg:Cmax:
2.700000
screening:Mg\PYGZus{}Mg\PYGZus{}Mg:Cmin:
0.400000
\end{sphinxVerbatim}

\sphinxAtStartPar
\sphinxstyleemphasis{networklayers} species the size of the neural network for each atom.
It counts both the input and output layer and so is 2 + \textless{}hidden layers\textgreater{}.

\begin{sphinxVerbatim}[commandchars=\\\{\}]
networklayers:Mg:
3
\end{sphinxVerbatim}

\sphinxAtStartPar
\sphinxstyleemphasis{layersize} specifies the length of each layer, including the input
layer and output layer.  The input layer is layer 0.  The size of the
input layer size must match the summed length of all the fingerprints
for that element, and the output layer size must be 1:

\begin{sphinxVerbatim}[commandchars=\\\{\}]
layersize:Mg:0:
14
layersize:Mg:1:
20
layersize:Mg:2:
1
\end{sphinxVerbatim}

\sphinxAtStartPar
\sphinxstyleemphasis{weight} specifies the weight for a given element and layer.  Weight
cannot be specified for the output layer.  The weight of layer i is a
\sphinxstyleemphasis{m} x \sphinxstyleemphasis{n} matrix where \sphinxstyleemphasis{m} is the layer size of \sphinxstyleemphasis{i} and \sphinxstyleemphasis{n} is the layer size of
\sphinxstyleemphasis{i}+1:

\begin{sphinxVerbatim}[commandchars=\\\{\}]
weight:Mg:0:
w11 w12 w13 ...
w21 w22 w23 ...
...
\end{sphinxVerbatim}

\sphinxAtStartPar
\sphinxstyleemphasis{bias} specifies the bias for a given element and layer.  Bias cannot be
specified for the output layer.  The bias of layer i is a nx1 vector
where n is the layer size of i+1:

\begin{sphinxVerbatim}[commandchars=\\\{\}]
bias:Mg:0:
b1
b2
b3
...
\end{sphinxVerbatim}

\sphinxAtStartPar
\sphinxstyleemphasis{activationfunctions} specifies the activation function for a given
element and layer.  Activation functions cannot be specified for the
output layer:

\begin{sphinxVerbatim}[commandchars=\\\{\}]
activationfunctions:Mg:0:
sigI
activationfunctions:Mg:1:
linear
\end{sphinxVerbatim}

\sphinxAtStartPar
The following activation styles are currently specified.  See the
{\hyperref[\detokenize{pair_rann:activations}]{\sphinxcrossref{\DUrole{std,std-ref}{formulation section}}}} below for their definitions.
\begin{itemize}
\item {} 
\sphinxAtStartPar
sigI

\item {} 
\sphinxAtStartPar
linear

\end{itemize}

\sphinxAtStartPar
\sphinxstyleemphasis{calibrationparameters} specifies a number of parameters used to calibrate the potential. These are ignored
by LAMMPS.


\subsection{Formulation}
\label{\detokenize{pair_rann:formulation}}
\sphinxAtStartPar
In the RANN formulation, the total energy of a system of atoms
is given by:
\begin{equation*}
\begin{split}E = \sum_{\alpha} E^{\alpha}\\\\
E^{\alpha} = {}^{N}\!A^{\alpha}\\\\
{}^{n+1}\!A_i^{\alpha} = {}^{n}\!F\left({}^{n}\!W_{ij}{\;}^{n}\!A_j^{\alpha}+{}^{n}\!B_i\right)\\\\
{}^{0}\!A_i^{\alpha} = \left[\begin{array}{c} {}^1\!S\!f^\alpha\\ {}^2\!S\!f^\alpha \\...\\\end{array}\right]\end{split}
\end{equation*}
\sphinxAtStartPar
Here \(E^\alpha\) is the energy of atom \(\alpha\),
\({}^n\!F()\), \({}^n\!W_{ij}\) and \({}^n\!B_i\) are the
activation function, weight matrix and bias vector of the n\sphinxhyphen{}th layer
respectively.  The inputs to the first layer are a collection of
structural fingerprints which are collected and reshaped into a single
long vector.  The individual fingerprints may be defined in any order
and have various shapes and sizes.  Multiple fingerprints of the same
type and varying parameters may also be defined in the input layer.

\sphinxAtStartPar
Eight types of structural fingerprints are currently defined.  In the
following, \(\beta\) and \(\gamma\) span the full neighbor list
of atom \(\alpha\).  \(\delta_i\) are decay meta\sphinxhyphen{}parameters, and
\(r_e\) is a meta\sphinxhyphen{}parameter roughly proportional to the first
neighbor distance.  \(r_c\) and \(dr\) are the neighbor cutoff
distance and cutoff smoothing distance respectively.
\(S^{\alpha\beta}\) is the MEAM screening function {\hyperref[\detokenize{pair_rann:baskes97}]{\sphinxcrossref{\DUrole{std,std-ref}{(Baskes)}}}}, \(s_i^\alpha\) and \(s_i^\beta\) are the atom spin
vectors {\hyperref[\detokenize{pair_rann:tranchida7}]{\sphinxcrossref{\DUrole{std,std-ref}{(Tranchida)}}}}.  \(r^{\alpha\beta}\) is the
distance from atom \(\alpha\) to atom \(\beta\), and
\(\theta^{\alpha\beta\gamma}\) is the bond angle:
\begin{equation*}
\begin{split}cos\left(\theta^{\alpha\beta\gamma}\right)=\frac{\mathbf{r}^{\alpha\beta} \cdot \mathbf{r}^{\alpha\gamma}}{r^{\alpha\beta}r^{\alpha\gamma}}\end{split}
\end{equation*}
\sphinxAtStartPar
\(S^{\alpha\beta}\) is defined as {\hyperref[\detokenize{pair_rann:baskes97}]{\sphinxcrossref{\DUrole{std,std-ref}{(Baskes)}}}}:
\begin{equation*}
\begin{split}X^{\gamma\beta} = \left(\frac{r^{\gamma\beta}}{r^{\alpha\beta}}\right)^2\\
\\
X^{\alpha\gamma} = \left(\frac{r^{\alpha\gamma}}{r^{\alpha\beta}}\right)^2\\
\\
C = \frac{2\left(X^{\alpha\gamma}+X^{\gamma\beta}\right)-\left(X^{\alpha\gamma}-X^{\gamma\beta}\right)^2-1}{1-\left(X^{\alpha\gamma}-X^{\gamma\beta}\right)^2}\\
\\
f_c(x) = \left[\begin{array}{l}  1 \; \: x \geq 1\\ \left(1-\left(1-x\right)^4\right)^2 \; \: 0<x<1\\0\; \; x\leq0\end{array}\right.\\
\\
S^{\alpha\beta\gamma} = f_c\left(\frac{C-C_{min}}{C_{max}-C_{min}}\right)\\
\\
S^{\alpha\beta} = \prod_\gamma S^{\alpha\beta\gamma}\\\end{split}
\end{equation*}
\sphinxAtStartPar
The structural fingerprints are computed as follows:

\phantomsection\label{\detokenize{pair_rann:fingerprints}}\begin{itemize}
\item {} 
\sphinxAtStartPar
\sphinxstylestrong{radial}

\end{itemize}
\begin{equation*}
\begin{split}{}^r\!S\!f_i^\alpha = \sum_{\beta} \left(\frac{r^{\alpha\beta}}{r_e}\right)^ie^{-\delta_i \frac{r^{\alpha\beta}}{r_e}}f_c\left(\frac{r_c-r^{\alpha\beta}}{dr}\right)\end{split}
\end{equation*}\begin{itemize}
\item {} 
\sphinxAtStartPar
\sphinxstylestrong{bond}

\end{itemize}
\begin{equation*}
\begin{split}{}^b\!S\!f_{ij}^\alpha = \sum_{\beta}\sum_{\gamma} \left(cos(\theta_{\alpha\beta\gamma})\right)^ie^{-\delta_j \frac{r^{\alpha\beta}}{r_e}}e^{-\delta_j \frac{r^{\alpha\gamma}}{r_e}}f_c\left(\frac{r_c-r^{\alpha\beta}}{dr}\right)f_c\left(\frac{r_c-r^{\alpha\gamma}}{dr}\right)\end{split}
\end{equation*}\begin{itemize}
\item {} 
\sphinxAtStartPar
\sphinxstylestrong{radialscreened}

\end{itemize}
\begin{equation*}
\begin{split}{}^{rsc}\!S\!f_i^\alpha = \sum_{\beta} \left(\frac{r^{\alpha\beta}}{r_e}\right)^ie^{-\delta_i \frac{r^{\alpha\beta}}{r_e}}S^{\alpha\beta}f_c\left(\frac{r_c-r^{\alpha\beta}}{dr}\right)\end{split}
\end{equation*}\begin{itemize}
\item {} 
\sphinxAtStartPar
\sphinxstylestrong{bondscreened}

\end{itemize}
\begin{equation*}
\begin{split}{}^{bsc}\!S\!f_{ij}^\alpha = \sum_{\beta}\sum_{\gamma} \left(cos(\theta_{\alpha\beta\gamma})\right)^ie^{-\delta_j \frac{r^{\alpha\beta}}{r_e}}e^{-\delta_j \frac{r^{\alpha\gamma}}{r_e}}S^{\alpha\beta}S^{\alpha\gamma}f_c\left(\frac{r_c-r^{\alpha\beta}}{dr}\right)f_c\left(\frac{r_c-r^{\alpha\gamma}}{dr}\right)\end{split}
\end{equation*}\begin{itemize}
\item {} 
\sphinxAtStartPar
\sphinxstylestrong{radialspin}

\end{itemize}
\begin{equation*}
\begin{split}{}^{rsp}\!S\!f_i^\alpha = \sum_{\beta} \left(\frac{r^{\alpha\beta}}{r_e}\right)^ie^{-\delta_i \frac{r^{\alpha\beta}}{r_e}}\left(\mathbf{s^\alpha \cdot s^\beta}\right)f_c\left(\frac{r_c-r^{\alpha\beta}}{dr}\right)\end{split}
\end{equation*}\begin{itemize}
\item {} 
\sphinxAtStartPar
\sphinxstylestrong{bondspin}

\end{itemize}
\begin{equation*}
\begin{split}{}^{bsp}\!S\!f_{ij}^\alpha = \sum_{\beta}\sum_{\gamma} \left(cos(\theta_{\alpha\beta\gamma})\right)^ie^{-\delta_j \frac{r^{\alpha\beta}}{r_e}}e^{-\delta_j \frac{r^{\alpha\gamma}}{r_e}}\left(\mathbf{s^\alpha \cdot s^\beta}\right)\left(\mathbf{s^\alpha \cdot s^\gamma}\right)f_c\left(\frac{r_c-r^{\alpha\beta}}{dr}\right)f_c\left(\frac{r_c-r^{\alpha\gamma}}{dr}\right)\end{split}
\end{equation*}\begin{itemize}
\item {} 
\sphinxAtStartPar
\sphinxstylestrong{radialscreenedspin}

\end{itemize}
\begin{equation*}
\begin{split}{}^{rscsp}\!S\!f_i^\alpha = \sum_{\beta} \left(\frac{r^{\alpha\beta}}{r_e}\right)^ie^{-\delta_i \frac{r^{\alpha\beta}}{r_e}}S^{\alpha\beta}\left(\mathbf{s^\alpha \cdot s^\beta}\right)f_c\left(\frac{r_c-r^{\alpha\beta}}{dr}\right)\end{split}
\end{equation*}\begin{itemize}
\item {} 
\sphinxAtStartPar
\sphinxstylestrong{bondscreenedspin}

\end{itemize}
\begin{equation*}
\begin{split}{}^{bscsp}\!S\!f_{ij}^\alpha = \sum_{\beta}\sum_{\gamma} \left(cos(\theta_{\alpha\beta\gamma})\right)^ie^{-\delta_j \frac{r^{\alpha\beta}}{r_e}}e^{-\delta_j \frac{r^{\alpha\gamma}}{r_e}}S^{\alpha\beta}S^{\alpha\gamma}\left(\mathbf{s^\alpha \cdot s^\beta}\right)\left(\mathbf{s^\alpha \cdot s^\gamma}\right)f_c\left(\frac{r_c-r^{\alpha\beta}}{dr}\right)f_c\left(\frac{r_c-r^{\alpha\gamma}}{dr}\right)\end{split}
\end{equation*}
\sphinxAtStartPar
The activation functions are computed as follows:

\phantomsection\label{\detokenize{pair_rann:activations}}\begin{itemize}
\item {} 
\sphinxAtStartPar
\sphinxstylestrong{sigI}

\end{itemize}
\begin{equation*}
\begin{split}F^{sigI}(x) = 0.1x+0.9ln\left(e^x+1\right)\end{split}
\end{equation*}\begin{itemize}
\item {} 
\sphinxAtStartPar
\sphinxstylestrong{linear}

\end{itemize}
\begin{equation*}
\begin{split}F^{linear}(x) = x\end{split}
\end{equation*}

\subsection{Restrictions}
\label{\detokenize{pair_rann:restrictions}}
\sphinxAtStartPar
Pair style \sphinxstyleemphasis{rann} is part of the ML\sphinxhyphen{}RANN package.  It is only enabled if LAMMPS was built with that
package.  Additionally, if any spin fingerprint styles are used LAMMPS must be built with the SPIN
package as well.

\sphinxAtStartPar
Pair style \sphinxstyleemphasis{rann} does not support computing per\sphinxhyphen{}atom stress or using {\hyperref[\detokenize{pair_modify::doc}]{\sphinxcrossref{\DUrole{doc}{pair\_modify nofdotr}}}}.


\subsection{Defaults}
\label{\detokenize{pair_rann:defaults}}
\sphinxAtStartPar
Cmin = 0.8, Cmax = 2.8.


\bigskip\hrule\bigskip

\phantomsection\label{\detokenize{pair_rann:baskes97}}
\sphinxAtStartPar
\sphinxstylestrong{(Baskes)} Baskes, Materials Chemistry and Physics, 50(2), 152\sphinxhyphen{}158, (1997).

\phantomsection\label{\detokenize{pair_rann:dickel}}
\sphinxAtStartPar
\sphinxstylestrong{(Dickel)} Dickel, Francis, and Barrett, Computational Materials Science 171 (2020): 109157.

\phantomsection\label{\detokenize{pair_rann:nitol}}
\sphinxAtStartPar
\sphinxstylestrong{(Nitol)} Nitol, Dickel, and Barrett, Computational Materials Science 188 (2021): 110207.

\phantomsection\label{\detokenize{pair_rann:tranchida7}}
\sphinxAtStartPar
\sphinxstylestrong{(Tranchida)} Tranchida, Plimpton, Thibaudeau and Thompson, Journal of Computational Physics, 372, 406\sphinxhyphen{}425, (2018).

\sphinxstepscope

\index{pair\_style reaxff@\spxentry{pair\_style reaxff}}\ignorespaces 
\index{pair\_style reaxff/kk@\spxentry{pair\_style reaxff/kk}}\ignorespaces 
\index{pair\_style reaxff/omp@\spxentry{pair\_style reaxff/omp}}\ignorespaces 

\section{pair\_style reaxff command}
\label{\detokenize{pair_reaxff:pair-style-reaxff-command}}\label{\detokenize{pair_reaxff:index-2}}\label{\detokenize{pair_reaxff:index-1}}\label{\detokenize{pair_reaxff:index-0}}\label{\detokenize{pair_reaxff::doc}}
\sphinxAtStartPar
Accelerator Variants: \sphinxstyleemphasis{reaxff/kk}, \sphinxstyleemphasis{reaxff/omp}


\subsection{Syntax}
\label{\detokenize{pair_reaxff:syntax}}
\begin{sphinxVerbatim}[commandchars=\\\{\}]
\PYG{k}{pair\PYGZus{}style}\PYG{+w}{ }\PYG{n}{reaxff}\PYG{+w}{ }\PYG{n}{cfile}\PYG{+w}{ }\PYG{n}{keyword}\PYG{+w}{ }\PYG{n}{value}
\end{sphinxVerbatim}
\begin{itemize}
\item {} 
\sphinxAtStartPar
cfile = NULL or name of a control file

\item {} 
\sphinxAtStartPar
zero or more keyword/value pairs may be appended
\begin{sphinxalltt}
keyword = \sphinxstyleemphasis{checkqeq} or \sphinxstyleemphasis{lgvdw} or \sphinxstyleemphasis{safezone} or \sphinxstyleemphasis{mincap} or \sphinxstyleemphasis{minhbonds} or \sphinxstyleemphasis{tabulate} or \sphinxstyleemphasis{list/blocking}
  \sphinxstyleemphasis{checkqeq} value = \sphinxstyleemphasis{yes} or \sphinxstyleemphasis{no} = whether or not to require qeq/reaxff or acks2/reaxff fix
  \sphinxstyleemphasis{enobonds} value = \sphinxstyleemphasis{yes} or \sphinxstyleemphasis{no} = whether or not to tally energy of atoms with no bonds
  \sphinxstyleemphasis{lgvdw} value = \sphinxstyleemphasis{yes} or \sphinxstyleemphasis{no} = whether or not to use a low gradient vdW correction
  \sphinxstyleemphasis{safezone} = factor used for array allocation
  \sphinxstyleemphasis{mincap} = minimum size for array allocation
  \sphinxstyleemphasis{minhbonds} = minimum size use for storing hydrogen bonds
  \sphinxstyleemphasis{tabulate} value = size of interpolation table for Lennard\sphinxhyphen{}Jones and Coulomb interactions
  \sphinxstyleemphasis{list/blocking} value = \sphinxstyleemphasis{yes} or \sphinxstyleemphasis{no} = whether or not to use "blocking" scheme for bond list build
\end{sphinxalltt}

\end{itemize}


\subsection{Examples}
\label{\detokenize{pair_reaxff:examples}}
\begin{sphinxVerbatim}[commandchars=\\\{\}]
\PYG{k}{pair\PYGZus{}style}\PYG{+w}{ }\PYG{n}{reaxff}\PYG{+w}{ }\PYG{n}{NULL}
\PYG{k}{pair\PYGZus{}style}\PYG{+w}{ }\PYG{n}{reaxff}\PYG{+w}{ }\PYG{n}{controlfile}\PYG{+w}{ }\PYG{n}{checkqeq}\PYG{+w}{ }\PYG{n}{no}
\PYG{k}{pair\PYGZus{}style}\PYG{+w}{ }\PYG{n}{reaxff}\PYG{+w}{ }\PYG{n}{NULL}\PYG{+w}{ }\PYG{n}{lgvdw}\PYG{+w}{ }\PYG{n}{yes}
\PYG{k}{pair\PYGZus{}style}\PYG{+w}{ }\PYG{n}{reaxff}\PYG{+w}{ }\PYG{n}{NULL}\PYG{+w}{ }\PYG{n}{safezone}\PYG{+w}{ }\PYG{l+m}{1.6}\PYG{+w}{ }\PYG{n}{mincap}\PYG{+w}{ }\PYG{l+m}{100}
\PYG{k}{pair\PYGZus{}coeff}\PYG{+w}{ }\PYG{o}{*}\PYG{+w}{ }\PYG{o}{*}\PYG{+w}{ }\PYG{n}{ffield.reax}\PYG{+w}{ }\PYG{n}{C}\PYG{+w}{ }\PYG{n}{H}\PYG{+w}{ }\PYG{n}{O}\PYG{+w}{ }\PYG{n}{N}
\end{sphinxVerbatim}


\subsection{Description}
\label{\detokenize{pair_reaxff:description}}
\sphinxAtStartPar
Pair style \sphinxstyleemphasis{reaxff} computes the ReaxFF potential of van Duin, Goddard
and co\sphinxhyphen{}workers.  ReaxFF uses distance\sphinxhyphen{}dependent bond\sphinxhyphen{}order functions to
represent the contributions of chemical bonding to the potential
energy.  There is more than one version of ReaxFF.  The version
implemented in LAMMPS uses the functional forms documented in the
supplemental information of the following paper:
{\hyperref[\detokenize{pair_reaxff:chenoweth-20082}]{\sphinxcrossref{\DUrole{std,std-ref}{(Chenoweth et al., 2008)}}}} and matches the
version of the reference ReaxFF implementation from Summer 2010.  For
more technical details about the implementation of ReaxFF in pair style
\sphinxstyleemphasis{reaxff}, see the {\hyperref[\detokenize{pair_reaxff:aktulga}]{\sphinxcrossref{\DUrole{std,std-ref}{(Aktulga)}}}} paper. The \sphinxstyleemphasis{reaxff} style
was initially implemented as a stand\sphinxhyphen{}alone C code and is now converted
to C++ and integrated into LAMMPS as a package.

\sphinxAtStartPar
The \sphinxstyleemphasis{reaxff/kk} style is a Kokkos version of the ReaxFF potential that
is derived from the \sphinxstyleemphasis{reaxff} style.  The Kokkos version can run on GPUs
and can also use OpenMP multithreading.  For more information about the
Kokkos package, see {\hyperref[\detokenize{Packages_details::doc}]{\sphinxcrossref{\DUrole{doc}{Packages details}}}} and
{\hyperref[\detokenize{Speed_kokkos::doc}]{\sphinxcrossref{\DUrole{doc}{Speed kokkos}}}} doc pages.  One important
consideration when using the \sphinxstyleemphasis{reaxff/kk} style is the choice of either
half or full neighbor lists. This setting can be changed using the
Kokkos {\hyperref[\detokenize{package::doc}]{\sphinxcrossref{\DUrole{doc}{package}}}} command.

\sphinxAtStartPar
The \sphinxstyleemphasis{reaxff} style differs from the (obsolete) “pair\_style reax”
command in the implementation details.  The \sphinxstyleemphasis{reax} style was a
Fortran library, linked to LAMMPS.  The \sphinxstyleemphasis{reax} style has been removed
from LAMMPS after the 12 December 2018 version.

\sphinxAtStartPar
LAMMPS provides several different versions of ffield.reax in its
potentials dir, each called potentials/ffield.reax.label.  These are
documented in potentials/README.reax.

\sphinxAtStartPar
The format of these files is identical to that used originally by van
Duin.  We have tested the accuracy of \sphinxstyleemphasis{pair\_style reaxff} potential
against the original ReaxFF code for the systems mentioned above.  You
can use other ffield files for specific chemical systems that may be
available elsewhere (but note that their accuracy may not have been
tested).

\begin{sphinxadmonition}{note}{Note:}
\sphinxAtStartPar
We do not distribute a wide variety of ReaxFF force field files
with LAMMPS.  Adri van Duin’s group at PSU is the central repository
for this kind of data as they are continuously deriving and updating
parameterizations for different classes of materials.  You can submit
a contact request at the Materials Computation Center (MCC) website
\sphinxurl{https://www.mri.psu.edu/materials-computation-center/connect-mcc},
describing the material(s) you are interested in modeling with ReaxFF.
They can tell you what is currently available or what it would take to
create a suitable ReaxFF parameterization.
\end{sphinxadmonition}

\sphinxAtStartPar
The \sphinxstyleemphasis{cfile} setting can be specified as NULL, in which case default
settings are used. A control file can be specified which defines values
of control variables. Some control variables are global parameters for
the ReaxFF potential. Others define certain performance and output
settings.  Each line in the control file specifies the value for a
control variable.  The format of the control file is described below.

\begin{sphinxadmonition}{note}{Note:}
\sphinxAtStartPar
The LAMMPS default values for the ReaxFF global parameters
correspond to those used by Adri van Duin’s stand\sphinxhyphen{}alone serial
code.  If these are changed by setting control variables in the
control file, the results from LAMMPS and the serial code will
not agree.
\end{sphinxadmonition}

\sphinxAtStartPar
Examples using \sphinxstyleemphasis{pair\_style reaxff} are provided in the examples/reax
directory and its subdirectories.

\sphinxAtStartPar
Use of this pair style requires using an {\hyperref[\detokenize{atom_style::doc}]{\sphinxcrossref{\DUrole{doc}{atom\_style}}}}
that includes a per\sphinxhyphen{}atom charge property \sphinxstyleemphasis{or} using
{\hyperref[\detokenize{fix_property_atom::doc}]{\sphinxcrossref{\DUrole{doc}{fix property/atom q}}}}.  Charges can be set
via {\hyperref[\detokenize{read_data::doc}]{\sphinxcrossref{\DUrole{doc}{read\_data}}}} or {\hyperref[\detokenize{set::doc}]{\sphinxcrossref{\DUrole{doc}{set}}}}.  Using an initial
charge that is close to the result of charge equilibration will speed
up that process.

\sphinxAtStartPar
The ReaxFF parameter files provided were created using a charge
equilibration (QEq) model for handling the electrostatic interactions.
Therefore, by default, LAMMPS requires that either the
{\hyperref[\detokenize{fix_qeq_reaxff::doc}]{\sphinxcrossref{\DUrole{doc}{fix qeq/reaxff}}}} or the
{\hyperref[\detokenize{fix_qeq::doc}]{\sphinxcrossref{\DUrole{doc}{fix qeq/shielded}}}} or {\hyperref[\detokenize{fix_acks2_reaxff::doc}]{\sphinxcrossref{\DUrole{doc}{fix acks2/reaxff}}}}
command be used with
\sphinxstyleemphasis{pair\_style reaxff} when simulating a ReaxFF model, to equilibrate
the charges each timestep.

\sphinxAtStartPar
Using the keyword \sphinxstyleemphasis{checkqeq} with the value \sphinxstyleemphasis{no} turns off the check
for the QEq fixes, allowing a simulation to be run without charge
equilibration. In this case, the static charges you assign to each
atom will be used for computing the electrostatic interactions in
the system. See the {\hyperref[\detokenize{fix_qeq_reaxff::doc}]{\sphinxcrossref{\DUrole{doc}{fix qeq/reaxff}}}} or
{\hyperref[\detokenize{fix_qeq::doc}]{\sphinxcrossref{\DUrole{doc}{fix qeq/shielded}}}} or {\hyperref[\detokenize{fix_acks2_reaxff::doc}]{\sphinxcrossref{\DUrole{doc}{fix acks2/reaxff}}}}
command documentation for more details.

\sphinxAtStartPar
Using the optional keyword \sphinxstyleemphasis{lgvdw} with the value \sphinxstyleemphasis{yes} turns on the
low\sphinxhyphen{}gradient correction of ReaxFF for long\sphinxhyphen{}range London Dispersion,
as described in the {\hyperref[\detokenize{pair_reaxff:liu-2011}]{\sphinxcrossref{\DUrole{std,std-ref}{(Liu)}}}} paper. The bundled force
field file \sphinxstyleemphasis{ffield.reax.lg} is designed for this correction, and is
trained for several energetic materials (see “Liu”). When using \sphinxstyleemphasis{lgvdw yes},
the recommended value for parameter \sphinxstyleemphasis{thb} is 0.01, which can be set in the
control file.  Note: Force field files are different for the original
or lg corrected pair styles, using the wrong ffield file generates an
error.

\sphinxAtStartPar
Using the optional keyword \sphinxstyleemphasis{enobonds} with the value \sphinxstyleemphasis{yes}, the energy
of atoms with no bonds (i.e. isolated atoms) is included in the total
potential energy and the per\sphinxhyphen{}atom energy of that atom.  If the value
\sphinxstyleemphasis{no} is specified then the energy of atoms with no bonds is set to
zero.  The latter behavior is usual not desired, as it causes
discontinuities in the potential energy when the bonding of an atom
drops to zero.

\sphinxAtStartPar
Optional keywords \sphinxstyleemphasis{safezone}, \sphinxstyleemphasis{mincap}, and \sphinxstyleemphasis{minhbonds} are used
for allocating reaxff arrays.  Increasing these values can avoid memory
problems, such as segmentation faults and bondchk failed errors, that
could occur under certain conditions. These keywords are not used by
the Kokkos version, which instead uses a more robust memory allocation
scheme that checks if the sizes of the arrays have been exceeded and
automatically allocates more memory.

\sphinxAtStartPar
The keyword \sphinxstyleemphasis{tabulate} controls the size of interpolation table for
Lennard\sphinxhyphen{}Jones and Coulomb interactions. Tabulation may also be set in the
control file (see below). If tabulation is set in both the input script and the
control file, the value in the control file will be ignored. A size of 10000 is
typically used for the interpolation table. A value of 0 means no tabulation
will be used.

\sphinxAtStartPar
The keyword \sphinxstyleemphasis{list/blocking} is only supported by the Kokkos version of
ReaxFF and ignored otherwise. Setting the value to \sphinxstyleemphasis{yes} enables the
“blocking” scheme (dynamically building interaction lists) for the
ReaxFF bond neighbor list. This reduces the number of empty
interactions and can improve performance in some cases (e.g. large
number of atoms/GPU on AMD hardware). It is also enabled by default
when running the CPU with Kokkos.

\sphinxAtStartPar
The thermo variable \sphinxstyleemphasis{evdwl} stores the sum of all the ReaxFF potential
energy contributions, with the exception of the Coulombic and charge
equilibration contributions which are stored in the thermo variable
\sphinxstyleemphasis{ecoul}.  The output of these quantities is controlled by the
{\hyperref[\detokenize{thermo::doc}]{\sphinxcrossref{\DUrole{doc}{thermo}}}} command.

\sphinxAtStartPar
This pair style tallies a breakdown of the total ReaxFF potential
energy into sub\sphinxhyphen{}categories, which can be accessed via the
{\hyperref[\detokenize{compute_pair::doc}]{\sphinxcrossref{\DUrole{doc}{compute pair}}}} command as a vector of values of length 14.
The 14 values correspond to the following sub\sphinxhyphen{}categories (the variable
names in italics match those used in the original FORTRAN ReaxFF
code):
\begin{enumerate}
\sphinxsetlistlabels{\arabic}{enumi}{enumii}{}{.}%
\item {} 
\sphinxAtStartPar
\sphinxstyleemphasis{eb} = bond energy

\item {} 
\sphinxAtStartPar
\sphinxstyleemphasis{ea} = atom energy

\item {} 
\sphinxAtStartPar
\sphinxstyleemphasis{elp} = lone\sphinxhyphen{}pair energy

\item {} 
\sphinxAtStartPar
\sphinxstyleemphasis{emol} = molecule energy (always 0.0)

\item {} 
\sphinxAtStartPar
\sphinxstyleemphasis{ev} = valence angle energy

\item {} 
\sphinxAtStartPar
\sphinxstyleemphasis{epen} = double\sphinxhyphen{}bond valence angle penalty

\item {} 
\sphinxAtStartPar
\sphinxstyleemphasis{ecoa} = valence angle conjugation energy

\item {} 
\sphinxAtStartPar
\sphinxstyleemphasis{ehb} = hydrogen bond energy

\item {} 
\sphinxAtStartPar
\sphinxstyleemphasis{et} = torsion energy

\item {} 
\sphinxAtStartPar
\sphinxstyleemphasis{eco} = conjugation energy

\item {} 
\sphinxAtStartPar
\sphinxstyleemphasis{ew} = van der Waals energy

\item {} 
\sphinxAtStartPar
\sphinxstyleemphasis{ep} = Coulomb energy

\item {} 
\sphinxAtStartPar
\sphinxstyleemphasis{efi} = electric field energy (always 0.0)

\item {} 
\sphinxAtStartPar
\sphinxstyleemphasis{eqeq} = charge equilibration energy

\end{enumerate}

\sphinxAtStartPar
To print these quantities to the log file (with descriptive column
headings) the following commands could be included in an input script:

\begin{sphinxVerbatim}[commandchars=\\\{\}]
\PYG{k}{compute }\PYG{n+nv+nvIdentifier}{reax}\PYG{+w}{ }\PYG{n+nv+nvIdentifier}{all}\PYG{+w}{ }\PYG{n}{pair}\PYG{+w}{ }\PYG{n}{reaxff}
\PYG{k}{variable }\PYG{n+nv+nvIdentifier}{eb}\PYG{+w}{      }\PYG{n}{equal}\PYG{+w}{ }\PYG{n}{c\PYGZus{}reax[1]}
\PYG{k}{variable }\PYG{n+nv+nvIdentifier}{ea}\PYG{+w}{      }\PYG{n}{equal}\PYG{+w}{ }\PYG{n}{c\PYGZus{}reax[2]}
\PYG{n}{[...]}
\PYG{k}{variable }\PYG{n+nv+nvIdentifier}{eqeq}\PYG{+w}{    }\PYG{n}{equal}\PYG{+w}{ }\PYG{n}{c\PYGZus{}reax[14]}
\PYG{k}{thermo\PYGZus{}style}\PYG{+w}{ }\PYG{n}{custom}\PYG{+w}{ }\PYG{n}{step}\PYG{+w}{ }\PYG{n}{temp}\PYG{+w}{ }\PYG{n}{epair}\PYG{+w}{ }\PYG{n}{v\PYGZus{}eb}\PYG{+w}{ }\PYG{n}{v\PYGZus{}ea}\PYG{+w}{ }\PYG{n}{[...]}\PYG{+w}{ }\PYG{n}{v\PYGZus{}eqeq}
\end{sphinxVerbatim}

\sphinxAtStartPar
Only a single pair\_coeff command is used with the \sphinxstyleemphasis{reaxff} style which
specifies a ReaxFF potential file with parameters for all needed
elements.  These are mapped to LAMMPS atom types by specifying N
additional arguments after the filename in the pair\_coeff command,
where N is the number of LAMMPS atom types:
\begin{itemize}
\item {} 
\sphinxAtStartPar
filename

\item {} 
\sphinxAtStartPar
N indices = ReaxFF elements

\end{itemize}

\sphinxAtStartPar
The filename is the ReaxFF potential file.

\sphinxAtStartPar
In the ReaxFF potential file, near the top, after the general
parameters, is the atomic parameters section that contains element
names, each with a couple dozen numeric parameters.  If there are M
elements specified in the \sphinxstyleemphasis{ffield} file, think of these as numbered 1
to M. Each of the N indices you specify for the N atom types of LAMMPS
atoms must be an integer from 1 to M.  Atoms with LAMMPS type 1 will
be mapped to whatever element you specify as the first index value,
etc.  If a mapping value is specified as NULL, the mapping is not
performed.  This can be used when the \sphinxstyleemphasis{reaxff} style is used as part
of the \sphinxstyleemphasis{hybrid} pair style.  The NULL values are placeholders for atom
types that will be used with other potentials.

\sphinxAtStartPar
As an example, say your LAMMPS simulation has 4 atom types and the
elements are ordered as C, H, O, N in the \sphinxstyleemphasis{ffield} file.  If you want
the LAMMPS atom type 1 and 2 to be C, type 3 to be N, and type 4 to be
H, you would use the following pair\_coeff command:

\begin{sphinxVerbatim}[commandchars=\\\{\}]
\PYG{k}{pair\PYGZus{}coeff}\PYG{+w}{ }\PYG{o}{*}\PYG{+w}{ }\PYG{o}{*}\PYG{+w}{ }\PYG{n}{ffield.reax}\PYG{+w}{ }\PYG{n}{C}\PYG{+w}{ }\PYG{n}{C}\PYG{+w}{ }\PYG{n}{N}\PYG{+w}{ }\PYG{n}{H}
\end{sphinxVerbatim}


\bigskip\hrule\bigskip



\subsection{Control file}
\label{\detokenize{pair_reaxff:control-file}}
\sphinxAtStartPar
The format of a line in the control file is as follows:

\begin{sphinxVerbatim}[commandchars=\\\{\}]
variable\PYGZus{}name value
\end{sphinxVerbatim}

\sphinxAtStartPar
and it may be followed by an “!” character and a trailing comment.

\sphinxAtStartPar
If the value of a control variable is not specified, then default
values are used.  What follows is the list of variables along with a
brief description of their use and default values.
\begin{description}
\sphinxlineitem{\sphinxstyleemphasis{simulation\_name}}
\sphinxAtStartPar
Output files produced by \sphinxstyleemphasis{pair\_style reaxff} carry
this name + extensions specific to their contents.  Partial energies
are reported with a “.pot” extension, while the trajectory file has
“.trj” extension.

\sphinxlineitem{\sphinxstyleemphasis{tabulate\_long\_range}}
\sphinxAtStartPar
To improve performance, long range interactions can optionally be
tabulated (0 means no tabulation). Value of this variable denotes the
size of the long range interaction table.  The range from 0 to long
range cutoff (defined in the \sphinxstyleemphasis{ffield} file) is divided into
\sphinxstyleemphasis{tabulate\_long\_range} points.  Then at the start of simulation, we
fill in the entries of the long range interaction table by computing
the energies and forces resulting from van der Waals and Coulomb
interactions between every possible atom type pairs present in the
input system.  During the simulation we consult to the long range
interaction table to estimate the energy and forces between a pair of
atoms. Linear interpolation is used for estimation. (default value = 0)

\sphinxlineitem{\sphinxstyleemphasis{energy\_update\_freq}}
\sphinxAtStartPar
Denotes the frequency (in number of steps) of writes into the partial
energies file. (default value = 0)

\sphinxlineitem{\sphinxstyleemphasis{nbrhood\_cutoff}}
\sphinxAtStartPar
Denotes the near neighbors cutoff (in Angstroms)
regarding the bonded interactions. (default value = 5.0)

\sphinxlineitem{\sphinxstyleemphasis{hbond\_cutoff}}
\sphinxAtStartPar
Denotes the cutoff distance (in Angstroms) for hydrogen
bond interactions.(default value = 7.5. A value of 0.0 turns off
hydrogen bonds)

\sphinxlineitem{\sphinxstyleemphasis{bond\_graph\_cutoff}}
\sphinxAtStartPar
is the threshold used in determining what is a
physical bond, what is not. Bonds and angles reported in the
trajectory file rely on this cutoff. (default value = 0.3)

\sphinxlineitem{\sphinxstyleemphasis{thb\_cutoff}}
\sphinxAtStartPar
cutoff value for the strength of bonds to be considered in
three body interactions. (default value = 0.001)

\sphinxlineitem{\sphinxstyleemphasis{thb\_cutoff\_sq}}
\sphinxAtStartPar
cutoff value for the strength of bond order products
to be considered in three body interactions. (default value = 0.00001)

\sphinxlineitem{\sphinxstyleemphasis{write\_freq}}
\sphinxAtStartPar
Frequency of writes into the trajectory file. (default value = 0)

\sphinxlineitem{\sphinxstyleemphasis{traj\_title}}
\sphinxAtStartPar
Title of the trajectory \sphinxhyphen{} not the name of the trajectory file.

\sphinxlineitem{\sphinxstyleemphasis{atom\_info}}
\sphinxAtStartPar
1 means print only atomic positions + charge (default = 0)

\sphinxlineitem{\sphinxstyleemphasis{atom\_forces}}
\sphinxAtStartPar
1 adds net forces to atom lines in the trajectory file (default = 0)

\sphinxlineitem{\sphinxstyleemphasis{atom\_velocities}}
\sphinxAtStartPar
1 adds atomic velocities to atoms line (default = 0)

\sphinxlineitem{\sphinxstyleemphasis{bond\_info}}
\sphinxAtStartPar
1 prints bonds in the trajectory file (default = 0)

\sphinxlineitem{\sphinxstyleemphasis{angle\_info}}
\sphinxAtStartPar
1 prints angles in the trajectory file (default = 0)

\end{description}


\bigskip\hrule\bigskip



\subsection{Mixing, shift, table, tail correction, restart, rRESPA info}
\label{\detokenize{pair_reaxff:mixing-shift-table-tail-correction-restart-rrespa-info}}
\sphinxAtStartPar
This pair style does not support the {\hyperref[\detokenize{pair_modify::doc}]{\sphinxcrossref{\DUrole{doc}{pair\_modify}}}}
mix, shift, table, and tail options.

\sphinxAtStartPar
This pair style does not write its information to {\hyperref[\detokenize{restart::doc}]{\sphinxcrossref{\DUrole{doc}{binary restart files}}}},
since it is stored in potential files.  Thus, you
need to re\sphinxhyphen{}specify the pair\_style and pair\_coeff commands in an input
script that reads a restart file.

\sphinxAtStartPar
This pair style can only be used via the \sphinxstyleemphasis{pair} keyword of the
{\hyperref[\detokenize{run_style::doc}]{\sphinxcrossref{\DUrole{doc}{run\_style respa}}}} command.  It does not support the
\sphinxstyleemphasis{inner}, \sphinxstyleemphasis{middle}, \sphinxstyleemphasis{outer} keywords.


\bigskip\hrule\bigskip


\sphinxAtStartPar
Styles with a \sphinxstyleemphasis{gpu}, \sphinxstyleemphasis{intel}, \sphinxstyleemphasis{kk}, \sphinxstyleemphasis{omp}, or \sphinxstyleemphasis{opt} suffix are
functionally the same as the corresponding style without the suffix.
They have been optimized to run faster, depending on your available
hardware, as discussed on the {\hyperref[\detokenize{Speed_packages::doc}]{\sphinxcrossref{\DUrole{doc}{Accelerator packages}}}}
page.  The accelerated styles take the same arguments and should
produce the same results, except for round\sphinxhyphen{}off and precision issues.

\sphinxAtStartPar
These accelerated styles are part of the GPU, INTEL, KOKKOS,
OPENMP, and OPT packages, respectively.  They are only enabled if
LAMMPS was built with those packages.  See the {\hyperref[\detokenize{Build_package::doc}]{\sphinxcrossref{\DUrole{doc}{Build package}}}} page for more info.

\sphinxAtStartPar
You can specify the accelerated styles explicitly in your input script
by including their suffix, or you can use the {\hyperref[\detokenize{Run_options::doc}]{\sphinxcrossref{\DUrole{doc}{\sphinxhyphen{}suffix command\sphinxhyphen{}line switch}}}} when you invoke LAMMPS, or you can use the
{\hyperref[\detokenize{suffix::doc}]{\sphinxcrossref{\DUrole{doc}{suffix}}}} command in your input script.

\sphinxAtStartPar
See the {\hyperref[\detokenize{Speed_packages::doc}]{\sphinxcrossref{\DUrole{doc}{Accelerator packages}}}} page for more
instructions on how to use the accelerated styles effectively.


\bigskip\hrule\bigskip



\subsection{Restrictions}
\label{\detokenize{pair_reaxff:restrictions}}
\sphinxAtStartPar
This pair style is part of the REAXFF package.  It is only enabled
if LAMMPS was built with that package.
See the {\hyperref[\detokenize{Build_package::doc}]{\sphinxcrossref{\DUrole{doc}{Build package}}}} page for more info.

\sphinxAtStartPar
The ReaxFF potential files provided with LAMMPS in the potentials
directory are parameterized for \sphinxstyleemphasis{real} {\hyperref[\detokenize{units::doc}]{\sphinxcrossref{\DUrole{doc}{units}}}}.  You can use
the ReaxFF pair style with any LAMMPS units, but you would need to
create your own potential file with coefficients listed in the
appropriate units if your simulation does not use “real” units.


\subsection{Related commands}
\label{\detokenize{pair_reaxff:related-commands}}
\sphinxAtStartPar
{\hyperref[\detokenize{pair_coeff::doc}]{\sphinxcrossref{\DUrole{doc}{pair\_coeff}}}}, {\hyperref[\detokenize{fix_qeq_reaxff::doc}]{\sphinxcrossref{\DUrole{doc}{fix qeq/reaxff}}}},
{\hyperref[\detokenize{fix_acks2_reaxff::doc}]{\sphinxcrossref{\DUrole{doc}{fix acks2/reaxff}}}}, {\hyperref[\detokenize{fix_reaxff_bonds::doc}]{\sphinxcrossref{\DUrole{doc}{fix reaxff/bonds}}}},
{\hyperref[\detokenize{fix_reaxff_species::doc}]{\sphinxcrossref{\DUrole{doc}{fix reaxff/species}}}},
{\hyperref[\detokenize{compute_reaxff_atom::doc}]{\sphinxcrossref{\DUrole{doc}{compute reaxff/atom}}}}


\subsection{Default}
\label{\detokenize{pair_reaxff:default}}
\sphinxAtStartPar
The keyword defaults are checkqeq = yes, enobonds = yes, lgvdw = no, safezone =
1.2, mincap = 50, minhbonds = 25, tabulate = 0, list/blocking = yes on CPU, no
on GPU.


\bigskip\hrule\bigskip

\phantomsection\label{\detokenize{pair_reaxff:chenoweth-20082}}
\sphinxAtStartPar
\sphinxstylestrong{(Chenoweth\_2008)} Chenoweth, van Duin and Goddard,
Journal of Physical Chemistry A, 112, 1040\sphinxhyphen{}1053 (2008).

\phantomsection\label{\detokenize{pair_reaxff:aktulga}}
\sphinxAtStartPar
(Aktulga) Aktulga, Fogarty, Pandit, Grama, Parallel Computing, 38,
245\sphinxhyphen{}259 (2012).

\phantomsection\label{\detokenize{pair_reaxff:liu-2011}}
\sphinxAtStartPar
\sphinxstylestrong{(Liu)} L. Liu, Y. Liu, S. V. Zybin, H. Sun and W. A. Goddard, Journal
of Physical Chemistry A, 115, 11016\sphinxhyphen{}11022 (2011).

\sphinxstepscope

\index{pair\_style rebomos@\spxentry{pair\_style rebomos}}\ignorespaces 
\index{pair\_style rebomos/omp@\spxentry{pair\_style rebomos/omp}}\ignorespaces 

\section{pair\_style rebomos command}
\label{\detokenize{pair_rebomos:pair-style-rebomos-command}}\label{\detokenize{pair_rebomos:index-1}}\label{\detokenize{pair_rebomos:index-0}}\label{\detokenize{pair_rebomos::doc}}
\sphinxAtStartPar
Accelerator Variants: \sphinxstyleemphasis{rebomos/omp}


\subsection{Syntax}
\label{\detokenize{pair_rebomos:syntax}}
\begin{sphinxVerbatim}[commandchars=\\\{\}]
\PYG{k}{pair\PYGZus{}style}\PYG{+w}{ }\PYG{n}{rebomos}
\end{sphinxVerbatim}
\begin{itemize}
\item {} 
\sphinxAtStartPar
rebomos = name of this pair style

\end{itemize}


\subsection{Examples}
\label{\detokenize{pair_rebomos:examples}}
\begin{sphinxVerbatim}[commandchars=\\\{\}]
\PYG{k}{pair\PYGZus{}style}\PYG{+w}{ }\PYG{n}{rebomos}
\PYG{k}{pair\PYGZus{}coeff}\PYG{+w}{ }\PYG{o}{*}\PYG{+w}{ }\PYG{o}{*}\PYG{+w}{ }\PYG{n}{..}\PYG{o}{/}\PYG{n}{potentials}\PYG{o}{/}\PYG{n}{MoS.rebomos}\PYG{+w}{ }\PYG{n}{Mo}\PYG{+w}{ }\PYG{n}{S}
\end{sphinxVerbatim}

\sphinxAtStartPar
Example input scripts available: examples/threebody/


\subsection{Description}
\label{\detokenize{pair_rebomos:description}}
\sphinxAtStartPar
\DUrole{versionmodified,added}{Added in version 17Apr2024.}

\sphinxAtStartPar
The \sphinxstyleemphasis{rebomos} pair style computes the interactions between molybdenum
and sulfur atoms {\hyperref[\detokenize{pair_rebomos:stewart}]{\sphinxcrossref{\DUrole{std,std-ref}{(Stewart)}}}} utilizing an adaptive
interatomic reactive empirical bond order potential that is similar in
form to the AIREBO potential {\hyperref[\detokenize{pair_rebomos:stuart2}]{\sphinxcrossref{\DUrole{std,std-ref}{(Stuart)}}}}.  The potential
is based on an earlier parameterizations for \(\text{MoS}_2\)
developed by {\hyperref[\detokenize{pair_rebomos:liang}]{\sphinxcrossref{\DUrole{std,std-ref}{(Liang)}}}}.

\sphinxAtStartPar
The REBOMoS potential consists of two terms:
\begin{equation*}
\begin{split}E & = \frac{1}{2} \sum_i \sum_{j \neq i}
\left[ E^{\text{REBO}}_{ij} + E^{\text{LJ}}_{ij}  \right] \\\end{split}
\end{equation*}
\sphinxAtStartPar
The \(E^{\text{REBO}}\) term describes the covalently bonded
interactions between Mo and S atoms while the \(E^{\text{LJ}}\) term
describes longer range dispersion forces between layers.  A cubic spline
function is applied to smoothly switch between covalent bonding at short
distances to dispersion interactions at longer distances. This allows
the model to capture bond formation and breaking events which may occur
between adjacent MoS2 layers, edges, defects, and more.


\bigskip\hrule\bigskip


\sphinxAtStartPar
Only a single pair\_coeff command is used with the \sphinxstyleemphasis{rebomos} pair style
which specifies an REBOMoS potential file with parameters for Mo and S.
These are mapped to LAMMPS atom types by specifying N additional
arguments after the filename in the pair\_coeff command, where N is the
number of LAMMPS atom types:
\begin{itemize}
\item {} 
\sphinxAtStartPar
filename

\item {} 
\sphinxAtStartPar
\(N\) element names = mapping of REBOMoS elements to atom types

\end{itemize}

\sphinxAtStartPar
See the {\hyperref[\detokenize{pair_coeff::doc}]{\sphinxcrossref{\DUrole{doc}{pair\_coeff}}}} page for alternate ways
to specify the path for the potential file.

\sphinxAtStartPar
As an example, if your LAMMPS simulation has three atom types and you want
the first two to be Mo, and the third to be S, you would use the following
pair\_coeff command:

\begin{sphinxVerbatim}[commandchars=\\\{\}]
\PYG{k}{pair\PYGZus{}coeff}\PYG{+w}{ }\PYG{o}{*}\PYG{+w}{ }\PYG{o}{*}\PYG{+w}{ }\PYG{n}{MoS.rebomos}\PYG{+w}{ }\PYG{n}{Mo}\PYG{+w}{ }\PYG{n}{Mo}\PYG{+w}{ }\PYG{n}{S}
\end{sphinxVerbatim}

\sphinxAtStartPar
The first 2 arguments must be * * so as to span all LAMMPS atom types.
The first two Mo arguments map LAMMPS atom types 1 and 2 to the Mo
element in the REBOMoS file.  The final S argument maps LAMMPS atom type
3 to the S element in the REBOMoS file.  If a mapping value is specified
as NULL, the mapping is not performed.  This can be used when a
\sphinxstyleemphasis{rebomos} potential is used as part of the \sphinxstyleemphasis{hybrid} pair style.  The
NULL values are placeholders for atom types that will be used with other
potentials.


\bigskip\hrule\bigskip


\sphinxAtStartPar
Styles with a \sphinxstyleemphasis{gpu}, \sphinxstyleemphasis{intel}, \sphinxstyleemphasis{kk}, \sphinxstyleemphasis{omp}, or \sphinxstyleemphasis{opt} suffix are
functionally the same as the corresponding style without the suffix.
They have been optimized to run faster, depending on your available
hardware, as discussed on the {\hyperref[\detokenize{Speed_packages::doc}]{\sphinxcrossref{\DUrole{doc}{Accelerator packages}}}}
page.  The accelerated styles take the same arguments and should
produce the same results, except for round\sphinxhyphen{}off and precision issues.

\sphinxAtStartPar
These accelerated styles are part of the GPU, INTEL, KOKKOS,
OPENMP, and OPT packages, respectively.  They are only enabled if
LAMMPS was built with those packages.  See the {\hyperref[\detokenize{Build_package::doc}]{\sphinxcrossref{\DUrole{doc}{Build package}}}} page for more info.

\sphinxAtStartPar
You can specify the accelerated styles explicitly in your input script
by including their suffix, or you can use the {\hyperref[\detokenize{Run_options::doc}]{\sphinxcrossref{\DUrole{doc}{\sphinxhyphen{}suffix command\sphinxhyphen{}line switch}}}} when you invoke LAMMPS, or you can use the
{\hyperref[\detokenize{suffix::doc}]{\sphinxcrossref{\DUrole{doc}{suffix}}}} command in your input script.

\sphinxAtStartPar
See the {\hyperref[\detokenize{Speed_packages::doc}]{\sphinxcrossref{\DUrole{doc}{Accelerator packages}}}} page for more
instructions on how to use the accelerated styles effectively.


\bigskip\hrule\bigskip



\subsection{Mixing, shift, table, tail correction, restart, rRESPA info}
\label{\detokenize{pair_rebomos:mixing-shift-table-tail-correction-restart-rrespa-info}}
\sphinxAtStartPar
This pair style does not support the {\hyperref[\detokenize{pair_modify::doc}]{\sphinxcrossref{\DUrole{doc}{pair\_modify}}}}
mix, shift, table, and tail options.

\sphinxAtStartPar
This pair style does not write their information to {\hyperref[\detokenize{restart::doc}]{\sphinxcrossref{\DUrole{doc}{binary restart
files}}}}, since it is stored in potential files.  Thus, you need
to re\sphinxhyphen{}specify the pair\_style and pair\_coeff commands in an input script
that reads a restart file.

\sphinxAtStartPar
This pair styles can only be used via the \sphinxstyleemphasis{pair} keyword of the
{\hyperref[\detokenize{run_style::doc}]{\sphinxcrossref{\DUrole{doc}{run\_style respa}}}} command.  It does not support the
\sphinxstyleemphasis{inner}, \sphinxstyleemphasis{middle}, \sphinxstyleemphasis{outer} keywords.


\subsection{Restrictions}
\label{\detokenize{pair_rebomos:restrictions}}
\sphinxAtStartPar
This pair style is part of the MANYBODY package.  It is only enabled if
LAMMPS was built with that package.  See the {\hyperref[\detokenize{Build_package::doc}]{\sphinxcrossref{\DUrole{doc}{Build package}}}} page for more info.

\sphinxAtStartPar
These pair potentials require the {\hyperref[\detokenize{newton::doc}]{\sphinxcrossref{\DUrole{doc}{newton}}}} setting to be
“on” for pair interactions.

\sphinxAtStartPar
The MoS.rebomos potential file provided with LAMMPS (see the potentials
directory) is parameterized for metal {\hyperref[\detokenize{units::doc}]{\sphinxcrossref{\DUrole{doc}{units}}}}.  You can use
the \sphinxstyleemphasis{rebomos} pair style with any LAMMPS units setting, but you would
need to create your own REBOMoS potential file with coefficients listed
in the appropriate units.

\sphinxAtStartPar
The pair style provided here \sphinxstylestrong{only} supports potential files parameterized
for the elements molybdenum and sulfur (designated with “Mo” and “S” in the
\sphinxstyleemphasis{pair\_coeff} command.  Using potential files for other elements will trigger
an error.


\subsection{Related commands}
\label{\detokenize{pair_rebomos:related-commands}}
\sphinxAtStartPar
{\hyperref[\detokenize{pair_coeff::doc}]{\sphinxcrossref{\DUrole{doc}{pair\_coeff}}}}, {\hyperref[\detokenize{pair_airebo::doc}]{\sphinxcrossref{\DUrole{doc}{pair style rebo}}}}


\subsection{Default}
\label{\detokenize{pair_rebomos:default}}
\sphinxAtStartPar
none


\bigskip\hrule\bigskip

\phantomsection\label{\detokenize{pair_rebomos:stewart}}
\sphinxAtStartPar
\sphinxstylestrong{(Steward)}  Stewart, Spearot, Modelling Simul. Mater. Sci. Eng. 21, 045003, (2013).

\phantomsection\label{\detokenize{pair_rebomos:stuart2}}
\sphinxAtStartPar
\sphinxstylestrong{(Stuart)} Stuart, Tutein, Harrison, J Chem Phys, 112, 6472\sphinxhyphen{}6486, (2000).

\phantomsection\label{\detokenize{pair_rebomos:liang}}
\sphinxAtStartPar
\sphinxstylestrong{(Liang)}  Liang, Phillpot, Sinnott Phys. Rev. B79 245110, (2009), Erratum: Phys. Rev. B85 199903(E), (2012)

\sphinxstepscope

\index{pair\_style resquared@\spxentry{pair\_style resquared}}\ignorespaces 
\index{pair\_style resquared/gpu@\spxentry{pair\_style resquared/gpu}}\ignorespaces 
\index{pair\_style resquared/omp@\spxentry{pair\_style resquared/omp}}\ignorespaces 

\section{pair\_style resquared command}
\label{\detokenize{pair_resquared:pair-style-resquared-command}}\label{\detokenize{pair_resquared:index-2}}\label{\detokenize{pair_resquared:index-1}}\label{\detokenize{pair_resquared:index-0}}\label{\detokenize{pair_resquared::doc}}
\sphinxAtStartPar
Accelerator Variants: \sphinxstyleemphasis{resquared/gpu}, \sphinxstyleemphasis{resquared/omp}


\subsection{Syntax}
\label{\detokenize{pair_resquared:syntax}}
\begin{sphinxVerbatim}[commandchars=\\\{\}]
\PYG{k}{pair\PYGZus{}style}\PYG{+w}{ }\PYG{n}{resquared}\PYG{+w}{ }\PYG{n}{cutoff}
\end{sphinxVerbatim}
\begin{itemize}
\item {} 
\sphinxAtStartPar
cutoff = global cutoff for interactions (distance units)

\end{itemize}


\subsection{Examples}
\label{\detokenize{pair_resquared:examples}}
\begin{sphinxVerbatim}[commandchars=\\\{\}]
\PYG{k}{pair\PYGZus{}style}\PYG{+w}{ }\PYG{n}{resquared}\PYG{+w}{ }\PYG{l+m}{10.0}
\PYG{k}{pair\PYGZus{}coeff}\PYG{+w}{ }\PYG{o}{*}\PYG{+w}{ }\PYG{o}{*}\PYG{+w}{ }\PYG{l+m}{1.0}\PYG{+w}{ }\PYG{l+m}{1.0}\PYG{+w}{ }\PYG{l+m}{1.7}\PYG{+w}{ }\PYG{l+m}{3.4}\PYG{+w}{ }\PYG{l+m}{3.4}\PYG{+w}{ }\PYG{l+m}{1.0}\PYG{+w}{ }\PYG{l+m}{1.0}\PYG{+w}{ }\PYG{l+m}{1.0}
\end{sphinxVerbatim}


\subsection{Description}
\label{\detokenize{pair_resquared:description}}
\sphinxAtStartPar
Style \sphinxstyleemphasis{resquared} computes the RE\sphinxhyphen{}squared anisotropic interaction
{\hyperref[\detokenize{pair_resquared:everaers3}]{\sphinxcrossref{\DUrole{std,std-ref}{(Everaers)}}}}, {\hyperref[\detokenize{pair_resquared:babadi}]{\sphinxcrossref{\DUrole{std,std-ref}{(Babadi)}}}} between pairs of
ellipsoidal and/or spherical Lennard\sphinxhyphen{}Jones particles.  For ellipsoidal
interactions, the potential considers the ellipsoid as being comprised
of small spheres of size \(\sigma\).  LJ particles are a single sphere of
size \(\sigma\).  The distinction is made to allow the pair style to make
efficient calculations of ellipsoid/solvent interactions.

\sphinxAtStartPar
Details for the equations used are given in the references below and
in \sphinxhref{PDF/pair\_resquared\_extra.pdf}{this supplementary document}.

\sphinxAtStartPar
Use of this pair style requires the NVE, NVT, or NPT fixes with the
\sphinxstyleemphasis{asphere} extension (e.g. {\hyperref[\detokenize{fix_nve_asphere::doc}]{\sphinxcrossref{\DUrole{doc}{fix nve/asphere}}}}) in
order to integrate particle rotation.  Additionally, {\hyperref[\detokenize{atom_style::doc}]{\sphinxcrossref{\DUrole{doc}{atom\_style ellipsoid}}}} should be used since it defines the
rotational state and the size and shape of each ellipsoidal particle.

\sphinxAtStartPar
The following coefficients must be defined for each pair of atoms
types via the {\hyperref[\detokenize{pair_coeff::doc}]{\sphinxcrossref{\DUrole{doc}{pair\_coeff}}}} command as in the examples
above, or in the data file or restart files read by the
{\hyperref[\detokenize{read_data::doc}]{\sphinxcrossref{\DUrole{doc}{read\_data}}}} or {\hyperref[\detokenize{read_restart::doc}]{\sphinxcrossref{\DUrole{doc}{read\_restart}}}}
commands:
\begin{itemize}
\item {} 
\sphinxAtStartPar
A12 = Energy Prefactor/Hamaker constant (energy units)

\item {} 
\sphinxAtStartPar
\(\sigma\) = atomic interaction diameter (distance units)

\item {} 
\sphinxAtStartPar
\(\epsilon_{i,a}\) = relative well depth of type I for side\sphinxhyphen{}to\sphinxhyphen{}side interactions

\item {} 
\sphinxAtStartPar
\(\epsilon_{i,b}\) = relative well depth of type I for face\sphinxhyphen{}to\sphinxhyphen{}face interactions

\item {} 
\sphinxAtStartPar
\(\epsilon_{i,c}\) = relative well depth of type I for end\sphinxhyphen{}to\sphinxhyphen{}end interactions

\item {} 
\sphinxAtStartPar
\(\epsilon_{j,a}\) = relative well depth of type J for side\sphinxhyphen{}to\sphinxhyphen{}side interactions

\item {} 
\sphinxAtStartPar
\(\epsilon_{j,b}\) = relative well depth of type J for face\sphinxhyphen{}to\sphinxhyphen{}face interactions

\item {} 
\sphinxAtStartPar
\(\epsilon_{j,c}\) = relative well depth of type J for end\sphinxhyphen{}to\sphinxhyphen{}end interactions

\item {} 
\sphinxAtStartPar
cutoff (distance units)

\end{itemize}

\sphinxAtStartPar
The parameters used depend on the type of the interacting particles,
i.e. ellipsoids or LJ spheres.  The type of a particle is determined
by the diameters specified for its 3 shape parameters.  If all 3 shape
parameters = 0.0, then the particle is treated as an LJ sphere.  The
\(\epsilon_{i,*}\) or \(\epsilon_{j,*}\) parameters are ignored
for LJ spheres.  If
the 3 shape parameters are \textgreater{} 0.0, then the particle is treated as an
ellipsoid (even if the 3 parameters are equal to each other).

\sphinxAtStartPar
A12 specifies the energy prefactor which depends on the types of the
two interacting particles.

\sphinxAtStartPar
For ellipsoid/ellipsoid interactions, the interaction is computed by
the formulas in the supplementary document referenced above.  A12 is
the Hamaker constant as described in {\hyperref[\detokenize{pair_resquared:everaers3}]{\sphinxcrossref{\DUrole{std,std-ref}{(Everaers)}}}}. In LJ
units:
\begin{equation*}
\begin{split}A_{12} = 4\pi^2\epsilon_{\mathrm{LJ}}(\rho\sigma^3)^2\end{split}
\end{equation*}
\sphinxAtStartPar
where \(\rho\) gives the number density of the spherical particles
composing the ellipsoids and \(\epsilon_{\mathrm{LJ}}\) determines
the interaction strength of the spherical particles.

\sphinxAtStartPar
For ellipsoid/LJ sphere interactions, the interaction is also computed
by the formulas in the supplementary document referenced above.  A12
has a modified form (see \sphinxhref{PDF/pair\_resquared\_extra.pdf}{here} for
details):
\begin{equation*}
\begin{split}A_{12} = 4\pi^2\epsilon_{\mathrm{LJ}}(\rho\sigma^3)\end{split}
\end{equation*}
\sphinxAtStartPar
For ellipsoid/LJ sphere interactions, a correction to the distance\sphinxhyphen{}
of\sphinxhyphen{}closest approach equation has been implemented to reduce the error
from two particles of disparate sizes; see \sphinxhref{PDF/pair\_resquared\_extra.pdf}{this supplementary document}.

\sphinxAtStartPar
For LJ sphere/LJ sphere interactions, the interaction is computed
using the standard Lennard\sphinxhyphen{}Jones formula, which is much cheaper to
compute than the ellipsoidal formulas.  A12 is used as epsilon in the
standard LJ formula:
\begin{equation*}
\begin{split}A_{12} = \epsilon_{\mathrm{LJ}}\end{split}
\end{equation*}
\sphinxAtStartPar
and the specified \(\sigma\) is used as the \(\sigma\) in the
standard LJ formula.

\sphinxAtStartPar
When one of both of the interacting particles are ellipsoids, then
\(\sigma\) specifies the diameter of the continuous distribution of
constituent particles within each ellipsoid used to model the RE\sphinxhyphen{}squared
potential.  Note that this is a different meaning for \(\sigma\)
than the {\hyperref[\detokenize{pair_gayberne::doc}]{\sphinxcrossref{\DUrole{doc}{pair\_style gayberne}}}} potential uses.

\sphinxAtStartPar
The \(\epsilon_i\) and \(\epsilon_j\) coefficients are defined
for atom types, not for pairs of atom types.  Thus, in a series of
pair\_coeff commands, they only need to be specified once for each atom
type.

\sphinxAtStartPar
Specifically, if any of \(\epsilon_{i,a}\), \(\epsilon_{i,b}\),
\(\epsilon_{i,c}\) are non\sphinxhyphen{}zero, the three values are assigned to
atom type I.  If all the \(\epsilon_i\) values are zero, they are
ignored.  If any of \(\epsilon_{j,a}\), \(\epsilon_{j,b}\),
\(\epsilon_{j,c}\) are non\sphinxhyphen{}zero, the three values are assigned to
atom type J.  If all three \(\epsilon_i\) values are zero, they are
ignored.  Thus the typical way to define the \(\epsilon_i\) and
\(\epsilon_j\) coefficients is to list their values in “pair\_coeff
I J” commands when I = J, but set them to 0.0 when I != J.  If you do
list them when I != J, you should ensure they are consistent with their
values in other pair\_coeff commands.

\sphinxAtStartPar
Note that if this potential is being used as a sub\sphinxhyphen{}style of
{\hyperref[\detokenize{pair_hybrid::doc}]{\sphinxcrossref{\DUrole{doc}{pair\_style hybrid}}}}, and there is no “pair\_coeff I I”
setting made for RE\sphinxhyphen{}squared for a particular type I (because I\sphinxhyphen{}I
interactions are computed by another hybrid pair potential), then you
still need to ensure the epsilon a,b,c coefficients are assigned to
that type in a “pair\_coeff I J” command.

\sphinxAtStartPar
For large uniform molecules it has been shown that the \(\epsilon_{*,*}\)
energy parameters are approximately representable in terms of local
contact curvatures {\hyperref[\detokenize{pair_resquared:everaers3}]{\sphinxcrossref{\DUrole{std,std-ref}{(Everaers)}}}}:
\begin{equation*}
\begin{split}\epsilon_a = \sigma \cdot { \frac{a}{ b \cdot c } }; \epsilon_b =
\sigma \cdot { \frac{b}{ a \cdot c } }; \epsilon_c = \sigma \cdot {
\frac{c}{ a \cdot b } }\end{split}
\end{equation*}
\sphinxAtStartPar
where a, b, and c give the particle diameters.

\sphinxAtStartPar
The last coefficient is optional.  If not specified, the global cutoff
specified in the pair\_style command is used.


\bigskip\hrule\bigskip


\sphinxAtStartPar
Styles with a \sphinxstyleemphasis{gpu}, \sphinxstyleemphasis{intel}, \sphinxstyleemphasis{kk}, \sphinxstyleemphasis{omp}, or \sphinxstyleemphasis{opt} suffix are
functionally the same as the corresponding style without the suffix.
They have been optimized to run faster, depending on your available
hardware, as discussed on the {\hyperref[\detokenize{Speed_packages::doc}]{\sphinxcrossref{\DUrole{doc}{Accelerator packages}}}}
page.  The accelerated styles take the same arguments and should
produce the same results, except for round\sphinxhyphen{}off and precision issues.

\sphinxAtStartPar
These accelerated styles are part of the GPU, INTEL, KOKKOS,
OPENMP, and OPT packages, respectively.  They are only enabled if
LAMMPS was built with those packages.  See the {\hyperref[\detokenize{Build_package::doc}]{\sphinxcrossref{\DUrole{doc}{Build package}}}} page for more info.

\sphinxAtStartPar
You can specify the accelerated styles explicitly in your input script
by including their suffix, or you can use the {\hyperref[\detokenize{Run_options::doc}]{\sphinxcrossref{\DUrole{doc}{\sphinxhyphen{}suffix command\sphinxhyphen{}line switch}}}} when you invoke LAMMPS, or you can use the
{\hyperref[\detokenize{suffix::doc}]{\sphinxcrossref{\DUrole{doc}{suffix}}}} command in your input script.

\sphinxAtStartPar
See the {\hyperref[\detokenize{Speed_packages::doc}]{\sphinxcrossref{\DUrole{doc}{Accelerator packages}}}} page for more
instructions on how to use the accelerated styles effectively.


\bigskip\hrule\bigskip



\subsection{Mixing, shift, table, tail correction, restart, rRESPA info}
\label{\detokenize{pair_resquared:mixing-shift-table-tail-correction-restart-rrespa-info}}
\sphinxAtStartPar
For atom type pairs I,J and I != J, the epsilon and sigma coefficients
and cutoff distance can be mixed, but only for sphere pairs.  The
default mix value is \sphinxstyleemphasis{geometric}.  See the “pair\_modify” command for
details.  Other type pairs cannot be mixed, due to the different
meanings of the energy prefactors used to calculate the interactions
and the implicit dependence of the ellipsoid\sphinxhyphen{}sphere interaction on the
equation for the Hamaker constant presented here.  Mixing of sigma and
epsilon followed by calculation of the energy prefactors using the
equations above is recommended.

\sphinxAtStartPar
This pair style supports the {\hyperref[\detokenize{pair_modify::doc}]{\sphinxcrossref{\DUrole{doc}{pair\_modify}}}} shift
option for the energy of the Lennard\sphinxhyphen{}Jones portion of the pair
interaction, but only for sphere\sphinxhyphen{}sphere interactions.  There is no
shifting performed for ellipsoidal interactions due to the anisotropic
dependence of the interaction.

\sphinxAtStartPar
The {\hyperref[\detokenize{pair_modify::doc}]{\sphinxcrossref{\DUrole{doc}{pair\_modify}}}} table option is not relevant
for this pair style.

\sphinxAtStartPar
This pair style does not support the {\hyperref[\detokenize{pair_modify::doc}]{\sphinxcrossref{\DUrole{doc}{pair\_modify}}}}
tail option for adding long\sphinxhyphen{}range tail corrections to energy and
pressure.

\sphinxAtStartPar
This pair style writes its information to {\hyperref[\detokenize{restart::doc}]{\sphinxcrossref{\DUrole{doc}{binary restart files}}}}, so pair\_style and pair\_coeff commands do not need
to be specified in an input script that reads a restart file.

\sphinxAtStartPar
This pair style can only be used via the \sphinxstyleemphasis{pair} keyword of the
{\hyperref[\detokenize{run_style::doc}]{\sphinxcrossref{\DUrole{doc}{run\_style respa}}}} command.  It does not support the
\sphinxstyleemphasis{inner}, \sphinxstyleemphasis{middle}, \sphinxstyleemphasis{outer} keywords of the {\hyperref[\detokenize{run_style::doc}]{\sphinxcrossref{\DUrole{doc}{run\_style command}}}}.


\bigskip\hrule\bigskip



\subsection{Restrictions}
\label{\detokenize{pair_resquared:restrictions}}
\sphinxAtStartPar
This style is part of the ASPHERE package.  It is only enabled if
LAMMPS was built with that package.  See the {\hyperref[\detokenize{Build_package::doc}]{\sphinxcrossref{\DUrole{doc}{Build package}}}} page for more info.

\sphinxAtStartPar
This pair style requires that atoms be ellipsoids as defined by the
{\hyperref[\detokenize{atom_style::doc}]{\sphinxcrossref{\DUrole{doc}{atom\_style ellipsoid}}}} command.

\sphinxAtStartPar
Particles acted on by the potential can be finite\sphinxhyphen{}size aspherical or
spherical particles, or point particles.  Spherical particles have all
3 of their shape parameters equal to each other.  Point particles have
all 3 of their shape parameters equal to 0.0.

\sphinxAtStartPar
The distance\sphinxhyphen{}of\sphinxhyphen{}closest\sphinxhyphen{}approach approximation used by LAMMPS becomes
less accurate when high\sphinxhyphen{}aspect ratio ellipsoids are used.


\subsection{Related commands}
\label{\detokenize{pair_resquared:related-commands}}
\sphinxAtStartPar
{\hyperref[\detokenize{pair_coeff::doc}]{\sphinxcrossref{\DUrole{doc}{pair\_coeff}}}}, {\hyperref[\detokenize{fix_nve_asphere::doc}]{\sphinxcrossref{\DUrole{doc}{fix nve/asphere}}}},
{\hyperref[\detokenize{compute_temp_asphere::doc}]{\sphinxcrossref{\DUrole{doc}{compute temp/asphere}}}}, {\hyperref[\detokenize{pair_gayberne::doc}]{\sphinxcrossref{\DUrole{doc}{pair\_style gayberne}}}}


\subsection{Default}
\label{\detokenize{pair_resquared:default}}
\sphinxAtStartPar
none


\bigskip\hrule\bigskip

\phantomsection\label{\detokenize{pair_resquared:everaers3}}
\sphinxAtStartPar
\sphinxstylestrong{(Everaers)} Everaers and Ejtehadi, Phys Rev E, 67, 041710 (2003).

\phantomsection\label{\detokenize{pair_resquared:babadi}}
\sphinxAtStartPar
\sphinxstylestrong{(Babadi)} Babadi, Ejtehadi, Everaers, J Comp Phys, 219, 770\sphinxhyphen{}779 (2006).

\sphinxstepscope

\index{pair\_style rheo@\spxentry{pair\_style rheo}}\ignorespaces 

\section{pair\_style rheo command}
\label{\detokenize{pair_rheo:pair-style-rheo-command}}\label{\detokenize{pair_rheo:index-0}}\label{\detokenize{pair_rheo::doc}}

\subsection{Syntax}
\label{\detokenize{pair_rheo:syntax}}
\begin{sphinxVerbatim}[commandchars=\\\{\}]
\PYG{k}{pair\PYGZus{}style}\PYG{+w}{ }\PYG{n}{rheo}\PYG{+w}{ }\PYG{n}{cutoff}\PYG{+w}{ }\PYG{n}{keyword}\PYG{+w}{ }\PYG{n}{values}
\end{sphinxVerbatim}
\begin{itemize}
\item {} 
\sphinxAtStartPar
cutoff = global cutoff for kernel (distance units)

\item {} 
\sphinxAtStartPar
zero or more keyword/value pairs may be appended to args

\item {} 
\sphinxAtStartPar
keyword = \sphinxstyleemphasis{rho/damp} or \sphinxstyleemphasis{artificial/visc} or \sphinxstyleemphasis{harmonic/means}

\end{itemize}
\begin{sphinxalltt}
\sphinxstyleemphasis{rho/damp} args = density damping prefactor \(\xi\)
\sphinxstyleemphasis{artificial/visc} args = artificial viscosity prefactor \(\zeta\)
\sphinxstyleemphasis{harmonic/means} args = none
\end{sphinxalltt}


\subsection{Examples}
\label{\detokenize{pair_rheo:examples}}
\begin{sphinxVerbatim}[commandchars=\\\{\}]
\PYG{k}{pair\PYGZus{}style}\PYG{+w}{ }\PYG{n}{rheo}\PYG{+w}{ }\PYG{l+m}{3.0}\PYG{+w}{ }\PYG{n}{rho}\PYG{o}{/}\PYG{n}{damp}\PYG{+w}{ }\PYG{l+m}{1.0}\PYG{+w}{ }\PYG{n}{artificial}\PYG{o}{/}\PYG{n}{visc}\PYG{+w}{ }\PYG{l+m}{2.0}
\PYG{k}{pair\PYGZus{}coeff}\PYG{+w}{ }\PYG{o}{*}\PYG{+w}{ }\PYG{o}{*}
\end{sphinxVerbatim}


\subsection{Description}
\label{\detokenize{pair_rheo:description}}
\sphinxAtStartPar
\DUrole{versionmodified,added}{Added in version 29Aug2024.}

\sphinxAtStartPar
Pair style \sphinxstyleemphasis{rheo} computes pressure and viscous forces between particles
in the {\hyperref[\detokenize{Howto_rheo::doc}]{\sphinxcrossref{\DUrole{doc}{rheo package}}}}. If thermal evolution is turned
on in {\hyperref[\detokenize{fix_rheo::doc}]{\sphinxcrossref{\DUrole{doc}{fix rheo}}}}, then the pair style also calculates
heat exchanged between particles.

\sphinxAtStartPar
The \sphinxstyleemphasis{artificial/viscosity} keyword is used to specify the magnitude
\(\zeta\) of an optional artificial viscosity contribution to forces.
This factor can help stabilize simulations by smoothing out small length
scale variations in velocity fields. Artificial viscous forces typically
are only exchanged by fluid particles. However, if interfaces are not
reconstructed in fix rheo, fluid particles will also exchange artificial
viscous forces with solid particles to improve stability.

\sphinxAtStartPar
The \sphinxstyleemphasis{rho/damp} keyword is used to specify the magnitude \(\xi\) of
an optional pairwise damping term between the density of particles. This
factor can help stabilize simulations by smoothing out small length
scale variations in density fields. However, in systems that develop
a density gradient in equilibrium (e.g. in a hydrostatic column underlying
gravity), this option may be inappropriate.

\sphinxAtStartPar
If particles have different viscosities or conductivities, the
\sphinxstyleemphasis{harmonic/means} keyword changes how they are averaged before calculating
pairwise forces or heat exchanges. By default, an arithmetic averaged is
used, however, a harmonic mean may improve stability in systems with multiple
fluid phases with large disparities in viscosities.

\sphinxAtStartPar
No coefficients are defined for each pair of atoms types via the
{\hyperref[\detokenize{pair_coeff::doc}]{\sphinxcrossref{\DUrole{doc}{pair\_coeff}}}} command as in the examples
above.


\bigskip\hrule\bigskip



\subsection{Mixing, shift, table, tail correction, restart, rRESPA info}
\label{\detokenize{pair_rheo:mixing-shift-table-tail-correction-restart-rrespa-info}}
\sphinxAtStartPar
This style does not support the {\hyperref[\detokenize{pair_modify::doc}]{\sphinxcrossref{\DUrole{doc}{pair\_modify}}}}
shift, table, and tail options.

\sphinxAtStartPar
This style does not write information to {\hyperref[\detokenize{restart::doc}]{\sphinxcrossref{\DUrole{doc}{binary restart files}}}}.
Thus, you need to re\sphinxhyphen{}specify the pair\_style and pair\_coeff commands in an input
script that reads a restart file.

\sphinxAtStartPar
This style can only be used via the \sphinxstyleemphasis{pair} keyword of the
{\hyperref[\detokenize{run_style::doc}]{\sphinxcrossref{\DUrole{doc}{run\_style respa}}}} command.  It does not support the \sphinxstyleemphasis{inner},
\sphinxstyleemphasis{middle}, \sphinxstyleemphasis{outer} keywords.


\subsection{Restrictions}
\label{\detokenize{pair_rheo:restrictions}}
\sphinxAtStartPar
This fix is part of the RHEO package.  It is only enabled if
LAMMPS was built with that package.  See the
{\hyperref[\detokenize{Build_package::doc}]{\sphinxcrossref{\DUrole{doc}{Build package}}}} page for more info.


\subsection{Related commands}
\label{\detokenize{pair_rheo:related-commands}}
\sphinxAtStartPar
{\hyperref[\detokenize{fix_rheo::doc}]{\sphinxcrossref{\DUrole{doc}{fix rheo}}}},
{\hyperref[\detokenize{fix_rheo_pressure::doc}]{\sphinxcrossref{\DUrole{doc}{fix rheo/pressure}}}},
{\hyperref[\detokenize{fix_rheo_thermal::doc}]{\sphinxcrossref{\DUrole{doc}{fix rheo/thermal}}}},
{\hyperref[\detokenize{fix_rheo_viscosity::doc}]{\sphinxcrossref{\DUrole{doc}{fix rheo/viscosity}}}},
{\hyperref[\detokenize{compute_rheo_property_atom::doc}]{\sphinxcrossref{\DUrole{doc}{compute rheo/property/atom}}}}


\subsection{Default}
\label{\detokenize{pair_rheo:default}}
\sphinxAtStartPar
Density damping and artificial viscous forces are not calculated.
Arithmetic means are used for mixing particle properties.

\sphinxstepscope

\index{pair\_style rheo/solid@\spxentry{pair\_style rheo/solid}}\ignorespaces 

\section{pair\_style rheo/solid command}
\label{\detokenize{pair_rheo_solid:pair-style-rheo-solid-command}}\label{\detokenize{pair_rheo_solid:index-0}}\label{\detokenize{pair_rheo_solid::doc}}

\subsection{Syntax}
\label{\detokenize{pair_rheo_solid:syntax}}
\begin{sphinxVerbatim}[commandchars=\\\{\}]
\PYG{k}{pair\PYGZus{}style}\PYG{+w}{ }\PYG{n}{rheo}\PYG{o}{/}\PYG{n}{solid}
\end{sphinxVerbatim}


\subsection{Examples}
\label{\detokenize{pair_rheo_solid:examples}}
\begin{sphinxVerbatim}[commandchars=\\\{\}]
\PYG{k}{pair\PYGZus{}style}\PYG{+w}{ }\PYG{n}{rheo}\PYG{o}{/}\PYG{n}{solid}
\PYG{k}{pair\PYGZus{}coeff}\PYG{+w}{ }\PYG{o}{*}\PYG{+w}{ }\PYG{o}{*}\PYG{+w}{ }\PYG{l+m}{1.0}\PYG{+w}{ }\PYG{l+m}{1.5}\PYG{+w}{ }\PYG{l+m}{1.0}
\end{sphinxVerbatim}


\subsection{Description}
\label{\detokenize{pair_rheo_solid:description}}
\sphinxAtStartPar
\DUrole{versionmodified,added}{Added in version 29Aug2024.}

\sphinxAtStartPar
Style \sphinxstyleemphasis{rheo/solid} is effectively a copy of pair style
{\hyperref[\detokenize{pair_bpm_spring::doc}]{\sphinxcrossref{\DUrole{doc}{bpm/spring}}}} except it only applies forces
between solid RHEO particles, determined by checking the status of
each pair of neighboring particles before calculating forces.

\sphinxAtStartPar
The style computes pairwise forces with the formula
\begin{equation*}
\begin{split}F = k (r - r_c)\end{split}
\end{equation*}
\sphinxAtStartPar
where \(k\) is a stiffness and \(r_c\) is the cutoff length.
An additional damping force is also applied to interacting
particles. The force is proportional to the difference in the
normal velocity of particles
\begin{equation*}
\begin{split}F_D = - \gamma w (\hat{r} \bullet \vec{v})\end{split}
\end{equation*}
\sphinxAtStartPar
where \(\gamma\) is the damping strength, \(\hat{r}\) is the
displacement normal vector, \(\vec{v}\) is the velocity difference
between the two particles, and \(w\) is a smoothing factor.
This smoothing factor is constructed such that damping forces go to zero
as particles come out of contact to avoid discontinuities. It is
given by
\begin{equation*}
\begin{split}w = 1.0 - \left( \frac{r}{r_c} \right)^8 .\end{split}
\end{equation*}
\sphinxAtStartPar
The following coefficients must be defined for each pair of atom types
via the {\hyperref[\detokenize{pair_coeff::doc}]{\sphinxcrossref{\DUrole{doc}{pair\_coeff}}}} command as in the examples
above, or in the data file or restart files read by the
{\hyperref[\detokenize{read_data::doc}]{\sphinxcrossref{\DUrole{doc}{read\_data}}}} or {\hyperref[\detokenize{read_restart::doc}]{\sphinxcrossref{\DUrole{doc}{read\_restart}}}}
commands, or by mixing as described below:
\begin{itemize}
\item {} 
\sphinxAtStartPar
\(k\)             (force/distance units)

\item {} 
\sphinxAtStartPar
\(r_c\)           (distance units)

\item {} 
\sphinxAtStartPar
\(\gamma\)        (force/velocity units)

\end{itemize}


\bigskip\hrule\bigskip



\subsection{Mixing, shift, table, tail correction, restart, rRESPA info}
\label{\detokenize{pair_rheo_solid:mixing-shift-table-tail-correction-restart-rrespa-info}}
\sphinxAtStartPar
For atom type pairs I,J and I != J, the A coefficient and cutoff
distance for this pair style can be mixed.  A is always mixed via a
\sphinxstyleemphasis{geometric} rule.  The cutoff is mixed according to the pair\_modify
mix value.  The default mix value is \sphinxstyleemphasis{geometric}.  See the
“pair\_modify” command for details.

\sphinxAtStartPar
This pair style does not support the {\hyperref[\detokenize{pair_modify::doc}]{\sphinxcrossref{\DUrole{doc}{pair\_modify}}}}
shift option, since the pair interaction goes to 0.0 at the cutoff.

\sphinxAtStartPar
The {\hyperref[\detokenize{pair_modify::doc}]{\sphinxcrossref{\DUrole{doc}{pair\_modify}}}} table and tail options are not
relevant for this pair style.

\sphinxAtStartPar
This pair style writes its information to {\hyperref[\detokenize{restart::doc}]{\sphinxcrossref{\DUrole{doc}{binary restart files}}}}, so pair\_style and pair\_coeff commands do not need to be
specified in an input script that reads a restart file.

\sphinxAtStartPar
This pair style can only be used via the \sphinxstyleemphasis{pair} keyword of the
{\hyperref[\detokenize{run_style::doc}]{\sphinxcrossref{\DUrole{doc}{run\_style respa}}}} command.  It does not support the
\sphinxstyleemphasis{inner}, \sphinxstyleemphasis{middle}, \sphinxstyleemphasis{outer} keywords.


\bigskip\hrule\bigskip



\subsection{Restrictions}
\label{\detokenize{pair_rheo_solid:restrictions}}
\sphinxAtStartPar
This pair style is part of the RHEO package.  It is only enabled if
LAMMPS was built with that package.  See the {\hyperref[\detokenize{Build_package::doc}]{\sphinxcrossref{\DUrole{doc}{Build package}}}} page for more info.


\subsection{Related commands}
\label{\detokenize{pair_rheo_solid:related-commands}}
\sphinxAtStartPar
{\hyperref[\detokenize{fix_rheo::doc}]{\sphinxcrossref{\DUrole{doc}{fix rheo}}}},
{\hyperref[\detokenize{fix_rheo_thermal::doc}]{\sphinxcrossref{\DUrole{doc}{fix rheo/thermal}}}},
{\hyperref[\detokenize{pair_bpm_spring::doc}]{\sphinxcrossref{\DUrole{doc}{pair bpm/spring}}}}


\subsection{Default}
\label{\detokenize{pair_rheo_solid:default}}
\sphinxAtStartPar
none

\sphinxstepscope

\index{pair\_style saip/metal@\spxentry{pair\_style saip/metal}}\ignorespaces 
\index{pair\_style saip/metal/opt@\spxentry{pair\_style saip/metal/opt}}\ignorespaces 

\section{pair\_style saip/metal command}
\label{\detokenize{pair_saip_metal:pair-style-saip-metal-command}}\label{\detokenize{pair_saip_metal:index-1}}\label{\detokenize{pair_saip_metal:index-0}}\label{\detokenize{pair_saip_metal::doc}}
\sphinxAtStartPar
Accelerator Variant: \sphinxstyleemphasis{saip/metal/opt}


\subsection{Syntax}
\label{\detokenize{pair_saip_metal:syntax}}
\begin{sphinxVerbatim}[commandchars=\\\{\}]
\PYG{k}{pair\PYGZus{}style}\PYG{+w}{ }\PYG{n}{[hybrid}\PYG{o}{/}\PYG{n}{overlay}\PYG{+w}{ }\PYG{n}{...]}\PYG{+w}{ }\PYG{n}{saip}\PYG{o}{/}\PYG{n}{metal}\PYG{+w}{ }\PYG{n}{cutoff}\PYG{+w}{ }\PYG{n}{tap\PYGZus{}flag}
\end{sphinxVerbatim}
\begin{itemize}
\item {} 
\sphinxAtStartPar
cutoff = global cutoff (distance units)

\item {} 
\sphinxAtStartPar
tap\_flag = 0/1 to turn off/on the taper function

\end{itemize}


\subsection{Examples}
\label{\detokenize{pair_saip_metal:examples}}
\begin{sphinxVerbatim}[commandchars=\\\{\}]
\PYG{k}{pair\PYGZus{}style}\PYG{+w}{  }\PYG{n}{hybrid}\PYG{o}{/}\PYG{n}{overlay}\PYG{+w}{ }\PYG{n}{saip}\PYG{o}{/}\PYG{n}{metal}\PYG{+w}{ }\PYG{l+m}{16.0}\PYG{+w}{ }\PYG{l+m}{1}
\PYG{k}{pair\PYGZus{}coeff}\PYG{+w}{  }\PYG{o}{*}\PYG{+w}{ }\PYG{o}{*}\PYG{+w}{ }\PYG{n}{saip}\PYG{o}{/}\PYG{n}{metal}\PYG{+w}{ }\PYG{n}{CHAu.ILP}\PYG{+w}{ }\PYG{n}{Au}\PYG{+w}{ }\PYG{n}{C}\PYG{+w}{ }\PYG{n}{H}

\PYG{k}{pair\PYGZus{}style}\PYG{+w}{  }\PYG{n}{hybrid}\PYG{o}{/}\PYG{n}{overlay}\PYG{+w}{ }\PYG{n}{eam}\PYG{+w}{ }\PYG{n}{rebo}\PYG{+w}{ }\PYG{n}{saip}\PYG{o}{/}\PYG{n}{metal}\PYG{+w}{ }\PYG{l+m}{16.0}
\PYG{k}{pair\PYGZus{}coeff}\PYG{+w}{  }\PYG{l+m}{1}\PYG{+w}{ }\PYG{l+m}{1}\PYG{+w}{ }\PYG{n}{eam}\PYG{+w}{  }\PYG{n}{Au\PYGZus{}u3.eam}\PYG{+w}{  }\PYG{n}{Au}\PYG{+w}{ }\PYG{n}{NULL}\PYG{+w}{ }\PYG{n}{NULL}
\PYG{k}{pair\PYGZus{}coeff}\PYG{+w}{  }\PYG{o}{*}\PYG{+w}{ }\PYG{o}{*}\PYG{+w}{ }\PYG{n}{rebo}\PYG{+w}{ }\PYG{n}{CH.rebo}\PYG{+w}{    }\PYG{n}{NULL}\PYG{+w}{  }\PYG{n}{C}\PYG{+w}{ }\PYG{n}{H}
\PYG{k}{pair\PYGZus{}coeff}\PYG{+w}{  }\PYG{o}{*}\PYG{+w}{ }\PYG{o}{*}\PYG{+w}{ }\PYG{n}{saip}\PYG{o}{/}\PYG{n}{metal}\PYG{+w}{  }\PYG{n}{CHAu.ILP}\PYG{+w}{  }\PYG{n}{Au}\PYG{+w}{ }\PYG{n}{C}\PYG{+w}{ }\PYG{n}{H}
\end{sphinxVerbatim}


\subsection{Description}
\label{\detokenize{pair_saip_metal:description}}
\sphinxAtStartPar
\DUrole{versionmodified,added}{Added in version 17Feb2022.}

\sphinxAtStartPar
The \sphinxstyleemphasis{saip/metal} style computes the registry\sphinxhyphen{}dependent interlayer
potential (ILP) potential for hetero\sphinxhyphen{}junctions formed with hexagonal
2D materials and metal surfaces, as described in {\hyperref[\detokenize{pair_saip_metal:ouyang6}]{\sphinxcrossref{\DUrole{std,std-ref}{(Ouyang6)}}}}.
\begin{equation*}
\begin{split}E  = & \frac{1}{2} \sum_i \sum_{j \neq i} V_{ij} \\
V_{ij}  = & {\rm Tap}(r_{ij})\left \{ e^{-\alpha (r_{ij}/\beta -1)}
             \left [ \epsilon + f(\rho_{ij}) + f(\rho_{ji})\right ] -
              \frac{1}{1+e^{-d\left [ \left ( r_{ij}/\left (s_R \cdot r^{eff} \right ) \right )-1 \right ]}}
              \cdot \frac{C_6}{r^6_{ij}} \right \}\\
\rho_{ij}^2 = & r_{ij}^2 - ({\bf r}_{ij} \cdot {\bf n}_i)^2 \\
\rho_{ji}^2  = & r_{ij}^2 - ({\bf r}_{ij} \cdot {\bf n}_j)^2 \\
f(\rho)  = &  C e^{ -( \rho / \delta )^2 } \\
{\rm Tap}(r_{ij})  = & 20\left ( \frac{r_{ij}}{R_{cut}} \right )^7 -
                        70\left ( \frac{r_{ij}}{R_{cut}} \right )^6 +
                        84\left ( \frac{r_{ij}}{R_{cut}} \right )^5 -
                        35\left ( \frac{r_{ij}}{R_{cut}} \right )^4 + 1\end{split}
\end{equation*}
\sphinxAtStartPar
Where \(\mathrm{Tap}(r_{ij})\) is the taper function which provides
a continuous cutoff (up to third derivative) for interatomic separations
larger than \(r_c\) {\hyperref[\detokenize{pair_ilp_graphene_hbn::doc}]{\sphinxcrossref{\DUrole{doc}{pair\_style ilp\_graphene\_hbn}}}}.

\sphinxAtStartPar
It is important to include all the pairs to build the neighbor list for
calculating the normals.

\begin{sphinxadmonition}{note}{Note:}
\sphinxAtStartPar
To account for the isotropic nature of the isolated gold atom
electron cloud, their corresponding normal vectors (\sphinxtitleref{\{bf n\}\_i}) are
assumed to lie along the interatomic vector \sphinxtitleref{\{bf r\}\_ij}. Notably, this
assumption is suitable for many bulk material surfaces, for
example, for systems possessing s\sphinxhyphen{}type valence orbitals or
metallic surfaces, whose valence electrons are mostly
delocalized, such that their Pauli repulsion with the electrons
of adjacent surfaces are isotropic. Caution should be used in
the case of very small gold contacts, for example, nano\sphinxhyphen{}clusters,
where edge effects may become relevant.
\end{sphinxadmonition}

\sphinxAtStartPar
The parameter file (e.g. CHAu.ILP), is intended for use with \sphinxstyleemphasis{metal}
{\hyperref[\detokenize{units::doc}]{\sphinxcrossref{\DUrole{doc}{units}}}}, with energies in meV. Two additional parameters,
\sphinxstyleemphasis{S}, and \sphinxstyleemphasis{rcut} are included in the parameter file. \sphinxstyleemphasis{S} is designed to
facilitate scaling of energies. \sphinxstyleemphasis{rcut} is designed to build the neighbor
list for calculating the normals for each atom pair.

\begin{sphinxadmonition}{note}{Note:}
\sphinxAtStartPar
The parameters presented in the parameter file (e.g. BNCH.ILP),
are fitted with taper function by setting the cutoff equal to 16.0
Angstrom.  Using different cutoff or taper function should be careful.
\end{sphinxadmonition}

\sphinxAtStartPar
This potential must be used in combination with hybrid/overlay.
Other interactions can be set to zero using pair\_style \sphinxstyleemphasis{none}.

\sphinxAtStartPar
This pair style tallies a breakdown of the total interlayer potential
energy into sub\sphinxhyphen{}categories, which can be accessed via the {\hyperref[\detokenize{compute_pair::doc}]{\sphinxcrossref{\DUrole{doc}{compute pair}}}} command as a vector of values of length 2.
The 2 values correspond to the following sub\sphinxhyphen{}categories:
\begin{enumerate}
\sphinxsetlistlabels{\arabic}{enumi}{enumii}{}{.}%
\item {} 
\sphinxAtStartPar
\sphinxstyleemphasis{E\_vdW} = vdW (attractive) energy

\item {} 
\sphinxAtStartPar
\sphinxstyleemphasis{E\_Rep} = Repulsive energy

\end{enumerate}

\sphinxAtStartPar
To print these quantities to the log file (with descriptive column
headings) the following commands could be included in an input script:

\begin{sphinxVerbatim}[commandchars=\\\{\}]
\PYG{k}{compute }\PYG{n+nv+nvIdentifier}{0}\PYG{+w}{ }\PYG{n+nv+nvIdentifier}{all}\PYG{+w}{ }\PYG{n}{pair}\PYG{+w}{ }\PYG{n}{saip}\PYG{o}{/}\PYG{n}{metal}
\PYG{k}{variable }\PYG{n+nv+nvIdentifier}{Evdw}\PYG{+w}{  }\PYG{n}{equal}\PYG{+w}{ }\PYG{n}{c\PYGZus{}0[1]}
\PYG{k}{variable }\PYG{n+nv+nvIdentifier}{Erep}\PYG{+w}{  }\PYG{n}{equal}\PYG{+w}{ }\PYG{n}{c\PYGZus{}0[2]}
\PYG{k}{thermo\PYGZus{}style}\PYG{+w}{ }\PYG{n}{custom}\PYG{+w}{ }\PYG{n}{step}\PYG{+w}{ }\PYG{n}{temp}\PYG{+w}{ }\PYG{n}{epair}\PYG{+w}{ }\PYG{n}{v\PYGZus{}Erep}\PYG{+w}{ }\PYG{n}{v\PYGZus{}Evdw}
\end{sphinxVerbatim}


\bigskip\hrule\bigskip


\sphinxAtStartPar
Styles with a \sphinxstyleemphasis{gpu}, \sphinxstyleemphasis{intel}, \sphinxstyleemphasis{kk}, \sphinxstyleemphasis{omp}, or \sphinxstyleemphasis{opt} suffix are
functionally the same as the corresponding style without the suffix.
They have been optimized to run faster, depending on your available
hardware, as discussed on the {\hyperref[\detokenize{Speed_packages::doc}]{\sphinxcrossref{\DUrole{doc}{Accelerator packages}}}}
page.  The accelerated styles take the same arguments and should
produce the same results, except for round\sphinxhyphen{}off and precision issues.

\sphinxAtStartPar
These accelerated styles are part of the GPU, INTEL, KOKKOS,
OPENMP, and OPT packages, respectively.  They are only enabled if
LAMMPS was built with those packages.  See the {\hyperref[\detokenize{Build_package::doc}]{\sphinxcrossref{\DUrole{doc}{Build package}}}} page for more info.

\sphinxAtStartPar
You can specify the accelerated styles explicitly in your input script
by including their suffix, or you can use the {\hyperref[\detokenize{Run_options::doc}]{\sphinxcrossref{\DUrole{doc}{\sphinxhyphen{}suffix command\sphinxhyphen{}line switch}}}} when you invoke LAMMPS, or you can use the
{\hyperref[\detokenize{suffix::doc}]{\sphinxcrossref{\DUrole{doc}{suffix}}}} command in your input script.

\sphinxAtStartPar
See the {\hyperref[\detokenize{Speed_packages::doc}]{\sphinxcrossref{\DUrole{doc}{Accelerator packages}}}} page for more
instructions on how to use the accelerated styles effectively.


\bigskip\hrule\bigskip



\subsection{Mixing, shift, table, tail correction, restart, rRESPA info}
\label{\detokenize{pair_saip_metal:mixing-shift-table-tail-correction-restart-rrespa-info}}
\sphinxAtStartPar
This pair style does not support the pair\_modify mix, shift, table, and
tail options.

\sphinxAtStartPar
This pair style does not write their information to binary restart
files, since it is stored in potential files. Thus, you need to
re\sphinxhyphen{}specify the pair\_style and pair\_coeff commands in an input script
that reads a restart file.


\subsection{Restrictions}
\label{\detokenize{pair_saip_metal:restrictions}}
\sphinxAtStartPar
This pair style is part of the INTERLAYER package.  It is only enabled
if LAMMPS was built with that package.  See the {\hyperref[\detokenize{Build_package::doc}]{\sphinxcrossref{\DUrole{doc}{Build package}}}} page for more info.

\sphinxAtStartPar
This pair style requires the newton setting to be \sphinxstyleemphasis{on} for pair
interactions.

\sphinxAtStartPar
The CHAu.ILP potential file provided with LAMMPS (see the potentials
directory) are parameterized for \sphinxstyleemphasis{metal} units.  You can use this
potential with any LAMMPS units, but you would need to create your own
custom CHAu.ILP potential file with coefficients listed in the appropriate
units, if your simulation does not use \sphinxstyleemphasis{metal} units.


\subsection{Related commands}
\label{\detokenize{pair_saip_metal:related-commands}}
\sphinxAtStartPar
{\hyperref[\detokenize{pair_coeff::doc}]{\sphinxcrossref{\DUrole{doc}{pair\_coeff}}}},
{\hyperref[\detokenize{pair_none::doc}]{\sphinxcrossref{\DUrole{doc}{pair\_none}}}},
{\hyperref[\detokenize{pair_hybrid::doc}]{\sphinxcrossref{\DUrole{doc}{pair\_style hybrid/overlay}}}},
{\hyperref[\detokenize{pair_drip::doc}]{\sphinxcrossref{\DUrole{doc}{pair\_style drip}}}},
{\hyperref[\detokenize{pair_ilp_tmd::doc}]{\sphinxcrossref{\DUrole{doc}{pair\_style ilp\_tmd}}}},
{\hyperref[\detokenize{pair_ilp_graphene_hbn::doc}]{\sphinxcrossref{\DUrole{doc}{pair\_style ilp\_graphene\_hbn}}}},
{\hyperref[\detokenize{pair_kolmogorov_crespi_z::doc}]{\sphinxcrossref{\DUrole{doc}{pair\_style pair\_kolmogorov\_crespi\_z}}}},
{\hyperref[\detokenize{pair_kolmogorov_crespi_full::doc}]{\sphinxcrossref{\DUrole{doc}{pair\_style pair\_kolmogorov\_crespi\_full}}}},
{\hyperref[\detokenize{pair_lebedeva_z::doc}]{\sphinxcrossref{\DUrole{doc}{pair\_style pair\_lebedeva\_z}}}},
{\hyperref[\detokenize{pair_coul_shield::doc}]{\sphinxcrossref{\DUrole{doc}{pair\_style pair\_coul\_shield}}}}.


\subsection{Default}
\label{\detokenize{pair_saip_metal:default}}
\sphinxAtStartPar
tap\_flag = 1


\bigskip\hrule\bigskip

\phantomsection\label{\detokenize{pair_saip_metal:ouyang6}}
\sphinxAtStartPar
\sphinxstylestrong{(Ouyang6)} W. Ouyang, O. Hod, and R. Guerra, J. Chem. Theory Comput. 17, 7215 (2021).

\sphinxstepscope

\index{pair\_style sdpd/taitwater/isothermal@\spxentry{pair\_style sdpd/taitwater/isothermal}}\ignorespaces 

\section{pair\_style sdpd/taitwater/isothermal command}
\label{\detokenize{pair_sdpd_taitwater_isothermal:pair-style-sdpd-taitwater-isothermal-command}}\label{\detokenize{pair_sdpd_taitwater_isothermal:index-0}}\label{\detokenize{pair_sdpd_taitwater_isothermal::doc}}

\subsection{Syntax}
\label{\detokenize{pair_sdpd_taitwater_isothermal:syntax}}
\begin{sphinxVerbatim}[commandchars=\\\{\}]
\PYG{k}{pair\PYGZus{}style}\PYG{+w}{ }\PYG{n}{sdpd}\PYG{o}{/}\PYG{n}{taitwater}\PYG{o}{/}\PYG{n}{isothermal}\PYG{+w}{ }\PYG{n}{temperature}\PYG{+w}{ }\PYG{n}{viscosity}\PYG{+w}{ }\PYG{n}{seed}
\end{sphinxVerbatim}
\begin{itemize}
\item {} 
\sphinxAtStartPar
temperature = temperature of the fluid (temperature units)

\item {} 
\sphinxAtStartPar
viscosity = dynamic viscosity of the fluid (mass*distance/time units)

\item {} 
\sphinxAtStartPar
seed = random number generator seed (positive integer, optional)

\end{itemize}


\subsection{Examples}
\label{\detokenize{pair_sdpd_taitwater_isothermal:examples}}
\begin{sphinxVerbatim}[commandchars=\\\{\}]
\PYG{k}{pair\PYGZus{}style}\PYG{+w}{ }\PYG{n}{sdpd}\PYG{o}{/}\PYG{n}{taitwater}\PYG{o}{/}\PYG{n}{isothermal}\PYG{+w}{ }\PYG{l+m}{300}\PYG{n}{.}\PYG{+w}{ }\PYG{l+m}{1}\PYG{n}{.}\PYG{+w}{ }\PYG{l+m}{28681}
\PYG{k}{pair\PYGZus{}coeff}\PYG{+w}{ }\PYG{o}{*}\PYG{+w}{ }\PYG{o}{*}\PYG{+w}{ }\PYG{l+m}{1000.0}\PYG{+w}{ }\PYG{l+m}{1430.0}\PYG{+w}{ }\PYG{l+m}{2.4}
\end{sphinxVerbatim}


\subsection{Description}
\label{\detokenize{pair_sdpd_taitwater_isothermal:description}}
\sphinxAtStartPar
The sdpd/taitwater/isothermal style computes forces between mesoscopic
particles according to the Smoothed Dissipative Particle Dynamics model
described in this paper by {\hyperref[\detokenize{pair_sdpd_taitwater_isothermal:espanol-revenga}]{\sphinxcrossref{\DUrole{std,std-ref}{(Espanol and Revenga)}}}} under
the following assumptions:
\begin{enumerate}
\sphinxsetlistlabels{\arabic}{enumi}{enumii}{}{.}%
\item {} 
\sphinxAtStartPar
The temperature is constant and uniform.

\item {} 
\sphinxAtStartPar
The shear viscosity is constant and uniform.

\item {} 
\sphinxAtStartPar
The volume viscosity is negligible before the shear viscosity.

\item {} 
\sphinxAtStartPar
The Boltzmann constant is negligible before the heat capacity of a
single mesoscopic particle of fluid.

\end{enumerate}

\sphinxAtStartPar
The third assumption is true for water in nearly incompressible flows.
The fourth holds true for water for any reasonable size one can
imagine for a mesoscopic particle.

\sphinxAtStartPar
The pressure forces between particles will be computed according to
Tait’s equation of state:
\begin{equation*}
\begin{split}p = B \left[(\frac{\rho}{\rho_0})^{\gamma} - 1\right]\end{split}
\end{equation*}
\sphinxAtStartPar
where \(\gamma = 7\) and \(B = c_0^2 \rho_0 / \gamma\), with
\(\rho_0\) being the reference density and \(c_0\) the reference
speed of sound.

\sphinxAtStartPar
The laminar viscosity and the random forces will be computed according
to formulas described in {\hyperref[\detokenize{pair_sdpd_taitwater_isothermal:espanol-revenga}]{\sphinxcrossref{\DUrole{std,std-ref}{(Espanol and Revenga)}}}}.

\begin{sphinxadmonition}{warning}{Warning:}
\sphinxAtStartPar
Similar to {\hyperref[\detokenize{pair_brownian::doc}]{\sphinxcrossref{\DUrole{doc}{brownian}}}} and
{\hyperref[\detokenize{pair_dpd::doc}]{\sphinxcrossref{\DUrole{doc}{dpd}}}} styles, the {\hyperref[\detokenize{newton::doc}]{\sphinxcrossref{\DUrole{doc}{newton}}}} setting for
pairwise interactions needs to be on when running LAMMPS in parallel
if you want to ensure linear momentum conservation. Otherwise random
forces generated for pairs straddling processor boundary will not be
equal and opposite.
\end{sphinxadmonition}

\begin{sphinxadmonition}{note}{Note:}
\sphinxAtStartPar
The actual random seed used will be a mix of what you specify
and other parameters like the MPI ranks. This is to ensure that
different MPI tasks have distinct seeds.
\end{sphinxadmonition}

\sphinxAtStartPar
The following coefficients must be defined for each pair of atoms
types via the {\hyperref[\detokenize{pair_coeff::doc}]{\sphinxcrossref{\DUrole{doc}{pair\_coeff}}}} command as in the examples
above.
\begin{itemize}
\item {} 
\sphinxAtStartPar
\(\rho_0\) reference density (mass/volume units)

\item {} 
\sphinxAtStartPar
\(c_0\) reference soundspeed (distance/time units)

\item {} 
\sphinxAtStartPar
h kernel function cutoff (distance units)

\end{itemize}


\bigskip\hrule\bigskip



\subsection{Mixing, shift, table, tail correction, restart, rRESPA info}
\label{\detokenize{pair_sdpd_taitwater_isothermal:mixing-shift-table-tail-correction-restart-rrespa-info}}
\sphinxAtStartPar
This style does not support mixing.  Thus, coefficients for all
I,J pairs must be specified explicitly.

\sphinxAtStartPar
This style does not support the {\hyperref[\detokenize{pair_modify::doc}]{\sphinxcrossref{\DUrole{doc}{pair\_modify}}}}
shift, table, and tail options.

\sphinxAtStartPar
This style does not write information to {\hyperref[\detokenize{restart::doc}]{\sphinxcrossref{\DUrole{doc}{binary restart files}}}}.  Thus, you need to re\sphinxhyphen{}specify the pair\_style and
pair\_coeff commands in an input script that reads a restart file.

\sphinxAtStartPar
This style can only be used via the \sphinxstyleemphasis{pair} keyword of the {\hyperref[\detokenize{run_style::doc}]{\sphinxcrossref{\DUrole{doc}{run\_style respa}}}} command.  It does not support the \sphinxstyleemphasis{inner},
\sphinxstyleemphasis{middle}, \sphinxstyleemphasis{outer} keywords.


\subsection{Restrictions}
\label{\detokenize{pair_sdpd_taitwater_isothermal:restrictions}}
\sphinxAtStartPar
This pair style is part of the DPD\sphinxhyphen{}SMOOTH package.  It is only enabled
if LAMMPS was built with that package.  See the {\hyperref[\detokenize{Build_package::doc}]{\sphinxcrossref{\DUrole{doc}{Build package}}}} page for more info.


\subsection{Related commands}
\label{\detokenize{pair_sdpd_taitwater_isothermal:related-commands}}
\sphinxAtStartPar
{\hyperref[\detokenize{pair_coeff::doc}]{\sphinxcrossref{\DUrole{doc}{pair coeff}}}}, {\hyperref[\detokenize{pair_sph_rhosum::doc}]{\sphinxcrossref{\DUrole{doc}{pair sph/rhosum}}}},
{\hyperref[\detokenize{pair_sph_taitwater::doc}]{\sphinxcrossref{\DUrole{doc}{pair sph/taitwater}}}}


\subsection{Default}
\label{\detokenize{pair_sdpd_taitwater_isothermal:default}}
\sphinxAtStartPar
The default seed is 0 (before mixing).


\bigskip\hrule\bigskip

\phantomsection\label{\detokenize{pair_sdpd_taitwater_isothermal:espanol-revenga}}
\sphinxAtStartPar
\sphinxstylestrong{(Espanol and Revenga)} Espanol, Revenga, Physical Review E, 67, 026705 (2003).

\sphinxstepscope

\index{pair\_style smatb@\spxentry{pair\_style smatb}}\ignorespaces 
\index{pair\_style smatb/single@\spxentry{pair\_style smatb/single}}\ignorespaces 

\section{pair\_style smatb command}
\label{\detokenize{pair_smatb:pair-style-smatb-command}}\label{\detokenize{pair_smatb:index-1}}\label{\detokenize{pair_smatb:index-0}}\label{\detokenize{pair_smatb::doc}}

\section{pair\_style smatb/single command}
\label{\detokenize{pair_smatb:pair-style-smatb-single-command}}

\subsection{Syntax}
\label{\detokenize{pair_smatb:syntax}}
\begin{sphinxVerbatim}[commandchars=\\\{\}]
\PYG{k}{pair\PYGZus{}style}\PYG{+w}{ }\PYG{n}{style}\PYG{+w}{ }\PYG{n}{args}
\end{sphinxVerbatim}
\begin{itemize}
\item {} 
\sphinxAtStartPar
style = \sphinxstyleemphasis{smatb} or \sphinxstyleemphasis{smatb/single}

\item {} 
\sphinxAtStartPar
args = none

\end{itemize}


\subsection{Examples}
\label{\detokenize{pair_smatb:examples}}
\begin{sphinxVerbatim}[commandchars=\\\{\}]
\PYG{k}{pair\PYGZus{}style}\PYG{+w}{ }\PYG{n}{smatb}
\PYG{k}{pair\PYGZus{}coeff}\PYG{+w}{ }\PYG{l+m}{1}\PYG{+w}{ }\PYG{l+m}{1}\PYG{+w}{ }\PYG{l+m}{2.88}\PYG{+w}{ }\PYG{l+m}{10.35}\PYG{+w}{ }\PYG{l+m}{4.178}\PYG{+w}{ }\PYG{l+m}{0.210}\PYG{+w}{ }\PYG{l+m}{1.818}\PYG{+w}{ }\PYG{l+m}{4.07293506}\PYG{+w}{ }\PYG{l+m}{4.9883063257983666}

\PYG{k}{pair\PYGZus{}style}\PYG{+w}{ }\PYG{n}{smatb}\PYG{o}{/}\PYG{n}{single}
\PYG{k}{pair\PYGZus{}coeff}\PYG{+w}{ }\PYG{l+m}{1}\PYG{+w}{ }\PYG{l+m}{1}\PYG{+w}{ }\PYG{l+m}{2.88}\PYG{+w}{ }\PYG{l+m}{10.35}\PYG{+w}{ }\PYG{l+m}{4.178}\PYG{+w}{ }\PYG{l+m}{0.210}\PYG{+w}{ }\PYG{l+m}{1.818}\PYG{+w}{ }\PYG{l+m}{4.07293506}\PYG{+w}{ }\PYG{l+m}{4.9883063257983666}
\end{sphinxVerbatim}


\subsection{Description}
\label{\detokenize{pair_smatb:description}}
\sphinxAtStartPar
\DUrole{versionmodified,added}{Added in version 4May2022.}

\sphinxAtStartPar
The \sphinxstyleemphasis{smatb} and \sphinxstyleemphasis{smatb/single} styles compute the Second Moment
Approximation to the Tight Binding {\hyperref[\detokenize{pair_smatb:cyrot}]{\sphinxcrossref{\DUrole{std,std-ref}{(Cyrot)}}}},
{\hyperref[\detokenize{pair_smatb:gupta}]{\sphinxcrossref{\DUrole{std,std-ref}{(Gupta)}}}}, {\hyperref[\detokenize{pair_smatb:rosato}]{\sphinxcrossref{\DUrole{std,std-ref}{(Rosato)}}}}, given by
\begin{equation*}
\begin{split}E_{i}  = \sum_{j,R_{ij}\leq R_{c}} \alpha(R_{ij}) - \sqrt{\sum_{j,R_{ij}\leq R_{c}}\Xi^2(R_{ij})}\end{split}
\end{equation*}
\sphinxAtStartPar
\(R_{ij}\) is the distance between the atom \(i\) and \(j\).
And the two functions \(\alpha\left(r\right)\) and \(\Xi\left(r\right)\) are:
\begin{equation*}
\begin{split}\alpha\left(r\right)=\left\lbrace\begin{array}{ll}
   A e^{-p \left(\frac{r}{R_{0}}-1\right)} & r < R_{sc}\\
   a_3\left(r-R_{c}\right)^3+a_4\left(r-R_{c}\right)^4
   +a_5\left(r-R_{c}\right)^5& R_{sc} < r < R_{c}
   \end{array}
   \right.\end{split}
\end{equation*}\begin{equation*}
\begin{split}\Xi\left(r\right)=\left\lbrace\begin{array}{ll}
\xi e^{-q \left(\frac{r}{R_{0}}-1\right)} & r < R_{sc}\\
x_3\left(r-R_{c}\right)^3+x_4\left(r-R_{c}\right)^4
+x_5\left(r-R_{c}\right)^5& R_{sc} < r < R_{c}
\end{array}
\right.\end{split}
\end{equation*}
\sphinxAtStartPar
The polynomial coefficients \(a_3\), \(a_4\), \(a_5\),
\(x_3\), \(x_4\), \(x_5\) are computed by LAMMPS: the two
exponential terms and their first and second derivatives are smoothly
reduced to zero, from the inner cutoff \(R_{sc}\) to the outer
cutoff \(R_{c}\).

\sphinxAtStartPar
The \sphinxstyleemphasis{smatb/single} style is an optimization when using only a single atom type.


\subsection{Coefficients}
\label{\detokenize{pair_smatb:coefficients}}
\sphinxAtStartPar
The following coefficients must be defined for each pair of atoms types via the
{\hyperref[\detokenize{pair_coeff::doc}]{\sphinxcrossref{\DUrole{doc}{pair\_coeff}}}} command as in the examples above, or in the data
file or restart files read by the {\hyperref[\detokenize{read_data::doc}]{\sphinxcrossref{\DUrole{doc}{read\_data}}}} or
{\hyperref[\detokenize{read_restart::doc}]{\sphinxcrossref{\DUrole{doc}{read\_restart}}}} commands, or by mixing as described below:
\begin{itemize}
\item {} 
\sphinxAtStartPar
\(R_{0}\) (distance units)

\item {} 
\sphinxAtStartPar
\(p\) (dimensionless)

\item {} 
\sphinxAtStartPar
\(q\) (dimensionless)

\item {} 
\sphinxAtStartPar
\(A\) (energy units)

\item {} 
\sphinxAtStartPar
\(\xi\) (energy units)

\item {} 
\sphinxAtStartPar
\(R_{cs}\) (distance units)

\item {} 
\sphinxAtStartPar
\(R_{c}\) (distance units)

\end{itemize}

\sphinxAtStartPar
Note that: \(R_{0}\) is the nearest neighbor distance, usually coincides
with the diameter of the atoms

\sphinxAtStartPar
See the {\hyperref[\detokenize{run_style::doc}]{\sphinxcrossref{\DUrole{doc}{run\_style}}}} command for details.


\bigskip\hrule\bigskip



\subsection{Mixing info}
\label{\detokenize{pair_smatb:mixing-info}}
\sphinxAtStartPar
For atom type pairs I,J and I != J the coefficients are not automatically mixed.


\bigskip\hrule\bigskip



\subsection{Restrictions}
\label{\detokenize{pair_smatb:restrictions}}
\sphinxAtStartPar
These pair styles are part of the SMTBQ package and are only enabled
if LAMMPS is built with that package.  See the {\hyperref[\detokenize{Build_package::doc}]{\sphinxcrossref{\DUrole{doc}{Build package}}}} page for more info.

\sphinxAtStartPar
These pair styles require the {\hyperref[\detokenize{newton::doc}]{\sphinxcrossref{\DUrole{doc}{newton}}}} setting to be “on” for pair interactions.


\subsection{Related commands}
\label{\detokenize{pair_smatb:related-commands}}\begin{itemize}
\item {} 
\sphinxAtStartPar
{\hyperref[\detokenize{pair_coeff::doc}]{\sphinxcrossref{\DUrole{doc}{pair\_coeff}}}}

\end{itemize}


\subsection{Default}
\label{\detokenize{pair_smatb:default}}
\sphinxAtStartPar
none


\bigskip\hrule\bigskip

\phantomsection\label{\detokenize{pair_smatb:cyrot}}
\sphinxAtStartPar
\sphinxstylestrong{(Cyrot)}  Cyrot\sphinxhyphen{}Lackmann and Ducastelle, Phys Rev. B, 4, 2406\sphinxhyphen{}2412 (1971).

\phantomsection\label{\detokenize{pair_smatb:gupta}}
\sphinxAtStartPar
\sphinxstylestrong{(Gupta)} Gupta ,Phys Rev. B, 23, 6265\sphinxhyphen{}6270 (1981).

\phantomsection\label{\detokenize{pair_smatb:rosato}}
\sphinxAtStartPar
\sphinxstylestrong{(Rosato)} Rosato and Guillope  and Legrand, Philosophical Magazine A, 59.2, 321\sphinxhyphen{}336 (1989).

\sphinxstepscope

\index{pair\_style smd/hertz@\spxentry{pair\_style smd/hertz}}\ignorespaces 

\section{pair\_style smd/hertz command}
\label{\detokenize{pair_smd_hertz:pair-style-smd-hertz-command}}\label{\detokenize{pair_smd_hertz:index-0}}\label{\detokenize{pair_smd_hertz::doc}}

\subsection{Syntax}
\label{\detokenize{pair_smd_hertz:syntax}}
\begin{sphinxVerbatim}[commandchars=\\\{\}]
\PYG{k}{pair\PYGZus{}style}\PYG{+w}{ }\PYG{n}{smd}\PYG{o}{/}\PYG{n}{hertz}\PYG{+w}{ }\PYG{n}{scale\PYGZus{}factor}
\end{sphinxVerbatim}


\subsection{Examples}
\label{\detokenize{pair_smd_hertz:examples}}
\begin{sphinxVerbatim}[commandchars=\\\{\}]
\PYG{k}{pair\PYGZus{}style}\PYG{+w}{ }\PYG{n}{smd}\PYG{o}{/}\PYG{n}{hertz}\PYG{+w}{ }\PYG{l+m}{1.0}
\PYG{k}{pair\PYGZus{}coeff}\PYG{+w}{ }\PYG{l+m}{1}\PYG{+w}{ }\PYG{l+m}{1}\PYG{+w}{ }\PYG{o}{\PYGZlt{}}\PYG{n}{contact\PYGZus{}stiffness}\PYG{o}{\PYGZgt{}}
\end{sphinxVerbatim}


\subsection{Description}
\label{\detokenize{pair_smd_hertz:description}}
\sphinxAtStartPar
The \sphinxstyleemphasis{smd/hertz} style calculates contact forces between SPH particles
belonging to different physical bodies.

\sphinxAtStartPar
The contact forces are calculated using a Hertz potential, which
evaluates the overlap between two particles (whose spatial extents are
defined via its contact radius).  The effect is that a particles
cannot penetrate into each other.  The parameter \textless{}contact\_stiffness\textgreater{}
has units of pressure and should equal roughly one half of the Young’s
modulus (or bulk modulus in the case of fluids) of the material model
associated with the SPH particles.

\sphinxAtStartPar
The parameter \sphinxstyleemphasis{scale\_factor} can be used to scale the particles’
contact radii. This can be useful to control how close particles can
approach each other. Usually, \sphinxstyleemphasis{scale\_factor} =1.0.


\bigskip\hrule\bigskip



\subsection{Mixing, shift, table, tail correction, restart, rRESPA info}
\label{\detokenize{pair_smd_hertz:mixing-shift-table-tail-correction-restart-rrespa-info}}
\sphinxAtStartPar
No mixing is performed automatically.  Currently, no part of MACHDYN
supports restarting nor minimization.  rRESPA does not apply to this
pair style.


\bigskip\hrule\bigskip



\subsection{Restrictions}
\label{\detokenize{pair_smd_hertz:restrictions}}
\sphinxAtStartPar
This fix is part of the MACHDYN package.  It is only enabled if
LAMMPS was built with that package.  See the {\hyperref[\detokenize{Build_package::doc}]{\sphinxcrossref{\DUrole{doc}{Build package}}}} page for more info.


\subsection{Related commands}
\label{\detokenize{pair_smd_hertz:related-commands}}
\sphinxAtStartPar
{\hyperref[\detokenize{pair_coeff::doc}]{\sphinxcrossref{\DUrole{doc}{pair\_coeff}}}}


\subsection{Default}
\label{\detokenize{pair_smd_hertz:default}}
\sphinxAtStartPar
none

\sphinxstepscope

\index{pair\_style smd/tlsph@\spxentry{pair\_style smd/tlsph}}\ignorespaces 

\section{pair\_style smd/tlsph command}
\label{\detokenize{pair_smd_tlsph:pair-style-smd-tlsph-command}}\label{\detokenize{pair_smd_tlsph:index-0}}\label{\detokenize{pair_smd_tlsph::doc}}

\subsection{Syntax}
\label{\detokenize{pair_smd_tlsph:syntax}}
\begin{sphinxVerbatim}[commandchars=\\\{\}]
\PYG{k}{pair\PYGZus{}style}\PYG{+w}{ }\PYG{n}{smd}\PYG{o}{/}\PYG{n}{tlsph}\PYG{+w}{ }\PYG{n}{args}
\end{sphinxVerbatim}


\subsection{Examples}
\label{\detokenize{pair_smd_tlsph:examples}}
\begin{sphinxVerbatim}[commandchars=\\\{\}]
\PYG{k}{pair\PYGZus{}style}\PYG{+w}{ }\PYG{n}{smd}\PYG{o}{/}\PYG{n}{tlsph}
\end{sphinxVerbatim}


\subsection{Description}
\label{\detokenize{pair_smd_tlsph:description}}
\sphinxAtStartPar
The \sphinxstyleemphasis{smd/tlsph} style computes particle interactions according to
continuum mechanics constitutive laws and a Total\sphinxhyphen{}Lagrangian
Smooth\sphinxhyphen{}Particle Hydrodynamics algorithm.

\sphinxAtStartPar
This pair style is invoked with the following command:

\begin{sphinxVerbatim}[commandchars=\\\{\}]
\PYG{k}{pair\PYGZus{}style}\PYG{+w}{ }\PYG{n}{smd}\PYG{o}{/}\PYG{n}{tlsph}
\PYG{k}{pair\PYGZus{}coeff}\PYG{+w}{ }\PYG{n}{i}\PYG{+w}{ }\PYG{n}{j}\PYG{+w}{ }\PYG{o}{*}\PYG{n}{COMMON}\PYG{+w}{ }\PYG{n}{rho0}\PYG{+w}{ }\PYG{n}{E}\PYG{+w}{ }\PYG{n}{nu}\PYG{+w}{ }\PYG{n}{Q1}\PYG{+w}{ }\PYG{n}{Q2}\PYG{+w}{ }\PYG{n}{hg}\PYG{+w}{ }\PYG{n}{Cp}\PYG{+w}{ }\PYG{o}{\PYGZam{}}
\PYG{+w}{               }\PYG{o}{*}\PYG{n}{END}
\end{sphinxVerbatim}

\sphinxAtStartPar
Here, \sphinxstyleemphasis{i} and \sphinxstyleemphasis{j} denote the \sphinxstyleemphasis{LAMMPS} particle types for which this
pair style is defined. Note that \sphinxstyleemphasis{i} and \sphinxstyleemphasis{j} must be equal, i.e., no
\sphinxstyleemphasis{tlsph} cross interactions between different particle types are
allowed.  In contrast to the usual \sphinxstyleemphasis{LAMMPS} \sphinxstyleemphasis{pair coeff} definitions,
which are given solely a number of floats and integers, the \sphinxstyleemphasis{tlsph}
\sphinxstyleemphasis{pair coeff} definition is organized using keywords. These keywords
mark the beginning of different sets of parameters for particle
properties, material constitutive models, and damage models. The \sphinxstyleemphasis{pair
coeff} line must be terminated with the \sphinxstyleemphasis{*END} keyword. The use the
line continuation operator \sphinxstyleemphasis{\&} is recommended. A typical invocation of
the \sphinxstyleemphasis{tlsph} for a solid body would consist of an equation of state for
computing the pressure (the diagonal components of the stress tensor),
and a material model to compute shear stresses (the off\sphinxhyphen{}diagonal
components of the stress tensor). Damage and failure models can also
be added.

\sphinxAtStartPar
Please see the \sphinxhref{PDF/MACHDYN\_LAMMPS\_userguide.pdf}{SMD user guide} for a
complete listing of the possible keywords and material models.


\bigskip\hrule\bigskip



\subsection{Mixing, shift, table, tail correction, restart, rRESPA info}
\label{\detokenize{pair_smd_tlsph:mixing-shift-table-tail-correction-restart-rrespa-info}}
\sphinxAtStartPar
No mixing is performed automatically.  Currently, no part of MACHDYN
supports restarting nor minimization.  rRESPA does not apply to this
pair style.


\bigskip\hrule\bigskip



\subsection{Restrictions}
\label{\detokenize{pair_smd_tlsph:restrictions}}
\sphinxAtStartPar
This fix is part of the MACHDYN package.  It is only enabled if
LAMMPS was built with that package.  See the {\hyperref[\detokenize{Build_package::doc}]{\sphinxcrossref{\DUrole{doc}{Build package}}}} page for more info.


\subsection{Related commands}
\label{\detokenize{pair_smd_tlsph:related-commands}}
\sphinxAtStartPar
{\hyperref[\detokenize{pair_coeff::doc}]{\sphinxcrossref{\DUrole{doc}{pair\_coeff}}}}


\subsection{Default}
\label{\detokenize{pair_smd_tlsph:default}}
\sphinxAtStartPar
none

\sphinxstepscope

\index{pair\_style smd/tri\_surface@\spxentry{pair\_style smd/tri\_surface}}\ignorespaces 

\section{pair\_style smd/tri\_surface command}
\label{\detokenize{pair_smd_triangulated_surface:pair-style-smd-tri-surface-command}}\label{\detokenize{pair_smd_triangulated_surface:index-0}}\label{\detokenize{pair_smd_triangulated_surface::doc}}

\subsection{Syntax}
\label{\detokenize{pair_smd_triangulated_surface:syntax}}
\begin{sphinxVerbatim}[commandchars=\\\{\}]
\PYG{k}{pair\PYGZus{}style}\PYG{+w}{ }\PYG{n}{smd}\PYG{o}{/}\PYG{n}{tri\PYGZus{}surface}\PYG{+w}{ }\PYG{n}{scale\PYGZus{}factor}
\end{sphinxVerbatim}


\subsection{Examples}
\label{\detokenize{pair_smd_triangulated_surface:examples}}
\begin{sphinxVerbatim}[commandchars=\\\{\}]
\PYG{k}{pair\PYGZus{}style}\PYG{+w}{ }\PYG{n}{smd}\PYG{o}{/}\PYG{n}{tri\PYGZus{}surface}\PYG{+w}{ }\PYG{l+m}{1.0}
\PYG{k}{pair\PYGZus{}coeff}\PYG{+w}{ }\PYG{l+m}{1}\PYG{+w}{ }\PYG{l+m}{1}\PYG{+w}{ }\PYG{o}{\PYGZlt{}}\PYG{n}{contact\PYGZus{}stiffness}\PYG{o}{\PYGZgt{}}
\end{sphinxVerbatim}


\subsection{Description}
\label{\detokenize{pair_smd_triangulated_surface:description}}
\sphinxAtStartPar
The \sphinxstyleemphasis{smd/tri\_surface} style calculates contact forces between SPH
particles and a rigid wall boundary defined via the
{\hyperref[\detokenize{fix_smd_wall_surface::doc}]{\sphinxcrossref{\DUrole{doc}{smd/wall\_surface}}}} fix.

\sphinxAtStartPar
The contact forces are calculated using a Hertz potential, which
evaluates the overlap between a particle (whose spatial extents are
defined via its contact radius) and the triangle.  The effect is that
a particle cannot penetrate into the triangular surface.  The
parameter \textless{}contact\_stiffness\textgreater{} has units of pressure and should equal
roughly one half of the Young’s modulus (or bulk modulus in the case
of fluids) of the material model associated with the SPH particle

\sphinxAtStartPar
The parameter \sphinxstyleemphasis{scale\_factor} can be used to scale the particles’
contact radii. This can be useful to control how close particles can
approach the triangulated surface. Usually, \sphinxstyleemphasis{scale\_factor} =1.0.


\bigskip\hrule\bigskip



\subsection{Mixing, shift, table, tail correction, restart, rRESPA info}
\label{\detokenize{pair_smd_triangulated_surface:mixing-shift-table-tail-correction-restart-rrespa-info}}
\sphinxAtStartPar
No mixing is performed automatically.
Currently, no part of MACHDYN supports restarting nor minimization.
rRESPA does not apply to this pair style.


\bigskip\hrule\bigskip



\subsection{Restrictions}
\label{\detokenize{pair_smd_triangulated_surface:restrictions}}
\sphinxAtStartPar
This fix is part of the MACHDYN package.  It is only enabled if
LAMMPS was built with that package.  See the {\hyperref[\detokenize{Build_package::doc}]{\sphinxcrossref{\DUrole{doc}{Build package}}}} page for more info.


\subsection{Related commands}
\label{\detokenize{pair_smd_triangulated_surface:related-commands}}
\sphinxAtStartPar
{\hyperref[\detokenize{pair_coeff::doc}]{\sphinxcrossref{\DUrole{doc}{pair\_coeff}}}}


\subsection{Default}
\label{\detokenize{pair_smd_triangulated_surface:default}}
\sphinxAtStartPar
none

\sphinxstepscope

\index{pair\_style smd/ulsph@\spxentry{pair\_style smd/ulsph}}\ignorespaces 

\section{pair\_style smd/ulsph command}
\label{\detokenize{pair_smd_ulsph:pair-style-smd-ulsph-command}}\label{\detokenize{pair_smd_ulsph:index-0}}\label{\detokenize{pair_smd_ulsph::doc}}

\subsection{Syntax}
\label{\detokenize{pair_smd_ulsph:syntax}}
\begin{sphinxVerbatim}[commandchars=\\\{\}]
\PYG{k}{pair\PYGZus{}style}\PYG{+w}{ }\PYG{n}{smd}\PYG{o}{/}\PYG{n}{ulsph}\PYG{+w}{ }\PYG{n}{args}
\end{sphinxVerbatim}
\begin{itemize}
\item {} 
\sphinxAtStartPar
these keywords must be given

\end{itemize}
\begin{sphinxalltt}
keyword = \sphinxstyleemphasis{*DENSITY\_SUMMATION} or \sphinxstyleemphasis{*DENSITY\_CONTINUITY} and \sphinxstyleemphasis{*VELOCITY\_GRADIENT} or \sphinxstyleemphasis{*NO\_VELOCITY\_GRADIENT} and \sphinxstyleemphasis{*GRADIENT\_CORRECTION} or \sphinxstyleemphasis{*NO\_GRADIENT\_CORRECTION}
\end{sphinxalltt}


\subsection{Examples}
\label{\detokenize{pair_smd_ulsph:examples}}
\begin{sphinxVerbatim}[commandchars=\\\{\}]
\PYG{k}{pair\PYGZus{}style}\PYG{+w}{ }\PYG{n}{smd}\PYG{o}{/}\PYG{n}{ulsph}\PYG{+w}{ }\PYG{o}{*}\PYG{n}{DENSITY\PYGZus{}CONTINUITY}\PYG{+w}{ }\PYG{o}{*}\PYG{n}{VELOCITY\PYGZus{}GRADIENT}\PYG{+w}{ }\PYG{o}{*}\PYG{n}{NO\PYGZus{}GRADIENT\PYGZus{}CORRECTION}
\end{sphinxVerbatim}


\subsection{Description}
\label{\detokenize{pair_smd_ulsph:description}}
\sphinxAtStartPar
The \sphinxstyleemphasis{smd/ulsph} style computes particle interactions according to
continuum mechanics constitutive laws and an updated Lagrangian
Smooth\sphinxhyphen{}Particle Hydrodynamics algorithm.

\sphinxAtStartPar
This pair style is invoked similar to the following command:

\begin{sphinxVerbatim}[commandchars=\\\{\}]
\PYG{k}{pair\PYGZus{}style}\PYG{+w}{ }\PYG{n}{smd}\PYG{o}{/}\PYG{n}{ulsph}\PYG{+w}{ }\PYG{o}{*}\PYG{n}{DENSITY\PYGZus{}CONTINUITY}\PYG{+w}{ }\PYG{o}{*}\PYG{n}{VELOCITY\PYGZus{}GRADIENT}\PYG{+w}{ }\PYG{o}{*}\PYG{n}{NO\PYGZus{}GRADIENT\PYGZus{}CORRECTION}
\PYG{k}{pair\PYGZus{}coeff}\PYG{+w}{ }\PYG{n}{i}\PYG{+w}{ }\PYG{n}{j}\PYG{+w}{ }\PYG{o}{*}\PYG{n}{COMMON}\PYG{+w}{ }\PYG{n}{rho0}\PYG{+w}{ }\PYG{n}{c0}\PYG{+w}{ }\PYG{n}{Q1}\PYG{+w}{ }\PYG{n}{Cp}\PYG{+w}{ }\PYG{n}{hg}\PYG{+w}{ }\PYG{o}{\PYGZam{}}
\PYG{+w}{               }\PYG{o}{*}\PYG{n}{END}
\end{sphinxVerbatim}

\sphinxAtStartPar
Here, \sphinxstyleemphasis{i} and \sphinxstyleemphasis{j} denote the \sphinxstyleemphasis{LAMMPS} particle types for which this
pair style is defined. Note that \sphinxstyleemphasis{i} and \sphinxstyleemphasis{j} can be different, i.e.,
\sphinxstyleemphasis{ulsph} cross interactions between different particle types are
allowed. However, \sphinxstyleemphasis{i}\textendash{}\sphinxstyleemphasis{i} respectively \sphinxstyleemphasis{j}\textendash{}\sphinxstyleemphasis{j} pair\_coeff lines have
to precede a cross interaction.  In contrast to the usual \sphinxstyleemphasis{LAMMPS}
\sphinxstyleemphasis{pair coeff} definitions, which are given solely a number of floats
and integers, the \sphinxstyleemphasis{ulsph} \sphinxstyleemphasis{pair coeff} definition is organized using
keywords. These keywords mark the beginning of different sets of
parameters for particle properties, material constitutive models, and
damage models. The \sphinxstyleemphasis{pair coeff} line must be terminated with the
\sphinxstyleemphasis{*END} keyword. The use the line continuation operator \sphinxstyleemphasis{\&} is
recommended. A typical invocation of the \sphinxstyleemphasis{ulsph} for a solid body
would consist of an equation of state for computing the pressure (the
diagonal components of the stress tensor), and a material model to
compute shear stresses (the off\sphinxhyphen{}diagonal components of the stress
tensor).

\sphinxAtStartPar
Note that the use of *GRADIENT\_CORRECTION can lead to severe numerical
instabilities. For a general fluid simulation, *NO\_GRADIENT\_CORRECTION
is recommended.

\sphinxAtStartPar
Please see the \sphinxhref{PDF/MACHDYN\_LAMMPS\_userguide.pdf}{SMD user guide} for a
complete listing of the possible keywords and material models.


\bigskip\hrule\bigskip



\subsection{Mixing, shift, table, tail correction, restart, rRESPA info}
\label{\detokenize{pair_smd_ulsph:mixing-shift-table-tail-correction-restart-rrespa-info}}
\sphinxAtStartPar
No mixing is performed automatically.  Currently, no part of MACHDYN
supports restarting nor minimization.  rRESPA does not apply to this
pair style.


\bigskip\hrule\bigskip



\subsection{Restrictions}
\label{\detokenize{pair_smd_ulsph:restrictions}}
\sphinxAtStartPar
This fix is part of the MACHDYN package.  It is only enabled if
LAMMPS was built with that package.  See the {\hyperref[\detokenize{Build_package::doc}]{\sphinxcrossref{\DUrole{doc}{Build package}}}} page for more info.


\subsection{Related commands}
\label{\detokenize{pair_smd_ulsph:related-commands}}
\sphinxAtStartPar
{\hyperref[\detokenize{pair_coeff::doc}]{\sphinxcrossref{\DUrole{doc}{pair\_coeff}}}}


\subsection{Default}
\label{\detokenize{pair_smd_ulsph:default}}
\sphinxAtStartPar
none

\sphinxstepscope

\index{pair\_style smtbq@\spxentry{pair\_style smtbq}}\ignorespaces 

\section{pair\_style smtbq command}
\label{\detokenize{pair_smtbq:pair-style-smtbq-command}}\label{\detokenize{pair_smtbq:index-0}}\label{\detokenize{pair_smtbq::doc}}

\subsection{Syntax}
\label{\detokenize{pair_smtbq:syntax}}
\begin{sphinxVerbatim}[commandchars=\\\{\}]
\PYG{k}{pair\PYGZus{}style}\PYG{+w}{ }\PYG{n}{smtbq}
\end{sphinxVerbatim}


\subsection{Examples}
\label{\detokenize{pair_smtbq:examples}}
\begin{sphinxVerbatim}[commandchars=\\\{\}]
\PYG{k}{pair\PYGZus{}style}\PYG{+w}{ }\PYG{n}{smtbq}
\PYG{k}{pair\PYGZus{}coeff}\PYG{+w}{ }\PYG{o}{*}\PYG{+w}{ }\PYG{o}{*}\PYG{+w}{ }\PYG{n}{ffield.smtbq.Al2O3}\PYG{+w}{ }\PYG{n}{O}\PYG{+w}{ }\PYG{n}{Al}
\end{sphinxVerbatim}


\subsection{Description}
\label{\detokenize{pair_smtbq:description}}
\sphinxAtStartPar
This pair style computes a variable charge SMTB\sphinxhyphen{}Q (Second\sphinxhyphen{}Moment
tight\sphinxhyphen{}Binding QEq) potential as described in {\hyperref[\detokenize{pair_smtbq:smtb-q-1}]{\sphinxcrossref{\DUrole{std,std-ref}{SMTB\sphinxhyphen{}Q\_1}}}} and
{\hyperref[\detokenize{pair_smtbq:smtb-q-2}]{\sphinxcrossref{\DUrole{std,std-ref}{SMTB\sphinxhyphen{}Q\_2}}}}.
This potential was first proposed in {\hyperref[\detokenize{pair_smtbq:smtb-q-0}]{\sphinxcrossref{\DUrole{std,std-ref}{SMTB\sphinxhyphen{}Q\_0}}}}.
Briefly, the energy of metallic\sphinxhyphen{}oxygen systems is given by three contributions:
\begin{equation*}
\begin{split}E_{tot} & =  E_{ES} + E_{OO} + E_{MO} \\
E_{ES}  & =  \sum_i{\biggl[ \chi_{i}^{0}Q_i + \frac{1}{2}J_{i}^{0}Q_{i}^{2} +
\frac{1}{2} \sum_{j\neq i}{ J_{ij}(r_{ij})f_{cut}^{R_{coul}}(r_{ij})Q_i Q_j } \biggr] } \\
E_{OO}  & =  \sum_{i,j}^{i,j = O}{\biggl[Cexp( -\frac{r_{ij}}{\rho} ) - Df_{cut}^{r_1^{OO}r_2^{OO}}(r_{ij}) exp(Br_{ij})\biggr]}  \\
E_{MO}  & =  \sum_i{E_{cov}^{i} + \sum_{j\neq i}{ Af_{cut}^{r_{c1}r_{c2}}(r_{ij})exp\bigl[-p(\frac{r_{ij}}{r_0} -1) \bigr] } }\end{split}
\end{equation*}
\sphinxAtStartPar
where \(E_{tot}\) is the total potential energy of the system,
\(E_{ES}\) is the electrostatic part of the total energy,
\(E_{OO}\) is the interaction between oxygen atoms and
\(E_{MO}\) is a short\sphinxhyphen{}range interaction between metal and oxygen
atoms. This interactions depend on interatomic distance \(r_{ij}\)
and/or the charge \(Q_{i}\) of atoms \sphinxstyleemphasis{i}. Cut\sphinxhyphen{}off function enables
smooth convergence to zero interaction.

\sphinxAtStartPar
The parameters appearing in the upper expressions are set in the
ffield.SMTBQ.Syst file where Syst corresponds to the selected system
(e.g. field.SMTBQ.Al2O3). Examples for \(\mathrm{TiO_2}\),
\(\mathrm{Al_2O_3}\) are provided.  A single pair\_coeff command
is used with the SMTBQ styles which provides the path to the potential
file with parameters for needed elements. These are mapped to LAMMPS
atom types by specifying additional arguments after the potential
filename in the pair\_coeff command. Note that atom type 1 must always
correspond to oxygen atoms. As an example, to simulate a \(\mathrm{TiO_2}\) system,
atom type 1 has to be oxygen and atom type 2 Ti. The following
pair\_coeff command should then be used:

\begin{sphinxVerbatim}[commandchars=\\\{\}]
\PYG{k}{pair\PYGZus{}coeff}\PYG{+w}{ }\PYG{o}{*}\PYG{+w}{ }\PYG{o}{*}\PYG{+w}{ }\PYG{n}{PathToLammps}\PYG{o}{/}\PYG{n}{potentials}\PYG{o}{/}\PYG{n}{ffield.smtbq.TiO2}\PYG{+w}{ }\PYG{n}{O}\PYG{+w}{ }\PYG{n}{Ti}
\end{sphinxVerbatim}

\sphinxAtStartPar
The electrostatic part of the energy consists of two components

\sphinxAtStartPar
self\sphinxhyphen{}energy of atom \sphinxstyleemphasis{i} in the form of a second order charge dependent
polynomial and a long\sphinxhyphen{}range Coulombic electrostatic interaction. The
latter uses the wolf summation method described in {\hyperref[\detokenize{pair_smtbq:wolf2}]{\sphinxcrossref{\DUrole{std,std-ref}{Wolf}}}},
spherically truncated at a longer cutoff, \(R_{coul}\). The
charge of each ion is modeled by an orbital Slater which depends on
the principal quantum number (\sphinxstyleemphasis{n}) of the outer orbital shared by the
ion.

\sphinxAtStartPar
Interaction between oxygen, \(E_{OO}\), consists of two parts,
an attractive and a repulsive part. The attractive part is effective
only at short range (\textless{} \(r_2^{OO}\)). The attractive
contribution was optimized to study surfaces reconstruction
(e.g. {\hyperref[\detokenize{pair_smtbq:smtb-q-2}]{\sphinxcrossref{\DUrole{std,std-ref}{SMTB\sphinxhyphen{}Q\_2}}}} in \(\mathrm{TiO_2}\)) and is not necessary
for oxide bulk modeling. The repulsive part is the Pauli interaction
between the electron clouds of oxygen. The Pauli repulsion and the
coulombic electrostatic interaction have same cut off value. In the
ffield.SMTBQ.Syst, the keyword \sphinxstyleemphasis{‘buck’} allows to consider only the
repulsive O\sphinxhyphen{}O interactions. The keyword \sphinxstyleemphasis{‘buckPlusAttr’} allows to
consider the repulsive and the attractive O\sphinxhyphen{}O interactions.

\sphinxAtStartPar
The short\sphinxhyphen{}range interaction between metal\sphinxhyphen{}oxygen, \(E_{MO}\) is
based on the second moment approximation of the density of states with
a N\sphinxhyphen{}body potential for the band energy term,
\(E^i_{cov}\), and a Born\sphinxhyphen{}Mayer type repulsive terms
as indicated by the keyword \sphinxstyleemphasis{‘second\_moment’} in the
ffield.SMTBQ.Syst. The energy band term is given by:
\begin{equation*}
\begin{split}E_{cov}^{i(i=M,O)} & = - \biggl\{\eta_i(\mu \xi^{0})^2 f_{cut}^{r_{c1}r_{c2}}(r_{ij})
\biggl( \sum_{j(j=O,M)}{ exp[ -2q(\frac{r_{ij}}{r_0} - 1)] } \biggr)
\delta Q_i \bigl( 2\frac{n_0}{\eta_i} - \delta Q_i \bigr) \biggr\}^{1/2} \\
\delta Q_i & =  | Q_i^{F} | - | Q_i |\end{split}
\end{equation*}
\sphinxAtStartPar
where \(\eta_i\) is the stoichiometry of atom \sphinxstyleemphasis{i},
\(\delta Q_i\) is the charge delocalization of atom \sphinxstyleemphasis{i},
compared to its formal charge
\(Q^F_i\). \(n_0\), the number of hybridized
orbitals, is calculated with to the atomic orbitals shared
\(d_i\) and the stoichiometry
\(\eta_i\). \(r_{c1}\) and \(r_{c2}\) are the two
cutoff radius around the fourth neighbors in the cutoff function.

\sphinxAtStartPar
In the formalism used here, \(\xi^0\) is the energy
parameter. \(\xi^0\) is in tight\sphinxhyphen{}binding approximation the
hopping integral between the hybridized orbitals of the cation and the
anion. In the literature we find many ways to write the hopping
integral depending on whether one takes the point of view of the anion
or cation. These are equivalent vision. The correspondence between the
two visions is explained in appendix A of the article in the
SrTiO3 {\hyperref[\detokenize{pair_smtbq:smtb-q-3}]{\sphinxcrossref{\DUrole{std,std-ref}{SMTB\sphinxhyphen{}Q\_3}}}} (parameter \(\beta\) shown in
this article is in fact the \(\beta_O\)). To summarize the
relationship between the hopping integral \(\xi^O\)  and the
others, we have in an oxide \(\mathrm{C_n O_m}\) the following
relationship:
\begin{equation*}
\begin{split}\xi^0 & = \frac{\xi_O}{m} = \frac{\xi_C}{n} \\
\frac{\beta_O}{\sqrt{m}} & = \frac{\beta_C}{\sqrt{n}} = \xi^0 \frac{\sqrt{m}+\sqrt{n}}{2}\end{split}
\end{equation*}
\sphinxAtStartPar
Thus parameter \(\mu\), indicated above, is given by \(\mu = \frac{1}{2}(\sqrt{n}+\sqrt{m})\)

\sphinxAtStartPar
The potential offers the possibility to consider the polarizability of
the electron clouds of oxygen by changing the slater radius of the
charge density around the oxygen atoms through the parameters \sphinxstyleemphasis{rBB, rB and
rS} in the ffield.SMTBQ.Syst. This change in radius is performed
according to the method developed by E. Maras
{\hyperref[\detokenize{pair_smtbq:smtb-q-2}]{\sphinxcrossref{\DUrole{std,std-ref}{SMTB\sphinxhyphen{}Q\_2}}}}. This method needs to determine the number of
nearest neighbors around the oxygen. This calculation is based on
first (\(r_{1n}\)) and second (\(r_{2n}\)) distances
neighbors.

\sphinxAtStartPar
The SMTB\sphinxhyphen{}Q potential is a variable charge potential. The equilibrium
charge on each atom is calculated by the electronegativity
equalization (QEq) method. See {\hyperref[\detokenize{pair_smtbq:rick3}]{\sphinxcrossref{\DUrole{std,std-ref}{Rick}}}} for further detail. One
can adjust the frequency, the maximum number of iterative loop and the
convergence of the equilibrium charge calculation. To obtain the
energy conservation in NVE thermodynamic ensemble, we recommend to use
a convergence parameter in the interval 10e\sphinxhyphen{}5 \sphinxhyphen{}
10e\sphinxhyphen{}6 eV.

\sphinxAtStartPar
The ffield.SMTBQ.Syst files are provided for few systems. They consist
of nine parts and the lines beginning with ‘\#’ are comments (note that
the number of comment lines matter). The first sections are on the
potential parameters and others are on the simulation options and
might be modified. Keywords are character type and must be enclosed in
quotation marks (‘’).
\begin{enumerate}
\sphinxsetlistlabels{\arabic}{enumi}{enumii}{}{)}%
\item {} 
\sphinxAtStartPar
Number of different element in the oxide:

\end{enumerate}
\begin{itemize}
\item {} 
\sphinxAtStartPar
N\_elem= 2 or 3

\item {} 
\sphinxAtStartPar
Divider line

\end{itemize}
\begin{enumerate}
\sphinxsetlistlabels{\arabic}{enumi}{enumii}{}{)}%
\setcounter{enumi}{1}
\item {} 
\sphinxAtStartPar
Atomic parameters

\end{enumerate}

\sphinxAtStartPar
For the anion (oxygen)
\begin{itemize}
\item {} 
\sphinxAtStartPar
Name of element (char) and stoichiometry in oxide

\item {} 
\sphinxAtStartPar
Formal charge and mass of element

\item {} 
\sphinxAtStartPar
Principal quantum number of outer orbital n), electronegativity (\(\chi^0_i\)) and hardness (\(J^0_i\))

\item {} 
\sphinxAtStartPar
Ionic radius parameters  : max coordination number (\sphinxstyleemphasis{coordBB} = 6 by default), bulk coordination number \sphinxstyleemphasis{(coordB)}, surface coordination number  \sphinxstyleemphasis{(coordS)} and \sphinxstyleemphasis{rBB, rB and rS}  the slater radius for each coordination number. (\sphinxstylestrong{note : If you don’t want to change the slater radius, use three identical radius values})

\item {} 
\sphinxAtStartPar
Number of orbital shared by the element in the oxide (\(d_i\))

\item {} 
\sphinxAtStartPar
Divider line

\end{itemize}

\sphinxAtStartPar
For each cations (metal):
\begin{itemize}
\item {} 
\sphinxAtStartPar
Name of element (char) and stoichiometry in oxide

\item {} 
\sphinxAtStartPar
Formal charge and mass of element

\item {} 
\sphinxAtStartPar
Number of electron in outer orbital \sphinxstyleemphasis{(ne)}, electronegativity (\(\chi^0_i\)), hardness (\(J^0_i\)) and \(r_{Slater}\) the slater radius for the cation.

\item {} 
\sphinxAtStartPar
Number of orbitals shared by the elements in the oxide (\(d_i\))

\item {} 
\sphinxAtStartPar
Divider line

\end{itemize}
\begin{enumerate}
\sphinxsetlistlabels{\arabic}{enumi}{enumii}{}{)}%
\setcounter{enumi}{2}
\item {} 
\sphinxAtStartPar
Potential parameters:

\end{enumerate}
\begin{itemize}
\item {} 
\sphinxAtStartPar
Keyword for element1, element2 and interaction potential
(‘second\_moment’ or ‘buck’ or ‘buckPlusAttr’) between element 1
and 2.  If the potential is ‘second\_moment’, specify ‘oxide’ or
‘metal’ for metal\sphinxhyphen{}oxygen or metal\sphinxhyphen{}metal interactions respectively.

\item {} 
\sphinxAtStartPar
Potential parameter:
\begin{itemize}
\item {} 
\sphinxAtStartPar
If type of potential is ‘second\_moment’ : A (eV), \sphinxstyleemphasis{p},
\(\zeta^0\) (eV) and \sphinxstyleemphasis{q}, \(r_{c1} (\AA)\), \(r_{c2}
(\AA)\) and \(r_0 (\AA)\)

\item {} 
\sphinxAtStartPar
If type of potential is ‘buck’ : \sphinxstyleemphasis{C} (eV) and \(\rho (\AA)\)

\item {} 
\sphinxAtStartPar
If type of potential is ‘buckPlusAttr’ : \sphinxstyleemphasis{C} (eV) and \(\rho
(\AA)\) \sphinxstyleemphasis{D} (eV), \sphinxstyleemphasis{B} \((\AA^{-1})\), \(r^{OO}_1 (\AA)\) and
\(r^{OO}_2 (\AA)\)

\end{itemize}

\item {} 
\sphinxAtStartPar
Divider line

\end{itemize}
\begin{enumerate}
\sphinxsetlistlabels{\arabic}{enumi}{enumii}{}{)}%
\setcounter{enumi}{3}
\item {} 
\sphinxAtStartPar
Tables parameters:

\end{enumerate}
\begin{itemize}
\item {} 
\sphinxAtStartPar
Cutoff radius for the Coulomb interaction (\(R_{coul}\))

\item {} 
\sphinxAtStartPar
Starting radius (\(r_{min} = 1,18845 \AA\)) and increments
(\(dr = 0.001 \AA\)) for creating the potential table.

\item {} 
\sphinxAtStartPar
Divider line

\end{itemize}
\begin{enumerate}
\sphinxsetlistlabels{\arabic}{enumi}{enumii}{}{)}%
\setcounter{enumi}{4}
\item {} 
\sphinxAtStartPar
Rick model parameter:

\end{enumerate}
\begin{itemize}
\item {} 
\sphinxAtStartPar
\sphinxstyleemphasis{Nevery} : parameter to set the frequency of the charge
resolution. The charges are evaluated each \sphinxstyleemphasis{Nevery} time steps.

\item {} 
\sphinxAtStartPar
Max number of iterative loop (\sphinxstyleemphasis{loopmax}) and convergence criterion
(\sphinxstyleemphasis{prec}) in eV of the charge resolution

\item {} 
\sphinxAtStartPar
Divider line

\end{itemize}
\begin{enumerate}
\sphinxsetlistlabels{\arabic}{enumi}{enumii}{}{)}%
\setcounter{enumi}{5}
\item {} 
\sphinxAtStartPar
Coordination parameter:

\end{enumerate}
\begin{itemize}
\item {} 
\sphinxAtStartPar
First (\(r_{1n}\)) and second (\(r_{2n}\)) neighbor distances
in angstroms

\item {} 
\sphinxAtStartPar
Divider line

\end{itemize}
\begin{enumerate}
\sphinxsetlistlabels{\arabic}{enumi}{enumii}{}{)}%
\setcounter{enumi}{6}
\item {} 
\sphinxAtStartPar
Charge initialization mode:

\end{enumerate}
\begin{itemize}
\item {} 
\sphinxAtStartPar
Keyword (\sphinxstyleemphasis{QInitMode}) and initial oxygen charge
(\(Q_{init}\)). If keyword = ‘true’, all oxygen charges are
initially set equal to \(Q_{init}\). The charges on the cations
are initially set in order to respect the neutrality of the box. If
keyword = ‘false’, all atom charges are initially set equal to 0 if
you use the {\hyperref[\detokenize{create_atoms::doc}]{\sphinxcrossref{\DUrole{doc}{create\_atoms}}}} command or the charge
specified in the file structure using {\hyperref[\detokenize{read_data::doc}]{\sphinxcrossref{\DUrole{doc}{read\_data}}}}
command.

\item {} 
\sphinxAtStartPar
Divider line

\end{itemize}
\begin{enumerate}
\sphinxsetlistlabels{\arabic}{enumi}{enumii}{}{)}%
\setcounter{enumi}{7}
\item {} 
\sphinxAtStartPar
Mode for the electronegativity equalization (Qeq)

\end{enumerate}
\begin{itemize}
\item {} 
\sphinxAtStartPar
Keyword (\sphinxstyleemphasis{mode}) followed by:
\begin{itemize}
\item {} 
\sphinxAtStartPar
QEqAll  (one QEq group) |   no parameters

\item {} 
\sphinxAtStartPar
QEqAllParallel (several QEq groups) |   no parameters

\item {} 
\sphinxAtStartPar
Surface |   zlim   (QEq only for z\textgreater{}zlim)

\end{itemize}

\item {} 
\sphinxAtStartPar
Parameter if necessary

\item {} 
\sphinxAtStartPar
Divider line

\end{itemize}
\begin{enumerate}
\sphinxsetlistlabels{\arabic}{enumi}{enumii}{}{)}%
\setcounter{enumi}{8}
\item {} 
\sphinxAtStartPar
Verbose

\end{enumerate}
\begin{itemize}
\item {} 
\sphinxAtStartPar
If you want the code to work in verbose mode or not : ‘true’ or ‘false’

\item {} 
\sphinxAtStartPar
If you want to print or not in the file ‘Energy\_component.txt’ the
three main contributions to the energy of the system according to the
description presented above : ‘true’ or ‘false’ and
\(N_{Energy}\). This option writes to the file every
\(N_{Energy}\) time steps. If the value is ‘false’ then
\(N_{Energy} = 0\). The file takes into account the possibility to
have several QEq groups \sphinxstyleemphasis{g} then it writes: time step, number of atoms
in group \sphinxstyleemphasis{g}, electrostatic part of energy, \(E_{ES}\), the
interaction between oxygen, \(E_{OO}\), and short range
metal\sphinxhyphen{}oxygen interaction, \(E_{MO}\).

\item {} 
\sphinxAtStartPar
If you want to print to the file ‘Electroneg\_component.txt’ the
electronegativity component (\(\frac{\partial E_{tot}}{\partial
Q_i}\)) or not: ‘true’ or ‘false’ and \(N_{Electroneg}\). This
option writes to the file every \(N_{Electroneg}\) time steps. If
the value is ‘false’ then \(N_{Electroneg} = 0\).  The file
consist of atom number \sphinxstyleemphasis{i}, atom type (1 for oxygen and \# higher
than 1 for metal), atom position: \sphinxstyleemphasis{x}, \sphinxstyleemphasis{y} and \sphinxstyleemphasis{z}, atomic charge
of atom \sphinxstyleemphasis{i}, electrostatic part of atom \sphinxstyleemphasis{i} electronegativity,
covalent part of atom \sphinxstyleemphasis{i} electronegativity, the hopping integral of
atom \sphinxstyleemphasis{i} \((Z\beta^2)_i\) and box electronegativity.

\end{itemize}

\begin{sphinxadmonition}{note}{Note:}
\sphinxAtStartPar
This last option slows down the calculation dramatically.  Use
only with a single processor simulation.
\end{sphinxadmonition}


\bigskip\hrule\bigskip



\subsection{Mixing, shift, table, tail correction, restart, rRESPA info}
\label{\detokenize{pair_smtbq:mixing-shift-table-tail-correction-restart-rrespa-info}}
\sphinxAtStartPar
This pair style does not support the {\hyperref[\detokenize{pair_modify::doc}]{\sphinxcrossref{\DUrole{doc}{pair\_modify}}}}
mix, shift, table, and tail options.

\sphinxAtStartPar
This pair style does not write its information to {\hyperref[\detokenize{restart::doc}]{\sphinxcrossref{\DUrole{doc}{binary restart files}}}}, since it is stored in potential files.  Thus, you
needs to re\sphinxhyphen{}specify the pair\_style and pair\_coeff commands in an input
script that reads a restart file.

\sphinxAtStartPar
This pair style can only be used via the \sphinxstyleemphasis{pair} keyword of the
{\hyperref[\detokenize{run_style::doc}]{\sphinxcrossref{\DUrole{doc}{run\_style respa}}}} command.  It does not support the
\sphinxstyleemphasis{inner}, \sphinxstyleemphasis{middle}, \sphinxstyleemphasis{outer} keywords.


\bigskip\hrule\bigskip



\subsection{Restrictions}
\label{\detokenize{pair_smtbq:restrictions}}
\sphinxAtStartPar
This pair style is part of the SMTBQ package and is only enabled
if LAMMPS is built with that package.  See the {\hyperref[\detokenize{Build_package::doc}]{\sphinxcrossref{\DUrole{doc}{Build package}}}} page for more info.

\sphinxAtStartPar
This potential requires using atom type 1 for oxygen and atom type
higher than 1 for metal atoms.

\sphinxAtStartPar
This pair style requires the {\hyperref[\detokenize{newton::doc}]{\sphinxcrossref{\DUrole{doc}{newton}}}} setting to be “on”
for pair interactions.

\sphinxAtStartPar
The SMTB\sphinxhyphen{}Q potential files provided with LAMMPS (see the potentials
directory) are parameterized for metal {\hyperref[\detokenize{units::doc}]{\sphinxcrossref{\DUrole{doc}{units}}}}.


\bigskip\hrule\bigskip



\subsection{Citing this work}
\label{\detokenize{pair_smtbq:citing-this-work}}
\sphinxAtStartPar
Please cite related publication: N. Salles, O. Politano, E. Amzallag
and R. Tetot, Comput. Mater. Sci. 111 (2016) 181\sphinxhyphen{}189


\bigskip\hrule\bigskip

\phantomsection\label{\detokenize{pair_smtbq:smtb-q-0}}
\sphinxAtStartPar
\sphinxstylestrong{(SMTB\sphinxhyphen{}Q\_0)} A. Hallil, E. Amzallag, S. Landron, R. Tetot,
Surface Science 605 738\sphinxhyphen{}745 (2011);
R. Tetot, A. Hallil, J. Creuze and I. Braems, EPL, 83 40001 (2008)

\phantomsection\label{\detokenize{pair_smtbq:smtb-q-1}}
\sphinxAtStartPar
\sphinxstylestrong{(SMTB\sphinxhyphen{}Q\_1)} N. Salles, O. Politano, E. Amzallag, R. Tetot,
Comput. Mater. Sci. 111 (2016) 181\sphinxhyphen{}189

\phantomsection\label{\detokenize{pair_smtbq:smtb-q-2}}
\sphinxAtStartPar
\sphinxstylestrong{(SMTB\sphinxhyphen{}Q\_2)} E. Maras, N. Salles, R. Tetot, T. Ala\sphinxhyphen{}Nissila,
H. Jonsson, J. Phys. Chem. C 2015, 119, 10391\sphinxhyphen{}10399

\phantomsection\label{\detokenize{pair_smtbq:smtb-q-3}}
\sphinxAtStartPar
\sphinxstylestrong{(SMTB\sphinxhyphen{}Q\_3)} R. Tetot, N. Salles, S. Landron, E. Amzallag, Surface
Science 616, 19\sphinxhyphen{}8722 28 (2013)

\phantomsection\label{\detokenize{pair_smtbq:wolf2}}
\sphinxAtStartPar
\sphinxstylestrong{(Wolf)} D. Wolf, P. Keblinski, S. R. Phillpot, J. Eggebrecht, J Chem
Phys, 110, 8254 (1999).

\phantomsection\label{\detokenize{pair_smtbq:rick3}}
\sphinxAtStartPar
\sphinxstylestrong{(Rick)} S. W. Rick, S. J. Stuart, B. J. Berne, J Chem Phys 101, 6141
(1994).

\sphinxstepscope

\index{pair\_style snap@\spxentry{pair\_style snap}}\ignorespaces 
\index{pair\_style snap/intel@\spxentry{pair\_style snap/intel}}\ignorespaces 
\index{pair\_style snap/kk@\spxentry{pair\_style snap/kk}}\ignorespaces 

\section{pair\_style snap command}
\label{\detokenize{pair_snap:pair-style-snap-command}}\label{\detokenize{pair_snap:index-2}}\label{\detokenize{pair_snap:index-1}}\label{\detokenize{pair_snap:index-0}}\label{\detokenize{pair_snap::doc}}
\sphinxAtStartPar
Accelerator Variants: \sphinxstyleemphasis{snap/intel}, \sphinxstyleemphasis{snap/kk}


\subsection{Syntax}
\label{\detokenize{pair_snap:syntax}}
\begin{sphinxVerbatim}[commandchars=\\\{\}]
\PYG{k}{pair\PYGZus{}style}\PYG{+w}{ }\PYG{n}{snap}
\end{sphinxVerbatim}


\subsection{Examples}
\label{\detokenize{pair_snap:examples}}
\begin{sphinxVerbatim}[commandchars=\\\{\}]
\PYG{k}{pair\PYGZus{}style}\PYG{+w}{ }\PYG{n}{snap}
\PYG{k}{pair\PYGZus{}coeff}\PYG{+w}{ }\PYG{o}{*}\PYG{+w}{ }\PYG{o}{*}\PYG{+w}{ }\PYG{n}{InP.snapcoeff}\PYG{+w}{ }\PYG{n}{InP.snapparam}\PYG{+w}{ }\PYG{n}{In}\PYG{+w}{ }\PYG{n}{In}\PYG{+w}{ }\PYG{n}{P}\PYG{+w}{ }\PYG{n}{P}
\end{sphinxVerbatim}


\subsection{Description}
\label{\detokenize{pair_snap:description}}
\sphinxAtStartPar
Pair style \sphinxstyleemphasis{snap} defines the spectral neighbor analysis potential
(SNAP), a machine\sphinxhyphen{}learning interatomic potential {\hyperref[\detokenize{pair_snap:thompson20142}]{\sphinxcrossref{\DUrole{std,std-ref}{(Thompson)}}}}.  Like the GAP framework of Bartok et
al. {\hyperref[\detokenize{pair_snap:bartok20102}]{\sphinxcrossref{\DUrole{std,std-ref}{(Bartok2010)}}}}, SNAP uses bispectrum components
to characterize the local neighborhood of each atom in a very general
way. The mathematical definition of the bispectrum calculation and its
derivatives w.r.t. atom positions is identical to that used by
{\hyperref[\detokenize{compute_sna_atom::doc}]{\sphinxcrossref{\DUrole{doc}{compute snap}}}}, which is used to fit SNAP
potentials to \sphinxstyleemphasis{ab initio} energy, force, and stress data.  In SNAP, the
total energy is decomposed into a sum over atom energies. The energy of
atom \sphinxstyleemphasis{i} is expressed as a weighted sum over bispectrum components.
\begin{equation*}
\begin{split}E^i_{SNAP}(B_1^i,...,B_K^i) = \beta^{\mu_i}_0 + \sum_{k=1}^K \beta_k^{\mu_i} B_k^i\end{split}
\end{equation*}
\sphinxAtStartPar
where \(B_k^i\) is the \sphinxstyleemphasis{k}\sphinxhyphen{}th bispectrum component of atom \sphinxstyleemphasis{i},
and \(\beta_k^{\mu_i}\) is the corresponding linear coefficient that
depends on \(\mu_i\), the SNAP element of atom \sphinxstyleemphasis{i}. The number of
bispectrum components used and their definitions depend on the value of
\sphinxstyleemphasis{twojmax} and other parameters defined in the SNAP parameter file
described below.  The bispectrum calculation is described in more detail
in {\hyperref[\detokenize{compute_sna_atom::doc}]{\sphinxcrossref{\DUrole{doc}{compute sna/atom}}}}.

\sphinxAtStartPar
Note that unlike for other potentials, cutoffs for SNAP potentials are
not set in the pair\_style or pair\_coeff command; they are specified in
the SNAP potential files themselves.

\sphinxAtStartPar
Only a single pair\_coeff command is used with the \sphinxstyleemphasis{snap} style which
specifies a SNAP coefficient file followed by a SNAP parameter file
and then N additional arguments specifying the mapping of SNAP
elements to LAMMPS atom types, where N is the number of
LAMMPS atom types:
\begin{itemize}
\item {} 
\sphinxAtStartPar
SNAP coefficient file

\item {} 
\sphinxAtStartPar
SNAP parameter file

\item {} 
\sphinxAtStartPar
N element names = mapping of SNAP elements to atom types

\end{itemize}

\sphinxAtStartPar
As an example, if a LAMMPS indium phosphide simulation has 4 atoms
types, with the first two being indium and the third and fourth being
phophorous, the pair\_coeff command would look like this:

\begin{sphinxVerbatim}[commandchars=\\\{\}]
\PYG{k}{pair\PYGZus{}coeff}\PYG{+w}{ }\PYG{o}{*}\PYG{+w}{ }\PYG{o}{*}\PYG{+w}{ }\PYG{n}{snap}\PYG{+w}{ }\PYG{n}{InP.snapcoeff}\PYG{+w}{ }\PYG{n}{InP.snapparam}\PYG{+w}{ }\PYG{n}{In}\PYG{+w}{ }\PYG{n}{In}\PYG{+w}{ }\PYG{n}{P}\PYG{+w}{ }\PYG{n}{P}
\end{sphinxVerbatim}

\sphinxAtStartPar
The first 2 arguments must be * * so as to span all LAMMPS atom types.
The two filenames are for the coefficient and parameter files, respectively.
The two trailing ‘In’ arguments map LAMMPS atom types 1 and 2 to the
SNAP ‘In’ element. The two trailing ‘P’ arguments map LAMMPS atom types
3 and 4 to the SNAP ‘P’ element.

\sphinxAtStartPar
If a SNAP mapping value is specified as NULL, the mapping is not
performed.  This can be used when a \sphinxstyleemphasis{snap} potential is used as part of
the \sphinxstyleemphasis{hybrid} pair style.  The NULL values are placeholders for atom
types that will be used with other potentials.

\sphinxAtStartPar
The name of the SNAP coefficient file usually ends in the “.snapcoeff”
extension. It may contain coefficients for many SNAP elements. The only
requirement is that each of the unique element names appearing in the
LAMMPS pair\_coeff command appear exactly once in the SNAP coefficient
file. It is okay if the SNAP coefficient file contains additional
elements not in the pair\_coeff command, except when using \sphinxstyleemphasis{chemflag}
(see below).  The name of the SNAP parameter file usually ends in the
“.snapparam” extension. It contains a small number of parameters that
define the overall form of the SNAP potential.  See the {\hyperref[\detokenize{pair_coeff::doc}]{\sphinxcrossref{\DUrole{doc}{pair\_coeff}}}} page for alternate ways to specify the path for these
files.

\sphinxAtStartPar
SNAP potentials are quite commonly combined with one or more other
LAMMPS pair styles using the \sphinxstyleemphasis{hybrid/overlay} pair style.  As an
example, the SNAP tantalum potential provided in the LAMMPS potentials
directory combines the \sphinxstyleemphasis{snap} and \sphinxstyleemphasis{zbl} pair styles. It is invoked by
the following commands:

\begin{sphinxVerbatim}[commandchars=\\\{\}]
\PYG{k}{variable }\PYG{n+nv+nvIdentifier}{zblcutinner}\PYG{+w}{ }\PYG{n}{equal}\PYG{+w}{ }\PYG{l+m}{4}
\PYG{k}{variable }\PYG{n+nv+nvIdentifier}{zblcutouter}\PYG{+w}{ }\PYG{n}{equal}\PYG{+w}{ }\PYG{l+m}{4.8}
\PYG{k}{variable }\PYG{n+nv+nvIdentifier}{zblz}\PYG{+w}{ }\PYG{n}{equal}\PYG{+w}{ }\PYG{l+m}{73}
\PYG{k}{pair\PYGZus{}style}\PYG{+w}{ }\PYG{n}{hybrid}\PYG{o}{/}\PYG{n}{overlay}\PYG{+w}{ }\PYG{o}{\PYGZam{}}
\PYG{n}{zbl}\PYG{+w}{ }\PYG{n+nv}{\PYGZdl{}\PYGZob{}}\PYG{n+nv}{zblcutinner}\PYG{n+nv}{\PYGZcb{}}\PYG{+w}{ }\PYG{n+nv}{\PYGZdl{}\PYGZob{}}\PYG{n+nv}{zblcutouter}\PYG{n+nv}{\PYGZcb{}}\PYG{+w}{ }\PYG{n}{snap}
\PYG{k}{pair\PYGZus{}coeff}\PYG{+w}{ }\PYG{o}{*}\PYG{+w}{ }\PYG{o}{*}\PYG{+w}{ }\PYG{n}{zbl}\PYG{+w}{ }\PYG{l+m}{0.0}
\PYG{k}{pair\PYGZus{}coeff}\PYG{+w}{ }\PYG{l+m}{1}\PYG{+w}{ }\PYG{l+m}{1}\PYG{+w}{ }\PYG{n}{zbl}\PYG{+w}{ }\PYG{n+nv}{\PYGZdl{}\PYGZob{}}\PYG{n+nv}{zblz}\PYG{n+nv}{\PYGZcb{}}
\PYG{k}{pair\PYGZus{}coeff}\PYG{+w}{ }\PYG{o}{*}\PYG{+w}{ }\PYG{o}{*}\PYG{+w}{ }\PYG{n}{snap}\PYG{+w}{ }\PYG{n}{Ta06A.snapcoeff}\PYG{+w}{ }\PYG{n}{Ta06A.snapparam}\PYG{+w}{ }\PYG{n}{Ta}
\end{sphinxVerbatim}

\sphinxAtStartPar
It is convenient to keep these commands in a separate file that can be
inserted in any LAMMPS input script using the {\hyperref[\detokenize{include::doc}]{\sphinxcrossref{\DUrole{doc}{include}}}}
command.

\sphinxAtStartPar
The top of the SNAP coefficient file can contain any number of blank and
comment lines (start with \#), but follows a strict format after
that. The first non\sphinxhyphen{}blank non\sphinxhyphen{}comment line must contain two integers:
\begin{itemize}
\item {} 
\sphinxAtStartPar
nelem  = Number of elements

\item {} 
\sphinxAtStartPar
ncoeff = Number of coefficients

\end{itemize}

\sphinxAtStartPar
This is followed by one block for each of the \sphinxstyleemphasis{nelem} elements.
The first line of each block contains three entries:
\begin{itemize}
\item {} 
\sphinxAtStartPar
Element name (text string)

\item {} 
\sphinxAtStartPar
R = Element radius (distance units)

\item {} 
\sphinxAtStartPar
w = Element weight (dimensionless)

\end{itemize}

\sphinxAtStartPar
This line is followed by \sphinxstyleemphasis{ncoeff} coefficients, one per line.

\sphinxAtStartPar
The SNAP parameter file can contain blank and comment lines (start
with \#) anywhere. Each non\sphinxhyphen{}blank non\sphinxhyphen{}comment line must contain one
keyword/value pair. The required keywords are \sphinxstyleemphasis{rcutfac} and
\sphinxstyleemphasis{twojmax}. Optional keywords are \sphinxstyleemphasis{rfac0}, \sphinxstyleemphasis{rmin0},
\sphinxstyleemphasis{switchflag}, \sphinxstyleemphasis{bzeroflag}, \sphinxstyleemphasis{quadraticflag}, \sphinxstyleemphasis{chemflag},
\sphinxstyleemphasis{bnormflag}, \sphinxstyleemphasis{wselfallflag},  \sphinxstyleemphasis{switchinnerflag},
\sphinxstyleemphasis{sinner}, \sphinxstyleemphasis{dinner}, \sphinxstyleemphasis{chunksize}, and \sphinxstyleemphasis{parallelthresh}.

\sphinxAtStartPar
The default values for these keywords are
\begin{itemize}
\item {} 
\sphinxAtStartPar
\sphinxstyleemphasis{rfac0} = 0.99363

\item {} 
\sphinxAtStartPar
\sphinxstyleemphasis{rmin0} = 0.0

\item {} 
\sphinxAtStartPar
\sphinxstyleemphasis{switchflag} = 1

\item {} 
\sphinxAtStartPar
\sphinxstyleemphasis{bzeroflag} = 1

\item {} 
\sphinxAtStartPar
\sphinxstyleemphasis{quadraticflag} = 0

\item {} 
\sphinxAtStartPar
\sphinxstyleemphasis{chemflag} = 0

\item {} 
\sphinxAtStartPar
\sphinxstyleemphasis{bnormflag} = 0

\item {} 
\sphinxAtStartPar
\sphinxstyleemphasis{wselfallflag} = 0

\item {} 
\sphinxAtStartPar
\sphinxstyleemphasis{switchinnerflag} = 0

\item {} 
\sphinxAtStartPar
\sphinxstyleemphasis{chunksize} = 32768

\item {} 
\sphinxAtStartPar
\sphinxstyleemphasis{parallelthresh} = 8192

\end{itemize}

\sphinxAtStartPar
For detailed definitions of all of these keywords,
see the {\hyperref[\detokenize{compute_sna_atom::doc}]{\sphinxcrossref{\DUrole{doc}{compute sna/atom}}}} doc page.

\sphinxAtStartPar
If \sphinxstyleemphasis{quadraticflag} is set to 1, then the SNAP energy expression includes
additional quadratic terms that have been shown to increase the overall
accuracy of the potential without much increase in computational cost
{\hyperref[\detokenize{pair_snap:wood20182}]{\sphinxcrossref{\DUrole{std,std-ref}{(Wood)}}}}.
\begin{equation*}
\begin{split}E^i_{SNAP}(\mathbf{B}^i) = \beta^{\mu_i}_0 + \boldsymbol{\beta}^{\mu_i} \cdot \mathbf{B}_i + \frac{1}{2}\mathbf{B}^t_i \cdot \boldsymbol{\alpha}^{\mu_i} \cdot \mathbf{B}_i\end{split}
\end{equation*}
\sphinxAtStartPar
where \(\mathbf{B}_i\) is the \sphinxstyleemphasis{K}\sphinxhyphen{}vector of bispectrum components,
\(\boldsymbol{\beta}^{\mu_i}\) is the \sphinxstyleemphasis{K}\sphinxhyphen{}vector of linear
coefficients for element \(\mu_i\), and
\(\boldsymbol{\alpha}^{\mu_i}\) is the symmetric \sphinxstyleemphasis{K} by \sphinxstyleemphasis{K} matrix
of quadratic coefficients.  The SNAP coefficient file should contain
\sphinxstyleemphasis{K}(\sphinxstyleemphasis{K}+1)/2 additional coefficients in each element block, the
upper\sphinxhyphen{}triangular elements of \(\boldsymbol{\alpha}^{\mu_i}\).

\sphinxAtStartPar
If \sphinxstyleemphasis{chemflag} is set to 1, then the energy expression is written in
terms of explicit multi\sphinxhyphen{}element bispectrum components indexed on ordered
triplets of elements, which has been shown to increase the ability of
the SNAP potential to capture energy differences in chemically complex
systems, at the expense of a significant increase in computational cost
{\hyperref[\detokenize{pair_snap:cusentino20202}]{\sphinxcrossref{\DUrole{std,std-ref}{(Cusentino)}}}}.
\begin{equation*}
\begin{split}E^i_{SNAP}(\mathbf{B}^i) = \beta^{\mu_i}_0 + \sum_{\kappa,\lambda,\mu} \boldsymbol{\beta}^{\kappa\lambda\mu}_{\mu_i} \cdot \mathbf{B}^{\kappa\lambda\mu}_i\end{split}
\end{equation*}
\sphinxAtStartPar
where \(\mathbf{B}^{\kappa\lambda\mu}_i\) is the \sphinxstyleemphasis{K}\sphinxhyphen{}vector of
bispectrum components for neighbors of elements \(\kappa\),
\(\lambda\), and \(\mu\) and
\(\boldsymbol{\beta}^{\kappa\lambda\mu}_{\mu_i}\) is the
corresponding \sphinxstyleemphasis{K}\sphinxhyphen{}vector of linear coefficients for element
\(\mu_i\). The SNAP coefficient file should contain a total of
\(K N_{elem}^3\) coefficients in each element block, where
\(N_{elem}\) is the number of elements in the SNAP coefficient file,
which must equal the number of unique elements appearing in the LAMMPS
pair\_coeff command, to avoid ambiguity in the number of coefficients.

\sphinxAtStartPar
The keyword \sphinxstyleemphasis{switchinnerflag} activates an additional switching function
that smoothly turns off contributions to the SNAP potential from neighbor
atoms at short separations. If \sphinxstyleemphasis{switchinnerflag} is set to 1 then
the additional keywords \sphinxstyleemphasis{sinner} and \sphinxstyleemphasis{dinner} must also be provided.
Each of these is followed by \sphinxstyleemphasis{nelements} values, where \sphinxstyleemphasis{nelements}
is the number of unique elements appearing in appearing in the LAMMPS
pair\_coeff command. The element order should correspond to the order
in which elements first appear in the pair\_coeff command reading from
left to right.

\sphinxAtStartPar
The keywords \sphinxstyleemphasis{chunksize} and \sphinxstyleemphasis{parallelthresh} are only applicable when
using the pair style \sphinxstyleemphasis{snap} with the KOKKOS package on GPUs and are
ignored otherwise.  The \sphinxstyleemphasis{chunksize} keyword controls the number of atoms
in each pass used to compute the bispectrum components and is used to
avoid running out of memory.  For example if there are 8192 atoms in the
simulation and the \sphinxstyleemphasis{chunksize} is set to 4096, the bispectrum
calculation will be broken up into two passes (running on a single GPU).
The \sphinxstyleemphasis{parallelthresh} keyword controls a crossover threshold for
performing extra parallelism.  For small systems, exposing additional
parallelism can be beneficial when there is not enough work to fully
saturate the GPU threads otherwise.  However, the extra parallelism also
leads to more divergence and can hurt performance when the system is
already large enough to saturate the GPU threads.  Extra parallelism
will be performed if the \sphinxstyleemphasis{chunksize} (or total number of atoms per GPU)
is smaller than \sphinxstyleemphasis{parallelthresh}.

\begin{sphinxadmonition}{note}{Note:}
\sphinxAtStartPar
The previously used \sphinxstyleemphasis{diagonalstyle} keyword was removed in 2019,
since all known SNAP potentials use the default value of 3.
\end{sphinxadmonition}


\bigskip\hrule\bigskip



\subsection{Mixing, shift, table, tail correction, restart, rRESPA info}
\label{\detokenize{pair_snap:mixing-shift-table-tail-correction-restart-rrespa-info}}
\sphinxAtStartPar
For atom type pairs I,J and I != J, where types I and J correspond to
two different element types, mixing is performed by LAMMPS with
user\sphinxhyphen{}specifiable parameters as described above.  You never need to
specify a pair\_coeff command with I != J arguments for this style.

\sphinxAtStartPar
This pair style does not support the {\hyperref[\detokenize{pair_modify::doc}]{\sphinxcrossref{\DUrole{doc}{pair\_modify}}}}
shift, table, and tail options.

\sphinxAtStartPar
This pair style does not write its information to {\hyperref[\detokenize{restart::doc}]{\sphinxcrossref{\DUrole{doc}{binary restart
files}}}}, since it is stored in potential files.  Thus, you need
to re\sphinxhyphen{}specify the pair\_style and pair\_coeff commands in an input script
that reads a restart file.

\sphinxAtStartPar
This pair style can only be used via the \sphinxstyleemphasis{pair} keyword of the
{\hyperref[\detokenize{run_style::doc}]{\sphinxcrossref{\DUrole{doc}{run\_style respa}}}} command.  It does not support the
\sphinxstyleemphasis{inner}, \sphinxstyleemphasis{middle}, \sphinxstyleemphasis{outer} keywords.


\bigskip\hrule\bigskip


\sphinxAtStartPar
Styles with a \sphinxstyleemphasis{gpu}, \sphinxstyleemphasis{intel}, \sphinxstyleemphasis{kk}, \sphinxstyleemphasis{omp}, or \sphinxstyleemphasis{opt} suffix are
functionally the same as the corresponding style without the suffix.
They have been optimized to run faster, depending on your available
hardware, as discussed on the {\hyperref[\detokenize{Speed_packages::doc}]{\sphinxcrossref{\DUrole{doc}{Accelerator packages}}}}
page.  The accelerated styles take the same arguments and should
produce the same results, except for round\sphinxhyphen{}off and precision issues.

\sphinxAtStartPar
These accelerated styles are part of the GPU, INTEL, KOKKOS,
OPENMP, and OPT packages, respectively.  They are only enabled if
LAMMPS was built with those packages.  See the {\hyperref[\detokenize{Build_package::doc}]{\sphinxcrossref{\DUrole{doc}{Build package}}}} page for more info.

\sphinxAtStartPar
You can specify the accelerated styles explicitly in your input script
by including their suffix, or you can use the {\hyperref[\detokenize{Run_options::doc}]{\sphinxcrossref{\DUrole{doc}{\sphinxhyphen{}suffix command\sphinxhyphen{}line switch}}}} when you invoke LAMMPS, or you can use the
{\hyperref[\detokenize{suffix::doc}]{\sphinxcrossref{\DUrole{doc}{suffix}}}} command in your input script.

\sphinxAtStartPar
See the {\hyperref[\detokenize{Speed_packages::doc}]{\sphinxcrossref{\DUrole{doc}{Accelerator packages}}}} page for more
instructions on how to use the accelerated styles effectively.


\bigskip\hrule\bigskip



\subsection{Restrictions}
\label{\detokenize{pair_snap:restrictions}}
\sphinxAtStartPar
This style is part of the ML\sphinxhyphen{}SNAP package.  It is only enabled if LAMMPS
was built with that package.  See the {\hyperref[\detokenize{Build_package::doc}]{\sphinxcrossref{\DUrole{doc}{Build package}}}} page for more info.

\sphinxAtStartPar
The \sphinxstyleemphasis{snap/intel} accelerator variant will \sphinxstyleemphasis{only} be available if LAMMPS
is built with Intel \sphinxstyleemphasis{compilers} and for CPUs with AVX\sphinxhyphen{}512 support.
While the INTEL package in general allows multiple floating point
precision modes to be selected, \sphinxstyleemphasis{snap/intel} will currently always use
full double precision regardless of the precision mode selected.
Additionally, the \sphinxstyleemphasis{intel} variant of snap will \sphinxstylestrong{NOT} use multiple
threads with OpenMP.


\subsection{Related commands}
\label{\detokenize{pair_snap:related-commands}}
\sphinxAtStartPar
{\hyperref[\detokenize{compute_sna_atom::doc}]{\sphinxcrossref{\DUrole{doc}{compute sna/atom}}}},
{\hyperref[\detokenize{compute_sna_atom::doc}]{\sphinxcrossref{\DUrole{doc}{compute snad/atom}}}},
{\hyperref[\detokenize{compute_sna_atom::doc}]{\sphinxcrossref{\DUrole{doc}{compute snav/atom}}}},
{\hyperref[\detokenize{compute_sna_atom::doc}]{\sphinxcrossref{\DUrole{doc}{compute snap}}}}


\subsection{Default}
\label{\detokenize{pair_snap:default}}
\sphinxAtStartPar
none


\bigskip\hrule\bigskip

\phantomsection\label{\detokenize{pair_snap:thompson20142}}
\sphinxAtStartPar
\sphinxstylestrong{(Thompson)} Thompson, Swiler, Trott, Foiles, Tucker, J Comp Phys, 285, 316 (2015).

\phantomsection\label{\detokenize{pair_snap:bartok20102}}
\sphinxAtStartPar
\sphinxstylestrong{(Bartok2010)} Bartok, Payne, Kondor, Csanyi, Phys Rev Lett, 104, 136403 (2010).

\phantomsection\label{\detokenize{pair_snap:wood20182}}
\sphinxAtStartPar
\sphinxstylestrong{(Wood)} Wood and Thompson, J Chem Phys, 148, 241721, (2018)

\phantomsection\label{\detokenize{pair_snap:cusentino20202}}
\sphinxAtStartPar
\sphinxstylestrong{(Cusentino)} Cusentino, Wood, Thompson, J Phys Chem A, 124, 5456, (2020)

\sphinxstepscope

\index{pair\_style soft@\spxentry{pair\_style soft}}\ignorespaces 
\index{pair\_style soft/gpu@\spxentry{pair\_style soft/gpu}}\ignorespaces 
\index{pair\_style soft/kk@\spxentry{pair\_style soft/kk}}\ignorespaces 
\index{pair\_style soft/omp@\spxentry{pair\_style soft/omp}}\ignorespaces 

\section{pair\_style soft command}
\label{\detokenize{pair_soft:pair-style-soft-command}}\label{\detokenize{pair_soft:index-3}}\label{\detokenize{pair_soft:index-2}}\label{\detokenize{pair_soft:index-1}}\label{\detokenize{pair_soft:index-0}}\label{\detokenize{pair_soft::doc}}
\sphinxAtStartPar
Accelerator Variants: \sphinxstyleemphasis{soft/gpu}, \sphinxstyleemphasis{soft/kk}, \sphinxstyleemphasis{soft/omp}


\subsection{Syntax}
\label{\detokenize{pair_soft:syntax}}
\begin{sphinxVerbatim}[commandchars=\\\{\}]
\PYG{k}{pair\PYGZus{}style}\PYG{+w}{ }\PYG{n}{soft}\PYG{+w}{ }\PYG{n}{cutoff}
\end{sphinxVerbatim}
\begin{itemize}
\item {} 
\sphinxAtStartPar
cutoff = global cutoff for soft interactions (distance units)

\end{itemize}


\subsection{Examples}
\label{\detokenize{pair_soft:examples}}
\begin{sphinxVerbatim}[commandchars=\\\{\}]
\PYG{k}{pair\PYGZus{}style}\PYG{+w}{ }\PYG{n}{soft}\PYG{+w}{ }\PYG{l+m}{1.0}
\PYG{k}{pair\PYGZus{}coeff}\PYG{+w}{ }\PYG{o}{*}\PYG{+w}{ }\PYG{o}{*}\PYG{+w}{ }\PYG{l+m}{10.0}
\PYG{k}{pair\PYGZus{}coeff}\PYG{+w}{ }\PYG{l+m}{1}\PYG{+w}{ }\PYG{l+m}{1}\PYG{+w}{ }\PYG{l+m}{10.0}\PYG{+w}{ }\PYG{l+m}{3.0}

\PYG{k}{pair\PYGZus{}style}\PYG{+w}{ }\PYG{n}{soft}\PYG{+w}{ }\PYG{l+m}{1.0}
\PYG{k}{pair\PYGZus{}coeff}\PYG{+w}{ }\PYG{o}{*}\PYG{+w}{ }\PYG{o}{*}\PYG{+w}{ }\PYG{l+m}{0.0}
\PYG{k}{variable }\PYG{n+nv+nvIdentifier}{prefactor}\PYG{+w}{ }\PYG{n}{equal}\PYG{+w}{ }\PYG{n}{ramp}\PYG{n+nv}{(}\PYG{n+nv}{0,30}\PYG{n+nv}{)}
\PYG{k}{fix }\PYG{n+nv+nvIdentifier}{1}\PYG{+w}{ }\PYG{n+nv+nvIdentifier}{all}\PYG{+w}{ }\PYG{n}{adapt}\PYG{+w}{ }\PYG{l+m}{1}\PYG{+w}{ }\PYG{n}{pair}\PYG{+w}{ }\PYG{n}{soft}\PYG{+w}{ }\PYG{n}{a}\PYG{+w}{ }\PYG{o}{*}\PYG{+w}{ }\PYG{o}{*}\PYG{+w}{ }\PYG{n}{v\PYGZus{}prefactor}
\end{sphinxVerbatim}


\subsection{Description}
\label{\detokenize{pair_soft:description}}
\sphinxAtStartPar
Style \sphinxstyleemphasis{soft} computes pairwise interactions with the formula
\begin{equation*}
\begin{split}E = A \left[ 1 + \cos\left(\frac{\pi r}{r_c}\right) \right]
\qquad r < r_c\end{split}
\end{equation*}
\sphinxAtStartPar
It is useful for pushing apart overlapping atoms, since it does not
blow up as r goes to 0.  A is a prefactor that can be made to vary in
time from the start to the end of the run (see discussion below),
e.g. to start with a very soft potential and slowly harden the
interactions over time.  \(r_c\) is the cutoff.
See the {\hyperref[\detokenize{fix_nve_limit::doc}]{\sphinxcrossref{\DUrole{doc}{fix nve/limit}}}} command for another way
to push apart overlapping atoms.

\sphinxAtStartPar
The following coefficients must be defined for each pair of atom types
via the {\hyperref[\detokenize{pair_coeff::doc}]{\sphinxcrossref{\DUrole{doc}{pair\_coeff}}}} command as in the examples above,
or in the data file or restart files read by the
{\hyperref[\detokenize{read_data::doc}]{\sphinxcrossref{\DUrole{doc}{read\_data}}}} or {\hyperref[\detokenize{read_restart::doc}]{\sphinxcrossref{\DUrole{doc}{read\_restart}}}}
commands, or by mixing as described below:
\begin{itemize}
\item {} 
\sphinxAtStartPar
A (energy units)

\item {} 
\sphinxAtStartPar
cutoff (distance units)

\end{itemize}

\sphinxAtStartPar
The last coefficient is optional.  If not specified, the global soft
cutoff is used.

\begin{sphinxadmonition}{note}{Note:}
\sphinxAtStartPar
The syntax for {\hyperref[\detokenize{pair_coeff::doc}]{\sphinxcrossref{\DUrole{doc}{pair\_coeff}}}} with a single A
coeff is different in the current version of LAMMPS than in older
versions which took two values, Astart and Astop, to ramp between
them.  This functionality is now available in a more general form
through the {\hyperref[\detokenize{fix_adapt::doc}]{\sphinxcrossref{\DUrole{doc}{fix adapt}}}} command, as explained below.
Note that if you use an old input script and specify Astart and Astop
without a cutoff, then LAMMPS will interpret that as A and a cutoff,
which is probably not what you want.
\end{sphinxadmonition}

\sphinxAtStartPar
The {\hyperref[\detokenize{fix_adapt::doc}]{\sphinxcrossref{\DUrole{doc}{fix adapt}}}} command can be used to vary A for one
or more pair types over the course of a simulation, in which case
pair\_coeff settings for A must still be specified, but will be
overridden.  For example these commands will vary the prefactor A for
all pairwise interactions from 0.0 at the beginning to 30.0 at the end
of a run:

\begin{sphinxVerbatim}[commandchars=\\\{\}]
\PYG{k}{variable }\PYG{n+nv+nvIdentifier}{prefactor}\PYG{+w}{ }\PYG{n}{equal}\PYG{+w}{ }\PYG{n}{ramp}\PYG{n+nv}{(}\PYG{n+nv}{0,30}\PYG{n+nv}{)}
\PYG{k}{fix }\PYG{n+nv+nvIdentifier}{1}\PYG{+w}{ }\PYG{n+nv+nvIdentifier}{all}\PYG{+w}{ }\PYG{n}{adapt}\PYG{+w}{ }\PYG{l+m}{1}\PYG{+w}{ }\PYG{n}{pair}\PYG{+w}{ }\PYG{n}{soft}\PYG{+w}{ }\PYG{n}{a}\PYG{+w}{ }\PYG{o}{*}\PYG{+w}{ }\PYG{o}{*}\PYG{+w}{ }\PYG{n}{v\PYGZus{}prefactor}
\end{sphinxVerbatim}

\sphinxAtStartPar
Note that a formula defined by an {\hyperref[\detokenize{variable::doc}]{\sphinxcrossref{\DUrole{doc}{equal\sphinxhyphen{}style variable}}}}
can use the current timestep, elapsed time in the current run, elapsed
time since the beginning of a series of runs, as well as access other
variables.


\bigskip\hrule\bigskip


\sphinxAtStartPar
Styles with a \sphinxstyleemphasis{gpu}, \sphinxstyleemphasis{intel}, \sphinxstyleemphasis{kk}, \sphinxstyleemphasis{omp}, or \sphinxstyleemphasis{opt} suffix are
functionally the same as the corresponding style without the suffix.
They have been optimized to run faster, depending on your available
hardware, as discussed on the {\hyperref[\detokenize{Speed_packages::doc}]{\sphinxcrossref{\DUrole{doc}{Accelerator packages}}}}
page.  The accelerated styles take the same arguments and should
produce the same results, except for round\sphinxhyphen{}off and precision issues.

\sphinxAtStartPar
These accelerated styles are part of the GPU, INTEL, KOKKOS,
OPENMP, and OPT packages, respectively.  They are only enabled if
LAMMPS was built with those packages.  See the {\hyperref[\detokenize{Build_package::doc}]{\sphinxcrossref{\DUrole{doc}{Build package}}}} page for more info.

\sphinxAtStartPar
You can specify the accelerated styles explicitly in your input script
by including their suffix, or you can use the {\hyperref[\detokenize{Run_options::doc}]{\sphinxcrossref{\DUrole{doc}{\sphinxhyphen{}suffix command\sphinxhyphen{}line switch}}}} when you invoke LAMMPS, or you can use the
{\hyperref[\detokenize{suffix::doc}]{\sphinxcrossref{\DUrole{doc}{suffix}}}} command in your input script.

\sphinxAtStartPar
See the {\hyperref[\detokenize{Speed_packages::doc}]{\sphinxcrossref{\DUrole{doc}{Accelerator packages}}}} page for more
instructions on how to use the accelerated styles effectively.


\bigskip\hrule\bigskip



\subsection{Mixing, shift, table, tail correction, restart, rRESPA info}
\label{\detokenize{pair_soft:mixing-shift-table-tail-correction-restart-rrespa-info}}
\sphinxAtStartPar
For atom type pairs I,J and I != J, the A coefficient and cutoff
distance for this pair style can be mixed.  A is always mixed via a
\sphinxstyleemphasis{geometric} rule.  The cutoff is mixed according to the pair\_modify
mix value.  The default mix value is \sphinxstyleemphasis{geometric}.  See the
“pair\_modify” command for details.

\sphinxAtStartPar
This pair style does not support the {\hyperref[\detokenize{pair_modify::doc}]{\sphinxcrossref{\DUrole{doc}{pair\_modify}}}}
shift option, since the pair interaction goes to 0.0 at the cutoff.

\sphinxAtStartPar
The {\hyperref[\detokenize{pair_modify::doc}]{\sphinxcrossref{\DUrole{doc}{pair\_modify}}}} table and tail options are not
relevant for this pair style.

\sphinxAtStartPar
This pair style writes its information to {\hyperref[\detokenize{restart::doc}]{\sphinxcrossref{\DUrole{doc}{binary restart files}}}}, so pair\_style and pair\_coeff commands do not need
to be specified in an input script that reads a restart file.

\sphinxAtStartPar
This pair style can only be used via the \sphinxstyleemphasis{pair} keyword of the
{\hyperref[\detokenize{run_style::doc}]{\sphinxcrossref{\DUrole{doc}{run\_style respa}}}} command.  It does not support the
\sphinxstyleemphasis{inner}, \sphinxstyleemphasis{middle}, \sphinxstyleemphasis{outer} keywords.


\bigskip\hrule\bigskip



\subsection{Restrictions}
\label{\detokenize{pair_soft:restrictions}}\begin{quote}

\sphinxAtStartPar
none
\end{quote}


\subsection{Related commands}
\label{\detokenize{pair_soft:related-commands}}
\sphinxAtStartPar
{\hyperref[\detokenize{pair_coeff::doc}]{\sphinxcrossref{\DUrole{doc}{pair\_coeff}}}}, {\hyperref[\detokenize{fix_nve_limit::doc}]{\sphinxcrossref{\DUrole{doc}{fix nve/limit}}}}, {\hyperref[\detokenize{fix_adapt::doc}]{\sphinxcrossref{\DUrole{doc}{fix adapt}}}}


\subsection{Default}
\label{\detokenize{pair_soft:default}}
\sphinxAtStartPar
none

\sphinxstepscope

\index{pair\_style sph/heatconduction@\spxentry{pair\_style sph/heatconduction}}\ignorespaces 
\index{pair\_style sph/heatconduction/gpu@\spxentry{pair\_style sph/heatconduction/gpu}}\ignorespaces 

\section{pair\_style sph/heatconduction command}
\label{\detokenize{pair_sph_heatconduction:pair-style-sph-heatconduction-command}}\label{\detokenize{pair_sph_heatconduction:index-1}}\label{\detokenize{pair_sph_heatconduction:index-0}}\label{\detokenize{pair_sph_heatconduction::doc}}
\sphinxAtStartPar
Accelerator Variants: \sphinxstyleemphasis{sph/heatconduction/gpu}


\subsection{Syntax}
\label{\detokenize{pair_sph_heatconduction:syntax}}
\begin{sphinxVerbatim}[commandchars=\\\{\}]
\PYG{k}{pair\PYGZus{}style}\PYG{+w}{ }\PYG{n}{sph}\PYG{o}{/}\PYG{n}{heatconduction}
\end{sphinxVerbatim}


\subsection{Examples}
\label{\detokenize{pair_sph_heatconduction:examples}}
\begin{sphinxVerbatim}[commandchars=\\\{\}]
\PYG{k}{pair\PYGZus{}style}\PYG{+w}{ }\PYG{n}{sph}\PYG{o}{/}\PYG{n}{heatconduction}
\PYG{k}{pair\PYGZus{}coeff}\PYG{+w}{ }\PYG{o}{*}\PYG{+w}{ }\PYG{o}{*}\PYG{+w}{ }\PYG{l+m}{1.0}\PYG{+w}{ }\PYG{l+m}{2.4}
\end{sphinxVerbatim}


\subsection{Description}
\label{\detokenize{pair_sph_heatconduction:description}}
\sphinxAtStartPar
The sph/heatconduction style computes heat transport between SPH particles.
The transport model is the diffusion equation for the internal energy.

\sphinxAtStartPar
See \sphinxhref{PDF/SPH\_LAMMPS\_userguide.pdf}{this PDF guide} to using SPH in
LAMMPS.

\sphinxAtStartPar
The following coefficients must be defined for each pair of atoms
types via the {\hyperref[\detokenize{pair_coeff::doc}]{\sphinxcrossref{\DUrole{doc}{pair\_coeff}}}} command as in the examples
above.
\begin{itemize}
\item {} 
\sphinxAtStartPar
D diffusion coefficient (length\textasciicircum{}2/time units)

\item {} 
\sphinxAtStartPar
h kernel function cutoff (distance units)

\end{itemize}


\bigskip\hrule\bigskip



\subsection{Mixing, shift, table, tail correction, restart, rRESPA info}
\label{\detokenize{pair_sph_heatconduction:mixing-shift-table-tail-correction-restart-rrespa-info}}
\sphinxAtStartPar
This style does not support mixing.  Thus, coefficients for all
I,J pairs must be specified explicitly.

\sphinxAtStartPar
This style does not support the {\hyperref[\detokenize{pair_modify::doc}]{\sphinxcrossref{\DUrole{doc}{pair\_modify}}}}
shift, table, and tail options.

\sphinxAtStartPar
This style does not write information to {\hyperref[\detokenize{restart::doc}]{\sphinxcrossref{\DUrole{doc}{binary restart files}}}}.  Thus, you need to re\sphinxhyphen{}specify the pair\_style and
pair\_coeff commands in an input script that reads a restart file.

\sphinxAtStartPar
This style can only be used via the \sphinxstyleemphasis{pair} keyword of the {\hyperref[\detokenize{run_style::doc}]{\sphinxcrossref{\DUrole{doc}{run\_style respa}}}} command.  It does not support the \sphinxstyleemphasis{inner},
\sphinxstyleemphasis{middle}, \sphinxstyleemphasis{outer} keywords.


\subsection{Restrictions}
\label{\detokenize{pair_sph_heatconduction:restrictions}}
\sphinxAtStartPar
This pair style is part of the SPH package.  It is only enabled
if LAMMPS was built with that package.  See the {\hyperref[\detokenize{Build_package::doc}]{\sphinxcrossref{\DUrole{doc}{Build package}}}} page for more info.


\subsection{Related commands}
\label{\detokenize{pair_sph_heatconduction:related-commands}}
\sphinxAtStartPar
{\hyperref[\detokenize{pair_coeff::doc}]{\sphinxcrossref{\DUrole{doc}{pair\_coeff}}}}, pair\_sph/rhosum


\subsection{Default}
\label{\detokenize{pair_sph_heatconduction:default}}
\sphinxAtStartPar
none

\sphinxstepscope

\index{pair\_style sph/idealgas@\spxentry{pair\_style sph/idealgas}}\ignorespaces 

\section{pair\_style sph/idealgas command}
\label{\detokenize{pair_sph_idealgas:pair-style-sph-idealgas-command}}\label{\detokenize{pair_sph_idealgas:index-0}}\label{\detokenize{pair_sph_idealgas::doc}}

\subsection{Syntax}
\label{\detokenize{pair_sph_idealgas:syntax}}
\begin{sphinxVerbatim}[commandchars=\\\{\}]
\PYG{k}{pair\PYGZus{}style}\PYG{+w}{ }\PYG{n}{sph}\PYG{o}{/}\PYG{n}{idealgas}
\end{sphinxVerbatim}


\subsection{Examples}
\label{\detokenize{pair_sph_idealgas:examples}}
\begin{sphinxVerbatim}[commandchars=\\\{\}]
\PYG{k}{pair\PYGZus{}style}\PYG{+w}{ }\PYG{n}{sph}\PYG{o}{/}\PYG{n}{idealgas}
\PYG{k}{pair\PYGZus{}coeff}\PYG{+w}{ }\PYG{o}{*}\PYG{+w}{ }\PYG{o}{*}\PYG{+w}{ }\PYG{l+m}{1.0}\PYG{+w}{ }\PYG{l+m}{2.4}
\end{sphinxVerbatim}


\subsection{Description}
\label{\detokenize{pair_sph_idealgas:description}}
\sphinxAtStartPar
The sph/idealgas style computes pressure forces between particles
according to the ideal gas equation of state:
\begin{equation*}
\begin{split}p = (\gamma - 1) \rho e\end{split}
\end{equation*}
\sphinxAtStartPar
where \(\gamma = 1.4\) is the heat capacity ratio, \(\rho\) is
the local density, and e is the internal energy per unit mass.  This
pair style also computes Monaghan’s artificial viscosity to prevent
particles from interpenetrating {\hyperref[\detokenize{pair_sph_idealgas:ideal-monoghan}]{\sphinxcrossref{\DUrole{std,std-ref}{(Monaghan)}}}}.

\sphinxAtStartPar
See \sphinxhref{PDF/SPH\_LAMMPS\_userguide.pdf}{this PDF guide} to using SPH in
LAMMPS.

\sphinxAtStartPar
The following coefficients must be defined for each pair of atoms
types via the {\hyperref[\detokenize{pair_coeff::doc}]{\sphinxcrossref{\DUrole{doc}{pair\_coeff}}}} command as in the examples
above.
\begin{itemize}
\item {} 
\sphinxAtStartPar
\(\nu\) artificial viscosity (no units)

\item {} 
\sphinxAtStartPar
h kernel function cutoff (distance units)

\end{itemize}


\bigskip\hrule\bigskip



\subsection{Mixing, shift, table, tail correction, restart, rRESPA info}
\label{\detokenize{pair_sph_idealgas:mixing-shift-table-tail-correction-restart-rrespa-info}}
\sphinxAtStartPar
This style does not support mixing.  Thus, coefficients for all
I,J pairs must be specified explicitly.

\sphinxAtStartPar
This style does not support the {\hyperref[\detokenize{pair_modify::doc}]{\sphinxcrossref{\DUrole{doc}{pair\_modify}}}}
shift, table, and tail options.

\sphinxAtStartPar
This style does not write information to {\hyperref[\detokenize{restart::doc}]{\sphinxcrossref{\DUrole{doc}{binary restart files}}}}.  Thus, you need to re\sphinxhyphen{}specify the pair\_style and
pair\_coeff commands in an input script that reads a restart file.

\sphinxAtStartPar
This style can only be used via the \sphinxstyleemphasis{pair} keyword of the {\hyperref[\detokenize{run_style::doc}]{\sphinxcrossref{\DUrole{doc}{run\_style respa}}}} command.  It does not support the \sphinxstyleemphasis{inner},
\sphinxstyleemphasis{middle}, \sphinxstyleemphasis{outer} keywords.


\subsection{Restrictions}
\label{\detokenize{pair_sph_idealgas:restrictions}}
\sphinxAtStartPar
This pair style is part of the SPH package.  It is only enabled
if LAMMPS was built with that package.  See the {\hyperref[\detokenize{Build_package::doc}]{\sphinxcrossref{\DUrole{doc}{Build package}}}} page for more info.


\subsection{Related commands}
\label{\detokenize{pair_sph_idealgas:related-commands}}
\sphinxAtStartPar
{\hyperref[\detokenize{pair_coeff::doc}]{\sphinxcrossref{\DUrole{doc}{pair\_coeff}}}}, pair\_sph/rhosum


\subsection{Default}
\label{\detokenize{pair_sph_idealgas:default}}
\sphinxAtStartPar
none


\bigskip\hrule\bigskip

\phantomsection\label{\detokenize{pair_sph_idealgas:ideal-monoghan}}
\sphinxAtStartPar
\sphinxstylestrong{(Monaghan)} Monaghan and Gingold, Journal of Computational Physics,
52, 374\sphinxhyphen{}389 (1983).

\sphinxstepscope

\index{pair\_style sph/lj@\spxentry{pair\_style sph/lj}}\ignorespaces 
\index{pair\_style sph/lj/gpu@\spxentry{pair\_style sph/lj/gpu}}\ignorespaces 

\section{pair\_style sph/lj command}
\label{\detokenize{pair_sph_lj:pair-style-sph-lj-command}}\label{\detokenize{pair_sph_lj:index-1}}\label{\detokenize{pair_sph_lj:index-0}}\label{\detokenize{pair_sph_lj::doc}}
\sphinxAtStartPar
Accelerator Variants: \sphinxstyleemphasis{sph/lj/gpu}


\subsection{Syntax}
\label{\detokenize{pair_sph_lj:syntax}}
\begin{sphinxVerbatim}[commandchars=\\\{\}]
\PYG{k}{pair\PYGZus{}style}\PYG{+w}{ }\PYG{n}{sph}\PYG{o}{/}\PYG{n}{lj}
\end{sphinxVerbatim}


\subsection{Examples}
\label{\detokenize{pair_sph_lj:examples}}
\begin{sphinxVerbatim}[commandchars=\\\{\}]
\PYG{k}{pair\PYGZus{}style}\PYG{+w}{ }\PYG{n}{sph}\PYG{o}{/}\PYG{n}{lj}
\PYG{k}{pair\PYGZus{}coeff}\PYG{+w}{ }\PYG{o}{*}\PYG{+w}{ }\PYG{o}{*}\PYG{+w}{ }\PYG{l+m}{1.0}\PYG{+w}{ }\PYG{l+m}{2.4}
\end{sphinxVerbatim}


\subsection{Description}
\label{\detokenize{pair_sph_lj:description}}
\sphinxAtStartPar
The sph/lj style computes pressure forces between particles according
to the Lennard\sphinxhyphen{}Jones equation of state, which is computed according to
Ree’s 1980 polynomial fit {\hyperref[\detokenize{pair_sph_lj:ree}]{\sphinxcrossref{\DUrole{std,std-ref}{(Ree)}}}}. The Lennard\sphinxhyphen{}Jones parameters
epsilon and sigma are set to unity.  This pair style also computes
Monaghan’s artificial viscosity to prevent particles from
interpenetrating {\hyperref[\detokenize{pair_sph_lj:monoghan}]{\sphinxcrossref{\DUrole{std,std-ref}{(Monaghan)}}}}.

\sphinxAtStartPar
See \sphinxhref{PDF/SPH\_LAMMPS\_userguide.pdf}{this PDF guide} to using SPH in
LAMMPS.

\sphinxAtStartPar
The following coefficients must be defined for each pair of atoms
types via the {\hyperref[\detokenize{pair_coeff::doc}]{\sphinxcrossref{\DUrole{doc}{pair\_coeff}}}} command as in the examples
above.
\begin{itemize}
\item {} 
\sphinxAtStartPar
\(\nu\) artificial viscosity (no units)

\item {} 
\sphinxAtStartPar
h kernel function cutoff (distance units)

\end{itemize}


\bigskip\hrule\bigskip



\subsection{Mixing, shift, table, tail correction, restart, rRESPA info}
\label{\detokenize{pair_sph_lj:mixing-shift-table-tail-correction-restart-rrespa-info}}
\sphinxAtStartPar
This style does not support mixing.  Thus, coefficients for all
I,J pairs must be specified explicitly.

\sphinxAtStartPar
This style does not support the {\hyperref[\detokenize{pair_modify::doc}]{\sphinxcrossref{\DUrole{doc}{pair\_modify}}}}
shift, table, and tail options.

\sphinxAtStartPar
This style does not write information to {\hyperref[\detokenize{restart::doc}]{\sphinxcrossref{\DUrole{doc}{binary restart files}}}}.  Thus, you need to re\sphinxhyphen{}specify the pair\_style and
pair\_coeff commands in an input script that reads a restart file.

\sphinxAtStartPar
This style can only be used via the \sphinxstyleemphasis{pair} keyword of the {\hyperref[\detokenize{run_style::doc}]{\sphinxcrossref{\DUrole{doc}{run\_style respa}}}} command.  It does not support the \sphinxstyleemphasis{inner},
\sphinxstyleemphasis{middle}, \sphinxstyleemphasis{outer} keywords.


\subsection{Restrictions}
\label{\detokenize{pair_sph_lj:restrictions}}
\sphinxAtStartPar
As noted above, the Lennard\sphinxhyphen{}Jones parameters epsilon and sigma are set
to unity.

\sphinxAtStartPar
This pair style is part of the SPH package.  It is only enabled
if LAMMPS was built with that package.  See the {\hyperref[\detokenize{Build_package::doc}]{\sphinxcrossref{\DUrole{doc}{Build package}}}} page for more info.


\subsection{Related commands}
\label{\detokenize{pair_sph_lj:related-commands}}
\sphinxAtStartPar
{\hyperref[\detokenize{pair_coeff::doc}]{\sphinxcrossref{\DUrole{doc}{pair\_coeff}}}}, pair\_sph/rhosum


\subsection{Default}
\label{\detokenize{pair_sph_lj:default}}
\sphinxAtStartPar
none


\bigskip\hrule\bigskip

\phantomsection\label{\detokenize{pair_sph_lj:ree}}
\sphinxAtStartPar
\sphinxstylestrong{(Ree)} Ree, Journal of Chemical Physics, 73, 5401 (1980).

\phantomsection\label{\detokenize{pair_sph_lj:monoghan}}
\sphinxAtStartPar
\sphinxstylestrong{(Monaghan)} Monaghan and Gingold, Journal of Computational Physics,
52, 374\sphinxhyphen{}389 (1983).

\sphinxstepscope

\index{pair\_style sph/rhosum@\spxentry{pair\_style sph/rhosum}}\ignorespaces 

\section{pair\_style sph/rhosum command}
\label{\detokenize{pair_sph_rhosum:pair-style-sph-rhosum-command}}\label{\detokenize{pair_sph_rhosum:index-0}}\label{\detokenize{pair_sph_rhosum::doc}}

\subsection{Syntax}
\label{\detokenize{pair_sph_rhosum:syntax}}
\begin{sphinxVerbatim}[commandchars=\\\{\}]
\PYG{k}{pair\PYGZus{}style}\PYG{+w}{ }\PYG{n}{sph}\PYG{o}{/}\PYG{n}{rhosum}\PYG{+w}{ }\PYG{n}{Nstep}
\end{sphinxVerbatim}
\begin{itemize}
\item {} 
\sphinxAtStartPar
Nstep = timestep interval

\end{itemize}


\subsection{Examples}
\label{\detokenize{pair_sph_rhosum:examples}}
\begin{sphinxVerbatim}[commandchars=\\\{\}]
\PYG{k}{pair\PYGZus{}style}\PYG{+w}{ }\PYG{n}{sph}\PYG{o}{/}\PYG{n}{rhosum}\PYG{+w}{ }\PYG{l+m}{10}
\PYG{k}{pair\PYGZus{}coeff}\PYG{+w}{ }\PYG{o}{*}\PYG{+w}{ }\PYG{o}{*}\PYG{+w}{ }\PYG{l+m}{2.4}
\end{sphinxVerbatim}


\subsection{Description}
\label{\detokenize{pair_sph_rhosum:description}}
\sphinxAtStartPar
The sph/rhosum style computes the local particle mass density rho for
SPH particles by kernel function interpolation, every Nstep timesteps.

\sphinxAtStartPar
See \sphinxhref{PDF/SPH\_LAMMPS\_userguide.pdf}{this PDF guide} to using SPH in
LAMMPS.

\sphinxAtStartPar
The following coefficients must be defined for each pair of atoms
types via the {\hyperref[\detokenize{pair_coeff::doc}]{\sphinxcrossref{\DUrole{doc}{pair\_coeff}}}} command as in the examples
above.
\begin{itemize}
\item {} 
\sphinxAtStartPar
h (distance units)

\end{itemize}


\bigskip\hrule\bigskip



\subsection{Mixing, shift, table, tail correction, restart, rRESPA info}
\label{\detokenize{pair_sph_rhosum:mixing-shift-table-tail-correction-restart-rrespa-info}}
\sphinxAtStartPar
This style does not support mixing.  Thus, coefficients for all
I,J pairs must be specified explicitly.

\sphinxAtStartPar
This style does not support the {\hyperref[\detokenize{pair_modify::doc}]{\sphinxcrossref{\DUrole{doc}{pair\_modify}}}}
shift, table, and tail options.

\sphinxAtStartPar
This style does not write information to {\hyperref[\detokenize{restart::doc}]{\sphinxcrossref{\DUrole{doc}{binary restart files}}}}.  Thus, you need to re\sphinxhyphen{}specify the pair\_style and
pair\_coeff commands in an input script that reads a restart file.

\sphinxAtStartPar
This style can only be used via the \sphinxstyleemphasis{pair} keyword of the {\hyperref[\detokenize{run_style::doc}]{\sphinxcrossref{\DUrole{doc}{run\_style respa}}}} command.  It does not support the \sphinxstyleemphasis{inner},
\sphinxstyleemphasis{middle}, \sphinxstyleemphasis{outer} keywords.


\subsection{Restrictions}
\label{\detokenize{pair_sph_rhosum:restrictions}}
\sphinxAtStartPar
This pair style is part of the SPH package.  It is only enabled
if LAMMPS was built with that package.  See the {\hyperref[\detokenize{Build_package::doc}]{\sphinxcrossref{\DUrole{doc}{Build package}}}} page for more info.


\subsection{Related commands}
\label{\detokenize{pair_sph_rhosum:related-commands}}
\sphinxAtStartPar
{\hyperref[\detokenize{pair_coeff::doc}]{\sphinxcrossref{\DUrole{doc}{pair\_coeff}}}}, pair\_sph/taitwater


\subsection{Default}
\label{\detokenize{pair_sph_rhosum:default}}
\sphinxAtStartPar
none

\sphinxstepscope

\index{pair\_style sph/taitwater@\spxentry{pair\_style sph/taitwater}}\ignorespaces 
\index{pair\_style sph/taitwater/gpu@\spxentry{pair\_style sph/taitwater/gpu}}\ignorespaces 

\section{pair\_style sph/taitwater command}
\label{\detokenize{pair_sph_taitwater:pair-style-sph-taitwater-command}}\label{\detokenize{pair_sph_taitwater:index-1}}\label{\detokenize{pair_sph_taitwater:index-0}}\label{\detokenize{pair_sph_taitwater::doc}}
\sphinxAtStartPar
Accelerator Variants: \sphinxstyleemphasis{sph/taitwater/gpu}


\subsection{Syntax}
\label{\detokenize{pair_sph_taitwater:syntax}}
\begin{sphinxVerbatim}[commandchars=\\\{\}]
\PYG{k}{pair\PYGZus{}style}\PYG{+w}{ }\PYG{n}{sph}\PYG{o}{/}\PYG{n}{taitwater}
\end{sphinxVerbatim}


\subsection{Examples}
\label{\detokenize{pair_sph_taitwater:examples}}
\begin{sphinxVerbatim}[commandchars=\\\{\}]
\PYG{k}{pair\PYGZus{}style}\PYG{+w}{ }\PYG{n}{sph}\PYG{o}{/}\PYG{n}{taitwater}
\PYG{k}{pair\PYGZus{}coeff}\PYG{+w}{ }\PYG{o}{*}\PYG{+w}{ }\PYG{o}{*}\PYG{+w}{ }\PYG{l+m}{1000.0}\PYG{+w}{ }\PYG{l+m}{1430.0}\PYG{+w}{ }\PYG{l+m}{1.0}\PYG{+w}{ }\PYG{l+m}{2.4}
\end{sphinxVerbatim}


\subsection{Description}
\label{\detokenize{pair_sph_taitwater:description}}
\sphinxAtStartPar
The sph/taitwater style computes pressure forces between SPH particles
according to Tait’s equation of state:
\begin{equation*}
\begin{split}p = B \biggl[\left(\frac{\rho}{\rho_0}\right)^{\gamma} - 1\biggr]\end{split}
\end{equation*}
\sphinxAtStartPar
where \(\gamma = 7\) and \(B = c_0^2 \rho_0 / \gamma\), with
\(\rho_0\) being the reference density and \(c_0\) the reference
speed of sound.

\sphinxAtStartPar
This pair style also computes Monaghan’s artificial viscosity to
prevent particles from interpenetrating {\hyperref[\detokenize{pair_sph_taitwater:monaghan}]{\sphinxcrossref{\DUrole{std,std-ref}{(Monaghan)}}}}.

\sphinxAtStartPar
See \sphinxhref{PDF/SPH\_LAMMPS\_userguide.pdf}{this PDF guide} to using SPH in
LAMMPS.

\sphinxAtStartPar
The following coefficients must be defined for each pair of atoms
types via the {\hyperref[\detokenize{pair_coeff::doc}]{\sphinxcrossref{\DUrole{doc}{pair\_coeff}}}} command as in the examples
above.
\begin{itemize}
\item {} 
\sphinxAtStartPar
\(\rho_0\) reference density (mass/volume units)

\item {} 
\sphinxAtStartPar
\(c_0\) reference soundspeed (distance/time units)

\item {} 
\sphinxAtStartPar
\(\nu\) artificial viscosity (no units)

\item {} 
\sphinxAtStartPar
h kernel function cutoff (distance units)

\end{itemize}


\bigskip\hrule\bigskip



\subsection{Mixing, shift, table, tail correction, restart, rRESPA info}
\label{\detokenize{pair_sph_taitwater:mixing-shift-table-tail-correction-restart-rrespa-info}}
\sphinxAtStartPar
This style does not support mixing.  Thus, coefficients for all
I,J pairs must be specified explicitly.

\sphinxAtStartPar
This style does not support the {\hyperref[\detokenize{pair_modify::doc}]{\sphinxcrossref{\DUrole{doc}{pair\_modify}}}}
shift, table, and tail options.

\sphinxAtStartPar
This style does not write information to {\hyperref[\detokenize{restart::doc}]{\sphinxcrossref{\DUrole{doc}{binary restart files}}}}.  Thus, you need to re\sphinxhyphen{}specify the pair\_style and
pair\_coeff commands in an input script that reads a restart file.

\sphinxAtStartPar
This style can only be used via the \sphinxstyleemphasis{pair} keyword of the {\hyperref[\detokenize{run_style::doc}]{\sphinxcrossref{\DUrole{doc}{run\_style respa}}}} command.  It does not support the \sphinxstyleemphasis{inner},
\sphinxstyleemphasis{middle}, \sphinxstyleemphasis{outer} keywords.


\subsection{Restrictions}
\label{\detokenize{pair_sph_taitwater:restrictions}}
\sphinxAtStartPar
This pair style is part of the SPH package.  It is only enabled
if LAMMPS was built with that package.  See the {\hyperref[\detokenize{Build_package::doc}]{\sphinxcrossref{\DUrole{doc}{Build package}}}} page for more info.


\subsection{Related commands}
\label{\detokenize{pair_sph_taitwater:related-commands}}
\sphinxAtStartPar
{\hyperref[\detokenize{pair_coeff::doc}]{\sphinxcrossref{\DUrole{doc}{pair\_coeff}}}}, pair\_sph/rhosum


\subsection{Default}
\label{\detokenize{pair_sph_taitwater:default}}
\sphinxAtStartPar
none


\bigskip\hrule\bigskip

\phantomsection\label{\detokenize{pair_sph_taitwater:monaghan}}
\sphinxAtStartPar
\sphinxstylestrong{(Monaghan)} Monaghan and Gingold, Journal of Computational Physics,
52, 374\sphinxhyphen{}389 (1983).

\sphinxstepscope

\index{pair\_style sph/taitwater/morris@\spxentry{pair\_style sph/taitwater/morris}}\ignorespaces 

\section{pair\_style sph/taitwater/morris command}
\label{\detokenize{pair_sph_taitwater_morris:pair-style-sph-taitwater-morris-command}}\label{\detokenize{pair_sph_taitwater_morris:index-0}}\label{\detokenize{pair_sph_taitwater_morris::doc}}

\subsection{Syntax}
\label{\detokenize{pair_sph_taitwater_morris:syntax}}
\begin{sphinxVerbatim}[commandchars=\\\{\}]
\PYG{k}{pair\PYGZus{}style}\PYG{+w}{ }\PYG{n}{sph}\PYG{o}{/}\PYG{n}{taitwater}\PYG{o}{/}\PYG{n}{morris}
\end{sphinxVerbatim}


\subsection{Examples}
\label{\detokenize{pair_sph_taitwater_morris:examples}}
\begin{sphinxVerbatim}[commandchars=\\\{\}]
\PYG{k}{pair\PYGZus{}style}\PYG{+w}{ }\PYG{n}{sph}\PYG{o}{/}\PYG{n}{taitwater}\PYG{o}{/}\PYG{n}{morris}
\PYG{k}{pair\PYGZus{}coeff}\PYG{+w}{ }\PYG{o}{*}\PYG{+w}{ }\PYG{o}{*}\PYG{+w}{ }\PYG{l+m}{1000.0}\PYG{+w}{ }\PYG{l+m}{1430.0}\PYG{+w}{ }\PYG{l+m}{1.0}\PYG{+w}{ }\PYG{l+m}{2.4}
\end{sphinxVerbatim}


\subsection{Description}
\label{\detokenize{pair_sph_taitwater_morris:description}}
\sphinxAtStartPar
The sph/taitwater/morris style computes pressure forces between SPH
particles according to Tait’s equation of state:
\begin{equation*}
\begin{split}p = B \biggl[\left(\frac{\rho}{\rho_0}\right)^{\gamma} - 1\biggr]\end{split}
\end{equation*}
\sphinxAtStartPar
where \(\gamma = 7\) and \(B = c_0^2 \rho_0 / \gamma\), with
\(\rho_0\) being the reference density and \(c_0\) the reference
speed of sound.

\sphinxAtStartPar
This pair style also computes laminar viscosity {\hyperref[\detokenize{pair_sph_taitwater_morris:morris}]{\sphinxcrossref{\DUrole{std,std-ref}{(Morris)}}}}.

\sphinxAtStartPar
See \sphinxhref{PDF/SPH\_LAMMPS\_userguide.pdf}{this PDF guide} to using SPH in
LAMMPS.

\sphinxAtStartPar
The following coefficients must be defined for each pair of atoms
types via the {\hyperref[\detokenize{pair_coeff::doc}]{\sphinxcrossref{\DUrole{doc}{pair\_coeff}}}} command as in the examples
above.
\begin{itemize}
\item {} 
\sphinxAtStartPar
\(\rho_0\) reference density (mass/volume units)

\item {} 
\sphinxAtStartPar
\(c_0\) reference soundspeed (distance/time units)

\item {} 
\sphinxAtStartPar
\(\nu\) dynamic viscosity (mass*distance/time units)

\item {} 
\sphinxAtStartPar
h kernel function cutoff (distance units)

\end{itemize}


\bigskip\hrule\bigskip



\subsection{Mixing, shift, table, tail correction, restart, rRESPA info}
\label{\detokenize{pair_sph_taitwater_morris:mixing-shift-table-tail-correction-restart-rrespa-info}}
\sphinxAtStartPar
This style does not support mixing.  Thus, coefficients for all
I,J pairs must be specified explicitly.

\sphinxAtStartPar
This style does not support the {\hyperref[\detokenize{pair_modify::doc}]{\sphinxcrossref{\DUrole{doc}{pair\_modify}}}}
shift, table, and tail options.

\sphinxAtStartPar
This style does not write information to {\hyperref[\detokenize{restart::doc}]{\sphinxcrossref{\DUrole{doc}{binary restart files}}}}.  Thus, you need to re\sphinxhyphen{}specify the pair\_style and
pair\_coeff commands in an input script that reads a restart file.

\sphinxAtStartPar
This style can only be used via the \sphinxstyleemphasis{pair} keyword of the {\hyperref[\detokenize{run_style::doc}]{\sphinxcrossref{\DUrole{doc}{run\_style respa}}}} command.  It does not support the \sphinxstyleemphasis{inner},
\sphinxstyleemphasis{middle}, \sphinxstyleemphasis{outer} keywords.


\subsection{Restrictions}
\label{\detokenize{pair_sph_taitwater_morris:restrictions}}
\sphinxAtStartPar
This pair style is part of the SPH package.  It is only enabled
if LAMMPS was built with that package.  See the {\hyperref[\detokenize{Build_package::doc}]{\sphinxcrossref{\DUrole{doc}{Build package}}}} page for more info.


\subsection{Related commands}
\label{\detokenize{pair_sph_taitwater_morris:related-commands}}
\sphinxAtStartPar
{\hyperref[\detokenize{pair_coeff::doc}]{\sphinxcrossref{\DUrole{doc}{pair\_coeff}}}}, pair\_sph/rhosum


\subsection{Default}
\label{\detokenize{pair_sph_taitwater_morris:default}}
\sphinxAtStartPar
none


\bigskip\hrule\bigskip

\phantomsection\label{\detokenize{pair_sph_taitwater_morris:morris}}
\sphinxAtStartPar
\sphinxstylestrong{(Morris)} Morris, Fox, Zhu, J Comp Physics, 136, 214\sphinxhyphen{}226 (1997).

\sphinxstepscope

\index{pair\_style lj/spica@\spxentry{pair\_style lj/spica}}\ignorespaces 
\index{pair\_style lj/spica/gpu@\spxentry{pair\_style lj/spica/gpu}}\ignorespaces 
\index{pair\_style lj/spica/kk@\spxentry{pair\_style lj/spica/kk}}\ignorespaces 
\index{pair\_style lj/spica/omp@\spxentry{pair\_style lj/spica/omp}}\ignorespaces 
\index{pair\_style lj/spica/coul/long@\spxentry{pair\_style lj/spica/coul/long}}\ignorespaces 
\index{pair\_style lj/spica/coul/long/gpu@\spxentry{pair\_style lj/spica/coul/long/gpu}}\ignorespaces 
\index{pair\_style lj/spica/coul/long/omp@\spxentry{pair\_style lj/spica/coul/long/omp}}\ignorespaces 
\index{pair\_style lj/spica/coul/long/kk@\spxentry{pair\_style lj/spica/coul/long/kk}}\ignorespaces 
\index{pair\_style lj/spica/coul/msm@\spxentry{pair\_style lj/spica/coul/msm}}\ignorespaces 
\index{pair\_style lj/spica/coul/msm/omp@\spxentry{pair\_style lj/spica/coul/msm/omp}}\ignorespaces 

\section{pair\_style lj/spica command}
\label{\detokenize{pair_spica:pair-style-lj-spica-command}}\label{\detokenize{pair_spica:index-9}}\label{\detokenize{pair_spica:index-8}}\label{\detokenize{pair_spica:index-7}}\label{\detokenize{pair_spica:index-6}}\label{\detokenize{pair_spica:index-5}}\label{\detokenize{pair_spica:index-4}}\label{\detokenize{pair_spica:index-3}}\label{\detokenize{pair_spica:index-2}}\label{\detokenize{pair_spica:index-1}}\label{\detokenize{pair_spica:index-0}}\label{\detokenize{pair_spica::doc}}
\sphinxAtStartPar
Accelerator Variants: \sphinxstyleemphasis{lj/spica/gpu}, \sphinxstyleemphasis{lj/spica/kk}, \sphinxstyleemphasis{lj/spica/omp}


\section{pair\_style lj/spica/coul/long command}
\label{\detokenize{pair_spica:pair-style-lj-spica-coul-long-command}}
\sphinxAtStartPar
Accelerator Variants: \sphinxstyleemphasis{lj/spica/coul/long/gpu}, \sphinxstyleemphasis{lj/spica/coul/long/omp}, \sphinxstyleemphasis{lj/spica/coul/long/kk}


\section{pair\_style lj/spica/coul/msm command}
\label{\detokenize{pair_spica:pair-style-lj-spica-coul-msm-command}}
\sphinxAtStartPar
Accelerator Variants: \sphinxstyleemphasis{lj/spica/coul/msm/omp}


\subsection{Syntax}
\label{\detokenize{pair_spica:syntax}}
\begin{sphinxVerbatim}[commandchars=\\\{\}]
\PYG{k}{pair\PYGZus{}style}\PYG{+w}{ }\PYG{n}{style}\PYG{+w}{ }\PYG{n}{args}
\end{sphinxVerbatim}
\begin{itemize}
\item {} 
\sphinxAtStartPar
style = \sphinxstyleemphasis{lj/spica} or \sphinxstyleemphasis{lj/spica/coul/long}

\item {} 
\sphinxAtStartPar
args = list of arguments for a particular style

\end{itemize}
\begin{sphinxalltt}
\sphinxstyleemphasis{lj/spica} args = cutoff
  cutoff = global cutoff for Lennard Jones interactions (distance units)
\sphinxstyleemphasis{lj/spica/coul/long} args = cutoff (cutoff2)
  cutoff = global cutoff for LJ (and Coulombic if only 1 arg) (distance units)
  cutoff2 = global cutoff for Coulombic (optional) (distance units)
\end{sphinxalltt}


\subsection{Examples}
\label{\detokenize{pair_spica:examples}}
\begin{sphinxVerbatim}[commandchars=\\\{\}]
\PYG{k}{pair\PYGZus{}style}\PYG{+w}{ }\PYG{n}{lj}\PYG{o}{/}\PYG{n}{spica}\PYG{+w}{ }\PYG{l+m}{2.5}
\PYG{k}{pair\PYGZus{}coeff}\PYG{+w}{ }\PYG{l+m}{1}\PYG{+w}{ }\PYG{l+m}{1}\PYG{+w}{ }\PYG{n}{lj12\PYGZus{}6}\PYG{+w}{ }\PYG{l+m}{1}\PYG{+w}{ }\PYG{l+m}{1.1}\PYG{+w}{ }\PYG{l+m}{2.8}

\PYG{k}{pair\PYGZus{}style}\PYG{+w}{ }\PYG{n}{lj}\PYG{o}{/}\PYG{n}{spica}\PYG{o}{/}\PYG{n}{coul}\PYG{o}{/}\PYG{n}{long}\PYG{+w}{ }\PYG{l+m}{10.0}
\PYG{k}{pair\PYGZus{}style}\PYG{+w}{ }\PYG{n}{lj}\PYG{o}{/}\PYG{n}{spica}\PYG{o}{/}\PYG{n}{coul}\PYG{o}{/}\PYG{n}{long}\PYG{+w}{ }\PYG{l+m}{10.0}\PYG{+w}{ }\PYG{l+m}{12.0}
\PYG{k}{pair\PYGZus{}coeff}\PYG{+w}{ }\PYG{l+m}{1}\PYG{+w}{ }\PYG{l+m}{1}\PYG{+w}{ }\PYG{n}{lj9\PYGZus{}6}\PYG{+w}{ }\PYG{l+m}{100.0}\PYG{+w}{ }\PYG{l+m}{3.5}\PYG{+w}{ }\PYG{l+m}{12.0}

\PYG{k}{pair\PYGZus{}style}\PYG{+w}{ }\PYG{n}{lj}\PYG{o}{/}\PYG{n}{spica}\PYG{o}{/}\PYG{n}{coul}\PYG{o}{/}\PYG{n}{msm}\PYG{+w}{ }\PYG{l+m}{10.0}
\PYG{k}{pair\PYGZus{}style}\PYG{+w}{ }\PYG{n}{lj}\PYG{o}{/}\PYG{n}{spica}\PYG{o}{/}\PYG{n}{coul}\PYG{o}{/}\PYG{n}{msm}\PYG{+w}{ }\PYG{l+m}{10.0}\PYG{+w}{ }\PYG{l+m}{12.0}
\PYG{k}{pair\PYGZus{}coeff}\PYG{+w}{ }\PYG{l+m}{1}\PYG{+w}{ }\PYG{l+m}{1}\PYG{+w}{ }\PYG{n}{lj9\PYGZus{}6}\PYG{+w}{ }\PYG{l+m}{100.0}\PYG{+w}{ }\PYG{l+m}{3.5}\PYG{+w}{ }\PYG{l+m}{12.0}
\end{sphinxVerbatim}


\subsection{Description}
\label{\detokenize{pair_spica:description}}
\sphinxAtStartPar
The \sphinxstyleemphasis{lj/spica} styles compute a 9/6, 12/4, 12/5, or 12/6 Lennard\sphinxhyphen{}Jones potential,
given by
\begin{equation*}
\begin{split}E = & \frac{27}{4} \epsilon \left[ \left(\frac{\sigma}{r}\right)^{9} -
                      \left(\frac{\sigma}{r}\right)^6 \right]
                      \qquad r < r_c \\
E = & \frac{3\sqrt{3}}{2} \epsilon \left[ \left(\frac{\sigma}{r}\right)^{12} -
                      \left(\frac{\sigma}{r}\right)^4 \right]
                      \qquad r < r_c \\
E = & \frac{12}{7}\left(\frac{12}{5}\right)^{\left(\frac{5}{7}\right)} \epsilon
                      \left[ \left(\frac{\sigma}{r}\right)^{12} -
                      \left(\frac{\sigma}{r}\right)^5 \right]
                      \qquad r < r_c \\
E = &  4 \epsilon  \left[ \left(\frac{\sigma}{r}\right)^{12} -
                      \left(\frac{\sigma}{r}\right)^6 \right]
                      \qquad r < r_c\end{split}
\end{equation*}
\sphinxAtStartPar
as required for the SPICA (formerly called SDK) and the pSPICA Coarse\sphinxhyphen{}grained MD parameterization discussed in
{\hyperref[\detokenize{pair_spica:shinoda3}]{\sphinxcrossref{\DUrole{std,std-ref}{(Shinoda)}}}}, {\hyperref[\detokenize{pair_spica:devane}]{\sphinxcrossref{\DUrole{std,std-ref}{(DeVane)}}}}, {\hyperref[\detokenize{pair_spica:seo}]{\sphinxcrossref{\DUrole{std,std-ref}{(Seo)}}}}, and {\hyperref[\detokenize{pair_spica:miyazaki}]{\sphinxcrossref{\DUrole{std,std-ref}{(Miyazaki)}}}}.
\(r_c\) is the cutoff.
Summary information on these force fields can be found at \sphinxurl{https://www.spica-ff.org}

\sphinxAtStartPar
Style \sphinxstyleemphasis{lj/spica/coul/long} computes the adds Coulombic interactions
with an additional damping factor applied so it can be used in
conjunction with the {\hyperref[\detokenize{kspace_style::doc}]{\sphinxcrossref{\DUrole{doc}{kspace\_style}}}} command and
its \sphinxstyleemphasis{ewald} or \sphinxstyleemphasis{pppm} or \sphinxstyleemphasis{pppm/cg} option.  The Coulombic cutoff
specified for this style means that pairwise interactions within
this distance are computed directly; interactions outside that
distance are computed in reciprocal space.

\sphinxAtStartPar
The following coefficients must be defined for each pair of atoms
types via the {\hyperref[\detokenize{pair_coeff::doc}]{\sphinxcrossref{\DUrole{doc}{pair\_coeff}}}} command as in the examples
above, or in the data file or restart files read by the
{\hyperref[\detokenize{read_data::doc}]{\sphinxcrossref{\DUrole{doc}{read\_data}}}} or {\hyperref[\detokenize{read_restart::doc}]{\sphinxcrossref{\DUrole{doc}{read\_restart}}}}
commands, or by mixing as described below:
\begin{itemize}
\item {} 
\sphinxAtStartPar
cg\_type (lj9\_6, lj12\_4, lj12\_5, or lj12\_6)

\item {} 
\sphinxAtStartPar
epsilon (energy units)

\item {} 
\sphinxAtStartPar
sigma (distance units)

\item {} 
\sphinxAtStartPar
cutoff1 (distance units)

\end{itemize}

\sphinxAtStartPar
Note that sigma is defined in the LJ formula as the zero\sphinxhyphen{}crossing
distance for the potential, not as the energy minimum. The prefactors
are chosen so that the potential minimum is at \sphinxhyphen{}epsilon.

\sphinxAtStartPar
The latter 2 coefficients are optional.  If not specified, the global
LJ and Coulombic cutoffs specified in the pair\_style command are used.
If only one cutoff is specified, it is used as the cutoff for both LJ
and Coulombic interactions for this type pair.  If both coefficients
are specified, they are used as the LJ and Coulombic cutoffs for this
type pair.

\sphinxAtStartPar
For \sphinxstyleemphasis{lj/spica/coul/long} and \sphinxstyleemphasis{lj/spica/coul/msm} only the LJ cutoff can be
specified since a Coulombic cutoff cannot be specified for an
individual I,J type pair.  All type pairs use the same global
Coulombic cutoff specified in the pair\_style command.

\sphinxAtStartPar
The original implementation of the above styles are
style \sphinxstyleemphasis{lj/sdk}, \sphinxstyleemphasis{lj/sdk/coul/long}, and \sphinxstyleemphasis{lj/sdk/coul/msm},
and available for backward compatibility.


\bigskip\hrule\bigskip


\sphinxAtStartPar
Styles with a \sphinxstyleemphasis{gpu}, \sphinxstyleemphasis{intel}, \sphinxstyleemphasis{kk}, \sphinxstyleemphasis{omp}, or \sphinxstyleemphasis{opt} suffix are
functionally the same as the corresponding style without the suffix.
They have been optimized to run faster, depending on your available
hardware, as discussed on the {\hyperref[\detokenize{Speed_packages::doc}]{\sphinxcrossref{\DUrole{doc}{Accelerator packages}}}}
page.  The accelerated styles take the same arguments and should
produce the same results, except for round\sphinxhyphen{}off and precision issues.

\sphinxAtStartPar
These accelerated styles are part of the GPU, INTEL, KOKKOS,
OPENMP, and OPT packages, respectively.  They are only enabled if
LAMMPS was built with those packages.  See the {\hyperref[\detokenize{Build_package::doc}]{\sphinxcrossref{\DUrole{doc}{Build package}}}} page for more info.

\sphinxAtStartPar
You can specify the accelerated styles explicitly in your input script
by including their suffix, or you can use the {\hyperref[\detokenize{Run_options::doc}]{\sphinxcrossref{\DUrole{doc}{\sphinxhyphen{}suffix command\sphinxhyphen{}line switch}}}} when you invoke LAMMPS, or you can use the
{\hyperref[\detokenize{suffix::doc}]{\sphinxcrossref{\DUrole{doc}{suffix}}}} command in your input script.

\sphinxAtStartPar
See the {\hyperref[\detokenize{Speed_packages::doc}]{\sphinxcrossref{\DUrole{doc}{Accelerator packages}}}} page for more
instructions on how to use the accelerated styles effectively.


\bigskip\hrule\bigskip



\subsection{Mixing, shift, table, tail correction, restart, rRESPA info}
\label{\detokenize{pair_spica:mixing-shift-table-tail-correction-restart-rrespa-info}}
\sphinxAtStartPar
For atom type pairs I,J and I != J, the epsilon and sigma coefficients
and cutoff distance for all of the lj/spica pair styles \sphinxstyleemphasis{cannot} be mixed,
since different pairs may have different exponents. So all parameters
for all pairs have to be specified explicitly through the “pair\_coeff”
command. Defining then in a data file is also not supported, due to
limitations of that file format.

\sphinxAtStartPar
All of the lj/spica pair styles support the
{\hyperref[\detokenize{pair_modify::doc}]{\sphinxcrossref{\DUrole{doc}{pair\_modify}}}} shift option for the energy of the
Lennard\sphinxhyphen{}Jones portion of the pair interaction.

\sphinxAtStartPar
The \sphinxstyleemphasis{lj/spica/coul/long} pair styles support the
{\hyperref[\detokenize{pair_modify::doc}]{\sphinxcrossref{\DUrole{doc}{pair\_modify}}}} table option since they can tabulate
the short\sphinxhyphen{}range portion of the long\sphinxhyphen{}range Coulombic interaction.

\sphinxAtStartPar
All of the lj/spica pair styles write their information to {\hyperref[\detokenize{restart::doc}]{\sphinxcrossref{\DUrole{doc}{binary restart files}}}}, so pair\_style and pair\_coeff commands do
not need to be specified in an input script that reads a restart file.

\sphinxAtStartPar
The lj/spica and lj/cut/coul/long pair styles do not support
the use of the \sphinxstyleemphasis{inner}, \sphinxstyleemphasis{middle}, and \sphinxstyleemphasis{outer} keywords of the {\hyperref[\detokenize{run_style::doc}]{\sphinxcrossref{\DUrole{doc}{run\_style respa}}}} command.


\bigskip\hrule\bigskip



\subsection{Restrictions}
\label{\detokenize{pair_spica:restrictions}}
\sphinxAtStartPar
All of the lj/spica pair styles are part of the CG\sphinxhyphen{}SPICA package.  The
\sphinxstyleemphasis{lj/spica/coul/long} style also requires the KSPACE package to be built
(which is enabled by default).  They are only enabled if LAMMPS was
built with that package.  See the {\hyperref[\detokenize{Build_package::doc}]{\sphinxcrossref{\DUrole{doc}{Build package}}}}
doc page for more info.


\subsection{Related commands}
\label{\detokenize{pair_spica:related-commands}}
\sphinxAtStartPar
{\hyperref[\detokenize{pair_coeff::doc}]{\sphinxcrossref{\DUrole{doc}{pair\_coeff}}}}, {\hyperref[\detokenize{angle_spica::doc}]{\sphinxcrossref{\DUrole{doc}{angle\_style spica}}}}


\subsection{Default}
\label{\detokenize{pair_spica:default}}
\sphinxAtStartPar
none


\bigskip\hrule\bigskip

\phantomsection\label{\detokenize{pair_spica:shinoda3}}
\sphinxAtStartPar
\sphinxstylestrong{(Shinoda)} Shinoda, DeVane, Klein, Mol Sim, 33, 27\sphinxhyphen{}36 (2007).

\phantomsection\label{\detokenize{pair_spica:devane}}
\sphinxAtStartPar
\sphinxstylestrong{(DeVane)}  Shinoda, DeVane, Klein, Soft Matter, 4, 2453\sphinxhyphen{}2462 (2008).

\phantomsection\label{\detokenize{pair_spica:seo}}
\sphinxAtStartPar
\sphinxstylestrong{(Seo)}  Seo, Shinoda, J Chem Theory Comput, 15, 762\sphinxhyphen{}774 (2019).

\phantomsection\label{\detokenize{pair_spica:miyazaki}}
\sphinxAtStartPar
\sphinxstylestrong{(Miyazaki)}  Miyazaki, Okazaki, Shinoda, J Chem Theory Comput, 16, 782\sphinxhyphen{}793 (2020).

\sphinxstepscope

\index{pair\_style spin/dipole/cut@\spxentry{pair\_style spin/dipole/cut}}\ignorespaces 
\index{pair\_style spin/dipole/long@\spxentry{pair\_style spin/dipole/long}}\ignorespaces 

\section{pair\_style spin/dipole/cut command}
\label{\detokenize{pair_spin_dipole:pair-style-spin-dipole-cut-command}}\label{\detokenize{pair_spin_dipole:index-1}}\label{\detokenize{pair_spin_dipole:index-0}}\label{\detokenize{pair_spin_dipole::doc}}

\section{pair\_style spin/dipole/long command}
\label{\detokenize{pair_spin_dipole:pair-style-spin-dipole-long-command}}

\subsection{Syntax}
\label{\detokenize{pair_spin_dipole:syntax}}
\begin{sphinxVerbatim}[commandchars=\\\{\}]
\PYG{k}{pair\PYGZus{}style}\PYG{+w}{ }\PYG{n}{spin}\PYG{o}{/}\PYG{n}{dipole}\PYG{o}{/}\PYG{n}{cut}\PYG{+w}{ }\PYG{n}{cutoff}
\PYG{k}{pair\PYGZus{}style}\PYG{+w}{ }\PYG{n}{spin}\PYG{o}{/}\PYG{n}{dipole}\PYG{o}{/}\PYG{n}{long}\PYG{+w}{ }\PYG{n}{cutoff}
\end{sphinxVerbatim}
\begin{itemize}
\item {} 
\sphinxAtStartPar
cutoff = global cutoff for magnetic dipole energy and forces
(optional) (distance units)

\end{itemize}


\subsection{Examples}
\label{\detokenize{pair_spin_dipole:examples}}
\begin{sphinxVerbatim}[commandchars=\\\{\}]
\PYG{k}{pair\PYGZus{}style}\PYG{+w}{ }\PYG{n}{spin}\PYG{o}{/}\PYG{n}{dipole}\PYG{o}{/}\PYG{n}{cut}\PYG{+w}{ }\PYG{l+m}{10.0}
\PYG{k}{pair\PYGZus{}coeff}\PYG{+w}{ }\PYG{o}{*}\PYG{+w}{ }\PYG{o}{*}\PYG{+w}{ }\PYG{l+m}{10.0}
\PYG{k}{pair\PYGZus{}coeff}\PYG{+w}{ }\PYG{l+m}{2}\PYG{+w}{ }\PYG{l+m}{3}\PYG{+w}{ }\PYG{l+m}{8.0}

\PYG{k}{pair\PYGZus{}style}\PYG{+w}{ }\PYG{n}{spin}\PYG{o}{/}\PYG{n}{dipole}\PYG{o}{/}\PYG{n}{long}\PYG{+w}{ }\PYG{l+m}{9.0}
\PYG{k}{pair\PYGZus{}coeff}\PYG{+w}{ }\PYG{o}{*}\PYG{+w}{ }\PYG{o}{*}\PYG{+w}{ }\PYG{l+m}{10.0}
\PYG{k}{pair\PYGZus{}coeff}\PYG{+w}{ }\PYG{l+m}{2}\PYG{+w}{ }\PYG{l+m}{3}\PYG{+w}{ }\PYG{l+m}{6.0}
\end{sphinxVerbatim}


\subsection{Description}
\label{\detokenize{pair_spin_dipole:description}}
\sphinxAtStartPar
Style \sphinxstyleemphasis{spin/dipole/cut} computes a short\sphinxhyphen{}range dipole\sphinxhyphen{}dipole
interaction between pairs of magnetic particles that each
have a magnetic spin.
The magnetic dipole\sphinxhyphen{}dipole interactions are computed by the
following formulas for the magnetic energy, magnetic precession
vector omega and mechanical force between particles I and J.
\begin{equation*}
\begin{split}\mathcal{H}_{\rm long} & =
-\frac{\mu_{0} \left( \mu_B\right)^2}{4\pi}
\sum_{i,j,i\neq j}^{N}
 \frac{g_i g_j}{r_{ij}^3}
 \biggl(3
 \left(\vec{e}_{ij}\cdot \vec{s}_{i}\right)
 \left(\vec{e}_{ij}\cdot \vec{s}_{j}\right)
 -\vec{s}_i\cdot\vec{s}_j \biggr) \\
 \mathbf{\omega}_i & =
 \frac{\mu_0 (\mu_B)^2}{4\pi\hbar}\sum_{j}
 \frac{g_i g_j}{r_{ij}^3}
 \, \biggl(
 3\,(\vec{e}_{ij}\cdot\vec{s}_{j})\vec{e}_{ij}
 -\vec{s}_{j} \biggr) \\
 \mathbf{F}_i & =
 \frac{3\, \mu_0 (\mu_B)^2}{4\pi} \sum_j
 \frac{g_i g_j}{r_{ij}^4}
 \biggl[\bigl( (\vec{s}_i\cdot\vec{s}_j)
 -5(\vec{e}_{ij}\cdot\vec{s}_i)
 (\vec{e}_{ij}\cdot\vec{s}_j)\bigr) \vec{e}_{ij}+
 \bigl(
 (\vec{e}_{ij}\cdot\vec{s}_i)\vec{s}_j+
 (\vec{e}_{ij}\cdot\vec{s}_j)\vec{s}_i
 \bigr)
 \biggr]\end{split}
\end{equation*}
\sphinxAtStartPar
where \(\vec{s}_i\) and \(\vec{s}_j\) are the spin on two magnetic
particles, r is their separation distance, and the vector \(\vec{e}_{ij}
= \frac{r_i - r_j}{\left| r_i - r_j \right|}\) is the direction vector
between the two particles.

\sphinxAtStartPar
Style \sphinxstyleemphasis{spin/dipole/long} computes long\sphinxhyphen{}range magnetic dipole\sphinxhyphen{}dipole
interaction.
A {\hyperref[\detokenize{kspace_style::doc}]{\sphinxcrossref{\DUrole{doc}{kspace\_style}}}} must be defined to
use this pair style.  Currently, {\hyperref[\detokenize{kspace_style::doc}]{\sphinxcrossref{\DUrole{doc}{kspace\_style ewald/dipole/spin}}}} and {\hyperref[\detokenize{kspace_style::doc}]{\sphinxcrossref{\DUrole{doc}{kspace\_style pppm/dipole/spin}}}} support long\sphinxhyphen{}range magnetic
dipole\sphinxhyphen{}dipole interactions.


\bigskip\hrule\bigskip


\sphinxAtStartPar
The {\hyperref[\detokenize{pair_modify::doc}]{\sphinxcrossref{\DUrole{doc}{pair\_modify}}}} table option is not relevant
for this pair style.

\sphinxAtStartPar
This pair style does not support the {\hyperref[\detokenize{pair_modify::doc}]{\sphinxcrossref{\DUrole{doc}{pair\_modify}}}}
tail option for adding long\sphinxhyphen{}range tail corrections to energy and
pressure.

\sphinxAtStartPar
This pair style writes its information to {\hyperref[\detokenize{restart::doc}]{\sphinxcrossref{\DUrole{doc}{binary restart files}}}}, so pair\_style and pair\_coeff commands do not need
to be specified in an input script that reads a restart file.


\subsection{Restrictions}
\label{\detokenize{pair_spin_dipole:restrictions}}
\sphinxAtStartPar
The \sphinxstyleemphasis{spin/dipole/cut} and \sphinxstyleemphasis{spin/dipole/long} styles are part of
the SPIN package.  They are only enabled if LAMMPS was built with that
package.  See the {\hyperref[\detokenize{Build_package::doc}]{\sphinxcrossref{\DUrole{doc}{Build package}}}} page for more
info.

\sphinxAtStartPar
Using dipole/spin pair styles with \sphinxstyleemphasis{electron} {\hyperref[\detokenize{units::doc}]{\sphinxcrossref{\DUrole{doc}{units}}}} is not
currently supported.


\subsection{Related commands}
\label{\detokenize{pair_spin_dipole:related-commands}}
\sphinxAtStartPar
{\hyperref[\detokenize{pair_coeff::doc}]{\sphinxcrossref{\DUrole{doc}{pair\_coeff}}}}, {\hyperref[\detokenize{kspace_style::doc}]{\sphinxcrossref{\DUrole{doc}{kspace\_style}}}}
{\hyperref[\detokenize{fix_nve_spin::doc}]{\sphinxcrossref{\DUrole{doc}{fix nve/spin}}}}


\subsection{Default}
\label{\detokenize{pair_spin_dipole:default}}
\sphinxAtStartPar
none

\sphinxstepscope

\index{pair\_style spin/dmi@\spxentry{pair\_style spin/dmi}}\ignorespaces 

\section{pair\_style spin/dmi command}
\label{\detokenize{pair_spin_dmi:pair-style-spin-dmi-command}}\label{\detokenize{pair_spin_dmi:index-0}}\label{\detokenize{pair_spin_dmi::doc}}

\subsection{Syntax}
\label{\detokenize{pair_spin_dmi:syntax}}
\begin{sphinxVerbatim}[commandchars=\\\{\}]
\PYG{k}{pair\PYGZus{}style}\PYG{+w}{ }\PYG{n}{spin}\PYG{o}{/}\PYG{n}{dmi}\PYG{+w}{ }\PYG{n}{cutoff}
\end{sphinxVerbatim}
\begin{itemize}
\item {} 
\sphinxAtStartPar
cutoff = global cutoff pair (distance in metal units)

\end{itemize}


\subsection{Examples}
\label{\detokenize{pair_spin_dmi:examples}}
\begin{sphinxVerbatim}[commandchars=\\\{\}]
\PYG{k}{pair\PYGZus{}style}\PYG{+w}{ }\PYG{n}{spin}\PYG{o}{/}\PYG{n}{dmi}\PYG{+w}{ }\PYG{l+m}{4.0}
\PYG{k}{pair\PYGZus{}coeff}\PYG{+w}{ }\PYG{o}{*}\PYG{+w}{ }\PYG{o}{*}\PYG{+w}{ }\PYG{n}{dmi}\PYG{+w}{ }\PYG{l+m}{2.6}\PYG{+w}{ }\PYG{l+m}{0.001}\PYG{+w}{ }\PYG{l+m}{1.0}\PYG{+w}{ }\PYG{l+m}{0.0}\PYG{+w}{ }\PYG{l+m}{0.0}
\PYG{k}{pair\PYGZus{}coeff}\PYG{+w}{ }\PYG{l+m}{1}\PYG{+w}{ }\PYG{l+m}{2}\PYG{+w}{ }\PYG{n}{dmi}\PYG{+w}{ }\PYG{l+m}{4.0}\PYG{+w}{ }\PYG{l+m}{0.00109}\PYG{+w}{ }\PYG{l+m}{0.0}\PYG{+w}{ }\PYG{l+m}{0.0}\PYG{+w}{ }\PYG{l+m}{1.0}
\end{sphinxVerbatim}


\subsection{Description}
\label{\detokenize{pair_spin_dmi:description}}
\sphinxAtStartPar
Style \sphinxstyleemphasis{spin/dmi} computes the Dzyaloshinskii\sphinxhyphen{}Moriya (DM) interaction
between pairs of magnetic spins.
According to the expression reported in {\hyperref[\detokenize{pair_spin_dmi:rohart}]{\sphinxcrossref{\DUrole{std,std-ref}{(Rohart)}}}}, one has
the following DM energy:
\begin{equation*}
\begin{split}\mathbf{H}_{dm} = \sum_{{ i,j}=1,i\neq j}^{N}
\left( \vec{e}_{ij} \times \vec{D} \right)
\cdot\left(\vec{s}_{i}\times \vec{s}_{j}\right),\end{split}
\end{equation*}
\sphinxAtStartPar
where \(\vec{s}_i\) and \(\vec{s}_j\) are two neighboring magnetic spins of
two particles, \(\vec{e}_ij = \frac{r_i - r_j}{\left| r_i - r_j \right|}\)
is the unit vector between sites \sphinxstyleemphasis{i} and \sphinxstyleemphasis{j}, and \(\vec{D}\) is the
DM vector defining the intensity (in eV) and the direction of the
interaction.

\sphinxAtStartPar
In {\hyperref[\detokenize{pair_spin_dmi:rohart}]{\sphinxcrossref{\DUrole{std,std-ref}{(Rohart)}}}}, \(\vec{D}\) is defined as the direction normal to the film oriented
from the high spin\sphinxhyphen{}orbit layer to the magnetic ultra\sphinxhyphen{}thin film.

\sphinxAtStartPar
The application of a spin\sphinxhyphen{}lattice Poisson bracket to this energy (as described
in {\hyperref[\detokenize{pair_spin_dmi:tranchida5}]{\sphinxcrossref{\DUrole{std,std-ref}{(Tranchida)}}}}) allows to derive a magnetic torque omega, and a
mechanical force F (for spin\sphinxhyphen{}lattice calculations only) for each magnetic
particle i:
\begin{equation*}
\begin{split}\vec{\omega}_i = -\frac{1}{\hbar} \sum_{j}^{Neighb} \vec{s}_{j}\times \left(\vec{e}_{ij}\times \vec{D} \right)
~~{\rm and}~~
\vec{F}_i = -\sum_{j}^{Neighb} \frac{1}{r_{ij}} \vec{D} \times \left( \vec{s}_{i}\times \vec{s}_{j} \right)\end{split}
\end{equation*}
\sphinxAtStartPar
More details about the derivation of these torques/forces are reported in
{\hyperref[\detokenize{pair_spin_dmi:tranchida5}]{\sphinxcrossref{\DUrole{std,std-ref}{(Tranchida)}}}}.

\sphinxAtStartPar
For the \sphinxstyleemphasis{spin/dmi} pair style, the following coefficients must be defined for
each pair of atoms types via the {\hyperref[\detokenize{pair_coeff::doc}]{\sphinxcrossref{\DUrole{doc}{pair\_coeff}}}} command as in
the examples above, or in the data file or restart files read by the
{\hyperref[\detokenize{read_data::doc}]{\sphinxcrossref{\DUrole{doc}{read\_data}}}} or {\hyperref[\detokenize{read_restart::doc}]{\sphinxcrossref{\DUrole{doc}{read\_restart}}}} commands, and
set in the following order:
\begin{itemize}
\item {} 
\sphinxAtStartPar
rc (distance units)

\item {} 
\sphinxAtStartPar
|D| (energy units)

\item {} 
\sphinxAtStartPar
Dx, Dy, Dz  (direction of D)

\end{itemize}

\sphinxAtStartPar
Note that rc is the radius cutoff of the considered DM interaction, |D| is
the norm of the DM vector (in eV), and Dx, Dy and Dz define its direction.

\sphinxAtStartPar
None of those coefficients is optional.  If not specified, the \sphinxstyleemphasis{spin/dmi}
pair style cannot be used.


\bigskip\hrule\bigskip



\subsection{Restrictions}
\label{\detokenize{pair_spin_dmi:restrictions}}
\sphinxAtStartPar
All the \sphinxstyleemphasis{pair/spin} styles are part of the SPIN package.  These styles
are only enabled if LAMMPS was built with this package, and if the
atom\_style “spin” was declared.  See the {\hyperref[\detokenize{Build_package::doc}]{\sphinxcrossref{\DUrole{doc}{Build package}}}} page for more info.


\subsection{Related commands}
\label{\detokenize{pair_spin_dmi:related-commands}}
\sphinxAtStartPar
{\hyperref[\detokenize{atom_style::doc}]{\sphinxcrossref{\DUrole{doc}{atom\_style spin}}}}, {\hyperref[\detokenize{pair_coeff::doc}]{\sphinxcrossref{\DUrole{doc}{pair\_coeff}}}},
{\hyperref[\detokenize{pair_eam::doc}]{\sphinxcrossref{\DUrole{doc}{pair\_eam}}}},


\subsection{Default}
\label{\detokenize{pair_spin_dmi:default}}
\sphinxAtStartPar
none


\bigskip\hrule\bigskip

\phantomsection\label{\detokenize{pair_spin_dmi:rohart}}\phantomsection\label{\detokenize{pair_spin_dmi:tranchida5}}
\sphinxAtStartPar
\sphinxstylestrong{(Rohart)} Rohart and Thiaville, Physical Review B, 88(18), 184422. (2013).

\sphinxAtStartPar
\sphinxstylestrong{(Tranchida)} Tranchida, Plimpton, Thibaudeau and Thompson, Journal of Computational Physics, 372, 406\sphinxhyphen{}425, (2018).

\sphinxstepscope

\index{pair\_style spin/exchange@\spxentry{pair\_style spin/exchange}}\ignorespaces 
\index{pair\_style spin/exchange/biquadratic@\spxentry{pair\_style spin/exchange/biquadratic}}\ignorespaces 

\section{pair\_style spin/exchange command}
\label{\detokenize{pair_spin_exchange:pair-style-spin-exchange-command}}\label{\detokenize{pair_spin_exchange:index-1}}\label{\detokenize{pair_spin_exchange:index-0}}\label{\detokenize{pair_spin_exchange::doc}}

\section{pair\_style spin/exchange/biquadratic command}
\label{\detokenize{pair_spin_exchange:pair-style-spin-exchange-biquadratic-command}}

\subsection{Syntax}
\label{\detokenize{pair_spin_exchange:syntax}}
\begin{sphinxVerbatim}[commandchars=\\\{\}]
\PYG{k}{pair\PYGZus{}style}\PYG{+w}{ }\PYG{n}{spin}\PYG{o}{/}\PYG{n}{exchange}\PYG{+w}{ }\PYG{n}{cutoff}
\PYG{k}{pair\PYGZus{}style}\PYG{+w}{ }\PYG{n}{spin}\PYG{o}{/}\PYG{n}{exchange}\PYG{o}{/}\PYG{n}{biquadratic}\PYG{+w}{ }\PYG{n}{cutoff}
\end{sphinxVerbatim}
\begin{itemize}
\item {} 
\sphinxAtStartPar
cutoff = global cutoff pair (distance in metal units)

\end{itemize}


\subsection{Examples}
\label{\detokenize{pair_spin_exchange:examples}}
\begin{sphinxVerbatim}[commandchars=\\\{\}]
\PYG{k}{pair\PYGZus{}style}\PYG{+w}{ }\PYG{n}{spin}\PYG{o}{/}\PYG{n}{exchange}\PYG{+w}{ }\PYG{l+m}{4.0}
\PYG{k}{pair\PYGZus{}coeff}\PYG{+w}{ }\PYG{o}{*}\PYG{+w}{ }\PYG{o}{*}\PYG{+w}{ }\PYG{n}{exchange}\PYG{+w}{ }\PYG{l+m}{4.0}\PYG{+w}{ }\PYG{l+m}{0.0446928}\PYG{+w}{ }\PYG{l+m}{0.003496}\PYG{+w}{ }\PYG{l+m}{1.4885}
\PYG{k}{pair\PYGZus{}coeff}\PYG{+w}{ }\PYG{l+m}{1}\PYG{+w}{ }\PYG{l+m}{2}\PYG{+w}{ }\PYG{n}{exchange}\PYG{+w}{ }\PYG{l+m}{6.0}\PYG{+w}{ }\PYG{o}{\PYGZhy{}}\PYG{l+m}{0.01575}\PYG{+w}{ }\PYG{l+m}{0.0}\PYG{+w}{ }\PYG{l+m}{1.965}\PYG{+w}{ }\PYG{n}{offset}\PYG{+w}{ }\PYG{n}{yes}

\PYG{k}{pair\PYGZus{}style}\PYG{+w}{ }\PYG{n}{spin}\PYG{o}{/}\PYG{n}{exchange}\PYG{o}{/}\PYG{n}{biquadratic}\PYG{+w}{ }\PYG{l+m}{4.0}
\PYG{k}{pair\PYGZus{}coeff}\PYG{+w}{ }\PYG{o}{*}\PYG{+w}{ }\PYG{o}{*}\PYG{+w}{ }\PYG{n}{biquadratic}\PYG{+w}{ }\PYG{l+m}{4.0}\PYG{+w}{ }\PYG{l+m}{0.05}\PYG{+w}{ }\PYG{l+m}{0.03}\PYG{+w}{ }\PYG{l+m}{1.48}\PYG{+w}{ }\PYG{l+m}{0.05}\PYG{+w}{ }\PYG{l+m}{0.03}\PYG{+w}{ }\PYG{l+m}{1.48}\PYG{+w}{ }\PYG{n}{offset}\PYG{+w}{ }\PYG{n}{no}
\PYG{k}{pair\PYGZus{}coeff}\PYG{+w}{ }\PYG{l+m}{1}\PYG{+w}{ }\PYG{l+m}{2}\PYG{+w}{ }\PYG{n}{biquadratic}\PYG{+w}{ }\PYG{l+m}{6.0}\PYG{+w}{ }\PYG{o}{\PYGZhy{}}\PYG{l+m}{0.01}\PYG{+w}{ }\PYG{l+m}{0.0}\PYG{+w}{ }\PYG{l+m}{1.9}\PYG{+w}{ }\PYG{l+m}{0.0}\PYG{+w}{ }\PYG{l+m}{0.1}\PYG{+w}{ }\PYG{l+m}{19}
\end{sphinxVerbatim}


\subsection{Description}
\label{\detokenize{pair_spin_exchange:description}}
\sphinxAtStartPar
Style \sphinxstyleemphasis{spin/exchange} computes the exchange interaction between
pairs of magnetic spins:
\begin{equation*}
\begin{split}H_{ex} = -\sum_{i,j}^N J_{ij} (r_{ij}) \,\vec{s}_i \cdot \vec{s}_j\end{split}
\end{equation*}
\sphinxAtStartPar
where \(\vec{s}_i\) and \(\vec{s}_j\) are two unit vectors representing
the magnetic spins of two particles (usually atoms), and
\(r_{ij} = \vert \vec{r}_i - \vec{r}_j \vert\) is the inter\sphinxhyphen{}atomic distance
between those two particles. The summation is over pairs of nearest neighbors.
\(J(r_{ij})\) is a function defining the intensity and the sign of the
exchange interaction for different neighboring shells.

\sphinxAtStartPar
Style \sphinxstyleemphasis{spin/exchange/biquadratic} computes a biquadratic exchange interaction
between pairs of magnetic spins:
\begin{equation*}
\begin{split}H_{bi} = -\sum_{i, j}^{N} {J}_{ij} \left(r_{ij} \right)\,
                   \vec{s}_{i}\cdot \vec{s}_{j}
                   -\sum_{i, j}^{N} {K}_{ij} \left(r_{ij} \right)\,
                   \left(\vec{s}_{i}\cdot
                   \vec{s}_{j}\right)^2\end{split}
\end{equation*}
\sphinxAtStartPar
where \(\vec{s}_i\),  \(\vec{s}_j\),  \(r_{ij}\) and
\(J(r_{ij})\) have the same definitions as above, and \(K(r_{ij})\) is
a second function, defining the intensity and the sign of the biquadratic term.

\sphinxAtStartPar
The interatomic dependence of \(J(r_{ij})\) and \(K(r_{ij})\) in both
interactions above is defined by the following function:
\begin{equation*}
\begin{split}{f}\left( r_{ij} \right) = 4 a \left( \frac{r_{ij}}{d}  \right)^2
\left( 1 - b \left( \frac{r_{ij}}{d}  \right)^2 \right)
e^{-\left( \frac{r_{ij}}{d} \right)^2 }\Theta (R_c - r_{ij})\end{split}
\end{equation*}
\sphinxAtStartPar
where \(a\), \(b\) and \(d\) are the three constant coefficients
defined in the associated “pair\_coeff” command, and \(R_c\) is the radius
cutoff associated to the pair interaction (see below for more explanations).

\sphinxAtStartPar
The coefficients \(a\), \(b\), and \(d\) need to be fitted so that
the function above matches with the value of the exchange interaction for the
\(N\) neighbor shells taken into account.
Examples and more explanations about this function and its parameterization
are reported in {\hyperref[\detokenize{pair_spin_exchange:tranchida3}]{\sphinxcrossref{\DUrole{std,std-ref}{(Tranchida)}}}}.

\sphinxAtStartPar
When a \sphinxstyleemphasis{spin/exchange/biquadratic} pair style is defined, six coefficients
(three for \(J(r_{ij})\), and three for \(K(r_{ij})\)) have to be
fitted.

\sphinxAtStartPar
From this exchange interaction, each spin \(i\) will be submitted
to a magnetic torque \(\vec{\omega}_{i}\), and its associated atom can be
submitted to a force \(\vec{F}_{i}\) for spin\sphinxhyphen{}lattice calculations (see
{\hyperref[\detokenize{fix_nve_spin::doc}]{\sphinxcrossref{\DUrole{doc}{fix nve/spin}}}}), such as:
\begin{equation*}
\begin{split}\vec{\omega}_{i} = \frac{1}{\hbar} \sum_{j}^{Neighb} {J}
\left(r_{ij} \right)\,\vec{s}_{j}
~~{\rm and}~~
\vec{F}_{i} = \sum_{j}^{Neighb} \frac{\partial {J} \left(r_{ij} \right)}{
\partial r_{ij}} \left( \vec{s}_{i}\cdot \vec{s}_{j} \right) \vec{e}_{ij}\end{split}
\end{equation*}
\sphinxAtStartPar
with \(\hbar\) the Planck constant (in metal units), and \(\vec{e}_{ij}
= \frac{\vec{r}_i - \vec{r}_j}{\vert \vec{r}_i-\vec{r}_j \vert}\) the unit
vector between sites \(i\) and \(j\).
Equivalent forces and magnetic torques are generated for the biquadratic term
when a \sphinxstyleemphasis{spin/exchange/biquadratic} pair style is defined.

\sphinxAtStartPar
More details about the derivation of these torques/forces are reported in
{\hyperref[\detokenize{pair_spin_exchange:tranchida3}]{\sphinxcrossref{\DUrole{std,std-ref}{(Tranchida)}}}}.

\sphinxAtStartPar
For the \sphinxstyleemphasis{spin/exchange} and \sphinxstyleemphasis{spin/exchange/biquadratic} pair styles, the
following coefficients must be defined for each pair of atoms types via the
{\hyperref[\detokenize{pair_coeff::doc}]{\sphinxcrossref{\DUrole{doc}{pair\_coeff}}}} command as in the examples above, or in the data
file or restart files read by the {\hyperref[\detokenize{read_data::doc}]{\sphinxcrossref{\DUrole{doc}{read\_data}}}} or
{\hyperref[\detokenize{read_restart::doc}]{\sphinxcrossref{\DUrole{doc}{read\_restart}}}} commands, and set in the following order:
\begin{itemize}
\item {} 
\sphinxAtStartPar
\(R_c\) (distance units)

\item {} 
\sphinxAtStartPar
\(a\)  (energy units)

\item {} 
\sphinxAtStartPar
\(b\)  (adim parameter)

\item {} 
\sphinxAtStartPar
\(d\)  (distance units)

\end{itemize}

\sphinxAtStartPar
for the \sphinxstyleemphasis{spin/exchange} pair style, and:
\begin{itemize}
\item {} 
\sphinxAtStartPar
\(R_c\) (distance units)

\item {} 
\sphinxAtStartPar
\(a_j\)  (energy units)

\item {} 
\sphinxAtStartPar
\(b_j\)  (adim parameter)

\item {} 
\sphinxAtStartPar
\(d_j\)  (distance units)

\item {} 
\sphinxAtStartPar
\(a_k\)  (energy units)

\item {} 
\sphinxAtStartPar
\(b_k\)  (adim parameter)

\item {} 
\sphinxAtStartPar
\(d_k\)  (distance units)

\end{itemize}

\sphinxAtStartPar
for the \sphinxstyleemphasis{spin/exchange/biquadratic} pair style.

\sphinxAtStartPar
Note that \(R_c\) is the radius cutoff of the considered exchange
interaction, and \(a\), \(b\) and \(d\) are the three coefficients
performing the parameterization of the function \(J(r_{ij})\) defined
above (in the \sphinxstyleemphasis{biquadratic} style, \(a_j\), \(b_j\), \(d_j\) and
\(a_k\), \(b_k\), \(d_k\) are the coefficients of \(J(r_{ij})\)
and \(K(r_{ij})\) respectively).

\sphinxAtStartPar
None of those coefficients is optional. If not specified, the
\sphinxstyleemphasis{spin/exchange} pair style cannot be used.


\bigskip\hrule\bigskip


\sphinxAtStartPar
\sphinxstylestrong{Offsetting magnetic forces and energies}:

\sphinxAtStartPar
For spin\sphinxhyphen{}lattice simulation, it can be useful to offset the
mechanical forces and energies generated by the exchange
interaction.
The \sphinxstyleemphasis{offset} keyword allows to apply this offset.
By setting \sphinxstyleemphasis{offset} to \sphinxstyleemphasis{yes}, the energy definitions above are
replaced by:
\begin{equation*}
\begin{split}H_{ex} = -\sum_{i,j}^N J_{ij} (r_{ij}) \,[ \vec{s}_i \cdot \vec{s}_j-1 ]\end{split}
\end{equation*}
\sphinxAtStartPar
for the \sphinxstyleemphasis{spin/exchange} pair style, and:
\begin{equation*}
\begin{split}H_{bi} = -\sum_{i, j}^{N} {J}_{ij} \left(r_{ij} \right)\,
                   [ \vec{s}_{i}\cdot \vec{s}_{j} -1 ]
                   -\sum_{i, j}^{N} {K}_{ij} \left(r_{ij} \right)\,
                   [ \left(\vec{s}_{i}\cdot
                   \vec{s}_{j}\right)^2 -1]\end{split}
\end{equation*}
\sphinxAtStartPar
for the \sphinxstyleemphasis{spin/exchange/biquadratic} pair style.

\sphinxAtStartPar
Note that this offset only affects the calculation of the energy
and mechanical forces. It does not modify the calculation of the
precession vectors (and thus does no impact the purely magnetic
properties).
This ensures that when all spins are aligned, the magnetic energy
and the associated mechanical forces (and thus the pressure
generated by the magnetic potential) are null.

\begin{sphinxadmonition}{note}{Note:}
\sphinxAtStartPar
This offset term can be very important when calculations such as
equations of state (energy vs volume, or energy vs pressure) are
being performed. Indeed, setting the \sphinxstyleemphasis{offset} term ensures that
at the ground state of the crystal and at the equilibrium magnetic
configuration (typically ferromagnetic), the pressure is null,
as expected.
Otherwise, magnetic forces could generate a residual pressure.
\end{sphinxadmonition}

\sphinxAtStartPar
When the \sphinxstyleemphasis{offset} option is set to \sphinxstyleemphasis{no}, no offset is applied
(also corresponding to the default option).


\bigskip\hrule\bigskip



\subsection{Restrictions}
\label{\detokenize{pair_spin_exchange:restrictions}}
\sphinxAtStartPar
All the \sphinxstyleemphasis{pair/spin} styles are part of the SPIN package.  These styles
are only enabled if LAMMPS was built with this package, and if the
atom\_style “spin” was declared.
See the {\hyperref[\detokenize{Build_package::doc}]{\sphinxcrossref{\DUrole{doc}{Build package}}}} page for more info.


\subsection{Related commands}
\label{\detokenize{pair_spin_exchange:related-commands}}
\sphinxAtStartPar
{\hyperref[\detokenize{atom_style::doc}]{\sphinxcrossref{\DUrole{doc}{atom\_style spin}}}}, {\hyperref[\detokenize{pair_coeff::doc}]{\sphinxcrossref{\DUrole{doc}{pair\_coeff}}}},
{\hyperref[\detokenize{pair_eam::doc}]{\sphinxcrossref{\DUrole{doc}{pair\_eam}}}},


\subsection{Default}
\label{\detokenize{pair_spin_exchange:default}}
\sphinxAtStartPar
The default \sphinxstyleemphasis{offset} keyword value is \sphinxstyleemphasis{no}.


\bigskip\hrule\bigskip

\phantomsection\label{\detokenize{pair_spin_exchange:tranchida3}}
\sphinxAtStartPar
\sphinxstylestrong{(Tranchida)} Tranchida, Plimpton, Thibaudeau and Thompson,
Journal of Computational Physics, 372, 406\sphinxhyphen{}425, (2018).

\sphinxstepscope

\index{pair\_style spin/magelec@\spxentry{pair\_style spin/magelec}}\ignorespaces 

\section{pair\_style spin/magelec command}
\label{\detokenize{pair_spin_magelec:pair-style-spin-magelec-command}}\label{\detokenize{pair_spin_magelec:index-0}}\label{\detokenize{pair_spin_magelec::doc}}

\subsection{Syntax}
\label{\detokenize{pair_spin_magelec:syntax}}
\begin{sphinxVerbatim}[commandchars=\\\{\}]
\PYG{k}{pair\PYGZus{}style}\PYG{+w}{ }\PYG{n}{spin}\PYG{o}{/}\PYG{n}{magelec}\PYG{+w}{ }\PYG{n}{cutoff}
\end{sphinxVerbatim}
\begin{itemize}
\item {} 
\sphinxAtStartPar
cutoff = global cutoff pair (distance in metal units)

\end{itemize}


\subsection{Examples}
\label{\detokenize{pair_spin_magelec:examples}}
\begin{sphinxVerbatim}[commandchars=\\\{\}]
\PYG{k}{pair\PYGZus{}style}\PYG{+w}{ }\PYG{n}{spin}\PYG{o}{/}\PYG{n}{magelec}\PYG{+w}{ }\PYG{l+m}{4.5}
\PYG{k}{pair\PYGZus{}coeff}\PYG{+w}{ }\PYG{o}{*}\PYG{+w}{ }\PYG{o}{*}\PYG{+w}{ }\PYG{n}{magelec}\PYG{+w}{ }\PYG{l+m}{4.5}\PYG{+w}{ }\PYG{l+m}{0.00109}\PYG{+w}{ }\PYG{l+m}{1.0}\PYG{+w}{ }\PYG{l+m}{1.0}\PYG{+w}{ }\PYG{l+m}{1.0}
\end{sphinxVerbatim}


\subsection{Description}
\label{\detokenize{pair_spin_magelec:description}}
\sphinxAtStartPar
Style \sphinxstyleemphasis{spin/me} computes a magneto\sphinxhyphen{}electric interaction between
pairs of magnetic spins. According to the derivation reported in
{\hyperref[\detokenize{pair_spin_magelec:katsura1}]{\sphinxcrossref{\DUrole{std,std-ref}{(Katsura)}}}}, this interaction is defined as:
\begin{equation*}
\begin{split}\vec{\omega}_i & = -\frac{1}{\hbar} \sum_{j}^{Neighb} \vec{s}_{j}\times\vec{D}(r_{ij}) \\
\vec{F}_i & = -\sum_{j}^{Neighb} \frac{\partial D(r_{ij})}{\partial r_{ij}} \left(\vec{s}_{i}\times \vec{s}_{j} \right) \cdot \vec{r}_{ij}\end{split}
\end{equation*}
\sphinxAtStartPar
where \(\vec{s}_i\) and \(\vec{s}_j\) are neighboring magnetic
spins of two particles.

\sphinxAtStartPar
From this magneto\sphinxhyphen{}electric interaction, each spin i will be submitted
to a magnetic torque omega, and its associated atom can be submitted to a
force F for spin\sphinxhyphen{}lattice calculations (see {\hyperref[\detokenize{fix_nve_spin::doc}]{\sphinxcrossref{\DUrole{doc}{fix nve/spin}}}}),
such as:
\begin{equation*}
\begin{split}\vec{F}^{i} & = -\sum_{j}^{Neighbor} \left( \vec{s}_{i}\times \vec{s}_{j} \right) \times \vec{E} \\
\vec{\omega}^{i} = -\frac{1}{\hbar} \sum_{j}^{Neighbor} \vec{s}_j \times \left(\vec{E}\times r_{ij} \right)\end{split}
\end{equation*}
\sphinxAtStartPar
with h the Planck constant (in metal units) and \(\vec{E}\) an
electric polarization vector.  The norm and direction of E are giving
the intensity and the direction of a screened dielectric atomic
polarization (in eV).

\sphinxAtStartPar
More details about the derivation of these torques/forces are reported in
{\hyperref[\detokenize{pair_spin_magelec:tranchida4}]{\sphinxcrossref{\DUrole{std,std-ref}{(Tranchida)}}}}.


\bigskip\hrule\bigskip



\subsection{Restrictions}
\label{\detokenize{pair_spin_magelec:restrictions}}
\sphinxAtStartPar
All the \sphinxstyleemphasis{pair/spin} styles are part of the SPIN package.  These styles
are only enabled if LAMMPS was built with this package, and if the
atom\_style “spin” was declared.  See the {\hyperref[\detokenize{Build_package::doc}]{\sphinxcrossref{\DUrole{doc}{Build package}}}} page for more info.


\subsection{Related commands}
\label{\detokenize{pair_spin_magelec:related-commands}}
\sphinxAtStartPar
{\hyperref[\detokenize{atom_style::doc}]{\sphinxcrossref{\DUrole{doc}{atom\_style spin}}}}, {\hyperref[\detokenize{pair_coeff::doc}]{\sphinxcrossref{\DUrole{doc}{pair\_coeff}}}},
{\hyperref[\detokenize{pair_spin_exchange::doc}]{\sphinxcrossref{\DUrole{doc}{pair\_style spin/exchange}}}}, {\hyperref[\detokenize{pair_eam::doc}]{\sphinxcrossref{\DUrole{doc}{pair\_eam}}}},


\subsection{Default}
\label{\detokenize{pair_spin_magelec:default}}
\sphinxAtStartPar
none


\bigskip\hrule\bigskip

\phantomsection\label{\detokenize{pair_spin_magelec:katsura1}}
\sphinxAtStartPar
\sphinxstylestrong{(Katsura)} H. Katsura, N. Nagaosa, A.V. Balatsky. Phys. Rev. Lett., 95(5), 057205. (2005)

\phantomsection\label{\detokenize{pair_spin_magelec:tranchida4}}
\sphinxAtStartPar
\sphinxstylestrong{(Tranchida)} Tranchida, Plimpton, Thibaudeau, and Thompson,
Journal of Computational Physics, 372, 406\sphinxhyphen{}425, (2018).

\sphinxstepscope

\index{pair\_style spin/neel@\spxentry{pair\_style spin/neel}}\ignorespaces 

\section{pair\_style spin/neel command}
\label{\detokenize{pair_spin_neel:pair-style-spin-neel-command}}\label{\detokenize{pair_spin_neel:index-0}}\label{\detokenize{pair_spin_neel::doc}}

\subsection{Syntax}
\label{\detokenize{pair_spin_neel:syntax}}
\begin{sphinxVerbatim}[commandchars=\\\{\}]
\PYG{k}{pair\PYGZus{}style}\PYG{+w}{ }\PYG{n}{spin}\PYG{o}{/}\PYG{n}{neel}\PYG{+w}{ }\PYG{n}{cutoff}
\end{sphinxVerbatim}
\begin{itemize}
\item {} 
\sphinxAtStartPar
cutoff = global cutoff pair (distance in metal units)

\end{itemize}


\subsection{Examples}
\label{\detokenize{pair_spin_neel:examples}}
\begin{sphinxVerbatim}[commandchars=\\\{\}]
\PYG{k}{pair\PYGZus{}style}\PYG{+w}{ }\PYG{n}{spin}\PYG{o}{/}\PYG{n}{neel}\PYG{+w}{ }\PYG{l+m}{4.0}
\PYG{k}{pair\PYGZus{}coeff}\PYG{+w}{ }\PYG{o}{*}\PYG{+w}{ }\PYG{o}{*}\PYG{+w}{ }\PYG{n}{neel}\PYG{+w}{ }\PYG{l+m}{4.0}\PYG{+w}{ }\PYG{l+m}{0.0048}\PYG{+w}{ }\PYG{l+m}{0.234}\PYG{+w}{ }\PYG{l+m}{1.168}\PYG{+w}{ }\PYG{l+m}{2.6905}\PYG{+w}{ }\PYG{l+m}{0.705}\PYG{+w}{ }\PYG{l+m}{0.652}
\PYG{k}{pair\PYGZus{}coeff}\PYG{+w}{ }\PYG{l+m}{1}\PYG{+w}{ }\PYG{l+m}{2}\PYG{+w}{ }\PYG{n}{neel}\PYG{+w}{ }\PYG{l+m}{4.0}\PYG{+w}{ }\PYG{l+m}{0.0048}\PYG{+w}{ }\PYG{l+m}{0.234}\PYG{+w}{ }\PYG{l+m}{1.168}\PYG{+w}{ }\PYG{l+m}{0.0}\PYG{+w}{ }\PYG{l+m}{0.0}\PYG{+w}{ }\PYG{l+m}{1.0}
\end{sphinxVerbatim}


\subsection{Description}
\label{\detokenize{pair_spin_neel:description}}
\sphinxAtStartPar
Style \sphinxstyleemphasis{spin/neel} computes the Neel pair anisotropy model
between pairs of magnetic spins:
\begin{equation*}
\begin{split}\mathcal{H}_{N\acute{e}el}=-\sum_{{ i,j=1,i\neq j}}^N g_1(r_{ij})\left(({\mathbf{e}}_{ij}\cdot {\mathbf{s}}_{i})({\mathbf{e}}_{ij}
\cdot {\mathbf{s}}_{j})-\frac{{\mathbf{s}}_{i}\cdot{\mathbf{s}}_{j}}{3} \right)
+q_1(r_{ij})\left( ({\mathbf{e}}_{ij}\cdot {\mathbf{s}}_{i})^2 -\frac{{\mathbf{s}}_{i}\cdot{\mathbf{s}}_{j}}{3}\right)
\left( ({\mathbf{e}}_{ij}\cdot {\mathbf{s}}_{i})^2 -\frac{{\mathbf{s}}_{i}\cdot{\mathbf{s}}_{j}}{3} \right)
+ q_2(r_{ij}) \Big( ({\mathbf{e}}_{ij}\cdot {\mathbf{s}}_{i}) ({\mathbf{e}}_{ij}\cdot {\mathbf{s}}_{j})^3 + ({\mathbf{e}}_{ij}\cdot
{\mathbf{s}}_{j}) ({\mathbf{e}}_{ij}\cdot {\mathbf{s}}_{i})^3\Big)\end{split}
\end{equation*}
\sphinxAtStartPar
where \(\mathbf{s}_i\) and \(\mathbf{s}_j\) are two neighboring magnetic spins of two particles,
\(r_{ij} = \vert \mathbf{r}_i - \mathbf{r}_j \vert\) is the inter\sphinxhyphen{}atomic distance between the two particles,
\(\mathbf{e}_{ij} = \frac{\mathbf{r}_i - \mathbf{r}_j}{\vert \mathbf{r}_i - \mathbf{r}_j\vert}\) is their normalized separation vector and \(g_1\),
\(q_1\) and \(q_2\) are three functions defining the intensity of the dipolar
and quadrupolar contributions, with:
\begin{equation*}
\begin{split}g_1(r_{ij}) &= g(r_{ij}) + \frac{12}{35} q(r_{ij}) \\
q_1(r_{ij}) &= \frac{9}{5} q(r_{ij}) \\
q_2(r_{ij}) &= - \frac{2}{5} q(r_{ij})\end{split}
\end{equation*}
\sphinxAtStartPar
With the functions \(g(r_{ij})\) and \(q(r_{ij})\) defined and fitted according to
the same Bethe\sphinxhyphen{}Slater function used to fit the exchange interaction:
\begin{equation*}
\begin{split}{J}\left( r_{ij} \right) = 4 a \left( \frac{r_{ij}}{d}  \right)^2 \left( 1 - b \left( \frac{r_{ij}}{d}  \right)^2 \right) e^{-\left( \frac{r_{ij}}{d} \right)^2 }\Theta (R_c - r_{ij})\end{split}
\end{equation*}
\sphinxAtStartPar
where \(a\), \(b\) and \(d\) are the three constant coefficients defined in the
associated “pair\_coeff” command.

\sphinxAtStartPar
The coefficients \(a\), \(b\), and \(d\) need to be fitted so that the function
above matches with the values of the magneto\sphinxhyphen{}elastic constant of the
materials at stake.

\sphinxAtStartPar
Examples and more explanations about this function and its
parameterization are reported in {\hyperref[\detokenize{pair_spin_neel:tranchida6}]{\sphinxcrossref{\DUrole{std,std-ref}{(Tranchida)}}}}. More
examples of parameterization will be provided in future work.

\sphinxAtStartPar
From this DM interaction, each spin \(i\) will be submitted to a magnetic
torque \(\mathbf{\omega}\) and its associated atom to a force \(\mathbf{F}\) (for spin\sphinxhyphen{}lattice
calculations only).

\sphinxAtStartPar
More details about the derivation of these torques/forces are reported
in {\hyperref[\detokenize{pair_spin_neel:tranchida6}]{\sphinxcrossref{\DUrole{std,std-ref}{(Tranchida)}}}}.


\bigskip\hrule\bigskip



\subsection{Restrictions}
\label{\detokenize{pair_spin_neel:restrictions}}
\sphinxAtStartPar
All the \sphinxstyleemphasis{pair/spin} styles are part of the SPIN package.  These styles
are only enabled if LAMMPS was built with this package, and if the
atom\_style “spin” was declared.  See the {\hyperref[\detokenize{Build_package::doc}]{\sphinxcrossref{\DUrole{doc}{Build package}}}} page for more info.


\subsection{Related commands}
\label{\detokenize{pair_spin_neel:related-commands}}
\sphinxAtStartPar
{\hyperref[\detokenize{atom_style::doc}]{\sphinxcrossref{\DUrole{doc}{atom\_style spin}}}}, {\hyperref[\detokenize{pair_coeff::doc}]{\sphinxcrossref{\DUrole{doc}{pair\_coeff}}}},
{\hyperref[\detokenize{pair_eam::doc}]{\sphinxcrossref{\DUrole{doc}{pair\_eam}}}},


\subsection{Default}
\label{\detokenize{pair_spin_neel:default}}
\sphinxAtStartPar
none


\bigskip\hrule\bigskip

\phantomsection\label{\detokenize{pair_spin_neel:tranchida6}}
\sphinxAtStartPar
\sphinxstylestrong{(Tranchida)} Tranchida, Plimpton, Thibaudeau and Thompson,
Journal of Computational Physics, 372, 406\sphinxhyphen{}425, (2018).

\sphinxstepscope

\index{pair\_style srp@\spxentry{pair\_style srp}}\ignorespaces 
\index{pair\_style srp/react@\spxentry{pair\_style srp/react}}\ignorespaces 

\section{pair\_style srp command}
\label{\detokenize{pair_srp:pair-style-srp-command}}\label{\detokenize{pair_srp:index-1}}\label{\detokenize{pair_srp:index-0}}\label{\detokenize{pair_srp::doc}}

\section{pair\_style srp/react command}
\label{\detokenize{pair_srp:pair-style-srp-react-command}}

\subsection{Syntax}
\label{\detokenize{pair_srp:syntax}}
\begin{sphinxVerbatim}[commandchars=\\\{\}]
\PYG{k}{pair\PYGZus{}style}\PYG{+w}{ }\PYG{n}{srp}\PYG{+w}{ }\PYG{n}{cutoff}\PYG{+w}{ }\PYG{n}{btype}\PYG{+w}{ }\PYG{n}{dist}\PYG{+w}{ }\PYG{n}{keyword}\PYG{+w}{ }\PYG{n}{value}\PYG{+w}{ }\PYG{n}{...}
\PYG{k}{pair\PYGZus{}style}\PYG{+w}{ }\PYG{n}{srp}\PYG{o}{/}\PYG{n}{react}\PYG{+w}{ }\PYG{n}{cutoff}\PYG{+w}{ }\PYG{n}{btype}\PYG{+w}{ }\PYG{n}{dist}\PYG{+w}{ }\PYG{n}{react}\PYG{o}{\PYGZhy{}}\PYG{n}{id}\PYG{+w}{ }\PYG{n}{keyword}\PYG{+w}{ }\PYG{n}{value}\PYG{+w}{ }\PYG{n}{...}
\end{sphinxVerbatim}
\begin{itemize}
\item {} 
\sphinxAtStartPar
cutoff = global cutoff for SRP interactions (distance units)

\item {} 
\sphinxAtStartPar
btype = bond type (numeric, type label, or wildcard) to apply SRP interactions to

\item {} 
\sphinxAtStartPar
distance = \sphinxstyleemphasis{min} or \sphinxstyleemphasis{mid}

\item {} 
\sphinxAtStartPar
react\sphinxhyphen{}id = id of either fix bond/break or fix bond/create

\item {} 
\sphinxAtStartPar
zero or more keyword/value pairs may be appended

\item {} 
\sphinxAtStartPar
keyword = \sphinxstyleemphasis{exclude}
\begin{sphinxalltt}
\sphinxstyleemphasis{bptype} value = atom type (numeric or type label) for bond particles
\sphinxstyleemphasis{exclude} value = \sphinxstyleemphasis{yes} or \sphinxstyleemphasis{no}
\end{sphinxalltt}

\end{itemize}


\subsection{Examples}
\label{\detokenize{pair_srp:examples}}
\begin{sphinxVerbatim}[commandchars=\\\{\}]
\PYG{k}{pair\PYGZus{}style}\PYG{+w}{ }\PYG{n}{hybrid}\PYG{+w}{ }\PYG{n}{dpd}\PYG{+w}{ }\PYG{l+m}{1.0}\PYG{+w}{ }\PYG{l+m}{1.0}\PYG{+w}{ }\PYG{l+m}{12345}\PYG{+w}{ }\PYG{n}{srp}\PYG{+w}{ }\PYG{l+m}{0.8}\PYG{+w}{ }\PYG{l+m}{1}\PYG{+w}{ }\PYG{n}{mid}\PYG{+w}{ }\PYG{n}{exclude}\PYG{+w}{ }\PYG{n}{yes}
\PYG{k}{pair\PYGZus{}coeff}\PYG{+w}{ }\PYG{l+m}{1}\PYG{+w}{ }\PYG{l+m}{1}\PYG{+w}{ }\PYG{n}{dpd}\PYG{+w}{ }\PYG{l+m}{60.0}\PYG{+w}{ }\PYG{l+m}{4.5}\PYG{+w}{ }\PYG{l+m}{1.0}
\PYG{k}{pair\PYGZus{}coeff}\PYG{+w}{ }\PYG{l+m}{1}\PYG{+w}{ }\PYG{l+m}{2}\PYG{+w}{ }\PYG{n}{none}
\PYG{k}{pair\PYGZus{}coeff}\PYG{+w}{ }\PYG{l+m}{2}\PYG{+w}{ }\PYG{l+m}{2}\PYG{+w}{ }\PYG{n}{srp}\PYG{+w}{ }\PYG{l+m}{100.0}\PYG{+w}{ }\PYG{l+m}{0.8}

\PYG{k}{pair\PYGZus{}style}\PYG{+w}{ }\PYG{n}{hybrid}\PYG{+w}{ }\PYG{n}{dpd}\PYG{+w}{ }\PYG{l+m}{1.0}\PYG{+w}{ }\PYG{l+m}{1.0}\PYG{+w}{ }\PYG{l+m}{12345}\PYG{+w}{ }\PYG{n}{srp}\PYG{+w}{ }\PYG{l+m}{0.8}\PYG{+w}{ }\PYG{o}{*}\PYG{+w}{ }\PYG{n}{min}\PYG{+w}{ }\PYG{n}{exclude}\PYG{+w}{ }\PYG{n}{yes}
\PYG{k}{pair\PYGZus{}coeff}\PYG{+w}{ }\PYG{l+m}{1}\PYG{+w}{ }\PYG{l+m}{1}\PYG{+w}{ }\PYG{n}{dpd}\PYG{+w}{ }\PYG{l+m}{60.0}\PYG{+w}{ }\PYG{l+m}{50}\PYG{+w}{ }\PYG{l+m}{1.0}
\PYG{k}{pair\PYGZus{}coeff}\PYG{+w}{ }\PYG{l+m}{1}\PYG{+w}{ }\PYG{l+m}{2}\PYG{+w}{ }\PYG{n}{none}
\PYG{k}{pair\PYGZus{}coeff}\PYG{+w}{ }\PYG{l+m}{2}\PYG{+w}{ }\PYG{l+m}{2}\PYG{+w}{ }\PYG{n}{srp}\PYG{+w}{ }\PYG{l+m}{40.0}

\PYG{k}{fix        }\PYG{n+nv+nvIdentifier}{create}\PYG{+w}{ }\PYG{n+nv+nvIdentifier}{all}\PYG{+w}{ }\PYG{n}{bond}\PYG{o}{/}\PYG{n}{create}\PYG{+w}{   }\PYG{l+m}{100}\PYG{+w}{  }\PYG{l+m}{1}\PYG{+w}{  }\PYG{l+m}{2}\PYG{+w}{  }\PYG{l+m}{1.0}\PYG{+w}{ }\PYG{l+m}{1}\PYG{+w}{ }\PYG{n}{prob}\PYG{+w}{  }\PYG{l+m}{0.2}\PYG{+w}{  }\PYG{l+m}{19852}
\PYG{k}{pair\PYGZus{}style}\PYG{+w}{ }\PYG{n}{hybrid}\PYG{+w}{ }\PYG{n}{dpd}\PYG{+w}{ }\PYG{l+m}{1.0}\PYG{+w}{ }\PYG{l+m}{1.0}\PYG{+w}{ }\PYG{l+m}{12345}\PYG{+w}{ }\PYG{n}{srp}\PYG{o}{/}\PYG{n}{react}\PYG{+w}{ }\PYG{l+m}{0.8}\PYG{+w}{ }\PYG{o}{*}\PYG{+w}{ }\PYG{n}{min}\PYG{+w}{ }\PYG{n}{create}\PYG{+w}{ }\PYG{n}{exclude}\PYG{+w}{ }\PYG{n}{yes}
\PYG{k}{pair\PYGZus{}coeff}\PYG{+w}{ }\PYG{l+m}{1}\PYG{+w}{ }\PYG{l+m}{1}\PYG{+w}{ }\PYG{n}{dpd}\PYG{+w}{ }\PYG{l+m}{60.0}\PYG{+w}{ }\PYG{l+m}{50}\PYG{+w}{ }\PYG{l+m}{1.0}
\PYG{k}{pair\PYGZus{}coeff}\PYG{+w}{ }\PYG{l+m}{1}\PYG{+w}{ }\PYG{l+m}{2}\PYG{+w}{ }\PYG{n}{none}
\PYG{k}{pair\PYGZus{}coeff}\PYG{+w}{ }\PYG{l+m}{2}\PYG{+w}{ }\PYG{l+m}{2}\PYG{+w}{ }\PYG{n}{srp}\PYG{o}{/}\PYG{n}{react}\PYG{+w}{ }\PYG{l+m}{40.0}

\PYG{k}{pair\PYGZus{}style}\PYG{+w}{ }\PYG{n}{hybrid}\PYG{+w}{ }\PYG{n}{srp}\PYG{+w}{ }\PYG{l+m}{0.8}\PYG{+w}{ }\PYG{l+m}{2}\PYG{+w}{ }\PYG{n}{mid}
\PYG{k}{pair\PYGZus{}coeff}\PYG{+w}{ }\PYG{l+m}{1}\PYG{+w}{ }\PYG{l+m}{1}\PYG{+w}{ }\PYG{n}{none}
\PYG{k}{pair\PYGZus{}coeff}\PYG{+w}{ }\PYG{l+m}{1}\PYG{+w}{ }\PYG{l+m}{2}\PYG{+w}{ }\PYG{n}{none}
\PYG{k}{pair\PYGZus{}coeff}\PYG{+w}{ }\PYG{l+m}{2}\PYG{+w}{ }\PYG{l+m}{2}\PYG{+w}{ }\PYG{n}{srp}\PYG{+w}{ }\PYG{l+m}{100.0}\PYG{+w}{ }\PYG{l+m}{0.8}

\PYG{k}{labelmap}\PYG{+w}{ }\PYG{n}{bond}\PYG{+w}{ }\PYG{l+m}{1}\PYG{+w}{ }\PYG{n}{C}\PYG{o}{\PYGZhy{}}\PYG{n}{C}
\PYG{k}{pair\PYGZus{}style}\PYG{+w}{ }\PYG{n}{hybrid}\PYG{+w}{ }\PYG{n}{srp}\PYG{+w}{ }\PYG{l+m}{0.8}\PYG{+w}{ }\PYG{n}{C}\PYG{o}{\PYGZhy{}}\PYG{n}{C}\PYG{+w}{ }\PYG{n}{mid}
\end{sphinxVerbatim}

\sphinxAtStartPar
Description

\sphinxAtStartPar
Style \sphinxstyleemphasis{srp} computes a soft segmental repulsive potential (SRP) that
acts between pairs of bonds. This potential is useful for preventing
bonds from passing through one another when a soft non\sphinxhyphen{}bonded
potential acts between beads in, for example, DPD polymer chains.  An
example input script that uses this command is provided in
examples/PACKAGES/srp.

\sphinxAtStartPar
Bonds of specified type \sphinxstyleemphasis{btype} interact with one another through a
bond\sphinxhyphen{}pairwise potential, such that the force on bond \sphinxstyleemphasis{i} due to bond
\sphinxstyleemphasis{j} is as follows
\begin{equation*}
\begin{split}F^{\mathrm{SRP}}_{ij} = C(1-r/r_c)\hat{r}_{ij} \qquad r < r_c\end{split}
\end{equation*}
\sphinxAtStartPar
where \sphinxstyleemphasis{r} and \(\hat{r}_{ij}\) are the distance and unit vector
between the two bonds.  Note that \sphinxstyleemphasis{btype} can be specified as an
asterisk “*”, which case the interaction is applied to all bond types.
The \sphinxstyleemphasis{mid} option computes \sphinxstyleemphasis{r} and \(\hat{r}_{ij}\) from the midpoint
distance between bonds. The \sphinxstyleemphasis{min} option computes \sphinxstyleemphasis{r} and
\(\hat{r}_{ij}\) from the minimum distance between bonds. The force
acting on a bond is mapped onto the two bond atoms according to the
lever rule,
\begin{equation*}
\begin{split}F_{i1}^{\mathrm{SRP}} & = F^{\mathrm{SRP}}_{ij}(L) \\
F_{i2}^{\mathrm{SRP}} & = F^{\mathrm{SRP}}_{ij}(1-L)\end{split}
\end{equation*}
\sphinxAtStartPar
where \sphinxstyleemphasis{L} is the normalized distance from the atom to the point of
closest approach of bond \sphinxstyleemphasis{i} and \sphinxstyleemphasis{j}. The \sphinxstyleemphasis{mid} option takes \sphinxstyleemphasis{L} as
0.5 for each interaction as described in {\hyperref[\detokenize{pair_srp:sirk2}]{\sphinxcrossref{\DUrole{std,std-ref}{(Sirk)}}}}.

\sphinxAtStartPar
The following coefficients must be defined via the
{\hyperref[\detokenize{pair_coeff::doc}]{\sphinxcrossref{\DUrole{doc}{pair\_coeff}}}} command as in the examples above, or in
the data file or restart file read by the {\hyperref[\detokenize{read_data::doc}]{\sphinxcrossref{\DUrole{doc}{read\_data}}}}
or {\hyperref[\detokenize{read_restart::doc}]{\sphinxcrossref{\DUrole{doc}{read\_restart}}}} commands:
\begin{itemize}
\item {} 
\sphinxAtStartPar
\sphinxstyleemphasis{C} (force units)

\item {} 
\sphinxAtStartPar
\(r_c\) (distance units)

\end{itemize}

\sphinxAtStartPar
The last coefficient is optional. If not specified, the global cutoff
is used.

\begin{sphinxadmonition}{note}{Note:}
\sphinxAtStartPar
Pair style srp considers each bond of type \sphinxstyleemphasis{btype} to be a
fictitious “particle” of type \sphinxstyleemphasis{bptype}, where \sphinxstyleemphasis{bptype} is either the
largest atom type in the system, or the type set by the \sphinxstyleemphasis{bptype} flag.
Any actual existing particles with this atom type will be deleted at
the beginning of a run. This means you must specify the number of
types in your system accordingly; usually to be one larger than what
would normally be the case, e.g. via the {\hyperref[\detokenize{create_box::doc}]{\sphinxcrossref{\DUrole{doc}{create\_box}}}}
or by changing the header in your {\hyperref[\detokenize{read_data::doc}]{\sphinxcrossref{\DUrole{doc}{data file}}}}.  The
fictitious “bond particles” are inserted at the beginning of the run,
and serve as placeholders that define the position of the bonds.  This
allows neighbor lists to be constructed and pairwise interactions to
be computed in almost the same way as is done for actual particles.
Because bonds interact only with other bonds, {\hyperref[\detokenize{pair_hybrid::doc}]{\sphinxcrossref{\DUrole{doc}{pair\_style hybrid}}}} should be used to turn off interactions
between atom type \sphinxstyleemphasis{bptype} and all other types of atoms.  An error
will be flagged if {\hyperref[\detokenize{pair_hybrid::doc}]{\sphinxcrossref{\DUrole{doc}{pair\_style hybrid}}}} is not used.
\end{sphinxadmonition}

\begin{sphinxadmonition}{note}{Note:}
\sphinxAtStartPar
If using type labels, the type labels must be defined before calling
the {\hyperref[\detokenize{pair_coeff::doc}]{\sphinxcrossref{\DUrole{doc}{pair\_coeff}}}} command.
\end{sphinxadmonition}

\sphinxAtStartPar
The optional \sphinxstyleemphasis{exclude} keyword determines if forces are computed
between first neighbor (directly connected) bonds.  For a setting of
\sphinxstyleemphasis{no}, first neighbor forces are computed; for \sphinxstyleemphasis{yes} they are not
computed. A setting of \sphinxstyleemphasis{no} cannot be used with the \sphinxstyleemphasis{min} option for
distance calculation because the minimum distance between directly
connected bonds is zero.

\sphinxAtStartPar
Pair style \sphinxstyleemphasis{srp} turns off normalization of thermodynamic properties
by particle number, as if the command {\hyperref[\detokenize{thermo_modify::doc}]{\sphinxcrossref{\DUrole{doc}{thermo\_modify norm no}}}} had been issued.

\sphinxAtStartPar
The pairwise energy associated with style \sphinxstyleemphasis{srp} is shifted to be zero
at the cutoff distance \(r_c\).


\bigskip\hrule\bigskip


\sphinxAtStartPar
\DUrole{versionmodified,added}{Added in version 3Aug2022.}

\sphinxAtStartPar
Pair style \sphinxstyleemphasis{srp/react} interfaces the pair style \sphinxstyleemphasis{srp} with the
bond breaking and formation mechanisms provided by fix \sphinxstyleemphasis{bond/break}
and fix \sphinxstyleemphasis{bond/create}, respectively. When using this pair style, whenever a
bond breaking (or formation) reaction occurs, the corresponding fictitious
particle is deleted (or inserted) during the same simulation time step as
the reaction. This is useful in the simulation of reactive systems involving
large polymeric molecules {\hyperref[\detokenize{pair_srp:palkar}]{\sphinxcrossref{\DUrole{std,std-ref}{(Palkar)}}}}  where the segmental repulsive
potential is necessary to minimize topological violations, and also needs to be
turned on and off according to the progress of the reaction.


\bigskip\hrule\bigskip



\subsection{Mixing, shift, table, tail correction, restart, rRESPA info}
\label{\detokenize{pair_srp:mixing-shift-table-tail-correction-restart-rrespa-info}}
\sphinxAtStartPar
This pair style does not support mixing.

\sphinxAtStartPar
This pair style does not support the {\hyperref[\detokenize{pair_modify::doc}]{\sphinxcrossref{\DUrole{doc}{pair\_modify}}}}
shift option for the energy of the pair interaction. Note that as
discussed above, the energy term is already shifted to be 0.0 at the
cutoff distance \(r_c\).

\sphinxAtStartPar
The {\hyperref[\detokenize{pair_modify::doc}]{\sphinxcrossref{\DUrole{doc}{pair\_modify}}}} table option is not relevant for
this pair style.

\sphinxAtStartPar
This pair style does not support the {\hyperref[\detokenize{pair_modify::doc}]{\sphinxcrossref{\DUrole{doc}{pair\_modify}}}}
tail option for adding long\sphinxhyphen{}range tail corrections to energy and
pressure.

\sphinxAtStartPar
This pair style writes global and per\sphinxhyphen{}atom information to {\hyperref[\detokenize{restart::doc}]{\sphinxcrossref{\DUrole{doc}{binary restart files}}}}. Pair srp should be used with {\hyperref[\detokenize{pair_hybrid::doc}]{\sphinxcrossref{\DUrole{doc}{pair\_style hybrid}}}}, thus the pair\_coeff commands need to be
specified in the input script when reading a restart file.

\sphinxAtStartPar
This pair style can only be used via the \sphinxstyleemphasis{pair} keyword of the
{\hyperref[\detokenize{run_style::doc}]{\sphinxcrossref{\DUrole{doc}{run\_style respa}}}} command.  It does not support the
\sphinxstyleemphasis{inner}, \sphinxstyleemphasis{middle}, \sphinxstyleemphasis{outer} keywords.


\bigskip\hrule\bigskip



\subsection{Restrictions}
\label{\detokenize{pair_srp:restrictions}}
\sphinxAtStartPar
This pair style is part of the MISC package. It is only enabled
if LAMMPS was built with that package. See the Making LAMMPS section
for more info.

\sphinxAtStartPar
This pair style must be used with {\hyperref[\detokenize{pair_hybrid::doc}]{\sphinxcrossref{\DUrole{doc}{pair\_style hybrid}}}}.

\sphinxAtStartPar
This pair style requires the {\hyperref[\detokenize{newton::doc}]{\sphinxcrossref{\DUrole{doc}{newton}}}} command to be \sphinxstyleemphasis{on}
for non\sphinxhyphen{}bonded interactions.

\sphinxAtStartPar
This pair style is not compatible with {\hyperref[\detokenize{fix_rigid::doc}]{\sphinxcrossref{\DUrole{doc}{rigid body integrators}}}}


\subsection{Related commands}
\label{\detokenize{pair_srp:related-commands}}
\sphinxAtStartPar
{\hyperref[\detokenize{pair_hybrid::doc}]{\sphinxcrossref{\DUrole{doc}{pair\_style hybrid}}}}, {\hyperref[\detokenize{pair_coeff::doc}]{\sphinxcrossref{\DUrole{doc}{pair\_coeff}}}},
{\hyperref[\detokenize{pair_dpd::doc}]{\sphinxcrossref{\DUrole{doc}{pair dpd}}}}


\subsection{Default}
\label{\detokenize{pair_srp:default}}
\sphinxAtStartPar
The default keyword value is exclude = yes.


\bigskip\hrule\bigskip

\phantomsection\label{\detokenize{pair_srp:sirk2}}
\sphinxAtStartPar
\sphinxstylestrong{(Sirk)} Sirk TW, Sliozberg YR, Brennan JK, Lisal M, Andzelm JW, J
Chem Phys, 136 (13) 134903, 2012.

\phantomsection\label{\detokenize{pair_srp:palkar}}
\sphinxAtStartPar
\sphinxstylestrong{(Palkar)} Palkar V, Kuksenok O, J. Phys. Chem. B, 126 (1), 336\sphinxhyphen{}346, 2022

\sphinxstepscope

\index{pair\_style sw@\spxentry{pair\_style sw}}\ignorespaces 
\index{pair\_style sw/gpu@\spxentry{pair\_style sw/gpu}}\ignorespaces 
\index{pair\_style sw/intel@\spxentry{pair\_style sw/intel}}\ignorespaces 
\index{pair\_style sw/kk@\spxentry{pair\_style sw/kk}}\ignorespaces 
\index{pair\_style sw/omp@\spxentry{pair\_style sw/omp}}\ignorespaces 
\index{pair\_style sw/mod@\spxentry{pair\_style sw/mod}}\ignorespaces 
\index{pair\_style sw/mod/omp@\spxentry{pair\_style sw/mod/omp}}\ignorespaces 

\section{pair\_style sw command}
\label{\detokenize{pair_sw:pair-style-sw-command}}\label{\detokenize{pair_sw:index-6}}\label{\detokenize{pair_sw:index-5}}\label{\detokenize{pair_sw:index-4}}\label{\detokenize{pair_sw:index-3}}\label{\detokenize{pair_sw:index-2}}\label{\detokenize{pair_sw:index-1}}\label{\detokenize{pair_sw:index-0}}\label{\detokenize{pair_sw::doc}}
\sphinxAtStartPar
Accelerator Variants: \sphinxstyleemphasis{sw/gpu}, \sphinxstyleemphasis{sw/intel}, \sphinxstyleemphasis{sw/kk}, \sphinxstyleemphasis{sw/omp}


\section{pair\_style sw/mod command}
\label{\detokenize{pair_sw:pair-style-sw-mod-command}}
\sphinxAtStartPar
Accelerator Variants: \sphinxstyleemphasis{sw/mod/omp}


\subsection{Syntax}
\label{\detokenize{pair_sw:syntax}}
\begin{sphinxVerbatim}[commandchars=\\\{\}]
\PYG{k}{pair\PYGZus{}style}\PYG{+w}{ }\PYG{n}{style}\PYG{+w}{ }\PYG{n}{keyword}\PYG{+w}{ }\PYG{n}{values}
\end{sphinxVerbatim}
\begin{itemize}
\item {} 
\sphinxAtStartPar
style = \sphinxstyleemphasis{sw} or \sphinxstyleemphasis{sw/mod}

\item {} 
\sphinxAtStartPar
keyword = \sphinxstyleemphasis{maxdelcs} or \sphinxstyleemphasis{threebody}
\begin{sphinxalltt}
\sphinxstyleemphasis{maxdelcs} value = delta1 delta2 (optional, sw/mod only)
  delta1 = The minimum thershold for the variation of cosine of three\sphinxhyphen{}body angle
  delta2 = The maximum threshold for the variation of cosine of three\sphinxhyphen{}body angle
\sphinxstyleemphasis{threebody} value = \sphinxstyleemphasis{on} or \sphinxstyleemphasis{off} (optional, sw only)
  on (default) = Compute both the three\sphinxhyphen{}body and two\sphinxhyphen{}body terms of the potential
  off = Compute only the two\sphinxhyphen{}body term of the potential
\end{sphinxalltt}

\end{itemize}


\subsection{Examples}
\label{\detokenize{pair_sw:examples}}
\begin{sphinxVerbatim}[commandchars=\\\{\}]
\PYG{k}{pair\PYGZus{}style}\PYG{+w}{ }\PYG{n}{sw}
\PYG{k}{pair\PYGZus{}coeff}\PYG{+w}{ }\PYG{o}{*}\PYG{+w}{ }\PYG{o}{*}\PYG{+w}{ }\PYG{n}{si.sw}\PYG{+w}{ }\PYG{n}{Si}
\PYG{k}{pair\PYGZus{}coeff}\PYG{+w}{ }\PYG{o}{*}\PYG{+w}{ }\PYG{o}{*}\PYG{+w}{ }\PYG{n}{GaN.sw}\PYG{+w}{ }\PYG{n}{Ga}\PYG{+w}{ }\PYG{n}{N}\PYG{+w}{ }\PYG{n}{Ga}

\PYG{k}{pair\PYGZus{}style}\PYG{+w}{ }\PYG{n}{sw}\PYG{o}{/}\PYG{n}{mod}\PYG{+w}{ }\PYG{n}{maxdelcs}\PYG{+w}{ }\PYG{l+m}{0.25}\PYG{+w}{ }\PYG{l+m}{0.35}
\PYG{k}{pair\PYGZus{}coeff}\PYG{+w}{ }\PYG{o}{*}\PYG{+w}{ }\PYG{o}{*}\PYG{+w}{ }\PYG{n}{tmd.sw.mod}\PYG{+w}{ }\PYG{n}{Mo}\PYG{+w}{ }\PYG{n}{S}\PYG{+w}{ }\PYG{n}{S}

\PYG{k}{pair\PYGZus{}style}\PYG{+w}{ }\PYG{n}{hybrid}\PYG{+w}{ }\PYG{n}{sw}\PYG{+w}{ }\PYG{n}{threebody}\PYG{+w}{ }\PYG{n}{on}\PYG{+w}{ }\PYG{n}{sw}\PYG{+w}{ }\PYG{n}{threebody}\PYG{+w}{ }\PYG{n}{off}
\PYG{k}{pair\PYGZus{}coeff}\PYG{+w}{ }\PYG{o}{*}\PYG{+w}{ }\PYG{o}{*}\PYG{+w}{ }\PYG{n}{sw}\PYG{+w}{ }\PYG{l+m}{1}\PYG{+w}{ }\PYG{n}{mW\PYGZus{}xL.sw}\PYG{+w}{ }\PYG{n}{mW}\PYG{+w}{ }\PYG{n}{NULL}
\PYG{k}{pair\PYGZus{}coeff}\PYG{+w}{ }\PYG{l+m}{1}\PYG{+w}{ }\PYG{l+m}{2}\PYG{+w}{ }\PYG{n}{sw}\PYG{+w}{ }\PYG{l+m}{2}\PYG{+w}{ }\PYG{n}{mW\PYGZus{}xL.sw}\PYG{+w}{ }\PYG{n}{mW}\PYG{+w}{ }\PYG{n}{xL}
\PYG{k}{pair\PYGZus{}coeff}\PYG{+w}{ }\PYG{l+m}{2}\PYG{+w}{ }\PYG{l+m}{2}\PYG{+w}{ }\PYG{n}{sw}\PYG{+w}{ }\PYG{l+m}{2}\PYG{+w}{ }\PYG{n}{mW\PYGZus{}xL.sw}\PYG{+w}{ }\PYG{n}{mW}\PYG{+w}{ }\PYG{n}{xL}
\end{sphinxVerbatim}


\subsection{Description}
\label{\detokenize{pair_sw:description}}
\sphinxAtStartPar
The \sphinxstyleemphasis{sw} style computes a 3\sphinxhyphen{}body {\hyperref[\detokenize{pair_sw:stillinger2}]{\sphinxcrossref{\DUrole{std,std-ref}{Stillinger\sphinxhyphen{}Weber}}}}
potential for the energy E of a system of atoms as
\begin{equation*}
\begin{split} E & =  \sum_i \sum_{j > i} \phi_2 (r_{ij}) +
        \sum_i \sum_{j \neq i} \sum_{k > j}
        \phi_3 (r_{ij}, r_{ik}, \theta_{ijk}) \\
\phi_2(r_{ij}) & =  A_{ij} \epsilon_{ij} \left[ B_{ij} (\frac{\sigma_{ij}}{r_{ij}})^{p_{ij}} -
                  (\frac{\sigma_{ij}}{r_{ij}})^{q_{ij}} \right]
                  \exp \left( \frac{\sigma_{ij}}{r_{ij} - a_{ij} \sigma_{ij}} \right) \\
\phi_3(r_{ij},r_{ik},\theta_{ijk}) & = \lambda_{ijk} \epsilon_{ijk} \left[ \cos \theta_{ijk} -
                  \cos \theta_{0ijk} \right]^2
                  \exp \left( \frac{\gamma_{ij} \sigma_{ij}}{r_{ij} - a_{ij} \sigma_{ij}} \right)
                  \exp \left( \frac{\gamma_{ik} \sigma_{ik}}{r_{ik} - a_{ik} \sigma_{ik}} \right)\end{split}
\end{equation*}
\sphinxAtStartPar
where \(\phi_2\) is a two\sphinxhyphen{}body term and \(\phi_3\) is a
three\sphinxhyphen{}body term.  The summations in the formula are over all neighbors J
and K of atom I within a cutoff distance \(a `\sigma\).

\sphinxAtStartPar
\DUrole{versionmodified,added}{Added in version 14Dec2021.}

\sphinxAtStartPar
The \sphinxstyleemphasis{sw/mod} style is designed for simulations of materials when
distinguishing three\sphinxhyphen{}body angles are necessary, such as borophene and
transition metal dichalcogenides, which cannot be described by the
original code for the Stillinger\sphinxhyphen{}Weber potential.  For instance, there
are several types of angles around each Mo atom in \sphinxtitleref{MoS\_2}, and some
unnecessary angle types should be excluded in the three\sphinxhyphen{}body
interaction.  Such exclusion may be realized by selecting proper angle
types directly.  The exclusion of unnecessary angles is achieved here by
the cut\sphinxhyphen{}off function (\sphinxtitleref{f\_C(delta)}), which induces only minimum
modifications for LAMMPS.

\sphinxAtStartPar
Validation, benchmark tests, and applications of the \sphinxstyleemphasis{sw/mod} style
can be found in {\hyperref[\detokenize{pair_sw:jiang2}]{\sphinxcrossref{\DUrole{std,std-ref}{(Jiang2)}}}} and {\hyperref[\detokenize{pair_sw:jiang3}]{\sphinxcrossref{\DUrole{std,std-ref}{(Jiang3)}}}}.

\sphinxAtStartPar
The \sphinxstyleemphasis{sw/mod} style computes the energy E of a system of atoms, whose
potential function is mostly the same as the Stillinger\sphinxhyphen{}Weber
potential. The only modification is in the three\sphinxhyphen{}body term, where the
value of \(\delta = \cos \theta_{ijk} - \cos \theta_{0ijk}\) used in
the original energy and force expression is scaled by a switching factor
\(f_C(\delta)\):
\begin{equation*}
\begin{split}f_C(\delta) & = \left\{ \begin{array} {r@{\quad:\quad}l}
  1 & \left| \delta \right| < \delta_1 \\
  \frac{1}{2} + \frac{1}{2} \cos \left( \pi \frac{\left| \delta \right| - \delta_1}{\delta_2 - \delta_1} \right) &
    \delta_1 < \left| \delta \right| < \delta_2 \\
  0 & \left| \delta \right| > \delta_2
  \end{array} \right. \\\end{split}
\end{equation*}
\sphinxAtStartPar
This cut\sphinxhyphen{}off function decreases smoothly from 1 to 0 over the range
\([\delta_1, \delta_2]\).  This smoothly turns off the energy and
force contributions for \(\left| \delta \right| > \delta_2\).  It is
suggested that \(\delta 1\) and \(\delta_2\) to be the value
around \(0.5 \left| \cos \theta_1 - \cos \theta_2 \right|\), with
\(\theta_1\) and \(\theta_2\) as the different types of angles
around an atom.  For borophene and transition metal dichalcogenides,
\(\delta_1 = 0.25\) and \(\delta_2 = 0.35\).  This value enables
the cut\sphinxhyphen{}off function to exclude unnecessary angles in the three\sphinxhyphen{}body SW
terms.

\begin{sphinxadmonition}{note}{Note:}
\sphinxAtStartPar
The cut\sphinxhyphen{}off function is just to be used as a technique to exclude
some unnecessary angles, and it has no physical meaning. It should be
noted that the force and potential are inconsistent with each other
in the decaying range of the cut\sphinxhyphen{}off function, as the angle
dependence for the cut\sphinxhyphen{}off function is not implemented in the force
(first derivation of potential).  However, the angle variation is
much smaller than the given threshold value for actual simulations,
so the inconsistency between potential and force can be neglected in
actual simulations.
\end{sphinxadmonition}

\sphinxAtStartPar
\DUrole{versionmodified,added}{Added in version 3Aug2022.}

\sphinxAtStartPar
The \sphinxstyleemphasis{threebody} keyword is optional and determines whether or not the
three\sphinxhyphen{}body term of the potential is calculated.  The default value is
“on” and it is only available for the plain \sphinxstyleemphasis{sw} pair style variants,
but not available for the \sphinxstyleemphasis{sw/mod} and {\hyperref[\detokenize{pair_sw_angle_table::doc}]{\sphinxcrossref{\DUrole{doc}{sw/angle/table}}}} pair style variants.  To turn off the threebody
contributions all \(\lambda_{ijk}\) parameters from the potential
file are forcibly set to 0.  In addition the pair style implementation
may employ code optimizations for the \sphinxstyleemphasis{threebody off} setting that can
result in significant speedups versus the default.  These code optimizations
are currently only available for the MANYBODY and OPENMP packages.

\sphinxAtStartPar
Only a single pair\_coeff command is used with the \sphinxstyleemphasis{sw} and \sphinxstyleemphasis{sw/mod}
styles which specifies a Stillinger\sphinxhyphen{}Weber potential file with parameters
for all needed elements, except for when the \sphinxstyleemphasis{threebody off} setting is
used (see note below).  These are mapped to LAMMPS atom types by
specifying N additional arguments after the filename in the pair\_coeff
command, where N is the number of LAMMPS atom types:
\begin{itemize}
\item {} 
\sphinxAtStartPar
filename

\item {} 
\sphinxAtStartPar
N element names = mapping of SW elements to atom types

\end{itemize}

\sphinxAtStartPar
See the {\hyperref[\detokenize{pair_coeff::doc}]{\sphinxcrossref{\DUrole{doc}{pair\_coeff}}}} page for alternate ways
to specify the path for the potential file.

\sphinxAtStartPar
As an example, imagine a file SiC.sw has Stillinger\sphinxhyphen{}Weber values for
Si and C.  If your LAMMPS simulation has 4 atoms types and you want
the first 3 to be Si, and the fourth to be C, you would use the following
pair\_coeff command:

\begin{sphinxVerbatim}[commandchars=\\\{\}]
\PYG{k}{pair\PYGZus{}style}\PYG{+w}{ }\PYG{n}{sw}
\PYG{k}{pair\PYGZus{}coeff}\PYG{+w}{ }\PYG{o}{*}\PYG{+w}{ }\PYG{o}{*}\PYG{+w}{ }\PYG{n}{SiC.sw}\PYG{+w}{ }\PYG{n}{Si}\PYG{+w}{ }\PYG{n}{Si}\PYG{+w}{ }\PYG{n}{Si}\PYG{+w}{ }\PYG{n}{C}
\end{sphinxVerbatim}

\sphinxAtStartPar
The first 2 arguments must be * * so as to span all LAMMPS atom types.
The first three Si arguments map LAMMPS atom types 1, 2, and 3 to the Si
element in the SW file.  The final C argument maps LAMMPS atom type 4 to
the C element in the SW file.  If an argument value is specified as
NULL, the mapping is not performed.  This can be used when an \sphinxstyleemphasis{sw}
potential is used as part of the \sphinxstyleemphasis{hybrid} pair style.  The NULL values
are placeholders for atom types that will be used with other potentials.

\begin{sphinxadmonition}{note}{Note:}
\sphinxAtStartPar
When the \sphinxstyleemphasis{threebody off} keyword is used, multiple pair\_coeff
commands may be used to specific the pairs of atoms which don’t
require three\sphinxhyphen{}body term.  In these cases, the first 2 arguments are
not required to be * *, the potential parameter file is only read
by the first {\hyperref[\detokenize{pair_coeff::doc}]{\sphinxcrossref{\DUrole{doc}{pair\_coeff command}}}} and the element
to atom type mappings must be consistent across all \sphinxstyleemphasis{pair\_coeff}
statements.  If not LAMMPS will abort with an error.
\end{sphinxadmonition}

\sphinxAtStartPar
Stillinger\sphinxhyphen{}Weber files in the \sphinxstyleemphasis{potentials} directory of the LAMMPS
distribution have a “.sw” suffix.  Lines that are not blank or
comments (starting with \#) define parameters for a triplet of
elements.  The parameters in a single entry correspond to the two\sphinxhyphen{}body
and three\sphinxhyphen{}body coefficients in the formula above:
\begin{itemize}
\item {} 
\sphinxAtStartPar
element 1 (the center atom in a 3\sphinxhyphen{}body interaction)

\item {} 
\sphinxAtStartPar
element 2

\item {} 
\sphinxAtStartPar
element 3

\item {} 
\sphinxAtStartPar
\(\epsilon\) (energy units)

\item {} 
\sphinxAtStartPar
\(\sigma\) (distance units)

\item {} 
\sphinxAtStartPar
a

\item {} 
\sphinxAtStartPar
\(\lambda\)

\item {} 
\sphinxAtStartPar
\(\gamma\)

\item {} 
\sphinxAtStartPar
\(\cos\theta_0\)

\item {} 
\sphinxAtStartPar
A

\item {} 
\sphinxAtStartPar
B

\item {} 
\sphinxAtStartPar
p

\item {} 
\sphinxAtStartPar
q

\item {} 
\sphinxAtStartPar
tol

\end{itemize}

\sphinxAtStartPar
The A, B, p, and q parameters are used only for two\sphinxhyphen{}body interactions.
The \(\lambda\) and \(\cos\theta_0\) parameters are used only
for three\sphinxhyphen{}body interactions. The \(\epsilon\), \(\sigma\) and
\sphinxstyleemphasis{a} parameters are used for both two\sphinxhyphen{}body and three\sphinxhyphen{}body
interactions. \(\gamma\) is used only in the three\sphinxhyphen{}body
interactions, but is defined for pairs of atoms.  The non\sphinxhyphen{}annotated
parameters are unitless.

\sphinxAtStartPar
LAMMPS introduces an additional performance\sphinxhyphen{}optimization parameter tol
that is used for both two\sphinxhyphen{}body and three\sphinxhyphen{}body interactions.  In the
Stillinger\sphinxhyphen{}Weber potential, the interaction energies become negligibly
small at atomic separations substantially less than the theoretical
cutoff distances.  LAMMPS therefore defines a virtual cutoff distance
based on a user defined tolerance tol.  The use of the virtual cutoff
distance in constructing atom neighbor lists can significantly reduce
the neighbor list sizes and therefore the computational cost.  LAMMPS
provides a \sphinxstyleemphasis{tol} value for each of the three\sphinxhyphen{}body entries so that they
can be separately controlled. If tol = 0.0, then the standard
Stillinger\sphinxhyphen{}Weber cutoff is used.

\sphinxAtStartPar
The Stillinger\sphinxhyphen{}Weber potential file must contain entries for all the
elements listed in the pair\_coeff command.  It can also contain
entries for additional elements not being used in a particular
simulation; LAMMPS ignores those entries.

\sphinxAtStartPar
For a single\sphinxhyphen{}element simulation, only a single entry is required
(e.g. SiSiSi).  For a two\sphinxhyphen{}element simulation, the file must contain 8
entries (for SiSiSi, SiSiC, SiCSi, SiCC, CSiSi, CSiC, CCSi, CCC), that
specify SW parameters for all permutations of the two elements
interacting in three\sphinxhyphen{}body configurations.  Thus for 3 elements, 27
entries would be required, etc.

\sphinxAtStartPar
As annotated above, the first element in the entry is the center atom
in a three\sphinxhyphen{}body interaction.  Thus an entry for SiCC means a Si atom
with 2 C atoms as neighbors.  The parameter values used for the
two\sphinxhyphen{}body interaction come from the entry where the second and third
elements are the same.  Thus the two\sphinxhyphen{}body parameters for Si
interacting with C, comes from the SiCC entry.  The three\sphinxhyphen{}body
parameters can in principle be specific to the three elements of the
configuration. In the literature, however, the three\sphinxhyphen{}body parameters
are usually defined by simple formulas involving two sets of pairwise
parameters, corresponding to the ij and ik pairs, where i is the
center atom. The user must ensure that the correct combining rule is
used to calculate the values of the three\sphinxhyphen{}body parameters for
alloys. Note also that the function \(\phi_3\) contains two exponential
screening factors with parameter values from the ij pair and ik
pairs. So \(\phi_3\) for a C atom bonded to a Si atom and a second C atom
will depend on the three\sphinxhyphen{}body parameters for the CSiC entry, and also
on the two\sphinxhyphen{}body parameters for the CCC and CSiSi entries. Since the
order of the two neighbors is arbitrary, the three\sphinxhyphen{}body parameters for
entries CSiC and CCSi should be the same.  Similarly, the two\sphinxhyphen{}body
parameters for entries SiCC and CSiSi should also be the same.  The
parameters used only for two\sphinxhyphen{}body interactions (A, B, p, and q) in
entries whose second and third element are different (e.g. SiCSi) are not
used for anything and can be set to 0.0 if desired.
This is also true for the parameters in \(\phi_3\) that are
taken from the ij and ik pairs (\(\sigma\), \sphinxstyleemphasis{a}, \(\gamma\))


\bigskip\hrule\bigskip


\sphinxAtStartPar
Styles with a \sphinxstyleemphasis{gpu}, \sphinxstyleemphasis{intel}, \sphinxstyleemphasis{kk}, \sphinxstyleemphasis{omp}, or \sphinxstyleemphasis{opt} suffix are
functionally the same as the corresponding style without the suffix.
They have been optimized to run faster, depending on your available
hardware, as discussed on the {\hyperref[\detokenize{Speed_packages::doc}]{\sphinxcrossref{\DUrole{doc}{Accelerator packages}}}}
page.  The accelerated styles take the same arguments and should
produce the same results, except for round\sphinxhyphen{}off and precision issues.

\sphinxAtStartPar
These accelerated styles are part of the GPU, INTEL, KOKKOS,
OPENMP, and OPT packages, respectively.  They are only enabled if
LAMMPS was built with those packages.  See the {\hyperref[\detokenize{Build_package::doc}]{\sphinxcrossref{\DUrole{doc}{Build package}}}} page for more info.

\sphinxAtStartPar
You can specify the accelerated styles explicitly in your input script
by including their suffix, or you can use the {\hyperref[\detokenize{Run_options::doc}]{\sphinxcrossref{\DUrole{doc}{\sphinxhyphen{}suffix command\sphinxhyphen{}line switch}}}} when you invoke LAMMPS, or you can use the
{\hyperref[\detokenize{suffix::doc}]{\sphinxcrossref{\DUrole{doc}{suffix}}}} command in your input script.

\sphinxAtStartPar
See the {\hyperref[\detokenize{Speed_packages::doc}]{\sphinxcrossref{\DUrole{doc}{Accelerator packages}}}} page for more
instructions on how to use the accelerated styles effectively.

\begin{sphinxadmonition}{note}{Note:}
\sphinxAtStartPar
When using the INTEL package with this style, there is an additional
5 to 10 percent performance improvement when the Stillinger\sphinxhyphen{}Weber
parameters p and q are set to 4 and 0 respectively.  These
parameters are common for modeling silicon and water.
\end{sphinxadmonition}


\bigskip\hrule\bigskip



\subsection{Mixing, shift, table, tail correction, restart, rRESPA info}
\label{\detokenize{pair_sw:mixing-shift-table-tail-correction-restart-rrespa-info}}
\sphinxAtStartPar
For atom type pairs I,J and I != J, where types I and J correspond to
two different element types, mixing is performed by LAMMPS as
described above from values in the potential file.

\sphinxAtStartPar
This pair style does not support the {\hyperref[\detokenize{pair_modify::doc}]{\sphinxcrossref{\DUrole{doc}{pair\_modify}}}}
shift, table, and tail options.

\sphinxAtStartPar
This pair style does not write its information to {\hyperref[\detokenize{restart::doc}]{\sphinxcrossref{\DUrole{doc}{binary restart
files}}}}, since it is stored in potential files.  Thus, you need
to re\sphinxhyphen{}specify the pair\_style and pair\_coeff commands in an input script
that reads a restart file.

\sphinxAtStartPar
This pair style can only be used via the \sphinxstyleemphasis{pair} keyword of the
{\hyperref[\detokenize{run_style::doc}]{\sphinxcrossref{\DUrole{doc}{run\_style respa}}}} command.  It does not support the
\sphinxstyleemphasis{inner}, \sphinxstyleemphasis{middle}, \sphinxstyleemphasis{outer} keywords.

\sphinxAtStartPar
The single() function of the \sphinxstyleemphasis{sw} pair style is only enabled and
supported for the case of the \sphinxstyleemphasis{threebody off} setting.


\bigskip\hrule\bigskip



\subsection{Restrictions}
\label{\detokenize{pair_sw:restrictions}}
\sphinxAtStartPar
This pair style is part of the MANYBODY package.  It is only enabled if
LAMMPS was built with that package.  See the {\hyperref[\detokenize{Build_package::doc}]{\sphinxcrossref{\DUrole{doc}{Build package}}}} page for more info.

\sphinxAtStartPar
This pair style requires the {\hyperref[\detokenize{newton::doc}]{\sphinxcrossref{\DUrole{doc}{newton}}}} setting to be “on”
for pair interactions.

\sphinxAtStartPar
The Stillinger\sphinxhyphen{}Weber potential files provided with LAMMPS (see the
potentials directory) are parameterized for metal {\hyperref[\detokenize{units::doc}]{\sphinxcrossref{\DUrole{doc}{units}}}}.
You can use the sw or sw/mod pair styles with any LAMMPS units, but you
would need to create your own SW potential file with coefficients listed
in the appropriate units if your simulation does not use “metal” units.
If the potential file contains a ‘UNITS:’ metadata tag in the first line
of the potential file, then LAMMPS can convert it transparently between
“metal” and “real” units.


\subsection{Related commands}
\label{\detokenize{pair_sw:related-commands}}
\sphinxAtStartPar
{\hyperref[\detokenize{pair_coeff::doc}]{\sphinxcrossref{\DUrole{doc}{pair\_coeff}}}}


\subsection{Default}
\label{\detokenize{pair_sw:default}}
\sphinxAtStartPar
The default value for the \sphinxstyleemphasis{threebody} setting of the “sw” pair style is
“on”, the default values for the “\sphinxstyleemphasis{maxdelcs} setting of the \sphinxstyleemphasis{sw/mod}
pair style are \sphinxstyleemphasis{delta1} = 0.25 and \sphinxstyleemphasis{delta2} = 0.35\textasciigrave{}.


\bigskip\hrule\bigskip

\phantomsection\label{\detokenize{pair_sw:stillinger2}}
\sphinxAtStartPar
\sphinxstylestrong{(Stillinger)} Stillinger and Weber, Phys Rev B, 31, 5262 (1985).

\phantomsection\label{\detokenize{pair_sw:jiang2}}
\sphinxAtStartPar
\sphinxstylestrong{(Jiang2)} J.\sphinxhyphen{}W. Jiang, Nanotechnology 26, 315706 (2015).

\phantomsection\label{\detokenize{pair_sw:jiang3}}
\sphinxAtStartPar
\sphinxstylestrong{(Jiang3)} J.\sphinxhyphen{}W. Jiang, Acta Mech. Solida. Sin 32, 17 (2019).

\sphinxstepscope

\index{pair\_style sw/angle/table@\spxentry{pair\_style sw/angle/table}}\ignorespaces 

\section{pair\_style sw/angle/table command}
\label{\detokenize{pair_sw_angle_table:pair-style-sw-angle-table-command}}\label{\detokenize{pair_sw_angle_table:index-0}}\label{\detokenize{pair_sw_angle_table::doc}}

\subsection{Syntax}
\label{\detokenize{pair_sw_angle_table:syntax}}
\begin{sphinxVerbatim}[commandchars=\\\{\}]
\PYG{k}{pair\PYGZus{}style}\PYG{+w}{ }\PYG{n}{style}
\end{sphinxVerbatim}
\begin{itemize}
\item {} 
\sphinxAtStartPar
style = \sphinxstyleemphasis{sw/angle/table}

\end{itemize}


\subsection{Examples}
\label{\detokenize{pair_sw_angle_table:examples}}
\begin{sphinxVerbatim}[commandchars=\\\{\}]
\PYG{k}{pair\PYGZus{}style}\PYG{+w}{ }\PYG{n}{sw}\PYG{o}{/}\PYG{n}{angle}\PYG{o}{/}\PYG{n}{table}
\PYG{k}{pair\PYGZus{}coeff}\PYG{+w}{ }\PYG{o}{*}\PYG{+w}{ }\PYG{o}{*}\PYG{+w}{ }\PYG{n}{spce.sw}\PYG{+w}{ }\PYG{n}{type}
\end{sphinxVerbatim}

\sphinxAtStartPar
Used in example input script:

\begin{sphinxVerbatim}[commandchars=\\\{\}]
examples/PACKAGES/manybody\PYGZus{}table/in.spce\PYGZus{}sw
\end{sphinxVerbatim}


\subsection{Description}
\label{\detokenize{pair_sw_angle_table:description}}
\sphinxAtStartPar
\DUrole{versionmodified,added}{Added in version 2Jun2022.}

\sphinxAtStartPar
The \sphinxstyleemphasis{sw/angle/table} style is a modification of the original
{\hyperref[\detokenize{pair_sw::doc}]{\sphinxcrossref{\DUrole{doc}{pair\_style sw}}}}. It has been developed for coarse\sphinxhyphen{}grained
simulations (of water) ({\hyperref[\detokenize{pair_sw_angle_table:scherer1}]{\sphinxcrossref{\DUrole{std,std-ref}{Scherer1}}}}), but can be employed
for all kinds of systems. It computes a modified 3\sphinxhyphen{}body
{\hyperref[\detokenize{pair_sw_angle_table:stillinger3}]{\sphinxcrossref{\DUrole{std,std-ref}{Stillinger\sphinxhyphen{}Weber}}}} potential for the energy E of a
system of atoms as
\begin{equation*}
\begin{split} E & =  \sum_i \sum_{j > i} \phi_2 (r_{ij}) +
        \sum_i \sum_{j \neq i} \sum_{k > j}
        \phi_3 (r_{ij}, r_{ik}, \theta_{ijk}) \\
\phi_2(r_{ij}) & =  A_{ij} \epsilon_{ij} \left[ B_{ij} (\frac{\sigma_{ij}}{r_{ij}})^{p_{ij}} -
                  (\frac{\sigma_{ij}}{r_{ij}})^{q_{ij}} \right]
                  \exp \left( \frac{\sigma_{ij}}{r_{ij} - a_{ij} \sigma_{ij}} \right) \\
\phi_3(r_{ij},r_{ik},\theta_{ijk}) & = f^{\textrm{3b}}\left(\theta_{ijk}\right)
                  \exp \left( \frac{\gamma_{ij} \sigma_{ij}}{r_{ij} - a_{ij} \sigma_{ij}} \right)
                  \exp \left( \frac{\gamma_{ik} \sigma_{ik}}{r_{ik} - a_{ik} \sigma_{ik}} \right)\end{split}
\end{equation*}
\sphinxAtStartPar
where \(\phi_2\) is a two\sphinxhyphen{}body term and \(\phi_3\) is a
three\sphinxhyphen{}body term.  The summations in the formula are over all neighbors J
and K of atom I within a cutoff distance \(a \sigma\).  In contrast
to the original \sphinxstyleemphasis{sw} style, \sphinxstyleemphasis{sw/angle/table} allows for a flexible
three\sphinxhyphen{}body term \(f^{\textrm{3b}}\left(\theta_{ijk}\right)\) which
is read in as a tabulated interaction. It can be parameterized with the
csg\_fmatch app of VOTCA as available at:
\sphinxurl{https://gitlab.mpcdf.mpg.de/votca/votca}.

\sphinxAtStartPar
Only a single pair\_coeff command is used with the \sphinxstyleemphasis{sw/angle/table} style
which specifies a modified Stillinger\sphinxhyphen{}Weber potential file with
parameters for all needed elements.  These are mapped to LAMMPS atom
types by specifying N\_el additional arguments after the “.sw” filename
in the pair\_coeff command, where N\_el is the number of LAMMPS atom
types:
\begin{itemize}
\item {} 
\sphinxAtStartPar
“.sw” filename

\item {} 
\sphinxAtStartPar
N\_el element names = mapping of SW elements to atom types

\end{itemize}

\sphinxAtStartPar
See the {\hyperref[\detokenize{pair_coeff::doc}]{\sphinxcrossref{\DUrole{doc}{pair\_coeff}}}} page for alternate ways to
specify the path for the potential file.

\sphinxAtStartPar
As an example, imagine a file SiC.sw has Stillinger\sphinxhyphen{}Weber values for Si
and C.  If your LAMMPS simulation has 4 atoms types and you want the
first 3 to be Si, and the fourth to be C, you would use the following
pair\_coeff command:

\begin{sphinxVerbatim}[commandchars=\\\{\}]
\PYG{k}{pair\PYGZus{}coeff}\PYG{+w}{ }\PYG{o}{*}\PYG{+w}{ }\PYG{o}{*}\PYG{+w}{ }\PYG{n}{SiC.sw}\PYG{+w}{ }\PYG{n}{Si}\PYG{+w}{ }\PYG{n}{Si}\PYG{+w}{ }\PYG{n}{Si}\PYG{+w}{ }\PYG{n}{C}
\end{sphinxVerbatim}

\sphinxAtStartPar
The first 2 arguments must be * * so as to span all LAMMPS atom types.
The first three Si arguments map LAMMPS atom types 1,2,3 to the Si
element in the SW file.  The final C argument maps LAMMPS atom type 4 to
the C element in the SW file.  If a mapping value is specified as NULL,
the mapping is not performed.  This can be used when a \sphinxstyleemphasis{sw/angle/table}
potential is used as part of the \sphinxstyleemphasis{hybrid} pair style.  The NULL values
are placeholders for atom types that will be used with other potentials.

\sphinxAtStartPar
The (modified) Stillinger\sphinxhyphen{}Weber files have a “.sw” suffix. Lines that
are not blank or comments (starting with \#) define parameters for a
triplet of elements. The parameters in a single entry correspond to the
two\sphinxhyphen{}body and three\sphinxhyphen{}body coefficients in the formula above. Here, also
the suffix “.sw” is used though the original Stillinger\sphinxhyphen{}Weber file
format is supplemented with four additional lines per parameter block to
specify the tabulated three\sphinxhyphen{}body interaction. A single entry then
contains:
\begin{itemize}
\item {} 
\sphinxAtStartPar
element 1 (the center atom in a 3\sphinxhyphen{}body interaction)

\item {} 
\sphinxAtStartPar
element 2

\item {} 
\sphinxAtStartPar
element 3

\item {} 
\sphinxAtStartPar
\(\epsilon\) (energy units)

\item {} 
\sphinxAtStartPar
\(\sigma\) (distance units)

\item {} 
\sphinxAtStartPar
a

\item {} 
\sphinxAtStartPar
\(\lambda\)

\item {} 
\sphinxAtStartPar
\(\gamma\)

\item {} 
\sphinxAtStartPar
\(\cos\theta_0\)

\item {} 
\sphinxAtStartPar
A

\item {} 
\sphinxAtStartPar
B

\item {} 
\sphinxAtStartPar
p

\item {} 
\sphinxAtStartPar
q

\item {} 
\sphinxAtStartPar
tol

\item {} 
\sphinxAtStartPar
filename

\item {} 
\sphinxAtStartPar
keyword

\item {} 
\sphinxAtStartPar
style

\item {} 
\sphinxAtStartPar
N

\end{itemize}

\sphinxAtStartPar
The A, B, p, and q parameters are used only for two\sphinxhyphen{}body interactions.
The \(\lambda\) and \(\cos\theta_0\) parameters, only used for
three\sphinxhyphen{}body interactions in the original Stillinger\sphinxhyphen{}Weber style, are read
in but ignored in this modified pair style. The \(\epsilon\)
parameter is only used for two\sphinxhyphen{}body interactions in this modified pair
style and not for the three\sphinxhyphen{}body terms. The \(\sigma\) and \sphinxstyleemphasis{a}
parameters are used for both two\sphinxhyphen{}body and three\sphinxhyphen{}body
interactions. \(\gamma\) is used only in the three\sphinxhyphen{}body
interactions, but is defined for pairs of atoms. The non\sphinxhyphen{}annotated
parameters are unitless.

\sphinxAtStartPar
LAMMPS introduces an additional performance\sphinxhyphen{}optimization parameter tol
that is used for both two\sphinxhyphen{}body and three\sphinxhyphen{}body interactions.  In the
Stillinger\sphinxhyphen{}Weber potential, the interaction energies become negligibly
small at atomic separations substantially less than the theoretical
cutoff distances.  LAMMPS therefore defines a virtual cutoff distance
based on a user defined tolerance tol.  The use of the virtual cutoff
distance in constructing atom neighbor lists can significantly reduce
the neighbor list sizes and therefore the computational cost.  LAMMPS
provides a \sphinxstyleemphasis{tol} value for each of the three\sphinxhyphen{}body entries so that they
can be separately controlled. If tol = 0.0, then the standard
Stillinger\sphinxhyphen{}Weber cutoff is used.

\sphinxAtStartPar
The additional parameters \sphinxstyleemphasis{filename}, \sphinxstyleemphasis{keyword}, \sphinxstyleemphasis{style}, and \sphinxstyleemphasis{N} refer
to the tabulated angular potential
\(f^{\textrm{3b}}\left(\theta_{ijk}\right)\).  The tabulated angular
potential has to be of the format as used in the {\hyperref[\detokenize{angle_table::doc}]{\sphinxcrossref{\DUrole{doc}{angle\_style table}}}} command:

\sphinxAtStartPar
An interpolation tables of length \sphinxstyleemphasis{N} is created. The interpolation is
done in one of 2 \sphinxstyleemphasis{styles}: \sphinxstyleemphasis{linear} or \sphinxstyleemphasis{spline}.  For the \sphinxstyleemphasis{linear}
style, the angle is used to find 2 surrounding table values from which
an energy or its derivative is computed by linear interpolation. For the
\sphinxstyleemphasis{spline} style, a cubic spline coefficients are computed and stored at
each of the \sphinxstyleemphasis{N} values in the table.  The angle is used to find the
appropriate set of coefficients which are used to evaluate a cubic
polynomial which computes the energy or derivative.

\sphinxAtStartPar
The \sphinxstyleemphasis{filename} specifies the file containing the tabulated energy and
derivative values of \(f^{\textrm{3b}}\left(\theta_{ijk}\right)\).
The \sphinxstyleemphasis{keyword} then specifies a section of the file.  The format of this
file is as follows (without the parenthesized comments):

\begin{sphinxVerbatim}[commandchars=\\\{\}]
\PYGZsh{} Angle potential for harmonic (one or more comment or blank lines)

HAM                           (keyword is the first text on line)
N 181 FP 0 0 EQ 90.0          (N, FP, EQ parameters)
                              (blank line)
1 0.0 200.5 2.5               (index, angle, energy, derivative)
2 1.0 198.0 2.5
...
181 180.0 0.0 0.0
\end{sphinxVerbatim}

\sphinxAtStartPar
A section begins with a non\sphinxhyphen{}blank line whose first character is not a
“\#”; blank lines or lines starting with “\#” can be used as comments
between sections.  The first line begins with a keyword which identifies
the section. The next line lists (in any order) one or more parameters
for the table.  Each parameter is a keyword followed by one or more
numeric values.

\sphinxAtStartPar
The parameter “N” is required and its value is the number of table
entries that follow.  Note that this may be different than the \sphinxstyleemphasis{N}
specified in the Stillinger\sphinxhyphen{}Weber potential file. Let Nsw = \sphinxstyleemphasis{N} in the
“.sw” file, and Nfile = “N” in the tabulated angular file.  What LAMMPS
does is a preliminary interpolation by creating splines using the Nfile
tabulated values as nodal points.  It uses these to interpolate as
needed to generate energy and derivative values at Ntable different
points.  The resulting tables of length Nsw are then used as described
above, when computing energy and force for individual angles and their
atoms.  This means that if you want the interpolation tables of length
Nsw to match exactly what is in the tabulated file (with effectively no
preliminary interpolation), you should set Nsw = Nfile.

\sphinxAtStartPar
The “FP” parameter is optional.  If used, it is followed by two values
fplo and fphi, which are the second derivatives at the innermost and
outermost angle settings.  These values are needed by the spline
construction routines.  If not specified by the “FP” parameter, they are
estimated (less accurately) by the first two and last two derivative
values in the table.

\sphinxAtStartPar
The “EQ” parameter is also optional.  If used, it is followed by a the
equilibrium angle value, which is used, for example, by the {\hyperref[\detokenize{fix_shake::doc}]{\sphinxcrossref{\DUrole{doc}{fix
shake}}}} command. If not used, the equilibrium angle is set to
180.0.

\sphinxAtStartPar
Following a blank line, the next N lines of the angular table file list
the tabulated values.  On each line, the first value is the index from 1
to N, the second value is the angle value (in degrees), the third value
is the energy (in energy units), and the fourth is \sphinxhyphen{}dE/d(theta) (also in
energy units).  The third term is the energy of the 3\sphinxhyphen{}atom configuration
for the specified angle.  The last term is the derivative of the energy
with respect to the angle (in degrees, not radians).  Thus the units of
the last term are still energy, not force.  The angle values must
increase from one line to the next.  The angle values must also begin
with 0.0 and end with 180.0, i.e. span the full range of possible
angles.

\sphinxAtStartPar
Note that one angular potential file can contain many sections, each
with a tabulated potential.  LAMMPS reads the file section by section
until it finds one that matches the specified \sphinxstyleemphasis{keyword} of appropriate
section of the “.sw” file.

\sphinxAtStartPar
The Stillinger\sphinxhyphen{}Weber potential file must contain entries for all the
elements listed in the pair\_coeff command.  It can also contain entries
for additional elements not being used in a particular simulation;
LAMMPS ignores those entries.

\sphinxAtStartPar
For a single\sphinxhyphen{}element simulation, only a single entry is required
(e.g. SiSiSi).  For a two\sphinxhyphen{}element simulation, the file must contain 8
entries (for SiSiSi, SiSiC, SiCSi, SiCC, CSiSi, CSiC, CCSi, CCC), that
specify SW parameters for all permutations of the two elements
interacting in three\sphinxhyphen{}body configurations.  Thus for 3 elements, 27
entries would be required, etc.

\sphinxAtStartPar
As annotated above, the first element in the entry is the center atom in
a three\sphinxhyphen{}body interaction.  Thus an entry for SiCC means a Si atom with 2
C atoms as neighbors.  The parameter values used for the two\sphinxhyphen{}body
interaction come from the entry where the second and third elements are
the same.  Thus the two\sphinxhyphen{}body parameters for Si interacting with C, comes
from the SiCC entry.  The three\sphinxhyphen{}body angular potential
\(f^{\textrm{3b}}\left(\theta_{ijk}\right)\) can in principle be
specific to the three elements of the configuration. However, the user
must ensure that it makes physically sense.  Note also that the function
\(\phi_3\) contains two exponential screening factors with parameter
values from the ij pair and ik pairs. So \(\phi_3\) for a C atom
bonded to a Si atom and a second C atom will depend on the three\sphinxhyphen{}body
parameters for the CSiC entry, and also on the two\sphinxhyphen{}body parameters for
the CCC and CSiSi entries. Since the order of the two neighbors is
arbitrary, the three\sphinxhyphen{}body parameters and the tabulated angular potential
for entries CSiC and CCSi should be the same.  Similarly, the two\sphinxhyphen{}body
parameters for entries SiCC and CSiSi should also be the same.  The
parameters used only for two\sphinxhyphen{}body interactions (A, B, p, and q) in
entries whose second and third element are different (e.g. SiCSi) are
not used for anything and can be set to 0.0 if desired.  This is also
true for the parameters in \(\phi_3\) that are taken from the ij and
ik pairs (\(\sigma\), \sphinxstyleemphasis{a}, \(\gamma\))

\sphinxAtStartPar
Additional input files and reference data can be found at:
\sphinxurl{https://gitlab.mpcdf.mpg.de/votca/votca/-/tree/master/csg-tutorials/spce/3body\_sw}


\bigskip\hrule\bigskip



\subsection{Mixing, shift, table, tail correction, restart, rRESPA info}
\label{\detokenize{pair_sw_angle_table:mixing-shift-table-tail-correction-restart-rrespa-info}}
\sphinxAtStartPar
For atom type pairs I,J and I != J, where types I and J correspond to
two different element types, mixing is performed by LAMMPS as described
above from values in the potential file, but not for the tabulated
angular potential file.

\sphinxAtStartPar
This pair style does not support the {\hyperref[\detokenize{pair_modify::doc}]{\sphinxcrossref{\DUrole{doc}{pair\_modify}}}}
shift, table, and tail options.

\sphinxAtStartPar
This pair style does not write its information to {\hyperref[\detokenize{restart::doc}]{\sphinxcrossref{\DUrole{doc}{binary restart
files}}}}, since it is stored in potential files.  Thus, you need
to re\sphinxhyphen{}specify the pair\_style and pair\_coeff commands in an input script
that reads a restart file.

\sphinxAtStartPar
This pair style can only be used via the \sphinxstyleemphasis{pair} keyword of the
{\hyperref[\detokenize{run_style::doc}]{\sphinxcrossref{\DUrole{doc}{run\_style respa}}}} command.  It does not support the
\sphinxstyleemphasis{inner}, \sphinxstyleemphasis{middle}, \sphinxstyleemphasis{outer} keywords.


\bigskip\hrule\bigskip



\subsection{Restrictions}
\label{\detokenize{pair_sw_angle_table:restrictions}}
\sphinxAtStartPar
This pair style is part of the MANYBODY package.  It is only enabled if
LAMMPS was built with that package. See the {\hyperref[\detokenize{Build_package::doc}]{\sphinxcrossref{\DUrole{doc}{Build package}}}} page for more info.

\sphinxAtStartPar
This pair style requires the {\hyperref[\detokenize{newton::doc}]{\sphinxcrossref{\DUrole{doc}{newton}}}} setting to be “on”
for pair interactions.


\subsection{Related commands}
\label{\detokenize{pair_sw_angle_table:related-commands}}
\sphinxAtStartPar
{\hyperref[\detokenize{pair_coeff::doc}]{\sphinxcrossref{\DUrole{doc}{pair\_coeff}}}}, {\hyperref[\detokenize{pair_sw::doc}]{\sphinxcrossref{\DUrole{doc}{pair\_style sw}}}},
{\hyperref[\detokenize{pair_threebody_table::doc}]{\sphinxcrossref{\DUrole{doc}{pair\_style threebody/table}}}}


\bigskip\hrule\bigskip

\phantomsection\label{\detokenize{pair_sw_angle_table:stillinger3}}
\sphinxAtStartPar
\sphinxstylestrong{(Stillinger)} Stillinger and Weber, Phys Rev B, 31, 5262 (1985).

\phantomsection\label{\detokenize{pair_sw_angle_table:scherer1}}
\sphinxAtStartPar
\sphinxstylestrong{(Scherer1)} C. Scherer and D. Andrienko, Phys. Chem. Chem. Phys. 20, 22387\sphinxhyphen{}22394 (2018).

\sphinxstepscope

\index{pair\_style table@\spxentry{pair\_style table}}\ignorespaces 
\index{pair\_style table/gpu@\spxentry{pair\_style table/gpu}}\ignorespaces 
\index{pair\_style table/kk@\spxentry{pair\_style table/kk}}\ignorespaces 
\index{pair\_style table/omp@\spxentry{pair\_style table/omp}}\ignorespaces 

\section{pair\_style table command}
\label{\detokenize{pair_table:pair-style-table-command}}\label{\detokenize{pair_table:index-3}}\label{\detokenize{pair_table:index-2}}\label{\detokenize{pair_table:index-1}}\label{\detokenize{pair_table:index-0}}\label{\detokenize{pair_table::doc}}
\sphinxAtStartPar
Accelerator Variants: \sphinxstyleemphasis{table/gpu}, \sphinxstyleemphasis{table/kk}, \sphinxstyleemphasis{table/omp}


\subsection{Syntax}
\label{\detokenize{pair_table:syntax}}
\begin{sphinxVerbatim}[commandchars=\\\{\}]
\PYG{k}{pair\PYGZus{}style}\PYG{+w}{ }\PYG{n}{table}\PYG{+w}{ }\PYG{n}{style}\PYG{+w}{ }\PYG{n}{N}\PYG{+w}{ }\PYG{n}{keyword}\PYG{+w}{ }\PYG{n}{...}
\end{sphinxVerbatim}
\begin{itemize}
\item {} 
\sphinxAtStartPar
style = \sphinxstyleemphasis{lookup} or \sphinxstyleemphasis{linear} or \sphinxstyleemphasis{spline} or \sphinxstyleemphasis{bitmap} = method of interpolation

\item {} 
\sphinxAtStartPar
N = use N values in \sphinxstyleemphasis{lookup}, \sphinxstyleemphasis{linear}, \sphinxstyleemphasis{spline} tables

\item {} 
\sphinxAtStartPar
N = use 2\textasciicircum{}N values in \sphinxstyleemphasis{bitmap} tables

\item {} 
\sphinxAtStartPar
zero or more keywords may be appended

\item {} 
\sphinxAtStartPar
keyword = \sphinxstyleemphasis{ewald} or \sphinxstyleemphasis{pppm} or \sphinxstyleemphasis{msm} or \sphinxstyleemphasis{dispersion} or \sphinxstyleemphasis{tip4p}

\end{itemize}


\subsection{Examples}
\label{\detokenize{pair_table:examples}}
\begin{sphinxVerbatim}[commandchars=\\\{\}]
\PYG{k}{pair\PYGZus{}style}\PYG{+w}{ }\PYG{n}{table}\PYG{+w}{ }\PYG{n}{linear}\PYG{+w}{ }\PYG{l+m}{1000}
\PYG{k}{pair\PYGZus{}style}\PYG{+w}{ }\PYG{n}{table}\PYG{+w}{ }\PYG{n}{linear}\PYG{+w}{ }\PYG{l+m}{1000}\PYG{+w}{ }\PYG{n}{pppm}
\PYG{k}{pair\PYGZus{}style}\PYG{+w}{ }\PYG{n}{table}\PYG{+w}{ }\PYG{n}{bitmap}\PYG{+w}{ }\PYG{l+m}{12}
\PYG{k}{pair\PYGZus{}coeff}\PYG{+w}{ }\PYG{o}{*}\PYG{+w}{ }\PYG{l+m}{3}\PYG{+w}{ }\PYG{n}{morse.table}\PYG{+w}{ }\PYG{n}{ENTRY1}
\PYG{k}{pair\PYGZus{}coeff}\PYG{+w}{ }\PYG{o}{*}\PYG{+w}{ }\PYG{l+m}{3}\PYG{+w}{ }\PYG{n}{morse.table}\PYG{+w}{ }\PYG{n}{ENTRY1}\PYG{+w}{ }\PYG{l+m}{7.0}
\end{sphinxVerbatim}


\subsection{Description}
\label{\detokenize{pair_table:description}}
\sphinxAtStartPar
Style \sphinxstyleemphasis{table} creates interpolation tables from potential energy and
force values listed in a file(s) as a function of distance.  When
performing dynamics or minimization, the interpolation tables are used
to evaluate energy and forces for pairwise interactions between
particles, similar to how analytic formulas are used for other pair
styles.

\sphinxAtStartPar
The interpolation tables are created as a pre\sphinxhyphen{}computation by fitting
cubic splines to the file values and interpolating energy and force
values at each of \sphinxstyleemphasis{N} distances.  During a simulation, the tables are
used to interpolate energy and force values as needed for each pair of
particles separated by a distance \sphinxstyleemphasis{R}.  The interpolation is done in
one of 4 styles: \sphinxstyleemphasis{lookup}, \sphinxstyleemphasis{linear}, \sphinxstyleemphasis{spline}, or \sphinxstyleemphasis{bitmap}.

\sphinxAtStartPar
For the \sphinxstyleemphasis{lookup} style, the distance \sphinxstyleemphasis{R} is used to find the nearest
table entry, which is the energy or force.

\sphinxAtStartPar
For the \sphinxstyleemphasis{linear} style, the distance \sphinxstyleemphasis{R} is used to find the 2
surrounding table values from which an energy or force is computed by
linear interpolation.

\sphinxAtStartPar
For the \sphinxstyleemphasis{spline} style, cubic spline coefficients are computed and
stored for each of the \sphinxstyleemphasis{N} values in the table, one set of splines for
energy, another for force.  Note that these splines are different than
the ones used to pre\sphinxhyphen{}compute the \sphinxstyleemphasis{N} values.  Those splines were fit
to the \sphinxstyleemphasis{Nfile} values in the tabulated file, where often \sphinxstyleemphasis{Nfile} \textless{}
\sphinxstyleemphasis{N}.  The distance \sphinxstyleemphasis{R} is used to find the appropriate set of spline
coefficients which are used to evaluate a cubic polynomial which
computes the energy or force.

\sphinxAtStartPar
For the \sphinxstyleemphasis{bitmap} style, the specified \sphinxstyleemphasis{N} is used to create
interpolation tables that are 2\textasciicircum{}N in length.  The distance \sphinxstyleemphasis{R} is used
to index into the table via a fast bit\sphinxhyphen{}mapping technique due to
{\hyperref[\detokenize{pair_table:wolff2}]{\sphinxcrossref{\DUrole{std,std-ref}{(Wolff)}}}}, and a linear interpolation is performed between
adjacent table values.

\sphinxAtStartPar
The following coefficients must be defined for each pair of atoms
types via the {\hyperref[\detokenize{pair_coeff::doc}]{\sphinxcrossref{\DUrole{doc}{pair\_coeff}}}} command as in the examples
above.
\begin{itemize}
\item {} 
\sphinxAtStartPar
filename

\item {} 
\sphinxAtStartPar
keyword

\item {} 
\sphinxAtStartPar
cutoff (distance units)

\end{itemize}

\sphinxAtStartPar
The filename specifies a file containing tabulated energy and force
values.  The keyword specifies a section of the file.  The cutoff is
an optional coefficient.  If not specified, the outer cutoff in the
table itself (see below) will be used to build an interpolation table
that extend to the largest tabulated distance.  If specified, only
file values up to the cutoff are used to create the interpolation
table.  The format of this file is described below.

\sphinxAtStartPar
If your tabulated potential(s) are designed to be used as the
short\sphinxhyphen{}range part of one of the long\sphinxhyphen{}range solvers specified by the
{\hyperref[\detokenize{kspace_style::doc}]{\sphinxcrossref{\DUrole{doc}{kspace\_style}}}} command, then you must use one or
more of the optional keywords listed above for the pair\_style command.
These are \sphinxstyleemphasis{ewald} or \sphinxstyleemphasis{pppm} or \sphinxstyleemphasis{msm} or \sphinxstyleemphasis{dispersion} or \sphinxstyleemphasis{tip4p}.  This
is so LAMMPS can ensure the short\sphinxhyphen{}range potential and long\sphinxhyphen{}range
solver are compatible with each other, as it does for other
short\sphinxhyphen{}range pair styles, such as {\hyperref[\detokenize{pair_lj_cut_coul::doc}]{\sphinxcrossref{\DUrole{doc}{pair\_style lj/cut/coul/long}}}}.  Note that it is up to you to ensure
the tabulated values for each pair of atom types has the correct
functional form to be compatible with the matching long\sphinxhyphen{}range solver.


\bigskip\hrule\bigskip


\sphinxAtStartPar
Here are some guidelines for using the pair\_style table command to
best effect:
\begin{itemize}
\item {} 
\sphinxAtStartPar
Vary the number of table points; you may need to use more than you think
to get good resolution.

\item {} 
\sphinxAtStartPar
Always use the {\hyperref[\detokenize{pair_write::doc}]{\sphinxcrossref{\DUrole{doc}{pair\_write}}}} command to produce a plot
of what the final interpolated potential looks like.  This can show up
interpolation “features” you may not like.

\item {} 
\sphinxAtStartPar
Start with the linear style; it’s the style least likely to have problems.

\item {} 
\sphinxAtStartPar
Use \sphinxstyleemphasis{N} in the pair\_style command equal to the “N” in the tabulation
file, and use the “RSQ” or “BITMAP” parameter, so additional interpolation
is not needed.  See discussion below.

\item {} 
\sphinxAtStartPar
Make sure that your tabulated forces and tabulated energies are
consistent (dE/dr = \sphinxhyphen{}F) over the entire range of r values.  LAMMPS
will warn if this is not the case.

\item {} 
\sphinxAtStartPar
Use as large an inner cutoff as possible.  This avoids fitting splines
to very steep parts of the potential.

\end{itemize}


\bigskip\hrule\bigskip


\sphinxAtStartPar
Suitable tables in the correct format for use with these pair styles can
be created by LAMMPS itself using the {\hyperref[\detokenize{pair_write::doc}]{\sphinxcrossref{\DUrole{doc}{pair\_write}}}}
command.  In combination with the pair styles {\hyperref[\detokenize{pair_python::doc}]{\sphinxcrossref{\DUrole{doc}{python}}}},
{\hyperref[\detokenize{pair_lepton::doc}]{\sphinxcrossref{\DUrole{doc}{lepton, or lepton/coul}}}} this can be a powerful
mechanism to implement and test tables for use with LAMMPS.  Another
option to generate tables is the Python code in the \sphinxcode{\sphinxupquote{tools/tabulate}}
folder of the LAMMPS source code distribution.

\sphinxAtStartPar
The format of a tabulated file has an (optional) header followed by a
series of one or more sections, defined as follows (without the
parenthesized comments). The header must start with a \sphinxtitleref{\#} character
and the DATE: and UNITS: tags will be parsed and used:

\begin{sphinxVerbatim}[commandchars=\\\{\}]
\PYGZsh{} DATE: 2020\PYGZhy{}06\PYGZhy{}10  UNITS: real  CONTRIBUTOR: ... (header line)
\PYGZsh{} Morse potential for Fe   (one or more comment or blank lines)

MORSE\PYGZus{}FE                   (keyword is first text on line)
N 500 R 1.0 10.0           (N, R, RSQ, BITMAP, FPRIME parameters)
                           (blank)
1 1.0 25.5 102.34          (index, r, energy, force)
2 1.02 23.4 98.5
...
500 10.0 0.001 0.003
\end{sphinxVerbatim}

\sphinxAtStartPar
A section begins with a non\sphinxhyphen{}blank line whose first character is not a
“\#”; blank lines or lines starting with “\#” can be used as comments
between sections.  The first line begins with a keyword which
identifies the section.  The line can contain additional text, but the
initial text must match the argument specified in the pair\_coeff
command.  The next line lists (in any order) one or more parameters
for the table.  Each parameter is a keyword followed by one or more
numeric values.

\sphinxAtStartPar
The parameter “N” is required and its value is the number of table
entries that follow.  Note that this may be different than the \sphinxstyleemphasis{N}
specified in the {\hyperref[\detokenize{pair_style::doc}]{\sphinxcrossref{\DUrole{doc}{pair\_style table}}}} command.  Let
Ntable = \sphinxstyleemphasis{N} in the pair\_style command, and Nfile = “N” in the
tabulated file.  What LAMMPS does is a preliminary interpolation by
creating splines using the Nfile tabulated values as nodal points.  It
uses these to interpolate energy and force values at Ntable different
points.  The resulting tables of length Ntable are then used as
described above, when computing energy and force for individual pair
distances.  This means that if you want the interpolation tables of
length Ntable to match exactly what is in the tabulated file (with
effectively no preliminary interpolation), you should set Ntable =
Nfile, and use the “RSQ” or “BITMAP” parameter.  This is because the
internal table abscissa is always RSQ (separation distance squared),
for efficient lookup.

\sphinxAtStartPar
All other parameters are optional.  If “R” or “RSQ” or “BITMAP” does
not appear, then the distances in each line of the table are used
as\sphinxhyphen{}is to perform spline interpolation.  In this case, the table values
can be spaced in \sphinxstyleemphasis{r} uniformly or however you wish to position table
values in regions of large gradients.

\sphinxAtStartPar
If used, the parameters “R” or “RSQ” are followed by 2 values \sphinxstyleemphasis{rlo}
and \sphinxstyleemphasis{rhi}.  If specified, the distance associated with each energy and
force value is computed from these 2 values (at high accuracy), rather
than using the (low\sphinxhyphen{}accuracy) value listed in each line of the table.
The distance values in the table file are ignored in this case.
For “R”, distances uniformly spaced between \sphinxstyleemphasis{rlo} and \sphinxstyleemphasis{rhi} are
computed; for “RSQ”, squared distances uniformly spaced between
\sphinxstyleemphasis{rlo*rlo} and \sphinxstyleemphasis{rhi*rhi} are computed.

\begin{sphinxadmonition}{note}{Note:}
\sphinxAtStartPar
If you use “R” or “RSQ”, the tabulated distance values in the
file are effectively ignored, and replaced by new values as described
in the previous paragraph.  If the distance value in the table is not
very close to the new value (i.e. round\sphinxhyphen{}off difference), then you will
be assigning energy/force values to a different distance, which is
probably not what you want.  LAMMPS will warn if this is occurring.
\end{sphinxadmonition}

\sphinxAtStartPar
If used, the parameter “BITMAP” is also followed by 2 values \sphinxstyleemphasis{rlo} and
\sphinxstyleemphasis{rhi}.  These values, along with the “N” value determine the ordering
of the N lines that follow and what distance is associated with each.
This ordering is complex, so it is not documented here, since this
file is typically produced by the {\hyperref[\detokenize{pair_write::doc}]{\sphinxcrossref{\DUrole{doc}{pair\_write}}}} command
with its \sphinxstyleemphasis{bitmap} option.  When the table is in BITMAP format, the “N”
parameter in the file must be equal to 2\textasciicircum{}M where M is the value
specified in the pair\_style command.  Also, a cutoff parameter cannot
be used as an optional third argument in the pair\_coeff command; the
entire table extent as specified in the file must be used.

\sphinxAtStartPar
If used, the parameter “FPRIME” is followed by 2 values \sphinxstyleemphasis{fplo} and
\sphinxstyleemphasis{fphi} which are the derivative of the force at the innermost and
outermost distances listed in the table.  These values are needed by
the spline construction routines.  If not specified by the “FPRIME”
parameter, they are estimated (less accurately) by the first 2 and
last 2 force values in the table.  This parameter is not used by
BITMAP tables.

\sphinxAtStartPar
Following a blank line, the next N lines list the tabulated values.
On each line, the first value is the index from 1 to N, the second value is
r (in distance units), the third value is the energy (in energy units),
and the fourth is the force (in force units).  The r values must increase
from one line to the next (unless the BITMAP parameter is specified).

\sphinxAtStartPar
Note that one file can contain many sections, each with a tabulated
potential.  LAMMPS reads the file section by section until it finds
one that matches the specified keyword.


\bigskip\hrule\bigskip


\sphinxAtStartPar
Styles with a \sphinxstyleemphasis{gpu}, \sphinxstyleemphasis{intel}, \sphinxstyleemphasis{kk}, \sphinxstyleemphasis{omp}, or \sphinxstyleemphasis{opt} suffix are
functionally the same as the corresponding style without the suffix.
They have been optimized to run faster, depending on your available
hardware, as discussed on the {\hyperref[\detokenize{Speed_packages::doc}]{\sphinxcrossref{\DUrole{doc}{Accelerator packages}}}}
page.  The accelerated styles take the same arguments and should
produce the same results, except for round\sphinxhyphen{}off and precision issues.

\sphinxAtStartPar
These accelerated styles are part of the GPU, INTEL, KOKKOS,
OPENMP, and OPT packages, respectively.  They are only enabled if
LAMMPS was built with those packages.  See the {\hyperref[\detokenize{Build_package::doc}]{\sphinxcrossref{\DUrole{doc}{Build package}}}} page for more info.

\sphinxAtStartPar
You can specify the accelerated styles explicitly in your input script
by including their suffix, or you can use the {\hyperref[\detokenize{Run_options::doc}]{\sphinxcrossref{\DUrole{doc}{\sphinxhyphen{}suffix command\sphinxhyphen{}line switch}}}} when you invoke LAMMPS, or you can use the
{\hyperref[\detokenize{suffix::doc}]{\sphinxcrossref{\DUrole{doc}{suffix}}}} command in your input script.

\sphinxAtStartPar
See the {\hyperref[\detokenize{Speed_packages::doc}]{\sphinxcrossref{\DUrole{doc}{Accelerator packages}}}} page for more
instructions on how to use the accelerated styles effectively.


\bigskip\hrule\bigskip



\subsection{Mixing, shift, table, tail correction, restart, rRESPA info}
\label{\detokenize{pair_table:mixing-shift-table-tail-correction-restart-rrespa-info}}
\sphinxAtStartPar
This pair style does not support mixing.  Thus, coefficients for all
I,J pairs must be specified explicitly.

\sphinxAtStartPar
The {\hyperref[\detokenize{pair_modify::doc}]{\sphinxcrossref{\DUrole{doc}{pair\_modify}}}} shift, table, and tail options are
not relevant for this pair style.

\sphinxAtStartPar
This pair style writes the settings for the “pair\_style table” command
to {\hyperref[\detokenize{restart::doc}]{\sphinxcrossref{\DUrole{doc}{binary restart files}}}}, so a pair\_style command does
not need to specified in an input script that reads a restart file.
However, the coefficient information is not stored in the restart
file, since it is tabulated in the potential files.  Thus, pair\_coeff
commands do need to be specified in the restart input script.

\sphinxAtStartPar
This pair style can only be used via the \sphinxstyleemphasis{pair} keyword of the
{\hyperref[\detokenize{run_style::doc}]{\sphinxcrossref{\DUrole{doc}{run\_style respa}}}} command.  It does not support the
\sphinxstyleemphasis{inner}, \sphinxstyleemphasis{middle}, \sphinxstyleemphasis{outer} keywords.


\bigskip\hrule\bigskip



\subsection{Restrictions}
\label{\detokenize{pair_table:restrictions}}
\sphinxAtStartPar
none


\subsection{Related commands}
\label{\detokenize{pair_table:related-commands}}
\sphinxAtStartPar
{\hyperref[\detokenize{pair_coeff::doc}]{\sphinxcrossref{\DUrole{doc}{pair\_coeff}}}}, {\hyperref[\detokenize{pair_write::doc}]{\sphinxcrossref{\DUrole{doc}{pair\_write}}}}


\subsection{Default}
\label{\detokenize{pair_table:default}}
\sphinxAtStartPar
none


\bigskip\hrule\bigskip

\phantomsection\label{\detokenize{pair_table:wolff2}}
\sphinxAtStartPar
\sphinxstylestrong{(Wolff)} Wolff and Rudd, Comp Phys Comm, 120, 200\sphinxhyphen{}32 (1999).

\sphinxstepscope

\index{pair\_style table/rx@\spxentry{pair\_style table/rx}}\ignorespaces 
\index{pair\_style table/rx/kk@\spxentry{pair\_style table/rx/kk}}\ignorespaces 

\section{pair\_style table/rx command}
\label{\detokenize{pair_table_rx:pair-style-table-rx-command}}\label{\detokenize{pair_table_rx:index-1}}\label{\detokenize{pair_table_rx:index-0}}\label{\detokenize{pair_table_rx::doc}}
\sphinxAtStartPar
Accelerator Variants: \sphinxstyleemphasis{table/rx/kk}


\subsection{Syntax}
\label{\detokenize{pair_table_rx:syntax}}
\begin{sphinxVerbatim}[commandchars=\\\{\}]
\PYG{k}{pair\PYGZus{}style}\PYG{+w}{ }\PYG{n}{table}\PYG{+w}{ }\PYG{n}{style}\PYG{+w}{ }\PYG{n}{N}\PYG{+w}{ }\PYG{n}{...}
\end{sphinxVerbatim}
\begin{itemize}
\item {} 
\sphinxAtStartPar
style = \sphinxstyleemphasis{lookup} or \sphinxstyleemphasis{linear} or \sphinxstyleemphasis{spline} or \sphinxstyleemphasis{bitmap} = method of interpolation

\item {} 
\sphinxAtStartPar
N = use N values in \sphinxstyleemphasis{lookup}, \sphinxstyleemphasis{linear}, \sphinxstyleemphasis{spline} tables

\item {} 
\sphinxAtStartPar
weighting = fractional or molecular (optional)

\end{itemize}


\subsection{Examples}
\label{\detokenize{pair_table_rx:examples}}
\begin{sphinxVerbatim}[commandchars=\\\{\}]
\PYG{k}{pair\PYGZus{}style}\PYG{+w}{ }\PYG{n}{table}\PYG{o}{/}\PYG{n}{rx}\PYG{+w}{ }\PYG{n}{linear}\PYG{+w}{ }\PYG{l+m}{1000}
\PYG{k}{pair\PYGZus{}style}\PYG{+w}{ }\PYG{n}{table}\PYG{o}{/}\PYG{n}{rx}\PYG{+w}{ }\PYG{n}{linear}\PYG{+w}{ }\PYG{l+m}{1000}\PYG{+w}{ }\PYG{n}{fractional}
\PYG{k}{pair\PYGZus{}style}\PYG{+w}{ }\PYG{n}{table}\PYG{o}{/}\PYG{n}{rx}\PYG{+w}{ }\PYG{n}{linear}\PYG{+w}{ }\PYG{l+m}{1000}\PYG{+w}{ }\PYG{n}{molecular}
\PYG{k}{pair\PYGZus{}coeff}\PYG{+w}{ }\PYG{o}{*}\PYG{+w}{ }\PYG{o}{*}\PYG{+w}{ }\PYG{n}{rxn.table}\PYG{+w}{ }\PYG{n}{ENTRY1}\PYG{+w}{ }\PYG{n}{h2o}\PYG{+w}{ }\PYG{n}{h2o}\PYG{+w}{ }\PYG{l+m}{10.0}
\PYG{k}{pair\PYGZus{}coeff}\PYG{+w}{ }\PYG{o}{*}\PYG{+w}{ }\PYG{o}{*}\PYG{+w}{ }\PYG{n}{rxn.table}\PYG{+w}{ }\PYG{n}{ENTRY1}\PYG{+w}{ }\PYG{l+m}{1}\PYG{n}{fluid}\PYG{+w}{ }\PYG{l+m}{1}\PYG{n}{fluid}\PYG{+w}{ }\PYG{l+m}{10.0}
\PYG{k}{pair\PYGZus{}coeff}\PYG{+w}{ }\PYG{o}{*}\PYG{+w}{ }\PYG{l+m}{3}\PYG{+w}{ }\PYG{n}{rxn.table}\PYG{+w}{ }\PYG{n}{ENTRY1}\PYG{+w}{ }\PYG{n}{h2o}\PYG{+w}{ }\PYG{n}{no2}\PYG{+w}{ }\PYG{l+m}{10.0}
\end{sphinxVerbatim}


\subsection{Description}
\label{\detokenize{pair_table_rx:description}}
\sphinxAtStartPar
Style \sphinxstyleemphasis{table/rx} is used in reaction DPD simulations,where the
coarse\sphinxhyphen{}grained (CG) particles are composed of \sphinxstyleemphasis{m} species whose
reaction rate kinetics are determined from a set of \sphinxstyleemphasis{n} reaction rate
equations through the {\hyperref[\detokenize{fix_rx::doc}]{\sphinxcrossref{\DUrole{doc}{fix rx}}}} command.  The species of
one CG particle can interact with a species in a neighboring CG
particle through a site\sphinxhyphen{}site interaction potential model. Style
\sphinxstyleemphasis{table/rx} creates interpolation tables of length \sphinxstyleemphasis{N} from pair
potential and force values listed in a file(s) as a function of
distance.  The files are read by the {\hyperref[\detokenize{pair_coeff::doc}]{\sphinxcrossref{\DUrole{doc}{pair\_coeff}}}}
command.

\sphinxAtStartPar
The interpolation tables are created by fitting cubic splines to the
file values and interpolating energy and force values at each of \sphinxstyleemphasis{N}
distances.  During a simulation, these tables are used to interpolate
energy and force values as needed.  The interpolation is done in one
of 4 styles: \sphinxstyleemphasis{lookup}, \sphinxstyleemphasis{linear}, \sphinxstyleemphasis{spline}, or \sphinxstyleemphasis{bitmap}.

\sphinxAtStartPar
For the \sphinxstyleemphasis{lookup} style, the distance between two atoms is used to find
the nearest table entry, which is the energy or force.

\sphinxAtStartPar
For the \sphinxstyleemphasis{linear} style, the pair distance is used to find 2
surrounding table values from which an energy or force is computed by
linear interpolation.

\sphinxAtStartPar
For the \sphinxstyleemphasis{spline} style, a cubic spline coefficients are computed and
stored at each of the \sphinxstyleemphasis{N} values in the table.  The pair distance is
used to find the appropriate set of coefficients which are used to
evaluate a cubic polynomial which computes the energy or force.

\sphinxAtStartPar
For the \sphinxstyleemphasis{bitmap} style, the N means to create interpolation tables
that are 2\textasciicircum{}N in length.  The pair distance is used to index into the
table via a fast bit\sphinxhyphen{}mapping technique {\hyperref[\detokenize{pair_table_rx:wolff}]{\sphinxcrossref{\DUrole{std,std-ref}{(Wolff)}}}} and a linear
interpolation is performed between adjacent table values.

\sphinxAtStartPar
The following coefficients must be defined for each pair of atoms
types via the {\hyperref[\detokenize{pair_coeff::doc}]{\sphinxcrossref{\DUrole{doc}{pair\_coeff}}}} command as in the examples
above.
\begin{itemize}
\item {} 
\sphinxAtStartPar
filename

\item {} 
\sphinxAtStartPar
keyword

\item {} 
\sphinxAtStartPar
species1

\item {} 
\sphinxAtStartPar
species2

\item {} 
\sphinxAtStartPar
cutoff (distance units)

\end{itemize}

\sphinxAtStartPar
The filename specifies a file containing tabulated energy and force
values.  The keyword specifies a section of the file.  The cutoff is
an optional coefficient.  If not specified, the outer cutoff in the
table itself (see below) will be used to build an interpolation table
that extend to the largest tabulated distance.  If specified, only
file values up to the cutoff are used to create the interpolation
table.  The format of this file is described below.

\sphinxAtStartPar
The species tags define the site\sphinxhyphen{}site interaction potential between
two species contained within two different particles.  The species
tags must either correspond to the species defined in the reaction
kinetics files specified with the {\hyperref[\detokenize{fix_rx::doc}]{\sphinxcrossref{\DUrole{doc}{fix rx}}}} command or they
must correspond to the tag “1fluid”, signifying interaction with a
product species mixture determined through a one\sphinxhyphen{}fluid approximation.
The interaction potential is weighted by the geometric average of
either the mole fraction concentrations or the number of molecules
associated with the interacting coarse\sphinxhyphen{}grained particles (see the
\sphinxstyleemphasis{fractional} or \sphinxstyleemphasis{molecular} weighting pair style options). The coarse\sphinxhyphen{}grained potential is
stored before and after the reaction kinetics solver is applied, where
the difference is defined to be the internal chemical energy (uChem).


\bigskip\hrule\bigskip


\sphinxAtStartPar
Here are some guidelines for using the pair\_style table/rx command to
best effect:
\begin{itemize}
\item {} 
\sphinxAtStartPar
Vary the number of table points; you may need to use more than you think
to get good resolution.

\item {} 
\sphinxAtStartPar
Always use the {\hyperref[\detokenize{pair_write::doc}]{\sphinxcrossref{\DUrole{doc}{pair\_write}}}} command to produce a plot
of what the final interpolated potential looks like.  This can show up
interpolation “features” you may not like.

\item {} 
\sphinxAtStartPar
Start with the linear style; it’s the style least likely to have problems.

\item {} 
\sphinxAtStartPar
Use \sphinxstyleemphasis{N} in the pair\_style command equal to the “N” in the tabulation
file, and use the “RSQ” or “BITMAP” parameter, so additional interpolation
is not needed.  See discussion below.

\item {} 
\sphinxAtStartPar
Make sure that your tabulated forces and tabulated energies are consistent
(dE/dr = \sphinxhyphen{}F) along the entire range of r values.

\item {} 
\sphinxAtStartPar
Use as large an inner cutoff as possible.  This avoids fitting splines
to very steep parts of the potential.

\end{itemize}


\bigskip\hrule\bigskip


\sphinxAtStartPar
The format of a tabulated file is a series of one or more sections,
defined as follows (without the parenthesized comments):

\begin{sphinxVerbatim}[commandchars=\\\{\}]
\PYGZsh{} Morse potential for Fe   (one or more comment or blank lines)

MORSE\PYGZus{}FE                   (keyword is first text on line)
N 500 R 1.0 10.0           (N, R, RSQ, BITMAP, FPRIME parameters)
                           (blank)
1 1.0 25.5 102.34          (index, r, energy, force)
2 1.02 23.4 98.5
...
500 10.0 0.001 0.003
\end{sphinxVerbatim}

\sphinxAtStartPar
A section begins with a non\sphinxhyphen{}blank line whose first character is not a
“\#”; blank lines or lines starting with “\#” can be used as comments
between sections.  The first line begins with a keyword which
identifies the section.  The line can contain additional text, but the
initial text must match the argument specified in the pair\_coeff
command.  The next line lists (in any order) one or more parameters
for the table.  Each parameter is a keyword followed by one or more
numeric values.

\sphinxAtStartPar
The parameter “N” is required and its value is the number of table
entries that follow.  Note that this may be different than the \sphinxstyleemphasis{N}
specified in the {\hyperref[\detokenize{pair_style::doc}]{\sphinxcrossref{\DUrole{doc}{pair\_style table/rx}}}} command.  Let
Ntable = \sphinxstyleemphasis{N} in the pair\_style command, and Nfile = “N” in the
tabulated file.  What LAMMPS does is a preliminary interpolation by
creating splines using the Nfile tabulated values as nodal points.  It
uses these to interpolate as needed to generate energy and force
values at Ntable different points.  The resulting tables of length
Ntable are then used as described above, when computing energy and
force for individual pair distances.  This means that if you want the
interpolation tables of length Ntable to match exactly what is in the
tabulated file (with effectively no preliminary interpolation), you
should set Ntable = Nfile, and use the “RSQ” or “BITMAP” parameter.
The internal table abscissa is RSQ (separation distance squared).

\sphinxAtStartPar
All other parameters are optional.  If “R” or “RSQ” or “BITMAP” does
not appear, then the distances in each line of the table are used
as\sphinxhyphen{}is to perform spline interpolation.  In this case, the table values
can be spaced in \sphinxstyleemphasis{r} uniformly or however you wish to position table
values in regions of large gradients.

\sphinxAtStartPar
If used, the parameters “R” or “RSQ” are followed by 2 values \sphinxstyleemphasis{rlo}
and \sphinxstyleemphasis{rhi}.  If specified, the distance associated with each energy and
force value is computed from these 2 values (at high accuracy), rather
than using the (low\sphinxhyphen{}accuracy) value listed in each line of the table.
The distance values in the table file are ignored in this case.
For “R”, distances uniformly spaced between \sphinxstyleemphasis{rlo} and \sphinxstyleemphasis{rhi} are
computed; for “RSQ”, squared distances uniformly spaced between
\sphinxstyleemphasis{rlo*rlo} and \sphinxstyleemphasis{rhi*rhi} are computed.

\sphinxAtStartPar
If used, the parameter “BITMAP” is also followed by 2 values \sphinxstyleemphasis{rlo} and
\sphinxstyleemphasis{rhi}.  These values, along with the “N” value determine the ordering
of the N lines that follow and what distance is associated with each.
This ordering is complex, so it is not documented here, since this
file is typically produced by the {\hyperref[\detokenize{pair_write::doc}]{\sphinxcrossref{\DUrole{doc}{pair\_write}}}} command
with its \sphinxstyleemphasis{bitmap} option.  When the table is in BITMAP format, the “N”
parameter in the file must be equal to 2\textasciicircum{}M where M is the value
specified in the pair\_style command.  Also, a cutoff parameter cannot
be used as an optional third argument in the pair\_coeff command; the
entire table extent as specified in the file must be used.

\sphinxAtStartPar
If used, the parameter “FPRIME” is followed by 2 values \sphinxstyleemphasis{fplo} and
\sphinxstyleemphasis{fphi} which are the derivative of the force at the innermost and
outermost distances listed in the table.  These values are needed by
the spline construction routines.  If not specified by the “FPRIME”
parameter, they are estimated (less accurately) by the first 2 and
last 2 force values in the table.  This parameter is not used by
BITMAP tables.

\sphinxAtStartPar
Following a blank line, the next N lines list the tabulated values.
On each line, the first value is the index from 1 to N, the second value is
r (in distance units), the third value is the energy (in energy units),
and the fourth is the force (in force units).  The r values must increase
from one line to the next (unless the BITMAP parameter is specified).

\sphinxAtStartPar
Note that one file can contain many sections, each with a tabulated
potential.  LAMMPS reads the file section by section until it finds
one that matches the specified keyword.


\bigskip\hrule\bigskip



\subsection{Mixing, shift, table, tail correction, restart, rRESPA info}
\label{\detokenize{pair_table_rx:mixing-shift-table-tail-correction-restart-rrespa-info}}
\sphinxAtStartPar
This pair style does not support mixing.  Thus, coefficients for all
I,J pairs must be specified explicitly.

\sphinxAtStartPar
The {\hyperref[\detokenize{pair_modify::doc}]{\sphinxcrossref{\DUrole{doc}{pair\_modify}}}} shift, table, and tail options are
not relevant for this pair style.

\sphinxAtStartPar
This pair style writes the settings for the “pair\_style table/rx” command
to {\hyperref[\detokenize{restart::doc}]{\sphinxcrossref{\DUrole{doc}{binary restart files}}}}, so a pair\_style command does
not need to specified in an input script that reads a restart file.
However, the coefficient information is not stored in the restart
file, since it is tabulated in the potential files.  Thus, pair\_coeff
commands do need to be specified in the restart input script.

\sphinxAtStartPar
This pair style can only be used via the \sphinxstyleemphasis{pair} keyword of the
{\hyperref[\detokenize{run_style::doc}]{\sphinxcrossref{\DUrole{doc}{run\_style respa}}}} command.  It does not support the
\sphinxstyleemphasis{inner}, \sphinxstyleemphasis{middle}, \sphinxstyleemphasis{outer} keywords.


\bigskip\hrule\bigskip


\sphinxAtStartPar
Styles with a \sphinxstyleemphasis{gpu}, \sphinxstyleemphasis{intel}, \sphinxstyleemphasis{kk}, \sphinxstyleemphasis{omp}, or \sphinxstyleemphasis{opt} suffix are
functionally the same as the corresponding style without the suffix.
They have been optimized to run faster, depending on your available
hardware, as discussed on the {\hyperref[\detokenize{Speed_packages::doc}]{\sphinxcrossref{\DUrole{doc}{Accelerator packages}}}}
page.  The accelerated styles take the same arguments and should
produce the same results, except for round\sphinxhyphen{}off and precision issues.

\sphinxAtStartPar
These accelerated styles are part of the GPU, INTEL, KOKKOS,
OPENMP, and OPT packages, respectively.  They are only enabled if
LAMMPS was built with those packages.  See the {\hyperref[\detokenize{Build_package::doc}]{\sphinxcrossref{\DUrole{doc}{Build package}}}} page for more info.

\sphinxAtStartPar
You can specify the accelerated styles explicitly in your input script
by including their suffix, or you can use the {\hyperref[\detokenize{Run_options::doc}]{\sphinxcrossref{\DUrole{doc}{\sphinxhyphen{}suffix command\sphinxhyphen{}line switch}}}} when you invoke LAMMPS, or you can use the
{\hyperref[\detokenize{suffix::doc}]{\sphinxcrossref{\DUrole{doc}{suffix}}}} command in your input script.

\sphinxAtStartPar
See the {\hyperref[\detokenize{Speed_packages::doc}]{\sphinxcrossref{\DUrole{doc}{Accelerator packages}}}} page for more
instructions on how to use the accelerated styles effectively.


\bigskip\hrule\bigskip



\subsection{Restrictions}
\label{\detokenize{pair_table_rx:restrictions}}
\sphinxAtStartPar
This command is part of the DPD\sphinxhyphen{}REACT package.  It is only enabled if
LAMMPS was built with that package.  See the {\hyperref[\detokenize{Build_package::doc}]{\sphinxcrossref{\DUrole{doc}{Build package}}}} page for more info.


\subsection{Related commands}
\label{\detokenize{pair_table_rx:related-commands}}
\sphinxAtStartPar
{\hyperref[\detokenize{pair_coeff::doc}]{\sphinxcrossref{\DUrole{doc}{pair\_coeff}}}}


\subsection{Default}
\label{\detokenize{pair_table_rx:default}}
\sphinxAtStartPar
fractional weighting


\bigskip\hrule\bigskip

\phantomsection\label{\detokenize{pair_table_rx:wolff}}
\sphinxAtStartPar
\sphinxstylestrong{(Wolff)} Wolff and Rudd, Comp Phys Comm, 120, 200\sphinxhyphen{}32 (1999).

\sphinxstepscope

\index{pair\_style tersoff@\spxentry{pair\_style tersoff}}\ignorespaces 
\index{pair\_style tersoff/table@\spxentry{pair\_style tersoff/table}}\ignorespaces 
\index{pair\_style tersoff/gpu@\spxentry{pair\_style tersoff/gpu}}\ignorespaces 
\index{pair\_style tersoff/intel@\spxentry{pair\_style tersoff/intel}}\ignorespaces 
\index{pair\_style tersoff/kk@\spxentry{pair\_style tersoff/kk}}\ignorespaces 
\index{pair\_style tersoff/omp@\spxentry{pair\_style tersoff/omp}}\ignorespaces 
\index{pair\_style tersoff/table/omp@\spxentry{pair\_style tersoff/table/omp}}\ignorespaces 

\section{pair\_style tersoff command}
\label{\detokenize{pair_tersoff:pair-style-tersoff-command}}\label{\detokenize{pair_tersoff:index-6}}\label{\detokenize{pair_tersoff:index-5}}\label{\detokenize{pair_tersoff:index-4}}\label{\detokenize{pair_tersoff:index-3}}\label{\detokenize{pair_tersoff:index-2}}\label{\detokenize{pair_tersoff:index-1}}\label{\detokenize{pair_tersoff:index-0}}\label{\detokenize{pair_tersoff::doc}}
\sphinxAtStartPar
Accelerator Variants: \sphinxstyleemphasis{tersoff/gpu}, \sphinxstyleemphasis{tersoff/intel}, \sphinxstyleemphasis{tersoff/kk}, \sphinxstyleemphasis{tersoff/omp}


\section{pair\_style tersoff/table command}
\label{\detokenize{pair_tersoff:pair-style-tersoff-table-command}}
\sphinxAtStartPar
Accelerator Variants: \sphinxstyleemphasis{tersoff/table/omp}


\subsection{Syntax}
\label{\detokenize{pair_tersoff:syntax}}
\begin{sphinxVerbatim}[commandchars=\\\{\}]
\PYG{k}{pair\PYGZus{}style}\PYG{+w}{ }\PYG{n}{style}\PYG{+w}{ }\PYG{n}{keywords}\PYG{+w}{ }\PYG{n}{values}
\end{sphinxVerbatim}
\begin{itemize}
\item {} 
\sphinxAtStartPar
style = \sphinxstyleemphasis{tersoff} or \sphinxstyleemphasis{tersoff/table}

\item {} 
\sphinxAtStartPar
keyword = \sphinxstyleemphasis{shift}
\begin{sphinxalltt}
\sphinxstyleemphasis{shift} value = delta
  delta = negative shift in equilibrium bond length
\end{sphinxalltt}

\end{itemize}


\subsection{Examples}
\label{\detokenize{pair_tersoff:examples}}
\begin{sphinxVerbatim}[commandchars=\\\{\}]
\PYG{k}{pair\PYGZus{}style}\PYG{+w}{ }\PYG{n}{tersoff}
\PYG{k}{pair\PYGZus{}coeff}\PYG{+w}{ }\PYG{o}{*}\PYG{+w}{ }\PYG{o}{*}\PYG{+w}{ }\PYG{n}{Si.tersoff}\PYG{+w}{ }\PYG{n}{Si}
\PYG{k}{pair\PYGZus{}coeff}\PYG{+w}{ }\PYG{o}{*}\PYG{+w}{ }\PYG{o}{*}\PYG{+w}{ }\PYG{n}{SiC.tersoff}\PYG{+w}{ }\PYG{n}{Si}\PYG{+w}{ }\PYG{n}{C}\PYG{+w}{ }\PYG{n}{Si}

\PYG{k}{pair\PYGZus{}style}\PYG{+w}{ }\PYG{n}{tersoff}\PYG{o}{/}\PYG{n}{table}
\PYG{k}{pair\PYGZus{}coeff}\PYG{+w}{ }\PYG{o}{*}\PYG{+w}{ }\PYG{o}{*}\PYG{+w}{ }\PYG{n}{SiCGe.tersoff}\PYG{+w}{ }\PYG{n}{Si}\PYG{n+nv}{(}\PYG{n+nv}{D}\PYG{n+nv}{)}

\PYG{k}{pair\PYGZus{}style}\PYG{+w}{ }\PYG{n}{tersoff}\PYG{+w}{ }\PYG{n}{shift}\PYG{+w}{ }\PYG{l+m}{0.05}
\PYG{k}{pair\PYGZus{}coeff}\PYG{+w}{ }\PYG{o}{*}\PYG{+w}{ }\PYG{o}{*}\PYG{+w}{ }\PYG{n}{Si.tersoff}\PYG{+w}{ }\PYG{n}{Si}
\end{sphinxVerbatim}


\subsection{Description}
\label{\detokenize{pair_tersoff:description}}
\sphinxAtStartPar
The \sphinxstyleemphasis{tersoff} style computes a 3\sphinxhyphen{}body Tersoff potential
{\hyperref[\detokenize{pair_tersoff:tersoff-11}]{\sphinxcrossref{\DUrole{std,std-ref}{(Tersoff\_1)}}}} for the energy E of a system of atoms as
\begin{equation*}
\begin{split}E & = \frac{1}{2} \sum_i \sum_{j \neq i} V_{ij} \\
V_{ij} & = f_C(r_{ij} + \delta) \left[ f_R(r_{ij} + \delta) + b_{ij} f_A(r_{ij} + \delta) \right] \\
f_C(r) & = \left\{ \begin{array} {r@{\quad:\quad}l}
  1 & r < R - D \\
  \frac{1}{2} - \frac{1}{2} \sin \left( \frac{\pi}{2} \frac{r-R}{D} \right) &
    R-D < r < R + D \\
  0 & r > R + D
  \end{array} \right. \\
f_R(r) & =  A \exp (-\lambda_1 r) \\
f_A(r) & =  -B \exp (-\lambda_2 r) \\
b_{ij} & =  \left( 1 + \beta^n {\zeta_{ij}}^n \right)^{-\frac{1}{2n}} \\
\zeta_{ij} & =  \sum_{k \neq i,j} f_C(r_{ik} + \delta) g \left[ \theta_{ijk}(r_{ij}, r_{ik}) \right]
                 \exp \left[ {\lambda_3}^m (r_{ij} - r_{ik})^m \right] \\
g(\theta) & =  \gamma_{ijk} \left( 1 + \frac{c^2}{d^2} -
                \frac{c^2}{\left[ d^2 + (\cos \theta - \cos \theta_0)^2\right]} \right)\end{split}
\end{equation*}
\sphinxAtStartPar
where \(f_R\) is a two\sphinxhyphen{}body term and \(f_A\) includes three\sphinxhyphen{}body
interactions.  The summations in the formula are over all neighbors
J and K of atom I within a cutoff distance = R + D. \(\delta\) is
an optional negative shift of the equilibrium bond length, as described below.

\sphinxAtStartPar
The \sphinxstyleemphasis{tersoff/table} style uses tabulated forms for the two\sphinxhyphen{}body,
environment and angular functions. Linear interpolation is performed
between adjacent table entries. The table length is chosen to be
accurate within 10\textasciicircum{}\sphinxhyphen{}6 with respect to the \sphinxstyleemphasis{tersoff} style energy.
The \sphinxstyleemphasis{tersoff/table} should give better performance in terms of speed.

\sphinxAtStartPar
Only a single pair\_coeff command is used with the \sphinxstyleemphasis{tersoff} style
which specifies a Tersoff potential file with parameters for all
needed elements.  These are mapped to LAMMPS atom types by specifying
N additional arguments after the filename in the pair\_coeff command,
where N is the number of LAMMPS atom types:
\begin{itemize}
\item {} 
\sphinxAtStartPar
filename

\item {} 
\sphinxAtStartPar
N element names = mapping of Tersoff elements to atom types

\end{itemize}

\sphinxAtStartPar
See the {\hyperref[\detokenize{pair_coeff::doc}]{\sphinxcrossref{\DUrole{doc}{pair\_coeff}}}} page for alternate ways
to specify the path for the potential file.

\sphinxAtStartPar
As an example, imagine the SiC.tersoff file has Tersoff values for Si
and C.  If your LAMMPS simulation has 4 atoms types and you want the
first 3 to be Si, and the fourth to be C, you would use the following
pair\_coeff command:

\begin{sphinxVerbatim}[commandchars=\\\{\}]
\PYG{k}{pair\PYGZus{}coeff}\PYG{+w}{ }\PYG{o}{*}\PYG{+w}{ }\PYG{o}{*}\PYG{+w}{ }\PYG{n}{SiC.tersoff}\PYG{+w}{ }\PYG{n}{Si}\PYG{+w}{ }\PYG{n}{Si}\PYG{+w}{ }\PYG{n}{Si}\PYG{+w}{ }\PYG{n}{C}
\end{sphinxVerbatim}

\sphinxAtStartPar
The first 2 arguments must be * * so as to span all LAMMPS atom types.
The first three Si arguments map LAMMPS atom types 1,2,3 to the Si
element in the Tersoff file.  The final C argument maps LAMMPS atom
type 4 to the C element in the Tersoff file.  If a mapping value is
specified as NULL, the mapping is not performed.  This can be used
when a \sphinxstyleemphasis{tersoff} potential is used as part of the \sphinxstyleemphasis{hybrid} pair style.
The NULL values are placeholders for atom types that will be used with
other potentials.

\sphinxAtStartPar
Tersoff files in the \sphinxstyleemphasis{potentials} directory of the LAMMPS distribution
have a “.tersoff” suffix.  Lines that are not blank or comments
(starting with \#) define parameters for a triplet of elements.  The
parameters in a single entry correspond to coefficients in the formula
above:
\begin{itemize}
\item {} 
\sphinxAtStartPar
element 1 (the center atom in a 3\sphinxhyphen{}body interaction)

\item {} 
\sphinxAtStartPar
element 2 (the atom bonded to the center atom)

\item {} 
\sphinxAtStartPar
element 3 (the atom influencing the 1\sphinxhyphen{}2 bond in a bond\sphinxhyphen{}order sense)

\item {} 
\sphinxAtStartPar
m

\item {} 
\sphinxAtStartPar
\(\gamma\)

\item {} 
\sphinxAtStartPar
\(\lambda_3\) (1/distance units)

\item {} 
\sphinxAtStartPar
c

\item {} 
\sphinxAtStartPar
d

\item {} 
\sphinxAtStartPar
\(\cos\theta_0\) (can be a value \textless{} \sphinxhyphen{}1 or \textgreater{} 1)

\item {} 
\sphinxAtStartPar
n

\item {} 
\sphinxAtStartPar
\(\beta\)

\item {} 
\sphinxAtStartPar
\(\lambda_2\) (1/distance units)

\item {} 
\sphinxAtStartPar
B (energy units)

\item {} 
\sphinxAtStartPar
R (distance units)

\item {} 
\sphinxAtStartPar
D (distance units)

\item {} 
\sphinxAtStartPar
\(\lambda_1\) (1/distance units)

\item {} 
\sphinxAtStartPar
A (energy units)

\end{itemize}

\sphinxAtStartPar
The n, \(\beta\), \(\lambda_2\), B, \(\lambda_1\), and A
parameters are only used for
two\sphinxhyphen{}body interactions.  The m, \(\gamma\), \(\lambda_3\), c, d,
and \(\cos\theta_0\)
parameters are only used for three\sphinxhyphen{}body interactions. The R and D
parameters are used for both two\sphinxhyphen{}body and three\sphinxhyphen{}body interactions. The
non\sphinxhyphen{}annotated parameters are unitless.  The value of m must be 3 or 1.

\sphinxAtStartPar
The Tersoff potential file must contain entries for all the elements
listed in the pair\_coeff command.  It can also contain entries for
additional elements not being used in a particular simulation; LAMMPS
ignores those entries.

\sphinxAtStartPar
For a single\sphinxhyphen{}element simulation, only a single entry is required
(e.g. SiSiSi).  For a two\sphinxhyphen{}element simulation, the file must contain 8
entries (for SiSiSi, SiSiC, SiCSi, SiCC, CSiSi, CSiC, CCSi, CCC), that
specify Tersoff parameters for all permutations of the two elements
interacting in three\sphinxhyphen{}body configurations.  Thus for 3 elements, 27
entries would be required, etc.

\sphinxAtStartPar
As annotated above, the first element in the entry is the center atom
in a three\sphinxhyphen{}body interaction and it is bonded to the second atom and the
bond is influenced by the third atom.  Thus an entry for SiCC means Si
bonded to a C with another C atom influencing the bond.  Thus
three\sphinxhyphen{}body parameters for SiCSi and SiSiC entries will not, in
general, be the same.  The parameters used for the two\sphinxhyphen{}body
interaction come from the entry where the second element is repeated.
Thus the two\sphinxhyphen{}body parameters for Si interacting with C, comes from the
SiCC entry.

\sphinxAtStartPar
The parameters used for a particular
three\sphinxhyphen{}body interaction come from the entry with the corresponding
three elements.  The parameters used only for two\sphinxhyphen{}body interactions
(n, \(\beta\), \(\lambda_2\), B, \(\lambda_1\), and A)
in entries whose second and third
element are different (e.g. SiCSi) are not used for anything and can
be set to 0.0 if desired.

\sphinxAtStartPar
Note that the twobody parameters in entries such as SiCC and CSiSi
are often the same, due to the common use of symmetric mixing rules,
but this is not always the case. For example, the beta and n parameters in
Tersoff\_2 {\hyperref[\detokenize{pair_tersoff:tersoff-21}]{\sphinxcrossref{\DUrole{std,std-ref}{(Tersoff\_2)}}}} are not symmetric.
Similarly, the threebody parameters in entries such as SiCSi and SiSiC
are often the same, but this is not always the case, particularly
the value of R, which is sometimes typed on the
first and second elements, sometimes on the first and third elements.
Hence the need to specify R and D explicitly for all element triples.
For example, while Tersoff’s notation
in Tersoff\_2 {\hyperref[\detokenize{pair_tersoff:tersoff-21}]{\sphinxcrossref{\DUrole{std,std-ref}{(Tersoff\_2)}}}} is ambiguous on this point,
and properties of the zincblende lattice are the same for either choice,
Tersoff’s results for rocksalt are consistent with typing on the first
and third elements. {\hyperref[\detokenize{pair_tersoff:albe}]{\sphinxcrossref{\DUrole{std,std-ref}{Albe et al.}}}} adopts the same convention.
Conversely, the potential for B/N/C from the Cagin group
uses the opposite convention, typing on the first and second elements.

\sphinxAtStartPar
We chose the above form so as to enable users to define all commonly
used variants of the Tersoff potential.  In particular, our form
reduces to the original Tersoff form when m = 3 and gamma = 1, while
it reduces to the form of {\hyperref[\detokenize{pair_tersoff:albe}]{\sphinxcrossref{\DUrole{std,std-ref}{Albe et al.}}}} when beta = 1 and m = 1.
Note that in the current Tersoff implementation in LAMMPS, m must be
specified as either 3 or 1.  Tersoff used a slightly different but
equivalent form for alloys, which we will refer to as Tersoff\_2
potential {\hyperref[\detokenize{pair_tersoff:tersoff-21}]{\sphinxcrossref{\DUrole{std,std-ref}{(Tersoff\_2)}}}}.
The \sphinxstyleemphasis{tersoff/table} style implements
Tersoff\_2 parameterization only.

\sphinxAtStartPar
LAMMPS parameter values for Tersoff\_2 can be obtained as follows:
\(\gamma_{ijk} = \omega_{ik}\), \(\lambda_3 = 0\) and the value of
m has no effect.  The parameters for species i and j can be calculated
using the Tersoff\_2 mixing rules:
\begin{equation*}
\begin{split}\lambda_1^{i,j} & = \frac{1}{2}(\lambda_1^i + \lambda_1^j)\\
\lambda_2^{i,j} & = \frac{1}{2}(\lambda_2^i + \lambda_2^j)\\
A_{i,j} & = (A_{i}A_{j})^{1/2}\\
B_{i,j} & = \chi_{ij}(B_{i}B_{j})^{1/2}\\
R_{i,j} & = (R_{i}R_{j})^{1/2}\\
S_{i,j} & = (S_{i}S_{j})^{1/2}\end{split}
\end{equation*}
\sphinxAtStartPar
Tersoff\_2 parameters R and S must be converted to the LAMMPS
parameters R and D (R is different in both forms), using the following
relations: R=(R’+S’)/2 and D=(S’\sphinxhyphen{}R’)/2, where the primes indicate the
Tersoff\_2 parameters.

\sphinxAtStartPar
In the potentials directory, the file SiCGe.tersoff provides the
LAMMPS parameters for Tersoff’s various versions of Si, as well as his
alloy parameters for Si, C, and Ge. This file can be used for pure Si,
(three different versions), pure C, pure Ge, binary SiC, and binary
SiGe.  LAMMPS will generate an error if this file is used with any
combination involving C and Ge, since there are no entries for the GeC
interactions (Tersoff did not publish parameters for this
cross\sphinxhyphen{}interaction.)  Tersoff files are also provided for the SiC alloy
(SiC.tersoff) and the GaN (GaN.tersoff) alloys.

\sphinxAtStartPar
Many thanks to Rutuparna Narulkar, David Farrell, and Xiaowang Zhou
for helping clarify how Tersoff parameters for alloys have been
defined in various papers.

\sphinxAtStartPar
The \sphinxstyleemphasis{shift} keyword computes the energy E of a system of atoms, whose
formula is the same as the Tersoff potential. The only modification is
that the original equilibrium bond length ( \(r_0\)) of the system
is shifted to \(r_0-\delta\).  The minus sign arises because each
radial distance \(r\) is replaced by \(r+\delta\).

\sphinxAtStartPar
The \sphinxstyleemphasis{shift} keyword is designed for simulations of closely matched van
der Waals heterostructures.  For instance, consider the case of a system
with few\sphinxhyphen{}layers graphene atop a thick hexagonal boron nitride (h\sphinxhyphen{}BN)
substrate simulated using periodic boundary conditions. The experimental
lattice mismatch of \textasciitilde{}1.8\% between graphene and h\sphinxhyphen{}BN is not well captured
by the equilibrium lattice constants of available potentials, thus a
small in\sphinxhyphen{}plane strain will be introduced in the system when building a
periodic supercell.  To minimize the effect of strain on simulation
results, the \sphinxstyleemphasis{shift} keyword allows adjusting the equilibrium bond
length of one of the two materials (e.g., h\sphinxhyphen{}BN). Validation, benchmark
tests, and applications of the \sphinxstyleemphasis{shift} keyword can be found in
{\hyperref[\detokenize{pair_tersoff:mandelli1}]{\sphinxcrossref{\DUrole{std,std-ref}{(Mandelli\_1)}}}} and {\hyperref[\detokenize{pair_tersoff:ouyang5}]{\sphinxcrossref{\DUrole{std,std-ref}{(Ouyang\_1)}}}}.

\sphinxAtStartPar
For the specific case discussed above, the force field can be defined as

\begin{sphinxVerbatim}[commandchars=\\\{\}]
\PYG{k}{pair\PYGZus{}style}\PYG{+w}{  }\PYG{n}{hybrid}\PYG{o}{/}\PYG{n}{overlay}\PYG{+w}{ }\PYG{n}{rebo}\PYG{+w}{ }\PYG{n}{tersoff}\PYG{+w}{ }\PYG{n}{shift}\PYG{+w}{ }\PYG{o}{\PYGZhy{}}\PYG{l+m}{0.00407}\PYG{+w}{ }\PYG{n}{ilp}\PYG{o}{/}\PYG{n}{graphene}\PYG{o}{/}\PYG{n}{hbn}\PYG{+w}{ }\PYG{l+m}{16.0}\PYG{+w}{ }\PYG{n}{coul}\PYG{o}{/}\PYG{n}{shield}\PYG{+w}{ }\PYG{l+m}{16.0}
\PYG{k}{pair\PYGZus{}coeff}\PYG{+w}{  }\PYG{o}{*}\PYG{+w}{ }\PYG{o}{*}\PYG{+w}{ }\PYG{n}{rebo}\PYG{+w}{              }\PYG{n}{CH.rebo}\PYG{+w}{     }\PYG{n}{NULL}\PYG{+w}{ }\PYG{n}{NULL}\PYG{+w}{ }\PYG{n}{C}
\PYG{k}{pair\PYGZus{}coeff}\PYG{+w}{  }\PYG{o}{*}\PYG{+w}{ }\PYG{o}{*}\PYG{+w}{ }\PYG{n}{tersoff}\PYG{+w}{           }\PYG{n}{BNC.tersoff}\PYG{+w}{ }\PYG{n}{B}\PYG{+w}{    }\PYG{n}{N}\PYG{+w}{    }\PYG{n}{NULL}
\PYG{k}{pair\PYGZus{}coeff}\PYG{+w}{  }\PYG{o}{*}\PYG{+w}{ }\PYG{o}{*}\PYG{+w}{ }\PYG{n}{ilp}\PYG{o}{/}\PYG{n}{graphene}\PYG{o}{/}\PYG{n}{hbn}\PYG{+w}{  }\PYG{n}{BNCH.ILP}\PYG{+w}{    }\PYG{n}{B}\PYG{+w}{    }\PYG{n}{N}\PYG{+w}{    }\PYG{n}{C}
\PYG{k}{pair\PYGZus{}coeff}\PYG{+w}{  }\PYG{l+m}{1}\PYG{+w}{ }\PYG{l+m}{1}\PYG{+w}{ }\PYG{n}{coul}\PYG{o}{/}\PYG{n}{shield}\PYG{+w}{ }\PYG{l+m}{0.70}
\PYG{k}{pair\PYGZus{}coeff}\PYG{+w}{  }\PYG{l+m}{1}\PYG{+w}{ }\PYG{l+m}{2}\PYG{+w}{ }\PYG{n}{coul}\PYG{o}{/}\PYG{n}{shield}\PYG{+w}{ }\PYG{l+m}{0.695}
\PYG{k}{pair\PYGZus{}coeff}\PYG{+w}{  }\PYG{l+m}{2}\PYG{+w}{ }\PYG{l+m}{2}\PYG{+w}{ }\PYG{n}{coul}\PYG{o}{/}\PYG{n}{shield}\PYG{+w}{ }\PYG{l+m}{0.69}
\end{sphinxVerbatim}


\bigskip\hrule\bigskip


\sphinxAtStartPar
Styles with a \sphinxstyleemphasis{gpu}, \sphinxstyleemphasis{intel}, \sphinxstyleemphasis{kk}, \sphinxstyleemphasis{omp}, or \sphinxstyleemphasis{opt} suffix are
functionally the same as the corresponding style without the suffix.
They have been optimized to run faster, depending on your available
hardware, as discussed on the {\hyperref[\detokenize{Speed_packages::doc}]{\sphinxcrossref{\DUrole{doc}{Accelerator packages}}}}
page.  The accelerated styles take the same arguments and should
produce the same results, except for round\sphinxhyphen{}off and precision issues.

\sphinxAtStartPar
These accelerated styles are part of the GPU, INTEL, KOKKOS,
OPENMP, and OPT packages, respectively.  They are only enabled if
LAMMPS was built with those packages.  See the {\hyperref[\detokenize{Build_package::doc}]{\sphinxcrossref{\DUrole{doc}{Build package}}}} page for more info.

\sphinxAtStartPar
You can specify the accelerated styles explicitly in your input script
by including their suffix, or you can use the {\hyperref[\detokenize{Run_options::doc}]{\sphinxcrossref{\DUrole{doc}{\sphinxhyphen{}suffix command\sphinxhyphen{}line switch}}}} when you invoke LAMMPS, or you can use the
{\hyperref[\detokenize{suffix::doc}]{\sphinxcrossref{\DUrole{doc}{suffix}}}} command in your input script.

\sphinxAtStartPar
See the {\hyperref[\detokenize{Speed_packages::doc}]{\sphinxcrossref{\DUrole{doc}{Accelerator packages}}}} page for more
instructions on how to use the accelerated styles effectively.


\bigskip\hrule\bigskip



\subsection{Mixing, shift, table, tail correction, restart, rRESPA info}
\label{\detokenize{pair_tersoff:mixing-shift-table-tail-correction-restart-rrespa-info}}
\sphinxAtStartPar
For atom type pairs I,J and I != J, where types I and J correspond to
two different element types, mixing is performed by LAMMPS as
described above from values in the potential file.

\sphinxAtStartPar
This pair style does not support the {\hyperref[\detokenize{pair_modify::doc}]{\sphinxcrossref{\DUrole{doc}{pair\_modify}}}}
shift, table, and tail options.

\sphinxAtStartPar
This pair style does not write its information to {\hyperref[\detokenize{restart::doc}]{\sphinxcrossref{\DUrole{doc}{binary restart
files}}}}, since it is stored in potential files.  Thus, you need
to re\sphinxhyphen{}specify the pair\_style and pair\_coeff commands in an input script
that reads a restart file.

\sphinxAtStartPar
This pair style can only be used via the \sphinxstyleemphasis{pair} keyword of the
{\hyperref[\detokenize{run_style::doc}]{\sphinxcrossref{\DUrole{doc}{run\_style respa}}}} command.  It does not support the
\sphinxstyleemphasis{inner}, \sphinxstyleemphasis{middle}, \sphinxstyleemphasis{outer} keywords.


\bigskip\hrule\bigskip



\subsection{Restrictions}
\label{\detokenize{pair_tersoff:restrictions}}
\sphinxAtStartPar
This pair style is part of the MANYBODY package.  It is only enabled if
LAMMPS was built with that package.  See the {\hyperref[\detokenize{Build_package::doc}]{\sphinxcrossref{\DUrole{doc}{Build package}}}} page for more info.

\sphinxAtStartPar
This pair style requires the {\hyperref[\detokenize{newton::doc}]{\sphinxcrossref{\DUrole{doc}{newton}}}} setting to be “on”
for pair interactions.

\sphinxAtStartPar
The \sphinxstyleemphasis{shift} keyword is not supported by the \sphinxstyleemphasis{tersoff/gpu},
\sphinxstyleemphasis{tersoff/intel}, \sphinxstyleemphasis{tersoff/kk}, \sphinxstyleemphasis{tersoff/table} or \sphinxstyleemphasis{tersoff/table/omp}
variants.

\sphinxAtStartPar
The \sphinxstyleemphasis{tersoff/intel} pair style is only available when compiling LAMMPS
with the Intel compilers.

\sphinxAtStartPar
The Tersoff potential files provided with LAMMPS (see the potentials
directory) are parameterized for {\hyperref[\detokenize{units::doc}]{\sphinxcrossref{\DUrole{doc}{“metal” units}}}}.  In addition
the pair style supports converting potential parameters on\sphinxhyphen{}the\sphinxhyphen{}fly between
“metal” and “real” units. You can use the \sphinxstyleemphasis{tersoff} pair style variants
with any LAMMPS units setting, but you would need to
create your own Tersoff potential file with coefficients listed in the
appropriate units if your simulation does not use “metal” or “real” units.


\subsection{Related commands}
\label{\detokenize{pair_tersoff:related-commands}}
\sphinxAtStartPar
{\hyperref[\detokenize{pair_coeff::doc}]{\sphinxcrossref{\DUrole{doc}{pair\_coeff}}}}


\subsection{Default}
\label{\detokenize{pair_tersoff:default}}
\sphinxAtStartPar
shift delta = 0.0


\bigskip\hrule\bigskip

\phantomsection\label{\detokenize{pair_tersoff:tersoff-11}}
\sphinxAtStartPar
\sphinxstylestrong{(Tersoff\_1)} J. Tersoff, Phys Rev B, 37, 6991 (1988).

\phantomsection\label{\detokenize{pair_tersoff:albe}}
\sphinxAtStartPar
\sphinxstylestrong{(Albe)} J. Nord, K. Albe, P. Erhart, and K. Nordlund, J. Phys.:
Condens. Matter, 15, 5649(2003).

\phantomsection\label{\detokenize{pair_tersoff:tersoff-21}}
\sphinxAtStartPar
\sphinxstylestrong{(Tersoff\_2)} J. Tersoff, Phys Rev B, 39, 5566 (1989); errata (PRB 41, 3248)

\phantomsection\label{\detokenize{pair_tersoff:mandelli1}}
\sphinxAtStartPar
\sphinxstylestrong{(Mandelli\_1)} D. Mandelli, W. Ouyang, M. Urbakh, and O. Hod, ACS Nano 13(7), 7603\sphinxhyphen{}7609 (2019).

\phantomsection\label{\detokenize{pair_tersoff:ouyang5}}
\sphinxAtStartPar
\sphinxstylestrong{(Ouyang\_1)} W. Ouyang et al., J. Chem. Theory Comput. 16(1), 666\sphinxhyphen{}676 (2020).

\sphinxstepscope

\index{pair\_style tersoff/mod@\spxentry{pair\_style tersoff/mod}}\ignorespaces 
\index{pair\_style tersoff/mod/c@\spxentry{pair\_style tersoff/mod/c}}\ignorespaces 
\index{pair\_style tersoff/mod/gpu@\spxentry{pair\_style tersoff/mod/gpu}}\ignorespaces 
\index{pair\_style tersoff/mod/kk@\spxentry{pair\_style tersoff/mod/kk}}\ignorespaces 
\index{pair\_style tersoff/mod/omp@\spxentry{pair\_style tersoff/mod/omp}}\ignorespaces 
\index{pair\_style tersoff/mod/c/omp@\spxentry{pair\_style tersoff/mod/c/omp}}\ignorespaces 

\section{pair\_style tersoff/mod command}
\label{\detokenize{pair_tersoff_mod:pair-style-tersoff-mod-command}}\label{\detokenize{pair_tersoff_mod:index-5}}\label{\detokenize{pair_tersoff_mod:index-4}}\label{\detokenize{pair_tersoff_mod:index-3}}\label{\detokenize{pair_tersoff_mod:index-2}}\label{\detokenize{pair_tersoff_mod:index-1}}\label{\detokenize{pair_tersoff_mod:index-0}}\label{\detokenize{pair_tersoff_mod::doc}}
\sphinxAtStartPar
Accelerator Variants: \sphinxstyleemphasis{tersoff/mod/gpu}, \sphinxstyleemphasis{tersoff/mod/kk}, \sphinxstyleemphasis{tersoff/mod/omp}


\section{pair\_style tersoff/mod/c command}
\label{\detokenize{pair_tersoff_mod:pair-style-tersoff-mod-c-command}}
\sphinxAtStartPar
Accelerator Variants: \sphinxstyleemphasis{tersoff/mod/c/omp}


\subsection{Syntax}
\label{\detokenize{pair_tersoff_mod:syntax}}
\begin{sphinxVerbatim}[commandchars=\\\{\}]
\PYG{k}{pair\PYGZus{}style}\PYG{+w}{ }\PYG{n}{style}\PYG{+w}{ }\PYG{n}{keywords}\PYG{+w}{ }\PYG{n}{values}
\end{sphinxVerbatim}
\begin{itemize}
\item {} 
\sphinxAtStartPar
style = \sphinxstyleemphasis{tersoff/mod} or \sphinxstyleemphasis{tersoff/mod/c}

\item {} 
\sphinxAtStartPar
keyword = \sphinxstyleemphasis{shift}
\begin{sphinxalltt}
\sphinxstyleemphasis{shift} value = delta
  delta = negative shift in equilibrium bond length
\end{sphinxalltt}

\end{itemize}


\subsection{Examples}
\label{\detokenize{pair_tersoff_mod:examples}}
\begin{sphinxVerbatim}[commandchars=\\\{\}]
\PYG{k}{pair\PYGZus{}style}\PYG{+w}{ }\PYG{n}{tersoff}\PYG{o}{/}\PYG{n}{mod}
\PYG{k}{pair\PYGZus{}coeff}\PYG{+w}{ }\PYG{o}{*}\PYG{+w}{ }\PYG{o}{*}\PYG{+w}{ }\PYG{n}{Si.tersoff.mod}\PYG{+w}{ }\PYG{n}{Si}\PYG{+w}{ }\PYG{n}{Si}

\PYG{k}{pair\PYGZus{}style}\PYG{+w}{ }\PYG{n}{tersoff}\PYG{o}{/}\PYG{n}{mod}\PYG{o}{/}\PYG{n}{c}
\PYG{k}{pair\PYGZus{}coeff}\PYG{+w}{ }\PYG{o}{*}\PYG{+w}{ }\PYG{o}{*}\PYG{+w}{ }\PYG{n}{Si.tersoff.modc}\PYG{+w}{ }\PYG{n}{Si}\PYG{+w}{ }\PYG{n}{Si}
\end{sphinxVerbatim}


\subsection{Description}
\label{\detokenize{pair_tersoff_mod:description}}
\sphinxAtStartPar
The \sphinxstyleemphasis{tersoff/mod} and \sphinxstyleemphasis{tersoff/mod/c} styles computes a bond\sphinxhyphen{}order type
interatomic potential {\hyperref[\detokenize{pair_tersoff_mod:kumagai}]{\sphinxcrossref{\DUrole{std,std-ref}{(Kumagai)}}}} based on a 3\sphinxhyphen{}body Tersoff
potential {\hyperref[\detokenize{pair_tersoff_mod:tersoff-12}]{\sphinxcrossref{\DUrole{std,std-ref}{(Tersoff\_1)}}}}, {\hyperref[\detokenize{pair_tersoff_mod:tersoff-22}]{\sphinxcrossref{\DUrole{std,std-ref}{(Tersoff\_2)}}}} with
modified cutoff function and angular\sphinxhyphen{}dependent term, giving the energy
E of a system of atoms as
\begin{equation*}
\begin{split}E & = \frac{1}{2} \sum_i \sum_{j \neq i} V_{ij} \\
V_{ij} & = f_C(r_{ij} + \delta) \left[ f_R(r_{ij} + \delta) + b_{ij} f_A(r_{ij} + \delta) \right] \\
f_C(r) & = \left\{ \begin{array} {r@{\quad:\quad}l}
  1 & r < R - D \\
  \frac{1}{2} - \frac{9}{16} \sin \left( \frac{\pi}{2} \frac{r-R}{D} \right) - \frac{1}{16} \sin \left( \frac{3\pi}{2} \frac{r-R}{D} \right) &
    R-D < r < R + D \\
  0 & r > R + D
  \end{array} \right. \\
f_R(r) & = A \exp (-\lambda_1 r) \\
f_A(r) & = -B \exp (-\lambda_2 r) \\
b_{ij} & = \left( 1 + {\zeta_{ij}}^\eta \right)^{-\frac{1}{2n}} \\
\zeta_{ij} & = \sum_{k \neq i,j} f_C(r_{ik} + \delta) g(\theta_{ijk})
                 \exp \left[ \alpha (r_{ij} - r_{ik})^\beta \right] \\
g(\theta) & = c_1 + g_o(\theta) g_a(\theta) \\
g_o(\theta) & = \frac{c_2 (h - \cos \theta)^2}{c_3 + (h - \cos \theta)^2} \\
g_a(\theta) & = 1 + c_4 \exp \left[ -c_5 (h - \cos \theta)^2 \right] \\\end{split}
\end{equation*}
\sphinxAtStartPar
where \(f_R\) is a two\sphinxhyphen{}body term and \(f_A\) includes three\sphinxhyphen{}body interactions.
\(\delta\) is an optional negative shift of the
equilibrium bond length, as described below.

\sphinxAtStartPar
The summations in the formula are over all neighbors J and K of atom I
within a cutoff distance = R + D.
The \sphinxstyleemphasis{tersoff/mod/c} style differs from \sphinxstyleemphasis{tersoff/mod} only in the
formulation of the V\_ij term, where it contains an additional c0 term.
\begin{equation*}
\begin{split}V_{ij} = f_C(r_{ij} + \delta) \left[ f_R(r_{ij} + \delta) + b_{ij} f_A(r_{ij} + \delta) + c_0 \right] \\\end{split}
\end{equation*}
\sphinxAtStartPar
The modified cutoff function \(f_C\) proposed by {\hyperref[\detokenize{pair_tersoff_mod:murty}]{\sphinxcrossref{\DUrole{std,std-ref}{(Murty)}}}} and
having a continuous second\sphinxhyphen{}order differential is employed. The
angular\sphinxhyphen{}dependent term \(g(\theta)\) was modified to increase the
flexibility of the potential.

\sphinxAtStartPar
The \sphinxstyleemphasis{tersoff/mod} potential is fitted to both the elastic constants
and melting point by employing the modified Tersoff potential function
form in which the angular\sphinxhyphen{}dependent term is improved. The model
performs extremely well in describing the crystalline, liquid, and
amorphous phases {\hyperref[\detokenize{pair_tersoff_mod:schelling}]{\sphinxcrossref{\DUrole{std,std-ref}{(Schelling)}}}}.

\sphinxAtStartPar
Only a single pair\_coeff command is used with the \sphinxstyleemphasis{tersoff/mod} style
which specifies a Tersoff/MOD potential file with parameters for all
needed elements.  These are mapped to LAMMPS atom types by specifying
N additional arguments after the filename in the pair\_coeff command,
where N is the number of LAMMPS atom types:
\begin{itemize}
\item {} 
\sphinxAtStartPar
filename

\item {} 
\sphinxAtStartPar
N element names = mapping of Tersoff/MOD elements to atom types

\end{itemize}

\sphinxAtStartPar
As an example, imagine the Si.tersoff\_mod file has Tersoff values for Si.
If your LAMMPS simulation has 3 Si atoms types, you would use the following
pair\_coeff command:

\begin{sphinxVerbatim}[commandchars=\\\{\}]
\PYG{k}{pair\PYGZus{}coeff}\PYG{+w}{ }\PYG{o}{*}\PYG{+w}{ }\PYG{o}{*}\PYG{+w}{ }\PYG{n}{Si.tersoff\PYGZus{}mod}\PYG{+w}{ }\PYG{n}{Si}\PYG{+w}{ }\PYG{n}{Si}\PYG{+w}{ }\PYG{n}{Si}
\end{sphinxVerbatim}

\sphinxAtStartPar
The first 2 arguments must be * * so as to span all LAMMPS atom types.
The three Si arguments map LAMMPS atom types 1,2,3 to the Si element
in the Tersoff/MOD file. If a mapping value is specified as NULL, the
mapping is not performed.  This can be used when a \sphinxstyleemphasis{tersoff/mod}
potential is used as part of the \sphinxstyleemphasis{hybrid} pair style. The NULL values
are placeholders for atom types that will be used with other
potentials.

\sphinxAtStartPar
Tersoff/MOD file in the \sphinxstyleemphasis{potentials} directory of the LAMMPS
distribution have a “.tersoff.mod” suffix. Potential files for the
\sphinxstyleemphasis{tersoff/mod/c} style have the suffix “.tersoff.modc”. Lines that are
not blank or comments (starting with \#) define parameters for a triplet
of elements.  The parameters in a single entry correspond to
coefficients in the formulae above:
\begin{itemize}
\item {} 
\sphinxAtStartPar
element 1 (the center atom in a 3\sphinxhyphen{}body interaction)

\item {} 
\sphinxAtStartPar
element 2 (the atom bonded to the center atom)

\item {} 
\sphinxAtStartPar
element 3 (the atom influencing the 1\sphinxhyphen{}2 bond in a bond\sphinxhyphen{}order sense)

\item {} 
\sphinxAtStartPar
\(\beta\)

\item {} 
\sphinxAtStartPar
\(\alpha\)

\item {} 
\sphinxAtStartPar
h

\item {} 
\sphinxAtStartPar
\(\eta\)

\item {} 
\sphinxAtStartPar
\(\beta_{ters}\) = 1 (dummy parameter)

\item {} 
\sphinxAtStartPar
\(\lambda_2\) (1/distance units)

\item {} 
\sphinxAtStartPar
B (energy units)

\item {} 
\sphinxAtStartPar
R (distance units)

\item {} 
\sphinxAtStartPar
D (distance units)

\item {} 
\sphinxAtStartPar
\(\lambda_1\) (1/distance units)

\item {} 
\sphinxAtStartPar
A (energy units)

\item {} 
\sphinxAtStartPar
n

\item {} 
\sphinxAtStartPar
c1

\item {} 
\sphinxAtStartPar
c2

\item {} 
\sphinxAtStartPar
c3

\item {} 
\sphinxAtStartPar
c4

\item {} 
\sphinxAtStartPar
c5

\item {} 
\sphinxAtStartPar
c0 (energy units, tersoff/mod/c only)

\end{itemize}

\sphinxAtStartPar
The n, \(\eta\), \(\lambda_2\), B, \(\lambda_1\), and A parameters are only used for
two\sphinxhyphen{}body interactions.  The \(\beta\), \(\alpha\), c1, c2, c3, c4, c5, h
parameters are only used for three\sphinxhyphen{}body interactions. The R and D
parameters are used for both two\sphinxhyphen{}body and three\sphinxhyphen{}body interactions.
The c0 term applies to \sphinxstyleemphasis{tersoff/mod/c} only. The non\sphinxhyphen{}annotated
parameters are unitless.

\sphinxAtStartPar
The Tersoff/MOD potential file must contain entries for all the elements
listed in the pair\_coeff command.  It can also contain entries for
additional elements not being used in a particular simulation; LAMMPS
ignores those entries.

\sphinxAtStartPar
For a single\sphinxhyphen{}element simulation, only a single entry is required
(e.g. SiSiSi). As annotated above, the first element in the entry is
the center atom in a three\sphinxhyphen{}body interaction and it is bonded to the
second atom and the bond is influenced by the third atom.  Thus an entry
for SiSiSi means Si bonded to a Si with another Si atom influencing the bond.

\sphinxAtStartPar
The \sphinxstyleemphasis{shift} keyword computes the energy E of a system of atoms, whose formula
is the same as the Tersoff potential. The only modification is that the original
equilibrium bond length ( \(r_0\)) of the system is shifted to \(r_0-\delta\).
The minus sign arises because each radial distance \(r\) is replaced by \(r+\delta\).
More information on this option is given on the main {\hyperref[\detokenize{pair_tersoff::doc}]{\sphinxcrossref{\DUrole{doc}{pair\_tersoff}}}} page.


\bigskip\hrule\bigskip


\sphinxAtStartPar
Styles with a \sphinxstyleemphasis{gpu}, \sphinxstyleemphasis{intel}, \sphinxstyleemphasis{kk}, \sphinxstyleemphasis{omp}, or \sphinxstyleemphasis{opt} suffix are
functionally the same as the corresponding style without the suffix.
They have been optimized to run faster, depending on your available
hardware, as discussed on the {\hyperref[\detokenize{Speed_packages::doc}]{\sphinxcrossref{\DUrole{doc}{Accelerator packages}}}}
page.  The accelerated styles take the same arguments and should
produce the same results, except for round\sphinxhyphen{}off and precision issues.

\sphinxAtStartPar
These accelerated styles are part of the GPU, INTEL, KOKKOS,
OPENMP, and OPT packages, respectively.  They are only enabled if
LAMMPS was built with those packages.  See the {\hyperref[\detokenize{Build_package::doc}]{\sphinxcrossref{\DUrole{doc}{Build package}}}} page for more info.

\sphinxAtStartPar
You can specify the accelerated styles explicitly in your input script
by including their suffix, or you can use the {\hyperref[\detokenize{Run_options::doc}]{\sphinxcrossref{\DUrole{doc}{\sphinxhyphen{}suffix command\sphinxhyphen{}line switch}}}} when you invoke LAMMPS, or you can use the
{\hyperref[\detokenize{suffix::doc}]{\sphinxcrossref{\DUrole{doc}{suffix}}}} command in your input script.

\sphinxAtStartPar
See the {\hyperref[\detokenize{Speed_packages::doc}]{\sphinxcrossref{\DUrole{doc}{Accelerator packages}}}} page for more
instructions on how to use the accelerated styles effectively.


\bigskip\hrule\bigskip



\subsection{Mixing, shift, table, tail correction, restart, rRESPA info}
\label{\detokenize{pair_tersoff_mod:mixing-shift-table-tail-correction-restart-rrespa-info}}
\sphinxAtStartPar
This pair style does not support the {\hyperref[\detokenize{pair_modify::doc}]{\sphinxcrossref{\DUrole{doc}{pair\_modify}}}}
shift, table, and tail options.

\sphinxAtStartPar
This pair style does not write its information to {\hyperref[\detokenize{restart::doc}]{\sphinxcrossref{\DUrole{doc}{binary restart files}}}}, since it is stored in potential files.  Thus, you
need to re\sphinxhyphen{}specify the pair\_style and pair\_coeff commands in an input
script that reads a restart file.

\sphinxAtStartPar
This pair style can only be used via the \sphinxstyleemphasis{pair} keyword of the
{\hyperref[\detokenize{run_style::doc}]{\sphinxcrossref{\DUrole{doc}{run\_style respa}}}} command.  It does not support the
\sphinxstyleemphasis{inner}, \sphinxstyleemphasis{middle}, \sphinxstyleemphasis{outer} keywords.


\bigskip\hrule\bigskip



\subsection{Restrictions}
\label{\detokenize{pair_tersoff_mod:restrictions}}
\sphinxAtStartPar
This pair style is part of the MANYBODY package.  It is only enabled
if LAMMPS was built with that package.  See the {\hyperref[\detokenize{Build_package::doc}]{\sphinxcrossref{\DUrole{doc}{Build package}}}} page for more info.

\sphinxAtStartPar
This pair style requires the {\hyperref[\detokenize{newton::doc}]{\sphinxcrossref{\DUrole{doc}{newton}}}} setting to be “on”
for pair interactions.

\sphinxAtStartPar
The \sphinxstyleemphasis{shift} keyword is not supported by the \sphinxstyleemphasis{tersoff/gpu},
\sphinxstyleemphasis{tersoff/intel}, \sphinxstyleemphasis{tersoff/kk}, \sphinxstyleemphasis{tersoff/table} or \sphinxstyleemphasis{tersoff/table/omp}
variants.

\sphinxAtStartPar
The \sphinxstyleemphasis{tersoff/mod} potential files provided with LAMMPS (see the potentials
directory) are parameterized for metal {\hyperref[\detokenize{units::doc}]{\sphinxcrossref{\DUrole{doc}{units}}}}.  You can
use the \sphinxstyleemphasis{tersoff/mod} pair style with any LAMMPS units, but you would need to
create your own Tersoff/MOD potential file with coefficients listed in the
appropriate units if your simulation does not use “metal” units.


\subsection{Related commands}
\label{\detokenize{pair_tersoff_mod:related-commands}}
\sphinxAtStartPar
{\hyperref[\detokenize{pair_coeff::doc}]{\sphinxcrossref{\DUrole{doc}{pair\_coeff}}}}


\subsection{Default}
\label{\detokenize{pair_tersoff_mod:default}}
\sphinxAtStartPar
none


\bigskip\hrule\bigskip

\phantomsection\label{\detokenize{pair_tersoff_mod:kumagai}}
\sphinxAtStartPar
\sphinxstylestrong{(Kumagai)} T. Kumagai, S. Izumi, S. Hara, S. Sakai,
Comp. Mat. Science, 39, 457 (2007).

\phantomsection\label{\detokenize{pair_tersoff_mod:tersoff-12}}
\sphinxAtStartPar
\sphinxstylestrong{(Tersoff\_1)} J. Tersoff, Phys Rev B, 37, 6991 (1988).

\phantomsection\label{\detokenize{pair_tersoff_mod:tersoff-22}}
\sphinxAtStartPar
\sphinxstylestrong{(Tersoff\_2)} J. Tersoff, Phys Rev B, 38, 9902 (1988).

\phantomsection\label{\detokenize{pair_tersoff_mod:murty}}
\sphinxAtStartPar
\sphinxstylestrong{(Murty)} M.V.R. Murty, H.A. Atwater, Phys Rev B, 51, 4889 (1995).

\phantomsection\label{\detokenize{pair_tersoff_mod:schelling}}
\sphinxAtStartPar
\sphinxstylestrong{(Schelling)} Patrick K. Schelling, Comp. Mat. Science, 44, 274 (2008).

\sphinxstepscope

\index{pair\_style tersoff/zbl@\spxentry{pair\_style tersoff/zbl}}\ignorespaces 
\index{pair\_style tersoff/zbl/gpu@\spxentry{pair\_style tersoff/zbl/gpu}}\ignorespaces 
\index{pair\_style tersoff/zbl/kk@\spxentry{pair\_style tersoff/zbl/kk}}\ignorespaces 
\index{pair\_style tersoff/zbl/omp@\spxentry{pair\_style tersoff/zbl/omp}}\ignorespaces 

\section{pair\_style tersoff/zbl command}
\label{\detokenize{pair_tersoff_zbl:pair-style-tersoff-zbl-command}}\label{\detokenize{pair_tersoff_zbl:index-3}}\label{\detokenize{pair_tersoff_zbl:index-2}}\label{\detokenize{pair_tersoff_zbl:index-1}}\label{\detokenize{pair_tersoff_zbl:index-0}}\label{\detokenize{pair_tersoff_zbl::doc}}
\sphinxAtStartPar
Accelerator Variants: \sphinxstyleemphasis{tersoff/zbl/gpu}, \sphinxstyleemphasis{tersoff/zbl/kk}, \sphinxstyleemphasis{tersoff/zbl/omp}


\subsection{Syntax}
\label{\detokenize{pair_tersoff_zbl:syntax}}
\begin{sphinxVerbatim}[commandchars=\\\{\}]
\PYG{k}{pair\PYGZus{}style}\PYG{+w}{ }\PYG{n}{tersoff}\PYG{o}{/}\PYG{n}{zbl}\PYG{+w}{ }\PYG{n}{keywords}\PYG{+w}{ }\PYG{n}{values}
\end{sphinxVerbatim}
\begin{itemize}
\item {} 
\sphinxAtStartPar
keyword = \sphinxstyleemphasis{shift}
\begin{sphinxalltt}
\sphinxstyleemphasis{shift} value = delta
  delta = negative shift in equilibrium bond length
\end{sphinxalltt}

\end{itemize}


\subsection{Examples}
\label{\detokenize{pair_tersoff_zbl:examples}}
\begin{sphinxVerbatim}[commandchars=\\\{\}]
\PYG{k}{pair\PYGZus{}style}\PYG{+w}{ }\PYG{n}{tersoff}\PYG{o}{/}\PYG{n}{zbl}
\PYG{k}{pair\PYGZus{}coeff}\PYG{+w}{ }\PYG{o}{*}\PYG{+w}{ }\PYG{o}{*}\PYG{+w}{ }\PYG{n}{SiC.tersoff.zbl}\PYG{+w}{ }\PYG{n}{Si}\PYG{+w}{ }\PYG{n}{C}\PYG{+w}{ }\PYG{n}{Si}
\end{sphinxVerbatim}


\subsection{Description}
\label{\detokenize{pair_tersoff_zbl:description}}
\sphinxAtStartPar
The \sphinxstyleemphasis{tersoff/zbl} style computes a 3\sphinxhyphen{}body Tersoff potential
{\hyperref[\detokenize{pair_tersoff_zbl:zbl-tersoff-1}]{\sphinxcrossref{\DUrole{std,std-ref}{(Tersoff\_1)}}}} with a close\sphinxhyphen{}separation pairwise modification
based on a Coulomb potential and the Ziegler\sphinxhyphen{}Biersack\sphinxhyphen{}Littmark
universal screening function {\hyperref[\detokenize{pair_tersoff_zbl:zbl-zbl}]{\sphinxcrossref{\DUrole{std,std-ref}{(ZBL)}}}}, giving the energy E of a
system of atoms as
\begin{equation*}
\begin{split} E & = \frac{1}{2} \sum_i \sum_{j \neq i} V_{ij} \\
 V_{ij} & =  (1 - f_F(r_{ij} + \delta)) V^{ZBL}(r_{ij} + \delta)
             + f_F(r_{ij} + \delta) V^{Tersoff}(r_{ij} + \delta) \\
 f_F(r) & =  \frac{1}{1 + e^{-A_F(r - r_C)}}\\
 \\
 \\
 V^{ZBL}(r) & = \frac{1}{4\pi\epsilon_0} \frac{Z_1 Z_2 \,e^2}{r} \phi(r/a) \\
a & = \frac{0.8854\,a_0}{Z_{1}^{0.23} + Z_{2}^{0.23}}\\
\phi(x) & =  0.1818e^{-3.2x} + 0.5099e^{-0.9423x} + 0.2802e^{-0.4029x} + 0.02817e^{-0.2016x}\\
\\
\\
V^{Tersoff}(r) & = f_C(r) \left[ f_R(r) + b_{ij} f_A(r) \right] \\
f_C(r) & = \left\{ \begin{array} {r@{\quad:\quad}l}
  1 & r < R - D \\
  \frac{1}{2} - \frac{1}{2} \sin \left( \frac{\pi}{2} \frac{r-R}{D} \right) &
    R-D < r < R + D \\
  0 & r > R + D
  \end{array} \right. \\
f_R(r) & = A \exp (-\lambda_1 r) \\
f_A(r) & = -B \exp (-\lambda_2 r) \\
b_{ij} & = \left( 1 + \beta^n {\zeta_{ij}}^n \right)^{-\frac{1}{2n}} \\
\zeta_{ij} & = \sum_{k \neq i,j} f_C(r_{ik} + \delta) g(\theta_{ijk})
                 \exp \left[ {\lambda_3}^m (r_{ij} - r_{ik})^m \right] \\
g(\theta) & =  \gamma_{ijk} \left( 1 + \frac{c^2}{d^2} -
                \frac{c^2}{\left[ d^2 + (\cos \theta - \cos \theta_0)^2\right]} \right)\end{split}
\end{equation*}
\sphinxAtStartPar
The \(f_F\) term is a fermi\sphinxhyphen{}like function used to smoothly connect the ZBL
repulsive potential with the Tersoff potential.  There are 2
parameters used to adjust it: \(A_F\) and \(r_C\).  \(A_F\)
controls how “sharp”
the transition is between the two, and \(r_C\) is essentially the cutoff
for the ZBL potential.

\sphinxAtStartPar
For the ZBL portion, there are two terms. The first is the Coulomb
repulsive term, with Z1, Z2 as the number of protons in each nucleus,
e as the electron charge (1 for metal and real units) and \(\epsilon_0\)
as the permittivity of vacuum.  The second part is the ZBL universal
screening function, with a0 being the Bohr radius (typically 0.529
Angstroms), and the remainder of the coefficients provided by the
original paper.  This screening function should be applicable to most
systems.  However, it is only accurate for small separations
(i.e. less than 1 Angstrom).

\sphinxAtStartPar
For the Tersoff portion, \(f_R\) is a two\sphinxhyphen{}body term and \(f_A\)
includes
three\sphinxhyphen{}body interactions. The summations in the formula are over all
neighbors J and K of atom I within a cutoff distance = R + D.

\sphinxAtStartPar
\(\delta\) is an optional negative shift of the
equilibrium bond length, as described below.

\sphinxAtStartPar
Only a single pair\_coeff command is used with the \sphinxstyleemphasis{tersoff/zbl} style
which specifies a Tersoff/ZBL potential file with parameters for all
needed elements.  These are mapped to LAMMPS atom types by specifying
N additional arguments after the filename in the pair\_coeff command,
where N is the number of LAMMPS atom types:
\begin{itemize}
\item {} 
\sphinxAtStartPar
filename

\item {} 
\sphinxAtStartPar
N element names = mapping of Tersoff/ZBL elements to atom types

\end{itemize}

\sphinxAtStartPar
See the {\hyperref[\detokenize{pair_coeff::doc}]{\sphinxcrossref{\DUrole{doc}{pair\_coeff}}}} page for alternate ways
to specify the path for the potential file.

\sphinxAtStartPar
As an example, imagine the SiC.tersoff.zbl file has Tersoff/ZBL values
for Si and C.  If your LAMMPS simulation has 4 atoms types and you
want the first 3 to be Si, and the fourth to be C, you would use the
following pair\_coeff command:

\begin{sphinxVerbatim}[commandchars=\\\{\}]
\PYG{k}{pair\PYGZus{}coeff}\PYG{+w}{ }\PYG{o}{*}\PYG{+w}{ }\PYG{o}{*}\PYG{+w}{ }\PYG{n}{SiC.tersoff}\PYG{+w}{ }\PYG{n}{Si}\PYG{+w}{ }\PYG{n}{Si}\PYG{+w}{ }\PYG{n}{Si}\PYG{+w}{ }\PYG{n}{C}
\end{sphinxVerbatim}

\sphinxAtStartPar
The first 2 arguments must be * * so as to span all LAMMPS atom types.
The first three Si arguments map LAMMPS atom types 1,2,3 to the Si
element in the Tersoff/ZBL file.  The final C argument maps LAMMPS
atom type 4 to the C element in the Tersoff/ZBL file.  If a mapping
value is specified as NULL, the mapping is not performed.  This can be
used when a \sphinxstyleemphasis{tersoff/zbl} potential is used as part of the \sphinxstyleemphasis{hybrid}
pair style.  The NULL values are placeholders for atom types that will
be used with other potentials.

\sphinxAtStartPar
Tersoff/ZBL files in the \sphinxstyleemphasis{potentials} directory of the LAMMPS
distribution have a “.tersoff.zbl” suffix.  Lines that are not blank
or comments (starting with \#) define parameters for a triplet of
elements.  The parameters in a single entry correspond to coefficients
in the formula above:
\begin{itemize}
\item {} 
\sphinxAtStartPar
element 1 (the center atom in a 3\sphinxhyphen{}body interaction)

\item {} 
\sphinxAtStartPar
element 2 (the atom bonded to the center atom)

\item {} 
\sphinxAtStartPar
element 3 (the atom influencing the 1\sphinxhyphen{}2 bond in a bond\sphinxhyphen{}order sense)

\item {} 
\sphinxAtStartPar
m

\item {} 
\sphinxAtStartPar
\(\gamma\)

\item {} 
\sphinxAtStartPar
\(\lambda_3\) (1/distance units)

\item {} 
\sphinxAtStartPar
c

\item {} 
\sphinxAtStartPar
d

\item {} 
\sphinxAtStartPar
\(\cos\theta_0\) (can be a value \textless{} \sphinxhyphen{}1 or \textgreater{} 1)

\item {} 
\sphinxAtStartPar
n

\item {} 
\sphinxAtStartPar
\(\beta\)

\item {} 
\sphinxAtStartPar
\(\lambda_2\) (1/distance units)

\item {} 
\sphinxAtStartPar
B (energy units)

\item {} 
\sphinxAtStartPar
R (distance units)

\item {} 
\sphinxAtStartPar
D (distance units)

\item {} 
\sphinxAtStartPar
\(\lambda_1\) (1/distance units)

\item {} 
\sphinxAtStartPar
A (energy units)

\item {} 
\sphinxAtStartPar
\(Z_i\)

\item {} 
\sphinxAtStartPar
\(Z_j\)

\item {} 
\sphinxAtStartPar
ZBLcut (distance units)

\item {} 
\sphinxAtStartPar
ZBLexpscale (1/distance units)

\end{itemize}

\sphinxAtStartPar
The n, \(\beta\), \(\lambda_2\), B, \(\lambda_1\), and A
parameters are only used for
two\sphinxhyphen{}body interactions.  The m, \(\gamma\), \(\lambda_3\), c, d,
and \(\cos\theta_0\)
parameters are only used for three\sphinxhyphen{}body interactions. The R and D
parameters are used for both two\sphinxhyphen{}body and three\sphinxhyphen{}body interactions. The
\(Z_i\), \(Z_j\), ZBLcut, ZBLexpscale parameters are used in the
ZBL repulsive
portion of the potential and in the Fermi\sphinxhyphen{}like function.  The
non\sphinxhyphen{}annotated parameters are unitless.  The value of m must be 3 or 1.

\sphinxAtStartPar
The Tersoff/ZBL potential file must contain entries for all the
elements listed in the pair\_coeff command.  It can also contain
entries for additional elements not being used in a particular
simulation; LAMMPS ignores those entries.

\sphinxAtStartPar
For a single\sphinxhyphen{}element simulation, only a single entry is required
(e.g. SiSiSi).  For a two\sphinxhyphen{}element simulation, the file must contain 8
entries (for SiSiSi, SiSiC, SiCSi, SiCC, CSiSi, CSiC, CCSi, CCC), that
specify Tersoff parameters for all permutations of the two elements
interacting in three\sphinxhyphen{}body configurations.  Thus for 3 elements, 27
entries would be required, etc.

\sphinxAtStartPar
As annotated above, the first element in the entry is the center atom
in a three\sphinxhyphen{}body interaction and it is bonded to the second atom and the
bond is influenced by the third atom.  Thus an entry for SiCC means Si
bonded to a C with another C atom influencing the bond.  Thus
three\sphinxhyphen{}body parameters for SiCSi and SiSiC entries will not, in
general, be the same.  The parameters used for the two\sphinxhyphen{}body
interaction come from the entry where the second element is repeated.
Thus the two\sphinxhyphen{}body parameters for Si interacting with C, comes from the
SiCC entry.

\sphinxAtStartPar
The parameters used for a particular
three\sphinxhyphen{}body interaction come from the entry with the corresponding
three elements.  The parameters used only for two\sphinxhyphen{}body interactions
(n, \(\beta\), \(\lambda_2\), B, \(\lambda_1\), and A)
in entries whose second and third
element are different (e.g. SiCSi) are not used for anything and can
be set to 0.0 if desired.

\sphinxAtStartPar
Note that the twobody parameters in entries such as SiCC and CSiSi
are often the same, due to the common use of symmetric mixing rules,
but this is not always the case. For example, the beta and n parameters in
Tersoff\_2 {\hyperref[\detokenize{pair_tersoff_zbl:zbl-tersoff-2}]{\sphinxcrossref{\DUrole{std,std-ref}{(Tersoff\_2)}}}} are not symmetric.

\sphinxAtStartPar
We chose the above form so as to enable users to define all commonly
used variants of the Tersoff portion of the potential.  In particular,
our form reduces to the original Tersoff form when m = 3 and gamma =
1, while it reduces to the form of {\hyperref[\detokenize{pair_tersoff_zbl:zbl-albe}]{\sphinxcrossref{\DUrole{std,std-ref}{Albe et al.}}}} when beta = 1
and m = 1.  Note that in the current Tersoff implementation in LAMMPS,
m must be specified as either 3 or 1.  Tersoff used a slightly
different but equivalent form for alloys, which we will refer to as
Tersoff\_2 potential {\hyperref[\detokenize{pair_tersoff_zbl:zbl-tersoff-2}]{\sphinxcrossref{\DUrole{std,std-ref}{(Tersoff\_2)}}}}.

\sphinxAtStartPar
LAMMPS parameter values for Tersoff\_2 can be obtained as follows:
\(\gamma = \omega_{ijk}\), \(\lambda_3 = 0\) and the value of
m has no effect.  The parameters for species i and j can be calculated
using the Tersoff\_2 mixing rules:
\begin{equation*}
\begin{split}\lambda_1^{i,j} & = \frac{1}{2}(\lambda_1^i + \lambda_1^j)\\
\lambda_2^{i,j} & = \frac{1}{2}(\lambda_2^i + \lambda_2^j)\\
A_{i,j} & = (A_{i}A_{j})^{1/2}\\
B_{i,j} & = \chi_{ij}(B_{i}B_{j})^{1/2}\\
R_{i,j} & = (R_{i}R_{j})^{1/2}\\
S_{i,j} & = (S_{i}S_{j})^{1/2}\\\end{split}
\end{equation*}
\sphinxAtStartPar
Tersoff\_2 parameters R and S must be converted to the LAMMPS
parameters R and D (R is different in both forms), using the following
relations: R=(R’+S’)/2 and D=(S’\sphinxhyphen{}R’)/2, where the primes indicate the
Tersoff\_2 parameters.

\sphinxAtStartPar
In the potentials directory, the file SiCGe.tersoff provides the
LAMMPS parameters for Tersoff’s various versions of Si, as well as his
alloy parameters for Si, C, and Ge. This file can be used for pure Si,
(three different versions), pure C, pure Ge, binary SiC, and binary
SiGe.  LAMMPS will generate an error if this file is used with any
combination involving C and Ge, since there are no entries for the GeC
interactions (Tersoff did not publish parameters for this
cross\sphinxhyphen{}interaction.)  Tersoff files are also provided for the SiC alloy
(SiC.tersoff) and the GaN (GaN.tersoff) alloys.

\sphinxAtStartPar
Many thanks to Rutuparna Narulkar, David Farrell, and Xiaowang Zhou
for helping clarify how Tersoff parameters for alloys have been
defined in various papers.  Also thanks to Ram Devanathan for
providing the base ZBL implementation.

\sphinxAtStartPar
The \sphinxstyleemphasis{shift} keyword computes the energy E of a system of atoms, whose
formula is the same as the Tersoff potential. The only modification is
that the original equilibrium bond length ( \(r_0\)) of the system
is shifted to \(r_0-\delta\).  The minus sign arises because each
radial distance \(r\) is replaced by \(r+\delta\).  More
information on this option is given on the main {\hyperref[\detokenize{pair_tersoff::doc}]{\sphinxcrossref{\DUrole{doc}{pair\_tersoff}}}} page.


\bigskip\hrule\bigskip


\sphinxAtStartPar
Styles with a \sphinxstyleemphasis{gpu}, \sphinxstyleemphasis{intel}, \sphinxstyleemphasis{kk}, \sphinxstyleemphasis{omp}, or \sphinxstyleemphasis{opt} suffix are
functionally the same as the corresponding style without the suffix.
They have been optimized to run faster, depending on your available
hardware, as discussed on the {\hyperref[\detokenize{Speed_packages::doc}]{\sphinxcrossref{\DUrole{doc}{Accelerator packages}}}}
page.  The accelerated styles take the same arguments and should
produce the same results, except for round\sphinxhyphen{}off and precision issues.

\sphinxAtStartPar
These accelerated styles are part of the GPU, INTEL, KOKKOS,
OPENMP, and OPT packages, respectively.  They are only enabled if
LAMMPS was built with those packages.  See the {\hyperref[\detokenize{Build_package::doc}]{\sphinxcrossref{\DUrole{doc}{Build package}}}} page for more info.

\sphinxAtStartPar
You can specify the accelerated styles explicitly in your input script
by including their suffix, or you can use the {\hyperref[\detokenize{Run_options::doc}]{\sphinxcrossref{\DUrole{doc}{\sphinxhyphen{}suffix command\sphinxhyphen{}line switch}}}} when you invoke LAMMPS, or you can use the
{\hyperref[\detokenize{suffix::doc}]{\sphinxcrossref{\DUrole{doc}{suffix}}}} command in your input script.

\sphinxAtStartPar
See the {\hyperref[\detokenize{Speed_packages::doc}]{\sphinxcrossref{\DUrole{doc}{Accelerator packages}}}} page for more
instructions on how to use the accelerated styles effectively.


\bigskip\hrule\bigskip



\subsection{Mixing, shift, table, tail correction, restart, rRESPA info}
\label{\detokenize{pair_tersoff_zbl:mixing-shift-table-tail-correction-restart-rrespa-info}}
\sphinxAtStartPar
For atom type pairs I,J and I != J, where types I and J correspond to
two different element types, mixing is performed by LAMMPS as
described above from values in the potential file.

\sphinxAtStartPar
This pair style does not support the {\hyperref[\detokenize{pair_modify::doc}]{\sphinxcrossref{\DUrole{doc}{pair\_modify}}}}
shift, table, and tail options.

\sphinxAtStartPar
This pair style does not write its information to {\hyperref[\detokenize{restart::doc}]{\sphinxcrossref{\DUrole{doc}{binary restart
files}}}}, since it is stored in potential files.  Thus, you need
to re\sphinxhyphen{}specify the pair\_style and pair\_coeff commands in an input script
that reads a restart file.

\sphinxAtStartPar
This pair style can only be used via the \sphinxstyleemphasis{pair} keyword of the
{\hyperref[\detokenize{run_style::doc}]{\sphinxcrossref{\DUrole{doc}{run\_style respa}}}} command.  It does not support the
\sphinxstyleemphasis{inner}, \sphinxstyleemphasis{middle}, \sphinxstyleemphasis{outer} keywords.


\bigskip\hrule\bigskip



\subsection{Restrictions}
\label{\detokenize{pair_tersoff_zbl:restrictions}}
\sphinxAtStartPar
This pair style is part of the MANYBODY package.  It is only enabled if
LAMMPS was built with that package.  See the {\hyperref[\detokenize{Build_package::doc}]{\sphinxcrossref{\DUrole{doc}{Build package}}}} page for more info.

\sphinxAtStartPar
This pair style requires the {\hyperref[\detokenize{newton::doc}]{\sphinxcrossref{\DUrole{doc}{newton}}}} setting to be “on”
for pair interactions.

\sphinxAtStartPar
The \sphinxstyleemphasis{shift} keyword is currently not supported for the \sphinxstyleemphasis{tersoff/gpu} and
\sphinxstyleemphasis{tersoff/kk} variants of this pair style.

\sphinxAtStartPar
The tersoff/zbl potential files provided with LAMMPS (see the potentials
directory) are parameterized for {\hyperref[\detokenize{units::doc}]{\sphinxcrossref{\DUrole{doc}{“metal” units}}}}.  Also the
pair style supports converting potential file parameters on\sphinxhyphen{}the\sphinxhyphen{}fly
between “metal” and “real” units.  You can use the tersoff/zbl pair
style with any LAMMPS units, but you would need to create your own
tersoff/zbl potential file with coefficients listed in the appropriate
units if your simulation does not use “metal” or “real” units.


\subsection{Related commands}
\label{\detokenize{pair_tersoff_zbl:related-commands}}
\sphinxAtStartPar
{\hyperref[\detokenize{pair_coeff::doc}]{\sphinxcrossref{\DUrole{doc}{pair\_coeff}}}}


\subsection{Default}
\label{\detokenize{pair_tersoff_zbl:default}}
\sphinxAtStartPar
none


\bigskip\hrule\bigskip

\phantomsection\label{\detokenize{pair_tersoff_zbl:zbl-tersoff-1}}
\sphinxAtStartPar
\sphinxstylestrong{(Tersoff\_1)} J. Tersoff, Phys Rev B, 37, 6991 (1988).

\phantomsection\label{\detokenize{pair_tersoff_zbl:zbl-zbl}}
\sphinxAtStartPar
\sphinxstylestrong{(ZBL)} J.F. Ziegler, J.P. Biersack, U. Littmark, ‘Stopping and Ranges
of Ions in Matter’ Vol 1, 1985, Pergamon Press.

\phantomsection\label{\detokenize{pair_tersoff_zbl:zbl-albe}}
\sphinxAtStartPar
\sphinxstylestrong{(Albe)} J. Nord, K. Albe, P. Erhart and K. Nordlund, J. Phys.:
Condens. Matter, 15, 5649(2003).

\phantomsection\label{\detokenize{pair_tersoff_zbl:zbl-tersoff-2}}
\sphinxAtStartPar
\sphinxstylestrong{(Tersoff\_2)} J. Tersoff, Phys Rev B, 39, 5566 (1989); errata (PRB 41, 3248)

\sphinxstepscope

\index{pair\_style thole@\spxentry{pair\_style thole}}\ignorespaces 
\index{pair\_style lj/cut/thole/long@\spxentry{pair\_style lj/cut/thole/long}}\ignorespaces 
\index{pair\_style lj/cut/thole/long/omp@\spxentry{pair\_style lj/cut/thole/long/omp}}\ignorespaces 

\section{pair\_style thole command}
\label{\detokenize{pair_thole:pair-style-thole-command}}\label{\detokenize{pair_thole:index-2}}\label{\detokenize{pair_thole:index-1}}\label{\detokenize{pair_thole:index-0}}\label{\detokenize{pair_thole::doc}}

\section{pair\_style lj/cut/thole/long command}
\label{\detokenize{pair_thole:pair-style-lj-cut-thole-long-command}}
\sphinxAtStartPar
Accelerator Variants: \sphinxstyleemphasis{lj/cut/thole/long/omp}


\subsection{Syntax}
\label{\detokenize{pair_thole:syntax}}
\begin{sphinxVerbatim}[commandchars=\\\{\}]
\PYG{k}{pair\PYGZus{}style}\PYG{+w}{ }\PYG{n}{style}\PYG{+w}{ }\PYG{n}{args}
\end{sphinxVerbatim}
\begin{itemize}
\item {} 
\sphinxAtStartPar
style = \sphinxstyleemphasis{thole} or \sphinxstyleemphasis{lj/cut/thole/long}

\item {} 
\sphinxAtStartPar
args = list of arguments for a particular style

\end{itemize}
\begin{sphinxalltt}
\sphinxstyleemphasis{thole} args = damp cutoff
  damp = global damping parameter
  cutoff = global cutoff (distance units)
\sphinxstyleemphasis{lj/cut/thole/long} args = damp cutoff (cutoff2)
  damp = global damping parameter
  cutoff = global cutoff for LJ (and Thole if only 1 arg) (distance units)
  cutoff2 = global cutoff for Thole (optional) (distance units)
\end{sphinxalltt}


\subsection{Examples}
\label{\detokenize{pair_thole:examples}}
\begin{sphinxVerbatim}[commandchars=\\\{\}]
\PYG{k}{pair\PYGZus{}style}\PYG{+w}{ }\PYG{n}{hybrid}\PYG{o}{/}\PYG{n}{overlay}\PYG{+w}{ }\PYG{n}{...}\PYG{+w}{ }\PYG{n}{thole}\PYG{+w}{ }\PYG{l+m}{2.6}\PYG{+w}{ }\PYG{l+m}{12.0}
\PYG{k}{pair\PYGZus{}coeff}\PYG{+w}{ }\PYG{l+m}{1}\PYG{+w}{ }\PYG{l+m}{1}\PYG{+w}{ }\PYG{n}{thole}\PYG{+w}{ }\PYG{l+m}{1.0}
\PYG{k}{pair\PYGZus{}coeff}\PYG{+w}{ }\PYG{l+m}{1}\PYG{+w}{ }\PYG{l+m}{2}\PYG{+w}{ }\PYG{n}{thole}\PYG{+w}{ }\PYG{l+m}{1.0}\PYG{+w}{ }\PYG{l+m}{2.6}\PYG{+w}{ }\PYG{l+m}{10.0}
\PYG{k}{pair\PYGZus{}coeff}\PYG{+w}{ }\PYG{o}{*}\PYG{+w}{ }\PYG{l+m}{2}\PYG{+w}{ }\PYG{n}{thole}\PYG{+w}{ }\PYG{l+m}{1.0}\PYG{+w}{ }\PYG{l+m}{2.6}

\PYG{k}{pair\PYGZus{}style}\PYG{+w}{ }\PYG{n}{lj}\PYG{o}{/}\PYG{n}{cut}\PYG{o}{/}\PYG{n}{thole}\PYG{o}{/}\PYG{n}{long}\PYG{+w}{ }\PYG{l+m}{2.6}\PYG{+w}{ }\PYG{l+m}{12.0}
\end{sphinxVerbatim}

\sphinxAtStartPar
Example input scripts available: examples/PACKAGES/drude


\subsection{Description}
\label{\detokenize{pair_thole:description}}
\sphinxAtStartPar
The \sphinxstyleemphasis{thole} pair styles are meant to be used with force fields that
include explicit polarization through Drude dipoles.  This link
describes how to use the {\hyperref[\detokenize{Howto_drude::doc}]{\sphinxcrossref{\DUrole{doc}{thermalized Drude oscillator model}}}} in LAMMPS and polarizable models in LAMMPS are
discussed on the {\hyperref[\detokenize{Howto_polarizable::doc}]{\sphinxcrossref{\DUrole{doc}{Howto polarizable}}}} doc page.

\sphinxAtStartPar
The \sphinxstyleemphasis{thole} pair style should be used as a sub\sphinxhyphen{}style within in the
{\hyperref[\detokenize{pair_hybrid::doc}]{\sphinxcrossref{\DUrole{doc}{pair\_style hybrid/overlay}}}} command, in conjunction with a
main pair style including Coulomb interactions, i.e. any pair style
containing \sphinxstyleemphasis{coul/cut} or \sphinxstyleemphasis{coul/long} in its style name.

\sphinxAtStartPar
The \sphinxstyleemphasis{lj/cut/thole/long} pair style is equivalent to, but more convenient that
the frequent combination \sphinxstyleemphasis{hybrid/overlay lj/cut/coul/long cutoff thole damp
cutoff2}. It is not only a shorthand for this pair\_style combination, but
it also allows for mixing pair coefficients instead of listing them all.
The \sphinxstyleemphasis{lj/cut/thole/long} pair style is also a bit faster because it avoids an
overlay and can benefit from OMP acceleration. Moreover, it uses a more
precise approximation of the direct Coulomb interaction at short range similar
to {\hyperref[\detokenize{pair_cs::doc}]{\sphinxcrossref{\DUrole{doc}{coul/long/cs}}}}, which stabilizes the temperature of
Drude particles.

\sphinxAtStartPar
The \sphinxstyleemphasis{thole} pair styles compute the Coulomb interaction damped at
short distances by a function
\begin{equation*}
\begin{split}T_{ij}(r_{ij}) = 1 - \left( 1 +
\frac{s_{ij} r_{ij} }{2} \right)
\exp \left( - s_{ij} r_{ij} \right)\end{split}
\end{equation*}
\sphinxAtStartPar
This function results from an adaptation to point charges
{\hyperref[\detokenize{pair_thole:noskov1}]{\sphinxcrossref{\DUrole{std,std-ref}{(Noskov)}}}} of the dipole screening scheme originally proposed
by {\hyperref[\detokenize{pair_thole:thole1}]{\sphinxcrossref{\DUrole{std,std-ref}{Thole}}}}. The scaling coefficient \(s_{ij}\) is determined
by the polarizability of the atoms, \(\alpha_i\), and by a Thole
damping parameter \(a\).  This Thole damping parameter usually takes
a value of 2.6, but in certain force fields the value can depend upon
the atom types. The mixing rule for Thole damping parameters is the
arithmetic average, and for polarizabilities the geometric average
between the atom\sphinxhyphen{}specific values.
\begin{equation*}
\begin{split}s_{ij} = \frac{ a_{ij} }{
(\alpha_{ij})^{1/3} } = \frac{ (a_i + a_j)/2 }{
[(\alpha_i\alpha_j)^{1/2}]^{1/3} }\end{split}
\end{equation*}
\sphinxAtStartPar
The damping function is only applied to the interactions between the
point charges representing the induced dipoles on polarizable sites,
that is, charges on Drude particles, \(q_{D,i}\), and opposite
charges, \(-q_{D,i}\), located on the respective core particles
(to which each Drude particle is bonded). Therefore, Thole screening
is not applied to the full charge of the core particle \(q_i\), but
only to the \(-q_{D,i}\) part of it.

\sphinxAtStartPar
The interactions between core charges are subject to the weighting
factors set by the {\hyperref[\detokenize{special_bonds::doc}]{\sphinxcrossref{\DUrole{doc}{special\_bonds}}}} command. The
interactions between Drude particles and core charges or
non\sphinxhyphen{}polarizable atoms are also subject to these weighting factors. The
Drude particles inherit the 1\sphinxhyphen{}2, 1\sphinxhyphen{}3 and 1\sphinxhyphen{}4 neighbor relations from
their respective cores.

\sphinxAtStartPar
For pair\_style \sphinxstyleemphasis{thole}, the following coefficients must be defined for
each pair of atoms types via the {\hyperref[\detokenize{pair_coeff::doc}]{\sphinxcrossref{\DUrole{doc}{pair\_coeff}}}} command
as in the example above.
\begin{itemize}
\item {} 
\sphinxAtStartPar
\(\alpha\) (distance units\textasciicircum{}3)

\item {} 
\sphinxAtStartPar
damp

\item {} 
\sphinxAtStartPar
cutoff (distance units)

\end{itemize}

\sphinxAtStartPar
The last two coefficients are optional.  If not specified the global
Thole damping parameter or global cutoff specified in the pair\_style
command are used. In order to specify a cutoff (third argument) a damp
parameter (second argument) must also be specified.

\sphinxAtStartPar
For pair style \sphinxstyleemphasis{lj/cut/thole/long}, the following coefficients must be
defined for each pair of atoms types via the {\hyperref[\detokenize{pair_coeff::doc}]{\sphinxcrossref{\DUrole{doc}{pair\_coeff}}}}
command.
\begin{itemize}
\item {} 
\sphinxAtStartPar
\(\epsilon\) (energy units)

\item {} 
\sphinxAtStartPar
\(\sigma\) (length units)

\item {} 
\sphinxAtStartPar
\(\alpha\) (distance units\textasciicircum{}3)

\item {} 
\sphinxAtStartPar
damp

\item {} 
\sphinxAtStartPar
LJ cutoff (distance units)

\end{itemize}

\sphinxAtStartPar
The last two coefficients are optional and default to the global values from
the \sphinxstyleemphasis{pair\_style} command line.


\bigskip\hrule\bigskip


\sphinxAtStartPar
Styles with a \sphinxstyleemphasis{gpu}, \sphinxstyleemphasis{intel}, \sphinxstyleemphasis{kk}, \sphinxstyleemphasis{omp}, or \sphinxstyleemphasis{opt} suffix are
functionally the same as the corresponding style without the suffix.
They have been optimized to run faster, depending on your available
hardware, as discussed on the {\hyperref[\detokenize{Speed_packages::doc}]{\sphinxcrossref{\DUrole{doc}{Accelerator packages}}}}
page.  The accelerated styles take the same arguments and should
produce the same results, except for round\sphinxhyphen{}off and precision issues.

\sphinxAtStartPar
These accelerated styles are part of the GPU, INTEL, KOKKOS,
OPENMP, and OPT packages, respectively.  They are only enabled if
LAMMPS was built with those packages.  See the {\hyperref[\detokenize{Build_package::doc}]{\sphinxcrossref{\DUrole{doc}{Build package}}}} page for more info.

\sphinxAtStartPar
You can specify the accelerated styles explicitly in your input script
by including their suffix, or you can use the {\hyperref[\detokenize{Run_options::doc}]{\sphinxcrossref{\DUrole{doc}{\sphinxhyphen{}suffix command\sphinxhyphen{}line switch}}}} when you invoke LAMMPS, or you can use the
{\hyperref[\detokenize{suffix::doc}]{\sphinxcrossref{\DUrole{doc}{suffix}}}} command in your input script.

\sphinxAtStartPar
See the {\hyperref[\detokenize{Speed_packages::doc}]{\sphinxcrossref{\DUrole{doc}{Accelerator packages}}}} page for more
instructions on how to use the accelerated styles effectively.


\subsection{Mixing, shift, table, tail correction, restart, rRESPA info}
\label{\detokenize{pair_thole:mixing-shift-table-tail-correction-restart-rrespa-info}}
\sphinxAtStartPar
The \sphinxstyleemphasis{thole} pair style does not support mixing.  Thus, coefficients
for all I,J pairs must be specified explicitly.

\sphinxAtStartPar
The \sphinxstyleemphasis{lj/cut/thole/long} pair style does support mixing. Mixed coefficients
are defined using
\begin{equation*}
\begin{split}\alpha_{ij} = & \sqrt{\alpha_i\alpha_j} \\
& \\
a_{ij} = & \frac 1 2 (a_i + a_j)\end{split}
\end{equation*}

\subsection{Restrictions}
\label{\detokenize{pair_thole:restrictions}}
\sphinxAtStartPar
These pair styles are part of the DRUDE package. They are only
enabled if LAMMPS was built with that package. See the {\hyperref[\detokenize{Build_package::doc}]{\sphinxcrossref{\DUrole{doc}{Build package}}}} page for more info.

\sphinxAtStartPar
This pair\_style should currently not be used with the {\hyperref[\detokenize{dihedral_charmm::doc}]{\sphinxcrossref{\DUrole{doc}{charmm dihedral style}}}} if the latter has non\sphinxhyphen{}zero 1\sphinxhyphen{}4 weighting
factors. This is because the \sphinxstyleemphasis{thole} pair style does not know which
pairs are 1\sphinxhyphen{}4 partners of which dihedrals.

\sphinxAtStartPar
The \sphinxstyleemphasis{lj/cut/thole/long} pair style should be used with a {\hyperref[\detokenize{kspace_style::doc}]{\sphinxcrossref{\DUrole{doc}{Kspace solver}}}}
like PPPM or Ewald, which is only enabled if LAMMPS was built with the kspace
package.


\subsection{Related commands}
\label{\detokenize{pair_thole:related-commands}}
\sphinxAtStartPar
{\hyperref[\detokenize{fix_drude::doc}]{\sphinxcrossref{\DUrole{doc}{fix drude}}}}, {\hyperref[\detokenize{fix_langevin_drude::doc}]{\sphinxcrossref{\DUrole{doc}{fix langevin/drude}}}}, {\hyperref[\detokenize{fix_drude_transform::doc}]{\sphinxcrossref{\DUrole{doc}{fix drude/transform}}}}, {\hyperref[\detokenize{compute_temp_drude::doc}]{\sphinxcrossref{\DUrole{doc}{compute temp/drude}}}}
{\hyperref[\detokenize{pair_lj_cut_coul::doc}]{\sphinxcrossref{\DUrole{doc}{pair\_style lj/cut/coul/long}}}}


\subsection{Default}
\label{\detokenize{pair_thole:default}}
\sphinxAtStartPar
none


\bigskip\hrule\bigskip

\phantomsection\label{\detokenize{pair_thole:noskov1}}
\sphinxAtStartPar
\sphinxstylestrong{(Noskov)} Noskov, Lamoureux and Roux, J Phys Chem B, 109, 6705 (2005).

\phantomsection\label{\detokenize{pair_thole:thole1}}
\sphinxAtStartPar
\sphinxstylestrong{(Thole)} Chem Phys, 59, 341 (1981).

\sphinxstepscope

\index{pair\_style threebody/table@\spxentry{pair\_style threebody/table}}\ignorespaces 

\section{pair\_style threebody/table command}
\label{\detokenize{pair_threebody_table:pair-style-threebody-table-command}}\label{\detokenize{pair_threebody_table:index-0}}\label{\detokenize{pair_threebody_table::doc}}

\subsection{Syntax}
\label{\detokenize{pair_threebody_table:syntax}}
\begin{sphinxVerbatim}[commandchars=\\\{\}]
\PYG{k}{pair\PYGZus{}style}\PYG{+w}{ }\PYG{n}{style}
\end{sphinxVerbatim}
\begin{itemize}
\item {} 
\sphinxAtStartPar
style = \sphinxstyleemphasis{threebody/table}

\end{itemize}


\subsection{Examples}
\label{\detokenize{pair_threebody_table:examples}}
\begin{sphinxVerbatim}[commandchars=\\\{\}]
\PYG{k}{pair\PYGZus{}style}\PYG{+w}{ }\PYG{n}{threebody}\PYG{o}{/}\PYG{n}{table}
\PYG{k}{pair\PYGZus{}coeff}\PYG{+w}{ }\PYG{o}{*}\PYG{+w}{ }\PYG{o}{*}\PYG{+w}{ }\PYG{n}{spce2.3b}\PYG{+w}{ }\PYG{n}{type1}\PYG{+w}{ }\PYG{n}{type2}

\PYG{k}{pair\PYGZus{}style}\PYG{+w}{ }\PYG{n}{hybrid}\PYG{o}{/}\PYG{n}{overlay}\PYG{+w}{ }\PYG{n}{table}\PYG{+w}{ }\PYG{n}{linear}\PYG{+w}{ }\PYG{l+m}{1200}\PYG{+w}{ }\PYG{n}{threebody}\PYG{o}{/}\PYG{n}{table}
\PYG{k}{pair\PYGZus{}coeff}\PYG{+w}{      }\PYG{l+m}{1}\PYG{+w}{ }\PYG{l+m}{1}\PYG{+w}{ }\PYG{n}{table}\PYG{+w}{ }\PYG{n}{table\PYGZus{}CG\PYGZus{}CG.txt}\PYG{+w}{ }\PYG{n}{VOTCA}
\PYG{k}{pair\PYGZus{}coeff}\PYG{+w}{      }\PYG{o}{*}\PYG{+w}{ }\PYG{o}{*}\PYG{+w}{ }\PYG{n}{threebody}\PYG{o}{/}\PYG{n}{table}\PYG{+w}{ }\PYG{n}{spce.3b}\PYG{+w}{ }\PYG{n}{type}
\end{sphinxVerbatim}

\sphinxAtStartPar
Used in example input scripts:

\begin{sphinxVerbatim}[commandchars=\\\{\}]
examples/PACKAGES/manybody\PYGZus{}table/in.spce
examples/PACKAGES/manybody\PYGZus{}table/in.spce2
\end{sphinxVerbatim}


\subsection{Description}
\label{\detokenize{pair_threebody_table:description}}
\sphinxAtStartPar
\DUrole{versionmodified,added}{Added in version 2Jun2022.}

\sphinxAtStartPar
The \sphinxstyleemphasis{threebody/table} style is a pair style for generic tabulated
three\sphinxhyphen{}body interactions.  It has been developed for (coarse\sphinxhyphen{}grained)
simulations (of water) with Kernel\sphinxhyphen{}based machine learning (ML)
potentials ({\hyperref[\detokenize{pair_threebody_table:scherer2}]{\sphinxcrossref{\DUrole{std,std-ref}{Scherer2}}}}).  As for many other MANYBODY
package pair styles the energy of a system is computed as a sum over
three\sphinxhyphen{}body terms:
\begin{equation*}
\begin{split}E =  \sum_i \sum_{j \neq i} \sum_{k > j} \phi_3 (r_{ij}, r_{ik}, \theta_{ijk})\end{split}
\end{equation*}
\sphinxAtStartPar
The summations in the formula are over all neighbors J and K of atom I
within a cutoff distance \(cut\).  In contrast to the
Stillinger\sphinxhyphen{}Weber potential, all forces are not calculated analytically,
but read in from a three\sphinxhyphen{}body force/energy table which can be generated
with the csg\_ml app of VOTCA as available at:
\sphinxurl{https://gitlab.mpcdf.mpg.de/votca/votca}.

\sphinxAtStartPar
Only a single pair\_coeff command is used with the \sphinxstyleemphasis{threebody/table}
style which specifies a threebody potential (“.3b”) file with parameters
for all needed elements. These are then mapped to LAMMPS atom types by
specifying N\_el additional arguments after the “.3b” filename in the
pair\_coeff command, where N\_el is the number of LAMMPS atom types:
\begin{itemize}
\item {} 
\sphinxAtStartPar
“.3b” filename

\item {} 
\sphinxAtStartPar
N\_el element names = mapping of threebody elements to atom types

\end{itemize}

\sphinxAtStartPar
See the {\hyperref[\detokenize{pair_coeff::doc}]{\sphinxcrossref{\DUrole{doc}{pair\_coeff}}}} page for alternate ways to
specify the path for the potential file.

\sphinxAtStartPar
As an example, imagine a file SiC.3b has three\sphinxhyphen{}body values for Si and C.
If your LAMMPS simulation has 4 atoms types and you want the first 3 to
be Si, and the fourth to be C, you would use the following pair\_coeff
command:

\begin{sphinxVerbatim}[commandchars=\\\{\}]
\PYG{k}{pair\PYGZus{}coeff}\PYG{+w}{ }\PYG{o}{*}\PYG{+w}{ }\PYG{o}{*}\PYG{+w}{ }\PYG{n}{SiC.3b}\PYG{+w}{ }\PYG{n}{Si}\PYG{+w}{ }\PYG{n}{Si}\PYG{+w}{ }\PYG{n}{Si}\PYG{+w}{ }\PYG{n}{C}
\end{sphinxVerbatim}

\sphinxAtStartPar
The first 2 arguments must be * * so as to span all LAMMPS atom types.
The first three Si arguments map LAMMPS atom types 1,2,3 to the Si
element in the “.3b” file.  The final C argument maps LAMMPS atom type 4
to the C element in the threebody file. If a mapping value is specified
as NULL, the mapping is not performed. This can be used when a
\sphinxstyleemphasis{threebody/table} potential is used as part of the \sphinxstyleemphasis{hybrid} pair style.
The NULL values are placeholders for atom types that will be used with
other potentials.

\sphinxAtStartPar
The three\sphinxhyphen{}body files have a “.3b” suffix. Lines that are not blank or
comments (starting with \#) define parameters for a triplet of
elements. The parameters in a single entry specify to the (three\sphinxhyphen{}body)
cutoff distance and the tabulated three\sphinxhyphen{}body interaction. A single entry
then contains:
\begin{itemize}
\item {} 
\sphinxAtStartPar
element 1 (the center atom in a 3\sphinxhyphen{}body interaction)

\item {} 
\sphinxAtStartPar
element 2

\item {} 
\sphinxAtStartPar
element 3

\item {} 
\sphinxAtStartPar
cut (distance units)

\item {} 
\sphinxAtStartPar
filename

\item {} 
\sphinxAtStartPar
keyword

\item {} 
\sphinxAtStartPar
style

\item {} 
\sphinxAtStartPar
N

\end{itemize}

\sphinxAtStartPar
The parameter \(cut\) is the (three\sphinxhyphen{}body) cutoff distance.  When set
to 0, no interaction is calculated for this element triplet.  The
parameters \sphinxstyleemphasis{filename}, \sphinxstyleemphasis{keyword}, \sphinxstyleemphasis{style}, and \sphinxstyleemphasis{N} refer to the
tabulated three\sphinxhyphen{}body potential.

\sphinxAtStartPar
The tabulation is done on a three\sphinxhyphen{}dimensional grid of the two distances
\(r_{ij}\) and \(r_{ik}\) as well as the angle
\(\theta_{ijk}\) which is constructed in the following way. There
are two different cases.  If element 2 and element 3 are of the same
type (e.g. SiCC), the distance \(r_{ij}\) is varied in “N” steps
from rmin to rmax and the distance \(r_{ik}\) is varied from
\(r_{ij}\) to rmax. This can be done, due to the symmetry of the
triplet. If element 2 and element 3 are not of the same type
(e.g. SiCSi), there is no additional symmetry and the distance
\(r_{ik}\) is also varied from rmin to rmax in “N” steps.  The angle
\(\theta_{ijk}\) is always varied in “2N” steps from (0.0 +
180.0/(4N)) to (180.0 \sphinxhyphen{} 180.0/(4N)). Therefore, the total number of
table entries is “M = N * N * (N+1)” for the symmetric (element 2 and
element 3 are of the same type) and “M = 2 * N * N * N” for the general
case (element 2 and element 3 are not of the same type).

\sphinxAtStartPar
The forces on all three particles I, J, and K of a triplet of this type
of three\sphinxhyphen{}body interaction potential (\(\phi_3 (r_{ij}, r_{ik},
\theta_{ijk})\)) lie within the plane defined by the three inter\sphinxhyphen{}particle
distance vectors \({\mathbf r}_{ij}\), \({\mathbf r}_{ik}\), and
\({\mathbf r}_{jk}\).  This property is used to project the forces
onto the inter\sphinxhyphen{}particle distance vectors as follows
\begin{equation*}
\begin{split}\begin{pmatrix}
   {\mathbf f}_{i} \\
   {\mathbf f}_{j} \\
   {\mathbf f}_{k} \\
\end{pmatrix} =
\begin{pmatrix}
   f_{i1} & f_{i2} & 0 \\
   f_{j1} & 0 & f_{j2} \\
   0 & f_{k1} & f_{k2} \\
\end{pmatrix}
\begin{pmatrix}
   {\mathbf r}_{ij} \\
   {\mathbf r}_{ik} \\
   {\mathbf r}_{jk} \\
\end{pmatrix}\end{split}
\end{equation*}
\sphinxAtStartPar
and then tabulate the 6 force constants \(f_{i1}\), \(f_{i2}\),
\(f_{j1}\), \(f_{j2}\), \(f_{k1}\), and \(f_{k2}\), as
well as the energy of a triplet e.  Due to symmetry reasons, the
following relations hold: \(f_{i1}=-f_{j1}\),
\(f_{i2}=-f_{k1}\), and \(f_{j2}=-f_{k2}\). As in this pair
style the forces are read in directly, a correct MD simulation is also
performed in the case that the triplet energies are set to e=0.

\sphinxAtStartPar
The \sphinxstyleemphasis{filename} specifies the file containing the tabulated energy and
derivative values of \(\phi_3 (r_{ij}, r_{ik}, \theta_{ijk})\).  The
\sphinxstyleemphasis{keyword} then specifies a section of the file. The format of this file
is as follows (without the parenthesized comments):

\begin{sphinxVerbatim}[commandchars=\\\{\}]
\PYGZsh{} Tabulated three\PYGZhy{}body potential for spce water (one or more comment or blank lines)

ENTRY1                                                                      (keyword is the first text on line)
N 12 rmin 2.55 rmax 3.65                                                    (N, rmin, rmax parameters)
                                                                            (blank line)
1 2.55 2.55 3.75 \PYGZhy{}867.212 \PYGZhy{}611.273 867.212 21386.8 611.273 \PYGZhy{}21386.8  0.0    (index, r\PYGZus{}ij, r\PYGZus{}ik, theta, f\PYGZus{}i1, f\PYGZus{}i2, f\PYGZus{}j1, f\PYGZus{}j2, f\PYGZus{}k1, f\PYGZus{}k2, e)
2 2.55 2.55 11.25 \PYGZhy{}621.539 \PYGZhy{}411.189 621.539 5035.95 411.189 \PYGZhy{}5035.95  0.0
...
1872 3.65 3.65 176.25 \PYGZhy{}0.00215132 \PYGZhy{}0.00412886 0.00215137 0.00111754 0.00412895 \PYGZhy{}0.00111757  0.0
\end{sphinxVerbatim}

\sphinxAtStartPar
A section begins with a non\sphinxhyphen{}blank line whose first character is not a
“\#”; blank lines or lines starting with “\#” can be used as comments
between sections.  The first line begins with a keyword which identifies
the section. The next line lists (in any order) one or more parameters
for the table.  Each parameter is a keyword followed by one or more
numeric values.

\sphinxAtStartPar
The parameter “N” is required. It should be the same than the parameter
“N” of the “.3b” file, otherwise its value is overwritten. “N”
determines the number of table entries “M” that follow: “M = N * N *
(N+1)” (symmetric triplet, e.g. SiCC) or “M = 2 * N * N * N” (asymmetric
triplet, e.g. SiCSi). Therefore “M = 12 * 12 * 13 = 1872” in the above
symmetric example. The parameters “rmin” and “rmax” are also required
and determine the minimum and maximum of the inter\sphinxhyphen{}particle distances
\(r_{ij}\) and \(r_{ik}\).

\sphinxAtStartPar
Following a blank line, the next M lines of the angular table file list
the tabulated values.  On each line, the first value is the index from 1
to M, the second value is the distance \(r_{ij}\), the third value
is the distance \(r_{ik}\), the fourth value is the angle
\(\theta_{ijk})\), the next six values are the force constants
\(f_{i1}\), \(f_{i2}\), \(f_{j1}\), \(f_{j2}\),
\(f_{k1}\), and \(f_{k2}\), and the last value is the energy e.

\sphinxAtStartPar
Note that one three\sphinxhyphen{}body potential file can contain many sections, each
with a tabulated potential. LAMMPS reads the file section by section
until it finds one that matches the specified \sphinxstyleemphasis{keyword} of appropriate
section of the “.3b” file.

\sphinxAtStartPar
At the moment, only the \sphinxstyleemphasis{style} \sphinxstyleemphasis{linear} is allowed and
implemented. After reading in the force table, it is internally stored
in LAMMPS as a lookup table. For each triplet configuration occurring in
the simulation within the cutoff distance, the next nearest tabulated
triplet configuration is looked up. No interpolation is done.  This
allows for a very efficient force calculation with the stored force
constants and energies. Due to the know table structure, the lookup can
be done efficiently. It has been tested ({\hyperref[\detokenize{pair_threebody_table:scherer2}]{\sphinxcrossref{\DUrole{std,std-ref}{Scherer2}}}})
that with a reasonably small bin size, the accuracy and speed is
comparable to that of a Stillinger\sphinxhyphen{}Weber potential with tabulated
three\sphinxhyphen{}body interactions ({\hyperref[\detokenize{pair_sw_angle_table::doc}]{\sphinxcrossref{\DUrole{doc}{pair\_style sw/angle/table}}}}) while the table format of this pair style allows
for more flexible three\sphinxhyphen{}body interactions.

\sphinxAtStartPar
As for the Stillinger\sphinxhyphen{}Weber potential, the three\sphinxhyphen{}body potential file
must contain entries for all the elements listed in the pair\_coeff
command.  It can also contain entries for additional elements not being
used in a particular simulation; LAMMPS ignores those entries.

\sphinxAtStartPar
For a single\sphinxhyphen{}element simulation, only a single entry is required
(e.g. SiSiSi).  For a two\sphinxhyphen{}element simulation, the file must contain 8
entries (for SiSiSi, SiSiC, SiCSi, SiCC, CSiSi, CSiC, CCSi, CCC), that
specify threebody parameters for all permutations of the two elements
interacting in three\sphinxhyphen{}body configurations.  Thus for 3 elements, 27
entries would be required, etc.

\sphinxAtStartPar
As annotated above, the first element in the entry is the center atom in
a three\sphinxhyphen{}body interaction. Thus an entry for SiCC means a Si atom with 2
C atoms as neighbors. The tabulated three\sphinxhyphen{}body forces can in principle
be specific to the three elements of the configuration.  However, the
user must ensure that it makes physically sense.  E.g., the tabulated
three\sphinxhyphen{}body forces for the entries CSiC and CCSi should be the same
exchanging \(r_{ij}\) with r\_\{ik\}, \(f_{j1}\) with
\(f_{k1}\), and \(f_{j2}\) with \(f_{k2}\).

\sphinxAtStartPar
Additional input files and reference data can be found at:
\sphinxurl{https://gitlab.mpcdf.mpg.de/votca/votca/-/tree/master/csg-tutorials/ml}


\bigskip\hrule\bigskip



\subsection{Mixing, shift, table, tail correction, restart, rRESPA info}
\label{\detokenize{pair_threebody_table:mixing-shift-table-tail-correction-restart-rrespa-info}}
\sphinxAtStartPar
As all interactions are tabulated, no mixing is performed.

\sphinxAtStartPar
This pair style does not support the {\hyperref[\detokenize{pair_modify::doc}]{\sphinxcrossref{\DUrole{doc}{pair\_modify}}}}
shift, table, and tail options.

\sphinxAtStartPar
This pair style does not write its information to {\hyperref[\detokenize{restart::doc}]{\sphinxcrossref{\DUrole{doc}{binary restart
files}}}}, since it is stored in potential files.  Thus, you need
to re\sphinxhyphen{}specify the pair\_style and pair\_coeff commands in an input script
that reads a restart file.

\sphinxAtStartPar
This pair style can only be used via the \sphinxstyleemphasis{pair} keyword of the
{\hyperref[\detokenize{run_style::doc}]{\sphinxcrossref{\DUrole{doc}{run\_style respa}}}} command.  It does not support the
\sphinxstyleemphasis{inner}, \sphinxstyleemphasis{middle}, \sphinxstyleemphasis{outer} keywords.


\bigskip\hrule\bigskip



\subsection{Restrictions}
\label{\detokenize{pair_threebody_table:restrictions}}
\sphinxAtStartPar
This pair style is part of the MANYBODY package.  It is only enabled if
LAMMPS was built with that package. See the {\hyperref[\detokenize{Build_package::doc}]{\sphinxcrossref{\DUrole{doc}{Build package}}}} page for more info.

\sphinxAtStartPar
This pair style requires the {\hyperref[\detokenize{newton::doc}]{\sphinxcrossref{\DUrole{doc}{newton}}}} setting to be “on”
for pair interactions.


\subsection{Related commands}
\label{\detokenize{pair_threebody_table:related-commands}}
\sphinxAtStartPar
{\hyperref[\detokenize{pair_coeff::doc}]{\sphinxcrossref{\DUrole{doc}{pair\_coeff}}}}, {\hyperref[\detokenize{pair_sw_angle_table::doc}]{\sphinxcrossref{\DUrole{doc}{pair sw/angle/table}}}}


\bigskip\hrule\bigskip

\phantomsection\label{\detokenize{pair_threebody_table:scherer2}}
\sphinxAtStartPar
\sphinxstylestrong{(Scherer2)} C. Scherer, R. Scheid, D. Andrienko, and T. Bereau, J. Chem. Theor. Comp. 16, 3194\sphinxhyphen{}3204 (2020).

\sphinxstepscope

\index{pair\_style tracker@\spxentry{pair\_style tracker}}\ignorespaces 

\section{pair\_style tracker command}
\label{\detokenize{pair_tracker:pair-style-tracker-command}}\label{\detokenize{pair_tracker:index-0}}\label{\detokenize{pair_tracker::doc}}

\subsection{Syntax}
\label{\detokenize{pair_tracker:syntax}}
\begin{sphinxVerbatim}[commandchars=\\\{\}]
\PYG{k}{pair\PYGZus{}style}\PYG{+w}{ }\PYG{n}{tracker}\PYG{+w}{ }\PYG{n}{fix\PYGZus{}ID}\PYG{+w}{ }\PYG{n}{N}\PYG{+w}{ }\PYG{n}{keyword}\PYG{+w}{ }\PYG{n}{values}\PYG{+w}{ }\PYG{n}{attribute1}\PYG{+w}{ }\PYG{n}{attribute2}\PYG{+w}{ }\PYG{n}{...}
\end{sphinxVerbatim}
\begin{itemize}
\item {} 
\sphinxAtStartPar
fix\_ID = ID of associated internal fix to store data

\item {} 
\sphinxAtStartPar
N = prepare data for output every this many timesteps

\item {} 
\sphinxAtStartPar
zero or more keywords may be appended

\item {} 
\sphinxAtStartPar
keyword = \sphinxstyleemphasis{finite} or \sphinxstyleemphasis{time/min} or \sphinxstyleemphasis{type/include}
\begin{sphinxalltt}
\sphinxstyleemphasis{finite} value = none
   pair style uses atomic diameters to identify contacts
\sphinxstyleemphasis{time/min} value = T
   T = minimum number of timesteps of interaction
\sphinxstyleemphasis{type/include} value = list1 list2
   list1,list2 = separate lists of types (see below)
\end{sphinxalltt}

\item {} 
\sphinxAtStartPar
one or more attributes may be appended

\begin{sphinxVerbatim}[commandchars=\\\{\}]
possible attributes = id1 id2 time/created time/broken time/total
                      r/min r/ave x y z
\end{sphinxVerbatim}

\begin{sphinxVerbatim}[commandchars=\\\{\}]
id1, id2 = IDs of the two atoms in each pair interaction
time/created = the timestep that the two atoms began interacting
time/broken = the timestep that the two atoms stopped interacting
time/total = the total number of timesteps the two atoms interacted
r/min = the minimum radial distance between the two atoms during the interaction (distance units)
r/ave = the average radial distance between the two atoms during the interaction (distance units)
x, y, z = the center of mass position of the two atoms when they stopped interacting (distance units)
\end{sphinxVerbatim}

\end{itemize}


\subsection{Examples}
\label{\detokenize{pair_tracker:examples}}
\begin{sphinxVerbatim}[commandchars=\\\{\}]
\PYG{k}{pair\PYGZus{}style}\PYG{+w}{ }\PYG{n}{hybrid}\PYG{o}{/}\PYG{n}{overlay}\PYG{+w}{ }\PYG{n}{tracker}\PYG{+w}{ }\PYG{n}{myfix}\PYG{+w}{ }\PYG{l+m}{1000}\PYG{+w}{ }\PYG{n}{id1}\PYG{+w}{ }\PYG{n}{id2}\PYG{+w}{ }\PYG{n}{type}\PYG{o}{/}\PYG{n}{include}\PYG{+w}{ }\PYG{l+m}{1}\PYG{+w}{ }\PYG{o}{*}\PYG{+w}{ }\PYG{n}{type}\PYG{o}{/}\PYG{n}{include}\PYG{+w}{ }\PYG{l+m}{2}\PYG{+w}{ }\PYG{l+m}{3},4\PYG{+w}{  }\PYG{n}{lj}\PYG{o}{/}\PYG{n}{cut}\PYG{+w}{ }\PYG{l+m}{2.5}
\PYG{k}{pair\PYGZus{}coeff}\PYG{+w}{ }\PYG{l+m}{1}\PYG{+w}{ }\PYG{l+m}{1}\PYG{+w}{ }\PYG{n}{tracker}\PYG{+w}{ }\PYG{l+m}{2.0}

\PYG{k}{pair\PYGZus{}style}\PYG{+w}{ }\PYG{n}{hybrid}\PYG{o}{/}\PYG{n}{overlay}\PYG{+w}{ }\PYG{n}{tracker}\PYG{+w}{ }\PYG{n}{myfix}\PYG{+w}{ }\PYG{l+m}{1000}\PYG{+w}{ }\PYG{n}{finite}\PYG{+w}{ }\PYG{n}{x}\PYG{+w}{ }\PYG{n}{y}\PYG{+w}{ }\PYG{n}{z}\PYG{+w}{ }\PYG{n}{time}\PYG{o}{/}\PYG{n}{min}\PYG{+w}{ }\PYG{l+m}{100}\PYG{+w}{ }\PYG{n}{granular}
\PYG{k}{pair\PYGZus{}coeff}\PYG{+w}{ }\PYG{o}{*}\PYG{+w}{ }\PYG{o}{*}\PYG{+w}{ }\PYG{n}{tracker}

\PYG{k}{dump }\PYG{n+nv+nvIdentifier}{1}\PYG{+w}{ }\PYG{n+nv+nvIdentifier}{all}\PYG{+w}{ }\PYG{n}{local}\PYG{+w}{ }\PYG{l+m}{1000}\PYG{+w}{ }\PYG{n}{dump.local}\PYG{+w}{ }\PYG{n}{f\PYGZus{}myfix[1]}\PYG{+w}{ }\PYG{n}{f\PYGZus{}myfix[2]}\PYG{+w}{ }\PYG{n}{f\PYGZus{}myfix[3]}
\PYG{k}{dump\PYGZus{}modify }\PYG{n+nv+nvIdentifier}{1}\PYG{+w}{ }\PYG{n}{write\PYGZus{}header}\PYG{+w}{ }\PYG{n}{no}
\end{sphinxVerbatim}


\subsection{Description}
\label{\detokenize{pair_tracker:description}}
\sphinxAtStartPar
Style \sphinxstyleemphasis{tracker} monitors information about pairwise interactions.  It
does not calculate any forces on atoms. {\hyperref[\detokenize{pair_hybrid::doc}]{\sphinxcrossref{\DUrole{doc}{Pair hybrid/overlay}}}} can be used to combine this pair style with any other
pair style, as shown in the examples above.

\sphinxAtStartPar
At each timestep, if two neighboring atoms move beyond the interaction
cutoff, pairwise data is processed and transferred to an internal fix
labeled \sphinxstyleemphasis{fix\_ID}. This allows the local data to be accessed by other
LAMMPS commands. Additional
filters can be applied using the \sphinxstyleemphasis{time/min} or \sphinxstyleemphasis{type/include} keywords
described below.  Note that this is the interaction cutoff defined by
this pair style, not the short\sphinxhyphen{}range cutoff defined by the pair style
that is calculating forces on atoms.

\sphinxAtStartPar
Following any optional keyword/value arguments, a list of one or more
attributes is specified.  These include the IDs of the two atoms in
the pair.  The other attributes for the pair of atoms are the
duration of time they were “interacting” or at the point in time they
started or stopped interacting.  In this context, “interacting” means
the time window during which the two atoms were closer than the
interaction cutoff distance.  The attributes for time/* refer to
timesteps.

\sphinxAtStartPar
Data is continuously accumulated by the internal fix over intervals of \sphinxstyleemphasis{N}
timesteps. At the end of each interval, all of the saved accumulated
data is deleted to make room for new data. Individual datum may
therefore persist anywhere between \sphinxstyleemphasis{1} to \sphinxstyleemphasis{N} timesteps depending on
when they are saved. This data can be accessed using the \sphinxstyleemphasis{fix\_ID} and a
{\hyperref[\detokenize{dump::doc}]{\sphinxcrossref{\DUrole{doc}{dump local}}}} command. To ensure all data is output,
the dump frequency should correspond to the same interval of \sphinxstyleemphasis{N}
timesteps. A dump frequency of an integer multiple of \sphinxstyleemphasis{N} can be used
to regularly output a sample of the accumulated data.


\bigskip\hrule\bigskip


\sphinxAtStartPar
The following optional keywords may be used.

\sphinxAtStartPar
If the \sphinxstyleemphasis{finite} keyword is not used, the following coefficients must
be defined for each pair of atom types via the {\hyperref[\detokenize{pair_coeff::doc}]{\sphinxcrossref{\DUrole{doc}{pair\_coeff}}}} command as in the examples above, or in the data file or
restart files read by the {\hyperref[\detokenize{read_data::doc}]{\sphinxcrossref{\DUrole{doc}{read\_data}}}} or
{\hyperref[\detokenize{read_restart::doc}]{\sphinxcrossref{\DUrole{doc}{read\_restart}}}} commands, or by mixing as described
below:
\begin{itemize}
\item {} 
\sphinxAtStartPar
cutoff (distance units)

\end{itemize}

\sphinxAtStartPar
If the \sphinxstyleemphasis{finite} keyword is used, there are no additional coefficients
to set for each pair of atom types via the
{\hyperref[\detokenize{pair_coeff::doc}]{\sphinxcrossref{\DUrole{doc}{pair\_coeff}}}} command. Interaction cutoffs are
instead calculated based on the diameter of finite particles. However
you must still use the {\hyperref[\detokenize{pair_coeff::doc}]{\sphinxcrossref{\DUrole{doc}{pair\_coeff}}}} for all atom
types. For example the command

\begin{sphinxVerbatim}[commandchars=\\\{\}]
\PYG{k}{pair\PYGZus{}coeff}\PYG{+w}{ }\PYG{o}{*}\PYG{+w}{ }\PYG{o}{*}
\end{sphinxVerbatim}

\sphinxAtStartPar
should be used.

\sphinxAtStartPar
The \sphinxstyleemphasis{time/min} keyword sets a minimum amount of time that an
interaction must persist to be included.  This setting can be used to
censor short\sphinxhyphen{}lived interactions.

\sphinxAtStartPar
The \sphinxstyleemphasis{type/include} keyword filters interactions based on the types of
the two atoms.  Data is only saved for interactions between atoms
whose two atom types appear in \sphinxstyleemphasis{list1} and \sphinxstyleemphasis{list2}.  Atom type 1 must
be in list1 and atom type 2 in list2.  Or vice versa.

\sphinxAtStartPar
Each type list consists of a series of type ranges separated by
commas.  Each range can be specified as a single numeric value, or a
wildcard asterisk can be used to specify a range of values.  This
takes the form “*” or “*n” or “n*” or “m*n”.  For example, if M =
the number of atom types, then an asterisk with no numeric values
means all types from 1 to M.  A leading asterisk means all types from
1 to n (inclusive).  A trailing asterisk means all types from n to M
(inclusive).  A middle asterisk means all types from m to n
(inclusive).  Note that the \sphinxstyleemphasis{type/include} keyword can be specified
multiple times.


\subsection{Mixing, shift, table, tail correction, restart, rRESPA info}
\label{\detokenize{pair_tracker:mixing-shift-table-tail-correction-restart-rrespa-info}}
\sphinxAtStartPar
For atom type pairs I,J and I != J, the cutoff coefficient and cutoff
distance for this pair style can be mixed.  The cutoff is always mixed
via a \sphinxstyleemphasis{geometric} rule.  The cutoff is mixed according to the
pair\_modify mix value.  The default mix value is \sphinxstyleemphasis{geometric}.  See
the “pair\_modify” command for details.

\sphinxAtStartPar
This pair style writes its information to {\hyperref[\detokenize{restart::doc}]{\sphinxcrossref{\DUrole{doc}{binary restart files}}}}, so pair\_style and pair\_coeff commands do not need to be
specified in an input script that reads a restart file.

\sphinxAtStartPar
The {\hyperref[\detokenize{pair_modify::doc}]{\sphinxcrossref{\DUrole{doc}{pair\_modify}}}} shift, table, and tail options
are not relevant for this pair style.

\sphinxAtStartPar
The accumulated data is not written to restart files and should be
output before a restart file is written to avoid missing data.

\sphinxAtStartPar
The internal fix calculates a local vector or local array depending on the
number of input values.  The length of the vector or number of rows in
the array is the number of recorded, lost interactions.  If a single
input is specified, a local vector is produced.  If two or more inputs
are specified, a local array is produced where the number of columns =
the number of inputs.  The vector or array can be accessed by any
command that uses local values from a compute as input.  See the
{\hyperref[\detokenize{Howto_output::doc}]{\sphinxcrossref{\DUrole{doc}{Howto output}}}} page for an overview of LAMMPS
output options.

\sphinxAtStartPar
The vector or array will be floating point values that correspond to
the specified attribute.


\bigskip\hrule\bigskip



\subsection{Restrictions}
\label{\detokenize{pair_tracker:restrictions}}
\sphinxAtStartPar
This pair style is part of the MISC package.  It is only enabled if
LAMMPS was built with that package.  See the {\hyperref[\detokenize{Build_package::doc}]{\sphinxcrossref{\DUrole{doc}{Build package}}}} page for more info.

\sphinxAtStartPar
This pair style is currently incompatible with granular pair styles
that extend beyond the contact (e.g. JKR and DMT).


\subsection{Related commands}
\label{\detokenize{pair_tracker:related-commands}}

\subsection{Default}
\label{\detokenize{pair_tracker:default}}
\sphinxAtStartPar
none

\sphinxstepscope

\index{pair\_style tri/lj@\spxentry{pair\_style tri/lj}}\ignorespaces 

\section{pair\_style tri/lj command}
\label{\detokenize{pair_tri_lj:pair-style-tri-lj-command}}\label{\detokenize{pair_tri_lj:index-0}}\label{\detokenize{pair_tri_lj::doc}}

\subsection{Syntax}
\label{\detokenize{pair_tri_lj:syntax}}
\begin{sphinxVerbatim}[commandchars=\\\{\}]
\PYG{k}{pair\PYGZus{}style}\PYG{+w}{ }\PYG{n}{tri}\PYG{o}{/}\PYG{n}{lj}\PYG{+w}{ }\PYG{n}{cutoff}
\end{sphinxVerbatim}

\sphinxAtStartPar
cutoff = global cutoff for interactions (distance units)


\subsection{Examples}
\label{\detokenize{pair_tri_lj:examples}}
\begin{sphinxVerbatim}[commandchars=\\\{\}]
\PYG{k}{pair\PYGZus{}style}\PYG{+w}{ }\PYG{n}{tri}\PYG{o}{/}\PYG{n}{lj}\PYG{+w}{ }\PYG{l+m}{3.0}
\PYG{k}{pair\PYGZus{}coeff}\PYG{+w}{ }\PYG{o}{*}\PYG{+w}{ }\PYG{o}{*}\PYG{+w}{ }\PYG{l+m}{1.0}\PYG{+w}{ }\PYG{l+m}{1.0}
\PYG{k}{pair\PYGZus{}coeff}\PYG{+w}{ }\PYG{l+m}{1}\PYG{+w}{ }\PYG{l+m}{1}\PYG{+w}{ }\PYG{l+m}{1.0}\PYG{+w}{ }\PYG{l+m}{1.5}\PYG{+w}{ }\PYG{l+m}{2.5}
\end{sphinxVerbatim}


\subsection{Description}
\label{\detokenize{pair_tri_lj:description}}
\sphinxAtStartPar
Style \sphinxstyleemphasis{tri/lj} treats particles which are triangles as a set of small
spherical particles that tile the triangle surface as explained below.
Interactions between two triangles, each with N1 and N2 spherical
particles, are calculated as the pairwise sum of N1*N2 Lennard\sphinxhyphen{}Jones
interactions.  Interactions between a triangle with N spherical
particles and a point particle are treated as the pairwise sum of N
Lennard\sphinxhyphen{}Jones interactions.  See the {\hyperref[\detokenize{pair_lj::doc}]{\sphinxcrossref{\DUrole{doc}{pair\_style lj/cut}}}}
doc page for the definition of Lennard\sphinxhyphen{}Jones interactions.

\sphinxAtStartPar
The cutoff distance for an interaction between 2 triangles, or between
a triangle and a point particle, is calculated from the position of
the triangle (its centroid), not between pairs of individual spheres
comprising the triangle.  Thus an interaction is either calculated in
its entirety or not at all.

\sphinxAtStartPar
The set of non\sphinxhyphen{}overlapping spherical particles that represent a
triangle, for purposes of this pair style, are generated in the
following manner.  Assume the triangle is of type I, and sigma\_II has
been specified.  We want a set of spheres with centers in the plane of
the triangle, none of them larger in diameter than sigma\_II, which
completely cover the triangle’s area, but with minimal overlap and a
minimal total number of spheres.  This is done in a recursive manner.
Place a sphere at the centroid of the original triangle.  Calculate
what diameter it must have to just cover all 3 corner points of the
triangle.  If that diameter is equal to or smaller than sigma\_II, then
include a sphere of the calculated diameter in the set of covering
spheres.  It the diameter is larger than sigma\_II, then split the
triangle into 2 triangles by bisecting its longest side.  Repeat the
process on each sub\sphinxhyphen{}triangle, recursing as far as needed to generate a
set of covering spheres.  When finished, the original criteria are
met, and the set of covering spheres should be near minimal in number
and overlap, at least for input triangles with a reasonable
aspect\sphinxhyphen{}ratio.

\sphinxAtStartPar
The LJ interaction between 2 spheres on different triangles of types
I,J is computed with an arithmetic mixing of the sigma values of the 2
spheres and using the specified epsilon value for I,J atom types.
Note that because the sigma values for triangles spheres is computed
using only sigma\_II values, specific to the triangles’s type, this
means that any specified sigma\_IJ values (for I != J) are effectively
ignored.

\sphinxAtStartPar
For style \sphinxstyleemphasis{tri/lj}, the following coefficients must be defined for
each pair of atoms types via the {\hyperref[\detokenize{pair_coeff::doc}]{\sphinxcrossref{\DUrole{doc}{pair\_coeff}}}} command
as in the examples above, or in the data file or restart files read by
the {\hyperref[\detokenize{read_data::doc}]{\sphinxcrossref{\DUrole{doc}{read\_data}}}} or {\hyperref[\detokenize{read_restart::doc}]{\sphinxcrossref{\DUrole{doc}{read\_restart}}}}
commands:
\begin{itemize}
\item {} 
\sphinxAtStartPar
epsilon (energy units)

\item {} 
\sphinxAtStartPar
sigma (distance units)

\item {} 
\sphinxAtStartPar
cutoff (distance units)

\end{itemize}

\sphinxAtStartPar
The last coefficient is optional.  If not specified, the global cutoff
is used.


\bigskip\hrule\bigskip



\subsection{Mixing, shift, table, tail correction, restart, rRESPA info}
\label{\detokenize{pair_tri_lj:mixing-shift-table-tail-correction-restart-rrespa-info}}
\sphinxAtStartPar
For atom type pairs I,J and I != J, the epsilon and sigma coefficients
and cutoff distance for all of this pair style can be mixed.  The
default mix value is \sphinxstyleemphasis{geometric}.  See the “pair\_modify” command for
details.

\sphinxAtStartPar
This pair style does not support the {\hyperref[\detokenize{pair_modify::doc}]{\sphinxcrossref{\DUrole{doc}{pair\_modify}}}}
shift, table, and tail options.

\sphinxAtStartPar
This pair style does not write its information to {\hyperref[\detokenize{restart::doc}]{\sphinxcrossref{\DUrole{doc}{binary restart files}}}}.

\sphinxAtStartPar
This pair style can only be used via the \sphinxstyleemphasis{pair} keyword of the
{\hyperref[\detokenize{run_style::doc}]{\sphinxcrossref{\DUrole{doc}{run\_style respa}}}} command.  It does not support the
\sphinxstyleemphasis{inner}, \sphinxstyleemphasis{middle}, \sphinxstyleemphasis{outer} keywords.


\bigskip\hrule\bigskip



\subsection{Restrictions}
\label{\detokenize{pair_tri_lj:restrictions}}
\sphinxAtStartPar
This style is part of the ASPHERE package.  It is only enabled if
LAMMPS was built with that package.  See the {\hyperref[\detokenize{Build_package::doc}]{\sphinxcrossref{\DUrole{doc}{Build package}}}} page for more info.

\sphinxAtStartPar
Defining particles to be triangles so they participate in tri/tri or
tri/particle interactions requires the use the {\hyperref[\detokenize{atom_style::doc}]{\sphinxcrossref{\DUrole{doc}{atom\_style tri}}}} command.


\subsection{Related commands}
\label{\detokenize{pair_tri_lj:related-commands}}
\sphinxAtStartPar
{\hyperref[\detokenize{pair_coeff::doc}]{\sphinxcrossref{\DUrole{doc}{pair\_coeff}}}}, {\hyperref[\detokenize{pair_line_lj::doc}]{\sphinxcrossref{\DUrole{doc}{pair\_style line/lj}}}}


\subsection{Default}
\label{\detokenize{pair_tri_lj:default}}
\sphinxAtStartPar
none

\sphinxstepscope

\index{pair\_style uf3@\spxentry{pair\_style uf3}}\ignorespaces 
\index{pair\_style uf3/kk@\spxentry{pair\_style uf3/kk}}\ignorespaces 

\section{pair\_style uf3 command}
\label{\detokenize{pair_uf3:pair-style-uf3-command}}\label{\detokenize{pair_uf3:index-1}}\label{\detokenize{pair_uf3:index-0}}\label{\detokenize{pair_uf3::doc}}
\sphinxAtStartPar
Accelerator Variants: \sphinxstyleemphasis{uf3/kk}


\subsection{Syntax}
\label{\detokenize{pair_uf3:syntax}}
\begin{sphinxVerbatim}[commandchars=\\\{\}]
\PYG{k}{pair\PYGZus{}style}\PYG{+w}{ }\PYG{n}{style}\PYG{+w}{ }\PYG{n}{BodyFlag}
\end{sphinxVerbatim}
\begin{itemize}
\item {} 
\sphinxAtStartPar
style = \sphinxstyleemphasis{uf3} or \sphinxstyleemphasis{uf3/kk}

\begin{sphinxVerbatim}[commandchars=\\\{\}]
BodyFlag = Indicates whether to calculate only 2\PYGZhy{}body or 2 and 3\PYGZhy{}body interactions. Possible values: 2 or 3
\end{sphinxVerbatim}

\end{itemize}


\subsection{Examples}
\label{\detokenize{pair_uf3:examples}}
\begin{sphinxVerbatim}[commandchars=\\\{\}]
\PYG{k}{pair\PYGZus{}style}\PYG{+w}{ }\PYG{n}{uf3}\PYG{+w}{ }\PYG{l+m}{3}
\PYG{k}{pair\PYGZus{}coeff}\PYG{+w}{ }\PYG{o}{*}\PYG{+w}{ }\PYG{o}{*}\PYG{+w}{ }\PYG{n}{Nb.uf3}\PYG{+w}{ }\PYG{n}{Nb}

\PYG{k}{pair\PYGZus{}style}\PYG{+w}{ }\PYG{n}{uf3}\PYG{+w}{ }\PYG{l+m}{2}
\PYG{k}{pair\PYGZus{}coeff}\PYG{+w}{ }\PYG{o}{*}\PYG{+w}{ }\PYG{o}{*}\PYG{+w}{ }\PYG{n}{NbSn.uf3}\PYG{+w}{ }\PYG{n}{Nb}\PYG{+w}{ }\PYG{n}{Sn}

\PYG{k}{pair\PYGZus{}style}\PYG{+w}{ }\PYG{n}{uf3}\PYG{+w}{ }\PYG{l+m}{3}
\PYG{k}{pair\PYGZus{}coeff}\PYG{+w}{ }\PYG{o}{*}\PYG{+w}{ }\PYG{o}{*}\PYG{+w}{ }\PYG{n}{NbSn.uf3}\PYG{+w}{ }\PYG{n}{Nb}\PYG{+w}{ }\PYG{n}{Sn}
\end{sphinxVerbatim}


\subsection{Description}
\label{\detokenize{pair_uf3:description}}
\sphinxAtStartPar
\DUrole{versionmodified,added}{Added in version 27June2024.}

\sphinxAtStartPar
The \sphinxstyleemphasis{uf3} style computes the {\hyperref[\detokenize{pair_uf3:xie23}]{\sphinxcrossref{\DUrole{std,std-ref}{Ultra\sphinxhyphen{}Fast Force Fields (UF3)}}}} potential, a machine\sphinxhyphen{}learning interatomic potential. In UF3,
the total energy of the system is defined via two\sphinxhyphen{} and three\sphinxhyphen{}body
interactions:
\begin{equation*}
\begin{split}E & = \sum_{i,j} V_2(r_{ij}) + \sum_{i,j,k} V_3 (r_{ij},r_{ik},r_{jk}) \\
V_2(r_{ij}) & = \sum_{n=0}^N c_n B_n(r_{ij}) \\
V_3 (r_{ij},r_{ik},r_{jk}) & = \sum_{l=0}^{N_l} \sum_{m=0}^{N_m} \sum_{n=0}^{N_n} c_{l,m,n} B_l(r_{ij}) B_m(r_{ik}) B_n(r_{jk})\end{split}
\end{equation*}
\sphinxAtStartPar
where \(V_2(r_{ij})\) and \(V_3 (r_{ij},r_{ik},r_{jk})\) are the
two\sphinxhyphen{} and three\sphinxhyphen{}body interactions, respectively. For the two\sphinxhyphen{}body the
summation is over all neighbors J and for the three\sphinxhyphen{}body the summation
is over all neighbors J and K of atom I within a cutoff distance
determined from the potential files. \(B_n(r_{ij})\) are the cubic
b\sphinxhyphen{}spline basis, \(c_n\) and \(c_{l,m,n}\) are the machine\sphinxhyphen{}learned
interaction parameters and \(N\), \(N_l\), \(N_m\), and
\(N_n\) denote the number of basis functions per spline or tensor
spline dimension.

\sphinxAtStartPar
With \sphinxstyleemphasis{uf3} style only a single pair\_coeff command is used to indicate the
UF3 LAMMPS potential file containing all the two\sphinxhyphen{} and three\sphinxhyphen{}body interactions
followed by N additional arguments specifying the mapping of UF3 elements to
LAMMPS atom types, where N is the number of LAMMPS atom types:
\begin{itemize}
\item {} 
\sphinxAtStartPar
UF3 LAMMPS potential file

\item {} 
\sphinxAtStartPar
N elements names = mapping of UF3 elements to atom types

\end{itemize}

\sphinxAtStartPar
As an example, if a LAMMPS simulation contains 2 atom types (elements
‘A’ and ‘B’), the pair\_coeff command will be:

\begin{sphinxVerbatim}[commandchars=\\\{\}]
\PYG{k}{pair\PYGZus{}style}\PYG{+w}{ }\PYG{n}{uf3}\PYG{+w}{ }\PYG{l+m}{3}
\PYG{k}{pair\PYGZus{}coeff}\PYG{+w}{ }\PYG{o}{*}\PYG{+w}{ }\PYG{o}{*}\PYG{+w}{ }\PYG{n}{AB.uf3}\PYG{+w}{ }\PYG{n}{A}\PYG{+w}{ }\PYG{n}{B}
\end{sphinxVerbatim}

\sphinxAtStartPar
The AB.uf3 file should contain all two\sphinxhyphen{}body (A\sphinxhyphen{}A, A\sphinxhyphen{}B, B\sphinxhyphen{}B) and three\sphinxhyphen{}body
(A\sphinxhyphen{}A\sphinxhyphen{}A, A\sphinxhyphen{}A\sphinxhyphen{}B, A\sphinxhyphen{}B\sphinxhyphen{}B, B\sphinxhyphen{}A\sphinxhyphen{}A, B\sphinxhyphen{}A\sphinxhyphen{}B, B\sphinxhyphen{}B\sphinxhyphen{}B).

\sphinxAtStartPar
If a value of “2” is specified in the \sphinxcode{\sphinxupquote{pair\_style uf3}} command,
only the two\sphinxhyphen{}body potentials are needed. For 3\sphinxhyphen{}body interaction the
first atom type is the central atom. We recommend using the
\sphinxcode{\sphinxupquote{generate\_uf3\_lammps\_pots.py}} script (found \sphinxhref{https://github.com/uf3/uf3/tree/develop/lammps\_plugin/scripts}{here}) for
generating the UF3 LAMMPS potential file from the UF3 JSON potentials.


\bigskip\hrule\bigskip


\sphinxAtStartPar
UF3 LAMMPS potential file in the \sphinxstyleemphasis{potentials} directory of the LAMMPS
distribution have a “.uf3” suffix. The interaction block in UF3 LAMMPS potential
file should start with \sphinxcode{\sphinxupquote{\#UF3 POT}} and end with \sphinxcode{\sphinxupquote{\#}} characters.
Following shows the format of a generic 2\sphinxhyphen{}body and 3\sphinxhyphen{}body potential block in
UF3 LAMMPS potential file\sphinxhyphen{}

\begin{sphinxVerbatim}[commandchars=\\\{\}]
\PYG{c}{\PYGZsh{}UF3 POT UNITS: units DATE: POT\PYGZus{}GEN\PYGZus{}DATE AUTHOR: AUTHOR\PYGZus{}NAME CITATION: CITE}
\PYG{l+m}{2}\PYG{n}{B}\PYG{+w}{ }\PYG{n}{ELEMENT1}\PYG{+w}{ }\PYG{n}{ELEMENT2}\PYG{+w}{ }\PYG{n}{LEADING\PYGZus{}TRIM}\PYG{+w}{ }\PYG{n}{TRAILING\PYGZus{}TRIM}
\PYG{n}{Rij\PYGZus{}CUTOFF}\PYG{+w}{ }\PYG{n}{NUM\PYGZus{}OF\PYGZus{}KNOTS}
\PYG{n}{BSPLINE\PYGZus{}KNOTS}
\PYG{n}{NUM\PYGZus{}OF\PYGZus{}COEFF}
\PYG{n}{COEFF}
\PYG{c}{\PYGZsh{}}
\PYG{c}{\PYGZsh{}UF3 POT UNITS: units DATE: POT\PYGZus{}GEN\PYGZus{}DATE AUTHOR: AUTHOR\PYGZus{}NAME CITATION: CITE}
\PYG{l+m}{3}\PYG{n}{B}\PYG{+w}{ }\PYG{n}{ELEMENT1}\PYG{+w}{ }\PYG{n}{ELEMENT2}\PYG{+w}{ }\PYG{n}{ELEMENT3}\PYG{+w}{ }\PYG{n}{LEADING\PYGZus{}TRIM}\PYG{+w}{ }\PYG{n}{TRAILING\PYGZus{}TRIM}
\PYG{n}{Rjk\PYGZus{}CUTOFF}\PYG{+w}{ }\PYG{n}{Rik\PYGZus{}CUTOFF}\PYG{+w}{ }\PYG{n}{Rij\PYGZus{}CUTOFF}\PYG{+w}{ }\PYG{n}{NUM\PYGZus{}OF\PYGZus{}KNOTS\PYGZus{}JK}\PYG{+w}{ }\PYG{n}{NUM\PYGZus{}OF\PYGZus{}KNOTS\PYGZus{}IK}\PYG{+w}{ }\PYG{n}{NUM\PYGZus{}OF\PYGZus{}KNOTS\PYGZus{}IJ}
\PYG{n}{BSPLINE\PYGZus{}KNOTS\PYGZus{}FOR\PYGZus{}JK}
\PYG{n}{BSPLINE\PYGZus{}KNOTS\PYGZus{}FOR\PYGZus{}IK}
\PYG{n}{BSPLINE\PYGZus{}KNOTS\PYGZus{}FOR\PYGZus{}IJ}
\PYG{n}{SHAPE\PYGZus{}OF\PYGZus{}COEFF\PYGZus{}MATRIX[I][J][K]}
\PYG{n}{COEFF\PYGZus{}MATRIX[0][0][K]}
\PYG{n}{COEFF\PYGZus{}MATRIX[0][1][K]}
\PYG{n}{COEFF\PYGZus{}MATRIX[0][2][K]}
\PYG{n}{.}
\PYG{n}{.}
\PYG{n}{.}
\PYG{n}{COEFF\PYGZus{}MATRIX[1][0][K]}
\PYG{n}{COEFF\PYGZus{}MATRIX[1][1][K]}
\PYG{n}{COEFF\PYGZus{}MATRIX[1][2][K]}
\PYG{n}{.}
\PYG{n}{.}
\PYG{n}{.}
\PYG{c}{\PYGZsh{}}
\end{sphinxVerbatim}

\sphinxAtStartPar
The second line indicates whether the block contains data for 2\sphinxhyphen{}body
(\sphinxcode{\sphinxupquote{2B}}) or 3\sphinxhyphen{}body (\sphinxcode{\sphinxupquote{3B}}) interaction. This is followed by element
combination interaction, \sphinxcode{\sphinxupquote{LEADING\_TRIM}} and \sphinxcode{\sphinxupquote{TRAILING\_TRIM}}
number on the same line. The current implementation is only tested for
\sphinxcode{\sphinxupquote{LEADING\_TRIM=0}} and \sphinxcode{\sphinxupquote{TRAILING\_TRIM=3}}.
If other values are used LAMMPS is terminated after issuing an error message.
The \sphinxcode{\sphinxupquote{Rij\_CUTOFF}} sets the 2\sphinxhyphen{}body cutoff for the interaction described
by the potential block. \sphinxcode{\sphinxupquote{NUM\_OF\_KNOTS}} is the number of knots
(or the length of the knot vector) present on the very next line. The
\sphinxcode{\sphinxupquote{BSPLINE\_KNOTS}} line should contain all the knots in ascending order.
\sphinxcode{\sphinxupquote{NUM\_OF\_COEFF}} is the number of coefficients in the \sphinxcode{\sphinxupquote{COEFF}} line.
All the numbers in the BSPLINE\_KNOTS and COEFF line should be space\sphinxhyphen{}separated.
Similar to the 2\sphinxhyphen{}body potential block, the third line sets the cutoffs and
length of the knots. The cutoff distance between atom\sphinxhyphen{}type I and J is
\sphinxcode{\sphinxupquote{Rij\_CUTOFF}}, atom\sphinxhyphen{}type I and K is \sphinxcode{\sphinxupquote{Rik\_CUTOFF}} and between
J and K is \sphinxcode{\sphinxupquote{Rjk\_CUTOFF}}.

\begin{sphinxadmonition}{note}{Note:}
\sphinxAtStartPar
The current implementation only works for UF3 potentials with cutoff
distances for 3\sphinxhyphen{}body interactions that follows
\sphinxcode{\sphinxupquote{2Rij\_CUTOFF=2Rik\_CUTOFF=Rjk\_CUTOFF}} relation.
\end{sphinxadmonition}

\sphinxAtStartPar
The \sphinxcode{\sphinxupquote{BSPLINE\_KNOTS\_FOR\_JK}}, \sphinxcode{\sphinxupquote{BSPLINE\_KNOTS\_FOR\_IK}}, and
\sphinxcode{\sphinxupquote{BSPLINE\_KNOTS\_FOR\_IJ}} lines (note the order) contain the knots in
increasing order for atoms J and K, I and K, and atoms I and J
respectively. The number of knots is defined by the
\sphinxcode{\sphinxupquote{NUM\_OF\_KNOTS\_*}} characters in the previous line.  The shape of
the coefficient matrix is defined on the
\sphinxcode{\sphinxupquote{SHAPE\_OF\_COEFF\_MATRIX{[}I{]}{[}J{]}{[}K{]}}} line followed by the columns of
the coefficient matrix, one per line, as shown above. For example, if
the coefficient matrix has the shape of 8x8x13, then
\sphinxcode{\sphinxupquote{SHAPE\_OF\_COEFF\_MATRIX{[}I{]}{[}J{]}{[}K{]}}} will be \sphinxcode{\sphinxupquote{8 8 13}} followed
by 64 (8x8) lines each containing 13 coefficients separated by space.


\bigskip\hrule\bigskip


\sphinxAtStartPar
Styles with a \sphinxstyleemphasis{gpu}, \sphinxstyleemphasis{intel}, \sphinxstyleemphasis{kk}, \sphinxstyleemphasis{omp}, or \sphinxstyleemphasis{opt} suffix are
functionally the same as the corresponding style without the suffix.
They have been optimized to run faster, depending on your available
hardware, as discussed on the {\hyperref[\detokenize{Speed_packages::doc}]{\sphinxcrossref{\DUrole{doc}{Accelerator packages}}}}
page.  The accelerated styles take the same arguments and should
produce the same results, except for round\sphinxhyphen{}off and precision issues.

\sphinxAtStartPar
These accelerated styles are part of the GPU, INTEL, KOKKOS,
OPENMP, and OPT packages, respectively.  They are only enabled if
LAMMPS was built with those packages.  See the {\hyperref[\detokenize{Build_package::doc}]{\sphinxcrossref{\DUrole{doc}{Build package}}}} page for more info.

\sphinxAtStartPar
You can specify the accelerated styles explicitly in your input script
by including their suffix, or you can use the {\hyperref[\detokenize{Run_options::doc}]{\sphinxcrossref{\DUrole{doc}{\sphinxhyphen{}suffix command\sphinxhyphen{}line switch}}}} when you invoke LAMMPS, or you can use the
{\hyperref[\detokenize{suffix::doc}]{\sphinxcrossref{\DUrole{doc}{suffix}}}} command in your input script.

\sphinxAtStartPar
See the {\hyperref[\detokenize{Speed_packages::doc}]{\sphinxcrossref{\DUrole{doc}{Accelerator packages}}}} page for more
instructions on how to use the accelerated styles effectively.


\bigskip\hrule\bigskip



\subsection{Mixing, shift, table, tail correction, restart, rRESPA info}
\label{\detokenize{pair_uf3:mixing-shift-table-tail-correction-restart-rrespa-info}}
\sphinxAtStartPar
For atom type pairs I,J and I != J, where types I and J correspond to
two different element types, mixing is performed by LAMMPS as described
above from values in the potential file.

\sphinxAtStartPar
This pair style does not support the {\hyperref[\detokenize{pair_modify::doc}]{\sphinxcrossref{\DUrole{doc}{pair\_modify}}}}
shift, table, and tail options.

\sphinxAtStartPar
This pair style does not write its information to {\hyperref[\detokenize{restart::doc}]{\sphinxcrossref{\DUrole{doc}{binary restart
files}}}}, since it is stored in potential file.

\sphinxAtStartPar
This pair style can only be used via the \sphinxstyleemphasis{pair} keyword of the
{\hyperref[\detokenize{run_style::doc}]{\sphinxcrossref{\DUrole{doc}{run\_style respa}}}} command.  It does not support the
\sphinxstyleemphasis{inner}, \sphinxstyleemphasis{middle}, \sphinxstyleemphasis{outer} keywords.


\subsection{Restrictions}
\label{\detokenize{pair_uf3:restrictions}}
\sphinxAtStartPar
The ‘uf3’ pair style is part of the ML\sphinxhyphen{}UF3 package. It is only enabled
if LAMMPS was built with that package. See the {\hyperref[\detokenize{Build_package::doc}]{\sphinxcrossref{\DUrole{doc}{Build package}}}} page for more info.

\sphinxAtStartPar
This pair style requires the {\hyperref[\detokenize{newton::doc}]{\sphinxcrossref{\DUrole{doc}{newton}}}} setting to be “on”.

\sphinxAtStartPar
The UF3 LAMMPS potential file provided with LAMMPS (see the potentials
directory) are parameterized for metal {\hyperref[\detokenize{units::doc}]{\sphinxcrossref{\DUrole{doc}{units}}}}.

\sphinxAtStartPar
The single() function of ‘uf3’ pair style only return the 2\sphinxhyphen{}body
interaction energy.


\subsection{Related commands}
\label{\detokenize{pair_uf3:related-commands}}
\sphinxAtStartPar
{\hyperref[\detokenize{pair_coeff::doc}]{\sphinxcrossref{\DUrole{doc}{pair\_coeff}}}}


\subsection{Default}
\label{\detokenize{pair_uf3:default}}
\sphinxAtStartPar
none


\bigskip\hrule\bigskip

\phantomsection\label{\detokenize{pair_uf3:xie23}}
\sphinxAtStartPar
\sphinxstylestrong{(Xie23)} Xie, S.R., Rupp, M. \& Hennig, R.G. Ultra\sphinxhyphen{}fast interpretable machine\sphinxhyphen{}learning potentials. npj Comput Mater 9, 162 (2023). \sphinxurl{https://doi.org/10.1038/s41524-023-01092-7}

\sphinxstepscope

\index{pair\_style ufm@\spxentry{pair\_style ufm}}\ignorespaces 
\index{pair\_style ufm/gpu@\spxentry{pair\_style ufm/gpu}}\ignorespaces 
\index{pair\_style ufm/omp@\spxentry{pair\_style ufm/omp}}\ignorespaces 
\index{pair\_style ufm/opt@\spxentry{pair\_style ufm/opt}}\ignorespaces 

\section{pair\_style ufm command}
\label{\detokenize{pair_ufm:pair-style-ufm-command}}\label{\detokenize{pair_ufm:index-3}}\label{\detokenize{pair_ufm:index-2}}\label{\detokenize{pair_ufm:index-1}}\label{\detokenize{pair_ufm:index-0}}\label{\detokenize{pair_ufm::doc}}
\sphinxAtStartPar
Accelerator Variants: \sphinxstyleemphasis{ufm/gpu}, \sphinxstyleemphasis{ufm/omp}, \sphinxstyleemphasis{ufm/opt}


\subsection{Syntax}
\label{\detokenize{pair_ufm:syntax}}
\begin{sphinxVerbatim}[commandchars=\\\{\}]
\PYG{k}{pair\PYGZus{}style}\PYG{+w}{ }\PYG{n}{ufm}\PYG{+w}{ }\PYG{n}{cutoff}
\end{sphinxVerbatim}
\begin{itemize}
\item {} 
\sphinxAtStartPar
cutoff = global cutoff for \sphinxstyleemphasis{ufm} interactions (distance units)

\end{itemize}


\subsection{Examples}
\label{\detokenize{pair_ufm:examples}}
\begin{sphinxVerbatim}[commandchars=\\\{\}]
\PYG{k}{pair\PYGZus{}style}\PYG{+w}{ }\PYG{n}{ufm}\PYG{+w}{ }\PYG{l+m}{4.0}
\PYG{k}{pair\PYGZus{}coeff}\PYG{+w}{ }\PYG{l+m}{1}\PYG{+w}{ }\PYG{l+m}{1}\PYG{+w}{ }\PYG{l+m}{100.0}\PYG{+w}{ }\PYG{l+m}{1.0}\PYG{+w}{ }\PYG{l+m}{2.5}
\PYG{k}{pair\PYGZus{}coeff}\PYG{+w}{ }\PYG{o}{*}\PYG{+w}{ }\PYG{o}{*}\PYG{+w}{ }\PYG{l+m}{100.0}\PYG{+w}{ }\PYG{l+m}{1.0}

\PYG{k}{pair\PYGZus{}style}\PYG{+w}{ }\PYG{n}{ufm}\PYG{+w}{ }\PYG{l+m}{4.0}
\PYG{k}{pair\PYGZus{}coeff}\PYG{+w}{ }\PYG{o}{*}\PYG{+w}{ }\PYG{o}{*}\PYG{+w}{ }\PYG{l+m}{10.0}\PYG{+w}{ }\PYG{l+m}{1.0}
\PYG{k}{variable }\PYG{n+nv+nvIdentifier}{prefactor}\PYG{+w}{ }\PYG{n}{equal}\PYG{+w}{ }\PYG{n}{ramp}\PYG{n+nv}{(}\PYG{n+nv}{10,100}\PYG{n+nv}{)}
\PYG{k}{fix }\PYG{n+nv+nvIdentifier}{1}\PYG{+w}{ }\PYG{n+nv+nvIdentifier}{all}\PYG{+w}{ }\PYG{n}{adapt}\PYG{+w}{ }\PYG{l+m}{1}\PYG{+w}{ }\PYG{n}{pair}\PYG{+w}{ }\PYG{n}{ufm}\PYG{+w}{ }\PYG{n}{epsilon}\PYG{+w}{ }\PYG{o}{*}\PYG{+w}{ }\PYG{o}{*}\PYG{+w}{ }\PYG{n}{v\PYGZus{}prefactor}
\end{sphinxVerbatim}


\subsection{Description}
\label{\detokenize{pair_ufm:description}}
\sphinxAtStartPar
Style \sphinxstyleemphasis{ufm} computes pairwise interactions using the Uhlenbeck\sphinxhyphen{}Ford model (UFM) potential {\hyperref[\detokenize{pair_ufm:pl2}]{\sphinxcrossref{\DUrole{std,std-ref}{(Paula Leite2016)}}}} which is given by
\begin{equation*}
\begin{split}E & = -\varepsilon\, \ln{\left[1-\exp{\left(-r^{2}/\sigma^{2}\right)}\right]} \qquad  r < r_c \\
\varepsilon & = p\,k_B\,T\end{split}
\end{equation*}
\sphinxAtStartPar
where \(r_c\) is the cutoff, \(\sigma\) is a distance\sphinxhyphen{}scale and
\(\epsilon\) is an energy\sphinxhyphen{}scale, i.e., a product of Boltzmann constant
\(k_B\), temperature \(T\) and the Uhlenbeck\sphinxhyphen{}Ford p\sphinxhyphen{}parameter which
is responsible
to control the softness of the interactions {\hyperref[\detokenize{pair_ufm:pl1}]{\sphinxcrossref{\DUrole{std,std-ref}{(Paula Leite2017)}}}}.
This model is useful as a reference system for fluid\sphinxhyphen{}phase free\sphinxhyphen{}energy calculations {\hyperref[\detokenize{pair_ufm:pl2}]{\sphinxcrossref{\DUrole{std,std-ref}{(Paula Leite2016)}}}}.

\sphinxAtStartPar
The following coefficients must be defined for each pair of atom types
via the {\hyperref[\detokenize{pair_coeff::doc}]{\sphinxcrossref{\DUrole{doc}{pair\_coeff}}}} command as in the examples above,
or in the data file or restart files read by the
{\hyperref[\detokenize{read_data::doc}]{\sphinxcrossref{\DUrole{doc}{read\_data}}}} or {\hyperref[\detokenize{read_restart::doc}]{\sphinxcrossref{\DUrole{doc}{read\_restart}}}}
commands, or by mixing as described below:
\begin{itemize}
\item {} 
\sphinxAtStartPar
\(\epsilon\) (energy units)

\item {} 
\sphinxAtStartPar
\(\sigma\) (distance units)

\item {} 
\sphinxAtStartPar
cutoff (distance units)

\end{itemize}

\sphinxAtStartPar
The last coefficient is optional.  If not specified, the global \sphinxstyleemphasis{ufm}
cutoff is used.

\sphinxAtStartPar
The {\hyperref[\detokenize{fix_adapt::doc}]{\sphinxcrossref{\DUrole{doc}{fix adapt}}}} command can be used to vary epsilon and sigma for this pair style over the course of a simulation, in which case
pair\_coeff settings for epsilon and sigma must still be specified, but will be
overridden.  For example these commands will vary the prefactor epsilon for
all pairwise interactions from 10.0 at the beginning to 100.0 at the end
of a run:

\begin{sphinxVerbatim}[commandchars=\\\{\}]
\PYG{k}{variable }\PYG{n+nv+nvIdentifier}{prefactor}\PYG{+w}{ }\PYG{n}{equal}\PYG{+w}{ }\PYG{n}{ramp}\PYG{n+nv}{(}\PYG{n+nv}{10,100}\PYG{n+nv}{)}
\PYG{k}{fix }\PYG{n+nv+nvIdentifier}{1}\PYG{+w}{ }\PYG{n+nv+nvIdentifier}{all}\PYG{+w}{ }\PYG{n}{adapt}\PYG{+w}{ }\PYG{l+m}{1}\PYG{+w}{ }\PYG{n}{pair}\PYG{+w}{ }\PYG{n}{ufm}\PYG{+w}{ }\PYG{n}{epsilon}\PYG{+w}{ }\PYG{o}{*}\PYG{+w}{ }\PYG{o}{*}\PYG{+w}{ }\PYG{n}{v\PYGZus{}prefactor}
\end{sphinxVerbatim}

\begin{sphinxadmonition}{note}{Note:}
\sphinxAtStartPar
The thermodynamic integration procedure can be performed with this
potential using {\hyperref[\detokenize{fix_adapt::doc}]{\sphinxcrossref{\DUrole{doc}{fix adapt}}}}. This command will
rescale the force on each atom by varying a scale variable, which
always starts with value 1.0. The syntax is the same described above,
however, changing epsilon to scale. A detailed explanation of how to
use this command and perform nonequilibrium thermodynamic integration
in LAMMPS is given in the paper by {\hyperref[\detokenize{pair_ufm:freitas2}]{\sphinxcrossref{\DUrole{std,std-ref}{(Freitas)}}}}.
\end{sphinxadmonition}


\bigskip\hrule\bigskip


\sphinxAtStartPar
Styles with a \sphinxstyleemphasis{gpu}, \sphinxstyleemphasis{intel}, \sphinxstyleemphasis{kk}, \sphinxstyleemphasis{omp}, or \sphinxstyleemphasis{opt} suffix are
functionally the same as the corresponding style without the suffix.
They have been optimized to run faster, depending on your available
hardware, as discussed on the {\hyperref[\detokenize{Speed_packages::doc}]{\sphinxcrossref{\DUrole{doc}{Accelerator packages}}}}
page.  The accelerated styles take the same arguments and should
produce the same results, except for round\sphinxhyphen{}off and precision issues.

\sphinxAtStartPar
These accelerated styles are part of the GPU, INTEL, KOKKOS,
OPENMP, and OPT packages, respectively.  They are only enabled if
LAMMPS was built with those packages.  See the {\hyperref[\detokenize{Build_package::doc}]{\sphinxcrossref{\DUrole{doc}{Build package}}}} page for more info.

\sphinxAtStartPar
You can specify the accelerated styles explicitly in your input script
by including their suffix, or you can use the {\hyperref[\detokenize{Run_options::doc}]{\sphinxcrossref{\DUrole{doc}{\sphinxhyphen{}suffix command\sphinxhyphen{}line switch}}}} when you invoke LAMMPS, or you can use the
{\hyperref[\detokenize{suffix::doc}]{\sphinxcrossref{\DUrole{doc}{suffix}}}} command in your input script.

\sphinxAtStartPar
See the {\hyperref[\detokenize{Speed_packages::doc}]{\sphinxcrossref{\DUrole{doc}{Accelerator packages}}}} page for more
instructions on how to use the accelerated styles effectively.


\bigskip\hrule\bigskip



\subsection{Mixing, shift, table, tail correction, restart, rRESPA info}
\label{\detokenize{pair_ufm:mixing-shift-table-tail-correction-restart-rrespa-info}}
\sphinxAtStartPar
For atom type pairs I,J and I != J, the epsilon and sigma coefficients and cutoff
distance for this pair style can be mixed. The default mix value is \sphinxstyleemphasis{geometric}.  See the
“pair\_modify” command for details.

\sphinxAtStartPar
This pair style support the {\hyperref[\detokenize{pair_modify::doc}]{\sphinxcrossref{\DUrole{doc}{pair\_modify}}}} shift option for the energy of the pair interaction.

\sphinxAtStartPar
The {\hyperref[\detokenize{pair_modify::doc}]{\sphinxcrossref{\DUrole{doc}{pair\_modify}}}} table and tail are not relevant for this
pair style.

\sphinxAtStartPar
This pair style does not support the {\hyperref[\detokenize{pair_modify::doc}]{\sphinxcrossref{\DUrole{doc}{pair\_modify}}}} tail option for adding long\sphinxhyphen{}range tail corrections to energy and pressure.

\sphinxAtStartPar
This pair style writes its information to {\hyperref[\detokenize{restart::doc}]{\sphinxcrossref{\DUrole{doc}{binary restart files}}}}, so pair\_style and pair\_coeff commands do not need
to be specified in an input script that reads a restart file.

\sphinxAtStartPar
This pair style can only be used via the \sphinxstyleemphasis{pair} keyword of the
{\hyperref[\detokenize{run_style::doc}]{\sphinxcrossref{\DUrole{doc}{run\_style respa}}}} command.  It does not support the
\sphinxstyleemphasis{inner}, \sphinxstyleemphasis{middle}, \sphinxstyleemphasis{outer} keywords.


\bigskip\hrule\bigskip



\subsection{Restrictions}
\label{\detokenize{pair_ufm:restrictions}}
\sphinxAtStartPar
This pair style is part of the EXTRA\sphinxhyphen{}PAIR package.  It is only enabled if
LAMMPS was built with that package.  See the
{\hyperref[\detokenize{Build_package::doc}]{\sphinxcrossref{\DUrole{doc}{Build package}}}} page for more info.


\subsection{Related commands}
\label{\detokenize{pair_ufm:related-commands}}
\sphinxAtStartPar
{\hyperref[\detokenize{pair_coeff::doc}]{\sphinxcrossref{\DUrole{doc}{pair\_coeff}}}}, {\hyperref[\detokenize{fix_adapt::doc}]{\sphinxcrossref{\DUrole{doc}{fix adapt}}}}


\subsection{Default}
\label{\detokenize{pair_ufm:default}}
\sphinxAtStartPar
none

\phantomsection\label{\detokenize{pair_ufm:pl1}}
\sphinxAtStartPar
\sphinxstylestrong{(Paula Leite2017)} Paula Leite, Santos\sphinxhyphen{}Florez, and de Koning, Phys Rev E, 96,
32115 (2017).

\phantomsection\label{\detokenize{pair_ufm:pl2}}
\sphinxAtStartPar
\sphinxstylestrong{(Paula Leite2016)} Paula Leite , Freitas, Azevedo, and de Koning, J Chem Phys, 126,
044509 (2016).

\phantomsection\label{\detokenize{pair_ufm:freitas2}}
\sphinxAtStartPar
\sphinxstylestrong{(Freitas)} Freitas, Asta, and de Koning, Computational Materials Science, 112, 333 (2016).

\sphinxstepscope

\index{pair\_style vashishta@\spxentry{pair\_style vashishta}}\ignorespaces 
\index{pair\_style vashishta/gpu@\spxentry{pair\_style vashishta/gpu}}\ignorespaces 
\index{pair\_style vashishta/omp@\spxentry{pair\_style vashishta/omp}}\ignorespaces 
\index{pair\_style vashishta/kk@\spxentry{pair\_style vashishta/kk}}\ignorespaces 
\index{pair\_style vashishta/table@\spxentry{pair\_style vashishta/table}}\ignorespaces 
\index{pair\_style vashishta/table/omp@\spxentry{pair\_style vashishta/table/omp}}\ignorespaces 

\section{pair\_style vashishta command}
\label{\detokenize{pair_vashishta:pair-style-vashishta-command}}\label{\detokenize{pair_vashishta:index-5}}\label{\detokenize{pair_vashishta:index-4}}\label{\detokenize{pair_vashishta:index-3}}\label{\detokenize{pair_vashishta:index-2}}\label{\detokenize{pair_vashishta:index-1}}\label{\detokenize{pair_vashishta:index-0}}\label{\detokenize{pair_vashishta::doc}}
\sphinxAtStartPar
Accelerator Variants: \sphinxstyleemphasis{vashishta/gpu}, \sphinxstyleemphasis{vashishta/omp}, \sphinxstyleemphasis{vashishta/kk}


\section{pair\_style vashishta/table command}
\label{\detokenize{pair_vashishta:pair-style-vashishta-table-command}}
\sphinxAtStartPar
Accelerator Variants: \sphinxstyleemphasis{vashishta/table/omp}


\subsection{Syntax}
\label{\detokenize{pair_vashishta:syntax}}
\begin{sphinxVerbatim}[commandchars=\\\{\}]
\PYG{k}{pair\PYGZus{}style}\PYG{+w}{ }\PYG{n}{style}\PYG{+w}{ }\PYG{n}{args}
\end{sphinxVerbatim}
\begin{itemize}
\item {} 
\sphinxAtStartPar
style = \sphinxstyleemphasis{vashishta} or \sphinxstyleemphasis{vashishta/table}

\item {} 
\sphinxAtStartPar
args = list of arguments for a particular style

\end{itemize}
\begin{sphinxalltt}
\sphinxstyleemphasis{vashishta} args = none
\sphinxstyleemphasis{vashishta/table} args = Ntable cutinner
  Ntable = \# of tabulation points
  cutinner = tablulate from cutinner to cutoff
\end{sphinxalltt}


\subsection{Examples}
\label{\detokenize{pair_vashishta:examples}}
\begin{sphinxVerbatim}[commandchars=\\\{\}]
\PYG{k}{pair\PYGZus{}style}\PYG{+w}{ }\PYG{n}{vashishta}
\PYG{k}{pair\PYGZus{}coeff}\PYG{+w}{ }\PYG{o}{*}\PYG{+w}{ }\PYG{o}{*}\PYG{+w}{ }\PYG{n}{SiC.vashishta}\PYG{+w}{ }\PYG{n}{Si}\PYG{+w}{ }\PYG{n}{C}

\PYG{k}{pair\PYGZus{}style}\PYG{+w}{ }\PYG{n}{vashishta}\PYG{o}{/}\PYG{n}{table}\PYG{+w}{ }\PYG{l+m}{100000}\PYG{+w}{ }\PYG{l+m}{0.2}
\PYG{k}{pair\PYGZus{}coeff}\PYG{+w}{ }\PYG{o}{*}\PYG{+w}{ }\PYG{o}{*}\PYG{+w}{ }\PYG{n}{SiC.vashishta}\PYG{+w}{ }\PYG{n}{Si}\PYG{+w}{ }\PYG{n}{C}
\end{sphinxVerbatim}


\subsection{Description}
\label{\detokenize{pair_vashishta:description}}
\sphinxAtStartPar
The \sphinxstyleemphasis{vashishta} and \sphinxstyleemphasis{vashishta/table} styles compute the combined
2\sphinxhyphen{}body and 3\sphinxhyphen{}body family of potentials developed in the group of Priya
Vashishta and collaborators.  By combining repulsive, screened
Coulombic, screened charge\sphinxhyphen{}dipole, and dispersion interactions with a
bond\sphinxhyphen{}angle energy based on the Stillinger\sphinxhyphen{}Weber potential, this
potential has been used to describe a variety of inorganic compounds,
including SiO2 {\hyperref[\detokenize{pair_vashishta:vashishta1990}]{\sphinxcrossref{\DUrole{std,std-ref}{Vashishta1990}}}}, SiC
{\hyperref[\detokenize{pair_vashishta:vashishta2007}]{\sphinxcrossref{\DUrole{std,std-ref}{Vashishta2007}}}}, and InP {\hyperref[\detokenize{pair_vashishta:branicio2009}]{\sphinxcrossref{\DUrole{std,std-ref}{Branicio2009}}}}.

\sphinxAtStartPar
The potential for the energy U of a system of atoms is
\begin{equation*}
\begin{split}U & =  \sum_i^N \sum_{j > i}^N U_{ij}^{(2)} (r_{ij}) + \sum_i^N \sum_{j \neq i}^N \sum_{k > j, k \neq i}^N U_{ijk}^{(3)} (r_{ij}, r_{ik}, \theta_{ijk}) \\
U_{ij}^{(2)} (r) & =   \frac{H_{ij}}{r^{\eta_{ij}}} + \frac{Z_i Z_j}{r}\exp(-r/\lambda_{1,ij}) - \frac{D_{ij}}{r^4}\exp(-r/\lambda_{4,ij}) - \frac{W_{ij}}{r^6}, r < r_{c,{ij}} \\
U_{ijk}^{(3)}(r_{ij},r_{ik},\theta_{ijk}) & =  B_{ijk} \frac{\left[ \cos \theta_{ijk} - \cos \theta_{0ijk} \right]^2} {1+C_{ijk}\left[ \cos \theta_{ijk} - \cos \theta_{0ijk} \right]^2} \times \\
                 &  \exp \left( \frac{\gamma_{ij}}{r_{ij} - r_{0,ij}} \right) \exp \left( \frac{\gamma_{ik}}{r_{ik} - r_{0,ik}} \right), r_{ij} < r_{0,ij}, r_{ik} < r_{0,ik}\end{split}
\end{equation*}
\sphinxAtStartPar
where we follow the notation used in {\hyperref[\detokenize{pair_vashishta:branicio2009}]{\sphinxcrossref{\DUrole{std,std-ref}{Branicio2009}}}}.
\(U^2\) is a two\sphinxhyphen{}body term and U3 is a three\sphinxhyphen{}body term.  The
summation over two\sphinxhyphen{}body terms is over all neighbors J within
a cutoff distance = \(r_c\).  The twobody terms are shifted and
tilted by a linear function so that the energy and force are
both zero at \(r_c\). The summation over three\sphinxhyphen{}body terms
is over all neighbors \sphinxstyleemphasis{i} and \sphinxstyleemphasis{k} within a cut\sphinxhyphen{}off distance \(= r_0\),
where the exponential screening function becomes zero.

\sphinxAtStartPar
The \sphinxstyleemphasis{vashishta} style computes these formulas analytically.  The
\sphinxstyleemphasis{vashishta/table} style tabulates the analytic values for \sphinxstyleemphasis{Ntable}
points from cutinner to the cutoff of the potential.  The points are
equally spaced in R\textasciicircum{}2 space from cutinner\textasciicircum{}2 to cutoff\textasciicircum{}2.  For the
two\sphinxhyphen{}body term in the above equation, a linear interpolation for each
pairwise distance between adjacent points in the table.  In practice
the tabulated version can run 3\sphinxhyphen{}5x faster than the analytic version
with moderate to little loss of accuracy for Ntable values
between 10000 and 1000000. It is not recommended to use less than
5000 tabulation points.

\sphinxAtStartPar
Only a single pair\_coeff command is used with either style which
specifies a Vashishta potential file with parameters for all needed
elements.  These are mapped to LAMMPS atom types by specifying N
additional arguments after the filename in the pair\_coeff command,
where N is the number of LAMMPS atom types:
\begin{itemize}
\item {} 
\sphinxAtStartPar
filename

\item {} 
\sphinxAtStartPar
N element names = mapping of Vashishta elements to atom types

\end{itemize}

\sphinxAtStartPar
See the {\hyperref[\detokenize{pair_coeff::doc}]{\sphinxcrossref{\DUrole{doc}{pair\_coeff}}}} page for alternate ways
to specify the path for the potential file.

\sphinxAtStartPar
As an example, imagine a file SiC.vashishta has parameters for
Si and C.  If your LAMMPS simulation has 4 atoms types and you want
the first 3 to be Si, and the fourth to be C, you would use the following
pair\_coeff command:

\begin{sphinxVerbatim}[commandchars=\\\{\}]
\PYG{k}{pair\PYGZus{}coeff}\PYG{+w}{ }\PYG{o}{*}\PYG{+w}{ }\PYG{o}{*}\PYG{+w}{ }\PYG{n}{SiC.vashishta}\PYG{+w}{ }\PYG{n}{Si}\PYG{+w}{ }\PYG{n}{Si}\PYG{+w}{ }\PYG{n}{Si}\PYG{+w}{ }\PYG{n}{C}
\end{sphinxVerbatim}

\sphinxAtStartPar
The first 2 arguments must be * * so as to span all LAMMPS atom types.
The first three Si arguments map LAMMPS atom types 1,2,3 to the Si
element in the file.  The final C argument maps LAMMPS atom type 4
to the C element in the file.  If a mapping value is specified as
NULL, the mapping is not performed.  This can be used when a \sphinxstyleemphasis{vashishta}
potential is used as part of the \sphinxstyleemphasis{hybrid} pair style.  The NULL values
are placeholders for atom types that will be used with other
potentials.

\sphinxAtStartPar
Vashishta files in the \sphinxstyleemphasis{potentials} directory of the LAMMPS
distribution have a “.vashishta” suffix.  Lines that are not blank or
comments (starting with \#) define parameters for a triplet of
elements.  The parameters in a single entry correspond to the two\sphinxhyphen{}body
and three\sphinxhyphen{}body coefficients in the formulae above:
\begin{itemize}
\item {} 
\sphinxAtStartPar
element 1 (the center atom in a 3\sphinxhyphen{}body interaction)

\item {} 
\sphinxAtStartPar
element 2

\item {} 
\sphinxAtStartPar
element 3

\item {} 
\sphinxAtStartPar
\sphinxstyleemphasis{H} (energy units)

\item {} 
\sphinxAtStartPar
\(\eta\)

\item {} 
\sphinxAtStartPar
\(Z_i\) (electron charge units)

\item {} 
\sphinxAtStartPar
\(Z_j\) (electron charge units)

\item {} 
\sphinxAtStartPar
\(\lambda_1\) (distance units)

\item {} 
\sphinxAtStartPar
\sphinxstyleemphasis{D} (energy units)

\item {} 
\sphinxAtStartPar
\(\lambda_4\) (distance units)

\item {} 
\sphinxAtStartPar
\sphinxstyleemphasis{W} (energy units)

\item {} 
\sphinxAtStartPar
\(r_c\) (distance units)

\item {} 
\sphinxAtStartPar
\sphinxstyleemphasis{B} (energy units)

\item {} 
\sphinxAtStartPar
\(\gamma\)

\item {} 
\sphinxAtStartPar
\(r_0\) (distance units)

\item {} 
\sphinxAtStartPar
\sphinxstyleemphasis{C}

\item {} 
\sphinxAtStartPar
\(\cos\theta_0\)

\end{itemize}

\sphinxAtStartPar
The non\sphinxhyphen{}annotated parameters are unitless.  The Vashishta potential
file must contain entries for all the elements listed in the
pair\_coeff command.  It can also contain entries for additional
elements not being used in a particular simulation; LAMMPS ignores
those entries.  For a single\sphinxhyphen{}element simulation, only a single entry
is required (e.g. SiSiSi).  For a two\sphinxhyphen{}element simulation, the file
must contain 8 entries (for SiSiSi, SiSiC, SiCSi, SiCC, CSiSi, CSiC,
CCSi, CCC), that specify parameters for all permutations of the two
elements interacting in three\sphinxhyphen{}body configurations.  Thus for 3
elements, 27 entries would be required, etc.

\sphinxAtStartPar
Depending on the particular version of the Vashishta potential, the
values of these parameters may be keyed to the identities of zero,
one, two, or three elements.  In order to make the input file format
unambiguous, general, and simple to code, LAMMPS uses a slightly
confusing method for specifying parameters.  All parameters are
divided into two classes: two\sphinxhyphen{}body and three\sphinxhyphen{}body.  Two\sphinxhyphen{}body and
three\sphinxhyphen{}body parameters are handled differently, as described below.
The two\sphinxhyphen{}body parameters are \sphinxstyleemphasis{H}, \(\eta\), \(\lambda_1\),
\sphinxstyleemphasis{D}, \(\lambda_4\), \sphinxstyleemphasis{W}, \(r_c\), \(\gamma\),
and \(r_0\).  They appear in the above formulae with two subscripts.
The parameters \(Z_i\) and \(Z_j\) are also classified
as two\sphinxhyphen{}body parameters, even
though they only have 1 subscript.  The three\sphinxhyphen{}body parameters are \sphinxstyleemphasis{B},
\sphinxstyleemphasis{C}, \(\cos\theta_0\).  They appear in the above formulae with
three subscripts.  Two\sphinxhyphen{}body and three\sphinxhyphen{}body parameters are handled
differently, as described below.

\sphinxAtStartPar
The first element in each entry is the center atom in a three\sphinxhyphen{}body
interaction, while the second and third elements are two neighbor
atoms. Three\sphinxhyphen{}body parameters for a central atom I and two neighbors J
and K are taken from the IJK entry.  Note that even though three\sphinxhyphen{}body
parameters do not depend on the order of J and K, LAMMPS stores
three\sphinxhyphen{}body parameters for both IJK and IKJ.  The user must ensure that
these values are equal.  Two\sphinxhyphen{}body parameters for an atom I interacting
with atom J are taken from the IJJ entry, where the second and third
elements are the same. Thus the two\sphinxhyphen{}body parameters for Si interacting
with C come from the SiCC entry. Note that even though two\sphinxhyphen{}body
parameters (except possibly gamma and r0 in U3) do not depend on the
order of the two elements, LAMMPS will get the Si\sphinxhyphen{}C value from the
SiCC entry and the C\sphinxhyphen{}Si value from the CSiSi entry. The user must
ensure that these values are equal. Two\sphinxhyphen{}body parameters appearing in
entries where the second and third elements are different are stored but
never used. It is good practice to enter zero for these values. Note
that the three\sphinxhyphen{}body function U3 above contains the two\sphinxhyphen{}body parameters
\(\gamma\) and \(r_0\). So U3 for a central C atom bonded to
an Si atom and a
second C atom will take three\sphinxhyphen{}body parameters from the CSiC entry, but
two\sphinxhyphen{}body parameters from the CCC and CSiSi entries.


\bigskip\hrule\bigskip


\sphinxAtStartPar
Styles with a \sphinxstyleemphasis{gpu}, \sphinxstyleemphasis{intel}, \sphinxstyleemphasis{kk}, \sphinxstyleemphasis{omp}, or \sphinxstyleemphasis{opt} suffix are
functionally the same as the corresponding style without the suffix.
They have been optimized to run faster, depending on your available
hardware, as discussed on the {\hyperref[\detokenize{Speed_packages::doc}]{\sphinxcrossref{\DUrole{doc}{Accelerator packages}}}}
page.  The accelerated styles take the same arguments and should
produce the same results, except for round\sphinxhyphen{}off and precision issues.

\sphinxAtStartPar
These accelerated styles are part of the GPU, INTEL, KOKKOS,
OPENMP, and OPT packages, respectively.  They are only enabled if
LAMMPS was built with those packages.  See the {\hyperref[\detokenize{Build_package::doc}]{\sphinxcrossref{\DUrole{doc}{Build package}}}} page for more info.

\sphinxAtStartPar
You can specify the accelerated styles explicitly in your input script
by including their suffix, or you can use the {\hyperref[\detokenize{Run_options::doc}]{\sphinxcrossref{\DUrole{doc}{\sphinxhyphen{}suffix command\sphinxhyphen{}line switch}}}} when you invoke LAMMPS, or you can use the
{\hyperref[\detokenize{suffix::doc}]{\sphinxcrossref{\DUrole{doc}{suffix}}}} command in your input script.

\sphinxAtStartPar
See the {\hyperref[\detokenize{Speed_packages::doc}]{\sphinxcrossref{\DUrole{doc}{Accelerator packages}}}} page for more
instructions on how to use the accelerated styles effectively.


\bigskip\hrule\bigskip



\subsection{Mixing, shift, table, tail correction, restart, rRESPA info}
\label{\detokenize{pair_vashishta:mixing-shift-table-tail-correction-restart-rrespa-info}}
\sphinxAtStartPar
For atom type pairs I,J and I != J, where types I and J correspond to
two different element types, mixing is performed by LAMMPS as
described above from values in the potential file.

\sphinxAtStartPar
This pair style does not support the {\hyperref[\detokenize{pair_modify::doc}]{\sphinxcrossref{\DUrole{doc}{pair\_modify}}}}
shift, table, and tail options.

\sphinxAtStartPar
This pair style does not write its information to {\hyperref[\detokenize{restart::doc}]{\sphinxcrossref{\DUrole{doc}{binary restart files}}}}, since it is stored in potential files.  Thus, you
need to re\sphinxhyphen{}specify the pair\_style and pair\_coeff commands in an input
script that reads a restart file.

\sphinxAtStartPar
This pair style can only be used via the \sphinxstyleemphasis{pair} keyword of the
{\hyperref[\detokenize{run_style::doc}]{\sphinxcrossref{\DUrole{doc}{run\_style respa}}}} command.  It does not support the
\sphinxstyleemphasis{inner}, \sphinxstyleemphasis{middle}, \sphinxstyleemphasis{outer} keywords.


\bigskip\hrule\bigskip



\subsection{Restrictions}
\label{\detokenize{pair_vashishta:restrictions}}
\sphinxAtStartPar
These pair styles are part of the MANYBODY package.  They are only
enabled if LAMMPS was built with that package.  See the {\hyperref[\detokenize{Build_package::doc}]{\sphinxcrossref{\DUrole{doc}{Build package}}}} page for more info.

\sphinxAtStartPar
These pair styles requires the {\hyperref[\detokenize{newton::doc}]{\sphinxcrossref{\DUrole{doc}{newton}}}} setting to be “on”
for pair interactions.

\sphinxAtStartPar
The Vashishta potential files provided with LAMMPS (see the potentials
directory) are parameterized for metal {\hyperref[\detokenize{units::doc}]{\sphinxcrossref{\DUrole{doc}{units}}}}.  You can
use the Vashishta potential with any LAMMPS units, but you would need
to create your own potential file with coefficients listed in the
appropriate units if your simulation does not use “metal” units.


\subsection{Related commands}
\label{\detokenize{pair_vashishta:related-commands}}
\sphinxAtStartPar
{\hyperref[\detokenize{pair_coeff::doc}]{\sphinxcrossref{\DUrole{doc}{pair\_coeff}}}}


\subsection{Default}
\label{\detokenize{pair_vashishta:default}}
\sphinxAtStartPar
none


\bigskip\hrule\bigskip

\phantomsection\label{\detokenize{pair_vashishta:vashishta1990}}
\sphinxAtStartPar
\sphinxstylestrong{(Vashishta1990)} P. Vashishta, R. K. Kalia, J. P. Rino, Phys. Rev. B
41, 12197 (1990).

\phantomsection\label{\detokenize{pair_vashishta:vashishta2007}}
\sphinxAtStartPar
\sphinxstylestrong{(Vashishta2007)} P. Vashishta, R. K. Kalia, A. Nakano,
J. P. Rino. J. Appl. Phys. 101, 103515 (2007).

\phantomsection\label{\detokenize{pair_vashishta:branicio2009}}
\sphinxAtStartPar
\sphinxstylestrong{(Branicio2009)} Branicio, Rino, Gan and Tsuzuki, J. Phys Condensed
Matter 21 (2009) 095002

\sphinxstepscope

\index{pair\_style wf/cut@\spxentry{pair\_style wf/cut}}\ignorespaces 

\section{pair\_style wf/cut command}
\label{\detokenize{pair_wf_cut:pair-style-wf-cut-command}}\label{\detokenize{pair_wf_cut:index-0}}\label{\detokenize{pair_wf_cut::doc}}

\subsection{Syntax}
\label{\detokenize{pair_wf_cut:syntax}}
\begin{sphinxVerbatim}[commandchars=\\\{\}]
\PYG{k}{pair\PYGZus{}style}\PYG{+w}{ }\PYG{n}{wf}\PYG{o}{/}\PYG{n}{cut}\PYG{+w}{ }\PYG{n}{cutoff}
\end{sphinxVerbatim}
\begin{itemize}
\item {} 
\sphinxAtStartPar
cutoff = cutoff for wf interactions (distance units)

\end{itemize}


\subsection{Examples}
\label{\detokenize{pair_wf_cut:examples}}
\begin{sphinxVerbatim}[commandchars=\\\{\}]
\PYG{k}{pair\PYGZus{}style}\PYG{+w}{        }\PYG{n}{wf}\PYG{o}{/}\PYG{n}{cut}\PYG{+w}{ }\PYG{l+m}{2.0}
\PYG{k}{pair\PYGZus{}coeff}\PYG{+w}{        }\PYG{l+m}{1}\PYG{+w}{ }\PYG{l+m}{1}\PYG{+w}{ }\PYG{l+m}{1.0}\PYG{+w}{ }\PYG{l+m}{1.0}\PYG{+w}{ }\PYG{l+m}{1}\PYG{+w}{ }\PYG{l+m}{1}\PYG{+w}{  }\PYG{l+m}{2.0}
\end{sphinxVerbatim}


\subsection{Description}
\label{\detokenize{pair_wf_cut:description}}
\sphinxAtStartPar
The \sphinxstyleemphasis{wf/cut} (Wang\sphinxhyphen{}Frenkel) style computes LJ\sphinxhyphen{}like potentials as
described in {\hyperref[\detokenize{pair_wf_cut:wang2020}]{\sphinxcrossref{\DUrole{std,std-ref}{Wang2020}}}}. This potential is by
construction finite ranged and it vanishes quadratically at the cutoff
distance, avoiding truncation, shifting, interpolation and other typical
procedures with the LJ potential. The \sphinxstyleemphasis{wf/cut} can be used when a
typical short\sphinxhyphen{}ranged potential with attraction is required. The
potential is given by which is given by:
\begin{equation*}
\begin{split}\phi(r)= \epsilon \alpha \left(\left[{\sigma\over r}\right]^{2\mu} -1 \right)\left(\left[{r_c\over r}\right]^{2\mu}-1\right)^{2\nu}\end{split}
\end{equation*}
\sphinxAtStartPar
with
\begin{equation*}
\begin{split}\alpha=2\nu\left(\frac{r_c}{\sigma}\right)^{2\mu}\left[\frac{1+2\nu}{2\nu\left[(r_c/\sigma)^{2\mu}-1\right]}\right]^{2\nu+1}\end{split}
\end{equation*}
\sphinxAtStartPar
and
\begin{equation*}
\begin{split}r_{min}=r_c\left[\frac{1+2\nu}{1+2\nu(r_c/\sigma)^{2\nu}}\right]^{1/{2\nu}}\end{split}
\end{equation*}
\sphinxAtStartPar
\(r_c\) is the cutoff.

\sphinxAtStartPar
Comparison of the non\sphinxhyphen{}truncated Lennard\sphinxhyphen{}Jones 12\sphinxhyphen{}6 potential (red curve),
and the WF potentials with \(\mu=1\) and \(\nu=1\) are shown in
the figure below. The blue curve has \(r_c =2.0\) and the green
curve has \(r_c =1.2\) and can be used to describe colloidal
interactions.

\noindent{\hspace*{\fill}\sphinxincludegraphics[scale=0.33]{{WF_LJ}.jpg}\hspace*{\fill}}

\sphinxAtStartPar
The following coefficients must be defined for each pair of atoms
types via the {\hyperref[\detokenize{pair_coeff::doc}]{\sphinxcrossref{\DUrole{doc}{pair\_coeff}}}} command as in the example
above, or in the data file or restart files read by the
{\hyperref[\detokenize{read_data::doc}]{\sphinxcrossref{\DUrole{doc}{read\_data}}}} or {\hyperref[\detokenize{read_restart::doc}]{\sphinxcrossref{\DUrole{doc}{read\_restart}}}}
commands:
\begin{itemize}
\item {} 
\sphinxAtStartPar
\(\epsilon\) (energy units)

\item {} 
\sphinxAtStartPar
\(\sigma\) (distance units)

\item {} 
\sphinxAtStartPar
\(\nu\)

\item {} 
\sphinxAtStartPar
\(\mu\)

\item {} 
\sphinxAtStartPar
\(r_c\) (distance units)

\end{itemize}

\sphinxAtStartPar
The last coefficient is optional. If not specified, the global cutoff
given in the pair\_style command is used.  The exponents \(\nu\) and
\(\mu\) are positive integers, usually set to 1. There is usually
little to be gained by choosing other values of \(\nu\) and
\(\mu\) (See discussion in {\hyperref[\detokenize{pair_wf_cut:wang2020}]{\sphinxcrossref{\DUrole{std,std-ref}{Wang2020}}}})


\bigskip\hrule\bigskip


\sphinxAtStartPar
\sphinxstylestrong{Mixing, shift, table, tail correction, restart, rRESPA info}:

\sphinxAtStartPar
This pair style does not support the {\hyperref[\detokenize{pair_modify::doc}]{\sphinxcrossref{\DUrole{doc}{pair\_modify}}}}
mixing and table options.

\sphinxAtStartPar
The {\hyperref[\detokenize{pair_modify::doc}]{\sphinxcrossref{\DUrole{doc}{pair\_modify}}}} tail and shift options are not
relevant for this pair style as it goes to zero at the cut\sphinxhyphen{}off radius.

\sphinxAtStartPar
This pair style writes its information to {\hyperref[\detokenize{restart::doc}]{\sphinxcrossref{\DUrole{doc}{binary restart files}}}}, so pair\_style and pair\_coeff commands do not need to be
specified in an input script that reads a restart file.

\sphinxAtStartPar
This pair style does not support the use of the \sphinxstyleemphasis{inner}, \sphinxstyleemphasis{middle},
and \sphinxstyleemphasis{outer} keywords of the {\hyperref[\detokenize{run_style::doc}]{\sphinxcrossref{\DUrole{doc}{run\_style respa}}}} command.


\bigskip\hrule\bigskip



\subsection{Restrictions}
\label{\detokenize{pair_wf_cut:restrictions}}
\sphinxAtStartPar
This pair style can only be used if LAMMPS was built with the
EXTRA\sphinxhyphen{}PAIR package.  See the {\hyperref[\detokenize{Build_package::doc}]{\sphinxcrossref{\DUrole{doc}{Build package}}}} doc
page for more info.


\subsection{Related commands}
\label{\detokenize{pair_wf_cut:related-commands}}
\sphinxAtStartPar
{\hyperref[\detokenize{pair_coeff::doc}]{\sphinxcrossref{\DUrole{doc}{pair\_coeff}}}}

\sphinxAtStartPar
\sphinxstylestrong{Default:} none


\bigskip\hrule\bigskip

\phantomsection\label{\detokenize{pair_wf_cut:wang2020}}
\sphinxAtStartPar
\sphinxstylestrong{(Wang2020)} X. Wang, S. Ramirez\sphinxhyphen{}Hinestrosa, J. Dobnikar, and D. Frenkel, Phys. Chem. Chem. Phys. 22, 10624 (2020).

\sphinxstepscope

\index{pair\_style ylz@\spxentry{pair\_style ylz}}\ignorespaces 

\section{pair\_style ylz command}
\label{\detokenize{pair_ylz:pair-style-ylz-command}}\label{\detokenize{pair_ylz:index-0}}\label{\detokenize{pair_ylz::doc}}

\subsection{Syntax}
\label{\detokenize{pair_ylz:syntax}}
\begin{sphinxVerbatim}[commandchars=\\\{\}]
\PYG{k}{pair\PYGZus{}style}\PYG{+w}{ }\PYG{n}{ylz}\PYG{+w}{ }\PYG{n}{cutoff}
\end{sphinxVerbatim}
\begin{itemize}
\item {} 
\sphinxAtStartPar
cutoff = global cutoff for interactions (distance units)

\end{itemize}


\subsection{Examples}
\label{\detokenize{pair_ylz:examples}}
\begin{sphinxVerbatim}[commandchars=\\\{\}]
\PYG{k}{pair\PYGZus{}style}\PYG{+w}{   }\PYG{n}{ylz}\PYG{+w}{  }\PYG{l+m}{2.6}
\PYG{k}{pair\PYGZus{}coeff}\PYG{+w}{   }\PYG{o}{*}\PYG{+w}{  }\PYG{o}{*}\PYG{+w}{  }\PYG{l+m}{1.0}\PYG{+w}{  }\PYG{l+m}{1.0}\PYG{+w}{  }\PYG{l+m}{4}\PYG{+w}{  }\PYG{l+m}{3}\PYG{+w}{  }\PYG{l+m}{0.0}\PYG{+w}{  }\PYG{l+m}{2.6}
\end{sphinxVerbatim}


\subsection{Description}
\label{\detokenize{pair_ylz:description}}
\sphinxAtStartPar
\DUrole{versionmodified,added}{Added in version 3Nov2022.}

\sphinxAtStartPar
The \sphinxstyleemphasis{ylz} (Yuan\sphinxhyphen{}Li\sphinxhyphen{}Zhang) style computes an anisotropic interaction
between pairs of coarse\sphinxhyphen{}grained particles considering the relative
particle orientations. This potential was originally developed as a
particle\sphinxhyphen{}based solvent\sphinxhyphen{}free model for biological membranes
{\hyperref[\detokenize{pair_ylz:yuan}]{\sphinxcrossref{\DUrole{std,std-ref}{(Yuan2010a)}}}}.  Unlike {\hyperref[\detokenize{pair_gayberne::doc}]{\sphinxcrossref{\DUrole{doc}{pair\_style gayberne}}}}, whose orientation dependence is strictly derived from
the closest distance between two ellipsoidal rigid bodies, the
orientation\sphinxhyphen{}dependence of this pair style is mathematically defined such
that the particles can self\sphinxhyphen{}assemble into one\sphinxhyphen{}particle\sphinxhyphen{}thick fluid
membranes.  The potential of this pair style is described by:
\begin{equation*}
\begin{split}U ( \mathbf{r}_{ij}, \mathbf{n}_i, \mathbf{n}_j ) =\left\{\begin{matrix} {u}_R(r)+\left [ 1-\phi (\mathbf{\hat{r}}_{ij}, \mathbf{n}_i, \mathbf{n}_j ) \right ]\epsilon, ~~ r<{r}_{min} \\ {u}_A(r)\phi (\mathbf{\hat{r}}_{ij}, \mathbf{n}_i, \mathbf{n}_j ),~~  {r}_{min}<r<{r}_{c} \\ \end{matrix}\right.\\\\ \phi (\mathbf{\hat{r}}_{ij}, \mathbf{n}_i, \mathbf{n}_j )=1+\left [  \mu (a(\mathbf{\hat{r}}_{ij}, \mathbf{n}_i, \mathbf{n}_j )-1) \right ] \\\\a(\mathbf{\hat{r}}_{ij}, \mathbf{n}_i, \mathbf{n}_j )=(\mathbf{n}_i\times\mathbf{\hat{r}}_{ij} )\cdot (\mathbf{n}_j\times\mathbf{\hat{r}}_{ij} )+{\beta}(\mathbf{n}_i-\mathbf{n}_j)\cdot \mathbf{\hat{r}}_{ij}-\beta^{2}\\\\  {u}_R(r)=\epsilon \left [ \left ( \frac{{r}_{min}}{r} \right )^{4}-2\left ( \frac{{r}_{min}}{r}\right )^{2} \right ] \\\\ {u}_A(r)=-\epsilon\;cos^{2\zeta }\left [ \frac{\pi}{2}\frac{\left ( {r}-{r}_{min} \right )}{\left ( {r}_{c}-{r}_{min} \right )} \right ]\\\end{split}
\end{equation*}
\sphinxAtStartPar
where \(\mathbf{r}_{i}\) and \(\mathbf{r}_{j}\) are the center
position vectors of particles i and j, respectively,
\(\mathbf{r}_{ij}=\mathbf{r}_{i}-\mathbf{r}_{j}\) is the
inter\sphinxhyphen{}particle distance vector, \(r=\left|\mathbf{r}_{ij} \right|\)
and \({\hat{\mathbf{r}}}_{ij}=\mathbf{r}_{ij}/r\).  The unit vectors
\(\mathbf{n}_{i}\) and \(\mathbf{n}_{j}\) represent the axes of
symmetry of particles i and j, respectively, \(u_R\) and \(u_A\)
are the repulsive and attractive potentials, \(\phi\) is an angular
function which depends on the relative orientation between pair
particles, \(\mu\) is the parameter related to the bending rigidity
of the membrane, \(\beta\) is the parameter related to the
spontaneous curvature, and \(\epsilon\) is the energy unit,
respectively.   The \(\zeta\) controls the slope of the attractive
branch and hence the diffusivity of the particles in the in\sphinxhyphen{}plane
direction of the membrane.  \({r}_{c}\) is the cutoff radius,
\(r_{min}\) is the distance which minimizes the potential energy
\(u_{A}(r)\) and \(r_{min}=2^{1/6}\sigma\), where \(\sigma\)
is the length unit.

\sphinxAtStartPar
This pair style is suited for solvent\sphinxhyphen{}free coarse\sphinxhyphen{}grained simulations of
biological systems involving lipid bilayer membranes, such as vesicle
shape transformations {\hyperref[\detokenize{pair_ylz:yuan}]{\sphinxcrossref{\DUrole{std,std-ref}{(Yuan2010b)}}}}, nanoparticle
endocytosis {\hyperref[\detokenize{pair_ylz:huang}]{\sphinxcrossref{\DUrole{std,std-ref}{(Huang)}}}}, modeling of red blood cell membranes
{\hyperref[\detokenize{pair_ylz:fu}]{\sphinxcrossref{\DUrole{std,std-ref}{(Fu)}}}}, {\hyperref[\detokenize{pair_ylz:appshaw}]{\sphinxcrossref{\DUrole{std,std-ref}{(Appshaw)}}}}, and modeling of cell
elasticity {\hyperref[\detokenize{pair_ylz:becton}]{\sphinxcrossref{\DUrole{std,std-ref}{(Becton)}}}}.

\sphinxAtStartPar
Use of this pair style requires the NVE, NVT, or NPT fixes with the
\sphinxstyleemphasis{asphere} extension (e.g. {\hyperref[\detokenize{fix_nve_asphere::doc}]{\sphinxcrossref{\DUrole{doc}{fix nve/asphere}}}}) in
order to integrate particle rotation.  Additionally, {\hyperref[\detokenize{atom_style::doc}]{\sphinxcrossref{\DUrole{doc}{atom\_style
ellipsoid}}}} should be used since it defines the rotational
state of each particle.

\sphinxAtStartPar
The following coefficients must be defined for each pair of atoms types
via the {\hyperref[\detokenize{pair_coeff::doc}]{\sphinxcrossref{\DUrole{doc}{pair\_coeff}}}} command as in the examples above,
or in the data file or restart files read by the {\hyperref[\detokenize{read_data::doc}]{\sphinxcrossref{\DUrole{doc}{read\_data}}}} or {\hyperref[\detokenize{read_restart::doc}]{\sphinxcrossref{\DUrole{doc}{read\_restart}}}} commands, or by
mixing as described below:
\begin{itemize}
\item {} 
\sphinxAtStartPar
\(\epsilon\) = well depth (energy units)

\item {} 
\sphinxAtStartPar
\(\sigma\) = minimum effective particle radii (distance units)

\item {} 
\sphinxAtStartPar
\(\zeta\) = tuning parameter for the slope of the attractive branch

\item {} 
\sphinxAtStartPar
\(\mu\) = parameter related to bending rigidity

\item {} 
\sphinxAtStartPar
\(\beta\) = parameter related to the spontaneous curvature

\item {} 
\sphinxAtStartPar
cutoff (distance units)

\end{itemize}

\sphinxAtStartPar
The last coefficient is optional.  If not specified, the global
cutoff specified in the pair\_style command is used.


\bigskip\hrule\bigskip



\subsection{Mixing, shift, table, tail correction, restart, rRESPA info}
\label{\detokenize{pair_ylz:mixing-shift-table-tail-correction-restart-rrespa-info}}
\sphinxAtStartPar
For atom type pairs I,J and I != J, the epsilon and sigma coefficients
and cutoff distance for this pair style can be mixed.  The default mix
value is \sphinxstyleemphasis{geometric}.  See the “pair\_modify” command for details.

\sphinxAtStartPar
The {\hyperref[\detokenize{pair_modify::doc}]{\sphinxcrossref{\DUrole{doc}{pair\_modify}}}} table option is not relevant for
this pair style.

\sphinxAtStartPar
This pair style does not support the {\hyperref[\detokenize{pair_modify::doc}]{\sphinxcrossref{\DUrole{doc}{pair\_modify}}}}
tail option for adding long\sphinxhyphen{}range tail corrections to energy and
pressure.

\sphinxAtStartPar
This pair style writes its information to {\hyperref[\detokenize{restart::doc}]{\sphinxcrossref{\DUrole{doc}{binary restart files}}}}, so pair\_style and pair\_coeff commands do not need to be
specified in an input script that reads a restart file.

\sphinxAtStartPar
This pair style can only be used via the \sphinxstyleemphasis{pair} keyword of the
{\hyperref[\detokenize{run_style::doc}]{\sphinxcrossref{\DUrole{doc}{run\_style respa}}}} command.  It does not support the
\sphinxstyleemphasis{inner}, \sphinxstyleemphasis{middle}, \sphinxstyleemphasis{outer} keywords.


\bigskip\hrule\bigskip



\subsection{Restrictions}
\label{\detokenize{pair_ylz:restrictions}}
\sphinxAtStartPar
The \sphinxstyleemphasis{ylz} style is part of the ASPHERE package.  It is only enabled if
LAMMPS was built with that package.  See the {\hyperref[\detokenize{Build_package::doc}]{\sphinxcrossref{\DUrole{doc}{Build package}}}} page for more info.

\sphinxAtStartPar
This pair style requires that atoms store torque and a quaternion to
represent their orientation, as defined by the {\hyperref[\detokenize{atom_style::doc}]{\sphinxcrossref{\DUrole{doc}{atom\_style}}}}.  It also requires they store a per\sphinxhyphen{}atom {\hyperref[\detokenize{set::doc}]{\sphinxcrossref{\DUrole{doc}{shape}}}}.  The particles cannot store a per\sphinxhyphen{}particle diameter.  To avoid
being mistakenly considered as point particles, the shape parameters ought
to be non\sphinxhyphen{}spherical, like {[}1 0.99 0.99{]}.  Unlike the {\hyperref[\detokenize{pair_resquared::doc}]{\sphinxcrossref{\DUrole{doc}{resquared}}}} pair style for which the shape directly determines the
mathematical expressions of the potential, the shape parameters for this
pair style is only involved in the computation of the moment of inertia
and thus only influences the rotational dynamics of individual
particles.

\sphinxAtStartPar
This pair style requires that \sphinxstylestrong{all} atoms are ellipsoids as defined by
the {\hyperref[\detokenize{atom_style::doc}]{\sphinxcrossref{\DUrole{doc}{atom\_style ellipsoid}}}} command.


\subsection{Related commands}
\label{\detokenize{pair_ylz:related-commands}}
\sphinxAtStartPar
{\hyperref[\detokenize{pair_coeff::doc}]{\sphinxcrossref{\DUrole{doc}{pair\_coeff}}}},
{\hyperref[\detokenize{fix_nve_asphere::doc}]{\sphinxcrossref{\DUrole{doc}{fix nve/asphere}}}},
{\hyperref[\detokenize{compute_temp_asphere::doc}]{\sphinxcrossref{\DUrole{doc}{compute temp/asphere}}}},
{\hyperref[\detokenize{pair_resquared::doc}]{\sphinxcrossref{\DUrole{doc}{pair\_style resquared}}}},
{\hyperref[\detokenize{pair_gayberne::doc}]{\sphinxcrossref{\DUrole{doc}{pair\_style gayberne}}}}


\subsection{Default}
\label{\detokenize{pair_ylz:default}}
\sphinxAtStartPar
none


\bigskip\hrule\bigskip

\phantomsection\label{\detokenize{pair_ylz:yuan}}
\sphinxAtStartPar
\sphinxstylestrong{(Yuan2010a)} Yuan, Huang, Li, Lykotrafitis, Zhang, Phys. Rev. E, 82, 011905(2010).

\sphinxAtStartPar
\sphinxstylestrong{(Yuan2010b)} Yuan, Huang, Zhang, Soft. Matter, 6, 4571(2010).

\phantomsection\label{\detokenize{pair_ylz:huang}}
\sphinxAtStartPar
\sphinxstylestrong{(Huang)} Huang, Zhang, Yuan, Gao, Zhang, Nano Lett. 13, 4546(2013).

\phantomsection\label{\detokenize{pair_ylz:fu}}
\sphinxAtStartPar
\sphinxstylestrong{(Fu)} Fu, Peng, Yuan, Kfoury, Young, Comput. Phys. Commun, 210, 193\sphinxhyphen{}203(2017).

\phantomsection\label{\detokenize{pair_ylz:appshaw}}
\sphinxAtStartPar
\sphinxstylestrong{(Appshaw)} Appshaw, Seddon, Hanna, Soft. Matter,18, 1747(2022).

\phantomsection\label{\detokenize{pair_ylz:becton}}
\sphinxAtStartPar
\sphinxstylestrong{(Becton)} Becton, Averett, Wang, Biomech. Model. Mechanobiology, 18, 425\sphinxhyphen{}433(2019).

\sphinxstepscope

\index{pair\_style yukawa@\spxentry{pair\_style yukawa}}\ignorespaces 
\index{pair\_style yukawa/gpu@\spxentry{pair\_style yukawa/gpu}}\ignorespaces 
\index{pair\_style yukawa/omp@\spxentry{pair\_style yukawa/omp}}\ignorespaces 
\index{pair\_style yukawa/kk@\spxentry{pair\_style yukawa/kk}}\ignorespaces 

\section{pair\_style yukawa command}
\label{\detokenize{pair_yukawa:pair-style-yukawa-command}}\label{\detokenize{pair_yukawa:index-3}}\label{\detokenize{pair_yukawa:index-2}}\label{\detokenize{pair_yukawa:index-1}}\label{\detokenize{pair_yukawa:index-0}}\label{\detokenize{pair_yukawa::doc}}
\sphinxAtStartPar
Accelerator Variants: \sphinxstyleemphasis{yukawa/gpu}, \sphinxstyleemphasis{yukawa/omp}, \sphinxstyleemphasis{yukawa/kk}


\subsection{Syntax}
\label{\detokenize{pair_yukawa:syntax}}
\begin{sphinxVerbatim}[commandchars=\\\{\}]
\PYG{k}{pair\PYGZus{}style}\PYG{+w}{ }\PYG{n}{yukawa}\PYG{+w}{ }\PYG{n}{kappa}\PYG{+w}{ }\PYG{n}{cutoff}
\end{sphinxVerbatim}
\begin{itemize}
\item {} 
\sphinxAtStartPar
kappa = screening length (inverse distance units)

\item {} 
\sphinxAtStartPar
cutoff = global cutoff for Yukawa interactions (distance units)

\end{itemize}


\subsection{Examples}
\label{\detokenize{pair_yukawa:examples}}
\begin{sphinxVerbatim}[commandchars=\\\{\}]
\PYG{k}{pair\PYGZus{}style}\PYG{+w}{ }\PYG{n}{yukawa}\PYG{+w}{ }\PYG{l+m}{2.0}\PYG{+w}{ }\PYG{l+m}{2.5}
\PYG{k}{pair\PYGZus{}coeff}\PYG{+w}{ }\PYG{l+m}{1}\PYG{+w}{ }\PYG{l+m}{1}\PYG{+w}{ }\PYG{l+m}{100.0}\PYG{+w}{ }\PYG{l+m}{2.3}
\PYG{k}{pair\PYGZus{}coeff}\PYG{+w}{ }\PYG{o}{*}\PYG{+w}{ }\PYG{o}{*}\PYG{+w}{ }\PYG{l+m}{100.0}
\end{sphinxVerbatim}


\subsection{Description}
\label{\detokenize{pair_yukawa:description}}
\sphinxAtStartPar
Style \sphinxstyleemphasis{yukawa} computes pairwise interactions with the formula
\begin{equation*}
\begin{split}E = A \frac{e^{- \kappa r}}{r} \qquad r < r_c\end{split}
\end{equation*}
\sphinxAtStartPar
\(r_c\) is the cutoff.

\sphinxAtStartPar
The following coefficients must be defined for each pair of atoms
types via the {\hyperref[\detokenize{pair_coeff::doc}]{\sphinxcrossref{\DUrole{doc}{pair\_coeff}}}} command as in the examples
above, or in the data file or restart files read by the
{\hyperref[\detokenize{read_data::doc}]{\sphinxcrossref{\DUrole{doc}{read\_data}}}} or {\hyperref[\detokenize{read_restart::doc}]{\sphinxcrossref{\DUrole{doc}{read\_restart}}}}
commands, or by mixing as described below:
\begin{itemize}
\item {} 
\sphinxAtStartPar
A (energy*distance units)

\item {} 
\sphinxAtStartPar
cutoff (distance units)

\end{itemize}

\sphinxAtStartPar
The last coefficient is optional.  If not specified, the global yukawa
cutoff is used.


\bigskip\hrule\bigskip


\sphinxAtStartPar
Styles with a \sphinxstyleemphasis{gpu}, \sphinxstyleemphasis{intel}, \sphinxstyleemphasis{kk}, \sphinxstyleemphasis{omp}, or \sphinxstyleemphasis{opt} suffix are
functionally the same as the corresponding style without the suffix.
They have been optimized to run faster, depending on your available
hardware, as discussed on the {\hyperref[\detokenize{Speed_packages::doc}]{\sphinxcrossref{\DUrole{doc}{Accelerator packages}}}}
page.  The accelerated styles take the same arguments and should
produce the same results, except for round\sphinxhyphen{}off and precision issues.

\sphinxAtStartPar
These accelerated styles are part of the GPU, INTEL, KOKKOS,
OPENMP, and OPT packages, respectively.  They are only enabled if
LAMMPS was built with those packages.  See the {\hyperref[\detokenize{Build_package::doc}]{\sphinxcrossref{\DUrole{doc}{Build package}}}} page for more info.

\sphinxAtStartPar
You can specify the accelerated styles explicitly in your input script
by including their suffix, or you can use the {\hyperref[\detokenize{Run_options::doc}]{\sphinxcrossref{\DUrole{doc}{\sphinxhyphen{}suffix command\sphinxhyphen{}line switch}}}} when you invoke LAMMPS, or you can use the
{\hyperref[\detokenize{suffix::doc}]{\sphinxcrossref{\DUrole{doc}{suffix}}}} command in your input script.

\sphinxAtStartPar
See the {\hyperref[\detokenize{Speed_packages::doc}]{\sphinxcrossref{\DUrole{doc}{Accelerator packages}}}} page for more
instructions on how to use the accelerated styles effectively.


\bigskip\hrule\bigskip



\subsection{Mixing, shift, table, tail correction, restart, rRESPA info}
\label{\detokenize{pair_yukawa:mixing-shift-table-tail-correction-restart-rrespa-info}}
\sphinxAtStartPar
For atom type pairs I,J and I != J, the A coefficient and cutoff
distance for this pair style can be mixed.  A is an energy value mixed
like a LJ epsilon.  The default mix value is \sphinxstyleemphasis{geometric}.  See the
“pair\_modify” command for details.

\sphinxAtStartPar
This pair style supports the {\hyperref[\detokenize{pair_modify::doc}]{\sphinxcrossref{\DUrole{doc}{pair\_modify}}}} shift
option for the energy of the pair interaction.

\sphinxAtStartPar
The {\hyperref[\detokenize{pair_modify::doc}]{\sphinxcrossref{\DUrole{doc}{pair\_modify}}}} table option is not relevant
for this pair style.

\sphinxAtStartPar
This pair style does not support the {\hyperref[\detokenize{pair_modify::doc}]{\sphinxcrossref{\DUrole{doc}{pair\_modify}}}}
tail option for adding long\sphinxhyphen{}range tail corrections to energy and
pressure.

\sphinxAtStartPar
This pair style writes its information to {\hyperref[\detokenize{restart::doc}]{\sphinxcrossref{\DUrole{doc}{binary restart files}}}}, so pair\_style and pair\_coeff commands do not need
to be specified in an input script that reads a restart file.

\sphinxAtStartPar
This pair style can only be used via the \sphinxstyleemphasis{pair} keyword of the
{\hyperref[\detokenize{run_style::doc}]{\sphinxcrossref{\DUrole{doc}{run\_style respa}}}} command.  It does not support the
\sphinxstyleemphasis{inner}, \sphinxstyleemphasis{middle}, \sphinxstyleemphasis{outer} keywords.


\bigskip\hrule\bigskip



\subsection{Restrictions}
\label{\detokenize{pair_yukawa:restrictions}}\begin{quote}

\sphinxAtStartPar
none
\end{quote}


\subsection{Related commands}
\label{\detokenize{pair_yukawa:related-commands}}
\sphinxAtStartPar
{\hyperref[\detokenize{pair_coeff::doc}]{\sphinxcrossref{\DUrole{doc}{pair\_coeff}}}}


\subsection{Default}
\label{\detokenize{pair_yukawa:default}}
\sphinxAtStartPar
none

\sphinxstepscope

\index{pair\_style yukawa/colloid@\spxentry{pair\_style yukawa/colloid}}\ignorespaces 
\index{pair\_style yukawa/colloid/gpu@\spxentry{pair\_style yukawa/colloid/gpu}}\ignorespaces 
\index{pair\_style yukawa/colloid/kk@\spxentry{pair\_style yukawa/colloid/kk}}\ignorespaces 
\index{pair\_style yukawa/colloid/omp@\spxentry{pair\_style yukawa/colloid/omp}}\ignorespaces 

\section{pair\_style yukawa/colloid command}
\label{\detokenize{pair_yukawa_colloid:pair-style-yukawa-colloid-command}}\label{\detokenize{pair_yukawa_colloid:index-3}}\label{\detokenize{pair_yukawa_colloid:index-2}}\label{\detokenize{pair_yukawa_colloid:index-1}}\label{\detokenize{pair_yukawa_colloid:index-0}}\label{\detokenize{pair_yukawa_colloid::doc}}
\sphinxAtStartPar
Accelerator Variants: \sphinxstyleemphasis{yukawa/colloid/gpu}, \sphinxstyleemphasis{yukawa/colloid/kk}, \sphinxstyleemphasis{yukawa/colloid/omp}


\subsection{Syntax}
\label{\detokenize{pair_yukawa_colloid:syntax}}
\begin{sphinxVerbatim}[commandchars=\\\{\}]
\PYG{k}{pair\PYGZus{}style}\PYG{+w}{ }\PYG{n}{yukawa}\PYG{o}{/}\PYG{n}{colloid}\PYG{+w}{ }\PYG{n}{kappa}\PYG{+w}{ }\PYG{n}{cutoff}
\end{sphinxVerbatim}
\begin{itemize}
\item {} 
\sphinxAtStartPar
kappa = screening length (inverse distance units)

\item {} 
\sphinxAtStartPar
cutoff = global cutoff for colloidal Yukawa interactions (distance units)

\end{itemize}


\subsection{Examples}
\label{\detokenize{pair_yukawa_colloid:examples}}
\begin{sphinxVerbatim}[commandchars=\\\{\}]
\PYG{k}{pair\PYGZus{}style}\PYG{+w}{ }\PYG{n}{yukawa}\PYG{o}{/}\PYG{n}{colloid}\PYG{+w}{ }\PYG{l+m}{2.0}\PYG{+w}{ }\PYG{l+m}{2.5}
\PYG{k}{pair\PYGZus{}coeff}\PYG{+w}{ }\PYG{l+m}{1}\PYG{+w}{ }\PYG{l+m}{1}\PYG{+w}{ }\PYG{l+m}{100.0}\PYG{+w}{ }\PYG{l+m}{2.3}
\PYG{k}{pair\PYGZus{}coeff}\PYG{+w}{ }\PYG{o}{*}\PYG{+w}{ }\PYG{o}{*}\PYG{+w}{ }\PYG{l+m}{100.0}
\end{sphinxVerbatim}


\subsection{Description}
\label{\detokenize{pair_yukawa_colloid:description}}
\sphinxAtStartPar
Style \sphinxstyleemphasis{yukawa/colloid} computes pairwise interactions with the formula
\begin{equation*}
\begin{split}E = \frac{A}{\kappa} e^{- \kappa (r - (r_i + r_j))} \qquad r < r_c\end{split}
\end{equation*}
\sphinxAtStartPar
where \(r_i\) and \(r_j\) are the radii of the two particles
and \(r_c\) is the cutoff.

\sphinxAtStartPar
In contrast to {\hyperref[\detokenize{pair_yukawa::doc}]{\sphinxcrossref{\DUrole{doc}{pair\_style yukawa}}}}, this functional
form arises from the Coulombic interaction between two colloid
particles, screened due to the presence of an electrolyte, see the
book by {\hyperref[\detokenize{pair_yukawa_colloid:safran}]{\sphinxcrossref{\DUrole{std,std-ref}{Safran}}}} for a derivation in the context of DLVO
theory.  {\hyperref[\detokenize{pair_yukawa::doc}]{\sphinxcrossref{\DUrole{doc}{Pair\_style yukawa}}}} is a screened Coulombic
potential between two point\sphinxhyphen{}charges and uses no such approximation.

\sphinxAtStartPar
This potential applies to nearby particle pairs for which the Derjagin
approximation holds, meaning \(h << r_i + r_j\), where \sphinxstyleemphasis{h} is the
surface\sphinxhyphen{}to\sphinxhyphen{}surface separation of the two particles.

\sphinxAtStartPar
When used in combination with {\hyperref[\detokenize{pair_colloid::doc}]{\sphinxcrossref{\DUrole{doc}{pair\_style colloid}}}},
the two terms become the so\sphinxhyphen{}called DLVO potential, which combines
electrostatic repulsion and van der Waals attraction.

\sphinxAtStartPar
The following coefficients must be defined for each pair of atoms
types via the {\hyperref[\detokenize{pair_coeff::doc}]{\sphinxcrossref{\DUrole{doc}{pair\_coeff}}}} command as in the examples
above, or in the data file or restart files read by the
{\hyperref[\detokenize{read_data::doc}]{\sphinxcrossref{\DUrole{doc}{read\_data}}}} or {\hyperref[\detokenize{read_restart::doc}]{\sphinxcrossref{\DUrole{doc}{read\_restart}}}}
commands, or by mixing as described below:
\begin{itemize}
\item {} 
\sphinxAtStartPar
A (energy/distance units)

\item {} 
\sphinxAtStartPar
cutoff (distance units)

\end{itemize}

\sphinxAtStartPar
The prefactor A is determined from the relationship between surface
charge and surface potential due to the presence of electrolyte.  Note
that the A for this potential style has different units than the A
used in {\hyperref[\detokenize{pair_yukawa::doc}]{\sphinxcrossref{\DUrole{doc}{pair\_style yukawa}}}}.  For low surface
potentials, i.e. less than about 25 mV, A can be written as:
\begin{equation*}
\begin{split}A = 2 \pi R\varepsilon\varepsilon_0 \kappa \psi^2\end{split}
\end{equation*}
\sphinxAtStartPar
where
\begin{itemize}
\item {} 
\sphinxAtStartPar
\sphinxstyleemphasis{R} = colloid radius (distance units)

\item {} 
\sphinxAtStartPar
\(\varepsilon_0\) = permittivity of free space (charge\textasciicircum{}2/energy/distance units)

\item {} 
\sphinxAtStartPar
\(\varepsilon\) = relative permittivity of fluid medium (dimensionless)

\item {} 
\sphinxAtStartPar
\(\kappa\) = inverse screening length (1/distance units)

\item {} 
\sphinxAtStartPar
\(\psi\) = surface potential (energy/charge units)

\end{itemize}

\sphinxAtStartPar
The last coefficient is optional.  If not specified, the global
yukawa/colloid cutoff is used.


\bigskip\hrule\bigskip


\sphinxAtStartPar
Styles with a \sphinxstyleemphasis{gpu}, \sphinxstyleemphasis{intel}, \sphinxstyleemphasis{kk}, \sphinxstyleemphasis{omp}, or \sphinxstyleemphasis{opt} suffix are
functionally the same as the corresponding style without the suffix.
They have been optimized to run faster, depending on your available
hardware, as discussed on the {\hyperref[\detokenize{Speed_packages::doc}]{\sphinxcrossref{\DUrole{doc}{Accelerator packages}}}}
page.  The accelerated styles take the same arguments and should
produce the same results, except for round\sphinxhyphen{}off and precision issues.

\sphinxAtStartPar
These accelerated styles are part of the GPU, INTEL, KOKKOS,
OPENMP, and OPT packages, respectively.  They are only enabled if
LAMMPS was built with those packages.  See the {\hyperref[\detokenize{Build_package::doc}]{\sphinxcrossref{\DUrole{doc}{Build package}}}} page for more info.

\sphinxAtStartPar
You can specify the accelerated styles explicitly in your input script
by including their suffix, or you can use the {\hyperref[\detokenize{Run_options::doc}]{\sphinxcrossref{\DUrole{doc}{\sphinxhyphen{}suffix command\sphinxhyphen{}line switch}}}} when you invoke LAMMPS, or you can use the
{\hyperref[\detokenize{suffix::doc}]{\sphinxcrossref{\DUrole{doc}{suffix}}}} command in your input script.

\sphinxAtStartPar
See the {\hyperref[\detokenize{Speed_packages::doc}]{\sphinxcrossref{\DUrole{doc}{Accelerator packages}}}} page for more
instructions on how to use the accelerated styles effectively.


\bigskip\hrule\bigskip



\subsection{Mixing, shift, table, tail correction, restart, rRESPA info}
\label{\detokenize{pair_yukawa_colloid:mixing-shift-table-tail-correction-restart-rrespa-info}}
\sphinxAtStartPar
For atom type pairs I,J and I != J, the A coefficient and cutoff
distance for this pair style can be mixed.  A is an energy value mixed
like a LJ epsilon.  The default mix value is \sphinxstyleemphasis{geometric}.  See the
“pair\_modify” command for details.

\sphinxAtStartPar
This pair style supports the {\hyperref[\detokenize{pair_modify::doc}]{\sphinxcrossref{\DUrole{doc}{pair\_modify}}}} shift
option for the energy of the pair interaction.

\sphinxAtStartPar
The {\hyperref[\detokenize{pair_modify::doc}]{\sphinxcrossref{\DUrole{doc}{pair\_modify}}}} table option is not relevant
for this pair style.

\sphinxAtStartPar
This pair style does not support the {\hyperref[\detokenize{pair_modify::doc}]{\sphinxcrossref{\DUrole{doc}{pair\_modify}}}}
tail option for adding long\sphinxhyphen{}range tail corrections to energy and
pressure.

\sphinxAtStartPar
This pair style writes its information to {\hyperref[\detokenize{restart::doc}]{\sphinxcrossref{\DUrole{doc}{binary restart files}}}}, so pair\_style and pair\_coeff commands do not need
to be specified in an input script that reads a restart file.

\sphinxAtStartPar
This pair style can only be used via the \sphinxstyleemphasis{pair} keyword of the
{\hyperref[\detokenize{run_style::doc}]{\sphinxcrossref{\DUrole{doc}{run\_style respa}}}} command.  It does not support the
\sphinxstyleemphasis{inner}, \sphinxstyleemphasis{middle}, \sphinxstyleemphasis{outer} keywords.


\bigskip\hrule\bigskip



\subsection{Restrictions}
\label{\detokenize{pair_yukawa_colloid:restrictions}}
\sphinxAtStartPar
This style is part of the COLLOID package.  It is only enabled if
LAMMPS was built with that package.  See the {\hyperref[\detokenize{Build_package::doc}]{\sphinxcrossref{\DUrole{doc}{Build package}}}} page for more info.

\sphinxAtStartPar
This pair style requires that atoms be finite\sphinxhyphen{}size spheres with a
diameter, as defined by the {\hyperref[\detokenize{atom_style::doc}]{\sphinxcrossref{\DUrole{doc}{atom\_style sphere}}}}
command.

\sphinxAtStartPar
Per\sphinxhyphen{}particle polydispersity is not yet supported by this pair style;
per\sphinxhyphen{}type polydispersity is allowed.  This means all particles of the
same type must have the same diameter.  Each type can have a different
diameter.


\bigskip\hrule\bigskip


\sphinxAtStartPar
Styles with a \sphinxstyleemphasis{gpu}, \sphinxstyleemphasis{intel}, \sphinxstyleemphasis{kk}, \sphinxstyleemphasis{omp}, or \sphinxstyleemphasis{opt} suffix are
functionally the same as the corresponding style without the suffix.
They have been optimized to run faster, depending on your available
hardware, as discussed on the {\hyperref[\detokenize{Speed_packages::doc}]{\sphinxcrossref{\DUrole{doc}{Accelerator packages}}}}
page.  The accelerated styles take the same arguments and should
produce the same results, except for round\sphinxhyphen{}off and precision issues.

\sphinxAtStartPar
These accelerated styles are part of the GPU, INTEL, KOKKOS,
OPENMP, and OPT packages, respectively.  They are only enabled if
LAMMPS was built with those packages.  See the {\hyperref[\detokenize{Build_package::doc}]{\sphinxcrossref{\DUrole{doc}{Build package}}}} page for more info.

\sphinxAtStartPar
You can specify the accelerated styles explicitly in your input script
by including their suffix, or you can use the {\hyperref[\detokenize{Run_options::doc}]{\sphinxcrossref{\DUrole{doc}{\sphinxhyphen{}suffix command\sphinxhyphen{}line switch}}}} when you invoke LAMMPS, or you can use the
{\hyperref[\detokenize{suffix::doc}]{\sphinxcrossref{\DUrole{doc}{suffix}}}} command in your input script.

\sphinxAtStartPar
See the {\hyperref[\detokenize{Speed_packages::doc}]{\sphinxcrossref{\DUrole{doc}{Accelerator packages}}}} page for more
instructions on how to use the accelerated styles effectively.


\bigskip\hrule\bigskip



\subsection{Related commands}
\label{\detokenize{pair_yukawa_colloid:related-commands}}
\sphinxAtStartPar
{\hyperref[\detokenize{pair_coeff::doc}]{\sphinxcrossref{\DUrole{doc}{pair\_coeff}}}}


\subsection{Default}
\label{\detokenize{pair_yukawa_colloid:default}}
\sphinxAtStartPar
none


\bigskip\hrule\bigskip

\phantomsection\label{\detokenize{pair_yukawa_colloid:safran}}
\sphinxAtStartPar
\sphinxstylestrong{(Safran)} Safran, Statistical Thermodynamics of Surfaces, Interfaces,
And Membranes, Westview Press, ISBN: 978\sphinxhyphen{}0813340791 (2003).

\sphinxstepscope

\index{pair\_style zbl@\spxentry{pair\_style zbl}}\ignorespaces 
\index{pair\_style zbl/gpu@\spxentry{pair\_style zbl/gpu}}\ignorespaces 
\index{pair\_style zbl/kk@\spxentry{pair\_style zbl/kk}}\ignorespaces 
\index{pair\_style zbl/omp@\spxentry{pair\_style zbl/omp}}\ignorespaces 

\section{pair\_style zbl command}
\label{\detokenize{pair_zbl:pair-style-zbl-command}}\label{\detokenize{pair_zbl:index-3}}\label{\detokenize{pair_zbl:index-2}}\label{\detokenize{pair_zbl:index-1}}\label{\detokenize{pair_zbl:index-0}}\label{\detokenize{pair_zbl::doc}}
\sphinxAtStartPar
Accelerator Variants: \sphinxstyleemphasis{zbl/gpu}, \sphinxstyleemphasis{zbl/kk}, \sphinxstyleemphasis{zbl/omp}


\subsection{Syntax}
\label{\detokenize{pair_zbl:syntax}}
\begin{sphinxVerbatim}[commandchars=\\\{\}]
\PYG{k}{pair\PYGZus{}style}\PYG{+w}{ }\PYG{n}{zbl}\PYG{+w}{ }\PYG{n}{inner}\PYG{+w}{ }\PYG{n}{outer}
\end{sphinxVerbatim}
\begin{itemize}
\item {} 
\sphinxAtStartPar
inner = distance where switching function begins

\item {} 
\sphinxAtStartPar
outer = global cutoff for ZBL interaction

\end{itemize}


\subsection{Examples}
\label{\detokenize{pair_zbl:examples}}
\begin{sphinxVerbatim}[commandchars=\\\{\}]
\PYG{k}{pair\PYGZus{}style}\PYG{+w}{ }\PYG{n}{zbl}\PYG{+w}{ }\PYG{l+m}{3.0}\PYG{+w}{ }\PYG{l+m}{4.0}
\PYG{k}{pair\PYGZus{}coeff}\PYG{+w}{ }\PYG{o}{*}\PYG{+w}{ }\PYG{o}{*}\PYG{+w}{ }\PYG{l+m}{73.0}\PYG{+w}{ }\PYG{l+m}{73.0}
\PYG{k}{pair\PYGZus{}coeff}\PYG{+w}{ }\PYG{l+m}{1}\PYG{+w}{ }\PYG{l+m}{1}\PYG{+w}{ }\PYG{l+m}{14.0}\PYG{+w}{ }\PYG{l+m}{14.0}
\end{sphinxVerbatim}


\subsection{Description}
\label{\detokenize{pair_zbl:description}}
\sphinxAtStartPar
Style \sphinxstyleemphasis{zbl} computes the Ziegler\sphinxhyphen{}Biersack\sphinxhyphen{}Littmark (ZBL) screened nuclear
repulsion for describing high\sphinxhyphen{}energy collisions between atoms.
{\hyperref[\detokenize{pair_zbl:ziegler}]{\sphinxcrossref{\DUrole{std,std-ref}{(Ziegler)}}}}. It includes an additional switching function
that ramps the energy, force, and curvature smoothly to zero
between an inner and outer cutoff. The potential
energy due to a pair of atoms at a distance r\_ij is given by:
\begin{equation*}
\begin{split}E^{ZBL}_{ij} & = \frac{1}{4\pi\epsilon_0} \frac{Z_i Z_j \,e^2}{r_{ij}} \phi(r_{ij}/a)+ S(r_{ij})\\
a & =  \frac{0.46850}{Z_{i}^{0.23} + Z_{j}^{0.23}}\\
\phi(x) & =  0.18175e^{-3.19980x} + 0.50986e^{-0.94229x} + 0.28022e^{-0.40290x} + 0.02817e^{-0.20162x}\\\end{split}
\end{equation*}
\sphinxAtStartPar
where \sphinxstyleemphasis{e} is the electron charge, \(\epsilon_0\) is the electrical
permittivity of vacuum, and \(Z_i\) and \(Z_j\) are the nuclear
charges of the
two atoms.  The switching function \(S(r)\) is identical to that used by
{\hyperref[\detokenize{pair_gromacs::doc}]{\sphinxcrossref{\DUrole{doc}{pair\_style lj/gromacs}}}}.  Here, the inner and outer
cutoff are the same for all pairs of atom types.

\sphinxAtStartPar
The following coefficients must be defined for each pair of atom types
via the {\hyperref[\detokenize{pair_coeff::doc}]{\sphinxcrossref{\DUrole{doc}{pair\_coeff}}}} command as in the examples above,
or in the LAMMPS data file.
\begin{itemize}
\item {} 
\sphinxAtStartPar
\(Z_i\) (atomic number for first atom type, e.g. 13.0 for aluminum)

\item {} 
\sphinxAtStartPar
\(Z_j\) (ditto for second atom type)

\end{itemize}

\sphinxAtStartPar
The values of \(Z_i\) and \(Z_j\) are normally equal to the atomic
numbers of the two atom types. Thus, the user may optionally
specify only the coefficients for each \(i==i\) pair, and rely
on the obvious mixing rule for cross interactions (see below).
Note that when \(i==i\) it is required that \(Z_i == Z_j\).
When used with {\hyperref[\detokenize{pair_hybrid::doc}]{\sphinxcrossref{\DUrole{doc}{hybrid/overlay}}}} and pairs are
assigned
to more than one sub\sphinxhyphen{}style, the mixing rule is not used and
each pair of types interacting with the ZBL sub\sphinxhyphen{}style must
be included in a pair\_coeff command.

\begin{sphinxadmonition}{note}{Note:}
\sphinxAtStartPar
The numerical values of the exponential decay constants in the
screening function depend on the unit of distance. In the above
equation they are given for units of Angstroms. LAMMPS will
automatically convert these values to the distance unit of the
specified LAMMPS {\hyperref[\detokenize{units::doc}]{\sphinxcrossref{\DUrole{doc}{units}}}} setting.  The values of Z should
always be given as multiples of a proton’s charge, e.g. 29.0 for
copper.
\end{sphinxadmonition}


\bigskip\hrule\bigskip


\sphinxAtStartPar
Styles with a \sphinxstyleemphasis{gpu}, \sphinxstyleemphasis{intel}, \sphinxstyleemphasis{kk}, \sphinxstyleemphasis{omp}, or \sphinxstyleemphasis{opt} suffix are
functionally the same as the corresponding style without the suffix.
They have been optimized to run faster, depending on your available
hardware, as discussed on the {\hyperref[\detokenize{Speed_packages::doc}]{\sphinxcrossref{\DUrole{doc}{Accelerator packages}}}}
page.  The accelerated styles take the same arguments and should
produce the same results, except for round\sphinxhyphen{}off and precision issues.

\sphinxAtStartPar
These accelerated styles are part of the GPU, INTEL, KOKKOS,
OPENMP, and OPT packages, respectively.  They are only enabled if
LAMMPS was built with those packages.  See the {\hyperref[\detokenize{Build_package::doc}]{\sphinxcrossref{\DUrole{doc}{Build package}}}} page for more info.

\sphinxAtStartPar
You can specify the accelerated styles explicitly in your input script
by including their suffix, or you can use the {\hyperref[\detokenize{Run_options::doc}]{\sphinxcrossref{\DUrole{doc}{\sphinxhyphen{}suffix command\sphinxhyphen{}line switch}}}} when you invoke LAMMPS, or you can use the
{\hyperref[\detokenize{suffix::doc}]{\sphinxcrossref{\DUrole{doc}{suffix}}}} command in your input script.

\sphinxAtStartPar
See the {\hyperref[\detokenize{Speed_packages::doc}]{\sphinxcrossref{\DUrole{doc}{Accelerator packages}}}} page for more
instructions on how to use the accelerated styles effectively.


\bigskip\hrule\bigskip



\subsection{Mixing, shift, table, tail correction, restart, rRESPA info}
\label{\detokenize{pair_zbl:mixing-shift-table-tail-correction-restart-rrespa-info}}
\sphinxAtStartPar
For atom type pairs \sphinxstyleemphasis{i,j} and \(i \neq i\), the \(Z_i\) and
\(Z_j\) coefficients
can be mixed by taking \(Z_i\) and \(Z_j\) from the values
specified for
\(i == i\) and \(j == j\) cases. When used
with {\hyperref[\detokenize{pair_hybrid::doc}]{\sphinxcrossref{\DUrole{doc}{hybrid/overlay}}}} and pairs are assigned
to more than one sub\sphinxhyphen{}style, the mixing rule is not used and
each pair of types interacting with the ZBL sub\sphinxhyphen{}style
must be included in a pair\_coeff command.
The {\hyperref[\detokenize{pair_modify::doc}]{\sphinxcrossref{\DUrole{doc}{pair\_modify}}}} mix option has no effect on
the mixing behavior

\sphinxAtStartPar
The ZBL pair style does not support the {\hyperref[\detokenize{pair_modify::doc}]{\sphinxcrossref{\DUrole{doc}{pair\_modify}}}}
shift option, since the ZBL interaction is already smoothed to 0.0 at
the cutoff.

\sphinxAtStartPar
The {\hyperref[\detokenize{pair_modify::doc}]{\sphinxcrossref{\DUrole{doc}{pair\_modify}}}} table option is not relevant for
this pair style.

\sphinxAtStartPar
This pair style does not support the {\hyperref[\detokenize{pair_modify::doc}]{\sphinxcrossref{\DUrole{doc}{pair\_modify}}}}
tail option for adding long\sphinxhyphen{}range tail corrections to energy and
pressure, since there are no corrections for a potential that goes to
0.0 at the cutoff.

\sphinxAtStartPar
This pair style does not write information to {\hyperref[\detokenize{restart::doc}]{\sphinxcrossref{\DUrole{doc}{binary restart files}}}}, so pair\_style and pair\_coeff commands must be
specified in an input script that reads a restart file.

\sphinxAtStartPar
This pair style can only be used via the \sphinxstyleemphasis{pair} keyword of the
{\hyperref[\detokenize{run_style::doc}]{\sphinxcrossref{\DUrole{doc}{run\_style respa}}}} command.  It does not support the
\sphinxstyleemphasis{inner}, \sphinxstyleemphasis{middle}, \sphinxstyleemphasis{outer} keywords.


\bigskip\hrule\bigskip



\subsection{Restrictions}
\label{\detokenize{pair_zbl:restrictions}}\begin{quote}

\sphinxAtStartPar
none
\end{quote}


\subsection{Related commands}
\label{\detokenize{pair_zbl:related-commands}}
\sphinxAtStartPar
{\hyperref[\detokenize{pair_coeff::doc}]{\sphinxcrossref{\DUrole{doc}{pair\_coeff}}}}


\subsection{Default}
\label{\detokenize{pair_zbl:default}}
\sphinxAtStartPar
none


\bigskip\hrule\bigskip

\phantomsection\label{\detokenize{pair_zbl:ziegler}}
\sphinxAtStartPar
\sphinxstylestrong{(Ziegler)} J.F. Ziegler, J. P. Biersack and U. Littmark, “The
Stopping and Range of Ions in Matter”, Volume 1, Pergamon, 1985.

\sphinxstepscope

\index{pair\_style zero@\spxentry{pair\_style zero}}\ignorespaces 

\section{pair\_style zero command}
\label{\detokenize{pair_zero:pair-style-zero-command}}\label{\detokenize{pair_zero:index-0}}\label{\detokenize{pair_zero::doc}}

\subsection{Syntax}
\label{\detokenize{pair_zero:syntax}}
\begin{sphinxVerbatim}[commandchars=\\\{\}]
\PYG{k}{pair\PYGZus{}style}\PYG{+w}{ }\PYG{n}{zero}\PYG{+w}{ }\PYG{n}{cutoff}\PYG{+w}{ }\PYG{n}{[nocoeff]}\PYG{+w}{ }\PYG{n}{[full]}
\end{sphinxVerbatim}
\begin{itemize}
\item {} 
\sphinxAtStartPar
zero = style name of this pair style

\item {} 
\sphinxAtStartPar
cutoff = global cutoff (distance units)

\item {} 
\sphinxAtStartPar
nocoeff = ignore all pair\_coeff parameters (optional)

\item {} 
\sphinxAtStartPar
full = build full neighbor list (optional)

\end{itemize}


\subsection{Examples}
\label{\detokenize{pair_zero:examples}}
\begin{sphinxVerbatim}[commandchars=\\\{\}]
\PYG{k}{pair\PYGZus{}style}\PYG{+w}{ }\PYG{n}{zero}\PYG{+w}{ }\PYG{l+m}{10.0}
\PYG{k}{pair\PYGZus{}style}\PYG{+w}{ }\PYG{n}{zero}\PYG{+w}{ }\PYG{l+m}{5.0}\PYG{+w}{ }\PYG{n}{nocoeff}
\PYG{k}{pair\PYGZus{}coeff}\PYG{+w}{ }\PYG{o}{*}\PYG{+w}{ }\PYG{o}{*}
\PYG{k}{pair\PYGZus{}coeff}\PYG{+w}{ }\PYG{l+m}{1}\PYG{+w}{ }\PYG{l+m}{2}\PYG{o}{*}\PYG{l+m}{4}\PYG{+w}{ }\PYG{l+m}{3.0}
\end{sphinxVerbatim}


\subsection{Description}
\label{\detokenize{pair_zero:description}}
\sphinxAtStartPar
Define a global or per\sphinxhyphen{}type cutoff length for the purpose of
building a neighbor list and acquiring ghost atoms, but do
not compute any pairwise forces or energies.

\sphinxAtStartPar
This can be useful for fixes or computes which require a neighbor list
to enumerate pairs of atoms within some cutoff distance, but when
pairwise forces are not otherwise needed.  Examples are the {\hyperref[\detokenize{fix_bond_create::doc}]{\sphinxcrossref{\DUrole{doc}{fix bond/create}}}}, {\hyperref[\detokenize{compute_rdf::doc}]{\sphinxcrossref{\DUrole{doc}{compute rdf}}}},
{\hyperref[\detokenize{compute_voronoi_atom::doc}]{\sphinxcrossref{\DUrole{doc}{compute voronoi/atom}}}} commands.

\sphinxAtStartPar
Note that the {\hyperref[\detokenize{comm_modify::doc}]{\sphinxcrossref{\DUrole{doc}{comm\_modify cutoff}}}} command can be
used to ensure communication of ghost atoms even when a pair style is
not defined, but it will not trigger neighbor list generation.

\sphinxAtStartPar
The optional \sphinxstyleemphasis{nocoeff} flag allows to read data files with a PairCoeff
section for any pair style. Similarly, any pair\_coeff commands
will only be checked for the atom type numbers and the rest ignored.
In this case, only the global cutoff will be used.

\sphinxAtStartPar
\DUrole{versionmodified,added}{Added in version 3Nov2022.}

\sphinxAtStartPar
The optional \sphinxstyleemphasis{full} flag builds a full neighbor list instead of the default
half neighbor list.

\sphinxAtStartPar
The following coefficients must be defined for each pair of atoms
types via the {\hyperref[\detokenize{pair_coeff::doc}]{\sphinxcrossref{\DUrole{doc}{pair\_coeff}}}} command as in the examples
above, or in the data file or restart files read by the
{\hyperref[\detokenize{read_data::doc}]{\sphinxcrossref{\DUrole{doc}{read\_data}}}} or {\hyperref[\detokenize{read_restart::doc}]{\sphinxcrossref{\DUrole{doc}{read\_restart}}}}
commands, or by mixing as described below:
\begin{itemize}
\item {} 
\sphinxAtStartPar
cutoff (distance units)

\end{itemize}

\sphinxAtStartPar
This coefficient is optional.  If not specified, the global cutoff
specified in the pair\_style command is used. If the pair\_style has
been specified with the optional \sphinxstyleemphasis{nocoeff} flag, then a cutoff
pair coefficient is ignored.


\bigskip\hrule\bigskip



\subsection{Mixing, shift, table, tail correction, restart, rRESPA info}
\label{\detokenize{pair_zero:mixing-shift-table-tail-correction-restart-rrespa-info}}
\sphinxAtStartPar
The cutoff distance for this pair style can be mixed.  The default mix
value is \sphinxstyleemphasis{geometric}.  See the “pair\_modify” command for details.

\sphinxAtStartPar
This pair style does not support the {\hyperref[\detokenize{pair_modify::doc}]{\sphinxcrossref{\DUrole{doc}{pair\_modify}}}}
shift, table, and tail options.

\sphinxAtStartPar
This pair style writes its information to {\hyperref[\detokenize{restart::doc}]{\sphinxcrossref{\DUrole{doc}{binary restart files}}}}, so pair\_style and pair\_coeff commands do not need
to be specified in an input script that reads a restart file.

\sphinxAtStartPar
This pair style supports the use of the \sphinxstyleemphasis{inner}, \sphinxstyleemphasis{middle},
and \sphinxstyleemphasis{outer} keywords of the {\hyperref[\detokenize{run_style::doc}]{\sphinxcrossref{\DUrole{doc}{run\_style respa}}}} command.


\bigskip\hrule\bigskip



\subsection{Restrictions}
\label{\detokenize{pair_zero:restrictions}}
\sphinxAtStartPar
none


\subsection{Related commands}
\label{\detokenize{pair_zero:related-commands}}
\sphinxAtStartPar
{\hyperref[\detokenize{pair_none::doc}]{\sphinxcrossref{\DUrole{doc}{pair\_style none}}}}


\subsection{Default}
\label{\detokenize{pair_zero:default}}
\sphinxAtStartPar
none

\sphinxstepscope


\chapter{Bond Styles}
\label{\detokenize{bonds:bond-styles}}\label{\detokenize{bonds::doc}}
\sphinxstepscope

\index{bond\_style bpm/rotational@\spxentry{bond\_style bpm/rotational}}\ignorespaces 

\section{bond\_style bpm/rotational command}
\label{\detokenize{bond_bpm_rotational:bond-style-bpm-rotational-command}}\label{\detokenize{bond_bpm_rotational:index-0}}\label{\detokenize{bond_bpm_rotational::doc}}

\subsection{Syntax}
\label{\detokenize{bond_bpm_rotational:syntax}}
\begin{sphinxVerbatim}[commandchars=\\\{\}]
\PYG{k}{bond\PYGZus{}style}\PYG{+w}{ }\PYG{n}{bpm}\PYG{o}{/}\PYG{n}{rotational}\PYG{+w}{ }\PYG{n}{keyword}\PYG{+w}{ }\PYG{n}{value}\PYG{+w}{ }\PYG{n}{attribute1}\PYG{+w}{ }\PYG{n}{attribute2}\PYG{+w}{ }\PYG{n}{...}
\end{sphinxVerbatim}
\begin{itemize}
\item {} 
\sphinxAtStartPar
optional keyword = \sphinxstyleemphasis{overlay/pair} or \sphinxstyleemphasis{store/local} or \sphinxstyleemphasis{smooth} or \sphinxstyleemphasis{break}
\begin{sphinxalltt}
\sphinxstyleemphasis{store/local} values = fix\_ID N attributes ...
   * fix\_ID = ID of associated internal fix to store data
   * N = prepare data for output every this many timesteps
   * attributes = zero or more of the below attributes may be appended

     \sphinxstyleemphasis{id1, id2} = IDs of two atoms in the bond
     \sphinxstyleemphasis{time} = the timestep the bond broke
     \sphinxstyleemphasis{x, y, z} = the center of mass position of the two atoms when the bond broke (distance units)
     \sphinxstyleemphasis{x/ref, y/ref, z/ref} = the initial center of mass position of the two atoms (distance units)

\sphinxstyleemphasis{overlay/pair} value = \sphinxstyleemphasis{yes} or \sphinxstyleemphasis{no}
   bonded particles will still interact with pair forces

\sphinxstyleemphasis{smooth} value = \sphinxstyleemphasis{yes} or \sphinxstyleemphasis{no}
   smooths bond forces near the breaking point

\sphinxstyleemphasis{normalize} value = \sphinxstyleemphasis{yes} or \sphinxstyleemphasis{no}
   normalizes normal and shear forces by the reference length

\sphinxstyleemphasis{break} value = \sphinxstyleemphasis{yes} or \sphinxstyleemphasis{no}
   indicates whether bonds break during a run
\end{sphinxalltt}

\end{itemize}


\subsection{Examples}
\label{\detokenize{bond_bpm_rotational:examples}}
\begin{sphinxVerbatim}[commandchars=\\\{\}]
\PYG{k}{bond\PYGZus{}style}\PYG{+w}{ }\PYG{n}{bpm}\PYG{o}{/}\PYG{n}{rotational}
\PYG{k}{bond\PYGZus{}coeff}\PYG{+w}{ }\PYG{l+m}{1}\PYG{+w}{ }\PYG{l+m}{1.0}\PYG{+w}{ }\PYG{l+m}{0.2}\PYG{+w}{ }\PYG{l+m}{0.02}\PYG{+w}{ }\PYG{l+m}{0.02}\PYG{+w}{ }\PYG{l+m}{0.20}\PYG{+w}{ }\PYG{l+m}{0.04}\PYG{+w}{ }\PYG{l+m}{0.04}\PYG{+w}{ }\PYG{l+m}{0.04}\PYG{+w}{ }\PYG{l+m}{0.1}\PYG{+w}{ }\PYG{l+m}{0.02}\PYG{+w}{ }\PYG{l+m}{0.002}\PYG{+w}{ }\PYG{l+m}{0.002}

\PYG{k}{bond\PYGZus{}style}\PYG{+w}{ }\PYG{n}{bpm}\PYG{o}{/}\PYG{n}{rotational}\PYG{+w}{ }\PYG{n}{store}\PYG{o}{/}\PYG{n}{local}\PYG{+w}{ }\PYG{n}{myfix}\PYG{+w}{ }\PYG{l+m}{1000}\PYG{+w}{ }\PYG{n}{time}\PYG{+w}{ }\PYG{n}{id1}\PYG{+w}{ }\PYG{n}{id2}
\PYG{k}{dump }\PYG{n+nv+nvIdentifier}{1}\PYG{+w}{ }\PYG{n+nv+nvIdentifier}{all}\PYG{+w}{ }\PYG{n}{local}\PYG{+w}{ }\PYG{l+m}{1000}\PYG{+w}{ }\PYG{n}{dump.broken}\PYG{+w}{ }\PYG{n}{f\PYGZus{}myfix[1]}\PYG{+w}{ }\PYG{n}{f\PYGZus{}myfix[2]}\PYG{+w}{ }\PYG{n}{f\PYGZus{}myfix[3]}
\PYG{k}{dump\PYGZus{}modify }\PYG{n+nv+nvIdentifier}{1}\PYG{+w}{ }\PYG{n}{write\PYGZus{}header}\PYG{+w}{ }\PYG{n}{no}
\end{sphinxVerbatim}


\subsection{Description}
\label{\detokenize{bond_bpm_rotational:description}}
\sphinxAtStartPar
\DUrole{versionmodified,added}{Added in version 4May2022.}

\sphinxAtStartPar
The \sphinxstyleemphasis{bpm/rotational} bond style computes forces and torques based on
deviations from the initial reference state of the two atoms.  The
reference state is stored by each bond when it is first computed in
the setup of a run. Data is then preserved across run commands and is
written to {\hyperref[\detokenize{restart::doc}]{\sphinxcrossref{\DUrole{doc}{binary restart files}}}} such that restarting
the system will not reset the reference state of a bond.

\sphinxAtStartPar
Forces include a normal and tangential component. The base normal force
has a magnitude of
\begin{equation*}
\begin{split}f_r = k_r (r - r_0)\end{split}
\end{equation*}
\sphinxAtStartPar
where \(k_r\) is a stiffness and \(r\) is the current distance and
\(r_0\) is the initial distance between the two particles.

\sphinxAtStartPar
A tangential force is applied perpendicular to the normal direction
which is proportional to the tangential shear displacement with a
stiffness of \(k_s\). This tangential force also induces a torque.
In addition, bending and twisting torques are also applied to
particles which are proportional to angular bending and twisting
displacements with stiffnesses of \(k_b\) and \(k_t\),
respectively.  Details on the calculations of shear displacements and
angular displacements can be found in {\hyperref[\detokenize{bond_bpm_rotational:wang2009}]{\sphinxcrossref{\DUrole{std,std-ref}{(Wang)}}}} and
{\hyperref[\detokenize{bond_bpm_rotational:wang2009b}]{\sphinxcrossref{\DUrole{std,std-ref}{(Wang and Mora)}}}}.

\sphinxAtStartPar
Bonds will break under sufficient stress. A breaking criterion is calculated
\begin{equation*}
\begin{split}B = \mathrm{max}\left\{0, \frac{f_r}{f_{r,c}} + \frac{|f_s|}{f_{s,c}} +
    \frac{|\tau_b|}{\tau_{b,c}} + \frac{|\tau_t|}{\tau_{t,c}} \right\}\end{split}
\end{equation*}
\sphinxAtStartPar
where \(|f_s|\) is the magnitude of the shear force and
\(|\tau_b|\) and \(|\tau_t|\) are the magnitudes of the
bending and twisting torques, respectively. The corresponding variables
\(f_{r,c}\) \(f_{s,c}\), \(\tau_{b,c}\), and
\(\tau_{t,c}\) are critical limits to each force or torque.  If
\(B\) is ever equal to or exceeds one, the bond will break.  This
is done by setting the bond type to 0 such that forces and
torques are no longer computed.

\sphinxAtStartPar
After computing the base magnitudes of the forces and torques, they
can be optionally multiplied by an extra factor \(w\) to smoothly
interpolate forces and torques to zero as the bond breaks. This term
is calculated as \(w = (1.0 - B^4)\). This smoothing factor can be added
or removed by setting the \sphinxstyleemphasis{smooth} keyword to \sphinxstyleemphasis{yes} or \sphinxstyleemphasis{no}, respectively.

\sphinxAtStartPar
Finally, additional damping forces and torques are applied to the two
particles. A force is applied proportional to the difference in the
normal velocity of particles using a similar construction as
dissipative particle dynamics {\hyperref[\detokenize{bond_bpm_rotational:groot3}]{\sphinxcrossref{\DUrole{std,std-ref}{(Groot)}}}}:
\begin{equation*}
\begin{split}F_D = - \gamma_n w (\hat{r} \bullet \vec{v})\end{split}
\end{equation*}
\sphinxAtStartPar
where \(\gamma_n\) is the damping strength, \(\hat{r}\) is the
radial normal vector, and \(\vec{v}\) is the velocity difference
between the two particles. Similarly, tangential forces are applied to
each atom proportional to the relative differences in sliding
velocities with a constant prefactor \(\gamma_s\) {\hyperref[\detokenize{bond_bpm_rotational:wang20152}]{\sphinxcrossref{\DUrole{std,std-ref}{(Wang et
al.)}}}} along with their associated torques. The rolling and
twisting components of the relative angular velocities of the two
atoms are also damped by applying torques with prefactors of
\(\gamma_r\) and \(\gamma_t\), respectively.

\sphinxAtStartPar
The following coefficients must be defined for each bond type via the
{\hyperref[\detokenize{bond_coeff::doc}]{\sphinxcrossref{\DUrole{doc}{bond\_coeff}}}} command as in the example above, or in
the data file or restart files read by the {\hyperref[\detokenize{read_data::doc}]{\sphinxcrossref{\DUrole{doc}{read\_data}}}}
or {\hyperref[\detokenize{read_restart::doc}]{\sphinxcrossref{\DUrole{doc}{read\_restart}}}} commands:
\begin{itemize}
\item {} 
\sphinxAtStartPar
\(k_r\)           (force/distance units)

\item {} 
\sphinxAtStartPar
\(k_s\)           (force/distance units)

\item {} 
\sphinxAtStartPar
\(k_t\)           (force*distance/radians units)

\item {} 
\sphinxAtStartPar
\(k_b\)           (force*distance/radians units)

\item {} 
\sphinxAtStartPar
\(f_{r,c}\)       (force units)

\item {} 
\sphinxAtStartPar
\(f_{s,c}\)       (force units)

\item {} 
\sphinxAtStartPar
\(\tau_{b,c}\)    (force*distance units)

\item {} 
\sphinxAtStartPar
\(\tau_{t,c}\)    (force*distance units)

\item {} 
\sphinxAtStartPar
\(\gamma_n\)      (force/velocity units)

\item {} 
\sphinxAtStartPar
\(\gamma_s\)      (force/velocity units)

\item {} 
\sphinxAtStartPar
\(\gamma_r\)      (force*distance/velocity units)

\item {} 
\sphinxAtStartPar
\(\gamma_t\)      (force*distance/velocity units)

\end{itemize}

\sphinxAtStartPar
If the \sphinxstyleemphasis{normalize} keyword is set to \sphinxstyleemphasis{yes}, the radial and shear forces
will be normalized by \(r_0\) such that \(k_r\) and \(k_s\)
must be given in force units.

\sphinxAtStartPar
By default, pair forces are not calculated between bonded particles.
Pair forces can alternatively be overlaid on top of bond forces by setting
the \sphinxstyleemphasis{overlay/pair} keyword to \sphinxstyleemphasis{yes}. These settings require specific
{\hyperref[\detokenize{special_bonds::doc}]{\sphinxcrossref{\DUrole{doc}{special\_bonds}}}} settings described in the
restrictions.  Further details can be found in the {\hyperref[\detokenize{Howto_bpm::doc}]{\sphinxcrossref{\DUrole{doc}{how to}}}}
page on BPMs.

\sphinxAtStartPar
\DUrole{versionmodified,added}{Added in version 28Mar2023.}

\sphinxAtStartPar
If the \sphinxstyleemphasis{break} keyword is set to \sphinxstyleemphasis{no}, LAMMPS assumes bonds should not break
during a simulation run. This will prevent some unnecessary calculation.
However, if a bond reaches a damage criterion greater than one,
it will trigger an error.

\sphinxAtStartPar
If the \sphinxstyleemphasis{store/local} keyword is used, an internal fix will track bonds that
break during the simulation. Whenever a bond breaks, data is processed
and transferred to an internal fix labeled \sphinxstyleemphasis{fix\_ID}. This allows the
local data to be accessed by other LAMMPS commands. Following this optional
keyword, a list of one or more attributes is specified.  These include the
IDs of the two atoms in the bond. The other attributes for the two atoms
include the timestep during which the bond broke and the current/initial
center of mass position of the two atoms.

\sphinxAtStartPar
Data is continuously accumulated over intervals of \sphinxstyleemphasis{N}
timesteps. At the end of each interval, all of the saved accumulated
data is deleted to make room for new data. Individual datum may
therefore persist anywhere between \sphinxstyleemphasis{1} to \sphinxstyleemphasis{N} timesteps depending on
when they are saved. This data can be accessed using the \sphinxstyleemphasis{fix\_ID} and a
{\hyperref[\detokenize{dump::doc}]{\sphinxcrossref{\DUrole{doc}{dump local}}}} command. To ensure all data is output,
the dump frequency should correspond to the same interval of \sphinxstyleemphasis{N}
timesteps. A dump frequency of an integer multiple of \sphinxstyleemphasis{N} can be used
to regularly output a sample of the accumulated data.

\sphinxAtStartPar
Note that when unbroken bonds are dumped to a file via the
{\hyperref[\detokenize{dump::doc}]{\sphinxcrossref{\DUrole{doc}{dump local}}}} command, bonds with type 0 (broken bonds)
are not included.
The {\hyperref[\detokenize{delete_bonds::doc}]{\sphinxcrossref{\DUrole{doc}{delete\_bonds}}}} command can also be used to
query the status of broken bonds or permanently delete them, e.g.:

\begin{sphinxVerbatim}[commandchars=\\\{\}]
\PYG{k}{delete\PYGZus{}bonds }\PYG{n+nv+nvIdentifier}{all}\PYG{+w}{ }\PYG{n}{stats}
\PYG{k}{delete\PYGZus{}bonds }\PYG{n+nv+nvIdentifier}{all}\PYG{+w}{ }\PYG{n}{bond}\PYG{+w}{ }\PYG{l+m}{0}\PYG{+w}{ }\PYG{n}{remove}
\end{sphinxVerbatim}


\bigskip\hrule\bigskip



\subsection{Restart and other info}
\label{\detokenize{bond_bpm_rotational:restart-and-other-info}}
\sphinxAtStartPar
This bond style writes the reference state of each bond to
{\hyperref[\detokenize{restart::doc}]{\sphinxcrossref{\DUrole{doc}{binary restart files}}}}. Loading a restart file will
properly resume bonds. However, the reference state is NOT
written to data files. Therefore reading a data file will not
restore bonds and will cause their reference states to be redefined.

\sphinxAtStartPar
If the \sphinxstyleemphasis{store/local} option is used, an internal fix will calculate
a local vector or local array depending on the number of input values.
The length of the vector or number of rows in the array is the number
of recorded, broken bonds.  If a single input is specified, a local
vector is produced. If two or more inputs are specified, a local array
is produced where the number of columns = the number of inputs.  The
vector or array can be accessed by any command that uses local values
from a compute as input. See the {\hyperref[\detokenize{Howto_output::doc}]{\sphinxcrossref{\DUrole{doc}{Howto output}}}} page
for an overview of LAMMPS output options.

\sphinxAtStartPar
The vector or array will be floating point values that correspond to
the specified attribute.

\sphinxAtStartPar
The single() function of this bond style returns 0.0 for the energy
of a bonded interaction, since energy is not conserved in these
dissipative potentials.  It also returns only the normal component of
the bonded interaction force.  However, the single() function also
calculates 7 extra bond quantities.  The first 4 are data from the
reference state of the bond including the initial distance between particles
\(r_0\) followed by the \(x\), \(y\), and \(z\) components
of the initial unit vector pointing to particle I from particle J. The next 3
quantities (5\sphinxhyphen{}7) are the  \(x\), \(y\), and \(z\) components
of the total force, including normal and tangential contributions, acting
on particle I.

\sphinxAtStartPar
These extra quantities can be accessed by the {\hyperref[\detokenize{compute_bond_local::doc}]{\sphinxcrossref{\DUrole{doc}{compute bond/local}}}}
command, as \sphinxstyleemphasis{b1}, \sphinxstyleemphasis{b2}, …, \sphinxstyleemphasis{b7}.


\subsection{Restrictions}
\label{\detokenize{bond_bpm_rotational:restrictions}}
\sphinxAtStartPar
This bond style is part of the BPM package.  It is only enabled if
LAMMPS was built with that package.  See the {\hyperref[\detokenize{Build_package::doc}]{\sphinxcrossref{\DUrole{doc}{Build package}}}} page for more info.

\sphinxAtStartPar
By default if pair interactions between bonded atoms are to be disabled,
this bond style requires setting

\begin{sphinxVerbatim}[commandchars=\\\{\}]
\PYG{k}{special\PYGZus{}bonds}\PYG{+w}{ }\PYG{n}{lj}\PYG{+w}{ }\PYG{l+m}{0}\PYG{+w}{ }\PYG{l+m}{1}\PYG{+w}{ }\PYG{l+m}{1}\PYG{+w}{ }\PYG{n}{coul}\PYG{+w}{ }\PYG{l+m}{1}\PYG{+w}{ }\PYG{l+m}{1}\PYG{+w}{ }\PYG{l+m}{1}
\end{sphinxVerbatim}

\sphinxAtStartPar
and {\hyperref[\detokenize{newton::doc}]{\sphinxcrossref{\DUrole{doc}{newton}}}} must be set to bond off.  If the \sphinxstyleemphasis{overlay/pair}
keyword is set to \sphinxstyleemphasis{yes}, this bond style alternatively requires setting

\begin{sphinxVerbatim}[commandchars=\\\{\}]
\PYG{k}{special\PYGZus{}bonds}\PYG{+w}{ }\PYG{n}{lj}\PYG{o}{/}\PYG{n}{coul}\PYG{+w}{ }\PYG{l+m}{1}\PYG{+w}{ }\PYG{l+m}{1}\PYG{+w}{ }\PYG{l+m}{1}
\end{sphinxVerbatim}

\sphinxAtStartPar
The \sphinxstyleemphasis{bpm/rotational} style requires {\hyperref[\detokenize{atom_style::doc}]{\sphinxcrossref{\DUrole{doc}{atom style bpm/sphere}}}}.


\subsection{Related commands}
\label{\detokenize{bond_bpm_rotational:related-commands}}
\sphinxAtStartPar
{\hyperref[\detokenize{bond_coeff::doc}]{\sphinxcrossref{\DUrole{doc}{bond\_coeff}}}}, {\hyperref[\detokenize{fix_nve_bpm_sphere::doc}]{\sphinxcrossref{\DUrole{doc}{fix nve/bpm/sphere}}}}


\subsection{Default}
\label{\detokenize{bond_bpm_rotational:default}}
\sphinxAtStartPar
The option defaults are \sphinxstyleemphasis{overlay/pair} = \sphinxstyleemphasis{no}, \sphinxstyleemphasis{smooth} = \sphinxstyleemphasis{yes}, \sphinxstyleemphasis{normalize} = \sphinxstyleemphasis{no}, and \sphinxstyleemphasis{break} = \sphinxstyleemphasis{yes}


\bigskip\hrule\bigskip

\phantomsection\label{\detokenize{bond_bpm_rotational:wang2009}}
\sphinxAtStartPar
\sphinxstylestrong{(Wang)} Wang, Acta Geotechnica, 4,
p 117\sphinxhyphen{}127 (2009).

\phantomsection\label{\detokenize{bond_bpm_rotational:wang2009b}}
\sphinxAtStartPar
\sphinxstylestrong{(Wang and Mora)} Wang, Mora, Advances in Geocomputing,
119, p 183\sphinxhyphen{}228 (2009).

\phantomsection\label{\detokenize{bond_bpm_rotational:groot3}}
\sphinxAtStartPar
\sphinxstylestrong{(Groot)} Groot and Warren, J Chem Phys, 107, 4423\sphinxhyphen{}35 (1997).

\phantomsection\label{\detokenize{bond_bpm_rotational:wang20152}}
\sphinxAtStartPar
\sphinxstylestrong{(Wang et al, 2015)} Wang, Y., Alonso\sphinxhyphen{}Marroquin, F., \& Guo,
W. W. (2015).  Rolling and sliding in 3\sphinxhyphen{}D discrete element
models. Particuology, 23, 49\sphinxhyphen{}55.

\sphinxstepscope

\index{bond\_style bpm/spring@\spxentry{bond\_style bpm/spring}}\ignorespaces 

\section{bond\_style bpm/spring command}
\label{\detokenize{bond_bpm_spring:bond-style-bpm-spring-command}}\label{\detokenize{bond_bpm_spring:index-0}}\label{\detokenize{bond_bpm_spring::doc}}

\subsection{Syntax}
\label{\detokenize{bond_bpm_spring:syntax}}
\begin{sphinxVerbatim}[commandchars=\\\{\}]
\PYG{k}{bond\PYGZus{}style}\PYG{+w}{ }\PYG{n}{bpm}\PYG{o}{/}\PYG{n}{spring}\PYG{+w}{ }\PYG{n}{keyword}\PYG{+w}{ }\PYG{n}{value}\PYG{+w}{ }\PYG{n}{attribute1}\PYG{+w}{ }\PYG{n}{attribute2}\PYG{+w}{ }\PYG{n}{...}
\end{sphinxVerbatim}
\begin{itemize}
\item {} 
\sphinxAtStartPar
optional keyword = \sphinxstyleemphasis{overlay/pair} or \sphinxstyleemphasis{store/local} or \sphinxstyleemphasis{smooth} or \sphinxstyleemphasis{break}
\begin{sphinxalltt}
\sphinxstyleemphasis{store/local} values = fix\_ID N attributes ...
   * fix\_ID = ID of associated internal fix to store data
   * N = prepare data for output every this many timesteps
   * attributes = zero or more of the below attributes may be appended

     \sphinxstyleemphasis{id1, id2} = IDs of two atoms in the bond
     \sphinxstyleemphasis{time} = the timestep the bond broke
     \sphinxstyleemphasis{x, y, z} = the center of mass position of the two atoms when the bond broke (distance units)
     \sphinxstyleemphasis{x/ref, y/ref, z/ref} = the initial center of mass position of the two atoms (distance units)

\sphinxstyleemphasis{overlay/pair} value = \sphinxstyleemphasis{yes} or \sphinxstyleemphasis{no}
   bonded particles will still interact with pair forces

\sphinxstyleemphasis{smooth} value = \sphinxstyleemphasis{yes} or \sphinxstyleemphasis{no}
   smooths bond forces near the breaking point

\sphinxstyleemphasis{normalize} value = \sphinxstyleemphasis{yes} or \sphinxstyleemphasis{no}
   normalizes bond forces by the reference length

\sphinxstyleemphasis{break} value = \sphinxstyleemphasis{yes} or \sphinxstyleemphasis{no}
   indicates whether bonds break during a run
\end{sphinxalltt}

\end{itemize}


\subsection{Examples}
\label{\detokenize{bond_bpm_spring:examples}}
\begin{sphinxVerbatim}[commandchars=\\\{\}]
\PYG{k}{bond\PYGZus{}style}\PYG{+w}{ }\PYG{n}{bpm}\PYG{o}{/}\PYG{n}{spring}
\PYG{k}{bond\PYGZus{}coeff}\PYG{+w}{ }\PYG{l+m}{1}\PYG{+w}{ }\PYG{l+m}{1.0}\PYG{+w}{ }\PYG{l+m}{0.05}\PYG{+w}{ }\PYG{l+m}{0.1}

\PYG{k}{bond\PYGZus{}style}\PYG{+w}{ }\PYG{n}{bpm}\PYG{o}{/}\PYG{n}{spring}\PYG{+w}{ }\PYG{n}{myfix}\PYG{+w}{ }\PYG{l+m}{1000}\PYG{+w}{ }\PYG{n}{time}\PYG{+w}{ }\PYG{n}{id1}\PYG{+w}{ }\PYG{n}{id2}
\PYG{k}{dump }\PYG{n+nv+nvIdentifier}{1}\PYG{+w}{ }\PYG{n+nv+nvIdentifier}{all}\PYG{+w}{ }\PYG{n}{local}\PYG{+w}{ }\PYG{l+m}{1000}\PYG{+w}{ }\PYG{n}{dump.broken}\PYG{+w}{ }\PYG{n}{f\PYGZus{}myfix[1]}\PYG{+w}{ }\PYG{n}{f\PYGZus{}myfix[2]}\PYG{+w}{ }\PYG{n}{f\PYGZus{}myfix[3]}
\PYG{k}{dump\PYGZus{}modify }\PYG{n+nv+nvIdentifier}{1}\PYG{+w}{ }\PYG{n}{write\PYGZus{}header}\PYG{+w}{ }\PYG{n}{no}
\end{sphinxVerbatim}


\subsection{Description}
\label{\detokenize{bond_bpm_spring:description}}
\sphinxAtStartPar
\DUrole{versionmodified,added}{Added in version 4May2022.}

\sphinxAtStartPar
The \sphinxstyleemphasis{bpm/spring} bond style computes forces based on
deviations from the initial reference state of the two atoms.  The
reference state is stored by each bond when it is first computed in
the setup of a run. Data is then preserved across run commands and is
written to {\hyperref[\detokenize{restart::doc}]{\sphinxcrossref{\DUrole{doc}{binary restart files}}}} such that restarting
the system will not reset the reference state of a bond.

\sphinxAtStartPar
This bond style only applies central\sphinxhyphen{}body forces which conserve the
translational and rotational degrees of freedom of a bonded set of
particles based on a model described by Clemmer and Robbins
{\hyperref[\detokenize{bond_bpm_spring:fragment-clemmer}]{\sphinxcrossref{\DUrole{std,std-ref}{(Clemmer)}}}}. The force has a magnitude of
\begin{equation*}
\begin{split}F = k (r - r_0) w\end{split}
\end{equation*}
\sphinxAtStartPar
where \(k\) is a stiffness, \(r\) is the current distance
and \(r_0\) is the initial distance between the two particles, and
\(w\) is an optional smoothing factor discussed below. Bonds will
break at a strain of \(\epsilon_c\).  This is done by setting
the bond type to 0 such that forces are no longer computed.

\sphinxAtStartPar
An additional damping force is applied to the bonded
particles.  This forces is proportional to the difference in the
normal velocity of particles using a similar construction as
dissipative particle dynamics {\hyperref[\detokenize{bond_bpm_spring:groot4}]{\sphinxcrossref{\DUrole{std,std-ref}{(Groot)}}}}:
\begin{equation*}
\begin{split}F_D = - \gamma w (\hat{r} \bullet \vec{v})\end{split}
\end{equation*}
\sphinxAtStartPar
where \(\gamma\) is the damping strength, \(\hat{r}\) is the
radial normal vector, and \(\vec{v}\) is the velocity difference
between the two particles.

\sphinxAtStartPar
The smoothing factor \(w\) can be added or removed by setting the
\sphinxstyleemphasis{smooth} keyword to \sphinxstyleemphasis{yes} or \sphinxstyleemphasis{no}, respectively. It is constructed such
that forces smoothly go to zero, avoiding discontinuities, as bonds
approach the critical strain
\begin{equation*}
\begin{split}w = 1.0 - \left( \frac{r - r_0}{r_0 \epsilon_c} \right)^8 .\end{split}
\end{equation*}
\sphinxAtStartPar
The following coefficients must be defined for each bond type via the
{\hyperref[\detokenize{bond_coeff::doc}]{\sphinxcrossref{\DUrole{doc}{bond\_coeff}}}} command as in the example above, or in
the data file or restart files read by the {\hyperref[\detokenize{read_data::doc}]{\sphinxcrossref{\DUrole{doc}{read\_data}}}} or {\hyperref[\detokenize{read_restart::doc}]{\sphinxcrossref{\DUrole{doc}{read\_restart}}}} commands:
\begin{itemize}
\item {} 
\sphinxAtStartPar
\(k\)             (force/distance units)

\item {} 
\sphinxAtStartPar
\(\epsilon_c\)    (unit less)

\item {} 
\sphinxAtStartPar
\(\gamma\)        (force/velocity units)

\end{itemize}

\sphinxAtStartPar
If the \sphinxstyleemphasis{normalize} keyword is set to \sphinxstyleemphasis{yes}, the elastic bond force will be
normalized by \(r_0\) such that \(k\) must be given in force units.

\sphinxAtStartPar
By default, pair forces are not calculated between bonded particles.
Pair forces can alternatively be overlaid on top of bond forces by setting
the \sphinxstyleemphasis{overlay/pair} keyword to \sphinxstyleemphasis{yes}. These settings require specific
{\hyperref[\detokenize{special_bonds::doc}]{\sphinxcrossref{\DUrole{doc}{special\_bonds}}}} settings described in the
restrictions.  Further details can be found in the {\hyperref[\detokenize{Howto_bpm::doc}]{\sphinxcrossref{\DUrole{doc}{how to}}}}
page on BPMs.

\sphinxAtStartPar
\DUrole{versionmodified,added}{Added in version 28Mar2023.}

\sphinxAtStartPar
If the \sphinxstyleemphasis{break} keyword is set to \sphinxstyleemphasis{no}, LAMMPS assumes bonds should not break
during a simulation run. This will prevent some unnecessary calculation.
However, if a bond reaches a strain greater than \(\epsilon_c\),
it will trigger an error.

\sphinxAtStartPar
If the \sphinxstyleemphasis{store/local} keyword is used, an internal fix will track bonds that
break during the simulation. Whenever a bond breaks, data is processed
and transferred to an internal fix labeled \sphinxstyleemphasis{fix\_ID}. This allows the
local data to be accessed by other LAMMPS commands. Following this optional
keyword, a list of one or more attributes is specified.  These include the
IDs of the two atoms in the bond. The other attributes for the two atoms
include the timestep during which the bond broke and the current/initial
center of mass position of the two atoms.

\sphinxAtStartPar
Data is continuously accumulated over intervals of \sphinxstyleemphasis{N}
timesteps. At the end of each interval, all of the saved accumulated
data is deleted to make room for new data. Individual datum may
therefore persist anywhere between \sphinxstyleemphasis{1} to \sphinxstyleemphasis{N} timesteps depending on
when they are saved. This data can be accessed using the \sphinxstyleemphasis{fix\_ID} and a
{\hyperref[\detokenize{dump::doc}]{\sphinxcrossref{\DUrole{doc}{dump local}}}} command. To ensure all data is output,
the dump frequency should correspond to the same interval of \sphinxstyleemphasis{N}
timesteps. A dump frequency of an integer multiple of \sphinxstyleemphasis{N} can be used
to regularly output a sample of the accumulated data.

\sphinxAtStartPar
Note that when unbroken bonds are dumped to a file via the
{\hyperref[\detokenize{dump::doc}]{\sphinxcrossref{\DUrole{doc}{dump local}}}} command, bonds with type 0 (broken bonds)
are not included.
The {\hyperref[\detokenize{delete_bonds::doc}]{\sphinxcrossref{\DUrole{doc}{delete\_bonds}}}} command can also be used to
query the status of broken bonds or permanently delete them, e.g.:

\begin{sphinxVerbatim}[commandchars=\\\{\}]
\PYG{k}{delete\PYGZus{}bonds }\PYG{n+nv+nvIdentifier}{all}\PYG{+w}{ }\PYG{n}{stats}
\PYG{k}{delete\PYGZus{}bonds }\PYG{n+nv+nvIdentifier}{all}\PYG{+w}{ }\PYG{n}{bond}\PYG{+w}{ }\PYG{l+m}{0}\PYG{+w}{ }\PYG{n}{remove}
\end{sphinxVerbatim}


\bigskip\hrule\bigskip



\subsection{Restart and other info}
\label{\detokenize{bond_bpm_spring:restart-and-other-info}}
\sphinxAtStartPar
This bond style writes the reference state of each bond to
{\hyperref[\detokenize{restart::doc}]{\sphinxcrossref{\DUrole{doc}{binary restart files}}}}. Loading a restart
file will properly restore bonds. However, the reference state is NOT
written to data files. Therefore reading a data file will not
restore bonds and will cause their reference states to be redefined.

\sphinxAtStartPar
If the \sphinxstyleemphasis{store/local} option is used, an internal fix will calculate
a local vector or local array depending on the number of input values.
The length of the vector or number of rows in the array is the number
of recorded, broken bonds.  If a single input is specified, a local
vector is produced. If two or more inputs are specified, a local array
is produced where the number of columns = the number of inputs.  The
vector or array can be accessed by any command that uses local values
from a compute as input. See the {\hyperref[\detokenize{Howto_output::doc}]{\sphinxcrossref{\DUrole{doc}{Howto output}}}} page
for an overview of LAMMPS output options.

\sphinxAtStartPar
The vector or array will be floating point values that correspond to
the specified attribute.

\sphinxAtStartPar
The single() function of this bond style returns 0.0 for the energy
of a bonded interaction, since energy is not conserved in these
dissipative potentials.  The single() function also calculates an
extra bond quantity, the initial distance \(r_0\). This
extra quantity can be accessed by the
{\hyperref[\detokenize{compute_bond_local::doc}]{\sphinxcrossref{\DUrole{doc}{compute bond/local}}}} command as \sphinxstyleemphasis{b1}.


\subsection{Restrictions}
\label{\detokenize{bond_bpm_spring:restrictions}}
\sphinxAtStartPar
This bond style is part of the BPM package.  It is only enabled if
LAMMPS was built with that package.  See the {\hyperref[\detokenize{Build_package::doc}]{\sphinxcrossref{\DUrole{doc}{Build package}}}} page for more info.

\sphinxAtStartPar
By default if pair interactions between bonded atoms are to be disabled,
this bond style requires setting

\begin{sphinxVerbatim}[commandchars=\\\{\}]
\PYG{k}{special\PYGZus{}bonds}\PYG{+w}{ }\PYG{n}{lj}\PYG{+w}{ }\PYG{l+m}{0}\PYG{+w}{ }\PYG{l+m}{1}\PYG{+w}{ }\PYG{l+m}{1}\PYG{+w}{ }\PYG{n}{coul}\PYG{+w}{ }\PYG{l+m}{1}\PYG{+w}{ }\PYG{l+m}{1}\PYG{+w}{ }\PYG{l+m}{1}
\end{sphinxVerbatim}

\sphinxAtStartPar
and {\hyperref[\detokenize{newton::doc}]{\sphinxcrossref{\DUrole{doc}{newton}}}} must be set to bond off.  If the \sphinxstyleemphasis{overlay/pair}
keyword is set to \sphinxstyleemphasis{yes}, this bond style alternatively requires setting

\begin{sphinxVerbatim}[commandchars=\\\{\}]
\PYG{k}{special\PYGZus{}bonds}\PYG{+w}{ }\PYG{n}{lj}\PYG{o}{/}\PYG{n}{coul}\PYG{+w}{ }\PYG{l+m}{1}\PYG{+w}{ }\PYG{l+m}{1}\PYG{+w}{ }\PYG{l+m}{1}
\end{sphinxVerbatim}


\subsection{Related commands}
\label{\detokenize{bond_bpm_spring:related-commands}}
\sphinxAtStartPar
{\hyperref[\detokenize{bond_coeff::doc}]{\sphinxcrossref{\DUrole{doc}{bond\_coeff}}}}, {\hyperref[\detokenize{pair_bpm_spring::doc}]{\sphinxcrossref{\DUrole{doc}{pair bpm/spring}}}}


\subsection{Default}
\label{\detokenize{bond_bpm_spring:default}}
\sphinxAtStartPar
The option defaults are \sphinxstyleemphasis{overlay/pair} = \sphinxstyleemphasis{no}, \sphinxstyleemphasis{smooth} = \sphinxstyleemphasis{yes}, \sphinxstyleemphasis{normalize} = \sphinxstyleemphasis{no}, and \sphinxstyleemphasis{break} = \sphinxstyleemphasis{yes}


\bigskip\hrule\bigskip

\phantomsection\label{\detokenize{bond_bpm_spring:fragment-clemmer}}
\sphinxAtStartPar
\sphinxstylestrong{(Clemmer)} Clemmer and Robbins, Phys. Rev. Lett. (2022).

\phantomsection\label{\detokenize{bond_bpm_spring:groot4}}
\sphinxAtStartPar
\sphinxstylestrong{(Groot)} Groot and Warren, J Chem Phys, 107, 4423\sphinxhyphen{}35 (1997).

\sphinxstepscope

\index{bond\_style class2@\spxentry{bond\_style class2}}\ignorespaces 
\index{bond\_style class2/omp@\spxentry{bond\_style class2/omp}}\ignorespaces 
\index{bond\_style class2/kk@\spxentry{bond\_style class2/kk}}\ignorespaces 

\section{bond\_style class2 command}
\label{\detokenize{bond_class2:bond-style-class2-command}}\label{\detokenize{bond_class2:index-2}}\label{\detokenize{bond_class2:index-1}}\label{\detokenize{bond_class2:index-0}}\label{\detokenize{bond_class2::doc}}
\sphinxAtStartPar
Accelerator Variants: \sphinxstyleemphasis{class2/omp}, \sphinxstyleemphasis{class2/kk}


\subsection{Syntax}
\label{\detokenize{bond_class2:syntax}}
\begin{sphinxVerbatim}[commandchars=\\\{\}]
\PYG{k}{bond\PYGZus{}style}\PYG{+w}{ }\PYG{n}{class2}
\end{sphinxVerbatim}


\subsection{Examples}
\label{\detokenize{bond_class2:examples}}
\begin{sphinxVerbatim}[commandchars=\\\{\}]
\PYG{k}{bond\PYGZus{}style}\PYG{+w}{ }\PYG{n}{class2}
\PYG{k}{bond\PYGZus{}coeff}\PYG{+w}{ }\PYG{l+m}{1}\PYG{+w}{ }\PYG{l+m}{1.0}\PYG{+w}{ }\PYG{l+m}{100.0}\PYG{+w}{ }\PYG{l+m}{80.0}\PYG{+w}{ }\PYG{l+m}{80.0}
\end{sphinxVerbatim}


\subsection{Description}
\label{\detokenize{bond_class2:description}}
\sphinxAtStartPar
The \sphinxstyleemphasis{class2} bond style uses the potential
\begin{equation*}
\begin{split}E = K_2 (r - r_0)^2 + K_3 (r - r_0)^3 + K_4 (r - r_0)^4\end{split}
\end{equation*}
\sphinxAtStartPar
where \(r_0\) is the equilibrium bond distance.

\sphinxAtStartPar
See {\hyperref[\detokenize{bond_class2:bond-sun}]{\sphinxcrossref{\DUrole{std,std-ref}{(Sun)}}}} for a description of the COMPASS class2 force field.

\sphinxAtStartPar
The following coefficients must be defined for each bond type via the
{\hyperref[\detokenize{bond_coeff::doc}]{\sphinxcrossref{\DUrole{doc}{bond\_coeff}}}} command as in the example above, or in
the data file or restart files read by the {\hyperref[\detokenize{read_data::doc}]{\sphinxcrossref{\DUrole{doc}{read\_data}}}}
or {\hyperref[\detokenize{read_restart::doc}]{\sphinxcrossref{\DUrole{doc}{read\_restart}}}} commands:
\begin{itemize}
\item {} 
\sphinxAtStartPar
\(r_0\) (distance)

\item {} 
\sphinxAtStartPar
\(K_2\) (energy/distance\textasciicircum{}2)

\item {} 
\sphinxAtStartPar
\(K_3\) (energy/distance\textasciicircum{}3)

\item {} 
\sphinxAtStartPar
\(K_4\) (energy/distance\textasciicircum{}4)

\end{itemize}


\bigskip\hrule\bigskip


\sphinxAtStartPar
Styles with a \sphinxstyleemphasis{gpu}, \sphinxstyleemphasis{intel}, \sphinxstyleemphasis{kk}, \sphinxstyleemphasis{omp}, or \sphinxstyleemphasis{opt} suffix are
functionally the same as the corresponding style without the suffix.
They have been optimized to run faster, depending on your available
hardware, as discussed on the {\hyperref[\detokenize{Speed_packages::doc}]{\sphinxcrossref{\DUrole{doc}{Accelerator packages}}}}
page.  The accelerated styles take the same arguments and should
produce the same results, except for round\sphinxhyphen{}off and precision issues.

\sphinxAtStartPar
These accelerated styles are part of the GPU, INTEL, KOKKOS,
OPENMP, and OPT packages, respectively.  They are only enabled if
LAMMPS was built with those packages.  See the {\hyperref[\detokenize{Build_package::doc}]{\sphinxcrossref{\DUrole{doc}{Build package}}}} page for more info.

\sphinxAtStartPar
You can specify the accelerated styles explicitly in your input script
by including their suffix, or you can use the {\hyperref[\detokenize{Run_options::doc}]{\sphinxcrossref{\DUrole{doc}{\sphinxhyphen{}suffix command\sphinxhyphen{}line switch}}}} when you invoke LAMMPS, or you can use the
{\hyperref[\detokenize{suffix::doc}]{\sphinxcrossref{\DUrole{doc}{suffix}}}} command in your input script.

\sphinxAtStartPar
See the {\hyperref[\detokenize{Speed_packages::doc}]{\sphinxcrossref{\DUrole{doc}{Accelerator packages}}}} page for more
instructions on how to use the accelerated styles effectively.


\bigskip\hrule\bigskip



\subsection{Restrictions}
\label{\detokenize{bond_class2:restrictions}}
\sphinxAtStartPar
This bond style can only be used if LAMMPS was built with the CLASS2
package.  See the {\hyperref[\detokenize{Build_package::doc}]{\sphinxcrossref{\DUrole{doc}{Build package}}}} page for more
info.


\subsection{Related commands}
\label{\detokenize{bond_class2:related-commands}}
\sphinxAtStartPar
{\hyperref[\detokenize{bond_coeff::doc}]{\sphinxcrossref{\DUrole{doc}{bond\_coeff}}}}, {\hyperref[\detokenize{delete_bonds::doc}]{\sphinxcrossref{\DUrole{doc}{delete\_bonds}}}}


\subsection{Default}
\label{\detokenize{bond_class2:default}}
\sphinxAtStartPar
none


\bigskip\hrule\bigskip

\phantomsection\label{\detokenize{bond_class2:bond-sun}}
\sphinxAtStartPar
\sphinxstylestrong{(Sun)} Sun, J Phys Chem B 102, 7338\sphinxhyphen{}7364 (1998).

\sphinxstepscope

\index{bond\_style fene@\spxentry{bond\_style fene}}\ignorespaces 
\index{bond\_style fene/nm@\spxentry{bond\_style fene/nm}}\ignorespaces 
\index{bond\_style fene/intel@\spxentry{bond\_style fene/intel}}\ignorespaces 
\index{bond\_style fene/kk@\spxentry{bond\_style fene/kk}}\ignorespaces 
\index{bond\_style fene/omp@\spxentry{bond\_style fene/omp}}\ignorespaces 

\section{bond\_style fene command}
\label{\detokenize{bond_fene:bond-style-fene-command}}\label{\detokenize{bond_fene:index-4}}\label{\detokenize{bond_fene:index-3}}\label{\detokenize{bond_fene:index-2}}\label{\detokenize{bond_fene:index-1}}\label{\detokenize{bond_fene:index-0}}\label{\detokenize{bond_fene::doc}}
\sphinxAtStartPar
Accelerator Variants: \sphinxstyleemphasis{fene/intel}, \sphinxstyleemphasis{fene/kk}, \sphinxstyleemphasis{fene/omp}


\section{bond\_style fene/nm command}
\label{\detokenize{bond_fene:bond-style-fene-nm-command}}

\subsection{Syntax}
\label{\detokenize{bond_fene:syntax}}
\begin{sphinxVerbatim}[commandchars=\\\{\}]
\PYG{k}{bond\PYGZus{}style}\PYG{+w}{ }\PYG{n}{fene}
\PYG{k}{bond\PYGZus{}style}\PYG{+w}{ }\PYG{n}{fene}\PYG{o}{/}\PYG{n}{nm}
\end{sphinxVerbatim}


\subsection{Examples}
\label{\detokenize{bond_fene:examples}}
\begin{sphinxVerbatim}[commandchars=\\\{\}]
\PYG{k}{bond\PYGZus{}style}\PYG{+w}{ }\PYG{n}{fene}
\PYG{k}{bond\PYGZus{}coeff}\PYG{+w}{ }\PYG{l+m}{1}\PYG{+w}{ }\PYG{l+m}{30.0}\PYG{+w}{ }\PYG{l+m}{1.5}\PYG{+w}{ }\PYG{l+m}{1.0}\PYG{+w}{ }\PYG{l+m}{1.0}

\PYG{k}{bond\PYGZus{}style}\PYG{+w}{ }\PYG{n}{fene}\PYG{o}{/}\PYG{n}{nm}
\PYG{k}{bond\PYGZus{}coeff}\PYG{+w}{ }\PYG{l+m}{1}\PYG{+w}{ }\PYG{l+m}{2.25344}\PYG{+w}{ }\PYG{l+m}{1.5}\PYG{+w}{ }\PYG{l+m}{1.0}\PYG{+w}{ }\PYG{l+m}{1.12246}\PYG{+w}{ }\PYG{l+m}{2}\PYG{+w}{ }\PYG{l+m}{6}
\end{sphinxVerbatim}


\subsection{Description}
\label{\detokenize{bond_fene:description}}
\sphinxAtStartPar
The \sphinxstyleemphasis{fene} bond style uses the potential
\begin{equation*}
\begin{split}E = -0.5 K R_0^2  \ln \left[ 1 - \left(\frac{r}{R_0}\right)^2\right] + 4 \epsilon \left[ \left(\frac{\sigma}{r}\right)^{12} - \left(\frac{\sigma}{r}\right)^6 \right] + \epsilon\end{split}
\end{equation*}
\sphinxAtStartPar
to define a finite extensible nonlinear elastic (FENE) potential
{\hyperref[\detokenize{bond_fene:fene-kremer}]{\sphinxcrossref{\DUrole{std,std-ref}{(Kremer)}}}}, used for bead\sphinxhyphen{}spring polymer models.  The first
term is attractive, the second Lennard\sphinxhyphen{}Jones term is repulsive.  The
first term extends to \(R_0\), the maximum extent of the bond.  The second
term is cutoff at \(2^\frac{1}{6} \sigma\), the minimum of the LJ potential.

\sphinxAtStartPar
The \sphinxstyleemphasis{fene/nm} bond style substitutes the standard LJ potential with the generalized LJ potential
in the same form as in pair style {\hyperref[\detokenize{pair_nm::doc}]{\sphinxcrossref{\DUrole{doc}{nm/cut}}}}. The bond energy is then given by
\begin{equation*}
\begin{split}E = -0.5 K R_0^2  \ln \left[ 1 - \left(\frac{r}{R_0}\right)^2\right] + \frac{E_0}{(n-m)} \left[ m \left(\frac{r_0}{r}\right)^n - n \left(\frac{r_0}{r}\right)^m \right]\end{split}
\end{equation*}
\sphinxAtStartPar
Similar to the \sphinxstyleemphasis{fene} style, the generalized Lennard\sphinxhyphen{}Jones is cut off at
the potential minimum, \(r_0\), to be repulsive only.  The following
coefficients must be defined for each bond type via the {\hyperref[\detokenize{bond_coeff::doc}]{\sphinxcrossref{\DUrole{doc}{bond\_coeff}}}} command as in the example above, or in the data file or
restart files read by the {\hyperref[\detokenize{read_data::doc}]{\sphinxcrossref{\DUrole{doc}{read\_data}}}} or
{\hyperref[\detokenize{read_restart::doc}]{\sphinxcrossref{\DUrole{doc}{read\_restart}}}} commands:
\begin{itemize}
\item {} 
\sphinxAtStartPar
\(K\) (energy/distance\textasciicircum{}2)

\item {} 
\sphinxAtStartPar
\(R_0\) (distance)

\item {} 
\sphinxAtStartPar
\(\epsilon\) (energy)

\item {} 
\sphinxAtStartPar
\(\sigma\) (distance)

\end{itemize}

\sphinxAtStartPar
For the \sphinxstyleemphasis{fene/nm} style, the following coefficients are used.  Please
note, that the standard LJ potential and thus the regular FENE potential
is recovered for (n=12 m=6) and \(r_0 = 2^\frac{1}{6} \sigma\).
\begin{itemize}
\item {} 
\sphinxAtStartPar
\(K\) (energy/distance\textasciicircum{}2)

\item {} 
\sphinxAtStartPar
\(R_0\) (distance)

\item {} 
\sphinxAtStartPar
\(E_0\) (energy)

\item {} 
\sphinxAtStartPar
\(r_0\) (distance)

\item {} 
\sphinxAtStartPar
\(n\) (unitless)

\item {} 
\sphinxAtStartPar
\(m\) (unitless)

\end{itemize}


\bigskip\hrule\bigskip


\sphinxAtStartPar
Styles with a \sphinxstyleemphasis{gpu}, \sphinxstyleemphasis{intel}, \sphinxstyleemphasis{kk}, \sphinxstyleemphasis{omp}, or \sphinxstyleemphasis{opt} suffix are
functionally the same as the corresponding style without the suffix.
They have been optimized to run faster, depending on your available
hardware, as discussed on the {\hyperref[\detokenize{Speed_packages::doc}]{\sphinxcrossref{\DUrole{doc}{Accelerator packages}}}}
page.  The accelerated styles take the same arguments and should
produce the same results, except for round\sphinxhyphen{}off and precision issues.

\sphinxAtStartPar
These accelerated styles are part of the GPU, INTEL, KOKKOS,
OPENMP, and OPT packages, respectively.  They are only enabled if
LAMMPS was built with those packages.  See the {\hyperref[\detokenize{Build_package::doc}]{\sphinxcrossref{\DUrole{doc}{Build package}}}} page for more info.

\sphinxAtStartPar
You can specify the accelerated styles explicitly in your input script
by including their suffix, or you can use the {\hyperref[\detokenize{Run_options::doc}]{\sphinxcrossref{\DUrole{doc}{\sphinxhyphen{}suffix command\sphinxhyphen{}line switch}}}} when you invoke LAMMPS, or you can use the
{\hyperref[\detokenize{suffix::doc}]{\sphinxcrossref{\DUrole{doc}{suffix}}}} command in your input script.

\sphinxAtStartPar
See the {\hyperref[\detokenize{Speed_packages::doc}]{\sphinxcrossref{\DUrole{doc}{Accelerator packages}}}} page for more
instructions on how to use the accelerated styles effectively.


\bigskip\hrule\bigskip



\subsection{Restrictions}
\label{\detokenize{bond_fene:restrictions}}
\sphinxAtStartPar
The \sphinxstyleemphasis{fene} bond style can only be used if LAMMPS was built with the MOLECULE
package; the \sphinxstyleemphasis{fene/nm} bond style can only be used if LAMMPS was built
with the EXTRA\sphinxhyphen{}MOLECULE package. See the {\hyperref[\detokenize{Build_package::doc}]{\sphinxcrossref{\DUrole{doc}{Build package}}}}
page for more info.

\sphinxAtStartPar
You typically should specify {\hyperref[\detokenize{special_bonds::doc}]{\sphinxcrossref{\DUrole{doc}{special\_bonds fene}}}}
or {\hyperref[\detokenize{special_bonds::doc}]{\sphinxcrossref{\DUrole{doc}{special\_bonds lj/coul 0 1 1}}}} to use this bond
style.  LAMMPS will issue a warning it that’s not the case.


\subsection{Related commands}
\label{\detokenize{bond_fene:related-commands}}
\sphinxAtStartPar
{\hyperref[\detokenize{bond_coeff::doc}]{\sphinxcrossref{\DUrole{doc}{bond\_coeff}}}}, {\hyperref[\detokenize{delete_bonds::doc}]{\sphinxcrossref{\DUrole{doc}{delete\_bonds}}}},
{\hyperref[\detokenize{pair_lj::doc}]{\sphinxcrossref{\DUrole{doc}{pair style lj/cut}}}}, {\hyperref[\detokenize{pair_nm::doc}]{\sphinxcrossref{\DUrole{doc}{pair style nm/cut}}}}.


\subsection{Default}
\label{\detokenize{bond_fene:default}}
\sphinxAtStartPar
none


\bigskip\hrule\bigskip

\phantomsection\label{\detokenize{bond_fene:fene-kremer}}
\sphinxAtStartPar
\sphinxstylestrong{(Kremer)} Kremer, Grest, J Chem Phys, 92, 5057 (1990).

\sphinxstepscope

\index{bond\_style fene/expand@\spxentry{bond\_style fene/expand}}\ignorespaces 
\index{bond\_style fene/expand/omp@\spxentry{bond\_style fene/expand/omp}}\ignorespaces 

\section{bond\_style fene/expand command}
\label{\detokenize{bond_fene_expand:bond-style-fene-expand-command}}\label{\detokenize{bond_fene_expand:index-1}}\label{\detokenize{bond_fene_expand:index-0}}\label{\detokenize{bond_fene_expand::doc}}
\sphinxAtStartPar
Accelerator Variants: \sphinxstyleemphasis{fene/expand/omp}


\subsection{Syntax}
\label{\detokenize{bond_fene_expand:syntax}}
\begin{sphinxVerbatim}[commandchars=\\\{\}]
\PYG{k}{bond\PYGZus{}style}\PYG{+w}{ }\PYG{n}{fene}\PYG{o}{/}\PYG{n}{expand}
\end{sphinxVerbatim}


\subsection{Examples}
\label{\detokenize{bond_fene_expand:examples}}
\begin{sphinxVerbatim}[commandchars=\\\{\}]
\PYG{k}{bond\PYGZus{}style}\PYG{+w}{ }\PYG{n}{fene}\PYG{o}{/}\PYG{n}{expand}
\PYG{k}{bond\PYGZus{}coeff}\PYG{+w}{ }\PYG{l+m}{1}\PYG{+w}{ }\PYG{l+m}{30.0}\PYG{+w}{ }\PYG{l+m}{1.5}\PYG{+w}{ }\PYG{l+m}{1.0}\PYG{+w}{ }\PYG{l+m}{1.0}\PYG{+w}{ }\PYG{l+m}{0.5}
\end{sphinxVerbatim}


\subsection{Description}
\label{\detokenize{bond_fene_expand:description}}
\sphinxAtStartPar
The \sphinxstyleemphasis{fene/expand} bond style uses the potential
\begin{equation*}
\begin{split}E = -0.5 K R_0^2 \ln \left[1 -\left( \frac{\left(r - \Delta\right)}{R_0}\right)^2 \right] + 4 \epsilon \left[ \left(\frac{\sigma}{\left(r - \Delta\right)}\right)^{12} - \left(\frac{\sigma}{\left(r - \Delta\right)}\right)^6 \right] + \epsilon\end{split}
\end{equation*}
\sphinxAtStartPar
to define a finite extensible nonlinear elastic (FENE) potential
{\hyperref[\detokenize{bond_fene_expand:feneexpand-kremer}]{\sphinxcrossref{\DUrole{std,std-ref}{(Kremer)}}}}, used for bead\sphinxhyphen{}spring polymer models.  The first
term is attractive, the second Lennard\sphinxhyphen{}Jones term is repulsive.

\sphinxAtStartPar
The \sphinxstyleemphasis{fene/expand} bond style is similar to \sphinxstyleemphasis{fene} except that an extra
shift factor of \(\Delta\) (positive or negative) is added to \(r\) to
effectively change the bead size of the bonded atoms.  The first term
now extends to \(R_0 + \Delta\) and the second term is cutoff at \(2^\frac{1}{6} \sigma + \Delta\).

\sphinxAtStartPar
The following coefficients must be defined for each bond type via the
{\hyperref[\detokenize{bond_coeff::doc}]{\sphinxcrossref{\DUrole{doc}{bond\_coeff}}}} command as in the example above, or in
the data file or restart files read by the {\hyperref[\detokenize{read_data::doc}]{\sphinxcrossref{\DUrole{doc}{read\_data}}}}
or {\hyperref[\detokenize{read_restart::doc}]{\sphinxcrossref{\DUrole{doc}{read\_restart}}}} commands:
\begin{itemize}
\item {} 
\sphinxAtStartPar
\(K\) (energy/distance\textasciicircum{}2)

\item {} 
\sphinxAtStartPar
\(R_0\) (distance)

\item {} 
\sphinxAtStartPar
\(\epsilon\) (energy)

\item {} 
\sphinxAtStartPar
\(\sigma\) (distance)

\item {} 
\sphinxAtStartPar
\(\Delta\) (distance)

\end{itemize}


\bigskip\hrule\bigskip


\sphinxAtStartPar
Styles with a \sphinxstyleemphasis{gpu}, \sphinxstyleemphasis{intel}, \sphinxstyleemphasis{kk}, \sphinxstyleemphasis{omp}, or \sphinxstyleemphasis{opt} suffix are
functionally the same as the corresponding style without the suffix.
They have been optimized to run faster, depending on your available
hardware, as discussed on the {\hyperref[\detokenize{Speed_packages::doc}]{\sphinxcrossref{\DUrole{doc}{Accelerator packages}}}}
page.  The accelerated styles take the same arguments and should
produce the same results, except for round\sphinxhyphen{}off and precision issues.

\sphinxAtStartPar
These accelerated styles are part of the GPU, INTEL, KOKKOS,
OPENMP, and OPT packages, respectively.  They are only enabled if
LAMMPS was built with those packages.  See the {\hyperref[\detokenize{Build_package::doc}]{\sphinxcrossref{\DUrole{doc}{Build package}}}} page for more info.

\sphinxAtStartPar
You can specify the accelerated styles explicitly in your input script
by including their suffix, or you can use the {\hyperref[\detokenize{Run_options::doc}]{\sphinxcrossref{\DUrole{doc}{\sphinxhyphen{}suffix command\sphinxhyphen{}line switch}}}} when you invoke LAMMPS, or you can use the
{\hyperref[\detokenize{suffix::doc}]{\sphinxcrossref{\DUrole{doc}{suffix}}}} command in your input script.

\sphinxAtStartPar
See the {\hyperref[\detokenize{Speed_packages::doc}]{\sphinxcrossref{\DUrole{doc}{Accelerator packages}}}} page for more
instructions on how to use the accelerated styles effectively.


\bigskip\hrule\bigskip



\subsection{Restrictions}
\label{\detokenize{bond_fene_expand:restrictions}}
\sphinxAtStartPar
This bond style can only be used if LAMMPS was built with the MOLECULE
package.  See the {\hyperref[\detokenize{Build_package::doc}]{\sphinxcrossref{\DUrole{doc}{Build package}}}} page for more
info.

\sphinxAtStartPar
You typically should specify {\hyperref[\detokenize{special_bonds::doc}]{\sphinxcrossref{\DUrole{doc}{special\_bonds fene}}}}
or {\hyperref[\detokenize{special_bonds::doc}]{\sphinxcrossref{\DUrole{doc}{special\_bonds lj/coul 0 1 1}}}} to use this bond
style.  LAMMPS will issue a warning it that’s not the case.


\subsection{Related commands}
\label{\detokenize{bond_fene_expand:related-commands}}
\sphinxAtStartPar
{\hyperref[\detokenize{bond_coeff::doc}]{\sphinxcrossref{\DUrole{doc}{bond\_coeff}}}}, {\hyperref[\detokenize{delete_bonds::doc}]{\sphinxcrossref{\DUrole{doc}{delete\_bonds}}}}


\subsection{Default}
\label{\detokenize{bond_fene_expand:default}}
\sphinxAtStartPar
none


\bigskip\hrule\bigskip

\phantomsection\label{\detokenize{bond_fene_expand:feneexpand-kremer}}
\sphinxAtStartPar
\sphinxstylestrong{(Kremer)} Kremer, Grest, J Chem Phys, 92, 5057 (1990).

\sphinxstepscope

\index{bond\_style gaussian@\spxentry{bond\_style gaussian}}\ignorespaces 

\section{bond\_style gaussian command}
\label{\detokenize{bond_gaussian:bond-style-gaussian-command}}\label{\detokenize{bond_gaussian:index-0}}\label{\detokenize{bond_gaussian::doc}}

\subsection{Syntax}
\label{\detokenize{bond_gaussian:syntax}}
\begin{sphinxVerbatim}[commandchars=\\\{\}]
\PYG{k}{bond\PYGZus{}style}\PYG{+w}{ }\PYG{n}{gaussian}
\end{sphinxVerbatim}


\subsection{Examples}
\label{\detokenize{bond_gaussian:examples}}
\begin{sphinxVerbatim}[commandchars=\\\{\}]
\PYG{k}{bond\PYGZus{}style}\PYG{+w}{ }\PYG{n}{gaussian}
\PYG{k}{bond\PYGZus{}coeff}\PYG{+w}{ }\PYG{l+m}{1}\PYG{+w}{ }\PYG{l+m}{300.0}\PYG{+w}{ }\PYG{l+m}{2}\PYG{+w}{ }\PYG{l+m}{0.0128}\PYG{+w}{ }\PYG{l+m}{0.375}\PYG{+w}{ }\PYG{l+m}{3.37}\PYG{+w}{ }\PYG{l+m}{0.0730}\PYG{+w}{ }\PYG{l+m}{0.148}\PYG{+w}{ }\PYG{l+m}{3.63}
\end{sphinxVerbatim}


\subsection{Description}
\label{\detokenize{bond_gaussian:description}}
\sphinxAtStartPar
The \sphinxstyleemphasis{gaussian} bond style uses the potential:
\begin{equation*}
\begin{split}E = -k_B T ln\left(\sum_{i=1}^{n} \frac{A_i}{w_i \sqrt{\pi/2}} exp\left( \frac{-2(r-r_{i})^2}{w_i^2}\right)\right)\end{split}
\end{equation*}
\sphinxAtStartPar
This analytical form is a suitable potential for obtaining mesoscale
effective force fields which can reproduce target atomistic
distributions {\hyperref[\detokenize{bond_gaussian:milano0}]{\sphinxcrossref{\DUrole{std,std-ref}{(Milano)}}}}

\sphinxAtStartPar
The following coefficients must be defined for each bond type via the
{\hyperref[\detokenize{bond_coeff::doc}]{\sphinxcrossref{\DUrole{doc}{bond\_coeff}}}} command as in the example above, or in
the data file or restart files read by the {\hyperref[\detokenize{read_data::doc}]{\sphinxcrossref{\DUrole{doc}{read\_data}}}}
or {\hyperref[\detokenize{read_restart::doc}]{\sphinxcrossref{\DUrole{doc}{read\_restart}}}} commands:
\begin{itemize}
\item {} 
\sphinxAtStartPar
\(T\) temperature at which the potential was derived

\item {} 
\sphinxAtStartPar
\(n\) (integer \textgreater{}=1)

\item {} 
\sphinxAtStartPar
\(A_1\) (\textgreater{} 0, distance)

\item {} 
\sphinxAtStartPar
\(w_1\) (\textgreater{} 0, distance)

\item {} 
\sphinxAtStartPar
\(r_1\) (\textgreater{}= 0, distance)

\item {} 
\sphinxAtStartPar
…

\item {} 
\sphinxAtStartPar
\(A_n\) (\textgreater{} 0, distance)

\item {} 
\sphinxAtStartPar
\(w_n\) (\textgreater{} 0, distance)

\item {} 
\sphinxAtStartPar
\(r_n\) (\textgreater{}= 0, distance)

\end{itemize}


\subsection{Restrictions}
\label{\detokenize{bond_gaussian:restrictions}}
\sphinxAtStartPar
This bond style can only be used if LAMMPS was built with the
EXTRA\sphinxhyphen{}MOLECULE package.  See the {\hyperref[\detokenize{Build_package::doc}]{\sphinxcrossref{\DUrole{doc}{Build package}}}}
doc page for more info.


\subsection{Related commands}
\label{\detokenize{bond_gaussian:related-commands}}
\sphinxAtStartPar
{\hyperref[\detokenize{bond_coeff::doc}]{\sphinxcrossref{\DUrole{doc}{bond\_coeff}}}}


\subsection{Default}
\label{\detokenize{bond_gaussian:default}}
\sphinxAtStartPar
none


\bigskip\hrule\bigskip

\phantomsection\label{\detokenize{bond_gaussian:milano0}}
\sphinxAtStartPar
\sphinxstylestrong{(Milano)} G. Milano, S. Goudeau, F. Mueller\sphinxhyphen{}Plathe, J. Polym. Sci. B Polym. Phys. 43, 871 (2005).

\sphinxstepscope

\index{bond\_style gromos@\spxentry{bond\_style gromos}}\ignorespaces 
\index{bond\_style gromos/omp@\spxentry{bond\_style gromos/omp}}\ignorespaces 

\section{bond\_style gromos command}
\label{\detokenize{bond_gromos:bond-style-gromos-command}}\label{\detokenize{bond_gromos:index-1}}\label{\detokenize{bond_gromos:index-0}}\label{\detokenize{bond_gromos::doc}}
\sphinxAtStartPar
Accelerator Variants: \sphinxstyleemphasis{gromos/omp}


\subsection{Syntax}
\label{\detokenize{bond_gromos:syntax}}
\begin{sphinxVerbatim}[commandchars=\\\{\}]
\PYG{k}{bond\PYGZus{}style}\PYG{+w}{ }\PYG{n}{gromos}
\end{sphinxVerbatim}


\subsection{Examples}
\label{\detokenize{bond_gromos:examples}}
\begin{sphinxVerbatim}[commandchars=\\\{\}]
\PYG{k}{bond\PYGZus{}style}\PYG{+w}{ }\PYG{n}{gromos}
\PYG{k}{bond\PYGZus{}coeff}\PYG{+w}{ }\PYG{l+m}{5}\PYG{+w}{ }\PYG{l+m}{80.0}\PYG{+w}{ }\PYG{l+m}{1.2}
\end{sphinxVerbatim}


\subsection{Description}
\label{\detokenize{bond_gromos:description}}
\sphinxAtStartPar
The \sphinxstyleemphasis{gromos} bond style uses the potential
\begin{equation*}
\begin{split}E = K (r^2 - r_0^2)^2\end{split}
\end{equation*}
\sphinxAtStartPar
where \(r_0\) is the equilibrium bond distance.  Note that the usual 1/4
factor is included in \(K\).

\sphinxAtStartPar
The following coefficients must be defined for each bond type via the
{\hyperref[\detokenize{bond_coeff::doc}]{\sphinxcrossref{\DUrole{doc}{bond\_coeff}}}} command as in the example above, or in
the data file or restart files read by the {\hyperref[\detokenize{read_data::doc}]{\sphinxcrossref{\DUrole{doc}{read\_data}}}}
or {\hyperref[\detokenize{read_restart::doc}]{\sphinxcrossref{\DUrole{doc}{read\_restart}}}} commands:
\begin{itemize}
\item {} 
\sphinxAtStartPar
\(K\) (energy/distance\textasciicircum{}4)

\item {} 
\sphinxAtStartPar
\(r_0\) (distance)

\end{itemize}


\bigskip\hrule\bigskip


\sphinxAtStartPar
Styles with a \sphinxstyleemphasis{gpu}, \sphinxstyleemphasis{intel}, \sphinxstyleemphasis{kk}, \sphinxstyleemphasis{omp}, or \sphinxstyleemphasis{opt} suffix are
functionally the same as the corresponding style without the suffix.
They have been optimized to run faster, depending on your available
hardware, as discussed on the {\hyperref[\detokenize{Speed_packages::doc}]{\sphinxcrossref{\DUrole{doc}{Accelerator packages}}}}
page.  The accelerated styles take the same arguments and should
produce the same results, except for round\sphinxhyphen{}off and precision issues.

\sphinxAtStartPar
These accelerated styles are part of the GPU, INTEL, KOKKOS,
OPENMP, and OPT packages, respectively.  They are only enabled if
LAMMPS was built with those packages.  See the {\hyperref[\detokenize{Build_package::doc}]{\sphinxcrossref{\DUrole{doc}{Build package}}}} page for more info.

\sphinxAtStartPar
You can specify the accelerated styles explicitly in your input script
by including their suffix, or you can use the {\hyperref[\detokenize{Run_options::doc}]{\sphinxcrossref{\DUrole{doc}{\sphinxhyphen{}suffix command\sphinxhyphen{}line switch}}}} when you invoke LAMMPS, or you can use the
{\hyperref[\detokenize{suffix::doc}]{\sphinxcrossref{\DUrole{doc}{suffix}}}} command in your input script.

\sphinxAtStartPar
See the {\hyperref[\detokenize{Speed_packages::doc}]{\sphinxcrossref{\DUrole{doc}{Accelerator packages}}}} page for more
instructions on how to use the accelerated styles effectively.


\bigskip\hrule\bigskip



\subsection{Restrictions}
\label{\detokenize{bond_gromos:restrictions}}
\sphinxAtStartPar
This bond style can only be used if LAMMPS was built with the MOLECULE
package.  See the {\hyperref[\detokenize{Build_package::doc}]{\sphinxcrossref{\DUrole{doc}{Build package}}}} page for more
info.


\subsection{Related commands}
\label{\detokenize{bond_gromos:related-commands}}
\sphinxAtStartPar
{\hyperref[\detokenize{bond_coeff::doc}]{\sphinxcrossref{\DUrole{doc}{bond\_coeff}}}}, {\hyperref[\detokenize{delete_bonds::doc}]{\sphinxcrossref{\DUrole{doc}{delete\_bonds}}}}


\subsection{Default}
\label{\detokenize{bond_gromos:default}}
\sphinxAtStartPar
none

\sphinxstepscope

\index{bond\_style harmonic@\spxentry{bond\_style harmonic}}\ignorespaces 
\index{bond\_style harmonic/intel@\spxentry{bond\_style harmonic/intel}}\ignorespaces 
\index{bond\_style harmonic/kk@\spxentry{bond\_style harmonic/kk}}\ignorespaces 
\index{bond\_style harmonic/omp@\spxentry{bond\_style harmonic/omp}}\ignorespaces 

\section{bond\_style harmonic command}
\label{\detokenize{bond_harmonic:bond-style-harmonic-command}}\label{\detokenize{bond_harmonic:index-3}}\label{\detokenize{bond_harmonic:index-2}}\label{\detokenize{bond_harmonic:index-1}}\label{\detokenize{bond_harmonic:index-0}}\label{\detokenize{bond_harmonic::doc}}
\sphinxAtStartPar
Accelerator Variants: \sphinxstyleemphasis{harmonic/intel}, \sphinxstyleemphasis{harmonic/kk}, \sphinxstyleemphasis{harmonic/omp}


\subsection{Syntax}
\label{\detokenize{bond_harmonic:syntax}}
\begin{sphinxVerbatim}[commandchars=\\\{\}]
\PYG{k}{bond\PYGZus{}style}\PYG{+w}{ }\PYG{n}{harmonic}
\end{sphinxVerbatim}


\subsection{Examples}
\label{\detokenize{bond_harmonic:examples}}
\begin{sphinxVerbatim}[commandchars=\\\{\}]
\PYG{k}{bond\PYGZus{}style}\PYG{+w}{ }\PYG{n}{harmonic}
\PYG{k}{bond\PYGZus{}coeff}\PYG{+w}{ }\PYG{l+m}{5}\PYG{+w}{ }\PYG{l+m}{80.0}\PYG{+w}{ }\PYG{l+m}{1.2}
\end{sphinxVerbatim}


\subsection{Description}
\label{\detokenize{bond_harmonic:description}}
\sphinxAtStartPar
The \sphinxstyleemphasis{harmonic} bond style uses the potential
\begin{equation*}
\begin{split}E = K (r - r_0)^2\end{split}
\end{equation*}
\sphinxAtStartPar
where \(r_0\) is the equilibrium bond distance.  Note that the usual 1/2
factor is included in \(K\).

\sphinxAtStartPar
The following coefficients must be defined for each bond type via the
{\hyperref[\detokenize{bond_coeff::doc}]{\sphinxcrossref{\DUrole{doc}{bond\_coeff}}}} command as in the example above, or in
the data file or restart files read by the {\hyperref[\detokenize{read_data::doc}]{\sphinxcrossref{\DUrole{doc}{read\_data}}}}
or {\hyperref[\detokenize{read_restart::doc}]{\sphinxcrossref{\DUrole{doc}{read\_restart}}}} commands:
\begin{itemize}
\item {} 
\sphinxAtStartPar
\(K\) (energy/distance\textasciicircum{}2)

\item {} 
\sphinxAtStartPar
\(r_0\) (distance)

\end{itemize}


\bigskip\hrule\bigskip


\sphinxAtStartPar
Styles with a \sphinxstyleemphasis{gpu}, \sphinxstyleemphasis{intel}, \sphinxstyleemphasis{kk}, \sphinxstyleemphasis{omp}, or \sphinxstyleemphasis{opt} suffix are
functionally the same as the corresponding style without the suffix.
They have been optimized to run faster, depending on your available
hardware, as discussed on the {\hyperref[\detokenize{Speed_packages::doc}]{\sphinxcrossref{\DUrole{doc}{Accelerator packages}}}}
page.  The accelerated styles take the same arguments and should
produce the same results, except for round\sphinxhyphen{}off and precision issues.

\sphinxAtStartPar
These accelerated styles are part of the GPU, INTEL, KOKKOS,
OPENMP, and OPT packages, respectively.  They are only enabled if
LAMMPS was built with those packages.  See the {\hyperref[\detokenize{Build_package::doc}]{\sphinxcrossref{\DUrole{doc}{Build package}}}} page for more info.

\sphinxAtStartPar
You can specify the accelerated styles explicitly in your input script
by including their suffix, or you can use the {\hyperref[\detokenize{Run_options::doc}]{\sphinxcrossref{\DUrole{doc}{\sphinxhyphen{}suffix command\sphinxhyphen{}line switch}}}} when you invoke LAMMPS, or you can use the
{\hyperref[\detokenize{suffix::doc}]{\sphinxcrossref{\DUrole{doc}{suffix}}}} command in your input script.

\sphinxAtStartPar
See the {\hyperref[\detokenize{Speed_packages::doc}]{\sphinxcrossref{\DUrole{doc}{Accelerator packages}}}} page for more
instructions on how to use the accelerated styles effectively.


\bigskip\hrule\bigskip



\subsection{Restrictions}
\label{\detokenize{bond_harmonic:restrictions}}
\sphinxAtStartPar
This bond style can only be used if LAMMPS was built with the MOLECULE
package.  See the {\hyperref[\detokenize{Build_package::doc}]{\sphinxcrossref{\DUrole{doc}{Build package}}}} page for more
info.


\subsection{Related commands}
\label{\detokenize{bond_harmonic:related-commands}}
\sphinxAtStartPar
{\hyperref[\detokenize{bond_coeff::doc}]{\sphinxcrossref{\DUrole{doc}{bond\_coeff}}}}, {\hyperref[\detokenize{delete_bonds::doc}]{\sphinxcrossref{\DUrole{doc}{delete\_bonds}}}}


\subsection{Default}
\label{\detokenize{bond_harmonic:default}}
\sphinxAtStartPar
none

\sphinxstepscope

\index{bond\_style harmonic/restrain@\spxentry{bond\_style harmonic/restrain}}\ignorespaces 

\section{bond\_style harmonic/restrain command}
\label{\detokenize{bond_harmonic_restrain:bond-style-harmonic-restrain-command}}\label{\detokenize{bond_harmonic_restrain:index-0}}\label{\detokenize{bond_harmonic_restrain::doc}}

\subsection{Syntax}
\label{\detokenize{bond_harmonic_restrain:syntax}}
\begin{sphinxVerbatim}[commandchars=\\\{\}]
\PYG{k}{bond\PYGZus{}style}\PYG{+w}{ }\PYG{n}{harmonic}\PYG{o}{/}\PYG{n}{restrain}
\end{sphinxVerbatim}


\subsection{Examples}
\label{\detokenize{bond_harmonic_restrain:examples}}
\begin{sphinxVerbatim}[commandchars=\\\{\}]
\PYG{k}{bond\PYGZus{}style}\PYG{+w}{ }\PYG{n}{harmonic}
\PYG{k}{bond\PYGZus{}coeff}\PYG{+w}{ }\PYG{l+m}{5}\PYG{+w}{ }\PYG{l+m}{80.0}
\end{sphinxVerbatim}


\subsection{Description}
\label{\detokenize{bond_harmonic_restrain:description}}
\sphinxAtStartPar
\DUrole{versionmodified,added}{Added in version 28Mar2023.}

\sphinxAtStartPar
The \sphinxstyleemphasis{harmonic/restrain} bond style uses the potential
\begin{equation*}
\begin{split}E = K (r - r_{t=0})^2\end{split}
\end{equation*}
\sphinxAtStartPar
where \(r_{t=0}\) is the distance between the bonded atoms at the
beginning of the first {\hyperref[\detokenize{run::doc}]{\sphinxcrossref{\DUrole{doc}{run}}}} or {\hyperref[\detokenize{minimize::doc}]{\sphinxcrossref{\DUrole{doc}{minimize}}}}
command after the bond style has been defined (\sphinxstyleemphasis{t=0}).  Note that the
usual 1/2 factor is included in \(K\).  This will effectively
restrain bonds to their initial length, whatever that is.  This is where
this bond style differs from {\hyperref[\detokenize{bond_harmonic::doc}]{\sphinxcrossref{\DUrole{doc}{bond style harmonic}}}}
where the bond length is set through the per bond type coefficients.

\sphinxAtStartPar
The following coefficient must be defined for each bond type via the
{\hyperref[\detokenize{bond_coeff::doc}]{\sphinxcrossref{\DUrole{doc}{bond\_coeff}}}} command as in the example above, or in
the data file or restart files read by the {\hyperref[\detokenize{read_data::doc}]{\sphinxcrossref{\DUrole{doc}{read\_data}}}}
or {\hyperref[\detokenize{read_restart::doc}]{\sphinxcrossref{\DUrole{doc}{read\_restart}}}} commands
\begin{itemize}
\item {} 
\sphinxAtStartPar
\(K\) (energy/distance\textasciicircum{}2)

\end{itemize}

\sphinxAtStartPar
This bond style differs from other options to add harmonic restraints
like {\hyperref[\detokenize{fix_restrain::doc}]{\sphinxcrossref{\DUrole{doc}{fix restrain}}}} or {\hyperref[\detokenize{pair_list::doc}]{\sphinxcrossref{\DUrole{doc}{pair style list}}}} or {\hyperref[\detokenize{fix_colvars::doc}]{\sphinxcrossref{\DUrole{doc}{fix colvars}}}} in that it requires a
bond topology, and thus the defined bonds will trigger exclusion of
special neighbors from the neighbor list according to the
{\hyperref[\detokenize{special_bonds::doc}]{\sphinxcrossref{\DUrole{doc}{special\_bonds}}}} settings.


\subsection{Restart info}
\label{\detokenize{bond_harmonic_restrain:restart-info}}
\sphinxAtStartPar
This bond style supports the {\hyperref[\detokenize{write_restart::doc}]{\sphinxcrossref{\DUrole{doc}{write\_restart}}}} and
{\hyperref[\detokenize{read_restart::doc}]{\sphinxcrossref{\DUrole{doc}{read\_restart}}}} commands. The state of the initial
bond lengths is stored with restart files and read back.


\subsection{Restrictions}
\label{\detokenize{bond_harmonic_restrain:restrictions}}
\sphinxAtStartPar
This bond style can only be used if LAMMPS was built with the
EXTRA\sphinxhyphen{}MOLECULE package.  See the {\hyperref[\detokenize{Build_package::doc}]{\sphinxcrossref{\DUrole{doc}{Build package}}}}
page for more info.

\sphinxAtStartPar
This bond style maintains internal data to determine the original bond
lengths \(r_{t=0}\).  This information will be written to
{\hyperref[\detokenize{write_restart::doc}]{\sphinxcrossref{\DUrole{doc}{binary restart files}}}} but \sphinxstylestrong{not} to {\hyperref[\detokenize{write_data::doc}]{\sphinxcrossref{\DUrole{doc}{data
files}}}}.  Thus, continuing a simulation is \sphinxstyleemphasis{only} possible
with {\hyperref[\detokenize{read_restart::doc}]{\sphinxcrossref{\DUrole{doc}{read\_restart}}}}.  When using the {\hyperref[\detokenize{read_data::doc}]{\sphinxcrossref{\DUrole{doc}{read\_data
command}}}}, the reference bond lengths \(r_{t=0}\) will be
re\sphinxhyphen{}initialized from the current geometry.

\sphinxAtStartPar
This bond style cannot be used with {\hyperref[\detokenize{fix_shake::doc}]{\sphinxcrossref{\DUrole{doc}{fix shake or fix rattle}}}}, with {\hyperref[\detokenize{fix_filter_corotate::doc}]{\sphinxcrossref{\DUrole{doc}{fix filter/corotate}}}}, or
any {\hyperref[\detokenize{pair_lj_cut_tip4p::doc}]{\sphinxcrossref{\DUrole{doc}{tip4p pair style}}}} since there is no specific
equilibrium distance for a given bond type.


\subsection{Related commands}
\label{\detokenize{bond_harmonic_restrain:related-commands}}
\sphinxAtStartPar
{\hyperref[\detokenize{bond_coeff::doc}]{\sphinxcrossref{\DUrole{doc}{bond\_coeff}}}},  {\hyperref[\detokenize{bond_harmonic::doc}]{\sphinxcrossref{\DUrole{doc}{bond\_harmonic}}}},
{\hyperref[\detokenize{fix_restrain::doc}]{\sphinxcrossref{\DUrole{doc}{fix restrain}}}}, {\hyperref[\detokenize{pair_list::doc}]{\sphinxcrossref{\DUrole{doc}{pair style list}}}}


\subsection{Default}
\label{\detokenize{bond_harmonic_restrain:default}}
\sphinxAtStartPar
none

\sphinxstepscope

\index{bond\_style harmonic/shift@\spxentry{bond\_style harmonic/shift}}\ignorespaces 
\index{bond\_style harmonic/shift/omp@\spxentry{bond\_style harmonic/shift/omp}}\ignorespaces 

\section{bond\_style harmonic/shift command}
\label{\detokenize{bond_harmonic_shift:bond-style-harmonic-shift-command}}\label{\detokenize{bond_harmonic_shift:index-1}}\label{\detokenize{bond_harmonic_shift:index-0}}\label{\detokenize{bond_harmonic_shift::doc}}
\sphinxAtStartPar
Accelerator Variants: \sphinxstyleemphasis{harmonic/shift/omp}


\subsection{Syntax}
\label{\detokenize{bond_harmonic_shift:syntax}}
\begin{sphinxVerbatim}[commandchars=\\\{\}]
\PYG{k}{bond\PYGZus{}style}\PYG{+w}{ }\PYG{n}{harmonic}\PYG{o}{/}\PYG{n}{shift}
\end{sphinxVerbatim}


\subsection{Examples}
\label{\detokenize{bond_harmonic_shift:examples}}
\begin{sphinxVerbatim}[commandchars=\\\{\}]
\PYG{k}{bond\PYGZus{}style}\PYG{+w}{ }\PYG{n}{harmonic}\PYG{o}{/}\PYG{n}{shift}
\PYG{k}{bond\PYGZus{}coeff}\PYG{+w}{ }\PYG{l+m}{5}\PYG{+w}{ }\PYG{l+m}{10.0}\PYG{+w}{ }\PYG{l+m}{0.5}\PYG{+w}{ }\PYG{l+m}{1.0}
\end{sphinxVerbatim}


\subsection{Description}
\label{\detokenize{bond_harmonic_shift:description}}
\sphinxAtStartPar
The \sphinxstyleemphasis{harmonic/shift} bond style is a shifted harmonic bond that uses
the potential
\begin{equation*}
\begin{split}E = \frac{U_{\text{min}}}{(r_0-r_c)^2} \left[ (r-r_0)^2-(r_c-r_0)^2 \right]\end{split}
\end{equation*}
\sphinxAtStartPar
where \(r_0\) is the equilibrium bond distance, and \(r_c\) the critical distance.
The potential is \(-U_{\text{min}}\) at \(r0\) and zero at \(r_c\). The spring constant is
\(k = U_{\text{min}} / [ 2 (r_0-r_c)^2]\).

\sphinxAtStartPar
The following coefficients must be defined for each bond type via the
{\hyperref[\detokenize{bond_coeff::doc}]{\sphinxcrossref{\DUrole{doc}{bond\_coeff}}}} command as in the example above, or in
the data file or restart files read by the {\hyperref[\detokenize{read_data::doc}]{\sphinxcrossref{\DUrole{doc}{read\_data}}}}
or {\hyperref[\detokenize{read_restart::doc}]{\sphinxcrossref{\DUrole{doc}{read\_restart}}}} commands:
\begin{itemize}
\item {} 
\sphinxAtStartPar
\(U_{\text{min}}\) (energy)

\item {} 
\sphinxAtStartPar
\(r_0\) (distance)

\item {} 
\sphinxAtStartPar
\(r_c\) (distance)

\end{itemize}


\bigskip\hrule\bigskip


\sphinxAtStartPar
Styles with a \sphinxstyleemphasis{gpu}, \sphinxstyleemphasis{intel}, \sphinxstyleemphasis{kk}, \sphinxstyleemphasis{omp}, or \sphinxstyleemphasis{opt} suffix are
functionally the same as the corresponding style without the suffix.
They have been optimized to run faster, depending on your available
hardware, as discussed on the {\hyperref[\detokenize{Speed_packages::doc}]{\sphinxcrossref{\DUrole{doc}{Accelerator packages}}}}
page.  The accelerated styles take the same arguments and should
produce the same results, except for round\sphinxhyphen{}off and precision issues.

\sphinxAtStartPar
These accelerated styles are part of the GPU, INTEL, KOKKOS,
OPENMP, and OPT packages, respectively.  They are only enabled if
LAMMPS was built with those packages.  See the {\hyperref[\detokenize{Build_package::doc}]{\sphinxcrossref{\DUrole{doc}{Build package}}}} page for more info.

\sphinxAtStartPar
You can specify the accelerated styles explicitly in your input script
by including their suffix, or you can use the {\hyperref[\detokenize{Run_options::doc}]{\sphinxcrossref{\DUrole{doc}{\sphinxhyphen{}suffix command\sphinxhyphen{}line switch}}}} when you invoke LAMMPS, or you can use the
{\hyperref[\detokenize{suffix::doc}]{\sphinxcrossref{\DUrole{doc}{suffix}}}} command in your input script.

\sphinxAtStartPar
See the {\hyperref[\detokenize{Speed_packages::doc}]{\sphinxcrossref{\DUrole{doc}{Accelerator packages}}}} page for more
instructions on how to use the accelerated styles effectively.


\bigskip\hrule\bigskip



\subsection{Restrictions}
\label{\detokenize{bond_harmonic_shift:restrictions}}
\sphinxAtStartPar
This bond style can only be used if LAMMPS was built with the
EXTRA\sphinxhyphen{}MOLECULE package.  See the {\hyperref[\detokenize{Build_package::doc}]{\sphinxcrossref{\DUrole{doc}{Build package}}}} doc
page for more info.


\subsection{Related commands}
\label{\detokenize{bond_harmonic_shift:related-commands}}
\sphinxAtStartPar
{\hyperref[\detokenize{bond_coeff::doc}]{\sphinxcrossref{\DUrole{doc}{bond\_coeff}}}}, {\hyperref[\detokenize{delete_bonds::doc}]{\sphinxcrossref{\DUrole{doc}{delete\_bonds}}}},
{\hyperref[\detokenize{bond_harmonic::doc}]{\sphinxcrossref{\DUrole{doc}{bond\_harmonic}}}}


\subsection{Default}
\label{\detokenize{bond_harmonic_shift:default}}
\sphinxAtStartPar
none

\sphinxstepscope

\index{bond\_style harmonic/shift/cut@\spxentry{bond\_style harmonic/shift/cut}}\ignorespaces 
\index{bond\_style harmonic/shift/cut/omp@\spxentry{bond\_style harmonic/shift/cut/omp}}\ignorespaces 

\section{bond\_style harmonic/shift/cut command}
\label{\detokenize{bond_harmonic_shift_cut:bond-style-harmonic-shift-cut-command}}\label{\detokenize{bond_harmonic_shift_cut:index-1}}\label{\detokenize{bond_harmonic_shift_cut:index-0}}\label{\detokenize{bond_harmonic_shift_cut::doc}}
\sphinxAtStartPar
Accelerator Variants: \sphinxstyleemphasis{harmonic/shift/cut/omp}


\subsection{Syntax}
\label{\detokenize{bond_harmonic_shift_cut:syntax}}
\begin{sphinxVerbatim}[commandchars=\\\{\}]
\PYG{k}{bond\PYGZus{}style}\PYG{+w}{ }\PYG{n}{harmonic}\PYG{o}{/}\PYG{n}{shift}\PYG{o}{/}\PYG{n}{cut}
\end{sphinxVerbatim}


\subsection{Examples}
\label{\detokenize{bond_harmonic_shift_cut:examples}}
\begin{sphinxVerbatim}[commandchars=\\\{\}]
\PYG{k}{bond\PYGZus{}style}\PYG{+w}{ }\PYG{n}{harmonic}\PYG{o}{/}\PYG{n}{shift}\PYG{o}{/}\PYG{n}{cut}
\PYG{k}{bond\PYGZus{}coeff}\PYG{+w}{ }\PYG{l+m}{5}\PYG{+w}{ }\PYG{l+m}{10.0}\PYG{+w}{ }\PYG{l+m}{0.5}\PYG{+w}{ }\PYG{l+m}{1.0}
\end{sphinxVerbatim}


\subsection{Description}
\label{\detokenize{bond_harmonic_shift_cut:description}}
\sphinxAtStartPar
The \sphinxstyleemphasis{harmonic/shift/cut} bond style is a shifted harmonic bond that
uses the potential
\begin{equation*}
\begin{split}E = \frac{U_{\text{min}}}{(r_0-r_c)^2} \left[ (r-r_0)^2-(r_c-r_0)^2 \right]\end{split}
\end{equation*}
\sphinxAtStartPar
where \(r_0\) is the equilibrium bond distance, and rc the critical distance.
The bond potential is zero for distances \(r > r_c\). The potential is \(-U_{\text{min}}\)
at \(r_0\) and zero at \(r_c\). The spring constant is \(k = U_{\text{min}} / [ 2 (r_0-r_c)^2]\).

\sphinxAtStartPar
The following coefficients must be defined for each bond type via the
{\hyperref[\detokenize{bond_coeff::doc}]{\sphinxcrossref{\DUrole{doc}{bond\_coeff}}}} command as in the example above, or in
the data file or restart files read by the {\hyperref[\detokenize{read_data::doc}]{\sphinxcrossref{\DUrole{doc}{read\_data}}}}
or {\hyperref[\detokenize{read_restart::doc}]{\sphinxcrossref{\DUrole{doc}{read\_restart}}}} commands:
\begin{itemize}
\item {} 
\sphinxAtStartPar
\(U_{\text{min}}\) (energy)

\item {} 
\sphinxAtStartPar
\(r_0\) (distance)

\item {} 
\sphinxAtStartPar
\(r_c\) (distance)

\end{itemize}


\bigskip\hrule\bigskip


\sphinxAtStartPar
Styles with a \sphinxstyleemphasis{gpu}, \sphinxstyleemphasis{intel}, \sphinxstyleemphasis{kk}, \sphinxstyleemphasis{omp}, or \sphinxstyleemphasis{opt} suffix are
functionally the same as the corresponding style without the suffix.
They have been optimized to run faster, depending on your available
hardware, as discussed on the {\hyperref[\detokenize{Speed_packages::doc}]{\sphinxcrossref{\DUrole{doc}{Accelerator packages}}}}
page.  The accelerated styles take the same arguments and should
produce the same results, except for round\sphinxhyphen{}off and precision issues.

\sphinxAtStartPar
These accelerated styles are part of the GPU, INTEL, KOKKOS,
OPENMP, and OPT packages, respectively.  They are only enabled if
LAMMPS was built with those packages.  See the {\hyperref[\detokenize{Build_package::doc}]{\sphinxcrossref{\DUrole{doc}{Build package}}}} page for more info.

\sphinxAtStartPar
You can specify the accelerated styles explicitly in your input script
by including their suffix, or you can use the {\hyperref[\detokenize{Run_options::doc}]{\sphinxcrossref{\DUrole{doc}{\sphinxhyphen{}suffix command\sphinxhyphen{}line switch}}}} when you invoke LAMMPS, or you can use the
{\hyperref[\detokenize{suffix::doc}]{\sphinxcrossref{\DUrole{doc}{suffix}}}} command in your input script.

\sphinxAtStartPar
See the {\hyperref[\detokenize{Speed_packages::doc}]{\sphinxcrossref{\DUrole{doc}{Accelerator packages}}}} page for more
instructions on how to use the accelerated styles effectively.


\bigskip\hrule\bigskip



\subsection{Restrictions}
\label{\detokenize{bond_harmonic_shift_cut:restrictions}}
\sphinxAtStartPar
This bond style can only be used if LAMMPS was built with the
EXTRA\sphinxhyphen{}MOLECULE package.  See the {\hyperref[\detokenize{Build_package::doc}]{\sphinxcrossref{\DUrole{doc}{Build package}}}} doc
page for more info.


\subsection{Related commands}
\label{\detokenize{bond_harmonic_shift_cut:related-commands}}
\sphinxAtStartPar
{\hyperref[\detokenize{bond_coeff::doc}]{\sphinxcrossref{\DUrole{doc}{bond\_coeff}}}}, {\hyperref[\detokenize{delete_bonds::doc}]{\sphinxcrossref{\DUrole{doc}{delete\_bonds}}}},
{\hyperref[\detokenize{bond_harmonic::doc}]{\sphinxcrossref{\DUrole{doc}{bond\_harmonic}}}},
{\hyperref[\detokenize{bond_harmonic_shift::doc}]{\sphinxcrossref{\DUrole{doc}{bond\_style harmonic/shift}}}}


\subsection{Default}
\label{\detokenize{bond_harmonic_shift_cut:default}}
\sphinxAtStartPar
none

\sphinxstepscope

\index{bond\_style hybrid@\spxentry{bond\_style hybrid}}\ignorespaces 
\index{bond\_style hybrid/kk@\spxentry{bond\_style hybrid/kk}}\ignorespaces 

\section{bond\_style hybrid command}
\label{\detokenize{bond_hybrid:bond-style-hybrid-command}}\label{\detokenize{bond_hybrid:index-1}}\label{\detokenize{bond_hybrid:index-0}}\label{\detokenize{bond_hybrid::doc}}
\sphinxAtStartPar
Accelerator Variants: \sphinxstyleemphasis{hybrid/kk}


\subsection{Syntax}
\label{\detokenize{bond_hybrid:syntax}}
\begin{sphinxVerbatim}[commandchars=\\\{\}]
\PYG{k}{bond\PYGZus{}style}\PYG{+w}{ }\PYG{n}{hybrid}\PYG{+w}{ }\PYG{n}{style1}\PYG{+w}{ }\PYG{n}{style2}\PYG{+w}{ }\PYG{n}{...}
\end{sphinxVerbatim}
\begin{itemize}
\item {} 
\sphinxAtStartPar
style1,style2 = list of one or more bond styles

\end{itemize}


\subsection{Examples}
\label{\detokenize{bond_hybrid:examples}}
\begin{sphinxVerbatim}[commandchars=\\\{\}]
\PYG{k}{bond\PYGZus{}style}\PYG{+w}{ }\PYG{n}{hybrid}\PYG{+w}{ }\PYG{n}{harmonic}\PYG{+w}{ }\PYG{n}{fene}
\PYG{k}{bond\PYGZus{}coeff}\PYG{+w}{ }\PYG{l+m}{1}\PYG{+w}{ }\PYG{n}{harmonic}\PYG{+w}{ }\PYG{l+m}{80.0}\PYG{+w}{ }\PYG{l+m}{1.2}
\PYG{k}{bond\PYGZus{}coeff}\PYG{+w}{ }\PYG{l+m}{2}\PYG{o}{*}\PYG{+w}{ }\PYG{n}{fene}\PYG{+w}{ }\PYG{l+m}{30.0}\PYG{+w}{ }\PYG{l+m}{1.5}\PYG{+w}{ }\PYG{l+m}{1.0}\PYG{+w}{ }\PYG{l+m}{1.0}
\end{sphinxVerbatim}


\subsection{Description}
\label{\detokenize{bond_hybrid:description}}
\sphinxAtStartPar
The \sphinxstyleemphasis{hybrid} style enables the use of multiple bond styles in one
simulation.  A bond style is assigned to each bond type.  For example,
bonds in a polymer flow (of bond type 1) could be computed with a
\sphinxstyleemphasis{fene} potential and bonds in the wall boundary (of bond type 2) could
be computed with a \sphinxstyleemphasis{harmonic} potential.  The assignment of bond type
to style is made via the {\hyperref[\detokenize{bond_coeff::doc}]{\sphinxcrossref{\DUrole{doc}{bond\_coeff}}}} command or in
the data file.

\sphinxAtStartPar
In the bond\_coeff commands, the name of a bond style must be added
after the bond type, with the remaining coefficients being those
appropriate to that style.  In the example above, the 2 bond\_coeff
commands set bonds of bond type 1 to be computed with a \sphinxstyleemphasis{harmonic}
potential with coefficients 80.0, 1.2 for \(K\), \(r_0\).  All other bond types
(2\sphinxhyphen{}N) are computed with a \sphinxstyleemphasis{fene} potential with coefficients 30.0,
1.5, 1.0, 1.0 for \(K\), \(R_0\), \(\epsilon\), \(\sigma\).

\sphinxAtStartPar
If bond coefficients are specified in the data file read via the
{\hyperref[\detokenize{read_data::doc}]{\sphinxcrossref{\DUrole{doc}{read\_data}}}} command, then the same rule applies.
E.g. “harmonic” or “fene” must be added after the bond type, for each
line in the “Bond Coeffs” section, e.g.

\begin{sphinxVerbatim}[commandchars=\\\{\}]
Bond Coeffs

1 harmonic 80.0 1.2
2 fene 30.0 1.5 1.0 1.0
...
\end{sphinxVerbatim}

\sphinxAtStartPar
A bond style of \sphinxstyleemphasis{none} with no additional coefficients can be used in
place of a bond style, either in a input script bond\_coeff command or
in the data file, if you desire to turn off interactions for specific
bond types.


\bigskip\hrule\bigskip


\sphinxAtStartPar
Styles with a \sphinxstyleemphasis{gpu}, \sphinxstyleemphasis{intel}, \sphinxstyleemphasis{kk}, \sphinxstyleemphasis{omp}, or \sphinxstyleemphasis{opt} suffix are
functionally the same as the corresponding style without the suffix.
They have been optimized to run faster, depending on your available
hardware, as discussed on the {\hyperref[\detokenize{Speed_packages::doc}]{\sphinxcrossref{\DUrole{doc}{Accelerator packages}}}}
page.  The accelerated styles take the same arguments and should
produce the same results, except for round\sphinxhyphen{}off and precision issues.

\sphinxAtStartPar
These accelerated styles are part of the GPU, INTEL, KOKKOS,
OPENMP, and OPT packages, respectively.  They are only enabled if
LAMMPS was built with those packages.  See the {\hyperref[\detokenize{Build_package::doc}]{\sphinxcrossref{\DUrole{doc}{Build package}}}} page for more info.

\sphinxAtStartPar
You can specify the accelerated styles explicitly in your input script
by including their suffix, or you can use the {\hyperref[\detokenize{Run_options::doc}]{\sphinxcrossref{\DUrole{doc}{\sphinxhyphen{}suffix command\sphinxhyphen{}line switch}}}} when you invoke LAMMPS, or you can use the
{\hyperref[\detokenize{suffix::doc}]{\sphinxcrossref{\DUrole{doc}{suffix}}}} command in your input script.

\sphinxAtStartPar
See the {\hyperref[\detokenize{Speed_packages::doc}]{\sphinxcrossref{\DUrole{doc}{Accelerator packages}}}} page for more
instructions on how to use the accelerated styles effectively.


\bigskip\hrule\bigskip



\subsection{Restrictions}
\label{\detokenize{bond_hybrid:restrictions}}
\sphinxAtStartPar
This bond style can only be used if LAMMPS was built with the MOLECULE
package.  See the {\hyperref[\detokenize{Build_package::doc}]{\sphinxcrossref{\DUrole{doc}{Build package}}}} page for more
info.

\sphinxAtStartPar
Unlike other bond styles, the hybrid bond style does not store bond
coefficient info for individual sub\sphinxhyphen{}styles in {\hyperref[\detokenize{restart::doc}]{\sphinxcrossref{\DUrole{doc}{binary restart files}}}} or {\hyperref[\detokenize{write_data::doc}]{\sphinxcrossref{\DUrole{doc}{data files}}}}.  Thus when restarting a
simulation, you need to re\sphinxhyphen{}specify the bond\_coeff commands.


\subsection{Related commands}
\label{\detokenize{bond_hybrid:related-commands}}
\sphinxAtStartPar
{\hyperref[\detokenize{bond_coeff::doc}]{\sphinxcrossref{\DUrole{doc}{bond\_coeff}}}}, {\hyperref[\detokenize{delete_bonds::doc}]{\sphinxcrossref{\DUrole{doc}{delete\_bonds}}}}


\subsection{Default}
\label{\detokenize{bond_hybrid:default}}
\sphinxAtStartPar
none

\sphinxstepscope

\index{bond\_style lepton@\spxentry{bond\_style lepton}}\ignorespaces 
\index{bond\_style lepton/omp@\spxentry{bond\_style lepton/omp}}\ignorespaces 

\section{bond\_style lepton command}
\label{\detokenize{bond_lepton:bond-style-lepton-command}}\label{\detokenize{bond_lepton:index-1}}\label{\detokenize{bond_lepton:index-0}}\label{\detokenize{bond_lepton::doc}}
\sphinxAtStartPar
Accelerator Variants: \sphinxstyleemphasis{lepton/omp}


\subsection{Syntax}
\label{\detokenize{bond_lepton:syntax}}
\begin{sphinxVerbatim}[commandchars=\\\{\}]
\PYG{k}{bond\PYGZus{}style}\PYG{+w}{ }\PYG{n}{style}\PYG{+w}{ }\PYG{n}{args}
\end{sphinxVerbatim}
\begin{itemize}
\item {} 
\sphinxAtStartPar
style = \sphinxstyleemphasis{lepton}

\item {} 
\sphinxAtStartPar
args = optional arguments

\end{itemize}
\begin{sphinxalltt}
args = \sphinxstyleemphasis{auto\_offset} or \sphinxstyleemphasis{no\_offset}
  \sphinxstyleemphasis{auto\_offset} = offset the potential energy so that the value at r0 is 0.0 (default)
  \sphinxstyleemphasis{no\_offset} = do not offset the potential energy
\end{sphinxalltt}


\subsection{Examples}
\label{\detokenize{bond_lepton:examples}}
\begin{sphinxVerbatim}[commandchars=\\\{\}]
\PYG{k}{bond\PYGZus{}style}\PYG{+w}{ }\PYG{n}{lepton}
\PYG{k}{bond\PYGZus{}style}\PYG{+w}{ }\PYG{n}{lepton}\PYG{+w}{ }\PYG{n}{no\PYGZus{}offset}

\PYG{k}{bond\PYGZus{}coeff}\PYG{+w}{  }\PYG{l+m}{1}\PYG{+w}{  }\PYG{l+m}{1.5}\PYG{+w}{ }\PYG{l+s}{\PYGZdq{}}\PYG{l+s}{k*r\PYGZca{}2; k=250.0}\PYG{l+s}{\PYGZdq{}}
\PYG{k}{bond\PYGZus{}coeff}\PYG{+w}{  }\PYG{l+m}{2}\PYG{+w}{  }\PYG{l+m}{1.1}\PYG{+w}{ }\PYG{l+s}{\PYGZdq{}}\PYG{l+s}{k2*r\PYGZca{}2 + k3*r\PYGZca{}3 + k4*r\PYGZca{}4; k2=300.0; k3=\PYGZhy{}100.0; k4=50.0}\PYG{l+s}{\PYGZdq{}}
\PYG{k}{bond\PYGZus{}coeff}\PYG{+w}{  }\PYG{l+m}{3}\PYG{+w}{  }\PYG{l+m}{1.3}\PYG{+w}{ }\PYG{l+s}{\PYGZdq{}}\PYG{l+s}{k*r\PYGZca{}2; k=350.0}\PYG{l+s}{\PYGZdq{}}
\end{sphinxVerbatim}


\subsection{Description}
\label{\detokenize{bond_lepton:description}}
\sphinxAtStartPar
\DUrole{versionmodified,added}{Added in version 8Feb2023.}

\sphinxAtStartPar
Bond style \sphinxstyleemphasis{lepton} computes bonded interactions between two atoms with
a custom function.  The potential function must be provided as an
expression string using “r” as the distance variable relative to the
reference distance \(r_0\) which is provided as a bond coefficient.
For example \sphinxtitleref{“200.0*r\textasciicircum{}2”} represents a harmonic potential with a force
constant \sphinxstyleemphasis{K} of 200.0 energy units:
\begin{equation*}
\begin{split}U_{bond,i} = K (r_i - r_0)^2 = K r^2 \qquad r = r_i - r_0\end{split}
\end{equation*}
\sphinxAtStartPar
\DUrole{versionmodified,changed}{Changed in version 7Feb2024.}

\sphinxAtStartPar
By default the potential energy U is shifted so that he value U is 0.0
for \$r = r\_0\$.  This is equivalent to using the optional keyword
\sphinxstyleemphasis{auto\_offset}.  When using the keyword \sphinxstyleemphasis{no\_offset} instead, the
potential energy is not shifted.

\sphinxAtStartPar
The \sphinxhref{https://simtk.org/projects/lepton}{Lepton library}, that the
\sphinxstyleemphasis{lepton} bond style interfaces with, evaluates this expression string at
run time to compute the pairwise energy.  It also creates an analytical
representation of the first derivative of this expression with respect to
“r” and then uses that to compute the force between the atom pairs forming
bonds as defined by the topology data.

\sphinxAtStartPar
The following coefficients must be defined for each bond type via the
{\hyperref[\detokenize{bond_coeff::doc}]{\sphinxcrossref{\DUrole{doc}{bond\_coeff}}}} command as in the examples above, or in
the data file or restart files read by the {\hyperref[\detokenize{read_data::doc}]{\sphinxcrossref{\DUrole{doc}{read\_data}}}}
or {\hyperref[\detokenize{read_restart::doc}]{\sphinxcrossref{\DUrole{doc}{read\_restart}}}} commands:
\begin{itemize}
\item {} 
\sphinxAtStartPar
Lepton expression (energy units)

\item {} 
\sphinxAtStartPar
\(r_0\) (distance)

\end{itemize}

\sphinxAtStartPar
The Lepton expression must be either enclosed in quotes or must not
contain any whitespace so that LAMMPS recognizes it as a single keyword.
More on valid Lepton expressions below.  The \(r_0\) is the
“equilibrium distance”.  The potential energy function in the Lepton
expression is shifted in such a way, that the potential energy is 0 for
a bond length \(r_i == r_0\).


\bigskip\hrule\bigskip



\subsection{Lepton expression syntax and features}
\label{\detokenize{bond_lepton:lepton-expression-syntax-and-features}}
\sphinxAtStartPar
Lepton supports the following operators in expressions:


\begin{savenotes}\sphinxattablestart
\sphinxthistablewithglobalstyle
\centering
\begin{tabulary}{\linewidth}[t]{TTTTTTTTTTTTTT}
\sphinxtoprule
\sphinxtableatstartofbodyhook
\sphinxAtStartPar
+
&
\sphinxAtStartPar
Add
&&
\sphinxAtStartPar
\sphinxhyphen{}
&
\sphinxAtStartPar
Subtract
&&
\sphinxAtStartPar
*
&
\sphinxAtStartPar
Multiply
&&
\sphinxAtStartPar
/
&
\sphinxAtStartPar
Divide
&&
\sphinxAtStartPar
\textasciicircum{}
&
\sphinxAtStartPar
Power
\\
\sphinxbottomrule
\end{tabulary}
\sphinxtableafterendhook\par
\sphinxattableend\end{savenotes}

\sphinxAtStartPar
The following mathematical functions are available:


\begin{savenotes}\sphinxattablestart
\sphinxthistablewithglobalstyle
\centering
\begin{tabulary}{\linewidth}[t]{TTTT}
\sphinxtoprule
\sphinxtableatstartofbodyhook
\sphinxAtStartPar
sqrt(x)
&
\sphinxAtStartPar
Square root
&
\sphinxAtStartPar
exp(x)
&
\sphinxAtStartPar
Exponential
\\
\sphinxhline
\sphinxAtStartPar
log(x)
&
\sphinxAtStartPar
Natural logarithm
&
\sphinxAtStartPar
sin(x)
&
\sphinxAtStartPar
Sine (angle in radians)
\\
\sphinxhline
\sphinxAtStartPar
cos(x)
&
\sphinxAtStartPar
Cosine (angle in radians)
&
\sphinxAtStartPar
sec(x)
&
\sphinxAtStartPar
Secant (angle in radians)
\\
\sphinxhline
\sphinxAtStartPar
csc(x)
&
\sphinxAtStartPar
Cosecant (angle in radians)
&
\sphinxAtStartPar
tan(x)
&
\sphinxAtStartPar
Tangent (angle in radians)
\\
\sphinxhline
\sphinxAtStartPar
cot(x)
&
\sphinxAtStartPar
Cotangent (angle in radians)
&
\sphinxAtStartPar
asin(x)
&
\sphinxAtStartPar
Inverse sine (in radians)
\\
\sphinxhline
\sphinxAtStartPar
acos(x)
&
\sphinxAtStartPar
Inverse cosine (in radians)
&
\sphinxAtStartPar
atan(x)
&
\sphinxAtStartPar
Inverse tangent (in radians)
\\
\sphinxhline
\sphinxAtStartPar
sinh(x)
&
\sphinxAtStartPar
Hyperbolic sine
&
\sphinxAtStartPar
cosh(x)
&
\sphinxAtStartPar
Hyperbolic cosine
\\
\sphinxhline
\sphinxAtStartPar
tanh(x)
&
\sphinxAtStartPar
Hyperbolic tangent
&
\sphinxAtStartPar
erf(x)
&
\sphinxAtStartPar
Error function
\\
\sphinxhline
\sphinxAtStartPar
erfc(x)
&
\sphinxAtStartPar
Complementary Error function
&
\sphinxAtStartPar
abs(x)
&
\sphinxAtStartPar
Absolute value
\\
\sphinxhline
\sphinxAtStartPar
min(x,y)
&
\sphinxAtStartPar
Minimum of two values
&
\sphinxAtStartPar
max(x,y)
&
\sphinxAtStartPar
Maximum of two values
\\
\sphinxhline
\sphinxAtStartPar
delta(x)
&
\sphinxAtStartPar
delta(x) is 1 for \sphinxtitleref{x = 0}, otherwise 0
&
\sphinxAtStartPar
step(x)
&
\sphinxAtStartPar
step(x) is 0 for \sphinxtitleref{x \textless{} 0}, otherwise 1
\\
\sphinxbottomrule
\end{tabulary}
\sphinxtableafterendhook\par
\sphinxattableend\end{savenotes}

\sphinxAtStartPar
Numbers may be given in either decimal or exponential form.  All of the
following are valid numbers: \sphinxtitleref{5}, \sphinxtitleref{\sphinxhyphen{}3.1}, \sphinxtitleref{1e6}, and \sphinxtitleref{3.12e\sphinxhyphen{}2}.

\sphinxAtStartPar
As an extension to the standard Lepton syntax, it is also possible to
use LAMMPS {\hyperref[\detokenize{variable::doc}]{\sphinxcrossref{\DUrole{doc}{variables}}}} in the format “v\_name”.  Before
evaluating the expression, “v\_name” will be replaced with the value of
the variable “name”.  This is compatible with all kinds of scalar
variables, but not with vectors, arrays, local, or per\sphinxhyphen{}atom
variables.  If necessary, a custom scalar variable needs to be defined
that can access the desired (single) item from a non\sphinxhyphen{}scalar variable.
As an example, the following lines will instruct LAMMPS to ramp
the force constant for a harmonic bond from 100.0 to 200.0 during the
next run:

\begin{sphinxVerbatim}[commandchars=\\\{\}]
\PYG{k}{variable }\PYG{n+nv+nvIdentifier}{fconst}\PYG{+w}{ }\PYG{n}{equal}\PYG{+w}{ }\PYG{n}{ramp}\PYG{n+nv}{(}\PYG{n+nv}{100.0, 200}\PYG{n+nv}{)}
\PYG{k}{bond\PYGZus{}style}\PYG{+w}{ }\PYG{n}{lepton}
\PYG{k}{bond\PYGZus{}coeff}\PYG{+w}{ }\PYG{l+m}{1}\PYG{+w}{ }\PYG{l+m}{1.5}\PYG{+w}{ }\PYG{l+s}{\PYGZdq{}}\PYG{l+s}{v\PYGZus{}fconst * (r\PYGZca{}2)}\PYG{l+s}{\PYGZdq{}}
\end{sphinxVerbatim}

\sphinxAtStartPar
An expression may be followed by definitions for intermediate values that appear in the
expression. A semicolon “;” is used as a delimiter between value definitions. For example,
the expression:

\begin{sphinxVerbatim}[commandchars=\\\{\}]
\PYG{n}{a}\PYG{o}{\PYGZca{}}\PYG{l+m+mi}{2}\PYG{o}{+}\PYG{n}{a}\PYG{o}{*}\PYG{n}{b}\PYG{o}{+}\PYG{n}{b}\PYG{o}{\PYGZca{}}\PYG{l+m+mi}{2}\PYG{p}{;}\PYG{+w}{ }\PYG{n}{a}\PYG{o}{=}\PYG{n}{a1}\PYG{o}{+}\PYG{n}{a2}\PYG{p}{;}\PYG{+w}{ }\PYG{n}{b}\PYG{o}{=}\PYG{n}{b1}\PYG{o}{+}\PYG{n}{b2}
\end{sphinxVerbatim}

\sphinxAtStartPar
is exactly equivalent to

\begin{sphinxVerbatim}[commandchars=\\\{\}]
\PYG{p}{(}\PYG{n}{a1}\PYG{o}{+}\PYG{n}{a2}\PYG{p}{)}\PYG{o}{\PYGZca{}}\PYG{l+m+mi}{2}\PYG{o}{+}\PYG{p}{(}\PYG{n}{a1}\PYG{o}{+}\PYG{n}{a2}\PYG{p}{)}\PYG{o}{*}\PYG{p}{(}\PYG{n}{b1}\PYG{o}{+}\PYG{n}{b2}\PYG{p}{)}\PYG{o}{+}\PYG{p}{(}\PYG{n}{b1}\PYG{o}{+}\PYG{n}{b2}\PYG{p}{)}\PYG{o}{\PYGZca{}}\PYG{l+m+mi}{2}
\end{sphinxVerbatim}

\sphinxAtStartPar
The definition of an intermediate value may itself involve other
intermediate values. Whitespace and quotation characters (’\textquotesingle{}’ and ‘”’)
are ignored.  All uses of a value must appear \sphinxstyleemphasis{before} that value’s
definition.  For efficiency reasons, the expression string is parsed,
optimized, and then stored in an internal, pre\sphinxhyphen{}parsed representation for
evaluation.

\sphinxAtStartPar
Evaluating a Lepton expression is typically between 2.5 and 5 times
slower than the corresponding compiled and optimized C++ code.  If
additional speed or GPU acceleration (via GPU or KOKKOS) is required,
the interaction can be represented as a table.  Suitable table files
can be created either internally using the {\hyperref[\detokenize{pair_write::doc}]{\sphinxcrossref{\DUrole{doc}{pair\_write}}}}
or {\hyperref[\detokenize{bond_write::doc}]{\sphinxcrossref{\DUrole{doc}{bond\_write}}}} command or through the Python scripts
in the {\hyperref[\detokenize{Tools:tabulate}]{\sphinxcrossref{\DUrole{std,std-ref}{tools/tabulate}}}} folder.


\bigskip\hrule\bigskip


\sphinxAtStartPar
Styles with a \sphinxstyleemphasis{gpu}, \sphinxstyleemphasis{intel}, \sphinxstyleemphasis{kk}, \sphinxstyleemphasis{omp}, or \sphinxstyleemphasis{opt} suffix are
functionally the same as the corresponding style without the suffix.
They have been optimized to run faster, depending on your available
hardware, as discussed on the {\hyperref[\detokenize{Speed_packages::doc}]{\sphinxcrossref{\DUrole{doc}{Accelerator packages}}}}
page.  The accelerated styles take the same arguments and should
produce the same results, except for round\sphinxhyphen{}off and precision issues.

\sphinxAtStartPar
These accelerated styles are part of the GPU, INTEL, KOKKOS,
OPENMP, and OPT packages, respectively.  They are only enabled if
LAMMPS was built with those packages.  See the {\hyperref[\detokenize{Build_package::doc}]{\sphinxcrossref{\DUrole{doc}{Build package}}}} page for more info.

\sphinxAtStartPar
You can specify the accelerated styles explicitly in your input script
by including their suffix, or you can use the {\hyperref[\detokenize{Run_options::doc}]{\sphinxcrossref{\DUrole{doc}{\sphinxhyphen{}suffix command\sphinxhyphen{}line switch}}}} when you invoke LAMMPS, or you can use the
{\hyperref[\detokenize{suffix::doc}]{\sphinxcrossref{\DUrole{doc}{suffix}}}} command in your input script.

\sphinxAtStartPar
See the {\hyperref[\detokenize{Speed_packages::doc}]{\sphinxcrossref{\DUrole{doc}{Accelerator packages}}}} page for more
instructions on how to use the accelerated styles effectively.


\bigskip\hrule\bigskip



\subsection{Restrictions}
\label{\detokenize{bond_lepton:restrictions}}
\sphinxAtStartPar
This bond style is part of the LEPTON package and only enabled if LAMMPS
was built with this package.  See the {\hyperref[\detokenize{Build_package::doc}]{\sphinxcrossref{\DUrole{doc}{Build package}}}} page for more info.


\subsection{Related commands}
\label{\detokenize{bond_lepton:related-commands}}
\sphinxAtStartPar
{\hyperref[\detokenize{bond_coeff::doc}]{\sphinxcrossref{\DUrole{doc}{bond\_coeff}}}}, {\hyperref[\detokenize{bond_table::doc}]{\sphinxcrossref{\DUrole{doc}{bond\_style table}}}},
{\hyperref[\detokenize{bond_write::doc}]{\sphinxcrossref{\DUrole{doc}{bond\_write}}}}, {\hyperref[\detokenize{angle_lepton::doc}]{\sphinxcrossref{\DUrole{doc}{angle\_style lepton}}}},
{\hyperref[\detokenize{dihedral_lepton::doc}]{\sphinxcrossref{\DUrole{doc}{dihedral\_style lepton}}}}


\subsection{Default}
\label{\detokenize{bond_lepton:default}}
\sphinxAtStartPar
none

\sphinxstepscope

\index{bond\_style mesocnt@\spxentry{bond\_style mesocnt}}\ignorespaces 

\section{bond\_style mesocnt command}
\label{\detokenize{bond_mesocnt:bond-style-mesocnt-command}}\label{\detokenize{bond_mesocnt:index-0}}\label{\detokenize{bond_mesocnt::doc}}

\subsection{Syntax}
\label{\detokenize{bond_mesocnt:syntax}}
\begin{sphinxVerbatim}[commandchars=\\\{\}]
\PYG{k}{bond\PYGZus{}style}\PYG{+w}{ }\PYG{n}{mesocnt}
\end{sphinxVerbatim}


\subsection{Examples}
\label{\detokenize{bond_mesocnt:examples}}
\begin{sphinxVerbatim}[commandchars=\\\{\}]
\PYG{k}{bond\PYGZus{}style}\PYG{+w}{ }\PYG{n}{mesocnt}
\PYG{k}{bond\PYGZus{}coeff}\PYG{+w}{ }\PYG{l+m}{1}\PYG{+w}{ }\PYG{n}{C}\PYG{+w}{ }\PYG{l+m}{10}\PYG{+w}{ }\PYG{l+m}{10}\PYG{+w}{ }\PYG{l+m}{20.0}
\PYG{k}{bond\PYGZus{}coeff}\PYG{+w}{ }\PYG{l+m}{4}\PYG{+w}{ }\PYG{n}{custom}\PYG{+w}{ }\PYG{l+m}{800.0}\PYG{+w}{ }\PYG{l+m}{10.0}
\end{sphinxVerbatim}


\subsection{Description}
\label{\detokenize{bond_mesocnt:description}}
\sphinxAtStartPar
\DUrole{versionmodified,added}{Added in version 15Sep2022.}

\sphinxAtStartPar
The \sphinxstyleemphasis{mesocnt} bond style is a wrapper for the {\hyperref[\detokenize{bond_harmonic::doc}]{\sphinxcrossref{\DUrole{doc}{harmonic}}}} style, and uses the potential
\begin{equation*}
\begin{split}E = K (r - r_0)^2\end{split}
\end{equation*}
\sphinxAtStartPar
where \(r_0\) is the equilibrium bond distance.  Note that the
usual 1/2 factor is included in \(K\).  The style implements
parameterization presets of \(K\) for mesoscopic simulations of
carbon nanotubes based on the atomistic simulations of
{\hyperref[\detokenize{bond_mesocnt:srivastava-1}]{\sphinxcrossref{\DUrole{std,std-ref}{(Srivastava)}}}}.

\sphinxAtStartPar
Other presets can be readily implemented in the future.

\sphinxAtStartPar
The following coefficients must be defined for each bond type via the
{\hyperref[\detokenize{bond_coeff::doc}]{\sphinxcrossref{\DUrole{doc}{bond\_coeff}}}} command as in the example above, or in
the data file or restart files read by the {\hyperref[\detokenize{read_data::doc}]{\sphinxcrossref{\DUrole{doc}{read\_data}}}} or {\hyperref[\detokenize{read_restart::doc}]{\sphinxcrossref{\DUrole{doc}{read\_restart}}}} commands:
\begin{itemize}
\item {} 
\sphinxAtStartPar
preset = \sphinxstyleemphasis{C} or \sphinxstyleemphasis{custom}

\item {} 
\sphinxAtStartPar
additional parameters depending on preset

\end{itemize}

\sphinxAtStartPar
Preset \sphinxstyleemphasis{C} is for carbon nanotubes, and the additional parameters are:
\begin{itemize}
\item {} 
\sphinxAtStartPar
chiral index \(n\) (unitless)

\item {} 
\sphinxAtStartPar
chiral index \(m\) (unitless)

\item {} 
\sphinxAtStartPar
\(r_0\) (distance)

\end{itemize}

\sphinxAtStartPar
Preset \sphinxstyleemphasis{custom} is simply a direct wrapper for the {\hyperref[\detokenize{bond_harmonic::doc}]{\sphinxcrossref{\DUrole{doc}{harmonic}}}} style, and the additional parameters are:
\begin{itemize}
\item {} 
\sphinxAtStartPar
\(K\) (energy/distance\textasciicircum{}2)

\item {} 
\sphinxAtStartPar
\(r_0\) (distance)

\end{itemize}


\subsection{Restrictions}
\label{\detokenize{bond_mesocnt:restrictions}}
\sphinxAtStartPar
This bond style can only be used if LAMMPS was built with the MOLECULE
and MESONT packages.  See the {\hyperref[\detokenize{Build_package::doc}]{\sphinxcrossref{\DUrole{doc}{Build package}}}}
page for more info.


\subsection{Related commands}
\label{\detokenize{bond_mesocnt:related-commands}}
\sphinxAtStartPar
{\hyperref[\detokenize{bond_coeff::doc}]{\sphinxcrossref{\DUrole{doc}{bond\_coeff}}}}, {\hyperref[\detokenize{delete_bonds::doc}]{\sphinxcrossref{\DUrole{doc}{delete\_bonds}}}}


\subsection{Default}
\label{\detokenize{bond_mesocnt:default}}
\sphinxAtStartPar
none


\bigskip\hrule\bigskip

\phantomsection\label{\detokenize{bond_mesocnt:srivastava-1}}
\sphinxAtStartPar
\sphinxstylestrong{(Srivastava)} Zhigilei, Wei and Srivastava, Phys. Rev. B 71, 165417
(2005).

\sphinxstepscope

\index{bond\_style mm3@\spxentry{bond\_style mm3}}\ignorespaces 

\section{bond\_style mm3 command}
\label{\detokenize{bond_mm3:bond-style-mm3-command}}\label{\detokenize{bond_mm3:index-0}}\label{\detokenize{bond_mm3::doc}}

\subsection{Syntax}
\label{\detokenize{bond_mm3:syntax}}
\begin{sphinxVerbatim}[commandchars=\\\{\}]
\PYG{k}{bond\PYGZus{}style}\PYG{+w}{ }\PYG{n}{mm3}
\end{sphinxVerbatim}


\subsection{Examples}
\label{\detokenize{bond_mm3:examples}}
\begin{sphinxVerbatim}[commandchars=\\\{\}]
\PYG{k}{bond\PYGZus{}style}\PYG{+w}{ }\PYG{n}{mm3}
\PYG{k}{bond\PYGZus{}coeff}\PYG{+w}{ }\PYG{l+m}{1}\PYG{+w}{ }\PYG{l+m}{100.0}\PYG{+w}{ }\PYG{l+m}{107.0}
\end{sphinxVerbatim}


\subsection{Description}
\label{\detokenize{bond_mm3:description}}
\sphinxAtStartPar
The \sphinxstyleemphasis{mm3} bond style uses the potential that is anharmonic in the bond
as defined in {\hyperref[\detokenize{bond_mm3:mm3-allinger1989}]{\sphinxcrossref{\DUrole{std,std-ref}{(Allinger)}}}}
\begin{equation*}
\begin{split}E = K (r - r_0)^2 \left[ 1 - 2.55(r-r_0) + \frac{7}{12} 2.55^2(r-r_0)^2 \right]\end{split}
\end{equation*}
\sphinxAtStartPar
where \(r_0\) is the equilibrium value of the bond, and \(K\) is a
prefactor. The anharmonic prefactors have units \(\AA^{-n}\):
\(-2.55 \AA^{-1}\) and \(\frac{7}{12} 2.55^2 \AA^{-2}\). The code takes
care of the necessary unit conversion for these factors internally.
Note that the MM3 papers contain an error in Eq (1):
\(\frac{7}{12} 2.55\) should be replaced with \(\frac{7}{12} 2.55^2\)

\sphinxAtStartPar
The following coefficients must be defined for each bond type via the
{\hyperref[\detokenize{bond_coeff::doc}]{\sphinxcrossref{\DUrole{doc}{bond\_coeff}}}} command as in the example above, or in
the data file or restart files read by the {\hyperref[\detokenize{read_data::doc}]{\sphinxcrossref{\DUrole{doc}{read\_data}}}}
or {\hyperref[\detokenize{read_restart::doc}]{\sphinxcrossref{\DUrole{doc}{read\_restart}}}} commands:
\begin{itemize}
\item {} 
\sphinxAtStartPar
\(K\) (energy/distance\textasciicircum{}2)

\item {} 
\sphinxAtStartPar
\(r_0\) (distance)

\end{itemize}


\subsection{Restrictions}
\label{\detokenize{bond_mm3:restrictions}}
\sphinxAtStartPar
This bond style can only be used if LAMMPS was built with the
YAFF package.  See the {\hyperref[\detokenize{Build_package::doc}]{\sphinxcrossref{\DUrole{doc}{Build package}}}} doc
page for more info.


\subsection{Related commands}
\label{\detokenize{bond_mm3:related-commands}}
\sphinxAtStartPar
{\hyperref[\detokenize{bond_coeff::doc}]{\sphinxcrossref{\DUrole{doc}{bond\_coeff}}}}


\subsection{Default}
\label{\detokenize{bond_mm3:default}}
\sphinxAtStartPar
none


\bigskip\hrule\bigskip

\phantomsection\label{\detokenize{bond_mm3:mm3-allinger1989}}
\sphinxAtStartPar
\sphinxstylestrong{(Allinger)} Allinger, Yuh, Lii, JACS, 111(23), 8551\sphinxhyphen{}8566
(1989),

\sphinxstepscope

\index{bond\_style morse@\spxentry{bond\_style morse}}\ignorespaces 
\index{bond\_style morse/omp@\spxentry{bond\_style morse/omp}}\ignorespaces 

\section{bond\_style morse command}
\label{\detokenize{bond_morse:bond-style-morse-command}}\label{\detokenize{bond_morse:index-1}}\label{\detokenize{bond_morse:index-0}}\label{\detokenize{bond_morse::doc}}
\sphinxAtStartPar
Accelerator Variants: \sphinxstyleemphasis{morse/omp}


\subsection{Syntax}
\label{\detokenize{bond_morse:syntax}}
\begin{sphinxVerbatim}[commandchars=\\\{\}]
\PYG{k}{bond\PYGZus{}style}\PYG{+w}{ }\PYG{n}{morse}
\end{sphinxVerbatim}


\subsection{Examples}
\label{\detokenize{bond_morse:examples}}
\begin{sphinxVerbatim}[commandchars=\\\{\}]
\PYG{k}{bond\PYGZus{}style}\PYG{+w}{ }\PYG{n}{morse}
\PYG{k}{bond\PYGZus{}coeff}\PYG{+w}{ }\PYG{l+m}{5}\PYG{+w}{ }\PYG{l+m}{1.0}\PYG{+w}{ }\PYG{l+m}{2.0}\PYG{+w}{ }\PYG{l+m}{1.2}
\end{sphinxVerbatim}


\subsection{Description}
\label{\detokenize{bond_morse:description}}
\sphinxAtStartPar
The \sphinxstyleemphasis{morse} bond style uses the potential
\begin{equation*}
\begin{split}E = D \left[ 1 - e^{-\alpha (r - r_0)} \right]^2\end{split}
\end{equation*}
\sphinxAtStartPar
where \(r_0\) is the equilibrium bond distance, \(\alpha\) is a stiffness
parameter, and \(D\) determines the depth of the potential well.

\sphinxAtStartPar
The following coefficients must be defined for each bond type via the
{\hyperref[\detokenize{bond_coeff::doc}]{\sphinxcrossref{\DUrole{doc}{bond\_coeff}}}} command as in the example above, or in
the data file or restart files read by the {\hyperref[\detokenize{read_data::doc}]{\sphinxcrossref{\DUrole{doc}{read\_data}}}}
or {\hyperref[\detokenize{read_restart::doc}]{\sphinxcrossref{\DUrole{doc}{read\_restart}}}} commands:
\begin{itemize}
\item {} 
\sphinxAtStartPar
\(D\) (energy)

\item {} 
\sphinxAtStartPar
\(\alpha\) (inverse distance)

\item {} 
\sphinxAtStartPar
\(r_0\) (distance)

\end{itemize}


\bigskip\hrule\bigskip


\sphinxAtStartPar
Styles with a \sphinxstyleemphasis{gpu}, \sphinxstyleemphasis{intel}, \sphinxstyleemphasis{kk}, \sphinxstyleemphasis{omp}, or \sphinxstyleemphasis{opt} suffix are
functionally the same as the corresponding style without the suffix.
They have been optimized to run faster, depending on your available
hardware, as discussed on the {\hyperref[\detokenize{Speed_packages::doc}]{\sphinxcrossref{\DUrole{doc}{Accelerator packages}}}}
page.  The accelerated styles take the same arguments and should
produce the same results, except for round\sphinxhyphen{}off and precision issues.

\sphinxAtStartPar
These accelerated styles are part of the GPU, INTEL, KOKKOS,
OPENMP, and OPT packages, respectively.  They are only enabled if
LAMMPS was built with those packages.  See the {\hyperref[\detokenize{Build_package::doc}]{\sphinxcrossref{\DUrole{doc}{Build package}}}} page for more info.

\sphinxAtStartPar
You can specify the accelerated styles explicitly in your input script
by including their suffix, or you can use the {\hyperref[\detokenize{Run_options::doc}]{\sphinxcrossref{\DUrole{doc}{\sphinxhyphen{}suffix command\sphinxhyphen{}line switch}}}} when you invoke LAMMPS, or you can use the
{\hyperref[\detokenize{suffix::doc}]{\sphinxcrossref{\DUrole{doc}{suffix}}}} command in your input script.

\sphinxAtStartPar
See the {\hyperref[\detokenize{Speed_packages::doc}]{\sphinxcrossref{\DUrole{doc}{Accelerator packages}}}} page for more
instructions on how to use the accelerated styles effectively.


\bigskip\hrule\bigskip



\subsection{Restrictions}
\label{\detokenize{bond_morse:restrictions}}
\sphinxAtStartPar
This bond style can only be used if LAMMPS was built with the MOLECULE
package.  See the {\hyperref[\detokenize{Build_package::doc}]{\sphinxcrossref{\DUrole{doc}{Build package}}}} page for more
info.


\subsection{Related commands}
\label{\detokenize{bond_morse:related-commands}}
\sphinxAtStartPar
{\hyperref[\detokenize{bond_coeff::doc}]{\sphinxcrossref{\DUrole{doc}{bond\_coeff}}}}, {\hyperref[\detokenize{delete_bonds::doc}]{\sphinxcrossref{\DUrole{doc}{delete\_bonds}}}}


\subsection{Default}
\label{\detokenize{bond_morse:default}}
\sphinxAtStartPar
none

\sphinxstepscope

\index{bond\_style none@\spxentry{bond\_style none}}\ignorespaces 

\section{bond\_style none command}
\label{\detokenize{bond_none:bond-style-none-command}}\label{\detokenize{bond_none:index-0}}\label{\detokenize{bond_none::doc}}

\subsection{Syntax}
\label{\detokenize{bond_none:syntax}}
\begin{sphinxVerbatim}[commandchars=\\\{\}]
\PYG{k}{bond\PYGZus{}style}\PYG{+w}{ }\PYG{n}{none}
\end{sphinxVerbatim}


\subsection{Examples}
\label{\detokenize{bond_none:examples}}
\begin{sphinxVerbatim}[commandchars=\\\{\}]
\PYG{k}{bond\PYGZus{}style}\PYG{+w}{ }\PYG{n}{none}
\end{sphinxVerbatim}


\subsection{Description}
\label{\detokenize{bond_none:description}}
\sphinxAtStartPar
Using a bond style of none means bond forces and energies are not
computed, even if pairs of bonded atoms were listed in the data file
read by the {\hyperref[\detokenize{read_data::doc}]{\sphinxcrossref{\DUrole{doc}{read\_data}}}} command.

\sphinxAtStartPar
See the {\hyperref[\detokenize{bond_zero::doc}]{\sphinxcrossref{\DUrole{doc}{bond\_style zero}}}} command for a way to
calculate bond statistics, but compute no bond interactions.


\subsection{Restrictions}
\label{\detokenize{bond_none:restrictions}}\begin{quote}

\sphinxAtStartPar
none
\end{quote}


\subsection{Related commands}
\label{\detokenize{bond_none:related-commands}}
\sphinxAtStartPar
none

\sphinxAtStartPar
{\hyperref[\detokenize{bond_zero::doc}]{\sphinxcrossref{\DUrole{doc}{bond\_style zero}}}}


\subsection{Default}
\label{\detokenize{bond_none:default}}
\sphinxAtStartPar
none

\sphinxstepscope

\index{bond\_style nonlinear@\spxentry{bond\_style nonlinear}}\ignorespaces 
\index{bond\_style nonlinear/omp@\spxentry{bond\_style nonlinear/omp}}\ignorespaces 

\section{bond\_style nonlinear command}
\label{\detokenize{bond_nonlinear:bond-style-nonlinear-command}}\label{\detokenize{bond_nonlinear:index-1}}\label{\detokenize{bond_nonlinear:index-0}}\label{\detokenize{bond_nonlinear::doc}}
\sphinxAtStartPar
Accelerator Variants: \sphinxstyleemphasis{nonlinear/omp}


\subsection{Syntax}
\label{\detokenize{bond_nonlinear:syntax}}
\begin{sphinxVerbatim}[commandchars=\\\{\}]
\PYG{k}{bond\PYGZus{}style}\PYG{+w}{ }\PYG{n}{nonlinear}
\end{sphinxVerbatim}


\subsection{Examples}
\label{\detokenize{bond_nonlinear:examples}}
\begin{sphinxVerbatim}[commandchars=\\\{\}]
\PYG{k}{bond\PYGZus{}style}\PYG{+w}{ }\PYG{n}{nonlinear}
\PYG{k}{bond\PYGZus{}coeff}\PYG{+w}{ }\PYG{l+m}{2}\PYG{+w}{ }\PYG{l+m}{100.0}\PYG{+w}{ }\PYG{l+m}{1.1}\PYG{+w}{ }\PYG{l+m}{1.4}
\end{sphinxVerbatim}


\subsection{Description}
\label{\detokenize{bond_nonlinear:description}}
\sphinxAtStartPar
The \sphinxstyleemphasis{nonlinear} bond style uses the potential
\begin{equation*}
\begin{split}E = \frac{\epsilon (r - r_0)^2}{ [ \lambda^2 - (r - r_0)^2 ]}\end{split}
\end{equation*}
\sphinxAtStartPar
to define an anharmonic spring {\hyperref[\detokenize{bond_nonlinear:rector}]{\sphinxcrossref{\DUrole{std,std-ref}{(Rector)}}}} of equilibrium
length \(r_0\) and maximum extension lamda.

\sphinxAtStartPar
The following coefficients must be defined for each bond type via the
{\hyperref[\detokenize{bond_coeff::doc}]{\sphinxcrossref{\DUrole{doc}{bond\_coeff}}}} command as in the example above, or in
the data file or restart files read by the {\hyperref[\detokenize{read_data::doc}]{\sphinxcrossref{\DUrole{doc}{read\_data}}}}
or {\hyperref[\detokenize{read_restart::doc}]{\sphinxcrossref{\DUrole{doc}{read\_restart}}}} commands:
\begin{itemize}
\item {} 
\sphinxAtStartPar
\(\epsilon\) (energy)

\item {} 
\sphinxAtStartPar
\(r_0\) (distance)

\item {} 
\sphinxAtStartPar
\(\lambda\) (distance)

\end{itemize}


\bigskip\hrule\bigskip


\sphinxAtStartPar
Styles with a \sphinxstyleemphasis{gpu}, \sphinxstyleemphasis{intel}, \sphinxstyleemphasis{kk}, \sphinxstyleemphasis{omp}, or \sphinxstyleemphasis{opt} suffix are
functionally the same as the corresponding style without the suffix.
They have been optimized to run faster, depending on your available
hardware, as discussed on the {\hyperref[\detokenize{Speed_packages::doc}]{\sphinxcrossref{\DUrole{doc}{Accelerator packages}}}}
page.  The accelerated styles take the same arguments and should
produce the same results, except for round\sphinxhyphen{}off and precision issues.

\sphinxAtStartPar
These accelerated styles are part of the GPU, INTEL, KOKKOS,
OPENMP, and OPT packages, respectively.  They are only enabled if
LAMMPS was built with those packages.  See the {\hyperref[\detokenize{Build_package::doc}]{\sphinxcrossref{\DUrole{doc}{Build package}}}} page for more info.

\sphinxAtStartPar
You can specify the accelerated styles explicitly in your input script
by including their suffix, or you can use the {\hyperref[\detokenize{Run_options::doc}]{\sphinxcrossref{\DUrole{doc}{\sphinxhyphen{}suffix command\sphinxhyphen{}line switch}}}} when you invoke LAMMPS, or you can use the
{\hyperref[\detokenize{suffix::doc}]{\sphinxcrossref{\DUrole{doc}{suffix}}}} command in your input script.

\sphinxAtStartPar
See the {\hyperref[\detokenize{Speed_packages::doc}]{\sphinxcrossref{\DUrole{doc}{Accelerator packages}}}} page for more
instructions on how to use the accelerated styles effectively.


\bigskip\hrule\bigskip



\subsection{Restrictions}
\label{\detokenize{bond_nonlinear:restrictions}}
\sphinxAtStartPar
This bond style can only be used if LAMMPS was built with the EXTRA\sphinxhyphen{}MOLECULE
package.  See the {\hyperref[\detokenize{Build_package::doc}]{\sphinxcrossref{\DUrole{doc}{Build package}}}} page for more
info.


\subsection{Related commands}
\label{\detokenize{bond_nonlinear:related-commands}}
\sphinxAtStartPar
{\hyperref[\detokenize{bond_coeff::doc}]{\sphinxcrossref{\DUrole{doc}{bond\_coeff}}}}, {\hyperref[\detokenize{delete_bonds::doc}]{\sphinxcrossref{\DUrole{doc}{delete\_bonds}}}}


\subsection{Default}
\label{\detokenize{bond_nonlinear:default}}
\sphinxAtStartPar
none


\bigskip\hrule\bigskip

\phantomsection\label{\detokenize{bond_nonlinear:rector}}
\sphinxAtStartPar
\sphinxstylestrong{(Rector)} Rector, Van Swol, Henderson, Molecular Physics, 82, 1009 (1994).

\sphinxstepscope

\index{bond\_style oxdna/fene@\spxentry{bond\_style oxdna/fene}}\ignorespaces 
\index{bond\_style oxdna2/fene@\spxentry{bond\_style oxdna2/fene}}\ignorespaces 
\index{bond\_style oxrna2/fene@\spxentry{bond\_style oxrna2/fene}}\ignorespaces 

\section{bond\_style oxdna/fene command}
\label{\detokenize{bond_oxdna:bond-style-oxdna-fene-command}}\label{\detokenize{bond_oxdna:index-2}}\label{\detokenize{bond_oxdna:index-1}}\label{\detokenize{bond_oxdna:index-0}}\label{\detokenize{bond_oxdna::doc}}

\section{bond\_style oxdna2/fene command}
\label{\detokenize{bond_oxdna:bond-style-oxdna2-fene-command}}

\section{bond\_style oxrna2/fene command}
\label{\detokenize{bond_oxdna:bond-style-oxrna2-fene-command}}

\subsection{Syntax}
\label{\detokenize{bond_oxdna:syntax}}
\begin{sphinxVerbatim}[commandchars=\\\{\}]
\PYG{k}{bond\PYGZus{}style}\PYG{+w}{ }\PYG{n}{oxdna}\PYG{o}{/}\PYG{n}{fene}

\PYG{k}{bond\PYGZus{}style}\PYG{+w}{ }\PYG{n}{oxdna2}\PYG{o}{/}\PYG{n}{fene}

\PYG{k}{bond\PYGZus{}style}\PYG{+w}{ }\PYG{n}{oxrna2}\PYG{o}{/}\PYG{n}{fene}
\end{sphinxVerbatim}


\subsection{Examples}
\label{\detokenize{bond_oxdna:examples}}
\begin{sphinxVerbatim}[commandchars=\\\{\}]
\PYG{c}{\PYGZsh{} LJ units}
\PYG{k}{bond\PYGZus{}style}\PYG{+w}{ }\PYG{n}{oxdna}\PYG{o}{/}\PYG{n}{fene}
\PYG{k}{bond\PYGZus{}coeff}\PYG{+w}{ }\PYG{o}{*}\PYG{+w}{ }\PYG{l+m}{2.0}\PYG{+w}{ }\PYG{l+m}{0.25}\PYG{+w}{ }\PYG{l+m}{0.7525}

\PYG{k}{bond\PYGZus{}style}\PYG{+w}{ }\PYG{n}{oxdna2}\PYG{o}{/}\PYG{n}{fene}
\PYG{k}{bond\PYGZus{}coeff}\PYG{+w}{ }\PYG{o}{*}\PYG{+w}{ }\PYG{l+m}{2.0}\PYG{+w}{ }\PYG{l+m}{0.25}\PYG{+w}{ }\PYG{l+m}{0.7564}

\PYG{k}{bond\PYGZus{}style}\PYG{+w}{ }\PYG{n}{oxrna2}\PYG{o}{/}\PYG{n}{fene}
\PYG{k}{bond\PYGZus{}coeff}\PYG{+w}{ }\PYG{o}{*}\PYG{+w}{ }\PYG{l+m}{2.0}\PYG{+w}{ }\PYG{l+m}{0.25}\PYG{+w}{ }\PYG{l+m}{0.76107}

\PYG{k}{bond\PYGZus{}style}\PYG{+w}{ }\PYG{n}{oxdna}\PYG{o}{/}\PYG{n}{fene}
\PYG{k}{bond\PYGZus{}coeff}\PYG{+w}{ }\PYG{o}{*}\PYG{+w}{ }\PYG{n}{oxdna\PYGZus{}lj.cgdna}

\PYG{c}{\PYGZsh{} Real units}
\PYG{k}{bond\PYGZus{}style}\PYG{+w}{ }\PYG{n}{oxdna}\PYG{o}{/}\PYG{n}{fene}
\PYG{k}{bond\PYGZus{}coeff}\PYG{+w}{ }\PYG{o}{*}\PYG{+w}{ }\PYG{l+m}{11.92337812042065}\PYG{+w}{ }\PYG{l+m}{2.1295}\PYG{+w}{ }\PYG{l+m}{6.409795}

\PYG{k}{bond\PYGZus{}style}\PYG{+w}{ }\PYG{n}{oxdna2}\PYG{o}{/}\PYG{n}{fene}
\PYG{k}{bond\PYGZus{}coeff}\PYG{+w}{ }\PYG{o}{*}\PYG{+w}{ }\PYG{l+m}{11.92337812042065}\PYG{+w}{ }\PYG{l+m}{2.1295}\PYG{+w}{ }\PYG{l+m}{6.4430152}

\PYG{k}{bond\PYGZus{}style}\PYG{+w}{ }\PYG{n}{oxrna2}\PYG{o}{/}\PYG{n}{fene}
\PYG{k}{bond\PYGZus{}coeff}\PYG{+w}{ }\PYG{o}{*}\PYG{+w}{ }\PYG{l+m}{11.92337812042065}\PYG{+w}{ }\PYG{l+m}{2.1295}\PYG{+w}{ }\PYG{l+m}{6.482800913}

\PYG{k}{bond\PYGZus{}style}\PYG{+w}{ }\PYG{n}{oxrna2}\PYG{o}{/}\PYG{n}{fene}
\PYG{k}{bond\PYGZus{}coeff}\PYG{+w}{ }\PYG{o}{*}\PYG{+w}{ }\PYG{n}{oxrna2\PYGZus{}real.cgdna}
\end{sphinxVerbatim}

\begin{sphinxadmonition}{note}{Note:}
\sphinxAtStartPar
The coefficients in the above examples have to be kept fixed and
cannot be changed without reparameterizing the entire model. They are
provided in forms compatible with both \sphinxstyleemphasis{units lj} and \sphinxstyleemphasis{units real}
(see documentation of {\hyperref[\detokenize{units::doc}]{\sphinxcrossref{\DUrole{doc}{units}}}}).  These can also be read
from a potential file with correct unit style by specifying the name
of the file. Several potential files for each unit style are included
in the \sphinxcode{\sphinxupquote{potentials}} directory of the LAMMPS distribution.
\end{sphinxadmonition}


\subsection{Description}
\label{\detokenize{bond_oxdna:description}}
\sphinxAtStartPar
The \sphinxstyleemphasis{oxdna/fene}, \sphinxstyleemphasis{oxdna2/fene}, and \sphinxstyleemphasis{oxrna2/fene} bond styles use the potential
\begin{equation*}
\begin{split}E = - \frac{\epsilon}{2} \ln \left[ 1 - \left(\frac{r-r_0}{\Delta}\right)^2\right]\end{split}
\end{equation*}
\sphinxAtStartPar
to define a modified finite extensible nonlinear elastic (FENE)
potential {\hyperref[\detokenize{bond_oxdna:ouldridge0}]{\sphinxcrossref{\DUrole{std,std-ref}{(Ouldridge)}}}} to model the connectivity of
the phosphate backbone in the oxDNA/oxRNA force field for coarse\sphinxhyphen{}grained
modelling of DNA/RNA.

\sphinxAtStartPar
The following coefficients must be defined for the bond type via the
{\hyperref[\detokenize{bond_coeff::doc}]{\sphinxcrossref{\DUrole{doc}{bond\_coeff}}}} command as given in the above example, or
in the data file or restart files read by the {\hyperref[\detokenize{read_data::doc}]{\sphinxcrossref{\DUrole{doc}{read\_data}}}} or {\hyperref[\detokenize{read_restart::doc}]{\sphinxcrossref{\DUrole{doc}{read\_restart}}}} commands:
\begin{itemize}
\item {} 
\sphinxAtStartPar
\(\epsilon\) (energy)

\item {} 
\sphinxAtStartPar
\(\Delta\) (distance)

\item {} 
\sphinxAtStartPar
\(r_0\) (distance)

\end{itemize}

\begin{sphinxadmonition}{note}{Note:}
\sphinxAtStartPar
The oxDNA bond style has to be used together with the corresponding
oxDNA pair styles for excluded volume interaction \sphinxstyleemphasis{oxdna/excv} ,
stacking \sphinxstyleemphasis{oxdna/stk} , cross\sphinxhyphen{}stacking \sphinxstyleemphasis{oxdna/xstk} and coaxial
stacking interaction \sphinxstyleemphasis{oxdna/coaxstk} as well as hydrogen\sphinxhyphen{}bonding
interaction \sphinxstyleemphasis{oxdna/hbond} (see also documentation of {\hyperref[\detokenize{pair_oxdna::doc}]{\sphinxcrossref{\DUrole{doc}{pair\_style
oxdna/excv}}}}). For the oxDNA2 {\hyperref[\detokenize{bond_oxdna:snodin0}]{\sphinxcrossref{\DUrole{std,std-ref}{(Snodin)}}}}
bond style the analogous pair styles \sphinxstyleemphasis{oxdna2/excv} , \sphinxstyleemphasis{oxdna2/stk} ,
\sphinxstyleemphasis{oxdna2/xstk} , \sphinxstyleemphasis{oxdna2/coaxstk} , \sphinxstyleemphasis{oxdna2/hbond} and an additional
Debye\sphinxhyphen{}Hueckel pair style \sphinxstyleemphasis{oxdna2/dh} have to be defined. The same
applies to the oxRNA2 {\hyperref[\detokenize{bond_oxdna:sulc01}]{\sphinxcrossref{\DUrole{std,std-ref}{(Sulc1)}}}} styles.
\end{sphinxadmonition}

\begin{sphinxadmonition}{note}{Note:}
\sphinxAtStartPar
This bond style has to be used with the \sphinxstyleemphasis{atom\_style hybrid bond
ellipsoid oxdna} (see documentation of {\hyperref[\detokenize{atom_style::doc}]{\sphinxcrossref{\DUrole{doc}{atom\_style}}}}). The \sphinxstyleemphasis{atom\_style oxdna} stores the 3’\sphinxhyphen{}to\sphinxhyphen{}5’ polarity
of the nucleotide strand, which is set through the bond topology in
the data file. The first (second) atom in a bond definition is
understood to point towards the 3’\sphinxhyphen{}end (5’\sphinxhyphen{}end) of the strand.
\end{sphinxadmonition}

\begin{sphinxadmonition}{warning}{Warning:}
\sphinxAtStartPar
If data files are produced with {\hyperref[\detokenize{write_data::doc}]{\sphinxcrossref{\DUrole{doc}{write\_data}}}}, then
the {\hyperref[\detokenize{newton::doc}]{\sphinxcrossref{\DUrole{doc}{newton}}}} command should be set to \sphinxstyleemphasis{newton on} or
\sphinxstyleemphasis{newton off on}.  Otherwise the data files will not have the same
3’\sphinxhyphen{}to\sphinxhyphen{}5’ polarity as the initial data file. This limitation does not
apply to binary restart files produced with {\hyperref[\detokenize{write_restart::doc}]{\sphinxcrossref{\DUrole{doc}{write\_restart}}}}.
\end{sphinxadmonition}

\sphinxAtStartPar
Example input and data files for DNA and RNA duplexes can be found in
\sphinxcode{\sphinxupquote{examples/PACKAGES/cgdna/examples/oxDNA/\textasciigrave{}, \textasciigrave{}.../oxDNA2/}} and
\sphinxcode{\sphinxupquote{.../oxRNA2/}}.  A simple python setup tool which creates single
straight or helical DNA strands, DNA/RNA duplexes or arrays of DNA/RNA
duplexes can be found in \sphinxcode{\sphinxupquote{examples/PACKAGES/cgdna/util/}}.

\sphinxAtStartPar
Please cite {\hyperref[\detokenize{bond_oxdna:henrich0}]{\sphinxcrossref{\DUrole{std,std-ref}{(Henrich)}}}} in any publication that uses
this implementation. An updated documentation that contains general information
on the model, its implementation and performance as well as the structure of
the data and input file can be found \sphinxhref{PDF/CG-DNA.pdf}{here}.

\sphinxAtStartPar
Please cite also the relevant oxDNA/oxRNA publications. These are
{\hyperref[\detokenize{bond_oxdna:ouldridge0}]{\sphinxcrossref{\DUrole{std,std-ref}{(Ouldridge)}}}} and
{\hyperref[\detokenize{bond_oxdna:ouldridge-dphil0}]{\sphinxcrossref{\DUrole{std,std-ref}{(Ouldridge\sphinxhyphen{}DPhil)}}}} for oxDNA,
{\hyperref[\detokenize{bond_oxdna:snodin0}]{\sphinxcrossref{\DUrole{std,std-ref}{(Snodin)}}}} for oxDNA2,
{\hyperref[\detokenize{bond_oxdna:sulc01}]{\sphinxcrossref{\DUrole{std,std-ref}{(Sulc1)}}}} for oxRNA2
and for sequence\sphinxhyphen{}specific hydrogen\sphinxhyphen{}bonding and stacking interactions
{\hyperref[\detokenize{bond_oxdna:sulc02}]{\sphinxcrossref{\DUrole{std,std-ref}{(Sulc2)}}}}.


\bigskip\hrule\bigskip



\subsection{Potential file reading}
\label{\detokenize{bond_oxdna:potential-file-reading}}
\sphinxAtStartPar
For each style oxdna, oxdna2 and oxrna2, the first parameter argument
can be a filename, and if it is, no further arguments should be
supplied. Therefore the following command:

\begin{sphinxVerbatim}[commandchars=\\\{\}]
\PYG{k}{bond\PYGZus{}style}\PYG{+w}{ }\PYG{n}{oxdna}\PYG{o}{/}\PYG{n}{fene}
\PYG{k}{bond\PYGZus{}coeff}\PYG{+w}{ }\PYG{o}{*}\PYG{+w}{ }\PYG{n}{oxdna\PYGZus{}lj.cgdna}
\end{sphinxVerbatim}

\sphinxAtStartPar
will be interpreted as a request to read the (FENE) potential
{\hyperref[\detokenize{bond_oxdna:ouldridge0}]{\sphinxcrossref{\DUrole{std,std-ref}{(Ouldridge)}}}} parameters from the file with the given
name.  The file can define multiple potential parameters for both bonded
and pair interactions, but for the above bonded interactions there must
exist in the file a line of the form:

\begin{sphinxVerbatim}[commandchars=\\\{\}]
\PYG{o}{*}\PYG{+w}{   }\PYG{n}{fene}\PYG{+w}{    }\PYG{n}{epsilon}\PYG{+w}{ }\PYG{n}{delta}\PYG{+w}{ }\PYG{n}{r0}
\end{sphinxVerbatim}

\sphinxAtStartPar
There are sample potential files for each unit style in the
\sphinxcode{\sphinxupquote{potentials}} directory of the LAMMPS distribution. The potential file
unit system must align with the units defined via the {\hyperref[\detokenize{units::doc}]{\sphinxcrossref{\DUrole{doc}{units}}}} command. For conversion between different \sphinxstyleemphasis{LJ} and \sphinxstyleemphasis{real} unit
systems for oxDNA, the python tool \sphinxstyleemphasis{lj2real.py} located in the
\sphinxcode{\sphinxupquote{examples/PACKAGES/cgdna/util/}} directory can be used. This tool
assumes similar file structure to the examples found in
\sphinxcode{\sphinxupquote{examples/PACKAGES/cgdna/examples/}}.


\bigskip\hrule\bigskip



\subsection{Restrictions}
\label{\detokenize{bond_oxdna:restrictions}}
\sphinxAtStartPar
This bond style can only be used if LAMMPS was built with the
CG\sphinxhyphen{}DNA package and the MOLECULE and ASPHERE package.  See the
{\hyperref[\detokenize{Build_package::doc}]{\sphinxcrossref{\DUrole{doc}{Build package}}}} page for more info.


\subsection{Related commands}
\label{\detokenize{bond_oxdna:related-commands}}
\sphinxAtStartPar
{\hyperref[\detokenize{pair_oxdna::doc}]{\sphinxcrossref{\DUrole{doc}{pair\_style oxdna/excv}}}}, {\hyperref[\detokenize{pair_oxdna2::doc}]{\sphinxcrossref{\DUrole{doc}{pair\_style oxdna2/excv}}}}, {\hyperref[\detokenize{pair_oxrna2::doc}]{\sphinxcrossref{\DUrole{doc}{pair\_style oxrna2/excv}}}},
{\hyperref[\detokenize{bond_coeff::doc}]{\sphinxcrossref{\DUrole{doc}{bond\_coeff}}}}, {\hyperref[\detokenize{atom_style::doc}]{\sphinxcrossref{\DUrole{doc}{atom\_style oxdna}}}}, {\hyperref[\detokenize{fix_nve_dotc_langevin::doc}]{\sphinxcrossref{\DUrole{doc}{fix nve/dotc/langevin}}}}


\subsection{Default}
\label{\detokenize{bond_oxdna:default}}
\sphinxAtStartPar
none


\bigskip\hrule\bigskip

\phantomsection\label{\detokenize{bond_oxdna:henrich0}}
\sphinxAtStartPar
\sphinxstylestrong{(Henrich)} O. Henrich, Y. A. Gutierrez\sphinxhyphen{}Fosado, T. Curk, T. E. Ouldridge, Eur. Phys. J. E 41, 57 (2018).

\phantomsection\label{\detokenize{bond_oxdna:ouldridge-dphil0}}
\sphinxAtStartPar
\sphinxstylestrong{(Ouldridge\sphinxhyphen{}DPhil)} T.E. Ouldridge, Coarse\sphinxhyphen{}grained modelling of DNA and DNA self\sphinxhyphen{}assembly, DPhil. University of Oxford (2011).

\phantomsection\label{\detokenize{bond_oxdna:ouldridge0}}
\sphinxAtStartPar
\sphinxstylestrong{(Ouldridge)} T.E. Ouldridge, A.A. Louis, J.P.K. Doye, J. Chem. Phys. 134, 085101 (2011).

\phantomsection\label{\detokenize{bond_oxdna:snodin0}}
\sphinxAtStartPar
\sphinxstylestrong{(Snodin)} B.E. Snodin, F. Randisi, M. Mosayebi, et al., J. Chem. Phys. 142, 234901 (2015).

\phantomsection\label{\detokenize{bond_oxdna:sulc01}}
\sphinxAtStartPar
\sphinxstylestrong{(Sulc1)} P. Sulc, F. Romano, T. E. Ouldridge, et al., J. Chem. Phys. 140, 235102 (2014).

\phantomsection\label{\detokenize{bond_oxdna:sulc02}}
\sphinxAtStartPar
\sphinxstylestrong{(Sulc2)} P. Sulc, F. Romano, T.E. Ouldridge, L. Rovigatti, J.P.K. Doye, A.A. Louis, J. Chem. Phys. 137, 135101 (2012).

\sphinxstepscope

\index{bond\_style quartic@\spxentry{bond\_style quartic}}\ignorespaces 
\index{bond\_style quartic/omp@\spxentry{bond\_style quartic/omp}}\ignorespaces 

\section{bond\_style quartic command}
\label{\detokenize{bond_quartic:bond-style-quartic-command}}\label{\detokenize{bond_quartic:index-1}}\label{\detokenize{bond_quartic:index-0}}\label{\detokenize{bond_quartic::doc}}
\sphinxAtStartPar
Accelerator Variants: \sphinxstyleemphasis{quartic/omp}


\subsection{Syntax}
\label{\detokenize{bond_quartic:syntax}}
\begin{sphinxVerbatim}[commandchars=\\\{\}]
\PYG{k}{bond\PYGZus{}style}\PYG{+w}{ }\PYG{n}{quartic}
\end{sphinxVerbatim}


\subsection{Examples}
\label{\detokenize{bond_quartic:examples}}
\begin{sphinxVerbatim}[commandchars=\\\{\}]
\PYG{k}{bond\PYGZus{}style}\PYG{+w}{ }\PYG{n}{quartic}
\PYG{k}{bond\PYGZus{}coeff}\PYG{+w}{ }\PYG{l+m}{2}\PYG{+w}{ }\PYG{l+m}{1200}\PYG{+w}{ }\PYG{o}{\PYGZhy{}}\PYG{l+m}{0.55}\PYG{+w}{ }\PYG{l+m}{0.25}\PYG{+w}{ }\PYG{l+m}{1.3}\PYG{+w}{ }\PYG{l+m}{34.6878}
\end{sphinxVerbatim}


\subsection{Description}
\label{\detokenize{bond_quartic:description}}
\sphinxAtStartPar
The \sphinxstyleemphasis{quartic} bond style uses the potential
\begin{equation*}
\begin{split}E      & = E_q + E_{LJ} \\
E_q    & = K (r - R_c)^ 2 (r - R_c - B_1) (r - R_c - B_2) + U_0 \\
E_{LJ} & = \left\{ \begin{array} {l@{\quad:\quad}l}
4 \epsilon \left[ \left(\frac{\sigma}{r}\right)^{12} - \left(\frac{\sigma}{r}\right)^6 \right] + \epsilon & r < 2^{\frac{1}{6}}, \epsilon = 1, \sigma = 1 \\
                                               0 & r >= 2^{\frac{1}{6}}
                      \end{array} \right.\end{split}
\end{equation*}
\sphinxAtStartPar
to define a bond that can be broken as the simulation proceeds (e.g.
due to a polymer being stretched).  The \(\sigma\) and
\(\epsilon\) used in the LJ portion of the formula are both set
equal to 1.0 by LAMMPS and the LJ portion is cut off at its minimum,
i.e. at \(r_c = 2^{\frac{1}{6}}\).

\sphinxAtStartPar
The following coefficients must be defined for each bond type via the
{\hyperref[\detokenize{bond_coeff::doc}]{\sphinxcrossref{\DUrole{doc}{bond\_coeff}}}} command as in the example above, or in
the data file or restart files read by the {\hyperref[\detokenize{read_data::doc}]{\sphinxcrossref{\DUrole{doc}{read\_data}}}}
or {\hyperref[\detokenize{read_restart::doc}]{\sphinxcrossref{\DUrole{doc}{read\_restart}}}} commands:
\begin{itemize}
\item {} 
\sphinxAtStartPar
\(K\) (energy/distance\textasciicircum{}4)

\item {} 
\sphinxAtStartPar
\(B_1\) (distance)

\item {} 
\sphinxAtStartPar
\(B_2\) (distance)

\item {} 
\sphinxAtStartPar
\(R_c\) (distance)

\item {} 
\sphinxAtStartPar
\(U_0\) (energy)

\end{itemize}

\sphinxAtStartPar
This potential was constructed to mimic the FENE bond potential for
coarse\sphinxhyphen{}grained polymer chains.  When monomers with \(\sigma =
\epsilon = 1.0\) are used, the following choice of parameters gives a
quartic potential that looks nearly like the FENE potential:
\begin{equation*}
\begin{split}K &= 1200 \\
B_1 &= -0.55 \\
B_2 &= 0.25 \\
R_c &= 1.3 \\
U_0 &= 34.6878\end{split}
\end{equation*}
\sphinxAtStartPar
Different parameters can be specified using the {\hyperref[\detokenize{bond_coeff::doc}]{\sphinxcrossref{\DUrole{doc}{bond\_coeff}}}}
command, but you will need to choose them carefully so they form a suitable
bond potential.

\sphinxAtStartPar
\(R_c\) is the cutoff length at which the bond potential goes smoothly to a
local maximum.  If a bond length ever becomes \(> R_c\), LAMMPS “breaks”
the bond, which means two things.  First, the bond potential is turned
off by setting its type to 0, and is no longer computed.  Second, a
pairwise interaction between the two atoms is turned on, since they
are no longer bonded. See the {\hyperref[\detokenize{Howto_broken_bonds::doc}]{\sphinxcrossref{\DUrole{doc}{Howto}}}} page
on broken bonds for more information.

\sphinxAtStartPar
LAMMPS does the second task via a computational sleight\sphinxhyphen{}of\sphinxhyphen{}hand.  It
subtracts the pairwise interaction as part of the bond computation.
When the bond breaks, the subtraction stops.  For this to work, the
pairwise interaction must always be computed by the
{\hyperref[\detokenize{pair_style::doc}]{\sphinxcrossref{\DUrole{doc}{pair\_style}}}} command, whether the bond is broken or
not.  This means that {\hyperref[\detokenize{special_bonds::doc}]{\sphinxcrossref{\DUrole{doc}{special\_bonds}}}} must be set
to 1,1,1, as indicated as a restriction below.

\sphinxAtStartPar
Note that when bonds are dumped to a file via the {\hyperref[\detokenize{dump::doc}]{\sphinxcrossref{\DUrole{doc}{dump local}}}} command, bonds with type 0 are not included.  The
{\hyperref[\detokenize{delete_bonds::doc}]{\sphinxcrossref{\DUrole{doc}{delete\_bonds}}}} command can also be used to query the
status of broken bonds or permanently delete them, e.g.:

\begin{sphinxVerbatim}[commandchars=\\\{\}]
\PYG{k}{delete\PYGZus{}bonds }\PYG{n+nv+nvIdentifier}{all}\PYG{+w}{ }\PYG{n}{stats}
\PYG{k}{delete\PYGZus{}bonds }\PYG{n+nv+nvIdentifier}{all}\PYG{+w}{ }\PYG{n}{bond}\PYG{+w}{ }\PYG{l+m}{0}\PYG{+w}{ }\PYG{n}{remove}
\end{sphinxVerbatim}


\bigskip\hrule\bigskip


\sphinxAtStartPar
Styles with a \sphinxstyleemphasis{gpu}, \sphinxstyleemphasis{intel}, \sphinxstyleemphasis{kk}, \sphinxstyleemphasis{omp}, or \sphinxstyleemphasis{opt} suffix are
functionally the same as the corresponding style without the suffix.
They have been optimized to run faster, depending on your available
hardware, as discussed on the {\hyperref[\detokenize{Speed_packages::doc}]{\sphinxcrossref{\DUrole{doc}{Accelerator packages}}}}
page.  The accelerated styles take the same arguments and should
produce the same results, except for round\sphinxhyphen{}off and precision issues.

\sphinxAtStartPar
These accelerated styles are part of the GPU, INTEL, KOKKOS,
OPENMP, and OPT packages, respectively.  They are only enabled if
LAMMPS was built with those packages.  See the {\hyperref[\detokenize{Build_package::doc}]{\sphinxcrossref{\DUrole{doc}{Build package}}}} page for more info.

\sphinxAtStartPar
You can specify the accelerated styles explicitly in your input script
by including their suffix, or you can use the {\hyperref[\detokenize{Run_options::doc}]{\sphinxcrossref{\DUrole{doc}{\sphinxhyphen{}suffix command\sphinxhyphen{}line switch}}}} when you invoke LAMMPS, or you can use the
{\hyperref[\detokenize{suffix::doc}]{\sphinxcrossref{\DUrole{doc}{suffix}}}} command in your input script.

\sphinxAtStartPar
See the {\hyperref[\detokenize{Speed_packages::doc}]{\sphinxcrossref{\DUrole{doc}{Accelerator packages}}}} page for more
instructions on how to use the accelerated styles effectively.


\bigskip\hrule\bigskip



\subsection{Restrictions}
\label{\detokenize{bond_quartic:restrictions}}
\sphinxAtStartPar
This bond style can only be used if LAMMPS was built with the MOLECULE
package.  See the {\hyperref[\detokenize{Build_package::doc}]{\sphinxcrossref{\DUrole{doc}{Build package}}}} page for more
info.

\sphinxAtStartPar
The \sphinxstyleemphasis{quartic} style requires that {\hyperref[\detokenize{special_bonds::doc}]{\sphinxcrossref{\DUrole{doc}{special\_bonds}}}}
parameters be set to 1,1,1.  Three\sphinxhyphen{} and four\sphinxhyphen{}body interactions (angle,
dihedral, etc) cannot be used with \sphinxstyleemphasis{quartic} bonds.


\subsection{Related commands}
\label{\detokenize{bond_quartic:related-commands}}
\sphinxAtStartPar
{\hyperref[\detokenize{bond_coeff::doc}]{\sphinxcrossref{\DUrole{doc}{bond\_coeff}}}}, {\hyperref[\detokenize{delete_bonds::doc}]{\sphinxcrossref{\DUrole{doc}{delete\_bonds}}}}


\subsection{Default}
\label{\detokenize{bond_quartic:default}}
\sphinxAtStartPar
none

\sphinxstepscope

\index{bond\_style rheo/shell@\spxentry{bond\_style rheo/shell}}\ignorespaces 

\section{bond\_style rheo/shell command}
\label{\detokenize{bond_rheo_shell:bond-style-rheo-shell-command}}\label{\detokenize{bond_rheo_shell:index-0}}\label{\detokenize{bond_rheo_shell::doc}}

\subsection{Syntax}
\label{\detokenize{bond_rheo_shell:syntax}}
\begin{sphinxVerbatim}[commandchars=\\\{\}]
\PYG{k}{bond\PYGZus{}style}\PYG{+w}{ }\PYG{n}{rheo}\PYG{o}{/}\PYG{n}{shell}\PYG{+w}{ }\PYG{n}{keyword}\PYG{+w}{ }\PYG{n}{value}\PYG{+w}{ }\PYG{n}{attribute1}\PYG{+w}{ }\PYG{n}{attribute2}\PYG{+w}{ }\PYG{n}{...}
\end{sphinxVerbatim}
\begin{itemize}
\item {} 
\sphinxAtStartPar
required keyword = \sphinxstyleemphasis{t/form}

\item {} 
\sphinxAtStartPar
optional keyword = \sphinxstyleemphasis{store/local}
\begin{sphinxalltt}
\sphinxstyleemphasis{t/form} value = formation time for a bond (time units)

\sphinxstyleemphasis{store/local} values = fix\_ID N attributes ...
   * fix\_ID = ID of associated internal fix to store data
   * N = prepare data for output every this many timesteps
   * attributes = zero or more of the below attributes may be appended

     \sphinxstyleemphasis{id1, id2} = IDs of two atoms in the bond
     \sphinxstyleemphasis{time} = the timestep the bond broke
     \sphinxstyleemphasis{x, y, z} = the center of mass position of the two atoms when the bond broke (distance units)
     \sphinxstyleemphasis{x/ref, y/ref, z/ref} = the initial center of mass position of the two atoms (distance units)
\end{sphinxalltt}

\end{itemize}


\subsection{Examples}
\label{\detokenize{bond_rheo_shell:examples}}
\begin{sphinxVerbatim}[commandchars=\\\{\}]
\PYG{k}{bond\PYGZus{}style}\PYG{+w}{ }\PYG{n}{rheo}\PYG{o}{/}\PYG{n}{shell}\PYG{+w}{ }\PYG{n}{t}\PYG{o}{/}\PYG{n}{form}\PYG{+w}{ }\PYG{l+m}{10.0}
\PYG{k}{bond\PYGZus{}coeff}\PYG{+w}{ }\PYG{l+m}{1}\PYG{+w}{ }\PYG{l+m}{1.0}\PYG{+w}{ }\PYG{l+m}{0.05}\PYG{+w}{ }\PYG{l+m}{0.1}
\end{sphinxVerbatim}


\subsection{Description}
\label{\detokenize{bond_rheo_shell:description}}
\sphinxAtStartPar
\DUrole{versionmodified,added}{Added in version 29Aug2024.}

\sphinxAtStartPar
The \sphinxstyleemphasis{rheo/shell} bond style is designed to work with
{\hyperref[\detokenize{fix_rheo_oxidation::doc}]{\sphinxcrossref{\DUrole{doc}{fix rheo/oxidation}}}} which creates candidate
bonds between eligible surface or near\sphinxhyphen{}surface particles. When a bond
is first created, it computes no forces and starts a timer. Forces are
not computed until the timer reaches the specified bond formation time,
\sphinxstyleemphasis{t/form}, and the bond is enabled and applies forces. If the two particles
move outside of the maximum bond distance or move into the bulk before
the timer reaches \sphinxstyleemphasis{t/form}, the bond automatically deletes itself. This
deletion is not recorded as a broken bond in the optional \sphinxstyleemphasis{store/local} fix.

\sphinxAtStartPar
Before bonds are enabled, they are still treated as regular bonds by
all other parts of LAMMPS. This means they are written to data files
and counted in computes such as {\hyperref[\detokenize{compute_nbond_atom::doc}]{\sphinxcrossref{\DUrole{doc}{nbond/atom}}}}.
To only count enabled bonds, use the \sphinxstyleemphasis{nbond/shell} attribute in
{\hyperref[\detokenize{compute_rheo_property_atom::doc}]{\sphinxcrossref{\DUrole{doc}{compute rheo/property/atom}}}}.

\sphinxAtStartPar
When enabled, the bond then computes forces based on deviations from
the initial reference state of the two atoms much like a BPM style
bond (as further discussed in the {\hyperref[\detokenize{Howto_bpm::doc}]{\sphinxcrossref{\DUrole{doc}{BPM howto page}}}}).
The reference state is stored by each bond when it is first enabled.
Data is then preserved across run commands and is written to
{\hyperref[\detokenize{restart::doc}]{\sphinxcrossref{\DUrole{doc}{binary restart files}}}} such that restarting the system
will not reset the reference state of a bond or the timer.

\sphinxAtStartPar
This bond style is based on a model described in
{\hyperref[\detokenize{bond_rheo_shell:rheo-clemmer}]{\sphinxcrossref{\DUrole{std,std-ref}{(Clemmer)}}}}. The force has a magnitude of
\begin{equation*}
\begin{split}F = 2 k (r - r_0) + \frac{2 k}{r_0^2 \epsilon_c^2} (r - r_0)^3\end{split}
\end{equation*}
\sphinxAtStartPar
where \(k\) is a stiffness, \(r\) is the current distance
and \(r_0\) is the initial distance between the two particles, and
\(\epsilon_c\) is maximum strain beyond which a bond breaks. This
is done by setting the bond type to 0 such that forces are no longer
computed.

\sphinxAtStartPar
A damping force proportional to the difference in the normal velocity
of particles is also applied to bonded particles:
\begin{equation*}
\begin{split}F_D = - \gamma w (\hat{r} \bullet \vec{v})\end{split}
\end{equation*}
\sphinxAtStartPar
where \(\gamma\) is the damping strength, \(\hat{r}\) is the
displacement normal vector, and \(\vec{v}\) is the velocity difference
between the two particles.

\sphinxAtStartPar
The following coefficients must be defined for each bond type via the
{\hyperref[\detokenize{bond_coeff::doc}]{\sphinxcrossref{\DUrole{doc}{bond\_coeff}}}} command as in the example above, or in
the data file or restart files read by the {\hyperref[\detokenize{read_data::doc}]{\sphinxcrossref{\DUrole{doc}{read\_data}}}} or {\hyperref[\detokenize{read_restart::doc}]{\sphinxcrossref{\DUrole{doc}{read\_restart}}}} commands:
\begin{itemize}
\item {} 
\sphinxAtStartPar
\(k\)             (force/distance units)

\item {} 
\sphinxAtStartPar
\(\epsilon_c\)    (unit less)

\item {} 
\sphinxAtStartPar
\(\gamma\)        (force/velocity units)

\end{itemize}

\sphinxAtStartPar
Unlike other BPM\sphinxhyphen{}style bonds, this bond style does not update special
bond settings when bonds are created or deleted. This bond style also
does not enforce specific {\hyperref[\detokenize{special_bonds::doc}]{\sphinxcrossref{\DUrole{doc}{special\_bonds}}}} settings.
This behavior is purposeful such {\hyperref[\detokenize{pair_rheo::doc}]{\sphinxcrossref{\DUrole{doc}{RHEO pair}}}} forces
and heat flows are still calculated.

\sphinxAtStartPar
If the \sphinxstyleemphasis{store/local} keyword is used, an internal fix will track bonds that
break during the simulation. Whenever a bond breaks, data is processed
and transferred to an internal fix labeled \sphinxstyleemphasis{fix\_ID}. This allows the
local data to be accessed by other LAMMPS commands. Following this optional
keyword, a list of one or more attributes is specified.  These include the
IDs of the two atoms in the bond. The other attributes for the two atoms
include the timestep during which the bond broke and the current/initial
center of mass position of the two atoms.

\sphinxAtStartPar
Data is continuously accumulated over intervals of \sphinxstyleemphasis{N}
timesteps. At the end of each interval, all of the saved accumulated
data is deleted to make room for new data. Individual datum may
therefore persist anywhere between \sphinxstyleemphasis{1} to \sphinxstyleemphasis{N} timesteps depending on
when they are saved. This data can be accessed using the \sphinxstyleemphasis{fix\_ID} and a
{\hyperref[\detokenize{dump::doc}]{\sphinxcrossref{\DUrole{doc}{dump local}}}} command. To ensure all data is output,
the dump frequency should correspond to the same interval of \sphinxstyleemphasis{N}
timesteps. A dump frequency of an integer multiple of \sphinxstyleemphasis{N} can be used
to regularly output a sample of the accumulated data.

\sphinxAtStartPar
Note that when unbroken bonds are dumped to a file via the
{\hyperref[\detokenize{dump::doc}]{\sphinxcrossref{\DUrole{doc}{dump local}}}} command, bonds with type 0 (broken bonds)
are not included.
The {\hyperref[\detokenize{delete_bonds::doc}]{\sphinxcrossref{\DUrole{doc}{delete\_bonds}}}} command can also be used to
query the status of broken bonds or permanently delete them, e.g.:

\begin{sphinxVerbatim}[commandchars=\\\{\}]
\PYG{k}{delete\PYGZus{}bonds }\PYG{n+nv+nvIdentifier}{all}\PYG{+w}{ }\PYG{n}{stats}
\PYG{k}{delete\PYGZus{}bonds }\PYG{n+nv+nvIdentifier}{all}\PYG{+w}{ }\PYG{n}{bond}\PYG{+w}{ }\PYG{l+m}{0}\PYG{+w}{ }\PYG{n}{remove}
\end{sphinxVerbatim}


\bigskip\hrule\bigskip



\subsection{Restart and other info}
\label{\detokenize{bond_rheo_shell:restart-and-other-info}}
\sphinxAtStartPar
This bond style writes the reference state of each bond to
{\hyperref[\detokenize{restart::doc}]{\sphinxcrossref{\DUrole{doc}{binary restart files}}}}. Loading a restart
file will properly restore bonds. However, the reference state is NOT
written to data files. Therefore reading a data file will not
restore bonds and will cause their reference states to be redefined.

\sphinxAtStartPar
If the \sphinxstyleemphasis{store/local} option is used, an internal fix will calculate
a local vector or local array depending on the number of input values.
The length of the vector or number of rows in the array is the number
of recorded, broken bonds.  If a single input is specified, a local
vector is produced. If two or more inputs are specified, a local array
is produced where the number of columns = the number of inputs.  The
vector or array can be accessed by any command that uses local values
from a compute as input. See the {\hyperref[\detokenize{Howto_output::doc}]{\sphinxcrossref{\DUrole{doc}{Howto output}}}} page
for an overview of LAMMPS output options.

\sphinxAtStartPar
The vector or array will be floating point values that correspond to
the specified attribute.

\sphinxAtStartPar
The single() function of this bond style returns 0.0 for the energy
of a bonded interaction, since energy is not conserved in these
dissipative potentials.  The single() function also calculates two
extra bond quantities, the initial distance \(r_0\) and a time.
These extra quantities can be accessed by the
{\hyperref[\detokenize{compute_bond_local::doc}]{\sphinxcrossref{\DUrole{doc}{compute bond/local}}}} command as \sphinxstyleemphasis{b1} and \sphinxstyleemphasis{b2}.


\subsection{Restrictions}
\label{\detokenize{bond_rheo_shell:restrictions}}
\sphinxAtStartPar
This bond style is part of the RHEO package.  It is only enabled if
LAMMPS was built with that package.  See the {\hyperref[\detokenize{Build_package::doc}]{\sphinxcrossref{\DUrole{doc}{Build package}}}} page for more info.


\subsection{Related commands}
\label{\detokenize{bond_rheo_shell:related-commands}}
\sphinxAtStartPar
{\hyperref[\detokenize{bond_coeff::doc}]{\sphinxcrossref{\DUrole{doc}{bond\_coeff}}}}, {\hyperref[\detokenize{fix_rheo_oxidation::doc}]{\sphinxcrossref{\DUrole{doc}{fix rheo/oxidation}}}}


\subsection{Default}
\label{\detokenize{bond_rheo_shell:default}}
\sphinxAtStartPar
NA


\bigskip\hrule\bigskip

\phantomsection\label{\detokenize{bond_rheo_shell:rheo-clemmer}}
\sphinxAtStartPar
\sphinxstylestrong{(Clemmer)} Clemmer, Pierce, O’Connor, Nevins, Jones, Lechman, Tencer, Appl. Math. Model., 130, 310\sphinxhyphen{}326 (2024).

\sphinxstepscope

\index{bond\_style special@\spxentry{bond\_style special}}\ignorespaces 

\section{bond\_style special command}
\label{\detokenize{bond_special:bond-style-special-command}}\label{\detokenize{bond_special:index-0}}\label{\detokenize{bond_special::doc}}

\subsection{Syntax}
\label{\detokenize{bond_special:syntax}}
\begin{sphinxVerbatim}[commandchars=\\\{\}]
\PYG{k}{bond\PYGZus{}style}\PYG{+w}{ }\PYG{n}{special}
\end{sphinxVerbatim}


\subsection{Examples}
\label{\detokenize{bond_special:examples}}
\begin{sphinxVerbatim}[commandchars=\\\{\}]
\PYG{k}{bond\PYGZus{}style}\PYG{+w}{ }\PYG{n}{special}
\PYG{k}{bond\PYGZus{}coeff}\PYG{+w}{ }\PYG{l+m}{0.5}\PYG{+w}{ }\PYG{l+m}{0.5}
\end{sphinxVerbatim}


\subsection{Description}
\label{\detokenize{bond_special:description}}
\sphinxAtStartPar
The \sphinxstyleemphasis{special} bond style can be used to create conceptual bonds which
effectively impose weightings on the pairwise Lennard Jones and/or
Coulombic interactions between selected pairs of particles in the
system.  The form of the pairwise interaction will be whatever is
computed by the {\hyperref[\detokenize{pair_style::doc}]{\sphinxcrossref{\DUrole{doc}{pair\_style}}}} command defined for the
system; this command defines the weightings for its two terms.

\sphinxAtStartPar
This command can thus be useful to apply weightings that cannot be
handled by the {\hyperref[\detokenize{special_bonds::doc}]{\sphinxcrossref{\DUrole{doc}{special\_bonds}}}} command, such as
on 1\sphinxhyphen{}5 or 1\sphinxhyphen{}6 interactions.  Or it can be used to add pairwise forces
between one or more pairs of atoms that otherwise would not be include
in the {\hyperref[\detokenize{pair_style::doc}]{\sphinxcrossref{\DUrole{doc}{pair\_style}}}} computation.

\sphinxAtStartPar
The potential for this bond style has the form
\begin{equation*}
\begin{split}E =  w_{LJ} E_{LJ} + w_{Coul} E_{Coul}\end{split}
\end{equation*}
\sphinxAtStartPar
The following coefficients must be defined for each bond type via the
{\hyperref[\detokenize{bond_coeff::doc}]{\sphinxcrossref{\DUrole{doc}{bond\_coeff}}}} command as in the example above, or in
the data file or restart files read by the {\hyperref[\detokenize{read_data::doc}]{\sphinxcrossref{\DUrole{doc}{read\_data}}}}
or {\hyperref[\detokenize{read_restart::doc}]{\sphinxcrossref{\DUrole{doc}{read\_restart}}}} commands:
\begin{itemize}
\item {} 
\sphinxAtStartPar
\(w_{LJ}\) weight (0.0 to 1.0) on pairwise Lennard\sphinxhyphen{}Jones interactions

\item {} 
\sphinxAtStartPar
\(w_{Coul}\) weight (0.0 to 1.0) on pairwise Coulombic interactions

\end{itemize}


\bigskip\hrule\bigskip


\sphinxAtStartPar
Normally this bond style should be used in conjunction with one (or
more) other bond styles which compute forces between atoms directly
bonded to each other in a molecule.  This means the {\hyperref[\detokenize{bond_hybrid::doc}]{\sphinxcrossref{\DUrole{doc}{bond\_style
hybrid}}}} command should be used with bond\_style special
as one of its sub\sphinxhyphen{}styles.

\sphinxAtStartPar
Note that the same as for any other bond style, pairs of bonded atoms
must be enumerated in the data file read by the {\hyperref[\detokenize{read_data::doc}]{\sphinxcrossref{\DUrole{doc}{read\_data}}}} command.  Thus if this command is used to weight all 1\sphinxhyphen{}5
interactions in the system, all the 1\sphinxhyphen{}5 pairs of atoms must be listed
in the “Bonds” section of the data file.

\sphinxAtStartPar
This bond style imposes strict requirements on settings made with the
{\hyperref[\detokenize{special_bonds::doc}]{\sphinxcrossref{\DUrole{doc}{special\_bonds}}}} command.  These requirements
ensure that the new bonds created by this style do not create spurious
1\sphinxhyphen{}2, 1\sphinxhyphen{}3, or 1\sphinxhyphen{}4 interactions within the molecular topology.

\sphinxAtStartPar
Specifically 1\sphinxhyphen{}2 interactions must have weights of zero, 1\sphinxhyphen{}3
interactions must either have weights of unity or {\hyperref[\detokenize{special_bonds::doc}]{\sphinxcrossref{\DUrole{doc}{special\_bonds
angle yes}}}} must be used, and 1\sphinxhyphen{}4 interactions must
have weights of unity or {\hyperref[\detokenize{special_bonds::doc}]{\sphinxcrossref{\DUrole{doc}{special\_bonds dihedral yes}}}}
must be used.

\sphinxAtStartPar
If this command is used to create bonded interactions between
particles that are further apart than usual (e.g. 1\sphinxhyphen{}5 or 1\sphinxhyphen{}6
interactions), this style may require an increase in the communication
cutoff via the {\hyperref[\detokenize{comm_modify::doc}]{\sphinxcrossref{\DUrole{doc}{comm\_modify cutoff}}}} command.  If
LAMMPS cannot find a partner atom in a bond, an error will be issued.


\subsection{Restrictions}
\label{\detokenize{bond_special:restrictions}}
\sphinxAtStartPar
This bond style can only be used if LAMMPS was built with the
MISC package.  See the {\hyperref[\detokenize{Build_package::doc}]{\sphinxcrossref{\DUrole{doc}{Build package}}}} doc
page for more info.

\sphinxAtStartPar
This bond style requires the use of a {\hyperref[\detokenize{pair_style::doc}]{\sphinxcrossref{\DUrole{doc}{pair\_style}}}} which
computes a pairwise additive interaction and provides the ability to
compute interactions for individual pairs of atoms.  Manybody potentials
are not compatible in general, but also some other pair styles are missing
the required functionality and thus will cause an error.

\sphinxAtStartPar
This command is not compatible with long\sphinxhyphen{}range Coulombic interactions. If a
\sphinxtitleref{kspace\_style \textless{}kspace\_style\textgreater{}} is declared, an error will be issued.


\subsection{Related commands}
\label{\detokenize{bond_special:related-commands}}
\sphinxAtStartPar
{\hyperref[\detokenize{bond_coeff::doc}]{\sphinxcrossref{\DUrole{doc}{bond\_coeff}}}}, {\hyperref[\detokenize{special_bonds::doc}]{\sphinxcrossref{\DUrole{doc}{special\_bonds}}}}


\subsection{Default}
\label{\detokenize{bond_special:default}}
\sphinxAtStartPar
none

\sphinxstepscope

\index{bond\_style table@\spxentry{bond\_style table}}\ignorespaces 
\index{bond\_style table/omp@\spxentry{bond\_style table/omp}}\ignorespaces 

\section{bond\_style table command}
\label{\detokenize{bond_table:bond-style-table-command}}\label{\detokenize{bond_table:index-1}}\label{\detokenize{bond_table:index-0}}\label{\detokenize{bond_table::doc}}
\sphinxAtStartPar
Accelerator Variants: \sphinxstyleemphasis{table/omp}


\subsection{Syntax}
\label{\detokenize{bond_table:syntax}}
\begin{sphinxVerbatim}[commandchars=\\\{\}]
\PYG{k}{bond\PYGZus{}style}\PYG{+w}{ }\PYG{n}{table}\PYG{+w}{ }\PYG{n}{style}\PYG{+w}{ }\PYG{n}{N}
\end{sphinxVerbatim}
\begin{itemize}
\item {} 
\sphinxAtStartPar
style = \sphinxstyleemphasis{linear} or \sphinxstyleemphasis{spline} = method of interpolation

\item {} 
\sphinxAtStartPar
N = use N values in table

\end{itemize}


\subsection{Examples}
\label{\detokenize{bond_table:examples}}
\begin{sphinxVerbatim}[commandchars=\\\{\}]
\PYG{k}{bond\PYGZus{}style}\PYG{+w}{ }\PYG{n}{table}\PYG{+w}{ }\PYG{n}{linear}\PYG{+w}{ }\PYG{l+m}{1000}
\PYG{k}{bond\PYGZus{}coeff}\PYG{+w}{ }\PYG{l+m}{1}\PYG{+w}{ }\PYG{n}{file.table}\PYG{+w}{ }\PYG{n}{ENTRY1}
\end{sphinxVerbatim}


\subsection{Description}
\label{\detokenize{bond_table:description}}
\sphinxAtStartPar
Style \sphinxstyleemphasis{table} creates interpolation tables of length \sphinxstyleemphasis{N} from bond
potential and force values listed in a file(s) as a function of bond
length.  The files are read by the {\hyperref[\detokenize{bond_coeff::doc}]{\sphinxcrossref{\DUrole{doc}{bond\_coeff}}}}
command.

\sphinxAtStartPar
The interpolation tables are created by fitting cubic splines to the
file values and interpolating energy and force values at each of \sphinxstyleemphasis{N}
distances.  During a simulation, these tables are used to interpolate
energy and force values as needed.  The interpolation is done in one
of 2 styles: \sphinxstyleemphasis{linear} or \sphinxstyleemphasis{spline}.

\sphinxAtStartPar
For the \sphinxstyleemphasis{linear} style, the bond length is used to find 2 surrounding
table values from which an energy or force is computed by linear
interpolation.

\sphinxAtStartPar
For the \sphinxstyleemphasis{spline} style, a cubic spline coefficients are computed and
stored at each of the \sphinxstyleemphasis{N} values in the table.  The bond length is
used to find the appropriate set of coefficients which are used to
evaluate a cubic polynomial which computes the energy or force.

\sphinxAtStartPar
The following coefficients must be defined for each bond type via the
{\hyperref[\detokenize{bond_coeff::doc}]{\sphinxcrossref{\DUrole{doc}{bond\_coeff}}}} command as in the example above.
\begin{itemize}
\item {} 
\sphinxAtStartPar
filename

\item {} 
\sphinxAtStartPar
keyword

\end{itemize}

\sphinxAtStartPar
The filename specifies a file containing tabulated energy and force
values.  The keyword specifies a section of the file.  The format of
this file is described below.


\bigskip\hrule\bigskip


\sphinxAtStartPar
Suitable tables for use with this bond style can be created by LAMMPS
itself from existing bond styles using the {\hyperref[\detokenize{bond_write::doc}]{\sphinxcrossref{\DUrole{doc}{bond\_write}}}} command.  This can be useful to have a template file for
testing the bond style settings and to build a compatible custom file.
Another option to generate tables is the Python code in the
\sphinxcode{\sphinxupquote{tools/tabulate}} folder of the LAMMPS source code distribution.

\sphinxAtStartPar
The format of a tabulated file is as follows (without the
parenthesized comments):

\begin{sphinxVerbatim}[commandchars=\\\{\}]
\PYGZsh{} Bond potential for harmonic (one or more comment or blank lines)

HAM                           (keyword is the first text on line)
N 101 FP 0 0 EQ 0.5           (N, FP, EQ  parameters)
                              (blank line)
1 0.00 338.0000 1352.0000     (index, bond\PYGZhy{}length, energy, force)
2 0.01 324.6152 1324.9600
...
101 1.00 338.0000 \PYGZhy{}1352.0000
\end{sphinxVerbatim}

\sphinxAtStartPar
A section begins with a non\sphinxhyphen{}blank line whose first character is not a
“\#”; blank lines or lines starting with “\#” can be used as comments
between sections.  The first line begins with a keyword which
identifies the section.  The line can contain additional text, but the
initial text must match the argument specified in the
{\hyperref[\detokenize{bond_coeff::doc}]{\sphinxcrossref{\DUrole{doc}{bond\_coeff}}}} command.  The next line lists (in any
order) one or more parameters for the table.  Each parameter is a
keyword followed by one or more numeric values.

\sphinxAtStartPar
The parameter “N” is required and its value is the number of table
entries that follow.  Note that this may be different than the \sphinxstyleemphasis{N}
specified in the {\hyperref[\detokenize{bond_style::doc}]{\sphinxcrossref{\DUrole{doc}{bond\_style table}}}} command.  Let
Ntable = \sphinxstyleemphasis{N} in the bond\_style command, and Nfile = “N” in the
tabulated file.  What LAMMPS does is a preliminary interpolation by
creating splines using the Nfile tabulated values as nodal points.  It
uses these to interpolate as needed to generate energy and force
values at Ntable different points.  The resulting tables of length
Ntable are then used as described above, when computing energy and
force for individual bond lengths.  This means that if you want the
interpolation tables of length Ntable to match exactly what is in the
tabulated file (with effectively no preliminary interpolation), you
should set Ntable = Nfile.

\sphinxAtStartPar
The “FP” parameter is optional.  If used, it is followed by two values
fplo and fphi, which are the derivatives of the force at the innermost
and outermost bond lengths.  These values are needed by the spline
construction routines.  If not specified by the “FP” parameter, they
are estimated (less accurately) by the first two and last two force
values in the table.

\sphinxAtStartPar
The “EQ” parameter is also optional.  If used, it is followed by a the
equilibrium bond length, which is used, for example, by the {\hyperref[\detokenize{fix_shake::doc}]{\sphinxcrossref{\DUrole{doc}{fix
shake}}}} command.  If not used, the equilibrium bond length is
to the distance in the table with the lowest potential energy.

\sphinxAtStartPar
Following a blank line, the next N lines list the tabulated values.
On each line, the first value is the index from 1 to N, the second value is
the bond length r (in distance units), the third value is the energy (in
energy units), and the fourth is the force (in force units).  The bond
lengths must range from a LO value to a HI value, and increase from
one line to the next.  If the actual bond length is ever smaller than
the LO value or larger than the HI value, then the calculation is
aborted with an error, so it is advisable to cover the whole range
of possible bond lengths.

\sphinxAtStartPar
Note that one file can contain many sections, each with a tabulated
potential.  LAMMPS reads the file section by section until it finds
one that matches the specified keyword.


\bigskip\hrule\bigskip


\sphinxAtStartPar
Styles with a \sphinxstyleemphasis{gpu}, \sphinxstyleemphasis{intel}, \sphinxstyleemphasis{kk}, \sphinxstyleemphasis{omp}, or \sphinxstyleemphasis{opt} suffix are
functionally the same as the corresponding style without the suffix.
They have been optimized to run faster, depending on your available
hardware, as discussed on the {\hyperref[\detokenize{Speed_packages::doc}]{\sphinxcrossref{\DUrole{doc}{Accelerator packages}}}}
page.  The accelerated styles take the same arguments and should
produce the same results, except for round\sphinxhyphen{}off and precision issues.

\sphinxAtStartPar
These accelerated styles are part of the GPU, INTEL, KOKKOS,
OPENMP, and OPT packages, respectively.  They are only enabled if
LAMMPS was built with those packages.  See the {\hyperref[\detokenize{Build_package::doc}]{\sphinxcrossref{\DUrole{doc}{Build package}}}} page for more info.

\sphinxAtStartPar
You can specify the accelerated styles explicitly in your input script
by including their suffix, or you can use the {\hyperref[\detokenize{Run_options::doc}]{\sphinxcrossref{\DUrole{doc}{\sphinxhyphen{}suffix command\sphinxhyphen{}line switch}}}} when you invoke LAMMPS, or you can use the
{\hyperref[\detokenize{suffix::doc}]{\sphinxcrossref{\DUrole{doc}{suffix}}}} command in your input script.

\sphinxAtStartPar
See the {\hyperref[\detokenize{Speed_packages::doc}]{\sphinxcrossref{\DUrole{doc}{Accelerator packages}}}} page for more
instructions on how to use the accelerated styles effectively.


\bigskip\hrule\bigskip



\subsection{Restart info}
\label{\detokenize{bond_table:restart-info}}
\sphinxAtStartPar
This bond style writes the settings for the “bond\_style table” command
to {\hyperref[\detokenize{restart::doc}]{\sphinxcrossref{\DUrole{doc}{binary restart files}}}}, so a bond\_style command does
not need to specified in an input script that reads a restart file.
However, the coefficient information is not stored in the restart file,
since it is tabulated in the potential files.  Thus, bond\_coeff commands
do need to be specified in the restart input script.


\subsection{Restrictions}
\label{\detokenize{bond_table:restrictions}}
\sphinxAtStartPar
This bond style can only be used if LAMMPS was built with the MOLECULE
package.  See the {\hyperref[\detokenize{Build_package::doc}]{\sphinxcrossref{\DUrole{doc}{Build package}}}} page for more
info.


\subsection{Related commands}
\label{\detokenize{bond_table:related-commands}}
\sphinxAtStartPar
{\hyperref[\detokenize{bond_coeff::doc}]{\sphinxcrossref{\DUrole{doc}{bond\_coeff}}}}, {\hyperref[\detokenize{delete_bonds::doc}]{\sphinxcrossref{\DUrole{doc}{delete\_bonds}}}},
{\hyperref[\detokenize{bond_write::doc}]{\sphinxcrossref{\DUrole{doc}{bond\_write}}}}


\subsection{Default}
\label{\detokenize{bond_table:default}}
\sphinxAtStartPar
none

\sphinxstepscope

\index{bond\_style zero@\spxentry{bond\_style zero}}\ignorespaces 

\section{bond\_style zero command}
\label{\detokenize{bond_zero:bond-style-zero-command}}\label{\detokenize{bond_zero:index-0}}\label{\detokenize{bond_zero::doc}}

\subsection{Syntax}
\label{\detokenize{bond_zero:syntax}}
\begin{sphinxVerbatim}[commandchars=\\\{\}]
\PYG{k}{bond\PYGZus{}style}\PYG{+w}{ }\PYG{n}{zero}\PYG{+w}{ }\PYG{n}{keyword}
\end{sphinxVerbatim}
\begin{itemize}
\item {} 
\sphinxAtStartPar
zero or more keywords may be appended

\item {} 
\sphinxAtStartPar
keyword = \sphinxstyleemphasis{nocoeff}

\end{itemize}


\subsection{Examples}
\label{\detokenize{bond_zero:examples}}
\begin{sphinxVerbatim}[commandchars=\\\{\}]
\PYG{k}{bond\PYGZus{}style}\PYG{+w}{ }\PYG{n}{zero}
\PYG{k}{bond\PYGZus{}style}\PYG{+w}{ }\PYG{n}{zero}\PYG{+w}{ }\PYG{n}{nocoeff}
\PYG{k}{bond\PYGZus{}coeff}\PYG{+w}{ }\PYG{o}{*}
\PYG{k}{bond\PYGZus{}coeff}\PYG{+w}{ }\PYG{o}{*}\PYG{+w}{ }\PYG{l+m}{2.14}
\end{sphinxVerbatim}


\subsection{Description}
\label{\detokenize{bond_zero:description}}
\sphinxAtStartPar
Using an bond style of zero means bond forces and energies are not
computed, but the geometry of bond pairs is still accessible to other
commands.

\sphinxAtStartPar
As an example, the {\hyperref[\detokenize{compute_bond_local::doc}]{\sphinxcrossref{\DUrole{doc}{compute bond/local}}}}
command can be used to compute distances for the list of pairs of bond
atoms listed in the data file read by the {\hyperref[\detokenize{read_data::doc}]{\sphinxcrossref{\DUrole{doc}{read\_data}}}}
command.  If no bond style is defined, this command cannot be used.

\sphinxAtStartPar
The optional \sphinxstyleemphasis{nocoeff} flag allows to read data files with a BondCoeff
section for any bond style. Similarly, any bond\_coeff commands
will only be checked for the bond type number and the rest ignored.

\sphinxAtStartPar
Note that the {\hyperref[\detokenize{bond_coeff::doc}]{\sphinxcrossref{\DUrole{doc}{bond\_coeff}}}} command must be used for
all bond types. If specified, there can be only one value, which is
going to be used to assign an equilibrium distance, e.g. for use with
{\hyperref[\detokenize{fix_shake::doc}]{\sphinxcrossref{\DUrole{doc}{fix shake}}}}.


\subsection{Restrictions}
\label{\detokenize{bond_zero:restrictions}}\begin{quote}

\sphinxAtStartPar
none
\end{quote}


\subsection{Related commands}
\label{\detokenize{bond_zero:related-commands}}
\sphinxAtStartPar
{\hyperref[\detokenize{bond_none::doc}]{\sphinxcrossref{\DUrole{doc}{bond\_style none}}}}


\subsection{Default}
\label{\detokenize{bond_zero:default}}
\sphinxAtStartPar
none

\sphinxstepscope


\chapter{Angle Styles}
\label{\detokenize{angles:angle-styles}}\label{\detokenize{angles::doc}}
\sphinxstepscope

\index{angle\_style amoeba@\spxentry{angle\_style amoeba}}\ignorespaces 

\section{angle\_style amoeba command}
\label{\detokenize{angle_amoeba:angle-style-amoeba-command}}\label{\detokenize{angle_amoeba:index-0}}\label{\detokenize{angle_amoeba::doc}}

\subsection{Syntax}
\label{\detokenize{angle_amoeba:syntax}}
\begin{sphinxVerbatim}[commandchars=\\\{\}]
\PYG{k}{angle\PYGZus{}style}\PYG{+w}{ }\PYG{n}{amoeba}
\end{sphinxVerbatim}


\subsection{Examples}
\label{\detokenize{angle_amoeba:examples}}
\begin{sphinxVerbatim}[commandchars=\\\{\}]
\PYG{k}{angle\PYGZus{}style}\PYG{+w}{ }\PYG{n}{amoeba}
\PYG{k}{angle\PYGZus{}coeff}\PYG{+w}{ }\PYG{o}{*}\PYG{+w}{ }\PYG{l+m}{75.0}\PYG{+w}{ }\PYG{o}{\PYGZhy{}}\PYG{l+m}{25.0}\PYG{+w}{ }\PYG{l+m}{1.0}\PYG{+w}{ }\PYG{l+m}{0.3}\PYG{+w}{ }\PYG{l+m}{0.02}\PYG{+w}{ }\PYG{l+m}{0.003}
\PYG{k}{angle\PYGZus{}coeff}\PYG{+w}{ }\PYG{o}{*}\PYG{+w}{ }\PYG{n}{ba}\PYG{+w}{ }\PYG{l+m}{3.6551}\PYG{+w}{ }\PYG{l+m}{24.895}\PYG{+w}{ }\PYG{l+m}{1.0119}\PYG{+w}{ }\PYG{l+m}{1.5228}
\PYG{k}{angle\PYGZus{}coeff}\PYG{+w}{ }\PYG{o}{*}\PYG{+w}{ }\PYG{n}{ub}\PYG{+w}{ }\PYG{o}{\PYGZhy{}}\PYG{l+m}{7.6}\PYG{+w}{ }\PYG{l+m}{1.5537}
\end{sphinxVerbatim}


\subsection{Description}
\label{\detokenize{angle_amoeba:description}}
\sphinxAtStartPar
The \sphinxstyleemphasis{amoeba} angle style uses the potential
\begin{equation*}
\begin{split}E & = E_a + E_{ba} + E_{ub} \\
E_a & = K_2\left(\theta - \theta_0\right)^2 + K_3\left(\theta - \theta_0\right)^3 + K_4\left(\theta - \theta_0\right)^4 + K_5\left(\theta - \theta_0\right)^5 + K_6\left(\theta - \theta_0\right)^6 \\
E_{ba} & = N_1 (r_{ij} - r_1) (\theta - \theta_0) + N_2(r_{jk} - r_2)(\theta - \theta_0) \\
E_{UB} & = K_{ub} (r_{ik} - r_{ub})^2\end{split}
\end{equation*}
\sphinxAtStartPar
where \(E_a\) is the angle term, \(E_{ba}\) is a bond\sphinxhyphen{}angle
term, \(E_{UB}\) is a Urey\sphinxhyphen{}Bradley bond term, \(\theta_0\) is
the equilibrium angle, \(r_1\) and \(r_2\) are the equilibrium
bond lengths, and \(r_{ub}\) is the equilibrium Urey\sphinxhyphen{}Bradley bond
length.

\sphinxAtStartPar
These formulas match how the Tinker MD code performs its angle
calculations for the AMOEBA and HIPPO force fields.  See the
{\hyperref[\detokenize{Howto_amoeba::doc}]{\sphinxcrossref{\DUrole{doc}{Howto amoeba}}}} page for more information about
the implementation of AMOEBA and HIPPO in LAMMPS.

\sphinxAtStartPar
Note that the \(E_a\) and \(E_{ba}\) formulas are identical to
those used for the {\hyperref[\detokenize{angle_class2::doc}]{\sphinxcrossref{\DUrole{doc}{angle\_style class2/p6}}}}
command, however there is no bond\sphinxhyphen{}bond cross term formula for
\(E_{bb}\).  Additionally, there is a \(E_{UB}\) term for a
Urey\sphinxhyphen{}Bradley bond.  It is effectively a harmonic bond between the I
and K atoms of angle IJK, even though that bond is not enumerated in
the “Bonds” section of the data file.

\sphinxAtStartPar
There are also two ways that Tinker computes the angle \(\theta\)
in the \(E_a\) formula.  The first is the standard way of treating
IJK as an “in\sphinxhyphen{}plane” angle.  The second is an “out\sphinxhyphen{}of\sphinxhyphen{}plane” method
which Tinker may use if the center atom J in the angle is bonded to
one additional atom in addition to I and K.  In this case, all 4 atoms
are used to compute the \(E_a\) formula, resulting in forces on
all 4 atoms.  In the Tinker PRM file, these 2 options are denoted by
\sphinxstyleemphasis{angle} versus \sphinxstyleemphasis{anglep} entries in the “Angle Bending Parameters”
section of the PRM force field file.  The \sphinxstyleemphasis{pflag} coefficient
described below selects between the 2 options.


\bigskip\hrule\bigskip


\sphinxAtStartPar
Coefficients for the \(E_a\), \(E_{bb}\), and \(E_{ub}\)
formulas must be defined for each angle type via the {\hyperref[\detokenize{angle_coeff::doc}]{\sphinxcrossref{\DUrole{doc}{angle\_coeff}}}} command as in the example above, or in the data file or
restart files read by the {\hyperref[\detokenize{read_data::doc}]{\sphinxcrossref{\DUrole{doc}{read\_data}}}} or
{\hyperref[\detokenize{read_restart::doc}]{\sphinxcrossref{\DUrole{doc}{read\_restart}}}} commands.

\sphinxAtStartPar
These are the 8 coefficients for the \(E_a\) formula:
\begin{itemize}
\item {} 
\sphinxAtStartPar
pflag = 0 or 1

\item {} 
\sphinxAtStartPar
ubflag = 0 or 1

\item {} 
\sphinxAtStartPar
\(\theta_0\) (degrees)

\item {} 
\sphinxAtStartPar
\(K_2\) (energy)

\item {} 
\sphinxAtStartPar
\(K_3\) (energy)

\item {} 
\sphinxAtStartPar
\(K_4\) (energy)

\item {} 
\sphinxAtStartPar
\(K_5\) (energy)

\item {} 
\sphinxAtStartPar
\(K_6\) (energy)

\end{itemize}

\sphinxAtStartPar
A pflag value of 0 vs 1 selects between the “in\sphinxhyphen{}plane” and
“out\sphinxhyphen{}of\sphinxhyphen{}plane” options described above.  Ubflag is 1 if there is a
Urey\sphinxhyphen{}Bradley term associated with this angle type, else it is 0.
\(\theta_0\) is specified in degrees, but LAMMPS converts it to
radians internally; hence the various \(K\) values are effectively
energy per radian\textasciicircum{}2 or radian\textasciicircum{}3 or radian\textasciicircum{}4 or radian\textasciicircum{}5 or
radian\textasciicircum{}6.

\sphinxAtStartPar
For the \(E_{ba}\) formula, each line in a {\hyperref[\detokenize{angle_coeff::doc}]{\sphinxcrossref{\DUrole{doc}{angle\_coeff}}}} command in the input script lists 5 coefficients, the
first of which is “ba” to indicate they are BondAngle coefficients.
In a data file, these coefficients should be listed under a “BondAngle
Coeffs” heading and you must leave out the “ba”, i.e. only list 4
coefficients after the angle type.
\begin{itemize}
\item {} 
\sphinxAtStartPar
ba

\item {} 
\sphinxAtStartPar
\(N_1\) (energy/distance\textasciicircum{}2)

\item {} 
\sphinxAtStartPar
\(N_2\) (energy/distance\textasciicircum{}2)

\item {} 
\sphinxAtStartPar
\(r_1\) (distance)

\item {} 
\sphinxAtStartPar
\(r_2\) (distance)

\end{itemize}

\sphinxAtStartPar
The \(\theta_0\) value in the \(E_{ba}\) formula is not specified,
since it is the same value from the \(E_a\) formula.

\sphinxAtStartPar
For the \(E_{ub}\) formula, each line in a {\hyperref[\detokenize{angle_coeff::doc}]{\sphinxcrossref{\DUrole{doc}{angle\_coeff}}}} command in the input script lists 3 coefficients, the
first of which is “ub” to indicate they are UreyBradley coefficients.
In a data file, these coefficients should be listed under a
“UreyBradley Coeffs” heading and you must leave out the “ub”,
i.e. only list 2 coefficients after the angle type.
\begin{itemize}
\item {} 
\sphinxAtStartPar
ub

\item {} 
\sphinxAtStartPar
\(K_{ub}\) (energy/distance\textasciicircum{}2)

\item {} 
\sphinxAtStartPar
\(r_{ub}\) (distance)

\end{itemize}


\bigskip\hrule\bigskip



\subsection{Restrictions}
\label{\detokenize{angle_amoeba:restrictions}}
\sphinxAtStartPar
This angle style can only be used if LAMMPS was built with the AMOEBA
package.  See the {\hyperref[\detokenize{Build_package::doc}]{\sphinxcrossref{\DUrole{doc}{Build package}}}} doc page for
more info.


\subsection{Related commands}
\label{\detokenize{angle_amoeba:related-commands}}
\sphinxAtStartPar
{\hyperref[\detokenize{angle_coeff::doc}]{\sphinxcrossref{\DUrole{doc}{angle\_coeff}}}}


\subsection{Default}
\label{\detokenize{angle_amoeba:default}}
\sphinxAtStartPar
none

\sphinxstepscope

\index{angle\_style charmm@\spxentry{angle\_style charmm}}\ignorespaces 
\index{angle\_style charmm/intel@\spxentry{angle\_style charmm/intel}}\ignorespaces 
\index{angle\_style charmm/kk@\spxentry{angle\_style charmm/kk}}\ignorespaces 
\index{angle\_style charmm/omp@\spxentry{angle\_style charmm/omp}}\ignorespaces 

\section{angle\_style charmm command}
\label{\detokenize{angle_charmm:angle-style-charmm-command}}\label{\detokenize{angle_charmm:index-3}}\label{\detokenize{angle_charmm:index-2}}\label{\detokenize{angle_charmm:index-1}}\label{\detokenize{angle_charmm:index-0}}\label{\detokenize{angle_charmm::doc}}
\sphinxAtStartPar
Accelerator Variants: \sphinxstyleemphasis{charmm/intel}, \sphinxstyleemphasis{charmm/kk}, \sphinxstyleemphasis{charmm/omp}


\subsection{Syntax}
\label{\detokenize{angle_charmm:syntax}}
\begin{sphinxVerbatim}[commandchars=\\\{\}]
\PYG{k}{angle\PYGZus{}style}\PYG{+w}{ }\PYG{n}{charmm}
\end{sphinxVerbatim}


\subsection{Examples}
\label{\detokenize{angle_charmm:examples}}
\begin{sphinxVerbatim}[commandchars=\\\{\}]
\PYG{k}{angle\PYGZus{}style}\PYG{+w}{ }\PYG{n}{charmm}
\PYG{k}{angle\PYGZus{}coeff}\PYG{+w}{ }\PYG{l+m}{1}\PYG{+w}{ }\PYG{l+m}{300.0}\PYG{+w}{ }\PYG{l+m}{107.0}\PYG{+w}{ }\PYG{l+m}{50.0}\PYG{+w}{ }\PYG{l+m}{3.0}
\end{sphinxVerbatim}


\subsection{Description}
\label{\detokenize{angle_charmm:description}}
\sphinxAtStartPar
The \sphinxstyleemphasis{charmm} angle style uses the potential
\begin{equation*}
\begin{split}E = K (\theta - \theta_0)^2 + K_{ub} (r - r_{ub})^2\end{split}
\end{equation*}
\sphinxAtStartPar
with an additional Urey\_Bradley term based on the distance \(r\) between
the first and third atoms in the angle.  \(K\), \(\theta_0\),
\(K_{ub}\), and \(R_{ub}\) are coefficients defined for each angle
type.

\sphinxAtStartPar
See {\hyperref[\detokenize{angle_charmm:angle-mackerell}]{\sphinxcrossref{\DUrole{std,std-ref}{(MacKerell)}}}} for a description of the CHARMM force
field.

\sphinxAtStartPar
The following coefficients must be defined for each angle type via the
{\hyperref[\detokenize{angle_coeff::doc}]{\sphinxcrossref{\DUrole{doc}{angle\_coeff}}}} command as in the example above, or in
the data file or restart files read by the {\hyperref[\detokenize{read_data::doc}]{\sphinxcrossref{\DUrole{doc}{read\_data}}}}
or {\hyperref[\detokenize{read_restart::doc}]{\sphinxcrossref{\DUrole{doc}{read\_restart}}}} commands:
\begin{itemize}
\item {} 
\sphinxAtStartPar
\(K\) (energy)

\item {} 
\sphinxAtStartPar
\(\theta_0\) (degrees)

\item {} 
\sphinxAtStartPar
\(K_{ub}\) (energy/distance\textasciicircum{}2)

\item {} 
\sphinxAtStartPar
\(r_{ub}\) (distance)

\end{itemize}

\sphinxAtStartPar
\(\theta_0\) is specified in degrees, but LAMMPS converts it to
radians internally; hence \(K\) is effectively energy per
radian\textasciicircum{}2.


\bigskip\hrule\bigskip


\sphinxAtStartPar
Styles with a \sphinxstyleemphasis{gpu}, \sphinxstyleemphasis{intel}, \sphinxstyleemphasis{kk}, \sphinxstyleemphasis{omp}, or \sphinxstyleemphasis{opt} suffix are
functionally the same as the corresponding style without the suffix.
They have been optimized to run faster, depending on your available
hardware, as discussed on the {\hyperref[\detokenize{Speed_packages::doc}]{\sphinxcrossref{\DUrole{doc}{Accelerator packages}}}}
page.  The accelerated styles take the same arguments and should
produce the same results, except for round\sphinxhyphen{}off and precision issues.

\sphinxAtStartPar
These accelerated styles are part of the GPU, INTEL, KOKKOS,
OPENMP, and OPT packages, respectively.  They are only enabled if
LAMMPS was built with those packages.  See the {\hyperref[\detokenize{Build_package::doc}]{\sphinxcrossref{\DUrole{doc}{Build package}}}} page for more info.

\sphinxAtStartPar
You can specify the accelerated styles explicitly in your input script
by including their suffix, or you can use the {\hyperref[\detokenize{Run_options::doc}]{\sphinxcrossref{\DUrole{doc}{\sphinxhyphen{}suffix command\sphinxhyphen{}line switch}}}} when you invoke LAMMPS, or you can use the
{\hyperref[\detokenize{suffix::doc}]{\sphinxcrossref{\DUrole{doc}{suffix}}}} command in your input script.

\sphinxAtStartPar
See the {\hyperref[\detokenize{Speed_packages::doc}]{\sphinxcrossref{\DUrole{doc}{Accelerator packages}}}} page for more
instructions on how to use the accelerated styles effectively.


\bigskip\hrule\bigskip



\subsection{Restrictions}
\label{\detokenize{angle_charmm:restrictions}}
\sphinxAtStartPar
This angle style can only be used if LAMMPS was built with the
MOLECULE package.  See the {\hyperref[\detokenize{Build_package::doc}]{\sphinxcrossref{\DUrole{doc}{Build package}}}} doc page
for more info.


\subsection{Related commands}
\label{\detokenize{angle_charmm:related-commands}}
\sphinxAtStartPar
{\hyperref[\detokenize{angle_coeff::doc}]{\sphinxcrossref{\DUrole{doc}{angle\_coeff}}}}, {\hyperref[\detokenize{pair_charmm::doc}]{\sphinxcrossref{\DUrole{doc}{pair\_style lj/charmm variants}}}},
{\hyperref[\detokenize{dihedral_charmm::doc}]{\sphinxcrossref{\DUrole{doc}{dihedral\_style charmm}}}},
{\hyperref[\detokenize{dihedral_charmm::doc}]{\sphinxcrossref{\DUrole{doc}{dihedral\_style charmmfsw}}}}, {\hyperref[\detokenize{fix_cmap::doc}]{\sphinxcrossref{\DUrole{doc}{fix cmap}}}}


\subsection{Default}
\label{\detokenize{angle_charmm:default}}
\sphinxAtStartPar
none


\bigskip\hrule\bigskip

\phantomsection\label{\detokenize{angle_charmm:angle-mackerell}}
\sphinxAtStartPar
\sphinxstylestrong{(MacKerell)} MacKerell, Bashford, Bellott, Dunbrack, Evanseck, Field,
Fischer, Gao, Guo, Ha, et al, J Phys Chem, 102, 3586 (1998).

\sphinxstepscope

\index{angle\_style class2@\spxentry{angle\_style class2}}\ignorespaces 
\index{angle\_style class2/kk@\spxentry{angle\_style class2/kk}}\ignorespaces 
\index{angle\_style class2/omp@\spxentry{angle\_style class2/omp}}\ignorespaces 
\index{angle\_style class2/p6@\spxentry{angle\_style class2/p6}}\ignorespaces 

\section{angle\_style class2 command}
\label{\detokenize{angle_class2:angle-style-class2-command}}\label{\detokenize{angle_class2:index-3}}\label{\detokenize{angle_class2:index-2}}\label{\detokenize{angle_class2:index-1}}\label{\detokenize{angle_class2:index-0}}\label{\detokenize{angle_class2::doc}}
\sphinxAtStartPar
Accelerator Variants: \sphinxstyleemphasis{class2/kk}, \sphinxstyleemphasis{class2/omp}


\section{angle\_style class2/p6 command}
\label{\detokenize{angle_class2:angle-style-class2-p6-command}}

\subsection{Syntax}
\label{\detokenize{angle_class2:syntax}}
\begin{sphinxVerbatim}[commandchars=\\\{\}]
\PYG{k}{angle\PYGZus{}style}\PYG{+w}{ }\PYG{n}{class2}
\end{sphinxVerbatim}


\subsection{Examples}
\label{\detokenize{angle_class2:examples}}
\begin{sphinxVerbatim}[commandchars=\\\{\}]
\PYG{k}{angle\PYGZus{}style}\PYG{+w}{ }\PYG{n}{class2}
\PYG{k}{angle\PYGZus{}coeff}\PYG{+w}{ }\PYG{o}{*}\PYG{+w}{ }\PYG{l+m}{75.0}\PYG{+w}{ }\PYG{l+m}{25.0}\PYG{+w}{ }\PYG{l+m}{0.3}\PYG{+w}{ }\PYG{l+m}{0.002}
\PYG{k}{angle\PYGZus{}coeff}\PYG{+w}{ }\PYG{l+m}{1}\PYG{+w}{ }\PYG{n}{bb}\PYG{+w}{ }\PYG{l+m}{10.5872}\PYG{+w}{ }\PYG{l+m}{1.0119}\PYG{+w}{ }\PYG{l+m}{1.5228}
\PYG{k}{angle\PYGZus{}coeff}\PYG{+w}{ }\PYG{o}{*}\PYG{+w}{ }\PYG{n}{ba}\PYG{+w}{ }\PYG{l+m}{3.6551}\PYG{+w}{ }\PYG{l+m}{24.895}\PYG{+w}{ }\PYG{l+m}{1.0119}\PYG{+w}{ }\PYG{l+m}{1.5228}
\end{sphinxVerbatim}


\subsection{Description}
\label{\detokenize{angle_class2:description}}
\sphinxAtStartPar
The \sphinxstyleemphasis{class2} angle style uses the potential
\begin{equation*}
\begin{split}E & = E_a + E_{bb} + E_{ba} \\
E_a & = K_2 (\theta - \theta_0)^2 + K_3 (\theta - \theta_0)^3 + K_4(\theta - \theta_0)^4 \\
E_{bb} & = M (r_{ij} - r_1) (r_{jk} - r_2) \\
E_{ba} & = N_1 (r_{ij} - r_1) (\theta - \theta_0) + N_2(r_{jk} - r_2)(\theta - \theta_0)\end{split}
\end{equation*}
\sphinxAtStartPar
where \(E_a\) is the angle term, \(E_{bb}\) is a bond\sphinxhyphen{}bond
term, and \(E_{ba}\) is a bond\sphinxhyphen{}angle term.  \(\theta_0\) is
the equilibrium angle and \(r_1\) and \(r_2\) are the
equilibrium bond lengths.

\sphinxAtStartPar
See {\hyperref[\detokenize{angle_class2:angle-sun}]{\sphinxcrossref{\DUrole{std,std-ref}{(Sun)}}}} for a description of the COMPASS class2 force field.

\sphinxAtStartPar
Coefficients for the \(E_a\), \(E_{bb}\), and \(E_{ba}\)
formulas must be defined for each angle type via the {\hyperref[\detokenize{angle_coeff::doc}]{\sphinxcrossref{\DUrole{doc}{angle\_coeff}}}} command as in the example above, or in the data file or
restart files read by the {\hyperref[\detokenize{read_data::doc}]{\sphinxcrossref{\DUrole{doc}{read\_data}}}} or
{\hyperref[\detokenize{read_restart::doc}]{\sphinxcrossref{\DUrole{doc}{read\_restart}}}} commands.

\sphinxAtStartPar
These are the 4 coefficients for the \(E_a\) formula:
\begin{itemize}
\item {} 
\sphinxAtStartPar
\(\theta_0\) (degrees)

\item {} 
\sphinxAtStartPar
\(K_2\) (energy)

\item {} 
\sphinxAtStartPar
\(K_3\) (energy)

\item {} 
\sphinxAtStartPar
\(K_4\) (energy)

\end{itemize}

\sphinxAtStartPar
\(\theta_0\) is specified in degrees, but LAMMPS converts it to
radians internally; hence the various \(K\) are effectively energy
per radian\textasciicircum{}2 or radian\textasciicircum{}3 or radian\textasciicircum{}4.

\sphinxAtStartPar
For the \(E_{bb}\) formula, each line in a {\hyperref[\detokenize{angle_coeff::doc}]{\sphinxcrossref{\DUrole{doc}{angle\_coeff}}}} command in the input script lists 4 coefficients, the
first of which is “bb” to indicate they are BondBond coefficients.  In
a data file, these coefficients should be listed under a “BondBond
Coeffs” heading and you must leave out the “bb”, i.e. only list 3
coefficients after the angle type.
\begin{itemize}
\item {} 
\sphinxAtStartPar
bb

\item {} 
\sphinxAtStartPar
\(M\) (energy/distance\textasciicircum{}2)

\item {} 
\sphinxAtStartPar
\(r_1\) (distance)

\item {} 
\sphinxAtStartPar
\(r_2\) (distance)

\end{itemize}

\sphinxAtStartPar
For the \(E_{ba}\) formula, each line in a {\hyperref[\detokenize{angle_coeff::doc}]{\sphinxcrossref{\DUrole{doc}{angle\_coeff}}}} command in the input script lists 5 coefficients, the
first of which is “ba” to indicate they are BondAngle coefficients.
In a data file, these coefficients should be listed under a “BondAngle
Coeffs” heading and you must leave out the “ba”, i.e. only list 4
coefficients after the angle type.
\begin{itemize}
\item {} 
\sphinxAtStartPar
ba

\item {} 
\sphinxAtStartPar
\(N_1\) (energy/distance)

\item {} 
\sphinxAtStartPar
\(N_2\) (energy/distance)

\item {} 
\sphinxAtStartPar
\(r_1\) (distance)

\item {} 
\sphinxAtStartPar
\(r_2\) (distance)

\end{itemize}

\sphinxAtStartPar
The \(\theta_0\) value in the \(E_{ba}\) formula is not specified,
since it is the same value from the \(E_a\) formula.

\begin{sphinxadmonition}{note}{Note:}
\sphinxAtStartPar
It is important that the order of the I,J,K atoms in each angle
listed in the Angles section of the data file read by the
{\hyperref[\detokenize{read_data::doc}]{\sphinxcrossref{\DUrole{doc}{read\_data}}}} command be consistent with the order
of the \(r_1\) and \(r_2\) BondBond and BondAngle
coefficients.  This is because the terms in the formulas for
\(E_{bb}\) and \(E_{ba}\) will use the I,J atoms to compute
\(r_{ij}\) and the J,K atoms to compute \(r_{jk}\).
\end{sphinxadmonition}


\bigskip\hrule\bigskip


\sphinxAtStartPar
Styles with a \sphinxstyleemphasis{gpu}, \sphinxstyleemphasis{intel}, \sphinxstyleemphasis{kk}, \sphinxstyleemphasis{omp}, or \sphinxstyleemphasis{opt} suffix are
functionally the same as the corresponding style without the suffix.
They have been optimized to run faster, depending on your available
hardware, as discussed on the {\hyperref[\detokenize{Speed_packages::doc}]{\sphinxcrossref{\DUrole{doc}{Accelerator packages}}}}
page.  The accelerated styles take the same arguments and should
produce the same results, except for round\sphinxhyphen{}off and precision issues.

\sphinxAtStartPar
These accelerated styles are part of the GPU, INTEL, KOKKOS,
OPENMP, and OPT packages, respectively.  They are only enabled if
LAMMPS was built with those packages.  See the {\hyperref[\detokenize{Build_package::doc}]{\sphinxcrossref{\DUrole{doc}{Build package}}}} page for more info.

\sphinxAtStartPar
You can specify the accelerated styles explicitly in your input script
by including their suffix, or you can use the {\hyperref[\detokenize{Run_options::doc}]{\sphinxcrossref{\DUrole{doc}{\sphinxhyphen{}suffix command\sphinxhyphen{}line switch}}}} when you invoke LAMMPS, or you can use the
{\hyperref[\detokenize{suffix::doc}]{\sphinxcrossref{\DUrole{doc}{suffix}}}} command in your input script.

\sphinxAtStartPar
See the {\hyperref[\detokenize{Speed_packages::doc}]{\sphinxcrossref{\DUrole{doc}{Accelerator packages}}}} page for more
instructions on how to use the accelerated styles effectively.


\bigskip\hrule\bigskip


\sphinxAtStartPar
The \sphinxstyleemphasis{class2/p6} angle style uses the \sphinxstyleemphasis{class2} potential expanded to sixth order:
\begin{equation*}
\begin{split}E_{a} = K_2\left(\theta - \theta_0\right)^2 + K_3\left(\theta - \theta_0\right)^3 + K_4\left(\theta - \theta_0\right)^4 + K_5\left(\theta - \theta_0\right)^5 + K_6\left(\theta - \theta_0\right)^6\end{split}
\end{equation*}
\sphinxAtStartPar
In this expanded term 6 coefficients for the \(E_a\) formula need to be set:
\begin{itemize}
\item {} 
\sphinxAtStartPar
\(\theta_0\) (degrees)

\item {} 
\sphinxAtStartPar
\(K_2\) (energy)

\item {} 
\sphinxAtStartPar
\(K_3\) (energy)

\item {} 
\sphinxAtStartPar
\(K_4\) (energy)

\item {} 
\sphinxAtStartPar
\(K_5\) (energy)

\item {} 
\sphinxAtStartPar
\(K_6\) (energy)

\end{itemize}

\sphinxAtStartPar
\(\theta_0\) is specified in degrees, but LAMMPS converts it to
radians internally; hence the various \(K\) are effectively energy
per radian\textasciicircum{}2 or radian\textasciicircum{}3 or radian\textasciicircum{}4 or radian\textasciicircum{}5 or radian\textasciicircum{}6.

\sphinxAtStartPar
The bond\sphinxhyphen{}bond and bond\sphinxhyphen{}angle terms remain unchanged.


\bigskip\hrule\bigskip



\subsection{Restrictions}
\label{\detokenize{angle_class2:restrictions}}
\sphinxAtStartPar
This angle style can only be used if LAMMPS was built with the CLASS2
package.  For the \sphinxstyleemphasis{class2/p6} style LAMMPS needs to be built with the
MOFFF package.  See the {\hyperref[\detokenize{Build_package::doc}]{\sphinxcrossref{\DUrole{doc}{Build package}}}} doc
page for more info.


\subsection{Related commands}
\label{\detokenize{angle_class2:related-commands}}
\sphinxAtStartPar
{\hyperref[\detokenize{angle_coeff::doc}]{\sphinxcrossref{\DUrole{doc}{angle\_coeff}}}}


\subsection{Default}
\label{\detokenize{angle_class2:default}}
\sphinxAtStartPar
none


\bigskip\hrule\bigskip

\phantomsection\label{\detokenize{angle_class2:angle-sun}}
\sphinxAtStartPar
\sphinxstylestrong{(Sun)} Sun, J Phys Chem B 102, 7338\sphinxhyphen{}7364 (1998).

\sphinxstepscope

\index{angle\_style cosine@\spxentry{angle\_style cosine}}\ignorespaces 
\index{angle\_style cosine/omp@\spxentry{angle\_style cosine/omp}}\ignorespaces 
\index{angle\_style cosine/kk@\spxentry{angle\_style cosine/kk}}\ignorespaces 

\section{angle\_style cosine command}
\label{\detokenize{angle_cosine:angle-style-cosine-command}}\label{\detokenize{angle_cosine:index-2}}\label{\detokenize{angle_cosine:index-1}}\label{\detokenize{angle_cosine:index-0}}\label{\detokenize{angle_cosine::doc}}
\sphinxAtStartPar
Accelerator Variants: \sphinxstyleemphasis{cosine/omp}, \sphinxstyleemphasis{cosine/kk}


\subsection{Syntax}
\label{\detokenize{angle_cosine:syntax}}
\begin{sphinxVerbatim}[commandchars=\\\{\}]
\PYG{k}{angle\PYGZus{}style}\PYG{+w}{ }\PYG{n}{cosine}
\end{sphinxVerbatim}


\subsection{Examples}
\label{\detokenize{angle_cosine:examples}}
\begin{sphinxVerbatim}[commandchars=\\\{\}]
\PYG{k}{angle\PYGZus{}style}\PYG{+w}{ }\PYG{n}{cosine}
\PYG{k}{angle\PYGZus{}coeff}\PYG{+w}{ }\PYG{o}{*}\PYG{+w}{ }\PYG{l+m}{75.0}
\end{sphinxVerbatim}


\subsection{Description}
\label{\detokenize{angle_cosine:description}}
\sphinxAtStartPar
The \sphinxstyleemphasis{cosine} angle style uses the potential
\begin{equation*}
\begin{split}E = K [1 + \cos(\theta)]\end{split}
\end{equation*}
\sphinxAtStartPar
where \(K\) is defined for each angle type.

\sphinxAtStartPar
The following coefficients must be defined for each angle type via the
{\hyperref[\detokenize{angle_coeff::doc}]{\sphinxcrossref{\DUrole{doc}{angle\_coeff}}}} command as in the example above, or in
the data file or restart files read by the {\hyperref[\detokenize{read_data::doc}]{\sphinxcrossref{\DUrole{doc}{read\_data}}}}
or {\hyperref[\detokenize{read_restart::doc}]{\sphinxcrossref{\DUrole{doc}{read\_restart}}}} commands:
\begin{itemize}
\item {} 
\sphinxAtStartPar
\(K\) (energy)

\end{itemize}


\bigskip\hrule\bigskip


\sphinxAtStartPar
Styles with a \sphinxstyleemphasis{gpu}, \sphinxstyleemphasis{intel}, \sphinxstyleemphasis{kk}, \sphinxstyleemphasis{omp}, or \sphinxstyleemphasis{opt} suffix are
functionally the same as the corresponding style without the suffix.
They have been optimized to run faster, depending on your available
hardware, as discussed on the {\hyperref[\detokenize{Speed_packages::doc}]{\sphinxcrossref{\DUrole{doc}{Accelerator packages}}}}
page.  The accelerated styles take the same arguments and should
produce the same results, except for round\sphinxhyphen{}off and precision issues.

\sphinxAtStartPar
These accelerated styles are part of the GPU, INTEL, KOKKOS,
OPENMP, and OPT packages, respectively.  They are only enabled if
LAMMPS was built with those packages.  See the {\hyperref[\detokenize{Build_package::doc}]{\sphinxcrossref{\DUrole{doc}{Build package}}}} page for more info.

\sphinxAtStartPar
You can specify the accelerated styles explicitly in your input script
by including their suffix, or you can use the {\hyperref[\detokenize{Run_options::doc}]{\sphinxcrossref{\DUrole{doc}{\sphinxhyphen{}suffix command\sphinxhyphen{}line switch}}}} when you invoke LAMMPS, or you can use the
{\hyperref[\detokenize{suffix::doc}]{\sphinxcrossref{\DUrole{doc}{suffix}}}} command in your input script.

\sphinxAtStartPar
See the {\hyperref[\detokenize{Speed_packages::doc}]{\sphinxcrossref{\DUrole{doc}{Accelerator packages}}}} page for more
instructions on how to use the accelerated styles effectively.


\bigskip\hrule\bigskip



\subsection{Restrictions}
\label{\detokenize{angle_cosine:restrictions}}
\sphinxAtStartPar
This angle style can only be used if LAMMPS was built with the
MOLECULE package.  See the {\hyperref[\detokenize{Build_package::doc}]{\sphinxcrossref{\DUrole{doc}{Build package}}}} doc page
for more info.


\subsection{Related commands}
\label{\detokenize{angle_cosine:related-commands}}
\sphinxAtStartPar
{\hyperref[\detokenize{angle_coeff::doc}]{\sphinxcrossref{\DUrole{doc}{angle\_coeff}}}}


\subsection{Default}
\label{\detokenize{angle_cosine:default}}
\sphinxAtStartPar
none

\sphinxstepscope

\index{angle\_style cosine/buck6d@\spxentry{angle\_style cosine/buck6d}}\ignorespaces 

\section{angle\_style cosine/buck6d command}
\label{\detokenize{angle_cosine_buck6d:angle-style-cosine-buck6d-command}}\label{\detokenize{angle_cosine_buck6d:index-0}}\label{\detokenize{angle_cosine_buck6d::doc}}

\subsection{Syntax}
\label{\detokenize{angle_cosine_buck6d:syntax}}
\begin{sphinxVerbatim}[commandchars=\\\{\}]
\PYG{k}{angle\PYGZus{}style}\PYG{+w}{ }\PYG{n}{cosine}\PYG{o}{/}\PYG{n}{buck6d}
\end{sphinxVerbatim}


\subsection{Examples}
\label{\detokenize{angle_cosine_buck6d:examples}}
\begin{sphinxVerbatim}[commandchars=\\\{\}]
\PYG{k}{angle\PYGZus{}style}\PYG{+w}{ }\PYG{n}{cosine}\PYG{o}{/}\PYG{n}{buck6d}
\PYG{k}{angle\PYGZus{}coeff}\PYG{+w}{ }\PYG{l+m}{1}\PYG{+w}{  }\PYG{n}{cosine}\PYG{o}{/}\PYG{n}{buck6d}\PYG{+w}{  }\PYG{l+m}{1.978350}\PYG{+w}{  }\PYG{l+m}{4}\PYG{+w}{  }\PYG{l+m}{180.000000}
\end{sphinxVerbatim}


\subsection{Description}
\label{\detokenize{angle_cosine_buck6d:description}}
\sphinxAtStartPar
The \sphinxstyleemphasis{cosine/buck6d} angle style uses the potential
\begin{equation*}
\begin{split}E = K \left[ 1 + \cos(n\theta - \theta_0)\right]\end{split}
\end{equation*}
\sphinxAtStartPar
where \(K\) is the energy constant, \(n\) is the periodic multiplicity and
\(\theta_0\) is the equilibrium angle.

\sphinxAtStartPar
The coefficients must be defined for each angle type via the
{\hyperref[\detokenize{angle_coeff::doc}]{\sphinxcrossref{\DUrole{doc}{angle\_coeff}}}} command as in the example above, or in
the data file or restart files read by the {\hyperref[\detokenize{read_data::doc}]{\sphinxcrossref{\DUrole{doc}{read\_data}}}}
or {\hyperref[\detokenize{read_restart::doc}]{\sphinxcrossref{\DUrole{doc}{read\_restart}}}} commands in the following order:
\begin{itemize}
\item {} 
\sphinxAtStartPar
\(K\) (energy)

\item {} 
\sphinxAtStartPar
\(n\)

\item {} 
\sphinxAtStartPar
\(\theta_0\) (degrees)

\end{itemize}

\sphinxAtStartPar
\(\theta_0\) is specified in degrees, but LAMMPS converts it to radians
internally.

\sphinxAtStartPar
Additional to the cosine term the \sphinxstyleemphasis{cosine/buck6d} angle style computes
the short range (vdW) interaction belonging to the
{\hyperref[\detokenize{pair_buck6d_coul_gauss::doc}]{\sphinxcrossref{\DUrole{doc}{pair\_style buck6d}}}} between the end atoms of the
angle.  For this reason this angle style only works in combination
with the {\hyperref[\detokenize{pair_buck6d_coul_gauss::doc}]{\sphinxcrossref{\DUrole{doc}{pair\_style buck6d}}}} styles and needs
the {\hyperref[\detokenize{special_bonds::doc}]{\sphinxcrossref{\DUrole{doc}{special\_bonds}}}} 1\sphinxhyphen{}3 interactions to be weighted
0.0 to prevent double counting.


\bigskip\hrule\bigskip



\subsection{Restrictions}
\label{\detokenize{angle_cosine_buck6d:restrictions}}
\sphinxAtStartPar
\sphinxstyleemphasis{cosine/buck6d} can only be used in combination with the
{\hyperref[\detokenize{pair_buck6d_coul_gauss::doc}]{\sphinxcrossref{\DUrole{doc}{pair\_style buck6d}}}} style and with a
{\hyperref[\detokenize{special_bonds::doc}]{\sphinxcrossref{\DUrole{doc}{special\_bonds}}}} 0.0 weighting of 1\sphinxhyphen{}3 interactions.

\sphinxAtStartPar
This angle style can only be used if LAMMPS was built with the
MOFFF package.  See the {\hyperref[\detokenize{Build_package::doc}]{\sphinxcrossref{\DUrole{doc}{Build package}}}} doc
page for more info.


\subsection{Related commands}
\label{\detokenize{angle_cosine_buck6d:related-commands}}
\sphinxAtStartPar
{\hyperref[\detokenize{angle_coeff::doc}]{\sphinxcrossref{\DUrole{doc}{angle\_coeff}}}}


\subsection{Default}
\label{\detokenize{angle_cosine_buck6d:default}}
\sphinxAtStartPar
none

\sphinxstepscope

\index{angle\_style cosine/delta@\spxentry{angle\_style cosine/delta}}\ignorespaces 
\index{angle\_style cosine/delta/omp@\spxentry{angle\_style cosine/delta/omp}}\ignorespaces 

\section{angle\_style cosine/delta command}
\label{\detokenize{angle_cosine_delta:angle-style-cosine-delta-command}}\label{\detokenize{angle_cosine_delta:index-1}}\label{\detokenize{angle_cosine_delta:index-0}}\label{\detokenize{angle_cosine_delta::doc}}
\sphinxAtStartPar
Accelerator Variants: \sphinxstyleemphasis{cosine/delta/omp}


\subsection{Syntax}
\label{\detokenize{angle_cosine_delta:syntax}}
\begin{sphinxVerbatim}[commandchars=\\\{\}]
\PYG{k}{angle\PYGZus{}style}\PYG{+w}{ }\PYG{n}{cosine}\PYG{o}{/}\PYG{n}{delta}
\end{sphinxVerbatim}


\subsection{Examples}
\label{\detokenize{angle_cosine_delta:examples}}
\begin{sphinxVerbatim}[commandchars=\\\{\}]
\PYG{k}{angle\PYGZus{}style}\PYG{+w}{ }\PYG{n}{cosine}\PYG{o}{/}\PYG{n}{delta}
\PYG{k}{angle\PYGZus{}coeff}\PYG{+w}{ }\PYG{l+m}{2}\PYG{o}{*}\PYG{l+m}{4}\PYG{+w}{ }\PYG{l+m}{75.0}\PYG{+w}{ }\PYG{l+m}{100.0}
\end{sphinxVerbatim}


\subsection{Description}
\label{\detokenize{angle_cosine_delta:description}}
\sphinxAtStartPar
The \sphinxstyleemphasis{cosine/delta} angle style uses the potential
\begin{equation*}
\begin{split}E = K [1 - \cos(\theta - \theta_0)]\end{split}
\end{equation*}
\sphinxAtStartPar
where \(\theta_0\) is the equilibrium value of the angle, and \(K\) is a
prefactor.  Note that the usual 1/2 factor is included in \(K\).

\sphinxAtStartPar
The following coefficients must be defined for each angle type via the
{\hyperref[\detokenize{angle_coeff::doc}]{\sphinxcrossref{\DUrole{doc}{angle\_coeff}}}} command as in the example above, or in
the data file or restart files read by the {\hyperref[\detokenize{read_data::doc}]{\sphinxcrossref{\DUrole{doc}{read\_data}}}}
or {\hyperref[\detokenize{read_restart::doc}]{\sphinxcrossref{\DUrole{doc}{read\_restart}}}} commands:
\begin{itemize}
\item {} 
\sphinxAtStartPar
\(K\) (energy)

\item {} 
\sphinxAtStartPar
\(\theta_0\) (degrees)

\end{itemize}

\sphinxAtStartPar
\(\theta_0\) is specified in degrees, but LAMMPS converts it to radians
internally.


\bigskip\hrule\bigskip


\sphinxAtStartPar
Styles with a \sphinxstyleemphasis{gpu}, \sphinxstyleemphasis{intel}, \sphinxstyleemphasis{kk}, \sphinxstyleemphasis{omp}, or \sphinxstyleemphasis{opt} suffix are
functionally the same as the corresponding style without the suffix.
They have been optimized to run faster, depending on your available
hardware, as discussed on the {\hyperref[\detokenize{Speed_packages::doc}]{\sphinxcrossref{\DUrole{doc}{Accelerator packages}}}}
page.  The accelerated styles take the same arguments and should
produce the same results, except for round\sphinxhyphen{}off and precision issues.

\sphinxAtStartPar
These accelerated styles are part of the GPU, INTEL, KOKKOS,
OPENMP, and OPT packages, respectively.  They are only enabled if
LAMMPS was built with those packages.  See the {\hyperref[\detokenize{Build_package::doc}]{\sphinxcrossref{\DUrole{doc}{Build package}}}} page for more info.

\sphinxAtStartPar
You can specify the accelerated styles explicitly in your input script
by including their suffix, or you can use the {\hyperref[\detokenize{Run_options::doc}]{\sphinxcrossref{\DUrole{doc}{\sphinxhyphen{}suffix command\sphinxhyphen{}line switch}}}} when you invoke LAMMPS, or you can use the
{\hyperref[\detokenize{suffix::doc}]{\sphinxcrossref{\DUrole{doc}{suffix}}}} command in your input script.

\sphinxAtStartPar
See the {\hyperref[\detokenize{Speed_packages::doc}]{\sphinxcrossref{\DUrole{doc}{Accelerator packages}}}} page for more
instructions on how to use the accelerated styles effectively.


\bigskip\hrule\bigskip



\subsection{Restrictions}
\label{\detokenize{angle_cosine_delta:restrictions}}
\sphinxAtStartPar
This angle style can only be used if LAMMPS was built with the
EXTRA\sphinxhyphen{}MOLECULE package.  See the {\hyperref[\detokenize{Build_package::doc}]{\sphinxcrossref{\DUrole{doc}{Build package}}}} doc page
for more info.


\subsection{Related commands}
\label{\detokenize{angle_cosine_delta:related-commands}}
\sphinxAtStartPar
{\hyperref[\detokenize{angle_coeff::doc}]{\sphinxcrossref{\DUrole{doc}{angle\_coeff}}}}, {\hyperref[\detokenize{angle_cosine_squared::doc}]{\sphinxcrossref{\DUrole{doc}{angle\_style cosine/squared}}}}


\subsection{Default}
\label{\detokenize{angle_cosine_delta:default}}
\sphinxAtStartPar
none

\sphinxstepscope

\index{angle\_style cosine/periodic@\spxentry{angle\_style cosine/periodic}}\ignorespaces 
\index{angle\_style cosine/periodic/omp@\spxentry{angle\_style cosine/periodic/omp}}\ignorespaces 

\section{angle\_style cosine/periodic command}
\label{\detokenize{angle_cosine_periodic:angle-style-cosine-periodic-command}}\label{\detokenize{angle_cosine_periodic:index-1}}\label{\detokenize{angle_cosine_periodic:index-0}}\label{\detokenize{angle_cosine_periodic::doc}}
\sphinxAtStartPar
Accelerator Variants: \sphinxstyleemphasis{cosine/periodic/omp}


\subsection{Syntax}
\label{\detokenize{angle_cosine_periodic:syntax}}
\begin{sphinxVerbatim}[commandchars=\\\{\}]
\PYG{k}{angle\PYGZus{}style}\PYG{+w}{ }\PYG{n}{cosine}\PYG{o}{/}\PYG{n}{periodic}
\end{sphinxVerbatim}


\subsection{Examples}
\label{\detokenize{angle_cosine_periodic:examples}}
\begin{sphinxVerbatim}[commandchars=\\\{\}]
\PYG{k}{angle\PYGZus{}style}\PYG{+w}{ }\PYG{n}{cosine}\PYG{o}{/}\PYG{n}{periodic}
\PYG{k}{angle\PYGZus{}coeff}\PYG{+w}{ }\PYG{o}{*}\PYG{+w}{ }\PYG{l+m}{75.0}\PYG{+w}{ }\PYG{l+m}{1}\PYG{+w}{ }\PYG{l+m}{6}
\end{sphinxVerbatim}


\subsection{Description}
\label{\detokenize{angle_cosine_periodic:description}}
\sphinxAtStartPar
The \sphinxstyleemphasis{cosine/periodic} angle style uses the following potential, which may be
particularly used for organometallic systems where \(n\) = 4 might be used
for an octahedral complex and \(n\) = 3 might be used for a trigonal
center:
\begin{equation*}
\begin{split}E = \frac{2.0}{n^2} * C \left[ 1 - B(-1)^n\cos\left( n\theta\right) \right]\end{split}
\end{equation*}
\sphinxAtStartPar
where \(C\), \(B\) and \(n\) are coefficients defined for each angle type.

\sphinxAtStartPar
The following coefficients must be defined for each angle type via the
{\hyperref[\detokenize{angle_coeff::doc}]{\sphinxcrossref{\DUrole{doc}{angle\_coeff}}}} command as in the example above, or in
the data file or restart files read by the {\hyperref[\detokenize{read_data::doc}]{\sphinxcrossref{\DUrole{doc}{read\_data}}}}
or {\hyperref[\detokenize{read_restart::doc}]{\sphinxcrossref{\DUrole{doc}{read\_restart}}}} commands:
\begin{itemize}
\item {} 
\sphinxAtStartPar
\(C\) (energy)

\item {} 
\sphinxAtStartPar
\(B\) = 1 or \sphinxhyphen{}1

\item {} 
\sphinxAtStartPar
\(n\) = 1, 2, 3, 4, 5 or 6 for periodicity

\end{itemize}

\sphinxAtStartPar
Note that the prefactor \(C\) is specified as coefficient and not the overall force
constant \(K = \frac{2 C}{n^2}\).  When \(B = 1\), it leads to a minimum for the
linear geometry.  When \(B = -1\), it leads to a maximum for the linear geometry.


\bigskip\hrule\bigskip


\sphinxAtStartPar
Styles with a \sphinxstyleemphasis{gpu}, \sphinxstyleemphasis{intel}, \sphinxstyleemphasis{kk}, \sphinxstyleemphasis{omp}, or \sphinxstyleemphasis{opt} suffix are
functionally the same as the corresponding style without the suffix.
They have been optimized to run faster, depending on your available
hardware, as discussed on the {\hyperref[\detokenize{Speed_packages::doc}]{\sphinxcrossref{\DUrole{doc}{Accelerator packages}}}}
page.  The accelerated styles take the same arguments and should
produce the same results, except for round\sphinxhyphen{}off and precision issues.

\sphinxAtStartPar
These accelerated styles are part of the GPU, INTEL, KOKKOS,
OPENMP, and OPT packages, respectively.  They are only enabled if
LAMMPS was built with those packages.  See the {\hyperref[\detokenize{Build_package::doc}]{\sphinxcrossref{\DUrole{doc}{Build package}}}} page for more info.

\sphinxAtStartPar
You can specify the accelerated styles explicitly in your input script
by including their suffix, or you can use the {\hyperref[\detokenize{Run_options::doc}]{\sphinxcrossref{\DUrole{doc}{\sphinxhyphen{}suffix command\sphinxhyphen{}line switch}}}} when you invoke LAMMPS, or you can use the
{\hyperref[\detokenize{suffix::doc}]{\sphinxcrossref{\DUrole{doc}{suffix}}}} command in your input script.

\sphinxAtStartPar
See the {\hyperref[\detokenize{Speed_packages::doc}]{\sphinxcrossref{\DUrole{doc}{Accelerator packages}}}} page for more
instructions on how to use the accelerated styles effectively.


\bigskip\hrule\bigskip



\subsection{Restrictions}
\label{\detokenize{angle_cosine_periodic:restrictions}}
\sphinxAtStartPar
This angle style can only be used if LAMMPS was built with the
EXTRA\sphinxhyphen{}MOLECULE package.  See the {\hyperref[\detokenize{Build_package::doc}]{\sphinxcrossref{\DUrole{doc}{Build package}}}} doc page
for more info.


\subsection{Related commands}
\label{\detokenize{angle_cosine_periodic:related-commands}}
\sphinxAtStartPar
{\hyperref[\detokenize{angle_coeff::doc}]{\sphinxcrossref{\DUrole{doc}{angle\_coeff}}}}


\subsection{Default}
\label{\detokenize{angle_cosine_periodic:default}}
\sphinxAtStartPar
none

\sphinxstepscope

\index{angle\_style cosine/shift@\spxentry{angle\_style cosine/shift}}\ignorespaces 
\index{angle\_style cosine/shift/omp@\spxentry{angle\_style cosine/shift/omp}}\ignorespaces 

\section{angle\_style cosine/shift command}
\label{\detokenize{angle_cosine_shift:angle-style-cosine-shift-command}}\label{\detokenize{angle_cosine_shift:index-1}}\label{\detokenize{angle_cosine_shift:index-0}}\label{\detokenize{angle_cosine_shift::doc}}
\sphinxAtStartPar
Accelerator Variants: \sphinxstyleemphasis{cosine/shift/omp}


\subsection{Syntax}
\label{\detokenize{angle_cosine_shift:syntax}}
\begin{sphinxVerbatim}[commandchars=\\\{\}]
\PYG{k}{angle\PYGZus{}style}\PYG{+w}{ }\PYG{n}{cosine}\PYG{o}{/}\PYG{n}{shift}
\end{sphinxVerbatim}


\subsection{Examples}
\label{\detokenize{angle_cosine_shift:examples}}
\begin{sphinxVerbatim}[commandchars=\\\{\}]
\PYG{k}{angle\PYGZus{}style}\PYG{+w}{ }\PYG{n}{cosine}\PYG{o}{/}\PYG{n}{shift}
\PYG{k}{angle\PYGZus{}coeff}\PYG{+w}{ }\PYG{o}{*}\PYG{+w}{ }\PYG{l+m}{10.0}\PYG{+w}{ }\PYG{l+m}{45.0}
\end{sphinxVerbatim}


\subsection{Description}
\label{\detokenize{angle_cosine_shift:description}}
\sphinxAtStartPar
The \sphinxstyleemphasis{cosine/shift} angle style uses the potential
\begin{equation*}
\begin{split}E = -\frac{U_{\text{min}}}{2} \left[ 1 + \cos(\theta-\theta_0) \right]\end{split}
\end{equation*}
\sphinxAtStartPar
where \(\theta_0\) is the equilibrium angle. The potential is bounded
between \(-U_{\text{min}}\) and zero. In the neighborhood of the minimum
\(E = - U_{\text{min}} + U_{\text{min}}/4(\theta - \theta_0)^2\) hence
the spring constant is \(\frac{U_{\text{min}}}{2}\).

\sphinxAtStartPar
The following coefficients must be defined for each angle type via the
{\hyperref[\detokenize{angle_coeff::doc}]{\sphinxcrossref{\DUrole{doc}{angle\_coeff}}}} command as in the example above, or in
the data file or restart files read by the {\hyperref[\detokenize{read_data::doc}]{\sphinxcrossref{\DUrole{doc}{read\_data}}}}
or {\hyperref[\detokenize{read_restart::doc}]{\sphinxcrossref{\DUrole{doc}{read\_restart}}}} commands:
\begin{itemize}
\item {} 
\sphinxAtStartPar
\(U_{\text{min}}\) (energy)

\item {} 
\sphinxAtStartPar
\(\theta\) (angle)

\end{itemize}


\bigskip\hrule\bigskip


\sphinxAtStartPar
Styles with a \sphinxstyleemphasis{gpu}, \sphinxstyleemphasis{intel}, \sphinxstyleemphasis{kk}, \sphinxstyleemphasis{omp}, or \sphinxstyleemphasis{opt} suffix are
functionally the same as the corresponding style without the suffix.
They have been optimized to run faster, depending on your available
hardware, as discussed on the {\hyperref[\detokenize{Speed_packages::doc}]{\sphinxcrossref{\DUrole{doc}{Accelerator packages}}}}
page.  The accelerated styles take the same arguments and should
produce the same results, except for round\sphinxhyphen{}off and precision issues.

\sphinxAtStartPar
These accelerated styles are part of the GPU, INTEL, KOKKOS,
OPENMP, and OPT packages, respectively.  They are only enabled if
LAMMPS was built with those packages.  See the {\hyperref[\detokenize{Build_package::doc}]{\sphinxcrossref{\DUrole{doc}{Build package}}}} page for more info.

\sphinxAtStartPar
You can specify the accelerated styles explicitly in your input script
by including their suffix, or you can use the {\hyperref[\detokenize{Run_options::doc}]{\sphinxcrossref{\DUrole{doc}{\sphinxhyphen{}suffix command\sphinxhyphen{}line switch}}}} when you invoke LAMMPS, or you can use the
{\hyperref[\detokenize{suffix::doc}]{\sphinxcrossref{\DUrole{doc}{suffix}}}} command in your input script.

\sphinxAtStartPar
See the {\hyperref[\detokenize{Speed_packages::doc}]{\sphinxcrossref{\DUrole{doc}{Accelerator packages}}}} page for more
instructions on how to use the accelerated styles effectively.


\bigskip\hrule\bigskip



\subsection{Restrictions}
\label{\detokenize{angle_cosine_shift:restrictions}}
\sphinxAtStartPar
This angle style can only be used if LAMMPS was built with the
EXTRA\sphinxhyphen{}MOLECULE package.  See the {\hyperref[\detokenize{Build_package::doc}]{\sphinxcrossref{\DUrole{doc}{Build package}}}} doc
page for more info.


\subsection{Related commands}
\label{\detokenize{angle_cosine_shift:related-commands}}
\sphinxAtStartPar
{\hyperref[\detokenize{angle_coeff::doc}]{\sphinxcrossref{\DUrole{doc}{angle\_coeff}}}},
{\hyperref[\detokenize{angle_cosine_shift_exp::doc}]{\sphinxcrossref{\DUrole{doc}{angle\_style cosine/shift/exp}}}}


\subsection{Default}
\label{\detokenize{angle_cosine_shift:default}}
\sphinxAtStartPar
none

\sphinxstepscope

\index{angle\_style cosine/shift/exp@\spxentry{angle\_style cosine/shift/exp}}\ignorespaces 
\index{angle\_style cosine/shift/exp/omp@\spxentry{angle\_style cosine/shift/exp/omp}}\ignorespaces 

\section{angle\_style cosine/shift/exp command}
\label{\detokenize{angle_cosine_shift_exp:angle-style-cosine-shift-exp-command}}\label{\detokenize{angle_cosine_shift_exp:index-1}}\label{\detokenize{angle_cosine_shift_exp:index-0}}\label{\detokenize{angle_cosine_shift_exp::doc}}
\sphinxAtStartPar
Accelerator Variants: \sphinxstyleemphasis{cosine/shift/exp/omp}


\subsection{Syntax}
\label{\detokenize{angle_cosine_shift_exp:syntax}}
\begin{sphinxVerbatim}[commandchars=\\\{\}]
\PYG{k}{angle\PYGZus{}style}\PYG{+w}{ }\PYG{n}{cosine}\PYG{o}{/}\PYG{n}{shift}\PYG{o}{/}\PYG{n}{exp}
\end{sphinxVerbatim}


\subsection{Examples}
\label{\detokenize{angle_cosine_shift_exp:examples}}
\begin{sphinxVerbatim}[commandchars=\\\{\}]
\PYG{k}{angle\PYGZus{}style}\PYG{+w}{ }\PYG{n}{cosine}\PYG{o}{/}\PYG{n}{shift}\PYG{o}{/}\PYG{n}{exp}
\PYG{k}{angle\PYGZus{}coeff}\PYG{+w}{ }\PYG{o}{*}\PYG{+w}{ }\PYG{l+m}{10.0}\PYG{+w}{ }\PYG{l+m}{45.0}\PYG{+w}{ }\PYG{l+m}{2.0}
\end{sphinxVerbatim}


\subsection{Description}
\label{\detokenize{angle_cosine_shift_exp:description}}
\sphinxAtStartPar
The \sphinxstyleemphasis{cosine/shift/exp} angle style uses the potential
\begin{equation*}
\begin{split}E = -U_{\text{min}} \frac{e^{-a U(\theta,\theta_0)}-1}{e^a-1} \quad \text{with} \quad U(\theta,\theta_0) = -0.5 \left(1+\cos(\theta-\theta_0) \right)\end{split}
\end{equation*}
\sphinxAtStartPar
where \(U_{\text{min}}\), \(\theta\), and \(a\) are defined for each angle type.

\sphinxAtStartPar
The potential is bounded between \([-U_{\text{min}}, 0]\) and the minimum is
located at the angle \(\theta_0\). The a parameter can be both positive or
negative and is used to control the spring constant at the
equilibrium.

\sphinxAtStartPar
The spring constant is given by \(k = A \exp(A) U_{\text{min}} / [2 (\exp(a)-1)]\).
For \(a > 3\), \(\frac{k}{U_{\text{min}}} = \frac{a}{2}\) to better than 5\% relative error. For negative
values of the \(a\) parameter, the spring constant is essentially zero,
and anharmonic terms takes over. The potential is furthermore well
behaved in the limit \(a \rightarrow 0\), where it has been implemented to linear
order in \(a\) for \(a < 0.001\). In this limit the potential reduces to the
cosineshifted potential.

\sphinxAtStartPar
The following coefficients must be defined for each angle type via the
{\hyperref[\detokenize{angle_coeff::doc}]{\sphinxcrossref{\DUrole{doc}{angle\_coeff}}}} command as in the example above, or in
the data file or restart files read by the {\hyperref[\detokenize{read_data::doc}]{\sphinxcrossref{\DUrole{doc}{read\_data}}}}
or {\hyperref[\detokenize{read_restart::doc}]{\sphinxcrossref{\DUrole{doc}{read\_restart}}}} commands:
\begin{itemize}
\item {} 
\sphinxAtStartPar
\(U_min\) (energy)

\item {} 
\sphinxAtStartPar
\(\theta\) (angle)

\item {} 
\sphinxAtStartPar
\(A\) (real number)

\end{itemize}


\bigskip\hrule\bigskip


\sphinxAtStartPar
Styles with a \sphinxstyleemphasis{gpu}, \sphinxstyleemphasis{intel}, \sphinxstyleemphasis{kk}, \sphinxstyleemphasis{omp}, or \sphinxstyleemphasis{opt} suffix are
functionally the same as the corresponding style without the suffix.
They have been optimized to run faster, depending on your available
hardware, as discussed on the {\hyperref[\detokenize{Speed_packages::doc}]{\sphinxcrossref{\DUrole{doc}{Accelerator packages}}}}
page.  The accelerated styles take the same arguments and should
produce the same results, except for round\sphinxhyphen{}off and precision issues.

\sphinxAtStartPar
These accelerated styles are part of the GPU, INTEL, KOKKOS,
OPENMP, and OPT packages, respectively.  They are only enabled if
LAMMPS was built with those packages.  See the {\hyperref[\detokenize{Build_package::doc}]{\sphinxcrossref{\DUrole{doc}{Build package}}}} page for more info.

\sphinxAtStartPar
You can specify the accelerated styles explicitly in your input script
by including their suffix, or you can use the {\hyperref[\detokenize{Run_options::doc}]{\sphinxcrossref{\DUrole{doc}{\sphinxhyphen{}suffix command\sphinxhyphen{}line switch}}}} when you invoke LAMMPS, or you can use the
{\hyperref[\detokenize{suffix::doc}]{\sphinxcrossref{\DUrole{doc}{suffix}}}} command in your input script.

\sphinxAtStartPar
See the {\hyperref[\detokenize{Speed_packages::doc}]{\sphinxcrossref{\DUrole{doc}{Accelerator packages}}}} page for more
instructions on how to use the accelerated styles effectively.


\bigskip\hrule\bigskip



\subsection{Restrictions}
\label{\detokenize{angle_cosine_shift_exp:restrictions}}
\sphinxAtStartPar
This angle style can only be used if LAMMPS was built with the
EXTRA\sphinxhyphen{}MOLECULE package.  See the {\hyperref[\detokenize{Build_package::doc}]{\sphinxcrossref{\DUrole{doc}{Build package}}}} doc
page for more info.


\subsection{Related commands}
\label{\detokenize{angle_cosine_shift_exp:related-commands}}
\sphinxAtStartPar
{\hyperref[\detokenize{angle_coeff::doc}]{\sphinxcrossref{\DUrole{doc}{angle\_coeff}}}},
{\hyperref[\detokenize{angle_cosine_shift::doc}]{\sphinxcrossref{\DUrole{doc}{angle\_style cosine/shift}}}},
{\hyperref[\detokenize{dihedral_cosine_shift_exp::doc}]{\sphinxcrossref{\DUrole{doc}{dihedral\_style cosine/shift/exp}}}}


\subsection{Default}
\label{\detokenize{angle_cosine_shift_exp:default}}
\sphinxAtStartPar
none

\sphinxstepscope

\index{angle\_style cosine/squared@\spxentry{angle\_style cosine/squared}}\ignorespaces 
\index{angle\_style cosine/squared/omp@\spxentry{angle\_style cosine/squared/omp}}\ignorespaces 

\section{angle\_style cosine/squared command}
\label{\detokenize{angle_cosine_squared:angle-style-cosine-squared-command}}\label{\detokenize{angle_cosine_squared:index-1}}\label{\detokenize{angle_cosine_squared:index-0}}\label{\detokenize{angle_cosine_squared::doc}}
\sphinxAtStartPar
Accelerator Variants: \sphinxstyleemphasis{cosine/squared/omp}


\subsection{Syntax}
\label{\detokenize{angle_cosine_squared:syntax}}
\begin{sphinxVerbatim}[commandchars=\\\{\}]
\PYG{k}{angle\PYGZus{}style}\PYG{+w}{ }\PYG{n}{cosine}\PYG{o}{/}\PYG{n}{squared}
\end{sphinxVerbatim}


\subsection{Examples}
\label{\detokenize{angle_cosine_squared:examples}}
\begin{sphinxVerbatim}[commandchars=\\\{\}]
\PYG{k}{angle\PYGZus{}style}\PYG{+w}{ }\PYG{n}{cosine}\PYG{o}{/}\PYG{n}{squared}
\PYG{k}{angle\PYGZus{}coeff}\PYG{+w}{ }\PYG{l+m}{2}\PYG{o}{*}\PYG{l+m}{4}\PYG{+w}{ }\PYG{l+m}{75.0}\PYG{+w}{ }\PYG{l+m}{100.0}
\end{sphinxVerbatim}


\subsection{Description}
\label{\detokenize{angle_cosine_squared:description}}
\sphinxAtStartPar
The \sphinxstyleemphasis{cosine/squared} angle style uses the potential
\begin{equation*}
\begin{split}E = K [\cos(\theta) - \cos(\theta_0)]^2\end{split}
\end{equation*}
\sphinxAtStartPar
, which is commonly used in the {\hyperref[\detokenize{Howto_bioFF::doc}]{\sphinxcrossref{\DUrole{doc}{DREIDING}}}} force field,
where \(\theta_0\) is the equilibrium value of the angle, and \(K\)
is a prefactor.  Note that the usual 1/2 factor is included in \(K\).

\sphinxAtStartPar
See {\hyperref[\detokenize{angle_cosine_squared:cosine-mayo}]{\sphinxcrossref{\DUrole{std,std-ref}{(Mayo)}}}} for a description of the DREIDING force field.

\sphinxAtStartPar
The following coefficients must be defined for each angle type via the
{\hyperref[\detokenize{angle_coeff::doc}]{\sphinxcrossref{\DUrole{doc}{angle\_coeff}}}} command as in the example above, or in
the data file or restart files read by the {\hyperref[\detokenize{read_data::doc}]{\sphinxcrossref{\DUrole{doc}{read\_data}}}}
or {\hyperref[\detokenize{read_restart::doc}]{\sphinxcrossref{\DUrole{doc}{read\_restart}}}} commands:
\begin{itemize}
\item {} 
\sphinxAtStartPar
\(K\) (energy)

\item {} 
\sphinxAtStartPar
\(\theta_0\) (degrees)

\end{itemize}

\sphinxAtStartPar
\(\theta_0\) is specified in degrees, but LAMMPS converts it to radians
internally.


\bigskip\hrule\bigskip


\sphinxAtStartPar
Styles with a \sphinxstyleemphasis{gpu}, \sphinxstyleemphasis{intel}, \sphinxstyleemphasis{kk}, \sphinxstyleemphasis{omp}, or \sphinxstyleemphasis{opt} suffix are
functionally the same as the corresponding style without the suffix.
They have been optimized to run faster, depending on your available
hardware, as discussed on the {\hyperref[\detokenize{Speed_packages::doc}]{\sphinxcrossref{\DUrole{doc}{Accelerator packages}}}}
page.  The accelerated styles take the same arguments and should
produce the same results, except for round\sphinxhyphen{}off and precision issues.

\sphinxAtStartPar
These accelerated styles are part of the GPU, INTEL, KOKKOS,
OPENMP, and OPT packages, respectively.  They are only enabled if
LAMMPS was built with those packages.  See the {\hyperref[\detokenize{Build_package::doc}]{\sphinxcrossref{\DUrole{doc}{Build package}}}} page for more info.

\sphinxAtStartPar
You can specify the accelerated styles explicitly in your input script
by including their suffix, or you can use the {\hyperref[\detokenize{Run_options::doc}]{\sphinxcrossref{\DUrole{doc}{\sphinxhyphen{}suffix command\sphinxhyphen{}line switch}}}} when you invoke LAMMPS, or you can use the
{\hyperref[\detokenize{suffix::doc}]{\sphinxcrossref{\DUrole{doc}{suffix}}}} command in your input script.

\sphinxAtStartPar
See the {\hyperref[\detokenize{Speed_packages::doc}]{\sphinxcrossref{\DUrole{doc}{Accelerator packages}}}} page for more
instructions on how to use the accelerated styles effectively.


\bigskip\hrule\bigskip



\subsection{Restrictions}
\label{\detokenize{angle_cosine_squared:restrictions}}
\sphinxAtStartPar
This angle style can only be used if LAMMPS was built with the
MOLECULE package.  See the {\hyperref[\detokenize{Build_package::doc}]{\sphinxcrossref{\DUrole{doc}{Build package}}}} doc page
for more info.


\subsection{Related commands}
\label{\detokenize{angle_cosine_squared:related-commands}}
\sphinxAtStartPar
{\hyperref[\detokenize{angle_coeff::doc}]{\sphinxcrossref{\DUrole{doc}{angle\_coeff}}}}


\subsection{Default}
\label{\detokenize{angle_cosine_squared:default}}
\sphinxAtStartPar
none


\bigskip\hrule\bigskip

\phantomsection\label{\detokenize{angle_cosine_squared:cosine-mayo}}
\sphinxAtStartPar
\sphinxstylestrong{(Mayo)} Mayo, Olfason, Goddard III, J Phys Chem, 94, 8897\sphinxhyphen{}8909
(1990).

\sphinxstepscope

\index{angle\_style cosine/squared/restricted@\spxentry{angle\_style cosine/squared/restricted}}\ignorespaces 
\index{angle\_style cosine/squared/restricted/omp@\spxentry{angle\_style cosine/squared/restricted/omp}}\ignorespaces 

\section{angle\_style cosine/squared/restricted command}
\label{\detokenize{angle_cosine_squared_restricted:angle-style-cosine-squared-restricted-command}}\label{\detokenize{angle_cosine_squared_restricted:index-1}}\label{\detokenize{angle_cosine_squared_restricted:index-0}}\label{\detokenize{angle_cosine_squared_restricted::doc}}
\sphinxAtStartPar
Accelerator Variants: \sphinxstyleemphasis{cosine/squared/restricted/omp}


\subsection{Syntax}
\label{\detokenize{angle_cosine_squared_restricted:syntax}}
\begin{sphinxVerbatim}[commandchars=\\\{\}]
\PYG{k}{angle\PYGZus{}style}\PYG{+w}{ }\PYG{n}{cosine}\PYG{o}{/}\PYG{n}{squared}\PYG{o}{/}\PYG{n}{restricted}
\end{sphinxVerbatim}


\subsection{Examples}
\label{\detokenize{angle_cosine_squared_restricted:examples}}
\begin{sphinxVerbatim}[commandchars=\\\{\}]
\PYG{k}{angle\PYGZus{}style}\PYG{+w}{ }\PYG{n}{cosine}\PYG{o}{/}\PYG{n}{squared}\PYG{o}{/}\PYG{n}{restricted}
\PYG{k}{angle\PYGZus{}coeff}\PYG{+w}{ }\PYG{l+m}{2}\PYG{o}{*}\PYG{l+m}{4}\PYG{+w}{ }\PYG{l+m}{75.0}\PYG{+w}{ }\PYG{l+m}{100.0}
\end{sphinxVerbatim}


\subsection{Description}
\label{\detokenize{angle_cosine_squared_restricted:description}}
\sphinxAtStartPar
\DUrole{versionmodified,added}{Added in version 17Apr2024.}

\sphinxAtStartPar
The \sphinxstyleemphasis{cosine/squared/restricted} angle style uses the potential
\begin{equation*}
\begin{split}E = K [\cos(\theta) - \cos(\theta_0)]^2 / \sin^2(\theta)\end{split}
\end{equation*}
\sphinxAtStartPar
, which is commonly used in the MARTINI force field,
where \(\theta_0\) is the equilibrium value of the angle, and \(K\)
is a prefactor.  Note that the usual 1/2 factor is included in \(K\).

\sphinxAtStartPar
See {\hyperref[\detokenize{angle_cosine_squared_restricted:restricted-bulacu}]{\sphinxcrossref{\DUrole{std,std-ref}{(Bulacu)}}}} for a description of the restricted angle for the MARTINI force field.

\sphinxAtStartPar
The following coefficients must be defined for each angle type via the
{\hyperref[\detokenize{angle_coeff::doc}]{\sphinxcrossref{\DUrole{doc}{angle\_coeff}}}} command as in the example above, or in
the data file or restart files read by the {\hyperref[\detokenize{read_data::doc}]{\sphinxcrossref{\DUrole{doc}{read\_data}}}}
or {\hyperref[\detokenize{read_restart::doc}]{\sphinxcrossref{\DUrole{doc}{read\_restart}}}} commands:
\begin{itemize}
\item {} 
\sphinxAtStartPar
\(K\) (energy)

\item {} 
\sphinxAtStartPar
\(\theta_0\) (degrees)

\end{itemize}

\sphinxAtStartPar
\(\theta_0\) is specified in degrees, but LAMMPS converts it to radians
internally.


\bigskip\hrule\bigskip


\sphinxAtStartPar
Styles with a \sphinxstyleemphasis{gpu}, \sphinxstyleemphasis{intel}, \sphinxstyleemphasis{kk}, \sphinxstyleemphasis{omp}, or \sphinxstyleemphasis{opt} suffix are
functionally the same as the corresponding style without the suffix.
They have been optimized to run faster, depending on your available
hardware, as discussed on the {\hyperref[\detokenize{Speed_packages::doc}]{\sphinxcrossref{\DUrole{doc}{Accelerator packages}}}}
page.  The accelerated styles take the same arguments and should
produce the same results, except for round\sphinxhyphen{}off and precision issues.

\sphinxAtStartPar
These accelerated styles are part of the GPU, INTEL, KOKKOS,
OPENMP, and OPT packages, respectively.  They are only enabled if
LAMMPS was built with those packages.  See the {\hyperref[\detokenize{Build_package::doc}]{\sphinxcrossref{\DUrole{doc}{Build package}}}} page for more info.

\sphinxAtStartPar
You can specify the accelerated styles explicitly in your input script
by including their suffix, or you can use the {\hyperref[\detokenize{Run_options::doc}]{\sphinxcrossref{\DUrole{doc}{\sphinxhyphen{}suffix command\sphinxhyphen{}line switch}}}} when you invoke LAMMPS, or you can use the
{\hyperref[\detokenize{suffix::doc}]{\sphinxcrossref{\DUrole{doc}{suffix}}}} command in your input script.

\sphinxAtStartPar
See the {\hyperref[\detokenize{Speed_packages::doc}]{\sphinxcrossref{\DUrole{doc}{Accelerator packages}}}} page for more
instructions on how to use the accelerated styles effectively.


\bigskip\hrule\bigskip



\subsection{Restrictions}
\label{\detokenize{angle_cosine_squared_restricted:restrictions}}
\sphinxAtStartPar
This angle style can only be used if LAMMPS was built with the
EXTRA\sphinxhyphen{}MOLECULE package.  See the {\hyperref[\detokenize{Build_package::doc}]{\sphinxcrossref{\DUrole{doc}{Build package}}}} doc page
for more info.


\subsection{Related commands}
\label{\detokenize{angle_cosine_squared_restricted:related-commands}}
\sphinxAtStartPar
{\hyperref[\detokenize{angle_coeff::doc}]{\sphinxcrossref{\DUrole{doc}{angle\_coeff}}}}


\subsection{Default}
\label{\detokenize{angle_cosine_squared_restricted:default}}
\sphinxAtStartPar
none


\bigskip\hrule\bigskip

\phantomsection\label{\detokenize{angle_cosine_squared_restricted:restricted-bulacu}}
\sphinxAtStartPar
\sphinxstylestrong{(Bulacu)} Bulacu, Goga, Zhao, Rossi, Monticelli, Periole, Tieleman, Marrink, J Chem Theory Comput, 9, 3282\sphinxhyphen{}3292
(2013).

\sphinxstepscope

\index{angle\_style cross@\spxentry{angle\_style cross}}\ignorespaces 

\section{angle\_style cross command}
\label{\detokenize{angle_cross:angle-style-cross-command}}\label{\detokenize{angle_cross:index-0}}\label{\detokenize{angle_cross::doc}}

\subsection{Syntax}
\label{\detokenize{angle_cross:syntax}}
\begin{sphinxVerbatim}[commandchars=\\\{\}]
\PYG{k}{angle\PYGZus{}style}\PYG{+w}{ }\PYG{n}{cross}
\end{sphinxVerbatim}


\subsection{Examples}
\label{\detokenize{angle_cross:examples}}
\begin{sphinxVerbatim}[commandchars=\\\{\}]
\PYG{k}{angle\PYGZus{}style}\PYG{+w}{ }\PYG{n}{cross}
\PYG{k}{angle\PYGZus{}coeff}\PYG{+w}{ }\PYG{l+m}{1}\PYG{+w}{ }\PYG{l+m}{200.0}\PYG{+w}{ }\PYG{l+m}{100.0}\PYG{+w}{ }\PYG{l+m}{100.0}\PYG{+w}{ }\PYG{l+m}{1.25}\PYG{+w}{ }\PYG{l+m}{1.25}\PYG{+w}{ }\PYG{l+m}{107.0}
\end{sphinxVerbatim}


\subsection{Description}
\label{\detokenize{angle_cross:description}}
\sphinxAtStartPar
The \sphinxstyleemphasis{cross} angle style uses a potential that couples the bond stretches of
a bend with the angle stretch of that bend:
\begin{equation*}
\begin{split}E = K_{SS} \left(r_{12}-r_{12,0}\right)\left(r_{32}-r_{32,0}\right) + K_{BS0}\left(r_{12}-r_{12,0}\right)\left(\theta-\theta_0\right) + K_{BS1}\left(r_{32}-r_{32,0}\right)\left(\theta-\theta_0\right)\end{split}
\end{equation*}
\sphinxAtStartPar
where \(r_{12,0}\) is the rest value of the bond length between atom 1 and 2,
\(r_{32,0}\) is the rest value of the bond length between atom 3 and 2,
and \(\theta_0\) is the rest value of the angle. \(K_{SS}\) is the force constant of
the bond stretch\sphinxhyphen{}bond stretch term and \(K_{BS0}\) and \(K_{BS1}\) are the force constants
of the bond stretch\sphinxhyphen{}angle stretch terms.

\sphinxAtStartPar
The following coefficients must be defined for each angle type via the
{\hyperref[\detokenize{angle_coeff::doc}]{\sphinxcrossref{\DUrole{doc}{angle\_coeff}}}} command as in the example above, or in
the data file or restart files read by the {\hyperref[\detokenize{read_data::doc}]{\sphinxcrossref{\DUrole{doc}{read\_data}}}}
or {\hyperref[\detokenize{read_restart::doc}]{\sphinxcrossref{\DUrole{doc}{read\_restart}}}} commands:
\begin{itemize}
\item {} 
\sphinxAtStartPar
\(K_{SS}\) (energy/distance\textasciicircum{}2)

\item {} 
\sphinxAtStartPar
\(K_{BS0}\) (energy/distance)

\item {} 
\sphinxAtStartPar
\(K_{BS1}\) (energy/distance)

\item {} 
\sphinxAtStartPar
\(r_{12,0}\) (distance)

\item {} 
\sphinxAtStartPar
\(r_{32,0}\) (distance)

\item {} 
\sphinxAtStartPar
\(\theta_0\) (degrees)

\end{itemize}

\sphinxAtStartPar
\(\theta_0\) is specified in degrees, but LAMMPS converts it to
radians internally; hence the \(K_{BS0}\) and \(K_{BS1}\) are
effectively energy/distance per radian.


\subsection{Restrictions}
\label{\detokenize{angle_cross:restrictions}}
\sphinxAtStartPar
This angle style can only be used if LAMMPS was built with the
YAFF package.  See the {\hyperref[\detokenize{Build_package::doc}]{\sphinxcrossref{\DUrole{doc}{Build package}}}} doc
page for more info.


\subsection{Related commands}
\label{\detokenize{angle_cross:related-commands}}
\sphinxAtStartPar
{\hyperref[\detokenize{angle_coeff::doc}]{\sphinxcrossref{\DUrole{doc}{angle\_coeff}}}}


\subsection{Default}
\label{\detokenize{angle_cross:default}}
\sphinxAtStartPar
none

\sphinxstepscope

\index{angle\_style dipole@\spxentry{angle\_style dipole}}\ignorespaces 
\index{angle\_style dipole/omp@\spxentry{angle\_style dipole/omp}}\ignorespaces 

\section{angle\_style dipole command}
\label{\detokenize{angle_dipole:angle-style-dipole-command}}\label{\detokenize{angle_dipole:index-1}}\label{\detokenize{angle_dipole:index-0}}\label{\detokenize{angle_dipole::doc}}
\sphinxAtStartPar
Accelerator Variants: \sphinxstyleemphasis{dipole/omp}


\subsection{Syntax}
\label{\detokenize{angle_dipole:syntax}}
\begin{sphinxVerbatim}[commandchars=\\\{\}]
\PYG{k}{angle\PYGZus{}style}\PYG{+w}{ }\PYG{n}{dipole}
\end{sphinxVerbatim}


\subsection{Examples}
\label{\detokenize{angle_dipole:examples}}
\begin{sphinxVerbatim}[commandchars=\\\{\}]
\PYG{k}{angle\PYGZus{}style}\PYG{+w}{ }\PYG{n}{dipole}
\PYG{k}{angle\PYGZus{}coeff}\PYG{+w}{ }\PYG{l+m}{6}\PYG{+w}{ }\PYG{l+m}{2.1}\PYG{+w}{ }\PYG{l+m}{180.0}
\end{sphinxVerbatim}


\subsection{Description}
\label{\detokenize{angle_dipole:description}}
\sphinxAtStartPar
The \sphinxstyleemphasis{dipole} angle style is used to control the orientation of a dipolar
atom within a molecule {\hyperref[\detokenize{angle_dipole:orsi}]{\sphinxcrossref{\DUrole{std,std-ref}{(Orsi)}}}}. Specifically, the \sphinxstyleemphasis{dipole} angle
style restrains the orientation of a point dipole \(\mu_j\) (embedded in atom
\(j\)) with respect to a reference (bond) vector
\(\vec{r}_{ij} = \vec{r}_i - \vec{r}_j\), where \(i\) is another atom of
the same molecule (typically, \(i\) and \(j\) are also covalently bonded).

\sphinxAtStartPar
It is convenient to define an angle gamma between the ‘free’ vector \(\vec{\mu}_j\)
and the reference (bond) vector \(\vec{r}_{ij}\):
\begin{equation*}
\begin{split}\cos\gamma = \frac{\vec{\mu}_j\cdot\vec{r}_{ij}}{\mu_j\,r_{ij}}\end{split}
\end{equation*}
\sphinxAtStartPar
The \sphinxstyleemphasis{dipole} angle style uses the potential:
\begin{equation*}
\begin{split}E = K (\cos\gamma - \cos\gamma_0)^2\end{split}
\end{equation*}
\sphinxAtStartPar
where \(K\) is a rigidity constant and gamma0 is an equilibrium (reference)
angle.

\sphinxAtStartPar
The torque on the dipole can be obtained by differentiating the
potential using the ‘chain rule’ as in appendix C.3 of
{\hyperref[\detokenize{angle_dipole:allen1}]{\sphinxcrossref{\DUrole{std,std-ref}{(AllenTildesley)}}}}:
\begin{equation*}
\begin{split}\vec{T}_j = \frac{2K(\cos\gamma - \cos\gamma_0)}{\mu_j\,r_{ij}}\, \vec{r}_{ij} \times \vec{\mu}_j\end{split}
\end{equation*}
\sphinxAtStartPar
Example: if \(\gamma_0\) is set to 0 degrees, the torque generated by
the potential will tend to align the dipole along the reference
direction defined by the (bond) vector \(\vec{r}_{ij}\) (in other words, \(\vec{\mu}_j\) is
restrained to point towards atom \(i\)).

\sphinxAtStartPar
The dipolar torque \(\vec{T}_j\) must be counterbalanced in order to conserve
the local angular momentum. This is achieved via an additional force
couple generating a torque equivalent to the opposite of \(\vec{T}_j\):
\begin{equation*}
\begin{split}-\vec{T}_j & = \vec{r}_{ij} \times \vec{F}_i \\
\vec{F}_j  & = -\vec{F}_i\end{split}
\end{equation*}
\sphinxAtStartPar
where \(\vec{F}_i\) and \(\vec{F}_j\) are applied on atoms \(i\)
and \(j\), respectively.

\sphinxAtStartPar
The following coefficients must be defined for each angle type via the
{\hyperref[\detokenize{angle_coeff::doc}]{\sphinxcrossref{\DUrole{doc}{angle\_coeff}}}} command as in the example above, or in
the data file or restart files read by the {\hyperref[\detokenize{read_data::doc}]{\sphinxcrossref{\DUrole{doc}{read\_data}}}}
or {\hyperref[\detokenize{read_restart::doc}]{\sphinxcrossref{\DUrole{doc}{read\_restart}}}} commands:
\begin{itemize}
\item {} 
\sphinxAtStartPar
\(K\) (energy)

\item {} 
\sphinxAtStartPar
\(\gamma_0\) (degrees)

\end{itemize}


\bigskip\hrule\bigskip


\sphinxAtStartPar
Styles with a \sphinxstyleemphasis{gpu}, \sphinxstyleemphasis{intel}, \sphinxstyleemphasis{kk}, \sphinxstyleemphasis{omp}, or \sphinxstyleemphasis{opt} suffix are
functionally the same as the corresponding style without the suffix.
They have been optimized to run faster, depending on your available
hardware, as discussed on the {\hyperref[\detokenize{Speed_packages::doc}]{\sphinxcrossref{\DUrole{doc}{Accelerator packages}}}}
page.  The accelerated styles take the same arguments and should
produce the same results, except for round\sphinxhyphen{}off and precision issues.

\sphinxAtStartPar
These accelerated styles are part of the GPU, INTEL, KOKKOS,
OPENMP, and OPT packages, respectively.  They are only enabled if
LAMMPS was built with those packages.  See the {\hyperref[\detokenize{Build_package::doc}]{\sphinxcrossref{\DUrole{doc}{Build package}}}} page for more info.

\sphinxAtStartPar
You can specify the accelerated styles explicitly in your input script
by including their suffix, or you can use the {\hyperref[\detokenize{Run_options::doc}]{\sphinxcrossref{\DUrole{doc}{\sphinxhyphen{}suffix command\sphinxhyphen{}line switch}}}} when you invoke LAMMPS, or you can use the
{\hyperref[\detokenize{suffix::doc}]{\sphinxcrossref{\DUrole{doc}{suffix}}}} command in your input script.

\sphinxAtStartPar
See the {\hyperref[\detokenize{Speed_packages::doc}]{\sphinxcrossref{\DUrole{doc}{Accelerator packages}}}} page for more
instructions on how to use the accelerated styles effectively.


\subsection{Restrictions}
\label{\detokenize{angle_dipole:restrictions}}
\sphinxAtStartPar
This angle style can only be used if LAMMPS was built with the
DIPOLE package.  See the {\hyperref[\detokenize{Build_package::doc}]{\sphinxcrossref{\DUrole{doc}{Build package}}}} doc
page for more info.

\begin{sphinxadmonition}{note}{Note:}
\sphinxAtStartPar
In the “Angles” section of the data file, the atom ID \(j\)
defining the direction of the dipole vector to restrain must come
before the atom ID of the reference atom \(i\). A third atom ID \(k\) must
also be provided to comply with the requirement of a valid angle
definition. This atom ID \(k\) should be chosen to be that of an atom
bonded to atom \(i\) to avoid errors with “lost angle atoms” when running
in parallel. Since the LAMMPS code checks for valid angle definitions,
cannot use the same atom ID of either \(i\) or \(j\) (this was allowed
and recommended with older LAMMPS versions).
\end{sphinxadmonition}

\sphinxAtStartPar
The {\hyperref[\detokenize{newton::doc}]{\sphinxcrossref{\DUrole{doc}{newton}}}} command for intramolecular interactions must be “on”
(which is the default except when using some accelerator packages).

\begin{sphinxadmonition}{note}{Note:}
\sphinxAtStartPar
This angle style should \sphinxstylestrong{NOT} be used with fix shake.
\end{sphinxadmonition}


\subsection{Related commands}
\label{\detokenize{angle_dipole:related-commands}}
\sphinxAtStartPar
{\hyperref[\detokenize{angle_coeff::doc}]{\sphinxcrossref{\DUrole{doc}{angle\_coeff}}}}, {\hyperref[\detokenize{angle_hybrid::doc}]{\sphinxcrossref{\DUrole{doc}{angle\_hybrid}}}}


\subsection{Default}
\label{\detokenize{angle_dipole:default}}
\sphinxAtStartPar
none


\bigskip\hrule\bigskip

\phantomsection\label{\detokenize{angle_dipole:orsi}}
\sphinxAtStartPar
\sphinxstylestrong{(Orsi)} Orsi \& Essex, The ELBA force field for coarse\sphinxhyphen{}grain modeling of
lipid membranes, PloS ONE 6(12): e28637, 2011.

\phantomsection\label{\detokenize{angle_dipole:allen1}}
\sphinxAtStartPar
\sphinxstylestrong{(AllenTildesley)} Allen \& Tildesley, Computer Simulation of Liquids,
Clarendon Press, Oxford, 1987.

\sphinxstepscope

\index{angle\_style fourier@\spxentry{angle\_style fourier}}\ignorespaces 
\index{angle\_style fourier/omp@\spxentry{angle\_style fourier/omp}}\ignorespaces 

\section{angle\_style fourier command}
\label{\detokenize{angle_fourier:angle-style-fourier-command}}\label{\detokenize{angle_fourier:index-1}}\label{\detokenize{angle_fourier:index-0}}\label{\detokenize{angle_fourier::doc}}
\sphinxAtStartPar
Accelerator Variants: \sphinxstyleemphasis{fourier/omp}


\subsection{Syntax}
\label{\detokenize{angle_fourier:syntax}}
\begin{sphinxVerbatim}[commandchars=\\\{\}]
\PYG{k}{angle\PYGZus{}style}\PYG{+w}{ }\PYG{n}{fourier}
\end{sphinxVerbatim}


\subsection{Examples}
\label{\detokenize{angle_fourier:examples}}
\begin{sphinxVerbatim}[commandchars=\\\{\}]
\PYG{k}{angle\PYGZus{}style}\PYG{+w}{ }\PYG{n}{fourier}
\PYG{k}{angle\PYGZus{}coeff}\PYG{+w}{ }\PYG{l+m}{75.0}\PYG{+w}{ }\PYG{l+m}{1.0}\PYG{+w}{ }\PYG{l+m}{1.0}\PYG{+w}{ }\PYG{l+m}{1.0}
\end{sphinxVerbatim}


\subsection{Description}
\label{\detokenize{angle_fourier:description}}
\sphinxAtStartPar
The \sphinxstyleemphasis{fourier} angle style uses the potential
\begin{equation*}
\begin{split}E = K [C_0 + C_1 \cos ( \theta) + C_2 \cos( 2 \theta) ]\end{split}
\end{equation*}
\sphinxAtStartPar
The following coefficients must be defined for each angle type via the
{\hyperref[\detokenize{angle_coeff::doc}]{\sphinxcrossref{\DUrole{doc}{angle\_coeff}}}} command as in the example above, or in
the data file or restart files read by the {\hyperref[\detokenize{read_data::doc}]{\sphinxcrossref{\DUrole{doc}{read\_data}}}}
or {\hyperref[\detokenize{read_restart::doc}]{\sphinxcrossref{\DUrole{doc}{read\_restart}}}} commands:
\begin{itemize}
\item {} 
\sphinxAtStartPar
\(K\) (energy)

\item {} 
\sphinxAtStartPar
\(C_0\) (real)

\item {} 
\sphinxAtStartPar
\(C_1\) (real)

\item {} 
\sphinxAtStartPar
\(C_2\) (real)

\end{itemize}


\bigskip\hrule\bigskip


\sphinxAtStartPar
Styles with a \sphinxstyleemphasis{gpu}, \sphinxstyleemphasis{intel}, \sphinxstyleemphasis{kk}, \sphinxstyleemphasis{omp}, or \sphinxstyleemphasis{opt} suffix are
functionally the same as the corresponding style without the suffix.
They have been optimized to run faster, depending on your available
hardware, as discussed on the {\hyperref[\detokenize{Speed_packages::doc}]{\sphinxcrossref{\DUrole{doc}{Accelerator packages}}}}
page.  The accelerated styles take the same arguments and should
produce the same results, except for round\sphinxhyphen{}off and precision issues.

\sphinxAtStartPar
These accelerated styles are part of the GPU, INTEL, KOKKOS,
OPENMP, and OPT packages, respectively.  They are only enabled if
LAMMPS was built with those packages.  See the {\hyperref[\detokenize{Build_package::doc}]{\sphinxcrossref{\DUrole{doc}{Build package}}}} page for more info.

\sphinxAtStartPar
You can specify the accelerated styles explicitly in your input script
by including their suffix, or you can use the {\hyperref[\detokenize{Run_options::doc}]{\sphinxcrossref{\DUrole{doc}{\sphinxhyphen{}suffix command\sphinxhyphen{}line switch}}}} when you invoke LAMMPS, or you can use the
{\hyperref[\detokenize{suffix::doc}]{\sphinxcrossref{\DUrole{doc}{suffix}}}} command in your input script.

\sphinxAtStartPar
See the {\hyperref[\detokenize{Speed_packages::doc}]{\sphinxcrossref{\DUrole{doc}{Accelerator packages}}}} page for more
instructions on how to use the accelerated styles effectively.


\bigskip\hrule\bigskip



\subsection{Restrictions}
\label{\detokenize{angle_fourier:restrictions}}
\sphinxAtStartPar
This angle style can only be used if LAMMPS was built with the
EXTRA\sphinxhyphen{}MOLECULE package.  See the {\hyperref[\detokenize{Build_package::doc}]{\sphinxcrossref{\DUrole{doc}{Build package}}}} doc
page for more info.


\subsection{Related commands}
\label{\detokenize{angle_fourier:related-commands}}
\sphinxAtStartPar
{\hyperref[\detokenize{angle_coeff::doc}]{\sphinxcrossref{\DUrole{doc}{angle\_coeff}}}}


\subsection{Default}
\label{\detokenize{angle_fourier:default}}
\sphinxAtStartPar
none

\sphinxstepscope

\index{angle\_style fourier/simple@\spxentry{angle\_style fourier/simple}}\ignorespaces 
\index{angle\_style fourier/simple/omp@\spxentry{angle\_style fourier/simple/omp}}\ignorespaces 

\section{angle\_style fourier/simple command}
\label{\detokenize{angle_fourier_simple:angle-style-fourier-simple-command}}\label{\detokenize{angle_fourier_simple:index-1}}\label{\detokenize{angle_fourier_simple:index-0}}\label{\detokenize{angle_fourier_simple::doc}}
\sphinxAtStartPar
Accelerator Variants: \sphinxstyleemphasis{fourier/simple/omp}


\subsection{Syntax}
\label{\detokenize{angle_fourier_simple:syntax}}
\begin{sphinxVerbatim}[commandchars=\\\{\}]
\PYG{k}{angle\PYGZus{}style}\PYG{+w}{ }\PYG{n}{fourier}\PYG{o}{/}\PYG{n}{simple}
\end{sphinxVerbatim}


\subsection{Examples}
\label{\detokenize{angle_fourier_simple:examples}}
\begin{sphinxVerbatim}[commandchars=\\\{\}]
\PYG{k}{angle\PYGZus{}style}\PYG{+w}{ }\PYG{n}{fourier}\PYG{o}{/}\PYG{n}{simple}
\PYG{k}{angle\PYGZus{}coeff}\PYG{+w}{ }\PYG{l+m}{100.0}\PYG{+w}{ }\PYG{o}{\PYGZhy{}}\PYG{l+m}{1.0}\PYG{+w}{ }\PYG{l+m}{1.0}
\end{sphinxVerbatim}


\subsection{Description}
\label{\detokenize{angle_fourier_simple:description}}
\sphinxAtStartPar
The \sphinxstyleemphasis{fourier/simple} angle style uses the potential
\begin{equation*}
\begin{split}E = K [ 1.0 + c \cos ( n \theta) ]\end{split}
\end{equation*}
\sphinxAtStartPar
The following coefficients must be defined for each angle type via the
{\hyperref[\detokenize{angle_coeff::doc}]{\sphinxcrossref{\DUrole{doc}{angle\_coeff}}}} command as in the example above, or in
the data file or restart files read by the {\hyperref[\detokenize{read_data::doc}]{\sphinxcrossref{\DUrole{doc}{read\_data}}}}
or {\hyperref[\detokenize{read_restart::doc}]{\sphinxcrossref{\DUrole{doc}{read\_restart}}}} commands:
\begin{itemize}
\item {} 
\sphinxAtStartPar
\(K\) (energy)

\item {} 
\sphinxAtStartPar
\(c\) (real)

\item {} 
\sphinxAtStartPar
\(n\) (real)

\end{itemize}


\bigskip\hrule\bigskip


\sphinxAtStartPar
Styles with a \sphinxstyleemphasis{gpu}, \sphinxstyleemphasis{intel}, \sphinxstyleemphasis{kk}, \sphinxstyleemphasis{omp}, or \sphinxstyleemphasis{opt} suffix are
functionally the same as the corresponding style without the suffix.
They have been optimized to run faster, depending on your available
hardware, as discussed on the {\hyperref[\detokenize{Speed_packages::doc}]{\sphinxcrossref{\DUrole{doc}{Accelerator packages}}}}
page.  The accelerated styles take the same arguments and should
produce the same results, except for round\sphinxhyphen{}off and precision issues.

\sphinxAtStartPar
These accelerated styles are part of the GPU, INTEL, KOKKOS,
OPENMP, and OPT packages, respectively.  They are only enabled if
LAMMPS was built with those packages.  See the {\hyperref[\detokenize{Build_package::doc}]{\sphinxcrossref{\DUrole{doc}{Build package}}}} page for more info.

\sphinxAtStartPar
You can specify the accelerated styles explicitly in your input script
by including their suffix, or you can use the {\hyperref[\detokenize{Run_options::doc}]{\sphinxcrossref{\DUrole{doc}{\sphinxhyphen{}suffix command\sphinxhyphen{}line switch}}}} when you invoke LAMMPS, or you can use the
{\hyperref[\detokenize{suffix::doc}]{\sphinxcrossref{\DUrole{doc}{suffix}}}} command in your input script.

\sphinxAtStartPar
See the {\hyperref[\detokenize{Speed_packages::doc}]{\sphinxcrossref{\DUrole{doc}{Accelerator packages}}}} page for more
instructions on how to use the accelerated styles effectively.


\bigskip\hrule\bigskip



\subsection{Restrictions}
\label{\detokenize{angle_fourier_simple:restrictions}}
\sphinxAtStartPar
This angle style can only be used if LAMMPS was built with the
EXTRA\sphinxhyphen{}MOLECULE package.  See the {\hyperref[\detokenize{Build_package::doc}]{\sphinxcrossref{\DUrole{doc}{Build package}}}} doc
page for more info.


\subsection{Related commands}
\label{\detokenize{angle_fourier_simple:related-commands}}
\sphinxAtStartPar
{\hyperref[\detokenize{angle_coeff::doc}]{\sphinxcrossref{\DUrole{doc}{angle\_coeff}}}}


\subsection{Default}
\label{\detokenize{angle_fourier_simple:default}}
\sphinxAtStartPar
none

\sphinxstepscope

\index{angle\_style gaussian@\spxentry{angle\_style gaussian}}\ignorespaces 

\section{angle\_style gaussian command}
\label{\detokenize{angle_gaussian:angle-style-gaussian-command}}\label{\detokenize{angle_gaussian:index-0}}\label{\detokenize{angle_gaussian::doc}}

\subsection{Syntax}
\label{\detokenize{angle_gaussian:syntax}}
\begin{sphinxVerbatim}[commandchars=\\\{\}]
\PYG{k}{angle\PYGZus{}style}\PYG{+w}{ }\PYG{n}{gaussian}
\end{sphinxVerbatim}


\subsection{Examples}
\label{\detokenize{angle_gaussian:examples}}
\begin{sphinxVerbatim}[commandchars=\\\{\}]
\PYG{k}{angle\PYGZus{}style}\PYG{+w}{ }\PYG{n}{gaussian}
\PYG{k}{angle\PYGZus{}coeff}\PYG{+w}{ }\PYG{l+m}{1}\PYG{+w}{ }\PYG{l+m}{300.0}\PYG{+w}{ }\PYG{l+m}{2}\PYG{+w}{ }\PYG{l+m}{0.0128}\PYG{+w}{ }\PYG{l+m}{0.375}\PYG{+w}{ }\PYG{l+m}{80.0}\PYG{+w}{ }\PYG{l+m}{0.0730}\PYG{+w}{ }\PYG{l+m}{0.148}\PYG{+w}{ }\PYG{l+m}{123.0}
\end{sphinxVerbatim}


\subsection{Description}
\label{\detokenize{angle_gaussian:description}}
\sphinxAtStartPar
The \sphinxstyleemphasis{gaussian} angle style uses the potential:
\begin{equation*}
\begin{split}E = -k_B T ln\left(\sum_{i=1}^{n} \frac{A_i}{w_i \sqrt{\pi/2}} exp\left( \frac{-2(\theta-\theta_{i})^2}{w_i^2}\right) \right)\end{split}
\end{equation*}
\sphinxAtStartPar
This analytical form is a suitable potential for obtaining mesoscale
effective force fields which can reproduce target atomistic
distributions {\hyperref[\detokenize{angle_gaussian:milano1}]{\sphinxcrossref{\DUrole{std,std-ref}{(Milano)}}}}.

\sphinxAtStartPar
The following coefficients must be defined for each angle type via the
{\hyperref[\detokenize{angle_coeff::doc}]{\sphinxcrossref{\DUrole{doc}{angle\_coeff}}}} command as in the example above, or in
the data file or restart files read by the {\hyperref[\detokenize{read_data::doc}]{\sphinxcrossref{\DUrole{doc}{read\_data}}}}
or {\hyperref[\detokenize{read_restart::doc}]{\sphinxcrossref{\DUrole{doc}{read\_restart}}}} commands:
\begin{itemize}
\item {} 
\sphinxAtStartPar
\(T\) temperature at which the potential was derived

\item {} 
\sphinxAtStartPar
\(n\) (integer \textgreater{}=1)

\item {} 
\sphinxAtStartPar
\(A_1\) (\textgreater{} 0, radians)

\item {} 
\sphinxAtStartPar
\(w_1\) (\textgreater{} 0, radians)

\item {} 
\sphinxAtStartPar
\(\theta_1\) (degrees)

\item {} 
\sphinxAtStartPar
…

\item {} 
\sphinxAtStartPar
\(A_n\) (\textgreater{} 0, radians)

\item {} 
\sphinxAtStartPar
\(w_n\) (\textgreater{} 0, radians)

\item {} 
\sphinxAtStartPar
\(\theta_n\) (degrees)

\end{itemize}


\subsection{Restrictions}
\label{\detokenize{angle_gaussian:restrictions}}
\sphinxAtStartPar
This angle style can only be used if LAMMPS was built with the
EXTRA\sphinxhyphen{}MOLECULE package.  See the {\hyperref[\detokenize{Build_package::doc}]{\sphinxcrossref{\DUrole{doc}{Build package}}}} doc
page for more info.


\subsection{Related commands}
\label{\detokenize{angle_gaussian:related-commands}}
\sphinxAtStartPar
{\hyperref[\detokenize{angle_coeff::doc}]{\sphinxcrossref{\DUrole{doc}{angle\_coeff}}}}


\subsection{Default}
\label{\detokenize{angle_gaussian:default}}
\sphinxAtStartPar
none


\bigskip\hrule\bigskip

\phantomsection\label{\detokenize{angle_gaussian:milano1}}
\sphinxAtStartPar
\sphinxstylestrong{(Milano)} G. Milano, S. Goudeau, F. Mueller\sphinxhyphen{}Plathe, J. Polym. Sci. B Polym. Phys. 43, 871 (2005).

\sphinxstepscope

\index{angle\_style harmonic@\spxentry{angle\_style harmonic}}\ignorespaces 
\index{angle\_style harmonic/intel@\spxentry{angle\_style harmonic/intel}}\ignorespaces 
\index{angle\_style harmonic/kk@\spxentry{angle\_style harmonic/kk}}\ignorespaces 
\index{angle\_style harmonic/omp@\spxentry{angle\_style harmonic/omp}}\ignorespaces 

\section{angle\_style harmonic command}
\label{\detokenize{angle_harmonic:angle-style-harmonic-command}}\label{\detokenize{angle_harmonic:index-3}}\label{\detokenize{angle_harmonic:index-2}}\label{\detokenize{angle_harmonic:index-1}}\label{\detokenize{angle_harmonic:index-0}}\label{\detokenize{angle_harmonic::doc}}
\sphinxAtStartPar
Accelerator Variants: \sphinxstyleemphasis{harmonic/intel}, \sphinxstyleemphasis{harmonic/kk}, \sphinxstyleemphasis{harmonic/omp}


\subsection{Syntax}
\label{\detokenize{angle_harmonic:syntax}}
\begin{sphinxVerbatim}[commandchars=\\\{\}]
\PYG{k}{angle\PYGZus{}style}\PYG{+w}{ }\PYG{n}{harmonic}
\end{sphinxVerbatim}


\subsection{Examples}
\label{\detokenize{angle_harmonic:examples}}
\begin{sphinxVerbatim}[commandchars=\\\{\}]
\PYG{k}{angle\PYGZus{}style}\PYG{+w}{ }\PYG{n}{harmonic}
\PYG{k}{angle\PYGZus{}coeff}\PYG{+w}{ }\PYG{l+m}{1}\PYG{+w}{ }\PYG{l+m}{300.0}\PYG{+w}{ }\PYG{l+m}{107.0}
\end{sphinxVerbatim}


\subsection{Description}
\label{\detokenize{angle_harmonic:description}}
\sphinxAtStartPar
The \sphinxstyleemphasis{harmonic} angle style uses the potential
\begin{equation*}
\begin{split}E = K (\theta - \theta_0)^2\end{split}
\end{equation*}
\sphinxAtStartPar
where \(\theta_0\) is the equilibrium value of the angle, and \(K\) is a
prefactor.  Note that the usual 1/2 factor is included in \(K\).

\sphinxAtStartPar
The following coefficients must be defined for each angle type via the
{\hyperref[\detokenize{angle_coeff::doc}]{\sphinxcrossref{\DUrole{doc}{angle\_coeff}}}} command as in the example above, or in
the data file or restart files read by the {\hyperref[\detokenize{read_data::doc}]{\sphinxcrossref{\DUrole{doc}{read\_data}}}}
or {\hyperref[\detokenize{read_restart::doc}]{\sphinxcrossref{\DUrole{doc}{read\_restart}}}} commands:
\begin{itemize}
\item {} 
\sphinxAtStartPar
\(K\) (energy)

\item {} 
\sphinxAtStartPar
\(\theta_0\) (degrees)

\end{itemize}

\sphinxAtStartPar
\(\theta_0\) is specified in degrees, but LAMMPS converts it to
radians internally; hence \(K\) is effectively energy per
radian\textasciicircum{}2.


\bigskip\hrule\bigskip


\sphinxAtStartPar
Styles with a \sphinxstyleemphasis{gpu}, \sphinxstyleemphasis{intel}, \sphinxstyleemphasis{kk}, \sphinxstyleemphasis{omp}, or \sphinxstyleemphasis{opt} suffix are
functionally the same as the corresponding style without the suffix.
They have been optimized to run faster, depending on your available
hardware, as discussed on the {\hyperref[\detokenize{Speed_packages::doc}]{\sphinxcrossref{\DUrole{doc}{Accelerator packages}}}}
page.  The accelerated styles take the same arguments and should
produce the same results, except for round\sphinxhyphen{}off and precision issues.

\sphinxAtStartPar
These accelerated styles are part of the GPU, INTEL, KOKKOS,
OPENMP, and OPT packages, respectively.  They are only enabled if
LAMMPS was built with those packages.  See the {\hyperref[\detokenize{Build_package::doc}]{\sphinxcrossref{\DUrole{doc}{Build package}}}} page for more info.

\sphinxAtStartPar
You can specify the accelerated styles explicitly in your input script
by including their suffix, or you can use the {\hyperref[\detokenize{Run_options::doc}]{\sphinxcrossref{\DUrole{doc}{\sphinxhyphen{}suffix command\sphinxhyphen{}line switch}}}} when you invoke LAMMPS, or you can use the
{\hyperref[\detokenize{suffix::doc}]{\sphinxcrossref{\DUrole{doc}{suffix}}}} command in your input script.

\sphinxAtStartPar
See the {\hyperref[\detokenize{Speed_packages::doc}]{\sphinxcrossref{\DUrole{doc}{Accelerator packages}}}} page for more
instructions on how to use the accelerated styles effectively.


\bigskip\hrule\bigskip



\subsection{Restrictions}
\label{\detokenize{angle_harmonic:restrictions}}
\sphinxAtStartPar
This angle style can only be used if LAMMPS was built with the
MOLECULE package.  See the {\hyperref[\detokenize{Build_package::doc}]{\sphinxcrossref{\DUrole{doc}{Build package}}}} doc
page for more info.


\subsection{Related commands}
\label{\detokenize{angle_harmonic:related-commands}}
\sphinxAtStartPar
{\hyperref[\detokenize{angle_coeff::doc}]{\sphinxcrossref{\DUrole{doc}{angle\_coeff}}}}


\subsection{Default}
\label{\detokenize{angle_harmonic:default}}
\sphinxAtStartPar
none

\sphinxstepscope

\index{angle\_style hybrid@\spxentry{angle\_style hybrid}}\ignorespaces 
\index{angle\_style hybrid/kk@\spxentry{angle\_style hybrid/kk}}\ignorespaces 

\section{angle\_style hybrid command}
\label{\detokenize{angle_hybrid:angle-style-hybrid-command}}\label{\detokenize{angle_hybrid:index-1}}\label{\detokenize{angle_hybrid:index-0}}\label{\detokenize{angle_hybrid::doc}}
\sphinxAtStartPar
Accelerator Variants: \sphinxstyleemphasis{hybrid/kk}


\subsection{Syntax}
\label{\detokenize{angle_hybrid:syntax}}
\begin{sphinxVerbatim}[commandchars=\\\{\}]
\PYG{k}{angle\PYGZus{}style}\PYG{+w}{ }\PYG{n}{hybrid}\PYG{+w}{ }\PYG{n}{style1}\PYG{+w}{ }\PYG{n}{style2}\PYG{+w}{ }\PYG{n}{...}
\end{sphinxVerbatim}
\begin{itemize}
\item {} 
\sphinxAtStartPar
style1,style2 = list of one or more angle styles

\end{itemize}


\subsection{Examples}
\label{\detokenize{angle_hybrid:examples}}
\begin{sphinxVerbatim}[commandchars=\\\{\}]
\PYG{k}{angle\PYGZus{}style}\PYG{+w}{ }\PYG{n}{hybrid}\PYG{+w}{ }\PYG{n}{harmonic}\PYG{+w}{ }\PYG{n}{cosine}
\PYG{k}{angle\PYGZus{}coeff}\PYG{+w}{ }\PYG{l+m}{1}\PYG{+w}{ }\PYG{n}{harmonic}\PYG{+w}{ }\PYG{l+m}{80.0}\PYG{+w}{ }\PYG{l+m}{30.0}
\PYG{k}{angle\PYGZus{}coeff}\PYG{+w}{ }\PYG{l+m}{2}\PYG{o}{*}\PYG{+w}{ }\PYG{n}{cosine}\PYG{+w}{ }\PYG{l+m}{50.0}
\end{sphinxVerbatim}


\subsection{Description}
\label{\detokenize{angle_hybrid:description}}
\sphinxAtStartPar
The \sphinxstyleemphasis{hybrid} style enables the use of multiple angle styles in one
simulation.  An angle style is assigned to each angle type.  For
example, angles in a polymer flow (of angle type 1) could be computed
with a \sphinxstyleemphasis{harmonic} potential and angles in the wall boundary (of angle
type 2) could be computed with a \sphinxstyleemphasis{cosine} potential.  The assignment
of angle type to style is made via the {\hyperref[\detokenize{angle_coeff::doc}]{\sphinxcrossref{\DUrole{doc}{angle\_coeff}}}}
command or in the data file.

\sphinxAtStartPar
In the {\hyperref[\detokenize{angle_coeff::doc}]{\sphinxcrossref{\DUrole{doc}{angle\_coeff}}}} commands, the name of an angle style must be added
after the angle type, with the remaining coefficients being those
appropriate to that style.  In the example above, the 2 angle\_coeff
commands set angles of angle type 1 to be computed with a \sphinxstyleemphasis{harmonic}
potential with coefficients 80.0, 30.0 for \(K\), \(\theta_0\).  All other angle
types \((2 - N)\) are computed with a \sphinxstyleemphasis{cosine} potential with coefficient
50.0 for \(K\).

\sphinxAtStartPar
If angle coefficients are specified in the data file read via the
{\hyperref[\detokenize{read_data::doc}]{\sphinxcrossref{\DUrole{doc}{read\_data}}}} command, then the same rule applies.
E.g. “harmonic” or “cosine”, must be added after the angle type, for each
line in the “Angle Coeffs” section, e.g.

\begin{sphinxVerbatim}[commandchars=\\\{\}]
Angle Coeffs

1 harmonic 80.0 30.0
2 cosine 50.0
...
\end{sphinxVerbatim}

\sphinxAtStartPar
If \sphinxstyleemphasis{class2} is one of the angle hybrid styles, the same rule holds for
specifying additional BondBond (and BondAngle) coefficients either via
the input script or in the data file.  I.e. \sphinxstyleemphasis{class2} must be added to
each line after the angle type.  For lines in the BondBond (or
BondAngle) section of the data file for angle types that are not
\sphinxstyleemphasis{class2}, you must use an angle style of \sphinxstyleemphasis{skip} as a placeholder, e.g.

\begin{sphinxVerbatim}[commandchars=\\\{\}]
BondBond Coeffs

1 skip
2 class2 3.6512 1.0119 1.0119
...
\end{sphinxVerbatim}

\sphinxAtStartPar
Note that it is not necessary to use the angle style \sphinxstyleemphasis{skip} in the
input script, since BondBond (or BondAngle) coefficients need not be
specified at all for angle types that are not \sphinxstyleemphasis{class2}.

\sphinxAtStartPar
An angle style of \sphinxstyleemphasis{none} with no additional coefficients can be used
in place of an angle style, either in a input script {\hyperref[\detokenize{angle_coeff::doc}]{\sphinxcrossref{\DUrole{doc}{angle\_coeff}}}}
command or in the data file, if you desire to turn off interactions
for specific angle types.


\bigskip\hrule\bigskip


\sphinxAtStartPar
Styles with a \sphinxstyleemphasis{gpu}, \sphinxstyleemphasis{intel}, \sphinxstyleemphasis{kk}, \sphinxstyleemphasis{omp}, or \sphinxstyleemphasis{opt} suffix are
functionally the same as the corresponding style without the suffix.
They have been optimized to run faster, depending on your available
hardware, as discussed on the {\hyperref[\detokenize{Speed_packages::doc}]{\sphinxcrossref{\DUrole{doc}{Accelerator packages}}}}
page.  The accelerated styles take the same arguments and should
produce the same results, except for round\sphinxhyphen{}off and precision issues.

\sphinxAtStartPar
These accelerated styles are part of the GPU, INTEL, KOKKOS,
OPENMP, and OPT packages, respectively.  They are only enabled if
LAMMPS was built with those packages.  See the {\hyperref[\detokenize{Build_package::doc}]{\sphinxcrossref{\DUrole{doc}{Build package}}}} page for more info.

\sphinxAtStartPar
You can specify the accelerated styles explicitly in your input script
by including their suffix, or you can use the {\hyperref[\detokenize{Run_options::doc}]{\sphinxcrossref{\DUrole{doc}{\sphinxhyphen{}suffix command\sphinxhyphen{}line switch}}}} when you invoke LAMMPS, or you can use the
{\hyperref[\detokenize{suffix::doc}]{\sphinxcrossref{\DUrole{doc}{suffix}}}} command in your input script.

\sphinxAtStartPar
See the {\hyperref[\detokenize{Speed_packages::doc}]{\sphinxcrossref{\DUrole{doc}{Accelerator packages}}}} page for more
instructions on how to use the accelerated styles effectively.


\bigskip\hrule\bigskip



\subsection{Restrictions}
\label{\detokenize{angle_hybrid:restrictions}}
\sphinxAtStartPar
This angle style can only be used if LAMMPS was built with the
MOLECULE package.  See the {\hyperref[\detokenize{Build_package::doc}]{\sphinxcrossref{\DUrole{doc}{Build package}}}} doc page
for more info.

\sphinxAtStartPar
Unlike other angle styles, the hybrid angle style does not store angle
coefficient info for individual sub\sphinxhyphen{}styles in {\hyperref[\detokenize{restart::doc}]{\sphinxcrossref{\DUrole{doc}{binary restart files}}}} or {\hyperref[\detokenize{write_data::doc}]{\sphinxcrossref{\DUrole{doc}{data files}}}}.  Thus when restarting a
simulation, you need to re\sphinxhyphen{}specify the angle\_coeff commands.


\subsection{Related commands}
\label{\detokenize{angle_hybrid:related-commands}}
\sphinxAtStartPar
{\hyperref[\detokenize{angle_coeff::doc}]{\sphinxcrossref{\DUrole{doc}{angle\_coeff}}}}


\subsection{Default}
\label{\detokenize{angle_hybrid:default}}
\sphinxAtStartPar
none

\sphinxstepscope

\index{angle\_style lepton@\spxentry{angle\_style lepton}}\ignorespaces 
\index{angle\_style lepton/omp@\spxentry{angle\_style lepton/omp}}\ignorespaces 

\section{angle\_style lepton command}
\label{\detokenize{angle_lepton:angle-style-lepton-command}}\label{\detokenize{angle_lepton:index-1}}\label{\detokenize{angle_lepton:index-0}}\label{\detokenize{angle_lepton::doc}}
\sphinxAtStartPar
Accelerator Variants: \sphinxstyleemphasis{lepton/omp}


\subsection{Syntax}
\label{\detokenize{angle_lepton:syntax}}
\begin{sphinxVerbatim}[commandchars=\\\{\}]
\PYG{k}{angle\PYGZus{}style}\PYG{+w}{ }\PYG{n}{style}\PYG{+w}{ }\PYG{n}{args}
\end{sphinxVerbatim}
\begin{itemize}
\item {} 
\sphinxAtStartPar
style = \sphinxstyleemphasis{lepton}

\item {} 
\sphinxAtStartPar
args = optional arguments

\end{itemize}
\begin{sphinxalltt}
args = \sphinxstyleemphasis{auto\_offset} or \sphinxstyleemphasis{no\_offset}
  \sphinxstyleemphasis{auto\_offset} = offset the potential energy so that the value at theta0 is 0.0 (default)
  \sphinxstyleemphasis{no\_offset} = do not offset the potential energy
\end{sphinxalltt}


\subsection{Examples}
\label{\detokenize{angle_lepton:examples}}
\begin{sphinxVerbatim}[commandchars=\\\{\}]
\PYG{k}{angle\PYGZus{}style}\PYG{+w}{ }\PYG{n}{lepton}
\PYG{k}{angle\PYGZus{}style}\PYG{+w}{ }\PYG{n}{lepton}\PYG{+w}{ }\PYG{n}{no\PYGZus{}offset}

\PYG{k}{angle\PYGZus{}coeff}\PYG{+w}{  }\PYG{l+m}{1}\PYG{+w}{  }\PYG{l+m}{120.0}\PYG{+w}{  }\PYG{l+s}{\PYGZdq{}}\PYG{l+s}{k*theta\PYGZca{}2; k=250.0}\PYG{l+s}{\PYGZdq{}}
\PYG{k}{angle\PYGZus{}coeff}\PYG{+w}{  }\PYG{l+m}{2}\PYG{+w}{   }\PYG{l+m}{90.0}\PYG{+w}{  }\PYG{l+s}{\PYGZdq{}}\PYG{l+s}{k2*theta\PYGZca{}2 + k3*theta\PYGZca{}3 + k4*theta\PYGZca{}4; k2=300.0; k3=\PYGZhy{}100.0; k4=50.0}\PYG{l+s}{\PYGZdq{}}
\PYG{k}{angle\PYGZus{}coeff}\PYG{+w}{  }\PYG{l+m}{3}\PYG{+w}{  }\PYG{l+m}{109.47}\PYG{+w}{ }\PYG{l+s}{\PYGZdq{}}\PYG{l+s}{k*theta\PYGZca{}2; k=350.0}\PYG{l+s}{\PYGZdq{}}
\end{sphinxVerbatim}


\subsection{Description}
\label{\detokenize{angle_lepton:description}}
\sphinxAtStartPar
\DUrole{versionmodified,added}{Added in version 8Feb2023.}

\sphinxAtStartPar
Angle style \sphinxstyleemphasis{lepton} computes angular interactions between three atoms
with a custom potential function.  The potential function must be
provided as an expression string using “theta” as the angle variable
relative to the reference angle \(\theta_0\) which is provided as an
angle coefficient.  For example \sphinxtitleref{“200.0*theta\textasciicircum{}2”} represents a
{\hyperref[\detokenize{angle_harmonic::doc}]{\sphinxcrossref{\DUrole{doc}{harmonic angle}}}} potential with a force constant
\sphinxstyleemphasis{K} of 200.0 energy units:
\begin{equation*}
\begin{split}U_{angle,i} = K (\theta_i - \theta_0)^2 = K \theta^2 \qquad \theta = \theta_i - \theta_0\end{split}
\end{equation*}
\sphinxAtStartPar
\DUrole{versionmodified,changed}{Changed in version 7Feb2024.}

\sphinxAtStartPar
By default the potential energy U is shifted so that the value U is 0.0
for \$theta = theta\_0\$.  This is equivalent to using the optional keyword
\sphinxstyleemphasis{auto\_offset}.  When using the keyword \sphinxstyleemphasis{no\_offset} instead, the
potential energy is not shifted.

\sphinxAtStartPar
The \sphinxhref{https://simtk.org/projects/lepton}{Lepton library}, that the
\sphinxstyleemphasis{lepton} angle style interfaces with, evaluates this expression string
at run time to compute the pairwise energy.  It also creates an
analytical representation of the first derivative of this expression
with respect to “theta” and then uses that to compute the force between
the angle atoms as defined by the topology data.

\sphinxAtStartPar
The following coefficients must be defined for each angle type via the
{\hyperref[\detokenize{angle_coeff::doc}]{\sphinxcrossref{\DUrole{doc}{angle\_coeff}}}} command as in the example above, or in
the data file or restart files read by the {\hyperref[\detokenize{read_data::doc}]{\sphinxcrossref{\DUrole{doc}{read\_data}}}}
or {\hyperref[\detokenize{read_restart::doc}]{\sphinxcrossref{\DUrole{doc}{read\_restart}}}} commands:
\begin{itemize}
\item {} 
\sphinxAtStartPar
Lepton expression (energy units)

\item {} 
\sphinxAtStartPar
\(\theta_0\) (degrees)

\end{itemize}

\sphinxAtStartPar
The Lepton expression must be either enclosed in quotes or must not
contain any whitespace so that LAMMPS recognizes it as a single keyword.
More on valid Lepton expressions below.  The \(\theta_0\)
coefficient is the “equilibrium angle”.  It is entered in degrees, but
internally converted to radians.  Thus the expression must assume
“theta” is in radians.  The potential energy function in the Lepton
expression is shifted in such a way, that the potential energy is 0 for
a angle \(\theta_i == \theta_0\).


\bigskip\hrule\bigskip



\subsection{Lepton expression syntax and features}
\label{\detokenize{angle_lepton:lepton-expression-syntax-and-features}}
\sphinxAtStartPar
Lepton supports the following operators in expressions:


\begin{savenotes}\sphinxattablestart
\sphinxthistablewithglobalstyle
\centering
\begin{tabulary}{\linewidth}[t]{TTTTTTTTTTTTTT}
\sphinxtoprule
\sphinxtableatstartofbodyhook
\sphinxAtStartPar
+
&
\sphinxAtStartPar
Add
&&
\sphinxAtStartPar
\sphinxhyphen{}
&
\sphinxAtStartPar
Subtract
&&
\sphinxAtStartPar
*
&
\sphinxAtStartPar
Multiply
&&
\sphinxAtStartPar
/
&
\sphinxAtStartPar
Divide
&&
\sphinxAtStartPar
\textasciicircum{}
&
\sphinxAtStartPar
Power
\\
\sphinxbottomrule
\end{tabulary}
\sphinxtableafterendhook\par
\sphinxattableend\end{savenotes}

\sphinxAtStartPar
The following mathematical functions are available:


\begin{savenotes}\sphinxattablestart
\sphinxthistablewithglobalstyle
\centering
\begin{tabulary}{\linewidth}[t]{TTTT}
\sphinxtoprule
\sphinxtableatstartofbodyhook
\sphinxAtStartPar
sqrt(x)
&
\sphinxAtStartPar
Square root
&
\sphinxAtStartPar
exp(x)
&
\sphinxAtStartPar
Exponential
\\
\sphinxhline
\sphinxAtStartPar
log(x)
&
\sphinxAtStartPar
Natural logarithm
&
\sphinxAtStartPar
sin(x)
&
\sphinxAtStartPar
Sine (angle in radians)
\\
\sphinxhline
\sphinxAtStartPar
cos(x)
&
\sphinxAtStartPar
Cosine (angle in radians)
&
\sphinxAtStartPar
sec(x)
&
\sphinxAtStartPar
Secant (angle in radians)
\\
\sphinxhline
\sphinxAtStartPar
csc(x)
&
\sphinxAtStartPar
Cosecant (angle in radians)
&
\sphinxAtStartPar
tan(x)
&
\sphinxAtStartPar
Tangent (angle in radians)
\\
\sphinxhline
\sphinxAtStartPar
cot(x)
&
\sphinxAtStartPar
Cotangent (angle in radians)
&
\sphinxAtStartPar
asin(x)
&
\sphinxAtStartPar
Inverse sine (in radians)
\\
\sphinxhline
\sphinxAtStartPar
acos(x)
&
\sphinxAtStartPar
Inverse cosine (in radians)
&
\sphinxAtStartPar
atan(x)
&
\sphinxAtStartPar
Inverse tangent (in radians)
\\
\sphinxhline
\sphinxAtStartPar
sinh(x)
&
\sphinxAtStartPar
Hyperbolic sine
&
\sphinxAtStartPar
cosh(x)
&
\sphinxAtStartPar
Hyperbolic cosine
\\
\sphinxhline
\sphinxAtStartPar
tanh(x)
&
\sphinxAtStartPar
Hyperbolic tangent
&
\sphinxAtStartPar
erf(x)
&
\sphinxAtStartPar
Error function
\\
\sphinxhline
\sphinxAtStartPar
erfc(x)
&
\sphinxAtStartPar
Complementary Error function
&
\sphinxAtStartPar
abs(x)
&
\sphinxAtStartPar
Absolute value
\\
\sphinxhline
\sphinxAtStartPar
min(x,y)
&
\sphinxAtStartPar
Minimum of two values
&
\sphinxAtStartPar
max(x,y)
&
\sphinxAtStartPar
Maximum of two values
\\
\sphinxhline
\sphinxAtStartPar
delta(x)
&
\sphinxAtStartPar
delta(x) is 1 for \sphinxtitleref{x = 0}, otherwise 0
&
\sphinxAtStartPar
step(x)
&
\sphinxAtStartPar
step(x) is 0 for \sphinxtitleref{x \textless{} 0}, otherwise 1
\\
\sphinxbottomrule
\end{tabulary}
\sphinxtableafterendhook\par
\sphinxattableend\end{savenotes}

\sphinxAtStartPar
Numbers may be given in either decimal or exponential form.  All of the
following are valid numbers: \sphinxtitleref{5}, \sphinxtitleref{\sphinxhyphen{}3.1}, \sphinxtitleref{1e6}, and \sphinxtitleref{3.12e\sphinxhyphen{}2}.

\sphinxAtStartPar
As an extension to the standard Lepton syntax, it is also possible to
use LAMMPS {\hyperref[\detokenize{variable::doc}]{\sphinxcrossref{\DUrole{doc}{variables}}}} in the format “v\_name”.  Before
evaluating the expression, “v\_name” will be replaced with the value of
the variable “name”.  This is compatible with all kinds of scalar
variables, but not with vectors, arrays, local, or per\sphinxhyphen{}atom
variables.  If necessary, a custom scalar variable needs to be defined
that can access the desired (single) item from a non\sphinxhyphen{}scalar variable.
As an example, the following lines will instruct LAMMPS to ramp
the force constant for a harmonic bond from 100.0 to 200.0 during the
next run:

\begin{sphinxVerbatim}[commandchars=\\\{\}]
\PYG{k}{variable }\PYG{n+nv+nvIdentifier}{fconst}\PYG{+w}{ }\PYG{n}{equal}\PYG{+w}{ }\PYG{n}{ramp}\PYG{n+nv}{(}\PYG{n+nv}{100.0, 200}\PYG{n+nv}{)}
\PYG{k}{bond\PYGZus{}style}\PYG{+w}{ }\PYG{n}{lepton}
\PYG{k}{bond\PYGZus{}coeff}\PYG{+w}{ }\PYG{l+m}{1}\PYG{+w}{ }\PYG{l+m}{1.5}\PYG{+w}{ }\PYG{l+s}{\PYGZdq{}}\PYG{l+s}{v\PYGZus{}fconst * (r\PYGZca{}2)}\PYG{l+s}{\PYGZdq{}}
\end{sphinxVerbatim}

\sphinxAtStartPar
An expression may be followed by definitions for intermediate values that appear in the
expression. A semicolon “;” is used as a delimiter between value definitions. For example,
the expression:

\begin{sphinxVerbatim}[commandchars=\\\{\}]
\PYG{n}{a}\PYG{o}{\PYGZca{}}\PYG{l+m+mi}{2}\PYG{o}{+}\PYG{n}{a}\PYG{o}{*}\PYG{n}{b}\PYG{o}{+}\PYG{n}{b}\PYG{o}{\PYGZca{}}\PYG{l+m+mi}{2}\PYG{p}{;}\PYG{+w}{ }\PYG{n}{a}\PYG{o}{=}\PYG{n}{a1}\PYG{o}{+}\PYG{n}{a2}\PYG{p}{;}\PYG{+w}{ }\PYG{n}{b}\PYG{o}{=}\PYG{n}{b1}\PYG{o}{+}\PYG{n}{b2}
\end{sphinxVerbatim}

\sphinxAtStartPar
is exactly equivalent to

\begin{sphinxVerbatim}[commandchars=\\\{\}]
\PYG{p}{(}\PYG{n}{a1}\PYG{o}{+}\PYG{n}{a2}\PYG{p}{)}\PYG{o}{\PYGZca{}}\PYG{l+m+mi}{2}\PYG{o}{+}\PYG{p}{(}\PYG{n}{a1}\PYG{o}{+}\PYG{n}{a2}\PYG{p}{)}\PYG{o}{*}\PYG{p}{(}\PYG{n}{b1}\PYG{o}{+}\PYG{n}{b2}\PYG{p}{)}\PYG{o}{+}\PYG{p}{(}\PYG{n}{b1}\PYG{o}{+}\PYG{n}{b2}\PYG{p}{)}\PYG{o}{\PYGZca{}}\PYG{l+m+mi}{2}
\end{sphinxVerbatim}

\sphinxAtStartPar
The definition of an intermediate value may itself involve other
intermediate values. Whitespace and quotation characters (’\textquotesingle{}’ and ‘”’)
are ignored.  All uses of a value must appear \sphinxstyleemphasis{before} that value’s
definition.  For efficiency reasons, the expression string is parsed,
optimized, and then stored in an internal, pre\sphinxhyphen{}parsed representation for
evaluation.

\sphinxAtStartPar
Evaluating a Lepton expression is typically between 2.5 and 5 times
slower than the corresponding compiled and optimized C++ code.  If
additional speed or GPU acceleration (via GPU or KOKKOS) is required,
the interaction can be represented as a table.  Suitable table files
can be created either internally using the {\hyperref[\detokenize{pair_write::doc}]{\sphinxcrossref{\DUrole{doc}{pair\_write}}}}
or {\hyperref[\detokenize{bond_write::doc}]{\sphinxcrossref{\DUrole{doc}{bond\_write}}}} command or through the Python scripts
in the {\hyperref[\detokenize{Tools:tabulate}]{\sphinxcrossref{\DUrole{std,std-ref}{tools/tabulate}}}} folder.


\bigskip\hrule\bigskip


\sphinxAtStartPar
Styles with a \sphinxstyleemphasis{gpu}, \sphinxstyleemphasis{intel}, \sphinxstyleemphasis{kk}, \sphinxstyleemphasis{omp}, or \sphinxstyleemphasis{opt} suffix are
functionally the same as the corresponding style without the suffix.
They have been optimized to run faster, depending on your available
hardware, as discussed on the {\hyperref[\detokenize{Speed_packages::doc}]{\sphinxcrossref{\DUrole{doc}{Accelerator packages}}}}
page.  The accelerated styles take the same arguments and should
produce the same results, except for round\sphinxhyphen{}off and precision issues.

\sphinxAtStartPar
These accelerated styles are part of the GPU, INTEL, KOKKOS,
OPENMP, and OPT packages, respectively.  They are only enabled if
LAMMPS was built with those packages.  See the {\hyperref[\detokenize{Build_package::doc}]{\sphinxcrossref{\DUrole{doc}{Build package}}}} page for more info.

\sphinxAtStartPar
You can specify the accelerated styles explicitly in your input script
by including their suffix, or you can use the {\hyperref[\detokenize{Run_options::doc}]{\sphinxcrossref{\DUrole{doc}{\sphinxhyphen{}suffix command\sphinxhyphen{}line switch}}}} when you invoke LAMMPS, or you can use the
{\hyperref[\detokenize{suffix::doc}]{\sphinxcrossref{\DUrole{doc}{suffix}}}} command in your input script.

\sphinxAtStartPar
See the {\hyperref[\detokenize{Speed_packages::doc}]{\sphinxcrossref{\DUrole{doc}{Accelerator packages}}}} page for more
instructions on how to use the accelerated styles effectively.


\bigskip\hrule\bigskip



\subsection{Restrictions}
\label{\detokenize{angle_lepton:restrictions}}
\sphinxAtStartPar
This angle style is part of the LEPTON package and only enabled if LAMMPS
was built with this package.  See the {\hyperref[\detokenize{Build_package::doc}]{\sphinxcrossref{\DUrole{doc}{Build package}}}} page for more info.


\subsection{Related commands}
\label{\detokenize{angle_lepton:related-commands}}
\sphinxAtStartPar
{\hyperref[\detokenize{angle_coeff::doc}]{\sphinxcrossref{\DUrole{doc}{angle\_coeff}}}}, {\hyperref[\detokenize{angle_table::doc}]{\sphinxcrossref{\DUrole{doc}{angle\_style table}}}},
{\hyperref[\detokenize{bond_lepton::doc}]{\sphinxcrossref{\DUrole{doc}{bond\_style lepton}}}},:doc:\sphinxtitleref{dihedral\_style lepton \textless{}dihedral\_lepton\textgreater{}}


\subsection{Default}
\label{\detokenize{angle_lepton:default}}
\sphinxAtStartPar
none

\sphinxstepscope

\index{angle\_style mesocnt@\spxentry{angle\_style mesocnt}}\ignorespaces 

\section{angle\_style mesocnt command}
\label{\detokenize{angle_mesocnt:angle-style-mesocnt-command}}\label{\detokenize{angle_mesocnt:index-0}}\label{\detokenize{angle_mesocnt::doc}}

\subsection{Syntax}
\label{\detokenize{angle_mesocnt:syntax}}
\begin{sphinxVerbatim}[commandchars=\\\{\}]
\PYG{k}{angle\PYGZus{}style}\PYG{+w}{ }\PYG{n}{mesocnt}
\end{sphinxVerbatim}


\subsection{Examples}
\label{\detokenize{angle_mesocnt:examples}}
\begin{sphinxVerbatim}[commandchars=\\\{\}]
\PYG{k}{angle\PYGZus{}style}\PYG{+w}{ }\PYG{n}{mesocnt}
\PYG{k}{angle\PYGZus{}coeff}\PYG{+w}{ }\PYG{l+m}{1}\PYG{+w}{ }\PYG{n}{buckling}\PYG{+w}{ }\PYG{n}{C}\PYG{+w}{ }\PYG{l+m}{10}\PYG{+w}{ }\PYG{l+m}{10}\PYG{+w}{ }\PYG{l+m}{20.0}
\PYG{k}{angle\PYGZus{}coeff}\PYG{+w}{ }\PYG{l+m}{4}\PYG{+w}{ }\PYG{n}{harmonic}\PYG{+w}{ }\PYG{n}{C}\PYG{+w}{ }\PYG{l+m}{8}\PYG{+w}{ }\PYG{l+m}{4}\PYG{+w}{ }\PYG{l+m}{10.0}
\PYG{k}{angle\PYGZus{}coeff}\PYG{+w}{ }\PYG{l+m}{2}\PYG{+w}{ }\PYG{n}{buckling}\PYG{+w}{ }\PYG{n}{custom}\PYG{+w}{ }\PYG{l+m}{400.0}\PYG{+w}{ }\PYG{l+m}{50.0}\PYG{+w}{ }\PYG{l+m}{5.0}
\PYG{k}{angle\PYGZus{}coeff}\PYG{+w}{ }\PYG{l+m}{1}\PYG{+w}{ }\PYG{n}{harmonic}\PYG{+w}{ }\PYG{n}{custom}\PYG{+w}{ }\PYG{l+m}{300.0}
\end{sphinxVerbatim}


\subsection{Description}
\label{\detokenize{angle_mesocnt:description}}
\sphinxAtStartPar
\DUrole{versionmodified,added}{Added in version 15Sep2022.}

\sphinxAtStartPar
The \sphinxstyleemphasis{mesocnt} angle style uses the potential
\begin{equation*}
\begin{split}E = K_\text{H} \Delta \theta^2, \qquad |\Delta \theta| < \Delta
\theta_\text{B} \\
E = K_\text{H} \Delta \theta_\text{B}^2 +
K_\text{B} (\Delta \theta - \Delta \theta_\text{B}), \qquad |\Delta
\theta| \geq \Delta \theta_\text{B}\end{split}
\end{equation*}
\sphinxAtStartPar
where \(\Delta \theta = \theta - \pi\) is the bending angle of the
nanotube, \(K_\text{H}\) and \(K_\text{B}\) are prefactors for
the harmonic and linear regime respectively and \(\Delta
\theta_\text{B}\) is the buckling angle. Note that the usual 1/2 factor
for the harmonic potential is included in \(K_\text{H}\).

\sphinxAtStartPar
The style implements parameterization presets of \(K_\text{H}\),
\(K_\text{B}\) and \(\Delta \theta_\text{B}\) for mesoscopic
simulations of carbon nanotubes based on the atomistic simulations of
{\hyperref[\detokenize{angle_mesocnt:srivastava-2}]{\sphinxcrossref{\DUrole{std,std-ref}{(Srivastava)}}}} and buckling considerations of
{\hyperref[\detokenize{angle_mesocnt:zhigilei1-1}]{\sphinxcrossref{\DUrole{std,std-ref}{(Zhigilei)}}}}.

\sphinxAtStartPar
The following coefficients must be defined for each angle type via the
{\hyperref[\detokenize{angle_coeff::doc}]{\sphinxcrossref{\DUrole{doc}{angle\_coeff}}}} command as in the examples above, or
in the data file or restart files read by the {\hyperref[\detokenize{read_data::doc}]{\sphinxcrossref{\DUrole{doc}{read\_data}}}} or {\hyperref[\detokenize{read_restart::doc}]{\sphinxcrossref{\DUrole{doc}{read\_restart}}}} commands:
\begin{itemize}
\item {} 
\sphinxAtStartPar
mode = \sphinxstyleemphasis{buckling} or \sphinxstyleemphasis{harmonic}

\item {} 
\sphinxAtStartPar
preset = \sphinxstyleemphasis{C} or \sphinxstyleemphasis{custom}

\item {} 
\sphinxAtStartPar
additional parameters depending on preset

\end{itemize}

\sphinxAtStartPar
If mode \sphinxstyleemphasis{harmonic} is chosen, the potential is simply harmonic and
does not switch to the linear term when the buckling angle is
reached. In \sphinxstyleemphasis{buckling} mode, the full piecewise potential is used.

\sphinxAtStartPar
Preset \sphinxstyleemphasis{C} is for carbon nanotubes, and the additional parameters are:
\begin{itemize}
\item {} 
\sphinxAtStartPar
chiral index \(n\) (unitless)

\item {} 
\sphinxAtStartPar
chiral index \(m\) (unitless)

\item {} 
\sphinxAtStartPar
\(r_0\) (distance)

\end{itemize}

\sphinxAtStartPar
Here, \(r_0\) is the equilibrium distance of the bonds included in
the angle, see {\hyperref[\detokenize{bond_mesocnt::doc}]{\sphinxcrossref{\DUrole{doc}{bond\_style mesocnt}}}}.

\sphinxAtStartPar
In harmonic mode with preset \sphinxstyleemphasis{custom}, the additional parameter is:
\begin{itemize}
\item {} 
\sphinxAtStartPar
\(K_\text{H}\) (energy)

\end{itemize}

\sphinxAtStartPar
Hence, this setting is simply a wrapper for {\hyperref[\detokenize{bond_harmonic::doc}]{\sphinxcrossref{\DUrole{doc}{bond\_style harmonic}}}} with an equilibrium angle of 180 degrees.

\sphinxAtStartPar
In harmonic mode with preset \sphinxstyleemphasis{custom}, the additional parameters are:
\begin{itemize}
\item {} 
\sphinxAtStartPar
\(K_\text{H}\) (energy)

\item {} 
\sphinxAtStartPar
\(K_\text{B}\) (energy)

\item {} 
\sphinxAtStartPar
\(\Delta \theta_\text{B}\) (degrees)

\end{itemize}

\sphinxAtStartPar
\(\Delta \theta_\text{B}\) is specified in degrees, but LAMMPS
converts it to radians internally; hence \(K_\text{H}\) is
effectively energy per radian\textasciicircum{}2 and \(K_\text{B}\) is energy per
radian.


\bigskip\hrule\bigskip


\sphinxAtStartPar
In \sphinxstyleemphasis{buckling} mode, this angle style adds the \sphinxstyleemphasis{buckled} property to
all atoms in the simulation, which is an integer flag indicating
whether the bending angle at a given atom has exceeded \(\Delta
\theta_\text{B}\). It can be accessed as an atomic variable, e.g. for
custom dump commands, as \sphinxstyleemphasis{i\_buckled}.

\begin{sphinxadmonition}{note}{Note:}
\sphinxAtStartPar
If the initial state of the simulation contains buckled nanotubes
and {\hyperref[\detokenize{pair_mesocnt::doc}]{\sphinxcrossref{\DUrole{doc}{pair\_style mesocnt}}}} is used, the
\sphinxstyleemphasis{i\_buckled} atomic variable needs to be initialized before the
pair\_style is defined by doing a \sphinxstyleemphasis{run 0} command straight after the
angle\_style command. See below for an example.
\end{sphinxadmonition}

\sphinxAtStartPar
If CNTs are already buckled at the start of the simulation, this
script will correctly initialize \sphinxstyleemphasis{i\_buckled}:

\begin{sphinxVerbatim}[commandchars=\\\{\}]
\PYG{k}{angle\PYGZus{}style}\PYG{+w}{ }\PYG{n}{mesocnt}
\PYG{k}{angle\PYGZus{}coeff}\PYG{+w}{ }\PYG{l+m}{1}\PYG{+w}{ }\PYG{n}{buckling}\PYG{+w}{ }\PYG{n}{C}\PYG{+w}{ }\PYG{l+m}{10}\PYG{+w}{ }\PYG{l+m}{10}\PYG{+w}{ }\PYG{l+m}{20.0}

\PYG{k}{run}\PYG{+w}{ }\PYG{l+m}{0}

\PYG{k}{pair\PYGZus{}style}\PYG{+w}{ }\PYG{n}{mesocnt}\PYG{+w}{ }\PYG{l+m}{60.0}
\PYG{k}{pair\PYGZus{}coeff}\PYG{+w}{ }\PYG{o}{*}\PYG{+w}{ }\PYG{o}{*}\PYG{+w}{ }\PYG{n}{C\PYGZus{}10\PYGZus{}10.mesocnt}\PYG{+w}{ }\PYG{l+m}{1}
\end{sphinxVerbatim}


\subsection{Restrictions}
\label{\detokenize{angle_mesocnt:restrictions}}
\sphinxAtStartPar
This angle style can only be used if LAMMPS was built with the
MOLECULE and MESONT packages.  See the {\hyperref[\detokenize{Build_package::doc}]{\sphinxcrossref{\DUrole{doc}{Build package}}}} doc page for more info.


\subsection{Related commands}
\label{\detokenize{angle_mesocnt:related-commands}}
\sphinxAtStartPar
{\hyperref[\detokenize{angle_coeff::doc}]{\sphinxcrossref{\DUrole{doc}{angle\_coeff}}}}


\subsection{Default}
\label{\detokenize{angle_mesocnt:default}}
\sphinxAtStartPar
none


\bigskip\hrule\bigskip

\phantomsection\label{\detokenize{angle_mesocnt:srivastava-2}}
\sphinxAtStartPar
\sphinxstylestrong{(Srivastava)} Zhigilei, Wei, Srivastava, Phys. Rev. B 71, 165417
(2005).

\phantomsection\label{\detokenize{angle_mesocnt:zhigilei1-1}}
\sphinxAtStartPar
\sphinxstylestrong{(Zhigilei)} Volkov and Zhigilei, ACS Nano 4, 6187 (2010).

\sphinxstepscope

\index{angle\_style mm3@\spxentry{angle\_style mm3}}\ignorespaces 

\section{angle\_style mm3 command}
\label{\detokenize{angle_mm3:angle-style-mm3-command}}\label{\detokenize{angle_mm3:index-0}}\label{\detokenize{angle_mm3::doc}}

\subsection{Syntax}
\label{\detokenize{angle_mm3:syntax}}
\begin{sphinxVerbatim}[commandchars=\\\{\}]
\PYG{k}{angle\PYGZus{}style}\PYG{+w}{ }\PYG{n}{mm3}
\end{sphinxVerbatim}


\subsection{Examples}
\label{\detokenize{angle_mm3:examples}}
\begin{sphinxVerbatim}[commandchars=\\\{\}]
\PYG{k}{angle\PYGZus{}style}\PYG{+w}{ }\PYG{n}{mm3}
\PYG{k}{angle\PYGZus{}coeff}\PYG{+w}{ }\PYG{l+m}{1}\PYG{+w}{ }\PYG{l+m}{100.0}\PYG{+w}{ }\PYG{l+m}{107.0}
\end{sphinxVerbatim}


\subsection{Description}
\label{\detokenize{angle_mm3:description}}
\sphinxAtStartPar
The \sphinxstyleemphasis{mm3} angle style uses the potential that is anharmonic in the angle
as defined in {\hyperref[\detokenize{bond_mm3:mm3-allinger1989}]{\sphinxcrossref{\DUrole{std,std-ref}{(Allinger)}}}}
\begin{equation*}
\begin{split}E = K (\theta - \theta_0)^2 \left[ 1 - 0.014(\theta - \theta_0) + 5.6(10)^{-5} (\theta - \theta_0)^2 - 7.0(10)^{-7} (\theta - \theta_0)^3 + 9(10)^{-10} (\theta - \theta_0)^4 \right]\end{split}
\end{equation*}
\sphinxAtStartPar
where \(\theta_0\) is the equilibrium value of the angle, and
\(K\) is a prefactor. The anharmonic prefactors have units
\(\deg^{-n}\), for example \(-0.014 \deg^{-1}\), \(5.6
\cdot 10^{-5} \deg^{-2}\), …

\sphinxAtStartPar
The following coefficients must be defined for each angle type via the
{\hyperref[\detokenize{angle_coeff::doc}]{\sphinxcrossref{\DUrole{doc}{angle\_coeff}}}} command as in the example above, or in
the data file or restart files read by the {\hyperref[\detokenize{read_data::doc}]{\sphinxcrossref{\DUrole{doc}{read\_data}}}}
or {\hyperref[\detokenize{read_restart::doc}]{\sphinxcrossref{\DUrole{doc}{read\_restart}}}} commands:
\begin{itemize}
\item {} 
\sphinxAtStartPar
\(K\) (energy)

\item {} 
\sphinxAtStartPar
\(\theta_0\) (degrees)

\end{itemize}

\sphinxAtStartPar
\(\theta_0\) is specified in degrees, but LAMMPS converts it to
radians internally; hence \(K\) is effectively energy per
radian\textasciicircum{}2.


\subsection{Restrictions}
\label{\detokenize{angle_mm3:restrictions}}
\sphinxAtStartPar
This angle style can only be used if LAMMPS was built with the
YAFF package.  See the {\hyperref[\detokenize{Build_package::doc}]{\sphinxcrossref{\DUrole{doc}{Build package}}}} doc
page for more info.


\subsection{Related commands}
\label{\detokenize{angle_mm3:related-commands}}
\sphinxAtStartPar
{\hyperref[\detokenize{angle_coeff::doc}]{\sphinxcrossref{\DUrole{doc}{angle\_coeff}}}}


\subsection{Default}
\label{\detokenize{angle_mm3:default}}
\sphinxAtStartPar
none

\sphinxstepscope

\index{angle\_style none@\spxentry{angle\_style none}}\ignorespaces 

\section{angle\_style none command}
\label{\detokenize{angle_none:angle-style-none-command}}\label{\detokenize{angle_none:index-0}}\label{\detokenize{angle_none::doc}}

\subsection{Syntax}
\label{\detokenize{angle_none:syntax}}
\begin{sphinxVerbatim}[commandchars=\\\{\}]
\PYG{k}{angle\PYGZus{}style}\PYG{+w}{ }\PYG{n}{none}
\end{sphinxVerbatim}


\subsection{Examples}
\label{\detokenize{angle_none:examples}}
\begin{sphinxVerbatim}[commandchars=\\\{\}]
\PYG{k}{angle\PYGZus{}style}\PYG{+w}{ }\PYG{n}{none}
\end{sphinxVerbatim}


\subsection{Description}
\label{\detokenize{angle_none:description}}
\sphinxAtStartPar
Using an angle style of none means angle forces and energies are not
computed, even if triplets of angle atoms were listed in the data file
read by the {\hyperref[\detokenize{read_data::doc}]{\sphinxcrossref{\DUrole{doc}{read\_data}}}} command.

\sphinxAtStartPar
See the {\hyperref[\detokenize{angle_zero::doc}]{\sphinxcrossref{\DUrole{doc}{angle\_style zero}}}} command for a way to
calculate angle statistics, but compute no angle interactions.


\subsection{Restrictions}
\label{\detokenize{angle_none:restrictions}}
\sphinxAtStartPar
none


\subsection{Related commands}
\label{\detokenize{angle_none:related-commands}}
\sphinxAtStartPar
{\hyperref[\detokenize{angle_zero::doc}]{\sphinxcrossref{\DUrole{doc}{angle\_style zero}}}}


\subsection{Default}
\label{\detokenize{angle_none:default}}
\sphinxAtStartPar
none

\sphinxstepscope

\index{angle\_style quartic@\spxentry{angle\_style quartic}}\ignorespaces 
\index{angle\_style quartic/omp@\spxentry{angle\_style quartic/omp}}\ignorespaces 

\section{angle\_style quartic command}
\label{\detokenize{angle_quartic:angle-style-quartic-command}}\label{\detokenize{angle_quartic:index-1}}\label{\detokenize{angle_quartic:index-0}}\label{\detokenize{angle_quartic::doc}}
\sphinxAtStartPar
Accelerator Variants: \sphinxstyleemphasis{quartic/omp}


\subsection{Syntax}
\label{\detokenize{angle_quartic:syntax}}
\begin{sphinxVerbatim}[commandchars=\\\{\}]
\PYG{k}{angle\PYGZus{}style}\PYG{+w}{ }\PYG{n}{quartic}
\end{sphinxVerbatim}


\subsection{Examples}
\label{\detokenize{angle_quartic:examples}}
\begin{sphinxVerbatim}[commandchars=\\\{\}]
\PYG{k}{angle\PYGZus{}style}\PYG{+w}{ }\PYG{n}{quartic}
\PYG{k}{angle\PYGZus{}coeff}\PYG{+w}{ }\PYG{l+m}{1}\PYG{+w}{ }\PYG{l+m}{129.1948}\PYG{+w}{ }\PYG{l+m}{56.8726}\PYG{+w}{ }\PYG{o}{\PYGZhy{}}\PYG{l+m}{25.9442}\PYG{+w}{ }\PYG{o}{\PYGZhy{}}\PYG{l+m}{14.2221}
\end{sphinxVerbatim}


\subsection{Description}
\label{\detokenize{angle_quartic:description}}
\sphinxAtStartPar
The \sphinxstyleemphasis{quartic} angle style uses the potential
\begin{equation*}
\begin{split}E = K_2 (\theta - \theta_0)^2 + K_3 (\theta - \theta_0)^3 + K_4 (\theta - \theta_0)^4\end{split}
\end{equation*}
\sphinxAtStartPar
where \(\theta_0\) is the equilibrium value of the angle, and \(K\) is a
prefactor.  Note that the usual 1/2 factor is included in \(K\).

\sphinxAtStartPar
The following coefficients must be defined for each angle type via the
{\hyperref[\detokenize{angle_coeff::doc}]{\sphinxcrossref{\DUrole{doc}{angle\_coeff}}}} command as in the example above, or in
the data file or restart files read by the {\hyperref[\detokenize{read_data::doc}]{\sphinxcrossref{\DUrole{doc}{read\_data}}}}
or {\hyperref[\detokenize{read_restart::doc}]{\sphinxcrossref{\DUrole{doc}{read\_restart}}}} commands:
\begin{itemize}
\item {} 
\sphinxAtStartPar
\(\theta_0\) (degrees)

\item {} 
\sphinxAtStartPar
\(K_2\) (energy)

\item {} 
\sphinxAtStartPar
\(K_3\) (energy)

\item {} 
\sphinxAtStartPar
\(K_4\) (energy)

\end{itemize}

\sphinxAtStartPar
\(\theta_0\) is specified in degrees, but LAMMPS converts it to
radians internally; hence the various \(K\) are effectively energy
per radian\textasciicircum{}2 or radian\textasciicircum{}3 or radian\textasciicircum{}4.


\bigskip\hrule\bigskip


\sphinxAtStartPar
Styles with a \sphinxstyleemphasis{gpu}, \sphinxstyleemphasis{intel}, \sphinxstyleemphasis{kk}, \sphinxstyleemphasis{omp}, or \sphinxstyleemphasis{opt} suffix are
functionally the same as the corresponding style without the suffix.
They have been optimized to run faster, depending on your available
hardware, as discussed on the {\hyperref[\detokenize{Speed_packages::doc}]{\sphinxcrossref{\DUrole{doc}{Accelerator packages}}}}
page.  The accelerated styles take the same arguments and should
produce the same results, except for round\sphinxhyphen{}off and precision issues.

\sphinxAtStartPar
These accelerated styles are part of the GPU, INTEL, KOKKOS,
OPENMP, and OPT packages, respectively.  They are only enabled if
LAMMPS was built with those packages.  See the {\hyperref[\detokenize{Build_package::doc}]{\sphinxcrossref{\DUrole{doc}{Build package}}}} page for more info.

\sphinxAtStartPar
You can specify the accelerated styles explicitly in your input script
by including their suffix, or you can use the {\hyperref[\detokenize{Run_options::doc}]{\sphinxcrossref{\DUrole{doc}{\sphinxhyphen{}suffix command\sphinxhyphen{}line switch}}}} when you invoke LAMMPS, or you can use the
{\hyperref[\detokenize{suffix::doc}]{\sphinxcrossref{\DUrole{doc}{suffix}}}} command in your input script.

\sphinxAtStartPar
See the {\hyperref[\detokenize{Speed_packages::doc}]{\sphinxcrossref{\DUrole{doc}{Accelerator packages}}}} page for more
instructions on how to use the accelerated styles effectively.


\bigskip\hrule\bigskip



\subsection{Restrictions}
\label{\detokenize{angle_quartic:restrictions}}
\sphinxAtStartPar
This angle style can only be used if LAMMPS was built with the
EXTRA\sphinxhyphen{}MOLECULE package.  See the {\hyperref[\detokenize{Build_package::doc}]{\sphinxcrossref{\DUrole{doc}{Build package}}}} doc
page for more info.


\subsection{Related commands}
\label{\detokenize{angle_quartic:related-commands}}
\sphinxAtStartPar
{\hyperref[\detokenize{angle_coeff::doc}]{\sphinxcrossref{\DUrole{doc}{angle\_coeff}}}}


\subsection{Default}
\label{\detokenize{angle_quartic:default}}
\sphinxAtStartPar
none

\sphinxstepscope

\index{angle\_style spica@\spxentry{angle\_style spica}}\ignorespaces 
\index{angle\_style spica/omp@\spxentry{angle\_style spica/omp}}\ignorespaces 
\index{angle\_style spica/kk@\spxentry{angle\_style spica/kk}}\ignorespaces 

\section{angle\_style spica command}
\label{\detokenize{angle_spica:angle-style-spica-command}}\label{\detokenize{angle_spica:index-2}}\label{\detokenize{angle_spica:index-1}}\label{\detokenize{angle_spica:index-0}}\label{\detokenize{angle_spica::doc}}
\sphinxAtStartPar
Accelerator Variants: \sphinxstyleemphasis{spica/omp}, \sphinxstyleemphasis{spica/kk}


\subsection{Syntax}
\label{\detokenize{angle_spica:syntax}}
\begin{sphinxVerbatim}[commandchars=\\\{\}]
\PYG{k}{angle\PYGZus{}style}\PYG{+w}{ }\PYG{n}{spica}

\PYG{k}{angle\PYGZus{}style}\PYG{+w}{ }\PYG{n}{spica}\PYG{o}{/}\PYG{n}{omp}
\end{sphinxVerbatim}


\subsection{Examples}
\label{\detokenize{angle_spica:examples}}
\begin{sphinxVerbatim}[commandchars=\\\{\}]
\PYG{k}{angle\PYGZus{}style}\PYG{+w}{ }\PYG{n}{spica}
\PYG{k}{angle\PYGZus{}coeff}\PYG{+w}{ }\PYG{l+m}{1}\PYG{+w}{ }\PYG{l+m}{300.0}\PYG{+w}{ }\PYG{l+m}{107.0}
\end{sphinxVerbatim}


\subsection{Description}
\label{\detokenize{angle_spica:description}}
\sphinxAtStartPar
The \sphinxstyleemphasis{spica} angle style is a combination of the harmonic angle potential,
\begin{equation*}
\begin{split}E = K (\theta - \theta_0)^2\end{split}
\end{equation*}
\sphinxAtStartPar
where \(\theta_0\) is the equilibrium value of the angle and
\(K\) a prefactor, with the \sphinxstyleemphasis{repulsive} part of the non\sphinxhyphen{}bonded
\sphinxstyleemphasis{lj/spica} pair style between the atoms 1 and 3.  This angle potential is
intended for coarse grained MD simulations with the SPICA (formerly called SDK) parameterization
using the {\hyperref[\detokenize{pair_spica::doc}]{\sphinxcrossref{\DUrole{doc}{pair\_style lj/spica}}}}.  Relative to the
pair\_style \sphinxstyleemphasis{lj/spica}, however, the energy is shifted by
\(\epsilon\), to avoid sudden jumps.  Note that the usual 1/2 factor
is included in \(K\).

\sphinxAtStartPar
The following coefficients must be defined for each angle type via the
{\hyperref[\detokenize{angle_coeff::doc}]{\sphinxcrossref{\DUrole{doc}{angle\_coeff}}}} command as in the example above:
\begin{itemize}
\item {} 
\sphinxAtStartPar
\(K\) (energy)

\item {} 
\sphinxAtStartPar
\(\theta_0\) (degrees)

\end{itemize}

\sphinxAtStartPar
\(\theta_0\) is specified in degrees, but LAMMPS converts it to
radians internally; hence \(K\) is effectively energy per
radian\textasciicircum{}2.

\sphinxAtStartPar
The required \sphinxstyleemphasis{lj/spica} parameters are extracted automatically from the
pair\_style.

\sphinxAtStartPar
Style \sphinxstyleemphasis{sdk}, the original implementation of style \sphinxstyleemphasis{spica}, is available
for backward compatibility.


\bigskip\hrule\bigskip


\sphinxAtStartPar
Styles with a \sphinxstyleemphasis{gpu}, \sphinxstyleemphasis{intel}, \sphinxstyleemphasis{kk}, \sphinxstyleemphasis{omp}, or \sphinxstyleemphasis{opt} suffix are
functionally the same as the corresponding style without the suffix.
They have been optimized to run faster, depending on your available
hardware, as discussed on the {\hyperref[\detokenize{Speed_packages::doc}]{\sphinxcrossref{\DUrole{doc}{Accelerator packages}}}}
page.  The accelerated styles take the same arguments and should
produce the same results, except for round\sphinxhyphen{}off and precision issues.

\sphinxAtStartPar
These accelerated styles are part of the GPU, INTEL, KOKKOS,
OPENMP, and OPT packages, respectively.  They are only enabled if
LAMMPS was built with those packages.  See the {\hyperref[\detokenize{Build_package::doc}]{\sphinxcrossref{\DUrole{doc}{Build package}}}} page for more info.

\sphinxAtStartPar
You can specify the accelerated styles explicitly in your input script
by including their suffix, or you can use the {\hyperref[\detokenize{Run_options::doc}]{\sphinxcrossref{\DUrole{doc}{\sphinxhyphen{}suffix command\sphinxhyphen{}line switch}}}} when you invoke LAMMPS, or you can use the
{\hyperref[\detokenize{suffix::doc}]{\sphinxcrossref{\DUrole{doc}{suffix}}}} command in your input script.

\sphinxAtStartPar
See the {\hyperref[\detokenize{Speed_packages::doc}]{\sphinxcrossref{\DUrole{doc}{Accelerator packages}}}} page for more
instructions on how to use the accelerated styles effectively.


\bigskip\hrule\bigskip



\subsection{Restrictions}
\label{\detokenize{angle_spica:restrictions}}
\sphinxAtStartPar
This angle style can only be used if LAMMPS was built with the
CG\sphinxhyphen{}SPICA package.  See the {\hyperref[\detokenize{Build_package::doc}]{\sphinxcrossref{\DUrole{doc}{Build package}}}} doc
page for more info.


\subsection{Related commands}
\label{\detokenize{angle_spica:related-commands}}
\sphinxAtStartPar
{\hyperref[\detokenize{angle_coeff::doc}]{\sphinxcrossref{\DUrole{doc}{angle\_coeff}}}}, {\hyperref[\detokenize{angle_harmonic::doc}]{\sphinxcrossref{\DUrole{doc}{angle\_style harmonic}}}}, {\hyperref[\detokenize{pair_spica::doc}]{\sphinxcrossref{\DUrole{doc}{pair\_style lj/spica}}}},
{\hyperref[\detokenize{pair_spica::doc}]{\sphinxcrossref{\DUrole{doc}{pair\_style lj/spica/coul/long}}}}


\subsection{Default}
\label{\detokenize{angle_spica:default}}
\sphinxAtStartPar
none

\sphinxstepscope

\index{angle\_style table@\spxentry{angle\_style table}}\ignorespaces 
\index{angle\_style table/omp@\spxentry{angle\_style table/omp}}\ignorespaces 

\section{angle\_style table command}
\label{\detokenize{angle_table:angle-style-table-command}}\label{\detokenize{angle_table:index-1}}\label{\detokenize{angle_table:index-0}}\label{\detokenize{angle_table::doc}}
\sphinxAtStartPar
Accelerator Variants: \sphinxstyleemphasis{table/omp}


\subsection{Syntax}
\label{\detokenize{angle_table:syntax}}
\begin{sphinxVerbatim}[commandchars=\\\{\}]
\PYG{k}{angle\PYGZus{}style}\PYG{+w}{ }\PYG{n}{table}\PYG{+w}{ }\PYG{n}{style}\PYG{+w}{ }\PYG{n}{N}
\end{sphinxVerbatim}
\begin{itemize}
\item {} 
\sphinxAtStartPar
style = \sphinxstyleemphasis{linear} or \sphinxstyleemphasis{spline} = method of interpolation

\item {} 
\sphinxAtStartPar
N = use N values in table

\end{itemize}


\subsection{Examples}
\label{\detokenize{angle_table:examples}}
\begin{sphinxVerbatim}[commandchars=\\\{\}]
\PYG{k}{angle\PYGZus{}style}\PYG{+w}{ }\PYG{n}{table}\PYG{+w}{ }\PYG{n}{linear}\PYG{+w}{ }\PYG{l+m}{1000}
\PYG{k}{angle\PYGZus{}coeff}\PYG{+w}{ }\PYG{l+m}{3}\PYG{+w}{ }\PYG{n}{file.table}\PYG{+w}{ }\PYG{n}{ENTRY1}
\end{sphinxVerbatim}


\subsection{Description}
\label{\detokenize{angle_table:description}}
\sphinxAtStartPar
Style \sphinxstyleemphasis{table} creates interpolation tables of length \sphinxstyleemphasis{N} from angle
potential and derivative values listed in a file(s) as a function of
angle The files are read by the {\hyperref[\detokenize{angle_coeff::doc}]{\sphinxcrossref{\DUrole{doc}{angle\_coeff}}}}
command.

\sphinxAtStartPar
The interpolation tables are created by fitting cubic splines to the
file values and interpolating energy and derivative values at each of
\sphinxstyleemphasis{N} angles.  During a simulation, these tables are used to interpolate
energy and force values on individual atoms as needed.  The
interpolation is done in one of 2 styles: \sphinxstyleemphasis{linear} or \sphinxstyleemphasis{spline}.

\sphinxAtStartPar
For the \sphinxstyleemphasis{linear} style, the angle is used to find 2 surrounding table
values from which an energy or its derivative is computed by linear
interpolation.

\sphinxAtStartPar
For the \sphinxstyleemphasis{spline} style, a cubic spline coefficients are computed and
stored at each of the \sphinxstyleemphasis{N} values in the table.  The angle is used to
find the appropriate set of coefficients which are used to evaluate a
cubic polynomial which computes the energy or derivative.

\sphinxAtStartPar
The following coefficients must be defined for each angle type via the
{\hyperref[\detokenize{angle_coeff::doc}]{\sphinxcrossref{\DUrole{doc}{angle\_coeff}}}} command as in the example above.
\begin{itemize}
\item {} 
\sphinxAtStartPar
filename

\item {} 
\sphinxAtStartPar
keyword

\end{itemize}

\sphinxAtStartPar
The filename specifies a file containing tabulated energy and
derivative values.  The keyword specifies a section of the file.  The
format of this file is described below.


\bigskip\hrule\bigskip


\sphinxAtStartPar
Suitable tables for use with this angle style can be created by LAMMPS
itself from existing angle styles using the {\hyperref[\detokenize{angle_write::doc}]{\sphinxcrossref{\DUrole{doc}{angle\_write}}}} command.  This can be useful to have a template file for
testing the angle style settings and to build a compatible custom file.
Another option to generate tables is the Python code in the
\sphinxcode{\sphinxupquote{tools/tabulate}} folder of the LAMMPS source code distribution.

\sphinxAtStartPar
The format of a tabulated file is as follows (without the
parenthesized comments):

\begin{sphinxVerbatim}[commandchars=\\\{\}]
\PYGZsh{} Angle potential for harmonic (one or more comment or blank lines)

HAM                           (keyword is the first text on line)
N 181 FP 0 0 EQ 90.0          (N, FP, EQ parameters)
                              (blank line)
1 0.0 200.5 2.5               (index, angle, energy, derivative)
2 1.0 198.0 2.5
...
181 180.0 0.0 0.0
\end{sphinxVerbatim}

\sphinxAtStartPar
A section begins with a non\sphinxhyphen{}blank line whose first character is not a
“\#”; blank lines or lines starting with “\#” can be used as comments
between sections.  The first line begins with a keyword which
identifies the section.  The line can contain additional text, but the
initial text must match the argument specified in the
{\hyperref[\detokenize{angle_coeff::doc}]{\sphinxcrossref{\DUrole{doc}{angle\_coeff}}}} command.  The next line lists (in any
order) one or more parameters for the table.  Each parameter is a
keyword followed by one or more numeric values.

\sphinxAtStartPar
The parameter “N” is required and its value is the number of table
entries that follow.  Note that this may be different than the \sphinxstyleemphasis{N}
specified in the {\hyperref[\detokenize{angle_style::doc}]{\sphinxcrossref{\DUrole{doc}{angle\_style table}}}} command.  Let
Ntable = \sphinxstyleemphasis{N} in the angle\_style command, and Nfile = “N” in the
tabulated file.  What LAMMPS does is a preliminary interpolation by
creating splines using the Nfile tabulated values as nodal points.  It
uses these to interpolate as needed to generate energy and derivative
values at Ntable different points.  The resulting tables of length
Ntable are then used as described above, when computing energy and
force for individual angles and their atoms.  This means that if you
want the interpolation tables of length Ntable to match exactly what
is in the tabulated file (with effectively no preliminary
interpolation), you should set Ntable = Nfile.

\sphinxAtStartPar
The “FP” parameter is optional.  If used, it is followed by two values
fplo and fphi, which are the second derivatives at the innermost and
outermost angle settings.  These values are needed by the spline
construction routines.  If not specified by the “FP” parameter, they
are estimated (less accurately) by the first two and last two
derivative values in the table.

\sphinxAtStartPar
The “EQ” parameter is also optional.  If used, it is followed by a the
equilibrium angle value, which is used, for example, by the {\hyperref[\detokenize{fix_shake::doc}]{\sphinxcrossref{\DUrole{doc}{fix shake}}}} command.  If not used, the equilibrium angle is
set to 180.0.

\sphinxAtStartPar
Following a blank line, the next N lines list the tabulated values.
On each line, the first value is the index from 1 to N, the second value is
the angle value (in degrees), the third value is the energy (in energy
units), and the fourth is \sphinxhyphen{}dE/d(theta) (also in energy units).  The third
term is the energy of the 3\sphinxhyphen{}atom configuration for the specified
angle.  The last term is the derivative of the energy with respect to
the angle (in degrees, not radians).  Thus the units of the last term
are still energy, not force.  The angle values must increase from one
line to the next.  The angle values must also begin with 0.0 and end
with 180.0, i.e. span the full range of possible angles.

\sphinxAtStartPar
Note that one file can contain many sections, each with a tabulated
potential.  LAMMPS reads the file section by section until it finds
one that matches the specified keyword.


\bigskip\hrule\bigskip


\sphinxAtStartPar
Styles with a \sphinxstyleemphasis{gpu}, \sphinxstyleemphasis{intel}, \sphinxstyleemphasis{kk}, \sphinxstyleemphasis{omp}, or \sphinxstyleemphasis{opt} suffix are
functionally the same as the corresponding style without the suffix.
They have been optimized to run faster, depending on your available
hardware, as discussed on the {\hyperref[\detokenize{Speed_packages::doc}]{\sphinxcrossref{\DUrole{doc}{Accelerator packages}}}}
page.  The accelerated styles take the same arguments and should
produce the same results, except for round\sphinxhyphen{}off and precision issues.

\sphinxAtStartPar
These accelerated styles are part of the GPU, INTEL, KOKKOS,
OPENMP, and OPT packages, respectively.  They are only enabled if
LAMMPS was built with those packages.  See the {\hyperref[\detokenize{Build_package::doc}]{\sphinxcrossref{\DUrole{doc}{Build package}}}} page for more info.

\sphinxAtStartPar
You can specify the accelerated styles explicitly in your input script
by including their suffix, or you can use the {\hyperref[\detokenize{Run_options::doc}]{\sphinxcrossref{\DUrole{doc}{\sphinxhyphen{}suffix command\sphinxhyphen{}line switch}}}} when you invoke LAMMPS, or you can use the
{\hyperref[\detokenize{suffix::doc}]{\sphinxcrossref{\DUrole{doc}{suffix}}}} command in your input script.

\sphinxAtStartPar
See the {\hyperref[\detokenize{Speed_packages::doc}]{\sphinxcrossref{\DUrole{doc}{Accelerator packages}}}} page for more
instructions on how to use the accelerated styles effectively.


\bigskip\hrule\bigskip



\subsection{Restart, fix\_modify, output, run start/stop, minimize info}
\label{\detokenize{angle_table:restart-fix-modify-output-run-start-stop-minimize-info}}
\sphinxAtStartPar
This angle style writes the settings for the “angle\_style table”
command to {\hyperref[\detokenize{restart::doc}]{\sphinxcrossref{\DUrole{doc}{binary restart files}}}}, so a angle\_style
command does not need to specified in an input script that reads a
restart file.  However, the coefficient information is not stored in
the restart file, since it is tabulated in the potential files.  Thus,
angle\_coeff commands do need to be specified in the restart input
script.


\subsection{Restrictions}
\label{\detokenize{angle_table:restrictions}}
\sphinxAtStartPar
This angle style can only be used if LAMMPS was built with the
MOLECULE package.  See the {\hyperref[\detokenize{Build_package::doc}]{\sphinxcrossref{\DUrole{doc}{Build package}}}} doc page
for more info.


\subsection{Related commands}
\label{\detokenize{angle_table:related-commands}}
\sphinxAtStartPar
{\hyperref[\detokenize{angle_coeff::doc}]{\sphinxcrossref{\DUrole{doc}{angle\_coeff}}}}, {\hyperref[\detokenize{angle_write::doc}]{\sphinxcrossref{\DUrole{doc}{angle\_write}}}}


\subsection{Default}
\label{\detokenize{angle_table:default}}
\sphinxAtStartPar
none

\sphinxstepscope

\index{angle\_style zero@\spxentry{angle\_style zero}}\ignorespaces 

\section{angle\_style zero command}
\label{\detokenize{angle_zero:angle-style-zero-command}}\label{\detokenize{angle_zero:index-0}}\label{\detokenize{angle_zero::doc}}

\subsection{Syntax}
\label{\detokenize{angle_zero:syntax}}
\begin{sphinxVerbatim}[commandchars=\\\{\}]
\PYG{k}{angle\PYGZus{}style}\PYG{+w}{ }\PYG{n}{zero}\PYG{+w}{ }\PYG{n}{keyword}
\end{sphinxVerbatim}
\begin{itemize}
\item {} 
\sphinxAtStartPar
zero or more keywords may be appended

\item {} 
\sphinxAtStartPar
keyword = \sphinxstyleemphasis{nocoeff}

\end{itemize}


\subsection{Examples}
\label{\detokenize{angle_zero:examples}}
\begin{sphinxVerbatim}[commandchars=\\\{\}]
\PYG{k}{angle\PYGZus{}style}\PYG{+w}{ }\PYG{n}{zero}
\PYG{k}{angle\PYGZus{}style}\PYG{+w}{ }\PYG{n}{zero}\PYG{+w}{ }\PYG{n}{nocoeff}
\PYG{k}{angle\PYGZus{}coeff}\PYG{+w}{ }\PYG{o}{*}
\PYG{k}{angle\PYGZus{}coeff}\PYG{+w}{ }\PYG{o}{*}\PYG{+w}{ }\PYG{l+m}{120.0}
\end{sphinxVerbatim}


\subsection{Description}
\label{\detokenize{angle_zero:description}}
\sphinxAtStartPar
Using an angle style of zero means angle forces and energies are not
computed, but the geometry of angle triplets is still accessible to
other commands.

\sphinxAtStartPar
As an example, the {\hyperref[\detokenize{compute_angle_local::doc}]{\sphinxcrossref{\DUrole{doc}{compute angle/local}}}}
command can be used to compute the theta values for the list of
triplets of angle atoms listed in the data file read by the
{\hyperref[\detokenize{read_data::doc}]{\sphinxcrossref{\DUrole{doc}{read\_data}}}} command.  If no angle style is defined,
this command cannot be used.

\sphinxAtStartPar
The optional \sphinxstyleemphasis{nocoeff} flag allows to read data files with AngleCoeff
section for any angle style. Similarly, any {\hyperref[\detokenize{angle_coeff::doc}]{\sphinxcrossref{\DUrole{doc}{angle\_coeff}}}} commands
will only be checked for the angle type number and the rest ignored.

\sphinxAtStartPar
Note that the {\hyperref[\detokenize{angle_coeff::doc}]{\sphinxcrossref{\DUrole{doc}{angle\_coeff}}}} command must be used for
all angle types. If specified, there can be only one value, which is
going to be used to assign an equilibrium angle, e.g. for use with
{\hyperref[\detokenize{fix_shake::doc}]{\sphinxcrossref{\DUrole{doc}{fix shake}}}}.


\subsection{Restrictions}
\label{\detokenize{angle_zero:restrictions}}\begin{quote}

\sphinxAtStartPar
none
\end{quote}


\subsection{Related commands}
\label{\detokenize{angle_zero:related-commands}}
\sphinxAtStartPar
{\hyperref[\detokenize{angle_none::doc}]{\sphinxcrossref{\DUrole{doc}{angle\_style none}}}}


\subsection{Default}
\label{\detokenize{angle_zero:default}}
\sphinxAtStartPar
none

\sphinxstepscope


\chapter{Dihedral Styles}
\label{\detokenize{dihedrals:dihedral-styles}}\label{\detokenize{dihedrals::doc}}
\sphinxstepscope

\index{dihedral\_style charmm@\spxentry{dihedral\_style charmm}}\ignorespaces 
\index{dihedral\_style charmm/intel@\spxentry{dihedral\_style charmm/intel}}\ignorespaces 
\index{dihedral\_style charmm/kk@\spxentry{dihedral\_style charmm/kk}}\ignorespaces 
\index{dihedral\_style charmm/omp@\spxentry{dihedral\_style charmm/omp}}\ignorespaces 
\index{dihedral\_style charmmfsw@\spxentry{dihedral\_style charmmfsw}}\ignorespaces 
\index{dihedral\_style charmmfsw/kk@\spxentry{dihedral\_style charmmfsw/kk}}\ignorespaces 

\section{dihedral\_style charmm command}
\label{\detokenize{dihedral_charmm:dihedral-style-charmm-command}}\label{\detokenize{dihedral_charmm:index-5}}\label{\detokenize{dihedral_charmm:index-4}}\label{\detokenize{dihedral_charmm:index-3}}\label{\detokenize{dihedral_charmm:index-2}}\label{\detokenize{dihedral_charmm:index-1}}\label{\detokenize{dihedral_charmm:index-0}}\label{\detokenize{dihedral_charmm::doc}}
\sphinxAtStartPar
Accelerator Variants: \sphinxstyleemphasis{charmm/intel}, \sphinxstyleemphasis{charmm/kk}, \sphinxstyleemphasis{charmm/omp}


\section{dihedral\_style charmmfsw command}
\label{\detokenize{dihedral_charmm:dihedral-style-charmmfsw-command}}
\sphinxAtStartPar
Accelerator Variants: \sphinxstyleemphasis{charmmfsw/kk}


\subsection{Syntax}
\label{\detokenize{dihedral_charmm:syntax}}
\begin{sphinxVerbatim}[commandchars=\\\{\}]
\PYG{k}{dihedral\PYGZus{}style}\PYG{+w}{ }\PYG{n}{style}
\end{sphinxVerbatim}
\begin{itemize}
\item {} 
\sphinxAtStartPar
style = \sphinxstyleemphasis{charmm} or \sphinxstyleemphasis{charmmfsw}

\end{itemize}


\subsection{Examples}
\label{\detokenize{dihedral_charmm:examples}}
\begin{sphinxVerbatim}[commandchars=\\\{\}]
\PYG{k}{dihedral\PYGZus{}style}\PYG{+w}{ }\PYG{n}{charmm}
\PYG{k}{dihedral\PYGZus{}style}\PYG{+w}{ }\PYG{n}{charmmfsw}
\PYG{k}{dihedral\PYGZus{}coeff}\PYG{+w}{  }\PYG{l+m}{1}\PYG{+w}{ }\PYG{l+m}{0.2}\PYG{+w}{ }\PYG{l+m}{1}\PYG{+w}{ }\PYG{l+m}{180}\PYG{+w}{ }\PYG{l+m}{1.0}
\PYG{k}{dihedral\PYGZus{}coeff}\PYG{+w}{  }\PYG{l+m}{2}\PYG{+w}{ }\PYG{l+m}{1.8}\PYG{+w}{ }\PYG{l+m}{1}\PYG{+w}{   }\PYG{l+m}{0}\PYG{+w}{ }\PYG{l+m}{1.0}
\PYG{k}{dihedral\PYGZus{}coeff}\PYG{+w}{  }\PYG{l+m}{1}\PYG{+w}{ }\PYG{l+m}{3.1}\PYG{+w}{ }\PYG{l+m}{2}\PYG{+w}{ }\PYG{l+m}{180}\PYG{+w}{ }\PYG{l+m}{0.5}
\end{sphinxVerbatim}


\subsection{Description}
\label{\detokenize{dihedral_charmm:description}}
\sphinxAtStartPar
The \sphinxstyleemphasis{charmm} and \sphinxstyleemphasis{charmmfsw} dihedral styles use the potential
\begin{equation*}
\begin{split}E = K [ 1 + \cos (n \phi - d) ]\end{split}
\end{equation*}
\sphinxAtStartPar
See {\hyperref[\detokenize{dihedral_charmm:dihedral-mackerell}]{\sphinxcrossref{\DUrole{std,std-ref}{(MacKerell)}}}} for a description of the CHARMM
force field.  This dihedral style can also be used for the AMBER force
field (see comment on weighting factors below).  See
{\hyperref[\detokenize{dihedral_charmm:dihedral-cornell}]{\sphinxcrossref{\DUrole{std,std-ref}{(Cornell)}}}} for a description of the AMBER force
field.

\begin{sphinxadmonition}{note}{Note:}
\sphinxAtStartPar
The newer \sphinxstyleemphasis{charmmfsw} style was released in March 2017.  We
recommend it be used instead of the older \sphinxstyleemphasis{charmm} style when running
a simulation with the CHARMM force field, either with long\sphinxhyphen{}range
Coulombics or a Coulombic cutoff, via the {\hyperref[\detokenize{pair_charmm::doc}]{\sphinxcrossref{\DUrole{doc}{pair\_style lj/charmmfsw/coul/long}}}} and {\hyperref[\detokenize{pair_charmm::doc}]{\sphinxcrossref{\DUrole{doc}{pair\_style lj/charmmfsw/coul/charmmfsh}}}} commands respectively.
Otherwise the older \sphinxstyleemphasis{charmm} style is fine to use.  See the discussion
below and more details on the {\hyperref[\detokenize{pair_charmm::doc}]{\sphinxcrossref{\DUrole{doc}{pair\_style charmm}}}} doc
page.
\end{sphinxadmonition}

\sphinxAtStartPar
The following coefficients must be defined for each dihedral type via the
{\hyperref[\detokenize{dihedral_coeff::doc}]{\sphinxcrossref{\DUrole{doc}{dihedral\_coeff}}}} command as in the example above, or in
the data file or restart files read by the {\hyperref[\detokenize{read_data::doc}]{\sphinxcrossref{\DUrole{doc}{read\_data}}}}
or {\hyperref[\detokenize{read_restart::doc}]{\sphinxcrossref{\DUrole{doc}{read\_restart}}}} commands:
\begin{itemize}
\item {} 
\sphinxAtStartPar
\(K\) (energy)

\item {} 
\sphinxAtStartPar
\(n\) (integer \textgreater{}= 0)

\item {} 
\sphinxAtStartPar
\(d\) (integer value of degrees)

\item {} 
\sphinxAtStartPar
weighting factor (1.0, 0.5, or 0.0)

\end{itemize}

\sphinxAtStartPar
The weighting factor is required to correct for double counting
pairwise non\sphinxhyphen{}bonded Lennard\sphinxhyphen{}Jones interactions in cyclic systems or
when using the CHARMM dihedral style with non\sphinxhyphen{}CHARMM force fields.
With the CHARMM dihedral style, interactions between the first and fourth
atoms in a dihedral are skipped during the normal non\sphinxhyphen{}bonded force
computation and instead evaluated as part of the dihedral using
special epsilon and sigma values specified with the
{\hyperref[\detokenize{pair_charmm::doc}]{\sphinxcrossref{\DUrole{doc}{pair\_coeff}}}} command of pair styles that contain
“lj/charmm” (e.g. {\hyperref[\detokenize{pair_charmm::doc}]{\sphinxcrossref{\DUrole{doc}{pair\_style lj/charmm/coul/long}}}})
In 6\sphinxhyphen{}membered rings, the same 1\sphinxhyphen{}4 interaction would be computed twice
(once for the clockwise 1\sphinxhyphen{}4 pair in dihedral 1\sphinxhyphen{}2\sphinxhyphen{}3\sphinxhyphen{}4 and once in the
counterclockwise dihedral 1\sphinxhyphen{}6\sphinxhyphen{}5\sphinxhyphen{}4) and thus the weighting factor has
to be 0.5 in this case.  In 4\sphinxhyphen{}membered or 5\sphinxhyphen{}membered rings, the 1\sphinxhyphen{}4
dihedral also is counted as a 1\sphinxhyphen{}2 or 1\sphinxhyphen{}3 interaction when going around
the ring in the opposite direction and thus the weighting factor is
0.0, as the 1\sphinxhyphen{}2 and 1\sphinxhyphen{}3 exclusions take precedence.

\sphinxAtStartPar
Note that this dihedral weighting factor is unrelated to the scaling
factor specified by the {\hyperref[\detokenize{special_bonds::doc}]{\sphinxcrossref{\DUrole{doc}{special bonds}}}} command
which applies to all 1\sphinxhyphen{}4 interactions in the system.  For CHARMM force
fields, the special\_bonds 1\sphinxhyphen{}4 interaction scaling factor should be set
to 0.0. Since the corresponding 1\sphinxhyphen{}4 non\sphinxhyphen{}bonded interactions are
computed with the dihedral.  This means that if any of the weighting
factors defined as dihedral coefficients (fourth coeff above) are
non\sphinxhyphen{}zero, then you must use a pair style with “lj/charmm” and set the
special\_bonds 1\sphinxhyphen{}4 scaling factor to 0.0 (which is the
default). Otherwise 1\sphinxhyphen{}4 non\sphinxhyphen{}bonded interactions in dihedrals will be
computed twice.

\sphinxAtStartPar
For simulations using the CHARMM force field with a Coulombic cutoff,
the difference between the \sphinxstyleemphasis{charmm} and \sphinxstyleemphasis{charmmfsw} styles is in the
computation of the 1\sphinxhyphen{}4 non\sphinxhyphen{}bond interactions, though only if the
distance between the two atoms is within the switching region of the
pairwise potential defined by the corresponding CHARMM pair style,
i.e. within the outer cutoff specified for the pair style.  The
\sphinxstyleemphasis{charmmfsw} style should only be used when using the corresponding
{\hyperref[\detokenize{pair_charmm::doc}]{\sphinxcrossref{\DUrole{doc}{pair\_style lj/charmmfsw/coul/charmmfsw}}}} or
{\hyperref[\detokenize{pair_charmm::doc}]{\sphinxcrossref{\DUrole{doc}{pair\_style lj/charmmfsw/coul/long}}}} commands.  Use
the \sphinxstyleemphasis{charmm} style with the older {\hyperref[\detokenize{pair_charmm::doc}]{\sphinxcrossref{\DUrole{doc}{pair\_style}}}}
commands that have just “charmm” in their style name.  See the
discussion on the {\hyperref[\detokenize{pair_charmm::doc}]{\sphinxcrossref{\DUrole{doc}{CHARMM pair\_style}}}} page for
details.

\sphinxAtStartPar
Note that for AMBER force fields, which use pair styles with “lj/cut”,
the special\_bonds 1\sphinxhyphen{}4 scaling factor should be set to the AMBER
defaults (1/2 and 5/6) and all the dihedral weighting factors (fourth
coeff above) must be set to 0.0. In this case, you can use any pair
style you wish, since the dihedral does not need any Lennard\sphinxhyphen{}Jones
parameter information and will not compute any 1\sphinxhyphen{}4 non\sphinxhyphen{}bonded
interactions.  Likewise the \sphinxstyleemphasis{charmm} or \sphinxstyleemphasis{charmmfsw} styles are
identical in this case since no 1\sphinxhyphen{}4 non\sphinxhyphen{}bonded interactions are
computed.


\bigskip\hrule\bigskip


\sphinxAtStartPar
Styles with a \sphinxstyleemphasis{gpu}, \sphinxstyleemphasis{intel}, \sphinxstyleemphasis{kk}, \sphinxstyleemphasis{omp}, or \sphinxstyleemphasis{opt} suffix are
functionally the same as the corresponding style without the suffix.
They have been optimized to run faster, depending on your available
hardware, as discussed on the {\hyperref[\detokenize{Speed_packages::doc}]{\sphinxcrossref{\DUrole{doc}{Accelerator packages}}}}
page.  The accelerated styles take the same arguments and should
produce the same results, except for round\sphinxhyphen{}off and precision issues.

\sphinxAtStartPar
These accelerated styles are part of the GPU, INTEL, KOKKOS,
OPENMP, and OPT packages, respectively.  They are only enabled if
LAMMPS was built with those packages.  See the {\hyperref[\detokenize{Build_package::doc}]{\sphinxcrossref{\DUrole{doc}{Build package}}}} page for more info.

\sphinxAtStartPar
You can specify the accelerated styles explicitly in your input script
by including their suffix, or you can use the {\hyperref[\detokenize{Run_options::doc}]{\sphinxcrossref{\DUrole{doc}{\sphinxhyphen{}suffix command\sphinxhyphen{}line switch}}}} when you invoke LAMMPS, or you can use the
{\hyperref[\detokenize{suffix::doc}]{\sphinxcrossref{\DUrole{doc}{suffix}}}} command in your input script.

\sphinxAtStartPar
See the {\hyperref[\detokenize{Speed_packages::doc}]{\sphinxcrossref{\DUrole{doc}{Accelerator packages}}}} page for more
instructions on how to use the accelerated styles effectively.


\bigskip\hrule\bigskip



\subsection{Restrictions}
\label{\detokenize{dihedral_charmm:restrictions}}
\sphinxAtStartPar
When using run\_style {\hyperref[\detokenize{run_style::doc}]{\sphinxcrossref{\DUrole{doc}{respa}}}}, these dihedral styles
must be assigned to the same r\sphinxhyphen{}RESPA level as \sphinxstyleemphasis{pair} or \sphinxstyleemphasis{outer}.

\sphinxAtStartPar
When used in combination with CHARMM pair styles, the 1\sphinxhyphen{}4
{\hyperref[\detokenize{special_bonds::doc}]{\sphinxcrossref{\DUrole{doc}{special\_bonds}}}} scaling factors must be set to 0.0.
Otherwise non\sphinxhyphen{}bonded contributions for these 1\sphinxhyphen{}4 pairs will be
computed multiple times.

\sphinxAtStartPar
These dihedral styles can only be used if LAMMPS was built with the
MOLECULE package.  See the {\hyperref[\detokenize{Build_package::doc}]{\sphinxcrossref{\DUrole{doc}{Build package}}}} doc page
for more info.


\subsection{Related commands}
\label{\detokenize{dihedral_charmm:related-commands}}
\sphinxAtStartPar
{\hyperref[\detokenize{dihedral_coeff::doc}]{\sphinxcrossref{\DUrole{doc}{dihedral\_coeff}}}},
{\hyperref[\detokenize{pair_charmm::doc}]{\sphinxcrossref{\DUrole{doc}{pair\_style lj/charmm variants}}}},
{\hyperref[\detokenize{angle_charmm::doc}]{\sphinxcrossref{\DUrole{doc}{angle\_style charmm}}}}, {\hyperref[\detokenize{fix_cmap::doc}]{\sphinxcrossref{\DUrole{doc}{fix cmap}}}}


\subsection{Default}
\label{\detokenize{dihedral_charmm:default}}
\sphinxAtStartPar
none


\bigskip\hrule\bigskip

\phantomsection\label{\detokenize{dihedral_charmm:dihedral-cornell}}
\sphinxAtStartPar
\sphinxstylestrong{(Cornell)} Cornell, Cieplak, Bayly, Gould, Merz, Ferguson,
Spellmeyer, Fox, Caldwell, Kollman, JACS 117, 5179\sphinxhyphen{}5197 (1995).

\phantomsection\label{\detokenize{dihedral_charmm:dihedral-mackerell}}
\sphinxAtStartPar
\sphinxstylestrong{(MacKerell)} MacKerell, Bashford, Bellott, Dunbrack, Evanseck, Field,
Fischer, Gao, Guo, Ha, et al, J Phys Chem B, 102, 3586 (1998).

\sphinxstepscope

\index{dihedral\_style class2@\spxentry{dihedral\_style class2}}\ignorespaces 
\index{dihedral\_style class2/omp@\spxentry{dihedral\_style class2/omp}}\ignorespaces 
\index{dihedral\_style class2/kk@\spxentry{dihedral\_style class2/kk}}\ignorespaces 

\section{dihedral\_style class2 command}
\label{\detokenize{dihedral_class2:dihedral-style-class2-command}}\label{\detokenize{dihedral_class2:index-2}}\label{\detokenize{dihedral_class2:index-1}}\label{\detokenize{dihedral_class2:index-0}}\label{\detokenize{dihedral_class2::doc}}
\sphinxAtStartPar
Accelerator Variants: \sphinxstyleemphasis{class2/omp}, \sphinxstyleemphasis{class2/kk}


\subsection{Syntax}
\label{\detokenize{dihedral_class2:syntax}}
\begin{sphinxVerbatim}[commandchars=\\\{\}]
\PYG{k}{dihedral\PYGZus{}style}\PYG{+w}{ }\PYG{n}{class2}
\end{sphinxVerbatim}


\subsection{Examples}
\label{\detokenize{dihedral_class2:examples}}
\begin{sphinxVerbatim}[commandchars=\\\{\}]
\PYG{k}{dihedral\PYGZus{}style}\PYG{+w}{ }\PYG{n}{class2}
\PYG{k}{dihedral\PYGZus{}coeff}\PYG{+w}{ }\PYG{l+m}{1}\PYG{+w}{ }\PYG{l+m}{100}\PYG{+w}{ }\PYG{l+m}{75}\PYG{+w}{ }\PYG{l+m}{100}\PYG{+w}{ }\PYG{l+m}{70}\PYG{+w}{ }\PYG{l+m}{80}\PYG{+w}{ }\PYG{l+m}{60}
\PYG{k}{dihedral\PYGZus{}coeff}\PYG{+w}{ }\PYG{o}{*}\PYG{+w}{ }\PYG{n}{mbt}\PYG{+w}{ }\PYG{l+m}{3.5945}\PYG{+w}{ }\PYG{l+m}{0.1704}\PYG{+w}{ }\PYG{o}{\PYGZhy{}}\PYG{l+m}{0.5490}\PYG{+w}{ }\PYG{l+m}{1.5228}
\PYG{k}{dihedral\PYGZus{}coeff}\PYG{+w}{ }\PYG{o}{*}\PYG{+w}{ }\PYG{n}{ebt}\PYG{+w}{ }\PYG{l+m}{0.3417}\PYG{+w}{ }\PYG{l+m}{0.3264}\PYG{+w}{ }\PYG{o}{\PYGZhy{}}\PYG{l+m}{0.9036}\PYG{+w}{ }\PYG{l+m}{0.1368}\PYG{+w}{ }\PYG{l+m}{0.0}\PYG{+w}{ }\PYG{o}{\PYGZhy{}}\PYG{l+m}{0.8080}\PYG{+w}{ }\PYG{l+m}{1.0119}\PYG{+w}{ }\PYG{l+m}{1.1010}
\PYG{k}{dihedral\PYGZus{}coeff}\PYG{+w}{ }\PYG{l+m}{2}\PYG{+w}{ }\PYG{n}{at}\PYG{+w}{ }\PYG{l+m}{0.0}\PYG{+w}{ }\PYG{o}{\PYGZhy{}}\PYG{l+m}{0.1850}\PYG{+w}{ }\PYG{o}{\PYGZhy{}}\PYG{l+m}{0.7963}\PYG{+w}{ }\PYG{o}{\PYGZhy{}}\PYG{l+m}{2.0220}\PYG{+w}{ }\PYG{l+m}{0.0}\PYG{+w}{ }\PYG{o}{\PYGZhy{}}\PYG{l+m}{0.3991}\PYG{+w}{ }\PYG{l+m}{110.2453}\PYG{+w}{ }\PYG{l+m}{105.1270}
\PYG{k}{dihedral\PYGZus{}coeff}\PYG{+w}{ }\PYG{o}{*}\PYG{+w}{ }\PYG{n}{aat}\PYG{+w}{ }\PYG{o}{\PYGZhy{}}\PYG{l+m}{13.5271}\PYG{+w}{ }\PYG{l+m}{110.2453}\PYG{+w}{ }\PYG{l+m}{105.1270}
\PYG{k}{dihedral\PYGZus{}coeff}\PYG{+w}{ }\PYG{o}{*}\PYG{+w}{ }\PYG{n}{bb13}\PYG{+w}{ }\PYG{l+m}{0.0}\PYG{+w}{ }\PYG{l+m}{1.0119}\PYG{+w}{ }\PYG{l+m}{1.1010}
\end{sphinxVerbatim}


\subsection{Description}
\label{\detokenize{dihedral_class2:description}}
\sphinxAtStartPar
The \sphinxstyleemphasis{class2} dihedral style uses the potential
\begin{equation*}
\begin{split}E        = & E_d + E_{mbt} + E_{ebt} + E_{at} + E_{aat} + E_{bb13} \\
E_d      = & \sum_{n=1}^{3} K_n [ 1 - \cos (n \phi - \phi_n) ] \\
E_{mbt}  = & (r_{jk} - r_2) [ A_1 \cos (\phi) + A_2 \cos (2\phi) + A_3 \cos (3\phi) ] \\
E_{ebt}  = & (r_{ij} - r_1) [ B_1 \cos (\phi) + B_2 \cos (2\phi) + B_3 \cos (3\phi) ] + \\
           & (r_{kl} - r_3) [ C_1 \cos (\phi) + C_2 \cos (2\phi) + C_3 \cos (3\phi) ] \\
E_{at}   = & (\theta_{ijk} - \theta_1) [ D_1 \cos (\phi) + D_2 \cos (2\phi) + D_3 \cos (3\phi) ] + \\
           & (\theta_{jkl} - \theta_2) [ E_1 \cos (\phi) + E_2 \cos (2\phi) + E_3 \cos (3\phi) ] \\
E_{aat}  = & M (\theta_{ijk} - \theta_1) (\theta_{jkl} - \theta_2) \cos (\phi) \\
E_{bb13} = & N (r_{ij} - r_1) (r_{kl} - r_3)\end{split}
\end{equation*}
\sphinxAtStartPar
where \(E_d\) is the dihedral term, \(E_{mbt}\) is a middle\sphinxhyphen{}bond\sphinxhyphen{}torsion term,
\(E_{ebt}\) is an end\sphinxhyphen{}bond\sphinxhyphen{}torsion term, \(E_{at}\) is an angle\sphinxhyphen{}torsion term, \(E_{aat}\)
is an angle\sphinxhyphen{}angle\sphinxhyphen{}torsion term, and \(E_{bb13}\) is a bond\sphinxhyphen{}bond\sphinxhyphen{}13 term.

\sphinxAtStartPar
\(\theta_1\) and \(\theta_2\) are equilibrium angles and \(r_1\), \(r_2\), and
\(r_3\) are equilibrium bond lengths.

\sphinxAtStartPar
See {\hyperref[\detokenize{dihedral_class2:dihedral-sun}]{\sphinxcrossref{\DUrole{std,std-ref}{(Sun)}}}} for a description of the COMPASS class2 force field.

\sphinxAtStartPar
Coefficients for the \(E_d\), \(E_{mbt}\), \(E_{ebt}\),
\(E_{at}\), \(E_{aat}\), and \(E_{bb13}\) formulas must be
defined for each dihedral type via the {\hyperref[\detokenize{dihedral_coeff::doc}]{\sphinxcrossref{\DUrole{doc}{dihedral\_coeff}}}}
command as in the example above, or in the data file
or restart files read by the {\hyperref[\detokenize{read_data::doc}]{\sphinxcrossref{\DUrole{doc}{read\_data}}}} or
{\hyperref[\detokenize{read_restart::doc}]{\sphinxcrossref{\DUrole{doc}{read\_restart}}}} commands.

\sphinxAtStartPar
These are the 6 coefficients for the \(E_d\) formula:
\begin{itemize}
\item {} 
\sphinxAtStartPar
\(K_1\) (energy)

\item {} 
\sphinxAtStartPar
\(\phi_1\) (degrees)

\item {} 
\sphinxAtStartPar
\(K_2\) (energy)

\item {} 
\sphinxAtStartPar
\(\phi_2\) (degrees)

\item {} 
\sphinxAtStartPar
\(K_3\) (energy)

\item {} 
\sphinxAtStartPar
\(phi_3\) (degrees)

\end{itemize}

\sphinxAtStartPar
For the \(E_{mbt}\) formula, each line in a
{\hyperref[\detokenize{dihedral_coeff::doc}]{\sphinxcrossref{\DUrole{doc}{dihedral\_coeff}}}} command in the input script lists
5 coefficients, the first of which is \sphinxstyleemphasis{mbt} to indicate they are
MiddleBondTorsion coefficients.  In a data file, these coefficients
should be listed under a \sphinxstyleemphasis{MiddleBondTorsion Coeffs} heading and you
must leave out the \sphinxstyleemphasis{mbt}, i.e. only list 4 coefficients after the
dihedral type.
\begin{itemize}
\item {} 
\sphinxAtStartPar
\sphinxstyleemphasis{mbt}

\item {} 
\sphinxAtStartPar
\(A_1\) (energy/distance)

\item {} 
\sphinxAtStartPar
\(A_2\) (energy/distance)

\item {} 
\sphinxAtStartPar
\(A_3\) (energy/distance)

\item {} 
\sphinxAtStartPar
\(r_2\) (distance)

\end{itemize}

\sphinxAtStartPar
For the \(E_{ebt}\) formula, each line in a
{\hyperref[\detokenize{dihedral_coeff::doc}]{\sphinxcrossref{\DUrole{doc}{dihedral\_coeff}}}} command in the input script lists
9 coefficients, the first of which is \sphinxstyleemphasis{ebt} to indicate they are
EndBondTorsion coefficients.  In a data file, these coefficients
should be listed under a \sphinxstyleemphasis{EndBondTorsion Coeffs} heading and you must
leave out the \sphinxstyleemphasis{ebt}, i.e. only list 8 coefficients after the dihedral
type.
\begin{itemize}
\item {} 
\sphinxAtStartPar
\sphinxstyleemphasis{ebt}

\item {} 
\sphinxAtStartPar
\(B_1\) (energy/distance)

\item {} 
\sphinxAtStartPar
\(B_2\) (energy/distance)

\item {} 
\sphinxAtStartPar
\(B_3\) (energy/distance)

\item {} 
\sphinxAtStartPar
\(C_1\) (energy/distance)

\item {} 
\sphinxAtStartPar
\(C_2\) (energy/distance)

\item {} 
\sphinxAtStartPar
\(C_3\) (energy/distance)

\item {} 
\sphinxAtStartPar
\(r_1\) (distance)

\item {} 
\sphinxAtStartPar
\(r_3\) (distance)

\end{itemize}

\sphinxAtStartPar
For the \(E_{at}\) formula, each line in a
{\hyperref[\detokenize{dihedral_coeff::doc}]{\sphinxcrossref{\DUrole{doc}{dihedral\_coeff}}}} command in the input script lists
9 coefficients, the first of which is \sphinxstyleemphasis{at} to indicate they are
AngleTorsion coefficients.  In a data file, these coefficients should
be listed under a \sphinxstyleemphasis{AngleTorsion Coeffs} heading and you must leave out
the \sphinxstyleemphasis{at}, i.e. only list 8 coefficients after the dihedral type.
\begin{itemize}
\item {} 
\sphinxAtStartPar
\sphinxstyleemphasis{at}

\item {} 
\sphinxAtStartPar
\(D_1\) (energy)

\item {} 
\sphinxAtStartPar
\(D_2\) (energy)

\item {} 
\sphinxAtStartPar
\(D_3\) (energy)

\item {} 
\sphinxAtStartPar
\(E_1\) (energy)

\item {} 
\sphinxAtStartPar
\(E_2\) (energy)

\item {} 
\sphinxAtStartPar
\(E_3\) (energy)

\item {} 
\sphinxAtStartPar
\(\theta_1\) (degrees)

\item {} 
\sphinxAtStartPar
\(\theta_2\) (degrees)

\end{itemize}

\sphinxAtStartPar
\(\theta_1\) and \(\theta_2\) are specified in degrees, but
LAMMPS converts them to radians internally; hence the various
\(D\) and \(E\) are effectively energy per radian.

\sphinxAtStartPar
For the \(E_{aat}\) formula, each line in a {\hyperref[\detokenize{dihedral_coeff::doc}]{\sphinxcrossref{\DUrole{doc}{dihedral\_coeff}}}} command in the input script lists 4 coefficients,
the first of which is \sphinxstyleemphasis{aat} to indicate they are AngleAngleTorsion
coefficients.  In a data file, these coefficients should be listed
under a \sphinxstyleemphasis{AngleAngleTorsion Coeffs} heading and you must leave out the
\sphinxstyleemphasis{aat}, i.e. only list 3 coefficients after the dihedral type.
\begin{itemize}
\item {} 
\sphinxAtStartPar
\sphinxstyleemphasis{aat}

\item {} 
\sphinxAtStartPar
\(M\) (energy)

\item {} 
\sphinxAtStartPar
\(\theta_1\) (degrees)

\item {} 
\sphinxAtStartPar
\(\theta_2\) (degrees)

\end{itemize}

\sphinxAtStartPar
\(\theta_1\) and \(\theta_2\) are specified in degrees, but
LAMMPS converts them to radians internally; hence \(M\) is
effectively energy per radian\textasciicircum{}2.

\sphinxAtStartPar
For the \(E_{bb13}\) formula, each line in a
{\hyperref[\detokenize{dihedral_coeff::doc}]{\sphinxcrossref{\DUrole{doc}{dihedral\_coeff}}}} command in the input script lists
4 coefficients, the first of which is \sphinxstyleemphasis{bb13} to indicate they are
BondBond13 coefficients.  In a data file, these coefficients should be
listed under a \sphinxstyleemphasis{BondBond13 Coeffs} heading and you must leave out the
\sphinxstyleemphasis{bb13}, i.e. only list 3 coefficients after the dihedral type.
\begin{itemize}
\item {} 
\sphinxAtStartPar
\sphinxstyleemphasis{bb13}

\item {} 
\sphinxAtStartPar
\(N\) (energy/distance\textasciicircum{}2)

\item {} 
\sphinxAtStartPar
\(r_1\) (distance)

\item {} 
\sphinxAtStartPar
\(r_3\) (distance)

\end{itemize}


\bigskip\hrule\bigskip


\sphinxAtStartPar
Styles with a \sphinxstyleemphasis{gpu}, \sphinxstyleemphasis{intel}, \sphinxstyleemphasis{kk}, \sphinxstyleemphasis{omp}, or \sphinxstyleemphasis{opt} suffix are
functionally the same as the corresponding style without the suffix.
They have been optimized to run faster, depending on your available
hardware, as discussed on the {\hyperref[\detokenize{Speed_packages::doc}]{\sphinxcrossref{\DUrole{doc}{Accelerator packages}}}}
page.  The accelerated styles take the same arguments and should
produce the same results, except for round\sphinxhyphen{}off and precision issues.

\sphinxAtStartPar
These accelerated styles are part of the GPU, INTEL, KOKKOS,
OPENMP, and OPT packages, respectively.  They are only enabled if
LAMMPS was built with those packages.  See the {\hyperref[\detokenize{Build_package::doc}]{\sphinxcrossref{\DUrole{doc}{Build package}}}} page for more info.

\sphinxAtStartPar
You can specify the accelerated styles explicitly in your input script
by including their suffix, or you can use the {\hyperref[\detokenize{Run_options::doc}]{\sphinxcrossref{\DUrole{doc}{\sphinxhyphen{}suffix command\sphinxhyphen{}line switch}}}} when you invoke LAMMPS, or you can use the
{\hyperref[\detokenize{suffix::doc}]{\sphinxcrossref{\DUrole{doc}{suffix}}}} command in your input script.

\sphinxAtStartPar
See the {\hyperref[\detokenize{Speed_packages::doc}]{\sphinxcrossref{\DUrole{doc}{Accelerator packages}}}} page for more
instructions on how to use the accelerated styles effectively.


\bigskip\hrule\bigskip



\subsection{Restrictions}
\label{\detokenize{dihedral_class2:restrictions}}
\sphinxAtStartPar
This dihedral style can only be used if LAMMPS was built with the
CLASS2 package.  See the {\hyperref[\detokenize{Build_package::doc}]{\sphinxcrossref{\DUrole{doc}{Build package}}}} doc
page for more info.


\subsection{Related commands}
\label{\detokenize{dihedral_class2:related-commands}}
\sphinxAtStartPar
{\hyperref[\detokenize{dihedral_coeff::doc}]{\sphinxcrossref{\DUrole{doc}{dihedral\_coeff}}}}


\subsection{Default}
\label{\detokenize{dihedral_class2:default}}
\sphinxAtStartPar
none


\bigskip\hrule\bigskip

\phantomsection\label{\detokenize{dihedral_class2:dihedral-sun}}
\sphinxAtStartPar
\sphinxstylestrong{(Sun)} Sun, J Phys Chem B 102, 7338\sphinxhyphen{}7364 (1998).

\sphinxstepscope

\index{dihedral\_style cosine/shift/exp@\spxentry{dihedral\_style cosine/shift/exp}}\ignorespaces 
\index{dihedral\_style cosine/shift/exp/omp@\spxentry{dihedral\_style cosine/shift/exp/omp}}\ignorespaces 

\section{dihedral\_style cosine/shift/exp command}
\label{\detokenize{dihedral_cosine_shift_exp:dihedral-style-cosine-shift-exp-command}}\label{\detokenize{dihedral_cosine_shift_exp:index-1}}\label{\detokenize{dihedral_cosine_shift_exp:index-0}}\label{\detokenize{dihedral_cosine_shift_exp::doc}}
\sphinxAtStartPar
Accelerator Variants: \sphinxstyleemphasis{cosine/shift/exp/omp}


\subsection{Syntax}
\label{\detokenize{dihedral_cosine_shift_exp:syntax}}
\begin{sphinxVerbatim}[commandchars=\\\{\}]
\PYG{k}{dihedral\PYGZus{}style}\PYG{+w}{ }\PYG{n}{cosine}\PYG{o}{/}\PYG{n}{shift}\PYG{o}{/}\PYG{n}{exp}
\end{sphinxVerbatim}


\subsection{Examples}
\label{\detokenize{dihedral_cosine_shift_exp:examples}}
\begin{sphinxVerbatim}[commandchars=\\\{\}]
\PYG{k}{dihedral\PYGZus{}style}\PYG{+w}{ }\PYG{n}{cosine}\PYG{o}{/}\PYG{n}{shift}\PYG{o}{/}\PYG{n}{exp}
\PYG{k}{dihedral\PYGZus{}coeff}\PYG{+w}{ }\PYG{l+m}{1}\PYG{+w}{ }\PYG{l+m}{10.0}\PYG{+w}{ }\PYG{l+m}{45.0}\PYG{+w}{ }\PYG{l+m}{2.0}
\end{sphinxVerbatim}


\subsection{Description}
\label{\detokenize{dihedral_cosine_shift_exp:description}}
\sphinxAtStartPar
The \sphinxstyleemphasis{cosine/shift/exp} dihedral style uses the potential
\begin{equation*}
\begin{split}E = -U_{min}\frac{e^{-a U(\theta,\theta_0)}-1}{e^a-1} \quad\mbox{with}\quad U(\theta,\theta_0)=-0.5 \left(1+\cos(\theta-\theta_0) \right)\end{split}
\end{equation*}
\sphinxAtStartPar
where \(U_{min}\), \(\theta\), and \(a\) are defined for
each dihedral type.

\sphinxAtStartPar
The potential is bounded between \(\left[-U_{min}:0\right]\) and the minimum is located
at the angle \(\theta_0\). The a parameter can be both positive or negative
and is used to control the spring constant at the equilibrium.

\sphinxAtStartPar
The spring constant is given by \(k=a e^a \frac{U_{min}}{2 \left(e^a-1\right)}\).
For \(a>3\) and  \(\frac{k}{U_{min}} = \frac{a}{2}\) to better than 5\% relative error. For negative
values of the a parameter, the spring constant is essentially zero,
and anharmonic terms takes over. The potential is furthermore well
behaved in the limit \(a \rightarrow 0\), where it has been implemented to linear
order in \(a\) for \(a < 0.001\).

\sphinxAtStartPar
The following coefficients must be defined for each dihedral type via
the {\hyperref[\detokenize{dihedral_coeff::doc}]{\sphinxcrossref{\DUrole{doc}{dihedral\_coeff}}}} command as in the example
above, or in the data file or restart files read by the
{\hyperref[\detokenize{read_data::doc}]{\sphinxcrossref{\DUrole{doc}{read\_data}}}} or {\hyperref[\detokenize{read_restart::doc}]{\sphinxcrossref{\DUrole{doc}{read\_restart}}}}
commands:
\begin{itemize}
\item {} 
\sphinxAtStartPar
\(U_{min}\) (energy)

\item {} 
\sphinxAtStartPar
\(\theta\) (angle)

\item {} 
\sphinxAtStartPar
\(a\) (real number)

\end{itemize}


\bigskip\hrule\bigskip


\sphinxAtStartPar
Styles with a \sphinxstyleemphasis{gpu}, \sphinxstyleemphasis{intel}, \sphinxstyleemphasis{kk}, \sphinxstyleemphasis{omp}, or \sphinxstyleemphasis{opt} suffix are
functionally the same as the corresponding style without the suffix.
They have been optimized to run faster, depending on your available
hardware, as discussed on the {\hyperref[\detokenize{Speed_packages::doc}]{\sphinxcrossref{\DUrole{doc}{Accelerator packages}}}}
page.  The accelerated styles take the same arguments and should
produce the same results, except for round\sphinxhyphen{}off and precision issues.

\sphinxAtStartPar
These accelerated styles are part of the GPU, INTEL, KOKKOS,
OPENMP, and OPT packages, respectively.  They are only enabled if
LAMMPS was built with those packages.  See the {\hyperref[\detokenize{Build_package::doc}]{\sphinxcrossref{\DUrole{doc}{Build package}}}} page for more info.

\sphinxAtStartPar
You can specify the accelerated styles explicitly in your input script
by including their suffix, or you can use the {\hyperref[\detokenize{Run_options::doc}]{\sphinxcrossref{\DUrole{doc}{\sphinxhyphen{}suffix command\sphinxhyphen{}line switch}}}} when you invoke LAMMPS, or you can use the
{\hyperref[\detokenize{suffix::doc}]{\sphinxcrossref{\DUrole{doc}{suffix}}}} command in your input script.

\sphinxAtStartPar
See the {\hyperref[\detokenize{Speed_packages::doc}]{\sphinxcrossref{\DUrole{doc}{Accelerator packages}}}} page for more
instructions on how to use the accelerated styles effectively.


\bigskip\hrule\bigskip



\subsection{Restrictions}
\label{\detokenize{dihedral_cosine_shift_exp:restrictions}}
\sphinxAtStartPar
This dihedral style can only be used if LAMMPS was built with the
MOLECULE package.  See the {\hyperref[\detokenize{Build_package::doc}]{\sphinxcrossref{\DUrole{doc}{Build package}}}} doc
page for more info.


\subsection{Related commands}
\label{\detokenize{dihedral_cosine_shift_exp:related-commands}}
\sphinxAtStartPar
{\hyperref[\detokenize{dihedral_coeff::doc}]{\sphinxcrossref{\DUrole{doc}{dihedral\_coeff}}}},
{\hyperref[\detokenize{angle_cosine_shift_exp::doc}]{\sphinxcrossref{\DUrole{doc}{angle\_style cosine/shift/exp}}}}


\subsection{Default}
\label{\detokenize{dihedral_cosine_shift_exp:default}}
\sphinxAtStartPar
none

\sphinxstepscope

\index{dihedral\_style cosine/squared/restricted@\spxentry{dihedral\_style cosine/squared/restricted}}\ignorespaces 

\section{dihedral\_style cosine/squared/restricted command}
\label{\detokenize{dihedral_cosine_squared_restricted:dihedral-style-cosine-squared-restricted-command}}\label{\detokenize{dihedral_cosine_squared_restricted:index-0}}\label{\detokenize{dihedral_cosine_squared_restricted::doc}}

\subsection{Syntax}
\label{\detokenize{dihedral_cosine_squared_restricted:syntax}}
\begin{sphinxVerbatim}[commandchars=\\\{\}]
\PYG{k}{dihedral\PYGZus{}style}\PYG{+w}{ }\PYG{n}{cosine}\PYG{o}{/}\PYG{n}{squared}\PYG{o}{/}\PYG{n}{restricted}
\end{sphinxVerbatim}


\subsection{Examples}
\label{\detokenize{dihedral_cosine_squared_restricted:examples}}
\begin{sphinxVerbatim}[commandchars=\\\{\}]
\PYG{k}{dihedral\PYGZus{}style}\PYG{+w}{ }\PYG{n}{cosine}\PYG{o}{/}\PYG{n}{squared}\PYG{o}{/}\PYG{n}{restricted}
\PYG{k}{dihedral\PYGZus{}coeff}\PYG{+w}{ }\PYG{l+m}{1}\PYG{+w}{ }\PYG{l+m}{10.0}\PYG{+w}{ }\PYG{l+m}{120}
\end{sphinxVerbatim}


\subsection{Description}
\label{\detokenize{dihedral_cosine_squared_restricted:description}}
\sphinxAtStartPar
\DUrole{versionmodified,added}{Added in version 17Apr2024.}

\sphinxAtStartPar
The \sphinxstyleemphasis{cosine/squared/restricted} dihedral style uses the potential
\begin{equation*}
\begin{split}E = K [\cos(\phi) - \cos(\phi_0)]^2 / \sin^2(\phi)\end{split}
\end{equation*}
\sphinxAtStartPar
, which is commonly used in the MARTINI force field.

\sphinxAtStartPar
See {\hyperref[\detokenize{dihedral_cosine_squared_restricted:restricted-bul}]{\sphinxcrossref{\DUrole{std,std-ref}{(Bulacu)}}}} for a description of the restricted dihedral for the MARTINI force field.

\sphinxAtStartPar
The following coefficients must be defined for each dihedral type via the
{\hyperref[\detokenize{dihedral_coeff::doc}]{\sphinxcrossref{\DUrole{doc}{dihedral\_coeff}}}} command as in the example above, or in
the data file or restart files read by the {\hyperref[\detokenize{read_data::doc}]{\sphinxcrossref{\DUrole{doc}{read\_data}}}}
or {\hyperref[\detokenize{read_restart::doc}]{\sphinxcrossref{\DUrole{doc}{read\_restart}}}} commands:
\begin{itemize}
\item {} 
\sphinxAtStartPar
\(K\) (energy)

\item {} 
\sphinxAtStartPar
\(\phi_0\) (degrees)

\end{itemize}

\sphinxAtStartPar
\(\phi_0\) is specified in degrees, but LAMMPS converts it to radians internally.


\bigskip\hrule\bigskip



\subsection{Restrictions}
\label{\detokenize{dihedral_cosine_squared_restricted:restrictions}}
\sphinxAtStartPar
This dihedral style can only be used if LAMMPS was built with the
EXTRA\sphinxhyphen{}MOLECULE package.  See the {\hyperref[\detokenize{Build_package::doc}]{\sphinxcrossref{\DUrole{doc}{Build package}}}} doc page
for more info.


\subsection{Related commands}
\label{\detokenize{dihedral_cosine_squared_restricted:related-commands}}
\sphinxAtStartPar
{\hyperref[\detokenize{dihedral_coeff::doc}]{\sphinxcrossref{\DUrole{doc}{dihedral\_coeff}}}}


\subsection{Default}
\label{\detokenize{dihedral_cosine_squared_restricted:default}}
\sphinxAtStartPar
none


\bigskip\hrule\bigskip

\phantomsection\label{\detokenize{dihedral_cosine_squared_restricted:restricted-bul}}
\sphinxAtStartPar
\sphinxstylestrong{(Bulacu)} Bulacu, Goga, Zhao, Rossi, Monticelli, Periole, Tieleman, Marrink, J Chem Theory Comput, 9, 3282\sphinxhyphen{}3292
(2013).

\sphinxstepscope

\index{dihedral\_style fourier@\spxentry{dihedral\_style fourier}}\ignorespaces 
\index{dihedral\_style fourier/intel@\spxentry{dihedral\_style fourier/intel}}\ignorespaces 
\index{dihedral\_style fourier/omp@\spxentry{dihedral\_style fourier/omp}}\ignorespaces 

\section{dihedral\_style fourier command}
\label{\detokenize{dihedral_fourier:dihedral-style-fourier-command}}\label{\detokenize{dihedral_fourier:index-2}}\label{\detokenize{dihedral_fourier:index-1}}\label{\detokenize{dihedral_fourier:index-0}}\label{\detokenize{dihedral_fourier::doc}}
\sphinxAtStartPar
Accelerator Variants: \sphinxstyleemphasis{fourier/intel}, \sphinxstyleemphasis{fourier/omp}


\subsection{Syntax}
\label{\detokenize{dihedral_fourier:syntax}}
\begin{sphinxVerbatim}[commandchars=\\\{\}]
\PYG{k}{dihedral\PYGZus{}style}\PYG{+w}{ }\PYG{n}{fourier}
\end{sphinxVerbatim}


\subsection{Examples}
\label{\detokenize{dihedral_fourier:examples}}
\begin{sphinxVerbatim}[commandchars=\\\{\}]
\PYG{k}{dihedral\PYGZus{}style}\PYG{+w}{ }\PYG{n}{fourier}
\PYG{k}{dihedral\PYGZus{}coeff}\PYG{+w}{ }\PYG{l+m}{1}\PYG{+w}{ }\PYG{l+m}{3}\PYG{+w}{ }\PYG{o}{\PYGZhy{}}\PYG{l+m}{0.846200}\PYG{+w}{ }\PYG{l+m}{3}\PYG{+w}{ }\PYG{l+m}{0.0}\PYG{+w}{ }\PYG{l+m}{7.578800}\PYG{+w}{ }\PYG{l+m}{1}\PYG{+w}{ }\PYG{l+m}{0}\PYG{+w}{ }\PYG{l+m}{0.138000}\PYG{+w}{ }\PYG{l+m}{2}\PYG{+w}{ }\PYG{o}{\PYGZhy{}}\PYG{l+m}{180.0}
\end{sphinxVerbatim}


\subsection{Description}
\label{\detokenize{dihedral_fourier:description}}
\sphinxAtStartPar
The \sphinxstyleemphasis{fourier} dihedral style uses the potential:
\begin{equation*}
\begin{split}E = \sum_{i=1,m} K_i  [ 1.0 + \cos ( n_i \phi - d_i ) ]\end{split}
\end{equation*}
\sphinxAtStartPar
The following coefficients must be defined for each dihedral type via the
{\hyperref[\detokenize{dihedral_coeff::doc}]{\sphinxcrossref{\DUrole{doc}{dihedral\_coeff}}}} command as in the example above, or in
the data file or restart files read by the {\hyperref[\detokenize{read_data::doc}]{\sphinxcrossref{\DUrole{doc}{read\_data}}}}
or {\hyperref[\detokenize{read_restart::doc}]{\sphinxcrossref{\DUrole{doc}{read\_restart}}}} commands:
\begin{itemize}
\item {} 
\sphinxAtStartPar
\(m\) (integer \textgreater{}=1)

\item {} 
\sphinxAtStartPar
\(K_1\) (energy)

\item {} 
\sphinxAtStartPar
\(n_1\) (integer \textgreater{}= 0)

\item {} 
\sphinxAtStartPar
\(d_1\) (degrees)

\item {} 
\sphinxAtStartPar
{[}…{]}

\item {} 
\sphinxAtStartPar
\(K_m\) (energy)

\item {} 
\sphinxAtStartPar
\(n_m\) (integer \textgreater{}= 0)

\item {} 
\sphinxAtStartPar
\(d_m\) (degrees)

\end{itemize}


\bigskip\hrule\bigskip


\sphinxAtStartPar
Styles with a \sphinxstyleemphasis{gpu}, \sphinxstyleemphasis{intel}, \sphinxstyleemphasis{kk}, \sphinxstyleemphasis{omp}, or \sphinxstyleemphasis{opt} suffix are
functionally the same as the corresponding style without the suffix.
They have been optimized to run faster, depending on your available
hardware, as discussed on the {\hyperref[\detokenize{Speed_packages::doc}]{\sphinxcrossref{\DUrole{doc}{Accelerator packages}}}}
page.  The accelerated styles take the same arguments and should
produce the same results, except for round\sphinxhyphen{}off and precision issues.

\sphinxAtStartPar
These accelerated styles are part of the GPU, INTEL, KOKKOS,
OPENMP, and OPT packages, respectively.  They are only enabled if
LAMMPS was built with those packages.  See the {\hyperref[\detokenize{Build_package::doc}]{\sphinxcrossref{\DUrole{doc}{Build package}}}} page for more info.

\sphinxAtStartPar
You can specify the accelerated styles explicitly in your input script
by including their suffix, or you can use the {\hyperref[\detokenize{Run_options::doc}]{\sphinxcrossref{\DUrole{doc}{\sphinxhyphen{}suffix command\sphinxhyphen{}line switch}}}} when you invoke LAMMPS, or you can use the
{\hyperref[\detokenize{suffix::doc}]{\sphinxcrossref{\DUrole{doc}{suffix}}}} command in your input script.

\sphinxAtStartPar
See the {\hyperref[\detokenize{Speed_packages::doc}]{\sphinxcrossref{\DUrole{doc}{Accelerator packages}}}} page for more
instructions on how to use the accelerated styles effectively.


\bigskip\hrule\bigskip



\subsection{Restrictions}
\label{\detokenize{dihedral_fourier:restrictions}}
\sphinxAtStartPar
This dihedral style can only be used if LAMMPS was built with the
EXTRA\sphinxhyphen{}MOLECULE package.  See the {\hyperref[\detokenize{Build_package::doc}]{\sphinxcrossref{\DUrole{doc}{Build package}}}} doc
page for more info.


\subsection{Related commands}
\label{\detokenize{dihedral_fourier:related-commands}}
\sphinxAtStartPar
{\hyperref[\detokenize{dihedral_coeff::doc}]{\sphinxcrossref{\DUrole{doc}{dihedral\_coeff}}}}


\subsection{Default}
\label{\detokenize{dihedral_fourier:default}}
\sphinxAtStartPar
none

\sphinxstepscope

\index{dihedral\_style harmonic@\spxentry{dihedral\_style harmonic}}\ignorespaces 
\index{dihedral\_style harmonic/intel@\spxentry{dihedral\_style harmonic/intel}}\ignorespaces 
\index{dihedral\_style harmonic/kk@\spxentry{dihedral\_style harmonic/kk}}\ignorespaces 
\index{dihedral\_style harmonic/omp@\spxentry{dihedral\_style harmonic/omp}}\ignorespaces 

\section{dihedral\_style harmonic command}
\label{\detokenize{dihedral_harmonic:dihedral-style-harmonic-command}}\label{\detokenize{dihedral_harmonic:index-3}}\label{\detokenize{dihedral_harmonic:index-2}}\label{\detokenize{dihedral_harmonic:index-1}}\label{\detokenize{dihedral_harmonic:index-0}}\label{\detokenize{dihedral_harmonic::doc}}
\sphinxAtStartPar
Accelerator Variants: \sphinxstyleemphasis{harmonic/intel}, \sphinxstyleemphasis{harmonic/kk}, \sphinxstyleemphasis{harmonic/omp}


\subsection{Syntax}
\label{\detokenize{dihedral_harmonic:syntax}}
\begin{sphinxVerbatim}[commandchars=\\\{\}]
\PYG{k}{dihedral\PYGZus{}style}\PYG{+w}{ }\PYG{n}{harmonic}
\end{sphinxVerbatim}


\subsection{Examples}
\label{\detokenize{dihedral_harmonic:examples}}
\begin{sphinxVerbatim}[commandchars=\\\{\}]
\PYG{k}{dihedral\PYGZus{}style}\PYG{+w}{ }\PYG{n}{harmonic}
\PYG{k}{dihedral\PYGZus{}coeff}\PYG{+w}{ }\PYG{l+m}{1}\PYG{+w}{ }\PYG{l+m}{80.0}\PYG{+w}{ }\PYG{l+m}{1}\PYG{+w}{ }\PYG{l+m}{2}
\end{sphinxVerbatim}


\subsection{Description}
\label{\detokenize{dihedral_harmonic:description}}
\sphinxAtStartPar
The \sphinxstyleemphasis{harmonic} dihedral style uses the potential
\begin{equation*}
\begin{split}E = K [ 1 + d  \cos (n \phi) ]\end{split}
\end{equation*}
\sphinxAtStartPar
The following coefficients must be defined for each dihedral type via the
{\hyperref[\detokenize{dihedral_coeff::doc}]{\sphinxcrossref{\DUrole{doc}{dihedral\_coeff}}}} command as in the example above, or in
the data file or restart files read by the {\hyperref[\detokenize{read_data::doc}]{\sphinxcrossref{\DUrole{doc}{read\_data}}}}
or {\hyperref[\detokenize{read_restart::doc}]{\sphinxcrossref{\DUrole{doc}{read\_restart}}}} commands:
\begin{itemize}
\item {} 
\sphinxAtStartPar
\(K\) (energy)

\item {} 
\sphinxAtStartPar
\(d\) (+1 or \sphinxhyphen{}1)

\item {} 
\sphinxAtStartPar
\(n\) (integer \textgreater{}= 0)

\end{itemize}

\begin{sphinxadmonition}{note}{Note:}
\sphinxAtStartPar
Here are important points to take note of when defining LAMMPS
dihedral coefficients for the harmonic style, so that they are
compatible with how harmonic dihedrals are defined by other force
fields:
\end{sphinxadmonition}
\begin{itemize}
\item {} 
\sphinxAtStartPar
The LAMMPS convention is that the trans position = 180 degrees, while
in some force fields trans = 0 degrees.

\item {} 
\sphinxAtStartPar
Some force fields reverse the sign convention on \(d\).

\item {} 
\sphinxAtStartPar
Some force fields let \(n\) be positive or negative which corresponds to
\(d = 1\) or \(d = -1\) for the harmonic style.

\end{itemize}


\bigskip\hrule\bigskip


\sphinxAtStartPar
Styles with a \sphinxstyleemphasis{gpu}, \sphinxstyleemphasis{intel}, \sphinxstyleemphasis{kk}, \sphinxstyleemphasis{omp}, or \sphinxstyleemphasis{opt} suffix are
functionally the same as the corresponding style without the suffix.
They have been optimized to run faster, depending on your available
hardware, as discussed on the {\hyperref[\detokenize{Speed_packages::doc}]{\sphinxcrossref{\DUrole{doc}{Accelerator packages}}}}
page.  The accelerated styles take the same arguments and should
produce the same results, except for round\sphinxhyphen{}off and precision issues.

\sphinxAtStartPar
These accelerated styles are part of the GPU, INTEL, KOKKOS,
OPENMP, and OPT packages, respectively.  They are only enabled if
LAMMPS was built with those packages.  See the {\hyperref[\detokenize{Build_package::doc}]{\sphinxcrossref{\DUrole{doc}{Build package}}}} page for more info.

\sphinxAtStartPar
You can specify the accelerated styles explicitly in your input script
by including their suffix, or you can use the {\hyperref[\detokenize{Run_options::doc}]{\sphinxcrossref{\DUrole{doc}{\sphinxhyphen{}suffix command\sphinxhyphen{}line switch}}}} when you invoke LAMMPS, or you can use the
{\hyperref[\detokenize{suffix::doc}]{\sphinxcrossref{\DUrole{doc}{suffix}}}} command in your input script.

\sphinxAtStartPar
See the {\hyperref[\detokenize{Speed_packages::doc}]{\sphinxcrossref{\DUrole{doc}{Accelerator packages}}}} page for more
instructions on how to use the accelerated styles effectively.


\bigskip\hrule\bigskip



\subsection{Restrictions}
\label{\detokenize{dihedral_harmonic:restrictions}}
\sphinxAtStartPar
This dihedral style can only be used if LAMMPS was built with the
MOLECULE package.  See the {\hyperref[\detokenize{Build_package::doc}]{\sphinxcrossref{\DUrole{doc}{Build package}}}} doc page
for more info.


\subsection{Related commands}
\label{\detokenize{dihedral_harmonic:related-commands}}
\sphinxAtStartPar
{\hyperref[\detokenize{dihedral_coeff::doc}]{\sphinxcrossref{\DUrole{doc}{dihedral\_coeff}}}}


\subsection{Default}
\label{\detokenize{dihedral_harmonic:default}}
\sphinxAtStartPar
none

\sphinxstepscope

\index{dihedral\_style helix@\spxentry{dihedral\_style helix}}\ignorespaces 
\index{dihedral\_style helix/omp@\spxentry{dihedral\_style helix/omp}}\ignorespaces 

\section{dihedral\_style helix command}
\label{\detokenize{dihedral_helix:dihedral-style-helix-command}}\label{\detokenize{dihedral_helix:index-1}}\label{\detokenize{dihedral_helix:index-0}}\label{\detokenize{dihedral_helix::doc}}
\sphinxAtStartPar
Accelerator Variants: \sphinxstyleemphasis{helix/omp}


\subsection{Syntax}
\label{\detokenize{dihedral_helix:syntax}}
\begin{sphinxVerbatim}[commandchars=\\\{\}]
\PYG{k}{dihedral\PYGZus{}style}\PYG{+w}{ }\PYG{n}{helix}
\end{sphinxVerbatim}


\subsection{Examples}
\label{\detokenize{dihedral_helix:examples}}
\begin{sphinxVerbatim}[commandchars=\\\{\}]
\PYG{k}{dihedral\PYGZus{}style}\PYG{+w}{ }\PYG{n}{helix}
\PYG{k}{dihedral\PYGZus{}coeff}\PYG{+w}{ }\PYG{l+m}{1}\PYG{+w}{ }\PYG{l+m}{80.0}\PYG{+w}{ }\PYG{l+m}{100.0}\PYG{+w}{ }\PYG{l+m}{40.0}
\end{sphinxVerbatim}


\subsection{Description}
\label{\detokenize{dihedral_helix:description}}
\sphinxAtStartPar
The \sphinxstyleemphasis{helix} dihedral style uses the potential
\begin{equation*}
\begin{split}E = A [1 - \cos(\theta)] + B [1 + \cos(3 \theta)] +
    C [1 + \cos(\theta + \frac{\pi}{4})]\end{split}
\end{equation*}
\sphinxAtStartPar
This coarse\sphinxhyphen{}grain dihedral potential is described in {\hyperref[\detokenize{dihedral_helix:guo}]{\sphinxcrossref{\DUrole{std,std-ref}{(Guo)}}}}.
For dihedral angles in the helical region, the energy function is
represented by a standard potential consisting of three minima, one
corresponding to the trans (t) state and the other to gauche states
(g+ and g\sphinxhyphen{}).  The paper describes how the \(A\), \(B\) and,
\(C\) parameters are chosen so as to balance secondary (largely
driven by local interactions) and
tertiary structure (driven by long\sphinxhyphen{}range interactions).

\sphinxAtStartPar
The following coefficients must be defined for each dihedral type via the
{\hyperref[\detokenize{dihedral_coeff::doc}]{\sphinxcrossref{\DUrole{doc}{dihedral\_coeff}}}} command as in the example above, or in
the data file or restart files read by the {\hyperref[\detokenize{read_data::doc}]{\sphinxcrossref{\DUrole{doc}{read\_data}}}}
or {\hyperref[\detokenize{read_restart::doc}]{\sphinxcrossref{\DUrole{doc}{read\_restart}}}} commands:
\begin{itemize}
\item {} 
\sphinxAtStartPar
\(A\) (energy)

\item {} 
\sphinxAtStartPar
\(B\) (energy)

\item {} 
\sphinxAtStartPar
\(C\) (energy)

\end{itemize}


\bigskip\hrule\bigskip


\sphinxAtStartPar
Styles with a \sphinxstyleemphasis{gpu}, \sphinxstyleemphasis{intel}, \sphinxstyleemphasis{kk}, \sphinxstyleemphasis{omp}, or \sphinxstyleemphasis{opt} suffix are
functionally the same as the corresponding style without the suffix.
They have been optimized to run faster, depending on your available
hardware, as discussed on the {\hyperref[\detokenize{Speed_packages::doc}]{\sphinxcrossref{\DUrole{doc}{Accelerator packages}}}}
page.  The accelerated styles take the same arguments and should
produce the same results, except for round\sphinxhyphen{}off and precision issues.

\sphinxAtStartPar
These accelerated styles are part of the GPU, INTEL, KOKKOS,
OPENMP, and OPT packages, respectively.  They are only enabled if
LAMMPS was built with those packages.  See the {\hyperref[\detokenize{Build_package::doc}]{\sphinxcrossref{\DUrole{doc}{Build package}}}} page for more info.

\sphinxAtStartPar
You can specify the accelerated styles explicitly in your input script
by including their suffix, or you can use the {\hyperref[\detokenize{Run_options::doc}]{\sphinxcrossref{\DUrole{doc}{\sphinxhyphen{}suffix command\sphinxhyphen{}line switch}}}} when you invoke LAMMPS, or you can use the
{\hyperref[\detokenize{suffix::doc}]{\sphinxcrossref{\DUrole{doc}{suffix}}}} command in your input script.

\sphinxAtStartPar
See the {\hyperref[\detokenize{Speed_packages::doc}]{\sphinxcrossref{\DUrole{doc}{Accelerator packages}}}} page for more
instructions on how to use the accelerated styles effectively.


\bigskip\hrule\bigskip



\subsection{Restrictions}
\label{\detokenize{dihedral_helix:restrictions}}
\sphinxAtStartPar
This dihedral style can only be used if LAMMPS was built with the
EXTRA\sphinxhyphen{}MOLECULE package.  See the {\hyperref[\detokenize{Build_package::doc}]{\sphinxcrossref{\DUrole{doc}{Build package}}}} doc page
for more info.


\subsection{Related commands}
\label{\detokenize{dihedral_helix:related-commands}}
\sphinxAtStartPar
{\hyperref[\detokenize{dihedral_coeff::doc}]{\sphinxcrossref{\DUrole{doc}{dihedral\_coeff}}}}


\subsection{Default}
\label{\detokenize{dihedral_helix:default}}
\sphinxAtStartPar
none


\bigskip\hrule\bigskip

\phantomsection\label{\detokenize{dihedral_helix:guo}}
\sphinxAtStartPar
\sphinxstylestrong{(Guo)} Guo and Thirumalai, Journal of Molecular Biology, 263, 323\sphinxhyphen{}43 (1996).

\sphinxstepscope

\index{dihedral\_style hybrid@\spxentry{dihedral\_style hybrid}}\ignorespaces 
\index{dihedral\_style hybrid/kk@\spxentry{dihedral\_style hybrid/kk}}\ignorespaces 

\section{dihedral\_style hybrid command}
\label{\detokenize{dihedral_hybrid:dihedral-style-hybrid-command}}\label{\detokenize{dihedral_hybrid:index-1}}\label{\detokenize{dihedral_hybrid:index-0}}\label{\detokenize{dihedral_hybrid::doc}}
\sphinxAtStartPar
Accelerator Variants: \sphinxstyleemphasis{hybrid/kk}


\subsection{Syntax}
\label{\detokenize{dihedral_hybrid:syntax}}
\begin{sphinxVerbatim}[commandchars=\\\{\}]
\PYG{k}{dihedral\PYGZus{}style}\PYG{+w}{ }\PYG{n}{hybrid}\PYG{+w}{ }\PYG{n}{style1}\PYG{+w}{ }\PYG{n}{style2}\PYG{+w}{ }\PYG{n}{...}
\end{sphinxVerbatim}
\begin{itemize}
\item {} 
\sphinxAtStartPar
style1,style2 = list of one or more dihedral styles

\end{itemize}


\subsection{Examples}
\label{\detokenize{dihedral_hybrid:examples}}
\begin{sphinxVerbatim}[commandchars=\\\{\}]
\PYG{k}{dihedral\PYGZus{}style}\PYG{+w}{ }\PYG{n}{hybrid}\PYG{+w}{ }\PYG{n}{harmonic}\PYG{+w}{ }\PYG{n}{helix}
\PYG{k}{dihedral\PYGZus{}coeff}\PYG{+w}{ }\PYG{l+m}{1}\PYG{+w}{ }\PYG{n}{harmonic}\PYG{+w}{ }\PYG{l+m}{6.0}\PYG{+w}{ }\PYG{l+m}{1}\PYG{+w}{ }\PYG{l+m}{3}
\PYG{k}{dihedral\PYGZus{}coeff}\PYG{+w}{ }\PYG{l+m}{2}\PYG{o}{*}\PYG{+w}{ }\PYG{n}{helix}\PYG{+w}{ }\PYG{l+m}{10}\PYG{+w}{ }\PYG{l+m}{10}\PYG{+w}{ }\PYG{l+m}{10}
\end{sphinxVerbatim}


\subsection{Description}
\label{\detokenize{dihedral_hybrid:description}}
\sphinxAtStartPar
The \sphinxstyleemphasis{hybrid} style enables the use of multiple dihedral styles in one
simulation.  An dihedral style is assigned to each dihedral type.  For
example, dihedrals in a polymer flow (of dihedral type 1) could be
computed with a \sphinxstyleemphasis{harmonic} potential and dihedrals in the wall
boundary (of dihedral type 2) could be computed with a \sphinxstyleemphasis{helix}
potential.  The assignment of dihedral type to style is made via the
{\hyperref[\detokenize{dihedral_coeff::doc}]{\sphinxcrossref{\DUrole{doc}{dihedral\_coeff}}}} command or in the data file.

\sphinxAtStartPar
In the dihedral\_coeff commands, the name of a dihedral style must be
added after the dihedral type, with the remaining coefficients being
those appropriate to that style.  In the example above, the 2
dihedral\_coeff commands set dihedrals of dihedral type 1 to be
computed with a \sphinxstyleemphasis{harmonic} potential with coefficients 6.0, 1, 3 for
K, d, n.  All other dihedral types (2\sphinxhyphen{}N) are computed with a \sphinxstyleemphasis{helix}
potential with coefficients 10, 10, 10 for A, B, C.

\sphinxAtStartPar
If dihedral coefficients are specified in the data file read via the
{\hyperref[\detokenize{read_data::doc}]{\sphinxcrossref{\DUrole{doc}{read\_data}}}} command, then the same rule applies.
E.g. “harmonic” or “helix”, must be added after the dihedral type, for
each line in the “Dihedral Coeffs” section, e.g.

\begin{sphinxVerbatim}[commandchars=\\\{\}]
Dihedral Coeffs

1 harmonic 6.0 1 3
2 helix 10 10 10
...
\end{sphinxVerbatim}

\sphinxAtStartPar
If \sphinxstyleemphasis{class2} is one of the dihedral hybrid styles, the same rule holds
for specifying additional AngleTorsion (and EndBondTorsion, etc)
coefficients either via the input script or in the data file.
I.e. \sphinxstyleemphasis{class2} must be added to each line after the dihedral type.  For
lines in the AngleTorsion (or EndBondTorsion, etc) Coeffs section of the data
file for dihedral types that are not \sphinxstyleemphasis{class2}, you must use an
dihedral style of \sphinxstyleemphasis{skip} as a placeholder, e.g.

\begin{sphinxVerbatim}[commandchars=\\\{\}]
AngleTorsion Coeffs

1 skip
2 class2 1.0 1.0 1.0 3.0 3.0 3.0 30.0 50.0
...
\end{sphinxVerbatim}

\sphinxAtStartPar
Note that it is not necessary to use the dihedral style \sphinxstyleemphasis{skip} in the
input script, since AngleTorsion (or EndBondTorsion, etc) coefficients
need not be specified at all for dihedral types that are not \sphinxstyleemphasis{class2}.

\sphinxAtStartPar
A dihedral style of \sphinxstyleemphasis{none} with no additional coefficients can be used
in place of a dihedral style, either in a input script dihedral\_coeff
command or in the data file, if you desire to turn off interactions
for specific dihedral types.


\bigskip\hrule\bigskip


\sphinxAtStartPar
Styles with a \sphinxstyleemphasis{gpu}, \sphinxstyleemphasis{intel}, \sphinxstyleemphasis{kk}, \sphinxstyleemphasis{omp}, or \sphinxstyleemphasis{opt} suffix are
functionally the same as the corresponding style without the suffix.
They have been optimized to run faster, depending on your available
hardware, as discussed on the {\hyperref[\detokenize{Speed_packages::doc}]{\sphinxcrossref{\DUrole{doc}{Accelerator packages}}}}
page.  The accelerated styles take the same arguments and should
produce the same results, except for round\sphinxhyphen{}off and precision issues.

\sphinxAtStartPar
These accelerated styles are part of the GPU, INTEL, KOKKOS,
OPENMP, and OPT packages, respectively.  They are only enabled if
LAMMPS was built with those packages.  See the {\hyperref[\detokenize{Build_package::doc}]{\sphinxcrossref{\DUrole{doc}{Build package}}}} page for more info.

\sphinxAtStartPar
You can specify the accelerated styles explicitly in your input script
by including their suffix, or you can use the {\hyperref[\detokenize{Run_options::doc}]{\sphinxcrossref{\DUrole{doc}{\sphinxhyphen{}suffix command\sphinxhyphen{}line switch}}}} when you invoke LAMMPS, or you can use the
{\hyperref[\detokenize{suffix::doc}]{\sphinxcrossref{\DUrole{doc}{suffix}}}} command in your input script.

\sphinxAtStartPar
See the {\hyperref[\detokenize{Speed_packages::doc}]{\sphinxcrossref{\DUrole{doc}{Accelerator packages}}}} page for more
instructions on how to use the accelerated styles effectively.


\bigskip\hrule\bigskip



\subsection{Restrictions}
\label{\detokenize{dihedral_hybrid:restrictions}}
\sphinxAtStartPar
This dihedral style can only be used if LAMMPS was built with the
MOLECULE package.  See the {\hyperref[\detokenize{Build_package::doc}]{\sphinxcrossref{\DUrole{doc}{Build package}}}} doc page
for more info.

\sphinxAtStartPar
Unlike other dihedral styles, the hybrid dihedral style does not store
dihedral coefficient info for individual sub\sphinxhyphen{}styles in {\hyperref[\detokenize{restart::doc}]{\sphinxcrossref{\DUrole{doc}{binary
restart files}}}} or {\hyperref[\detokenize{write_data::doc}]{\sphinxcrossref{\DUrole{doc}{data files}}}}.  Thus when
restarting a simulation, you need to re\sphinxhyphen{}specify the dihedral\_coeff
commands.


\subsection{Related commands}
\label{\detokenize{dihedral_hybrid:related-commands}}
\sphinxAtStartPar
{\hyperref[\detokenize{dihedral_coeff::doc}]{\sphinxcrossref{\DUrole{doc}{dihedral\_coeff}}}}


\subsection{Default}
\label{\detokenize{dihedral_hybrid:default}}
\sphinxAtStartPar
none

\sphinxstepscope

\index{dihedral\_style lepton@\spxentry{dihedral\_style lepton}}\ignorespaces 
\index{dihedral\_style lepton/omp@\spxentry{dihedral\_style lepton/omp}}\ignorespaces 

\section{dihedral\_style lepton command}
\label{\detokenize{dihedral_lepton:dihedral-style-lepton-command}}\label{\detokenize{dihedral_lepton:index-1}}\label{\detokenize{dihedral_lepton:index-0}}\label{\detokenize{dihedral_lepton::doc}}
\sphinxAtStartPar
Accelerator Variants: \sphinxstyleemphasis{lepton/omp}


\subsection{Syntax}
\label{\detokenize{dihedral_lepton:syntax}}
\begin{sphinxVerbatim}[commandchars=\\\{\}]
\PYG{k}{dihedral\PYGZus{}style}\PYG{+w}{ }\PYG{n}{lepton}
\end{sphinxVerbatim}


\subsection{Examples}
\label{\detokenize{dihedral_lepton:examples}}
\begin{sphinxVerbatim}[commandchars=\\\{\}]
\PYG{k}{dihedral\PYGZus{}style}\PYG{+w}{ }\PYG{n}{lepton}

\PYG{k}{dihedral\PYGZus{}coeff}\PYG{+w}{  }\PYG{l+m}{1}\PYG{+w}{ }\PYG{l+s}{\PYGZdq{}}\PYG{l+s}{k*(1 + d*cos(n*phi)); k=75.0; d=1; n=2}\PYG{l+s}{\PYGZdq{}}
\PYG{k}{dihedral\PYGZus{}coeff}\PYG{+w}{  }\PYG{l+m}{2}\PYG{+w}{ }\PYG{l+s}{\PYGZdq{}}\PYG{l+s}{45*(1\PYGZhy{}cos(4*phi))}\PYG{l+s}{\PYGZdq{}}
\PYG{k}{dihedral\PYGZus{}coeff}\PYG{+w}{  }\PYG{l+m}{2}\PYG{+w}{ }\PYG{l+s}{\PYGZdq{}}\PYG{l+s}{k2*cos(phi) + k3*cos(phi)\PYGZca{}2; k2=100.0}\PYG{l+s}{\PYGZdq{}}
\PYG{k}{dihedral\PYGZus{}coeff}\PYG{+w}{  }\PYG{l+m}{3}\PYG{+w}{ }\PYG{l+s}{\PYGZdq{}}\PYG{l+s}{k*(phi\PYGZhy{}phi0)\PYGZca{}2; k=85.0; phi0=120.0}\PYG{l+s}{\PYGZdq{}}
\end{sphinxVerbatim}


\subsection{Description}
\label{\detokenize{dihedral_lepton:description}}
\sphinxAtStartPar
\DUrole{versionmodified,added}{Added in version 8Feb2023.}

\sphinxAtStartPar
Dihedral style \sphinxstyleemphasis{lepton} computes dihedral interactions between four
atoms forming a dihedral angle with a custom potential function.  The
potential function must be provided as an expression string using “phi”
as the dihedral angle variable.  For example \sphinxtitleref{“200.0*(phi\sphinxhyphen{}120.0)\textasciicircum{}2”}
represents a {\hyperref[\detokenize{dihedral_quadratic::doc}]{\sphinxcrossref{\DUrole{doc}{quadratic dihedral}}}} potential
around a 120 degree dihedral angle with a force constant \sphinxstyleemphasis{K} of 200.0
energy units:
\begin{equation*}
\begin{split}U_{dihedral,i} = K (\phi_i - \phi_0)^2\end{split}
\end{equation*}
\sphinxAtStartPar
The \sphinxhref{https://simtk.org/projects/lepton}{Lepton library}, that the
\sphinxstyleemphasis{lepton} dihedral style interfaces with, evaluates this expression
string at run time to compute the pairwise energy.  It also creates an
analytical representation of the first derivative of this expression
with respect to “phi” and then uses that to compute the force between
the dihedral atoms as defined by the topology data.

\sphinxAtStartPar
The potential function expression for each dihedral type is provided via the
{\hyperref[\detokenize{dihedral_coeff::doc}]{\sphinxcrossref{\DUrole{doc}{dihedral\_coeff}}}} command as in the example above, or in
the data file or restart files read by the {\hyperref[\detokenize{read_data::doc}]{\sphinxcrossref{\DUrole{doc}{read\_data}}}}
or {\hyperref[\detokenize{read_restart::doc}]{\sphinxcrossref{\DUrole{doc}{read\_restart}}}} commands.  The expression is in energy units.

\sphinxAtStartPar
The Lepton expression must be either enclosed in quotes or must not
contain any whitespace so that LAMMPS recognizes it as a single keyword.
More on valid Lepton expressions below.  Dihedral angles are internally
computed in radians and thus the expression must assume “phi” is in
radians.


\bigskip\hrule\bigskip



\subsection{Lepton expression syntax and features}
\label{\detokenize{dihedral_lepton:lepton-expression-syntax-and-features}}
\sphinxAtStartPar
Lepton supports the following operators in expressions:


\begin{savenotes}\sphinxattablestart
\sphinxthistablewithglobalstyle
\centering
\begin{tabulary}{\linewidth}[t]{TTTTTTTTTTTTTT}
\sphinxtoprule
\sphinxtableatstartofbodyhook
\sphinxAtStartPar
+
&
\sphinxAtStartPar
Add
&&
\sphinxAtStartPar
\sphinxhyphen{}
&
\sphinxAtStartPar
Subtract
&&
\sphinxAtStartPar
*
&
\sphinxAtStartPar
Multiply
&&
\sphinxAtStartPar
/
&
\sphinxAtStartPar
Divide
&&
\sphinxAtStartPar
\textasciicircum{}
&
\sphinxAtStartPar
Power
\\
\sphinxbottomrule
\end{tabulary}
\sphinxtableafterendhook\par
\sphinxattableend\end{savenotes}

\sphinxAtStartPar
The following mathematical functions are available:


\begin{savenotes}\sphinxattablestart
\sphinxthistablewithglobalstyle
\centering
\begin{tabulary}{\linewidth}[t]{TTTT}
\sphinxtoprule
\sphinxtableatstartofbodyhook
\sphinxAtStartPar
sqrt(x)
&
\sphinxAtStartPar
Square root
&
\sphinxAtStartPar
exp(x)
&
\sphinxAtStartPar
Exponential
\\
\sphinxhline
\sphinxAtStartPar
log(x)
&
\sphinxAtStartPar
Natural logarithm
&
\sphinxAtStartPar
sin(x)
&
\sphinxAtStartPar
Sine (angle in radians)
\\
\sphinxhline
\sphinxAtStartPar
cos(x)
&
\sphinxAtStartPar
Cosine (angle in radians)
&
\sphinxAtStartPar
sec(x)
&
\sphinxAtStartPar
Secant (angle in radians)
\\
\sphinxhline
\sphinxAtStartPar
csc(x)
&
\sphinxAtStartPar
Cosecant (angle in radians)
&
\sphinxAtStartPar
tan(x)
&
\sphinxAtStartPar
Tangent (angle in radians)
\\
\sphinxhline
\sphinxAtStartPar
cot(x)
&
\sphinxAtStartPar
Cotangent (angle in radians)
&
\sphinxAtStartPar
asin(x)
&
\sphinxAtStartPar
Inverse sine (in radians)
\\
\sphinxhline
\sphinxAtStartPar
acos(x)
&
\sphinxAtStartPar
Inverse cosine (in radians)
&
\sphinxAtStartPar
atan(x)
&
\sphinxAtStartPar
Inverse tangent (in radians)
\\
\sphinxhline
\sphinxAtStartPar
sinh(x)
&
\sphinxAtStartPar
Hyperbolic sine
&
\sphinxAtStartPar
cosh(x)
&
\sphinxAtStartPar
Hyperbolic cosine
\\
\sphinxhline
\sphinxAtStartPar
tanh(x)
&
\sphinxAtStartPar
Hyperbolic tangent
&
\sphinxAtStartPar
erf(x)
&
\sphinxAtStartPar
Error function
\\
\sphinxhline
\sphinxAtStartPar
erfc(x)
&
\sphinxAtStartPar
Complementary Error function
&
\sphinxAtStartPar
abs(x)
&
\sphinxAtStartPar
Absolute value
\\
\sphinxhline
\sphinxAtStartPar
min(x,y)
&
\sphinxAtStartPar
Minimum of two values
&
\sphinxAtStartPar
max(x,y)
&
\sphinxAtStartPar
Maximum of two values
\\
\sphinxhline
\sphinxAtStartPar
delta(x)
&
\sphinxAtStartPar
delta(x) is 1 for \sphinxtitleref{x = 0}, otherwise 0
&
\sphinxAtStartPar
step(x)
&
\sphinxAtStartPar
step(x) is 0 for \sphinxtitleref{x \textless{} 0}, otherwise 1
\\
\sphinxbottomrule
\end{tabulary}
\sphinxtableafterendhook\par
\sphinxattableend\end{savenotes}

\sphinxAtStartPar
Numbers may be given in either decimal or exponential form.  All of the
following are valid numbers: \sphinxtitleref{5}, \sphinxtitleref{\sphinxhyphen{}3.1}, \sphinxtitleref{1e6}, and \sphinxtitleref{3.12e\sphinxhyphen{}2}.

\sphinxAtStartPar
As an extension to the standard Lepton syntax, it is also possible to
use LAMMPS {\hyperref[\detokenize{variable::doc}]{\sphinxcrossref{\DUrole{doc}{variables}}}} in the format “v\_name”.  Before
evaluating the expression, “v\_name” will be replaced with the value of
the variable “name”.  This is compatible with all kinds of scalar
variables, but not with vectors, arrays, local, or per\sphinxhyphen{}atom
variables.  If necessary, a custom scalar variable needs to be defined
that can access the desired (single) item from a non\sphinxhyphen{}scalar variable.
As an example, the following lines will instruct LAMMPS to ramp
the force constant for a harmonic bond from 100.0 to 200.0 during the
next run:

\begin{sphinxVerbatim}[commandchars=\\\{\}]
\PYG{k}{variable }\PYG{n+nv+nvIdentifier}{fconst}\PYG{+w}{ }\PYG{n}{equal}\PYG{+w}{ }\PYG{n}{ramp}\PYG{n+nv}{(}\PYG{n+nv}{100.0, 200}\PYG{n+nv}{)}
\PYG{k}{bond\PYGZus{}style}\PYG{+w}{ }\PYG{n}{lepton}
\PYG{k}{bond\PYGZus{}coeff}\PYG{+w}{ }\PYG{l+m}{1}\PYG{+w}{ }\PYG{l+m}{1.5}\PYG{+w}{ }\PYG{l+s}{\PYGZdq{}}\PYG{l+s}{v\PYGZus{}fconst * (r\PYGZca{}2)}\PYG{l+s}{\PYGZdq{}}
\end{sphinxVerbatim}

\sphinxAtStartPar
An expression may be followed by definitions for intermediate values that appear in the
expression. A semicolon “;” is used as a delimiter between value definitions. For example,
the expression:

\begin{sphinxVerbatim}[commandchars=\\\{\}]
\PYG{n}{a}\PYG{o}{\PYGZca{}}\PYG{l+m+mi}{2}\PYG{o}{+}\PYG{n}{a}\PYG{o}{*}\PYG{n}{b}\PYG{o}{+}\PYG{n}{b}\PYG{o}{\PYGZca{}}\PYG{l+m+mi}{2}\PYG{p}{;}\PYG{+w}{ }\PYG{n}{a}\PYG{o}{=}\PYG{n}{a1}\PYG{o}{+}\PYG{n}{a2}\PYG{p}{;}\PYG{+w}{ }\PYG{n}{b}\PYG{o}{=}\PYG{n}{b1}\PYG{o}{+}\PYG{n}{b2}
\end{sphinxVerbatim}

\sphinxAtStartPar
is exactly equivalent to

\begin{sphinxVerbatim}[commandchars=\\\{\}]
\PYG{p}{(}\PYG{n}{a1}\PYG{o}{+}\PYG{n}{a2}\PYG{p}{)}\PYG{o}{\PYGZca{}}\PYG{l+m+mi}{2}\PYG{o}{+}\PYG{p}{(}\PYG{n}{a1}\PYG{o}{+}\PYG{n}{a2}\PYG{p}{)}\PYG{o}{*}\PYG{p}{(}\PYG{n}{b1}\PYG{o}{+}\PYG{n}{b2}\PYG{p}{)}\PYG{o}{+}\PYG{p}{(}\PYG{n}{b1}\PYG{o}{+}\PYG{n}{b2}\PYG{p}{)}\PYG{o}{\PYGZca{}}\PYG{l+m+mi}{2}
\end{sphinxVerbatim}

\sphinxAtStartPar
The definition of an intermediate value may itself involve other
intermediate values. Whitespace and quotation characters (’\textquotesingle{}’ and ‘”’)
are ignored.  All uses of a value must appear \sphinxstyleemphasis{before} that value’s
definition.  For efficiency reasons, the expression string is parsed,
optimized, and then stored in an internal, pre\sphinxhyphen{}parsed representation for
evaluation.

\sphinxAtStartPar
Evaluating a Lepton expression is typically between 2.5 and 5 times
slower than the corresponding compiled and optimized C++ code.  If
additional speed or GPU acceleration (via GPU or KOKKOS) is required,
the interaction can be represented as a table.  Suitable table files
can be created either internally using the {\hyperref[\detokenize{pair_write::doc}]{\sphinxcrossref{\DUrole{doc}{pair\_write}}}}
or {\hyperref[\detokenize{bond_write::doc}]{\sphinxcrossref{\DUrole{doc}{bond\_write}}}} command or through the Python scripts
in the {\hyperref[\detokenize{Tools:tabulate}]{\sphinxcrossref{\DUrole{std,std-ref}{tools/tabulate}}}} folder.


\bigskip\hrule\bigskip


\sphinxAtStartPar
Styles with a \sphinxstyleemphasis{gpu}, \sphinxstyleemphasis{intel}, \sphinxstyleemphasis{kk}, \sphinxstyleemphasis{omp}, or \sphinxstyleemphasis{opt} suffix are
functionally the same as the corresponding style without the suffix.
They have been optimized to run faster, depending on your available
hardware, as discussed on the {\hyperref[\detokenize{Speed_packages::doc}]{\sphinxcrossref{\DUrole{doc}{Accelerator packages}}}}
page.  The accelerated styles take the same arguments and should
produce the same results, except for round\sphinxhyphen{}off and precision issues.

\sphinxAtStartPar
These accelerated styles are part of the GPU, INTEL, KOKKOS,
OPENMP, and OPT packages, respectively.  They are only enabled if
LAMMPS was built with those packages.  See the {\hyperref[\detokenize{Build_package::doc}]{\sphinxcrossref{\DUrole{doc}{Build package}}}} page for more info.

\sphinxAtStartPar
You can specify the accelerated styles explicitly in your input script
by including their suffix, or you can use the {\hyperref[\detokenize{Run_options::doc}]{\sphinxcrossref{\DUrole{doc}{\sphinxhyphen{}suffix command\sphinxhyphen{}line switch}}}} when you invoke LAMMPS, or you can use the
{\hyperref[\detokenize{suffix::doc}]{\sphinxcrossref{\DUrole{doc}{suffix}}}} command in your input script.

\sphinxAtStartPar
See the {\hyperref[\detokenize{Speed_packages::doc}]{\sphinxcrossref{\DUrole{doc}{Accelerator packages}}}} page for more
instructions on how to use the accelerated styles effectively.


\bigskip\hrule\bigskip



\subsection{Restrictions}
\label{\detokenize{dihedral_lepton:restrictions}}
\sphinxAtStartPar
This dihedral style is part of the LEPTON package and only enabled if LAMMPS
was built with this package.  See the {\hyperref[\detokenize{Build_package::doc}]{\sphinxcrossref{\DUrole{doc}{Build package}}}} page for more info.


\subsection{Related commands}
\label{\detokenize{dihedral_lepton:related-commands}}
\sphinxAtStartPar
{\hyperref[\detokenize{dihedral_coeff::doc}]{\sphinxcrossref{\DUrole{doc}{dihedral\_coeff}}}}, {\hyperref[\detokenize{dihedral_table::doc}]{\sphinxcrossref{\DUrole{doc}{dihedral\_style table}}}},
{\hyperref[\detokenize{bond_lepton::doc}]{\sphinxcrossref{\DUrole{doc}{bond\_style lepton}}}}, {\hyperref[\detokenize{angle_lepton::doc}]{\sphinxcrossref{\DUrole{doc}{angle\_style lepton}}}}


\subsection{Default}
\label{\detokenize{dihedral_lepton:default}}
\sphinxAtStartPar
none

\sphinxstepscope

\index{dihedral\_style multi/harmonic@\spxentry{dihedral\_style multi/harmonic}}\ignorespaces 
\index{dihedral\_style multi/harmonic/omp@\spxentry{dihedral\_style multi/harmonic/omp}}\ignorespaces 

\section{dihedral\_style multi/harmonic command}
\label{\detokenize{dihedral_multi_harmonic:dihedral-style-multi-harmonic-command}}\label{\detokenize{dihedral_multi_harmonic:index-1}}\label{\detokenize{dihedral_multi_harmonic:index-0}}\label{\detokenize{dihedral_multi_harmonic::doc}}
\sphinxAtStartPar
Accelerator Variants: \sphinxstyleemphasis{multi/harmonic/omp}


\subsection{Syntax}
\label{\detokenize{dihedral_multi_harmonic:syntax}}
\begin{sphinxVerbatim}[commandchars=\\\{\}]
\PYG{k}{dihedral\PYGZus{}style}\PYG{+w}{ }\PYG{n}{multi}\PYG{o}{/}\PYG{n}{harmonic}
\end{sphinxVerbatim}


\subsection{Examples}
\label{\detokenize{dihedral_multi_harmonic:examples}}
\begin{sphinxVerbatim}[commandchars=\\\{\}]
\PYG{k}{dihedral\PYGZus{}style}\PYG{+w}{ }\PYG{n}{multi}\PYG{o}{/}\PYG{n}{harmonic}
\PYG{k}{dihedral\PYGZus{}coeff}\PYG{+w}{ }\PYG{l+m}{1}\PYG{+w}{ }\PYG{l+m}{20}\PYG{+w}{ }\PYG{l+m}{20}\PYG{+w}{ }\PYG{l+m}{20}\PYG{+w}{ }\PYG{l+m}{20}\PYG{+w}{ }\PYG{l+m}{20}
\end{sphinxVerbatim}


\subsection{Description}
\label{\detokenize{dihedral_multi_harmonic:description}}
\sphinxAtStartPar
The \sphinxstyleemphasis{multi/harmonic} dihedral style uses the potential
\begin{equation*}
\begin{split}E = \sum_{n=1,5} A_n  \cos^{n-1}(\phi)\end{split}
\end{equation*}
\sphinxAtStartPar
The following coefficients must be defined for each dihedral type via the
{\hyperref[\detokenize{dihedral_coeff::doc}]{\sphinxcrossref{\DUrole{doc}{dihedral\_coeff}}}} command as in the example above, or in
the data file or restart files read by the {\hyperref[\detokenize{read_data::doc}]{\sphinxcrossref{\DUrole{doc}{read\_data}}}}
or {\hyperref[\detokenize{read_restart::doc}]{\sphinxcrossref{\DUrole{doc}{read\_restart}}}} commands:
\begin{itemize}
\item {} 
\sphinxAtStartPar
\(A_1\) (energy)

\item {} 
\sphinxAtStartPar
\(A_2\) (energy)

\item {} 
\sphinxAtStartPar
\(A_3\) (energy)

\item {} 
\sphinxAtStartPar
\(A_4\) (energy)

\item {} 
\sphinxAtStartPar
\(A_5\) (energy)

\end{itemize}


\bigskip\hrule\bigskip


\sphinxAtStartPar
Styles with a \sphinxstyleemphasis{gpu}, \sphinxstyleemphasis{intel}, \sphinxstyleemphasis{kk}, \sphinxstyleemphasis{omp}, or \sphinxstyleemphasis{opt} suffix are
functionally the same as the corresponding style without the suffix.
They have been optimized to run faster, depending on your available
hardware, as discussed on the {\hyperref[\detokenize{Speed_packages::doc}]{\sphinxcrossref{\DUrole{doc}{Accelerator packages}}}}
page.  The accelerated styles take the same arguments and should
produce the same results, except for round\sphinxhyphen{}off and precision issues.

\sphinxAtStartPar
These accelerated styles are part of the GPU, INTEL, KOKKOS,
OPENMP, and OPT packages, respectively.  They are only enabled if
LAMMPS was built with those packages.  See the {\hyperref[\detokenize{Build_package::doc}]{\sphinxcrossref{\DUrole{doc}{Build package}}}} page for more info.

\sphinxAtStartPar
You can specify the accelerated styles explicitly in your input script
by including their suffix, or you can use the {\hyperref[\detokenize{Run_options::doc}]{\sphinxcrossref{\DUrole{doc}{\sphinxhyphen{}suffix command\sphinxhyphen{}line switch}}}} when you invoke LAMMPS, or you can use the
{\hyperref[\detokenize{suffix::doc}]{\sphinxcrossref{\DUrole{doc}{suffix}}}} command in your input script.

\sphinxAtStartPar
See the {\hyperref[\detokenize{Speed_packages::doc}]{\sphinxcrossref{\DUrole{doc}{Accelerator packages}}}} page for more
instructions on how to use the accelerated styles effectively.


\bigskip\hrule\bigskip



\subsection{Restrictions}
\label{\detokenize{dihedral_multi_harmonic:restrictions}}
\sphinxAtStartPar
This dihedral style can only be used if LAMMPS was built with the
MOLECULE package.  See the {\hyperref[\detokenize{Build_package::doc}]{\sphinxcrossref{\DUrole{doc}{Build package}}}} doc page
for more info.


\subsection{Related commands}
\label{\detokenize{dihedral_multi_harmonic:related-commands}}
\sphinxAtStartPar
{\hyperref[\detokenize{dihedral_coeff::doc}]{\sphinxcrossref{\DUrole{doc}{dihedral\_coeff}}}}


\subsection{Default}
\label{\detokenize{dihedral_multi_harmonic:default}}
\sphinxAtStartPar
none

\sphinxstepscope

\index{dihedral\_style nharmonic@\spxentry{dihedral\_style nharmonic}}\ignorespaces 
\index{dihedral\_style nharmonic/omp@\spxentry{dihedral\_style nharmonic/omp}}\ignorespaces 

\section{dihedral\_style nharmonic command}
\label{\detokenize{dihedral_nharmonic:dihedral-style-nharmonic-command}}\label{\detokenize{dihedral_nharmonic:index-1}}\label{\detokenize{dihedral_nharmonic:index-0}}\label{\detokenize{dihedral_nharmonic::doc}}
\sphinxAtStartPar
Accelerator Variants: \sphinxstyleemphasis{nharmonic/omp}


\subsection{Syntax}
\label{\detokenize{dihedral_nharmonic:syntax}}
\begin{sphinxVerbatim}[commandchars=\\\{\}]
\PYG{k}{dihedral\PYGZus{}style}\PYG{+w}{ }\PYG{n}{nharmonic}
\end{sphinxVerbatim}


\subsection{Examples}
\label{\detokenize{dihedral_nharmonic:examples}}
\begin{sphinxVerbatim}[commandchars=\\\{\}]
\PYG{k}{dihedral\PYGZus{}style}\PYG{+w}{ }\PYG{n}{nharmonic}
\PYG{k}{dihedral\PYGZus{}coeff}\PYG{+w}{ }\PYG{o}{*}\PYG{+w}{ }\PYG{l+m}{3}\PYG{+w}{ }\PYG{l+m}{10.0}\PYG{+w}{ }\PYG{l+m}{20.0}\PYG{+w}{ }\PYG{l+m}{30.0}
\end{sphinxVerbatim}


\subsection{Description}
\label{\detokenize{dihedral_nharmonic:description}}
\sphinxAtStartPar
The \sphinxstyleemphasis{nharmonic} dihedral style uses the potential:
\begin{equation*}
\begin{split}E = \sum_{i=1,n} A_i  \cos^{i-1}(\phi)\end{split}
\end{equation*}
\sphinxAtStartPar
The following coefficients must be defined for each dihedral type via the
{\hyperref[\detokenize{dihedral_coeff::doc}]{\sphinxcrossref{\DUrole{doc}{dihedral\_coeff}}}} command as in the example above, or in
the data file or restart files read by the {\hyperref[\detokenize{read_data::doc}]{\sphinxcrossref{\DUrole{doc}{read\_data}}}}
or {\hyperref[\detokenize{read_restart::doc}]{\sphinxcrossref{\DUrole{doc}{read\_restart}}}} commands:
\begin{itemize}
\item {} 
\sphinxAtStartPar
\(n\) (integer \textgreater{}=1)

\item {} 
\sphinxAtStartPar
\(A_1\) (energy)

\item {} 
\sphinxAtStartPar
\(A_2\) (energy)

\item {} 
\sphinxAtStartPar
…

\item {} 
\sphinxAtStartPar
\(A_n\) (energy)

\end{itemize}


\bigskip\hrule\bigskip


\sphinxAtStartPar
Styles with a \sphinxstyleemphasis{gpu}, \sphinxstyleemphasis{intel}, \sphinxstyleemphasis{kk}, \sphinxstyleemphasis{omp}, or \sphinxstyleemphasis{opt} suffix are
functionally the same as the corresponding style without the suffix.
They have been optimized to run faster, depending on your available
hardware, as discussed on the {\hyperref[\detokenize{Speed_packages::doc}]{\sphinxcrossref{\DUrole{doc}{Accelerator packages}}}}
page.  The accelerated styles take the same arguments and should
produce the same results, except for round\sphinxhyphen{}off and precision issues.

\sphinxAtStartPar
These accelerated styles are part of the GPU, INTEL, KOKKOS,
OPENMP, and OPT packages, respectively.  They are only enabled if
LAMMPS was built with those packages.  See the {\hyperref[\detokenize{Build_package::doc}]{\sphinxcrossref{\DUrole{doc}{Build package}}}} page for more info.

\sphinxAtStartPar
You can specify the accelerated styles explicitly in your input script
by including their suffix, or you can use the {\hyperref[\detokenize{Run_options::doc}]{\sphinxcrossref{\DUrole{doc}{\sphinxhyphen{}suffix command\sphinxhyphen{}line switch}}}} when you invoke LAMMPS, or you can use the
{\hyperref[\detokenize{suffix::doc}]{\sphinxcrossref{\DUrole{doc}{suffix}}}} command in your input script.

\sphinxAtStartPar
See the {\hyperref[\detokenize{Speed_packages::doc}]{\sphinxcrossref{\DUrole{doc}{Accelerator packages}}}} page for more
instructions on how to use the accelerated styles effectively.


\bigskip\hrule\bigskip



\subsection{Restrictions}
\label{\detokenize{dihedral_nharmonic:restrictions}}
\sphinxAtStartPar
This dihedral style can only be used if LAMMPS was built with the
EXTRA\sphinxhyphen{}MOLECULE package.  See the {\hyperref[\detokenize{Build_package::doc}]{\sphinxcrossref{\DUrole{doc}{Build package}}}} doc
page for more info.


\subsection{Related commands}
\label{\detokenize{dihedral_nharmonic:related-commands}}
\sphinxAtStartPar
{\hyperref[\detokenize{dihedral_coeff::doc}]{\sphinxcrossref{\DUrole{doc}{dihedral\_coeff}}}}


\subsection{Default}
\label{\detokenize{dihedral_nharmonic:default}}
\sphinxAtStartPar
none

\sphinxstepscope

\index{dihedral\_style none@\spxentry{dihedral\_style none}}\ignorespaces 

\section{dihedral\_style none command}
\label{\detokenize{dihedral_none:dihedral-style-none-command}}\label{\detokenize{dihedral_none:index-0}}\label{\detokenize{dihedral_none::doc}}

\subsection{Syntax}
\label{\detokenize{dihedral_none:syntax}}
\begin{sphinxVerbatim}[commandchars=\\\{\}]
\PYG{k}{dihedral\PYGZus{}style}\PYG{+w}{ }\PYG{n}{none}
\end{sphinxVerbatim}


\subsection{Examples}
\label{\detokenize{dihedral_none:examples}}
\begin{sphinxVerbatim}[commandchars=\\\{\}]
\PYG{k}{dihedral\PYGZus{}style}\PYG{+w}{ }\PYG{n}{none}
\end{sphinxVerbatim}


\subsection{Description}
\label{\detokenize{dihedral_none:description}}
\sphinxAtStartPar
Using a dihedral style of none means dihedral forces and energies are
not computed, even if quadruplets of dihedral atoms were listed in the
data file read by the {\hyperref[\detokenize{read_data::doc}]{\sphinxcrossref{\DUrole{doc}{read\_data}}}} command.

\sphinxAtStartPar
See the {\hyperref[\detokenize{dihedral_zero::doc}]{\sphinxcrossref{\DUrole{doc}{dihedral\_style zero}}}} command for a way to
calculate dihedral statistics, but compute no dihedral interactions.


\subsection{Restrictions}
\label{\detokenize{dihedral_none:restrictions}}\begin{quote}

\sphinxAtStartPar
none
\end{quote}


\subsection{Related commands}
\label{\detokenize{dihedral_none:related-commands}}
\sphinxAtStartPar
{\hyperref[\detokenize{dihedral_zero::doc}]{\sphinxcrossref{\DUrole{doc}{dihedral\_style zero}}}}


\subsection{Default}
\label{\detokenize{dihedral_none:default}}
\sphinxAtStartPar
none

\sphinxstepscope

\index{dihedral\_style opls@\spxentry{dihedral\_style opls}}\ignorespaces 
\index{dihedral\_style opls/intel@\spxentry{dihedral\_style opls/intel}}\ignorespaces 
\index{dihedral\_style opls/kk@\spxentry{dihedral\_style opls/kk}}\ignorespaces 
\index{dihedral\_style opls/omp@\spxentry{dihedral\_style opls/omp}}\ignorespaces 

\section{dihedral\_style opls command}
\label{\detokenize{dihedral_opls:dihedral-style-opls-command}}\label{\detokenize{dihedral_opls:index-3}}\label{\detokenize{dihedral_opls:index-2}}\label{\detokenize{dihedral_opls:index-1}}\label{\detokenize{dihedral_opls:index-0}}\label{\detokenize{dihedral_opls::doc}}
\sphinxAtStartPar
Accelerator Variants: \sphinxstyleemphasis{opls/intel}, \sphinxstyleemphasis{opls/kk}, \sphinxstyleemphasis{opls/omp}


\subsection{Syntax}
\label{\detokenize{dihedral_opls:syntax}}
\begin{sphinxVerbatim}[commandchars=\\\{\}]
\PYG{k}{dihedral\PYGZus{}style}\PYG{+w}{ }\PYG{n}{opls}
\end{sphinxVerbatim}


\subsection{Examples}
\label{\detokenize{dihedral_opls:examples}}
\begin{sphinxVerbatim}[commandchars=\\\{\}]
\PYG{k}{dihedral\PYGZus{}style}\PYG{+w}{ }\PYG{n}{opls}
\PYG{k}{dihedral\PYGZus{}coeff}\PYG{+w}{ }\PYG{l+m}{1}\PYG{+w}{ }\PYG{l+m}{1.740}\PYG{+w}{ }\PYG{o}{\PYGZhy{}}\PYG{l+m}{0.157}\PYG{+w}{ }\PYG{l+m}{0.279}\PYG{+w}{ }\PYG{l+m}{0.00}\PYG{+w}{   }\PYG{c}{\PYGZsh{} CT\PYGZhy{}CT\PYGZhy{}CT\PYGZhy{}CT}
\PYG{k}{dihedral\PYGZus{}coeff}\PYG{+w}{ }\PYG{l+m}{2}\PYG{+w}{ }\PYG{l+m}{0.000}\PYG{+w}{ }\PYG{l+m}{0.000}\PYG{+w}{ }\PYG{l+m}{0.366}\PYG{+w}{ }\PYG{l+m}{0.000}\PYG{+w}{   }\PYG{c}{\PYGZsh{} CT\PYGZhy{}CT\PYGZhy{}CT\PYGZhy{}HC}
\PYG{k}{dihedral\PYGZus{}coeff}\PYG{+w}{ }\PYG{l+m}{3}\PYG{+w}{ }\PYG{l+m}{0.000}\PYG{+w}{ }\PYG{l+m}{0.000}\PYG{+w}{ }\PYG{l+m}{0.318}\PYG{+w}{ }\PYG{l+m}{0.000}\PYG{+w}{   }\PYG{c}{\PYGZsh{} HC\PYGZhy{}CT\PYGZhy{}CT\PYGZhy{}HC}
\end{sphinxVerbatim}


\subsection{Description}
\label{\detokenize{dihedral_opls:description}}
\sphinxAtStartPar
The \sphinxstyleemphasis{opls} dihedral style uses the potential
\begin{equation*}
\begin{split}E = & \frac{1}{2} K_1 [1 + \cos(\phi)] + \frac{1}{2} K_2 [1 - \cos(2 \phi)] + \\
    & \frac{1}{2} K_3 [1 + \cos(3 \phi)] + \frac{1}{2} K_4 [1 - \cos(4 \phi)]\end{split}
\end{equation*}
\sphinxAtStartPar
Note that the usual 1/2 factor is not included in the K values.

\sphinxAtStartPar
This dihedral potential is used in the OPLS force field and is
described in {\hyperref[\detokenize{dihedral_opls:watkins}]{\sphinxcrossref{\DUrole{std,std-ref}{(Watkins)}}}}.

\sphinxAtStartPar
The following coefficients must be defined for each dihedral type via the
{\hyperref[\detokenize{dihedral_coeff::doc}]{\sphinxcrossref{\DUrole{doc}{dihedral\_coeff}}}} command as in the example above, or in
the data file or restart files read by the {\hyperref[\detokenize{read_data::doc}]{\sphinxcrossref{\DUrole{doc}{read\_data}}}}
or {\hyperref[\detokenize{read_restart::doc}]{\sphinxcrossref{\DUrole{doc}{read\_restart}}}} commands:
\begin{itemize}
\item {} 
\sphinxAtStartPar
\(K_1\) (energy)

\item {} 
\sphinxAtStartPar
\(K_2\) (energy)

\item {} 
\sphinxAtStartPar
\(K_3\) (energy)

\item {} 
\sphinxAtStartPar
\(K_4\) (energy)

\end{itemize}


\bigskip\hrule\bigskip


\sphinxAtStartPar
Styles with a \sphinxstyleemphasis{gpu}, \sphinxstyleemphasis{intel}, \sphinxstyleemphasis{kk}, \sphinxstyleemphasis{omp}, or \sphinxstyleemphasis{opt} suffix are
functionally the same as the corresponding style without the suffix.
They have been optimized to run faster, depending on your available
hardware, as discussed on the {\hyperref[\detokenize{Speed_packages::doc}]{\sphinxcrossref{\DUrole{doc}{Accelerator packages}}}}
page.  The accelerated styles take the same arguments and should
produce the same results, except for round\sphinxhyphen{}off and precision issues.

\sphinxAtStartPar
These accelerated styles are part of the GPU, INTEL, KOKKOS,
OPENMP, and OPT packages, respectively.  They are only enabled if
LAMMPS was built with those packages.  See the {\hyperref[\detokenize{Build_package::doc}]{\sphinxcrossref{\DUrole{doc}{Build package}}}} page for more info.

\sphinxAtStartPar
You can specify the accelerated styles explicitly in your input script
by including their suffix, or you can use the {\hyperref[\detokenize{Run_options::doc}]{\sphinxcrossref{\DUrole{doc}{\sphinxhyphen{}suffix command\sphinxhyphen{}line switch}}}} when you invoke LAMMPS, or you can use the
{\hyperref[\detokenize{suffix::doc}]{\sphinxcrossref{\DUrole{doc}{suffix}}}} command in your input script.

\sphinxAtStartPar
See the {\hyperref[\detokenize{Speed_packages::doc}]{\sphinxcrossref{\DUrole{doc}{Accelerator packages}}}} page for more
instructions on how to use the accelerated styles effectively.


\bigskip\hrule\bigskip



\subsection{Restrictions}
\label{\detokenize{dihedral_opls:restrictions}}
\sphinxAtStartPar
This dihedral style can only be used if LAMMPS was built with the
MOLECULE package.  See the {\hyperref[\detokenize{Build_package::doc}]{\sphinxcrossref{\DUrole{doc}{Build package}}}} doc page
for more info.


\subsection{Related commands}
\label{\detokenize{dihedral_opls:related-commands}}
\sphinxAtStartPar
{\hyperref[\detokenize{dihedral_coeff::doc}]{\sphinxcrossref{\DUrole{doc}{dihedral\_coeff}}}}


\subsection{Default}
\label{\detokenize{dihedral_opls:default}}
\sphinxAtStartPar
none


\bigskip\hrule\bigskip

\phantomsection\label{\detokenize{dihedral_opls:watkins}}
\sphinxAtStartPar
\sphinxstylestrong{(Watkins)} Watkins and Jorgensen, J Phys Chem A, 105, 4118\sphinxhyphen{}4125 (2001).

\sphinxstepscope

\index{dihedral\_style quadratic@\spxentry{dihedral\_style quadratic}}\ignorespaces 
\index{dihedral\_style quadratic/omp@\spxentry{dihedral\_style quadratic/omp}}\ignorespaces 

\section{dihedral\_style quadratic command}
\label{\detokenize{dihedral_quadratic:dihedral-style-quadratic-command}}\label{\detokenize{dihedral_quadratic:index-1}}\label{\detokenize{dihedral_quadratic:index-0}}\label{\detokenize{dihedral_quadratic::doc}}
\sphinxAtStartPar
Accelerator Variants: \sphinxstyleemphasis{quadratic/omp}


\subsection{Syntax}
\label{\detokenize{dihedral_quadratic:syntax}}
\begin{sphinxVerbatim}[commandchars=\\\{\}]
\PYG{k}{dihedral\PYGZus{}style}\PYG{+w}{ }\PYG{n}{quadratic}
\end{sphinxVerbatim}


\subsection{Examples}
\label{\detokenize{dihedral_quadratic:examples}}
\begin{sphinxVerbatim}[commandchars=\\\{\}]
\PYG{k}{dihedral\PYGZus{}style}\PYG{+w}{ }\PYG{n}{quadratic}
\PYG{k}{dihedral\PYGZus{}coeff}\PYG{+w}{ }\PYG{l+m}{100.0}\PYG{+w}{ }\PYG{l+m}{80.0}
\end{sphinxVerbatim}


\subsection{Description}
\label{\detokenize{dihedral_quadratic:description}}
\sphinxAtStartPar
The \sphinxstyleemphasis{quadratic} dihedral style uses the potential:
\begin{equation*}
\begin{split}E = K (\phi - \phi_0)^2\end{split}
\end{equation*}
\sphinxAtStartPar
This dihedral potential can be used to keep a dihedral in a predefined
value (cis=zero, right\sphinxhyphen{}hand convention is used).

\sphinxAtStartPar
The following coefficients must be defined for each dihedral type via
the {\hyperref[\detokenize{dihedral_coeff::doc}]{\sphinxcrossref{\DUrole{doc}{dihedral\_coeff}}}} command as in the example
above, or in the data file or restart files read by the
{\hyperref[\detokenize{read_data::doc}]{\sphinxcrossref{\DUrole{doc}{read\_data}}}} or {\hyperref[\detokenize{read_restart::doc}]{\sphinxcrossref{\DUrole{doc}{read\_restart}}}}
commands:
\begin{itemize}
\item {} 
\sphinxAtStartPar
\(K\) (energy)

\item {} 
\sphinxAtStartPar
\(\phi_0\) (degrees)

\end{itemize}

\sphinxAtStartPar
\(\phi_0\) is specified in degrees, but LAMMPS converts it to
radians internally; hence \(K\) is effectively energy per
radian\textasciicircum{}2.


\bigskip\hrule\bigskip


\sphinxAtStartPar
Styles with a \sphinxstyleemphasis{gpu}, \sphinxstyleemphasis{intel}, \sphinxstyleemphasis{kk}, \sphinxstyleemphasis{omp}, or \sphinxstyleemphasis{opt} suffix are
functionally the same as the corresponding style without the suffix.
They have been optimized to run faster, depending on your available
hardware, as discussed on the {\hyperref[\detokenize{Speed_packages::doc}]{\sphinxcrossref{\DUrole{doc}{Accelerator packages}}}}
page.  The accelerated styles take the same arguments and should
produce the same results, except for round\sphinxhyphen{}off and precision issues.

\sphinxAtStartPar
These accelerated styles are part of the GPU, INTEL, KOKKOS,
OPENMP, and OPT packages, respectively.  They are only enabled if
LAMMPS was built with those packages.  See the {\hyperref[\detokenize{Build_package::doc}]{\sphinxcrossref{\DUrole{doc}{Build package}}}} page for more info.

\sphinxAtStartPar
You can specify the accelerated styles explicitly in your input script
by including their suffix, or you can use the {\hyperref[\detokenize{Run_options::doc}]{\sphinxcrossref{\DUrole{doc}{\sphinxhyphen{}suffix command\sphinxhyphen{}line switch}}}} when you invoke LAMMPS, or you can use the
{\hyperref[\detokenize{suffix::doc}]{\sphinxcrossref{\DUrole{doc}{suffix}}}} command in your input script.

\sphinxAtStartPar
See the {\hyperref[\detokenize{Speed_packages::doc}]{\sphinxcrossref{\DUrole{doc}{Accelerator packages}}}} page for more
instructions on how to use the accelerated styles effectively.


\bigskip\hrule\bigskip



\subsection{Restrictions}
\label{\detokenize{dihedral_quadratic:restrictions}}
\sphinxAtStartPar
This dihedral style can only be used if LAMMPS was built with the
EXTRA\sphinxhyphen{}MOLECULE package.  See the {\hyperref[\detokenize{Build_package::doc}]{\sphinxcrossref{\DUrole{doc}{Build package}}}} doc
page for more info.


\subsection{Related commands}
\label{\detokenize{dihedral_quadratic:related-commands}}
\sphinxAtStartPar
{\hyperref[\detokenize{dihedral_coeff::doc}]{\sphinxcrossref{\DUrole{doc}{dihedral\_coeff}}}}


\subsection{Default}
\label{\detokenize{dihedral_quadratic:default}}
\sphinxAtStartPar
none

\sphinxstepscope

\index{dihedral\_style spherical@\spxentry{dihedral\_style spherical}}\ignorespaces 

\section{dihedral\_style spherical command}
\label{\detokenize{dihedral_spherical:dihedral-style-spherical-command}}\label{\detokenize{dihedral_spherical:index-0}}\label{\detokenize{dihedral_spherical::doc}}

\subsection{Syntax}
\label{\detokenize{dihedral_spherical:syntax}}
\begin{sphinxVerbatim}[commandchars=\\\{\}]
\PYG{k}{dihedral\PYGZus{}style}\PYG{+w}{ }\PYG{n}{spherical}
\end{sphinxVerbatim}


\subsection{Examples}
\label{\detokenize{dihedral_spherical:examples}}
\begin{sphinxVerbatim}[commandchars=\\\{\}]
\PYG{k}{dihedral\PYGZus{}coeff}\PYG{+w}{ }\PYG{l+m}{1}\PYG{+w}{ }\PYG{l+m}{1}\PYG{+w}{  }\PYG{l+m}{286.1}\PYG{+w}{  }\PYG{l+m}{1}\PYG{+w}{ }\PYG{l+m}{124}\PYG{+w}{  }\PYG{l+m}{1}\PYG{+w}{    }\PYG{l+m}{1}\PYG{+w}{ }\PYG{l+m}{90.0}\PYG{+w}{ }\PYG{l+m}{0}\PYG{+w}{    }\PYG{l+m}{1}\PYG{+w}{ }\PYG{l+m}{90.0}\PYG{+w}{ }\PYG{l+m}{0}
\PYG{k}{dihedral\PYGZus{}coeff}\PYG{+w}{ }\PYG{l+m}{1}\PYG{+w}{ }\PYG{l+m}{3}\PYG{+w}{  }\PYG{l+m}{69.3}\PYG{+w}{   }\PYG{l+m}{1}\PYG{+w}{ }\PYG{l+m}{93.9}\PYG{+w}{ }\PYG{l+m}{1}\PYG{+w}{    }\PYG{l+m}{1}\PYG{+w}{ }\PYG{l+m}{90}\PYG{+w}{   }\PYG{l+m}{0}\PYG{+w}{    }\PYG{l+m}{1}\PYG{+w}{ }\PYG{l+m}{90}\PYG{+w}{   }\PYG{l+m}{0}\PYG{+w}{  }\PYG{o}{\PYGZam{}}
\PYG{+w}{                    }\PYG{l+m}{49.1}\PYG{+w}{   }\PYG{l+m}{0}\PYG{+w}{ }\PYG{l+m}{0.00}\PYG{+w}{ }\PYG{l+m}{0}\PYG{+w}{    }\PYG{l+m}{1}\PYG{+w}{ }\PYG{l+m}{74.4}\PYG{+w}{ }\PYG{l+m}{1}\PYG{+w}{    }\PYG{l+m}{0}\PYG{+w}{ }\PYG{l+m}{0.00}\PYG{+w}{ }\PYG{l+m}{0}\PYG{+w}{  }\PYG{o}{\PYGZam{}}
\PYG{+w}{                    }\PYG{l+m}{25.2}\PYG{+w}{   }\PYG{l+m}{0}\PYG{+w}{ }\PYG{l+m}{0.00}\PYG{+w}{ }\PYG{l+m}{0}\PYG{+w}{    }\PYG{l+m}{0}\PYG{+w}{ }\PYG{l+m}{0.00}\PYG{+w}{ }\PYG{l+m}{0}\PYG{+w}{    }\PYG{l+m}{1}\PYG{+w}{ }\PYG{l+m}{48.1}\PYG{+w}{ }\PYG{l+m}{1}
\end{sphinxVerbatim}


\subsection{Description}
\label{\detokenize{dihedral_spherical:description}}
\sphinxAtStartPar
The \sphinxstyleemphasis{spherical} dihedral style uses the potential:

\noindent{\hspace*{\fill}\sphinxincludegraphics{{dihedral_spherical_angles}.jpg}\hspace*{\fill}}
\begin{equation*}
\begin{split}E(\phi,\theta_1,\theta_2) & = \sum_{i=1}^N\nolimits\ C_i\ \Phi_i(\phi)\ \Theta_{1i}(\theta_1)\ \Theta_{2i}(\theta_2) \\
\Phi_{i}(\phi)            & = u_i - \mathrm{cos}((\phi   - a_i)K_i) \\
\Theta_{1i}(\theta_1)     & = v_i - \mathrm{cos}((\theta_1-b_i)L_i) \\
\Theta_{2i}(\theta_2)     & = w_i - \mathrm{cos}((\theta_2-c_i)M_i)\end{split}
\end{equation*}
\sphinxAtStartPar
For this dihedral style, the energy can be any function that combines the
4\sphinxhyphen{}body dihedral\sphinxhyphen{}angle (\(\phi\)) and the two 3\sphinxhyphen{}body bond\sphinxhyphen{}angles
(\(\theta_1\), \(\theta_2\)).
For this reason, there is usually no need to define 3\sphinxhyphen{}body “angle” forces
separately for the atoms participating in these interactions.
It is probably more efficient to incorporate 3\sphinxhyphen{}body angle forces into
the dihedral interaction even if it requires adding additional terms to
the expansion (as was done in the second example).  A careful choice of
parameters can prevent singularities that occur with traditional
force\sphinxhyphen{}fields whenever theta1 or theta2 approach 0 or 180 degrees.

\sphinxAtStartPar
The last example above corresponds to an interaction with a single energy
minima located near \(\phi=93.9\), \(\theta_1=74.4\),
\(\theta_2=48.1\) degrees, and it remains numerically stable at all
angles (\(\phi\), \(\theta_1\), \(\theta_2\)). In this example,
the coefficients 49.1, and 25.2 can be physically interpreted as the
harmonic spring constants for theta1 and theta2 around their minima.
The coefficient 69.3 is the harmonic spring constant for phi after
division by sin(74.4)*sin(48.1) (the minima positions for theta1 and theta2).

\sphinxAtStartPar
The following coefficients must be defined for each dihedral type via the
{\hyperref[\detokenize{dihedral_coeff::doc}]{\sphinxcrossref{\DUrole{doc}{dihedral\_coeff}}}} command as in the example above, or in
the Dihedral Coeffs section of a data file read by the
{\hyperref[\detokenize{read_data::doc}]{\sphinxcrossref{\DUrole{doc}{read\_data}}}} command:
\begin{itemize}
\item {} 
\sphinxAtStartPar
\(n\) (integer \textgreater{}= 1)

\item {} 
\sphinxAtStartPar
\(C_1\) (energy)

\item {} 
\sphinxAtStartPar
\(K_1\) (typically an integer)

\item {} 
\sphinxAtStartPar
\(a_1\) (degrees)

\item {} 
\sphinxAtStartPar
\(u_1\) (typically 0.0 or 1.0)

\item {} 
\sphinxAtStartPar
\(L_1\) (typically an integer)

\item {} 
\sphinxAtStartPar
\(b_1\) (degrees, typically 0.0 or 90.0)

\item {} 
\sphinxAtStartPar
\(v_1\) (typically 0.0 or 1.0)

\item {} 
\sphinxAtStartPar
\(M_1\) (typically an integer)

\item {} 
\sphinxAtStartPar
\(c_1\) (degrees, typically 0.0 or 90.0)

\item {} 
\sphinxAtStartPar
\(w_1\) (typically 0.0 or 1.0)

\item {} 
\sphinxAtStartPar
{[}…{]}

\item {} 
\sphinxAtStartPar
\(C_n\) (energy)

\item {} 
\sphinxAtStartPar
\(K_n\) (typically an integer)

\item {} 
\sphinxAtStartPar
\(a_n\) (degrees)

\item {} 
\sphinxAtStartPar
\(u_n\) (typically 0.0 or 1.0)

\item {} 
\sphinxAtStartPar
\(L_n\) (typically an integer)

\item {} 
\sphinxAtStartPar
\(b_n\) (degrees, typically 0.0 or 90.0)

\item {} 
\sphinxAtStartPar
\(v_n\) (typically 0.0 or 1.0)

\item {} 
\sphinxAtStartPar
\(M_n\) (typically an integer)

\item {} 
\sphinxAtStartPar
\(c_n\) (degrees, typically 0.0 or 90.0)

\item {} 
\sphinxAtStartPar
\(w_n\) (typically 0.0 or 1.0)

\end{itemize}


\bigskip\hrule\bigskip



\subsection{Restrictions}
\label{\detokenize{dihedral_spherical:restrictions}}
\sphinxAtStartPar
This dihedral style can only be used if LAMMPS was built with the
EXTRA\sphinxhyphen{}MOLECULE package.  See the {\hyperref[\detokenize{Build_package::doc}]{\sphinxcrossref{\DUrole{doc}{Build package}}}} doc
page for more info.


\subsection{Related commands}
\label{\detokenize{dihedral_spherical:related-commands}}
\sphinxAtStartPar
{\hyperref[\detokenize{dihedral_coeff::doc}]{\sphinxcrossref{\DUrole{doc}{dihedral\_coeff}}}}


\subsection{Default}
\label{\detokenize{dihedral_spherical:default}}
\sphinxAtStartPar
none

\sphinxstepscope

\index{dihedral\_style table@\spxentry{dihedral\_style table}}\ignorespaces 
\index{dihedral\_style table/omp@\spxentry{dihedral\_style table/omp}}\ignorespaces 
\index{dihedral\_style table/cut@\spxentry{dihedral\_style table/cut}}\ignorespaces 

\section{dihedral\_style table command}
\label{\detokenize{dihedral_table:dihedral-style-table-command}}\label{\detokenize{dihedral_table:index-2}}\label{\detokenize{dihedral_table:index-1}}\label{\detokenize{dihedral_table:index-0}}\label{\detokenize{dihedral_table::doc}}
\sphinxAtStartPar
Accelerator Variants: \sphinxstyleemphasis{table/omp}


\section{dihedral\_style table/cut command}
\label{\detokenize{dihedral_table:dihedral-style-table-cut-command}}

\subsection{Syntax}
\label{\detokenize{dihedral_table:syntax}}
\begin{sphinxVerbatim}[commandchars=\\\{\}]
\PYG{k}{dihedral\PYGZus{}style}\PYG{+w}{ }\PYG{n}{style}\PYG{+w}{ }\PYG{n}{interpolation}\PYG{+w}{ }\PYG{n}{Ntable}
\end{sphinxVerbatim}
\begin{itemize}
\item {} 
\sphinxAtStartPar
style = \sphinxstyleemphasis{table} or \sphinxstyleemphasis{table/cut}

\item {} 
\sphinxAtStartPar
interpolation = \sphinxstyleemphasis{linear} or \sphinxstyleemphasis{spline} = method of interpolation

\item {} 
\sphinxAtStartPar
Ntable = size of the internal lookup table

\end{itemize}


\subsection{Examples}
\label{\detokenize{dihedral_table:examples}}
\begin{sphinxVerbatim}[commandchars=\\\{\}]
\PYG{k}{dihedral\PYGZus{}style}\PYG{+w}{ }\PYG{n}{table}\PYG{+w}{ }\PYG{n}{spline}\PYG{+w}{ }\PYG{l+m}{400}
\PYG{k}{dihedral\PYGZus{}style}\PYG{+w}{ }\PYG{n}{table}\PYG{+w}{ }\PYG{n}{linear}\PYG{+w}{ }\PYG{l+m}{1000}
\PYG{k}{dihedral\PYGZus{}coeff}\PYG{+w}{ }\PYG{l+m}{1}\PYG{+w}{ }\PYG{n}{file.table}\PYG{+w}{ }\PYG{n}{DIH\PYGZus{}TABLE1}
\PYG{k}{dihedral\PYGZus{}coeff}\PYG{+w}{ }\PYG{l+m}{2}\PYG{+w}{ }\PYG{n}{file.table}\PYG{+w}{ }\PYG{n}{DIH\PYGZus{}TABLE2}

\PYG{k}{dihedral\PYGZus{}style}\PYG{+w}{ }\PYG{n}{table}\PYG{o}{/}\PYG{n}{cut}\PYG{+w}{ }\PYG{n}{spline}\PYG{+w}{ }\PYG{l+m}{400}
\PYG{k}{dihedral\PYGZus{}style}\PYG{+w}{ }\PYG{n}{table}\PYG{o}{/}\PYG{n}{cut}\PYG{+w}{ }\PYG{n}{linear}\PYG{+w}{ }\PYG{l+m}{1000}
\PYG{k}{dihedral\PYGZus{}coeff}\PYG{+w}{ }\PYG{l+m}{1}\PYG{+w}{ }\PYG{n}{aat}\PYG{+w}{ }\PYG{l+m}{1.0}\PYG{+w}{ }\PYG{l+m}{177}\PYG{+w}{ }\PYG{l+m}{180}\PYG{+w}{ }\PYG{n}{file.table}\PYG{+w}{ }\PYG{n}{DIH\PYGZus{}TABLE1}
\PYG{k}{dihedral\PYGZus{}coeff}\PYG{+w}{ }\PYG{l+m}{2}\PYG{+w}{ }\PYG{n}{aat}\PYG{+w}{ }\PYG{l+m}{0.5}\PYG{+w}{ }\PYG{l+m}{170}\PYG{+w}{ }\PYG{l+m}{180}\PYG{+w}{ }\PYG{n}{file.table}\PYG{+w}{ }\PYG{n}{DIH\PYGZus{}TABLE2}
\end{sphinxVerbatim}


\subsection{Description}
\label{\detokenize{dihedral_table:description}}
\sphinxAtStartPar
The \sphinxstyleemphasis{table} and \sphinxstyleemphasis{table/cut} dihedral styles create interpolation tables
of length \sphinxstyleemphasis{Ntable} from dihedral potential and derivative values listed
in a file(s) as a function of the dihedral angle “phi”.  The files are
read by the {\hyperref[\detokenize{dihedral_coeff::doc}]{\sphinxcrossref{\DUrole{doc}{dihedral\_coeff}}}} command.  For
dihedral style \sphinxstyleemphasis{table/cut} additionally an analytic cutoff that is
quadratic in the bond\sphinxhyphen{}angle (theta) is applied in order to regularize
the dihedral interaction.

\sphinxAtStartPar
The interpolation tables are created by fitting cubic splines to the
file values and interpolating energy and derivative values at each of
\sphinxstyleemphasis{Ntable} dihedral angles. During a simulation, these tables are used
to interpolate energy and force values on individual atoms as
needed. The interpolation is done in one of 2 styles: \sphinxstyleemphasis{linear} or
\sphinxstyleemphasis{spline}.

\sphinxAtStartPar
For the \sphinxstyleemphasis{linear} style, the dihedral angle (phi) is used to find 2
surrounding table values from which an energy or its derivative is
computed by linear interpolation.

\sphinxAtStartPar
For the \sphinxstyleemphasis{spline} style, cubic spline coefficients are computed and
stored at each of the \sphinxstyleemphasis{Ntable} evenly\sphinxhyphen{}spaced values in the
interpolated table.  For a given dihedral angle (phi), the appropriate
coefficients are chosen from this list, and a cubic polynomial is used
to compute the energy and the derivative at this angle.

\sphinxAtStartPar
For dihedral style \sphinxstyleemphasis{table} the following coefficients must be defined
for each dihedral type via the {\hyperref[\detokenize{dihedral_coeff::doc}]{\sphinxcrossref{\DUrole{doc}{dihedral\_coeff}}}}
command as in the example above.
\begin{itemize}
\item {} 
\sphinxAtStartPar
filename

\item {} 
\sphinxAtStartPar
keyword

\end{itemize}

\sphinxAtStartPar
The filename specifies a file containing tabulated energy and derivative
values. The keyword specifies which section of the file to read.  The
format of this file is the same for both dihedral styles and described
below.

\sphinxAtStartPar
For dihedral style \sphinxstyleemphasis{table/cut} the following coefficients must be
defined for each dihedral type via the {\hyperref[\detokenize{dihedral_coeff::doc}]{\sphinxcrossref{\DUrole{doc}{dihedral\_coeff}}}} command as in the example above.
\begin{itemize}
\item {} 
\sphinxAtStartPar
style (= aat)

\item {} 
\sphinxAtStartPar
cutoff prefactor (unitless)

\item {} 
\sphinxAtStartPar
cutoff angle1 (degrees)

\item {} 
\sphinxAtStartPar
cutoff angle2 (degrees)

\item {} 
\sphinxAtStartPar
filename

\item {} 
\sphinxAtStartPar
keyword

\end{itemize}

\sphinxAtStartPar
The cutoff dihedral style uses a tabulated dihedral interaction with a
cutoff function:
\begin{equation*}
\begin{split}f(\theta) & = K \qquad\qquad\qquad\qquad\qquad\qquad \theta < \theta_1 \\
f(\theta) & = K \left(1-\frac{(\theta - \theta_1)^2}{(\theta_2 - \theta_1)^2}\right) \qquad \theta_1 < \theta < \theta_2\end{split}
\end{equation*}
\sphinxAtStartPar
The cutoff includes a prefactor \sphinxstyleemphasis{K} to the cutoff function \(f(\theta)\).
While this value would ordinarily be 1, there may be situations where the value
could be different.

\sphinxAtStartPar
The cutoff \(\theta_1\) specifies the angle (in degrees) below which
the dihedral interaction is unmodified, i.e. the cutoff function is 1.

\sphinxAtStartPar
The cutoff function is applied between \(\theta_1\) and
\(\theta_2\), which is the angle at which the cutoff function drops
to zero.  The value of zero effectively “turns off” the dihedral
interaction.

\sphinxAtStartPar
The filename specifies a file containing tabulated energy and derivative
values. The keyword specifies which section of the file to read.  The
format of this file is the same for both dihedral styles and described
below.


\bigskip\hrule\bigskip


\sphinxAtStartPar
Suitable tables for use with this dihedral style can be created using
the Python code in the \sphinxcode{\sphinxupquote{tools/tabulate}} folder of the LAMMPS source
code distribution.

\sphinxAtStartPar
The format of a tabulated file is as follows (without the
parenthesized comments).  It can begin with one or more comment
or blank lines.

\begin{sphinxVerbatim}[commandchars=\\\{\}]
\PYGZsh{} Table of the potential and its negative derivative

DIH\PYGZus{}TABLE1                   (keyword is the first text on line)
N 30 DEGREES                 (N, NOF, DEGREES, RADIANS, CHECKU/F)
                             (blank line)
1 \PYGZhy{}168.0 \PYGZhy{}1.40351172223 0.0423346818422
2 \PYGZhy{}156.0 \PYGZhy{}1.70447981034 0.00811786522531
3 \PYGZhy{}144.0 \PYGZhy{}1.62956100432 \PYGZhy{}0.0184129719987
...
30 180.0 \PYGZhy{}0.707106781187 0.0719306095245

\PYGZsh{} Example 2: table of the potential. Forces omitted

DIH\PYGZus{}TABLE2
N 30 NOF CHECKU testU.dat CHECKF testF.dat

1 \PYGZhy{}168.0 \PYGZhy{}1.40351172223
2 \PYGZhy{}156.0 \PYGZhy{}1.70447981034
3 \PYGZhy{}144.0 \PYGZhy{}1.62956100432
...
30 180.0 \PYGZhy{}0.707106781187
\end{sphinxVerbatim}

\sphinxAtStartPar
A section begins with a non\sphinxhyphen{}blank line whose first character is not a
“\#”; blank lines or lines starting with “\#” can be used as comments
between sections. The first line begins with a keyword which
identifies the section. The line can contain additional text, but the
initial text must match the argument specified in the
{\hyperref[\detokenize{dihedral_coeff::doc}]{\sphinxcrossref{\DUrole{doc}{dihedral\_coeff}}}} command. The next line lists (in
any order) one or more parameters for the table. Each parameter is a
keyword followed by one or more numeric values.

\sphinxAtStartPar
Following a blank line, the next N lines list the tabulated values. On
each line, the first value is the index from 1 to N, the second value is
the angle value, the third value is the energy (in energy units), and
the fourth is \sphinxhyphen{}dE/d(phi) also in energy units). The third term is the
energy of the 4\sphinxhyphen{}atom configuration for the specified angle.  The fourth
term (when present) is the negative derivative of the energy with
respect to the angle (in degrees, or radians depending on whether the
user selected DEGREES or RADIANS).  Thus the units of the last term
are still energy, not force. The dihedral angle values must increase
from one line to the next.

\sphinxAtStartPar
Dihedral table splines are cyclic.  There is no discontinuity at 180
degrees (or at any other angle).  Although in the examples above, the
angles range from \sphinxhyphen{}180 to 180 degrees, in general, the first angle in
the list can have any value (positive, zero, or negative).  However
the \sphinxstyleemphasis{range} of angles represented in the table must be \sphinxstyleemphasis{strictly} less
than 360 degrees (2pi radians) to avoid angle overlap.  (You may not
supply entries in the table for both 180 and \sphinxhyphen{}180, for example.)  If
the user’s table covers only a narrow range of dihedral angles,
strange numerical behavior can occur in the large remaining gap.

\sphinxAtStartPar
\sphinxstylestrong{Parameters:}

\sphinxAtStartPar
The parameter “N” is required and its value is the number of table
entries that follow. Note that this may be different than the N
specified in the {\hyperref[\detokenize{dihedral_style::doc}]{\sphinxcrossref{\DUrole{doc}{dihedral\_style table}}}} command.
Let \sphinxstyleemphasis{Ntable} is the number of table entries requested dihedral\_style
command, and let \sphinxstyleemphasis{Nfile} be the parameter following “N” in the
tabulated file (“30” in the sparse example above).  What LAMMPS does
is a preliminary interpolation by creating splines using the \sphinxstyleemphasis{Nfile}
tabulated values as nodal points.  It uses these to interpolate as
needed to generate energy and derivative values at \sphinxstyleemphasis{Ntable} different
points (which are evenly spaced over a 360 degree range, even if the
angles in the file are not).  The resulting tables of length \sphinxstyleemphasis{Ntable}
are then used as described above, when computing energy and force for
individual dihedral angles and their atoms.  This means that if you
want the interpolation tables of length \sphinxstyleemphasis{Ntable} to match exactly what
is in the tabulated file (with effectively nopreliminary
interpolation), you should set \sphinxstyleemphasis{Ntable} = \sphinxstyleemphasis{Nfile}.  To ensure the
nodal points in the user’s file are aligned with the interpolated
table entries, the angles in the table should be integer multiples of
360/\sphinxstyleemphasis{Ntable} degrees, or 2*PI/\sphinxstyleemphasis{Ntable} radians (depending on your
choice of angle units).

\sphinxAtStartPar
The optional “NOF” keyword allows the user to omit the forces
(negative energy derivatives) from the table file (normally located in
the fourth column).  In their place, forces will be calculated
automatically by differentiating the potential energy function
indicated by the third column of the table (using either linear or
spline interpolation).

\sphinxAtStartPar
The optional “DEGREES” keyword allows the user to specify angles in
degrees instead of radians (default).

\sphinxAtStartPar
The optional “RADIANS” keyword allows the user to specify angles in
radians instead of degrees.  (Note: This changes the way the forces
are scaled in the fourth column of the data file.)

\sphinxAtStartPar
The optional “CHECKU” keyword is followed by a filename.  This allows
the user to save all of the \sphinxstyleemphasis{Ntable} different entries in the
interpolated energy table to a file to make sure that the interpolated
function agrees with the user’s expectations.  (Note: You can
temporarily increase the \sphinxstyleemphasis{Ntable} parameter to a high value for this
purpose.  “\sphinxstyleemphasis{Ntable}“ is explained above.)

\sphinxAtStartPar
The optional “CHECKF” keyword is analogous to the “CHECKU” keyword.
It is followed by a filename, and it allows the user to check the
interpolated force table.  This option is available even if the user
selected the “NOF” option.

\sphinxAtStartPar
Note that one file can contain many sections, each with a tabulated
potential. LAMMPS reads the file section by section until it finds one
that matches the specified keyword.


\bigskip\hrule\bigskip


\sphinxAtStartPar
Styles with a \sphinxstyleemphasis{gpu}, \sphinxstyleemphasis{intel}, \sphinxstyleemphasis{kk}, \sphinxstyleemphasis{omp}, or \sphinxstyleemphasis{opt} suffix are
functionally the same as the corresponding style without the suffix.
They have been optimized to run faster, depending on your available
hardware, as discussed on the {\hyperref[\detokenize{Speed_packages::doc}]{\sphinxcrossref{\DUrole{doc}{Accelerator packages}}}}
page.  The accelerated styles take the same arguments and should
produce the same results, except for round\sphinxhyphen{}off and precision issues.

\sphinxAtStartPar
These accelerated styles are part of the GPU, INTEL, KOKKOS,
OPENMP, and OPT packages, respectively.  They are only enabled if
LAMMPS was built with those packages.  See the {\hyperref[\detokenize{Build_package::doc}]{\sphinxcrossref{\DUrole{doc}{Build package}}}} page for more info.

\sphinxAtStartPar
You can specify the accelerated styles explicitly in your input script
by including their suffix, or you can use the {\hyperref[\detokenize{Run_options::doc}]{\sphinxcrossref{\DUrole{doc}{\sphinxhyphen{}suffix command\sphinxhyphen{}line switch}}}} when you invoke LAMMPS, or you can use the
{\hyperref[\detokenize{suffix::doc}]{\sphinxcrossref{\DUrole{doc}{suffix}}}} command in your input script.

\sphinxAtStartPar
See the {\hyperref[\detokenize{Speed_packages::doc}]{\sphinxcrossref{\DUrole{doc}{Accelerator packages}}}} page for more
instructions on how to use the accelerated styles effectively.


\subsection{Restart, fix\_modify, output, run start/stop, minimize info}
\label{\detokenize{dihedral_table:restart-fix-modify-output-run-start-stop-minimize-info}}
\sphinxAtStartPar
These dihedral styles write the settings for the “dihedral\_style table”
or “dihedral\_style table/cut” command to {\hyperref[\detokenize{restart::doc}]{\sphinxcrossref{\DUrole{doc}{binary restart files}}}}, so a dihedral\_style command does not need to specified in an
input script that reads a restart file.  However, the coefficient
information loaded from the table file(s) is not stored in the restart
file, since it is tabulated in the potential files.  Thus, suitable
dihedral\_coeff commands do need to be specified in the restart input
script after reading the restart file.


\subsection{Restrictions}
\label{\detokenize{dihedral_table:restrictions}}
\sphinxAtStartPar
The \sphinxstyleemphasis{table} dihedral style can only be used if LAMMPS was built with the
MOLECULE package.  The \sphinxstyleemphasis{table/cut} dihedral style can only be used if
LAMMPS was built with the EXTRA\sphinxhyphen{}MOLECULE package. See the
{\hyperref[\detokenize{Build_package::doc}]{\sphinxcrossref{\DUrole{doc}{Build package}}}} doc page for more info.


\subsection{Related commands}
\label{\detokenize{dihedral_table:related-commands}}
\sphinxAtStartPar
{\hyperref[\detokenize{dihedral_coeff::doc}]{\sphinxcrossref{\DUrole{doc}{dihedral\_coeff}}}}


\subsection{Default}
\label{\detokenize{dihedral_table:default}}
\sphinxAtStartPar
none

\sphinxstepscope

\index{dihedral\_style zero@\spxentry{dihedral\_style zero}}\ignorespaces 

\section{dihedral\_style zero command}
\label{\detokenize{dihedral_zero:dihedral-style-zero-command}}\label{\detokenize{dihedral_zero:index-0}}\label{\detokenize{dihedral_zero::doc}}

\subsection{Syntax}
\label{\detokenize{dihedral_zero:syntax}}
\begin{sphinxVerbatim}[commandchars=\\\{\}]
\PYG{k}{dihedral\PYGZus{}style}\PYG{+w}{ }\PYG{n}{zero}\PYG{+w}{ }\PYG{n}{keyword}
\end{sphinxVerbatim}
\begin{itemize}
\item {} 
\sphinxAtStartPar
zero or more keywords may be appended

\item {} 
\sphinxAtStartPar
keyword = \sphinxstyleemphasis{nocoeff}

\end{itemize}


\subsection{Examples}
\label{\detokenize{dihedral_zero:examples}}
\begin{sphinxVerbatim}[commandchars=\\\{\}]
\PYG{k}{dihedral\PYGZus{}style}\PYG{+w}{ }\PYG{n}{zero}
\PYG{k}{dihedral\PYGZus{}style}\PYG{+w}{ }\PYG{n}{zero}\PYG{+w}{ }\PYG{n}{nocoeff}
\PYG{k}{dihedral\PYGZus{}coeff}\PYG{+w}{ }\PYG{o}{*}
\end{sphinxVerbatim}


\subsection{Description}
\label{\detokenize{dihedral_zero:description}}
\sphinxAtStartPar
Using a dihedral style of zero means dihedral forces and energies are
not computed, but the geometry of dihedral quadruplets is still
accessible to other commands.

\sphinxAtStartPar
As an example, the {\hyperref[\detokenize{compute_dihedral_local::doc}]{\sphinxcrossref{\DUrole{doc}{compute dihedral/local}}}} command can be used to
compute the theta values for the list of quadruplets of dihedral atoms
listed in the data file read by the {\hyperref[\detokenize{read_data::doc}]{\sphinxcrossref{\DUrole{doc}{read\_data}}}}
command.  If no dihedral style is defined, this command cannot be
used.

\sphinxAtStartPar
The optional \sphinxstyleemphasis{nocoeff} flag allows to read data files with a DihedralCoeff
section for any dihedral style. Similarly, any dihedral\_coeff commands
will only be checked for the dihedral type number and the rest ignored.

\sphinxAtStartPar
Note that the {\hyperref[\detokenize{dihedral_coeff::doc}]{\sphinxcrossref{\DUrole{doc}{dihedral\_coeff}}}} command must be
used for all dihedral types, though no additional values are
specified.


\subsection{Restrictions}
\label{\detokenize{dihedral_zero:restrictions}}\begin{quote}

\sphinxAtStartPar
none
\end{quote}


\subsection{Related commands}
\label{\detokenize{dihedral_zero:related-commands}}
\sphinxAtStartPar
none

\sphinxAtStartPar
{\hyperref[\detokenize{dihedral_none::doc}]{\sphinxcrossref{\DUrole{doc}{dihedral\_style none}}}}


\subsection{Default}
\label{\detokenize{dihedral_zero:default}}
\sphinxAtStartPar
none

\sphinxstepscope


\chapter{Improper Styles}
\label{\detokenize{impropers:improper-styles}}\label{\detokenize{impropers::doc}}
\sphinxstepscope

\index{improper\_style amoeba@\spxentry{improper\_style amoeba}}\ignorespaces 

\section{improper\_style amoeba command}
\label{\detokenize{improper_amoeba:improper-style-amoeba-command}}\label{\detokenize{improper_amoeba:index-0}}\label{\detokenize{improper_amoeba::doc}}

\subsection{Syntax}
\label{\detokenize{improper_amoeba:syntax}}
\begin{sphinxVerbatim}[commandchars=\\\{\}]
\PYG{k}{improper\PYGZus{}style}\PYG{+w}{ }\PYG{n}{amoeba}
\end{sphinxVerbatim}


\subsection{Examples}
\label{\detokenize{improper_amoeba:examples}}
\begin{sphinxVerbatim}[commandchars=\\\{\}]
\PYG{k}{improper\PYGZus{}style}\PYG{+w}{ }\PYG{n}{amoeba}
\PYG{k}{improper\PYGZus{}coeff}\PYG{+w}{ }\PYG{l+m}{1}\PYG{+w}{ }\PYG{l+m}{49.6}
\end{sphinxVerbatim}


\subsection{Description}
\label{\detokenize{improper_amoeba:description}}
\sphinxAtStartPar
The \sphinxstyleemphasis{amoeba} improper style uses the potential
\begin{equation*}
\begin{split}E = K (\chi)^2\end{split}
\end{equation*}
\sphinxAtStartPar
where \(\chi\) is the improper angle and \(K\) is a prefactor.
Note that the usual 1/2 factor is included in \(K\).

\sphinxAtStartPar
This formula seems like a simplified version of the formula for the
{\hyperref[\detokenize{improper_harmonic::doc}]{\sphinxcrossref{\DUrole{doc}{improper\_style harmonic}}}} command with
\(\chi_0\) = 0.0.  However the computation of the angle
\(\chi\) is done differently to match how the Tinker MD code
computes its out\sphinxhyphen{}of\sphinxhyphen{}plane improper for the AMOEBA and HIPPO force
fields.  See the {\hyperref[\detokenize{Howto_amoeba::doc}]{\sphinxcrossref{\DUrole{doc}{Howto amoeba}}}} doc page for more
information about the implementation of AMOEBA and HIPPO in LAMMPS.

\sphinxAtStartPar
If the 4 atoms in an improper quadruplet (listed in the data file read
by the {\hyperref[\detokenize{read_data::doc}]{\sphinxcrossref{\DUrole{doc}{read\_data}}}} command are ordered I,J,K,L then
atoms I,K,L are considered to lie in a plane and atom J is
out\sphinxhyphen{}of\sphinxhyphen{}place.  The angle \(\chi_0\) is computed as the Allinger
angle which is defined as the angle between the plane of I,K,L, and
the vector from atom I to atom J.

\sphinxAtStartPar
The following coefficient must be defined for each improper type via
the {\hyperref[\detokenize{improper_coeff::doc}]{\sphinxcrossref{\DUrole{doc}{improper\_coeff}}}} command as in the example
above, or in the data file or restart files read by the
{\hyperref[\detokenize{read_data::doc}]{\sphinxcrossref{\DUrole{doc}{read\_data}}}} or {\hyperref[\detokenize{read_restart::doc}]{\sphinxcrossref{\DUrole{doc}{read\_restart}}}}
commands:
\begin{itemize}
\item {} 
\sphinxAtStartPar
\(K\) (energy)

\end{itemize}

\sphinxAtStartPar
Note that the angle \(\chi\) is computed in radians; hence
\(K\) is effectively energy per radian\textasciicircum{}2.


\bigskip\hrule\bigskip



\subsection{Restrictions}
\label{\detokenize{improper_amoeba:restrictions}}
\sphinxAtStartPar
This improper style can only be used if LAMMPS was built with the
AMOEBA package.  See the {\hyperref[\detokenize{Build_package::doc}]{\sphinxcrossref{\DUrole{doc}{Build package}}}} doc page
for more info.


\subsection{Related commands}
\label{\detokenize{improper_amoeba:related-commands}}
\sphinxAtStartPar
{\hyperref[\detokenize{improper_coeff::doc}]{\sphinxcrossref{\DUrole{doc}{improper\_coeff}}}},
{\hyperref[\detokenize{improper_harmonic::doc}]{\sphinxcrossref{\DUrole{doc}{improper\_harmonic}}}}


\subsection{Default}
\label{\detokenize{improper_amoeba:default}}
\sphinxAtStartPar
none

\sphinxstepscope

\index{improper\_style class2@\spxentry{improper\_style class2}}\ignorespaces 
\index{improper\_style class2/omp@\spxentry{improper\_style class2/omp}}\ignorespaces 
\index{improper\_style class2/kk@\spxentry{improper\_style class2/kk}}\ignorespaces 

\section{improper\_style class2 command}
\label{\detokenize{improper_class2:improper-style-class2-command}}\label{\detokenize{improper_class2:index-2}}\label{\detokenize{improper_class2:index-1}}\label{\detokenize{improper_class2:index-0}}\label{\detokenize{improper_class2::doc}}
\sphinxAtStartPar
Accelerator Variants: \sphinxstyleemphasis{class2/omp}, \sphinxstyleemphasis{class2/kk}


\subsection{Syntax}
\label{\detokenize{improper_class2:syntax}}
\begin{sphinxVerbatim}[commandchars=\\\{\}]
\PYG{k}{improper\PYGZus{}style}\PYG{+w}{ }\PYG{n}{class2}
\end{sphinxVerbatim}


\subsection{Examples}
\label{\detokenize{improper_class2:examples}}
\begin{sphinxVerbatim}[commandchars=\\\{\}]
\PYG{k}{improper\PYGZus{}style}\PYG{+w}{ }\PYG{n}{class2}
\PYG{k}{improper\PYGZus{}coeff}\PYG{+w}{ }\PYG{l+m}{1}\PYG{+w}{ }\PYG{l+m}{100.0}\PYG{+w}{ }\PYG{l+m}{0}
\PYG{k}{improper\PYGZus{}coeff}\PYG{+w}{ }\PYG{o}{*}\PYG{+w}{ }\PYG{n}{aa}\PYG{+w}{ }\PYG{l+m}{0.0}\PYG{+w}{ }\PYG{l+m}{0.0}\PYG{+w}{ }\PYG{l+m}{0.0}\PYG{+w}{ }\PYG{l+m}{115.06}\PYG{+w}{ }\PYG{l+m}{130.01}\PYG{+w}{ }\PYG{l+m}{115.06}
\end{sphinxVerbatim}


\subsection{Description}
\label{\detokenize{improper_class2:description}}
\sphinxAtStartPar
The \sphinxstyleemphasis{class2} improper style uses the potential
\begin{equation*}
\begin{split}E      = & E_i + E_{aa} \\
E_i    = & K [ \frac{\chi_{ijkl} + \chi_{kjli} + \chi_{ljik}}{3} - \chi_0 ]^2 \\
E_{aa} = & M_1 (\theta_{ijk} - \theta_1) (\theta_{kjl} - \theta_3) + \\
         & M_2 (\theta_{ijk} - \theta_1) (\theta_{ijl} - \theta_2) + \\
         & M_3 (\theta_{ijl} - \theta_2) (\theta_{kjl} - \theta_3)\end{split}
\end{equation*}
\sphinxAtStartPar
where \(E_i\) is the improper term and \(E_{aa}\) is an
angle\sphinxhyphen{}angle term.  The 3 \(\chi\) terms in \(E_i\) are an
average over 3 out\sphinxhyphen{}of\sphinxhyphen{}plane angles.

\sphinxAtStartPar
The 4 atoms in an improper quadruplet (listed in the data file read by
the {\hyperref[\detokenize{read_data::doc}]{\sphinxcrossref{\DUrole{doc}{read\_data}}}} command) are ordered I,J,K,L.
\(\chi_{ijkl}\) refers to the angle between the plane of I,J,K and
the plane of J,K,L, and the bond JK lies in both planes.  Similarly for
\(\chi_{kjli}\) and \(\chi_{ljik}\).
Note that atom J appears in the common bonds (JI, JK, JL) of all 3 X
terms.  Thus J (the second atom in the quadruplet) is the atom of
symmetry in the 3 \(\chi\) angles.

\sphinxAtStartPar
The subscripts on the various \(\theta\)s refer to different
combinations of three atoms (I,J,K,L) used to form a particular angle.
E.g. \(\theta_{ijl}\) is the angle formed by atoms I,J,L with J
in the middle.  \(\theta_1\), \(\theta_2\), \(\theta_3\)
are the equilibrium positions of those angles.  Again,
atom J (the second atom in the quadruplet) is the atom of symmetry in the
theta angles, since it is always the center atom.

\sphinxAtStartPar
Since atom J is the atom of symmetry, normally the bonds J\sphinxhyphen{}I, J\sphinxhyphen{}K, J\sphinxhyphen{}L
would exist for an improper to be defined between the 4 atoms, but
this is not required.

\sphinxAtStartPar
See {\hyperref[\detokenize{improper_class2:improper-sun}]{\sphinxcrossref{\DUrole{std,std-ref}{(Sun)}}}} for a description of the COMPASS class2 force field.

\sphinxAtStartPar
Coefficients for the \(E_i\) and \(E_{aa}\) formulas must be
defined for each
improper type via the {\hyperref[\detokenize{improper_coeff::doc}]{\sphinxcrossref{\DUrole{doc}{improper\_coeff}}}} command as
in the example above, or in the data file or restart files read by the
{\hyperref[\detokenize{read_data::doc}]{\sphinxcrossref{\DUrole{doc}{read\_data}}}} or {\hyperref[\detokenize{read_restart::doc}]{\sphinxcrossref{\DUrole{doc}{read\_restart}}}}
commands.

\sphinxAtStartPar
These are the 2 coefficients for the \(E_i\) formula:
\begin{itemize}
\item {} 
\sphinxAtStartPar
\(K\) (energy)

\item {} 
\sphinxAtStartPar
\(\chi_0\) (degrees)

\end{itemize}

\sphinxAtStartPar
\(\chi_0\) is specified in degrees, but LAMMPS converts it to
radians internally; hence \(K\) is effectively energy per
radian\textasciicircum{}2.

\sphinxAtStartPar
For the \(E_{aa}\) formula, each line in a {\hyperref[\detokenize{improper_coeff::doc}]{\sphinxcrossref{\DUrole{doc}{improper\_coeff}}}} command in the input script lists 7 coefficients,
the first of which is \sphinxstyleemphasis{aa} to indicate they are AngleAngle
coefficients.  In a data file, these coefficients should be listed
under a \sphinxstyleemphasis{AngleAngle Coeffs} heading and you must leave out the \sphinxstyleemphasis{aa},
i.e. only list 6 coefficients after the improper type.
\begin{itemize}
\item {} 
\sphinxAtStartPar
\sphinxstyleemphasis{aa}

\item {} 
\sphinxAtStartPar
\(M_1\) (energy)

\item {} 
\sphinxAtStartPar
\(M_2\) (energy)

\item {} 
\sphinxAtStartPar
\(M_3\) (energy)

\item {} 
\sphinxAtStartPar
\(\theta_1\) (degrees)

\item {} 
\sphinxAtStartPar
\(\theta_2\) (degrees)

\item {} 
\sphinxAtStartPar
\(\theta_3\) (degrees)

\end{itemize}

\sphinxAtStartPar
The \(\theta\) values are specified in degrees, but LAMMPS
converts them to radians internally; hence the hence the various
\(M\) are effectively energy per radian\textasciicircum{}2.


\bigskip\hrule\bigskip


\sphinxAtStartPar
Styles with a \sphinxstyleemphasis{gpu}, \sphinxstyleemphasis{intel}, \sphinxstyleemphasis{kk}, \sphinxstyleemphasis{omp}, or \sphinxstyleemphasis{opt} suffix are
functionally the same as the corresponding style without the suffix.
They have been optimized to run faster, depending on your available
hardware, as discussed on the {\hyperref[\detokenize{Speed_packages::doc}]{\sphinxcrossref{\DUrole{doc}{Accelerator packages}}}}
page.  The accelerated styles take the same arguments and should
produce the same results, except for round\sphinxhyphen{}off and precision issues.

\sphinxAtStartPar
These accelerated styles are part of the GPU, INTEL, KOKKOS,
OPENMP, and OPT packages, respectively.  They are only enabled if
LAMMPS was built with those packages.  See the {\hyperref[\detokenize{Build_package::doc}]{\sphinxcrossref{\DUrole{doc}{Build package}}}} page for more info.

\sphinxAtStartPar
You can specify the accelerated styles explicitly in your input script
by including their suffix, or you can use the {\hyperref[\detokenize{Run_options::doc}]{\sphinxcrossref{\DUrole{doc}{\sphinxhyphen{}suffix command\sphinxhyphen{}line switch}}}} when you invoke LAMMPS, or you can use the
{\hyperref[\detokenize{suffix::doc}]{\sphinxcrossref{\DUrole{doc}{suffix}}}} command in your input script.

\sphinxAtStartPar
See the {\hyperref[\detokenize{Speed_packages::doc}]{\sphinxcrossref{\DUrole{doc}{Accelerator packages}}}} page for more
instructions on how to use the accelerated styles effectively.


\bigskip\hrule\bigskip



\subsection{Restrictions}
\label{\detokenize{improper_class2:restrictions}}
\sphinxAtStartPar
This improper style can only be used if LAMMPS was built with the
CLASS2 package.  See the {\hyperref[\detokenize{Build_package::doc}]{\sphinxcrossref{\DUrole{doc}{Build package}}}} doc
page for more info.


\subsection{Related commands}
\label{\detokenize{improper_class2:related-commands}}
\sphinxAtStartPar
{\hyperref[\detokenize{improper_coeff::doc}]{\sphinxcrossref{\DUrole{doc}{improper\_coeff}}}}


\subsection{Default}
\label{\detokenize{improper_class2:default}}
\sphinxAtStartPar
none


\bigskip\hrule\bigskip

\phantomsection\label{\detokenize{improper_class2:improper-sun}}
\sphinxAtStartPar
\sphinxstylestrong{(Sun)} Sun, J Phys Chem B 102, 7338\sphinxhyphen{}7364 (1998).

\sphinxstepscope

\index{improper\_style cossq@\spxentry{improper\_style cossq}}\ignorespaces 
\index{improper\_style cossq/omp@\spxentry{improper\_style cossq/omp}}\ignorespaces 

\section{improper\_style cossq command}
\label{\detokenize{improper_cossq:improper-style-cossq-command}}\label{\detokenize{improper_cossq:index-1}}\label{\detokenize{improper_cossq:index-0}}\label{\detokenize{improper_cossq::doc}}
\sphinxAtStartPar
Accelerator Variants: \sphinxstyleemphasis{cossq/omp}


\subsection{Syntax}
\label{\detokenize{improper_cossq:syntax}}
\begin{sphinxVerbatim}[commandchars=\\\{\}]
\PYG{k}{improper\PYGZus{}style}\PYG{+w}{ }\PYG{n}{cossq}
\end{sphinxVerbatim}


\subsection{Examples}
\label{\detokenize{improper_cossq:examples}}
\begin{sphinxVerbatim}[commandchars=\\\{\}]
\PYG{k}{improper\PYGZus{}style}\PYG{+w}{ }\PYG{n}{cossq}
\PYG{k}{improper\PYGZus{}coeff}\PYG{+w}{ }\PYG{l+m}{1}\PYG{+w}{ }\PYG{l+m}{4.0}\PYG{+w}{ }\PYG{l+m}{0.0}
\end{sphinxVerbatim}


\subsection{Description}
\label{\detokenize{improper_cossq:description}}
\sphinxAtStartPar
The \sphinxstyleemphasis{cossq} improper style uses the potential
\begin{equation*}
\begin{split}E = \frac{1}{2} K \cos^2{\left(\chi - \chi_0\right)}\end{split}
\end{equation*}
\sphinxAtStartPar
where \(\chi\) is the improper angle, \(\chi_0\) is its
equilibrium value, and \(K\) is a prefactor.

\sphinxAtStartPar
If the 4 atoms in an improper quadruplet (listed in the data file read
by the {\hyperref[\detokenize{read_data::doc}]{\sphinxcrossref{\DUrole{doc}{read\_data}}}} command) are ordered I,J,K,L then
\(\chi\) is the angle between the plane of I,J,K and the plane of J,K,L.
Alternatively, you can think of atoms J,K,L as being in a plane, and
atom I above the plane, and \(\chi\) as a measure of how far
out\sphinxhyphen{}of\sphinxhyphen{}plane I is with respect to the other 3 atoms.

\sphinxAtStartPar
Note that defining 4 atoms to interact in this way, does not mean that
bonds necessarily exist between I\sphinxhyphen{}J, J\sphinxhyphen{}K, or K\sphinxhyphen{}L, as they would in a
linear dihedral.  Normally, the bonds I\sphinxhyphen{}J, I\sphinxhyphen{}K, I\sphinxhyphen{}L would exist for an
improper to be defined between the 4 atoms.

\sphinxAtStartPar
The following coefficients must be defined for each improper type via
the {\hyperref[\detokenize{improper_coeff::doc}]{\sphinxcrossref{\DUrole{doc}{improper\_coeff}}}} command as in the example
above, or in the data file or restart files read by the
{\hyperref[\detokenize{read_data::doc}]{\sphinxcrossref{\DUrole{doc}{read\_data}}}} or {\hyperref[\detokenize{read_restart::doc}]{\sphinxcrossref{\DUrole{doc}{read\_restart}}}}
commands:
\begin{itemize}
\item {} 
\sphinxAtStartPar
\(K\) (energy)

\item {} 
\sphinxAtStartPar
\(\chi_0\) (degrees)

\end{itemize}


\bigskip\hrule\bigskip


\sphinxAtStartPar
Styles with a \sphinxstyleemphasis{gpu}, \sphinxstyleemphasis{intel}, \sphinxstyleemphasis{kk}, \sphinxstyleemphasis{omp}, or \sphinxstyleemphasis{opt} suffix are
functionally the same as the corresponding style without the suffix.
They have been optimized to run faster, depending on your available
hardware, as discussed on the {\hyperref[\detokenize{Speed_packages::doc}]{\sphinxcrossref{\DUrole{doc}{Accelerator packages}}}}
page.  The accelerated styles take the same arguments and should
produce the same results, except for round\sphinxhyphen{}off and precision issues.

\sphinxAtStartPar
These accelerated styles are part of the GPU, INTEL, KOKKOS,
OPENMP, and OPT packages, respectively.  They are only enabled if
LAMMPS was built with those packages.  See the {\hyperref[\detokenize{Build_package::doc}]{\sphinxcrossref{\DUrole{doc}{Build package}}}} page for more info.

\sphinxAtStartPar
You can specify the accelerated styles explicitly in your input script
by including their suffix, or you can use the {\hyperref[\detokenize{Run_options::doc}]{\sphinxcrossref{\DUrole{doc}{\sphinxhyphen{}suffix command\sphinxhyphen{}line switch}}}} when you invoke LAMMPS, or you can use the
{\hyperref[\detokenize{suffix::doc}]{\sphinxcrossref{\DUrole{doc}{suffix}}}} command in your input script.

\sphinxAtStartPar
See the {\hyperref[\detokenize{Speed_packages::doc}]{\sphinxcrossref{\DUrole{doc}{Accelerator packages}}}} page for more
instructions on how to use the accelerated styles effectively.


\bigskip\hrule\bigskip



\subsection{Restrictions}
\label{\detokenize{improper_cossq:restrictions}}
\sphinxAtStartPar
This improper style can only be used if LAMMPS was built with the
EXTRA\sphinxhyphen{}MOLECULE package.  See the {\hyperref[\detokenize{Build_package::doc}]{\sphinxcrossref{\DUrole{doc}{Build package}}}}
doc page for more info.


\subsection{Related commands}
\label{\detokenize{improper_cossq:related-commands}}
\sphinxAtStartPar
{\hyperref[\detokenize{improper_coeff::doc}]{\sphinxcrossref{\DUrole{doc}{improper\_coeff}}}}


\subsection{Default}
\label{\detokenize{improper_cossq:default}}
\sphinxAtStartPar
none

\sphinxstepscope

\index{improper\_style cvff@\spxentry{improper\_style cvff}}\ignorespaces 
\index{improper\_style cvff/intel@\spxentry{improper\_style cvff/intel}}\ignorespaces 
\index{improper\_style cvff/omp@\spxentry{improper\_style cvff/omp}}\ignorespaces 

\section{improper\_style cvff command}
\label{\detokenize{improper_cvff:improper-style-cvff-command}}\label{\detokenize{improper_cvff:index-2}}\label{\detokenize{improper_cvff:index-1}}\label{\detokenize{improper_cvff:index-0}}\label{\detokenize{improper_cvff::doc}}
\sphinxAtStartPar
Accelerator Variants: \sphinxstyleemphasis{cvff/intel}, \sphinxstyleemphasis{cvff/omp}


\subsection{Syntax}
\label{\detokenize{improper_cvff:syntax}}
\begin{sphinxVerbatim}[commandchars=\\\{\}]
\PYG{k}{improper\PYGZus{}style}\PYG{+w}{ }\PYG{n}{cvff}
\end{sphinxVerbatim}


\subsection{Examples}
\label{\detokenize{improper_cvff:examples}}
\begin{sphinxVerbatim}[commandchars=\\\{\}]
\PYG{k}{improper\PYGZus{}style}\PYG{+w}{ }\PYG{n}{cvff}
\PYG{k}{improper\PYGZus{}coeff}\PYG{+w}{ }\PYG{l+m}{1}\PYG{+w}{ }\PYG{l+m}{80.0}\PYG{+w}{ }\PYG{o}{\PYGZhy{}}\PYG{l+m}{1}\PYG{+w}{ }\PYG{l+m}{4}
\end{sphinxVerbatim}


\subsection{Description}
\label{\detokenize{improper_cvff:description}}
\sphinxAtStartPar
The \sphinxstyleemphasis{cvff} improper style uses the potential
\begin{equation*}
\begin{split}E = K [1 + d  \cos (n \phi) ]\end{split}
\end{equation*}
\sphinxAtStartPar
where phi is the improper dihedral angle.

\sphinxAtStartPar
If the 4 atoms in an improper quadruplet (listed in the data file read
by the {\hyperref[\detokenize{read_data::doc}]{\sphinxcrossref{\DUrole{doc}{read\_data}}}} command) are ordered I,J,K,L then
the improper dihedral angle is between the plane of I,J,K and the
plane of J,K,L.  Note that because this is effectively a dihedral
angle, the formula for this improper style is the same as for
{\hyperref[\detokenize{dihedral_harmonic::doc}]{\sphinxcrossref{\DUrole{doc}{dihedral\_style harmonic}}}}.

\sphinxAtStartPar
Note that defining 4 atoms to interact in this way, does not mean that
bonds necessarily exist between I\sphinxhyphen{}J, J\sphinxhyphen{}K, or K\sphinxhyphen{}L, as they would in a
linear dihedral.  Normally, the bonds I\sphinxhyphen{}J, I\sphinxhyphen{}K, I\sphinxhyphen{}L would exist for an
improper to be defined between the 4 atoms.

\sphinxAtStartPar
The following coefficients must be defined for each improper type via
the {\hyperref[\detokenize{improper_coeff::doc}]{\sphinxcrossref{\DUrole{doc}{improper\_coeff}}}} command as in the example
above, or in the data file or restart files read by the
{\hyperref[\detokenize{read_data::doc}]{\sphinxcrossref{\DUrole{doc}{read\_data}}}} or {\hyperref[\detokenize{read_restart::doc}]{\sphinxcrossref{\DUrole{doc}{read\_restart}}}}
commands:
\begin{itemize}
\item {} 
\sphinxAtStartPar
\(K\) (energy)

\item {} 
\sphinxAtStartPar
\(d\) (+1 or \sphinxhyphen{}1)

\item {} 
\sphinxAtStartPar
\(n\) (0,1,2,3,4,6)

\end{itemize}


\bigskip\hrule\bigskip


\sphinxAtStartPar
Styles with a \sphinxstyleemphasis{gpu}, \sphinxstyleemphasis{intel}, \sphinxstyleemphasis{kk}, \sphinxstyleemphasis{omp}, or \sphinxstyleemphasis{opt} suffix are
functionally the same as the corresponding style without the suffix.
They have been optimized to run faster, depending on your available
hardware, as discussed on the {\hyperref[\detokenize{Speed_packages::doc}]{\sphinxcrossref{\DUrole{doc}{Accelerator packages}}}}
page.  The accelerated styles take the same arguments and should
produce the same results, except for round\sphinxhyphen{}off and precision issues.

\sphinxAtStartPar
These accelerated styles are part of the GPU, INTEL, KOKKOS,
OPENMP, and OPT packages, respectively.  They are only enabled if
LAMMPS was built with those packages.  See the {\hyperref[\detokenize{Build_package::doc}]{\sphinxcrossref{\DUrole{doc}{Build package}}}} page for more info.

\sphinxAtStartPar
You can specify the accelerated styles explicitly in your input script
by including their suffix, or you can use the {\hyperref[\detokenize{Run_options::doc}]{\sphinxcrossref{\DUrole{doc}{\sphinxhyphen{}suffix command\sphinxhyphen{}line switch}}}} when you invoke LAMMPS, or you can use the
{\hyperref[\detokenize{suffix::doc}]{\sphinxcrossref{\DUrole{doc}{suffix}}}} command in your input script.

\sphinxAtStartPar
See the {\hyperref[\detokenize{Speed_packages::doc}]{\sphinxcrossref{\DUrole{doc}{Accelerator packages}}}} page for more
instructions on how to use the accelerated styles effectively.


\bigskip\hrule\bigskip



\subsection{Restrictions}
\label{\detokenize{improper_cvff:restrictions}}
\sphinxAtStartPar
This improper style can only be used if LAMMPS was built with the
MOLECULE package.  See the {\hyperref[\detokenize{Build_package::doc}]{\sphinxcrossref{\DUrole{doc}{Build package}}}} doc page
for more info.


\subsection{Related commands}
\label{\detokenize{improper_cvff:related-commands}}
\sphinxAtStartPar
{\hyperref[\detokenize{improper_coeff::doc}]{\sphinxcrossref{\DUrole{doc}{improper\_coeff}}}}


\subsection{Default}
\label{\detokenize{improper_cvff:default}}
\sphinxAtStartPar
none

\sphinxstepscope

\index{improper\_style distance@\spxentry{improper\_style distance}}\ignorespaces 

\section{improper\_style distance command}
\label{\detokenize{improper_distance:improper-style-distance-command}}\label{\detokenize{improper_distance:index-0}}\label{\detokenize{improper_distance::doc}}

\subsection{Syntax}
\label{\detokenize{improper_distance:syntax}}
\begin{sphinxVerbatim}[commandchars=\\\{\}]
\PYG{k}{improper\PYGZus{}style}\PYG{+w}{ }\PYG{n}{distance}
\end{sphinxVerbatim}


\subsection{Examples}
\label{\detokenize{improper_distance:examples}}
\begin{sphinxVerbatim}[commandchars=\\\{\}]
\PYG{k}{improper\PYGZus{}style}\PYG{+w}{ }\PYG{n}{distance}
\PYG{k}{improper\PYGZus{}coeff}\PYG{+w}{ }\PYG{l+m}{1}\PYG{+w}{ }\PYG{l+m}{80.0}\PYG{+w}{ }\PYG{l+m}{100.0}
\end{sphinxVerbatim}


\subsection{Description}
\label{\detokenize{improper_distance:description}}
\sphinxAtStartPar
The \sphinxstyleemphasis{distance} improper style uses the potential
\begin{equation*}
\begin{split}E = K_2 d^2 + K_4 d^4\end{split}
\end{equation*}
\sphinxAtStartPar
where \(d\) is the distance between the central atom and the plane formed
by the other three atoms.  If the 4 atoms in an improper quadruplet
(listed in the data file read by the {\hyperref[\detokenize{read_data::doc}]{\sphinxcrossref{\DUrole{doc}{read\_data}}}}
command) are ordered I,J,K,L then the I\sphinxhyphen{}atom is assumed to be the
central atom.

\noindent{\hspace*{\fill}\sphinxincludegraphics{{improper_distance}.jpg}\hspace*{\fill}}

\sphinxAtStartPar
Note that defining 4 atoms to interact in this way, does not mean that
bonds necessarily exist between I\sphinxhyphen{}J, J\sphinxhyphen{}K, or K\sphinxhyphen{}L, as they would in a
linear dihedral. Normally, the bonds I\sphinxhyphen{}J, I\sphinxhyphen{}K, I\sphinxhyphen{}L would exist for an
improper to be defined between the 4 atoms.

\sphinxAtStartPar
The following coefficients must be defined for each improper type via
the improper\_coeff command as in the example above, or in the data
file or restart files read by the read\_data or read\_restart commands:
\begin{itemize}
\item {} 
\sphinxAtStartPar
\(K_2\) (energy/distance\textasciicircum{}2)

\item {} 
\sphinxAtStartPar
\(K_4\) (energy/distance\textasciicircum{}4)

\end{itemize}


\bigskip\hrule\bigskip



\subsection{Restrictions}
\label{\detokenize{improper_distance:restrictions}}
\sphinxAtStartPar
This improper style can only be used if LAMMPS was built with the
EXTRA\sphinxhyphen{}MOLECULE package.  See the {\hyperref[\detokenize{Build_package::doc}]{\sphinxcrossref{\DUrole{doc}{Build package}}}}
doc page for more info.


\subsection{Related commands}
\label{\detokenize{improper_distance:related-commands}}
\sphinxAtStartPar
{\hyperref[\detokenize{improper_coeff::doc}]{\sphinxcrossref{\DUrole{doc}{improper\_coeff}}}}


\subsection{Default}
\label{\detokenize{improper_distance:default}}
\sphinxAtStartPar
none

\sphinxstepscope

\index{improper\_style distharm@\spxentry{improper\_style distharm}}\ignorespaces 

\section{improper\_style distharm command}
\label{\detokenize{improper_distharm:improper-style-distharm-command}}\label{\detokenize{improper_distharm:index-0}}\label{\detokenize{improper_distharm::doc}}

\subsection{Syntax}
\label{\detokenize{improper_distharm:syntax}}
\begin{sphinxVerbatim}[commandchars=\\\{\}]
\PYG{k}{improper\PYGZus{}style}\PYG{+w}{ }\PYG{n}{distharm}
\end{sphinxVerbatim}


\subsection{Examples}
\label{\detokenize{improper_distharm:examples}}
\begin{sphinxVerbatim}[commandchars=\\\{\}]
\PYG{k}{improper\PYGZus{}style}\PYG{+w}{ }\PYG{n}{distharm}
\PYG{k}{improper\PYGZus{}coeff}\PYG{+w}{ }\PYG{l+m}{1}\PYG{+w}{ }\PYG{l+m}{25.0}\PYG{+w}{ }\PYG{l+m}{0.5}
\end{sphinxVerbatim}


\subsection{Description}
\label{\detokenize{improper_distharm:description}}
\sphinxAtStartPar
The \sphinxstyleemphasis{distharm} improper style uses the potential
\begin{equation*}
\begin{split}E = K (d - d_0)^2\end{split}
\end{equation*}
\sphinxAtStartPar
where \(d\) is the oriented distance between the central atom and the plane formed
by the other three atoms.  If the 4 atoms in an improper quadruplet
(listed in the data file read by the {\hyperref[\detokenize{read_data::doc}]{\sphinxcrossref{\DUrole{doc}{read\_data}}}}
command) are ordered I,J,K,L then the L\sphinxhyphen{}atom is assumed to be the
central atom. Note that this is different from the convention used
in the improper\_style distance. The distance \(d\) is oriented and can take
on negative values. This may lead to unwanted behavior if \(d_0\) is not equal to zero.

\sphinxAtStartPar
The following coefficients must be defined for each improper type via
the improper\_coeff command as in the example above, or in the data
file or restart files read by the read\_data or read\_restart commands:
\begin{itemize}
\item {} 
\sphinxAtStartPar
\(K\) (energy/distance\textasciicircum{}2)

\item {} 
\sphinxAtStartPar
\(d_0\) (distance)

\end{itemize}


\bigskip\hrule\bigskip



\subsection{Restrictions}
\label{\detokenize{improper_distharm:restrictions}}
\sphinxAtStartPar
This improper style can only be used if LAMMPS was built with the
YAFF package.  See the {\hyperref[\detokenize{Build_package::doc}]{\sphinxcrossref{\DUrole{doc}{Build package}}}} doc
page for more info.


\subsection{Related commands}
\label{\detokenize{improper_distharm:related-commands}}
\sphinxAtStartPar
{\hyperref[\detokenize{improper_coeff::doc}]{\sphinxcrossref{\DUrole{doc}{improper\_coeff}}}}


\subsection{Default}
\label{\detokenize{improper_distharm:default}}
\sphinxAtStartPar
none

\sphinxstepscope

\index{improper\_style fourier@\spxentry{improper\_style fourier}}\ignorespaces 
\index{improper\_style fourier/omp@\spxentry{improper\_style fourier/omp}}\ignorespaces 

\section{improper\_style fourier command}
\label{\detokenize{improper_fourier:improper-style-fourier-command}}\label{\detokenize{improper_fourier:index-1}}\label{\detokenize{improper_fourier:index-0}}\label{\detokenize{improper_fourier::doc}}
\sphinxAtStartPar
Accelerator Variants: \sphinxstyleemphasis{fourier/omp}


\subsection{Syntax}
\label{\detokenize{improper_fourier:syntax}}
\begin{sphinxVerbatim}[commandchars=\\\{\}]
\PYG{k}{improper\PYGZus{}style}\PYG{+w}{ }\PYG{n}{fourier}
\end{sphinxVerbatim}


\subsection{Examples}
\label{\detokenize{improper_fourier:examples}}
\begin{sphinxVerbatim}[commandchars=\\\{\}]
\PYG{k}{improper\PYGZus{}style}\PYG{+w}{ }\PYG{n}{fourier}
\PYG{k}{improper\PYGZus{}coeff}\PYG{+w}{ }\PYG{l+m}{1}\PYG{+w}{ }\PYG{l+m}{100.0}\PYG{+w}{ }\PYG{l+m}{0.0}\PYG{+w}{ }\PYG{l+m}{1.0}\PYG{+w}{ }\PYG{l+m}{0.5}\PYG{+w}{ }\PYG{l+m}{1}
\end{sphinxVerbatim}


\subsection{Description}
\label{\detokenize{improper_fourier:description}}
\sphinxAtStartPar
The \sphinxstyleemphasis{fourier} improper style uses the following potential:
\begin{equation*}
\begin{split}E = K [C_0 + C_1 \cos ( \omega) + C_2 \cos( 2 \omega) ]\end{split}
\end{equation*}
\sphinxAtStartPar
where K is the force constant, C0, C1, C2 are dimensionless coefficients,
and omega is the angle between the IL axis and the IJK plane:

\noindent{\hspace*{\fill}\sphinxincludegraphics{{umbrella}.jpg}\hspace*{\fill}}

\sphinxAtStartPar
If all parameter (see below) is not zero, the all the three possible angles will taken in account.

\sphinxAtStartPar
The following coefficients must be defined for each improper type via
the {\hyperref[\detokenize{improper_coeff::doc}]{\sphinxcrossref{\DUrole{doc}{improper\_coeff}}}} command as in the example
above, or in the data file or restart files read by the
{\hyperref[\detokenize{read_data::doc}]{\sphinxcrossref{\DUrole{doc}{read\_data}}}} or {\hyperref[\detokenize{read_restart::doc}]{\sphinxcrossref{\DUrole{doc}{read\_restart}}}}
commands:
\begin{itemize}
\item {} 
\sphinxAtStartPar
\(K\) (energy)

\item {} 
\sphinxAtStartPar
\(C_0\) (unitless)

\item {} 
\sphinxAtStartPar
\(C_1\) (unitless)

\item {} 
\sphinxAtStartPar
\(C_2\) (unitless)

\item {} 
\sphinxAtStartPar
all  (0 or 1, optional)

\end{itemize}


\bigskip\hrule\bigskip


\sphinxAtStartPar
Styles with a \sphinxstyleemphasis{gpu}, \sphinxstyleemphasis{intel}, \sphinxstyleemphasis{kk}, \sphinxstyleemphasis{omp}, or \sphinxstyleemphasis{opt} suffix are
functionally the same as the corresponding style without the suffix.
They have been optimized to run faster, depending on your available
hardware, as discussed on the {\hyperref[\detokenize{Speed_packages::doc}]{\sphinxcrossref{\DUrole{doc}{Accelerator packages}}}}
page.  The accelerated styles take the same arguments and should
produce the same results, except for round\sphinxhyphen{}off and precision issues.

\sphinxAtStartPar
These accelerated styles are part of the GPU, INTEL, KOKKOS,
OPENMP, and OPT packages, respectively.  They are only enabled if
LAMMPS was built with those packages.  See the {\hyperref[\detokenize{Build_package::doc}]{\sphinxcrossref{\DUrole{doc}{Build package}}}} page for more info.

\sphinxAtStartPar
You can specify the accelerated styles explicitly in your input script
by including their suffix, or you can use the {\hyperref[\detokenize{Run_options::doc}]{\sphinxcrossref{\DUrole{doc}{\sphinxhyphen{}suffix command\sphinxhyphen{}line switch}}}} when you invoke LAMMPS, or you can use the
{\hyperref[\detokenize{suffix::doc}]{\sphinxcrossref{\DUrole{doc}{suffix}}}} command in your input script.

\sphinxAtStartPar
See the {\hyperref[\detokenize{Speed_packages::doc}]{\sphinxcrossref{\DUrole{doc}{Accelerator packages}}}} page for more
instructions on how to use the accelerated styles effectively.


\bigskip\hrule\bigskip



\subsection{Restrictions}
\label{\detokenize{improper_fourier:restrictions}}
\sphinxAtStartPar
This angle style can only be used if LAMMPS was built with the
EXTRA\sphinxhyphen{}MOLECULE package.  See the {\hyperref[\detokenize{Build_package::doc}]{\sphinxcrossref{\DUrole{doc}{Build package}}}}
doc page for more info.


\subsection{Related commands}
\label{\detokenize{improper_fourier:related-commands}}
\sphinxAtStartPar
{\hyperref[\detokenize{improper_coeff::doc}]{\sphinxcrossref{\DUrole{doc}{improper\_coeff}}}}


\subsection{Default}
\label{\detokenize{improper_fourier:default}}
\sphinxAtStartPar
none

\sphinxstepscope

\index{improper\_style harmonic@\spxentry{improper\_style harmonic}}\ignorespaces 
\index{improper\_style harmonic/intel@\spxentry{improper\_style harmonic/intel}}\ignorespaces 
\index{improper\_style harmonic/kk@\spxentry{improper\_style harmonic/kk}}\ignorespaces 
\index{improper\_style harmonic/omp@\spxentry{improper\_style harmonic/omp}}\ignorespaces 

\section{improper\_style harmonic command}
\label{\detokenize{improper_harmonic:improper-style-harmonic-command}}\label{\detokenize{improper_harmonic:index-3}}\label{\detokenize{improper_harmonic:index-2}}\label{\detokenize{improper_harmonic:index-1}}\label{\detokenize{improper_harmonic:index-0}}\label{\detokenize{improper_harmonic::doc}}
\sphinxAtStartPar
Accelerator Variants: \sphinxstyleemphasis{harmonic/intel}, \sphinxstyleemphasis{harmonic/kk}, \sphinxstyleemphasis{harmonic/omp}


\subsection{Syntax}
\label{\detokenize{improper_harmonic:syntax}}
\begin{sphinxVerbatim}[commandchars=\\\{\}]
\PYG{k}{improper\PYGZus{}style}\PYG{+w}{ }\PYG{n}{harmonic}
\end{sphinxVerbatim}


\subsection{Examples}
\label{\detokenize{improper_harmonic:examples}}
\begin{sphinxVerbatim}[commandchars=\\\{\}]
\PYG{k}{improper\PYGZus{}style}\PYG{+w}{ }\PYG{n}{harmonic}
\PYG{k}{improper\PYGZus{}coeff}\PYG{+w}{ }\PYG{l+m}{1}\PYG{+w}{ }\PYG{l+m}{100.0}\PYG{+w}{ }\PYG{l+m}{0}
\end{sphinxVerbatim}


\subsection{Description}
\label{\detokenize{improper_harmonic:description}}
\sphinxAtStartPar
The \sphinxstyleemphasis{harmonic} improper style uses the potential
\begin{equation*}
\begin{split}E = K (\chi - \chi_0)^2\end{split}
\end{equation*}
\sphinxAtStartPar
where \(\chi\) is the improper angle, \(\chi_0\) is its equilibrium
value, and \(K\) is a prefactor.  Note that the usual 1/2 factor is
included in \(K\).

\sphinxAtStartPar
If the 4 atoms in an improper quadruplet (listed in the data file read
by the {\hyperref[\detokenize{read_data::doc}]{\sphinxcrossref{\DUrole{doc}{read\_data}}}} command) are ordered I,J,K,L then
\(\chi\)
is the angle between the plane of I,J,K and the plane of J,K,L.
Alternatively, you can think of atoms J,K,L as being in a plane, and
atom I above the plane, and \(\chi\) as a measure of how far out\sphinxhyphen{}of\sphinxhyphen{}plane
I is with respect to the other 3 atoms.

\sphinxAtStartPar
Note that defining 4 atoms to interact in this way, does not mean that
bonds necessarily exist between I\sphinxhyphen{}J, J\sphinxhyphen{}K, or K\sphinxhyphen{}L, as they would in a
linear dihedral.  Normally, the bonds I\sphinxhyphen{}J, I\sphinxhyphen{}K, I\sphinxhyphen{}L would exist for an
improper to be defined between the 4 atoms.

\sphinxAtStartPar
The following coefficients must be defined for each improper type via
the {\hyperref[\detokenize{improper_coeff::doc}]{\sphinxcrossref{\DUrole{doc}{improper\_coeff}}}} command as in the example
above, or in the data file or restart files read by the
{\hyperref[\detokenize{read_data::doc}]{\sphinxcrossref{\DUrole{doc}{read\_data}}}} or {\hyperref[\detokenize{read_restart::doc}]{\sphinxcrossref{\DUrole{doc}{read\_restart}}}}
commands:
\begin{itemize}
\item {} 
\sphinxAtStartPar
\(K\) (energy)

\item {} 
\sphinxAtStartPar
\(\chi_0\) (degrees)

\end{itemize}

\sphinxAtStartPar
\(\chi_0\) is specified in degrees, but LAMMPS converts it to
radians internally; hence \(K\) is effectively energy per
radian\textasciicircum{}2.


\bigskip\hrule\bigskip


\sphinxAtStartPar
Styles with a \sphinxstyleemphasis{gpu}, \sphinxstyleemphasis{intel}, \sphinxstyleemphasis{kk}, \sphinxstyleemphasis{omp}, or \sphinxstyleemphasis{opt} suffix are
functionally the same as the corresponding style without the suffix.
They have been optimized to run faster, depending on your available
hardware, as discussed on the {\hyperref[\detokenize{Speed_packages::doc}]{\sphinxcrossref{\DUrole{doc}{Accelerator packages}}}}
page.  The accelerated styles take the same arguments and should
produce the same results, except for round\sphinxhyphen{}off and precision issues.

\sphinxAtStartPar
These accelerated styles are part of the GPU, INTEL, KOKKOS,
OPENMP, and OPT packages, respectively.  They are only enabled if
LAMMPS was built with those packages.  See the {\hyperref[\detokenize{Build_package::doc}]{\sphinxcrossref{\DUrole{doc}{Build package}}}} page for more info.

\sphinxAtStartPar
You can specify the accelerated styles explicitly in your input script
by including their suffix, or you can use the {\hyperref[\detokenize{Run_options::doc}]{\sphinxcrossref{\DUrole{doc}{\sphinxhyphen{}suffix command\sphinxhyphen{}line switch}}}} when you invoke LAMMPS, or you can use the
{\hyperref[\detokenize{suffix::doc}]{\sphinxcrossref{\DUrole{doc}{suffix}}}} command in your input script.

\sphinxAtStartPar
See the {\hyperref[\detokenize{Speed_packages::doc}]{\sphinxcrossref{\DUrole{doc}{Accelerator packages}}}} page for more
instructions on how to use the accelerated styles effectively.


\bigskip\hrule\bigskip



\subsection{Restrictions}
\label{\detokenize{improper_harmonic:restrictions}}
\sphinxAtStartPar
This improper style can only be used if LAMMPS was built with the
MOLECULE package.  See the {\hyperref[\detokenize{Build_package::doc}]{\sphinxcrossref{\DUrole{doc}{Build package}}}} doc page
for more info.


\subsection{Related commands}
\label{\detokenize{improper_harmonic:related-commands}}
\sphinxAtStartPar
{\hyperref[\detokenize{improper_coeff::doc}]{\sphinxcrossref{\DUrole{doc}{improper\_coeff}}}}


\subsection{Default}
\label{\detokenize{improper_harmonic:default}}
\sphinxAtStartPar
none

\sphinxstepscope

\index{improper\_style hybrid@\spxentry{improper\_style hybrid}}\ignorespaces 
\index{improper\_style hybrid/kk@\spxentry{improper\_style hybrid/kk}}\ignorespaces 

\section{improper\_style hybrid command}
\label{\detokenize{improper_hybrid:improper-style-hybrid-command}}\label{\detokenize{improper_hybrid:index-1}}\label{\detokenize{improper_hybrid:index-0}}\label{\detokenize{improper_hybrid::doc}}
\sphinxAtStartPar
Accelerator Variants: \sphinxstyleemphasis{hybrid/kk}


\subsection{Syntax}
\label{\detokenize{improper_hybrid:syntax}}
\begin{sphinxVerbatim}[commandchars=\\\{\}]
\PYG{k}{improper\PYGZus{}style}\PYG{+w}{ }\PYG{n}{hybrid}\PYG{+w}{ }\PYG{n}{style1}\PYG{+w}{ }\PYG{n}{style2}\PYG{+w}{ }\PYG{n}{...}
\end{sphinxVerbatim}
\begin{itemize}
\item {} 
\sphinxAtStartPar
style1,style2 = list of one or more improper styles

\end{itemize}


\subsection{Examples}
\label{\detokenize{improper_hybrid:examples}}
\begin{sphinxVerbatim}[commandchars=\\\{\}]
\PYG{k}{improper\PYGZus{}style}\PYG{+w}{ }\PYG{n}{hybrid}\PYG{+w}{ }\PYG{n}{harmonic}\PYG{+w}{ }\PYG{n}{cvff}
\PYG{k}{improper\PYGZus{}coeff}\PYG{+w}{ }\PYG{l+m}{1}\PYG{+w}{ }\PYG{n}{harmonic}\PYG{+w}{ }\PYG{l+m}{120.0}\PYG{+w}{ }\PYG{l+m}{30}
\PYG{k}{improper\PYGZus{}coeff}\PYG{+w}{ }\PYG{l+m}{2}\PYG{+w}{ }\PYG{n}{cvff}\PYG{+w}{ }\PYG{l+m}{20.0}\PYG{+w}{ }\PYG{o}{\PYGZhy{}}\PYG{l+m}{1}\PYG{+w}{ }\PYG{l+m}{2}
\end{sphinxVerbatim}


\subsection{Description}
\label{\detokenize{improper_hybrid:description}}
\sphinxAtStartPar
The \sphinxstyleemphasis{hybrid} style enables the use of multiple improper styles in one
simulation.  An improper style is assigned to each improper type.  For
example, impropers in a polymer flow (of improper type 1) could be
computed with a \sphinxstyleemphasis{harmonic} potential and impropers in the wall
boundary (of improper type 2) could be computed with a \sphinxstyleemphasis{cvff}
potential.  The assignment of improper type to style is made via the
{\hyperref[\detokenize{improper_coeff::doc}]{\sphinxcrossref{\DUrole{doc}{improper\_coeff}}}} command or in the data file.

\sphinxAtStartPar
In the improper\_coeff command, the first coefficient sets the improper
style and the remaining coefficients are those appropriate to that
style.  In the example above, the 2 improper\_coeff commands would set
impropers of improper type 1 to be computed with a \sphinxstyleemphasis{harmonic}
potential with coefficients 120.0, 30 for \(K\), \(\chi_0\).
Improper type 2 would be computed with a \sphinxstyleemphasis{cvff} potential with coefficients
20.0, \sphinxhyphen{}1, 2 for K, d, and n, respectively.

\sphinxAtStartPar
If improper coefficients are specified in the data file read via the
{\hyperref[\detokenize{read_data::doc}]{\sphinxcrossref{\DUrole{doc}{read\_data}}}} command, then the same rule applies.
E.g. “harmonic” or “cvff”, must be added after the improper type, for
each line in the “Improper Coeffs” section, e.g.

\begin{sphinxVerbatim}[commandchars=\\\{\}]
Improper Coeffs

1 harmonic 120.0 30
2 cvff 20.0 \PYGZhy{}1 2
...
\end{sphinxVerbatim}

\sphinxAtStartPar
If \sphinxstyleemphasis{class2} is one of the improper hybrid styles, the same rule holds
for specifying additional AngleAngle coefficients either via the input
script or in the data file. I.e. \sphinxstyleemphasis{class2} must be added to each line
after the improper type.  For
lines in the AngleAngle Coeffs section of the data
file for dihedral types that are not \sphinxstyleemphasis{class2}, you must use an
improper style of \sphinxstyleemphasis{skip} as a placeholder, e.g.

\begin{sphinxVerbatim}[commandchars=\\\{\}]
AngleAngle Coeffs

1 skip
2 class2 0.0 0.0 0.0 115.06 130.01 115.06
...
\end{sphinxVerbatim}

\sphinxAtStartPar
Note that it is not necessary to use the improper style \sphinxstyleemphasis{skip} in the
input script, since AngleAngle coefficients
need not be specified at all for improper types that are not \sphinxstyleemphasis{class2}.

\sphinxAtStartPar
An improper style of \sphinxstyleemphasis{none} can be specified as the second argument to
the improper\_coeff command, if you desire to turn off certain improper
types.


\bigskip\hrule\bigskip


\sphinxAtStartPar
Styles with a \sphinxstyleemphasis{gpu}, \sphinxstyleemphasis{intel}, \sphinxstyleemphasis{kk}, \sphinxstyleemphasis{omp}, or \sphinxstyleemphasis{opt} suffix are
functionally the same as the corresponding style without the suffix.
They have been optimized to run faster, depending on your available
hardware, as discussed on the {\hyperref[\detokenize{Speed_packages::doc}]{\sphinxcrossref{\DUrole{doc}{Accelerator packages}}}}
page.  The accelerated styles take the same arguments and should
produce the same results, except for round\sphinxhyphen{}off and precision issues.

\sphinxAtStartPar
These accelerated styles are part of the GPU, INTEL, KOKKOS,
OPENMP, and OPT packages, respectively.  They are only enabled if
LAMMPS was built with those packages.  See the {\hyperref[\detokenize{Build_package::doc}]{\sphinxcrossref{\DUrole{doc}{Build package}}}} page for more info.

\sphinxAtStartPar
You can specify the accelerated styles explicitly in your input script
by including their suffix, or you can use the {\hyperref[\detokenize{Run_options::doc}]{\sphinxcrossref{\DUrole{doc}{\sphinxhyphen{}suffix command\sphinxhyphen{}line switch}}}} when you invoke LAMMPS, or you can use the
{\hyperref[\detokenize{suffix::doc}]{\sphinxcrossref{\DUrole{doc}{suffix}}}} command in your input script.

\sphinxAtStartPar
See the {\hyperref[\detokenize{Speed_packages::doc}]{\sphinxcrossref{\DUrole{doc}{Accelerator packages}}}} page for more
instructions on how to use the accelerated styles effectively.


\bigskip\hrule\bigskip



\subsection{Restrictions}
\label{\detokenize{improper_hybrid:restrictions}}
\sphinxAtStartPar
This improper style can only be used if LAMMPS was built with the
MOLECULE package.  See the {\hyperref[\detokenize{Build_package::doc}]{\sphinxcrossref{\DUrole{doc}{Build package}}}} doc page
for more info.

\sphinxAtStartPar
Unlike other improper styles, the hybrid improper style does not store
improper coefficient info for individual sub\sphinxhyphen{}styles in {\hyperref[\detokenize{restart::doc}]{\sphinxcrossref{\DUrole{doc}{binary
restart files}}}} or {\hyperref[\detokenize{write_data::doc}]{\sphinxcrossref{\DUrole{doc}{data files}}}}.  Thus when
restarting a simulation, you need to re\sphinxhyphen{}specify the improper\_coeff
commands.


\subsection{Related commands}
\label{\detokenize{improper_hybrid:related-commands}}
\sphinxAtStartPar
{\hyperref[\detokenize{improper_coeff::doc}]{\sphinxcrossref{\DUrole{doc}{improper\_coeff}}}}


\subsection{Default}
\label{\detokenize{improper_hybrid:default}}
\sphinxAtStartPar
none

\sphinxstepscope

\index{improper\_style inversion/harmonic@\spxentry{improper\_style inversion/harmonic}}\ignorespaces 

\section{improper\_style inversion/harmonic command}
\label{\detokenize{improper_inversion_harmonic:improper-style-inversion-harmonic-command}}\label{\detokenize{improper_inversion_harmonic:index-0}}\label{\detokenize{improper_inversion_harmonic::doc}}

\subsection{Syntax}
\label{\detokenize{improper_inversion_harmonic:syntax}}
\begin{sphinxVerbatim}[commandchars=\\\{\}]
\PYG{k}{improper\PYGZus{}style}\PYG{+w}{ }\PYG{n}{inversion}\PYG{o}{/}\PYG{n}{harmonic}
\end{sphinxVerbatim}


\subsection{Examples}
\label{\detokenize{improper_inversion_harmonic:examples}}
\begin{sphinxVerbatim}[commandchars=\\\{\}]
\PYG{k}{improper\PYGZus{}style}\PYG{+w}{ }\PYG{n}{inversion}\PYG{o}{/}\PYG{n}{harmonic}
\PYG{k}{improper\PYGZus{}coeff}\PYG{+w}{ }\PYG{l+m}{1}\PYG{+w}{ }\PYG{l+m}{18.776340}\PYG{+w}{ }\PYG{l+m}{0.000000}
\end{sphinxVerbatim}


\subsection{Description}
\label{\detokenize{improper_inversion_harmonic:description}}
\sphinxAtStartPar
The \sphinxstyleemphasis{inversion/harmonic} improper style follows the Wilson\sphinxhyphen{}Decius
out\sphinxhyphen{}of\sphinxhyphen{}plane angle definition and uses an harmonic potential:
\begin{equation*}
\begin{split}E = K \left(\omega - \omega_0\right)^2\end{split}
\end{equation*}
\sphinxAtStartPar
where \(K\) is the force constant and \(\omega\) is the angle
evaluated for all three axis\sphinxhyphen{}plane combinations centered around the atom I.
For the IL axis and the IJK plane \(\omega\) looks as follows:

\noindent{\hspace*{\fill}\sphinxincludegraphics{{umbrella}.jpg}\hspace*{\fill}}

\sphinxAtStartPar
Note that the \sphinxstyleemphasis{inversion/harmonic} angle term evaluation differs to
the {\hyperref[\detokenize{improper_umbrella::doc}]{\sphinxcrossref{\DUrole{doc}{improper\_umbrella}}}} due to the cyclic
evaluation of all possible angles \(\omega\).

\sphinxAtStartPar
The following coefficients must be defined for each improper type via
the {\hyperref[\detokenize{improper_coeff::doc}]{\sphinxcrossref{\DUrole{doc}{improper\_coeff}}}} command as in the example
above, or in the data file or restart files read by the
{\hyperref[\detokenize{read_data::doc}]{\sphinxcrossref{\DUrole{doc}{read\_data}}}} or {\hyperref[\detokenize{read_restart::doc}]{\sphinxcrossref{\DUrole{doc}{read\_restart}}}}
commands:
\begin{itemize}
\item {} 
\sphinxAtStartPar
\(K\) (energy)

\item {} 
\sphinxAtStartPar
\(\omega_0\) (degrees)

\end{itemize}

\sphinxAtStartPar
If \(\omega_0 = 0\) the potential term has a single minimum for
the planar structure.  Otherwise it has two minima at +/\sphinxhyphen{} \(\omega_0\),
with a barrier in between.


\bigskip\hrule\bigskip



\subsection{Restrictions}
\label{\detokenize{improper_inversion_harmonic:restrictions}}
\sphinxAtStartPar
This improper style can only be used if LAMMPS was built with the
MOFFF package.  See the {\hyperref[\detokenize{Build_package::doc}]{\sphinxcrossref{\DUrole{doc}{Build package}}}} doc
page for more info.


\subsection{Related commands}
\label{\detokenize{improper_inversion_harmonic:related-commands}}
\sphinxAtStartPar
{\hyperref[\detokenize{improper_coeff::doc}]{\sphinxcrossref{\DUrole{doc}{improper\_coeff}}}}


\subsection{Default}
\label{\detokenize{improper_inversion_harmonic:default}}
\sphinxAtStartPar
none

\sphinxstepscope

\index{improper\_style none@\spxentry{improper\_style none}}\ignorespaces 

\section{improper\_style none command}
\label{\detokenize{improper_none:improper-style-none-command}}\label{\detokenize{improper_none:index-0}}\label{\detokenize{improper_none::doc}}

\subsection{Syntax}
\label{\detokenize{improper_none:syntax}}
\begin{sphinxVerbatim}[commandchars=\\\{\}]
\PYG{k}{improper\PYGZus{}style}\PYG{+w}{ }\PYG{n}{none}
\end{sphinxVerbatim}


\subsection{Examples}
\label{\detokenize{improper_none:examples}}
\begin{sphinxVerbatim}[commandchars=\\\{\}]
\PYG{k}{improper\PYGZus{}style}\PYG{+w}{ }\PYG{n}{none}
\end{sphinxVerbatim}


\subsection{Description}
\label{\detokenize{improper_none:description}}
\sphinxAtStartPar
Using an improper style of none means improper forces and energies are
not computed, even if quadruplets of improper atoms were listed in the
data file read by the {\hyperref[\detokenize{read_data::doc}]{\sphinxcrossref{\DUrole{doc}{read\_data}}}} command.

\sphinxAtStartPar
See the {\hyperref[\detokenize{improper_zero::doc}]{\sphinxcrossref{\DUrole{doc}{improper\_style zero}}}} command for a way to
calculate improper statistics, but compute no improper interactions.


\subsection{Restrictions}
\label{\detokenize{improper_none:restrictions}}\begin{quote}

\sphinxAtStartPar
none
\end{quote}


\subsection{Related commands}
\label{\detokenize{improper_none:related-commands}}
\sphinxAtStartPar
{\hyperref[\detokenize{improper_zero::doc}]{\sphinxcrossref{\DUrole{doc}{improper\_style zero}}}}


\subsection{Default}
\label{\detokenize{improper_none:default}}
\sphinxAtStartPar
none

\sphinxstepscope

\index{improper\_style ring@\spxentry{improper\_style ring}}\ignorespaces 
\index{improper\_style ring/omp@\spxentry{improper\_style ring/omp}}\ignorespaces 

\section{improper\_style ring command}
\label{\detokenize{improper_ring:improper-style-ring-command}}\label{\detokenize{improper_ring:index-1}}\label{\detokenize{improper_ring:index-0}}\label{\detokenize{improper_ring::doc}}
\sphinxAtStartPar
Accelerator Variants: \sphinxstyleemphasis{ring/omp}


\subsection{Syntax}
\label{\detokenize{improper_ring:syntax}}
\begin{sphinxVerbatim}[commandchars=\\\{\}]
\PYG{k}{improper\PYGZus{}style}\PYG{+w}{ }\PYG{n}{ring}
\end{sphinxVerbatim}


\subsection{Examples}
\label{\detokenize{improper_ring:examples}}
\begin{sphinxVerbatim}[commandchars=\\\{\}]
\PYG{k}{improper\PYGZus{}style}\PYG{+w}{ }\PYG{n}{ring}
\PYG{k}{improper\PYGZus{}coeff}\PYG{+w}{ }\PYG{l+m}{1}\PYG{+w}{ }\PYG{l+m}{8000}\PYG{+w}{ }\PYG{l+m}{70.5}
\end{sphinxVerbatim}


\subsection{Description}
\label{\detokenize{improper_ring:description}}
\sphinxAtStartPar
The \sphinxstyleemphasis{ring} improper style uses the potential
\begin{equation*}
\begin{split}E = &\frac{1}{6} K \left(\Delta_{ijl} + \Delta_{ijk} + \Delta_{kjl} \right)^6 \\
\Delta_{ijl} = & \cos{\theta_{ijl} - \cos{\theta_0}} \\
\Delta_{ijk} = & \cos{\theta_{ijk} - \cos{\theta_0}} \\
\Delta_{kjl} = & \cos{\theta_{kjl} - \cos{\theta_0}}\end{split}
\end{equation*}
\sphinxAtStartPar
where \(K\) is a prefactor, \(\theta\) is the angle formed by
the atoms specified by (i,j,k,l) indices and \(\theta_0\) its
equilibrium value.

\sphinxAtStartPar
If the 4 atoms in an improper quadruplet (listed in the data file read
by the {\hyperref[\detokenize{read_data::doc}]{\sphinxcrossref{\DUrole{doc}{read\_data}}}} command) are ordered i,j,k,l then
\(\theta_{ijl}\) is the angle between atoms i,j and l,
\(\theta_{ijk}\) is the angle between atoms i,j and k,
\(\theta_{kjl}\) is the angle between atoms j,k, and l.

\sphinxAtStartPar
The “ring” improper style implements the improper potential introduced
by Destree et al., in Equation (9) of {\hyperref[\detokenize{improper_ring:destree}]{\sphinxcrossref{\DUrole{std,std-ref}{(Destree)}}}}.  This
potential does not affect small amplitude vibrations but is used in an
ad\sphinxhyphen{}hoc way to prevent the onset of accidentally large amplitude
fluctuations leading to the occurrence of a planar conformation of the
three bonds i\sphinxhyphen{}j, j\sphinxhyphen{}k and j\sphinxhyphen{}l, an intermediate conformation toward the
chiral inversion of a methine carbon.  In the “Impropers” section of
data file four atoms: i, j, k and l are specified with i,j and l lying
on the backbone of the chain and k specifying the chirality of j.

\sphinxAtStartPar
The following coefficients must be defined for each improper type via
the {\hyperref[\detokenize{improper_coeff::doc}]{\sphinxcrossref{\DUrole{doc}{improper\_coeff}}}} command as in the example
above, or in the data file or restart files read by the
{\hyperref[\detokenize{read_data::doc}]{\sphinxcrossref{\DUrole{doc}{read\_data}}}} or {\hyperref[\detokenize{read_restart::doc}]{\sphinxcrossref{\DUrole{doc}{read\_restart}}}}
commands:
\begin{itemize}
\item {} 
\sphinxAtStartPar
\(K\) (energy)

\item {} 
\sphinxAtStartPar
\(\theta_0\) (degrees)

\end{itemize}


\bigskip\hrule\bigskip


\sphinxAtStartPar
Styles with a \sphinxstyleemphasis{gpu}, \sphinxstyleemphasis{intel}, \sphinxstyleemphasis{kk}, \sphinxstyleemphasis{omp}, or \sphinxstyleemphasis{opt} suffix are
functionally the same as the corresponding style without the suffix.
They have been optimized to run faster, depending on your available
hardware, as discussed on the {\hyperref[\detokenize{Speed_packages::doc}]{\sphinxcrossref{\DUrole{doc}{Accelerator packages}}}}
page.  The accelerated styles take the same arguments and should
produce the same results, except for round\sphinxhyphen{}off and precision issues.

\sphinxAtStartPar
These accelerated styles are part of the GPU, INTEL, KOKKOS,
OPENMP, and OPT packages, respectively.  They are only enabled if
LAMMPS was built with those packages.  See the {\hyperref[\detokenize{Build_package::doc}]{\sphinxcrossref{\DUrole{doc}{Build package}}}} page for more info.

\sphinxAtStartPar
You can specify the accelerated styles explicitly in your input script
by including their suffix, or you can use the {\hyperref[\detokenize{Run_options::doc}]{\sphinxcrossref{\DUrole{doc}{\sphinxhyphen{}suffix command\sphinxhyphen{}line switch}}}} when you invoke LAMMPS, or you can use the
{\hyperref[\detokenize{suffix::doc}]{\sphinxcrossref{\DUrole{doc}{suffix}}}} command in your input script.

\sphinxAtStartPar
See the {\hyperref[\detokenize{Speed_packages::doc}]{\sphinxcrossref{\DUrole{doc}{Accelerator packages}}}} page for more
instructions on how to use the accelerated styles effectively.


\bigskip\hrule\bigskip



\subsection{Restrictions}
\label{\detokenize{improper_ring:restrictions}}
\sphinxAtStartPar
This improper style can only be used if LAMMPS was built with the
EXTRA\sphinxhyphen{}MOLECULE package.  See the {\hyperref[\detokenize{Build_package::doc}]{\sphinxcrossref{\DUrole{doc}{Build package}}}}
doc page for more info.


\subsection{Related commands}
\label{\detokenize{improper_ring:related-commands}}
\sphinxAtStartPar
{\hyperref[\detokenize{improper_coeff::doc}]{\sphinxcrossref{\DUrole{doc}{improper\_coeff}}}}

\phantomsection\label{\detokenize{improper_ring:destree}}
\sphinxAtStartPar
\sphinxstylestrong{(Destree)} M. Destree, F. Laupretre, A. Lyulin, and J.\sphinxhyphen{}P.  Ryckaert,
J Chem Phys, 112, 9632 (2000).

\sphinxstepscope

\index{improper\_style sqdistharm@\spxentry{improper\_style sqdistharm}}\ignorespaces 

\section{improper\_style sqdistharm command}
\label{\detokenize{improper_sqdistharm:improper-style-sqdistharm-command}}\label{\detokenize{improper_sqdistharm:index-0}}\label{\detokenize{improper_sqdistharm::doc}}

\subsection{Syntax}
\label{\detokenize{improper_sqdistharm:syntax}}
\begin{sphinxVerbatim}[commandchars=\\\{\}]
\PYG{k}{improper\PYGZus{}style}\PYG{+w}{ }\PYG{n}{sqdistharm}
\end{sphinxVerbatim}


\subsection{Examples}
\label{\detokenize{improper_sqdistharm:examples}}
\begin{sphinxVerbatim}[commandchars=\\\{\}]
\PYG{k}{improper\PYGZus{}style}\PYG{+w}{ }\PYG{n}{sqdistharm}
\PYG{k}{improper\PYGZus{}coeff}\PYG{+w}{ }\PYG{l+m}{1}\PYG{+w}{ }\PYG{l+m}{50.0}\PYG{+w}{ }\PYG{l+m}{0.1}
\end{sphinxVerbatim}


\subsection{Description}
\label{\detokenize{improper_sqdistharm:description}}
\sphinxAtStartPar
The \sphinxstyleemphasis{sqdistharm} improper style uses the potential
\begin{equation*}
\begin{split}E = K (d^2 - {d_0}^2)^2\end{split}
\end{equation*}
\sphinxAtStartPar
where \(d\) is the distance between the central atom and the plane formed
by the other three atoms.  If the 4 atoms in an improper quadruplet
(listed in the data file read by the {\hyperref[\detokenize{read_data::doc}]{\sphinxcrossref{\DUrole{doc}{read\_data}}}}
command) are ordered I,J,K,L then the L\sphinxhyphen{}atom is assumed to be the
central atom. Note that this is different from the convention used
in the improper\_style distance.

\sphinxAtStartPar
The following coefficients must be defined for each improper type via
the improper\_coeff command as in the example above, or in the data
file or restart files read by the read\_data or read\_restart commands:
\begin{itemize}
\item {} 
\sphinxAtStartPar
\(K\) (energy/distance\textasciicircum{}4)

\item {} 
\sphinxAtStartPar
\({d_0}^2\) (distance\textasciicircum{}2)

\end{itemize}

\sphinxAtStartPar
Note that \({d_0}^2\) (in units distance\textasciicircum{}2) has be provided and not \(d_0\).


\bigskip\hrule\bigskip



\subsection{Restrictions}
\label{\detokenize{improper_sqdistharm:restrictions}}
\sphinxAtStartPar
This improper style can only be used if LAMMPS was built with the
MOLECULE package.  See the {\hyperref[\detokenize{Build_package::doc}]{\sphinxcrossref{\DUrole{doc}{Build package}}}} doc
page for more info.


\subsection{Related commands}
\label{\detokenize{improper_sqdistharm:related-commands}}
\sphinxAtStartPar
{\hyperref[\detokenize{improper_coeff::doc}]{\sphinxcrossref{\DUrole{doc}{improper\_coeff}}}}


\subsection{Default}
\label{\detokenize{improper_sqdistharm:default}}
\sphinxAtStartPar
none

\sphinxstepscope

\index{improper\_style umbrella@\spxentry{improper\_style umbrella}}\ignorespaces 
\index{improper\_style umbrella/omp@\spxentry{improper\_style umbrella/omp}}\ignorespaces 

\section{improper\_style umbrella command}
\label{\detokenize{improper_umbrella:improper-style-umbrella-command}}\label{\detokenize{improper_umbrella:index-1}}\label{\detokenize{improper_umbrella:index-0}}\label{\detokenize{improper_umbrella::doc}}
\sphinxAtStartPar
Accelerator Variants: \sphinxstyleemphasis{umbrella/omp}


\subsection{Syntax}
\label{\detokenize{improper_umbrella:syntax}}
\begin{sphinxVerbatim}[commandchars=\\\{\}]
\PYG{k}{improper\PYGZus{}style}\PYG{+w}{ }\PYG{n}{umbrella}
\end{sphinxVerbatim}


\subsection{Examples}
\label{\detokenize{improper_umbrella:examples}}
\begin{sphinxVerbatim}[commandchars=\\\{\}]
\PYG{k}{improper\PYGZus{}style}\PYG{+w}{ }\PYG{n}{umbrella}
\PYG{k}{improper\PYGZus{}coeff}\PYG{+w}{ }\PYG{l+m}{1}\PYG{+w}{ }\PYG{l+m}{100.0}\PYG{+w}{ }\PYG{l+m}{180.0}
\end{sphinxVerbatim}


\subsection{Description}
\label{\detokenize{improper_umbrella:description}}
\sphinxAtStartPar
The \sphinxstyleemphasis{umbrella} improper style uses the following potential, which is
commonly referred to as a classic inversion and used in the
{\hyperref[\detokenize{Howto_bioFF::doc}]{\sphinxcrossref{\DUrole{doc}{DREIDING}}}} force field:
\begin{equation*}
\begin{split}E = & \frac{1}{2}K\left( \frac{1}{\sin\omega_0}\right) ^2 \left( \cos\omega - \cos\omega_0\right) ^2 \qquad \omega_0 \neq 0^o \\
E = & K\left( 1-cos\omega\right)  \qquad \omega_0 = 0^o\end{split}
\end{equation*}
\sphinxAtStartPar
where \(K\) is the force constant and \(\omega\) is the angle between the IL
axis and the IJK plane:

\noindent{\hspace*{\fill}\sphinxincludegraphics{{umbrella}.jpg}\hspace*{\fill}}

\sphinxAtStartPar
If \(\omega_0 = 0\) the potential term has a minimum for the planar
structure.  Otherwise it has two minima at \(\omega +/- \omega_0\),
with a barrier in between.

\sphinxAtStartPar
See {\hyperref[\detokenize{improper_umbrella:umbrella-mayo}]{\sphinxcrossref{\DUrole{std,std-ref}{(Mayo)}}}} for a description of the DREIDING force field.

\sphinxAtStartPar
The following coefficients must be defined for each improper type via
the {\hyperref[\detokenize{improper_coeff::doc}]{\sphinxcrossref{\DUrole{doc}{improper\_coeff}}}} command as in the example
above, or in the data file or restart files read by the
{\hyperref[\detokenize{read_data::doc}]{\sphinxcrossref{\DUrole{doc}{read\_data}}}} or {\hyperref[\detokenize{read_restart::doc}]{\sphinxcrossref{\DUrole{doc}{read\_restart}}}}
commands:
\begin{itemize}
\item {} 
\sphinxAtStartPar
\(K\) (energy)

\item {} 
\sphinxAtStartPar
\(\omega_0\) (degrees)

\end{itemize}


\bigskip\hrule\bigskip


\sphinxAtStartPar
Styles with a \sphinxstyleemphasis{gpu}, \sphinxstyleemphasis{intel}, \sphinxstyleemphasis{kk}, \sphinxstyleemphasis{omp}, or \sphinxstyleemphasis{opt} suffix are
functionally the same as the corresponding style without the suffix.
They have been optimized to run faster, depending on your available
hardware, as discussed on the {\hyperref[\detokenize{Speed_packages::doc}]{\sphinxcrossref{\DUrole{doc}{Accelerator packages}}}}
page.  The accelerated styles take the same arguments and should
produce the same results, except for round\sphinxhyphen{}off and precision issues.

\sphinxAtStartPar
These accelerated styles are part of the GPU, INTEL, KOKKOS,
OPENMP, and OPT packages, respectively.  They are only enabled if
LAMMPS was built with those packages.  See the {\hyperref[\detokenize{Build_package::doc}]{\sphinxcrossref{\DUrole{doc}{Build package}}}} page for more info.

\sphinxAtStartPar
You can specify the accelerated styles explicitly in your input script
by including their suffix, or you can use the {\hyperref[\detokenize{Run_options::doc}]{\sphinxcrossref{\DUrole{doc}{\sphinxhyphen{}suffix command\sphinxhyphen{}line switch}}}} when you invoke LAMMPS, or you can use the
{\hyperref[\detokenize{suffix::doc}]{\sphinxcrossref{\DUrole{doc}{suffix}}}} command in your input script.

\sphinxAtStartPar
See the {\hyperref[\detokenize{Speed_packages::doc}]{\sphinxcrossref{\DUrole{doc}{Accelerator packages}}}} page for more
instructions on how to use the accelerated styles effectively.


\bigskip\hrule\bigskip



\subsection{Restrictions}
\label{\detokenize{improper_umbrella:restrictions}}
\sphinxAtStartPar
This improper style can only be used if LAMMPS was built with the
MOLECULE package.  See the {\hyperref[\detokenize{Build_package::doc}]{\sphinxcrossref{\DUrole{doc}{Build package}}}} doc page
for more info.


\subsection{Related commands}
\label{\detokenize{improper_umbrella:related-commands}}
\sphinxAtStartPar
{\hyperref[\detokenize{improper_coeff::doc}]{\sphinxcrossref{\DUrole{doc}{improper\_coeff}}}}


\subsection{Default}
\label{\detokenize{improper_umbrella:default}}
\sphinxAtStartPar
none


\bigskip\hrule\bigskip

\phantomsection\label{\detokenize{improper_umbrella:umbrella-mayo}}
\sphinxAtStartPar
\sphinxstylestrong{(Mayo)} Mayo, Olfason, Goddard III, J Phys Chem, 94, 8897\sphinxhyphen{}8909
(1990),

\sphinxstepscope

\index{improper\_style zero@\spxentry{improper\_style zero}}\ignorespaces 

\section{improper\_style zero command}
\label{\detokenize{improper_zero:improper-style-zero-command}}\label{\detokenize{improper_zero:index-0}}\label{\detokenize{improper_zero::doc}}

\subsection{Syntax}
\label{\detokenize{improper_zero:syntax}}
\begin{sphinxVerbatim}[commandchars=\\\{\}]
\PYG{k}{improper\PYGZus{}style}\PYG{+w}{ }\PYG{n}{zero}\PYG{+w}{ }\PYG{n}{[nocoeff]}
\end{sphinxVerbatim}


\subsection{Examples}
\label{\detokenize{improper_zero:examples}}
\begin{sphinxVerbatim}[commandchars=\\\{\}]
\PYG{k}{improper\PYGZus{}style}\PYG{+w}{ }\PYG{n}{zero}
\PYG{k}{improper\PYGZus{}style}\PYG{+w}{ }\PYG{n}{zero}\PYG{+w}{ }\PYG{n}{nocoeff}
\PYG{k}{improper\PYGZus{}coeff}\PYG{+w}{ }\PYG{o}{*}
\end{sphinxVerbatim}


\subsection{Description}
\label{\detokenize{improper_zero:description}}
\sphinxAtStartPar
Using an improper style of zero means improper forces and energies are
not computed, but the geometry of improper quadruplets is still
accessible to other commands.

\sphinxAtStartPar
As an example, the {\hyperref[\detokenize{compute_improper_local::doc}]{\sphinxcrossref{\DUrole{doc}{compute improper/local}}}} command can be used to
compute the chi values for the list of quadruplets of improper atoms
listed in the data file read by the {\hyperref[\detokenize{read_data::doc}]{\sphinxcrossref{\DUrole{doc}{read\_data}}}}
command.  If no improper style is defined, this command cannot be
used.

\sphinxAtStartPar
The optional \sphinxstyleemphasis{nocoeff} flag allows to read data files with a ImproperCoeff
section for any improper style. Similarly, any improper\_coeff commands
will only be checked for the improper type number and the rest ignored.

\sphinxAtStartPar
Note that the {\hyperref[\detokenize{improper_coeff::doc}]{\sphinxcrossref{\DUrole{doc}{improper\_coeff}}}} command must be
used for all improper types, though no additional values are
specified.


\subsection{Restrictions}
\label{\detokenize{improper_zero:restrictions}}\begin{quote}

\sphinxAtStartPar
none
\end{quote}


\subsection{Related commands}
\label{\detokenize{improper_zero:related-commands}}
\sphinxAtStartPar
none

\sphinxAtStartPar
{\hyperref[\detokenize{improper_none::doc}]{\sphinxcrossref{\DUrole{doc}{improper\_style none}}}}


\subsection{Default}
\label{\detokenize{improper_zero:default}}
\sphinxAtStartPar
none

\sphinxstepscope


\chapter{Dump Styles}
\label{\detokenize{dumps:dump-styles}}\label{\detokenize{dumps::doc}}
\sphinxstepscope

\index{dump@\spxentry{dump}}\ignorespaces 
\index{dump atom@\spxentry{dump atom}}\ignorespaces 
\index{dump cfg@\spxentry{dump cfg}}\ignorespaces 
\index{dump custom@\spxentry{dump custom}}\ignorespaces 
\index{dump dcd@\spxentry{dump dcd}}\ignorespaces 
\index{dump grid@\spxentry{dump grid}}\ignorespaces 
\index{dump grid/vtk@\spxentry{dump grid/vtk}}\ignorespaces 
\index{dump local@\spxentry{dump local}}\ignorespaces 
\index{dump xtc@\spxentry{dump xtc}}\ignorespaces 
\index{dump yaml@\spxentry{dump yaml}}\ignorespaces 
\index{dump xyz@\spxentry{dump xyz}}\ignorespaces 
\index{dump atom/gz@\spxentry{dump atom/gz}}\ignorespaces 
\index{dump cfg/gz@\spxentry{dump cfg/gz}}\ignorespaces 
\index{dump custom/gz@\spxentry{dump custom/gz}}\ignorespaces 
\index{dump local/gz@\spxentry{dump local/gz}}\ignorespaces 
\index{dump xyz/gz@\spxentry{dump xyz/gz}}\ignorespaces 
\index{dump atom/zstd@\spxentry{dump atom/zstd}}\ignorespaces 
\index{dump cfg/zstd@\spxentry{dump cfg/zstd}}\ignorespaces 
\index{dump custom/zstd@\spxentry{dump custom/zstd}}\ignorespaces 
\index{dump xyz/zstd@\spxentry{dump xyz/zstd}}\ignorespaces 
\index{dump local/zstd@\spxentry{dump local/zstd}}\ignorespaces 

\section{dump command}
\label{\detokenize{dump:dump-command}}\label{\detokenize{dump:index-20}}\label{\detokenize{dump:index-19}}\label{\detokenize{dump:index-18}}\label{\detokenize{dump:index-17}}\label{\detokenize{dump:index-16}}\label{\detokenize{dump:index-15}}\label{\detokenize{dump:index-14}}\label{\detokenize{dump:index-13}}\label{\detokenize{dump:index-12}}\label{\detokenize{dump:index-11}}\label{\detokenize{dump:index-10}}\label{\detokenize{dump:index-9}}\label{\detokenize{dump:index-8}}\label{\detokenize{dump:index-7}}\label{\detokenize{dump:index-6}}\label{\detokenize{dump:index-5}}\label{\detokenize{dump:index-4}}\label{\detokenize{dump:index-3}}\label{\detokenize{dump:index-2}}\label{\detokenize{dump:index-1}}\label{\detokenize{dump:index-0}}\label{\detokenize{dump::doc}}

\section{dump vtk command}
\label{\detokenize{dump:dump-vtk-command}}

\section{dump h5md command}
\label{\detokenize{dump:dump-h5md-command}}

\section{dump molfile command}
\label{\detokenize{dump:dump-molfile-command}}

\section{dump netcdf command}
\label{\detokenize{dump:dump-netcdf-command}}

\section{dump image command}
\label{\detokenize{dump:dump-image-command}}

\section{dump movie command}
\label{\detokenize{dump:dump-movie-command}}

\section{dump atom/adios command}
\label{\detokenize{dump:dump-atom-adios-command}}

\section{dump custom/adios command}
\label{\detokenize{dump:dump-custom-adios-command}}

\section{dump cfg/uef command}
\label{\detokenize{dump:dump-cfg-uef-command}}

\subsection{Syntax}
\label{\detokenize{dump:syntax}}
\begin{sphinxVerbatim}[commandchars=\\\{\}]
\PYG{k}{dump }\PYG{n+nv+nvIdentifier}{ID}\PYG{+w}{ }\PYG{n+nv+nvIdentifier}{group\PYGZhy{}ID}\PYG{+w}{ }\PYG{n}{style}\PYG{+w}{ }\PYG{n}{N}\PYG{+w}{ }\PYG{n}{file}\PYG{+w}{ }\PYG{n}{attribute1}\PYG{+w}{ }\PYG{n}{attribute2}\PYG{+w}{ }\PYG{n}{...}
\end{sphinxVerbatim}
\begin{itemize}
\item {} 
\sphinxAtStartPar
ID = user\sphinxhyphen{}assigned name for the dump

\item {} 
\sphinxAtStartPar
group\sphinxhyphen{}ID = ID of the group of atoms to be dumped

\item {} 
\sphinxAtStartPar
style = \sphinxstyleemphasis{atom} or \sphinxstyleemphasis{atom/adios} or \sphinxstyleemphasis{atom/gz} or \sphinxstyleemphasis{atom/zstd} or \sphinxstyleemphasis{cfg} or \sphinxstyleemphasis{cfg/gz} or \sphinxstyleemphasis{cfg/zstd} or \sphinxstyleemphasis{cfg/uef} or \sphinxstyleemphasis{custom} or \sphinxstyleemphasis{custom/gz} or \sphinxstyleemphasis{custom/zstd} or \sphinxstyleemphasis{custom/adios} or \sphinxstyleemphasis{dcd} or \sphinxstyleemphasis{grid} or \sphinxstyleemphasis{grid/vtk} or \sphinxstyleemphasis{h5md} or \sphinxstyleemphasis{image} or \sphinxstyleemphasis{local} or \sphinxstyleemphasis{local/gz} or \sphinxstyleemphasis{local/zstd} or \sphinxstyleemphasis{molfile} or \sphinxstyleemphasis{movie} or \sphinxstyleemphasis{netcdf} or \sphinxstyleemphasis{netcdf/mpiio} or \sphinxstyleemphasis{vtk} or \sphinxstyleemphasis{xtc} or \sphinxstyleemphasis{xyz} or \sphinxstyleemphasis{xyz/gz} or \sphinxstyleemphasis{xyz/zstd} or \sphinxstyleemphasis{yaml}

\item {} 
\sphinxAtStartPar
N = dump on timesteps which are multiples of N

\item {} 
\sphinxAtStartPar
file = name of file to write dump info to

\item {} 
\sphinxAtStartPar
attribute1,attribute2,… = list of attributes for a particular style
\begin{sphinxalltt}
\sphinxstyleemphasis{atom} attributes = none
\sphinxstyleemphasis{atom/adios} attributes = none,  discussed on {\hyperref[\detokenize{dump_adios::doc}]{\sphinxcrossref{\DUrole{doc}{dump atom/adios}}}} page
\sphinxstyleemphasis{atom/gz} attributes = none
\sphinxstyleemphasis{atom/zstd} attributes = none
\sphinxstyleemphasis{cfg} attributes = same as \sphinxstyleemphasis{custom} attributes, see below
\sphinxstyleemphasis{cfg/gz} attributes = same as \sphinxstyleemphasis{custom} attributes, see below
\sphinxstyleemphasis{cfg/zstd} attributes = same as \sphinxstyleemphasis{custom} attributes, see below
\sphinxstyleemphasis{cfg/uef} attributes = same as \sphinxstyleemphasis{custom} attributes, discussed on {\hyperref[\detokenize{dump_cfg_uef::doc}]{\sphinxcrossref{\DUrole{doc}{dump cfg/uef}}}} page
\sphinxstyleemphasis{custom}, \sphinxstyleemphasis{custom/gz}, \sphinxstyleemphasis{custom/zstd} attributes = see below
\sphinxstyleemphasis{custom/adios} attributes = same as \sphinxstyleemphasis{custom} attributes, discussed on {\hyperref[\detokenize{dump_adios::doc}]{\sphinxcrossref{\DUrole{doc}{dump custom/adios}}}} page
\sphinxstyleemphasis{dcd} attributes = none
\sphinxstyleemphasis{h5md} attributes = discussed on {\hyperref[\detokenize{dump_h5md::doc}]{\sphinxcrossref{\DUrole{doc}{dump h5md}}}} page
\sphinxstyleemphasis{grid} attributes = see below
\sphinxstyleemphasis{grid/vtk} attributes = see below
\sphinxstyleemphasis{image} attributes = discussed on {\hyperref[\detokenize{dump_image::doc}]{\sphinxcrossref{\DUrole{doc}{dump image}}}} page
\sphinxstyleemphasis{local}, \sphinxstyleemphasis{local/gz}, \sphinxstyleemphasis{local/zstd} attributes = see below
\sphinxstyleemphasis{molfile} attributes = discussed on {\hyperref[\detokenize{dump_molfile::doc}]{\sphinxcrossref{\DUrole{doc}{dump molfile}}}} page
\sphinxstyleemphasis{movie} attributes = discussed on {\hyperref[\detokenize{dump_image::doc}]{\sphinxcrossref{\DUrole{doc}{dump image}}}} page
\sphinxstyleemphasis{netcdf} attributes = discussed on {\hyperref[\detokenize{dump_netcdf::doc}]{\sphinxcrossref{\DUrole{doc}{dump netcdf}}}} page
\sphinxstyleemphasis{netcdf/mpiio} attributes = discussed on {\hyperref[\detokenize{dump_netcdf::doc}]{\sphinxcrossref{\DUrole{doc}{dump netcdf}}}} page
\sphinxstyleemphasis{vtk} attributes = same as \sphinxstyleemphasis{custom} attributes, see below, also {\hyperref[\detokenize{dump_vtk::doc}]{\sphinxcrossref{\DUrole{doc}{dump vtk}}}} page
\sphinxstyleemphasis{xtc} attributes = none
\sphinxstyleemphasis{xyz} attributes = none
\sphinxstyleemphasis{xyz/gz} attributes = none
\sphinxstyleemphasis{xyz/zstd} attributes = none
\sphinxstyleemphasis{yaml} attributes = same as \sphinxstyleemphasis{custom} attributes, see below
\end{sphinxalltt}

\item {} 
\sphinxAtStartPar
\sphinxstyleemphasis{custom} or \sphinxstyleemphasis{custom/gz} or \sphinxstyleemphasis{custom/zstd} or \sphinxstyleemphasis{cfg} or \sphinxstyleemphasis{cfg/gz} or \sphinxstyleemphasis{cfg/zstd} or \sphinxstyleemphasis{cfg/uef} or \sphinxstyleemphasis{netcdf} or \sphinxstyleemphasis{netcdf/mpiio} or \sphinxstyleemphasis{yaml} attributes:

\begin{sphinxVerbatim}[commandchars=\\\{\}]
possible attributes = id, mol, proc, procp1, type, element, mass,
                      x, y, z, xs, ys, zs, xu, yu, zu,
                      xsu, ysu, zsu, ix, iy, iz,
                      vx, vy, vz, fx, fy, fz,
                      q, mux, muy, muz, mu,
                      radius, diameter, omegax, omegay, omegaz,
                      angmomx, angmomy, angmomz, tqx, tqy, tqz,
                      c\PYGZus{}ID, c\PYGZus{}ID[I], f\PYGZus{}ID, f\PYGZus{}ID[I], v\PYGZus{}name,
                      i\PYGZus{}name, d\PYGZus{}name, i2\PYGZus{}name[I], d2\PYGZus{}name[I]
\end{sphinxVerbatim}
\begin{sphinxalltt}
id = atom ID
mol = molecule ID
proc = ID of processor that owns atom
procp1 = ID+1 of processor that owns atom
type = atom type
typelabel = atom {\hyperref[\detokenize{Howto_type_labels::doc}]{\sphinxcrossref{\DUrole{doc}{type label}}}}
element = name of atom element, as defined by {\hyperref[\detokenize{dump_modify::doc}]{\sphinxcrossref{\DUrole{doc}{dump\_modify}}}} command
mass = atom mass
x,y,z = unscaled atom coordinates
xs,ys,zs = scaled atom coordinates
xu,yu,zu = unwrapped atom coordinates
xsu,ysu,zsu = scaled unwrapped atom coordinates
ix,iy,iz = box image that the atom is in
vx,vy,vz = atom velocities
fx,fy,fz = forces on atoms
q = atom charge
mux,muy,muz = orientation of dipole moment of atom
mu = magnitude of dipole moment of atom
radius,diameter = radius, diameter of spherical particle
omegax,omegay,omegaz = angular velocity of spherical particle
angmomx,angmomy,angmomz = angular momentum of aspherical particle
tqx,tqy,tqz = torque on finite\sphinxhyphen{}size particles
c\_ID = per\sphinxhyphen{}atom vector calculated by a compute with ID
c\_ID{[}I{]} = Ith column of per\sphinxhyphen{}atom array calculated by a compute with ID, I can include wildcard (see below)
f\_ID = per\sphinxhyphen{}atom vector calculated by a fix with ID
f\_ID{[}I{]} = Ith column of per\sphinxhyphen{}atom array calculated by a fix with ID, I can include wildcard (see below)
v\_name = per\sphinxhyphen{}atom vector calculated by an atom\sphinxhyphen{}style variable with name
i\_name = custom integer vector with name
d\_name = custom floating point vector with name
i2\_name{[}I{]} = Ith column of custom integer array with name, I can include wildcard (see below)
d2\_name{[}I{]} = Ith column of custom floating point vector with name, I can include wildcard (see below)
\end{sphinxalltt}

\item {} 
\sphinxAtStartPar
\sphinxstyleemphasis{local} or \sphinxstyleemphasis{local/gz} or \sphinxstyleemphasis{local/zstd} attributes:

\begin{sphinxVerbatim}[commandchars=\\\{\}]
possible attributes = index, c\PYGZus{}ID, c\PYGZus{}ID[I], f\PYGZus{}ID, f\PYGZus{}ID[I]
  index = enumeration of local values
  c\PYGZus{}ID = local vector calculated by a compute with ID
  c\PYGZus{}ID[I] = Ith column of local array calculated by a compute with ID, I can include wildcard (see below)
  f\PYGZus{}ID = local vector calculated by a fix with ID
  f\PYGZus{}ID[I] = Ith column of local array calculated by a fix with ID, I can include wildcard (see below)
\end{sphinxVerbatim}

\item {} 
\sphinxAtStartPar
\sphinxstyleemphasis{grid} or \sphinxstyleemphasis{grid/vtk} attributes:

\begin{sphinxVerbatim}[commandchars=\\\{\}]
possible attributes = c\PYGZus{}ID:gname:dname, c\PYGZus{}ID:gname:dname[I], f\PYGZus{}ID:gname:dname, f\PYGZus{}ID:gname:dname[I]
  gname = name of grid defined by compute or fix
  dname = name of data field defined by compute or fix
  c\PYGZus{}ID = per\PYGZhy{}grid vector calculated by a compute with ID
  c\PYGZus{}ID[I] = Ith column of per\PYGZhy{}grid array calculated by a compute with ID, I can include wildcard (see below)
  f\PYGZus{}ID = per\PYGZhy{}grid vector calculated by a fix with ID
  f\PYGZus{}ID[I] = Ith column of per\PYGZhy{}grid array calculated by a fix with ID, I can include wildcard (see below)
\end{sphinxVerbatim}

\end{itemize}


\subsection{Examples}
\label{\detokenize{dump:examples}}
\begin{sphinxVerbatim}[commandchars=\\\{\}]
\PYG{k}{dump }\PYG{n+nv+nvIdentifier}{myDump}\PYG{+w}{ }\PYG{n+nv+nvIdentifier}{all}\PYG{+w}{ }\PYG{n}{atom}\PYG{+w}{ }\PYG{l+m}{100}\PYG{+w}{ }\PYG{n}{dump.lammpstrj}
\PYG{k}{dump }\PYG{n+nv+nvIdentifier}{myDump}\PYG{+w}{ }\PYG{n+nv+nvIdentifier}{all}\PYG{+w}{ }\PYG{n}{atom}\PYG{o}{/}\PYG{n}{gz}\PYG{+w}{ }\PYG{l+m}{100}\PYG{+w}{ }\PYG{n}{dump.atom.gz}
\PYG{k}{dump }\PYG{n+nv+nvIdentifier}{myDump}\PYG{+w}{ }\PYG{n+nv+nvIdentifier}{all}\PYG{+w}{ }\PYG{n}{atom}\PYG{o}{/}\PYG{n}{zstd}\PYG{+w}{ }\PYG{l+m}{100}\PYG{+w}{ }\PYG{n}{dump.atom.zst}
\PYG{k}{dump }\PYG{n+nv+nvIdentifier}{2}\PYG{+w}{ }\PYG{n+nv+nvIdentifier}{subgroup}\PYG{+w}{ }\PYG{n}{atom}\PYG{+w}{ }\PYG{l+m}{50}\PYG{+w}{ }\PYG{n}{dump.run.bin}
\PYG{k}{dump }\PYG{n+nv+nvIdentifier}{4a}\PYG{+w}{ }\PYG{n+nv+nvIdentifier}{all}\PYG{+w}{ }\PYG{n}{custom}\PYG{+w}{ }\PYG{l+m}{100}\PYG{+w}{ }\PYG{n}{dump.myforce.}\PYG{o}{*}\PYG{+w}{ }\PYG{n}{id}\PYG{+w}{ }\PYG{n}{type}\PYG{+w}{ }\PYG{n}{x}\PYG{+w}{ }\PYG{n}{y}\PYG{+w}{ }\PYG{n}{vx}\PYG{+w}{ }\PYG{n}{fx}
\PYG{k}{dump }\PYG{n+nv+nvIdentifier}{4a}\PYG{+w}{ }\PYG{n+nv+nvIdentifier}{all}\PYG{+w}{ }\PYG{n}{custom}\PYG{+w}{ }\PYG{l+m}{100}\PYG{+w}{ }\PYG{n}{dump.myvel.lammpsbin}\PYG{+w}{ }\PYG{n}{id}\PYG{+w}{ }\PYG{n}{type}\PYG{+w}{ }\PYG{n}{x}\PYG{+w}{ }\PYG{n}{y}\PYG{+w}{ }\PYG{n}{z}\PYG{+w}{ }\PYG{n}{vx}\PYG{+w}{ }\PYG{n}{vy}\PYG{+w}{ }\PYG{n}{vz}
\PYG{k}{dump }\PYG{n+nv+nvIdentifier}{4b}\PYG{+w}{ }\PYG{n+nv+nvIdentifier}{flow}\PYG{+w}{ }\PYG{n}{custom}\PYG{+w}{ }\PYG{l+m}{100}\PYG{+w}{ }\PYG{n}{dump.}\PYG{o}{\PYGZpc{}}\PYG{n}{.myforce}\PYG{+w}{ }\PYG{n}{id}\PYG{+w}{ }\PYG{n}{type}\PYG{+w}{ }\PYG{n}{c\PYGZus{}myF[3]}\PYG{+w}{ }\PYG{n}{v\PYGZus{}ke}
\PYG{k}{dump }\PYG{n+nv+nvIdentifier}{4b}\PYG{+w}{ }\PYG{n+nv+nvIdentifier}{flow}\PYG{+w}{ }\PYG{n}{custom}\PYG{+w}{ }\PYG{l+m}{100}\PYG{+w}{ }\PYG{n}{dump.}\PYG{o}{\PYGZpc{}}\PYG{n}{.myforce}\PYG{+w}{ }\PYG{n}{id}\PYG{+w}{ }\PYG{n}{type}\PYG{+w}{ }\PYG{n}{c\PYGZus{}myF[}\PYG{o}{*}\PYG{n}{]}\PYG{+w}{ }\PYG{n}{v\PYGZus{}ke}
\PYG{k}{dump }\PYG{n+nv+nvIdentifier}{2}\PYG{+w}{ }\PYG{n+nv+nvIdentifier}{inner}\PYG{+w}{ }\PYG{n}{cfg}\PYG{+w}{ }\PYG{l+m}{10}\PYG{+w}{ }\PYG{n}{dump.snap.}\PYG{o}{*}\PYG{n}{.cfg}\PYG{+w}{ }\PYG{n}{mass}\PYG{+w}{ }\PYG{n}{type}\PYG{+w}{ }\PYG{n}{xs}\PYG{+w}{ }\PYG{n}{ys}\PYG{+w}{ }\PYG{n}{zs}\PYG{+w}{ }\PYG{n}{vx}\PYG{+w}{ }\PYG{n}{vy}\PYG{+w}{ }\PYG{n}{vz}
\PYG{k}{dump }\PYG{n+nv+nvIdentifier}{snap}\PYG{+w}{ }\PYG{n+nv+nvIdentifier}{all}\PYG{+w}{ }\PYG{n}{cfg}\PYG{+w}{ }\PYG{l+m}{100}\PYG{+w}{ }\PYG{n}{dump.config.}\PYG{o}{*}\PYG{n}{.cfg}\PYG{+w}{ }\PYG{n}{mass}\PYG{+w}{ }\PYG{n}{type}\PYG{+w}{ }\PYG{n}{xs}\PYG{+w}{ }\PYG{n}{ys}\PYG{+w}{ }\PYG{n}{zs}\PYG{+w}{ }\PYG{n}{id}\PYG{+w}{ }\PYG{n}{type}\PYG{+w}{ }\PYG{n}{c\PYGZus{}Stress[2]}
\PYG{k}{dump }\PYG{n+nv+nvIdentifier}{1}\PYG{+w}{ }\PYG{n+nv+nvIdentifier}{all}\PYG{+w}{ }\PYG{n}{xtc}\PYG{+w}{ }\PYG{l+m}{1000}\PYG{+w}{ }\PYG{n}{file.xtc}
\end{sphinxVerbatim}


\subsection{Description}
\label{\detokenize{dump:description}}
\sphinxAtStartPar
Dump a snapshot of quantities to one or more files once every
\(N\) timesteps in one of several styles.  The timesteps on which
dump output is written can also be controlled by a variable.  See the
{\hyperref[\detokenize{dump_modify::doc}]{\sphinxcrossref{\DUrole{doc}{dump\_modify every}}}} command.

\sphinxAtStartPar
Almost all the styles output per\sphinxhyphen{}atom data, i.e. one or more values
per atom.  The exceptions are as follows.  The \sphinxstyleemphasis{local} styles output
one or more values per bond (angle, dihedral, improper) or per pair of
interacting atoms (force or neighbor interactions).  The \sphinxstyleemphasis{grid} styles
output one or more values per grid cell, which are produced by other
commands which overlay the simulation domain with a regular grid.  See
the {\hyperref[\detokenize{Howto_grid::doc}]{\sphinxcrossref{\DUrole{doc}{Howto grid}}}} doc page for details.  The \sphinxstyleemphasis{image}
style renders a JPG, PNG, or PPM image file of the system for each
snapshot, while the \sphinxstyleemphasis{movie} style combines and compresses the series
of images into a movie file; both styles are discussed in detail on
the {\hyperref[\detokenize{dump_image::doc}]{\sphinxcrossref{\DUrole{doc}{dump image}}}} page.

\sphinxAtStartPar
Only information for atoms in the specified group is dumped.  The
{\hyperref[\detokenize{dump_modify::doc}]{\sphinxcrossref{\DUrole{doc}{dump\_modify thresh and region and refresh}}}}
commands can also alter what atoms are included.  Not all styles
support these options; see details on the {\hyperref[\detokenize{dump_modify::doc}]{\sphinxcrossref{\DUrole{doc}{dump\_modify}}}} doc page.

\sphinxAtStartPar
As described below, the filename determines the kind of output: text
or binary or gzipped, one big file or one per timestep, one file for
all the processors or multiple smaller files.

\begin{sphinxadmonition}{note}{Note:}
\sphinxAtStartPar
Because periodic boundary conditions are enforced only on timesteps
when neighbor lists are rebuilt, the coordinates of an atom written
to a dump file may be slightly outside the simulation box.
Re\sphinxhyphen{}neighbor timesteps will not typically coincide with the
timesteps dump snapshots are written.  See the {\hyperref[\detokenize{dump_modify::doc}]{\sphinxcrossref{\DUrole{doc}{dump\_modify
pbc}}}} command if you wish to force coordinates to be
strictly inside the simulation box.
\end{sphinxadmonition}

\begin{sphinxadmonition}{note}{Note:}
\sphinxAtStartPar
Unless the {\hyperref[\detokenize{dump_modify::doc}]{\sphinxcrossref{\DUrole{doc}{dump\_modify sort}}}} option is invoked,
the lines of atom or grid information written to dump files
(typically one line per atom or grid cell) will be in an
indeterminate order for each snapshot.  This is even true when
running on a single processor, if the {\hyperref[\detokenize{atom_modify::doc}]{\sphinxcrossref{\DUrole{doc}{atom\_modify sort}}}} option is on, which it is by default.  In this case
atoms are re\sphinxhyphen{}ordered periodically during a simulation, due to
spatial sorting.  It is also true when running in parallel, because
data for a single snapshot is collected from multiple processors,
each of which owns a subset of the atoms.
\end{sphinxadmonition}

\begin{sphinxadmonition}{warning}{Warning:}
\sphinxAtStartPar
Without either including atom IDs or using the {\hyperref[\detokenize{dump_modify::doc}]{\sphinxcrossref{\DUrole{doc}{dump\_modify sort}}}} option, it is impossible for visualization programs
(e.g. OVITO or VMD) or analysis tools to assign data in different
frames consistently to the same atom.  This can lead to incorrect
visualizations or results.  LAMMPS will print a warning in such cases.
\end{sphinxadmonition}

\sphinxAtStartPar
For the \sphinxstyleemphasis{atom}, \sphinxstyleemphasis{custom}, \sphinxstyleemphasis{cfg}, \sphinxstyleemphasis{grid}, and \sphinxstyleemphasis{local} styles, sorting
is off by default.  For the \sphinxstyleemphasis{dcd}, \sphinxstyleemphasis{grid/vtk}, \sphinxstyleemphasis{xtc}, \sphinxstyleemphasis{xyz}, and
\sphinxstyleemphasis{molfile} styles, sorting by atom ID or grid ID is on by default. See
the {\hyperref[\detokenize{dump_modify::doc}]{\sphinxcrossref{\DUrole{doc}{dump\_modify}}}} page for details.

\sphinxAtStartPar
The \sphinxstyleemphasis{style} keyword determines what kind of data is written to the
dump file(s) and in what format.

\sphinxAtStartPar
Note that \sphinxstyleemphasis{atom}, \sphinxstyleemphasis{custom}, \sphinxstyleemphasis{dcd}, \sphinxstyleemphasis{xtc}, \sphinxstyleemphasis{xyz}, and \sphinxstyleemphasis{yaml} style dump
files can be read directly by \sphinxhref{https://www.ks.uiuc.edu/Research/vmd}{VMD},
a popular tool for visualizing and analyzing trajectories from atomic
and molecular systems.  For reading \sphinxstyleemphasis{netcdf} style dump files, the
netcdf plugin needs to be recompiled from source using a NetCDF version
compatible with the one used by LAMMPS.  The bundled plugin binary
uses a very old version of NetCDF that is not compatible with LAMMPS.

\sphinxAtStartPar
Likewise the \sphinxhref{https://www.ovito.org}{OVITO visualization package},
popular for materials modeling, can read the \sphinxstyleemphasis{atom}, \sphinxstyleemphasis{custom},
\sphinxstyleemphasis{local}, \sphinxstyleemphasis{xtc}, \sphinxstyleemphasis{cfg}, \sphinxstyleemphasis{netcdf}, and \sphinxstyleemphasis{xyz} style atom dump files
directly.  With version 3.8 and above, OVITO can also read and
visualize \sphinxstyleemphasis{grid} style dump files with grid cell data, including
iso\sphinxhyphen{}surface images of the grid cell values.

\sphinxAtStartPar
Note that settings made via the {\hyperref[\detokenize{dump_modify::doc}]{\sphinxcrossref{\DUrole{doc}{dump\_modify}}}}
command can also alter the format of individual values and content of
the dump file itself.  This includes the precision of values output to
text\sphinxhyphen{}based dump files which is controlled by the {\hyperref[\detokenize{dump_modify::doc}]{\sphinxcrossref{\DUrole{doc}{dump\_modify
format}}}} command and its options.


\bigskip\hrule\bigskip


\sphinxAtStartPar
Format of native LAMMPS format dump files:

\sphinxAtStartPar
The \sphinxstyleemphasis{atom}, \sphinxstyleemphasis{custom}, \sphinxstyleemphasis{grid}, and \sphinxstyleemphasis{local} styles create files in a
simple LAMMPS\sphinxhyphen{}specific text format that is mostly self\sphinxhyphen{}explanatory
when viewing a dump file.  Many post\sphinxhyphen{}processing tools either included
with LAMMPS or third\sphinxhyphen{}party tools can read this format, as does the
{\hyperref[\detokenize{rerun::doc}]{\sphinxcrossref{\DUrole{doc}{rerun}}}} command.  See tools described on the {\hyperref[\detokenize{Tools::doc}]{\sphinxcrossref{\DUrole{doc}{Tools}}}} doc page for examples, including \sphinxhref{https://lammps.github.io/pizza}{Pizza.py}.

\sphinxAtStartPar
For all these styles, the dimensions of the simulation box are
included in each snapshot.  The simulation box in LAMMPS can be
defined in one of 3 ways: orthogonal, restricted triclinic, and
general triclinic.  See the {\hyperref[\detokenize{Howto_triclinic::doc}]{\sphinxcrossref{\DUrole{doc}{Howto triclinic}}}}
doc page for a detailed description of all 3 options.

\sphinxAtStartPar
For an orthogonal simulation box the box information is formatted as:

\begin{sphinxVerbatim}[commandchars=\\\{\}]
ITEM: BOX BOUNDS xx yy zz
xlo xhi
ylo yhi
zlo zhi
\end{sphinxVerbatim}

\sphinxAtStartPar
where xlo,xhi are the maximum extents of the simulation box in the
\(x\)\sphinxhyphen{}dimension, and similarly for \(y\) and \(z\).  The
“xx yy zz” terms are six characters that encode the style of boundary for each
of the six simulation box boundaries (xlo,xhi; ylo,yhi; and zlo,zhi).  Each of
the six characters is one of \sphinxstyleemphasis{p} (periodic), \sphinxstyleemphasis{f} (fixed), \sphinxstyleemphasis{s} (shrink wrap),
or \sphinxstyleemphasis{m} (shrink wrapped with a minimum value).  See the
{\hyperref[\detokenize{boundary::doc}]{\sphinxcrossref{\DUrole{doc}{boundary}}}} command for details.

\sphinxAtStartPar
For a restricted triclinic simulation box, an orthogonal bounding box
which encloses the restricted triclinic simulation box is output,
along with the three tilt factors (\sphinxstyleemphasis{xy}, \sphinxstyleemphasis{xz}, \sphinxstyleemphasis{yz}) of the triclinic
box, formatted as follows:

\begin{sphinxVerbatim}[commandchars=\\\{\}]
ITEM: BOX BOUNDS xy xz yz xx yy zz
xlo\PYGZus{}bound xhi\PYGZus{}bound xy
ylo\PYGZus{}bound yhi\PYGZus{}bound xz
zlo\PYGZus{}bound zhi\PYGZus{}bound yz
\end{sphinxVerbatim}

\sphinxAtStartPar
The presence of the text “xy xz yz” in the ITEM line indicates that
the three tilt factors will be included on each of the three following lines.
This bounding box is convenient for many visualization programs.  The
meaning of the six character flags for “xx yy zz” is the same as above.

\sphinxAtStartPar
Note that the first two numbers on each line are now xlo\_bound instead
of xlo, etc. because they represent a bounding box.  See the {\hyperref[\detokenize{Howto_triclinic::doc}]{\sphinxcrossref{\DUrole{doc}{Howto
triclinic}}}} page for a geometric description of
triclinic boxes, as defined by LAMMPS, simple formulas for how the six
bounding box extents (xlo\_bound, xhi\_bound, etc.) are calculated from the
triclinic parameters, and how to transform those parameters to and
from other commonly used triclinic representations.

\sphinxAtStartPar
For a general triclinic simulation box, see the “General triclinic”
section below for a description of the ITEM: BOX BOUNDS format as well
as how per\sphinxhyphen{}atom coordinates and per\sphinxhyphen{}atom vector quantities are output.

\sphinxAtStartPar
The \sphinxstyleemphasis{atom} and \sphinxstyleemphasis{custom} styles output a “ITEM: NUMBER OF ATOMS” line
with the count of atoms in the snapshot.  Likewise they output an
“ITEM: ATOMS” line which includes column descriptors for the per\sphinxhyphen{}atom
lines that follow.  For example, the descriptors would be “id type xs
ys zs” for the default \sphinxstyleemphasis{atom} style, and would be the atom attributes
you specify in the dump command for the \sphinxstyleemphasis{custom} style.  Each
subsequent line will list the data for a single atom.

\sphinxAtStartPar
For style \sphinxstyleemphasis{atom}, atom coordinates are written to the file, along with
the atom ID and atom type.  By default, atom coords are written in a
scaled format (from 0 to 1).  That is, an \(x\) value of 0.25 means the
atom is at a location 1/4 of the distance from \sphinxstyleemphasis{xlo} to \sphinxstyleemphasis{xhi} of the box
boundaries.  The format can be changed to unscaled coords via the
{\hyperref[\detokenize{dump_modify::doc}]{\sphinxcrossref{\DUrole{doc}{dump\_modify}}}} settings.  Image flags can also be
added for each atom via dump\_modify.

\sphinxAtStartPar
Style \sphinxstyleemphasis{custom} allows you to specify a list of atom attributes to be
written to the dump file for each atom.  Possible attributes are
listed above and will appear in the order specified.  You cannot
specify a quantity that is not defined for a particular
simulation—such as \sphinxstyleemphasis{q} for atom style \sphinxstyleemphasis{bond}, since that atom style
does not assign charges.  Dumps occur at the very end of a timestep,
so atom attributes will include effects due to fixes that are applied
during the timestep.  An explanation of the possible dump custom
attributes is given below.

\sphinxAtStartPar
\DUrole{versionmodified,added}{Added in version 22Dec2022.}

\sphinxAtStartPar
For style \sphinxstyleemphasis{grid} the dimension of the simulation domain and size of
the Nx by Ny by Nz grid that overlays the simulation domain are also
output with each snapshot:

\begin{sphinxVerbatim}[commandchars=\\\{\}]
ITEM: DIMENSION
dim
ITEM: GRID SIZE
nx ny nz
\end{sphinxVerbatim}

\sphinxAtStartPar
The value dim will be 2 or 3 for 2d or 3d simulations.  It is included
so that post\sphinxhyphen{}processing tools like \sphinxhref{https://www.ovito.org}{OVITO},
which can visualize grid\sphinxhyphen{}based quantities know how to draw each grid
cell.  The grid size will match the input script parameters for
grid(s) created by the computes or fixes which are referenced by the
the dump command.  For 2d simulations (and grids), nz will always
be 1.

\sphinxAtStartPar
There will also be an “ITEM: GRID DATA” line which includes column
descriptors for the per grid cell data.  Each subsequent line (Nx *
Ny * Nz lines) will list the data for a single grid cell.  If grid
cell IDs are included in the output via the {\hyperref[\detokenize{compute_property_grid::doc}]{\sphinxcrossref{\DUrole{doc}{compute
property/grid}}}} command, then the IDs will
range from 1 to N = Nx*Ny*Nz.  The ordering of IDs is with the x index
varying fastest, then the y index, and the z index varying slowest.

\sphinxAtStartPar
For style \sphinxstyleemphasis{local}, local output generated by {\hyperref[\detokenize{compute::doc}]{\sphinxcrossref{\DUrole{doc}{computes}}}}
and {\hyperref[\detokenize{fix::doc}]{\sphinxcrossref{\DUrole{doc}{fixes}}}} is used to generate lines of output that is
written to the dump file.  This local data is typically calculated by
each processor based on the atoms it owns, but there may be zero or
more entities per atom (e.g., a list of bond distances).  An explanation
of the possible dump local attributes is given below.  Note that by
using input from the {\hyperref[\detokenize{compute_property_local::doc}]{\sphinxcrossref{\DUrole{doc}{compute property/local}}}} command with dump local, it is possible to
generate information on bonds, angles, etc. that can be cut and pasted
directly into a data file read by the {\hyperref[\detokenize{read_data::doc}]{\sphinxcrossref{\DUrole{doc}{read\_data}}}}
command.


\bigskip\hrule\bigskip


\sphinxAtStartPar
Dump files in other popular formats:

\begin{sphinxadmonition}{note}{Note:}
\sphinxAtStartPar
This section only discusses file formats relevant to this doc page.
The top of this page has links to other dump commands (with their
own pages) which write files in additional popular formats.
\end{sphinxadmonition}

\sphinxAtStartPar
Style \sphinxstyleemphasis{cfg} has the same command syntax as style \sphinxstyleemphasis{custom} and writes
extended CFG format files, as used by the \sphinxhref{http://li.mit.edu/Archive/Graphics/A/}{AtomEye} visualization package.
Since the extended CFG format uses a single snapshot of the system per
file, a wildcard “*” must be included in the filename, as discussed
below.  The list of atom attributes for style \sphinxstyleemphasis{cfg} must begin with
either “mass type xs ys zs” or “mass type xsu ysu zsu” since these
quantities are needed to write the CFG files in the appropriate format
(though the “mass” and “type” fields do not appear explicitly in the
file).  Any remaining attributes will be stored as “auxiliary
properties” in the CFG files.  Note that you will typically want to
use the {\hyperref[\detokenize{dump_modify::doc}]{\sphinxcrossref{\DUrole{doc}{dump\_modify element}}}} command with
CFG\sphinxhyphen{}formatted files, to associate element names with atom types, so
that AtomEye can render atoms appropriately. When unwrapped
coordinates \sphinxstyleemphasis{xsu}, \sphinxstyleemphasis{ysu}, and \sphinxstyleemphasis{zsu} are requested, the nominal AtomEye
periodic cell dimensions are expanded by a large factor UNWRAPEXPAND =
10.0, which ensures atoms that are displayed correctly for up to
UNWRAPEXPAND/2 periodic boundary crossings in any direction.  Beyond
this, AtomEye will rewrap the unwrapped coordinates.  The expansion
causes the atoms to be drawn farther away from the viewer, but it is
easy to zoom the atoms closer, and the interatomic distances are
unaffected.

\sphinxAtStartPar
The \sphinxstyleemphasis{dcd} style writes DCD files, a standard atomic trajectory format
used by the CHARMM, NAMD, and XPlor molecular dynamics packages.  DCD
files are binary and thus may not be portable to different machines.
The number of atoms per snapshot cannot change with the \sphinxstyleemphasis{dcd} style.
The \sphinxstyleemphasis{unwrap} option of the {\hyperref[\detokenize{dump_modify::doc}]{\sphinxcrossref{\DUrole{doc}{dump\_modify}}}} command
allows DCD coordinates to be written “unwrapped” by the image flags
for each atom.  Unwrapped means that if the atom has passed through
a periodic boundary one or more times, the value is printed for what
the coordinate would be if it had not been wrapped back into the
periodic box.  Note that these coordinates may thus be far outside
the box size stored with the snapshot.

\sphinxAtStartPar
The \sphinxstyleemphasis{xtc} style writes XTC files, a compressed trajectory format used
by the GROMACS molecular dynamics package, and described \sphinxhref{https://manual.gromacs.org/current/reference-manual/file-formats.html\#xtc}{here}.
The precision used in XTC files can be adjusted via the
{\hyperref[\detokenize{dump_modify::doc}]{\sphinxcrossref{\DUrole{doc}{dump\_modify}}}} command.  The default value of 1000
means that coordinates are stored to 1/1000 nanometer accuracy.  XTC
files are portable binary files written in the NFS XDR data format, so
that any machine which supports XDR should be able to read them.  The
number of atoms per snapshot cannot change with the \sphinxstyleemphasis{xtc} style.  The
\sphinxstyleemphasis{unwrap} option of the {\hyperref[\detokenize{dump_modify::doc}]{\sphinxcrossref{\DUrole{doc}{dump\_modify}}}} command allows
XTC coordinates to be written “unwrapped” by the image flags for each
atom.  Unwrapped means that if the atom has passed through a periodic
boundary one or more times, the value is printed for what the
coordinate would be if it had not been wrapped back into the periodic
box.  Note that these coordinates may thus be far outside the box size
stored with the snapshot.

\sphinxAtStartPar
The \sphinxstyleemphasis{xyz} style writes XYZ files, which is a simple text\sphinxhyphen{}based
coordinate format that many codes can read. Specifically it has a line
with the number of atoms, then a comment line that is usually ignored
followed by one line per atom with the atom type and the \(x\)\sphinxhyphen{},
\(y\)\sphinxhyphen{}, and \(z\)\sphinxhyphen{}coordinate of that atom.  You can use the
{\hyperref[\detokenize{dump_modify::doc}]{\sphinxcrossref{\DUrole{doc}{dump\_modify element}}}} option to change the output
from using the (numerical) atom type to an element name (or some other
label). This option will help many visualization programs to guess bonds
and colors. You can use the {\hyperref[\detokenize{dump_modify::doc}]{\sphinxcrossref{\DUrole{doc}{dump\_modify types labels}}}}
option to replace numeric atom types with {\hyperref[\detokenize{Howto_type_labels::doc}]{\sphinxcrossref{\DUrole{doc}{type labels}}}}.

\sphinxAtStartPar
\DUrole{versionmodified,added}{Added in version 22Dec2022.}

\sphinxAtStartPar
The \sphinxstyleemphasis{grid/vtk} style writes VTK files for grid data on a regular
rectilinear grid.  Its content is conceptually similar to that of the
text file produced by the \sphinxstyleemphasis{grid} style, except that it in an XML\sphinxhyphen{}based
format which visualization programs which support the VTK format can
read, e.g. the \sphinxhref{https://www.paraview.org}{ParaView tool}.  For this
style, there can only be 1 or 3 per grid cell attributes specified.
If it is a single value, it is a scalar quantity.  If 3 values are
specified it is encoded in the VTK file as a vector quantity (for each
grid cell).  The filename for this style must include a “*” wildcard
character to produce one file per snapshot; see details below.

\sphinxAtStartPar
\DUrole{versionmodified,added}{Added in version 4May2022.}

\sphinxAtStartPar
Dump style \sphinxstyleemphasis{yaml} has the same command syntax as style \sphinxstyleemphasis{custom} and
writes YAML format files that can be easily parsed by a variety of
data processing tools and programming languages.  Each timestep will
be written as a YAML “document” (i.e., starts with “—” and ends with
“…”).  The style supports writing one file per timestep through the
“*” wildcard but not multi\sphinxhyphen{}processor outputs with the “\%” token in
the filename.  In addition to per\sphinxhyphen{}atom data, {\hyperref[\detokenize{thermo::doc}]{\sphinxcrossref{\DUrole{doc}{thermo}}}}
data can be included in the \sphinxstyleemphasis{yaml} style dump file using the
{\hyperref[\detokenize{dump_modify::doc}]{\sphinxcrossref{\DUrole{doc}{dump\_modify thermo yes}}}}. The data included in the
dump file uses the “thermo” tag and is otherwise identical to data
specified by the {\hyperref[\detokenize{thermo_style::doc}]{\sphinxcrossref{\DUrole{doc}{thermo\_style}}}} command.

\sphinxAtStartPar
Below is an example for a YAML format dump created by the following commands.

\begin{sphinxVerbatim}[commandchars=\\\{\}]
\PYG{k}{dump }\PYG{n+nv+nvIdentifier}{out}\PYG{+w}{ }\PYG{n+nv+nvIdentifier}{all}\PYG{+w}{ }\PYG{n}{yaml}\PYG{+w}{ }\PYG{l+m}{100}\PYG{+w}{ }\PYG{n}{dump.yaml}\PYG{+w}{ }\PYG{n}{id}\PYG{+w}{ }\PYG{n}{type}\PYG{+w}{ }\PYG{n}{x}\PYG{+w}{ }\PYG{n}{y}\PYG{+w}{ }\PYG{n}{z}\PYG{+w}{ }\PYG{n}{vx}\PYG{+w}{ }\PYG{n}{vy}\PYG{+w}{ }\PYG{n}{vz}\PYG{+w}{ }\PYG{n}{ix}\PYG{+w}{ }\PYG{n}{iy}\PYG{+w}{ }\PYG{n}{iz}
\PYG{k}{dump\PYGZus{}modify }\PYG{n+nv+nvIdentifier}{out}\PYG{+w}{ }\PYG{n}{time}\PYG{+w}{ }\PYG{n}{yes}\PYG{+w}{ }\PYG{n}{units}\PYG{+w}{ }\PYG{n}{yes}\PYG{+w}{ }\PYG{n}{thermo}\PYG{+w}{ }\PYG{n}{yes}\PYG{+w}{ }\PYG{n}{format}\PYG{+w}{ }\PYG{l+m}{1}\PYG{+w}{ }\PYG{o}{\PYGZpc{}}\PYG{l+m}{5}\PYG{n}{d}\PYG{+w}{ }\PYG{n}{format}\PYG{+w}{ }\PYG{l+s}{\PYGZdq{}}\PYG{l+s}{\PYGZpc{} 10.6e}\PYG{l+s}{\PYGZdq{}}
\end{sphinxVerbatim}

\sphinxAtStartPar
The tags “time”, “units”, and “thermo” are optional and enabled by the
dump\_modify command. The list under the “box” tag has three lines for
orthogonal boxes and four lines for triclinic boxes, where the first
three are the box boundaries and the fourth the three tilt factors
(\(xy\), \(xz\), \(yz\)).  The “thermo” data follows the
format of the \sphinxstyleemphasis{yaml} thermo style.  The “keywords” tag lists the
per\sphinxhyphen{}atom properties contained in the “data” columns, which contain a
list with one line per atom.  The keywords may be renamed using the
dump\_modify command same as for the \sphinxstyleemphasis{custom} dump style.

\begin{sphinxVerbatim}[commandchars=\\\{\}]
\PYG{n+nn}{\PYGZhy{}\PYGZhy{}\PYGZhy{}}
\PYG{n+nt}{creator}\PYG{p}{:}\PYG{+w}{ }\PYG{l+lScalar+lScalarPlain}{LAMMPS}
\PYG{n+nt}{timestep}\PYG{p}{:}\PYG{+w}{ }\PYG{l+lScalar+lScalarPlain}{0}
\PYG{n+nt}{units}\PYG{p}{:}\PYG{+w}{ }\PYG{l+lScalar+lScalarPlain}{lj}
\PYG{n+nt}{time}\PYG{p}{:}\PYG{+w}{ }\PYG{l+lScalar+lScalarPlain}{0}
\PYG{n+nt}{natoms}\PYG{p}{:}\PYG{+w}{ }\PYG{l+lScalar+lScalarPlain}{4000}
\PYG{n+nt}{boundary}\PYG{p}{:}\PYG{+w}{ }\PYG{p+pIndicator}{[}\PYG{+w}{ }\PYG{n+nv}{p}\PYG{p+pIndicator}{,}\PYG{+w}{ }\PYG{n+nv}{p}\PYG{p+pIndicator}{,}\PYG{+w}{ }\PYG{n+nv}{p}\PYG{p+pIndicator}{,}\PYG{+w}{ }\PYG{n+nv}{p}\PYG{p+pIndicator}{,}\PYG{+w}{ }\PYG{n+nv}{p}\PYG{p+pIndicator}{,}\PYG{+w}{ }\PYG{n+nv}{p}\PYG{p+pIndicator}{,}\PYG{+w}{ }\PYG{p+pIndicator}{]}
\PYG{n+nt}{thermo}\PYG{p}{:}
\PYG{+w}{  }\PYG{p+pIndicator}{\PYGZhy{}}\PYG{+w}{ }\PYG{n+nt}{keywords}\PYG{p}{:}\PYG{+w}{ }\PYG{p+pIndicator}{[}\PYG{+w}{ }\PYG{n+nv}{Step}\PYG{p+pIndicator}{,}\PYG{+w}{ }\PYG{n+nv}{Temp}\PYG{p+pIndicator}{,}\PYG{+w}{ }\PYG{n+nv}{E\PYGZus{}pair}\PYG{p+pIndicator}{,}\PYG{+w}{ }\PYG{n+nv}{E\PYGZus{}mol}\PYG{p+pIndicator}{,}\PYG{+w}{ }\PYG{n+nv}{TotEng}\PYG{p+pIndicator}{,}\PYG{+w}{ }\PYG{n+nv}{Press}\PYG{p+pIndicator}{,}\PYG{+w}{ }\PYG{p+pIndicator}{]}
\PYG{+w}{  }\PYG{p+pIndicator}{\PYGZhy{}}\PYG{+w}{ }\PYG{n+nt}{data}\PYG{p}{:}\PYG{+w}{ }\PYG{p+pIndicator}{[}\PYG{+w}{ }\PYG{n+nv}{0}\PYG{p+pIndicator}{,}\PYG{+w}{ }\PYG{n+nv}{0}\PYG{p+pIndicator}{,}\PYG{+w}{ }\PYG{n+nv}{\PYGZhy{}27093.472213010766}\PYG{p+pIndicator}{,}\PYG{+w}{ }\PYG{n+nv}{0}\PYG{p+pIndicator}{,}\PYG{+w}{ }\PYG{n+nv}{0}\PYG{p+pIndicator}{,}\PYG{+w}{ }\PYG{n+nv}{0}\PYG{p+pIndicator}{,}\PYG{+w}{ }\PYG{p+pIndicator}{]}
\PYG{n+nt}{box}\PYG{p}{:}
\PYG{+w}{  }\PYG{p+pIndicator}{\PYGZhy{}}\PYG{+w}{ }\PYG{p+pIndicator}{[}\PYG{+w}{ }\PYG{n+nv}{0}\PYG{p+pIndicator}{,}\PYG{+w}{ }\PYG{n+nv}{16.795961913825074}\PYG{+w}{ }\PYG{p+pIndicator}{]}
\PYG{+w}{  }\PYG{p+pIndicator}{\PYGZhy{}}\PYG{+w}{ }\PYG{p+pIndicator}{[}\PYG{+w}{ }\PYG{n+nv}{0}\PYG{p+pIndicator}{,}\PYG{+w}{ }\PYG{n+nv}{16.795961913825074}\PYG{+w}{ }\PYG{p+pIndicator}{]}
\PYG{+w}{  }\PYG{p+pIndicator}{\PYGZhy{}}\PYG{+w}{ }\PYG{p+pIndicator}{[}\PYG{+w}{ }\PYG{n+nv}{0}\PYG{p+pIndicator}{,}\PYG{+w}{ }\PYG{n+nv}{16.795961913825074}\PYG{+w}{ }\PYG{p+pIndicator}{]}
\PYG{+w}{  }\PYG{p+pIndicator}{\PYGZhy{}}\PYG{+w}{ }\PYG{p+pIndicator}{[}\PYG{+w}{ }\PYG{n+nv}{0}\PYG{p+pIndicator}{,}\PYG{+w}{ }\PYG{n+nv}{0}\PYG{p+pIndicator}{,}\PYG{+w}{ }\PYG{n+nv}{0}\PYG{+w}{ }\PYG{p+pIndicator}{]}
\PYG{n+nt}{keywords}\PYG{p}{:}\PYG{+w}{ }\PYG{p+pIndicator}{[}\PYG{+w}{ }\PYG{n+nv}{id}\PYG{p+pIndicator}{,}\PYG{+w}{ }\PYG{n+nv}{type}\PYG{p+pIndicator}{,}\PYG{+w}{ }\PYG{n+nv}{x}\PYG{p+pIndicator}{,}\PYG{+w}{ }\PYG{n+nv}{y}\PYG{p+pIndicator}{,}\PYG{+w}{ }\PYG{n+nv}{z}\PYG{p+pIndicator}{,}\PYG{+w}{ }\PYG{n+nv}{vx}\PYG{p+pIndicator}{,}\PYG{+w}{ }\PYG{n+nv}{vy}\PYG{p+pIndicator}{,}\PYG{+w}{ }\PYG{n+nv}{vz}\PYG{p+pIndicator}{,}\PYG{+w}{ }\PYG{n+nv}{ix}\PYG{p+pIndicator}{,}\PYG{+w}{ }\PYG{n+nv}{iy}\PYG{p+pIndicator}{,}\PYG{+w}{ }\PYG{n+nv}{iz}\PYG{p+pIndicator}{,}\PYG{+w}{  }\PYG{p+pIndicator}{]}
\PYG{n+nt}{data}\PYG{p}{:}
\PYG{+w}{  }\PYG{p+pIndicator}{\PYGZhy{}}\PYG{+w}{ }\PYG{p+pIndicator}{[}\PYG{+w}{     }\PYG{n+nv}{1}\PYG{+w}{ }\PYG{p+pIndicator}{,}\PYG{+w}{ }\PYG{n+nv}{1}\PYG{+w}{ }\PYG{p+pIndicator}{,}\PYG{+w}{  }\PYG{n+nv}{0.000000e+00}\PYG{+w}{ }\PYG{p+pIndicator}{,}\PYG{+w}{  }\PYG{n+nv}{0.000000e+00}\PYG{+w}{ }\PYG{p+pIndicator}{,}\PYG{+w}{  }\PYG{n+nv}{0.000000e+00}\PYG{+w}{ }\PYG{p+pIndicator}{,}\PYG{+w}{  }\PYG{n+nv}{\PYGZhy{}1.841579e\PYGZhy{}01}\PYG{+w}{ }\PYG{p+pIndicator}{,}\PYG{+w}{ }\PYG{n+nv}{\PYGZhy{}9.710036e\PYGZhy{}01}\PYG{+w}{ }\PYG{p+pIndicator}{,}\PYG{+w}{ }\PYG{n+nv}{\PYGZhy{}2.934617e+00}\PYG{+w}{ }\PYG{p+pIndicator}{,}\PYG{+w}{ }\PYG{n+nv}{0}\PYG{+w}{ }\PYG{p+pIndicator}{,}\PYG{+w}{ }\PYG{n+nv}{0}\PYG{+w}{ }\PYG{p+pIndicator}{,}\PYG{+w}{ }\PYG{n+nv}{0}\PYG{p+pIndicator}{,}\PYG{+w}{ }\PYG{p+pIndicator}{]}
\PYG{+w}{  }\PYG{p+pIndicator}{\PYGZhy{}}\PYG{+w}{ }\PYG{p+pIndicator}{[}\PYG{+w}{     }\PYG{n+nv}{2}\PYG{+w}{ }\PYG{p+pIndicator}{,}\PYG{+w}{ }\PYG{n+nv}{1}\PYG{+w}{ }\PYG{p+pIndicator}{,}\PYG{+w}{  }\PYG{n+nv}{8.397981e\PYGZhy{}01}\PYG{+w}{ }\PYG{p+pIndicator}{,}\PYG{+w}{  }\PYG{n+nv}{8.397981e\PYGZhy{}01}\PYG{+w}{ }\PYG{p+pIndicator}{,}\PYG{+w}{  }\PYG{n+nv}{0.000000e+00}\PYG{+w}{ }\PYG{p+pIndicator}{,}\PYG{+w}{  }\PYG{n+nv}{\PYGZhy{}1.799591e+00}\PYG{+w}{ }\PYG{p+pIndicator}{,}\PYG{+w}{  }\PYG{n+nv}{2.127197e+00}\PYG{+w}{ }\PYG{p+pIndicator}{,}\PYG{+w}{  }\PYG{n+nv}{2.298572e+00}\PYG{+w}{ }\PYG{p+pIndicator}{,}\PYG{+w}{ }\PYG{n+nv}{0}\PYG{+w}{ }\PYG{p+pIndicator}{,}\PYG{+w}{ }\PYG{n+nv}{0}\PYG{+w}{ }\PYG{p+pIndicator}{,}\PYG{+w}{ }\PYG{n+nv}{0}\PYG{p+pIndicator}{,}\PYG{+w}{ }\PYG{p+pIndicator}{]}
\PYG{+w}{  }\PYG{p+pIndicator}{\PYGZhy{}}\PYG{+w}{ }\PYG{p+pIndicator}{[}\PYG{+w}{     }\PYG{n+nv}{3}\PYG{+w}{ }\PYG{p+pIndicator}{,}\PYG{+w}{ }\PYG{n+nv}{1}\PYG{+w}{ }\PYG{p+pIndicator}{,}\PYG{+w}{  }\PYG{n+nv}{8.397981e\PYGZhy{}01}\PYG{+w}{ }\PYG{p+pIndicator}{,}\PYG{+w}{  }\PYG{n+nv}{0.000000e+00}\PYG{+w}{ }\PYG{p+pIndicator}{,}\PYG{+w}{  }\PYG{n+nv}{8.397981e\PYGZhy{}01}\PYG{+w}{ }\PYG{p+pIndicator}{,}\PYG{+w}{  }\PYG{n+nv}{\PYGZhy{}1.807682e+00}\PYG{+w}{ }\PYG{p+pIndicator}{,}\PYG{+w}{ }\PYG{n+nv}{\PYGZhy{}9.585130e\PYGZhy{}01}\PYG{+w}{ }\PYG{p+pIndicator}{,}\PYG{+w}{  }\PYG{n+nv}{1.605884e+00}\PYG{+w}{ }\PYG{p+pIndicator}{,}\PYG{+w}{ }\PYG{n+nv}{0}\PYG{+w}{ }\PYG{p+pIndicator}{,}\PYG{+w}{ }\PYG{n+nv}{0}\PYG{+w}{ }\PYG{p+pIndicator}{,}\PYG{+w}{ }\PYG{n+nv}{0}\PYG{p+pIndicator}{,}\PYG{+w}{ }\PYG{p+pIndicator}{]}

\PYG{+w}{  }\PYG{p+pIndicator}{[}\PYG{n+nv}{...}\PYG{p+pIndicator}{]}
\PYG{n+nn}{...}
\PYG{n+nn}{\PYGZhy{}\PYGZhy{}\PYGZhy{}}
\PYG{n+nt}{timestep}\PYG{p}{:}\PYG{+w}{ }\PYG{l+lScalar+lScalarPlain}{100}
\PYG{n+nt}{units}\PYG{p}{:}\PYG{+w}{ }\PYG{l+lScalar+lScalarPlain}{lj}
\PYG{n+nt}{time}\PYG{p}{:}\PYG{+w}{ }\PYG{l+lScalar+lScalarPlain}{0.5}

\PYG{+w}{  }\PYG{l+lScalar+lScalarPlain}{[...]}

\PYG{n+nn}{...}
\end{sphinxVerbatim}


\bigskip\hrule\bigskip


\sphinxAtStartPar
Frequency of dump output:

\sphinxAtStartPar
Dumps are performed on timesteps that are a multiple of \(N\)
(including timestep 0) and on the last timestep of a minimization if
the minimization converges.  Note that this means a dump will not be
performed on the initial timestep after the dump command is invoked,
if the current timestep is not a multiple of \(N\).  This behavior
can be changed via the {\hyperref[\detokenize{dump_modify::doc}]{\sphinxcrossref{\DUrole{doc}{dump\_modify first}}}} command,
which can also be useful if the dump command is invoked after a
minimization ended on an arbitrary timestep.

\sphinxAtStartPar
The value of \(N\) can be changed between runs by using the
{\hyperref[\detokenize{dump_modify::doc}]{\sphinxcrossref{\DUrole{doc}{dump\_modify every}}}} command (not allowed for \sphinxstyleemphasis{dcd}
style).  The {\hyperref[\detokenize{dump_modify::doc}]{\sphinxcrossref{\DUrole{doc}{dump\_modify every}}}} command also
allows a variable to be used to determine the sequence of timesteps on
which dump files are written.  In this mode a dump on the first
timestep of a run will also not be written unless the
{\hyperref[\detokenize{dump_modify::doc}]{\sphinxcrossref{\DUrole{doc}{dump\_modify first}}}} command is used.

\sphinxAtStartPar
If you instead want to dump snapshots based on simulation time (in
time units of the {\hyperref[\detokenize{units::doc}]{\sphinxcrossref{\DUrole{doc}{units command}}}} command), the {\hyperref[\detokenize{dump_modify::doc}]{\sphinxcrossref{\DUrole{doc}{dump\_modify
every/time}}}} command can be used.  This can be useful
when the timestep size varies during a simulation run, e.g. by use of
the {\hyperref[\detokenize{fix_dt_reset::doc}]{\sphinxcrossref{\DUrole{doc}{fix dt/reset}}}} command.


\bigskip\hrule\bigskip


\sphinxAtStartPar
Dump filenames:

\sphinxAtStartPar
The specified dump filename determines how the dump file(s) is
written.  The default is to write one large text file, which is opened
when the dump command is invoked and closed when an {\hyperref[\detokenize{undump::doc}]{\sphinxcrossref{\DUrole{doc}{undump}}}} command is used or when LAMMPS exits.  For the \sphinxstyleemphasis{dcd} and
\sphinxstyleemphasis{xtc} styles, this is a single large binary file.

\sphinxAtStartPar
Many of the styles allow dump filenames to contain either or both of
two wildcard characters.  If a “*” character appears in the filename,
then one file per snapshot is written and the “*” character is
replaced with the timestep value.  For example, tmp.dump.* becomes
tmp.dump.0, tmp.dump.10000, tmp.dump.20000, etc.  This option is not
available for the \sphinxstyleemphasis{dcd} and \sphinxstyleemphasis{xtc} styles.  Note that the
{\hyperref[\detokenize{dump_modify::doc}]{\sphinxcrossref{\DUrole{doc}{dump\_modify pad}}}} command can be used to ensure all
timestep numbers are the same length (e.g., 00010), which can make it
easier to read a series of dump files in order with some
post\sphinxhyphen{}processing tools.

\sphinxAtStartPar
If a “\%” character appears in the filename, then each of P processors
writes a portion of the dump file, and the “\%” character is replaced
with the processor ID from \(0\) to \(P-1\).  For example,
tmp.dump.\% becomes tmp.dump.0, tmp.dump.1, … tmp.dump.:math:\sphinxtitleref{P\sphinxhyphen{}1},
etc.  This creates smaller files and can be a fast mode of output on
parallel machines that support parallel I/O for output. This option is
\sphinxstylestrong{not} available for the \sphinxstyleemphasis{dcd}, \sphinxstyleemphasis{xtc}, \sphinxstyleemphasis{xyz}, \sphinxstyleemphasis{grid/vtk}, and \sphinxstyleemphasis{yaml}
styles.

\sphinxAtStartPar
By default, \(P\) is the the number of processors, meaning one file per
processor, but \(P\) can be set to a smaller value via the \sphinxstyleemphasis{nfile} or
\sphinxstyleemphasis{fileper} keywords of the {\hyperref[\detokenize{dump_modify::doc}]{\sphinxcrossref{\DUrole{doc}{dump\_modify}}}} command.
These options can be the most efficient way of writing out dump files
when running on large numbers of processors.

\sphinxAtStartPar
Note that using the “*” and “\%” characters together can produce a
large number of small dump files!

\sphinxAtStartPar
\DUrole{versionmodified,deprecated}{Deprecated since version 21Nov2023.}

\sphinxAtStartPar
The MPIIO package and the the corresponding “/mpiio” dump styles, except
for the unrelated “netcdf/mpiio” style were removed from LAMMPS.


\bigskip\hrule\bigskip


\sphinxAtStartPar
Compression of dump file data:

\sphinxAtStartPar
If the specified filename ends with “.bin” or “.lammpsbin”, the dump
file (or files, if “*” or “\%” is also used) is written in binary
format.  A binary dump file will be about the same size as a text
version, but will typically write out much faster.  Of course, when
post\sphinxhyphen{}processing, you will need to convert it back to text format (see
the {\hyperref[\detokenize{Tools:binary}]{\sphinxcrossref{\DUrole{std,std-ref}{binary2txt tool}}}}) or write your own code to read
the binary file.  The format of the binary file can be understood by
looking at the \sphinxcode{\sphinxupquote{tools/binary2txt.cpp}} file.  This option is only
available for the \sphinxstyleemphasis{atom} and \sphinxstyleemphasis{custom} styles.

\sphinxAtStartPar
If the filename ends with “.gz”, the dump file (or files, if “*” or “\%”
is also used) is written in gzipped format.  A gzipped dump file will be
about \(3\times\) smaller than the text version, but will also take
longer to write.  This option is not available for the \sphinxstyleemphasis{dcd} and \sphinxstyleemphasis{xtc}
styles.

\sphinxAtStartPar
Note that styles that end with \sphinxstyleemphasis{gz} are identical in command syntax to
the corresponding styles without “gz”, however, they generate
compressed files using the zlib library. Thus the filename suffix
“.gz” is mandatory. This is an alternative approach to writing
compressed files via a pipe, as done by the regular dump styles, which
may be required on clusters where the interface to the high\sphinxhyphen{}speed
network disallows using the fork() library call (which is needed for a
pipe).  For the remainder of this page, you should thus consider the
\sphinxstyleemphasis{atom} and \sphinxstyleemphasis{atom/gz} styles (etc.) to be inter\sphinxhyphen{}changeable, with the
exception of the required filename suffix.

\sphinxAtStartPar
Similarly, styles that end with \sphinxstyleemphasis{zstd} are identical to the gz styles,
but use the Zstd compression library instead and require a “.zst”
suffix. See the {\hyperref[\detokenize{dump_modify::doc}]{\sphinxcrossref{\DUrole{doc}{dump\_modify}}}} page for details on
how to control the compression level in both variants.


\bigskip\hrule\bigskip


\sphinxAtStartPar
General triclinic simulation box output for the \sphinxstyleemphasis{atom} and \sphinxstyleemphasis{custom} styles:

\sphinxAtStartPar
As mentioned above, the simulation box can be defined as a general
triclinic box, which means that 3 arbitrary box edge vectors \sphinxstylestrong{A},
\sphinxstylestrong{B}, \sphinxstylestrong{C} can be specified.  See the {\hyperref[\detokenize{Howto_triclinic::doc}]{\sphinxcrossref{\DUrole{doc}{Howto triclinic}}}} doc page for a detailed description of general
triclinic boxes.

\sphinxAtStartPar
This option is provided as a convenience for users who may be
converting data from solid\sphinxhyphen{}state crystallographic representations or
from DFT codes for input to LAMMPS.  However, as explained on the
{\hyperref[\detokenize{Howto_triclinic::doc}]{\sphinxcrossref{\DUrole{doc}{Howto\_triclinic}}}} doc page, internally, LAMMPS
only uses restricted triclinic simulation boxes.  This means the box
and per\sphinxhyphen{}atom information (e.g. coordinates, velocities) LAMMPS stores
are converted (rotated) from general to restricted triclinic form when
the system is created.

\sphinxAtStartPar
For dump output, if the {\hyperref[\detokenize{dump_modify::doc}]{\sphinxcrossref{\DUrole{doc}{dump\_modify triclinic/general}}}} command is used, the box description and per\sphinxhyphen{}atom
coordinates and other per\sphinxhyphen{}atom vectors will be converted (rotated)
from restricted to general form when each dump file snapshots is
output.  This option can only be used if the simulation box was
initially created as general triclinic.  If the option is not used,
and the simulation box is general triclinic, then the dump file
snapshots will reflect the internal restricted triclinic geometry.

\sphinxAtStartPar
The dump\_modify triclinic/general option affects 3 aspects of the dump
file output.

\sphinxAtStartPar
First, the format for the BOX BOUNDS is as follows

\begin{sphinxVerbatim}[commandchars=\\\{\}]
ITEM: BOX BOUNDS abc origin
ax ay az originx
bx by bz originy
cx cy cz originz
\end{sphinxVerbatim}

\sphinxAtStartPar
where the \sphinxstylestrong{A} edge vector of the box is (ax,ay,az) and similarly
for \sphinxstylestrong{B} and \sphinxstylestrong{C}.  The origin of all 3 edge vectors is (originx,
originy, originz).

\sphinxAtStartPar
Second, the coordinates of each atom are converted (rotated) so that
the atom is inside (or near) the general triclinic box defined by the
\sphinxstylestrong{A}, \sphinxstylestrong{B}, \sphinxstylestrong{C} edge vectors.  For style \sphinxstyleemphasis{atom}, this only alters
output for unscaled atom coords, via the {\hyperref[\detokenize{dump_modify::doc}]{\sphinxcrossref{\DUrole{doc}{dump\_modify scaled no}}}} setting. For style \sphinxstyleemphasis{custom}, this alters output for
either unscaled or unwrapped output of atom coords, via the \sphinxstyleemphasis{x,y,z} or
\sphinxstyleemphasis{xu,yu,zu} attributes.  For output of scaled atom coords by both
styles, there is no difference between restricted and general
triclinic values.

\sphinxAtStartPar
Third, the output for any attribute of the \sphinxstyleemphasis{custom} style which
represents a per\sphinxhyphen{}atom vector quantity will be converted (rotated) to
be oriented consistent with the general triclinic box and its
orientation relative to the standard xyz coordinate axes.

\sphinxAtStartPar
This applies to the following \sphinxstyleemphasis{custom} style attributes:
\begin{itemize}
\item {} 
\sphinxAtStartPar
vx,vy,vz = atom velocities

\item {} 
\sphinxAtStartPar
fx,fy,fz = forces on atoms

\item {} 
\sphinxAtStartPar
mux,muy,muz = orientation of dipole moment of atom

\item {} 
\sphinxAtStartPar
omegax,omegay,omegaz = angular velocity of spherical particle

\item {} 
\sphinxAtStartPar
angmomx,angmomy,angmomz = angular momentum of aspherical particle

\item {} 
\sphinxAtStartPar
tqx,tqy,tqz = torque on finite\sphinxhyphen{}size particles

\end{itemize}

\sphinxAtStartPar
For example, if the velocity of an atom in a restricted triclinic box
is along the x\sphinxhyphen{}axis, then it will be output for a general triclinic
box as a vector along the \sphinxstylestrong{A} edge vector of the box.

\begin{sphinxadmonition}{note}{Note:}
\sphinxAtStartPar
For style \sphinxstyleemphasis{custom}, the {\hyperref[\detokenize{dump_modify::doc}]{\sphinxcrossref{\DUrole{doc}{dump\_modify thresh}}}}
command may access per\sphinxhyphen{}atom attributes either directly or
indirectly through a compute or variable.  If the attribute is an
atom coordinate or one of the vectors mentioned above, its value
will \sphinxstyleemphasis{NOT} be a general triclinic (rotated) value.  Rather it will
be a restricted triclinic value.
\end{sphinxadmonition}


\bigskip\hrule\bigskip


\sphinxAtStartPar
Arguments for different styles:

\sphinxAtStartPar
The sections below describe per\sphinxhyphen{}atom, local, and per grid cell
attributes which can be used as arguments to the various styles.

\sphinxAtStartPar
Note that in the discussion below, for styles which can reference
values from a compute or fix or custom atom property, like the
\sphinxstyleemphasis{custom}, \sphinxstyleemphasis{cfg}, \sphinxstyleemphasis{grid} or \sphinxstyleemphasis{local} styles, the bracketed index
\(i\) can be specified using a wildcard asterisk with the index to
effectively specify multiple values.  This takes the form “*” or
“*n” or “m*” or “m*n”.  If \(N\) is the number of columns in
the array, then an asterisk with no numeric values means all column
indices from 1 to \(N\).  A leading asterisk means all indices
from 1 to n (inclusive).  A trailing asterisk means all indices from m
to \(N\) (inclusive).  A middle asterisk means all indices from m
to n (inclusive).

\sphinxAtStartPar
Using a wildcard is the same as if the individual columns of the array
had been listed one by one.  For example, these two dump commands are
equivalent, since the {\hyperref[\detokenize{compute_stress_atom::doc}]{\sphinxcrossref{\DUrole{doc}{compute stress/atom}}}}
command creates a per\sphinxhyphen{}atom array with six columns:

\begin{sphinxVerbatim}[commandchars=\\\{\}]
\PYG{k}{compute }\PYG{n+nv+nvIdentifier}{myPress}\PYG{+w}{ }\PYG{n+nv+nvIdentifier}{all}\PYG{+w}{ }\PYG{n}{stress}\PYG{o}{/}\PYG{n}{atom}\PYG{+w}{ }\PYG{n}{NULL}
\PYG{k}{dump }\PYG{n+nv+nvIdentifier}{2}\PYG{+w}{ }\PYG{n+nv+nvIdentifier}{all}\PYG{+w}{ }\PYG{n}{custom}\PYG{+w}{ }\PYG{l+m}{100}\PYG{+w}{ }\PYG{n}{tmp.dump}\PYG{+w}{ }\PYG{n}{id}\PYG{+w}{ }\PYG{n}{myPress[}\PYG{o}{*}\PYG{n}{]}
\PYG{k}{dump }\PYG{n+nv+nvIdentifier}{2}\PYG{+w}{ }\PYG{n+nv+nvIdentifier}{all}\PYG{+w}{ }\PYG{n}{custom}\PYG{+w}{ }\PYG{l+m}{100}\PYG{+w}{ }\PYG{n}{tmp.dump}\PYG{+w}{ }\PYG{n}{id}\PYG{+w}{ }\PYG{n}{myPress[1]}\PYG{+w}{ }\PYG{n}{myPress[2]}\PYG{+w}{ }\PYG{n}{myPress[3]}\PYG{+w}{ }\PYG{o}{\PYGZam{}}
\PYG{+w}{                                  }\PYG{n}{myPress[4]}\PYG{+w}{ }\PYG{n}{myPress[5]}\PYG{+w}{ }\PYG{n}{myPress[6]}
\end{sphinxVerbatim}


\bigskip\hrule\bigskip


\sphinxAtStartPar
Per\sphinxhyphen{}atom attributes used as arguments to the \sphinxstyleemphasis{custom} and \sphinxstyleemphasis{cfg} styles:

\sphinxAtStartPar
The \sphinxstyleemphasis{id}, \sphinxstyleemphasis{mol}, \sphinxstyleemphasis{proc}, \sphinxstyleemphasis{procp1}, \sphinxstyleemphasis{type}, \sphinxstyleemphasis{typelabel}, \sphinxstyleemphasis{element}, \sphinxstyleemphasis{mass},
\sphinxstyleemphasis{vx}, \sphinxstyleemphasis{vy}, \sphinxstyleemphasis{vz}, \sphinxstyleemphasis{fx}, \sphinxstyleemphasis{fy}, \sphinxstyleemphasis{fz}, \sphinxstyleemphasis{q} attributes are self\sphinxhyphen{}explanatory.

\sphinxAtStartPar
\sphinxstyleemphasis{Id} is the atom ID.  \sphinxstyleemphasis{Mol} is the molecule ID, included in the data file
for molecular systems.  \sphinxstyleemphasis{Proc} is the ID of the processor (0 to
\(N_\text{procs}-1\)) that currently owns the atom.  \sphinxstyleemphasis{Procp1} is the
proc ID+1, which can be convenient in place of a \sphinxstyleemphasis{type} attribute (1 to
\(N_\text{types}\)) for coloring atoms in a visualization program.
\sphinxstyleemphasis{Type} is the atom type (1 to \(N_\text{types}\)).  \sphinxstyleemphasis{Typelabel} is the
atom {\hyperref[\detokenize{Howto_type_labels::doc}]{\sphinxcrossref{\DUrole{doc}{type label}}}}.  \sphinxstyleemphasis{Element} is typically the
chemical name of an element, which you must assign to each type via the
{\hyperref[\detokenize{dump_modify::doc}]{\sphinxcrossref{\DUrole{doc}{dump\_modify element}}}} command.  More generally, it can
be any string you wish to associated with an atom type.  \sphinxstyleemphasis{Mass} is the atom
mass.  The quantities \sphinxstyleemphasis{vx}, \sphinxstyleemphasis{vy}, \sphinxstyleemphasis{vz}, \sphinxstyleemphasis{fx}, \sphinxstyleemphasis{fy}, \sphinxstyleemphasis{fz}, and \sphinxstyleemphasis{q} are
components of atom velocity and force and atomic charge.

\sphinxAtStartPar
There are several options for outputting atom coordinates.  The \sphinxstyleemphasis{x},
\sphinxstyleemphasis{y}, and \sphinxstyleemphasis{z} attributes write atom coordinates “unscaled”, in the
appropriate distance {\hyperref[\detokenize{units::doc}]{\sphinxcrossref{\DUrole{doc}{units}}}} (\(\AA\),
\(\sigma\), etc.).  Use \sphinxstyleemphasis{xs}, \sphinxstyleemphasis{ys}, and \sphinxstyleemphasis{zs} if you want the
coordinates “scaled” to the box size so that each value is 0.0 to 1.0.
If the simulation box is triclinic (tilted), then all atom coords will
still be between 0.0 and 1.0.  The actual unscaled \((x,y,z)\)
coordinate is \(x_s a + y_s b + z_s c\), where \((a,b,c)\) are
the non\sphinxhyphen{}orthogonal vectors of the simulation box edges, as discussed on
the {\hyperref[\detokenize{Howto_triclinic::doc}]{\sphinxcrossref{\DUrole{doc}{Howto triclinic}}}} page.

\sphinxAtStartPar
Use \sphinxstyleemphasis{xu}, \sphinxstyleemphasis{yu}, and \sphinxstyleemphasis{zu} if you want the coordinates “unwrapped” by the
image flags for each atom.  Unwrapped means that if the atom has passed
through a periodic boundary one or more times, the value is printed for
what the coordinate would be if it had not been wrapped back into the
periodic box.  Note that using \sphinxstyleemphasis{xu}, \sphinxstyleemphasis{yu}, and \sphinxstyleemphasis{zu} means that the
coordinate values may be far outside the box bounds printed with the
snapshot.  Using \sphinxstyleemphasis{xsu}, \sphinxstyleemphasis{ysu}, and \sphinxstyleemphasis{zsu} is similar to using \sphinxstyleemphasis{xu}, \sphinxstyleemphasis{yu},
and \sphinxstyleemphasis{zu}, except that the unwrapped coordinates are scaled by the box
size. Atoms that have passed through a periodic boundary will have the
corresponding coordinate increased or decreased by 1.0.

\sphinxAtStartPar
The image flags can be printed directly using the \sphinxstyleemphasis{ix}, \sphinxstyleemphasis{iy}, and \sphinxstyleemphasis{iz}
attributes.  For periodic dimensions, they specify which image of the
simulation box the atom is considered to be in.  An image of 0 means
it is inside the box as defined.  A value of 2 means add 2 box lengths
to get the true value.  A value of \(-1\) means subtract 1 box length to
get the true value.  LAMMPS updates these flags as atoms cross
periodic boundaries during the simulation.

\sphinxAtStartPar
The \sphinxstyleemphasis{mux}, \sphinxstyleemphasis{muy}, and \sphinxstyleemphasis{muz} attributes are specific to dipolar systems
defined with an atom style of \sphinxstyleemphasis{dipole}.  They give the orientation of
the atom’s point dipole moment.  The \sphinxstyleemphasis{mu} attribute gives the magnitude
of the atom’s dipole moment.

\sphinxAtStartPar
The \sphinxstyleemphasis{radius} and \sphinxstyleemphasis{diameter} attributes are specific to spherical
particles that have a finite size, such as those defined with an atom
style of \sphinxstyleemphasis{sphere}.

\sphinxAtStartPar
The \sphinxstyleemphasis{omegax}, \sphinxstyleemphasis{omegay}, and \sphinxstyleemphasis{omegaz} attributes are specific to
finite\sphinxhyphen{}size spherical particles that have an angular velocity.  Only
certain atom styles, such as \sphinxstyleemphasis{sphere}, define this quantity.

\sphinxAtStartPar
The \sphinxstyleemphasis{angmomx}, \sphinxstyleemphasis{angmomy}, and \sphinxstyleemphasis{angmomz} attributes are specific to
finite\sphinxhyphen{}size aspherical particles that have an angular momentum.  Only
the \sphinxstyleemphasis{ellipsoid} atom style defines this quantity.

\sphinxAtStartPar
The \sphinxstyleemphasis{tqx}, \sphinxstyleemphasis{tqy}, and \sphinxstyleemphasis{tqz} attributes are for finite\sphinxhyphen{}size particles
that can sustain a rotational torque due to interactions with other
particles.

\sphinxAtStartPar
The \sphinxstyleemphasis{c\_ID} and \sphinxstyleemphasis{c\_ID{[}I{]}} attributes allow per\sphinxhyphen{}atom vectors or arrays
calculated by a {\hyperref[\detokenize{compute::doc}]{\sphinxcrossref{\DUrole{doc}{compute}}}} to be output.  The ID in the
attribute should be replaced by the actual ID of the compute that has
been defined previously in the input script.  See the {\hyperref[\detokenize{compute::doc}]{\sphinxcrossref{\DUrole{doc}{compute}}}} command for details.  There are computes for calculating the
per\sphinxhyphen{}atom energy, stress, centro\sphinxhyphen{}symmetry parameter, and coordination
number of individual atoms.

\sphinxAtStartPar
Note that computes which calculate global or local quantities, as
opposed to per\sphinxhyphen{}atom quantities, cannot be output in a dump custom
command.  Instead, global quantities can be output by the
{\hyperref[\detokenize{thermo_style::doc}]{\sphinxcrossref{\DUrole{doc}{thermo\_style custom}}}} command, and local quantities
can be output by the dump local command.

\sphinxAtStartPar
If \sphinxstyleemphasis{c\_ID} is used as a attribute, then the per\sphinxhyphen{}atom vector calculated by
the compute is printed.  If \sphinxstyleemphasis{c\_ID{[}i{]}} is used, then \(i\) must be in
the range from 1 to \(M\), which will print the \(i\)th column
of the per\sphinxhyphen{}atom array with \(M\) columns calculated by the compute.
See the discussion above for how \(i\) can be specified with a
wildcard asterisk to effectively specify multiple values.

\sphinxAtStartPar
The \sphinxstyleemphasis{f\_ID} and \sphinxstyleemphasis{f\_ID{[}I{]}} attributes allow vector or array per\sphinxhyphen{}atom
quantities calculated by a {\hyperref[\detokenize{fix::doc}]{\sphinxcrossref{\DUrole{doc}{fix}}}} to be output.  The ID in the
attribute should be replaced by the actual ID of the fix that has been
defined previously in the input script.  The {\hyperref[\detokenize{fix_ave_atom::doc}]{\sphinxcrossref{\DUrole{doc}{fix ave/atom}}}} command is one that calculates per\sphinxhyphen{}atom quantities.
Since it can time\sphinxhyphen{}average per\sphinxhyphen{}atom quantities produced by any
{\hyperref[\detokenize{compute::doc}]{\sphinxcrossref{\DUrole{doc}{compute}}}}, {\hyperref[\detokenize{fix::doc}]{\sphinxcrossref{\DUrole{doc}{fix}}}}, or atom\sphinxhyphen{}style {\hyperref[\detokenize{variable::doc}]{\sphinxcrossref{\DUrole{doc}{variable}}}}, this allows those time\sphinxhyphen{}averaged results to be written to a
dump file.

\sphinxAtStartPar
If \sphinxstyleemphasis{f\_ID} is used as a attribute, then the per\sphinxhyphen{}atom vector calculated by
the fix is printed.  If \sphinxstyleemphasis{f\_ID{[}i{]}} is used, then \(i\) must be in the
range from 1 to \(M\), which will print the \(i\)th column of
the per\sphinxhyphen{}atom array with \(M\) columns calculated by the fix.  See
the discussion above for how \(i\) can be specified with a wildcard
asterisk to effectively specify multiple values.

\sphinxAtStartPar
The \sphinxstyleemphasis{v\_name} attribute allows per\sphinxhyphen{}atom vectors calculated by a
{\hyperref[\detokenize{variable::doc}]{\sphinxcrossref{\DUrole{doc}{variable}}}} to be output.  The name in the attribute
should be replaced by the actual name of the variable that has been
defined previously in the input script.  Only an atom\sphinxhyphen{}style variable can
be referenced, since it is the only style that generates per\sphinxhyphen{}atom
values.  Variables of style \sphinxstyleemphasis{atom} can reference individual atom
attributes, per\sphinxhyphen{}atom attributes, thermodynamic keywords, or invoke other
computes, fixes, or variables when they are evaluated, so this is a very
general means of creating quantities to output to a dump file.

\sphinxAtStartPar
The \sphinxstyleemphasis{i\_name}, \sphinxstyleemphasis{d\_name}, \sphinxstyleemphasis{i2\_name}, \sphinxstyleemphasis{d2\_name} attributes refer to
custom per\sphinxhyphen{}atom integer and floating\sphinxhyphen{}point vectors or arrays that have
been added via the {\hyperref[\detokenize{fix_property_atom::doc}]{\sphinxcrossref{\DUrole{doc}{fix property/atom}}}}
command.  When that command is used specific names are given to each
attribute which are the “name” portion of these keywords.  For arrays
\sphinxstyleemphasis{i2\_name} and \sphinxstyleemphasis{d2\_name}, the column of the array must also be included
following the name in brackets (e.g., d2\_xyz{[}i{]}, i2\_mySpin{[}i{]}, where
\(i\) is in the range from 1 to \(M\), where \(M\) is the
number of columns in the custom array).  See the discussion above for
how \(i\) can be specified with a wildcard asterisk to effectively
specify multiple values.

\sphinxAtStartPar
See the {\hyperref[\detokenize{Modify::doc}]{\sphinxcrossref{\DUrole{doc}{Modify}}}} page for information on how to add
new compute and fix styles to LAMMPS to calculate per\sphinxhyphen{}atom quantities
which could then be output into dump files.


\bigskip\hrule\bigskip


\sphinxAtStartPar
Attributes used as arguments to the \sphinxstyleemphasis{local} style:

\sphinxAtStartPar
The \sphinxstyleemphasis{index} attribute can be used to generate an index number from 1
to N for each line written into the dump file, where N is the total
number of local datums from all processors, or lines of output that
will appear in the snapshot.  Note that because data from different
processors depend on what atoms they currently own, and atoms migrate
between processor, there is no guarantee that the same index will be
used for the same info (e.g. a particular bond) in successive
snapshots.

\sphinxAtStartPar
The \sphinxstyleemphasis{c\_ID} and \sphinxstyleemphasis{c\_ID{[}I{]}} attributes allow local vectors or arrays
calculated by a {\hyperref[\detokenize{compute::doc}]{\sphinxcrossref{\DUrole{doc}{compute}}}} to be output.  The ID in the
attribute should be replaced by the actual ID of the compute that has
been defined previously in the input script.  See the
{\hyperref[\detokenize{compute::doc}]{\sphinxcrossref{\DUrole{doc}{compute}}}} command for details.  There are computes for
calculating local information such as indices, types, and energies for
bonds and angles.

\sphinxAtStartPar
Note that computes which calculate global or per\sphinxhyphen{}atom quantities, as
opposed to local quantities, cannot be output in a dump local command.
Instead, global quantities can be output by the {\hyperref[\detokenize{thermo_style::doc}]{\sphinxcrossref{\DUrole{doc}{thermo\_style
custom}}}} command, and per\sphinxhyphen{}atom quantities can be output
by the dump custom command.

\sphinxAtStartPar
If \sphinxstyleemphasis{c\_ID} is used as a attribute, then the local vector calculated by
the compute is printed.  If \sphinxstyleemphasis{c\_ID{[}I{]}} is used, then I must be in the
range from 1\sphinxhyphen{}M, which will print the Ith column of the local array
with M columns calculated by the compute.  See the discussion above
for how I can be specified with a wildcard asterisk to effectively
specify multiple values.

\sphinxAtStartPar
The \sphinxstyleemphasis{f\_ID} and \sphinxstyleemphasis{f\_ID{[}I{]}} attributes allow local vectors or arrays
calculated by a {\hyperref[\detokenize{fix::doc}]{\sphinxcrossref{\DUrole{doc}{fix}}}} to be output.  The ID in the attribute
should be replaced by the actual ID of the fix that has been defined
previously in the input script.

\sphinxAtStartPar
If \sphinxstyleemphasis{f\_ID} is used as a attribute, then the local vector calculated by
the fix is printed.  If \sphinxstyleemphasis{f\_ID{[}I{]}} is used, then I must be in the
range from 1\sphinxhyphen{}M, which will print the Ith column of the local with M
columns calculated by the fix.  See the discussion above for how I can
be specified with a wildcard asterisk to effectively specify multiple
values.

\sphinxAtStartPar
Here is an example of how to dump bond info for a system, including
the distance and energy of each bond:

\begin{sphinxVerbatim}[commandchars=\\\{\}]
\PYG{k}{compute }\PYG{n+nv+nvIdentifier}{1}\PYG{+w}{ }\PYG{n+nv+nvIdentifier}{all}\PYG{+w}{ }\PYG{n}{property}\PYG{o}{/}\PYG{n}{local}\PYG{+w}{ }\PYG{n}{batom1}\PYG{+w}{ }\PYG{n}{batom2}\PYG{+w}{ }\PYG{n}{btype}
\PYG{k}{compute }\PYG{n+nv+nvIdentifier}{2}\PYG{+w}{ }\PYG{n+nv+nvIdentifier}{all}\PYG{+w}{ }\PYG{n}{bond}\PYG{o}{/}\PYG{n}{local}\PYG{+w}{ }\PYG{n}{dist}\PYG{+w}{ }\PYG{n}{eng}
\PYG{k}{dump }\PYG{n+nv+nvIdentifier}{1}\PYG{+w}{ }\PYG{n+nv+nvIdentifier}{all}\PYG{+w}{ }\PYG{n}{local}\PYG{+w}{ }\PYG{l+m}{1000}\PYG{+w}{ }\PYG{n}{tmp.dump}\PYG{+w}{ }\PYG{n}{index}\PYG{+w}{ }\PYG{n}{c\PYGZus{}1[1]}\PYG{+w}{ }\PYG{n}{c\PYGZus{}1[2]}\PYG{+w}{ }\PYG{n}{c\PYGZus{}1[3]}\PYG{+w}{ }\PYG{n}{c\PYGZus{}2[1]}\PYG{+w}{ }\PYG{n}{c\PYGZus{}2[2]}
\end{sphinxVerbatim}


\bigskip\hrule\bigskip


\sphinxAtStartPar
Attributes used as arguments to the \sphinxstyleemphasis{grid} and \sphinxstyleemphasis{grid/vtk} styles:

\sphinxAtStartPar
The attributes that begin with \sphinxstyleemphasis{c\_ID} and \sphinxstyleemphasis{f\_ID} both take
colon\sphinxhyphen{}separated fields \sphinxstyleemphasis{gname} and \sphinxstyleemphasis{dname}.  These refer to a grid
name and data field name which is defined by the compute or fix.  Note
that a compute or fix can define one or more grids (of different
sizes) and one or more data fields for each of those grids.  The sizes
of all grids output in a single dump grid command must be the same.

\sphinxAtStartPar
The \sphinxstyleemphasis{c\_ID:gname:dname} and \sphinxstyleemphasis{c\_ID:gname:dname{[}I{]}} attributes allow
per\sphinxhyphen{}grid vectors or arrays calculated by a {\hyperref[\detokenize{compute::doc}]{\sphinxcrossref{\DUrole{doc}{compute}}}} to
be output.  The ID in the attribute should be replaced by the actual
ID of the compute that has been defined previously in the input
script.

\sphinxAtStartPar
If \sphinxstyleemphasis{c\_ID:gname:dname} is used as a attribute, then the per\sphinxhyphen{}grid vector
calculated by the compute is printed.  If \sphinxstyleemphasis{c\_ID:gname:dname{[}I{]}} is
used, then I must be in the range from 1\sphinxhyphen{}M, which will print the Ith
column of the per\sphinxhyphen{}grid array with M columns calculated by the compute.
See the discussion above for how I can be specified with a wildcard
asterisk to effectively specify multiple values.

\sphinxAtStartPar
The \sphinxstyleemphasis{f\_ID:gname:dname} and \sphinxstyleemphasis{f\_ID:gname:dname{[}I{]}} attributes allow
per\sphinxhyphen{}grid vectors or arrays calculated by a {\hyperref[\detokenize{fix::doc}]{\sphinxcrossref{\DUrole{doc}{fix}}}} to be
output.  The ID in the attribute should be replaced by the actual ID
of the fix that has been defined previously in the input script.

\sphinxAtStartPar
If \sphinxstyleemphasis{f\_ID:gname:dname} is used as a attribute, then the per\sphinxhyphen{}grid vector
calculated by the fix is printed.  If \sphinxstyleemphasis{f\_ID:gname:dname{[}I{]}} is used,
then I must be in the range from 1\sphinxhyphen{}M, which will print the Ith column
of the per\sphinxhyphen{}grid with M columns calculated by the fix.  See the
discussion above for how I can be specified with a wildcard asterisk
to effectively specify multiple values.


\bigskip\hrule\bigskip



\subsection{Restrictions}
\label{\detokenize{dump:restrictions}}
\sphinxAtStartPar
To write gzipped dump files, you must either compile LAMMPS with the
\sphinxhyphen{}DLAMMPS\_GZIP option or use the styles from the COMPRESS package.
See the {\hyperref[\detokenize{Build_settings::doc}]{\sphinxcrossref{\DUrole{doc}{Build settings}}}} page for details.

\sphinxAtStartPar
While a dump command is active (i.e., has not been stopped by using
the {\hyperref[\detokenize{undump::doc}]{\sphinxcrossref{\DUrole{doc}{undump command}}}}), no commands may be used that will
change the timestep (e.g., {\hyperref[\detokenize{reset_timestep::doc}]{\sphinxcrossref{\DUrole{doc}{reset\_timestep}}}}).
LAMMPS will terminate with an error otherwise.

\sphinxAtStartPar
The \sphinxstyleemphasis{atom/gz}, \sphinxstyleemphasis{cfg/gz}, \sphinxstyleemphasis{custom/gz}, and \sphinxstyleemphasis{xyz/gz} styles are part of
the COMPRESS package.  They are only enabled if LAMMPS was built with
that package.  See the {\hyperref[\detokenize{Build_package::doc}]{\sphinxcrossref{\DUrole{doc}{Build package}}}} page for
more info.

\sphinxAtStartPar
The \sphinxstyleemphasis{xtc}, \sphinxstyleemphasis{dcd}, and \sphinxstyleemphasis{yaml} styles are part of the EXTRA\sphinxhyphen{}DUMP package.
They are only enabled if LAMMPS was built with that package.  See the
{\hyperref[\detokenize{Build_package::doc}]{\sphinxcrossref{\DUrole{doc}{Build package}}}} page for more info.


\subsection{Related commands}
\label{\detokenize{dump:related-commands}}
\sphinxAtStartPar
{\hyperref[\detokenize{dump_adios::doc}]{\sphinxcrossref{\DUrole{doc}{dump atom/adios}}}}, {\hyperref[\detokenize{dump_adios::doc}]{\sphinxcrossref{\DUrole{doc}{dump custom/adios}}}},
{\hyperref[\detokenize{dump_cfg_uef::doc}]{\sphinxcrossref{\DUrole{doc}{dump cfg/uef}}}}, {\hyperref[\detokenize{dump_h5md::doc}]{\sphinxcrossref{\DUrole{doc}{dump h5md}}}},
{\hyperref[\detokenize{dump_image::doc}]{\sphinxcrossref{\DUrole{doc}{dump image}}}}, {\hyperref[\detokenize{dump_molfile::doc}]{\sphinxcrossref{\DUrole{doc}{dump molfile}}}},
{\hyperref[\detokenize{dump_netcdf::doc}]{\sphinxcrossref{\DUrole{doc}{dump netcdf}}}}, {\hyperref[\detokenize{dump_netcdf::doc}]{\sphinxcrossref{\DUrole{doc}{dump netcdf/mpiio}}}},
{\hyperref[\detokenize{dump_modify::doc}]{\sphinxcrossref{\DUrole{doc}{dump\_modify}}}}, {\hyperref[\detokenize{undump::doc}]{\sphinxcrossref{\DUrole{doc}{undump}}}},
{\hyperref[\detokenize{write_dump::doc}]{\sphinxcrossref{\DUrole{doc}{write\_dump}}}}


\subsection{Default}
\label{\detokenize{dump:default}}
\sphinxAtStartPar
The defaults for the \sphinxstyleemphasis{image} and \sphinxstyleemphasis{movie} styles are listed on the
{\hyperref[\detokenize{dump_image::doc}]{\sphinxcrossref{\DUrole{doc}{dump image}}}} page.

\sphinxstepscope

\index{dump atom/adios@\spxentry{dump atom/adios}}\ignorespaces 
\index{dump custom/adios@\spxentry{dump custom/adios}}\ignorespaces 

\section{dump atom/adios  command}
\label{\detokenize{dump_adios:dump-atom-adios-command}}\label{\detokenize{dump_adios:index-1}}\label{\detokenize{dump_adios:index-0}}\label{\detokenize{dump_adios::doc}}

\section{dump custom/adios command}
\label{\detokenize{dump_adios:dump-custom-adios-command}}

\subsection{Syntax}
\label{\detokenize{dump_adios:syntax}}
\begin{sphinxVerbatim}[commandchars=\\\{\}]
\PYG{k}{dump }\PYG{n+nv+nvIdentifier}{ID}\PYG{+w}{ }\PYG{n+nv+nvIdentifier}{group\PYGZhy{}ID}\PYG{+w}{ }\PYG{n}{atom}\PYG{o}{/}\PYG{n}{adios}\PYG{+w}{ }\PYG{n}{N}\PYG{+w}{ }\PYG{n}{file.bp}
\PYG{k}{dump }\PYG{n+nv+nvIdentifier}{ID}\PYG{+w}{ }\PYG{n+nv+nvIdentifier}{group\PYGZhy{}ID}\PYG{+w}{ }\PYG{n}{custom}\PYG{o}{/}\PYG{n}{adios}\PYG{+w}{ }\PYG{n}{N}\PYG{+w}{ }\PYG{n}{file.bp}\PYG{+w}{ }\PYG{n}{args}
\end{sphinxVerbatim}
\begin{itemize}
\item {} 
\sphinxAtStartPar
ID = user\sphinxhyphen{}assigned name for the dump

\item {} 
\sphinxAtStartPar
group\sphinxhyphen{}ID = ID of the group of atoms to be imaged

\item {} 
\sphinxAtStartPar
adios = style of dump command (other styles \sphinxstyleemphasis{atom} or \sphinxstyleemphasis{cfg} or \sphinxstyleemphasis{dcd} or \sphinxstyleemphasis{xtc} or \sphinxstyleemphasis{xyz} or \sphinxstyleemphasis{local} or \sphinxstyleemphasis{custom} are discussed on the {\hyperref[\detokenize{dump::doc}]{\sphinxcrossref{\DUrole{doc}{dump}}}} doc page)

\item {} 
\sphinxAtStartPar
N = dump every this many timesteps

\item {} 
\sphinxAtStartPar
file.bp = name of file/stream to write to

\item {} 
\sphinxAtStartPar
args = same options as in {\hyperref[\detokenize{dump::doc}]{\sphinxcrossref{\DUrole{doc}{dump custom}}}} command

\end{itemize}


\subsection{Examples}
\label{\detokenize{dump_adios:examples}}
\begin{sphinxVerbatim}[commandchars=\\\{\}]
\PYG{k}{dump }\PYG{n+nv+nvIdentifier}{adios1}\PYG{+w}{ }\PYG{n+nv+nvIdentifier}{all}\PYG{+w}{ }\PYG{n}{atom}\PYG{o}{/}\PYG{n}{adios}\PYG{+w}{   }\PYG{l+m}{100}\PYG{+w}{ }\PYG{n}{atoms.bp}
\PYG{k}{dump }\PYG{n+nv+nvIdentifier}{4a}\PYG{+w}{     }\PYG{n+nv+nvIdentifier}{all}\PYG{+w}{ }\PYG{n}{custom}\PYG{o}{/}\PYG{n}{adios}\PYG{+w}{ }\PYG{l+m}{100}\PYG{+w}{ }\PYG{n}{dump\PYGZus{}adios.bp}\PYG{+w}{ }\PYG{n}{id}\PYG{+w}{ }\PYG{n}{v\PYGZus{}p}\PYG{+w}{ }\PYG{n}{x}\PYG{+w}{ }\PYG{n}{y}\PYG{+w}{ }\PYG{n}{z}
\PYG{k}{dump }\PYG{n+nv+nvIdentifier}{2}\PYG{+w}{ }\PYG{n+nv+nvIdentifier}{subgroup}\PYG{+w}{ }\PYG{n}{custom}\PYG{o}{/}\PYG{n}{adios}\PYG{+w}{ }\PYG{l+m}{100}\PYG{+w}{ }\PYG{n}{dump\PYGZus{}adios.bp}\PYG{+w}{ }\PYG{n}{mass}\PYG{+w}{ }\PYG{n}{type}\PYG{+w}{ }\PYG{n}{xs}\PYG{+w}{ }\PYG{n}{ys}\PYG{+w}{ }\PYG{n}{zs}\PYG{+w}{ }\PYG{n}{vx}\PYG{+w}{ }\PYG{n}{vy}\PYG{+w}{ }\PYG{n}{vz}
\end{sphinxVerbatim}


\subsection{Description}
\label{\detokenize{dump_adios:description}}
\sphinxAtStartPar
Dump a snapshot of atom coordinates every \(N\) timesteps in the \sphinxhref{https://github.com/ornladios/ADIOS2}{ADIOS}\sphinxhyphen{}based “BP” file format, or using different I/O solutions in
ADIOS, to a stream that can be read on\sphinxhyphen{}line by another program.
ADIOS\sphinxhyphen{}BP files are binary, portable, and self\sphinxhyphen{}describing.

\begin{sphinxadmonition}{note}{Note:}
\sphinxAtStartPar
To be able to use ADIOS, a file \sphinxcode{\sphinxupquote{adios2\_config.xml}} with specific
configuration settings is expected in the current working directory.
If the file is not present, LAMMPS will try to create a minimal
default file.  Please refer to the ADIOS documentation for details on
how to adjust this file for optimal performance and desired features.
\end{sphinxadmonition}

\sphinxAtStartPar
\sphinxstylestrong{Use from write\_dump:}

\sphinxAtStartPar
It is possible to use these dump styles with the
{\hyperref[\detokenize{write_dump::doc}]{\sphinxcrossref{\DUrole{doc}{write\_dump}}}} command.  In this case, the sub\sphinxhyphen{}intervals
must not be set at all.  The write\_dump command can be used to
create a new file at each individual dump.

\begin{sphinxVerbatim}[commandchars=\\\{\}]
\PYG{k}{dump }\PYG{n+nv+nvIdentifier}{4}\PYG{+w}{     }\PYG{n+nv+nvIdentifier}{all}\PYG{+w}{ }\PYG{n}{atom}\PYG{o}{/}\PYG{n}{adios}\PYG{+w}{ }\PYG{l+m}{100}\PYG{+w}{ }\PYG{n}{dump.bp}
\PYG{k}{write\PYGZus{}dump }\PYG{n+nv+nvIdentifier}{all}\PYG{+w}{ }\PYG{n}{atom}\PYG{o}{/}\PYG{n}{adios}\PYG{+w}{ }\PYG{n}{singledump.bp}
\end{sphinxVerbatim}


\bigskip\hrule\bigskip



\subsection{Restrictions}
\label{\detokenize{dump_adios:restrictions}}
\sphinxAtStartPar
The number of atoms per snapshot \sphinxstylestrong{can} change with the adios style.
When using the ADIOS tool ‘bpls’ to list the content of a .bp file,
bpls will print \sphinxstyleemphasis{\_\_} for the size of the output table indicating that
its size is changing every step.

\sphinxAtStartPar
The \sphinxstyleemphasis{atom/adios} and \sphinxstyleemphasis{custom/adios} dump styles are part of the ADIOS
package.  They are only enabled if LAMMPS was built with that package.
See the {\hyperref[\detokenize{Build_package::doc}]{\sphinxcrossref{\DUrole{doc}{Build package}}}} page for more info.


\bigskip\hrule\bigskip



\subsection{Related commands}
\label{\detokenize{dump_adios:related-commands}}
\sphinxAtStartPar
{\hyperref[\detokenize{dump::doc}]{\sphinxcrossref{\DUrole{doc}{dump}}}}, {\hyperref[\detokenize{dump_modify::doc}]{\sphinxcrossref{\DUrole{doc}{dump\_modify}}}}, {\hyperref[\detokenize{undump::doc}]{\sphinxcrossref{\DUrole{doc}{undump}}}}

\sphinxstepscope

\index{dump cfg/uef@\spxentry{dump cfg/uef}}\ignorespaces 

\section{dump cfg/uef command}
\label{\detokenize{dump_cfg_uef:dump-cfg-uef-command}}\label{\detokenize{dump_cfg_uef:index-0}}\label{\detokenize{dump_cfg_uef::doc}}

\subsection{Syntax}
\label{\detokenize{dump_cfg_uef:syntax}}
\begin{sphinxVerbatim}[commandchars=\\\{\}]
\PYG{k}{dump }\PYG{n+nv+nvIdentifier}{ID}\PYG{+w}{ }\PYG{n+nv+nvIdentifier}{group\PYGZhy{}ID}\PYG{+w}{ }\PYG{n}{cfg}\PYG{o}{/}\PYG{n}{uef}\PYG{+w}{ }\PYG{n}{N}\PYG{+w}{ }\PYG{n}{file}\PYG{+w}{ }\PYG{n}{mass}\PYG{+w}{ }\PYG{n}{type}\PYG{+w}{ }\PYG{n}{xs}\PYG{+w}{ }\PYG{n}{ys}\PYG{+w}{ }\PYG{n}{zs}\PYG{+w}{ }\PYG{n}{args}
\end{sphinxVerbatim}
\begin{itemize}
\item {} 
\sphinxAtStartPar
ID = user\sphinxhyphen{}assigned name for the dump

\item {} 
\sphinxAtStartPar
group\sphinxhyphen{}ID = ID of the group of atoms to be dumped

\item {} 
\sphinxAtStartPar
N = dump every this many timesteps

\item {} 
\sphinxAtStartPar
file = name of file to write dump info to
\begin{sphinxalltt}
args = same as args for {\hyperref[\detokenize{dump::doc}]{\sphinxcrossref{\DUrole{doc}{dump custom}}}}
\end{sphinxalltt}

\end{itemize}


\subsection{Examples}
\label{\detokenize{dump_cfg_uef:examples}}
\begin{sphinxVerbatim}[commandchars=\\\{\}]
\PYG{k}{dump }\PYG{n+nv+nvIdentifier}{1}\PYG{+w}{ }\PYG{n+nv+nvIdentifier}{all}\PYG{+w}{ }\PYG{n}{cfg}\PYG{o}{/}\PYG{n}{uef}\PYG{+w}{ }\PYG{l+m}{10}\PYG{+w}{ }\PYG{n}{dump.}\PYG{o}{*}\PYG{n}{.cfg}\PYG{+w}{ }\PYG{n}{mass}\PYG{+w}{ }\PYG{n}{type}\PYG{+w}{ }\PYG{n}{xs}\PYG{+w}{ }\PYG{n}{ys}\PYG{+w}{ }\PYG{n}{zs}
\PYG{k}{dump }\PYG{n+nv+nvIdentifier}{2}\PYG{+w}{ }\PYG{n+nv+nvIdentifier}{all}\PYG{+w}{ }\PYG{n}{cfg}\PYG{o}{/}\PYG{n}{uef}\PYG{+w}{ }\PYG{l+m}{100}\PYG{+w}{ }\PYG{n}{dump.}\PYG{o}{*}\PYG{n}{.cfg}\PYG{+w}{ }\PYG{n}{mass}\PYG{+w}{ }\PYG{n}{type}\PYG{+w}{ }\PYG{n}{xs}\PYG{+w}{ }\PYG{n}{ys}\PYG{+w}{ }\PYG{n}{zs}\PYG{+w}{ }\PYG{n}{id}\PYG{+w}{ }\PYG{n}{c\PYGZus{}stress}
\end{sphinxVerbatim}


\subsection{Description}
\label{\detokenize{dump_cfg_uef:description}}
\sphinxAtStartPar
This command is used to dump atomic coordinates in the
reference frame of the applied flow field when
{\hyperref[\detokenize{fix_nh_uef::doc}]{\sphinxcrossref{\DUrole{doc}{fix nvt/uef}}}} or {\hyperref[\detokenize{fix_nh_uef::doc}]{\sphinxcrossref{\DUrole{doc}{fix npt/uef}}}} is used.
Only the atomic coordinates and frame\sphinxhyphen{}invariant scalar quantities
will be in the flow frame. If velocities are selected
as output, for example, they will not be in the same
reference frame as the atomic positions.


\subsection{Restrictions}
\label{\detokenize{dump_cfg_uef:restrictions}}
\sphinxAtStartPar
This fix is part of the UEF package. It is only enabled if LAMMPS
was built with that package. See the {\hyperref[\detokenize{Build_package::doc}]{\sphinxcrossref{\DUrole{doc}{Build package}}}}
page for more info.

\sphinxAtStartPar
This command can only be used when {\hyperref[\detokenize{fix_nh_uef::doc}]{\sphinxcrossref{\DUrole{doc}{fix nvt/uef}}}}
or {\hyperref[\detokenize{fix_nh_uef::doc}]{\sphinxcrossref{\DUrole{doc}{fix npt/uef}}}} is active.


\subsection{Related commands}
\label{\detokenize{dump_cfg_uef:related-commands}}
\sphinxAtStartPar
{\hyperref[\detokenize{dump::doc}]{\sphinxcrossref{\DUrole{doc}{dump}}}},
{\hyperref[\detokenize{fix_nh_uef::doc}]{\sphinxcrossref{\DUrole{doc}{fix nvt/uef}}}}


\subsection{Default}
\label{\detokenize{dump_cfg_uef:default}}
\sphinxAtStartPar
none

\sphinxstepscope

\index{dump h5md@\spxentry{dump h5md}}\ignorespaces 

\section{dump h5md command}
\label{\detokenize{dump_h5md:dump-h5md-command}}\label{\detokenize{dump_h5md:index-0}}\label{\detokenize{dump_h5md::doc}}

\subsection{Syntax}
\label{\detokenize{dump_h5md:syntax}}
\begin{sphinxVerbatim}[commandchars=\\\{\}]
\PYG{k}{dump }\PYG{n+nv+nvIdentifier}{ID}\PYG{+w}{ }\PYG{n+nv+nvIdentifier}{group\PYGZhy{}ID}\PYG{+w}{ }\PYG{n}{h5md}\PYG{+w}{ }\PYG{n}{N}\PYG{+w}{ }\PYG{n}{file.h5}\PYG{+w}{ }\PYG{n}{args}
\end{sphinxVerbatim}
\begin{itemize}
\item {} 
\sphinxAtStartPar
ID = user\sphinxhyphen{}assigned name for the dump

\item {} 
\sphinxAtStartPar
group\sphinxhyphen{}ID = ID of the group of atoms to be imaged

\item {} 
\sphinxAtStartPar
\sphinxstyleemphasis{h5md} = style of dump command (other styles \sphinxstyleemphasis{atom} or \sphinxstyleemphasis{cfg} or \sphinxstyleemphasis{dcd} or \sphinxstyleemphasis{xtc} or \sphinxstyleemphasis{xyz} or \sphinxstyleemphasis{local} or \sphinxstyleemphasis{custom} are discussed on the {\hyperref[\detokenize{dump::doc}]{\sphinxcrossref{\DUrole{doc}{dump}}}} doc page)

\item {} 
\sphinxAtStartPar
N = dump every this many timesteps

\item {} 
\sphinxAtStartPar
file.h5 = name of file to write to

\item {} 
\sphinxAtStartPar
args = \sphinxstyleemphasis{position} options or \sphinxstyleemphasis{image} or \sphinxstyleemphasis{velocity} options or \sphinxstyleemphasis{force} options or \sphinxstyleemphasis{species} options or \sphinxstyleemphasis{file\_from} ID or \sphinxstyleemphasis{box} value or \sphinxstyleemphasis{create\_group} value or \sphinxstyleemphasis{author} value = list of data elements to dump, with their dump “sub\sphinxhyphen{}intervals”
\begin{sphinxalltt}
\sphinxstyleemphasis{position} options
\sphinxstyleemphasis{image}
\sphinxstyleemphasis{velocity} options
\sphinxstyleemphasis{force} options
\sphinxstyleemphasis{species} options
\sphinxstyleemphasis{file\_from} ID = do not open a new file, re\sphinxhyphen{}use the already opened file from dump ID
\sphinxstyleemphasis{box} value = \sphinxstyleemphasis{yes} or \sphinxstyleemphasis{no}
\sphinxstyleemphasis{create\_group} value = \sphinxstyleemphasis{yes} or \sphinxstyleemphasis{no}
\sphinxstyleemphasis{author} value = quoted string
\end{sphinxalltt}

\end{itemize}

\sphinxAtStartPar
Note that at least one element must be specified and that \sphinxstyleemphasis{image} may only be
present if \sphinxstyleemphasis{position} is specified first.

\sphinxAtStartPar
For the elements \sphinxstyleemphasis{position}, \sphinxstyleemphasis{velocity}, \sphinxstyleemphasis{force} and \sphinxstyleemphasis{species}, a
sub\sphinxhyphen{}interval may be specified to write the data only every N\_element
iterations of the dump (i.e. every N*N\_element time steps). This is
specified by this option directly following the element declaration:
\begin{sphinxalltt}
options = \sphinxstyleemphasis{every} N\_element
\end{sphinxalltt}


\subsection{Examples}
\label{\detokenize{dump_h5md:examples}}
\begin{sphinxVerbatim}[commandchars=\\\{\}]
\PYG{k}{dump }\PYG{n+nv+nvIdentifier}{h5md1}\PYG{+w}{ }\PYG{n+nv+nvIdentifier}{all}\PYG{+w}{ }\PYG{n}{h5md}\PYG{+w}{ }\PYG{l+m}{100}\PYG{+w}{ }\PYG{n}{dump\PYGZus{}h5md.h5}\PYG{+w}{ }\PYG{n}{position}\PYG{+w}{ }\PYG{n}{image}
\PYG{k}{dump }\PYG{n+nv+nvIdentifier}{h5md1}\PYG{+w}{ }\PYG{n+nv+nvIdentifier}{all}\PYG{+w}{ }\PYG{n}{h5md}\PYG{+w}{ }\PYG{l+m}{100}\PYG{+w}{ }\PYG{n}{dump\PYGZus{}h5md.h5}\PYG{+w}{ }\PYG{n}{position}\PYG{+w}{ }\PYG{k}{velocity }\PYG{n+nv+nvIdentifier}{every}\PYG{+w}{ }\PYG{l+m}{10}
\PYG{k}{dump }\PYG{n+nv+nvIdentifier}{h5md1}\PYG{+w}{ }\PYG{n+nv+nvIdentifier}{all}\PYG{+w}{ }\PYG{n}{h5md}\PYG{+w}{ }\PYG{l+m}{100}\PYG{+w}{ }\PYG{n}{dump\PYGZus{}h5md.h5}\PYG{+w}{ }\PYG{k}{velocity }\PYG{n+nv+nvIdentifier}{author}\PYG{+w}{ }\PYG{l+s}{\PYGZdq{}}\PYG{l+s}{John Doe}\PYG{l+s}{\PYGZdq{}}
\end{sphinxVerbatim}


\subsection{Description}
\label{\detokenize{dump_h5md:description}}
\sphinxAtStartPar
Dump a snapshot of atom coordinates every N timesteps in the
\sphinxhref{https://www.hdfgroup.org/solutions/hdf5/}{HDF5} based \sphinxhref{https://nongnu.org/h5md/}{H5MD} file format {\hyperref[\detokenize{dump_h5md:h5md-cpc}]{\sphinxcrossref{\DUrole{std,std-ref}{(de Buyl)}}}}.
HDF5 files are binary, portable and self\sphinxhyphen{}describing.  This dump style
will write only one file, on the root node.

\sphinxAtStartPar
Several dumps may write to the same file, by using file\_from and
referring to a previously defined dump.  Several groups may also be
stored within the same file by defining several dumps.  A dump that
refers (via \sphinxstyleemphasis{file\_from}) to an already open dump ID and that concerns
another particle group must specify \sphinxstyleemphasis{create\_group yes}.

\sphinxAtStartPar
Each data element is written every N*N\_element steps. For \sphinxstyleemphasis{image}, no
sub\sphinxhyphen{}interval is needed as it must be present at the same interval as
\sphinxstyleemphasis{position}.  \sphinxstyleemphasis{image} must be given after \sphinxstyleemphasis{position} in any case.  The
box information (edges in each dimension) is stored at the same
interval than the \sphinxstyleemphasis{position} element, if present. Else it is stored
every N steps.

\begin{sphinxadmonition}{note}{Note:}
\sphinxAtStartPar
Because periodic boundary conditions are enforced only on
timesteps when neighbor lists are rebuilt, the coordinates of an atom
written to a dump file may be slightly outside the simulation box.
\end{sphinxadmonition}

\sphinxAtStartPar
\sphinxstylestrong{Use from write\_dump:}

\sphinxAtStartPar
It is possible to use this dump style with the
{\hyperref[\detokenize{write_dump::doc}]{\sphinxcrossref{\DUrole{doc}{write\_dump}}}} command.  In this case, the sub\sphinxhyphen{}intervals
must not be set at all.  The write\_dump command can be used either to
create a new file or to add current data to an existing dump file by
using the \sphinxstyleemphasis{file\_from} keyword.

\sphinxAtStartPar
Typically, the \sphinxstyleemphasis{species} data is fixed. The following two commands
store the position data every 100 timesteps, with the image data, and
store once the species data in the same file.

\begin{sphinxVerbatim}[commandchars=\\\{\}]
\PYG{k}{dump }\PYG{n+nv+nvIdentifier}{h5md1}\PYG{+w}{ }\PYG{n+nv+nvIdentifier}{all}\PYG{+w}{ }\PYG{n}{h5md}\PYG{+w}{ }\PYG{l+m}{100}\PYG{+w}{ }\PYG{n}{dump.h5}\PYG{+w}{ }\PYG{n}{position}\PYG{+w}{ }\PYG{n}{image}
\PYG{k}{write\PYGZus{}dump }\PYG{n+nv+nvIdentifier}{all}\PYG{+w}{ }\PYG{n}{h5md}\PYG{+w}{ }\PYG{n}{dump.h5}\PYG{+w}{ }\PYG{n}{file\PYGZus{}from}\PYG{+w}{ }\PYG{n}{h5md1}\PYG{+w}{ }\PYG{n}{species}
\end{sphinxVerbatim}


\bigskip\hrule\bigskip



\subsection{Restrictions}
\label{\detokenize{dump_h5md:restrictions}}
\sphinxAtStartPar
The number of atoms per snapshot cannot change with the h5md style.
The position data is stored wrapped (box boundaries not enforced, see
note above).  Only orthogonal domains are currently supported. This is
a limitation of the present dump h5md command and not of H5MD itself.

\sphinxAtStartPar
The \sphinxstyleemphasis{h5md} dump style is part of the H5MD package. It is only
enabled if LAMMPS was built with that package. See the {\hyperref[\detokenize{Build_package::doc}]{\sphinxcrossref{\DUrole{doc}{Build package}}}} page for more info. It also requires
(i) building the ch5md library provided with LAMMPS (See the {\hyperref[\detokenize{Build_package::doc}]{\sphinxcrossref{\DUrole{doc}{Build package}}}} page for more info.) and (ii) having
the \sphinxhref{https://www.hdfgroup.org/solutions/hdf5/}{HDF5} library installed (C bindings are sufficient) on
your system.  The library ch5md is compiled with the h5cc wrapper
provided by the HDF5 library.


\bigskip\hrule\bigskip



\subsection{Related commands}
\label{\detokenize{dump_h5md:related-commands}}
\sphinxAtStartPar
{\hyperref[\detokenize{dump::doc}]{\sphinxcrossref{\DUrole{doc}{dump}}}}, {\hyperref[\detokenize{dump_modify::doc}]{\sphinxcrossref{\DUrole{doc}{dump\_modify}}}}, {\hyperref[\detokenize{undump::doc}]{\sphinxcrossref{\DUrole{doc}{undump}}}}


\bigskip\hrule\bigskip

\phantomsection\label{\detokenize{dump_h5md:h5md-cpc}}
\sphinxAtStartPar
\sphinxstylestrong{(de Buyl)} de Buyl, Colberg and Hofling, H5MD: A structured,
efficient, and portable file format for molecular data,
Comp. Phys. Comm. 185(6), 1546\sphinxhyphen{}1553 (2014) \sphinxhyphen{}
\sphinxhref{https://arxiv.org/abs/1308.6382/}{{[}arXiv:1308.6382{]}}.

\sphinxstepscope

\index{dump image@\spxentry{dump image}}\ignorespaces 
\index{dump movie@\spxentry{dump movie}}\ignorespaces 

\section{dump image command}
\label{\detokenize{dump_image:dump-image-command}}\label{\detokenize{dump_image:index-1}}\label{\detokenize{dump_image:index-0}}\label{\detokenize{dump_image::doc}}

\section{dump movie command}
\label{\detokenize{dump_image:dump-movie-command}}
\sphinxAtStartPar
(see below for {\hyperref[\detokenize{dump_image:dump-modify-image}]{\sphinxcrossref{\DUrole{std,std-ref}{dump\_modify options}}}} specific to dump image/movie)


\subsection{Syntax}
\label{\detokenize{dump_image:syntax}}
\begin{sphinxVerbatim}[commandchars=\\\{\}]
\PYG{k}{dump }\PYG{n+nv+nvIdentifier}{ID}\PYG{+w}{ }\PYG{n+nv+nvIdentifier}{group\PYGZhy{}ID}\PYG{+w}{ }\PYG{n}{style}\PYG{+w}{ }\PYG{n}{N}\PYG{+w}{ }\PYG{n}{file}\PYG{+w}{ }\PYG{n}{color}\PYG{+w}{ }\PYG{n}{diameter}\PYG{+w}{ }\PYG{n}{keyword}\PYG{+w}{ }\PYG{n}{value}\PYG{+w}{ }\PYG{n}{...}
\end{sphinxVerbatim}
\begin{itemize}
\item {} 
\sphinxAtStartPar
ID = user\sphinxhyphen{}assigned name for the dump

\item {} 
\sphinxAtStartPar
group\sphinxhyphen{}ID = ID of the group of atoms to be imaged

\item {} 
\sphinxAtStartPar
style = \sphinxstyleemphasis{image} or \sphinxstyleemphasis{movie} = style of dump command (other styles such as \sphinxstyleemphasis{atom} or \sphinxstyleemphasis{cfg} or \sphinxstyleemphasis{dcd} or \sphinxstyleemphasis{xtc} or \sphinxstyleemphasis{xyz} or \sphinxstyleemphasis{local} or \sphinxstyleemphasis{custom} are discussed on the {\hyperref[\detokenize{dump::doc}]{\sphinxcrossref{\DUrole{doc}{dump}}}} doc page)

\item {} 
\sphinxAtStartPar
N = dump every this many timesteps

\item {} 
\sphinxAtStartPar
file = name of file to write image to

\item {} 
\sphinxAtStartPar
color = atom attribute that determines color of each atom

\item {} 
\sphinxAtStartPar
diameter = atom attribute that determines size of each atom

\item {} 
\sphinxAtStartPar
zero or more keyword/value pairs may be appended

\item {} 
\sphinxAtStartPar
keyword = \sphinxstyleemphasis{atom} or \sphinxstyleemphasis{adiam} or \sphinxstyleemphasis{bond} or \sphinxstyleemphasis{grid} or \sphinxstyleemphasis{line} or \sphinxstyleemphasis{tri} or \sphinxstyleemphasis{body} or \sphinxstyleemphasis{fix} or \sphinxstyleemphasis{size} or \sphinxstyleemphasis{view} or \sphinxstyleemphasis{center} or \sphinxstyleemphasis{up} or \sphinxstyleemphasis{zoom} or \sphinxstyleemphasis{box} or \sphinxstyleemphasis{axes} or \sphinxstyleemphasis{subbox} or \sphinxstyleemphasis{shiny} or \sphinxstyleemphasis{fsaa} or \sphinxstyleemphasis{ssao}
\begin{sphinxalltt}
\sphinxstyleemphasis{atom} = \sphinxstyleemphasis{yes} or \sphinxstyleemphasis{no} = do or do not draw atoms
\sphinxstyleemphasis{adiam} size = numeric value for atom diameter (distance units)
\sphinxstyleemphasis{bond} values = color width = color and width of bonds
  color = \sphinxstyleemphasis{atom} or \sphinxstyleemphasis{type} or \sphinxstyleemphasis{none}
  width = number or \sphinxstyleemphasis{atom} or \sphinxstyleemphasis{type} or \sphinxstyleemphasis{none}
    number = numeric value for bond width (distance units)
\sphinxstyleemphasis{grid} = per\sphinxhyphen{}grid value to use when coloring each grid cell
  per\sphinxhyphen{}grid value = c\_ID:gname:dname, c\_ID:gname:dname{[}I{]}, f\_ID:gname:dname, f\_ID:gname:dname{[}I{]}
    gname = name of grid defined by compute or fix
    dname = name of data field defined by compute or fix
    c\_ID = per\sphinxhyphen{}grid vector calculated by a compute with ID
    c\_ID{[}I{]} = Ith column of per\sphinxhyphen{}grid array calculated by a compute with ID
    f\_ID = per\sphinxhyphen{}grid vector calculated by a fix with ID
    f\_ID{[}I{]} = Ith column of per\sphinxhyphen{}grid array calculated by a fix with ID
\sphinxstyleemphasis{line} = color width
  color = \sphinxstyleemphasis{type}
  width = numeric value for line width (distance units)
\sphinxstyleemphasis{tri} = color tflag width
  color = \sphinxstyleemphasis{type}
  tflag = 1 for just triangle, 2 for just tri edges, 3 for both
  width = numeric value for tringle edge width (distance units)
\sphinxstyleemphasis{body} = color bflag1 bflag2
  color = \sphinxstyleemphasis{type}
  bflag1,bflag2 = 2 numeric flags to affect how bodies are drawn
\sphinxstyleemphasis{fix} = fixID color fflag1 fflag2
  fixID = ID of fix that generates objects to draw
  color = \sphinxstyleemphasis{type}
  fflag1,fflag2 = 2 numeric flags to affect how fix objects are drawn
\sphinxstyleemphasis{size} values = width height = size of images
  width = width of image in \# of pixels
  height = height of image in \# of pixels
\sphinxstyleemphasis{view} values = theta phi = view of simulation box
  theta = view angle from +z axis (degrees)
  phi = azimuthal view angle (degrees)
  theta or phi can be a variable (see below)
\sphinxstyleemphasis{center} values = flag Cx Cy Cz = center point of image
  flag = \sphinxstyleemphasis{s} for static, \sphinxstyleemphasis{d} for dynamic
  Cx,Cy,Cz = center point of image as fraction of box dimension (0.5 = center of box)
  Cx,Cy,Cz can be variables (see below)
\sphinxstyleemphasis{up} values = Ux Uy Uz = direction that is "up" in image
  Ux,Uy,Uz = components of up vector
  Ux,Uy,Uz can be variables (see below)
\sphinxstyleemphasis{zoom} value = zfactor = size that simulation box appears in image
  zfactor = scale image size by factor \textgreater{} 1 to enlarge, factor \textless{} 1 to shrink
  zfactor can be a variable (see below)
\sphinxstyleemphasis{box} values = yes/no diam = draw outline of simulation box
  yes/no = do or do not draw simulation box lines
  diam = diameter of box lines as fraction of shortest box length
\sphinxstyleemphasis{axes} values = axes length diam = draw xyz axes
  axes = \sphinxstyleemphasis{yes} or \sphinxstyleemphasis{no} = do or do not draw xyz axes lines next to simulation box
  length = length of axes lines as fraction of respective box lengths
  diam = diameter of axes lines as fraction of shortest box length
\sphinxstyleemphasis{subbox} values = lines diam = draw outline of processor subdomains
  lines = \sphinxstyleemphasis{yes} or \sphinxstyleemphasis{no} = do or do not draw subdomain lines
  diam = diameter of subdomain lines as fraction of shortest box length
\sphinxstyleemphasis{shiny} value = sfactor = shinyness of spheres and cylinders
  sfactor = shinyness of spheres and cylinders from 0.0 to 1.0
\sphinxstyleemphasis{fsaa} arg = yes/no
  yes/no = do or do not apply anti\sphinxhyphen{}aliasing
\sphinxstyleemphasis{ssao} value = shading seed dfactor = SSAO depth shading
  shading = \sphinxstyleemphasis{yes} or \sphinxstyleemphasis{no} = turn depth shading on/off
  seed = random \# seed (positive integer)
  dfactor = strength of shading from 0.0 to 1.0
\end{sphinxalltt}

\end{itemize}


\section{dump\_modify options for dump image/movie}
\label{\detokenize{dump_image:dump-modify-options-for-dump-image-movie}}\label{\detokenize{dump_image:dump-modify-image}}

\subsection{Syntax}
\label{\detokenize{dump_image:id1}}
\begin{sphinxVerbatim}[commandchars=\\\{\}]
\PYG{k}{dump\PYGZus{}modify }\PYG{n+nv+nvIdentifier}{dump\PYGZhy{}ID}\PYG{+w}{ }\PYG{n}{keyword}\PYG{+w}{ }\PYG{n}{values}\PYG{+w}{ }\PYG{n}{...}
\end{sphinxVerbatim}
\begin{itemize}
\item {} 
\sphinxAtStartPar
these keywords apply only to the \sphinxstyleemphasis{image} and \sphinxstyleemphasis{movie} styles and are documented on this page

\item {} 
\sphinxAtStartPar
keyword = \sphinxstyleemphasis{acolor} or \sphinxstyleemphasis{adiam} or \sphinxstyleemphasis{amap} or \sphinxstyleemphasis{gmap} or \sphinxstyleemphasis{backcolor} or \sphinxstyleemphasis{bcolor} or \sphinxstyleemphasis{bdiam} or \sphinxstyleemphasis{bitrate} or \sphinxstyleemphasis{boxcolor} or \sphinxstyleemphasis{color} or \sphinxstyleemphasis{framerate} or \sphinxstyleemphasis{gmap}

\item {} 
\sphinxAtStartPar
see the {\hyperref[\detokenize{dump_modify::doc}]{\sphinxcrossref{\DUrole{doc}{dump modify}}}} doc page for more general keywords
\begin{sphinxalltt}
\sphinxstyleemphasis{acolor} args = type color
  type = atom type (numeric or type label) or range of numeric types (see below)
  color = name of color or color1/color2/...
\sphinxstyleemphasis{adiam} args = type diam
  type = atom type (numeric or type label) or range of numeric types (see below)
  diam = diameter of atoms of that type (distance units)
\sphinxstyleemphasis{amap} args = lo hi style delta N entry1 entry2 ... entryN
  lo = number or \sphinxstyleemphasis{min} = lower bound of range of color map
  hi = number or \sphinxstyleemphasis{max} = upper bound of range of color map
  style = 2 letters = \sphinxstyleemphasis{c} or \sphinxstyleemphasis{d} or \sphinxstyleemphasis{s} plus \sphinxstyleemphasis{a} or \sphinxstyleemphasis{f}
    \sphinxstyleemphasis{c} for continuous
    \sphinxstyleemphasis{d} for discrete
    \sphinxstyleemphasis{s} for sequential
    \sphinxstyleemphasis{a} for absolute
    \sphinxstyleemphasis{f} for fractional
  delta = binsize (only used for style \sphinxstyleemphasis{s}, otherwise ignored)
    binsize = range is divided into bins of this width
  N = \# of subsequent entries
  entry = value color (for continuous style)
    value = number or \sphinxstyleemphasis{min} or \sphinxstyleemphasis{max} = single value within range
    color = name of color used for that value
  entry = lo hi color (for discrete style)
    lo/hi = number or \sphinxstyleemphasis{min} or \sphinxstyleemphasis{max} = lower/upper bound of subset of range
    color = name of color used for that subset of values
  entry = color (for sequential style)
    color = name of color used for a bin of values
\sphinxstyleemphasis{backcolor} arg = color
  color = name of color for background
\sphinxstyleemphasis{bcolor} args = type color
  type = bond type (numeric or type label) or range of numeric types (see below)
  color = name of color or color1/color2/...
\sphinxstyleemphasis{bdiam} args = type diam
  type = bond type (numeric or type label) or range of numeric types (see below)
  diam = diameter of bonds of that type (distance units)
\sphinxstyleemphasis{bitrate} arg = rate
  rate = target bitrate for movie in kbps
\sphinxstyleemphasis{boxcolor} arg = color
  color = name of color for simulation box lines and processor subdomain lines
\sphinxstyleemphasis{color} args = name R G B
  name = name of color
  R,G,B = red/green/blue numeric values from 0.0 to 1.0
\sphinxstyleemphasis{framerate} arg = fps
  fps = frames per second for movie
\sphinxstyleemphasis{gmap} args = identical to \sphinxstyleemphasis{amap} args
\end{sphinxalltt}

\end{itemize}


\subsection{Examples}
\label{\detokenize{dump_image:examples}}
\begin{sphinxVerbatim}[commandchars=\\\{\}]
\PYG{k}{dump }\PYG{n+nv+nvIdentifier}{d0}\PYG{+w}{ }\PYG{n+nv+nvIdentifier}{all}\PYG{+w}{ }\PYG{n}{image}\PYG{+w}{ }\PYG{l+m}{100}\PYG{+w}{ }\PYG{n}{dump.}\PYG{o}{*}\PYG{n}{.jpg}\PYG{+w}{ }\PYG{n}{type}\PYG{+w}{ }\PYG{n}{type}
\PYG{k}{dump }\PYG{n+nv+nvIdentifier}{d1}\PYG{+w}{ }\PYG{n+nv+nvIdentifier}{mobile}\PYG{+w}{ }\PYG{n}{image}\PYG{+w}{ }\PYG{l+m}{500}\PYG{+w}{ }\PYG{n}{snap.}\PYG{o}{*}\PYG{n}{.png}\PYG{+w}{ }\PYG{n}{element}\PYG{+w}{ }\PYG{n}{element}\PYG{+w}{ }\PYG{n}{ssao}\PYG{+w}{ }\PYG{n}{yes}\PYG{+w}{ }\PYG{l+m}{4539}\PYG{+w}{ }\PYG{l+m}{0.6}
\PYG{k}{dump }\PYG{n+nv+nvIdentifier}{d2}\PYG{+w}{ }\PYG{n+nv+nvIdentifier}{all}\PYG{+w}{ }\PYG{n}{image}\PYG{+w}{ }\PYG{l+m}{200}\PYG{+w}{ }\PYG{n}{img}\PYG{o}{\PYGZhy{}}\PYG{o}{*}\PYG{n}{.ppm}\PYG{+w}{ }\PYG{n}{type}\PYG{+w}{ }\PYG{n}{type}\PYG{+w}{ }\PYG{n}{zoom}\PYG{+w}{ }\PYG{l+m}{2.5}\PYG{+w}{ }\PYG{n}{adiam}\PYG{+w}{ }\PYG{l+m}{1.5}\PYG{+w}{ }\PYG{n}{size}\PYG{+w}{ }\PYG{l+m}{1280}\PYG{+w}{ }\PYG{l+m}{720}
\PYG{k}{dump }\PYG{n+nv+nvIdentifier}{m0}\PYG{+w}{ }\PYG{n+nv+nvIdentifier}{all}\PYG{+w}{ }\PYG{n}{movie}\PYG{+w}{ }\PYG{l+m}{1000}\PYG{+w}{ }\PYG{n}{movie.mpg}\PYG{+w}{ }\PYG{n}{type}\PYG{+w}{ }\PYG{n}{type}\PYG{+w}{ }\PYG{n}{size}\PYG{+w}{ }\PYG{l+m}{640}\PYG{+w}{ }\PYG{l+m}{480}
\PYG{k}{dump }\PYG{n+nv+nvIdentifier}{m1}\PYG{+w}{ }\PYG{n+nv+nvIdentifier}{all}\PYG{+w}{ }\PYG{n}{movie}\PYG{+w}{ }\PYG{l+m}{1000}\PYG{+w}{ }\PYG{n}{movie.avi}\PYG{+w}{ }\PYG{n}{type}\PYG{+w}{ }\PYG{n}{type}\PYG{+w}{ }\PYG{n}{size}\PYG{+w}{ }\PYG{l+m}{640}\PYG{+w}{ }\PYG{l+m}{480}
\PYG{k}{dump }\PYG{n+nv+nvIdentifier}{m2}\PYG{+w}{ }\PYG{n+nv+nvIdentifier}{all}\PYG{+w}{ }\PYG{n}{movie}\PYG{+w}{ }\PYG{l+m}{100}\PYG{+w}{ }\PYG{n}{movie.m4v}\PYG{+w}{ }\PYG{n}{type}\PYG{+w}{ }\PYG{n}{type}\PYG{+w}{ }\PYG{n}{zoom}\PYG{+w}{ }\PYG{l+m}{1.8}\PYG{+w}{ }\PYG{n}{adiam}\PYG{+w}{ }\PYG{n}{v\PYGZus{}value}\PYG{+w}{ }\PYG{n}{size}\PYG{+w}{ }\PYG{l+m}{1280}\PYG{+w}{ }\PYG{l+m}{720}

\PYG{k}{dump\PYGZus{}modify }\PYG{n+nv+nvIdentifier}{1}\PYG{+w}{ }\PYG{n}{amap}\PYG{+w}{ }\PYG{n}{min}\PYG{+w}{ }\PYG{n}{max}\PYG{+w}{ }\PYG{n}{cf}\PYG{+w}{ }\PYG{l+m}{0.0}\PYG{+w}{ }\PYG{l+m}{3}\PYG{+w}{ }\PYG{n}{min}\PYG{+w}{ }\PYG{n}{green}\PYG{+w}{ }\PYG{l+m}{0.5}\PYG{+w}{ }\PYG{n}{yellow}\PYG{+w}{ }\PYG{n}{max}\PYG{+w}{ }\PYG{n}{blue}\PYG{+w}{ }\PYG{n}{boxcolor}\PYG{+w}{ }\PYG{n}{red}

\PYG{k}{labelmap}\PYG{+w}{ }\PYG{n}{atom}\PYG{+w}{ }\PYG{l+m}{1}\PYG{+w}{ }\PYG{n}{C}\PYG{+w}{ }\PYG{l+m}{2}\PYG{+w}{ }\PYG{n}{H}\PYG{+w}{ }\PYG{l+m}{3}\PYG{+w}{ }\PYG{n}{O}\PYG{+w}{ }\PYG{l+m}{4}\PYG{+w}{ }\PYG{n}{N}
\PYG{k}{dump\PYGZus{}modify }\PYG{n+nv+nvIdentifier}{1}\PYG{+w}{ }\PYG{n}{acolor}\PYG{+w}{ }\PYG{n}{C}\PYG{+w}{ }\PYG{n}{gray}\PYG{+w}{ }\PYG{n}{acolor}\PYG{+w}{ }\PYG{n}{H}\PYG{+w}{ }\PYG{n}{white}\PYG{+w}{ }\PYG{n}{acolor}\PYG{+w}{ }\PYG{n}{O}\PYG{+w}{ }\PYG{n}{red}\PYG{+w}{ }\PYG{n}{acolor}\PYG{+w}{ }\PYG{n}{N}\PYG{+w}{ }\PYG{n}{blue}
\end{sphinxVerbatim}


\subsection{Description}
\label{\detokenize{dump_image:description}}
\sphinxAtStartPar
Dump a high\sphinxhyphen{}quality rendered image of the atom configuration every \(N\)
timesteps and save the images either as a sequence of JPEG or PNG or
PPM files, or as a single movie file.  The options for this command as
well as the {\hyperref[\detokenize{dump_modify::doc}]{\sphinxcrossref{\DUrole{doc}{dump\_modify}}}} command control what is
included in the image or movie and how it appears.  A series of such
images can easily be manually converted into an animated movie of your
simulation or the process can be automated without writing the
intermediate files using the dump movie style; see further details
below.  Other dump styles store snapshots of numerical data associated
with atoms in various formats, as discussed on the {\hyperref[\detokenize{dump::doc}]{\sphinxcrossref{\DUrole{doc}{dump}}}}
doc page.

\sphinxAtStartPar
Note that a set of images or a movie can be made after a simulation
has been run, using the {\hyperref[\detokenize{rerun::doc}]{\sphinxcrossref{\DUrole{doc}{rerun}}}} command to read snapshots
from an existing dump file, and using these dump commands in the rerun
script to generate the images/movie.

\sphinxAtStartPar
Here are two sample images, rendered as \(1024\times 1024\) JPEG files.

\sphinxAtStartPar
\sphinxincludegraphics[width=0.480\linewidth]{{dump1}.jpg}  \sphinxincludegraphics[width=0.480\linewidth]{{dump2}.jpg}



\sphinxAtStartPar
Only atoms in the specified group are rendered in the image.  The
{\hyperref[\detokenize{dump_modify::doc}]{\sphinxcrossref{\DUrole{doc}{dump\_modify region and thresh}}}} commands can also
alter what atoms are included in the image.
The filename suffix determines whether a JPEG, PNG, or PPM file is
created with the \sphinxstyleemphasis{image} dump style.  If the suffix is “.jpg” or
“.jpeg”, then a \sphinxhref{https://jpeg.org/jpeg/}{JPEG format} file is created, if the
suffix is “.png”, then a \sphinxhref{https://en.wikipedia.org/wiki/Portable\_Network\_Graphics}{PNG format} is created, else
a \sphinxhref{https://en.wikipedia.org/wiki/Netpbm}{PPM (aka NETPBM) format} file is created.
The JPEG and PNG files are binary; PPM has a text mode header followed
by binary data. JPEG images have lossy compression, PNG has lossless
compression, and PPM files are uncompressed but can be compressed with
gzip, if LAMMPS has been compiled with \sphinxhyphen{}DLAMMPS\_GZIP and a “.gz” suffix
is used.

\sphinxAtStartPar
Similarly, the format of the resulting movie is chosen with the \sphinxstyleemphasis{movie}
dump style. This is handled by the underlying FFmpeg converter and thus
details have to be looked up in the \sphinxhref{https://ffmpeg.org/}{FFmpeg documentation}.  Typical examples are: .avi, .mpg, .m4v, .mp4,
.mkv, .flv, .mov, .gif Additional settings of the movie compression like
\sphinxstyleemphasis{bitrate} and \sphinxstyleemphasis{framerate} can be set using the dump\_modify command as
described below.

\sphinxAtStartPar
To write out JPEG and PNG format files, you must build LAMMPS with
support for the corresponding JPEG or PNG library.  To convert images
into movies, LAMMPS has to be compiled with the \sphinxhyphen{}DLAMMPS\_FFMPEG
flag. See the {\hyperref[\detokenize{Build_settings::doc}]{\sphinxcrossref{\DUrole{doc}{Build settings}}}} page for
details.

\begin{sphinxadmonition}{note}{Note:}
\sphinxAtStartPar
Because periodic boundary conditions are enforced only on
timesteps when neighbor lists are rebuilt, the coordinates of an atom
in the image may be slightly outside the simulation box.
\end{sphinxadmonition}


\bigskip\hrule\bigskip


\sphinxAtStartPar
Dumps are performed on timesteps that are a multiple of \(N\) (including
timestep 0) and on the last timestep of a minimization if the
minimization converges.  Note that this means a dump will not be
performed on the initial timestep after the dump command is invoked,
if the current timestep is not a multiple of \(N\).  This behavior can be
changed via the {\hyperref[\detokenize{dump_modify::doc}]{\sphinxcrossref{\DUrole{doc}{dump\_modify first}}}} command, which
can be useful if the dump command is invoked after a minimization
ended on an arbitrary timestep. \(N\) can be changed between runs by
using the {\hyperref[\detokenize{dump_modify::doc}]{\sphinxcrossref{\DUrole{doc}{dump\_modify every}}}} command.

\sphinxAtStartPar
Dump \sphinxstyleemphasis{image} filenames must contain a wildcard character “*” so that
one image file per snapshot is written.  The “*” character is replaced
with the timestep value.  For example, tmp.dump.*.jpg becomes
tmp.dump.0.jpg, tmp.dump.10000.jpg, tmp.dump.20000.jpg, etc.  Note
that the {\hyperref[\detokenize{dump_modify::doc}]{\sphinxcrossref{\DUrole{doc}{dump\_modify pad}}}} command can be used to
ensure all timestep numbers are the same length (e.g., 00010), which
can make it easier to convert a series of images into a movie in the
correct ordering.

\sphinxAtStartPar
Dump \sphinxstyleemphasis{movie} filenames on the other hand, must not have any wildcard
character since only one file combining all images into a single
movie will be written by the movie encoder.


\bigskip\hrule\bigskip


\sphinxAtStartPar
The \sphinxstyleemphasis{color} and \sphinxstyleemphasis{diameter} settings determine the color and size of
atoms rendered in the image.  They can be any atom attribute defined
for the {\hyperref[\detokenize{dump::doc}]{\sphinxcrossref{\DUrole{doc}{dump custom}}}} command, including \sphinxstyleemphasis{type} and
\sphinxstyleemphasis{element}.  This includes per\sphinxhyphen{}atom quantities calculated by a
{\hyperref[\detokenize{compute::doc}]{\sphinxcrossref{\DUrole{doc}{compute}}}}, {\hyperref[\detokenize{fix::doc}]{\sphinxcrossref{\DUrole{doc}{fix}}}}, or {\hyperref[\detokenize{variable::doc}]{\sphinxcrossref{\DUrole{doc}{variable}}}},
which are prefixed by “c\_”, “f\_”, or “v\_”, respectively.  Note that the
\sphinxstyleemphasis{diameter} setting can be overridden with a numeric value applied to
all atoms by the optional \sphinxstyleemphasis{adiam} keyword.

\sphinxAtStartPar
If \sphinxstyleemphasis{type} is specified for the \sphinxstyleemphasis{color} setting, then the color of each
atom is determined by its atom type.  By default the mapping of types
to colors is as follows:
\begin{itemize}
\item {} 
\sphinxAtStartPar
type 1 = red

\item {} 
\sphinxAtStartPar
type 2 = green

\item {} 
\sphinxAtStartPar
type 3 = blue

\item {} 
\sphinxAtStartPar
type 4 = yellow

\item {} 
\sphinxAtStartPar
type 5 = aqua

\item {} 
\sphinxAtStartPar
type 6 = cyan

\end{itemize}

\sphinxAtStartPar
and repeats itself for types \(> 6\).  This mapping can be changed by the
“dump\_modify acolor” command, as described below.

\sphinxAtStartPar
If \sphinxstyleemphasis{type} is specified for the \sphinxstyleemphasis{diameter} setting then the diameter of
each atom is determined by its atom type.  By default all types have
diameter 1.0.  This mapping can be changed by the “dump\_modify adiam”
command, as described below.

\sphinxAtStartPar
If \sphinxstyleemphasis{element} is specified for the \sphinxstyleemphasis{color} and/or \sphinxstyleemphasis{diameter} setting,
then the color and/or diameter of each atom is determined by which
element it is, which in turn is specified by the element\sphinxhyphen{}to\sphinxhyphen{}type
mapping specified by the “dump\_modify element” command, as described
below.  By default every atom type is C (carbon).  Every element has a
color and diameter associated with it, which is the same as the colors
and sizes used by the \sphinxhref{http://li.mit.edu/Archive/Graphics/A/}{AtomEye} visualization package.

\sphinxAtStartPar
If other atom attributes are used for the \sphinxstyleemphasis{color} or \sphinxstyleemphasis{diameter}
settings, they are interpreted in the following way.

\sphinxAtStartPar
If “vx”, for example, is used as the \sphinxstyleemphasis{color} setting, then the color
of the atom will depend on the x\sphinxhyphen{}component of its velocity.  The
association of a per\sphinxhyphen{}atom value with a specific color is determined by
a “color map”, which can be specified via the dump\_modify amap
command, as described below.  The basic idea is that the
atom\sphinxhyphen{}attribute will be within a range of values, and every value
within the range is mapped to a specific color.  Depending on how the
color map is defined, that mapping can take place via interpolation so
that a value of \sphinxhyphen{}3.2 is halfway between “red” and “blue”, or
discretely so that the value of \sphinxhyphen{}3.2 is “orange”.

\sphinxAtStartPar
If “vx”, for example, is used as the \sphinxstyleemphasis{diameter} setting, then the atom
will be rendered using the x\sphinxhyphen{}component of its velocity as the
diameter.  If the per\sphinxhyphen{}atom value \textless{}= 0.0, them the atom will not be
drawn.  Note that finite\sphinxhyphen{}size spherical particles, as defined by
{\hyperref[\detokenize{atom_style::doc}]{\sphinxcrossref{\DUrole{doc}{atom\_style sphere}}}} define a per\sphinxhyphen{}particle radius or
diameter, which can be used as the \sphinxstyleemphasis{diameter} setting.


\bigskip\hrule\bigskip


\sphinxAtStartPar
The various keywords listed above control how the image is rendered.
As listed below, all of the keywords have defaults, most of which you
will likely not need to change.  As described below, the dump modify
command also has options specific to the dump image style,
particularly for assigning colors to atoms, bonds, and other image
features.


\bigskip\hrule\bigskip


\sphinxAtStartPar
The \sphinxstyleemphasis{atom} keyword allow you to turn off the drawing of all atoms, if
the specified value is \sphinxstyleemphasis{no}.  Note that this will not turn off the
drawing of particles that are represented as lines, triangles, or
bodies, as discussed below.  These particles can be drawn separately
if the \sphinxstyleemphasis{line}, \sphinxstyleemphasis{tri}, or \sphinxstyleemphasis{body} keywords are used.

\sphinxAtStartPar
The \sphinxstyleemphasis{adiam} keyword allows you to override the \sphinxstyleemphasis{diameter} setting to
set a single numeric \sphinxstyleemphasis{size}.  All atoms will be drawn with that
diameter, e.g. 1.5, which is in whatever distance {\hyperref[\detokenize{units::doc}]{\sphinxcrossref{\DUrole{doc}{units}}}}
the input script defines, e.g. Angstroms.


\bigskip\hrule\bigskip


\sphinxAtStartPar
The \sphinxstyleemphasis{bond} keyword allows to you to alter how bonds are drawn.  A bond
is only drawn if both atoms in the bond are being drawn due to being
in the specified group and due to other selection criteria
(e.g. region, threshold settings of the
{\hyperref[\detokenize{dump_modify::doc}]{\sphinxcrossref{\DUrole{doc}{dump\_modify}}}} command).  By default, bonds are drawn
if they are defined in the input data file as read by the
{\hyperref[\detokenize{read_data::doc}]{\sphinxcrossref{\DUrole{doc}{read\_data}}}} command.  Using \sphinxstyleemphasis{none} for both the bond
\sphinxstyleemphasis{color} and \sphinxstyleemphasis{width} value will turn off the drawing of all bonds.

\sphinxAtStartPar
If \sphinxstyleemphasis{atom} is specified for the bond \sphinxstyleemphasis{color} value, then each bond is
drawn in 2 halves, with the color of each half being the color of the
atom at that end of the bond.

\sphinxAtStartPar
If \sphinxstyleemphasis{type} is specified for the \sphinxstyleemphasis{color} value, then the color of each
bond is determined by its bond type.  By default the mapping of bond
types to colors is as follows:
\begin{itemize}
\item {} 
\sphinxAtStartPar
type 1 = red

\item {} 
\sphinxAtStartPar
type 2 = green

\item {} 
\sphinxAtStartPar
type 3 = blue

\item {} 
\sphinxAtStartPar
type 4 = yellow

\item {} 
\sphinxAtStartPar
type 5 = aqua

\item {} 
\sphinxAtStartPar
type 6 = cyan

\end{itemize}

\sphinxAtStartPar
and repeats itself for bond types \textgreater{} 6.  This mapping can be changed by
the “dump\_modify bcolor” command, as described below.

\sphinxAtStartPar
The bond \sphinxstyleemphasis{width} value can be a numeric value or \sphinxstyleemphasis{atom} or \sphinxstyleemphasis{type} (or
\sphinxstyleemphasis{none} as indicated above).

\sphinxAtStartPar
If a numeric value is specified, then all bonds will be drawn as
cylinders with that diameter, e.g. 1.0, which is in whatever distance
{\hyperref[\detokenize{units::doc}]{\sphinxcrossref{\DUrole{doc}{units}}}} the input script defines, e.g. Angstroms.

\sphinxAtStartPar
If \sphinxstyleemphasis{atom} is specified for the \sphinxstyleemphasis{width} value, then each bond
will be drawn with a width corresponding to the minimum diameter
of the two atoms in the bond.

\sphinxAtStartPar
If \sphinxstyleemphasis{type} is specified for the \sphinxstyleemphasis{width} value then the diameter of each
bond is determined by its bond type.  By default all types have
diameter 0.5.  This mapping can be changed by the “dump\_modify bdiam” command,
as described below.


\bigskip\hrule\bigskip


\sphinxAtStartPar
The \sphinxstyleemphasis{line} keyword can be used when {\hyperref[\detokenize{atom_style::doc}]{\sphinxcrossref{\DUrole{doc}{atom\_style line}}}}
is used to define particles as line segments, and will draw them as
lines.  If this keyword is not used, such particles will be drawn as
spheres, the same as if they were regular atoms.  The only setting
currently allowed for the \sphinxstyleemphasis{color} value is \sphinxstyleemphasis{type}, which will color
the lines according to the atom type of the particle.  By default the
mapping of types to colors is as follows:
\begin{itemize}
\item {} 
\sphinxAtStartPar
type 1 = red

\item {} 
\sphinxAtStartPar
type 2 = green

\item {} 
\sphinxAtStartPar
type 3 = blue

\item {} 
\sphinxAtStartPar
type 4 = yellow

\item {} 
\sphinxAtStartPar
type 5 = aqua

\item {} 
\sphinxAtStartPar
type 6 = cyan

\end{itemize}

\sphinxAtStartPar
and repeats itself for types \textgreater{} 6.  There is not yet an option to
change this via the dump\_modify command.

\sphinxAtStartPar
The line \sphinxstyleemphasis{width} can only be a numeric value, which specifies that all
lines will be drawn as cylinders with that diameter, e.g. 1.0, which
is in whatever distance {\hyperref[\detokenize{units::doc}]{\sphinxcrossref{\DUrole{doc}{units}}}} the input script defines,
e.g. Angstroms.


\bigskip\hrule\bigskip


\sphinxAtStartPar
The \sphinxstyleemphasis{tri} keyword can be used when {\hyperref[\detokenize{atom_style::doc}]{\sphinxcrossref{\DUrole{doc}{atom\_style tri}}}} is
used to define particles as triangles, and will draw them as triangles
or edges (3 lines) or both, depending on the setting for \sphinxstyleemphasis{tflag}.  If
edges are drawn, the \sphinxstyleemphasis{width} setting determines the diameters of the
line segments.  If this keyword is not used, triangle particles will
be drawn as spheres, the same as if they were regular atoms.  The only
setting currently allowed for the \sphinxstyleemphasis{color} value is \sphinxstyleemphasis{type}, which will
color the triangles according to the atom type of the particle.  By
default the mapping of types to colors is as follows:
\begin{itemize}
\item {} 
\sphinxAtStartPar
type 1 = red

\item {} 
\sphinxAtStartPar
type 2 = green

\item {} 
\sphinxAtStartPar
type 3 = blue

\item {} 
\sphinxAtStartPar
type 4 = yellow

\item {} 
\sphinxAtStartPar
type 5 = aqua

\item {} 
\sphinxAtStartPar
type 6 = cyan

\end{itemize}

\sphinxAtStartPar
and repeats itself for types \textgreater{} 6.  There is not yet an option to
change this via the dump\_modify command.


\bigskip\hrule\bigskip


\sphinxAtStartPar
The \sphinxstyleemphasis{body} keyword can be used when {\hyperref[\detokenize{atom_style::doc}]{\sphinxcrossref{\DUrole{doc}{atom\_style body}}}}
is used to define body particles with internal state
(e.g. sub\sphinxhyphen{}particles), and will drawn them in a manner specific to the
body style.  If this keyword is not used, such particles will be drawn
as spheres, the same as if they were regular atoms.

\sphinxAtStartPar
The {\hyperref[\detokenize{Howto_body::doc}]{\sphinxcrossref{\DUrole{doc}{Howto body}}}} page describes the body styles
LAMMPS currently supports, and provides more details as to the kind of
body particles they represent and how they are drawn by this dump
image command.  For all the body styles, individual atoms can be
either a body particle or a usual point (non\sphinxhyphen{}body) particle.  Non\sphinxhyphen{}body
particles will be drawn the same way they would be as a regular atom.
The \sphinxstyleemphasis{bflag1} and \sphinxstyleemphasis{bflag2} settings are numerical values which are
passed to the body style to affect how the drawing of a body particle
is done.  See the {\hyperref[\detokenize{Howto_body::doc}]{\sphinxcrossref{\DUrole{doc}{Howto body}}}} page for a
description of what these parameters mean for each body style.

\sphinxAtStartPar
The only setting currently allowed for the \sphinxstyleemphasis{color} value is \sphinxstyleemphasis{type},
which will color the body particles according to the atom type of the
particle.  By default the mapping of types to colors is as follows:
\begin{itemize}
\item {} 
\sphinxAtStartPar
type 1 = red

\item {} 
\sphinxAtStartPar
type 2 = green

\item {} 
\sphinxAtStartPar
type 3 = blue

\item {} 
\sphinxAtStartPar
type 4 = yellow

\item {} 
\sphinxAtStartPar
type 5 = aqua

\item {} 
\sphinxAtStartPar
type 6 = cyan

\end{itemize}

\sphinxAtStartPar
and repeats itself for types \textgreater{} 6.  There is not yet an option to
change this via the dump\_modify command.


\bigskip\hrule\bigskip


\sphinxAtStartPar
The \sphinxstyleemphasis{fix} keyword can be used with a {\hyperref[\detokenize{fix::doc}]{\sphinxcrossref{\DUrole{doc}{fix}}}} that produces
objects to be drawn.

\sphinxAtStartPar
The \sphinxstyleemphasis{fflag1} and \sphinxstyleemphasis{fflag2} settings are numerical values which are
passed to the fix to affect how the drawing of its objects is done.
See the individual fix page for a description of what these
parameters mean for a particular fix.

\sphinxAtStartPar
The only setting currently allowed for the \sphinxstyleemphasis{color} value is \sphinxstyleemphasis{type},
which will color the fix objects according to their type.  By default
the mapping of types to colors is as follows:
\begin{itemize}
\item {} 
\sphinxAtStartPar
type 1 = red

\item {} 
\sphinxAtStartPar
type 2 = green

\item {} 
\sphinxAtStartPar
type 3 = blue

\item {} 
\sphinxAtStartPar
type 4 = yellow

\item {} 
\sphinxAtStartPar
type 5 = aqua

\item {} 
\sphinxAtStartPar
type 6 = cyan

\end{itemize}

\sphinxAtStartPar
and repeats itself for types \textgreater{} 6.  There is not yet an option to
change this via the dump\_modify command.


\bigskip\hrule\bigskip


\sphinxAtStartPar
The \sphinxstyleemphasis{size} keyword sets the width and height of the created images,
i.e. the number of pixels in each direction.


\bigskip\hrule\bigskip


\sphinxAtStartPar
The \sphinxstyleemphasis{view}, \sphinxstyleemphasis{center}, \sphinxstyleemphasis{up}, and \sphinxstyleemphasis{zoom} values determine how
3d simulation space is mapped to the 2d plane of the image.  Basically
they control how the simulation box appears in the image.

\sphinxAtStartPar
All of the \sphinxstyleemphasis{view}, \sphinxstyleemphasis{center}, \sphinxstyleemphasis{up}, and \sphinxstyleemphasis{zoom} values can be
specified as numeric quantities, whose meaning is explained below.
Any of them can also be specified as an {\hyperref[\detokenize{variable::doc}]{\sphinxcrossref{\DUrole{doc}{equal\sphinxhyphen{}style variable}}}},
by using v\_name as the value, where “name” is
the variable name.  In this case the variable will be evaluated on the
timestep each image is created to create a new value.  If the
equal\sphinxhyphen{}style variable is time\sphinxhyphen{}dependent, this is a means of changing
the way the simulation box appears from image to image, effectively
doing a pan or fly\sphinxhyphen{}by view of your simulation.

\sphinxAtStartPar
The \sphinxstyleemphasis{view} keyword determines the viewpoint from which the simulation
box is viewed, looking towards the \sphinxstyleemphasis{center} point.  The \sphinxstyleemphasis{theta} value
is the vertical angle from the +z axis, and must be an angle from 0 to
180 degrees.  The \sphinxstyleemphasis{phi} value is an azimuthal angle around the z axis
and can be positive or negative.  A value of 0.0 is a view along the
+x axis, towards the \sphinxstyleemphasis{center} point.  If \sphinxstyleemphasis{theta} or \sphinxstyleemphasis{phi} are
specified via variables, then the variable values should be in
degrees.

\sphinxAtStartPar
The \sphinxstyleemphasis{center} keyword determines the point in simulation space that
will be at the center of the image.  \sphinxstyleemphasis{Cx}, \sphinxstyleemphasis{Cy}, and \sphinxstyleemphasis{Cz} are
specified as fractions of the box dimensions, so that (0.5,0.5,0.5) is
the center of the simulation box.  These values do not have to be
between 0.0 and 1.0, if you want the simulation box to be offset from
the center of the image.  Note, however, that if you choose strange
values for \sphinxstyleemphasis{Cx}, \sphinxstyleemphasis{Cy}, or \sphinxstyleemphasis{Cz} you may get a blank image.  Internally,
\sphinxstyleemphasis{Cx}, \sphinxstyleemphasis{Cy}, and \sphinxstyleemphasis{Cz} are converted into a point in simulation space.
If \sphinxstyleemphasis{flag} is set to “s” for static, then this conversion is done once,
at the time the dump command is issued.  If \sphinxstyleemphasis{flag} is set to “d” for
dynamic then the conversion is performed every time a new image is
created.  If the box size or shape is changing, this will adjust the
center point in simulation space.

\sphinxAtStartPar
The \sphinxstyleemphasis{up} keyword determines what direction in simulation space will be
“up” in the image.  Internally it is stored as a vector that is in the
plane perpendicular to the view vector implied by the \sphinxstyleemphasis{theta} and
\sphinxstyleemphasis{pni} values, and which is also in the plane defined by the view
vector and user\sphinxhyphen{}specified up vector.  Thus this internal vector is
computed from the user\sphinxhyphen{}specified \sphinxstyleemphasis{up} vector as

\begin{sphinxVerbatim}[commandchars=\\\{\}]
up\PYGZus{}internal = view cross (up cross view)
\end{sphinxVerbatim}

\sphinxAtStartPar
This means the only restriction on the specified \sphinxstyleemphasis{up} vector is that
it cannot be parallel to the \sphinxstyleemphasis{view} vector, implied by the \sphinxstyleemphasis{theta} and
\sphinxstyleemphasis{phi} values.

\sphinxAtStartPar
The \sphinxstyleemphasis{zoom} keyword scales the size of the simulation box as it appears
in the image.  The default \sphinxstyleemphasis{zfactor} value of 1 should display an
image mostly filled by the atoms in the simulation box.  A \sphinxstyleemphasis{zfactor} \textgreater{}
1 will make the simulation box larger; a \sphinxstyleemphasis{zfactor} \textless{} 1 will make it
smaller.  \sphinxstyleemphasis{Zfactor} must be a value \textgreater{} 0.0.


\bigskip\hrule\bigskip


\sphinxAtStartPar
The \sphinxstyleemphasis{box} keyword determines if and how the simulation box boundaries
are rendered as thin cylinders in the image.  If \sphinxstyleemphasis{no} is set, then the
box boundaries are not drawn and the \sphinxstyleemphasis{diam} setting is ignored.  If
\sphinxstyleemphasis{yes} is set, the 12 edges of the box are drawn, with a diameter that
is a fraction of the shortest box length in x,y,z (for 3d) or x,y (for
2d).  The color of the box boundaries can be set with the “dump\_modify
boxcolor” command.

\sphinxAtStartPar
The \sphinxstyleemphasis{axes} keyword determines if and how the coordinate axes are
rendered as thin cylinders in the image.  If \sphinxstyleemphasis{no} is set, then the
axes are not drawn and the \sphinxstyleemphasis{length} and \sphinxstyleemphasis{diam} settings are ignored.
If \sphinxstyleemphasis{yes} is set, 3 thin cylinders are drawn to represent the x,y,z
axes in colors red,green,blue.  The origin of these cylinders will be
offset from the lower left corner of the box by 10\%.  The \sphinxstyleemphasis{length}
setting determines how long the cylinders will be as a fraction of the
respective box lengths.  The \sphinxstyleemphasis{diam} setting determines their thickness
as a fraction of the shortest box length in x,y,z (for 3d) or x,y (for
2d).

\sphinxAtStartPar
The \sphinxstyleemphasis{subbox} keyword determines if and how processor subdomain
boundaries are rendered as thin cylinders in the image.  If \sphinxstyleemphasis{no} is
set (default), then the subdomain boundaries are not drawn and the
\sphinxstyleemphasis{diam} setting is ignored.  If \sphinxstyleemphasis{yes} is set, the 12 edges of each
processor subdomain are drawn, with a diameter that is a fraction of
the shortest box length in x,y,z (for 3d) or x,y (for 2d).  The color
of the subdomain boundaries can be set with the “dump\_modify
boxcolor” command.


\bigskip\hrule\bigskip


\sphinxAtStartPar
The \sphinxstyleemphasis{shiny} keyword determines how shiny the objects rendered in the
image will appear.  The \sphinxstyleemphasis{sfactor} value must be a value 0.0 \textless{}=
\sphinxstyleemphasis{sfactor} \textless{}= 1.0, where \sphinxstyleemphasis{sfactor} = 1 is a highly reflective surface
and \sphinxstyleemphasis{sfactor} = 0 is a rough non\sphinxhyphen{}shiny surface.

\sphinxAtStartPar
\DUrole{versionmodified,added}{Added in version 21Nov2023.}

\sphinxAtStartPar
The \sphinxstyleemphasis{fsaa} keyword can be used with the dump image command to improve
the image quality by enabling full scene anti\sphinxhyphen{}aliasing.  Internally the
image is rendered at twice the width and height and then scaled down by
computing the average of each 2x2 block of pixels to produce a single
pixel in the final image at the original size. This produces images with
smoother, less ragged edges.  The application of this algorithm can
increase the cost of computing the image by about 3x or more.

\sphinxAtStartPar
The \sphinxstyleemphasis{ssao} keyword turns on/off a screen space ambient occlusion (SSAO)
model for depth shading.  If \sphinxstyleemphasis{yes} is set, then atoms further away from
the viewer are darkened via a randomized process, which is perceived as
depth.  The strength of the effect can be scaled by the \sphinxstyleemphasis{dfactor}
parameter.  If \sphinxstyleemphasis{no} is set, no depth shading is performed.  The
calculation of this effect can increase the cost of computing the image
substantially by 5x or more, especially with larger images.  When used
in combination with the \sphinxstyleemphasis{fsaa} keyword the computational cost of depth
shading is particularly large.


\bigskip\hrule\bigskip



\subsection{Image Quality Settings}
\label{\detokenize{dump_image:image-quality-settings}}
\sphinxAtStartPar
The two keywords \sphinxstyleemphasis{fsaa} and \sphinxstyleemphasis{ssao} can be used to improve the image
quality at the expense of additional computational cost to render the
images. The images below show from left to right the same render with
default settings, with \sphinxstyleemphasis{fsaa} added, with \sphinxstyleemphasis{ssao} added, and with both
keywords added.

\sphinxAtStartPar
\sphinxincludegraphics[width=0.240\linewidth]{{image.default}.png}  \sphinxincludegraphics[width=0.240\linewidth]{{image.fsaa}.png}  \sphinxincludegraphics[width=0.240\linewidth]{{image.ssao}.png}  \sphinxincludegraphics[width=0.240\linewidth]{{image.both}.png}


\bigskip\hrule\bigskip


\sphinxAtStartPar
A series of JPEG, PNG, or PPM images can be converted into a movie
file and then played as a movie using commonly available tools. Using
dump style \sphinxstyleemphasis{movie} automates this step and avoids the intermediate
step of writing (many) image snapshot file. But LAMMPS has to be
compiled with \sphinxhyphen{}DLAMMPS\_FFMPEG and an FFmpeg executable have to be
installed.

\sphinxAtStartPar
To manually convert JPEG, PNG or PPM files into an animated GIF or
MPEG or other movie file you can use:
\begin{itemize}
\item {} \begin{enumerate}
\sphinxsetlistlabels{\alph}{enumi}{enumii}{}{)}%
\item {} 
\sphinxAtStartPar
Use the ImageMagick convert program.

\end{enumerate}

\begin{sphinxVerbatim}[commandchars=\\\{\}]
convert\PYG{+w}{ }*.jpg\PYG{+w}{ }foo.gif
convert\PYG{+w}{ }\PYGZhy{}loop\PYG{+w}{ }\PYG{l+m}{1}\PYG{+w}{ }*.ppm\PYG{+w}{ }foo.mpg
\end{sphinxVerbatim}

\sphinxAtStartPar
Animated GIF files from ImageMagick are not optimized. You can use
a program like gifsicle to optimize and thus massively shrink them.
MPEG files created by ImageMagick are in MPEG\sphinxhyphen{}1 format with a rather
inefficient compression and low quality compared to more modern
compression styles like MPEG\sphinxhyphen{}4, H.264, VP8, VP9, H.265 and so on.

\item {} \begin{enumerate}
\sphinxsetlistlabels{\alph}{enumi}{enumii}{}{)}%
\setcounter{enumi}{1}
\item {} 
\sphinxAtStartPar
Use QuickTime.

\end{enumerate}

\sphinxAtStartPar
Select “Open Image Sequence” under the File menu Load the images into
QuickTime to animate them Select “Export” under the File menu Save the
movie as a QuickTime movie (*.mov) or in another format.  QuickTime
can generate very high quality and efficiently compressed movie
files. Some of the supported formats require to buy a license and some
are not readable on all platforms until specific runtime libraries are
installed.

\item {} \begin{enumerate}
\sphinxsetlistlabels{\alph}{enumi}{enumii}{}{)}%
\setcounter{enumi}{2}
\item {} 
\sphinxAtStartPar
Use FFmpeg

\end{enumerate}

\sphinxAtStartPar
FFmpeg is a command line tool that is available on many platforms and
allows extremely flexible encoding and decoding of movies.

\begin{sphinxVerbatim}[commandchars=\\\{\}]
cat\PYG{+w}{ }snap.*.jpg\PYG{+w}{ }\PYG{p}{|}\PYG{+w}{ }ffmpeg\PYG{+w}{ }\PYGZhy{}y\PYG{+w}{ }\PYGZhy{}f\PYG{+w}{ }image2pipe\PYG{+w}{ }\PYGZhy{}c:v\PYG{+w}{ }mjpeg\PYG{+w}{ }\PYGZhy{}i\PYG{+w}{ }\PYGZhy{}\PYG{+w}{ }\PYGZhy{}b:v\PYG{+w}{ }2000k\PYG{+w}{ }movie.m4v
cat\PYG{+w}{ }snap.*.ppm\PYG{+w}{ }\PYG{p}{|}\PYG{+w}{ }ffmpeg\PYG{+w}{ }\PYGZhy{}y\PYG{+w}{ }\PYGZhy{}f\PYG{+w}{ }image2pipe\PYG{+w}{ }\PYGZhy{}c:v\PYG{+w}{ }ppm\PYG{+w}{ }\PYGZhy{}i\PYG{+w}{ }\PYGZhy{}\PYG{+w}{ }\PYGZhy{}b:v\PYG{+w}{ }2400k\PYG{+w}{ }movie.avi
\end{sphinxVerbatim}

\sphinxAtStartPar
Front ends for FFmpeg exist for multiple platforms. For more
information see the \sphinxhref{https://ffmpeg.org/}{FFmpeg homepage}

\end{itemize}


\bigskip\hrule\bigskip


\sphinxAtStartPar
Play the movie:
\begin{itemize}
\item {} \begin{enumerate}
\sphinxsetlistlabels{\alph}{enumi}{enumii}{}{)}%
\item {} 
\sphinxAtStartPar
Use your browser to view an animated GIF movie.

\end{enumerate}

\sphinxAtStartPar
Select “Open File” under the File menu
Load the animated GIF file

\item {} 
\sphinxAtStartPar
b) Use the freely available mplayer or ffplay tool to view a
movie. Both are available for multiple OSes and support a large
variety of file formats and decoders.

\begin{sphinxVerbatim}[commandchars=\\\{\}]
mplayer\PYG{+w}{ }foo.mpg
ffplay\PYG{+w}{ }bar.avi
\end{sphinxVerbatim}

\item {} 
\sphinxAtStartPar
c) Use the \sphinxhref{https://lammps.github.io/pizza}{Pizza.py}
\sphinxhref{https://lammps.github.io/pizza/doc/animate.html}{animate tool},
which works directly on a series of image files.

\begin{sphinxVerbatim}[commandchars=\\\{\}]
\PYG{n}{a} \PYG{o}{=} \PYG{n}{animate}\PYG{p}{(}\PYG{l+s+s2}{\PYGZdq{}}\PYG{l+s+s2}{foo*.jpg}\PYG{l+s+s2}{\PYGZdq{}}\PYG{p}{)}
\end{sphinxVerbatim}

\item {} 
\sphinxAtStartPar
d) QuickTime and other Windows\sphinxhyphen{} or macOS\sphinxhyphen{}based media players can
obviously play movie files directly. Similarly for corresponding tools
bundled with Linux desktop environments.  However, due to licensing
issues with some file formats, the formats may require installing
additional libraries, purchasing a license, or may not be
supported.

\end{itemize}


\bigskip\hrule\bigskip



\subsection{Dump\_modify keywords for dump image and dump movie}
\label{\detokenize{dump_image:dump-modify-keywords-for-dump-image-and-dump-movie}}
\sphinxAtStartPar
The following dump\_modify keywords apply only to the dump image and
dump movie styles.  Any keyword that works with dump image also works
with dump movie, since the movie is simply a collection of images.
Some of the keywords only affect the dump movie style.  The
descriptions give details.


\bigskip\hrule\bigskip


\sphinxAtStartPar
The \sphinxstyleemphasis{acolor} keyword can be used with the dump image command, when its
atom color setting is \sphinxstyleemphasis{type}, to set the color that atoms of each type
will be drawn in the image.

\sphinxAtStartPar
The specified \sphinxstyleemphasis{type} should be a type label or integer from 1 to Ntypes
= the number of atom types.  For numeric types, a wildcard asterisk can
be used in place of or in conjunction with the \sphinxstyleemphasis{type} argument to
specify a range of atom types.  This takes the form “*” or “*n” or
“n*” or “m*n”. If N = the number of atom types, then an asterisk with
no numeric values means all types from 1 to N.  A leading asterisk
means all types from 1 to n (inclusive).  A trailing asterisk means all
types from n to N (inclusive).  A middle asterisk means all types from
m to n (inclusive).

\sphinxAtStartPar
The specified \sphinxstyleemphasis{color} can be a single color which is any of the 140
pre\sphinxhyphen{}defined colors (see below) or a color name defined by the
“dump\_modify color” command, as described below.  Or it can be two or
more colors separated by a “/” character, e.g. red/green/blue.  In the
former case, that color is assigned to all the specified atom types.
In the latter case, the list of colors are assigned in a round\sphinxhyphen{}robin
fashion to each of the specified atom types.


\bigskip\hrule\bigskip


\sphinxAtStartPar
The \sphinxstyleemphasis{adiam} keyword can be used with the dump image command, when its
atom diameter setting is \sphinxstyleemphasis{type}, to set the size that atoms of each
type will be drawn in the image.  The specified \sphinxstyleemphasis{type} should be a type
label or integer from 1 to Ntypes.  As with the \sphinxstyleemphasis{acolor} keyword, a
wildcard asterisk can be used as part of the \sphinxstyleemphasis{type} argument to specify
a range of numeric atom types.  The specified \sphinxstyleemphasis{diam} is the size in
whatever distance {\hyperref[\detokenize{units::doc}]{\sphinxcrossref{\DUrole{doc}{units}}}} the input script is using, e.g.
Angstroms.


\bigskip\hrule\bigskip


\sphinxAtStartPar
The \sphinxstyleemphasis{amap} keyword can be used with the dump image command, with its
\sphinxstyleemphasis{atom} keyword, when its atom setting is an atom\sphinxhyphen{}attribute, to setup a
color map.  The color map is used to assign a specific RGB
(red/green/blue) color value to an individual atom when it is drawn,
based on the atom’s attribute, which is a numeric value, e.g. its
x\sphinxhyphen{}component of velocity if the atom\sphinxhyphen{}attribute “vx” was specified.

\sphinxAtStartPar
The basic idea of a color map is that the atom\sphinxhyphen{}attribute will be
within a range of values, and that range is associated with a series
of colors (e.g. red, blue, green).  An atom’s specific value (vx =
\sphinxhyphen{}3.2) can then mapped to the series of colors (e.g. halfway between
red and blue), and a specific color is determined via an interpolation
procedure.

\sphinxAtStartPar
There are many possible options for the color map, enabled by the
\sphinxstyleemphasis{amap} keyword.  Here are the details.

\sphinxAtStartPar
The \sphinxstyleemphasis{lo} and \sphinxstyleemphasis{hi} settings determine the range of values allowed for
the atom attribute.  If numeric values are used for \sphinxstyleemphasis{lo} and/or \sphinxstyleemphasis{hi},
then values that are lower/higher than that value are set to the
value.  I.e. the range is static.  If \sphinxstyleemphasis{lo} is specified as \sphinxstyleemphasis{min} or
\sphinxstyleemphasis{hi} as \sphinxstyleemphasis{max} then the range is dynamic, and the lower and/or
upper bound will be calculated each time an image is drawn, based
on the set of atoms being visualized.

\sphinxAtStartPar
The \sphinxstyleemphasis{style} setting is two letters, such as “ca”.  The first letter is
either “c” for continuous, “d” for discrete, or “s” for sequential.
The second letter is either “a” for absolute, or “f” for fractional.

\sphinxAtStartPar
A continuous color map is one in which the color changes continuously
from value to value within the range.  A discrete color map is one in
which discrete colors are assigned to sub\sphinxhyphen{}ranges of values within the
range.  A sequential color map is one in which discrete colors are
assigned to a sequence of sub\sphinxhyphen{}ranges of values covering the entire
range.

\sphinxAtStartPar
An absolute color map is one in which the values to which colors are
assigned are specified explicitly as values within the range.  A
fractional color map is one in which the values to which colors are
assigned are specified as a fractional portion of the range.  For
example if the range is from \sphinxhyphen{}10.0 to 10.0, and the color red is to be
assigned to atoms with a value of 5.0, then for an absolute color map
the number 5.0 would be used.  But for a fractional map, the number
0.75 would be used since 5.0 is 3/4 of the way from \sphinxhyphen{}10.0 to 10.0.

\sphinxAtStartPar
The \sphinxstyleemphasis{delta} setting must be specified for all styles, but is only used
for the sequential style; otherwise the value is ignored.  It
specifies the bin size to use within the range for assigning
consecutive colors to.  For example, if the range is from \(-10.0\) to
\(10.0\) and a \sphinxstyleemphasis{delta} of \(1.0\) is used, then 20 colors will be
assigned to the range.  The first will be from
\(-10.0 \le \text{color1} < -9.0\), then second from
\(-9.0 \le color2 < -8.0\), etc.

\sphinxAtStartPar
The \sphinxstyleemphasis{N} setting is how many entries follow.  The format of the entries
depends on whether the color map style is continuous, discrete or
sequential.  In all cases the \sphinxstyleemphasis{color} setting can be any of the 140
pre\sphinxhyphen{}defined colors (see below) or a color name defined by the
dump\_modify color option.

\sphinxAtStartPar
For continuous color maps, each entry has a \sphinxstyleemphasis{value} and a \sphinxstyleemphasis{color}.
The \sphinxstyleemphasis{value} is either a number within the range of values or \sphinxstyleemphasis{min} or
\sphinxstyleemphasis{max}.  The \sphinxstyleemphasis{value} of the first entry must be \sphinxstyleemphasis{min} and the \sphinxstyleemphasis{value}
of the last entry must be \sphinxstyleemphasis{max}.  Any entries in between must have
increasing values.  Note that numeric values can be specified either
as absolute numbers or as fractions (0.0 to 1.0) of the range,
depending on the “a” or “f” in the style setting for the color map.

\sphinxAtStartPar
Here is how the entries are used to determine the color of an individual
atom, given the value \(X\) of its atom attribute.  \(X\) will
fall between 2 of the entry values.  The color of the atom is linearly
interpolated (in each of the RGB values) between the 2 colors associated
with those entries.  For example, if \(X = -5.0\) and the two
surrounding entries are “red” at \(-10.0\) and “blue” at
\(0.0\), then the atom’s color will be halfway between “red” and
“blue”, which happens to be “purple”.

\sphinxAtStartPar
For discrete color maps, each entry has a \sphinxstyleemphasis{lo} and \sphinxstyleemphasis{hi} value and a
\sphinxstyleemphasis{color}.  The \sphinxstyleemphasis{lo} and \sphinxstyleemphasis{hi} settings are either numbers within the
range of values or \sphinxstyleemphasis{lo} can be \sphinxstyleemphasis{min} or \sphinxstyleemphasis{hi} can be \sphinxstyleemphasis{max}.  The \sphinxstyleemphasis{lo}
and \sphinxstyleemphasis{hi} settings of the last entry must be \sphinxstyleemphasis{min} and \sphinxstyleemphasis{max}.  Other
entries can have any \sphinxstyleemphasis{lo} and \sphinxstyleemphasis{hi} values and the sub\sphinxhyphen{}ranges of
different values can overlap.  Note that numeric \sphinxstyleemphasis{lo} and \sphinxstyleemphasis{hi} values
can be specified either as absolute numbers or as fractions (0.0 to 1.0)
of the range, depending on the “a” or “f” in the style setting for the
color map.

\sphinxAtStartPar
Here is how the entries are used to determine the color of an individual
atom, given the value X of its atom attribute.  The entries are scanned
from first to last.  The first time that \sphinxstyleemphasis{lo} \textless{}= X \textless{}= \sphinxstyleemphasis{hi}, X is
assigned the color associated with that entry.  You can think of the
last entry as assigning a default color (since it will always be matched
by X), and the earlier entries as colors that override the default.
Also note that no interpolation of a color RGB is done.  All atoms will
be drawn with one of the colors in the list of entries.

\sphinxAtStartPar
For sequential color maps, each entry has only a \sphinxstyleemphasis{color}.  Here is how
the entries are used to determine the color of an individual atom,
given the value X of its atom attribute.  The range is partitioned
into N bins of width \sphinxstyleemphasis{binsize}.  Thus X will fall in a specific bin
from 1 to N, say the Mth bin.  If it falls on a boundary between 2
bins, it is considered to be in the higher of the 2 bins.  Each bin is
assigned a color from the E entries.  If E \textless{} N, then the colors are
repeated.  For example if 2 entries with colors red and green are
specified, then the odd numbered bins will be red and the even bins
green.  The color of the atom is the color of its bin.  Note that the
sequential color map is really a shorthand way of defining a discrete
color map without having to specify where all the bin boundaries are.

\sphinxAtStartPar
Here is an example of using a sequential color map to color all the
atoms in individual molecules with a different color.  See the
examples/pour/in.pour.2d.molecule input script for an example of how
this is used.

\begin{sphinxVerbatim}[commandchars=\\\{\}]
\PYG{k}{variable        }\PYG{n+nv+nvIdentifier}{colors}\PYG{+w}{ }\PYG{n}{string}\PYG{+w}{ }\PYG{o}{\PYGZam{}}
\PYG{+w}{                }\PYG{l+s}{\PYGZdq{}}\PYG{l+s}{red green blue yellow white \PYGZam{}}
\PYG{l+s}{                purple pink orange lime gray}\PYG{l+s}{\PYGZdq{}}
\PYG{k}{variable        }\PYG{n+nv+nvIdentifier}{mol}\PYG{+w}{ }\PYG{n}{atom}\PYG{+w}{ }\PYG{n}{mol}\PYG{o}{\PYGZpc{}}\PYG{l+m}{10}
\PYG{k}{dump            }\PYG{n+nv+nvIdentifier}{1}\PYG{+w}{ }\PYG{n+nv+nvIdentifier}{all}\PYG{+w}{ }\PYG{n}{image}\PYG{+w}{ }\PYG{l+m}{250}\PYG{+w}{ }\PYG{n}{image.}\PYG{o}{*}\PYG{n}{.jpg}\PYG{+w}{ }\PYG{n}{v\PYGZus{}mol}\PYG{+w}{ }\PYG{n}{type}\PYG{+w}{ }\PYG{o}{\PYGZam{}}
\PYG{+w}{                }\PYG{n}{zoom}\PYG{+w}{ }\PYG{l+m}{1.6}\PYG{+w}{ }\PYG{n}{adiam}\PYG{+w}{ }\PYG{l+m}{1.5}
\PYG{k}{dump\PYGZus{}modify     }\PYG{n+nv+nvIdentifier}{1}\PYG{+w}{ }\PYG{n}{pad}\PYG{+w}{ }\PYG{l+m}{5}\PYG{+w}{ }\PYG{n}{amap}\PYG{+w}{ }\PYG{l+m}{0}\PYG{+w}{ }\PYG{l+m}{10}\PYG{+w}{ }\PYG{n}{sa}\PYG{+w}{ }\PYG{l+m}{1}\PYG{+w}{ }\PYG{l+m}{10}\PYG{+w}{ }\PYG{n+nv}{\PYGZdl{}\PYGZob{}}\PYG{n+nv}{colors}\PYG{n+nv}{\PYGZcb{}}
\end{sphinxVerbatim}

\sphinxAtStartPar
In this case, 10 colors are defined, and molecule IDs are
mapped to one of the colors, even if there are 1000s of molecules.


\bigskip\hrule\bigskip


\sphinxAtStartPar
The \sphinxstyleemphasis{backcolor} sets the background color of the images.  The color
name can be any of the 140 pre\sphinxhyphen{}defined colors (see below) or a color
name defined by the dump\_modify color option.


\bigskip\hrule\bigskip


\sphinxAtStartPar
The \sphinxstyleemphasis{bcolor} keyword can be used with the dump image command, with its
\sphinxstyleemphasis{bond} keyword, when its color setting is \sphinxstyleemphasis{type}, to set the color
that bonds of each type will be drawn in the image.

\sphinxAtStartPar
The specified \sphinxstyleemphasis{type} should be a type label or integer from 1 to
\(N\), where \(N\) is the number of bond types.  For numeric
types, a wildcard asterisk can be used in place of or in conjunction
with the \sphinxstyleemphasis{type} argument to specify a range of bond types.  This takes
the form “*” or “*n” or “m*” or “m*n”.  If \(N\) is the number
of bond types, then an asterisk with no numerical values means all
types from 1 to \(N\).  A leading asterisk means all types from 1
to n (inclusive).  A trailing asterisk means all types from m to
\(N\) (inclusive).  A middle asterisk means all types from m to n
(inclusive).

\sphinxAtStartPar
The specified \sphinxstyleemphasis{color} can be a single color which is any of the 140
pre\sphinxhyphen{}defined colors (see below) or a color name defined by the
dump\_modify color option.  Or it can be two or more colors separated
by a “/” character (e.g., red/green/blue).  In the former case, that
color is assigned to all the specified bond types.  In the latter
case, the list of colors are assigned in a round\sphinxhyphen{}robin fashion to each
of the specified bond types.


\bigskip\hrule\bigskip


\sphinxAtStartPar
The \sphinxstyleemphasis{bdiam} keyword can be used with the dump image command, with its
\sphinxstyleemphasis{bond} keyword, when its \sphinxstyleemphasis{diam} setting is \sphinxstyleemphasis{type}, to set the diameter
that bonds of each type will be drawn in the image.  The specified
\sphinxstyleemphasis{type} should be a type label or integer from 1 to Nbondtypes.  As with
the \sphinxstyleemphasis{bcolor} keyword, a wildcard asterisk can be used as part of the
\sphinxstyleemphasis{type} argument to specify a range of numeric bond types.  The
specified \sphinxstyleemphasis{diam} is the size in whatever distance {\hyperref[\detokenize{units::doc}]{\sphinxcrossref{\DUrole{doc}{units}}}}
you are using (e.g., Angstroms).


\bigskip\hrule\bigskip


\sphinxAtStartPar
The \sphinxstyleemphasis{bitrate} keyword can be used with the {\hyperref[\detokenize{dump_image::doc}]{\sphinxcrossref{\DUrole{doc}{dump movie}}}} command to define the size of the resulting movie file
and its quality via setting how many kbits per second are to be used
for the movie file. Higher bitrates require less compression and will
result in higher quality movies.  The quality is also determined by
the compression format and encoder.  The default setting is 2000
kbit/s, which will result in average quality with older compression
formats.

\begin{sphinxadmonition}{note}{Note:}
\sphinxAtStartPar
Not all movie file formats supported by dump movie allow the
bitrate to be set.  If not, the setting is silently ignored.
\end{sphinxadmonition}


\bigskip\hrule\bigskip


\sphinxAtStartPar
The \sphinxstyleemphasis{boxcolor} keyword sets the color of the simulation box drawn
around the atoms in each image as well as the color of processor
subdomain boundaries.  See the “dump image box” command for how to
specify that a box be drawn via the \sphinxstyleemphasis{box} keyword, and the subdomain
boundaries via the \sphinxstyleemphasis{subbox} keyword.  The color name can be any of the
140 pre\sphinxhyphen{}defined colors (see below) or a color name defined by the
dump\_modify color option.


\bigskip\hrule\bigskip


\sphinxAtStartPar
The \sphinxstyleemphasis{color} keyword allows definition of a new color name, in addition
to the 140\sphinxhyphen{}predefined colors (see below), and associates three
red/green/blue RGB values with that color name.  The color name can
then be used with any other dump\_modify keyword that takes a color
name as a value.  The RGB values should each be floating point values
between 0.0 and 1.0 inclusive.

\sphinxAtStartPar
When a color name is converted to RGB values, the user\sphinxhyphen{}defined color
names are searched first, then the 140 pre\sphinxhyphen{}defined color names.  This
means you can also use the \sphinxstyleemphasis{color} keyword to overwrite one of the
pre\sphinxhyphen{}defined color names with new RBG values.


\bigskip\hrule\bigskip


\sphinxAtStartPar
The \sphinxstyleemphasis{framerate} keyword can be used with the {\hyperref[\detokenize{dump_image::doc}]{\sphinxcrossref{\DUrole{doc}{dump movie}}}} command to define the duration of the resulting movie
file.  Movie files written by the dump \sphinxstyleemphasis{movie} command have a default
frame rate of 24 frames per second and the images generated will be
converted at that rate.  Thus a sequence of 1000 dump images will
result in a movie of about 42 seconds.  To make a movie run longer you
can either generate images more frequently or lower the frame rate.
To speed a movie up, you can do the inverse.  Using a frame rate
higher than 24 is not recommended, as it will result in simply
dropping the rendered images. It is more efficient to dump images less
frequently.


\bigskip\hrule\bigskip


\sphinxAtStartPar
The \sphinxstyleemphasis{gmap} keyword can be used with the dump image command, with its
\sphinxstyleemphasis{grid} keyword, to setup a color map.  The color map is used to assign
a specific RGB (red/green/blue) color value to an individual grid cell
when it is drawn, based on the grid cell value, which is a numeric
quantity specified with the \sphinxstyleemphasis{grid} keyword.

\sphinxAtStartPar
The arguments for the \sphinxstyleemphasis{gmap} keyword are identical to those for the
\sphinxstyleemphasis{amap} keyword (for atom coloring) described above.


\bigskip\hrule\bigskip



\subsection{Restrictions}
\label{\detokenize{dump_image:restrictions}}
\sphinxAtStartPar
To write JPEG images, you must use the \sphinxhyphen{}DLAMMPS\_JPEG switch when
building LAMMPS and link with a JPEG library. To write PNG images, you
must use the \sphinxhyphen{}DLAMMPS\_PNG switch when building LAMMPS and link with a
PNG library.

\sphinxAtStartPar
To write \sphinxstyleemphasis{movie} dumps, you must use the \sphinxhyphen{}DLAMMPS\_FFMPEG switch when
building LAMMPS and have the FFmpeg executable available on the
machine where LAMMPS is being run.  Typically its name is lowercase
(i.e., “ffmpeg”).

\sphinxAtStartPar
See the {\hyperref[\detokenize{Build_settings::doc}]{\sphinxcrossref{\DUrole{doc}{Build settings}}}} page for details.

\sphinxAtStartPar
Note that since FFmpeg is run as an external program via a pipe,
LAMMPS has limited control over its execution and no knowledge about
errors and warnings printed by it. Those warnings and error messages
will be printed to the screen only. Due to the way image data are
communicated to FFmpeg, it will often print the message

\begin{sphinxVerbatim}[commandchars=\\\{\}]
pipe:: Input/output error
\end{sphinxVerbatim}

\sphinxAtStartPar
which can be safely ignored. Other warnings
and errors have to be addressed according to the FFmpeg documentation.
One known issue is that certain movie file formats (e.g., MPEG level 1
and 2 format streams) have video bandwidth limits that can be crossed
when rendering too large of image sizes. Typical warnings look like
this:

\begin{sphinxVerbatim}[commandchars=\\\{\}]
[mpeg @ 0x98b5e0] packet too large, ignoring buffer limits to mux it
[mpeg @ 0x98b5e0] buffer underflow st=0 bufi=281407 size=285018
[mpeg @ 0x98b5e0] buffer underflow st=0 bufi=283448 size=285018
\end{sphinxVerbatim}

\sphinxAtStartPar
In this case it is recommended either to reduce the size of the image
or to encode in a different format that is also supported by your copy of
FFmpeg and which does not have this limitation (e.g., .avi, .mkv, mp4).


\subsection{Related commands}
\label{\detokenize{dump_image:related-commands}}
\sphinxAtStartPar
{\hyperref[\detokenize{dump::doc}]{\sphinxcrossref{\DUrole{doc}{dump}}}}, {\hyperref[\detokenize{dump_modify::doc}]{\sphinxcrossref{\DUrole{doc}{dump\_modify}}}}, {\hyperref[\detokenize{undump::doc}]{\sphinxcrossref{\DUrole{doc}{undump}}}}


\subsection{Default}
\label{\detokenize{dump_image:default}}
\sphinxAtStartPar
The defaults for the dump image and dump movie keywords are as follows:
\begin{itemize}
\item {} 
\sphinxAtStartPar
adiam = not specified (use diameter setting)

\item {} 
\sphinxAtStartPar
atom = yes

\item {} 
\sphinxAtStartPar
bond = none none (if no bonds in system)

\item {} 
\sphinxAtStartPar
bond = atom 0.5 (if bonds in system)

\item {} 
\sphinxAtStartPar
size = 512 512

\item {} 
\sphinxAtStartPar
view = 60 30 (for 3d)

\item {} 
\sphinxAtStartPar
view = 0 0 (for 2d)

\item {} 
\sphinxAtStartPar
center = s 0.5 0.5 0.5

\item {} 
\sphinxAtStartPar
up = 0 0 1 (for 3d)

\item {} 
\sphinxAtStartPar
up = 0 1 0 (for 2d)

\item {} 
\sphinxAtStartPar
zoom = 1.0

\item {} 
\sphinxAtStartPar
box = yes 0.02

\item {} 
\sphinxAtStartPar
axes = no 0.0 0.0

\item {} 
\sphinxAtStartPar
subbox no 0.0

\item {} 
\sphinxAtStartPar
shiny = 1.0

\item {} 
\sphinxAtStartPar
ssao = no

\end{itemize}


\bigskip\hrule\bigskip


\sphinxAtStartPar
The defaults for the dump\_modify keywords specific to dump image and dump movie are as follows:
\begin{itemize}
\item {} 
\sphinxAtStartPar
acolor = * red/green/blue/yellow/aqua/cyan

\item {} 
\sphinxAtStartPar
adiam = * 1.0

\item {} 
\sphinxAtStartPar
amap = min max cf 0.0 2 min blue max red

\item {} 
\sphinxAtStartPar
backcolor = black

\item {} 
\sphinxAtStartPar
bcolor = * red/green/blue/yellow/aqua/cyan

\item {} 
\sphinxAtStartPar
bdiam = * 0.5

\item {} 
\sphinxAtStartPar
bitrate = 2000

\item {} 
\sphinxAtStartPar
boxcolor = yellow

\item {} 
\sphinxAtStartPar
color = 140 color names are pre\sphinxhyphen{}defined as listed below

\item {} 
\sphinxAtStartPar
framerate = 24

\item {} 
\sphinxAtStartPar
fsaa = no

\item {} 
\sphinxAtStartPar
gmap = min max cf 0.0 2 min blue max red

\end{itemize}


\bigskip\hrule\bigskip


\sphinxAtStartPar
These are the standard 109 element names that LAMMPS pre\sphinxhyphen{}defines for
use with the dump image and dump\_modify commands.
\begin{itemize}
\item {} 
\sphinxAtStartPar
1\sphinxhyphen{}10 = “H”, “He”, “Li”, “Be”, “B”, “C”, “N”, “O”, “F”, “Ne”

\item {} 
\sphinxAtStartPar
11\sphinxhyphen{}20 = “Na”, “Mg”, “Al”, “Si”, “P”, “S”, “Cl”, “Ar”, “K”, “Ca”

\item {} 
\sphinxAtStartPar
21\sphinxhyphen{}30 = “Sc”, “Ti”, “V”, “Cr”, “Mn”, “Fe”, “Co”, “Ni”, “Cu”, “Zn”

\item {} 
\sphinxAtStartPar
31\sphinxhyphen{}40 = “Ga”, “Ge”, “As”, “Se”, “Br”, “Kr”, “Rb”, “Sr”, “Y”, “Zr”

\item {} 
\sphinxAtStartPar
41\sphinxhyphen{}50 = “Nb”, “Mo”, “Tc”, “Ru”, “Rh”, “Pd”, “Ag”, “Cd”, “In”, “Sn”

\item {} 
\sphinxAtStartPar
51\sphinxhyphen{}60 = “Sb”, “Te”, “I”, “Xe”, “Cs”, “Ba”, “La”, “Ce”, “Pr”, “Nd”

\item {} 
\sphinxAtStartPar
61\sphinxhyphen{}70 = “Pm”, “Sm”, “Eu”, “Gd”, “Tb”, “Dy”, “Ho”, “Er”, “Tm”, “Yb”

\item {} 
\sphinxAtStartPar
71\sphinxhyphen{}80 = “Lu”, “Hf”, “Ta”, “W”, “Re”, “Os”, “Ir”, “Pt”, “Au”, “Hg”

\item {} 
\sphinxAtStartPar
81\sphinxhyphen{}90 = “Tl”, “Pb”, “Bi”, “Po”, “At”, “Rn”, “Fr”, “Ra”, “Ac”, “Th”

\item {} 
\sphinxAtStartPar
91\sphinxhyphen{}100 = “Pa”, “U”, “Np”, “Pu”, “Am”, “Cm”, “Bk”, “Cf”, “Es”, “Fm”

\item {} 
\sphinxAtStartPar
101\sphinxhyphen{}109 = “Md”, “No”, “Lr”, “Rf”, “Db”, “Sg”, “Bh”, “Hs”, “Mt”

\end{itemize}


\bigskip\hrule\bigskip


\sphinxAtStartPar
These are the 140 colors that LAMMPS pre\sphinxhyphen{}defines for use with the dump
image and dump\_modify commands.  Additional colors can be defined with
the dump\_modify color command.  The 3 numbers listed for each name are
the RGB (red/green/blue) values.  Divide each value by 255 to get the
equivalent 0.0 to 1.0 value.


\begin{savenotes}\sphinxattablestart
\sphinxthistablewithglobalstyle
\centering
\begin{tabulary}{\linewidth}[t]{TTTTT}
\sphinxtoprule
\sphinxtableatstartofbodyhook
\sphinxAtStartPar
aliceblue = 240, 248, 255
&
\sphinxAtStartPar
antiquewhite = 250, 235, 215
&
\sphinxAtStartPar
aqua = 0, 255, 255
&
\sphinxAtStartPar
aquamarine = 127, 255, 212
&
\sphinxAtStartPar
azure = 240, 255, 255
\\
\sphinxhline
\sphinxAtStartPar
beige = 245, 245, 220
&
\sphinxAtStartPar
bisque = 255, 228, 196
&
\sphinxAtStartPar
black = 0, 0, 0
&
\sphinxAtStartPar
blanchedalmond = 255, 255, 205
&
\sphinxAtStartPar
blue = 0, 0, 255
\\
\sphinxhline
\sphinxAtStartPar
blueviolet = 138, 43, 226
&
\sphinxAtStartPar
brown = 165, 42, 42
&
\sphinxAtStartPar
burlywood = 222, 184, 135
&
\sphinxAtStartPar
cadetblue = 95, 158, 160
&
\sphinxAtStartPar
chartreuse = 127, 255, 0
\\
\sphinxhline
\sphinxAtStartPar
chocolate = 210, 105, 30
&
\sphinxAtStartPar
coral = 255, 127, 80
&
\sphinxAtStartPar
cornflowerblue = 100, 149, 237
&
\sphinxAtStartPar
cornsilk = 255, 248, 220
&
\sphinxAtStartPar
crimson = 220, 20, 60
\\
\sphinxhline
\sphinxAtStartPar
cyan = 0, 255, 255
&
\sphinxAtStartPar
darkblue = 0, 0, 139
&
\sphinxAtStartPar
darkcyan = 0, 139, 139
&
\sphinxAtStartPar
darkgoldenrod = 184, 134, 11
&
\sphinxAtStartPar
darkgray = 169, 169, 169
\\
\sphinxhline
\sphinxAtStartPar
darkgreen = 0, 100, 0
&
\sphinxAtStartPar
darkkhaki = 189, 183, 107
&
\sphinxAtStartPar
darkmagenta = 139, 0, 139
&
\sphinxAtStartPar
darkolivegreen = 85, 107, 47
&
\sphinxAtStartPar
darkorange = 255, 140, 0
\\
\sphinxhline
\sphinxAtStartPar
darkorchid = 153, 50, 204
&
\sphinxAtStartPar
darkred = 139, 0, 0
&
\sphinxAtStartPar
darksalmon = 233, 150, 122
&
\sphinxAtStartPar
darkseagreen = 143, 188, 143
&
\sphinxAtStartPar
darkslateblue = 72, 61, 139
\\
\sphinxhline
\sphinxAtStartPar
darkslategray = 47, 79, 79
&
\sphinxAtStartPar
darkturquoise = 0, 206, 209
&
\sphinxAtStartPar
darkviolet = 148, 0, 211
&
\sphinxAtStartPar
deeppink = 255, 20, 147
&
\sphinxAtStartPar
deepskyblue = 0, 191, 255
\\
\sphinxhline
\sphinxAtStartPar
dimgray = 105, 105, 105
&
\sphinxAtStartPar
dodgerblue = 30, 144, 255
&
\sphinxAtStartPar
firebrick = 178, 34, 34
&
\sphinxAtStartPar
floralwhite = 255, 250, 240
&
\sphinxAtStartPar
forestgreen = 34, 139, 34
\\
\sphinxhline
\sphinxAtStartPar
fuchsia = 255, 0, 255
&
\sphinxAtStartPar
gainsboro = 220, 220, 220
&
\sphinxAtStartPar
ghostwhite = 248, 248, 255
&
\sphinxAtStartPar
gold = 255, 215, 0
&
\sphinxAtStartPar
goldenrod = 218, 165, 32
\\
\sphinxhline
\sphinxAtStartPar
gray = 128, 128, 128
&
\sphinxAtStartPar
green = 0, 128, 0
&
\sphinxAtStartPar
greenyellow = 173, 255, 47
&
\sphinxAtStartPar
honeydew = 240, 255, 240
&
\sphinxAtStartPar
hotpink = 255, 105, 180
\\
\sphinxhline
\sphinxAtStartPar
indianred = 205, 92, 92
&
\sphinxAtStartPar
indigo = 75, 0, 130
&
\sphinxAtStartPar
ivory = 255, 240, 240
&
\sphinxAtStartPar
khaki = 240, 230, 140
&
\sphinxAtStartPar
lavender = 230, 230, 250
\\
\sphinxhline
\sphinxAtStartPar
lavenderblush = 255, 240, 245
&
\sphinxAtStartPar
lawngreen = 124, 252, 0
&
\sphinxAtStartPar
lemonchiffon = 255, 250, 205
&
\sphinxAtStartPar
lightblue = 173, 216, 230
&
\sphinxAtStartPar
lightcoral = 240, 128, 128
\\
\sphinxhline
\sphinxAtStartPar
lightcyan = 224, 255, 255
&
\sphinxAtStartPar
lightgoldenrodyellow = 250, 250, 210
&
\sphinxAtStartPar
lightgreen = 144, 238, 144
&
\sphinxAtStartPar
lightgrey = 211, 211, 211
&
\sphinxAtStartPar
lightpink = 255, 182, 193
\\
\sphinxhline
\sphinxAtStartPar
lightsalmon = 255, 160, 122
&
\sphinxAtStartPar
lightseagreen = 32, 178, 170
&
\sphinxAtStartPar
lightskyblue = 135, 206, 250
&
\sphinxAtStartPar
lightslategray = 119, 136, 153
&
\sphinxAtStartPar
lightsteelblue = 176, 196, 222
\\
\sphinxhline
\sphinxAtStartPar
lightyellow = 255, 255, 224
&
\sphinxAtStartPar
lime = 0, 255, 0
&
\sphinxAtStartPar
limegreen = 50, 205, 50
&
\sphinxAtStartPar
linen = 250, 240, 230
&
\sphinxAtStartPar
magenta = 255, 0, 255
\\
\sphinxhline
\sphinxAtStartPar
maroon = 128, 0, 0
&
\sphinxAtStartPar
mediumaquamarine = 102, 205, 170
&
\sphinxAtStartPar
mediumblue = 0, 0, 205
&
\sphinxAtStartPar
mediumorchid = 186, 85, 211
&
\sphinxAtStartPar
mediumpurple = 147, 112, 219
\\
\sphinxhline
\sphinxAtStartPar
mediumseagreen = 60, 179, 113
&
\sphinxAtStartPar
mediumslateblue = 123, 104, 238
&
\sphinxAtStartPar
mediumspringgreen = 0, 250, 154
&
\sphinxAtStartPar
mediumturquoise = 72, 209, 204
&
\sphinxAtStartPar
mediumvioletred = 199, 21, 133
\\
\sphinxhline
\sphinxAtStartPar
midnightblue = 25, 25, 112
&
\sphinxAtStartPar
mintcream = 245, 255, 250
&
\sphinxAtStartPar
mistyrose = 255, 228, 225
&
\sphinxAtStartPar
moccasin = 255, 228, 181
&
\sphinxAtStartPar
navajowhite = 255, 222, 173
\\
\sphinxhline
\sphinxAtStartPar
navy = 0, 0, 128
&
\sphinxAtStartPar
oldlace = 253, 245, 230
&
\sphinxAtStartPar
olive = 128, 128, 0
&
\sphinxAtStartPar
olivedrab = 107, 142, 35
&
\sphinxAtStartPar
orange = 255, 165, 0
\\
\sphinxhline
\sphinxAtStartPar
orangered = 255, 69, 0
&
\sphinxAtStartPar
orchid = 218, 112, 214
&
\sphinxAtStartPar
palegoldenrod = 238, 232, 170
&
\sphinxAtStartPar
palegreen = 152, 251, 152
&
\sphinxAtStartPar
paleturquoise = 175, 238, 238
\\
\sphinxhline
\sphinxAtStartPar
palevioletred = 219, 112, 147
&
\sphinxAtStartPar
papayawhip = 255, 239, 213
&
\sphinxAtStartPar
peachpuff = 255, 239, 213
&
\sphinxAtStartPar
peru = 205, 133, 63
&
\sphinxAtStartPar
pink = 255, 192, 203
\\
\sphinxhline
\sphinxAtStartPar
plum = 221, 160, 221
&
\sphinxAtStartPar
powderblue = 176, 224, 230
&
\sphinxAtStartPar
purple = 128, 0, 128
&
\sphinxAtStartPar
red = 255, 0, 0
&
\sphinxAtStartPar
rosybrown = 188, 143, 143
\\
\sphinxhline
\sphinxAtStartPar
royalblue = 65, 105, 225
&
\sphinxAtStartPar
saddlebrown = 139, 69, 19
&
\sphinxAtStartPar
salmon = 250, 128, 114
&
\sphinxAtStartPar
sandybrown = 244, 164, 96
&
\sphinxAtStartPar
seagreen = 46, 139, 87
\\
\sphinxhline
\sphinxAtStartPar
seashell = 255, 245, 238
&
\sphinxAtStartPar
sienna = 160, 82, 45
&
\sphinxAtStartPar
silver = 192, 192, 192
&
\sphinxAtStartPar
skyblue = 135, 206, 235
&
\sphinxAtStartPar
slateblue = 106, 90, 205
\\
\sphinxhline
\sphinxAtStartPar
slategray = 112, 128, 144
&
\sphinxAtStartPar
snow = 255, 250, 250
&
\sphinxAtStartPar
springgreen = 0, 255, 127
&
\sphinxAtStartPar
steelblue = 70, 130, 180
&
\sphinxAtStartPar
tan = 210, 180, 140
\\
\sphinxhline
\sphinxAtStartPar
teal = 0, 128, 128
&
\sphinxAtStartPar
thistle = 216, 191, 216
&
\sphinxAtStartPar
tomato = 253, 99, 71
&
\sphinxAtStartPar
turquoise = 64, 224, 208
&
\sphinxAtStartPar
violet = 238, 130, 238
\\
\sphinxhline
\sphinxAtStartPar
wheat = 245, 222, 179
&
\sphinxAtStartPar
white = 255, 255, 255
&
\sphinxAtStartPar
whitesmoke = 245, 245, 245
&
\sphinxAtStartPar
yellow = 255, 255, 0
&
\sphinxAtStartPar
yellowgreen = 154, 205, 50
\\
\sphinxbottomrule
\end{tabulary}
\sphinxtableafterendhook\par
\sphinxattableend\end{savenotes}

\sphinxstepscope

\index{dump\_modify@\spxentry{dump\_modify}}\ignorespaces 

\section{dump\_modify command}
\label{\detokenize{dump_modify:dump-modify-command}}\label{\detokenize{dump_modify:index-0}}\label{\detokenize{dump_modify::doc}}

\section{dump\_modify command for image/movie options}
\label{\detokenize{dump_modify:dump-modify-command-for-image-movie-options}}

\subsection{Syntax}
\label{\detokenize{dump_modify:syntax}}
\begin{sphinxVerbatim}[commandchars=\\\{\}]
\PYG{k}{dump\PYGZus{}modify }\PYG{n+nv+nvIdentifier}{dump\PYGZhy{}ID}\PYG{+w}{ }\PYG{n}{keyword}\PYG{+w}{ }\PYG{n}{values}\PYG{+w}{ }\PYG{n}{...}
\end{sphinxVerbatim}
\begin{itemize}
\item {} 
\sphinxAtStartPar
dump\sphinxhyphen{}ID = ID of dump to modify

\item {} 
\sphinxAtStartPar
one or more keyword/value pairs may be appended

\item {} 
\sphinxAtStartPar
these keywords apply to various dump styles

\item {} 
\sphinxAtStartPar
keyword = \sphinxstyleemphasis{append} or \sphinxstyleemphasis{at} or \sphinxstyleemphasis{balance} or \sphinxstyleemphasis{buffer} or \sphinxstyleemphasis{colname} or \sphinxstyleemphasis{delay} or \sphinxstyleemphasis{element} or \sphinxstyleemphasis{every} or \sphinxstyleemphasis{every/time} or \sphinxstyleemphasis{fileper} or \sphinxstyleemphasis{first} or \sphinxstyleemphasis{flush} or \sphinxstyleemphasis{format} or \sphinxstyleemphasis{header} or \sphinxstyleemphasis{image} or \sphinxstyleemphasis{label} or \sphinxstyleemphasis{maxfiles} or \sphinxstyleemphasis{nfile} or \sphinxstyleemphasis{pad} or \sphinxstyleemphasis{pbc} or \sphinxstyleemphasis{precision} or \sphinxstyleemphasis{region} or \sphinxstyleemphasis{refresh} or \sphinxstyleemphasis{scale} or \sphinxstyleemphasis{sfactor} or \sphinxstyleemphasis{skip} or \sphinxstyleemphasis{sort} or \sphinxstyleemphasis{tfactor} or \sphinxstyleemphasis{thermo} or \sphinxstyleemphasis{thresh} or \sphinxstyleemphasis{time} or \sphinxstyleemphasis{triclinic/general} or \sphinxstyleemphasis{types} or \sphinxstyleemphasis{units} or \sphinxstyleemphasis{unwrap}
\begin{sphinxalltt}
\sphinxstyleemphasis{append} arg = \sphinxstyleemphasis{yes} or \sphinxstyleemphasis{no}
\sphinxstyleemphasis{at} arg = N
  N = index of frame written upon first dump
\sphinxstyleemphasis{balance} arg = \sphinxstyleemphasis{yes} or \sphinxstyleemphasis{no}
\sphinxstyleemphasis{buffer} arg = \sphinxstyleemphasis{yes} or \sphinxstyleemphasis{no}
\sphinxstyleemphasis{colname} values =  ID string, or \sphinxstyleemphasis{default}
  string = new column header name
  ID = integer from 1 to N, or integer from \sphinxhyphen{}1 to \sphinxhyphen{}N, where N = \# of quantities being output
       \sphinxstyleemphasis{or} a custom dump keyword or reference to compute, fix, property or variable.
\sphinxstyleemphasis{delay} arg = Dstep
  Dstep = delay output until this timestep
\sphinxstyleemphasis{element} args = E1 E2 ... EN, where N = \# of atom types
  E1,...,EN = element name (e.g., C or Fe or Ga)
\sphinxstyleemphasis{every} arg = N
  N = dump on timesteps which are a multiple of N
  N can be a variable (see below)
\sphinxstyleemphasis{every/time} arg = Delta
  Delta = dump once every Delta interval of simulation time (time units)
  Delta can be a variable (see below)
\sphinxstyleemphasis{fileper} arg = Np
  Np = write one file for every this many processors
\sphinxstyleemphasis{first} arg = \sphinxstyleemphasis{yes} or \sphinxstyleemphasis{no}
\sphinxstyleemphasis{flush} arg = \sphinxstyleemphasis{yes} or \sphinxstyleemphasis{no}
\sphinxstyleemphasis{format} args = \sphinxstyleemphasis{line} string, \sphinxstyleemphasis{int} string, \sphinxstyleemphasis{float} string, ID string, or \sphinxstyleemphasis{none}
  string = C\sphinxhyphen{}style format string
  ID = integer from 1 to N, or integer from \sphinxhyphen{}1 to \sphinxhyphen{}N, where N = \# of quantities being output
       \sphinxstyleemphasis{or} a custom dump keyword or reference to compute, fix, property or variable.
\sphinxstyleemphasis{header} arg = \sphinxstyleemphasis{yes} or \sphinxstyleemphasis{no}
  \sphinxstyleemphasis{yes} to write the header
  \sphinxstyleemphasis{no} to not write the header
\sphinxstyleemphasis{image} arg = \sphinxstyleemphasis{yes} or \sphinxstyleemphasis{no}
\sphinxstyleemphasis{label} arg = string
  string = character string (e.g., BONDS) to use in header of dump local file
\sphinxstyleemphasis{maxfiles} arg = Fmax
  Fmax = keep only the most recent \sphinxstyleemphasis{Fmax} snapshots (one snapshot per file)
\sphinxstyleemphasis{nfile} arg = Nf
  Nf = write this many files, one from each of Nf processors
\sphinxstyleemphasis{pad} arg = Nchar = \# of characters to convert timestep to
\sphinxstyleemphasis{pbc} arg = \sphinxstyleemphasis{yes} or \sphinxstyleemphasis{no} = remap atoms via periodic boundary conditions
\sphinxstyleemphasis{precision} arg = power\sphinxhyphen{}of\sphinxhyphen{}10 value from 10 to 1000000
\sphinxstyleemphasis{region} arg = region\sphinxhyphen{}ID or "none"
\sphinxstyleemphasis{refresh} arg = c\_ID = compute ID that supports a refresh operation
\sphinxstyleemphasis{scale} arg = \sphinxstyleemphasis{yes} or \sphinxstyleemphasis{no}
\sphinxstyleemphasis{sfactor} arg = coordinate scaling factor (\textgreater{} 0.0)
\sphinxstyleemphasis{skip} arg = v\_name
  v\_name = variable with name which evaluates to non\sphinxhyphen{}zero (skip) or 0
\sphinxstyleemphasis{sort} arg = \sphinxstyleemphasis{off} or \sphinxstyleemphasis{id} or N or \sphinxhyphen{}N
   off = no sorting of per\sphinxhyphen{}atom lines within a snapshot
   id = sort per\sphinxhyphen{}atom lines by atom ID
   N = sort per\sphinxhyphen{}atom lines in ascending order by the Nth column
   \sphinxhyphen{}N = sort per\sphinxhyphen{}atom lines in descending order by the Nth column
\sphinxstyleemphasis{tfactor} arg = time scaling factor (\textgreater{} 0.0)
\sphinxstyleemphasis{thermo} arg = \sphinxstyleemphasis{yes} or \sphinxstyleemphasis{no}
\sphinxstyleemphasis{thresh} args = attribute operator value
  attribute = same attributes (x,fy,etotal,sxx,etc) used by dump custom style
  operator = "\textless{}" or "\textless{}=" or "\textgreater{}" or "\textgreater{}=" or "==" or "!=" or "|\textasciicircum{}"
  value = numeric value to compare to, or LAST
  these 3 args can be replaced by the word "none" to turn off thresholding
\sphinxstyleemphasis{time} arg = \sphinxstyleemphasis{yes} or \sphinxstyleemphasis{no}
\sphinxstyleemphasis{triclinic/general} arg = \sphinxstyleemphasis{yes} or \sphinxstyleemphasis{no}
\sphinxstyleemphasis{types} value = \sphinxstyleemphasis{numeric} or \sphinxstyleemphasis{labels}
\sphinxstyleemphasis{units} arg = \sphinxstyleemphasis{yes} or \sphinxstyleemphasis{no}
\sphinxstyleemphasis{unwrap} arg = \sphinxstyleemphasis{yes} or \sphinxstyleemphasis{no}
\end{sphinxalltt}

\item {} 
\sphinxAtStartPar
these keywords apply only to the \sphinxstyleemphasis{image} and \sphinxstyleemphasis{movie} {\hyperref[\detokenize{dump_image::doc}]{\sphinxcrossref{\DUrole{doc}{styles}}}}

\item {} 
\sphinxAtStartPar
keyword = \sphinxstyleemphasis{acolor} or \sphinxstyleemphasis{adiam} or \sphinxstyleemphasis{amap} or \sphinxstyleemphasis{backcolor} or \sphinxstyleemphasis{bcolor} or \sphinxstyleemphasis{bdiam} or \sphinxstyleemphasis{boxcolor} or \sphinxstyleemphasis{color} or \sphinxstyleemphasis{bitrate} or \sphinxstyleemphasis{framerate}
\begin{sphinxalltt}
see the {\hyperref[\detokenize{dump_image::doc}]{\sphinxcrossref{\DUrole{doc}{dump image}}}} doc page for details
\end{sphinxalltt}

\item {} 
\sphinxAtStartPar
these keywords apply only to the \sphinxstyleemphasis{/gz} and \sphinxstyleemphasis{/zstd} dump styles

\item {} 
\sphinxAtStartPar
keyword = \sphinxstyleemphasis{compression\_level}
\begin{sphinxalltt}
\sphinxstyleemphasis{compression\_level} args = level
  level = integer specifying the compression level that should be used (see below for supported levels)
\end{sphinxalltt}

\item {} 
\sphinxAtStartPar
these keywords apply only to the \sphinxstyleemphasis{/zstd} dump styles

\item {} 
\sphinxAtStartPar
keyword = \sphinxstyleemphasis{checksum}
\begin{sphinxalltt}
\sphinxstyleemphasis{checksum} args = \sphinxstyleemphasis{yes} or \sphinxstyleemphasis{no} (add checksum at end of zst file)
\end{sphinxalltt}

\item {} 
\sphinxAtStartPar
these keywords apply only to the vtk* dump style

\item {} 
\sphinxAtStartPar
keyword = \sphinxstyleemphasis{binary}
\begin{sphinxalltt}
\sphinxstyleemphasis{binary} args = \sphinxstyleemphasis{yes} or \sphinxstyleemphasis{no} (select between binary and text mode VTK files)
\end{sphinxalltt}

\end{itemize}


\subsection{Examples}
\label{\detokenize{dump_modify:examples}}
\begin{sphinxVerbatim}[commandchars=\\\{\}]
\PYG{k}{dump\PYGZus{}modify }\PYG{n+nv+nvIdentifier}{1}\PYG{+w}{ }\PYG{n}{format}\PYG{+w}{ }\PYG{n}{line}\PYG{+w}{ }\PYG{l+s}{\PYGZdq{}}\PYG{l+s}{\PYGZpc{}d \PYGZpc{}d \PYGZpc{}20.15g \PYGZpc{}g \PYGZpc{}g}\PYG{l+s}{\PYGZdq{}}\PYG{+w}{ }\PYG{n}{scale}\PYG{+w}{ }\PYG{n}{yes}
\PYG{k}{dump\PYGZus{}modify }\PYG{n+nv+nvIdentifier}{1}\PYG{+w}{ }\PYG{n}{format}\PYG{+w}{ }\PYG{n}{float}\PYG{+w}{ }\PYG{o}{\PYGZpc{}}\PYG{l+m}{20.15}\PYG{n}{g}\PYG{+w}{ }\PYG{n}{scale}\PYG{+w}{ }\PYG{n}{yes}
\PYG{k}{dump\PYGZus{}modify }\PYG{n+nv+nvIdentifier}{myDump}\PYG{+w}{ }\PYG{n}{image}\PYG{+w}{ }\PYG{n}{yes}\PYG{+w}{ }\PYG{n}{scale}\PYG{+w}{ }\PYG{n}{no}\PYG{+w}{ }\PYG{n}{flush}\PYG{+w}{ }\PYG{n}{yes}
\PYG{k}{dump\PYGZus{}modify }\PYG{n+nv+nvIdentifier}{1}\PYG{+w}{ }\PYG{k}{region }\PYG{n+nv+nvIdentifier}{mySphere}\PYG{+w}{ }\PYG{n}{thresh}\PYG{+w}{ }\PYG{n}{x}\PYG{+w}{ }\PYG{o}{\PYGZlt{}}\PYG{+w}{ }\PYG{l+m}{0.0}\PYG{+w}{ }\PYG{n}{thresh}\PYG{+w}{ }\PYG{n}{fx}\PYG{+w}{ }\PYG{o}{\PYGZgt{}}\PYG{o}{=}\PYG{+w}{ }\PYG{l+m}{3.2}
\PYG{k}{dump\PYGZus{}modify }\PYG{n+nv+nvIdentifier}{xtcdump}\PYG{+w}{ }\PYG{n}{precision}\PYG{+w}{ }\PYG{l+m}{10000}\PYG{+w}{ }\PYG{n}{sfactor}\PYG{+w}{ }\PYG{l+m}{0.1}
\PYG{k}{dump\PYGZus{}modify }\PYG{n+nv+nvIdentifier}{1}\PYG{+w}{ }\PYG{n}{every}\PYG{+w}{ }\PYG{l+m}{1000}\PYG{+w}{ }\PYG{n}{nfile}\PYG{+w}{ }\PYG{l+m}{20}
\PYG{k}{dump\PYGZus{}modify }\PYG{n+nv+nvIdentifier}{1}\PYG{+w}{ }\PYG{n}{every}\PYG{+w}{ }\PYG{n}{v\PYGZus{}myVar}
\end{sphinxVerbatim}


\subsection{Description}
\label{\detokenize{dump_modify:description}}
\sphinxAtStartPar
Modify the parameters of a previously defined dump command.  Not all
parameters are relevant to all dump styles.

\sphinxAtStartPar
Unless otherwise noted, the following keywords apply to all the
various dump styles, including the {\hyperref[\detokenize{dump_image::doc}]{\sphinxcrossref{\DUrole{doc}{dump image}}}} and
{\hyperref[\detokenize{dump_image::doc}]{\sphinxcrossref{\DUrole{doc}{dump movie}}}} styles.


\bigskip\hrule\bigskip


\sphinxAtStartPar
The \sphinxstyleemphasis{append} keyword applies to all dump styles except \sphinxstyleemphasis{cfg} and \sphinxstyleemphasis{xtc}
and \sphinxstyleemphasis{dcd}.  It also applies only to text output files, not to binary
or gzipped or image/movie files.  If specified as \sphinxstyleemphasis{yes}, then dump
snapshots are appended to the end of an existing dump file.  If
specified as \sphinxstyleemphasis{no}, then a new dump file will be created which will
overwrite an existing file with the same name.


\bigskip\hrule\bigskip


\sphinxAtStartPar
The \sphinxstyleemphasis{at} keyword only applies to the \sphinxstyleemphasis{netcdf} dump style.  It can only
be used if the \sphinxstyleemphasis{append yes} keyword is also used.  The \sphinxstyleemphasis{N} argument is
the index of which frame to append to.  A negative value can be
specified for \sphinxstyleemphasis{N}, which means a frame counted from the end of the
file.  The \sphinxstyleemphasis{at} keyword can only be used if the dump\_modify command is
before the first command that causes dump snapshots to be output
(e.g., a {\hyperref[\detokenize{run::doc}]{\sphinxcrossref{\DUrole{doc}{run}}}} or {\hyperref[\detokenize{minimize::doc}]{\sphinxcrossref{\DUrole{doc}{minimize}}}} command).  Once the
dump file has been opened, this keyword has no further effect.


\bigskip\hrule\bigskip


\sphinxAtStartPar
The \sphinxstyleemphasis{buffer} keyword applies only to dump styles \sphinxstyleemphasis{atom}, \sphinxstyleemphasis{cfg},
\sphinxstyleemphasis{custom}, \sphinxstyleemphasis{local}, and \sphinxstyleemphasis{xyz}.  It also applies only to text output
files, not to binary or gzipped files.  If specified as \sphinxstyleemphasis{yes}, which
is the default, then each processor writes its output into an internal
text buffer, which is then sent to the processor(s) which perform file
writes, and written by those processors(s) as one large chunk of text.
If specified as \sphinxstyleemphasis{no}, each processor sends its per\sphinxhyphen{}atom data in binary
format to the processor(s) which perform file wirtes, and those
processor(s) format and write it line by line into the output file.

\sphinxAtStartPar
The buffering mode is typically faster since each processor does the
relatively expensive task of formatting the output for its own atoms.
However it requires about twice the memory (per processor) for the
extra buffering.


\bigskip\hrule\bigskip


\sphinxAtStartPar
\DUrole{versionmodified,added}{Added in version 4May2022.}

\sphinxAtStartPar
The \sphinxstyleemphasis{colname} keyword can be used to change the default header keyword
for dump styles: \sphinxstyleemphasis{atom}, \sphinxstyleemphasis{custom}, \sphinxstyleemphasis{cfg}, and \sphinxstyleemphasis{local} and their
compressed, ADIOS variants.  The setting for \sphinxstyleemphasis{ID string} replaces the
default text with the provided string.  \sphinxstyleemphasis{ID} can be a positive integer
when it represents the column number counting from the left, a negative
integer when it represents the column number from the right (i.e. \sphinxhyphen{}1 is
the last column/keyword), or a custom dump keyword (or compute, fix,
property, or variable reference) and then it replaces the string for
that specific keyword. For \sphinxstyleemphasis{atom} dump styles only the keywords “id”,
“type”, “x”, “y”, “z”, “ix”, “iy”, “iz” can be accessed via string
regardless of whether scaled or unwrapped coordinates were enabled or
disabled, and it always assumes 8 columns for indexing regardless of
whether image flags are enabled or not.  For dump style \sphinxstyleemphasis{cfg} only
changes to the “auxiliary” keywords (6th or later keyword) will become
visible.

\sphinxAtStartPar
The \sphinxstyleemphasis{colname} keyword can be used multiple times. If multiple \sphinxstyleemphasis{colname}
settings refer to the same keyword, the last setting has precedence.  A
setting of \sphinxstyleemphasis{default} clears all previous settings, reverting all values
to their default names. Using the \sphinxstyleemphasis{scale} or \sphinxstyleemphasis{image} keyword will also
reset all header keywords to their default values.


\bigskip\hrule\bigskip


\sphinxAtStartPar
The \sphinxstyleemphasis{delay} keyword applies to all dump styles.  No snapshots will be
output until the specified \sphinxstyleemphasis{Dstep} timestep or later.  Specifying
\sphinxstyleemphasis{Dstep} \textless{} 0 is the same as turning off the delay setting.  This is a
way to turn off unwanted output early in a simulation, for example,
during an equilibration phase.


\bigskip\hrule\bigskip


\sphinxAtStartPar
The \sphinxstyleemphasis{element} keyword applies only to the dump \sphinxstyleemphasis{cfg}, \sphinxstyleemphasis{xyz}, and
\sphinxstyleemphasis{image} styles.  It associates element names (e.g., H, C, Fe) with
LAMMPS atom types.  See the list of element names at the bottom of
this page.

\sphinxAtStartPar
In the case of dump \sphinxstyleemphasis{cfg}, this allows the \sphinxhref{http://li.mit.edu/Archive/Graphics/A/}{AtomEye}
visualization package to read the dump file and render atoms with the
appropriate size and color.

\sphinxAtStartPar
In the case of dump \sphinxstyleemphasis{image}, the output images will follow the same
\sphinxhref{http://li.mit.edu/Archive/Graphics/A/}{AtomEye} convention.  An element name is specified for each
atom type (1 to Ntype) in the simulation.  The same element name can
be given to multiple atom types.

\sphinxAtStartPar
In the case of \sphinxstyleemphasis{xyz} format dumps, there are no restrictions to what
label can be used as an element name.  Any white\sphinxhyphen{}space separated text
will be accepted.


\bigskip\hrule\bigskip


\sphinxAtStartPar
The \sphinxstyleemphasis{every} keyword can be used with any dump style except the \sphinxstyleemphasis{dcd}
and \sphinxstyleemphasis{xtc} styles.  It specifies that the output of dump snapshots will
now be performed on timesteps which are a multiple of a new \(N\)
value, This overrides the dump frequency originally specified by the
{\hyperref[\detokenize{dump::doc}]{\sphinxcrossref{\DUrole{doc}{dump}}}} command.

\sphinxAtStartPar
The \sphinxstyleemphasis{every} keyword can be specified in one of two ways.  It can be a
numeric value in which case it must be \textgreater{} 0.  Or it can be an
{\hyperref[\detokenize{variable::doc}]{\sphinxcrossref{\DUrole{doc}{equal\sphinxhyphen{}style variable}}}}, which should be specified as
v\_name, where name is the variable name.

\sphinxAtStartPar
In this case, the variable is evaluated at the beginning of a run to
determine the next timestep at which a dump snapshot will be written
out.  On that timestep the variable will be evaluated again to
determine the next timestep, etc.  Thus the variable should return
timestep values.  See the stagger() and logfreq() and stride() math
functions for {\hyperref[\detokenize{variable::doc}]{\sphinxcrossref{\DUrole{doc}{equal\sphinxhyphen{}style variables}}}}, as examples of
useful functions to use in this context.  Other similar math functions
could easily be added as options for {\hyperref[\detokenize{variable::doc}]{\sphinxcrossref{\DUrole{doc}{equal\sphinxhyphen{}style variables}}}}.  Also see the next() function, which allows use of a
file\sphinxhyphen{}style variable which reads successive values from a file, each
time the variable is evaluated.  Used with the \sphinxstyleemphasis{every} keyword, if the
file contains a list of ascending timesteps, you can output snapshots
whenever you wish.

\sphinxAtStartPar
Note that when using the variable option with the \sphinxstyleemphasis{every} keyword, you
need to use the \sphinxstyleemphasis{first} option if you want an initial snapshot written
to the dump file.  The \sphinxstyleemphasis{every} keyword cannot be used with the dump
\sphinxstyleemphasis{dcd} style.

\sphinxAtStartPar
For example, the following commands will
write snapshots at timesteps 0,10,20,30,100,200,300,1000,2000,etc:

\begin{sphinxVerbatim}[commandchars=\\\{\}]
\PYG{k}{variable        }\PYG{n+nv+nvIdentifier}{s}\PYG{+w}{ }\PYG{n}{equal}\PYG{+w}{ }\PYG{n}{logfreq}\PYG{n+nv}{(}\PYG{n+nv}{10,3,10}\PYG{n+nv}{)}
\PYG{k}{dump            }\PYG{n+nv+nvIdentifier}{1}\PYG{+w}{ }\PYG{n+nv+nvIdentifier}{all}\PYG{+w}{ }\PYG{n}{atom}\PYG{+w}{ }\PYG{l+m}{100}\PYG{+w}{ }\PYG{n}{tmp.dump}
\PYG{k}{dump\PYGZus{}modify     }\PYG{n+nv+nvIdentifier}{1}\PYG{+w}{ }\PYG{n}{every}\PYG{+w}{ }\PYG{n}{v\PYGZus{}s}\PYG{+w}{ }\PYG{n}{first}\PYG{+w}{ }\PYG{n}{yes}
\end{sphinxVerbatim}

\sphinxAtStartPar
The following commands would write snapshots at the timesteps listed
in file tmp.times:

\begin{sphinxVerbatim}[commandchars=\\\{\}]
\PYG{k}{variable        }\PYG{n+nv+nvIdentifier}{f}\PYG{+w}{ }\PYG{n}{file}\PYG{+w}{ }\PYG{n}{tmp.times}
\PYG{k}{variable        }\PYG{n+nv+nvIdentifier}{s}\PYG{+w}{ }\PYG{n}{equal}\PYG{+w}{ }\PYG{n}{next}\PYG{n+nv}{(}\PYG{n+nv}{f}\PYG{n+nv}{)}
\PYG{k}{dump            }\PYG{n+nv+nvIdentifier}{1}\PYG{+w}{ }\PYG{n+nv+nvIdentifier}{all}\PYG{+w}{ }\PYG{n}{atom}\PYG{+w}{ }\PYG{l+m}{100}\PYG{+w}{ }\PYG{n}{tmp.dump}
\PYG{k}{dump\PYGZus{}modify     }\PYG{n+nv+nvIdentifier}{1}\PYG{+w}{ }\PYG{n}{every}\PYG{+w}{ }\PYG{n}{v\PYGZus{}s}
\end{sphinxVerbatim}

\begin{sphinxadmonition}{note}{Note:}
\sphinxAtStartPar
When using a file\sphinxhyphen{}style variable with the \sphinxstyleemphasis{every} keyword, the
file of timesteps must list a first timestep that is beyond the
current timestep (e.g., it cannot be 0).  And it must list one or more
timesteps beyond the length of the run you perform.  This is because
the dump command will generate an error if the next timestep it reads
from the file is not a value greater than the current timestep.  Thus
if you wanted output on steps 0,15,100 of a 100\sphinxhyphen{}timestep run, the file
should contain the values 15,100,101 and you should also use the
dump\_modify first command.  Any final value \textgreater{} 100 could be used in
place of 101.
\end{sphinxadmonition}


\bigskip\hrule\bigskip


\sphinxAtStartPar
\DUrole{versionmodified,added}{Added in version 7Jan2022.}

\sphinxAtStartPar
The \sphinxstyleemphasis{every/time} keyword can be used with any dump style except the
\sphinxstyleemphasis{dcd} and \sphinxstyleemphasis{xtc} styles.  It changes the frequency of dump snapshots
from being based on the current timestep to being determined by
elapsed simulation time, i.e. in time units of the {\hyperref[\detokenize{units::doc}]{\sphinxcrossref{\DUrole{doc}{units}}}} command, and specifies \sphinxstyleemphasis{Delta} for the interval between
snapshots.  This can be useful when the timestep size varies during a
simulation run, e.g. by use of the {\hyperref[\detokenize{fix_dt_reset::doc}]{\sphinxcrossref{\DUrole{doc}{fix dt/reset}}}}
command.  The default is to perform output on timesteps which a
multiples of specified timestep value \(N\); see the \sphinxstyleemphasis{every}
keyword.

\sphinxAtStartPar
The \sphinxstyleemphasis{every/time} keyword can be used with any dump style except the
\sphinxstyleemphasis{dcd} and \sphinxstyleemphasis{xtc} styles.  It does two things.  It specifies that the
interval between dump snapshots will be set in simulation time
(i.e. in time units of the {\hyperref[\detokenize{units::doc}]{\sphinxcrossref{\DUrole{doc}{units}}}} command).  This can be
useful when the timestep size varies during a simulation run (e.g., by
use of the {\hyperref[\detokenize{fix_dt_reset::doc}]{\sphinxcrossref{\DUrole{doc}{fix dt/reset}}}} command).  The default is
to specify the interval in timesteps; see the \sphinxstyleemphasis{every} keyword.  The
\sphinxstyleemphasis{every/time} command also sets the interval value.

\begin{sphinxadmonition}{note}{Note:}
\sphinxAtStartPar
If you wish dump styles \sphinxstyleemphasis{atom}, \sphinxstyleemphasis{custom}, \sphinxstyleemphasis{local}, or \sphinxstyleemphasis{xyz} to
include the simulation time as a field in the header portion of
each snapshot, you also need to use the dump\_modify \sphinxstyleemphasis{time} keyword
with a setting of \sphinxstyleemphasis{yes}.  See its documentation below.
\end{sphinxadmonition}

\sphinxAtStartPar
Note that since snapshots are output on simulation steps, each
snapshot will be written on the first timestep whose associated
simulation time is \textgreater{}= the exact snapshot time value.

\sphinxAtStartPar
As with the \sphinxstyleemphasis{every} option, the \sphinxstyleemphasis{Delta} value can be specified in one
of two ways.  It can be a numeric value in which case it must be \textgreater{}
0.0.  Or it can be an {\hyperref[\detokenize{variable::doc}]{\sphinxcrossref{\DUrole{doc}{equal\sphinxhyphen{}style variable}}}}, which
should be specified as v\_name, where name is the variable name.

\sphinxAtStartPar
In this case, the variable is evaluated at the beginning of a run to
determine the next simulation time at which a dump snapshot will be
written out.  On that timestep the variable will be evaluated again to
determine the next simulation time, etc.  Thus the variable should
return values in time units.  Note the current timestep or simulation
time can be used in an {\hyperref[\detokenize{variable::doc}]{\sphinxcrossref{\DUrole{doc}{equal\sphinxhyphen{}style variables}}}} since
they are both thermodynamic keywords.  Also see the next() function,
which allows use of a file\sphinxhyphen{}style variable which reads successive
values from a file, each time the variable is evaluated.  Used with
the \sphinxstyleemphasis{every/time} keyword, if the file contains a list of ascending
simulation times, you can output snapshots whenever you wish.

\sphinxAtStartPar
Note that when using the variable option with the \sphinxstyleemphasis{every/time}
keyword, you need to use the \sphinxstyleemphasis{first} option if you want an initial
snapshot written to the dump file.  The \sphinxstyleemphasis{every/time} keyword cannot be
used with the dump \sphinxstyleemphasis{dcd} style.

\sphinxAtStartPar
For example, the following commands will write snapshots at successive
simulation times which grow by a factor of 1.5 with each interval.
The dt value used in the variable is to avoid a zero result when the
initial simulation time is 0.0.

\begin{sphinxVerbatim}[commandchars=\\\{\}]
\PYG{k}{variable        }\PYG{n+nv+nvIdentifier}{increase}\PYG{+w}{ }\PYG{n}{equal}\PYG{+w}{ }\PYG{l+m}{1.5}\PYG{o}{*}\PYG{n+nv}{(}\PYG{n+nv}{time+dt}\PYG{n+nv}{)}
\PYG{k}{dump            }\PYG{n+nv+nvIdentifier}{1}\PYG{+w}{ }\PYG{n+nv+nvIdentifier}{all}\PYG{+w}{ }\PYG{n}{atom}\PYG{+w}{ }\PYG{l+m}{100}\PYG{+w}{ }\PYG{n}{tmp.dump}
\PYG{k}{dump\PYGZus{}modify     }\PYG{n+nv+nvIdentifier}{1}\PYG{+w}{ }\PYG{n}{every}\PYG{o}{/}\PYG{n}{time}\PYG{+w}{ }\PYG{n}{v\PYGZus{}increase}\PYG{+w}{ }\PYG{n}{first}\PYG{+w}{ }\PYG{n}{yes}
\end{sphinxVerbatim}

\sphinxAtStartPar
The following commands would write snapshots at the times listed in
file tmp.times:

\begin{sphinxVerbatim}[commandchars=\\\{\}]
\PYG{k}{variable        }\PYG{n+nv+nvIdentifier}{f}\PYG{+w}{ }\PYG{n}{file}\PYG{+w}{ }\PYG{n}{tmp.times}
\PYG{k}{variable        }\PYG{n+nv+nvIdentifier}{s}\PYG{+w}{ }\PYG{n}{equal}\PYG{+w}{ }\PYG{n}{next}\PYG{n+nv}{(}\PYG{n+nv}{f}\PYG{n+nv}{)}
\PYG{k}{dump            }\PYG{n+nv+nvIdentifier}{1}\PYG{+w}{ }\PYG{n+nv+nvIdentifier}{all}\PYG{+w}{ }\PYG{n}{atom}\PYG{+w}{ }\PYG{l+m}{100}\PYG{+w}{ }\PYG{n}{tmp.dump}
\PYG{k}{dump\PYGZus{}modify     }\PYG{n+nv+nvIdentifier}{1}\PYG{+w}{ }\PYG{n}{every}\PYG{o}{/}\PYG{n}{time}\PYG{+w}{ }\PYG{n}{v\PYGZus{}s}
\end{sphinxVerbatim}

\begin{sphinxadmonition}{note}{Note:}
\sphinxAtStartPar
When using a file\sphinxhyphen{}style variable with the \sphinxstyleemphasis{every/time} keyword, the
file of timesteps must list a first time that is beyond the time
associated with the current timestep (e.g., it cannot be 0.0).  And
it must list one or more times beyond the length of the run you
perform.  This is because the dump command will generate an error
if the next time it reads from the file is not a value greater than
the current time.  Thus if you wanted output at times 0,15,100 of a
run of length 100 in simulation time, the file should contain the
values 15,100,101 and you should also use the dump\_modify first
command.  Any final value \textgreater{} 100 could be used in place of 101.
\end{sphinxadmonition}


\bigskip\hrule\bigskip


\sphinxAtStartPar
The \sphinxstyleemphasis{first} keyword determines whether a dump snapshot is written on
the very first timestep after the dump command is invoked.  This will
always occur if the current timestep is a multiple of \$N\$, the
frequency specified in the {\hyperref[\detokenize{dump::doc}]{\sphinxcrossref{\DUrole{doc}{dump}}}} command or
{\hyperref[\detokenize{dump_modify::doc}]{\sphinxcrossref{\DUrole{doc}{dump\_modify every}}}} command, including timestep 0.
It will also always occur if the current simulation time is a multiple
of \sphinxstyleemphasis{Delta}, the time interval specified in the {\hyperref[\detokenize{dump_modify::doc}]{\sphinxcrossref{\DUrole{doc}{dump\_modify
every/time}}}} command.

\sphinxAtStartPar
But if this is not the case, a dump snapshot will only be written if
the setting of this keyword is \sphinxstyleemphasis{yes}.  If it is \sphinxstyleemphasis{no}, which is the
default, then it will not be written.

\sphinxAtStartPar
Note that if the argument to the {\hyperref[\detokenize{dump_modify::doc}]{\sphinxcrossref{\DUrole{doc}{dump\_modify every}}}} {\hyperref[\detokenize{dump_modify::doc}]{\sphinxcrossref{\DUrole{doc}{dump\_modify every/time}}}} commands is
a variable and not a numeric value, then specifying \sphinxstyleemphasis{first yes} is the
only way to write a dump snapshot on the first timestep after the dump
command is invoked.


\bigskip\hrule\bigskip


\sphinxAtStartPar
The \sphinxstyleemphasis{flush} keyword determines whether a flush operation is invoked
after a dump snapshot is written to the dump file.  A flush ensures
the output in that file is current (no buffering by the OS), even if
LAMMPS halts before the simulation completes.  Flushes cannot be
performed with dump style \sphinxstyleemphasis{xtc}.


\bigskip\hrule\bigskip


\sphinxAtStartPar
The \sphinxstyleemphasis{format} keyword can be used to change the default numeric format
output by the text\sphinxhyphen{}based dump styles: \sphinxstyleemphasis{atom}, \sphinxstyleemphasis{local}, \sphinxstyleemphasis{custom}, \sphinxstyleemphasis{cfg},
and \sphinxstyleemphasis{xyz} styles. Only the \sphinxstyleemphasis{line} or \sphinxstyleemphasis{none} options can be used with the
\sphinxstyleemphasis{atom} and \sphinxstyleemphasis{xyz} styles.

\sphinxAtStartPar
All the specified format strings are C\sphinxhyphen{}style formats, such as used by
the C/C++ printf() command.  The \sphinxstyleemphasis{line} keyword takes a single
argument which is the format string for an entire line of output for
each atom (do not include a trailing “n”), with \(N\) fields, which you
must enclose in quotes if there is more than one field.  The \sphinxstyleemphasis{int} and
\sphinxstyleemphasis{float} keywords take a single format argument and are applied to all
integer or floating\sphinxhyphen{}point quantities output.  The setting for \sphinxstyleemphasis{M string}
also takes a single format argument which is used for the \(M\)th
value output in each line (e.g., the fifth column is output in high
precision by “format 5 \%20.15g”).

\begin{sphinxadmonition}{note}{Note:}
\sphinxAtStartPar
When using the \sphinxstyleemphasis{line} keyword for the \sphinxstyleemphasis{cfg} style, the first two
fields (atom ID and type) are not actually written into the CFG file,
however you must include formats for them in the format string.
\end{sphinxadmonition}

\sphinxAtStartPar
The \sphinxstyleemphasis{format} keyword can be used multiple times.  The precedence is
that for each value in a line of output, the \sphinxstyleemphasis{M} format (if specified)
is used, else the \sphinxstyleemphasis{int} or \sphinxstyleemphasis{float} setting (if specified) is used,
else the \sphinxstyleemphasis{line} setting (if specified) for that value is used, else
the default setting is used.  A setting of \sphinxstyleemphasis{none} clears all previous
settings, reverting all values to their default format.

\begin{sphinxadmonition}{note}{Note:}
\sphinxAtStartPar
Atom and molecule IDs are stored internally as 4\sphinxhyphen{}byte or 8\sphinxhyphen{}byte
signed integers, depending on how LAMMPS was compiled.  When
specifying the \sphinxstyleemphasis{format int} option you can use a “\%d”\sphinxhyphen{}style format
identifier in the format string and LAMMPS will convert this to the
corresponding 8\sphinxhyphen{}byte form if it is needed when outputting those
values.  However, when specifying the \sphinxstyleemphasis{line} option or \sphinxstyleemphasis{format M
string} option for those values, you should specify a format string
appropriate for an 8\sphinxhyphen{}byte signed integer (e.g., one with “\%ld”) if
LAMMPS was compiled with the \sphinxhyphen{}DLAMMPS\_BIGBIG option for 8\sphinxhyphen{}byte IDs.
\end{sphinxadmonition}

\begin{sphinxadmonition}{note}{Note:}
\sphinxAtStartPar
Any value written to a text\sphinxhyphen{}based dump file that is a per\sphinxhyphen{}atom
quantity calculated by a {\hyperref[\detokenize{compute::doc}]{\sphinxcrossref{\DUrole{doc}{compute}}}} or {\hyperref[\detokenize{fix::doc}]{\sphinxcrossref{\DUrole{doc}{fix}}}} is
stored internally as a floating\sphinxhyphen{}point value.  If the value is actually
an integer and you wish it to appear in the text dump file as a
(large) integer, then you need to use an appropriate format.  For
example, these commands:
\end{sphinxadmonition}

\begin{sphinxVerbatim}[commandchars=\\\{\}]
\PYG{k}{compute     }\PYG{n+nv+nvIdentifier}{1}\PYG{+w}{ }\PYG{n+nv+nvIdentifier}{all}\PYG{+w}{ }\PYG{n}{property}\PYG{o}{/}\PYG{n}{local}\PYG{+w}{ }\PYG{n}{batom1}\PYG{+w}{ }\PYG{n}{batom2}
\PYG{k}{dump        }\PYG{n+nv+nvIdentifier}{1}\PYG{+w}{ }\PYG{n+nv+nvIdentifier}{all}\PYG{+w}{ }\PYG{n}{local}\PYG{+w}{ }\PYG{l+m}{100}\PYG{+w}{ }\PYG{n}{tmp.bonds}\PYG{+w}{ }\PYG{n}{index}\PYG{+w}{ }\PYG{n}{c\PYGZus{}1[1]}\PYG{+w}{ }\PYG{n}{c\PYGZus{}1[2]}
\PYG{k}{dump\PYGZus{}modify }\PYG{n+nv+nvIdentifier}{1}\PYG{+w}{ }\PYG{n}{format}\PYG{+w}{ }\PYG{n}{line}\PYG{+w}{ }\PYG{l+s}{\PYGZdq{}}\PYG{l+s}{\PYGZpc{}d \PYGZpc{}0.0f \PYGZpc{}0.0f}\PYG{l+s}{\PYGZdq{}}
\end{sphinxVerbatim}

\sphinxAtStartPar
will output the two atom IDs for atoms in each bond as integers.  If
the dump\_modify command were omitted, they would appear as
floating\sphinxhyphen{}point values, assuming they were large integers (more than six
digits).  The “index” keyword should use the “\%d” format since it is
not generated by a compute or fix, and is stored internally as an
integer.


\bigskip\hrule\bigskip


\sphinxAtStartPar
The \sphinxstyleemphasis{fileper} keyword is documented below with the \sphinxstyleemphasis{nfile} keyword.


\bigskip\hrule\bigskip


\sphinxAtStartPar
The \sphinxstyleemphasis{header} keyword toggles whether the dump file will include a
header.  Excluding a header will reduce the size of the dump file for
data produced by {\hyperref[\detokenize{pair_tracker::doc}]{\sphinxcrossref{\DUrole{doc}{pair tracker}}}} or
{\hyperref[\detokenize{Howto_bpm::doc}]{\sphinxcrossref{\DUrole{doc}{bpm bond styles}}}} which may not require the
information typically written to the header.


\bigskip\hrule\bigskip


\sphinxAtStartPar
The \sphinxstyleemphasis{image} keyword applies only to the dump \sphinxstyleemphasis{atom} style.  If the
image value is \sphinxstyleemphasis{yes}, three flags are appended to each atom’s coords which
are the absolute box image of the atom in each dimension.  For
example, an \(x\) image flag of \(-2\) with a normalized coord of 0.5
means the atom is in the center of the box, but has passed through the box
boundary twice and is really two box lengths to the left of its
current coordinate.  Note that for dump style \sphinxstyleemphasis{custom} these various
values can be printed in the dump file by using the appropriate atom
attributes in the dump command itself.
Using this keyword will reset all custom header names set with
\sphinxstyleemphasis{dump\_modify colname} to their respective default values.


\bigskip\hrule\bigskip


\sphinxAtStartPar
The \sphinxstyleemphasis{label} keyword applies only to the dump \sphinxstyleemphasis{local} style.
When it writes local information, such as bond or angle topology
to a dump file, it will use the specified \sphinxstyleemphasis{label} to format the header.
By default this includes two lines:

\begin{sphinxVerbatim}[commandchars=\\\{\}]
ITEM: NUMBER OF ENTRIES
ITEM: ENTRIES ...
\end{sphinxVerbatim}

\sphinxAtStartPar
The word “ENTRIES” will be replaced with the string specified
(e.g., BONDS or ANGLES).


\bigskip\hrule\bigskip


\sphinxAtStartPar
The \sphinxstyleemphasis{maxfiles} keyword can only be used when a ‘*’ wildcard is
included in the dump file name (i.e., when writing a new file(s) for
each snapshot).  The specified \sphinxstyleemphasis{Fmax} is how many snapshots will be
kept.  Once this number is reached, the file(s) containing the oldest
snapshot is deleted before a new dump file is written.  If the
specified \(\text{Fmax} \le 0\), then all files are retained.

\sphinxAtStartPar
This can be useful for debugging, especially if you do not know on what
timestep something bad will happen (e.g., when LAMMPS will exit with an
error).  You can dump every time step and limit the number of dump
files produced, even if you run for thousands of steps.


\bigskip\hrule\bigskip


\sphinxAtStartPar
The \sphinxstyleemphasis{nfile} or \sphinxstyleemphasis{fileper} keywords can be used in conjunction with the
“\%” wildcard character in the specified dump file name, for all dump
styles except the \sphinxstyleemphasis{dcd}, \sphinxstyleemphasis{image}, \sphinxstyleemphasis{movie}, \sphinxstyleemphasis{xtc}, and \sphinxstyleemphasis{xyz} styles
(for which “\%” is not allowed).  As explained on the {\hyperref[\detokenize{dump::doc}]{\sphinxcrossref{\DUrole{doc}{dump}}}}
command doc page, the “\%” character causes the dump file to be written
in pieces, one piece for each of \(P\) processors.  By default, \(P\)
is the number of processors the simulation is running on.  The \sphinxstyleemphasis{nfile} or
\sphinxstyleemphasis{fileper} keyword can be used to set \(P\) to a smaller value, which can
be more efficient when running on a large number of processors.

\sphinxAtStartPar
The \sphinxstyleemphasis{nfile} keyword sets \(P\) to the specified \(N_f\) value.
For example, if \(N_f = 4\), and the simulation is running on 100
processors, four files will be written by processors 0, 25, 50, and 75.
Each will collect information from itself and the next 24 processors and write
it to a dump file.

\sphinxAtStartPar
For the \sphinxstyleemphasis{fileper} keyword, the specified value of \(N_p\) means write one
file for every \(N_p\) processors.  For example, if \(N_p = 4\),
every fourth processor (0, 4, 8, 12, etc.) will collect information from itself
and the next three processors and write it to a dump file.


\bigskip\hrule\bigskip


\sphinxAtStartPar
The \sphinxstyleemphasis{pad} keyword only applies when the dump filename is specified
with a wildcard “*” character which becomes the timestep.  If \sphinxstyleemphasis{pad} is
0, which is the default, the timestep is converted into a string of
unpadded length (e.g., 100 or 12000 or 2000000).  When \sphinxstyleemphasis{pad} is
specified with \sphinxstyleemphasis{Nchar} \(>\) 0, the string is padded with leading zeroes
so they are all the same length = \sphinxstyleemphasis{Nchar}.  For example, pad 7 would
yield 0000100, 0012000, 2000000.  This can be useful so that
post\sphinxhyphen{}processing programs can easily read the files in ascending
timestep order.


\bigskip\hrule\bigskip


\sphinxAtStartPar
The \sphinxstyleemphasis{pbc} keyword applies to all the dump styles.  As explained on the
{\hyperref[\detokenize{dump::doc}]{\sphinxcrossref{\DUrole{doc}{dump}}}} doc page, atom coordinates in a dump file may be
slightly outside the simulation box.  This is because periodic
boundary conditions are enforced only on timesteps when neighbor lists
are rebuilt, which will not typically coincide with the timesteps dump
snapshots are written.  If the setting of this keyword is set to
\sphinxstyleemphasis{yes}, then all atoms will be remapped to the periodic box before the
snapshot is written, then restored to their original position.  If it
is set to \sphinxstyleemphasis{no} they will not be.  The \sphinxstyleemphasis{no} setting is the default
because it requires no extra computation.


\bigskip\hrule\bigskip


\sphinxAtStartPar
The \sphinxstyleemphasis{precision} keyword only applies to the dump \sphinxstyleemphasis{xtc} style.  A
specified value of \(N\) means that coordinates are stored to \(1/N\)
nanometer accuracy (e.g., for \(N = 1000\), the coordinates are written to
\(1/1000\) nanometer accuracy).


\bigskip\hrule\bigskip


\sphinxAtStartPar
The \sphinxstyleemphasis{refresh} keyword only applies to the dump \sphinxstyleemphasis{custom}, \sphinxstyleemphasis{cfg},
\sphinxstyleemphasis{image}, and \sphinxstyleemphasis{movie} styles.  It allows an “incremental” dump file to
be written, by refreshing a compute that is used as a threshold for
determining which atoms are included in a dump snapshot.  The
specified \sphinxstyleemphasis{c\_ID} gives the ID of the compute.  It is prefixed by “c\_”
to indicate a compute, which is the only current option.  At some
point, other options may be added (e.g., fixes or variables).

\begin{sphinxadmonition}{note}{Note:}
\sphinxAtStartPar
This keyword can only be specified once for a dump.  Refreshes
of multiple computes cannot yet be performed.
\end{sphinxadmonition}

\sphinxAtStartPar
The definition and motivation of an incremental dump file is as
follows.  Instead of outputting all atoms at each snapshot (with some
associated values), you may only wish to output the subset of atoms
with a value that has changed in some way compared to the value the
last time that atom was output.  In some scenarios this can result in
a dramatically smaller dump file.  If desired, by post\sphinxhyphen{}processing the
sequence of snapshots, the values for all atoms at all timesteps can
be inferred.

\sphinxAtStartPar
A concrete example is a simulation of atom diffusion in a solid,
represented as atoms on a lattice.  Diffusive hops are rare.  Imagine
that when a hop occurs an atom moves more than a distance \sphinxstyleemphasis{Dhop}.  For
any snapshot we only want to output atoms that have hopped since the
last snapshot.  This can be accomplished with something the following
commands:

\begin{sphinxVerbatim}[commandchars=\\\{\}]
\PYG{k}{variable        }\PYG{n+nv+nvIdentifier}{Dhop}\PYG{+w}{ }\PYG{n}{equal}\PYG{+w}{ }\PYG{l+m}{0.6}
\PYG{k}{variable        }\PYG{n+nv+nvIdentifier}{check}\PYG{+w}{ }\PYG{n}{atom}\PYG{+w}{ }\PYG{l+s}{\PYGZdq{}}\PYG{l+s}{c\PYGZus{}dsp[4] \PYGZgt{} v\PYGZus{}Dhop}\PYG{l+s}{\PYGZdq{}}
\PYG{k}{compute         }\PYG{n+nv+nvIdentifier}{dsp}\PYG{+w}{ }\PYG{n+nv+nvIdentifier}{all}\PYG{+w}{ }\PYG{n}{displace}\PYG{o}{/}\PYG{n}{atom}\PYG{+w}{ }\PYG{n}{refresh}\PYG{+w}{ }\PYG{n}{check}
\PYG{k}{dump            }\PYG{n+nv+nvIdentifier}{1}\PYG{+w}{ }\PYG{n+nv+nvIdentifier}{all}\PYG{+w}{ }\PYG{n}{custom}\PYG{+w}{ }\PYG{l+m}{20}\PYG{+w}{ }\PYG{n}{tmp.dump}\PYG{+w}{ }\PYG{n}{id}\PYG{+w}{ }\PYG{n}{type}\PYG{+w}{ }\PYG{n}{x}\PYG{+w}{ }\PYG{n}{y}\PYG{+w}{ }\PYG{n}{z}
\PYG{k}{dump\PYGZus{}modify     }\PYG{n+nv+nvIdentifier}{1}\PYG{+w}{ }\PYG{n}{append}\PYG{+w}{ }\PYG{n}{yes}\PYG{+w}{ }\PYG{n}{thresh}\PYG{+w}{ }\PYG{n}{c\PYGZus{}dsp[4]}\PYG{+w}{ }\PYG{o}{\PYGZgt{}}\PYG{+w}{ }\PYG{n+nv}{\PYGZdl{}\PYGZob{}}\PYG{n+nv}{Dhop}\PYG{n+nv}{\PYGZcb{}}\PYG{+w}{ }\PYG{n}{refresh}\PYG{+w}{ }\PYG{n}{c\PYGZus{}dsp}
\end{sphinxVerbatim}

\sphinxAtStartPar
The {\hyperref[\detokenize{compute_displace_atom::doc}]{\sphinxcrossref{\DUrole{doc}{compute displace/atom}}}} command
calculates the displacement of each atom from its reference position.
The “4” index is the scalar displacement; 1, 2, and 3 are the \(xyz\)
components of the displacement.  The {\hyperref[\detokenize{dump_modify::doc}]{\sphinxcrossref{\DUrole{doc}{dump\_modify thresh}}}}
command will cause only atoms that have displaced more than
\(0.6~\AA\) to be output on a given snapshot (assuming
metal units).  However, note that when an atom is output, we also need to
update the reference position for that atom to its new coordinates.  So that it
will not be output in every snapshot thereafter.  That reference position is
stored by {\hyperref[\detokenize{compute_displace_atom::doc}]{\sphinxcrossref{\DUrole{doc}{compute displace/atom}}}}.  So the
dump\_modify \sphinxstyleemphasis{refresh} option triggers a call to compute displace/atom at the
end of every dump to perform that update.  The \sphinxstyleemphasis{refresh check} option
shown as part of the {\hyperref[\detokenize{compute_displace_atom::doc}]{\sphinxcrossref{\DUrole{doc}{compute displace/atom}}}}
command enables the compute to respond to the call from the dump command, and
update the appropriate reference positions.  This is done be defining an
{\hyperref[\detokenize{variable::doc}]{\sphinxcrossref{\DUrole{doc}{atom\sphinxhyphen{}style variable}}}}, \sphinxstyleemphasis{check} in this example, which
calculates a Boolean value (0 or 1) for each atom, based on the same
criterion used by dump\_modify thresh.

\sphinxAtStartPar
See the {\hyperref[\detokenize{compute_displace_atom::doc}]{\sphinxcrossref{\DUrole{doc}{compute displace/atom}}}} command for
more details, including an example of how to produce output that
includes an initial snapshot with the reference position of all atoms.

\sphinxAtStartPar
Note that only computes with a \sphinxstyleemphasis{refresh} option will work with
dump\_modify refresh.  See individual compute doc pages for details.
Currently, only compute displace/atom supports this option.  Others
may be added at some point.  If you use a compute that does not support
refresh operations, LAMMPS will not complain; dump\_modify refresh will
simply do nothing.


\bigskip\hrule\bigskip


\sphinxAtStartPar
The \sphinxstyleemphasis{region} keyword only applies to the dump \sphinxstyleemphasis{custom}, \sphinxstyleemphasis{cfg},
\sphinxstyleemphasis{image}, and \sphinxstyleemphasis{movie} styles.  If specified, only atoms in the region
will be written to the dump file or included in the image/movie.  Only
one region can be applied as a filter (the last one specified).  See
the {\hyperref[\detokenize{region::doc}]{\sphinxcrossref{\DUrole{doc}{region}}}} command for more details.  Note that a region
can be defined as the “inside” or “outside” of a geometric shape, and
it can be the “union” or “intersection” of a series of simpler
regions.


\bigskip\hrule\bigskip


\sphinxAtStartPar
The \sphinxstyleemphasis{scale} keyword applies only to the dump \sphinxstyleemphasis{atom} style.  A scale
value of \sphinxstyleemphasis{yes} means atom coords are written in normalized units from
0.0 to 1.0 in each box dimension.  If the simulation box is triclinic
(tilted), then all atom coords will still be between 0.0 and 1.0.  A
value of \sphinxstyleemphasis{no} means they are written in absolute distance units
(e.g., \(\AA\) or \(\sigma\)).
Using this keyword will reset all custom header names set with
\sphinxstyleemphasis{dump\_modify colname} to their respective default values.


\bigskip\hrule\bigskip


\sphinxAtStartPar
The \sphinxstyleemphasis{sfactor} and \sphinxstyleemphasis{tfactor} keywords only apply to the dump \sphinxstyleemphasis{xtc}
style.  They allow customization of the unit conversion factors used
when writing to XTC files.  By default, they are initialized for
whatever {\hyperref[\detokenize{units::doc}]{\sphinxcrossref{\DUrole{doc}{units}}}} style is being used, to write out
coordinates in nanometers and time in picoseconds.  For example, for \sphinxstyleemphasis{real}
units, LAMMPS defines \sphinxstyleemphasis{sfactor} = 0.1 and \sphinxstyleemphasis{tfactor} = 0.001, since the
\(\AA\) and fs used by \sphinxstyleemphasis{real} units are 0.1 nm and
0.001 ps, respectively.  If you are using a units system with distance and time
units far from nm and ps, you may wish to write XTC files with
different units, since the compression algorithm used in XTC files is
most effective when the typical magnitude of position data is between
10.0 and 0.1.


\bigskip\hrule\bigskip


\sphinxAtStartPar
\DUrole{versionmodified,added}{Added in version 15Sep2022.}

\sphinxAtStartPar
The \sphinxstyleemphasis{skip} keyword can be used with all dump styles.  It allows a dump
snapshot to be skipped (not written to the dump file), if a condition
is met.  The condition is computed by an {\hyperref[\detokenize{variable::doc}]{\sphinxcrossref{\DUrole{doc}{equal\sphinxhyphen{}style variable}}}}, which should be specified as v\_name, where name is the
variable name.  If the variable evaluation returns a non\sphinxhyphen{}zero value,
then the dump snapshot is skipped.  If it returns zero, the dump
proceeds as usual.  Note that {\hyperref[\detokenize{variable::doc}]{\sphinxcrossref{\DUrole{doc}{equal\sphinxhyphen{}style variable}}}}
can contain Boolean operators which effectively evaluate as a true
(non\sphinxhyphen{}zero) or false (zero) result.

\sphinxAtStartPar
The \sphinxstyleemphasis{skip} keyword can be useful for debugging purposes, e.g. to dump
only on a particular timestep.  Or to limit output to conditions of
interest, e.g. only when the force on some atom exceeds a threshold
value.


\bigskip\hrule\bigskip


\sphinxAtStartPar
The \sphinxstyleemphasis{sort} keyword determines whether lines of per\sphinxhyphen{}atom output in a
snapshot are sorted or not.  A sort value of \sphinxstyleemphasis{off} means they will
typically be written in indeterminate order, either in serial or
parallel.  This is the case even in serial if the {\hyperref[\detokenize{atom_modify::doc}]{\sphinxcrossref{\DUrole{doc}{atom\_modify sort}}}} option is turned on, which it is by default, to improve
performance.  A sort value of \sphinxstyleemphasis{id} means sort the output by atom ID.  A
sort value of \(N\) or \(-N\) means sort the output by the value
in the \(N\)th column of per\sphinxhyphen{}atom info in either ascending or
descending order.

\sphinxAtStartPar
The dump \sphinxstyleemphasis{local} style cannot be sorted by atom ID, since there are
typically multiple lines of output per atom.  Some dump styles, such
as \sphinxstyleemphasis{dcd} and \sphinxstyleemphasis{xtc}, require sorting by atom ID to format the output
file correctly.  If multiple processors are writing the dump file, via
the “\%” wildcard in the dump filename and the \sphinxstyleemphasis{nfile} or \sphinxstyleemphasis{fileper}
keywords are set to non\sphinxhyphen{}default values (i.e., the number of dump file
pieces is not equal to the number of procs), then sorting cannot be
performed.

\sphinxAtStartPar
In a parallel run, the per\sphinxhyphen{}processor dump file pieces can have
significant imbalance in number of lines of per\sphinxhyphen{}atom info. The \sphinxstyleemphasis{balance}
keyword determines whether the number of lines in each processor
snapshot are balanced to be nearly the same. A balance value of \sphinxstyleemphasis{no}
means no balancing will be done, while \sphinxstyleemphasis{yes} means balancing will be
performed. This balancing preserves dump sorting order. For a serial
run, this option is ignored since the output is already balanced.

\begin{sphinxadmonition}{note}{Note:}
\sphinxAtStartPar
Unless it is required by the dump style, sorting dump file
output requires extra overhead in terms of CPU and communication cost,
as well as memory, versus unsorted output.
\end{sphinxadmonition}


\bigskip\hrule\bigskip


\sphinxAtStartPar
The \sphinxstyleemphasis{thermo} keyword only applies the dump styles \sphinxstyleemphasis{netcdf} and \sphinxstyleemphasis{yaml}.
It triggers writing of {\hyperref[\detokenize{thermo::doc}]{\sphinxcrossref{\DUrole{doc}{thermo}}}} information to the dump
file alongside per\sphinxhyphen{}atom data.  The values included in the dump file are
cached values from the last thermo output and include the exact same the
values as specified by the {\hyperref[\detokenize{thermo_style::doc}]{\sphinxcrossref{\DUrole{doc}{thermo\_style}}}} command.
Because these are cached values, they are only up\sphinxhyphen{}to\sphinxhyphen{}date when dump
output is on a timestep that also has thermo output. Dump style \sphinxstyleemphasis{yaml}
will skip thermo output on incompatible steps.


\bigskip\hrule\bigskip


\sphinxAtStartPar
The \sphinxstyleemphasis{thresh} keyword only applies to the dump \sphinxstyleemphasis{custom}, \sphinxstyleemphasis{cfg},
\sphinxstyleemphasis{image}, and \sphinxstyleemphasis{movie} styles.  Multiple thresholds can be specified.
Specifying \sphinxstyleemphasis{none} turns off all threshold criteria.  If thresholds are
specified, only atoms whose attributes meet all the threshold criteria
are written to the dump file or included in the image.  The possible
attributes that can be tested for are the same as those that can be
specified in the {\hyperref[\detokenize{dump::doc}]{\sphinxcrossref{\DUrole{doc}{dump custom}}}} command, with the exception
of the \sphinxstyleemphasis{element} attribute, since it is not a numeric value.  Note
that a different attributes can be used than those output by the
{\hyperref[\detokenize{dump::doc}]{\sphinxcrossref{\DUrole{doc}{dump custom}}}} command.  For example, you can output the
coordinates and stress of atoms whose energy is above some threshold.

\sphinxAtStartPar
If an atom\sphinxhyphen{}style variable is used as the attribute, then it can
produce continuous numeric values or effective Boolean 0/1 values,
which may be useful for the comparison operator.  Boolean values can
be generated by variable formulas that use comparison or Boolean math
operators or special functions like gmask() and rmask() and grmask().
See the {\hyperref[\detokenize{variable::doc}]{\sphinxcrossref{\DUrole{doc}{variable}}}} command page for details.

\sphinxAtStartPar
The specified value must be a simple numeric value or the word LAST.
If LAST is used, it refers to the value of the attribute the last time
the dump command was invoked to produce a snapshot.  This is a way to
only dump atoms whose attribute has changed (or not changed).
Three examples follow.

\begin{sphinxVerbatim}[commandchars=\\\{\}]
\PYG{k}{dump\PYGZus{}modify }\PYG{n+nv+nvIdentifier}{...}\PYG{+w}{ }\PYG{n}{thresh}\PYG{+w}{ }\PYG{n}{ix}\PYG{+w}{ }\PYG{o}{!}\PYG{o}{=}\PYG{+w}{ }\PYG{n}{LAST}
\end{sphinxVerbatim}

\sphinxAtStartPar
This will dump atoms which have crossed the periodic \(x\) boundary of the
simulation box since the last dump.  (Note that atoms that crossed
once and then crossed back between the two dump timesteps would not be
included.)

\begin{sphinxVerbatim}[commandchars=\\\{\}]
\PYG{k}{region }\PYG{n+nv+nvIdentifier}{foo}\PYG{+w}{ }\PYG{n}{sphere}\PYG{+w}{ }\PYG{l+m}{10}\PYG{+w}{ }\PYG{l+m}{20}\PYG{+w}{ }\PYG{l+m}{10}\PYG{+w}{ }\PYG{l+m}{15}
\PYG{k}{variable }\PYG{n+nv+nvIdentifier}{inregion}\PYG{+w}{ }\PYG{n}{atom}\PYG{+w}{ }\PYG{n}{rmask}\PYG{n+nv}{(}\PYG{n+nv}{foo}\PYG{n+nv}{)}
\PYG{k}{dump\PYGZus{}modify }\PYG{n+nv+nvIdentifier}{...}\PYG{+w}{ }\PYG{n}{thresh}\PYG{+w}{ }\PYG{n}{v\PYGZus{}inregion}\PYG{+w}{ }\PYG{o}{|}\PYG{o}{\PYGZca{}}\PYG{+w}{ }\PYG{n}{LAST}
\end{sphinxVerbatim}

\sphinxAtStartPar
This will dump atoms which crossed the boundary of the spherical
region since the last dump.

\begin{sphinxVerbatim}[commandchars=\\\{\}]
\PYG{k}{variable }\PYG{n+nv+nvIdentifier}{charge}\PYG{+w}{ }\PYG{n}{atom}\PYG{+w}{ }\PYG{l+s}{\PYGZdq{}}\PYG{l+s}{(q \PYGZgt{} 0.5) || (q \PYGZlt{} \PYGZhy{}0.5)}\PYG{l+s}{\PYGZdq{}}
\PYG{k}{dump\PYGZus{}modify }\PYG{n+nv+nvIdentifier}{...}\PYG{+w}{ }\PYG{n}{thresh}\PYG{+w}{ }\PYG{n}{v\PYGZus{}charge}\PYG{+w}{ }\PYG{o}{|}\PYG{o}{\PYGZca{}}\PYG{+w}{ }\PYG{n}{LAST}
\end{sphinxVerbatim}

\sphinxAtStartPar
This will dump atoms whose charge has changed from an absolute value
less than \(\frac12\) to greater than \(\frac12\) (or vice
versa) since the last dump (e.g., due to reactions and subsequent
charge equilibration in a reactive force field).

\sphinxAtStartPar
The choice of operators listed above are the usual comparison
operators.  The XOR operation (exclusive or) is also included as “|\textasciicircum{}”.
In this context, XOR means that if either the attribute or value is
0.0 and the other is non\sphinxhyphen{}zero, then the result is “true” and the
threshold criterion is met.  Otherwise it is not met.

\begin{sphinxadmonition}{note}{Note:}
\sphinxAtStartPar
For style \sphinxstyleemphasis{custom}, the \sphinxstyleemphasis{triclinic/general} keyword can alter dump
output for general triclinic simulation boxes and their atoms.  See
the {\hyperref[\detokenize{dump::doc}]{\sphinxcrossref{\DUrole{doc}{dump}}}} command for details of how this changes the
format of dump file snapshots.  The thresh keyword may access
per\sphinxhyphen{}atom attributes either directly or indirectly through a compute
or variable.  If the attribute is an atom coordinate or a per\sphinxhyphen{}atom
vector (such as velocity, force, or dipole moment), its value will
\sphinxstyleemphasis{NOT} be a general triclinic (rotated) value.  Rather it will be a
restricted triclinic value.
\end{sphinxadmonition}


\bigskip\hrule\bigskip


\sphinxAtStartPar
The \sphinxstyleemphasis{time} keyword only applies to the dump \sphinxstyleemphasis{atom}, \sphinxstyleemphasis{custom}, \sphinxstyleemphasis{local},
and \sphinxstyleemphasis{xyz} styles (and their COMPRESS package versions \sphinxstyleemphasis{atom/gz},
\sphinxstyleemphasis{custom/gz} and \sphinxstyleemphasis{local/gz}).  For the first three styles, if set to
\sphinxstyleemphasis{yes}, each frame will will contain two extra lines before the “ITEM:
TIMESTEP” entry:
\begin{sphinxalltt}
ITEM: TIME
\textless{}elapsed time\textgreater{}
\end{sphinxalltt}

\sphinxAtStartPar
For the \sphinxstyleemphasis{xyz} style, the simulation time is included on the same line
as the timestep value.

\sphinxAtStartPar
This will output the current elapsed simulation time in current
time units equivalent to the {\hyperref[\detokenize{thermo_style::doc}]{\sphinxcrossref{\DUrole{doc}{thermo keyword}}}} \sphinxstyleemphasis{time}.
This is to simplify post\sphinxhyphen{}processing of trajectories using a variable time
step (e.g., when using {\hyperref[\detokenize{fix_dt_reset::doc}]{\sphinxcrossref{\DUrole{doc}{fix dt/reset}}}}).
The default setting is \sphinxstyleemphasis{no}.


\bigskip\hrule\bigskip


\sphinxAtStartPar
The \sphinxstyleemphasis{types} keyword applies only to the dump xyz style. If this keyword is
used with a value of \sphinxstyleemphasis{numeric}, then numeric atom types are printed in the
xyz file (default). If the value \sphinxstyleemphasis{labels} is specified, then
{\hyperref[\detokenize{Howto_type_labels::doc}]{\sphinxcrossref{\DUrole{doc}{type labels}}}} are printed for atom types.


\bigskip\hrule\bigskip


\sphinxAtStartPar
The \sphinxstyleemphasis{triclinic/general} keyword only applies to the dump \sphinxstyleemphasis{atom} and
\sphinxstyleemphasis{custom} styles.  It can only be used with a value of \sphinxstyleemphasis{yes} if the
simulation box was created as a general triclinic box.  See the
{\hyperref[\detokenize{Howto_triclinic::doc}]{\sphinxcrossref{\DUrole{doc}{Howto\_triclinic}}}} doc page for a detailed
explanation of orthogonal, restricted triclinic, and general triclinic
simulation boxes.

\sphinxAtStartPar
If this keyword is used with a value of \sphinxstyleemphasis{yes}, the box information at
the beginning of each snapshot will include information about the 3
arbitrary edge vectors \sphinxstylestrong{A}, \sphinxstylestrong{B}, \sphinxstylestrong{C} that define the general
triclinic box as well as their origin.  The format is described on the
{\hyperref[\detokenize{dump::doc}]{\sphinxcrossref{\DUrole{doc}{dump}}}} doc page.

\sphinxAtStartPar
The coordinates of each atom will likewise be output as values in (or
near) the general triclinic box.  Likewise, per\sphinxhyphen{}atom vector quantities
such as velocity, omega, dipole moment, etc will have orientations
consistent with the general triclinic box, meaning they will be
rotated relative to the standard xyz coordinate axes.  See the
{\hyperref[\detokenize{dump::doc}]{\sphinxcrossref{\DUrole{doc}{dump}}}} doc page for a full list of which dump attributes
this affects.


\bigskip\hrule\bigskip


\sphinxAtStartPar
The \sphinxstyleemphasis{units} keyword only applies to the dump \sphinxstyleemphasis{atom}, \sphinxstyleemphasis{custom}, and
\sphinxstyleemphasis{local} styles (and their COMPRESS package versions \sphinxstyleemphasis{atom/gz},
\sphinxstyleemphasis{custom/gz} and \sphinxstyleemphasis{local/gz}). If set to \sphinxstyleemphasis{yes}, each individual dump
file will contain two extra lines at the very beginning with:
\begin{sphinxalltt}
ITEM: UNITS
\textless{}units style\textgreater{}
\end{sphinxalltt}

\sphinxAtStartPar
This will output the current selected {\hyperref[\detokenize{units::doc}]{\sphinxcrossref{\DUrole{doc}{units}}}} style
to the dump file and thus allows visualization and post\sphinxhyphen{}processing
tools to determine the choice of units of the data in the dump file.
The default setting is \sphinxstyleemphasis{no}.


\bigskip\hrule\bigskip


\sphinxAtStartPar
The \sphinxstyleemphasis{unwrap} keyword only applies to the dump \sphinxstyleemphasis{dcd} and \sphinxstyleemphasis{xtc} styles.
If set to \sphinxstyleemphasis{yes}, coordinates will be written “unwrapped” by the image
flags for each atom.  Unwrapped means that if the atom has passed through
a periodic boundary one or more times, the value is printed for what
the coordinate would be if it had not been wrapped back into the
periodic box.  Note that these coordinates may thus be far outside the
box size stored with the snapshot.


\bigskip\hrule\bigskip


\sphinxAtStartPar
The {\hyperref[\detokenize{Packages_details:pkg-compress}]{\sphinxcrossref{\DUrole{std,std-ref}{COMPRESS package}}}} offers both GZ and Zstd
compression variants of styles atom, custom, local, cfg, and xyz. When
using these styles the compression level can be controlled by the
\sphinxcode{\sphinxupquote{compression\_level}} keyword. File names with these styles have to
end in either \sphinxcode{\sphinxupquote{.gz}} or \sphinxcode{\sphinxupquote{.zst}}.

\sphinxAtStartPar
GZ supports compression levels from \(-1\) (default), 0 (no compression),
and 1 to 9, 9 being the best compression. The COMPRESS \sphinxcode{\sphinxupquote{/gz}} styles use 9
as default compression level.

\sphinxAtStartPar
Zstd offers a wider range of compression levels, including negative
levels that sacrifice compression for performance. 0 is the default,
positive levels are 1 to 22, with 22 being the most expensive
compression. Zstd promises higher compression/decompression speeds for
similar compression ratios. For more details see
\sphinxtitleref{https://facebook.github.io/zstd/}.

\sphinxAtStartPar
In addition, Zstd compressed files can include a checksum of the
entire contents. The Zstd enabled dump styles enable this feature by
default and it can be disabled with the \sphinxcode{\sphinxupquote{checksum}} keyword.


\bigskip\hrule\bigskip


\sphinxAtStartPar
The {\hyperref[\detokenize{Packages_details:pkg-vtk}]{\sphinxcrossref{\DUrole{std,std-ref}{VTK package}}}} offers writing dump files in \sphinxhref{https://www.vtk.org/}{VTK file
formats} that can be read by a variety of
visualization tools based on the VTK library.  These VTK files follow
naming conventions that collide with the LAMMPS convention to append
“.bin” to a file name in order to switch to a binary output.  Thus for
{\hyperref[\detokenize{dump_vtk::doc}]{\sphinxcrossref{\DUrole{doc}{vtk style dumps}}}} the dump\_modify command supports the
keyword \sphinxstyleemphasis{binary} which selects between generating text mode and binary
style VTK files.


\bigskip\hrule\bigskip



\subsection{Restrictions}
\label{\detokenize{dump_modify:restrictions}}
\sphinxAtStartPar
Not all \sphinxstyleemphasis{dump\_modify} options can be applied to all dump styles.
Details are in the discussions of the individual options.


\subsection{Related commands}
\label{\detokenize{dump_modify:related-commands}}
\sphinxAtStartPar
{\hyperref[\detokenize{dump::doc}]{\sphinxcrossref{\DUrole{doc}{dump}}}}, {\hyperref[\detokenize{dump_image::doc}]{\sphinxcrossref{\DUrole{doc}{dump image}}}}, {\hyperref[\detokenize{undump::doc}]{\sphinxcrossref{\DUrole{doc}{undump}}}}


\subsection{Default}
\label{\detokenize{dump_modify:default}}
\sphinxAtStartPar
The option defaults are
\begin{itemize}
\item {} 
\sphinxAtStartPar
append = no

\item {} 
\sphinxAtStartPar
balance = no

\item {} 
\sphinxAtStartPar
buffer = yes for dump styles \sphinxstyleemphasis{atom}, \sphinxstyleemphasis{custom}, \sphinxstyleemphasis{loca}, and \sphinxstyleemphasis{xyz}

\item {} 
\sphinxAtStartPar
element = “C” for every atom type

\item {} 
\sphinxAtStartPar
every = whatever it was set to via the {\hyperref[\detokenize{dump::doc}]{\sphinxcrossref{\DUrole{doc}{dump}}}} command

\item {} 
\sphinxAtStartPar
fileper = \# of processors

\item {} 
\sphinxAtStartPar
first = no

\item {} 
\sphinxAtStartPar
flush = yes

\item {} 
\sphinxAtStartPar
format = \%d and \%g for each integer or floating point value

\item {} 
\sphinxAtStartPar
image = no

\item {} 
\sphinxAtStartPar
label = ENTRIES

\item {} 
\sphinxAtStartPar
maxfiles = \sphinxhyphen{}1

\item {} 
\sphinxAtStartPar
nfile = 1

\item {} 
\sphinxAtStartPar
pad = 0

\item {} 
\sphinxAtStartPar
pbc = no

\item {} 
\sphinxAtStartPar
precision = 1000

\item {} 
\sphinxAtStartPar
region = none

\item {} 
\sphinxAtStartPar
scale = yes

\item {} 
\sphinxAtStartPar
sort = off for dump styles \sphinxstyleemphasis{atom}, \sphinxstyleemphasis{custom}, \sphinxstyleemphasis{cfg}, and \sphinxstyleemphasis{local}

\item {} 
\sphinxAtStartPar
sort = id for dump styles \sphinxstyleemphasis{dcd}, \sphinxstyleemphasis{xtc}, and \sphinxstyleemphasis{xyz}

\item {} 
\sphinxAtStartPar
thresh = none

\item {} 
\sphinxAtStartPar
time = no

\item {} 
\sphinxAtStartPar
triclinic/general = no

\item {} 
\sphinxAtStartPar
types = numeric

\item {} 
\sphinxAtStartPar
units = no

\item {} 
\sphinxAtStartPar
unwrap = no

\item {} 
\sphinxAtStartPar
compression\_level = 9 (gz variants)

\item {} 
\sphinxAtStartPar
compression\_level = 0 (zstd variants)

\item {} 
\sphinxAtStartPar
checksum = yes (zstd variants)

\end{itemize}

\sphinxstepscope

\index{dump molfile@\spxentry{dump molfile}}\ignorespaces 

\section{dump molfile command}
\label{\detokenize{dump_molfile:dump-molfile-command}}\label{\detokenize{dump_molfile:index-0}}\label{\detokenize{dump_molfile::doc}}

\subsection{Syntax}
\label{\detokenize{dump_molfile:syntax}}
\begin{sphinxVerbatim}[commandchars=\\\{\}]
\PYG{k}{dump }\PYG{n+nv+nvIdentifier}{ID}\PYG{+w}{ }\PYG{n+nv+nvIdentifier}{group\PYGZhy{}ID}\PYG{+w}{ }\PYG{n}{molfile}\PYG{+w}{ }\PYG{n}{N}\PYG{+w}{ }\PYG{n}{file}\PYG{+w}{ }\PYG{n}{format}\PYG{+w}{ }\PYG{n}{path}
\end{sphinxVerbatim}
\begin{itemize}
\item {} 
\sphinxAtStartPar
ID = user\sphinxhyphen{}assigned name for the dump

\item {} 
\sphinxAtStartPar
group\sphinxhyphen{}ID = ID of the group of atoms to be imaged

\item {} 
\sphinxAtStartPar
molfile = style of dump command (other styles \sphinxstyleemphasis{atom} or \sphinxstyleemphasis{cfg} or \sphinxstyleemphasis{dcd} or \sphinxstyleemphasis{xtc} or \sphinxstyleemphasis{xyz} or \sphinxstyleemphasis{local} or \sphinxstyleemphasis{custom} are discussed on the {\hyperref[\detokenize{dump::doc}]{\sphinxcrossref{\DUrole{doc}{dump}}}} doc page)

\item {} 
\sphinxAtStartPar
N = dump every this many timesteps

\item {} 
\sphinxAtStartPar
file = name of file to write to

\item {} 
\sphinxAtStartPar
format = file format to be used

\item {} 
\sphinxAtStartPar
path = file path with plugins (optional)

\end{itemize}


\subsection{Examples}
\label{\detokenize{dump_molfile:examples}}
\begin{sphinxVerbatim}[commandchars=\\\{\}]
\PYG{k}{dump }\PYG{n+nv+nvIdentifier}{mf1}\PYG{+w}{ }\PYG{n+nv+nvIdentifier}{all}\PYG{+w}{ }\PYG{n}{molfile}\PYG{+w}{ }\PYG{l+m}{10}\PYG{+w}{ }\PYG{n}{melt1.xml}\PYG{+w}{ }\PYG{n}{hoomd}
\PYG{k}{dump }\PYG{n+nv+nvIdentifier}{mf2}\PYG{+w}{ }\PYG{n+nv+nvIdentifier}{all}\PYG{+w}{ }\PYG{n}{molfile}\PYG{+w}{ }\PYG{l+m}{10}\PYG{+w}{ }\PYG{n}{melt2}\PYG{o}{\PYGZhy{}}\PYG{o}{*}\PYG{n}{.pdb}\PYG{+w}{ }\PYG{n}{pdb}\PYG{+w}{ }\PYG{n}{.}
\PYG{k}{dump }\PYG{n+nv+nvIdentifier}{mf3}\PYG{+w}{ }\PYG{n+nv+nvIdentifier}{all}\PYG{+w}{ }\PYG{n}{molfile}\PYG{+w}{ }\PYG{l+m}{50}\PYG{+w}{ }\PYG{n}{melt3.xyz}\PYG{+w}{ }\PYG{n}{xyz}\PYG{+w}{ }\PYG{n}{.}:/home/akohlmey/vmd/plugins/LINUX/molfile
\end{sphinxVerbatim}


\subsection{Description}
\label{\detokenize{dump_molfile:description}}
\sphinxAtStartPar
Dump a snapshot of atom coordinates and selected additional quantities
to one or more files every N timesteps in one of several formats.
Only information for atoms in the specified group is dumped.  This
specific dump style uses molfile plugins that are bundled with the
\sphinxhref{https://www.ks.uiuc.edu/Research/vmd}{VMD} molecular visualization and
analysis program.

\sphinxAtStartPar
Unless the filename contains a * character, the output will be written
to one single file with the specified format. Otherwise there will be
one file per snapshot and the * will be replaced by the time step number
when the snapshot is written.

\begin{sphinxadmonition}{note}{Note:}
\sphinxAtStartPar
Because periodic boundary conditions are enforced only on
timesteps when neighbor lists are rebuilt, the coordinates of an atom
written to a dump file may be slightly outside the simulation box.
\end{sphinxadmonition}

\sphinxAtStartPar
The molfile plugin API has a few restrictions that have to be honored
by this dump style: the number of atoms must not change, the atoms
must be sorted, outside of the coordinates no change in atom properties
(like type, mass, charge) will be recorded.


\bigskip\hrule\bigskip


\sphinxAtStartPar
The \sphinxstyleemphasis{format} keyword determines what format is used to write out the
dump. For this to work, LAMMPS must be able to find and load a
compatible molfile plugin that supports this format.  Settings made via
the {\hyperref[\detokenize{dump_modify::doc}]{\sphinxcrossref{\DUrole{doc}{dump\_modify}}}} command can alter per atom properties
like element names.

\sphinxAtStartPar
The \sphinxstyleemphasis{path} keyword determines which in directories. This is a “path”
like other search paths, i.e. it can contain multiple directories
separated by a colon (or semicolon on Windows). This keyword is
optional and default to “.”, the current directory.

\sphinxAtStartPar
The \sphinxstyleemphasis{unwrap} option of the {\hyperref[\detokenize{dump_modify::doc}]{\sphinxcrossref{\DUrole{doc}{dump\_modify}}}} command allows
coordinates to be written “unwrapped” by the image flags for each atom.
Unwrapped means that if the atom has passed through a periodic boundary
one or more times, the value is printed for what the coordinate would be
if it had not been wrapped back into the periodic box.  Note that these
coordinates may thus be far outside the box size stored with the
snapshot.


\bigskip\hrule\bigskip


\sphinxAtStartPar
Dumps are performed on timesteps that are a multiple of N (including
timestep 0) and on the last timestep of a minimization if the
minimization converges.  Note that this means a dump will not be
performed on the initial timestep after the dump command is invoked,
if the current timestep is not a multiple of N.  This behavior can be
changed via the {\hyperref[\detokenize{dump_modify::doc}]{\sphinxcrossref{\DUrole{doc}{dump\_modify first}}}} command, which can
be useful if the dump command is invoked after a minimization ended on
an arbitrary timestep.  N can be changed between runs by using the
{\hyperref[\detokenize{dump_modify::doc}]{\sphinxcrossref{\DUrole{doc}{dump\_modify every}}}} command. The {\hyperref[\detokenize{dump_modify::doc}]{\sphinxcrossref{\DUrole{doc}{dump\_modify every}}}} command also allows a variable to be used to
determine the sequence of timesteps on which dump files are written.


\bigskip\hrule\bigskip



\subsection{Restrictions}
\label{\detokenize{dump_molfile:restrictions}}
\sphinxAtStartPar
The \sphinxstyleemphasis{molfile} dump style is part of the MOLFILE package.  It is
only enabled if LAMMPS was built with that package.  See the {\hyperref[\detokenize{Build_package::doc}]{\sphinxcrossref{\DUrole{doc}{Build package}}}} page for more info.

\sphinxAtStartPar
Molfile plugins provide a consistent programming interface to read and
write file formats commonly used in molecular simulations. The
MOLFILE package only provides the interface code, not the plugins.
These can be obtained from a VMD installation which has to match the
platform that you are using to compile LAMMPS for. By adding plugins
to VMD, support for new file formats can be added to LAMMPS (or VMD
or other programs that use them) without having to re\sphinxhyphen{}compile the
application itself.  The plugins are installed in the directory:
\textless{}VMDHOME\textgreater{}/plugins/\textless{}VMDARCH\textgreater{}/molfile

\begin{sphinxadmonition}{note}{Note:}
\sphinxAtStartPar
while the programming interface (API) to the plugins is backward
compatible, the binary interface (ABI) has been changing over time, so
it is necessary to compile this package with the plugin header files
from VMD that match the binary plugins.  These header files in the
directory: \textless{}VMDHOME\textgreater{}/plugins/include For convenience, the package ships
with a set of header files that are compatible with VMD 1.9 and 1.9.1
(June 2012)
\end{sphinxadmonition}


\bigskip\hrule\bigskip



\subsection{Related commands}
\label{\detokenize{dump_molfile:related-commands}}
\sphinxAtStartPar
{\hyperref[\detokenize{dump::doc}]{\sphinxcrossref{\DUrole{doc}{dump}}}}, {\hyperref[\detokenize{dump_modify::doc}]{\sphinxcrossref{\DUrole{doc}{dump\_modify}}}}, {\hyperref[\detokenize{undump::doc}]{\sphinxcrossref{\DUrole{doc}{undump}}}}


\subsection{Default}
\label{\detokenize{dump_molfile:default}}
\sphinxAtStartPar
The default path is “.”. All other properties have to be specified.

\sphinxstepscope

\index{dump netcdf@\spxentry{dump netcdf}}\ignorespaces 
\index{dump netcdf/mpiio@\spxentry{dump netcdf/mpiio}}\ignorespaces 

\section{dump netcdf command}
\label{\detokenize{dump_netcdf:dump-netcdf-command}}\label{\detokenize{dump_netcdf:index-1}}\label{\detokenize{dump_netcdf:index-0}}\label{\detokenize{dump_netcdf::doc}}

\section{dump netcdf/mpiio command}
\label{\detokenize{dump_netcdf:dump-netcdf-mpiio-command}}

\subsection{Syntax}
\label{\detokenize{dump_netcdf:syntax}}
\begin{sphinxVerbatim}[commandchars=\\\{\}]
\PYG{k}{dump }\PYG{n+nv+nvIdentifier}{ID}\PYG{+w}{ }\PYG{n+nv+nvIdentifier}{group\PYGZhy{}ID}\PYG{+w}{ }\PYG{n}{netcdf}\PYG{+w}{ }\PYG{n}{N}\PYG{+w}{ }\PYG{n}{file}\PYG{+w}{ }\PYG{n}{args}
\PYG{k}{dump }\PYG{n+nv+nvIdentifier}{ID}\PYG{+w}{ }\PYG{n+nv+nvIdentifier}{group\PYGZhy{}ID}\PYG{+w}{ }\PYG{n}{netcdf}\PYG{o}{/}\PYG{n}{mpiio}\PYG{+w}{ }\PYG{n}{N}\PYG{+w}{ }\PYG{n}{file}\PYG{+w}{ }\PYG{n}{args}
\end{sphinxVerbatim}
\begin{itemize}
\item {} 
\sphinxAtStartPar
ID = user\sphinxhyphen{}assigned name for the dump

\item {} 
\sphinxAtStartPar
group\sphinxhyphen{}ID = ID of the group of atoms to be imaged

\item {} 
\sphinxAtStartPar
\sphinxstyleemphasis{netcdf} or \sphinxstyleemphasis{netcdf/mpiio}  = style of dump command (other styles \sphinxstyleemphasis{atom} or \sphinxstyleemphasis{cfg} or \sphinxstyleemphasis{dcd} or \sphinxstyleemphasis{xtc} or \sphinxstyleemphasis{xyz} or \sphinxstyleemphasis{local} or \sphinxstyleemphasis{custom} are discussed on the {\hyperref[\detokenize{dump::doc}]{\sphinxcrossref{\DUrole{doc}{dump}}}} doc page)

\item {} 
\sphinxAtStartPar
N = dump every this many timesteps

\item {} 
\sphinxAtStartPar
file = name of file to write dump info to

\item {} 
\sphinxAtStartPar
args = list of atom attributes, same as for {\hyperref[\detokenize{dump::doc}]{\sphinxcrossref{\DUrole{doc}{dump\_style custom}}}}

\end{itemize}


\subsection{Examples}
\label{\detokenize{dump_netcdf:examples}}
\begin{sphinxVerbatim}[commandchars=\\\{\}]
\PYG{k}{dump }\PYG{n+nv+nvIdentifier}{1}\PYG{+w}{ }\PYG{n+nv+nvIdentifier}{all}\PYG{+w}{ }\PYG{n}{netcdf}\PYG{+w}{ }\PYG{l+m}{100}\PYG{+w}{ }\PYG{n}{traj.nc}\PYG{+w}{ }\PYG{n}{type}\PYG{+w}{ }\PYG{n}{x}\PYG{+w}{ }\PYG{n}{y}\PYG{+w}{ }\PYG{n}{z}\PYG{+w}{ }\PYG{n}{vx}\PYG{+w}{ }\PYG{n}{vy}\PYG{+w}{ }\PYG{n}{vz}
\PYG{k}{dump\PYGZus{}modify }\PYG{n+nv+nvIdentifier}{1}\PYG{+w}{ }\PYG{n}{append}\PYG{+w}{ }\PYG{n}{yes}\PYG{+w}{ }\PYG{n}{at}\PYG{+w}{ }\PYG{o}{\PYGZhy{}}\PYG{l+m}{1}\PYG{+w}{ }\PYG{n}{thermo}\PYG{+w}{ }\PYG{n}{yes}
\PYG{k}{dump }\PYG{n+nv+nvIdentifier}{1}\PYG{+w}{ }\PYG{n+nv+nvIdentifier}{all}\PYG{+w}{ }\PYG{n}{netcdf}\PYG{o}{/}\PYG{n}{mpiio}\PYG{+w}{ }\PYG{l+m}{1000}\PYG{+w}{ }\PYG{n}{traj.nc}\PYG{+w}{ }\PYG{n}{id}\PYG{+w}{ }\PYG{n}{type}\PYG{+w}{ }\PYG{n}{x}\PYG{+w}{ }\PYG{n}{y}\PYG{+w}{ }\PYG{n}{z}
\PYG{k}{dump }\PYG{n+nv+nvIdentifier}{1}\PYG{+w}{ }\PYG{n+nv+nvIdentifier}{all}\PYG{+w}{ }\PYG{n}{netcdf}\PYG{+w}{ }\PYG{l+m}{1000}\PYG{+w}{ }\PYG{n}{traj.}\PYG{o}{*}\PYG{n}{.nc}\PYG{+w}{ }\PYG{n}{id}\PYG{+w}{ }\PYG{n}{type}\PYG{+w}{ }\PYG{n}{x}\PYG{+w}{ }\PYG{n}{y}\PYG{+w}{ }\PYG{n}{z}
\end{sphinxVerbatim}


\subsection{Description}
\label{\detokenize{dump_netcdf:description}}
\sphinxAtStartPar
Dump a snapshot of atom coordinates every N timesteps in Amber\sphinxhyphen{}style
NetCDF file format.  NetCDF files are binary, portable and
self\sphinxhyphen{}describing.  This dump style will write only one file on the root
node.  The dump style \sphinxstyleemphasis{netcdf} uses the \sphinxhref{https://www.unidata.ucar.edu/software/netcdf/}{standard NetCDF library}.  All data is collected on one processor and then
written to the dump file.  Dump style \sphinxstyleemphasis{netcdf/mpiio} uses the
\sphinxhref{https://trac.mcs.anl.gov/projects/parallel-netcdf/}{parallel NetCDF library} and MPI\sphinxhyphen{}IO to write to the dump
file in parallel; it has better performance on a larger number of
processors.  Note that style \sphinxstyleemphasis{netcdf} outputs all atoms sorted by atom
tag while style \sphinxstyleemphasis{netcdf/mpiio} outputs atoms in order of their MPI
rank.

\sphinxAtStartPar
NetCDF files can be directly visualized via the following tools:
\begin{itemize}
\item {} 
\sphinxAtStartPar
Ovito (\sphinxurl{https://www.ovito.org/}). Ovito supports the AMBER convention and
all extensions of this dump style.

\item {} 
\sphinxAtStartPar
VMD (\sphinxurl{https://www.ks.uiuc.edu/Research/vmd/}).

\end{itemize}

\sphinxAtStartPar
In addition to per\sphinxhyphen{}atom data, {\hyperref[\detokenize{thermo::doc}]{\sphinxcrossref{\DUrole{doc}{thermo}}}} data can be included in the
dump file. The data included in the dump file is identical to the data specified
by {\hyperref[\detokenize{thermo_style::doc}]{\sphinxcrossref{\DUrole{doc}{thermo\_style}}}}.


\bigskip\hrule\bigskip



\subsection{Restrictions}
\label{\detokenize{dump_netcdf:restrictions}}
\sphinxAtStartPar
The \sphinxstyleemphasis{netcdf} and \sphinxstyleemphasis{netcdf/mpiio} dump styles are part of the
NETCDF package.  They are only enabled if LAMMPS was built with
that package. See the {\hyperref[\detokenize{Build_package::doc}]{\sphinxcrossref{\DUrole{doc}{Build package}}}} page for
more info.

\sphinxAtStartPar
The \sphinxstyleemphasis{netcdf} and \sphinxstyleemphasis{netcdf/mpiio} dump styles currently cannot dump
string properties or properties from variables.


\bigskip\hrule\bigskip



\subsection{Related commands}
\label{\detokenize{dump_netcdf:related-commands}}
\sphinxAtStartPar
{\hyperref[\detokenize{dump::doc}]{\sphinxcrossref{\DUrole{doc}{dump}}}}, {\hyperref[\detokenize{dump_modify::doc}]{\sphinxcrossref{\DUrole{doc}{dump\_modify}}}}, {\hyperref[\detokenize{undump::doc}]{\sphinxcrossref{\DUrole{doc}{undump}}}}

\sphinxstepscope

\index{dump vtk@\spxentry{dump vtk}}\ignorespaces 

\section{dump vtk command}
\label{\detokenize{dump_vtk:dump-vtk-command}}\label{\detokenize{dump_vtk:index-0}}\label{\detokenize{dump_vtk::doc}}

\subsection{Syntax}
\label{\detokenize{dump_vtk:syntax}}
\begin{sphinxVerbatim}[commandchars=\\\{\}]
\PYG{k}{dump }\PYG{n+nv+nvIdentifier}{ID}\PYG{+w}{ }\PYG{n+nv+nvIdentifier}{group\PYGZhy{}ID}\PYG{+w}{ }\PYG{n}{vtk}\PYG{+w}{ }\PYG{n}{N}\PYG{+w}{ }\PYG{n}{file}\PYG{+w}{ }\PYG{n}{args}
\end{sphinxVerbatim}
\begin{itemize}
\item {} 
\sphinxAtStartPar
ID = user\sphinxhyphen{}assigned name for the dump

\item {} 
\sphinxAtStartPar
group\sphinxhyphen{}ID = ID of the group of atoms to be dumped

\item {} 
\sphinxAtStartPar
vtk = style of dump command (other styles such as \sphinxstyleemphasis{atom} or \sphinxstyleemphasis{cfg} or \sphinxstyleemphasis{dcd} or \sphinxstyleemphasis{xtc} or \sphinxstyleemphasis{xyz} or \sphinxstyleemphasis{local} or \sphinxstyleemphasis{custom} are discussed on the {\hyperref[\detokenize{dump::doc}]{\sphinxcrossref{\DUrole{doc}{dump}}}} doc page)

\item {} 
\sphinxAtStartPar
N = dump every this many timesteps

\item {} 
\sphinxAtStartPar
file = name of file to write dump info to

\item {} 
\sphinxAtStartPar
args = same as arguments for {\hyperref[\detokenize{dump::doc}]{\sphinxcrossref{\DUrole{doc}{dump\_style custom}}}}

\end{itemize}


\subsection{Examples}
\label{\detokenize{dump_vtk:examples}}
\begin{sphinxVerbatim}[commandchars=\\\{\}]
\PYG{k}{dump }\PYG{n+nv+nvIdentifier}{dmpvtk}\PYG{+w}{ }\PYG{n+nv+nvIdentifier}{all}\PYG{+w}{ }\PYG{n}{vtk}\PYG{+w}{ }\PYG{l+m}{100}\PYG{+w}{ }\PYG{n}{dump}\PYG{o}{*}\PYG{n}{.myforce.vtk}\PYG{+w}{ }\PYG{n}{id}\PYG{+w}{ }\PYG{n}{type}\PYG{+w}{ }\PYG{n}{vx}\PYG{+w}{ }\PYG{n}{fx}
\PYG{k}{dump }\PYG{n+nv+nvIdentifier}{dmpvtp}\PYG{+w}{ }\PYG{n+nv+nvIdentifier}{flow}\PYG{+w}{ }\PYG{n}{vtk}\PYG{+w}{ }\PYG{l+m}{100}\PYG{+w}{ }\PYG{n}{dump}\PYG{o}{*}\PYG{n}{.}\PYG{o}{\PYGZpc{}}\PYG{n}{.displace.vtp}\PYG{+w}{ }\PYG{n}{id}\PYG{+w}{ }\PYG{n}{type}\PYG{+w}{ }\PYG{n}{c\PYGZus{}myD[1]}\PYG{+w}{ }\PYG{n}{c\PYGZus{}myD[2]}\PYG{+w}{ }\PYG{n}{c\PYGZus{}myD[3]}\PYG{+w}{ }\PYG{n}{v\PYGZus{}ke}
\end{sphinxVerbatim}


\subsection{Description}
\label{\detokenize{dump_vtk:description}}
\sphinxAtStartPar
Dump a snapshot of atom quantities to one or more files every \(N\)
timesteps in a format readable by the \sphinxhref{https://www.vtk.org}{VTK visualization toolkit} or other visualization tools that use it, such
as \sphinxhref{https://www.paraview.org}{ParaView}.  The time steps on which dump
output is written can also be controlled by a variable; see the
{\hyperref[\detokenize{dump_modify::doc}]{\sphinxcrossref{\DUrole{doc}{dump\_modify every}}}} command for details.

\sphinxAtStartPar
This dump style is similar to {\hyperref[\detokenize{dump::doc}]{\sphinxcrossref{\DUrole{doc}{dump\_style custom}}}} but uses
the VTK library to write data to VTK simple legacy or XML format,
depending on the filename extension specified for the dump file.  This
can be either \sphinxstyleemphasis{*.vtk} for the legacy format or \sphinxstyleemphasis{*.vtp} and \sphinxstyleemphasis{*.vtu},
respectively, for XML format; see the \sphinxhref{https://www.vtk.org/VTK/img/file-formats.pdf}{VTK homepage} for a detailed
description of these formats.  Since this naming convention conflicts
with the way binary output is usually specified (see below), the
{\hyperref[\detokenize{dump_modify::doc}]{\sphinxcrossref{\DUrole{doc}{dump\_modify binary}}}} command allows setting of a
binary option for this dump style explicitly.

\sphinxAtStartPar
Only information for atoms in the specified group is dumped.  The
{\hyperref[\detokenize{dump_modify::doc}]{\sphinxcrossref{\DUrole{doc}{dump\_modify thresh and region}}}} commands can also
alter what atoms are included; see details below.

\sphinxAtStartPar
As described below, special characters (”*”, “\%”) in the filename
determine the kind of output.

\begin{sphinxadmonition}{warning}{Warning:}
\sphinxAtStartPar
Because periodic boundary conditions are enforced only
on timesteps when neighbor lists are rebuilt, the coordinates of an
atom written to a dump file may be slightly outside the simulation
box.
\end{sphinxadmonition}

\begin{sphinxadmonition}{warning}{Warning:}
\sphinxAtStartPar
Unless the {\hyperref[\detokenize{dump_modify::doc}]{\sphinxcrossref{\DUrole{doc}{dump\_modify sort}}}} option is invoked,
the lines of atom information written to dump files will be in an
indeterminate order for each snapshot.  This is even true when
running on a single processor, if the {\hyperref[\detokenize{atom_modify::doc}]{\sphinxcrossref{\DUrole{doc}{atom\_modify sort}}}} option is on, which it is by default.  In this case
atoms are re\sphinxhyphen{}ordered periodically during a simulation, due to spatial
sorting.  It is also true when running in parallel, because data for
a single snapshot is collected from multiple processors, each of
which owns a subset of the atoms.
\end{sphinxadmonition}

\sphinxAtStartPar
For the \sphinxstyleemphasis{vtk} style, sorting is off by default. See the
{\hyperref[\detokenize{dump_modify::doc}]{\sphinxcrossref{\DUrole{doc}{dump\_modify}}}} page for details.


\bigskip\hrule\bigskip


\sphinxAtStartPar
The dimensions of the simulation box are written to a separate file
for each snapshot (either in legacy VTK or XML format depending on the
format of the main dump file) with the suffix \sphinxstyleemphasis{\_boundingBox} appended
to the given dump filename.

\sphinxAtStartPar
For an orthogonal simulation box this information is saved as a
rectilinear grid (legacy .vtk or .vtr XML format).

\sphinxAtStartPar
Triclinic simulation boxes (non\sphinxhyphen{}orthogonal) are saved as
hexahedrons in either legacy .vtk or .vtu XML format.

\sphinxAtStartPar
Style \sphinxstyleemphasis{vtk} allows you to specify a list of atom attributes to be
written to the dump file for each atom.  The list of possible attributes
is the same as for the {\hyperref[\detokenize{dump::doc}]{\sphinxcrossref{\DUrole{doc}{dump\_style custom}}}} command; see its
documentation page for a listing and an explanation of each attribute.

\begin{sphinxadmonition}{note}{Note:}
\sphinxAtStartPar
Since position data is required to write VTK files the atom
attributes “x y z” do not have to be specified explicitly; they will
be included in the dump file regardless.  Also, in contrast to the
\sphinxstyleemphasis{custom} style, the specified \sphinxstyleemphasis{vtk} attributes are rearranged to
ensure correct ordering of vector components (except for computes and
fixes \sphinxhyphen{} these have to be given in the right order) and duplicate
entries are removed.
\end{sphinxadmonition}

\sphinxAtStartPar
The VTK format uses a single snapshot of the system per file, thus
a wildcard “*” must be included in the filename, as discussed below.
Otherwise the dump files will get overwritten with the new snapshot
each time.


\bigskip\hrule\bigskip


\sphinxAtStartPar
Dumps are performed on timesteps that are a multiple of N (including
timestep 0) and on the last timestep of a minimization if the
minimization converges.  Note that this means a dump will not be
performed on the initial timestep after the dump command is invoked,
if the current timestep is not a multiple of N.  This behavior can be
changed via the {\hyperref[\detokenize{dump_modify::doc}]{\sphinxcrossref{\DUrole{doc}{dump\_modify first}}}} command, which
can also be useful if the dump command is invoked after a minimization
ended on an arbitrary timestep.  N can be changed between runs by
using the {\hyperref[\detokenize{dump_modify::doc}]{\sphinxcrossref{\DUrole{doc}{dump\_modify every}}}} command.
The {\hyperref[\detokenize{dump_modify::doc}]{\sphinxcrossref{\DUrole{doc}{dump\_modify every}}}} command
also allows a variable to be used to determine the sequence of
timesteps on which dump files are written.  In this mode a dump on the
first timestep of a run will also not be written unless the
{\hyperref[\detokenize{dump_modify::doc}]{\sphinxcrossref{\DUrole{doc}{dump\_modify first}}}} command is used.

\sphinxAtStartPar
Dump filenames can contain two wildcard characters.  If a “*”
character appears in the filename, then one file per snapshot is
written and the “*” character is replaced with the timestep value.
For example, tmp.dump*.vtk becomes tmp.dump0.vtk, tmp.dump10000.vtk,
tmp.dump20000.vtk, etc.  Note that the {\hyperref[\detokenize{dump_modify::doc}]{\sphinxcrossref{\DUrole{doc}{dump\_modify pad}}}}
command can be used to ensure all timestep numbers are the same length
(e.g. 00010), which can make it easier to read a series of dump files
in order with some post\sphinxhyphen{}processing tools.

\sphinxAtStartPar
If a “\%” character appears in the filename, then each of P processors
writes a portion of the dump file, and the “\%” character is replaced
with the processor ID from 0 to P\sphinxhyphen{}1 preceded by an underscore character.
For example, tmp.dump\%.vtp becomes tmp.dump\_0.vtp, tmp.dump\_1.vtp, …
tmp.dump\_P\sphinxhyphen{}1.vtp, etc.  This creates smaller files and can be a fast
mode of output on parallel machines that support parallel I/O for output.

\sphinxAtStartPar
By default, P = the number of processors meaning one file per
processor, but P can be set to a smaller value via the \sphinxstyleemphasis{nfile} or
\sphinxstyleemphasis{fileper} keywords of the {\hyperref[\detokenize{dump_modify::doc}]{\sphinxcrossref{\DUrole{doc}{dump\_modify}}}} command.
These options can be the most efficient way of writing out dump files
when running on large numbers of processors.

\sphinxAtStartPar
For the legacy VTK format “\%” is ignored and P = 1, i.e., only
processor 0 does write files.

\sphinxAtStartPar
Note that using the “*” and “\%” characters together can produce a
large number of small dump files!

\sphinxAtStartPar
If \sphinxstyleemphasis{dump\_modify binary} is used, the dump file (or files, if “*” or
“\%” is also used) is written in binary format.  A binary dump file
will be about the same size as a text version, but will typically
write out much faster.


\bigskip\hrule\bigskip



\subsection{Restrictions}
\label{\detokenize{dump_vtk:restrictions}}
\sphinxAtStartPar
The \sphinxstyleemphasis{vtk} style does not support writing of gzipped dump files.

\sphinxAtStartPar
The \sphinxstyleemphasis{vtk} dump style is part of the VTK package. It is only
enabled if LAMMPS was built with that package. See the {\hyperref[\detokenize{Build_package::doc}]{\sphinxcrossref{\DUrole{doc}{Build package}}}} page for more info.

\sphinxAtStartPar
To use this dump style, you also must link to the VTK library.  See
the info in lib/vtk/README and ensure the Makefile.lammps file in that
directory is appropriate for your machine.

\sphinxAtStartPar
The \sphinxstyleemphasis{vtk} dump style supports neither buffering or custom format
strings.


\subsection{Related commands}
\label{\detokenize{dump_vtk:related-commands}}
\sphinxAtStartPar
{\hyperref[\detokenize{dump::doc}]{\sphinxcrossref{\DUrole{doc}{dump}}}}, {\hyperref[\detokenize{dump_image::doc}]{\sphinxcrossref{\DUrole{doc}{dump image}}}},
{\hyperref[\detokenize{dump_modify::doc}]{\sphinxcrossref{\DUrole{doc}{dump\_modify}}}}, {\hyperref[\detokenize{undump::doc}]{\sphinxcrossref{\DUrole{doc}{undump}}}}


\subsection{Default}
\label{\detokenize{dump_vtk:default}}
\sphinxAtStartPar
By default, files are written in ASCII format. If the file extension
is not one of .vtk, .vtp or .vtu, the legacy VTK file format is used.

\sphinxstepscope


\chapter{Bibliography}
\label{\detokenize{Bibliography:bibliography}}\label{\detokenize{Bibliography::doc}}\begin{description}
\sphinxlineitem{\sphinxstylestrong{(Abascal1)}}
\sphinxAtStartPar
Abascal, Sanz, Fernandez, Vega, J Chem Phys, 122, 234511 (2005)

\sphinxlineitem{\sphinxstylestrong{(Abascal2)}}
\sphinxAtStartPar
Abascal, J Chem Phys, 123, 234505 (2005)

\sphinxlineitem{\sphinxstylestrong{(Ackland)}}
\sphinxAtStartPar
Ackland, Jones, Phys Rev B, 73, 054104 (2006).

\sphinxlineitem{\sphinxstylestrong{(Ackland1)}}
\sphinxAtStartPar
Ackland, Condensed Matter (2005).

\sphinxlineitem{\sphinxstylestrong{(Ackland2)}}
\sphinxAtStartPar
Ackland, Mendelev, Srolovitz, Han and Barashev, Journal of Physics: Condensed Matter, 16, S2629 (2004).

\sphinxlineitem{\sphinxstylestrong{(Addington)}}
\sphinxAtStartPar
Addington, Long, Gubbins, J Chem Phys, 149, 084109 (2018).

\sphinxlineitem{\sphinxstylestrong{(Adhikari et al.)}}
\sphinxAtStartPar
Adhikari, R., Stratford, K.,  Cates, M. E., and Wagner, A. J., Fluctuating lattice Boltzmann, Europhys. Lett. 71 (2005) 473\sphinxhyphen{}479.

\sphinxlineitem{\sphinxstylestrong{(Afshar)}}
\sphinxAtStartPar
Afshar, F. Schmid, A. Pishevar, S. Worley, Comput Phys Comm, 184, 1119\sphinxhyphen{}1128 (2013).

\sphinxlineitem{\sphinxstylestrong{(Agnolin and Roux 2007)}}
\sphinxAtStartPar
Agnolin, I. \& Roux, J\sphinxhyphen{}N. (2007). Internal states of model isotropic granular packings. I. Assembling process, geometry, and contact networks. Phys. Rev. E, 76, 061302.

\sphinxlineitem{\sphinxstylestrong{(Ahrens\sphinxhyphen{}Iwers2022)}}
\sphinxAtStartPar
Ahrens\sphinxhyphen{}Iwers \sphinxstyleemphasis{et al.}, J. Chem. Phys. 157, 084801 (2022).

\sphinxlineitem{\sphinxstylestrong{(Ahrens\sphinxhyphen{}Iwers)}}
\sphinxAtStartPar
Ahrens\sphinxhyphen{}Iwers and Meissner, J. Chem. Phys. 155, 104104 (2021).

\sphinxlineitem{\sphinxstylestrong{(Aktulga)}}
\sphinxAtStartPar
Aktulga, Fogarty, Pandit, Grama, Parallel Computing, 38, 245\sphinxhyphen{}259 (2012).

\sphinxlineitem{\sphinxstylestrong{(Albe)}}
\sphinxAtStartPar
J. Nord, K. Albe, P. Erhart, and K. Nordlund, J. Phys.: Condens. Matter, 15, 5649(2003).

\sphinxlineitem{\sphinxstylestrong{(Albe1)}}
\sphinxAtStartPar
K. Albe, K. Nordlund, J. Nord, and A. Kuronen, Phys. Rev. B, 66, 035205 (2002).

\sphinxlineitem{\sphinxstylestrong{(Allen)}}
\sphinxAtStartPar
Allen and Germano, Mol Phys 104, 3225\sphinxhyphen{}3235 (2006).

\sphinxlineitem{\sphinxstylestrong{(AllenTildesley)}}
\sphinxAtStartPar
Allen and Tildesley, Computer Simulation of Liquids, Oxford University Press (1987)

\sphinxlineitem{\sphinxstylestrong{(Allinger)}}
\sphinxAtStartPar
Allinger, Yuh, Lii, JACS, 111(23), 8551\sphinxhyphen{}8566 (1989),

\sphinxlineitem{\sphinxstylestrong{(Andersen)}}
\sphinxAtStartPar
H. Andersen, J of Comp Phys, 52, 24\sphinxhyphen{}34 (1983).

\sphinxlineitem{\sphinxstylestrong{(Anderson)}}
\sphinxAtStartPar
Anderson, Mukherjee, Critchley, Ziegler, and Lipton “POEMS: Parallelizable Open\sphinxhyphen{}source Efficient Multibody Software “, Engineering With Computers (2006).

\sphinxlineitem{\sphinxstylestrong{(Appshaw)}}
\sphinxAtStartPar
Appshaw, Seddon, Hanna, Soft. Matter,18, 1747(2022).

\sphinxlineitem{\sphinxstylestrong{(Avendano)}}
\sphinxAtStartPar
C. Avendano, T. Lafitte, A. Galindo, C. S. Adjiman, G. Jackson, E. Muller, J Phys Chem B, 115, 11154 (2011).

\sphinxlineitem{\sphinxstylestrong{(Axilrod)}}
\sphinxAtStartPar
Axilrod and Teller, J Chem Phys, 11, 299 (1943); Muto, Nippon Sugaku\sphinxhyphen{}Buturigakkwaishi 17, 629 (1943).

\sphinxlineitem{\sphinxstylestrong{(Babadi)}}
\sphinxAtStartPar
Babadi, Ejtehadi, Everaers, J Comp Phys, 219, 770\sphinxhyphen{}779 (2006).

\sphinxlineitem{\sphinxstylestrong{(Babadi2)}}
\sphinxAtStartPar
Babadi and Ejtehadi, EPL, 77 (2007) 23002.

\sphinxlineitem{\sphinxstylestrong{(Baczewski)}}
\sphinxAtStartPar
A.D. Baczewski and S.D. Bond, J. Chem. Phys. 139, 044107 (2013).

\sphinxlineitem{\sphinxstylestrong{(Bal)}}
\sphinxAtStartPar
K. M Bal and E. C. Neyts, J. Chem. Phys. 141, 204104 (2014).

\sphinxlineitem{\sphinxstylestrong{(Ball)}}
\sphinxAtStartPar
Ball and Melrose, Physica A, 247, 444\sphinxhyphen{}472 (1997).

\sphinxlineitem{\sphinxstylestrong{(Ballenegger)}}
\sphinxAtStartPar
Ballenegger, Arnold, Cerda, J Chem Phys, 131, 094107 (2009).

\sphinxlineitem{\sphinxstylestrong{(Barrat)}}
\sphinxAtStartPar
Barrat and Rodney, J. Stat. Phys, 144, 679 (2011).

\sphinxlineitem{\sphinxstylestrong{(Barrett)}}
\sphinxAtStartPar
Barrett, Tschopp, El Kadiri, Scripta Mat. 66, p.666 (2012).

\sphinxlineitem{\sphinxstylestrong{(Barros)}}
\sphinxAtStartPar
Barros, Sinkovits, Luijten, J. Chem. Phys, 140, 064903 (2014)

\sphinxlineitem{\sphinxstylestrong{(Bartok)}}
\sphinxAtStartPar
Bartok, Payne, Risi, Csanyi, Phys Rev Lett, 104, 136403 (2010).

\sphinxlineitem{\sphinxstylestrong{(Bartok2010)}}
\sphinxAtStartPar
Bartok, Payne, Risi, Csanyi, Phys Rev Lett, 104, 136403 (2010).

\sphinxlineitem{\sphinxstylestrong{(Bartok2013)}}
\sphinxAtStartPar
Bartok, Kondor, Csanyi, Phys Rev B, 87, 184115 (2013).

\sphinxlineitem{\sphinxstylestrong{(Bartok\_PhD)}}
\sphinxAtStartPar
A Bartok\sphinxhyphen{}Partay, PhD Thesis, University of Cambridge, (2010).

\sphinxlineitem{\sphinxstylestrong{(Baskes)}}
\sphinxAtStartPar
Baskes, Phys Rev B, 46, 2727\sphinxhyphen{}2742 (1992).

\sphinxlineitem{\sphinxstylestrong{(Baskes2)}}
\sphinxAtStartPar
Baskes, Phys Rev B, 75, 094113 (2007).

\sphinxlineitem{\sphinxstylestrong{(Beck)}}
\sphinxAtStartPar
Beck, Molecular Physics, 14, 311 (1968).

\sphinxlineitem{\sphinxstylestrong{(Becton)}}
\sphinxAtStartPar
Becton, Averett, Wang, Biomech. Model. Mechanobiology, 18, 425\sphinxhyphen{}433(2019).

\sphinxlineitem{\sphinxstylestrong{(Behler and Parrinello 2007)}}
\sphinxAtStartPar
Behler, J.; Parrinello, M. Phys. Rev. Lett.  2007, 98 (14), 146401.

\sphinxlineitem{\sphinxstylestrong{(Bennet)}}
\sphinxAtStartPar
Bennet, J Comput Phys, 22, 245 (1976)

\sphinxlineitem{\sphinxstylestrong{(Berardi)}}
\sphinxAtStartPar
Berardi, Fava, Zannoni, Chem Phys Lett, 297, 8\sphinxhyphen{}14 (1998). Berardi, Muccioli, Zannoni, J Chem Phys, 128, 024905 (2008).

\sphinxlineitem{\sphinxstylestrong{(Berendsen)}}
\sphinxAtStartPar
Berendsen, Postma, van Gunsteren, DiNola, Haak, J Chem Phys, 81, 3684 (1984).

\sphinxlineitem{\sphinxstylestrong{(Berendsen2)}}
\sphinxAtStartPar
Berendsen, Grigera, Straatsma, J Phys Chem, 91, 6269\sphinxhyphen{}6271 (1987).

\sphinxlineitem{\sphinxstylestrong{(Bessarab)}}
\sphinxAtStartPar
Bessarab, Uzdin, Jonsson, Comp Phys Comm, 196, 335\sphinxhyphen{}347 (2015).

\sphinxlineitem{\sphinxstylestrong{(Beutler)}}
\sphinxAtStartPar
Beutler, Mark, van Schaik, Gerber, van Gunsteren, Chem Phys Lett, 222, 529 (1994).

\sphinxlineitem{\sphinxstylestrong{(Bird)}}
\sphinxAtStartPar
G. A. Bird, “Molecular Gas Dynamics and the Direct Simulation of Gas Flows” (1994).

\sphinxlineitem{\sphinxstylestrong{(Bitzek)}}
\sphinxAtStartPar
Bitzek, Koskinen, Gahler, Moseler, Gumbsch, Phys Rev Lett, 97, 170201 (2006).

\sphinxlineitem{\sphinxstylestrong{(Bolintineanu1)}}
\sphinxAtStartPar
Bolintineanu, Lechman, Plimpton, Grest, Phys Rev E, 86, 066703 (2012).

\sphinxlineitem{\sphinxstylestrong{(Bolintineanu2)}}
\sphinxAtStartPar
Bolintineanu, Grest, Lechman, Pierce, Plimpton, Schunk, Comp Particle Mechanics, 1, 321\sphinxhyphen{}356 (2014).

\sphinxlineitem{\sphinxstylestrong{(Bomont)}}
\sphinxAtStartPar
Bomont, Bretonnet, J. Chem. Phys. 124, 054504 (2006).

\sphinxlineitem{\sphinxstylestrong{(Bond)}}
\sphinxAtStartPar
Bond and Leimkuhler, SIAM J Sci Comput, 30, p 134 (2007).

\sphinxlineitem{\sphinxstylestrong{(Boone)}}
\sphinxAtStartPar
Boone, Babaei, Wilmer, J Chem Theory Comput, 15, 5579\textendash{}5587 (2019).

\sphinxlineitem{\sphinxstylestrong{(BoreschKarplus)}}
\sphinxAtStartPar
Boresch and Karplus, J Phys Chem A, 103, 103 (1999).

\sphinxlineitem{\sphinxstylestrong{(Botu1)}}
\sphinxAtStartPar
V. Botu and R. Ramprasad, Int. J. Quant. Chem., 115(16), 1074 (2015).

\sphinxlineitem{\sphinxstylestrong{(Botu2)}}
\sphinxAtStartPar
V. Botu and R. Ramprasad, Phys. Rev. B, 92(9), 094306 (2015).

\sphinxlineitem{\sphinxstylestrong{(Botu3)}}
\sphinxAtStartPar
V. Botu, R. Batra, J. Chapman and R. Ramprasad, \sphinxurl{https://arxiv.org/abs/1610.02098} (2016).

\sphinxlineitem{\sphinxstylestrong{(Branicio2009)}}
\sphinxAtStartPar
Branicio, Rino, Gan and Tsuzuki, J. Phys Condensed Matter 21 (2009) 095002

\sphinxlineitem{\sphinxstylestrong{(Brennan)}}
\sphinxAtStartPar
Brennan, J Chem Phys Lett, 5, 2144\sphinxhyphen{}2149 (2014).

\sphinxlineitem{\sphinxstylestrong{(Brenner)}}
\sphinxAtStartPar
Brenner, Shenderova, Harrison, Stuart, Ni, Sinnott, J Physics: Condensed Matter, 14, 783\sphinxhyphen{}802 (2002).

\sphinxlineitem{\sphinxstylestrong{(Brilliantov)}}
\sphinxAtStartPar
Brilliantov, Spahn, Hertzsch, Poschel, Phys Rev E, 53, p 5382\sphinxhyphen{}5392 (1996).

\sphinxlineitem{\sphinxstylestrong{(Brooks)}}
\sphinxAtStartPar
Brooks, Brooks, MacKerell Jr., J Comput Chem, 30, 1545 (2009).

\sphinxlineitem{\sphinxstylestrong{(Brooks)}}
\sphinxAtStartPar
Brooks, et al, J Comput Chem, 30, 1545 (2009).

\sphinxlineitem{\sphinxstylestrong{(Brown)}}
\sphinxAtStartPar
Brown et al. International Tables for Crystallography Volume C: Mathematical and Chemical Tables, 554\sphinxhyphen{}95 (2004).

\sphinxlineitem{\sphinxstylestrong{(Buck)}}
\sphinxAtStartPar
Buck, Bouguet\sphinxhyphen{}Bonnet, Pastor, MacKerell Jr., Biophys J, 90, L36 (2006).

\sphinxlineitem{\sphinxstylestrong{(Bulacu)}}
\sphinxAtStartPar
Bulacu, Goga, Zhao, Rossi, Monticelli, Periole, Tieleman, Marrink, J Chem Theory Comput, 9, 3282\sphinxhyphen{}3292

\sphinxlineitem{\sphinxstylestrong{(Bussi)}}\begin{enumerate}
\sphinxsetlistlabels{\Alph}{enumi}{enumii}{}{.}%
\setcounter{enumi}{6}
\item {} 
\sphinxAtStartPar
Bussi, T. Zykova\sphinxhyphen{}Timan, M. Parrinello, J Chem Phys, 130, 074101 (2009).

\end{enumerate}

\sphinxlineitem{\sphinxstylestrong{(Bussi1)}}
\sphinxAtStartPar
Bussi, Donadio and Parrinello, J. Chem. Phys. 126, 014101(2007)

\sphinxlineitem{\sphinxstylestrong{(Bussi2)}}
\sphinxAtStartPar
Bussi and Parrinello, Phys. Rev. E 75, 056707 (2007)

\sphinxlineitem{\sphinxstylestrong{(COMB\_1)}}
\sphinxAtStartPar
J. Yu, S. B. Sinnott, S. R. Phillpot, Phys Rev B, 75, 085311 (2007),

\sphinxlineitem{\sphinxstylestrong{(COMB\_2)}}
\sphinxAtStartPar
T.\sphinxhyphen{}R. Shan, B. D. Devine, T. W. Kemper, S. B. Sinnott, and S. R. Phillpot, Phys. Rev. B 81, 125328 (2010)

\sphinxlineitem{\sphinxstylestrong{(COMB3)}}
\sphinxAtStartPar
T. Liang, T.\sphinxhyphen{}R. Shan, Y.\sphinxhyphen{}T. Cheng, B. D. Devine, M. Noordhoek, Y. Li, Z. Lu, S. R. Phillpot, and S. B. Sinnott, Mat. Sci. \& Eng: R 74, 255\sphinxhyphen{}279 (2013).

\sphinxlineitem{\sphinxstylestrong{(Calhoun)}}
\sphinxAtStartPar
A. Calhoun, M. Pavese, G. Voth, Chem Phys Letters, 262, 415 (1996).

\sphinxlineitem{\sphinxstylestrong{(Campana)}}
\sphinxAtStartPar
C. Campana and M. H. Muser, Phys. Rev. B {[}74{]}, 075420 (2006).

\sphinxlineitem{\sphinxstylestrong{(Cao1)}}
\sphinxAtStartPar
J. Cao and B. Berne, J Chem Phys, 99, 2902 (1993).

\sphinxlineitem{\sphinxstylestrong{(Cao2)}}
\sphinxAtStartPar
J. Cao and G. Voth, J Chem Phys, 100, 5093 (1994).

\sphinxlineitem{\sphinxstylestrong{(Caro)}}
\sphinxAtStartPar
A Caro, DA Crowson, M Caro; Phys Rev Lett, 95, 075702 (2005)

\sphinxlineitem{\sphinxstylestrong{(CasP)}}
\sphinxAtStartPar
CasP webpage: \sphinxurl{http://www.casp-program.org/}

\sphinxlineitem{\sphinxstylestrong{(Cawkwell2012)}}
\sphinxAtStartPar
A. M. N. Niklasson, M. J. Cawkwell, Phys. Rev. B, 86 (17), 174308 (2012).

\sphinxlineitem{\sphinxstylestrong{(Cercignani)}}
\sphinxAtStartPar
C. Cercignani and M. Lampis. Trans. Theory Stat. Phys. 1, 2, 101 (1971).

\sphinxlineitem{\sphinxstylestrong{(Cerda)}}
\sphinxAtStartPar
Cerda, Ballenegger, Lenz, Holm, J Chem Phys 129, 234104 (2008)

\sphinxlineitem{\sphinxstylestrong{(Ceriotti)}}\begin{enumerate}
\sphinxsetlistlabels{\Alph}{enumi}{enumii}{}{.}%
\setcounter{enumi}{12}
\item {} 
\sphinxAtStartPar
Ceriotti, M. Parrinello, T. Markland, D. Manolopoulos, J. Chem. Phys. 133, 124104 (2010).

\end{enumerate}

\sphinxlineitem{\sphinxstylestrong{(Ceriotti1)}}
\sphinxAtStartPar
Ceriotti, Bussi and Parrinello, J Chem Theory Comput 6, 1170\sphinxhyphen{}80 (2010)

\sphinxlineitem{\sphinxstylestrong{(Ceriotti2)}}
\sphinxAtStartPar
Ceriotti, Bussi and Parrinello, Phys Rev Lett 103, 030603 (2009)

\sphinxlineitem{\sphinxstylestrong{(Cerutti)}}
\sphinxAtStartPar
Cerutti, Duke, Darden, Lybrand, Journal of Chemical Theory and Computation 5, 2322 (2009)

\sphinxlineitem{\sphinxstylestrong{(Chen)}}
\sphinxAtStartPar
J Chen, D Tzou and J Beraun, Int. J. Heat Mass Transfer, 49, 307\sphinxhyphen{}316 (2006).

\sphinxlineitem{\sphinxstylestrong{(Chenoweth\_2008)}}
\sphinxAtStartPar
Chenoweth, van Duin and Goddard, Journal of Physical Chemistry A, 112, 1040\sphinxhyphen{}1053 (2008).

\sphinxlineitem{\sphinxstylestrong{(Clarke)}}
\sphinxAtStartPar
Clarke and Smith, J Chem Phys, 84, 2290 (1986).

\sphinxlineitem{\sphinxstylestrong{(Clavier)}}\begin{enumerate}
\sphinxsetlistlabels{\Alph}{enumi}{enumii}{}{.}%
\setcounter{enumi}{6}
\item {} 
\sphinxAtStartPar
Clavier, N. Desbiens, E. Bourasseau, V. Lachet, N. Brusselle\sphinxhyphen{}Dupend and B. Rousseau, Mol Sim, 43, 1413 (2017).

\end{enumerate}

\sphinxlineitem{\sphinxstylestrong{(Clemmer)}}
\sphinxAtStartPar
Clemmer and Robbins, Phys. Rev. Lett. (2022).

\sphinxlineitem{\sphinxstylestrong{(Clemmer1)}}
\sphinxAtStartPar
Clemmer, Monti, Lechman, Soft Matter, 20, 1702 (2024).

\sphinxlineitem{\sphinxstylestrong{(Clemmer2)}}
\sphinxAtStartPar
Clemmer, Pierce, O’Connor, Nevins, Jones, Lechman, Tencer, Appl. Math. Model., 130, 310\sphinxhyphen{}326 (2024).

\sphinxlineitem{\sphinxstylestrong{(Coleman)}}
\sphinxAtStartPar
Coleman, Spearot, Capolungo, MSMSE, 21, 055020 (2013).

\sphinxlineitem{\sphinxstylestrong{(Colliex)}}
\sphinxAtStartPar
Colliex et al. International Tables for Crystallography Volume C: Mathematical and Chemical Tables, 249\sphinxhyphen{}429 (2004).

\sphinxlineitem{\sphinxstylestrong{(Cooke)}}
\sphinxAtStartPar
“Cooke, Kremer and Deserno, Phys. Rev. E, 72, 011506 (2005)”

\sphinxlineitem{\sphinxstylestrong{(Cornell)}}
\sphinxAtStartPar
Cornell, Cieplak, Bayly, Gould, Merz, Ferguson, Spellmeyer, Fox, Caldwell, Kollman, JACS 117, 5179\sphinxhyphen{}5197 (1995).

\sphinxlineitem{\sphinxstylestrong{(Cundall, 1987)}}
\sphinxAtStartPar
Cundall, P. A. Distinct Element Models of Rock and Soil

\sphinxlineitem{\sphinxstylestrong{(Curk1)}}\begin{enumerate}
\sphinxsetlistlabels{\Alph}{enumi}{enumii}{}{.}%
\setcounter{enumi}{19}
\item {} 
\sphinxAtStartPar
Curk, J. Yuan, and E. Luijten, JCP 156 (2022).

\end{enumerate}

\sphinxlineitem{\sphinxstylestrong{(Curk2)}}\begin{enumerate}
\sphinxsetlistlabels{\Alph}{enumi}{enumii}{}{.}%
\setcounter{enumi}{19}
\item {} 
\sphinxAtStartPar
Curk and E. Luijten, PRL 126 (2021)

\end{enumerate}

\sphinxlineitem{\sphinxstylestrong{(Cusentino)}}
\sphinxAtStartPar
Cusentino, Wood, Thompson, J Phys Chem A, 124, 5456, (2020)

\sphinxlineitem{\sphinxstylestrong{(Daivis and Todd)}}
\sphinxAtStartPar
Daivis and Todd, J Chem Phys, 124, 194103 (2006).

\sphinxlineitem{\sphinxstylestrong{(Daivis and Todd)}}
\sphinxAtStartPar
Daivis and Todd, Nonequilibrium Molecular Dynamics (book), Cambridge University Press, \sphinxurl{https://doi.org/10.1017/9781139017848}, (2017).

\sphinxlineitem{\sphinxstylestrong{(Dammak)}}
\sphinxAtStartPar
Dammak, Chalopin, Laroche, Hayoun, and Greffet, Phys Rev Lett, 103, 190601 (2009).

\sphinxlineitem{\sphinxstylestrong{(Darden)}}
\sphinxAtStartPar
Darden, York, Pedersen, J Chem Phys, 98, 10089 (1993).

\sphinxlineitem{\sphinxstylestrong{(Davidchack)}}
\sphinxAtStartPar
R.L Davidchack, T.E. Ouldridge, and M.V. Tretyakov. J. Chem. Phys. 142, 144114 (2015).

\sphinxlineitem{\sphinxstylestrong{(Daw1)}}
\sphinxAtStartPar
Daw, Baskes, Phys Rev Lett, 50, 1285 (1983). Daw, Baskes, Phys Rev B, 29, 6443 (1984).

\sphinxlineitem{\sphinxstylestrong{(Daw2)}}
\sphinxAtStartPar
M. S. Daw, and M. I. Baskes, Phys. Rev. B, 29, 6443 (1984).

\sphinxlineitem{\sphinxstylestrong{(de Buyl)}}
\sphinxAtStartPar
de Buyl, Colberg and Hofling, Comp. Phys. Comm. 185(6), 1546\sphinxhyphen{}1553 (2014) \sphinxhyphen{}

\sphinxlineitem{\sphinxstylestrong{(Deissenbeck)}}
\sphinxAtStartPar
Deissenbeck \sphinxstyleemphasis{et al.}, Phys. Rev. Letters 126, 136803 (2021).

\sphinxlineitem{\sphinxstylestrong{(de Koning)}}
\sphinxAtStartPar
de Koning and Antonelli, Phys Rev E, 53, 465 (1996).

\sphinxlineitem{\sphinxstylestrong{(DeVane)}}
\sphinxAtStartPar
Shinoda, DeVane, Klein, Soft Matter, 4, 2453\sphinxhyphen{}2462 (2008).

\sphinxlineitem{\sphinxstylestrong{(Deserno)}}
\sphinxAtStartPar
Deserno and Holm, J Chem Phys, 109, 7694 (1998).

\sphinxlineitem{\sphinxstylestrong{(Destree)}}
\sphinxAtStartPar
M. Destree, F. Laupretre, A. Lyulin, and J.\sphinxhyphen{}P.  Ryckaert, J Chem Phys, 112, 9632 (2000).

\sphinxlineitem{\sphinxstylestrong{(Dickel)}}
\sphinxAtStartPar
Dickel, Francis, and Barrett, Computational Materials Science 171 (2020): 109157.

\sphinxlineitem{\sphinxstylestrong{(Dietz)}}
\sphinxAtStartPar
Dietz and Hoy, J. Chem Phys, 156, 014103 (2022).

\sphinxlineitem{\sphinxstylestrong{(Dobson)}}
\sphinxAtStartPar
Dobson, J Chem Phys, 141, 184103 (2014).

\sphinxlineitem{\sphinxstylestrong{(Drautz)}}
\sphinxAtStartPar
Drautz, Phys Rev B, 99, 014104 (2019).

\sphinxlineitem{\sphinxstylestrong{(Duffy)}}
\sphinxAtStartPar
D M Duffy and A M Rutherford, J. Phys.: Condens. Matter, 19, 016207\sphinxhyphen{}016218 (2007).

\sphinxlineitem{\sphinxstylestrong{(Dufils)}}
\sphinxAtStartPar
Dufils \sphinxstyleemphasis{et al.}, Phys. Rev. Letters 123, 195501 (2019).

\sphinxlineitem{\sphinxstylestrong{(Dullweber)}}
\sphinxAtStartPar
Dullweber, Leimkuhler and McLachlan, J Chem Phys, 107, 5840 (1997).

\sphinxlineitem{\sphinxstylestrong{(Dunn1)}}
\sphinxAtStartPar
Dunn and Noid, J Chem Phys, 143, 243148 (2015).

\sphinxlineitem{\sphinxstylestrong{(Dunn2)}}
\sphinxAtStartPar
Dunn, Lebold, DeLyser, Rudzinski, and Noid, J. Phys. Chem. B, 122, 3363 (2018).

\sphinxlineitem{\sphinxstylestrong{(Dunweg)}}
\sphinxAtStartPar
Dunweg and Paul, Int J of Modern Physics C, 2, 817\sphinxhyphen{}27 (1991).

\sphinxlineitem{\sphinxstylestrong{(EcheverriRestrepo)}}
\sphinxAtStartPar
Echeverri Restrepo, Andric, Comput Mater Sci, 218, 111978 (2023).

\sphinxlineitem{\sphinxstylestrong{(EDIP)}}
\sphinxAtStartPar
J F Justo et al, Phys Rev B 58, 2539 (1998).

\sphinxlineitem{\sphinxstylestrong{(Eike)}}
\sphinxAtStartPar
Eike and Maginn, Journal of Chemical Physics, 124, 164503 (2006).

\sphinxlineitem{\sphinxstylestrong{(Elliott)}}
\sphinxAtStartPar
Elliott, Tadmor and Bernstein, \sphinxurl{https://openkim.org/kim-api} (2011) doi: \sphinxurl{https://doi.org/10.25950/FF8F563A}

\sphinxlineitem{\sphinxstylestrong{(Ellis)}}
\sphinxAtStartPar
Ellis, Fiedler, Popoola, Modine, Stephens, Thompson, Cangi, Rajamanickam,  Phys Rev B, 104, 035120, (2021)

\sphinxlineitem{\sphinxstylestrong{(Emmrich)}}
\sphinxAtStartPar
Emmrich, Weckner, Commun. Math. Sci., 5, 851\sphinxhyphen{}864 (2007),

\sphinxlineitem{\sphinxstylestrong{(Elstner)}}
\sphinxAtStartPar
M. Elstner, D. Poresag, G. Jungnickel, J. Elsner, M. Haugk, T. Frauenheim, S. Suhai, and G. Seifert, Phys. Rev. B, 58, 7260 (1998).

\sphinxlineitem{\sphinxstylestrong{(Erdmann)}}
\sphinxAtStartPar
U. Erdmann , W. Ebeling, L. Schimansky\sphinxhyphen{}Geier, and F. Schweitzer, Eur. Phys. J. B 15, 105\sphinxhyphen{}113, 2000.

\sphinxlineitem{\sphinxstylestrong{(Eshelby)}}
\sphinxAtStartPar
J.D. Eshelby, Philos. Trans. Royal Soc. London A, Math. Phys. Sci., Vol. 244, No. 877 (1951) pp. 87\sphinxhyphen{}112; J. Elasticity, Vol. 5, Nos. 3\sphinxhyphen{}4 (1975) pp. 321\sphinxhyphen{}335{]}

\sphinxlineitem{\sphinxstylestrong{(Espanol and Revenga)}}
\sphinxAtStartPar
Espanol, Revenga, Physical Review E, 67, 026705 (2003).

\sphinxlineitem{\sphinxstylestrong{(Espanol1997)}}
\sphinxAtStartPar
Espanol, Europhys Lett, 40(6): 631\sphinxhyphen{}636 (1997).  DOI:10.1209/epl/i1997\sphinxhyphen{}00515\sphinxhyphen{}8

\sphinxlineitem{\sphinxstylestrong{(Evans and Morriss)}}
\sphinxAtStartPar
Evans and Morriss, Phys Rev A, 30, 1528 (1984).

\sphinxlineitem{\sphinxstylestrong{(Evans)}}
\sphinxAtStartPar
Evans and Morriss, Phys. Rev. Lett. 56, 2172 (1986).

\sphinxlineitem{\sphinxstylestrong{(Everaers)}}
\sphinxAtStartPar
Everaers and Ejtehadi, Phys Rev E, 67, 041710 (2003).

\sphinxlineitem{\sphinxstylestrong{(Faken)}}
\sphinxAtStartPar
Faken, Jonsson, Comput Mater Sci, 2, 279 (1994).

\sphinxlineitem{\sphinxstylestrong{(Falk)}}
\sphinxAtStartPar
Falk and Langer PRE, 57, 7192 (1998).

\sphinxlineitem{\sphinxstylestrong{(Fath)}}
\sphinxAtStartPar
Fath, Hochbruck, Singh, J Comp Phys, 333, 180\sphinxhyphen{}198 (2017).

\sphinxlineitem{\sphinxstylestrong{(Feng1)}}\begin{enumerate}
\sphinxsetlistlabels{\arabic}{enumi}{enumii}{}{.}%
\setcounter{enumi}{25}
\item {} 
\sphinxAtStartPar
Feng, …, and W. Ouyang, J. Phys. Chem. C. 127(18), 8704\sphinxhyphen{}8713 (2023).

\end{enumerate}

\sphinxlineitem{\sphinxstylestrong{(Feng2)}}\begin{enumerate}
\sphinxsetlistlabels{\arabic}{enumi}{enumii}{}{.}%
\setcounter{enumi}{25}
\item {} 
\sphinxAtStartPar
Feng, …, and W. Ouyang, Langmuir 39(50), 18198\sphinxhyphen{}18207 (2023).

\end{enumerate}

\sphinxlineitem{\sphinxstylestrong{(Fennell)}}
\sphinxAtStartPar
C. J. Fennell, J. D. Gezelter, J Chem Phys, 124, 234104 (2006).

\sphinxlineitem{\sphinxstylestrong{(Feynman)}}
\sphinxAtStartPar
R. Feynman and A. Hibbs, Chapter 7, Quantum Mechanics and Path Integrals, McGraw\sphinxhyphen{}Hill, New York (1965).

\sphinxlineitem{\sphinxstylestrong{(Fichthorn)}}
\sphinxAtStartPar
Fichthorn, Balankura, Qi, CrystEngComm, 18(29), 5410\sphinxhyphen{}5417 (2016).

\sphinxlineitem{\sphinxstylestrong{(Fily)}}
\sphinxAtStartPar
Y. Fily and M.C. Marchetti, Phys. Rev. Lett. 108, 235702, 2012.  Default

\sphinxlineitem{\sphinxstylestrong{(Fincham)}}
\sphinxAtStartPar
Fincham, Mackrodt and Mitchell, J Phys Condensed Matter, 6, 393\sphinxhyphen{}404 (1994).

\sphinxlineitem{\sphinxstylestrong{(Finnis1)}}
\sphinxAtStartPar
Finnis, Sinclair, Philosophical Magazine A, 50, 45 (1984).

\sphinxlineitem{\sphinxstylestrong{(Finnis2)}}
\sphinxAtStartPar
M. W. Finnis, A. T. Paxton, M. Methfessel, and M. van Schilfgarde, Phys. Rev. Lett., 81, 5149 (1998).

\sphinxlineitem{\sphinxstylestrong{(Fiorin)}}
\sphinxAtStartPar
Fiorin, Klein, Henin, Mol. Phys., DOI:10.1080/00268976.2013.813594

\sphinxlineitem{\sphinxstylestrong{(Fox)}}
\sphinxAtStartPar
Fox, O’Keefe, Tabbernor, Acta Crystallogr. A, 45, 786\sphinxhyphen{}93 (1989).

\sphinxlineitem{\sphinxstylestrong{(Fraige)}}
\sphinxAtStartPar
F. Y. Fraige, P. A. Langston, A. J. Matchett, J. Dodds, Particuology, 6, 455 (2008).

\sphinxlineitem{\sphinxstylestrong{(Freitas)}}
\sphinxAtStartPar
Freitas, Asta, and de Koning, Computational Materials Science, 112, 333 (2016).

\sphinxlineitem{\sphinxstylestrong{(Frenkel)}}
\sphinxAtStartPar
Frenkel and Smit, Understanding Molecular Simulation, Academic Press, London, 2002.

\sphinxlineitem{\sphinxstylestrong{(Fu)}}
\sphinxAtStartPar
Fu, Peng, Yuan, Kfoury, Young, Comput. Phys. Commun, 210, 193\sphinxhyphen{}203(2017).

\sphinxlineitem{\sphinxstylestrong{(Gao)}}
\sphinxAtStartPar
Gao and Weber, Nuclear Instruments and Methods in Physics Research B 191 (2012) 504.

\sphinxlineitem{\sphinxstylestrong{(Gingrich)}}
\sphinxAtStartPar
Gingrich, \sphinxtitleref{MSc thesis} \textless{}\sphinxurl{https://gingrich.chem.northwestern.edu/papers/ThesiswCorrections.pdf}\textgreater{}\textasciigrave{} (2010).

\sphinxlineitem{\sphinxstylestrong{(Gissinger2017)}}
\sphinxAtStartPar
Gissinger, Jensen and Wise, Polymer, 128, 211\sphinxhyphen{}217 (2017).

\sphinxlineitem{\sphinxstylestrong{(Gissinger2020)}}
\sphinxAtStartPar
Gissinger, Jensen and Wise, Macromolecules, 53, 22, 9953\sphinxhyphen{}9961 (2020).

\sphinxlineitem{\sphinxstylestrong{(Gissinger)}}
\sphinxAtStartPar
Jacob R. Gissinger, Scott R. Zavada, Joseph G. Smith, Josh Kemppainen, Ivan Gallegos, Gregory M. Odegard, Emilie J. Siochi, and Kristopher E. Wise, Carbon, 202, 336\sphinxhyphen{}347 (2023).

\sphinxlineitem{\sphinxstylestrong{(Gissinger2024)}}\begin{enumerate}
\sphinxsetlistlabels{\Alph}{enumi}{enumii}{}{.}%
\setcounter{enumi}{9}
\item {} \begin{enumerate}
\sphinxsetlistlabels{\Alph}{enumii}{enumiii}{}{.}%
\setcounter{enumii}{17}
\item {} 
\sphinxAtStartPar
Gissinger, I. Nikiforov, Y. Afshar, B. Waters, M. Choi, D. S. Karls, A. Stukowski, W. Im, H. Heinz, A. Kohlmeyer, and E. B. Tadmor, J Phys Chem B, 128, 3282\sphinxhyphen{}3297 (2024).

\end{enumerate}

\end{enumerate}

\sphinxlineitem{\sphinxstylestrong{(Gloor)}}
\sphinxAtStartPar
Gloor, J Chem Phys, 123, 134703 (2005)

\sphinxlineitem{\sphinxstylestrong{(Glosli)}}
\sphinxAtStartPar
Glosli, unpublished, 2005. Streitz, Glosli, Patel, Chan, Yates, de Supinski, Sexton and Gunnels, Journal of Physics: Conference Series, 46, 254 (2006).

\sphinxlineitem{\sphinxstylestrong{(Goff)}}
\sphinxAtStartPar
Goff, Zhang, Negre, Rohskopf, Niklasson, Journal of Chemical Theory and Computation 19, no. 13 (2023).

\sphinxlineitem{\sphinxstylestrong{(Goldman1)}}
\sphinxAtStartPar
Goldman, Reed and Fried, J. Chem. Phys. 131, 204103 (2009)

\sphinxlineitem{\sphinxstylestrong{(Goldman2)}}
\sphinxAtStartPar
Goldman, Srinivasan, Hamel, Fried, Gaus, and Elstner, J. Phys. Chem. C, 117, 7885 (2013).

\sphinxlineitem{\sphinxstylestrong{(Grime)}}
\sphinxAtStartPar
Grime and Voth, to appear in J Chem Theory \& Computation (2014).

\sphinxlineitem{\sphinxstylestrong{(Grimme)}}
\sphinxAtStartPar
Grimme, J Comput Chem, 27(15), 1787\sphinxhyphen{}1799 (2006).

\sphinxlineitem{\sphinxstylestrong{(Gronbech\sphinxhyphen{}Jensen1)}}
\sphinxAtStartPar
Gronbech Jensen and Gronbech\sphinxhyphen{}Jensen, Mol Phys, 117, 2511 (2019)

\sphinxlineitem{\sphinxstylestrong{(Gronbech\sphinxhyphen{}Jensen2)}}
\sphinxAtStartPar
Gronbech\sphinxhyphen{}Jensen and Farago, Mol Phys, 111, 983 (2013)

\sphinxlineitem{\sphinxstylestrong{(Gronbech\sphinxhyphen{}Jensen3)}}
\sphinxAtStartPar
Hayre, and Farago, Comp Phys Comm, 185, 524 (2014)

\sphinxlineitem{\sphinxstylestrong{(Groot)}}
\sphinxAtStartPar
Groot and Warren, J Chem Phys, 107: 4423\sphinxhyphen{}4435 (1997).  DOI:10.1063/1.474784

\sphinxlineitem{\sphinxstylestrong{(Guenole)}}
\sphinxAtStartPar
Guenole, Noehring, Vaid, Houlle, Xie, Prakash, Bitzek, Comput Mater Sci, 175, 109584 (2020).

\sphinxlineitem{\sphinxstylestrong{(Gullet)}}
\sphinxAtStartPar
Gullet, Wagner, Slepoy, SANDIA Report 2003\sphinxhyphen{}8782 (2003). DOI:10.2172/918395

\sphinxlineitem{\sphinxstylestrong{(Guo)}}
\sphinxAtStartPar
Guo and Thirumalai, Journal of Molecular Biology, 263, 323\sphinxhyphen{}43 (1996).

\sphinxlineitem{\sphinxstylestrong{(Gupta)}}
\sphinxAtStartPar
Gupta ,Phys Rev. B, 23, 6265\sphinxhyphen{}6270 (1981).

\sphinxlineitem{\sphinxstylestrong{(Hardy)}}
\sphinxAtStartPar
David Hardy thesis: Multilevel Summation for the Fast Evaluation of Forces for the Simulation of Biomolecules, University of Illinois at Urbana\sphinxhyphen{}Champaign, (2006).

\sphinxlineitem{\sphinxstylestrong{(Hardy2)}}
\sphinxAtStartPar
Hardy, Stone, Schulten, Parallel Computing, 35, 164\sphinxhyphen{}177 (2009).

\sphinxlineitem{\sphinxstylestrong{(Hecht)}}
\sphinxAtStartPar
Hecht, Harting, Ihle, Herrmann, Phys Rev E, 72, 011408 (2005).

\sphinxlineitem{\sphinxstylestrong{(Henkelman1)}}
\sphinxAtStartPar
Henkelman and Jonsson, J Chem Phys, 113, 9978\sphinxhyphen{}9985 (2000).

\sphinxlineitem{\sphinxstylestrong{(Henkelman2)}}
\sphinxAtStartPar
Henkelman, Uberuaga, Jonsson, J Chem Phys, 113, 9901\sphinxhyphen{}9904 (2000).

\sphinxlineitem{\sphinxstylestrong{(Henkes)}}
\sphinxAtStartPar
Henkes, S, Fily, Y., and Marchetti, M. C. Phys. Rev. E, 84, 040301(R), 2011.

\sphinxlineitem{\sphinxstylestrong{(Henrich)}}
\sphinxAtStartPar
O. Henrich, Y. A. Gutierrez\sphinxhyphen{}Fosado, T. Curk, T. E. Ouldridge, Eur. Phys. J. E 41, 57 (2018).

\sphinxlineitem{\sphinxstylestrong{(Herman)}}\begin{enumerate}
\sphinxsetlistlabels{\Alph}{enumi}{enumii}{}{.}%
\setcounter{enumi}{12}
\item {} \begin{enumerate}
\sphinxsetlistlabels{\Alph}{enumii}{enumiii}{}{.}%
\setcounter{enumii}{5}
\item {} 
\sphinxAtStartPar
Herman, E. J. Bruskin, B. J. Berne, J Chem Phys, 76, 5150 (1982).

\end{enumerate}

\end{enumerate}

\sphinxlineitem{\sphinxstylestrong{(Hess)}}
\sphinxAtStartPar
Hess, B. The Journal of Chemical Physics 2002, 116 (1), 209\sphinxhyphen{}217.

\sphinxlineitem{\sphinxstylestrong{(Heyes)}}
\sphinxAtStartPar
Heyes, Phys Rev B, 49, 755 (1994).

\sphinxlineitem{\sphinxstylestrong{(Hijazi)}}
\sphinxAtStartPar
M. Hijazi, D. M. Wilkins, M. Ceriotti, J. Chem. Phys. 148, 184109 (2018)

\sphinxlineitem{\sphinxstylestrong{(Hockney)}}
\sphinxAtStartPar
Hockney and Eastwood, Computer Simulation Using Particles, Adam Hilger, NY (1989).

\sphinxlineitem{\sphinxstylestrong{(Holian)}}
\sphinxAtStartPar
Holian and Ravelo, Phys Rev B, 51, 11275 (1995).

\sphinxlineitem{\sphinxstylestrong{(Hone)}}
\sphinxAtStartPar
T. Hone, P. Rossky, G. Voth, J Chem Phys, 124, 154103 (2006).

\sphinxlineitem{\sphinxstylestrong{(Hoover)}}
\sphinxAtStartPar
Hoover, Phys Rev A, 31, 1695 (1985).

\sphinxlineitem{\sphinxstylestrong{(Huang)}}
\sphinxAtStartPar
Huang, Zhang, Yuan, Gao, Zhang, Nano Lett. 13, 4546(2013).

\sphinxlineitem{\sphinxstylestrong{(Huang2014)}}\begin{enumerate}
\sphinxsetlistlabels{\Alph}{enumi}{enumii}{}{.}%
\setcounter{enumi}{23}
\item {} 
\sphinxAtStartPar
Huang, “Exploring critical\sphinxhyphen{}state behavior using DEM”, Doctoral dissertation, Imperial College. (2014). \sphinxurl{https://doi.org/10.25560/25316}

\end{enumerate}

\sphinxlineitem{\sphinxstylestrong{(Hu)}}
\sphinxAtStartPar
Hu, and Adams J. Comp. Physics, 213, 844\sphinxhyphen{}861 (2006).

\sphinxlineitem{\sphinxstylestrong{(Hu)}}
\sphinxAtStartPar
Hu, J. Chem. Theory Comput. 10, 5254 (2014).

\sphinxlineitem{\sphinxstylestrong{(Hummer)}}
\sphinxAtStartPar
Hummer, Gronbech\sphinxhyphen{}Jensen, Neumann, J Chem Phys, 109, 2791 (1998)

\sphinxlineitem{\sphinxstylestrong{(Hunt)}}
\sphinxAtStartPar
Hunt, Mol Simul, 42, 347 (2016).

\sphinxlineitem{\sphinxstylestrong{(Ikeshoji)}}
\sphinxAtStartPar
Ikeshoji and Hafskjold, Molecular Physics, 81, 251\sphinxhyphen{}261 (1994).

\sphinxlineitem{\sphinxstylestrong{(Ikeshoji2)}}
\sphinxAtStartPar
Ikeshoji, Hafskjold, Furuholt, Mol Sim, 29, 101\sphinxhyphen{}109, (2003).

\sphinxlineitem{\sphinxstylestrong{(Ilie)}}
\sphinxAtStartPar
Ilie, Briels, den Otter, Journal of Chemical Physics, 142, 114103 (2015).

\sphinxlineitem{\sphinxstylestrong{(In ‘t Veld)}}
\sphinxAtStartPar
In ‘t Veld, Ismail, Grest, J Chem Phys (accepted) (2007).

\sphinxlineitem{\sphinxstylestrong{(IPI)}}
\sphinxAtStartPar
\sphinxtitleref{https://ipi\sphinxhyphen{}code.org/ \textless{}https://ipi\sphinxhyphen{}code.org/\textgreater{}}

\sphinxlineitem{\sphinxstylestrong{(IPI\sphinxhyphen{}CPC)}}
\sphinxAtStartPar
Ceriotti, More and Manolopoulos, Comp Phys Comm, 185, 1019\sphinxhyphen{}1026 (2014).

\sphinxlineitem{\sphinxstylestrong{(Isele\sphinxhyphen{}Holder)}}
\sphinxAtStartPar
Isele\sphinxhyphen{}Holder, Mitchell, Ismail, J Chem Phys, 137, 174107 (2012).

\sphinxlineitem{\sphinxstylestrong{(Isele\sphinxhyphen{}Holder2)}}
\sphinxAtStartPar
Isele\sphinxhyphen{}Holder, Mitchell, Hammond, Kohlmeyer, Ismail, J Chem Theory Comput 9, 5412 (2013).

\sphinxlineitem{\sphinxstylestrong{(Ismail)}}
\sphinxAtStartPar
Ismail, Tsige, In ‘t Veld, Grest, Molecular Physics (accepted) (2007).

\sphinxlineitem{\sphinxstylestrong{(Ivanov)}}
\sphinxAtStartPar
Ivanov, Uzdin, Jonsson. arXiv preprint arXiv:1904.02669, (2019).

\sphinxlineitem{\sphinxstylestrong{(Izrailev)}}
\sphinxAtStartPar
Izrailev, Stepaniants, Isralewitz, Kosztin, Lu, Molnar, Wriggers, Schulten. Computational Molecular Dynamics: Challenges, Methods, Ideas, volume 4 of Lecture Notes in Computational Science and

\sphinxlineitem{\sphinxstylestrong{(Izvekov)}}
\sphinxAtStartPar
Izvekov, Voth, J Chem Phys 123, 134105 (2005).

\sphinxlineitem{\sphinxstylestrong{(Jadhao)}}
\sphinxAtStartPar
Jadhao, Solis, Olvera de la Cruz, J Chem Phys, 138, 054119 (2013)

\sphinxlineitem{\sphinxstylestrong{(Janssens)}}
\sphinxAtStartPar
Janssens, Olmsted, Holm, Foiles, Plimpton, Derlet, Nature Materials, 5, 124\sphinxhyphen{}127 (2006).

\sphinxlineitem{\sphinxstylestrong{(Jaramillo\sphinxhyphen{}Botero)}}
\sphinxAtStartPar
Jaramillo\sphinxhyphen{}Botero, Su, Qi, Goddard, Large\sphinxhyphen{}scale, Long\sphinxhyphen{}term Non\sphinxhyphen{}adiabatic Electron Molecular Dynamics for Describing Material Properties and Phenomena in Extreme Environments, J Comp

\sphinxlineitem{\sphinxstylestrong{(Jarzynski)}}
\sphinxAtStartPar
Jarzynski, Phys. Rev. Lett. 78, 2690 (1997)

\sphinxlineitem{\sphinxstylestrong{(Jiang)}}
\sphinxAtStartPar
Jiang, Hardy, Phillips, MacKerell, Schulten, and Roux,  J Phys Chem Lett, 2, 87\sphinxhyphen{}92 (2011).

\sphinxlineitem{\sphinxstylestrong{(Jiang1)}}
\sphinxAtStartPar
Jiang, Hardy, Phillips, MacKerell, Schulten, and Roux, J Phys Chem Lett, 2, 87\sphinxhyphen{}92 (2011).

\sphinxlineitem{\sphinxstylestrong{(Jiang2)}}
\sphinxAtStartPar
J.\sphinxhyphen{}W. Jiang, Nanotechnology 26, 315706 (2015).

\sphinxlineitem{\sphinxstylestrong{(Jiang3)}}
\sphinxAtStartPar
J.\sphinxhyphen{}W. Jiang, Acta Mech. Solida. Sin 32, 17 (2019).

\sphinxlineitem{\sphinxstylestrong{(Johnson et al, 1971)}}
\sphinxAtStartPar
Johnson, K. L., Kendall, K., \& Roberts, A. D. (1971).  Surface energy and the contact of elastic solids. Proc. R. Soc. Lond. A, 324(1558), 301\sphinxhyphen{}313.

\sphinxlineitem{\sphinxstylestrong{(Jones)}}
\sphinxAtStartPar
Jones, RE; Templeton, JA; Wagner, GJ; Olmsted, D; Modine, JA, “Electron transport enhanced molecular dynamics for metals and semi\sphinxhyphen{}metals.” International Journal for Numerical Methods in Engineering (2010), 83:940.

\sphinxlineitem{\sphinxstylestrong{(Jonsson)}}
\sphinxAtStartPar
Jonsson, Mills and Jacobsen, in Classical and Quantum Dynamics in Condensed Phase Simulations, edited by Berne, Ciccotti, and Coker World Scientific, Singapore, 1998, p 385.

\sphinxlineitem{\sphinxstylestrong{(Jorgensen)}}
\sphinxAtStartPar
Jorgensen, Chandrasekhar, Madura, Impey, Klein, J Chem Phys, 79, 926 (1983).

\sphinxlineitem{\sphinxstylestrong{(Jusufi)}}
\sphinxAtStartPar
Jusufi, Hynninen, and Panagiotopoulos, J Phys Chem B, 112, 13783 (2008).

\sphinxlineitem{\sphinxstylestrong{(Kamberaj)}}
\sphinxAtStartPar
Kamberaj, Low, Neal, J Chem Phys, 122, 224114 (2005).

\sphinxlineitem{\sphinxstylestrong{(Katsura)}}
\sphinxAtStartPar
H. Katsura, N. Nagaosa, A.V. Balatsky. Phys. Rev. Lett., 95(5), 057205. (2005)

\sphinxlineitem{\sphinxstylestrong{(Kelchner)}}
\sphinxAtStartPar
Kelchner, Plimpton, Hamilton, Phys Rev B, 58, 11085 (1998).

\sphinxlineitem{\sphinxstylestrong{(Khrapak)}}
\sphinxAtStartPar
Khrapak, Chaudhuri, and Morfill, J Chem Phys, 134, 054120 (2011).

\sphinxlineitem{\sphinxstylestrong{(Kim)}}
\sphinxAtStartPar
Kim, Keyes, Straub, J Chem. Phys, 132, 224107 (2010).

\sphinxlineitem{\sphinxstylestrong{(Klapp)}}
\sphinxAtStartPar
Klapp, Schoen, J Chem Phys, 117, 8050 (2002).

\sphinxlineitem{\sphinxstylestrong{(Kolafa)}}
\sphinxAtStartPar
Kolafa and Perram, Molecular Simulation, 9, 351 (1992).

\sphinxlineitem{\sphinxstylestrong{(Kolmogorov)}}
\sphinxAtStartPar
A. N. Kolmogorov, V. H. Crespi, Phys. Rev. B 71, 235415 (2005).

\sphinxlineitem{\sphinxstylestrong{(Kong)}}
\sphinxAtStartPar
L.T. Kong, G. Bartels, C. Campana, C. Denniston, and Martin H. Muser, \sphinxhref{https://doi.org/10.1016/j.cpc.2008.12.035}{Computer Physics Communications {[}180{]}(6):1004\sphinxhyphen{}1010 (2009).}

\sphinxlineitem{\sphinxstylestrong{(Kong2011)}}
\sphinxAtStartPar
L.T. Kong, \sphinxhref{https://doi.org/10.1016/j.cpc.2011.04.019}{Computer Physics Communications {[}182{]}(10):2201\sphinxhyphen{}2207, (2011).}

\sphinxlineitem{\sphinxstylestrong{(Kremer)}}
\sphinxAtStartPar
Kremer, Grest, J Chem Phys, 92, 5057 (1990).

\sphinxlineitem{\sphinxstylestrong{(Kuhn and Bagi, 2005)}}
\sphinxAtStartPar
Kuhn, M. R., \& Bagi, K. (2004). Contact rolling and deformation in granular media.  International journal of solids and structures, 41(21), 5793\sphinxhyphen{}5820.

\sphinxlineitem{\sphinxstylestrong{(Kumagai)}}
\sphinxAtStartPar
T. Kumagai, S. Izumi, S. Hara, S. Sakai, Comp. Mat. Science, 39, 457 (2007).

\sphinxlineitem{\sphinxstylestrong{(Kumar)}}
\sphinxAtStartPar
Kumar and Higdon, Phys Rev E, 82, 051401 (2010).

\sphinxlineitem{\sphinxstylestrong{(Kumar)}}
\sphinxAtStartPar
Kumar and Skinner, J. Phys. Chem. B, 112, 8311 (2008)

\sphinxlineitem{\sphinxstylestrong{(Lafourcade)}}
\sphinxAtStartPar
Lafourcade, Maillet, Denoual, Duval, Allera, Goryaeva, and Marinica, \sphinxhref{https://doi.org/10.1016/j.commatsci.2023.112534}{Comp. Mat. Science, 230, 112534 (2023)}

\sphinxlineitem{\sphinxstylestrong{(Lamoureux and Roux)}}
\sphinxAtStartPar
G. Lamoureux, B. Roux, J. Chem. Phys 119, 3025 (2003)

\sphinxlineitem{\sphinxstylestrong{(Lamoureux)}}
\sphinxAtStartPar
Lamoureux and Roux, J Chem Phys, 119, 3025\sphinxhyphen{}3039 (2003).

\sphinxlineitem{\sphinxstylestrong{(Landsgesell)}}\begin{enumerate}
\sphinxsetlistlabels{\Alph}{enumi}{enumii}{}{.}%
\setcounter{enumi}{9}
\item {} 
\sphinxAtStartPar
Landsgesell, P. Hebbeker, O. Rud, R. Lunkad, P. Kosovan, and C. Holm, Macromolecules 53, 3007\sphinxhyphen{}3020 (2020).

\end{enumerate}

\sphinxlineitem{\sphinxstylestrong{(Larentzos1)}}
\sphinxAtStartPar
J.P. Larentzos, J.K. Brennan, J.D. Moore, M. Lisal and W.D. Mattson, Comput. Phys. Commun., 185, 1987\sphinxhyphen{}1998 (2014).

\sphinxlineitem{\sphinxstylestrong{(Larentzos2)}}
\sphinxAtStartPar
J.P. Larentzos, J.K. Brennan, J.D. Moore, and W.D. Mattson, ARL\sphinxhyphen{}TR\sphinxhyphen{}6863, U.S. Army Research Laboratory, Aberdeen Proving Ground, MD (2014).

\sphinxlineitem{\sphinxstylestrong{(Larsen)}}
\sphinxAtStartPar
Larsen, Schmidt, Schiotz, Modelling Simul Mater Sci Eng, 24, 055007 (2016).

\sphinxlineitem{\sphinxstylestrong{(Lebedeva1)}}
\sphinxAtStartPar
I. V. Lebedeva, A. A. Knizhnik, A. M. Popov, Y. E. Lozovik, B. V. Potapkin, Phys. Rev. B, 84, 245437 (2011)

\sphinxlineitem{\sphinxstylestrong{(Lebedeva2)}}\begin{enumerate}
\sphinxsetlistlabels{\Roman}{enumi}{enumii}{}{.}%
\item {} \begin{enumerate}
\sphinxsetlistlabels{\Alph}{enumii}{enumiii}{}{.}%
\setcounter{enumii}{21}
\item {} 
\sphinxAtStartPar
Lebedeva, A. A. Knizhnik, A. M. Popov, Y. E. Lozovik, B. V. Potapkin, Physica E: 44, 949\sphinxhyphen{}954 (2012)

\end{enumerate}

\end{enumerate}

\sphinxlineitem{\sphinxstylestrong{(Lechman)}}
\sphinxAtStartPar
Lechman, et al, in preparation (2010).

\sphinxlineitem{\sphinxstylestrong{(Lee)}}
\sphinxAtStartPar
Lee, Baskes, Phys. Rev. B, 62, 8564\sphinxhyphen{}8567 (2000).

\sphinxlineitem{\sphinxstylestrong{(Lee2)}}
\sphinxAtStartPar
Lee, Baskes, Kim, Cho.  Phys. Rev. B, 64, 184102 (2001).

\sphinxlineitem{\sphinxstylestrong{(Lee2020)}}
\sphinxAtStartPar
C.W. Lee, et al. (2020) Physical Review B, 102(2), 024107.

\sphinxlineitem{\sphinxstylestrong{(Lenart)}}
\sphinxAtStartPar
Lenart , Jusufi, and Panagiotopoulos, J Chem Phys, 126, 044509 (2007).

\sphinxlineitem{\sphinxstylestrong{(Lenosky)}}
\sphinxAtStartPar
Lenosky, Sadigh, Alonso, Bulatov, de la Rubia, Kim, Voter, Kress, Modelling Simulation Materials Science Engineering, 8, 825 (2000).

\sphinxlineitem{\sphinxstylestrong{(Leven1)}}
\sphinxAtStartPar
I. Leven, I. Azuri, L. Kronik and O. Hod, J. Chem. Phys. 140, 104106 (2014).

\sphinxlineitem{\sphinxstylestrong{(Leven2)}}
\sphinxAtStartPar
I. Leven et al, J. Chem.Theory Comput. 12, 2896\sphinxhyphen{}905 (2016).

\sphinxlineitem{\sphinxstylestrong{(Li2013\_POF)}}
\sphinxAtStartPar
Li, Hu, Wang, Ma, Zhou, Phys Fluids, 25: 072103 (2013). DOI:10.1063/1.4812366.

\sphinxlineitem{\sphinxstylestrong{(Li2014\_JCP)}}
\sphinxAtStartPar
Li, Tang, Lei, Caswell, Karniadakis, J Comput Phys, 265: 113\sphinxhyphen{}127 (2014).  DOI:10.1016/j.jcp.2014.02.003.

\sphinxlineitem{\sphinxstylestrong{(Li2015\_CC)}}
\sphinxAtStartPar
Li, Tang, Li, Karniadakis, Chem Commun, 51: 11038\sphinxhyphen{}11040 (2015).  DOI:10.1039/C5CC01684C.

\sphinxlineitem{\sphinxstylestrong{(Li2015\_JCP)}}
\sphinxAtStartPar
Li, Yazdani, Tartakovsky, Karniadakis, J Chem Phys, 143: 014101 (2015).  DOI:10.1063/1.4923254.

\sphinxlineitem{\sphinxstylestrong{(Liang)}}
\sphinxAtStartPar
Liang, Phillpot, Sinnott Phys. Rev. B79 245110, (2009), Erratum: Phys. Rev. B85 199903(E), (2012)

\sphinxlineitem{\sphinxstylestrong{(Lisal)}}
\sphinxAtStartPar
M. Lisal, J.K. Brennan, J. Bonet Avalos, J. Chem. Phys., 135, 204105 (2011).

\sphinxlineitem{\sphinxstylestrong{(Liu1)}}
\sphinxAtStartPar
L. Liu, Y. Liu, S. V. Zybin, H. Sun and W. A. Goddard, Journal of Physical Chemistry A, 115, 11016\sphinxhyphen{}11022 (2011).

\sphinxlineitem{\sphinxstylestrong{(Liu2)}}
\sphinxAtStartPar
Liu, Bryantsev, Diallo, Goddard III, J. Am. Chem. Soc 131 (8) 2798 (2009)

\sphinxlineitem{\sphinxstylestrong{(Los and Fasolino)}}
\sphinxAtStartPar
J. H. Los and A. Fasolino, Phys. Rev. B 68, 024107 (2003).

\sphinxlineitem{\sphinxstylestrong{(Los2017)}}
\sphinxAtStartPar
J. H. Los et al. “Extended Tersoff potential for boron nitride: Energetics and elastic properties of pristine and defective h\sphinxhyphen{}BN”, Phys. Rev. B 96 (184108), 2017.

\sphinxlineitem{\sphinxstylestrong{(Luding, 2008)}}
\sphinxAtStartPar
Luding, S. (2008). Cohesive, frictional powders: contact models for tension. Granular matter, 10(4), 235.

\sphinxlineitem{\sphinxstylestrong{(Lysogorskiy)}}
\sphinxAtStartPar
Lysogorskiy, van der Oord, Bochkarev, Menon, Rinaldi, Hammerschmidt, Mrovec, Thompson, Csanyi, Ortner, Drautz, npj Comp Mat, 7, 97 (2021).

\sphinxlineitem{\sphinxstylestrong{(Lysogorskiy21)}}
\sphinxAtStartPar
Lysogorskiy, van der Oord, Bochkarev, Menon, Rinaldi, Hammerschmidt, Mrovec, Thompson, Csanyi, Ortner, Drautz, npj Comp Mat, 7, 97 (2021).

\sphinxlineitem{\sphinxstylestrong{(Lysogorskiy23)}}
\sphinxAtStartPar
Lysogorskiy, Bochkarev, Mrovec, Drautz, Phys Rev Mater, 7, 043801 (2023) / arXiv:2212.08716 (2022).

\sphinxlineitem{\sphinxstylestrong{(Maaravi)}}
\sphinxAtStartPar
T. Maaravi et al, J. Phys. Chem. C 121, 22826\sphinxhyphen{}22835 (2017).

\sphinxlineitem{\sphinxstylestrong{(MacKerell)}}
\sphinxAtStartPar
MacKerell, Bashford, Bellott, Dunbrack, Evanseck, Field, Fischer, Gao, Guo, Ha, et al, J Phys Chem B, 102, 3586 (1998).

\sphinxlineitem{\sphinxstylestrong{(Mackay and Denniston)}}
\sphinxAtStartPar
Mackay, F. E., and Denniston, C., Coupling MD particles to a lattice\sphinxhyphen{}Boltzmann fluid through the use of conservative forces, J. Comput. Phys. 237 (2013) 289\sphinxhyphen{}298.

\sphinxlineitem{\sphinxstylestrong{(Mackay et al.)}}
\sphinxAtStartPar
Mackay, F. E., Ollila, S.T.T., and Denniston, C., Hydrodynamic Forces Implemented into LAMMPS through a lattice\sphinxhyphen{}Boltzmann fluid, Computer Physics Communications 184 (2013) 2021\sphinxhyphen{}2031.

\sphinxlineitem{\sphinxstylestrong{(Magda)}}
\sphinxAtStartPar
Magda, Tirrell, Davis, J Chem Phys, 83, 1888\sphinxhyphen{}1901 (1985); erratum in JCP 84, 2901 (1986).

\sphinxlineitem{\sphinxstylestrong{(Maginn)}}
\sphinxAtStartPar
Kelkar, Rafferty, Maginn, Siepmann, Fluid Phase Equilibria, 260, 218\sphinxhyphen{}231 (2007).

\sphinxlineitem{\sphinxstylestrong{(Mahoney)}}
\sphinxAtStartPar
Mahoney, Jorgensen, J Chem Phys 112, 8910 (2000)

\sphinxlineitem{\sphinxstylestrong{(Malolepsza)}}
\sphinxAtStartPar
Malolepsza, Secor, Keyes, J Phys Chem B 119 (42), 13379\sphinxhyphen{}13384 (2015).

\sphinxlineitem{\sphinxstylestrong{(Mandadapu)}}
\sphinxAtStartPar
Mandadapu, KK; Templeton, JA; Lee, JW, “Polarization as a field variable from molecular dynamics simulations.” Journal of Chemical Physics (2013), 139:054115.  Please refer to the standard finite element (FE) texts, e.g. T.J.R Hughes “ The finite element method “, Dover 2003, for the basics of FE simulation.

\sphinxlineitem{\sphinxstylestrong{(Mandelli\_1)}}\begin{enumerate}
\sphinxsetlistlabels{\Alph}{enumi}{enumii}{}{.}%
\setcounter{enumi}{3}
\item {} 
\sphinxAtStartPar
Mandelli, W. Ouyang, M. Urbakh, and O. Hod, ACS Nano 13(7), 7603\sphinxhyphen{}7609 (2019).

\end{enumerate}

\sphinxlineitem{\sphinxstylestrong{(Maras)}}
\sphinxAtStartPar
Maras, Trushin, Stukowski, Ala\sphinxhyphen{}Nissila, Jonsson, Comp Phys Comm, 205, 13\sphinxhyphen{}21 (2016).

\sphinxlineitem{\sphinxstylestrong{(Marrink)}}
\sphinxAtStartPar
Marrink, de Vries, Mark, J Phys Chem B, 108, 750\sphinxhyphen{}760 (2004).

\sphinxlineitem{\sphinxstylestrong{(Marshall, 2009)}}
\sphinxAtStartPar
Marshall, J. S. (2009). Discrete\sphinxhyphen{}element modeling of particulate aerosol flows.  Journal of Computational Physics, 228(5), 1541\sphinxhyphen{}1561.

\sphinxlineitem{\sphinxstylestrong{(Martyna1992)}}
\sphinxAtStartPar
Martyna, Klein, Tuckerman, J Chem Phys, 97, 2635 (1992); Martyna, Tuckerman, Tobias, Klein, Mol Phys, 87, 1117.

\sphinxlineitem{\sphinxstylestrong{(Martyna1994)}}
\sphinxAtStartPar
Martyna, Tobias and Klein, J Chem Phys, 101, 4177 (1994).

\sphinxlineitem{\sphinxstylestrong{(Martyna2)}}\begin{enumerate}
\sphinxsetlistlabels{\Alph}{enumi}{enumii}{}{.}%
\setcounter{enumi}{6}
\item {} 
\sphinxAtStartPar
Martyna, A. Hughes, M. Tuckerman, J. Chem. Phys. 110, 3275 (1999).

\end{enumerate}

\sphinxlineitem{\sphinxstylestrong{(Mason)}}
\sphinxAtStartPar
J. K. Mason, Acta Cryst A65, 259 (2009).

\sphinxlineitem{\sphinxstylestrong{(Mattice)}}
\sphinxAtStartPar
Mattice, Suter, Conformational Theory of Large Molecules, Wiley, New York, 1994.

\sphinxlineitem{\sphinxstylestrong{(Maxwell)}}
\sphinxAtStartPar
J.C. Maxwell, Philos. Tans. Royal Soc. London, 157: 49\sphinxhyphen{}88 (1867).

\sphinxlineitem{\sphinxstylestrong{(Mayergoyz)}}
\sphinxAtStartPar
I.D. Mayergoyz, G. Bertotti, C. Serpico (2009). Elsevier (2009)

\sphinxlineitem{\sphinxstylestrong{(Mayo)}}
\sphinxAtStartPar
Mayo, Olfason, Goddard III, J Phys Chem, 94, 8897\sphinxhyphen{}8909 (1990).

\sphinxlineitem{\sphinxstylestrong{(Mees)}}
\sphinxAtStartPar
M. J. Mees, G. Pourtois, E. C. Neyts, B. J. Thijsse, and A. Stesmans, Phys. Rev. B 85, 134301 (2012).

\sphinxlineitem{\sphinxstylestrong{(Mei)}}
\sphinxAtStartPar
Mei, Davenport, Fernando, Phys Rev B, 43 4653 (1991)

\sphinxlineitem{\sphinxstylestrong{(Melchor)}}
\sphinxAtStartPar
Gonzalez\sphinxhyphen{}Melchor, Mayoral, Velazquez, and Alejandre, J Chem Phys, 125, 224107 (2006).

\sphinxlineitem{\sphinxstylestrong{(Meloni)}}
\sphinxAtStartPar
Meloni, Rosati and Colombo, J Chem Phys, 126, 121102 (2007).

\sphinxlineitem{\sphinxstylestrong{(Meremianin)}}
\sphinxAtStartPar
Meremianin, J. Phys. A,  39, 3099 (2006).

\sphinxlineitem{\sphinxstylestrong{(Mezei)}}
\sphinxAtStartPar
Mezei, J Chem Phys, 86, 7084 (1987)

\sphinxlineitem{\sphinxstylestrong{(Mickel)}}
\sphinxAtStartPar
W. Mickel, S. C. Kapfer, G. E. Schroeder\sphinxhyphen{}Turkand, K. Mecke, J. Chem. Phys. 138, 044501 (2013).

\sphinxlineitem{\sphinxstylestrong{(Mie)}}
\sphinxAtStartPar
G. Mie, Ann Phys, 316, 657 (1903).

\sphinxlineitem{\sphinxstylestrong{(Milano)}}\begin{enumerate}
\sphinxsetlistlabels{\Alph}{enumi}{enumii}{}{.}%
\setcounter{enumi}{6}
\item {} 
\sphinxAtStartPar
Milano, S. Goudeau, F. Mueller\sphinxhyphen{}Plathe, J. Polym. Sci. B Polym. Phys. 43, 871 (2005).

\end{enumerate}

\sphinxlineitem{\sphinxstylestrong{(Miller1)}}
\sphinxAtStartPar
T. F. Miller III, M. Eleftheriou, P. Pattnaik, A. Ndirango, G. J. Martyna, J. Chem. Phys., 116, 8649\sphinxhyphen{}8659 (2002).

\sphinxlineitem{\sphinxstylestrong{(Miller2)}}
\sphinxAtStartPar
Miller, Tadmor, Gibson, Bernstein and Pavia, J Chem Phys, 144, 184107 (2016).

\sphinxlineitem{\sphinxstylestrong{(Minary)}}
\sphinxAtStartPar
Minary, Martyna, and Tuckerman, J Chem Phys, 18, 2510 (2003).

\sphinxlineitem{\sphinxstylestrong{(Mindlin and Deresiewicz, 1953)}}
\sphinxAtStartPar
Mindlin, R.D., \& Deresiewicz, H (1953). Elastic Spheres in Contact under Varying Oblique Force. J. Appl. Mech., ASME 20, 327\sphinxhyphen{}344.

\sphinxlineitem{\sphinxstylestrong{(Mindlin, 1949)}}
\sphinxAtStartPar
Mindlin, R. D. (1949). Compliance of elastic bodies in contact.  J. Appl. Mech., ASME 16, 259\sphinxhyphen{}268.

\sphinxlineitem{\sphinxstylestrong{(Miron)}}
\sphinxAtStartPar
R. A. Miron and K. A. Fichthorn, J Chem Phys, 119, 6210 (2003).

\sphinxlineitem{\sphinxstylestrong{(Mishin)}}
\sphinxAtStartPar
Mishin, Mehl, and Papaconstantopoulos, Acta Mater, 53, 4029 (2005).

\sphinxlineitem{\sphinxstylestrong{(Mitchell and Fincham)}}
\sphinxAtStartPar
Mitchell, Fincham, J Phys Condensed Matter, 5, 1031\sphinxhyphen{}1038 (1993).

\sphinxlineitem{\sphinxstylestrong{(Mitchell2011)}}
\sphinxAtStartPar
Mitchell. A non\sphinxhyphen{}local, ordinary\sphinxhyphen{}state\sphinxhyphen{}based viscoelasticity model for peridynamics. Sandia National Lab Report, 8064:1\sphinxhyphen{}28 (2011).

\sphinxlineitem{\sphinxstylestrong{(Mitchell2011a)}}
\sphinxAtStartPar
Mitchell. A Nonlocal, Ordinary, State\sphinxhyphen{}Based Plasticity Model for Peridynamics. Sandia National Lab Report, 3166:1\sphinxhyphen{}34 (2011).

\sphinxlineitem{\sphinxstylestrong{(Miyazaki)}}
\sphinxAtStartPar
Miyazaki, Okazaki, Shinoda, J Chem Theory Comput, 16, 782\sphinxhyphen{}793 (2020).

\sphinxlineitem{\sphinxstylestrong{(Mniszewski)}}
\sphinxAtStartPar
S. M. Mniszewski, M. J. Cawkwell, M. E. Wall, J. Mohd\sphinxhyphen{}Yusof, N. Bock, T. C.  Germann, and A. M. N. Niklasson, J. Chem. Theory Comput., 11, 4644 (2015).

\sphinxlineitem{\sphinxstylestrong{(Monaghan)}}
\sphinxAtStartPar
Monaghan and Gingold, Journal of Computational Physics, 52, 374\sphinxhyphen{}389 (1983).

\sphinxlineitem{\sphinxstylestrong{(Monti)}}
\sphinxAtStartPar
Monti, Clemmer, Srivastava, Silbert, Grest, and Lechman, Phys. Rev. E, (2022).

\sphinxlineitem{\sphinxstylestrong{(Moore)}}
\sphinxAtStartPar
Moore, J Chem Phys, 144, 104501 (2016).

\sphinxlineitem{\sphinxstylestrong{(Mori)}}
\sphinxAtStartPar
Y. Mori, Y. Okamoto, J. Phys. Soc. Jpn., 7, 074003 (2010).

\sphinxlineitem{\sphinxstylestrong{(Moriarty1)}}
\sphinxAtStartPar
Moriarty, Physical Review B, 38, 3199 (1988).

\sphinxlineitem{\sphinxstylestrong{(Moriarty2)}}
\sphinxAtStartPar
Moriarty, Physical Review B, 42, 1609 (1990). Moriarty, Physical Review B 49, 12431 (1994).

\sphinxlineitem{\sphinxstylestrong{(Moriarty3)}}
\sphinxAtStartPar
Moriarty, Benedict, Glosli, Hood, Orlikowski, Patel, Soderlind, Streitz, Tang, and Yang, Journal of Materials Research, 21, 563 (2006).

\sphinxlineitem{\sphinxstylestrong{(Morris)}}
\sphinxAtStartPar
Morris, Fox, Zhu, J Comp Physics, 136, 214\sphinxhyphen{}226 (1997).

\sphinxlineitem{\sphinxstylestrong{(Moustafa)}}
\sphinxAtStartPar
Sabry G. Moustafa, Andrew J. Schultz, and David A. Kofke, \sphinxhref{https://link.aps.org/doi/10.1103/PhysRevE.92.043303}{Phys. Rev. E {[}92{]}, 043303 (2015)}

\sphinxlineitem{\sphinxstylestrong{(Muller\sphinxhyphen{}Plathe1)}}
\sphinxAtStartPar
Muller\sphinxhyphen{}Plathe, J Chem Phys, 106, 6082 (1997).

\sphinxlineitem{\sphinxstylestrong{(Muller\sphinxhyphen{}Plathe2)}}
\sphinxAtStartPar
Muller\sphinxhyphen{}Plathe, Phys Rev E, 59, 4894\sphinxhyphen{}4898 (1999).

\sphinxlineitem{\sphinxstylestrong{(Murdick)}}
\sphinxAtStartPar
D.A. Murdick, X.W. Zhou, H.N.G. Wadley, D. Nguyen\sphinxhyphen{}Manh, R. Drautz, and D.G. Pettifor, Phys. Rev. B, 73, 45206 (2006).

\sphinxlineitem{\sphinxstylestrong{(Murty)}}
\sphinxAtStartPar
M.V.R. Murty, H.A. Atwater, Phys Rev B, 51, 4889 (1995).

\sphinxlineitem{\sphinxstylestrong{(Nakano)}}
\sphinxAtStartPar
A. Nakano, Computer Physics Communications, 104, 59\sphinxhyphen{}69 (1997).

\sphinxlineitem{\sphinxstylestrong{(Neelov)}}
\sphinxAtStartPar
Neelov, Holm, J Chem Phys 132, 234103 (2010)

\sphinxlineitem{\sphinxstylestrong{(Nelson)}}
\sphinxAtStartPar
Nelson, Halperin, Phys Rev B, 19, 2457 (1979).

\sphinxlineitem{\sphinxstylestrong{(Nettleton)}}
\sphinxAtStartPar
Nettleton and Green, J Chem Phys, 29, 6 (1958).

\sphinxlineitem{\sphinxstylestrong{(Neyts)}}
\sphinxAtStartPar
E. C. Neyts and A. Bogaerts, Theor. Chem. Acc. 132, 1320 (2013).

\sphinxlineitem{\sphinxstylestrong{(Nguyen2023)}}
\sphinxAtStartPar
Nguyen, Physical Review B, 107(14), 144103, (2023).

\sphinxlineitem{\sphinxstylestrong{(Nguyen2024)}}
\sphinxAtStartPar
Nguyen, Journal of Computational Physics, 113102, (2024).

\sphinxlineitem{\sphinxstylestrong{(Nguyen and Rohskopf)}}
\sphinxAtStartPar
Nguyen and Rohskopf,  Journal of Computational Physics, 480, 112030, (2023).

\sphinxlineitem{\sphinxstylestrong{(Nguyen and Sema)}}
\sphinxAtStartPar
Nguyen and Sema, \sphinxurl{https://arxiv.org/abs/2405.00306}, (2024).

\sphinxlineitem{\sphinxstylestrong{(NguyenTD)}}
\sphinxAtStartPar
Nguyen, Li, Bagchi, Solis, Olvera de la Cruz, Comput Phys Commun 241, 80\sphinxhyphen{}19 (2019)

\sphinxlineitem{\sphinxstylestrong{(Nicholson and Rutledge)}}
\sphinxAtStartPar
Nicholson and Rutledge, J Chem Phys, 145, 244903 (2016).

\sphinxlineitem{\sphinxstylestrong{(Niklasson2002)}}
\sphinxAtStartPar
A. M. N. Niklasson, Phys. Rev. B, 66, 155115 (2002).

\sphinxlineitem{\sphinxstylestrong{(Niklasson2008)}}
\sphinxAtStartPar
A. M. N. Niklasson, Phys. Rev. Lett., 100, 123004 (2008).

\sphinxlineitem{\sphinxstylestrong{(Niklasson2014)}}
\sphinxAtStartPar
A. M. N. Niklasson and M. Cawkwell, J. Chem. Phys., 141, 164123, (2014).

\sphinxlineitem{\sphinxstylestrong{(Niklasson2017)}}
\sphinxAtStartPar
A. M. N. Niklasson, J. Chem. Phys., 147, 054103 (2017).

\sphinxlineitem{\sphinxstylestrong{(Nitol)}}
\sphinxAtStartPar
Nitol, Dickel, and Barrett, Computational Materials Science 188 (2021): 110207.

\sphinxlineitem{\sphinxstylestrong{(Noid)}}
\sphinxAtStartPar
Noid, Chu, Ayton, Krishna, Izvekov, Voth, Das, Andersen, J Chem Phys 128, 134105 (2008).

\sphinxlineitem{\sphinxstylestrong{(Nordlund95)}}
\sphinxAtStartPar
Nordlund, Kai. Computational materials science 3.4 (1995): 448\sphinxhyphen{}456.

\sphinxlineitem{\sphinxstylestrong{(Nordlund98)}}
\sphinxAtStartPar
Nordlund, Kai, et al.  Physical Review B 57.13 (1998): 7556.

\sphinxlineitem{\sphinxstylestrong{(Norman)}}
\sphinxAtStartPar
G E Norman, S V Starikov, V V Stegailov et al., Contrib. Plasma Phys., 53, 129\sphinxhyphen{}139 (2013).

\sphinxlineitem{\sphinxstylestrong{(Noskov)}}
\sphinxAtStartPar
Noskov, Lamoureux and Roux, J Phys Chem B, 109, 6705 (2005).

\sphinxlineitem{\sphinxstylestrong{(Nurdin)}}
\sphinxAtStartPar
Nurdin and Schotte Phys Rev E, 61(4), 3579 (2000)

\sphinxlineitem{\sphinxstylestrong{(O’Connor)}}
\sphinxAtStartPar
O’Connor et al., J. Chem. Phys. 142, 024903 (2015).

\sphinxlineitem{\sphinxstylestrong{(O’Hearn)}}
\sphinxAtStartPar
O’Hearn, Alperen, Aktulga, SIAM J. Sci. Comput., 42(1), C1\textendash{}C22 (2020).

\sphinxlineitem{\sphinxstylestrong{(Okabe)}}
\sphinxAtStartPar
T. Okabe, M. Kawata, Y. Okamoto, M. Masuhiro, Chem. Phys. Lett., 335, 435\sphinxhyphen{}439 (2001).

\sphinxlineitem{\sphinxstylestrong{(Ollila et al.)}}
\sphinxAtStartPar
Ollila, S.T.T., Denniston, C., Karttunen, M., and Ala\sphinxhyphen{}Nissila, T., Fluctuating lattice\sphinxhyphen{}Boltzmann model for complex fluids, J. Chem. Phys. 134 (2011) 064902.

\sphinxlineitem{\sphinxstylestrong{(Omelyan)}}
\sphinxAtStartPar
Omelyan, Mryglod, and Folk. Phys. Rev. Lett. 86(5), 898. (2001).

\end{description}

\sphinxAtStartPar
\sphinxstylestrong{(OPLS\sphinxhyphen{}AA96)}  Jorgensen, Maxwell, Tirado\sphinxhyphen{}Rives, J Am Chem Soc, 118(45), 11225\sphinxhyphen{}11236 (1996).
\begin{description}
\sphinxlineitem{\sphinxstylestrong{(Oppelstrup)}}
\sphinxAtStartPar
Oppelstrup, unpublished, 2015. Oppelstrup and Moriarty, to be published.

\sphinxlineitem{\sphinxstylestrong{(Orsi)}}
\sphinxAtStartPar
Orsi \& Essex, The ELBA force field for coarse\sphinxhyphen{}grain modeling of lipid membranes, PloS ONE 6(12): e28637, 2011.

\sphinxlineitem{\sphinxstylestrong{(Otis R. Walton)}}
\sphinxAtStartPar
Walton, O.R., Personal Communication

\sphinxlineitem{\sphinxstylestrong{(Ouldridge)}}
\sphinxAtStartPar
T.E. Ouldridge, A.A. Louis, J.P.K. Doye, J. Chem. Phys. 134, 085101 (2011).

\sphinxlineitem{\sphinxstylestrong{(Ouldridge\sphinxhyphen{}DPhil)}}
\sphinxAtStartPar
T.E. Ouldridge, Coarse\sphinxhyphen{}grained modelling of DNA and DNA self\sphinxhyphen{}assembly, DPhil. University of Oxford (2011).

\sphinxlineitem{\sphinxstylestrong{(Ouyang1)}}
\sphinxAtStartPar
W. Ouyang, D. Mandelli, M. Urbakh and O. Hod, Nano Lett. 18, 6009\sphinxhyphen{}6016 (2018).

\sphinxlineitem{\sphinxstylestrong{(Ouyang2)}}
\sphinxAtStartPar
W. Ouyang et al., J. Chem. Theory Comput. 16(1), 666\sphinxhyphen{}676 (2020).

\sphinxlineitem{\sphinxstylestrong{(Ouyang\_1)}}\begin{enumerate}
\sphinxsetlistlabels{\Alph}{enumi}{enumii}{}{.}%
\setcounter{enumi}{22}
\item {} 
\sphinxAtStartPar
Ouyang et al., J. Chem. Theory Comput. 16(1), 666\sphinxhyphen{}676 (2020).

\end{enumerate}

\sphinxlineitem{\sphinxstylestrong{(Ouyang6)}}\begin{enumerate}
\sphinxsetlistlabels{\Alph}{enumi}{enumii}{}{.}%
\setcounter{enumi}{22}
\item {} 
\sphinxAtStartPar
Ouyang, O. Hod, and R. Guerra, J. Chem. Theory Comput. 17, 7215 (2021).

\end{enumerate}

\sphinxlineitem{\sphinxstylestrong{(Ouyang7)}}\begin{enumerate}
\sphinxsetlistlabels{\Alph}{enumi}{enumii}{}{.}%
\setcounter{enumi}{22}
\item {} 
\sphinxAtStartPar
Ouyang, et al., J. Chem. Theory Comput. 17, 7237 (2021).

\end{enumerate}

\sphinxlineitem{\sphinxstylestrong{(Palkar)}}
\sphinxAtStartPar
Palkar V, Kuksenok O, J. Phys. Chem. B, 126 (1), 336\sphinxhyphen{}346, 2022

\sphinxlineitem{\sphinxstylestrong{(Paquay)}}
\sphinxAtStartPar
Paquay and Kusters, Biophys. J., 110, 6, (2016). preprint available at \sphinxhref{https://arxiv.org/abs/1411.3019/}{arXiv:1411.3019}.

\sphinxlineitem{\sphinxstylestrong{(Park)}}
\sphinxAtStartPar
Park, Schulten, J. Chem. Phys. 120 (13), 5946 (2004)

\sphinxlineitem{\sphinxstylestrong{(Parks)}}
\sphinxAtStartPar
Parks, Lehoucq, Plimpton, Silling, Comp Phys Comm, 179(11), 777\sphinxhyphen{}783 (2008).

\sphinxlineitem{\sphinxstylestrong{(Parrinello)}}
\sphinxAtStartPar
Parrinello and Rahman, J Appl Phys, 52, 7182 (1981).

\sphinxlineitem{\sphinxstylestrong{(PASS)}}
\sphinxAtStartPar
PASS webpage: \sphinxurl{https://www.sdu.dk/en/DPASS}

\sphinxlineitem{\sphinxstylestrong{(Paula Leite2016)}}
\sphinxAtStartPar
Paula Leite , Freitas, Azevedo, and de Koning, J Chem Phys, 126, 044509 (2016).

\sphinxlineitem{\sphinxstylestrong{(Paula Leite2017)}}
\sphinxAtStartPar
Paula Leite, Santos\sphinxhyphen{}Florez, and de Koning, Phys Rev E, 96, 32115 (2017).

\sphinxlineitem{\sphinxstylestrong{(Pavlov1)}}
\sphinxAtStartPar
D Pavlov, V Galigerov, D Kolotinskii, V Nikolskiy, V Stegailov, International Journal of High Performance Computing Applications, 38, 34\sphinxhyphen{}49 (2024).

\sphinxlineitem{\sphinxstylestrong{(Pavlov2)}}
\sphinxAtStartPar
Pavlov, Galigerov, Kolotinskii, Nikolskiy, Stegailov, “GPU\sphinxhyphen{}based Molecular Dynamics of Fluid Flows: Reaching for Turbulence”, Int. J. High Perf. Comp. Appl., (2024)

\sphinxlineitem{\sphinxstylestrong{(Pearlman)}}
\sphinxAtStartPar
Pearlman, J Chem Phys, 98, 1487 (1994)

\sphinxlineitem{\sphinxstylestrong{(Pedersen)}}
\sphinxAtStartPar
Pedersen, J. Chem. Phys., 139, 104102 (2013).

\sphinxlineitem{\sphinxstylestrong{(Pedone)}}\begin{enumerate}
\sphinxsetlistlabels{\Alph}{enumi}{enumii}{}{.}%
\item {} 
\sphinxAtStartPar
Pedone, G. Malavasi, M. C. Menziani, A. N. Cormack, and U. Segre, J. Phys. Chem. B, 110, 11780 (2006)

\end{enumerate}

\sphinxlineitem{\sphinxstylestrong{(Peng)}}
\sphinxAtStartPar
Peng, Ren, Dudarev, Whelan, Acta Crystallogr. A, 52, 257\sphinxhyphen{}76 (1996).

\sphinxlineitem{\sphinxstylestrong{(Perram)}}
\sphinxAtStartPar
Perram and Rasmussen, Phys Rev E, 54, 6565\sphinxhyphen{}6572 (1996).

\sphinxlineitem{\sphinxstylestrong{(Petersen)}}
\sphinxAtStartPar
Petersen, J Chem Phys, 103, 3668 (1995).

\sphinxlineitem{\sphinxstylestrong{(Petersen)}}
\sphinxAtStartPar
Petersen, Lechman, Plimpton, Grest, in’ t Veld, Schunk, J Chem Phys, 132, 174106 (2010).

\sphinxlineitem{\sphinxstylestrong{(Pettifor\_1)}}
\sphinxAtStartPar
D.G. Pettifor and I.I. Oleinik, Phys. Rev. B, 59, 8487 (1999).

\sphinxlineitem{\sphinxstylestrong{(Pettifor\_2)}}
\sphinxAtStartPar
D.G. Pettifor and I.I. Oleinik, Phys. Rev. Lett., 84, 4124 (2000).

\sphinxlineitem{\sphinxstylestrong{(Pettifor\_3)}}
\sphinxAtStartPar
D.G. Pettifor and I.I. Oleinik, Phys. Rev. B, 65, 172103 (2002).

\sphinxlineitem{\sphinxstylestrong{(PFC)}}
\sphinxAtStartPar
PFC Particle Flow Code 6.0 Documentation. Itasca Consulting Group.

\sphinxlineitem{\sphinxstylestrong{(Phillips)}}
\sphinxAtStartPar
C. L. Phillips, J. A. Anderson, S. C. Glotzer, Comput Phys Comm, 230, 7191\sphinxhyphen{}7201 (2011).

\sphinxlineitem{\sphinxstylestrong{(Piaggi)}}
\sphinxAtStartPar
Piaggi and Parrinello, J Chem Phys, 147, 114112 (2017).

\sphinxlineitem{\sphinxstylestrong{(Pisarev)}}
\sphinxAtStartPar
V V Pisarev and S V Starikov, J. Phys.: Condens. Matter, 26, 475401 (2014).

\sphinxlineitem{\sphinxstylestrong{(Plimpton)}}
\sphinxAtStartPar
Plimpton and Knight, JPDC, 147, 184\sphinxhyphen{}195 (2021).

\sphinxlineitem{\sphinxstylestrong{(PLUMED)}}
\sphinxAtStartPar
G.A. Tribello, M. Bonomi, D. Branduardi, C. Camilloni and G. Bussi, Comp. Phys. Comm 185, 604 (2014)

\sphinxlineitem{\sphinxstylestrong{(Pollock)}}
\sphinxAtStartPar
Pollock and Glosli, Comp Phys Comm, 95, 93 (1996).

\sphinxlineitem{\sphinxstylestrong{(Ponder)}}
\sphinxAtStartPar
Ponder, Wu, Ren, Pande, Chodera, Schnieders, Haque,  Mobley, Lambrecht, DiStasio Jr, M. Head\sphinxhyphen{}Gordon, Clark,  Johnson, T. Head\sphinxhyphen{}Gordon, J Phys Chem B, 114, 2549\sphinxhyphen{}2564 (2010).

\sphinxlineitem{\sphinxstylestrong{(Popov1)}}
\sphinxAtStartPar
A.M. Popov, I. V. Lebedeva, A. A. Knizhnik, Y. E. Lozovik and B. V. Potapkin, Chem. Phys. Lett. 536, 82\sphinxhyphen{}86 (2012).

\sphinxlineitem{\sphinxstylestrong{(Price1)}}
\sphinxAtStartPar
Price and Brooks, J Chem Phys, 121, 10096 (2004).

\sphinxlineitem{\sphinxstylestrong{(Price2)}}
\sphinxAtStartPar
Price, Stone and Alderton, Mol Phys, 52, 987 (1984).

\sphinxlineitem{\sphinxstylestrong{(QEq/Fire)}}
\sphinxAtStartPar
T.\sphinxhyphen{}R. Shan, A. P. Thompson, S. J. Plimpton, in preparation

\sphinxlineitem{\sphinxstylestrong{(Qi)}}
\sphinxAtStartPar
Qi and Reed, J. Phys. Chem. A 116, 10451 (2012).

\sphinxlineitem{\sphinxstylestrong{(Ramirez)}}
\sphinxAtStartPar
J. Ramirez, S.K. Sukumaran, B. Vorselaars and A.E. Likhtman, J. Chem. Phys. 133, 154103 (2010).

\sphinxlineitem{\sphinxstylestrong{(Rappe)}}
\sphinxAtStartPar
Rappe and Goddard III, Journal of Physical Chemistry, 95, 3358\sphinxhyphen{}3363 (1991).

\sphinxlineitem{\sphinxstylestrong{(Ravelo)}}
\sphinxAtStartPar
Ravelo, Holian, Germann and Lomdahl, Phys Rev B, 70, 014103 (2004).

\sphinxlineitem{\sphinxstylestrong{(ReaxFF)}}
\sphinxAtStartPar
A. C. T. van Duin, S. Dasgupta, F. Lorant, W. A. Goddard III, J Physical Chemistry, 105, 9396\sphinxhyphen{}9049 (2001)

\sphinxlineitem{\sphinxstylestrong{(Rector)}}
\sphinxAtStartPar
Rector, Van Swol, Henderson, Molecular Physics, 82, 1009 (1994).

\sphinxlineitem{\sphinxstylestrong{(Ree)}}
\sphinxAtStartPar
Ree, Journal of Chemical Physics, 73, 5401 (1980).

\sphinxlineitem{\sphinxstylestrong{(Reed)}}
\sphinxAtStartPar
Reed, Fried, and Joannopoulos, Phys. Rev. Lett., 90, 235503 (2003).

\sphinxlineitem{\sphinxstylestrong{(Reed2)}}
\sphinxAtStartPar
Reed, J. Phys. Chem. C, 116, 2205 (2012).

\sphinxlineitem{\sphinxstylestrong{(Rick)}}
\sphinxAtStartPar
S. W. Rick, S. J. Stuart, B. J. Berne, J Chem Phys 101, 16141 (1994).

\sphinxlineitem{\sphinxstylestrong{(Rick2)}}\begin{enumerate}
\sphinxsetlistlabels{\Alph}{enumi}{enumii}{}{.}%
\setcounter{enumi}{18}
\item {} \begin{enumerate}
\sphinxsetlistlabels{\Alph}{enumii}{enumiii}{}{.}%
\setcounter{enumii}{22}
\item {} 
\sphinxAtStartPar
Rick, S. J. Stuart, B. J. Berne, J Chem Phys 101, 6141

\end{enumerate}

\end{enumerate}

\sphinxlineitem{\sphinxstylestrong{(Roberts)}}\begin{enumerate}
\sphinxsetlistlabels{\Alph}{enumi}{enumii}{}{.}%
\setcounter{enumi}{17}
\item {} 
\sphinxAtStartPar
Roberts (2019) “Evenly Distributing Points in a Triangle.” Extreme Learning.  \sphinxurl{http://extremelearning.com.au/evenly-distributing-points-in-a-triangle/}

\end{enumerate}

\sphinxlineitem{\sphinxstylestrong{(Rohart)}}
\sphinxAtStartPar
Rohart and Thiaville, Physical Review B, 88(18), 184422. (2013).

\sphinxlineitem{\sphinxstylestrong{(Rosenberger)}}
\sphinxAtStartPar
Rosenberger, Sanyal, Shell and van der Vegt,  Journal of Chemical Physics, 2019, 151 (4), 044111.

\sphinxlineitem{\sphinxstylestrong{(Rubensson)}}
\sphinxAtStartPar
E. H. Rubensson, A. M. N. Niklasson, SIAM J. Sci. Comput. 36 (2), 147\sphinxhyphen{}170, (2014).

\sphinxlineitem{\sphinxstylestrong{(Rutherford)}}
\sphinxAtStartPar
A M Rutherford and D M Duffy, J. Phys.: Condens. Matter, 19, 496201\sphinxhyphen{}496210 (2007).

\sphinxlineitem{\sphinxstylestrong{(Ryckaert)}}
\sphinxAtStartPar
J.\sphinxhyphen{}P. Ryckaert, G. Ciccotti and H. J. C. Berendsen, J of Comp Phys, 23, 327\sphinxhyphen{}341 (1977).

\sphinxlineitem{\sphinxstylestrong{(SMTB\sphinxhyphen{}Q\_1)}}
\sphinxAtStartPar
N. Salles, O. Politano, E. Amzallag, R. Tetot, Comput. Mater. Sci. 111 (2016) 181\sphinxhyphen{}189

\sphinxlineitem{\sphinxstylestrong{(SMTB\sphinxhyphen{}Q\_2)}}
\sphinxAtStartPar
E. Maras, N. Salles, R. Tetot, T. Ala\sphinxhyphen{}Nissila, H. Jonsson, J. Phys. Chem. C 2015, 119, 10391\sphinxhyphen{}10399

\sphinxlineitem{\sphinxstylestrong{(SMTB\sphinxhyphen{}Q\_3)}}
\sphinxAtStartPar
R. Tetot, N. Salles, S. Landron, E. Amzallag, Surface Science 616, 19\sphinxhyphen{}8722 28 (2013)

\sphinxlineitem{\sphinxstylestrong{(SRIM)}}
\sphinxAtStartPar
SRIM webpage: \sphinxurl{http://www.srim.org/}

\sphinxlineitem{\sphinxstylestrong{(SW)}}
\sphinxAtStartPar
F. H. Stillinger, and T. A. Weber, Phys. Rev. B, 31, 5262 (1985).

\sphinxlineitem{\sphinxstylestrong{(SWM4\sphinxhyphen{}NDP)}}
\sphinxAtStartPar
Lamoureux, Harder, Vorobyov, Roux, MacKerell, Chem Phys Let, 418, 245\sphinxhyphen{}249 (2006)

\sphinxlineitem{\sphinxstylestrong{(Sadigh)}}
\sphinxAtStartPar
B Sadigh, P Erhart, A Stukowski, A Caro, E Martinez, and L Zepeda\sphinxhyphen{}Ruiz, Phys. Rev. B, 85, 184203 (2012).

\sphinxlineitem{\sphinxstylestrong{(Sadigh1)}}\begin{enumerate}
\sphinxsetlistlabels{\Alph}{enumi}{enumii}{}{.}%
\setcounter{enumi}{1}
\item {} 
\sphinxAtStartPar
Sadigh, P. Erhart, A. Stukowski, A. Caro, E. Martinez, and L. Zepeda\sphinxhyphen{}Ruiz, Phys. Rev. B \sphinxstylestrong{85}, 184203 (2012)

\end{enumerate}

\sphinxlineitem{\sphinxstylestrong{(Sadigh2)}}\begin{enumerate}
\sphinxsetlistlabels{\Alph}{enumi}{enumii}{}{.}%
\setcounter{enumi}{1}
\item {} 
\sphinxAtStartPar
Sadigh and P. Erhart, Phys. Rev. B \sphinxstylestrong{86}, 134204 (2012)

\end{enumerate}

\sphinxlineitem{\sphinxstylestrong{(Safran)}}
\sphinxAtStartPar
Safran, Statistical Thermodynamics of Surfaces, Interfaces, And Membranes, Westview Press, ISBN: 978\sphinxhyphen{}0813340791 (2003).

\sphinxlineitem{\sphinxstylestrong{(Salanne)}}
\sphinxAtStartPar
Salanne, Rotenberg, Jahn, Vuilleumier, Simon, Christian and Madden, Theor Chem Acc, 131, 1143 (2012).

\sphinxlineitem{\sphinxstylestrong{(Salerno)}}
\sphinxAtStartPar
Salerno, Bernstein, J Chem Theory Comput, \textendash{}, —\sphinxhyphen{} (2018).

\sphinxlineitem{\sphinxstylestrong{(Sanyal1)}}
\sphinxAtStartPar
Sanyal and Shell, Journal of Chemical Physics, 2016, 145 (3), 034109.

\sphinxlineitem{\sphinxstylestrong{(Sanyal2)}}
\sphinxAtStartPar
Sanyal and Shell, Journal of Physical Chemistry B, 122 (21), 5678\sphinxhyphen{}5693.

\sphinxlineitem{\sphinxstylestrong{(Scalfi)}}
\sphinxAtStartPar
Scalfi \sphinxstyleemphasis{et al.}, J. Chem. Phys., 153, 174704 (2020).

\sphinxlineitem{\sphinxstylestrong{(Schelling)}}
\sphinxAtStartPar
Patrick K. Schelling, Comp. Mat. Science, 44, 274 (2008).

\sphinxlineitem{\sphinxstylestrong{(Scherer1)}}\begin{enumerate}
\sphinxsetlistlabels{\Alph}{enumi}{enumii}{}{.}%
\setcounter{enumi}{2}
\item {} 
\sphinxAtStartPar
Scherer and D. Andrienko, Phys. Chem. Chem. Phys. 20, 22387\sphinxhyphen{}22394 (2018).

\end{enumerate}

\sphinxlineitem{\sphinxstylestrong{(Scherer2)}}\begin{enumerate}
\sphinxsetlistlabels{\Alph}{enumi}{enumii}{}{.}%
\setcounter{enumi}{2}
\item {} 
\sphinxAtStartPar
Scherer, R. Scheid, D. Andrienko, and T. Bereau, J. Chem. Theor. Comp. 16, 3194\sphinxhyphen{}3204 (2020).

\end{enumerate}

\sphinxlineitem{\sphinxstylestrong{(Schlitter1)}}
\sphinxAtStartPar
Schlitter, Swegat, Mulders, “Distance\sphinxhyphen{}type reaction coordinates for modelling activated processes”, J Molecular Modeling, 7, 171\sphinxhyphen{}177 (2001).

\sphinxlineitem{\sphinxstylestrong{(Schlitter2)}}
\sphinxAtStartPar
Schlitter and Klahn, “The free energy of a reaction coordinate at multiple constraints: a concise formulation”, Molecular Physics, 101, 3439\sphinxhyphen{}3443 (2003).

\sphinxlineitem{\sphinxstylestrong{(Schmid)}}
\sphinxAtStartPar
S. Bureekaew, S. Amirjalayer, M. Tafipolsky, C. Spickermann, T.K. Roy and R. Schmid, Phys. Status Solidi B, 6, 1128 (2013).

\sphinxlineitem{\sphinxstylestrong{(Schneider)}}
\sphinxAtStartPar
Schneider and Stoll, Phys Rev B, 17, 1302 (1978).

\sphinxlineitem{\sphinxstylestrong{(Schratt \& Mohles)}}
\sphinxAtStartPar
Schratt, Mohles. Comp. Mat. Sci. 182 (2020) 109774  ———\sphinxhyphen{}

\sphinxlineitem{\sphinxstylestrong{(Schroeder)}}
\sphinxAtStartPar
Schroeder and Steinhauser, J Chem Phys, 133, 154511 (2010).

\sphinxlineitem{\sphinxstylestrong{(Seleson 2010)}}
\sphinxAtStartPar
Seleson, Parks, Int J Mult Comp Eng 9(6), pp. 689\sphinxhyphen{}706, 2011.

\sphinxlineitem{\sphinxstylestrong{(Semaev)}}
\sphinxAtStartPar
Semaev, Cryptography and Lattices, 181 (2001).

\sphinxlineitem{\sphinxstylestrong{(Seo)}}
\sphinxAtStartPar
Seo, Shinoda, J Chem Theory Comput, 15, 762\sphinxhyphen{}774 (2019).

\sphinxlineitem{\sphinxstylestrong{(Sheppard)}}
\sphinxAtStartPar
Sheppard, Terrell, Henkelman, J Chem Phys, 128, 134106 (2008).  See ref 1 in this paper for original reference to Qmin in Jonsson, Mills, Jacobsen.

\sphinxlineitem{\sphinxstylestrong{(Shi)}}
\sphinxAtStartPar
Shi, Xia, Zhang, Best, Wu, Ponder, Ren, J Chem Theory Comp, 9, 4046, 2013.

\sphinxlineitem{\sphinxstylestrong{(Shinoda)}}
\sphinxAtStartPar
Shinoda, DeVane, Klein, Mol Sim, 33, 27 (2007).

\sphinxlineitem{\sphinxstylestrong{(Shinoda)}}
\sphinxAtStartPar
Shinoda, Shiga, and Mikami, Phys Rev B, 69, 134103 (2004).

\sphinxlineitem{\sphinxstylestrong{(Shire)}}
\sphinxAtStartPar
Shire, Hanley and Stratford, Comp. Part. Mech., (2020).

\sphinxlineitem{\sphinxstylestrong{(Sides)}}
\sphinxAtStartPar
Sides, Grest, Stevens, Plimpton, J Polymer Science B, 42, 199\sphinxhyphen{}208 (2004).

\sphinxlineitem{\sphinxstylestrong{(Siepmann)}}
\sphinxAtStartPar
Siepmann and Sprik, J. Chem. Phys. 102, 511 (1995).

\sphinxlineitem{\sphinxstylestrong{(Silbert)}}
\sphinxAtStartPar
Silbert, Ertas, Grest, Halsey, Levine, Plimpton, Phys Rev E, 64, p 051302 (2001).

\sphinxlineitem{\sphinxstylestrong{(Silbert, 2001)}}
\sphinxAtStartPar
Silbert, L. E., Ertas, D., Grest, G. S., Halsey, T. C., Levine, D., \& Plimpton, S. J. (2001).  Granular flow down an inclined plane: Bagnold scaling and rheology. Physical Review E,

\sphinxlineitem{\sphinxstylestrong{(Silling 2000)}}
\sphinxAtStartPar
Silling, J Mech Phys Solids, 48, 175\sphinxhyphen{}209 (2000).

\sphinxlineitem{\sphinxstylestrong{(Silling 2005)}}
\sphinxAtStartPar
Silling Askari, Computer and Structures, 83, 1526\sphinxhyphen{}1535 (2005).

\sphinxlineitem{\sphinxstylestrong{(Silling 2007)}}
\sphinxAtStartPar
Silling, Epton, Weckner, Xu, Askari, J Elasticity, 88, 151\sphinxhyphen{}184 (2007).

\sphinxlineitem{\sphinxstylestrong{(Singh)}}
\sphinxAtStartPar
Singh and Warner, Acta Mater, 58, 5797\sphinxhyphen{}5805 (2010),

\sphinxlineitem{\sphinxstylestrong{(Singraber, Behler and Dellago 2019)}}
\sphinxAtStartPar
Singraber, A.; Behler, J.; Dellago, C. J., Chem.  Theory Comput. 2019, 15 (3), 1827\sphinxhyphen{}1840

\sphinxlineitem{\sphinxstylestrong{(Singraber et al 2019)}}
\sphinxAtStartPar
Singraber, A.; Morawietz, T.; Behler, J.; Dellago, C., J. Chem. Theory Comput.  2019, 15 (5), 3075\sphinxhyphen{}3092.

\sphinxlineitem{\sphinxstylestrong{(Sirk1)}}
\sphinxAtStartPar
Sirk TW, Sliozberg YR, Brennan JK, Lisal M, Andzelm JW, J Chem Phys, 136 (13) 134903, 2012.

\sphinxlineitem{\sphinxstylestrong{(Sirk2)}}
\sphinxAtStartPar
Sirk, Moore, Brown, J Chem Phys, 138, 064505 (2013).

\sphinxlineitem{\sphinxstylestrong{(Skomski)}}
\sphinxAtStartPar
Skomski, R. (2008). Simple models of magnetism. Oxford University Press.

\sphinxlineitem{\sphinxstylestrong{(Snodin)}}
\sphinxAtStartPar
B.E. Snodin, F. Randisi, M. Mosayebi, et al., J. Chem. Phys. 142, 234901 (2015).

\sphinxlineitem{\sphinxstylestrong{(Son)}}
\sphinxAtStartPar
Son, McDaniel, Cui and Yethiraj, J Phys Chem Lett, 10, 7523 (2019).

\sphinxlineitem{\sphinxstylestrong{(Srivastava)}}
\sphinxAtStartPar
Zhigilei, Wei, Srivastava, Phys. Rev. B 71, 165417 (2005).

\sphinxlineitem{\sphinxstylestrong{(Steinbach)}}
\sphinxAtStartPar
Steinbach, Brooks, J Comput Chem, 15, 667 (1994).

\sphinxlineitem{\sphinxstylestrong{(Steinhardt)}}
\sphinxAtStartPar
P. Steinhardt, D. Nelson, and M. Ronchetti, Phys. Rev. B 28, 784 (1983).

\sphinxlineitem{\sphinxstylestrong{(Steward)}}
\sphinxAtStartPar
Stewart, Spearot, Modelling Simul. Mater. Sci. Eng. 21, 045003, (2013).

\sphinxlineitem{\sphinxstylestrong{(Stewart2018)}}
\sphinxAtStartPar
J.A. Stewart, et al. (2018) Journal of Applied Physics, 123(16), 165902.

\sphinxlineitem{\sphinxstylestrong{(Stiles)}}
\sphinxAtStartPar
Stiles , Hubbard, and Kayser, J Chem Phys, 77, 6189 (1982).

\sphinxlineitem{\sphinxstylestrong{(Stillinger)}}
\sphinxAtStartPar
Stillinger, Weber, Phys. Rev. B 31, 5262 (1985).

\sphinxlineitem{\sphinxstylestrong{(Stoddard)}}
\sphinxAtStartPar
Stoddard and Ford, Phys Rev A, 8, 1504 (1973).

\sphinxlineitem{\sphinxstylestrong{(Streitz)}}
\sphinxAtStartPar
F. H. Streitz, J. W. Mintmire, Phys Rev B, 50, 11996\sphinxhyphen{}12003 (1994).

\sphinxlineitem{\sphinxstylestrong{(Strong)}}
\sphinxAtStartPar
Strong and Eaves, J. Phys. Chem. B 121, 189 (2017).

\sphinxlineitem{\sphinxstylestrong{(Stuart)}}
\sphinxAtStartPar
Stuart, Tutein, Harrison, J Chem Phys, 112, 6472\sphinxhyphen{}6486 (2000).

\sphinxlineitem{\sphinxstylestrong{(Stukowski)}}
\sphinxAtStartPar
Stukowski, Sadigh, Erhart, Caro; Modeling Simulation Materials Science \& Engineering, 7, 075005 (2009).

\sphinxlineitem{\sphinxstylestrong{(Su)}}
\sphinxAtStartPar
Su and Goddard, Excited Electron Dynamics Modeling of Warm Dense Matter, Phys Rev Lett, 99:185003 (2007).

\sphinxlineitem{\sphinxstylestrong{(Sulc1)}}
\sphinxAtStartPar
P. Sulc, F. Romano, T. E. Ouldridge, et al., J. Chem. Phys. 140, 235102 (2014).

\sphinxlineitem{\sphinxstylestrong{(Sulc2)}}
\sphinxAtStartPar
P. Sulc, F. Romano, T.E. Ouldridge, L. Rovigatti, J.P.K. Doye, A.A. Louis, J. Chem. Phys. 137, 135101 (2012).

\sphinxlineitem{\sphinxstylestrong{(Sun)}}
\sphinxAtStartPar
Sun, J. Phys. Chem. B, 102, 7338\sphinxhyphen{}7364 (1998).

\sphinxlineitem{\sphinxstylestrong{(Surblys2019)}}
\sphinxAtStartPar
Surblys, Matsubara, Kikugawa, Ohara, Phys Rev E, 99, 051301(R) (2019).

\sphinxlineitem{\sphinxstylestrong{(Surblys2021)}}
\sphinxAtStartPar
Surblys, Matsubara, Kikugawa, Ohara, J Appl Phys 130, 215104 (2021).

\sphinxlineitem{\sphinxstylestrong{(Sutmann)}}
\sphinxAtStartPar
Sutmann, Arnold, Fahrenberger, et. al., Physical review / E 88(6), 063308 (2013)

\sphinxlineitem{\sphinxstylestrong{(Sutmann)} G. Sutmann. ScaFaCoS \sphinxhyphen{} a Scalable library of Fast Coulomb Solvers for particle Systems.}
\sphinxAtStartPar
In Bajaj, Zavattieri, Koslowski, Siegmund, Proceedings of the Society of Engineering Science 51st Annual Technical Meeting. 2014.

\sphinxlineitem{\sphinxstylestrong{(Swinburne)}}
\sphinxAtStartPar
Swinburne and Marinica, Physical Review Letters, 120, 1 (2018)

\sphinxlineitem{\sphinxstylestrong{(Tadmor)}}
\sphinxAtStartPar
Tadmor, Elliott, Sethna, Miller and Becker, JOM, 63, 17 (2011). doi: \sphinxurl{https://doi.org/10.1007/s11837-011-0102-6}

\sphinxlineitem{\sphinxstylestrong{(Tainter 2011)}}
\sphinxAtStartPar
Tainter, Pieniazek, Lin, and Skinner, J. Chem. Phys., 134, 184501 (2011)

\sphinxlineitem{\sphinxstylestrong{(Tainter 2015)}}
\sphinxAtStartPar
Tainter, Shi, and Skinner, 11, 2268 (2015)

\sphinxlineitem{\sphinxstylestrong{(Tang and Toennies)}}
\sphinxAtStartPar
J Chem Phys, 80, 3726 (1984).

\sphinxlineitem{\sphinxstylestrong{(Tee)}}
\sphinxAtStartPar
Tee and Searles, J. Chem. Phys. 156, 184101 (2022).

\sphinxlineitem{\sphinxstylestrong{(Templeton2010)}}
\sphinxAtStartPar
Templeton, JA; Jones, RE; Wagner, GJ, “Application of a field\sphinxhyphen{}based method to spatially varying thermal transport problems in molecular dynamics.” Modelling and Simulation in Materials Science and Engineering (2010), 18:085007.

\sphinxlineitem{\sphinxstylestrong{(Templeton2011)}}
\sphinxAtStartPar
Templeton, JA; Jones, RE; Lee, JW; Zimmerman, JA; Wong, BM, “A long\sphinxhyphen{}range electric field solver for molecular dynamics based on atomistic\sphinxhyphen{}to\sphinxhyphen{}continuum modeling.” Journal of Chemical Theory and Computation (2011), 7:1736.

\sphinxlineitem{\sphinxstylestrong{(tenWolde)}}
\sphinxAtStartPar
P. R. ten Wolde, M. J. Ruiz\sphinxhyphen{}Montero, D. Frenkel, J. Chem. Phys. 104, 9932 (1996).

\sphinxlineitem{\sphinxstylestrong{(Tersoff\_1)}}
\sphinxAtStartPar
J. Tersoff, Phys Rev B, 37, 6991 (1988).

\sphinxlineitem{\sphinxstylestrong{(Tersoff\_2)}}
\sphinxAtStartPar
J. Tersoff, Phys Rev B, 38, 9902 (1988).

\sphinxlineitem{\sphinxstylestrong{(Tersoff\_3)}}
\sphinxAtStartPar
J. Tersoff, Phys Rev B, 39, 5566 (1989); errata (PRB 41, 3248)

\sphinxlineitem{\sphinxstylestrong{(Theodorou)}}
\sphinxAtStartPar
Theodorou, Suter, Macromolecules, 18, 1206 (1985).

\sphinxlineitem{\sphinxstylestrong{(Thole)}}
\sphinxAtStartPar
Chem Phys, 59, 341 (1981).

\sphinxlineitem{\sphinxstylestrong{(Thompson1)}}
\sphinxAtStartPar
Thompson, Plimpton, Mattson, J Chem Phys, 131, 154107 (2009).

\sphinxlineitem{\sphinxstylestrong{(Thompson2)}}
\sphinxAtStartPar
Thompson, Swiler, Trott, Foiles, Tucker, J Comp Phys, 285, 316 (2015).

\sphinxlineitem{\sphinxstylestrong{(Thornton et al, 2013)}}
\sphinxAtStartPar
Thornton, C., Cummins, S. J., \& Cleary, P. W. (2013).  An investigation of the comparative behavior of alternative contact force models during inelastic collisions. Powder

\sphinxlineitem{\sphinxstylestrong{(Thornton, 1991)}}
\sphinxAtStartPar
Thornton, C. (1991). Interparticle sliding in the presence of adhesion.  J. Phys. D: Appl. Phys. 24 1942

\sphinxlineitem{\sphinxstylestrong{(To)}}
\sphinxAtStartPar
Q.D. To, V.H. Vu, G. Lauriat, and C. Leonard. J. Math. Phys. 56, 103101 (2015).

\sphinxlineitem{\sphinxstylestrong{(Todd)}}
\sphinxAtStartPar
B. D. Todd, Denis J. Evans, and Peter J. Daivis: “Pressure tensor for inhomogeneous fluids”, Phys. Rev. E 52, 1627 (1995).

\sphinxlineitem{\sphinxstylestrong{(Toukmaji)}}
\sphinxAtStartPar
Toukmaji, Sagui, Board, and Darden, J Chem Phys, 113, 10913 (2000).

\sphinxlineitem{\sphinxstylestrong{(Toxvaerd)}}
\sphinxAtStartPar
Toxvaerd, Dyre, J Chem Phys, 134, 081102 (2011).

\sphinxlineitem{\sphinxstylestrong{(Tranchida)}}
\sphinxAtStartPar
Tranchida, Plimpton, Thibaudeau and Thompson, Journal of Computational Physics, 372, 406\sphinxhyphen{}425, (2018).

\sphinxlineitem{\sphinxstylestrong{(Tribello)}}
\sphinxAtStartPar
G.A. Tribello, M. Bonomi, D. Branduardi, C. Camilloni and G. Bussi, Comp. Phys. Comm 185, 604 (2014)

\sphinxlineitem{\sphinxstylestrong{(Tsuji et al, 1992)}}
\sphinxAtStartPar
Tsuji, Y., Tanaka, T., \& Ishida, T. (1992). Lagrangian numerical simulation of plug flow of cohesionless particles in a horizontal pipe. Powder technology, 71(3),

\sphinxlineitem{\sphinxstylestrong{(Tsuzuki)}}
\sphinxAtStartPar
Tsuzuki, Branicio, Rino, Comput Phys Comm, 177, 518 (2007).

\sphinxlineitem{\sphinxstylestrong{(Tuckerman1)}}
\sphinxAtStartPar
M. Tuckerman and B. Berne, J Chem Phys, 99, 2796 (1993).

\sphinxlineitem{\sphinxstylestrong{(Tuckerman2)}}
\sphinxAtStartPar
Tuckerman, Alejandre, Lopez\sphinxhyphen{}Rendon, Jochim, and Martyna, J Phys A: Math Gen, 39, 5629 (2006).

\sphinxlineitem{\sphinxstylestrong{(Tuckerman3)}}
\sphinxAtStartPar
Tuckerman, Berne and Martyna, J Chem Phys, 97, p 1990 (1992).

\sphinxlineitem{\sphinxstylestrong{(Tuckerman4)}}
\sphinxAtStartPar
Tuckerman, Mundy, Balasubramanian, Klein, J Chem Phys, 106, 5615 (1997).

\sphinxlineitem{\sphinxstylestrong{(Tyagi)}}
\sphinxAtStartPar
Tyagi, Suzen, Sega, Barbosa, Kantorovich, Holm, J Chem Phys, 132, 154112 (2010)

\sphinxlineitem{\sphinxstylestrong{(Ulomek)}}
\sphinxAtStartPar
Ulomek, Brien, Foiles, Mohles, Modelling Simul. Mater. Sci. Eng. 23 (2015) 025007

\sphinxlineitem{\sphinxstylestrong{(Vaiwala)}}
\sphinxAtStartPar
Vaiwala, Jadhav, and Thaokar, J Chem Phys, 146, 124904 (2017).

\sphinxlineitem{\sphinxstylestrong{(Valone)}}
\sphinxAtStartPar
Valone, Baskes, Martin, Phys. Rev. B, 73, 214209 (2006).

\sphinxlineitem{\sphinxstylestrong{(vanWijk)}}\begin{enumerate}
\sphinxsetlistlabels{\Alph}{enumi}{enumii}{}{.}%
\setcounter{enumi}{12}
\item {} \begin{enumerate}
\sphinxsetlistlabels{\Alph}{enumii}{enumiii}{}{.}%
\setcounter{enumii}{12}
\item {} 
\sphinxAtStartPar
van Wijk, A. Schuring, M. I. Katsnelson, and A. Fasolino, Physical Review Letters, 113, 135504 (2014)

\end{enumerate}

\end{enumerate}

\sphinxlineitem{\sphinxstylestrong{(Van Workum)}}\begin{enumerate}
\sphinxsetlistlabels{\Alph}{enumi}{enumii}{}{.}%
\setcounter{enumi}{10}
\item {} 
\sphinxAtStartPar
Van Workum et al., J. Chem. Phys. 125 144506 (2006)

\end{enumerate}

\sphinxlineitem{\sphinxstylestrong{(Vargas and McCarthy 2001)}}
\sphinxAtStartPar
Vargas, W.L. and McCarthy, J.J. (2001).

\sphinxlineitem{\sphinxstylestrong{(Varshalovich)}}
\sphinxAtStartPar
Varshalovich, Moskalev, Khersonskii, Quantum Theory of Angular Momentum, World Scientific, Singapore (1987).

\sphinxlineitem{\sphinxstylestrong{(Vashishta1990)}}
\sphinxAtStartPar
P. Vashishta, R. K. Kalia, J. P. Rino, Phys. Rev. B 41, 12197 (1990).

\sphinxlineitem{\sphinxstylestrong{(Vashishta2007)}}
\sphinxAtStartPar
P. Vashishta, R. K. Kalia, A. Nakano, J. P. Rino. J. Appl. Phys. 101, 103515 (2007).

\sphinxlineitem{\sphinxstylestrong{(Veld)}}
\sphinxAtStartPar
In ‘t Veld, Ismail, Grest, J Chem Phys, 127, 144711 (2007).

\sphinxlineitem{\sphinxstylestrong{(Verstraelen)}}
\sphinxAtStartPar
Verstraelen, Ayers, Speybroeck, Waroquier, J. Chem. Phys. 138, 074108 (2013).

\sphinxlineitem{\sphinxstylestrong{(Volkov1)}}
\sphinxAtStartPar
Volkov and Zhigilei, J Phys Chem C, 114, 5513 (2010).

\sphinxlineitem{\sphinxstylestrong{(Volkov2)}}
\sphinxAtStartPar
Volkov, Simov and Zhigilei, APS Meeting Abstracts, Q31.013 (2008).

\sphinxlineitem{\sphinxstylestrong{(Voter1998)}}
\sphinxAtStartPar
Voter, Phys Rev B, 57, 13985 (1998).

\sphinxlineitem{\sphinxstylestrong{(Voter2000)}}
\sphinxAtStartPar
Sorensen and Voter, J Chem Phys, 112, 9599 (2000)

\sphinxlineitem{\sphinxstylestrong{(Voter2002)}}
\sphinxAtStartPar
Voter, Montalenti, Germann, Annual Review of Materials Research 32, 321 (2002).

\sphinxlineitem{\sphinxstylestrong{(Voter2013)}}
\sphinxAtStartPar
S. Y. Kim, D. Perez, A. F. Voter, J Chem Phys, 139, 144110 (2013).

\sphinxlineitem{\sphinxstylestrong{(Wagner)}}
\sphinxAtStartPar
Wagner, GJ; Jones, RE; Templeton, JA; Parks, MA, “An atomistic\sphinxhyphen{}to\sphinxhyphen{}continuum coupling method for heat transfer in solids.” Special Issue of Computer Methods and Applied Mechanics (2008) 197:3351.

\sphinxlineitem{\sphinxstylestrong{(Wang et al, 2015)}}
\sphinxAtStartPar
Wang, Y., Alonso\sphinxhyphen{}Marroquin, F., \& Guo, W. W. (2015).  Rolling and sliding in 3\sphinxhyphen{}D discrete element models. Particuology, 23, 49\sphinxhyphen{}55.

\sphinxlineitem{\sphinxstylestrong{(Wang2020)}}\begin{enumerate}
\sphinxsetlistlabels{\Alph}{enumi}{enumii}{}{.}%
\setcounter{enumi}{23}
\item {} 
\sphinxAtStartPar
Wang, S. Ramirez\sphinxhyphen{}Hinestrosa, J. Dobnikar, and D. Frenkel, Phys. Chem. Chem. Phys. 22, 10624 (2020).

\end{enumerate}

\sphinxlineitem{\sphinxstylestrong{(Wang1)}}
\sphinxAtStartPar
J. Wang, H. S. Yu, P. A. Langston, F. Y. Fraige, Granular Matter, 13, 1 (2011).

\sphinxlineitem{\sphinxstylestrong{(Wang2)}}
\sphinxAtStartPar
J. Wang, and A. Rockett, Phys. Rev. B, 43, 12571 (1991).

\sphinxlineitem{\sphinxstylestrong{(Wang3)}}
\sphinxAtStartPar
Wang and Holm, J Chem Phys, 115, 6277 (2001).

\sphinxlineitem{\sphinxstylestrong{(Wang4)}}
\sphinxAtStartPar
Wang, Van Hove, Ross, Baskes, J. Chem. Phys., 121, 5410 (2004).

\sphinxlineitem{\sphinxstylestrong{(Ward)}}
\sphinxAtStartPar
D.K. Ward, X.W. Zhou, B.M. Wong, F.P. Doty, and J.A. Zimmerman, Phys. Rev. B, 85,115206 (2012).

\sphinxlineitem{\sphinxstylestrong{(Warren)}}
\sphinxAtStartPar
Warren, Phys Rev E, 68, 066702 (2003).

\sphinxlineitem{\sphinxstylestrong{(Watkins)}}
\sphinxAtStartPar
Watkins and Jorgensen, J Phys Chem A, 105, 4118\sphinxhyphen{}4125 (2001).

\sphinxlineitem{\sphinxstylestrong{(Weeks)}}
\sphinxAtStartPar
Weeks, Chandler and Andersen, J. Chem. Phys., 54, 5237 (1971)

\sphinxlineitem{\sphinxstylestrong{(WeinanE)}}
\sphinxAtStartPar
E, Ren, Vanden\sphinxhyphen{}Eijnden, Phys Rev B, 66, 052301 (2002).

\sphinxlineitem{\sphinxstylestrong{(Wen)}}
\sphinxAtStartPar
M. Wen, S. Carr, S. Fang, E. Kaxiras, and E. B. Tadmor, Phys. Rev. B, 98, 235404 (2018)

\sphinxlineitem{\sphinxstylestrong{(Wennberg)}}
\sphinxAtStartPar
Wennberg, Murtola, Hess, Lindahl, J Chem Theory Comput, 9, 3527 (2013).

\sphinxlineitem{\sphinxstylestrong{(Wicaksono1)}}
\sphinxAtStartPar
Wicaksono, Sinclair, Militzer, Computational Materials Science, 117, 397\sphinxhyphen{}405 (2016).

\sphinxlineitem{\sphinxstylestrong{(Wicaksono2)}}
\sphinxAtStartPar
Wicaksono, figshare, \sphinxurl{https://doi.org/10.6084/m9.figshare.1488628.v1} (2015).

\sphinxlineitem{\sphinxstylestrong{(Winkler)}}
\sphinxAtStartPar
Winkler, Wysocki, and Gompper, Soft Matter, 11, 6680 (2015).

\sphinxlineitem{\sphinxstylestrong{(Wirnsberger)}}
\sphinxAtStartPar
Wirnsberger, Frenkel, and Dellago, J Chem Phys, 143, 124104 (2015).

\sphinxlineitem{\sphinxstylestrong{(Wolf)}}
\sphinxAtStartPar
D. Wolf, P. Keblinski, S. R. Phillpot, J. Eggebrecht, J Chem Phys, 110, 8254 (1999).

\sphinxlineitem{\sphinxstylestrong{(Wolff)}}
\sphinxAtStartPar
Wolff and Rudd, Comp Phys Comm, 120, 200\sphinxhyphen{}32 (1999).

\sphinxlineitem{\sphinxstylestrong{(Wood)}}
\sphinxAtStartPar
Wood and Thompson, J Chem Phys, 148, 241721, (2018)

\sphinxlineitem{\sphinxstylestrong{(Xie23)}}
\sphinxAtStartPar
Xie, S.R., Rupp, M. \& Hennig, R.G. Ultra\sphinxhyphen{}fast interpretable machine\sphinxhyphen{}learning potentials. npj Comput Mater 9, 162 (2023). \sphinxurl{https://doi.org/10.1038/s41524-023-01092-7}

\sphinxlineitem{\sphinxstylestrong{(Yade\sphinxhyphen{}DEM)}}\begin{enumerate}
\sphinxsetlistlabels{\Alph}{enumi}{enumii}{}{.}%
\setcounter{enumi}{21}
\item {} 
\sphinxAtStartPar
Smilauer et al. (2021), Yade Documentation 3rd ed.

\end{enumerate}

\sphinxlineitem{\sphinxstylestrong{(Yanxon2020)}}
\sphinxAtStartPar
Yanxon, Zagaceta, Tang, Matteson, Zhu, Mach. Learn.: Sci. Technol. 2, 027001 (2020).

\sphinxlineitem{\sphinxstylestrong{(Yeh)}}
\sphinxAtStartPar
Yeh and Berkowitz, J Chem Phys, 111, 3155 (1999).

\sphinxlineitem{\sphinxstylestrong{(Yuan2010a)}}
\sphinxAtStartPar
Yuan, Huang, Li, Lykotrafitis, Zhang, Phys. Rev. E, 82, 011905(2010).

\sphinxlineitem{\sphinxstylestrong{(Yuan2010b)}}
\sphinxAtStartPar
Yuan, Huang, Zhang, Soft. Matter, 6, 4571(2010).

\sphinxlineitem{\sphinxstylestrong{(Zagaceta2020)}}
\sphinxAtStartPar
Zagaceta, Yanxon, Zhu, J Appl Phys, 128, 045113 (2020).

\sphinxlineitem{\sphinxstylestrong{(ZBL)}}
\sphinxAtStartPar
J.F. Ziegler, J.P. Biersack, U. Littmark, ‘Stopping and Ranges of Ions in Matter’ Vol 1, 1985, Pergamon Press.

\sphinxlineitem{\sphinxstylestrong{(Zhang1)}}
\sphinxAtStartPar
Zhang and Makse, Phys Rev E, 72, p 011301 (2005).

\sphinxlineitem{\sphinxstylestrong{(Zhang2)}}
\sphinxAtStartPar
Zhang and Trinkle, Computational Materials Science, 124, 204\sphinxhyphen{}210 (2016).

\sphinxlineitem{\sphinxstylestrong{(Zhang3)}}
\sphinxAtStartPar
Zhang, Glotzer, Nanoletters, 4, 1407\sphinxhyphen{}1413 (2004).

\sphinxlineitem{\sphinxstylestrong{(Zhang4)}}
\sphinxAtStartPar
Zhang, J Chem Phys, 106, 6102 (1997).

\sphinxlineitem{\sphinxstylestrong{(Zhang5)}}
\sphinxAtStartPar
Zhang, Lussetti, de Souza, Muller\sphinxhyphen{}Plathe, J Phys Chem B, 109, 15060\sphinxhyphen{}15067 (2005).

\sphinxlineitem{\sphinxstylestrong{(Zhigilei1)}}
\sphinxAtStartPar
Volkov and Zhigilei, ACS Nano 4, 6187 (2010).

\sphinxlineitem{\sphinxstylestrong{(Zhigilei2)}}
\sphinxAtStartPar
Volkov, Simov, Zhigilei, ASME paper IMECE2008, 68021 (2008).

\sphinxlineitem{\sphinxstylestrong{(Zhigilei3)}}
\sphinxAtStartPar
Volkov, Zhigilei, J. Phys. Chem. C 114, 5513 (2010).

\sphinxlineitem{\sphinxstylestrong{(Zhigilei4)}}
\sphinxAtStartPar
Wittmaack, Banna, Volkov, Zhigilei, Carbon 130, 69 (2018).

\sphinxlineitem{\sphinxstylestrong{(Zhigilei5)}}
\sphinxAtStartPar
Wittmaack, Volkov, Zhigilei, Compos. Sci. Technol. 166, 66 (2018).

\sphinxlineitem{\sphinxstylestrong{(Zhigilei6)}}
\sphinxAtStartPar
Wittmaack, Volkov, Zhigilei, Carbon 143, 587 (2019).

\sphinxlineitem{\sphinxstylestrong{(Zhigilei7)}}
\sphinxAtStartPar
Volkov, Zhigilei, Phys. Rev. Lett. 104, 215902 (2010).

\sphinxlineitem{\sphinxstylestrong{(Zhigilei8)}}
\sphinxAtStartPar
Volkov, Shiga, Nicholson, Shiomi, Zhigilei, J. Appl. Phys. 111, 053501 (2012).

\sphinxlineitem{\sphinxstylestrong{(Zhigilei9)}}
\sphinxAtStartPar
Volkov, Zhigilei, Appl. Phys. Lett. 101, 043113 (2012).

\sphinxlineitem{\sphinxstylestrong{(Zhigilei10)}}
\sphinxAtStartPar
Jacobs, Nicholson, Zemer, Volkov, Zhigilei, Phys. Rev. B 86, 165414 (2012).

\sphinxlineitem{\sphinxstylestrong{(Zhou1)}}
\sphinxAtStartPar
Zhou, Saidi, Fichthorn, J Phys Chem C, 118(6), 3366\sphinxhyphen{}3374 (2014).

\sphinxlineitem{\sphinxstylestrong{(Zhou3)}}
\sphinxAtStartPar
X. W. Zhou, M. E. Foster, R. E. Jones, P. Yang, H. Fan, and F. P. Doty, J. Mater. Sci. Res., 4, 15 (2015).

\sphinxlineitem{\sphinxstylestrong{(Zhou4)}}
\sphinxAtStartPar
X. W. Zhou, M. E. Foster, J. A. Ronevich, and C. W. San Marchi, J. Comp. Chem., 41, 1299 (2020).

\sphinxlineitem{\sphinxstylestrong{(Zhu)}}
\sphinxAtStartPar
Zhu, Tajkhorshid, and Schulten, Biophys. J. 83, 154 (2002).

\sphinxlineitem{\sphinxstylestrong{(Ziegler)}}
\sphinxAtStartPar
J.F. Ziegler, J. P. Biersack and U. Littmark, “The Stopping and Range of Ions in Matter”, Volume 1, Pergamon, 1985.

\sphinxlineitem{\sphinxstylestrong{(Zimmerman2004)}}
\sphinxAtStartPar
Zimmerman, JA; Webb, EB; Hoyt, JJ;. Jones, RE; Klein, PA; Bammann, DJ, “Calculation of stress in atomistic simulation.” Special Issue of Modelling and Simulation in Materials Science and Engineering (2004),12:S319.

\sphinxlineitem{\sphinxstylestrong{(Zimmerman2010)}}
\sphinxAtStartPar
Zimmerman, JA; Jones, RE; Templeton, JA, “A material frame approach for evaluating continuum variables in atomistic simulations.” Journal of Computational Physics (2010), 229:2364.

\sphinxlineitem{\sphinxstylestrong{(electronic stopping)}}
\sphinxAtStartPar
Wikipedia \sphinxhyphen{} Electronic Stopping Power: \sphinxurl{https://en.wikipedia.org/wiki/Stopping\_power\_\%28particle\_radiation\%29}

\end{description}


\part{Indices and tables}
\label{\detokenize{Manual:indices-and-tables}}


\renewcommand{\indexname}{Index}
\printindex
\end{document}